% Options for packages loaded elsewhere
% Options for packages loaded elsewhere
\PassOptionsToPackage{unicode}{hyperref}
\PassOptionsToPackage{hyphens}{url}
\PassOptionsToPackage{dvipsnames,svgnames,x11names}{xcolor}
%
\documentclass[
  letterpaper,
  DIV=11,
  numbers=noendperiod]{scrreprt}
\usepackage{xcolor}
\usepackage[margin=1in]{geometry}
\usepackage{amsmath,amssymb}
\setcounter{secnumdepth}{5}
\usepackage{iftex}
\ifPDFTeX
  \usepackage[T1]{fontenc}
  \usepackage[utf8]{inputenc}
  \usepackage{textcomp} % provide euro and other symbols
\else % if luatex or xetex
  \usepackage{unicode-math} % this also loads fontspec
  \defaultfontfeatures{Scale=MatchLowercase}
  \defaultfontfeatures[\rmfamily]{Ligatures=TeX,Scale=1}
\fi
\usepackage{lmodern}
\ifPDFTeX\else
  % xetex/luatex font selection
\fi
% Use upquote if available, for straight quotes in verbatim environments
\IfFileExists{upquote.sty}{\usepackage{upquote}}{}
\IfFileExists{microtype.sty}{% use microtype if available
  \usepackage[]{microtype}
  \UseMicrotypeSet[protrusion]{basicmath} % disable protrusion for tt fonts
}{}
\makeatletter
\@ifundefined{KOMAClassName}{% if non-KOMA class
  \IfFileExists{parskip.sty}{%
    \usepackage{parskip}
  }{% else
    \setlength{\parindent}{0pt}
    \setlength{\parskip}{6pt plus 2pt minus 1pt}}
}{% if KOMA class
  \KOMAoptions{parskip=half}}
\makeatother
% Make \paragraph and \subparagraph free-standing
\makeatletter
\ifx\paragraph\undefined\else
  \let\oldparagraph\paragraph
  \renewcommand{\paragraph}{
    \@ifstar
      \xxxParagraphStar
      \xxxParagraphNoStar
  }
  \newcommand{\xxxParagraphStar}[1]{\oldparagraph*{#1}\mbox{}}
  \newcommand{\xxxParagraphNoStar}[1]{\oldparagraph{#1}\mbox{}}
\fi
\ifx\subparagraph\undefined\else
  \let\oldsubparagraph\subparagraph
  \renewcommand{\subparagraph}{
    \@ifstar
      \xxxSubParagraphStar
      \xxxSubParagraphNoStar
  }
  \newcommand{\xxxSubParagraphStar}[1]{\oldsubparagraph*{#1}\mbox{}}
  \newcommand{\xxxSubParagraphNoStar}[1]{\oldsubparagraph{#1}\mbox{}}
\fi
\makeatother


\usepackage{longtable,booktabs,array}
\usepackage{calc} % for calculating minipage widths
% Correct order of tables after \paragraph or \subparagraph
\usepackage{etoolbox}
\makeatletter
\patchcmd\longtable{\par}{\if@noskipsec\mbox{}\fi\par}{}{}
\makeatother
% Allow footnotes in longtable head/foot
\IfFileExists{footnotehyper.sty}{\usepackage{footnotehyper}}{\usepackage{footnote}}
\makesavenoteenv{longtable}
\usepackage{graphicx}
\makeatletter
\newsavebox\pandoc@box
\newcommand*\pandocbounded[1]{% scales image to fit in text height/width
  \sbox\pandoc@box{#1}%
  \Gscale@div\@tempa{\textheight}{\dimexpr\ht\pandoc@box+\dp\pandoc@box\relax}%
  \Gscale@div\@tempb{\linewidth}{\wd\pandoc@box}%
  \ifdim\@tempb\p@<\@tempa\p@\let\@tempa\@tempb\fi% select the smaller of both
  \ifdim\@tempa\p@<\p@\scalebox{\@tempa}{\usebox\pandoc@box}%
  \else\usebox{\pandoc@box}%
  \fi%
}
% Set default figure placement to htbp
\def\fps@figure{htbp}
\makeatother





\setlength{\emergencystretch}{3em} % prevent overfull lines

\providecommand{\tightlist}{%
  \setlength{\itemsep}{0pt}\setlength{\parskip}{0pt}}



 


\usepackage{fvextra}
\DefineVerbatimEnvironment{Highlighting}{Verbatim}{breaklines,commandchars=\\\{\}}
\KOMAoption{captions}{tableheading}
\makeatletter
\@ifpackageloaded{tcolorbox}{}{\usepackage[skins,breakable]{tcolorbox}}
\@ifpackageloaded{fontawesome5}{}{\usepackage{fontawesome5}}
\definecolor{quarto-callout-color}{HTML}{909090}
\definecolor{quarto-callout-note-color}{HTML}{0758E5}
\definecolor{quarto-callout-important-color}{HTML}{CC1914}
\definecolor{quarto-callout-warning-color}{HTML}{EB9113}
\definecolor{quarto-callout-tip-color}{HTML}{00A047}
\definecolor{quarto-callout-caution-color}{HTML}{FC5300}
\definecolor{quarto-callout-color-frame}{HTML}{acacac}
\definecolor{quarto-callout-note-color-frame}{HTML}{4582ec}
\definecolor{quarto-callout-important-color-frame}{HTML}{d9534f}
\definecolor{quarto-callout-warning-color-frame}{HTML}{f0ad4e}
\definecolor{quarto-callout-tip-color-frame}{HTML}{02b875}
\definecolor{quarto-callout-caution-color-frame}{HTML}{fd7e14}
\makeatother
\makeatletter
\@ifpackageloaded{bookmark}{}{\usepackage{bookmark}}
\makeatother
\makeatletter
\@ifpackageloaded{caption}{}{\usepackage{caption}}
\AtBeginDocument{%
\ifdefined\contentsname
  \renewcommand*\contentsname{Table of contents}
\else
  \newcommand\contentsname{Table of contents}
\fi
\ifdefined\listfigurename
  \renewcommand*\listfigurename{List of Figures}
\else
  \newcommand\listfigurename{List of Figures}
\fi
\ifdefined\listtablename
  \renewcommand*\listtablename{List of Tables}
\else
  \newcommand\listtablename{List of Tables}
\fi
\ifdefined\figurename
  \renewcommand*\figurename{Figure}
\else
  \newcommand\figurename{Figure}
\fi
\ifdefined\tablename
  \renewcommand*\tablename{Table}
\else
  \newcommand\tablename{Table}
\fi
}
\@ifpackageloaded{float}{}{\usepackage{float}}
\floatstyle{ruled}
\@ifundefined{c@chapter}{\newfloat{codelisting}{h}{lop}}{\newfloat{codelisting}{h}{lop}[chapter]}
\floatname{codelisting}{Listing}
\newcommand*\listoflistings{\listof{codelisting}{List of Listings}}
\makeatother
\makeatletter
\makeatother
\makeatletter
\@ifpackageloaded{caption}{}{\usepackage{caption}}
\@ifpackageloaded{subcaption}{}{\usepackage{subcaption}}
\makeatother
\usepackage{bookmark}
\IfFileExists{xurl.sty}{\usepackage{xurl}}{} % add URL line breaks if available
\urlstyle{same}
\hypersetup{
  pdftitle={OPM Flow Reference Manual},
  pdfauthor={Open Porous Media},
  colorlinks=true,
  linkcolor={blue},
  filecolor={Maroon},
  citecolor={Blue},
  urlcolor={Blue},
  pdfcreator={LaTeX via pandoc}}


\title{OPM Flow Reference Manual}
\usepackage{etoolbox}
\makeatletter
\providecommand{\subtitle}[1]{% add subtitle to \maketitle
  \apptocmd{\@title}{\par {\large #1 \par}}{}{}
}
\makeatother
\subtitle{2025-04}
\author{Open Porous Media}
\date{2025-01-01}
\begin{document}
\maketitle

\renewcommand*\contentsname{Table of contents}
{
\hypersetup{linkcolor=}
\setcounter{tocdepth}{2}
\tableofcontents
}

\bookmarksetup{startatroot}

\chapter{OPM Flow Reference Manual}\label{opm-flow-reference-manual}

2025-04

\hfill\break

\section{About This Manual}\label{about-this-manual}

The \textbf{OPM Flow Reference Manual} provides comprehensive
documentation for the Open Porous Media (OPM) Flow reservoir simulator.
OPM Flow is a fully-implicit, black-oil and compositional reservoir
simulator capable of running industry-standard simulation models.

This manual covers:

\begin{itemize}
\tightlist
\item
  \textbf{Installation and Setup} -- How to install and run OPM Flow
\item
  \textbf{Keyword Reference} -- Complete documentation of all supported
  keywords organized by input deck section (GLOBAL, RUNSPEC, GRID, EDIT,
  PROPS, REGIONS, SOLUTION, SUMMARY, SCHEDULE)
\item
  \textbf{Appendices} -- Keyword summaries, release notes, the OPMRUN
  graphical interface, Python scripting interface, command line options,
  and output file formats
\end{itemize}

\subsection{How to Use This Manual}\label{how-to-use-this-manual}

Use the table of contents or the search bar to navigate to specific
keywords or topics. Keywords are organized by their input deck section,
matching the structure of OPM Flow simulation input files.

\begin{tcolorbox}[enhanced jigsaw, title=\textcolor{quarto-callout-note-color}{\faInfo}\hspace{0.5em}{Note}, coltitle=black, breakable, toprule=.15mm, colbacktitle=quarto-callout-note-color!10!white, leftrule=.75mm, opacitybacktitle=0.6, colback=white, bottomtitle=1mm, toptitle=1mm, titlerule=0mm, left=2mm, colframe=quarto-callout-note-color-frame, arc=.35mm, rightrule=.15mm, bottomrule=.15mm, opacityback=0]

This manual corresponds to \textbf{OPM Flow version 2025-04}.

\end{tcolorbox}

\bookmarksetup{startatroot}

\chapter{INTRODUCTION}\label{introduction}

\section{Overview}\label{overview}

The Open Porous Media (``OPM'') initiative was started in 2009 to
encourage open innovation and reproducible research on modeling and
simulation of porous media processes. OPM was initially founded as a
collaboration between groups at Equinor (formerly Statoil), SINTEF, the
University of Stuttgart, and the University of Bergen, but over time,
several other groups and individuals have joined and contributed. What
today forms the OPM suite of software, has mainly been developed by
SINTEF, NORCE (formerly IRIS), Equinor, Ceetron Solutions (OPM
Resinsight), Poware Software Solutions, Dr.~Blatt
HPC-Simulation-Software \& Services, and OPM-OP. The initial vision was
to create long-lasting, efficient, and well-maintained, open-source
software for simulating flow and transport in porous media. The scope
has later been extended to also provide open data sets, thus making it
easier to benchmark, compare, and test different mathematical models,
computational methods, and software implementations.

All software in OPM is distributed under the GNU General Public License
(``GPL''), \href{http://www.gnu.org/licenses/gpl-3.0.en.html}{version
3}, with the option of using newer versions of the license. Datasets are
distributed under the Open Data Commons Open Database License
(``ODbL''), \href{http://opendatacommons.org/licenses/odbl/1.0/}{version
1.0}, with individual content under the Database Contents License
(``DbCL''), \href{http://opendatacommons.org/licenses/dbcl/1.0/}{version
1.0.}

The project's numerical simulator, OPM Flow, is a fully-implicit,
black-oil simulator capable of running industry-standard simulation
models {[}Rasmussen, A.F., Sandve, T.H., Bao, K., Lauser, A., Hove, J.,
Skaflestad, B., Klöfkorn, R., Blatt, M., Rustad, A.B., Sævareid, O. et
al.~{[}2019{]}. The Open Porous Media Flow Reservoir Simulator.
https://doi.org/10.1016/j.camwa.2020.05.014.{]}. The simulator's input
file is compatible with the commercial simulator, although OPM Flow has
extra keywords for supporting additional functionality not found in the
commercial simulator. The simulator's output files, except for the print
and debug files, are compatible with the commercial simulator, and are
readable by commercial post-processing software as well as by OPM
ResInsight. OPM ResInsight is the Open Porous Media's post-processing
software.

The simulator is implemented using automatic differentiation to enable
rapid development of new fluid models. In most simulators the
coefficients of the linearized systems of equations (Jacobian matrix)
required by the non-linear Newton Solver, have traditionally been
computed by evaluating the closed-form expressions obtained by
differentiating the discretized flow equations analytically.
Differentiating these flow equations manually and programming the
resulting formulas is time-consuming and error-prone. Thus, even simple
extensions in functional dependencies can be difficult to accommodate
into existing simulators. Automatic differentiation eliminates these
drawbacks.

OPM Flow runs under Linux and Windows Subsystem for Linux, and the
program is in active development with new features added in each
bi-annual release. Currently the program has the following
functionality:

\subsubsection{Model Formulation}\label{model-formulation}

\begin{itemize}
\tightlist
\item
  Black-oil with dissolved gas and vaporized oil.
\item
  Rock-dependent capillary and relative-permeability curves.
\item
  End-point scaling and hysteresis.
\item
  Oil vaporization controls (VAPPARS).
\item
  Rock compaction.
\end{itemize}

\subsubsection{Enhance Oil Recovery
Models}\label{enhance-oil-recovery-models}

Enhanced Oil Recovery (``EOR'') is a term applied to methods used for
recovering oil from a petroleum reservoir beyond that recoverable by
primary and secondary methods. Primary recovery normally refers to
production using the energy inherent in the reservoir from gas under
pressure or a natural water drive. Secondary recovery usually refers to
injection of water or water flooding. Thus, Enhanced Oil Recovery is
often synonymous with tertiary recovery. Improved Oil Recovery (``IOR'')
and Advanced Oil Recovery (``AOR'') have similar meaning, except they
also apply to primary and secondary methods, and sometimes EOR methods
can be used earlier in the sequence. Originally water flooding was
considered as Enhanced Oil Recovery, but now EOR is generally considered
to follow water flooding.

There are four major EOR groups of methods:

\begin{itemize}
\tightlist
\item
  thermal recovery,
\item
  gas miscible recovery,
\item
  chemical flooding, and
\item
  microbial flooding.
\end{itemize}

These main groups are further subdivided into individual processes as
illustrated in Figure 1.1. For example the thermal recovery methods
consist of steam flooding and cyclic steam stimulation, as well as
in-situ combustion. Figure 1.1 also includes hot water as a process for
thermal EOR, but this process is rarely used in the industry as a drive
mechanism, as steam contains significantly more heat than hot water and
therefore steam is more effective in heating up the in situ oil. Note
that microbial EOR methods, are not shown in Figure 1.1, due to there
limited use in the industry.

\pandocbounded{\includegraphics[keepaspectratio]{chapters/images/Frame31_751f1e6cbafb.png}}
\pandocbounded{\includegraphics[keepaspectratio]{chapters/images/Image157_751f1e6cbafb.png}}

Chemical flooding methods include polymer flooding (including polymer
gels), micellar-polymer flooding, and alkaline flooding. Quite often
these type of floods are combined to form a more efficient displacement
process, for example in an Alkaline Surfactant Polymer (``ASP'')
process.

Conventional gas flooding is an immiscible process, that is the oil and
gas phases remain separate during the displacement drive. For the EOR
gas drive process, the displacement is miscible, here the gas ``mixes''
with the oil to form a single phase fluid in the reservoir. Gas miscible
processes are subdivided into carbon dioxide flooding, cyclic carbon
dioxide stimulation, nitrogen flooding, and nitrogen-CO2 flooding.

OPM Flow has several EOR models based on extending the black-oil model
with a fourth component in the simulator:

\begin{itemize}
\item
  CO2 Standard EOR Model: OPM Flow's CO2 Standard EOR model is based on
  the simulators Solvent model by extending the black-oil formulation
  with a fourth component by adding a CO2 component to the gas phase.
  This is the standard formulation used by most black-oil simulators.
\item
  CO2 Dynamic EOR Model: Although the CO2 Standard EOR model is commonly
  used, the formulation cannot account for the mass transfer of the
  various components and phases. The extended black-oil formulations
  often poorly represent the PVT properties of the oil-CO2 mixtures,
  resulting in poor agreement with the compositional formulation. OPM
  Flow's CO2 Dynamic EOR Model black-oil formulation attempts to
  overcome this limitation by making the black-oil properties dependent
  on the fraction of CO2 in the cell {[}T. H. Sandve, O. Sævareid and I.
  Aavatsmark: ``Improved Extended Blackoil Formulation -- for CO2 EOR
  Simulations.'' in ECMOR XVII -- The 17th European Conference on the --
  Mathematics of Oil Recovery, September 2020.{]}. This approach models
  the oil-CO2 mixture more accurately, and thus give results closer to
  the compositional simulator. Note the model is not restricted to CO2
  utilization, as long as the appropriate oil-injection-gas dependent
  properties are entered in the model.
\item
  Foam Model: An experimental foam module has been added to OPM Flow
  2019-10 release. With this it is possible to simulate certain types of
  surfactant injection. Such injection stimulates formation of foam to
  change mobility ratios, and give better reservoir sweep. The
  implemented foam model treats surfactant transported in the gas or
  water phase, and reduces the mobility of that phase depending on the
  surfactant concentration. In addition to mobility reduction, foam
  adsorption to the reservoir rock is also included in the model.
\item
  Polymer Model (Standard): OPM Flow's polymer model is based on a
  black-oil polymer formulation, which is developed by extending the
  black-oil model with a polymer component. The effects of the polymer
  mixing are simulated based on the Todd-Longstaff {[}M. R. Todd and W.
  J Longstaff, The Development, Testing, and Application Of a Numerical
  Simulator for Predicting Miscible Flood Performance''. In: J. Petrol.
  Tech. 24.7 (1972), pages 874-882.{]} mixing model, and adsorption,
  dead pore space, and permeability reduction effects are also
  considered. A logarithmic shear thinning/thickening model has also
  been incorporated since the 2015-10 Release (see
  \href{http://opm-project.org/?p=592}{Flow-polymer}). Note that the
  Polymer model has now been incorporated into the main OPM Flow
  simulator and is no longer a separate simulator.
\item
  Polymer Model (Molecular Weight Transport): In addition to the
  standard polymer model, the simulator also has the Polymer Molecular
  Weight Transport option, that uses the polymer molecular weight in
  calculating the polymer viscosity. Here, the Mark-Howink equation,
  also known as the Mark-Howink-Staudinger equation {[}H. Mark, in R.
  Saenger, Der feste Koerper, Hirzel, Leipzig, 1938.{]}, {[}R. Houwink ,
  J. Prakt. Chem., Vol. 157, Issue 1-3, p.~15 (1940).{]}, and {[}H.
  Staudinger, Die Hochmolekulare Organischen Verbindungen, Julius
  Springer, Berlin 1932.{]} is used to calculate the intrinsic viscosity
  as a function of the molecular weight of the polymer.
\item
  Solvent Model: Similar to the polymer model, OPM Flow's solvent model
  again uses an extra component by extending the black-oil formulation
  with a fourth component in the simulator by adding a solvent component
  to the gas phase (see
  \href{http://opm-project.org/?page_id=565}{Flow-solvent}). Note that
  Solvent model has now been incorporated into the main OPM Flow
  simulator and is no longer a separate simulator.
\end{itemize}

\subsubsection{CO2 Storage:}\label{co2-storage}

\begin{itemize}
\tightlist
\item
  Brine-CO2 PVT module that generates brine and CO2 PVT properties
  directly eliminating user input for PVT.
\item
  Thermal properties generated by OPM Flow include the CO2 enthalpy
  represented as a table internally in the simulator. The liquid
  enthalpy depends on the dissolved CO2 and salinity as well as pressure
  and temperature. The water enthalpy is automatically calculated and
  modified to account for salinity and for CO2 following.
\item
  Diffusion for fine grid modeling is implemented via coefficients that
  depend on temperature, pressure, and salinity. The diffusion
  coefficient is computed internally for pure water and modified to
  account for salinity. The effect of the porous media on the diffusion
  is modeled using an industry standard relationship. In addition, the
  coefficient can also be given as an input parameter using a keyword.
\item
  Convective dissolution of CO2 into in situ brine within a grid cell is
  implemented together with a convective dissolution parameter for
  controlling the rate of dissolution.
\item
  Microbially Induced Calcite Precipitation (``MICP'') model for the
  modeling of biochemical processes in creating barriers by calcium
  carbonate cementation, the technology has the potential to be used for
  sealing leakage zones in geological formations {[}Landa-Marbán, D.,
  Tveit, S., Kumar, K., Gasda, S.E., 2021. Practical approaches to study
  microbially induced calcite precipitation at the eld scale. Int. J.
  Greenh. Gas Control 106, 103256.
  https://doi.org/10.1016/j.ijggc.2021.103256.{]} and {[}Landa-Marbán,
  D., Kumar, K., Tveit, S., Gasda, S.E., 2021. Numerical studies of CO2
  leakage remediation by micp-based plugging technology. In: Røkke, N.A.
  and Knuutila, H.K. (Eds) Short Papers from the 11th International
  Trondheim CCS conference, ISBN: 978-82-536-1714-5, 284-290.{]}.
\end{itemize}

\subsubsection{Other OPM Flow Specific
Enhancements:}\label{other-opm-flow-specific-enhancements}

\begin{itemize}
\tightlist
\item
  The Salt Precipitation and Water Evaporation Model: The model accounts
  for water evaporation of the in situ brine and the resulting
  increasing dissolved salt concentration. Once the salt solubility
  limit is exceeded, salt precipitation occurs, that results in a
  decrease in porosity and permeability, and consequently a decrease in
  gas production. This feature is primary used in modeling gas fields
  under depletion.
\item
  Thermal Model Enhancements: The energy black-oil implementation in OPM
  Flow is a mixture of the commercial simulators black-oil temperature
  and compositional thermal options. These commercial simulator options
  are two separate modeling facilities in the commercial simulator. OPM
  Flow's thermal implementation is based on solving the energy equation
  fully coupled with the black-oil equations so the results are not
  directly equivalent to the commercial simulator's black-oil
  temperature or compositional thermal formulations. In addition, OPM
  Flow can take into account the Joule--Thomson effect for oil, gas, and
  water. This behavior is especially important in gas wells, when a real
  gas expands, resulting in the temperature of the gas dropping.
\item
  Biofilm Effects in Underground Applications: Biofilm-related effects
  in subsurface applications such as hydrogen storage include reduced
  injectivity and hydrogen loss. The conceptual model includes the
  following mechanisms: (1) Biofilm is present in the storage site prior
  to injection and (2) the biofilm consumes injected H2/CO2, leading to
  clogging effects.
\end{itemize}

\subsubsection{Description of Geology:}\label{description-of-geology}

\begin{itemize}
\tightlist
\item
  Rectilinear and fully-unstructured grid.
\item
  Corner-point grids from the commercial simulator input, including
  fault and region multipliers, minimum pore volume, pinch outs, etc.
\item
  Spider grids that are similar to radial grids for modeling near well
  bore effects.
\item
  Analytical and numerical aquifers.
\end{itemize}

\subsubsection{Well and Group Controls:}\label{well-and-group-controls}

\begin{itemize}
\item
  Bottom-hole pressure and surface/reservoir rate.
\item
  Group controls, including voidage replacement, network model, gas lift
  optimization, guide rate control for producers and injection groups
  and well, etc.
\item
  Shut/stop/open individual completions.
\item
  History-matching wells.
\item
  Multi-segment wells together with various inflow control devices
  (``ICDs'').
\item
  Well lists.
\item
  Action conditions and command processing via the ACTIONX and UDQ
  keywords employed in the commercial simulator.
\item
  Python interface to control and process simulator data.
\end{itemize}

\subsubsection{Input and Output:}\label{input-and-output}

\begin{itemize}
\tightlist
\item
  General reader/parser for Eclipse input decks.
\item
  XML-based or simple text-format input of additional parameters.
\item
  Flexible output of summary and restart files in he commercial
  simulator's format.
\item
  Logging to terminal and print file.
\item
  Front-end graphical interface for running the simulator with various
  tools, including: case compression and uncompression, keyword
  generation etc.
\end{itemize}

\subsubsection{Simulation Technology:}\label{simulation-technology}

\begin{itemize}
\tightlist
\item
  Fully-implicit in time.
\item
  Automatic differentiation.
\item
  Two-point flux approximation in space with upstream-mobility
  weighting.
\item
  Flexible assembly through the use of automatic differentiation.
\item
  Bi-conjugate Gradient method (BiCG-stab) and a restarted Generalized
  Minimal Residual (``GMRes'') iterative linear solvers with ILU0 and
  Constraint Pressure Residual (``CPR'') preconditioners.
\item
  Adaptive step-size controls.
\item
  Parallel processing.
\end{itemize}

The following section briefly describe the simulator's formulation,
equation, and the various numerical procedures available to solve the
equations.

\begin{tcolorbox}[enhanced jigsaw, title=\textcolor{quarto-callout-note-color}{\faInfo}\hspace{0.5em}{Note}, coltitle=black, breakable, toprule=.15mm, colbacktitle=quarto-callout-note-color!10!white, leftrule=.75mm, opacitybacktitle=0.6, colback=white, bottomtitle=1mm, toptitle=1mm, titlerule=0mm, left=2mm, colframe=quarto-callout-note-color-frame, arc=.35mm, rightrule=.15mm, bottomrule=.15mm, opacityback=0]

The material in the following sections is taken from Rasmussen et
al.~{[}Rasmussen, A.F., Sandve, T.H., Bao, K., Lauser, A., Hove, J.,
Skaflestad, B., Klöfkorn, R., Blatt, M., Rustad, A.B., Sævareid, O. et
al.~{[}2019{]}. The Open Porous Media Flow Reservoir Simulator.
https://doi.org/10.1016/j.camwa.2020.05.014.{]}, and the reader is
encouraged to read the full paper to obtain a fuller understanding of
the simulator's formulation and the resulting equations.

\end{tcolorbox}

\section{Black-Oil Model Equations}\label{black-oil-model-equations}

The black-oil equations constitute the most widely used flow model in
the simulation of hydrocarbon reservoirs. The model is based on the
premise that we have three different fluid phases (aqueous, oleic, and
gaseous), and three (pseudo) components: water, oil, and gas. Oil and
gas are not single hydrocarbon species, but represent all hydrocarbon
species that exist in liquid and vapor form at surface conditions.
Mixing is allowed, in the sense that both oil and gas can be found in
the oleic phase, in the gaseous phase, or in both. The amount of
dissolved gas in the oleic phase and vaporized oil in the gaseous phase
must be kept track of on a grid cell basis. In the following discussions
the subscripts w represents the water, o the oil, and g the gas
quantities related to each of the phases or components.

The black-oil equations can be deduced from conservation of mass for
each component with suitable closure relationships such as Darcy's law
and initial and boundary conditions. The equations are discretized in
space with an upstream finite-volume scheme using a two-point flux
approximation, and in time using an implicit (backward) Euler scheme.
The resulting equations are solved simultaneously in a fully implicit
formulation by a Newton-type linearization with a properly
preconditioned, iterative linear solver (see Rasmussen et
al.~{[}Rasmussen, A.F., Sandve, T.H., Bao, K., Lauser, A., Hove, J.,
Skaflestad, B., Klöfkorn, R., Blatt, M., Rustad, A.B., Sævareid, O. et
al.~{[}2019{]}. The Open Porous Media Flow Reservoir Simulator.
https://doi.org/10.1016/j.camwa.2020.05.014.{]}).

The conservation laws form a system of partial differential equations,
one for each pseudo component α:

\begin{equation}\phantomsection\label{eq-1-1}{
\frac{\partial }{\partial t}\left({ϕ}_{\mathit{ref}}{A}_{α}\right) + \nabla \cdot \mathbf{{\text{u}}_{\mathrm{α}}} + {q}_{α} = 0
}\end{equation}

where the accumulation terms and fluxes are given by:

\begin{equation}\phantomsection\label{eq-1-2}{
{A}_{w}={m}_{ϕ}{b}_{w}{s}_{w}
}\end{equation}

\begin{equation}\phantomsection\label{eq-1-3}{
{A}_{o}={m}_{ϕ}\left({b}_{o}{s}_{o}+{r}_{\mathit{og}}{b}_{g}{s}_{g}\right)
}\end{equation}

\begin{equation}\phantomsection\label{eq-1-4}{
{A}_{g}={m}_{ϕ}\left({b}_{g}{s}_{g}+{r}_{\mathit{go}}{b}_{o}{s}_{o}\right)
}\end{equation}

and the phase fluxes are given by Darcy's law:

\begin{equation}\phantomsection\label{eq-1-5}{
\mathbf{{\text{v}}_{\mathrm{α}}} = -{λ}_{α}\mathbf{\text{K}}\left(\nabla {p}_{α} - {ρ}_{α}\mathbf{\text{g}}\right)
}\end{equation}

In addition, the following closure relations should hold:

\begin{equation}\phantomsection\label{eq-1-6}{
{s}_{w} + {s}_{o} + {s}_{g} = 1
}\end{equation}

\begin{equation}\phantomsection\label{eq-1-7}{
{p}_{c,\mathit{ow}} = {p}_{o} - {p}_{w}
}\end{equation}

\begin{equation}\phantomsection\label{eq-1-8}{
{p}_{c,\mathit{og}} = {p}_{o} - {p}_{g}
}\end{equation}

where:

\(ϕ\) = porosity, can be a function of (oil)
pressure:\(ϕ = {m}_{ϕ}\left({p}_{o}\right){ϕ}_{\mathit{ref}}\).

\({m}_{ϕ}\) = pore volume multiplier as function of pressure.

\({ϕ}_{\mathit{ref}}\) = reference porosity, a constant in time but
varying in space.

\({p}_{α}\) = pressure for phase\(α\).

\({b}_{α}\) = shrinkage/expansion factor for phase\(α\)defined as the
ratio of the surface

volume at standard conditions to the reservoir volume for a given amount
of
fluid:\({b}_{α}={V}_{\mathit{surface};α} / {V}_{\mathit{reservoir};α}\).
For the oil phase,\({b}_{o}\)is called the shrinkage factor, whereas for
gas\({b}_{g}\) is called the expansion factor. The reciprocal quantity
is called the formation volume factor, and is usually denoted with a
capital B:\({B}_{α}=1 / {b}_{α}\). The variable is usually a function of
phase pressure and
composition:\({b}_{α}={b}_{α}\left({p}_{α},{r}_{\mathit{go}},{r}_{\mathit{og}}\right)\).

\({r}_{\mathit{go}}\) = gas-oil ratio (``GOR''), ratio of dissolved gas
to oil in the oleic phase,

commonly called \({r}_{s}\)or \({R}_{s}\) in other literature.

\({r}_{\mathit{og}}\) = oil-gas ratio or condensate-gas ratio (``CGR''),
ratio of vaporized oil to gas

in the gaseous phase, commonly called\({r}_{v}\)or \({R}_{v}\) in other
literature.

\({s}_{α}\) = saturation of phase \(α\), the pore volume fraction
occupied by the phase.

The saturations of all phases sum to one.

\({p}_{c,αβ}\) = capillary pressure between phases\(α\)and\(β\).
Typically a function of

saturation:\({p}_{c,αβ}={p}_{c,αβ}\left({S}_{α}\right)\).

\(\mathbf{\text{K}}\) = permeability of the porous medium.

\({k}_{r;α}\) = relative permeability for phase\(α\)modeling the
reduction in effective

permeability for a fluid phase in the presence of other phases.
Typically a function of saturation.

\({μ}_{α}\) = viscosity of phase\(α\), typically a function of phase
pressure and

composition.

\({λ}_{α}\) = mobility of phase\(α\), given
by\({λ}_{α}={k}_{r,α} / {μ}_{α}\).

\(\mathbf{{\text{v}}_{\mathrm{α}}}\) = velocity of phase\(α\).

\(\mathbf{{\text{u}}_{\mathrm{α}}}\) = velocity of component\(α\).

\({ρ}_{S,α}\) = surface density of phase\(α\)at one atmosphere, a given
constant.

\({ρ}_{α}\) = density of phase \(α\)in the reservoir. For
water,\({ρ}_{w}={b}_{w}{ρ}_{S,w}\). For the oleic

phase the relationship is more complex since it must include dissolved
gas:
\({ρ}_{o}={b}_{o}\left({ρ}_{S,o}+{r}_{\mathit{go}}{ρ}_{S,g}\right)\),
and similar for the gaseous phase we have:
\({ρ}_{g}={b}_{g}\left({ρ}_{S,g}+{r}_{\mathit{og}}{ρ}_{S,o}\right)\).

\(\mathbf{\text{g}}\) = gravitational acceleration vector.

\({q}_{α}\) = well out flux density of pseudo component\(α\), (negative
for well in-flow).

The form of this term depends on the well model used, either the
standard well model or the multi-segment model.

Equations (1.1) to (1.8) are then transformed to the discrete form using
a forward finite difference scheme. In the following, let
subscript\(i\)denote a discrete quantity defined in cell\(i\)and
subscript\(\mathit{ij}\)denote a discrete quantity defined at the
connection between two cells\(i\)and\(j\). Two cells share a connection
if they are adjacent geometrically in the computational grid or share an
explicit Non-Neighbor Connection (``NNC'') {[}Connections between cells
due to geometric adjacency are automatically calculated by the
simulator, including across faults in the grid. In addition, the user
can specify extra connections between arbitrary cells using the NNC
keyword in the GRID section.{]}. For example\({υ}_{o,\mathit{ij}}\)is
the oil flux from cell\(i\)to cell\(j\). For oriented quantities such as
fluxes, the orientation is taken to be from cell\(i\) to cell\(j\), and
the quantity is skew-symmetric
(\({υ}_{o,\mathit{ij}}=-{υ}_{o,\mathit{ij}}\)), whereas non-oriented
quantities such as the transmissibility\({T}_{\mathit{ij}}\)are
symmetric (\({T}_{\mathit{ij}}={T}_{\mathit{ji}}\)). Quantities with
superscript 0 are taken at the start of the discrete time step, other
quantities are at the end of the time step. Superscripts or subscripts
applied to an expression in parenthesis apply to each element in the
expression.

The discretized equations and residuals are, for each
pseudo-component\(α\) and cell\(i\):

\begin{equation}\phantomsection\label{eq-1-9}{
{R}_{α,i}=\frac{{ϕ}_{\mathit{ref},i}{V}_{i}}{\mathrm{Δ}t}\left({A}_{α,i}-{{A}^{0}}_{α,i}\right) +\underset{jϵC(i)}{\sum }{u}_{α,\mathit{ij}}+{q}_{α,i}=0
}\end{equation}

where the accumulation terms are the same as equations (1.2) to (1.4)
and fluxes \({u}_{α}\)are of similar form to the velocities,
\(\mathbf{{\text{u}}_{\mathrm{α}}}\) in the aforementioned equations;
thus we have

\begin{equation}\phantomsection\label{eq-1-10}{
{A}_{w i}={\left({m}_{ϕ}{b}_{w}{s}_{w}\right)}_{i}
}\end{equation}

\begin{equation}\phantomsection\label{eq-1-11}{
{A}_{o i}={\left({m}_{ϕ}\left({b}_{o}{s}_{o}+{r}_{\mathit{og}}{b}_{g}{s}_{g}\right)\right)}_{i}
}\end{equation}

\begin{equation}\phantomsection\label{eq-1-12}{
{A}_{g i}={\left({m}_{ϕ}\left({b}_{g}{s}_{g}+{r}_{\mathit{go}}{b}_{o}{s}_{o}\right)\right)}_{i}
}\end{equation}

In addition, equations (1.6) to (1.8) also hold for each cell i, that
is:

\begin{equation}\phantomsection\label{eq-1-13}{
{\left( {s}_{w} + {s}_{o} + {s}_{g} \right)}_{i} = 1
}\end{equation}

\begin{equation}\phantomsection\label{eq-1-14}{
{\left( {p}_{c,\mathit{ow}} \right)}_{i} = {\left( {p}_{o} - {p}_{w} \right)}_{i}
}\end{equation}

\begin{equation}\phantomsection\label{eq-1-15}{
{\left( {p}_{c,\mathit{og}} \right)}_{i} = {\left( {p}_{o} - {p}_{g} \right)}_{i}
}\end{equation}

Finally, the fluxes are given for each connection\(\mathit{ij}\)by the
following relationships:

\begin{equation}\phantomsection\label{eq-1-16}{
{\left({b}_{α}{υ}_{α}\right)}_{\mathit{ij}} = {\left({b}_{α}{λ}_{α}{m}_{T}\right)}_{{U}_{(α,\mathit{ij})}}{T}_{\mathit{ij}}\mathrm{Δ}{\mathrm{Φ}}_{α,\mathit{ij}}
}\end{equation}

\begin{equation}\phantomsection\label{eq-1-17}{
{\left({r}_{βα}{b}_{α}{υ}_{α}\right)}_{\mathit{ij}} = {\left({r}_{βα}{b}_{α}{λ}_{α}{m}_{T}\right)}_{{U}_{(α,\mathit{ij})}}{T}_{\mathit{ij}}\mathrm{Δ}{\mathrm{Φ}}_{α,\mathit{ij}}
}\end{equation}

\begin{equation}\phantomsection\label{eq-1-18}{
\mathrm{Δ}{\mathrm{Φ}}_{α,\mathit{ij}} = {p}_{α,i}-{p}_{α,j}-g{ρ}_{α,\mathit{ij}}\left({z}_{i}-{z}_{j}\right)
}\end{equation}

\begin{equation}\phantomsection\label{eq-1-19}{
{ρ}_{α,\mathit{ij}} = \frac{\left({ρ}_{α,i}+{ρ}_{α,j}\right)}{2}
}\end{equation}

\begin{equation}\phantomsection\label{eq-1-20}{
U\left(α,\mathit{ij}\right) = \left{\begin{matrix}i, \mathrm{Δ}{\mathrm{Φ}}_{α,\mathit{ij}}\geq 0 \\ j, \mathrm{Δ}{\mathrm{Φ}}_{α,\mathit{ij}}\geq 0\end{matrix}
}\end{equation}

where:

\(V\) = cell volume.

\(\mathrm{Δ}t\) = time step length for the current Euler step.

\({υ}_{α}\) = volume flux of phase \(α\)(oriented quantity).

\({u}_{α}\) = surface volume flux of pseudo component\(α\) (oriented
quantity).

\({T}_{\mathit{ij}}\) = transmissibility factor for a connection,
derived from permeability.

\({m}_{T}\) = transmissibility multiplier as function of pressure.

\(C(i)\) = connections from cell \(i\), i. e. the set of cells connected
to it.

\(U(α,\mathit{ij})\) = upwind cell for phase\(α\)for the connection
between cells\(i\)and\(j\).

\(\mathrm{Δ}{\mathrm{Φ}}_{(α,\mathit{ij})}\) = potential difference for
phase\(α\)for the connection between cells \(i\)

and\(j\)(oriented quantity).

\(g\) = gravitational acceleration in the z-direction.

\({z}_{i}\) = depth of center of cell \(i\).

For a grid containing a total of n active cells, the material balance
relationship expressed by equation (1.9) results in a system of 3n
equations for the three phase formulation (oil, water, and gas).
Similarly, for the two phase model, oil-water or gas-water, there are 2n
equations to be solved. In addition, there are also the sink and source
terms associated with the wells. The solution variables for the system
of equations is dependent on the state of a cell, basically, for:

\begin{verbatim}
- Non-miscible: Here pressure of one of the phases (which we take to be oil pressure${p}_{o}$), ${s}_{w},$ and ${s}_{g}$ as our primary variables. The pressure ${p}_{o}$ will then behave mostly in an elliptic fashion, whereas the saturations ${s}_{w}$ and ${s}_{g}$are more hyperbolic.
- Miscible flow: Since the gaseous phase may disappear if all the gas dissolves into the oleic phase, and similarly the oleic phase may disappear if all the oil vaporizes into the gaseous phase, one cannot always use ${s}_{g}$ for our third variable. Instead we use the third variable to track the composition of the phase that has not disappeared. We therefore use either ${s}_{g}$,${r}_{\mathit{go}}$, or ${r}_{\mathit{og}}$ as our third primary variable, depending on the fluid state, and we call that variable x:
\end{verbatim}

The next two sections describe the formulation of the two well models
implemented in OPM flow: the Standard Well Model and the Multi-Segment
Well Model.

\section{Standard Well Model}\label{standard-well-model}

For the standard well model, in a three phase black-oil system, there
are four primary variables, three fluid variables:\({Q}_{t}\),
\({F}_{w}\), and\({F}_{g}\), which describe the fluid composition within
the wellbore {[}J. Holmes, Enhancements to the strongly coupled, fully
implicit well model: wellbore crossflow modeling and collective well
control, in: SPE Reservoir Simulation Symposium, Society of Petroleum
Engineers, 1983.{]}, as depicted in the following equation:

\begin{equation}\phantomsection\label{eq-1-22}{
{Q}_{t} = \underset{α\in {o,g,w}}{\sum }{g}_{α}{Q}_{α}, {F}_{w} = \frac{{g}_{w}{Q}_{w}}{{Q}_{t}}, {F}_{g} = \frac{{g}_{g}{Q}_{g}}{{Q}_{t}}
}\end{equation}

where:

\({Q}_{t}\) = the weighted total flow rate,

\({F}_{w}\) = the weighted fraction of water, and

\({F}_{g}\) = the weighted fraction of gas.

Here, \({Q}_{α}\) is the flow rate of component \(α\)under surface
conditions and \({g}_{α}\)is a weighting factor. For gas, this factor is
typically chosen to be a small value, for example, 0.01 to avoid gas
fractions close to unity12.

The fourth variable is the bottom-hole pressure
(\({p}_{\mathit{bhp}}\)), that is the pressure in the wellbore at the
datum depth.

Volumetric inflow rates at reservoir conditions are calculated as:

\begin{equation}\phantomsection\label{eq-1-23}{
{{q}^{r}}_{α,j}={T}_{w,j} {M}_{α,j}\left[{p}_{j}-({p}_{\mathit{bhp},w} + {h}_{w,j})\right]
}\end{equation}

where:

\({{q}^{r}}_{α,j}\) = is the flow rate of phase \(α\) through connection
\(j\). The rate is negative

for flow from the wellbore to reservoir, and positive for flow into the
opposite direction.

\({T}_{w,j}\) = connection transmissibility factor.

\({M}_{α,j}\) = the mobility for phase \(α\) at the connection j. For
injection wells, the total

mobility of the injecting phase is used, as per Holmes12.

\({p}_{j}\) = pressure of the grid block that contains the connection j.

\({p}_{\mathit{bhp},w}\) = the bottom-hole pressure of well\(w\).

\({h}_{w,j}\) = pressure difference within the wellbore between
connection \(j\) and the

well's bottom-hole datum depth. The pressure differences \({h}_{w,j}\)
between connection\(j\) and the datum point are computed explicitly
based on the fluid composition in the wellbore at the start of the time
step.

To keep the system closed, we introduce conservation equations for each
component using:

\begin{equation}\phantomsection\label{eq-1-24}{
{R}_{α,w}=\frac{{A}_{α,w}-{A}_{α,w}^{0}}{\mathrm{Δ}t} + {Q}_{α} - \underset{j\in C(w)}{\sum }{q}_{α,j} = 0
}\end{equation}

Equation (1.24) is solved in a fully implicit and coupled way with the
black-oil equations, as per equation (1.9)). Here, \(C(w)\) is the set
of connections of the well \(w\), \({q}_{α,j}\) is the flow rate of
phase \(α\) through connection \(j\) under surface condition, which can
be calculated from \({{q}^{r}}_{α,j}\)defined in equation (1.23); the
relationship is similar to the one between component and mass fluxes, as
per equations (1.2) through (1.4). The storage term \({A}_{α,w}\)
describes the amount of component \(α\) in the wellbore (a small volume,
that is 0.1 ft3 as per Holmes12), and it is introduced for better
stability of the well solution.

In addition, we need equations that model how the wells are controlled.
For wells controlled by a prescribed bottom-hole pressure
(\(\mathit{bhp}\)) target, we have:

\begin{equation}\phantomsection\label{eq-1-25}{
{R}_{c,w} = {p}_{\mathit{bhp},w}-{p}_{\mathit{bhp},w}^{\mathit{target}} =0
}\end{equation}

whereas for a rate-controlled well we have:

\begin{equation}\phantomsection\label{eq-1-26}{
{R}_{c,w} ={Q}_{α} - {Q}_{α}^{\mathit{target}} = 0
}\end{equation}

where\({Q}_{α}^{\mathit{target}}\) is the desired surface-volume rate of
the controlled component \(α\), typically oil or gas rate for a
production well. The Standard Well model handles cross-flow, that is
flow exiting from one well connection and entering another connection;
however, this model assumes that the fluid composition is uniform
throughout the wellbore (given by Fg and Fw), and thus the same
composition will be injected through all injecting connections. When
this is not a reasonable assumption, then the Multi-Segment model should
be used instead, which is the topic of the next section.

\section{Multi-Segment Well Model}\label{multi-segment-well-model}

\pandocbounded{\includegraphics[keepaspectratio]{chapters/images/Frame70_8c99a481697c.png}}
\pandocbounded{\includegraphics[keepaspectratio]{chapters/images/Image33_8c99a481697c.png}}

Multi-segment wells were introduced to handle more complex well
configurations, including for example horizontal, multi-lateral and fish
bone well types, as well Inflow Control Devices (``ICD'') {[}J. Holmes,
T. Barkve, O. Lund, Application of a multisegment well model to simulate
ow in advanced wells, in: European Petroleum Conference, Society of
Petroleum Engineers, 1998.{]}. This model divides the wellbore into a
number of arbitrary segments. Each segment has a segment node and a flow
path to the neighboring segment in the direction towards the well head,
that defines as the outlet segment (Figure 1.2). Most segments also have
inlet segment neighbors in the direction away from the well head. For
multi-lateral wells, a segment node must be placed at the branch
junction. Segments at branch junction points have additional inlet
segments and segments at the end of branches do not have inlet segments.
The top segment of each well does not have an outlet segment. The
pressure at this segment is the bottom-hole pressure of the well,
whereas the component rates equal the component rates for the well as a
whole. For example, Figure 1.3 on the following page, shows the main
wellbore branch together with four offsetting branches.

Each segment has the node pressure p as a primary variable in addition
to the primary variables \({Q}_{t}\), \({F}_{w}\), and \({F}_{g}\)
introduced for the standard well model, as per equation (1.22).
Similarly, the component equations, (1.24), are extended to include the
additional \({Q}_{α,i}\) terms that represent the flow rate from the
inlet segments \(I(n)\)of segment \(n\), as well as the inflow from the
reservoir through the segment's connections \(C(n)\). Thus, the
equivalent standard well equation (1.24) for multi-segment wells is:

\begin{equation}\phantomsection\label{eq-1-27}{
{R}_{α,n}=\frac{{A}_{α,n}-{A}_{α,n}^{0}}{\mathrm{Δ}t} + \underset{i\in I(n)}{\sum }{Q}_{α,i} - \underset{j\in C(w)}{\sum }{q}_{α,j} + {Q}_{α,n} = 0
}\end{equation}

Each segment can thus receive inflow from more than one connection or no
connection through the segment node, whereas each connection can only
contribute to one segment. In the standard well model, the connection is
located at the centroid of the grid block. In the multi-segment well
model, the connection and

segment node can be located at any place within the grid block. The
inflow calculation of the connection in the multi-segment well model is
as follows:

\begin{equation}\phantomsection\label{eq-1-28}{
{{q}^{r}}_{α,j} = {T}_{w,j} {M}_{α,j}\left({p}_{j} + {H}_{c,j} - {p}_{n} - {H}_{\mathit{nc}}\right)
}\end{equation}

where:

\({p}_{j}\) = the pressure at the cell center of the grid block that
contains connection

\(j\),

\({H}_{c,j}\) = the hydrostatic pressure difference between the cell
center and the

connection,

\({p}_{n}\) = the pressure of segment \(n\), and

\({H}_{\mathit{nc}}\) = the hydrostatic pressure difference between the
connection and the

segment.

The last equation for each segment describes the pressure relationships
between the segment \(n\) and its outlet segment \(m\):

\begin{equation}\phantomsection\label{eq-1-29}{
{R}_{p,n} = {p}_{n} - {p}_{m} - {H}_{h} - {H}_{f} - {H}_{a} = 0
}\end{equation}

Here, the \(H\) terms represent the hydrostatic, frictional, and
acceleration pressure drops between the segments.

The top segment does not have an outlet segment, and hence the pressure
equation is replaced by a well control equation, which is the same as
for the standard well model, that is either equation (1.25) or (1.26).

As with the standard well model, the well equations are solved with the
reservoir mass conservation equations in a coupled, fully implicit way.
The well equations are added to the global system using a Schur
{[}Zhang, Fuzhen (2005). Zhang, Fuzhen (ed.). The Schur Complement and
Its Applications. Numerical Methods and Algorithms. Vol. 4. Springer.
doi:10.1007/b105056. ISBN 0-387-24271-6.{]} complement approach:

\begin{equation}\phantomsection\label{eq-1-30}{
(A - \underset{w}{\sum }{C}_{w}{D}_{w}^{-1}{B}_{w}){x}_{r} = {R}_{r} - \underset{w}{\sum }{C}_{w}{{D}_{w}}^{-1}{R}_{w}
}\end{equation}

\[
{x}_{w} = {{D}_{w}}^{-1}({R}_{w} - {B}_{w}{x}_{r})
\]

Figure 1.3 depicts a multi-segment well with the main wellbore branch as
the vertical segment, together with four offsetting horizontal branches.

\pandocbounded{\includegraphics[keepaspectratio]{chapters/images/Frame125_827ae00c3b64.png}}
\pandocbounded{\includegraphics[keepaspectratio]{chapters/images/Image34_827ae00c3b64.png}}

\section{Salt Precipitation Model}\label{salt-precipitation-model}

OPM Flow's salt precipitation and water evaporation model is an
extension of the black-oil model with a water in gas component, plus
addition of a solid salt phase component to the salt transport equation.
The effects of porosity and permeability reduction effects are also
considered. The oil, gas, and water and salt conservation equations
including these extensions are given by equations (1.31) to (1.34).

Oil:

\begin{equation}\phantomsection\label{eq-1-31}{
\frac{\partial }{\partial t}\left[ϕ\left({b}_{o}{s}_{o} + {r}_{\mathit{og}}{b}_{g}{s}_{g}\right)\right] + \nabla \cdot \left({b}_{o}{ν}_{o} + {r}_{\mathit{og}}{b}_{g}{ν}_{g}\right) + {q}_{o} = 0
}\end{equation}

Gas:

\begin{equation}\phantomsection\label{eq-1-32}{
\frac{\partial }{\partial t}\left[ϕ\left({b}_{g}{s}_{g} + {r}_{\mathit{go}}{b}_{o}{s}_{o}\right)\right] + \nabla \cdot \left({b}_{g}{ν}_{g} + {r}_{\mathit{go}}{b}_{o}{ν}_{o}\right) + {q}_{g} = 0
}\end{equation}

Water:

\begin{equation}\phantomsection\label{eq-1-33}{
\frac{\partial }{\partial t}\left[ϕ\left({b}_{w}{s}_{w} + {r}_{\mathit{wg}}{b}_{g}{s}_{g}\right)\right] + \nabla \cdot \left({b}_{w}{ν}_{w} + {r}_{\mathit{wg}}{b}_{g}{ν}_{g}\right) + {q}_{w} = 0
}\end{equation}

Salt:

\begin{equation}\phantomsection\label{eq-1-34}{
\frac{\partial }{\partial t}\left[ϕ{b}_{w}{s}_{w}{c}_{w}^{\mathit{salt}} + {m}_{ϕ}{ϕ}_{0}{s}_{s}{ρ}^{\mathit{salt}}\right] + \nabla \cdot \left({c}_{w}^{\mathit{salt}}{b}_{w}{ν}_{w}\right) + {c}_{w}^{\mathit{salt}}{q}_{w} = 0
}\end{equation}

where:

\(ϕ\) = porosity.

\({m}_{ϕ}\) = pore volume multiplier as function of pressure.

\({ϕ}_{0}\) = reference porosity, a constant in time but varying in
space.

\({p}_{α}\) = pressure for phase\(α\).

\({b}_{α}\) = shrinkage/expansion factor for phase\(α\)defined as the
ratio of the surface

volume at standard conditions to the reservoir volume for a given amount
of
fluid:\({b}_{α}={V}_{\mathit{surface};α} / {V}_{\mathit{reservoir};α}\).
For the oil phase,\({b}_{o}\)is called the shrinkage factor, whereas for
gas \({b}_{g}\) is called the expansion factor. The reciprocal quantity
is called the formation volume factor, and is usually denoted with a
capital B: \({B}_{α}=1 / {b}_{α}\). The variable is usually a function
of phase pressure and composition.

\({r}_{\mathit{go}}\) = gas-oil ratio (``GOR''), ratio of dissolved gas
to oil in the oleic phase,

commonly called \({r}_{s}\)or \({R}_{s}\) in other literature.

\({r}_{\mathit{og}}\) = oil-gas ratio or condensate-gas ratio (``CGR''),
ratio of vaporized oil to gas

in the gaseous phase, commonly called\({r}_{v}\)or \({R}_{v}\) in other
literature.

\({r}_{\mathit{wg}}\) = evaporated water in gas ratio.

\({s}_{α}\) = saturation of phase \(α\), the pore volume fraction
occupied by the phase.

\({s}_{s}\) = saturation of precipitated salt (volume fraction).

\({ρ}_{α}\) = surface density of phase\(α\)at one atmosphere, a given
constant.

\({ρ}^{\mathit{salt}}\) = density of the solid salt.

\({q}_{α}\) = well out flow of component\(α\), (negative for well
in-flow).

The form of this term depends on the well model used, either the
standard well model or the multi-segment model.

\({c}_{w}^{\mathit{salt}}\) = salt concentration in water (mass per unit
volume).

The extension has two extra primary variable switching mechanisms. One
between \({s}_{w}\) and \({r}_{\mathit{wg}}\), when water disappears and
reappears, and one between \({c}_{w}^{\mathit{salt}}\) and \({s}_{s}\)
when the salt concentration reaches the (user defined) solubility limit
and salt precipitates or when solid salt dissolves. Furthermore, in the
new formulation porosity becomes:

\begin{equation}\phantomsection\label{eq-1-35}{
ϕ = \left(1 - {s}_{s}\right){m}_{ϕ}{ϕ}_{0}
}\end{equation}

The resulting (reduced) permeability, \(k\), is based on a user
specified permeability-porosity relation \(k=k(ϕ)\).

Both for dry and humid gas and wet and humid gas the PVT data can be
provided using dedicated PVT tables. For the capillary pressure data it
is recommended to extend its values towards zero water saturation to
permit water saturations below the irreducible water saturation due to
water evaporation.

For further details see Machado et al.~{[}1C. G. Machado, P. Egberts, J.
Alvestad, O. S. Hustad. Salt Precipitation and Water Evaporation
Modelling in a Black-Oil Reservoir Simulator. SPE Reservoir Simulation
Conference 2023 (SPE-212257).{]}

\section{Numerical Solution of
Equations}\label{numerical-solution-of-equations}

Efficiently solving the aforementioned equations is non-trivial, and
involves a step iterative process that consists of an outer,
Newton-Raphson type method, of a series of iterations, and an inner set
of linear iterations, as depicted in Figure 1.4, which shows a
simplified schematic of the iterative scheme.

\pandocbounded{\includegraphics[keepaspectratio]{chapters/images/Frame126_8bd907266a91.png}}
\pandocbounded{\includegraphics[keepaspectratio]{chapters/images/Image35_8bd907266a91.png}}

The next two sections attempt to give an overview of the two numerical
schemes.

\subsection{Newton-Raphson Solver (Outer
Iteration)}\label{newton-raphson-solver-outer-iteration}

The aforementioned reservoir and well equations define a large system of
(fully implicit) non-linear equations, which one can write in compact
residual form as R(y) = 0, where y is the vector of primary variables,
\({p}_{o}, {s}_{w}, x\). This system is solved using a Newton-Raphson
type method.

Let \({y}_{n}=({p}_{o}, {s}_{w}, x)\) be the primary variables after
\(n\) Newton iterations, and given the initial state \({y}_{0}\), the
solution of \(R(y)=0\) can be found by iteratively solving:

\begin{equation}\phantomsection\label{eq-1-36}{
J({y}_{n})\left({y}_{n+1} - {y}_{n}\right) = -R({y}_{n})
}\end{equation}

until \(R({y}_{n})\) satisfies a given tolerance. Creating the Jacobian
matrix \(J ({y}_{n})\) and solving this linearized problem is the core
computational task of the simulator, and its performance depends largely
on how this part of the code is programmed.

The coefficients of the linearized systems required by the non-linear
Newton solver,\(J ({y}_{n})\) in equation (1.36), have traditionally
been computed by evaluating closed-form expressions obtained by
differentiating the discretized flow equations analytically.
Differentiating these flow equations manually and programming the
resulting formulas is both time-consuming and error-prone. The problem
is particularly pronounced when extending an existing simulator with new
functional relationships. If a quantity is set to depend on a different
or new (primary) variable, the changes in derivatives will propagate
through the chain rule to the Jacobians of all quantities that depend on
the modified quantity. Thus, even very simple extensions in functional
dependencies may cause large modifications throughout the simulator
code.

Automatic Differentiation (``AD'') is a way to mitigate this problem
{[}R. D. Neidinger, Introduction to automatic differentiation and matlab
object-oriented programming, SIAM Rev.~52 (3) (2010) 545\{563.
doi:10.1137/080743627.{]} and {[}A. Griewank, A. Walther, Evaluating
Derivatives, Principles and Techniques of Algorithmic Differentiation,
second edition, SIAM, Philadelphia, 2008.{]}, and there are many
examples of AD tools used for scientific computing, such as Sacado
{[}Eric T. Phipps, Roscoe A. Bartlett, David M. Gay, Robert J. Hoekstra,
Large-Scale Transient Sensitivity Analysis of a Radiation-Damaged
Bipolar Junction Transistor via Automatic Differentiation Advances in
Automatic Differentiation, Springer, 2008.{]}, ADOL-C {[}A. Walther, A.
Griewank, Getting started with ADOL-C, in: U. Naumann, O. Schenk (Eds.),
Combinatorial Scientific Computing, Chapman-Hall CRC Computational
Science, 2012, Ch. 7, pp.~181-202.{]} or ADETL {[}Rami M. Younis, Modern
Advances In Software And Solution Algorithms For Reservoir Simulation,
PhD thesis, Stanford University, (August 2011){]}; with the latter used
for the AD-GPRS {[}AD-GRPS, Stanford Earth, Stanford University,
https://suetri-b.stanford.edu/research/ad-gprs.{]} reservoir simulator
at Stanford University.

The OPM Flow's implementation of AD introduces a class that mimics the
behavior of the built-in floating point types of C++ (double and float)
as closely as possible. The AD object contains a value and a fixed
number of derivatives and defines all basic arithmetic operators as well
as common mathematical functions. Also, to allow easy comparisons with
code that uses standard floating-point objects, all functions that work
on AD-objects can also be used with objects of built-in floating-point
types. For further detailed information on OPM Flow's implementation of
AD, again see Rasmussen et al.~{[}Rasmussen, A.F., Sandve, T.H., Bao,
K., Lauser, A., Hove, J., Skaflestad, B., Klöfkorn, R., Blatt, M.,
Rustad, A.B., Sævareid, O. et al.~{[}2019{]}. The Open Porous Media Flow
Reservoir Simulator. https://doi.org/10.1016/j.camwa.2020.05.014.{]}.

When solving the Newton-Raphson equation (1.36), the first iteration at
the new time step uses the solution at the end of the previous time step
(\({t}_{n}\)), as an initial estimate for the new time step
(\({t}_{n + 1}\)). Thus, the Newton-Raphson iteration is expected to
converge rapidly for sufficiently small time steps. For large time steps
and/or large changes in the solution, typically caused by changing well
controls, the Newton-Raphson method may converge poorly or not at all.
There are several means in which one can increase the probability of
convergence, namely:

\begin{itemize}
\tightlist
\item
  Restrict the change allowed per iteration and prevent the solution
  from jumping far across the boundary between saturated and
  undersaturated states in a single iteration. Thus, like the commercial
  simulator, OPM Flow implements the Appleyard {[}J. Appleyard, I. M.
  Cheshire, Nested factorization, in: 7th SPE Symposium on Reservoir
  Simulation, San Francisco, USA, Society of Petroleum Engineers, 1983.
  doi:10.2118/12264-MS.{]} chop technique, which is often sufficient for
  the method to converge. However, time step chops {[}A time step chop
  occurs when convergence is not achieved within the allowable number of
  iterations. For example, if the simulator was attempting to take a 30
  day time step and failed, then the time step would be cutback, to say
  ten days, and the iterations restarted. If this 10 day time step also
  failed to converge, then a second time step chop would occur reducing
  the new time step, to say three days. This process is repeated until
  either convergence is achieved, or the time step is too small to chop
  and the simulation stops.{]} may still occur in order to obtain
  convergence for some difficult cases, particularly for models with
  very fine grids.
\item
  Secondly, since the third primary variable,\(x\),may change solution
  state from one iteration to the next, this may disturb the
  convergence, and therefore a small threshold is applied (\(ε\)) to the
  variable-switching logic for equation (1.21), in order to prevent
  oscillation between states.
\item
  Thirdly, and frequently the main cause of non-convergence and the lack
  of smoothness of the residuals, is the discontinuity or inconsistent
  input data, particularly the relative permeability and capillary
  pressure functions. Thus, using Corey {[}Corey, A. T. : ``The
  Interrelation Between gas and Oil Relative Permeabilities'',
  Production Mon., 19. 38. (1954).{]} or Lomeland et al.~{[}Lomeland F.,
  2018.Overview Of The Let Family Of Versatile Correlations For Flow
  Functions. Paper SCA2018-056 presented at the International Symposium
  of the Society of Core Analysts held in Trondheim, Norway, 27-30
  August 2018.{]} type of curves for these data may well improve
  convergence performance, see Error: Reference source not found, Error:
  Reference source not found, and Error: Reference source not found, for
  OPM Flow's implementation of the Lomeland type of curves.
\item
  The final point that influences the convergence of the non-linear
  solver is the accuracy of the linear solvers, as discussed in the next
  section.
\end{itemize}

For more information about run time options and numerical controls, see
section Error: Reference source not found Error: Reference source not
found on how to invoke various numerical schemes via the OPM Flow
command line interface.

\subsection{Linear Solver (Inner
Iteration)}\label{linear-solver-inner-iteration}

As mentioned previously, efficiently solving the aforementioned
equations is non-trivial, and involves a step iterative process that
consists of an outer, Newton-Raphson type method, series of iterations,
and an inner set of linear iterations, as depicted in the detailed work
flow of the iterative process in Figure 1.5.

\pandocbounded{\includegraphics[keepaspectratio]{chapters/images/Frame131_3abbcb7f49bb.png}}
\pandocbounded{\includegraphics[keepaspectratio]{chapters/images/Image78_3abbcb7f49bb.png}}

Depending on the fluid properties and other reservoir specific
parameters, the system of equations as described by equation (1.9), that
is:

\begin{equation}\phantomsection\label{eq-1-9}{
{R}_{α,i}=\frac{{ϕ}_{\mathit{ref},i}{V}_{i}}{\mathrm{Δ}t}\left({A}_{α,i}-{{A}^{0}}_{α,i}\right) +\underset{jϵC(i)}{\sum }{u}_{α,\mathit{ij}}+{q}_{α,i}=0
}\end{equation}

could exhibit elliptic or degenerate parabolic behavior, or could even
be hyperbolic for specif setups. Therefore, solving the linearized
system stemming from these equations is challenging, since it is likely
non-symmetric and ill-conditioned. Thus, instead of solving equation
(1.36), that is:

\[
J({y}_{n})\left({y}_{n+1} - {y}_{n}\right) = -R({y}_{n})
\]

directly, the solution is approximated by an iterative linear solver.
Two variants are available in OPM Flow:

\begin{verbatim}
- a stabilized Bi-conjugate Gradient method (BiCG-stab), and
- a restarted Generalized Minimal Residual (GMRes) solver.
\end{verbatim}

By default, BiCG-stab is used. In addition, there are several options
available for preconditioning the linear system, such as the incomplete
lower triangular-upper triangular factorization with level of fill in 0
(ILU0). The default option is based on an algebraic multi-grid (AMG)
method for preconditioning, named the Constraint Pressure Residual (CPR)
preconditioner (originally introduced by Wallis {[}J. Wallis, Incomplete
Gaussian Elimination as a Preconditioning for Generalized Conjugate
Gradient Acceleration, in: 7th SPE Symposium on Reservoir Simulation,
San Francisco, USA, Society of Petroleum Engineers, 1983.
doi:10.2118/12265-MS.{]} using the approach of Scheichl et al.~{[}R.
Scheichl, M. Roland, J. Wendebourg, Decoupling and block preconditioning
for sedimentary basin simulations, Computational Geosciences 7 (2003)
295-318.{]}). The solvers and preconditioners are implemented in the
Dune module dune-istl and described by Blatt et al.~{[}M. Blatt, A
parallel algebraic multigrid method for elliptic problems with highly
discontinuous coefficients, Ph.D.~thesis, Ruprecht-Karls-Universitat
Heidelberg (2010).{]}, {[}M. Blatt, P. Bastian, The Iterative Solver
Template Library, in: B. Kåagström, E. Elmroth, J. Dongarra, J.
Waśniewski (Eds.), Applied Parallel Computing. State of the Art in
Scientific Computing, Vol. 4699 of Lecture Notes in Computer Science,
Springer, 2007, pp.~666-675.{]}, and {[}M. Blatt, P. Bastian, On the
generic parallelisation of iterative solvers for the finite element
method, Int. J. Comput. Sci. Engrg. 4 (1) (2008) 56-69.
doi:10.1504/IJCSE.2008.021112.{]}. The interface between OPM Flow and
the linear solver packages has been designed in a flexible way so that
besides dune-istl other linear algebra packages could be used.

For more information about run time options and numerical controls, see
section Error: Reference source not found Error: Reference source not
found on how to invoke various numerical schemes via the OPM Flow
command line interface. In addition, section Error: Reference source not
found Error: Reference source not found provides some practical advice
on how to address numerical performance issues.

\section{Formation Damage Models}\label{formation-damage-models}

\subsection{Incompressible External Filter Cake Model -- Linear
Geometry}\label{incompressible-external-filter-cake-model-linear-geometry}

This simple model assumes that all of the filtrate is captured on the
formation surface to form a filter cake of constant porosity and
permeability. Hence, the filter cake will simply grow in thickness as a
function of the volume of deposited material, and a simple volume
balance assuming linear flow yields,

\begin{equation}\phantomsection\label{eq-1-36}{
{A}_{f}{l}_{c}\left(t\right)\left(1-{ϕ}_{c}\right)={f}_{s}{\int }_{0}^{t}q(τ)dτ={f}_{s}{Q}_{c}\left(t\right)
}\end{equation}

Here \({A}_{f}\) is the area of flow, \({l}_{c}\) is the filter cake
thickness, \({ϕ}_{c}\) is the filter cake porosity, \({f}_{s}\) is the
volume concentration of suspended material, \(q\) is the injection rate,
\({Q}_{c}\) is the cumulative injected volume, and \(t\) is time. The
filter cake thickness is hence,

\begin{longtable}[]{@{}
  >{\raggedright\arraybackslash}p{(\linewidth - 2\tabcolsep) * \real{0.5000}}
  >{\raggedright\arraybackslash}p{(\linewidth - 2\tabcolsep) * \real{0.5000}}@{}}
\toprule\noalign{}
\begin{minipage}[b]{\linewidth}\raggedright
\({l}_{c}\left(t\right)=\frac{{f}_{s}{Q}_{c}\left(t\right)}{{A}_{f}\left(1-{ϕ}_{c}\right)}\).
\end{minipage} & \begin{minipage}[b]{\linewidth}\raggedright
(1.36)
\end{minipage} \\
\midrule\noalign{}
\endhead
\bottomrule\noalign{}
\endlastfoot
\end{longtable}

Assuming a constant filter cake permeability \({k}_{c}\), and viscosity
\(μ\), Darcy's law for the pressure drop across the filtrate gives the
damage pressure drop \(Δ{P}_{d}\) as,

\begin{equation}\phantomsection\label{eq-1-36}{
\mathrm{\mathrm{Δ}}{P}_{d}\left(t\right)=\frac{qμ{f}_{s}{Q}_{c}\left(t\right)}{{{k}_{c}A}_{f}^{2}\left(1-{ϕ}_{c}\right)}
}\end{equation}

and the skin factor as,

\begin{equation}\phantomsection\label{eq-1-36}{
S=\frac{{f}_{s}}{1-{ϕ}_{c}}\frac{k}{{k}_{c}}\frac{{Q}_{c}\left(t\right)}{{A}_{f}{r}_{w}}=\frac{{l}_{c}}{{r}_{w}}\frac{k}{{k}_{c}}
}\end{equation}

Here \(k\) is the rock permeability. The linear flow assumption is valid
for flow near a fracture surface, and it is also a decent approximation
for radial flow assuming that the filter cake is thin compared to the
well radius.

\subsection{Incompressible External Filter Cake Model -- Radial
Geometry}\label{incompressible-external-filter-cake-model-radial-geometry}

\subsubsection{Deposition Inside the
Wellbore}\label{deposition-inside-the-wellbore}

Simple volume balance on deposited filtrate inside a well in radial
geometry yields, for a 2D cross section perpendicular to the well,

\begin{equation}\phantomsection\label{eq-1-36}{
π{r}_{w}^{2}-π{r}_{c}^{2}=\frac{{f}_{s}}{1-{ϕ}_{c}}{\int }_{0}^{t}{q}_{h}\left(t\right)dt=\frac{{f}_{s}}{1-{ϕ}_{c}}{Q}_{h}(t)
}\end{equation}

Here \({q}_{h}\) and \({Q}_{h}\) are rate and cumulative injected volume
per length, respectively. This may be solved for the filtrate inner
radius \({r}_{c}\) to compute the cake thickness \({r}_{w}-{r}_{c}\) as,

\begin{equation}\phantomsection\label{eq-1-36}{
\mathrm{\mathrm{Δ}}{r}_{c}(t)={r}_{w}-{r}_{c}={r}_{w}-\sqrt{{r}_{w}^{2}-\frac{{f}_{s}{Q}_{h}\left(t\right)}{π(1-{ϕ}_{c})}}
}\end{equation}

The pressure drop across the filter cake is then given as,

\begin{equation}\phantomsection\label{eq-1-36}{
\mathrm{\mathrm{Δ}}{P}_{c}=\frac{qμ}{2π{k}_{c}}\mathrm{ln}⁡\frac{{r}_{c}}{{r}_{w}}
}\end{equation}

The skin is simply,

\begin{equation}\phantomsection\label{eq-1-36}{
S=\frac{k}{{k}_{c}}\mathrm{ln}⁡\frac{{r}_{w}}{{r}_{c}}
}\end{equation}

\subsubsection{Deposition in the
Reservoir}\label{deposition-in-the-reservoir}

If the filtrate is instead deposited in the near-well reservoir (i.e.,
outside the well), we will have \({r}_{c}>{r}_{w}\), and the skin factor
will be,

\begin{equation}\phantomsection\label{eq-1-36}{
S=\frac{k}{{k}_{c}}\mathrm{ln}⁡\frac{{r}_{c}}{{r}_{w}}
}\end{equation}

For implementation convenience and consistency with the linear geometry
case the filter cake thickness may be expressed in terms of a connection
area, \({A}_{f}=2π{r}_{w}h\), so that with \(Q\) being the cumulative
injected volume, the filter cake radius follows from,

\begin{equation}\phantomsection\label{eq-1-36}{
{r}_{c}^{2}={r}_{w}^{2}+\frac{2{r}_{w}{f}_{s}Q\left(t\right)}{{A}_{f}(1-{ϕ}_{c})}
}\end{equation}

\bookmarksetup{startatroot}

\chapter{INSTALLING AND RUNNING FLOW}\label{installing-and-running-flow}

\section{Installing OPM Flow}\label{installing-opm-flow}

OPM Flow can be installed and used in a variety of ways. Perhaps the
simplest way is to use the operating system's package manager to install
the OPM modules as binary packages, including OPM Flow and other OPM
programs. That way all prerequisite libraries will automatically be
installed on the system. This is supported on Ubuntu Linux version 22.04
or 24.04 and Red Hat Enterprise (or CentOS) version 7 or 8.

OPM Flow can also be used via a Docker container. This allows OPM Flow
to run on any system that supports Docker, including Microsoft Windows,
without installing prerequisite libraries or affecting any other
software on your system.

Finally, OPM Flow can be installed by compiling from source on Linux or
macOS systems.

\subsection{Ubuntu Linux 22.04 or 24.04 (64-bit version
only)}\label{ubuntu-linux-22.04-or-24.04-64-bit-version-only}

The easiest way to install the Ubuntu packages is to first add the OPM
personal package archive (ppa). In order to do that we need to install
the apt-add-repository command.

\begin{verbatim}
sudo apt-get update

sudo apt-get install software-properties-common
\end{verbatim}

Then we add the repository, and run update again:

\begin{verbatim}
sudo apt-add-repository ppa:opm/ppa

sudo apt-get update
\end{verbatim}

At this point, all the OPM modules should be available to install. To
see a list of (for example) the opm-simulators packages:

\begin{verbatim}
apt-cache search opm-simulators
\end{verbatim}

Then to install the opm-simulators programs (including OPM Flow) and
their dependencies:

\begin{verbatim}
sudo apt-get install mpi-default-bin

sudo apt-get install libopm-simulators-bin python3-opm-common
\end{verbatim}

To optionally install the OPM Python bindings use:

\begin{verbatim}
sudo apt-get install python3-opm-simulators
\end{verbatim}

Notes

\begin{verbatim}
- The mpi install above is required for the mpi libraries which the OPM executables are linked to, even those that are not intended to be run with mpirun. The reason the OPM packages do not depend on that package is that by Debian policy the user should be allowed to choose which MPI implementation to use (above we picked the default set by Ubuntu).
- If an old versions of prerequisite libraries have already been installed (for example from installing a previous release of OPM Flow) it might be necessary to upgrade them via the following commands:
\end{verbatim}

\begin{verbatim}
sudo apt-get update

sudo apt-get upgrade
\end{verbatim}

\begin{verbatim}
- Some users have experienced trouble with OPM Flow simply aborting with no error message. In some cases that can be related to locale issues. Check by running the “locale” command. OPM Flow requires the “C” or an English locale to run correctly. One way to fix this can be to put the following in your “.bash_profile” (note the initial period or full stop in the filename), and open a new terminal (command) window:
\end{verbatim}

\begin{verbatim}

        LANG="en_US.UTF-8"

        export LANG

        LC_ALL="en_US.UTF-8"

            export LC_ALL
\end{verbatim}

\subsection{Red Hat Enterprise or CentOS (version 7 or
8)}\label{red-hat-enterprise-or-centos-version-7-or-8}

First add the OPM package repository:

\begin{verbatim}
sudo yum-config-manager --add-repo \
http://www.opm-project.org/package/opm.repo
\end{verbatim}

The OPM software is split in several packages. To list all available OPM
packages one can use:

\begin{verbatim}
sudo yum search opm-
\end{verbatim}

For example you can install the opm-simulators binary package to get
access to the OPM Flow reservoir simulator:

\begin{verbatim}
sudo yum install opm-simulators-bin python3-opm-common
\end{verbatim}

Development packages are available using the -devel suffix. For example,
to install the opm-upscaling development package use:

\begin{verbatim}
sudo yum install opm-upscaling-devel
\end{verbatim}

Notes

\begin{itemize}
\tightlist
\item
  If you wish to install the MPI parallel version of Flow, you must
  choose an MPI implementation, either mpich or openmpi, and instead use
  the following command:
\end{itemize}

\begin{verbatim}
sudo yum install opm-simulators-openmpi-bin
\end{verbatim}

\begin{itemize}
\tightlist
\item
  For mpich, just replace openmpi with mpich in the above command. In
  order to run Flow in parallel using mpirun you must also load the
  required module first:
\end{itemize}

\begin{verbatim}

module load mpi/openmpi-x86_64
\end{verbatim}

To use mpich instead, replace openmpi with mpich in the above command.

\begin{verbatim}
\end{verbatim}

\subsection{Windows 7 and 10 Using
VirtualBox}\label{windows-7-and-10-using-virtualbox}

Oracle VM VirtualBox's (formerly Sun VirtualBox, Sun xVM VirtualBox and
Innotek VirtualBox) is a free and open source hosted hypervisor for x86
virtualization, developed by Oracle Corporation. The software was
originally created by Innotek, who were acquired by Sun Microsystems in
2008, which was in turn acquired by Oracle in 2010. VirtualBox can be
installed on Windows, macOS, Linux, Solaris and Open Solaris operating
systems. The software supports the creation and management of guest
virtual machines running Windows, Linux, BSD, OS/2, Solaris, Haiku, and
OSx86, as well as limited virtualization of macOS guests on Apple
hardware For some guest operating systems, a ``Guest Additions'' package
of device drivers and system applications is available, which typically
improves performance, especially that of graphics {[}After
\url{https://en.wikipedia.org/wiki/VirtualBox}.{]}.

The OPM web site outlines the install process using VirtualBox, Vagrant
{[}See \url{https://www.vagrantup.com/} for details{]} and a Vagrant
file to automatically set up a virtual environment -- see
\url{https://opm-project.org/?page_id=294&page=2} for details on this
approach.

Alternatively, one can use the more standard approach of installing
virtual machines using VirtualBox. The basic outline for this method to
use VirtualBox to run a guest Linux distribution and install OPM Flow in
the guest Linux virtual machine. Note only the Linux distributions
mentioned previously should be used as a guest virtual machine, although
any of the Ubuntu ``flavors'' (Kubuntu, Lubuntu, Ubuntu Mate, etc.) will
work as well. The steps are:

\begin{verbatim}
- Installing VirtualBox [See also [https://www.virtualbox.org/wiki/Documentation](https://www.virtualbox.org/wiki/Documentation) for further and more detailed information.]
\end{verbatim}

\begin{itemize}
\item
  First download the latest VirtualBox software from Oracle's web site
  at \url{https://www.virtualbox.org/} for Windows host machines.
\item
  Install the software via double clicking on the executable and
  following the on screen instructions.

  \begin{itemize}
  \tightlist
  \item
    Create a Guest Linux Virtual Machine.
  \end{itemize}
\item
  Download a suitable Linux distribution in ISO format that supports
  running OPM Flow, for example Ubuntu 24.04 LTS (64-bit version only)
  from \url{https://ubuntu.com/download/desktop} or CentOS Linux
  (version 7 or 8) from \url{https://www.centos.org/download/}.
\item
  Start VirtualBox, which should display a screen similar to Figure 2.1
  without all the guest virtual machines.
\end{itemize}

\pandocbounded{\includegraphics[keepaspectratio]{chapters/images/Frame54_1264b6329c73.png}}
\pandocbounded{\includegraphics[keepaspectratio]{chapters/images/Image16_1264b6329c73.png}}

\begin{itemize}
\tightlist
\item
  Then click the ``New'' icon button to install a new guest virtual
  machine.
\item
  Identify the operating system type (Linux) and the Linux distribution
  that one wishes to install (Ubuntu 64-bit or RedHat for one of the
  CentOS distributions).
\item
  Set the amount of RAM, should be 4GB or more; note the guest machine
  properties can subsequently be changed after the guest machine has
  been created.
\item
  Create a virtual hard drive, the default virtual disk of VDI is
  normally used, and select either the dynamic or the fixed option. Next
  select the size of the virtual disk, the size should be sufficient to
  allow for the operating system and any other software, including OPM
  Flow, to be installed. For Ubuntu 24.04 LTS the operating system needs
  at least 10 GB. VirtualBox will then create the guest virtual machine.
\item
  The next step is to start the guest virtual machine and installing the
  Linux operating system. This can be done by double clicking on the
  guest virtual machine and then selecting the previously downloaded IOS
  file. For Ubuntu 20.04 LTS the machine will boot into a live
  distribution, that is the guest virtual machine is running from a
  ``live CD'', shown below.
\end{itemize}

\pandocbounded{\includegraphics[keepaspectratio]{chapters/images/Frame55_c3788dc4a8a0.png}}
\pandocbounded{\includegraphics[keepaspectratio]{chapters/images/Image44_c3788dc4a8a0.png}}

\begin{itemize}
\item
  Select the ``Install Ubuntu'' option to install the operating system.
  The ``Try Ubuntu'' option means the guest machine will run from a
  ``Live CD''.
\item
  Follow the on screen instructions to install the operating system,
  together with any updates suggested by the installer. Once the
  installation is complete the guest machine will re-boot as a normal
  guest Linux machine.

  \begin{itemize}
  \tightlist
  \item
    Installing OPM Flow in the Guest Virtual Machine.
  \end{itemize}
\item
  Boot up your virtual machine by selecting the machine and selecting
  the ``Start'' icon on the main VirtualBox screen (Figure 2.1). One
  should then be presented with the newly created guest virtual machine,
  similar to one shown in Figure 2.3, that displays the standard Ubuntu
  desktop.
\item
  In order to install OPM Flow in the guest virtual machine one first
  needs to open a command terminal using ALT+CNTL+T key sequence (Figure
  2.4).
\end{itemize}

\pandocbounded{\includegraphics[keepaspectratio]{chapters/images/Frame57_a643df6e083e.png}}
\pandocbounded{\includegraphics[keepaspectratio]{chapters/images/Image51_a643df6e083e.png}}

\begin{itemize}
\tightlist
\item
  One can then use the same commands as for a host machine, as described
  earlier in Sections 2.1.1 Ubuntu Linux 22.04 or 24.04 (64-bit version
  only) and 2.1.2 Red Hat Enterprise or CentOS (version 7 or 8) to
  install the software. To check if OPM Flow is installed, execute the
  following command in the terminal window:
\end{itemize}

\begin{verbatim}
flow –-help
\end{verbatim}

which should then display the command line parameters for the installed
version of OPM Flow, as depicted in Figure 2.5.

\begin{verbatim}
- Using OPM Flow in a Linux Virtual Machine.
\end{verbatim}

\begin{itemize}
\tightlist
\item
  The guest virtual machine only has access to files within the guest
  machine. To enable the guest to have access to the host machine files;
  on the main VirtualBox screen (Figure 2.1), select the ``Setting''
  icon and then the ``Shared Folders'' option to set a share directory
  between the guest and the host machines. One can then have all the
  input and output files on the host machine and run OPM Flow in the
  guest virtual machine.
\item
  Note one should shut down the guest virtual machine in the normal way
  to avoid potential file corruptions.
\item
  One may also take snapshots of the current state of the virtual
  machine, as a form of backup if one changes the configuration of the
  virtual machine, or if one updates to the latest release of OPM Flow.
\end{itemize}

Historically this has been one of the preferred methods to run OPM Flow
on Windows machines and still has several advantages, including being
able to export and import various guest virtual machines for testing
various versions of the simulator. However, Microsoft, with the
introduction of Windows Subsystem for Linux (``WSL'') has create a more
seamless method for running OPM Flow on Windows 10 and higher machines.
This option is outlined in the next section.

\subsection{Windows 10 and 11 -- Windows Subsystem for Linux
(WSL)}\label{windows-10-and-11-windows-subsystem-for-linux-wsl}

The first release of WSL was on August 2nd, 2016 and there has been
various updates to this version which is known as WSL 1. On May 6th,
2019 WSL 2 was announced, that was a complete re-write of the Windows
Subsystem for Linux, and offers significant performance improvements
over WSL 1. Note that WSL 2 requires Windows 10 version 1903 or higher,
with Build 18362 or higher, for x64 systems. WSL 2 is based on a virtual
machine with an actual Linux Kernel that can immediately react to system
calls. Such an architecture enables Full System call capabilities.

To run OPM Flow under WSL 1 or WSL 2 one has to first activate WSL, then
install a Linux distribution from the Application store, and then
finally install OPM Flow. The detail instruction are as follows:

\begin{verbatim}
- Activating Windows Subsystem for Linux.
\end{verbatim}

\begin{itemize}
\tightlist
\item
  On the Windows machine, Click Start and navigate to Control Panel,
  then Select the Programs Category. Under Programs and Features, click
  on Turn Windows Features on or off and click the check box next to the
  Windows Subsystem for Linux option, to enable WSL.
\end{itemize}

\pandocbounded{\includegraphics[keepaspectratio]{chapters/images/Frame58_584c66751192.png}}
\pandocbounded{\includegraphics[keepaspectratio]{chapters/images/Image52_584c66751192.png}}

You may also have to ensure that CPU virtualization. has been activated
in your computer's BIOS.

\begin{verbatim}
- Installing a Linux Distribution.
\end{verbatim}

\begin{itemize}
\tightlist
\item
  To install the Linux distribution in Windows, click Start and then
  select the Microsoft Store.
\item
  Using the search box available in the Microsoft Store, pick one of the
  distributions that is compatible with OPM Flow (Figure 2.7). Note that
  for Ubuntu these are usually the two latest LTS versions before the
  release, i.e.~version ``20XX.04 LTS'' where XX is a multiple of 4.
\end{itemize}

\pandocbounded{\includegraphics[keepaspectratio]{chapters/images/Frame59_c73b5cb709ae.png}}
\pandocbounded{\includegraphics[keepaspectratio]{chapters/images/Image53_c73b5cb709ae.png}}

\begin{itemize}
\tightlist
\item
  After selecting the distribution, click Install or Get to install the
  Linux Distribution. Note the installer will ask you to sign into your
  account, but you can close the pop-up window if you don't want to sign
  in at this time.
\item
  On completion of the installation process the Linux distribution will
  be installed on Windows and ready for use. Furthermore, the
  distribution will be available directly in the Start Menu itself in
  the form of a standard everyday application (Figure 2.8).
\end{itemize}

\pandocbounded{\includegraphics[keepaspectratio]{chapters/images/Frame60_5e4b9d976c6a.png}}
\pandocbounded{\includegraphics[keepaspectratio]{chapters/images/Image54_5e4b9d976c6a.png}}

\begin{itemize}
\tightlist
\item
  When starting the Linux Distribution for the first time, a brief setup
  process will take place. Basically, it will ask for the username and
  password you want to set for the Linux environment, a similar process
  as when you install a Linux distribution via VirtualBox.
\item
  To check the version of WSL that will be used to run the Linux
  distribution, in PowerShell use the following command:
\end{itemize}

\begin{verbatim}
wsl -l -v
\end{verbatim}

\begin{itemize}
\tightlist
\item
  To set the WSL for a given distribution, for example Ubuntu-24-04 to
  WSL 2, in PowerShell use the following command:
\end{itemize}

\begin{verbatim}
wsl --set-version Ubuntu-24-04 2
\end{verbatim}

\begin{itemize}
\tightlist
\item
  And finally to set the default version for WSL for all Linux
  distributions, use the following PowerShell command:
\end{itemize}

\begin{verbatim}
wsl --set-default-version 2
\end{verbatim}

\begin{verbatim}
- Installing OPM Flow in a WSL Linux Distribution:
\end{verbatim}

\begin{itemize}
\tightlist
\item
  Click on installed Linux distribution in the Start Menu to open a
  Linux terminal session, and then execute the Linux commands as
  described earlier in Sections 2.1.1 Ubuntu Linux 22.04 or 24.04
  (64-bit version only) and 2.1.2 Red Hat Enterprise or CentOS (version
  7 or 8) to install the software.
\item
  To check if OPM Flow is installed, execute the following command in
  the Linux terminal window:
\end{itemize}

\begin{verbatim}
flow –-help
\end{verbatim}

which should then display the command line parameters for the installed
version of OPM Flow, as depicted in Figure 2.9.

\pandocbounded{\includegraphics[keepaspectratio]{chapters/images/Frame62_d816261771f0.png}}
\pandocbounded{\includegraphics[keepaspectratio]{chapters/images/Image56_d816261771f0.png}}

One can also use the Windows command terminal, instead of the WSL Linux
Terminal, to do the same thing by adding WSL character string before the
command, that is:

\begin{verbatim}
wsl flow –-help

\end{verbatim}

\begin{verbatim}
- Using OPM Flow via the Windows  WSL Linux Virtual Machine.
\end{verbatim}

\begin{itemize}
\tightlist
\item
  Similar to running OPM Flow using a Linux distribution installed in
  VirtualBox, the WSL Linux virtual machine only has access to files
  within WSL virtual machine. To access the Windows file system, one has
  to execute the Linux CD command in the WSL Linux Terminal, the command
  has the form:
\end{itemize}

\begin{verbatim}
cd /mnt/<drive letter>
\end{verbatim}

For example, to access the Windows system drive C one would use the
following command:

\begin{verbatim}
cd /mnt/c
\end{verbatim}

One can then run OPM Flow from the desired directory in the WSL Linux
virtual machine and use OPM ResInsight under Windows to view the results
of the simulation run.

User experience indicates this approach of running OPM Flow under
Windows is the most integrated and convenient method, compared to the
VirtualBox approach, especially for the less technically advance users.

\subsection{Using a Docker Container}\label{using-a-docker-container}

See the tutorial on the OPM website
\href{http://opm-project.org/?page_id=853&page=2}{Running Flow in
Docker}.

\subsection{Installing from Source}\label{installing-from-source}

See instructions on the OPM website
\href{http://opm-project.org/?page_id=231}{Building from Source}.

\section{Running OPM Flow From The Command
Line}\label{running-opm-flow-from-the-command-line}

This section describes how to run the OPM Flow simulator and the various
command line options associated with this release, the command line
syntax for running OPM Flow is:

\begin{verbatim}

flow [OPTIONS] CASENAME
\end{verbatim}

and typing the following command from your terminal:

\begin{verbatim}

flow CASENAME
\end{verbatim}

will start the simulator and run the case specified by CASENAME.DATA.
Alternatively, you can type the path to the binary and the *.DATA file,
as per the example below:

\begin{verbatim}

path_to_binary/flow path_to_data/CASENAME
\end{verbatim}

The simulator can also be run using command line parameters for example
to run the previous CASENAME one could use:

\begin{verbatim}

flow --ecl-deck-file-name=path_to_data/CASENAME
\end{verbatim}

Note that there must be no spaces around the equals sign and there
should be no spaces in CASENAME. If there are spaces in CASENAME, which
can typically occur on Windows based systems, then the CASENAME should
be enclosed in quotes; thus instead of CASE NAME use ``CASE NAME''.

By default output files are generated in the same folder as the *.DATA
file. To override this one can use the --output-dir command line
parameter, for instance:

\begin{verbatim}

flow --output-dir=foo CASENAME
\end{verbatim}

will send the output files to the foo directory.

It is also possible to put multiple command line options together in a
parameter file. To do so, put one option on each line of the file with
``.param'' as the extension and pass that filename as a command line
parameter to OPM Flow using:

\begin{verbatim}

flow --parameter-file=CASENAME.param
\end{verbatim}

to instruct OPM Flow to read the parameter file. Note, it is good
practice to name both the *.DATA and parameter file with the same
CASENAME, to avoid confusion. However, if one wishes to use the same
parameter file for various cases one can specify the parameter file
followed by the DATA file, for example:

\begin{verbatim}

flow --parameter-file=path_to_params/ALLCASES.param CASENAME.DATA
\end{verbatim}

If OPM Flow is installed with parallel capabilities then:

\begin{verbatim}

mpirun -np 4 flow --parameter-file=path_to_data/CASENAME.param
\end{verbatim}

will run OPM Flow on four nodes etc.

\begin{tcolorbox}[enhanced jigsaw, title=\textcolor{quarto-callout-note-color}{\faInfo}\hspace{0.5em}{Note}, coltitle=black, breakable, toprule=.15mm, colbacktitle=quarto-callout-note-color!10!white, leftrule=.75mm, opacitybacktitle=0.6, colback=white, bottomtitle=1mm, toptitle=1mm, titlerule=0mm, left=2mm, colframe=quarto-callout-note-color-frame, arc=.35mm, rightrule=.15mm, bottomrule=.15mm, opacityback=0]

On Red Hat based distributions, including centOS distributions, one
needs to install the openmpi version of the simulator,
opm-simulators-openmpi-bin, via: sudo yum install
opm-simulators-openmpi-bin and then set the MPI version to use. The way
this is done is by first making the module command available, by running
the following command: . /usr/share/Modules/init/bash Then query for the
installed modules with: module avail And lastly, you tell the system it
to use openmpi with: module add mpi/openmpi-x86\_64 On Ubuntu based
distributions, there is no need for this as openmpi is installed and
mpirun just works.

\end{tcolorbox}

A complete list of the active user facing command line options and their
function is given in the Table 2.1 for the current release. These
commands are the options listed with command line --help option, there
is also a --help-all option that lists all the command line options
included in the release, including experimental, obsolete, hidden and
deprecated options.

Additional tutorials for running OPM Flow are available on the OPM
website in the \href{http://opm-project.org/?page_id=43}{Tutorials}
section.

\begin{longtable}[]{@{}
  >{\raggedright\arraybackslash}p{(\linewidth - 6\tabcolsep) * \real{0.2500}}
  >{\raggedright\arraybackslash}p{(\linewidth - 6\tabcolsep) * \real{0.2500}}
  >{\raggedright\arraybackslash}p{(\linewidth - 6\tabcolsep) * \real{0.2500}}
  >{\raggedright\arraybackslash}p{(\linewidth - 6\tabcolsep) * \real{0.2500}}@{}}
\toprule\noalign{}
\begin{minipage}[b]{\linewidth}\raggedright
OPM Flow Command Line Options
\end{minipage} & \begin{minipage}[b]{\linewidth}\raggedright
\end{minipage} & \begin{minipage}[b]{\linewidth}\raggedright
\end{minipage} & \begin{minipage}[b]{\linewidth}\raggedright
\end{minipage} \\
\midrule\noalign{}
\endhead
\bottomrule\noalign{}
\endlastfoot
No. & Variable Name & Description & Default \\
1 & -h or --help & A character string. Print the help message that gives
a brief description of the available command line parameters. & N/A \\
2 & --help-all & Prints all the command line options included in the
release, including obsolete, hidden and deprecated options. & N/A \\
3 & --accelerator-mode & A character string. Choose a linear solver,
usage:
`--accelerator-mode={[}none\textbar cusparse\textbar opencl\textbar amgcl\textbar rocalution\textbar rocsparse{]}'.
& ``none'' \\
4 & --action-parsing-strictness & A character string. Set strictness of
parsing process for ActionX and PyAction. Available options are normal
(do not apply keywords that have not been tested for ActionX or
PyAction) and low (try to apply all keywords, beware: the simulation
outcome might be incorrect). & ``normal'' \\
5 & --add-corners & A boolean value. Add corners to partition. &
false \\
6 & --allow-distributed-wells & A boolean value. Allow the perforations
of a well to be distributed to interior of multiple processes. &
false \\
7 & --allow-splitting-inactive-wells & A boolean value. Allow inactive
(never non-shut) wells to be split across multiple domains. & true \\
8 & --alternative-well-rate-init & A boolean value. Use alternative well
rate initialization procedure. & true \\
9 & --check-group-constraints-inner-well-iterations & A boolean value.
Allow checking of group constraints during inner well iterations. &
true \\
10 & --check-satfunc-consistency & Saturation function consistency
checking is now enabled by default. A boolean value. Whether or not to
check saturation function consistency requirements. & true \\
11 & --continue-on-convergence-error & A boolean value. Continue with a
non-converged solution instead of giving up if we encounter a time step
size smaller than the minimum time step size. & false \\
12 & --convergence-monitoring & A boolean value. Enable convergence
monitoring. & false \\
13 & --convergence-monitoring-cut-off & An integer value. Cut off limit
for convergence monitoring. & 6 \\
14 & --convergence-monitoring-decay-factor & A positive real value.
Decay factor for convergence monitoring. & 0.75 \\
15 & --cpr-reuse-interval & An integer value. Reuse preconditioner
interval. Used when CprReuseSetup is set to 4, then the preconditioner
will be fully recreated instead of reused every N linear solve, where N
is this parameter. & 30 \\
16 & --cpr-reuse-setup & A positive integer. Specify if the CPR solver
should re-use the AMG setup. Valid options are: & 4 \\
17 & --cpr-weights-thread-parallel & A boolean value. Enable OpenMP
thread parallelization of CPR weight calculation. This can improve
performance for large models but is disabled by default. Usage:
--cpr-weights-thread-parallel={[}true\textbar false{]}. & false \\
18 & --dbhp-max-rel & A positive real value. Maximum relative change of
the bottom-hole pressure in a single iteration. & 1 \\
19 & --debug-emit-cell-partition & A boolean value. Whether or not to
emit cell partitions as a debugging aid. & false \\
20 & --debug-verbosity-level & An integer value. Set debug verbosity
level globally. Default is 1, increasing values give additional output
and 0 disables most messages to the .DBG file. & 1 \\
21 & --dp-max-rel & A positive real value. Maximum relative change of
pressure in a single iteration. & 0.3 \\
22 & --ds-max & A positive real value. Maximum absolute change of any
saturation in a single iteration. & 0.2 \\
23 & --dwell-fraction-max & A positive real value. Maximum absolute
change of a well's volume fraction in a single iteration. & 0.2 \\
24 & --ecl-deck-file-name & A character string. The name of the file
which contains the ECL deck to be simulated. & ``\,'' \\
25 & --ecl-output-double-precision & A boolean value. Tell the output
writer to use double precision. Useful for `perfect' restarts. &
false \\
26 & --ecl-output-interval & An integer value. The number of report
steps that ought to be skipped between two writes of ECL results. &
-1 \\
27 & --edge-conformal & A boolean value. Edge conformal cornerpoint
processing. & false \\
28 & --edge-weights-method & This option is now specified by a string
rather than an integer. A defined character string. Specify the
edge-weighting strategy: & ``transmissibility'' \\
29 & --enable-adaptive-time-stepping & A boolean value. Use adaptive
time stepping between report steps. & true \\
30 & --enable-async-ecl-output & A boolean value. Write the ECL-formated
results in a non-blocking way (i.e., using a separate thread). & true \\
31 & --enable-async-vtk-output & A boolean value. Dispatch a separate
thread to write the VTK output. & true \\
32 & --enable-drift-compensation & A boolean value. Enable partial
compensation of systematic mass losses via the source term of the next
time step. & false \\
33 & --enable-dry-run & A boolean value. Specify if the simulation
should actually run (true), or just check the input deck (false). This
option is equivalent to activating the NOSIM keyword in the RUNSPEC
section of the input deck. The default value of ``auto'' will use
whatever is stipulated in the input deck via the NOSIM keyword. &
``auto'' \\
34 & --enable-ecl-output & A boolean value. Write binary output which is
compatible with the commercial Eclipse simulator. & true \\
35 & --enable-esmry & A boolean value. Write ESMRY file for fast loading
of summary data. & false \\
36 & --enable-gravity & A boolean value. Use the gravity correction for
the pressure gradients. & true \\
37 & --enable-grid-adaptation & A boolean value. Enable adaptive grid
refinement/coarsening. & false \\
38 & --enable-intensive-quantity-cache & A boolean value. Turn on
caching of intensive quantities. & true \\
39 & --enable-logging-fallout-warning & A boolean value. Developer
option to see whether logging was on non-root processors. In that case
it will be appended to the \emph{.DBG or }.PRT files. & false \\
40 & --enable-opm-rst-file & A boolean value. Include OPM-specific
keywords in the ECL restart file to enable restart of OPM simulators
from these files. & false \\
41 & --enable-storage-cache & A boolean value. Store previous storage
terms and avoid re-calculating them. & true \\
42 & --enable-terminal-output & A boolean value. Print high-level
information about the simulation's progress to the terminal. & true \\
43 & --enable-thermodynamic-hints & A boolean value. Enable
thermodynamic hints. & false \\
44 & --enable-tuning & A boolean value. Read the time stepping and
convergence parameters from the TUNING keyword in the input deck, if set
to true. Note that only some of the TUNING keyword parameters are
processed. & false \\
45 & --enable-vtk-output & A boolean value. Global switch for turning on
writing VTK files. & false \\
46 & --enable-well-operability-check & A boolean value. Enable the well
operability checking. & true \\
47 & --enable-well-operability-check-iter & A boolean value. Enable the
well operability checking during iterations. & false \\
48 & --enable-write-all-solutions & A boolean value. Write all solutions
to disk instead of only the ones for the report steps. & false \\
49 & --end-time & A positive real value. The simulation time at which
the simulation is finished {[}s{]}. & 1e+100 \\
50 & --explicit-rock-compaction & A boolean value. Use pressure from end
of the last time step when evaluating rock compaction. & false \\
51 & --external-partition & A character string. Name of file from which
to load an externally generated partitioning of the model's active cells
for MPI distribution purposes. If empty, the built-in partitioning
method will be employed. & ``\,'' \\
52 & --force-disable-fluid-in-place-output & A boolean value. Do not
print fluid-in-place values after each report step even if requested by
the deck. Note this parameter will override the print request in the
input deck. & false \\
53 & --force-disable-resv-fluid-in-place-output & A boolean value. Do
not print reservoir volumes values after each report step even if
requested by the deck. Note this parameter will override the print
request in the input deck. & false \\
54 & --full-time-step-initially & A boolean value. Always attempt to
finish a report step using a single substep. & false \\
55 & --gpu-aware-mpi & A boolean value. MPI communication use GPU aware
MPI in the sense that it will use GPU direct communication. Setting this
to true will require that the MPI implementation supports GPU direct
communication. If you are unsure, set this to false. Usage:
--gpu-aware-mpi={[}true\textbar false{]}. . & false \\
56 & --gpu-device-id & An integer value. Choose device ID for
cusparseSolver or openclSolver, use `nvidia-smi' or `clinfo' to
determine valid IDs. & 0 \\
57 & --hy-ne-config-file & A character string. Use config files for
Hybrid Newton. & ``hybridNewtonConfig.json'' \\
58 & --ignore-keywords & A character string. List of Eclipse keywords
which should be ignored. As a `:' separated string. & ``\,'' \\
59 & --ilu-fillin-level & An integer value. The fill-in level of the
linear solver's ILU preconditioner. & 0 \\
60 & --ilu-redblack & A boolean value. Use red-black partitioning for
the ILU preconditioner. & false \\
61 & --ilu-relaxation & A real value ranging from zero to one. The
relaxation factor of the linear solver's ILU preconditioner. & 0.9 \\
62 & --ilu-reorder-spheres & A boolean value. Whether to reorder the
entries of the matrix in the red-black ILU preconditioner in spheres
starting at an edge. If false the original ordering is preserved in each
color. Otherwise why try to ensure D4 ordering (in a 2D structured grid,
the diagonal elements are consecutive). & false \\
63 & --imbalance-tol & A positive real value. Tolerable imbalance of the
loadbalancing. & 1.1 \\
64 & --initial-time-step-in-days & A positive real value. The size of
the initial time step in days. & 1 \\
65 & --initial-time-step-size & A positive real value. The size of the
initial time step {[}s{]}. & 86400 \\
66 & --inj-mult-damp-mult & A real value ranging from zero to one.
Injection multiplier dampening factor (dampening multiplied by this each
time oscillation is detected). & 0.9 \\
67 & --inj-mult-min-damp-factor & A real value ranging from zero to one.
Minimum injection multiplier dampening factor (maximum dampening level).
& 0.05 \\
68 & --inj-mult-osc-threshold & A positive real value. Injection
multiplier oscillation threshold (used for multiplier dampening). &
0.1 \\
69 & --input-skip-mode & A defined character string. Set compatibility
mode for the \hyperref[REF_HEADING_KEYWORD_SKIP100_11_3]{SKIP100} and
\hyperref[REF_HEADING_KEYWORD_SKIP300_11_3]{SKIP300} keywords. Options
are: & ``100'' \\
70 & --linear-solver & A defined character string. Specify the
configuration of the linear solver, valid values are: Option (1) extends
the existing Constrained Pressure Residual (``CPR'') preconditioner to
include wells. This option can also be invoked via the CPR keyword in
the RUNSPEC section; however, the command line parameter takes
precedence. For option (10) one enters a character string enclosed in
quotes that defines the filename of a JSON configuration file for the
flexible linear solver system. In this case the file extension should be
`.json'. Note that the Print (PRT) file contains the ``Property tree for
the linear solver'' listing, which is the JSON specification of the
current case, and can be used to configure a user specific linear solver
JSON file. & ``cprw'' \\
71 & --linear-solver-accelerator & A character string. Choose the
backend for the linear solver, usage:
`--linear-solver-accelerator={[}cpu\textbar gpu{]}'. & ``cpu'' \\
72 & --linear-solver-ignore-convergence-failure & A boolean value.
Continue with the simulation like nothing happened after the linear
solver did not converge. The default value for ``CPR'' solvers is 20. &
false \\
73 & --linear-solver-max-iter & An integer value. The maximum number of
iterations of the linear solver. & 200 \\
74 & --linear-solver-print-json-definition & A boolean value. Write the
JSON definition of the linear solver setup to the DBG file. The default
value for ``CPR'' solvers is 0.005. & true \\
75 & --linear-solver-reduction & A positive real value. The minimum
reduction of the residual which the linear solver must achieve. &
0.01 \\
76 & --linear-solver-restart & An integer value. The number of
iterations after which GMRES is restarted. & 40 \\
77 & --linear-solver-verbosity & An integer value. Specify the output
from the linear solver: & 0 \\
78 & --load-file & A character string. FileName for .OPMRST file used to
load serialized state. If empty, CASENAME.OPMRST is used. & ``\,'' \\
79 & --load-step & An integer value. Specify if the simulator should
load the serialized state from OPM Flow's version of the restart file,
and should be set to either a specific report step, or 0 to load the
last stored report step. The default value of -1 does not load the data
from the OPM Flow specific restart file. OPM Flow's version of the
restart file, is written using the --save-step=N option. & -1 \\
80 & --local-domains-ordering-measure & A character string. Subdomain
ordering measure. Allowed values are `maxpressure', `averagepressure'
and `residual'. & ``maxpressure'' \\
81 & --local-domains-partition-well-neighbor-levels & An integer value.
Number of neighbor levels around wells to include in the same domain
during NLDD partitioning. & 1 \\
82 & --local-domains-partitioning-imbalance & A real value greater than
one. Subdomain partitioning imbalance tolerance. 1.03 is 3 percent
imbalance. & 1.03 \\
83 & --local-domains-partitioning-method & A character string. Subdomain
partitioning method. Allowed values are `zoltan', `simple', and the name
of a partition file ending with `.partition'. & ``zoltan'' \\
84 & --local-solve-approach & A character string. Choose local solve
approach. Valid choices are jacobi and gauss-seidel. &
``gauss-seidel'' \\
85 & --local-tolerance-scaling-cnv & A positive real value. Set lower
than 1.0 to use stricter convergence tolerance for local solves. &
0.1 \\
86 & --local-tolerance-scaling-mb & A positive real value. Set lower
than 1.0 to use stricter convergence tolerance for local solves. & 1 \\
87 & --local-well-solve-control-switching & A boolean value. Allow
control switching during local well solutions. & true \\
88 & --matrix-add-well-contributions & A boolean value. Explicitly
specify the influences of wells between cells in the Jacobian and
preconditioner matrices. & false \\
89 & --max-inner-iter-ms-wells & An integer value. Maximum number of
inner iterations for multi-segment wells. & 100 \\
90 & --max-inner-iter-wells & An integer value. Maximum number of inner
iterations for standard wells. & 50 \\
91 & --max-local-solve-iterations & An integer value. Max iterations for
local solves with NLDD nonlinear solver. & 20 \\
92 & --max-newton-iterations-with-inner-well-iterations & An integer
value. Maximum newton iterations with inner well iterations. & 8 \\
93 & --max-pressure-change-ms-wells & A positive real value. Maximum
relative pressure change for a single iteration of the multi-segment
well model. & 1e+06 \\
94 & --max-residual-allowed & A positive real value. Absolute maximum
tolerated for residuals without cutting the time step size. & 1e+07 \\
95 & --max-single-precision-days & A positive real value. Maximum time
step size where single precision floating point arithmetic can be used
solving for the linear systems of equations. & 20 \\
96 & --max-temperature-change & A positive real value. Maximum absolute
change of temperature in a single iteration. & 5 \\
97 & --max-time-step-divisions & An integer value. The maximum number of
divisions by two of the timestep size before the simulation bails out. &
10 \\
98 & --max-time-step-size & A positive real value. The maximum size to
which all time steps are limited to {[}s{]}. & inf \\
99 & --max-well-status-switch-for-wells & An integer value. Maximum
number of status switching (stop\textless-\textgreater open) for a well
during a time-step. & 99 \\
100 & --max-well-status-switch-in-inner-iter-wells & An integer value.
Maximum number of status switching (stop\textless-\textgreater open) for
a well during inner iterations. & 99 \\
101 & --max-welleq-iter & An integer value. Maximum number of iterations
to determine solution the well equations. & 30 \\
102 & --maximum-number-of-group-switches & An integer value. Maximum
number of times a group can switch to the same control. & 3 \\
103 & --maximum-number-of-well-switches & An integer value. Maximum
number of times a well can switch to the same control. & 3 \\
104 & --maximum-water-saturation & A positive real value. Maximum water
saturation. & 1 \\
105 & --metis-params & A character string. Configuration of Metis
partitioner. You can request a configuration to be read from a JSON file
by giving the filename here, with extension `.json'. See
http://glaros.dtc.umn.edu/gkhome/fetch/sw/metis/manual.pdf for available
METIS options. & ``default'' \\
106 & --milu-variant & Option strings have been converted to lower case.
A defined character string. Specify which variant of the modified ILU
preconditioner ought to be used. Possible variants are: & ``ilu'' \\
107 & --min-strict-cnv-iter & An integer value. Minimum number of Newton
iterations before relaxed tolerances can be used for the CNV convergence
criterion. & -1 \\
108 & --min-strict-mb-iter & An integer value. Minimum number of Newton
iterations before relaxed tolerances can be used for the MB convergence
criterion. Default -1 means that the relaxed tolerance is used when
maximum number of Newton iterations are reached. & -1 \\
109 & --min-time-step-based-on-newton-iterations & A real value greater
than or equal to zero. The minimum time step size (in days for field and
metric unit and hours for lab unit) can be reduced to based on newton
iteration counts. & 0 \\
110 & --min-time-step-before-shutting-problematic-wells-in-days & A real
value greater than or equal to zero. The minimum time step size in days
for which problematic wells are not shut. & 0.01 \\
111 & --min-time-step-size & A real value greater than or equal to zero.
The minimum size to which all time steps are limited to {[}s{]}. & 0 \\
112 & --network-max-outer-iterations & An integer value. Maximum outer
number of iterations in the network solver before giving up. & 10 \\
113 & --network-max-pressure-update-in-bars & A positive real value.
Maximum pressure update in the inner network pressure update iterations.
& 5 \\
114 & --network-max-strict-outer-iterations & An integer value. Maximum
outer iterations in network solver before relaxing tolerance. & 10 \\
115 & --network-max-sub-iterations & An integer value. Maximum number of
sub-iterations to update network pressures (within a single well/group
control update). & 20 \\
116 & --network-pressure-update-damping-factor & A real value ranging
from zero to one. Damping factor in the inner network pressure update
iterations. & 0.1 \\
117 & --newton-max-error & A positive real value. The maximum error
tolerated by the Newton method to which does not cause an abort. &
1e+100 \\
118 & --newton-max-iterations & An integer value. The maximum number of
Newton iterations per time step. & 20 \\
119 & --newton-max-relax & A positive real value. The maximum relaxation
factor of a Newton iteration. & 0.5 \\
120 & --newton-min-iterations & An integer value. The minimum number of
Newton iterations per time step. & 2 \\
121 & --newton-relaxation-type & A character string. The type of
relaxation used by the Newton method. Valid options are: dampen or sor.
& ``dampen'' \\
122 & --newton-target-iterations & An integer value. The `optimum'
number of Newton iterations per time step. & 10 \\
123 & --newton-tolerance & A positive real value. The maximum raw error
tolerated by the Newton method for considering a solution to be
converged. & 0.01 \\
124 & --newton-verbose & A boolean value. Specify whether the Newton
method should inform the user about its progress or not. & true \\
125 & --newton-write-convergence & A boolean value. Write the
convergence behaviour of the Newton method to a VTK file. & false \\
126 & --nldd-local-linear-solver & A character string. Configuration of
NLDD local linear solver. Valid options are: ilu0 (default), dilu,
cpr\_quasiimpes and amg. Alternatively, you can request a configuration
to be read from a JSON file by giving the filename here, ending with
`.json.'. & ``ilu0'' \\
127 & --nldd-local-linear-solver-max-iter & An integer value. The
maximum number of iterations of the NLDD local linear solver. & 200 \\
128 & --nldd-local-linear-solver-reduction & A positive real value. The
minimum reduction of the residual which the NLDD local linear solver
must achieve. & 0.01 \\
129 & --nldd-num-initial-newton-iter & An integer value. Number of
initial global Newton iterations when running the NLDD nonlinear solver.
& 1 \\
130 & --nldd-relative-mobility-change-tol & A positive real value.
Threshold for single cell relative mobility change in the NLDD solver. &
0.1 \\
131 & --nonlinear-solver & A character string. Choose nonlinear solver.
Valid choices are newton or nldd. & ``newton'' \\
132 & --num-local-domains & An integer value. Number of local domains
for NLDD nonlinear solver. & 0 \\
133 & --num-overlap & An integer value. Numbers of layers overlap in
parallel partition. & 1 \\
134 & --num-pressure-points-equil & An integer value. Number of pressure
points (in each direction) in tables used for equilibration. & 2000 \\
135 & --num-satfunc-consistency-sample-points & An integer value.
Maximum number of reported failures for each individual saturation
function consistency check. & 5 \\
136 & --nupcol-group-rate-tolerance & A positive real value. Tolerance
for acceptable changes in VREP/REIN group rates. & 0.001 \\
137 & --opencl-ilu-parallel & A boolean value. Parallelize ILU
decomposition and application on GPU. & true \\
138 & --opencl-platform-id & An integer value. Choose platform ID for
openclSolver, use `clinfo' to determine valid platform IDs. & 0 \\
139 & --output-dir & A character string. The directory to which result
files are written. & ``\,'' \\
140 & --output-extra-convergence-info & A defined character string.
Specify whether to provide additional convergence output to separate
files for diagnostic purposes. The available options are: Options can be
combined with commas, e.g.''steps,iterations'' for multiple outputs. The
default value of ``none'' prevents the two files from being written out,
for better compatibility with the commercial simulator. & ``none'' \\
141 & --output-interval & An integer value. Specify the number of report
steps between two consecutive writes of restart data. & 1 \\
142 & --output-mode & A defined character string. Specify the output to
the \emph{.PRT and }.DBG files: For example to just output logging
information use: --output-mode=``log'' or --output-mode=false &
``all'' \\
143 & --owner-cells-first & A boolean value. Order cells owned by rank
before ghost/overlap cells. & true \\
144 & --parameter-file & A character string. An .ini file which contains
a set of run-time parameters. & ``\,'' \\
145 & --parsing-strictness & A character string. Set strictness of
parsing process. Available options are: & ``normal'' \\
146 & --partition-method & Options are now specified by strings rather
than integers. A character string. Choose partitioning strategy: &
``zoltanwell'' \\
147 & --pre-solve-network & A boolean value. Pre solve and iterate the
network model at start-up. & true \\
148 & --predetermined-time-steps-file & A character string. A file with
a list of predetermined time step sizes (one time step per line). &
``\,'' \\
149 & --pressure-max & A positive real value. Maximum absolute pressure.
& 1e+99 \\
150 & --pressure-min & A real value. Minimum absolute pressure. &
-1e+99 \\
151 & --pressure-scale & A positive real value. Scaling of pressure
primary variable. & 1 \\
152 & --pri-var-oscilation-threshold & A positive real value. The
threshold value for the primary variable switching conditions after its
meaning has switched to hinder oscillations. & 1e-05 \\
153 & --print-parameters & An integer value. Specify the run time
parameters to be printed at the start of the run: & 2 \\
154 & --project-saturations & A boolean value. Option for doing
saturation projection. & false \\
155 & --regularization-factor-wells & A positive real value.
Regularization factor for wells. & 100 \\
156 & --relaxed-linear-solver-reduction & A positive real value. The
minimum reduction of the residual which the linear solver need to
achieve for the solution to be accepted. & 0.01 \\
157 & --relaxed-max-pv-fraction & A positive real value. The fraction of
the pore volume of the reservoir where the volumetric error (CNV) may be
violated during strict Newton iterations. & 0.03 \\
158 & --relaxed-pressure-tol-msw & A positive real value. Relaxed
tolerance for the MSW pressure solution. & 10000 \\
159 & --relaxed-well-flow-tol & A positive real value. Relaxed tolerance
for the well flow residual. & 0.001 \\
160 & --restart-time & A real value. The simulation time at which a
restart should be attempted {[}s{]}. & -1e+35 \\
161 & --restart-writing-interval & An integer value. The frequencies of
which time steps are serialized to disk. & 16777215 \\
162 & --save-file & A character string. FileName for .OPMRST file used
for saving serialized state. If empty, CASENAME.OPMRST is used. &
``\,'' \\
163 & --save-step & A character string. Specify if the simulator should
save the serialized state of the OPM Flow simulator at one or more
report steps to a special *.OPMRST file. This is in addition to the
normal restart files written, and consumes significantly more space than
the normal restart files, but restarting OPM Flow from this file using
the --load-step option deviates less from the original run, compared to
restarting using the normal restart features. The files produced are not
compatible with other simulators, and also will not be compatible
between different releases of OPM Flow. The parameter should be set to
one of the following: Use negative values of ``x'' to keep only the last
written step. The default value of ``\,'' does not write anything to the
OPM Flow specific restart file. & ``\,'' \\
164 & --scale-linear-system & A boolean value. Scale linear system
according to equation scale and primary variable types. & false \\
165 & --sched-restart & A boolean value. Determine for a restart run if
the case should initialize wells and groups from the historical SCHEDULE
section (true), or from the well and group data on the restart file
(false). Note that the commercial simulator always uses data from the
restart file (false). & false \\
166 & --serial-partitioning & A boolean value. Perform partitioning for
parallel runs on a single process. & false \\
167 & --shut-unsolvable-wells & A boolean value. Shut unsolvable wells.
& true \\
168 & --slave & A boolean value. Specify if the simulation is a slave
simulation in a master-slave simulation. Note that reservoir coupling is
not currently supported. This is a developer option that should not be
used. & false \\
169 & --solve-welleq-initially & A boolean value. Fully solve the well
equations before each iteration of the reservoir model. Note that the
well equations are always added to the full system and solved until
convergence. & true \\
170 & --solver-continue-on-convergence-failure & A boolean value.
Continue instead of stop when minimum solver time step is reached. &
false \\
171 & --solver-growth-factor & A real value greater than one. Specify
the growth factor a time step can be increased by when recovering from
one or more time step chops, subject to the maximum allowable time step
size set by the --solver-max-time-step-in-days parameter. For example,
if the current time step has converged at 5 days after a time step chop,
and --solver-growth-factor is set to the default value of 2.0, then the
next time step will be 2.0 * 5 days, that is at 10 days. & 2 \\
172 & --solver-max-growth & A real value greater than one. Specify the
maximum growth factor a time step can be increased by after a report
time step, subject to the maximum allowable time step size set by the
--solver-max-time-step-in-days parameter. Thus, if the current time step
has converged at 5 days after at a report step, and --solver-max-growth
is set to the default value of 3.0, then the next time step will be 3.0
* 5 days, that is at 15 days. & 3 \\
173 & --solver-max-restarts & An integer value. The maximum number of
breakdowns before a substep is given up and the simulator is terminated.
& 10 \\
174 & --solver-max-time-step-in-days & A positive real value. The
maximum size of a time step in days. & 365 \\
175 & --solver-min-time-step & A positive real value. The minimum size
of a time step in days for field and metric and hours for lab. If a step
cannot converge without getting cut below this step size the simulator
will stop. & 1e-12 \\
176 & --solver-restart-factor & A real value between zero and one.
Specify the time step chop factor of the time step after a convergence
failure. For example, if the current non-convergent time step is 30 days
and --solver-restart-factor is set to the default value, then the time
step will be repeated using 0.33 * 30 days as the time step, that is 9.9
days. & 0.33 \\
177 & --solver-verbosity & An integer value. Specify the ``chattiness''
of the non-linear solver itself. & 1 \\
178 & --strict-inner-iter-wells & An integer value. Number of inner well
iterations with strict tolerance. & 40 \\
179 & --strict-outer-iter-wells & An integer value. Number of newton
iterations for which wells are checked with strict tolerance. & 6 \\
180 & --temperature-max & A positive real value. Maximum absolute
temperature. & 1e+09 \\
181 & --temperature-min & A real value. Minimum absolute temperature. &
0 \\
182 & --threads-per-process & An integer value. The maximum number of
threads to be instantiated per process (`-1' means `automatic'). & 2 \\
183 & --time-step-after-event-in-days & A real value. Specify the
maximum allowed time step after an event; for example, when a well is
open or closed etc. The default value of -1 means that events do not
affect the time stepping. & -1 \\
184 & --time-step-control & A character string. The algorithm used to
determine time-step sizes. Valid options are: &
``pid+newtoniteration'' \\
185 & --time-step-control-decay-damping-factor & A positive real value
less than or equal to one. The decay rate of the time step decrease when
the target iterations is exceeded. & 1 \\
186 & --time-step-control-decay-rate & A positive real value less than
one. The decay rate of the time step size of the number of target
iterations is exceeded. & 0.75 \\
187 & --time-step-control-file-name & A character string. Specify the
name of a file where time steps are specified. For instance time steps
can be generated by the summary application as per the following UNIX
command line: path\_to\_application/summary DECK TIME \textgreater{}
timesteps Where: DECK is the name of the data deck you want to get the
time steps from, TIME tells the application to return the timing for the
run, and ``filename'' is the name of the file the times are piped to. &
``timesteps'' \\
188 & --time-step-control-growth-damping-factor & A real value greater
than one. The growth rate of the time step increase when the target
iterations is undercut. & 3.2 \\
189 & --time-step-control-growth-rate & A real value greater than one.
The growth rate of the time step size of the number of target iterations
is undercut. & 1.25 \\
190 & --time-step-control-max-reduction-time-step & A positive real
value (only applicable for the general 3rd order controller, using
`control-error-filtering' as time-step-control-tolerance-test-version).
If the (proposed) relative change in time step size is larger than this
parameter, the time step will be rejected. & 0.1 \\
191 & --time-step-control-parameters & A character string (only
applicable for the general 3rd order controller). Parameters for the
general 3rd order controller. Should be given as
`beta\_1;beta\_2;beta\_3;alpha\_2;alpha\_3'. &
``0.125;0.25;0.125;0.75;0.25'' \\
192 & --time-step-control-reject-completed-step & A boolean value (only
applicable for the general 3rd order controller). Include rejection of
completed time steps if the relative change is larger than the time step
control tolerance. & false \\
193 & --time-step-control-safety-factor & A positive real value less
than or equal to one. Value to be multiplied with the time step control
tolerance to ensure that the target relative change is lower than the
tolerance. & 0.8 \\
194 & --time-step-control-target-iterations & An integer value. The
number of linear iterations which the time step control scheme should
aim for (if applicable). & 30 \\
195 & --time-step-control-target-newton-iterations & An integer value.
The number of Newton iterations which the time step control scheme
should aim for (if applicable). & 8 \\
196 & --time-step-control-tolerance & A positive real value. The
tolerance used by the time step size control algorithm. & 0.1 \\
197 & --time-step-control-tolerance-test-version & A character string
(only applicable for the general 3rd order controller). Ways to decide
if the time step should be rejected. Options: `standard' and
`control-error-filtering'. The standard version compares relative change
to tolerance directly to decide if the time step should be rejected,
while the control-error-filtering version compares the relative change
in time step size to the max-reduction-time-step parameter. &
``standard'' \\
198 & --time-step-verbosity & An integer value. Specify the
``chattiness'' during the time integration. & 1 \\
199 & --tolerance-cnv & A positive real value. Local convergence
tolerance (maximum of local saturation errors). & 0.01 \\
200 & --tolerance-cnv-energy & A positive real value. Local energy
convergence tolerance (maximum of local energy errors). & 0.01 \\
201 & --tolerance-cnv-energy-relaxed & A positive real value. Relaxed
local energy convergence tolerance that applies for iterations after the
iterations with the strict tolerance. & 1 \\
202 & --tolerance-cnv-relaxed & A positive real value. Relaxed local
convergence tolerance that applies for iterations after the iterations
with the strict tolerance. & 1 \\
203 & --tolerance-energy-balance & A positive real value. Tolerated
energy balance error relative to (scaled) total energy present. &
1e-07 \\
204 & --tolerance-energy-balance-relaxed & A positive real value.
Relaxed tolerated energy balance error that applies for iterations after
the iterations with the strict tolerance. & 1e-06 \\
205 & --tolerance-mb & A positive real value. Tolerated mass balance
error relative to total mass present. & 1e-07 \\
206 & --tolerance-mb-relaxed & A positive real value. Relaxed tolerated
mass balance error that applies for iterations after the iterations with
the strict tolerance. & 1e-06 \\
207 & --tolerance-pressure-ms-wells & A positive real value. Tolerance
for the pressure equations for multi-segment wells. & 1000 \\
208 & --tolerance-well-control & A positive real value. Tolerance for
the well control equations. & 1e-07 \\
209 & --tolerance-wells & A positive real value. Well convergence
tolerance. & 0.0001 \\
210 & --update-equations-scaling & A boolean value. Update scaling
factors for mass balance equations during the run. & false \\
211 & --use-average-density-ms-wells & A boolean value. Approximate
segment densitities by averaging over segment and its outlet. & false \\
212 & --use-gmres & A boolean value. Use the Generalized Minimal
Residual (``GMRES'') {[}Y. Saad, A flexible inner-outer preconditioned
GMRES algorithm, SIAM J. Sci. Statist. Comput.,14, (1993).{]} and {[}Y.
Saad and M.H. Schultz, GMRES: A generalized minimum residual algorithm
for solving nonsymmetric linear systems, SIAM J. Sci. Statist. Comput.,
7 (1986), pp.~856\{869.{]} solver instead of Biconjugate Gradient
Stabilized (``BiCGSTAB'') {[}Van der Vorst, H. A. (1992). ``Bi-CGSTAB: A
Fast and Smoothly Converging Variant of Bi-CG for the Solution of
Nonsymmetric Linear Systems''. SIAM J. Sci. Stat. Comput. 13 (2):
631--644. doi:10.1137/0913035. hdl:10338.dmlcz/104566{]} and
{[}Sleijpen, G. L. G.; Fokkema, D. R. (November 1993). ``BiCGstab(l) for
linear equations involving unsymmetric matrices with complex spectrum''
(PDF). Electronic Transactions on Numerical Analysis. Kent, OH: Kent
State University. 1: 11--32. ISSN 1068-9613.{]} as the linear solver
within the Newton iterations. & false \\
213 & --use-hy-ne & A boolean value. Wheter or not to use Hybrid Newton.
& false \\
214 & --use-implicit-ipr & A boolean value. Compute implict IPR for
stability checks and stable solution search. & true \\
215 & --use-multisegment-well & A boolean value. Use the well model for
multi-segment wells instead of the one for single-segment wells. &
true \\
216 & --use-update-stabilization & A boolean value. Try to detect and
correct oscillations or stagnation during the Newton method. & true \\
217 & --verify-gpu-aware-mpi & A boolean value. Verify that the MPI
implementation supports GPU aware MPI. If this is set to true \emph{and}
--gpu-aware-mpi=true, the simulation will fail if the MPI implementation
does not support GPU aware MPI. Note that the verification is not
exhaustive, and some configurations will not verify, but will work in
practice. Usage: --verify-gpu-aware-mpi={[}true\textbar false{]}. . &
false \\
218 & --water-only-threshold & A positive real value. Cells with water
saturation above or equal is considered one-phase water only. & 1 \\
219 & --well-group-constraints-max-iterations & An integer value.
Maximum number of iterations in the well/group switching algorithm. &
1 \\
220 & --zoltan-imbalance-tol & A real value. Tolerable imbalance of the
loadbalancing provided by Zoltan. DEPRECATED: Use --imbalance-tol
instead. & 1.1 \\
221 & --zoltan-params & A character string. Configuration of Zoltan
partitioner. Valid options are: graph, hypergraph or scotch.
Alternatively, you can request a configuration to be read from a JSON
file by giving the filename here, with extension `.json'. See
https://sandialabs.github.io/Zoltan/ug\_html/ug.html for available
Zoltan options. & ``graph'' \\
222 & --zoltan-phg-edge-size-threshold & A real value ranging from zero
to one. Low-level threshold fraction in the range {[}0,1{]} controlling
which hypergraph edge to omit. Used if --zoltan-params=``graph'' or if
--zoltan-params=``hypergraph''. & 0.35 \\
Notes: & & & \\
\end{longtable}

\emph{Table 2.1: OPM Flow Command Line Options}

As per all UNIX and LINUX based systems the input is case sensitive. The
command line parameters must either be entered as lower case or
CamelCase, for example:

\begin{verbatim}
flow --enable-dry-run=false CASENAME.DATA
\end{verbatim}

or:

\begin{verbatim}
flow --EnableDryRun=false CASENAME.DATA
\end{verbatim}

If a parameter file is used to set one of the above parameters, the
leading ``--'' should not be incorporated in the parameter file, for
example on the command line use:

\begin{verbatim}
flow --enable-dry-run=false CASENAME.DATA
\end{verbatim}

but in the parameter file use:

\begin{verbatim}
ecl-deck-file-name=CASENAME.DATA

enable-dry-run=false
\end{verbatim}

or:

\begin{verbatim}
EclDeckFileName=CASENAME.DATA

EnableDryRun=false
\end{verbatim}

OPM Flow prints out the command line run time and compile time
parameters in CamelCase in both the \emph{.PRT and }.DBG files for
reference, one can therefore use this as basis for creating additional
run specific parameter files.

\begin{tcolorbox}[enhanced jigsaw, title=\textcolor{quarto-callout-note-color}{\faInfo}\hspace{0.5em}{Note}, coltitle=black, breakable, toprule=.15mm, colbacktitle=quarto-callout-note-color!10!white, leftrule=.75mm, opacitybacktitle=0.6, colback=white, bottomtitle=1mm, toptitle=1mm, titlerule=0mm, left=2mm, colframe=quarto-callout-note-color-frame, arc=.35mm, rightrule=.15mm, bottomrule=.15mm, opacityback=0]

It is worth considering having a parameter file for each run of the form
CASENAME.param, in order to re-run a case using the same parameters or
to use the same parameters for other cases.

\end{tcolorbox}

\subsubsection{Example}\label{example}

The following example shows the parameter file NORNE\_ATW2013.PARAM
taken from the Norne project, where the comments below (preceded with
``\#'') explain the parameter settings used in the file.

\begin{verbatim}
# ----------------------------------------------------------------------------
# INPUT AND OUTPUT OPTIONS
# ----------------------------------------------------------------------------
#
# Input File
#
ecl-deck-file-name=NORNE_ATW2013.DATA
#
# Output and Output Directory
#
output-dir=’\$HOME\OPM\NORNE\’
# ----------------------------------------------------------------------------
# NEWTON SOLVER PARAMETER
# ----------------------------------------------------------------------------
#
# Define Numerical Tolerances
#
tolerance-cnv=1e-2
tolerance-mb=1e-5
tolerance-wells=1e-2
#
# Set Min Newton Solver iterations to 1 and Max to 15
#
newton-min-iterations=1
newton-max-iterations=15
#
# ----------------------------------------------------------------------------
\end{verbatim}

Notice that the leading ``--'' have not be incorporated in the parameter
file, as per the notes in Table 2.1.

In order to run the deck file NORNE\_ATW2013.DATA in OPM Flow using the
above parameter file one could use either of the following commands:

\begin{verbatim}

flow --parameter-file=NORNE_ATW2013.PARAM
\end{verbatim}

or:

\begin{verbatim}
flow --ParameterFile=NORNE_ATW2013.PARAM
\end{verbatim}

\section{Running OPM Flow Using
OPMRUN}\label{running-opm-flow-using-opmrun}

This section briefly describes how to run the OPM Flow simulator using
the Graphical User Interface (``GUI'') program called OPMRUN, a more
detailed description is given in Error: Reference source not found
Error: Reference source not found, that also describes the various tools
included with OPMRUN. The program is written in Python 3 and has been
tested under Ubuntu-Mate 24.04 TLS. The software can be downloaded from
the following link:

\url{https://github.com/OPM/opm-utilities/tree/master/opmrun}

The intent is to develop a graphical user interface to OPM Flow that has
similar functionality to the commercial simulator's program, with the
targeted audience being Reservoir Engineers in a production environment.
Developers and experienced Linux users will already have compatible work
flows. OPMRUN enables the editing and management of OPM Flow's run time
parameters, setting up job queues to run a series of simulation jobs
sequentially, as well as the management of the job queues. Figure 2.10
shows the initial display.

\pandocbounded{\includegraphics[keepaspectratio]{chapters/images/Frame19_f5dfac27d930.png}}
\pandocbounded{\includegraphics[keepaspectratio]{chapters/images/Image17_f5dfac27d930.png}}

Upon launch the program runs OPM Flow to get a list of command line
parameters from the current version of OPM Flow. These default
parameters can be edited for each case, or alternative default parameter
sets can be loaded from an existing parameter file from another job, or
a *.PRT file from a completed simulation.

As can be seen in Figure 2.11 the program has upper and lower display
elements. The upper element shows a list of simulation jobs that are in
the job queue and the lower element consists of two elements, one for
the OPM Flow Output (the terminal output from OPM Flow) and a second
element (OPM Run Log) that is a session log of the jobs run by OPMRUN.
Clicking the OPM Flow Output and OPM Run Log tabs switches the display
on the lowered element between two display types.

\pandocbounded{\includegraphics[keepaspectratio]{chapters/images/Frame20_fc8a971c9ec3.png}}
\pandocbounded{\includegraphics[keepaspectratio]{chapters/images/Image18_fc8a971c9ec3.png}}

To add jobs to the queue use the Add Job button or load an existing job
queue using the Load Queue button. Jobs can be edited or deleted from
the queue using the Edit Job and Delete Job buttons, and a series of
jobs can be saved as a job queue by using the Save Queue button. The
Clear Queue button deletes all jobs from the queue.

Pressing the Add Job button will display the following dialogue box:

\pandocbounded{\includegraphics[keepaspectratio]{chapters/images/Frame21_557019afb7dc.png}}
\pandocbounded{\includegraphics[keepaspectratio]{chapters/images/Image19_557019afb7dc.png}}

Use the Browse button to select the input file to add to the queue, then
select the Run Parameters for this input file, then press the Submit
button to add the input file to the job queue.

The Overwrite Parameter File Options allow for different default
treatments of existing \emph{.PARAM files, which is particularly useful
when adding multiple jobs at the same time. The reason for this is
because different versions of OPM Flow have different parameter sets and
if a newer version of OPM Flow runs with a previous version's }.PARAM
file then the simulator will stop with an error message for the various
invalid parameters for the current version of the simulator. To load a
previously saved job queue, press the Load Queue button this will
display a dialog box allowing the user to select a queue file (*.que),
after pressing the OK button the jobs will be displayed in the Job List
Element as illustrated in Figure 2.13.

\pandocbounded{\includegraphics[keepaspectratio]{chapters/images/Frame22_0d01d65c6763.png}}
\pandocbounded{\includegraphics[keepaspectratio]{chapters/images/Image20_0d01d65c6763.png}}

Jobs in the queue can be edited by selecting the Edit Job button that
will display two options: one to edit the input file using the defined
editor and the second to edit the OPM Flow Parameter File. If the second
option is selected OPMRUN will display a dialog box that shows a list of
the OPM Flow command line parameters together with the parameter help
information (Figure 2.14). Selecting a parameter from the list and
selecting the Edit button will display the setting for the selected
parameter (alternatively one can double click the required entry). One
can then edit the parameter as required. Use the Save button to save the
change and use the Exit but to save all the changes to the parameter
file. The Cancel button will cancel all changes to the parameter file.

Alternatively one can use the:

\begin{verbatim}
- Edit OPM Flow Parameter menu option to edit the parameter file for a job.
- List OPM Flow Parameters menu option to list the commands in the parameter file for a job.
- Set OPM Flow Default Parameters to set the default parameters for all subsequent jobs added to the queue. This option allows the user to load a default set of parameters from (1) OPM Flow, (2) an OPM Flow Parameter File, or (3) an OPM Flow print file (*.PRT).
\end{verbatim}

One can also right-click on a job in the Job List Element and select one
of the available options.

\pandocbounded{\includegraphics[keepaspectratio]{chapters/images/Frame23_30b2cadad0e0.png}}
\pandocbounded{\includegraphics[keepaspectratio]{chapters/images/Image21_30b2cadad0e0.png}}

Selecting the Run Jobs button displays the Select Run Option dialog box
shown in Figure 2.15.

\pandocbounded{\includegraphics[keepaspectratio]{chapters/images/Frame124_e58f1a43b3ac.png}}
\pandocbounded{\includegraphics[keepaspectratio]{chapters/images/Image124_e58f1a43b3ac.png}}

The Run in No Simulation Mode option is equivalent to setting the NOSIM
option in the input deck for all jobs in the queue (see see section
Error: Reference source not found Error: Reference source not found and
the --enable-dry-run command line parameter in Table 2.1 in section 2.2
Running OPM Flow From The Command Line. This allows for checking all the
jobs at once.

Selecting Run in Standard Simulation Mode will run all the jobs in the
queue sequentially, with the OPM Flow terminal output directed to OPM
Flow Output Element, as shown in Figure 2.16. The terminal output is
also directed to a *.LOG file as well, similar to what the commercial
simulator does.

\pandocbounded{\includegraphics[keepaspectratio]{chapters/images/Frame25_601e0d9ae485.png}}
\pandocbounded{\includegraphics[keepaspectratio]{chapters/images/Image23_601e0d9ae485.png}}

Clicking the OPM Run Log tab displays the OPMRUN's session log file that
records the time and date of the major events that have occurred,
including the start and end times of each run. Notice also how OPMRUN
deletes all the existing output files for a given job, if they exists,
before running OPM Flow.

\pandocbounded{\includegraphics[keepaspectratio]{chapters/images/Frame26_6dee5c16282c.png}}
\pandocbounded{\includegraphics[keepaspectratio]{chapters/images/Image24_6dee5c16282c.png}}

The Kill button will ask the user if the current running job should be
killed, and if the job is to be killed, the application will prompt as
whether or not all the jobs in the queue should be killed.

The Clear clears the OPM Flow Output Element from the currently
displayed tab (Output or Log) and the Copy button copies the data to the
clipboard.

Finally, an example *.param file generated by OPMRUN for the 2022-10
release is shown below and on the following pages for reference:

\begin{verbatim}
#
# OPMRUN Parameter File
#
# File Name   : "/media/Fields/Norne/Model/NOR-BASE-A00.param"
# Created By  : opmuser
# Date Created: 2022-11-01 13:12:44
#
accelerator-mode="none"
allow-distributed-wells=false
alternative-well-rate-init=true
bda-device-id=0
continue-on-convergence-error=false
cpr-max-ell-iter=20
cpr-reuse-interval=10
cpr-reuse-setup=3
dbhp-max-rel=1
dp-max-rel=0.3
ds-max=0.2
dwell-fraction-max=0.2
ecl-deck-file-name="NOR-BASE-A00.DATA"
ecl-enable-drift-compensation=true
ecl-output-double-precision=false
ecl-output-interval=-1
ecl-strict-parsing=false
edge-weights-method=1
enable-adaptive-time-stepping=true
enable-async-ecl-output=true
enable-async-vtk-output=true
enable-dry-run="auto"
enable-ecl-output=true
enable-esmry=false
enable-logging-fallout-warning=false
enable-opm-rst-file=false
enable-storage-cache=true
enable-terminal-output=true
enable-tuning=false
enable-vtk-output=false
enable-well-operability-check=true
enable-well-operability-check-iter=false
enable-write-all-solutions=false
force-disable-fluid-in-place-output=false
force-disable-resv-fluid-in-place-output=false
fpga-bitstream=""
full-time-step-initially=false
ignore-keywords=""
ilu-fillin-level=0
ilu-redblack=false
ilu-relaxation=0.9
ilu-reorder-spheres=false
initial-time-step-in-days=1
linear-solver="ilu0"
linear-solver-ignore-convergence-failure=false
linear-solver-max-iter=200
linear-solver-reduction=0.01
linear-solver-require-full-sparsity-pattern=false
linear-solver-restart=40
linear-solver-verbosity=0
matrix-add-well-contributions=false
max-inner-iter-ms-wells=100
max-inner-iter-wells=50
max-newton-iterations-with-inner-well-iterations=8
max-pressure-change-ms-wells=1e+06
max-residual-allowed=1e+07
max-single-precision-days=20
max-temperature-change=5
max-welleq-iter=30
maximum-number-of-well-switches=3
milu-variant="ILU"
min-strict-cnv-iter=0
min-time-step-based-on-newton-iterations=0
min-time-step-before-shutting-problematic-wells-in-days=0.01
newton-max-iterations=20
newton-max-relax=0.5
newton-min-iterations=1
newton-relaxation-type="dampen"
opencl-ilu-parallel=true
opencl-platform-id=0
output-dir=""
output-interval=1
output-mode="all"
owner-cells-first=true
parameter-file=""
pri-var-oscilation-threshold=1e-05
print-parameters=2
print-properties=2
project-saturations=false
regularization-factor-wells=100
relaxed-max-pv-fraction=0.03
relaxed-pressure-tol-msw=10000
relaxed-well-flow-tol=0.001
scale-linear-system=false
sched-restart=false
serial-partitioning=false
shut-unsolvable-wells=true
solve-welleq-initially=true
solver-continue-on-convergence-failure=false
solver-growth-factor=2
solver-max-growth=3
solver-max-restarts=10
solver-max-time-step-in-days=365
solver-min-time-step=1e-12
solver-restart-factor=0.33
solver-verbosity=1
strict-inner-iter-wells=40
strict-outer-iter-wells=6
temperature-max=1e+09
temperature-min=0
threads-per-process=-1
time-step-after-event-in-days=-1
time-step-control="pid+newtoniteration"
time-step-control-decay-damping-factor=1
time-step-control-decay-rate=0.75
time-step-control-file-name="timesteps"
time-step-control-growth-damping-factor=3.2
time-step-control-growth-rate=1.25
time-step-control-target-iterations=30
time-step-control-target-newton-iterations=8
time-step-control-tolerance=0.1
time-step-verbosity=1
tolerance-cnv=0.01
tolerance-cnv-relaxed=1
tolerance-mb=1e-06
tolerance-pressure-ms-wells=1000
tolerance-well-control=1e-07
tolerance-wells=0.0001
update-equations-scaling=false
use-gmres=false
use-multisegment-well=true
use-update-stabilization=true
zoltan-imbalance-tol=1.1
#
# End of Parameter File
\end{verbatim}

In the above example, all the default parameters have be used.

In addition to the functionality described in this section, the
application has various additional features accessible through the
application's menus including:

\begin{verbatim}
- Configuring OPMRUN various options, default editor, default terminal console for background jobs, setting default project directories, etc.
- Various additional simulation input generation and conversion utilities are available including:

  - Compressing a job to save space (DATA, and all OPM Flow output files) and uncompressing previously compressed jobs.
  - Keywords, a keyword generator based on the Apache Velocity Template Language (“VTL”). The templates can therefore also be used with any editor that supports VTL, jEdit for example. There is one template per keyword, with the formatting the same as the OPM Flow manual. Over 450 templates are currently implemented. One can also customize the existing templates as well as creating User defined templates. The keywords are examples, one still has to edit the resulting deck with the actual required data, but the format with comments should make this a straight forward process.
  - A Production Schedule application that takes a comma delimited CSV file containing historical production and injection data and converts the data to an OPM Flow SCHEDULE file using the WCONHIST series of keywords. Currently only production data is supported.
  - Sensitivities application that generates sensitivity cases based on a "Base" case file. The Base file contains "Factors" (variable names), $X01, $X02, etc., that are substituted with user defined values using the data entered and the type of Sensitivity Scenario selected.
  - A Well Specification application that uses the standard well export files from OPM ResInsight to reformat the data in a more user-friendly manner for the WELSPECS and COMPDAT keywords. Optionally, the application can generate the COMPLUMP keyword based on the OPM ResInsight layers file, with one completion per defined reservoir layer.
  - Calling OPM ResInsight and loading the currently selected job into OPM ResInsight for viewing.
  - A Well Trajectory Conversion application that converts a Schlumberger Petrel exported well trajectory file into an OPM ResInsight file, containing all the wells.
\end{verbatim}

All the aforementioned options and tools are described in detail in
Error: Reference source not foundError: Reference source not found,

\begin{tcolorbox}[enhanced jigsaw, title=\textcolor{quarto-callout-note-color}{\faInfo}\hspace{0.5em}{Note}, coltitle=black, breakable, toprule=.15mm, colbacktitle=quarto-callout-note-color!10!white, leftrule=.75mm, opacitybacktitle=0.6, colback=white, bottomtitle=1mm, toptitle=1mm, titlerule=0mm, left=2mm, colframe=quarto-callout-note-color-frame, arc=.35mm, rightrule=.15mm, bottomrule=.15mm, opacityback=0]

OPMRUN can be run as either a Linux or a Windows application. If the
software is installed and run under Windows, then the program
automatically handles the file naming conventions of the two operating
systems. That is under Windows, the application uses the Windows file
name to select the files, and then automatically converts the file name
to the equivalent Linux file name to run under WSL.

\end{tcolorbox}

\section{Improving Simulator Convergence and Numerical
Performance}\label{improving-simulator-convergence-and-numerical-performance}

In general, most simulator performance issues are caused by the data
used by the simulator, rather than the numerical algorithms used to
solve the equations presented in sections Error: Reference source not
found Error: Reference source not found through Error: Reference source
not found Error: Reference source not found. Thus, it is important to
verify all warning messages and to address the issues raised by these
warnings. Naturally, some warnings can be safely ignored, while others
should not. Just because the simulator runs with warning messages does
not mean that the results are reasonable, or that numerical performance
has not been impacted.

This section is therefore broken into two parts, the first part
discusses potential data errors and how to address them, and the second
section outlines the numerical controls available to improve the
simulator's performance. The latter should only be done after the data
quality issues have been rectified, although, it is not uncommon for
users to skip the data quality step altogether, and jump to ``tuning''
the model's numerical controls - to be clear, skipping the data quality
step is strongly not recommended.

\subsection{Data Quality and
Assurance}\label{data-quality-and-assurance}

\subsubsection{GRID Section}\label{grid-section}

Grid Configuration

In full field models grid distortions may occur due to faulting and the
structural complexity of the geological model, and ideally one wishes to
have all the cells to be orthogonal. Grid cells are orthogonal if all
grid lines intersect at a right angles, that is at 90o. This is because
orthogonal grids offer significant advantages in solving a system of
partial differential equations used computational fluid dynamics. In
addition, they ensure that the fluid flow is not being unduly influenced
by the grid orientation effects, as cell to cell flow only occurs across
connected cell faces.

Most earth modeling software can produce a histogram of the cell angles
as an aide to checking the orthogonality of the grid. A common cause for
cells not being orthogonal, as well as creating other grid artifacts, is
the grid cells following actual the faulting traces, as oppose to using
``zig-zag'' fault traces that follow the grid. A typical fault trace
grid is illustrated in Figure 2.18.

\pandocbounded{\includegraphics[keepaspectratio]{chapters/images/Frame148_7c86ced6c66a.png}}
\pandocbounded{\includegraphics[keepaspectratio]{chapters/images/Image141_7c86ced6c66a.png}}

Notice in Figure 2.18 how the cells along the fault plane are severely
distorted, with some cells having only three faces instead of four in
the x-y plane. Now compare this to the Norne model shown in Figure 2.19,
which has more complex faulting, and yet the grid along the fault planes
is mainly orthogonal.

\pandocbounded{\includegraphics[keepaspectratio]{chapters/images/Frame146_942fda1d79a8.png}}
\pandocbounded{\includegraphics[keepaspectratio]{chapters/images/Image139_942fda1d79a8.png}}

Grid Property Data and Active Cells

There is a current trend in the industry for not applying cut-offs to
the petrophysical parameters used to generate the grid property data
(porosity, net-to-gross, etc.). The philosophy behind this approach is
that the simulator should determine what is reservoir and is capable of
contributing to flow within the reservoir, and what is non-reservoir. At
first glance this would appear reasonable, but unfortunately the
inclusion of low pore volume and low permeability grid cells will have a
dramatic negative impact on the model's numerical convergence and run
times. This is because one ends up with many more active cells than
otherwise, plus the fact that low pore volumes grid blocks will be
connected to high pore volume cells. Because of this
inter-connectedness, a change in a large pore volume grid block will
have disproportional change on a low pore volume cell, which will result
in smaller time step sizes in order to satisfy the numerical controls.
{[}The large pore volume differences make the solution of the linear
system of equations more challenging, due to the residual errors in
certain grid cells or areas within the grid, that are much harder to
reduce with iterative solvers. This results in more convergence
failures, smaller time steps, etc. Similarly, for highly varying cell
sizes.{]} That is, the low pore volume cell throughput will influence
the simulator's selected time step size disproportionately.

Nevertheless, even if the more conventional approach to the model build
is undertaken by applying some form of petrophysical cut-offs, there may
well still be grid blocks containing small pore volumes. In this case
the MINPV keyword in the GRID section can be used to set a minimum pore
volume threshold in which cells below the threshold value will be
automatically made inactive. It is not possible to give absolute
guidance for a suitable pore volume value to be used with MINPV, the
best approach is to run some sensitivities to investigate the impact of
different pore volume thresholds on the reduction in the number of cells
and the consequent change in the fluid in-place volumes. One should also
be mindful not to remove thin high permeability cells that have small
pore volumes due to their size, as this will effect the model's flow
behavior.

Finally, if the PINCH keyword in the GRID section is also in the input
deck, then one must be cognizant of the interaction between the two
keywords.

Grid Coarsening

As the number of active grid blocks directly effect a model's run time,
reducing the number of active cells by coarsening areas of model that
are of secondary importance, can be beneficial to the overall run time
of a model. An example of where this approach should be considered is
the modeling of aquifers. It is not necessary to fully model an aquifer,
by having the model cells cover the whole area of the aquifer. Instead,
incorporating either analytical aquifers via the Error: Reference source
not found and Error: Reference source not found keywords, or numerical
aquifers using the Error: Reference source not found keyword, can
effectively re-produce the aquifer response. Both types of aquifers have
advantages and disadvantages.

Analytical aquifers have the advantage that their response and behavior
can be matched by using material balance techniques to quantify the
aquifer properties, by history matching the measured pressure response
in the hydrocarbon zones. This is an efficient analytical modeling work
flow to identify the aquifer response, the properties derived can
therefore be more or less directly incorporated into the dynamic model.
The dis-advantage is that analytical aquifers are not directly
comparable with geological interpretations, which may, or may not, be
important, depending upon the circumstances.

The opposite is true for numerical aquifers in that an approximate
aquifer volume and property data can be matched with current geological
interpretation, but cannot be history matched in most material balance
software, as numerical aquifers in general are not supported.

Nevertheless, no mater which type of aquifer is employed, they should
not be directly connected to hydrocarbon bearing cells, ideally there
should be between three to five water bearing cells between the aquifer
connected cells and the hydrocarbon cells. However, the water cells may
vary in size, such that water cells connected to the hydrocarbon zone
are of similar size, and then increase in size towards the edge of the
model. The reason for this is that we wish to limit ``small''
hydrocarbon cells being directly coupled to ``large'' water cells, that
will hinder the length of time steps selected by the simulator due to
the convergence criteria.

Grid Isolated Cells

Simulation models should only include grid blocks that influence the
flow and recovery in the model. One common issue with some of the
geo-statistical formulations used to generate the grid and the
associated grid property data, is the tendency to create isolated grid
blocks. These isolated grid blocks are not connected to the rest of the
model and therefore do not influence the behavior of the model. Figure
2.20 shows a typical example, with the cells circled in red being the
isolated cells.

\pandocbounded{\includegraphics[keepaspectratio]{chapters/images/Frame147_2323a867c0db.png}}
\pandocbounded{\includegraphics[keepaspectratio]{chapters/images/Image140_2323a867c0db.png}}

Unfortunately, these cells will be included in all the calculations, and
may cause numerical issues due to their unchanging state, both pressure
and saturation. They will also unnecessarily add to the active cell
count, and therefore effect the simulation run time. In addition, they
will give a false impression of the average reservoir pressure, as the
isolated cell pressures will remain at initial conditions throughout the
simulation run. If model convergence is proving difficult, then this may
be a potential cause of the problem. In any case, it is recommended that
the identified isolated cells be removed by making them inactive via the
ACTNUM keyword in the GRID section.

If there are only a few isolated cells then they can manually be made
inactive; however, in many cases this not possible due to the number of
isolated grid blocks. In this case, it is better to revert back to the
static model and fix the issue in the earth modeling software. Indeed,
if there are numerous isolated grid blocks this may be indicative that
an inappropriate approach has been used to generate the static model,
although it may be also due to the depositional environment as well.

\subsubsection{PROPS Section}\label{props-section}

This section describes potential issues associated with the property
data and is split into two separate sub-sections:

\begin{verbatim}
- The first sub-section covers the rock properties, and in particular the saturation functions (relative permeability curves), that can have a dramatic impact on reservoir performance, especially for the non-linear solver. As will be shown later, it is important that the derivatives of the relative permeability curves are smooth and monotonically increasing or decreasing depending on the phase of the relative permeability fluid. Particular care must be taken if end-point scaling is being used, as it is not uncommon for the resulting de-normalized curves to be inappropriate, even though the normalized curves are reasonable and correct.
- The second sub-section addresses issues that may arise with the fluid properties, and inconsistent PVT data that may result in the simulation struggling to converge in the linear solver.
\end{verbatim}

PROPS Section: Saturation Function (Derivative)

In general, one should not enter laboratory derived relative
permeability data directly into the simulator. This is because it is
dubious that laboratory measurements on 1½ inch core plugs can represent
the flow in 100 x 100 x 1.0 m grid cell in the model. There are various
upscaling techniques that can account for the scaling issue, using
various forms of relative permeability curves, including the laboratory
curves, piston-like curves and thickness average curves etc. However,
the resulting pseudo relative permeability curves tend to not have
smooth derivatives, which will have a detrimental impact on convergence
efficiency. Thus, it is common practice to use the laboratory derived
curves, either in an upscaling workflow or as part of the input
workflow, and to fit a Corey17 type curve to the ``rock'' curves in
order to smooth out any discontinuities.

Unfortunately, neither OPM Flow nor the commercial simulator, unlike
some other simulators, support the direct entry of Corey type curves;
however, OPM Flow does support the more advanced and flexible LET family
of models instead. See section Error: Reference source not foundError:
Reference source not found and Lomeland et al {[}Lomeland F., Ebeltoft
E. and Thomas W.H., 2005. A New Versatile Relative Permeability
Correlation. Paper SCA2005-32 presented at the International Symposium
of the Society of Core Analysts held in Toronto, Canada, 21-25 August,
2005.{]} for further information on the model, as well as the Error:
Reference source not found, Error: Reference source not found, and the
Error: Reference source not found keywords in the PROPS section.

Corey {[}Corey, A. T. : ``The Interrelation Between gas and Oil Relative
Permeabilities'', Production Mon., 19. 38. (1954).{]} combined the work
of Purcell {[}Purcell, W. R., ``Capillary Pressures- Their Measurement
Using Mercury and the Calculation of Permeability Therefrom'',
Transactions AIME, 186, 39 (1949).{]} and Burdine {[}Burdine, N. T.,
``Relative Permeability Calculations from Pore Size Distribution Data'',
Transactions AIME, 198, 71 (1953).{]} that was widely accepted for its
simplicity. His original equations were developed for the drainage cycle
in water-wet sandstones, but have also been used in carbonate
formations. Corey's original water-oil equations were as follows:

\begin{equation}\phantomsection\label{eq-2-1}{
{k}_{\mathit{ro}} ({S}_{w}) = {\left[\frac{1 - {S}_{w }}{1 - {S}_{\mathit{wcr} }}\right]}^{{n}_{o}}
}\end{equation}

\begin{equation}\phantomsection\label{eq-2-2}{
{k}_{\mathit{rw}} ({S}_{w}) = {\left[\frac{{S}_{w} - {S}_{\mathit{wcr} }}{1 - {S}_{\mathit{wcr} }}\right]}^{{n}_{w}}
}\end{equation}

where:

kro(Sw) = relative permeability to oil,

krw(Sw) = relative permeability to water,

no = Corey oil exponent, set to four in the original paper,

nw = Corey water exponent, set to four in the original paper,

Sw = water saturation, and

Swcr = critical water saturation.

The denominator in equations (2.1) and (2.2) scales the water saturation
to the mobile water phase. There are several forms of these equations,
with the most common normalizing the saturation over the mobile
hydrocarbon phase, as depicted in equations (2.3) and (2.4).

\begin{equation}\phantomsection\label{eq-2-3}{
{k}_{\mathit{ro}} ({S}_{w}) = {k}_{\mathit{row}}{\left[\frac{1 - {S}_{w} - {S}_{\mathit{orw}}}{1 - {S}_{\mathit{orw}} - {S}_{\mathit{wcr}}}\right]}^{{n}_{o}}
}\end{equation}

\begin{equation}\phantomsection\label{eq-2-4}{
{k}_{\mathit{rw}} ({S}_{w}) = {k}_{\mathit{rww}}{\left[\frac{{S}_{w} - {S}_{\mathit{wc}}}{1 - {S}_{\mathit{orw}} - {S}_{\mathit{wcr}}}\right]}^{{n}_{w}}
}\end{equation}

where:

krow = maximum oil relative permeability at Swc,

krw = maximum water relative permeability st Sorw,

Swcr = critical water saturation, and

Sorw = residual oil saturation under a water flood (SOWCR).

\pandocbounded{\includegraphics[keepaspectratio]{chapters/images/Frame149_9902c207158d.png}}
\pandocbounded{\includegraphics[keepaspectratio]{chapters/images/Image142_9902c207158d.png}}

Figure 2.21 shows a typical Corey plot, with the plot at the top of the
figure showing the Corey plot and the bottom plot showing the derivative
plot, the data is also tabulated in Table 2.2. Notice in Figure 2.21
that the water curve does not extend to 100\% water saturation, this is
an error as the water below the oil-water contact would have restrictive
water flow.

\begin{longtable}[]{@{}
  >{\raggedright\arraybackslash}p{(\linewidth - 12\tabcolsep) * \real{0.1429}}
  >{\raggedright\arraybackslash}p{(\linewidth - 12\tabcolsep) * \real{0.1429}}
  >{\raggedright\arraybackslash}p{(\linewidth - 12\tabcolsep) * \real{0.1429}}
  >{\raggedright\arraybackslash}p{(\linewidth - 12\tabcolsep) * \real{0.1429}}
  >{\raggedright\arraybackslash}p{(\linewidth - 12\tabcolsep) * \real{0.1429}}
  >{\raggedright\arraybackslash}p{(\linewidth - 12\tabcolsep) * \real{0.1429}}
  >{\raggedright\arraybackslash}p{(\linewidth - 12\tabcolsep) * \real{0.1429}}@{}}
\toprule\noalign{}
\begin{minipage}[b]{\linewidth}\raggedright
Oil-Water Saturation Function and Derivative Table
\end{minipage} & \begin{minipage}[b]{\linewidth}\raggedright
\end{minipage} & \begin{minipage}[b]{\linewidth}\raggedright
\end{minipage} & \begin{minipage}[b]{\linewidth}\raggedright
\end{minipage} & \begin{minipage}[b]{\linewidth}\raggedright
\end{minipage} & \begin{minipage}[b]{\linewidth}\raggedright
\end{minipage} & \begin{minipage}[b]{\linewidth}\raggedright
\end{minipage} \\
\midrule\noalign{}
\endhead
\bottomrule\noalign{}
\endlastfoot
Sw & Krw & Krow & Pcow (psia) &
\(\frac{{\mathit{dk}}_{\mathit{rw}}}{{\mathit{dS}}_{w}}\) &
\(\frac{{\mathit{dk}}_{\mathit{row}}}{{\mathit{dS}}_{w}}\) &
\(\frac{{\mathit{dP}}_{\mathit{cow}}}{{\mathit{dS}}_{w}}\) \\
0.2600 & 0.0000 & 0.8750 & 16.4040 & 0.00 & 0.00 & 0.00 \\
0.2740 & 0.0000 & 0.7737 & 13.4170 & 0.00 & -7.23 & -213.13 \\
0.2889 & 0.0000 & 0.6814 & 11.0910 & 0.00 & -6.20 & -156.36 \\
0.3038 & 0.0000 & 0.5976 & 9.2570 & 0.00 & -5.63 & -123.34 \\
0.3335 & 0.0002 & 0.4535 & 6.6120 & 0.01 & -4.84 & -88.90 \\
0.3633 & 0.0008 & 0.3372 & 4.8610 & 0.02 & -3.91 & -58.87 \\
0.3930 & 0.0022 & 0.2448 & 3.6610 & 0.05 & -3.11 & -40.33 \\
0.4228 & 0.0049 & 0.1728 & 2.8150 & 0.09 & -2.42 & -28.44 \\
0.4525 & 0.0096 & 0.1181 & 2.2030 & 0.16 & -1.84 & -20.55 \\
0.4823 & 0.0171 & 0.0775 & 1.7520 & 0.25 & -1.36 & -15.18 \\
0.5120 & 0.0282 & 0.0484 & 1.4120 & 0.37 & -0.98 & -11.42 \\
0.5418 & 0.0440 & 0.0283 & 1.1520 & 0.53 & -0.67 & -8.74 \\
0.5715 & 0.0656 & 0.0153 & 0.9500 & 0.73 & -0.44 & -6.78 \\
0.6013 & 0.0944 & 0.0074 & 0.7920 & 0.97 & -0.27 & -5.34 \\
0.6310 & 0.1318 & 0.0030 & 0.6650 & 1.26 & -0.15 & -4.25 \\
0.6608 & 0.1793 & 0.0010 & 0.5640 & 1.60 & -0.07 & -3.42 \\
0.6905 & 0.2386 & 0.0002 & 0.4810 & 1.99 & -0.03 & -2.78 \\
0.7203 & 0.3115 & 0.0000 & 0.4130 & 2.45 & -0.01 & -2.28 \\
0.7351 & 0.3537 & 0.0000 & 0.3840 & 2.84 & 0.00 & -1.97 \\
0.7500 & 0.4000 & 0.0000 & 0.3570 & 3.11 & 0.00 & -1.80 \\
1.0000 & 0.9000 & 0.0000 & 0.1270 & 2.00 & 0.00 & -0.92 \\
\end{longtable}

\emph{Table 2.2: Oil-Water Saturation Function and Derivative Table}

If we correct the data shown in Figure 2.21 by extending the water curve
to KRW, then we obtain the result shown in Figure 2.22, but notice what
has happened to the water derivative plot, which now has a discontinuity
because of the extension. Thus, small changes to the relative
permeability curves can have an undue influence on the simulator's
performance. In this particular case, the effect is minor as it is only
one point that has the issue, but it is not uncommon to have multiple
discontinuities in numerous curves that will have an impact on
convergence.

\pandocbounded{\includegraphics[keepaspectratio]{chapters/images/Frame150_bf9e161b7fc2.png}}
\pandocbounded{\includegraphics[keepaspectratio]{chapters/images/Image143_bf9e161b7fc2.png}}

To summarize, if convergence is an issue verify that both the entered
relative permeability data and capillary pressure do not have numerous
discontinuities in the derivative plots.

The next section discusses the impact and relationship of lowest
saturation and critical saturation, in terms of pressure depletion.
However, again, care should be taken to avoid any potential
discontinuities in the derivatives between these to data points.

PROPS Section: Saturation Function (Lowest Saturation and Critical
Saturation)

In most cases, time step chops caused by saturation changes may be
reduced by avoiding having the lowest and critical saturations be the
same. For example for the water phase, this would mean having different
values for SWL and SWCR. The reason for this due to the compressibility
difference between the gas and liquids (oil and water), as illustrated
Figure 2.23. Here we assume that initially 80\% of cell's volume of 100
m3 is occupied by gas and the remaining 20\% by water.

\pandocbounded{\includegraphics[keepaspectratio]{chapters/images/Frame152_09713ad1f09e.png}}
\pandocbounded{\includegraphics[keepaspectratio]{chapters/images/Image145_09713ad1f09e.png}}

Assuming there is no flow in or out of the grid block and the cell
pressure decreases due to production elsewhere in the model, that
results in a small reduction (1.0\%) in the cell's pore volume. And
since gas is much more compressible than oil and water, then the gas
volume will decrease by approximately 1.0\%; whereas the water will
increase slightly to approximately 20.2\% {[}Assuming that the initial
pore volume of the grid block was 100 m3, then a 1.0\% reduction would
result in a pore volume 99 m3, and 20 m3 of water in 99 m3 pore volume
results in 20.2\% of the grid block being occupied by water.{]}. Thus,
if SWL was set to 20\%, then the 0.2\% increase in water saturation
would become mobile and start to flow to the other connect grid blocks.
This can be a cause of instability.

The solution is to enter values for SWL and SWCR, such that there is a
small difference between them, around 2\% to 5\% for example, instead of
this:

\begin{verbatim}
--
--       WATER RELATIVE PERMEABILITY TABLES (SWFN)
--
SWFN
--       SWAT       KRW        PCOW
--       FRAC       FRAC       PSIA
--       --------   --------   -------
          0.30      0.0000     0.6190
          0.35      0.0005     0.3891
          0.40      0.0028     0.2279
          0.50      0.0237     0.0602
          0.60      0.0931     0.0077
          0.80      0.5000     0.0000
          0.90      0.7000     0.0000
          1.00      0.9000     0.0000                      / TABLE NO. 1
\end{verbatim}

Amend the table as follows:

\begin{verbatim}
--
--       WATER RELATIVE PERMEABILITY TABLES (SWFN)
--
SWFN
--       SWAT       KRW        PCOW
--       FRAC       FRAC       PSIA
--       --------   --------   -------
          0.30      0.0000     0.6190
          0.31      0.0000     1*
          0.35      0.0005     0.3891
          0.40      0.0028     0.2279
          0.50      0.0237     0.0602
          0.60      0.0931     0.0077
          0.80      0.5000     0.0000
          0.90      0.7000     0.0000
          1.00      0.9000     0.0000                      / TABLE NO. 1
\end{verbatim}

The incorporation of the row for Sw equal to 0.31 should assist
convergence and will not influence the oil in-place, although it will
slightly reduce the water mobility between the water saturations of 0.30
and 0.31. Secondly, as one can see from Table 2.3 the additional water
saturation value did not cause any derivative discontinuities.

\begin{longtable}[]{@{}
  >{\raggedright\arraybackslash}p{(\linewidth - 12\tabcolsep) * \real{0.1429}}
  >{\raggedright\arraybackslash}p{(\linewidth - 12\tabcolsep) * \real{0.1429}}
  >{\raggedright\arraybackslash}p{(\linewidth - 12\tabcolsep) * \real{0.1429}}
  >{\raggedright\arraybackslash}p{(\linewidth - 12\tabcolsep) * \real{0.1429}}
  >{\raggedright\arraybackslash}p{(\linewidth - 12\tabcolsep) * \real{0.1429}}
  >{\raggedright\arraybackslash}p{(\linewidth - 12\tabcolsep) * \real{0.1429}}
  >{\raggedright\arraybackslash}p{(\linewidth - 12\tabcolsep) * \real{0.1429}}@{}}
\toprule\noalign{}
\begin{minipage}[b]{\linewidth}\raggedright
Oil-Water Saturation Function and Derivative Table
\end{minipage} & \begin{minipage}[b]{\linewidth}\raggedright
\end{minipage} & \begin{minipage}[b]{\linewidth}\raggedright
\end{minipage} & \begin{minipage}[b]{\linewidth}\raggedright
\end{minipage} & \begin{minipage}[b]{\linewidth}\raggedright
\end{minipage} & \begin{minipage}[b]{\linewidth}\raggedright
\end{minipage} & \begin{minipage}[b]{\linewidth}\raggedright
\end{minipage} \\
\midrule\noalign{}
\endhead
\bottomrule\noalign{}
\endlastfoot
Sw & Krw & Krow & Pcow (psia) &
\(\frac{{\mathit{dk}}_{\mathit{rw}}}{{\mathit{dS}}_{w}}\) &
\(\frac{{\mathit{dk}}_{\mathit{row}}}{{\mathit{dS}}_{w}}\) &
\(\frac{{\mathit{dP}}_{\mathit{cow}}}{{\mathit{dS}}_{w}}\) \\
0.30 & 0.0000 & 0.6190 & 9.7247 & & & \\
0.31 & 0.0000 & 0.5918 & 9.1501 & 0.00 & -5.44 & -114.92 \\
0.35 & 0.0005 & 0.3891 & 5.6425 & 0.01 & -4.50 & -77.95 \\
0.40 & 0.0028 & 0.2279 & 3.4623 & 0.05 & -3.22 & -43.60 \\
0.50 & 0.0237 & 0.0602 & 1.5494 & 0.21 & -1.68 & -19.13 \\
0.60 & 0.0931 & 0.0077 & 0.7989 & 0.69 & -0.52 & -7.50 \\
0.80 & 0.5000 & 0.0000 & 0.3110 & 2.03 & -0.04 & -2.44 \\
0.90 & 0.7000 & 0.0000 & 0.2190 & 2.00 & 0.00 & -0.92 \\
1.00 & 0.9000 & 0.0000 & 0.1270 & 2.00 & 0.00 & -0.92 \\
\end{longtable}

\emph{Table 2.3: Oil-Water Saturation Function and Derivative Table
(Modified)}

Note also that same approach can also be applied to the oil phase.

PROPS Section: Saturation Table End-point Scaling

The end-point scaling functionality in OPM Flow gives tremendous
flexibility, allowing one to use a limited set of relative permeability
functions by accounting for rock property variation, such that each cell
in the model has its own relative permeability function. However, both
the original (normalized) input functions and the grid cell's scaled (or
de-noramlized) functions should be checked for consistency. In general,
it is best to input the original curves as standardized curves that have
distinct end-points, rather than fully normalized curves, in the
classical sense. The reason for this is shown in Figure 2.24 and Figure
2.25.

\pandocbounded{\includegraphics[keepaspectratio]{chapters/images/Frame142_ef371648a600.png}}
\pandocbounded{\includegraphics[keepaspectratio]{chapters/images/Image135_ef371648a600.png}}
\pandocbounded{\includegraphics[keepaspectratio]{chapters/images/Frame144_40faed7f4b95.png}}
\pandocbounded{\includegraphics[keepaspectratio]{chapters/images/Image137_40faed7f4b95.png}}

Notice how the distinction between KRO and KRORW is no longer present in
the normalized curve in Figure 2.25, and similarly for KRW and KRWR, and
SWL and SWCR. Thus, the simulator will ignore one of each pair of
end-points - this can lead to unexpected results for the de-normalized
grid cell based saturation functions. An example is shown in Figure 2.26
for a gas-water saturation table. Notice how the de-normalized end-point
water relative permeability, KRW, is only 0.20 versus 1.00 for the
normalized value of KRW.

\pandocbounded{\includegraphics[keepaspectratio]{chapters/images/Frame145_0475708f3372.png}}
\pandocbounded{\includegraphics[keepaspectratio]{chapters/images/Image138_0475708f3372.png}}

Consequently, water influx will be inhibited leading to unreliable
results in terms of water production, and potentially gas production if
aquifers are present in the model.

As mentioned previously, entering a correct original curve that is not
normalized will automatically resolve these type of issues. If a
normalized curve is entered, then ensure that simulator can distinguish
between SWL and SWCR etc., based on the variable definitions provided in
Error: Reference source not found, Error: Reference source not found and
Error: Reference source not found.

Finally, see the SCALCERS keyword in the PROPS section, that defines the
options used in the re-scaling process, including the options for two
point and three point scaling. Whatever option is selected, always
verify that the de-normalized curves that will be used in the model are
correct and reasonable.

PROPS Section: Fluid Properties

It is important that the fluid properties are consistent to avoid
convergence issues, and thus the simulator performs a series of
validation checks to ensure that data is consistent. The following
example indicates the oil formation volume factor is is not
monotonically increasing for increasing GOR.

\begin{verbatim}
--
--       OIL PVT TABLE FOR LIVE OIL
--
PVTO
--       RS        PSAT       BO        VISC
--       MSCF/STB  PSIA       RB/STB    CPOISE
--       --------  --------   -------   ------
...............................................
                                                           /  TABLE NO. 43
         0.000    1.010      1.0388000  120924.000
                  2510.0     0.9801748  529773.445         /
         3.909      10.0     1.0563000   45850.000
                  2510.0     0.9969008  200315.440         /
         5.725      20.0     1.0617000   34198.000
                  2510.0     1.0022360  148947.825         /
         6.787      30.0     1.0648000   28808.000
                  2510.0     1.0054019  125083.718         /
         7.540      40.0     1.0671000   25508.000
                  2510.0     1.0078136  110411.450         /
         8.125      50.0     1.0688000   23210.000
                  2510.0     1.0096595  100151.779         /
         9.940     100.0     1.0742000   17312.000
                  2510.0     1.0159689   73535.291         /
        11.002     150.0     1.0774000   14583.000
                  2510.0     1.0202071   60960.876         /
        11.756     200.0     1.0797000   12912.000
                  2510.0     1.0235993   53105.652         /
        12.340     250.0     1.0814000   11749.000
                  2510.0     1.0264272   47530.728         /
        12.818     300.0     1.0829000   10877.000
                  2510.0     1.0290688   43270.164         /
        13.221     350.0     1.0841000   10191.000
                  2510.0     1.0314284   39854.509         /
        13.571     400.0     1.0851000    9631.000
                  2510.0     1.0336002   37015.563         /
        13.880     450.0     1.0861000    9163.000
                  2510.0     1.0357743   34599.472         /
        14.156     500.0     1.0869000    8763.000
                  2510.0     1.0377596   32498.634         /
        14.405     550.0     1.0877000    8417.000
                  2510.0     1.0397467   30648.326         /
        14.633     600.0     1.0884000    8113.000
                  2510.0     1.0416399   28994.746         /
        14.843     650.0     1.0890000    7843.000
                  2510.0     1.0434389   27501.353         /
        15.037     700.0     1.0896000    7601.000
                  2510.0     1.0452393   26140.639         /
        17.800    2500.0     1.0222000    4436.000
                  2510.0     1.0219701    4495.778         /
                                                           / TABLE NO. 44
\end{verbatim}

The data is displayed above and the offending entries are shown in bold.
These type of PVT warning messages are important and should be addressed
before continuing with the simulation run.

\subsubsection{SOLUTION Section:
Initialization}\label{solution-section-initialization}

The first quality control check that should be performed once the model
has been initialized, is to check that the pressure distribution matches
the observed data, this is especially important when there are multiple
stacked reservoirs in different hydraulic pressures zones, as per the
likes of the oil and gas fields in the Malay Basin and the deep water
fields in Indonesia and the Gulf of Mexico. One effective way to do this
is to plot the initial pressure data versus depth from exploration
wells, and compared this to the equivalent model pressures at the same
depth, as shown in Figure 2.27. Here the measured data is from the RFT
{[}Repeat Formation Tester (RFT) tool is an open hole device which is an
updated version of the Formation Interval Tester (FIT), both of which
are run on wire line. Both tools take multiple pressure readings (at
various depths) thus enabling a pressure depth profile to be obtained
from the formation, and, in addition, they can also take fluid samples
from the formation. The latest tool available is the Modular Formation
Dynamics Tester (MDT), which, as its name suggests, is a modular tool
that can be assembled in different configurations depending on what are
the objectives for running the tool. Thus, the term RFT measurements
applies to all tools that measure a pressure profiles versus depth
(RFT/FIT/MDT etc.).{]} (``Repeat Formation Tester'') tool that measures
the observed pressure profile in the open-hole well bore.

\pandocbounded{\includegraphics[keepaspectratio]{chapters/images/Frame141_27e19fd666cb.png}}
\pandocbounded{\includegraphics[keepaspectratio]{chapters/images/Image134_27e19fd666cb.png}}

Figure 2.27 was created in OPM ResInsight, and shows good agreement
between the two data sets. This type of plot is commonly constructed
when assessing the history matching field performance, but is also
useful in verifying the initial pressure distribution. Note that for
simple models with only a few reservoir zones, then this plot can be
replaced by simply checking the three dimensional pressure distribution.

If the SWATINIT keyword in the PROPS section is used to initialize the
model by setting each grid block's initial water saturation (``Sw''),
then OPM Flow will re-scale the water-oil capillary pressure curves
entered via the SWFN saturation functions in the PROPS section, so that
the resulting initialized Sw matches the values in the SWATINIT array.

Normally the SWATINIT array is generated in the static earth model when
calculating the hydrocarbons in-place volumes using a Saturation Height
Functions (``SHF'') derived from capillary pressure functions. Static
earth models do not directly use capillary pressure in these type of
calculations as individual cell pressures are not available. There is
therefore some potential for inconsistencies to arise between the two
sets of formulations. This is normally manifested by extreme scaling in
the scaled capillary pressure values calculated by the simulator. These
large scaled capillary pressures values can result in numerical issues,
and should be addressed, by verifying that the SHF used in the earth
model to derive the SWATINIT array and the unscaled capillary pressure
values on the SWFN keyword are consistent. Note in general, the scaled
capillary pressure values calculated by the simulator, should be roughly
in the same range as the entered unscaled data, as shown in Figure 2.28
for the Norne field.

\pandocbounded{\includegraphics[keepaspectratio]{chapters/images/Frame143_1e27793ec858.png}}
\pandocbounded{\includegraphics[keepaspectratio]{chapters/images/Image136_1e27793ec858.png}}

Note that although the PPCWMAX keyword can be used to set a maximum
scaled capillary pressure value, a more technically sound approach is to
resolve these inconsistencies before continuing with the model build.

Once a model is built and can be initialized, it is good practice to run
the model with no production for several years to ensure that the
initialization is stable. Under these circumstances, there should be no
fluid movement, or only very limited movement, and the pressure should
remain stable, that is constant, although a small drift may be
observable. An example for the Norne model is shown in Figure 2.29 on
the field level, although it can, and should be performed on a region by
region basis if the model is not stable.

\pandocbounded{\includegraphics[keepaspectratio]{chapters/images/Frame137_d09fef45212a.png}}
\pandocbounded{\includegraphics[keepaspectratio]{chapters/images/Image130_d09fef45212a.png}}

Potential causes for the instability are:

\begin{verbatim}
- Aquifer Not in Equilibrium with the Connected Hydrocarbon Zones. One can check for this by switching off all the aquifer connections and see if the instability is still present. If this corrects the issue then verify the aquifer connections by reviewing the AQUIFERA array in the *.INIT file in OPM ResInsight. OPM Flow writes out the AQUIFERA value as 2(AQUID-1) for cells connected to aquifer AQUID. If a cell is connected to multiple analytical aquifers then AQUIFERA is summed for all aquifers connected to a cell. Note that connecting cells to multiple aquifers should be avoided.
\end{verbatim}

Note aquifers should only be connected to grid blocks containing 100\%
water, or at the very least any hydrocarbons in the connected cells
should be immobile.

\begin{verbatim}
- Different PVT Regions in Communication. Recall that fluid proprieties are associated with a grid cell, so if two different PVT regions are in communication, then a m3 of fluid from PVT region one flowing into PVT region two, will automatically assume the properties of PVT region two. This creates a material balance error that will result in numerical difficulties.  One may be able to identify this by reviewing the inter-region flows to identify the source of the problem, either via the *.PRT file or the summary file out. It typically occurs across fault planes, when one PVT region butts against a different PVT region along the fault face.
\end{verbatim}

There are several solutions to this, firstly, one can set the
transmissibility along the fault plane to zero using the MULTFLT keyword
in the GRID or EDIT sections.

Secondly, create a MULTNUM array based on the PVTNUM array, as the
PVTNUM keyword is not available in the GRID and EDIT sections. Then use
the MULTREGT keyword with MULTNUM array in the GRID or EDIT sections to
set the transmissibility between the offending regions to zero.

Finally, one can use the THPRES keyword in the SOLUTION section combined
with the EQLNUM array to isolate the offending regions. Normally, but
not always, the PVTNUM and EQLNUM arrays are identical, if not, then it
is relatively easy to copy the aforementioned MULTNUM array to the
EQLNUM array, then add the additional equilibration data to the EQUIL
keyword. One advantage of this option is that it will cater for the
scenario when the PVT regions are initially isolated, but after a
certain production period, the faults break down and there is
communication between the two regions. In this scenario, one still has
the material balance error as different fluids inter-mix. If this is the
case then the API Tracking option should be used to account for this
behavior. See the API keyword in the RUNSPEC on how to activate this
option.

\begin{verbatim}
- Different Equilibration Regions in Communication. This is similar to (2) above and the remedies are also similar, so the solution is not repeated here.


- PVT Data Not Consistent with Equilibration Data. In an oil-gas system the saturation pressure, that is the bubble-point pressure should be the pressure at the GOC, and similarly for the dew point. Conveniently, OPM ResInsight can produce a plot, as depicted in Figure 2.30. The figure illustrates that gas dew point pressure gradient (shown in red), and the oil bubble-point pressure gradient (shown in green) from the PVT data intersect at the GOC, which matches the initial pressure profile (shown in blue) from the grid cells. Thus, confirming the consistency of the two data sets.
\end{verbatim}

\pandocbounded{\includegraphics[keepaspectratio]{chapters/images/Frame138_f190d694e1a5.png}}
\pandocbounded{\includegraphics[keepaspectratio]{chapters/images/Image131_f190d694e1a5.png}}

\subsubsection{SCHEDULE Section}\label{schedule-section}

A common issue in most simulation models that employ Vertical Lift
Performance (``VLP'') tables to relate a well's Tubing Head Pressure
(``THP'') to its Bottom-Hole Pressure (``BHP''), is the consistency of
the data. Note that this is not just limited to wells, but also the
network pipelines, that describe the pressure drop between the inlet
node and the outlet node.

\begin{longtable}[]{@{}
  >{\raggedright\arraybackslash}p{(\linewidth - 0\tabcolsep) * \real{1.0000}}@{}}
\toprule\noalign{}
\begin{minipage}[b]{\linewidth}\raggedright
Warning: VFPPROD table 39 has 5 non-monotonic points of BHP(THP) In
/media/OPM/Fields/Norne/Model/INCLUDE/VFP/B3H.Ecl line 68 This may cause
convergence issues due to switching between BHP and THP control.
\end{minipage} \\
\midrule\noalign{}
\endhead
\bottomrule\noalign{}
\endlastfoot
\end{longtable}

Generally, the cause of this warning is because the VFP tables have been
generated with too much fidelity or resolution. Users often consider
that having more detailed and higher resolution VFP tables will give
more accurate results. This is not always the case and in this
particular example, may very well end up causing the wells using this
particular VFP table to oscillate between control modes. Figure 2.31
shows the VFP header section of the offending table, that declares the
values used to generate the table.

\pandocbounded{\includegraphics[keepaspectratio]{chapters/images/Frame139_6302045ea559.png}}
\pandocbounded{\includegraphics[keepaspectratio]{chapters/images/Image132_6302045ea559.png}}

At first glance this looks reasonable, but the problem is that the
variables (LIQ, THP, WCT and GOR) are at a too high resolution, which
causes some of the resulting curves to cross, due to the numerical
accuracy when generating the curves. In addition, note that the maximum
water cut is defined as 100\%, which is unlikely to occur in practice,
perhaps 98\% would be more appropriate as the maximum value. OPM
ResInsight is able to load VFPPROD and VFPINJ tables and to plot the
data. Here, Figure 2.32 shows a plot of BHP versus THP with the various
water cut values represented by the individual curves. The figure
clearly illustrates the crossing of the water cut line of 80\% over the
70\% water cut curve, as well as the issue with the 100\% water cut
curve.

\pandocbounded{\includegraphics[keepaspectratio]{chapters/images/Frame140_9ca57b60112c.png}}
\pandocbounded{\includegraphics[keepaspectratio]{chapters/images/Image133_9ca57b60112c.png}}

The solution is to re-generate the offending VFP tables using a coarser
set of variable values. The Petroleum Experts' PROSPER {[}IPM Prosper
Version 16.5 User Manual - November 2020, Petroleum Experts, Edinburgh,
Scotland (www.petex.com).{]} program, as well as other nodal analysis
software programs used to generate VFP tables for simulators, have
options to generate the parameters based on the number of entries, and
the minimum and maximum values for a variable, combined with various
interval spacing options. Keeping the same range and selecting four or
five entries and using the geometric spacing option normally resolves
this type of issue. The simulator will then interpolate the coarser
curves and avoid any issues with crossing curves. Note that it is
important to ensure that the variable ranges cover the full range of
values expected to occur during the simulation, as extrapolation, unlike
interpolation of the curves, is more prone to be problematic.

Note that this warning message can be safely ignored:

\begin{longtable}[]{@{}
  >{\raggedright\arraybackslash}p{(\linewidth - 0\tabcolsep) * \real{1.0000}}@{}}
\toprule\noalign{}
\begin{minipage}[b]{\linewidth}\raggedright
Warning: LIQ preferred phase not supported for well OP-01, using OIL
instead
\end{minipage} \\
\midrule\noalign{}
\endhead
\bottomrule\noalign{}
\endlastfoot
\end{longtable}

as it refers to the reference phase in a well's productivity index
calculation. This is because OPM Flow only currently supports OIL, GAS
and WATER well type fluids, as defined by the WELSPECS(TYPE) parameter
in the SCHEDULE section, i.e., liquid (LIQ) is not supported. For
producing, wells this mostly matters if one plots the WPI summary vector
(productivity index for well's preferred phase). In the current
treatment WPI will not have contributions from the water phase if the
declared preferred phase is LIQ. Thus, if history matching a well's
actual productivity index, this should be done on one of the recognized
fluid well types. For injection wells the preferred phase defined by the
WELSPECS keyword does not matter, since the preferred phase is
(typically) reset to the injected phase via the WCONINJE or WCONINJH
keywords.

\subsection{Convergence and Numerical
Performance}\label{convergence-and-numerical-performance}

Given that all the potential issues with the input data have been
successfully resolved, there may still be problems with convergence and
numerical performance, due to the nature of the problem the simulator is
trying to solve. In Figure 2.33, at first glance, the non-linear
iterations would appear to be reasonable, ranging from one to four,
combined with linear iterations averaging around ten per time step.
However, the time steps are very small resulting in nearly 12,500 Newton
iterations and over 60,000 linear iterations in the commercial
simulator, for this relatively simple problem. In fact, at the time of
the study {[}Christie, M. A., and M. J. Blunt. ``Tenth SPE Comparative
Solution Project: A Comparison of Upscaling Techniques.'' SPE Res Eval
\& Eng 4 (2001): 308--317. doi: https://doi.org/10.2118/72469-PA.{]},
the commercial simulator was unable to complete this simulation.
Re-running the model with the latest version of the commercial simulator
(Figure 2.33), still gives incorrect results compared to the published
data.

\pandocbounded{\includegraphics[keepaspectratio]{chapters/images/Frame132_45cd1251bc35.png}}
\pandocbounded{\includegraphics[keepaspectratio]{chapters/images/Image125_45cd1251bc35.png}}

Fortunately, there are several ways to remedy poor performance,
including:

\begin{verbatim}
- The simplest and common approach is to manually restrict the time step size to reduce the number of time steps chops. Time step chops occurs when the simulator fails to converge at a given time step (tn+1), and therefore restarts at the last successful time step (tn) and advances to the next time step using a smaller time step size (∆tn+1) than before (∆tn.). This was the approach used for the example in Figure 2.33, which although the model now runs, the results are incorrect.
- Alternatively, one can change the numerical controls for the Newton-Raphson (“Newton” or non-linear solver) iteration, by changing the minimum and/or maximum number of Newton iterations. Although in general this may resolve some time step chops it can also result in a significant increase in run times, and should therefore be used with caution.
- The last option is to change the numerical parameters for the linear solver, which is generally preferable to modifying the Newton iteration parameters. OPM Flow contains various options in terms of the type of linear solver and the linear solver preconditioner to be employed, in addition to various standard linear solver parameters employed by the commercial simulator (minimum,  maximum number of iterations etc.).
\end{verbatim}

For example, Figure 2.34 illustrates the impact of selecting the more
appropriate numerical controls. Here, the number of Newton iterations
are greatly reduced to around 2,500 at the expense of increasing the
number of linear iterations from \textasciitilde60,000 to over 100,000.
More importantly, the simulator results are correct, in that they are
close to the published oil, gas, and pressure values23.

\pandocbounded{\includegraphics[keepaspectratio]{chapters/images/Frame133_f21c643a8896.png}}
\pandocbounded{\includegraphics[keepaspectratio]{chapters/images/Image126_f21c643a8896.png}}

Note that the simulator has various other numerical performance
parameters than can use used to identify the cause for under
performance.

\begin{tcolorbox}[enhanced jigsaw, title=\textcolor{quarto-callout-note-color}{\faInfo}\hspace{0.5em}{Note}, coltitle=black, breakable, toprule=.15mm, colbacktitle=quarto-callout-note-color!10!white, leftrule=.75mm, opacitybacktitle=0.6, colback=white, bottomtitle=1mm, toptitle=1mm, titlerule=0mm, left=2mm, colframe=quarto-callout-note-color-frame, arc=.35mm, rightrule=.15mm, bottomrule=.15mm, opacityback=0]

To obtain the numerical performance data, the easiest way is to add the
PERFORMA keyword in the SUMMARY section, that enables the writing out of
a series of performance keywords to the summary file, which can then be
plotted in OPM ResInsight as per Figure 2.33 and Figure 2.34.

\end{tcolorbox}

In addition to the aforementioned PERFORMA keyword in the SUMMARY
section, which is available for both OPM Flow and the commercial
simulator, OPM Flow has two additional numerical performance reports
that may help in identifying numerical issues. Both of these OPM Flow
specific reports tabulate the numerical performance data available in
the print file, in an easily readable format.

The first is the INFOSTEP report illustrated in Figure 2.35, which
summarizes the numerical performance for each time step. The Assembly,
LSetup, LSolve, LocSol, Update and Output header columns represent the
time in seconds taken to complete the particular work item.

\pandocbounded{\includegraphics[keepaspectratio]{chapters/images/Frame153_6c7d73102f69.png}}
\pandocbounded{\includegraphics[keepaspectratio]{chapters/images/Image144_6c7d73102f69.png}}

Here, the Assembly column represents the time taken to assemble the
equations, LSetup is the time taken to setup the linear solver
equations, LSolve is the linear solver time, LocSol is the local solver
time, Update is the time taken in updating the results, and Output is
the time used for writing out the results, that is the output to the
print, debug, summary, and restart files.

The WellIt, Lins, NewtIt, and LinIt column headers represent the number
of iterations used to complete the time step for a given work item. Thus
for time step one, the number of well iterations was zero (WellIt), the
number of times the linear equations were setup was 16 (Lins), the
number of Newton iterations used to complete the time step was 15
(NewIt), and the total number of linear iterations to completed the time
step was 19 (LinIt). The last column, Conv, indicates if the time step
had successfully converged given the current convergence criteria, with
a one indicating success and a zero indicating a failure.

By default the report is not generated to maintain compatibility with
commercial simulator. To generate the report one can use the
--output-extra-convergence-info command line parameter, for instance:

\begin{verbatim}
flow --output-extra-convergence-info="steps" CASENAME
\end{verbatim}

will write the report to the CASENAME.INFOSTEP file in the same
directory as the other output files.

The second report, INFOITER, contains more detailed information by
reporting the numerical performance per iteration, as can be seen in
Figure 2.36.

\pandocbounded{\includegraphics[keepaspectratio]{chapters/images/Frame154_980aec28e2bf.png}}
\pandocbounded{\includegraphics[keepaspectratio]{chapters/images/Image146_980aec28e2bf.png}}

Here, the Iteration column contains the number of Newton iterations,
same as the NewIt column in the INFOSTEP report. Note that adding one to
the values in the Iteration column gives the number of times the linear
equations were setup. Columns starting CnvPvFrac give the fractions of
the pore volume that have converged under the standard convergence
criteria, have converged under the relaxed convergence criteria, and
have not converged. Similarly, columns starting CnvCellCnt give the
number of cells that have converged under the standard convergence
criteria, have converged under the relaxed convergence criteria, and
have not converged. Columns starting Penalty indicate the penalty values
for non-convergence, distance decay, and large well residuals. Columns
containing MB indicate a phase's normalized material balance error, and
columns containing CNV show a phase's normalized residual error. The
last column, WelllStatus, indicates if the iteration converged (CONV) or
failed (FAIL).

Like the INFOSTEP report, the INFOITER report by default is not
generated in order to maintain compatibility with commercial simulator.
To generate the report one can use the --output-extra-convergence-info
command line parameter, for instance:

\begin{verbatim}
flow --output-extra-convergence-info="iterations" CASENAME
\end{verbatim}

will write the report to the CASENAME.INFOITER file in the same
directory as the other output files. Finally to generate both reports
use the following command line option.

\begin{verbatim}
flow –output-extra-convergence-info="steps,iterations" CASENAME
\end{verbatim}

The following sections now attempt to give some guidance on how to use
the numerical controls to resolve the simulator's performance, assuming
that all the issues with the input data have been successfully resolved.

\subsection{Time Step Control
Guidance}\label{time-step-control-guidance}

Figure 2.37 illustrates how the simulator advances from the start of the
simulation to the end, and shows two user defined report steps, each
containing two simulator selected time steps. Each time step within a
report time step contains one to three Newton iterations, and each
Newton iteration contains a various number of linear iterations.

\pandocbounded{\includegraphics[keepaspectratio]{chapters/images/Frame129_69393748d817.png}}
\pandocbounded{\includegraphics[keepaspectratio]{chapters/images/Image47_69393748d817.png}}

In an actual model the number of time steps are determined by the
simulator's internal adaptive time stepping algorithm, combined with any
user settings overriding the internal stepping algorithm and parameters.
The default initial time step of one day can be overwritten via either
the command line parameter:

\begin{verbatim}
--initial-time-step-in-days
\end{verbatim}

or via the Error: Reference source not found keyword in the SCHEDULE
section, that offers more flexibility, as the initial time step size can
be reset at various report time steps, by repeating the keyword and
setting TUNING(TSINIT) to a given value. Alternatively, TUNING(TSMINZ)
may be used to define the minimum length of all time steps. Note,
however, the simulator only reads some of the parameters from the TUNING
keyword if, and only if, the following command line parameter has been
set equal to true:

\begin{verbatim}
--enable-tuning=true
\end{verbatim}

One can also set the maximum time step size via both the command line
parameter:

\begin{verbatim}
--solver-max-time-step-in-days
\end{verbatim}

or via the TUNING(TSMAXZ) parameter, that defines the maximum length of
the next time step following TUNING(TSINIT).

There are many other time stepping command line parameters, including
those that are equivalent to all the first record parameters on the
TUNING keyword; however, provided there are no data inconsistencies
within the input deck, resetting the maximum time step size from the
default value of 365 days may improve numerical performance. Typically
the default value of 365 days is not binding, as the report time steps
are more restrictive. So for example, in a history matching model, the
production data is normally entered on a monthly basis, so the maximum
report time step size would be 31 days. Even so, time step chops may
still occur even in this scenario, and therefore reducing the maximum
time step to half a month, that is 16 days may avoid time step chops.
Sixteen days is preferred over 15 days as the maximum time step, because
it enables only two successful time steps to complete a 31 day month (a
16 day plus a 15 day time step). If 15 days was used for the maximum
time step size, then one would need three successful time steps to
complete a 31 day month. In some cases using 11, or even a five day
maximum time step may increase computational efficiency by avoiding
expensive computational time step chops.

Basically, one has to balance the maximum length of time steps with the
required number of time steps to reach the end of the run, versus the
total elapsed time {[}Elapsed real time, real time, wall-clock time, or
wall time is the actual time taken from the start of a simulation to the
end. In other words, it is the difference between the time at which the
simulation finishes and the time at which the simulation started.{]}.
Time steps chops are computationally expensive and thus should therefore
be minimized in order to reduce run times.

Figure 2.38 shows a typical report step and time step report from the
simulator that contains two time step chops shown in the red outline
boxes. At the start of the report:

\begin{verbatim}
- The simulation has advanced to 00:00:00 on October 9th, 1967 at report step 15 (Current Step) out of 33 steps (End Step) and after an elasped time of 250 days out of a total of 2000 days needed to complete the run.
- As requested, the simulator writes out a restart record at report step 15, that is October 9th, 1967.
- The simulator attempts to advance the simulation from October 9th, 1967 by taking a time step of 50 days, as determined by the internal adaptive time stepping algorithm. If successful, this would advance the simulator from the current report step 15 at an elapsed time of 250 days to the next report step 16 at an elapsed time of 300 days.
- However, the time step fails due to a Numerical problem encountered, that is either the non-linear or linear solver failed to convergence within the maximum number of iterations.  Thus, the time step was chopped from 50 days to 16.5 days.
- The simulator restarts again from report step 15, that is October 9th, 1967 and attempts to advance 16.5 days to October 25th, 1967. If successful, this would advance the simulator by 16.5 days from the current elapsed time of 250 days towards the next report step at 300 days.
- The 16.5 day time step converges in three Newton iterations using a total of 342 linear iterations, and the simulator has successfully advanced to October 25th, 1967.
- The simulator then attempts to advance 33.5 days, but again fails due to a Numerical problem encountered, that is either the non-linear or the linear solver again failed to convergence within the maximum number of iterations.  Thus, the time step was chopped form 33.5 days to 11.055 days.
- The simulator restarts again from the last successful time step, that is October 25th, 1967 and attempts to advance to November 5th, 1967. The 11.055 day time step converges in three Newton iterations using a total of 317 linear iteration, and the simulator has advanced to November 5th, 1967.
- The simulator's internal adaptive time stepping algorithm has determined that the next time step should be 22.445 days to coincide with the next report time step (report step 16) at November 28th, 1967. Again, the 22.445 day time step is successful and converges in three Newton iteration using a total of 318 linear iterations.
- The simulation has now advanced to 00:00:00 on November 28th, 1967 at report step 16 (Current Step) out of 33 steps (End Step) and after an elasped time of 300 days out of a total of 2000 days needed to complete the run.
\end{verbatim}

\pandocbounded{\includegraphics[keepaspectratio]{chapters/images/Frame130_8707dd32f779.png}}
\pandocbounded{\includegraphics[keepaspectratio]{chapters/images/Image63_8707dd32f779.png}}

If the maximum time step was set to 16 days then advancing 50 days to
November 28th, 1967 would take four time steps with a good chance of
avoiding time step chops and thus reducing run times.

Various command line parameters can be used to control the behavior of
the time stepping algorithm, as summarized in the Table 2.4, that lists
the important time stepping control parameters for the simulator,
including the TUNING equivalent parameters where applicable.

\begin{longtable}[]{@{}
  >{\raggedright\arraybackslash}p{(\linewidth - 6\tabcolsep) * \real{0.2500}}
  >{\raggedright\arraybackslash}p{(\linewidth - 6\tabcolsep) * \real{0.2500}}
  >{\raggedright\arraybackslash}p{(\linewidth - 6\tabcolsep) * \real{0.2500}}
  >{\raggedright\arraybackslash}p{(\linewidth - 6\tabcolsep) * \real{0.2500}}@{}}
\toprule\noalign{}
\begin{minipage}[b]{\linewidth}\raggedright
Major Solver Time Stepping Controls
\end{minipage} & \begin{minipage}[b]{\linewidth}\raggedright
\end{minipage} & \begin{minipage}[b]{\linewidth}\raggedright
\end{minipage} & \begin{minipage}[b]{\linewidth}\raggedright
\end{minipage} \\
\midrule\noalign{}
\endhead
\bottomrule\noalign{}
\endlastfoot
No. & Command Line Parameter & TUNING Keyword Equivalent &
Description \\
1 & --initial-time-step-in-days & TSINIT & A real positive value. Set
the size of the initial time step in days. The default value is one day.
Note that as the TUNING keyword can be entered multiple times, then
TUNING(TSINIT) can be reset multiple times, whereas as
--initial-time-step-in-days, is only set once at the beginning of the
run. \\
2 & --solver-growth-factor & TSDIFF & A real positive value. Set the
growth factor a time step can be increased by when recovering from one
or more time step chops, subject to the maximum allowable time step size
set by the --solver-max-time-step-in-days parameter. For example, if the
current time step has converged at 5 days after a time step chop, and
--solver-growth-factor is set to the default value of 2.0, then the next
time step will be 2.0 * 5 days, that is at 10 days. \\
3 & --solver-max-growth & TSFMAX & A real positive value. Set the
maximum growth factor a time step can be increased by after a report
time step, subject to the maximum allowable time step size set by the
--solver-max-time-step-in-days parameter. Thus, if the current time step
has converged at 5 days after at a report step, and --solver-max-growth
is set to the default value of 3.0, then the next time step will be 3.0
* 5 days, that is at 15 days. \\
4 & --solver-max-restarts & & A positive integer. Set the number of
allowed consecutive restarts (or time step chops) before the simulation
is terminated. Here the default value is 10. TUNING (TSMCHP) sets the
minimum length of all chopped time steps. \\
5 & --solver-max-time-step-in-days & TSMAXZ & A real positive value. Set
the maximum allowed time step size in days. The default value of 365
days is rarely used in practise due to the report time steps being more
restrictive, although in some specialized models (CO2 long term
prediction cases) this may indeed be the upper limit. As mentioned
previously, if the time step chops are mostly occuring when the time
step size is greater than ∆t, then setting this parameter to a value
less than ∆t may reduce the number of time step chops. Provided of cause
the value of ∆t is reasonable. \\
6 & --solver-min-time-step & TSMINZ & A real positive value. Set the
minimum size of a time step in days for field and metric units, and
hours for laboratory units If a time step cannot converge without
getting cut below this time step size the simulator will stop, unless
overwritten by the --solver-continue-on-convergence-failure command line
parameter (see Table 2.1 for further details). Note if the current time
step is smaller than this minimum and has not been subject to a time
step chop, then in some circumstances it may still be accepted due to
other numeric controls. \\
7 & --solver-restart-factor & TSFCNV & A real positive value. Set the
time step chop factor of the time step after a convergence failure. For
example, if the current non-convergent time step is 30 days and this
parameter is set to the default value of 0.33, then the time step will
be repeated using 0.33 * 30 days as the time step, that is 9.9 days. \\
Notes: & & & \\
\end{longtable}

\emph{Table 2.4: Major Solver Time Stepping Controls}

\begin{tcolorbox}[enhanced jigsaw, title=\textcolor{quarto-callout-note-color}{\faInfo}\hspace{0.5em}{Note}, coltitle=black, breakable, toprule=.15mm, colbacktitle=quarto-callout-note-color!10!white, leftrule=.75mm, opacitybacktitle=0.6, colback=white, bottomtitle=1mm, toptitle=1mm, titlerule=0mm, left=2mm, colframe=quarto-callout-note-color-frame, arc=.35mm, rightrule=.15mm, bottomrule=.15mm, opacityback=0]

OPM Flow also has a series of command line parameters of the form
--time-step* that are subordinate to the --solver* series of command
line parameters. The former, --time-step*, are used to select the time
step control algorithm and associated parameters, and in general should
only be used by users with appropriate experience.

\end{tcolorbox}

As mentioned above, the --time-step* series of command line parameters
should in general not be used; however, the following parameter may be
useful if group and well events are resulting in time step chops:

\begin{verbatim}
--time-step-after-event-in-days
\end{verbatim}

If the simulation is running smoothly with monthly historical report
time steps, except when a new well is bought on stream in a given month,
or a group control is reset resulting in a significant change in
production or injection, and the event causes time step chops, then
using --time-step-after-event-in-days may improve the situation. The
parameter can be used to automatically enforce a small time step when an
event occurs, thus reducing the probability, or even eliminating a time
step chop. Typically, a one or five day time step is used, as per:

\begin{verbatim}
--time-step-after-event-in-days=5.0
\end{verbatim}

The drawback is that the control is applied to the complete run, which
again needs to be taken into consideration.

Alternatively, one could use the NEXT or NEXTSTEP keywords in the
SCHEDULE section, to incorporate similar behavior at specific times in
the SCHEDULE section. See Error: Reference source not found or Error:
Reference source not found for further information.

\subsection{Newton-Raphson Solver Iteration Controls (Non-Linear
Iterations)}\label{newton-raphson-solver-iteration-controls-non-linear-iterations}

Figure 2.39 shows typical time step output for two report time steps,
March 1st and April 1st, 2022; with both solving the equations with two
Newton iterations per time step, and between four to six linear
iterations per time step. The numerical count for the number of
linearizations is the number of times the Jacobian and the residuals
have been constructed, and is generally one more than the number of
Newton iterations. Figure 2.39, thus indicates very good performance,
especially as the first iteration is taken from the previous successful
time step. Note however, that the model is only using five day time
steps in this case, which significantly contributes to the convergence
stability of both the non-linear and linear solvers.

\pandocbounded{\includegraphics[keepaspectratio]{chapters/images/Frame127_65f003e88fab.png}}
\pandocbounded{\includegraphics[keepaspectratio]{chapters/images/Image36_65f003e88fab.png}}

The number of Newton iterations taken to solve a time step is indicative
of the degree of ease the simulator has in converging to a solution,
with higher numbers of iterations indicating an increasingly difficult
problem. Note, unlike Figure 2.39, the number of non-linear iterations
will vary as the simulation progresses, with periods of relatively
efficient convergence and periods where the simulator is struggling to
reach non-linear convergence. The latter is normally associated with
operational events, for example new wells being brought on stream, major
changes to group controls etc.

\begin{tcolorbox}[enhanced jigsaw, title=\textcolor{quarto-callout-note-color}{\faInfo}\hspace{0.5em}{Note}, coltitle=black, breakable, toprule=.15mm, colbacktitle=quarto-callout-note-color!10!white, leftrule=.75mm, opacitybacktitle=0.6, colback=white, bottomtitle=1mm, toptitle=1mm, titlerule=0mm, left=2mm, colframe=quarto-callout-note-color-frame, arc=.35mm, rightrule=.15mm, bottomrule=.15mm, opacityback=0]

As mentioned previously, in general the number of Newton iterations is
indicative of the ease of convergence, with one to three iterations
indicating that converge was relatively easy and efficient, four to nine
iterations indicating an increasingly difficult problem, and greater
than ten would tend to suggest that there may be an issue with the
model.

\end{tcolorbox}

It is important to stress that the convergence of the non-linear solver
is very dependent on the accuracy of the linear solver. If the accuracy
of the linear solver is too relaxed or is not fully converged, the
non-linear solver may fail to converge. For example, if the linear
solver solution is not sufficiently converged, then the non-linear
solution may no longer be well-posed {[}A problem in differential
equations is said to be well-posed if: (1) a solution exists; (2) that
the solution is unique; (3) and the solution changes continuously with
the changes in the data.{]}. That is the solution starts digressing from
the true solution, and additional iterations only exasperate the
problem. Normally, under these circumstances OPM Flow will chop the time
step, which should correct the divergence by reverting to previous
converged time step. However, if the time step is too small to chop and
the run continues with the minimum time step, then the divergence from
the true solution will continue. This may manifest itself with various
error messages and NaNs (``Not a Number''), or eventually with the
simulator aborting. Figure 2.40 shows an example of this behavior.

\pandocbounded{\includegraphics[keepaspectratio]{chapters/images/Frame134_16c3503b2239.png}}
\pandocbounded{\includegraphics[keepaspectratio]{chapters/images/Image127_16c3503b2239.png}}

Fortunately, OPM Flow has various controls that can resolve these type
of issues similar to the commercial simulator, and in addition, and
unlike the commercial simulator, OPM Flow has several linear solvers and
preconditioners that increase the likelihood a solution may be found for
various non-convergence issues.

As mentioned previously in section Error: Reference source not found
Error: Reference source not found, convergence is achieved when
\(R({y}_{n})\), the residual, in equation (2.5), satisfies a given
tolerance.

\begin{equation}\phantomsection\label{eq-2-5}{
J({y}_{n})\left({y}_{n+1} - {y}_{n}\right) = -R({y}_{n})
}\end{equation}

where:

\(J ({y}_{n})\) = the Jacobian matrix,

\({y}_{n}\) = the solution of the non-linear equations at the nth
non-linear

iteration, and

\({y}_{n+1}\) = the solution of the non-linear equations iteration at
the n+1

iteration.

However, as the iterations proceed during a time step the residual
convergence criteria is reduced in an effort to avoid a time step chop.
Thus, by relaxing the convergence criteria as the non-linear iterations
progress, the simulator is attempting to avoid throwing away all the
previous work performed in trying to reach a solution for the time step.

From the user prospective, there are several command-line options that
can be used to control the Newton-Raphson solver, in particular:

\begin{verbatim}
- To modify the maximum number of non-linear iterations per time step, the following command line parameter can be used:
\end{verbatim}

\begin{verbatim}
--newton-max-iterations
\end{verbatim}

Although, the default value of 20 is sufficiently large that it should
not be necessary to increase this value.

\begin{verbatim}
- There is also a command that sets the minimum number of Newton iterations per time step used by the simulator.
\end{verbatim}

\begin{verbatim}
--newton-min-iterations
\end{verbatim}

The default value of two ensures that at least two Newton iterations are
performed after the previous time step. Increasing this parameter, may
help convergence in difficult time steps, by ensuring slightly more
accurate results in preceding ``easy'' time steps.

\begin{verbatim}
- Thirdly, the mass balance error as a reservoir average saturation error, can be use to slackened or tighten the material balance tolerance.
\end{verbatim}

\begin{verbatim}
--tolerance-mb
\end{verbatim}

In general, if the non-linear solver is having difficulty in converging,
then tightening the linear mass balance error as a reservoir average
saturation error, may help the non-linear solver converge more rapidly.

\begin{verbatim}
- Fourthly,  maximal local residual error can be slackened or tighten using:
\end{verbatim}

\begin{verbatim}
--tolerance-cnv
\end{verbatim}

command line parameter. And may be used to set the maximal local
residual error.

For completeness, Table 2.5 lists the major non-linear command line
parameters used to control the behavior of the non-linear solver,
including those that relate to relaxed convergence criteria. In general,
it is recommended not to modify these parameters, unless one has the
sufficient expertise in this area.

\begin{longtable}[]{@{}
  >{\raggedright\arraybackslash}p{(\linewidth - 6\tabcolsep) * \real{0.2500}}
  >{\raggedright\arraybackslash}p{(\linewidth - 6\tabcolsep) * \real{0.2500}}
  >{\raggedright\arraybackslash}p{(\linewidth - 6\tabcolsep) * \real{0.2500}}
  >{\raggedright\arraybackslash}p{(\linewidth - 6\tabcolsep) * \real{0.2500}}@{}}
\toprule\noalign{}
\begin{minipage}[b]{\linewidth}\raggedright
Major Newton-Raphson Solver Controls
\end{minipage} & \begin{minipage}[b]{\linewidth}\raggedright
\end{minipage} & \begin{minipage}[b]{\linewidth}\raggedright
\end{minipage} & \begin{minipage}[b]{\linewidth}\raggedright
\end{minipage} \\
\midrule\noalign{}
\endhead
\bottomrule\noalign{}
\endlastfoot
No. & Command Line Parameter & TUNING Keyword Equivalent &
Description \\
1 & --newton-max-iterations & NEWTMX & A positive integer. Set the
maximum number of Newton iterations per time step used by the simulator.
The default value is 20 which compares to the commercial simulator's
value of 12, so there should not really be a requirement to increase
this parameter. \\
2 & --newton-min-iterations & NEWTMN & A positive integer. Set the
minimum number of Newton iterations per time step used by the simulator.
The default values for --newton-min-iterations and NEWTMN are 2 and 1,
respectively. The value of one ensures that at least one Newton
iteration is performed after the previous time step. Increasing this
parameter, to say three for example, may help convergence in difficult
time steps, by ensuring slightly more accurate results in preceding
``easy'' time steps. \\
3 & --newton-max-relax & & A real value ranging from zero to one. Set
the maximum relaxation factor of a Newton iteration used by the
simulator. The default is 0.5. In general, it is not recommended to
modify this parameter. \\
4 & --newton-relaxation-type & & A character string. The type of
relaxation used by the Newton method. Valid options are: dampen
(default) or sor. In general, it is not recommended to modify this
parameter. \\
5 & --max-newton-iterations-with-inner-welliterations & MXWSIT & A
positive integer. Set the maximum Newton iterations with inner well
iterations. The default value is eight, the same as the commercial
simulator. If there are issues with well convergence, then sometimes
increasing this value, to say 12, may improve the situation. \\
6 & --min-strict-cnv-iter & & An integer. Set the minimum number of
Newton iterations before relaxed tolerances are used for the CNV
convergence criterion. The default value -1 means that relaxed
convergence criterion is applied after the initial iteration, that is
iteration zero in Figure 2.42 in section 2.4.5 Linear Solver Controls
and Options. Increasing this parameter may improve convergence by
applying stricter convergence controls over additional Newton
iterations, but may also increase the number of Newton iterations over a
time step before the time step is accepted. \\
7 & --tolerance-cnv & TRGCNV & A real positive value. Set the maximum
non-linear tolerance error. This is the local convergence tolerance
(maximum of local saturation errors). The default value for
--tolerance-cnv is 0.01 and for TRGCNV is 0.001. \\
8 & --tolerance-cnv-relaxed & XXXCNV & A real positive value. Set the
maximum non-linear relaxed tolerance error. The default value for
--tolerance-cnv-relaxed is 1 and for XXXCNV is 0.01. \\
9 & --tolerance-mb & TRGMBE & A real positive value. Set the maximum
mass balance error, that is the tolerated mass balance error relative to
total mass present. The default value is 1.0 x 10-7. \\
10 & --tolerance-mb-relaxed & XXXMBE & A real positive value. Set the
maximum mass balance relaxed error. The default value is 1.0 x 10-6. \\
Notes: & & & \\
\end{longtable}

\emph{Table 2.5: Newton-Raphson Solver Controls}

Detailed descriptions of all of the command line parameters are given in
Table 2.1 in section 2.2 Running OPM Flow From The Command Line.

\subsection{Linear Solver Controls and
Options}\label{linear-solver-controls-and-options}

There are numerous options for the linear solver, which can be
overwhelming even for the most experienced users, given the rich feature
set enabled by the various command line parameters. Fortunately, the
simulator summarizes the options in use in the Linear Solver Property
Tree printed in both the \emph{.PRT and }.DBG files, right after
processing the input data and just before the simulation starts.

\pandocbounded{\includegraphics[keepaspectratio]{chapters/images/Frame135_e66a670170e0.png}}
\pandocbounded{\includegraphics[keepaspectratio]{chapters/images/Image128_e66a670170e0.png}}

Figure 2.41 above shows an example of the Linear Property Tree, based on
using mostly the linear solver default parameters. Here the tol value of
0.005 is equivalent to the --tolerance-cnv command line parameter, and
maxiter is equivalent to --linear-solver-max-iter. The key parameters
are summarized in Table 2.6. See also
\href{https://opm-project.org/?page_id=1798}{Customizing the Linear
Solver in OPM Flow}.

\begin{longtable}[]{@{}lll@{}}
\toprule\noalign{}
No. & Property Tree Parameter & Equivalent Command Line Parameter \\
\midrule\noalign{}
\endhead
\bottomrule\noalign{}
\endlastfoot
1 & tol & --tolerance-cnv \\
2 & maxiter & --linear-solver-max-iter \\
3 & verbosity & --linear-solver-verbosity \\
4 & solver & --use-gmres=false \\
5 & preconditioner & \\
6 & type & --linear-solver \\
Notes: & & \\
\end{longtable}

\emph{Table 2.6: Linear Solver Property Tree versus Command Line
Parameters}

Table 2.7 lists the major linear solver command line parameters
discussed in this section.

Figure 2.42 shows typical output for the linear solver for the March 1st
report step advancing to the next report step. Naturally, the Newton and
linear iteration counts are the same as for Figure 2.39 for the common
periods; however, oil, gas, and water material balance values are
printed for each non-linear iteration, together with the respective
convergence values.

\pandocbounded{\includegraphics[keepaspectratio]{chapters/images/Frame128_036d4e418156.png}}
\pandocbounded{\includegraphics[keepaspectratio]{chapters/images/Image45_036d4e418156.png}}

The total number of linear iterations displayed in Figure 2.42 is the
total number of iterations for all the non-linear iterations, not for a
particular non-linear iteration. To get the number of linear iterations
per non-linear iteration use the following command line parameter:

\begin{verbatim}
--linear-solver-verbosity=1
\end{verbatim}

which will produce a report similar to that shown in Figure 2.43, with
the IT parameter indicating the number of linear iterations for the
given Newton iteration.

\pandocbounded{\includegraphics[keepaspectratio]{chapters/images/Frame136_74a8e7b3224d.png}}
\pandocbounded{\includegraphics[keepaspectratio]{chapters/images/Image129_74a8e7b3224d.png}}

As mentioned previously, the convergence of the non-linear solver is
dependent on the accuracy of the linear solvers. Solving the linear
system to a loose accuracy may lead to poor convergence of the
non-linear solver. The tuning of the time stepping procedure, and the
tolerance used in the linear and non-linear solvers are highly related,
and optimal choices are case and problem dependent.

From the user prospective, there are several command-line options that
can be used to configure and control the linear solver, in particular:

\begin{verbatim}
- Using an alternative linear solver by changing
\end{verbatim}

\begin{verbatim}
--use-gmres
\end{verbatim}

The default value of false sets the linear solver to the Biconjugate
Gradient Stabilized (``BiCGSTAB'') linear solver, whereas a value of
true will set the linear solver to use the Generalized Minimal Residual
linear solver (``GMRES''). If the default linear solver is not
converging then changing to the GMRES solver may help. Note, that the
GMRES linear solver has a slightly higher memory consumption, as it
stores some of the previous search iterations as part of the
optimization process. The default number of search iterations stored is
15, which means 14 more vectors than BiCGSTAB solver.

\begin{verbatim}
- Changing the linear solver preconditioner can dramatically improve the performance of the linear solver, which is accomplished by the following command line parameter:
\end{verbatim}

\begin{verbatim}
--linear-solver
\end{verbatim}

In many cases using the default Constrained Pressure Residual (``CPR'')
preconditioner, instead of the widely used Incomplete Lower-triangular
Upper-triangular factorization with level of fill in 0 (``ILU0''),
results in difficult problems being solved at the expense of additional
work. For example, model two in the Tenth SPE Comparative Solution
Project {[}Christie, M. A., and M. J. Blunt. ``Tenth SPE Comparative
Solution Project: A Comparison of Upscaling Techniques.'' SPE Res Eval
\& Eng 4 (2001): 308--317. doi: https://doi.org/10.2118/72469-PA.{]},
fails to run using the ILU0 preconditioner, but runs satisfactory with
the CPR preconditioner. Note that the commercial simulator also has a
CPR option that is implemented via the CPR keyword in the RUNSPEC, but
the implementation is different to that of OPM Flow.

\begin{verbatim}
- And finally increasing the number of linear iterations may be appropriate, via:
\end{verbatim}

\begin{verbatim}
--linear-solver-max-iter
\end{verbatim}

that defines the maximum number of linear iterations. Sometimes
increasing this value may improve performance. Note that for the default
linear solver CPRW the flag is ignored, but one could set this and
additional values using a JSON file, see
\href{https://opm-project.org/?page_id=1798}{Customizing the Linear
Solver in OPM Flow} for details.

\begin{longtable}[]{@{}
  >{\raggedright\arraybackslash}p{(\linewidth - 6\tabcolsep) * \real{0.2500}}
  >{\raggedright\arraybackslash}p{(\linewidth - 6\tabcolsep) * \real{0.2500}}
  >{\raggedright\arraybackslash}p{(\linewidth - 6\tabcolsep) * \real{0.2500}}
  >{\raggedright\arraybackslash}p{(\linewidth - 6\tabcolsep) * \real{0.2500}}@{}}
\toprule\noalign{}
\begin{minipage}[b]{\linewidth}\raggedright
Major Linear Solver Controls
\end{minipage} & \begin{minipage}[b]{\linewidth}\raggedright
\end{minipage} & \begin{minipage}[b]{\linewidth}\raggedright
\end{minipage} & \begin{minipage}[b]{\linewidth}\raggedright
\end{minipage} \\
\midrule\noalign{}
\endhead
\bottomrule\noalign{}
\endlastfoot
No. & Command Line Parameter & TUNING Keyword Equivalent &
Description \\
1 & --linear-solver & & A defined string. Specify the configuration of
the linear solver, valid values are: Option (1) extends the existing
Constrained Pressure Residual (``CPR'') preconditioner to include wells.
This option can also be invoked via the CPR keyword in the RUNSPEC
section; however, the command line parameter takes precedence. For
option (10) one enters a character string enclosed in quotes that
defines the filename of a JSON configuration file for the flexible
linear solver system. In this case the file extension should be `.json'.
Note that the *.PRT file contains the ``Property tree for the linear
solver'' listing, which is the JSON specification of the current case,
and can be used to configure a user specific linear solver JSON file. \\
2 & --linear-solver-ignore-convergence-failure & & A boolean value set
to true or false. If set to true convergences failures in the linear
solver are ignored. This option should be used with care, as the results
may be unreliable. \\
3 & --linear-solver-max-iter & LITMAX & A positive integer value. Set
the maximum number of linear iterations. \\
4 & --linear-solver-reduction & & A real positive value. Set the minimum
reduction of the residual for the linear solver. The linear solver
convergences when the residual is reduced sufficiently. In general, it
is not recommended to modify this parameter. \\
5 & --linear-solver-restart & & A positive integer value. Set the number
of iterations after which GMRES is restarted. In general, it is not
recommended to modify this parameter. \\
6 & --use-gmres & & A boolean value. Use Generalized Minimal Residual
(``GMRES'') solver instead of Biconjugate Gradient Stabilized
(``BiCGSTAB'') as the linear solver within the Newton iterations. \\
Notes: & & & \\
\end{longtable}

\emph{Table 2.7: Linear Solver Controls}

Again detailed descriptions of all of the command line parameters are
given in Table 2.1 in section 2.2 Running OPM Flow From The Command
Line.

\bookmarksetup{startatroot}

\chapter{KEYWORD DOCUMENTATION
STRUCTURE}\label{keyword-documentation-structure}

\section{Introduction}\label{introduction-1}

The OPM Flow manual is constructed in a manner to enable the reader to
reference various parts of the document by using the table contents or
by simply pressing on a link embedded in the text. This automatic cross
referencing has been extensively employed to ensure effective
documentation of the keywords used by the simulator.

There are several key features that can be used to navigate the manual
in an efficient manner. The first is the ``footer'' which can be used to
move to various sections in the manual. For example, the ``Table of
Contents'' footer shown below:

\pandocbounded{\includegraphics[keepaspectratio]{chapters/images/Frame63_62664cb5267c.png}}
\pandocbounded{\includegraphics[keepaspectratio]{chapters/images/Image57_62664cb5267c.png}}

allows the reader to move to the Table of Contents by clicking on the
link (highlighted in italic blue text on a gray background in Figure
3.1). Note also that the entries in the table of contents are also
``clickable'' enabling the reader to move to the desired entry.

The second type of footer is the ``Keyword Index'' footer that also
contains the Table of Contents link mentioned above. This footer is
illustrated in Figure 3.2 and is displayed on a keyword definition page.

\pandocbounded{\includegraphics[keepaspectratio]{chapters/images/Frame64_b74b8c203553.png}}
\pandocbounded{\includegraphics[keepaspectratio]{chapters/images/Image58_b74b8c203553.png}}

Clicking on a letter (highlighted in italic blue text on a gray
background in) takes the reader to an alphabetic listing of all the
keywords beginning with the selected letter (Figure 3.3).

\pandocbounded{\includegraphics[keepaspectratio]{chapters/images/Frame65_d621d11b4bea.png}}
\pandocbounded{\includegraphics[keepaspectratio]{chapters/images/Image59_d621d11b4bea.png}}

The list is color coded, so one instantly knows what keywords are
implemented, with green colored cells indicating the keyword is
available and is mostly or fully implemented. Cells colored orange show
that keyword is recognized but not implemented in OPM Flow. Finally,
cells colored red mean that keyword is available in the commercial
simulator but has not been implemented in OPM Flow, and may cause an
error if used in the input deck. Clicking on a keyword in the list will
move the reader to the keyword definition. Note also that clicking one
of the section names (RUNSPEC, GRID, etc.) in Figure 3.3, will take the
reader to the beginning of the selected section.

Finally, in the PDF version of the manual if one displays the
``bookmarks'' in the PDF reading software one can jump to a particular
keyword or section without having to scroll up or down.

\section{Keyword Definitions}\label{keyword-definitions}

Each keyword is defined in it's own section that contains a section
header, that contains the keyword name in capital letters followed by a
brief description of the keyword's function. This is then followed by
the Keyword Table Section which defines the status of the keyword and in
which sections of the input deck the keyword can be utilized. Table 3.1
illustrates a typical Keyword Table Section defining the keyword status
with the various OPM Flow sections.

\begin{longtable}[]{@{}
  >{\raggedright\arraybackslash}p{(\linewidth - 14\tabcolsep) * \real{0.1250}}
  >{\raggedright\arraybackslash}p{(\linewidth - 14\tabcolsep) * \real{0.1250}}
  >{\raggedright\arraybackslash}p{(\linewidth - 14\tabcolsep) * \real{0.1250}}
  >{\raggedright\arraybackslash}p{(\linewidth - 14\tabcolsep) * \real{0.1250}}
  >{\raggedright\arraybackslash}p{(\linewidth - 14\tabcolsep) * \real{0.1250}}
  >{\raggedright\arraybackslash}p{(\linewidth - 14\tabcolsep) * \real{0.1250}}
  >{\raggedright\arraybackslash}p{(\linewidth - 14\tabcolsep) * \real{0.1250}}
  >{\raggedright\arraybackslash}p{(\linewidth - 14\tabcolsep) * \real{0.1250}}@{}}
\toprule\noalign{}
\begin{minipage}[b]{\linewidth}\raggedright
RUNSPEC
\end{minipage} & \begin{minipage}[b]{\linewidth}\raggedright
GRID
\end{minipage} & \begin{minipage}[b]{\linewidth}\raggedright
EDIT
\end{minipage} & \begin{minipage}[b]{\linewidth}\raggedright
PROPS
\end{minipage} & \begin{minipage}[b]{\linewidth}\raggedright
REGIONS
\end{minipage} & \begin{minipage}[b]{\linewidth}\raggedright
SOLUTION
\end{minipage} & \begin{minipage}[b]{\linewidth}\raggedright
SUMMARY
\end{minipage} & \begin{minipage}[b]{\linewidth}\raggedright
SCHEDULE
\end{minipage} \\
\midrule\noalign{}
\endhead
\bottomrule\noalign{}
\endlastfoot
\end{longtable}

\emph{Table 3.1: Example Keyword Table Section}

The cells are colored in such a manner as to quickly indicate to the
reader the keyword's section availability and function availability,
with green colored cells indicating the keyword is available for this
section and is mostly or fully implemented. Cells colored gray indicate
that keyword cannot be used in that particular section, and cells
colored orange show that the keyword is only partially implemented
within OPM Flow, for example OPM Flow may just recognize the keyword and
ignore the keyword's function. Finally, cells colored red mean that the
keyword is available in the commercial simulator but has not been
implemented in OPM Flow. In this scenario the keyword should not be used
in OPM Flow as it will result in unpredictable results, including
causing the simulator to abort or throw an exception.

\section{Multi-Section Keywords}\label{multi-section-keywords}

As there are numerous keywords that can be used within multiple OPM Flow
sections of the input file, for example the ADD and EQUALS keywords,
there is a need to avoid duplication of the keyword definitions but at
the same time attempt to define only those keywords for a given section.
Thus for multi-section keywords, the keyword is defined in the first
available section that the keyword can be found. The Keyword Table
Section as shown below for the ADD keyword below, indicates in which
sections the keyword can be utilized.

\begin{longtable}[]{@{}
  >{\raggedright\arraybackslash}p{(\linewidth - 14\tabcolsep) * \real{0.1250}}
  >{\raggedright\arraybackslash}p{(\linewidth - 14\tabcolsep) * \real{0.1250}}
  >{\raggedright\arraybackslash}p{(\linewidth - 14\tabcolsep) * \real{0.1250}}
  >{\raggedright\arraybackslash}p{(\linewidth - 14\tabcolsep) * \real{0.1250}}
  >{\raggedright\arraybackslash}p{(\linewidth - 14\tabcolsep) * \real{0.1250}}
  >{\raggedright\arraybackslash}p{(\linewidth - 14\tabcolsep) * \real{0.1250}}
  >{\raggedright\arraybackslash}p{(\linewidth - 14\tabcolsep) * \real{0.1250}}
  >{\raggedright\arraybackslash}p{(\linewidth - 14\tabcolsep) * \real{0.1250}}@{}}
\toprule\noalign{}
\begin{minipage}[b]{\linewidth}\raggedright
RUNSPEC
\end{minipage} & \begin{minipage}[b]{\linewidth}\raggedright
GRID
\end{minipage} & \begin{minipage}[b]{\linewidth}\raggedright
EDIT
\end{minipage} & \begin{minipage}[b]{\linewidth}\raggedright
PROPS
\end{minipage} & \begin{minipage}[b]{\linewidth}\raggedright
REGIONS
\end{minipage} & \begin{minipage}[b]{\linewidth}\raggedright
SOLUTION
\end{minipage} & \begin{minipage}[b]{\linewidth}\raggedright
SUMMARY
\end{minipage} & \begin{minipage}[b]{\linewidth}\raggedright
SCHEDULE
\end{minipage} \\
\midrule\noalign{}
\endhead
\bottomrule\noalign{}
\endlastfoot
\end{longtable}

\emph{Table 3.2: ADD Keyword Table Section}

Here the ADD keyword can be used in the GRID, EDIT, PROPS, REGIONS, and
SOLUTION sections as indicated by those cells colored green and not for
the cells colored in light gray. A full description of the ADD keyword
is in the GRID section. In the EDIT, PROPS, REGIONS, and SOLUTION
sections there is a brief description of the ADD keyword with a link to
the full description, as shown below.

The ADD keyword adds a constant to a specified array or part of an
array. The constant can be real or integer depending on the array type;
however, the arrays that can be operated on is dependent on which
section the ADD keyword is being applied. See ADD -- Add a Constant to a
Specified Array in the GRID section for a full description.

A complete list of keywords in alphabetic order is given in APPENDIX A:
KEYWORD INDEX - ALPHABETIC LISTING , and clicking on a specific keyword
in the appendix will take the reader to the keyword definition in a
particular section.

\section{Keyword Formats}\label{keyword-formats}

All keywords in OPM Flow should be entered in capital case and start in
column one, lowercase entry of keywords will produce errors, and
keywords not starting in column one will not be recognized. There are
five types of keyword format types used by OPM Flow for data input. The
description of the five types is given in the next three sections
together with some examples.

\subsection{Keyword Format Type --
Comment}\label{keyword-format-type-comment}

Comments in the input deck can occur anywhere in the file are preceded
by ``--'' in columns one and two, for example for the EQLDIMS keyword:

\begin{verbatim}
--
--       MAX     MAX     RSVD    TVDP    TVDP
--       EQLNUM  DEPTH   NODES   TABLE   NODES
EQLDIMS
         9       1*      20      1*      1*                                    /
\end{verbatim}

In addition, comments can be place after ``/'' that terminates a record
entry as shown below;

\begin{verbatim}
--
-- -- ARRAY    CONSTANT --  ---------- BOX ---------
--                          I1  I2   J1  J2   K1  K2
MULTIPLY
   'PERMZ'     0.50000      1*  1*   1*  1*   1*  1* / PERMZ * 0.5
/
\end{verbatim}

\subsection{Keyword Format Type --
Activation}\label{keyword-format-type-activation}

This type of keyword format only consists of the keyword itself and is
usually used to invoke a feature or to switch on or off a processing
feature. The keyword is documented by describing the functionality or
action the keyword performs, followed by an example. Examples of this
type of keyword include API (to switch on API tracking), GAS (to
activate the gas phase in the model), ECHO (to switching echoing of the
input file to the output file), and SKIP (for skipping parts of the
input deck). For example the GAS keyword in the RUNSPEC section would be
described as:

Description

This keyword indicate that the gas phase is present in the model and
must be used for oil-gas, gas-water, oil-water-gas input decks that
contain the gas phase. The keyword will also invoke the data input file
checking to ensure that all the required gas phase input parameters are
defined in the input deck.

There is no data required for this keyword.

Example

\begin{verbatim}
--
--       GAS PHASE IS PRESENT IN THE RUN
--
GAS
\end{verbatim}

The above example declares that the gas phase is active in the model.

\subsection{Keyword Format Type - Vector (Row
Vector)}\label{keyword-format-type---vector-row-vector}

Vector based keywords consist of the keyword followed by a vector of
parameters on a separate line and may consists of multiple lines of
vectors with each line representing a data set (see the second example
for this type of vector keyword). The vector may contain integer, real,
and character parameters depending on the keywords requirements. This
type of keyword is documented by describing the functionality or action
the keyword performs, a table describing the parameters associated with
the keyword, followed by one or two examples on how to use the keyword.
For example the DIMENS keyword in the RUNSPEC section would be described
as:

Description

DIMENS defines the dimensions of the model entered as integer vector.
The keyword can be used with all grid types.

\begin{longtable}[]{@{}
  >{\raggedright\arraybackslash}p{(\linewidth - 6\tabcolsep) * \real{0.2500}}
  >{\raggedright\arraybackslash}p{(\linewidth - 6\tabcolsep) * \real{0.2500}}
  >{\raggedright\arraybackslash}p{(\linewidth - 6\tabcolsep) * \real{0.2500}}
  >{\raggedright\arraybackslash}p{(\linewidth - 6\tabcolsep) * \real{0.2500}}@{}}
\toprule\noalign{}
\begin{minipage}[b]{\linewidth}\raggedright
No.
\end{minipage} & \begin{minipage}[b]{\linewidth}\raggedright
Name
\end{minipage} & \begin{minipage}[b]{\linewidth}\raggedright
Description
\end{minipage} & \begin{minipage}[b]{\linewidth}\raggedright
Default
\end{minipage} \\
\midrule\noalign{}
\endhead
\bottomrule\noalign{}
\endlastfoot
1 & NX & The number of grid blocks in the x direction for Cartesian
grids or the number of grid blocks in the r direction for radial grids &
None \\
2 & NY & The number of grid blocks in the y direction for Cartesian
grids or the number of grid blocks in the theta direction for radial
grids. & None \\
3 & NZ & The number of grid blocks in the z direction for both Cartesian
and radial grids. & None \\
Notes: & & & \\
\end{longtable}

\emph{Table 3.3: DIMENS Keyword Description}

Note that NX, NY, and NZ are not maximum values but the actual size of
the grid. OPM Flow applies these parameters when reading in particular
data sets. For example if NX, NY, and NZ are set to 10, 10, and 10
respectively, then for the grid property data like PORO; OPM Flow
expects to read in 10 x 10 x 10 or 1,000 porosity values for the PORO
array. If the number of porosity values is not equal to 1,000 then OPM
Flow will produce an error.

Example

\begin{verbatim}
--
--       MAX     MAX     MAX
--       NDIVIX  NDIVIY  NDIVIZ
DIMENS
         46      112     22                                                    /
\end{verbatim}

The above example defines the dimensions for the Norne model of 36 cells
in the x direction, 122 cells in the y direction, and 22 cells in the z
direction.

For vector keywords that have parameters associated with units, then
there is a slightly different table format from the one used above to
take into account documenting the default values for the three sets of
units supported by OPM Flow, for example for the ROCK keyword is
describe as follows:

Description

ROCK defines the rock compressibility for various regions in the model.
The number of ROCK vector data sets is defined by the NTPVT parameter on
the TABDIMS keyword in the RUNSPEC section and the allocation of the
ROCK tables to different grid blocks in the model is done via the PVTNUM
keyword in the REGIONS section. One data set consists of one record or
line which is terminated by a ``/''.

This keyword must be defined in the OPM Flow input deck.

\begin{longtable}[]{@{}
  >{\raggedright\arraybackslash}p{(\linewidth - 6\tabcolsep) * \real{0.2500}}
  >{\raggedright\arraybackslash}p{(\linewidth - 6\tabcolsep) * \real{0.2500}}
  >{\raggedright\arraybackslash}p{(\linewidth - 6\tabcolsep) * \real{0.2500}}
  >{\raggedright\arraybackslash}p{(\linewidth - 6\tabcolsep) * \real{0.2500}}@{}}
\toprule\noalign{}
\begin{minipage}[b]{\linewidth}\raggedright
No.
\end{minipage} & \begin{minipage}[b]{\linewidth}\raggedright
Name
\end{minipage} & \begin{minipage}[b]{\linewidth}\raggedright
Description
\end{minipage} & \begin{minipage}[b]{\linewidth}\raggedright
Default
\end{minipage} \\
\midrule\noalign{}
\endhead
\bottomrule\noalign{}
\endlastfoot
Field & Metric & Laboratory & \\
1 & Pref & Pref is a real number defining the reference pressure for the
other parameters for this data set. & Default \\
psia 1.032 & barsa 1.032 & atma 1.032 & \\
2 & Cf & Cf is a real number defining the rock compressibility at the
reference pressure, Cf(Pref) and is defined as:
\({C}_{f} = -\frac{1}{V}\left(\frac{\mathit{dV}}{\mathit{dP}}\right)\) &
Defined \\
1/psia 0.0 & 1/barsa 0.0 & 1/atma 0.0 & \\
Notes: & & & \\
\end{longtable}

\emph{Table 3.4: ROCK Keyword Description}

Example

The following shows the PVTW keyword for when NTPVT on the TABDIMS
keyword in the RUNSPEC section is set to one.

\begin{verbatim}
--
-- ROCK COMPRESSIBILITY
--
-- (1) REFERENCE PRESSURE IS TAKEN FROM THE HCPV WEIGHTED RESERVOIR PRESSURE
--     AS THE PORV IS ALREADY AT RESERVOIR CONDITIONS (FLOW USES THE REFERENCE
--     PRESSURE) TO CONVERT THE GIVEN PORV TO RESERVOIR CONDITIONS USING THE DATA
--     ON THE ROCK KEYWORD)
--
ROCK
    3566.9    5.0E-06                             / ROCK COMPRESSIBILITY REGION 1
    3966.9    5.5E-06                             / ROCK COMPRESSIBILITY REGION 2
    4566.9    6.0E-06                             / ROCK COMPRESSIBILITY REGION 3
\end{verbatim}

There is no terminating ``/'' for this keyword.

In this case the example shows a multiple data set entry of the vector
format keyword, with three ROCK data sets being defined by the keyword.

\subsection{Keyword Format Type -- Vector (Columnar
Vector)}\label{keyword-format-type-vector-columnar-vector}

Columnar vector based keywords consist of the keyword followed by a
columnar vector of parameters in a separate column for each parameter.
The vector may contain integer, real, and character parameters depending
on the keywords requirements. This type of keyword is documented by
describing the functionality or action the keyword performs, a table
describing the parameters associated with the keyword, followed by one
or two examples on how to use the keyword. For example the SWFN keyword
in the PROPS section would be described as:

Description

The SWFN keyword defines the water relative permeability and water-oil
capillary pressure data versus water saturation tables for when water is
present in the input deck. This keyword should only be used if water is
present in the run.

\begin{longtable}[]{@{}
  >{\raggedright\arraybackslash}p{(\linewidth - 6\tabcolsep) * \real{0.2500}}
  >{\raggedright\arraybackslash}p{(\linewidth - 6\tabcolsep) * \real{0.2500}}
  >{\raggedright\arraybackslash}p{(\linewidth - 6\tabcolsep) * \real{0.2500}}
  >{\raggedright\arraybackslash}p{(\linewidth - 6\tabcolsep) * \real{0.2500}}@{}}
\toprule\noalign{}
\begin{minipage}[b]{\linewidth}\raggedright
No.
\end{minipage} & \begin{minipage}[b]{\linewidth}\raggedright
Name
\end{minipage} & \begin{minipage}[b]{\linewidth}\raggedright
Description
\end{minipage} & \begin{minipage}[b]{\linewidth}\raggedright
Default
\end{minipage} \\
\midrule\noalign{}
\endhead
\bottomrule\noalign{}
\endlastfoot
Field & Metric & Laboratory & \\
1 & SWAT & A columnar vector of real monotonically increasing down the
column values starting from zero and terminating at one, that defines
the water saturation. & None \\
dimensionless & dimensionless & dimensionless & \\
2 & KRW & A columnar vector of real values that are either equal or
increasing down the column and that are greater than or equal to zero
and less than or equal to one that defines the water relative
permeability with respect to gas saturation. The first value in the
column should be zero. & None \\
dimensionless & dimensionless & dimensionless & \\
3 & PCWO & A columnar vector of real values that are either equal or
increasing down the column that defines the water-oil relative capillary
pressure. If the SWATINIT keyword has been used to initialize the model
then columnar vector has to be strictly monotonically increasing. &
None \\
psia & bars & atm & \\
Notes: & & & \\
\end{longtable}

\emph{Table 3.3: SWFN Keyword Description}

Example

\begin{verbatim}
--
--       WATER RELATIVE PERMEABILITY TABLES (SWFN)
--
SWFN
--       SWAT       KRW        PCOW
--       FRAC       FRAC       PSIA
--       --------   --------   -------
            0.15     0.00000     1*
            0.30     0.00050     1*
            0.40     0.00390     1*
            0.50     0.01500     1*
            0.60     0.04100     1*
            0.65     0.06250     1*
            0.70     0.09150     1*
            0.80     0.17850     1*
            0.90     0.31640     1*
            0.95     0.40960     1*
            1.00     0.52200     1*                        / TABLE NO. 1
\end{verbatim}

The example defines two SWFN tables for use when water is present in the
run. In the tables the water-oil capillary pressure data has been
defaulted with ``1*'' and will be set to zero as there are no other
values for the water-oil capillary pressure columns.

\subsection{Keyword Format Type --
Array}\label{keyword-format-type-array}

This type of keyword defines a property for the grid or an area of the
grid using a previously entered BOX keyword to define the area where the
property will be defined. For array data a full set of values for each
element in the array is required. For example, the documentation for the
PORO array would be:

Description

PORO defines the porosity for all the cells in the model via an array.
The keyword can be used with all grid types.

\begin{longtable}[]{@{}
  >{\raggedright\arraybackslash}p{(\linewidth - 6\tabcolsep) * \real{0.2500}}
  >{\raggedright\arraybackslash}p{(\linewidth - 6\tabcolsep) * \real{0.2500}}
  >{\raggedright\arraybackslash}p{(\linewidth - 6\tabcolsep) * \real{0.2500}}
  >{\raggedright\arraybackslash}p{(\linewidth - 6\tabcolsep) * \real{0.2500}}@{}}
\toprule\noalign{}
\begin{minipage}[b]{\linewidth}\raggedright
No.
\end{minipage} & \begin{minipage}[b]{\linewidth}\raggedright
Name
\end{minipage} & \begin{minipage}[b]{\linewidth}\raggedright
Description
\end{minipage} & \begin{minipage}[b]{\linewidth}\raggedright
Default
\end{minipage} \\
\midrule\noalign{}
\endhead
\bottomrule\noalign{}
\endlastfoot
Field & Metric & Laboratory & \\
1 & PORO & PORO is an array of real numbers assigning the porosity
values to each cell in the model. The number of entries should
correspond to the NX x NY x NZ parameters on the DIMENS keyword. Repeat
counts may be used, for example 30*100.0. & None \\
dimensionless & dimensionless & dimensionless & \\
Notes: & & & \\
\end{longtable}

\emph{Table 3.3: PORO Keyword Description}

See also the DX, DY, and TOPS keywords to fully define a Cartesian
Regular Grid.

Example

\begin{verbatim}
--
--  DEFINE GRID BLOCK POROSITY DATA FOR ALL CELLS (BASED ON NX x NY x NZ = 300)
--
PORO
  300*0.300                                                                       /
\end{verbatim}

\section{Input File Structure}\label{input-file-structure}

OPM Flow input files are similar to commercial simulators that are used
in the oil and gas industry, that is the input file is separated into
sections in an effort to avoid errors in an engineer's input data and a
computer programmer's code to interpret the data. OPM Flow has been
designed from an engineer's prospective with an input structure similar
to Schlumberger's industry wide ECLIPSE 100 {[}ECLIPSE
Industry-Reference Reservoir Simulator -- Reference Manual 2019.1,
Schlumberger.{]} simulator. Table 3.4 lists the various sections
together with a brief description of purpose of each section, as well as
declaring if a section is mandatory or not for a run to form a valid
input deck.

\begin{longtable}[]{@{}
  >{\raggedright\arraybackslash}p{(\linewidth - 4\tabcolsep) * \real{0.3333}}
  >{\raggedright\arraybackslash}p{(\linewidth - 4\tabcolsep) * \real{0.3333}}
  >{\raggedright\arraybackslash}p{(\linewidth - 4\tabcolsep) * \real{0.3333}}@{}}
\toprule\noalign{}
\begin{minipage}[b]{\linewidth}\raggedright
Section Name
\end{minipage} & \begin{minipage}[b]{\linewidth}\raggedright
Description
\end{minipage} & \begin{minipage}[b]{\linewidth}\raggedright
Required Optional
\end{minipage} \\
\midrule\noalign{}
\endhead
\bottomrule\noalign{}
\endlastfoot
RUNSPEC & This is the first section in the OPM Flow input file and
defines the key parameters for the simulator including the dimensions of
the model, phases present in the model (oil, gas, and water for
example), number of tables for a given property, and the maximum number
of rows for each table, the maximum number of groups, wells and well
completions, as well as various options to be invoked by OPM Flow. &
Required \\
GRID & Defines the basic grid properties, including structure, faults,
and various static rock properties (porosity, permeability, etc.). Upon
completion of reading this section, the software calculates the pore
volume (PORV) for each cell and the transmissibilities (TRANX, TRANY,
and TRANZ) between all neighbouring cells, as well as calculating the
transmissibilities of the Non-Neighbor Connections (``NNC''). &
Required \\
EDIT & The properties calculated by OPM Flow in the GRID section are
available for editing in this section (PORV, TRANX, etc.). & Optional \\
PROPS & This section defines the fluid properties for all the phases
present in the run, for example oil viscosity, oil formation volume
factor, etc. The section also defines the rock flow properties such as
the relative permeabilities and the capillary pressure functions. &
Required \\
REGIONS & This section allows the engineer to define various regions in
the model for reporting purposes and to define how the fluid and rock
properties defined in the PROPS section are allocated throughout the
model. & Optional \\
SOLUTION & Defines the parameters to initialize the model, fluid
contacts, reservoir pressures, etc., together with the data from the
previous sections. This section, if requested, reports the initial
in-place volumes for phases present in the model, as well as the average
pressure for the various defined regions. & Required \\
SUMMARY1 & Defines the vectors of summary results versus time to be
written to various output files for reviewing the results of the
simulation. This includes field, group, well and well completion
production, and injection results, for example field oil rate versus
time. Grid block property results can also be reported versus time, for
example grid block pressure versus time2. & Optional \\
SCHEDULE & This is the final section that defines the time dependent
data such as the field, group and well parameters, targets and
constraints, the numerical controls, the operating schedule, and the
reporting requirements. & Required \\
Notes: & & \\
\end{longtable}

\emph{Table 3.4: OPM Flow Input Deck Sections}

\bookmarksetup{startatroot}

\chapter{GLOBAL SECTION KEYWORDS}\label{global-section-keywords}

\section{Introduction}\label{introduction-2}

Keywords used in this section can be used in all input file sections,
that is the RUNSPEC, GRID, EDIT PROPS, REGIONS, SOLUTION, SUMMARY, and
SCHEDULE sections.

\section{Keyword Definitions}\label{keyword-definitions-1}

\section{Data Requirements}\label{data-requirements}

There are no data requirements for the GLOBAL section keywords, as all
the keywords are optional and considered to be utility keywords that can
be used throughout the input deck. A complete list of these global
keywords is given in Table 4.1.

\begin{longtable}[]{@{}
  >{\raggedright\arraybackslash}p{(\linewidth - 6\tabcolsep) * \real{0.2500}}
  >{\raggedright\arraybackslash}p{(\linewidth - 6\tabcolsep) * \real{0.2500}}
  >{\raggedright\arraybackslash}p{(\linewidth - 6\tabcolsep) * \real{0.2500}}
  >{\raggedright\arraybackslash}p{(\linewidth - 6\tabcolsep) * \real{0.2500}}@{}}
\toprule\noalign{}
\begin{minipage}[b]{\linewidth}\raggedright
No.
\end{minipage} & \begin{minipage}[b]{\linewidth}\raggedright
GLOBAL Keyword
\end{minipage} & \begin{minipage}[b]{\linewidth}\raggedright
Description
\end{minipage} & \begin{minipage}[b]{\linewidth}\raggedright
Section Applied
\end{minipage} \\
\midrule\noalign{}
\endhead
\bottomrule\noalign{}
\endlastfoot
1 & COLUMNS & COLUMNS -- Define Input File Column Margins. The COLUMNS
keyword defines the input file column margins; characters outside the
margins are ignored by the input parser. & ALL \\
2 & DEBUG & DEBUG -- Define the Debug Data to be Printed to File. This
keyword defines the debug data to be written to the debug file (*.DBG),
it is ignored by OPM Flow. & ALL \\
3 & ECHO & ECHO -- Activate Echoing of User Input Files to the Print
File. Turns on echoing of all the input files to the print file. &
ALL \\
4 & END & END -- Define the End of the Input File. This keyword marks
the end of the input file and can occur in any section. Any keywords and
data after the END keyword are ignored. & ALL \\
5 & ENDINC & ENDINC -- Define the End of an Include File. This keyword
marks the end of an include file specified on the INCLUDE keyword. When
the ENDINC keyword is encountered in the INCLUDE file, input data is
read from the next keyword in the current file. & ALL \\
6 & ENDSKIP & ENDSKIP -- Deactivate Skipping of Keywords and Input Data.
Turns off skipping of keywords that was activated by the SKIP, SKIP100,
or SKIP300 keywords. & ALL \\
7 & EXTRAPMS & EXTRAPMS -- Activate Extrapolation Warning Messages. The
EXTRAPMS keyword activates extrapolation warning messages for when
simulator extrapolates the PVT or VFP tables. & ALL \\
8 & FORMFEED & FORMFEED -- Defined the Print File Form-Feed Character.
The keyword defines the form-feed character, or carriage control
character, for the output print (\emph{.PRT) run summary (}.RSM) files.
& ALL \\
9 & INCLUDE & INCLUDE -- Load Another Data File at the Current Position.
The INCLUDE keyword informs OPM Flow to continue reading input data from
the specified INCLUDE file. & ALL \\
10 & MESSAGE & MESSAGE -- Output User Message. The MESSAGE keyword
outputs a user message to the terminal, as well as to the print
(\emph{.PRT) and debug (}.DBG) files. & ALL \\
11 & MESSAGES & MESSAGES -- Define Message Print Limits and Stop Limits.
The MESSAGES keyword defines the print and stops levels for various
messages. & ALL \\
12 & NOECHO & NOECHO -- Deactivate Echoing of User Input Files to the
Print File. Turns off echoing of all the input files to the print file.
& ALL \\
13 & NOWARN & NOWARN -- Deactivate Warning Messages. Turns off warning
messages to be printed to the print file. & ALL \\
14 & SKIP & SKIP -- Activate Skipping of All Keywords and Input Data.
The keyword turns on skipping of keywords until the ENDSKIP activation
keyword is encountered. & ALL \\
15 & SKIP100 & SKIP100 -- Activate Skipping of Keywords by Black-Oil
Simulator. This keyword turns on skipping of all keywords and input data
by the commercial black-oil simulator until the ENDSKIP keyword is
encountered. & ALL \\
16 & SKIP300 & SKIP300 -- Activate Skipping of Compositional Keywords.
This keyword turns on skipping of all keywords and input data by the
commercial compositional simulator until the ENDSKIP keyword is
encountered. & ALL \\
17 & WARN & WARN -- Activate Warning Messages. Turns on warning messages
to be printed to the print file (*.PRT); note that this keyword is
activated by default. & ALL \\
Notes: & & & \\
\end{longtable}

\emph{Table 4.1: Alphabetic List of GLOBAL Section Keywords}

A more detailed description of the keywords is presented in the
following sections.

\section{Keyword Definitions}\label{keyword-definitions-2}

\subsection{COLUMNS -- Define Input File Column
Margins}\label{columns-define-input-file-column-margins}

\begin{longtable}[]{@{}
  >{\raggedright\arraybackslash}p{(\linewidth - 14\tabcolsep) * \real{0.1250}}
  >{\raggedright\arraybackslash}p{(\linewidth - 14\tabcolsep) * \real{0.1250}}
  >{\raggedright\arraybackslash}p{(\linewidth - 14\tabcolsep) * \real{0.1250}}
  >{\raggedright\arraybackslash}p{(\linewidth - 14\tabcolsep) * \real{0.1250}}
  >{\raggedright\arraybackslash}p{(\linewidth - 14\tabcolsep) * \real{0.1250}}
  >{\raggedright\arraybackslash}p{(\linewidth - 14\tabcolsep) * \real{0.1250}}
  >{\raggedright\arraybackslash}p{(\linewidth - 14\tabcolsep) * \real{0.1250}}
  >{\raggedright\arraybackslash}p{(\linewidth - 14\tabcolsep) * \real{0.1250}}@{}}
\toprule\noalign{}
\begin{minipage}[b]{\linewidth}\raggedright
RUNSPEC
\end{minipage} & \begin{minipage}[b]{\linewidth}\raggedright
GRID
\end{minipage} & \begin{minipage}[b]{\linewidth}\raggedright
EDIT
\end{minipage} & \begin{minipage}[b]{\linewidth}\raggedright
PROPS
\end{minipage} & \begin{minipage}[b]{\linewidth}\raggedright
REGIONS
\end{minipage} & \begin{minipage}[b]{\linewidth}\raggedright
SOLUTION
\end{minipage} & \begin{minipage}[b]{\linewidth}\raggedright
SUMMARY
\end{minipage} & \begin{minipage}[b]{\linewidth}\raggedright
SCHEDULE
\end{minipage} \\
\midrule\noalign{}
\endhead
\bottomrule\noalign{}
\endlastfoot
\end{longtable}

\subsubsection{Description}\label{description}

The COLUMNS keyword defines the input file column margins; characters
outside the margins are ignored by the input parser.

This keyword is not supported by OPM Flow but would change the results
if supported so the simulation will be stopped.

\subsection{DEBUG -- Define the Debug Data to be Printed to
File}\label{debug-define-the-debug-data-to-be-printed-to-file}

\begin{longtable}[]{@{}
  >{\raggedright\arraybackslash}p{(\linewidth - 14\tabcolsep) * \real{0.1250}}
  >{\raggedright\arraybackslash}p{(\linewidth - 14\tabcolsep) * \real{0.1250}}
  >{\raggedright\arraybackslash}p{(\linewidth - 14\tabcolsep) * \real{0.1250}}
  >{\raggedright\arraybackslash}p{(\linewidth - 14\tabcolsep) * \real{0.1250}}
  >{\raggedright\arraybackslash}p{(\linewidth - 14\tabcolsep) * \real{0.1250}}
  >{\raggedright\arraybackslash}p{(\linewidth - 14\tabcolsep) * \real{0.1250}}
  >{\raggedright\arraybackslash}p{(\linewidth - 14\tabcolsep) * \real{0.1250}}
  >{\raggedright\arraybackslash}p{(\linewidth - 14\tabcolsep) * \real{0.1250}}@{}}
\toprule\noalign{}
\begin{minipage}[b]{\linewidth}\raggedright
RUNSPEC
\end{minipage} & \begin{minipage}[b]{\linewidth}\raggedright
GRID
\end{minipage} & \begin{minipage}[b]{\linewidth}\raggedright
EDIT
\end{minipage} & \begin{minipage}[b]{\linewidth}\raggedright
PROPS
\end{minipage} & \begin{minipage}[b]{\linewidth}\raggedright
REGIONS
\end{minipage} & \begin{minipage}[b]{\linewidth}\raggedright
SOLUTION
\end{minipage} & \begin{minipage}[b]{\linewidth}\raggedright
SUMMARY
\end{minipage} & \begin{minipage}[b]{\linewidth}\raggedright
SCHEDULE
\end{minipage} \\
\midrule\noalign{}
\endhead
\bottomrule\noalign{}
\endlastfoot
\end{longtable}

\subsubsection{Description}\label{description-1}

This keyword defines the debug data to be written to the debug file
(*.DBG). This keyword is not supported by OPM Flow but has no effect on
the results so it will be ignored.

\subsection{ECHO -- Activate Echoing of User Input Files to the Print
File}\label{echo-activate-echoing-of-user-input-files-to-the-print-file}

\begin{longtable}[]{@{}
  >{\raggedright\arraybackslash}p{(\linewidth - 14\tabcolsep) * \real{0.1250}}
  >{\raggedright\arraybackslash}p{(\linewidth - 14\tabcolsep) * \real{0.1250}}
  >{\raggedright\arraybackslash}p{(\linewidth - 14\tabcolsep) * \real{0.1250}}
  >{\raggedright\arraybackslash}p{(\linewidth - 14\tabcolsep) * \real{0.1250}}
  >{\raggedright\arraybackslash}p{(\linewidth - 14\tabcolsep) * \real{0.1250}}
  >{\raggedright\arraybackslash}p{(\linewidth - 14\tabcolsep) * \real{0.1250}}
  >{\raggedright\arraybackslash}p{(\linewidth - 14\tabcolsep) * \real{0.1250}}
  >{\raggedright\arraybackslash}p{(\linewidth - 14\tabcolsep) * \real{0.1250}}@{}}
\toprule\noalign{}
\begin{minipage}[b]{\linewidth}\raggedright
RUNSPEC
\end{minipage} & \begin{minipage}[b]{\linewidth}\raggedright
GRID
\end{minipage} & \begin{minipage}[b]{\linewidth}\raggedright
EDIT
\end{minipage} & \begin{minipage}[b]{\linewidth}\raggedright
PROPS
\end{minipage} & \begin{minipage}[b]{\linewidth}\raggedright
REGIONS
\end{minipage} & \begin{minipage}[b]{\linewidth}\raggedright
SOLUTION
\end{minipage} & \begin{minipage}[b]{\linewidth}\raggedright
SUMMARY
\end{minipage} & \begin{minipage}[b]{\linewidth}\raggedright
SCHEDULE
\end{minipage} \\
\midrule\noalign{}
\endhead
\bottomrule\noalign{}
\endlastfoot
\end{longtable}

\subsubsection{Description}\label{description-2}

Turns on echoing of all the input files to the print file; note that
this keyword is activated by default and can subsequently be switched
off by the NOECHO activation keyword.

There is no data required for this keyword.

This keyword is not supported by OPM Flow but has no effect on the
results so it will be ignored.

\subsubsection{Example}\label{example-1}

\begin{verbatim}
--
--        SWITCH OFF ECHOING OF INPUT FILES
--
NOECHO
--
--        INCLUDE SIMULATION GRID WITH SLOPING FAULTS
--
INCLUDE
          './INCLUDE/GRID/IRAP_1005.GRDECL' /
--
--        SWITCH ON ECHOING OF INPUT FILES
--
ECHO
\end{verbatim}

The examples deactivates the echoing of the input files, reads in the
grid geometry data using the INCLUDE keyword, and then activates the
echoing of the input files again.

\begin{tcolorbox}[enhanced jigsaw, title=\textcolor{quarto-callout-note-color}{\faInfo}\hspace{0.5em}{Note}, coltitle=black, breakable, toprule=.15mm, colbacktitle=quarto-callout-note-color!10!white, leftrule=.75mm, opacitybacktitle=0.6, colback=white, bottomtitle=1mm, toptitle=1mm, titlerule=0mm, left=2mm, colframe=quarto-callout-note-color-frame, arc=.35mm, rightrule=.15mm, bottomrule=.15mm, opacityback=0]

Especially for the large voluminous data sets in the GRID section, it is
good practice to deactivate the echoing of the input files when loading
this data to avoid the print output file becoming too large to view in a
text editor.

\end{tcolorbox}

\subsection{END -- Define the End of the Input
File}\label{end-define-the-end-of-the-input-file}

\begin{longtable}[]{@{}
  >{\raggedright\arraybackslash}p{(\linewidth - 14\tabcolsep) * \real{0.1250}}
  >{\raggedright\arraybackslash}p{(\linewidth - 14\tabcolsep) * \real{0.1250}}
  >{\raggedright\arraybackslash}p{(\linewidth - 14\tabcolsep) * \real{0.1250}}
  >{\raggedright\arraybackslash}p{(\linewidth - 14\tabcolsep) * \real{0.1250}}
  >{\raggedright\arraybackslash}p{(\linewidth - 14\tabcolsep) * \real{0.1250}}
  >{\raggedright\arraybackslash}p{(\linewidth - 14\tabcolsep) * \real{0.1250}}
  >{\raggedright\arraybackslash}p{(\linewidth - 14\tabcolsep) * \real{0.1250}}
  >{\raggedright\arraybackslash}p{(\linewidth - 14\tabcolsep) * \real{0.1250}}@{}}
\toprule\noalign{}
\begin{minipage}[b]{\linewidth}\raggedright
RUNSPEC
\end{minipage} & \begin{minipage}[b]{\linewidth}\raggedright
GRID
\end{minipage} & \begin{minipage}[b]{\linewidth}\raggedright
EDIT
\end{minipage} & \begin{minipage}[b]{\linewidth}\raggedright
PROPS
\end{minipage} & \begin{minipage}[b]{\linewidth}\raggedright
REGIONS
\end{minipage} & \begin{minipage}[b]{\linewidth}\raggedright
SOLUTION
\end{minipage} & \begin{minipage}[b]{\linewidth}\raggedright
SUMMARY
\end{minipage} & \begin{minipage}[b]{\linewidth}\raggedright
SCHEDULE
\end{minipage} \\
\midrule\noalign{}
\endhead
\bottomrule\noalign{}
\endlastfoot
\end{longtable}

\subsubsection{Description}\label{description-3}

This keyword marks the end of the input file and can occur in any
section. Any keywords and data after the END keyword are ignored.

There is no data required for this keyword.

\subsubsection{Example}\label{example-2}

\begin{verbatim}
-- ------------------------------------------------------------------------------
-- SCHEDULE SECTION - 2006-01-01
-- ------------------------------------------------------------------------------
RPTSCHED
'WELLS=2'    'WELSPECS'    'CPU=2'     'FIP=2'                                 /

DATES
 1  JAN   2006  /
/

RPTSCHED
'NOTHING'                                                                      /

DATES
 1  APR   2006  /
 1  JLY   2006  /
 1  OCT   2006  /
/
       ECHO
--
-- ******************************************************************************
-- END OF FILE
-- ******************************************************************************
END
-- ------------------------------------------------------------------------------
-- SCHEDULE SECTION - 2007-01-01
-- ------------------------------------------------------------------------------
RPTSCHED
'WELLS=2'    'WELSPECS'    'CPU=2'     'FIP=2'                                 /

DATES
 1  JAN   2007  /
/
\end{verbatim}

In the above example OPM Flow will process the data up to October 1,
2006 only, and then start to run the simulation. All keywords after the
END file keyword will not be read or processed.

\subsection{ENDINC -- Define the End of an Include
File}\label{endinc-define-the-end-of-an-include-file}

\begin{longtable}[]{@{}
  >{\raggedright\arraybackslash}p{(\linewidth - 14\tabcolsep) * \real{0.1250}}
  >{\raggedright\arraybackslash}p{(\linewidth - 14\tabcolsep) * \real{0.1250}}
  >{\raggedright\arraybackslash}p{(\linewidth - 14\tabcolsep) * \real{0.1250}}
  >{\raggedright\arraybackslash}p{(\linewidth - 14\tabcolsep) * \real{0.1250}}
  >{\raggedright\arraybackslash}p{(\linewidth - 14\tabcolsep) * \real{0.1250}}
  >{\raggedright\arraybackslash}p{(\linewidth - 14\tabcolsep) * \real{0.1250}}
  >{\raggedright\arraybackslash}p{(\linewidth - 14\tabcolsep) * \real{0.1250}}
  >{\raggedright\arraybackslash}p{(\linewidth - 14\tabcolsep) * \real{0.1250}}@{}}
\toprule\noalign{}
\begin{minipage}[b]{\linewidth}\raggedright
RUNSPEC
\end{minipage} & \begin{minipage}[b]{\linewidth}\raggedright
GRID
\end{minipage} & \begin{minipage}[b]{\linewidth}\raggedright
EDIT
\end{minipage} & \begin{minipage}[b]{\linewidth}\raggedright
PROPS
\end{minipage} & \begin{minipage}[b]{\linewidth}\raggedright
REGIONS
\end{minipage} & \begin{minipage}[b]{\linewidth}\raggedright
SOLUTION
\end{minipage} & \begin{minipage}[b]{\linewidth}\raggedright
SUMMARY
\end{minipage} & \begin{minipage}[b]{\linewidth}\raggedright
SCHEDULE
\end{minipage} \\
\midrule\noalign{}
\endhead
\bottomrule\noalign{}
\endlastfoot
\end{longtable}

\subsubsection{Description}\label{description-4}

This keyword marks the end of an include file specified on the INCLUDE
keyword. When the ENDINC keyword is encountered in the INCLUDE file,
input data is read from the next keyword in the current file. Any
keywords and data after the ENDINC keyword in the INCLUDE file are
ignored.

There is no data required for this keyword.

\subsubsection{Example}\label{example-3}

\begin{verbatim}
-- ------------------------------------------------------------------------------
-- SCHEDULE SECTION - 2006-01-01
-- ------------------------------------------------------------------------------
RPTSCHED
'WELLS=2'    'WELSPECS'    'CPU=2'     'FIP=2'                                 /

DATES
 1  JAN   2006  /
/

RPTSCHED
'NOTHING'                                                                      /

DATES
 1  APR   2006  /
 1  JLY   2006  /
 1  OCT   2006  /
/
       ECHO
--
-- ******************************************************************************
-- END OF INCLUDE FILE PROCESSING
-- ******************************************************************************
ENDINC
-- ------------------------------------------------------------------------------
-- SCHEDULE SECTION - 2007-01-01
-- ------------------------------------------------------------------------------
RPTSCHED
'WELLS=2'    'WELSPECS'    'CPU=2'     'FIP=2'                                 /

DATES
 1  JAN   2007  /
/
\end{verbatim}

In the above example OPM Flow will process the data up to October 1,
2006 only, and return control to the file that called the INCLUDE
keyword, and then continue processing the input files. All keywords
after the ENDINC keyword in the INCLUDE FILE will not be read or
processed.

\subsection{ENDSKIP -- Deactivate Skipping of Keywords and Input
Data}\label{endskip-deactivate-skipping-of-keywords-and-input-data}

\begin{longtable}[]{@{}
  >{\raggedright\arraybackslash}p{(\linewidth - 14\tabcolsep) * \real{0.1250}}
  >{\raggedright\arraybackslash}p{(\linewidth - 14\tabcolsep) * \real{0.1250}}
  >{\raggedright\arraybackslash}p{(\linewidth - 14\tabcolsep) * \real{0.1250}}
  >{\raggedright\arraybackslash}p{(\linewidth - 14\tabcolsep) * \real{0.1250}}
  >{\raggedright\arraybackslash}p{(\linewidth - 14\tabcolsep) * \real{0.1250}}
  >{\raggedright\arraybackslash}p{(\linewidth - 14\tabcolsep) * \real{0.1250}}
  >{\raggedright\arraybackslash}p{(\linewidth - 14\tabcolsep) * \real{0.1250}}
  >{\raggedright\arraybackslash}p{(\linewidth - 14\tabcolsep) * \real{0.1250}}@{}}
\toprule\noalign{}
\begin{minipage}[b]{\linewidth}\raggedright
RUNSPEC
\end{minipage} & \begin{minipage}[b]{\linewidth}\raggedright
GRID
\end{minipage} & \begin{minipage}[b]{\linewidth}\raggedright
EDIT
\end{minipage} & \begin{minipage}[b]{\linewidth}\raggedright
PROPS
\end{minipage} & \begin{minipage}[b]{\linewidth}\raggedright
REGIONS
\end{minipage} & \begin{minipage}[b]{\linewidth}\raggedright
SOLUTION
\end{minipage} & \begin{minipage}[b]{\linewidth}\raggedright
SUMMARY
\end{minipage} & \begin{minipage}[b]{\linewidth}\raggedright
SCHEDULE
\end{minipage} \\
\midrule\noalign{}
\endhead
\bottomrule\noalign{}
\endlastfoot
\end{longtable}

\subsubsection{Description}\label{description-5}

The ENDSKIP keyword deactivates the skipping of keywords that was
activated by the SKIP, SKIP100 or SKIP300 keywords. Each SKIP, SKIP100
or SKIP300 keyword should be paired with an ENDSKIP keyword.

There is no data required for this keyword.

The nesting of pairs of either the SKIP, SKIP100, or SKIP300 keyword and
the ENDSKIP keyword is not supported by OPM Flow.

\subsubsection{Example}\label{example-4}

\begin{verbatim}
--
--       SWITCH ON SKIPPING OF KEYWORDS AND DATA
--
SKIP
--
--       INCLUDE SIMULATION GRID WITH SLOPING FAULTS
--
INCLUDE
         './INCLUDE/GRID/IRAP_1005.GRDECL' /
--
--       SWITCH ON READING OF KEYWORDS AND DATA
--
ENDSKIP
\end{verbatim}

The example skips reading of the grid geometry data input using the
INCLUDE keyword, and then reverts back to reading the remainder of the
input file.

\subsection{EXTRAPMS -- Activate Extrapolation Warning
Messages}\label{extrapms-activate-extrapolation-warning-messages}

\begin{longtable}[]{@{}
  >{\raggedright\arraybackslash}p{(\linewidth - 14\tabcolsep) * \real{0.1250}}
  >{\raggedright\arraybackslash}p{(\linewidth - 14\tabcolsep) * \real{0.1250}}
  >{\raggedright\arraybackslash}p{(\linewidth - 14\tabcolsep) * \real{0.1250}}
  >{\raggedright\arraybackslash}p{(\linewidth - 14\tabcolsep) * \real{0.1250}}
  >{\raggedright\arraybackslash}p{(\linewidth - 14\tabcolsep) * \real{0.1250}}
  >{\raggedright\arraybackslash}p{(\linewidth - 14\tabcolsep) * \real{0.1250}}
  >{\raggedright\arraybackslash}p{(\linewidth - 14\tabcolsep) * \real{0.1250}}
  >{\raggedright\arraybackslash}p{(\linewidth - 14\tabcolsep) * \real{0.1250}}@{}}
\toprule\noalign{}
\begin{minipage}[b]{\linewidth}\raggedright
RUNSPEC
\end{minipage} & \begin{minipage}[b]{\linewidth}\raggedright
GRID
\end{minipage} & \begin{minipage}[b]{\linewidth}\raggedright
EDIT
\end{minipage} & \begin{minipage}[b]{\linewidth}\raggedright
PROPS
\end{minipage} & \begin{minipage}[b]{\linewidth}\raggedright
REGIONS
\end{minipage} & \begin{minipage}[b]{\linewidth}\raggedright
SOLUTION
\end{minipage} & \begin{minipage}[b]{\linewidth}\raggedright
SUMMARY
\end{minipage} & \begin{minipage}[b]{\linewidth}\raggedright
SCHEDULE
\end{minipage} \\
\midrule\noalign{}
\endhead
\bottomrule\noalign{}
\endlastfoot
\end{longtable}

\subsubsection{Description}\label{description-6}

The EXTRAPMS keyword activates extrapolation warning messages for when
OPM Flow extrapolates the PVT or VFP tables. Frequent extrapolation
warning messages should be investigated and resolved as this would
indicate possible incorrect data and may result in the simulator
extrapolating to unrealistic values.

\begin{longtable}[]{@{}
  >{\raggedright\arraybackslash}p{(\linewidth - 6\tabcolsep) * \real{0.2500}}
  >{\raggedright\arraybackslash}p{(\linewidth - 6\tabcolsep) * \real{0.2500}}
  >{\raggedright\arraybackslash}p{(\linewidth - 6\tabcolsep) * \real{0.2500}}
  >{\raggedright\arraybackslash}p{(\linewidth - 6\tabcolsep) * \real{0.2500}}@{}}
\toprule\noalign{}
\begin{minipage}[b]{\linewidth}\raggedright
No.
\end{minipage} & \begin{minipage}[b]{\linewidth}\raggedright
Name
\end{minipage} & \begin{minipage}[b]{\linewidth}\raggedright
Description
\end{minipage} & \begin{minipage}[b]{\linewidth}\raggedright
Default
\end{minipage} \\
\midrule\noalign{}
\endhead
\bottomrule\noalign{}
\endlastfoot
1 & EXTRAP & Defines a single integer that activates the extrapolation
warning message options for PVT and VFP tables. EXTRAP can have the
following values: & 0 \\
Notes: & & & \\
\end{longtable}

\emph{Table 4.2: EXTRAPMS Keyword Description}

This keyword is not supported by OPM Flow but has no effect on the
results so it will be ignored.

\subsubsection{Example}\label{example-5}

\begin{verbatim}
--
--       ACTIVATE EXTRAPOLATION MESSAGES
--
EXTRAPMS
         2                                                                    /
\end{verbatim}

The above example activates the default the VFP table extrapolation
warnings option.

\subsection{FORMFEED -- Defined the Print File Form-Feed
Character}\label{formfeed-defined-the-print-file-form-feed-character}

\begin{longtable}[]{@{}
  >{\raggedright\arraybackslash}p{(\linewidth - 14\tabcolsep) * \real{0.1250}}
  >{\raggedright\arraybackslash}p{(\linewidth - 14\tabcolsep) * \real{0.1250}}
  >{\raggedright\arraybackslash}p{(\linewidth - 14\tabcolsep) * \real{0.1250}}
  >{\raggedright\arraybackslash}p{(\linewidth - 14\tabcolsep) * \real{0.1250}}
  >{\raggedright\arraybackslash}p{(\linewidth - 14\tabcolsep) * \real{0.1250}}
  >{\raggedright\arraybackslash}p{(\linewidth - 14\tabcolsep) * \real{0.1250}}
  >{\raggedright\arraybackslash}p{(\linewidth - 14\tabcolsep) * \real{0.1250}}
  >{\raggedright\arraybackslash}p{(\linewidth - 14\tabcolsep) * \real{0.1250}}@{}}
\toprule\noalign{}
\begin{minipage}[b]{\linewidth}\raggedright
RUNSPEC
\end{minipage} & \begin{minipage}[b]{\linewidth}\raggedright
GRID
\end{minipage} & \begin{minipage}[b]{\linewidth}\raggedright
EDIT
\end{minipage} & \begin{minipage}[b]{\linewidth}\raggedright
PROPS
\end{minipage} & \begin{minipage}[b]{\linewidth}\raggedright
REGIONS
\end{minipage} & \begin{minipage}[b]{\linewidth}\raggedright
SOLUTION
\end{minipage} & \begin{minipage}[b]{\linewidth}\raggedright
SUMMARY
\end{minipage} & \begin{minipage}[b]{\linewidth}\raggedright
SCHEDULE
\end{minipage} \\
\midrule\noalign{}
\endhead
\bottomrule\noalign{}
\endlastfoot
\end{longtable}

\subsubsection{Description}\label{description-7}

The FORMFEED keyword defines the form-feed character, or carriage
control character, for the output print (\emph{.PRT) run summary (}.RSM)
files. The keyword should be place at the very top of the input file.

\begin{longtable}[]{@{}
  >{\raggedright\arraybackslash}p{(\linewidth - 6\tabcolsep) * \real{0.2500}}
  >{\raggedright\arraybackslash}p{(\linewidth - 6\tabcolsep) * \real{0.2500}}
  >{\raggedright\arraybackslash}p{(\linewidth - 6\tabcolsep) * \real{0.2500}}
  >{\raggedright\arraybackslash}p{(\linewidth - 6\tabcolsep) * \real{0.2500}}@{}}
\toprule\noalign{}
\begin{minipage}[b]{\linewidth}\raggedright
No.
\end{minipage} & \begin{minipage}[b]{\linewidth}\raggedright
Name
\end{minipage} & \begin{minipage}[b]{\linewidth}\raggedright
Description
\end{minipage} & \begin{minipage}[b]{\linewidth}\raggedright
Default
\end{minipage} \\
\midrule\noalign{}
\endhead
\bottomrule\noalign{}
\endlastfoot
1 & FORMFEED & Defines a single integer that defines the carriage
control character activates, and should be set to: & 0 \\
Notes: & & & \\
\end{longtable}

\emph{Table 4.3: FORMFEED Keyword Description}

This keyword is not supported by OPM Flow but would change the results
if supported so the simulation will be stopped.

\subsubsection{Example}\label{example-6}

\begin{verbatim}
--
--       ACTIVATE EXTRAPOLATION MESSAGES
--
FORMFEED
         3                                                                      /
\end{verbatim}

The above example sets the carriage return character to no form-feed
character.

\subsection{INCLUDE -- Load Another Data File at the Current
Position}\label{include-load-another-data-file-at-the-current-position}

\begin{longtable}[]{@{}
  >{\raggedright\arraybackslash}p{(\linewidth - 14\tabcolsep) * \real{0.1250}}
  >{\raggedright\arraybackslash}p{(\linewidth - 14\tabcolsep) * \real{0.1250}}
  >{\raggedright\arraybackslash}p{(\linewidth - 14\tabcolsep) * \real{0.1250}}
  >{\raggedright\arraybackslash}p{(\linewidth - 14\tabcolsep) * \real{0.1250}}
  >{\raggedright\arraybackslash}p{(\linewidth - 14\tabcolsep) * \real{0.1250}}
  >{\raggedright\arraybackslash}p{(\linewidth - 14\tabcolsep) * \real{0.1250}}
  >{\raggedright\arraybackslash}p{(\linewidth - 14\tabcolsep) * \real{0.1250}}
  >{\raggedright\arraybackslash}p{(\linewidth - 14\tabcolsep) * \real{0.1250}}@{}}
\toprule\noalign{}
\begin{minipage}[b]{\linewidth}\raggedright
RUNSPEC
\end{minipage} & \begin{minipage}[b]{\linewidth}\raggedright
GRID
\end{minipage} & \begin{minipage}[b]{\linewidth}\raggedright
EDIT
\end{minipage} & \begin{minipage}[b]{\linewidth}\raggedright
PROPS
\end{minipage} & \begin{minipage}[b]{\linewidth}\raggedright
REGIONS
\end{minipage} & \begin{minipage}[b]{\linewidth}\raggedright
SOLUTION
\end{minipage} & \begin{minipage}[b]{\linewidth}\raggedright
SUMMARY
\end{minipage} & \begin{minipage}[b]{\linewidth}\raggedright
SCHEDULE
\end{minipage} \\
\midrule\noalign{}
\endhead
\bottomrule\noalign{}
\endlastfoot
\end{longtable}

\subsubsection{Description}\label{description-8}

The INCLUDE keyword informs OPM Flow to continue reading input data from
the specified INCLUDE file. When the end of the INCLUDE file is reached,
or the ENDINC keyword is encountered in the included file, input data is
read from the next keyword in the current file. Although INCLUDE files
can be nested, that is INCLUDE files within INCLUDE files etc., in
practice this should be avoided due to the complexity of tracking the
files.

\begin{longtable}[]{@{}
  >{\raggedright\arraybackslash}p{(\linewidth - 6\tabcolsep) * \real{0.2500}}
  >{\raggedright\arraybackslash}p{(\linewidth - 6\tabcolsep) * \real{0.2500}}
  >{\raggedright\arraybackslash}p{(\linewidth - 6\tabcolsep) * \real{0.2500}}
  >{\raggedright\arraybackslash}p{(\linewidth - 6\tabcolsep) * \real{0.2500}}@{}}
\toprule\noalign{}
\begin{minipage}[b]{\linewidth}\raggedright
No.
\end{minipage} & \begin{minipage}[b]{\linewidth}\raggedright
Name
\end{minipage} & \begin{minipage}[b]{\linewidth}\raggedright
Description
\end{minipage} & \begin{minipage}[b]{\linewidth}\raggedright
Default
\end{minipage} \\
\midrule\noalign{}
\endhead
\bottomrule\noalign{}
\endlastfoot
1 & FILENAME & A character string enclosed in quotes that defines a file
to read in and be processed by OPM Flow. & None \\
Notes: & & & \\
\end{longtable}

\emph{Table 4.4: INCLUDE Keyword Description}

\subsubsection{Examples}\label{examples}

The first example shown below loads the grid file from the same
directory as the data file.

\begin{verbatim}
--
--       LOAD INCLUDE FILE
--
INCLUDE
         'NOR-OPM-A00-GRID.inc'  /
\end{verbatim}

The next example loads the same file one directory above from where the
data file is located.

\begin{verbatim}
--
--       LOAD INCLUDE FILE
--
INCLUDE
         '../NOR-OPM-A00-FAULTS.inc'  /
\end{verbatim}

The final example loads the same file from a separate include directory
found in the parent directory relative to where the data file is
located.

\begin{verbatim}
--
--       LOAD INCLUDE FILE
--
INCLUDE
         '../INCLUDE/NOR-OPM-A00-FAULTS.inc'  /
\end{verbatim}

\subsection{MESSAGE -- Output User
Message}\label{message-output-user-message}

\begin{longtable}[]{@{}
  >{\raggedright\arraybackslash}p{(\linewidth - 14\tabcolsep) * \real{0.1250}}
  >{\raggedright\arraybackslash}p{(\linewidth - 14\tabcolsep) * \real{0.1250}}
  >{\raggedright\arraybackslash}p{(\linewidth - 14\tabcolsep) * \real{0.1250}}
  >{\raggedright\arraybackslash}p{(\linewidth - 14\tabcolsep) * \real{0.1250}}
  >{\raggedright\arraybackslash}p{(\linewidth - 14\tabcolsep) * \real{0.1250}}
  >{\raggedright\arraybackslash}p{(\linewidth - 14\tabcolsep) * \real{0.1250}}
  >{\raggedright\arraybackslash}p{(\linewidth - 14\tabcolsep) * \real{0.1250}}
  >{\raggedright\arraybackslash}p{(\linewidth - 14\tabcolsep) * \real{0.1250}}@{}}
\toprule\noalign{}
\begin{minipage}[b]{\linewidth}\raggedright
RUNSPEC
\end{minipage} & \begin{minipage}[b]{\linewidth}\raggedright
GRID
\end{minipage} & \begin{minipage}[b]{\linewidth}\raggedright
EDIT
\end{minipage} & \begin{minipage}[b]{\linewidth}\raggedright
PROPS
\end{minipage} & \begin{minipage}[b]{\linewidth}\raggedright
REGIONS
\end{minipage} & \begin{minipage}[b]{\linewidth}\raggedright
SOLUTION
\end{minipage} & \begin{minipage}[b]{\linewidth}\raggedright
SUMMARY
\end{minipage} & \begin{minipage}[b]{\linewidth}\raggedright
SCHEDULE
\end{minipage} \\
\midrule\noalign{}
\endhead
\bottomrule\noalign{}
\endlastfoot
\end{longtable}

\subsubsection{Description}\label{description-9}

The MESSAGE keyword outputs a user message to the terminal, as well as
to the print (\emph{.PRT) and debug (}.DBG) files. Note this is
different to the MESSAGES keyword, that defines OPM Flows message print
limits and stop limits generated by the simulator.

This keyword is not supported by OPM Flow but has no effect on the
results so it will be ignored.

\subsection{MESSAGES -- Define Message Print Limits and Stop
Limits}\label{messages-define-message-print-limits-and-stop-limits}

\begin{longtable}[]{@{}
  >{\raggedright\arraybackslash}p{(\linewidth - 14\tabcolsep) * \real{0.1250}}
  >{\raggedright\arraybackslash}p{(\linewidth - 14\tabcolsep) * \real{0.1250}}
  >{\raggedright\arraybackslash}p{(\linewidth - 14\tabcolsep) * \real{0.1250}}
  >{\raggedright\arraybackslash}p{(\linewidth - 14\tabcolsep) * \real{0.1250}}
  >{\raggedright\arraybackslash}p{(\linewidth - 14\tabcolsep) * \real{0.1250}}
  >{\raggedright\arraybackslash}p{(\linewidth - 14\tabcolsep) * \real{0.1250}}
  >{\raggedright\arraybackslash}p{(\linewidth - 14\tabcolsep) * \real{0.1250}}
  >{\raggedright\arraybackslash}p{(\linewidth - 14\tabcolsep) * \real{0.1250}}@{}}
\toprule\noalign{}
\begin{minipage}[b]{\linewidth}\raggedright
RUNSPEC
\end{minipage} & \begin{minipage}[b]{\linewidth}\raggedright
GRID
\end{minipage} & \begin{minipage}[b]{\linewidth}\raggedright
EDIT
\end{minipage} & \begin{minipage}[b]{\linewidth}\raggedright
PROPS
\end{minipage} & \begin{minipage}[b]{\linewidth}\raggedright
REGIONS
\end{minipage} & \begin{minipage}[b]{\linewidth}\raggedright
SOLUTION
\end{minipage} & \begin{minipage}[b]{\linewidth}\raggedright
SUMMARY
\end{minipage} & \begin{minipage}[b]{\linewidth}\raggedright
SCHEDULE
\end{minipage} \\
\midrule\noalign{}
\endhead
\bottomrule\noalign{}
\endlastfoot
\end{longtable}

\subsubsection{Description}\label{description-10}

The MESSAGES keyword defines the print and stops levels for various
messages. The ``print limits'' set the maximum number of messages that
will be printed, after which no more messages will be printed and the
``stop limits'' terminate the run when these limits are exceeded. There
are six levels of message that increase in severity from informative all
the way to programming errors, as outlined in Table 4.5.

\begin{longtable}[]{@{}
  >{\raggedright\arraybackslash}p{(\linewidth - 6\tabcolsep) * \real{0.2500}}
  >{\raggedright\arraybackslash}p{(\linewidth - 6\tabcolsep) * \real{0.2500}}
  >{\raggedright\arraybackslash}p{(\linewidth - 6\tabcolsep) * \real{0.2500}}
  >{\raggedright\arraybackslash}p{(\linewidth - 6\tabcolsep) * \real{0.2500}}@{}}
\toprule\noalign{}
\begin{minipage}[b]{\linewidth}\raggedright
No.
\end{minipage} & \begin{minipage}[b]{\linewidth}\raggedright
Name
\end{minipage} & \begin{minipage}[b]{\linewidth}\raggedright
Description
\end{minipage} & \begin{minipage}[b]{\linewidth}\raggedright
Default
\end{minipage} \\
\midrule\noalign{}
\endhead
\bottomrule\noalign{}
\endlastfoot
1 & PRTSMESG & An integer defining the maximum number of MESSAGE type
messages to be printed after which this type of message stops printing.
& 1,000,000 \\
2 & PRTSCOMT & An integer defining the maximum number of COMMENT type
messages be printed after which this type of message stops printing. &
1,000,000 \\
3 & PRTSWARN & An integer defining the maximum number of WARNING type
messages be printed after which this type of message stops printing. &
10,000 \\
4 & PRTSPROB & An integer defining the maximum number of PROBLEM type
messages to be printed after which this type of message stops printing.
& 100 \\
5 & PRTSERRS & An integer defining the maximum number of ERROR type
messages to be printed after which this type of message stops printing.
& 100 \\
6 & PRTSBUGS & An integer defining the maximum number of BUG type
messages to be printed after which this type of message stops printing.
& 100 \\
7 & STOPMESG & An integer defining the maximum number of MESSAGE type
messages to be printed after which OPM Flow terminates the run. Not used
by OPM Flow. & 1,000,000 \\
8 & STOPCOMT & An integer defining the maximum number of COMMENT type
messages to be printed after which OPM Flow terminates the run. Not used
by OPM Flow. & 1,000,000 \\
9 & STOPWARN & An integer defining the maximum number of WARNING type
messages to be printed after which OPM Flow terminates the run. Not used
by OPM Flow. & 10,000 \\
10 & STOPPROB & An integer defining the maximum number of PROBLEM type
messages to be printed after which OPM Flow terminates the run. Not used
by OPM Flow. & 100 \\
11 & STOPERRS & An integer defining the maximum number of ERROR type
messages to be printed after which OPM Flow terminates the run. Not used
by OPM Flow. & 10 \\
12 & STOPBUGS & An integer defining the maximum number of BUG type
messages to be printed after which OPM Flow terminates the run. Not used
by OPM Flow. & 1 \\
13 & PRTGRPMS & An integer defining the maximum number of GROUP MESSAGE
type messages to be printed after which this type of message stops
printing. Not used by OPM Flow. & 10 \\
Notes: & & & \\
\end{longtable}

\emph{Table 4.5: MESSAGES Keyword Description}

\subsubsection{Example}\label{example-7}

\begin{verbatim}
--
--       MESS  COMMT WARN  PROBL ERROR BUG   MESS COMMT WARN  PROBL ERROR BUG
--       LIMIT LIMIT LIMIT LIMIT LIMIT LIMIT STOP STOP  STOP  STOP  STOP  STOP
MESSAGES
         1*    1*    1*    1500  1*    1*    1*   1*    1*    1000  1*    1*   /
\end{verbatim}

The above example sets the PROBLEM type message print limit to 1,500 and
the stop limit to 1,000.

\subsection{NOECHO -- Deactivate Echoing of User Input Files to the
Print
File}\label{noecho-deactivate-echoing-of-user-input-files-to-the-print-file}

\begin{longtable}[]{@{}
  >{\raggedright\arraybackslash}p{(\linewidth - 14\tabcolsep) * \real{0.1250}}
  >{\raggedright\arraybackslash}p{(\linewidth - 14\tabcolsep) * \real{0.1250}}
  >{\raggedright\arraybackslash}p{(\linewidth - 14\tabcolsep) * \real{0.1250}}
  >{\raggedright\arraybackslash}p{(\linewidth - 14\tabcolsep) * \real{0.1250}}
  >{\raggedright\arraybackslash}p{(\linewidth - 14\tabcolsep) * \real{0.1250}}
  >{\raggedright\arraybackslash}p{(\linewidth - 14\tabcolsep) * \real{0.1250}}
  >{\raggedright\arraybackslash}p{(\linewidth - 14\tabcolsep) * \real{0.1250}}
  >{\raggedright\arraybackslash}p{(\linewidth - 14\tabcolsep) * \real{0.1250}}@{}}
\toprule\noalign{}
\begin{minipage}[b]{\linewidth}\raggedright
RUNSPEC
\end{minipage} & \begin{minipage}[b]{\linewidth}\raggedright
GRID
\end{minipage} & \begin{minipage}[b]{\linewidth}\raggedright
EDIT
\end{minipage} & \begin{minipage}[b]{\linewidth}\raggedright
PROPS
\end{minipage} & \begin{minipage}[b]{\linewidth}\raggedright
REGIONS
\end{minipage} & \begin{minipage}[b]{\linewidth}\raggedright
SOLUTION
\end{minipage} & \begin{minipage}[b]{\linewidth}\raggedright
SUMMARY
\end{minipage} & \begin{minipage}[b]{\linewidth}\raggedright
SCHEDULE
\end{minipage} \\
\midrule\noalign{}
\endhead
\bottomrule\noalign{}
\endlastfoot
\end{longtable}

\subsubsection{Description}\label{description-11}

Turns off echoing of all the input files to the print file. Note by
default echoing of the inputs files is active, but can subsequently be
switched off by the NOECHO activation keyword.

There is no data required for this keyword.

This keyword is not supported by OPM Flow but has no effect on the
results so it will be ignored.

\subsubsection{Example}\label{example-8}

\begin{verbatim}
--
--       SWITCH OFF ECHOING OF INPUT FILES
--
NOECHO
--
--       INCLUDE SIMULATION GRID WITH SLOPING FAULTS
--
INCLUDE
         './INCLUDE/GRID/IRAP_1005.GRDECL' /
--
--       SWITCH ON ECHOING OF INPUT FILES
--
ECHO
\end{verbatim}

The example deactivates the echoing of the input files, reads in the
grid geometry data using the INCLUDE keyword, and then activates the
echoing of the input files again.

\begin{tcolorbox}[enhanced jigsaw, title=\textcolor{quarto-callout-note-color}{\faInfo}\hspace{0.5em}{Note}, coltitle=black, breakable, toprule=.15mm, colbacktitle=quarto-callout-note-color!10!white, leftrule=.75mm, opacitybacktitle=0.6, colback=white, bottomtitle=1mm, toptitle=1mm, titlerule=0mm, left=2mm, colframe=quarto-callout-note-color-frame, arc=.35mm, rightrule=.15mm, bottomrule=.15mm, opacityback=0]

Especially for the large voluminous data sets in the GRID section, it is
good practice to deactivate the echoing of the input files when loading
this data to avoid the print output file becoming too large to view in a
text editor.

\end{tcolorbox}

\subsection{NOWARN -- Deactivate Warning
Messages}\label{nowarn-deactivate-warning-messages}

\begin{longtable}[]{@{}
  >{\raggedright\arraybackslash}p{(\linewidth - 14\tabcolsep) * \real{0.1250}}
  >{\raggedright\arraybackslash}p{(\linewidth - 14\tabcolsep) * \real{0.1250}}
  >{\raggedright\arraybackslash}p{(\linewidth - 14\tabcolsep) * \real{0.1250}}
  >{\raggedright\arraybackslash}p{(\linewidth - 14\tabcolsep) * \real{0.1250}}
  >{\raggedright\arraybackslash}p{(\linewidth - 14\tabcolsep) * \real{0.1250}}
  >{\raggedright\arraybackslash}p{(\linewidth - 14\tabcolsep) * \real{0.1250}}
  >{\raggedright\arraybackslash}p{(\linewidth - 14\tabcolsep) * \real{0.1250}}
  >{\raggedright\arraybackslash}p{(\linewidth - 14\tabcolsep) * \real{0.1250}}@{}}
\toprule\noalign{}
\begin{minipage}[b]{\linewidth}\raggedright
RUNSPEC
\end{minipage} & \begin{minipage}[b]{\linewidth}\raggedright
GRID
\end{minipage} & \begin{minipage}[b]{\linewidth}\raggedright
EDIT
\end{minipage} & \begin{minipage}[b]{\linewidth}\raggedright
PROPS
\end{minipage} & \begin{minipage}[b]{\linewidth}\raggedright
REGIONS
\end{minipage} & \begin{minipage}[b]{\linewidth}\raggedright
SOLUTION
\end{minipage} & \begin{minipage}[b]{\linewidth}\raggedright
SUMMARY
\end{minipage} & \begin{minipage}[b]{\linewidth}\raggedright
SCHEDULE
\end{minipage} \\
\midrule\noalign{}
\endhead
\bottomrule\noalign{}
\endlastfoot
\end{longtable}

\subsubsection{Description}\label{description-12}

Turns off warning messages to be printed to the print file; note that
this keyword is deactivated by default and can subsequently be switched
off by the WARN activation keyword. The warning messages may be turned
on and off using keywords WARN and NOWARN.

It is recommended that WARN should always be used and action taken if
necessary. For subsequent runs, the warning messages can be turned off.

There is no data required for this keyword.

This keyword is not supported by OPM Flow but has no effect on the
results so it will be ignored.

\subsubsection{Example}\label{example-9}

\begin{verbatim}
--
--      SWITCH OFF WARNING MESSAGES
NOWARN
--
--      INCLUDE SIMULATION GRID WITH SLOPING FAULTS
--
INCLUDE
     './INCLUDE/GRID/IRAP_1005.GRDECL' /
--
--      SWITCH ON WARNING MESSAGES
--
WARN
\end{verbatim}

The example deactivates the warning messages before reading the grid
geometry data using the INCLUDE keyword, and then activates the warning
messages after reading the INCLUDE file.

\subsection{SKIP -- Activate Skipping of All Keywords and Input
Data}\label{skip-activate-skipping-of-all-keywords-and-input-data}

\begin{longtable}[]{@{}
  >{\raggedright\arraybackslash}p{(\linewidth - 14\tabcolsep) * \real{0.1250}}
  >{\raggedright\arraybackslash}p{(\linewidth - 14\tabcolsep) * \real{0.1250}}
  >{\raggedright\arraybackslash}p{(\linewidth - 14\tabcolsep) * \real{0.1250}}
  >{\raggedright\arraybackslash}p{(\linewidth - 14\tabcolsep) * \real{0.1250}}
  >{\raggedright\arraybackslash}p{(\linewidth - 14\tabcolsep) * \real{0.1250}}
  >{\raggedright\arraybackslash}p{(\linewidth - 14\tabcolsep) * \real{0.1250}}
  >{\raggedright\arraybackslash}p{(\linewidth - 14\tabcolsep) * \real{0.1250}}
  >{\raggedright\arraybackslash}p{(\linewidth - 14\tabcolsep) * \real{0.1250}}@{}}
\toprule\noalign{}
\begin{minipage}[b]{\linewidth}\raggedright
RUNSPEC
\end{minipage} & \begin{minipage}[b]{\linewidth}\raggedright
GRID
\end{minipage} & \begin{minipage}[b]{\linewidth}\raggedright
EDIT
\end{minipage} & \begin{minipage}[b]{\linewidth}\raggedright
PROPS
\end{minipage} & \begin{minipage}[b]{\linewidth}\raggedright
REGIONS
\end{minipage} & \begin{minipage}[b]{\linewidth}\raggedright
SOLUTION
\end{minipage} & \begin{minipage}[b]{\linewidth}\raggedright
SUMMARY
\end{minipage} & \begin{minipage}[b]{\linewidth}\raggedright
SCHEDULE
\end{minipage} \\
\midrule\noalign{}
\endhead
\bottomrule\noalign{}
\endlastfoot
\end{longtable}

\subsubsection{Description}\label{description-13}

The SKIP keyword activates skipping of all keywords and input data until
the ENDSKIP keyword is encountered. All keywords between the SKIP and
ENDSKIP keywords are ignored. Each SKIP keyword should be paired with an
ENDSKIP keyword. See also the SKIP100 and SKIP300 keywords.

There is no data required for this keyword.

The nesting of pairs of the SKIP, SKIP100, or SKIP300 keywords and the
ENDSKIP keyword is not supported by OPM Flow.

\subsubsection{Example}\label{example-10}

\begin{verbatim}
--
--       SWITCH ON SKIPPING OF ALL KEYWORDS AND DATA
--
SKIP
--
--       INCLUDE SIMULATION GRID WITH SLOPING FAULTS
--
INCLUDE
         './INCLUDE/GRID/IRAP_1005.GRDECL' /
--
--       SWITCH ON READING OF ALL KEYWORDS AND DATA
--
ENDSKIP
\end{verbatim}

The example skips reading of the grid geometry data input using the
INCLUDE keyword, and then reverts back to reading the remainder of the
input file.

\subsection{SKIP100 -- Activate Skipping of Keywords by Black-Oil
Simulator}\label{skip100-activate-skipping-of-keywords-by-black-oil-simulator}

\begin{longtable}[]{@{}
  >{\raggedright\arraybackslash}p{(\linewidth - 14\tabcolsep) * \real{0.1250}}
  >{\raggedright\arraybackslash}p{(\linewidth - 14\tabcolsep) * \real{0.1250}}
  >{\raggedright\arraybackslash}p{(\linewidth - 14\tabcolsep) * \real{0.1250}}
  >{\raggedright\arraybackslash}p{(\linewidth - 14\tabcolsep) * \real{0.1250}}
  >{\raggedright\arraybackslash}p{(\linewidth - 14\tabcolsep) * \real{0.1250}}
  >{\raggedright\arraybackslash}p{(\linewidth - 14\tabcolsep) * \real{0.1250}}
  >{\raggedright\arraybackslash}p{(\linewidth - 14\tabcolsep) * \real{0.1250}}
  >{\raggedright\arraybackslash}p{(\linewidth - 14\tabcolsep) * \real{0.1250}}@{}}
\toprule\noalign{}
\begin{minipage}[b]{\linewidth}\raggedright
RUNSPEC
\end{minipage} & \begin{minipage}[b]{\linewidth}\raggedright
GRID
\end{minipage} & \begin{minipage}[b]{\linewidth}\raggedright
EDIT
\end{minipage} & \begin{minipage}[b]{\linewidth}\raggedright
PROPS
\end{minipage} & \begin{minipage}[b]{\linewidth}\raggedright
REGIONS
\end{minipage} & \begin{minipage}[b]{\linewidth}\raggedright
SOLUTION
\end{minipage} & \begin{minipage}[b]{\linewidth}\raggedright
SUMMARY
\end{minipage} & \begin{minipage}[b]{\linewidth}\raggedright
SCHEDULE
\end{minipage} \\
\midrule\noalign{}
\endhead
\bottomrule\noalign{}
\endlastfoot
\end{longtable}

\subsubsection{Description}\label{description-14}

The SKIP100 keyword activates skipping of all keywords and input data by
the commercial black-oil simulator until the ENDSKIP keyword is
encountered. All keywords between the SKIP100 and ENDSKIP keywords are
ignored by the commercial black-oil simulator. The SKIP100 keyword is
ignored by the commercial compositional simulator. Each SKIP100 keyword
should be paired with an ENDSKIP keyword. See also the SKIP and SKIP300
keywords.

There is no data required for this keyword.

By default OPM Flow behaves like the commercial black-oil simulator.
However, this behaviour can be changed using the --input-skip-mode
command line parameter.

The nesting of pairs of either the SKIP, SKIP100, or SKIP300 keyword and
the ENDSKIP keyword is not supported by OPM Flow.

\subsubsection{Example}\label{example-11}

If the following example is read using the commercial black-oil
simulator then the SKIP300 keyword will be ignored and the black-oil PVT
data will be read, the SKIP100 keyword will cause the compositional PVT
data to be ignored, and then the remainder of the input will be read
after the next ENDSKIP keyword is encountered.

\begin{verbatim}
--
--       ACTIVATE SKIPPING OF KEYWORDS BY COMPOSITIONAL SIMULATOR
--
SKIP300
--
--       INCLUDE BLACK-OIL PVT DATA
--
INCLUDE
         './INCLUDE/’BLACK-OIL-PVT'         /
--
--       DEACTIVATE SKIPPING OF KEYWORDS
--
ENDSKIP
--
--       ACTIVATE SKIPPING OF KEYWORDS BY BLACK-OIL SIMULATOR
--
SKIP100
--
--       INCLUDE COMPOSITIONAL PVT DATA
--
INCLUDE
         './INCLUDE/’COMPOSITION-PVT-EOS'   /
--
--       SWITCH ON READING OF ALL KEYWORDS AND DATA
--
ENDSKIP
--
--       WATER PVT TABLE
--
PVTW
--       REF PRES  BW        CW        VISC     VISC
--       PSIA      RB/STB   1/PSIA    CPOISE   GRAD
--       --------  ------   -------   ------   ------
          4840.0    1.019    2.7E-6    0.370    1*         / WATER DATA REGION 1
--
--       OIL      WAT        GAS
--       DENSITY  DENSITY    DENSITY
--       -------  -------    -------
DENSITY
          39.0     62.37     0.04520                       / PVT DATA REGION 1
--       ROCK COMPRESSIBILITY
--
--       REF PRES  CF
--       PSIA      1/PSIA
--       --------  -------
ROCK
          3966.9   5.0E-06                                 / ROCK COMPRESSIBILITY
\end{verbatim}

\subsection{SKIP300 -- Activate Skipping of Keywords by Compositional
Simulator}\label{skip300-activate-skipping-of-keywords-by-compositional-simulator}

\begin{longtable}[]{@{}
  >{\raggedright\arraybackslash}p{(\linewidth - 14\tabcolsep) * \real{0.1250}}
  >{\raggedright\arraybackslash}p{(\linewidth - 14\tabcolsep) * \real{0.1250}}
  >{\raggedright\arraybackslash}p{(\linewidth - 14\tabcolsep) * \real{0.1250}}
  >{\raggedright\arraybackslash}p{(\linewidth - 14\tabcolsep) * \real{0.1250}}
  >{\raggedright\arraybackslash}p{(\linewidth - 14\tabcolsep) * \real{0.1250}}
  >{\raggedright\arraybackslash}p{(\linewidth - 14\tabcolsep) * \real{0.1250}}
  >{\raggedright\arraybackslash}p{(\linewidth - 14\tabcolsep) * \real{0.1250}}
  >{\raggedright\arraybackslash}p{(\linewidth - 14\tabcolsep) * \real{0.1250}}@{}}
\toprule\noalign{}
\begin{minipage}[b]{\linewidth}\raggedright
RUNSPEC
\end{minipage} & \begin{minipage}[b]{\linewidth}\raggedright
GRID
\end{minipage} & \begin{minipage}[b]{\linewidth}\raggedright
EDIT
\end{minipage} & \begin{minipage}[b]{\linewidth}\raggedright
PROPS
\end{minipage} & \begin{minipage}[b]{\linewidth}\raggedright
REGIONS
\end{minipage} & \begin{minipage}[b]{\linewidth}\raggedright
SOLUTION
\end{minipage} & \begin{minipage}[b]{\linewidth}\raggedright
SUMMARY
\end{minipage} & \begin{minipage}[b]{\linewidth}\raggedright
SCHEDULE
\end{minipage} \\
\midrule\noalign{}
\endhead
\bottomrule\noalign{}
\endlastfoot
\end{longtable}

\subsubsection{Description}\label{description-15}

The SKIP300 keyword activates skipping of all keywords and input data by
the commercial compositional simulator until the ENDSKIP keyword is
encountered. All keywords between the SKIP300 and ENDSKIP keywords are
ignored by the commercial compositional simulator. The SKIP300 keyword
is ignored by the commercial black-oil simulator. Each SKIP300 keyword
should be paired with an ENDSKIP keyword. See also the SKIP and SKIP100
keywords.

There is no data required for this keyword.

By default OPM Flow behaves like the commercial black-oil simulator.
However, this behaviour can be changed using the --input-skip-mode
command line parameter.

The nesting of pairs of either the SKIP, SKIP100, or SKIP300 keyword and
the ENDSKIP keyword is not supported by OPM Flow.

\subsubsection{Example}\label{example-12}

If the following example is read using the commercial compositional
simulator then the SKIP300 keyword will cause the black-oil PVT data to
be ignored, the input will continue to be read after the next ENDSKIP
keyword is encountered, and the SKIP100 keyword will be ignored and the
compositional PVT data will be read along with the remainder of the
input.

\begin{verbatim}
--
--       ACTIVATE SKIPPING OF KEYWORDS BY COMPOSITIONAL SIMULATOR
--
SKIP300
--
--       INCLUDE BLACK-OIL PVT DATA
--
INCLUDE
         './INCLUDE/’BLACK-OIL-PVT'         /
--
--       DEACTIVATE SKIPPING OF KEYWORDS
--
ENDSKIP
--
--       ACTIVATE SKIPPING OF KEYWORDS BY BLACK-OIL SIMULATOR
--
SKIP100
--
--       INCLUDE COMPOSITIONAL PVT DATA
--
INCLUDE
         './INCLUDE/’COMPOSITION-PVT-EOS'   /
--
--       SWITCH ON READING OF ALL KEYWORDS AND DATA
--
ENDSKIP
--
--       WATER PVT TABLE
--
PVTW
--       REF PRES  BW        CW        VISC     VISC
--       PSIA      RB/STB   1/PSIA    CPOISE   GRAD
--       --------  ------   -------   ------   ------
          4840.0    1.019    2.7E-6    0.370    1*         / WATER DATA REGION 1
--
--       OIL      WAT        GAS
--       DENSITY  DENSITY    DENSITY
--       -------  -------    -------
DENSITY
          39.0     62.37     0.04520                       / PVT DATA REGION 1
--       ROCK COMPRESSIBILITY
--
--       REF PRES  CF
--       PSIA      1/PSIA
--       --------  -------
ROCK
          3966.9   5.0E-06                                 / ROCK COMPRESSIBILITY
\end{verbatim}

\subsection{WARN -- Activate Warning
Messages}\label{warn-activate-warning-messages}

\begin{longtable}[]{@{}
  >{\raggedright\arraybackslash}p{(\linewidth - 14\tabcolsep) * \real{0.1250}}
  >{\raggedright\arraybackslash}p{(\linewidth - 14\tabcolsep) * \real{0.1250}}
  >{\raggedright\arraybackslash}p{(\linewidth - 14\tabcolsep) * \real{0.1250}}
  >{\raggedright\arraybackslash}p{(\linewidth - 14\tabcolsep) * \real{0.1250}}
  >{\raggedright\arraybackslash}p{(\linewidth - 14\tabcolsep) * \real{0.1250}}
  >{\raggedright\arraybackslash}p{(\linewidth - 14\tabcolsep) * \real{0.1250}}
  >{\raggedright\arraybackslash}p{(\linewidth - 14\tabcolsep) * \real{0.1250}}
  >{\raggedright\arraybackslash}p{(\linewidth - 14\tabcolsep) * \real{0.1250}}@{}}
\toprule\noalign{}
\begin{minipage}[b]{\linewidth}\raggedright
RUNSPEC
\end{minipage} & \begin{minipage}[b]{\linewidth}\raggedright
GRID
\end{minipage} & \begin{minipage}[b]{\linewidth}\raggedright
EDIT
\end{minipage} & \begin{minipage}[b]{\linewidth}\raggedright
PROPS
\end{minipage} & \begin{minipage}[b]{\linewidth}\raggedright
REGIONS
\end{minipage} & \begin{minipage}[b]{\linewidth}\raggedright
SOLUTION
\end{minipage} & \begin{minipage}[b]{\linewidth}\raggedright
SUMMARY
\end{minipage} & \begin{minipage}[b]{\linewidth}\raggedright
SCHEDULE
\end{minipage} \\
\midrule\noalign{}
\endhead
\bottomrule\noalign{}
\endlastfoot
\end{longtable}

\subsubsection{Description}\label{description-16}

Turns on warning messages to be printed to the print file (*.PRT); note
that this keyword is activated by default and can subsequently be
switched off by the NOWARN activation keyword. The warning messages may
be turned on and off using keywords WARN and NOWARN. OPM Flow always
prints error messages.

It is recommended that WARN should always be used and action taken if
necessary for the initial runs, once the run has been ``cleaned up'' the
warning messages can be turned off.

There is no data required for this keyword.

This keyword is not supported by OPM Flow but has no effect on the
results so it will be ignored.

\subsubsection{Example}\label{example-13}

\begin{verbatim}
--
--       SWITCH OFF WARNING MESSAGES
--
NOWARN

--
--       INCLUDE SIMULATION GRID WITH SLOPING FAULTS
--
INCLUDE
         './INCLUDE/GRID/IRAP_1005.GRDECL' /

--
--       SWITCH ON WARNING MESSAGES
--
WARN
\end{verbatim}

The example deactivates the warning messages before reading the grid
geometry data using the INCLUDE keyword, and then activates the warning
messages after reading the INCLUDE file.

\bookmarksetup{startatroot}

\chapter{RUNSPEC SECTION}\label{runspec-section}

\section{Introduction}\label{introduction-3}

This is the first section in the OPM Flow input file and defines the key
parameters for the simulator including the dimensions of the model,
phases present in the model (oil, gas, and water for example), number of
tables for a given property and the maximum number of rows for each
table, the maximum number of groups, wells and well completions, as well
as various options to be invoked by the simulator.

The section is defined by the RUNSPEC keyword and is terminated by the
GRID keyword.

Keywords associated with the RUNSPEC section are either Activation type
keywords to switch on a feature or declare a phase, or a Vector (Row
Vector) type of keyword that declares various options or array
dimensions. Examples of Vector type keywords are the SATOPTS keyword
that activates relative permeability assignment options, and the
WELLDIMS keyword that defines the maximum number of wells, completions
per well, maximum number of groups, etc. for the run. The keywords can
be in any order, but the values entered impact the data the simulator
expects in subsequent sections. So for example if the VAPOIL keyword has
been invoked to declare that the vaporized oil phase (condensate) in wet
gas {[}Natural gas that contains significant heavy hydrocarbons such as
propane, butane, and other liquid hydrocarbons is known as wet gas or
rich gas.{]} is present in the run, then the simulator expects the PVTG
keyword in the PROPS section to define the wet gas PVT properties. If
the PVDG keyword (dry gas {[}Natural gas that occurs in the absence of
condensate or liquid hydrocarbons, or gas that had condensable
hydrocarbons removed, is called dry gas. It is primarily methane with
some intermediates{]} PVT data set) has been entered instead in the
PROPS section, then the simulator will report this as an error when it
reads the PVDG keyword or after processing the PROPS section.

Note that in some instances the numerical values entered need to be the
actual values, as per the DIMENS keyword, or maximum values, for example
the previously mentioned WELLDIMS keyword. If the requirement for the
keyword is for the actual value and an incorrect value has been entered
then this may or may not be directly trapped by the simulator, and may
result in numerous error message associated with the data instead. Thus,
if the error reporting messages in a given section are associated with
data that is required to be declared via a RUNSPEC keyword, one should
also verify that the RUNSPEC keyword has been correctly defined with the
associated problematic data error messages. This is particularly true
for the TABDIMS keyword in the RUNSPEC section where some of the
variables are defined as the actual number of tables, whereas the length
of the tables are a maximum value. To be clear, if NTSFUN on the TABDIMS
keyword has been set to five, then there must be exactly five sets of
relative permeability data entered in the PROPS section, otherwise data
input errors will occur.

\section{Keyword Definitions}\label{keyword-definitions-3}

\section{Data Requirements}\label{data-requirements-1}

Table 5.2.1 outlines the minimum set of keywords required by OPM Flow in
order for the simulator to successfully parse the RUNSPEC section.

\begin{longtable}[]{@{}
  >{\raggedright\arraybackslash}p{(\linewidth - 6\tabcolsep) * \real{0.2500}}
  >{\raggedright\arraybackslash}p{(\linewidth - 6\tabcolsep) * \real{0.2500}}
  >{\raggedright\arraybackslash}p{(\linewidth - 6\tabcolsep) * \real{0.2500}}
  >{\raggedright\arraybackslash}p{(\linewidth - 6\tabcolsep) * \real{0.2500}}@{}}
\toprule\noalign{}
\begin{minipage}[b]{\linewidth}\raggedright
No.
\end{minipage} & \begin{minipage}[b]{\linewidth}\raggedright
RUNSPEC Keyword
\end{minipage} & \begin{minipage}[b]{\linewidth}\raggedright
Description
\end{minipage} & \begin{minipage}[b]{\linewidth}\raggedright
Section Applied
\end{minipage} \\
\midrule\noalign{}
\endhead
\bottomrule\noalign{}
\endlastfoot
1 & TITLE & Error: Reference source not found - this is printed on most
reports routed to the print file. & RUNSPEC \\
2 & DIMENS & Error: Reference source not found - the values entered must
be the exact dimensions of the grid, otherwise the simulator will report
errors. & GRID \\
3 & OIL, DISGAS, GAS, VAPOIL, and/or WATER & Error: Reference source not
found Error: Reference source not found Error: Reference source not
found Error: Reference source not found Error: Reference source not
found Only the phases present in the run should be declared, so for a
dry gas run then only GAS and WATER keywords should be declared, since
water is always present in hydrocarbon accumulations. & PROPS \\
4 & FIELD, LAB, or METRIC & Error: Reference source not found Error:
Reference source not found Error: Reference source not found States the
system of units to be used for the input deck. & ALL \\
5 & WELLDIMS & Error: Reference source not found Sets the well and group
dimensions etc. & SCHEDULE \\
Notes: & & & \\
\end{longtable}

\emph{Table 5.2.1: Minimum Set of RUNSPEC Keywords Required by OPM Flow}

If the minimum set of RUNSPEC keywords is only entered, then the default
values of the missing keywords will apply. This means for example only
one set of PVT, relative permeability tables, and one equilibrium region
will be assumed by the simulator.

A complete list of RUNSPEC keywords in alphabetic order is shown in
Table 5.2.2 together with a generalized Topic column that classifies the
functionality of the keyword. Note that not all keywords and features
listed in Table 5.2.2 are implemented in OPM Flow. Cells not colored in
the No.~column indicate the keyword is supported, cells colored gray
indicate that the keyword is not applicable, and finally, cells colored
in orange indicate keywords that are not supported by OPM Flow.

\begin{longtable}[]{@{}
  >{\raggedright\arraybackslash}p{(\linewidth - 6\tabcolsep) * \real{0.2500}}
  >{\raggedright\arraybackslash}p{(\linewidth - 6\tabcolsep) * \real{0.2500}}
  >{\raggedright\arraybackslash}p{(\linewidth - 6\tabcolsep) * \real{0.2500}}
  >{\raggedright\arraybackslash}p{(\linewidth - 6\tabcolsep) * \real{0.2500}}@{}}
\toprule\noalign{}
\begin{minipage}[b]{\linewidth}\raggedright
No.
\end{minipage} & \begin{minipage}[b]{\linewidth}\raggedright
RUNSPEC Keyword
\end{minipage} & \begin{minipage}[b]{\linewidth}\raggedright
Description
\end{minipage} & \begin{minipage}[b]{\linewidth}\raggedright
Topic
\end{minipage} \\
\midrule\noalign{}
\endhead
\bottomrule\noalign{}
\endlastfoot
1 & ACTDIMS & Error: Reference source not found. & Action \\
2 & ACTPARAM & Error: Reference source not found. & Action \\
3 & AITS & Error: Reference source not found. & Numeric \\
4 & AITSOFF & Error: Reference source not found. & Numeric \\
5 & ALKALINE & Error: Reference source not found. & Alkaline \\
6 & API & Error: Reference source not found. & API \\
7 & AQUDIMS & Error: Reference source not found. & Aquifer \\
8 & AUTOREF & Error: Reference source not found. & LGR \\
9 & BIGMODEL & Error: Reference source not found. & N/A \\
10 & \hyperref[REF_HEADING_KEYWORD_BIOFILM]{BIOFILM} & Error: Reference
source not found. This is an OPM Flow specific keyword used to
investigate biofilm effects in underground applications. The module
requires that either the CO2STORE or
\hyperref[REF_HEADING_KEYWORD_H2STORE]{H2STORE} keywords to be active. &
Biofilm \\
11 & BLACKOIL & Error: Reference source not found. & Phases \\
12 & BPARA & Error: Reference source not found. & N/A \\
13 & BPIDIMS & Error: Reference source not found. & Output \\
14 & BRINE & Error: Reference source not found. & Brine \\
15 & CART & Error: Reference source not found. This the default grid
type that is normally not required to be entered. & Input \\
16 & CBMOPTS & Error: Reference source not found. & CBM \\
17 & \hyperref[REF_HEADING_KEYWORD_CO2SOL]{CO2SOL} &
\hyperref[0.0.0.CO2SOL]{CO2SOL -- Activate Dissolved CO2 in the Water
Phase}. This is an OPM Flow implementation that is different to the
commercial simulator. The keyword is a compositional keyword in the
commercial simulator but has been implemented in OPM Flow's black-oil
model. The keyword actives the carbon dioxide storage model for the run
to account for both carbon dioxide and water phase solubility the
simulators CO2-Brine PVT model. & Phases \\
18 & CO2STORE & Error: Reference source not found. This is an OPM Flow
implementation that is different to the commercial simulator, the
keyword is a compositional keyword in the commercial simulator but has
been implemented in OPM Flow's black-oil model. The keyword actives the
carbon dioxide storage model for the run to account for both carbon
dioxide and water phase solubility the simulators CO2-Brine PVT model. &
Phases \\
19 & COAL & Error: Reference source not found. & CBM \\
20 & COLUMNS & Error: Reference source not found. & Input \\
21 & COMPS & COMPS -- Activate Compositional Modeling Formulation. OPM
Flow does not currently support the general compositional modelling
formulation. This keyword is only supported by OPM Flow when the two
component gas-water CO2 storage model has been activated using the
CO2STORE keyword. & Phases \\
22 & CPR & Error: Reference source not found. See section Error:
Reference source not found Error: Reference source not found on how to
invoke various numerical schemes via the OPM Flow command line
interface. & Numeric \\
23 & DEBUG & Error: Reference source not found. & Output \\
24 & DIFFUSE & Error: Reference source not found. & Diffusion \\
25 & DIMENS & Error: Reference source not found. The values entered must
be the exact dimensions of the grid, otherwise the simulator will report
errors. & Input \\
26 & DISGAS & Error: Reference source not found. & Phases \\
27 & DISGASW & Error: Reference source not found This is an OPM Flow
specific keyword that activates dissolved gas in the water phase in the
model. & Phases \\
28 & DISPDIMS & Error: Reference source not found. & Dispersion \\
29 & DUALPERM & Error: Reference source not found. & Dual-Porosity \\
30 & DUALPORO & Error: Reference source not found. & Dual-Porosity \\
31 & DYNRDIMS & Error: Reference source not found. & Input \\
32 & ECHO & Error: Reference source not found. & Output \\
33 & ECLMC & Error: Reference source not found. & Brine \\
34 & END & Error: Reference source not found. & Input \\
35 & ENDINC & Error: Reference source not found. & Input \\
36 & ENDSCALE & Error: Reference source not found. & End-Point \\
37 & ENDSKIP & Error: Reference source not found. & Input \\
38 & \hyperref[REF_HEADING_KEYWORD_EOS_5_3]{EOS} &
\hyperref[0.0.0.EOS]{EOS -- Specify Equation of State}. OPM Flow does
not currently support the general compositional modelling formulation. &
Phases \\
39 & EQLDIMS & Error: Reference source not found. & Equilibration \\
40 & EQLOPTS & Error: Reference source not found. & Equilibration \\
41 & EXTRAPMS & Error: Reference source not found. & Output \\
42 & FAULTDIM & Error: Reference source not found. & Input \\
43 & FIELD & Error: Reference source not found. & Input \\
44 & FMTHMD & Error: Reference source not found. & Output \\
45 & FMTIN & Error: Reference source not found. & Input \\
46 & FMTOUT & Error: Reference source not found. & Output \\
47 & FOAM & Error: Reference source not found. & Foam \\
48 & FORMFEED & Error: Reference source not found. & Output \\
49 & FRICTION & Error: Reference source not found. & Well \\
50 & FULLIMP & Error: Reference source not found. See section Error:
Reference source not found Error: Reference source not found on how to
invoke various numerical schemes via the OPM Flow command line
interface. & Numeric \\
51 & GAS & Error: Reference source not found. & Phases \\
52 & GASFIELD & Error: Reference source not found. & Gas Field \\
53 & GASWAT & GASWAT -- Activate the Gas-Water Model Formulation. &
Phases \\
54 & GDIMS & Error: Reference source not found. & Gradient \\
55 & GIMODEL & Error: Reference source not found. & Phases \\
56 & GRAVDR & Error: Reference source not found. & Dual-Porosity \\
57 & GRAVDRB & Error: Reference source not found. & Dual-Porosity \\
58 & GRAVDRM & Error: Reference source not found. & Dual-Porosity \\
59 & GRIDOPTS & Error: Reference source not found. & Input \\
60 & \hyperref[REF_HEADING_KEYWORD_H2SOL]{H2SOL} &
\hyperref[0.0.0.H2SOL]{H2SOL -- Activate Dissolved H2 in the Water
Phase}. The keyword actives the dissolved hydrogen in the water phase,
where hydrogen is represented by the SOLVENT pseudo component. &
Phases \\
61 & \hyperref[REF_HEADING_KEYWORD_H2STORE]{H2STORE} & Error: Reference
source not found. The keyword actives the hydrogen storage model for the
run to account for both hydrogen and water phase solubility using the
simulator's H2-Brine PVT model. & Phases \\
62 & HMDIMS & Error: Reference source not found. & Gradient \\
63 & HYST & Error: Reference source not found. The HYST keyword
activates the hysteresis option, the keyword should be avoided and the
hysteresis option should be enabled by the HYSTER parameter on the
SATOTPS keyword see - Error: Reference source not found. & Hysteresis \\
64 & IMPES & Error: Reference source not found. See section Error:
Reference source not found Error: Reference source not found on how to
invoke various numerical schemes via the OPM Flow command line
interface. & Numeric \\
65 & IMPLICIT & Error: Reference source not found. See section Error:
Reference source not found Error: Reference source not found on how to
invoke various numerical schemes via the OPM Flow command line
interface. & Numeric \\
66 & INCLUDE & Error: Reference source not found. & Input \\
67 & INSPEC & Error: Reference source not found. & Output \\
68 & LAB & Error: Reference source not found. & Input \\
69 & LGR & Error: Reference source not found. & LGR \\
70 & LGRCOPY & Error: Reference source not found. & LGR \\
71 & LICENSES & Error: Reference source not found. & N/A \\
72 & LIVEOIL & Error: Reference source not found. & Phases \\
73 & LOAD & Error: Reference source not found. & Input \\
74 & LOWSALT & Error: Reference source not found. & Brine \\
75 & MEMORY & Error: Reference source not found. & N/A \\
76 & MESSAGE & Error: Reference source not found. The MESSAGE keyword
outputs a user message to the terminal, as well as to the print
(\emph{.PRT) and debug (}.DBG) files. & Output \\
77 & MESSAGES & Error: Reference source not found. The MESSAGES keyword
defines the print and stops levels for various messages. & Output \\
78 & MESSSRVC & Error: Reference source not found. & Output \\
79 & METRIC & Error: Reference source not found. & Input \\
80 & MICP & Error: Reference source not found. This is an OPM Flow
specific keyword that activates the Microbial Induced Calcite
Precipitation (``MICP'') Model used to investigate leakage remediation.
The module requires that both the MICP and WATER keywords in the RUNSPEC
section to be active. & MICP \\
81 & MINNPCOL & Error: Reference source not found. This is an OPM Flow
specific keyword that sets the minimum number of Newton iterations, as
oppose to the commercial simulator's NUPCOL keyword that defines the
maximum number of Newton iterations within a time step, after which well
targets are frozen. & Numeric \\
82 & MISCIBLE & Error: Reference source not found. & Miscible \\
83 & MONITOR & Error: Reference source not found. & Output \\
84 & MSGFILE & Error: Reference source not found. & Output \\
85 & MULTIN & Error: Reference source not found. & Input \\
86 & MULTOUT & Error: Reference source not found. & Output \\
87 & MULTOUTS & Error: Reference source not found. & Output \\
88 & MULTREAL & Error: Reference source not found. & N/A \\
89 & NETWORK & Error: Reference source not found. & Network \\
90 & NINEPOIN & Error: Reference source not found. & Numeric \\
91 & NMATRIX & Error: Reference source not found. & Dual-Porosity \\
92 & NNEWTF & Error: Reference source not found. & Numeric \\
93 & NOCASC & Error: Reference source not found. See section Error:
Reference source not found Error: Reference source not found on how to
invoke various numerical schemes via the OPM Flow command line
interface. & Tracers \\
94 & NODPPM & Error: Reference source not found & Dual-Porosity \\
95 & NOECHO & Error: Reference source not found. & \\
96 & NOHYST & Error: Reference source not found. & Hysteresis \\
97 & NOINSPEC & Error: Reference source not found. & Output \\
98 & NOMONITO & Error: Reference source not found. & Output \\
99 & NONNC & Error: Reference source not found. & Input \\
100 & NORSSPEC & Error: Reference source not found. & Output \\
101 & NOSIM & Error: Reference source not found. & Numeric \\
102 & NOWARN & Error: Reference source not found. & Output \\
103 & NRSOUT & Error: Reference source not found. & Output \\
104 & NSTACK & Error: Reference source not found. & Numeric \\
105 & NUMRES & Error: Reference source not found. Currently, this should
only be set to one in OPM Flow. & Input \\
106 & NUPCOL & Error: Reference source not found. & Numeric \\
107 & OIL & Error: Reference source not found. & Phases \\
108 & OPTIONS & Error: Reference source not found. & Options \\
109 & PARALLEL & Error: Reference source not found. The keyword defines
the run to use parallel processing and sets the domain decomposition
options. See section Error: Reference source not found Error: Reference
source not found on how to run OPM Flow in parallel mode. & Numeric \\
110 & PARTTRAC & Error: Reference source not found. & Tracers \\
111 & PATHS & Error: Reference source not found. & Input \\
112 & PEDIMS & Error: Reference source not found. & Petro-Elastic \\
113 & PETOPTS & Error: Reference source not found. & Input \\
114 & PIMTDIMS & Error: Reference source not found. & Wells \\
115 & PINTDIMS & Error: Reference source not found. This is an OPM Flow
specific keyword for the simulator's Polymer Molecular Weight Transport
option that uses the polymer molecular weight in calculating the polymer
viscosity, as well as accounting for formation damage due to the water
and polymer injection, by adjusting the wellbore skin pressure. &
Polymer \\
116 & POLYMER & Error: Reference source not found. & Polymer \\
117 & POLYMW & Error: Reference source not found This is an OPM Flow
specific keyword for the simulator's Polymer Molecular Weight Transport
option that uses the polymer molecular weight in calculating the polymer
viscosity, as well as accounting for formation damage due to the water
and polymer injection, by adjusting the wellbore skin pressure. &
Polymer \\
118 & PRECSALT & Error: Reference source not found. This is an OPM Flow
specific keyword for the simulator's Salt Precipitation Model, note that
this is an extension to the commercial simulator's Brine Model. & Water
Vaporization \\
119 & PSTEADY & Error: Reference source not found. & Numeric \\
120 & RADIAL & Error: Reference source not found. & Radial \\
121 & REGDIMS & Error: Reference source not found. & Input \\
122 & RIVRDIMS & Error: Reference source not found. & Rivers \\
123 & ROCKCOMP & Error: Reference source not found. & Compaction \\
124 & RPTCPL & Error: Reference source not found. & Output \\
125 & RPTHMD & Error: Reference source not found. & Gradient \\
126 & RPTRUNSP & Error: Reference source not found. & Output \\
127 & RSSPEC & Error: Reference source not found. & Output \\
128 & RUNSPEC & Error: Reference source not found & \\
129 & SAMG & Error: Reference source not found. See section Error:
Reference source not found Error: Reference source not found on how to
invoke various numerical schemes via the OPM Flow command line
interface. & Numeric \\
130 & SATOPTS & Error: Reference source not found. & Input \\
131 & SAVE & Error: Reference source not found. & Output \\
132 & SCDPDIMS & Error: Reference source not found. & Scale
Deposition \\
133 & SKIP & Error: Reference source not found. & Input \\
134 & SKIP100 & Error: Reference source not found. & Input \\
135 & SKIP300 & Error: Reference source not found. & Input \\
136 & SMRYDIMS & Error: Reference source not found. & Output \\
137 & SOLVDIMS & Error: Reference source not found. See section Error:
Reference source not found Error: Reference source not found on how to
invoke various numerical schemes via the OPM Flow command line
interface. & Numeric \\
138 & SOLVENT & Error: Reference source not found & Solvent \\
139 & SPIDER & Error: Reference source not found. This is an OPM Flow
specific keyword for the simulator's Spider Grid option, that emulates a
Radial Grid via corner-point geometry. The option employs the standard
GRID section radial grid keywords to construct the grid - see the
example in the section on Spider Grids (Error: Reference source not
found). & Radial \\
140 & START & Error: Reference source not found. & Numeric \\
141 & SURFACT & Error: Reference source not found. & Surfactant \\
142 & SURFACTW & Error: Reference source not found. & Surfactant \\
143 & TABDIMS & Error: Reference source not found. & Input \\
144 & TEMP & Error: Reference source not found. The temperature option
(TEMP keyword) and the thermal option (THERMAL keyword) are two separate
modeling facilities in the commercial simulator. See Error: Reference
source not found for additional information. & Thermal \\
145 & THERMAL & Error: Reference source not found. This keyword
activates OPM Flow's thermal modeling option. The energy black-oil
implementation in OPM Flow is a mixture of the commercial simulators
black-oil and the commercial simulator's ``compositional thermal''
keywords, as well as some OPM Flow specific keywords. See the keyword
for further information. & Thermal \\
146 & TITLE & Error: Reference source not found. & Input \\
147 & TRACERS & Error: Reference source not found. & Tracers \\
148 & TRPLPORO & Error: Reference source not found. & Triple-Porosity \\
149 & UDADIMS & Error: Reference source not found. & Action \\
150 & UDQDIMS & Error: Reference source not found. & Action \\
151 & UDQPARAM & Error: Reference source not found. & Action \\
152 & UDTDIMS & Error: Reference source not found. & Action \\
153 & UNCODHMD & Error: Reference source not found. & Gradient \\
154 & UNIFIN & Error: Reference source not found. & Input \\
155 & UNIFOUT & Error: Reference source not found. & Output \\
156 & UNIFOUTS & Error: Reference source not found. & Output \\
157 & UNIFSAVE & Error: Reference source not found. & Output \\
158 & VAPOIL & Error: Reference source not found & Phases \\
159 & VAPWAT & Error: Reference source not found. This is an OPM Flow
specific keyword for the simulator's Water Vaporization Model that is
activated by declaring that vaporized water is present in the run. &
Water Vaporization \\
160 & VE & Error: Reference source not found. & Input \\
161 & VFPIDIMS & Error: Reference source not found. & Well \\
162 & VFPPDIMS & Error: Reference source not found. & Well \\
163 & VISAGE & Error: Reference source not found. & N/A \\
164 & VISCD & Error: Reference source not found. & Dual-Porosity \\
165 & WARN & Error: Reference source not found. & Output \\
166 & WATER & Error: Reference source not found. & Phases \\
167 & WELLDIMS & Error: Reference source not found. & Well \\
168 & WPOTCALC & Error: Reference source not found. & Well \\
169 & WSEGDIMS & Error: Reference source not found. & Multi-Segment
Well \\
Notes: & & & \\
\end{longtable}

\emph{Table 5.2.2: Alphabetic List of RUNSPEC Keywords}

Note that a number of keywords can be classified under several topics,
so the Topic column should only be used as a general reference guide.

\subsubsection{Example}\label{example-14}

A typical RUNSPEC section using various keywords is shown below and on
the next few pages for reference.

\begin{verbatim}
-- ==============================================================================
--
-- RUNSPEC SECTION
--
-- ==============================================================================
RUNSPEC
--
--       DEFINE THE TITLE FOR THE RUN
--
TITLE
CASE CASENAME - SIMULATION RUN TITLE
--
--       DEFINE THE START DATE FOR THE RUN
--
START
         01 'JAN' 2021                                                         /
--
--       SWITCH NO SIMULATION MODE FOR DATA CHECKING COMMENT OUT TO RUN THE MODEL
--
NOSIM

-- ------------------------------------------------------------------------------
-- FLUID TYPES AND TRACER OPTIONS
-- ------------------------------------------------------------------------------
--
--       OIL PHASE IS PRESENT IN THE RUN
--
OIL
--
--       WATER PHASE IS PRESENT IN THE RUN
--
WATER
--
--       GAS PHASE IS PRESENT IN THE RUN
--
GAS
--
--       DISSOLVED GAS IN LIVE OIL IS PRESENT IN THE RUN
--
DISGAS
--
--       VAPORIZED OIL IN WET GAS IS PRESENT IN THE RUN
--
VAPOIL
--
--       NUMBER AND TYPE OF TRACERS
--       NO OIL  NO WAT  NO GAS  NO ENV  DIFF    MAX    MIN    TRACER
--       TRACERS TRACERS TRACERS TRACERS CONTL   NONLIN NONLIN NONLIN
TRACERS
         0       0       1       0      'NODIFF' 1*     1*     1*              /
\end{verbatim}

\begin{verbatim}
-- ------------------------------------------------------------------------------
-- GRID AND EQUILBRATION DIMENSIONS AND OPTIONS
-- ------------------------------------------------------------------------------
--
--       MAX     MAX     MAX
--       NDIVIX  NDIVIY  NDIVIZ
DIMENS
         46      112     22                                                    /
--
--       FAULT
--       SEGMS
FAULTDIM
         10000                                                                 /
--
--       MAX     MAX     RSVD    TVDP    TVDP
--       EQLNUM  DEPTH   NODES   TABLE   NODES
EQLDIMS
         9       1*      20      1       1                                     /
--
--       MAX     TOTAL   INDEP   FLUX    TRACK  CBM    OPERN  WORK  WORK  POLY
--       FIPNUM  REGNS   REGNS   REGNS   REGNS  REGNS  REGNS  REAL  INTG  REGNS
REGDIMS
         9       9       1*      1*      1*     1*     1*     1*    1*    1*   /
--
--       NEG      MAX     MAX
--       MULTS    MULTNUM PINCHNUM
GRIDOPTS
         YES      9       1*                                                   /
--
--       ACTIVATE EQUILIBRATION OPTIONS
--       MOBILE ENDPOINT(MOBILE) STEADY STATE(QUIESC) THRESHOLD(THPRES)
--       IRREVERSIBLE THRESHOLD(IRREVERS)
EQLOPTS
         'MOBILE'   'QUIESC'   'THPRES'   'IRREVERS'                           /

-- ------------------------------------------------------------------------------
-- ROCK AND SATURATION TABLES DIMENSIONS AND OPTIONS
-- ------------------------------------------------------------------------------
--
--       MAX     MAX     MAX     MAX     MAX     MAX    E300
--       NTSFUN  NTPVT   NSSFUN  NPPVT   NTFIP   NRPVT  BLANK  NTEND
TABDIMS
         15      9       40      30      9       25     1*     1               /
--
--       DIRC   REVERSE  MAX     MAX
--       SCALE  SCALE    TABLES  NODES
ENDSCALE
         NODIR  REVERS   1*      1*                                            /
--
--       ACTIVATE RELATIVE PERMEABILITY ASSIGNMENT HYSTERESIS OPTIONS
--       DIRECTTIONAL(DIRECT) IRREVERSIBLE(IRREVERS) HYSTERESIS(HYSTER)
SATOPTS
         'DIRECT'   'IRREVERS'   'HYSTER'                                      /
\end{verbatim}

\begin{verbatim}
-- ------------------------------------------------------------------------------
-- GROUP, WELL AND VFP TABLE DIMENSIONS
-- ------------------------------------------------------------------------------
--
--       WELL    WELL    GRUPS   GRUPS
--       MXWELS  MXCONS  MXGRPS  MXGRPW
WELLDIMS
         60      110     18      40                                            /
--
--       DEFINE WELL PRODUCTIVITY SCALING TABLE DIMENSIONS
--       MAX     MAX
--       TABLES  ENTRIES
PIMTDIMS
         1       51                                                            /
--
--       INJECTING VFP TABLES
--       VFP     VFP     VFP
--       MXMFLO  MXMTHP  NMMVFT
VFPIDIMS
         10      10      12                                                    /
--
--       PRODUCING VFP TABLES
--       VFP     VFP     VFP     VFP     VFP     VFP
--       MXMFLO  MXMTHP  MXMWFR  MXMGFR  MXMALQ  NMMVFT
VFPPDIMS
         20      10      10      10      6       9                             /

-- ------------------------------------------------------------------------------
-- MISCELLANEOUS OPTIONS
-- ------------------------------------------------------------------------------
--
--       USER DEFINED ARGUMENT DIMENSIONS FACILITY
--       MAX     MAX     MAX    MAX    MAX    MAX   MAX   MAX   MAX  MAX    RAND
--       FUNCS   ITEMS   CONNS  FIELD  GROUP  REGS  SEGTM WELL  AQUF BLCKS  OPT
UDQDIMS
         16      16      0      0      0      0     0     0     0    0      N  /
--
--       USER DEFINED ARGUMENT DIMENSIONS
--       NO.     NOT     TOTAL
--       ARGS    USED    UDQ
UDADIMS
         2       1*      1*                                                    /
--
--       USER DEFINED DEFAULT VALUES
--       SEED    RANGE    UNDEFINED   COMPARISON
--       INTG    -AND+    VALUE       TOLERANCE
UDQPARAM
         1       1*       0.0         1.0E-4                                   /

-- ------------------------------------------------------------------------------
-- NUMERICAL AND RUN CONTROL OPTIONS
-- ------------------------------------------------------------------------------
--
--       SET STACK SIZE FOR LINEAR SOLVER
--
NSTACK
         30                                                                    /
--
--       DEFINE THE NUMBER OF ITERATIONS TO UPDATE WELL FLOW TARGETS
--
NUPCOL
         3                                                                     /

--
--       SKIP      ACTIVATE
--       OPTIONS   OPTION
OPTIONS
         77*0      1                                                           /

-- ------------------------------------------------------------------------------
-- INPUT AND OUTPUT OPTIONS
-- ------------------------------------------------------------------------------
--
--       METRIC SYSTEM OF UNITS FOR BOTH INPUT AND OUTPUT
--
METRIC
--
--       SWITCH ON THE UNIFIED INPUT FILES OPTION
--
UNIFIN
--
--       SWITCH ON THE UNIFIED OUTPUT FILES OPTION
--
UNIFOUT
--
--       MESS  COMMT WARN  PROBL ERROR BUG   MESS COMMT WARN  PROBL ERROR BUG
--       LIMIT LIMIT LIMIT LIMIT LIMIT LIMIT STOP STOP  STOP  STOP  STOP  STOP
MESSAGES
         3000  1*    1000  1000  1*    1*    1*   1*    9000  1*    9000  1*   /
--
--       DEBUG PRINTING OPTIONS
--
DEBUG
         8*0   1     11*0  1     30*0                                          /
-- ==============================================================================
\end{verbatim}

The above example is rather extensive but is not unrealistic in terms of
the features being activated and the array size declarations.

On 32-bit systems, array sizes had to be carefully managed to fit within
limited addressable memory. On modern 64-bit systems such as Linux, this
limitation is no longer a concern.

\section{Keyword Definitions}\label{keyword-definitions-4}

\subsection{ACTDIMS -- ACTION Keyword
Dimensions}\label{actdims-action-keyword-dimensions}

\begin{longtable}[]{@{}
  >{\raggedright\arraybackslash}p{(\linewidth - 14\tabcolsep) * \real{0.1250}}
  >{\raggedright\arraybackslash}p{(\linewidth - 14\tabcolsep) * \real{0.1250}}
  >{\raggedright\arraybackslash}p{(\linewidth - 14\tabcolsep) * \real{0.1250}}
  >{\raggedright\arraybackslash}p{(\linewidth - 14\tabcolsep) * \real{0.1250}}
  >{\raggedright\arraybackslash}p{(\linewidth - 14\tabcolsep) * \real{0.1250}}
  >{\raggedright\arraybackslash}p{(\linewidth - 14\tabcolsep) * \real{0.1250}}
  >{\raggedright\arraybackslash}p{(\linewidth - 14\tabcolsep) * \real{0.1250}}
  >{\raggedright\arraybackslash}p{(\linewidth - 14\tabcolsep) * \real{0.1250}}@{}}
\toprule\noalign{}
\begin{minipage}[b]{\linewidth}\raggedright
RUNSPEC
\end{minipage} & \begin{minipage}[b]{\linewidth}\raggedright
GRID
\end{minipage} & \begin{minipage}[b]{\linewidth}\raggedright
EDIT
\end{minipage} & \begin{minipage}[b]{\linewidth}\raggedright
PROPS
\end{minipage} & \begin{minipage}[b]{\linewidth}\raggedright
REGIONS
\end{minipage} & \begin{minipage}[b]{\linewidth}\raggedright
SOLUTION
\end{minipage} & \begin{minipage}[b]{\linewidth}\raggedright
SUMMARY
\end{minipage} & \begin{minipage}[b]{\linewidth}\raggedright
SCHEDULE
\end{minipage} \\
\midrule\noalign{}
\endhead
\bottomrule\noalign{}
\endlastfoot
\end{longtable}

\subsubsection{Description}\label{description-17}

The ACTDIMS keyword defines the maximum number of properties associated
with the ACTION keyword. The ACTION keyword allows the user to enter
computational logic and calculations to the simulation run based on the
how the simulation run is proceeding.

\begin{longtable}[]{@{}
  >{\raggedright\arraybackslash}p{(\linewidth - 6\tabcolsep) * \real{0.2500}}
  >{\raggedright\arraybackslash}p{(\linewidth - 6\tabcolsep) * \real{0.2500}}
  >{\raggedright\arraybackslash}p{(\linewidth - 6\tabcolsep) * \real{0.2500}}
  >{\raggedright\arraybackslash}p{(\linewidth - 6\tabcolsep) * \real{0.2500}}@{}}
\toprule\noalign{}
\begin{minipage}[b]{\linewidth}\raggedright
No.
\end{minipage} & \begin{minipage}[b]{\linewidth}\raggedright
Name
\end{minipage} & \begin{minipage}[b]{\linewidth}\raggedright
Description
\end{minipage} & \begin{minipage}[b]{\linewidth}\raggedright
Default
\end{minipage} \\
\midrule\noalign{}
\endhead
\bottomrule\noalign{}
\endlastfoot
1 & MXACTNS & A positive integer value that defines the maximum number
of ACTION keywords defined in the input deck. & 2 \\
2 & MXLINES & A positive integer value that defines the maximum number
of lines in an ACTION statement. & 50 \\
3 & MXCHARS & A positive integer value that defines the maximum
characters in an ACTION statement. & 80 \\
4 & MXSTATMS & A positive integer value that defines the maximum number
of conditional statements in the ACTION statement. & 3 \\
Notes: & & & \\
\end{longtable}

\emph{Table 5.3: ACTDIMS Keyword Description}

Only the ACTIONX keyword and computational logic have been implemented
in OPM Flow.

\subsubsection{Example}\label{example-15}

\begin{verbatim}
--       ACTION    ACTION   ACTION   ACTION
--       MXACTNS   MXLINES  MXCHARS  MXSTATMS
ACTDIMS
         2         50       80       3                                         /
\end{verbatim}

The above example defines the default values for the ACTDIMS keyword.

\subsection{ACTPARAM -- Define Action Facility Target and Tolerance
Parameters}\label{actparam-define-action-facility-target-and-tolerance-parameters}

\begin{longtable}[]{@{}
  >{\raggedright\arraybackslash}p{(\linewidth - 14\tabcolsep) * \real{0.1250}}
  >{\raggedright\arraybackslash}p{(\linewidth - 14\tabcolsep) * \real{0.1250}}
  >{\raggedright\arraybackslash}p{(\linewidth - 14\tabcolsep) * \real{0.1250}}
  >{\raggedright\arraybackslash}p{(\linewidth - 14\tabcolsep) * \real{0.1250}}
  >{\raggedright\arraybackslash}p{(\linewidth - 14\tabcolsep) * \real{0.1250}}
  >{\raggedright\arraybackslash}p{(\linewidth - 14\tabcolsep) * \real{0.1250}}
  >{\raggedright\arraybackslash}p{(\linewidth - 14\tabcolsep) * \real{0.1250}}
  >{\raggedright\arraybackslash}p{(\linewidth - 14\tabcolsep) * \real{0.1250}}@{}}
\toprule\noalign{}
\begin{minipage}[b]{\linewidth}\raggedright
RUNSPEC
\end{minipage} & \begin{minipage}[b]{\linewidth}\raggedright
GRID
\end{minipage} & \begin{minipage}[b]{\linewidth}\raggedright
EDIT
\end{minipage} & \begin{minipage}[b]{\linewidth}\raggedright
PROPS
\end{minipage} & \begin{minipage}[b]{\linewidth}\raggedright
REGIONS
\end{minipage} & \begin{minipage}[b]{\linewidth}\raggedright
SOLUTION
\end{minipage} & \begin{minipage}[b]{\linewidth}\raggedright
SUMMARY
\end{minipage} & \begin{minipage}[b]{\linewidth}\raggedright
SCHEDULE
\end{minipage} \\
\midrule\noalign{}
\endhead
\bottomrule\noalign{}
\endlastfoot
\end{longtable}

\subsubsection{Description}\label{description-18}

The ACTPARAM keyword defines the maximum target percent value for the
ACTION series of keywords and the fractional equality tolerance for
determining if two numbers are numerically equal when comparing values
using the ACTION series of keywords. The ACTION keyword allows the user
to enter computational logic and calculations to the simulation run
based on the how the simulation run is proceeding.

\begin{longtable}[]{@{}
  >{\raggedright\arraybackslash}p{(\linewidth - 6\tabcolsep) * \real{0.2500}}
  >{\raggedright\arraybackslash}p{(\linewidth - 6\tabcolsep) * \real{0.2500}}
  >{\raggedright\arraybackslash}p{(\linewidth - 6\tabcolsep) * \real{0.2500}}
  >{\raggedright\arraybackslash}p{(\linewidth - 6\tabcolsep) * \real{0.2500}}@{}}
\toprule\noalign{}
\begin{minipage}[b]{\linewidth}\raggedright
No.
\end{minipage} & \begin{minipage}[b]{\linewidth}\raggedright
Name
\end{minipage} & \begin{minipage}[b]{\linewidth}\raggedright
Description
\end{minipage} & \begin{minipage}[b]{\linewidth}\raggedright
Default
\end{minipage} \\
\midrule\noalign{}
\endhead
\bottomrule\noalign{}
\endlastfoot
Field & Metric & Laboratory & \\
1 & MXTOLS & A positive real value that defines the maximum target
percent number for the ACTION series of keywords. The default value of
100 means the target is not applied. & Defined \\
percent 100.0 & percent 100.0 & percent 100.0 & \\
2 & MXEQLS & MXEQLS a real positive number greater than zero and less
than one that defines the tolerance used to determine if two real values
are equal for comparing values in the ACTION series of keywords.
Floating-point numbers (as implemented in computers) are never exact,
one cannot compare floating point numbers for exact equality. Thus,
MXEQLS defines a tolerance. For example, the default value of 1 x 10-4
means that if the difference between two real values is less than 1 x
10-4 then the values are considered equal. & Defined \\
fraction 1 x 10-4 & fraction 1 x 10-4 & fraction 1 x 10-4 & \\
Notes: & & & \\
\end{longtable}

\emph{Table 5.4: ACTPARAM Keyword Description}

Although this keyword is read by OPM Flow, the ACTION and UDQ
computational logic and calculations have not been fully implemented and
therefore this keyword should not be used as it may result in OPM Flow
terminating.

\subsubsection{Example}\label{example-16}

\begin{verbatim}
--
--       ACTION    ACTION
--       MXTOLS    MXEQLS
ACTPARAM
         5.0       1.0E-4                                                      /
\end{verbatim}

The above example defines the maximum tolerance to be 5\% and the
equality tolerance to be the default value of 1.0 x 10-4.

\subsection{AITS -- Activate Intelligent Time
Stepping}\label{aits-activate-intelligent-time-stepping}

\begin{longtable}[]{@{}
  >{\raggedright\arraybackslash}p{(\linewidth - 14\tabcolsep) * \real{0.1250}}
  >{\raggedright\arraybackslash}p{(\linewidth - 14\tabcolsep) * \real{0.1250}}
  >{\raggedright\arraybackslash}p{(\linewidth - 14\tabcolsep) * \real{0.1250}}
  >{\raggedright\arraybackslash}p{(\linewidth - 14\tabcolsep) * \real{0.1250}}
  >{\raggedright\arraybackslash}p{(\linewidth - 14\tabcolsep) * \real{0.1250}}
  >{\raggedright\arraybackslash}p{(\linewidth - 14\tabcolsep) * \real{0.1250}}
  >{\raggedright\arraybackslash}p{(\linewidth - 14\tabcolsep) * \real{0.1250}}
  >{\raggedright\arraybackslash}p{(\linewidth - 14\tabcolsep) * \real{0.1250}}@{}}
\toprule\noalign{}
\begin{minipage}[b]{\linewidth}\raggedright
RUNSPEC
\end{minipage} & \begin{minipage}[b]{\linewidth}\raggedright
GRID
\end{minipage} & \begin{minipage}[b]{\linewidth}\raggedright
EDIT
\end{minipage} & \begin{minipage}[b]{\linewidth}\raggedright
PROPS
\end{minipage} & \begin{minipage}[b]{\linewidth}\raggedright
REGIONS
\end{minipage} & \begin{minipage}[b]{\linewidth}\raggedright
SOLUTION
\end{minipage} & \begin{minipage}[b]{\linewidth}\raggedright
SUMMARY
\end{minipage} & \begin{minipage}[b]{\linewidth}\raggedright
SCHEDULE
\end{minipage} \\
\midrule\noalign{}
\endhead
\bottomrule\noalign{}
\endlastfoot
\end{longtable}

\subsubsection{Description}\label{description-19}

Turns on the commercial simulator's intelligent time stepping.

This keyword is not supported by OPM Flow but has no effect on the
results so it will be ignored.

\subsection{AITSOFF -- Deactivate Intelligent Time
Stepping}\label{aitsoff-deactivate-intelligent-time-stepping}

\begin{longtable}[]{@{}
  >{\raggedright\arraybackslash}p{(\linewidth - 14\tabcolsep) * \real{0.1250}}
  >{\raggedright\arraybackslash}p{(\linewidth - 14\tabcolsep) * \real{0.1250}}
  >{\raggedright\arraybackslash}p{(\linewidth - 14\tabcolsep) * \real{0.1250}}
  >{\raggedright\arraybackslash}p{(\linewidth - 14\tabcolsep) * \real{0.1250}}
  >{\raggedright\arraybackslash}p{(\linewidth - 14\tabcolsep) * \real{0.1250}}
  >{\raggedright\arraybackslash}p{(\linewidth - 14\tabcolsep) * \real{0.1250}}
  >{\raggedright\arraybackslash}p{(\linewidth - 14\tabcolsep) * \real{0.1250}}
  >{\raggedright\arraybackslash}p{(\linewidth - 14\tabcolsep) * \real{0.1250}}@{}}
\toprule\noalign{}
\begin{minipage}[b]{\linewidth}\raggedright
RUNSPEC
\end{minipage} & \begin{minipage}[b]{\linewidth}\raggedright
GRID
\end{minipage} & \begin{minipage}[b]{\linewidth}\raggedright
EDIT
\end{minipage} & \begin{minipage}[b]{\linewidth}\raggedright
PROPS
\end{minipage} & \begin{minipage}[b]{\linewidth}\raggedright
REGIONS
\end{minipage} & \begin{minipage}[b]{\linewidth}\raggedright
SOLUTION
\end{minipage} & \begin{minipage}[b]{\linewidth}\raggedright
SUMMARY
\end{minipage} & \begin{minipage}[b]{\linewidth}\raggedright
SCHEDULE
\end{minipage} \\
\midrule\noalign{}
\endhead
\bottomrule\noalign{}
\endlastfoot
\end{longtable}

\subsubsection{Description}\label{description-20}

Turns off the commercial simulator's intelligent time stepping.

This keyword is not supported by OPM Flow but has no effect on the
results so it will be ignored.

\subsection{ALKALINE -- Activate the Alkaline Phase and
Model}\label{alkaline-activate-the-alkaline-phase-and-model}

\begin{longtable}[]{@{}
  >{\raggedright\arraybackslash}p{(\linewidth - 14\tabcolsep) * \real{0.1250}}
  >{\raggedright\arraybackslash}p{(\linewidth - 14\tabcolsep) * \real{0.1250}}
  >{\raggedright\arraybackslash}p{(\linewidth - 14\tabcolsep) * \real{0.1250}}
  >{\raggedright\arraybackslash}p{(\linewidth - 14\tabcolsep) * \real{0.1250}}
  >{\raggedright\arraybackslash}p{(\linewidth - 14\tabcolsep) * \real{0.1250}}
  >{\raggedright\arraybackslash}p{(\linewidth - 14\tabcolsep) * \real{0.1250}}
  >{\raggedright\arraybackslash}p{(\linewidth - 14\tabcolsep) * \real{0.1250}}
  >{\raggedright\arraybackslash}p{(\linewidth - 14\tabcolsep) * \real{0.1250}}@{}}
\toprule\noalign{}
\begin{minipage}[b]{\linewidth}\raggedright
RUNSPEC
\end{minipage} & \begin{minipage}[b]{\linewidth}\raggedright
GRID
\end{minipage} & \begin{minipage}[b]{\linewidth}\raggedright
EDIT
\end{minipage} & \begin{minipage}[b]{\linewidth}\raggedright
PROPS
\end{minipage} & \begin{minipage}[b]{\linewidth}\raggedright
REGIONS
\end{minipage} & \begin{minipage}[b]{\linewidth}\raggedright
SOLUTION
\end{minipage} & \begin{minipage}[b]{\linewidth}\raggedright
SUMMARY
\end{minipage} & \begin{minipage}[b]{\linewidth}\raggedright
SCHEDULE
\end{minipage} \\
\midrule\noalign{}
\endhead
\bottomrule\noalign{}
\endlastfoot
\end{longtable}

\subsubsection{Description}\label{description-21}

This keyword indicates that an alkaline phase is present in the model
and to activate the Alkaline Model in the run. The keyword will also
invoke data input file checking to ensure that all the required Alkaline
Model input parameters are defined in the input deck.

There is no data required for this keyword and there is no terminating
``/'' for this keyword.

This keyword is not supported by OPM Flow but would change the results
if supported so the simulation will be stopped.

\subsubsection{Example}\label{example-17}

\begin{verbatim}
--
--       ALKALINE PHASE IS PRESENT IN THE RUN
--
ALKALINE
\end{verbatim}

The above example declares that the alkaline phase is active in the
model to activate Alkaline Model.

\subsection{API -- Activate API
Tracking}\label{api-activate-api-tracking}

\begin{longtable}[]{@{}
  >{\raggedright\arraybackslash}p{(\linewidth - 14\tabcolsep) * \real{0.1250}}
  >{\raggedright\arraybackslash}p{(\linewidth - 14\tabcolsep) * \real{0.1250}}
  >{\raggedright\arraybackslash}p{(\linewidth - 14\tabcolsep) * \real{0.1250}}
  >{\raggedright\arraybackslash}p{(\linewidth - 14\tabcolsep) * \real{0.1250}}
  >{\raggedright\arraybackslash}p{(\linewidth - 14\tabcolsep) * \real{0.1250}}
  >{\raggedright\arraybackslash}p{(\linewidth - 14\tabcolsep) * \real{0.1250}}
  >{\raggedright\arraybackslash}p{(\linewidth - 14\tabcolsep) * \real{0.1250}}
  >{\raggedright\arraybackslash}p{(\linewidth - 14\tabcolsep) * \real{0.1250}}@{}}
\toprule\noalign{}
\begin{minipage}[b]{\linewidth}\raggedright
RUNSPEC
\end{minipage} & \begin{minipage}[b]{\linewidth}\raggedright
GRID
\end{minipage} & \begin{minipage}[b]{\linewidth}\raggedright
EDIT
\end{minipage} & \begin{minipage}[b]{\linewidth}\raggedright
PROPS
\end{minipage} & \begin{minipage}[b]{\linewidth}\raggedright
REGIONS
\end{minipage} & \begin{minipage}[b]{\linewidth}\raggedright
SOLUTION
\end{minipage} & \begin{minipage}[b]{\linewidth}\raggedright
SUMMARY
\end{minipage} & \begin{minipage}[b]{\linewidth}\raggedright
SCHEDULE
\end{minipage} \\
\midrule\noalign{}
\endhead
\bottomrule\noalign{}
\endlastfoot
\end{longtable}

\subsubsection{Description}\label{description-22}

This keyword switches on API tracking so that the various ``oil types''
are tracked in the model.

In many reservoirs the initial API gravity of oil varies with depth due
to the heavy viscous fractions occupying the deepest part of the
reservoir whilst the lighter more mobile fractions will occupy the upper
part of the reservoir. As a reservoir is depleted the API gravity of oil
in a cell will gradually change as the different fluids mix.

In OPM Flow it is possible to define different PVT regions in a
reservoir, as in all finite difference formulated simulators, oil moving
from one region to another will suddenly assume the properties of that
region it has moved to. The fluid type tracking option allows the smooth
change of PVT properties in a cell to be simulated by correlating PVT
properties against the API gravity of oil.

There is no data required for this keyword and there is no terminating
``/'' for this keyword.

This keyword is not supported by OPM Flow but would change the results
if supported so the simulation will be stopped.

\subsubsection{Example}\label{example-18}

\begin{verbatim}
--
--       ACTIVATE THE API TRACKING OPTION
--
API
\end{verbatim}

The above example switches on the API tracking facility.

\subsection{AQUDIMS -- Define Aquifer
Dimensions}\label{aqudims-define-aquifer-dimensions}

\begin{longtable}[]{@{}
  >{\raggedright\arraybackslash}p{(\linewidth - 14\tabcolsep) * \real{0.1250}}
  >{\raggedright\arraybackslash}p{(\linewidth - 14\tabcolsep) * \real{0.1250}}
  >{\raggedright\arraybackslash}p{(\linewidth - 14\tabcolsep) * \real{0.1250}}
  >{\raggedright\arraybackslash}p{(\linewidth - 14\tabcolsep) * \real{0.1250}}
  >{\raggedright\arraybackslash}p{(\linewidth - 14\tabcolsep) * \real{0.1250}}
  >{\raggedright\arraybackslash}p{(\linewidth - 14\tabcolsep) * \real{0.1250}}
  >{\raggedright\arraybackslash}p{(\linewidth - 14\tabcolsep) * \real{0.1250}}
  >{\raggedright\arraybackslash}p{(\linewidth - 14\tabcolsep) * \real{0.1250}}@{}}
\toprule\noalign{}
\begin{minipage}[b]{\linewidth}\raggedright
RUNSPEC
\end{minipage} & \begin{minipage}[b]{\linewidth}\raggedright
GRID
\end{minipage} & \begin{minipage}[b]{\linewidth}\raggedright
EDIT
\end{minipage} & \begin{minipage}[b]{\linewidth}\raggedright
PROPS
\end{minipage} & \begin{minipage}[b]{\linewidth}\raggedright
REGIONS
\end{minipage} & \begin{minipage}[b]{\linewidth}\raggedright
SOLUTION
\end{minipage} & \begin{minipage}[b]{\linewidth}\raggedright
SUMMARY
\end{minipage} & \begin{minipage}[b]{\linewidth}\raggedright
SCHEDULE
\end{minipage} \\
\midrule\noalign{}
\endhead
\bottomrule\noalign{}
\endlastfoot
\end{longtable}

\subsubsection{Description}\label{description-23}

The AQUDIMS keyword defines the dimensions of the various aquifer
property data. The data is normally entered on a single line and is
terminated by a ``/''.

\begin{longtable}[]{@{}
  >{\raggedright\arraybackslash}p{(\linewidth - 6\tabcolsep) * \real{0.2500}}
  >{\raggedright\arraybackslash}p{(\linewidth - 6\tabcolsep) * \real{0.2500}}
  >{\raggedright\arraybackslash}p{(\linewidth - 6\tabcolsep) * \real{0.2500}}
  >{\raggedright\arraybackslash}p{(\linewidth - 6\tabcolsep) * \real{0.2500}}@{}}
\toprule\noalign{}
\begin{minipage}[b]{\linewidth}\raggedright
No.
\end{minipage} & \begin{minipage}[b]{\linewidth}\raggedright
Name
\end{minipage} & \begin{minipage}[b]{\linewidth}\raggedright
Description
\end{minipage} & \begin{minipage}[b]{\linewidth}\raggedright
Default
\end{minipage} \\
\midrule\noalign{}
\endhead
\bottomrule\noalign{}
\endlastfoot
1 & MXNAQN & A positive integer value that defines the AQUNUM keyword
maximum number of lines associated with this keyword, that is the
maximum number of numerical aquifers & 1 \\
2 & MXNAQC & A positive integer value that defines the AQUCON keyword
maximum number of lines of connection data associated with this keyword,
that is the maximum number of lines of connection data for numerical
aquifers. & 1 \\
3 & NIFTBL & A positive integer value that defines the AQUTAB keyword
maximum number of Carter-Tracy aquifer tables associated with this
keyword. & 1 \\
4 & NRIFTB & A positive integer value that defines the AQUTAB keyword
maximum number of rows in the Carter-Tracy aquifer tables associated
with this keyword. NRIFTB must not be less than 36 in order to
accommodate the default infinite acting Carter-Tracy aquifer influence
function. & 36 \\
5 & NANAQ & A positive integer value that defines the AQUFETP, AQUFLUX
and AQUCT maximum number of analytical aquifers defined by these three
keywords. & 1 \\
6 & NCAMAX & A positive integer value that defines the maximum number of
cells connected to an analytical aquifer. & 1 \\
7 & MXNALI & A positive integer value that defines the maximum number of
aquifer lists. & 0 \\
8 & MXAAQL & A positive integer value that defines the maximum number of
analytical aquifers in any single aquifer list as defined by (7). & 0 \\
Notes: & & & \\
\end{longtable}

\emph{Table 5.5: AQUDIMS Keyword Description}

\subsubsection{Example}\label{example-19}

\begin{verbatim}
--
--       AQF     AQF     AQF     AQF     AQF     AQF    AQF    AQF
--       MXAQN   MXNAQC  NIFTBL  NRIFTB  NANAQU  NCAMAX MXNALI MXAAQL
AQUDIMS
         1*      1*      1*      1*      1*      1*     1*     1*              /
\end{verbatim}

The above example defines the default values for the AQUDIMS keyword.

\subsection{AUTOREF -- Define Auto Refinement
Options}\label{autoref-define-auto-refinement-options}

\begin{longtable}[]{@{}
  >{\raggedright\arraybackslash}p{(\linewidth - 14\tabcolsep) * \real{0.1250}}
  >{\raggedright\arraybackslash}p{(\linewidth - 14\tabcolsep) * \real{0.1250}}
  >{\raggedright\arraybackslash}p{(\linewidth - 14\tabcolsep) * \real{0.1250}}
  >{\raggedright\arraybackslash}p{(\linewidth - 14\tabcolsep) * \real{0.1250}}
  >{\raggedright\arraybackslash}p{(\linewidth - 14\tabcolsep) * \real{0.1250}}
  >{\raggedright\arraybackslash}p{(\linewidth - 14\tabcolsep) * \real{0.1250}}
  >{\raggedright\arraybackslash}p{(\linewidth - 14\tabcolsep) * \real{0.1250}}
  >{\raggedright\arraybackslash}p{(\linewidth - 14\tabcolsep) * \real{0.1250}}@{}}
\toprule\noalign{}
\begin{minipage}[b]{\linewidth}\raggedright
RUNSPEC
\end{minipage} & \begin{minipage}[b]{\linewidth}\raggedright
GRID
\end{minipage} & \begin{minipage}[b]{\linewidth}\raggedright
EDIT
\end{minipage} & \begin{minipage}[b]{\linewidth}\raggedright
PROPS
\end{minipage} & \begin{minipage}[b]{\linewidth}\raggedright
REGIONS
\end{minipage} & \begin{minipage}[b]{\linewidth}\raggedright
SOLUTION
\end{minipage} & \begin{minipage}[b]{\linewidth}\raggedright
SUMMARY
\end{minipage} & \begin{minipage}[b]{\linewidth}\raggedright
SCHEDULE
\end{minipage} \\
\midrule\noalign{}
\endhead
\bottomrule\noalign{}
\endlastfoot
\end{longtable}

\subsubsection{Description}\label{description-24}

The AUTOREF keyword activates the Auto Refinement option and defines the
parameters for this feature.

This keyword is not supported by OPM Flow but would change the results
if supported so the simulation will be stopped.

\subsection{BIGMODEL -- Activate Big Model Option
(Retired)}\label{bigmodel-activate-big-model-option-retired}

\begin{longtable}[]{@{}
  >{\raggedright\arraybackslash}p{(\linewidth - 14\tabcolsep) * \real{0.1250}}
  >{\raggedright\arraybackslash}p{(\linewidth - 14\tabcolsep) * \real{0.1250}}
  >{\raggedright\arraybackslash}p{(\linewidth - 14\tabcolsep) * \real{0.1250}}
  >{\raggedright\arraybackslash}p{(\linewidth - 14\tabcolsep) * \real{0.1250}}
  >{\raggedright\arraybackslash}p{(\linewidth - 14\tabcolsep) * \real{0.1250}}
  >{\raggedright\arraybackslash}p{(\linewidth - 14\tabcolsep) * \real{0.1250}}
  >{\raggedright\arraybackslash}p{(\linewidth - 14\tabcolsep) * \real{0.1250}}
  >{\raggedright\arraybackslash}p{(\linewidth - 14\tabcolsep) * \real{0.1250}}@{}}
\toprule\noalign{}
\begin{minipage}[b]{\linewidth}\raggedright
RUNSPEC
\end{minipage} & \begin{minipage}[b]{\linewidth}\raggedright
GRID
\end{minipage} & \begin{minipage}[b]{\linewidth}\raggedright
EDIT
\end{minipage} & \begin{minipage}[b]{\linewidth}\raggedright
PROPS
\end{minipage} & \begin{minipage}[b]{\linewidth}\raggedright
REGIONS
\end{minipage} & \begin{minipage}[b]{\linewidth}\raggedright
SOLUTION
\end{minipage} & \begin{minipage}[b]{\linewidth}\raggedright
SUMMARY
\end{minipage} & \begin{minipage}[b]{\linewidth}\raggedright
SCHEDULE
\end{minipage} \\
\midrule\noalign{}
\endhead
\bottomrule\noalign{}
\endlastfoot
\end{longtable}

\subsubsection{Description}\label{description-25}

The original intention in the commercial simulator was to define an
optimized memory allocation method to handle large models; this has
since become redundant and has been retired in the commercial simulator.

This keyword is ignored by both OPM Flow and the commercial simulator
and has no effect on the simulation but is documented here for
completeness.

\subsection{BIOFILM -- Activate the Biofilm
Model}\label{biofilm-activate-the-biofilm-model}

\begin{longtable}[]{@{}
  >{\raggedright\arraybackslash}p{(\linewidth - 14\tabcolsep) * \real{0.1250}}
  >{\raggedright\arraybackslash}p{(\linewidth - 14\tabcolsep) * \real{0.1250}}
  >{\raggedright\arraybackslash}p{(\linewidth - 14\tabcolsep) * \real{0.1250}}
  >{\raggedright\arraybackslash}p{(\linewidth - 14\tabcolsep) * \real{0.1250}}
  >{\raggedright\arraybackslash}p{(\linewidth - 14\tabcolsep) * \real{0.1250}}
  >{\raggedright\arraybackslash}p{(\linewidth - 14\tabcolsep) * \real{0.1250}}
  >{\raggedright\arraybackslash}p{(\linewidth - 14\tabcolsep) * \real{0.1250}}
  >{\raggedright\arraybackslash}p{(\linewidth - 14\tabcolsep) * \real{0.1250}}@{}}
\toprule\noalign{}
\begin{minipage}[b]{\linewidth}\raggedright
RUNSPEC
\end{minipage} & \begin{minipage}[b]{\linewidth}\raggedright
GRID
\end{minipage} & \begin{minipage}[b]{\linewidth}\raggedright
EDIT
\end{minipage} & \begin{minipage}[b]{\linewidth}\raggedright
PROPS
\end{minipage} & \begin{minipage}[b]{\linewidth}\raggedright
REGIONS
\end{minipage} & \begin{minipage}[b]{\linewidth}\raggedright
SOLUTION
\end{minipage} & \begin{minipage}[b]{\linewidth}\raggedright
SUMMARY
\end{minipage} & \begin{minipage}[b]{\linewidth}\raggedright
SCHEDULE
\end{minipage} \\
\midrule\noalign{}
\endhead
\bottomrule\noalign{}
\endlastfoot
\end{longtable}

\subsubsection{Description}\label{description-26}

The \hyperref[REF_HEADING_KEYWORD_BIOFILM]{BIOFILM} keyword activates
the biofilm model for the run. Biofilm-related effects in subsurface
applications such as hydrogen storage include reduced injectivity and
hydrogen loss. The conceptual model includes the following mechanisms:
(1) Biofilm is present in the storage site prior to injection and (2)
the biofilm consumes injected H2/CO2, leading to clogging effects.

\begin{tcolorbox}[enhanced jigsaw, title=\textcolor{quarto-callout-note-color}{\faInfo}\hspace{0.5em}{Note}, coltitle=black, breakable, toprule=.15mm, colbacktitle=quarto-callout-note-color!10!white, leftrule=.75mm, opacitybacktitle=0.6, colback=white, bottomtitle=1mm, toptitle=1mm, titlerule=0mm, left=2mm, colframe=quarto-callout-note-color-frame, arc=.35mm, rightrule=.15mm, bottomrule=.15mm, opacityback=0]

This is an OPM Flow specific keyword used to investigate biofilm effects
in underground applications. The module requires that either the
CO2STORE or \hyperref[REF_HEADING_KEYWORD_H2STORE]{H2STORE} keywords in
the RUNSPEC to be active.

\end{tcolorbox}

There is no data required for this keyword and there is no terminating
``/'' for this keyword.

\subsubsection{Example}\label{example-20}

\begin{verbatim}
        --
--       ACTIVATE THE BIOFILM MODEL
--
BIOFILM
--
--       ACTIVATE THE H2STORE MODEL
--
WATER
GAS
H2STORE
\end{verbatim}

The above example declares that the
\hyperref[REF_HEADING_KEYWORD_BIOFILM]{BIOFILM} model is active for the
run and activates the \hyperref[REF_HEADING_KEYWORD_H2STORE]{H2STORE}
model. For a complete example of this model, see
\href{https://github.com/OPM/opm-tests/blob/master/h2store/H2STORE_BIOFILM.DATA}{H2STORE\_BIOFILM.DATA}.

\subsection{BLACKOIL -- Activate Black-Oil
Phases}\label{blackoil-activate-black-oil-phases}

\begin{longtable}[]{@{}
  >{\raggedright\arraybackslash}p{(\linewidth - 14\tabcolsep) * \real{0.1250}}
  >{\raggedright\arraybackslash}p{(\linewidth - 14\tabcolsep) * \real{0.1250}}
  >{\raggedright\arraybackslash}p{(\linewidth - 14\tabcolsep) * \real{0.1250}}
  >{\raggedright\arraybackslash}p{(\linewidth - 14\tabcolsep) * \real{0.1250}}
  >{\raggedright\arraybackslash}p{(\linewidth - 14\tabcolsep) * \real{0.1250}}
  >{\raggedright\arraybackslash}p{(\linewidth - 14\tabcolsep) * \real{0.1250}}
  >{\raggedright\arraybackslash}p{(\linewidth - 14\tabcolsep) * \real{0.1250}}
  >{\raggedright\arraybackslash}p{(\linewidth - 14\tabcolsep) * \real{0.1250}}@{}}
\toprule\noalign{}
\begin{minipage}[b]{\linewidth}\raggedright
RUNSPEC
\end{minipage} & \begin{minipage}[b]{\linewidth}\raggedright
GRID
\end{minipage} & \begin{minipage}[b]{\linewidth}\raggedright
EDIT
\end{minipage} & \begin{minipage}[b]{\linewidth}\raggedright
PROPS
\end{minipage} & \begin{minipage}[b]{\linewidth}\raggedright
REGIONS
\end{minipage} & \begin{minipage}[b]{\linewidth}\raggedright
SOLUTION
\end{minipage} & \begin{minipage}[b]{\linewidth}\raggedright
SUMMARY
\end{minipage} & \begin{minipage}[b]{\linewidth}\raggedright
SCHEDULE
\end{minipage} \\
\midrule\noalign{}
\endhead
\bottomrule\noalign{}
\endlastfoot
\end{longtable}

\subsubsection{Description}\label{description-27}

This keyword switches on the black-oil formulation, and is equivalent to
setting the phases present in the model to be oil, vaporized oil, gas,
and dissolved gas. Note if water is present in the model this needs to
be explicitly stated via the WATER keyword in the RUNSPEC section (see
also the LIVEOIL keywords in the RUNSPEC section). The keyword is used
by the commercial simulator's compositional THERMAL option to set the
phases present in the model.

There is no data required for this keyword and there is no terminating
``/'' for this keyword.

\subsubsection{Example}\label{example-21}

The following example activates the black-oil phases in the model.

\begin{verbatim}
--
--       ACTIVATE BLACK-OIL PHASES
--
BLACKOIL
\end{verbatim}

Alternatively one could explicitly declare the phases using the
following keywords in the RUNSPEC section.

\begin{verbatim}
--
--       OIL PHASE IS PRESENT IN THE RUN
--
OIL
--
--       VAPORIZED OIL IN WET GAS IS PRESENT IN THE RUN
--
VAPOIL
--
--       GAS PHASE IS PRESENT IN THE RUN
--
GAS
--
--       DISSOLVED GAS IN LIVE OIL IS PRESENT IN THE RUN
--
DISGAS
\end{verbatim}

The above example switches on the black-oil phases in the model.

\subsection{BPARA -- Activate Block Parallel
License}\label{bpara-activate-block-parallel-license}

\begin{longtable}[]{@{}
  >{\raggedright\arraybackslash}p{(\linewidth - 14\tabcolsep) * \real{0.1250}}
  >{\raggedright\arraybackslash}p{(\linewidth - 14\tabcolsep) * \real{0.1250}}
  >{\raggedright\arraybackslash}p{(\linewidth - 14\tabcolsep) * \real{0.1250}}
  >{\raggedright\arraybackslash}p{(\linewidth - 14\tabcolsep) * \real{0.1250}}
  >{\raggedright\arraybackslash}p{(\linewidth - 14\tabcolsep) * \real{0.1250}}
  >{\raggedright\arraybackslash}p{(\linewidth - 14\tabcolsep) * \real{0.1250}}
  >{\raggedright\arraybackslash}p{(\linewidth - 14\tabcolsep) * \real{0.1250}}
  >{\raggedright\arraybackslash}p{(\linewidth - 14\tabcolsep) * \real{0.1250}}@{}}
\toprule\noalign{}
\begin{minipage}[b]{\linewidth}\raggedright
RUNSPEC
\end{minipage} & \begin{minipage}[b]{\linewidth}\raggedright
GRID
\end{minipage} & \begin{minipage}[b]{\linewidth}\raggedright
EDIT
\end{minipage} & \begin{minipage}[b]{\linewidth}\raggedright
PROPS
\end{minipage} & \begin{minipage}[b]{\linewidth}\raggedright
REGIONS
\end{minipage} & \begin{minipage}[b]{\linewidth}\raggedright
SOLUTION
\end{minipage} & \begin{minipage}[b]{\linewidth}\raggedright
SUMMARY
\end{minipage} & \begin{minipage}[b]{\linewidth}\raggedright
SCHEDULE
\end{minipage} \\
\midrule\noalign{}
\endhead
\bottomrule\noalign{}
\endlastfoot
\end{longtable}

\subsubsection{Description}\label{description-28}

The BPARA keyword activates the block parallel license in the commercial
simulator. There is no data required for this keyword; however the
keyword should be followed by the PARALLEL keyword in the RUNSPEC
section, as illustrated in the example below.

There is no data required for this keyword and there is no terminating
``/'' for this keyword.

OPM Flow is an open source project and therefore there is no license
management of the various implemented options or the number of
cores/threads that can be utilized; hence this keyword is not supported
by OPM Flow but would change the results if supported so the simulation
will be stopped.

See section 2.2 Running OPM Flow 2023-04 From The Command Line on how to
run OPM Flow in parallel mode.

\subsubsection{Example}\label{example-22}

\begin{verbatim}
--
--       ACTIVATE BLOCK PARALLEL LICENSE
--
BPARA
--
--       PARALLEL MULTI-CORE OPTIONS
--       NDMAIN     MACHINE TYPE
PARALLEL
           8      DISTRIBUTED                                                /
\end{verbatim}

The above example sets the number of domains (or processors) to eight
and for the simulation to run in block parallel mode. This has no effect
in OPM Flow input decks.

\subsection{BPIDIMS -- Define the Dimensions of the Interpolated Block
Quantities}\label{bpidims-define-the-dimensions-of-the-interpolated-block-quantities}

\begin{longtable}[]{@{}
  >{\raggedright\arraybackslash}p{(\linewidth - 14\tabcolsep) * \real{0.1250}}
  >{\raggedright\arraybackslash}p{(\linewidth - 14\tabcolsep) * \real{0.1250}}
  >{\raggedright\arraybackslash}p{(\linewidth - 14\tabcolsep) * \real{0.1250}}
  >{\raggedright\arraybackslash}p{(\linewidth - 14\tabcolsep) * \real{0.1250}}
  >{\raggedright\arraybackslash}p{(\linewidth - 14\tabcolsep) * \real{0.1250}}
  >{\raggedright\arraybackslash}p{(\linewidth - 14\tabcolsep) * \real{0.1250}}
  >{\raggedright\arraybackslash}p{(\linewidth - 14\tabcolsep) * \real{0.1250}}
  >{\raggedright\arraybackslash}p{(\linewidth - 14\tabcolsep) * \real{0.1250}}@{}}
\toprule\noalign{}
\begin{minipage}[b]{\linewidth}\raggedright
RUNSPEC
\end{minipage} & \begin{minipage}[b]{\linewidth}\raggedright
GRID
\end{minipage} & \begin{minipage}[b]{\linewidth}\raggedright
EDIT
\end{minipage} & \begin{minipage}[b]{\linewidth}\raggedright
PROPS
\end{minipage} & \begin{minipage}[b]{\linewidth}\raggedright
REGIONS
\end{minipage} & \begin{minipage}[b]{\linewidth}\raggedright
SOLUTION
\end{minipage} & \begin{minipage}[b]{\linewidth}\raggedright
SUMMARY
\end{minipage} & \begin{minipage}[b]{\linewidth}\raggedright
SCHEDULE
\end{minipage} \\
\midrule\noalign{}
\endhead
\bottomrule\noalign{}
\endlastfoot
\end{longtable}

\subsubsection{Description}\label{description-29}

The BPDIMS keyword defines the dimensions of the interpolated grid block
quantities for the BPR\_X, BHD\_X,BHDF\_X, BSCN\_X and BCTRA\_X, etc.
variables declared in the SUMMARY section.

This keyword is not supported by OPM Flow but would change the results
if supported so the simulation will be stopped.

\subsection{BRINE -- Activate Brine Tracking
Option}\label{brine-activate-brine-tracking-option}

\begin{longtable}[]{@{}
  >{\raggedright\arraybackslash}p{(\linewidth - 14\tabcolsep) * \real{0.1250}}
  >{\raggedright\arraybackslash}p{(\linewidth - 14\tabcolsep) * \real{0.1250}}
  >{\raggedright\arraybackslash}p{(\linewidth - 14\tabcolsep) * \real{0.1250}}
  >{\raggedright\arraybackslash}p{(\linewidth - 14\tabcolsep) * \real{0.1250}}
  >{\raggedright\arraybackslash}p{(\linewidth - 14\tabcolsep) * \real{0.1250}}
  >{\raggedright\arraybackslash}p{(\linewidth - 14\tabcolsep) * \real{0.1250}}
  >{\raggedright\arraybackslash}p{(\linewidth - 14\tabcolsep) * \real{0.1250}}
  >{\raggedright\arraybackslash}p{(\linewidth - 14\tabcolsep) * \real{0.1250}}@{}}
\toprule\noalign{}
\begin{minipage}[b]{\linewidth}\raggedright
RUNSPEC
\end{minipage} & \begin{minipage}[b]{\linewidth}\raggedright
GRID
\end{minipage} & \begin{minipage}[b]{\linewidth}\raggedright
EDIT
\end{minipage} & \begin{minipage}[b]{\linewidth}\raggedright
PROPS
\end{minipage} & \begin{minipage}[b]{\linewidth}\raggedright
REGIONS
\end{minipage} & \begin{minipage}[b]{\linewidth}\raggedright
SOLUTION
\end{minipage} & \begin{minipage}[b]{\linewidth}\raggedright
SUMMARY
\end{minipage} & \begin{minipage}[b]{\linewidth}\raggedright
SCHEDULE
\end{minipage} \\
\midrule\noalign{}
\endhead
\bottomrule\noalign{}
\endlastfoot
\end{longtable}

\subsubsection{Description}\label{description-30}

The BRINE keyword activates the standard Brine Tracking model and
optionally defines the water phase to have various salinities if the
ECLMC keyword in the RUNSPEC section has been used to activate the
Multi-Component Brine model, that allows for the water phase to have
multiple water salinities. Note that the Multi-Component Brine model is
not supported by OPM Flow.

\begin{longtable}[]{@{}
  >{\raggedright\arraybackslash}p{(\linewidth - 6\tabcolsep) * \real{0.2500}}
  >{\raggedright\arraybackslash}p{(\linewidth - 6\tabcolsep) * \real{0.2500}}
  >{\raggedright\arraybackslash}p{(\linewidth - 6\tabcolsep) * \real{0.2500}}
  >{\raggedright\arraybackslash}p{(\linewidth - 6\tabcolsep) * \real{0.2500}}@{}}
\toprule\noalign{}
\begin{minipage}[b]{\linewidth}\raggedright
No.
\end{minipage} & \begin{minipage}[b]{\linewidth}\raggedright
Name
\end{minipage} & \begin{minipage}[b]{\linewidth}\raggedright
Description
\end{minipage} & \begin{minipage}[b]{\linewidth}\raggedright
Default
\end{minipage} \\
\midrule\noalign{}
\endhead
\bottomrule\noalign{}
\endlastfoot
1 & SALTS & An optional character vector string that defines the salts
to be tracked for when the Multi-Component Brine model has been
activated by the ECLMC keyword in the RUNSPEC section. SALTS should be
set to one or more of the following salt chemical formulae: & None \\
Salt Name & Salt Chemical Formulae & & \\
Sodium Chloride & NaCl & & \\
Potassium Chloride & KCl & & \\
Calcium Chloride & CaCl2 & & \\
Magnesium Chloride & MgCl2 & & \\
Sodium Carbonate & Na2CO3 & & \\
Potassium Carbonate & K2CO3 & & \\
Calcium Carbonate & CaCO3 & & \\
Magnesium Carbonate & MgCO3 & & \\
Sodium Sulfate & Na2SO4 & & \\
Potassium Sulfate & K2SO4 & & \\
Calcium Sulfate & CaSO4 & & \\
Magnesium Sulfate & MgSO4 & & \\
Note that the ECLMC option is currently not available in OPM Flow, so
only the BRINE keyword without the optional SALT variables should be
declared in the input deck. & & & \\
Notes: & & & \\
\end{longtable}

\emph{Table 5.6: BRINE Keyword Description}

See also the PRECSALT and VAPWAT keywords in the RUNSPEC section that
activates OPM Flow's Salt Precipitation model, and the PVTWSALT keyword
in the PROPS section to define the water properties with respect to salt
concentration.

\subsubsection{Example}\label{example-23}

The first example actives the standard Brine model and has no
terminating ``/''.

\begin{verbatim}
--
--       ACTIVATE STANDARD BRINE MODEL IN THE RUN
--
BRINE
\end{verbatim}

The second example illustrates how to activate OPM Flow's Salt
Precipitation model.

\begin{verbatim}
--
--       ACTIVATE STANDARD BRINE MODEL IN THE RUN
--
BRINE
--
--       ACTIVATE THE OPM FLOW SALT PRECIPITATION MODEL (OPM FLOW KEYWORD)
--
PRECSALT
--
--       VAPORIZED WATER IN DRY/WET GAS IS PRESENT IN THE RUN (OPM FLOW KEYWORD)
--
VAPWAT
\end{verbatim}

The third and final example activates the Multi-Component brine model
with four different salts.

\begin{verbatim}
--
--       ACTIVATE MULTI-COMPONENT BRINE MODEL
--
ECLMC
--
--       DEFINE WATER PHASE MULTI-COMPONENT BRINE COMPONENTS
--
--       SALT1   SALT2   SALT3   SALT4   SALT5
BRINE
         NACL    CACL2   MGC03   K2CO3                                         /
\end{verbatim}

This option is currently not available in OPM Flow.

\subsection{CART -- Activate Cartesian
Geometry}\label{cart-activate-cartesian-geometry}

\begin{longtable}[]{@{}
  >{\raggedright\arraybackslash}p{(\linewidth - 14\tabcolsep) * \real{0.1250}}
  >{\raggedright\arraybackslash}p{(\linewidth - 14\tabcolsep) * \real{0.1250}}
  >{\raggedright\arraybackslash}p{(\linewidth - 14\tabcolsep) * \real{0.1250}}
  >{\raggedright\arraybackslash}p{(\linewidth - 14\tabcolsep) * \real{0.1250}}
  >{\raggedright\arraybackslash}p{(\linewidth - 14\tabcolsep) * \real{0.1250}}
  >{\raggedright\arraybackslash}p{(\linewidth - 14\tabcolsep) * \real{0.1250}}
  >{\raggedright\arraybackslash}p{(\linewidth - 14\tabcolsep) * \real{0.1250}}
  >{\raggedright\arraybackslash}p{(\linewidth - 14\tabcolsep) * \real{0.1250}}@{}}
\toprule\noalign{}
\begin{minipage}[b]{\linewidth}\raggedright
RUNSPEC
\end{minipage} & \begin{minipage}[b]{\linewidth}\raggedright
GRID
\end{minipage} & \begin{minipage}[b]{\linewidth}\raggedright
EDIT
\end{minipage} & \begin{minipage}[b]{\linewidth}\raggedright
PROPS
\end{minipage} & \begin{minipage}[b]{\linewidth}\raggedright
REGIONS
\end{minipage} & \begin{minipage}[b]{\linewidth}\raggedright
SOLUTION
\end{minipage} & \begin{minipage}[b]{\linewidth}\raggedright
SUMMARY
\end{minipage} & \begin{minipage}[b]{\linewidth}\raggedright
SCHEDULE
\end{minipage} \\
\midrule\noalign{}
\endhead
\bottomrule\noalign{}
\endlastfoot
\end{longtable}

\subsubsection{Description}\label{description-31}

CART activates the Cartesian grid geometry for the main model, as oppose
to a radial geometry. This is the default geometry and therefore the
keyword does have to be used to activate this type of geometry.

This keyword is not supported by OPM Flow but would change the results
if supported so the simulation will be stopped.

\subsection{CBMOPTS -- Define Coal Bed Methane
Options}\label{cbmopts-define-coal-bed-methane-options}

\begin{longtable}[]{@{}
  >{\raggedright\arraybackslash}p{(\linewidth - 14\tabcolsep) * \real{0.1250}}
  >{\raggedright\arraybackslash}p{(\linewidth - 14\tabcolsep) * \real{0.1250}}
  >{\raggedright\arraybackslash}p{(\linewidth - 14\tabcolsep) * \real{0.1250}}
  >{\raggedright\arraybackslash}p{(\linewidth - 14\tabcolsep) * \real{0.1250}}
  >{\raggedright\arraybackslash}p{(\linewidth - 14\tabcolsep) * \real{0.1250}}
  >{\raggedright\arraybackslash}p{(\linewidth - 14\tabcolsep) * \real{0.1250}}
  >{\raggedright\arraybackslash}p{(\linewidth - 14\tabcolsep) * \real{0.1250}}
  >{\raggedright\arraybackslash}p{(\linewidth - 14\tabcolsep) * \real{0.1250}}@{}}
\toprule\noalign{}
\begin{minipage}[b]{\linewidth}\raggedright
RUNSPEC
\end{minipage} & \begin{minipage}[b]{\linewidth}\raggedright
GRID
\end{minipage} & \begin{minipage}[b]{\linewidth}\raggedright
EDIT
\end{minipage} & \begin{minipage}[b]{\linewidth}\raggedright
PROPS
\end{minipage} & \begin{minipage}[b]{\linewidth}\raggedright
REGIONS
\end{minipage} & \begin{minipage}[b]{\linewidth}\raggedright
SOLUTION
\end{minipage} & \begin{minipage}[b]{\linewidth}\raggedright
SUMMARY
\end{minipage} & \begin{minipage}[b]{\linewidth}\raggedright
SCHEDULE
\end{minipage} \\
\midrule\noalign{}
\endhead
\bottomrule\noalign{}
\endlastfoot
\end{longtable}

\subsubsection{Description}\label{description-32}

This keyword sets the options for the Coal Bed Methane model which is
activated via the COAL keyword in the RUNSPEC section.

This keyword is not supported by OPM Flow but would change the results
if supported so the simulation will be stopped.

\subsection{CO2SOL -- Activate Dissolved CO2 in the Water
Phase}\label{co2sol-activate-dissolved-co2-in-the-water-phase}

\begin{longtable}[]{@{}
  >{\raggedright\arraybackslash}p{(\linewidth - 14\tabcolsep) * \real{0.1250}}
  >{\raggedright\arraybackslash}p{(\linewidth - 14\tabcolsep) * \real{0.1250}}
  >{\raggedright\arraybackslash}p{(\linewidth - 14\tabcolsep) * \real{0.1250}}
  >{\raggedright\arraybackslash}p{(\linewidth - 14\tabcolsep) * \real{0.1250}}
  >{\raggedright\arraybackslash}p{(\linewidth - 14\tabcolsep) * \real{0.1250}}
  >{\raggedright\arraybackslash}p{(\linewidth - 14\tabcolsep) * \real{0.1250}}
  >{\raggedright\arraybackslash}p{(\linewidth - 14\tabcolsep) * \real{0.1250}}
  >{\raggedright\arraybackslash}p{(\linewidth - 14\tabcolsep) * \real{0.1250}}@{}}
\toprule\noalign{}
\begin{minipage}[b]{\linewidth}\raggedright
RUNSPEC
\end{minipage} & \begin{minipage}[b]{\linewidth}\raggedright
GRID
\end{minipage} & \begin{minipage}[b]{\linewidth}\raggedright
EDIT
\end{minipage} & \begin{minipage}[b]{\linewidth}\raggedright
PROPS
\end{minipage} & \begin{minipage}[b]{\linewidth}\raggedright
REGIONS
\end{minipage} & \begin{minipage}[b]{\linewidth}\raggedright
SOLUTION
\end{minipage} & \begin{minipage}[b]{\linewidth}\raggedright
SUMMARY
\end{minipage} & \begin{minipage}[b]{\linewidth}\raggedright
SCHEDULE
\end{minipage} \\
\midrule\noalign{}
\endhead
\bottomrule\noalign{}
\endlastfoot
\end{longtable}

\subsubsection{Description}\label{description-33}

The \hyperref[REF_HEADING_KEYWORD_CO2SOL]{CO2SOL} keyword activates
dissolved carbon dioxide (CO2) in the water phase, where CO2 is
represented by the SOLVENT pseudo component, using the simulator's
CO2-Brine PVT model. This keyword is a compositional keyword in the
commercial simulator but has been implemented in OPM Flow's black-oil
model.

The \hyperref[REF_HEADING_KEYWORD_CO2SOL]{CO2SOL} keyword can be used
when modelling CO2 injection in depleted hydrocarbon reservoirs. See
also the CO2STORE keyword in the RUNSPEC section which can be used when
modelling CO2 injection in saline aquifers.

The CO2-Brine PVT model computes the PVT properties such as density,
viscosity, and enthalpy internally as functions of pressure,
temperature, and composition by using analytic correlations and models
from the literature rather than by interpolation of tabulated values in
the input deck. These values are transformed to the standard black-oil
equivalent PVT tables internally by the simulator. A full description of
the underlying PVT models is described by Sandve et al.~{[}Tor Harald
Sandve1, Sarah E. Gasda, Atgeirr Rasmussen, and Alf Birger Rustad.
Convective dissolution in field scale CO2 storage simulation using the
OPM Flow simulator. Submitted to TCCS 11 -- Trondheim Conference on CO2
Capture, Transport and Storage Trondheim, Norway -- June 21-23,
2021.{]}. Note that the CO2-Brine PVT properties depend on the
temperature and salinity and these must therefore be entered in the
PROPS section. The reservoir temperature can be defined using, e.g., the
RTEMP keyword. Region based salinity can be provided using the SALINITY
keyword.

The \hyperref[REF_HEADING_KEYWORD_CO2SOL]{CO2SOL} keyword must be used
together with the SOLVENT keyword in the RUNSPEC section. The SOLVENT
keyword activates the four component black-oil model: oil, water and
gas, plus a solvent (in this case CO2).

The DISGASW keyword in the RUNSPEC section can be used to model
dissolution of CO2 in the Brine. Note that hydrocarbon gas is not
allowed to dissolve in the Brine (this is considered a reasonable
assumption for most hydrocarbon gases).

The hydrocarbon reservoir fluid properties should be set up as normal
apart from the water PVT properties, which will be defined by the
simulators' CO2-Brine PVT model, so the water PVT keywords (e.g.~PVTW)
are not required. The CO2 PVT properties will also be defined by the
CO2-Brine PVT model so the solvent PVT keywords (e.g.~SDENSITY, PVTSOL)
are also not required.

The oil-water and gas-water saturation functions should be set up as
normal. The CO2 and gas miscible relative permeability tables should be
defined using the SSFN keyword in the PROPS section.

The initial CO2 saturations in the model can be specified by enumeration
using the SSOL keyword in the SOLUTION section.

The CO2 fraction in a well's injection stream can be defined using the
WSOLVENT keyword in the SCHEDULE section.

There is no data required for this keyword and there is no terminating
``/'' for this keyword.

\subsubsection{Example}\label{example-24}

The example shows the usage of
\hyperref[REF_HEADING_KEYWORD_CO2SOL]{CO2SOL} to model CO2 injection in
a live-oil reservoir. The RUNSPEC section includes the
\hyperref[REF_HEADING_KEYWORD_CO2SOL]{CO2SOL} keyword, which activates
dissolved carbon dioxide (CO2) in the water phase. The SOLVENT keyword
activates the solvent pseudo component, which represents CO2. The
DISGASW keyword has been used to activate dissolved CO2 in the brine.

\begin{verbatim}
-- ==============================================================================
--
-- RUNSPEC SECTION
--
-- ==============================================================================
RUNSPEC
-- ------------------------------------------------------------------------------
-- FLUID TYPES AND TRACER OPTIONS
-- ------------------------------------------------------------------------------
--
--       ACTIVATE DISSOLVED CO2 IN WATER (OPM FLOW KEYWORD)
--
CO2SOL
--
--       SOLVENT PHASE IS PRESENT IN THE RUN
--
SOLVENT
--
--       DISSOLVED CO2 IN WATER IS PRESENT IN THE RUN (OPM FLOW KEYWORD)
--
DISGASW
\end{verbatim}

The second part of the example covers the additional data required in
the PROPS section, in which the CO2-gas miscible relative permeabilities
are defined using the SSFN keyword. The initial temperature and salinity
are defined using the RTEMP and SALINITY keywords. The oil and gas PVT
properties should be defined in the normal way. No water PVT data is
required.

\begin{verbatim}
-- ==============================================================================
--
-- PROPS SECTION
--
-- ==============================================================================
PROPS
--
--       RESERVOIR
--       TEMPERATURE
--       -----------
RTEMP
         80.0                                              / RESERVOIR TEMP
--
--       SET SALINITY FOR ALL CELLS (OPM FLOW KEYWORD)
--
SALINITY
         0.7                                                       /
--
-- SOLVENT RELATIVE PERMEABILITY TABLES
--
SSFN
-- SGAS KRGT KRST
-- FRAC
-- -------- -------- --------
 0.0000 0.0000 0.0000
 1.0000 1.0000 1.0000 /
\end{verbatim}

The third part of the example covers initialising the model in the
SOLUTION section. Here the EQUIL keyword is used as normal to
equilibrate the oil, gas and water phases, and the SSOL keyword to
specify that there is initially no CO2 present in the gas phase.

\begin{verbatim}
-- ==============================================================================
--
-- SOLUTION SECTION
--
-- ==============================================================================
SOLUTION
--
--       DATUM   DATUM   OWC     PCOW   GOC    PCGO   RS   RV   N    E300  RVW
--       DEPTH   PRESS   DEPTH   ----   DEPTH  ----   OPT  OPT  OPT  OPT   OPT
EQUIL
         2000.0  200.0   1900.0  0.00  2000.0  0.00   1*   1*   1*   2*    1*  /
--
-- DEFINE INITIAL EQUILIBRATION SOLVENT SAT VALUES FOR ALL CELLS IN THE MODEL
-- BASED ON NX = 100, NY = 100 AND NZ = 3
--
SSOL
         30000*0.0 /
\end{verbatim}

For the summary section, the CO2 storage quantity can be tracked using
the solvent injection and production volumes as shown below.

\begin{verbatim}
-- ==============================================================================
--
-- SUMMARY SECTION
--
-- ==============================================================================
SUMMARY
--
--       EXPORT STANDARD SUMMARY VARIABLE VECTORS TO FILE
--
ALL
--
--       FIELD SOLVENT INJECTION RATE
FNIR
--
--       FIELD SOLVENT PRODUCTION RATE
FNPR
--
--       FIELD SOLVENT INJECTION TOTAL
FNIT
--
--       FIELD SOLVENT PRODUCTION TOTAL
FNPT
\end{verbatim}

The final part of the example covers the SCHEDULE section. The WSOLVENT
keyword is used to specify the injected solvent fraction (in this case
pure CO2).

\begin{verbatim}
-- ==============================================================================
--
-- SCHEDULE SECTION
--
-- ==============================================================================
SCHEDULE
--
--       RESTART CONTROL BASIC = 1 (LAST REPORT TIME ONLY)
--
RPTRST
        'BASIC=1'  'ALLPROPS'   /
--
-- DEFINE GAS INJECTION WELL SOLVENT FRACTION
--
-- WELL SOLVENT
-- NAME FRACTION
-- --------
WSOLVENT
   GI01 1.0000 /
/
\end{verbatim}

\subsection{CO2STORE -- Activate the CO2 Storage
Model}\label{co2store-activate-the-co2-storage-model}

\begin{longtable}[]{@{}
  >{\raggedright\arraybackslash}p{(\linewidth - 14\tabcolsep) * \real{0.1250}}
  >{\raggedright\arraybackslash}p{(\linewidth - 14\tabcolsep) * \real{0.1250}}
  >{\raggedright\arraybackslash}p{(\linewidth - 14\tabcolsep) * \real{0.1250}}
  >{\raggedright\arraybackslash}p{(\linewidth - 14\tabcolsep) * \real{0.1250}}
  >{\raggedright\arraybackslash}p{(\linewidth - 14\tabcolsep) * \real{0.1250}}
  >{\raggedright\arraybackslash}p{(\linewidth - 14\tabcolsep) * \real{0.1250}}
  >{\raggedright\arraybackslash}p{(\linewidth - 14\tabcolsep) * \real{0.1250}}
  >{\raggedright\arraybackslash}p{(\linewidth - 14\tabcolsep) * \real{0.1250}}@{}}
\toprule\noalign{}
\begin{minipage}[b]{\linewidth}\raggedright
RUNSPEC
\end{minipage} & \begin{minipage}[b]{\linewidth}\raggedright
GRID
\end{minipage} & \begin{minipage}[b]{\linewidth}\raggedright
EDIT
\end{minipage} & \begin{minipage}[b]{\linewidth}\raggedright
PROPS
\end{minipage} & \begin{minipage}[b]{\linewidth}\raggedright
REGIONS
\end{minipage} & \begin{minipage}[b]{\linewidth}\raggedright
SOLUTION
\end{minipage} & \begin{minipage}[b]{\linewidth}\raggedright
SUMMARY
\end{minipage} & \begin{minipage}[b]{\linewidth}\raggedright
SCHEDULE
\end{minipage} \\
\midrule\noalign{}
\endhead
\bottomrule\noalign{}
\endlastfoot
\end{longtable}

\subsubsection{Description}\label{description-34}

The CO2STORE keyword activates the carbon dioxide (CO2) storage model
for the run to account for both carbon dioxide and water phase
solubility, via the simulator's CO2-Brine PVT model. This keyword is a
compositional keyword in the commercial simulator but has been
implemented in OPM Flow's black-oil model.

The CO2-Brine PVT model computes the PVT properties such as density,
viscosity, and enthalpy internally as functions of pressure,
temperature, and composition by using analytic correlations and models
from the literature rather than by interpolation from tabulated values.
These values are transformed to the standard black-oil equivalent PVT
tables internally by the simulator. A full description of the underlying
PVT models is described by Sandve et al.~{[}Tor Harald Sandve1, Sarah E.
Gasda, Atgeirr Rasmussen, and Alf Birger Rustad. Convective dissolution
in field scale CO2 storage simulation using the OPM Flow simulator.
Submitted to TCCS 11 -- Trondheim Conference on CO2 Capture, Transport
and Storage Trondheim, Norway -- June 21-23, 2021.{]}. This means that
the normal PVT keywords like DENSITY, PVTO, PVDG etc. are not required
by OPM Flow when this model is activated, and if entered will be ignored
by the simulator. Note that the CO2-Brine PVT properties depend on the
temperature and salinity and these must therefore be entered in the
PROPS or SOLUTION sections. The reservoir temperature can be defined
using, e.g., the RTEMP keyword. Region based salinity can be provided
using the SALINITY keyword.

The CO2STORE keyword must be used together with either: (1) the GAS and
WATER keywords (or alternatively the GASWAT keyword), or (2) the GAS and
OIL keywords in the RUNSPEC section. It is recommended that the standard
method option (1) is used.

The DISGASW keyword in the RUNSPEC section can be used with option (1)
to model dissolution of CO2 in the Brine.

Option (1) has the advantage that it can be used with the VAPWAT and
PRECSALT keywords in the RUNSPEC section to model the impact of both
vaporization of residual water and salt precipitation in the near
wellbore region on injectivity of CO2 injection wells.

In option (1), the GAS and WATER (or GASWAT) keywords declare that the
gas and water phases are present in the model. When the CO2STORE option
is used the water phase represents the brine and the gas phase
represents CO2. Note that the input and output keywords need to be
consistent with this assumption, e.g., GSF (gas saturation function) and
WSF (water saturation function) should be used for the CO2-Brine
relative permeability, etc.

Although, the DISGAS and VAPOIL keywords can be used with option (2) to
model CO2 dissolution and water vaporization, salt precipitation is not
supported with option (2).

In option (2), the GAS and OIL keywords declare that the gas and oil
phases are present in the model. Internally when CO2STORE is used the
oil phase refers to the brine and the gas phase to CO2. Again, the input
and output keywords need to be consistent with this assumption, e.g.,
SGOF (gas-oil relative permeability) is used for the CO2-Brine relative
permeability, FOIP (Field Oil-In-Place) shows the total amount of brine
in the reservoir, etc.

There is no data required for this keyword and there is no terminating
``/'' for this keyword.

\subsubsection{Examples}\label{examples-1}

The first example shows the standard usage of CO2STORE with Option (1)
the Gas-Water model (GASWAT). Here we also activate the dissolved gas in
water (DISGASW) and vaporized water in gas (VAPWAT) options in the
RUNSPEC section.

\begin{verbatim}
-- ==============================================================================
--
-- RUNSPEC SECTION
--
-- ==============================================================================
RUNSPEC
       -- ------------------------------------------------------------------------------
-- FLUID TYPES AND TRACER OPTIONS
-- ------------------------------------------------------------------------------
--
--       ACTIVATE CO2 STORAGE IN THE MODEL (OPM FLOW CO2 STORAGE KEYWORD)
--
CO2STORE
--
--       ACTIVATE GAS-WATER THE MODEL (OPM FLOW KEYWORD)
--
GASWAT
--
--       DISSOLVED GAS IN WATER IS PRESENT IN THE RUN (OPM FLOW KEYWORD)
--
DISGASW
--
--       VAPORIZED WATER IN DRY/WET GAS IS PRESENT IN THE RUN (OPM FLOW KEYWORD)
--
VAPWAT
\end{verbatim}

The second part of the example covers the data required for the PROPS
section, in which the two-phase relative permeability functions are set
using GSF and WSF keywords.

\begin{verbatim}
-- ==============================================================================
--
-- PROPS SECTION
--
-- ==============================================================================
PROPS
--
--       RESERVOIR
--       TEMPERATURE
--       -----------
RTEMP
         80.0                                              / RESERVOIR TEMP
--
--       ROCK COMPRESSIBILITY
--
--       REFERENCE PRESSURE IS TAKEN FROM THE HCPV WEIGHTED RESERVOIR PRESSURE
--
--       REF PRES  CF
--       BARSA     1/BARSA
--       --------  --------
ROCK
         200.0     5.0E-05                                 / ROCK COMPRESSIBILITY
--
--       GAS RELATIVE PERMEABILITY TABLES (OPM FLOW KEYWORD)
--
GSF
--       SGAS       KRG        PCGW
--       FRAC                  PSIA
--       --------  --------   -------
            0.000    0.000       0.0
            0.080    0.001       0.0
            0.170    0.010       0.0
            0.350    0.050       0.0
            0.530    0.200       0.0
            0.620    0.350       0.0
            0.650    0.390       0.0
            0.710    0.560       0.0
            0.800    1.000       0.0   / TABLE NO. 01
--
--       WATER RELATIVE PERMEABILITY TABLES (OPM FLOW KEYWORD)
--
WSF
--       SWAT       KRW
--       FRAC
--       --------   --------
            0.200   0.0000
            0.400   0.1000
            0.800   0.5000
            1.000   1.0000             / TABLE NO. 01
--
--       SET SALINITY FOR ALL CELLS (OPM FLOW KEYWORD)
--
SALINITY
0.7                                                       /
\end{verbatim}

No other data is required to define the fluid and rock properties in the
PROPS section as the data is generated from internal analytic
correlations and models by the simulator. Finally, note that units for
salinity are to the 10-3, thus for metric units we have 10-3 x kg-M/kg.

The third part of the example covers initializing the model in the
SOLUTION section. Here we set the EQUIL(EQLOPT6) parameter equal to one,
to use table number one of the RVWVD keyword, in order to set the
vaporized water versus depth distribution for the model.

\begin{verbatim}
-- ==============================================================================
--
-- SOLUTION SECTION
--
-- ==============================================================================
SOLUTION
--
--       SYSTEM IS SATURATED WITH WATER
--
--       DATUM   DATUM   OWC     PCOW   GOC    PCGO   RS   RV   N    E300  RVW
--       DEPTH   PRESS   DEPTH   ----   DEPTH  ----   OPT  OPT  OPT  OPT   OPT
EQUIL
         2000.0  200.0   1800.0  0.00  1800.0  0.00   1*   1*   1*   2*    1   /
--
--       WATER VAPOR RATIO VS DEPTH (OPM FLOW KEYWORD)
--
--       DEPTH    RVW
--                STB/MSCF
--       ------   --------
RVWVD
         1000.0   0.000
         3000.0   0.000                              / RVW VS DEPTH EQUIL REGN 01
\end{verbatim}

In the SUMMARY section, the simulator supports summary vectors specific
to CO2 storage (see Section 11.1.12 Option Specific Variables -
CO2STORE/H2STORE Model) including those shown below.

\begin{verbatim}
-- ==============================================================================
--
-- SUMMARY SECTION
--
-- ==============================================================================
SUMMARY
--
--       EXPORT STANDARD SUMMARY VARIABLE VECTORS TO FILE
--
ALL
--
--       FIELD CO2 DISOLVED IN WATER PHASE
FWCD
--
--       FIELD CO2 TRAPPED IN GAS PHASE
FGCDI
--
--       FIELD CO2 MOBILE IN GAS PHASE
FGCDM
\end{verbatim}

The final part of the example covers the SCHEDULE section. The standard
WCONINJE keyword is used to set the gas injection rate, in this case
100,000 sm3/day of CO2.

\begin{verbatim}
-- ==============================================================================
--
-- SCHEDULE SECTION
--
-- ==============================================================================
SCHEDULE
--
--       RESTART CONTROL BASIC (LAST=1, ALL=2, YEARLY=4, MONTHLY=5, TSTEP=6)
--
RPTRST
        'BASIC=2' ALLPROPS TEMP RSWSAT RVWSAT                                   /
--
--       WELL SPECIFICATION DATA
--
-- WELL  GROUP     LOCATION  BHP    PHASE  DRAIN  INFLOW  OPEN  CROSS  PVT
-- NAME  NAME        I    J  DEPTH  FLUID  AREA   EQUANS  SHUT  FLOW   TABLE
WELSPECS
INJ1     'G1'       25   25  1*     GAS                                         /
/
--
--       WELL CONNECTION DATA
--
-- WELL  --- LOCATION ---  OPEN   SAT   CONN   WELL   KH    SKIN   D     DIR
-- NAME   II  JJ  K1  K2   SHUT   TAB   FACT   DIA    FACT  FACT   FACT  PEN
COMPDAT
INJ1      25  25   1  10  OPEN    1*    1*    0.3                              /
/
--
--       WELL INJECTION CONTROLS
--
-- WELL  FLUID  OPEN/  CNTL  SURF   RESV   BHP   THP   VFP
-- NAME  TYPE   SHUT   MODE  RATE   RATE   PRES  PRES  TABLE
WCONINJE
INJ1     GAS    OPEN   RATE 100000  1*     300                                 /
/
--
--       ADVANCE SIMULATION BY REPORTING TIME
--
TSTEP
         36*30.4375
/
\end{verbatim}

Note in order to get the liquid phase mole fractions of CO2, that is,
the mole fractions of CO2 in the water phase (XMFCO2), and the vapor
phase mole fractions (YMFWAT) to the restart file, one must use the
command line parameter enable-opm-rst-file set equal to true.

The second example shows how to use CO2STORE with the alternative option
(2). The example below declares that the carbon dioxide storage model is
active for the run to account for both carbon dioxide and water phase
solubility using OPM Flow's CO2-Brine PVT model. Option (2) is used
where the OIL phase refers to the brine and the GAS phase to CO2.

\begin{verbatim}
-- ==============================================================================
--
-- RUNSPEC SECTION
--
-- ==============================================================================
RUNSPEC
-- ------------------------------------------------------------------------------ -- FLUID TYPES AND TRACER OPTIONS
-- ------------------------------------------------------------------------------ --
--       ACTIVATE CO2 STORAGE IN THE MODEL (OPM FLOW CO2 STORAGE KEYWORD)
--
CO2STORE
--
--       OIL PHASE IS PRESENT IN THE RUN BUT IS THE BRINE PHASE FOR CO2STORE
--
OIL
--
--       GAS PHASE IS PRESENT IN THE RUN BUT IS THE CO2 PHASE FOR CO2STORE
--
GAS
--
--       DISSOLVED GAS IN LIVE OIL IS PRESENT IN THE RUN (CO2 IN BRINE)
--
DISGAS
--
--       VAPORIZED OIL IN WET GAS IS PRESENT IN THE RUN (WATER IN CO2)
--
VAPOIL
\end{verbatim}

The second part of the example covers the data required for the PROPS
section, in which the input keywords need to be consistent with the OIL
phase referring to the Brine and the GAS to CO2; that is SGOF (gas-oil
relative permeability) is used to define the CO2-Brine relative
permeability table.

\begin{verbatim}
-- ==============================================================================
--
-- PROPS SECTION
--
-- ==============================================================================
PROPS
--
--       RESERVOIR
--       TEMPERATURE
--       -----------
RTEMP
90.0                                              / RESERVOIR TEMPERATURE
--
--       ROCK COMPRESSIBILITY
--
--       REFERENCE PRESSURE IS TAKEN FROM THE HCPV WEIGHTED RESERVOIR PRESSURE
--
--       REF PRES  CF
--       BARSA     1/BARSA
--       --------  --------
ROCK
         200.0     5.0E-05                                 / ROCK COMPRESSIBILITY
--
--       GAS-OIL RELATIVE PERMEABILITY TABLES (SGOF) – CO2STORE PHASES
SGOF
--       SG         KRG       KROG      PCOG
--       FRAC                           PSIA
--       -------   --------  -------   -------
         0.00000   0.000000  1.00000    0.0000
         1.00000   1.000000  0.00000    0.0000             / TABLE No. 01
--
--       SET SALINITY FOR ALL CELLS (OPM FLOW KEYWORD)
--
SALINITY
0.7                                                       /
\end{verbatim}

The third part of the example covers initializing the model in the
SOLUTION section. Here we set the EQUIL(EQLOPT1 and EQLOPT2) parameters
equal to one, to use table number one of the RSVD and RVVD keywords, in
order to set the initial dissolved CO2 and vaporized water versus depth
distribution for the model.

\begin{verbatim}
-- ==============================================================================
--
-- SOLUTION SECTION
--
-- ==============================================================================
SOLUTION
--
--       SYSTEM IS SATURATED WITH OIL (BRINE)
--
--       DATUM   DATUM   OWC     PCOW   GOC    PCGO   RS   RV   N    E300  RVW
--       DEPTH   PRESS   DEPTH   ----   DEPTH  ----   OPT  OPT  OPT  OPT   OPT
EQUIL
         2000.0  200.0   2200.0  0.00  1800.0  0.00   1    1    1*   2*    1*   /
--
--       DEPTH    RS
--                MSCF/STB
--       ------   --------
RSVD
         1000.0   0.000
         3000.0   0.000                                      /
--
--       DEPTH    RVW
--                STB/MSCF
--       ------   --------
RVVD
         1000.0   0.000
         3000.0   0.000                                      /
\end{verbatim}

See also the DRSDTCON keyword in the SCHEDULE section for details of how
to control convective dissolution of carbon dioxide into the in-situ
brine.

\subsection{COAL -- Activate the Coal Phase (CBM
Model)}\label{coal-activate-the-coal-phase-cbm-model}

\begin{longtable}[]{@{}
  >{\raggedright\arraybackslash}p{(\linewidth - 14\tabcolsep) * \real{0.1250}}
  >{\raggedright\arraybackslash}p{(\linewidth - 14\tabcolsep) * \real{0.1250}}
  >{\raggedright\arraybackslash}p{(\linewidth - 14\tabcolsep) * \real{0.1250}}
  >{\raggedright\arraybackslash}p{(\linewidth - 14\tabcolsep) * \real{0.1250}}
  >{\raggedright\arraybackslash}p{(\linewidth - 14\tabcolsep) * \real{0.1250}}
  >{\raggedright\arraybackslash}p{(\linewidth - 14\tabcolsep) * \real{0.1250}}
  >{\raggedright\arraybackslash}p{(\linewidth - 14\tabcolsep) * \real{0.1250}}
  >{\raggedright\arraybackslash}p{(\linewidth - 14\tabcolsep) * \real{0.1250}}@{}}
\toprule\noalign{}
\begin{minipage}[b]{\linewidth}\raggedright
RUNSPEC
\end{minipage} & \begin{minipage}[b]{\linewidth}\raggedright
GRID
\end{minipage} & \begin{minipage}[b]{\linewidth}\raggedright
EDIT
\end{minipage} & \begin{minipage}[b]{\linewidth}\raggedright
PROPS
\end{minipage} & \begin{minipage}[b]{\linewidth}\raggedright
REGIONS
\end{minipage} & \begin{minipage}[b]{\linewidth}\raggedright
SOLUTION
\end{minipage} & \begin{minipage}[b]{\linewidth}\raggedright
SUMMARY
\end{minipage} & \begin{minipage}[b]{\linewidth}\raggedright
SCHEDULE
\end{minipage} \\
\midrule\noalign{}
\endhead
\bottomrule\noalign{}
\endlastfoot
\end{longtable}

\subsubsection{Description}\label{description-35}

The COAL keyword actives the coal phase and the Coal Bed Methane
(``CBM'') model for the run.

This keyword is not supported by OPM Flow but would change the results
if supported so the simulation will be stopped.

There is no data required for this keyword and there is no terminating
``/'' for this keyword.

\subsubsection{Example}\label{example-25}

\begin{verbatim}
--
--       ACTIVATE THE COAL PHASE (CBM MODEL) IN THE MODEL
--
COAL
\end{verbatim}

The above example declares that the Coal phase is active in the run and
activates the CBM model option.

\subsection{COLUMNS -- Define Input File Column
Margins}\label{columns-define-input-file-column-margins-1}

The COLUMNS keyword defines the input file column margins; characters
outside the margins are ignored by the input parser.

See COLUMNS -- Define Input File Column Margins in the GLOBAL section
for a full description.

\subsection{COMPS -- Activate Compositional Modeling
Formulation}\label{comps-activate-compositional-modeling-formulation}

\begin{longtable}[]{@{}
  >{\raggedright\arraybackslash}p{(\linewidth - 14\tabcolsep) * \real{0.1250}}
  >{\raggedright\arraybackslash}p{(\linewidth - 14\tabcolsep) * \real{0.1250}}
  >{\raggedright\arraybackslash}p{(\linewidth - 14\tabcolsep) * \real{0.1250}}
  >{\raggedright\arraybackslash}p{(\linewidth - 14\tabcolsep) * \real{0.1250}}
  >{\raggedright\arraybackslash}p{(\linewidth - 14\tabcolsep) * \real{0.1250}}
  >{\raggedright\arraybackslash}p{(\linewidth - 14\tabcolsep) * \real{0.1250}}
  >{\raggedright\arraybackslash}p{(\linewidth - 14\tabcolsep) * \real{0.1250}}
  >{\raggedright\arraybackslash}p{(\linewidth - 14\tabcolsep) * \real{0.1250}}@{}}
\toprule\noalign{}
\begin{minipage}[b]{\linewidth}\raggedright
RUNSPEC
\end{minipage} & \begin{minipage}[b]{\linewidth}\raggedright
GRID
\end{minipage} & \begin{minipage}[b]{\linewidth}\raggedright
EDIT
\end{minipage} & \begin{minipage}[b]{\linewidth}\raggedright
PROPS
\end{minipage} & \begin{minipage}[b]{\linewidth}\raggedright
REGIONS
\end{minipage} & \begin{minipage}[b]{\linewidth}\raggedright
SOLUTION
\end{minipage} & \begin{minipage}[b]{\linewidth}\raggedright
SUMMARY
\end{minipage} & \begin{minipage}[b]{\linewidth}\raggedright
SCHEDULE
\end{minipage} \\
\midrule\noalign{}
\endhead
\bottomrule\noalign{}
\endlastfoot
\end{longtable}

\subsubsection{Description}\label{description-36}

The COMPS keyword activates the Compositional Modeling Formulation, and
declares the number of components active in the model.

OPM Flow does not currently support the general compositional modeling
formulation.

\begin{tcolorbox}[enhanced jigsaw, title=\textcolor{quarto-callout-note-color}{\faInfo}\hspace{0.5em}{Note}, coltitle=black, breakable, toprule=.15mm, colbacktitle=quarto-callout-note-color!10!white, leftrule=.75mm, opacitybacktitle=0.6, colback=white, bottomtitle=1mm, toptitle=1mm, titlerule=0mm, left=2mm, colframe=quarto-callout-note-color-frame, arc=.35mm, rightrule=.15mm, bottomrule=.15mm, opacityback=0]

This keyword is only supported by OPM Flow when the two component
gas-water CO2 storage model has been activated using the CO2STORE
keyword and either the GASWAT or the GAS and WATER keywords in the
RUNSPEC section. Secondly, although OPM Flow parses the keyword, the
simulator currently ignores the data for this keyword.

\end{tcolorbox}

\begin{longtable}[]{@{}
  >{\raggedright\arraybackslash}p{(\linewidth - 6\tabcolsep) * \real{0.2500}}
  >{\raggedright\arraybackslash}p{(\linewidth - 6\tabcolsep) * \real{0.2500}}
  >{\raggedright\arraybackslash}p{(\linewidth - 6\tabcolsep) * \real{0.2500}}
  >{\raggedright\arraybackslash}p{(\linewidth - 6\tabcolsep) * \real{0.2500}}@{}}
\toprule\noalign{}
\begin{minipage}[b]{\linewidth}\raggedright
No.
\end{minipage} & \begin{minipage}[b]{\linewidth}\raggedright
Name
\end{minipage} & \begin{minipage}[b]{\linewidth}\raggedright
Description
\end{minipage} & \begin{minipage}[b]{\linewidth}\raggedright
Default
\end{minipage} \\
\midrule\noalign{}
\endhead
\bottomrule\noalign{}
\endlastfoot
1 & COMPS & A positive integer defining the number of compositional
components active in the model. Only the default value of two is
currently supported by OPM Flow. & 2 \\
Notes: & & & \\
\end{longtable}

\emph{Table 5.7: COMPS Keyword Description}

\subsubsection{Example}\label{example-26}

The following example shows how to request a two component compositional
modeling formulation to be used with the CO2STORE and GASWAT options.

\begin{verbatim}
-- ==============================================================================
--
-- RUNSPEC SECTION
--
-- ==============================================================================
RUNSPEC

--
--       ACTIVATE CO2 STORAGE IN THE MODEL (OPM FLOW CO2 STORAGE KEYWORD)
--
CO2STORE
--
--       ACTIVATE COMPOSITIONAL MODELING FORMULATION (OPM FLOW KEYWORD)
--
COMPS
         2                                                                     /
--
--       ACTIVATE GAS-WATER THE MODEL (OPM FLOW KEYWORD)
--
GASWAT
\end{verbatim}

\subsection{CPR -- Activate Constrained Pressure Residual (``CPR'')
Linear
Solver}\label{cpr-activate-constrained-pressure-residual-cpr-linear-solver}

\begin{longtable}[]{@{}
  >{\raggedright\arraybackslash}p{(\linewidth - 14\tabcolsep) * \real{0.1250}}
  >{\raggedright\arraybackslash}p{(\linewidth - 14\tabcolsep) * \real{0.1250}}
  >{\raggedright\arraybackslash}p{(\linewidth - 14\tabcolsep) * \real{0.1250}}
  >{\raggedright\arraybackslash}p{(\linewidth - 14\tabcolsep) * \real{0.1250}}
  >{\raggedright\arraybackslash}p{(\linewidth - 14\tabcolsep) * \real{0.1250}}
  >{\raggedright\arraybackslash}p{(\linewidth - 14\tabcolsep) * \real{0.1250}}
  >{\raggedright\arraybackslash}p{(\linewidth - 14\tabcolsep) * \real{0.1250}}
  >{\raggedright\arraybackslash}p{(\linewidth - 14\tabcolsep) * \real{0.1250}}@{}}
\toprule\noalign{}
\begin{minipage}[b]{\linewidth}\raggedright
RUNSPEC
\end{minipage} & \begin{minipage}[b]{\linewidth}\raggedright
GRID
\end{minipage} & \begin{minipage}[b]{\linewidth}\raggedright
EDIT
\end{minipage} & \begin{minipage}[b]{\linewidth}\raggedright
PROPS
\end{minipage} & \begin{minipage}[b]{\linewidth}\raggedright
REGIONS
\end{minipage} & \begin{minipage}[b]{\linewidth}\raggedright
SOLUTION
\end{minipage} & \begin{minipage}[b]{\linewidth}\raggedright
SUMMARY
\end{minipage} & \begin{minipage}[b]{\linewidth}\raggedright
SCHEDULE
\end{minipage} \\
\midrule\noalign{}
\endhead
\bottomrule\noalign{}
\endlastfoot
\end{longtable}

\subsubsection{Description}\label{description-37}

Turns on the Constrained Pressure Residual (``CPR'') {[}Wallis, J. R.,
Little, T. E., and Nolen, J. S.: ``Constrained Residual Acceleration of
Conjugate Residual Methods,'' paper SPE 13536 presented at the SPE
Reservoir Simulation Symposium, Dallas, Texas, USA (February 10-13,
1985).{]}, {[}R. Scheichl, M. Roland, J. Wendebourg, Decoupling and
block preconditioning for sedimentary basin simulations, Computational
Geosciences 7 (2003) 295\{318.{]} and {[}Klemetsdal, Ø.S., Møyner, O. \&
Lie, KA. Accelerating multiscale simulation of complex geomodels by use
of dynamically adapted basis functions. Comput Geosci 24, 459--476
(2020). https://doi.org/10.1007/s10596-019-9827-z.{]} preconditioner
linear solver option, and declares how the solver should be applied. The
keyword is equivalent to using the OPM Flow command line parameter
--linear-solver= ``cprw''. Note that if the command line has been used,
then this will take precedence over the CPR keyword.

\begin{longtable}[]{@{}
  >{\raggedright\arraybackslash}p{(\linewidth - 6\tabcolsep) * \real{0.2500}}
  >{\raggedright\arraybackslash}p{(\linewidth - 6\tabcolsep) * \real{0.2500}}
  >{\raggedright\arraybackslash}p{(\linewidth - 6\tabcolsep) * \real{0.2500}}
  >{\raggedright\arraybackslash}p{(\linewidth - 6\tabcolsep) * \real{0.2500}}@{}}
\toprule\noalign{}
\begin{minipage}[b]{\linewidth}\raggedright
No.
\end{minipage} & \begin{minipage}[b]{\linewidth}\raggedright
Name
\end{minipage} & \begin{minipage}[b]{\linewidth}\raggedright
Description
\end{minipage} & \begin{minipage}[b]{\linewidth}\raggedright
Default
\end{minipage} \\
\midrule\noalign{}
\endhead
\bottomrule\noalign{}
\endlastfoot
1 & CPROPTN & A defined character string that determines how the CPR
linear solver should be applied, and should be set to one of the
following: Note that OPM Flow only supports the ORIGINAL option, which
is the default value in OPM Flow, whereas the default value in the
commercial simulator is ADAPTIVE. . & ORIGINAL \\
Notes: & & & \\
\end{longtable}

\emph{Table 5.8: CPR Keyword Description}

See section 2.2 Running OPM Flow 2023-04 From The Command Line on how to
invoke various numerical schemes via the OPM Flow command line
interface.

\subsubsection{Example}\label{example-27}

\begin{verbatim}
--
--       ACTIVATE CONSTRAINED PRESSURE RESIDUAL LINEAR SOLVER FOR THE RUN
--
CPR
/
\end{verbatim}

The above example activates the Constrained Pressure Residual
preconditioner linear solver using the default option, that is the
ORIGINAL option for OPM Flow.

To enable the CPR preconditioner linear solver option using the command
line parameter, use:

\begin{verbatim}
flow --linear-solver="cpr" CASENAME.DATA
\end{verbatim}

which should usually be combined with the
-matrix-add-well-contributions=true option, that is:

\begin{verbatim}
flow --linear-solver=cpr --matrix-add-well-contributions=true CASENAME.DATA
\end{verbatim}

However, the new CPRW preconditioner, that includes wells does not need
the latter option, so in this instance use:

\begin{verbatim}
flow --linear-solver=cprw CASENAME.DATA
\end{verbatim}

As of this release the CPRW preconditioner, should be considered
experimental, and the recommended configuration for running this is:

\begin{verbatim}
flow --linear-solver=cprw --linear-solver-reduction=0.005 --cpr-reuse-setup=4
     --cpr-reuse-interval=10 CASENAME.DATA
\end{verbatim}

See new Rasmussen et al.~{[}Atgeirr Flø Rasmussen, Tor Harald Sandve,
Kai Bao, Andreas Lauser, Joakim Hove, Bård Skaflestad, Robert Klöfkorn,
Markus Blatt, Alf Birger Rustad, Ove Sævareid, Knut-Andreas Lie, Andreas
Thune, The Open Porous Media Flow reservoir simulator, Computers \&
Mathematics with Applications, Volume 81, 2021, Pages 159-185, ISSN
0898-1221, \url{https://doi.org/10.1016/j.camwa.2020.05.014}.
(https://www.sciencedirect.com/science/article/pii/S0898122120302182).{]}
for further information on the available numerical algorithms available
in OPM Flow.

\begin{verbatim}
\end{verbatim}

\subsection{DEBUG -- Define the Debug Data to be Printed to
File}\label{debug-define-the-debug-data-to-be-printed-to-file-1}

This keyword defines the debug data to be written to the debug file
(*.DBG). This keyword is not supported by OPM Flow but has no effect on
the results so it will be ignored.

See DEBUG -- Define the Debug Data to be Printed to File in the GLOBAL
section for a full description.

\subsection{DIFFUSE -- Activate Molecular Diffusion
Option}\label{diffuse-activate-molecular-diffusion-option}

\begin{longtable}[]{@{}
  >{\raggedright\arraybackslash}p{(\linewidth - 14\tabcolsep) * \real{0.1250}}
  >{\raggedright\arraybackslash}p{(\linewidth - 14\tabcolsep) * \real{0.1250}}
  >{\raggedright\arraybackslash}p{(\linewidth - 14\tabcolsep) * \real{0.1250}}
  >{\raggedright\arraybackslash}p{(\linewidth - 14\tabcolsep) * \real{0.1250}}
  >{\raggedright\arraybackslash}p{(\linewidth - 14\tabcolsep) * \real{0.1250}}
  >{\raggedright\arraybackslash}p{(\linewidth - 14\tabcolsep) * \real{0.1250}}
  >{\raggedright\arraybackslash}p{(\linewidth - 14\tabcolsep) * \real{0.1250}}
  >{\raggedright\arraybackslash}p{(\linewidth - 14\tabcolsep) * \real{0.1250}}@{}}
\toprule\noalign{}
\begin{minipage}[b]{\linewidth}\raggedright
RUNSPEC
\end{minipage} & \begin{minipage}[b]{\linewidth}\raggedright
GRID
\end{minipage} & \begin{minipage}[b]{\linewidth}\raggedright
EDIT
\end{minipage} & \begin{minipage}[b]{\linewidth}\raggedright
PROPS
\end{minipage} & \begin{minipage}[b]{\linewidth}\raggedright
REGIONS
\end{minipage} & \begin{minipage}[b]{\linewidth}\raggedright
SOLUTION
\end{minipage} & \begin{minipage}[b]{\linewidth}\raggedright
SUMMARY
\end{minipage} & \begin{minipage}[b]{\linewidth}\raggedright
SCHEDULE
\end{minipage} \\
\midrule\noalign{}
\endhead
\bottomrule\noalign{}
\endlastfoot
\end{longtable}

\subsubsection{Description}\label{description-38}

The DIFFUSE keyword activates OPM Flow's Molecular Diffusion option
based on fluid and grid data (Sandve et al.~{[}Tor Harald Sandve1, Sarah
E. Gasda, Atgeirr Rasmussen, and Alf Birger Rustad. Convective
dissolution in field scale CO2 storage simulation using the OPM Flow
simulator. Submitted to TCCS 11 -- Trondheim Conference on CO2 Capture,
Transport and Storage Trondheim, Norway -- June 21-23, 2021.{]}),
similar to the commercial simulator. For field-scale simulations
diffusion is a sub-grid phenomenon and is typically not explicitly
represented in the flow equations. However, for simulations on the
laboratory scale diffusion plays a direct role and therefore needs to be
explicitly represented in the flow equations. Diffusion coefficients
that control the diffusion depend on temperature, pressure, and
salinity. In OPM Flow, the diffusion coefficients are computed
internally for pure water using McLachlan and Danckwerts {[}McLachlan,
C. N. S., \& Danckwerts, P. V. (1972). Desorption of carbon dioxide from
aqueous potash solutions with and without the addition of arsenite as a
catalyst. Trans. Inst. Chem. Eng, 50, 300-309.{]} and modified to
account for salinity using Ratcliff and Holdcroft {[}Ratcliff, G. A., \&
Holdcroft, J. G. (1963). Diffusivities of gases in aqueous electrolyte
solutions. Trans. Inst. Chem. Eng, 41(10), 315-319.{]}. The effect of
the porous media on the diffusion is modeled using the relation
suggested by Millington and Quirk {[}Millington, R. J., \& Quirk, J. P.
(1961). Permeability of porous solids. Transactions of the Faraday
Society, 57, 1200-1207{]}. The coefficients can also be given for each
PVT region as an input parameter using the DIFFC keyword in the PROPS
section.

There is no data required for this keyword and there is no terminating
``/'' for this keyword.

\begin{tcolorbox}[enhanced jigsaw, title=\textcolor{quarto-callout-note-color}{\faInfo}\hspace{0.5em}{Note}, coltitle=black, breakable, toprule=.15mm, colbacktitle=quarto-callout-note-color!10!white, leftrule=.75mm, opacitybacktitle=0.6, colback=white, bottomtitle=1mm, toptitle=1mm, titlerule=0mm, left=2mm, colframe=quarto-callout-note-color-frame, arc=.35mm, rightrule=.15mm, bottomrule=.15mm, opacityback=0]

The option has been tested in combination with the CO2STORE keyword, but
not for the general case at this point.

\end{tcolorbox}

\subsubsection{Example}\label{example-28}

\begin{verbatim}
--
--       ACTIVATE MOLECULAR DIFFUSION OPTION
--
DIFFUSE
\end{verbatim}

The above example switches on the molecular diffusion facility.

\subsection{DIMENS -- Define the Dimensions of the
Model}\label{dimens-define-the-dimensions-of-the-model}

\begin{longtable}[]{@{}
  >{\raggedright\arraybackslash}p{(\linewidth - 14\tabcolsep) * \real{0.1250}}
  >{\raggedright\arraybackslash}p{(\linewidth - 14\tabcolsep) * \real{0.1250}}
  >{\raggedright\arraybackslash}p{(\linewidth - 14\tabcolsep) * \real{0.1250}}
  >{\raggedright\arraybackslash}p{(\linewidth - 14\tabcolsep) * \real{0.1250}}
  >{\raggedright\arraybackslash}p{(\linewidth - 14\tabcolsep) * \real{0.1250}}
  >{\raggedright\arraybackslash}p{(\linewidth - 14\tabcolsep) * \real{0.1250}}
  >{\raggedright\arraybackslash}p{(\linewidth - 14\tabcolsep) * \real{0.1250}}
  >{\raggedright\arraybackslash}p{(\linewidth - 14\tabcolsep) * \real{0.1250}}@{}}
\toprule\noalign{}
\begin{minipage}[b]{\linewidth}\raggedright
RUNSPEC
\end{minipage} & \begin{minipage}[b]{\linewidth}\raggedright
GRID
\end{minipage} & \begin{minipage}[b]{\linewidth}\raggedright
EDIT
\end{minipage} & \begin{minipage}[b]{\linewidth}\raggedright
PROPS
\end{minipage} & \begin{minipage}[b]{\linewidth}\raggedright
REGIONS
\end{minipage} & \begin{minipage}[b]{\linewidth}\raggedright
SOLUTION
\end{minipage} & \begin{minipage}[b]{\linewidth}\raggedright
SUMMARY
\end{minipage} & \begin{minipage}[b]{\linewidth}\raggedright
SCHEDULE
\end{minipage} \\
\midrule\noalign{}
\endhead
\bottomrule\noalign{}
\endlastfoot
\end{longtable}

\subsubsection{Description}\label{description-39}

DIMENS defines the dimensions of the model entered as integer vector.
The keyword can be used with all grid types.

\begin{longtable}[]{@{}
  >{\raggedright\arraybackslash}p{(\linewidth - 6\tabcolsep) * \real{0.2500}}
  >{\raggedright\arraybackslash}p{(\linewidth - 6\tabcolsep) * \real{0.2500}}
  >{\raggedright\arraybackslash}p{(\linewidth - 6\tabcolsep) * \real{0.2500}}
  >{\raggedright\arraybackslash}p{(\linewidth - 6\tabcolsep) * \real{0.2500}}@{}}
\toprule\noalign{}
\begin{minipage}[b]{\linewidth}\raggedright
No.
\end{minipage} & \begin{minipage}[b]{\linewidth}\raggedright
Name
\end{minipage} & \begin{minipage}[b]{\linewidth}\raggedright
Description
\end{minipage} & \begin{minipage}[b]{\linewidth}\raggedright
Default
\end{minipage} \\
\midrule\noalign{}
\endhead
\bottomrule\noalign{}
\endlastfoot
1 & NX & A positive integer value that defines the number of grid blocks
in the x direction for Cartesian grids or the number of grid blocks in
the r direction for radial grids & None \\
2 & NY & A positive integer value that defines the number of grid blocks
in the y direction for Cartesian grids or the number of grid blocks in
the theta direction for radial grids. & None \\
3 & NZ & A positive integer value that defines the number of grid blocks
in the z direction for both Cartesian and radial grids. & None \\
Notes: & & & \\
\end{longtable}

\emph{Table 5.9: DIMENS Keyword Description}

Note that NX, NY and NZ are not maximum values but the actual size of
the grid. OPM Flow applies these parameters when reading in particular
data sets. For example if NX, NY, and NZ are set to 10, 10 and 10
respectively, then for the grid property data like PORO. OPM Flow
expects to read in 10 x 10 x 10 or 1,000 porosity values for the PORO
array. If the number of porosity values is not equal to 1,000 then OPM
Flow will produce an error.

\subsubsection{Example}\label{example-29}

\begin{verbatim}
--
--       MAX     MAX     MAX
--       NDIVIX  NDIVIY  NDIVIZ
DIMENS
         46      112     22                                                    /
\end{verbatim}

The above example defines the dimensions for the Norne model of 46 cells
in the x direction, 112 cells in the y direction and 22 cells in the z
direction.

\subsection{DISGAS -- Activate the Dissolved Gas Phase in the
Model}\label{disgas-activate-the-dissolved-gas-phase-in-the-model}

\begin{longtable}[]{@{}
  >{\raggedright\arraybackslash}p{(\linewidth - 14\tabcolsep) * \real{0.1250}}
  >{\raggedright\arraybackslash}p{(\linewidth - 14\tabcolsep) * \real{0.1250}}
  >{\raggedright\arraybackslash}p{(\linewidth - 14\tabcolsep) * \real{0.1250}}
  >{\raggedright\arraybackslash}p{(\linewidth - 14\tabcolsep) * \real{0.1250}}
  >{\raggedright\arraybackslash}p{(\linewidth - 14\tabcolsep) * \real{0.1250}}
  >{\raggedright\arraybackslash}p{(\linewidth - 14\tabcolsep) * \real{0.1250}}
  >{\raggedright\arraybackslash}p{(\linewidth - 14\tabcolsep) * \real{0.1250}}
  >{\raggedright\arraybackslash}p{(\linewidth - 14\tabcolsep) * \real{0.1250}}@{}}
\toprule\noalign{}
\begin{minipage}[b]{\linewidth}\raggedright
RUNSPEC
\end{minipage} & \begin{minipage}[b]{\linewidth}\raggedright
GRID
\end{minipage} & \begin{minipage}[b]{\linewidth}\raggedright
EDIT
\end{minipage} & \begin{minipage}[b]{\linewidth}\raggedright
PROPS
\end{minipage} & \begin{minipage}[b]{\linewidth}\raggedright
REGIONS
\end{minipage} & \begin{minipage}[b]{\linewidth}\raggedright
SOLUTION
\end{minipage} & \begin{minipage}[b]{\linewidth}\raggedright
SUMMARY
\end{minipage} & \begin{minipage}[b]{\linewidth}\raggedright
SCHEDULE
\end{minipage} \\
\midrule\noalign{}
\endhead
\bottomrule\noalign{}
\endlastfoot
\end{longtable}

\subsubsection{Description}\label{description-40}

This keyword indicates that dissolved gas is present in live {[}``Live''
oil is oil that contains gas in solution, which is normally the case for
most conventional oil reservoirs. However, for oil reservoirs classified
as heavy oil reservoirs, the in situ dissolved gas may be negligible and
oil would then be classified as gas-free oil which is commonly referred
to as ``dead'' oil.{]} oil in the model and the keyword should only be
used if the there is both oil and gas phases in the model. The keyword
may be used for oil-water and oil-water-gas input decks that contain the
oil and gas phases. The keyword will also invoke data input file
checking to ensure that all the required oil and gas phase input
parameters are defined in the input deck.

If the oil has a constant and uniform dissolved gas concentration,
Gas-Oil Ratio (``GOR''), and if the reservoir pressure never drops below
the saturation pressure (bubble point pressure), then the model can be
run more efficiently by omitting the GAS and DISGAS keywords from the
RUNSPEC section, treating the oil as a dead oil {[}``Dead'' oil is oil
that it contains no dissolved gas or a relatively thick oil or residue
that has lost its volatile components.{]}, and defining a constant Rs
(GOR) value with keyword RSCONST or RSCONSTT in the PROPS section. This
results in the model being run with as a dead oil problem with no active
gas phase. However, OPM Flow takes into account the constant Rs in the
calculations and reporting.

There is no data required for this keyword and there is no terminating
``/'' for this keyword.

\subsubsection{Example}\label{example-30}

\begin{verbatim}
--
--       DISSOLVED GAS IN LIVE OIL IS PRESENT IN THE RUN
--
DISGAS

\end{verbatim}

The above example declares that the dissolved gas in the oil phase is
active in the model.

\subsection{DISGASW -- Activate Dissolved Gas in the Water Phase in the
Model}\label{disgasw-activate-dissolved-gas-in-the-water-phase-in-the-model}

\begin{longtable}[]{@{}
  >{\raggedright\arraybackslash}p{(\linewidth - 14\tabcolsep) * \real{0.1250}}
  >{\raggedright\arraybackslash}p{(\linewidth - 14\tabcolsep) * \real{0.1250}}
  >{\raggedright\arraybackslash}p{(\linewidth - 14\tabcolsep) * \real{0.1250}}
  >{\raggedright\arraybackslash}p{(\linewidth - 14\tabcolsep) * \real{0.1250}}
  >{\raggedright\arraybackslash}p{(\linewidth - 14\tabcolsep) * \real{0.1250}}
  >{\raggedright\arraybackslash}p{(\linewidth - 14\tabcolsep) * \real{0.1250}}
  >{\raggedright\arraybackslash}p{(\linewidth - 14\tabcolsep) * \real{0.1250}}
  >{\raggedright\arraybackslash}p{(\linewidth - 14\tabcolsep) * \real{0.1250}}@{}}
\toprule\noalign{}
\begin{minipage}[b]{\linewidth}\raggedright
RUNSPEC
\end{minipage} & \begin{minipage}[b]{\linewidth}\raggedright
GRID
\end{minipage} & \begin{minipage}[b]{\linewidth}\raggedright
EDIT
\end{minipage} & \begin{minipage}[b]{\linewidth}\raggedright
PROPS
\end{minipage} & \begin{minipage}[b]{\linewidth}\raggedright
REGIONS
\end{minipage} & \begin{minipage}[b]{\linewidth}\raggedright
SOLUTION
\end{minipage} & \begin{minipage}[b]{\linewidth}\raggedright
SUMMARY
\end{minipage} & \begin{minipage}[b]{\linewidth}\raggedright
SCHEDULE
\end{minipage} \\
\midrule\noalign{}
\endhead
\bottomrule\noalign{}
\endlastfoot
\end{longtable}

\subsubsection{Description}\label{description-41}

This keyword indicates that dissolved gas is present in the water phase
in the model. The keyword may only be used for gas-water input decks
that contain just the gas and the water phases, which must also be
declared. The keyword will also invoke data input file checking to
ensure that all the required gas and water phase input parameters are
defined in the input deck.

\begin{tcolorbox}[enhanced jigsaw, title=\textcolor{quarto-callout-note-color}{\faInfo}\hspace{0.5em}{Note}, coltitle=black, breakable, toprule=.15mm, colbacktitle=quarto-callout-note-color!10!white, leftrule=.75mm, opacitybacktitle=0.6, colback=white, bottomtitle=1mm, toptitle=1mm, titlerule=0mm, left=2mm, colframe=quarto-callout-note-color-frame, arc=.35mm, rightrule=.15mm, bottomrule=.15mm, opacityback=0]

This is an OPM Flow specific keyword for the simulator's Dissolved Gas
in Water Model that is activated by declaring that this phase is present
in the run.

\end{tcolorbox}

The activation of this phase may be used for modeling the production of
gas from the water in gas fields, where gas has dissolved into the in
situ water phase. Although this is always the case, the volumes are
usually insignificant compared with the volumes in the hydrocarbon
zones, and therefore in nearly all cases the dissolved gas volumes are
usually ignored in terms of resource volumes. For gas fields, the
interest is normally the amount of water in the gas phase, as this
influences the liquid loading rate as the water is liberated from the
gas as the fluid transverses up the wellbore.

The second, and primary, application of this black-oil formulation is
for CO2 storage, where injection of CO2 injection is being conducted in
either depleted gas fields or in water filled aquifer zones. Activating
the dissolved gas in the water phase under these circumstances will
allow the CO2 to dissolve into the water phase based on the entered PVT
properties in the PROPS section. In this case the CO2STORE keyword in
the RUNSPEC section should also be present in the deck.

There is no data required for this keyword and there is no terminating
``/'' for this keyword.

\subsubsection{Example}\label{example-31}

\begin{verbatim}
--
--       GAS PHASE IS PRESENT IN THE RUN
--
GAS
--
--       WATER PHASE IS PRESENT IN THE RUN
--
WATER
--
--       DISSOLVED GAS IN WATER IS PRESENT IN THE RUN (OPM FLOW KEYWORD)
--
DISGASW
--
--       ACTIVATE CO2 STORAGE IN THE MODEL (OPM FLOW CO2 STORAGE KEYWORD)
--
CO2STORE
\end{verbatim}

The above example declares that the gas and water phases are present,
and that the dissolved gas in the water phase is also active in the
model, for modeling CO2 storage.

\subsection{DISPDIMS -- Define the Maximum Number of Dispersion
Tables}\label{dispdims-define-the-maximum-number-of-dispersion-tables}

\begin{longtable}[]{@{}
  >{\raggedright\arraybackslash}p{(\linewidth - 14\tabcolsep) * \real{0.1250}}
  >{\raggedright\arraybackslash}p{(\linewidth - 14\tabcolsep) * \real{0.1250}}
  >{\raggedright\arraybackslash}p{(\linewidth - 14\tabcolsep) * \real{0.1250}}
  >{\raggedright\arraybackslash}p{(\linewidth - 14\tabcolsep) * \real{0.1250}}
  >{\raggedright\arraybackslash}p{(\linewidth - 14\tabcolsep) * \real{0.1250}}
  >{\raggedright\arraybackslash}p{(\linewidth - 14\tabcolsep) * \real{0.1250}}
  >{\raggedright\arraybackslash}p{(\linewidth - 14\tabcolsep) * \real{0.1250}}
  >{\raggedright\arraybackslash}p{(\linewidth - 14\tabcolsep) * \real{0.1250}}@{}}
\toprule\noalign{}
\begin{minipage}[b]{\linewidth}\raggedright
RUNSPEC
\end{minipage} & \begin{minipage}[b]{\linewidth}\raggedright
GRID
\end{minipage} & \begin{minipage}[b]{\linewidth}\raggedright
EDIT
\end{minipage} & \begin{minipage}[b]{\linewidth}\raggedright
PROPS
\end{minipage} & \begin{minipage}[b]{\linewidth}\raggedright
REGIONS
\end{minipage} & \begin{minipage}[b]{\linewidth}\raggedright
SOLUTION
\end{minipage} & \begin{minipage}[b]{\linewidth}\raggedright
SUMMARY
\end{minipage} & \begin{minipage}[b]{\linewidth}\raggedright
SCHEDULE
\end{minipage} \\
\midrule\noalign{}
\endhead
\bottomrule\noalign{}
\endlastfoot
\end{longtable}

\subsubsection{Description}\label{description-42}

The DISPDIMS key defines the maximum number of dispersion tables, and
the maximum number of velocity and concentration elements per table.

This keyword is not supported by OPM Flow but would change the results
if supported so the simulation will be stopped.

\subsection{DUALPERM -- Activate Dual Permeability
Model}\label{dualperm-activate-dual-permeability-model}

\begin{longtable}[]{@{}
  >{\raggedright\arraybackslash}p{(\linewidth - 14\tabcolsep) * \real{0.1250}}
  >{\raggedright\arraybackslash}p{(\linewidth - 14\tabcolsep) * \real{0.1250}}
  >{\raggedright\arraybackslash}p{(\linewidth - 14\tabcolsep) * \real{0.1250}}
  >{\raggedright\arraybackslash}p{(\linewidth - 14\tabcolsep) * \real{0.1250}}
  >{\raggedright\arraybackslash}p{(\linewidth - 14\tabcolsep) * \real{0.1250}}
  >{\raggedright\arraybackslash}p{(\linewidth - 14\tabcolsep) * \real{0.1250}}
  >{\raggedright\arraybackslash}p{(\linewidth - 14\tabcolsep) * \real{0.1250}}
  >{\raggedright\arraybackslash}p{(\linewidth - 14\tabcolsep) * \real{0.1250}}@{}}
\toprule\noalign{}
\begin{minipage}[b]{\linewidth}\raggedright
RUNSPEC
\end{minipage} & \begin{minipage}[b]{\linewidth}\raggedright
GRID
\end{minipage} & \begin{minipage}[b]{\linewidth}\raggedright
EDIT
\end{minipage} & \begin{minipage}[b]{\linewidth}\raggedright
PROPS
\end{minipage} & \begin{minipage}[b]{\linewidth}\raggedright
REGIONS
\end{minipage} & \begin{minipage}[b]{\linewidth}\raggedright
SOLUTION
\end{minipage} & \begin{minipage}[b]{\linewidth}\raggedright
SUMMARY
\end{minipage} & \begin{minipage}[b]{\linewidth}\raggedright
SCHEDULE
\end{minipage} \\
\midrule\noalign{}
\endhead
\bottomrule\noalign{}
\endlastfoot
\end{longtable}

\subsubsection{Description}\label{description-43}

The DUALPERM keyword activates the Dual Permeability option and the Dual
Porosity option for the run. In a dual porosity system flow occurs
between the matrix and the fracture only, whereas in a dual permeability
system flow also occurs between the matrix grid blocks {[}Warren, J.E.
and Root, P.J. 1963. The Behavior of Naturally Fractured Reservoirs. SPE
J. 3 (3): 245--255. SPE-426-PA. http://dx.doi.org/10.2118/426-PA.{]},
{[}Gringarten, A.C. 1984. Interpretation of Tests in Fissured and
Multilayered Reservoirs With Double-Porosity Behavior: Theory and
Practice. J Pet Technol 36 (4): 549-564. SPE-10044-PA.
http://dx.doi.org/10.2118/10044-PA.{]}, {[}Serra, K., Reynolds, A.C.,
and Raghavan, R. 1983. New Pressure Transient Analysis Methods for
Naturally Fractured Reservoirs(includes associated papers 12940 and
13014 ). J Pet Technol 35 (12): 2271-2283. SPE-10780-PA.
http://dx.doi.org/10.2118/10780-PA{]}, {[}Barenblatt, G.E., Zheltov,
I.P., and Kochina, I.N. 1960. Basic Concepts in the Theory of
Homogeneous Liquids in Fissured Rocks. J. Appl. Math. Mech. 24:
1286-1303.{]} and {[}Kazemi, H., Merrill JR., L. S., Porterfield, K. L.,
and Zeman, P. R. ``Numerical Simulation of Water-Oil Flow in Naturally
Fractured Reservoirs,'' paper SPE 5719, Society of Petroleum Engineers
Journal(1976) 16, No.~6, 317-326.{]}.

This keyword is not supported by OPM Flow but would change the results
if supported so the simulation will be stopped.

There is no data required for this keyword and there is no terminating
``/'' for this keyword.

\subsubsection{Example}\label{example-32}

\begin{verbatim}
--
--       ACTIVATE DUAL PERMEABILITY MODEL
--
DUALPERM

\end{verbatim}

The above example declares that the Dual Permeability option is active
for the run.

\subsection{DUALPORO -- Activate Dual Porosity
Model}\label{dualporo-activate-dual-porosity-model}

\begin{longtable}[]{@{}
  >{\raggedright\arraybackslash}p{(\linewidth - 14\tabcolsep) * \real{0.1250}}
  >{\raggedright\arraybackslash}p{(\linewidth - 14\tabcolsep) * \real{0.1250}}
  >{\raggedright\arraybackslash}p{(\linewidth - 14\tabcolsep) * \real{0.1250}}
  >{\raggedright\arraybackslash}p{(\linewidth - 14\tabcolsep) * \real{0.1250}}
  >{\raggedright\arraybackslash}p{(\linewidth - 14\tabcolsep) * \real{0.1250}}
  >{\raggedright\arraybackslash}p{(\linewidth - 14\tabcolsep) * \real{0.1250}}
  >{\raggedright\arraybackslash}p{(\linewidth - 14\tabcolsep) * \real{0.1250}}
  >{\raggedright\arraybackslash}p{(\linewidth - 14\tabcolsep) * \real{0.1250}}@{}}
\toprule\noalign{}
\begin{minipage}[b]{\linewidth}\raggedright
RUNSPEC
\end{minipage} & \begin{minipage}[b]{\linewidth}\raggedright
GRID
\end{minipage} & \begin{minipage}[b]{\linewidth}\raggedright
EDIT
\end{minipage} & \begin{minipage}[b]{\linewidth}\raggedright
PROPS
\end{minipage} & \begin{minipage}[b]{\linewidth}\raggedright
REGIONS
\end{minipage} & \begin{minipage}[b]{\linewidth}\raggedright
SOLUTION
\end{minipage} & \begin{minipage}[b]{\linewidth}\raggedright
SUMMARY
\end{minipage} & \begin{minipage}[b]{\linewidth}\raggedright
SCHEDULE
\end{minipage} \\
\midrule\noalign{}
\endhead
\bottomrule\noalign{}
\endlastfoot
\end{longtable}

\subsubsection{Description}\label{description-44}

The DUALPORO keyword activates the Dual Porosity option for the run. In
a dual porosity system flow occurs between the matrix and the fracture
only, whereas in a dual permeability system flow also occurs between the
matrix grid blocks {[}Warren, J.E. and Root, P.J. 1963. The Behavior of
Naturally Fractured Reservoirs. SPE J. 3 (3): 245--255. SPE-426-PA.
http://dx.doi.org/10.2118/426-PA.{]}, {[}Gringarten, A.C. 1984.
Interpretation of Tests in Fissured and Multilayered Reservoirs With
Double-Porosity Behavior: Theory and Practice. J Pet Technol 36 (4):
549-564. SPE-10044-PA. http://dx.doi.org/10.2118/10044-PA.{]}, {[}Serra,
K., Reynolds, A.C., and Raghavan, R. 1983. New Pressure Transient
Analysis Methods for Naturally Fractured Reservoirs(includes associated
papers 12940 and 13014 ). J Pet Technol 35 (12): 2271-2283.
SPE-10780-PA. http://dx.doi.org/10.2118/10780-PA{]}, {[}Barenblatt,
G.E., Zheltov, I.P., and Kochina, I.N. 1960. Basic Concepts in the
Theory of Homogeneous Liquids in Fissured Rocks. J. Appl. Math. Mech.
24: 1286-1303.{]} and {[}Kazemi, H., Merrill JR., L. S., Porterfield, K.
L., and Zeman, P. R. ``Numerical Simulation of Water-Oil Flow in
Naturally Fractured Reservoirs,'' paper SPE 5719, Society of Petroleum
Engineers Journal(1976) 16, No.~6, 317-326.{]}.

This keyword is not supported by OPM Flow but would change the results
if supported so the simulation will be stopped.

There is no data required for this keyword and there is no terminating
``/'' for this keyword.

\subsubsection{Example}\label{example-33}

\begin{verbatim}
--
--       ACTIVATE DUAL POROSITY MODEL
--
DUALPORO
\end{verbatim}

The above example declares that the Dual Porosity option is active for
the run.

\subsection{DYNRDIMS -- Define Dynamic Region
Dimensions}\label{dynrdims-define-dynamic-region-dimensions}

\begin{longtable}[]{@{}
  >{\raggedright\arraybackslash}p{(\linewidth - 14\tabcolsep) * \real{0.1250}}
  >{\raggedright\arraybackslash}p{(\linewidth - 14\tabcolsep) * \real{0.1250}}
  >{\raggedright\arraybackslash}p{(\linewidth - 14\tabcolsep) * \real{0.1250}}
  >{\raggedright\arraybackslash}p{(\linewidth - 14\tabcolsep) * \real{0.1250}}
  >{\raggedright\arraybackslash}p{(\linewidth - 14\tabcolsep) * \real{0.1250}}
  >{\raggedright\arraybackslash}p{(\linewidth - 14\tabcolsep) * \real{0.1250}}
  >{\raggedright\arraybackslash}p{(\linewidth - 14\tabcolsep) * \real{0.1250}}
  >{\raggedright\arraybackslash}p{(\linewidth - 14\tabcolsep) * \real{0.1250}}@{}}
\toprule\noalign{}
\begin{minipage}[b]{\linewidth}\raggedright
RUNSPEC
\end{minipage} & \begin{minipage}[b]{\linewidth}\raggedright
GRID
\end{minipage} & \begin{minipage}[b]{\linewidth}\raggedright
EDIT
\end{minipage} & \begin{minipage}[b]{\linewidth}\raggedright
PROPS
\end{minipage} & \begin{minipage}[b]{\linewidth}\raggedright
REGIONS
\end{minipage} & \begin{minipage}[b]{\linewidth}\raggedright
SOLUTION
\end{minipage} & \begin{minipage}[b]{\linewidth}\raggedright
SUMMARY
\end{minipage} & \begin{minipage}[b]{\linewidth}\raggedright
SCHEDULE
\end{minipage} \\
\midrule\noalign{}
\endhead
\bottomrule\noalign{}
\endlastfoot
\end{longtable}

\subsubsection{Description}\label{description-45}

The DYNRDIMS keyword defines the dimensions for the parameters used by
the Dynamic Regions facility, including the maximum number of dynamic
regions. The Dynamic Regions facility allows for property and reporting
regions to vary as the run progresses, based on the parameters and logic
defined by the DYNAMICR keyword in the SOLUTION and PROPS section.

This keyword is not supported by OPM Flow but would change the results
if supported so the simulation will be stopped.

\subsection{ECHO -- Activate Echoing of User Input Files to the Print
File}\label{echo-activate-echoing-of-user-input-files-to-the-print-file-1}

Turns on echoing of all the input files to the print file; note that
this keyword is activated by default and can subsequently be switched
off by the NOECHO activation keyword.

See ECHO -- Activate Echoing of User Input Files to the Print File in
the GLOBAL section for a full description.

\subsection{ECLMC -- Activate Multi-Component Brine
Model}\label{eclmc-activate-multi-component-brine-model}

\begin{longtable}[]{@{}
  >{\raggedright\arraybackslash}p{(\linewidth - 14\tabcolsep) * \real{0.1250}}
  >{\raggedright\arraybackslash}p{(\linewidth - 14\tabcolsep) * \real{0.1250}}
  >{\raggedright\arraybackslash}p{(\linewidth - 14\tabcolsep) * \real{0.1250}}
  >{\raggedright\arraybackslash}p{(\linewidth - 14\tabcolsep) * \real{0.1250}}
  >{\raggedright\arraybackslash}p{(\linewidth - 14\tabcolsep) * \real{0.1250}}
  >{\raggedright\arraybackslash}p{(\linewidth - 14\tabcolsep) * \real{0.1250}}
  >{\raggedright\arraybackslash}p{(\linewidth - 14\tabcolsep) * \real{0.1250}}
  >{\raggedright\arraybackslash}p{(\linewidth - 14\tabcolsep) * \real{0.1250}}@{}}
\toprule\noalign{}
\begin{minipage}[b]{\linewidth}\raggedright
RUNSPEC
\end{minipage} & \begin{minipage}[b]{\linewidth}\raggedright
GRID
\end{minipage} & \begin{minipage}[b]{\linewidth}\raggedright
EDIT
\end{minipage} & \begin{minipage}[b]{\linewidth}\raggedright
PROPS
\end{minipage} & \begin{minipage}[b]{\linewidth}\raggedright
REGIONS
\end{minipage} & \begin{minipage}[b]{\linewidth}\raggedright
SOLUTION
\end{minipage} & \begin{minipage}[b]{\linewidth}\raggedright
SUMMARY
\end{minipage} & \begin{minipage}[b]{\linewidth}\raggedright
SCHEDULE
\end{minipage} \\
\midrule\noalign{}
\endhead
\bottomrule\noalign{}
\endlastfoot
\end{longtable}

\subsubsection{Description}\label{description-46}

The ECLMC keyword activates the Multi-Component Brine model that allows
for the water phase to have multiple water salinities. The keyword
should be used in conjunction with the BRINE keyword in the RUNSPEC.
Both keywords must be specified to activate the Multi-Component Brine
model, whereas the BRINE keyword only is required to activate the
standard brine tracking model.

There is no data required for this keyword and there is no terminating
``/'' for this keyword.

This keyword is not supported by OPM Flow but would change the results
if supported so the simulation will be stopped.

\subsubsection{Example}\label{example-34}

The first example activates the standard Brine Tracking model.

\begin{verbatim}
--
--       ACTIVATE STANDARD BRINE MODEL
--
BRINE
\end{verbatim}

The next example shows the ECLMC and BRINE keywords for when the
Multi-Component Brine model is required.

\begin{verbatim}
--
--       ACTIVATE MULTI-COMPONENT BRINE MODEL
--
ECLMC
--
--       DEFINE WATER PHASE MULTI-COMPONENT BRINE COMPONENTS
--
--       SALT1   SALT2   SALT3   SALT4   SALT5
BRINE
         NACL    CACL2   MGC03                                                 /
\end{verbatim}

The above example activates the Multi-Component Brine model with three
different water salinities.

\subsection{END -- Define the End of the Input
File}\label{end-define-the-end-of-the-input-file-1}

This keyword marks the end of the input file and can occur in any
section. Any keywords and data after the END keyword are ignored.

See END -- Define the End of the Input File in the GLOBAL section for a
full description.

\subsection{ENDINC -- Define the End of an Include
File}\label{endinc-define-the-end-of-an-include-file-1}

This keyword marks the end of an include file specified on the INCLUDE
keyword. When the ENDINC keyword is encountered in the INCLUDE file,
input data is read from the next keyword in the current file. Any
keywords and data after the ENDINC keyword in the INCLUDE file are
ignored.

See ENDINC -- Define the End of an Include File in the GLOBAL section
for a full description.

\subsection{ENDSCALE -- Activate Relative Permeability End-Point Scaling
Option}\label{endscale-activate-relative-permeability-end-point-scaling-option}

\begin{longtable}[]{@{}
  >{\raggedright\arraybackslash}p{(\linewidth - 14\tabcolsep) * \real{0.1250}}
  >{\raggedright\arraybackslash}p{(\linewidth - 14\tabcolsep) * \real{0.1250}}
  >{\raggedright\arraybackslash}p{(\linewidth - 14\tabcolsep) * \real{0.1250}}
  >{\raggedright\arraybackslash}p{(\linewidth - 14\tabcolsep) * \real{0.1250}}
  >{\raggedright\arraybackslash}p{(\linewidth - 14\tabcolsep) * \real{0.1250}}
  >{\raggedright\arraybackslash}p{(\linewidth - 14\tabcolsep) * \real{0.1250}}
  >{\raggedright\arraybackslash}p{(\linewidth - 14\tabcolsep) * \real{0.1250}}
  >{\raggedright\arraybackslash}p{(\linewidth - 14\tabcolsep) * \real{0.1250}}@{}}
\toprule\noalign{}
\begin{minipage}[b]{\linewidth}\raggedright
RUNSPEC
\end{minipage} & \begin{minipage}[b]{\linewidth}\raggedright
GRID
\end{minipage} & \begin{minipage}[b]{\linewidth}\raggedright
EDIT
\end{minipage} & \begin{minipage}[b]{\linewidth}\raggedright
PROPS
\end{minipage} & \begin{minipage}[b]{\linewidth}\raggedright
REGIONS
\end{minipage} & \begin{minipage}[b]{\linewidth}\raggedright
SOLUTION
\end{minipage} & \begin{minipage}[b]{\linewidth}\raggedright
SUMMARY
\end{minipage} & \begin{minipage}[b]{\linewidth}\raggedright
SCHEDULE
\end{minipage} \\
\midrule\noalign{}
\endhead
\bottomrule\noalign{}
\endlastfoot
\end{longtable}

\subsubsection{Description}\label{description-47}

The ENDSCALE keyword activates OPM Flow's relative permeability
end-point scaling option. The relative permeability functions are
defined using the either the:

\begin{itemize}
\tightlist
\item
  SWOF, SGOF, SLGOF series of saturation functions, or the
\item
  SWFN, SGFN, SGWFN, SOF2, SOF3, SOF32D series of functions.
\end{itemize}

And are allocated to the grid cells via the SATNUM keyword.

End-point scaling allows the entered relative permeability functions to
be re-scaled based on the saturation end-points allocated to each cell
by the SWL, SWCR, SWU, SGL, SGCR, SGU, SOWCR, and SOGCR saturation grid
arrays for the saturation end-points, and the KRG, KRORG, KRORW and KRW
relative permeability grid cell arrays for the relative permeability
end-point data. In addition end-point scaling may be directional
dependent in which case the directional dependent versions of the
aforementioned arrays should be used, that is SWLX, SWLY and SWLZ
instead of SWL etc. There is also the facility to make the directional
end-point scaling reversible or non-reversible and if the non-reversible
option is selected the non-reversible versions of the aforementioned
arrays should be used, that is SWLX, SWLX-, SWLY, SWLY- SWLZ and SWLZ-,
instead of SWL or the SWLX, SWLY and SWLZ set of keywords.

The keyword also defines the number of saturation end-point tables that
allows for the re-scaling of the saturation functions to be a function
of depth as opposed to being a grid property array. This is accomplished
via the ENKRVD and ENPTVD keywords in the PROPS section.

\begin{longtable}[]{@{}
  >{\raggedright\arraybackslash}p{(\linewidth - 6\tabcolsep) * \real{0.2500}}
  >{\raggedright\arraybackslash}p{(\linewidth - 6\tabcolsep) * \real{0.2500}}
  >{\raggedright\arraybackslash}p{(\linewidth - 6\tabcolsep) * \real{0.2500}}
  >{\raggedright\arraybackslash}p{(\linewidth - 6\tabcolsep) * \real{0.2500}}@{}}
\toprule\noalign{}
\begin{minipage}[b]{\linewidth}\raggedright
No.
\end{minipage} & \begin{minipage}[b]{\linewidth}\raggedright
Name
\end{minipage} & \begin{minipage}[b]{\linewidth}\raggedright
Description
\end{minipage} & \begin{minipage}[b]{\linewidth}\raggedright
Default
\end{minipage} \\
\midrule\noalign{}
\endhead
\bottomrule\noalign{}
\endlastfoot
1 & DIRECT & A character string that activates or deactivates
directional end-point scaling option. If DIRECT is set to NODIR then
directional end-point scaling is switched off and the same saturation
function is used in the x, y and z directions (unless activated
otherwise by the SATOPTS keyword in the RUNSPEC section). In this case
the SWL, SWCR, SWU, SGL, SGCR, SGU, SOWCR and SOGCR saturation grid
arrays and the KRG, KRORG, KRORW and KRW relative permeability grid cell
arrays should be used to enter the grid block end-point data. If DIRECT
is set to DIRECT then directional end-point scaling is switched on and
different saturation functions are used in the x, y and z directions
(unless activated otherwise by the SATOPTS keyword in the RUNSPEC
section). Here the directional form of the SWL, SWCR, SWU, SGL, SGCR,
SGU, SOWCR and SOGCR saturation grid arrays and the KRG, KRORG, KRORW
and KRW relative permeability grid cell arrays should be use to enter
the grid block end-point data. For example SWLX, SWLY and SWLZ for SWL.
Only the default option is supported by OPM Flow. & NODIR \\
2 & IRREVERS & A character string that activates or deactivates
non-reversible end-point scaling option. If IRREVERS is set to REVERS
then the end-point scaling is set to reversible and results in the same
set of end-point arrays being used for flow from the xI to xI + 1
direction as for the flow from the xI to the xI -- 1 for all directions
(x, y and z). Here the SWLX, SWLY and SWLZ series of keywords should be
used instead of SWL type of keywords. Alternatively, if IRREVERS is set
to IRREVERS then the end-point scaling is set to non-reversible and
results in different sets of end-point arrays being applied for flow
from the xI to xI + 1 direction and the xI to the xI -- 1 direction, for
all directions (x, y, z). in this case the SWLX+, SWLX-, SWLY+, SWLY-
SWLZ+ and SWLZ- series of keywords should be utilized instead of SWL or
the SWLX, SWLY and SWLZ set of keywords. Only the default option is
supported by OPM Flow. & REVERS \\
3 & NTENDP & A positive integer that defines the maximum number of
saturation end-point depth tables. The end-point depth tables are used
to re-scale the saturation tables as a function of depth as opposed to
being a grid block property. NTENDP may also be specified on the TABDIMS
keyword, and if specified on both here and on the TABDIMS keyword the
maximum value of the two is used. Only the default option is supported
by OPM Flow. & 1 \\
4 & NNODES & A positive integer the defines the maximum number entries
for saturation end-point depth tables. & 20 \\
5 & MODE & A positive integer that activates the options for temperature
dependent saturation end-point scaling in the commercial compositional
simulator. MODE should be defaulted with either 1* or zero, which means
that scaling can only be performed by grid block end-point scaling
properties or via saturation end-point depth tables. & 0 \\
Notes: & & & \\
\end{longtable}

\emph{Table 5.10: ENDSCALE Keyword Description}

\subsubsection{Example}\label{example-35}

\begin{verbatim}
--       DIRC   REVERSE  MAX     MAX
--       SCALE  SCALE    TABLES  NODES
ENDSCALE
         NODIR  REVERS   1*      1*                                            /
\end{verbatim}

The above example invokes the end-point scaling option with end-point
scaling being non-directional and reversible with the default number of
saturation end-point depth tables (one) with 20 entries per table.

\subsection{ENDSKIP -- Deactivate Skipping of Keywords and Input
Data}\label{endskip-deactivate-skipping-of-keywords-and-input-data-1}

The ENDSKIP keyword deactivates the skipping of keywords that was
activated by the SKIP, SKIP100, or SKIP300 keywords. Each SKIP keyword
should be paired with an ENDSKIP keyword.

See ENDSKIP -- Deactivate Skipping of Keywords and Input Data in the
GLOBAL section for a full description.

\subsection{EOS -- Specify Equation of
State}\label{eos-specify-equation-of-state}

\begin{longtable}[]{@{}
  >{\raggedright\arraybackslash}p{(\linewidth - 14\tabcolsep) * \real{0.1250}}
  >{\raggedright\arraybackslash}p{(\linewidth - 14\tabcolsep) * \real{0.1250}}
  >{\raggedright\arraybackslash}p{(\linewidth - 14\tabcolsep) * \real{0.1250}}
  >{\raggedright\arraybackslash}p{(\linewidth - 14\tabcolsep) * \real{0.1250}}
  >{\raggedright\arraybackslash}p{(\linewidth - 14\tabcolsep) * \real{0.1250}}
  >{\raggedright\arraybackslash}p{(\linewidth - 14\tabcolsep) * \real{0.1250}}
  >{\raggedright\arraybackslash}p{(\linewidth - 14\tabcolsep) * \real{0.1250}}
  >{\raggedright\arraybackslash}p{(\linewidth - 14\tabcolsep) * \real{0.1250}}@{}}
\toprule\noalign{}
\begin{minipage}[b]{\linewidth}\raggedright
RUNSPEC
\end{minipage} & \begin{minipage}[b]{\linewidth}\raggedright
GRID
\end{minipage} & \begin{minipage}[b]{\linewidth}\raggedright
EDIT
\end{minipage} & \begin{minipage}[b]{\linewidth}\raggedright
PROPS
\end{minipage} & \begin{minipage}[b]{\linewidth}\raggedright
REGIONS
\end{minipage} & \begin{minipage}[b]{\linewidth}\raggedright
SOLUTION
\end{minipage} & \begin{minipage}[b]{\linewidth}\raggedright
SUMMARY
\end{minipage} & \begin{minipage}[b]{\linewidth}\raggedright
SCHEDULE
\end{minipage} \\
\midrule\noalign{}
\endhead
\bottomrule\noalign{}
\endlastfoot
\end{longtable}

\subsubsection{Description}\label{description-48}

The \hyperref[REF_HEADING_KEYWORD_EOS_5_3]{EOS} keyword specifies the
Equation Of State (EOS) to be used for each EOS region. The keyword
should only be used if the compositional mode has been requested using
the COMPS keyword in the RUNSPEC section.

See \hyperref[0.0.0.EOS]{EOS -- Specify Equation of State} in the PROPS
section for a full description.

\subsection{EQLDIMS -- Define the Equilibration Data
Dimensions}\label{eqldims-define-the-equilibration-data-dimensions}

\begin{longtable}[]{@{}
  >{\raggedright\arraybackslash}p{(\linewidth - 14\tabcolsep) * \real{0.1250}}
  >{\raggedright\arraybackslash}p{(\linewidth - 14\tabcolsep) * \real{0.1250}}
  >{\raggedright\arraybackslash}p{(\linewidth - 14\tabcolsep) * \real{0.1250}}
  >{\raggedright\arraybackslash}p{(\linewidth - 14\tabcolsep) * \real{0.1250}}
  >{\raggedright\arraybackslash}p{(\linewidth - 14\tabcolsep) * \real{0.1250}}
  >{\raggedright\arraybackslash}p{(\linewidth - 14\tabcolsep) * \real{0.1250}}
  >{\raggedright\arraybackslash}p{(\linewidth - 14\tabcolsep) * \real{0.1250}}
  >{\raggedright\arraybackslash}p{(\linewidth - 14\tabcolsep) * \real{0.1250}}@{}}
\toprule\noalign{}
\begin{minipage}[b]{\linewidth}\raggedright
RUNSPEC
\end{minipage} & \begin{minipage}[b]{\linewidth}\raggedright
GRID
\end{minipage} & \begin{minipage}[b]{\linewidth}\raggedright
EDIT
\end{minipage} & \begin{minipage}[b]{\linewidth}\raggedright
PROPS
\end{minipage} & \begin{minipage}[b]{\linewidth}\raggedright
REGIONS
\end{minipage} & \begin{minipage}[b]{\linewidth}\raggedright
SOLUTION
\end{minipage} & \begin{minipage}[b]{\linewidth}\raggedright
SUMMARY
\end{minipage} & \begin{minipage}[b]{\linewidth}\raggedright
SCHEDULE
\end{minipage} \\
\midrule\noalign{}
\endhead
\bottomrule\noalign{}
\endlastfoot
\end{longtable}

\subsubsection{Description}\label{description-49}

The EQLDIMS keyword defines the maximum number of properties associated
with equilibrating the model, that is initializing the model. A
reservoir grid can be separated into separate regions in order to apply
different pressure regimes and/or fluid contacts. Care should be taken
that the different regions are not in communication if the pressures or
fluid contacts are different for the various regions, as this would lead
to an unstable initialization and would also imply errors in the model
description as implemented.

\begin{longtable}[]{@{}
  >{\raggedright\arraybackslash}p{(\linewidth - 6\tabcolsep) * \real{0.2500}}
  >{\raggedright\arraybackslash}p{(\linewidth - 6\tabcolsep) * \real{0.2500}}
  >{\raggedright\arraybackslash}p{(\linewidth - 6\tabcolsep) * \real{0.2500}}
  >{\raggedright\arraybackslash}p{(\linewidth - 6\tabcolsep) * \real{0.2500}}@{}}
\toprule\noalign{}
\begin{minipage}[b]{\linewidth}\raggedright
No.
\end{minipage} & \begin{minipage}[b]{\linewidth}\raggedright
Name
\end{minipage} & \begin{minipage}[b]{\linewidth}\raggedright
Description
\end{minipage} & \begin{minipage}[b]{\linewidth}\raggedright
Default
\end{minipage} \\
\midrule\noalign{}
\endhead
\bottomrule\noalign{}
\endlastfoot
1 & NTEQUL & A positive integer value that defines the number of
equilibration regions entered using the EQLNUM keyword in the REGIONS
section and the number of entries associated with the EQUIL keyword in
the SOLUTION section. & 1 \\
2 & NPRSVD & A positive integer value setting the number of pressure
versus depth entries used by OPM Flow in determining equilibration
parameters. Unless there is a requirement for a very fine equilibration
this parameter should be defaulted. & 100 \\
3 & NDRXVD & A positive integer value that defines the maximum number of
depth entries in equilibration property versus depth tables (RSVD, RVVD,
PBVD or PDVD etc.) as defined in the SOLUTION section. & 20 \\
4 & NTTRVD & A positive integer that defines the maximum number of TVDP
tables that describe the initial tracer concentration versus depth. &
1 \\
5 & NSTRVD & A positive integer that defines the maximum number of depth
entries in the TVDP tables as described in (4) & 20 \\
Notes: & & & \\
\end{longtable}

\emph{Table 5.11: EQLDIMS Keyword Description}

It is common that the EQLNUM and FIPNUM arrays are identical so that the
fluid in-place reporting matches the equilibration regions. Thus, in
order to avoid errors in this case, one should just use one array (say
the FIPNUM property array) and use the COPY keyword to generate the
EQLNUM array.

\subsubsection{Example}\label{example-36}

\begin{verbatim}
--
--       MAX     MAX     RSVD    TVDP    TVDP
--       EQLNUM  DEPTH   NODES   TABLE   NODES
EQLDIMS
         9       1*      20      1*      1*                                    /
\end{verbatim}

The above example defines nine equilibration regions the default values
for the remaining parameters on the EQLDIMS keyword.

\subsection{EQLOPTS -- Activate the Equilibration
Options}\label{eqlopts-activate-the-equilibration-options}

\begin{longtable}[]{@{}
  >{\raggedright\arraybackslash}p{(\linewidth - 14\tabcolsep) * \real{0.1250}}
  >{\raggedright\arraybackslash}p{(\linewidth - 14\tabcolsep) * \real{0.1250}}
  >{\raggedright\arraybackslash}p{(\linewidth - 14\tabcolsep) * \real{0.1250}}
  >{\raggedright\arraybackslash}p{(\linewidth - 14\tabcolsep) * \real{0.1250}}
  >{\raggedright\arraybackslash}p{(\linewidth - 14\tabcolsep) * \real{0.1250}}
  >{\raggedright\arraybackslash}p{(\linewidth - 14\tabcolsep) * \real{0.1250}}
  >{\raggedright\arraybackslash}p{(\linewidth - 14\tabcolsep) * \real{0.1250}}
  >{\raggedright\arraybackslash}p{(\linewidth - 14\tabcolsep) * \real{0.1250}}@{}}
\toprule\noalign{}
\begin{minipage}[b]{\linewidth}\raggedright
RUNSPEC
\end{minipage} & \begin{minipage}[b]{\linewidth}\raggedright
GRID
\end{minipage} & \begin{minipage}[b]{\linewidth}\raggedright
EDIT
\end{minipage} & \begin{minipage}[b]{\linewidth}\raggedright
PROPS
\end{minipage} & \begin{minipage}[b]{\linewidth}\raggedright
REGIONS
\end{minipage} & \begin{minipage}[b]{\linewidth}\raggedright
SOLUTION
\end{minipage} & \begin{minipage}[b]{\linewidth}\raggedright
SUMMARY
\end{minipage} & \begin{minipage}[b]{\linewidth}\raggedright
SCHEDULE
\end{minipage} \\
\midrule\noalign{}
\endhead
\bottomrule\noalign{}
\endlastfoot
\end{longtable}

\subsubsection{Description}\label{description-50}

The EQLOPTS keyword defines the equilibration options by stating the
character command to activate an option to be used for initializing the
model. Multiple commands may be utilized to activate several
equilibration options following the keyword.

\begin{longtable}[]{@{}
  >{\raggedright\arraybackslash}p{(\linewidth - 6\tabcolsep) * \real{0.2500}}
  >{\raggedright\arraybackslash}p{(\linewidth - 6\tabcolsep) * \real{0.2500}}
  >{\raggedright\arraybackslash}p{(\linewidth - 6\tabcolsep) * \real{0.2500}}
  >{\raggedright\arraybackslash}p{(\linewidth - 6\tabcolsep) * \real{0.2500}}@{}}
\toprule\noalign{}
\begin{minipage}[b]{\linewidth}\raggedright
No.
\end{minipage} & \begin{minipage}[b]{\linewidth}\raggedright
Name
\end{minipage} & \begin{minipage}[b]{\linewidth}\raggedright
Description
\end{minipage} & \begin{minipage}[b]{\linewidth}\raggedright
Default
\end{minipage} \\
\midrule\noalign{}
\endhead
\bottomrule\noalign{}
\endlastfoot
1 & MOBILE & A character string that activates the mobile fluid critical
saturation end point correction. If the MOBILE command is stated then
this option is activated. This option is not supported by OPM Flow but
would change the results if supported so the simulation will be stopped.
Therefore this character string should be omitted. & None \\
2 & QUIESC & A character string that activates the initial quiescence
option that modifies the equilibrium calculated phase pressures to
ensure that a steady state solution is obtained. This options ensures
that there is no flow potential between the grid blocks in a given
region, which is the normal case when block-centered equilibration is
used by setting BOINIT on the EQUIL keyword to zero in the SOLUTION
section. If the QUIESC command is stated then this option is activated.
This option is not supported by OPM Flow but would change the results if
supported so the simulation will be stopped. Therefore this character
string should be omitted. & None \\
3 & THPRES & A character string that activates the inter-region
equilibration flow option. This option allows for a threshold pressure
variable entered via the THPRES keyword to define a pressure which
prevents flow between regions until the THPRES value between regions is
exceeded. If the THPRES command is stated then this option is activated.
& None \\
4 & IRREVERS & A character string that activates the irreversible
inter-region equilibration flow option. This option can only be invoked
if the THPRES command has been stated. The option allows for different
THPRES values for different directions. If the IRREVERS command is
stated then this option is activated. This option is not supported by
OPM Flow but would change the results if supported so the simulation
will be stopped. Therefore this character string should be omitted. &
None \\
Notes: & & & \\
\end{longtable}

\emph{Table 5.3.37.1: EQLOPTS Keyword Description}

\subsubsection{Example}\label{example-37}

\begin{verbatim}
--
--       ACTIVATE EQUILIBRATION OPTIONS
--       MOBILE END-POINT(MOBILE) STEADY STATE(QUIESC) THRESHOLD(THPRES)
--       IRREVERSIBLE THRESHOLD(IRREVERS)
EQLOPTS
         'THPRES'  'IRREVERS'                                                   /
\end{verbatim}

The above example activates the threshold pressure option with different
threshold pressure for different directions.

\subsection{EXTRAPMS -- Activate Extrapolation Warning
Messages}\label{extrapms-activate-extrapolation-warning-messages-1}

The EXTRAPMS keyword activates extrapolation warning messages for when
OPM Flow extrapolates the PVT or VFP tables. Frequent extrapolation
warning messages should be investigated and resolved as this would
indicate possible incorrect data and may result in the simulator
extrapolating to unrealistic values.

See EXTRAPMS -- Activate Extrapolation Warning Messages in the GLOBAL
section for a full description.

\subsection{FAULTDIM -- Define the Number of Fault
Segments}\label{faultdim-define-the-number-of-fault-segments}

\begin{longtable}[]{@{}
  >{\raggedright\arraybackslash}p{(\linewidth - 14\tabcolsep) * \real{0.1250}}
  >{\raggedright\arraybackslash}p{(\linewidth - 14\tabcolsep) * \real{0.1250}}
  >{\raggedright\arraybackslash}p{(\linewidth - 14\tabcolsep) * \real{0.1250}}
  >{\raggedright\arraybackslash}p{(\linewidth - 14\tabcolsep) * \real{0.1250}}
  >{\raggedright\arraybackslash}p{(\linewidth - 14\tabcolsep) * \real{0.1250}}
  >{\raggedright\arraybackslash}p{(\linewidth - 14\tabcolsep) * \real{0.1250}}
  >{\raggedright\arraybackslash}p{(\linewidth - 14\tabcolsep) * \real{0.1250}}
  >{\raggedright\arraybackslash}p{(\linewidth - 14\tabcolsep) * \real{0.1250}}@{}}
\toprule\noalign{}
\begin{minipage}[b]{\linewidth}\raggedright
RUNSPEC
\end{minipage} & \begin{minipage}[b]{\linewidth}\raggedright
GRID
\end{minipage} & \begin{minipage}[b]{\linewidth}\raggedright
EDIT
\end{minipage} & \begin{minipage}[b]{\linewidth}\raggedright
PROPS
\end{minipage} & \begin{minipage}[b]{\linewidth}\raggedright
REGIONS
\end{minipage} & \begin{minipage}[b]{\linewidth}\raggedright
SOLUTION
\end{minipage} & \begin{minipage}[b]{\linewidth}\raggedright
SUMMARY
\end{minipage} & \begin{minipage}[b]{\linewidth}\raggedright
SCHEDULE
\end{minipage} \\
\midrule\noalign{}
\endhead
\bottomrule\noalign{}
\endlastfoot
\end{longtable}

\subsubsection{Description}\label{description-51}

The FAULTDIM keyword defines the maximum number of records (or segments)
that can be entered with the FAULTS keyword. The FAULTS keyword defines
the faults in the grid than can be used for setting (or re-setting)
transmissibility barriers across the fault planes.

\begin{longtable}[]{@{}
  >{\raggedright\arraybackslash}p{(\linewidth - 6\tabcolsep) * \real{0.2500}}
  >{\raggedright\arraybackslash}p{(\linewidth - 6\tabcolsep) * \real{0.2500}}
  >{\raggedright\arraybackslash}p{(\linewidth - 6\tabcolsep) * \real{0.2500}}
  >{\raggedright\arraybackslash}p{(\linewidth - 6\tabcolsep) * \real{0.2500}}@{}}
\toprule\noalign{}
\begin{minipage}[b]{\linewidth}\raggedright
No.
\end{minipage} & \begin{minipage}[b]{\linewidth}\raggedright
Name
\end{minipage} & \begin{minipage}[b]{\linewidth}\raggedright
Description
\end{minipage} & \begin{minipage}[b]{\linewidth}\raggedright
Default
\end{minipage} \\
\midrule\noalign{}
\endhead
\bottomrule\noalign{}
\endlastfoot
1 & MFSEGS & A positive integer value that defines the maximum number of
records (segments) for the FAULTS keyword. & 0 \\
Notes: & & & \\
\end{longtable}

\emph{Table 5.13: FAULTDIM Keyword Description}

\subsubsection{Example}\label{example-38}

\begin{verbatim}
--
--       FAULT
--       SEGMS
--
FAULTDIM
         10000                                                                 /
\end{verbatim}

The above example defines the maximum number of records that can be
entered using the FAULT keyword to be 10,000 segments.

\subsection{FIELD -- Activate the Oil Field System of Units for the
Model}\label{field-activate-the-oil-field-system-of-units-for-the-model}

\begin{longtable}[]{@{}
  >{\raggedright\arraybackslash}p{(\linewidth - 14\tabcolsep) * \real{0.1250}}
  >{\raggedright\arraybackslash}p{(\linewidth - 14\tabcolsep) * \real{0.1250}}
  >{\raggedright\arraybackslash}p{(\linewidth - 14\tabcolsep) * \real{0.1250}}
  >{\raggedright\arraybackslash}p{(\linewidth - 14\tabcolsep) * \real{0.1250}}
  >{\raggedright\arraybackslash}p{(\linewidth - 14\tabcolsep) * \real{0.1250}}
  >{\raggedright\arraybackslash}p{(\linewidth - 14\tabcolsep) * \real{0.1250}}
  >{\raggedright\arraybackslash}p{(\linewidth - 14\tabcolsep) * \real{0.1250}}
  >{\raggedright\arraybackslash}p{(\linewidth - 14\tabcolsep) * \real{0.1250}}@{}}
\toprule\noalign{}
\begin{minipage}[b]{\linewidth}\raggedright
RUNSPEC
\end{minipage} & \begin{minipage}[b]{\linewidth}\raggedright
GRID
\end{minipage} & \begin{minipage}[b]{\linewidth}\raggedright
EDIT
\end{minipage} & \begin{minipage}[b]{\linewidth}\raggedright
PROPS
\end{minipage} & \begin{minipage}[b]{\linewidth}\raggedright
REGIONS
\end{minipage} & \begin{minipage}[b]{\linewidth}\raggedright
SOLUTION
\end{minipage} & \begin{minipage}[b]{\linewidth}\raggedright
SUMMARY
\end{minipage} & \begin{minipage}[b]{\linewidth}\raggedright
SCHEDULE
\end{minipage} \\
\midrule\noalign{}
\endhead
\bottomrule\noalign{}
\endlastfoot
\end{longtable}

\subsubsection{Description}\label{description-52}

This keyword switches on the oil FIELD system of units for the model.

OPM Flow has three sets of units, namely: METRIC, FIELD and LAB and one
of these keyword should be invoked in the RUNSPEC section to avoid any
ambiguity. Both the simulator input and output units are controlled by
including one of the METRIC, FIELD or LAB keywords in the RUNSPEC
section of the input file.

There is no data required for this keyword and there is no terminating
``/'' for this keyword.

\subsubsection{Example}\label{example-39}

\begin{verbatim}
--
--       SWITCH ON THE FIELD SYSTEM OF UNITS FOR BOTH INPUT AND OUTPUT
--
FIELD
\end{verbatim}

The above example switches on the FIELD system of units for the model.

\subsection{FMTHMD -- Activate The Format History Match Gradient File
Option}\label{fmthmd-activate-the-format-history-match-gradient-file-option}

\begin{longtable}[]{@{}
  >{\raggedright\arraybackslash}p{(\linewidth - 14\tabcolsep) * \real{0.1250}}
  >{\raggedright\arraybackslash}p{(\linewidth - 14\tabcolsep) * \real{0.1250}}
  >{\raggedright\arraybackslash}p{(\linewidth - 14\tabcolsep) * \real{0.1250}}
  >{\raggedright\arraybackslash}p{(\linewidth - 14\tabcolsep) * \real{0.1250}}
  >{\raggedright\arraybackslash}p{(\linewidth - 14\tabcolsep) * \real{0.1250}}
  >{\raggedright\arraybackslash}p{(\linewidth - 14\tabcolsep) * \real{0.1250}}
  >{\raggedright\arraybackslash}p{(\linewidth - 14\tabcolsep) * \real{0.1250}}
  >{\raggedright\arraybackslash}p{(\linewidth - 14\tabcolsep) * \real{0.1250}}@{}}
\toprule\noalign{}
\begin{minipage}[b]{\linewidth}\raggedright
RUNSPEC
\end{minipage} & \begin{minipage}[b]{\linewidth}\raggedright
GRID
\end{minipage} & \begin{minipage}[b]{\linewidth}\raggedright
EDIT
\end{minipage} & \begin{minipage}[b]{\linewidth}\raggedright
PROPS
\end{minipage} & \begin{minipage}[b]{\linewidth}\raggedright
REGIONS
\end{minipage} & \begin{minipage}[b]{\linewidth}\raggedright
SOLUTION
\end{minipage} & \begin{minipage}[b]{\linewidth}\raggedright
SUMMARY
\end{minipage} & \begin{minipage}[b]{\linewidth}\raggedright
SCHEDULE
\end{minipage} \\
\midrule\noalign{}
\endhead
\bottomrule\noalign{}
\endlastfoot
\end{longtable}

\subsubsection{Description}\label{description-53}

This keyword switches on formatted output for the History Match Gradient
output file, for when the History Match Gradient option has been
activated by the HMDIMS keyword in the RUNSPEC section.

This keyword is not supported by OPM Flow but would change the results
if supported so the simulation will be stopped.

\subsection{FMTIN -- Activate The Format Input File
Option}\label{fmtin-activate-the-format-input-file-option}

\begin{longtable}[]{@{}
  >{\raggedright\arraybackslash}p{(\linewidth - 14\tabcolsep) * \real{0.1250}}
  >{\raggedright\arraybackslash}p{(\linewidth - 14\tabcolsep) * \real{0.1250}}
  >{\raggedright\arraybackslash}p{(\linewidth - 14\tabcolsep) * \real{0.1250}}
  >{\raggedright\arraybackslash}p{(\linewidth - 14\tabcolsep) * \real{0.1250}}
  >{\raggedright\arraybackslash}p{(\linewidth - 14\tabcolsep) * \real{0.1250}}
  >{\raggedright\arraybackslash}p{(\linewidth - 14\tabcolsep) * \real{0.1250}}
  >{\raggedright\arraybackslash}p{(\linewidth - 14\tabcolsep) * \real{0.1250}}
  >{\raggedright\arraybackslash}p{(\linewidth - 14\tabcolsep) * \real{0.1250}}@{}}
\toprule\noalign{}
\begin{minipage}[b]{\linewidth}\raggedright
RUNSPEC
\end{minipage} & \begin{minipage}[b]{\linewidth}\raggedright
GRID
\end{minipage} & \begin{minipage}[b]{\linewidth}\raggedright
EDIT
\end{minipage} & \begin{minipage}[b]{\linewidth}\raggedright
PROPS
\end{minipage} & \begin{minipage}[b]{\linewidth}\raggedright
REGIONS
\end{minipage} & \begin{minipage}[b]{\linewidth}\raggedright
SOLUTION
\end{minipage} & \begin{minipage}[b]{\linewidth}\raggedright
SUMMARY
\end{minipage} & \begin{minipage}[b]{\linewidth}\raggedright
SCHEDULE
\end{minipage} \\
\midrule\noalign{}
\endhead
\bottomrule\noalign{}
\endlastfoot
\end{longtable}

\subsubsection{Description}\label{description-54}

This keyword switches on the Format Input Files option for all input
files. Similar to the commercial simulator, OPM Flow has various options
for reading various input files and writing the resulting OPM Flow
output files as described in Table 5.14.

\begin{longtable}[]{@{}
  >{\raggedright\arraybackslash}p{(\linewidth - 6\tabcolsep) * \real{0.2500}}
  >{\raggedright\arraybackslash}p{(\linewidth - 6\tabcolsep) * \real{0.2500}}
  >{\raggedright\arraybackslash}p{(\linewidth - 6\tabcolsep) * \real{0.2500}}
  >{\raggedright\arraybackslash}p{(\linewidth - 6\tabcolsep) * \real{0.2500}}@{}}
\toprule\noalign{}
\begin{minipage}[b]{\linewidth}\raggedright
Process
\end{minipage} & \begin{minipage}[b]{\linewidth}\raggedright
Keyword
\end{minipage} & \begin{minipage}[b]{\linewidth}\raggedright
Description
\end{minipage} & \begin{minipage}[b]{\linewidth}\raggedright
Files
\end{minipage} \\
\midrule\noalign{}
\endhead
\bottomrule\noalign{}
\endlastfoot
Input & FMTIN & A character string that defines the input files to be
formatted as ASCII i.e.~text files, as oppose to binary files. The input
deck file is always of this type. The option relates to the OPM Flow
derived files that used as input, for for example when restarting from
another case. If the keyword is omitted then the default is for binary
file input. & \emph{.FEGRID }.FINSPEC \emph{.FINIT }.FRSSPEC
\emph{.FUNRST }.FSMSPEC *.FUNSMRY \\
MULTIN & A character string that defines the input files to be
non-unified multiple files, as opposed to unified files. In this case,
one file is read in per reporting time step, as opposed to all time
steps reports being read from one file. If the keyword is omitted then
the default is for one file per reporting time step. & \emph{.RSSPEC
}.X0001 \emph{.SMSPEC }.S0001 & \\
UNIFIN & A character string that defines the input files to be unified
files, as opposed to non-unified multiple files. A unified file is a
single file containing output for each reporting time step. For this
option a single summary file and a single restart file will be read. If
the keyword is omitted then the default is for one file per reporting
time step. & \emph{.RSSPEC }.UNRST \emph{.SMSPEC }.UNSMRY & \\
Output & FMTOUT & A character string that sets all output files to be
formatted as ASCII i.e.~text files, as opposed to binary files. The
\emph{.PRT, }.LOG and *.DBG files are always of this type. The option
relates to the OPM Flow output files only. In this case the files will
be portable across operating systems, but will also be very large in
terms of hard disk space. For this reason it is recommend that the
default option is used so that binary files are outputted. If the
keyword is omitted then the default is for binary file input. &
\emph{.FEGRID }.FINSPEC \emph{.FINIT }.FRSSPEC \emph{.FUNRST }.FSMSPEC
*.FUNSMRY \\
MULTOUT & A character string that defines the output files to be
non-unified multiple files, as opposed to unified files. In this case,
one file is written for each reporting time step, as opposed to all time
steps reports being written in one file. If the keyword is omitted then
the default is for one file per reporting time step. & \emph{.RSSPEC
}.X0001 \emph{.SMSPEC }.S0001 & \\
UNIFOUT & A character string that defines the output files to be unified
files, as opposed to non-unified multiple files. A unified file is a
single file containing output for each reporting time step. Here a
single summary file and a single restart file will be generated, as
opposed to one file per report time step. If the keyword is omitted then
the default is for one file per reporting time step. & \emph{.RSSPEC
}.UNRST \emph{.SMSPEC }.UNSMRY & \\
Notes: & & & \\
\end{longtable}

\emph{Table 5.14: FMTIN Keyword Description}

There is no data required for this keyword.

See also OPM FLOW OUTPUT FILE FORMATS for a more detailed description of
the various file types (ASCII or binary) and file structure formats
(unified or non-unified formats).

\subsubsection{Example}\label{example-40}

\begin{verbatim}
--
--       SWITCH ON THE FORMAT INPUT FILES OPTION
--
FMTIN
\end{verbatim}

The above example switches on the format input file option.

\subsection{FMTOUT -- Activate The Format Output File
Option}\label{fmtout-activate-the-format-output-file-option}

\begin{longtable}[]{@{}
  >{\raggedright\arraybackslash}p{(\linewidth - 14\tabcolsep) * \real{0.1250}}
  >{\raggedright\arraybackslash}p{(\linewidth - 14\tabcolsep) * \real{0.1250}}
  >{\raggedright\arraybackslash}p{(\linewidth - 14\tabcolsep) * \real{0.1250}}
  >{\raggedright\arraybackslash}p{(\linewidth - 14\tabcolsep) * \real{0.1250}}
  >{\raggedright\arraybackslash}p{(\linewidth - 14\tabcolsep) * \real{0.1250}}
  >{\raggedright\arraybackslash}p{(\linewidth - 14\tabcolsep) * \real{0.1250}}
  >{\raggedright\arraybackslash}p{(\linewidth - 14\tabcolsep) * \real{0.1250}}
  >{\raggedright\arraybackslash}p{(\linewidth - 14\tabcolsep) * \real{0.1250}}@{}}
\toprule\noalign{}
\begin{minipage}[b]{\linewidth}\raggedright
RUNSPEC
\end{minipage} & \begin{minipage}[b]{\linewidth}\raggedright
GRID
\end{minipage} & \begin{minipage}[b]{\linewidth}\raggedright
EDIT
\end{minipage} & \begin{minipage}[b]{\linewidth}\raggedright
PROPS
\end{minipage} & \begin{minipage}[b]{\linewidth}\raggedright
REGIONS
\end{minipage} & \begin{minipage}[b]{\linewidth}\raggedright
SOLUTION
\end{minipage} & \begin{minipage}[b]{\linewidth}\raggedright
SUMMARY
\end{minipage} & \begin{minipage}[b]{\linewidth}\raggedright
SCHEDULE
\end{minipage} \\
\midrule\noalign{}
\endhead
\bottomrule\noalign{}
\endlastfoot
\end{longtable}

\subsubsection{Description}\label{description-55}

This keyword switches on the Format Output Files option for all output
files. Similar to the commercial simulator, OPM Flow has various options
for reading various input files and writing the resulting OPM Flow
output files as described in Table 5.15.

\begin{longtable}[]{@{}
  >{\raggedright\arraybackslash}p{(\linewidth - 6\tabcolsep) * \real{0.2500}}
  >{\raggedright\arraybackslash}p{(\linewidth - 6\tabcolsep) * \real{0.2500}}
  >{\raggedright\arraybackslash}p{(\linewidth - 6\tabcolsep) * \real{0.2500}}
  >{\raggedright\arraybackslash}p{(\linewidth - 6\tabcolsep) * \real{0.2500}}@{}}
\toprule\noalign{}
\begin{minipage}[b]{\linewidth}\raggedright
Process
\end{minipage} & \begin{minipage}[b]{\linewidth}\raggedright
Keyword
\end{minipage} & \begin{minipage}[b]{\linewidth}\raggedright
Description
\end{minipage} & \begin{minipage}[b]{\linewidth}\raggedright
Files
\end{minipage} \\
\midrule\noalign{}
\endhead
\bottomrule\noalign{}
\endlastfoot
Input & FMTIN & A character string that defines the input files to be
formatted as ASCII i.e.~text files, as oppose to binary files. The input
deck file is always of this type. The option relates to the OPM Flow
derived files that used as input, for for example when restarting from
another case. If the keyword is omitted then the default is for binary
file input. & \emph{.FEGRID }.FINSPEC \emph{.FINIT }.FRSSPEC
\emph{.FUNRST }.FSMSPEC *.FUNSMRY \\
MULTIN & A character string that defines the input files to be
non-unified multiple files, as opposed to unified files. In this case,
one file is read in per reporting time step, as opposed to all time
steps reports being read from one file. If the keyword is omitted then
the default is for one file per reporting time step. & \emph{.RSSPEC
}.X0001 \emph{.SMSPEC }.S0001 & \\
UNIFIN & A character string that defines the input files to be unified
files, as opposed to non-unified multiple files. A unified file is a
single file containing output for each reporting time step. For this
option a single summary file and a single restart file will be read. If
the keyword is omitted then the default is for one file per reporting
time step. & \emph{.RSSPEC }.UNRST \emph{.SMSPEC }.UNSMRY & \\
Output & FMTOUT & A character string that sets all output files to be
formatted as ASCII i.e.~text files, as opposed to binary files. The
\emph{.PRT, }.LOG and *.DBG files are always of this type. The option
relates to the OPM Flow output files only. In this case the files will
be portable across operating systems, but will also be very large in
terms of hard disk space. For this reason it is recommend that the
default option is used so that binary files are outputted. If the
keyword is omitted then the default is for binary file input. &
\emph{.FEGRID }.FINSPEC \emph{.FINIT }.FRSSPEC \emph{.FUNRST }.FSMSPEC
*.FUNSMRY \\
MULTOUT & A character string that defines the output files to be
non-unified multiple files, as opposed to unified files. In this case,
one file is written for each reporting time step, as opposed to all time
steps reports being written in one file. If the keyword is omitted then
the default is for one file per reporting time step. & \emph{.RSSPEC
}.X0001 \emph{.SMSPEC }.S0001 & \\
UNIFOUT & A character string that defines the output files to be unified
files, as opposed to non-unified multiple files. A unified file is a
single file containing output for each reporting time step. Here a
single summary file and a single restart file will be generated, as
opposed to one file per report time step. If the keyword is omitted then
the default is for one file per reporting time step. & \emph{.RSSPEC
}.UNRST \emph{.SMSPEC }.UNSMRY & \\
Notes: & & & \\
\end{longtable}

\emph{Table 5.15: FMOUT Keyword Description}

There is no data required for this keyword.

See also OPM FLOW OUTPUT FILE FORMATS for a more detailed description of
the various file types (ASCII or binary) and file structure formats
(unified or non-unified formats).

\subsubsection{Example}\label{example-41}

\begin{verbatim}
--
--       SWITCH ON THE FORMAT OUTPUT FILES OPTION
--
FMTOUT
\end{verbatim}

The above example switches on the format output file option.

\subsection{FOAM -- Activate the Foam Phase and
Model}\label{foam-activate-the-foam-phase-and-model}

\begin{longtable}[]{@{}
  >{\raggedright\arraybackslash}p{(\linewidth - 14\tabcolsep) * \real{0.1250}}
  >{\raggedright\arraybackslash}p{(\linewidth - 14\tabcolsep) * \real{0.1250}}
  >{\raggedright\arraybackslash}p{(\linewidth - 14\tabcolsep) * \real{0.1250}}
  >{\raggedright\arraybackslash}p{(\linewidth - 14\tabcolsep) * \real{0.1250}}
  >{\raggedright\arraybackslash}p{(\linewidth - 14\tabcolsep) * \real{0.1250}}
  >{\raggedright\arraybackslash}p{(\linewidth - 14\tabcolsep) * \real{0.1250}}
  >{\raggedright\arraybackslash}p{(\linewidth - 14\tabcolsep) * \real{0.1250}}
  >{\raggedright\arraybackslash}p{(\linewidth - 14\tabcolsep) * \real{0.1250}}@{}}
\toprule\noalign{}
\begin{minipage}[b]{\linewidth}\raggedright
RUNSPEC
\end{minipage} & \begin{minipage}[b]{\linewidth}\raggedright
GRID
\end{minipage} & \begin{minipage}[b]{\linewidth}\raggedright
EDIT
\end{minipage} & \begin{minipage}[b]{\linewidth}\raggedright
PROPS
\end{minipage} & \begin{minipage}[b]{\linewidth}\raggedright
REGIONS
\end{minipage} & \begin{minipage}[b]{\linewidth}\raggedright
SOLUTION
\end{minipage} & \begin{minipage}[b]{\linewidth}\raggedright
SUMMARY
\end{minipage} & \begin{minipage}[b]{\linewidth}\raggedright
SCHEDULE
\end{minipage} \\
\midrule\noalign{}
\endhead
\bottomrule\noalign{}
\endlastfoot
\end{longtable}

\subsubsection{Description}\label{description-56}

This keyword activates the foam phase and modeling option. The keyword
will also invoke data input file checking to ensure that all the
required foam phase input parameters are defined in the input deck. Note
in the commercial simulator the FOAM phase and model can be used in
conjunction with the POLYMER and SURFACT phases; this is not the case
for OPM Flow. OPM Flow's FOAM phase and model is a standalone
implementation and cannot be used in conjunction with either the POLYMER
or SURFACT phases.

Foam flooding is an enhanced oil recovery flood process that attempts to
control injected gas breakthrough in an oil reservoir by changing the
mobility of the injected fluid.

There is no data required for this keyword and there is no terminating
``/'' for this keyword.

\subsubsection{Example}\label{example-42}

\begin{verbatim}
--
--       ACTIVATE THE FOAM PHASE IN THE MODEL
--
FOAM
\end{verbatim}

The above example declares that the foam phase is active in the model.

\subsection{FORMFEED -- Defined the Print File Form-Feed
Character}\label{formfeed-defined-the-print-file-form-feed-character-1}

The FORMFEED keyword defines the form-feed character, or carriage
control character, for the output print (\emph{.PRT) run summary (}.RSM)
files. The keyword should be place at the very top of the input file.

See FORMFEED -- Defined the Print File Form-Feed Character in the GLOBAL
section for a dull description.

\subsection{FRICTION -- Activate Wellbore Friction
Option}\label{friction-activate-wellbore-friction-option}

\begin{longtable}[]{@{}
  >{\raggedright\arraybackslash}p{(\linewidth - 14\tabcolsep) * \real{0.1250}}
  >{\raggedright\arraybackslash}p{(\linewidth - 14\tabcolsep) * \real{0.1250}}
  >{\raggedright\arraybackslash}p{(\linewidth - 14\tabcolsep) * \real{0.1250}}
  >{\raggedright\arraybackslash}p{(\linewidth - 14\tabcolsep) * \real{0.1250}}
  >{\raggedright\arraybackslash}p{(\linewidth - 14\tabcolsep) * \real{0.1250}}
  >{\raggedright\arraybackslash}p{(\linewidth - 14\tabcolsep) * \real{0.1250}}
  >{\raggedright\arraybackslash}p{(\linewidth - 14\tabcolsep) * \real{0.1250}}
  >{\raggedright\arraybackslash}p{(\linewidth - 14\tabcolsep) * \real{0.1250}}@{}}
\toprule\noalign{}
\begin{minipage}[b]{\linewidth}\raggedright
RUNSPEC
\end{minipage} & \begin{minipage}[b]{\linewidth}\raggedright
GRID
\end{minipage} & \begin{minipage}[b]{\linewidth}\raggedright
EDIT
\end{minipage} & \begin{minipage}[b]{\linewidth}\raggedright
PROPS
\end{minipage} & \begin{minipage}[b]{\linewidth}\raggedright
REGIONS
\end{minipage} & \begin{minipage}[b]{\linewidth}\raggedright
SOLUTION
\end{minipage} & \begin{minipage}[b]{\linewidth}\raggedright
SUMMARY
\end{minipage} & \begin{minipage}[b]{\linewidth}\raggedright
SCHEDULE
\end{minipage} \\
\midrule\noalign{}
\endhead
\bottomrule\noalign{}
\endlastfoot
\end{longtable}

\subsubsection{Description}\label{description-57}

The FRICTION keyword activates the Wellbore Friction option and defines
the maximum number of wellbore friction wells together with the maximum
number of well branches.

This keyword is not supported by OPM Flow but would change the results
if supported so the simulation will be stopped.

\begin{longtable}[]{@{}
  >{\raggedright\arraybackslash}p{(\linewidth - 6\tabcolsep) * \real{0.2500}}
  >{\raggedright\arraybackslash}p{(\linewidth - 6\tabcolsep) * \real{0.2500}}
  >{\raggedright\arraybackslash}p{(\linewidth - 6\tabcolsep) * \real{0.2500}}
  >{\raggedright\arraybackslash}p{(\linewidth - 6\tabcolsep) * \real{0.2500}}@{}}
\toprule\noalign{}
\begin{minipage}[b]{\linewidth}\raggedright
No.
\end{minipage} & \begin{minipage}[b]{\linewidth}\raggedright
Name
\end{minipage} & \begin{minipage}[b]{\linewidth}\raggedright
Description
\end{minipage} & \begin{minipage}[b]{\linewidth}\raggedright
Default
\end{minipage} \\
\midrule\noalign{}
\endhead
\bottomrule\noalign{}
\endlastfoot
1 & MXWELS & A positive integer defining the maximum number of wellbore
friction wells for this model. & 0 \\
3 & MXBRAN & A positive integer defining the maximum number of branches
per well. The default value of one implies a standard well with no
branches. & 1 \\
Notes: & & & \\
\end{longtable}

\emph{Table 5.16: FRICTION Keyword Description}

\subsubsection{Example}\label{example-43}

\begin{verbatim}
--
--       WELL    BRANCH
--       MXWELS  MXBRAN
FRICTION
         5       1                                                             /
\end{verbatim}

The above example defines the maximum number of wellbore friction wells
to be five and the maximum number of branches set to one, for standard
wells.

\subsection{FULLIMP -- Activate Fully Implicit Solution
Option}\label{fullimp-activate-fully-implicit-solution-option}

\begin{longtable}[]{@{}
  >{\raggedright\arraybackslash}p{(\linewidth - 14\tabcolsep) * \real{0.1250}}
  >{\raggedright\arraybackslash}p{(\linewidth - 14\tabcolsep) * \real{0.1250}}
  >{\raggedright\arraybackslash}p{(\linewidth - 14\tabcolsep) * \real{0.1250}}
  >{\raggedright\arraybackslash}p{(\linewidth - 14\tabcolsep) * \real{0.1250}}
  >{\raggedright\arraybackslash}p{(\linewidth - 14\tabcolsep) * \real{0.1250}}
  >{\raggedright\arraybackslash}p{(\linewidth - 14\tabcolsep) * \real{0.1250}}
  >{\raggedright\arraybackslash}p{(\linewidth - 14\tabcolsep) * \real{0.1250}}
  >{\raggedright\arraybackslash}p{(\linewidth - 14\tabcolsep) * \real{0.1250}}@{}}
\toprule\noalign{}
\begin{minipage}[b]{\linewidth}\raggedright
RUNSPEC
\end{minipage} & \begin{minipage}[b]{\linewidth}\raggedright
GRID
\end{minipage} & \begin{minipage}[b]{\linewidth}\raggedright
EDIT
\end{minipage} & \begin{minipage}[b]{\linewidth}\raggedright
PROPS
\end{minipage} & \begin{minipage}[b]{\linewidth}\raggedright
REGIONS
\end{minipage} & \begin{minipage}[b]{\linewidth}\raggedright
SOLUTION
\end{minipage} & \begin{minipage}[b]{\linewidth}\raggedright
SUMMARY
\end{minipage} & \begin{minipage}[b]{\linewidth}\raggedright
SCHEDULE
\end{minipage} \\
\midrule\noalign{}
\endhead
\bottomrule\noalign{}
\endlastfoot
\end{longtable}

\subsubsection{Description}\label{description-58}

The FULLIMP keyword activates the Fully Implicit Solution formulation
and solution options. OPM Flow uses a different numerical scheme which
makes this keyword redundant; hence, OPM Flow ignores this keyword. It
is documented here for completeness. The keyword as the same function as
the IMPLICIT keyword in the RUNSPEC section.

There is no data required for this keyword and there is no terminating
``/'' for this keyword.

See section 2.2 Running OPM Flow 2023-04 From The Command Line on how to
invoke various numerical schemes via the OPM Flow command line
interface.

\subsubsection{Example}\label{example-44}

\begin{verbatim}
--
--       ACTIVATES THE FULLY IMPLICIT SOLUTION OPTION
--
FULLIMP
\end{verbatim}

The above example switches on the fully implicit solution option;
however, this has no effect in OPM Flow input decks.

\subsection{GAS -- Activate the Gas Phase in the
Model}\label{gas-activate-the-gas-phase-in-the-model}

\begin{longtable}[]{@{}
  >{\raggedright\arraybackslash}p{(\linewidth - 14\tabcolsep) * \real{0.1250}}
  >{\raggedright\arraybackslash}p{(\linewidth - 14\tabcolsep) * \real{0.1250}}
  >{\raggedright\arraybackslash}p{(\linewidth - 14\tabcolsep) * \real{0.1250}}
  >{\raggedright\arraybackslash}p{(\linewidth - 14\tabcolsep) * \real{0.1250}}
  >{\raggedright\arraybackslash}p{(\linewidth - 14\tabcolsep) * \real{0.1250}}
  >{\raggedright\arraybackslash}p{(\linewidth - 14\tabcolsep) * \real{0.1250}}
  >{\raggedright\arraybackslash}p{(\linewidth - 14\tabcolsep) * \real{0.1250}}
  >{\raggedright\arraybackslash}p{(\linewidth - 14\tabcolsep) * \real{0.1250}}@{}}
\toprule\noalign{}
\begin{minipage}[b]{\linewidth}\raggedright
RUNSPEC
\end{minipage} & \begin{minipage}[b]{\linewidth}\raggedright
GRID
\end{minipage} & \begin{minipage}[b]{\linewidth}\raggedright
EDIT
\end{minipage} & \begin{minipage}[b]{\linewidth}\raggedright
PROPS
\end{minipage} & \begin{minipage}[b]{\linewidth}\raggedright
REGIONS
\end{minipage} & \begin{minipage}[b]{\linewidth}\raggedright
SOLUTION
\end{minipage} & \begin{minipage}[b]{\linewidth}\raggedright
SUMMARY
\end{minipage} & \begin{minipage}[b]{\linewidth}\raggedright
SCHEDULE
\end{minipage} \\
\midrule\noalign{}
\endhead
\bottomrule\noalign{}
\endlastfoot
\end{longtable}

\subsubsection{Description}\label{description-59}

This keyword indicate that the gas phase is present in the model and
must be used for oil-gas, gas-water, oil-water-gas input decks that
contain the gas phase. The keyword will also invoke data input file
checking to ensure that all the required gas phase input parameters are
defined in the input deck.

There is no data required for this keyword and there is no terminating
``/'' for this keyword.

\subsubsection{Example}\label{example-45}

\begin{verbatim}
--
--       GAS PHASE IS PRESENT IN THE RUN
--
GAS
\end{verbatim}

The above example declares that the gas phase is active in the model.

\subsection{GASFIELD -- Define Gas Field Operations
Options}\label{gasfield-define-gas-field-operations-options}

\begin{longtable}[]{@{}
  >{\raggedright\arraybackslash}p{(\linewidth - 14\tabcolsep) * \real{0.1250}}
  >{\raggedright\arraybackslash}p{(\linewidth - 14\tabcolsep) * \real{0.1250}}
  >{\raggedright\arraybackslash}p{(\linewidth - 14\tabcolsep) * \real{0.1250}}
  >{\raggedright\arraybackslash}p{(\linewidth - 14\tabcolsep) * \real{0.1250}}
  >{\raggedright\arraybackslash}p{(\linewidth - 14\tabcolsep) * \real{0.1250}}
  >{\raggedright\arraybackslash}p{(\linewidth - 14\tabcolsep) * \real{0.1250}}
  >{\raggedright\arraybackslash}p{(\linewidth - 14\tabcolsep) * \real{0.1250}}
  >{\raggedright\arraybackslash}p{(\linewidth - 14\tabcolsep) * \real{0.1250}}@{}}
\toprule\noalign{}
\begin{minipage}[b]{\linewidth}\raggedright
RUNSPEC
\end{minipage} & \begin{minipage}[b]{\linewidth}\raggedright
GRID
\end{minipage} & \begin{minipage}[b]{\linewidth}\raggedright
EDIT
\end{minipage} & \begin{minipage}[b]{\linewidth}\raggedright
PROPS
\end{minipage} & \begin{minipage}[b]{\linewidth}\raggedright
REGIONS
\end{minipage} & \begin{minipage}[b]{\linewidth}\raggedright
SOLUTION
\end{minipage} & \begin{minipage}[b]{\linewidth}\raggedright
SUMMARY
\end{minipage} & \begin{minipage}[b]{\linewidth}\raggedright
SCHEDULE
\end{minipage} \\
\midrule\noalign{}
\endhead
\bottomrule\noalign{}
\endlastfoot
\end{longtable}

\subsubsection{Description}\label{description-60}

The GASFIELD keyword activates and specifies the Gas Field Operations
options and determines if extended compressors are present in the run
and if the expedited first pass DCQ calculation should be used.

This keyword is not supported by OPM Flow but would change the results
if supported so the simulation will be stopped.

\subsection{GASWAT -- Activate the Gas-Water Model
Formulation}\label{gaswat-activate-the-gas-water-model-formulation}

\begin{longtable}[]{@{}
  >{\raggedright\arraybackslash}p{(\linewidth - 14\tabcolsep) * \real{0.1250}}
  >{\raggedright\arraybackslash}p{(\linewidth - 14\tabcolsep) * \real{0.1250}}
  >{\raggedright\arraybackslash}p{(\linewidth - 14\tabcolsep) * \real{0.1250}}
  >{\raggedright\arraybackslash}p{(\linewidth - 14\tabcolsep) * \real{0.1250}}
  >{\raggedright\arraybackslash}p{(\linewidth - 14\tabcolsep) * \real{0.1250}}
  >{\raggedright\arraybackslash}p{(\linewidth - 14\tabcolsep) * \real{0.1250}}
  >{\raggedright\arraybackslash}p{(\linewidth - 14\tabcolsep) * \real{0.1250}}
  >{\raggedright\arraybackslash}p{(\linewidth - 14\tabcolsep) * \real{0.1250}}@{}}
\toprule\noalign{}
\begin{minipage}[b]{\linewidth}\raggedright
RUNSPEC
\end{minipage} & \begin{minipage}[b]{\linewidth}\raggedright
GRID
\end{minipage} & \begin{minipage}[b]{\linewidth}\raggedright
EDIT
\end{minipage} & \begin{minipage}[b]{\linewidth}\raggedright
PROPS
\end{minipage} & \begin{minipage}[b]{\linewidth}\raggedright
REGIONS
\end{minipage} & \begin{minipage}[b]{\linewidth}\raggedright
SOLUTION
\end{minipage} & \begin{minipage}[b]{\linewidth}\raggedright
SUMMARY
\end{minipage} & \begin{minipage}[b]{\linewidth}\raggedright
SCHEDULE
\end{minipage} \\
\midrule\noalign{}
\endhead
\bottomrule\noalign{}
\endlastfoot
\end{longtable}

\subsubsection{Description}\label{description-61}

This keyword activates the two-phase Gas-Water model, as such it is
equivalent to using both the GAS and WATER keywords in the RUNSPEC
section..

\begin{tcolorbox}[enhanced jigsaw, title=\textcolor{quarto-callout-note-color}{\faInfo}\hspace{0.5em}{Note}, coltitle=black, breakable, toprule=.15mm, colbacktitle=quarto-callout-note-color!10!white, leftrule=.75mm, opacitybacktitle=0.6, colback=white, bottomtitle=1mm, toptitle=1mm, titlerule=0mm, left=2mm, colframe=quarto-callout-note-color-frame, arc=.35mm, rightrule=.15mm, bottomrule=.15mm, opacityback=0]

This is an OPM Flow keyword, and should not be confused with the more
general version of the GASWAT keyword used in the commercial
compositional simulator.

\end{tcolorbox}

There is no data required for this keyword and there is no terminating
``/'' for this keyword.

\subsubsection{Example}\label{example-46}

\begin{verbatim}
-- ==============================================================================
--
-- RUNSPEC SECTION
--
-- ==============================================================================
RUNSPEC
--
--       ACTIVATE CO2 STORAGE IN THE MODEL (OPM FLOW CO2 STORAGE KEYWORD)
--
CO2STORE
--
--       ACTIVATE COMPOSITIONAL MODELING FORMULATION (OPM FLOW KEYWORD)
--
COMPS
         2                                                                     /
--
--       ACTIVATE GAS-WATER THE MODEL (OPM FLOW KEYWORD)
--
GASWAT
\end{verbatim}

The above example declares that the run should use the Gas-Water model,
together with the CO2STORE option, and with two components.

\subsection{GDIMS -- Activate Instantaneous Gradient Option and Define
Dimensions}\label{gdims-activate-instantaneous-gradient-option-and-define-dimensions}

\begin{longtable}[]{@{}
  >{\raggedright\arraybackslash}p{(\linewidth - 14\tabcolsep) * \real{0.1250}}
  >{\raggedright\arraybackslash}p{(\linewidth - 14\tabcolsep) * \real{0.1250}}
  >{\raggedright\arraybackslash}p{(\linewidth - 14\tabcolsep) * \real{0.1250}}
  >{\raggedright\arraybackslash}p{(\linewidth - 14\tabcolsep) * \real{0.1250}}
  >{\raggedright\arraybackslash}p{(\linewidth - 14\tabcolsep) * \real{0.1250}}
  >{\raggedright\arraybackslash}p{(\linewidth - 14\tabcolsep) * \real{0.1250}}
  >{\raggedright\arraybackslash}p{(\linewidth - 14\tabcolsep) * \real{0.1250}}
  >{\raggedright\arraybackslash}p{(\linewidth - 14\tabcolsep) * \real{0.1250}}@{}}
\toprule\noalign{}
\begin{minipage}[b]{\linewidth}\raggedright
RUNSPEC
\end{minipage} & \begin{minipage}[b]{\linewidth}\raggedright
GRID
\end{minipage} & \begin{minipage}[b]{\linewidth}\raggedright
EDIT
\end{minipage} & \begin{minipage}[b]{\linewidth}\raggedright
PROPS
\end{minipage} & \begin{minipage}[b]{\linewidth}\raggedright
REGIONS
\end{minipage} & \begin{minipage}[b]{\linewidth}\raggedright
SOLUTION
\end{minipage} & \begin{minipage}[b]{\linewidth}\raggedright
SUMMARY
\end{minipage} & \begin{minipage}[b]{\linewidth}\raggedright
SCHEDULE
\end{minipage} \\
\midrule\noalign{}
\endhead
\bottomrule\noalign{}
\endlastfoot
\end{longtable}

\subsubsection{Description}\label{description-62}

The GDIMS keyword activates the Instantaneous Gradient option and
defines the maximum dimensions as used by the GWRTWCV keyword in the
SCHEDULE section. The Instantaneous Gradient option calculates
derivatives of solution quantities at the current time step with respect
to variations in the variables at the current time step. This is
different to Gradient option that calculates the derivatives of solution
quantities at the current time step with respect to variations in the
variables at the initial time step, that is a time equal to zero.
Consequently, the Instantaneous Gradient option can be switched on and
off by the GUPFREQ keyword in the SCHEDULE section, whereas the Gradient
option cannot.

This keyword is not supported by OPM Flow but would change the results
if supported so the simulation will be stopped.

\subsection{GIMODEL -- Activate Gi Pseudo Compositional
Option}\label{gimodel-activate-gi-pseudo-compositional-option}

\begin{longtable}[]{@{}
  >{\raggedright\arraybackslash}p{(\linewidth - 14\tabcolsep) * \real{0.1250}}
  >{\raggedright\arraybackslash}p{(\linewidth - 14\tabcolsep) * \real{0.1250}}
  >{\raggedright\arraybackslash}p{(\linewidth - 14\tabcolsep) * \real{0.1250}}
  >{\raggedright\arraybackslash}p{(\linewidth - 14\tabcolsep) * \real{0.1250}}
  >{\raggedright\arraybackslash}p{(\linewidth - 14\tabcolsep) * \real{0.1250}}
  >{\raggedright\arraybackslash}p{(\linewidth - 14\tabcolsep) * \real{0.1250}}
  >{\raggedright\arraybackslash}p{(\linewidth - 14\tabcolsep) * \real{0.1250}}
  >{\raggedright\arraybackslash}p{(\linewidth - 14\tabcolsep) * \real{0.1250}}@{}}
\toprule\noalign{}
\begin{minipage}[b]{\linewidth}\raggedright
RUNSPEC
\end{minipage} & \begin{minipage}[b]{\linewidth}\raggedright
GRID
\end{minipage} & \begin{minipage}[b]{\linewidth}\raggedright
EDIT
\end{minipage} & \begin{minipage}[b]{\linewidth}\raggedright
PROPS
\end{minipage} & \begin{minipage}[b]{\linewidth}\raggedright
REGIONS
\end{minipage} & \begin{minipage}[b]{\linewidth}\raggedright
SOLUTION
\end{minipage} & \begin{minipage}[b]{\linewidth}\raggedright
SUMMARY
\end{minipage} & \begin{minipage}[b]{\linewidth}\raggedright
SCHEDULE
\end{minipage} \\
\midrule\noalign{}
\endhead
\bottomrule\noalign{}
\endlastfoot
\end{longtable}

\subsubsection{Description}\label{description-63}

This keyword, GIMODEL, activates the Gi Pseudo Compositional option for
gas condensate and volatile oil fluids.

This keyword is not supported by OPM Flow but would change the results
if supported so the simulation will be stopped.

The accuracy of gas condensate and volatile oil modeling using a
black-oil reservoir simulator is now firmly accepted in the industry;
both depletion and gas cycling above the saturation point can be modeled
and yield an acceptable match compared with compositional simulators.
The main deficiency with the existing black-oil formulation is the
treatment of gas injection below the saturation point, where the
compositional effect of the stripping of liquid components is inversely
proportional to their molecular weights, is not considered. This is
because the black-oil model assumes that the saturated hydrocarbon fluid
properties are only functions of pressure. Thus, when dry gas is
injected into a condensate reservoir below the saturation pressure the
injected gas continues to re-vaporize liquid at a rate governed only by
the cell pressures. Compositional modeling indicate that this not the
case.

The GI option attempts to overcome the limitation of the standard
black-oil approach by extending the black-oil model using the method of
Cook et al.~{[}Cook, R. E., Jacoby, R. H., and Ramesh, A. B. ``A
Beta-type Reservoir Simulator for Approximating Compositional Effects
During Gas Injection'' paper SPE 4272, Society of Petroleum Engineers
Journal (1974) 14, No.~5, 471-481.{]} to take into account the fluid
property changes occurring during gas injection. This is done by
extending the fluid property treatment so that the saturated fluid
properties depend on pressure (as per the standard black-oil
formulation) and also on an additional parameter which characterizes the
compositional changes in the reservoir liquid and vapor phases at
constant pressure.

There is no data required for this keyword.

\subsubsection{Example}\label{example-47}

\begin{verbatim}
--
--       ACTIVATE THE GI PSEUDO COMPOSITIONAL OPTION
--
GIMODEL
\end{verbatim}

The above example switches on the Gi Pseudo Compositional option.

\subsection{GRAVDR -- Activate Gravity Drainage and Imbibition for Dual
Porosity
Model}\label{gravdr-activate-gravity-drainage-and-imbibition-for-dual-porosity-model}

\begin{longtable}[]{@{}
  >{\raggedright\arraybackslash}p{(\linewidth - 14\tabcolsep) * \real{0.1250}}
  >{\raggedright\arraybackslash}p{(\linewidth - 14\tabcolsep) * \real{0.1250}}
  >{\raggedright\arraybackslash}p{(\linewidth - 14\tabcolsep) * \real{0.1250}}
  >{\raggedright\arraybackslash}p{(\linewidth - 14\tabcolsep) * \real{0.1250}}
  >{\raggedright\arraybackslash}p{(\linewidth - 14\tabcolsep) * \real{0.1250}}
  >{\raggedright\arraybackslash}p{(\linewidth - 14\tabcolsep) * \real{0.1250}}
  >{\raggedright\arraybackslash}p{(\linewidth - 14\tabcolsep) * \real{0.1250}}
  >{\raggedright\arraybackslash}p{(\linewidth - 14\tabcolsep) * \real{0.1250}}@{}}
\toprule\noalign{}
\begin{minipage}[b]{\linewidth}\raggedright
RUNSPEC
\end{minipage} & \begin{minipage}[b]{\linewidth}\raggedright
GRID
\end{minipage} & \begin{minipage}[b]{\linewidth}\raggedright
EDIT
\end{minipage} & \begin{minipage}[b]{\linewidth}\raggedright
PROPS
\end{minipage} & \begin{minipage}[b]{\linewidth}\raggedright
REGIONS
\end{minipage} & \begin{minipage}[b]{\linewidth}\raggedright
SOLUTION
\end{minipage} & \begin{minipage}[b]{\linewidth}\raggedright
SUMMARY
\end{minipage} & \begin{minipage}[b]{\linewidth}\raggedright
SCHEDULE
\end{minipage} \\
\midrule\noalign{}
\endhead
\bottomrule\noalign{}
\endlastfoot
\end{longtable}

\subsubsection{Description}\label{description-64}

This keyword switches on gravity drainage and imbibition modeling
between the matrix and the fracture grid blocks in dual porosity and
dual permeability runs. Note that either DZMTRX or DZMTRXV keywords in
the GRID section should be used to set the matrix vertical dimensions if
this option is activated.

There is no data required for this keyword and there is no terminating
``/'' for this keyword.

This keyword is not supported by OPM Flow but would change the results
if supported so the simulation will be stopped.

\subsubsection{Example}\label{example-48}

\begin{verbatim}
--
--       ACTIVATE GRAVITY DRAINAGE AND IMBIBITION FOR DUAL POROSITY MODEL
--
GRAVDR
\end{verbatim}

The above example switches on the gravity drainage and imbibition option
for the run.

\subsection{GRAVDRB -- Activate Vertical Discretized Gravity Drainage
and Imbibition for Dual Porosity
Model}\label{gravdrb-activate-vertical-discretized-gravity-drainage-and-imbibition-for-dual-porosity-model}

\begin{longtable}[]{@{}
  >{\raggedright\arraybackslash}p{(\linewidth - 14\tabcolsep) * \real{0.1250}}
  >{\raggedright\arraybackslash}p{(\linewidth - 14\tabcolsep) * \real{0.1250}}
  >{\raggedright\arraybackslash}p{(\linewidth - 14\tabcolsep) * \real{0.1250}}
  >{\raggedright\arraybackslash}p{(\linewidth - 14\tabcolsep) * \real{0.1250}}
  >{\raggedright\arraybackslash}p{(\linewidth - 14\tabcolsep) * \real{0.1250}}
  >{\raggedright\arraybackslash}p{(\linewidth - 14\tabcolsep) * \real{0.1250}}
  >{\raggedright\arraybackslash}p{(\linewidth - 14\tabcolsep) * \real{0.1250}}
  >{\raggedright\arraybackslash}p{(\linewidth - 14\tabcolsep) * \real{0.1250}}@{}}
\toprule\noalign{}
\begin{minipage}[b]{\linewidth}\raggedright
RUNSPEC
\end{minipage} & \begin{minipage}[b]{\linewidth}\raggedright
GRID
\end{minipage} & \begin{minipage}[b]{\linewidth}\raggedright
EDIT
\end{minipage} & \begin{minipage}[b]{\linewidth}\raggedright
PROPS
\end{minipage} & \begin{minipage}[b]{\linewidth}\raggedright
REGIONS
\end{minipage} & \begin{minipage}[b]{\linewidth}\raggedright
SOLUTION
\end{minipage} & \begin{minipage}[b]{\linewidth}\raggedright
SUMMARY
\end{minipage} & \begin{minipage}[b]{\linewidth}\raggedright
SCHEDULE
\end{minipage} \\
\midrule\noalign{}
\endhead
\bottomrule\noalign{}
\endlastfoot
\end{longtable}

\subsubsection{Description}\label{description-65}

This keyword switches on vertical discretized gravity drainage and
imbibition modeling between the matrix and the fracture grid blocks in
dual porosity and dual permeability runs. Note that the geometry of the
matrix sub-cells should be set to VERTICAL on the NMATOPS keyword in the
GRID section if this option is activated.

There is no data required for this keyword and there is no terminating
``/'' for this keyword.

This keyword is not supported by OPM Flow but would change the results
if supported so the simulation will be stopped.

\subsubsection{Example}\label{example-49}

\begin{verbatim}
--
--       ACTIVATE VERTICAL DISCRETIZED GRAVITY DRAINAGE AND IMBIBITION
--
GRAVDRB
\end{verbatim}

The above example switches on the vertical discretized gravity drainage
and imbibition option for the run.

\subsection{GRAVDRM -- Activate Alternative Gravity Drainage and
Imbibition for Dual Porosity
Model}\label{gravdrm-activate-alternative-gravity-drainage-and-imbibition-for-dual-porosity-model}

\begin{longtable}[]{@{}
  >{\raggedright\arraybackslash}p{(\linewidth - 14\tabcolsep) * \real{0.1250}}
  >{\raggedright\arraybackslash}p{(\linewidth - 14\tabcolsep) * \real{0.1250}}
  >{\raggedright\arraybackslash}p{(\linewidth - 14\tabcolsep) * \real{0.1250}}
  >{\raggedright\arraybackslash}p{(\linewidth - 14\tabcolsep) * \real{0.1250}}
  >{\raggedright\arraybackslash}p{(\linewidth - 14\tabcolsep) * \real{0.1250}}
  >{\raggedright\arraybackslash}p{(\linewidth - 14\tabcolsep) * \real{0.1250}}
  >{\raggedright\arraybackslash}p{(\linewidth - 14\tabcolsep) * \real{0.1250}}
  >{\raggedright\arraybackslash}p{(\linewidth - 14\tabcolsep) * \real{0.1250}}@{}}
\toprule\noalign{}
\begin{minipage}[b]{\linewidth}\raggedright
RUNSPEC
\end{minipage} & \begin{minipage}[b]{\linewidth}\raggedright
GRID
\end{minipage} & \begin{minipage}[b]{\linewidth}\raggedright
EDIT
\end{minipage} & \begin{minipage}[b]{\linewidth}\raggedright
PROPS
\end{minipage} & \begin{minipage}[b]{\linewidth}\raggedright
REGIONS
\end{minipage} & \begin{minipage}[b]{\linewidth}\raggedright
SOLUTION
\end{minipage} & \begin{minipage}[b]{\linewidth}\raggedright
SUMMARY
\end{minipage} & \begin{minipage}[b]{\linewidth}\raggedright
SCHEDULE
\end{minipage} \\
\midrule\noalign{}
\endhead
\bottomrule\noalign{}
\endlastfoot
\end{longtable}

\subsubsection{Description}\label{description-66}

This keyword switches on the alternative gravity drainage and imbibition
modeling between the matrix and the fracture grid blocks in dual
porosity and dual permeability runs. Either the GRAVDRM or GRAVDR
keywords should be used to activate this standard or alternative type of
formulation.

There is no data required for this keyword.

This keyword is not supported by OPM Flow but would change the results
if supported so the simulation will be stopped.

\begin{longtable}[]{@{}
  >{\raggedright\arraybackslash}p{(\linewidth - 6\tabcolsep) * \real{0.2500}}
  >{\raggedright\arraybackslash}p{(\linewidth - 6\tabcolsep) * \real{0.2500}}
  >{\raggedright\arraybackslash}p{(\linewidth - 6\tabcolsep) * \real{0.2500}}
  >{\raggedright\arraybackslash}p{(\linewidth - 6\tabcolsep) * \real{0.2500}}@{}}
\toprule\noalign{}
\begin{minipage}[b]{\linewidth}\raggedright
No.
\end{minipage} & \begin{minipage}[b]{\linewidth}\raggedright
Name
\end{minipage} & \begin{minipage}[b]{\linewidth}\raggedright
Description
\end{minipage} & \begin{minipage}[b]{\linewidth}\raggedright
Default
\end{minipage} \\
\midrule\noalign{}
\endhead
\bottomrule\noalign{}
\endlastfoot
1 & OPTION1 & A defined character string that sets the matrix flow in
and out of the matrix block option, and should be set to one of the
following: & YES \\
Notes: & & & \\
\end{longtable}

\emph{Table 5.17: GRAVDRM Keyword Description}

\subsubsection{Example}\label{example-50}

\begin{verbatim}
--
--       ACTIVATE ALTERNATIVE GRAVITY DRAINAGE AND IMBIBITION MODEL
--
--       MATRIX
--       OPTION
GRAVDRM
         YES                                                                   /
\end{verbatim}

The above example switches on the alternative gravity drainage and
imbibition option for the run and sets oil flow to be bi-directional,
that is oil can flow into and out of the matrix block.

\subsection{GRIDOPTS -- Grid Processing
Options}\label{gridopts-grid-processing-options}

\begin{longtable}[]{@{}
  >{\raggedright\arraybackslash}p{(\linewidth - 14\tabcolsep) * \real{0.1250}}
  >{\raggedright\arraybackslash}p{(\linewidth - 14\tabcolsep) * \real{0.1250}}
  >{\raggedright\arraybackslash}p{(\linewidth - 14\tabcolsep) * \real{0.1250}}
  >{\raggedright\arraybackslash}p{(\linewidth - 14\tabcolsep) * \real{0.1250}}
  >{\raggedright\arraybackslash}p{(\linewidth - 14\tabcolsep) * \real{0.1250}}
  >{\raggedright\arraybackslash}p{(\linewidth - 14\tabcolsep) * \real{0.1250}}
  >{\raggedright\arraybackslash}p{(\linewidth - 14\tabcolsep) * \real{0.1250}}
  >{\raggedright\arraybackslash}p{(\linewidth - 14\tabcolsep) * \real{0.1250}}@{}}
\toprule\noalign{}
\begin{minipage}[b]{\linewidth}\raggedright
RUNSPEC
\end{minipage} & \begin{minipage}[b]{\linewidth}\raggedright
GRID
\end{minipage} & \begin{minipage}[b]{\linewidth}\raggedright
EDIT
\end{minipage} & \begin{minipage}[b]{\linewidth}\raggedright
PROPS
\end{minipage} & \begin{minipage}[b]{\linewidth}\raggedright
REGIONS
\end{minipage} & \begin{minipage}[b]{\linewidth}\raggedright
SOLUTION
\end{minipage} & \begin{minipage}[b]{\linewidth}\raggedright
SUMMARY
\end{minipage} & \begin{minipage}[b]{\linewidth}\raggedright
SCHEDULE
\end{minipage} \\
\midrule\noalign{}
\endhead
\bottomrule\noalign{}
\endlastfoot
\end{longtable}

\subsubsection{Description}\label{description-67}

GRIDOPTS activates the negative directional dependent transmissibility
multipliers option, defines the maximum number of MULTNUM regions and
the number of PINCHNUM regions for the model.

\begin{longtable}[]{@{}
  >{\raggedright\arraybackslash}p{(\linewidth - 6\tabcolsep) * \real{0.2500}}
  >{\raggedright\arraybackslash}p{(\linewidth - 6\tabcolsep) * \real{0.2500}}
  >{\raggedright\arraybackslash}p{(\linewidth - 6\tabcolsep) * \real{0.2500}}
  >{\raggedright\arraybackslash}p{(\linewidth - 6\tabcolsep) * \real{0.2500}}@{}}
\toprule\noalign{}
\begin{minipage}[b]{\linewidth}\raggedright
No.
\end{minipage} & \begin{minipage}[b]{\linewidth}\raggedright
Name
\end{minipage} & \begin{minipage}[b]{\linewidth}\raggedright
Description
\end{minipage} & \begin{minipage}[b]{\linewidth}\raggedright
Default
\end{minipage} \\
\midrule\noalign{}
\endhead
\bottomrule\noalign{}
\endlastfoot
1 & TRANMULT & A character string that activates the negative
directional dependent transmissibility multipliers option by setting
TRANMULT to YES. Setting the value to NO switches off this option. OPM
Flow uses a positive directional dependent transmissibility formulation
to describe the flow between two cells, that is for cell (I, J, K) OPM
Flow calculates the x face transmissibility between (I, J, K) and (I +1,
J, K) cell face. Modification to the transmissibilities in this case is
accomplished by the MULTX, MULTY and MULTZ keywords. Setting TRANMULT to
YES invokes the option to use a negative directional dependent
multiplier scheme using the MULTX-, MULTY- and MULTZ- keywords. In this
case OPM Flow applies the x face transmissibility between (I - 1, J, K)
and (I, J, K) cell face when using the MULTX-, MULTY- and MULTZ-
keywords. Note that if TRANMULT is defaulted, and there are negative
directional dependent multiplier keywords in the input deck, then OPM
Flow will continue to process the MULTX-, MULTY- and MULTZ keywords
correctly. Whereas, the commercial simulator will terminate with an
error. & NO \\
2 & NRMULT & A positive integer value that defines the maximum number of
MULTNUM regions for the MULTNUM array. The MULTNUM array is used in the
GRID section to define various inter-region transmissibility regions in
the model and NRMULT sets the maximum number of regions which is the
maximum value of an element in the MULTNUM array. Inter-region MULTNUM
transmissibility multipliers can be defined using the MULTREGT and
regional pore volumes multipliers can be set using the MULTREGP keyword.
& 0 \\
3 & NRPINC & A positive integer value that defines the maximum number of
PINCHNUM regions for the PINCHNUM array. The PINCHNUM array is used in
the GRID section to define various regions in the model and NRPINC sets
the maximum of regions which is the maximum value of an element in the
PINCHNUM array. Each regions criteria for setting the pinch out criteria
is set by the PINCHREG keyword. & 0 \\
Notes: & & & \\
\end{longtable}

\emph{Table 5.18: GRIDOPTS Keyword Description}

See also the MULTNUM, MULTREGP, MULTREGT, PINCHNUM, and PINCHREG
keywords.

\subsubsection{Example}\label{example-51}

\begin{verbatim}
–-
--      NEG      MAX     MAX
--      MULTS    MULTNUM PINCHNUM
GRIDOPTS
        NO       9       1*                                                    /
\end{verbatim}

The above example switches off the negative directional dependent
transmissibility multipliers option and defines the maximum of MULTNUM
regions to be nine,. The NRPINC parameter is defaulted which means there
the maximum number of PINCHREG regions is zero.

\subsection{H2SOL -- Activate Dissolved H2 in the Water
Phase}\label{h2sol-activate-dissolved-h2-in-the-water-phase}

\begin{longtable}[]{@{}
  >{\raggedright\arraybackslash}p{(\linewidth - 14\tabcolsep) * \real{0.1250}}
  >{\raggedright\arraybackslash}p{(\linewidth - 14\tabcolsep) * \real{0.1250}}
  >{\raggedright\arraybackslash}p{(\linewidth - 14\tabcolsep) * \real{0.1250}}
  >{\raggedright\arraybackslash}p{(\linewidth - 14\tabcolsep) * \real{0.1250}}
  >{\raggedright\arraybackslash}p{(\linewidth - 14\tabcolsep) * \real{0.1250}}
  >{\raggedright\arraybackslash}p{(\linewidth - 14\tabcolsep) * \real{0.1250}}
  >{\raggedright\arraybackslash}p{(\linewidth - 14\tabcolsep) * \real{0.1250}}
  >{\raggedright\arraybackslash}p{(\linewidth - 14\tabcolsep) * \real{0.1250}}@{}}
\toprule\noalign{}
\begin{minipage}[b]{\linewidth}\raggedright
RUNSPEC
\end{minipage} & \begin{minipage}[b]{\linewidth}\raggedright
GRID
\end{minipage} & \begin{minipage}[b]{\linewidth}\raggedright
EDIT
\end{minipage} & \begin{minipage}[b]{\linewidth}\raggedright
PROPS
\end{minipage} & \begin{minipage}[b]{\linewidth}\raggedright
REGIONS
\end{minipage} & \begin{minipage}[b]{\linewidth}\raggedright
SOLUTION
\end{minipage} & \begin{minipage}[b]{\linewidth}\raggedright
SUMMARY
\end{minipage} & \begin{minipage}[b]{\linewidth}\raggedright
SCHEDULE
\end{minipage} \\
\midrule\noalign{}
\endhead
\bottomrule\noalign{}
\endlastfoot
\end{longtable}

\subsubsection{Description}\label{description-68}

The \hyperref[REF_HEADING_KEYWORD_H2SOL]{H2SOL} keyword activates
dissolved hydrogen (H2) in the water phase, where H2 is represented by
the SOLVENT pseudo component, using the simulator's H2-Brine PVT model.
This keyword is a compositional keyword in the commercial simulator but
has been implemented in OPM Flow's black-oil model.

The \hyperref[REF_HEADING_KEYWORD_H2SOL]{H2SOL} keyword can be used when
modelling H2 injection in depleted hydrocarbon reservoirs. See also the
\hyperref[REF_HEADING_KEYWORD_H2STORE]{H2STORE} keyword in the RUNSPEC
section which can be used when modelling H2 injection in saline
aquifers.

The H2-Brine PVT model computes the PVT properties such as density,
viscosity, and enthalpy internally as functions of pressure,
temperature, and composition by using analytic correlations and models
from the literature rather than by interpolation of tabulated values in
the input deck. These values are transformed to the standard black-oil
equivalent PVT tables internally by the simulator. Dissolved hydrogen in
brine is modeled using Li et al.~{[}Dedong Li, Christof Beyer, Sebastian
Bauer, A unified phase equilibrium model for hydrogen solubility and
solution density, International Journal of Hydrogen Energy, Volume 43,
Issue 1, 2018, Pages 512-529.{]}. Hydrogen gas density was modeled using
Leachman et al.~{[}J. W. Leachman, R. T Jacobsen, S. G. Penoncello, E.
W. Lemmon; Fundamental Equations of State for Parahydrogen, Normal
Hydrogen, and Orthohydrogen. J. Phys. Chem. Ref. Data 1 September 2009;
38 (3): 721--748.{]}. Other properties are modeled the same way as in
CO2STORE but modified for hydrogen-brine (a similar approach was
recently reported by Raad et al.~{[}Seyed Mostafa Jafari Raad, Ehsan
Ranjbar, Hassan Hassanzadeh, Yuri Leonenko, Hydrogen-brine mixture PVT
data for reservoir simulation of hydrogen storage in deep saline
aquifers, International Journal of Hydrogen Energy, Volume 48, Issue 2,
2023, Pages 696-708.{]}). A full description of the underlying PVT
models used by CO2STORE is described by Sandve et al.~{[}Tor Harald
Sandve, Sarah E. Gasda, Atgeirr Rasmussen, and Alf Birger Rustad.
Convective dissolution in field scale CO2 storage simulation using the
OPM Flow simulator. Submitted to TCCS 11 -- Trondheim Conference on CO2
Capture, Transport and Storage Trondheim, Norway -- June 21-23,
2021.{]}. Note that the H2-Brine PVT properties depend on the
temperature and salinity and these must therefore be entered in the
PROPS section. The reservoir temperature can be defined using, e.g., the
RTEMP keyword. Region based salinity can be provided using the SALINITY
keyword.

The \hyperref[REF_HEADING_KEYWORD_H2SOL]{H2SOL} keyword must be used
together with the SOLVENT keyword in the RUNSPEC section. The SOLVENT
keyword activates the four component black-oil model: oil, water and
gas, plus a solvent (in this case H2).

The DISGASW keyword in the RUNSPEC section can be used to model
dissolution of H2 in the Brine. Note that hydrocarbon gas is not allowed
to dissolve in the Brine (this is considered a reasonable assumption for
most hydrocarbon gases).

The hydrocarbon reservoir fluid properties should be set up as normal
apart from the water PVT properties, which will be defined by the
simulators' H2-Brine PVT model, so the water PVT keywords (e.g.~PVTW)
are not required. The H2 PVT properties will also be defined by the
H2-Brine PVT model so the solvent PVT keywords (e.g.~SDENSITY, PVTSOL)
are also not required.

The oil-water and gas-water saturation functions should be set up as
normal. The H2 and gas miscible relative permeability tables should be
defined using the SSFN keyword in the PROPS section.

The initial H2 saturations in the model can be specified by enumeration
using the SSOL keyword in the SOLUTION section.

The H2 fraction in a well's injection stream can be defined using the
WSOLVENT keyword in the SCHEDULE section.

There is no data required for this keyword and there is no terminating
``/'' for this keyword.

\subsubsection{Example}\label{example-52}

The example shows the usage of
\hyperref[REF_HEADING_KEYWORD_H2SOL]{H2SOL} to model H2 injection in a
live-oil reservoir. The RUNSPEC section includes the
\hyperref[REF_HEADING_KEYWORD_H2SOL]{H2SOL} keyword, which activates
dissolved hydrogen (H2) in the water phase. The SOLVENT keyword
activates the solvent pseudo component, which represents H2. The DISGASW
keyword has been used to activate dissolved H2 in the brine.

\begin{verbatim}
-- ==============================================================================
--
-- RUNSPEC SECTION
--
-- ==============================================================================
RUNSPEC
-- ------------------------------------------------------------------------------
-- FLUID TYPES AND TRACER OPTIONS
-- ------------------------------------------------------------------------------
--
--       ACTIVATE DISSOLVED H2 IN WATER (OPM FLOW KEYWORD)
--
H2SOL
--
--       SOLVENT PHASE IS PRESENT IN THE RUN
--
SOLVENT
--
--       DISSOLVED H2 IN WATER IS PRESENT IN THE RUN (OPM FLOW KEYWORD)
--
DISGASW
\end{verbatim}

The second part of the example covers the additional data required in
the PROPS section, in which the H2-gas miscible relative permeabilities
are defined using the SSFN keyword. The initial temperature and salinity
are defined using the RTEMP and SALINITY keywords. The oil and gas PVT
properties should be defined in the normal way. No water PVT data is
required.

\begin{verbatim}
-- ==============================================================================
--
-- PROPS SECTION
--
-- ==============================================================================
PROPS
--
--       RESERVOIR
--       TEMPERATURE
--       -----------
RTEMP
         80.0                                              / RESERVOIR TEMP
--
--       SET SALINITY FOR ALL CELLS (OPM FLOW KEYWORD)
--
SALINITY
         0.7                                                       /
--
-- SOLVENT RELATIVE PERMEABILITY TABLES
--
SSFN
-- SGAS KRGT KRST
-- FRAC
-- -------- -------- --------
 0.0000 0.0000 0.0000
 1.0000 1.0000 1.0000 /
\end{verbatim}

The third part of the example covers initialising the model in the
SOLUTION section. Here the EQUIL keyword is used as normal to
equilibrate the oil, gas and water phases, and the SSOL keyword to
specify that there is initially no H2 present in the gas phase.

\begin{verbatim}
-- ==============================================================================
--
-- SOLUTION SECTION
--
-- ==============================================================================
SOLUTION
--
--       DATUM   DATUM   OWC     PCOW   GOC    PCGO   RS   RV   N    E300  RVW
--       DEPTH   PRESS   DEPTH   ----   DEPTH  ----   OPT  OPT  OPT  OPT   OPT
EQUIL
         2000.0  200.0   2100.0  0.00  2000.0  0.00   1*   1*   1*   2*    1*  /
--
-- DEFINE INITIAL EQUILIBRATION SOLVENT SAT VALUES FOR ALL CELLS IN THE MODEL
-- BASED ON NX = 100, NY = 100 AND NZ = 3
--
SSOL
         30000*0.0 /
\end{verbatim}

For the summary section, the H2 storage quantity can be tracked using
the solvent injection and production volumes as shown below.

\begin{verbatim}
-- ==============================================================================
--
-- SUMMARY SECTION
--
-- ==============================================================================
SUMMARY
--
--       EXPORT STANDARD SUMMARY VARIABLE VECTORS TO FILE
--
ALL
--
--       FIELD SOLVENT INJECTION RATE
FNIR
--
--       FIELD SOLVENT PRODUCTION RATE
FNPR
--
--       FIELD SOLVENT INJECTION TOTAL
FNIT
--
--       FIELD SOLVENT PRODUCTION TOTAL
FNPT
\end{verbatim}

The final part of the example covers the SCHEDULE section. The WSOLVENT
keyword is used to specify the injected solvent fraction (in this case
pure H2).

\begin{verbatim}
-- ==============================================================================
--
-- SCHEDULE SECTION
--
-- ==============================================================================
SCHEDULE
--
--       RESTART CONTROL BASIC = 4 (ALL=2, YEARLY=4, MONTHLY=5, TSTEP=6)
--
RPTRST
        'BASIC=1'  'ALLPROPS'   /
--
-- DEFINE GAS INJECTION WELL SOLVENT FRACTION
--
-- WELL SOLVENT
-- NAME FRACTION
-- --------
WSOLVENT
   GI01 1.0000 /
/
\end{verbatim}

\subsection{H2STORE -- Activate the H2 Storage
Model}\label{h2store-activate-the-h2-storage-model}

\begin{longtable}[]{@{}
  >{\raggedright\arraybackslash}p{(\linewidth - 14\tabcolsep) * \real{0.1250}}
  >{\raggedright\arraybackslash}p{(\linewidth - 14\tabcolsep) * \real{0.1250}}
  >{\raggedright\arraybackslash}p{(\linewidth - 14\tabcolsep) * \real{0.1250}}
  >{\raggedright\arraybackslash}p{(\linewidth - 14\tabcolsep) * \real{0.1250}}
  >{\raggedright\arraybackslash}p{(\linewidth - 14\tabcolsep) * \real{0.1250}}
  >{\raggedright\arraybackslash}p{(\linewidth - 14\tabcolsep) * \real{0.1250}}
  >{\raggedright\arraybackslash}p{(\linewidth - 14\tabcolsep) * \real{0.1250}}
  >{\raggedright\arraybackslash}p{(\linewidth - 14\tabcolsep) * \real{0.1250}}@{}}
\toprule\noalign{}
\begin{minipage}[b]{\linewidth}\raggedright
RUNSPEC
\end{minipage} & \begin{minipage}[b]{\linewidth}\raggedright
GRID
\end{minipage} & \begin{minipage}[b]{\linewidth}\raggedright
EDIT
\end{minipage} & \begin{minipage}[b]{\linewidth}\raggedright
PROPS
\end{minipage} & \begin{minipage}[b]{\linewidth}\raggedright
REGIONS
\end{minipage} & \begin{minipage}[b]{\linewidth}\raggedright
SOLUTION
\end{minipage} & \begin{minipage}[b]{\linewidth}\raggedright
SUMMARY
\end{minipage} & \begin{minipage}[b]{\linewidth}\raggedright
SCHEDULE
\end{minipage} \\
\midrule\noalign{}
\endhead
\bottomrule\noalign{}
\endlastfoot
\end{longtable}

\subsubsection{Description}\label{description-69}

The \hyperref[REF_HEADING_KEYWORD_H2STORE]{H2STORE} keyword activates
the hydrogen (H2) storage model for the run to account for both hydrogen
and water phase solubility.
\hyperref[REF_HEADING_KEYWORD_H2STORE]{H2STORE} is similar to CO2STORE,
which activates the carbon dioxide (CO2) storage model.

The H2-Brine PVT model computes the PVT properties such as density,
viscosity, and enthalpy internally as functions of pressure,
temperature, and composition by using analytic correlations and models
from the literature rather than by interpolation from tabulated values.
These values are transformed to the standard black-oil equivalent PVT
tables internally by the simulator. Dissolved hydrogen in brine is
modeled using Li et al.~{[}Dedong Li, Christof Beyer, Sebastian Bauer, A
unified phase equilibrium model for hydrogen solubility and solution
density, International Journal of Hydrogen Energy, Volume 43, Issue 1,
2018, Pages 512-529.{]}. Hydrogen gas density was modeled using Leachman
et al.~{[}J. W. Leachman, R. T Jacobsen, S. G. Penoncello, E. W. Lemmon;
Fundamental Equations of State for Parahydrogen, Normal Hydrogen, and
Orthohydrogen. J. Phys. Chem. Ref. Data 1 September 2009; 38 (3):
721--748.{]}. Other properties are modeled the same way as in CO2STORE
but modified for hydrogen-brine (a similar approach was recently
reported by Raad et al.~{[}Seyed Mostafa Jafari Raad, Ehsan Ranjbar,
Hassan Hassanzadeh, Yuri Leonenko, Hydrogen-brine mixture PVT data for
reservoir simulation of hydrogen storage in deep saline aquifers,
International Journal of Hydrogen Energy, Volume 48, Issue 2, 2023,
Pages 696-708.{]}). A full description of the underlying PVT models used
by CO2STORE is described by Sandve et al.~{[}Tor Harald Sandve, Sarah E.
Gasda, Atgeirr Rasmussen, and Alf Birger Rustad. Convective dissolution
in field scale CO2 storage simulation using the OPM Flow simulator.
Submitted to TCCS 11 -- Trondheim Conference on CO2 Capture, Transport
and Storage Trondheim, Norway -- June 21-23, 2021.{]}. This means that
the normal PVT keywords like DENSITY, PVTO, PVDG etc. are not required
by OPM Flow when this model is activated, and if entered will be ignored
by the simulator. Note that the H2-Brine PVT properties depend on the
temperature and salinity and these must therefore be entered in the
PROPS or SOLUTION sections. The reservoir temperature can be defined
using, e.g., the RTEMP keyword. Region based salinity can be provided
using the SALINITY keyword.

The \hyperref[REF_HEADING_KEYWORD_H2STORE]{H2STORE} keyword must be used
with either: (1) the GAS and WATER keywords (or alternatively the GASWAT
keyword), or (2) the GAS and OIL keywords in the RUNSPEC section. It is
recommended that the standard option (1) is used.

The DISGASW keyword can also be used with option (1) to model
dissolution of H2 in the Brine.

Option (1) has the advantage that it can be used with the VAPWAT and
PRECSALT keywords to model the impact of both vaporization of residual
water and salt precipitation in the near wellbore region on injectivity
of H2 injection wells.

In option (1), the GAS and WATER (or GASWAT) keywords declare that the
gas and water phases are present in the model. When the
\hyperref[REF_HEADING_KEYWORD_H2STORE]{H2STORE} option is used the water
phase represents the brine and the gas phase represents H2. Note that
the input and output keywords need to be consistent with this
assumption, e.g., GSF (gas saturation function) and WSF (water
saturation function) should be used for the H2-Brine relative
permeability, etc.

Although, the DISGAS and VAPOIL keywords can be used with option (2) to
model water vaporization and H2 dissolution, salt precipitation is not
currently supported with option (2).

In option (2), the GAS and OIL keywords declare that the gas and oil
phases are present in the model. Internally when
\hyperref[REF_HEADING_KEYWORD_H2STORE]{H2STORE} is used the oil phase
refers to the brine and the gas phase to H2. Again, the input and output
keywords need to be consistent with this assumption, e.g., SGOF (gas-oil
relative permeability) is used for the H2-Brine relative permeability,
FOIP (Field Oil-In-Place) shows the total amount of brine in the
reservoir, etc.

There is no data required for this keyword and there is no terminating
``/'' for this keyword.

\subsubsection{Examples}\label{examples-2}

The first example shows the standard useage of
\hyperref[REF_HEADING_KEYWORD_H2STORE]{H2STORE} with Option (1) the
Gas-Water model (GASWAT). Here we also activate the dissolved gas in
water (DISGASW) and vaporized water in gas (VAPWAT) options.

\begin{verbatim}
-- ==============================================================================
--
-- RUNSPEC SECTION
--
-- ==============================================================================
RUNSPEC
-- ------------------------------------------------------------------------------ -- FLUID TYPES AND TRACER OPTIONS
-- ------------------------------------------------------------------------------
--
--       ACTIVATE H2 STORAGE IN THE MODEL (OPM FLOW H2 STORAGE KEYWORD)
--
H2STORE
--
--       ACTIVATE GAS-WATER THE MODEL (OPM FLOW KEYWORD)
--
GASWAT
--
--       DISSOLVED GAS IN WATER IS PRESENT IN THE RUN (OPM FLOW KEYWORD)
--
DISGASW
--
--       VAPORIZED WATER IN DRY/WET GAS IS PRESENT IN THE RUN (OPM FLOW KEYWORD)
--
VAPWAT
\end{verbatim}

The second part of the example covers the data required for the PROPS
section, in which the two-phase relative permeability functions are set
using GSF and WSF keywords.

\begin{verbatim}
-- ==============================================================================
--
-- PROPS SECTION
--
-- ==============================================================================
PROPS
--
--       RESERVOIR
--       TEMPERATURE
--       -----------
RTEMP
         80.0                                              / RESERVOIR TEMP
--
--       ROCK COMPRESSIBILITY
--
--       REFERENCE PRESSURE IS TAKEN FROM THE HCPV WEIGHTED RESERVOIR PRESSURE
--
--       REF PRES  CF
--       BARSA     1/BARSA
--       --------  --------
ROCK
         200.0     5.0E-05                                 / ROCK COMPRESSIBILITY
--
--       GAS RELATIVE PERMEABILITY TABLES (OPM FLOW KEYWORD)
--
GSF
--       SGAS       KRG        PCGW
--       FRAC                  PSIA
--       --------  --------   -------
            0.000    0.000       0.0
            0.080    0.001       0.0
            0.170    0.010       0.0
            0.350    0.050       0.0
            0.530    0.200       0.0
            0.620    0.350       0.0
            0.650    0.390       0.0
            0.710    0.560       0.0
            0.800    1.000       0.0   / TABLE NO. 01
--
--       WATER RELATIVE PERMEABILITY TABLES (OPM FLOW KEYWORD)
--
WSF
--       SWAT       KRW
--       FRAC
--       --------   --------
            0.200   0.0000
            0.400   0.1000
            0.800   0.5000
            1.000   1.0000             / TABLE NO. 01
--
--       SET SALINITY FOR ALL CELLS (OPM FLOW KEYWORD)
--
SALINITY
0.7                                                       /
\end{verbatim}

No other data is required to define the fluid and rock properties in the
PROPS section as the data is generated from internal analytic
correlations and models by the simulator. Finally, note that units for
salinity are to the 10-3, thus for metric units we have 10-3 x kg-M/kg.

The third part of the example covers initializing the model in the
SOLUTION section. Here we use the EQUIL(EQLOPT6) parameter equal to one,
to use table number one of the RVWVD keyword, in order to set the
vaporized water versus depth distribution for the model.

\begin{verbatim}
-- ==============================================================================
--
-- SOLUTION SECTION
--
-- ==============================================================================
SOLUTION
--
--       SYSTEM IS SATURATED WITH WATER
--
--       DATUM   DATUM   OWC     PCOW   GOC    PCGO   RS   RV   N    E300  RVW
--       DEPTH   PRESS   DEPTH   ----   DEPTH  ----   OPT  OPT  OPT  OPT   OPT
EQUIL
         2000.0  200.0   1800.0  0.00  1800.0  0.00   1*   1*   1*   2*    1   /
--
--       WATER VAPOR RATIO VS DEPTH (OPM FLOW KEYWORD)
--
--       DEPTH    RVW
--                STB/MSCF
--       ------   --------
RVWVD
         1000.0   0.000
         3000.0   0.000                              / RVW VS DEPTH EQUIL REGN 01
\end{verbatim}

In the SUMMARY section, the simulator supports summary vectors specific
to CO2 storage (see Section 11.1.12 Option Specific Variables -
CO2STORE/H2STORE Model) many of these can also be used for H2 storage
including those shown below.

\begin{verbatim}
-- ==============================================================================
--
-- SUMMARY SECTION
--
-- ==============================================================================
SUMMARY
--
--       EXPORT STANDARD SUMMARY VARIABLE VECTORS TO FILE
--
ALL
--
--       FIELD H2 DISOLVED IN WATER PHASE
FWCD
--
--       FIELD H2 TRAPPED IN GAS PHASE
FGCDI
--
--       FIELD H2 MOBILE IN GAS PHASE
FGCDM
\end{verbatim}

The final part of the example covers the SCHEDULE section. The standard
WCONINJE keyword is then used to set the gas injection rate, in this
case 100,000 sm3/day of H2.

\begin{verbatim}
-- ==============================================================================
--
-- SCHEDULE SECTION
--
-- ==============================================================================
SCHEDULE
--
--       RESTART CONTROL BASIC (LAST=1, ALL=2, YEARLY=4, MONTHLY=5, TSTEP=6)
--
RPTRST
        'BASIC=1' ALLPROPS TEMP RSWSAT RVWSAT                                   /
--
--       WELL SPECIFICATION DATA
--
-- WELL  GROUP     LOCATION  BHP    PHASE  DRAIN  INFLOW  OPEN  CROSS  PVT
-- NAME  NAME        I    J  DEPTH  FLUID  AREA   EQUANS  SHUT  FLOW   TABLE
WELSPECS
INJ1     'G1'       25   25  1*     GAS                                         /
/
--
--       WELL CONNECTION DATA
--
-- WELL  --- LOCATION ---  OPEN   SAT   CONN   WELL   KH    SKIN   D     DIR
-- NAME   II  JJ  K1  K2   SHUT   TAB   FACT   DIA    FACT  FACT   FACT  PEN
COMPDAT
INJ1      25  25   1  10  OPEN    1*    1*    0.3                              /
/
--
--       WELL INJECTION CONTROLS
--
-- WELL  FLUID  OPEN/  CNTL  SURF   RESV   BHP   THP   VFP
-- NAME  TYPE   SHUT   MODE  RATE   RATE   PRES  PRES  TABLE
WCONINJE
INJ1     GAS    OPEN   RATE 100000  1*     300                                 /
/
--
--       ADVANCE SIMULATION BY REPORTING TIME
--
TSTEP
         36*30.4375
/
\end{verbatim}

Note in order to get the liquid phase mole fractions of H2, that is, the
mole fractions of H2 in the water phase (XMFH2), and the vapor phase
mole fractions (YMFWAT) to the restart file, one must use the command
line parameter enable-opm-rst-file set equal to true.

The second example shows how to use
\hyperref[REF_HEADING_KEYWORD_H2STORE]{H2STORE} with the alternative
option (2). The example below declares that the hydrogen storage model
is active for the run to account for both hydrogen and water phase
solubility using OPM Flow's H2-Brine PVT model. Option (2) is used where
the OIL phase refers to the brine and the GAS phase to H2.

\begin{verbatim}
-- ==============================================================================
--
-- RUNSPEC SECTION
--
-- ==============================================================================
RUNSPEC
-- ------------------------------------------------------------------------------ -- FLUID TYPES AND TRACER OPTIONS
-- ------------------------------------------------------------------------------ --
--       ACTIVATE H2 STORAGE IN THE MODEL (OPM FLOW H2 STORAGE KEYWORD)
--
H2STORE
--
--       OIL PHASE IS PRESENT IN THE RUN BUT IS THE BRINE PHASE FOR H2STORE
--
OIL
--
--       GAS PHASE IS PRESENT IN THE RUN BUT IS THE H2 PHASE FOR H2STORE
--
GAS
--
--       DISSOLVED GAS IN LIVE OIL IS PRESENT IN THE RUN (H2 IN BRINE)
--
DISGAS
--
--       VAPORIZED OIL IN WET GAS IS PRESENT IN THE RUN (WATER IN H2)
--
VAPOIL
\end{verbatim}

The second part of the example covers the data required for the PROPS
section, in which the input keywords need to be consistent with the OIL
phase referring to the Brine and the GAS to H2; that is SGOF (gas-oil
relative permeability) is used to define the H2-Brine relative
permeability table.

\begin{verbatim}
-- ==============================================================================
--
-- PROPS SECTION
--
-- ==============================================================================
PROPS
--
--       RESERVOIR
--       TEMPERATURE
--       -----------
RTEMP
90.0                                              / RESERVOIR TEMPERATURE
--
--       ROCK COMPRESSIBILITY
--
--       REFERENCE PRESSURE IS TAKEN FROM THE HCPV WEIGHTED RESERVOIR PRESSURE
--
--       REF PRES  CF
--       BARSA     1/BARSA
--       --------  --------
ROCK
         200.0     5.0E-05                                 / ROCK COMPRESSIBILITY
--
--       GAS-OIL RELATIVE PERMEABILITY TABLES (SGOF) – H2STORE PHASES
SGOF
--       SG         KRG       KROG      PCOG
--       FRAC                           PSIA
--       -------   --------  -------   -------
         0.00000   0.000000  1.00000    0.0000
         1.00000   1.000000  0.00000    0.0000             / TABLE No. 01
--
--       SET SALINITY FOR ALL CELLS (OPM FLOW KEYWORD)
--
SALINITY
0.7                                                       /
\end{verbatim}

The third part and final part of the example covers initializing the
model in the SOLUTION section. Here we set the EQUIL(EQLOPT1 and
EQLOPT2) parameters equal to one, to use table number one of the RSVD
and RVVD keywords, in order to set the initial dissolved H2 and
vaporised water versus depth distribution for the model.

\begin{verbatim}
-- ==============================================================================
--
-- SOLUTION SECTION
--
-- ==============================================================================
SOLUTION
--
--       SYSTEM IS SATURATED WITH OIL (BRINE)
--
--       DATUM   DATUM   OWC     PCOW   GOC    PCGO   RS   RV   N    E300  RVW
--       DEPTH   PRESS   DEPTH   ----   DEPTH  ----   OPT  OPT  OPT  OPT   OPT
EQUIL
         2000.0  200.0   2200.0  0.00  1800.0  0.00   1    1    1*   2*    1*   /
--
--       DEPTH    RS
--                MSCF/STB
--       ------   --------
RSVD
         1000.0   0.000
         3000.0   0.000                                      /
--
--       DEPTH    RV
--                STB/MSCF
--       ------   --------
RVVD
         1000.0   0.000
         3000.0   0.000                                       /
\end{verbatim}

\subsection{HMDIMS -- Define History Match Gradient Parameter
Dimensions}\label{hmdims-define-history-match-gradient-parameter-dimensions}

\begin{longtable}[]{@{}
  >{\raggedright\arraybackslash}p{(\linewidth - 14\tabcolsep) * \real{0.1250}}
  >{\raggedright\arraybackslash}p{(\linewidth - 14\tabcolsep) * \real{0.1250}}
  >{\raggedright\arraybackslash}p{(\linewidth - 14\tabcolsep) * \real{0.1250}}
  >{\raggedright\arraybackslash}p{(\linewidth - 14\tabcolsep) * \real{0.1250}}
  >{\raggedright\arraybackslash}p{(\linewidth - 14\tabcolsep) * \real{0.1250}}
  >{\raggedright\arraybackslash}p{(\linewidth - 14\tabcolsep) * \real{0.1250}}
  >{\raggedright\arraybackslash}p{(\linewidth - 14\tabcolsep) * \real{0.1250}}
  >{\raggedright\arraybackslash}p{(\linewidth - 14\tabcolsep) * \real{0.1250}}@{}}
\toprule\noalign{}
\begin{minipage}[b]{\linewidth}\raggedright
RUNSPEC
\end{minipage} & \begin{minipage}[b]{\linewidth}\raggedright
GRID
\end{minipage} & \begin{minipage}[b]{\linewidth}\raggedright
EDIT
\end{minipage} & \begin{minipage}[b]{\linewidth}\raggedright
PROPS
\end{minipage} & \begin{minipage}[b]{\linewidth}\raggedright
REGIONS
\end{minipage} & \begin{minipage}[b]{\linewidth}\raggedright
SOLUTION
\end{minipage} & \begin{minipage}[b]{\linewidth}\raggedright
SUMMARY
\end{minipage} & \begin{minipage}[b]{\linewidth}\raggedright
SCHEDULE
\end{minipage} \\
\midrule\noalign{}
\endhead
\bottomrule\noalign{}
\endlastfoot
\end{longtable}

\subsubsection{Description}\label{description-70}

This keyword, HMDIMS, defines the maximum parameter dimensions for the
History Match Gradient option.

This keyword is not supported by OPM Flow but would change the results
if supported so the simulation will be stopped.

\subsection{HYST -- Activate the Hysteresis Option
(Retired)}\label{hyst-activate-the-hysteresis-option-retired}

\begin{longtable}[]{@{}
  >{\raggedright\arraybackslash}p{(\linewidth - 14\tabcolsep) * \real{0.1250}}
  >{\raggedright\arraybackslash}p{(\linewidth - 14\tabcolsep) * \real{0.1250}}
  >{\raggedright\arraybackslash}p{(\linewidth - 14\tabcolsep) * \real{0.1250}}
  >{\raggedright\arraybackslash}p{(\linewidth - 14\tabcolsep) * \real{0.1250}}
  >{\raggedright\arraybackslash}p{(\linewidth - 14\tabcolsep) * \real{0.1250}}
  >{\raggedright\arraybackslash}p{(\linewidth - 14\tabcolsep) * \real{0.1250}}
  >{\raggedright\arraybackslash}p{(\linewidth - 14\tabcolsep) * \real{0.1250}}
  >{\raggedright\arraybackslash}p{(\linewidth - 14\tabcolsep) * \real{0.1250}}@{}}
\toprule\noalign{}
\begin{minipage}[b]{\linewidth}\raggedright
RUNSPEC
\end{minipage} & \begin{minipage}[b]{\linewidth}\raggedright
GRID
\end{minipage} & \begin{minipage}[b]{\linewidth}\raggedright
EDIT
\end{minipage} & \begin{minipage}[b]{\linewidth}\raggedright
PROPS
\end{minipage} & \begin{minipage}[b]{\linewidth}\raggedright
REGIONS
\end{minipage} & \begin{minipage}[b]{\linewidth}\raggedright
SOLUTION
\end{minipage} & \begin{minipage}[b]{\linewidth}\raggedright
SUMMARY
\end{minipage} & \begin{minipage}[b]{\linewidth}\raggedright
SCHEDULE
\end{minipage} \\
\midrule\noalign{}
\endhead
\bottomrule\noalign{}
\endlastfoot
\end{longtable}

\subsubsection{Description}\label{description-71}

The HYST keyword activates the hysteresis option, the keyword should be
avoided and the hysteresis option should be enabled by the HYSER
parameter on the SATOTPS keyword in the RUNSPEC section.

This keyword is not supported by OPM Flow but would change the results
if supported so the simulation will be stopped.

\subsection{IMPES -- Activate Implicit Pressure Explicit Saturation
Solution
Option}\label{impes-activate-implicit-pressure-explicit-saturation-solution-option}

\begin{longtable}[]{@{}
  >{\raggedright\arraybackslash}p{(\linewidth - 14\tabcolsep) * \real{0.1250}}
  >{\raggedright\arraybackslash}p{(\linewidth - 14\tabcolsep) * \real{0.1250}}
  >{\raggedright\arraybackslash}p{(\linewidth - 14\tabcolsep) * \real{0.1250}}
  >{\raggedright\arraybackslash}p{(\linewidth - 14\tabcolsep) * \real{0.1250}}
  >{\raggedright\arraybackslash}p{(\linewidth - 14\tabcolsep) * \real{0.1250}}
  >{\raggedright\arraybackslash}p{(\linewidth - 14\tabcolsep) * \real{0.1250}}
  >{\raggedright\arraybackslash}p{(\linewidth - 14\tabcolsep) * \real{0.1250}}
  >{\raggedright\arraybackslash}p{(\linewidth - 14\tabcolsep) * \real{0.1250}}@{}}
\toprule\noalign{}
\begin{minipage}[b]{\linewidth}\raggedright
RUNSPEC
\end{minipage} & \begin{minipage}[b]{\linewidth}\raggedright
GRID
\end{minipage} & \begin{minipage}[b]{\linewidth}\raggedright
EDIT
\end{minipage} & \begin{minipage}[b]{\linewidth}\raggedright
PROPS
\end{minipage} & \begin{minipage}[b]{\linewidth}\raggedright
REGIONS
\end{minipage} & \begin{minipage}[b]{\linewidth}\raggedright
SOLUTION
\end{minipage} & \begin{minipage}[b]{\linewidth}\raggedright
SUMMARY
\end{minipage} & \begin{minipage}[b]{\linewidth}\raggedright
SCHEDULE
\end{minipage} \\
\midrule\noalign{}
\endhead
\bottomrule\noalign{}
\endlastfoot
\end{longtable}

\subsubsection{Description}\label{description-72}

The IMPES keyword activates the Implicit Pressure Explicit Saturation
formulation and solution options, commonly know as IMPES. OPM Flow uses
a different numerical scheme which makes this keyword redundant; hence,
OPM Flow ignores this keyword. It is documented here for completeness.

There is no data required for this keyword and there is no terminating
``/'' for this keyword.

See section 2.2 Running OPM Flow 2023-04 From The Command Line on how to
invoke various numerical schemes via the OPM Flow command line
interface.

\subsubsection{Example}\label{example-53}

\begin{verbatim}
--
--       ACTIVATE THE IMPES SOLUTION OPTION
--
IMPES
\end{verbatim}

The above example switches on the IMPES solution option; however, this
has no effect in OPM Flow input decks.

\subsection{IMPLICIT -- Activate Fully Implicit Solution
Option}\label{implicit-activate-fully-implicit-solution-option}

\begin{longtable}[]{@{}
  >{\raggedright\arraybackslash}p{(\linewidth - 14\tabcolsep) * \real{0.1250}}
  >{\raggedright\arraybackslash}p{(\linewidth - 14\tabcolsep) * \real{0.1250}}
  >{\raggedright\arraybackslash}p{(\linewidth - 14\tabcolsep) * \real{0.1250}}
  >{\raggedright\arraybackslash}p{(\linewidth - 14\tabcolsep) * \real{0.1250}}
  >{\raggedright\arraybackslash}p{(\linewidth - 14\tabcolsep) * \real{0.1250}}
  >{\raggedright\arraybackslash}p{(\linewidth - 14\tabcolsep) * \real{0.1250}}
  >{\raggedright\arraybackslash}p{(\linewidth - 14\tabcolsep) * \real{0.1250}}
  >{\raggedright\arraybackslash}p{(\linewidth - 14\tabcolsep) * \real{0.1250}}@{}}
\toprule\noalign{}
\begin{minipage}[b]{\linewidth}\raggedright
RUNSPEC
\end{minipage} & \begin{minipage}[b]{\linewidth}\raggedright
GRID
\end{minipage} & \begin{minipage}[b]{\linewidth}\raggedright
EDIT
\end{minipage} & \begin{minipage}[b]{\linewidth}\raggedright
PROPS
\end{minipage} & \begin{minipage}[b]{\linewidth}\raggedright
REGIONS
\end{minipage} & \begin{minipage}[b]{\linewidth}\raggedright
SOLUTION
\end{minipage} & \begin{minipage}[b]{\linewidth}\raggedright
SUMMARY
\end{minipage} & \begin{minipage}[b]{\linewidth}\raggedright
SCHEDULE
\end{minipage} \\
\midrule\noalign{}
\endhead
\bottomrule\noalign{}
\endlastfoot
\end{longtable}

\subsubsection{Description}\label{description-73}

The IMPLICIT keyword activates the Fully Implicit Solution formulation
and solution options. OPM Flow uses a different numerical scheme which
makes this keyword redundant; hence, OPM Flow ignores this keyword. It
is documented here for completeness. The keyword as the same function as
the FULLIMP keyword in the RUNSPEC section.

There is no data required for this keyword and there is no terminating
``/'' for this keyword.

See section 2.2 Running OPM Flow 2023-04 From The Command Line on how to
invoke various numerical schemes via the OPM Flow command line
interface.

\subsubsection{Example}\label{example-54}

\begin{verbatim}
--
--       ACTIVATES THE FULLY IMPLICIT SOLUTION OPTION
--
IMPLICIT
\end{verbatim}

The above example switches on the fully implicit solution option;
however, this has no effect in OPM Flow input decks.

\subsection{INCLUDE -- Load Another Data File at the Current
Position}\label{include-load-another-data-file-at-the-current-position-1}

The INCLUDE keyword informs OPM Flow to continue reading input data from
the specified INCLUDE file. When the end of the INCLUDE file is reached,
or the ENDINC keyword is encountered in the included file, input data is
read from the next keyword in the current file.

See INCLUDE -- Load Another Data File at the Current Position in the
GLOBAL section for a full description.

\subsection{INSPEC -- Activate the INSPEC File
Option}\label{inspec-activate-the-inspec-file-option}

\begin{longtable}[]{@{}
  >{\raggedright\arraybackslash}p{(\linewidth - 14\tabcolsep) * \real{0.1250}}
  >{\raggedright\arraybackslash}p{(\linewidth - 14\tabcolsep) * \real{0.1250}}
  >{\raggedright\arraybackslash}p{(\linewidth - 14\tabcolsep) * \real{0.1250}}
  >{\raggedright\arraybackslash}p{(\linewidth - 14\tabcolsep) * \real{0.1250}}
  >{\raggedright\arraybackslash}p{(\linewidth - 14\tabcolsep) * \real{0.1250}}
  >{\raggedright\arraybackslash}p{(\linewidth - 14\tabcolsep) * \real{0.1250}}
  >{\raggedright\arraybackslash}p{(\linewidth - 14\tabcolsep) * \real{0.1250}}
  >{\raggedright\arraybackslash}p{(\linewidth - 14\tabcolsep) * \real{0.1250}}@{}}
\toprule\noalign{}
\begin{minipage}[b]{\linewidth}\raggedright
RUNSPEC
\end{minipage} & \begin{minipage}[b]{\linewidth}\raggedright
GRID
\end{minipage} & \begin{minipage}[b]{\linewidth}\raggedright
EDIT
\end{minipage} & \begin{minipage}[b]{\linewidth}\raggedright
PROPS
\end{minipage} & \begin{minipage}[b]{\linewidth}\raggedright
REGIONS
\end{minipage} & \begin{minipage}[b]{\linewidth}\raggedright
SOLUTION
\end{minipage} & \begin{minipage}[b]{\linewidth}\raggedright
SUMMARY
\end{minipage} & \begin{minipage}[b]{\linewidth}\raggedright
SCHEDULE
\end{minipage} \\
\midrule\noalign{}
\endhead
\bottomrule\noalign{}
\endlastfoot
\end{longtable}

\subsubsection{Description}\label{description-74}

This keyword switches on the writing of the INIT Index file that
specifies and defines the format and data type written to the
\emph{.INIT data file. The }.INIT data file contains the static data
specified in the GRID, PROPS and REGIONS sections. For example, the
PORO, PERM and NTG arrays from the GRID section. The data is used in
post-processing software, for example OPM ResInsight, to visualize the
static grid properties.

The INIT Index file can either be written out in formatted form as ASCII
i.e.~text files, if the FMTOUT keyword has been activated
(\emph{.FINSPEC), or binary format (}.INSPEC) if the FMTOUT keyword has
not been activated. If the INIT keyword in the GRID section has been
used to switch on the writing of the *. INIT data file then a binary
INIT Index file is automatically written out as well, unless the
NOINSPEC keyword in the RUNSPEC section has been used to switch off the
writing of the INIT Index file.

Note that most post-processing software require the \emph{.INSPEC file
to load the }.INIT data set, although OPM ResInsight does not require
this file to be able to load the *.INIT data file.

There is no data required for this keyword and there is no terminating
``/'' for this keyword.

This keyword is not supported by OPM Flow but would change the results
if supported so the simulation will be stopped.

\subsubsection{Example}\label{example-55}

\begin{verbatim}
--
--       ACTIVATE WRITING THE INIT INDEX FILE FOR POST-PROCESSING
--
INSPEC
\end{verbatim}

The above example switches on the writing of the INIT Index file for
post-processing in ResInsight.

\subsection{LAB -- Activate the Laboratory System of Units for the
Model}\label{lab-activate-the-laboratory-system-of-units-for-the-model}

\begin{longtable}[]{@{}
  >{\raggedright\arraybackslash}p{(\linewidth - 14\tabcolsep) * \real{0.1250}}
  >{\raggedright\arraybackslash}p{(\linewidth - 14\tabcolsep) * \real{0.1250}}
  >{\raggedright\arraybackslash}p{(\linewidth - 14\tabcolsep) * \real{0.1250}}
  >{\raggedright\arraybackslash}p{(\linewidth - 14\tabcolsep) * \real{0.1250}}
  >{\raggedright\arraybackslash}p{(\linewidth - 14\tabcolsep) * \real{0.1250}}
  >{\raggedright\arraybackslash}p{(\linewidth - 14\tabcolsep) * \real{0.1250}}
  >{\raggedright\arraybackslash}p{(\linewidth - 14\tabcolsep) * \real{0.1250}}
  >{\raggedright\arraybackslash}p{(\linewidth - 14\tabcolsep) * \real{0.1250}}@{}}
\toprule\noalign{}
\begin{minipage}[b]{\linewidth}\raggedright
RUNSPEC
\end{minipage} & \begin{minipage}[b]{\linewidth}\raggedright
GRID
\end{minipage} & \begin{minipage}[b]{\linewidth}\raggedright
EDIT
\end{minipage} & \begin{minipage}[b]{\linewidth}\raggedright
PROPS
\end{minipage} & \begin{minipage}[b]{\linewidth}\raggedright
REGIONS
\end{minipage} & \begin{minipage}[b]{\linewidth}\raggedright
SOLUTION
\end{minipage} & \begin{minipage}[b]{\linewidth}\raggedright
SUMMARY
\end{minipage} & \begin{minipage}[b]{\linewidth}\raggedright
SCHEDULE
\end{minipage} \\
\midrule\noalign{}
\endhead
\bottomrule\noalign{}
\endlastfoot
\end{longtable}

\subsubsection{Description}\label{description-75}

This keyword switches on the LABORATORY system of units for the model.

OPM Flow has three sets of units, namely: METRIC, FIELD and LAB and one
of these keyword should be invoked in the RUNSPEC section to avoid any
ambiguity. Both the simulator input and output units are controlled by
including one of the METRIC, FIELD or LAB keywords in the RUNSPEC
section of the input file.

There is no data required for this keyword and there is no terminating
``/'' for this keyword.

\subsubsection{Example}\label{example-56}

\begin{verbatim}
--
--       SWITCH ON THE LABORATORY SYSTEM OF UNITS FOR BOTH INPUT AND OUTPUT
--
LAB
\end{verbatim}

The above example switches on the LABORATORY system of units for the
model.

\subsection{LGR -- Define Local Grid Refinement Dimensions and
Parameters}\label{lgr-define-local-grid-refinement-dimensions-and-parameters}

\begin{longtable}[]{@{}
  >{\raggedright\arraybackslash}p{(\linewidth - 14\tabcolsep) * \real{0.1250}}
  >{\raggedright\arraybackslash}p{(\linewidth - 14\tabcolsep) * \real{0.1250}}
  >{\raggedright\arraybackslash}p{(\linewidth - 14\tabcolsep) * \real{0.1250}}
  >{\raggedright\arraybackslash}p{(\linewidth - 14\tabcolsep) * \real{0.1250}}
  >{\raggedright\arraybackslash}p{(\linewidth - 14\tabcolsep) * \real{0.1250}}
  >{\raggedright\arraybackslash}p{(\linewidth - 14\tabcolsep) * \real{0.1250}}
  >{\raggedright\arraybackslash}p{(\linewidth - 14\tabcolsep) * \real{0.1250}}
  >{\raggedright\arraybackslash}p{(\linewidth - 14\tabcolsep) * \real{0.1250}}@{}}
\toprule\noalign{}
\begin{minipage}[b]{\linewidth}\raggedright
RUNSPEC
\end{minipage} & \begin{minipage}[b]{\linewidth}\raggedright
GRID
\end{minipage} & \begin{minipage}[b]{\linewidth}\raggedright
EDIT
\end{minipage} & \begin{minipage}[b]{\linewidth}\raggedright
PROPS
\end{minipage} & \begin{minipage}[b]{\linewidth}\raggedright
REGIONS
\end{minipage} & \begin{minipage}[b]{\linewidth}\raggedright
SOLUTION
\end{minipage} & \begin{minipage}[b]{\linewidth}\raggedright
SUMMARY
\end{minipage} & \begin{minipage}[b]{\linewidth}\raggedright
SCHEDULE
\end{minipage} \\
\midrule\noalign{}
\endhead
\bottomrule\noalign{}
\endlastfoot
\end{longtable}

\subsubsection{Description}\label{description-76}

The LGR keyword defines the maximum dimensions and parameters for the
Local Grid Refinement (``LGR'') option.

Currently, OPM Flow does not support the local grid refinement feature
and therefore this keyword is ignored by the simulator, but is
documented here for completeness.

\begin{longtable}[]{@{}
  >{\raggedright\arraybackslash}p{(\linewidth - 6\tabcolsep) * \real{0.2500}}
  >{\raggedright\arraybackslash}p{(\linewidth - 6\tabcolsep) * \real{0.2500}}
  >{\raggedright\arraybackslash}p{(\linewidth - 6\tabcolsep) * \real{0.2500}}
  >{\raggedright\arraybackslash}p{(\linewidth - 6\tabcolsep) * \real{0.2500}}@{}}
\toprule\noalign{}
\begin{minipage}[b]{\linewidth}\raggedright
No.
\end{minipage} & \begin{minipage}[b]{\linewidth}\raggedright
Name
\end{minipage} & \begin{minipage}[b]{\linewidth}\raggedright
Description
\end{minipage} & \begin{minipage}[b]{\linewidth}\raggedright
Default
\end{minipage} \\
\midrule\noalign{}
\endhead
\bottomrule\noalign{}
\endlastfoot
1 & MAXLGR & A positive integer value that defines the maximum number of
LGRs in the model. & 0 \\
2 & MAXCLS & A positive integer value that defines the maximum number of
grid blocks in all the LGRs. & 0 \\
3 & MCOARS & A positive integer value that defines the maximum number of
amalgamated coarse grid blocks in the model. & 0 \\
4 & MAMALG & A positive integer value that defines the maximum number of
LGR amalgamations in the model. & 0 \\
5 & MXLALG & A positive integer value that defines the maximum number of
LGRs in any amalgamation in the model. & 0 \\
6 & LSTACK & A positive integer that defines the maximum number of
previous search directions stored by the linear solver for the LGR. See
the NSTACK keyword in the RUNSPEC section for a full description. &
10 \\
7 & INTOPT & A character string set to either INTERP to activate the
Quandalle {[}Quandalle, Philippe \& Besset, P. (1985). Reduction of Grid
Effects Due to Local Sub-Gridding in Simulations Using a Composite Grid.
10.2118/13527-MS.{]} pressure correction, or NOINTERP to deactivate this
option. The option applies bi-linear interpolation to the global cells
surrounding an LGR in order to improve the accuracy of the flow
calculations between the LGR and the host cells. & NOINTERP \\
8 & NCHCOR & A positive integer value that defines the maximum number of
grid blocks within a coarsened grid that overlap parallel domain
boundaries for when the Parallel option has been invoked by the PARALLEL
keyword in the RUNSPEC section. OPM Flow uses a different numerical
scheme which makes this parameter redundant, see section 2.2 Running OPM
Flow 2023-04 From The Command Line on how to run OPM Flow in parallel
mode. & 0 \\
Notes: & & & \\
\end{longtable}

\emph{Table 5.19: LGR Keyword Description}

\subsubsection{Example}\label{example-57}

\begin{verbatim}
--
--       LOCAL GRID REFINEMENT DIMENSIONS AND PARAMETERS
--
--       LGR     LGR     LGR     LGR     LGR     LGR    LGR     LGR
--       MAXLGR  MAXCLS  MCOARS  MAMALG  MXLALG  LSTACK INTOPT  NCHCOR
LGR
         10      1000    1*      1*      1*      1*     INTERP  1*             /
\end{verbatim}

The above example sets the maximum number of LGRs to 10 and the maximum
number of grid blocks a LGR may contain to 1,000, and that the Quandalle
pressure correction should be used to improve the flow calculation.

\subsection{LGRCOPY -- Activate Local Grid Refinement
Inheritance}\label{lgrcopy-activate-local-grid-refinement-inheritance}

\begin{longtable}[]{@{}
  >{\raggedright\arraybackslash}p{(\linewidth - 14\tabcolsep) * \real{0.1250}}
  >{\raggedright\arraybackslash}p{(\linewidth - 14\tabcolsep) * \real{0.1250}}
  >{\raggedright\arraybackslash}p{(\linewidth - 14\tabcolsep) * \real{0.1250}}
  >{\raggedright\arraybackslash}p{(\linewidth - 14\tabcolsep) * \real{0.1250}}
  >{\raggedright\arraybackslash}p{(\linewidth - 14\tabcolsep) * \real{0.1250}}
  >{\raggedright\arraybackslash}p{(\linewidth - 14\tabcolsep) * \real{0.1250}}
  >{\raggedright\arraybackslash}p{(\linewidth - 14\tabcolsep) * \real{0.1250}}
  >{\raggedright\arraybackslash}p{(\linewidth - 14\tabcolsep) * \real{0.1250}}@{}}
\toprule\noalign{}
\begin{minipage}[b]{\linewidth}\raggedright
RUNSPEC
\end{minipage} & \begin{minipage}[b]{\linewidth}\raggedright
GRID
\end{minipage} & \begin{minipage}[b]{\linewidth}\raggedright
EDIT
\end{minipage} & \begin{minipage}[b]{\linewidth}\raggedright
PROPS
\end{minipage} & \begin{minipage}[b]{\linewidth}\raggedright
REGIONS
\end{minipage} & \begin{minipage}[b]{\linewidth}\raggedright
SOLUTION
\end{minipage} & \begin{minipage}[b]{\linewidth}\raggedright
SUMMARY
\end{minipage} & \begin{minipage}[b]{\linewidth}\raggedright
SCHEDULE
\end{minipage} \\
\midrule\noalign{}
\endhead
\bottomrule\noalign{}
\endlastfoot
\end{longtable}

\subsubsection{Description}\label{description-77}

The LGRCOPY keyword actives the Local Grid Refinement (``LGR'')
Inheritance option that allows the LGR to inherit the properties of the
global or host cell containing an LGR grid block when it is defined, as
opposed to the normal process of applying this transform at the end of
the GRID section. LGRCOPY can be used in the RUNSPEC, GRID and EDIT
sections. If used in the RUNSPEC section then the option is applied to
all LGRs defined in the input file, whereas if used in the GRID or EDIT
sections the keyword must be placed inside a LGR definition section,
that is between a CARFIN (Cartesian LGR grid) or RADFIN/RADFIN4 (radial
LGR grid) keyword and an ENDFIN keyword. In the latter case inheritance
is applied on an individual LGR basis.

Currently, OPM Flow does not support the local grid refinement feature
and therefore this keyword is ignored by the simulator.

\subsubsection{Example}\label{example-58}

The following example activates the LGR Inheritance option for all LGRs
in the model.

\begin{verbatim}
-- ==============================================================================
--
-- RUNSPEC SECTION
--
-- ==============================================================================
RUNSPEC
--
--       ACTIVATE LOCAL GRID REFINEMENT INHERITANCE
--
LGRCOPY
\end{verbatim}

\subsection{LICENSES -- Define Required Licenses for
Run}\label{licenses-define-required-licenses-for-run}

\begin{longtable}[]{@{}
  >{\raggedright\arraybackslash}p{(\linewidth - 14\tabcolsep) * \real{0.1250}}
  >{\raggedright\arraybackslash}p{(\linewidth - 14\tabcolsep) * \real{0.1250}}
  >{\raggedright\arraybackslash}p{(\linewidth - 14\tabcolsep) * \real{0.1250}}
  >{\raggedright\arraybackslash}p{(\linewidth - 14\tabcolsep) * \real{0.1250}}
  >{\raggedright\arraybackslash}p{(\linewidth - 14\tabcolsep) * \real{0.1250}}
  >{\raggedright\arraybackslash}p{(\linewidth - 14\tabcolsep) * \real{0.1250}}
  >{\raggedright\arraybackslash}p{(\linewidth - 14\tabcolsep) * \real{0.1250}}
  >{\raggedright\arraybackslash}p{(\linewidth - 14\tabcolsep) * \real{0.1250}}@{}}
\toprule\noalign{}
\begin{minipage}[b]{\linewidth}\raggedright
RUNSPEC
\end{minipage} & \begin{minipage}[b]{\linewidth}\raggedright
GRID
\end{minipage} & \begin{minipage}[b]{\linewidth}\raggedright
EDIT
\end{minipage} & \begin{minipage}[b]{\linewidth}\raggedright
PROPS
\end{minipage} & \begin{minipage}[b]{\linewidth}\raggedright
REGIONS
\end{minipage} & \begin{minipage}[b]{\linewidth}\raggedright
SOLUTION
\end{minipage} & \begin{minipage}[b]{\linewidth}\raggedright
SUMMARY
\end{minipage} & \begin{minipage}[b]{\linewidth}\raggedright
SCHEDULE
\end{minipage} \\
\midrule\noalign{}
\endhead
\bottomrule\noalign{}
\endlastfoot
\end{longtable}

\subsubsection{Description}\label{description-78}

This keyword defines the additional software licenses that are required
to invoke various licensed options in the commercial simulator at the
start of the run. The commercial simulator requests a license when
keywords associated with a licensed option is encountered in the input
deck, this may result in the license being unavailable at the time of
request and after the simulation has been initiated, resulting in the
run terminating. This keyword avoids this scenario by reserving the
license at the start of the run.

OPM Flow is an open source project and therefore there is no license
management of the various implemented options; hence this keyword is not
supported by OPM Flow but would change the results if supported so the
simulation will be stopped.

\subsection{LIVEOIL -- Activate the Live Oil Phase (Oil with Free and
Dissolved
Gas)}\label{liveoil-activate-the-live-oil-phase-oil-with-free-and-dissolved-gas}

\begin{longtable}[]{@{}
  >{\raggedright\arraybackslash}p{(\linewidth - 14\tabcolsep) * \real{0.1250}}
  >{\raggedright\arraybackslash}p{(\linewidth - 14\tabcolsep) * \real{0.1250}}
  >{\raggedright\arraybackslash}p{(\linewidth - 14\tabcolsep) * \real{0.1250}}
  >{\raggedright\arraybackslash}p{(\linewidth - 14\tabcolsep) * \real{0.1250}}
  >{\raggedright\arraybackslash}p{(\linewidth - 14\tabcolsep) * \real{0.1250}}
  >{\raggedright\arraybackslash}p{(\linewidth - 14\tabcolsep) * \real{0.1250}}
  >{\raggedright\arraybackslash}p{(\linewidth - 14\tabcolsep) * \real{0.1250}}
  >{\raggedright\arraybackslash}p{(\linewidth - 14\tabcolsep) * \real{0.1250}}@{}}
\toprule\noalign{}
\begin{minipage}[b]{\linewidth}\raggedright
RUNSPEC
\end{minipage} & \begin{minipage}[b]{\linewidth}\raggedright
GRID
\end{minipage} & \begin{minipage}[b]{\linewidth}\raggedright
EDIT
\end{minipage} & \begin{minipage}[b]{\linewidth}\raggedright
PROPS
\end{minipage} & \begin{minipage}[b]{\linewidth}\raggedright
REGIONS
\end{minipage} & \begin{minipage}[b]{\linewidth}\raggedright
SOLUTION
\end{minipage} & \begin{minipage}[b]{\linewidth}\raggedright
SUMMARY
\end{minipage} & \begin{minipage}[b]{\linewidth}\raggedright
SCHEDULE
\end{minipage} \\
\midrule\noalign{}
\endhead
\bottomrule\noalign{}
\endlastfoot
\end{longtable}

\subsubsection{Description}\label{description-79}

This keyword activates oil, free and dissolved gas in the model and
therefore makes the oil phase live oil {[}``Live'' oil is oil that
contains gas in solution, which is normally the case for most
conventional oil reservoirs. However, for oil reservoirs classified as
heavy oil reservoirs, the in situ dissolved gas may be negligible and
oil would then be classified as gas-free oil which is commonly referred
to as ``dead'' oil.{]} in the black-oil formulation, and is equivalent
to setting the phases present in the model to be oil, dissolved gas, gas
and water phases. Note if water is present in the model this needs to be
explicitly stated via the WATER keyword in the RUNSPEC section (see also
the BLACKOIL and DEADOIL keywords in the RUNSPEC section). The keyword
is used by the commercial simulator's compositional THERMAL option to
set the phases present in the model.

There is no data required for this keyword and there is no terminating
``/'' for this keyword.

This keyword is not supported by OPM Flow but has no effect on the
results so it will be ignored.

\subsubsection{Example}\label{example-59}

The following example activates the black-oil phases in the model.

\begin{verbatim}
--
--       ACTIVATE LIVE-OIL PHASE
--
LIVEOIL
\end{verbatim}

Alternatively one could explicitly declare the phases using the
following keywords in the RUNSPEC section.

\begin{verbatim}
--
--       OIL PHASE IS PRESENT IN THE RUN
--
OIL
--
--       DISSOLVED GAS IN LIVE OIL IS PRESENT IN THE RUN
--
DISGAS
--
--       GAS PHASE IS PRESENT IN THE RUN
--
GAS
--
--       WATER PHASE IS PRESENT IN THE RUN
--
WATER
\end{verbatim}

The above example switches on the oil, dissolved gas, gas and water
phases in the model.

\subsection{LOAD -- Load a SAVE File for a Fast
Restart}\label{load-load-a-save-file-for-a-fast-restart}

\begin{longtable}[]{@{}
  >{\raggedright\arraybackslash}p{(\linewidth - 14\tabcolsep) * \real{0.1250}}
  >{\raggedright\arraybackslash}p{(\linewidth - 14\tabcolsep) * \real{0.1250}}
  >{\raggedright\arraybackslash}p{(\linewidth - 14\tabcolsep) * \real{0.1250}}
  >{\raggedright\arraybackslash}p{(\linewidth - 14\tabcolsep) * \real{0.1250}}
  >{\raggedright\arraybackslash}p{(\linewidth - 14\tabcolsep) * \real{0.1250}}
  >{\raggedright\arraybackslash}p{(\linewidth - 14\tabcolsep) * \real{0.1250}}
  >{\raggedright\arraybackslash}p{(\linewidth - 14\tabcolsep) * \real{0.1250}}
  >{\raggedright\arraybackslash}p{(\linewidth - 14\tabcolsep) * \real{0.1250}}@{}}
\toprule\noalign{}
\begin{minipage}[b]{\linewidth}\raggedright
RUNSPEC
\end{minipage} & \begin{minipage}[b]{\linewidth}\raggedright
GRID
\end{minipage} & \begin{minipage}[b]{\linewidth}\raggedright
EDIT
\end{minipage} & \begin{minipage}[b]{\linewidth}\raggedright
PROPS
\end{minipage} & \begin{minipage}[b]{\linewidth}\raggedright
REGIONS
\end{minipage} & \begin{minipage}[b]{\linewidth}\raggedright
SOLUTION
\end{minipage} & \begin{minipage}[b]{\linewidth}\raggedright
SUMMARY
\end{minipage} & \begin{minipage}[b]{\linewidth}\raggedright
SCHEDULE
\end{minipage} \\
\midrule\noalign{}
\endhead
\bottomrule\noalign{}
\endlastfoot
\end{longtable}

\subsubsection{Description}\label{description-80}

The LOAD keyword loads a previously generated SAVE file to enable a fast
restart. A SAVE file contains all the data from a previous run's
RUNSPEC, GRID, EDIT, PROPS and REGIONS sections, and thus there is no
need for the simulator to calculate various parameters, including grid
block transmissibilities etc. This allows for the current run to restart
quicker than a conventional restart run using the RESTART keyword in the
SOLUTION section via a RESTART file (\emph{.UNRST or }.FUNRST etc.). The
keyword should be the first keyword in the input deck and the RUNSPEC,
GRID, EDIT, PROPS and REGIONS sections should be deleted from the input
deck.

This keyword is not supported by OPM Flow but would change the results
if supported so the simulation will be stopped.

\subsection{LOWSALT -- Activate the Low Salt Brine Phase in the Brine
Model}\label{lowsalt-activate-the-low-salt-brine-phase-in-the-brine-model}

\begin{longtable}[]{@{}
  >{\raggedright\arraybackslash}p{(\linewidth - 14\tabcolsep) * \real{0.1250}}
  >{\raggedright\arraybackslash}p{(\linewidth - 14\tabcolsep) * \real{0.1250}}
  >{\raggedright\arraybackslash}p{(\linewidth - 14\tabcolsep) * \real{0.1250}}
  >{\raggedright\arraybackslash}p{(\linewidth - 14\tabcolsep) * \real{0.1250}}
  >{\raggedright\arraybackslash}p{(\linewidth - 14\tabcolsep) * \real{0.1250}}
  >{\raggedright\arraybackslash}p{(\linewidth - 14\tabcolsep) * \real{0.1250}}
  >{\raggedright\arraybackslash}p{(\linewidth - 14\tabcolsep) * \real{0.1250}}
  >{\raggedright\arraybackslash}p{(\linewidth - 14\tabcolsep) * \real{0.1250}}@{}}
\toprule\noalign{}
\begin{minipage}[b]{\linewidth}\raggedright
RUNSPEC
\end{minipage} & \begin{minipage}[b]{\linewidth}\raggedright
GRID
\end{minipage} & \begin{minipage}[b]{\linewidth}\raggedright
EDIT
\end{minipage} & \begin{minipage}[b]{\linewidth}\raggedright
PROPS
\end{minipage} & \begin{minipage}[b]{\linewidth}\raggedright
REGIONS
\end{minipage} & \begin{minipage}[b]{\linewidth}\raggedright
SOLUTION
\end{minipage} & \begin{minipage}[b]{\linewidth}\raggedright
SUMMARY
\end{minipage} & \begin{minipage}[b]{\linewidth}\raggedright
SCHEDULE
\end{minipage} \\
\midrule\noalign{}
\endhead
\bottomrule\noalign{}
\endlastfoot
\end{longtable}

\subsubsection{Description}\label{description-81}

This keyword, LOWSALT, activates the low salt brine phase for the Brine
option and also activates the Brine option. See also the BRINE keyword
in the RUNSPEC section.

There is no data required for this keyword and there is no terminating
``/'' for this keyword.

This keyword is not supported by OPM Flow but would change the results
if supported so the simulation will be stopped.

\subsubsection{Example}\label{example-60}

\begin{verbatim}
--
--       ACTIVATE THE LOW SALT BRINE PHASE FOR THE BRINE OPTION
--
LOWSALT
\end{verbatim}

The above example declares that the low salt brine phase is active in
the model for the Brine option.

\subsection{MEMORY -- Define Allocated Memory
(Retired)}\label{memory-define-allocated-memory-retired}

\begin{longtable}[]{@{}
  >{\raggedright\arraybackslash}p{(\linewidth - 14\tabcolsep) * \real{0.1250}}
  >{\raggedright\arraybackslash}p{(\linewidth - 14\tabcolsep) * \real{0.1250}}
  >{\raggedright\arraybackslash}p{(\linewidth - 14\tabcolsep) * \real{0.1250}}
  >{\raggedright\arraybackslash}p{(\linewidth - 14\tabcolsep) * \real{0.1250}}
  >{\raggedright\arraybackslash}p{(\linewidth - 14\tabcolsep) * \real{0.1250}}
  >{\raggedright\arraybackslash}p{(\linewidth - 14\tabcolsep) * \real{0.1250}}
  >{\raggedright\arraybackslash}p{(\linewidth - 14\tabcolsep) * \real{0.1250}}
  >{\raggedright\arraybackslash}p{(\linewidth - 14\tabcolsep) * \real{0.1250}}@{}}
\toprule\noalign{}
\begin{minipage}[b]{\linewidth}\raggedright
RUNSPEC
\end{minipage} & \begin{minipage}[b]{\linewidth}\raggedright
GRID
\end{minipage} & \begin{minipage}[b]{\linewidth}\raggedright
EDIT
\end{minipage} & \begin{minipage}[b]{\linewidth}\raggedright
PROPS
\end{minipage} & \begin{minipage}[b]{\linewidth}\raggedright
REGIONS
\end{minipage} & \begin{minipage}[b]{\linewidth}\raggedright
SOLUTION
\end{minipage} & \begin{minipage}[b]{\linewidth}\raggedright
SUMMARY
\end{minipage} & \begin{minipage}[b]{\linewidth}\raggedright
SCHEDULE
\end{minipage} \\
\midrule\noalign{}
\endhead
\bottomrule\noalign{}
\endlastfoot
\end{longtable}

\subsubsection{Description}\label{description-82}

This keyword defines the memory allocation for the run.

OPM Flow and now the commercial simulator uses dynamic memory allocation
and therefore the keyword has no effect and is ignored by both
simulators.

\subsection{MESSAGE -- Output User
Message}\label{message-output-user-message-1}

The MESSAGE keyword outputs a user message to the terminal, as well as
to the print (\emph{.PRT) and debug (}.DBG) files. Note this is
different to the MESSAGES keyword, that defines OPM Flows message print
limits and stop limits generated by the simulator.

See MESSAGE -- Output User Message in the GLOBAL section for a full
description.

\subsection{MESSAGES -- Define Message Print Limits and Stop
Limits}\label{messages-define-message-print-limits-and-stop-limits-1}

The MESSAGES keyword defines the print and stops levels for various
messages. The ``print limits'' set the maximum number of messages that
will be printed, after which no more messages will be printed and the
``stop limits'' terminate the run when these limits are exceeded. There
are six levels of message that increase in severity from informative all
the way to programming errors, as outlined in Table 4.5.

See MESSAGES -- Define Message Print Limits and Stop Limits in the
GLOBAL section for a full description.

\subsection{MESSSRVC -- Activate or Deactivate Database Message File
Output}\label{messsrvc-activate-or-deactivate-database-message-file-output}

\begin{longtable}[]{@{}
  >{\raggedright\arraybackslash}p{(\linewidth - 14\tabcolsep) * \real{0.1250}}
  >{\raggedright\arraybackslash}p{(\linewidth - 14\tabcolsep) * \real{0.1250}}
  >{\raggedright\arraybackslash}p{(\linewidth - 14\tabcolsep) * \real{0.1250}}
  >{\raggedright\arraybackslash}p{(\linewidth - 14\tabcolsep) * \real{0.1250}}
  >{\raggedright\arraybackslash}p{(\linewidth - 14\tabcolsep) * \real{0.1250}}
  >{\raggedright\arraybackslash}p{(\linewidth - 14\tabcolsep) * \real{0.1250}}
  >{\raggedright\arraybackslash}p{(\linewidth - 14\tabcolsep) * \real{0.1250}}
  >{\raggedright\arraybackslash}p{(\linewidth - 14\tabcolsep) * \real{0.1250}}@{}}
\toprule\noalign{}
\begin{minipage}[b]{\linewidth}\raggedright
RUNSPEC
\end{minipage} & \begin{minipage}[b]{\linewidth}\raggedright
GRID
\end{minipage} & \begin{minipage}[b]{\linewidth}\raggedright
EDIT
\end{minipage} & \begin{minipage}[b]{\linewidth}\raggedright
PROPS
\end{minipage} & \begin{minipage}[b]{\linewidth}\raggedright
REGIONS
\end{minipage} & \begin{minipage}[b]{\linewidth}\raggedright
SOLUTION
\end{minipage} & \begin{minipage}[b]{\linewidth}\raggedright
SUMMARY
\end{minipage} & \begin{minipage}[b]{\linewidth}\raggedright
SCHEDULE
\end{minipage} \\
\midrule\noalign{}
\endhead
\bottomrule\noalign{}
\endlastfoot
\end{longtable}

\subsubsection{Description}\label{description-83}

The MESSSRVC keyword activates or deactivates output to the database
message file (*.DBPRTX). The file contains all the messages from run in
binary format and is used in some post-processing software to annotate
production line plots from the run.

This keyword is not supported by OPM Flow but would change the results
if supported so the simulation will be stopped.

\subsection{METRIC -- Activate the Metric System of Units for the
Model}\label{metric-activate-the-metric-system-of-units-for-the-model}

\begin{longtable}[]{@{}
  >{\raggedright\arraybackslash}p{(\linewidth - 14\tabcolsep) * \real{0.1250}}
  >{\raggedright\arraybackslash}p{(\linewidth - 14\tabcolsep) * \real{0.1250}}
  >{\raggedright\arraybackslash}p{(\linewidth - 14\tabcolsep) * \real{0.1250}}
  >{\raggedright\arraybackslash}p{(\linewidth - 14\tabcolsep) * \real{0.1250}}
  >{\raggedright\arraybackslash}p{(\linewidth - 14\tabcolsep) * \real{0.1250}}
  >{\raggedright\arraybackslash}p{(\linewidth - 14\tabcolsep) * \real{0.1250}}
  >{\raggedright\arraybackslash}p{(\linewidth - 14\tabcolsep) * \real{0.1250}}
  >{\raggedright\arraybackslash}p{(\linewidth - 14\tabcolsep) * \real{0.1250}}@{}}
\toprule\noalign{}
\begin{minipage}[b]{\linewidth}\raggedright
RUNSPEC
\end{minipage} & \begin{minipage}[b]{\linewidth}\raggedright
GRID
\end{minipage} & \begin{minipage}[b]{\linewidth}\raggedright
EDIT
\end{minipage} & \begin{minipage}[b]{\linewidth}\raggedright
PROPS
\end{minipage} & \begin{minipage}[b]{\linewidth}\raggedright
REGIONS
\end{minipage} & \begin{minipage}[b]{\linewidth}\raggedright
SOLUTION
\end{minipage} & \begin{minipage}[b]{\linewidth}\raggedright
SUMMARY
\end{minipage} & \begin{minipage}[b]{\linewidth}\raggedright
SCHEDULE
\end{minipage} \\
\midrule\noalign{}
\endhead
\bottomrule\noalign{}
\endlastfoot
\end{longtable}

\subsubsection{Description}\label{description-84}

This keyword switches on the METRIC system of units for the model.

OPM Flow has three sets of units, namely: METRIC, FIELD and LAB and one
of these keyword should be invoked in the RUNSPEC section to avoid any
ambiguity. Both the simulator input and output units are controlled by
including one of the METRIC, FIELD or LAB keywords in the RUNSPEC
section of the input file.

There is no data required for this keyword and there is no terminating
``/'' for this keyword.

\subsubsection{Example}\label{example-61}

\begin{verbatim}
--
--       SWITCH ON THE METRIC SYSTEM OF UNITS FOR BOTH INPUT AND OUTPUT
--
METRIC
\end{verbatim}

The above example switches on the METRIC system of units for the model.

\subsection{MICP -- Activate the Microbially Induced Calcite
Precipitation
Model}\label{micp-activate-the-microbially-induced-calcite-precipitation-model}

\begin{longtable}[]{@{}
  >{\raggedright\arraybackslash}p{(\linewidth - 14\tabcolsep) * \real{0.1250}}
  >{\raggedright\arraybackslash}p{(\linewidth - 14\tabcolsep) * \real{0.1250}}
  >{\raggedright\arraybackslash}p{(\linewidth - 14\tabcolsep) * \real{0.1250}}
  >{\raggedright\arraybackslash}p{(\linewidth - 14\tabcolsep) * \real{0.1250}}
  >{\raggedright\arraybackslash}p{(\linewidth - 14\tabcolsep) * \real{0.1250}}
  >{\raggedright\arraybackslash}p{(\linewidth - 14\tabcolsep) * \real{0.1250}}
  >{\raggedright\arraybackslash}p{(\linewidth - 14\tabcolsep) * \real{0.1250}}
  >{\raggedright\arraybackslash}p{(\linewidth - 14\tabcolsep) * \real{0.1250}}@{}}
\toprule\noalign{}
\begin{minipage}[b]{\linewidth}\raggedright
RUNSPEC
\end{minipage} & \begin{minipage}[b]{\linewidth}\raggedright
GRID
\end{minipage} & \begin{minipage}[b]{\linewidth}\raggedright
EDIT
\end{minipage} & \begin{minipage}[b]{\linewidth}\raggedright
PROPS
\end{minipage} & \begin{minipage}[b]{\linewidth}\raggedright
REGIONS
\end{minipage} & \begin{minipage}[b]{\linewidth}\raggedright
SOLUTION
\end{minipage} & \begin{minipage}[b]{\linewidth}\raggedright
SUMMARY
\end{minipage} & \begin{minipage}[b]{\linewidth}\raggedright
SCHEDULE
\end{minipage} \\
\midrule\noalign{}
\endhead
\bottomrule\noalign{}
\endlastfoot
\end{longtable}

\subsubsection{Description}\label{description-85}

The MICP keyword activates the Microbially Induced Calcite Precipitation
(``MICP'') model for the run. MICP is a new and sustainable technology
which utilizes biochemical processes to create barriers by calcium
carbonate cementation, the technology has the potential to be used for
sealing leakage zones in geological formations. Further information on
the mathematical model can be found in the open-access publications
Landa-Marbán et al {[}Landa-Marbán, D., Tveit, S., Kumar, K., Gasda,
S.E., 2021. Practical approaches to study microbially induced calcite
precipitation at the eld scale. Int. J. Greenh. Gas Control 106, 103256.
https://doi.org/10.1016/j.ijggc.2021.103256.{]} and {[}Landa-Marbán, D.,
Kumar, K., Tveit, S., Gasda, S.E., 2021. Numerical studies of CO2
leakage remediation by micp-based plugging technology. In: Røkke, N.A.
and Knuutila, H.K. (Eds) Short Papers from the 11th International
Trondheim CCS conference, ISBN: 978-82-536-1714-5, 284-290.{]}.

\begin{tcolorbox}[enhanced jigsaw, title=\textcolor{quarto-callout-note-color}{\faInfo}\hspace{0.5em}{Note}, coltitle=black, breakable, toprule=.15mm, colbacktitle=quarto-callout-note-color!10!white, leftrule=.75mm, opacitybacktitle=0.6, colback=white, bottomtitle=1mm, toptitle=1mm, titlerule=0mm, left=2mm, colframe=quarto-callout-note-color-frame, arc=.35mm, rightrule=.15mm, bottomrule=.15mm, opacityback=0]

This is an OPM Flow specific keyword used to investigate leakage
remediation. The module requires that both the MICP and WATER keywords
in the RUNSPEC to be active.

\end{tcolorbox}

There is no data required for this keyword and there is no terminating
``/'' for this keyword.

\subsubsection{Example}\label{example-62}

\begin{verbatim}
        --
--       ACTIVATE THE MICROBIAL INDUCED CALCITE PRECIPITATION MODEL
--
MICP
--
--       WATER PHASE IS PRESENT IN THE RUN
--
WATER
\end{verbatim}

The above example declares that the MICP model is active for the run and
activates the water phase for the MICP model. For a complete example of
this model, see
\href{https://github.com/OPM/opm-tests/blob/master/micp/MICP.DATA}{MICP.DATA}.

\subsection{MINNPCOL -- Define the Minimum Number of Newton Iterations
Used to Update Well
Targets}\label{minnpcol-define-the-minimum-number-of-newton-iterations-used-to-update-well-targets}

\begin{longtable}[]{@{}
  >{\raggedright\arraybackslash}p{(\linewidth - 14\tabcolsep) * \real{0.1250}}
  >{\raggedright\arraybackslash}p{(\linewidth - 14\tabcolsep) * \real{0.1250}}
  >{\raggedright\arraybackslash}p{(\linewidth - 14\tabcolsep) * \real{0.1250}}
  >{\raggedright\arraybackslash}p{(\linewidth - 14\tabcolsep) * \real{0.1250}}
  >{\raggedright\arraybackslash}p{(\linewidth - 14\tabcolsep) * \real{0.1250}}
  >{\raggedright\arraybackslash}p{(\linewidth - 14\tabcolsep) * \real{0.1250}}
  >{\raggedright\arraybackslash}p{(\linewidth - 14\tabcolsep) * \real{0.1250}}
  >{\raggedright\arraybackslash}p{(\linewidth - 14\tabcolsep) * \real{0.1250}}@{}}
\toprule\noalign{}
\begin{minipage}[b]{\linewidth}\raggedright
RUNSPEC
\end{minipage} & \begin{minipage}[b]{\linewidth}\raggedright
GRID
\end{minipage} & \begin{minipage}[b]{\linewidth}\raggedright
EDIT
\end{minipage} & \begin{minipage}[b]{\linewidth}\raggedright
PROPS
\end{minipage} & \begin{minipage}[b]{\linewidth}\raggedright
REGIONS
\end{minipage} & \begin{minipage}[b]{\linewidth}\raggedright
SOLUTION
\end{minipage} & \begin{minipage}[b]{\linewidth}\raggedright
SUMMARY
\end{minipage} & \begin{minipage}[b]{\linewidth}\raggedright
SCHEDULE
\end{minipage} \\
\midrule\noalign{}
\endhead
\bottomrule\noalign{}
\endlastfoot
\end{longtable}

\subsubsection{Description}\label{description-86}

The MINNPCOL keyword defines the minimum number of Newton iterations
within a time step where the well production and injection targets may
be updated, after which the well targets will be frozen until the time
step calculations have converged and the time step is complete.

\begin{tcolorbox}[enhanced jigsaw, title=\textcolor{quarto-callout-note-color}{\faInfo}\hspace{0.5em}{Note}, coltitle=black, breakable, toprule=.15mm, colbacktitle=quarto-callout-note-color!10!white, leftrule=.75mm, opacitybacktitle=0.6, colback=white, bottomtitle=1mm, toptitle=1mm, titlerule=0mm, left=2mm, colframe=quarto-callout-note-color-frame, arc=.35mm, rightrule=.15mm, bottomrule=.15mm, opacityback=0]

This is an OPM Flow specific keyword that sets the minimum number of
Newton iterations, as opposed to the NUPCOL keyword that defines the
maximum number of Newton iterations within a time step, after which well
targets are frozen.

\end{tcolorbox}

The MINNPCOL keyword has been deprecated and the NUPCOL keyword should
be used instead.

Wells under group control may suffer from some dependency with other
wells in the same group that are also under group control. This may
cause oscillations in the production and injection well rates within the
group. In order to avoid this, the NUPCOL keyword in the RUNSPEC section
can be used to set the maximum number of Newton iterations within a time
step, after which the group well rates are frozen until the time step
has converged. Reducing the potential for well rate oscillations within
the time step may result in the group targets and limits not being
exactly met in this case. Increasing the value of NUPCOL, will improve
the accuracy of the group targets and limits at the expense of
computational efficiency. Here, the MINNPCOL keyword sets the minimum
number of Newton iterations within a timestep where the well production
and injection targets may be updated.

See also section 2.2 Running OPM Flow 2023-04 From The Command Line on
how to set various other numerical control parameters for OPM Flow.

\begin{longtable}[]{@{}
  >{\raggedright\arraybackslash}p{(\linewidth - 6\tabcolsep) * \real{0.2500}}
  >{\raggedright\arraybackslash}p{(\linewidth - 6\tabcolsep) * \real{0.2500}}
  >{\raggedright\arraybackslash}p{(\linewidth - 6\tabcolsep) * \real{0.2500}}
  >{\raggedright\arraybackslash}p{(\linewidth - 6\tabcolsep) * \real{0.2500}}@{}}
\toprule\noalign{}
\begin{minipage}[b]{\linewidth}\raggedright
No.
\end{minipage} & \begin{minipage}[b]{\linewidth}\raggedright
Name
\end{minipage} & \begin{minipage}[b]{\linewidth}\raggedright
Description
\end{minipage} & \begin{minipage}[b]{\linewidth}\raggedright
Default
\end{minipage} \\
\midrule\noalign{}
\endhead
\bottomrule\noalign{}
\endlastfoot
1 & MINNPCOL & A positive integer that defines the minimum number of
Newton iterations within a timestep where well targets may be updated. &
3 \\
Notes: & & & \\
\end{longtable}

\emph{Table 5.20: MINNPCOL Keyword Description}

\subsubsection{Example}\label{example-63}

\begin{verbatim}
--
--       DEFINE THE MIN NUMBER OF ITERATIONS TO UPDATE WELL FLOW TARGETS
--
MINNPCOL
         3                                                                    /
\end{verbatim}

The above example sets the MINNPCOL value to its default value of three.

\subsection{MISCIBLE -- Define Miscibility Todd-Longstaff
Parameters}\label{miscible-define-miscibility-todd-longstaff-parameters}

\begin{longtable}[]{@{}
  >{\raggedright\arraybackslash}p{(\linewidth - 14\tabcolsep) * \real{0.1250}}
  >{\raggedright\arraybackslash}p{(\linewidth - 14\tabcolsep) * \real{0.1250}}
  >{\raggedright\arraybackslash}p{(\linewidth - 14\tabcolsep) * \real{0.1250}}
  >{\raggedright\arraybackslash}p{(\linewidth - 14\tabcolsep) * \real{0.1250}}
  >{\raggedright\arraybackslash}p{(\linewidth - 14\tabcolsep) * \real{0.1250}}
  >{\raggedright\arraybackslash}p{(\linewidth - 14\tabcolsep) * \real{0.1250}}
  >{\raggedright\arraybackslash}p{(\linewidth - 14\tabcolsep) * \real{0.1250}}
  >{\raggedright\arraybackslash}p{(\linewidth - 14\tabcolsep) * \real{0.1250}}@{}}
\toprule\noalign{}
\begin{minipage}[b]{\linewidth}\raggedright
RUNSPEC
\end{minipage} & \begin{minipage}[b]{\linewidth}\raggedright
GRID
\end{minipage} & \begin{minipage}[b]{\linewidth}\raggedright
EDIT
\end{minipage} & \begin{minipage}[b]{\linewidth}\raggedright
PROPS
\end{minipage} & \begin{minipage}[b]{\linewidth}\raggedright
REGIONS
\end{minipage} & \begin{minipage}[b]{\linewidth}\raggedright
SOLUTION
\end{minipage} & \begin{minipage}[b]{\linewidth}\raggedright
SUMMARY
\end{minipage} & \begin{minipage}[b]{\linewidth}\raggedright
SCHEDULE
\end{minipage} \\
\midrule\noalign{}
\endhead
\bottomrule\noalign{}
\endlastfoot
\end{longtable}

\subsubsection{Description}\label{description-87}

The MISCIBLE keyword defines the options associated with the
Todd-Longstaff {[}M. R. Todd and W. J Longstaff, The Development,
Testing, and Application Of a Numerical Simulator for Predicting
Miscible Flood Performance''. In: J. Petrol. Tech. 24.7 (1972), pages
874-882.{]} mixing parameters used for when polymer flooding or CO2 EOR
simulation cases are being run.

\begin{longtable}[]{@{}
  >{\raggedright\arraybackslash}p{(\linewidth - 6\tabcolsep) * \real{0.2500}}
  >{\raggedright\arraybackslash}p{(\linewidth - 6\tabcolsep) * \real{0.2500}}
  >{\raggedright\arraybackslash}p{(\linewidth - 6\tabcolsep) * \real{0.2500}}
  >{\raggedright\arraybackslash}p{(\linewidth - 6\tabcolsep) * \real{0.2500}}@{}}
\toprule\noalign{}
\begin{minipage}[b]{\linewidth}\raggedright
No.
\end{minipage} & \begin{minipage}[b]{\linewidth}\raggedright
Name
\end{minipage} & \begin{minipage}[b]{\linewidth}\raggedright
Description
\end{minipage} & \begin{minipage}[b]{\linewidth}\raggedright
Default
\end{minipage} \\
\midrule\noalign{}
\endhead
\bottomrule\noalign{}
\endlastfoot
1 & NTMISC & A positive integer value that declares the number miscible
residual oil saturations versus water saturations tables for SORWMIS
keyword and the number Todd-Longstaff mixing parameters entries on the
TLMIXPAR keyword. & 1 \\
2 & NSMISC & A positive integer value that sets the maximum number of
entries (or rows) for each SORWMIS table defined by the SORWMIS keyword.
& 20 \\
3 & MISOPT & A character string that defines the numerical dispersion
control options for the oil and gas relative permeability curves, set to
either NONE or TWOPOINT: Only the default value of NONE is supported. &
NONE \\
Notes: & & & \\
\end{longtable}

\emph{Table 5.21: MISCIBLE Keyword Description}

\subsubsection{Example}\label{example-64}

\begin{verbatim}
--
--       NTAB    MAX     UPSTRM
--       NTMISC  NSMISC  MISOPT
MISCIBLE
         1       20      NONE                                                  /
\end{verbatim}

The above example defines the default values for the MISCIBLE keyword,
that is one table with a maximum of 20 rows per table using the standard
one cell upstream option.

\subsection{MONITOR -- Activate Output of the Monitoring Data and
File}\label{monitor-activate-output-of-the-monitoring-data-and-file}

\begin{longtable}[]{@{}
  >{\raggedright\arraybackslash}p{(\linewidth - 14\tabcolsep) * \real{0.1250}}
  >{\raggedright\arraybackslash}p{(\linewidth - 14\tabcolsep) * \real{0.1250}}
  >{\raggedright\arraybackslash}p{(\linewidth - 14\tabcolsep) * \real{0.1250}}
  >{\raggedright\arraybackslash}p{(\linewidth - 14\tabcolsep) * \real{0.1250}}
  >{\raggedright\arraybackslash}p{(\linewidth - 14\tabcolsep) * \real{0.1250}}
  >{\raggedright\arraybackslash}p{(\linewidth - 14\tabcolsep) * \real{0.1250}}
  >{\raggedright\arraybackslash}p{(\linewidth - 14\tabcolsep) * \real{0.1250}}
  >{\raggedright\arraybackslash}p{(\linewidth - 14\tabcolsep) * \real{0.1250}}@{}}
\toprule\noalign{}
\begin{minipage}[b]{\linewidth}\raggedright
RUNSPEC
\end{minipage} & \begin{minipage}[b]{\linewidth}\raggedright
GRID
\end{minipage} & \begin{minipage}[b]{\linewidth}\raggedright
EDIT
\end{minipage} & \begin{minipage}[b]{\linewidth}\raggedright
PROPS
\end{minipage} & \begin{minipage}[b]{\linewidth}\raggedright
REGIONS
\end{minipage} & \begin{minipage}[b]{\linewidth}\raggedright
SOLUTION
\end{minipage} & \begin{minipage}[b]{\linewidth}\raggedright
SUMMARY
\end{minipage} & \begin{minipage}[b]{\linewidth}\raggedright
SCHEDULE
\end{minipage} \\
\midrule\noalign{}
\endhead
\bottomrule\noalign{}
\endlastfoot
\end{longtable}

\subsubsection{Description}\label{description-88}

The MONITOR keyword activates the writing out of the run time monitoring
information used by post-processing graphics software to display run
time information, for example the simulated production and injection
rates and cumulative values. OPM Flow does not have this functionality.

Hence, OPM Flow ignores this keyword but it is documented here for
completeness.

There is no data required for this keyword and there is no terminating
``/'' for this keyword.

\subsubsection{Example}\label{example-65}

\begin{verbatim}
--
--       ACTIVATE MONITORING OUTPUT DATA AND FILES
--
MONITOR
\end{verbatim}

The above example switches on the output required for run time
monitoring required by post-processing graphics software to review the
simulation results in real time as the run progresses; however, this has
no effect in OPM Flow input decks.

\subsection{MSGFILE -- Activate or Deactivate Message File
Output}\label{msgfile-activate-or-deactivate-message-file-output}

\begin{longtable}[]{@{}
  >{\raggedright\arraybackslash}p{(\linewidth - 14\tabcolsep) * \real{0.1250}}
  >{\raggedright\arraybackslash}p{(\linewidth - 14\tabcolsep) * \real{0.1250}}
  >{\raggedright\arraybackslash}p{(\linewidth - 14\tabcolsep) * \real{0.1250}}
  >{\raggedright\arraybackslash}p{(\linewidth - 14\tabcolsep) * \real{0.1250}}
  >{\raggedright\arraybackslash}p{(\linewidth - 14\tabcolsep) * \real{0.1250}}
  >{\raggedright\arraybackslash}p{(\linewidth - 14\tabcolsep) * \real{0.1250}}
  >{\raggedright\arraybackslash}p{(\linewidth - 14\tabcolsep) * \real{0.1250}}
  >{\raggedright\arraybackslash}p{(\linewidth - 14\tabcolsep) * \real{0.1250}}@{}}
\toprule\noalign{}
\begin{minipage}[b]{\linewidth}\raggedright
RUNSPEC
\end{minipage} & \begin{minipage}[b]{\linewidth}\raggedright
GRID
\end{minipage} & \begin{minipage}[b]{\linewidth}\raggedright
EDIT
\end{minipage} & \begin{minipage}[b]{\linewidth}\raggedright
PROPS
\end{minipage} & \begin{minipage}[b]{\linewidth}\raggedright
REGIONS
\end{minipage} & \begin{minipage}[b]{\linewidth}\raggedright
SOLUTION
\end{minipage} & \begin{minipage}[b]{\linewidth}\raggedright
SUMMARY
\end{minipage} & \begin{minipage}[b]{\linewidth}\raggedright
SCHEDULE
\end{minipage} \\
\midrule\noalign{}
\endhead
\bottomrule\noalign{}
\endlastfoot
\end{longtable}

\subsubsection{Description}\label{description-89}

MSGFILE keyword activates or deactivates the message file output used by
pre- and post-processing software. Note that message file processing is
not available in OPM Flow.

This keyword is not supported by OPM Flow but has no effect on the
results so it will be ignored.

\begin{longtable}[]{@{}
  >{\raggedright\arraybackslash}p{(\linewidth - 6\tabcolsep) * \real{0.2500}}
  >{\raggedright\arraybackslash}p{(\linewidth - 6\tabcolsep) * \real{0.2500}}
  >{\raggedright\arraybackslash}p{(\linewidth - 6\tabcolsep) * \real{0.2500}}
  >{\raggedright\arraybackslash}p{(\linewidth - 6\tabcolsep) * \real{0.2500}}@{}}
\toprule\noalign{}
\begin{minipage}[b]{\linewidth}\raggedright
No.
\end{minipage} & \begin{minipage}[b]{\linewidth}\raggedright
Name
\end{minipage} & \begin{minipage}[b]{\linewidth}\raggedright
Description
\end{minipage} & \begin{minipage}[b]{\linewidth}\raggedright
Default
\end{minipage} \\
\midrule\noalign{}
\endhead
\bottomrule\noalign{}
\endlastfoot
1 & MSGOPT & A positive integer set to 0 for to deactivate message file
output or 1 to activate message file output. & 1 \\
Notes: & & & \\
\end{longtable}

\emph{Table 5.22: MSGFILE Keyword Description}

\subsubsection{Example}\label{example-66}

\begin{verbatim}
--
--       OUTPUT
--       OPTN
MSGFILE
         0                                                                     /
\end{verbatim}

The above example deactivates the message file output, but the keyword
is ignored by OPM Flow.

\subsection{MULTIN -- Activate the Non-Unified Multiple Input File
Option}\label{multin-activate-the-non-unified-multiple-input-file-option}

\begin{longtable}[]{@{}
  >{\raggedright\arraybackslash}p{(\linewidth - 14\tabcolsep) * \real{0.1250}}
  >{\raggedright\arraybackslash}p{(\linewidth - 14\tabcolsep) * \real{0.1250}}
  >{\raggedright\arraybackslash}p{(\linewidth - 14\tabcolsep) * \real{0.1250}}
  >{\raggedright\arraybackslash}p{(\linewidth - 14\tabcolsep) * \real{0.1250}}
  >{\raggedright\arraybackslash}p{(\linewidth - 14\tabcolsep) * \real{0.1250}}
  >{\raggedright\arraybackslash}p{(\linewidth - 14\tabcolsep) * \real{0.1250}}
  >{\raggedright\arraybackslash}p{(\linewidth - 14\tabcolsep) * \real{0.1250}}
  >{\raggedright\arraybackslash}p{(\linewidth - 14\tabcolsep) * \real{0.1250}}@{}}
\toprule\noalign{}
\begin{minipage}[b]{\linewidth}\raggedright
RUNSPEC
\end{minipage} & \begin{minipage}[b]{\linewidth}\raggedright
GRID
\end{minipage} & \begin{minipage}[b]{\linewidth}\raggedright
EDIT
\end{minipage} & \begin{minipage}[b]{\linewidth}\raggedright
PROPS
\end{minipage} & \begin{minipage}[b]{\linewidth}\raggedright
REGIONS
\end{minipage} & \begin{minipage}[b]{\linewidth}\raggedright
SOLUTION
\end{minipage} & \begin{minipage}[b]{\linewidth}\raggedright
SUMMARY
\end{minipage} & \begin{minipage}[b]{\linewidth}\raggedright
SCHEDULE
\end{minipage} \\
\midrule\noalign{}
\endhead
\bottomrule\noalign{}
\endlastfoot
\end{longtable}

\subsubsection{Description}\label{description-90}

This keyword switches on the Multiple Input Files option for all input
files. Similar to the commercial simulator, OPM Flow has various options
for reading various input files and writing the resulting OPM Flow
output files as described in Table 5.23.

\begin{longtable}[]{@{}
  >{\raggedright\arraybackslash}p{(\linewidth - 6\tabcolsep) * \real{0.2500}}
  >{\raggedright\arraybackslash}p{(\linewidth - 6\tabcolsep) * \real{0.2500}}
  >{\raggedright\arraybackslash}p{(\linewidth - 6\tabcolsep) * \real{0.2500}}
  >{\raggedright\arraybackslash}p{(\linewidth - 6\tabcolsep) * \real{0.2500}}@{}}
\toprule\noalign{}
\begin{minipage}[b]{\linewidth}\raggedright
Process
\end{minipage} & \begin{minipage}[b]{\linewidth}\raggedright
Keyword
\end{minipage} & \begin{minipage}[b]{\linewidth}\raggedright
Description
\end{minipage} & \begin{minipage}[b]{\linewidth}\raggedright
Files
\end{minipage} \\
\midrule\noalign{}
\endhead
\bottomrule\noalign{}
\endlastfoot
Input & FMTIN & A character string that defines the input files to be
formatted as ASCII i.e.~text files, as oppose to binary files. The input
deck file is always of this type. The option relates to the OPM Flow
derived files that used as input, for for example when restarting from
another case. If the keyword is omitted then the default is for binary
file input. & \emph{.FEGRID }.FINSPEC \emph{.FINIT }.FRSSPEC
\emph{.FUNRST }.FSMSPEC *.FUNSMRY \\
MULTIN & A character string that defines the input files to be
non-unified multiple files, as opposed to unified files. In this case,
one file is read in per reporting time step, as opposed to all time
steps reports being read from one file. If the keyword is omitted then
the default is for one file per reporting time step. & \emph{.RSSPEC
}.X0001 \emph{.SMSPEC }.S0001 & \\
UNIFIN & A character string that defines the input files to be unified
files, as opposed to non-unified multiple files. A unified file is a
single file containing output for each reporting time step. For this
option a single summary file and a single restart file will be read. If
the keyword is omitted then the default is for one file per reporting
time step. & \emph{.RSSPEC }.UNRST \emph{.SMSPEC }.UNSMRY & \\
Output & FMTOUT & A character string that sets all output files to be
formatted as ASCII i.e.~text files, as opposed to binary files. The
\emph{.PRT, }.LOG and *.DBG files are always of this type. The option
relates to the OPM Flow output files only. In this case the files will
be portable across operating systems, but will also be very large in
terms of hard disk space. For this reason it is recommend that the
default option is used so that binary files are outputted. If the
keyword is omitted then the default is for binary file input. &
\emph{.FEGRID }.FINSPEC \emph{.FINIT }.FRSSPEC \emph{.FUNRST }.FSMSPEC
*.FUNSMRY \\
MULTOUT & A character string that defines the output files to be
non-unified multiple files, as opposed to unified files. In this case,
one file is written for each reporting time step, as opposed to all time
steps reports being written in one file. If the keyword is omitted then
the default is for one file per reporting time step. & \emph{.RSSPEC
}.X0001 \emph{.SMSPEC }.S0001 & \\
UNIFOUT & A character string that defines the output files to be unified
files, as opposed to non-unified multiple files. A unified file is a
single file containing output for each reporting time step. Here a
single summary file and a single restart file will be generated, as
opposed to one file per report time step. If the keyword is omitted then
the default is for one file per reporting time step. & \emph{.RSSPEC
}.UNRST \emph{.SMSPEC }.UNSMRY & \\
Notes: & & & \\
\end{longtable}

\emph{Table 5.23: MULTIN Keyword Description}

There is no data required for this keyword.

See also OPM FLOW OUTPUT FILE FORMATS for a more detailed description of
the various file types (ASCII or binary) and file structure formats
(unified or non-unified formats).

\subsubsection{Example}\label{example-67}

\begin{verbatim}
--
--       ACTIVATE THE MULTIPLE INPUT FILES OPTION
--
MULTIN
\end{verbatim}

The above example switches on the multiple input file option.

\subsection{MULTOUT -- Activate the Non-Unified Multiple Output File
Option}\label{multout-activate-the-non-unified-multiple-output-file-option}

\begin{longtable}[]{@{}
  >{\raggedright\arraybackslash}p{(\linewidth - 14\tabcolsep) * \real{0.1250}}
  >{\raggedright\arraybackslash}p{(\linewidth - 14\tabcolsep) * \real{0.1250}}
  >{\raggedright\arraybackslash}p{(\linewidth - 14\tabcolsep) * \real{0.1250}}
  >{\raggedright\arraybackslash}p{(\linewidth - 14\tabcolsep) * \real{0.1250}}
  >{\raggedright\arraybackslash}p{(\linewidth - 14\tabcolsep) * \real{0.1250}}
  >{\raggedright\arraybackslash}p{(\linewidth - 14\tabcolsep) * \real{0.1250}}
  >{\raggedright\arraybackslash}p{(\linewidth - 14\tabcolsep) * \real{0.1250}}
  >{\raggedright\arraybackslash}p{(\linewidth - 14\tabcolsep) * \real{0.1250}}@{}}
\toprule\noalign{}
\begin{minipage}[b]{\linewidth}\raggedright
RUNSPEC
\end{minipage} & \begin{minipage}[b]{\linewidth}\raggedright
GRID
\end{minipage} & \begin{minipage}[b]{\linewidth}\raggedright
EDIT
\end{minipage} & \begin{minipage}[b]{\linewidth}\raggedright
PROPS
\end{minipage} & \begin{minipage}[b]{\linewidth}\raggedright
REGIONS
\end{minipage} & \begin{minipage}[b]{\linewidth}\raggedright
SOLUTION
\end{minipage} & \begin{minipage}[b]{\linewidth}\raggedright
SUMMARY
\end{minipage} & \begin{minipage}[b]{\linewidth}\raggedright
SCHEDULE
\end{minipage} \\
\midrule\noalign{}
\endhead
\bottomrule\noalign{}
\endlastfoot
\end{longtable}

\subsubsection{Description}\label{description-91}

This keyword switches on the Multiple Output Files option for all output
files. Similar to the commercial simulator, OPM Flow has various options
for reading various input files and writing the resulting OPM Flow
output files as described in Table 5.24.

\begin{longtable}[]{@{}
  >{\raggedright\arraybackslash}p{(\linewidth - 6\tabcolsep) * \real{0.2500}}
  >{\raggedright\arraybackslash}p{(\linewidth - 6\tabcolsep) * \real{0.2500}}
  >{\raggedright\arraybackslash}p{(\linewidth - 6\tabcolsep) * \real{0.2500}}
  >{\raggedright\arraybackslash}p{(\linewidth - 6\tabcolsep) * \real{0.2500}}@{}}
\toprule\noalign{}
\begin{minipage}[b]{\linewidth}\raggedright
Process
\end{minipage} & \begin{minipage}[b]{\linewidth}\raggedright
Keyword
\end{minipage} & \begin{minipage}[b]{\linewidth}\raggedright
Description
\end{minipage} & \begin{minipage}[b]{\linewidth}\raggedright
Files
\end{minipage} \\
\midrule\noalign{}
\endhead
\bottomrule\noalign{}
\endlastfoot
Input & FMTIN & A character string that defines the input files to be
formatted as ASCII i.e.~text files, as oppose to binary files. The input
deck file is always of this type. The option relates to the OPM Flow
derived files that used as input, for for example when restarting from
another case. If the keyword is omitted then the default is for binary
file input. & \emph{.FEGRID }.FINSPEC \emph{.FINIT }.FRSSPEC
\emph{.FUNRST }.FSMSPEC *.FUNSMRY \\
MULTIN & A character string that defines the input files to be
non-unified multiple files, as opposed to unified files. In this case,
one file is read in per reporting time step, as opposed to all time
steps reports being read from one file. If the keyword is omitted then
the default is for one file per reporting time step. & \emph{.RSSPEC
}.X0001 \emph{.SMSPEC }.S0001 & \\
UNIFIN & A character string that defines the input files to be unified
files, as opposed to non-unified multiple files. A unified file is a
single file containing output for each reporting time step. For this
option a single summary file and a single restart file will be read. If
the keyword is omitted then the default is for one file per reporting
time step. & \emph{.RSSPEC }.UNRST \emph{.SMSPEC }.UNSMRY & \\
Output & FMTOUT & A character string that sets all output files to be
formatted as ASCII i.e.~text files, as opposed to binary files. The
\emph{.PRT, }.LOG and *.DBG files are always of this type. The option
relates to the OPM Flow output files only. In this case the files will
be portable across operating systems, but will also be very large in
terms of hard disk space. For this reason it is recommend that the
default option is used so that binary files are outputted. If the
keyword is omitted then the default is for binary file input. &
\emph{.FEGRID }.FINSPEC \emph{.FINIT }.FRSSPEC \emph{.FUNRST }.FSMSPEC
*.FUNSMRY \\
MULTOUT & A character string that defines the output files to be
non-unified multiple files, as opposed to unified files. In this case,
one file is written for each reporting time step, as opposed to all time
steps reports being written in one file. If the keyword is omitted then
the default is for one file per reporting time step. & \emph{.RSSPEC
}.X0001 \emph{.SMSPEC }.S0001 & \\
UNIFOUT & A character string that defines the output files to be unified
files, as opposed to non-unified multiple files. A unified file is a
single file containing output for each reporting time step. Here a
single summary file and a single restart file will be generated, as
opposed to one file per report time step. If the keyword is omitted then
the default is for one file per reporting time step. & \emph{.RSSPEC
}.UNRST \emph{.SMSPEC }.UNSMRY & \\
Notes: & & & \\
\end{longtable}

\emph{Table 5.24: MULTOUT Keyword Description}

There is no data required for this keyword.

See also OPM FLOW OUTPUT FILE FORMATS for a more detailed description of
the various file types (ASCII or binary) and file structure formats
(unified or non-unified formats).

\subsubsection{Example}\label{example-68}

\begin{verbatim}
--
--       ACTIVATE THE MULTIPLE OUTPUT FILES OPTION
--
MULTOUT
\end{verbatim}

The above example switches on the multiple output file option.

\subsection{MULTOUTS -- Activate Non-Unified Multiple Summary Output
File
Option}\label{multouts-activate-non-unified-multiple-summary-output-file-option}

\begin{longtable}[]{@{}
  >{\raggedright\arraybackslash}p{(\linewidth - 14\tabcolsep) * \real{0.1250}}
  >{\raggedright\arraybackslash}p{(\linewidth - 14\tabcolsep) * \real{0.1250}}
  >{\raggedright\arraybackslash}p{(\linewidth - 14\tabcolsep) * \real{0.1250}}
  >{\raggedright\arraybackslash}p{(\linewidth - 14\tabcolsep) * \real{0.1250}}
  >{\raggedright\arraybackslash}p{(\linewidth - 14\tabcolsep) * \real{0.1250}}
  >{\raggedright\arraybackslash}p{(\linewidth - 14\tabcolsep) * \real{0.1250}}
  >{\raggedright\arraybackslash}p{(\linewidth - 14\tabcolsep) * \real{0.1250}}
  >{\raggedright\arraybackslash}p{(\linewidth - 14\tabcolsep) * \real{0.1250}}@{}}
\toprule\noalign{}
\begin{minipage}[b]{\linewidth}\raggedright
RUNSPEC
\end{minipage} & \begin{minipage}[b]{\linewidth}\raggedright
GRID
\end{minipage} & \begin{minipage}[b]{\linewidth}\raggedright
EDIT
\end{minipage} & \begin{minipage}[b]{\linewidth}\raggedright
PROPS
\end{minipage} & \begin{minipage}[b]{\linewidth}\raggedright
REGIONS
\end{minipage} & \begin{minipage}[b]{\linewidth}\raggedright
SOLUTION
\end{minipage} & \begin{minipage}[b]{\linewidth}\raggedright
SUMMARY
\end{minipage} & \begin{minipage}[b]{\linewidth}\raggedright
SCHEDULE
\end{minipage} \\
\midrule\noalign{}
\endhead
\bottomrule\noalign{}
\endlastfoot
\end{longtable}

\subsubsection{Description}\label{description-92}

This keyword switches on the Multiple Output Files option for SUMMARY
files only, and overwrites the UNIFOUT keyword in the RUNSPEC section
that activates the Unified Output Files option for all output files.

There is no data required for this keyword and there is no terminating
``/'' for this keyword.

This keyword is not supported by OPM Flow but would change the results
if supported so the simulation will be stopped.

See also OPM FLOW OUTPUT FILE FORMATS for a more detailed description of
the various file types (ASCII or binary) and file structure formats
(unified or non-unified formats).

\subsubsection{Example}\label{example-69}

\begin{verbatim}
--
--       ACTIVATE MULTIPLE OUTPUT SUMMARY FILES ONLY OPTION
--
MULTOUTS
\end{verbatim}

The above example switches on the multiple output file option.

\subsection{MULTREAL -- Activate Commercial Simulator's
Multi-Realization
License}\label{multreal-activate-commercial-simulators-multi-realization-license}

\begin{longtable}[]{@{}
  >{\raggedright\arraybackslash}p{(\linewidth - 14\tabcolsep) * \real{0.1250}}
  >{\raggedright\arraybackslash}p{(\linewidth - 14\tabcolsep) * \real{0.1250}}
  >{\raggedright\arraybackslash}p{(\linewidth - 14\tabcolsep) * \real{0.1250}}
  >{\raggedright\arraybackslash}p{(\linewidth - 14\tabcolsep) * \real{0.1250}}
  >{\raggedright\arraybackslash}p{(\linewidth - 14\tabcolsep) * \real{0.1250}}
  >{\raggedright\arraybackslash}p{(\linewidth - 14\tabcolsep) * \real{0.1250}}
  >{\raggedright\arraybackslash}p{(\linewidth - 14\tabcolsep) * \real{0.1250}}
  >{\raggedright\arraybackslash}p{(\linewidth - 14\tabcolsep) * \real{0.1250}}@{}}
\toprule\noalign{}
\begin{minipage}[b]{\linewidth}\raggedright
RUNSPEC
\end{minipage} & \begin{minipage}[b]{\linewidth}\raggedright
GRID
\end{minipage} & \begin{minipage}[b]{\linewidth}\raggedright
EDIT
\end{minipage} & \begin{minipage}[b]{\linewidth}\raggedright
PROPS
\end{minipage} & \begin{minipage}[b]{\linewidth}\raggedright
REGIONS
\end{minipage} & \begin{minipage}[b]{\linewidth}\raggedright
SOLUTION
\end{minipage} & \begin{minipage}[b]{\linewidth}\raggedright
SUMMARY
\end{minipage} & \begin{minipage}[b]{\linewidth}\raggedright
SCHEDULE
\end{minipage} \\
\midrule\noalign{}
\endhead
\bottomrule\noalign{}
\endlastfoot
\end{longtable}

\subsubsection{Description}\label{description-93}

The MULTREAL keyword activates the commercial simulator's
Multi-Realization License option.

OPM Flow is an open source project and therefore there is no license
management of the various implemented options; hence this keyword is not
supported by OPM Flow but would change the results if supported so the
simulation will be stopped.

\subsection{NETWORK -- Activate the Extended Network Option and Define
Parameters}\label{network-activate-the-extended-network-option-and-define-parameters}

\begin{longtable}[]{@{}
  >{\raggedright\arraybackslash}p{(\linewidth - 14\tabcolsep) * \real{0.1250}}
  >{\raggedright\arraybackslash}p{(\linewidth - 14\tabcolsep) * \real{0.1250}}
  >{\raggedright\arraybackslash}p{(\linewidth - 14\tabcolsep) * \real{0.1250}}
  >{\raggedright\arraybackslash}p{(\linewidth - 14\tabcolsep) * \real{0.1250}}
  >{\raggedright\arraybackslash}p{(\linewidth - 14\tabcolsep) * \real{0.1250}}
  >{\raggedright\arraybackslash}p{(\linewidth - 14\tabcolsep) * \real{0.1250}}
  >{\raggedright\arraybackslash}p{(\linewidth - 14\tabcolsep) * \real{0.1250}}
  >{\raggedright\arraybackslash}p{(\linewidth - 14\tabcolsep) * \real{0.1250}}@{}}
\toprule\noalign{}
\begin{minipage}[b]{\linewidth}\raggedright
RUNSPEC
\end{minipage} & \begin{minipage}[b]{\linewidth}\raggedright
GRID
\end{minipage} & \begin{minipage}[b]{\linewidth}\raggedright
EDIT
\end{minipage} & \begin{minipage}[b]{\linewidth}\raggedright
PROPS
\end{minipage} & \begin{minipage}[b]{\linewidth}\raggedright
REGIONS
\end{minipage} & \begin{minipage}[b]{\linewidth}\raggedright
SOLUTION
\end{minipage} & \begin{minipage}[b]{\linewidth}\raggedright
SUMMARY
\end{minipage} & \begin{minipage}[b]{\linewidth}\raggedright
SCHEDULE
\end{minipage} \\
\midrule\noalign{}
\endhead
\bottomrule\noalign{}
\endlastfoot
\end{longtable}

\subsubsection{Description}\label{description-94}

This keyword activates the Extended Network option and defines the
maximum number of nodes and links (branches) in the network. The
Extended Network option is a different facility to the Standard Network
facility, as such, this keyword should only be used if the former
network is required for the run.

\begin{longtable}[]{@{}
  >{\raggedright\arraybackslash}p{(\linewidth - 6\tabcolsep) * \real{0.2500}}
  >{\raggedright\arraybackslash}p{(\linewidth - 6\tabcolsep) * \real{0.2500}}
  >{\raggedright\arraybackslash}p{(\linewidth - 6\tabcolsep) * \real{0.2500}}
  >{\raggedright\arraybackslash}p{(\linewidth - 6\tabcolsep) * \real{0.2500}}@{}}
\toprule\noalign{}
\begin{minipage}[b]{\linewidth}\raggedright
No.
\end{minipage} & \begin{minipage}[b]{\linewidth}\raggedright
Name
\end{minipage} & \begin{minipage}[b]{\linewidth}\raggedright
Description
\end{minipage} & \begin{minipage}[b]{\linewidth}\raggedright
Default
\end{minipage} \\
\midrule\noalign{}
\endhead
\bottomrule\noalign{}
\endlastfoot
1 & NODMAX & NODMAX is a positive integer that defines the maximum
number of nodes in the Extended Network model. & None \\
2 & NBRMAX & NBRMAX is a positive integer that defines the maximum
number of links in the Extended Network model. & None \\
3 & NBCMAX & NBCMAX is a positive integer that defines the maximum
number of branches that can be connected to a node in the Extended
Network model, used in the commercial compositional simulator. The
parameter is ignored by OPM Flow and should be defaulted or set to the
default value of 20. & 20 \\
Notes: & & & \\
\end{longtable}

\emph{Table 5.25: NETWORK Keyword Description}

\subsubsection{Example}\label{example-70}

\begin{verbatim}
--
--       ACTIVATE THE EXTENDED NETWORK OPTION AND DEFINE PARAMETERS
--
--       MAX.    MAX     NOT
--       NODE    LINK    USED
NETWORK
         10       12     1*                                                    /
\end{verbatim}

In the above example the maximum number of nodes is set equal to ten and
the maximum number of links (or branches) is set equal to 12, for the
Extended Network option.

\subsection{NINEPOIN -- Activate the Nine-Point Discretization
Option}\label{ninepoin-activate-the-nine-point-discretization-option}

\begin{longtable}[]{@{}
  >{\raggedright\arraybackslash}p{(\linewidth - 14\tabcolsep) * \real{0.1250}}
  >{\raggedright\arraybackslash}p{(\linewidth - 14\tabcolsep) * \real{0.1250}}
  >{\raggedright\arraybackslash}p{(\linewidth - 14\tabcolsep) * \real{0.1250}}
  >{\raggedright\arraybackslash}p{(\linewidth - 14\tabcolsep) * \real{0.1250}}
  >{\raggedright\arraybackslash}p{(\linewidth - 14\tabcolsep) * \real{0.1250}}
  >{\raggedright\arraybackslash}p{(\linewidth - 14\tabcolsep) * \real{0.1250}}
  >{\raggedright\arraybackslash}p{(\linewidth - 14\tabcolsep) * \real{0.1250}}
  >{\raggedright\arraybackslash}p{(\linewidth - 14\tabcolsep) * \real{0.1250}}@{}}
\toprule\noalign{}
\begin{minipage}[b]{\linewidth}\raggedright
RUNSPEC
\end{minipage} & \begin{minipage}[b]{\linewidth}\raggedright
GRID
\end{minipage} & \begin{minipage}[b]{\linewidth}\raggedright
EDIT
\end{minipage} & \begin{minipage}[b]{\linewidth}\raggedright
PROPS
\end{minipage} & \begin{minipage}[b]{\linewidth}\raggedright
REGIONS
\end{minipage} & \begin{minipage}[b]{\linewidth}\raggedright
SOLUTION
\end{minipage} & \begin{minipage}[b]{\linewidth}\raggedright
SUMMARY
\end{minipage} & \begin{minipage}[b]{\linewidth}\raggedright
SCHEDULE
\end{minipage} \\
\midrule\noalign{}
\endhead
\bottomrule\noalign{}
\endlastfoot
\end{longtable}

\subsubsection{Description}\label{description-95}

The NINEPOIN keyword activates the Nine-Point Discretization formulation
for the whole grid. If the keyword is absent from the run then the
conventional standard five-point discretization formulation is used for
the model. The nine-point scheme is based on adding additional
non-neighbor connections between the diagonal neighbors in the areal
plane, in order to reduce grid orientation effects {[}Yanosik, J. L. and
McCracken, T. A. ``A Nine-Point, Finite-Difference Reservoir Simulator
for Realistic Prediction of Adverse Mobility Ratio Displacements,''
paper SPE 5734, Society of Petroleum Engineers Journal (1979) 19, No.~4,
253-262.{]}.

This keyword is not supported by OPM Flow but would change the results
if supported so the simulation will be stopped.

In none Local Grid Refinement runs the NINENUM keyword in the GRID
section may be use to optionally set parts of the grid to use nine-point
discretization and the remaining regions to use the conventional
standard five-point discretization formulation.

There is no data required for this keyword and there is no terminating
``/'' for this keyword.

\subsubsection{Example}\label{example-71}

\begin{verbatim}
--
--       ACTIVATE THE NINE-POINT DISCRETIZATION OPTION
--
NINEPOIN
\end{verbatim}

The above example switches on the Nine-Point Discretization option for
the whole grid.

\subsection{NMATRIX -- Activate the Discretized Matrix Dual Porosity
Option}\label{nmatrix-activate-the-discretized-matrix-dual-porosity-option}

\begin{longtable}[]{@{}
  >{\raggedright\arraybackslash}p{(\linewidth - 14\tabcolsep) * \real{0.1250}}
  >{\raggedright\arraybackslash}p{(\linewidth - 14\tabcolsep) * \real{0.1250}}
  >{\raggedright\arraybackslash}p{(\linewidth - 14\tabcolsep) * \real{0.1250}}
  >{\raggedright\arraybackslash}p{(\linewidth - 14\tabcolsep) * \real{0.1250}}
  >{\raggedright\arraybackslash}p{(\linewidth - 14\tabcolsep) * \real{0.1250}}
  >{\raggedright\arraybackslash}p{(\linewidth - 14\tabcolsep) * \real{0.1250}}
  >{\raggedright\arraybackslash}p{(\linewidth - 14\tabcolsep) * \real{0.1250}}
  >{\raggedright\arraybackslash}p{(\linewidth - 14\tabcolsep) * \real{0.1250}}@{}}
\toprule\noalign{}
\begin{minipage}[b]{\linewidth}\raggedright
RUNSPEC
\end{minipage} & \begin{minipage}[b]{\linewidth}\raggedright
GRID
\end{minipage} & \begin{minipage}[b]{\linewidth}\raggedright
EDIT
\end{minipage} & \begin{minipage}[b]{\linewidth}\raggedright
PROPS
\end{minipage} & \begin{minipage}[b]{\linewidth}\raggedright
REGIONS
\end{minipage} & \begin{minipage}[b]{\linewidth}\raggedright
SOLUTION
\end{minipage} & \begin{minipage}[b]{\linewidth}\raggedright
SUMMARY
\end{minipage} & \begin{minipage}[b]{\linewidth}\raggedright
SCHEDULE
\end{minipage} \\
\midrule\noalign{}
\endhead
\bottomrule\noalign{}
\endlastfoot
\end{longtable}

\subsubsection{Description}\label{description-96}

The NMATRIX keyword activates the Discretized Matrix Dual Porosity
option and specifies the number of sub-grid blocks in the actual matrix
grid blocks. See also the NMATOPS keyword in the GRID section that
defines various parameters for this option.

This keyword is not supported by OPM Flow but would change the results
if supported so the simulation will be stopped.

\begin{longtable}[]{@{}
  >{\raggedright\arraybackslash}p{(\linewidth - 6\tabcolsep) * \real{0.2500}}
  >{\raggedright\arraybackslash}p{(\linewidth - 6\tabcolsep) * \real{0.2500}}
  >{\raggedright\arraybackslash}p{(\linewidth - 6\tabcolsep) * \real{0.2500}}
  >{\raggedright\arraybackslash}p{(\linewidth - 6\tabcolsep) * \real{0.2500}}@{}}
\toprule\noalign{}
\begin{minipage}[b]{\linewidth}\raggedright
No.
\end{minipage} & \begin{minipage}[b]{\linewidth}\raggedright
Name
\end{minipage} & \begin{minipage}[b]{\linewidth}\raggedright
Description
\end{minipage} & \begin{minipage}[b]{\linewidth}\raggedright
Default
\end{minipage} \\
\midrule\noalign{}
\endhead
\bottomrule\noalign{}
\endlastfoot
1 & NMATRIX & A positive integer value that specifies the number of
sub-grid blocks in the actual matrix grid blocks. & 1 \\
Notes: & & & \\
\end{longtable}

\emph{Table 5.26: NMATRIX Keyword Description}

Note the keyword cannot be used in conjunction with the TRPLPORO
keyword, which is also in the RUNSPEC section.

\subsubsection{Example}\label{example-72}

\begin{verbatim}
--
--       SUB-GRIDS
--       NMATRIX
NMATRIX
         4                                                        /
\end{verbatim}

The above example activates the Discretized Matrix Dual Porosity option
and specifies the number of sub-grid blocks in the actual matrix grid
block to be four.

\subsection{NNEWTF -- Activate the Non-Newtonian Fluid
Model}\label{nnewtf-activate-the-non-newtonian-fluid-model}

\begin{longtable}[]{@{}
  >{\raggedright\arraybackslash}p{(\linewidth - 14\tabcolsep) * \real{0.1250}}
  >{\raggedright\arraybackslash}p{(\linewidth - 14\tabcolsep) * \real{0.1250}}
  >{\raggedright\arraybackslash}p{(\linewidth - 14\tabcolsep) * \real{0.1250}}
  >{\raggedright\arraybackslash}p{(\linewidth - 14\tabcolsep) * \real{0.1250}}
  >{\raggedright\arraybackslash}p{(\linewidth - 14\tabcolsep) * \real{0.1250}}
  >{\raggedright\arraybackslash}p{(\linewidth - 14\tabcolsep) * \real{0.1250}}
  >{\raggedright\arraybackslash}p{(\linewidth - 14\tabcolsep) * \real{0.1250}}
  >{\raggedright\arraybackslash}p{(\linewidth - 14\tabcolsep) * \real{0.1250}}@{}}
\toprule\noalign{}
\begin{minipage}[b]{\linewidth}\raggedright
RUNSPEC
\end{minipage} & \begin{minipage}[b]{\linewidth}\raggedright
GRID
\end{minipage} & \begin{minipage}[b]{\linewidth}\raggedright
EDIT
\end{minipage} & \begin{minipage}[b]{\linewidth}\raggedright
PROPS
\end{minipage} & \begin{minipage}[b]{\linewidth}\raggedright
REGIONS
\end{minipage} & \begin{minipage}[b]{\linewidth}\raggedright
SOLUTION
\end{minipage} & \begin{minipage}[b]{\linewidth}\raggedright
SUMMARY
\end{minipage} & \begin{minipage}[b]{\linewidth}\raggedright
SCHEDULE
\end{minipage} \\
\midrule\noalign{}
\endhead
\bottomrule\noalign{}
\endlastfoot
\end{longtable}

\subsubsection{Description}\label{description-97}

This keyword activates the Non-Newtonian Fluid phase and model for when
the polymer phase is present in the model, as indicated by the POLYMER
keyword in the RUNSPEC section.

This keyword is not supported by OPM Flow but would change the results
if supported so the simulation will be stopped.

\begin{longtable}[]{@{}
  >{\raggedright\arraybackslash}p{(\linewidth - 6\tabcolsep) * \real{0.2500}}
  >{\raggedright\arraybackslash}p{(\linewidth - 6\tabcolsep) * \real{0.2500}}
  >{\raggedright\arraybackslash}p{(\linewidth - 6\tabcolsep) * \real{0.2500}}
  >{\raggedright\arraybackslash}p{(\linewidth - 6\tabcolsep) * \real{0.2500}}@{}}
\toprule\noalign{}
\begin{minipage}[b]{\linewidth}\raggedright
No.
\end{minipage} & \begin{minipage}[b]{\linewidth}\raggedright
Name
\end{minipage} & \begin{minipage}[b]{\linewidth}\raggedright
Description
\end{minipage} & \begin{minipage}[b]{\linewidth}\raggedright
Default
\end{minipage} \\
\midrule\noalign{}
\endhead
\bottomrule\noalign{}
\endlastfoot
1 & NTHRBL & A positive integer that defines the maximum number of
Herschel-Bulkley versus polymer concentration tables to be used with the
polymer model, as entered via the FHERCHBL keyword in the PROPS section.
The tables are allocated to different parts of the grid by the HBNUM
keyword in the REGIONS section & NTPVT \\
2 & NLNHBL & A positive integer that defines the maximum number of rows
for each table entered by the FHERCHBL keyword in the PROPS section. &
2 \\
Notes: & & & \\
\end{longtable}

\emph{Table 5.27: NNEWTF Keyword Description}

\subsubsection{Example}\label{example-73}

\begin{verbatim}
--
--       MAX     MAX
--       NTHRBL  NLNHBL
NNEWTF
          3      5                                                             /
\end{verbatim}

The above example defines maximum number of Herschel-Bulkley tables to
be three with a maximum number of rows for each table set to five.

\subsection{NOCASC -- Activate Linear Solver Tracer
Algorithm}\label{nocasc-activate-linear-solver-tracer-algorithm}

\begin{longtable}[]{@{}
  >{\raggedright\arraybackslash}p{(\linewidth - 14\tabcolsep) * \real{0.1250}}
  >{\raggedright\arraybackslash}p{(\linewidth - 14\tabcolsep) * \real{0.1250}}
  >{\raggedright\arraybackslash}p{(\linewidth - 14\tabcolsep) * \real{0.1250}}
  >{\raggedright\arraybackslash}p{(\linewidth - 14\tabcolsep) * \real{0.1250}}
  >{\raggedright\arraybackslash}p{(\linewidth - 14\tabcolsep) * \real{0.1250}}
  >{\raggedright\arraybackslash}p{(\linewidth - 14\tabcolsep) * \real{0.1250}}
  >{\raggedright\arraybackslash}p{(\linewidth - 14\tabcolsep) * \real{0.1250}}
  >{\raggedright\arraybackslash}p{(\linewidth - 14\tabcolsep) * \real{0.1250}}@{}}
\toprule\noalign{}
\begin{minipage}[b]{\linewidth}\raggedright
RUNSPEC
\end{minipage} & \begin{minipage}[b]{\linewidth}\raggedright
GRID
\end{minipage} & \begin{minipage}[b]{\linewidth}\raggedright
EDIT
\end{minipage} & \begin{minipage}[b]{\linewidth}\raggedright
PROPS
\end{minipage} & \begin{minipage}[b]{\linewidth}\raggedright
REGIONS
\end{minipage} & \begin{minipage}[b]{\linewidth}\raggedright
SOLUTION
\end{minipage} & \begin{minipage}[b]{\linewidth}\raggedright
SUMMARY
\end{minipage} & \begin{minipage}[b]{\linewidth}\raggedright
SCHEDULE
\end{minipage} \\
\midrule\noalign{}
\endhead
\bottomrule\noalign{}
\endlastfoot
\end{longtable}

\subsubsection{Description}\label{description-98}

NOCASC keyword activates the linear solver tracer algorithm for single
phase tracers.

OPM Flow uses a different numerical scheme which makes this keyword
redundant; hence, OPM Flow ignores this keyword. It is documented here
for completeness.

There is no data required for this keyword and there is no terminating
``/'' for this keyword.

\subsubsection{Example}\label{example-74}

\begin{verbatim}
--
--       TRACER SOLVER OPTION
--
NOCASC
\end{verbatim}

The above example switches on the linear solver tracer algorithm;
however, this has no effect in OPM Flow input decks.

\subsection{NODPPM -- Deactivate Fracture Porosity-Permeability
Calculation}\label{nodppm-deactivate-fracture-porosity-permeability-calculation}

\begin{longtable}[]{@{}
  >{\raggedright\arraybackslash}p{(\linewidth - 14\tabcolsep) * \real{0.1250}}
  >{\raggedright\arraybackslash}p{(\linewidth - 14\tabcolsep) * \real{0.1250}}
  >{\raggedright\arraybackslash}p{(\linewidth - 14\tabcolsep) * \real{0.1250}}
  >{\raggedright\arraybackslash}p{(\linewidth - 14\tabcolsep) * \real{0.1250}}
  >{\raggedright\arraybackslash}p{(\linewidth - 14\tabcolsep) * \real{0.1250}}
  >{\raggedright\arraybackslash}p{(\linewidth - 14\tabcolsep) * \real{0.1250}}
  >{\raggedright\arraybackslash}p{(\linewidth - 14\tabcolsep) * \real{0.1250}}
  >{\raggedright\arraybackslash}p{(\linewidth - 14\tabcolsep) * \real{0.1250}}@{}}
\toprule\noalign{}
\begin{minipage}[b]{\linewidth}\raggedright
RUNSPEC
\end{minipage} & \begin{minipage}[b]{\linewidth}\raggedright
GRID
\end{minipage} & \begin{minipage}[b]{\linewidth}\raggedright
EDIT
\end{minipage} & \begin{minipage}[b]{\linewidth}\raggedright
PROPS
\end{minipage} & \begin{minipage}[b]{\linewidth}\raggedright
REGIONS
\end{minipage} & \begin{minipage}[b]{\linewidth}\raggedright
SOLUTION
\end{minipage} & \begin{minipage}[b]{\linewidth}\raggedright
SUMMARY
\end{minipage} & \begin{minipage}[b]{\linewidth}\raggedright
SCHEDULE
\end{minipage} \\
\midrule\noalign{}
\endhead
\bottomrule\noalign{}
\endlastfoot
\end{longtable}

\subsubsection{Description}\label{description-99}

The NODPPM keyword deactivates the default behavior of multiplying the
fracture porosity by the fracture permeability to calculate the
effective fracture permeability in dual porosity and dual permeability
runs. Either the DUALPORO or DUALPERM keywords in the RUNSPEC section
must be declared in the input file in order to use this keyword. If the
default calculation is switched off by this keyword, then the effective
fracture permeability is taken to be those entered for the fracture
using the PERMX, PERMY and PERMZ keywords in the GRID section. If the
keyword is absent from the input deck, then the entered PERMX, PERMY and
PERMZ arrays for the fractures are multiplied by fracture PORO array
values in order to obtain the effective fracture permeability.

This keyword is not supported by OPM Flow but would change the results
if supported so the simulation will be stopped.

There is no data required for this keyword and there is no terminating
``/'' for this keyword.

\subsubsection{Example}\label{example-75}

\begin{verbatim}
--
--       DEACTIVATE FRACTURE POROSITY-PERMEABILITY CALCULATION
--
NODPPM
\end{verbatim}

The above example switches off the default behavior of multiplying the
fracture porosity by the fracture permeability to calculate the
effective fracture permeability.

\subsection{NOECHO -- Deactivate Echoing of User Input Files to the
Print
File}\label{noecho-deactivate-echoing-of-user-input-files-to-the-print-file-1}

Turns off echoing of all the input files to the print file. Note by
default echoing of the inputs files is active. but can subsequently be
switched off by the NOECHO activation keyword.

See NOECHO -- Deactivate Echoing of User Input Files to the Print File
in the GLOBAL section for a full description.

\subsection{NOHYST -- Deactivate the Hysteresis
Option}\label{nohyst-deactivate-the-hysteresis-option}

\begin{longtable}[]{@{}
  >{\raggedright\arraybackslash}p{(\linewidth - 14\tabcolsep) * \real{0.1250}}
  >{\raggedright\arraybackslash}p{(\linewidth - 14\tabcolsep) * \real{0.1250}}
  >{\raggedright\arraybackslash}p{(\linewidth - 14\tabcolsep) * \real{0.1250}}
  >{\raggedright\arraybackslash}p{(\linewidth - 14\tabcolsep) * \real{0.1250}}
  >{\raggedright\arraybackslash}p{(\linewidth - 14\tabcolsep) * \real{0.1250}}
  >{\raggedright\arraybackslash}p{(\linewidth - 14\tabcolsep) * \real{0.1250}}
  >{\raggedright\arraybackslash}p{(\linewidth - 14\tabcolsep) * \real{0.1250}}
  >{\raggedright\arraybackslash}p{(\linewidth - 14\tabcolsep) * \real{0.1250}}@{}}
\toprule\noalign{}
\begin{minipage}[b]{\linewidth}\raggedright
RUNSPEC
\end{minipage} & \begin{minipage}[b]{\linewidth}\raggedright
GRID
\end{minipage} & \begin{minipage}[b]{\linewidth}\raggedright
EDIT
\end{minipage} & \begin{minipage}[b]{\linewidth}\raggedright
PROPS
\end{minipage} & \begin{minipage}[b]{\linewidth}\raggedright
REGIONS
\end{minipage} & \begin{minipage}[b]{\linewidth}\raggedright
SOLUTION
\end{minipage} & \begin{minipage}[b]{\linewidth}\raggedright
SUMMARY
\end{minipage} & \begin{minipage}[b]{\linewidth}\raggedright
SCHEDULE
\end{minipage} \\
\midrule\noalign{}
\endhead
\bottomrule\noalign{}
\endlastfoot
\end{longtable}

\subsubsection{Description}\label{description-100}

The NOHYST keyword deactivates the Hysteresis option and informs the
simulator to ignore the IMBNUM array in the REGIONS section.

This keyword is not supported by OPM Flow but would change the results
if supported so the simulation will be stopped.

There is no data required for this keyword and there is no terminating
``/'' for this keyword.

\subsubsection{Example}\label{example-76}

\begin{verbatim}
--
--       DEACTIVATE THE HYSTERESIS OPTION
--
NOHYST
\end{verbatim}

The above example switches off the default behavior of multiplying the
fracture porosity by the fracture permeability to calculate the
effective fracture permeability.

\subsection{NOINSPEC -- Deactivate Output of the INIT Index
File}\label{noinspec-deactivate-output-of-the-init-index-file}

\begin{longtable}[]{@{}
  >{\raggedright\arraybackslash}p{(\linewidth - 14\tabcolsep) * \real{0.1250}}
  >{\raggedright\arraybackslash}p{(\linewidth - 14\tabcolsep) * \real{0.1250}}
  >{\raggedright\arraybackslash}p{(\linewidth - 14\tabcolsep) * \real{0.1250}}
  >{\raggedright\arraybackslash}p{(\linewidth - 14\tabcolsep) * \real{0.1250}}
  >{\raggedright\arraybackslash}p{(\linewidth - 14\tabcolsep) * \real{0.1250}}
  >{\raggedright\arraybackslash}p{(\linewidth - 14\tabcolsep) * \real{0.1250}}
  >{\raggedright\arraybackslash}p{(\linewidth - 14\tabcolsep) * \real{0.1250}}
  >{\raggedright\arraybackslash}p{(\linewidth - 14\tabcolsep) * \real{0.1250}}@{}}
\toprule\noalign{}
\begin{minipage}[b]{\linewidth}\raggedright
RUNSPEC
\end{minipage} & \begin{minipage}[b]{\linewidth}\raggedright
GRID
\end{minipage} & \begin{minipage}[b]{\linewidth}\raggedright
EDIT
\end{minipage} & \begin{minipage}[b]{\linewidth}\raggedright
PROPS
\end{minipage} & \begin{minipage}[b]{\linewidth}\raggedright
REGIONS
\end{minipage} & \begin{minipage}[b]{\linewidth}\raggedright
SOLUTION
\end{minipage} & \begin{minipage}[b]{\linewidth}\raggedright
SUMMARY
\end{minipage} & \begin{minipage}[b]{\linewidth}\raggedright
SCHEDULE
\end{minipage} \\
\midrule\noalign{}
\endhead
\bottomrule\noalign{}
\endlastfoot
\end{longtable}

\subsubsection{Description}\label{description-101}

The NOINSPEC keyword deactivates the writing out of the INIT index file
(\emph{.INSPEC). The initialization data (or static data) is written out
to two files one file contains the data, }.INIT, and the second file
contains an index of the data (\emph{.INSPEC) stored in the }.INIT file.
This functionality is redundant as most post-processing software require
the \emph{.INSPEC file to load the }.INIT data set.

Hence, OPM Flow ignores this keyword. It is documented here for
completeness.

There is no data required for this keyword and there is no terminating
``/'' for this keyword.

\subsubsection{Example}\label{example-77}

\begin{verbatim}
--
--       DEACTIVATE OUTPUT OF THE INIT INDEX FILE *.INSPEC
--
NOINSPEC
\end{verbatim}

The above example switches off the writing of the INIT index file
(*.INSPEC); however, this has no effect in OPM Flow input decks.

\subsection{NOMONITO -- Deactivate Output of the Monitoring Data and
File}\label{nomonito-deactivate-output-of-the-monitoring-data-and-file}

\begin{longtable}[]{@{}
  >{\raggedright\arraybackslash}p{(\linewidth - 14\tabcolsep) * \real{0.1250}}
  >{\raggedright\arraybackslash}p{(\linewidth - 14\tabcolsep) * \real{0.1250}}
  >{\raggedright\arraybackslash}p{(\linewidth - 14\tabcolsep) * \real{0.1250}}
  >{\raggedright\arraybackslash}p{(\linewidth - 14\tabcolsep) * \real{0.1250}}
  >{\raggedright\arraybackslash}p{(\linewidth - 14\tabcolsep) * \real{0.1250}}
  >{\raggedright\arraybackslash}p{(\linewidth - 14\tabcolsep) * \real{0.1250}}
  >{\raggedright\arraybackslash}p{(\linewidth - 14\tabcolsep) * \real{0.1250}}
  >{\raggedright\arraybackslash}p{(\linewidth - 14\tabcolsep) * \real{0.1250}}@{}}
\toprule\noalign{}
\begin{minipage}[b]{\linewidth}\raggedright
RUNSPEC
\end{minipage} & \begin{minipage}[b]{\linewidth}\raggedright
GRID
\end{minipage} & \begin{minipage}[b]{\linewidth}\raggedright
EDIT
\end{minipage} & \begin{minipage}[b]{\linewidth}\raggedright
PROPS
\end{minipage} & \begin{minipage}[b]{\linewidth}\raggedright
REGIONS
\end{minipage} & \begin{minipage}[b]{\linewidth}\raggedright
SOLUTION
\end{minipage} & \begin{minipage}[b]{\linewidth}\raggedright
SUMMARY
\end{minipage} & \begin{minipage}[b]{\linewidth}\raggedright
SCHEDULE
\end{minipage} \\
\midrule\noalign{}
\endhead
\bottomrule\noalign{}
\endlastfoot
\end{longtable}

\subsubsection{Description}\label{description-102}

The NOMONITO keyword deactivates the writing out of the run time
monitoring information used by post-processing graphics software to
display run time information, for example the simulated production and
injection rates and cumulative values. OPM Flow does not have this
functionality.

Hence, OPM Flow ignores this keyword but it is documented here for
completeness.

There is no data required for this keyword and there is no terminating
``/'' for this keyword.

\subsubsection{Example}\label{example-78}

\begin{verbatim}
--
--       DEACTIVATE MONITORING OUTPUT DATA AND FILES
--
NOMONITO
\end{verbatim}

The above example switches off the output required for run time
monitoring required by post-processing graphics software to review the
simulation results in real time as the run progresses; however, this has
no effect in OPM Flow input decks.

\subsection{NONNC -- Deactivate Non-Neighbor
Connections}\label{nonnc-deactivate-non-neighbor-connections}

\begin{longtable}[]{@{}
  >{\raggedright\arraybackslash}p{(\linewidth - 14\tabcolsep) * \real{0.1250}}
  >{\raggedright\arraybackslash}p{(\linewidth - 14\tabcolsep) * \real{0.1250}}
  >{\raggedright\arraybackslash}p{(\linewidth - 14\tabcolsep) * \real{0.1250}}
  >{\raggedright\arraybackslash}p{(\linewidth - 14\tabcolsep) * \real{0.1250}}
  >{\raggedright\arraybackslash}p{(\linewidth - 14\tabcolsep) * \real{0.1250}}
  >{\raggedright\arraybackslash}p{(\linewidth - 14\tabcolsep) * \real{0.1250}}
  >{\raggedright\arraybackslash}p{(\linewidth - 14\tabcolsep) * \real{0.1250}}
  >{\raggedright\arraybackslash}p{(\linewidth - 14\tabcolsep) * \real{0.1250}}@{}}
\toprule\noalign{}
\begin{minipage}[b]{\linewidth}\raggedright
RUNSPEC
\end{minipage} & \begin{minipage}[b]{\linewidth}\raggedright
GRID
\end{minipage} & \begin{minipage}[b]{\linewidth}\raggedright
EDIT
\end{minipage} & \begin{minipage}[b]{\linewidth}\raggedright
PROPS
\end{minipage} & \begin{minipage}[b]{\linewidth}\raggedright
REGIONS
\end{minipage} & \begin{minipage}[b]{\linewidth}\raggedright
SOLUTION
\end{minipage} & \begin{minipage}[b]{\linewidth}\raggedright
SUMMARY
\end{minipage} & \begin{minipage}[b]{\linewidth}\raggedright
SCHEDULE
\end{minipage} \\
\midrule\noalign{}
\endhead
\bottomrule\noalign{}
\endlastfoot
\end{longtable}

\subsubsection{Description}\label{description-103}

The NONNC keyword deactivates non-neighbor connections (``NNCs'') in the
current run. NNCs create off-diagonal elements in the Jacobi matrix that
impact the numerical efficiency of the solution algorithms, and thus if
the run does not contain NNCs then there is the potential for greater
computation efficiency. Unfortunately, nearly all models, except for the
most simple models, generate NNCs via for example:

\begin{itemize}
\tightlist
\item
  aquifer connections,
\item
  faults, and
\item
  manually entered NNCs, including those automatically generate by
  pre-processing software.
\end{itemize}

There is no data required for this keyword and there is no terminating
``/'' for this keyword.

\subsubsection{Example}\label{example-79}

\begin{verbatim}
--
--       DEACTIVATE NON-NEIGHBOR CONNECTIONS
--
NONNC
\end{verbatim}

The above example switches off the NNCs.

\subsection{NORSSPEC -- Deactivate Output of the RESTART Index
File}\label{norsspec-deactivate-output-of-the-restart-index-file}

\begin{longtable}[]{@{}
  >{\raggedright\arraybackslash}p{(\linewidth - 14\tabcolsep) * \real{0.1250}}
  >{\raggedright\arraybackslash}p{(\linewidth - 14\tabcolsep) * \real{0.1250}}
  >{\raggedright\arraybackslash}p{(\linewidth - 14\tabcolsep) * \real{0.1250}}
  >{\raggedright\arraybackslash}p{(\linewidth - 14\tabcolsep) * \real{0.1250}}
  >{\raggedright\arraybackslash}p{(\linewidth - 14\tabcolsep) * \real{0.1250}}
  >{\raggedright\arraybackslash}p{(\linewidth - 14\tabcolsep) * \real{0.1250}}
  >{\raggedright\arraybackslash}p{(\linewidth - 14\tabcolsep) * \real{0.1250}}
  >{\raggedright\arraybackslash}p{(\linewidth - 14\tabcolsep) * \real{0.1250}}@{}}
\toprule\noalign{}
\begin{minipage}[b]{\linewidth}\raggedright
RUNSPEC
\end{minipage} & \begin{minipage}[b]{\linewidth}\raggedright
GRID
\end{minipage} & \begin{minipage}[b]{\linewidth}\raggedright
EDIT
\end{minipage} & \begin{minipage}[b]{\linewidth}\raggedright
PROPS
\end{minipage} & \begin{minipage}[b]{\linewidth}\raggedright
REGIONS
\end{minipage} & \begin{minipage}[b]{\linewidth}\raggedright
SOLUTION
\end{minipage} & \begin{minipage}[b]{\linewidth}\raggedright
SUMMARY
\end{minipage} & \begin{minipage}[b]{\linewidth}\raggedright
SCHEDULE
\end{minipage} \\
\midrule\noalign{}
\endhead
\bottomrule\noalign{}
\endlastfoot
\end{longtable}

\subsubsection{Description}\label{description-104}

The NORSSPEC keyword deactivates the writing out of the RESTART index
file (\emph{.RSSPEC). The restart data (pressure, saturations etc.
through time for each active cell) are written out to two files one file
contains the data, }.UNRST for example, and the second file contains an
index of the data (\emph{.RSSPEC) stored in the }.UNRST file. This
functionality is redundant as most post-processing software require the
\emph{.RSSPEC file to load the }.UNRST data set.

Hence, OPM Flow ignores this keyword. It is documented here for
completeness.

There is no data required for this keyword and there is no terminating
``/'' for this keyword.

\subsubsection{Example}\label{example-80}

\begin{verbatim}
--
--       DEACTIVATE OUTPUT OF THE RESTART INDEX FILE *.RSSPEC
--
NORSSPEC
\end{verbatim}

The above example switches off the writing of the restart index file
(*.RSSPEC); however, this has no effect in OPM Flow input decks.

\subsection{NOSIM -- Activate the No Simulation Mode for Data File
Checking}\label{nosim-activate-the-no-simulation-mode-for-data-file-checking}

\begin{longtable}[]{@{}
  >{\raggedright\arraybackslash}p{(\linewidth - 14\tabcolsep) * \real{0.1250}}
  >{\raggedright\arraybackslash}p{(\linewidth - 14\tabcolsep) * \real{0.1250}}
  >{\raggedright\arraybackslash}p{(\linewidth - 14\tabcolsep) * \real{0.1250}}
  >{\raggedright\arraybackslash}p{(\linewidth - 14\tabcolsep) * \real{0.1250}}
  >{\raggedright\arraybackslash}p{(\linewidth - 14\tabcolsep) * \real{0.1250}}
  >{\raggedright\arraybackslash}p{(\linewidth - 14\tabcolsep) * \real{0.1250}}
  >{\raggedright\arraybackslash}p{(\linewidth - 14\tabcolsep) * \real{0.1250}}
  >{\raggedright\arraybackslash}p{(\linewidth - 14\tabcolsep) * \real{0.1250}}@{}}
\toprule\noalign{}
\begin{minipage}[b]{\linewidth}\raggedright
RUNSPEC
\end{minipage} & \begin{minipage}[b]{\linewidth}\raggedright
GRID
\end{minipage} & \begin{minipage}[b]{\linewidth}\raggedright
EDIT
\end{minipage} & \begin{minipage}[b]{\linewidth}\raggedright
PROPS
\end{minipage} & \begin{minipage}[b]{\linewidth}\raggedright
REGIONS
\end{minipage} & \begin{minipage}[b]{\linewidth}\raggedright
SOLUTION
\end{minipage} & \begin{minipage}[b]{\linewidth}\raggedright
SUMMARY
\end{minipage} & \begin{minipage}[b]{\linewidth}\raggedright
SCHEDULE
\end{minipage} \\
\midrule\noalign{}
\endhead
\bottomrule\noalign{}
\endlastfoot
\end{longtable}

\subsubsection{Description}\label{description-105}

NOSIM switches the mode of OPM Flow to data input checking mode. In this
mode the input file is read and all messages and print instructions are
sent to the respective output files. The SCHEDULE section is read but
the simulation is not performed.

There is no data required for this keyword and there is no terminating
``/'' for this keyword.

\subsubsection{Example}\label{example-81}

The example below switches OPM Flow to no simulation mode for data
checking of the input deck.

\begin{verbatim}
--
--       SWITCH NO SIMULATION MODE FOR DATA CHECKING COMMENT OUT TO RUN THE MODEL
--
NOSIM
\end{verbatim}

And the next example shows how to commented out the NOSIM activation
keyword so that the simulation will proceed.

\begin{verbatim}
--
--       SWITCH NO SIMULATION MODE FOR DATA CHECKING COMMENT OUT TO RUN THE MODEL
--
-- NOSIM
\end{verbatim}

\begin{tcolorbox}[enhanced jigsaw, title=\textcolor{quarto-callout-note-color}{\faInfo}\hspace{0.5em}{Note}, coltitle=black, breakable, toprule=.15mm, colbacktitle=quarto-callout-note-color!10!white, leftrule=.75mm, opacitybacktitle=0.6, colback=white, bottomtitle=1mm, toptitle=1mm, titlerule=0mm, left=2mm, colframe=quarto-callout-note-color-frame, arc=.35mm, rightrule=.15mm, bottomrule=.15mm, opacityback=0]

Simulation input decks are complex and are therefore prone to typing
errors, thus before submitting a run that will take over 15 minutes or
so, it is a good idea to run the model with the NOSIM option. If no
errors are found then the NOSIM keyword should be commented out by
placing ``--'' before the keyword, and then re-running the model.
Alternatively, one could use OPMRUN to run all the jobs in the queue in
NOSIM mode and have the software re-run jobs in simulation mode if there
are no errors.

\end{tcolorbox}

\subsection{NOWARN -- Deactivate Warning
Messages}\label{nowarn-deactivate-warning-messages-1}

Turns off warning messages to be printed to the print file; note that
this keyword is deactivated by default and can subsequently be switched
off by the WARN activation keyword. The warning messages may be turned
on and off using keywords WARN and NOWARN.

See NOWARN -- Deactivate Warning Messages in the GLOBAL section for a
full description.

\subsection{NRSOUT -- Defined Maximum Number of RESTART
Elements}\label{nrsout-defined-maximum-number-of-restart-elements}

\begin{longtable}[]{@{}
  >{\raggedright\arraybackslash}p{(\linewidth - 14\tabcolsep) * \real{0.1250}}
  >{\raggedright\arraybackslash}p{(\linewidth - 14\tabcolsep) * \real{0.1250}}
  >{\raggedright\arraybackslash}p{(\linewidth - 14\tabcolsep) * \real{0.1250}}
  >{\raggedright\arraybackslash}p{(\linewidth - 14\tabcolsep) * \real{0.1250}}
  >{\raggedright\arraybackslash}p{(\linewidth - 14\tabcolsep) * \real{0.1250}}
  >{\raggedright\arraybackslash}p{(\linewidth - 14\tabcolsep) * \real{0.1250}}
  >{\raggedright\arraybackslash}p{(\linewidth - 14\tabcolsep) * \real{0.1250}}
  >{\raggedright\arraybackslash}p{(\linewidth - 14\tabcolsep) * \real{0.1250}}@{}}
\toprule\noalign{}
\begin{minipage}[b]{\linewidth}\raggedright
RUNSPEC
\end{minipage} & \begin{minipage}[b]{\linewidth}\raggedright
GRID
\end{minipage} & \begin{minipage}[b]{\linewidth}\raggedright
EDIT
\end{minipage} & \begin{minipage}[b]{\linewidth}\raggedright
PROPS
\end{minipage} & \begin{minipage}[b]{\linewidth}\raggedright
REGIONS
\end{minipage} & \begin{minipage}[b]{\linewidth}\raggedright
SOLUTION
\end{minipage} & \begin{minipage}[b]{\linewidth}\raggedright
SUMMARY
\end{minipage} & \begin{minipage}[b]{\linewidth}\raggedright
SCHEDULE
\end{minipage} \\
\midrule\noalign{}
\endhead
\bottomrule\noalign{}
\endlastfoot
\end{longtable}

\subsubsection{Description}\label{description-106}

The NRSOUT keyword specifies the maximum number of elements that can be
written to the RESTART file at each reporting time step.

This keyword is not supported by OPM Flow but would change the results
if supported so the simulation will be stopped.

\begin{longtable}[]{@{}
  >{\raggedright\arraybackslash}p{(\linewidth - 6\tabcolsep) * \real{0.2500}}
  >{\raggedright\arraybackslash}p{(\linewidth - 6\tabcolsep) * \real{0.2500}}
  >{\raggedright\arraybackslash}p{(\linewidth - 6\tabcolsep) * \real{0.2500}}
  >{\raggedright\arraybackslash}p{(\linewidth - 6\tabcolsep) * \real{0.2500}}@{}}
\toprule\noalign{}
\begin{minipage}[b]{\linewidth}\raggedright
No.
\end{minipage} & \begin{minipage}[b]{\linewidth}\raggedright
Name
\end{minipage} & \begin{minipage}[b]{\linewidth}\raggedright
Description
\end{minipage} & \begin{minipage}[b]{\linewidth}\raggedright
Default
\end{minipage} \\
\midrule\noalign{}
\endhead
\bottomrule\noalign{}
\endlastfoot
1 & NRSOUT & A positive integer value that specifies the maximum number
of elements that can be written to the RESTART file at each reporting
time step. & 3600 \\
Notes: & & & \\
\end{longtable}

\emph{Table 5.28: NRSOUT Keyword Description}

\subsubsection{Example}\label{example-82}

\begin{verbatim}
--
--       MAX
--       NRSOUT
NRSOUT
         6000                                                     /
\end{verbatim}

The above example sets the maximum number of elements that can be
written to the RESTART file at each reporting time step to 6000.

\subsection{NSTACK -- Define the Stack Length for the Iterative Linear
Solver}\label{nstack-define-the-stack-length-for-the-iterative-linear-solver}

\begin{longtable}[]{@{}
  >{\raggedright\arraybackslash}p{(\linewidth - 14\tabcolsep) * \real{0.1250}}
  >{\raggedright\arraybackslash}p{(\linewidth - 14\tabcolsep) * \real{0.1250}}
  >{\raggedright\arraybackslash}p{(\linewidth - 14\tabcolsep) * \real{0.1250}}
  >{\raggedright\arraybackslash}p{(\linewidth - 14\tabcolsep) * \real{0.1250}}
  >{\raggedright\arraybackslash}p{(\linewidth - 14\tabcolsep) * \real{0.1250}}
  >{\raggedright\arraybackslash}p{(\linewidth - 14\tabcolsep) * \real{0.1250}}
  >{\raggedright\arraybackslash}p{(\linewidth - 14\tabcolsep) * \real{0.1250}}
  >{\raggedright\arraybackslash}p{(\linewidth - 14\tabcolsep) * \real{0.1250}}@{}}
\toprule\noalign{}
\begin{minipage}[b]{\linewidth}\raggedright
RUNSPEC
\end{minipage} & \begin{minipage}[b]{\linewidth}\raggedright
GRID
\end{minipage} & \begin{minipage}[b]{\linewidth}\raggedright
EDIT
\end{minipage} & \begin{minipage}[b]{\linewidth}\raggedright
PROPS
\end{minipage} & \begin{minipage}[b]{\linewidth}\raggedright
REGIONS
\end{minipage} & \begin{minipage}[b]{\linewidth}\raggedright
SOLUTION
\end{minipage} & \begin{minipage}[b]{\linewidth}\raggedright
SUMMARY
\end{minipage} & \begin{minipage}[b]{\linewidth}\raggedright
SCHEDULE
\end{minipage} \\
\midrule\noalign{}
\endhead
\bottomrule\noalign{}
\endlastfoot
\end{longtable}

\subsubsection{Description}\label{description-107}

The NSTACK keyword defines the maximum number of previous search
directions stored by the linear solver. Increasing the value of NSTACK
may improve the efficiency of the solver on difficult problems, but will
increase the memory requirements of the simulator. The default value of
10 should be sufficient for most problems; however, if OPM Flow is
having issues with the convergence of the linear questions then
increasing NSTACK and the value of LITMAX on the TUNING keyword may
improve performance.

OPM Flow uses a different numerical scheme which makes this keyword
redundant; see section 2.2 Running OPM Flow 2023-04 From The Command
Line on how to invoke various numerical schemes via the OPM Flow command
line interface.

This keyword is not supported by OPM Flow but has no effect on the
results so it will be ignored.

\begin{longtable}[]{@{}
  >{\raggedright\arraybackslash}p{(\linewidth - 6\tabcolsep) * \real{0.2500}}
  >{\raggedright\arraybackslash}p{(\linewidth - 6\tabcolsep) * \real{0.2500}}
  >{\raggedright\arraybackslash}p{(\linewidth - 6\tabcolsep) * \real{0.2500}}
  >{\raggedright\arraybackslash}p{(\linewidth - 6\tabcolsep) * \real{0.2500}}@{}}
\toprule\noalign{}
\begin{minipage}[b]{\linewidth}\raggedright
No.
\end{minipage} & \begin{minipage}[b]{\linewidth}\raggedright
Name
\end{minipage} & \begin{minipage}[b]{\linewidth}\raggedright
Description
\end{minipage} & \begin{minipage}[b]{\linewidth}\raggedright
Default
\end{minipage} \\
\midrule\noalign{}
\endhead
\bottomrule\noalign{}
\endlastfoot
1 & NSTACK & A positive integer that defines the maximum number of
previous search directions stored by the linear solver. & 10 \\
Notes: & & & \\
\end{longtable}

\emph{Table 5.29: NSTACK Keyword Description}

\subsubsection{Example}\label{example-83}

\begin{verbatim}
--
--         SET STACK SIZE FOR LINEAR SOLVER
--
NSTACK
           30                                                                  /
\end{verbatim}

The above example sets maximum number of previous search directions
stored by the linear solver to 30, this has no effect in OPM Flow input
decks.

\begin{tcolorbox}[enhanced jigsaw, title=\textcolor{quarto-callout-note-color}{\faInfo}\hspace{0.5em}{Note}, coltitle=black, breakable, toprule=.15mm, colbacktitle=quarto-callout-note-color!10!white, leftrule=.75mm, opacitybacktitle=0.6, colback=white, bottomtitle=1mm, toptitle=1mm, titlerule=0mm, left=2mm, colframe=quarto-callout-note-color-frame, arc=.35mm, rightrule=.15mm, bottomrule=.15mm, opacityback=0]

If the run is suffering from linear convergence problems, then check the
data first for any data issues before manipulating the numerical control
parameters. For example, if OPM Flow has written some WARNING messages
with respect to end-point scaling, etc., then resolve these messages
first before adjusting the numerical controls.

\end{tcolorbox}

\subsection{NUMRES -- Define the Number of Reservoir
Grids}\label{numres-define-the-number-of-reservoir-grids}

\begin{longtable}[]{@{}
  >{\raggedright\arraybackslash}p{(\linewidth - 14\tabcolsep) * \real{0.1250}}
  >{\raggedright\arraybackslash}p{(\linewidth - 14\tabcolsep) * \real{0.1250}}
  >{\raggedright\arraybackslash}p{(\linewidth - 14\tabcolsep) * \real{0.1250}}
  >{\raggedright\arraybackslash}p{(\linewidth - 14\tabcolsep) * \real{0.1250}}
  >{\raggedright\arraybackslash}p{(\linewidth - 14\tabcolsep) * \real{0.1250}}
  >{\raggedright\arraybackslash}p{(\linewidth - 14\tabcolsep) * \real{0.1250}}
  >{\raggedright\arraybackslash}p{(\linewidth - 14\tabcolsep) * \real{0.1250}}
  >{\raggedright\arraybackslash}p{(\linewidth - 14\tabcolsep) * \real{0.1250}}@{}}
\toprule\noalign{}
\begin{minipage}[b]{\linewidth}\raggedright
RUNSPEC
\end{minipage} & \begin{minipage}[b]{\linewidth}\raggedright
GRID
\end{minipage} & \begin{minipage}[b]{\linewidth}\raggedright
EDIT
\end{minipage} & \begin{minipage}[b]{\linewidth}\raggedright
PROPS
\end{minipage} & \begin{minipage}[b]{\linewidth}\raggedright
REGIONS
\end{minipage} & \begin{minipage}[b]{\linewidth}\raggedright
SOLUTION
\end{minipage} & \begin{minipage}[b]{\linewidth}\raggedright
SUMMARY
\end{minipage} & \begin{minipage}[b]{\linewidth}\raggedright
SCHEDULE
\end{minipage} \\
\midrule\noalign{}
\endhead
\bottomrule\noalign{}
\endlastfoot
\end{longtable}

\subsubsection{Description}\label{description-108}

The NUMRES keyword defines the number of reservoir grids (COORD data
sets) that the simulator should process. OPM Flow currently only
supports a single reservoir grid.

\begin{longtable}[]{@{}
  >{\raggedright\arraybackslash}p{(\linewidth - 6\tabcolsep) * \real{0.2500}}
  >{\raggedright\arraybackslash}p{(\linewidth - 6\tabcolsep) * \real{0.2500}}
  >{\raggedright\arraybackslash}p{(\linewidth - 6\tabcolsep) * \real{0.2500}}
  >{\raggedright\arraybackslash}p{(\linewidth - 6\tabcolsep) * \real{0.2500}}@{}}
\toprule\noalign{}
\begin{minipage}[b]{\linewidth}\raggedright
No.
\end{minipage} & \begin{minipage}[b]{\linewidth}\raggedright
Name
\end{minipage} & \begin{minipage}[b]{\linewidth}\raggedright
Description
\end{minipage} & \begin{minipage}[b]{\linewidth}\raggedright
Default
\end{minipage} \\
\midrule\noalign{}
\endhead
\bottomrule\noalign{}
\endlastfoot
1 & NUMRES & A positive integer that defines the number of COORD data
sets to be processed by OPM Flow. OPM Flow currently only supports a
single reservoir grid and so this item should be defaulted (1*) or set
to one. & 1 \\
Notes: & & & \\
\end{longtable}

\emph{Table 5.30: NUMRES Keyword Description}

\subsubsection{Example}\label{example-84}

\begin{verbatim}
--
--       DEFINE THE NUMBER OF RESERVOIR GRIDS (COORD DATA SETS)
--
NUMRES
         1                                                                    /
\end{verbatim}

The above example sets the maximum number of COORD data sets to be
processed to one.

\subsection{NUPCOL -- Define the Maximum Number of Newton Iterations
Used to Update Well
Targets}\label{nupcol-define-the-maximum-number-of-newton-iterations-used-to-update-well-targets}

\begin{longtable}[]{@{}
  >{\raggedright\arraybackslash}p{(\linewidth - 14\tabcolsep) * \real{0.1250}}
  >{\raggedright\arraybackslash}p{(\linewidth - 14\tabcolsep) * \real{0.1250}}
  >{\raggedright\arraybackslash}p{(\linewidth - 14\tabcolsep) * \real{0.1250}}
  >{\raggedright\arraybackslash}p{(\linewidth - 14\tabcolsep) * \real{0.1250}}
  >{\raggedright\arraybackslash}p{(\linewidth - 14\tabcolsep) * \real{0.1250}}
  >{\raggedright\arraybackslash}p{(\linewidth - 14\tabcolsep) * \real{0.1250}}
  >{\raggedright\arraybackslash}p{(\linewidth - 14\tabcolsep) * \real{0.1250}}
  >{\raggedright\arraybackslash}p{(\linewidth - 14\tabcolsep) * \real{0.1250}}@{}}
\toprule\noalign{}
\begin{minipage}[b]{\linewidth}\raggedright
RUNSPEC
\end{minipage} & \begin{minipage}[b]{\linewidth}\raggedright
GRID
\end{minipage} & \begin{minipage}[b]{\linewidth}\raggedright
EDIT
\end{minipage} & \begin{minipage}[b]{\linewidth}\raggedright
PROPS
\end{minipage} & \begin{minipage}[b]{\linewidth}\raggedright
REGIONS
\end{minipage} & \begin{minipage}[b]{\linewidth}\raggedright
SOLUTION
\end{minipage} & \begin{minipage}[b]{\linewidth}\raggedright
SUMMARY
\end{minipage} & \begin{minipage}[b]{\linewidth}\raggedright
SCHEDULE
\end{minipage} \\
\midrule\noalign{}
\endhead
\bottomrule\noalign{}
\endlastfoot
\end{longtable}

\subsubsection{Description}\label{description-109}

The NUPCOL keyword defines the maximum number of Newton iterations
within a time step that may be used to update the well production and
injection targets, after which the well targets will be frozen until the
time step calculations have converged and the time step is complete.

Wells under group control may suffer from some dependency with other
wells in the same group that are also under group control. This can
cause oscillations in the production and injection well rates within the
group. In order to avoid this, after the number Newton iterations within
a time step exceeds NUPCOL, the group well rates are frozen until the
time step has converged. Reducing the potential for well rate
oscillations within the time step may result in the group targets and
limits not being exactly met in this case. Increasing the value of
NUPCOL will improve the accuracy of the group targets and limits at the
expense of computational efficiency.

See also section 2.2 Running OPM Flow 2023-04 From The Command Line on
how to set various other numerical control parameters for OPM Flow.

\begin{longtable}[]{@{}
  >{\raggedright\arraybackslash}p{(\linewidth - 6\tabcolsep) * \real{0.2500}}
  >{\raggedright\arraybackslash}p{(\linewidth - 6\tabcolsep) * \real{0.2500}}
  >{\raggedright\arraybackslash}p{(\linewidth - 6\tabcolsep) * \real{0.2500}}
  >{\raggedright\arraybackslash}p{(\linewidth - 6\tabcolsep) * \real{0.2500}}@{}}
\toprule\noalign{}
\begin{minipage}[b]{\linewidth}\raggedright
No.
\end{minipage} & \begin{minipage}[b]{\linewidth}\raggedright
Name
\end{minipage} & \begin{minipage}[b]{\linewidth}\raggedright
Description
\end{minipage} & \begin{minipage}[b]{\linewidth}\raggedright
Default
\end{minipage} \\
\midrule\noalign{}
\endhead
\bottomrule\noalign{}
\endlastfoot
1 & NUPCOL & A positive integer that defines the maximum number of
Newton iterations used to update well targets within a time step. & 3 \\
Notes: & & & \\
\end{longtable}

\emph{Table 5.31: NUPCOL Keyword Description}

See also the MINNPCOL keyword in the RUNSPEC section that sets the
minimum number of Newton iterations within a time step that may be used
to update the well production and injection targets, after which the
well targets will be frozen until the time step calculations have
converged and the time step is complete. Note that MINNPCOL is an OPM
Flow specific keyword.

\subsubsection{Example}\label{example-85}

\begin{verbatim}
--
--       DEFINE THE MAX NUMBER OF ITERATIONS TO UPDATE WELL FLOW TARGETS
--
NUPCOL
         3                                                                    /
\end{verbatim}

The above example sets the NUPCOL value to its default value of 3.

\subsection{OIL -- Activate the Oil Phase in the
Model}\label{oil-activate-the-oil-phase-in-the-model}

\begin{longtable}[]{@{}
  >{\raggedright\arraybackslash}p{(\linewidth - 14\tabcolsep) * \real{0.1250}}
  >{\raggedright\arraybackslash}p{(\linewidth - 14\tabcolsep) * \real{0.1250}}
  >{\raggedright\arraybackslash}p{(\linewidth - 14\tabcolsep) * \real{0.1250}}
  >{\raggedright\arraybackslash}p{(\linewidth - 14\tabcolsep) * \real{0.1250}}
  >{\raggedright\arraybackslash}p{(\linewidth - 14\tabcolsep) * \real{0.1250}}
  >{\raggedright\arraybackslash}p{(\linewidth - 14\tabcolsep) * \real{0.1250}}
  >{\raggedright\arraybackslash}p{(\linewidth - 14\tabcolsep) * \real{0.1250}}
  >{\raggedright\arraybackslash}p{(\linewidth - 14\tabcolsep) * \real{0.1250}}@{}}
\toprule\noalign{}
\begin{minipage}[b]{\linewidth}\raggedright
RUNSPEC
\end{minipage} & \begin{minipage}[b]{\linewidth}\raggedright
GRID
\end{minipage} & \begin{minipage}[b]{\linewidth}\raggedright
EDIT
\end{minipage} & \begin{minipage}[b]{\linewidth}\raggedright
PROPS
\end{minipage} & \begin{minipage}[b]{\linewidth}\raggedright
REGIONS
\end{minipage} & \begin{minipage}[b]{\linewidth}\raggedright
SOLUTION
\end{minipage} & \begin{minipage}[b]{\linewidth}\raggedright
SUMMARY
\end{minipage} & \begin{minipage}[b]{\linewidth}\raggedright
SCHEDULE
\end{minipage} \\
\midrule\noalign{}
\endhead
\bottomrule\noalign{}
\endlastfoot
\end{longtable}

\subsubsection{Description}\label{description-110}

This keyword indicates that the oil phase is present in the model and
must be used for oil-gas, oil-water, oil-water-gas input decks that
contain the oil phase. The keyword will also invoke data input file
checking to ensure that all the required oil phase input parameters are
defined in the input deck.

There is no data required for this keyword and there is no terminating
``/'' for this keyword.

\subsubsection{Example}\label{example-86}

\begin{verbatim}
--
--       OIL PHASE IS PRESENT IN THE RUN
--
OIL
\end{verbatim}

The above example declares that the oil phase is active in the model.

\subsection{OPTIONS -- Activate Various Program
Options}\label{options-activate-various-program-options}

\begin{longtable}[]{@{}
  >{\raggedright\arraybackslash}p{(\linewidth - 14\tabcolsep) * \real{0.1250}}
  >{\raggedright\arraybackslash}p{(\linewidth - 14\tabcolsep) * \real{0.1250}}
  >{\raggedright\arraybackslash}p{(\linewidth - 14\tabcolsep) * \real{0.1250}}
  >{\raggedright\arraybackslash}p{(\linewidth - 14\tabcolsep) * \real{0.1250}}
  >{\raggedright\arraybackslash}p{(\linewidth - 14\tabcolsep) * \real{0.1250}}
  >{\raggedright\arraybackslash}p{(\linewidth - 14\tabcolsep) * \real{0.1250}}
  >{\raggedright\arraybackslash}p{(\linewidth - 14\tabcolsep) * \real{0.1250}}
  >{\raggedright\arraybackslash}p{(\linewidth - 14\tabcolsep) * \real{0.1250}}@{}}
\toprule\noalign{}
\begin{minipage}[b]{\linewidth}\raggedright
RUNSPEC
\end{minipage} & \begin{minipage}[b]{\linewidth}\raggedright
GRID
\end{minipage} & \begin{minipage}[b]{\linewidth}\raggedright
EDIT
\end{minipage} & \begin{minipage}[b]{\linewidth}\raggedright
PROPS
\end{minipage} & \begin{minipage}[b]{\linewidth}\raggedright
REGIONS
\end{minipage} & \begin{minipage}[b]{\linewidth}\raggedright
SOLUTION
\end{minipage} & \begin{minipage}[b]{\linewidth}\raggedright
SUMMARY
\end{minipage} & \begin{minipage}[b]{\linewidth}\raggedright
SCHEDULE
\end{minipage} \\
\midrule\noalign{}
\endhead
\bottomrule\noalign{}
\endlastfoot
\end{longtable}

\subsubsection{Description}\label{description-111}

The OPTIONS keyword activates various program options in the commercial
simulator. Currently, none of the options available in the commercial
simulator are implemented in OPM Flow, and it is unlikely that this
keyword will be supported in the future releases of OPM Flow.

This keyword is not supported by OPM Flow but would change the results
if supported so the simulation will be stopped.

\begin{longtable}[]{@{}
  >{\raggedright\arraybackslash}p{(\linewidth - 4\tabcolsep) * \real{0.3333}}
  >{\raggedright\arraybackslash}p{(\linewidth - 4\tabcolsep) * \real{0.3333}}
  >{\raggedright\arraybackslash}p{(\linewidth - 4\tabcolsep) * \real{0.3333}}@{}}
\toprule\noalign{}
\begin{minipage}[b]{\linewidth}\raggedright
No.
\end{minipage} & \begin{minipage}[b]{\linewidth}\raggedright
Description
\end{minipage} & \begin{minipage}[b]{\linewidth}\raggedright
Default
\end{minipage} \\
\midrule\noalign{}
\endhead
\bottomrule\noalign{}
\endlastfoot
1 - 273 & Commercial simulator options keyword, used to to switch on or
off specific features. The keyword is commonly used to revert the
simulator's behavior to past functionality that has been depreciated in
the current version of the commercial simulator, for backward
compatibility with previous models. & 0 \\
Notes: & & \\
\end{longtable}

\emph{Table 5.32: OPTIONS Keyword Description}

\subsubsection{Examples}\label{examples-3}

\begin{verbatim}
--
--       SKIP      ACTIVATE
--       OPTIONS   OPTION
OPTIONS
         77*0      1                                                           /
\end{verbatim}

The above example activates the use of scratch files for pre-processing
grid geometry data for non-neighbor connections. Note if multiple
options are required then one can just repeat the format of the example
to activate multiple options as the keyword does not overwrite previous
entries. So for example:

\begin{verbatim}
--
--       SKIP      ACTIVATE
--       OPTIONS   OPTION
OPTIONS
          7*0      1                                                           /
--
--       SKIP      ACTIVATE
--       OPTIONS   OPTION
OPTIONS
         77*0      1                                                           /
--
--       SKIP      ACTIVATE
--       OPTIONS   OPTION
OPTIONS
         177*0     1                                                           /
\end{verbatim}

Could be used to activate the 8, 78 and 178 options if they were
available.

\subsection{PARALLEL -- Define Parallel Run
Configuration}\label{parallel-define-parallel-run-configuration}

\begin{longtable}[]{@{}
  >{\raggedright\arraybackslash}p{(\linewidth - 14\tabcolsep) * \real{0.1250}}
  >{\raggedright\arraybackslash}p{(\linewidth - 14\tabcolsep) * \real{0.1250}}
  >{\raggedright\arraybackslash}p{(\linewidth - 14\tabcolsep) * \real{0.1250}}
  >{\raggedright\arraybackslash}p{(\linewidth - 14\tabcolsep) * \real{0.1250}}
  >{\raggedright\arraybackslash}p{(\linewidth - 14\tabcolsep) * \real{0.1250}}
  >{\raggedright\arraybackslash}p{(\linewidth - 14\tabcolsep) * \real{0.1250}}
  >{\raggedright\arraybackslash}p{(\linewidth - 14\tabcolsep) * \real{0.1250}}
  >{\raggedright\arraybackslash}p{(\linewidth - 14\tabcolsep) * \real{0.1250}}@{}}
\toprule\noalign{}
\begin{minipage}[b]{\linewidth}\raggedright
RUNSPEC
\end{minipage} & \begin{minipage}[b]{\linewidth}\raggedright
GRID
\end{minipage} & \begin{minipage}[b]{\linewidth}\raggedright
EDIT
\end{minipage} & \begin{minipage}[b]{\linewidth}\raggedright
PROPS
\end{minipage} & \begin{minipage}[b]{\linewidth}\raggedright
REGIONS
\end{minipage} & \begin{minipage}[b]{\linewidth}\raggedright
SOLUTION
\end{minipage} & \begin{minipage}[b]{\linewidth}\raggedright
SUMMARY
\end{minipage} & \begin{minipage}[b]{\linewidth}\raggedright
SCHEDULE
\end{minipage} \\
\midrule\noalign{}
\endhead
\bottomrule\noalign{}
\endlastfoot
\end{longtable}

\subsubsection{Description}\label{description-112}

The PARALLEL keyword defines the run to use parallel processing and sets
the domain decomposition options. See section 2.2 Running OPM Flow
2023-04 From The Command Line on how to run OPM Flow in parallel mode.

\begin{longtable}[]{@{}
  >{\raggedright\arraybackslash}p{(\linewidth - 6\tabcolsep) * \real{0.2500}}
  >{\raggedright\arraybackslash}p{(\linewidth - 6\tabcolsep) * \real{0.2500}}
  >{\raggedright\arraybackslash}p{(\linewidth - 6\tabcolsep) * \real{0.2500}}
  >{\raggedright\arraybackslash}p{(\linewidth - 6\tabcolsep) * \real{0.2500}}@{}}
\toprule\noalign{}
\begin{minipage}[b]{\linewidth}\raggedright
No.
\end{minipage} & \begin{minipage}[b]{\linewidth}\raggedright
Name
\end{minipage} & \begin{minipage}[b]{\linewidth}\raggedright
Description
\end{minipage} & \begin{minipage}[b]{\linewidth}\raggedright
Default
\end{minipage} \\
\midrule\noalign{}
\endhead
\bottomrule\noalign{}
\endlastfoot
1 & NPROCS & A positive integer that defines the number of domains or
parallel processors to use for this run. & 1 \\
2 & RTYPE & A character string set to either SERIAL to run the parallel
code in serial mode for testing the code, or DISTRIBUTED to full utilize
parallel processing. & DISTRIBUTED \\
Notes: & & & \\
\end{longtable}

\emph{Table 5.33: PARALLEL Keyword Description}

OPM Flow uses a different numerical scheme which makes this keyword
redundant; hence, OPM Flow ignores this keyword. It is documented here
for completeness.

\subsubsection{Example}\label{example-87}

\begin{verbatim}
--
--       PARALLEL MULTI-CORE OPTIONS
--       NDMAIN     MACHINE TYPE
PARALLEL
           2        DISTRIBUTED                                                /
\end{verbatim}

The above example sets the number of domains (or processors) to two and
for the simulation to run in parallel mode. This has no effect in OPM
Flow input decks.

\subsection{PARTTRAC -- Activate and Define Partitioned Tracer
Option}\label{parttrac-activate-and-define-partitioned-tracer-option}

\begin{longtable}[]{@{}
  >{\raggedright\arraybackslash}p{(\linewidth - 14\tabcolsep) * \real{0.1250}}
  >{\raggedright\arraybackslash}p{(\linewidth - 14\tabcolsep) * \real{0.1250}}
  >{\raggedright\arraybackslash}p{(\linewidth - 14\tabcolsep) * \real{0.1250}}
  >{\raggedright\arraybackslash}p{(\linewidth - 14\tabcolsep) * \real{0.1250}}
  >{\raggedright\arraybackslash}p{(\linewidth - 14\tabcolsep) * \real{0.1250}}
  >{\raggedright\arraybackslash}p{(\linewidth - 14\tabcolsep) * \real{0.1250}}
  >{\raggedright\arraybackslash}p{(\linewidth - 14\tabcolsep) * \real{0.1250}}
  >{\raggedright\arraybackslash}p{(\linewidth - 14\tabcolsep) * \real{0.1250}}@{}}
\toprule\noalign{}
\begin{minipage}[b]{\linewidth}\raggedright
RUNSPEC
\end{minipage} & \begin{minipage}[b]{\linewidth}\raggedright
GRID
\end{minipage} & \begin{minipage}[b]{\linewidth}\raggedright
EDIT
\end{minipage} & \begin{minipage}[b]{\linewidth}\raggedright
PROPS
\end{minipage} & \begin{minipage}[b]{\linewidth}\raggedright
REGIONS
\end{minipage} & \begin{minipage}[b]{\linewidth}\raggedright
SOLUTION
\end{minipage} & \begin{minipage}[b]{\linewidth}\raggedright
SUMMARY
\end{minipage} & \begin{minipage}[b]{\linewidth}\raggedright
SCHEDULE
\end{minipage} \\
\midrule\noalign{}
\endhead
\bottomrule\noalign{}
\endlastfoot
\end{longtable}

\subsubsection{Description}\label{description-113}

The PARTTRAC keyword activates the Partitioned Tracer option and defines
the maximum number of partitioned tracers, the number of TRACERKP or
TRACERKM partitioning tables in the PROPS section, and the maximum
number of number of rows in the TRACERKP or TRACERKM partitioning
tables.

This keyword is not supported by OPM Flow but would change the results
if supported so the simulation will be stopped.

\subsection{PATHS -- Define Filename Directory Path
Aliases}\label{paths-define-filename-directory-path-aliases}

\begin{longtable}[]{@{}
  >{\raggedright\arraybackslash}p{(\linewidth - 14\tabcolsep) * \real{0.1250}}
  >{\raggedright\arraybackslash}p{(\linewidth - 14\tabcolsep) * \real{0.1250}}
  >{\raggedright\arraybackslash}p{(\linewidth - 14\tabcolsep) * \real{0.1250}}
  >{\raggedright\arraybackslash}p{(\linewidth - 14\tabcolsep) * \real{0.1250}}
  >{\raggedright\arraybackslash}p{(\linewidth - 14\tabcolsep) * \real{0.1250}}
  >{\raggedright\arraybackslash}p{(\linewidth - 14\tabcolsep) * \real{0.1250}}
  >{\raggedright\arraybackslash}p{(\linewidth - 14\tabcolsep) * \real{0.1250}}
  >{\raggedright\arraybackslash}p{(\linewidth - 14\tabcolsep) * \real{0.1250}}@{}}
\toprule\noalign{}
\begin{minipage}[b]{\linewidth}\raggedright
RUNSPEC
\end{minipage} & \begin{minipage}[b]{\linewidth}\raggedright
GRID
\end{minipage} & \begin{minipage}[b]{\linewidth}\raggedright
EDIT
\end{minipage} & \begin{minipage}[b]{\linewidth}\raggedright
PROPS
\end{minipage} & \begin{minipage}[b]{\linewidth}\raggedright
REGIONS
\end{minipage} & \begin{minipage}[b]{\linewidth}\raggedright
SOLUTION
\end{minipage} & \begin{minipage}[b]{\linewidth}\raggedright
SUMMARY
\end{minipage} & \begin{minipage}[b]{\linewidth}\raggedright
SCHEDULE
\end{minipage} \\
\midrule\noalign{}
\endhead
\bottomrule\noalign{}
\endlastfoot
\end{longtable}

\subsubsection{Description}\label{description-114}

PATHS allows the user to define alias directory filenames to avoid long
filenames with the INCLUDE, IMPORT, RESTART or GDFILE keywords. To use
the alias with the aforementioned keywords PATHS should be prefixed with
the \$ symbol.

\begin{longtable}[]{@{}
  >{\raggedright\arraybackslash}p{(\linewidth - 6\tabcolsep) * \real{0.2500}}
  >{\raggedright\arraybackslash}p{(\linewidth - 6\tabcolsep) * \real{0.2500}}
  >{\raggedright\arraybackslash}p{(\linewidth - 6\tabcolsep) * \real{0.2500}}
  >{\raggedright\arraybackslash}p{(\linewidth - 6\tabcolsep) * \real{0.2500}}@{}}
\toprule\noalign{}
\begin{minipage}[b]{\linewidth}\raggedright
No.
\end{minipage} & \begin{minipage}[b]{\linewidth}\raggedright
Name
\end{minipage} & \begin{minipage}[b]{\linewidth}\raggedright
Description
\end{minipage} & \begin{minipage}[b]{\linewidth}\raggedright
Default
\end{minipage} \\
\midrule\noalign{}
\endhead
\bottomrule\noalign{}
\endlastfoot
1 & ALIAS & A character string enclosed in quotes defining the alias. &
None \\
2 & DIRC & A character string enclosed in quotes defining the directory
filename. & None \\
Notes: & & & \\
\end{longtable}

\emph{Table 5.34: PATHS Keyword Description}

\subsubsection{Examples}\label{examples-4}

\begin{verbatim}
--
--       PATH       PATH
--       ALIAS      DIRECTORY FILENAME
PATHS
        'GRID'      '/DISK1/NORNE/2017/GRID-INCLUDES'                          /
        'SCHD'      '/DISK1/NORNE/2017/SCHD-INCLUDES'                          /
/
\end{verbatim}

The above example defines ``GRID'' and ``SCHD'' aliases in the RUNSPEC
section than can be used in the GRID and SCHEDULE sections of the input
deck. The next example shows how to use the ``GRID'' alias with the
INCLUDE keyword in the GRID section.

\begin{verbatim}
--
--       LOAD INCLUDE FILES
--
INCLUDE
         '$GRID/PORO.INC'                                                      /

INCLUDE
         '$GRID/PERMX.INC'                                                     /

INCLUDE
         '$GRID/NTG.INC'                                                       /
\end{verbatim}

Here the porosity, permeability and net-to-gross arrays are loaded in
the GRID section using the directory filename aliases declared in the
RUNSPEC section.

\subsection{PEDIMS -- Define Petro-Elastic Model Regions and Table
Dimensions}\label{pedims-define-petro-elastic-model-regions-and-table-dimensions}

\begin{longtable}[]{@{}
  >{\raggedright\arraybackslash}p{(\linewidth - 14\tabcolsep) * \real{0.1250}}
  >{\raggedright\arraybackslash}p{(\linewidth - 14\tabcolsep) * \real{0.1250}}
  >{\raggedright\arraybackslash}p{(\linewidth - 14\tabcolsep) * \real{0.1250}}
  >{\raggedright\arraybackslash}p{(\linewidth - 14\tabcolsep) * \real{0.1250}}
  >{\raggedright\arraybackslash}p{(\linewidth - 14\tabcolsep) * \real{0.1250}}
  >{\raggedright\arraybackslash}p{(\linewidth - 14\tabcolsep) * \real{0.1250}}
  >{\raggedright\arraybackslash}p{(\linewidth - 14\tabcolsep) * \real{0.1250}}
  >{\raggedright\arraybackslash}p{(\linewidth - 14\tabcolsep) * \real{0.1250}}@{}}
\toprule\noalign{}
\begin{minipage}[b]{\linewidth}\raggedright
RUNSPEC
\end{minipage} & \begin{minipage}[b]{\linewidth}\raggedright
GRID
\end{minipage} & \begin{minipage}[b]{\linewidth}\raggedright
EDIT
\end{minipage} & \begin{minipage}[b]{\linewidth}\raggedright
PROPS
\end{minipage} & \begin{minipage}[b]{\linewidth}\raggedright
REGIONS
\end{minipage} & \begin{minipage}[b]{\linewidth}\raggedright
SOLUTION
\end{minipage} & \begin{minipage}[b]{\linewidth}\raggedright
SUMMARY
\end{minipage} & \begin{minipage}[b]{\linewidth}\raggedright
SCHEDULE
\end{minipage} \\
\midrule\noalign{}
\endhead
\bottomrule\noalign{}
\endlastfoot
\end{longtable}

\subsubsection{Description}\label{description-115}

The PEDIMS keyword defines the number of petro-elastic regions to be
used with the PENUM keyword in the REGIONS section and the number of
rows in the PEGTAB0 to PEGTAB7 keywords, as well as the number of rows
in the PEKTAB0 to PEKTAB7 keywords in the PROPS section.

This keyword is not supported by OPM Flow but would change the results
if supported so the simulation will be stopped.

\subsection{PETOPTS -- Define Petrel and Generic Simulation File
Options}\label{petopts-define-petrel-and-generic-simulation-file-options}

\begin{longtable}[]{@{}
  >{\raggedright\arraybackslash}p{(\linewidth - 14\tabcolsep) * \real{0.1250}}
  >{\raggedright\arraybackslash}p{(\linewidth - 14\tabcolsep) * \real{0.1250}}
  >{\raggedright\arraybackslash}p{(\linewidth - 14\tabcolsep) * \real{0.1250}}
  >{\raggedright\arraybackslash}p{(\linewidth - 14\tabcolsep) * \real{0.1250}}
  >{\raggedright\arraybackslash}p{(\linewidth - 14\tabcolsep) * \real{0.1250}}
  >{\raggedright\arraybackslash}p{(\linewidth - 14\tabcolsep) * \real{0.1250}}
  >{\raggedright\arraybackslash}p{(\linewidth - 14\tabcolsep) * \real{0.1250}}
  >{\raggedright\arraybackslash}p{(\linewidth - 14\tabcolsep) * \real{0.1250}}@{}}
\toprule\noalign{}
\begin{minipage}[b]{\linewidth}\raggedright
RUNSPEC
\end{minipage} & \begin{minipage}[b]{\linewidth}\raggedright
GRID
\end{minipage} & \begin{minipage}[b]{\linewidth}\raggedright
EDIT
\end{minipage} & \begin{minipage}[b]{\linewidth}\raggedright
PROPS
\end{minipage} & \begin{minipage}[b]{\linewidth}\raggedright
REGIONS
\end{minipage} & \begin{minipage}[b]{\linewidth}\raggedright
SOLUTION
\end{minipage} & \begin{minipage}[b]{\linewidth}\raggedright
SUMMARY
\end{minipage} & \begin{minipage}[b]{\linewidth}\raggedright
SCHEDULE
\end{minipage} \\
\midrule\noalign{}
\endhead
\bottomrule\noalign{}
\endlastfoot
\end{longtable}

\subsubsection{Description}\label{description-116}

The PETOPTS keyword defines various Petrel and Generic Simulation
(*.GSG) file options.

This keyword is not supported by OPM Flow but would change the results
if supported so the simulation will be stopped.

\subsection{PIMTDIMS -- Define Well Productivity Scaling Table
Dimensions}\label{pimtdims-define-well-productivity-scaling-table-dimensions}

\begin{longtable}[]{@{}
  >{\raggedright\arraybackslash}p{(\linewidth - 14\tabcolsep) * \real{0.1250}}
  >{\raggedright\arraybackslash}p{(\linewidth - 14\tabcolsep) * \real{0.1250}}
  >{\raggedright\arraybackslash}p{(\linewidth - 14\tabcolsep) * \real{0.1250}}
  >{\raggedright\arraybackslash}p{(\linewidth - 14\tabcolsep) * \real{0.1250}}
  >{\raggedright\arraybackslash}p{(\linewidth - 14\tabcolsep) * \real{0.1250}}
  >{\raggedright\arraybackslash}p{(\linewidth - 14\tabcolsep) * \real{0.1250}}
  >{\raggedright\arraybackslash}p{(\linewidth - 14\tabcolsep) * \real{0.1250}}
  >{\raggedright\arraybackslash}p{(\linewidth - 14\tabcolsep) * \real{0.1250}}@{}}
\toprule\noalign{}
\begin{minipage}[b]{\linewidth}\raggedright
RUNSPEC
\end{minipage} & \begin{minipage}[b]{\linewidth}\raggedright
GRID
\end{minipage} & \begin{minipage}[b]{\linewidth}\raggedright
EDIT
\end{minipage} & \begin{minipage}[b]{\linewidth}\raggedright
PROPS
\end{minipage} & \begin{minipage}[b]{\linewidth}\raggedright
REGIONS
\end{minipage} & \begin{minipage}[b]{\linewidth}\raggedright
SOLUTION
\end{minipage} & \begin{minipage}[b]{\linewidth}\raggedright
SUMMARY
\end{minipage} & \begin{minipage}[b]{\linewidth}\raggedright
SCHEDULE
\end{minipage} \\
\midrule\noalign{}
\endhead
\bottomrule\noalign{}
\endlastfoot
\end{longtable}

\subsubsection{Description}\label{description-117}

PIMTDIMS keyword defines the maximum number of PIMULTAB tables and the
maximum number of entries (or rows) per PIMULTAB table. The PIMULTAB
keyword is used to define a well's productivity index factor as a
function of a well's producing water cut.

\begin{longtable}[]{@{}
  >{\raggedright\arraybackslash}p{(\linewidth - 6\tabcolsep) * \real{0.2500}}
  >{\raggedright\arraybackslash}p{(\linewidth - 6\tabcolsep) * \real{0.2500}}
  >{\raggedright\arraybackslash}p{(\linewidth - 6\tabcolsep) * \real{0.2500}}
  >{\raggedright\arraybackslash}p{(\linewidth - 6\tabcolsep) * \real{0.2500}}@{}}
\toprule\noalign{}
\begin{minipage}[b]{\linewidth}\raggedright
No.
\end{minipage} & \begin{minipage}[b]{\linewidth}\raggedright
Name
\end{minipage} & \begin{minipage}[b]{\linewidth}\raggedright
Description
\end{minipage} & \begin{minipage}[b]{\linewidth}\raggedright
Default
\end{minipage} \\
\midrule\noalign{}
\endhead
\bottomrule\noalign{}
\endlastfoot
1 & NTPIMT & A positive integer value that defines the maximum number of
PIMULTAB keywords defined in the input deck. & 0 \\
2 & NRPIMT & A positive integer value defining the maximum number of
entries (rows) in the PIMULTAB keyword. & 0 \\
Notes: & & & \\
\end{longtable}

\emph{Table 5.35: PIMTDIMS Keyword Description}

\subsubsection{Example}\label{example-88}

\begin{verbatim}
--
--       MAX     MAX
--       TABLES  ENTRIES
PIMTDIMS
         1       51                                                            /
\end{verbatim}

The above example defines that there is one PIMULTAB table with a
maximum number of 51 rows.

\subsection{PINTDIMS -- Define Polymer Molecular Weight Model Table
Dimensions}\label{pintdims-define-polymer-molecular-weight-model-table-dimensions}

\begin{longtable}[]{@{}
  >{\raggedright\arraybackslash}p{(\linewidth - 14\tabcolsep) * \real{0.1250}}
  >{\raggedright\arraybackslash}p{(\linewidth - 14\tabcolsep) * \real{0.1250}}
  >{\raggedright\arraybackslash}p{(\linewidth - 14\tabcolsep) * \real{0.1250}}
  >{\raggedright\arraybackslash}p{(\linewidth - 14\tabcolsep) * \real{0.1250}}
  >{\raggedright\arraybackslash}p{(\linewidth - 14\tabcolsep) * \real{0.1250}}
  >{\raggedright\arraybackslash}p{(\linewidth - 14\tabcolsep) * \real{0.1250}}
  >{\raggedright\arraybackslash}p{(\linewidth - 14\tabcolsep) * \real{0.1250}}
  >{\raggedright\arraybackslash}p{(\linewidth - 14\tabcolsep) * \real{0.1250}}@{}}
\toprule\noalign{}
\begin{minipage}[b]{\linewidth}\raggedright
RUNSPEC
\end{minipage} & \begin{minipage}[b]{\linewidth}\raggedright
GRID
\end{minipage} & \begin{minipage}[b]{\linewidth}\raggedright
EDIT
\end{minipage} & \begin{minipage}[b]{\linewidth}\raggedright
PROPS
\end{minipage} & \begin{minipage}[b]{\linewidth}\raggedright
REGIONS
\end{minipage} & \begin{minipage}[b]{\linewidth}\raggedright
SOLUTION
\end{minipage} & \begin{minipage}[b]{\linewidth}\raggedright
SUMMARY
\end{minipage} & \begin{minipage}[b]{\linewidth}\raggedright
SCHEDULE
\end{minipage} \\
\midrule\noalign{}
\endhead
\bottomrule\noalign{}
\endlastfoot
\end{longtable}

\subsubsection{Description}\label{description-118}

The PINTDIMS keyword defines the number of property tables used in the
OPM Flow's Polymer Molecular Weight Transport option, that uses the
polymer molecular weight in calculating the polymer viscosity, as well
as accounting for formation damage due to the water and polymer
injection, by adjusting the wellbore skin pressure. This keyword should
only be used if the POLYMER and POLYMW keywords in the RUNSPEC section
are also activated. The PINTDIMS keyword defines the maximum number of
tables for the SKPRWAT, SKPRPOLY, and PLYMWINJ keywords, and the number
of entries in the PLYVMH keyword. All the aforementioned keywords are in
the PROPS section.

\begin{tcolorbox}[enhanced jigsaw, title=\textcolor{quarto-callout-note-color}{\faInfo}\hspace{0.5em}{Note}, coltitle=black, breakable, toprule=.15mm, colbacktitle=quarto-callout-note-color!10!white, leftrule=.75mm, opacitybacktitle=0.6, colback=white, bottomtitle=1mm, toptitle=1mm, titlerule=0mm, left=2mm, colframe=quarto-callout-note-color-frame, arc=.35mm, rightrule=.15mm, bottomrule=.15mm, opacityback=0]

This is an OPM Flow specific keyword that employs an alternative polymer
flood model based on a Polymer Molecular Weight Transport equation, that
is not available in the commercial simulator.

\end{tcolorbox}

\begin{longtable}[]{@{}
  >{\raggedright\arraybackslash}p{(\linewidth - 6\tabcolsep) * \real{0.2500}}
  >{\raggedright\arraybackslash}p{(\linewidth - 6\tabcolsep) * \real{0.2500}}
  >{\raggedright\arraybackslash}p{(\linewidth - 6\tabcolsep) * \real{0.2500}}
  >{\raggedright\arraybackslash}p{(\linewidth - 6\tabcolsep) * \real{0.2500}}@{}}
\toprule\noalign{}
\begin{minipage}[b]{\linewidth}\raggedright
No.
\end{minipage} & \begin{minipage}[b]{\linewidth}\raggedright
Name
\end{minipage} & \begin{minipage}[b]{\linewidth}\raggedright
Description
\end{minipage} & \begin{minipage}[b]{\linewidth}\raggedright
Default
\end{minipage} \\
\midrule\noalign{}
\endhead
\bottomrule\noalign{}
\endlastfoot
1 & NTSKWAT & NTSKWAT is a positive integer that defines the number of
SKPRWAT tables in the PROPS section, used to describe the relationship
of wellbore skin pressure as a function of water throughput and water
velocity, for the simulator's Polymer Molecular Weight Transport option.
& 1 \\
2 & NTSKPOLY & NTSKPOLY is a positive integer that defines the number of
SKPRPOLY tables in the PROPS section, used to describe the relationship
of wellbore skin pressure as a function of polymer throughput and
polymer velocity, for the simulator's Polymer Molecular Weight Transport
option. & 1 \\
3 & NTPMWINJ & NTPMWINJ is a positive integer that defines the number of
PLYMWINJ tables in the PROPS section, used to describe the relationship
of the injected polymer molecular weight as a function of polymer
throughput and polymer velocity, for the simulator's Polymer Molecular
Weight Transport option. & 1 \\
4 & NPLYVMH & NPLYVMH is a positive integer that defines the maximum
number of entries (rows) in the PLYVMH table in the PROPS section, used
to describe the relationship of the injected polymer viscosity as a
function of polymer molecular weight and polymer concentration, for the
simulator's Polymer Molecular Weight Transport option. & 1 \\
Notes: & & & \\
\end{longtable}

\emph{Table 5.36: PINTDIMS Keyword Description}

The SKPRWAT, SKPRPOLY, PLYMWINJ, and PLYVMH keywords in the PROPS
section, are the additional keywords required for the Polymer Molecular
Weight Transport option. Note that the standard polymer property data
keywords: PLYROCK, PLYADS, PLYMAX, etc., are still required to fully
describe the polymer fluid.

\subsubsection{Example}\label{example-89}

\begin{verbatim}
--
--       POLYMER MOLECULAR WEIGHT TRANSPORT TABLES (OPM FLOW RUNSPEC KEYWORD)
--
--       NO.      NO.       NO.       NO.
--       NTSKWAT  NTSKPOLY  NTPMWINJ  NPLYVMH
PINTDIMS
         2        2         2         1                                        /
\end{verbatim}

The above example declares two SKPRWAT, SKPRPOLY, and PLYMWINJ keywords
in the PROPS section will be used, as well as the default value of one
for the number of rows in the PLYVMH keyword.

\subsection{POLYMER -- Activate the Polymer Phase in the
Model}\label{polymer-activate-the-polymer-phase-in-the-model}

\begin{longtable}[]{@{}
  >{\raggedright\arraybackslash}p{(\linewidth - 14\tabcolsep) * \real{0.1250}}
  >{\raggedright\arraybackslash}p{(\linewidth - 14\tabcolsep) * \real{0.1250}}
  >{\raggedright\arraybackslash}p{(\linewidth - 14\tabcolsep) * \real{0.1250}}
  >{\raggedright\arraybackslash}p{(\linewidth - 14\tabcolsep) * \real{0.1250}}
  >{\raggedright\arraybackslash}p{(\linewidth - 14\tabcolsep) * \real{0.1250}}
  >{\raggedright\arraybackslash}p{(\linewidth - 14\tabcolsep) * \real{0.1250}}
  >{\raggedright\arraybackslash}p{(\linewidth - 14\tabcolsep) * \real{0.1250}}
  >{\raggedright\arraybackslash}p{(\linewidth - 14\tabcolsep) * \real{0.1250}}@{}}
\toprule\noalign{}
\begin{minipage}[b]{\linewidth}\raggedright
RUNSPEC
\end{minipage} & \begin{minipage}[b]{\linewidth}\raggedright
GRID
\end{minipage} & \begin{minipage}[b]{\linewidth}\raggedright
EDIT
\end{minipage} & \begin{minipage}[b]{\linewidth}\raggedright
PROPS
\end{minipage} & \begin{minipage}[b]{\linewidth}\raggedright
REGIONS
\end{minipage} & \begin{minipage}[b]{\linewidth}\raggedright
SOLUTION
\end{minipage} & \begin{minipage}[b]{\linewidth}\raggedright
SUMMARY
\end{minipage} & \begin{minipage}[b]{\linewidth}\raggedright
SCHEDULE
\end{minipage} \\
\midrule\noalign{}
\endhead
\bottomrule\noalign{}
\endlastfoot
\end{longtable}

\subsubsection{Description}\label{description-119}

This keyword indicates that the polymer phase is present in the model
and to activate the polymer flooding model. The keyword will also invoke
data input file checking to ensure that all the required polymer phase
input parameters are defined in the input deck.

There is no data required for this keyword and there is no terminating
``/'' for this keyword.

\subsubsection{Example}\label{example-90}

\begin{verbatim}
--
--       ACTIVATE THE POLYMER PHASE IN THE MODEL
--
POLYMER
\end{verbatim}

The above example declares that the polymer phase is active in the
model.

\subsection{POLYMW -- Activate the Polymer Molecular Weight Transport
Option}\label{polymw-activate-the-polymer-molecular-weight-transport-option}

\begin{longtable}[]{@{}
  >{\raggedright\arraybackslash}p{(\linewidth - 14\tabcolsep) * \real{0.1250}}
  >{\raggedright\arraybackslash}p{(\linewidth - 14\tabcolsep) * \real{0.1250}}
  >{\raggedright\arraybackslash}p{(\linewidth - 14\tabcolsep) * \real{0.1250}}
  >{\raggedright\arraybackslash}p{(\linewidth - 14\tabcolsep) * \real{0.1250}}
  >{\raggedright\arraybackslash}p{(\linewidth - 14\tabcolsep) * \real{0.1250}}
  >{\raggedright\arraybackslash}p{(\linewidth - 14\tabcolsep) * \real{0.1250}}
  >{\raggedright\arraybackslash}p{(\linewidth - 14\tabcolsep) * \real{0.1250}}
  >{\raggedright\arraybackslash}p{(\linewidth - 14\tabcolsep) * \real{0.1250}}@{}}
\toprule\noalign{}
\begin{minipage}[b]{\linewidth}\raggedright
RUNSPEC
\end{minipage} & \begin{minipage}[b]{\linewidth}\raggedright
GRID
\end{minipage} & \begin{minipage}[b]{\linewidth}\raggedright
EDIT
\end{minipage} & \begin{minipage}[b]{\linewidth}\raggedright
PROPS
\end{minipage} & \begin{minipage}[b]{\linewidth}\raggedright
REGIONS
\end{minipage} & \begin{minipage}[b]{\linewidth}\raggedright
SOLUTION
\end{minipage} & \begin{minipage}[b]{\linewidth}\raggedright
SUMMARY
\end{minipage} & \begin{minipage}[b]{\linewidth}\raggedright
SCHEDULE
\end{minipage} \\
\midrule\noalign{}
\endhead
\bottomrule\noalign{}
\endlastfoot
\end{longtable}

\subsubsection{Description}\label{description-120}

This keyword activates OPM Flow's Polymer Molecular Weight Transport
option that uses the polymer molecular weight in calculating the polymer
viscosity, as well as accounting for formation damage due to the water
and polymer injection, by adjusting the wellbore skin pressure. This
keyword should only be used if the POLYMER keyword in the RUNSPEC
section is also activated. The keyword will also invoke data input file
checking to ensure that all the required polymer phase input parameters
are defined in the input deck.

In the PROPS section the PLYVMH, PLYMWINJ, SKPRWAT and SKPRPOLY keywords
are used to define the additional polymer and water properties for the
model. And the SPOLYMW keyword in the SOLUTION section may be used to
set the initial molecular weight for each grid block in the model, in
order to initialize the model. In the SCHEDULE section, the WPMITAB
keyword should be used to associate a water injection well with the
appropriate PLYMWINJ table. Finally, the WSKPTAB keyword, also in the
SCHEDULE section, may be used to define the SKPRPOLY and SKPRWAT tables
associated with the skin pressure for polymer and water injection. All
the aforementioned keywords are OPM Flow specific keywords, that are not
available in the commercial simulator.

\begin{tcolorbox}[enhanced jigsaw, title=\textcolor{quarto-callout-note-color}{\faInfo}\hspace{0.5em}{Note}, coltitle=black, breakable, toprule=.15mm, colbacktitle=quarto-callout-note-color!10!white, leftrule=.75mm, opacitybacktitle=0.6, colback=white, bottomtitle=1mm, toptitle=1mm, titlerule=0mm, left=2mm, colframe=quarto-callout-note-color-frame, arc=.35mm, rightrule=.15mm, bottomrule=.15mm, opacityback=0]

This is an OPM Flow specific keyword that employs an alternative polymer
flood model based on a Polymer Molecular Weight Transport equation, that
is not available in the commercial simulator.

\end{tcolorbox}

There is no data required for this keyword and there is no terminating
``/'' for this keyword.

\subsubsection{Example}\label{example-91}

\begin{verbatim}
--
--       ACTIVATE THE POLYMER PHASE IN THE MODEL
--
POLYMER
--
--       ACTIVATE THE POLYMER MOLECULAR WEIGHT TRANSPORT OPTION
--
POLYMW
\end{verbatim}

The above example declares that the polymer phase is present, and that
the simulator's Polymer Molecular Weight Transport option should be used
instead of the standard polymer model.

\subsection{PRECSALT -- Activate the OPM Flow Salt Precipitation
Model}\label{precsalt-activate-the-opm-flow-salt-precipitation-model}

\begin{longtable}[]{@{}
  >{\raggedright\arraybackslash}p{(\linewidth - 14\tabcolsep) * \real{0.1250}}
  >{\raggedright\arraybackslash}p{(\linewidth - 14\tabcolsep) * \real{0.1250}}
  >{\raggedright\arraybackslash}p{(\linewidth - 14\tabcolsep) * \real{0.1250}}
  >{\raggedright\arraybackslash}p{(\linewidth - 14\tabcolsep) * \real{0.1250}}
  >{\raggedright\arraybackslash}p{(\linewidth - 14\tabcolsep) * \real{0.1250}}
  >{\raggedright\arraybackslash}p{(\linewidth - 14\tabcolsep) * \real{0.1250}}
  >{\raggedright\arraybackslash}p{(\linewidth - 14\tabcolsep) * \real{0.1250}}
  >{\raggedright\arraybackslash}p{(\linewidth - 14\tabcolsep) * \real{0.1250}}@{}}
\toprule\noalign{}
\begin{minipage}[b]{\linewidth}\raggedright
RUNSPEC
\end{minipage} & \begin{minipage}[b]{\linewidth}\raggedright
GRID
\end{minipage} & \begin{minipage}[b]{\linewidth}\raggedright
EDIT
\end{minipage} & \begin{minipage}[b]{\linewidth}\raggedright
PROPS
\end{minipage} & \begin{minipage}[b]{\linewidth}\raggedright
REGIONS
\end{minipage} & \begin{minipage}[b]{\linewidth}\raggedright
SOLUTION
\end{minipage} & \begin{minipage}[b]{\linewidth}\raggedright
SUMMARY
\end{minipage} & \begin{minipage}[b]{\linewidth}\raggedright
SCHEDULE
\end{minipage} \\
\midrule\noalign{}
\endhead
\bottomrule\noalign{}
\endlastfoot
\end{longtable}

\subsubsection{Description}\label{description-121}

This keyword activates the OPM Flow Salt Precipitation model that
accounts for salt precipitating out of the water phase when the water is
being vaporized into the gas phase and the dissolved salt reaches the
solubility limit as the pressure in the reservoir is being depleted (see
the VAPWAT keyword in the RUNSPEC section). This facility is an
extension to the standard Brine model, and as such the BRINE keyword in
the RUNSPEC must also be present in the input deck. In general, if the
PRECSALT keyword has been activated in the input deck then the VAPWAT
keyword should also be activated. The keyword should only be used if
both water and gas phases are active in the model.

\begin{tcolorbox}[enhanced jigsaw, title=\textcolor{quarto-callout-note-color}{\faInfo}\hspace{0.5em}{Note}, coltitle=black, breakable, toprule=.15mm, colbacktitle=quarto-callout-note-color!10!white, leftrule=.75mm, opacitybacktitle=0.6, colback=white, bottomtitle=1mm, toptitle=1mm, titlerule=0mm, left=2mm, colframe=quarto-callout-note-color-frame, arc=.35mm, rightrule=.15mm, bottomrule=.15mm, opacityback=0]

This is an OPM Flow specific keyword for the simulator's Salt
Precipitation model, note that this is an extension to the commercial
simulator's Brine model.

\end{tcolorbox}

If the keyword is present in the input deck then the SALTSOL keyword in
the PROPS section also needs to be present in the input deck to define
the salt solubility. In addition, either the SALTPVD or SALTP keywords
in the SOLUTION section should be used to define the initial salt
precipitated saturation.

There is no data required for this keyword and there is no terminating
``/'' for this keyword.

\subsubsection{Example}\label{example-92}

The first part of the example shows the keywords for the RUNSPEC
section.

\begin{verbatim}
-- ==============================================================================
--
-- RUNSPEC SECTION
--
-- ==============================================================================
RUNSPEC
       -- ------------------------------------------------------------------------------
-- FLUID TYPES AND TRACER OPTIONS
-- ------------------------------------------------------------------------------
--
--       OIL PHASE IS PRESENT IN THE RUN
--
OIL
--
--       WATER PHASE IS PRESENT IN THE RUN
--
WATER
--
--       GAS PHASE IS PRESENT IN THE RUN
--
GAS
--
--       DISSOLVED GAS IN LIVE OIL IS PRESENT IN THE RUN
--
DISGAS

--
--       ACTIVATE STANDARD BRINE MODEL IN THE RUN
--
BRINE
--
--       ACTIVATE THE OPM FLOW SALT PRECIPITATION MODEL (OPM FLOW KEYWORD)
--
PRECSALT
\end{verbatim}

The above example declares that the oil, water, gas, dissolved gas, and
brine phases are present in the model, and activates the Brine (BRINE
keyword) and Salt Precipitation (PRECSALT keyword) models.

The second part of the example shows the additional PVT fluid keywords
require for the Salt Precipitation model in the PROPS section. Note that
in addition to the Salt Precipitation model specific keywords, the
standard DENSITY (or GRAVITY), PVDG, and PVTO keywords are required.

\begin{verbatim}
-- ==============================================================================
--
-- PROPS SECTION
--
-- ==============================================================================
PROPS
--
--       WATER SALT PVT TABLE
--
PVTWSALT
--       REF PRES  REF SALT
--       PSIA      LB/STB
--       --------  --------
         3000.0    0.000                                   / REFERENCE DATA
--
--       SALTCONC  BW         CW        VISC     VISC
--       LB/STB    RB/STB     1/PSIA    CPOISE   GRAD
--       --------  --------   -------   ------   ------
            0.0    1.0        3.0E-6    1.0      1.0E-6
           10.0    0.8        3.0E-6    1.0      1.0E-6    / SALT CONCENTRATION
--
--       SET SALT SOLUBILITY LIMIT FOR ALL CELLS (OPM FLOW KEYWORD)
--
SALTSOL
--       MAX       SALT
--       SALTSOL   DENSITY
         0.500     135.46867445023383                      /
\end{verbatim}

In addition, to the above the PERMFACT keyword may be used to account
for the reduction in porosity and permeability due to salt precipitating
into the pore space. Again, in addition to the Salt Precipitation model
specific keywords, the standard ROCK, SGOF, and SWOF keywords are
required.

\begin{verbatim}
--
--       PERMEABILITY FACTOR REDUCTION DUE TO SALT (OPM FLOW KEYWORD)
--
PERMFACT
--       PORO       PERM
--       FACTOR     FACTOR
--       -------    --------
         0.0         0.000
         0.4         0.005
         0.6         0.010
         0.9         0.100
         1.0         1.000                                 / TABLE NO. 01
/
\end{verbatim}

In order to initialize the model in the SOLUTION section, apart from the
standard equilibration keywords, the following Salt Precipitation model
specific keywords should be included.

\begin{verbatim}
-- ==============================================================================
--
-- SOLUTION SECTION
--
-- ==============================================================================
SOLUTION
--
--       PRECIPITATED SALT VOLUME FRACTION VERSUS DEPTH (OPM FLOW KEYWORD)
--
--       DEPTH    SALTPSAT
--       ------   --------
SALTPVD
         8300.0    0.010
         8450.0    0.010                        / SALT VS DEPTH EQUIL REGN 01
--
--       DEPTH    SALT-1    SALT-2    SALT-3
--                SALTCON   SALTCON   SALTCON
--       ------   -------   -------   -------
SALTVD
         8300.0    0.500
         8450.0    0.500                         / SALT VS DEPTH EQUIL REGN 01
\end{verbatim}

Note that OPM Flow does not support Multi-Component Brine model, thus
there should only one column of salt concentrations.

Finally, in the SCHEDULE section, one may optionally one can add the
injections well's injection salt concentrations, as shown below.

\begin{verbatim}
-- ==============================================================================
--
-- SCHEDULE SECTION
--
-- ==============================================================================
SCHEDULE
--
--       DEFINE WATER INJECTION WELL SALT CONCENTRATIONS
--
-- WELL  SALT-1     SALT-2     SALT-3     SALT-4
-- NAME  SALTCON    SALTCON    SALTCON    SALTCON
--       -------    -------    -------    --------
WSALT
'INJ'    0.0000                                            /
/
\end{verbatim}

Normally, with this option one would also include vaporized water via
the VAPWAT keyword in the RUNSPEC section, as well as the PVTGW keyword
to define gas PVT properties for dry gas, or the PVTGWO keyword that
declares the gas PVT properties for wet gas. Both keywords are in the
PROPS section. The RWGSALT keyword in the PROPS section, may be used to
define the relationship of water vaporization versus pressure and salt
concentration. Again, the keyword is in the PROPS section. Secondly, in
the SOLUTION section, the RVW keyword can be used to set the initial
equilibration vaporized water in gas ratio values for all grid cells in
the model for Enumeration Initialization, or the RVWVD keyword that
declares the vaporized water-gas ratio (Rvw) versus depth tables for
each equilibration region, for Equilibrium Initialization.

\subsection{PSTEADY -- Activate Pseudo Steady State Flow Calculation
Option}\label{psteady-activate-pseudo-steady-state-flow-calculation-option}

\begin{longtable}[]{@{}
  >{\raggedright\arraybackslash}p{(\linewidth - 14\tabcolsep) * \real{0.1250}}
  >{\raggedright\arraybackslash}p{(\linewidth - 14\tabcolsep) * \real{0.1250}}
  >{\raggedright\arraybackslash}p{(\linewidth - 14\tabcolsep) * \real{0.1250}}
  >{\raggedright\arraybackslash}p{(\linewidth - 14\tabcolsep) * \real{0.1250}}
  >{\raggedright\arraybackslash}p{(\linewidth - 14\tabcolsep) * \real{0.1250}}
  >{\raggedright\arraybackslash}p{(\linewidth - 14\tabcolsep) * \real{0.1250}}
  >{\raggedright\arraybackslash}p{(\linewidth - 14\tabcolsep) * \real{0.1250}}
  >{\raggedright\arraybackslash}p{(\linewidth - 14\tabcolsep) * \real{0.1250}}@{}}
\toprule\noalign{}
\begin{minipage}[b]{\linewidth}\raggedright
RUNSPEC
\end{minipage} & \begin{minipage}[b]{\linewidth}\raggedright
GRID
\end{minipage} & \begin{minipage}[b]{\linewidth}\raggedright
EDIT
\end{minipage} & \begin{minipage}[b]{\linewidth}\raggedright
PROPS
\end{minipage} & \begin{minipage}[b]{\linewidth}\raggedright
REGIONS
\end{minipage} & \begin{minipage}[b]{\linewidth}\raggedright
SOLUTION
\end{minipage} & \begin{minipage}[b]{\linewidth}\raggedright
SUMMARY
\end{minipage} & \begin{minipage}[b]{\linewidth}\raggedright
SCHEDULE
\end{minipage} \\
\midrule\noalign{}
\endhead
\bottomrule\noalign{}
\endlastfoot
\end{longtable}

\subsubsection{Description}\label{description-122}

The PSTEADY keyword activates Pseudo Steady State Flow Calculation
option by advancing the simulator until it reaches a pseudo steady state
flow and then sets the date to the date defined on this keyword, that is
written to the RESTART file. Keyword also includes parameters defining
the conditions for pseudo steady flow state.

This keyword is not supported by OPM Flow but would change the results
if supported so the simulation will be stopped.

\subsection{RADIAL -- Radial Grid Activation
Option}\label{radial-radial-grid-activation-option}

\begin{longtable}[]{@{}
  >{\raggedright\arraybackslash}p{(\linewidth - 14\tabcolsep) * \real{0.1250}}
  >{\raggedright\arraybackslash}p{(\linewidth - 14\tabcolsep) * \real{0.1250}}
  >{\raggedright\arraybackslash}p{(\linewidth - 14\tabcolsep) * \real{0.1250}}
  >{\raggedright\arraybackslash}p{(\linewidth - 14\tabcolsep) * \real{0.1250}}
  >{\raggedright\arraybackslash}p{(\linewidth - 14\tabcolsep) * \real{0.1250}}
  >{\raggedright\arraybackslash}p{(\linewidth - 14\tabcolsep) * \real{0.1250}}
  >{\raggedright\arraybackslash}p{(\linewidth - 14\tabcolsep) * \real{0.1250}}
  >{\raggedright\arraybackslash}p{(\linewidth - 14\tabcolsep) * \real{0.1250}}@{}}
\toprule\noalign{}
\begin{minipage}[b]{\linewidth}\raggedright
RUNSPEC
\end{minipage} & \begin{minipage}[b]{\linewidth}\raggedright
GRID
\end{minipage} & \begin{minipage}[b]{\linewidth}\raggedright
EDIT
\end{minipage} & \begin{minipage}[b]{\linewidth}\raggedright
PROPS
\end{minipage} & \begin{minipage}[b]{\linewidth}\raggedright
REGIONS
\end{minipage} & \begin{minipage}[b]{\linewidth}\raggedright
SOLUTION
\end{minipage} & \begin{minipage}[b]{\linewidth}\raggedright
SUMMARY
\end{minipage} & \begin{minipage}[b]{\linewidth}\raggedright
SCHEDULE
\end{minipage} \\
\midrule\noalign{}
\endhead
\bottomrule\noalign{}
\endlastfoot
\end{longtable}

\subsubsection{Description}\label{description-123}

RADIAL activates the radial grid geometry option for the model, if this
keyword and the SPIDER keyword are omitted then Cartesian geometry is
assumed by OPM Flow.

See also the SPIDER keyword in the RUNSPEC section (SPIDER -- Spider
Grid Activation Option). This is an OPM Flow specific keyword for the
simulator's Spider Grid option, that emulates a Radial Grid via
corner-point geometry. Both Spider and Radial grids will be displayed
the same in OPM ResInsight, as they both use Irregular Corner-Point Grid
geometry. The difference is that the radial model's pore volumes are
adjusted to match the radial geometry pore volumes, whereas, the Spider
grid volumes are not adjusted. See the examples in the section on Radial
Grids (Radial Grids) and Spider Grids (Spider Grids) for further
information.

\subsubsection{Example}\label{example-93}

\begin{verbatim}
--
--       DEFINE RADIAL GRID GEOMETRY
--
RADIAL
\end{verbatim}

The example activates the radial grid geometry option.

\subsection{REGDIMS -- Define the Maximum Number of Regions for a Region
Array}\label{regdims-define-the-maximum-number-of-regions-for-a-region-array}

\begin{longtable}[]{@{}
  >{\raggedright\arraybackslash}p{(\linewidth - 14\tabcolsep) * \real{0.1250}}
  >{\raggedright\arraybackslash}p{(\linewidth - 14\tabcolsep) * \real{0.1250}}
  >{\raggedright\arraybackslash}p{(\linewidth - 14\tabcolsep) * \real{0.1250}}
  >{\raggedright\arraybackslash}p{(\linewidth - 14\tabcolsep) * \real{0.1250}}
  >{\raggedright\arraybackslash}p{(\linewidth - 14\tabcolsep) * \real{0.1250}}
  >{\raggedright\arraybackslash}p{(\linewidth - 14\tabcolsep) * \real{0.1250}}
  >{\raggedright\arraybackslash}p{(\linewidth - 14\tabcolsep) * \real{0.1250}}
  >{\raggedright\arraybackslash}p{(\linewidth - 14\tabcolsep) * \real{0.1250}}@{}}
\toprule\noalign{}
\begin{minipage}[b]{\linewidth}\raggedright
RUNSPEC
\end{minipage} & \begin{minipage}[b]{\linewidth}\raggedright
GRID
\end{minipage} & \begin{minipage}[b]{\linewidth}\raggedright
EDIT
\end{minipage} & \begin{minipage}[b]{\linewidth}\raggedright
PROPS
\end{minipage} & \begin{minipage}[b]{\linewidth}\raggedright
REGIONS
\end{minipage} & \begin{minipage}[b]{\linewidth}\raggedright
SOLUTION
\end{minipage} & \begin{minipage}[b]{\linewidth}\raggedright
SUMMARY
\end{minipage} & \begin{minipage}[b]{\linewidth}\raggedright
SCHEDULE
\end{minipage} \\
\midrule\noalign{}
\endhead
\bottomrule\noalign{}
\endlastfoot
\end{longtable}

\subsubsection{Description}\label{description-124}

The REGDIMS keyword defines the maximum number of regions for various
region arrays used in the model. Note that the maximum number of FIPNUM
regions can be defined both on this keyword and the TABDIMS keyword, if
it set in both locations the maximum value is used. The reason for this
type of inconsistency is due to the commercial simulator evolving with
time as new features were added, but at the same time having to maintain
backward input deck compatibility.

\begin{longtable}[]{@{}
  >{\raggedright\arraybackslash}p{(\linewidth - 6\tabcolsep) * \real{0.2500}}
  >{\raggedright\arraybackslash}p{(\linewidth - 6\tabcolsep) * \real{0.2500}}
  >{\raggedright\arraybackslash}p{(\linewidth - 6\tabcolsep) * \real{0.2500}}
  >{\raggedright\arraybackslash}p{(\linewidth - 6\tabcolsep) * \real{0.2500}}@{}}
\toprule\noalign{}
\begin{minipage}[b]{\linewidth}\raggedright
No.
\end{minipage} & \begin{minipage}[b]{\linewidth}\raggedright
Name
\end{minipage} & \begin{minipage}[b]{\linewidth}\raggedright
Description
\end{minipage} & \begin{minipage}[b]{\linewidth}\raggedright
Default
\end{minipage} \\
\midrule\noalign{}
\endhead
\bottomrule\noalign{}
\endlastfoot
1 & NTFIP & A positive integer defining the maximum number of regions in
the FIPNUM region array. If additional sets of fluid in-place regions
have been defined, as per the FIPxxxxx series of fluid in-place region
keywords, then NTFIP should be set to the maximum number of regions in
either the FIPNUM or FIPxxxxx associated arrays. Thus, if the maximum
number of regions in the FIPNUM array is 12, and the maximum value in
the FIPxxxxx series of arrays is 20, then 20 should be entered for
NTFIP. Note that this parameter may also be set on the TABDIMS keyword
as well. If NTFIP is set in both places, then the maximum value is used.
& 1 \\
2 & NMFIPR & A positive integer defining the total maximum number of
fluid in-place regions. The number of FIPNUM regions are defined by
NTFIP. However, if additional sets of fluid in-place regions are
required, as per the FIPxxxxx series of fluid in-place region keywords,
then these are to be defined here by adding the number of FIPxxxxx
arrays to the value NTFIP. So for example, if NTFIP equals five and the
number of distinct FIPxxxxx regions is three, then the value to enter
for NMFIPR is eight. & 1 \\
3 & NRFREG & A positive integer defining the maximum number of
independent reservoir regions in the ISOLNUM region array. & 0 \\
4 & MXNFLX & A positive integer defining the maximum number of flux
regions in the FLUXNUM region array. MXNFLX can also be defined on the
TABDIMS keywords as well. If MXNFLX is defined both here and on the
TABDIMS keyword then the maximum value of the two is used. & 0 \\
5 & NUSREG & A positive integer defining the maximum user defined
regions in a commercial simulator's compositional model. This parameter
is included for compatibility and should be defaulted as it is not used
in OPM Flow. & 0 \\
6 & NTCREG & A positive integer defining the maximum number of regions
in the COALNUM region array. This parameter is included for
compatibility and should be defaulted as it is not used in OPM Flow. &
1 \\
7 & NOPREG & A positive integer defining the maximum number of regions
in the OPERNUM region array. & 0 \\
8 & NWKDREG & A positive integer defining the maximum number of real
double-precision work arrays for use with the OPERATE and OPERATER
keywords. This parameter is included for compatibility and should be
defaulted as it is not used in OPM Flow. & 0 \\
9 & NWKIREG & A positive integer defining the maximum number of integer
work arrays for use with the OPERATE and OPERATER keywords. This
parameter is included for compatibility and should be defaulted as it is
not used in OPM Flow. & 0 \\
10 & NPLMIX & A positive integer defining the maximum number of regions
in the PLMIXNUM region array. & 1 \\
Notes: & & & \\
\end{longtable}

\emph{Table 5.37: REGDIMS Keyword Description}

\subsubsection{Example}\label{example-94}

\begin{verbatim}
--
--       MAX     TOTAL   INDEP   FLUX    TRACK  CBM    OPERN  WORK  WORK  POLY
--       FIPNUM  REGNS   REGNS   REGNS   REGNS  REGNS  REGNS  REAL  INTG  REGNS
REGDIMS
         9       12      1*      1*      1*     1*     1*     1*    1*    1*   /
\end{verbatim}

The above example defines the number of FIPNUM regions to be nine and
the number of FIPxxx type of regions to be three (12 -- 9), the rest of
the region sizes are set to the default values.

\subsection{RIVRDIMS -- Define the River Dimensions and Associated
Data}\label{rivrdims-define-the-river-dimensions-and-associated-data}

\begin{longtable}[]{@{}
  >{\raggedright\arraybackslash}p{(\linewidth - 14\tabcolsep) * \real{0.1250}}
  >{\raggedright\arraybackslash}p{(\linewidth - 14\tabcolsep) * \real{0.1250}}
  >{\raggedright\arraybackslash}p{(\linewidth - 14\tabcolsep) * \real{0.1250}}
  >{\raggedright\arraybackslash}p{(\linewidth - 14\tabcolsep) * \real{0.1250}}
  >{\raggedright\arraybackslash}p{(\linewidth - 14\tabcolsep) * \real{0.1250}}
  >{\raggedright\arraybackslash}p{(\linewidth - 14\tabcolsep) * \real{0.1250}}
  >{\raggedright\arraybackslash}p{(\linewidth - 14\tabcolsep) * \real{0.1250}}
  >{\raggedright\arraybackslash}p{(\linewidth - 14\tabcolsep) * \real{0.1250}}@{}}
\toprule\noalign{}
\begin{minipage}[b]{\linewidth}\raggedright
RUNSPEC
\end{minipage} & \begin{minipage}[b]{\linewidth}\raggedright
GRID
\end{minipage} & \begin{minipage}[b]{\linewidth}\raggedright
EDIT
\end{minipage} & \begin{minipage}[b]{\linewidth}\raggedright
PROPS
\end{minipage} & \begin{minipage}[b]{\linewidth}\raggedright
REGIONS
\end{minipage} & \begin{minipage}[b]{\linewidth}\raggedright
SOLUTION
\end{minipage} & \begin{minipage}[b]{\linewidth}\raggedright
SUMMARY
\end{minipage} & \begin{minipage}[b]{\linewidth}\raggedright
SCHEDULE
\end{minipage} \\
\midrule\noalign{}
\endhead
\bottomrule\noalign{}
\endlastfoot
\end{longtable}

\subsubsection{Description}\label{description-125}

RIVRDIMS defines the river system array dimensions used with the REACHES
keyword and other river keywords in the SOLUTION and SCHEDULE sections.
The keyword also enables the River option.

This keyword is not supported by OPM Flow but would change the results
if supported so the simulation will be stopped.

\subsection{ROCKCOMP -- Activate Rock
Compaction}\label{rockcomp-activate-rock-compaction}

\begin{longtable}[]{@{}
  >{\raggedright\arraybackslash}p{(\linewidth - 14\tabcolsep) * \real{0.1250}}
  >{\raggedright\arraybackslash}p{(\linewidth - 14\tabcolsep) * \real{0.1250}}
  >{\raggedright\arraybackslash}p{(\linewidth - 14\tabcolsep) * \real{0.1250}}
  >{\raggedright\arraybackslash}p{(\linewidth - 14\tabcolsep) * \real{0.1250}}
  >{\raggedright\arraybackslash}p{(\linewidth - 14\tabcolsep) * \real{0.1250}}
  >{\raggedright\arraybackslash}p{(\linewidth - 14\tabcolsep) * \real{0.1250}}
  >{\raggedright\arraybackslash}p{(\linewidth - 14\tabcolsep) * \real{0.1250}}
  >{\raggedright\arraybackslash}p{(\linewidth - 14\tabcolsep) * \real{0.1250}}@{}}
\toprule\noalign{}
\begin{minipage}[b]{\linewidth}\raggedright
RUNSPEC
\end{minipage} & \begin{minipage}[b]{\linewidth}\raggedright
GRID
\end{minipage} & \begin{minipage}[b]{\linewidth}\raggedright
EDIT
\end{minipage} & \begin{minipage}[b]{\linewidth}\raggedright
PROPS
\end{minipage} & \begin{minipage}[b]{\linewidth}\raggedright
REGIONS
\end{minipage} & \begin{minipage}[b]{\linewidth}\raggedright
SOLUTION
\end{minipage} & \begin{minipage}[b]{\linewidth}\raggedright
SUMMARY
\end{minipage} & \begin{minipage}[b]{\linewidth}\raggedright
SCHEDULE
\end{minipage} \\
\midrule\noalign{}
\endhead
\bottomrule\noalign{}
\endlastfoot
\end{longtable}

\subsubsection{Description}\label{description-126}

The ROCKCOMP keyword activates rock compaction and defines various rock
compaction options for the run. By default OPM Flow models rock
compaction via pore volume compressibility as entered on the ROCK
keyword in the PROPS section. This keyword enables pressure dependent
pore volume and transmissibility multipliers for rock compaction that
are entered in the PROPS section using the ROCKTAB keyword. Currently
OPM Flow only supports the default options for rock compaction.

\begin{longtable}[]{@{}
  >{\raggedright\arraybackslash}p{(\linewidth - 6\tabcolsep) * \real{0.2500}}
  >{\raggedright\arraybackslash}p{(\linewidth - 6\tabcolsep) * \real{0.2500}}
  >{\raggedright\arraybackslash}p{(\linewidth - 6\tabcolsep) * \real{0.2500}}
  >{\raggedright\arraybackslash}p{(\linewidth - 6\tabcolsep) * \real{0.2500}}@{}}
\toprule\noalign{}
\begin{minipage}[b]{\linewidth}\raggedright
No.
\end{minipage} & \begin{minipage}[b]{\linewidth}\raggedright
Name
\end{minipage} & \begin{minipage}[b]{\linewidth}\raggedright
Description
\end{minipage} & \begin{minipage}[b]{\linewidth}\raggedright
Default
\end{minipage} \\
\midrule\noalign{}
\endhead
\bottomrule\noalign{}
\endlastfoot
1 & ROCKOPT & A character string that defines the rock compaction option
based on one of the following character strings: Only the REVERS and
IRREVERS options are supported by OPM Flow. & REVERS \\
2 & NTROCC & A positive integer that defines the number of rock
compaction tables, that is the number of ROCKTAB tables to be used by
OPM Flow. & 1 \\
3 & WATINOPT & A character string that states if the water induced rock
compaction option should be used (YES) or not (NO). If set to YES then
either the ROCKTABW or the ROCK2D and ROCKWNOD keywords should be
entered in the PROPS section. & NO \\
4 & PORTXROP & A character string that specifies the model to be used
for when transmissibility is dependent on porosity, and should be set to
either: This option is used in the commercial compositional simulator
and is therefore ignored by OPM Flow. & 1* \\
5 & CARKZEXP & The exponent constant in the Carmen-Kozeny
porosity-transmissibility equation for when PORTXROP has been set to CZ.
This option is used in the commercial compositional simulator and is
therefore ignored by OPM Flow. & 0.0 \\
Notes: & & & \\
\end{longtable}

\emph{Table 5.38: ROCKCOMP Keyword Description}

\subsubsection{Example}\label{example-95}

\begin{verbatim}
--
--       ROCK   NUMBER   WAT     POR-TRAN
--       OPTN   TABLES   INDUCE  OPTION
ROCKCOMP
         REVERS 5        NO      1*                                            /
\end{verbatim}

The above example defines the default values for the ROCKCOMP keyword
with five rock compaction tables.

\subsection{RPTCPL -- Activate Couple Simulation
Reporting}\label{rptcpl-activate-couple-simulation-reporting}

\begin{longtable}[]{@{}
  >{\raggedright\arraybackslash}p{(\linewidth - 14\tabcolsep) * \real{0.1250}}
  >{\raggedright\arraybackslash}p{(\linewidth - 14\tabcolsep) * \real{0.1250}}
  >{\raggedright\arraybackslash}p{(\linewidth - 14\tabcolsep) * \real{0.1250}}
  >{\raggedright\arraybackslash}p{(\linewidth - 14\tabcolsep) * \real{0.1250}}
  >{\raggedright\arraybackslash}p{(\linewidth - 14\tabcolsep) * \real{0.1250}}
  >{\raggedright\arraybackslash}p{(\linewidth - 14\tabcolsep) * \real{0.1250}}
  >{\raggedright\arraybackslash}p{(\linewidth - 14\tabcolsep) * \real{0.1250}}
  >{\raggedright\arraybackslash}p{(\linewidth - 14\tabcolsep) * \real{0.1250}}@{}}
\toprule\noalign{}
\begin{minipage}[b]{\linewidth}\raggedright
FRUNSPEC
\end{minipage} & \begin{minipage}[b]{\linewidth}\raggedright
GRID
\end{minipage} & \begin{minipage}[b]{\linewidth}\raggedright
EDIT
\end{minipage} & \begin{minipage}[b]{\linewidth}\raggedright
PROPS
\end{minipage} & \begin{minipage}[b]{\linewidth}\raggedright
REGIONS
\end{minipage} & \begin{minipage}[b]{\linewidth}\raggedright
SOLUTION
\end{minipage} & \begin{minipage}[b]{\linewidth}\raggedright
SUMMARY
\end{minipage} & \begin{minipage}[b]{\linewidth}\raggedright
SCHEDULE
\end{minipage} \\
\midrule\noalign{}
\endhead
\bottomrule\noalign{}
\endlastfoot
\end{longtable}

\subsubsection{Description}\label{description-127}

This keyword activates the couple simulation reporting, that results in
the simulator writing out various initialization data and simulation
data in order for external ``controlling programs'' to interactively
manage the simulation. There is no data required for this keyword but
the keyword should be terminated by a ``/''.

This keyword is not supported by OPM Flow but has no effect on the
results so it will be ignored.

\subsubsection{Example}\label{example-96}

\begin{verbatim}
--
--       ACTIVATE COUPLE SIMULATION REPORTING
--
RPTCPL
/
\end{verbatim}

The above example switches on couple simulation reporting; however, this
has no effect in OPM Flow input decks.

\subsection{RPTHMD -- Define Well History Match Gradient Reporting
Options}\label{rpthmd-define-well-history-match-gradient-reporting-options}

\begin{longtable}[]{@{}
  >{\raggedright\arraybackslash}p{(\linewidth - 14\tabcolsep) * \real{0.1250}}
  >{\raggedright\arraybackslash}p{(\linewidth - 14\tabcolsep) * \real{0.1250}}
  >{\raggedright\arraybackslash}p{(\linewidth - 14\tabcolsep) * \real{0.1250}}
  >{\raggedright\arraybackslash}p{(\linewidth - 14\tabcolsep) * \real{0.1250}}
  >{\raggedright\arraybackslash}p{(\linewidth - 14\tabcolsep) * \real{0.1250}}
  >{\raggedright\arraybackslash}p{(\linewidth - 14\tabcolsep) * \real{0.1250}}
  >{\raggedright\arraybackslash}p{(\linewidth - 14\tabcolsep) * \real{0.1250}}
  >{\raggedright\arraybackslash}p{(\linewidth - 14\tabcolsep) * \real{0.1250}}@{}}
\toprule\noalign{}
\begin{minipage}[b]{\linewidth}\raggedright
RUNSPEC
\end{minipage} & \begin{minipage}[b]{\linewidth}\raggedright
GRID
\end{minipage} & \begin{minipage}[b]{\linewidth}\raggedright
EDIT
\end{minipage} & \begin{minipage}[b]{\linewidth}\raggedright
PROPS
\end{minipage} & \begin{minipage}[b]{\linewidth}\raggedright
REGIONS
\end{minipage} & \begin{minipage}[b]{\linewidth}\raggedright
SOLUTION
\end{minipage} & \begin{minipage}[b]{\linewidth}\raggedright
SUMMARY
\end{minipage} & \begin{minipage}[b]{\linewidth}\raggedright
SCHEDULE
\end{minipage} \\
\midrule\noalign{}
\endhead
\bottomrule\noalign{}
\endlastfoot
\end{longtable}

\subsubsection{Description}\label{description-128}

This keyword, RPTHMD, defines the options and level of history match
output that should be written to history match file (*.HMD), for when
the History Match Gradient option has been activated by the HMDIMS
keyword in the RUNSPEC section.

This keyword is not supported by OPM Flow but has no effect on the
results so it will be ignored.

\subsection{RPTRUNSP -- Activate RUNSPEC
Reporting}\label{rptrunsp-activate-runspec-reporting}

\begin{longtable}[]{@{}
  >{\raggedright\arraybackslash}p{(\linewidth - 14\tabcolsep) * \real{0.1250}}
  >{\raggedright\arraybackslash}p{(\linewidth - 14\tabcolsep) * \real{0.1250}}
  >{\raggedright\arraybackslash}p{(\linewidth - 14\tabcolsep) * \real{0.1250}}
  >{\raggedright\arraybackslash}p{(\linewidth - 14\tabcolsep) * \real{0.1250}}
  >{\raggedright\arraybackslash}p{(\linewidth - 14\tabcolsep) * \real{0.1250}}
  >{\raggedright\arraybackslash}p{(\linewidth - 14\tabcolsep) * \real{0.1250}}
  >{\raggedright\arraybackslash}p{(\linewidth - 14\tabcolsep) * \real{0.1250}}
  >{\raggedright\arraybackslash}p{(\linewidth - 14\tabcolsep) * \real{0.1250}}@{}}
\toprule\noalign{}
\begin{minipage}[b]{\linewidth}\raggedright
RUNSPEC
\end{minipage} & \begin{minipage}[b]{\linewidth}\raggedright
GRID
\end{minipage} & \begin{minipage}[b]{\linewidth}\raggedright
EDIT
\end{minipage} & \begin{minipage}[b]{\linewidth}\raggedright
PROPS
\end{minipage} & \begin{minipage}[b]{\linewidth}\raggedright
REGIONS
\end{minipage} & \begin{minipage}[b]{\linewidth}\raggedright
SOLUTION
\end{minipage} & \begin{minipage}[b]{\linewidth}\raggedright
SUMMARY
\end{minipage} & \begin{minipage}[b]{\linewidth}\raggedright
SCHEDULE
\end{minipage} \\
\midrule\noalign{}
\endhead
\bottomrule\noalign{}
\endlastfoot
\end{longtable}

\subsubsection{Description}\label{description-129}

This keyword activates reporting of all the RUNSPEC options utilized in
the run. There is no data required for this keyword.

This keyword is not supported by OPM Flow but has no effect on the
results so it will be ignored.

\subsubsection{Example}\label{example-97}

\begin{verbatim}
--
--       ACTIVATE RUNSPEC SECTION REPORTING
--
RPTRUNSP
\end{verbatim}

The above example switches on RUNSPEC reporting; however, this has no
effect in OPM Flow input decks.

\subsection{RSSPEC -- Activate Output of the RESTART Index
File}\label{rsspec-activate-output-of-the-restart-index-file}

\begin{longtable}[]{@{}
  >{\raggedright\arraybackslash}p{(\linewidth - 14\tabcolsep) * \real{0.1250}}
  >{\raggedright\arraybackslash}p{(\linewidth - 14\tabcolsep) * \real{0.1250}}
  >{\raggedright\arraybackslash}p{(\linewidth - 14\tabcolsep) * \real{0.1250}}
  >{\raggedright\arraybackslash}p{(\linewidth - 14\tabcolsep) * \real{0.1250}}
  >{\raggedright\arraybackslash}p{(\linewidth - 14\tabcolsep) * \real{0.1250}}
  >{\raggedright\arraybackslash}p{(\linewidth - 14\tabcolsep) * \real{0.1250}}
  >{\raggedright\arraybackslash}p{(\linewidth - 14\tabcolsep) * \real{0.1250}}
  >{\raggedright\arraybackslash}p{(\linewidth - 14\tabcolsep) * \real{0.1250}}@{}}
\toprule\noalign{}
\begin{minipage}[b]{\linewidth}\raggedright
RUNSPEC
\end{minipage} & \begin{minipage}[b]{\linewidth}\raggedright
GRID
\end{minipage} & \begin{minipage}[b]{\linewidth}\raggedright
EDIT
\end{minipage} & \begin{minipage}[b]{\linewidth}\raggedright
PROPS
\end{minipage} & \begin{minipage}[b]{\linewidth}\raggedright
REGIONS
\end{minipage} & \begin{minipage}[b]{\linewidth}\raggedright
SOLUTION
\end{minipage} & \begin{minipage}[b]{\linewidth}\raggedright
SUMMARY
\end{minipage} & \begin{minipage}[b]{\linewidth}\raggedright
SCHEDULE
\end{minipage} \\
\midrule\noalign{}
\endhead
\bottomrule\noalign{}
\endlastfoot
\end{longtable}

\subsubsection{Description}\label{description-130}

The RSSPEC keyword activates the writing out of the RESTART index file
(\emph{.RSSPEC). The restart data (pressure, saturations etc. through
time for each active cell) are written out to two files one file
contains the data, }.UNRST for example, and the second file contains an
index of the data (\emph{.RSSPEC) stored in the }.UNRST file. This
keyword is somewhat redundant as the RESTART index file is written out
by default. See the NORSSPEC keyword in the RUNSPEC section that
deactivates the writing out of the file.

Hence, OPM Flow ignores this keyword. It is documented here for
completeness.

There is no data required for this keyword and there is no terminating
``/'' for this keyword.

\subsubsection{Example}\label{example-98}

\begin{verbatim}
--
--       ACTIVATE OUTPUT OF THE RESTART INDEX FILE *.RSSPEC
--
RSSPEC
\end{verbatim}

The above example switches on the writing of the restart index file
(*.RSSPEC); however, this has no effect in OPM Flow input decks.

\subsection{RUNSPEC -- Define the Start of the RUNSPEC Section of
Keywords}\label{runspec-define-the-start-of-the-runspec-section-of-keywords}

\begin{longtable}[]{@{}
  >{\raggedright\arraybackslash}p{(\linewidth - 14\tabcolsep) * \real{0.1250}}
  >{\raggedright\arraybackslash}p{(\linewidth - 14\tabcolsep) * \real{0.1250}}
  >{\raggedright\arraybackslash}p{(\linewidth - 14\tabcolsep) * \real{0.1250}}
  >{\raggedright\arraybackslash}p{(\linewidth - 14\tabcolsep) * \real{0.1250}}
  >{\raggedright\arraybackslash}p{(\linewidth - 14\tabcolsep) * \real{0.1250}}
  >{\raggedright\arraybackslash}p{(\linewidth - 14\tabcolsep) * \real{0.1250}}
  >{\raggedright\arraybackslash}p{(\linewidth - 14\tabcolsep) * \real{0.1250}}
  >{\raggedright\arraybackslash}p{(\linewidth - 14\tabcolsep) * \real{0.1250}}@{}}
\toprule\noalign{}
\begin{minipage}[b]{\linewidth}\raggedright
RUNSPEC
\end{minipage} & \begin{minipage}[b]{\linewidth}\raggedright
GRID
\end{minipage} & \begin{minipage}[b]{\linewidth}\raggedright
EDIT
\end{minipage} & \begin{minipage}[b]{\linewidth}\raggedright
PROPS
\end{minipage} & \begin{minipage}[b]{\linewidth}\raggedright
REGIONS
\end{minipage} & \begin{minipage}[b]{\linewidth}\raggedright
SOLUTION
\end{minipage} & \begin{minipage}[b]{\linewidth}\raggedright
SUMMARY
\end{minipage} & \begin{minipage}[b]{\linewidth}\raggedright
SCHEDULE
\end{minipage} \\
\midrule\noalign{}
\endhead
\bottomrule\noalign{}
\endlastfoot
\end{longtable}

\subsubsection{Description}\label{description-131}

The RUNSPEC activation keyword marks the start of the RUNSPEC section
that defines the key parameters for the simulator including the
dimensions of the model, phases present in the model (oil, gas and water
for example), number of tables for a given property and the maximum
number of rows for each table, the maximum number of groups, wells and
well completions, as well as various options to be invoked by OPM Flow.

Apart from COMMENTS entered by ``--'' in columns one and two, this
keyword should be the first keyword in the input deck.

There is no data required for this keyword and there is no keyword
terminating ``/''

\subsubsection{Example}\label{example-99}

\begin{verbatim}
-- ==============================================================================
--
-- RUNSPEC SECTION
--
-- ==============================================================================
RUNSPEC
\end{verbatim}

The above example marks the start of the RUNSPEC section in the OPM Flow
data input file.

\subsection{SAMG -- Activate Algebraic Multi-Grid Linear
Solver}\label{samg-activate-algebraic-multi-grid-linear-solver}

\begin{longtable}[]{@{}
  >{\raggedright\arraybackslash}p{(\linewidth - 14\tabcolsep) * \real{0.1250}}
  >{\raggedright\arraybackslash}p{(\linewidth - 14\tabcolsep) * \real{0.1250}}
  >{\raggedright\arraybackslash}p{(\linewidth - 14\tabcolsep) * \real{0.1250}}
  >{\raggedright\arraybackslash}p{(\linewidth - 14\tabcolsep) * \real{0.1250}}
  >{\raggedright\arraybackslash}p{(\linewidth - 14\tabcolsep) * \real{0.1250}}
  >{\raggedright\arraybackslash}p{(\linewidth - 14\tabcolsep) * \real{0.1250}}
  >{\raggedright\arraybackslash}p{(\linewidth - 14\tabcolsep) * \real{0.1250}}
  >{\raggedright\arraybackslash}p{(\linewidth - 14\tabcolsep) * \real{0.1250}}@{}}
\toprule\noalign{}
\begin{minipage}[b]{\linewidth}\raggedright
RUNSPEC
\end{minipage} & \begin{minipage}[b]{\linewidth}\raggedright
GRID
\end{minipage} & \begin{minipage}[b]{\linewidth}\raggedright
EDIT
\end{minipage} & \begin{minipage}[b]{\linewidth}\raggedright
PROPS
\end{minipage} & \begin{minipage}[b]{\linewidth}\raggedright
REGIONS
\end{minipage} & \begin{minipage}[b]{\linewidth}\raggedright
SOLUTION
\end{minipage} & \begin{minipage}[b]{\linewidth}\raggedright
SUMMARY
\end{minipage} & \begin{minipage}[b]{\linewidth}\raggedright
SCHEDULE
\end{minipage} \\
\midrule\noalign{}
\endhead
\bottomrule\noalign{}
\endlastfoot
\end{longtable}

\subsubsection{Description}\label{description-132}

This keyword actives the algebraic multi-grid linear solve; note this
solver is not available to the general public in the commercial
simulator.

This keyword is not supported by OPM Flow but would change the results
if supported so the simulation will be stopped.

See section 2.2 Running OPM Flow 2023-04 From The Command Line on how to
invoke various numerical schemes via the OPM Flow command line
interface.

\subsection{SATOPTS -- Activate Relative Permeability Assignment
Options}\label{satopts-activate-relative-permeability-assignment-options}

\begin{longtable}[]{@{}
  >{\raggedright\arraybackslash}p{(\linewidth - 14\tabcolsep) * \real{0.1250}}
  >{\raggedright\arraybackslash}p{(\linewidth - 14\tabcolsep) * \real{0.1250}}
  >{\raggedright\arraybackslash}p{(\linewidth - 14\tabcolsep) * \real{0.1250}}
  >{\raggedright\arraybackslash}p{(\linewidth - 14\tabcolsep) * \real{0.1250}}
  >{\raggedright\arraybackslash}p{(\linewidth - 14\tabcolsep) * \real{0.1250}}
  >{\raggedright\arraybackslash}p{(\linewidth - 14\tabcolsep) * \real{0.1250}}
  >{\raggedright\arraybackslash}p{(\linewidth - 14\tabcolsep) * \real{0.1250}}
  >{\raggedright\arraybackslash}p{(\linewidth - 14\tabcolsep) * \real{0.1250}}@{}}
\toprule\noalign{}
\begin{minipage}[b]{\linewidth}\raggedright
RUNSPEC
\end{minipage} & \begin{minipage}[b]{\linewidth}\raggedright
GRID
\end{minipage} & \begin{minipage}[b]{\linewidth}\raggedright
EDIT
\end{minipage} & \begin{minipage}[b]{\linewidth}\raggedright
PROPS
\end{minipage} & \begin{minipage}[b]{\linewidth}\raggedright
REGIONS
\end{minipage} & \begin{minipage}[b]{\linewidth}\raggedright
SOLUTION
\end{minipage} & \begin{minipage}[b]{\linewidth}\raggedright
SUMMARY
\end{minipage} & \begin{minipage}[b]{\linewidth}\raggedright
SCHEDULE
\end{minipage} \\
\midrule\noalign{}
\endhead
\bottomrule\noalign{}
\endlastfoot
\end{longtable}

\subsubsection{Description}\label{description-133}

The SATOPTS keyword activates OPM Flow's relative permeability
assignment options. The relative permeability functions are defined
using the either the:

\begin{itemize}
\tightlist
\item
  SWOF, SGOF, SLGOF series of saturation functions, or the
\item
  SWFN, SGFN, SGWFN, SOF2, SOF3, SOF32D series of functions.
\end{itemize}

The allocation of the relative permeability tables to the grid cells is
dependent on the options selected on this keyword (SATOPTS).

If the DIRECT option has been activated and the IRREVERS option has not
been invoked on the SATOPTS keyword, then different relative
permeability functions are used for the x, y, and z directions. Here the
KRNUMX, KRNUMY and KRNUMZ keywords are used for Cartesian grids to
allocate the relative permeability tables to the cells. For Radial grids
the KRNUMR, KRNUMT and KRNUMZ keywords should be used. This results in
the same relative permeability curves being used in both the xi to xi --
1 and the xi to xi + 1 flow directions. Similarly for the y direction
the same curves are used for the yi to yi -- 1 and the yi to yi + 1 flow
directions. And again for the z direction, the same relative
permeability function is used for flow in the zi to zi -- 1 and the zi
to zi + 1 flow directions.

If the DIRECT option has been activated and the IRREVERS option has also
been invoked on the SATOPTS keyword, then KRNUMX, KRNUMY and KRNUMZ
keywords are used for Cartesian grids to allocate the relative
permeability tables in the xi to xi + 1, yi to yi + 1, zi to zi + 1,
flow directions, respectively. For Radial grids the KRNUMR, KRNUMT and
KRNUMZ keywords should be used. For flow in the xi to xi -- 1 flow
directions, etc., the KRNUMX-, KRNUMY- and KRNUMZ- keywords are used for
Cartesian grids and the KRNUMR-, KRNUMT- and KRNUMZ- are used for radial
grids.

The HYSTER option activates relative permeability hysteresis for the
non-wetting phases (liquid and vapour) such that their values depend on
whether the non-wetting phase saturation is increasing or decreasing.
For this option the user specifies two sets of saturation functions, one
for a drainage process (decreasing wetting phase saturation) and one for
an imbibition process (increasing wetting phase saturation). The SATNUM
and IMBNUM keywords should then be used to identify which saturation
table applies to drainage and imbibition respectively. For a process
starting at the maximum wetting phase saturation, on the drainage curve,
with the wetting phase saturation decreasing, the drainage curve is
followed. Similarly for a process starting at the minimum wetting phase
saturation with the wetting phase saturation increasing, the imbibition
curve is followed. If the drainage or imbibition process is reversed at
some point, then the curve does not necessarily run back over its
previous values. In OPM Flow the Carlson {[}Carlson, F. M. (1981) SPE
10157, presented at the 56th Annual SPE Fall Meeting, San Antonio,
1981{]} model may be used to describe relative permeability hysteresis
and the Killough {[}Killough, J. E. ``Reservoir Simulation with
History-dependent Saturation Functions,'' paper SPE 5106, Society of
Petroleum Engineers Journal (1976) 16, No.~1, 37-48.{]} model may be
used for relative permeability or capillary pressure. See also the
EHYSTR keyword used to define the hysteresis model that is applied.

If the DIRECT option has been activated and the IRREVERS option has not
been invoked on the SATOPTS keyword, then the same set of keywords as
for the DIRECT only option are used to assign the drainage relative
permeability curves, that is: KRNUMX, KRNUMY and KRNUMZ, plus the
IMBNUMX, IMBNUMY, and IMBNUMZ, keywords for the imbibition curves. If
the DIRECT option has been activated and the IRREVERS option has also
been invoked on the SATOPTS keyword, then the same set are keywords as
for the DIRECT and IRREVERS option are used to assign the drainage
relative permeability curves, that is: KRNUMX, KRNUMX-, etc., plus the
IMBNUMX, IMBNUMY, IMBNUMZ, IMBNUMX-, IMBNUMY-, IMBNUMZ- keywords for the
imbibition curves. See Table 5.40 for the various relative permeability
table allocation keywords required for different combinations of the
DIRECT, IRREVERS and HYSTER options.

The keyword should be followed by one or more of the following keyword
options.

\begin{longtable}[]{@{}
  >{\raggedright\arraybackslash}p{(\linewidth - 6\tabcolsep) * \real{0.2500}}
  >{\raggedright\arraybackslash}p{(\linewidth - 6\tabcolsep) * \real{0.2500}}
  >{\raggedright\arraybackslash}p{(\linewidth - 6\tabcolsep) * \real{0.2500}}
  >{\raggedright\arraybackslash}p{(\linewidth - 6\tabcolsep) * \real{0.2500}}@{}}
\toprule\noalign{}
\begin{minipage}[b]{\linewidth}\raggedright
No.
\end{minipage} & \begin{minipage}[b]{\linewidth}\raggedright
Name
\end{minipage} & \begin{minipage}[b]{\linewidth}\raggedright
Description
\end{minipage} & \begin{minipage}[b]{\linewidth}\raggedright
Default
\end{minipage} \\
\midrule\noalign{}
\endhead
\bottomrule\noalign{}
\endlastfoot
1 & DIRECT & A character string that activates the directional relative
permeability assignment option. If the DIRECT option is specified then
directional relative permeability assignment is activated and different
relative permeability functions are assigned to the x, y and z
directions. In this case the KRNUMX, KRNUMY and KRNUMZ keywords are used
for Cartesian grids to allocate the relative permeability tables. For
Radial grids the KRNUMR, KRNUMT and KRNUMZ keywords should be used. &
None \\
2 & IRREVERS & A character string that activates the irreversible
directional relative permeability assignment option. If the IRREVERS
option is specified then the relative permeability assignment is set to
irreversible and results in different sets of relative permeability
tables being applied for flow in the xi to xi + 1 direction and the xi
to the xi -- 1 direction, for all directions (x, y, z). In this case the
KRNUMX, KRNUMY and KRNUMZ keywords are used for Cartesian grids to
allocate the relative permeability tables in the xi to xi + 1 flow
directions etc. For Radial grids the KRNUMR, KRNUMT and KRNUMZ keywords
should be used. For flow in the xi to xi -- 1 flow directions, etc., the
KRNUMX-, KRNUMY- and KRNUMZ- keywords are used for Cartesian grids and
the KRNUMR-, KRNUMT- and KRNUMZ- are used for radial grids. Only the
default option is supported by OPM Flow, that is irreversible
permeability tables are not supported. & None \\
3 & HYSTER & A character string that activates the hysteresis option. If
the HYSTER and DIRECT options have been activated and the IRREVERS has
not been invoked on the SATOPTS keyword, then different relative
permeability functions are used for the x, y, and z directions and for
the drainage and imbibition processes. Here the drainage relative
permeability curves are allocated via the KRNUMX, KRNUMY and KRNUMZ
keywords for Cartesian grids and the KRNUMR, KRNUMT and KRNUMZ keywords
for radial grids. The imbibition relative permeability curves are
allocated via the IMBNUMX, IMBNUMY and IMBNUMZ keywords for Cartesian
grids and the IMBNUMR, IMBNUMT and IMBNUMZ keywords for radial grids. If
the HYSTER, DIRECT and IRREVERS options have been activated, then
different relative permeability functions are used for the x, y, and z
directions, each flow direction and for the drainage and imbibition
processes. In addition to the aforementioned relative permeability
allocation keywords for the xi to xi + 1 flow direction etc., the xi to
xi -- 1 flow directions keywords, KRNUMX-, KRNUMY- and KRNUMZ- are used
for Cartesian grids and the KRNUMR-, KRNUMT- and KRNUMZ- are used for
radial grids. The imbibition relative permeability curves are allocated
via the IMBNUMX-, IMBNUMY- and IMBNUMZ- keywords for Cartesian grids and
the IMBNUMR, IMBNUMT and IMBNUMZ- keywords for radial grids. & None \\
4 & SURFTENS & A character string that activates the capillary pressure
surface tension pressure dependency option. Only the default option is
supported by OPM Flow, that is capillary pressure surface tension
pressure dependency option is not supported. & None \\
Notes: & & & \\
\end{longtable}

\emph{Table 5.39: SATOPTS Keyword Description}

For clarity the following table outlines the keywords that should be
used in allocating the relative permeability tables for the various
SATOPTS options.

\begin{longtable}[]{@{}
  >{\raggedright\arraybackslash}p{(\linewidth - 8\tabcolsep) * \real{0.2000}}
  >{\raggedright\arraybackslash}p{(\linewidth - 8\tabcolsep) * \real{0.2000}}
  >{\raggedright\arraybackslash}p{(\linewidth - 8\tabcolsep) * \real{0.2000}}
  >{\raggedright\arraybackslash}p{(\linewidth - 8\tabcolsep) * \real{0.2000}}
  >{\raggedright\arraybackslash}p{(\linewidth - 8\tabcolsep) * \real{0.2000}}@{}}
\toprule\noalign{}
\begin{minipage}[b]{\linewidth}\raggedright
Option
\end{minipage} & \begin{minipage}[b]{\linewidth}\raggedright
Cartesian
\end{minipage} & \begin{minipage}[b]{\linewidth}\raggedright
Radial
\end{minipage} & \begin{minipage}[b]{\linewidth}\raggedright
\end{minipage} & \begin{minipage}[b]{\linewidth}\raggedright
\end{minipage} \\
\midrule\noalign{}
\endhead
\bottomrule\noalign{}
\endlastfoot
DIRECT Flow in all directions & KRNUMX KRNUMY KRNUMZ & KRNUMR KRNUMT
KRNUMZ & & \\
DIRECT and IRREVERS Flow in the i to i +1 directions. Flow in the i to i
-1 directions. & KRNUMX, KRNUMY KRNUMZ KRNUMX- KRNUMY- KRNUMZ- & KRNUMR
KRNUMT KRNUMZ KRNUMR- KRNUMT- KRNUMZ- & & \\
DIRECT and HYSTER Flow in all directions. & Drainage KRNUMX KRNUMY
KRNUMZ & Imbibition IMBNUMX IMBNUMY IMBNUMZ & Drainage KRNUMR KRNUMT
KRNUMZ & Imbibition IMBNUMR IMBNUMT IMBNUMZ \\
DIRECT, IRREVERS and HYSTER Flow in the i to i +1 directions. Flow in
the i to i -1 directions. & Drainage KRNUMX KRNUMY KRNUMZ KRNUMX-,
KRNUMY- KRNUMZ- & Imbibition IMBNUMX IMBNUMY IMBNUMZ IMBNUMX- IMBNUMY-
IMBNUMZ- & Drainage KRNUMR KRNUMT KRNUMZ KRNUMR- KRNUMT- KRNUMZ- &
Imbibition IMBNUMR IMBNUMT IMBNUMZ IMBNUMR- IMBNUMT- IMBNUMZ- \\
Notes: & & & & \\
\end{longtable}

\emph{Table 5.40: SATOPTS Relative Permeability Function Allocation
Keywords}

\subsubsection{Examples}\label{examples-5}

The first example actives the directional relative permeability
assignment option only and hence the following keywords are used to
allocate the relative permeability arrays for Cartesian grids: KRNUMX,
KRNUMY, and KRNUMZ.

\begin{verbatim}
--
--       ACTIVATE RELATIVE PERMEABILITY ASSIGNMENT HYSTERESIS OPTIONS
--       DIRECTTIONAL(DIRECT) IRREVERSIBLE(IRREVERS) HYSTERESIS(HYSTER)
SATOPTS
         'DIRECT'                                                              /
\end{verbatim}

The next example actives the directional and irreversible relative
permeability assignment options, and hence the following keywords are
used to allocate the relative permeability arrays for Cartesian grids:
KRNUMX, KRNUMY, KRNUMZ, KRNUMX-, KRNUMY- and KRNUMZ-.

\begin{verbatim}
--
--       ACTIVATE RELATIVE PERMEABILITY ASSIGNMENT HYSTERESIS OPTIONS
--       DIRECTTIONAL(DIRECT) IRREVERSIBLE(IRREVERS) HYSTERESIS(HYSTER)
SATOPTS
         'DIRECT'   'IRREVERS'                                                 /
\end{verbatim}

Finally, the last example invokes the directional, irreversible and
hysteresis assignment options.

\begin{verbatim}
--
--       ACTIVATE RELATIVE PERMEABILITY ASSIGNMENT HYSTERESIS OPTIONS
--       DIRECTTIONAL(DIRECT) IRREVERSIBLE(IRREVERS) HYSTERESIS(HYSTER)
SATOPTS
         'DIRECT'   'IRREVERS'   'HYSTER'                                      /
\end{verbatim}

In this case the drainage relative permeability curves are allocated by
the KRNUMX, KRNUMY, KRNUMZ, KRNUMX-, KRNUMY-, KRNUMZ- keywords, and the
imbibition relative permeability curves are allocated by the IMBNUMX,
IMBNUMY, IMBNUMZ, IMBNUMX-, IMBNUMY-, IMBNUMZ- keywords.

\begin{tcolorbox}[enhanced jigsaw, title=\textcolor{quarto-callout-note-color}{\faInfo}\hspace{0.5em}{Note}, coltitle=black, breakable, toprule=.15mm, colbacktitle=quarto-callout-note-color!10!white, leftrule=.75mm, opacitybacktitle=0.6, colback=white, bottomtitle=1mm, toptitle=1mm, titlerule=0mm, left=2mm, colframe=quarto-callout-note-color-frame, arc=.35mm, rightrule=.15mm, bottomrule=.15mm, opacityback=0]

This keyword activates how relative permeability curves are assigned in
the model. The ENDSCALE keyword allows the end-point scaling also to
vary with direction, flow direction and hysteresis process, resulting in
a great deal of flexibility. Whether or not all these features should be
used though is another question.

\end{tcolorbox}

\subsection{SAVE -- Activate Output of a SAVE File for Fast
Restarts}\label{save-activate-output-of-a-save-file-for-fast-restarts}

\begin{longtable}[]{@{}
  >{\raggedright\arraybackslash}p{(\linewidth - 14\tabcolsep) * \real{0.1250}}
  >{\raggedright\arraybackslash}p{(\linewidth - 14\tabcolsep) * \real{0.1250}}
  >{\raggedright\arraybackslash}p{(\linewidth - 14\tabcolsep) * \real{0.1250}}
  >{\raggedright\arraybackslash}p{(\linewidth - 14\tabcolsep) * \real{0.1250}}
  >{\raggedright\arraybackslash}p{(\linewidth - 14\tabcolsep) * \real{0.1250}}
  >{\raggedright\arraybackslash}p{(\linewidth - 14\tabcolsep) * \real{0.1250}}
  >{\raggedright\arraybackslash}p{(\linewidth - 14\tabcolsep) * \real{0.1250}}
  >{\raggedright\arraybackslash}p{(\linewidth - 14\tabcolsep) * \real{0.1250}}@{}}
\toprule\noalign{}
\begin{minipage}[b]{\linewidth}\raggedright
RUNSPEC
\end{minipage} & \begin{minipage}[b]{\linewidth}\raggedright
GRID
\end{minipage} & \begin{minipage}[b]{\linewidth}\raggedright
EDIT
\end{minipage} & \begin{minipage}[b]{\linewidth}\raggedright
PROPS
\end{minipage} & \begin{minipage}[b]{\linewidth}\raggedright
REGIONS
\end{minipage} & \begin{minipage}[b]{\linewidth}\raggedright
SOLUTION
\end{minipage} & \begin{minipage}[b]{\linewidth}\raggedright
SUMMARY
\end{minipage} & \begin{minipage}[b]{\linewidth}\raggedright
SCHEDULE
\end{minipage} \\
\midrule\noalign{}
\endhead
\bottomrule\noalign{}
\endlastfoot
\end{longtable}

\subsubsection{Description}\label{description-134}

This keyword activates output of a SAVE file for fast restarts which are
read in by the LOAD keyword in the RUNSPEC section. No data is required
for this keyword.

OPM Flow does not support the SAVE file format for fast restarts like
the commercial simulator, but instead writes a standard RESTART file at
the requested time step in the SCHEDULE section, which can then be used
to restart the simulation at a given point in time via the RESTART
keyword in the RUNSPEC section.

\subsubsection{Example}\label{example-100}

The first example requests that a SAVE file be written out in the
RUNSPEC section; however, OPM Flow will not write a RESTART record if
the SAVE keyword is encountered in the RUNSPEC section.

\begin{verbatim}
-- ==============================================================================
--
-- RUNSPEC SECTION
--
-- ==============================================================================
RUNSPEC
--
--       WRITE OUT SAVE FILE FOR FAST RESTARTS
--
SAVE
\end{verbatim}

The second example shows how the keyword is used in the SCHEDULE
section.

\begin{verbatim}
-- ==============================================================================
--
-- SCHEDULE SECTION
--
-- ==============================================================================
SCHEDULE

-- ------------------------------------------------------------------------------
-- SCHEDULE SECTION - 2020-01-01
-- ------------------------------------------------------------------------------
RPTSCHED
          'WELLS=2'    'WELSPECS'    'CPU=2'     FIP=2'                         /

DATES
          1  APR   2021  /
/
--
--       WRITE OUT SAVE FILE FOR FAST RESTARTS
--
SAVE
\end{verbatim}

\begin{verbatim}
DATES
          1  JLY   2021  /
          1  OCT   2021  /
/
\end{verbatim}

Here OPM Flow will write out a RESTART file instead of a SAVE file at
April 1st, 2021.

\subsection{SCDPDIMS -- Define Scale Deposition and Damage Table
Dimensions}\label{scdpdims-define-scale-deposition-and-damage-table-dimensions}

\begin{longtable}[]{@{}
  >{\raggedright\arraybackslash}p{(\linewidth - 14\tabcolsep) * \real{0.1250}}
  >{\raggedright\arraybackslash}p{(\linewidth - 14\tabcolsep) * \real{0.1250}}
  >{\raggedright\arraybackslash}p{(\linewidth - 14\tabcolsep) * \real{0.1250}}
  >{\raggedright\arraybackslash}p{(\linewidth - 14\tabcolsep) * \real{0.1250}}
  >{\raggedright\arraybackslash}p{(\linewidth - 14\tabcolsep) * \real{0.1250}}
  >{\raggedright\arraybackslash}p{(\linewidth - 14\tabcolsep) * \real{0.1250}}
  >{\raggedright\arraybackslash}p{(\linewidth - 14\tabcolsep) * \real{0.1250}}
  >{\raggedright\arraybackslash}p{(\linewidth - 14\tabcolsep) * \real{0.1250}}@{}}
\toprule\noalign{}
\begin{minipage}[b]{\linewidth}\raggedright
RUNSPEC
\end{minipage} & \begin{minipage}[b]{\linewidth}\raggedright
GRID
\end{minipage} & \begin{minipage}[b]{\linewidth}\raggedright
EDIT
\end{minipage} & \begin{minipage}[b]{\linewidth}\raggedright
PROPS
\end{minipage} & \begin{minipage}[b]{\linewidth}\raggedright
REGIONS
\end{minipage} & \begin{minipage}[b]{\linewidth}\raggedright
SOLUTION
\end{minipage} & \begin{minipage}[b]{\linewidth}\raggedright
SUMMARY
\end{minipage} & \begin{minipage}[b]{\linewidth}\raggedright
SCHEDULE
\end{minipage} \\
\midrule\noalign{}
\endhead
\bottomrule\noalign{}
\endlastfoot
\end{longtable}

\subsubsection{Description}\label{description-135}

The SCDPDIMS keyword defines the number of tables used in the Scale
Deposition option and the maximum number of entries for the various
tables.

This keyword is not supported by OPM Flow but would change the results
if supported so the simulation will be stopped.

\begin{longtable}[]{@{}
  >{\raggedright\arraybackslash}p{(\linewidth - 6\tabcolsep) * \real{0.2500}}
  >{\raggedright\arraybackslash}p{(\linewidth - 6\tabcolsep) * \real{0.2500}}
  >{\raggedright\arraybackslash}p{(\linewidth - 6\tabcolsep) * \real{0.2500}}
  >{\raggedright\arraybackslash}p{(\linewidth - 6\tabcolsep) * \real{0.2500}}@{}}
\toprule\noalign{}
\begin{minipage}[b]{\linewidth}\raggedright
No.
\end{minipage} & \begin{minipage}[b]{\linewidth}\raggedright
Name
\end{minipage} & \begin{minipage}[b]{\linewidth}\raggedright
Description
\end{minipage} & \begin{minipage}[b]{\linewidth}\raggedright
Default
\end{minipage} \\
\midrule\noalign{}
\endhead
\bottomrule\noalign{}
\endlastfoot
1 & NTSCDP & NTSCDP is a positive integer that defines the number of
SCDPTAB scale deposition tables used in the Scale Deposition option. &
0 \\
2 & NPSCDP & NPSCDP is a positive integer that defines the maximum
number of entries (or rows) in any one SCDPTAB scale deposition table
defined in the input deck. & 0 \\
3 & NTSCDA & NTSCDA is a positive integer that defines the number of
SCDATAB scale damage tables used in the Scale Deposition option. & 0 \\
4 & NPSCDA & NPSCDA is a positive integer that defines the maximum
number of entries (or rows) in any one SCDATAB scale damage table
defined in the input deck. & 0 \\
5 & Not Used & & 1* \\
6 & Not Used & & 1* \\
7 & NTSCDE & NTSCDE is a positive integer that defines the number of
SCDETAB karst aquifer dissolution tables used in the Scale Deposition
option. & 0 \\
Notes: & & & \\
\end{longtable}

\emph{Table 5.41: SCDPDIMS Keyword Description}

\subsubsection{Example}\label{example-101}

\begin{verbatim}
--
--       NO.     MAX     NO.     MAX     NOT   NOT   NO.
--       NTSCDP  NPSCDP  NTSCDA  NPSCDA  USED  USED  NTSCDE
SCDPDIMS
         5       10      4       10      1*    1*    3                         /
\end{verbatim}

The above example defines the number of SCDPTAB scale deposition tables
to be five with a maximum number of rows for each table set to 10, the
maximum number of SCDATAB scale damage tables to be four with a maximum
number of 10 rows per table, and the maximum number of SCDETAB karst
aquifer dissolution tables to be three.

\subsection{SKIP -- Activate Skipping of All Keywords and Input
Data}\label{skip-activate-skipping-of-all-keywords-and-input-data-1}

The SKIP keyword activates skipping of all keywords and input data until
the ENDSKIP keyword is encountered. All keywords between the SKIP and
ENDSKIP keywords are ignored.

See SKIP -- Activate Skipping of All Keywords and Input Data in the
GLOBAL section for a full description.

\subsection{SKIP100 -- Activate Skipping of Black-Oil Keywords and Input
Data}\label{skip100-activate-skipping-of-black-oil-keywords-and-input-data}

The SKIP100 keyword activates skipping of all keywords and input data by
the commercial black-oil simulator until the ENDSKIP keyword is
encountered. All keywords between the SKIP100 and ENDSKIP keywords are
ignored by the commercial black-oil simulator. The SKIP100 keyword is
ignored by the commercial compositional simulator.

See SKIP100 -- Activate Skipping of Keywords by Black-Oil Simulator in
the GLOBAL section for a full description.

\subsection{SKIP300 -- Activate Skipping of Keywords by Compositional
Simulator}\label{skip300-activate-skipping-of-keywords-by-compositional-simulator-1}

The SKIP300 keyword activates skipping of all keywords and input data by
the commercial compositional simulator until the ENDSKIP keyword is
encountered. All keywords between the SKIP300 and ENDSKIP keywords are
ignored by the commercial compositional simulator. The SKIP300 keyword
is ignored by the commercial black-oil simulator.

See Error: Reference source not found in the GLOBAL section for a full
description.

\subsection{SMRYDIMS -- Define Maximum Number of Summary Vectors to be
Written}\label{smrydims-define-maximum-number-of-summary-vectors-to-be-written}

\begin{longtable}[]{@{}
  >{\raggedright\arraybackslash}p{(\linewidth - 14\tabcolsep) * \real{0.1250}}
  >{\raggedright\arraybackslash}p{(\linewidth - 14\tabcolsep) * \real{0.1250}}
  >{\raggedright\arraybackslash}p{(\linewidth - 14\tabcolsep) * \real{0.1250}}
  >{\raggedright\arraybackslash}p{(\linewidth - 14\tabcolsep) * \real{0.1250}}
  >{\raggedright\arraybackslash}p{(\linewidth - 14\tabcolsep) * \real{0.1250}}
  >{\raggedright\arraybackslash}p{(\linewidth - 14\tabcolsep) * \real{0.1250}}
  >{\raggedright\arraybackslash}p{(\linewidth - 14\tabcolsep) * \real{0.1250}}
  >{\raggedright\arraybackslash}p{(\linewidth - 14\tabcolsep) * \real{0.1250}}@{}}
\toprule\noalign{}
\begin{minipage}[b]{\linewidth}\raggedright
RUNSPEC
\end{minipage} & \begin{minipage}[b]{\linewidth}\raggedright
GRID
\end{minipage} & \begin{minipage}[b]{\linewidth}\raggedright
EDIT
\end{minipage} & \begin{minipage}[b]{\linewidth}\raggedright
PROPS
\end{minipage} & \begin{minipage}[b]{\linewidth}\raggedright
REGIONS
\end{minipage} & \begin{minipage}[b]{\linewidth}\raggedright
SOLUTION
\end{minipage} & \begin{minipage}[b]{\linewidth}\raggedright
SUMMARY
\end{minipage} & \begin{minipage}[b]{\linewidth}\raggedright
SCHEDULE
\end{minipage} \\
\midrule\noalign{}
\endhead
\bottomrule\noalign{}
\endlastfoot
\end{longtable}

\subsubsection{Description}\label{description-136}

The SMRYDIMS keyword defines the maximum number of summary vectors to be
written out to the SUMMARY file (*.SUMMARY).

OPM Flow uses dynamic memory allocation and therefore the keyword has no
effect and is ignored by the simulator, but is documented here for
completeness.

\begin{longtable}[]{@{}
  >{\raggedright\arraybackslash}p{(\linewidth - 6\tabcolsep) * \real{0.2500}}
  >{\raggedright\arraybackslash}p{(\linewidth - 6\tabcolsep) * \real{0.2500}}
  >{\raggedright\arraybackslash}p{(\linewidth - 6\tabcolsep) * \real{0.2500}}
  >{\raggedright\arraybackslash}p{(\linewidth - 6\tabcolsep) * \real{0.2500}}@{}}
\toprule\noalign{}
\begin{minipage}[b]{\linewidth}\raggedright
No.
\end{minipage} & \begin{minipage}[b]{\linewidth}\raggedright
Name
\end{minipage} & \begin{minipage}[b]{\linewidth}\raggedright
Description
\end{minipage} & \begin{minipage}[b]{\linewidth}\raggedright
Default
\end{minipage} \\
\midrule\noalign{}
\endhead
\bottomrule\noalign{}
\endlastfoot
1 & NSUMMX & A positive integer that defines the maximum number of
summary vectors to be written out to the SUMMARY file (*.SUMMARY). &
10000 \\
Notes: & & & \\
\end{longtable}

\emph{Table 5.42: SMRYDIMS Keyword Description}

\subsubsection{Example}\label{example-102}

\begin{verbatim}
--
--       SET THE MAXIMUM NUMBER OF SUMMARY VECTORS THAT CAN BE WRITTEN OUT
--
SMRYDIMS
         10000                                                                 /
\end{verbatim}

The above example sets maximum number of summary vectors that can be
written out to the SUMMARY file to the default value of 10,000; however,
this has no effect in OPM Flow input decks.

\subsection{SOLVDIMS -- Define PEBI Grid Nested Factorization Solver
Dimensions}\label{solvdims-define-pebi-grid-nested-factorization-solver-dimensions}

\begin{longtable}[]{@{}
  >{\raggedright\arraybackslash}p{(\linewidth - 14\tabcolsep) * \real{0.1250}}
  >{\raggedright\arraybackslash}p{(\linewidth - 14\tabcolsep) * \real{0.1250}}
  >{\raggedright\arraybackslash}p{(\linewidth - 14\tabcolsep) * \real{0.1250}}
  >{\raggedright\arraybackslash}p{(\linewidth - 14\tabcolsep) * \real{0.1250}}
  >{\raggedright\arraybackslash}p{(\linewidth - 14\tabcolsep) * \real{0.1250}}
  >{\raggedright\arraybackslash}p{(\linewidth - 14\tabcolsep) * \real{0.1250}}
  >{\raggedright\arraybackslash}p{(\linewidth - 14\tabcolsep) * \real{0.1250}}
  >{\raggedright\arraybackslash}p{(\linewidth - 14\tabcolsep) * \real{0.1250}}@{}}
\toprule\noalign{}
\begin{minipage}[b]{\linewidth}\raggedright
RUNSPEC
\end{minipage} & \begin{minipage}[b]{\linewidth}\raggedright
GRID
\end{minipage} & \begin{minipage}[b]{\linewidth}\raggedright
EDIT
\end{minipage} & \begin{minipage}[b]{\linewidth}\raggedright
PROPS
\end{minipage} & \begin{minipage}[b]{\linewidth}\raggedright
REGIONS
\end{minipage} & \begin{minipage}[b]{\linewidth}\raggedright
SOLUTION
\end{minipage} & \begin{minipage}[b]{\linewidth}\raggedright
SUMMARY
\end{minipage} & \begin{minipage}[b]{\linewidth}\raggedright
SCHEDULE
\end{minipage} \\
\midrule\noalign{}
\endhead
\bottomrule\noalign{}
\endlastfoot
\end{longtable}

\subsubsection{Description}\label{description-137}

The SOLVDIMS defines the unstructured Perpendicular Bisector (``PEBI'')
{[}Heinemann, Z.E. and Brand, C.W. 1988. Gridding Techniques in
Reservoir Simulation. Proc., First Intl. Forum on Reservoir Simulation,
Alpbach, Austria, 339.{]} and {[}Heinemann, Z.E., Brand, C.W., Munka, M.
et al.~1991. Modeling Reservoir Geometry With Irregular Grids. SPE Res
Eng 6 (2): 225--232. SPE-18412-PA. http://dx.doi.org/10.2118/18412-PA{]}
grid nested factorization solver dimensions. This keyword is generated
by an external pre-processing program for generating simulation grids.

This keyword is not supported by OPM Flow but would change the results
if supported so the simulation will be stopped.

\subsection{SOLVENT -- Activate the SOLVENT Phase in the
Model}\label{solvent-activate-the-solvent-phase-in-the-model}

\begin{longtable}[]{@{}
  >{\raggedright\arraybackslash}p{(\linewidth - 14\tabcolsep) * \real{0.1250}}
  >{\raggedright\arraybackslash}p{(\linewidth - 14\tabcolsep) * \real{0.1250}}
  >{\raggedright\arraybackslash}p{(\linewidth - 14\tabcolsep) * \real{0.1250}}
  >{\raggedright\arraybackslash}p{(\linewidth - 14\tabcolsep) * \real{0.1250}}
  >{\raggedright\arraybackslash}p{(\linewidth - 14\tabcolsep) * \real{0.1250}}
  >{\raggedright\arraybackslash}p{(\linewidth - 14\tabcolsep) * \real{0.1250}}
  >{\raggedright\arraybackslash}p{(\linewidth - 14\tabcolsep) * \real{0.1250}}
  >{\raggedright\arraybackslash}p{(\linewidth - 14\tabcolsep) * \real{0.1250}}@{}}
\toprule\noalign{}
\begin{minipage}[b]{\linewidth}\raggedright
RUNSPEC
\end{minipage} & \begin{minipage}[b]{\linewidth}\raggedright
GRID
\end{minipage} & \begin{minipage}[b]{\linewidth}\raggedright
EDIT
\end{minipage} & \begin{minipage}[b]{\linewidth}\raggedright
PROPS
\end{minipage} & \begin{minipage}[b]{\linewidth}\raggedright
REGIONS
\end{minipage} & \begin{minipage}[b]{\linewidth}\raggedright
SOLUTION
\end{minipage} & \begin{minipage}[b]{\linewidth}\raggedright
SUMMARY
\end{minipage} & \begin{minipage}[b]{\linewidth}\raggedright
SCHEDULE
\end{minipage} \\
\midrule\noalign{}
\endhead
\bottomrule\noalign{}
\endlastfoot
\end{longtable}

\subsubsection{Description}\label{description-138}

This keyword indicates that the solvent phase is present in the model
and to activate the four component solvent model for this run. In
addition to this keyword, the oil, water and gases phases should also be
declared for the run using the OIL, WATER and GAS keywords. The keyword
will also invoke data input file checking to ensure that all the
required Solvent phase input parameters are defined in the input deck.

There is no data required for this keyword and there is no terminating
``/'' for this keyword.

\subsubsection{Example}\label{example-103}

\begin{verbatim}
--
--       SOLVENT PHASE IS PRESENT IN THE RUN
--
SOLVENT
\end{verbatim}

The above example declares that the solvent phase is active in the
model.

\subsection{SPIDER -- Spider Grid Activation
Option}\label{spider-spider-grid-activation-option}

\begin{longtable}[]{@{}
  >{\raggedright\arraybackslash}p{(\linewidth - 14\tabcolsep) * \real{0.1250}}
  >{\raggedright\arraybackslash}p{(\linewidth - 14\tabcolsep) * \real{0.1250}}
  >{\raggedright\arraybackslash}p{(\linewidth - 14\tabcolsep) * \real{0.1250}}
  >{\raggedright\arraybackslash}p{(\linewidth - 14\tabcolsep) * \real{0.1250}}
  >{\raggedright\arraybackslash}p{(\linewidth - 14\tabcolsep) * \real{0.1250}}
  >{\raggedright\arraybackslash}p{(\linewidth - 14\tabcolsep) * \real{0.1250}}
  >{\raggedright\arraybackslash}p{(\linewidth - 14\tabcolsep) * \real{0.1250}}
  >{\raggedright\arraybackslash}p{(\linewidth - 14\tabcolsep) * \real{0.1250}}@{}}
\toprule\noalign{}
\begin{minipage}[b]{\linewidth}\raggedright
RUNSPEC
\end{minipage} & \begin{minipage}[b]{\linewidth}\raggedright
GRID
\end{minipage} & \begin{minipage}[b]{\linewidth}\raggedright
EDIT
\end{minipage} & \begin{minipage}[b]{\linewidth}\raggedright
PROPS
\end{minipage} & \begin{minipage}[b]{\linewidth}\raggedright
REGIONS
\end{minipage} & \begin{minipage}[b]{\linewidth}\raggedright
SOLUTION
\end{minipage} & \begin{minipage}[b]{\linewidth}\raggedright
SUMMARY
\end{minipage} & \begin{minipage}[b]{\linewidth}\raggedright
SCHEDULE
\end{minipage} \\
\midrule\noalign{}
\endhead
\bottomrule\noalign{}
\endlastfoot
\end{longtable}

\subsubsection{Description}\label{description-139}

SPIDER activates the OPM Flow spider grid geometry option for the model,
if this keyword and the RADIAL keyword are omitted then Cartesian
geometry is assumed by OPM Flow. Note that is an OPM Flow specific
keyword and option.

Spider Grids basically take the standard radial geometry keywords in the
grid section and transform the grid specification to an Irregular
Corner-Point Grid which can then be viewed in OPM ResInsight in a more
intuitive form. The only difference is that in the RUNSPEC section the
SPIDER keyword is used instead of the RADIAL keyword.

See also the RADIAL keyword in the RUNSPEC section (RADIAL -- Radial
Grid Activation Option). Both Spider and Radial grids created by OPM
Flow will be displayed the same in OPM ResInsight, as they are both
output as Irregular Corner-Point Grid files. The difference is that the
radial model's pore volumes are adjusted to match the radial geometry
pore volumes, whereas, the Spider grid volumes are not adjusted. See the
examples in the section on Radial Grids (Radial Grids) and Spider Grids
(Spider Grids) for further information

\subsubsection{Example}\label{example-104}

\begin{verbatim}
--
--       DEFINE SPIDER GRID GEOMETRY (OPM FLOW RADIAL GRID KEYWORD)
--
SPIDER
\end{verbatim}

The above example actives OPM Flow's spider grid geometry option.

\subsection{START -- Simulation Start
Date}\label{start-simulation-start-date}

\begin{longtable}[]{@{}
  >{\raggedright\arraybackslash}p{(\linewidth - 14\tabcolsep) * \real{0.1250}}
  >{\raggedright\arraybackslash}p{(\linewidth - 14\tabcolsep) * \real{0.1250}}
  >{\raggedright\arraybackslash}p{(\linewidth - 14\tabcolsep) * \real{0.1250}}
  >{\raggedright\arraybackslash}p{(\linewidth - 14\tabcolsep) * \real{0.1250}}
  >{\raggedright\arraybackslash}p{(\linewidth - 14\tabcolsep) * \real{0.1250}}
  >{\raggedright\arraybackslash}p{(\linewidth - 14\tabcolsep) * \real{0.1250}}
  >{\raggedright\arraybackslash}p{(\linewidth - 14\tabcolsep) * \real{0.1250}}
  >{\raggedright\arraybackslash}p{(\linewidth - 14\tabcolsep) * \real{0.1250}}@{}}
\toprule\noalign{}
\begin{minipage}[b]{\linewidth}\raggedright
RUNSPEC
\end{minipage} & \begin{minipage}[b]{\linewidth}\raggedright
GRID
\end{minipage} & \begin{minipage}[b]{\linewidth}\raggedright
EDIT
\end{minipage} & \begin{minipage}[b]{\linewidth}\raggedright
PROPS
\end{minipage} & \begin{minipage}[b]{\linewidth}\raggedright
REGIONS
\end{minipage} & \begin{minipage}[b]{\linewidth}\raggedright
SOLUTION
\end{minipage} & \begin{minipage}[b]{\linewidth}\raggedright
SUMMARY
\end{minipage} & \begin{minipage}[b]{\linewidth}\raggedright
SCHEDULE
\end{minipage} \\
\midrule\noalign{}
\endhead
\bottomrule\noalign{}
\endlastfoot
\end{longtable}

\subsubsection{Description}\label{description-140}

This keyword sets the start date for the simulation switches. If the
DATES keyword is to be used during the simulation, then a start date
should be entered.

\begin{longtable}[]{@{}
  >{\raggedright\arraybackslash}p{(\linewidth - 6\tabcolsep) * \real{0.2500}}
  >{\raggedright\arraybackslash}p{(\linewidth - 6\tabcolsep) * \real{0.2500}}
  >{\raggedright\arraybackslash}p{(\linewidth - 6\tabcolsep) * \real{0.2500}}
  >{\raggedright\arraybackslash}p{(\linewidth - 6\tabcolsep) * \real{0.2500}}@{}}
\toprule\noalign{}
\begin{minipage}[b]{\linewidth}\raggedright
No.
\end{minipage} & \begin{minipage}[b]{\linewidth}\raggedright
Name
\end{minipage} & \begin{minipage}[b]{\linewidth}\raggedright
Description
\end{minipage} & \begin{minipage}[b]{\linewidth}\raggedright
Default
\end{minipage} \\
\midrule\noalign{}
\endhead
\bottomrule\noalign{}
\endlastfoot
1 & DAY & A positive integer that defines the day of the month, the
value should be greater than or equal to one and less than or equal to
31. & None \\
2 & MONTH & Character string for the month and should be one of the
following `JAN', `FEB', `MAR', `APR', `MAY', `JUN', `JUL' (or `JLY'),
`AUG', `SEP', `OCT', `NOV', or `DEC' & None \\
3 & YEAR & A positive four digit integer value of the start year, which
must be specified fully by four digits, that is 1986. & None \\
Notes: & & & \\
\end{longtable}

\emph{Table 5.43: START Keyword Description}

\subsubsection{Example}\label{example-105}

\begin{verbatim}
--
--       DEFINE THE START DATE FOR THE RUN
--
START
         01 'JAN' 2014                                                         /
\end{verbatim}

The above example sets the start date for the run to be January 1, 2014.

\begin{tcolorbox}[enhanced jigsaw, title=\textcolor{quarto-callout-note-color}{\faInfo}\hspace{0.5em}{Note}, coltitle=black, breakable, toprule=.15mm, colbacktitle=quarto-callout-note-color!10!white, leftrule=.75mm, opacitybacktitle=0.6, colback=white, bottomtitle=1mm, toptitle=1mm, titlerule=0mm, left=2mm, colframe=quarto-callout-note-color-frame, arc=.35mm, rightrule=.15mm, bottomrule=.15mm, opacityback=0]

Whenever possible it is a good idea to always set the start date to be
at the beginning of the year as per the example. As like most
simulators, OPM Flow reports are always stated at the number of days
from the start date (and sometimes at a given date). If the start date
is at the beginning of the year, then calculating the actual date is
relatively straight forward and simple.

\end{tcolorbox}

\subsection{SURFACT -- Activate the Surfactant Phase in the
Model}\label{surfact-activate-the-surfactant-phase-in-the-model}

\begin{longtable}[]{@{}
  >{\raggedright\arraybackslash}p{(\linewidth - 14\tabcolsep) * \real{0.1250}}
  >{\raggedright\arraybackslash}p{(\linewidth - 14\tabcolsep) * \real{0.1250}}
  >{\raggedright\arraybackslash}p{(\linewidth - 14\tabcolsep) * \real{0.1250}}
  >{\raggedright\arraybackslash}p{(\linewidth - 14\tabcolsep) * \real{0.1250}}
  >{\raggedright\arraybackslash}p{(\linewidth - 14\tabcolsep) * \real{0.1250}}
  >{\raggedright\arraybackslash}p{(\linewidth - 14\tabcolsep) * \real{0.1250}}
  >{\raggedright\arraybackslash}p{(\linewidth - 14\tabcolsep) * \real{0.1250}}
  >{\raggedright\arraybackslash}p{(\linewidth - 14\tabcolsep) * \real{0.1250}}@{}}
\toprule\noalign{}
\begin{minipage}[b]{\linewidth}\raggedright
RUNSPEC
\end{minipage} & \begin{minipage}[b]{\linewidth}\raggedright
GRID
\end{minipage} & \begin{minipage}[b]{\linewidth}\raggedright
EDIT
\end{minipage} & \begin{minipage}[b]{\linewidth}\raggedright
PROPS
\end{minipage} & \begin{minipage}[b]{\linewidth}\raggedright
REGIONS
\end{minipage} & \begin{minipage}[b]{\linewidth}\raggedright
SOLUTION
\end{minipage} & \begin{minipage}[b]{\linewidth}\raggedright
SUMMARY
\end{minipage} & \begin{minipage}[b]{\linewidth}\raggedright
SCHEDULE
\end{minipage} \\
\midrule\noalign{}
\endhead
\bottomrule\noalign{}
\endlastfoot
\end{longtable}

\subsubsection{Description}\label{description-141}

This keyword indicates that the surfactant phase is present in the model
and to activate the surfactant flooding model. The keyword will also
invoke data input file checking to ensure that all the required
surfactant phase input parameters are defined in the input deck. See
also the SURFACTW keyword in the RUNSPEC section that actives the
surfactant phase, but with the changes to the wettability option
activated as well.

There is no data required for this keyword and there is no terminating
``/'' for this keyword.

This keyword is not supported by OPM Flow but would change the results
if supported so the simulation will be stopped.

\subsubsection{Example}\label{example-106}

\begin{verbatim}
--
--       ACTIVATE THE SURFACTANT PHASE IN THE MODEL
--
SURFACT
\end{verbatim}

The above example declares that the surfactant phase is active in the
model.

\subsection{SURFACTW -- Activate the Surfactant Phase with Wettability
Changes in the
Model}\label{surfactw-activate-the-surfactant-phase-with-wettability-changes-in-the-model}

\begin{longtable}[]{@{}
  >{\raggedright\arraybackslash}p{(\linewidth - 14\tabcolsep) * \real{0.1250}}
  >{\raggedright\arraybackslash}p{(\linewidth - 14\tabcolsep) * \real{0.1250}}
  >{\raggedright\arraybackslash}p{(\linewidth - 14\tabcolsep) * \real{0.1250}}
  >{\raggedright\arraybackslash}p{(\linewidth - 14\tabcolsep) * \real{0.1250}}
  >{\raggedright\arraybackslash}p{(\linewidth - 14\tabcolsep) * \real{0.1250}}
  >{\raggedright\arraybackslash}p{(\linewidth - 14\tabcolsep) * \real{0.1250}}
  >{\raggedright\arraybackslash}p{(\linewidth - 14\tabcolsep) * \real{0.1250}}
  >{\raggedright\arraybackslash}p{(\linewidth - 14\tabcolsep) * \real{0.1250}}@{}}
\toprule\noalign{}
\begin{minipage}[b]{\linewidth}\raggedright
RUNSPEC
\end{minipage} & \begin{minipage}[b]{\linewidth}\raggedright
GRID
\end{minipage} & \begin{minipage}[b]{\linewidth}\raggedright
EDIT
\end{minipage} & \begin{minipage}[b]{\linewidth}\raggedright
PROPS
\end{minipage} & \begin{minipage}[b]{\linewidth}\raggedright
REGIONS
\end{minipage} & \begin{minipage}[b]{\linewidth}\raggedright
SOLUTION
\end{minipage} & \begin{minipage}[b]{\linewidth}\raggedright
SUMMARY
\end{minipage} & \begin{minipage}[b]{\linewidth}\raggedright
SCHEDULE
\end{minipage} \\
\midrule\noalign{}
\endhead
\bottomrule\noalign{}
\endlastfoot
\end{longtable}

\subsubsection{Description}\label{description-142}

This keyword indicates that the surfactant phase is present in the model
and to activate the surfactant flooding mode with Changes to Wettability
option activated as well. The keyword will also invoke data input file
checking to ensure that all the required surfactant phase input
parameters are defined in the input deck. See also the SURFACT keyword
in the RUNSPEC section that actives the surfactant phase only, that is
without the Changes to the Wettability option.

There is no data required for this keyword and there is no terminating
``/'' for this keyword.

This keyword is not supported by OPM Flow but would change the results
if supported so the simulation will be stopped.

\subsubsection{Example}\label{example-107}

\begin{verbatim}
--
--       ACTIVATE THE SURFACTANT PHASE WITH WETTABILITY CHANGES IN THE MODEL
--
SURFACTW
\end{verbatim}

The above example declares that the surfactant phase is active in the
model together with the wettability changes.

\subsection{TABDIMS -- Define the Number of Tables and the Table
Dimensions}\label{tabdims-define-the-number-of-tables-and-the-table-dimensions}

\begin{longtable}[]{@{}
  >{\raggedright\arraybackslash}p{(\linewidth - 14\tabcolsep) * \real{0.1250}}
  >{\raggedright\arraybackslash}p{(\linewidth - 14\tabcolsep) * \real{0.1250}}
  >{\raggedright\arraybackslash}p{(\linewidth - 14\tabcolsep) * \real{0.1250}}
  >{\raggedright\arraybackslash}p{(\linewidth - 14\tabcolsep) * \real{0.1250}}
  >{\raggedright\arraybackslash}p{(\linewidth - 14\tabcolsep) * \real{0.1250}}
  >{\raggedright\arraybackslash}p{(\linewidth - 14\tabcolsep) * \real{0.1250}}
  >{\raggedright\arraybackslash}p{(\linewidth - 14\tabcolsep) * \real{0.1250}}
  >{\raggedright\arraybackslash}p{(\linewidth - 14\tabcolsep) * \real{0.1250}}@{}}
\toprule\noalign{}
\begin{minipage}[b]{\linewidth}\raggedright
RUNSPEC
\end{minipage} & \begin{minipage}[b]{\linewidth}\raggedright
GRID
\end{minipage} & \begin{minipage}[b]{\linewidth}\raggedright
EDIT
\end{minipage} & \begin{minipage}[b]{\linewidth}\raggedright
PROPS
\end{minipage} & \begin{minipage}[b]{\linewidth}\raggedright
REGIONS
\end{minipage} & \begin{minipage}[b]{\linewidth}\raggedright
SOLUTION
\end{minipage} & \begin{minipage}[b]{\linewidth}\raggedright
SUMMARY
\end{minipage} & \begin{minipage}[b]{\linewidth}\raggedright
SCHEDULE
\end{minipage} \\
\midrule\noalign{}
\endhead
\bottomrule\noalign{}
\endlastfoot
\end{longtable}

\subsubsection{Description}\label{description-143}

The TABDIMS keyword defines the maximum number of tables for a given
table type dataset and the maximum number of entries for the various
tables. The commercial simulator combines both the black-oil and
compositional simulator variables on this keyword; however, although all
the parameters are explained below only the black-oil parameters are
used by OPM Flow.

\begin{longtable}[]{@{}
  >{\raggedright\arraybackslash}p{(\linewidth - 6\tabcolsep) * \real{0.2500}}
  >{\raggedright\arraybackslash}p{(\linewidth - 6\tabcolsep) * \real{0.2500}}
  >{\raggedright\arraybackslash}p{(\linewidth - 6\tabcolsep) * \real{0.2500}}
  >{\raggedright\arraybackslash}p{(\linewidth - 6\tabcolsep) * \real{0.2500}}@{}}
\toprule\noalign{}
\begin{minipage}[b]{\linewidth}\raggedright
No.
\end{minipage} & \begin{minipage}[b]{\linewidth}\raggedright
Name
\end{minipage} & \begin{minipage}[b]{\linewidth}\raggedright
Description
\end{minipage} & \begin{minipage}[b]{\linewidth}\raggedright
Default
\end{minipage} \\
\midrule\noalign{}
\endhead
\bottomrule\noalign{}
\endlastfoot
1 & NTSFUN & A positive integer that defines the number of relative
permeability table sets defined in the input deck. The tables are
allocated to different parts of the grid by the SATNUM keyword. & 1 \\
2 & NTPVT & A positive integer that defines the number of fluid property
table sets defined in the input deck. The tables are allocated to
different parts of the grid by the PVTNUM keyword. & 1 \\
3 & NSSFUN & A positive integer that defines the maximum number of
saturation entries in the relative permeability tables defined in the
input deck. & 20 \\
4 & NPPVT & A positive integer that defines the maximum number of
pressure entries in the PVT tables. & 20 \\
5 & NTFIP & A positive integer defining the maximum number of regions in
the FIPNUM region array. Note that this parameter may also be set on the
REGDIMS keyword as well. If NTFIP is set in both places then the maximum
value is used. & 1 \\
6 & NRPVT & A positive integer that defines the maximum number of Rs and
Rv entries in the PVT tables. If the DISGAS and VAPOIL options have not
been activated then this parameter is ignored. & 20 \\
7 & NRVPVT & A positive integer that defines the maximum number of Rv
entries in the PVT tables for the commercial compositional simulator. &
20 \\
8 & NTENDP & A positive integer that defines the maximum number of
saturation end-point depth tables. The end-point depth tables are used
to re-scale the saturation tables as a function of depth as oppose to
being a grid block property. NTENDP may also be specified on the
ENDSCALE keyword, and if specified on both here and on the ENDSCALE
keyword the maximum value of the two is used. & 1 \\
9 & NMEOSR & A positive integer that defines the maximum number of
reservoir equations of states for the commercial compositional
simulator. & 1 \\
10 & NMEOSS & A positive integer that defines the maximum number of
separator or surface equations of states for the commercial
compositional simulator. & 1 \\
11 & MXNFLX & A positive integer defining the maximum number flux
regions in the FLUXNUM region array. MXNFLX can also be defined on the
REGDIMS keywords as well. If MXNFLX is defined both here and on the
REGDIMS keyword then the maximum value of the two is used. & 10 \\
12 & MXNTHR & A positive integer that defines the maximum number of
thermal regions for the commercial compositional simulator. & 1 \\
13 & NTROCC & A positive integer that defines the number of rock
compressibility entries enter by the ROCK keyword defined in the input
deck. The ROCK data is allocated to different parts of the grid by the
either the PVTNUM or ROCKNUM keywords in the REGIONS section. If NTROCC
is defaulted then the PVTNUM array will be used to allocate the ROCK
keyword data to the grid blocks; whereas, if a value is entered then the
ROCKNUM array will be used instead. See also the ROCKOPTS keyword in the
PROPS section that can be used to redefine if the PVTNUM, ROCKNUM or
SATNUM arrays should be employed to allocate the ROCK keyword data. &
1* \\
14 & MXNPMR & A positive integer that defines the maximum number of
pressure maintenance regions for the commercial compositional simulator.
& 0 \\
15 & NTABKT & A positive integer that defines the maximum number of
temperature dependent K-value tables for the when the thermal option is
activated in the commercial compositional simulator. & 0 \\
16 & NTALPHA & A positive integer that defines the maximum number of
transport coefficient tables for the commercial compositional simulator.
& 0 \\
17 & NASPKA & A positive integer that defines the maximum number of
maximum number of entries in the ASPKDAM keyword tables for the
commercial compositional simulator. & 10 \\
18 & MXRAWG & A positive integer that defines the maximum number of
maximum number of entries in the ASPREWG keyword tables for the
commercial compositional simulator. & 10 \\
19 & MXRASO & A positive integer that defines the maximum number of
pressure maintenance regions for the commercial compositional simulator.
& 10 \\
20 & & Not Used & 1* \\
21 & MCASPP & A positive integer that defines the maximum number of
column entries in the ASPPW2D keyword tables for the commercial
compositional simulator. & 5 \\
22 & MRASPP & A positive integer that defines the maximum number of row
entries in the ASPPW2D keyword tables for the commercial compositional
simulator. & 5 \\
23 & MXRATF & A positive integer that defines the maximum number of
entries in the ASPWETF table for the commercial compositional simulator.
& 5 \\
24 & MXNKVT & A positive integer that defines the maximum number of
composition dependent K-value tables for the commercial compositional
simulator. & 0 \\
25 & RESVED & Not Used & 1* \\
Notes: & & & \\
\end{longtable}

\emph{Table 5.44: TABDIMS Keyword Description}

\subsubsection{Example}\label{example-108}

\begin{verbatim}
--
--       NO.     NO.     MAX     MAX     MAX     MAX    E300
--       NTSFUN  NTPVT   NSSFUN  NPPVT   NTFIP   NRPVT  BLANK  NTEND
TABDIMS
         15      9       40      30      1*      1*     1*     1               /
\end{verbatim}

The above example defines number of relative permeability tables to be
15 with a maximum number of rows for each table set to 40, and the
number of PVT tables to be nine with a maximum number of 30 rows per
table.

\subsection{TEMP -- Activate the Temperature Modeling
Option}\label{temp-activate-the-temperature-modeling-option}

\begin{longtable}[]{@{}
  >{\raggedright\arraybackslash}p{(\linewidth - 14\tabcolsep) * \real{0.1250}}
  >{\raggedright\arraybackslash}p{(\linewidth - 14\tabcolsep) * \real{0.1250}}
  >{\raggedright\arraybackslash}p{(\linewidth - 14\tabcolsep) * \real{0.1250}}
  >{\raggedright\arraybackslash}p{(\linewidth - 14\tabcolsep) * \real{0.1250}}
  >{\raggedright\arraybackslash}p{(\linewidth - 14\tabcolsep) * \real{0.1250}}
  >{\raggedright\arraybackslash}p{(\linewidth - 14\tabcolsep) * \real{0.1250}}
  >{\raggedright\arraybackslash}p{(\linewidth - 14\tabcolsep) * \real{0.1250}}
  >{\raggedright\arraybackslash}p{(\linewidth - 14\tabcolsep) * \real{0.1250}}@{}}
\toprule\noalign{}
\begin{minipage}[b]{\linewidth}\raggedright
RUNSPEC
\end{minipage} & \begin{minipage}[b]{\linewidth}\raggedright
GRID
\end{minipage} & \begin{minipage}[b]{\linewidth}\raggedright
EDIT
\end{minipage} & \begin{minipage}[b]{\linewidth}\raggedright
PROPS
\end{minipage} & \begin{minipage}[b]{\linewidth}\raggedright
REGIONS
\end{minipage} & \begin{minipage}[b]{\linewidth}\raggedright
SOLUTION
\end{minipage} & \begin{minipage}[b]{\linewidth}\raggedright
SUMMARY
\end{minipage} & \begin{minipage}[b]{\linewidth}\raggedright
SCHEDULE
\end{minipage} \\
\midrule\noalign{}
\endhead
\bottomrule\noalign{}
\endlastfoot
\end{longtable}

\subsubsection{Description}\label{description-144}

This keyword activates the temperature modeling option.

In OPM Flow's thermal model (THERMAL keyword) enthalpy is conserved and
the energy equations are solved fully implicitly with the black-oil
equations. Whereas, in OPM Flow's temperature model (TEMP keyword)
internal energy is conserved (the work term is neglected and internal
energy is equal to enthalpy) and the energy equations are solved
sequentially after the black-oil equations. Note that the thermal model
can be modified so that internal energy rather than enthalpy is
conserved by setting the command line option
--conserve-inner-energy-thermal=true (the default is false). If
properties like density or viscosity ``strongly'' depend on temperature,
the thermal model is recommended.

There is no data required for this keyword and there is no terminating
``/'' for this keyword.

The temperature option (TEMP keyword) and the thermal option (THERMAL
keyword) are two separate modeling facilities in the commercial
simulator, although some keywords can be used by both options, for
example the RTEMP keyword. OPM Flow's temperature model is not directly
equivalent to either the commercial simulator's black-oil TEMP or
compositional THERMAL formulation.

The energy black-oil implementation in OPM Flow uses a mixture of the
commercial simulators black-oil and the commercial simulators
``compositional thermal'' keywords, as well as some OPM Flow specific
keywords.

\subsubsection{Example}\label{example-109}

\begin{verbatim}
--
--       ACTIVATE THE TEMPERATURE MODELING OPTION (OPM FLOW TEMPERATURE OPTION ONLY)
--
TEMP
\end{verbatim}

The above example activates the temperature modeling option.

\subsection{THERMAL -- Activate the Thermal Modeling
Option}\label{thermal-activate-the-thermal-modeling-option}

\begin{longtable}[]{@{}
  >{\raggedright\arraybackslash}p{(\linewidth - 14\tabcolsep) * \real{0.1250}}
  >{\raggedright\arraybackslash}p{(\linewidth - 14\tabcolsep) * \real{0.1250}}
  >{\raggedright\arraybackslash}p{(\linewidth - 14\tabcolsep) * \real{0.1250}}
  >{\raggedright\arraybackslash}p{(\linewidth - 14\tabcolsep) * \real{0.1250}}
  >{\raggedright\arraybackslash}p{(\linewidth - 14\tabcolsep) * \real{0.1250}}
  >{\raggedright\arraybackslash}p{(\linewidth - 14\tabcolsep) * \real{0.1250}}
  >{\raggedright\arraybackslash}p{(\linewidth - 14\tabcolsep) * \real{0.1250}}
  >{\raggedright\arraybackslash}p{(\linewidth - 14\tabcolsep) * \real{0.1250}}@{}}
\toprule\noalign{}
\begin{minipage}[b]{\linewidth}\raggedright
RUNSPEC
\end{minipage} & \begin{minipage}[b]{\linewidth}\raggedright
GRID
\end{minipage} & \begin{minipage}[b]{\linewidth}\raggedright
EDIT
\end{minipage} & \begin{minipage}[b]{\linewidth}\raggedright
PROPS
\end{minipage} & \begin{minipage}[b]{\linewidth}\raggedright
REGIONS
\end{minipage} & \begin{minipage}[b]{\linewidth}\raggedright
SOLUTION
\end{minipage} & \begin{minipage}[b]{\linewidth}\raggedright
SUMMARY
\end{minipage} & \begin{minipage}[b]{\linewidth}\raggedright
SCHEDULE
\end{minipage} \\
\midrule\noalign{}
\endhead
\bottomrule\noalign{}
\endlastfoot
\end{longtable}

\subsubsection{Description}\label{description-145}

This keyword activates the thermal modeling option.

In OPM Flow's thermal model (THERMAL keyword) enthalpy is conserved and
the energy equations are solved fully implicitly with the black-oil
equations. Whereas, in OPM Flow's temperature model (TEMP keyword)
internal energy is conserved (the work term is neglected and internal
energy is equal to enthalpy) and the energy equations are solved
sequentially after the black-oil equations. Note that the thermal model
can be modified so that internal energy rather than enthalpy is
conserved by setting the command line option
--conserve-inner-energy-thermal=true (the default is false). If
properties like density or viscosity ``strongly'' depend on temperature,
the thermal model is recommended.

There is no data required for this keyword and there is no terminating
``/'' for this keyword.

The temperature option (TEMP keyword) and the thermal option (THERMAL
keyword) are two separate modeling facilities in the commercial
simulator, although some keywords can be used by both options, for
example the RTEMP keyword. OPM Flow's thermal model is not directly
equivalent to either the commercial simulator's black-oil TEMP or
compositional THERMAL formulation.

The energy black-oil implementation in OPM Flow uses a mixture of the
commercial simulators black-oil and the commercial simulators
``compositional thermal'' keywords, as well as some OPM Flow specific
keywords. Keywords specifically associated with both OPM Flow's THERMAL
and the commercial simulators TEMP and THERMAL options are listed in
Table 5.46 for ease of reference.

\begin{longtable}[]{@{}
  >{\raggedright\arraybackslash}p{(\linewidth - 10\tabcolsep) * \real{0.1667}}
  >{\raggedright\arraybackslash}p{(\linewidth - 10\tabcolsep) * \real{0.1667}}
  >{\raggedright\arraybackslash}p{(\linewidth - 10\tabcolsep) * \real{0.1667}}
  >{\raggedright\arraybackslash}p{(\linewidth - 10\tabcolsep) * \real{0.1667}}
  >{\raggedright\arraybackslash}p{(\linewidth - 10\tabcolsep) * \real{0.1667}}
  >{\raggedright\arraybackslash}p{(\linewidth - 10\tabcolsep) * \real{0.1667}}@{}}
\toprule\noalign{}
\begin{minipage}[b]{\linewidth}\raggedright
Section
\end{minipage} & \begin{minipage}[b]{\linewidth}\raggedright
Keyword
\end{minipage} & \begin{minipage}[b]{\linewidth}\raggedright
Function
\end{minipage} & \begin{minipage}[b]{\linewidth}\raggedright
OPM Flow
\end{minipage} & \begin{minipage}[b]{\linewidth}\raggedright
Commercial Simulator
\end{minipage} & \begin{minipage}[b]{\linewidth}\raggedright
\end{minipage} \\
\midrule\noalign{}
\endhead
\bottomrule\noalign{}
\endlastfoot
THERMAL MODEL & TEMP MODEL & THERMAL MODEL & & & \\
RUNSPEC & ROCKDIMS & Thermal Rock Dimensions of Over and Underburden
Rock Types & & & \\
TEMP & Activate the Temperature Modeling Option & & & Black-Oil & \\
THERMAL & Activate the Thermal Modeling Option & & & & \\
GRID & ROCKCON & Thermal Rock Over and Underburden Connection Data & &
& \\
ROCKPROP & Thermal Rock Over and Underburden Property Data & & & & \\
HEATCR & Rock Heat Capacity. & & & & \\
HEATCRT & Rock Heat Capacity Temperature. & & & & \\
THCGAS & Gas Phase Thermal Conductivity. & & & & \\
THCOIL & Oil Phase Thermal Conductivity. & & & & \\
THCONR & Thermal Conductivity of liquids and reservoir rock. & & & & \\
THCONSF & Thermal Conductivity of liquids and reservoir rock scaling
factor applied to THCONR to account for gas saturation. & & & & \\
THCROCK & Rock Thermal Conductivity. & & & & \\
THCSOLID & Solid Phase Thermal Conductivity. & & & & \\
THCWATER & Water Thermal Conductivity. & & & & \\
PROPS & HEATVAP & Thermal Oil Component Heat of Vaporization & & & \\
GASDENT & Gas Density Temperature Coefficients (OPM Flow keyword). & & &
& \\
GASJT & Gas Joule-Thomson Coefficient (OPM Flow keyword). & & & & \\
GASVISCT & Gas Viscosity versus Temperature Functions (OPM Flow
black-oil keyword). & & & & \\
OILDENT & Oil Density Temperature Coefficients (OPM Flow keyword). & & &
& \\
OILJT & Oil Joule-Thomson Coefficient (OPM Flow keyword). & & & & \\
OILVISCT & Oil Viscosity versus Temperature Functions (OPM Flow
black-oil keyword). & & & & \\
RTEMP & Constant Initial Reservoir Temperature. & & & & \\
RTEMPA & Constant Initial Reservoir Temperature. & & & & \\
RTEMPVD & Initial Reservoir Temperature versus Depth. & & & & \\
SPECHA & Thermal Specific Heat of Oil Component & & & & \\
SPECHEAT & Specific Heat of Oil, Water and Gas & & & & \\
SPECROCK & Specific Heat of the Reservoir Rock & & & & \\
TEMPVD & Initial Reservoir Temperature versus Depth. & & & & \\
THANALB & Activate Thermal Analytic Water Density Option & & & & \\
THERMEX1 & Liquid Components Thermal Expansion Coefficient & & & & \\
WATDENT & Oil Density Temperature Coefficients. & & & & \\
WATJT & Water Joule-Thomson Coefficient (OPM Flow keyword). & & & & \\
WATVISCT & Oil Viscosity versus Temperature Function. & & & & \\
REGIONS & THERMNUM & Thermal Region Numbers. & & & \\
SOLUTION & RTEMP & Constant Initial Reservoir Temperature. & & & \\
RTEMPA & Constant Initial Reservoir Temperature. & & & & \\
RTEMPVD & Initial Reservoir Temperature versus Depth. & & & & \\
TEMPI & Initial Reservoir Temperature for All Cells. & & & & \\
TEMPVD & Initial Reservoir Temperature versus Depth. & & & & \\
SCHEDULE & WTEMP & Set An Injection Well's Fluid Temperature & & & \\
WINJTEMP & Define Injection Fluid Thermal Properties & & & & \\
Notes: & & & & & \\
\end{longtable}

Table 5.46: OPM Flow's THERMAL Option Associated Keywords

In thermal runs a producing well's bottom-hole temperature is calculated
based on a weighted average of the temperature in the grid cell
connections open to flow in the producing well, that is the reported
bottom-hole temperature, TBHT, is calculated as:

\begin{equation}\phantomsection\label{eq-5-1}{
{T}_{\mathit{BHT}} = \frac{{\sum }_{i=1}^{M}{W}_{i}{T}_{i}}{{\sum }_{i=1}^{M}{W}_{i}}
}\end{equation}

with

\begin{equation}\phantomsection\label{eq-5-2}{
{W}_{i} = {\sum }_{p}^{N}\left({{ρ}^{r}}_{\mathit{pi}}\right) \left({{q}^{r}}_{\mathit{pi}}\right) {c}_{p}
}\end{equation}

The term
\(\left({{ρ}^{r}}_{\mathit{pi}}\right) \left({{q}^{r}}_{\mathit{pi}}\right) {c}_{p}\)is
the energy rate density (J/(K s)) of phase p,

where:

\(N\) = number of phases,

\(M\) = number of open connections in the well,

\(i\) = the open connection index,

\(p\) = the phase index, oil, water and gas.

\(r\) = indicating that the parameter is evaluated at reservoir
conditions,

\(T\) = temperature (K),

\(q\) = connection flow rate (m3/s),

\(ρ\) = fluid density (kg/m3), and

\(c\) = specific heat capacity (J/(K kg))

(see the SPECHEAT keyword in the PROPS section).

The current implementation makes use of the specific internal energy:

\begin{equation}\phantomsection\label{eq-5-3}{
{e}_{\mathit{pi}} = {c}_{p} ⋅ {T}_{i}
}\end{equation}

derived from the specific enthalpy, hpi

\begin{equation}\phantomsection\label{eq-5-4}{
{e}_{\mathit{pi}} = {h}_{\mathit{pi}} - \frac{{P}_{\mathit{pi}}}{{{ρ}^{r}}_{\mathit{pi}}}
}\end{equation}

where \({P}_{\mathit{pi}}\) is connection grid block pressure of phase
p.

The phase rates at surface conditions\(\left({{q}^{s}}_{pi}\right)\)are
converted to reservoir in situ rates\(\left({{q}^{r}}_{pi}\right)\)using
the phase formation volume factor, Bpi, via:

\begin{equation}\phantomsection\label{eq-5-5}{
{{q}^{r}}_{\mathit{pi}} = \frac{{{q}^{s}}_{\mathit{pi}} }{{{B}^{-1}}_{\mathit{pi}}}
}\end{equation}

And thus equation (5.2) can be simplified to:

\begin{equation}\phantomsection\label{eq-5-6}{
{W}_{i} = {\sum }_{p}^{N}\left({{ρ}^{r}}_{\mathit{pi}}\right) \left(\frac{{{q}^{s}}_{\mathit{pi}} {e}_{\mathit{pi}}}{{{B}^{-1}}_{\mathit{pi}} ⋅ {T}_{i}}\right)
}\end{equation}

\subsubsection{Example}\label{example-110}

\begin{verbatim}
--
--       ACTIVATE THE THERMAL MODELING OPTION (OPM FLOW THERMAL OPTION ONLY)
--
THERMAL
\end{verbatim}

The above example activates the thermal modeling option.

\subsection{TITLE -- Define the Title for the Input
Deck}\label{title-define-the-title-for-the-input-deck}

\begin{longtable}[]{@{}
  >{\raggedright\arraybackslash}p{(\linewidth - 14\tabcolsep) * \real{0.1250}}
  >{\raggedright\arraybackslash}p{(\linewidth - 14\tabcolsep) * \real{0.1250}}
  >{\raggedright\arraybackslash}p{(\linewidth - 14\tabcolsep) * \real{0.1250}}
  >{\raggedright\arraybackslash}p{(\linewidth - 14\tabcolsep) * \real{0.1250}}
  >{\raggedright\arraybackslash}p{(\linewidth - 14\tabcolsep) * \real{0.1250}}
  >{\raggedright\arraybackslash}p{(\linewidth - 14\tabcolsep) * \real{0.1250}}
  >{\raggedright\arraybackslash}p{(\linewidth - 14\tabcolsep) * \real{0.1250}}
  >{\raggedright\arraybackslash}p{(\linewidth - 14\tabcolsep) * \real{0.1250}}@{}}
\toprule\noalign{}
\begin{minipage}[b]{\linewidth}\raggedright
RUNSPEC
\end{minipage} & \begin{minipage}[b]{\linewidth}\raggedright
GRID
\end{minipage} & \begin{minipage}[b]{\linewidth}\raggedright
EDIT
\end{minipage} & \begin{minipage}[b]{\linewidth}\raggedright
PROPS
\end{minipage} & \begin{minipage}[b]{\linewidth}\raggedright
REGIONS
\end{minipage} & \begin{minipage}[b]{\linewidth}\raggedright
SOLUTION
\end{minipage} & \begin{minipage}[b]{\linewidth}\raggedright
SUMMARY
\end{minipage} & \begin{minipage}[b]{\linewidth}\raggedright
SCHEDULE
\end{minipage} \\
\midrule\noalign{}
\endhead
\bottomrule\noalign{}
\endlastfoot
\end{longtable}

\subsubsection{Description}\label{description-146}

The TITLE keyword defines the title for the input deck. The title text
will be printed on all reports so as to act as a reference for the run.

\begin{longtable}[]{@{}
  >{\raggedright\arraybackslash}p{(\linewidth - 6\tabcolsep) * \real{0.2500}}
  >{\raggedright\arraybackslash}p{(\linewidth - 6\tabcolsep) * \real{0.2500}}
  >{\raggedright\arraybackslash}p{(\linewidth - 6\tabcolsep) * \real{0.2500}}
  >{\raggedright\arraybackslash}p{(\linewidth - 6\tabcolsep) * \real{0.2500}}@{}}
\toprule\noalign{}
\begin{minipage}[b]{\linewidth}\raggedright
No.
\end{minipage} & \begin{minipage}[b]{\linewidth}\raggedright
Name
\end{minipage} & \begin{minipage}[b]{\linewidth}\raggedright
Description
\end{minipage} & \begin{minipage}[b]{\linewidth}\raggedright
Default
\end{minipage} \\
\midrule\noalign{}
\endhead
\bottomrule\noalign{}
\endlastfoot
1 & TITLE & A character string that defines the title for the input deck
& None \\
Notes: & & & \\
\end{longtable}

\emph{Table 5.47: TITLE Keyword Description}

\begin{tcolorbox}[enhanced jigsaw, title=\textcolor{quarto-callout-note-color}{\faInfo}\hspace{0.5em}{Note}, coltitle=black, breakable, toprule=.15mm, colbacktitle=quarto-callout-note-color!10!white, leftrule=.75mm, opacitybacktitle=0.6, colback=white, bottomtitle=1mm, toptitle=1mm, titlerule=0mm, left=2mm, colframe=quarto-callout-note-color-frame, arc=.35mm, rightrule=.15mm, bottomrule=.15mm, opacityback=0]

It is good practice to include the name of the input file in the tittle
(without the extension) for when cross checking results from multiple
cases.

\end{tcolorbox}

\subsubsection{Example}\label{example-111}

\begin{verbatim}
--
--       DEFINE THE TITLE FOR THE RUN
TITLE
SPE01-THEM01-OPM1810-R01 - OPM THERMAL OPTION RUN
\end{verbatim}

The above example defines the title for the run to be
``SPE01-THEM01-OPM1810-R01 - OPM THERMAL OPTION RUN''.

\subsection{TRACERS -- Activate Tracer Options and Set Tracer Array
Dimensions}\label{tracers-activate-tracer-options-and-set-tracer-array-dimensions}

\begin{longtable}[]{@{}
  >{\raggedright\arraybackslash}p{(\linewidth - 14\tabcolsep) * \real{0.1250}}
  >{\raggedright\arraybackslash}p{(\linewidth - 14\tabcolsep) * \real{0.1250}}
  >{\raggedright\arraybackslash}p{(\linewidth - 14\tabcolsep) * \real{0.1250}}
  >{\raggedright\arraybackslash}p{(\linewidth - 14\tabcolsep) * \real{0.1250}}
  >{\raggedright\arraybackslash}p{(\linewidth - 14\tabcolsep) * \real{0.1250}}
  >{\raggedright\arraybackslash}p{(\linewidth - 14\tabcolsep) * \real{0.1250}}
  >{\raggedright\arraybackslash}p{(\linewidth - 14\tabcolsep) * \real{0.1250}}
  >{\raggedright\arraybackslash}p{(\linewidth - 14\tabcolsep) * \real{0.1250}}@{}}
\toprule\noalign{}
\begin{minipage}[b]{\linewidth}\raggedright
RUNSPEC
\end{minipage} & \begin{minipage}[b]{\linewidth}\raggedright
GRID
\end{minipage} & \begin{minipage}[b]{\linewidth}\raggedright
EDIT
\end{minipage} & \begin{minipage}[b]{\linewidth}\raggedright
PROPS
\end{minipage} & \begin{minipage}[b]{\linewidth}\raggedright
REGIONS
\end{minipage} & \begin{minipage}[b]{\linewidth}\raggedright
SOLUTION
\end{minipage} & \begin{minipage}[b]{\linewidth}\raggedright
SUMMARY
\end{minipage} & \begin{minipage}[b]{\linewidth}\raggedright
SCHEDULE
\end{minipage} \\
\midrule\noalign{}
\endhead
\bottomrule\noalign{}
\endlastfoot
\end{longtable}

\subsubsection{Description}\label{description-147}

The TRACERS keyword defines the number of tracers in the model and the
various passive tracer tracking options.

\begin{longtable}[]{@{}
  >{\raggedright\arraybackslash}p{(\linewidth - 6\tabcolsep) * \real{0.2500}}
  >{\raggedright\arraybackslash}p{(\linewidth - 6\tabcolsep) * \real{0.2500}}
  >{\raggedright\arraybackslash}p{(\linewidth - 6\tabcolsep) * \real{0.2500}}
  >{\raggedright\arraybackslash}p{(\linewidth - 6\tabcolsep) * \real{0.2500}}@{}}
\toprule\noalign{}
\begin{minipage}[b]{\linewidth}\raggedright
No.
\end{minipage} & \begin{minipage}[b]{\linewidth}\raggedright
Name
\end{minipage} & \begin{minipage}[b]{\linewidth}\raggedright
Description
\end{minipage} & \begin{minipage}[b]{\linewidth}\raggedright
Default
\end{minipage} \\
\midrule\noalign{}
\endhead
\bottomrule\noalign{}
\endlastfoot
1 & MXOILTR & A positive integer defining the maximum number of passive
oil tracers defined using the TRACER keyword. & 0 \\
2 & MXWATTR & A positive integer defining the maximum number of passive
water tracers defined using the TRACER keyword. & 0 \\
3 & MXGASTR & A positive integer defining the maximum number of passive
gas tracers defined using the TRACER keyword. & 0 \\
4 & MXENVTR & A positive integer defining the maximum number of passive
environmental tracers defined using the TRACER keyword. Environmental
tracers are used with the Environment Tracer model that takes into
account tracer adsorption and decay. Only the default value of zero is
currently supported. & 0 \\
5 & DIFFOPT & A character string defining the numerical diffusion option
for tracer tracking runs that should be set to: Only the default value
of NODIFF is supported & NODIFF \\
6 & MXITRTR & A positive integer defining the maximum number of
non-linear iterations to be used when the tracer option is activated. &
12 \\
7 & MNITRTR & A positive integer defining the minimum number of
non-linear iterations to be used when the tracer option is activated. &
1 \\
8 & NONLIN & A character string stating if passive tracers as should be
linear (NO) or non-linear (YES). Only the default value of NO is
supported. & No \\
9 & LNCONFAC & A real value defining the initial linear convergence
factor. The default value of 1* means the parameter will not be
utilized. & 1* \\
10 & NLCONFAC & A real value defining the initial non-linear convergence
factor. The default value of 1* means the parameter will not be
utilized. & 1* \\
11 & CONFAC & A real value defining the LNCONFAC and NLCONFAC
convergence factors to be used after the initial convergence factor has
been applied. & 1.0 \\
12 & NUMCONF & A positive integer defining the maximum number of times
CONFAC can be used. & 0 \\
Notes: & & & \\
\end{longtable}

\emph{Table 5.48: TRACERS Keyword Description}

See also the TRACER keyword in the PROPS section that defines the
individual tracers.

\subsubsection{Example}\label{example-112}

\begin{verbatim}
--
--       NO OIL  NO WAT  NO GAS  NO ENV  DIFF    MAX    MIN    TRACER
--       TRACERS TRACERS TRACERS TRACERS CONTL   NONLIN NONLIN NONLIN
TRACERS
         0       7       1       0      'NODIFF' 1*     1*     1*              /
\end{verbatim}

The above example defines seven tracers in the water phase and one
tracer in the gas phase.

\subsection{TRPLPORO -- Activate the Triple Porosity Model
Option}\label{trplporo-activate-the-triple-porosity-model-option}

\begin{longtable}[]{@{}
  >{\raggedright\arraybackslash}p{(\linewidth - 14\tabcolsep) * \real{0.1250}}
  >{\raggedright\arraybackslash}p{(\linewidth - 14\tabcolsep) * \real{0.1250}}
  >{\raggedright\arraybackslash}p{(\linewidth - 14\tabcolsep) * \real{0.1250}}
  >{\raggedright\arraybackslash}p{(\linewidth - 14\tabcolsep) * \real{0.1250}}
  >{\raggedright\arraybackslash}p{(\linewidth - 14\tabcolsep) * \real{0.1250}}
  >{\raggedright\arraybackslash}p{(\linewidth - 14\tabcolsep) * \real{0.1250}}
  >{\raggedright\arraybackslash}p{(\linewidth - 14\tabcolsep) * \real{0.1250}}
  >{\raggedright\arraybackslash}p{(\linewidth - 14\tabcolsep) * \real{0.1250}}@{}}
\toprule\noalign{}
\begin{minipage}[b]{\linewidth}\raggedright
RUNSPEC
\end{minipage} & \begin{minipage}[b]{\linewidth}\raggedright
GRID
\end{minipage} & \begin{minipage}[b]{\linewidth}\raggedright
EDIT
\end{minipage} & \begin{minipage}[b]{\linewidth}\raggedright
PROPS
\end{minipage} & \begin{minipage}[b]{\linewidth}\raggedright
REGIONS
\end{minipage} & \begin{minipage}[b]{\linewidth}\raggedright
SOLUTION
\end{minipage} & \begin{minipage}[b]{\linewidth}\raggedright
SUMMARY
\end{minipage} & \begin{minipage}[b]{\linewidth}\raggedright
SCHEDULE
\end{minipage} \\
\midrule\noalign{}
\endhead
\bottomrule\noalign{}
\endlastfoot
\end{longtable}

\subsubsection{Description}\label{description-148}

The TRPLPORO keyword activates the Triple Porosity Model option that
models matrix, fractures and vuggy porosity for carbonate reservoirs,
and specifies the number of matrix porosity systems

This keyword is not supported by OPM Flow but would change the results
if supported so the simulation will be stopped.

\begin{longtable}[]{@{}
  >{\raggedright\arraybackslash}p{(\linewidth - 6\tabcolsep) * \real{0.2500}}
  >{\raggedright\arraybackslash}p{(\linewidth - 6\tabcolsep) * \real{0.2500}}
  >{\raggedright\arraybackslash}p{(\linewidth - 6\tabcolsep) * \real{0.2500}}
  >{\raggedright\arraybackslash}p{(\linewidth - 6\tabcolsep) * \real{0.2500}}@{}}
\toprule\noalign{}
\begin{minipage}[b]{\linewidth}\raggedright
No.
\end{minipage} & \begin{minipage}[b]{\linewidth}\raggedright
Name
\end{minipage} & \begin{minipage}[b]{\linewidth}\raggedright
Description
\end{minipage} & \begin{minipage}[b]{\linewidth}\raggedright
Default
\end{minipage} \\
\midrule\noalign{}
\endhead
\bottomrule\noalign{}
\endlastfoot
1 & TRPLPORO & A positive integer value that specifies the number of
matrix porosity systems in the model. TRPLPORO should be set to either:
& 1 \\
Notes: & & & \\
\end{longtable}

\emph{Table 5.49: TRPLPORO Keyword Description}

Note the keyword cannot be used in conjunction with the NMATRIX keyword,
which is also in the RUNSPEC section.

\subsubsection{Example}\label{example-113}

\begin{verbatim}
--
--       TRPLPORO
--       OPTION
TRPLPORO
         3                                                        /
\end{verbatim}

The above example activates the Triple Porosity Model option and
specifies the porosity system is matrix, connected vugs, and isolated
vugs.

\subsection{UDADIMS -- Define the Dimensions of the User Defined
Arguments}\label{udadims-define-the-dimensions-of-the-user-defined-arguments}

\begin{longtable}[]{@{}
  >{\raggedright\arraybackslash}p{(\linewidth - 14\tabcolsep) * \real{0.1250}}
  >{\raggedright\arraybackslash}p{(\linewidth - 14\tabcolsep) * \real{0.1250}}
  >{\raggedright\arraybackslash}p{(\linewidth - 14\tabcolsep) * \real{0.1250}}
  >{\raggedright\arraybackslash}p{(\linewidth - 14\tabcolsep) * \real{0.1250}}
  >{\raggedright\arraybackslash}p{(\linewidth - 14\tabcolsep) * \real{0.1250}}
  >{\raggedright\arraybackslash}p{(\linewidth - 14\tabcolsep) * \real{0.1250}}
  >{\raggedright\arraybackslash}p{(\linewidth - 14\tabcolsep) * \real{0.1250}}
  >{\raggedright\arraybackslash}p{(\linewidth - 14\tabcolsep) * \real{0.1250}}@{}}
\toprule\noalign{}
\begin{minipage}[b]{\linewidth}\raggedright
RUNSPEC
\end{minipage} & \begin{minipage}[b]{\linewidth}\raggedright
GRID
\end{minipage} & \begin{minipage}[b]{\linewidth}\raggedright
EDIT
\end{minipage} & \begin{minipage}[b]{\linewidth}\raggedright
PROPS
\end{minipage} & \begin{minipage}[b]{\linewidth}\raggedright
REGIONS
\end{minipage} & \begin{minipage}[b]{\linewidth}\raggedright
SOLUTION
\end{minipage} & \begin{minipage}[b]{\linewidth}\raggedright
SUMMARY
\end{minipage} & \begin{minipage}[b]{\linewidth}\raggedright
SCHEDULE
\end{minipage} \\
\midrule\noalign{}
\endhead
\bottomrule\noalign{}
\endlastfoot
\end{longtable}

\subsubsection{Description}\label{description-149}

This keyword defines the dimensions of the User Defined Arguments
(``UDA'') used by OPM Flow that can be applied to various connection,
group, and well keywords in the SCHEDULE section. UDAs are defined by
the UDQ keyword that is used to specify values to be constants, SUMMARY
variables, as defined in SUMMARY section, or a formula using various
mathematical functions together with constants and SUMMARY variables.

Although this keyword is read by OPM Flow and the ACTION and UDQ
computational logic and calculations have been implemented, one should
used caution using this facility as it may result in OPM Flow aborting.

\begin{longtable}[]{@{}
  >{\raggedright\arraybackslash}p{(\linewidth - 6\tabcolsep) * \real{0.2500}}
  >{\raggedright\arraybackslash}p{(\linewidth - 6\tabcolsep) * \real{0.2500}}
  >{\raggedright\arraybackslash}p{(\linewidth - 6\tabcolsep) * \real{0.2500}}
  >{\raggedright\arraybackslash}p{(\linewidth - 6\tabcolsep) * \real{0.2500}}@{}}
\toprule\noalign{}
\begin{minipage}[b]{\linewidth}\raggedright
No.
\end{minipage} & \begin{minipage}[b]{\linewidth}\raggedright
Name
\end{minipage} & \begin{minipage}[b]{\linewidth}\raggedright
Description
\end{minipage} & \begin{minipage}[b]{\linewidth}\raggedright
Default
\end{minipage} \\
\midrule\noalign{}
\endhead
\bottomrule\noalign{}
\endlastfoot
1 & NMUDA & NMUDA is a positive integer that defines the number of
arguments in a SCHEDULE section keyword, that are replaced by numeric
UDQ values. & 0 \\
2 & IGNORED & Not used and should be defaulted. & 1* \\
3 & MXUDA & MXUDA is a positive integer that defines the maximum number
of unique arguments in a keyword that are replaced by numeric UDA
values. Note that MXUDA differs from NMUDA, for example: As MXUDA's
default value is 100 then this only needs to be increased where the same
UDA is used more than 100 times. & 100 \\
Notes: & & & \\
\end{longtable}

\emph{Table 5.50: UDADIMS Keyword Description}

Note that OPM Flow has a more restricted UDQ feature set than the
commercial simulator, so not all options and functions are currently
available.

\subsubsection{Example}\label{example-114}

\begin{verbatim}
--
--       USER DEFINED ARGUMENT DIMENSIONS
--       NO.     NOT     TOTAL
--       ARGS    USED    UDQ
UDADIMS
         10       1*     10                                                    /
\end{verbatim}

In the above example both NMUDA and MXUDA are set equal to ten.

\subsection{UDQDIMS -- Define the Dimensions of the User Defined UDQ
Feature}\label{udqdims-define-the-dimensions-of-the-user-defined-udq-feature}

\begin{longtable}[]{@{}
  >{\raggedright\arraybackslash}p{(\linewidth - 14\tabcolsep) * \real{0.1250}}
  >{\raggedright\arraybackslash}p{(\linewidth - 14\tabcolsep) * \real{0.1250}}
  >{\raggedright\arraybackslash}p{(\linewidth - 14\tabcolsep) * \real{0.1250}}
  >{\raggedright\arraybackslash}p{(\linewidth - 14\tabcolsep) * \real{0.1250}}
  >{\raggedright\arraybackslash}p{(\linewidth - 14\tabcolsep) * \real{0.1250}}
  >{\raggedright\arraybackslash}p{(\linewidth - 14\tabcolsep) * \real{0.1250}}
  >{\raggedright\arraybackslash}p{(\linewidth - 14\tabcolsep) * \real{0.1250}}
  >{\raggedright\arraybackslash}p{(\linewidth - 14\tabcolsep) * \real{0.1250}}@{}}
\toprule\noalign{}
\begin{minipage}[b]{\linewidth}\raggedright
RUNSPEC
\end{minipage} & \begin{minipage}[b]{\linewidth}\raggedright
GRID
\end{minipage} & \begin{minipage}[b]{\linewidth}\raggedright
EDIT
\end{minipage} & \begin{minipage}[b]{\linewidth}\raggedright
PROPS
\end{minipage} & \begin{minipage}[b]{\linewidth}\raggedright
REGIONS
\end{minipage} & \begin{minipage}[b]{\linewidth}\raggedright
SOLUTION
\end{minipage} & \begin{minipage}[b]{\linewidth}\raggedright
SUMMARY
\end{minipage} & \begin{minipage}[b]{\linewidth}\raggedright
SCHEDULE
\end{minipage} \\
\midrule\noalign{}
\endhead
\bottomrule\noalign{}
\endlastfoot
\end{longtable}

\subsubsection{Description}\label{description-150}

This keyword defines the dimensions associated with the UDQ keyword used
in OPM Flow to calculate various user defined values in the SCHEDULE
section. The UDQ keyword defined variables can be constants, SUMMARY
variables, as defined in the SUMMARY section, or a formula using various
mathematical functions together with constants and SUMMARY variables.

Although this keyword is read by OPM Flow and the ACTION and UDQ
computational logic and calculations have been implemented, one should
used caution using this facility as it may result in OPM Flow aborting.

\begin{longtable}[]{@{}
  >{\raggedright\arraybackslash}p{(\linewidth - 6\tabcolsep) * \real{0.2500}}
  >{\raggedright\arraybackslash}p{(\linewidth - 6\tabcolsep) * \real{0.2500}}
  >{\raggedright\arraybackslash}p{(\linewidth - 6\tabcolsep) * \real{0.2500}}
  >{\raggedright\arraybackslash}p{(\linewidth - 6\tabcolsep) * \real{0.2500}}@{}}
\toprule\noalign{}
\begin{minipage}[b]{\linewidth}\raggedright
No.
\end{minipage} & \begin{minipage}[b]{\linewidth}\raggedright
Name
\end{minipage} & \begin{minipage}[b]{\linewidth}\raggedright
Description
\end{minipage} & \begin{minipage}[b]{\linewidth}\raggedright
Default
\end{minipage} \\
\midrule\noalign{}
\endhead
\bottomrule\noalign{}
\endlastfoot
1 & MXFUNS & A positive integer that defines the maximum number of
functions that can be included when defining a UDQ definition. This
should also include any brackets that will be used in the UDQ
definition. & 16 \\
2 & MXITEMS & MXITEMS is a positive integer that defines the maximum
number of ITEMS allowed in an UDQ definition. & 16 \\
3 & MXUDC & MXUDC is a positive integer that defines the maximum number
of user defined CONNECTION quantities allowed in an UDQ definition. &
0 \\
4 & MXUDF & MXUDF is a positive integer that defines the maximum number
of user defined FIELD quantities allowed in an UDQ definition. & 0 \\
5 & MXUDG & MXUDG is a positive integer that defines the maximum number
of user defined GROUP quantities allowed in an UDQ definition. & 0 \\
6 & MXUDR & MXUDR is a positive integer that defines the maximum number
of user defined REGION quantities allowed in an UDQ definition. & 0 \\
7 & MXUDS & MXUDS is a positive integer that defines the maximum number
of user defined SEGMENT quantities allowed in an UDQ definition. & 0 \\
8 & MXUDW & MXUDW is a positive integer that defines the maximum number
of user defined WELL quantities allowed in an UDQ definition. & 0 \\
9 & MXUDA & MXUDA is a positive integer that defines the maximum number
of user defined AQUIFER quantities allowed in an UDQ definition. & 0 \\
10 & MXUDB & MXUDB is a positive integer that defines the maximum number
of user defined BLOCK quantities allowed in an UDQ definition. & 0 \\
11 & RSEED & RSEED is a character string that determines if a new random
number seed should be generated for restart runs for use in the UDQ
functions RANDN, RANDU RRNDN and RRNDU. If RSEED is set to Y than a new
seed will be generated and if set to the default value of N or 1* then
the same seed of the ``base'' simulation will be employed. See also the
RSEED integer variable on the UDQPARAM keyword in the RUNSPEC section to
set the random number seed for the current run. This feature is not
supported by OPM Flow. & N \\
Notes: & & & \\
\end{longtable}

\emph{Table 5.51: UDQDIMS Keyword Description}

Note that OPM Flow has a more restricted UDQ feature set than the
commercial simulator, so not all options and functions are available.

\subsubsection{Example}\label{example-115}

\begin{verbatim}
--
--       USER DEFINED ARGUMENT DIMENSIONS FACILITY
--       MAX     MAX     MAX    MAX    MAX    MAX   MAX   MAX   MAX  MAX    RAND
--       FUNCS   ITEMS   CONNS  FIELD  GROUP  REGS  SEGTM WELL  AQUF BLCKS  OPT
UDQDIMS
         50      25      0      50     50     0     0     0     0    0      N  /
\end{verbatim}

In this case the maximum number of functions that can be included when
defining a UDQ definition is set to 50, maximum number of items allowed
in an UDQ definition is 25, the maximum number of user defined field
quantities allowed in an UDQ definition is 50, and the maximum number of
user defined group quantities allowed in an UDQ definition is also 50.
All other parameters are defaulted including the RSEED variable (the
same seed of the ``base'' simulation will be employed).

\subsection{UDQPARAM -- Define Parameters for the User Defined Quantity
Feature}\label{udqparam-define-parameters-for-the-user-defined-quantity-feature}

\begin{longtable}[]{@{}
  >{\raggedright\arraybackslash}p{(\linewidth - 14\tabcolsep) * \real{0.1250}}
  >{\raggedright\arraybackslash}p{(\linewidth - 14\tabcolsep) * \real{0.1250}}
  >{\raggedright\arraybackslash}p{(\linewidth - 14\tabcolsep) * \real{0.1250}}
  >{\raggedright\arraybackslash}p{(\linewidth - 14\tabcolsep) * \real{0.1250}}
  >{\raggedright\arraybackslash}p{(\linewidth - 14\tabcolsep) * \real{0.1250}}
  >{\raggedright\arraybackslash}p{(\linewidth - 14\tabcolsep) * \real{0.1250}}
  >{\raggedright\arraybackslash}p{(\linewidth - 14\tabcolsep) * \real{0.1250}}
  >{\raggedright\arraybackslash}p{(\linewidth - 14\tabcolsep) * \real{0.1250}}@{}}
\toprule\noalign{}
\begin{minipage}[b]{\linewidth}\raggedright
RUNSPEC
\end{minipage} & \begin{minipage}[b]{\linewidth}\raggedright
GRID
\end{minipage} & \begin{minipage}[b]{\linewidth}\raggedright
EDIT
\end{minipage} & \begin{minipage}[b]{\linewidth}\raggedright
PROPS
\end{minipage} & \begin{minipage}[b]{\linewidth}\raggedright
REGIONS
\end{minipage} & \begin{minipage}[b]{\linewidth}\raggedright
SOLUTION
\end{minipage} & \begin{minipage}[b]{\linewidth}\raggedright
SUMMARY
\end{minipage} & \begin{minipage}[b]{\linewidth}\raggedright
SCHEDULE
\end{minipage} \\
\midrule\noalign{}
\endhead
\bottomrule\noalign{}
\endlastfoot
\end{longtable}

\subsubsection{Description}\label{description-151}

This keyword defines the dimensions of the User Defined Arguments
(``UDA'') used by OPM Flow that can be applied to various connection,
group, and well keywords in the SCHEDULE section. UDAs are defined by
the UDQ keyword that is used to specify values to be constants, SUMMARY
variables, as defined in SUMMARY section, or a formula using various
mathematical functions together with constants and SUMMARY variables.

Although this keyword is read by OPM Flow and the ACTION and UDQ
computational logic and calculations have been implemented, one should
used caution using this facility as it may result in OPM Flow aborting.

\begin{longtable}[]{@{}
  >{\raggedright\arraybackslash}p{(\linewidth - 6\tabcolsep) * \real{0.2500}}
  >{\raggedright\arraybackslash}p{(\linewidth - 6\tabcolsep) * \real{0.2500}}
  >{\raggedright\arraybackslash}p{(\linewidth - 6\tabcolsep) * \real{0.2500}}
  >{\raggedright\arraybackslash}p{(\linewidth - 6\tabcolsep) * \real{0.2500}}@{}}
\toprule\noalign{}
\begin{minipage}[b]{\linewidth}\raggedright
No.
\end{minipage} & \begin{minipage}[b]{\linewidth}\raggedright
Name
\end{minipage} & \begin{minipage}[b]{\linewidth}\raggedright
Description
\end{minipage} & \begin{minipage}[b]{\linewidth}\raggedright
Default
\end{minipage} \\
\midrule\noalign{}
\endhead
\bottomrule\noalign{}
\endlastfoot
1 & RSEED & RSEED is a positive integer greater than zero that sets a
new random number seed for use in the UDQ functions RANDN, RANDU RRNDN
and RRNDU. See also the RSEED character variable on the UDQDIMS keyword
in the RUNSPEC section to default the random number seed for a restart
run. This feature is not supported by OPM Flow. & 1 \\
2 & RANGE & RANGE is a real positive value greater than or equal to one
and less than or equal to 1.0 x 1020, that sets the absolute range for
the user defined quantities. The default value of 1 x 1020 sets the
range from -1 x 1020 to +1 x 1020. & 1 x 1020 \\
3 & DEFAULT & DEFAULT is real value that is the default numerical value
given to undefined UDQ variables and should be in the same range as
RANGE. & 0.0 \\
4 & TOLUDQ & TOLUDQ a real positive number greater than zero and less
than one that defines the tolerance used to determine if two real values
are equal. Floating-point numbers (as implemented in computers) are
never exact, one cannot compare floating point numbers for exact
equality. Thus, TOLUDQ defines a tolerance. For example, the default
value of 1 x 10-4 means that if the difference between two real values
is less than 1 x 10-4 then the values are considered equal. & 1 x
10-4 \\
Notes: & & & \\
\end{longtable}

\emph{Table 5.52: UDQPARAM Keyword Description}

Note that OPM Flow has a more restricted UDQ feature set than the
commercial simulator, so not all options and functions are available.

\subsubsection{Example}\label{example-116}

\begin{verbatim}
--
--       USER DEFINED DEFAULT VALUES
--       SEED    RANGE    UNDEFINED   COMPARISON
--       INTG    -AND+    VALUE       TOLERANCE
UDQPARAM
         1      1.0E20    0.0         1.0E-4                                   /
\end{verbatim}

The example explicitly sets the default values for all four variables on
the UDAPARAM keyword, namely the random seed to one, the range to 1 x
1020, the undefined UDQ variables to zero, and the comparison tolerance
to 1.0 x 10-4.

\subsection{UDTDIMS -- Define the Dimensions of the User Defined
Tables}\label{udtdims-define-the-dimensions-of-the-user-defined-tables}

\begin{longtable}[]{@{}
  >{\raggedright\arraybackslash}p{(\linewidth - 14\tabcolsep) * \real{0.1250}}
  >{\raggedright\arraybackslash}p{(\linewidth - 14\tabcolsep) * \real{0.1250}}
  >{\raggedright\arraybackslash}p{(\linewidth - 14\tabcolsep) * \real{0.1250}}
  >{\raggedright\arraybackslash}p{(\linewidth - 14\tabcolsep) * \real{0.1250}}
  >{\raggedright\arraybackslash}p{(\linewidth - 14\tabcolsep) * \real{0.1250}}
  >{\raggedright\arraybackslash}p{(\linewidth - 14\tabcolsep) * \real{0.1250}}
  >{\raggedright\arraybackslash}p{(\linewidth - 14\tabcolsep) * \real{0.1250}}
  >{\raggedright\arraybackslash}p{(\linewidth - 14\tabcolsep) * \real{0.1250}}@{}}
\toprule\noalign{}
\begin{minipage}[b]{\linewidth}\raggedright
RUNSPEC
\end{minipage} & \begin{minipage}[b]{\linewidth}\raggedright
GRID
\end{minipage} & \begin{minipage}[b]{\linewidth}\raggedright
EDIT
\end{minipage} & \begin{minipage}[b]{\linewidth}\raggedright
PROPS
\end{minipage} & \begin{minipage}[b]{\linewidth}\raggedright
REGIONS
\end{minipage} & \begin{minipage}[b]{\linewidth}\raggedright
SOLUTION
\end{minipage} & \begin{minipage}[b]{\linewidth}\raggedright
SUMMARY
\end{minipage} & \begin{minipage}[b]{\linewidth}\raggedright
SCHEDULE
\end{minipage} \\
\midrule\noalign{}
\endhead
\bottomrule\noalign{}
\endlastfoot
\end{longtable}

\subsubsection{Description}\label{description-152}

This keyword defines the dimensions of the User Defined Tables (``UDT'')
used by OPM Flow that can be used as lookup tables when assigning values
to User Defined Quantities (``UDQ'') using the UDQ keyword in the
SCHEDULE section. UDTs are defined by the UDT keyword in the SCHEDULE
section.

\begin{longtable}[]{@{}
  >{\raggedright\arraybackslash}p{(\linewidth - 6\tabcolsep) * \real{0.2500}}
  >{\raggedright\arraybackslash}p{(\linewidth - 6\tabcolsep) * \real{0.2500}}
  >{\raggedright\arraybackslash}p{(\linewidth - 6\tabcolsep) * \real{0.2500}}
  >{\raggedright\arraybackslash}p{(\linewidth - 6\tabcolsep) * \real{0.2500}}@{}}
\toprule\noalign{}
\begin{minipage}[b]{\linewidth}\raggedright
No.
\end{minipage} & \begin{minipage}[b]{\linewidth}\raggedright
Name
\end{minipage} & \begin{minipage}[b]{\linewidth}\raggedright
Description
\end{minipage} & \begin{minipage}[b]{\linewidth}\raggedright
Default
\end{minipage} \\
\midrule\noalign{}
\endhead
\bottomrule\noalign{}
\endlastfoot
1 & MXUDT & MXUDT is a positive integer that defines the maximum number
of User Defined Tables & 0 \\
2 & NUDT & NUDT is a positive integer that defines the maximum number of
rows in any given User Defined Table. & 0 \\
3 & MXINTP & MXINTP is a positive integer that defines the maximum
number of interpolation points allowed in any given dimension. & 0 \\
4 & MXDIMS & MXDIMS is a positive integer that defines the maximum
number of dimensions in any given User Defined Table. Only one
dimensional tables are currently supported by OPM Flow. & 0 \\
Notes: & & & \\
\end{longtable}

\emph{Table 5.53: UDTDIMS Keyword Description}

\subsubsection{Example}\label{example-117}

\begin{verbatim}
--
--       USER DEFINED TABLE DIMENSIONS
--
--       MAX     MAX   MAX     MAX
--       TABLES  ROWS  INTPOL  DIMS
UDTDIMS
         3       20    3       2                                               /
\end{verbatim}

In the above example the maximum number of UDT tables is set to three
and the maximum number of rows for each table is 20, the maximum number
of interpolation points in any given dimension is set to three and the
maximum number of dimensions is defined as two.

\subsection{UNCODHMD -- Activate History Match Gradient Unencoded
Output}\label{uncodhmd-activate-history-match-gradient-unencoded-output}

\begin{longtable}[]{@{}
  >{\raggedright\arraybackslash}p{(\linewidth - 14\tabcolsep) * \real{0.1250}}
  >{\raggedright\arraybackslash}p{(\linewidth - 14\tabcolsep) * \real{0.1250}}
  >{\raggedright\arraybackslash}p{(\linewidth - 14\tabcolsep) * \real{0.1250}}
  >{\raggedright\arraybackslash}p{(\linewidth - 14\tabcolsep) * \real{0.1250}}
  >{\raggedright\arraybackslash}p{(\linewidth - 14\tabcolsep) * \real{0.1250}}
  >{\raggedright\arraybackslash}p{(\linewidth - 14\tabcolsep) * \real{0.1250}}
  >{\raggedright\arraybackslash}p{(\linewidth - 14\tabcolsep) * \real{0.1250}}
  >{\raggedright\arraybackslash}p{(\linewidth - 14\tabcolsep) * \real{0.1250}}@{}}
\toprule\noalign{}
\begin{minipage}[b]{\linewidth}\raggedright
RUNSPEC
\end{minipage} & \begin{minipage}[b]{\linewidth}\raggedright
GRID
\end{minipage} & \begin{minipage}[b]{\linewidth}\raggedright
EDIT
\end{minipage} & \begin{minipage}[b]{\linewidth}\raggedright
PROPS
\end{minipage} & \begin{minipage}[b]{\linewidth}\raggedright
REGIONS
\end{minipage} & \begin{minipage}[b]{\linewidth}\raggedright
SOLUTION
\end{minipage} & \begin{minipage}[b]{\linewidth}\raggedright
SUMMARY
\end{minipage} & \begin{minipage}[b]{\linewidth}\raggedright
SCHEDULE
\end{minipage} \\
\midrule\noalign{}
\endhead
\bottomrule\noalign{}
\endlastfoot
\end{longtable}

\subsubsection{Description}\label{description-153}

UNCODHMD activates the history match gradient unencoded output for the
history match gradient output file. Unencoded files allows external
programs to read this file type.

There is no data required for this keyword and there is no terminating
``/'' for this keyword.

This keyword is not supported by OPM Flow but would change the results
if supported so the simulation will be stopped.

\subsubsection{Example}\label{example-118}

\begin{verbatim}
--
--       ACTIVATE HISTORY MATCH GRADIENT UNENCODED OUTPUT
--
UNCODHMD
\end{verbatim}

The above example switches on the unified output file option.

\subsection{UNIFIN -- Activate The Unified Input File
Option}\label{unifin-activate-the-unified-input-file-option}

\begin{longtable}[]{@{}
  >{\raggedright\arraybackslash}p{(\linewidth - 14\tabcolsep) * \real{0.1250}}
  >{\raggedright\arraybackslash}p{(\linewidth - 14\tabcolsep) * \real{0.1250}}
  >{\raggedright\arraybackslash}p{(\linewidth - 14\tabcolsep) * \real{0.1250}}
  >{\raggedright\arraybackslash}p{(\linewidth - 14\tabcolsep) * \real{0.1250}}
  >{\raggedright\arraybackslash}p{(\linewidth - 14\tabcolsep) * \real{0.1250}}
  >{\raggedright\arraybackslash}p{(\linewidth - 14\tabcolsep) * \real{0.1250}}
  >{\raggedright\arraybackslash}p{(\linewidth - 14\tabcolsep) * \real{0.1250}}
  >{\raggedright\arraybackslash}p{(\linewidth - 14\tabcolsep) * \real{0.1250}}@{}}
\toprule\noalign{}
\begin{minipage}[b]{\linewidth}\raggedright
RUNSPEC
\end{minipage} & \begin{minipage}[b]{\linewidth}\raggedright
GRID
\end{minipage} & \begin{minipage}[b]{\linewidth}\raggedright
EDIT
\end{minipage} & \begin{minipage}[b]{\linewidth}\raggedright
PROPS
\end{minipage} & \begin{minipage}[b]{\linewidth}\raggedright
REGIONS
\end{minipage} & \begin{minipage}[b]{\linewidth}\raggedright
SOLUTION
\end{minipage} & \begin{minipage}[b]{\linewidth}\raggedright
SUMMARY
\end{minipage} & \begin{minipage}[b]{\linewidth}\raggedright
SCHEDULE
\end{minipage} \\
\midrule\noalign{}
\endhead
\bottomrule\noalign{}
\endlastfoot
\end{longtable}

\subsubsection{Description}\label{description-154}

This keyword switches on the Unified Input Files option for all input
files. Similar to the commercial simulator, OPM Flow has various options
for reading various input files and writing the resulting OPM Flow
output files as described in Table 5.54.

\begin{longtable}[]{@{}
  >{\raggedright\arraybackslash}p{(\linewidth - 6\tabcolsep) * \real{0.2500}}
  >{\raggedright\arraybackslash}p{(\linewidth - 6\tabcolsep) * \real{0.2500}}
  >{\raggedright\arraybackslash}p{(\linewidth - 6\tabcolsep) * \real{0.2500}}
  >{\raggedright\arraybackslash}p{(\linewidth - 6\tabcolsep) * \real{0.2500}}@{}}
\toprule\noalign{}
\begin{minipage}[b]{\linewidth}\raggedright
Process
\end{minipage} & \begin{minipage}[b]{\linewidth}\raggedright
Keyword
\end{minipage} & \begin{minipage}[b]{\linewidth}\raggedright
Description
\end{minipage} & \begin{minipage}[b]{\linewidth}\raggedright
Files
\end{minipage} \\
\midrule\noalign{}
\endhead
\bottomrule\noalign{}
\endlastfoot
Input & FMTIN & A character string that defines the input files to be
formatted as ASCII i.e.~text files, as oppose to binary files. The input
deck file is always of this type. The option relates to the OPM Flow
derived files that used as input, for for example when restarting from
another case. If the keyword is omitted then the default is for binary
file input. & \emph{.FEGRID }.FINSPEC \emph{.FINIT }.FRSSPEC
\emph{.FUNRST }.FSMSPEC *.FUNSMRY \\
MULTIN & A character string that defines the input files to be
non-unified multiple files, as opposed to unified files. In this case,
one file is read in per reporting time step, as opposed to all time
steps reports being read from one file. If the keyword is omitted then
the default is for one file per reporting time step. & \emph{.RSSPEC
}.X0001 \emph{.SMSPEC }.S0001 & \\
UNIFIN & A character string that defines the input files to be unified
files, as opposed to non-unified multiple files. A unified file is a
single file containing output for each reporting time step. For this
option a single summary file and a single restart file will be read. If
the keyword is omitted then the default is for one file per reporting
time step. & \emph{.RSSPEC }.UNRST \emph{.SMSPEC }.UNSMRY & \\
Output & FMTOUT & A character string that sets all output files to be
formatted as ASCII i.e.~text files, as opposed to binary files. The
\emph{.PRT, }.LOG and *.DBG files are always of this type. The option
relates to the OPM Flow output files only. In this case the files will
be portable across operating systems, but will also be very large in
terms of hard disk space. For this reason it is recommend that the
default option is used so that binary files are outputted. If the
keyword is omitted then the default is for binary file input. &
\emph{.FEGRID }.FINSPEC \emph{.FINIT }.FRSSPEC \emph{.FUNRST }.FSMSPEC
*.FUNSMRY \\
MULTOUT & A character string that defines the output files to be
non-unified multiple files, as opposed to unified files. In this case,
one file is written for each reporting time step, as opposed to all time
steps reports being written in one file. If the keyword is omitted then
the default is for one file per reporting time step. & \emph{.RSSPEC
}.X0001 \emph{.SMSPEC }.S0001 & \\
UNIFOUT & A character string that defines the output files to be unified
files, as opposed to non-unified multiple files. A unified file is a
single file containing output for each reporting time step. Here a
single summary file and a single restart file will be generated, as
opposed to one file per report time step. If the keyword is omitted then
the default is for one file per reporting time step. & \emph{.RSSPEC
}.UNRST \emph{.SMSPEC }.UNSMRY & \\
Notes: & & & \\
\end{longtable}

\emph{Table 5.54: UNIFIN Keyword Description}

There is no data required for this keyword and there is no terminating
``/'' for this keyword.

See also OPM FLOW OUTPUT FILE FORMATS for a more detailed description of
the various file types (ASCII or binary) and file structure formats
(unified or non-unified formats).

\subsubsection{Example}\label{example-119}

\begin{verbatim}
--
--       SWITCH ON THE UNIFIED INPUT FILES OPTION
--
UNIFIN
\end{verbatim}

The above example switches on the unified input file option.

\subsection{UNIFOUT -- Activate The Unified Output File
Option}\label{unifout-activate-the-unified-output-file-option}

\begin{longtable}[]{@{}
  >{\raggedright\arraybackslash}p{(\linewidth - 14\tabcolsep) * \real{0.1250}}
  >{\raggedright\arraybackslash}p{(\linewidth - 14\tabcolsep) * \real{0.1250}}
  >{\raggedright\arraybackslash}p{(\linewidth - 14\tabcolsep) * \real{0.1250}}
  >{\raggedright\arraybackslash}p{(\linewidth - 14\tabcolsep) * \real{0.1250}}
  >{\raggedright\arraybackslash}p{(\linewidth - 14\tabcolsep) * \real{0.1250}}
  >{\raggedright\arraybackslash}p{(\linewidth - 14\tabcolsep) * \real{0.1250}}
  >{\raggedright\arraybackslash}p{(\linewidth - 14\tabcolsep) * \real{0.1250}}
  >{\raggedright\arraybackslash}p{(\linewidth - 14\tabcolsep) * \real{0.1250}}@{}}
\toprule\noalign{}
\begin{minipage}[b]{\linewidth}\raggedright
RUNSPEC
\end{minipage} & \begin{minipage}[b]{\linewidth}\raggedright
GRID
\end{minipage} & \begin{minipage}[b]{\linewidth}\raggedright
EDIT
\end{minipage} & \begin{minipage}[b]{\linewidth}\raggedright
PROPS
\end{minipage} & \begin{minipage}[b]{\linewidth}\raggedright
REGIONS
\end{minipage} & \begin{minipage}[b]{\linewidth}\raggedright
SOLUTION
\end{minipage} & \begin{minipage}[b]{\linewidth}\raggedright
SUMMARY
\end{minipage} & \begin{minipage}[b]{\linewidth}\raggedright
SCHEDULE
\end{minipage} \\
\midrule\noalign{}
\endhead
\bottomrule\noalign{}
\endlastfoot
\end{longtable}

\subsubsection{Description}\label{description-155}

This keyword switches on the Unified Output Files option for all output
files. Similar to the commercial simulator, OPM Flow has various options
for reading various input files and writing the resulting OPM Flow
output files as described in Table 5.55.

\begin{longtable}[]{@{}
  >{\raggedright\arraybackslash}p{(\linewidth - 6\tabcolsep) * \real{0.2500}}
  >{\raggedright\arraybackslash}p{(\linewidth - 6\tabcolsep) * \real{0.2500}}
  >{\raggedright\arraybackslash}p{(\linewidth - 6\tabcolsep) * \real{0.2500}}
  >{\raggedright\arraybackslash}p{(\linewidth - 6\tabcolsep) * \real{0.2500}}@{}}
\toprule\noalign{}
\begin{minipage}[b]{\linewidth}\raggedright
Process
\end{minipage} & \begin{minipage}[b]{\linewidth}\raggedright
Keyword
\end{minipage} & \begin{minipage}[b]{\linewidth}\raggedright
Description
\end{minipage} & \begin{minipage}[b]{\linewidth}\raggedright
Files
\end{minipage} \\
\midrule\noalign{}
\endhead
\bottomrule\noalign{}
\endlastfoot
Input & FMTIN & A character string that defines the input files to be
formatted as ASCII i.e.~text files, as oppose to binary files. The input
deck file is always of this type. The option relates to the OPM Flow
derived files that used as input, for for example when restarting from
another case. If the keyword is omitted then the default is for binary
file input. & \emph{.FEGRID }.FINSPEC \emph{.FINIT }.FRSSPEC
\emph{.FUNRST }.FSMSPEC *.FUNSMRY \\
MULTIN & A character string that defines the input files to be
non-unified multiple files, as opposed to unified files. In this case,
one file is read in per reporting time step, as opposed to all time
steps reports being read from one file. If the keyword is omitted then
the default is for one file per reporting time step. & \emph{.RSSPEC
}.X0001 \emph{.SMSPEC }.S0001 & \\
UNIFIN & A character string that defines the input files to be unified
files, as opposed to non-unified multiple files. A unified file is a
single file containing output for each reporting time step. For this
option a single summary file and a single restart file will be read. If
the keyword is omitted then the default is for one file per reporting
time step. & \emph{.RSSPEC }.UNRST \emph{.SMSPEC }.UNSMRY & \\
Output & FMTOUT & A character string that sets all output files to be
formatted as ASCII i.e.~text files, as opposed to binary files. The
\emph{.PRT, }.LOG and *.DBG files are always of this type. The option
relates to the OPM Flow output files only. In this case the files will
be portable across operating systems, but will also be very large in
terms of hard disk space. For this reason it is recommend that the
default option is used so that binary files are outputted. If the
keyword is omitted then the default is for binary file input. &
\emph{.FEGRID }.FINSPEC \emph{.FINIT }.FRSSPEC \emph{.FUNRST }.FSMSPEC
*.FUNSMRY \\
MULTOUT & A character string that defines the output files to be
non-unified multiple files, as opposed to unified files. In this case,
one file is written for each reporting time step, as opposed to all time
steps reports being written in one file. If the keyword is omitted then
the default is for one file per reporting time step. & \emph{.RSSPEC
}.X0001 \emph{.SMSPEC }.S0001 & \\
UNIFOUT & A character string that defines the output files to be unified
files, as opposed to non-unified multiple files. A unified file is a
single file containing output for each reporting time step. Here a
single summary file and a single restart file will be generated, as
opposed to one file per report time step. If the keyword is omitted then
the default is for one file per reporting time step. & \emph{.RSSPEC
}.UNRST \emph{.SMSPEC }.UNSMRY & \\
Notes: & & & \\
\end{longtable}

\emph{Table 5.55: UNIFOUT Keyword Description}

There is no data required for this keyword and there is no terminating
``/'' for this keyword.

See also OPM FLOW OUTPUT FILE FORMATS for a more detailed description of
the various file types (ASCII or binary) and file structure formats
(unified or non-unified formats).

\subsubsection{Example}\label{example-120}

\begin{verbatim}
--
--       SWITCH ON THE UNIFIED OUTPUT FILES OPTION
--
UNIFOUT
\end{verbatim}

The above example switches on the unified output file option.

\subsection{UNIFOUTS -- Activate The Unified Output Summary File
Option}\label{unifouts-activate-the-unified-output-summary-file-option}

\begin{longtable}[]{@{}
  >{\raggedright\arraybackslash}p{(\linewidth - 14\tabcolsep) * \real{0.1250}}
  >{\raggedright\arraybackslash}p{(\linewidth - 14\tabcolsep) * \real{0.1250}}
  >{\raggedright\arraybackslash}p{(\linewidth - 14\tabcolsep) * \real{0.1250}}
  >{\raggedright\arraybackslash}p{(\linewidth - 14\tabcolsep) * \real{0.1250}}
  >{\raggedright\arraybackslash}p{(\linewidth - 14\tabcolsep) * \real{0.1250}}
  >{\raggedright\arraybackslash}p{(\linewidth - 14\tabcolsep) * \real{0.1250}}
  >{\raggedright\arraybackslash}p{(\linewidth - 14\tabcolsep) * \real{0.1250}}
  >{\raggedright\arraybackslash}p{(\linewidth - 14\tabcolsep) * \real{0.1250}}@{}}
\toprule\noalign{}
\begin{minipage}[b]{\linewidth}\raggedright
RUNSPEC
\end{minipage} & \begin{minipage}[b]{\linewidth}\raggedright
GRID
\end{minipage} & \begin{minipage}[b]{\linewidth}\raggedright
EDIT
\end{minipage} & \begin{minipage}[b]{\linewidth}\raggedright
PROPS
\end{minipage} & \begin{minipage}[b]{\linewidth}\raggedright
REGIONS
\end{minipage} & \begin{minipage}[b]{\linewidth}\raggedright
SOLUTION
\end{minipage} & \begin{minipage}[b]{\linewidth}\raggedright
SUMMARY
\end{minipage} & \begin{minipage}[b]{\linewidth}\raggedright
SCHEDULE
\end{minipage} \\
\midrule\noalign{}
\endhead
\bottomrule\noalign{}
\endlastfoot
\end{longtable}

\subsubsection{Description}\label{description-156}

The UNIFOUTS keyword causes the SUMMARY file output files to be a
unified file, as opposed to non-unified multiple files. A unified file
is a single file containing output for each reporting time step. Here a
single SUMMARY file will be generated, as opposed to one file per report
time step. See also the MULTOUT keyword in the RUNSPEC section that sets
both the SUMMARY and RESTART files to be non-unified multiple files, as
opposed to unified files. Note also that UNIFOUTS keyword has precedence
over the MULTOUT keyword for SUMMARY files.

There is no data required for this keyword and there is no terminating
``/'' for this keyword.

This keyword is not supported by OPM Flow but has no effect on the
results so it will be ignored.

\subsubsection{Example}\label{example-121}

\begin{verbatim}
--
--       ACTIVATE THE UNIFIED OUTPUT SUMMARY FILE OPTION
--
UNIFOUTS
\end{verbatim}

The above example switches on writing of unified SUMMARY output files.

\subsection{UNIFSAVE -- Activate The Unified Output Save File
Option}\label{unifsave-activate-the-unified-output-save-file-option}

\begin{longtable}[]{@{}
  >{\raggedright\arraybackslash}p{(\linewidth - 14\tabcolsep) * \real{0.1250}}
  >{\raggedright\arraybackslash}p{(\linewidth - 14\tabcolsep) * \real{0.1250}}
  >{\raggedright\arraybackslash}p{(\linewidth - 14\tabcolsep) * \real{0.1250}}
  >{\raggedright\arraybackslash}p{(\linewidth - 14\tabcolsep) * \real{0.1250}}
  >{\raggedright\arraybackslash}p{(\linewidth - 14\tabcolsep) * \real{0.1250}}
  >{\raggedright\arraybackslash}p{(\linewidth - 14\tabcolsep) * \real{0.1250}}
  >{\raggedright\arraybackslash}p{(\linewidth - 14\tabcolsep) * \real{0.1250}}
  >{\raggedright\arraybackslash}p{(\linewidth - 14\tabcolsep) * \real{0.1250}}@{}}
\toprule\noalign{}
\begin{minipage}[b]{\linewidth}\raggedright
RUNSPEC
\end{minipage} & \begin{minipage}[b]{\linewidth}\raggedright
GRID
\end{minipage} & \begin{minipage}[b]{\linewidth}\raggedright
EDIT
\end{minipage} & \begin{minipage}[b]{\linewidth}\raggedright
PROPS
\end{minipage} & \begin{minipage}[b]{\linewidth}\raggedright
REGIONS
\end{minipage} & \begin{minipage}[b]{\linewidth}\raggedright
SOLUTION
\end{minipage} & \begin{minipage}[b]{\linewidth}\raggedright
SUMMARY
\end{minipage} & \begin{minipage}[b]{\linewidth}\raggedright
SCHEDULE
\end{minipage} \\
\midrule\noalign{}
\endhead
\bottomrule\noalign{}
\endlastfoot
\end{longtable}

\subsubsection{Description}\label{description-157}

The UNIFSAVE keyword causes the SAVE file output file to be a unified
file, as opposed to non-unified multiple files. A unified file is a
single file containing output for each reporting time step. Here a
single SAVE file will be generated, as opposed to one file per report
time step. See also the MULTOUT keyword in the RUNSPEC section that sets
both the SUMMARY and RESTART files to be non-unified multiple files, as
opposed to unified files.

There is no data required for this keyword and there is no terminating
``/'' for this keyword.

This keyword is not supported by OPM Flow but has no effect on the
results so it will be ignored.

\subsubsection{Example}\label{example-122}

\begin{verbatim}
--
--      ACTIVATE THE UNIFIED OUTPUT SAVE FILE OPTION
--
UNIFSAVE
\end{verbatim}

The above example switches on writing of unified SAVE output files.

\subsection{VAPOIL -- Activate the Vaporized Oil in Wet Gas Phase in the
Model}\label{vapoil-activate-the-vaporized-oil-in-wet-gas-phase-in-the-model}

\begin{longtable}[]{@{}
  >{\raggedright\arraybackslash}p{(\linewidth - 14\tabcolsep) * \real{0.1250}}
  >{\raggedright\arraybackslash}p{(\linewidth - 14\tabcolsep) * \real{0.1250}}
  >{\raggedright\arraybackslash}p{(\linewidth - 14\tabcolsep) * \real{0.1250}}
  >{\raggedright\arraybackslash}p{(\linewidth - 14\tabcolsep) * \real{0.1250}}
  >{\raggedright\arraybackslash}p{(\linewidth - 14\tabcolsep) * \real{0.1250}}
  >{\raggedright\arraybackslash}p{(\linewidth - 14\tabcolsep) * \real{0.1250}}
  >{\raggedright\arraybackslash}p{(\linewidth - 14\tabcolsep) * \real{0.1250}}
  >{\raggedright\arraybackslash}p{(\linewidth - 14\tabcolsep) * \real{0.1250}}@{}}
\toprule\noalign{}
\begin{minipage}[b]{\linewidth}\raggedright
RUNSPEC
\end{minipage} & \begin{minipage}[b]{\linewidth}\raggedright
GRID
\end{minipage} & \begin{minipage}[b]{\linewidth}\raggedright
EDIT
\end{minipage} & \begin{minipage}[b]{\linewidth}\raggedright
PROPS
\end{minipage} & \begin{minipage}[b]{\linewidth}\raggedright
REGIONS
\end{minipage} & \begin{minipage}[b]{\linewidth}\raggedright
SOLUTION
\end{minipage} & \begin{minipage}[b]{\linewidth}\raggedright
SUMMARY
\end{minipage} & \begin{minipage}[b]{\linewidth}\raggedright
SCHEDULE
\end{minipage} \\
\midrule\noalign{}
\endhead
\bottomrule\noalign{}
\endlastfoot
\end{longtable}

\subsubsection{Description}\label{description-158}

This keyword indicates that vaporized oil (more commonly referred to as
condensate) is present in wet {[}Natural gas that contains significant
heavy hydrocarbons such as propane, butane and other liquid hydrocarbons
is known as wet gas or rich gas. The general rule of thumb is if the gas
contains less methane (typically less than 85\% methane) and more
ethane, and other more complex hydrocarbons, it is labeled as wet gas.
Wet gas normally has GOR's less than 100,000 scf/stb or 18,000 Sm3/m3,
with the condensate having a gravity greater than 50 oAPI.{]} gas in the
model and the keyword should only be used if the there is both oil and
gas phases in the model. The keyword may be used for gas-water and
oil-water-gas input decks that contain the oil and gas phases. The
keyword will also invoke data input file checking to ensure that all the
required oil and gas phase input parameters are defined in the input
deck.

If the gas has a constant and uniform vaporized oil concentration,
Condensate-Gas Ratio (``CGR''), and if the reservoir pressure never
drops below the saturation pressure (dew point pressure), then the model
can be run more efficiently by omitting the OIL and VAPOIL keywords from
the RUNSPEC section, treating the gas as a dry gas {[}Natural gas that
occurs in the absence of condensate or liquid hydrocarbons, or gas that
had condensable hydrocarbons removed, is called dry gas. It is primarily
methane with some intermediates. The hydrocarbon mixture is solely gas
in the reservoir and there is no liquid (condensate surface liquid)
formed either in the reservoir or at surface. The term dry indicates
that the gas does not contain heavier hydrocarbons to form liquids at
the surface conditions. Dry gas typically has GOR's greater than 100,000
scf/stb or 18,000 Sm3/m3.{]}, and defining a constant Rv (CGR) value
with keyword RVCONST or RVCONSTT in the PROPS section. This results in
the model being run with as a dry gas problem with no active oil
(condensate) phase. However, OPM Flow takes into account the constant Rv
in the calculations and reporting.

There is no data required for this keyword and there is no terminating
``/'' for this keyword.

\subsubsection{Example}\label{example-123}

\begin{verbatim}
--
--       VAPORIZED OIL IN WET GAS IS PRESENT IN THE RUN
--
VAPOIL
\end{verbatim}

The above example declares that the vaporized oil, i.e.~condensate, in
the gas phase is active in the model.

\subsection{VAPWAT -- Activate Vaporized Water in the Dry and Wet Gas
Phases}\label{vapwat-activate-vaporized-water-in-the-dry-and-wet-gas-phases}

\begin{longtable}[]{@{}
  >{\raggedright\arraybackslash}p{(\linewidth - 14\tabcolsep) * \real{0.1250}}
  >{\raggedright\arraybackslash}p{(\linewidth - 14\tabcolsep) * \real{0.1250}}
  >{\raggedright\arraybackslash}p{(\linewidth - 14\tabcolsep) * \real{0.1250}}
  >{\raggedright\arraybackslash}p{(\linewidth - 14\tabcolsep) * \real{0.1250}}
  >{\raggedright\arraybackslash}p{(\linewidth - 14\tabcolsep) * \real{0.1250}}
  >{\raggedright\arraybackslash}p{(\linewidth - 14\tabcolsep) * \real{0.1250}}
  >{\raggedright\arraybackslash}p{(\linewidth - 14\tabcolsep) * \real{0.1250}}
  >{\raggedright\arraybackslash}p{(\linewidth - 14\tabcolsep) * \real{0.1250}}@{}}
\toprule\noalign{}
\begin{minipage}[b]{\linewidth}\raggedright
RUNSPEC
\end{minipage} & \begin{minipage}[b]{\linewidth}\raggedright
GRID
\end{minipage} & \begin{minipage}[b]{\linewidth}\raggedright
EDIT
\end{minipage} & \begin{minipage}[b]{\linewidth}\raggedright
PROPS
\end{minipage} & \begin{minipage}[b]{\linewidth}\raggedright
REGIONS
\end{minipage} & \begin{minipage}[b]{\linewidth}\raggedright
SOLUTION
\end{minipage} & \begin{minipage}[b]{\linewidth}\raggedright
SUMMARY
\end{minipage} & \begin{minipage}[b]{\linewidth}\raggedright
SCHEDULE
\end{minipage} \\
\midrule\noalign{}
\endhead
\bottomrule\noalign{}
\endlastfoot
\end{longtable}

\subsubsection{Description}\label{description-159}

This keyword indicates that vaporized water is present in the gas phase
and the keyword should only be used if both water and gas phases are
present in the model. VAPWAT should also be used in conjunction with the
PRECSALT keyword in the RUNSPEC section in order to activate OPM Flow's
Salt Precipitation model. VAPWAT may be used for gas-water and
oil-water-gas input decks that contain the oil, gas and water phases.

\begin{tcolorbox}[enhanced jigsaw, title=\textcolor{quarto-callout-note-color}{\faInfo}\hspace{0.5em}{Note}, coltitle=black, breakable, toprule=.15mm, colbacktitle=quarto-callout-note-color!10!white, leftrule=.75mm, opacitybacktitle=0.6, colback=white, bottomtitle=1mm, toptitle=1mm, titlerule=0mm, left=2mm, colframe=quarto-callout-note-color-frame, arc=.35mm, rightrule=.15mm, bottomrule=.15mm, opacityback=0]

This is an OPM Flow specific keyword for the simulator's Water
Vaporization Model that is activated by declaring that vaporized water
is present in the run.

\end{tcolorbox}

Note that if the VAPWAT keyword is in the input deck then either the
PVTGW or PVTGWO keywords in PROPS section should be used to defined the
gas and water PVT properties.

Secondly, if both the VAPWAT keyword and the PRECSALT keyword (used to
activate the OPM Flow's Salt Precipitation model) are present in the
input deck, then the RWGSALT keyword in the PROPS section also needs to
be present.

There is no data required for this keyword and there is no terminating
``/'' for this keyword.

\subsubsection{Example}\label{example-124}

\begin{verbatim}
--
--       VAPORIZED WATER IN DRY/WET GAS IS PRESENT IN THE RUN (OPM FLOW KEYWORD)
--
VAPWAT
\end{verbatim}

The above example declares that the vaporized water is present in the
gas phase and is active in the model.

\subsection{VE -- Activate Vertical Equilibrium Model
(Global)}\label{ve-activate-vertical-equilibrium-model-global}

\begin{longtable}[]{@{}
  >{\raggedright\arraybackslash}p{(\linewidth - 14\tabcolsep) * \real{0.1250}}
  >{\raggedright\arraybackslash}p{(\linewidth - 14\tabcolsep) * \real{0.1250}}
  >{\raggedright\arraybackslash}p{(\linewidth - 14\tabcolsep) * \real{0.1250}}
  >{\raggedright\arraybackslash}p{(\linewidth - 14\tabcolsep) * \real{0.1250}}
  >{\raggedright\arraybackslash}p{(\linewidth - 14\tabcolsep) * \real{0.1250}}
  >{\raggedright\arraybackslash}p{(\linewidth - 14\tabcolsep) * \real{0.1250}}
  >{\raggedright\arraybackslash}p{(\linewidth - 14\tabcolsep) * \real{0.1250}}
  >{\raggedright\arraybackslash}p{(\linewidth - 14\tabcolsep) * \real{0.1250}}@{}}
\toprule\noalign{}
\begin{minipage}[b]{\linewidth}\raggedright
RUNSPEC
\end{minipage} & \begin{minipage}[b]{\linewidth}\raggedright
GRID
\end{minipage} & \begin{minipage}[b]{\linewidth}\raggedright
EDIT
\end{minipage} & \begin{minipage}[b]{\linewidth}\raggedright
PROPS
\end{minipage} & \begin{minipage}[b]{\linewidth}\raggedright
REGIONS
\end{minipage} & \begin{minipage}[b]{\linewidth}\raggedright
SOLUTION
\end{minipage} & \begin{minipage}[b]{\linewidth}\raggedright
SUMMARY
\end{minipage} & \begin{minipage}[b]{\linewidth}\raggedright
SCHEDULE
\end{minipage} \\
\midrule\noalign{}
\endhead
\bottomrule\noalign{}
\endlastfoot
\end{longtable}

\subsubsection{Description}\label{description-160}

This keyword actives the Vertical Equilibrium (``VE'') model for the
global grid and optionally specifies the type of VE model. The VE model
type can either be compressed for merging columns of grid blocks into a
single grid block, or uncompressed for the standard VE model.

This keyword is not supported by OPM Flow but would change the results
if supported so the simulation will be stopped.

\subsection{VFPIDIMS -- Injection Vertical Flow Performance Table
Dimensions}\label{vfpidims-injection-vertical-flow-performance-table-dimensions}

\begin{longtable}[]{@{}
  >{\raggedright\arraybackslash}p{(\linewidth - 14\tabcolsep) * \real{0.1250}}
  >{\raggedright\arraybackslash}p{(\linewidth - 14\tabcolsep) * \real{0.1250}}
  >{\raggedright\arraybackslash}p{(\linewidth - 14\tabcolsep) * \real{0.1250}}
  >{\raggedright\arraybackslash}p{(\linewidth - 14\tabcolsep) * \real{0.1250}}
  >{\raggedright\arraybackslash}p{(\linewidth - 14\tabcolsep) * \real{0.1250}}
  >{\raggedright\arraybackslash}p{(\linewidth - 14\tabcolsep) * \real{0.1250}}
  >{\raggedright\arraybackslash}p{(\linewidth - 14\tabcolsep) * \real{0.1250}}
  >{\raggedright\arraybackslash}p{(\linewidth - 14\tabcolsep) * \real{0.1250}}@{}}
\toprule\noalign{}
\begin{minipage}[b]{\linewidth}\raggedright
RUNSPEC
\end{minipage} & \begin{minipage}[b]{\linewidth}\raggedright
GRID
\end{minipage} & \begin{minipage}[b]{\linewidth}\raggedright
EDIT
\end{minipage} & \begin{minipage}[b]{\linewidth}\raggedright
PROPS
\end{minipage} & \begin{minipage}[b]{\linewidth}\raggedright
REGIONS
\end{minipage} & \begin{minipage}[b]{\linewidth}\raggedright
SOLUTION
\end{minipage} & \begin{minipage}[b]{\linewidth}\raggedright
SUMMARY
\end{minipage} & \begin{minipage}[b]{\linewidth}\raggedright
SCHEDULE
\end{minipage} \\
\midrule\noalign{}
\endhead
\bottomrule\noalign{}
\endlastfoot
\end{longtable}

\subsubsection{Description}\label{description-161}

VFPIDIMS keyword defines the maximum dimensions of the injection well
Vertical Lift Performance (``VFP'') tables defined by VFPINJ keyword.
The VFP tables for the producing wells are defined by the VFPPDIMS
keyword.

\begin{longtable}[]{@{}
  >{\raggedright\arraybackslash}p{(\linewidth - 6\tabcolsep) * \real{0.2500}}
  >{\raggedright\arraybackslash}p{(\linewidth - 6\tabcolsep) * \real{0.2500}}
  >{\raggedright\arraybackslash}p{(\linewidth - 6\tabcolsep) * \real{0.2500}}
  >{\raggedright\arraybackslash}p{(\linewidth - 6\tabcolsep) * \real{0.2500}}@{}}
\toprule\noalign{}
\begin{minipage}[b]{\linewidth}\raggedright
No.
\end{minipage} & \begin{minipage}[b]{\linewidth}\raggedright
Name
\end{minipage} & \begin{minipage}[b]{\linewidth}\raggedright
Description
\end{minipage} & \begin{minipage}[b]{\linewidth}\raggedright
Default
\end{minipage} \\
\midrule\noalign{}
\endhead
\bottomrule\noalign{}
\endlastfoot
1 & MXMFLO & A positive integer that defines the maximum number of
injection rate entries for the VFPINJ keyword. & 0 \\
2 & MXMTHP & A positive integer that defines the maximum number of THP
entries for the VFPINJ keyword. & 0 \\
3 & MXVFPTAB & A positive integer that defines the maximum number of
VFPINJ tables entered through the VFPINJ keyword. & 0 \\
Notes: & & & \\
\end{longtable}

\emph{Table 5.56: VFPIDIMS Keyword Description}

\subsubsection{Example}\label{example-125}

\begin{verbatim}
--       INJECTING VFP TABLES
--       VFP     VFP     VFP
--       MXMFLO  MXMTHP  NMMVFT
VFPIDIMS
         10      10      12                                                    /
\end{verbatim}

The above example defines that the maximum number of injection rates and
THP entries on the VFPINJ keyword is 10, and the number of VFPINJ tables
is 12.

\subsection{VFPPDIMS -- Production Vertical Flow Performance Table
Dimensions}\label{vfppdims-production-vertical-flow-performance-table-dimensions}

\begin{longtable}[]{@{}
  >{\raggedright\arraybackslash}p{(\linewidth - 14\tabcolsep) * \real{0.1250}}
  >{\raggedright\arraybackslash}p{(\linewidth - 14\tabcolsep) * \real{0.1250}}
  >{\raggedright\arraybackslash}p{(\linewidth - 14\tabcolsep) * \real{0.1250}}
  >{\raggedright\arraybackslash}p{(\linewidth - 14\tabcolsep) * \real{0.1250}}
  >{\raggedright\arraybackslash}p{(\linewidth - 14\tabcolsep) * \real{0.1250}}
  >{\raggedright\arraybackslash}p{(\linewidth - 14\tabcolsep) * \real{0.1250}}
  >{\raggedright\arraybackslash}p{(\linewidth - 14\tabcolsep) * \real{0.1250}}
  >{\raggedright\arraybackslash}p{(\linewidth - 14\tabcolsep) * \real{0.1250}}@{}}
\toprule\noalign{}
\begin{minipage}[b]{\linewidth}\raggedright
RUNSPEC
\end{minipage} & \begin{minipage}[b]{\linewidth}\raggedright
GRID
\end{minipage} & \begin{minipage}[b]{\linewidth}\raggedright
EDIT
\end{minipage} & \begin{minipage}[b]{\linewidth}\raggedright
PROPS
\end{minipage} & \begin{minipage}[b]{\linewidth}\raggedright
REGIONS
\end{minipage} & \begin{minipage}[b]{\linewidth}\raggedright
SOLUTION
\end{minipage} & \begin{minipage}[b]{\linewidth}\raggedright
SUMMARY
\end{minipage} & \begin{minipage}[b]{\linewidth}\raggedright
SCHEDULE
\end{minipage} \\
\midrule\noalign{}
\endhead
\bottomrule\noalign{}
\endlastfoot
\end{longtable}

\subsubsection{Description}\label{description-162}

VFPPDIMS keyword defines the maximum dimensions of the production well
Vertical Lift Performance (``VFP'') tables defined by VFPPROD keyword.
The VFP tables for the injection wells are defined by the VFPIDIMS
keyword.

\begin{longtable}[]{@{}
  >{\raggedright\arraybackslash}p{(\linewidth - 6\tabcolsep) * \real{0.2500}}
  >{\raggedright\arraybackslash}p{(\linewidth - 6\tabcolsep) * \real{0.2500}}
  >{\raggedright\arraybackslash}p{(\linewidth - 6\tabcolsep) * \real{0.2500}}
  >{\raggedright\arraybackslash}p{(\linewidth - 6\tabcolsep) * \real{0.2500}}@{}}
\toprule\noalign{}
\begin{minipage}[b]{\linewidth}\raggedright
No.
\end{minipage} & \begin{minipage}[b]{\linewidth}\raggedright
Name
\end{minipage} & \begin{minipage}[b]{\linewidth}\raggedright
Description
\end{minipage} & \begin{minipage}[b]{\linewidth}\raggedright
Default
\end{minipage} \\
\midrule\noalign{}
\endhead
\bottomrule\noalign{}
\endlastfoot
1 & MXMFLO & A positive integer that defines the maximum number of
production flow rate entries for the VFPPROD keyword. & 0 \\
2 & MXMTHP & A positive integer that defines the maximum number of THP
entries for the VFPPROD keyword. & 0 \\
3 & MXMWFR & A positive integer that defines the maximum number of water
fraction entries (WOR, WCUT, GWR etc.) for the VFPPROD keyword. & 0 \\
4 & MXMGFR & A positive integer that defines the maximum number of gas
fraction entries (GOR, GLR, OGR etc.) for the VFPPROD keyword. & 0 \\
5 & MXMALQ & A positive integer that defines the maximum number of
artificial lift quantity entries for the VFPPROD keyword. & 0 \\
6 & MXVFPTAB & A positive integer that defines the maximum number of
VFPPROD tables entered through the VFPPROD keyword. & 0 \\
Notes: & & & \\
\end{longtable}

\emph{Table 5.57: VFPPDIMS Keyword Description}

\subsubsection{Example}\label{example-126}

\begin{verbatim}
--       PRODUCING VFP TABLES
--       VFP     VFP     VFP     VFP     VFP     VFP
--       MXMFLO  MXMTHP  MXMWFR  MXMGFR  MXMALQ  NMMVFT
VFPPDIMS
         20      10      10      10      6       9                             /
\end{verbatim}

Here the example shows that there are a maximum of 20 flow rates, 10 THP
entries, 10 water and gas fraction entries, and six artificial lift
entries for the nine VFPPROD VFP production tables.

\subsection{VISAGE -- Activate External Reservoir Geo-Mechanics VISAGE
Option}\label{visage-activate-external-reservoir-geo-mechanics-visage-option}

\begin{longtable}[]{@{}
  >{\raggedright\arraybackslash}p{(\linewidth - 14\tabcolsep) * \real{0.1250}}
  >{\raggedright\arraybackslash}p{(\linewidth - 14\tabcolsep) * \real{0.1250}}
  >{\raggedright\arraybackslash}p{(\linewidth - 14\tabcolsep) * \real{0.1250}}
  >{\raggedright\arraybackslash}p{(\linewidth - 14\tabcolsep) * \real{0.1250}}
  >{\raggedright\arraybackslash}p{(\linewidth - 14\tabcolsep) * \real{0.1250}}
  >{\raggedright\arraybackslash}p{(\linewidth - 14\tabcolsep) * \real{0.1250}}
  >{\raggedright\arraybackslash}p{(\linewidth - 14\tabcolsep) * \real{0.1250}}
  >{\raggedright\arraybackslash}p{(\linewidth - 14\tabcolsep) * \real{0.1250}}@{}}
\toprule\noalign{}
\begin{minipage}[b]{\linewidth}\raggedright
RUNSPEC
\end{minipage} & \begin{minipage}[b]{\linewidth}\raggedright
GRID
\end{minipage} & \begin{minipage}[b]{\linewidth}\raggedright
EDIT
\end{minipage} & \begin{minipage}[b]{\linewidth}\raggedright
PROPS
\end{minipage} & \begin{minipage}[b]{\linewidth}\raggedright
REGIONS
\end{minipage} & \begin{minipage}[b]{\linewidth}\raggedright
SOLUTION
\end{minipage} & \begin{minipage}[b]{\linewidth}\raggedright
SUMMARY
\end{minipage} & \begin{minipage}[b]{\linewidth}\raggedright
SCHEDULE
\end{minipage} \\
\midrule\noalign{}
\endhead
\bottomrule\noalign{}
\endlastfoot
\end{longtable}

\subsubsection{Description}\label{description-163}

The VISAGE keyword activates the External Reservoir Geo-Mechanics VISAGE
option. The keyword should not be used in input decks as the associated
data is generated by an external program.

This keyword is not supported by OPM Flow but would change the results
if supported so the simulation will be stopped.

\subsection{VISCD -- Activate Dual Porosity Viscous Displacement
Option}\label{viscd-activate-dual-porosity-viscous-displacement-option}

\begin{longtable}[]{@{}
  >{\raggedright\arraybackslash}p{(\linewidth - 14\tabcolsep) * \real{0.1250}}
  >{\raggedright\arraybackslash}p{(\linewidth - 14\tabcolsep) * \real{0.1250}}
  >{\raggedright\arraybackslash}p{(\linewidth - 14\tabcolsep) * \real{0.1250}}
  >{\raggedright\arraybackslash}p{(\linewidth - 14\tabcolsep) * \real{0.1250}}
  >{\raggedright\arraybackslash}p{(\linewidth - 14\tabcolsep) * \real{0.1250}}
  >{\raggedright\arraybackslash}p{(\linewidth - 14\tabcolsep) * \real{0.1250}}
  >{\raggedright\arraybackslash}p{(\linewidth - 14\tabcolsep) * \real{0.1250}}
  >{\raggedright\arraybackslash}p{(\linewidth - 14\tabcolsep) * \real{0.1250}}@{}}
\toprule\noalign{}
\begin{minipage}[b]{\linewidth}\raggedright
RUNSPEC
\end{minipage} & \begin{minipage}[b]{\linewidth}\raggedright
GRID
\end{minipage} & \begin{minipage}[b]{\linewidth}\raggedright
EDIT
\end{minipage} & \begin{minipage}[b]{\linewidth}\raggedright
PROPS
\end{minipage} & \begin{minipage}[b]{\linewidth}\raggedright
REGIONS
\end{minipage} & \begin{minipage}[b]{\linewidth}\raggedright
SOLUTION
\end{minipage} & \begin{minipage}[b]{\linewidth}\raggedright
SUMMARY
\end{minipage} & \begin{minipage}[b]{\linewidth}\raggedright
SCHEDULE
\end{minipage} \\
\midrule\noalign{}
\endhead
\bottomrule\noalign{}
\endlastfoot
\end{longtable}

\subsubsection{Description}\label{description-164}

The VISCD keyword activates the Dual Porosity Viscous Displacement
option for dual porosity and dual permeability models, and therefore
requires either the DUALPORO or DUALPERM keyword to be entered in the
RUNSPEC section to activate either one of these options. The VISCD
option is used to model the viscous displacement of fluids from the
matrix by the fracture pressure gradient, for when the fracture system
has a more moderate permeability, and flow to and from the matrix caused
by the fracture pressure gradient acts as an additional production
mechanism {[}Gilman, J. R. and Kazemi, H. ``Improved Calculation for
Viscous and Gravity Displacement in Matrix Blocks in Dual-Porosity
Simulators,'' paper SPE 16010 (includes a number of associated papers),
Journal of Petroleum Technology (1988) 40, No.~1, 60-70.{]}. Normally
this mechanism is ignored as the pressure gradient in the fracture
system is small due to the very high permeability of the fracture
system. See the LX, LY and LZ keywords in the GRID section that define
representative matrix grid block sizes.

There is no data required for this keyword and there is no terminating
``/'' for this keyword.

This keyword is not supported by OPM Flow but would change the results
if supported so the simulation will be stopped.

\subsubsection{Example}\label{example-127}

\begin{verbatim}
--
--       ACTIVATE DUAL POROSITY VISCOUS DISPLACEMENT OPTION
--
VISCD
\end{verbatim}

The above example activates the dual porosity viscous displacement
option.

\subsection{WARN -- Activate Warning
Messages}\label{warn-activate-warning-messages-1}

Turns on warning messages to be printed to the print file (*.PRT); note
that this keyword is activated by default and can subsequently be
switched off by the NOWARN activation keyword. The warning messages may
be turned on and off using keywords WARN and NOWARN. OPM Flow always
prints error messages.

See WARN -- Activate Warning Messages in the GLOBAL section for a full
description.

\subsection{WATER -- Activate the Water Phase in the
Model}\label{water-activate-the-water-phase-in-the-model}

\begin{longtable}[]{@{}
  >{\raggedright\arraybackslash}p{(\linewidth - 14\tabcolsep) * \real{0.1250}}
  >{\raggedright\arraybackslash}p{(\linewidth - 14\tabcolsep) * \real{0.1250}}
  >{\raggedright\arraybackslash}p{(\linewidth - 14\tabcolsep) * \real{0.1250}}
  >{\raggedright\arraybackslash}p{(\linewidth - 14\tabcolsep) * \real{0.1250}}
  >{\raggedright\arraybackslash}p{(\linewidth - 14\tabcolsep) * \real{0.1250}}
  >{\raggedright\arraybackslash}p{(\linewidth - 14\tabcolsep) * \real{0.1250}}
  >{\raggedright\arraybackslash}p{(\linewidth - 14\tabcolsep) * \real{0.1250}}
  >{\raggedright\arraybackslash}p{(\linewidth - 14\tabcolsep) * \real{0.1250}}@{}}
\toprule\noalign{}
\begin{minipage}[b]{\linewidth}\raggedright
RUNSPEC
\end{minipage} & \begin{minipage}[b]{\linewidth}\raggedright
GRID
\end{minipage} & \begin{minipage}[b]{\linewidth}\raggedright
EDIT
\end{minipage} & \begin{minipage}[b]{\linewidth}\raggedright
PROPS
\end{minipage} & \begin{minipage}[b]{\linewidth}\raggedright
REGIONS
\end{minipage} & \begin{minipage}[b]{\linewidth}\raggedright
SOLUTION
\end{minipage} & \begin{minipage}[b]{\linewidth}\raggedright
SUMMARY
\end{minipage} & \begin{minipage}[b]{\linewidth}\raggedright
SCHEDULE
\end{minipage} \\
\midrule\noalign{}
\endhead
\bottomrule\noalign{}
\endlastfoot
\end{longtable}

\subsubsection{Description}\label{description-165}

This keyword indicate that the water phase is present in the model and
must be used for gas-water, oil-water, oil-water-gas input decks that
contain the water phase. The keyword will also invoke data input file
checking to ensure that all the required water phase input parameters
are defined in the input deck.

There is no data required for this keyword and there is no terminating
``/'' for this keyword.

\subsubsection{Example}\label{example-128}

\begin{verbatim}
--
--       WATER PHASE IS PRESENT IN THE RUN
--
WATER
\end{verbatim}

The above example declares that the water phase is active in the model.

\subsection{WELLDIMS -- Define the Wells and Group
Dimensions}\label{welldims-define-the-wells-and-group-dimensions}

\begin{longtable}[]{@{}
  >{\raggedright\arraybackslash}p{(\linewidth - 14\tabcolsep) * \real{0.1250}}
  >{\raggedright\arraybackslash}p{(\linewidth - 14\tabcolsep) * \real{0.1250}}
  >{\raggedright\arraybackslash}p{(\linewidth - 14\tabcolsep) * \real{0.1250}}
  >{\raggedright\arraybackslash}p{(\linewidth - 14\tabcolsep) * \real{0.1250}}
  >{\raggedright\arraybackslash}p{(\linewidth - 14\tabcolsep) * \real{0.1250}}
  >{\raggedright\arraybackslash}p{(\linewidth - 14\tabcolsep) * \real{0.1250}}
  >{\raggedright\arraybackslash}p{(\linewidth - 14\tabcolsep) * \real{0.1250}}
  >{\raggedright\arraybackslash}p{(\linewidth - 14\tabcolsep) * \real{0.1250}}@{}}
\toprule\noalign{}
\begin{minipage}[b]{\linewidth}\raggedright
RUNSPEC
\end{minipage} & \begin{minipage}[b]{\linewidth}\raggedright
GRID
\end{minipage} & \begin{minipage}[b]{\linewidth}\raggedright
EDIT
\end{minipage} & \begin{minipage}[b]{\linewidth}\raggedright
PROPS
\end{minipage} & \begin{minipage}[b]{\linewidth}\raggedright
REGIONS
\end{minipage} & \begin{minipage}[b]{\linewidth}\raggedright
SOLUTION
\end{minipage} & \begin{minipage}[b]{\linewidth}\raggedright
SUMMARY
\end{minipage} & \begin{minipage}[b]{\linewidth}\raggedright
SCHEDULE
\end{minipage} \\
\midrule\noalign{}
\endhead
\bottomrule\noalign{}
\endlastfoot
\end{longtable}

\subsubsection{Description}\label{description-166}

WELLDIMS defines various well and group dimensions for the run. The
commercial simulator combines both the black-oil and compositional
simulator variables on this keyword; however, although all the
parameters are explained below only the black-oil parameters are used by
OPM Flow.

\begin{longtable}[]{@{}
  >{\raggedright\arraybackslash}p{(\linewidth - 6\tabcolsep) * \real{0.2500}}
  >{\raggedright\arraybackslash}p{(\linewidth - 6\tabcolsep) * \real{0.2500}}
  >{\raggedright\arraybackslash}p{(\linewidth - 6\tabcolsep) * \real{0.2500}}
  >{\raggedright\arraybackslash}p{(\linewidth - 6\tabcolsep) * \real{0.2500}}@{}}
\toprule\noalign{}
\begin{minipage}[b]{\linewidth}\raggedright
No.
\end{minipage} & \begin{minipage}[b]{\linewidth}\raggedright
Name
\end{minipage} & \begin{minipage}[b]{\linewidth}\raggedright
Description
\end{minipage} & \begin{minipage}[b]{\linewidth}\raggedright
Default
\end{minipage} \\
\midrule\noalign{}
\endhead
\bottomrule\noalign{}
\endlastfoot
1 & MXWELS & A positive integer defining the maximum number of wells for
this model. & 0 \\
2 & MXCONS & A positive integer defining the maximum number of grid
block connections per well for this model. & 0 \\
3 & MXGRPS & A positive integer defining the maximum number of groups
for this model. & 0 \\
4 & MXGRPW & A positive integer defining the maximum number of wells
that can belong to a group in the model and the maximum number of child
groups in a group. Note that MXGRPW sets both the maximum number of
wells in a group and the maximum number of child groups in a group. The
former applies to groups that contain wells and the latter applies to
groups that contain other groups. See also the GRUPTREE keyword in the
SCHEDULE section to define group hierarchy. & 0 \\
5 & MXSTAGE & A positive integer defining the maximum number of stages
per separator for this model. This option is ignored by OPM Flow. & 5 \\
6 & MXSTRMS & A positive integer defining the maximum number of well
streams for this model. This option is ignored by OPM Flow. & 10 \\
7 & MXMIXS & A positive integer defining the maximum number of mixtures
for this model. This option is ignored by OPM Flow. & 5 \\
8 & MXSEPS & A positive integer defining the maximum number of
separators for this model. This option is ignored by OPM Flow. & 4 \\
9 & MXCOMPS & A positive integer defining the maximum number of mixture
components in a mixture for the model. This option is ignored by OPM
Flow. & 3 \\
10 & MXDOCOMP & A positive integer defining the maximum number of well
completions that can cross a parallel run domain boundary when the
PARALLEL option has been activated. This option is ignored by OPM Flow.
& 0 \\
11 & MXWSLIST & A positive integer defining the maximum number of well
lists that a well may be concurrent belong to at one time for this
model. This option is ignored by OPM Flow. & 1 \\
12 & MXWLISTS & A positive integer defining the maximum number of
dynamic well lists for this model. This option is ignored by OPM Flow. &
1 \\
13 & MXWSECD & A positive integer defining the maximum number of
secondary wells for this model. This option is ignored by OPM Flow. &
10 \\
14 & MXNGPP & A positive integer defining the maximum number of row
entries per completion in the generalized pseudo-pressure tables used
for to calculate the blocking factor associated with condensate drop-out
in gas condensate reservoirs. If the generalized pseudo-pressure option
has not been activated then this is ignored. This option is ignored by
OPM Flow. & 201 \\
Notes: & & & \\
\end{longtable}

\emph{Table 5.58: WELLDIMS Keyword Description}

\subsubsection{Example}\label{example-129}

\begin{verbatim}
--
--       WELL    WELL    GRUPS   GRUPS
--       MXWELS  MXCONS  MXGRPS  MXGRPW
WELLDIMS
         60      110     18      40                                            /
\end{verbatim}

The above example defines the maximum number of wells to be 60 with 110
completions per well, and maximum number of groups to be 18 with maximum
number of wells per group of 40. All other parameters are defaulted.

\subsection{WPOTCALC -- Well Potential Calculation
Options}\label{wpotcalc-well-potential-calculation-options}

\begin{longtable}[]{@{}
  >{\raggedright\arraybackslash}p{(\linewidth - 14\tabcolsep) * \real{0.1250}}
  >{\raggedright\arraybackslash}p{(\linewidth - 14\tabcolsep) * \real{0.1250}}
  >{\raggedright\arraybackslash}p{(\linewidth - 14\tabcolsep) * \real{0.1250}}
  >{\raggedright\arraybackslash}p{(\linewidth - 14\tabcolsep) * \real{0.1250}}
  >{\raggedright\arraybackslash}p{(\linewidth - 14\tabcolsep) * \real{0.1250}}
  >{\raggedright\arraybackslash}p{(\linewidth - 14\tabcolsep) * \real{0.1250}}
  >{\raggedright\arraybackslash}p{(\linewidth - 14\tabcolsep) * \real{0.1250}}
  >{\raggedright\arraybackslash}p{(\linewidth - 14\tabcolsep) * \real{0.1250}}@{}}
\toprule\noalign{}
\begin{minipage}[b]{\linewidth}\raggedright
RUNSPEC
\end{minipage} & \begin{minipage}[b]{\linewidth}\raggedright
GRID
\end{minipage} & \begin{minipage}[b]{\linewidth}\raggedright
EDIT
\end{minipage} & \begin{minipage}[b]{\linewidth}\raggedright
PROPS
\end{minipage} & \begin{minipage}[b]{\linewidth}\raggedright
REGIONS
\end{minipage} & \begin{minipage}[b]{\linewidth}\raggedright
SOLUTION
\end{minipage} & \begin{minipage}[b]{\linewidth}\raggedright
SUMMARY
\end{minipage} & \begin{minipage}[b]{\linewidth}\raggedright
SCHEDULE
\end{minipage} \\
\midrule\noalign{}
\endhead
\bottomrule\noalign{}
\endlastfoot
\end{longtable}

\subsubsection{Description}\label{description-167}

WPOTCALC defines how shut-in and stopped wells should have their well
potentials calculated. Well potentials for wells under these conditions
need to have their potentials calculated if they are in a Priority
Drilling Queue via the WDRILPRI keyword in the SCHEDULE section, or the
Prioritization option has been enabled by the PRIORITY keyword in the
SCHEDULE section.

This keyword is not supported by OPM Flow but has no effect on the
results so it will be ignored.

\subsection{WSEGDIMS -- Define Multi-Segment Well
Dimensions}\label{wsegdims-define-multi-segment-well-dimensions}

\begin{longtable}[]{@{}
  >{\raggedright\arraybackslash}p{(\linewidth - 14\tabcolsep) * \real{0.1250}}
  >{\raggedright\arraybackslash}p{(\linewidth - 14\tabcolsep) * \real{0.1250}}
  >{\raggedright\arraybackslash}p{(\linewidth - 14\tabcolsep) * \real{0.1250}}
  >{\raggedright\arraybackslash}p{(\linewidth - 14\tabcolsep) * \real{0.1250}}
  >{\raggedright\arraybackslash}p{(\linewidth - 14\tabcolsep) * \real{0.1250}}
  >{\raggedright\arraybackslash}p{(\linewidth - 14\tabcolsep) * \real{0.1250}}
  >{\raggedright\arraybackslash}p{(\linewidth - 14\tabcolsep) * \real{0.1250}}
  >{\raggedright\arraybackslash}p{(\linewidth - 14\tabcolsep) * \real{0.1250}}@{}}
\toprule\noalign{}
\begin{minipage}[b]{\linewidth}\raggedright
RUNSPEC
\end{minipage} & \begin{minipage}[b]{\linewidth}\raggedright
GRID
\end{minipage} & \begin{minipage}[b]{\linewidth}\raggedright
EDIT
\end{minipage} & \begin{minipage}[b]{\linewidth}\raggedright
PROPS
\end{minipage} & \begin{minipage}[b]{\linewidth}\raggedright
REGIONS
\end{minipage} & \begin{minipage}[b]{\linewidth}\raggedright
SOLUTION
\end{minipage} & \begin{minipage}[b]{\linewidth}\raggedright
SUMMARY
\end{minipage} & \begin{minipage}[b]{\linewidth}\raggedright
SCHEDULE
\end{minipage} \\
\midrule\noalign{}
\endhead
\bottomrule\noalign{}
\endlastfoot
\end{longtable}

\subsubsection{Description}\label{description-168}

The WSEGSDIMS keyword defines the multi-segment well dimensions for the
multi-segment well model and the keyword is obligatory if multi-segment
wells are being employed in the model.

\begin{longtable}[]{@{}
  >{\raggedright\arraybackslash}p{(\linewidth - 6\tabcolsep) * \real{0.2500}}
  >{\raggedright\arraybackslash}p{(\linewidth - 6\tabcolsep) * \real{0.2500}}
  >{\raggedright\arraybackslash}p{(\linewidth - 6\tabcolsep) * \real{0.2500}}
  >{\raggedright\arraybackslash}p{(\linewidth - 6\tabcolsep) * \real{0.2500}}@{}}
\toprule\noalign{}
\begin{minipage}[b]{\linewidth}\raggedright
No.
\end{minipage} & \begin{minipage}[b]{\linewidth}\raggedright
Name
\end{minipage} & \begin{minipage}[b]{\linewidth}\raggedright
Description
\end{minipage} & \begin{minipage}[b]{\linewidth}\raggedright
Default
\end{minipage} \\
\midrule\noalign{}
\endhead
\bottomrule\noalign{}
\endlastfoot
1 & MXWELS & A positive integer defining the maximum number of
multi-segment wells for this model. & 0 \\
2 & MXSEGS & A positive integer defining the maximum number of segments
per well for this model. & 1 \\
3 & MXBRAN & A positive integer defining the maximum number of branches
per multi-segment well, including the main branch groups for this model.
& 1 \\
4 & MXLINKS & A positive integer defining the maximum number of segment
links per multi-segment well. The simulator does not currently support
multi-segment chord links, and therefore this parameter is ignored. &
0 \\
Notes: & & & \\
\end{longtable}

\emph{Table 5.59: WSEGDIMS Keyword Description}

\subsubsection{Example}\label{example-130}

\begin{verbatim}
--
--       WELL    WELL    BRANCH  SEGMENT
--       MXWELS  MXSEGS  MXBRAN  MXLINKS
WSEGDIMS
         5       100     10      10                                            /
\end{verbatim}

The above example defines the maximum number of multi-segment wells to
be five with up to 100 segments per multi-segment well, a maximum number
of 10 branches per multi-segment well, and up to 10 segment links per
multi-segment well.

\bookmarksetup{startatroot}

\chapter{GRID SECTION}\label{grid-section-1}

\section{Introduction}\label{introduction-4}

The GRID section defines the basic grid properties, including structure,
faults, and various static rock properties (porosity, permeability,
etc.). The information in this section will be used by the software to
calculate the pore volume (PORV) for each cell, the cell mid-point
depths, and the regular transmissibilities (TRANX, TRANY, and TRANZ)
between all the cells, as well as across faults. The OPM Flow calculated
parameters can then be edited in the EDIT section.

\pandocbounded{\includegraphics[keepaspectratio]{chapters/images/Frame9_53a6f211bd32.png}}
\pandocbounded{\includegraphics[keepaspectratio]{chapters/images/graphics9_53a6f211bd32.png}}

All models can be classified by the number of dimensions as show in
Figure 6.1 (after Mattax {[}Mattax, C.C. and Dalton R.L. 1990. Reservoir
Simulation. Society of Petroleum Engineers, Henry L. Doherty Series,
Monograph Vol. 13{]}). The zero and one dimension models are employed in
analytical modeling, while the higher dimensions are used in numerical
modeling. The term 4D modeling refers to a 3D model with the fourth
dimension being the time domain derived from time-lapse seismic, that is
the comparison of 3D seismic surveys at two or more points in time.

OPM Flow enables the user to define 1D, 2D and 3D models using three
types of grids: Cartesian Regular Grid, Radial Grid, and Irregular
Corner-Point Grids. The first two type of grids are rather limited in
their ability to describe the structural complexity of oil and gas
reservoirs; however, this simplicity allows the engineer to quickly
build simple models to investigate reservoir performance. Indeed in the
early days of numerical modeling back in the late 1970's two-dimensional
cross-section and radial models were the main models used to predict
reservoir performance due to limited computer resources at the time.
That is not to say that full field models were not developed, but that
these full field models were very coarse in comparison to what is
designed and built today using static earth modeling software.

A brief introduction to the three types of grid and the data required to
fully define the structural element of the grid together with the rock
properties necessary to complete the GRID section data requirements is
outlined in the following section. This is then followed by the keyword
definitions applicable to this section.

\section{Keyword Definitions}\label{keyword-definitions-5}

\section{Data Requirements}\label{data-requirements-2}

\subsection{Cartesian Regular Grids}\label{cartesian-regular-grids}

This type of grid defines a regular orthogonal grid based on defining
the sizes of all the cells in the x, y, and z directions and is normally
employed when a complex structural model is not required. Figure 6.2
shows the SPE Comparative Solution Project Number 1 (``SPE-CSP01'') as
documented by Odeh {[}Odeh, A. ``Comparison of Solutions to a Three
Dimensional Black‐Oil Reservoir Simulation Problem.'' JPT 33
(1981):13‐25.{]}.

\pandocbounded{\includegraphics[keepaspectratio]{chapters/images/Frame1_884b5a2d6253.png}}
\pandocbounded{\includegraphics[keepaspectratio]{chapters/images/Image2_884b5a2d6253.png}}

The model consists of a simple 10 x 10 x 3 (NX, NY, NZ) grid and is
defined using the following GRID section keywords to define the grid
geometry:

\begin{verbatim}
-- ==============================================================================
--
-- GRID SECTION
--
-- ==============================================================================
GRID
--
--       DEFINE GRID BLOCK SIZES IN THE X DIRECTION BASED ON NX x NY x NZ = 300)
--       (There Are In Total 300 Cells With Length 1000ft In X-Direction)
DX
         300*1000                                                               /
--
--       DEFINE GRID BLOCK SIZES IN THE Y DIRECTION (BASED ON NX x NY x NZ = 300)
--       (There Are In Total 300 Cells With Length 1000ft In Y-Direction)
DY
         300*1000                                                               /
--
--       DEFINE GRID BLOCK SIZES IN THE Z DIRECTION (BASED ON NX x NY x NZ = 300)
--
DZ
         100*20.0   100*30.0   100*50.0                                         /
--
--       DEFINE GRID BLOCK TOPS FOR THE TOP LAYER (BASED ON NX = 100, NY = 100)
--       (Layer 2 and 3 TOPS Calculated by Simulator)
TOPS
         25*3100  25*3105  25*3110                                              /
\end{verbatim}

The rock property data required to complete the GRID section is as
follows:

\begin{verbatim}
--
--      DEFINE POROSITY DATA FOR ALL CELLS (BASED ON NX x NY x NZ = 300)
--
PORO
        300*0.300                                                               /
--
--      DEFINE PERMY DATA FOR ALL CELLS (BASED ON NX x NY x NZ = 300)
--
PERMX
        100*500.0   100*50.0   100*200.0                                        /                                                                                 --
--      DEFINE PERMY DATA FOR ALL CELLS (BASED ON NX x NY x NZ = 300)
--
PERMY
        100*500.0   100*50.0   100*200.0                                        /
--
--      DEFINE PERMZ DATA FOR ALL CELLS (BASED ON NX x NY x NZ = 300)
--      (Not Defined in Original Paper So Assume That PERMX = PERMY = PERMZ)
PERMZ
        100*500.0   100*50.0   100*200.0                                        /
\end{verbatim}

The above keywords define all the properties required for the GRID
section for this type of grid geometry.

\subsection{Irregular Corner-Point
Grids}\label{irregular-corner-point-grids}

This type of grid is an industry standard grid used to represent the
structure of complex reservoirs {[}K. Ponting, D, Corner Point Geometry
in Reservoir Simulation, Conference Proceedings, ECMOR I - 1st European
Conference on the Mathematics of Oil Recovery, Jul 1989, cp-234-00003
dDOI= ``https://doi.org/10.3997/2214-4609.201411305'', European
Association of Geoscientists \& Engineers.{]}. Here a static modeling
application is used to build the grid which is then exported and
imported into a dynamic simulation model. Figure 6.3 illustrates the
skeleton grid for the Norne Field which has dimensions of 46 x 112 x 22
in the x, y, and z dimensions respectively. This results in a total of
113,344 cells although not all of these cells will be active in the
model.

\pandocbounded{\includegraphics[keepaspectratio]{chapters/images/Frame2_450d297a615e.png}}
\pandocbounded{\includegraphics[keepaspectratio]{chapters/images/Image4_450d297a615e.png}}

Similar to a Cartesian Regular Grid the grid geometry and properties
must be defined for each cell. The formulation of the grid geometry is
based on a corner-point geometry, where basically the coordinate lines
or pillars are given, then the top and bottom surfaces for the cell are
defined by specifying the depth (z-coordinates) of the cell's corner
points along each of the four adjacent pillars. The cell then forms an
irregular hexahedron as depicted in Figure 6.4. Note that the figure
shows a corner-point cell which is more or less orthogonal, which is
ideally what we want to minimize grid orientation effects.

The data required to define this type of grid consists of the SPECGRID
keyword to define the dimensions of the grid, that is:

\begin{verbatim}
--       MAX     MAX     MAX     MAX     GRID
--       NDIVIX  NDIVIY  NDIVIZ  NUMRES  TYPE
SPECGRID
         46      112     22      1       F                                     /
\end{verbatim}

A portion of the coordinate line data defined by the COORD keyword from
the Norne model is shown on the next page.

\pandocbounded{\includegraphics[keepaspectratio]{chapters/images/Frame3_106d554ea496.png}}
\pandocbounded{\includegraphics[keepaspectratio]{chapters/images/Image5_106d554ea496.png}}

\begin{verbatim}
COORD
-- X1           Y1             Z1          X2           Y2             Z2
-- ---          ---            ----        ---          ---            ----
   453114.000   7319921.000    3037.473    453114.000   7319921.000    3132.831
   453155.031   7319840.000    2983.933    453142.750   7319864.000    3173.572
   453196.094   7319759.000    3005.969    453171.500   7319807.500    3215.836
   453237.156   7319678.000    3000.265    453200.250   7319751.000    3217.250
   453278.188   7319597.000    2989.348    453229.031   7319694.000    3213.951
   453319.250   7319516.500    2995.680    453257.781   7319637.500    3215.323
   453356.250   7319443.500    3000.855    453308.750   7319537.000    3220.549
   453393.250   7319370.500    3005.252    453359.688   7319436.500    3210.393
   453423.969   7319310.000    3030.862    453394.219   7319368.500    3203.438
   453454.688   7319249.500    3036.870    453428.719   7319300.500    3190.770
   453485.406   7319189.000    3038.017    453463.219   7319232.500    3190.660
   453516.125   7319128.000    3045.027    453497.750   7319164.500    3188.813
   453546.844   7319067.500    3055.410    453532.250   7319096.500    3185.966
   453577.562   7319007.000    3066.541    453566.750   7319028.500    3184.325
   453608.281   7318946.500    3076.624    453601.250   7318960.500    3183.584
   453639.000   7318886.000    3086.938    453635.781   7318892.500    3184.057
   453669.719   7318825.500    3096.153    453670.281   7318824.500    3185.988
   453700.438   7318765.000    3104.703    453704.781   7318756.500    3188.598
   453731.156   7318704.500    3097.016    453739.281   7318688.500    3180.484
   453761.875   7318644.000    3088.539    453773.812   7318620.500    3177.091
   453780.000   7318608.000    3098.118    453796.562   7318575.500    3176.401
   453798.125   7318572.500    3096.691    453819.344   7318530.500    3172.299
…………………………………….
/
\end{verbatim}

The final keyword to define an Irregular Corner-Point geometry grid is
the ZCORN keyword that defines the depths of the cell corners. A portion
of the ZCORN data from the Norne model is shown below.

\begin{verbatim}
ZCORN
   3037.473    2983.933    2983.933    3005.969    3005.969    3000.265
   3000.265    2989.348    2989.348    2995.680    2995.680    3000.855
   3000.855    3005.252    3005.252    3030.862    3030.862    3036.870
   3036.870    3038.017    3038.017    3045.027    3045.027    3055.410
   3055.410    3066.541    3066.541    3076.624    3076.624    3086.938
   3086.938    3096.153    3096.153    3104.703    3104.703    3097.016
   3097.016    3088.539    3088.539    3098.118    3098.118    3096.691
   3096.691    3093.886    3093.886    3085.393    3085.393    3081.957
   3081.957    3080.645    3080.645    3115.021    3115.021    3130.474
   3130.474    3204.674    3204.674    3193.187    3193.187    3169.512
   3169.512    3101.928    3101.928    3044.277    3044.277    3023.930
   3023.930    2964.244    2964.244    2900.178    2900.178    2875.715
   2875.715    2864.913    2864.913    2855.256    2855.256    2841.119
   2841.119    2826.261    2826.261    2806.556    2806.556    2781.052
   2781.052    2791.720    2791.720    2817.940    2817.940    2813.308
   2813.308    2788.492
…………………………………………………………………………………………………
/
\end{verbatim}

The rock property data required to complete the GRID section is the same
as for a Cartesian Regular grid, as defined in section 6.2.1 Cartesian
Regular Grids and the data is defined using the same keywords. The
resulting Norne model showing the ternary solution variable is displayed
in Figure 6.5.

\pandocbounded{\includegraphics[keepaspectratio]{chapters/images/Frame4_1e6c544544c0.png}}
\pandocbounded{\includegraphics[keepaspectratio]{chapters/images/Image3_1e6c544544c0.png}}

\subsection{Radial Grids}\label{radial-grids}

Radial flow exists near the wellbore and linear or Cartesian flow occurs
within the reservoir away from the wellbore. For field-wide studies the
growth or decay of reservoir volumes activated by unsteady-state flow is
small and hence linear flow is modeled via a linear grid, using either a
Cartesian Regular Grid or an Irregular Corner-Point Grid, in these type
of studies. On the other hand, for converging flow near wells the growth
or decay of reservoir volumes activated by unsteady-state flow is large
and radial flow should modeled. Thus, radial flow is important in only
very localized areas around the wells and linear flow is modeled in most
reservoir studies, for example in full field or sector models.
Consequently for investigating near well bore effects, radial models are
used to investigate the behavior around the wellbore, for example in gas
cusping and water conning studies, as one can see from Figure 6.6.

\pandocbounded{\includegraphics[keepaspectratio]{chapters/images/Frame68_4c2e14c3afcb.png}}
\pandocbounded{\includegraphics[keepaspectratio]{chapters/images/graphics10_4c2e14c3afcb.png}}

In both Cartesian Regular and Irregular Corner-Point grids the dimension
nomenclature is (x, y, z) for the x, y, and z planes. The corresponding
nomenclature for Radial Grids is (r, theta, and z) for r or radius
plane, theta for the angular plane measured in degrees, and z for the z
plane.

In Figure 6.6 the radial grid has been labeled as a 2-D Cylindrical grid
because there is no flow in the theta direction only in the r and z
planes, whereas the 3-D Reservoir grid in the figure has flow in both
the x and y planes as well as the z plane. Figure 6.7 shows a typical
3-D radial grid with dimensions (10, 10, 3) in the (r, theta, z)
dimensions taken from Odeh {[}Comparison of Solutions to a
Three-Dimensional Black-Oil Reservoir Simulation Problem by Aziz S.
Odeh,-- Journal of Petroleum Technology, January 1981.{]}.

\pandocbounded{\includegraphics[keepaspectratio]{chapters/images/Frame69_f0eff26a3ebd.png}}
\pandocbounded{\includegraphics[keepaspectratio]{chapters/images/Image62_f0eff26a3ebd.png}}

Radial grids normally have a fine grid near the wellbore which then
expands logarithmically away from the wellbore. This is illustrated in
Figure 6.7 where one can clearly see the four most outer rings, but not
the inner six. This type of grid makes solving the flow equations more
challenging than conventional Cartesian and Irregular Corner-Point
grids, where the cell pore volume distribution is relatively similar for
neighboring cells, compared with Radial grids. For comparison, a typical
conventional grid will have grid blocks size of 300, 300, 3 feet (or
approximately 100m x 100m x 1m) spacing in the (x, y, z) plane; whereas
a radial grid will utilize an inner most radius of between 0.25 to 1.0
foot in the R direction, and then logarithmically expand from the inner
radius. The Odeh110 example employs 0.25 ft. for the inner radius, and
1.75, 2.32, 5.01, 10.84, 23.39, 50.55, 109.21, 235.92, 509.68, and
1101.08 ft. for the outer radii.

Figure 6.7 is based on a 10 x 10 x 3 grid, where NX is the R dimensions,
NY the THETA dimensions, and NZ the z dimensions on the DIMENS keyword
in the RUNSPEC section. Thus, in order to fully define this radial grid
the following RUNSPEC and GRID section keywords are required:

\begin{verbatim}
-- ==============================================================================
--
-- RUNSPEC SECTION
--
-- ==============================================================================
RUNSPEC
--
--       MAX     MAX     MAX
--       NDIVIX  NDIVIY  NDIVIZ
DIMENS
         10      10      3                                                     /
--
--       DEFINE RADIAL GRID GEOMETRY
--
RADIAL
-- ==============================================================================
--
-- GRID SECTION
--
-- ==============================================================================
GRID
--
--       INNER RADIUS OF FIRST GRID BLOCK IN THE RADIAL DIRECTION
--
INRAD
         0.25
/
--
--       DEFINE GRID BLOCK SIZES IN THE R DIRECTION
--
DRV
         1.75  2.32  5.01  10.84  23.39  50.55  109.21  235.92  509.68  1101.08 /
--
--       DEFINE GRID BLOCK SIZES IN THE THETA DIRECTION (BASED ON NY = 10)
--
DTHETAV
         10*36                                                                  /
--
--       DEFINE GRID BLOCK Z DIRECTION CELL SIZE (BASED ON NX x NY x NZ = 300)
--
DZ
         100*20.0   100*30.0   100*50.0                                         /
--
--       DEFINE GRID BLOCK TOPS FOR THE TOP LAYER (NX=10,NY=10,and NZ=3)
--
TOPS
         100*8325                                                               /
\end{verbatim}

The rock property data required to complete the GRID section is as
follows:

\begin{verbatim}
--
--       DEFINE POROSITY DATA FOR ALL CELLS (BASED ON NX x NY x NZ = 300)
--
PORO
         300*0.300                                                              /
--
--       DEFINE GRID BLOCK PERMR DATA FOR ALL CELLS (BASED ON NR x NY x NZ = 300)
--
PERMR
         100*500.0   100*50.0   100*200.0                                       /
--
--       DEFINE GRID BLOCK PERMTHT DATA FOR ALL CELLS (NX x NY x NZ = 300)
--
PERMTHT
         100*500.0   100*50.0   100*200.0                                       /
--
--        DEFINE GRID BLOCK Z DIRECTION CELL SIZE (BASED ON NX x NY x NZ = 300)
--
PERMZ
         100*500.0   100*50.0   100*200.0                                       /
\end{verbatim}

The above keywords define all the properties required for the RUNSPEC
and GRID sections for this type of grid geometry.

\begin{tcolorbox}[enhanced jigsaw, title=\textcolor{quarto-callout-note-color}{\faInfo}\hspace{0.5em}{Note}, coltitle=black, breakable, toprule=.15mm, colbacktitle=quarto-callout-note-color!10!white, leftrule=.75mm, opacitybacktitle=0.6, colback=white, bottomtitle=1mm, toptitle=1mm, titlerule=0mm, left=2mm, colframe=quarto-callout-note-color-frame, arc=.35mm, rightrule=.15mm, bottomrule=.15mm, opacityback=0]

Radial grids are currently not fully implemented in OPM Flow as the
completion of the circle in the THETA direction is not supported. For
reference only, this feature is activated via the COMPLETE parameter on
the COORDSYS keyword in the GRID section.

\end{tcolorbox}

OPM ResInsight is the 3D visualization software that is used to view the
grid block property data and OPM Flow results (grid block pressures,
fluid saturations, etc.). If the radial grid has a low number of angular
cells, OPM ResInsight will automatically create a temporary Local Grid
Refinement (``LGR'') with higher angular refinement for improved
visualization. Figure 6.8 shows how the Odeh {[}Comparison of Solutions
to a Three-Dimensional Black-Oil Reservoir Simulation Problem by Aziz S.
Odeh,-- Journal of Petroleum Technology, January 1981.{]} radial model
is rendered in OPM ResInsight.

\pandocbounded{\includegraphics[keepaspectratio]{chapters/images/Frame29_704598d1522c.png}}
\pandocbounded{\includegraphics[keepaspectratio]{chapters/images/Image26_704598d1522c.png}}

The previous figure (Figure 6.8) shows how the model would be displayed
if the commercial simulator run was loaded into OPM ResInsight. Support
for displaying radial grids was added to OPM ResInsight in version
2025.12.

To overcome the display limitation in versions of OPM ResInsight prior
to version 2025.12, OPM Flow converts radial grids to Irregular
Corner-Point Grids (Cartesian coordinates) and adjusts the model to
reflect radial coordinates. Thus, the OPM Flow radial model has the same
pore volumes as the commercial simulator radial model. However, since
this conversion is no longer necessary it is anticipated that it will be
removed in a future version of OPM Flow.

\pandocbounded{\includegraphics[keepaspectratio]{chapters/images/Frame30_d2b36a02f78d.png}}
\pandocbounded{\includegraphics[keepaspectratio]{chapters/images/Image27_d2b36a02f78d.png}}

Figure 6.9 shows how the OPM Flow run is displayed in OPM ResInsight. As
will be shown in section 6.2.4 Spider Grids, both Spider and Radial
grids generated by OPM Flow will be displayed the same, as they both use
Irregular Corner-Point Grid geometry. The difference is that the radial
model's pore volumes are adjusted to match the radial geometry pore
volumes, whereas, the Spider grid volumes are not adjusted.

\subsection{Spider Grids}\label{spider-grids}

Spider Grids have been implemented as an alternative to Radial Grids.
They convert standard radial geometry keywords in the grid section into
an Irregular Corner-Point Grid, which can then be visualized in OPM
ResInsight. The only difference is that in the RUNSPEC section the
SPIDER keyword is used instead of the RADIAL keyword.

\pandocbounded{\includegraphics[keepaspectratio]{chapters/images/Frame71_a8a1634db92f.png}}
\pandocbounded{\includegraphics[keepaspectratio]{chapters/images/Image64_a8a1634db92f.png}}

Figure 6.10 illustrates the spider grid for the Odeh111 example, which
is not dissimilar to the conventional radial model shown in Figure 6.7
and is identical to OPM Flow's radial model implementation shown in
Figure 6.9.

Naturally, there will be differences in the results between radial and
spider grid formulations, but the user should use their own judgement if
the differences are relevant.

Again, Figure 6.10 is based on a 10 x 10 x 3 grid, where NX is the R
dimensions, NY the THETA dimensions and NZ the z dimensions on the
DIMENS keyword in the RUNSPEC section. The only difference between this
example and the previous example for radial grids is that the RADIAL
keyword has been replaced by the SPIDER keyword in the RUNSPEC section,
all the other keywords are exactly the same.

\begin{verbatim}
-- ==============================================================================
--
-- RUNSPEC SECTION
--
-- ==============================================================================
RUNSPEC
--
--       MAX     MAX     MAX
--       NDIVIX  NDIVIY  NDIVIZ
DIMENS
         10      1       3                                                     /
--
--       DEFINE SPIDER GRID GEOMETRY (OPM FLOW RADIAL GRID KEYWORD)
--
SPIDER
-- ==============================================================================
--
-- GRID SECTION
--
-- ==============================================================================
GRID
--
--       INNER RADIUS OF FIRST GRID BLOCK IN THE RADIAL DIRECTION
--
INRAD
         0.25
/
--
--       DEFINE GRID BLOCK SIZES IN THE R DIRECTION
--
DRV
         1.75  2.32  5.01  10.84  23.39  50.55  109.21  235.92  509.68  1101.08 /
--
--       DEFINE GRID BLOCK SIZES IN THE THETA DIRECTION (BASED ON NY = 10)
--
DTHETAV
         1*45                                                                   /
--
--       DEFINE GRID BLOCK Z DIRECTION CELL SIZE (BASED ON NX x NY x NZ = 30)
--
DZ
         10*20.0   10*30.0   10*50.0                                            /
--
--       DEFINE GRID BLOCK TOPS FOR THE TOP LAYER (NX=10,NY=1,and NZ=3)
--
TOPS
         10*8325                                                               /
\end{verbatim}

The rock property data required to complete the GRID section is as
follows:

\begin{verbatim}
--
--       DEFINE POROSITY DATA FOR ALL CELLS (BASED ON NX x NY x NZ = 30)
--
PORO
         30*0.300                                                              /
--
--       DEFINE GRID BLOCK PERMR DATA FOR ALL CELLS (BASED ON NR x NY x NZ = 30)
--
PERMR
         10*500.0   10*50.0   10*200.0                                         /

--
--       DEFINE GRID BLOCK PERMTHT DATA FOR ALL CELLS (NX x NY x NZ = 30)
--
PERMTHT
         10*500.0   10*50.0   10*200.0                                        /
--
--        DEFINE GRID BLOCK Z DIRECTION CELL SIZE (BASED ON NX x NY x NZ = 300)
--
PERMZ
         10*500.0   10*50.0   10*200.0                                       /
\end{verbatim}

The above keywords define all the properties required for the GRID
section for this type of grid geometry.

\subsection{Rock Properties}\label{rock-properties}

Regardless of the grid type used to define the structural component of
the model, certain static properties must be specified to complete the
grid definition. These include identifying active and inactive grid
blocks, porosity, permeability, and reservoir quality through the
net-to-gross fraction (``NTG''). These parameters must be set for each
cell in the model

\begin{longtable}[]{@{}
  >{\raggedright\arraybackslash}p{(\linewidth - 6\tabcolsep) * \real{0.2500}}
  >{\raggedright\arraybackslash}p{(\linewidth - 6\tabcolsep) * \real{0.2500}}
  >{\raggedright\arraybackslash}p{(\linewidth - 6\tabcolsep) * \real{0.2500}}
  >{\raggedright\arraybackslash}p{(\linewidth - 6\tabcolsep) * \real{0.2500}}@{}}
\toprule\noalign{}
\begin{minipage}[b]{\linewidth}\raggedright
Property
\end{minipage} & \begin{minipage}[b]{\linewidth}\raggedright
Description
\end{minipage} & \begin{minipage}[b]{\linewidth}\raggedright
Cartesian and Irregular Corner-Point Grids Keywords
\end{minipage} & \begin{minipage}[b]{\linewidth}\raggedright
Radial Grid Keywords
\end{minipage} \\
\midrule\noalign{}
\endhead
\bottomrule\noalign{}
\endlastfoot
Active and Inactive cells & Defines if a cell in the model is active by
setting the ACTNUM property for a cell to either one or inactive by
setting the value to zero. Cells that are inactive in the model are
ignored computationally and can act as barriers to flow. Thus, a shale
in a conventional reservoir is normally treated as non-reservoir and is
made inactive either by setting the ACTNUM, PORO, or NTG to zero for the
cells representing the shale. & ACTNUM & \\
Porosity & Porosity is a measure of the space in a reservoir rock. It is
defined as the fraction of the total bulk volume of the rock not
occupied by solids, that is it is the fraction of the cell that is
porous and contains the reservoir fluids. & PORO & \\
Reservoir Quality & Reservoir quality of the cell in terms of the gross
volume derived from the structural grid and the net volume available for
fluid flow in the model is expressed as a fraction from zero to one. A
zero values means the cell does not contribute to flow and therefore is
made inactive. A value of one means the gross and net volumes are
identical for the cell & NTG & \\
Permeability & Permeability is a measure of the ease with which a fluid
will flow through a porous medium. In numerical models permeability is
dependent on the direction of flow, that is x, y, and z directions in
Cartesian and Irregular Corner-Point Grids, and the radial, theta, and z
directions in radial grids. There are various formulations for
permeability, absolute permeability, effective permeability, gas
permeability, liquid permeability, etc., and the values are saturation
dependent. Thus, values entered should be consistent with the relative
permeability entered in the PROPS section. Normally Kair (Sg=1.0) should
be entered for the cell permeability and the values may or not be
corrected for overburden or humidity drying effects. Correcting for
liquid flow and saturation end points etc., is accomplished by the
relative permeability curves. For example, if Kair (Sg=1.0) has been
entered for the cell permeability when Krg (Sg=1-Swc) should be less
than one. & PERMX PERMY PERMZ & PERMR PERMTHT PERMZ \\
\end{longtable}

\emph{Table 6.1: Key Static Grid Properties}

\begin{tcolorbox}[enhanced jigsaw, title=\textcolor{quarto-callout-note-color}{\faInfo}\hspace{0.5em}{Note}, coltitle=black, breakable, toprule=.15mm, colbacktitle=quarto-callout-note-color!10!white, leftrule=.75mm, opacitybacktitle=0.6, colback=white, bottomtitle=1mm, toptitle=1mm, titlerule=0mm, left=2mm, colframe=quarto-callout-note-color-frame, arc=.35mm, rightrule=.15mm, bottomrule=.15mm, opacityback=0]

Static grid properties are frequently generated from a static earth
model using petrophysical evaluation of the well logs and propagated
through to the model based on a variety of geostatistical techniques.
Petrophysical evaluations are conducted in either the ``Total'' or the
``Effective'' porosity domain, and it is important that all the rock
property data entered into the model is of the same basis. It is not
important which porosity domain is used, as long as all the data is in
the same domain.

\end{tcolorbox}

Pore volume and transmissibility are common terms in the reservoir
simulation vernacular. Pore volume is self-explanatory, that is, given
the grid property data the pore volume for each cell is calculated
using:

\begin{equation}\phantomsection\label{eq-6-1}{
\mathit{PV} = \mathit{Cell}\mathit{Gross}\mathit{Volume}\times \mathit{PORO}\times \mathit{NTG}\times \mathit{ACTNUM}
}\end{equation}

Where

PV = the pore volume of a cell,

Cell Gross Volume = the gross volume (or bulk volume) calculated from
the structural

parameters of the cell,

PORO = cell porosity,

NTG = cell net-to-gross ratio, and

ACTNUM = active and inactive cell indicator.

Any cell with a pore volume equal to zero is made inactive automatically
in the model. However, there may be some cells that have small pore
volumes that may negatively impact computational performance of the
model. If this is the case then the MINPV and MINPORV keywords in the
GRID section can be used to make these cells inactive.

There has been a trend in the industry in recent years to not apply
petrophysical cut-offs to determine net volumes in static models. This
results in large models with numerous cells with very low porosity
values (less than 0.01 for example) and corresponding very low
permeabilities. The theory behind this approach is that the numerical
model will determine the effective (or net) reservoir. However, the
approach is questionable as by not applying net pay cut-offs, the
in-place and recoverable volumes may not satisfy Reserve reporting
requirements and guidelines to various agencies. Secondly, although this
may be appropriate in unconventional reservoirs, as all the cells in the
model will have similar values of porosity and permeability, but in
conventional reservoirs this methodology will lead to severe
computational issues when attempting to run the model, due to very tight
cells being next to relative high permeability cells. Again, the MINPV
and MINPORV keywords can be used to resolve this issue.

Transmissibility on the other hand is more complex as it relates the
flow from one cell face to another cell face and is a function of the
area open to flow, the direction of flow, the permeability, saturation,
and viscosity of the phases flowing between the cells. For single phase
flow in a Cartesian grid the x-direction transmissibility is of the
form:

\begin{equation}\phantomsection\label{eq-6-2}{
{T}_{{x}_{i+½, j}} = {\left[\frac{{k}_{x} h (\mathit{Δy})}{μ (\mathit{Δx})}\right] }_{i+½, j}
}\end{equation}

As transmissibility is a property of the flow between two cell faces,
not a block centered grid cell property like porosity or permeability,
then the nomenclature for transmissibility is different. In OPM Flow,
the transmissible of cell face Tx(i, j, k) is the transmissibility
between cells (i, j, k) and (i+1, j, k). In some simulators it would be
between (i, j, k) and (i-1, j, k). This is important to note if manual
modifications to cell connections are to be made in the model.

Modifications to basic grid property data (porosity, permeability, etc.)
can only be done in the GRID section, thereafter only the calculated
pore volumes and transmissibilities are available for adjustment in the
EDIT section. Some engineers use the GRID section for the basic model
formulation, and the EDIT section to document the changes to the
``base'' model using the PORV, TRANX, TRANY, TRANZ, etc. keywords.

\section{Keyword Definitions}\label{keyword-definitions-6}

\subsection{ACTNUM -- Set the Status of a Grid Block To Active or
Inactive}\label{actnum-set-the-status-of-a-grid-block-to-active-or-inactive}

\begin{longtable}[]{@{}
  >{\raggedright\arraybackslash}p{(\linewidth - 14\tabcolsep) * \real{0.1250}}
  >{\raggedright\arraybackslash}p{(\linewidth - 14\tabcolsep) * \real{0.1250}}
  >{\raggedright\arraybackslash}p{(\linewidth - 14\tabcolsep) * \real{0.1250}}
  >{\raggedright\arraybackslash}p{(\linewidth - 14\tabcolsep) * \real{0.1250}}
  >{\raggedright\arraybackslash}p{(\linewidth - 14\tabcolsep) * \real{0.1250}}
  >{\raggedright\arraybackslash}p{(\linewidth - 14\tabcolsep) * \real{0.1250}}
  >{\raggedright\arraybackslash}p{(\linewidth - 14\tabcolsep) * \real{0.1250}}
  >{\raggedright\arraybackslash}p{(\linewidth - 14\tabcolsep) * \real{0.1250}}@{}}
\toprule\noalign{}
\begin{minipage}[b]{\linewidth}\raggedright
RUNSPEC
\end{minipage} & \begin{minipage}[b]{\linewidth}\raggedright
GRID
\end{minipage} & \begin{minipage}[b]{\linewidth}\raggedright
EDIT
\end{minipage} & \begin{minipage}[b]{\linewidth}\raggedright
PROPS
\end{minipage} & \begin{minipage}[b]{\linewidth}\raggedright
REGIONS
\end{minipage} & \begin{minipage}[b]{\linewidth}\raggedright
SOLUTION
\end{minipage} & \begin{minipage}[b]{\linewidth}\raggedright
SUMMARY
\end{minipage} & \begin{minipage}[b]{\linewidth}\raggedright
SCHEDULE
\end{minipage} \\
\midrule\noalign{}
\endhead
\bottomrule\noalign{}
\endlastfoot
\end{longtable}

\subsubsection{Description}\label{description-169}

The ACTNUM keyword specifies which grid blocks are either active or
inactive. A grid block is automatically set to inactive if its pore
volume is less than the value entered using the MINPV keyword. The
ACTNUM keyword can be used to also make blocks with pore volumes greater
than MINPV inactive.

\begin{longtable}[]{@{}
  >{\raggedright\arraybackslash}p{(\linewidth - 6\tabcolsep) * \real{0.2500}}
  >{\raggedright\arraybackslash}p{(\linewidth - 6\tabcolsep) * \real{0.2500}}
  >{\raggedright\arraybackslash}p{(\linewidth - 6\tabcolsep) * \real{0.2500}}
  >{\raggedright\arraybackslash}p{(\linewidth - 6\tabcolsep) * \real{0.2500}}@{}}
\toprule\noalign{}
\begin{minipage}[b]{\linewidth}\raggedright
No.
\end{minipage} & \begin{minipage}[b]{\linewidth}\raggedright
Name
\end{minipage} & \begin{minipage}[b]{\linewidth}\raggedright
Description
\end{minipage} & \begin{minipage}[b]{\linewidth}\raggedright
Default
\end{minipage} \\
\midrule\noalign{}
\endhead
\bottomrule\noalign{}
\endlastfoot
1 & ACTNUM & An array of integers equal to either 0 or 1 that define the
activity of each cell in the model. A value of 0 indicates the cell is
inactive. Grid blocks are ordered with the I index cycling fastest,
followed by the J and K indices. Repeat counts may be used, for example
20*1. & 1 \\
Notes: & & & \\
\end{longtable}

\emph{Table 6.2: ACTNUM Keyword Description}

\subsubsection{Examples}\label{examples-6}

The example below sets several cells to be inactive for a 4 x 5 x 2
model.

\begin{verbatim}
ACTNUM
-- Layer 1
0 0 1 1
0 0 1 1
1 1 1 1
1 1 1 1
1 1 1 1
-- Layer 2
1 1 1 1
1 1 1 1
1 1 1 1
1 1 1 1
0 0 0 0
/
\end{verbatim}

Alternatively the EQUALS keyword could be employed to accomplish the
same task, that is:

\begin{verbatim}
--       -- ARRAY    CONSTANT --  ---------- BOX ---------
--                                I1  I2   J1  J2   K1  K2
EQUALS
         'ACTNUM’    1.0000       1*  1*   1*  1*   1*  1* / SET ACTIVE CELLS
         'ACTNUM’    0.0000       1   2    1   2    1   1  / SET INACTIVE CELLS
         'ACTNUM’    0.0000       1   4    5   5    2   2  / SET INACTIVE CELLS
/
\end{verbatim}

\subsection{ADD -- Add a Constant to a Specified
Array}\label{add-add-a-constant-to-a-specified-array}

\begin{longtable}[]{@{}
  >{\raggedright\arraybackslash}p{(\linewidth - 14\tabcolsep) * \real{0.1250}}
  >{\raggedright\arraybackslash}p{(\linewidth - 14\tabcolsep) * \real{0.1250}}
  >{\raggedright\arraybackslash}p{(\linewidth - 14\tabcolsep) * \real{0.1250}}
  >{\raggedright\arraybackslash}p{(\linewidth - 14\tabcolsep) * \real{0.1250}}
  >{\raggedright\arraybackslash}p{(\linewidth - 14\tabcolsep) * \real{0.1250}}
  >{\raggedright\arraybackslash}p{(\linewidth - 14\tabcolsep) * \real{0.1250}}
  >{\raggedright\arraybackslash}p{(\linewidth - 14\tabcolsep) * \real{0.1250}}
  >{\raggedright\arraybackslash}p{(\linewidth - 14\tabcolsep) * \real{0.1250}}@{}}
\toprule\noalign{}
\begin{minipage}[b]{\linewidth}\raggedright
RUNSPEC
\end{minipage} & \begin{minipage}[b]{\linewidth}\raggedright
GRID
\end{minipage} & \begin{minipage}[b]{\linewidth}\raggedright
EDIT
\end{minipage} & \begin{minipage}[b]{\linewidth}\raggedright
PROPS
\end{minipage} & \begin{minipage}[b]{\linewidth}\raggedright
REGIONS
\end{minipage} & \begin{minipage}[b]{\linewidth}\raggedright
SOLUTION
\end{minipage} & \begin{minipage}[b]{\linewidth}\raggedright
SUMMARY
\end{minipage} & \begin{minipage}[b]{\linewidth}\raggedright
SCHEDULE
\end{minipage} \\
\midrule\noalign{}
\endhead
\bottomrule\noalign{}
\endlastfoot
\end{longtable}

\subsubsection{Description}\label{description-170}

The ADD keyword adds a constant to a specified array or part of an
array. The constant can be an integer or real value depending on the
array type; however, the arrays that can be operated on are dependent on
which section the ADD keyword is being applied in.

\begin{longtable}[]{@{}
  >{\raggedright\arraybackslash}p{(\linewidth - 6\tabcolsep) * \real{0.2500}}
  >{\raggedright\arraybackslash}p{(\linewidth - 6\tabcolsep) * \real{0.2500}}
  >{\raggedright\arraybackslash}p{(\linewidth - 6\tabcolsep) * \real{0.2500}}
  >{\raggedright\arraybackslash}p{(\linewidth - 6\tabcolsep) * \real{0.2500}}@{}}
\toprule\noalign{}
\begin{minipage}[b]{\linewidth}\raggedright
No.
\end{minipage} & \begin{minipage}[b]{\linewidth}\raggedright
Name
\end{minipage} & \begin{minipage}[b]{\linewidth}\raggedright
Description
\end{minipage} & \begin{minipage}[b]{\linewidth}\raggedright
Default
\end{minipage} \\
\midrule\noalign{}
\endhead
\bottomrule\noalign{}
\endlastfoot
1 & ARRAY & A character string of up to eight characters in length that
defines the keyword identifying the property to be modified. & None \\
2 & CONSTANT & An integer or real value to be added to the ARRAY in the
same units as the ARRAY property. & 0 \\
3 & I1 & A positive integer that defines the lower bound of the array in
the I-direction to be modified must be greater than or equal 1 and less
than or equal to I2 and NX. & 1 \\
4 & I2 & A positive integer that defines the upper bound of the array in
the I-direction to be modified must be greater than or equal to II and
less than or equal to NX & NX \\
5 & J1 & A positive integer that defines the lower bound of the array in
the J-direction to be modified must be greater than or equal 1 and less
than or equal to J2 and NY. & 1 \\
6 & J2 & A positive integer that defines the upper bound of the array in
the J-direction to be modified must be greater than or equal to JI and
less than or equal to NY. & NY \\
7 & K1 & A positive integer that defines the lower bound of the array in
the K-direction to be modified must be greater than or equal to one and
less than or equal to K2 and NZ. & 1 \\
8 & K2 & A positive integer that defines the upper bound of the array in
the K-direction to be modified must be greater than or equal to KI and
less than or equal to NZ. & NZ \\
Notes: & & & \\
\end{longtable}

\emph{Table 6.3: ADD Keyword Description}

Examples of the arrays most commonly operated on in each section are
given in Table 6.4. Cells colored red indicate arrays that are not
currently supported by OPM Flow operations.

\begin{longtable}[]{@{}lllllll@{}}
\toprule\noalign{}
ADD Keyword and Variable Options by Section & & & & & & \\
\midrule\noalign{}
\endhead
\bottomrule\noalign{}
\endlastfoot
GRID & EDIT & PROPS & REGIONS & SOLUTION & SUMMARY & SCHEDULE \\
DX & DEPTH & SWL & ENDNUM & PRESSURE & & \\
DY & PORV & SWCR & EQLNUM & SWAT & & \\
DZ & TRANX & SWU & FIPNUM & SGAS & & \\
PERMX & TRANY & SGL & IMBNUM & RV & & \\
PERMY & TRANZ & SGCR & MISCNUM & RS & & \\
PERMZ & DIFFX & SGU & PVTNUM & TBLK & & \\
MULTX & DIFFY & KRW & ROCKNUM & GI & & \\
MULTY & DIFFZ & KRO & SATNUM & OILAPI & & \\
MULTZ & TRANR & KRG & WH2NUM & SALT & & \\
DR & TRANTHT & PCG & & GASCONC & & \\
DTHETA & DIFFR & PCW & & SOLVCONC & & \\
PERMR & DIFFTHT & & & SOLVFRAC & & \\
PERMTHT & & & & SFOAM & & \\
DZNET & & & & SPOLY & & \\
PORO & & & & & & \\
NTG & & & & & & \\
FLUXNUM & & & & & & \\
MULTNUM & & & & & & \\
MPFANUM & & & & & & \\
DIFFX & & & & & & \\
DIFFY & & & & & & \\
DIFFZ & & & & & & \\
DIFFR & & & & & & \\
DIFFTHT & & & & & & \\
\end{longtable}

\emph{Table 6.4: ADD Keyword Applicable Arrays by Section}

\subsubsection{Example}\label{example-131}

\begin{verbatim}
--
--       ARRAY      CONSTANT     ---------- BOX ---------
--                                I1  I2   J1  J2   K1  K2
ADD
         PERMX      20.000        1*  1*   1*  1*   1*  1* / ADD 20 mD TO PERMX
/
\end{verbatim}

The above example adds 20 md to the PERMX property array throughout the
model.

\subsection{ADDREG -- Add a Constant to an Array based on a Region
Number}\label{addreg-add-a-constant-to-an-array-based-on-a-region-number}

\begin{longtable}[]{@{}
  >{\raggedright\arraybackslash}p{(\linewidth - 14\tabcolsep) * \real{0.1250}}
  >{\raggedright\arraybackslash}p{(\linewidth - 14\tabcolsep) * \real{0.1250}}
  >{\raggedright\arraybackslash}p{(\linewidth - 14\tabcolsep) * \real{0.1250}}
  >{\raggedright\arraybackslash}p{(\linewidth - 14\tabcolsep) * \real{0.1250}}
  >{\raggedright\arraybackslash}p{(\linewidth - 14\tabcolsep) * \real{0.1250}}
  >{\raggedright\arraybackslash}p{(\linewidth - 14\tabcolsep) * \real{0.1250}}
  >{\raggedright\arraybackslash}p{(\linewidth - 14\tabcolsep) * \real{0.1250}}
  >{\raggedright\arraybackslash}p{(\linewidth - 14\tabcolsep) * \real{0.1250}}@{}}
\toprule\noalign{}
\begin{minipage}[b]{\linewidth}\raggedright
RUNSPEC
\end{minipage} & \begin{minipage}[b]{\linewidth}\raggedright
GRID
\end{minipage} & \begin{minipage}[b]{\linewidth}\raggedright
EDIT
\end{minipage} & \begin{minipage}[b]{\linewidth}\raggedright
PROPS
\end{minipage} & \begin{minipage}[b]{\linewidth}\raggedright
REGIONS
\end{minipage} & \begin{minipage}[b]{\linewidth}\raggedright
SOLUTION
\end{minipage} & \begin{minipage}[b]{\linewidth}\raggedright
SUMMARY
\end{minipage} & \begin{minipage}[b]{\linewidth}\raggedright
SCHEDULE
\end{minipage} \\
\midrule\noalign{}
\endhead
\bottomrule\noalign{}
\endlastfoot
\end{longtable}

\subsubsection{Description}\label{description-171}

The ADDREG keyword adds a constant to a specified array or part of an
array based on cells with a specific region number. The region number
array can be FLUXNUM, MULTNUM or OPERNUM and these arrays must be
defined and be available before the ADDREG keyword is read by the
simulator. The constant can be an integer or real value depending on the
array type; however, the arrays that can be operated on are dependent on
which section the ADDREG keyword is being applied in.

\begin{longtable}[]{@{}
  >{\raggedright\arraybackslash}p{(\linewidth - 6\tabcolsep) * \real{0.2500}}
  >{\raggedright\arraybackslash}p{(\linewidth - 6\tabcolsep) * \real{0.2500}}
  >{\raggedright\arraybackslash}p{(\linewidth - 6\tabcolsep) * \real{0.2500}}
  >{\raggedright\arraybackslash}p{(\linewidth - 6\tabcolsep) * \real{0.2500}}@{}}
\toprule\noalign{}
\begin{minipage}[b]{\linewidth}\raggedright
No.
\end{minipage} & \begin{minipage}[b]{\linewidth}\raggedright
Name
\end{minipage} & \begin{minipage}[b]{\linewidth}\raggedright
Description
\end{minipage} & \begin{minipage}[b]{\linewidth}\raggedright
Default
\end{minipage} \\
\midrule\noalign{}
\endhead
\bottomrule\noalign{}
\endlastfoot
1 & ARRAY & A character string of up to eight characters in length that
defines the keyword identifying the array to be modified. & None \\
2 & CONSTANT & An integer or real value to be added to the ARRAY in the
same units as the ARRAY property for a given REGION. & 0 \\
3 & REGION NUMBER & REGION NUMBER is a positive integer representing the
region for which the CONSTANT in (2) should be applied. & None \\
4 & REGION ARRAY & The REGION ARRAY to use for applying the CONSTANT in
(2) based on the REGION NUMBER in (3). REGION ARRAY can have the
following values: & M \\
Notes: & & & \\
\end{longtable}

\emph{Table 6.5: ADDREG Keyword Description}

Examples of the arrays most commonly operated on in each section are
given in Table 6.6. Cells colored red indicate arrays that are not
supported by OPM Flow operations.

\begin{longtable}[]{@{}lllllll@{}}
\toprule\noalign{}
ADDREG Keyword and Variable Options by Section & & & & & & \\
\midrule\noalign{}
\endhead
\bottomrule\noalign{}
\endlastfoot
GRID & EDIT & PROPS & REGIONS & SOLUTION & SUMMARY & SCHEDULE \\
DX & DEPTH & SWL & ENDNUM & PRESSURE & & \\
DY & PORV & SWCR & EQLNUM & SWAT & & \\
DZ & TRANX & SWU & FIPNUM & SGAS & & \\
PERMX & TRANY & SGL & IMBNUM & RV & & \\
PERMY & TRANZ & SGCR & MISCNUM & RS & & \\
PERMZ & DIFFX & SGU & PVTNUM & TBLK & & \\
MULTX & DIFFY & KRW & ROCKNUM & GI & & \\
MULTY & DIFFZ & KRO & SATNUM & OILAPI & & \\
MULTZ & TRANR & KRG & WH2NUM & SALT & & \\
DR & TRANTHT & PCG & & GASCONC & & \\
DTHETA & DIFFR & PCW & & SOLVCONC & & \\
PERMR & DIFFTHT & & & SOLVFRAC & & \\
PERMTHT & & & & SFOAM & & \\
DZNET & & & & SPOLY & & \\
PORO & & & & & & \\
NTG & & & & & & \\
FLUXNUM & & & & & & \\
MULTNUM & & & & & & \\
MPFANUM & & & & & & \\
DIFFX & & & & & & \\
DIFFY & & & & & & \\
DIFFZ & & & & & & \\
DIFFR & & & & & & \\
DIFFTHT & & & & & & \\
\end{longtable}

\emph{Table 6.6: ADDREG Keyword Applicable Arrays by Section}

\subsubsection{Example}\label{example-132}

\begin{verbatim}
--
-- FIRST DEFINE THE PROPERTY ARRAYS AND MULTNUM ARRAYS FOR 10 X 10 X 20 MODEL
--
--  ARRAY    CONSTANT --  ---------- BOX ---------
--                          I1  I2   J1  J2   K1  K2
EQUALS
    PORO      0.2000        1*  1*   1*  1*   1*  1*  / PORO  TO 0.20 IN MODEL
    PERMX     100.00        1*  1*   1*  1*   1*  1*  / PERMX TO 0.10 IN MODEL
    MULTNUM   1             1*  1*   1*  1*   1*  1*  / MULTNUM IN MODEL
    MULTNUM   2             1*  5    1   5    6   6   / MULTNUM IN MODEL
    MULTNUM   3             1*  1*   1*  1*   10  10  / MULTNUM IN MODEL
/
--
–        NOW RESET PORO AND PERMX BASED ON THE MULTNUM REGION NUMBER
--
--       ADD A CONSTANT TO AN ARRAY BASED ON A REGION NUMBER
--
--       ARRAY     CONSTANT  REGION   REGION ARRAY
--                 VALUE     NUMBER    M / F / O                                           ADDREG
      PORO         0.050     1         M                /
      PORO         0.100     2         M                /
      PORO        -0.050     3         M                /
      PERMX        25.00     1         M                /
      PERMX        100.0     2         M                /
      PERMX       -50.00     3         M                /
/
\end{verbatim}

The example first defines the PORO and PERMX property arrays for the
model and then sets the MULTNUM array to 1 for all cells in the model,
after which selected areas of model are assigned various MULTNUM integer
values. The ADDREG can then be invoked to add or subtract constant
values from the PORO and PERMX arrays for the various MULTNUM regions.

\subsection{ADDZCORN -- Add a Constant to the ZCORN Depth
Array}\label{addzcorn-add-a-constant-to-the-zcorn-depth-array}

\begin{longtable}[]{@{}
  >{\raggedright\arraybackslash}p{(\linewidth - 14\tabcolsep) * \real{0.1250}}
  >{\raggedright\arraybackslash}p{(\linewidth - 14\tabcolsep) * \real{0.1250}}
  >{\raggedright\arraybackslash}p{(\linewidth - 14\tabcolsep) * \real{0.1250}}
  >{\raggedright\arraybackslash}p{(\linewidth - 14\tabcolsep) * \real{0.1250}}
  >{\raggedright\arraybackslash}p{(\linewidth - 14\tabcolsep) * \real{0.1250}}
  >{\raggedright\arraybackslash}p{(\linewidth - 14\tabcolsep) * \real{0.1250}}
  >{\raggedright\arraybackslash}p{(\linewidth - 14\tabcolsep) * \real{0.1250}}
  >{\raggedright\arraybackslash}p{(\linewidth - 14\tabcolsep) * \real{0.1250}}@{}}
\toprule\noalign{}
\begin{minipage}[b]{\linewidth}\raggedright
RUNSPEC
\end{minipage} & \begin{minipage}[b]{\linewidth}\raggedright
GRID
\end{minipage} & \begin{minipage}[b]{\linewidth}\raggedright
EDIT
\end{minipage} & \begin{minipage}[b]{\linewidth}\raggedright
PROPS
\end{minipage} & \begin{minipage}[b]{\linewidth}\raggedright
REGIONS
\end{minipage} & \begin{minipage}[b]{\linewidth}\raggedright
SOLUTION
\end{minipage} & \begin{minipage}[b]{\linewidth}\raggedright
SUMMARY
\end{minipage} & \begin{minipage}[b]{\linewidth}\raggedright
SCHEDULE
\end{minipage} \\
\midrule\noalign{}
\endhead
\bottomrule\noalign{}
\endlastfoot
\end{longtable}

\subsubsection{Description}\label{description-172}

The ADDZCORN keyword adds a constant to the ZCORN array or part of the
array based on cells defined in the specified input box. The constant
can be real or integer and can be negative or positive.

This keyword is not supported by OPM Flow but would change the results
if supported so the simulation will be stopped.

\subsection{AMALGAM -- Define LGR
Amalgamations}\label{amalgam-define-lgr-amalgamations}

\begin{longtable}[]{@{}
  >{\raggedright\arraybackslash}p{(\linewidth - 14\tabcolsep) * \real{0.1250}}
  >{\raggedright\arraybackslash}p{(\linewidth - 14\tabcolsep) * \real{0.1250}}
  >{\raggedright\arraybackslash}p{(\linewidth - 14\tabcolsep) * \real{0.1250}}
  >{\raggedright\arraybackslash}p{(\linewidth - 14\tabcolsep) * \real{0.1250}}
  >{\raggedright\arraybackslash}p{(\linewidth - 14\tabcolsep) * \real{0.1250}}
  >{\raggedright\arraybackslash}p{(\linewidth - 14\tabcolsep) * \real{0.1250}}
  >{\raggedright\arraybackslash}p{(\linewidth - 14\tabcolsep) * \real{0.1250}}
  >{\raggedright\arraybackslash}p{(\linewidth - 14\tabcolsep) * \real{0.1250}}@{}}
\toprule\noalign{}
\begin{minipage}[b]{\linewidth}\raggedright
RUNSPEC
\end{minipage} & \begin{minipage}[b]{\linewidth}\raggedright
GRID
\end{minipage} & \begin{minipage}[b]{\linewidth}\raggedright
EDIT
\end{minipage} & \begin{minipage}[b]{\linewidth}\raggedright
PROPS
\end{minipage} & \begin{minipage}[b]{\linewidth}\raggedright
REGIONS
\end{minipage} & \begin{minipage}[b]{\linewidth}\raggedright
SOLUTION
\end{minipage} & \begin{minipage}[b]{\linewidth}\raggedright
SUMMARY
\end{minipage} & \begin{minipage}[b]{\linewidth}\raggedright
SCHEDULE
\end{minipage} \\
\midrule\noalign{}
\endhead
\bottomrule\noalign{}
\endlastfoot
\end{longtable}

\subsubsection{Description}\label{description-173}

The AMALGAM keyword defines a Cartesian Local Grid Refinements (``LGR'')
amalgamations, that is merging several LGRs into one amalgamated LGR.

This keyword is not supported by OPM Flow but would change the results
if supported so the simulation will be stopped.

\subsection{AQUANCON -- Define Analytical Connections to the
Grid}\label{aquancon-define-analytical-connections-to-the-grid}

\begin{longtable}[]{@{}
  >{\raggedright\arraybackslash}p{(\linewidth - 14\tabcolsep) * \real{0.1250}}
  >{\raggedright\arraybackslash}p{(\linewidth - 14\tabcolsep) * \real{0.1250}}
  >{\raggedright\arraybackslash}p{(\linewidth - 14\tabcolsep) * \real{0.1250}}
  >{\raggedright\arraybackslash}p{(\linewidth - 14\tabcolsep) * \real{0.1250}}
  >{\raggedright\arraybackslash}p{(\linewidth - 14\tabcolsep) * \real{0.1250}}
  >{\raggedright\arraybackslash}p{(\linewidth - 14\tabcolsep) * \real{0.1250}}
  >{\raggedright\arraybackslash}p{(\linewidth - 14\tabcolsep) * \real{0.1250}}
  >{\raggedright\arraybackslash}p{(\linewidth - 14\tabcolsep) * \real{0.1250}}@{}}
\toprule\noalign{}
\begin{minipage}[b]{\linewidth}\raggedright
RUNSPEC
\end{minipage} & \begin{minipage}[b]{\linewidth}\raggedright
GRID
\end{minipage} & \begin{minipage}[b]{\linewidth}\raggedright
EDIT
\end{minipage} & \begin{minipage}[b]{\linewidth}\raggedright
PROPS
\end{minipage} & \begin{minipage}[b]{\linewidth}\raggedright
REGIONS
\end{minipage} & \begin{minipage}[b]{\linewidth}\raggedright
SOLUTION
\end{minipage} & \begin{minipage}[b]{\linewidth}\raggedright
SUMMARY
\end{minipage} & \begin{minipage}[b]{\linewidth}\raggedright
SCHEDULE
\end{minipage} \\
\midrule\noalign{}
\endhead
\bottomrule\noalign{}
\endlastfoot
\end{longtable}

\subsubsection{Description}\label{description-174}

The AQUANCON keyword defines how analytical aquifers are connected to
the simulation grid, this includes the Carter-Tracy, Fetkovich and
Constant Flux analytical aquifers, all of which are implemented in OPM
Flow. Carter-Tracy analytical aquifers are characterized by the AQUCT
keyword in the GRID section and Fetkovich analytical aquifers are
defined by the AQUFETP keyword in the SOLUTION section. Finally, the
Constant Flux aquifer is defined by the AQUFLUX keyword in SOLUTION
section.

Note that numerical aquifers are connected to the grid using the AQUCON
keyword in the GRID section and that both aquifer types dimensions are
declared by the AQUDIMS keyword in the RUNSPEC section.

\begin{longtable}[]{@{}
  >{\raggedright\arraybackslash}p{(\linewidth - 6\tabcolsep) * \real{0.2500}}
  >{\raggedright\arraybackslash}p{(\linewidth - 6\tabcolsep) * \real{0.2500}}
  >{\raggedright\arraybackslash}p{(\linewidth - 6\tabcolsep) * \real{0.2500}}
  >{\raggedright\arraybackslash}p{(\linewidth - 6\tabcolsep) * \real{0.2500}}@{}}
\toprule\noalign{}
\begin{minipage}[b]{\linewidth}\raggedright
No.
\end{minipage} & \begin{minipage}[b]{\linewidth}\raggedright
Name
\end{minipage} & \begin{minipage}[b]{\linewidth}\raggedright
Description
\end{minipage} & \begin{minipage}[b]{\linewidth}\raggedright
Default
\end{minipage} \\
\midrule\noalign{}
\endhead
\bottomrule\noalign{}
\endlastfoot
Field & Metric & Laboratory & \\
1 & AQUID & AQUID is a positive integer greater than or equal to one and
less than the maximum number of analytical aquifers as defined by the
NANAQU variable on the AQUDIMS keyword in the RUNSPEC section, that
defines the aquifer to be connected to the grid. & None \\
2 & I1 & A positive integer that defines the lower bound of the cells in
the I-direction to be connected to the aquifer and must be greater than
or equal to one and less than or equal to I2 and NX. & 1 \\
3 & I2 & A positive integer that defines the upper bound of the cells in
the I-direction to be connected to the aquifer and must be greater than
or equal to I1 and less than or equal to NX & NX \\
4 & J1 & A positive integer that defines the lower bound of the cells in
the J-direction to be connected to the aquifer and must be greater than
or equal to one and less than or equal to J2 and NY. & 1 \\
5 & J2 & A positive integer that defines the upper bound of the cells in
the J-direction to be connected to the aquifer and must be greater than
or equal to JI and less than or equal to NY. & NY \\
6 & K1 & A positive integer that defines the lower bound of the cells in
the K-direction to be to be connected to the aquifer and must be greater
than or equal to one and less than or equal to K2 and NZ. & 1 \\
7 & K2 & A positive integer that defines the upper bound of the cells in
the K-direction to be connected to the aquifer and must be greater than
or equal to KI and less than or equal to NZ. & NZ \\
8 & AQUFACE & AQUFACE is a character string that sets the connection
``face'' of the cells declared by this record and should be set to one
of the following: & None \\
9 & AQUFLUX & AQUFLUX is a positive real value that sets the fraction of
the total influx between the aquifer and the defined cells declared on
this keyword. If defaulted the cell face area for each cell is applied
and if a value is declared then this value is applied to all cells
declared by this record. & 1* \\
ft2 & m2 & cm2 & \\
10 & AQUCOEF & AQUCOEF is a real positive values that scales the
calculated connection between the aquifer and the cells declared on this
record. & 1.0 \\
dimensionless & dimensionless & dimensionless & \\
11 & AQUOPT & AQUOPT is a character string that sets the cell face
connection and should be set to one of the following: & NO \\
Notes: & & & \\
\end{longtable}

\emph{Table 6.7: AQUANCON Keyword Description}

\begin{tcolorbox}[enhanced jigsaw, title=\textcolor{quarto-callout-note-color}{\faInfo}\hspace{0.5em}{Note}, coltitle=black, breakable, toprule=.15mm, colbacktitle=quarto-callout-note-color!10!white, leftrule=.75mm, opacitybacktitle=0.6, colback=white, bottomtitle=1mm, toptitle=1mm, titlerule=0mm, left=2mm, colframe=quarto-callout-note-color-frame, arc=.35mm, rightrule=.15mm, bottomrule=.15mm, opacityback=0]

If the AQUANCON keyword has been utilized in the run deck then OPM Flow
will write the AQUIFERA array to the *.INIT file in order to visualize
the aquifer connections in OPM ResInsight. This is accomplished by
setting the AQUIFERA value to 2\^{}(AQUID-1) for cells connected to
aquifer AQUID. If a cell is connected to multiple analytical aquifers
then AQUIFERA is summed for all aquifers connected to a cell. Note that
connecting cells to multiple aquifers is best avoided.

\end{tcolorbox}

\subsubsection{Example}\label{example-133}

The following example defines aquifer number one connected to the I+
face of various cells in the model.

\begin{verbatim}
--
--                      ANALYTIC AQUIFER CONNECTION
--
--       ID     ---------- BOX ---------   CONNECT  AQF    AQF      ADJOIN
--       NUMBER I1  I2   J1  J2   K1  K2   FACE     INFLX  MULTI    CELLS
AQUANCON
         1     57  57   28  36   46  58    'I+'    1*     1*       'NO'        /
         1    111 111   38  41   22  31    'I+'    1*     1*       'NO'        /
         1     96  96   44  49   22  31    'I+'    1*     1*       'NO'        /
         1     43  43   28  35   54  58    'I+'    1*     1*       'NO'        /
         1     98  98   38  42   32  40    'I+'    1*     1*       'NO'        /
         1     79  79   41  67    5  11    'I+'    1*     1*       'NO'        /
         1     61  61   48  72   12  17    'I+'    1*     1*       'NO'        /
/
\end{verbatim}

See the AQUCT keyword in the GRID section for a complete example on
defining and connecting a Carter-Tracy aquifer to a simulation grid.

\subsection{AQUCON -- Define Numerical Aquifer Connections to the
Grid}\label{aqucon-define-numerical-aquifer-connections-to-the-grid}

\begin{longtable}[]{@{}
  >{\raggedright\arraybackslash}p{(\linewidth - 14\tabcolsep) * \real{0.1250}}
  >{\raggedright\arraybackslash}p{(\linewidth - 14\tabcolsep) * \real{0.1250}}
  >{\raggedright\arraybackslash}p{(\linewidth - 14\tabcolsep) * \real{0.1250}}
  >{\raggedright\arraybackslash}p{(\linewidth - 14\tabcolsep) * \real{0.1250}}
  >{\raggedright\arraybackslash}p{(\linewidth - 14\tabcolsep) * \real{0.1250}}
  >{\raggedright\arraybackslash}p{(\linewidth - 14\tabcolsep) * \real{0.1250}}
  >{\raggedright\arraybackslash}p{(\linewidth - 14\tabcolsep) * \real{0.1250}}
  >{\raggedright\arraybackslash}p{(\linewidth - 14\tabcolsep) * \real{0.1250}}@{}}
\toprule\noalign{}
\begin{minipage}[b]{\linewidth}\raggedright
RUNSPEC
\end{minipage} & \begin{minipage}[b]{\linewidth}\raggedright
GRID
\end{minipage} & \begin{minipage}[b]{\linewidth}\raggedright
EDIT
\end{minipage} & \begin{minipage}[b]{\linewidth}\raggedright
PROPS
\end{minipage} & \begin{minipage}[b]{\linewidth}\raggedright
REGIONS
\end{minipage} & \begin{minipage}[b]{\linewidth}\raggedright
SOLUTION
\end{minipage} & \begin{minipage}[b]{\linewidth}\raggedright
SUMMARY
\end{minipage} & \begin{minipage}[b]{\linewidth}\raggedright
SCHEDULE
\end{minipage} \\
\midrule\noalign{}
\endhead
\bottomrule\noalign{}
\endlastfoot
\end{longtable}

\subsubsection{Description}\label{description-175}

AQUCON keyword defines how numerical aquifers are connected to the
simulation grid and these type of aquifers are characterized by the
AQUNUM keyword in the GRID section. Analytical aquifers are connected to
the simulation grid by the AQUANCON keyword in the GRID section, this
includes the Carter-Tracy and Fetkovich analytical aquifers, both of
which are implemented in OPM Flow. Both aquifer types dimensions are
declared by the AQUDIMS keyword in the RUNSPEC section.

\begin{longtable}[]{@{}
  >{\raggedright\arraybackslash}p{(\linewidth - 6\tabcolsep) * \real{0.2500}}
  >{\raggedright\arraybackslash}p{(\linewidth - 6\tabcolsep) * \real{0.2500}}
  >{\raggedright\arraybackslash}p{(\linewidth - 6\tabcolsep) * \real{0.2500}}
  >{\raggedright\arraybackslash}p{(\linewidth - 6\tabcolsep) * \real{0.2500}}@{}}
\toprule\noalign{}
\begin{minipage}[b]{\linewidth}\raggedright
No.
\end{minipage} & \begin{minipage}[b]{\linewidth}\raggedright
Name
\end{minipage} & \begin{minipage}[b]{\linewidth}\raggedright
Description
\end{minipage} & \begin{minipage}[b]{\linewidth}\raggedright
Default
\end{minipage} \\
\midrule\noalign{}
\endhead
\bottomrule\noalign{}
\endlastfoot
Field & Metric & Laboratory & \\
1 & AQUID & AQUID is a positive integer greater than or equal to one and
less than or equal to the maximum number of numerical aquifers as
defined by the MXNAQN variable on the AQUDIMS keyword in the RUNSPEC
section, that defines the aquifer to be connected to the grid. & None \\
2 & I1 & A positive integer that defines the lower bound of the cells in
the I-direction to be connected to the aquifer and must be greater than
or equal to one and less than or equal to I2 and NX. & 1 \\
3 & I2 & A positive integer that defines the upper bound of the cells in
the I-direction to be connected to the aquifer and must be greater than
or equal to I1 and less than or equal to NX & NX \\
4 & J1 & A positive integer that defines the lower bound of the cells in
the J-direction to be connected to the aquifer and must be greater than
or equal to one and less than or equal to J2 and NY. & 1 \\
5 & J2 & A positive integer that defines the upper bound of the cells in
the J-direction to be connected to the aquifer and must be greater than
or equal to JI and less than or equal to NY. & NY \\
6 & K1 & A positive integer that defines the lower bound of the cells in
the K-direction to be to be connected to the aquifer and must be greater
than or equal to one and less than or equal to K2 and NZ. & 1 \\
7 & K2 & A positive integer that defines the upper bound of the cells in
the K-direction to be connected to the aquifer and must be greater than
or equal to KI and less than or equal to NZ. & NZ \\
8 & AQUFACE & AQUFACE is a character string that sets the connection
``face'' of the cells declared by this record and should be set to one
of the following: & None \\
9 & AQUMULT & AQUMULT is a positive real number greater than or equal to
zero that scales the OPM Flow calculated transmissibility between the
AQUID aquifer and the grid cell connections defined by this record. The
default value of one sets the transmissibility between the aquifer and
grid cells to the OPM Flow calculated value. & 1.0 \\
dimensionless & dimensionless & dimensionless & \\
10 & AQUOPT1 & AQUOPT1 is a defined integer value set to either zero or
one, that defines the area to be used in calculating the connection
transmissibility between the aquifer and the grid cells: & 0 \\
dimensionless & dimensionless & dimensionless & \\
11 & AQUOPT2 & AQUOPT2 is a character string that sets the cell face
connection and should be set to one of the following: & NO \\
12 & VEOPT1 & Vertical Equilibrium Option Number 1-- Not Used & 1 \\
13 & VEOPT2 & Vertical Equilibrium Option Number 2-- Not Used & 1 \\
Notes: & & & \\
\end{longtable}

\emph{Table 6.8: AQUCON Keyword Description}

\begin{tcolorbox}[enhanced jigsaw, title=\textcolor{quarto-callout-note-color}{\faInfo}\hspace{0.5em}{Note}, coltitle=black, breakable, toprule=.15mm, colbacktitle=quarto-callout-note-color!10!white, leftrule=.75mm, opacitybacktitle=0.6, colback=white, bottomtitle=1mm, toptitle=1mm, titlerule=0mm, left=2mm, colframe=quarto-callout-note-color-frame, arc=.35mm, rightrule=.15mm, bottomrule=.15mm, opacityback=0]

If the AQUCON keyword has been utilized in the run deck then OPM Flow
will write the AQUIFERN array to the *.INIT file in order to visualize
the aquifer connections in OPM ResInsight. This is accomplished by
setting the AQUIFERN value to 2\^{}(AQUID-1) for cells connected to
aquifer AQUID. If a cell is connected to multiple numerical aquifers
then AQUIFERN is summed for all aquifers connected to a cell. Note that
connecting cells to multiple aquifers is best avoided. Finally, for
cells representing the numerical aquifers themselves, AQUIFERN is set to
minus AQUID.

\end{tcolorbox}

\subsubsection{Example}\label{example-134}

The following example defines numerical aquifer number one connected to
the I+ face of various cells in the model.

\begin{verbatim}
--
--        NUMERICAL AQUIFER CONNECTIONS
--
--       ID     ---------- BOX ---------   CONNECT  TRANS  TRANS  ADJOIN
--       NUMBER I1  I2   J1  J2   K1  K2   FACE     MULT   OPTN   CELLS
AQUCON
         1     57  57   28  36   46  58    'I+'     1*     1*     'NO'        /
         1    111 111   38  41   22  31    'I+'     1*     1*     'NO'        /
         1     96  96   44  49   22  31    'I+'     1*     1*     'NO'        /
         1     43  43   28  35   54  58    'I+'     1*     1*     'NO'        /
         1     98  98   38  42   32  40    'I+'     1*     1*     'NO'        /
         1     79  79   41  67    5  11    'I+'     1*     1*     'NO'        /
         1     61  61   48  72   12  17    'I+'     1*     1*     'NO'        /
/
\end{verbatim}

See the AQUNUM keyword in the GRID section for a complete example on
defining and connecting a numerical aquifer to a simulation grid.

\subsection{AQUCT -- Define Carter-Tracy Analytical
Aquifers}\label{aquct-define-carter-tracy-analytical-aquifers}

\begin{longtable}[]{@{}
  >{\raggedright\arraybackslash}p{(\linewidth - 14\tabcolsep) * \real{0.1250}}
  >{\raggedright\arraybackslash}p{(\linewidth - 14\tabcolsep) * \real{0.1250}}
  >{\raggedright\arraybackslash}p{(\linewidth - 14\tabcolsep) * \real{0.1250}}
  >{\raggedright\arraybackslash}p{(\linewidth - 14\tabcolsep) * \real{0.1250}}
  >{\raggedright\arraybackslash}p{(\linewidth - 14\tabcolsep) * \real{0.1250}}
  >{\raggedright\arraybackslash}p{(\linewidth - 14\tabcolsep) * \real{0.1250}}
  >{\raggedright\arraybackslash}p{(\linewidth - 14\tabcolsep) * \real{0.1250}}
  >{\raggedright\arraybackslash}p{(\linewidth - 14\tabcolsep) * \real{0.1250}}@{}}
\toprule\noalign{}
\begin{minipage}[b]{\linewidth}\raggedright
RUNSPEC
\end{minipage} & \begin{minipage}[b]{\linewidth}\raggedright
GRID
\end{minipage} & \begin{minipage}[b]{\linewidth}\raggedright
EDIT
\end{minipage} & \begin{minipage}[b]{\linewidth}\raggedright
PROPS
\end{minipage} & \begin{minipage}[b]{\linewidth}\raggedright
REGIONS
\end{minipage} & \begin{minipage}[b]{\linewidth}\raggedright
SOLUTION
\end{minipage} & \begin{minipage}[b]{\linewidth}\raggedright
SUMMARY
\end{minipage} & \begin{minipage}[b]{\linewidth}\raggedright
SCHEDULE
\end{minipage} \\
\midrule\noalign{}
\endhead
\bottomrule\noalign{}
\endlastfoot
\end{longtable}

\subsubsection{Description}\label{description-176}

The AQUCT keyword defines Carter-Tracy {[}Carter, R. D. and Tracy, G. W.
``An Improved Method for Calculating Water Influx,'' Transactions of
AIME (1960) 219, 215-417.{]} {[}Van Everdingen, A. \& Hurst, W.,.The
Application of the Laplace Transformation to Flow Problems in
Reservoirs. Petroleum Transactions, AIME (December, 1949).{]} analytical
aquifers, the properties of the aquifer, including the Carter-Tracy
aquifer influence function associated with the aquifer, defined by the
AQUTAB keyword in the PROPS section.

Each row entry in the AQUCT keyword defines one Carter-Tracy aquifer.

\begin{longtable}[]{@{}
  >{\raggedright\arraybackslash}p{(\linewidth - 6\tabcolsep) * \real{0.2500}}
  >{\raggedright\arraybackslash}p{(\linewidth - 6\tabcolsep) * \real{0.2500}}
  >{\raggedright\arraybackslash}p{(\linewidth - 6\tabcolsep) * \real{0.2500}}
  >{\raggedright\arraybackslash}p{(\linewidth - 6\tabcolsep) * \real{0.2500}}@{}}
\toprule\noalign{}
\begin{minipage}[b]{\linewidth}\raggedright
No.
\end{minipage} & \begin{minipage}[b]{\linewidth}\raggedright
Name
\end{minipage} & \begin{minipage}[b]{\linewidth}\raggedright
Description
\end{minipage} & \begin{minipage}[b]{\linewidth}\raggedright
Default
\end{minipage} \\
\midrule\noalign{}
\endhead
\bottomrule\noalign{}
\endlastfoot
Field & Metric & Laboratory & \\
1 & AQUID & A positive integer greater than or equal to one and less
than or equal to NANAQU on the AQUDIMS keyword in the RUNSPEC section,
that defines the Carter-Tracy aquifer number. & 1 \\
2 & DATUM & DATUM is a single positive value that defines the
Carter-Tracy reference datum depth for PRESS. & None \\
feet & m & cm & \\
3 & PRESS & PRESS is a single positive value that defines the aquifer
pressure at DATUM. If PRESS is defaulted then the simulator will set the
aquifer's initial reservoir pressure to be in equilibrium with the cells
the aquifer is contacted to. & 1* \\
psia & barsa & atma & \\
4 & PERM & PERM is a real positive number that assigns the permeability
to the aquifer. & None \\
mD & mD & mD & \\
5 & PORO & PORO is a real positive number greater than zero and less
than or equal to one that assigns the porosity to the aquifer. & None \\
dimensionless & dimensionless & dimensionless & \\
6 & RCOMP & RCOMP is a real number defining the total (rock and water)
compressibility (Ct) at the DATUM pressure. & None \\
1/psia & 1/barsa & 1/atma & \\
7 & RE & RE is a real positive number that defines the Carter-Tracy
aquifer external radius. & None \\
feet & m & cm & \\
8 & DZ & DZ is a real positive number that defines the Carter-Tracy
aquifer average net thickness. & None \\
feet & m & cm & \\
9 & ANGLE & ANGLE is a real positive number that defines the angle of
influence, that is the angular connection between the aquifer and the
hydrocarbon reservoir. A value of 360o degrees, the default value,
indicates that the aquifer completely surrounds the hydrocarbon
reservoir. & 360.0 \\
degrees & degrees & degrees & \\
10 & PVTNUM & PVTNUM is positive integer greater than zero and less than
the NTPVT variable on the TABDIMS keyword in the RUNSPEC section, that
defines the PVTW table allocated to the Carter-Tracy aquifer. & 1 \\
11 & AQUTAB & AQUTAB is positive integer greater than zero and less than
the NIFTBL variable as declared on the AQUDIMS keyword in the RUNSPEC
section, that defines the AQUTAB table allocated to this Carter-Tracy
aquifer. The default value of one sets the internal infinite acting
Carter-Tracy aquifer influence table not the first table in the AQUTAB
keyword in the PROPS section The first table in the AQUTAB keyword is
considered to be table number two. The internal table, AQUTAB equal to
one, is based on the Radial Flow, Constant Pressure and Constant
Terminal Rate Cases for Infinite Reservoirs (Table I) in Van Everdingen
and Hurst's {[}Van Everdingen, A. \& Hurst, W.,.The Application of the
Laplace Transformation to Flow Problems in Reservoirs. Petroleum
Transactions, AIME (December, 1949).{]} paper. & 1 \\
12 & SALTCON & SALTCON is a real positive number that defines the
initial salt concentration in the aquifer, for when with simulator's
Brine Model has been activated via the BRINE keyword in the RUNSPEC
section. This variable is ignored by OPM Flow. & 0.0 \\
lb/stb & kg/sm3 & gm/scc & \\
13 & TEMP & TEMP is a real positive number that defines the initial
temperature of the aquifer at DATUM for use with OPM Flow's thermal
option. The THERMAL keyword in the RUNSPEC section must be activated to
use this option. & 1* \\
oF & oC & oC & \\
Notes: & & & \\
\end{longtable}

\emph{Table 6.9: AQUCT Keyword Description}

\begin{tcolorbox}[enhanced jigsaw, title=\textcolor{quarto-callout-note-color}{\faInfo}\hspace{0.5em}{Note}, coltitle=black, breakable, toprule=.15mm, colbacktitle=quarto-callout-note-color!10!white, leftrule=.75mm, opacitybacktitle=0.6, colback=white, bottomtitle=1mm, toptitle=1mm, titlerule=0mm, left=2mm, colframe=quarto-callout-note-color-frame, arc=.35mm, rightrule=.15mm, bottomrule=.15mm, opacityback=0]

OPM Flow includes the infinite acting Carter-Tracy aquifer influence
table as a default for table number one; thus data entered on AQUTAB
keyword starts from table number two.

\end{tcolorbox}

In order to full define a Carter-Tracy aquifer one has to define the
aquifer properties via the AQUCT keyword, the Carter-Tracy influence
function via the AQUTAB keyword in the PROPS section (AQUTAB -- Define
Carter-Tracy Aquifer Influence Functions), if the default infinite
acting table is not being employed, and finally, how the aquifer is
connected to the reservoir using the AQUANCON keyword in the GRID or
SOLUTION sections.

\subsubsection{Example}\label{example-135}

Given the following grid and aquifer dimensions in the RUNSPEC section:

\begin{verbatim}
--
--       MAX     MAX     MAX
--       NDIVIX  NDIVIY  NDIVIZ
DIMENS
         20      1       5                                                    /

--       AQF     AQF     AQF     AQF     AQF     AQF    AQF    AQF
--       MXAQN   MXNAQC  NIFTBL  NRIFTB  NANAQU  NCAMAX MXNALI MXAAQL
AQUDIMS
         1*      1*      5       100     1       1*     1*     1*              /
\end{verbatim}

And the AQUTAB keyword in the PROPS section

\begin{verbatim}
--
--       CARTER-TRACY AQUIFER INFLUENCE TABLES
--       (STARTS FROM TABLE NO. 2, AS DEFAULT IS TABLE NO. 1)
--
AQUTAB
--       DIMLESS    DIMLESS
--       TIME       PRESSURE
--       -------    ---------
          0.01        0.112
          0.05        0.229
          0.10        0.315
          0.15        0.376
          0.20        0.424
          0.22        0.443
          0.24        0.459
          0.26        0.476
          0.28        0.492
          0.30        0.507
          0.32        0.522
          0.34        0.536
          0.36        0.551
          0.38        0.565
          0.40        0.579
          0.42        0.593
          0.44        0.607
          0.46        0.621
          0.48        0.634
          0.50        0.648
          0.60        0.715
          0.70        0.782
          0.80        0.849
          0.90        0.915
          1.00        0.982
          2.00        1.649
          3.00        2.316
          5.00        3.649
         10.00        6.982
         20.00       13.649
         30.00       20.316
         50.00       33.649
        100.00       66.982
        200.00      133.649
        300.00      200.316
        500.00      333.649
       1000.00      666.982 /
\end{verbatim}

The Carter-Tracy aquifer is defined in the GRID or SOLUTION sections as:

\begin{verbatim}
-- ==============================================================================
--
-- SOLUTION SECTION
--
-- ==============================================================================
SOLUTION
--
--                      CARTER-TRACY AQUIFER DESCRIPTION
--
--      ID   DATUM   AQF    AQF    AQF    AQF     AQF   AQF   INFL   PVT  AQU
--      NUM  DEPTH   PRESS  PERM   PORO   RCOMP   RE    DZ    ANGLE  NUM  TAB
--
AQUCT
         1   2000.0  269    100.0  0.30   3.0e-5  330   10.0  360.0   1    2   /
/
\end{verbatim}

And the connection of the aquifer is set in the GRID or SOLUTION
sections as:

\begin{verbatim}
--
--                      ANALYTIC AQUIFER CONNECTION
--
--       ID     ---------- BOX ---------   CONNECT  AQF    AQF      ADJOIN
--       NUMBER I1  I2   J1  J2   K1  K2   FACE     INFLX  MULTI    CELLS
AQUANCON
         1      1   1    1   1    1   1     J-      1.0    1.0      'NO'        /
/
\end{verbatim}

Here one Carter-Tracy aquifer is connected to a single cell (1, 1, 1) at
the J- face (or X- face) of the cell.

\subsection{AQUNNC -- Define Numerical Aquifer Non-Neighbor
Connections}\label{aqunnc-define-numerical-aquifer-non-neighbor-connections}

\begin{longtable}[]{@{}
  >{\raggedright\arraybackslash}p{(\linewidth - 14\tabcolsep) * \real{0.1250}}
  >{\raggedright\arraybackslash}p{(\linewidth - 14\tabcolsep) * \real{0.1250}}
  >{\raggedright\arraybackslash}p{(\linewidth - 14\tabcolsep) * \real{0.1250}}
  >{\raggedright\arraybackslash}p{(\linewidth - 14\tabcolsep) * \real{0.1250}}
  >{\raggedright\arraybackslash}p{(\linewidth - 14\tabcolsep) * \real{0.1250}}
  >{\raggedright\arraybackslash}p{(\linewidth - 14\tabcolsep) * \real{0.1250}}
  >{\raggedright\arraybackslash}p{(\linewidth - 14\tabcolsep) * \real{0.1250}}
  >{\raggedright\arraybackslash}p{(\linewidth - 14\tabcolsep) * \real{0.1250}}@{}}
\toprule\noalign{}
\begin{minipage}[b]{\linewidth}\raggedright
RUNSPEC
\end{minipage} & \begin{minipage}[b]{\linewidth}\raggedright
GRID
\end{minipage} & \begin{minipage}[b]{\linewidth}\raggedright
EDIT
\end{minipage} & \begin{minipage}[b]{\linewidth}\raggedright
PROPS
\end{minipage} & \begin{minipage}[b]{\linewidth}\raggedright
REGIONS
\end{minipage} & \begin{minipage}[b]{\linewidth}\raggedright
SOLUTION
\end{minipage} & \begin{minipage}[b]{\linewidth}\raggedright
SUMMARY
\end{minipage} & \begin{minipage}[b]{\linewidth}\raggedright
SCHEDULE
\end{minipage} \\
\midrule\noalign{}
\endhead
\bottomrule\noalign{}
\endlastfoot
\end{longtable}

\subsubsection{Description}\label{description-177}

The AQUNNC keyword defines Numerical Aquifer Non-Neighbor Connections.

This keyword is not supported by OPM Flow but would change the results
if supported so the simulation will be stopped.

\subsection{AQUNUM -- Define Numerical Aquifer
Properties}\label{aqunum-define-numerical-aquifer-properties}

\begin{longtable}[]{@{}
  >{\raggedright\arraybackslash}p{(\linewidth - 14\tabcolsep) * \real{0.1250}}
  >{\raggedright\arraybackslash}p{(\linewidth - 14\tabcolsep) * \real{0.1250}}
  >{\raggedright\arraybackslash}p{(\linewidth - 14\tabcolsep) * \real{0.1250}}
  >{\raggedright\arraybackslash}p{(\linewidth - 14\tabcolsep) * \real{0.1250}}
  >{\raggedright\arraybackslash}p{(\linewidth - 14\tabcolsep) * \real{0.1250}}
  >{\raggedright\arraybackslash}p{(\linewidth - 14\tabcolsep) * \real{0.1250}}
  >{\raggedright\arraybackslash}p{(\linewidth - 14\tabcolsep) * \real{0.1250}}
  >{\raggedright\arraybackslash}p{(\linewidth - 14\tabcolsep) * \real{0.1250}}@{}}
\toprule\noalign{}
\begin{minipage}[b]{\linewidth}\raggedright
RUNSPEC
\end{minipage} & \begin{minipage}[b]{\linewidth}\raggedright
GRID
\end{minipage} & \begin{minipage}[b]{\linewidth}\raggedright
EDIT
\end{minipage} & \begin{minipage}[b]{\linewidth}\raggedright
PROPS
\end{minipage} & \begin{minipage}[b]{\linewidth}\raggedright
REGIONS
\end{minipage} & \begin{minipage}[b]{\linewidth}\raggedright
SOLUTION
\end{minipage} & \begin{minipage}[b]{\linewidth}\raggedright
SUMMARY
\end{minipage} & \begin{minipage}[b]{\linewidth}\raggedright
SCHEDULE
\end{minipage} \\
\midrule\noalign{}
\endhead
\bottomrule\noalign{}
\endlastfoot
\end{longtable}

\subsubsection{Description}\label{description-178}

The AQUNUM keyword defines the properties of numerical aquifers,
including which grid blocks in the model should be utilized as part of
the numerical aquifer. Each row entry in the AQUNUM keyword defines one
numerical aquifer. Note that a numerical aquifer may consists of more
than one grid cell, in order to better describe the water influx from
the aquifer to the grid.

\begin{longtable}[]{@{}
  >{\raggedright\arraybackslash}p{(\linewidth - 6\tabcolsep) * \real{0.2500}}
  >{\raggedright\arraybackslash}p{(\linewidth - 6\tabcolsep) * \real{0.2500}}
  >{\raggedright\arraybackslash}p{(\linewidth - 6\tabcolsep) * \real{0.2500}}
  >{\raggedright\arraybackslash}p{(\linewidth - 6\tabcolsep) * \real{0.2500}}@{}}
\toprule\noalign{}
\begin{minipage}[b]{\linewidth}\raggedright
No.
\end{minipage} & \begin{minipage}[b]{\linewidth}\raggedright
Name
\end{minipage} & \begin{minipage}[b]{\linewidth}\raggedright
Description
\end{minipage} & \begin{minipage}[b]{\linewidth}\raggedright
Default
\end{minipage} \\
\midrule\noalign{}
\endhead
\bottomrule\noalign{}
\endlastfoot
Field & Metric & Laboratory & \\
1 & AQUID & AQUID is a positive integer greater than or equal to one and
less than or equal to the maximum number of numerical aquifers as
defined by the MXNAQN variable on the AQUDIMS keyword in the RUNSPEC
section, that defines the aquifer to be connected to the grid. & None \\
2 & I & A positive integer that defines the cell in the I-direction that
represents the AQUID aquifer, and which must be greater than or equal to
one and less than or equal to NX. & None \\
3 & J & A positive integer that defines the cell in the J-direction that
represents the AQUID aquifer, and which must be greater than or equal to
one and less than or equal to NY. & None \\
4 & K & A positive integer that defines the cell in the K-direction that
represents the AQUID aquifer, and which must be greater than or equal to
one and less than or equal to NZ. & None \\
5 & AREA & AREA is a real positive value that defines the
cross-sectional area of the aquifer used in calculating the aquifer
connection transmissibility. The actual transmissibility between the
numerical aquifer and the connected grid cell is the harmonic average of
the aquifer connection transmissibility and the calculated connected
cell transmissibility. The value entered for AREA does not effect the
visualization of the cell in OPM ResInsight, as it is only used in
calculating the transmissibility. Note that the AQUOPT1 variable on the
AQUCON keyword in the GRID section allows one to use the value entered
for AREA or to use the grid cell cross-sectional area instead for the
transmissibility calculation. & None \\
ft2 & m2 & cm2 & \\
6 & LENGTH & LENGTH is a real positive value that defines the length of
the numerical aquifer. Similar to the AREA variable, LENGTH is not
constrained by the original host cell size and the value entered does
not effect the visualization of the cell in OPM ResInsight. & None \\
feet & m & cm & \\
7 & PORO & PORO is a real positive number greater than zero and less
than or equal to one that assigns the porosity to the numerical aquifer.
& None \\
dimensionless & dimensionless & dimensionless & \\
8 & PERM & PERM is a real positive number that assigns the permeability
to the numerical aquifer. & None \\
mD & mD & mD & \\
9 & DATUM & DATUM is a real positive number that sets the reference
datum depth of the numerical aquifer. Similar to the AREA variable,
DATUM is not constrained by the original host cell depth and the value
entered does not effect the visualization of the cell in OPM ResInsight.
If defaulted then the depth of cell (I, J, K) will be used & Cell
Depth \\
feet & m & cm & \\
10 & PRESS & PRESS is a single positive value that defines the numerical
aquifer pressure at DATUM. If PRESS is defaulted then the simulator will
set the aquifer's initial reservoir pressure to be in equilibrium with
the cells the aquifer is contacted to. This is the preferred manner to
initialize the numerical aquifer. & 1* \\
psia & barsa & atma & \\
11 & PVTNUM & PVTNUM is positive integer greater than zero and less than
the NTPVT variable on the TABDIMS keyword in the RUNSPEC section, that
defines the PVTW table allocated to the numerical aquifer. If defaulted
then the PVT tables allocated to cell (I, J, K) will be used. & Cell
PVTNUM \\
12 & SATNUM & SATNUM is positive integer greater than zero and less than
the NTSFUN variable on the TABDIMS keyword in the RUNSPEC section, that
defines the saturation tables allocated to the numerical aquifer. If
defaulted then the saturation tables allocated to cell (I, J, K) will be
used. & Cell SATNUM \\
Notes: & & & \\
\end{longtable}

\emph{Table 6.10: AQUNUM Keyword Description}

Numerical aquifers are modeled as one-dimensional, with aquifer flow
assumed to be in the direction defined by LENGTH, and flux out of the
aquifer is through the cross sectional area defined by AREA. Note that,
all the aquifer cells must be isolated from the reservoir cells, with
only the AQUCON connections connecting to the actual reservoir cells.

The values entered on the AQUNUM keyword are used to calculate the
aquifer's pore volume and the transmissibility between the aquifer and
the connected cell faces defined on the AQUCON keyword. Thus:

\begin{verbatim}
- The aquifer’s pore volume is always calculated from the data entered on the AQUNUM keyword using$\mathit{Pore} \mathit{Volume}=\mathit{AREA}\times \mathit{LENGTH}\times \mathit{PORO}$and any modifications to the host cell values performed in either the GRID or EDIT sections are always ignored for cells declared as numerical aquifers cells.
- For the transmissibility calculation either the cross-sectional area (AREA) defined on the AQUNUM keyword may be used or the connecting cell cross-sectional area by setting the AQUOPT1 variable on the AQUCON keyword.
\end{verbatim}

In order to fully define a numerical aquifer one has to define the
aquifer properties via the AQUNUM keyword, and how the aquifer is
connected to the reservoir using the AQUCON keyword in the GRID or
SOLUTION sections.

\begin{tcolorbox}[enhanced jigsaw, title=\textcolor{quarto-callout-note-color}{\faInfo}\hspace{0.5em}{Note}, coltitle=black, breakable, toprule=.15mm, colbacktitle=quarto-callout-note-color!10!white, leftrule=.75mm, opacitybacktitle=0.6, colback=white, bottomtitle=1mm, toptitle=1mm, titlerule=0mm, left=2mm, colframe=quarto-callout-note-color-frame, arc=.35mm, rightrule=.15mm, bottomrule=.15mm, opacityback=0]

If the AQUCON keyword has been utilized in the run deck then OPM Flow
will write the AQUIFERN array to the *.INIT file in order to visualize
the aquifer connections in OPM ResInsight. This is accomplished by
setting the AQUIFERN value to 2\^{}(AQUID-1) for cells connected to
aquifer AQUID. If a cell is connected to multiple numerical aquifers
then AQUIFERN is summed for all aquifers connected to a cell. Note that
connecting cells to multiple aquifers is best avoided. Finally for cells
representing the numerical aquifers themselves, AQUIFERN is set to minus
AQUID.

\end{tcolorbox}

Using one aquifer cell should generally be sufficient, provided the
aquifer properties are constant. However, employing multiple cells may
be appropriate if the aquifer properties vary with distance or depth,
which is not uncommon. Secondly, having multiple cells may help to
minimize throughput-related convergence problems by increasing the pore
volumes exponentially away from the reservoir cells. Typically, three to
five aquifer cells are employed under theses circumstances.

The aquifer pore volume can be used to define the amount of pressure
support; whereas, the aquifer transmissibility will influence the
responsiveness of the aquifer.

See also the AQUDIMS keyword in the RUNSPEC section that defines the
numerical aquifer dimensions.

\subsubsection{Example}\label{example-136}

Given the following grid and aquifer dimensions in the RUNSPEC section:

\begin{verbatim}
--
--       MAX     MAX     MAX
--       NDIVIX  NDIVIY  NDIVIZ
DIMENS
         96      58      28                                                    /

--       AQF     AQF     AQF     AQF     AQF     AQF    AQF    AQF
--       MXAQN   MXNAQC  NIFTBL  NRIFTB  NANAQU  NCAMAX MXNALI MXAAQL
AQUDIMS
         3       3       1*       1*      1*      1*     1*     1*             /
\end{verbatim}

The following numerical aquifer definition:

\begin{verbatim}
--
--       NUMERICAL AQUIFER DESCRIPTION
--
--       AQUID  - LOCATION - AQF      AQF     AQF   AQF  AQF    AQF  PVT SATNUM
--       NUMBER   I   J   K  AREA     LENGTH  PORO  PERM DATUM  PRES TAB TAB
--
AQUNUM
         1        1   3  28  5.20E03  2.0E3   0.05  0.3  1*     1*   1*  1*    /
         1        1   2  28  5.20E06  2.0E6   0.05  0.3  1*     1*   1*  1*    /
         1        1   1  28  5.20E09  2.0E9   0.05  0.3  1*     1*   1*  1*    /
/
\end{verbatim}

defines one numerical aquifer consisting of three cells. The connection
to the grid cells would take the form of:

\begin{verbatim}
--
--        NUMERICAL AQUIFER CONNECTIONS
--
--       ID     ---------- BOX ---------   CONNECT  TRANS  TRANS  ADJOIN
--       NUMBER I1  I2   J1  J2   K1  K2   FACE     MULT   OPTN   CELLS
AQUCON
         1       1   1    4  58   28   28  'K+'     1.3    1*     1*          /
         1       2  96    1  58   28   28  'K+'     1.3    1*     1*          /
/
\end{verbatim}

that creates a basal aquifer {[}Basal Aquifer: An aquifer located at the
bottom of a geologic unit.{]}.

\subsection{AUTOCOAR -- Define Auto Refinement Grid Coarsen
Area}\label{autocoar-define-auto-refinement-grid-coarsen-area}

\begin{longtable}[]{@{}
  >{\raggedright\arraybackslash}p{(\linewidth - 14\tabcolsep) * \real{0.1250}}
  >{\raggedright\arraybackslash}p{(\linewidth - 14\tabcolsep) * \real{0.1250}}
  >{\raggedright\arraybackslash}p{(\linewidth - 14\tabcolsep) * \real{0.1250}}
  >{\raggedright\arraybackslash}p{(\linewidth - 14\tabcolsep) * \real{0.1250}}
  >{\raggedright\arraybackslash}p{(\linewidth - 14\tabcolsep) * \real{0.1250}}
  >{\raggedright\arraybackslash}p{(\linewidth - 14\tabcolsep) * \real{0.1250}}
  >{\raggedright\arraybackslash}p{(\linewidth - 14\tabcolsep) * \real{0.1250}}
  >{\raggedright\arraybackslash}p{(\linewidth - 14\tabcolsep) * \real{0.1250}}@{}}
\toprule\noalign{}
\begin{minipage}[b]{\linewidth}\raggedright
RUNSPEC
\end{minipage} & \begin{minipage}[b]{\linewidth}\raggedright
GRID
\end{minipage} & \begin{minipage}[b]{\linewidth}\raggedright
EDIT
\end{minipage} & \begin{minipage}[b]{\linewidth}\raggedright
PROPS
\end{minipage} & \begin{minipage}[b]{\linewidth}\raggedright
REGIONS
\end{minipage} & \begin{minipage}[b]{\linewidth}\raggedright
SOLUTION
\end{minipage} & \begin{minipage}[b]{\linewidth}\raggedright
SUMMARY
\end{minipage} & \begin{minipage}[b]{\linewidth}\raggedright
SCHEDULE
\end{minipage} \\
\midrule\noalign{}
\endhead
\bottomrule\noalign{}
\endlastfoot
\end{longtable}

\subsubsection{Description}\label{description-179}

The AUTOCOAR keyword defines an area in the global grid that should be
coarsen for when the AUTOREF keyword has been declared in the RUNSPEC
section.

This keyword is not supported by OPM Flow but would change the results
if supported so the simulation will be stopped.

\subsection{BCCON -- Define Boundary Conditions
Connections}\label{bccon-define-boundary-conditions-connections}

\begin{longtable}[]{@{}
  >{\raggedright\arraybackslash}p{(\linewidth - 14\tabcolsep) * \real{0.1250}}
  >{\raggedright\arraybackslash}p{(\linewidth - 14\tabcolsep) * \real{0.1250}}
  >{\raggedright\arraybackslash}p{(\linewidth - 14\tabcolsep) * \real{0.1250}}
  >{\raggedright\arraybackslash}p{(\linewidth - 14\tabcolsep) * \real{0.1250}}
  >{\raggedright\arraybackslash}p{(\linewidth - 14\tabcolsep) * \real{0.1250}}
  >{\raggedright\arraybackslash}p{(\linewidth - 14\tabcolsep) * \real{0.1250}}
  >{\raggedright\arraybackslash}p{(\linewidth - 14\tabcolsep) * \real{0.1250}}
  >{\raggedright\arraybackslash}p{(\linewidth - 14\tabcolsep) * \real{0.1250}}@{}}
\toprule\noalign{}
\begin{minipage}[b]{\linewidth}\raggedright
RUNSPEC
\end{minipage} & \begin{minipage}[b]{\linewidth}\raggedright
GRID
\end{minipage} & \begin{minipage}[b]{\linewidth}\raggedright
EDIT
\end{minipage} & \begin{minipage}[b]{\linewidth}\raggedright
PROPS
\end{minipage} & \begin{minipage}[b]{\linewidth}\raggedright
REGIONS
\end{minipage} & \begin{minipage}[b]{\linewidth}\raggedright
SOLUTION
\end{minipage} & \begin{minipage}[b]{\linewidth}\raggedright
SUMMARY
\end{minipage} & \begin{minipage}[b]{\linewidth}\raggedright
SCHEDULE
\end{minipage} \\
\midrule\noalign{}
\endhead
\bottomrule\noalign{}
\endlastfoot
\end{longtable}

\subsubsection{Description}\label{description-180}

The \hyperref[REF_HEADING_KEYWORD_BCCON]{BCCON} keyword defines the grid
connections for the boundary conditions.

Together the \hyperref[REF_HEADING_KEYWORD_BCCON]{BCCON} and
\hyperref[REF_HEADING_KEYWORD_BCPROP]{BCPROP} keywords define the
boundary conditions for the model, and can be used to set boundary
conditions for when external influx or efflux volumes are influencing
the reservoir pressure and production history. For example, when the
average reservoir pressure remains constant throughout the production
period due to water influx, or gas migration from an external source.

\begin{tcolorbox}[enhanced jigsaw, title=\textcolor{quarto-callout-note-color}{\faInfo}\hspace{0.5em}{Note}, coltitle=black, breakable, toprule=.15mm, colbacktitle=quarto-callout-note-color!10!white, leftrule=.75mm, opacitybacktitle=0.6, colback=white, bottomtitle=1mm, toptitle=1mm, titlerule=0mm, left=2mm, colframe=quarto-callout-note-color-frame, arc=.35mm, rightrule=.15mm, bottomrule=.15mm, opacityback=0]

This is an OPM Flow specific keyword and will therefore cause an error
in the commercial simulator.

\end{tcolorbox}

\begin{longtable}[]{@{}
  >{\raggedright\arraybackslash}p{(\linewidth - 6\tabcolsep) * \real{0.2500}}
  >{\raggedright\arraybackslash}p{(\linewidth - 6\tabcolsep) * \real{0.2500}}
  >{\raggedright\arraybackslash}p{(\linewidth - 6\tabcolsep) * \real{0.2500}}
  >{\raggedright\arraybackslash}p{(\linewidth - 6\tabcolsep) * \real{0.2500}}@{}}
\toprule\noalign{}
\begin{minipage}[b]{\linewidth}\raggedright
No.
\end{minipage} & \begin{minipage}[b]{\linewidth}\raggedright
Name
\end{minipage} & \begin{minipage}[b]{\linewidth}\raggedright
Description
\end{minipage} & \begin{minipage}[b]{\linewidth}\raggedright
Default
\end{minipage} \\
\midrule\noalign{}
\endhead
\bottomrule\noalign{}
\endlastfoot
Field & Metric & Laboratory & \\
1 & INDEX & A positive integer that identifies the boundary condition. &
None \\
2 & I1 & A positive integer that defines the lower bound of the grid in
the I-direction for which the boundary conditions are to be applied,
must be greater than or equal 1 and less than or equal to I2 and NX. &
1 \\
2 & I2 & A positive integer that defines the upper bound of the grid in
the I-direction for which the boundary conditions are to be applied,
must be greater than or equal to II and less than or equal to NX & NX \\
3 & J1 & A positive integer that defines the lower bound of the grid in
the J-direction for which the boundary conditions are to be applied,
must be greater than or equal 1 and less than or equal to J2 and NY. &
1 \\
4 & J2 & A positive integer that defines the upper bound of the grid in
the J-direction for which the boundary conditions are to be applied,
must be greater than or equal to JI and less than or equal to NY. &
NY \\
5 & K1 & A positive integer that defines the lower bound of the grid in
the K-direction for which the boundary conditions are to be applied,
must be greater than or equal to one and less than or equal to K2 and
NZ. & 1 \\
6 & K2 & A positive integer that defines the upper bound of the grid in
the K-direction for which the boundary conditions are to be applied,
must be greater than or equal to KI and less than or equal to NZ. &
NZ \\
7 & DIRECT & A character string that defines the direction to apply the
boundary conditions, and should be set to one of the following X, Y, or
Z for the positive direction, or X-, Y- or Z- for the negative
direction. & None \\
Notes: & & & \\
\end{longtable}

\emph{Table 6.3.12.1: BCCON Keyword Description}

See also the AQUFLUX keyword that is supported by OPM Flow in both the
SOLUTION and SCHEDULE sections, to define a constant flux analytical
aquifer.

If the \hyperref[REF_HEADING_KEYWORD_BCCON]{BCCON} and
\hyperref[REF_HEADING_KEYWORD_BCPROP]{BCPROP} keywords are not present
in the input deck, then the boundary conditions for the model are set to
be no flow, which is the normal behavior in both OPM Flow and the
commercial simulator.

The BC keyword has been replaced by the
\hyperref[REF_HEADING_KEYWORD_BCCON]{BCCON} and
\hyperref[REF_HEADING_KEYWORD_BCPROP]{BCPROP} keywords.

\subsubsection{Examples}\label{examples-7}

The first example shows how to set a constant pressure boundary using
TYPE equal to FREE:

\begin{verbatim}
--
--       DEFINE BOUNDARY CONDITIONS CONNECTION (OPM FLOW GRID KEYWORD)
--
--INDEX  ---------- BOX ---------   BC
--       I1  I2   J1  J2   K1  K2   DIRC
BCCON
1        1   1    1   1*   1   1*   X-    /
2        1   1*   1   1    1   1*   Y     /
/

--
--       DEFINE BOUNDARY CONDITIONS PROPERTIES (OPM FLOW SCHEDULE KEYWORD)
--
--INDEX  BC         BC     BC    BC     BC
--       TYPE       PHASE  RATE  PRESS  TEMP
BCPROP
1        FREE       1*     1*    1*     1*    /
2        FREE                                 /
/
\end{verbatim}

With this option it is only necessary to define the boundary cells and
all the other parameters (COMPONENT, RATE, PRESS, and TEMP) can be
defaulted, as they are ignored when TYPE equals FREE.

The next example is based on NX, NY and NZ equal to 20, 1, 10
respectively, on the DIMENS keyword in the RUNSPEC section, and shows
how different boundary types can be assigned to different parts of the
model.

\begin{verbatim}
--
--       DEFINE BOUNDARY CONDITIONS CONNECTION (OPM FLOW GRID KEYWORD)
--
--INDEX  ---------- BOX ---------   BC
--       I1  I2   J1  J2   K1  K2   DIRC
BCCON
1        1   1    1   1    1   10   X-    /
2        20  20   1   1    1   10   X     /
/

--
--       DEFINE BOUNDARY CONDITIONS PROPERTIES (OPM FLOW SCHEDULE KEYWORD)
--
--INDEX  BC         BC     BC    BC     BC
--       TYPE       PHASE  RATE  PRESS  TEMP
BCPROP
1        DIRICHLET       GAS    1*    256.0  100.0 /
2        FREE       4*                        /
/
\end{verbatim}

The last example shows how the DIRICHLET boundary condition option may
be used:

\begin{verbatim}
--
--       DEFINE BOUNDARY CONDITIONS CONNECTION (OPM FLOW GRID KEYWORD)
--
--INDEX  ---------- BOX ---------   BC
--       I1  I2   J1  J2   K1  K2   DIRC
BCCON
1        1   1    1   1*   1   1*   X-    /
2        1   1*   1   1    1   1*   Y     /
/

--
--       DEFINE BOUNDARY CONDITIONS PROPERTIES (OPM FLOW SCHEDULE KEYWORD)
--
--INDEX  BC         BC     BC    BC     BC
--       TYPE       PHASE  RATE  PRESS  TEMP
BCPROP
1        DIRICHLET  WAT    1*    256.0  100.0 /
2        DIRICHLET  WAT    1*    1*     100.0 /
/
\end{verbatim}

Here, the first line sets both the pressure and temperature at the
boundary, and the second line defaults the pressure entry, so that the
simulator calculated initial boundary pressure will be used.

\subsection{BOUNDARY -- Define a Boundary Box for
Printing}\label{boundary-define-a-boundary-box-for-printing}

\begin{longtable}[]{@{}
  >{\raggedright\arraybackslash}p{(\linewidth - 14\tabcolsep) * \real{0.1250}}
  >{\raggedright\arraybackslash}p{(\linewidth - 14\tabcolsep) * \real{0.1250}}
  >{\raggedright\arraybackslash}p{(\linewidth - 14\tabcolsep) * \real{0.1250}}
  >{\raggedright\arraybackslash}p{(\linewidth - 14\tabcolsep) * \real{0.1250}}
  >{\raggedright\arraybackslash}p{(\linewidth - 14\tabcolsep) * \real{0.1250}}
  >{\raggedright\arraybackslash}p{(\linewidth - 14\tabcolsep) * \real{0.1250}}
  >{\raggedright\arraybackslash}p{(\linewidth - 14\tabcolsep) * \real{0.1250}}
  >{\raggedright\arraybackslash}p{(\linewidth - 14\tabcolsep) * \real{0.1250}}@{}}
\toprule\noalign{}
\begin{minipage}[b]{\linewidth}\raggedright
RUNSPEC
\end{minipage} & \begin{minipage}[b]{\linewidth}\raggedright
GRID
\end{minipage} & \begin{minipage}[b]{\linewidth}\raggedright
EDIT
\end{minipage} & \begin{minipage}[b]{\linewidth}\raggedright
PROPS
\end{minipage} & \begin{minipage}[b]{\linewidth}\raggedright
REGIONS
\end{minipage} & \begin{minipage}[b]{\linewidth}\raggedright
SOLUTION
\end{minipage} & \begin{minipage}[b]{\linewidth}\raggedright
SUMMARY
\end{minipage} & \begin{minipage}[b]{\linewidth}\raggedright
SCHEDULE
\end{minipage} \\
\midrule\noalign{}
\endhead
\bottomrule\noalign{}
\endlastfoot
\end{longtable}

\subsubsection{Description}\label{description-181}

The BOUNDARY keyword defines a rectangular grid for printing various
arrays to the output print file (*.PRT); thus, avoiding printing all the
elements in the selected array.

This keyword is not supported by OPM Flow but would change the results
if supported so the simulation will be stopped.

\subsection{BOX -- Define a Range of Grid Blocks to Enter Property
Data}\label{box-define-a-range-of-grid-blocks-to-enter-property-data}

\begin{longtable}[]{@{}
  >{\raggedright\arraybackslash}p{(\linewidth - 14\tabcolsep) * \real{0.1250}}
  >{\raggedright\arraybackslash}p{(\linewidth - 14\tabcolsep) * \real{0.1250}}
  >{\raggedright\arraybackslash}p{(\linewidth - 14\tabcolsep) * \real{0.1250}}
  >{\raggedright\arraybackslash}p{(\linewidth - 14\tabcolsep) * \real{0.1250}}
  >{\raggedright\arraybackslash}p{(\linewidth - 14\tabcolsep) * \real{0.1250}}
  >{\raggedright\arraybackslash}p{(\linewidth - 14\tabcolsep) * \real{0.1250}}
  >{\raggedright\arraybackslash}p{(\linewidth - 14\tabcolsep) * \real{0.1250}}
  >{\raggedright\arraybackslash}p{(\linewidth - 14\tabcolsep) * \real{0.1250}}@{}}
\toprule\noalign{}
\begin{minipage}[b]{\linewidth}\raggedright
RUNSPEC
\end{minipage} & \begin{minipage}[b]{\linewidth}\raggedright
GRID
\end{minipage} & \begin{minipage}[b]{\linewidth}\raggedright
EDIT
\end{minipage} & \begin{minipage}[b]{\linewidth}\raggedright
PROPS
\end{minipage} & \begin{minipage}[b]{\linewidth}\raggedright
REGIONS
\end{minipage} & \begin{minipage}[b]{\linewidth}\raggedright
SOLUTION
\end{minipage} & \begin{minipage}[b]{\linewidth}\raggedright
SUMMARY
\end{minipage} & \begin{minipage}[b]{\linewidth}\raggedright
SCHEDULE
\end{minipage} \\
\midrule\noalign{}
\endhead
\bottomrule\noalign{}
\endlastfoot
\end{longtable}

\subsubsection{Description}\label{description-182}

BOX defines a range of grid blocks for which subsequent data is assigned
for all the cells in the defined BOX. Values are set for cells within
the defined box grid using natural reading order, initially along the
I-direction then J-direction and finally the K-direction. If fewer
values are assigned than exist within the defined block space, then
subsequent values are set starting from the next block that was not
previously assigned for that property. This is the same behavior as
applies to setting grid properties for an unboxed grid.

Note that the BOX grid is reset by the keyword ENDBOX by resetting the
current defined BOX to be the whole grid. The keyword can be used for
any array and for all grid types.

\begin{longtable}[]{@{}
  >{\raggedright\arraybackslash}p{(\linewidth - 6\tabcolsep) * \real{0.2500}}
  >{\raggedright\arraybackslash}p{(\linewidth - 6\tabcolsep) * \real{0.2500}}
  >{\raggedright\arraybackslash}p{(\linewidth - 6\tabcolsep) * \real{0.2500}}
  >{\raggedright\arraybackslash}p{(\linewidth - 6\tabcolsep) * \real{0.2500}}@{}}
\toprule\noalign{}
\begin{minipage}[b]{\linewidth}\raggedright
No.
\end{minipage} & \begin{minipage}[b]{\linewidth}\raggedright
Name
\end{minipage} & \begin{minipage}[b]{\linewidth}\raggedright
Description
\end{minipage} & \begin{minipage}[b]{\linewidth}\raggedright
Default
\end{minipage} \\
\midrule\noalign{}
\endhead
\bottomrule\noalign{}
\endlastfoot
1 & I1 & A positive integer that defines the lower bound of the array in
the I-direction to be modified must be greater than or equal 1 and less
than or equal to I2 and NX. & 1 \\
2 & I2 & A positive integer that defines the upper bound of the array in
the I-direction to be modified must be greater than or equal to II and
less than or equal to NX & NX \\
3 & J1 & A positive integer that defines the lower bound of the array in
the J-direction to be modified must be greater than or equal 1 and less
than or equal to J2 and NY. & 1 \\
4 & J2 & A positive integer that defines the upper bound of the array in
the J-direction to be modified must be greater than or equal to JI and
less than or equal to NY. & NY \\
5 & K1 & A positive integer that defines the lower bound of the array in
the K-direction to be modified must be greater than or equal to one and
less than or equal to K2 and NZ. & 1 \\
6 & K2 & A positive integer that defines the upper bound of the array in
the K-direction to be modified must be greater than or equal to KI and
less than or equal to NZ. & NZ \\
Notes: & & & \\
\end{longtable}

\emph{Table 6.11: BOX Keyword Description}

See also the ADD, COPY, ENDBOX, EQUALS, and MULTIPLY keywords can also
be used to enter data in a subset of the model.

\subsubsection{Examples}\label{examples-8}

\begin{verbatim}
--
--        DEFINE A BOX GRID FOR THE BOTTOM LAYER OF A 100 X 100 X 20 MODEL
--
--        ---------- BOX ---------
--        I1  I2   J1  J2   K1  K2
BOX
          1*  1*   1*  1*   20  20 / SELECT THE BOTTOM LAYER
--
--        DEFINE THE POROSITY AND OTHER PROPERTIES ON THE BOX GRID
--
PORO
          10000*0.300                                                                     /
PERMX
          5000*100.0   5000*75.0                                                                     /
NTG
          10000*0.500                                                                     /
--
--        RESET THE INPUT BOX TO BE THE FULL MODEL
--
ENDBOX
\end{verbatim}

The above example set the BOX grid to be the last layer in the model
which means that 100 x 100, that is 10,000 values need to entered for
each property.

Alternatively, one could use the EQUALS keyword to accomplish the same
thing.

\begin{verbatim}
-- -- ARRAY    CONSTANT --  ---------- BOX ---------
--                          I1  I2   J1  J2   K1  K2
EQUALS
   'PORO’      0.3000       1*  1*   1*  1*   20  20 / PORO  TO 0.30 IN LAYER 20
   'PERMX'     100.00       1*  1*   1   50   20  20 / PERMX TO 100. IN LAYER 20
   'PERMX'      75.00       1*  1*   51  100  20  20 / PERMX TO 75.0 IN LAYER 20
   'NTG’       0.5000       1*  1*   1*  1*   20  20 / NTG   TO 0.50 IN LAYER 20
/
\end{verbatim}

\begin{tcolorbox}[enhanced jigsaw, title=\textcolor{quarto-callout-note-color}{\faInfo}\hspace{0.5em}{Note}, coltitle=black, breakable, toprule=.15mm, colbacktitle=quarto-callout-note-color!10!white, leftrule=.75mm, opacitybacktitle=0.6, colback=white, bottomtitle=1mm, toptitle=1mm, titlerule=0mm, left=2mm, colframe=quarto-callout-note-color-frame, arc=.35mm, rightrule=.15mm, bottomrule=.15mm, opacityback=0]

It is good practice to always use the ENDBOX keyword to reset the input
back to the full grid when all the modifications for a sub-grid have
been completed.

\end{tcolorbox}

\subsection{BTOBALFA -- Dual Porosity Matrix to Fracture Multiplier (All
Cells)}\label{btobalfa-dual-porosity-matrix-to-fracture-multiplier-all-cells}

\begin{longtable}[]{@{}
  >{\raggedright\arraybackslash}p{(\linewidth - 14\tabcolsep) * \real{0.1250}}
  >{\raggedright\arraybackslash}p{(\linewidth - 14\tabcolsep) * \real{0.1250}}
  >{\raggedright\arraybackslash}p{(\linewidth - 14\tabcolsep) * \real{0.1250}}
  >{\raggedright\arraybackslash}p{(\linewidth - 14\tabcolsep) * \real{0.1250}}
  >{\raggedright\arraybackslash}p{(\linewidth - 14\tabcolsep) * \real{0.1250}}
  >{\raggedright\arraybackslash}p{(\linewidth - 14\tabcolsep) * \real{0.1250}}
  >{\raggedright\arraybackslash}p{(\linewidth - 14\tabcolsep) * \real{0.1250}}
  >{\raggedright\arraybackslash}p{(\linewidth - 14\tabcolsep) * \real{0.1250}}@{}}
\toprule\noalign{}
\begin{minipage}[b]{\linewidth}\raggedright
RUNSPEC
\end{minipage} & \begin{minipage}[b]{\linewidth}\raggedright
GRID
\end{minipage} & \begin{minipage}[b]{\linewidth}\raggedright
EDIT
\end{minipage} & \begin{minipage}[b]{\linewidth}\raggedright
PROPS
\end{minipage} & \begin{minipage}[b]{\linewidth}\raggedright
REGIONS
\end{minipage} & \begin{minipage}[b]{\linewidth}\raggedright
SOLUTION
\end{minipage} & \begin{minipage}[b]{\linewidth}\raggedright
SUMMARY
\end{minipage} & \begin{minipage}[b]{\linewidth}\raggedright
SCHEDULE
\end{minipage} \\
\midrule\noalign{}
\endhead
\bottomrule\noalign{}
\endlastfoot
\end{longtable}

\subsubsection{Description}\label{description-183}

The BTOBALFA keyword defines a dual porosity matrix to fracture
multiplier that is applied to all cells, for when the Dual Porosity
model has been activated by either the DUALPORO or the DUALPERM keywords
in the RUNSPEC section.

See also the BTOBALFV keyword in the GRID section that applies a
multipliers to individual cells.

This keyword is not supported by OPM Flow but would change the results
if supported so the simulation will be stopped.

\subsection{BTOBALFV -- Dual Porosity Matrix to Fracture Multiplier
(Individual
Cells)}\label{btobalfv-dual-porosity-matrix-to-fracture-multiplier-individual-cells}

\begin{longtable}[]{@{}
  >{\raggedright\arraybackslash}p{(\linewidth - 14\tabcolsep) * \real{0.1250}}
  >{\raggedright\arraybackslash}p{(\linewidth - 14\tabcolsep) * \real{0.1250}}
  >{\raggedright\arraybackslash}p{(\linewidth - 14\tabcolsep) * \real{0.1250}}
  >{\raggedright\arraybackslash}p{(\linewidth - 14\tabcolsep) * \real{0.1250}}
  >{\raggedright\arraybackslash}p{(\linewidth - 14\tabcolsep) * \real{0.1250}}
  >{\raggedright\arraybackslash}p{(\linewidth - 14\tabcolsep) * \real{0.1250}}
  >{\raggedright\arraybackslash}p{(\linewidth - 14\tabcolsep) * \real{0.1250}}
  >{\raggedright\arraybackslash}p{(\linewidth - 14\tabcolsep) * \real{0.1250}}@{}}
\toprule\noalign{}
\begin{minipage}[b]{\linewidth}\raggedright
RUNSPEC
\end{minipage} & \begin{minipage}[b]{\linewidth}\raggedright
GRID
\end{minipage} & \begin{minipage}[b]{\linewidth}\raggedright
EDIT
\end{minipage} & \begin{minipage}[b]{\linewidth}\raggedright
PROPS
\end{minipage} & \begin{minipage}[b]{\linewidth}\raggedright
REGIONS
\end{minipage} & \begin{minipage}[b]{\linewidth}\raggedright
SOLUTION
\end{minipage} & \begin{minipage}[b]{\linewidth}\raggedright
SUMMARY
\end{minipage} & \begin{minipage}[b]{\linewidth}\raggedright
SCHEDULE
\end{minipage} \\
\midrule\noalign{}
\endhead
\bottomrule\noalign{}
\endlastfoot
\end{longtable}

\subsubsection{Description}\label{description-184}

The BTOBALFV keyword defines a dual porosity matrix to fracture
multiplier that is applied to individual cells, for when the Dual
Porosity model has been invoked by either the DUALPORO or the DUALPERM
keywords in the RUNSPEC section.

See also the BTOBALFAV keyword in the GRID section that applies a
constant multiplier to all cells.

This keyword is not supported by OPM Flow but would change the results
if supported so the simulation will be stopped.

\subsection{CARFIN -- Define a Cartesian Local Grid
Refinement}\label{carfin-define-a-cartesian-local-grid-refinement}

\begin{longtable}[]{@{}
  >{\raggedright\arraybackslash}p{(\linewidth - 14\tabcolsep) * \real{0.1250}}
  >{\raggedright\arraybackslash}p{(\linewidth - 14\tabcolsep) * \real{0.1250}}
  >{\raggedright\arraybackslash}p{(\linewidth - 14\tabcolsep) * \real{0.1250}}
  >{\raggedright\arraybackslash}p{(\linewidth - 14\tabcolsep) * \real{0.1250}}
  >{\raggedright\arraybackslash}p{(\linewidth - 14\tabcolsep) * \real{0.1250}}
  >{\raggedright\arraybackslash}p{(\linewidth - 14\tabcolsep) * \real{0.1250}}
  >{\raggedright\arraybackslash}p{(\linewidth - 14\tabcolsep) * \real{0.1250}}
  >{\raggedright\arraybackslash}p{(\linewidth - 14\tabcolsep) * \real{0.1250}}@{}}
\toprule\noalign{}
\begin{minipage}[b]{\linewidth}\raggedright
RUNSPEC
\end{minipage} & \begin{minipage}[b]{\linewidth}\raggedright
GRID
\end{minipage} & \begin{minipage}[b]{\linewidth}\raggedright
EDIT
\end{minipage} & \begin{minipage}[b]{\linewidth}\raggedright
PROPS
\end{minipage} & \begin{minipage}[b]{\linewidth}\raggedright
REGIONS
\end{minipage} & \begin{minipage}[b]{\linewidth}\raggedright
SOLUTION
\end{minipage} & \begin{minipage}[b]{\linewidth}\raggedright
SUMMARY
\end{minipage} & \begin{minipage}[b]{\linewidth}\raggedright
SCHEDULE
\end{minipage} \\
\midrule\noalign{}
\endhead
\bottomrule\noalign{}
\endlastfoot
\end{longtable}

\subsubsection{Description}\label{description-185}

CARFIN defines a Cartesian Local Grid Refinement (``LGR'') in a cell or
a group of cells in the host grid, for when LGRs have been activated for
the input deck using the LGR keyword in the RUNSPEC section. The keyword
marks the start of an LGR description section and all subsequent
keywords between the CARFIN and ENDFIN keywords are deemed to be
associated with the current LGR and not the host grid.

This keyword is not supported by OPM Flow but would change the results
if supported so the simulation will be stopped.

\begin{longtable}[]{@{}
  >{\raggedright\arraybackslash}p{(\linewidth - 6\tabcolsep) * \real{0.2500}}
  >{\raggedright\arraybackslash}p{(\linewidth - 6\tabcolsep) * \real{0.2500}}
  >{\raggedright\arraybackslash}p{(\linewidth - 6\tabcolsep) * \real{0.2500}}
  >{\raggedright\arraybackslash}p{(\linewidth - 6\tabcolsep) * \real{0.2500}}@{}}
\toprule\noalign{}
\begin{minipage}[b]{\linewidth}\raggedright
No.
\end{minipage} & \begin{minipage}[b]{\linewidth}\raggedright
Name
\end{minipage} & \begin{minipage}[b]{\linewidth}\raggedright
Description
\end{minipage} & \begin{minipage}[b]{\linewidth}\raggedright
Default
\end{minipage} \\
\midrule\noalign{}
\endhead
\bottomrule\noalign{}
\endlastfoot
1 & LGRNAME & A character string of up to eight characters in length
that defines the LGR name for which the LGR is being being defined. &
None \\
2 & I1 & A positive integer that defines the lower index of the global
or host grid in the I-direction to be refined; must be greater than or
equal 1 and less than or equal to I2 and NX on the DIMENS keyword in the
RUNSPEC section. & None \\
3 & I2 & A positive integer that defines the upper index of the global
or host grid in the I-direction to be refined; must be greater than or
equal 1 and II, and less than or equal to NX on the DIMENS keyword in
the RUNSPEC section. & None \\
4 & J1 & A positive integer that defines the lower index of the global
or host grid in the J-direction to be refined; must be greater than or
equal 1 and less than or equal to J2 and NY on the DIMENS keyword in the
RUNSPEC section. & None \\
5 & J2 & A positive integer that defines the upper index of the global
or host grid in the J-direction to be refined; must be greater than or
equal 1 and JI, and less than or equal to NY on the DIMENS keyword in
the RUNSPEC section. & None \\
6 & K1 & A positive integer that defines the lower index of the global
or host grid in the K-direction to be refined; must be greater than or
equal 1 and less than or equal to K2 and NZ on the DIMENS keyword in the
RUNSPEC section. & None \\
7 & K2 & A positive integer that defines the upper index of the global
or host grid in the K-direction to be refined; must be greater than or
equal 1 and KI, and less than or equal to NZ on the DIMENS keyword in
the RUNSPEC section. & None \\
8 & NX & A positive integer value that defines the number of LGR grid
blocks in the x direction for Cartesian grids or the number of grid
blocks in the r direction for radial grids & None \\
9 & NY & A positive integer value that defines the number of LGR grid
blocks in the y direction for Cartesian grids or the number of grid
blocks in the theta direction for radial grids. & None \\
10 & NZ & A positive integer value that defines the number of LGR grid
blocks in the z direction for both Cartesian and radial grids. & None \\
11 & MXWELS & A positive integer defining the maximum number of wells
contained in this LGR. & None \\
12 & HOSTNAME & A character string of up to eight characters in length
that defines the host grid name for nested refinements. The default
value of ``GLOBAL'' sets the host name to the global grid, that is for a
conventional LGR. A nested refinement is when the HOSTNAME is a
previously declared LGR for which the current LGR is specifying a
further LGR refinement. & GLOBAL \\
Notes: & & & \\
\end{longtable}

\emph{Table 6.12: CARFIN Keyword Description}

Note that if the Dual Porosity option has be activated by either the
DUALPORO or DUALPERM keywords in the RUNSPEC section, then the host grid
definition (I1-I2, J1-J2, K1-K2) applies only to the matrix cells;
however, the LGR NZ parameter in this case must include the fracture
blocks, similar to the NZ parameter on the DIMENS keyword. This means
that all property data should be entered for both the matrix and
fracture cells in the LGR description.

\subsubsection{Example}\label{example-137}

The example below defines an LGR in the global grid, named LGR-OP01 with
a maximum of one well allowed in the LGR.

\begin{verbatim}
--
--       CARFIN LGR GRID COMMANDS
--
--       LGR        ----- HOST GRID ------   -- CARFIN GRID --  MAX     HOST
--       NAME       I1  I2  J1  J2  K1  K2     NX    NY    NZ   WELLS   NAME
CARFIN
         LGR-OP01   24  24  87  87   1  50      3     3    50     1     GLOBAL /


ENDFIN
\end{verbatim}

Here the one global cell in the areal plane (24, 87) is divided into
three LGR cells in the x-direction and three cells in the y-direction.
Since no other property data is given, then the LGR cells take their
properties from the host grid, that is the global grid.

\subsection{COALNUM -- Define the Coal Region
Numbers}\label{coalnum-define-the-coal-region-numbers}

\begin{longtable}[]{@{}
  >{\raggedright\arraybackslash}p{(\linewidth - 14\tabcolsep) * \real{0.1250}}
  >{\raggedright\arraybackslash}p{(\linewidth - 14\tabcolsep) * \real{0.1250}}
  >{\raggedright\arraybackslash}p{(\linewidth - 14\tabcolsep) * \real{0.1250}}
  >{\raggedright\arraybackslash}p{(\linewidth - 14\tabcolsep) * \real{0.1250}}
  >{\raggedright\arraybackslash}p{(\linewidth - 14\tabcolsep) * \real{0.1250}}
  >{\raggedright\arraybackslash}p{(\linewidth - 14\tabcolsep) * \real{0.1250}}
  >{\raggedright\arraybackslash}p{(\linewidth - 14\tabcolsep) * \real{0.1250}}
  >{\raggedright\arraybackslash}p{(\linewidth - 14\tabcolsep) * \real{0.1250}}@{}}
\toprule\noalign{}
\begin{minipage}[b]{\linewidth}\raggedright
RUNSPEC
\end{minipage} & \begin{minipage}[b]{\linewidth}\raggedright
GRID
\end{minipage} & \begin{minipage}[b]{\linewidth}\raggedright
EDIT
\end{minipage} & \begin{minipage}[b]{\linewidth}\raggedright
PROPS
\end{minipage} & \begin{minipage}[b]{\linewidth}\raggedright
REGIONS
\end{minipage} & \begin{minipage}[b]{\linewidth}\raggedright
SOLUTION
\end{minipage} & \begin{minipage}[b]{\linewidth}\raggedright
SUMMARY
\end{minipage} & \begin{minipage}[b]{\linewidth}\raggedright
SCHEDULE
\end{minipage} \\
\midrule\noalign{}
\endhead
\bottomrule\noalign{}
\endlastfoot
\end{longtable}

\subsubsection{Description}\label{description-186}

The COALNUM keyword defines the coal region numbers for each grid block
used with the Coal Bed Methane option (``CBM''). OPM Flow does not have
a CBM option; however, the keyword is documented here for completeness.

This keyword is not supported by OPM Flow but would change the results
if supported so the simulation will be stopped.

\begin{longtable}[]{@{}
  >{\raggedright\arraybackslash}p{(\linewidth - 6\tabcolsep) * \real{0.2500}}
  >{\raggedright\arraybackslash}p{(\linewidth - 6\tabcolsep) * \real{0.2500}}
  >{\raggedright\arraybackslash}p{(\linewidth - 6\tabcolsep) * \real{0.2500}}
  >{\raggedright\arraybackslash}p{(\linewidth - 6\tabcolsep) * \real{0.2500}}@{}}
\toprule\noalign{}
\begin{minipage}[b]{\linewidth}\raggedright
No.
\end{minipage} & \begin{minipage}[b]{\linewidth}\raggedright
Name
\end{minipage} & \begin{minipage}[b]{\linewidth}\raggedright
Description
\end{minipage} & \begin{minipage}[b]{\linewidth}\raggedright
Default
\end{minipage} \\
\midrule\noalign{}
\endhead
\bottomrule\noalign{}
\endlastfoot
1 & COALNUM & COALNUM defines an array of positive integers assigning a
grid cell to a particular coal region. The maximum number of COALNUM
regions is set by the NTCREG variable on REGDIMS keywords in the RUNSPEC
section. & 1 \\
Notes: & & & \\
\end{longtable}

\emph{Table 6.13: COALNUM Keyword Description}

\subsubsection{Example}\label{example-138}

The example below sets three COALNUM regions for a 4 x 5 x 2 model.

\begin{verbatim}
--
--       DEFINE COALNUM REGIONS FOR ALL CELLS
--
COALNUM
         2  2  1  1  2  2  1  1  1  1  1  1  1  1  1  1  1  1  1  1
         3  3  1  1  3  3  1  1  1  1  1  1  1  1  1  1  1  1  1  1
/
\end{verbatim}

The above will no effect in an OPM Flow input deck.

\subsection{COARSEN -- Define Grid Coarsening
Cells}\label{coarsen-define-grid-coarsening-cells}

\begin{longtable}[]{@{}
  >{\raggedright\arraybackslash}p{(\linewidth - 14\tabcolsep) * \real{0.1250}}
  >{\raggedright\arraybackslash}p{(\linewidth - 14\tabcolsep) * \real{0.1250}}
  >{\raggedright\arraybackslash}p{(\linewidth - 14\tabcolsep) * \real{0.1250}}
  >{\raggedright\arraybackslash}p{(\linewidth - 14\tabcolsep) * \real{0.1250}}
  >{\raggedright\arraybackslash}p{(\linewidth - 14\tabcolsep) * \real{0.1250}}
  >{\raggedright\arraybackslash}p{(\linewidth - 14\tabcolsep) * \real{0.1250}}
  >{\raggedright\arraybackslash}p{(\linewidth - 14\tabcolsep) * \real{0.1250}}
  >{\raggedright\arraybackslash}p{(\linewidth - 14\tabcolsep) * \real{0.1250}}@{}}
\toprule\noalign{}
\begin{minipage}[b]{\linewidth}\raggedright
RUNSPEC
\end{minipage} & \begin{minipage}[b]{\linewidth}\raggedright
GRID
\end{minipage} & \begin{minipage}[b]{\linewidth}\raggedright
EDIT
\end{minipage} & \begin{minipage}[b]{\linewidth}\raggedright
PROPS
\end{minipage} & \begin{minipage}[b]{\linewidth}\raggedright
REGIONS
\end{minipage} & \begin{minipage}[b]{\linewidth}\raggedright
SOLUTION
\end{minipage} & \begin{minipage}[b]{\linewidth}\raggedright
SUMMARY
\end{minipage} & \begin{minipage}[b]{\linewidth}\raggedright
SCHEDULE
\end{minipage} \\
\midrule\noalign{}
\endhead
\bottomrule\noalign{}
\endlastfoot
\end{longtable}

\subsubsection{Description}\label{description-187}

The COARSEN keyword defines how a set of cells should be coarsened for
when the Local Grid Refinement (``LGR'') option has been activated by
LGR keyword in the RUNSPEC section.

This keyword is not supported by OPM Flow but would change the results
if supported so the simulation will be stopped.

\subsection{COLLAPSE -- Define Compressed Vertical Equilibrium
Cells}\label{collapse-define-compressed-vertical-equilibrium-cells}

\begin{longtable}[]{@{}
  >{\raggedright\arraybackslash}p{(\linewidth - 14\tabcolsep) * \real{0.1250}}
  >{\raggedright\arraybackslash}p{(\linewidth - 14\tabcolsep) * \real{0.1250}}
  >{\raggedright\arraybackslash}p{(\linewidth - 14\tabcolsep) * \real{0.1250}}
  >{\raggedright\arraybackslash}p{(\linewidth - 14\tabcolsep) * \real{0.1250}}
  >{\raggedright\arraybackslash}p{(\linewidth - 14\tabcolsep) * \real{0.1250}}
  >{\raggedright\arraybackslash}p{(\linewidth - 14\tabcolsep) * \real{0.1250}}
  >{\raggedright\arraybackslash}p{(\linewidth - 14\tabcolsep) * \real{0.1250}}
  >{\raggedright\arraybackslash}p{(\linewidth - 14\tabcolsep) * \real{0.1250}}@{}}
\toprule\noalign{}
\begin{minipage}[b]{\linewidth}\raggedright
RUNSPEC
\end{minipage} & \begin{minipage}[b]{\linewidth}\raggedright
GRID
\end{minipage} & \begin{minipage}[b]{\linewidth}\raggedright
EDIT
\end{minipage} & \begin{minipage}[b]{\linewidth}\raggedright
PROPS
\end{minipage} & \begin{minipage}[b]{\linewidth}\raggedright
REGIONS
\end{minipage} & \begin{minipage}[b]{\linewidth}\raggedright
SOLUTION
\end{minipage} & \begin{minipage}[b]{\linewidth}\raggedright
SUMMARY
\end{minipage} & \begin{minipage}[b]{\linewidth}\raggedright
SCHEDULE
\end{minipage} \\
\midrule\noalign{}
\endhead
\bottomrule\noalign{}
\endlastfoot
\end{longtable}

\subsubsection{Description}\label{description-188}

The COLLAPSE keyword defines the which cells can be collapsed in a
collapsed Vertical Equilibrium (``VE'') run when the VE option has been
invoked via the VE keyword in the RUNSPEC section.

This keyword is not supported by OPM Flow but would change the results
if supported so the simulation will be stopped.

\subsection{COLUMNS -- Define Input File Column
Margins}\label{columns-define-input-file-column-margins-2}

The COLUMNS keyword defines the input file column margins; characters
outside the margins are ignored by the input parser.

See COLUMNS -- Define Input File Column Margins in the GLOBAL section
for a full description.

\subsection{COORD -- Define a Set of Coordinates Lines for a Reservoir
Grid}\label{coord-define-a-set-of-coordinates-lines-for-a-reservoir-grid}

\begin{longtable}[]{@{}
  >{\raggedright\arraybackslash}p{(\linewidth - 14\tabcolsep) * \real{0.1250}}
  >{\raggedright\arraybackslash}p{(\linewidth - 14\tabcolsep) * \real{0.1250}}
  >{\raggedright\arraybackslash}p{(\linewidth - 14\tabcolsep) * \real{0.1250}}
  >{\raggedright\arraybackslash}p{(\linewidth - 14\tabcolsep) * \real{0.1250}}
  >{\raggedright\arraybackslash}p{(\linewidth - 14\tabcolsep) * \real{0.1250}}
  >{\raggedright\arraybackslash}p{(\linewidth - 14\tabcolsep) * \real{0.1250}}
  >{\raggedright\arraybackslash}p{(\linewidth - 14\tabcolsep) * \real{0.1250}}
  >{\raggedright\arraybackslash}p{(\linewidth - 14\tabcolsep) * \real{0.1250}}@{}}
\toprule\noalign{}
\begin{minipage}[b]{\linewidth}\raggedright
RUNSPEC
\end{minipage} & \begin{minipage}[b]{\linewidth}\raggedright
GRID
\end{minipage} & \begin{minipage}[b]{\linewidth}\raggedright
EDIT
\end{minipage} & \begin{minipage}[b]{\linewidth}\raggedright
PROPS
\end{minipage} & \begin{minipage}[b]{\linewidth}\raggedright
REGIONS
\end{minipage} & \begin{minipage}[b]{\linewidth}\raggedright
SOLUTION
\end{minipage} & \begin{minipage}[b]{\linewidth}\raggedright
SUMMARY
\end{minipage} & \begin{minipage}[b]{\linewidth}\raggedright
SCHEDULE
\end{minipage} \\
\midrule\noalign{}
\endhead
\bottomrule\noalign{}
\endlastfoot
\end{longtable}

\subsubsection{Description}\label{description-189}

COORD defines a set of coordinate lines or pillars for a reservoir grid
via an array. A total of 6 x (NX+1) x (NY+1) values must be specified
for each coordinate data set (or reservoir). For multiple reservoirs,
where NUMRES is greater than one, there must be 6 x (NX+1) x (NY+1) x
NUMRES values. In OPM Flow NUMRES can only be set to one.

For Cartesian geometry, each line is defined by the (x, y, z)
coordinates of two distinct points on the line. The lines are entered
with I cycling fastest then J. For radial geometry, each line is defined
by the (r, theta, z) coordinates of two distinct points on the line. The
lines are entered with R cycling fastest then THETA.

The keyword can only be used with Irregular Corner-Point Grids.

\begin{longtable}[]{@{}llll@{}}
\toprule\noalign{}
No. & Name & Description & Default \\
\midrule\noalign{}
\endhead
\bottomrule\noalign{}
\endlastfoot
Field & Metric & Laboratory & \\
1 & X1-Column & Top X coordinate & None \\
2 & Y1-Column & Top Y coordinate & \\
3 & Z1-Column & Top Z coordinate & \\
4 & X2-Column & Base X coordinate & \\
5 & Y2-Column & Base Y coordinate & \\
6 & Z2-Column & Base Z coordinate & \\
& feet & metres & cm \\
Notes: & & & \\
\end{longtable}

\emph{Table 6.14: COORD Keyword Description}

See also the SPECGRID, COORDSYS and ZCORN keywords to fully define an
Irregular Corner-Point Grid.

\subsubsection{Example}\label{example-139}

\begin{verbatim}
--
--       SPECIFY VERTICAL COORDINATE LINES FOR A REGULAR 3 x 2 GRID
--       (DX = 100 and DY = 200)
--
--       X1     Y1    Z1      X2     Y2    Z2
--       ---    ---   ----    ---    ---   ----
COORD
           0      0   1000      0      0   5000
         100      0   1000    100      0   5000
         200      0   1000    200      0   5000
         300      0   1000    300      0   5000
           0    200   1000      0    200   5000
         100    200   1000    100    200   5000
         200    200   1000    200    200   5000
         300    200   1000    300    200   5000
           0    400   1000      0    400   5000
         100    400   1000    100    400   5000
         200    400   1000    200    400   5000
         300    400   1000    300    400   5000
/
\end{verbatim}

The above example defines vertical coordinate lines for a regular 3 by 2
grid with DX equal to 100 and DY equal to 200.

\subsection{COORDSYS -- Define Coordinate Grid
Options}\label{coordsys-define-coordinate-grid-options}

\begin{longtable}[]{@{}
  >{\raggedright\arraybackslash}p{(\linewidth - 14\tabcolsep) * \real{0.1250}}
  >{\raggedright\arraybackslash}p{(\linewidth - 14\tabcolsep) * \real{0.1250}}
  >{\raggedright\arraybackslash}p{(\linewidth - 14\tabcolsep) * \real{0.1250}}
  >{\raggedright\arraybackslash}p{(\linewidth - 14\tabcolsep) * \real{0.1250}}
  >{\raggedright\arraybackslash}p{(\linewidth - 14\tabcolsep) * \real{0.1250}}
  >{\raggedright\arraybackslash}p{(\linewidth - 14\tabcolsep) * \real{0.1250}}
  >{\raggedright\arraybackslash}p{(\linewidth - 14\tabcolsep) * \real{0.1250}}
  >{\raggedright\arraybackslash}p{(\linewidth - 14\tabcolsep) * \real{0.1250}}@{}}
\toprule\noalign{}
\begin{minipage}[b]{\linewidth}\raggedright
RUNSPEC
\end{minipage} & \begin{minipage}[b]{\linewidth}\raggedright
GRID
\end{minipage} & \begin{minipage}[b]{\linewidth}\raggedright
EDIT
\end{minipage} & \begin{minipage}[b]{\linewidth}\raggedright
PROPS
\end{minipage} & \begin{minipage}[b]{\linewidth}\raggedright
REGIONS
\end{minipage} & \begin{minipage}[b]{\linewidth}\raggedright
SOLUTION
\end{minipage} & \begin{minipage}[b]{\linewidth}\raggedright
SUMMARY
\end{minipage} & \begin{minipage}[b]{\linewidth}\raggedright
SCHEDULE
\end{minipage} \\
\midrule\noalign{}
\endhead
\bottomrule\noalign{}
\endlastfoot
\end{longtable}

\subsubsection{Description}\label{description-190}

This keyword sets various options for when multiple grid systems are
being used, as declared by the NUMRES keyword in the RUNSPEC section.
OPM Flow does not support multiple grid systems. The keyword is also
used to stipulate for radial grids if the completion of the circle in
the THETA direction should be implemented using non-neighbor
connections.

This keyword is not supported by OPM Flow but has no effect on the
results so it will be ignored.

\begin{longtable}[]{@{}
  >{\raggedright\arraybackslash}p{(\linewidth - 6\tabcolsep) * \real{0.2500}}
  >{\raggedright\arraybackslash}p{(\linewidth - 6\tabcolsep) * \real{0.2500}}
  >{\raggedright\arraybackslash}p{(\linewidth - 6\tabcolsep) * \real{0.2500}}
  >{\raggedright\arraybackslash}p{(\linewidth - 6\tabcolsep) * \real{0.2500}}@{}}
\toprule\noalign{}
\begin{minipage}[b]{\linewidth}\raggedright
No.
\end{minipage} & \begin{minipage}[b]{\linewidth}\raggedright
Name
\end{minipage} & \begin{minipage}[b]{\linewidth}\raggedright
Description
\end{minipage} & \begin{minipage}[b]{\linewidth}\raggedright
Default
\end{minipage} \\
\midrule\noalign{}
\endhead
\bottomrule\noalign{}
\endlastfoot
1 & K1 & A positive integer that defines the lower bound of the array in
the K-direction for the given grid system. & None \\
2 & K2 & A positive integer that defines the upper bound of the array in
the K-direction for the given grid system. & None \\
3 & COMPLETE & COMPLETE is a defined character string that determines
for radial grids if the circle should be completed in THETA direction,
and should be set to COMP to complete the circle, or INCOMP for not
completing the circle. & INCOMP \\
4 & CONNECT & A defined character string that declares how the reservoir
below should be connected to the given reservoir, and should be set to
JOIN to connect the two reservoirs by calculating the inter-reservoir
transmissibilities, or SEPARATE to isolate the reservoirs. & SEPARATE \\
5 & R1 & R1 is a positive integer defining the lower reservoir unit that
is is connected to the given reservoir unit. & Current Reservoir
Record \\
6 & R2 & R2 is a positive integer defining the upper reservoir unit that
is is connected to the given reservoir unit. & \\
Notes: & & & \\
\end{longtable}

\emph{Table 6.15: COORDSYS Keyword Description}

\subsubsection{Example}\label{example-140}

\begin{verbatim}
--
--       DEFINE COORDINATE GRID OPTIONS
--
--       K1     K2     COMP    CONNECT  LOWER  UPPER
--       Layer  Layer  CIRCLE  RES      RES    RES
COORDSYS
         1      1      COMP                                /
/
\end{verbatim}

The above example connects the circle in the THETA direction for the
RADIAL model, for when the number of grids have been set to one via the
NUMRES keyword in the SCHEDULE section.

\subsection{COPY -- Copy Array Data to Another
Array}\label{copy-copy-array-data-to-another-array}

\begin{longtable}[]{@{}
  >{\raggedright\arraybackslash}p{(\linewidth - 14\tabcolsep) * \real{0.1250}}
  >{\raggedright\arraybackslash}p{(\linewidth - 14\tabcolsep) * \real{0.1250}}
  >{\raggedright\arraybackslash}p{(\linewidth - 14\tabcolsep) * \real{0.1250}}
  >{\raggedright\arraybackslash}p{(\linewidth - 14\tabcolsep) * \real{0.1250}}
  >{\raggedright\arraybackslash}p{(\linewidth - 14\tabcolsep) * \real{0.1250}}
  >{\raggedright\arraybackslash}p{(\linewidth - 14\tabcolsep) * \real{0.1250}}
  >{\raggedright\arraybackslash}p{(\linewidth - 14\tabcolsep) * \real{0.1250}}
  >{\raggedright\arraybackslash}p{(\linewidth - 14\tabcolsep) * \real{0.1250}}@{}}
\toprule\noalign{}
\begin{minipage}[b]{\linewidth}\raggedright
RUNSPEC
\end{minipage} & \begin{minipage}[b]{\linewidth}\raggedright
GRID
\end{minipage} & \begin{minipage}[b]{\linewidth}\raggedright
EDIT
\end{minipage} & \begin{minipage}[b]{\linewidth}\raggedright
PROPS
\end{minipage} & \begin{minipage}[b]{\linewidth}\raggedright
REGIONS
\end{minipage} & \begin{minipage}[b]{\linewidth}\raggedright
SOLUTION
\end{minipage} & \begin{minipage}[b]{\linewidth}\raggedright
SUMMARY
\end{minipage} & \begin{minipage}[b]{\linewidth}\raggedright
SCHEDULE
\end{minipage} \\
\midrule\noalign{}
\endhead
\bottomrule\noalign{}
\endlastfoot
\end{longtable}

\subsubsection{Description}\label{description-191}

The COPY keyword copies an array (or part of an array) to another array
or part of an array. The arrays can be integer or real valued; however,
the arrays that can be operated on are dependent on which section the
COPY keyword is being applied in.

\begin{longtable}[]{@{}
  >{\raggedright\arraybackslash}p{(\linewidth - 6\tabcolsep) * \real{0.2500}}
  >{\raggedright\arraybackslash}p{(\linewidth - 6\tabcolsep) * \real{0.2500}}
  >{\raggedright\arraybackslash}p{(\linewidth - 6\tabcolsep) * \real{0.2500}}
  >{\raggedright\arraybackslash}p{(\linewidth - 6\tabcolsep) * \real{0.2500}}@{}}
\toprule\noalign{}
\begin{minipage}[b]{\linewidth}\raggedright
No.
\end{minipage} & \begin{minipage}[b]{\linewidth}\raggedright
Name
\end{minipage} & \begin{minipage}[b]{\linewidth}\raggedright
Description
\end{minipage} & \begin{minipage}[b]{\linewidth}\raggedright
Default
\end{minipage} \\
\midrule\noalign{}
\endhead
\bottomrule\noalign{}
\endlastfoot
1 & ARRAY-1 & A character string of up to eight characters in length
that defines the keyword identifying the array to be copied from. &
None \\
2 & ARRAY-2 & A character string of up to eight characters in length
that defines the keyword identifying the array to be copied to. &
None \\
3 & I1 & A positive integer that defines the lower bound of the array in
the I-direction to be modified must be greater than or equal 1 and less
than or equal to I2 and NX. & 1 \\
4 & I2 & A positive integer that defines the upper bound of the array in
the I-direction to be modified must be greater than or equal to II and
less than or equal to NX & NX \\
5 & J1 & A positive integer that defines the lower bound of the array in
the J-direction to be modified must be greater than or equal 1 and less
than or equal to J2 and NY. & 1 \\
6 & J2 & A positive integer that defines the upper bound of the array in
the J-direction to be modified must be greater than or equal to JI and
less than or equal to NY. & NY \\
7 & K1 & A positive integer that defines the lower bound of the array in
the K-direction to be modified must be greater than or equal to one and
less than or equal to K2 and NZ. & 1 \\
8 & K2 & A positive integer that defines the upper bound of the array in
the K-direction to be modified must be greater than or equal to KI and
less than or equal to NZ. & NZ \\
Notes: & & & \\
\end{longtable}

\emph{Table 6.16: COPY Keyword Description}

Examples of the arrays most commonly operated on in each section are
given in Table 6.17. Cells colored red indicate arrays that are not
supported by OPM Flow operations.

\begin{longtable}[]{@{}lllllll@{}}
\toprule\noalign{}
COPY Keyword and Variable Options by Section & & & & & & \\
\midrule\noalign{}
\endhead
\bottomrule\noalign{}
\endlastfoot
GRID & EDIT & PROPS & REGIONS & SOLUTION & SUMMARY & SCHEDULE \\
DX & DEPTH & SWL & ENDNUM & PRESSURE & & \\
DY & PORV & SWCR & EQLNUM & SWAT & & \\
DZ & TRANX & SWU & FIPNUM & SGAS & & \\
PERMX & TRANY & SGL & IMBNUM & RV & & \\
PERMY & TRANZ & SGCR & MISCNUM & RS & & \\
PERMZ & DIFFX & SGU & PVTNUM & TBLK & & \\
MULTX & DIFFY & KRW & ROCKNUM & GI & & \\
MULTY & DIFFZ & KRO & SATNUM & OILAPI & & \\
MULTZ & TRANR & KRG & WH2NUM & SALT & & \\
DR & TRANTHT & PCG & & GASCONC & & \\
DTHETA & DIFFR & PCW & & SOLVCONC & & \\
PERMR & DIFFTHT & & & SOLVFRAC & & \\
PERMTHT & & & & SFOAM & & \\
DZNET & & & & SPOLY & & \\
PORO & & & & & & \\
NTG & & & & & & \\
FLUXNUM & & & & & & \\
MULTNUM & & & & & & \\
MPFANUM & & & & & & \\
DIFFX & & & & & & \\
DIFFY & & & & & & \\
DIFFZ & & & & & & \\
DIFFR & & & & & & \\
DIFFTHT & & & & & & \\
\end{longtable}

\emph{Table 6.17: COPY Keyword Applicable Arrays by Section}

\subsubsection{Example}\label{example-141}

\begin{verbatim}
--
--       SOURCE      DESTIN.      ---------- BOX ---------
--                                I1  I2   J1  J2   K1  K2
COPY
         PERMX       PERMY        1*  1*   1*  1*   1*  1* / CREATE PERMY
         PERMX       PERMZ        1*  1*   1*  1*   1*  1* / CREATE PERMZ
/

--
--       ARRAY       CONSTANT     ---------- BOX ---------
--                                I1  I2   J1  J2   K1  K2
MULTIPLY
         PERMZ       0.50000      1*  1*   1*  1*   1*  1* / PERMZ * 0.5
/
\end{verbatim}

The above example copies PERMX array to the PERMY and PERMZ arrays in
the GRID section for all grid blocks in the model. The PERMZ array is
then multiplied by 0.5 for all grid blocks in the model.

\subsection{COPYBOX -- Copy Array Data Defined by a
Box}\label{copybox-copy-array-data-defined-by-a-box}

\begin{longtable}[]{@{}
  >{\raggedright\arraybackslash}p{(\linewidth - 14\tabcolsep) * \real{0.1250}}
  >{\raggedright\arraybackslash}p{(\linewidth - 14\tabcolsep) * \real{0.1250}}
  >{\raggedright\arraybackslash}p{(\linewidth - 14\tabcolsep) * \real{0.1250}}
  >{\raggedright\arraybackslash}p{(\linewidth - 14\tabcolsep) * \real{0.1250}}
  >{\raggedright\arraybackslash}p{(\linewidth - 14\tabcolsep) * \real{0.1250}}
  >{\raggedright\arraybackslash}p{(\linewidth - 14\tabcolsep) * \real{0.1250}}
  >{\raggedright\arraybackslash}p{(\linewidth - 14\tabcolsep) * \real{0.1250}}
  >{\raggedright\arraybackslash}p{(\linewidth - 14\tabcolsep) * \real{0.1250}}@{}}
\toprule\noalign{}
\begin{minipage}[b]{\linewidth}\raggedright
RUNSPEC
\end{minipage} & \begin{minipage}[b]{\linewidth}\raggedright
GRID
\end{minipage} & \begin{minipage}[b]{\linewidth}\raggedright
EDIT
\end{minipage} & \begin{minipage}[b]{\linewidth}\raggedright
PROPS
\end{minipage} & \begin{minipage}[b]{\linewidth}\raggedright
REGIONS
\end{minipage} & \begin{minipage}[b]{\linewidth}\raggedright
SOLUTION
\end{minipage} & \begin{minipage}[b]{\linewidth}\raggedright
SUMMARY
\end{minipage} & \begin{minipage}[b]{\linewidth}\raggedright
SCHEDULE
\end{minipage} \\
\midrule\noalign{}
\endhead
\bottomrule\noalign{}
\endlastfoot
\end{longtable}

\subsubsection{Description}\label{description-192}

The COPYBOX keyword copies an array (or part of an array) to another
part of the same array. The array can be real or integer depending on
the array type; however, the array that can be operated on is dependent
on which section the COPYBOX keyword is being used.

This keyword is not supported by OPM Flow but would change the results
if supported so the simulation will be stopped.

\begin{longtable}[]{@{}
  >{\raggedright\arraybackslash}p{(\linewidth - 6\tabcolsep) * \real{0.2500}}
  >{\raggedright\arraybackslash}p{(\linewidth - 6\tabcolsep) * \real{0.2500}}
  >{\raggedright\arraybackslash}p{(\linewidth - 6\tabcolsep) * \real{0.2500}}
  >{\raggedright\arraybackslash}p{(\linewidth - 6\tabcolsep) * \real{0.2500}}@{}}
\toprule\noalign{}
\begin{minipage}[b]{\linewidth}\raggedright
No.
\end{minipage} & \begin{minipage}[b]{\linewidth}\raggedright
Name
\end{minipage} & \begin{minipage}[b]{\linewidth}\raggedright
Description
\end{minipage} & \begin{minipage}[b]{\linewidth}\raggedright
Default
\end{minipage} \\
\midrule\noalign{}
\endhead
\bottomrule\noalign{}
\endlastfoot
1 & ARRAY-1 & The name of the array to be copied This is the keyword
name identifying the property and is up to eight characters in length
and enclosed in quotes. & None \\
2 & I1 & A positive integer that defines the
\hyperref[REF_HEADING_KEYWORD_SOURCE_12_3]{SOURCE} lower bound of the
array in the I-direction to be modified must be greater than or equal 1
and less than or equal to I2 and NX. & 1 \\
3 & I2 & A positive integer that defines the
\hyperref[REF_HEADING_KEYWORD_SOURCE_12_3]{SOURCE} upper bound of the
array in the I-direction to be modified must be greater than or equal to
II and less than or equal to NX & NX \\
4 & J1 & A positive integer that defines the
\hyperref[REF_HEADING_KEYWORD_SOURCE_12_3]{SOURCE} lower bound of the
array in the J-direction to be modified must be greater than or equal 1
and less than or equal to J2 and NY. & 1 \\
5 & J2 & A positive integer that defines the
\hyperref[REF_HEADING_KEYWORD_SOURCE_12_3]{SOURCE} upper bound of the
array in the J-direction to be modified must be greater than or equal to
JI and less than or equal to NY. & NY \\
6 & K1 & A positive integer that defines the
\hyperref[REF_HEADING_KEYWORD_SOURCE_12_3]{SOURCE} lower bound of the
array in the K-direction to be modified must be greater than or equal to
one and less than or equal to K2 and NZ. & 1 \\
7 & K2 & A positive integer that defines the
\hyperref[REF_HEADING_KEYWORD_SOURCE_12_3]{SOURCE} upper bound of the
array in the K-direction to be modified must be greater than or equal to
KI and less than or equal to NZ. & NZ \\
8 & I3 & A positive integer that defines the DESTINATION lower bound of
the array in the I-direction to be modified must be greater than or
equal 1 and less than or equal to I2 and NX. & 1 \\
9 & I4 & A positive integer that defines the DESTINATION upper bound of
the array in the I-direction to be modified must be greater than or
equal to II and less than or equal to NX & NX \\
10 & J3 & A positive integer that defines the DESTINATION lower bound of
the array in the J-direction to be modified must be greater than or
equal 1 and less than or equal to J2 and NY. & 1 \\
11 & J4 & A positive integer that defines the DESTINATION upper bound of
the array in the J-direction to be modified must be greater than or
equal to JI and less than or equal to NY. & NY \\
12 & K3 & A positive integer that defines the DESTINATION lower bound of
the array in the K-direction to be modified must be greater than or
equal to one and less than or equal to K2 and NZ. & 1 \\
13 & K4 & A positive integer that defines the DESTINATION upper bound of
the array in the K-direction to be modified must be greater than or
equal to KI and less than or equal to NZ. & NZ \\
Notes: & & & \\
\end{longtable}

\emph{Table 6.18: COPYBOX Keyword Description}

Note that the \hyperref[REF_HEADING_KEYWORD_SOURCE_12_3]{SOURCE} and
DESTINATION arrays must be of the same size in all dimensions and the
applicable arrays for each section are defined in Table 6.19.

\begin{longtable}[]{@{}lllllll@{}}
\toprule\noalign{}
COPYBOX Keyword and Variable Options by Section & & & & & & \\
\midrule\noalign{}
\endhead
\bottomrule\noalign{}
\endlastfoot
GRID & EDIT & PROPS & REGIONS & SOLUTION & SUMMARY & SCHEDULE \\
DX & & SWL & ENDNUM & & & \\
DY & & SWCR & EQLNUM & & & \\
DZ & & SWU & FIPNUM & & & \\
PERMX & & SGL & IMBNUM & & & \\
PERMY & & SGCR & MISCNUM & & & \\
PERMZ & & SGU & PVTNUM & & & \\
MULTX & & KRW & ROCKNUM & & & \\
MULTY & & KRO & SATNUM & & & \\
MULTZ & & KRG & WH2NUM & & & \\
DR & & PCG & & & & \\
DTHETA & & PCW & & & & \\
PERMR & & & & & & \\
PERMTHT & & & & & & \\
DZNET & & & & & & \\
PORO & & & & & & \\
NTG & & & & & & \\
FLUXNUM & & & & & & \\
MULTNUM & & & & & & \\
MPFANUM & & & & & & \\
DIFFX & & & & & & \\
DIFFY & & & & & & \\
DIFFZ & & & & & & \\
DIFFR & & & & & & \\
DIFFTHT & & & & & & \\
\end{longtable}

\emph{Table 6.19: COPYBOX Keyword Applicable Arrays by Section}

\subsubsection{Example}\label{example-142}

\begin{verbatim}
--
--       SOURCE   ------ SOURCE BOX ------   ---- DESTINATION BOX ----
--       ARRAY    I1  I2   J1  J2   K1  K2   I1  I2   J1  J2   K1  K2
COPYBOX
          PORO    1*  1*   1*  1*   12  14   1*  1*   1*  1*   15  17 / PORO
          PERMX   1*  1*   1*  1*   12  14   1*  1*   1*  1*   15  17 / PERMX
/
\end{verbatim}

The above example copies all the PORO and PERMX values in layers 12 to
14 to layers 15 and 17.

\subsection{COPYREG -- Copy an Array to Another Array based on a Region
Number}\label{copyreg-copy-an-array-to-another-array-based-on-a-region-number}

\begin{longtable}[]{@{}
  >{\raggedright\arraybackslash}p{(\linewidth - 14\tabcolsep) * \real{0.1250}}
  >{\raggedright\arraybackslash}p{(\linewidth - 14\tabcolsep) * \real{0.1250}}
  >{\raggedright\arraybackslash}p{(\linewidth - 14\tabcolsep) * \real{0.1250}}
  >{\raggedright\arraybackslash}p{(\linewidth - 14\tabcolsep) * \real{0.1250}}
  >{\raggedright\arraybackslash}p{(\linewidth - 14\tabcolsep) * \real{0.1250}}
  >{\raggedright\arraybackslash}p{(\linewidth - 14\tabcolsep) * \real{0.1250}}
  >{\raggedright\arraybackslash}p{(\linewidth - 14\tabcolsep) * \real{0.1250}}
  >{\raggedright\arraybackslash}p{(\linewidth - 14\tabcolsep) * \real{0.1250}}@{}}
\toprule\noalign{}
\begin{minipage}[b]{\linewidth}\raggedright
RUNSPEC
\end{minipage} & \begin{minipage}[b]{\linewidth}\raggedright
GRID
\end{minipage} & \begin{minipage}[b]{\linewidth}\raggedright
EDIT
\end{minipage} & \begin{minipage}[b]{\linewidth}\raggedright
PROPS
\end{minipage} & \begin{minipage}[b]{\linewidth}\raggedright
REGIONS
\end{minipage} & \begin{minipage}[b]{\linewidth}\raggedright
SOLUTION
\end{minipage} & \begin{minipage}[b]{\linewidth}\raggedright
SUMMARY
\end{minipage} & \begin{minipage}[b]{\linewidth}\raggedright
SCHEDULE
\end{minipage} \\
\midrule\noalign{}
\endhead
\bottomrule\noalign{}
\endlastfoot
\end{longtable}

\subsubsection{Description}\label{description-193}

The COPYREG keyword copies a specified array or part of an array based
on cells with a specific region number to another array. The region
number array can be FLUXNUM, MULTNUM or OPERNUM and these arrays must be
defined and be available before the COPYREG keyword is read by the
simulator. The property arrays can be real or integer valued depending
on the property array type; however, the property arrays that can be
operated on are dependent on which section the COPYREG keyword is being
applied in.

\begin{longtable}[]{@{}
  >{\raggedright\arraybackslash}p{(\linewidth - 6\tabcolsep) * \real{0.2500}}
  >{\raggedright\arraybackslash}p{(\linewidth - 6\tabcolsep) * \real{0.2500}}
  >{\raggedright\arraybackslash}p{(\linewidth - 6\tabcolsep) * \real{0.2500}}
  >{\raggedright\arraybackslash}p{(\linewidth - 6\tabcolsep) * \real{0.2500}}@{}}
\toprule\noalign{}
\begin{minipage}[b]{\linewidth}\raggedright
No.
\end{minipage} & \begin{minipage}[b]{\linewidth}\raggedright
Name
\end{minipage} & \begin{minipage}[b]{\linewidth}\raggedright
Description
\end{minipage} & \begin{minipage}[b]{\linewidth}\raggedright
Default
\end{minipage} \\
\midrule\noalign{}
\endhead
\bottomrule\noalign{}
\endlastfoot
1 & ARRAY-1 & A character string of up to eight characters in length
that defines the keyword identifying the array to be copied from. &
None \\
2 & ARRAY-2 & A character string of up to eight characters in length
that defines the keyword identifying the array to be copied to. &
None \\
3 & REGION NUMBER & Integer REGION NUMBER is the region for which the
array data in (1) should be copied to array data in (2). & None \\
4 & REGION ARRAY & The REGION ARRAY to use for selecting the REGION
NUMBER in (3) for selecting the data to be copied. REGION ARRAY can have
the following values: & M \\
Notes: & & & \\
\end{longtable}

\emph{Table 6.20: COPYREG Keyword Description}

Examples of the arrays most commonly operated on in each section are
given in Table 6.21. Cells colored red indicate arrays that are not
supported by OPM Flow operations.

\begin{longtable}[]{@{}lllllll@{}}
\toprule\noalign{}
COPYREG Keyword and Variable Options by Section & & & & & & \\
\midrule\noalign{}
\endhead
\bottomrule\noalign{}
\endlastfoot
GRID & EDIT & PROPS & REGIONS & SOLUTION & SUMMARY & SCHEDULE \\
DX & DEPTH & SWL & ENDNUM & PRESSURE & & \\
DY & PORV & SWCR & EQLNUM & SWAT & & \\
DZ & TRANX & SWU & FIPNUM & SGAS & & \\
PERMX & TRANY & SGL & IMBNUM & RV & & \\
PERMY & TRANZ & SGCR & MISCNUM & RS & & \\
PERMZ & DIFFX & SGU & PVTNUM & TBLK & & \\
MULTX & DIFFY & KRW & ROCKNUM & GI & & \\
MULTY & DIFFZ & KRO & SATNUM & OILAPI & & \\
MULTZ & TRANR & KRG & WH2NUM & SALT & & \\
DR & TRANTHT & PCG & & GASCONC & & \\
DTHETA & DIFFR & PCW & & SOLVCONC & & \\
PERMR & DIFFTHT & & & SOLVFRAC & & \\
PERMTHT & & & & SFOAM & & \\
DZNET & & & & SPOLY & & \\
PORO & & & & & & \\
NTG & & & & & & \\
FLUXNUM & & & & & & \\
MULTNUM & & & & & & \\
MPFANUM & & & & & & \\
DIFFX & & & & & & \\
DIFFY & & & & & & \\
DIFFZ & & & & & & \\
DIFFR & & & & & & \\
DIFFTHT & & & & & & \\
\end{longtable}

\emph{Table 6.21: COPYREG Keyword Applicable Arrays by Section}

\subsubsection{Example}\label{example-143}

\begin{verbatim}
--
--       COPY AN ARRAY TO ANOTHER ARRAY BASED ON A REGION NUMBER
--
--       ARRAY     ARRAY     REGION   REGION ARRAY
--       FROM      TO        NUMBER    M / F / O
COPYREG
         PERMX     PERMY     1         M                   / COPY PERMX TO PERMY
         PERMX     PERMZ     1         M                   / COPY PERMX TO PERMZ
/
--
--       NOW RESET PERMZ BASED ON THE MULTNUM REGION NUMBER
--
--
--       MULTIPLY AN ARRAY BY A CONSTANT BASED ON A REGION NUMBER
--
--       ARRAY     CONSTANT  REGION   REGION ARRAY
--                 VALUE     NUMBER    M / F / O
MULTIREG
         PERMX     0.95      1         M                   /
/
\end{verbatim}

The above example first copies the PERMX property array for region
number one to the PERMY and PERMZ property arrays for region one using
the MULTNUM array to define the region numbers. After which PERMZ
property array for region one is multiplied by 0.5 using the MULTIREG
keyword.

\subsection{CRITPERM -- Define Minimum Permeability for Vertical
Equilibrium Grid Cell
Compression}\label{critperm-define-minimum-permeability-for-vertical-equilibrium-grid-cell-compression}

\begin{longtable}[]{@{}
  >{\raggedright\arraybackslash}p{(\linewidth - 14\tabcolsep) * \real{0.1250}}
  >{\raggedright\arraybackslash}p{(\linewidth - 14\tabcolsep) * \real{0.1250}}
  >{\raggedright\arraybackslash}p{(\linewidth - 14\tabcolsep) * \real{0.1250}}
  >{\raggedright\arraybackslash}p{(\linewidth - 14\tabcolsep) * \real{0.1250}}
  >{\raggedright\arraybackslash}p{(\linewidth - 14\tabcolsep) * \real{0.1250}}
  >{\raggedright\arraybackslash}p{(\linewidth - 14\tabcolsep) * \real{0.1250}}
  >{\raggedright\arraybackslash}p{(\linewidth - 14\tabcolsep) * \real{0.1250}}
  >{\raggedright\arraybackslash}p{(\linewidth - 14\tabcolsep) * \real{0.1250}}@{}}
\toprule\noalign{}
\begin{minipage}[b]{\linewidth}\raggedright
RUNSPEC
\end{minipage} & \begin{minipage}[b]{\linewidth}\raggedright
GRID
\end{minipage} & \begin{minipage}[b]{\linewidth}\raggedright
EDIT
\end{minipage} & \begin{minipage}[b]{\linewidth}\raggedright
PROPS
\end{minipage} & \begin{minipage}[b]{\linewidth}\raggedright
REGIONS
\end{minipage} & \begin{minipage}[b]{\linewidth}\raggedright
SOLUTION
\end{minipage} & \begin{minipage}[b]{\linewidth}\raggedright
SUMMARY
\end{minipage} & \begin{minipage}[b]{\linewidth}\raggedright
SCHEDULE
\end{minipage} \\
\midrule\noalign{}
\endhead
\bottomrule\noalign{}
\endlastfoot
\end{longtable}

\subsubsection{Description}\label{description-194}

The CRITPERM keyword is used to define the minimum permeability for
Vertical Equilibrium(``VE'') grid cell compression, for when the
Vertical Equilibrium formulation has been activated by the VE keyword in
the RUNSPEC section.

This keyword is not supported by OPM Flow but would change the results
if supported so the simulation will be stopped.

\subsection{DEBUG -- Define the Debug Data to be Printed to
File}\label{debug-define-the-debug-data-to-be-printed-to-file-2}

This keyword defines the debug data to be written to the debug file
(*.DBG). This keyword is not supported by OPM Flow but has no effect on
the results so it will be ignored.

See DEBUG -- Define the Debug Data to be Printed to File in the GLOBAL
section for a full description.

\subsection{DIFFMR- -- Define Grid Block Negative Radial Direction
Diffusivity
Multipliers}\label{diffmr--define-grid-block-negative-radial-direction-diffusivity-multipliers}

\begin{longtable}[]{@{}
  >{\raggedright\arraybackslash}p{(\linewidth - 14\tabcolsep) * \real{0.1250}}
  >{\raggedright\arraybackslash}p{(\linewidth - 14\tabcolsep) * \real{0.1250}}
  >{\raggedright\arraybackslash}p{(\linewidth - 14\tabcolsep) * \real{0.1250}}
  >{\raggedright\arraybackslash}p{(\linewidth - 14\tabcolsep) * \real{0.1250}}
  >{\raggedright\arraybackslash}p{(\linewidth - 14\tabcolsep) * \real{0.1250}}
  >{\raggedright\arraybackslash}p{(\linewidth - 14\tabcolsep) * \real{0.1250}}
  >{\raggedright\arraybackslash}p{(\linewidth - 14\tabcolsep) * \real{0.1250}}
  >{\raggedright\arraybackslash}p{(\linewidth - 14\tabcolsep) * \real{0.1250}}@{}}
\toprule\noalign{}
\begin{minipage}[b]{\linewidth}\raggedright
RUNSPEC
\end{minipage} & \begin{minipage}[b]{\linewidth}\raggedright
GRID
\end{minipage} & \begin{minipage}[b]{\linewidth}\raggedright
EDIT
\end{minipage} & \begin{minipage}[b]{\linewidth}\raggedright
PROPS
\end{minipage} & \begin{minipage}[b]{\linewidth}\raggedright
REGIONS
\end{minipage} & \begin{minipage}[b]{\linewidth}\raggedright
SOLUTION
\end{minipage} & \begin{minipage}[b]{\linewidth}\raggedright
SUMMARY
\end{minipage} & \begin{minipage}[b]{\linewidth}\raggedright
SCHEDULE
\end{minipage} \\
\midrule\noalign{}
\endhead
\bottomrule\noalign{}
\endlastfoot
\end{longtable}

\subsubsection{Description}\label{description-195}

The DIFFMR- keyword defines the negative radial direction diffusivity
multipliers for cells in the current input box for when the Diffusivity
option has been activated by the DIFFUSE keyword in the RUNSPEC section.

This keyword is not supported by OPM Flow but would change the results
if supported so the simulation will be stopped.

\subsection{DIFFMR -- Define Grid Block Radial Direction Diffusivity
Multipliers}\label{diffmr-define-grid-block-radial-direction-diffusivity-multipliers}

\begin{longtable}[]{@{}
  >{\raggedright\arraybackslash}p{(\linewidth - 14\tabcolsep) * \real{0.1250}}
  >{\raggedright\arraybackslash}p{(\linewidth - 14\tabcolsep) * \real{0.1250}}
  >{\raggedright\arraybackslash}p{(\linewidth - 14\tabcolsep) * \real{0.1250}}
  >{\raggedright\arraybackslash}p{(\linewidth - 14\tabcolsep) * \real{0.1250}}
  >{\raggedright\arraybackslash}p{(\linewidth - 14\tabcolsep) * \real{0.1250}}
  >{\raggedright\arraybackslash}p{(\linewidth - 14\tabcolsep) * \real{0.1250}}
  >{\raggedright\arraybackslash}p{(\linewidth - 14\tabcolsep) * \real{0.1250}}
  >{\raggedright\arraybackslash}p{(\linewidth - 14\tabcolsep) * \real{0.1250}}@{}}
\toprule\noalign{}
\begin{minipage}[b]{\linewidth}\raggedright
RUNSPEC
\end{minipage} & \begin{minipage}[b]{\linewidth}\raggedright
GRID
\end{minipage} & \begin{minipage}[b]{\linewidth}\raggedright
EDIT
\end{minipage} & \begin{minipage}[b]{\linewidth}\raggedright
PROPS
\end{minipage} & \begin{minipage}[b]{\linewidth}\raggedright
REGIONS
\end{minipage} & \begin{minipage}[b]{\linewidth}\raggedright
SOLUTION
\end{minipage} & \begin{minipage}[b]{\linewidth}\raggedright
SUMMARY
\end{minipage} & \begin{minipage}[b]{\linewidth}\raggedright
SCHEDULE
\end{minipage} \\
\midrule\noalign{}
\endhead
\bottomrule\noalign{}
\endlastfoot
\end{longtable}

\subsubsection{Description}\label{description-196}

The DIFFMR keyword defines the radial direction diffusivity multipliers
for cells in the current input box for when the Diffusivity option has
been activated by the DIFFUSE keyword in the RUNSPEC section.

This keyword is not supported by OPM Flow but would change the results
if supported so the simulation will be stopped.

\subsection{DIFFMTH- -- Define Grid Block Negative Theta Direction
Diffusivity
Multipliers}\label{diffmth--define-grid-block-negative-theta-direction-diffusivity-multipliers}

\begin{longtable}[]{@{}
  >{\raggedright\arraybackslash}p{(\linewidth - 14\tabcolsep) * \real{0.1250}}
  >{\raggedright\arraybackslash}p{(\linewidth - 14\tabcolsep) * \real{0.1250}}
  >{\raggedright\arraybackslash}p{(\linewidth - 14\tabcolsep) * \real{0.1250}}
  >{\raggedright\arraybackslash}p{(\linewidth - 14\tabcolsep) * \real{0.1250}}
  >{\raggedright\arraybackslash}p{(\linewidth - 14\tabcolsep) * \real{0.1250}}
  >{\raggedright\arraybackslash}p{(\linewidth - 14\tabcolsep) * \real{0.1250}}
  >{\raggedright\arraybackslash}p{(\linewidth - 14\tabcolsep) * \real{0.1250}}
  >{\raggedright\arraybackslash}p{(\linewidth - 14\tabcolsep) * \real{0.1250}}@{}}
\toprule\noalign{}
\begin{minipage}[b]{\linewidth}\raggedright
RUNSPEC
\end{minipage} & \begin{minipage}[b]{\linewidth}\raggedright
GRID
\end{minipage} & \begin{minipage}[b]{\linewidth}\raggedright
EDIT
\end{minipage} & \begin{minipage}[b]{\linewidth}\raggedright
PROPS
\end{minipage} & \begin{minipage}[b]{\linewidth}\raggedright
REGIONS
\end{minipage} & \begin{minipage}[b]{\linewidth}\raggedright
SOLUTION
\end{minipage} & \begin{minipage}[b]{\linewidth}\raggedright
SUMMARY
\end{minipage} & \begin{minipage}[b]{\linewidth}\raggedright
SCHEDULE
\end{minipage} \\
\midrule\noalign{}
\endhead
\bottomrule\noalign{}
\endlastfoot
\end{longtable}

\subsubsection{Description}\label{description-197}

The DIFFMR- keyword defines the negative theta direction diffusivity
multipliers for cells in the current input box for when the Diffusivity
option has been activated by the DIFFUSE keyword in the RUNSPEC section.

This keyword is not supported by OPM Flow but would change the results
if supported so the simulation will be stopped.

\subsection{DIFFMTHT -- Define Grid Block Theta Direction Diffusivity
Multipliers}\label{diffmtht-define-grid-block-theta-direction-diffusivity-multipliers}

\begin{longtable}[]{@{}
  >{\raggedright\arraybackslash}p{(\linewidth - 14\tabcolsep) * \real{0.1250}}
  >{\raggedright\arraybackslash}p{(\linewidth - 14\tabcolsep) * \real{0.1250}}
  >{\raggedright\arraybackslash}p{(\linewidth - 14\tabcolsep) * \real{0.1250}}
  >{\raggedright\arraybackslash}p{(\linewidth - 14\tabcolsep) * \real{0.1250}}
  >{\raggedright\arraybackslash}p{(\linewidth - 14\tabcolsep) * \real{0.1250}}
  >{\raggedright\arraybackslash}p{(\linewidth - 14\tabcolsep) * \real{0.1250}}
  >{\raggedright\arraybackslash}p{(\linewidth - 14\tabcolsep) * \real{0.1250}}
  >{\raggedright\arraybackslash}p{(\linewidth - 14\tabcolsep) * \real{0.1250}}@{}}
\toprule\noalign{}
\begin{minipage}[b]{\linewidth}\raggedright
RUNSPEC
\end{minipage} & \begin{minipage}[b]{\linewidth}\raggedright
GRID
\end{minipage} & \begin{minipage}[b]{\linewidth}\raggedright
EDIT
\end{minipage} & \begin{minipage}[b]{\linewidth}\raggedright
PROPS
\end{minipage} & \begin{minipage}[b]{\linewidth}\raggedright
REGIONS
\end{minipage} & \begin{minipage}[b]{\linewidth}\raggedright
SOLUTION
\end{minipage} & \begin{minipage}[b]{\linewidth}\raggedright
SUMMARY
\end{minipage} & \begin{minipage}[b]{\linewidth}\raggedright
SCHEDULE
\end{minipage} \\
\midrule\noalign{}
\endhead
\bottomrule\noalign{}
\endlastfoot
\end{longtable}

\subsubsection{Description}\label{description-198}

The DIFFMTHT keyword defines the theta direction diffusivity multipliers
for cells in the current input box for when the Diffusivity option has
been activated by the DIFFUSE keyword in the RUNSPEC section.

This keyword is not supported by OPM Flow but would change the results
if supported so the simulation will be stopped.

\subsection{DIFFMX- -- Define Grid Block Negative X-Direction
Diffusivity
Multipliers}\label{diffmx--define-grid-block-negative-x-direction-diffusivity-multipliers}

\begin{longtable}[]{@{}
  >{\raggedright\arraybackslash}p{(\linewidth - 14\tabcolsep) * \real{0.1250}}
  >{\raggedright\arraybackslash}p{(\linewidth - 14\tabcolsep) * \real{0.1250}}
  >{\raggedright\arraybackslash}p{(\linewidth - 14\tabcolsep) * \real{0.1250}}
  >{\raggedright\arraybackslash}p{(\linewidth - 14\tabcolsep) * \real{0.1250}}
  >{\raggedright\arraybackslash}p{(\linewidth - 14\tabcolsep) * \real{0.1250}}
  >{\raggedright\arraybackslash}p{(\linewidth - 14\tabcolsep) * \real{0.1250}}
  >{\raggedright\arraybackslash}p{(\linewidth - 14\tabcolsep) * \real{0.1250}}
  >{\raggedright\arraybackslash}p{(\linewidth - 14\tabcolsep) * \real{0.1250}}@{}}
\toprule\noalign{}
\begin{minipage}[b]{\linewidth}\raggedright
RUNSPEC
\end{minipage} & \begin{minipage}[b]{\linewidth}\raggedright
GRID
\end{minipage} & \begin{minipage}[b]{\linewidth}\raggedright
EDIT
\end{minipage} & \begin{minipage}[b]{\linewidth}\raggedright
PROPS
\end{minipage} & \begin{minipage}[b]{\linewidth}\raggedright
REGIONS
\end{minipage} & \begin{minipage}[b]{\linewidth}\raggedright
SOLUTION
\end{minipage} & \begin{minipage}[b]{\linewidth}\raggedright
SUMMARY
\end{minipage} & \begin{minipage}[b]{\linewidth}\raggedright
SCHEDULE
\end{minipage} \\
\midrule\noalign{}
\endhead
\bottomrule\noalign{}
\endlastfoot
\end{longtable}

\subsubsection{Description}\label{description-199}

The DIFFMX- keyword defines the negative x-direction diffusivity
multipliers for cells in the current input box for when the Diffusivity
option has been activated by the DIFFUSE keyword in the RUNSPEC section.

This keyword is not supported by OPM Flow but would change the results
if supported so the simulation will be stopped.

\subsection{DIFFMX -- Define Grid Block X-Direction Diffusivity
Multipliers}\label{diffmx-define-grid-block-x-direction-diffusivity-multipliers}

\begin{longtable}[]{@{}
  >{\raggedright\arraybackslash}p{(\linewidth - 14\tabcolsep) * \real{0.1250}}
  >{\raggedright\arraybackslash}p{(\linewidth - 14\tabcolsep) * \real{0.1250}}
  >{\raggedright\arraybackslash}p{(\linewidth - 14\tabcolsep) * \real{0.1250}}
  >{\raggedright\arraybackslash}p{(\linewidth - 14\tabcolsep) * \real{0.1250}}
  >{\raggedright\arraybackslash}p{(\linewidth - 14\tabcolsep) * \real{0.1250}}
  >{\raggedright\arraybackslash}p{(\linewidth - 14\tabcolsep) * \real{0.1250}}
  >{\raggedright\arraybackslash}p{(\linewidth - 14\tabcolsep) * \real{0.1250}}
  >{\raggedright\arraybackslash}p{(\linewidth - 14\tabcolsep) * \real{0.1250}}@{}}
\toprule\noalign{}
\begin{minipage}[b]{\linewidth}\raggedright
RUNSPEC
\end{minipage} & \begin{minipage}[b]{\linewidth}\raggedright
GRID
\end{minipage} & \begin{minipage}[b]{\linewidth}\raggedright
EDIT
\end{minipage} & \begin{minipage}[b]{\linewidth}\raggedright
PROPS
\end{minipage} & \begin{minipage}[b]{\linewidth}\raggedright
REGIONS
\end{minipage} & \begin{minipage}[b]{\linewidth}\raggedright
SOLUTION
\end{minipage} & \begin{minipage}[b]{\linewidth}\raggedright
SUMMARY
\end{minipage} & \begin{minipage}[b]{\linewidth}\raggedright
SCHEDULE
\end{minipage} \\
\midrule\noalign{}
\endhead
\bottomrule\noalign{}
\endlastfoot
\end{longtable}

\subsubsection{Description}\label{description-200}

The DIFFMX keyword defines the x-direction diffusivity multipliers for
cells in the current input box for when the Diffusivity option has been
activated by the DIFFUSE keyword in the RUNSPEC section.

This keyword is not supported by OPM Flow but would change the results
if supported so the simulation will be stopped.

\subsection{DIFFMY- -- Define Grid Block Negative Y-Direction
Diffusivity
Multipliers}\label{diffmy--define-grid-block-negative-y-direction-diffusivity-multipliers}

\begin{longtable}[]{@{}
  >{\raggedright\arraybackslash}p{(\linewidth - 14\tabcolsep) * \real{0.1250}}
  >{\raggedright\arraybackslash}p{(\linewidth - 14\tabcolsep) * \real{0.1250}}
  >{\raggedright\arraybackslash}p{(\linewidth - 14\tabcolsep) * \real{0.1250}}
  >{\raggedright\arraybackslash}p{(\linewidth - 14\tabcolsep) * \real{0.1250}}
  >{\raggedright\arraybackslash}p{(\linewidth - 14\tabcolsep) * \real{0.1250}}
  >{\raggedright\arraybackslash}p{(\linewidth - 14\tabcolsep) * \real{0.1250}}
  >{\raggedright\arraybackslash}p{(\linewidth - 14\tabcolsep) * \real{0.1250}}
  >{\raggedright\arraybackslash}p{(\linewidth - 14\tabcolsep) * \real{0.1250}}@{}}
\toprule\noalign{}
\begin{minipage}[b]{\linewidth}\raggedright
RUNSPEC
\end{minipage} & \begin{minipage}[b]{\linewidth}\raggedright
GRID
\end{minipage} & \begin{minipage}[b]{\linewidth}\raggedright
EDIT
\end{minipage} & \begin{minipage}[b]{\linewidth}\raggedright
PROPS
\end{minipage} & \begin{minipage}[b]{\linewidth}\raggedright
REGIONS
\end{minipage} & \begin{minipage}[b]{\linewidth}\raggedright
SOLUTION
\end{minipage} & \begin{minipage}[b]{\linewidth}\raggedright
SUMMARY
\end{minipage} & \begin{minipage}[b]{\linewidth}\raggedright
SCHEDULE
\end{minipage} \\
\midrule\noalign{}
\endhead
\bottomrule\noalign{}
\endlastfoot
\end{longtable}

\subsubsection{Description}\label{description-201}

The DIFFMY- keyword defines the negative y-direction diffusivity
multipliers for cells in the current input box for when the Diffusivity
option has been activated by the DIFFUSE keyword in the RUNSPEC section.

This keyword is not supported by OPM Flow but would change the results
if supported so the simulation will be stopped.

\subsection{DIFFMY -- Define Grid Block Y-Direction Diffusivity
Multipliers}\label{diffmy-define-grid-block-y-direction-diffusivity-multipliers}

\begin{longtable}[]{@{}
  >{\raggedright\arraybackslash}p{(\linewidth - 14\tabcolsep) * \real{0.1250}}
  >{\raggedright\arraybackslash}p{(\linewidth - 14\tabcolsep) * \real{0.1250}}
  >{\raggedright\arraybackslash}p{(\linewidth - 14\tabcolsep) * \real{0.1250}}
  >{\raggedright\arraybackslash}p{(\linewidth - 14\tabcolsep) * \real{0.1250}}
  >{\raggedright\arraybackslash}p{(\linewidth - 14\tabcolsep) * \real{0.1250}}
  >{\raggedright\arraybackslash}p{(\linewidth - 14\tabcolsep) * \real{0.1250}}
  >{\raggedright\arraybackslash}p{(\linewidth - 14\tabcolsep) * \real{0.1250}}
  >{\raggedright\arraybackslash}p{(\linewidth - 14\tabcolsep) * \real{0.1250}}@{}}
\toprule\noalign{}
\begin{minipage}[b]{\linewidth}\raggedright
RUNSPEC
\end{minipage} & \begin{minipage}[b]{\linewidth}\raggedright
GRID
\end{minipage} & \begin{minipage}[b]{\linewidth}\raggedright
EDIT
\end{minipage} & \begin{minipage}[b]{\linewidth}\raggedright
PROPS
\end{minipage} & \begin{minipage}[b]{\linewidth}\raggedright
REGIONS
\end{minipage} & \begin{minipage}[b]{\linewidth}\raggedright
SOLUTION
\end{minipage} & \begin{minipage}[b]{\linewidth}\raggedright
SUMMARY
\end{minipage} & \begin{minipage}[b]{\linewidth}\raggedright
SCHEDULE
\end{minipage} \\
\midrule\noalign{}
\endhead
\bottomrule\noalign{}
\endlastfoot
\end{longtable}

\subsubsection{Description}\label{description-202}

The DIFFMY keyword defines the y-direction diffusivity multipliers for
cells in the current input box for when the Diffusivity option has been
activated by the DIFFUSE keyword in the RUNSPEC section.

This keyword is not supported by OPM Flow but would change the results
if supported so the simulation will be stopped.

\subsection{DIFFMZ- -- Define Grid Block Negative Z-Direction
Diffusivity
Multipliers}\label{diffmz--define-grid-block-negative-z-direction-diffusivity-multipliers}

\begin{longtable}[]{@{}
  >{\raggedright\arraybackslash}p{(\linewidth - 14\tabcolsep) * \real{0.1250}}
  >{\raggedright\arraybackslash}p{(\linewidth - 14\tabcolsep) * \real{0.1250}}
  >{\raggedright\arraybackslash}p{(\linewidth - 14\tabcolsep) * \real{0.1250}}
  >{\raggedright\arraybackslash}p{(\linewidth - 14\tabcolsep) * \real{0.1250}}
  >{\raggedright\arraybackslash}p{(\linewidth - 14\tabcolsep) * \real{0.1250}}
  >{\raggedright\arraybackslash}p{(\linewidth - 14\tabcolsep) * \real{0.1250}}
  >{\raggedright\arraybackslash}p{(\linewidth - 14\tabcolsep) * \real{0.1250}}
  >{\raggedright\arraybackslash}p{(\linewidth - 14\tabcolsep) * \real{0.1250}}@{}}
\toprule\noalign{}
\begin{minipage}[b]{\linewidth}\raggedright
RUNSPEC
\end{minipage} & \begin{minipage}[b]{\linewidth}\raggedright
GRID
\end{minipage} & \begin{minipage}[b]{\linewidth}\raggedright
EDIT
\end{minipage} & \begin{minipage}[b]{\linewidth}\raggedright
PROPS
\end{minipage} & \begin{minipage}[b]{\linewidth}\raggedright
REGIONS
\end{minipage} & \begin{minipage}[b]{\linewidth}\raggedright
SOLUTION
\end{minipage} & \begin{minipage}[b]{\linewidth}\raggedright
SUMMARY
\end{minipage} & \begin{minipage}[b]{\linewidth}\raggedright
SCHEDULE
\end{minipage} \\
\midrule\noalign{}
\endhead
\bottomrule\noalign{}
\endlastfoot
\end{longtable}

\subsubsection{Description}\label{description-203}

The DIFFMZ- keyword defines the negative z-direction diffusivity
multipliers for cells in the current input box for when the Diffusivity
option has been activated by the DIFFUSE keyword in the RUNSPEC section.

This keyword is not supported by OPM Flow but would change the results
if supported so the simulation will be stopped.

\subsection{DIFFMZ -- Define Grid Block Z-Direction Diffusivity
Multipliers}\label{diffmz-define-grid-block-z-direction-diffusivity-multipliers}

\begin{longtable}[]{@{}
  >{\raggedright\arraybackslash}p{(\linewidth - 14\tabcolsep) * \real{0.1250}}
  >{\raggedright\arraybackslash}p{(\linewidth - 14\tabcolsep) * \real{0.1250}}
  >{\raggedright\arraybackslash}p{(\linewidth - 14\tabcolsep) * \real{0.1250}}
  >{\raggedright\arraybackslash}p{(\linewidth - 14\tabcolsep) * \real{0.1250}}
  >{\raggedright\arraybackslash}p{(\linewidth - 14\tabcolsep) * \real{0.1250}}
  >{\raggedright\arraybackslash}p{(\linewidth - 14\tabcolsep) * \real{0.1250}}
  >{\raggedright\arraybackslash}p{(\linewidth - 14\tabcolsep) * \real{0.1250}}
  >{\raggedright\arraybackslash}p{(\linewidth - 14\tabcolsep) * \real{0.1250}}@{}}
\toprule\noalign{}
\begin{minipage}[b]{\linewidth}\raggedright
RUNSPEC
\end{minipage} & \begin{minipage}[b]{\linewidth}\raggedright
GRID
\end{minipage} & \begin{minipage}[b]{\linewidth}\raggedright
EDIT
\end{minipage} & \begin{minipage}[b]{\linewidth}\raggedright
PROPS
\end{minipage} & \begin{minipage}[b]{\linewidth}\raggedright
REGIONS
\end{minipage} & \begin{minipage}[b]{\linewidth}\raggedright
SOLUTION
\end{minipage} & \begin{minipage}[b]{\linewidth}\raggedright
SUMMARY
\end{minipage} & \begin{minipage}[b]{\linewidth}\raggedright
SCHEDULE
\end{minipage} \\
\midrule\noalign{}
\endhead
\bottomrule\noalign{}
\endlastfoot
\end{longtable}

\subsubsection{Description}\label{description-204}

The DIFFMZ keyword defines the z-direction diffusivity multipliers for
cells in the current input box for when the Diffusivity option has been
activated by the DIFFUSE keyword in the RUNSPEC section.

This keyword is not supported by OPM Flow but would change the results
if supported so the simulation will be stopped.

\subsection{DISPERC -- Define the Mechanical Dispersivity for All
Cells}\label{disperc-define-the-mechanical-dispersivity-for-all-cells}

\begin{longtable}[]{@{}
  >{\raggedright\arraybackslash}p{(\linewidth - 14\tabcolsep) * \real{0.1250}}
  >{\raggedright\arraybackslash}p{(\linewidth - 14\tabcolsep) * \real{0.1250}}
  >{\raggedright\arraybackslash}p{(\linewidth - 14\tabcolsep) * \real{0.1250}}
  >{\raggedright\arraybackslash}p{(\linewidth - 14\tabcolsep) * \real{0.1250}}
  >{\raggedright\arraybackslash}p{(\linewidth - 14\tabcolsep) * \real{0.1250}}
  >{\raggedright\arraybackslash}p{(\linewidth - 14\tabcolsep) * \real{0.1250}}
  >{\raggedright\arraybackslash}p{(\linewidth - 14\tabcolsep) * \real{0.1250}}
  >{\raggedright\arraybackslash}p{(\linewidth - 14\tabcolsep) * \real{0.1250}}@{}}
\toprule\noalign{}
\begin{minipage}[b]{\linewidth}\raggedright
RUNSPEC
\end{minipage} & \begin{minipage}[b]{\linewidth}\raggedright
GRID
\end{minipage} & \begin{minipage}[b]{\linewidth}\raggedright
EDIT
\end{minipage} & \begin{minipage}[b]{\linewidth}\raggedright
PROPS
\end{minipage} & \begin{minipage}[b]{\linewidth}\raggedright
REGIONS
\end{minipage} & \begin{minipage}[b]{\linewidth}\raggedright
SOLUTION
\end{minipage} & \begin{minipage}[b]{\linewidth}\raggedright
SUMMARY
\end{minipage} & \begin{minipage}[b]{\linewidth}\raggedright
SCHEDULE
\end{minipage} \\
\midrule\noalign{}
\endhead
\bottomrule\noalign{}
\endlastfoot
\end{longtable}

\subsubsection{Description}\label{description-205}

The \hyperref[REF_HEADING_KEYWORD_DISPERC_6_3]{DISPERC} keyword defines
the mechanical dispersivity for all cells in the model via an array.

\begin{longtable}[]{@{}
  >{\raggedright\arraybackslash}p{(\linewidth - 6\tabcolsep) * \real{0.2500}}
  >{\raggedright\arraybackslash}p{(\linewidth - 6\tabcolsep) * \real{0.2500}}
  >{\raggedright\arraybackslash}p{(\linewidth - 6\tabcolsep) * \real{0.2500}}
  >{\raggedright\arraybackslash}p{(\linewidth - 6\tabcolsep) * \real{0.2500}}@{}}
\toprule\noalign{}
\begin{minipage}[b]{\linewidth}\raggedright
No.
\end{minipage} & \begin{minipage}[b]{\linewidth}\raggedright
Name
\end{minipage} & \begin{minipage}[b]{\linewidth}\raggedright
Description
\end{minipage} & \begin{minipage}[b]{\linewidth}\raggedright
Default
\end{minipage} \\
\midrule\noalign{}
\endhead
\bottomrule\noalign{}
\endlastfoot
& Metric & Laboratory & \\
1 & \hyperref[REF_HEADING_KEYWORD_DISPERC_6_3]{DISPERC} &
\hyperref[REF_HEADING_KEYWORD_DISPERC_6_3]{DISPERC} is an array of real
positive values that defines the mechanical dispersivity for each cell
in the model. Repeat counts may be used, for example 20*1.0. & None \\
feet & m & cm & \\
Notes: & & & \\
\end{longtable}

\emph{Table 6.3.39.1:
\hyperref[REF_HEADING_KEYWORD_DISPERC_6_3]{DISPERC} Keyword Description}

\begin{tcolorbox}[enhanced jigsaw, title=\textcolor{quarto-callout-note-color}{\faInfo}\hspace{0.5em}{Note}, coltitle=black, breakable, toprule=.15mm, colbacktitle=quarto-callout-note-color!10!white, leftrule=.75mm, opacitybacktitle=0.6, colback=white, bottomtitle=1mm, toptitle=1mm, titlerule=0mm, left=2mm, colframe=quarto-callout-note-color-frame, arc=.35mm, rightrule=.15mm, bottomrule=.15mm, opacityback=0]

The option has been tested in combination with the CO2STORE,
\hyperref[REF_HEADING_KEYWORD_H2STORE]{H2STORE},
\hyperref[REF_HEADING_KEYWORD_BIOFILM]{BIOFILM}, or MICP keywords, but
not for the general case at this point.

\end{tcolorbox}

See also the CO2STORE and
\hyperref[REF_HEADING_KEYWORD_H2STORE]{H2STORE} keywords in the RUNSPEC
section that active OPM Flow's CO2 and H2 storage models respectively.

\subsubsection{Example}\label{example-144}

The example sets the mechanical dispersivity for all cells in the model
to 1.0.

\begin{verbatim}
--
--       SET MECHANICAL DISPERSIVITY FOR ALL CELLS (OPM FLOW KEYWORD)
--
DISPERC
         1000*1.0                                           /
\end{verbatim}

\subsection{DOMAINS -- Define the Parallel Domain
Properties}\label{domains-define-the-parallel-domain-properties}

\begin{longtable}[]{@{}
  >{\raggedright\arraybackslash}p{(\linewidth - 14\tabcolsep) * \real{0.1250}}
  >{\raggedright\arraybackslash}p{(\linewidth - 14\tabcolsep) * \real{0.1250}}
  >{\raggedright\arraybackslash}p{(\linewidth - 14\tabcolsep) * \real{0.1250}}
  >{\raggedright\arraybackslash}p{(\linewidth - 14\tabcolsep) * \real{0.1250}}
  >{\raggedright\arraybackslash}p{(\linewidth - 14\tabcolsep) * \real{0.1250}}
  >{\raggedright\arraybackslash}p{(\linewidth - 14\tabcolsep) * \real{0.1250}}
  >{\raggedright\arraybackslash}p{(\linewidth - 14\tabcolsep) * \real{0.1250}}
  >{\raggedright\arraybackslash}p{(\linewidth - 14\tabcolsep) * \real{0.1250}}@{}}
\toprule\noalign{}
\begin{minipage}[b]{\linewidth}\raggedright
RUNSPEC
\end{minipage} & \begin{minipage}[b]{\linewidth}\raggedright
GRID
\end{minipage} & \begin{minipage}[b]{\linewidth}\raggedright
EDIT
\end{minipage} & \begin{minipage}[b]{\linewidth}\raggedright
PROPS
\end{minipage} & \begin{minipage}[b]{\linewidth}\raggedright
REGIONS
\end{minipage} & \begin{minipage}[b]{\linewidth}\raggedright
SOLUTION
\end{minipage} & \begin{minipage}[b]{\linewidth}\raggedright
SUMMARY
\end{minipage} & \begin{minipage}[b]{\linewidth}\raggedright
SCHEDULE
\end{minipage} \\
\midrule\noalign{}
\endhead
\bottomrule\noalign{}
\endlastfoot
\end{longtable}

\subsubsection{Description}\label{description-206}

The DOMAINS keyword defines the parallel domain properties for when
parallel processing has been invoked by the PARALLEL keyword in the
RUNSPEC section.

This keyword is not supported by OPM Flow but would change the results
if supported so the simulation will be stopped.

See section 2.2 Running OPM Flow 2023-04 From The Command Line on how to
run OPM Flow in parallel mode.

\subsection{DPGRID -- Activate The Matrix Cell to Fracture Cell
Option}\label{dpgrid-activate-the-matrix-cell-to-fracture-cell-option}

\begin{longtable}[]{@{}
  >{\raggedright\arraybackslash}p{(\linewidth - 14\tabcolsep) * \real{0.1250}}
  >{\raggedright\arraybackslash}p{(\linewidth - 14\tabcolsep) * \real{0.1250}}
  >{\raggedright\arraybackslash}p{(\linewidth - 14\tabcolsep) * \real{0.1250}}
  >{\raggedright\arraybackslash}p{(\linewidth - 14\tabcolsep) * \real{0.1250}}
  >{\raggedright\arraybackslash}p{(\linewidth - 14\tabcolsep) * \real{0.1250}}
  >{\raggedright\arraybackslash}p{(\linewidth - 14\tabcolsep) * \real{0.1250}}
  >{\raggedright\arraybackslash}p{(\linewidth - 14\tabcolsep) * \real{0.1250}}
  >{\raggedright\arraybackslash}p{(\linewidth - 14\tabcolsep) * \real{0.1250}}@{}}
\toprule\noalign{}
\begin{minipage}[b]{\linewidth}\raggedright
RUNSPEC
\end{minipage} & \begin{minipage}[b]{\linewidth}\raggedright
GRID
\end{minipage} & \begin{minipage}[b]{\linewidth}\raggedright
EDIT
\end{minipage} & \begin{minipage}[b]{\linewidth}\raggedright
PROPS
\end{minipage} & \begin{minipage}[b]{\linewidth}\raggedright
REGIONS
\end{minipage} & \begin{minipage}[b]{\linewidth}\raggedright
SOLUTION
\end{minipage} & \begin{minipage}[b]{\linewidth}\raggedright
SUMMARY
\end{minipage} & \begin{minipage}[b]{\linewidth}\raggedright
SCHEDULE
\end{minipage} \\
\midrule\noalign{}
\endhead
\bottomrule\noalign{}
\endlastfoot
\end{longtable}

\subsubsection{Description}\label{description-207}

The DPGRID keyword activates the matrix cell to fracture cell option for
dual porosity runs for when a Dual Porosity model has been activated by
either the DUALPORO or DUALPERM keywords in the RUNSPEC section. The
keyword allows for only the matrix grid data to be entered and the
missing fracture cells are set to the inputted matrix cells.

This keyword is not supported by OPM Flow but would change the results
if supported so the simulation will be stopped.

\subsection{DPNUM -- Define Dual and Single Porosity Grid Block
Array}\label{dpnum-define-dual-and-single-porosity-grid-block-array}

\begin{longtable}[]{@{}
  >{\raggedright\arraybackslash}p{(\linewidth - 14\tabcolsep) * \real{0.1250}}
  >{\raggedright\arraybackslash}p{(\linewidth - 14\tabcolsep) * \real{0.1250}}
  >{\raggedright\arraybackslash}p{(\linewidth - 14\tabcolsep) * \real{0.1250}}
  >{\raggedright\arraybackslash}p{(\linewidth - 14\tabcolsep) * \real{0.1250}}
  >{\raggedright\arraybackslash}p{(\linewidth - 14\tabcolsep) * \real{0.1250}}
  >{\raggedright\arraybackslash}p{(\linewidth - 14\tabcolsep) * \real{0.1250}}
  >{\raggedright\arraybackslash}p{(\linewidth - 14\tabcolsep) * \real{0.1250}}
  >{\raggedright\arraybackslash}p{(\linewidth - 14\tabcolsep) * \real{0.1250}}@{}}
\toprule\noalign{}
\begin{minipage}[b]{\linewidth}\raggedright
RUNSPEC
\end{minipage} & \begin{minipage}[b]{\linewidth}\raggedright
GRID
\end{minipage} & \begin{minipage}[b]{\linewidth}\raggedright
EDIT
\end{minipage} & \begin{minipage}[b]{\linewidth}\raggedright
PROPS
\end{minipage} & \begin{minipage}[b]{\linewidth}\raggedright
REGIONS
\end{minipage} & \begin{minipage}[b]{\linewidth}\raggedright
SOLUTION
\end{minipage} & \begin{minipage}[b]{\linewidth}\raggedright
SUMMARY
\end{minipage} & \begin{minipage}[b]{\linewidth}\raggedright
SCHEDULE
\end{minipage} \\
\midrule\noalign{}
\endhead
\bottomrule\noalign{}
\endlastfoot
\end{longtable}

\subsubsection{Description}\label{description-208}

In dual porosity runs only, that is not dual permeability runs, the
DPNUM keyword defines which wells should be treated as single porosity
cells and which cells should be treated as dual porosity cells, for when
the Dual Porosity model has been activated by the DUALPORO keyword in
the RUNSPEC section.

This keyword is not supported by OPM Flow but would change the results
if supported so the simulation will be stopped.

\subsection{DR -- Define the Size of Grid Blocks in the R Direction for
All
Cells}\label{dr-define-the-size-of-grid-blocks-in-the-r-direction-for-all-cells}

\begin{longtable}[]{@{}
  >{\raggedright\arraybackslash}p{(\linewidth - 14\tabcolsep) * \real{0.1250}}
  >{\raggedright\arraybackslash}p{(\linewidth - 14\tabcolsep) * \real{0.1250}}
  >{\raggedright\arraybackslash}p{(\linewidth - 14\tabcolsep) * \real{0.1250}}
  >{\raggedright\arraybackslash}p{(\linewidth - 14\tabcolsep) * \real{0.1250}}
  >{\raggedright\arraybackslash}p{(\linewidth - 14\tabcolsep) * \real{0.1250}}
  >{\raggedright\arraybackslash}p{(\linewidth - 14\tabcolsep) * \real{0.1250}}
  >{\raggedright\arraybackslash}p{(\linewidth - 14\tabcolsep) * \real{0.1250}}
  >{\raggedright\arraybackslash}p{(\linewidth - 14\tabcolsep) * \real{0.1250}}@{}}
\toprule\noalign{}
\begin{minipage}[b]{\linewidth}\raggedright
RUNSPEC
\end{minipage} & \begin{minipage}[b]{\linewidth}\raggedright
GRID
\end{minipage} & \begin{minipage}[b]{\linewidth}\raggedright
EDIT
\end{minipage} & \begin{minipage}[b]{\linewidth}\raggedright
PROPS
\end{minipage} & \begin{minipage}[b]{\linewidth}\raggedright
REGIONS
\end{minipage} & \begin{minipage}[b]{\linewidth}\raggedright
SOLUTION
\end{minipage} & \begin{minipage}[b]{\linewidth}\raggedright
SUMMARY
\end{minipage} & \begin{minipage}[b]{\linewidth}\raggedright
SCHEDULE
\end{minipage} \\
\midrule\noalign{}
\endhead
\bottomrule\noalign{}
\endlastfoot
\end{longtable}

\subsubsection{Description}\label{description-209}

DR defines the size of all grid blocks in the R direction via an array
for each cell in the model. The RADIAL or SPIDER keyword in the RUNSPEC
section should be activated to indicate that radial or spider geometry
is being used.

This keyword is not supported by OPM Flow but would change the results
if supported so the simulation will be stopped.

\begin{longtable}[]{@{}
  >{\raggedright\arraybackslash}p{(\linewidth - 6\tabcolsep) * \real{0.2500}}
  >{\raggedright\arraybackslash}p{(\linewidth - 6\tabcolsep) * \real{0.2500}}
  >{\raggedright\arraybackslash}p{(\linewidth - 6\tabcolsep) * \real{0.2500}}
  >{\raggedright\arraybackslash}p{(\linewidth - 6\tabcolsep) * \real{0.2500}}@{}}
\toprule\noalign{}
\begin{minipage}[b]{\linewidth}\raggedright
No.
\end{minipage} & \begin{minipage}[b]{\linewidth}\raggedright
Name
\end{minipage} & \begin{minipage}[b]{\linewidth}\raggedright
Description
\end{minipage} & \begin{minipage}[b]{\linewidth}\raggedright
Default
\end{minipage} \\
\midrule\noalign{}
\endhead
\bottomrule\noalign{}
\endlastfoot
Field & Metric & Laboratory & \\
1 & DR & DR is an array of real numbers describing the cell size in the
R direction for each cell in the model in a radial grid. Repeat counts
may be used, for example 10*100.0. & None \\
feet & m & cm & \\
Notes: & & & \\
\end{longtable}

\emph{Table 6.22: DR Keyword Description}

Note that the SPIDER keyword activates OPM Flow's radial grid geometry
option for the model. This keyword will create a spiderweb-shaped grid
based on a corner-point grid using the standard radial grid keywords:
INRAD, DRV, DTHETAV, DZ/DZV etc. in the GRID the section. A spider grid
can be viewed in 3D in OPM ResInsight unlike radial grids that cannot be
viewed in the software. To overcome this, the simulator now converts
radial grids to Irregular Corner-Point Grids and adjusts the model's
pore volume to reflect radial coordinates; thus, overcoming the display
limitation.

See also the DRV, DTHETAV, DZ and TOPS keywords in the GRID section to
fully define a radial or spider grid model.

\subsubsection{Example}\label{example-145}

Given the dimensions of the grid in the RUNSPEC section to be 10,1, 8
for NX, NY and NZ respectively, we first define the inner radius of the
radial model,

\begin{verbatim}
--
--       INNER RADIUS OF FIRST GRID BLOCK IN THE RADIAL DIRECTION
--
INRAD
         0.25
/
\end{verbatim}

and then DR should be defined as:

\begin{verbatim}
--
--       DEFINE GRID BLOCK R DIRECTION CELL SIZE
--
DR
         1.75  2.32  5.01  10.84  23.39  50.55  109.21  235.92  509.68  1101.0
       1.75  2.32  5.01  10.84  23.39  50.55  109.21  235.92  509.68  1101.0
         1.75  2.32  5.01  10.84  23.39  50.55  109.21  235.92  509.68  1101.0
         1.75  2.32  5.01  10.84  23.39  50.55  109.21  235.92  509.68  1101.0
         1.75  2.32  5.01  10.84  23.39  50.55  109.21  235.92  509.68  1101.0
         1.75  2.32  5.01  10.84  23.39  50.55  109.21  235.92  509.68  1101.0
         1.75  2.32  5.01  10.84  23.39  50.55  109.21  235.92  509.68  1101.0
         1.75  2.32  5.01  10.84  23.39  50.55  109.21  235.92  509.68  1101.0
    /
\end{verbatim}

The above example defines the size of the cells in the R direction based
on 80 cells in the model as defined by the DIMENS keyword in the RUNSPEC
section.

Note that since the first layer (K=1) must be defined and subsequent
layers default to the layer above then:

\begin{verbatim}
--
--       INNER RADIUS OF FIRST GRID BLOCK IN THE RADIAL DIRECTION
--
INRAD
         0.25
/
--
--       DEFINE GRID BLOCK R DIRECTION CELL SIZE
--
DR
         1.75  2.32  5.01  10.84  23.39  50.55  109.21  235.92  509.68  1101.0
    /
\end{verbatim}

is equivalent to previous example.

\subsection{DRV -- Define the Size of Grid Blocks in the R Direction via
a
Vector}\label{drv-define-the-size-of-grid-blocks-in-the-r-direction-via-a-vector}

\begin{longtable}[]{@{}
  >{\raggedright\arraybackslash}p{(\linewidth - 14\tabcolsep) * \real{0.1250}}
  >{\raggedright\arraybackslash}p{(\linewidth - 14\tabcolsep) * \real{0.1250}}
  >{\raggedright\arraybackslash}p{(\linewidth - 14\tabcolsep) * \real{0.1250}}
  >{\raggedright\arraybackslash}p{(\linewidth - 14\tabcolsep) * \real{0.1250}}
  >{\raggedright\arraybackslash}p{(\linewidth - 14\tabcolsep) * \real{0.1250}}
  >{\raggedright\arraybackslash}p{(\linewidth - 14\tabcolsep) * \real{0.1250}}
  >{\raggedright\arraybackslash}p{(\linewidth - 14\tabcolsep) * \real{0.1250}}
  >{\raggedright\arraybackslash}p{(\linewidth - 14\tabcolsep) * \real{0.1250}}@{}}
\toprule\noalign{}
\begin{minipage}[b]{\linewidth}\raggedright
RUNSPEC
\end{minipage} & \begin{minipage}[b]{\linewidth}\raggedright
GRID
\end{minipage} & \begin{minipage}[b]{\linewidth}\raggedright
EDIT
\end{minipage} & \begin{minipage}[b]{\linewidth}\raggedright
PROPS
\end{minipage} & \begin{minipage}[b]{\linewidth}\raggedright
REGIONS
\end{minipage} & \begin{minipage}[b]{\linewidth}\raggedright
SOLUTION
\end{minipage} & \begin{minipage}[b]{\linewidth}\raggedright
SUMMARY
\end{minipage} & \begin{minipage}[b]{\linewidth}\raggedright
SCHEDULE
\end{minipage} \\
\midrule\noalign{}
\endhead
\bottomrule\noalign{}
\endlastfoot
\end{longtable}

\subsubsection{Description}\label{description-210}

DRV defines the size of grid blocks in the R direction via a vector as
opposed to defining the property for each cell for a Radial or Spider
Grid. The RADIAL or SPIDER keyword in the RUNSPEC section should be
activated to indicate that radial or spider geometry is being used.

This keyword is not supported by OPM Flow but has no effect on the
results so it will be ignored.

\begin{longtable}[]{@{}
  >{\raggedright\arraybackslash}p{(\linewidth - 6\tabcolsep) * \real{0.2500}}
  >{\raggedright\arraybackslash}p{(\linewidth - 6\tabcolsep) * \real{0.2500}}
  >{\raggedright\arraybackslash}p{(\linewidth - 6\tabcolsep) * \real{0.2500}}
  >{\raggedright\arraybackslash}p{(\linewidth - 6\tabcolsep) * \real{0.2500}}@{}}
\toprule\noalign{}
\begin{minipage}[b]{\linewidth}\raggedright
No.
\end{minipage} & \begin{minipage}[b]{\linewidth}\raggedright
Name
\end{minipage} & \begin{minipage}[b]{\linewidth}\raggedright
Description
\end{minipage} & \begin{minipage}[b]{\linewidth}\raggedright
Default
\end{minipage} \\
\midrule\noalign{}
\endhead
\bottomrule\noalign{}
\endlastfoot
Field & Metric & Laboratory & \\
1 & DRV & DRV is a vector of real numbers describing the cell size for
the grid blocks in the R direction in a radial grid. Repeat counts may
be used, for example 10*100.0. & None \\
feet & m & cm & \\
Notes: & & & \\
\end{longtable}

\emph{Table 6.23: DRV Keyword Description}

Note that the SPIDER keyword activates OPM Flow's radial grid geometry
option for the model. This keyword will create a spiderweb-shaped grid
based on a corner-point grid using the standard radial grid keywords:
INRAD, DRV, DTHETAV, DZ/DZV etc. in the GRID the section. A spider grid
can be viewed in 3D in OPM ResInsight unlike radial grids that cannot be
viewed in the software. To overcome this, the simulator now converts
radial grids to Irregular Corner-Point Grids and adjusts the model's
pore volume to reflect radial coordinates; thus, overcoming the display
limitation.

See also the DR, DTHETAV, DZ and TOPS keywords in the GRID section to
fully define a radial or spider grid model.

\subsubsection{Example}\label{example-146}

\begin{verbatim}
--
--       INNER RADIUS OF FIRST GRID BLOCK IN THE RADIAL DIRECTION
--
INRAD
         0.25
/
--
--       DEFINE GRID BLOCK SIZES IN THE R DIRECTION
--
DRV
         1.75  2.32  5.01  10.84  23.39  50.55  109.21  235.92  509.68  1101.08 /
\end{verbatim}

The above example defines the size of the cells in the R direction based
on NX equals 10 on the DIMENS keyword in the RUNSPEC section. Note the
INRAD keyword to define the inner radius of the radial grid.

\subsection{DTHETA -- Define the Size of Grid Blocks in the THETA
Direction for All
Cells}\label{dtheta-define-the-size-of-grid-blocks-in-the-theta-direction-for-all-cells}

\begin{longtable}[]{@{}
  >{\raggedright\arraybackslash}p{(\linewidth - 14\tabcolsep) * \real{0.1250}}
  >{\raggedright\arraybackslash}p{(\linewidth - 14\tabcolsep) * \real{0.1250}}
  >{\raggedright\arraybackslash}p{(\linewidth - 14\tabcolsep) * \real{0.1250}}
  >{\raggedright\arraybackslash}p{(\linewidth - 14\tabcolsep) * \real{0.1250}}
  >{\raggedright\arraybackslash}p{(\linewidth - 14\tabcolsep) * \real{0.1250}}
  >{\raggedright\arraybackslash}p{(\linewidth - 14\tabcolsep) * \real{0.1250}}
  >{\raggedright\arraybackslash}p{(\linewidth - 14\tabcolsep) * \real{0.1250}}
  >{\raggedright\arraybackslash}p{(\linewidth - 14\tabcolsep) * \real{0.1250}}@{}}
\toprule\noalign{}
\begin{minipage}[b]{\linewidth}\raggedright
RUNSPEC
\end{minipage} & \begin{minipage}[b]{\linewidth}\raggedright
GRID
\end{minipage} & \begin{minipage}[b]{\linewidth}\raggedright
EDIT
\end{minipage} & \begin{minipage}[b]{\linewidth}\raggedright
PROPS
\end{minipage} & \begin{minipage}[b]{\linewidth}\raggedright
REGIONS
\end{minipage} & \begin{minipage}[b]{\linewidth}\raggedright
SOLUTION
\end{minipage} & \begin{minipage}[b]{\linewidth}\raggedright
SUMMARY
\end{minipage} & \begin{minipage}[b]{\linewidth}\raggedright
SCHEDULE
\end{minipage} \\
\midrule\noalign{}
\endhead
\bottomrule\noalign{}
\endlastfoot
\end{longtable}

\subsubsection{Description}\label{description-211}

DTHETA defines the size of all grid blocks in the Theta direction via an
array for each cell in model. The RADIAL or SPIDER keyword in the
RUNSPEC section should be activated to indicate that radial or spider
geometry is being used.

This keyword is not supported by OPM Flow but would change the results
if supported so the simulation will be stopped.

\begin{longtable}[]{@{}
  >{\raggedright\arraybackslash}p{(\linewidth - 6\tabcolsep) * \real{0.2500}}
  >{\raggedright\arraybackslash}p{(\linewidth - 6\tabcolsep) * \real{0.2500}}
  >{\raggedright\arraybackslash}p{(\linewidth - 6\tabcolsep) * \real{0.2500}}
  >{\raggedright\arraybackslash}p{(\linewidth - 6\tabcolsep) * \real{0.2500}}@{}}
\toprule\noalign{}
\begin{minipage}[b]{\linewidth}\raggedright
No.
\end{minipage} & \begin{minipage}[b]{\linewidth}\raggedright
Name
\end{minipage} & \begin{minipage}[b]{\linewidth}\raggedright
Description
\end{minipage} & \begin{minipage}[b]{\linewidth}\raggedright
Default
\end{minipage} \\
\midrule\noalign{}
\endhead
\bottomrule\noalign{}
\endlastfoot
Field & Metric & Laboratory & \\
1 & DTHETA & DTHETA is an array of real numbers describing the cell size
in the THETA direction in radial grids for each cell in the model.
Repeat counts may be used, for example 10*25.0 & None \\
degrees & degrees & degrees & \\
Notes: & & & \\
\end{longtable}

\emph{Table 6.24: DTHETA Keyword Description}

Note that the SPIDER keyword activates OPM Flow's radial grid geometry
option for the model. This keyword will create a spiderweb-shaped grid
based on a corner-point grid using the standard radial grid keywords:
INRAD, DRV, DTHETAV, DZ/DZV etc. in the GRID the section. A spider grid
can be viewed in 3D in OPM ResInsight unlike radial grids that cannot be
viewed in the software. To overcome this, the simulator now converts
radial grids to Irregular Corner-Point Grids and adjusts the model's
pore volume to reflect radial coordinates; thus, overcoming the display
limitation. See also the DRV, DTHETAV, DZ and TOPS keywords in the GRID
section to fully define a radial or spider grid model.

\subsubsection{Example}\label{example-147}

Given the dimensions of the grid in the RUNSPEC section to be 10, 6, 1
for NX, NY and NZ respectively, then DTHETA should be defined as:

\begin{verbatim}
--
--       DEFINE GRID BLOCK SIZES IN THE THETA DIRECTION
--
DTHETA
         10*60.0
         10*60.0
         10*60.0
         10*60.0
         10*60.0
         10*60.0
/
\end{verbatim}

The above example defines the size of the cells in the R direction based
on 60 cells in the model as defined by the DIMENS keyword in the RUNSPEC
section.

\subsection{DTHETAV -- Sets the Size of Grid Blocks in THETA Direction
via a
Vector}\label{dthetav-sets-the-size-of-grid-blocks-in-theta-direction-via-a-vector}

\begin{longtable}[]{@{}
  >{\raggedright\arraybackslash}p{(\linewidth - 14\tabcolsep) * \real{0.1250}}
  >{\raggedright\arraybackslash}p{(\linewidth - 14\tabcolsep) * \real{0.1250}}
  >{\raggedright\arraybackslash}p{(\linewidth - 14\tabcolsep) * \real{0.1250}}
  >{\raggedright\arraybackslash}p{(\linewidth - 14\tabcolsep) * \real{0.1250}}
  >{\raggedright\arraybackslash}p{(\linewidth - 14\tabcolsep) * \real{0.1250}}
  >{\raggedright\arraybackslash}p{(\linewidth - 14\tabcolsep) * \real{0.1250}}
  >{\raggedright\arraybackslash}p{(\linewidth - 14\tabcolsep) * \real{0.1250}}
  >{\raggedright\arraybackslash}p{(\linewidth - 14\tabcolsep) * \real{0.1250}}@{}}
\toprule\noalign{}
\begin{minipage}[b]{\linewidth}\raggedright
RUNSPEC
\end{minipage} & \begin{minipage}[b]{\linewidth}\raggedright
GRID
\end{minipage} & \begin{minipage}[b]{\linewidth}\raggedright
EDIT
\end{minipage} & \begin{minipage}[b]{\linewidth}\raggedright
PROPS
\end{minipage} & \begin{minipage}[b]{\linewidth}\raggedright
REGIONS
\end{minipage} & \begin{minipage}[b]{\linewidth}\raggedright
SOLUTION
\end{minipage} & \begin{minipage}[b]{\linewidth}\raggedright
SUMMARY
\end{minipage} & \begin{minipage}[b]{\linewidth}\raggedright
SCHEDULE
\end{minipage} \\
\midrule\noalign{}
\endhead
\bottomrule\noalign{}
\endlastfoot
\end{longtable}

\subsubsection{Description}\label{description-212}

DTHETAV defines the size of grid blocks in the THETA direction via a
vector as opposed to defining the property for each cell for a Radial
Grid. The RADIAL or SPIDER keyword in the RUNSPEC should be activated to
indicate that radial geometry is being used.

\begin{longtable}[]{@{}
  >{\raggedright\arraybackslash}p{(\linewidth - 6\tabcolsep) * \real{0.2500}}
  >{\raggedright\arraybackslash}p{(\linewidth - 6\tabcolsep) * \real{0.2500}}
  >{\raggedright\arraybackslash}p{(\linewidth - 6\tabcolsep) * \real{0.2500}}
  >{\raggedright\arraybackslash}p{(\linewidth - 6\tabcolsep) * \real{0.2500}}@{}}
\toprule\noalign{}
\begin{minipage}[b]{\linewidth}\raggedright
No.
\end{minipage} & \begin{minipage}[b]{\linewidth}\raggedright
Name
\end{minipage} & \begin{minipage}[b]{\linewidth}\raggedright
Description
\end{minipage} & \begin{minipage}[b]{\linewidth}\raggedright
Default
\end{minipage} \\
\midrule\noalign{}
\endhead
\bottomrule\noalign{}
\endlastfoot
Field & Metric & Laboratory & \\
1 & DTHETAV & DTHETAV is a vector of real numbers describing the cell
size for the grid blocks in the THETA direction in a radial grid. Repeat
counts may be used, for example 10*100.0. & None \\
degrees & degrees & degrees & \\
Notes: & & & \\
\end{longtable}

\emph{Table 6.25: DTHETA Keyword Description}

Note that the SPIDER keyword activates OPM Flow's radial grid geometry
option for the model. This keyword will create a spiderweb-shaped grid
based on a corner-point grid using the standard radial grid keywords:
INRAD, DRV, DTHETAV, DZ/DZV etc. in the GRID the section. A spider grid
can be viewed in 3D in OPM ResInsight unlike radial grids that cannot be
viewed in the software. To overcome this, the simulator now converts
radial grids to Irregular Corner-Point Grids and adjusts the model's
pore volume to reflect radial coordinates; thus, overcoming the display
limitation.

See also the DRV, DZV and TOPS keywords to fully define a radial or
spider grid model.

\subsubsection{Example}\label{example-148}

\begin{verbatim}
--
--       DEFINE GRID BLOCK SIZES IN THE THETA DIRECTION (BASED ON NY = 6)
--
DTHETAV
         60.0  60.0  60.0  60.0  60.0  60.0                                    /
\end{verbatim}

The above example defines the size of the cells in the THETA direction
based on NY equals six in the DIMENS keyword in the RUNSPEC section.

\subsection{DUMPFLUX -- Activate Writing Out of a Flux
File}\label{dumpflux-activate-writing-out-of-a-flux-file}

\begin{longtable}[]{@{}
  >{\raggedright\arraybackslash}p{(\linewidth - 14\tabcolsep) * \real{0.1250}}
  >{\raggedright\arraybackslash}p{(\linewidth - 14\tabcolsep) * \real{0.1250}}
  >{\raggedright\arraybackslash}p{(\linewidth - 14\tabcolsep) * \real{0.1250}}
  >{\raggedright\arraybackslash}p{(\linewidth - 14\tabcolsep) * \real{0.1250}}
  >{\raggedright\arraybackslash}p{(\linewidth - 14\tabcolsep) * \real{0.1250}}
  >{\raggedright\arraybackslash}p{(\linewidth - 14\tabcolsep) * \real{0.1250}}
  >{\raggedright\arraybackslash}p{(\linewidth - 14\tabcolsep) * \real{0.1250}}
  >{\raggedright\arraybackslash}p{(\linewidth - 14\tabcolsep) * \real{0.1250}}@{}}
\toprule\noalign{}
\begin{minipage}[b]{\linewidth}\raggedright
RUNSPEC
\end{minipage} & \begin{minipage}[b]{\linewidth}\raggedright
GRID
\end{minipage} & \begin{minipage}[b]{\linewidth}\raggedright
EDIT
\end{minipage} & \begin{minipage}[b]{\linewidth}\raggedright
PROPS
\end{minipage} & \begin{minipage}[b]{\linewidth}\raggedright
REGIONS
\end{minipage} & \begin{minipage}[b]{\linewidth}\raggedright
SOLUTION
\end{minipage} & \begin{minipage}[b]{\linewidth}\raggedright
SUMMARY
\end{minipage} & \begin{minipage}[b]{\linewidth}\raggedright
SCHEDULE
\end{minipage} \\
\midrule\noalign{}
\endhead
\bottomrule\noalign{}
\endlastfoot
\end{longtable}

\subsubsection{Description}\label{description-213}

This keyword activates the writing out of a full field (the full grid)
FLUX file for later processing in a Flux Boundary run. The Flux Boundary
feature allows for the segmentation of the full grid into flux boundary
areas which allow for a sub-area of the grid to be run and at the same
time model the flux across the boundary derived from the main grid. The
object of this feature is to be able to investigate the performance of
various areas of the model without having to run the full field, thus
improving computational efficiency and run times, but still obtain
``reasonable'' results due to the incorporation of the fluxes across the
boundary.

This feature is not available in OPM Flow; however it is documented here
for completeness.

There is no data required for this keyword and there is no terminating
``/'' for this keyword.

\subsubsection{Example}\label{example-149}

\begin{verbatim}
--
--       ACTIVATE WRITING OUT OF A FLUX FILE
--
DUMPFLUX                                                                        /
\end{verbatim}

The above example switches on the writing of the FLUX output file; the
keyword has no effect and is ignored by the simulator.

\subsection{DX -- Define the Size of Grid Blocks in the X Direction for
All
Cells}\label{dx-define-the-size-of-grid-blocks-in-the-x-direction-for-all-cells}

\begin{longtable}[]{@{}
  >{\raggedright\arraybackslash}p{(\linewidth - 14\tabcolsep) * \real{0.1250}}
  >{\raggedright\arraybackslash}p{(\linewidth - 14\tabcolsep) * \real{0.1250}}
  >{\raggedright\arraybackslash}p{(\linewidth - 14\tabcolsep) * \real{0.1250}}
  >{\raggedright\arraybackslash}p{(\linewidth - 14\tabcolsep) * \real{0.1250}}
  >{\raggedright\arraybackslash}p{(\linewidth - 14\tabcolsep) * \real{0.1250}}
  >{\raggedright\arraybackslash}p{(\linewidth - 14\tabcolsep) * \real{0.1250}}
  >{\raggedright\arraybackslash}p{(\linewidth - 14\tabcolsep) * \real{0.1250}}
  >{\raggedright\arraybackslash}p{(\linewidth - 14\tabcolsep) * \real{0.1250}}@{}}
\toprule\noalign{}
\begin{minipage}[b]{\linewidth}\raggedright
RUNSPEC
\end{minipage} & \begin{minipage}[b]{\linewidth}\raggedright
GRID
\end{minipage} & \begin{minipage}[b]{\linewidth}\raggedright
EDIT
\end{minipage} & \begin{minipage}[b]{\linewidth}\raggedright
PROPS
\end{minipage} & \begin{minipage}[b]{\linewidth}\raggedright
REGIONS
\end{minipage} & \begin{minipage}[b]{\linewidth}\raggedright
SOLUTION
\end{minipage} & \begin{minipage}[b]{\linewidth}\raggedright
SUMMARY
\end{minipage} & \begin{minipage}[b]{\linewidth}\raggedright
SCHEDULE
\end{minipage} \\
\midrule\noalign{}
\endhead
\bottomrule\noalign{}
\endlastfoot
\end{longtable}

\subsubsection{Description}\label{description-214}

DX defines the size of all grid blocks in the X direction via an array
for each cell in a Cartesian Regular Grid model.

\begin{longtable}[]{@{}
  >{\raggedright\arraybackslash}p{(\linewidth - 6\tabcolsep) * \real{0.2500}}
  >{\raggedright\arraybackslash}p{(\linewidth - 6\tabcolsep) * \real{0.2500}}
  >{\raggedright\arraybackslash}p{(\linewidth - 6\tabcolsep) * \real{0.2500}}
  >{\raggedright\arraybackslash}p{(\linewidth - 6\tabcolsep) * \real{0.2500}}@{}}
\toprule\noalign{}
\begin{minipage}[b]{\linewidth}\raggedright
No.
\end{minipage} & \begin{minipage}[b]{\linewidth}\raggedright
Name
\end{minipage} & \begin{minipage}[b]{\linewidth}\raggedright
Description
\end{minipage} & \begin{minipage}[b]{\linewidth}\raggedright
Default
\end{minipage} \\
\midrule\noalign{}
\endhead
\bottomrule\noalign{}
\endlastfoot
Field & Metric & Laboratory & \\
1 & DX & DX is an array of real numbers describing the cell size in the
X direction for each cell in the model. Repeat counts may be used, for
example 10*100.0. & None \\
feet & m & cm & \\
Notes: & & & \\
\end{longtable}

\emph{Table 6.26: DX Keyword Description}

See also the DY, DZ and TOPS keywords to fully define a Cartesian
Regular Grid.

\subsubsection{Example}\label{example-150}

\begin{verbatim}
--
--       DEFINE GRID BLOCK X DIRECTION CELL SIZE (BASED ON NX x NY x NZ = 300)
--
DX
         300*1000                                                              /
\end{verbatim}

The above example defines the size of the cells in the X direction based
on 300 cells in the model as defined by the DIMENS keyword in the
RUNSPEC section.

\subsection{DXV -- Define the Size of Grid Blocks in the X Direction via
a
Vector}\label{dxv-define-the-size-of-grid-blocks-in-the-x-direction-via-a-vector}

\begin{longtable}[]{@{}
  >{\raggedright\arraybackslash}p{(\linewidth - 14\tabcolsep) * \real{0.1250}}
  >{\raggedright\arraybackslash}p{(\linewidth - 14\tabcolsep) * \real{0.1250}}
  >{\raggedright\arraybackslash}p{(\linewidth - 14\tabcolsep) * \real{0.1250}}
  >{\raggedright\arraybackslash}p{(\linewidth - 14\tabcolsep) * \real{0.1250}}
  >{\raggedright\arraybackslash}p{(\linewidth - 14\tabcolsep) * \real{0.1250}}
  >{\raggedright\arraybackslash}p{(\linewidth - 14\tabcolsep) * \real{0.1250}}
  >{\raggedright\arraybackslash}p{(\linewidth - 14\tabcolsep) * \real{0.1250}}
  >{\raggedright\arraybackslash}p{(\linewidth - 14\tabcolsep) * \real{0.1250}}@{}}
\toprule\noalign{}
\begin{minipage}[b]{\linewidth}\raggedright
RUNSPEC
\end{minipage} & \begin{minipage}[b]{\linewidth}\raggedright
GRID
\end{minipage} & \begin{minipage}[b]{\linewidth}\raggedright
EDIT
\end{minipage} & \begin{minipage}[b]{\linewidth}\raggedright
PROPS
\end{minipage} & \begin{minipage}[b]{\linewidth}\raggedright
REGIONS
\end{minipage} & \begin{minipage}[b]{\linewidth}\raggedright
SOLUTION
\end{minipage} & \begin{minipage}[b]{\linewidth}\raggedright
SUMMARY
\end{minipage} & \begin{minipage}[b]{\linewidth}\raggedright
SCHEDULE
\end{minipage} \\
\midrule\noalign{}
\endhead
\bottomrule\noalign{}
\endlastfoot
\end{longtable}

\subsubsection{Description}\label{description-215}

DXV defines the size of grid blocks in the X direction via a vector as
opposed to defining the X direction cell size for each cell for a
Cartesian Regular Grid.

\begin{longtable}[]{@{}
  >{\raggedright\arraybackslash}p{(\linewidth - 6\tabcolsep) * \real{0.2500}}
  >{\raggedright\arraybackslash}p{(\linewidth - 6\tabcolsep) * \real{0.2500}}
  >{\raggedright\arraybackslash}p{(\linewidth - 6\tabcolsep) * \real{0.2500}}
  >{\raggedright\arraybackslash}p{(\linewidth - 6\tabcolsep) * \real{0.2500}}@{}}
\toprule\noalign{}
\begin{minipage}[b]{\linewidth}\raggedright
No.
\end{minipage} & \begin{minipage}[b]{\linewidth}\raggedright
Name
\end{minipage} & \begin{minipage}[b]{\linewidth}\raggedright
Description
\end{minipage} & \begin{minipage}[b]{\linewidth}\raggedright
Default
\end{minipage} \\
\midrule\noalign{}
\endhead
\bottomrule\noalign{}
\endlastfoot
Field & Metric & Laboratory & \\
1 & DXV & DXV is a vector of real numbers describing the cell size for
the grid blocks in the X direction. Repeat counts may be used, for
example 10*100.0. & None \\
feet & m & cm & \\
Notes: & & & \\
\end{longtable}

\emph{Table 6.27: DXV Keyword Description}

See also the DYV, DZV and TOPS keywords to fully define a Cartesian
Regular Grid.

\subsubsection{Example}\label{example-151}

\begin{verbatim}
--
--       DEFINE GRID BLOCK X DIRECTION CELL SIZE (BASED ON NX = 5)
--
DXV
         5*100                                                                 /
\end{verbatim}

The above example defines the size of the cells in the X direction based
on NX equals 5 on the DIMENS keyword in the RUNSPEC section.

\subsection{DY -- Define the Size of Grid Blocks in the Y Direction for
All
Cells}\label{dy-define-the-size-of-grid-blocks-in-the-y-direction-for-all-cells}

\begin{longtable}[]{@{}
  >{\raggedright\arraybackslash}p{(\linewidth - 14\tabcolsep) * \real{0.1250}}
  >{\raggedright\arraybackslash}p{(\linewidth - 14\tabcolsep) * \real{0.1250}}
  >{\raggedright\arraybackslash}p{(\linewidth - 14\tabcolsep) * \real{0.1250}}
  >{\raggedright\arraybackslash}p{(\linewidth - 14\tabcolsep) * \real{0.1250}}
  >{\raggedright\arraybackslash}p{(\linewidth - 14\tabcolsep) * \real{0.1250}}
  >{\raggedright\arraybackslash}p{(\linewidth - 14\tabcolsep) * \real{0.1250}}
  >{\raggedright\arraybackslash}p{(\linewidth - 14\tabcolsep) * \real{0.1250}}
  >{\raggedright\arraybackslash}p{(\linewidth - 14\tabcolsep) * \real{0.1250}}@{}}
\toprule\noalign{}
\begin{minipage}[b]{\linewidth}\raggedright
RUNSPEC
\end{minipage} & \begin{minipage}[b]{\linewidth}\raggedright
GRID
\end{minipage} & \begin{minipage}[b]{\linewidth}\raggedright
EDIT
\end{minipage} & \begin{minipage}[b]{\linewidth}\raggedright
PROPS
\end{minipage} & \begin{minipage}[b]{\linewidth}\raggedright
REGIONS
\end{minipage} & \begin{minipage}[b]{\linewidth}\raggedright
SOLUTION
\end{minipage} & \begin{minipage}[b]{\linewidth}\raggedright
SUMMARY
\end{minipage} & \begin{minipage}[b]{\linewidth}\raggedright
SCHEDULE
\end{minipage} \\
\midrule\noalign{}
\endhead
\bottomrule\noalign{}
\endlastfoot
\end{longtable}

\subsubsection{Description}\label{description-216}

DY defines the size of all grid blocks in the Y direction via an array
for each cell in a Cartesian Regular Grid model.

\begin{longtable}[]{@{}
  >{\raggedright\arraybackslash}p{(\linewidth - 6\tabcolsep) * \real{0.2500}}
  >{\raggedright\arraybackslash}p{(\linewidth - 6\tabcolsep) * \real{0.2500}}
  >{\raggedright\arraybackslash}p{(\linewidth - 6\tabcolsep) * \real{0.2500}}
  >{\raggedright\arraybackslash}p{(\linewidth - 6\tabcolsep) * \real{0.2500}}@{}}
\toprule\noalign{}
\begin{minipage}[b]{\linewidth}\raggedright
No.
\end{minipage} & \begin{minipage}[b]{\linewidth}\raggedright
Name
\end{minipage} & \begin{minipage}[b]{\linewidth}\raggedright
Description
\end{minipage} & \begin{minipage}[b]{\linewidth}\raggedright
Default
\end{minipage} \\
\midrule\noalign{}
\endhead
\bottomrule\noalign{}
\endlastfoot
Field & Metric & Laboratory & \\
1 & DY & DY is an array of real numbers describing the cell size in the
Y direction for each cell in the model. Repeat counts may be used, for
example 10*100.0. & None \\
feet & m & cm & \\
Notes: & & & \\
\end{longtable}

\emph{Table 6.28: DY Keyword Description}

See also the DX, DZ and TOPS keywords to fully define a Cartesian
Regular Grid.

\subsubsection{Example}\label{example-152}

\begin{verbatim}
--
--        DEFINE GRID BLOCK Y DIRECTION CELL SIZE (BASED ON NX x NY x NZ = 300)
--
DY
          300*1000                                                              /
\end{verbatim}

The above example defines the size of the cells in the Y direction based
on 300 cells in the model as defined by the DIMENS keyword in the
RUNSPEC section.

\subsection{DYV -- Define the Size of Grid Blocks in the Y Direction via
a
Vector}\label{dyv-define-the-size-of-grid-blocks-in-the-y-direction-via-a-vector}

\begin{longtable}[]{@{}
  >{\raggedright\arraybackslash}p{(\linewidth - 14\tabcolsep) * \real{0.1250}}
  >{\raggedright\arraybackslash}p{(\linewidth - 14\tabcolsep) * \real{0.1250}}
  >{\raggedright\arraybackslash}p{(\linewidth - 14\tabcolsep) * \real{0.1250}}
  >{\raggedright\arraybackslash}p{(\linewidth - 14\tabcolsep) * \real{0.1250}}
  >{\raggedright\arraybackslash}p{(\linewidth - 14\tabcolsep) * \real{0.1250}}
  >{\raggedright\arraybackslash}p{(\linewidth - 14\tabcolsep) * \real{0.1250}}
  >{\raggedright\arraybackslash}p{(\linewidth - 14\tabcolsep) * \real{0.1250}}
  >{\raggedright\arraybackslash}p{(\linewidth - 14\tabcolsep) * \real{0.1250}}@{}}
\toprule\noalign{}
\begin{minipage}[b]{\linewidth}\raggedright
RUNSPEC
\end{minipage} & \begin{minipage}[b]{\linewidth}\raggedright
GRID
\end{minipage} & \begin{minipage}[b]{\linewidth}\raggedright
EDIT
\end{minipage} & \begin{minipage}[b]{\linewidth}\raggedright
PROPS
\end{minipage} & \begin{minipage}[b]{\linewidth}\raggedright
REGIONS
\end{minipage} & \begin{minipage}[b]{\linewidth}\raggedright
SOLUTION
\end{minipage} & \begin{minipage}[b]{\linewidth}\raggedright
SUMMARY
\end{minipage} & \begin{minipage}[b]{\linewidth}\raggedright
SCHEDULE
\end{minipage} \\
\midrule\noalign{}
\endhead
\bottomrule\noalign{}
\endlastfoot
\end{longtable}

\subsubsection{Description}\label{description-217}

DYV defines the size of grid blocks in the Y direction via a vector as
opposed to defining the Y direction cell size for each cell for a
Cartesian Regular Grid.

\begin{longtable}[]{@{}
  >{\raggedright\arraybackslash}p{(\linewidth - 6\tabcolsep) * \real{0.2500}}
  >{\raggedright\arraybackslash}p{(\linewidth - 6\tabcolsep) * \real{0.2500}}
  >{\raggedright\arraybackslash}p{(\linewidth - 6\tabcolsep) * \real{0.2500}}
  >{\raggedright\arraybackslash}p{(\linewidth - 6\tabcolsep) * \real{0.2500}}@{}}
\toprule\noalign{}
\begin{minipage}[b]{\linewidth}\raggedright
No.
\end{minipage} & \begin{minipage}[b]{\linewidth}\raggedright
Name
\end{minipage} & \begin{minipage}[b]{\linewidth}\raggedright
Description
\end{minipage} & \begin{minipage}[b]{\linewidth}\raggedright
Default
\end{minipage} \\
\midrule\noalign{}
\endhead
\bottomrule\noalign{}
\endlastfoot
Field & Metric & Laboratory & \\
1 & DYV & DYV is a vector of real numbers describing the cell size for
the grid blocks in the Y direction. Repeat counts may be used, for
example 10*100.0. & None \\
feet & m & cm & \\
Notes: & & & \\
\end{longtable}

\emph{Table 6.29: DYV Keyword Description}

See also the DXV, DZV and TOPS keywords to fully define a Cartesian
Regular Grid.

\subsubsection{Example}\label{example-153}

\begin{verbatim}
--
--        DEFINE GRID BLOCK Y DIRECTION CELL SIZE (BASED ON NY = 5)
--
DYV
          5*100                                                                 /
\end{verbatim}

The above example defines the size of the cells in the Y direction based
on NY equals 5 on the DIMENS keyword in the RUNSPEC section.

\subsection{DZ -- Define the Size of Grid Blocks in the Z Direction for
All
Cells}\label{dz-define-the-size-of-grid-blocks-in-the-z-direction-for-all-cells}

\begin{longtable}[]{@{}
  >{\raggedright\arraybackslash}p{(\linewidth - 14\tabcolsep) * \real{0.1250}}
  >{\raggedright\arraybackslash}p{(\linewidth - 14\tabcolsep) * \real{0.1250}}
  >{\raggedright\arraybackslash}p{(\linewidth - 14\tabcolsep) * \real{0.1250}}
  >{\raggedright\arraybackslash}p{(\linewidth - 14\tabcolsep) * \real{0.1250}}
  >{\raggedright\arraybackslash}p{(\linewidth - 14\tabcolsep) * \real{0.1250}}
  >{\raggedright\arraybackslash}p{(\linewidth - 14\tabcolsep) * \real{0.1250}}
  >{\raggedright\arraybackslash}p{(\linewidth - 14\tabcolsep) * \real{0.1250}}
  >{\raggedright\arraybackslash}p{(\linewidth - 14\tabcolsep) * \real{0.1250}}@{}}
\toprule\noalign{}
\begin{minipage}[b]{\linewidth}\raggedright
RUNSPEC
\end{minipage} & \begin{minipage}[b]{\linewidth}\raggedright
GRID
\end{minipage} & \begin{minipage}[b]{\linewidth}\raggedright
EDIT
\end{minipage} & \begin{minipage}[b]{\linewidth}\raggedright
PROPS
\end{minipage} & \begin{minipage}[b]{\linewidth}\raggedright
REGIONS
\end{minipage} & \begin{minipage}[b]{\linewidth}\raggedright
SOLUTION
\end{minipage} & \begin{minipage}[b]{\linewidth}\raggedright
SUMMARY
\end{minipage} & \begin{minipage}[b]{\linewidth}\raggedright
SCHEDULE
\end{minipage} \\
\midrule\noalign{}
\endhead
\bottomrule\noalign{}
\endlastfoot
\end{longtable}

\subsubsection{Description}\label{description-218}

DZ defines the size of all grid blocks in the Z direction via an array
for each cell in a Cartesian Regular Grid model.

\begin{longtable}[]{@{}
  >{\raggedright\arraybackslash}p{(\linewidth - 6\tabcolsep) * \real{0.2500}}
  >{\raggedright\arraybackslash}p{(\linewidth - 6\tabcolsep) * \real{0.2500}}
  >{\raggedright\arraybackslash}p{(\linewidth - 6\tabcolsep) * \real{0.2500}}
  >{\raggedright\arraybackslash}p{(\linewidth - 6\tabcolsep) * \real{0.2500}}@{}}
\toprule\noalign{}
\begin{minipage}[b]{\linewidth}\raggedright
No.
\end{minipage} & \begin{minipage}[b]{\linewidth}\raggedright
Name
\end{minipage} & \begin{minipage}[b]{\linewidth}\raggedright
Description
\end{minipage} & \begin{minipage}[b]{\linewidth}\raggedright
Default
\end{minipage} \\
\midrule\noalign{}
\endhead
\bottomrule\noalign{}
\endlastfoot
Field & Metric & Laboratory & \\
1 & DZ & DZ is an array of real numbers describing the cell size in the
Z direction for each cell in the model. Repeat counts may be used, for
example 10*100.0. & None \\
feet & m & cm & \\
Notes: & & & \\
\end{longtable}

\emph{Table 6.30: DZ Keyword Description}

See also the DX, DY and TOPS keywords to fully define a Cartesian
Regular Grid.

\subsubsection{Example}\label{example-154}

\begin{verbatim}
--
--        DEFINE GRID BLOCK Z DIRECTION CELL SIZE (BASED ON NX x NY x NZ = 300)
--
DZ
          100*20.0   100*30.0   100*50.0                                        /
\end{verbatim}

The above example defines the size of the cells in the Z direction based
on 300 cells in the model as defined by the DIMENS keyword in the
RUNSPEC section.

\subsection{DZMATRIX -- Matrix Block Height for Gravity Drainage Model
For All
Cells}\label{dzmatrix-matrix-block-height-for-gravity-drainage-model-for-all-cells}

\begin{longtable}[]{@{}
  >{\raggedright\arraybackslash}p{(\linewidth - 14\tabcolsep) * \real{0.1250}}
  >{\raggedright\arraybackslash}p{(\linewidth - 14\tabcolsep) * \real{0.1250}}
  >{\raggedright\arraybackslash}p{(\linewidth - 14\tabcolsep) * \real{0.1250}}
  >{\raggedright\arraybackslash}p{(\linewidth - 14\tabcolsep) * \real{0.1250}}
  >{\raggedright\arraybackslash}p{(\linewidth - 14\tabcolsep) * \real{0.1250}}
  >{\raggedright\arraybackslash}p{(\linewidth - 14\tabcolsep) * \real{0.1250}}
  >{\raggedright\arraybackslash}p{(\linewidth - 14\tabcolsep) * \real{0.1250}}
  >{\raggedright\arraybackslash}p{(\linewidth - 14\tabcolsep) * \real{0.1250}}@{}}
\toprule\noalign{}
\begin{minipage}[b]{\linewidth}\raggedright
RUNSPEC
\end{minipage} & \begin{minipage}[b]{\linewidth}\raggedright
GRID
\end{minipage} & \begin{minipage}[b]{\linewidth}\raggedright
EDIT
\end{minipage} & \begin{minipage}[b]{\linewidth}\raggedright
PROPS
\end{minipage} & \begin{minipage}[b]{\linewidth}\raggedright
REGIONS
\end{minipage} & \begin{minipage}[b]{\linewidth}\raggedright
SOLUTION
\end{minipage} & \begin{minipage}[b]{\linewidth}\raggedright
SUMMARY
\end{minipage} & \begin{minipage}[b]{\linewidth}\raggedright
SCHEDULE
\end{minipage} \\
\midrule\noalign{}
\endhead
\bottomrule\noalign{}
\endlastfoot
\end{longtable}

\subsubsection{Description}\label{description-219}

The DZMATRIX keyword defines the matrix block height for the gravity
drainage model by grid block for when the Dual Permeability or Dual
Porosity models are activated by the DUALPERM and DUALPORO keywords and
the Gravity Drainage option is invoked via the GRAVDR and GRAVDRM
keywords. All keywords are in the RUNSPEC section.

The keyword is identical to the DZMTRXV keyword in the GRID section.

This keyword is not supported by OPM Flow but would change the results
if supported so the simulation will be stopped.

\subsection{DZMTRX -- Matrix Block Height for Gravity Drainage Model for
the
Grid}\label{dzmtrx-matrix-block-height-for-gravity-drainage-model-for-the-grid}

\begin{longtable}[]{@{}
  >{\raggedright\arraybackslash}p{(\linewidth - 14\tabcolsep) * \real{0.1250}}
  >{\raggedright\arraybackslash}p{(\linewidth - 14\tabcolsep) * \real{0.1250}}
  >{\raggedright\arraybackslash}p{(\linewidth - 14\tabcolsep) * \real{0.1250}}
  >{\raggedright\arraybackslash}p{(\linewidth - 14\tabcolsep) * \real{0.1250}}
  >{\raggedright\arraybackslash}p{(\linewidth - 14\tabcolsep) * \real{0.1250}}
  >{\raggedright\arraybackslash}p{(\linewidth - 14\tabcolsep) * \real{0.1250}}
  >{\raggedright\arraybackslash}p{(\linewidth - 14\tabcolsep) * \real{0.1250}}
  >{\raggedright\arraybackslash}p{(\linewidth - 14\tabcolsep) * \real{0.1250}}@{}}
\toprule\noalign{}
\begin{minipage}[b]{\linewidth}\raggedright
RUNSPEC
\end{minipage} & \begin{minipage}[b]{\linewidth}\raggedright
GRID
\end{minipage} & \begin{minipage}[b]{\linewidth}\raggedright
EDIT
\end{minipage} & \begin{minipage}[b]{\linewidth}\raggedright
PROPS
\end{minipage} & \begin{minipage}[b]{\linewidth}\raggedright
REGIONS
\end{minipage} & \begin{minipage}[b]{\linewidth}\raggedright
SOLUTION
\end{minipage} & \begin{minipage}[b]{\linewidth}\raggedright
SUMMARY
\end{minipage} & \begin{minipage}[b]{\linewidth}\raggedright
SCHEDULE
\end{minipage} \\
\midrule\noalign{}
\endhead
\bottomrule\noalign{}
\endlastfoot
\end{longtable}

\subsubsection{Description}\label{description-220}

The DZMTRX keyword defines a constant matrix block height for the
gravity drainage model for the entire grid for when the Dual
Permeability or Dual Porosity models are activated by the DUALPERM and
DUALPORO keywords and the Gravity Drainage option is invoked via the
GRAVDR and GRAVDRM keywords. All keywords are in the RUNSPEC section.

This keyword is not supported by OPM Flow but would change the results
if supported so the simulation will be stopped.

\subsection{DZMTRXV -- Matrix Block Height for Gravity Drainage Model
For All
Cells}\label{dzmtrxv-matrix-block-height-for-gravity-drainage-model-for-all-cells}

\begin{longtable}[]{@{}
  >{\raggedright\arraybackslash}p{(\linewidth - 14\tabcolsep) * \real{0.1250}}
  >{\raggedright\arraybackslash}p{(\linewidth - 14\tabcolsep) * \real{0.1250}}
  >{\raggedright\arraybackslash}p{(\linewidth - 14\tabcolsep) * \real{0.1250}}
  >{\raggedright\arraybackslash}p{(\linewidth - 14\tabcolsep) * \real{0.1250}}
  >{\raggedright\arraybackslash}p{(\linewidth - 14\tabcolsep) * \real{0.1250}}
  >{\raggedright\arraybackslash}p{(\linewidth - 14\tabcolsep) * \real{0.1250}}
  >{\raggedright\arraybackslash}p{(\linewidth - 14\tabcolsep) * \real{0.1250}}
  >{\raggedright\arraybackslash}p{(\linewidth - 14\tabcolsep) * \real{0.1250}}@{}}
\toprule\noalign{}
\begin{minipage}[b]{\linewidth}\raggedright
RUNSPEC
\end{minipage} & \begin{minipage}[b]{\linewidth}\raggedright
GRID
\end{minipage} & \begin{minipage}[b]{\linewidth}\raggedright
EDIT
\end{minipage} & \begin{minipage}[b]{\linewidth}\raggedright
PROPS
\end{minipage} & \begin{minipage}[b]{\linewidth}\raggedright
REGIONS
\end{minipage} & \begin{minipage}[b]{\linewidth}\raggedright
SOLUTION
\end{minipage} & \begin{minipage}[b]{\linewidth}\raggedright
SUMMARY
\end{minipage} & \begin{minipage}[b]{\linewidth}\raggedright
SCHEDULE
\end{minipage} \\
\midrule\noalign{}
\endhead
\bottomrule\noalign{}
\endlastfoot
\end{longtable}

\subsubsection{Description}\label{description-221}

The DZMATRIX keyword defines the matrix block height for the gravity
drainage model by grid block for when the Dual Permeability or Dual
Porosity models are activated by the DUALPERM and DUALPORO keywords and
the Gravity Drainage option is invoked via the GRAVDR and GRAVDRM
keywords. All keywords are in the RUNSPEC section.

The keyword is identical to the DZMATRIX keyword in the GRID section.

This keyword is not supported by OPM Flow but would change the results
if supported so the simulation will be stopped.

\subsection{DZNET -- Define Grid Block Net Thickness for All
Cells}\label{dznet-define-grid-block-net-thickness-for-all-cells}

\begin{longtable}[]{@{}
  >{\raggedright\arraybackslash}p{(\linewidth - 14\tabcolsep) * \real{0.1250}}
  >{\raggedright\arraybackslash}p{(\linewidth - 14\tabcolsep) * \real{0.1250}}
  >{\raggedright\arraybackslash}p{(\linewidth - 14\tabcolsep) * \real{0.1250}}
  >{\raggedright\arraybackslash}p{(\linewidth - 14\tabcolsep) * \real{0.1250}}
  >{\raggedright\arraybackslash}p{(\linewidth - 14\tabcolsep) * \real{0.1250}}
  >{\raggedright\arraybackslash}p{(\linewidth - 14\tabcolsep) * \real{0.1250}}
  >{\raggedright\arraybackslash}p{(\linewidth - 14\tabcolsep) * \real{0.1250}}
  >{\raggedright\arraybackslash}p{(\linewidth - 14\tabcolsep) * \real{0.1250}}@{}}
\toprule\noalign{}
\begin{minipage}[b]{\linewidth}\raggedright
RUNSPEC
\end{minipage} & \begin{minipage}[b]{\linewidth}\raggedright
GRID
\end{minipage} & \begin{minipage}[b]{\linewidth}\raggedright
EDIT
\end{minipage} & \begin{minipage}[b]{\linewidth}\raggedright
PROPS
\end{minipage} & \begin{minipage}[b]{\linewidth}\raggedright
REGIONS
\end{minipage} & \begin{minipage}[b]{\linewidth}\raggedright
SOLUTION
\end{minipage} & \begin{minipage}[b]{\linewidth}\raggedright
SUMMARY
\end{minipage} & \begin{minipage}[b]{\linewidth}\raggedright
SCHEDULE
\end{minipage} \\
\midrule\noalign{}
\endhead
\bottomrule\noalign{}
\endlastfoot
\end{longtable}

\subsubsection{Description}\label{description-222}

DZNET defines the net thickness of all grid blocks in the Z direction
via an array for each cell in a Cartesian Regular Grid model.

\begin{longtable}[]{@{}
  >{\raggedright\arraybackslash}p{(\linewidth - 6\tabcolsep) * \real{0.2500}}
  >{\raggedright\arraybackslash}p{(\linewidth - 6\tabcolsep) * \real{0.2500}}
  >{\raggedright\arraybackslash}p{(\linewidth - 6\tabcolsep) * \real{0.2500}}
  >{\raggedright\arraybackslash}p{(\linewidth - 6\tabcolsep) * \real{0.2500}}@{}}
\toprule\noalign{}
\begin{minipage}[b]{\linewidth}\raggedright
No.
\end{minipage} & \begin{minipage}[b]{\linewidth}\raggedright
Name
\end{minipage} & \begin{minipage}[b]{\linewidth}\raggedright
Description
\end{minipage} & \begin{minipage}[b]{\linewidth}\raggedright
Default
\end{minipage} \\
\midrule\noalign{}
\endhead
\bottomrule\noalign{}
\endlastfoot
Field & Metric & Laboratory & \\
1 & DZNET & DZNET is an array of real numbers describing the net
thickness in the Z direction for each cell in the model. Repeat counts
may be used, for example 10*100.0. If the value for a grid block is not
defined then the grid block size (DZ) is assigned to the missing values.
& DZ \\
feet & m & cm & \\
Notes: & & & \\
\end{longtable}

\emph{Table 6.31: DZNET Keyword Description}

See also the DX, DY, DZ, NTG and TOPS keywords to fully define a
Cartesian Regular Grid.

\subsubsection{Example}\label{example-155}

\begin{verbatim}
--
--       DEFINE GRID BLOCK Z DIRECTION NET THICKNESS(BASED ON NX x NY x NZ = 300)
--
DZNET
         100*15.0   100*25.0   00*45.0                                          /
\end{verbatim}

The above example defines the net thickness of the cells in the Z
direction based on 300 cells in the model as defined by the DIMENS
keyword in the RUNSPEC section.

\subsection{DZV -- Define the Size of Grid Blocks in the Z Direction via
a
Vector}\label{dzv-define-the-size-of-grid-blocks-in-the-z-direction-via-a-vector}

\begin{longtable}[]{@{}
  >{\raggedright\arraybackslash}p{(\linewidth - 14\tabcolsep) * \real{0.1250}}
  >{\raggedright\arraybackslash}p{(\linewidth - 14\tabcolsep) * \real{0.1250}}
  >{\raggedright\arraybackslash}p{(\linewidth - 14\tabcolsep) * \real{0.1250}}
  >{\raggedright\arraybackslash}p{(\linewidth - 14\tabcolsep) * \real{0.1250}}
  >{\raggedright\arraybackslash}p{(\linewidth - 14\tabcolsep) * \real{0.1250}}
  >{\raggedright\arraybackslash}p{(\linewidth - 14\tabcolsep) * \real{0.1250}}
  >{\raggedright\arraybackslash}p{(\linewidth - 14\tabcolsep) * \real{0.1250}}
  >{\raggedright\arraybackslash}p{(\linewidth - 14\tabcolsep) * \real{0.1250}}@{}}
\toprule\noalign{}
\begin{minipage}[b]{\linewidth}\raggedright
RUNSPEC
\end{minipage} & \begin{minipage}[b]{\linewidth}\raggedright
GRID
\end{minipage} & \begin{minipage}[b]{\linewidth}\raggedright
EDIT
\end{minipage} & \begin{minipage}[b]{\linewidth}\raggedright
PROPS
\end{minipage} & \begin{minipage}[b]{\linewidth}\raggedright
REGIONS
\end{minipage} & \begin{minipage}[b]{\linewidth}\raggedright
SOLUTION
\end{minipage} & \begin{minipage}[b]{\linewidth}\raggedright
SUMMARY
\end{minipage} & \begin{minipage}[b]{\linewidth}\raggedright
SCHEDULE
\end{minipage} \\
\midrule\noalign{}
\endhead
\bottomrule\noalign{}
\endlastfoot
\end{longtable}

\subsubsection{Description}\label{description-223}

DZV defines the size of grid blocks in the Z direction via a vector as
opposed to defining the thickness property for each cell. The keyword is
used for both Cartesian Regular Grids and Radial Grids.

\begin{longtable}[]{@{}
  >{\raggedright\arraybackslash}p{(\linewidth - 6\tabcolsep) * \real{0.2500}}
  >{\raggedright\arraybackslash}p{(\linewidth - 6\tabcolsep) * \real{0.2500}}
  >{\raggedright\arraybackslash}p{(\linewidth - 6\tabcolsep) * \real{0.2500}}
  >{\raggedright\arraybackslash}p{(\linewidth - 6\tabcolsep) * \real{0.2500}}@{}}
\toprule\noalign{}
\begin{minipage}[b]{\linewidth}\raggedright
No.
\end{minipage} & \begin{minipage}[b]{\linewidth}\raggedright
Name
\end{minipage} & \begin{minipage}[b]{\linewidth}\raggedright
Description
\end{minipage} & \begin{minipage}[b]{\linewidth}\raggedright
Default
\end{minipage} \\
\midrule\noalign{}
\endhead
\bottomrule\noalign{}
\endlastfoot
Field & Metric & Laboratory & \\
1 & DZV & DZV is a vector of real numbers describing the cell size for
the grid blocks in the Z direction. Repeat counts may be used, for
example 10*20.0. & None \\
feet & m & cm & \\
Notes: & & & \\
\end{longtable}

\emph{Table 6.32: DZV Keyword Description}

See also the DXV, DYV and TOPS keywords for a Cartesian Regular Grid and
DRV, DTHETAV and TOPS keywords to fully define a Radial Grid model.

\subsubsection{Example}\label{example-156}

\begin{verbatim}
--
--       DEFINE GRID BLOCK SIZES IN THE Z DIRECTION (BASED ON NZ = 20)
--
DZV
         3.0   5.0   3.0   2.0   5.0  15*3.0                                    /
\end{verbatim}

The above example defines the size of the cells in the Z direction based
on NZ equals 20 on the DIMENS keyword in the RUNSPEC section.

\subsection{ECHO -- Activate Echoing of User Input Files to the Print
File}\label{echo-activate-echoing-of-user-input-files-to-the-print-file-2}

Turns on echoing of all the input files to the print file; note that
this keyword is activated by default and can subsequently be switched
off by the NOECHO activation keyword.

See ECHO -- Activate Echoing of User Input Files to the Print File in
the GLOBAL section for a full description.

\subsection{END -- Define the End of the Input
File}\label{end-define-the-end-of-the-input-file-2}

This keyword marks the end of the input file and can occur in any
section. Any keywords and data after the END keyword are ignored.

See END -- Define the End of the Input File in the GLOBAL section for a
full description.

\subsection{ENDBOX -- Define the End of the BOX Defined
Grid}\label{endbox-define-the-end-of-the-box-defined-grid}

\begin{longtable}[]{@{}
  >{\raggedright\arraybackslash}p{(\linewidth - 14\tabcolsep) * \real{0.1250}}
  >{\raggedright\arraybackslash}p{(\linewidth - 14\tabcolsep) * \real{0.1250}}
  >{\raggedright\arraybackslash}p{(\linewidth - 14\tabcolsep) * \real{0.1250}}
  >{\raggedright\arraybackslash}p{(\linewidth - 14\tabcolsep) * \real{0.1250}}
  >{\raggedright\arraybackslash}p{(\linewidth - 14\tabcolsep) * \real{0.1250}}
  >{\raggedright\arraybackslash}p{(\linewidth - 14\tabcolsep) * \real{0.1250}}
  >{\raggedright\arraybackslash}p{(\linewidth - 14\tabcolsep) * \real{0.1250}}
  >{\raggedright\arraybackslash}p{(\linewidth - 14\tabcolsep) * \real{0.1250}}@{}}
\toprule\noalign{}
\begin{minipage}[b]{\linewidth}\raggedright
RUNSPEC
\end{minipage} & \begin{minipage}[b]{\linewidth}\raggedright
GRID
\end{minipage} & \begin{minipage}[b]{\linewidth}\raggedright
EDIT
\end{minipage} & \begin{minipage}[b]{\linewidth}\raggedright
PROPS
\end{minipage} & \begin{minipage}[b]{\linewidth}\raggedright
REGIONS
\end{minipage} & \begin{minipage}[b]{\linewidth}\raggedright
SOLUTION
\end{minipage} & \begin{minipage}[b]{\linewidth}\raggedright
SUMMARY
\end{minipage} & \begin{minipage}[b]{\linewidth}\raggedright
SCHEDULE
\end{minipage} \\
\midrule\noalign{}
\endhead
\bottomrule\noalign{}
\endlastfoot
\end{longtable}

\subsubsection{Description}\label{description-224}

This keyword marks the end of a previously defined BOX sub-grid as
defined by a previously entered BOX keyword. The keyword resets the
input grid to be the full grid as defined by the NX, NY, and NZ
variables on the DIMENS keyword in the RUNSPEC section.

There is no data required for this keyword.

\subsubsection{Example}\label{example-157}

\begin{verbatim}
--
--       DEFINE INPUT BOX FOR EDITING INPUT ARRAYS
--
--       ---------- BOX ---------
--       I1  I2   J1  J2   K1  K2
BOX
         10  10   1   6    1   1                           / DEFINE BOX AREA
--
--
--       DEFINE GRID BLOCK PERMZ DATA FOR THE INPUT BOX
--
PERMZ
         6*0.01                                            /
--
--       DEFINE END OF INPUT BOX EDITING OF INPUT ARRAYS
--
ENDBOX
\end{verbatim}

The above example defines a subset of the grid and sets the cells PERMZ
values to 0.01 for that area. After which the ENDBOX keyword resets the
input to be the full grid.

\begin{tcolorbox}[enhanced jigsaw, title=\textcolor{quarto-callout-note-color}{\faInfo}\hspace{0.5em}{Note}, coltitle=black, breakable, toprule=.15mm, colbacktitle=quarto-callout-note-color!10!white, leftrule=.75mm, opacitybacktitle=0.6, colback=white, bottomtitle=1mm, toptitle=1mm, titlerule=0mm, left=2mm, colframe=quarto-callout-note-color-frame, arc=.35mm, rightrule=.15mm, bottomrule=.15mm, opacityback=0]

It is good practice to always use the ENDBOX keyword to reset the input
back to the full grid when all the modifications for a sub-grid have
been completed.

\end{tcolorbox}

\subsection{ENDFIN -- End the Definition of a Local Grid
Refinement}\label{endfin-end-the-definition-of-a-local-grid-refinement}

\begin{longtable}[]{@{}
  >{\raggedright\arraybackslash}p{(\linewidth - 14\tabcolsep) * \real{0.1250}}
  >{\raggedright\arraybackslash}p{(\linewidth - 14\tabcolsep) * \real{0.1250}}
  >{\raggedright\arraybackslash}p{(\linewidth - 14\tabcolsep) * \real{0.1250}}
  >{\raggedright\arraybackslash}p{(\linewidth - 14\tabcolsep) * \real{0.1250}}
  >{\raggedright\arraybackslash}p{(\linewidth - 14\tabcolsep) * \real{0.1250}}
  >{\raggedright\arraybackslash}p{(\linewidth - 14\tabcolsep) * \real{0.1250}}
  >{\raggedright\arraybackslash}p{(\linewidth - 14\tabcolsep) * \real{0.1250}}
  >{\raggedright\arraybackslash}p{(\linewidth - 14\tabcolsep) * \real{0.1250}}@{}}
\toprule\noalign{}
\begin{minipage}[b]{\linewidth}\raggedright
RUNSPEC
\end{minipage} & \begin{minipage}[b]{\linewidth}\raggedright
GRID
\end{minipage} & \begin{minipage}[b]{\linewidth}\raggedright
EDIT
\end{minipage} & \begin{minipage}[b]{\linewidth}\raggedright
PROPS
\end{minipage} & \begin{minipage}[b]{\linewidth}\raggedright
REGIONS
\end{minipage} & \begin{minipage}[b]{\linewidth}\raggedright
SOLUTION
\end{minipage} & \begin{minipage}[b]{\linewidth}\raggedright
SUMMARY
\end{minipage} & \begin{minipage}[b]{\linewidth}\raggedright
SCHEDULE
\end{minipage} \\
\midrule\noalign{}
\endhead
\bottomrule\noalign{}
\endlastfoot
\end{longtable}

\subsubsection{Description}\label{description-225}

The ENDFIN keyword defines the end of a Cartesian or radial local grid
refinement (``LGR'') definition and a LGR property definition data set.
In the GRID section the CARFIN, RADFIN, and RADFIN4 keywords defines the
start of an LGR description section, whereas the REFINE keyword in the
EDIT, PROPS, REGIONS, SOLUTION and SCHEDULE section defines the start.
The REFINE keyword can also be used in the GRID section provided the LGR
has been previously specified by the CARFIN, RADFIN, or RADFIN4
keywords.

This keyword is not supported by OPM Flow but would change the results
if supported so the simulation will be stopped.

There is no data required for this keyword and there is no terminating
``/'' for this keyword.

\subsubsection{Example}\label{example-158}

The example below is based on using the CARFIN keyword in the GRID
section to define an LGR in the global grid, named LGR-OP01 with a
maximum of one well allowed in the LGR.

\begin{verbatim}
--
--       CARFIN LGR GRID COMMANDS
--
--       LGR        ----- FINE GRID ------   -- CARFIN GRID --  MAX     HOST
--       NAME       I1  I2  J1  J2  K1  K2     NX    NY    NZ   WELLS   NAME
CARFIN
         LGR-OP01   24  24  87  87   1  50      3     3    50     1     GLOBAL /


ENDFIN
\end{verbatim}

Here the one global cell in the areal plane (24, 87) is divided into
three LGR cells in the x-direction and three cells in the y-direction.
Since no other property data is given, then the LGR cells take their
properties from the host grid, that is the global grid.

\subsection{ENDINC -- Define the End of an Include
File}\label{endinc-define-the-end-of-an-include-file-2}

This keyword marks the end of an include file specified on the INCLUDE
keyword. When the ENDINC keyword is encountered in the INCLUDE file,
input data is read from the next keyword in the current file. Any
keywords and data after the ENDINC keyword in the INCLUDE file are
ignored.

See ENDINC -- Define the End of an Include File in the GLOBAL section
for a full description.

\subsection{ENDSKIP -- Deactivate Skipping of Keywords and Input
Data}\label{endskip-deactivate-skipping-of-keywords-and-input-data-2}

The ENDSKIP keyword deactivates the skipping of keywords that was
activated by the SKIP, SKIP100, or SKIP300 keywords. Each SKIP keyword
should be paired with an ENDSKIP keyword.

See ENDSKIP -- Deactivate Skipping of Keywords and Input Data in the
GLOBAL section for a full description.

\subsection{EQLZCORN -- Modify the Depth of the Corner-Point Depth
Array}\label{eqlzcorn-modify-the-depth-of-the-corner-point-depth-array}

\begin{longtable}[]{@{}
  >{\raggedright\arraybackslash}p{(\linewidth - 14\tabcolsep) * \real{0.1250}}
  >{\raggedright\arraybackslash}p{(\linewidth - 14\tabcolsep) * \real{0.1250}}
  >{\raggedright\arraybackslash}p{(\linewidth - 14\tabcolsep) * \real{0.1250}}
  >{\raggedright\arraybackslash}p{(\linewidth - 14\tabcolsep) * \real{0.1250}}
  >{\raggedright\arraybackslash}p{(\linewidth - 14\tabcolsep) * \real{0.1250}}
  >{\raggedright\arraybackslash}p{(\linewidth - 14\tabcolsep) * \real{0.1250}}
  >{\raggedright\arraybackslash}p{(\linewidth - 14\tabcolsep) * \real{0.1250}}
  >{\raggedright\arraybackslash}p{(\linewidth - 14\tabcolsep) * \real{0.1250}}@{}}
\toprule\noalign{}
\begin{minipage}[b]{\linewidth}\raggedright
RUNSPEC
\end{minipage} & \begin{minipage}[b]{\linewidth}\raggedright
GRID
\end{minipage} & \begin{minipage}[b]{\linewidth}\raggedright
EDIT
\end{minipage} & \begin{minipage}[b]{\linewidth}\raggedright
PROPS
\end{minipage} & \begin{minipage}[b]{\linewidth}\raggedright
REGIONS
\end{minipage} & \begin{minipage}[b]{\linewidth}\raggedright
SOLUTION
\end{minipage} & \begin{minipage}[b]{\linewidth}\raggedright
SUMMARY
\end{minipage} & \begin{minipage}[b]{\linewidth}\raggedright
SCHEDULE
\end{minipage} \\
\midrule\noalign{}
\endhead
\bottomrule\noalign{}
\endlastfoot
\end{longtable}

\subsubsection{Description}\label{description-226}

The EQLZCORN keyword modifies the depth of a corner point of a grid
block on the pillars defining the reservoir grid. The keyword can be
only used be used with Irregular Corner-Point Grids.

This keyword is not supported by OPM Flow but would change the results
if supported so the simulation will be stopped.

\subsection{EQUALREG -- Sets an Array to a Constant by Region
Number}\label{equalreg-sets-an-array-to-a-constant-by-region-number}

\begin{longtable}[]{@{}
  >{\raggedright\arraybackslash}p{(\linewidth - 14\tabcolsep) * \real{0.1250}}
  >{\raggedright\arraybackslash}p{(\linewidth - 14\tabcolsep) * \real{0.1250}}
  >{\raggedright\arraybackslash}p{(\linewidth - 14\tabcolsep) * \real{0.1250}}
  >{\raggedright\arraybackslash}p{(\linewidth - 14\tabcolsep) * \real{0.1250}}
  >{\raggedright\arraybackslash}p{(\linewidth - 14\tabcolsep) * \real{0.1250}}
  >{\raggedright\arraybackslash}p{(\linewidth - 14\tabcolsep) * \real{0.1250}}
  >{\raggedright\arraybackslash}p{(\linewidth - 14\tabcolsep) * \real{0.1250}}
  >{\raggedright\arraybackslash}p{(\linewidth - 14\tabcolsep) * \real{0.1250}}@{}}
\toprule\noalign{}
\begin{minipage}[b]{\linewidth}\raggedright
RUNSPEC
\end{minipage} & \begin{minipage}[b]{\linewidth}\raggedright
GRID
\end{minipage} & \begin{minipage}[b]{\linewidth}\raggedright
EDIT
\end{minipage} & \begin{minipage}[b]{\linewidth}\raggedright
PROPS
\end{minipage} & \begin{minipage}[b]{\linewidth}\raggedright
REGIONS
\end{minipage} & \begin{minipage}[b]{\linewidth}\raggedright
SOLUTION
\end{minipage} & \begin{minipage}[b]{\linewidth}\raggedright
SUMMARY
\end{minipage} & \begin{minipage}[b]{\linewidth}\raggedright
SCHEDULE
\end{minipage} \\
\midrule\noalign{}
\endhead
\bottomrule\noalign{}
\endlastfoot
\end{longtable}

\subsubsection{Description}\label{description-227}

The EQUALREG keyword sets a specified array to a constant for cells with
a specific region number. The region number array can be FLUXNUM,
MULTNUM or OPERNUM and these arrays must be defined and be available
before the EQUALREG keyword is read by the simulator. The constant can
be an integer or real value depending on the array type; however, the
arrays that can be operated on are dependent on which section the
EQUALREG keyword is being applied in.

This keyword is not supported by OPM Flow but would change the results
if supported so the simulation will be stopped.

\begin{longtable}[]{@{}
  >{\raggedright\arraybackslash}p{(\linewidth - 6\tabcolsep) * \real{0.2500}}
  >{\raggedright\arraybackslash}p{(\linewidth - 6\tabcolsep) * \real{0.2500}}
  >{\raggedright\arraybackslash}p{(\linewidth - 6\tabcolsep) * \real{0.2500}}
  >{\raggedright\arraybackslash}p{(\linewidth - 6\tabcolsep) * \real{0.2500}}@{}}
\toprule\noalign{}
\begin{minipage}[b]{\linewidth}\raggedright
No.
\end{minipage} & \begin{minipage}[b]{\linewidth}\raggedright
Name
\end{minipage} & \begin{minipage}[b]{\linewidth}\raggedright
Description
\end{minipage} & \begin{minipage}[b]{\linewidth}\raggedright
Default
\end{minipage} \\
\midrule\noalign{}
\endhead
\bottomrule\noalign{}
\endlastfoot
1 & ARRAY & A character string of up to eight characters in length that
defines the keyword identifying the array to be modified. & None \\
2 & CONSTANT & An integer or real value to be assigned to the ARRAY in
the same units as the ARRAY property for a given REGION & 0 \\
3 & REGION NUMBER & REGION NUMBER is a positive integer representing the
region for which the CONSTANT in (2) should be applied & None \\
4 & REGION ARRAY & The REGION ARRAY to use for applying the CONSTANT in
(2) based on the REGION NUMBER in (3). REGION ARRAY can have the
following values: & M \\
Notes: & & & \\
\end{longtable}

\emph{Table 6.33: EQUALREG Keyword Description}

Examples of the arrays most commonly operated on in each section are
given in Table 6.34. Cells colored red indicate arrays that are not
supported by OPM Flow operations.

\begin{longtable}[]{@{}lllllll@{}}
\toprule\noalign{}
EQUALREG Keyword and Variable Options by Section & & & & & & \\
\midrule\noalign{}
\endhead
\bottomrule\noalign{}
\endlastfoot
GRID & EDIT & PROPS & REGIONS & SOLUTION & SUMMARY & SCHEDULE \\
DX & DEPTH & SWL & ENDNUM & PRESSURE & & \\
DY & PORV & SWCR & EQLNUM & SWAT & & \\
DZ & TRANX & SWU & FIPNUM & SGAS & & \\
PERMX & TRANY & SGL & IMBNUM & RV & & \\
PERMY & TRANZ & SGCR & MISCNUM & RS & & \\
PERMZ & DIFFX & SGU & PVTNUM & TBLK & & \\
MULTX & DIFFY & KRW & ROCKNUM & GI & & \\
MULTY & DIFFZ & KRO & SATNUM & OILAPI & & \\
MULTZ & TRANR & KRG & WH2NUM & SALT & & \\
DR & TRANTHT & PCG & & GASCONC & & \\
DTHETA & DIFFR & PCW & & SOLVCONC & & \\
PERMR & DIFFTHT & & & SOLVFRAC & & \\
PERMTHT & & & & SFOAM & & \\
DZNET & & & & SPOLY & & \\
PORO & & & & & & \\
NTG & & & & & & \\
FLUXNUM & & & & & & \\
MULTNUM & & & & & & \\
MPFANUM & & & & & & \\
DIFFX & & & & & & \\
DIFFY & & & & & & \\
DIFFZ & & & & & & \\
DIFFR & & & & & & \\
DIFFTHT & & & & & & \\
\end{longtable}

\emph{Table 6.34: EQUALREG Keyword Applicable Arrays by Section}

\subsubsection{Example}\label{example-159}

\begin{verbatim}
-- FIRST DEFINE MULTNUM ARRAYS FOR 10 X 10 X 20 MODEL
--
--      ARRAY     CONSTANT    ---------- BOX ---------
--                            I1  I2   J1  J2   K1  K2
EQUALS
        MULTNUM    1          1*  1*   1*  1*   1*  1*     / MULTNUM IN MODEL
        MULTNUM    2          1*  1*   1*  1*   6   6      / MULTNUM IN MODEL
        MULTNUM    3          1*  1*   1*  1*   10  10     / MULTNUM IN MODEL
/
--    NOW SET PORO AND PERMX BASED ON THE MULTNUM REGION NUMBER
--
--       SETS A CONSTANT TO AN ARRAY BASED ON A REGION NUMBER
--
--       ARRAY     CONSTANT  REGION   REGION ARRAY
--                 VALUE     NUMBER    M / F / O                                           EQUALREG
         PORO      0.200     1         M                   /
         PORO      0.150     2         M                   /
         PORO      0.120     3         M                   /
         PERMX    100.00     1         M                   /
         PERMX     75.00     2         M                   /
         PERMX     50.00     3         M                   /
/
\end{verbatim}

The example first defines the MULTNUM array to 1 for all cells in the
model, after which selected areas of model are assigned various MULTNUM
integer values. The EQUALREG can then be invoked to set constant values
for the PORO and PERMX arrays for the various MULTNUM regions.

\subsection{EQUALS -- Sets a Specified Array to a
Constant}\label{equals-sets-a-specified-array-to-a-constant}

\begin{longtable}[]{@{}
  >{\raggedright\arraybackslash}p{(\linewidth - 14\tabcolsep) * \real{0.1250}}
  >{\raggedright\arraybackslash}p{(\linewidth - 14\tabcolsep) * \real{0.1250}}
  >{\raggedright\arraybackslash}p{(\linewidth - 14\tabcolsep) * \real{0.1250}}
  >{\raggedright\arraybackslash}p{(\linewidth - 14\tabcolsep) * \real{0.1250}}
  >{\raggedright\arraybackslash}p{(\linewidth - 14\tabcolsep) * \real{0.1250}}
  >{\raggedright\arraybackslash}p{(\linewidth - 14\tabcolsep) * \real{0.1250}}
  >{\raggedright\arraybackslash}p{(\linewidth - 14\tabcolsep) * \real{0.1250}}
  >{\raggedright\arraybackslash}p{(\linewidth - 14\tabcolsep) * \real{0.1250}}@{}}
\toprule\noalign{}
\begin{minipage}[b]{\linewidth}\raggedright
RUNSPEC
\end{minipage} & \begin{minipage}[b]{\linewidth}\raggedright
GRID
\end{minipage} & \begin{minipage}[b]{\linewidth}\raggedright
EDIT
\end{minipage} & \begin{minipage}[b]{\linewidth}\raggedright
PROPS
\end{minipage} & \begin{minipage}[b]{\linewidth}\raggedright
REGIONS
\end{minipage} & \begin{minipage}[b]{\linewidth}\raggedright
SOLUTION
\end{minipage} & \begin{minipage}[b]{\linewidth}\raggedright
SUMMARY
\end{minipage} & \begin{minipage}[b]{\linewidth}\raggedright
SCHEDULE
\end{minipage} \\
\midrule\noalign{}
\endhead
\bottomrule\noalign{}
\endlastfoot
\end{longtable}

\subsubsection{Description}\label{description-228}

The EQUALS keyword sets a specified array or part of an array to a
constant. The constant can be an integer or real value depending on the
array type; however, the arrays that can be operated on are dependent on
which section the EQUALS keyword is being applied in.

\begin{longtable}[]{@{}
  >{\raggedright\arraybackslash}p{(\linewidth - 6\tabcolsep) * \real{0.2500}}
  >{\raggedright\arraybackslash}p{(\linewidth - 6\tabcolsep) * \real{0.2500}}
  >{\raggedright\arraybackslash}p{(\linewidth - 6\tabcolsep) * \real{0.2500}}
  >{\raggedright\arraybackslash}p{(\linewidth - 6\tabcolsep) * \real{0.2500}}@{}}
\toprule\noalign{}
\begin{minipage}[b]{\linewidth}\raggedright
No.
\end{minipage} & \begin{minipage}[b]{\linewidth}\raggedright
Name
\end{minipage} & \begin{minipage}[b]{\linewidth}\raggedright
Description
\end{minipage} & \begin{minipage}[b]{\linewidth}\raggedright
Default
\end{minipage} \\
\midrule\noalign{}
\endhead
\bottomrule\noalign{}
\endlastfoot
1 & ARRAY & A character string of up to eight characters in length that
defines the keyword identifying the property to be modified. & None \\
2 & CONSTANT & An integer or real value that the ARRAY will be set to in
the same units as the ARRAY property. & None \\
3 & I1 & A positive integer that defines the lower bound of the array in
the I-direction to be modified must be greater than or equal to one and
less than or equal to I2 and NX. & 1 \\
4 & I2 & A positive integer that defines the upper bound of the array in
the I-direction to be modified must be greater than or equal to II and
less than or equal to NX & NX \\
5 & J1 & A positive integer that defines the lower bound of the array in
the J-direction to be modified must be greater than or equal to one and
less than or equal to J2 and NY. & 1 \\
6 & J2 & A positive integer that defines the upper bound of the array in
the J-direction to be modified must be greater than or equal to JI and
less than or equal to NY. & NY \\
7 & K1 & A positive integer that defines the lower bound of the array in
the K-direction to be modified must be greater than or equal to one and
less than or equal to K2 and NZ. & 1 \\
8 & K2 & A positive integer that defines the upper bound of the array in
the K-direction to be modified must be greater than or equal to KI and
less than or equal to NZ. & NZ \\
Notes: & & & \\
\end{longtable}

\emph{Table 6.35: EQUALS Keyword Description}

Examples of the arrays most commonly operated on in each section are
given in Table 6.36. Cells colored red indicate arrays that are not
supported by OPM Flow operations.

\begin{longtable}[]{@{}lllllll@{}}
\toprule\noalign{}
EQUALS Keyword and Variable Options by Section & & & & & & \\
\midrule\noalign{}
\endhead
\bottomrule\noalign{}
\endlastfoot
GRID & EDIT & PROPS & REGIONS & SOLUTION & SUMMARY & SCHEDULE \\
DX & DEPTH & SWL & ENDNUM & PRESSURE & & \\
DY & PORV & SWCR & EQLNUM & SWAT & & \\
DZ & TRANX & SWU & FIPNUM & SGAS & & \\
PERMX & TRANY & SGL & IMBNUM & RV & & \\
PERMY & TRANZ & SGCR & MISCNUM & RS & & \\
PERMZ & DIFFX & SGU & PVTNUM & TBLK & & \\
MULTX & DIFFY & KRW & ROCKNUM & GI & & \\
MULTY & DIFFZ & KRO & SATNUM & OILAPI & & \\
MULTZ & TRANR & KRG & WH2NUM & SALT & & \\
DR & TRANTHT & PCG & & GASCONC & & \\
DTHETA & DIFFR & PCW & & SOLVCONC & & \\
PERMR & DIFFTHT & & & SOLVFRAC & & \\
PERMTHT & & & & SFOAM & & \\
DZNET & & & & SPOLY & & \\
PORO & & & & & & \\
NTG & & & & & & \\
FLUXNUM & & & & & & \\
MULTNUM & & & & & & \\
MPFANUM & & & & & & \\
DIFFX & & & & & & \\
DIFFY & & & & & & \\
DIFFZ & & & & & & \\
DIFFR & & & & & & \\
DIFFTHT & & & & & & \\
\end{longtable}

\emph{Table 6.36: EQUALS Keyword Applicable Arrays by Section}

\begin{tcolorbox}[enhanced jigsaw, title=\textcolor{quarto-callout-note-color}{\faInfo}\hspace{0.5em}{Note}, coltitle=black, breakable, toprule=.15mm, colbacktitle=quarto-callout-note-color!10!white, leftrule=.75mm, opacitybacktitle=0.6, colback=white, bottomtitle=1mm, toptitle=1mm, titlerule=0mm, left=2mm, colframe=quarto-callout-note-color-frame, arc=.35mm, rightrule=.15mm, bottomrule=.15mm, opacityback=0]

Unlike the commercial simulator, using the EQUALS keyword to setup the
structure of the grid using the DX, DY, DZ and TOPS keywords should be
avoided as it may cause OPM Flow to prematurely fail during the
initialization. See the second example on the correct way to setup this
type of grid.

\end{tcolorbox}

\subsubsection{Examples}\label{examples-9}

The first example resets the PERMX, PERMY and PERMZ, arrays to 0.10,
0.10, and 0.01 for all cells in layer five, respectively.

\begin{verbatim}
--
--  ARRAY     CONSTANT      ---------- BOX ---------
-                           I1  I2   J1  J2   K1  K2
EQUALS
    PERMX     0.1000        1*  1*   1*  1*   5   5  / PERMX TO 0.10 IN LAYER 5
    PERMY     0.1000        1*  1*   1*  1*   5   5  / PERMY TO 0.10 IN LAYER 5
    PERMZ     0.0100        1*  1*   1*  1*   5   5  / PERMZ TO 0.01 IN LAYER 5
/
\end{verbatim}

The second example illustrates how to correctly setup a Cartesian
Regular Grid in OPM Flow, given the DIMENS keyword in the RUNSPEC
section is set to:

\begin{verbatim}
--
--       MAX     MAX     MAX
--       NDIVIX  NDIVIY  NDIVIZ
DIMENS
         10       10      3                                                    /
\end{verbatim}

and the following keywords in the GRID section:

\begin{verbatim}
--
--       ACTIVATE IRREGULAR CORNER-POINT GRID TRANSMISSIBILITIES
--
OLDTRAN
--
--       DEFINE GRID BLOCK X DIRECTION CELL SIZE (BASED ON NX = 5)
--
DX
         300*1000                                                              /
--
--       DEFINE GRID BLOCK Y DIRECTION CELL SIZE (BASED ON NY = 5)
--
DY
         300*1000                                                              /
--
--       DEFINE GRID BLOCK SIZES IN THE Z DIRECTION
--
DZ
         100*20  100*30  100*50                                                /
--
--       DEFINE GRID BLOCK TOPS FOR THE TOP LAYER
--
TOPS
         100*8325                                                              /
--
--       ARRAY    CONSTANT        ---------- BOX ---------
--                                I1  I2   J1  J2   K1  K2
EQUALS
         PERMX     500.000        1*  1*   1*  1*    1   1 / Layer #01 Properties
         PERMY     500.000                                 /
         PERMZ      20.000                                 /
         PORO        0.300                                 /
         NTG         1.000                                 /

         PERMX      50.000        1*  1*   1*  1*    2   2 / Layer #02 Properties
         PERMY      50.000                                 /
         PERMZ      50.000                                 /
         PORO        0.300                                 /
         NTG         1.000                                 /

         PERMX     200.000        1*  1*   1*  1*    3   3 / Layer #03 Properties
         PERMY     200.000                                 /
         PERMZ     200.000                                 /
         PORO        0.300                                 /
         NTG         1.000                                 /
/
\end{verbatim}

Notice that the DX, DY, DZ and TOPS keywords are defined separately,
that is they are not included in the EQUALS keyword.

\subsection{EXTFIN -- Define an External Unstructured Local Grid
Refinement}\label{extfin-define-an-external-unstructured-local-grid-refinement}

\begin{longtable}[]{@{}
  >{\raggedright\arraybackslash}p{(\linewidth - 14\tabcolsep) * \real{0.1250}}
  >{\raggedright\arraybackslash}p{(\linewidth - 14\tabcolsep) * \real{0.1250}}
  >{\raggedright\arraybackslash}p{(\linewidth - 14\tabcolsep) * \real{0.1250}}
  >{\raggedright\arraybackslash}p{(\linewidth - 14\tabcolsep) * \real{0.1250}}
  >{\raggedright\arraybackslash}p{(\linewidth - 14\tabcolsep) * \real{0.1250}}
  >{\raggedright\arraybackslash}p{(\linewidth - 14\tabcolsep) * \real{0.1250}}
  >{\raggedright\arraybackslash}p{(\linewidth - 14\tabcolsep) * \real{0.1250}}
  >{\raggedright\arraybackslash}p{(\linewidth - 14\tabcolsep) * \real{0.1250}}@{}}
\toprule\noalign{}
\begin{minipage}[b]{\linewidth}\raggedright
RUNSPEC
\end{minipage} & \begin{minipage}[b]{\linewidth}\raggedright
GRID
\end{minipage} & \begin{minipage}[b]{\linewidth}\raggedright
EDIT
\end{minipage} & \begin{minipage}[b]{\linewidth}\raggedright
PROPS
\end{minipage} & \begin{minipage}[b]{\linewidth}\raggedright
REGIONS
\end{minipage} & \begin{minipage}[b]{\linewidth}\raggedright
SOLUTION
\end{minipage} & \begin{minipage}[b]{\linewidth}\raggedright
SUMMARY
\end{minipage} & \begin{minipage}[b]{\linewidth}\raggedright
SCHEDULE
\end{minipage} \\
\midrule\noalign{}
\endhead
\bottomrule\noalign{}
\endlastfoot
\end{longtable}

\subsubsection{Description}\label{description-229}

The EXTFIN keyword defines an external Unstructured Local Grid
Refinement (``LGR'') in a cell or a group of cells in the global grid,
and for when LGRs have been activated for the model using the LGR
keyword in the RUNSPEC section. Note the global grid can be either
structured, see the EXTREPGL keyword in the GRID section for global
structure grids, or unstructured, see the EXTHOST keyword in the GRID
section for unstructured global grids.

This keyword is not supported by OPM Flow but would change the results
if supported so the simulation will be stopped.

\subsection{EXTHOST -- Define Host Cells for External
LGRs}\label{exthost-define-host-cells-for-external-lgrs}

\begin{longtable}[]{@{}
  >{\raggedright\arraybackslash}p{(\linewidth - 14\tabcolsep) * \real{0.1250}}
  >{\raggedright\arraybackslash}p{(\linewidth - 14\tabcolsep) * \real{0.1250}}
  >{\raggedright\arraybackslash}p{(\linewidth - 14\tabcolsep) * \real{0.1250}}
  >{\raggedright\arraybackslash}p{(\linewidth - 14\tabcolsep) * \real{0.1250}}
  >{\raggedright\arraybackslash}p{(\linewidth - 14\tabcolsep) * \real{0.1250}}
  >{\raggedright\arraybackslash}p{(\linewidth - 14\tabcolsep) * \real{0.1250}}
  >{\raggedright\arraybackslash}p{(\linewidth - 14\tabcolsep) * \real{0.1250}}
  >{\raggedright\arraybackslash}p{(\linewidth - 14\tabcolsep) * \real{0.1250}}@{}}
\toprule\noalign{}
\begin{minipage}[b]{\linewidth}\raggedright
RUNSPEC
\end{minipage} & \begin{minipage}[b]{\linewidth}\raggedright
GRID
\end{minipage} & \begin{minipage}[b]{\linewidth}\raggedright
EDIT
\end{minipage} & \begin{minipage}[b]{\linewidth}\raggedright
PROPS
\end{minipage} & \begin{minipage}[b]{\linewidth}\raggedright
REGIONS
\end{minipage} & \begin{minipage}[b]{\linewidth}\raggedright
SOLUTION
\end{minipage} & \begin{minipage}[b]{\linewidth}\raggedright
SUMMARY
\end{minipage} & \begin{minipage}[b]{\linewidth}\raggedright
SCHEDULE
\end{minipage} \\
\midrule\noalign{}
\endhead
\bottomrule\noalign{}
\endlastfoot
\end{longtable}

\subsubsection{Description}\label{description-230}

The EXTHOST keyword defines the host global grid blocks for an external
Local Grid Refinement (``LGR'') for when LGRs have been activated for
the model using the LGR keyword in the RUNSPEC section, and the global
grid is an unstructured grid.

This keyword is not supported by OPM Flow but would change the results
if supported so the simulation will be stopped.

\subsection{EXTRAPMS -- Activate Extrapolation Warning
Messages}\label{extrapms-activate-extrapolation-warning-messages-2}

The EXTRAPMS keyword activates extrapolation warning messages for when
OPM Flow extrapolates the PVT or VFP tables. Frequent extrapolation
warning messages should be investigated and resolved as this would
indicate possible incorrect data and may result in the simulator
extrapolating to unrealistic values.

See EXTRAPMS -- Activate Extrapolation Warning Messages in the GLOBAL
section for a full description.

\subsection{EXTREPGL -- Define Host Cells for External Unstructured
LGRs}\label{extrepgl-define-host-cells-for-external-unstructured-lgrs}

\begin{longtable}[]{@{}
  >{\raggedright\arraybackslash}p{(\linewidth - 14\tabcolsep) * \real{0.1250}}
  >{\raggedright\arraybackslash}p{(\linewidth - 14\tabcolsep) * \real{0.1250}}
  >{\raggedright\arraybackslash}p{(\linewidth - 14\tabcolsep) * \real{0.1250}}
  >{\raggedright\arraybackslash}p{(\linewidth - 14\tabcolsep) * \real{0.1250}}
  >{\raggedright\arraybackslash}p{(\linewidth - 14\tabcolsep) * \real{0.1250}}
  >{\raggedright\arraybackslash}p{(\linewidth - 14\tabcolsep) * \real{0.1250}}
  >{\raggedright\arraybackslash}p{(\linewidth - 14\tabcolsep) * \real{0.1250}}
  >{\raggedright\arraybackslash}p{(\linewidth - 14\tabcolsep) * \real{0.1250}}@{}}
\toprule\noalign{}
\begin{minipage}[b]{\linewidth}\raggedright
RUNSPEC
\end{minipage} & \begin{minipage}[b]{\linewidth}\raggedright
GRID
\end{minipage} & \begin{minipage}[b]{\linewidth}\raggedright
EDIT
\end{minipage} & \begin{minipage}[b]{\linewidth}\raggedright
PROPS
\end{minipage} & \begin{minipage}[b]{\linewidth}\raggedright
REGIONS
\end{minipage} & \begin{minipage}[b]{\linewidth}\raggedright
SOLUTION
\end{minipage} & \begin{minipage}[b]{\linewidth}\raggedright
SUMMARY
\end{minipage} & \begin{minipage}[b]{\linewidth}\raggedright
SCHEDULE
\end{minipage} \\
\midrule\noalign{}
\endhead
\bottomrule\noalign{}
\endlastfoot
\end{longtable}

\subsubsection{Description}\label{description-231}

The EXTREPGL keyword defines the host global grid blocks for an external
Unstructured Local Grid Refinement (``LGR'') for when LGRs have been
activated for the model using the LGR keyword in the RUNSPEC section,
and the global grid is a structured grid.

This keyword is not supported by OPM Flow but would change the results
if supported so the simulation will be stopped.

\subsection{FAULTS -- Define Faults in the Grid
Geometry}\label{faults-define-faults-in-the-grid-geometry}

\begin{longtable}[]{@{}
  >{\raggedright\arraybackslash}p{(\linewidth - 14\tabcolsep) * \real{0.1250}}
  >{\raggedright\arraybackslash}p{(\linewidth - 14\tabcolsep) * \real{0.1250}}
  >{\raggedright\arraybackslash}p{(\linewidth - 14\tabcolsep) * \real{0.1250}}
  >{\raggedright\arraybackslash}p{(\linewidth - 14\tabcolsep) * \real{0.1250}}
  >{\raggedright\arraybackslash}p{(\linewidth - 14\tabcolsep) * \real{0.1250}}
  >{\raggedright\arraybackslash}p{(\linewidth - 14\tabcolsep) * \real{0.1250}}
  >{\raggedright\arraybackslash}p{(\linewidth - 14\tabcolsep) * \real{0.1250}}
  >{\raggedright\arraybackslash}p{(\linewidth - 14\tabcolsep) * \real{0.1250}}@{}}
\toprule\noalign{}
\begin{minipage}[b]{\linewidth}\raggedright
RUNSPEC
\end{minipage} & \begin{minipage}[b]{\linewidth}\raggedright
GRID
\end{minipage} & \begin{minipage}[b]{\linewidth}\raggedright
EDIT
\end{minipage} & \begin{minipage}[b]{\linewidth}\raggedright
PROPS
\end{minipage} & \begin{minipage}[b]{\linewidth}\raggedright
REGIONS
\end{minipage} & \begin{minipage}[b]{\linewidth}\raggedright
SOLUTION
\end{minipage} & \begin{minipage}[b]{\linewidth}\raggedright
SUMMARY
\end{minipage} & \begin{minipage}[b]{\linewidth}\raggedright
SCHEDULE
\end{minipage} \\
\midrule\noalign{}
\endhead
\bottomrule\noalign{}
\endlastfoot
\end{longtable}

\subsubsection{Description}\label{description-232}

The FAULTS keyword defines the faults in the grid geometry and the
keyword is normally exported with the grid geometry COORD and ZCORN data
sets from static earth modeling software. Note that the FAULTS keyword
is not required to describe the structural geometry as this is already
accounted for in the COORD and ZCORN data sets, but instead lists the
fault traces with respect to the grid. Once the fault traces have been
defined with the FAULTS keyword then the fault transmissibilities can be
modified by the MULTFLT keyword. Note that without the FAULTS keyword
one would still get proper cross-fault transmissibilities but they would
not be modifiable using MULTFLT keyword.

\begin{longtable}[]{@{}
  >{\raggedright\arraybackslash}p{(\linewidth - 6\tabcolsep) * \real{0.2500}}
  >{\raggedright\arraybackslash}p{(\linewidth - 6\tabcolsep) * \real{0.2500}}
  >{\raggedright\arraybackslash}p{(\linewidth - 6\tabcolsep) * \real{0.2500}}
  >{\raggedright\arraybackslash}p{(\linewidth - 6\tabcolsep) * \real{0.2500}}@{}}
\toprule\noalign{}
\begin{minipage}[b]{\linewidth}\raggedright
No.
\end{minipage} & \begin{minipage}[b]{\linewidth}\raggedright
Name
\end{minipage} & \begin{minipage}[b]{\linewidth}\raggedright
Description
\end{minipage} & \begin{minipage}[b]{\linewidth}\raggedright
Default
\end{minipage} \\
\midrule\noalign{}
\endhead
\bottomrule\noalign{}
\endlastfoot
1 & FLTNAME & FLTNAME is a character string enclosed in quotes with a
maximum length of eight characters, that defines the name of the fault.
& None \\
2 & I1 & The lower bound of the fault's I-direction range must be
greater than or equal to one and less than or equal to I2 and NX. &
None \\
3 & I2 & The upper bound of the fault's I-direction range must be
greater than or equal to II and less than or equal to NX & None \\
4 & J1 & The lower bound of the fault's J-direction range must be
greater than or equal to one and less than or equal to J2 and NY. &
None \\
5 & J2 & The upper bound of the fault's J-direction range must be
greater than or equal to JI and less than or equal to NY. & None \\
6 & K1 & The lower bound of the fault's K-direction range must be
greater than or equal to one and less than or equal to K2 and NZ. &
None \\
7 & K2 & The upper bound of the fault's K-direction range must be
greater than or equal to KI and less than or equal to NZ. & None \\
8 & FLTFACE & Unknown Author 2017-06-27T13:59:19 I am not able to review
this properly, as I am not familiar enough with what is supposed to
happen. I do not know if we support X-, Y-, Z- transmissibilities.
FLTFACE is a character string enclosed in quotes with a maximum length
of two characters, that classifies the fault face. & None \\
Notes: & & & \\
\end{longtable}

\emph{Table 6.37: FAULTS Keyword Description}

\subsubsection{Example}\label{example-160}

The example below defines two fault traces, the first being the
`M\_WEST' fault and the second the `BC' fault trace.

\begin{verbatim}
--
--       DEFINE FAULTS IN THE GRID GEOMETRY
--
--       FAULT       ------------ FAULT TRACE -------------
--       NAME         I1   I2    J1   J2    K1   K2    FACE
FAULTS
         'M_WEST'      5    5     3    3    1    22    'X' /
         'M_WEST'      5    5     4    4    1    22    'X' /
         'M_WEST'      5    5     5    5    1    22    'X' /
         'M_WEST'      5    5     6    6    1    22    'X' /
         'M_WEST'      5    5     7    7    1    22    'X' /
         'M_WEST'      5    5     8    8    1    22    'X' /
         'M_WEST'      5    5     9    9    1    22    'X' /
         'M_WEST'      5    5    10   10    1    22    'X' /
         'M_WEST'      5    5    11   11    1    22    'X' /
……………………………………………..

         'BC'         43   43     8    8    1    22   'Y'  /
         'BC'         42   42     9    9    1    22   'X'  /
         'BC'         44   44     8    8    1    22   'Y'  /
         'BC'         45   45     8    8    1    22   'Y'  /
         'BC'         46   46     8    8    1    22   'Y'  /
         'BC'         31   31     9    9    1    22   'Y'  /
         'BC'         30   30    10   10    1    22   'X'  /
         'BC'         32   32     9    9    1    22   'Y'  /
         'BC'         33   33     9    9    1    22   'Y'  /
         'BC'         34   34     9    9    1    22   'Y'  /
         'BC'         35   35     9    9    1    22   'Y'  /
         'BC'         36   36     9    9    1    22   'Y'  /
         'BC'         37   37     9    9    1    22   'Y'  /
         'BC'         38   38     9    9    1    22   'Y'  /
         'BC'         39   39     9    9    1    22   'Y'  /
         'BC'         40   40     9    9    1    22   'Y'  /
……………………………………………..
/
\end{verbatim}

\subsection{FILEUNIT -- Activate Unit Consistency
Checking}\label{fileunit-activate-unit-consistency-checking}

\begin{longtable}[]{@{}
  >{\raggedright\arraybackslash}p{(\linewidth - 14\tabcolsep) * \real{0.1250}}
  >{\raggedright\arraybackslash}p{(\linewidth - 14\tabcolsep) * \real{0.1250}}
  >{\raggedright\arraybackslash}p{(\linewidth - 14\tabcolsep) * \real{0.1250}}
  >{\raggedright\arraybackslash}p{(\linewidth - 14\tabcolsep) * \real{0.1250}}
  >{\raggedright\arraybackslash}p{(\linewidth - 14\tabcolsep) * \real{0.1250}}
  >{\raggedright\arraybackslash}p{(\linewidth - 14\tabcolsep) * \real{0.1250}}
  >{\raggedright\arraybackslash}p{(\linewidth - 14\tabcolsep) * \real{0.1250}}
  >{\raggedright\arraybackslash}p{(\linewidth - 14\tabcolsep) * \real{0.1250}}@{}}
\toprule\noalign{}
\begin{minipage}[b]{\linewidth}\raggedright
RUNSPEC
\end{minipage} & \begin{minipage}[b]{\linewidth}\raggedright
GRID
\end{minipage} & \begin{minipage}[b]{\linewidth}\raggedright
EDIT
\end{minipage} & \begin{minipage}[b]{\linewidth}\raggedright
PROPS
\end{minipage} & \begin{minipage}[b]{\linewidth}\raggedright
REGIONS
\end{minipage} & \begin{minipage}[b]{\linewidth}\raggedright
SOLUTION
\end{minipage} & \begin{minipage}[b]{\linewidth}\raggedright
SUMMARY
\end{minipage} & \begin{minipage}[b]{\linewidth}\raggedright
SCHEDULE
\end{minipage} \\
\midrule\noalign{}
\endhead
\bottomrule\noalign{}
\endlastfoot
\end{longtable}

\subsubsection{Description}\label{description-233}

The FILEUNIT keyword defines the units of the data set, and is used to
verify that the units in the input deck and any associated include files
are consistent. The keyword does not provide for the conversion between
different sets of units.

\begin{longtable}[]{@{}
  >{\raggedright\arraybackslash}p{(\linewidth - 6\tabcolsep) * \real{0.2500}}
  >{\raggedright\arraybackslash}p{(\linewidth - 6\tabcolsep) * \real{0.2500}}
  >{\raggedright\arraybackslash}p{(\linewidth - 6\tabcolsep) * \real{0.2500}}
  >{\raggedright\arraybackslash}p{(\linewidth - 6\tabcolsep) * \real{0.2500}}@{}}
\toprule\noalign{}
\begin{minipage}[b]{\linewidth}\raggedright
No.
\end{minipage} & \begin{minipage}[b]{\linewidth}\raggedright
Name
\end{minipage} & \begin{minipage}[b]{\linewidth}\raggedright
Description
\end{minipage} & \begin{minipage}[b]{\linewidth}\raggedright
Default
\end{minipage} \\
\midrule\noalign{}
\endhead
\bottomrule\noalign{}
\endlastfoot
1 & FILEUNIT & A character string that defines the units of the data
set, and should be set to: & None \\
Notes: & & & \\
\end{longtable}

\emph{Table 6.38: FILEUNIT Keyword Description}

OPM Flow's behavior is controllable through the
``UNIT\_SYSTEM\_MISMATCH'' environment variable. The default behavior if
the check fails (i.e., if one of the INCLUDE files has a unit system
different from the main run specification) is to terminate the
simulation with an error.

\subsubsection{Example}\label{example-161}

\begin{verbatim}
--
--       ACTIVATE UNIT CONSISTENCY CHECKING
--
FILEUNIT
         FIELD                                                                 /
\end{verbatim}

The above example defines the data set units to be FIELD units.

\subsection{FLUXNUM -- Define the Flux
Regions}\label{fluxnum-define-the-flux-regions}

\begin{longtable}[]{@{}
  >{\raggedright\arraybackslash}p{(\linewidth - 14\tabcolsep) * \real{0.1250}}
  >{\raggedright\arraybackslash}p{(\linewidth - 14\tabcolsep) * \real{0.1250}}
  >{\raggedright\arraybackslash}p{(\linewidth - 14\tabcolsep) * \real{0.1250}}
  >{\raggedright\arraybackslash}p{(\linewidth - 14\tabcolsep) * \real{0.1250}}
  >{\raggedright\arraybackslash}p{(\linewidth - 14\tabcolsep) * \real{0.1250}}
  >{\raggedright\arraybackslash}p{(\linewidth - 14\tabcolsep) * \real{0.1250}}
  >{\raggedright\arraybackslash}p{(\linewidth - 14\tabcolsep) * \real{0.1250}}
  >{\raggedright\arraybackslash}p{(\linewidth - 14\tabcolsep) * \real{0.1250}}@{}}
\toprule\noalign{}
\begin{minipage}[b]{\linewidth}\raggedright
RUNSPEC
\end{minipage} & \begin{minipage}[b]{\linewidth}\raggedright
GRID
\end{minipage} & \begin{minipage}[b]{\linewidth}\raggedright
EDIT
\end{minipage} & \begin{minipage}[b]{\linewidth}\raggedright
PROPS
\end{minipage} & \begin{minipage}[b]{\linewidth}\raggedright
REGIONS
\end{minipage} & \begin{minipage}[b]{\linewidth}\raggedright
SOLUTION
\end{minipage} & \begin{minipage}[b]{\linewidth}\raggedright
SUMMARY
\end{minipage} & \begin{minipage}[b]{\linewidth}\raggedright
SCHEDULE
\end{minipage} \\
\midrule\noalign{}
\endhead
\bottomrule\noalign{}
\endlastfoot
\end{longtable}

\subsubsection{Description}\label{description-234}

The FLUXNUM keyword defines the flux region numbers for each grid block,
as such there must be one entry for each cell in the model. The array is
used with the Flux Boundary option to define the various flux regions;
however, the Flux Boundary option has not been implemented in OPM Flow.
In addition, the array can be used with the EQUALREG, ADDREG, COPYREG,
MULTIREG, MULTREGP and MULTREGT keywords in calculating various grid
properties in the GRID section. This facility has been implemented in
OPM Flow.

\begin{longtable}[]{@{}
  >{\raggedright\arraybackslash}p{(\linewidth - 6\tabcolsep) * \real{0.2500}}
  >{\raggedright\arraybackslash}p{(\linewidth - 6\tabcolsep) * \real{0.2500}}
  >{\raggedright\arraybackslash}p{(\linewidth - 6\tabcolsep) * \real{0.2500}}
  >{\raggedright\arraybackslash}p{(\linewidth - 6\tabcolsep) * \real{0.2500}}@{}}
\toprule\noalign{}
\begin{minipage}[b]{\linewidth}\raggedright
No.
\end{minipage} & \begin{minipage}[b]{\linewidth}\raggedright
Name
\end{minipage} & \begin{minipage}[b]{\linewidth}\raggedright
Description
\end{minipage} & \begin{minipage}[b]{\linewidth}\raggedright
Default
\end{minipage} \\
\midrule\noalign{}
\endhead
\bottomrule\noalign{}
\endlastfoot
1 & FLUXNUM & FLUXNUM defines an array of positive integers assigning a
grid cell to a particular flux region. The maximum number of flux
regions is set by the MXNFLX variable on the REGDIMS keyword in the
RUNSPEC section. & 1 \\
Notes: & & & \\
\end{longtable}

\emph{Table 6.39: FLUXNUM Keyword Description}

\subsubsection{Examples}\label{examples-10}

The example below sets three FLUXNUM regions for a 4 x 5 x 2 model.

\begin{verbatim}
--
--       DEFINE FLUXNUM REGIONS FOR ALL CELLS
--
FLUXNUM
         2  2  1  1  2  2  1  1  1  1  1  1  1  1  1  1  1  1  1  1
         3  3  1  1  3  3  1  1  1  1  1  1  1  1  1  1  1  1  1  1
/
\end{verbatim}

Alternatively the EQUALS keyword could be employed to accomplish the
same task, that is:

\begin{verbatim}
--
--       ARRAY     CONSTANT       ---------- BOX ---------
--                                I1  I2   J1  J2   K1  K2
EQUALS
         FLUXNUM    1             1*  1*   1*  1*   1*  1* / SET REGION 1
         FLUXNUM    2             1   2    1   2    1   1  / SET REGION 2
         FLUXNUM    3             1   2    1   2    2   2  / SET REGION 3
/

\end{verbatim}

\subsection{FLUXREG -- Define Active Flux
Regions}\label{fluxreg-define-active-flux-regions}

\begin{longtable}[]{@{}
  >{\raggedright\arraybackslash}p{(\linewidth - 14\tabcolsep) * \real{0.1250}}
  >{\raggedright\arraybackslash}p{(\linewidth - 14\tabcolsep) * \real{0.1250}}
  >{\raggedright\arraybackslash}p{(\linewidth - 14\tabcolsep) * \real{0.1250}}
  >{\raggedright\arraybackslash}p{(\linewidth - 14\tabcolsep) * \real{0.1250}}
  >{\raggedright\arraybackslash}p{(\linewidth - 14\tabcolsep) * \real{0.1250}}
  >{\raggedright\arraybackslash}p{(\linewidth - 14\tabcolsep) * \real{0.1250}}
  >{\raggedright\arraybackslash}p{(\linewidth - 14\tabcolsep) * \real{0.1250}}
  >{\raggedright\arraybackslash}p{(\linewidth - 14\tabcolsep) * \real{0.1250}}@{}}
\toprule\noalign{}
\begin{minipage}[b]{\linewidth}\raggedright
RUNSPEC
\end{minipage} & \begin{minipage}[b]{\linewidth}\raggedright
GRID
\end{minipage} & \begin{minipage}[b]{\linewidth}\raggedright
EDIT
\end{minipage} & \begin{minipage}[b]{\linewidth}\raggedright
PROPS
\end{minipage} & \begin{minipage}[b]{\linewidth}\raggedright
REGIONS
\end{minipage} & \begin{minipage}[b]{\linewidth}\raggedright
SOLUTION
\end{minipage} & \begin{minipage}[b]{\linewidth}\raggedright
SUMMARY
\end{minipage} & \begin{minipage}[b]{\linewidth}\raggedright
SCHEDULE
\end{minipage} \\
\midrule\noalign{}
\endhead
\bottomrule\noalign{}
\endlastfoot
\end{longtable}

\subsubsection{Description}\label{description-235}

The FLUXREG is used in conjunction with the USEFLUX keyword in runs with
have multiple flux regions, to reduce the number of flux regions, that
is the keyword specifies which flux regions are active and which are not
in the current run.

This keyword is not supported by OPM Flow but would change the results
if supported so the simulation will be stopped.

\subsection{FLUXTYPE -- Defines the Flux Boundary
Type}\label{fluxtype-defines-the-flux-boundary-type}

\begin{longtable}[]{@{}
  >{\raggedright\arraybackslash}p{(\linewidth - 14\tabcolsep) * \real{0.1250}}
  >{\raggedright\arraybackslash}p{(\linewidth - 14\tabcolsep) * \real{0.1250}}
  >{\raggedright\arraybackslash}p{(\linewidth - 14\tabcolsep) * \real{0.1250}}
  >{\raggedright\arraybackslash}p{(\linewidth - 14\tabcolsep) * \real{0.1250}}
  >{\raggedright\arraybackslash}p{(\linewidth - 14\tabcolsep) * \real{0.1250}}
  >{\raggedright\arraybackslash}p{(\linewidth - 14\tabcolsep) * \real{0.1250}}
  >{\raggedright\arraybackslash}p{(\linewidth - 14\tabcolsep) * \real{0.1250}}
  >{\raggedright\arraybackslash}p{(\linewidth - 14\tabcolsep) * \real{0.1250}}@{}}
\toprule\noalign{}
\begin{minipage}[b]{\linewidth}\raggedright
RUNSPEC
\end{minipage} & \begin{minipage}[b]{\linewidth}\raggedright
GRID
\end{minipage} & \begin{minipage}[b]{\linewidth}\raggedright
EDIT
\end{minipage} & \begin{minipage}[b]{\linewidth}\raggedright
PROPS
\end{minipage} & \begin{minipage}[b]{\linewidth}\raggedright
REGIONS
\end{minipage} & \begin{minipage}[b]{\linewidth}\raggedright
SOLUTION
\end{minipage} & \begin{minipage}[b]{\linewidth}\raggedright
SUMMARY
\end{minipage} & \begin{minipage}[b]{\linewidth}\raggedright
SCHEDULE
\end{minipage} \\
\midrule\noalign{}
\endhead
\bottomrule\noalign{}
\endlastfoot
\end{longtable}

\subsubsection{Description}\label{description-236}

The FLUXTYPE keyword defines the type of flux boundary to be used in the
run.

This keyword is not supported by OPM Flow but would change the results
if supported so the simulation will be stopped.

\subsection{FORMFEED -- Defined the Print File Form-Feed
Character}\label{formfeed-defined-the-print-file-form-feed-character-2}

The FORMFEED keyword defines the form-feed character, or carriage
control character, for the output print (\emph{.PRT) run summary (}.RSM)
files. The keyword should be place at the very top of the input file.

See FORMFEED -- Defined the Print File Form-Feed Character in the GLOBAL
section for a dull description.

\subsection{GDFILE -- Load a Grid File}\label{gdfile-load-a-grid-file}

\begin{longtable}[]{@{}
  >{\raggedright\arraybackslash}p{(\linewidth - 14\tabcolsep) * \real{0.1250}}
  >{\raggedright\arraybackslash}p{(\linewidth - 14\tabcolsep) * \real{0.1250}}
  >{\raggedright\arraybackslash}p{(\linewidth - 14\tabcolsep) * \real{0.1250}}
  >{\raggedright\arraybackslash}p{(\linewidth - 14\tabcolsep) * \real{0.1250}}
  >{\raggedright\arraybackslash}p{(\linewidth - 14\tabcolsep) * \real{0.1250}}
  >{\raggedright\arraybackslash}p{(\linewidth - 14\tabcolsep) * \real{0.1250}}
  >{\raggedright\arraybackslash}p{(\linewidth - 14\tabcolsep) * \real{0.1250}}
  >{\raggedright\arraybackslash}p{(\linewidth - 14\tabcolsep) * \real{0.1250}}@{}}
\toprule\noalign{}
\begin{minipage}[b]{\linewidth}\raggedright
RUNSPEC
\end{minipage} & \begin{minipage}[b]{\linewidth}\raggedright
GRID
\end{minipage} & \begin{minipage}[b]{\linewidth}\raggedright
EDIT
\end{minipage} & \begin{minipage}[b]{\linewidth}\raggedright
PROPS
\end{minipage} & \begin{minipage}[b]{\linewidth}\raggedright
REGIONS
\end{minipage} & \begin{minipage}[b]{\linewidth}\raggedright
SOLUTION
\end{minipage} & \begin{minipage}[b]{\linewidth}\raggedright
SUMMARY
\end{minipage} & \begin{minipage}[b]{\linewidth}\raggedright
SCHEDULE
\end{minipage} \\
\midrule\noalign{}
\endhead
\bottomrule\noalign{}
\endlastfoot
\end{longtable}

\subsubsection{Description}\label{description-237}

The GDFILE keyword loads a GRID file that contains the structural data
for the grid as a set of topological cuboidal cells, and EGRID files
that contain structural and property data. Note OPM Flow only supports
reading in EGRID files at this time.

\begin{longtable}[]{@{}
  >{\raggedright\arraybackslash}p{(\linewidth - 6\tabcolsep) * \real{0.2500}}
  >{\raggedright\arraybackslash}p{(\linewidth - 6\tabcolsep) * \real{0.2500}}
  >{\raggedright\arraybackslash}p{(\linewidth - 6\tabcolsep) * \real{0.2500}}
  >{\raggedright\arraybackslash}p{(\linewidth - 6\tabcolsep) * \real{0.2500}}@{}}
\toprule\noalign{}
\begin{minipage}[b]{\linewidth}\raggedright
No.
\end{minipage} & \begin{minipage}[b]{\linewidth}\raggedright
Name
\end{minipage} & \begin{minipage}[b]{\linewidth}\raggedright
Description
\end{minipage} & \begin{minipage}[b]{\linewidth}\raggedright
Default
\end{minipage} \\
\midrule\noalign{}
\endhead
\bottomrule\noalign{}
\endlastfoot
1 & GRIDFILE & A character string enclosed in quotes that defines the
GRID or EGRID file to be read in and be processed by OPM Flow. Again,
OPM Flow only supports reading in EGRID files. & None \\
2 & FMTOPT & A defined character string that defines the format of the
GRID or EGRID file to be read and should be set to one of the following:
If the variable FMTOPT is omitted then the default is for binary file
input for the commercial simulator; whereas, OPM Flow derives FMTOPT
from the file extension (\emph{.EGRID or }.FEGRID), making FMTOPT
superfluous. & U \\
Notes: & & & \\
\end{longtable}

\emph{Table 6.40: GDFILE Keyword Description}

See also the GRIDFILE keyword in the GRID section for exporting the GRID
and EGRID files from the current simulation run.

\subsubsection{Examples}\label{examples-11}

The first example shown below loads the NOR-OPM-A00-GRID.EGRID file in
binary format from the same directory as the data file.

\begin{verbatim}
--
--       LOAD A GRID FILE
--
GDFILE
         'NOR-OPM-A00-GRID.EGRID'      /
\end{verbatim}

The next example loads the same EGRID file one directory above from
where the data file is located.

\begin{verbatim}

--
--       LOAD a GRID FILE
--
GDFILE
         '../NOR-OPM-A00-GRID.EGRID'  /

\end{verbatim}

\subsection{GDORIENT -- Define Grid Orientation
Parameters}\label{gdorient-define-grid-orientation-parameters}

\begin{longtable}[]{@{}
  >{\raggedright\arraybackslash}p{(\linewidth - 14\tabcolsep) * \real{0.1250}}
  >{\raggedright\arraybackslash}p{(\linewidth - 14\tabcolsep) * \real{0.1250}}
  >{\raggedright\arraybackslash}p{(\linewidth - 14\tabcolsep) * \real{0.1250}}
  >{\raggedright\arraybackslash}p{(\linewidth - 14\tabcolsep) * \real{0.1250}}
  >{\raggedright\arraybackslash}p{(\linewidth - 14\tabcolsep) * \real{0.1250}}
  >{\raggedright\arraybackslash}p{(\linewidth - 14\tabcolsep) * \real{0.1250}}
  >{\raggedright\arraybackslash}p{(\linewidth - 14\tabcolsep) * \real{0.1250}}
  >{\raggedright\arraybackslash}p{(\linewidth - 14\tabcolsep) * \real{0.1250}}@{}}
\toprule\noalign{}
\begin{minipage}[b]{\linewidth}\raggedright
RUNSPEC
\end{minipage} & \begin{minipage}[b]{\linewidth}\raggedright
GRID
\end{minipage} & \begin{minipage}[b]{\linewidth}\raggedright
EDIT
\end{minipage} & \begin{minipage}[b]{\linewidth}\raggedright
PROPS
\end{minipage} & \begin{minipage}[b]{\linewidth}\raggedright
REGIONS
\end{minipage} & \begin{minipage}[b]{\linewidth}\raggedright
SOLUTION
\end{minipage} & \begin{minipage}[b]{\linewidth}\raggedright
SUMMARY
\end{minipage} & \begin{minipage}[b]{\linewidth}\raggedright
SCHEDULE
\end{minipage} \\
\midrule\noalign{}
\endhead
\bottomrule\noalign{}
\endlastfoot
\end{longtable}

\subsubsection{Description}\label{description-238}

This keyword defines the grid orientation parameters for post-processing
applications.

This keyword is not supported by OPM Flow but has no effect on the
results so it will be ignored.

\subsection{GETDATA -- Load and Assign Data Array from INIT or RESTART
Files}\label{getdata-load-and-assign-data-array-from-init-or-restart-files}

\begin{longtable}[]{@{}
  >{\raggedright\arraybackslash}p{(\linewidth - 14\tabcolsep) * \real{0.1250}}
  >{\raggedright\arraybackslash}p{(\linewidth - 14\tabcolsep) * \real{0.1250}}
  >{\raggedright\arraybackslash}p{(\linewidth - 14\tabcolsep) * \real{0.1250}}
  >{\raggedright\arraybackslash}p{(\linewidth - 14\tabcolsep) * \real{0.1250}}
  >{\raggedright\arraybackslash}p{(\linewidth - 14\tabcolsep) * \real{0.1250}}
  >{\raggedright\arraybackslash}p{(\linewidth - 14\tabcolsep) * \real{0.1250}}
  >{\raggedright\arraybackslash}p{(\linewidth - 14\tabcolsep) * \real{0.1250}}
  >{\raggedright\arraybackslash}p{(\linewidth - 14\tabcolsep) * \real{0.1250}}@{}}
\toprule\noalign{}
\begin{minipage}[b]{\linewidth}\raggedright
RUNSPEC
\end{minipage} & \begin{minipage}[b]{\linewidth}\raggedright
GRID
\end{minipage} & \begin{minipage}[b]{\linewidth}\raggedright
EDIT
\end{minipage} & \begin{minipage}[b]{\linewidth}\raggedright
PROPS
\end{minipage} & \begin{minipage}[b]{\linewidth}\raggedright
REGIONS
\end{minipage} & \begin{minipage}[b]{\linewidth}\raggedright
SOLUTION
\end{minipage} & \begin{minipage}[b]{\linewidth}\raggedright
SUMMARY
\end{minipage} & \begin{minipage}[b]{\linewidth}\raggedright
SCHEDULE
\end{minipage} \\
\midrule\noalign{}
\endhead
\bottomrule\noalign{}
\endlastfoot
\end{longtable}

\subsubsection{Description}\label{description-239}

The GETDATA keyword loads a data array from a previously generated INIT
or RESTART file and assigns the loaded array to either same array in the
run or another property array.

This keyword is not supported by OPM Flow but would change the results
if supported so the simulation will be stopped.

\subsection{GRID -- Define the Start of the GRID Section of
Keywords}\label{grid-define-the-start-of-the-grid-section-of-keywords}

\begin{longtable}[]{@{}
  >{\raggedright\arraybackslash}p{(\linewidth - 14\tabcolsep) * \real{0.1250}}
  >{\raggedright\arraybackslash}p{(\linewidth - 14\tabcolsep) * \real{0.1250}}
  >{\raggedright\arraybackslash}p{(\linewidth - 14\tabcolsep) * \real{0.1250}}
  >{\raggedright\arraybackslash}p{(\linewidth - 14\tabcolsep) * \real{0.1250}}
  >{\raggedright\arraybackslash}p{(\linewidth - 14\tabcolsep) * \real{0.1250}}
  >{\raggedright\arraybackslash}p{(\linewidth - 14\tabcolsep) * \real{0.1250}}
  >{\raggedright\arraybackslash}p{(\linewidth - 14\tabcolsep) * \real{0.1250}}
  >{\raggedright\arraybackslash}p{(\linewidth - 14\tabcolsep) * \real{0.1250}}@{}}
\toprule\noalign{}
\begin{minipage}[b]{\linewidth}\raggedright
RUNSPEC
\end{minipage} & \begin{minipage}[b]{\linewidth}\raggedright
GRID
\end{minipage} & \begin{minipage}[b]{\linewidth}\raggedright
EDIT
\end{minipage} & \begin{minipage}[b]{\linewidth}\raggedright
PROPS
\end{minipage} & \begin{minipage}[b]{\linewidth}\raggedright
REGIONS
\end{minipage} & \begin{minipage}[b]{\linewidth}\raggedright
SOLUTION
\end{minipage} & \begin{minipage}[b]{\linewidth}\raggedright
SUMMARY
\end{minipage} & \begin{minipage}[b]{\linewidth}\raggedright
SCHEDULE
\end{minipage} \\
\midrule\noalign{}
\endhead
\bottomrule\noalign{}
\endlastfoot
\end{longtable}

\subsubsection{Description}\label{description-240}

The GRID activation keyword marks the end of the RUNSPEC section and the
start of the GRID section that defines the key grid property data for
the simulator including the grid structure, porosity, permeability and
other relevant grid property data.

There is no data required for this keyword.

\subsubsection{Example}\label{example-162}

\begin{verbatim}
-- ==============================================================================
--
-- GRID SECTION
--
-- ==============================================================================
GRID
\end{verbatim}

The above example marks the end of the RUNSPEC section and the start of
the GRID section in the OPM Flow data input file.

\subsection{GRIDFILE -- Set the Grid File Output
Options}\label{gridfile-set-the-grid-file-output-options}

\begin{longtable}[]{@{}
  >{\raggedright\arraybackslash}p{(\linewidth - 14\tabcolsep) * \real{0.1250}}
  >{\raggedright\arraybackslash}p{(\linewidth - 14\tabcolsep) * \real{0.1250}}
  >{\raggedright\arraybackslash}p{(\linewidth - 14\tabcolsep) * \real{0.1250}}
  >{\raggedright\arraybackslash}p{(\linewidth - 14\tabcolsep) * \real{0.1250}}
  >{\raggedright\arraybackslash}p{(\linewidth - 14\tabcolsep) * \real{0.1250}}
  >{\raggedright\arraybackslash}p{(\linewidth - 14\tabcolsep) * \real{0.1250}}
  >{\raggedright\arraybackslash}p{(\linewidth - 14\tabcolsep) * \real{0.1250}}
  >{\raggedright\arraybackslash}p{(\linewidth - 14\tabcolsep) * \real{0.1250}}@{}}
\toprule\noalign{}
\begin{minipage}[b]{\linewidth}\raggedright
RUNSPEC
\end{minipage} & \begin{minipage}[b]{\linewidth}\raggedright
GRID
\end{minipage} & \begin{minipage}[b]{\linewidth}\raggedright
EDIT
\end{minipage} & \begin{minipage}[b]{\linewidth}\raggedright
PROPS
\end{minipage} & \begin{minipage}[b]{\linewidth}\raggedright
REGIONS
\end{minipage} & \begin{minipage}[b]{\linewidth}\raggedright
SOLUTION
\end{minipage} & \begin{minipage}[b]{\linewidth}\raggedright
SUMMARY
\end{minipage} & \begin{minipage}[b]{\linewidth}\raggedright
SCHEDULE
\end{minipage} \\
\midrule\noalign{}
\endhead
\bottomrule\noalign{}
\endlastfoot
\end{longtable}

\subsubsection{Description}\label{description-241}

This keyword controls the output of a standard GRID or extended GRID
file, as well as the extensible EGRID file for post-processing
applications. The extended and extensible GRID formats are comparable;
however, the extensible GRID format is more compact and is the only
format supported by OPM Flow.

\begin{longtable}[]{@{}
  >{\raggedright\arraybackslash}p{(\linewidth - 6\tabcolsep) * \real{0.2500}}
  >{\raggedright\arraybackslash}p{(\linewidth - 6\tabcolsep) * \real{0.2500}}
  >{\raggedright\arraybackslash}p{(\linewidth - 6\tabcolsep) * \real{0.2500}}
  >{\raggedright\arraybackslash}p{(\linewidth - 6\tabcolsep) * \real{0.2500}}@{}}
\toprule\noalign{}
\begin{minipage}[b]{\linewidth}\raggedright
No.
\end{minipage} & \begin{minipage}[b]{\linewidth}\raggedright
Name
\end{minipage} & \begin{minipage}[b]{\linewidth}\raggedright
Description
\end{minipage} & \begin{minipage}[b]{\linewidth}\raggedright
Default
\end{minipage} \\
\midrule\noalign{}
\endhead
\bottomrule\noalign{}
\endlastfoot
1 & NGRID & A positive integer that controls the output of the GRID
geometry file: Only the default value of zero is supported. & 0 \\
2 & NEGRID & A positive integer that controls the output of the EGRID
geometry file: Only the default value of one is supported. & 1 \\
Notes: & & & \\
\end{longtable}

\emph{Table 6.41: GRIDFILE Keyword Description}

\subsubsection{Example}\label{example-163}

\begin{verbatim}
--
--       GRID FILE OUTPUT OPTIONS
--       GRID    EGRID
--       OPTN    OPTN
GRIDFILE
         0       1                                                             /
\end{verbatim}

The above example defines that no GRID file will be written out and that
the extensible GRID (that is the EGRID geometry format) file will be
produced. This is the only configuration that OPM Flow supports

\subsection{GRIDUNIT -- Define the Grid
Units}\label{gridunit-define-the-grid-units}

\begin{longtable}[]{@{}
  >{\raggedright\arraybackslash}p{(\linewidth - 14\tabcolsep) * \real{0.1250}}
  >{\raggedright\arraybackslash}p{(\linewidth - 14\tabcolsep) * \real{0.1250}}
  >{\raggedright\arraybackslash}p{(\linewidth - 14\tabcolsep) * \real{0.1250}}
  >{\raggedright\arraybackslash}p{(\linewidth - 14\tabcolsep) * \real{0.1250}}
  >{\raggedright\arraybackslash}p{(\linewidth - 14\tabcolsep) * \real{0.1250}}
  >{\raggedright\arraybackslash}p{(\linewidth - 14\tabcolsep) * \real{0.1250}}
  >{\raggedright\arraybackslash}p{(\linewidth - 14\tabcolsep) * \real{0.1250}}
  >{\raggedright\arraybackslash}p{(\linewidth - 14\tabcolsep) * \real{0.1250}}@{}}
\toprule\noalign{}
\begin{minipage}[b]{\linewidth}\raggedright
RUNSPEC
\end{minipage} & \begin{minipage}[b]{\linewidth}\raggedright
GRID
\end{minipage} & \begin{minipage}[b]{\linewidth}\raggedright
EDIT
\end{minipage} & \begin{minipage}[b]{\linewidth}\raggedright
PROPS
\end{minipage} & \begin{minipage}[b]{\linewidth}\raggedright
REGIONS
\end{minipage} & \begin{minipage}[b]{\linewidth}\raggedright
SOLUTION
\end{minipage} & \begin{minipage}[b]{\linewidth}\raggedright
SUMMARY
\end{minipage} & \begin{minipage}[b]{\linewidth}\raggedright
SCHEDULE
\end{minipage} \\
\midrule\noalign{}
\endhead
\bottomrule\noalign{}
\endlastfoot
\end{longtable}

\subsubsection{Description}\label{description-242}

The GRIDUNIT keyword defines the units of the grid data. It is usually
output by pre-processing software when exporting the grid geometry. The
data is not used by OPM Flow intrinsically, but is merely written to the
output EGRID file, as specified by the GRIDFILE keyword, for the use of
post-processing software like OPM ResInsight.

\begin{longtable}[]{@{}
  >{\raggedright\arraybackslash}p{(\linewidth - 6\tabcolsep) * \real{0.2500}}
  >{\raggedright\arraybackslash}p{(\linewidth - 6\tabcolsep) * \real{0.2500}}
  >{\raggedright\arraybackslash}p{(\linewidth - 6\tabcolsep) * \real{0.2500}}
  >{\raggedright\arraybackslash}p{(\linewidth - 6\tabcolsep) * \real{0.2500}}@{}}
\toprule\noalign{}
\begin{minipage}[b]{\linewidth}\raggedright
No.
\end{minipage} & \begin{minipage}[b]{\linewidth}\raggedright
Name
\end{minipage} & \begin{minipage}[b]{\linewidth}\raggedright
Description
\end{minipage} & \begin{minipage}[b]{\linewidth}\raggedright
Default
\end{minipage} \\
\midrule\noalign{}
\endhead
\bottomrule\noalign{}
\endlastfoot
1 & GRIDUNIT & A character string that defines the units of the
coordinates stated on the MAPAXES keyword, and should be set to: &
METRES \\
2 & MAPOPT & A character string that defines if the grid data are
measured relative to the map, or relative to the origin as stated on the
MAPAXES keyword. MAPOPT should either be left blank (the default)
indicating the origin is relative to the origin on the MAPAXES keyword,
or set equal to MAP measured relative to the map. & 1* \\
Notes: & & & \\
\end{longtable}

\emph{Table 6.42: GRIDUNIT Keyword Description}

\subsubsection{Example}\label{example-164}

\begin{verbatim}
--
--       SET THE GRID UNITS FOR THE GRID
--
GRIDUNIT
         METRES                                                                /
\end{verbatim}

The above example defines that the GRID units to be metric.

\subsection{HALFTRAN -- Activate Half Block Transmissibility
Calculations}\label{halftran-activate-half-block-transmissibility-calculations}

\begin{longtable}[]{@{}
  >{\raggedright\arraybackslash}p{(\linewidth - 14\tabcolsep) * \real{0.1250}}
  >{\raggedright\arraybackslash}p{(\linewidth - 14\tabcolsep) * \real{0.1250}}
  >{\raggedright\arraybackslash}p{(\linewidth - 14\tabcolsep) * \real{0.1250}}
  >{\raggedright\arraybackslash}p{(\linewidth - 14\tabcolsep) * \real{0.1250}}
  >{\raggedright\arraybackslash}p{(\linewidth - 14\tabcolsep) * \real{0.1250}}
  >{\raggedright\arraybackslash}p{(\linewidth - 14\tabcolsep) * \real{0.1250}}
  >{\raggedright\arraybackslash}p{(\linewidth - 14\tabcolsep) * \real{0.1250}}
  >{\raggedright\arraybackslash}p{(\linewidth - 14\tabcolsep) * \real{0.1250}}@{}}
\toprule\noalign{}
\begin{minipage}[b]{\linewidth}\raggedright
RUNSPEC
\end{minipage} & \begin{minipage}[b]{\linewidth}\raggedright
GRID
\end{minipage} & \begin{minipage}[b]{\linewidth}\raggedright
EDIT
\end{minipage} & \begin{minipage}[b]{\linewidth}\raggedright
PROPS
\end{minipage} & \begin{minipage}[b]{\linewidth}\raggedright
REGIONS
\end{minipage} & \begin{minipage}[b]{\linewidth}\raggedright
SOLUTION
\end{minipage} & \begin{minipage}[b]{\linewidth}\raggedright
SUMMARY
\end{minipage} & \begin{minipage}[b]{\linewidth}\raggedright
SCHEDULE
\end{minipage} \\
\midrule\noalign{}
\endhead
\bottomrule\noalign{}
\endlastfoot
\end{longtable}

\subsubsection{Description}\label{description-243}

The HALFTRAN keyword activates the half block transmissibility
calculation option.

This keyword is not supported by OPM Flow but would change the results
if supported so the simulation will be stopped.

\subsection{HEATCR -- Define Reservoir Rock Heat Capacity for All
Cells}\label{heatcr-define-reservoir-rock-heat-capacity-for-all-cells}

\begin{longtable}[]{@{}
  >{\raggedright\arraybackslash}p{(\linewidth - 14\tabcolsep) * \real{0.1250}}
  >{\raggedright\arraybackslash}p{(\linewidth - 14\tabcolsep) * \real{0.1250}}
  >{\raggedright\arraybackslash}p{(\linewidth - 14\tabcolsep) * \real{0.1250}}
  >{\raggedright\arraybackslash}p{(\linewidth - 14\tabcolsep) * \real{0.1250}}
  >{\raggedright\arraybackslash}p{(\linewidth - 14\tabcolsep) * \real{0.1250}}
  >{\raggedright\arraybackslash}p{(\linewidth - 14\tabcolsep) * \real{0.1250}}
  >{\raggedright\arraybackslash}p{(\linewidth - 14\tabcolsep) * \real{0.1250}}
  >{\raggedright\arraybackslash}p{(\linewidth - 14\tabcolsep) * \real{0.1250}}@{}}
\toprule\noalign{}
\begin{minipage}[b]{\linewidth}\raggedright
RUNSPEC
\end{minipage} & \begin{minipage}[b]{\linewidth}\raggedright
GRID
\end{minipage} & \begin{minipage}[b]{\linewidth}\raggedright
EDIT
\end{minipage} & \begin{minipage}[b]{\linewidth}\raggedright
PROPS
\end{minipage} & \begin{minipage}[b]{\linewidth}\raggedright
REGIONS
\end{minipage} & \begin{minipage}[b]{\linewidth}\raggedright
SOLUTION
\end{minipage} & \begin{minipage}[b]{\linewidth}\raggedright
SUMMARY
\end{minipage} & \begin{minipage}[b]{\linewidth}\raggedright
SCHEDULE
\end{minipage} \\
\midrule\noalign{}
\endhead
\bottomrule\noalign{}
\endlastfoot
\end{longtable}

\subsubsection{Description}\label{description-244}

The HEATCR keyword defines the reservoir rock volumetric heat capacity
for all cells for when OPM Flow's thermal calculation is activated by
the THERMAL keywords in the RUNSPEC section.

This keyword can only be used if OPM Flow's thermal option has been
activated by the THERMAL keyword in the RUNSPEC section. Note this is
different to the commercial simulator that uses the TEMP keyword in the
RUNSPEC section to activate the black-oil thermal model.

\begin{longtable}[]{@{}
  >{\raggedright\arraybackslash}p{(\linewidth - 6\tabcolsep) * \real{0.2500}}
  >{\raggedright\arraybackslash}p{(\linewidth - 6\tabcolsep) * \real{0.2500}}
  >{\raggedright\arraybackslash}p{(\linewidth - 6\tabcolsep) * \real{0.2500}}
  >{\raggedright\arraybackslash}p{(\linewidth - 6\tabcolsep) * \real{0.2500}}@{}}
\toprule\noalign{}
\begin{minipage}[b]{\linewidth}\raggedright
No.
\end{minipage} & \begin{minipage}[b]{\linewidth}\raggedright
Name
\end{minipage} & \begin{minipage}[b]{\linewidth}\raggedright
Description
\end{minipage} & \begin{minipage}[b]{\linewidth}\raggedright
Default
\end{minipage} \\
\midrule\noalign{}
\endhead
\bottomrule\noalign{}
\endlastfoot
Field & Metric & Laboratory & \\
1 & HEATCR & HEATCR is an array of real positive numbers that define
reservoir rock volumetric heat capacity of a grid block. Repeat counts
may be used, for example 3000*25.0 & None \\
Btu/ft3/°R & kJ/m3/K & J/cm3/K & \\
Notes: & & & \\
\end{longtable}

\emph{Table 6.43: HEATCR Keyword Description}

Note this keyword is incompatible with SPECROCK keyword in the PROPS
section.

\subsubsection{Example}\label{example-165}

\begin{verbatim}
--
--       DEFINE GRID BLOCK RESERVOIR ROCK HEAT CAPACITY
--       FOR ALL CELLS (BASED ON NX x NY x NZ = 300)
--       KEYWORD IS INCOMPATIBLE WITH THE SPECROCK KEYWORD
--       (OPM FLOW THERMAL OPTION ONLY)
--
HEATCR
         300*32.0                                                              /
\end{verbatim}

The above example defines the reservoir rock volumetric heat capacity of
32.0 for each cell in the 300 grid block model.

\subsection{HEATCRT -- Define Reservoir Rock Heat Capacity Temperature
Dependence for All
Cells}\label{heatcrt-define-reservoir-rock-heat-capacity-temperature-dependence-for-all-cells}

\begin{longtable}[]{@{}
  >{\raggedright\arraybackslash}p{(\linewidth - 14\tabcolsep) * \real{0.1250}}
  >{\raggedright\arraybackslash}p{(\linewidth - 14\tabcolsep) * \real{0.1250}}
  >{\raggedright\arraybackslash}p{(\linewidth - 14\tabcolsep) * \real{0.1250}}
  >{\raggedright\arraybackslash}p{(\linewidth - 14\tabcolsep) * \real{0.1250}}
  >{\raggedright\arraybackslash}p{(\linewidth - 14\tabcolsep) * \real{0.1250}}
  >{\raggedright\arraybackslash}p{(\linewidth - 14\tabcolsep) * \real{0.1250}}
  >{\raggedright\arraybackslash}p{(\linewidth - 14\tabcolsep) * \real{0.1250}}
  >{\raggedright\arraybackslash}p{(\linewidth - 14\tabcolsep) * \real{0.1250}}@{}}
\toprule\noalign{}
\begin{minipage}[b]{\linewidth}\raggedright
RUNSPEC
\end{minipage} & \begin{minipage}[b]{\linewidth}\raggedright
GRID
\end{minipage} & \begin{minipage}[b]{\linewidth}\raggedright
EDIT
\end{minipage} & \begin{minipage}[b]{\linewidth}\raggedright
PROPS
\end{minipage} & \begin{minipage}[b]{\linewidth}\raggedright
REGIONS
\end{minipage} & \begin{minipage}[b]{\linewidth}\raggedright
SOLUTION
\end{minipage} & \begin{minipage}[b]{\linewidth}\raggedright
SUMMARY
\end{minipage} & \begin{minipage}[b]{\linewidth}\raggedright
SCHEDULE
\end{minipage} \\
\midrule\noalign{}
\endhead
\bottomrule\noalign{}
\endlastfoot
\end{longtable}

\subsubsection{Description}\label{description-245}

The HEATCRT keyword defines the reservoir rock volumetric heat capacity
temperature dependence for all cells for when OPM Flow's thermal
calculation is activated by the THERMAL keywords in the RUNSPEC section.

This keyword can only be used if OPM Flow's thermal option has been
activated by the THERMAL keyword in the RUNSPEC section. Note this is
different to the commercial simulator that uses the TEMP keyword in the
RUNSPEC section to activate the black-oil thermal model.

\begin{longtable}[]{@{}
  >{\raggedright\arraybackslash}p{(\linewidth - 6\tabcolsep) * \real{0.2500}}
  >{\raggedright\arraybackslash}p{(\linewidth - 6\tabcolsep) * \real{0.2500}}
  >{\raggedright\arraybackslash}p{(\linewidth - 6\tabcolsep) * \real{0.2500}}
  >{\raggedright\arraybackslash}p{(\linewidth - 6\tabcolsep) * \real{0.2500}}@{}}
\toprule\noalign{}
\begin{minipage}[b]{\linewidth}\raggedright
No.
\end{minipage} & \begin{minipage}[b]{\linewidth}\raggedright
Name
\end{minipage} & \begin{minipage}[b]{\linewidth}\raggedright
Description
\end{minipage} & \begin{minipage}[b]{\linewidth}\raggedright
Default
\end{minipage} \\
\midrule\noalign{}
\endhead
\bottomrule\noalign{}
\endlastfoot
Field & Metric & Laboratory & \\
1 & HEATCRT & HEATCRT is an array of real positive numbers that define
reservoir rock volumetric heat capacity temperature dependence of a grid
block. Repeat counts may be used, for example 3000*0.05 & None \\
Btu/ft3/°R2 & kJ/m3/K2 & J/cm3/K2 & \\
Notes: & & & \\
\end{longtable}

\emph{Table 6.44: HEATCRT Keyword Description}

Note this keyword is incompatible with SPECROCK keyword in the PROPS
section.

The data for this keyword and the HEATCR keyword are used to calculate
the reservoir rock volumetric heat capacity temperature dependence using
the following relationship:

\begin{equation}\phantomsection\label{eq-6-3}{
\text{Heat Capacity of Rock }=\text{HEATCR}\left(\mathit{Temp} - {\mathit{Temp}}_{\mathit{ref}}\right) + \frac{\text{HEATCRT}{\left(\mathit{Temp} - {\mathit{Temp}}_{\mathit{ref}}\right)}^{2}}{2}
}\end{equation}

\subsubsection{Example}\label{example-166}

\begin{verbatim}
--
--       DEFINE RESERVOIR ROCK HEAT CAPACITY TEMPERATURE DEPENDENCE
--       FOR ALL CELLS (BASED ON NX x NY x NZ = 300)
--       KEYWORD IS INCOMPATIBLE WITH THE SPECROCK KEYWORD
--       (OPM FLOW THERMAL OPTION ONLY)
--
HEATCRT
         300*0.05                                                              /
\end{verbatim}

The above example defines the reservoir rock volumetric heat capacity
temperature dependence of 0.05 for each cell in the 300 grid block
model.

\subsection{HMAQUNUM -- History Match Numerical Aquifer Gradient
Parameters}\label{hmaqunum-history-match-numerical-aquifer-gradient-parameters}

\begin{longtable}[]{@{}
  >{\raggedright\arraybackslash}p{(\linewidth - 14\tabcolsep) * \real{0.1250}}
  >{\raggedright\arraybackslash}p{(\linewidth - 14\tabcolsep) * \real{0.1250}}
  >{\raggedright\arraybackslash}p{(\linewidth - 14\tabcolsep) * \real{0.1250}}
  >{\raggedright\arraybackslash}p{(\linewidth - 14\tabcolsep) * \real{0.1250}}
  >{\raggedright\arraybackslash}p{(\linewidth - 14\tabcolsep) * \real{0.1250}}
  >{\raggedright\arraybackslash}p{(\linewidth - 14\tabcolsep) * \real{0.1250}}
  >{\raggedright\arraybackslash}p{(\linewidth - 14\tabcolsep) * \real{0.1250}}
  >{\raggedright\arraybackslash}p{(\linewidth - 14\tabcolsep) * \real{0.1250}}@{}}
\toprule\noalign{}
\begin{minipage}[b]{\linewidth}\raggedright
RUNSPEC
\end{minipage} & \begin{minipage}[b]{\linewidth}\raggedright
GRID
\end{minipage} & \begin{minipage}[b]{\linewidth}\raggedright
EDIT
\end{minipage} & \begin{minipage}[b]{\linewidth}\raggedright
PROPS
\end{minipage} & \begin{minipage}[b]{\linewidth}\raggedright
REGIONS
\end{minipage} & \begin{minipage}[b]{\linewidth}\raggedright
SOLUTION
\end{minipage} & \begin{minipage}[b]{\linewidth}\raggedright
SUMMARY
\end{minipage} & \begin{minipage}[b]{\linewidth}\raggedright
SCHEDULE
\end{minipage} \\
\midrule\noalign{}
\endhead
\bottomrule\noalign{}
\endlastfoot
\end{longtable}

\subsubsection{Description}\label{description-246}

The HMAQUNUM keyword defines the history match numerical aquifer
gradient parameters for when the History Match Gradient option has been
activated by the HMDIMS keyword in the RUNSPEC section, and numerical
aquifers have been specified in the model via the AQUNUM keyword and
connected to the grid using AQUCON keyword. All keywords are in the GRID
section.

See also the HMDIMS keyword in the RUNSPEC section that specifies the
dimensions for the gradient option, including the maximum number of
aquifers that can be used with the History Match Gradient option.

This keyword is not supported by OPM Flow but would change the results
if supported so the simulation will be stopped.

\subsection{HMFAULTS -- History Match Fault Gradient
Parameters}\label{hmfaults-history-match-fault-gradient-parameters}

\begin{longtable}[]{@{}
  >{\raggedright\arraybackslash}p{(\linewidth - 14\tabcolsep) * \real{0.1250}}
  >{\raggedright\arraybackslash}p{(\linewidth - 14\tabcolsep) * \real{0.1250}}
  >{\raggedright\arraybackslash}p{(\linewidth - 14\tabcolsep) * \real{0.1250}}
  >{\raggedright\arraybackslash}p{(\linewidth - 14\tabcolsep) * \real{0.1250}}
  >{\raggedright\arraybackslash}p{(\linewidth - 14\tabcolsep) * \real{0.1250}}
  >{\raggedright\arraybackslash}p{(\linewidth - 14\tabcolsep) * \real{0.1250}}
  >{\raggedright\arraybackslash}p{(\linewidth - 14\tabcolsep) * \real{0.1250}}
  >{\raggedright\arraybackslash}p{(\linewidth - 14\tabcolsep) * \real{0.1250}}@{}}
\toprule\noalign{}
\begin{minipage}[b]{\linewidth}\raggedright
RUNSPEC
\end{minipage} & \begin{minipage}[b]{\linewidth}\raggedright
GRID
\end{minipage} & \begin{minipage}[b]{\linewidth}\raggedright
EDIT
\end{minipage} & \begin{minipage}[b]{\linewidth}\raggedright
PROPS
\end{minipage} & \begin{minipage}[b]{\linewidth}\raggedright
REGIONS
\end{minipage} & \begin{minipage}[b]{\linewidth}\raggedright
SOLUTION
\end{minipage} & \begin{minipage}[b]{\linewidth}\raggedright
SUMMARY
\end{minipage} & \begin{minipage}[b]{\linewidth}\raggedright
SCHEDULE
\end{minipage} \\
\midrule\noalign{}
\endhead
\bottomrule\noalign{}
\endlastfoot
\end{longtable}

\subsubsection{Description}\label{description-247}

The HMFAULTS keyword defines the history match faults gradient
parameters for when the History Match Gradient option has been activated
by the HMDIMS keyword in the RUNSPEC section, and for when the FAULTS
keyword in the GRID section has been used to define faults in the model.

See also the HMDIMS keyword in the RUNSPEC section that specifies the
dimensions for the gradient option, including the maximum number of
faults that can be used with the History Match Gradient option.

This keyword is not supported by OPM Flow but would change the results
if supported so the simulation will be stopped.

\subsection{HMMLAQUN -- History Match Numerical Aquifer Gradient
Multipliers}\label{hmmlaqun-history-match-numerical-aquifer-gradient-multipliers}

\begin{longtable}[]{@{}
  >{\raggedright\arraybackslash}p{(\linewidth - 14\tabcolsep) * \real{0.1250}}
  >{\raggedright\arraybackslash}p{(\linewidth - 14\tabcolsep) * \real{0.1250}}
  >{\raggedright\arraybackslash}p{(\linewidth - 14\tabcolsep) * \real{0.1250}}
  >{\raggedright\arraybackslash}p{(\linewidth - 14\tabcolsep) * \real{0.1250}}
  >{\raggedright\arraybackslash}p{(\linewidth - 14\tabcolsep) * \real{0.1250}}
  >{\raggedright\arraybackslash}p{(\linewidth - 14\tabcolsep) * \real{0.1250}}
  >{\raggedright\arraybackslash}p{(\linewidth - 14\tabcolsep) * \real{0.1250}}
  >{\raggedright\arraybackslash}p{(\linewidth - 14\tabcolsep) * \real{0.1250}}@{}}
\toprule\noalign{}
\begin{minipage}[b]{\linewidth}\raggedright
RUNSPEC
\end{minipage} & \begin{minipage}[b]{\linewidth}\raggedright
GRID
\end{minipage} & \begin{minipage}[b]{\linewidth}\raggedright
EDIT
\end{minipage} & \begin{minipage}[b]{\linewidth}\raggedright
PROPS
\end{minipage} & \begin{minipage}[b]{\linewidth}\raggedright
REGIONS
\end{minipage} & \begin{minipage}[b]{\linewidth}\raggedright
SOLUTION
\end{minipage} & \begin{minipage}[b]{\linewidth}\raggedright
SUMMARY
\end{minipage} & \begin{minipage}[b]{\linewidth}\raggedright
SCHEDULE
\end{minipage} \\
\midrule\noalign{}
\endhead
\bottomrule\noalign{}
\endlastfoot
\end{longtable}

\subsubsection{Description}\label{description-248}

The HMMLAQUN keyword defines the history match numerical aquifer
gradient multipliers for when the History Match Gradient option has been
activated by the HMDIMS keyword in the RUNSPEC section, and numerical
aquifers have been specified in the model via the AQUNUM keyword and
connected to the grid using the AQUCON keyword. All keywords are in the
GRID section.

Multipliers can be declared for numerical aquifers' pore volume,
permeability, and aquifer to grid connection factors.

This keyword is not supported by OPM Flow but would change the results
if supported so the simulation will be stopped.

\subsection{HMMLT -- History Match Grid Permeability Gradient Cumulative
Multipliers}\label{hmmlt-history-match-grid-permeability-gradient-cumulative-multipliers}

\begin{longtable}[]{@{}
  >{\raggedright\arraybackslash}p{(\linewidth - 14\tabcolsep) * \real{0.1250}}
  >{\raggedright\arraybackslash}p{(\linewidth - 14\tabcolsep) * \real{0.1250}}
  >{\raggedright\arraybackslash}p{(\linewidth - 14\tabcolsep) * \real{0.1250}}
  >{\raggedright\arraybackslash}p{(\linewidth - 14\tabcolsep) * \real{0.1250}}
  >{\raggedright\arraybackslash}p{(\linewidth - 14\tabcolsep) * \real{0.1250}}
  >{\raggedright\arraybackslash}p{(\linewidth - 14\tabcolsep) * \real{0.1250}}
  >{\raggedright\arraybackslash}p{(\linewidth - 14\tabcolsep) * \real{0.1250}}
  >{\raggedright\arraybackslash}p{(\linewidth - 14\tabcolsep) * \real{0.1250}}@{}}
\toprule\noalign{}
\begin{minipage}[b]{\linewidth}\raggedright
RUNSPEC
\end{minipage} & \begin{minipage}[b]{\linewidth}\raggedright
GRID
\end{minipage} & \begin{minipage}[b]{\linewidth}\raggedright
EDIT
\end{minipage} & \begin{minipage}[b]{\linewidth}\raggedright
PROPS
\end{minipage} & \begin{minipage}[b]{\linewidth}\raggedright
REGIONS
\end{minipage} & \begin{minipage}[b]{\linewidth}\raggedright
SOLUTION
\end{minipage} & \begin{minipage}[b]{\linewidth}\raggedright
SUMMARY
\end{minipage} & \begin{minipage}[b]{\linewidth}\raggedright
SCHEDULE
\end{minipage} \\
\midrule\noalign{}
\endhead
\bottomrule\noalign{}
\endlastfoot
\end{longtable}

\subsubsection{Description}\label{description-249}

The HMMLT series of keywords defines the history match gradient
cumulative permeability multipliers, for when the History Match Gradient
option has been activated by the HMDIMS keyword in the RUNSPEC section.
The keyword consists of the first five characters of ``HMMLT'' followed
by a two or three character string shown in Table 6.45, that determines
the permeability direction, for example, HMMLTPX.

This keyword is not supported by OPM Flow but has no effect on the
results so it will be ignored.

\begin{longtable}[]{@{}lllll@{}}
\toprule\noalign{}
Mnemonic & Cartesian Grid & Radial Grid & & \\
\midrule\noalign{}
\endhead
\bottomrule\noalign{}
\endlastfoot
Grid Keyword & HMMULT Keyword & Grid Keyword & HMMULT Keyword & \\
PX/PR & PERMX & HMMLTPX & PERMR & HMMLTPR \\
PXY & PERMXY & HMMLTPXY & & \\
PY/THT & PERMY & hMMLTPY & PERMTHT & HMMLTTH \\
PZ & PERMZ & HMMLTPZ & PERMZ & HMMLTPZ \\
\end{longtable}

\emph{Table 6.45: HMMLT Keyword List}

See also the HMMULT keyword in the EDIT section

\subsection{HMMMREGT -- History Match Region Transmissibility Gradient
Cumulative
Multipliers}\label{hmmmregt-history-match-region-transmissibility-gradient-cumulative-multipliers}

\begin{longtable}[]{@{}
  >{\raggedright\arraybackslash}p{(\linewidth - 14\tabcolsep) * \real{0.1250}}
  >{\raggedright\arraybackslash}p{(\linewidth - 14\tabcolsep) * \real{0.1250}}
  >{\raggedright\arraybackslash}p{(\linewidth - 14\tabcolsep) * \real{0.1250}}
  >{\raggedright\arraybackslash}p{(\linewidth - 14\tabcolsep) * \real{0.1250}}
  >{\raggedright\arraybackslash}p{(\linewidth - 14\tabcolsep) * \real{0.1250}}
  >{\raggedright\arraybackslash}p{(\linewidth - 14\tabcolsep) * \real{0.1250}}
  >{\raggedright\arraybackslash}p{(\linewidth - 14\tabcolsep) * \real{0.1250}}
  >{\raggedright\arraybackslash}p{(\linewidth - 14\tabcolsep) * \real{0.1250}}@{}}
\toprule\noalign{}
\begin{minipage}[b]{\linewidth}\raggedright
RUNSPEC
\end{minipage} & \begin{minipage}[b]{\linewidth}\raggedright
GRID
\end{minipage} & \begin{minipage}[b]{\linewidth}\raggedright
EDIT
\end{minipage} & \begin{minipage}[b]{\linewidth}\raggedright
PROPS
\end{minipage} & \begin{minipage}[b]{\linewidth}\raggedright
REGIONS
\end{minipage} & \begin{minipage}[b]{\linewidth}\raggedright
SOLUTION
\end{minipage} & \begin{minipage}[b]{\linewidth}\raggedright
SUMMARY
\end{minipage} & \begin{minipage}[b]{\linewidth}\raggedright
SCHEDULE
\end{minipage} \\
\midrule\noalign{}
\endhead
\bottomrule\noalign{}
\endlastfoot
\end{longtable}

\subsubsection{Description}\label{description-250}

The HMMMREGT keyword multiplies the transmissibility between two regions
by a constant, for when the History Match Gradient option has been
activated by the HMDIMS keyword in the RUNSPEC section. The constant
should be a real number. Unlike the MULTREGT keyword in the GRID
section, the HMMMREGT keyword modifications are cumulative.

Note that the HMMMREGT keyword only declares the two regions and the
multiplier between those regions, the transmissibility direction (DIR on
the MULTREGT keyword), type of transmissibility multiplier (TYPE on the
MULTREGT keyword), and the region number array to use (ARRAY on the
MULTREGT keyword), are all taken from the MULTREGY keyword. For example,
the region number array can be FLUXNUM, MULTNUM or OPERNUM and these
arrays must be defined and be available before the MULTREGT keyword is
read by the simulator, and before the HMMMREGT keyword is used.

This keyword is not supported by OPM Flow but has no effect on the
results so it will be ignored.

\subsection{HMMULRGT -- History Match Region Transmissibility
Parameters}\label{hmmulrgt-history-match-region-transmissibility-parameters}

\begin{longtable}[]{@{}
  >{\raggedright\arraybackslash}p{(\linewidth - 14\tabcolsep) * \real{0.1250}}
  >{\raggedright\arraybackslash}p{(\linewidth - 14\tabcolsep) * \real{0.1250}}
  >{\raggedright\arraybackslash}p{(\linewidth - 14\tabcolsep) * \real{0.1250}}
  >{\raggedright\arraybackslash}p{(\linewidth - 14\tabcolsep) * \real{0.1250}}
  >{\raggedright\arraybackslash}p{(\linewidth - 14\tabcolsep) * \real{0.1250}}
  >{\raggedright\arraybackslash}p{(\linewidth - 14\tabcolsep) * \real{0.1250}}
  >{\raggedright\arraybackslash}p{(\linewidth - 14\tabcolsep) * \real{0.1250}}
  >{\raggedright\arraybackslash}p{(\linewidth - 14\tabcolsep) * \real{0.1250}}@{}}
\toprule\noalign{}
\begin{minipage}[b]{\linewidth}\raggedright
RUNSPEC
\end{minipage} & \begin{minipage}[b]{\linewidth}\raggedright
GRID
\end{minipage} & \begin{minipage}[b]{\linewidth}\raggedright
EDIT
\end{minipage} & \begin{minipage}[b]{\linewidth}\raggedright
PROPS
\end{minipage} & \begin{minipage}[b]{\linewidth}\raggedright
REGIONS
\end{minipage} & \begin{minipage}[b]{\linewidth}\raggedright
SOLUTION
\end{minipage} & \begin{minipage}[b]{\linewidth}\raggedright
SUMMARY
\end{minipage} & \begin{minipage}[b]{\linewidth}\raggedright
SCHEDULE
\end{minipage} \\
\midrule\noalign{}
\endhead
\bottomrule\noalign{}
\endlastfoot
\end{longtable}

\subsubsection{Description}\label{description-251}

HMMULRGT defines the transmissibility between two regions gradient
parameters, for when the History Match Gradient option has been
activated by the HMDIMS keyword in the RUNSPEC section.

See also the HMDIMS keyword in the RUNSPEC section that specifies the
dimensions for the gradient option, including the maximum number of
gradient regions that can be used with the History Match Gradient
option.

This keyword is not supported by OPM Flow but has no effect on the
results so it will be ignored.

\subsection{HMMULTFT -- History Match Fault Transmissibility Gradient
Cumulative
Multipliers}\label{hmmultft-history-match-fault-transmissibility-gradient-cumulative-multipliers}

\begin{longtable}[]{@{}
  >{\raggedright\arraybackslash}p{(\linewidth - 14\tabcolsep) * \real{0.1250}}
  >{\raggedright\arraybackslash}p{(\linewidth - 14\tabcolsep) * \real{0.1250}}
  >{\raggedright\arraybackslash}p{(\linewidth - 14\tabcolsep) * \real{0.1250}}
  >{\raggedright\arraybackslash}p{(\linewidth - 14\tabcolsep) * \real{0.1250}}
  >{\raggedright\arraybackslash}p{(\linewidth - 14\tabcolsep) * \real{0.1250}}
  >{\raggedright\arraybackslash}p{(\linewidth - 14\tabcolsep) * \real{0.1250}}
  >{\raggedright\arraybackslash}p{(\linewidth - 14\tabcolsep) * \real{0.1250}}
  >{\raggedright\arraybackslash}p{(\linewidth - 14\tabcolsep) * \real{0.1250}}@{}}
\toprule\noalign{}
\begin{minipage}[b]{\linewidth}\raggedright
RUNSPEC
\end{minipage} & \begin{minipage}[b]{\linewidth}\raggedright
GRID
\end{minipage} & \begin{minipage}[b]{\linewidth}\raggedright
EDIT
\end{minipage} & \begin{minipage}[b]{\linewidth}\raggedright
PROPS
\end{minipage} & \begin{minipage}[b]{\linewidth}\raggedright
REGIONS
\end{minipage} & \begin{minipage}[b]{\linewidth}\raggedright
SOLUTION
\end{minipage} & \begin{minipage}[b]{\linewidth}\raggedright
SUMMARY
\end{minipage} & \begin{minipage}[b]{\linewidth}\raggedright
SCHEDULE
\end{minipage} \\
\midrule\noalign{}
\endhead
\bottomrule\noalign{}
\endlastfoot
\end{longtable}

\subsubsection{Description}\label{description-252}

HMMULTFT defines the history match fault transmissibility gradient
cumulative multipliers to be applied to the fault transmissibilities for
faults declared by the FAULT keyword in the GRID section, for when the
History Match Gradient option has been activated by the HMDIMS keyword
in the RUNSPEC section. The keyword allows for the re-scaling of the
existing fault transmissibilities calculated by OPM Flow, or if the
MULTFLT keyword has been entered, then HMMULTFT is applied to the
existing MULTFLT multipliers.

This keyword is not supported by OPM Flow but would change the results
if supported so the simulation will be stopped.

\subsection{HMMULTSG -- History Match Dual porosity Sigma Gradient
Cumulative
Multipliers}\label{hmmultsg-history-match-dual-porosity-sigma-gradient-cumulative-multipliers}

\begin{longtable}[]{@{}
  >{\raggedright\arraybackslash}p{(\linewidth - 14\tabcolsep) * \real{0.1250}}
  >{\raggedright\arraybackslash}p{(\linewidth - 14\tabcolsep) * \real{0.1250}}
  >{\raggedright\arraybackslash}p{(\linewidth - 14\tabcolsep) * \real{0.1250}}
  >{\raggedright\arraybackslash}p{(\linewidth - 14\tabcolsep) * \real{0.1250}}
  >{\raggedright\arraybackslash}p{(\linewidth - 14\tabcolsep) * \real{0.1250}}
  >{\raggedright\arraybackslash}p{(\linewidth - 14\tabcolsep) * \real{0.1250}}
  >{\raggedright\arraybackslash}p{(\linewidth - 14\tabcolsep) * \real{0.1250}}
  >{\raggedright\arraybackslash}p{(\linewidth - 14\tabcolsep) * \real{0.1250}}@{}}
\toprule\noalign{}
\begin{minipage}[b]{\linewidth}\raggedright
RUNSPEC
\end{minipage} & \begin{minipage}[b]{\linewidth}\raggedright
GRID
\end{minipage} & \begin{minipage}[b]{\linewidth}\raggedright
EDIT
\end{minipage} & \begin{minipage}[b]{\linewidth}\raggedright
PROPS
\end{minipage} & \begin{minipage}[b]{\linewidth}\raggedright
REGIONS
\end{minipage} & \begin{minipage}[b]{\linewidth}\raggedright
SOLUTION
\end{minipage} & \begin{minipage}[b]{\linewidth}\raggedright
SUMMARY
\end{minipage} & \begin{minipage}[b]{\linewidth}\raggedright
SCHEDULE
\end{minipage} \\
\midrule\noalign{}
\endhead
\bottomrule\noalign{}
\endlastfoot
\end{longtable}

\subsubsection{Description}\label{description-253}

HMMULTSG defines the history match dual porosity sigma parameter
gradient cumulative multipliers applied to the dual porosity sigma value
declared by the SIGMAV and SIGMAGDV keywords in the PROPS section, for
when the History Match Gradient option has been activated by the HMDIMS
keyword in the RUNSPEC section. In addition to the HMDIMS keyword,
either the DUALPERM keyword that activates the Dual Permeability option,
or the DUALPORO keyword that activates the Dual Porosity option for the
run, must be declared in the RUNSPEC section.

This keyword is not supported by OPM Flow but would change the results
if supported so the simulation will be stopped.

\subsection{HRFIN -- Define the Ratio of LGR Grid Blocks in the
R-Direction}\label{hrfin-define-the-ratio-of-lgr-grid-blocks-in-the-r-direction}

\begin{longtable}[]{@{}
  >{\raggedright\arraybackslash}p{(\linewidth - 14\tabcolsep) * \real{0.1250}}
  >{\raggedright\arraybackslash}p{(\linewidth - 14\tabcolsep) * \real{0.1250}}
  >{\raggedright\arraybackslash}p{(\linewidth - 14\tabcolsep) * \real{0.1250}}
  >{\raggedright\arraybackslash}p{(\linewidth - 14\tabcolsep) * \real{0.1250}}
  >{\raggedright\arraybackslash}p{(\linewidth - 14\tabcolsep) * \real{0.1250}}
  >{\raggedright\arraybackslash}p{(\linewidth - 14\tabcolsep) * \real{0.1250}}
  >{\raggedright\arraybackslash}p{(\linewidth - 14\tabcolsep) * \real{0.1250}}
  >{\raggedright\arraybackslash}p{(\linewidth - 14\tabcolsep) * \real{0.1250}}@{}}
\toprule\noalign{}
\begin{minipage}[b]{\linewidth}\raggedright
RUNSPEC
\end{minipage} & \begin{minipage}[b]{\linewidth}\raggedright
GRID
\end{minipage} & \begin{minipage}[b]{\linewidth}\raggedright
EDIT
\end{minipage} & \begin{minipage}[b]{\linewidth}\raggedright
PROPS
\end{minipage} & \begin{minipage}[b]{\linewidth}\raggedright
REGIONS
\end{minipage} & \begin{minipage}[b]{\linewidth}\raggedright
SOLUTION
\end{minipage} & \begin{minipage}[b]{\linewidth}\raggedright
SUMMARY
\end{minipage} & \begin{minipage}[b]{\linewidth}\raggedright
SCHEDULE
\end{minipage} \\
\midrule\noalign{}
\endhead
\bottomrule\noalign{}
\endlastfoot
\end{longtable}

\subsubsection{Description}\label{description-254}

HRFIN {[}Radial grids are not currently implemented in this version of
OPM Flow, but is expected to be incorporated in a future release.{]}
defines the ratio of grid blocks for the DRV keyword in the r-direction
via a vector within a Local Grid Refinement (``LGR'') as opposed to
defining the size for each cell for a Radial LGR Grid. The LGR keyword
in the RUNSPEC section should be activated to indicate an LGR is being
used, and the keyword HRFIN should be placed in between the RADIN (or
RAFDIN4) and ENDFIN keywords in the GRID section. The DRV keyword in the
GRID section defines the radial grid size in terms of the length, that
is feet for field units, this keyword defines the length as the ratio of
the previous cell size, staring with the inner radius (INRAD).

This keyword is not supported by OPM Flow but would change the results
if supported so the simulation will be stopped.

\begin{longtable}[]{@{}
  >{\raggedright\arraybackslash}p{(\linewidth - 6\tabcolsep) * \real{0.2500}}
  >{\raggedright\arraybackslash}p{(\linewidth - 6\tabcolsep) * \real{0.2500}}
  >{\raggedright\arraybackslash}p{(\linewidth - 6\tabcolsep) * \real{0.2500}}
  >{\raggedright\arraybackslash}p{(\linewidth - 6\tabcolsep) * \real{0.2500}}@{}}
\toprule\noalign{}
\begin{minipage}[b]{\linewidth}\raggedright
No.
\end{minipage} & \begin{minipage}[b]{\linewidth}\raggedright
Name
\end{minipage} & \begin{minipage}[b]{\linewidth}\raggedright
Description
\end{minipage} & \begin{minipage}[b]{\linewidth}\raggedright
Default
\end{minipage} \\
\midrule\noalign{}
\endhead
\bottomrule\noalign{}
\endlastfoot
Field & Metric & Laboratory & \\
1 & HRFIN & HRFIN is a vector of real numbers describing the ratio of
cell size for the grid blocks in the r-direction in a radial LGR for the
DRV keyword. Repeat counts may be used, for example 2*1.5. & None \\
dimensionless & dimensionless & dimensionless & \\
Notes: & & & \\
\end{longtable}

\emph{Table 6.46: HRFIN Keyword Description}

See also the DR, DRV, DTHETAV, and DZ keywords in the GRID section to
fully define a radial LGR model.

\subsubsection{Example}\label{example-167}

\begin{verbatim}
--
--       INNER RADIUS OF FIRST GRID BLOCK IN THE RADIAL DIRECTION
--
INRAD    0.25                                                                 /
--
--       DEFINE GRID BLOCK DRV RATIOS IN THE R DIRECTION
--
HRFIN
         1.50  2.00  3.00  5.00  7.00  10.00                                  /
\end{verbatim}

The above example defines the size of the cells in the R direction based
on NR equals 7, resulting in NR-1 values on the RADFIN keyword in the
GRID section. Note the INRAD keyword to define the inner radius of the
radial grid.

\subsection{HXFIN -- Define the Ratio of LGR Grid Blocks in the
X-Direction}\label{hxfin-define-the-ratio-of-lgr-grid-blocks-in-the-x-direction}

\begin{longtable}[]{@{}
  >{\raggedright\arraybackslash}p{(\linewidth - 14\tabcolsep) * \real{0.1250}}
  >{\raggedright\arraybackslash}p{(\linewidth - 14\tabcolsep) * \real{0.1250}}
  >{\raggedright\arraybackslash}p{(\linewidth - 14\tabcolsep) * \real{0.1250}}
  >{\raggedright\arraybackslash}p{(\linewidth - 14\tabcolsep) * \real{0.1250}}
  >{\raggedright\arraybackslash}p{(\linewidth - 14\tabcolsep) * \real{0.1250}}
  >{\raggedright\arraybackslash}p{(\linewidth - 14\tabcolsep) * \real{0.1250}}
  >{\raggedright\arraybackslash}p{(\linewidth - 14\tabcolsep) * \real{0.1250}}
  >{\raggedright\arraybackslash}p{(\linewidth - 14\tabcolsep) * \real{0.1250}}@{}}
\toprule\noalign{}
\begin{minipage}[b]{\linewidth}\raggedright
RUNSPEC
\end{minipage} & \begin{minipage}[b]{\linewidth}\raggedright
GRID
\end{minipage} & \begin{minipage}[b]{\linewidth}\raggedright
EDIT
\end{minipage} & \begin{minipage}[b]{\linewidth}\raggedright
PROPS
\end{minipage} & \begin{minipage}[b]{\linewidth}\raggedright
REGIONS
\end{minipage} & \begin{minipage}[b]{\linewidth}\raggedright
SOLUTION
\end{minipage} & \begin{minipage}[b]{\linewidth}\raggedright
SUMMARY
\end{minipage} & \begin{minipage}[b]{\linewidth}\raggedright
SCHEDULE
\end{minipage} \\
\midrule\noalign{}
\endhead
\bottomrule\noalign{}
\endlastfoot
\end{longtable}

\subsubsection{Description}\label{description-255}

HXFIN defines the split ratio of grid blocks for the DXV keyword in the
x-direction via a vector within a Local Grid Refinement (``LGR'') as
opposed to defining the size for each cell for a Cartesian LGR Grid. The
LGR keyword in the RUNSPEC section should be activated to indicate an
LGR is being used, and the keyword HXFIN should be placed in between the
CARFIN and ENDFIN keywords in the GRID section. The DXV keyword in the
GRID section defines the grid size in terms of the length, that is feet
for field units, this keyword defines the length as the ratio of the
coarse cells.

This keyword is not supported by OPM Flow but would change the results
if supported so the simulation will be stopped.

\begin{longtable}[]{@{}
  >{\raggedright\arraybackslash}p{(\linewidth - 6\tabcolsep) * \real{0.2500}}
  >{\raggedright\arraybackslash}p{(\linewidth - 6\tabcolsep) * \real{0.2500}}
  >{\raggedright\arraybackslash}p{(\linewidth - 6\tabcolsep) * \real{0.2500}}
  >{\raggedright\arraybackslash}p{(\linewidth - 6\tabcolsep) * \real{0.2500}}@{}}
\toprule\noalign{}
\begin{minipage}[b]{\linewidth}\raggedright
No.
\end{minipage} & \begin{minipage}[b]{\linewidth}\raggedright
Name
\end{minipage} & \begin{minipage}[b]{\linewidth}\raggedright
Description
\end{minipage} & \begin{minipage}[b]{\linewidth}\raggedright
Default
\end{minipage} \\
\midrule\noalign{}
\endhead
\bottomrule\noalign{}
\endlastfoot
Field & Metric & Laboratory & \\
1 & HXFIN & HXFIN is a vector of real numbers describing the ratio of
cell size for the grid blocks in the x-direction in a Cartesian LGR
grid. Repeat counts may be used, for example 2*2.0. & None \\
dimensionless & dimensionless & dimensionless & \\
Notes: & & & \\
\end{longtable}

\emph{Table 6.47: HXFIN Keyword Description}

See also the CARFIN, ENDFIN, HYFIN, and HZFIN keywords in the GRID
section to fully define a Cartesian LGR grid model.

\subsubsection{Example}\label{example-168}

\begin{verbatim}
--
--       CARFIN LGR GRID COMMANDS
--
--       LGR        ----- HOST GRID ------   -- CARFIN GRID --  MAX     HOST
--       NAME       I1  I2  J1  J2  K1  K2     NX    NY    NZ   WELLS   NAME
CARFIN
         LGR-OP01   24  25  86  87   1  50      5     3    50     1     GLOBAL /
--
--       DEFINE LGR GRID BLOCK IN THE X-DIRECTION
NXFIN
         3  2                                                                  /
--
--       DEFINE GRID BLOCK LGR RATIOS IN THE X-DIRECTION
--
HXFIN
         1.00  2.00  3.00  2.00  1.00                                          /
ENDFIN
\end{verbatim}

The above example defines the size of the cells in the x-direction based
on NX equals five on the CARFIN keyword in the GRID section.

\subsection{HYFIN -- Define the Ratio of LGR Grid Blocks in the
Y-Direction}\label{hyfin-define-the-ratio-of-lgr-grid-blocks-in-the-y-direction}

\begin{longtable}[]{@{}
  >{\raggedright\arraybackslash}p{(\linewidth - 14\tabcolsep) * \real{0.1250}}
  >{\raggedright\arraybackslash}p{(\linewidth - 14\tabcolsep) * \real{0.1250}}
  >{\raggedright\arraybackslash}p{(\linewidth - 14\tabcolsep) * \real{0.1250}}
  >{\raggedright\arraybackslash}p{(\linewidth - 14\tabcolsep) * \real{0.1250}}
  >{\raggedright\arraybackslash}p{(\linewidth - 14\tabcolsep) * \real{0.1250}}
  >{\raggedright\arraybackslash}p{(\linewidth - 14\tabcolsep) * \real{0.1250}}
  >{\raggedright\arraybackslash}p{(\linewidth - 14\tabcolsep) * \real{0.1250}}
  >{\raggedright\arraybackslash}p{(\linewidth - 14\tabcolsep) * \real{0.1250}}@{}}
\toprule\noalign{}
\begin{minipage}[b]{\linewidth}\raggedright
RUNSPEC
\end{minipage} & \begin{minipage}[b]{\linewidth}\raggedright
GRID
\end{minipage} & \begin{minipage}[b]{\linewidth}\raggedright
EDIT
\end{minipage} & \begin{minipage}[b]{\linewidth}\raggedright
PROPS
\end{minipage} & \begin{minipage}[b]{\linewidth}\raggedright
REGIONS
\end{minipage} & \begin{minipage}[b]{\linewidth}\raggedright
SOLUTION
\end{minipage} & \begin{minipage}[b]{\linewidth}\raggedright
SUMMARY
\end{minipage} & \begin{minipage}[b]{\linewidth}\raggedright
SCHEDULE
\end{minipage} \\
\midrule\noalign{}
\endhead
\bottomrule\noalign{}
\endlastfoot
\end{longtable}

\subsubsection{Description}\label{description-256}

HYFIN defines the split ratio of grid blocks for the DYV keyword in the
y-direction via a vector within a Local Grid Refinement (``LGR'') as
opposed to defining the size for each cell for a Cartesian LGR Grid. The
LGR keyword in the RUNSPEC section should be activated to indicate an
LGR is being used, and the keyword HYFIN should be placed in between the
CARFIN and ENDFIN keywords in the GRID section. The DYV keyword in the
GRID section defines the grid size in terms of the length, that is feet
for field units, this keyword defines the length as the ratio of the
coarse cells.

This keyword is not supported by OPM Flow but would change the results
if supported so the simulation will be stopped.

\begin{longtable}[]{@{}
  >{\raggedright\arraybackslash}p{(\linewidth - 6\tabcolsep) * \real{0.2500}}
  >{\raggedright\arraybackslash}p{(\linewidth - 6\tabcolsep) * \real{0.2500}}
  >{\raggedright\arraybackslash}p{(\linewidth - 6\tabcolsep) * \real{0.2500}}
  >{\raggedright\arraybackslash}p{(\linewidth - 6\tabcolsep) * \real{0.2500}}@{}}
\toprule\noalign{}
\begin{minipage}[b]{\linewidth}\raggedright
No.
\end{minipage} & \begin{minipage}[b]{\linewidth}\raggedright
Name
\end{minipage} & \begin{minipage}[b]{\linewidth}\raggedright
Description
\end{minipage} & \begin{minipage}[b]{\linewidth}\raggedright
Default
\end{minipage} \\
\midrule\noalign{}
\endhead
\bottomrule\noalign{}
\endlastfoot
Field & Metric & Laboratory & \\
1 & HYFIN & HYFIN is a vector of real numbers describing the ratio of
cell size for the grid blocks in the y-direction in a Cartesian LGR
grid. Repeat counts may be used, for example 2*2.0. & None \\
dimensionless & dimensionless & dimensionless & \\
Notes: & & & \\
\end{longtable}

\emph{Table 6.48: HYFIN Keyword Description}

See also the CARFIN, ENDFIN, HXFIN, and HZFIN keywords in the GRID
section to fully define a Cartesian LGR grid model.

\subsubsection{Example}\label{example-169}

\begin{verbatim}
--
--       CARFIN LGR GRID COMMANDS
--
--       LGR        ----- HOST GRID ------   -- CARFIN GRID --  MAX     HOST
--       NAME       I1  I2  J1  J2  K1  K2     NX    NY    NZ   WELLS   NAME
CARFIN
         LGR-OP01   24  25  86  87   1  50      3     5    50     1     GLOBAL /
--
--       DEFINE LGR GRID BLOCK IN THE Y-DIRECTION
NYFIN
         3  2                                                                  /
--
--       DEFINE GRID BLOCK LGR RATIOS IN THE Y-DIRECTION
--
HYFIN
         1.00  2.00  3.00  2.00  1.00                                          /
ENDFIN
\end{verbatim}

The above example defines the size of the cells in the y-direction based
on NY equals five on the CARFIN keyword in the GRID section.

\subsection{HZFIN -- Define the Ratio of LGR Grid Blocks in the
Z-Direction}\label{hzfin-define-the-ratio-of-lgr-grid-blocks-in-the-z-direction}

\begin{longtable}[]{@{}
  >{\raggedright\arraybackslash}p{(\linewidth - 14\tabcolsep) * \real{0.1250}}
  >{\raggedright\arraybackslash}p{(\linewidth - 14\tabcolsep) * \real{0.1250}}
  >{\raggedright\arraybackslash}p{(\linewidth - 14\tabcolsep) * \real{0.1250}}
  >{\raggedright\arraybackslash}p{(\linewidth - 14\tabcolsep) * \real{0.1250}}
  >{\raggedright\arraybackslash}p{(\linewidth - 14\tabcolsep) * \real{0.1250}}
  >{\raggedright\arraybackslash}p{(\linewidth - 14\tabcolsep) * \real{0.1250}}
  >{\raggedright\arraybackslash}p{(\linewidth - 14\tabcolsep) * \real{0.1250}}
  >{\raggedright\arraybackslash}p{(\linewidth - 14\tabcolsep) * \real{0.1250}}@{}}
\toprule\noalign{}
\begin{minipage}[b]{\linewidth}\raggedright
RUNSPEC
\end{minipage} & \begin{minipage}[b]{\linewidth}\raggedright
GRID
\end{minipage} & \begin{minipage}[b]{\linewidth}\raggedright
EDIT
\end{minipage} & \begin{minipage}[b]{\linewidth}\raggedright
PROPS
\end{minipage} & \begin{minipage}[b]{\linewidth}\raggedright
REGIONS
\end{minipage} & \begin{minipage}[b]{\linewidth}\raggedright
SOLUTION
\end{minipage} & \begin{minipage}[b]{\linewidth}\raggedright
SUMMARY
\end{minipage} & \begin{minipage}[b]{\linewidth}\raggedright
SCHEDULE
\end{minipage} \\
\midrule\noalign{}
\endhead
\bottomrule\noalign{}
\endlastfoot
\end{longtable}

\subsubsection{Description}\label{description-257}

HZFIN defines the split ratio of grid blocks for the DZV keyword in the
z-direction via a vector within a Local Grid Refinement (``LGR'') as
opposed to defining the size for each cell for a Cartesian LGR Grid. The
LGR keyword in the RUNSPEC section should be activated to indicate an
LGR is being used, and the keyword HYFIN should be placed in between the
CARFIN and ENDFIN keywords in the GRID section. The DZV keyword in the
GRID section defines the grid size in terms of the length, that is feet
for field units, this keyword defines the length as the ratio of the
coarse cells.

This keyword is not supported by OPM Flow but would change the results
if supported so the simulation will be stopped.

\begin{longtable}[]{@{}
  >{\raggedright\arraybackslash}p{(\linewidth - 6\tabcolsep) * \real{0.2500}}
  >{\raggedright\arraybackslash}p{(\linewidth - 6\tabcolsep) * \real{0.2500}}
  >{\raggedright\arraybackslash}p{(\linewidth - 6\tabcolsep) * \real{0.2500}}
  >{\raggedright\arraybackslash}p{(\linewidth - 6\tabcolsep) * \real{0.2500}}@{}}
\toprule\noalign{}
\begin{minipage}[b]{\linewidth}\raggedright
No.
\end{minipage} & \begin{minipage}[b]{\linewidth}\raggedright
Name
\end{minipage} & \begin{minipage}[b]{\linewidth}\raggedright
Description
\end{minipage} & \begin{minipage}[b]{\linewidth}\raggedright
Default
\end{minipage} \\
\midrule\noalign{}
\endhead
\bottomrule\noalign{}
\endlastfoot
Field & Metric & Laboratory & \\
1 & HYZIN & HZFIN is a vector of real numbers describing the ratio of
cell size for the grid blocks in the z-direction in a Cartesian LGR
grid. Repeat counts may be used, for example 2*2.0. & None \\
dimensionless & dimensionless & dimensionless & \\
Notes: & & & \\
\end{longtable}

\emph{Table 6.49: HZFIN Keyword Description}

See also the CARFIN, ENDFIN, HXFIN, and HYFIN keywords in the GRID
section to fully define a Cartesian LGR grid model.

\subsubsection{Example}\label{example-170}

\begin{verbatim}
--
--       CARFIN LGR GRID COMMANDS
--
--       LGR        ----- HOST GRID ------   -- CARFIN GRID --  MAX     HOST
--       NAME       I1  I2  J1  J2  K1  K2     NX    NY    NZ   WELLS   NAME
CARFIN
         LGR-OP01   24  25  86  87   1  50      5     3    100    1     GLOBAL /
--
--       DEFINE LGR GRID BLOCK IN THE Z-DIRECTION
NZFIN
         50*2                                                                  /
--
--       DEFINE GRID BLOCK LGR RATIOS IN THE Z-DIRECTION
--
HZFIN
         50*2.0                                                                /
ENDFIN
\end{verbatim}

The above example defines the size of the cells in the z-direction based
on NZ equals 100 on the CARFIN keyword in the GRID section.

\subsection{IHOST -- Assign LGRs to Parallel Process
Number}\label{ihost-assign-lgrs-to-parallel-process-number}

\begin{longtable}[]{@{}
  >{\raggedright\arraybackslash}p{(\linewidth - 14\tabcolsep) * \real{0.1250}}
  >{\raggedright\arraybackslash}p{(\linewidth - 14\tabcolsep) * \real{0.1250}}
  >{\raggedright\arraybackslash}p{(\linewidth - 14\tabcolsep) * \real{0.1250}}
  >{\raggedright\arraybackslash}p{(\linewidth - 14\tabcolsep) * \real{0.1250}}
  >{\raggedright\arraybackslash}p{(\linewidth - 14\tabcolsep) * \real{0.1250}}
  >{\raggedright\arraybackslash}p{(\linewidth - 14\tabcolsep) * \real{0.1250}}
  >{\raggedright\arraybackslash}p{(\linewidth - 14\tabcolsep) * \real{0.1250}}
  >{\raggedright\arraybackslash}p{(\linewidth - 14\tabcolsep) * \real{0.1250}}@{}}
\toprule\noalign{}
\begin{minipage}[b]{\linewidth}\raggedright
RUNSPEC
\end{minipage} & \begin{minipage}[b]{\linewidth}\raggedright
GRID
\end{minipage} & \begin{minipage}[b]{\linewidth}\raggedright
EDIT
\end{minipage} & \begin{minipage}[b]{\linewidth}\raggedright
PROPS
\end{minipage} & \begin{minipage}[b]{\linewidth}\raggedright
REGIONS
\end{minipage} & \begin{minipage}[b]{\linewidth}\raggedright
SOLUTION
\end{minipage} & \begin{minipage}[b]{\linewidth}\raggedright
SUMMARY
\end{minipage} & \begin{minipage}[b]{\linewidth}\raggedright
SCHEDULE
\end{minipage} \\
\midrule\noalign{}
\endhead
\bottomrule\noalign{}
\endlastfoot
\end{longtable}

\subsubsection{Description}\label{description-258}

The IHIST keyword assigns Local Grid Refinements (``LGR'') to a parallel
process number, for when the PARALLEL keyword has been invoked in the
RUNSPEC section.

This keyword is not supported by OPM Flow but would change the results
if supported so the simulation will be stopped.

\subsection{IMPORT -- Import Grid File Data at the Current
Position}\label{import-import-grid-file-data-at-the-current-position}

\begin{longtable}[]{@{}
  >{\raggedright\arraybackslash}p{(\linewidth - 14\tabcolsep) * \real{0.1250}}
  >{\raggedright\arraybackslash}p{(\linewidth - 14\tabcolsep) * \real{0.1250}}
  >{\raggedright\arraybackslash}p{(\linewidth - 14\tabcolsep) * \real{0.1250}}
  >{\raggedright\arraybackslash}p{(\linewidth - 14\tabcolsep) * \real{0.1250}}
  >{\raggedright\arraybackslash}p{(\linewidth - 14\tabcolsep) * \real{0.1250}}
  >{\raggedright\arraybackslash}p{(\linewidth - 14\tabcolsep) * \real{0.1250}}
  >{\raggedright\arraybackslash}p{(\linewidth - 14\tabcolsep) * \real{0.1250}}
  >{\raggedright\arraybackslash}p{(\linewidth - 14\tabcolsep) * \real{0.1250}}@{}}
\toprule\noalign{}
\begin{minipage}[b]{\linewidth}\raggedright
RUNSPEC
\end{minipage} & \begin{minipage}[b]{\linewidth}\raggedright
GRID
\end{minipage} & \begin{minipage}[b]{\linewidth}\raggedright
EDIT
\end{minipage} & \begin{minipage}[b]{\linewidth}\raggedright
PROPS
\end{minipage} & \begin{minipage}[b]{\linewidth}\raggedright
REGIONS
\end{minipage} & \begin{minipage}[b]{\linewidth}\raggedright
SOLUTION
\end{minipage} & \begin{minipage}[b]{\linewidth}\raggedright
SUMMARY
\end{minipage} & \begin{minipage}[b]{\linewidth}\raggedright
SCHEDULE
\end{minipage} \\
\midrule\noalign{}
\endhead
\bottomrule\noalign{}
\endlastfoot
\end{longtable}

\subsubsection{Description}\label{description-259}

The IMPORT keyword informs the simulator to import data from the
specified IMPORT file. When the end of the IMPORT file is reached, input
data is read from the next keyword in the current file. Normally IMPORT
files are generated by grid pre-processing software and the keyword
allows for both formatted and unformatted (binary) files to be loaded.
The keyword can be used to import any valid grid arrays within a
section, for example the EQLNUM array in the REGIONS section

\begin{longtable}[]{@{}
  >{\raggedright\arraybackslash}p{(\linewidth - 6\tabcolsep) * \real{0.2500}}
  >{\raggedright\arraybackslash}p{(\linewidth - 6\tabcolsep) * \real{0.2500}}
  >{\raggedright\arraybackslash}p{(\linewidth - 6\tabcolsep) * \real{0.2500}}
  >{\raggedright\arraybackslash}p{(\linewidth - 6\tabcolsep) * \real{0.2500}}@{}}
\toprule\noalign{}
\begin{minipage}[b]{\linewidth}\raggedright
No.
\end{minipage} & \begin{minipage}[b]{\linewidth}\raggedright
Name
\end{minipage} & \begin{minipage}[b]{\linewidth}\raggedright
Description
\end{minipage} & \begin{minipage}[b]{\linewidth}\raggedright
Default
\end{minipage} \\
\midrule\noalign{}
\endhead
\bottomrule\noalign{}
\endlastfoot
1 & FILENAME & A character string enclosed in quotes that defines a file
to be imported and to be processed by OPM Flow. & None \\
2 & FMTOPT & & U \\
Notes: & & & \\
\end{longtable}

\emph{Table 6.50: IMPORT Keyword Description}

\subsubsection{Examples}\label{examples-12}

The first example shown below loads the grid file from the same
directory as the data file.

\begin{verbatim}
--
--       LOAD A IMPORT FILE
--
IMPORT
         'NOR-OPM-A00-GRID.EGRID'  /
\end{verbatim}

The next example loads the same file one directory above from where the
data file is located.

\begin{verbatim}
--
--       LOAD A IMPORT FILE
--
IMPORT
         '../NOR-OPM-A00-GRID.EGRID'  /
\end{verbatim}

The final example loads the same file from a separate include
subdirectory in the parent directory of the data file.

\begin{verbatim}
--
--       LOAD A IMPORT FILE
--
IMPORT
         '../INCLUDE/NOR-OPM-A00-GRID.EGRID'  /
\end{verbatim}

\subsection{INCLUDE -- Load Another Data File at the Current
Position}\label{include-load-another-data-file-at-the-current-position-2}

The INCLUDE keyword informs OPM Flow to continue reading input data from
the specified INCLUDE file. When the end of the INCLUDE file is reached,
or the ENDINC keyword is encountered in the included file, input data is
read from the next keyword in the current file.

See INCLUDE -- Load Another Data File at the Current Position in the
GLOBAL section for a full description.

\subsection{INIT -- Activate the INIT File
Option}\label{init-activate-the-init-file-option}

\begin{longtable}[]{@{}
  >{\raggedright\arraybackslash}p{(\linewidth - 14\tabcolsep) * \real{0.1250}}
  >{\raggedright\arraybackslash}p{(\linewidth - 14\tabcolsep) * \real{0.1250}}
  >{\raggedright\arraybackslash}p{(\linewidth - 14\tabcolsep) * \real{0.1250}}
  >{\raggedright\arraybackslash}p{(\linewidth - 14\tabcolsep) * \real{0.1250}}
  >{\raggedright\arraybackslash}p{(\linewidth - 14\tabcolsep) * \real{0.1250}}
  >{\raggedright\arraybackslash}p{(\linewidth - 14\tabcolsep) * \real{0.1250}}
  >{\raggedright\arraybackslash}p{(\linewidth - 14\tabcolsep) * \real{0.1250}}
  >{\raggedright\arraybackslash}p{(\linewidth - 14\tabcolsep) * \real{0.1250}}@{}}
\toprule\noalign{}
\begin{minipage}[b]{\linewidth}\raggedright
RUNSPEC
\end{minipage} & \begin{minipage}[b]{\linewidth}\raggedright
GRID
\end{minipage} & \begin{minipage}[b]{\linewidth}\raggedright
EDIT
\end{minipage} & \begin{minipage}[b]{\linewidth}\raggedright
PROPS
\end{minipage} & \begin{minipage}[b]{\linewidth}\raggedright
REGIONS
\end{minipage} & \begin{minipage}[b]{\linewidth}\raggedright
SOLUTION
\end{minipage} & \begin{minipage}[b]{\linewidth}\raggedright
SUMMARY
\end{minipage} & \begin{minipage}[b]{\linewidth}\raggedright
SCHEDULE
\end{minipage} \\
\midrule\noalign{}
\endhead
\bottomrule\noalign{}
\endlastfoot
\end{longtable}

\subsubsection{Description}\label{description-260}

This keyword switches on the writing of the INIT file that contains the
static data specified in the GRID, PROPS and REGIONS sections. For
example, the PORO, PERM and NTG arrays from the GRID section. The data
is used in post-processing software, for example OPM ResInsight, to
visualize the static grid properties.

The INIT file can either be written out in formatted form as ASCII
i.e.~text files, if the FMTOUT keyword has been activated, or binary
format if the FMTOUT keyword has not been activated. Normally, this
option is always activated by the user and when activated the binary
form of the file is used.

There is no data required for this keyword and there is no terminating
``/'' for this keyword.

\subsubsection{Example}\label{example-171}

\begin{verbatim}
--
--       ACTIVATE WRITING THE INIT FILE FOR POST-PROCESSING
INIT
\end{verbatim}

The above example switches on the writing of the INIT file for
post-processing in ResInsight.

\subsection{INRAD -- Define the Inner Radius of a Radial
Grid}\label{inrad-define-the-inner-radius-of-a-radial-grid}

\begin{longtable}[]{@{}
  >{\raggedright\arraybackslash}p{(\linewidth - 14\tabcolsep) * \real{0.1250}}
  >{\raggedright\arraybackslash}p{(\linewidth - 14\tabcolsep) * \real{0.1250}}
  >{\raggedright\arraybackslash}p{(\linewidth - 14\tabcolsep) * \real{0.1250}}
  >{\raggedright\arraybackslash}p{(\linewidth - 14\tabcolsep) * \real{0.1250}}
  >{\raggedright\arraybackslash}p{(\linewidth - 14\tabcolsep) * \real{0.1250}}
  >{\raggedright\arraybackslash}p{(\linewidth - 14\tabcolsep) * \real{0.1250}}
  >{\raggedright\arraybackslash}p{(\linewidth - 14\tabcolsep) * \real{0.1250}}
  >{\raggedright\arraybackslash}p{(\linewidth - 14\tabcolsep) * \real{0.1250}}@{}}
\toprule\noalign{}
\begin{minipage}[b]{\linewidth}\raggedright
RUNSPEC
\end{minipage} & \begin{minipage}[b]{\linewidth}\raggedright
GRID
\end{minipage} & \begin{minipage}[b]{\linewidth}\raggedright
EDIT
\end{minipage} & \begin{minipage}[b]{\linewidth}\raggedright
PROPS
\end{minipage} & \begin{minipage}[b]{\linewidth}\raggedright
REGIONS
\end{minipage} & \begin{minipage}[b]{\linewidth}\raggedright
SOLUTION
\end{minipage} & \begin{minipage}[b]{\linewidth}\raggedright
SUMMARY
\end{minipage} & \begin{minipage}[b]{\linewidth}\raggedright
SCHEDULE
\end{minipage} \\
\midrule\noalign{}
\endhead
\bottomrule\noalign{}
\endlastfoot
\end{longtable}

\subsubsection{Description}\label{description-261}

INRAD defines the inner radius of the reservoir model for a radial grid
geometry. The RADIAL or SPIDER keyword in the RUNSPEC should be
activated to indicate that radial geometry is being used.

\begin{longtable}[]{@{}
  >{\raggedright\arraybackslash}p{(\linewidth - 6\tabcolsep) * \real{0.2500}}
  >{\raggedright\arraybackslash}p{(\linewidth - 6\tabcolsep) * \real{0.2500}}
  >{\raggedright\arraybackslash}p{(\linewidth - 6\tabcolsep) * \real{0.2500}}
  >{\raggedright\arraybackslash}p{(\linewidth - 6\tabcolsep) * \real{0.2500}}@{}}
\toprule\noalign{}
\begin{minipage}[b]{\linewidth}\raggedright
No.
\end{minipage} & \begin{minipage}[b]{\linewidth}\raggedright
Name
\end{minipage} & \begin{minipage}[b]{\linewidth}\raggedright
Description
\end{minipage} & \begin{minipage}[b]{\linewidth}\raggedright
Default
\end{minipage} \\
\midrule\noalign{}
\endhead
\bottomrule\noalign{}
\endlastfoot
Field & Metric & Laboratory & \\
1 & INRAD & A single real positive number defining the inner radius of a
radial grid. & None \\
feet & m & cm & \\
Notes: & & & \\
\end{longtable}

\emph{Table 6.51: INRAD Keyword Description}

Note that the SPIDER keyword activates OPM Flow's radial grid geometry
option for the model. This keyword will create a spiderweb-shaped grid
based on a corner-point grid using the standard radial grid keywords:
INRAD, DRV, DTHETAV, DZ/DZV etc. in the GRID the section. A spider grid
can be viewed in 3D in OPM ResInsight unlike radial grids that cannot be
viewed in the software. To overcome this, the simulator now converts
radial grids to Irregular Corner-Point Grids and adjusts the model's
pore volume to reflect radial coordinates; thus, overcoming the display
limitation.

See also the DR, DRV, DTHETA, DTHETAV and TOPS keywords to fully define
a radial or spider grid in the model.

\subsubsection{Example}\label{example-172}

\begin{verbatim}
--
--       INNER RADIUS OF FIRST GRID BLOCK IN THE RADIAL DIRECTION
--
INRAD
         0.25                                                                   /
\end{verbatim}

The above example defines the inner radius of a radial grid to be 0.25
feet.

\subsection{IONROCK -- Define the Ion Exchange Capacity for All the
Cells}\label{ionrock-define-the-ion-exchange-capacity-for-all-the-cells}

\begin{longtable}[]{@{}
  >{\raggedright\arraybackslash}p{(\linewidth - 14\tabcolsep) * \real{0.1250}}
  >{\raggedright\arraybackslash}p{(\linewidth - 14\tabcolsep) * \real{0.1250}}
  >{\raggedright\arraybackslash}p{(\linewidth - 14\tabcolsep) * \real{0.1250}}
  >{\raggedright\arraybackslash}p{(\linewidth - 14\tabcolsep) * \real{0.1250}}
  >{\raggedright\arraybackslash}p{(\linewidth - 14\tabcolsep) * \real{0.1250}}
  >{\raggedright\arraybackslash}p{(\linewidth - 14\tabcolsep) * \real{0.1250}}
  >{\raggedright\arraybackslash}p{(\linewidth - 14\tabcolsep) * \real{0.1250}}
  >{\raggedright\arraybackslash}p{(\linewidth - 14\tabcolsep) * \real{0.1250}}@{}}
\toprule\noalign{}
\begin{minipage}[b]{\linewidth}\raggedright
RUNSPEC
\end{minipage} & \begin{minipage}[b]{\linewidth}\raggedright
GRID
\end{minipage} & \begin{minipage}[b]{\linewidth}\raggedright
EDIT
\end{minipage} & \begin{minipage}[b]{\linewidth}\raggedright
PROPS
\end{minipage} & \begin{minipage}[b]{\linewidth}\raggedright
REGIONS
\end{minipage} & \begin{minipage}[b]{\linewidth}\raggedright
SOLUTION
\end{minipage} & \begin{minipage}[b]{\linewidth}\raggedright
SUMMARY
\end{minipage} & \begin{minipage}[b]{\linewidth}\raggedright
SCHEDULE
\end{minipage} \\
\midrule\noalign{}
\endhead
\bottomrule\noalign{}
\endlastfoot
\end{longtable}

\subsubsection{Description}\label{description-262}

The IONROCK keyword defines the ion exchange capacity for all the cells
in the model, for when the brine phase has been activated by the BRINE
keyword and the Multi-Component Brine model, that allows for the water
phase to have multiple water salinities, has been activated by the ECLMC
keyword. Both keywords are in the RUNSPEC section.

This keyword is not supported by OPM Flow but would change the results
if supported so the simulation will be stopped.

\subsection{ISOLNUM -- Define the Independent Reservoir
Regions}\label{isolnum-define-the-independent-reservoir-regions}

\begin{longtable}[]{@{}
  >{\raggedright\arraybackslash}p{(\linewidth - 14\tabcolsep) * \real{0.1250}}
  >{\raggedright\arraybackslash}p{(\linewidth - 14\tabcolsep) * \real{0.1250}}
  >{\raggedright\arraybackslash}p{(\linewidth - 14\tabcolsep) * \real{0.1250}}
  >{\raggedright\arraybackslash}p{(\linewidth - 14\tabcolsep) * \real{0.1250}}
  >{\raggedright\arraybackslash}p{(\linewidth - 14\tabcolsep) * \real{0.1250}}
  >{\raggedright\arraybackslash}p{(\linewidth - 14\tabcolsep) * \real{0.1250}}
  >{\raggedright\arraybackslash}p{(\linewidth - 14\tabcolsep) * \real{0.1250}}
  >{\raggedright\arraybackslash}p{(\linewidth - 14\tabcolsep) * \real{0.1250}}@{}}
\toprule\noalign{}
\begin{minipage}[b]{\linewidth}\raggedright
RUNSPEC
\end{minipage} & \begin{minipage}[b]{\linewidth}\raggedright
GRID
\end{minipage} & \begin{minipage}[b]{\linewidth}\raggedright
EDIT
\end{minipage} & \begin{minipage}[b]{\linewidth}\raggedright
PROPS
\end{minipage} & \begin{minipage}[b]{\linewidth}\raggedright
REGIONS
\end{minipage} & \begin{minipage}[b]{\linewidth}\raggedright
SOLUTION
\end{minipage} & \begin{minipage}[b]{\linewidth}\raggedright
SUMMARY
\end{minipage} & \begin{minipage}[b]{\linewidth}\raggedright
SCHEDULE
\end{minipage} \\
\midrule\noalign{}
\endhead
\bottomrule\noalign{}
\endlastfoot
\end{longtable}

\subsubsection{Description}\label{description-263}

The ISOLNUM keyword defines areas of the grid that consists of isolated
reservoirs where the only form of communication between the reservoirs
is via wellbore connections This enables the reservoir flow equations to
be solved independently for greater computational efficiency.

\begin{longtable}[]{@{}
  >{\raggedright\arraybackslash}p{(\linewidth - 6\tabcolsep) * \real{0.2500}}
  >{\raggedright\arraybackslash}p{(\linewidth - 6\tabcolsep) * \real{0.2500}}
  >{\raggedright\arraybackslash}p{(\linewidth - 6\tabcolsep) * \real{0.2500}}
  >{\raggedright\arraybackslash}p{(\linewidth - 6\tabcolsep) * \real{0.2500}}@{}}
\toprule\noalign{}
\begin{minipage}[b]{\linewidth}\raggedright
No.
\end{minipage} & \begin{minipage}[b]{\linewidth}\raggedright
Name
\end{minipage} & \begin{minipage}[b]{\linewidth}\raggedright
Description
\end{minipage} & \begin{minipage}[b]{\linewidth}\raggedright
Default
\end{minipage} \\
\midrule\noalign{}
\endhead
\bottomrule\noalign{}
\endlastfoot
1 & ISOLNUM & ISOLNUM defines an array of positive integers assigning a
grid cell to a particular isolated reservoir region. The maximum number
of ISOLNUM regions is set by the NRFREG variable on the REGDIMS keyword
in the RUNSPEC section. & 1 \\
Notes: & & & \\
\end{longtable}

\emph{Table 6.52: ISOLNUM Keyword Description}

This keyword is not supported by OPM Flow but would change the results
if supported so the simulation will be stopped.

\subsubsection{Example}\label{example-173}

The example below defines three separate independent reservoirs; the
first reservoir covers the whole grid and layers 1 to 50, reservoir two
cover the whole grid and layers 52 to 150, and finally the third
reservoir again covers the whole grid but with layers 152 to 300. The
layers 51 and 151 are shale layers made inactive by setting ISOLNUM to
zero.

\begin{verbatim}
--
--       ARRAY     CONSTANT       ---------- BOX ---------
--                                I1  I2   J1  J2   K1  K2
EQUALS
         ISOLNUM    1             1*  1*   1*  1*   1   50 / DEFINED RESERVOIR 1
         ISOLNUM    0             1*  1*   1*  1*   51  51 / DEFINED A SHALE
         ISOLNUM    2             1*  1*   1*  1*   52 150 / DEFINED RESERVOIR 2
         ISOLNUM    0             1*  1*   1*  1*  151 151 / DEFINED A SHALE
         ISOLNUM    3             1*  1*   1*  1*  152 300 / DEFINED RESERVOIR 3
/
\end{verbatim}

Note the above example has no effect as the keyword is ignored by the
simulator.

\subsection{JFUNC -- Activate the Leverett J-function
Option}\label{jfunc-activate-the-leverett-j-function-option}

\begin{longtable}[]{@{}
  >{\raggedright\arraybackslash}p{(\linewidth - 14\tabcolsep) * \real{0.1250}}
  >{\raggedright\arraybackslash}p{(\linewidth - 14\tabcolsep) * \real{0.1250}}
  >{\raggedright\arraybackslash}p{(\linewidth - 14\tabcolsep) * \real{0.1250}}
  >{\raggedright\arraybackslash}p{(\linewidth - 14\tabcolsep) * \real{0.1250}}
  >{\raggedright\arraybackslash}p{(\linewidth - 14\tabcolsep) * \real{0.1250}}
  >{\raggedright\arraybackslash}p{(\linewidth - 14\tabcolsep) * \real{0.1250}}
  >{\raggedright\arraybackslash}p{(\linewidth - 14\tabcolsep) * \real{0.1250}}
  >{\raggedright\arraybackslash}p{(\linewidth - 14\tabcolsep) * \real{0.1250}}@{}}
\toprule\noalign{}
\begin{minipage}[b]{\linewidth}\raggedright
RUNSPEC
\end{minipage} & \begin{minipage}[b]{\linewidth}\raggedright
GRID
\end{minipage} & \begin{minipage}[b]{\linewidth}\raggedright
EDIT
\end{minipage} & \begin{minipage}[b]{\linewidth}\raggedright
PROPS
\end{minipage} & \begin{minipage}[b]{\linewidth}\raggedright
REGIONS
\end{minipage} & \begin{minipage}[b]{\linewidth}\raggedright
SOLUTION
\end{minipage} & \begin{minipage}[b]{\linewidth}\raggedright
SUMMARY
\end{minipage} & \begin{minipage}[b]{\linewidth}\raggedright
SCHEDULE
\end{minipage} \\
\midrule\noalign{}
\endhead
\bottomrule\noalign{}
\endlastfoot
\end{longtable}

\subsubsection{Description}\label{description-264}

The JFUNC keyword activates the Leverett-J-Function {[}Leverett, M. C.;
``Capillary Behaviour in Porous Solids'', Trans. AIME (1941) 142,
152-168.{]} option which is a commonly used technique to normalize
capillary pressure based on laboratory measured core plugs porosity and
permeability values and the resulting capillary pressure data. The
keyword performs the calculation based on the parameters on the this
keyword combined with a cells porosity and permeability to perform the
scaling globally.

The keyword should only be used if end-point scaling is switched on
using the ENDSCALE keyword in the RUNSPEC section.

\begin{longtable}[]{@{}
  >{\raggedright\arraybackslash}p{(\linewidth - 6\tabcolsep) * \real{0.2500}}
  >{\raggedright\arraybackslash}p{(\linewidth - 6\tabcolsep) * \real{0.2500}}
  >{\raggedright\arraybackslash}p{(\linewidth - 6\tabcolsep) * \real{0.2500}}
  >{\raggedright\arraybackslash}p{(\linewidth - 6\tabcolsep) * \real{0.2500}}@{}}
\toprule\noalign{}
\begin{minipage}[b]{\linewidth}\raggedright
No.
\end{minipage} & \begin{minipage}[b]{\linewidth}\raggedright
Name
\end{minipage} & \begin{minipage}[b]{\linewidth}\raggedright
Description
\end{minipage} & \begin{minipage}[b]{\linewidth}\raggedright
Default
\end{minipage} \\
\midrule\noalign{}
\endhead
\bottomrule\noalign{}
\endlastfoot
Field & Metric & Laboratory & \\
1 & JFOPT & A character string that defines which capillary data sets
the J-Function option should be applied to, based on the following
options: & BOTH \\
2 & OWSTEN & A positive real number that defines oil-water surface
tension used to de-normalized J-Function data entered in the PROPS
section. & None \\
dynes/cm & dynes/cm & dynes/cm & \\
3 & OGSTEN & A positive real number that defines oil-gas surface tension
used to de-normalized J-Function data entered in the PROPS section. &
None \\
dynes/cm & dynes/cm & dynes/cm & \\
4 & ALPHA & A positive real value that defines an alternative power
value for the porosity term in the J-Function equation, that is instead
of \(\sqrt{\frac{k}{\mathrm{φ}}}\)use
\(\frac{{k}^{0.5}}{{\mathrm{φ}}^{\mathrm{α}}}\)instead in the
transformation. & 0.5 \\
5 & BETA & A positive real number that defines an alternative power
value for the permeability term in the J-Function equation, that is
instead of \(\sqrt{\frac{k}{\mathrm{φ}}}\)use
\(\frac{{k}^{\mathrm{β}}}{{\mathrm{φ}}^{0.5}}\)instead in the
transformation. & 0.5 \\
6 & PERM & PERM is a character string that sets the permeability array
to be used in the transform, based on the following options: & XY \\
Notes: & & & \\
\end{longtable}

\emph{Table 6.53: JFUNC Keyword Description}

Just like the relative permeability data capillary pressure data are
measured on core plugs with varying quality and perhaps from different
reservoirs. It is therefore necessary to determine averaged data, before
employing the data in engineering calculations. This is commonly done by
using the Leverett J-function {[}Leverett, M. C.; ``Capillary Behaviour
in Porous Solids'', Trans. AIME (1941) 142, 152-168.{]}, which is
defined as:

\begin{equation}\phantomsection\label{eq-6-4}{
J ({S}_{w}) = \frac{{P}_{c,\mathit{res}}({S}_{w}) \sqrt{\frac{k}{φ}}}{σ}
}\end{equation}

Where:

J (Sw) = dimensionless function of water saturation

Pc (Sw) = capillary pressure (kPa)

k = permeability, (m2)

φ = porosity (fraction)

σ = interfacial tension (mN/m)

Θ = contact angle

Sometimes the equation is stated with the cos θ term included, that is:

\begin{equation}\phantomsection\label{eq-6-5}{
J ({S}_{w}) = \frac{{P}_{c,\mathit{res}}({S}_{w}) \sqrt{\frac{k}{φ}}}{σ cos Θ}
}\end{equation}

Since the above function is just a normalizing function, then units are
not important, as long as when we de-normalize the average curve we use
the same unit set. Secondly, if all the capillary pressure data has been
converted to reservoir conditions, we actually ignore the denominator as
it is a constant, and we can therefore just use:

\begin{equation}\phantomsection\label{eq-6-6}{
J ({S}_{w}) = {P}_{c,\mathit{res}}({S}_{w}) \sqrt{\frac{k}{φ}}
}\end{equation}

However, in the simulator it is necessary to use the formal definition
as outlined in equation (6.4). In addition to the standard the equation
the keyword allows for de-normalizing the curve to use alternative power
functions instead of the standard 0.5 used in equation (6.4), that is:

\begin{equation}\phantomsection\label{eq-6-7}{
J ({S}_{w}) = \frac{{P}_{c,\mathit{res}}({S}_{w}) \left(\frac{{k}^{\mathrm{β}}}{{φ}^{\mathrm{α}}}\right)}{σ}
}\end{equation}

Where:

J (Sw) = dimensionless function of water saturation

Pc (Sw) = capillary pressure (kPa)

k = permeability, (m2)

φ = porosity (fraction)

σ = interfacial tension (mN/m)

Θ = contact angle

α = porosity power value

β = permeability value

The JFUNC keyword allows the data entered as capillary pressure in the
saturation tables, for example, by using the SGFN and SWFN keywords in
the PROPS section to be treated as J-functions instead, and to
de-normalize these curves for each active cell in the model using the
options and values defined with the JFUNC keyword combined with a cells
porosity and permeability values.

\begin{tcolorbox}[enhanced jigsaw, title=\textcolor{quarto-callout-note-color}{\faInfo}\hspace{0.5em}{Note}, coltitle=black, breakable, toprule=.15mm, colbacktitle=quarto-callout-note-color!10!white, leftrule=.75mm, opacitybacktitle=0.6, colback=white, bottomtitle=1mm, toptitle=1mm, titlerule=0mm, left=2mm, colframe=quarto-callout-note-color-frame, arc=.35mm, rightrule=.15mm, bottomrule=.15mm, opacityback=0]

If either the JFUNC or JFUNCR keywords are used to activate J-Function
scaling then the ENDSCALE keyword in the RUNSPEC section must also be
present in the input deck, in order for the dimensionless J-function
values entered on the SWFN, SGFN or the SWOF, SGOF, SLGOF keywords to be
re-scaled to capillary pressure data. Note if the ENDSCALE keyword is
absent, then like the commercial simulator, J-Function scaling is not
performed, and the values entered on the SWFN, SGFN or the SWOF, SGOF,
SLGOF keywords are used as entered.

\end{tcolorbox}

See also the JFUNCR keyword in the GRID section that performs similar
calculations based on the J-Function parameters being declared by
saturation table number.

\subsubsection{Example}\label{example-174}

\begin{verbatim}
--
--       DEFINE LEVERETT J-FUNCTION PARAMETERS
--       JFUN    OILWAT  GASOIL  PORO    PERM    PERM
--       OPTN    SDENS   SDEN    ALPHA   BETA    OPTN
JFUNC
         WATER   22.5    1*      0.5     0.5     XY                            /
\end{verbatim}

The above example results in the oil-water capillary pressure data
entered on the SWFN keyword in the PROPS section being treated as
J-Functions, and that the J-Function should be de-normalized using an
oil-water surface density of 22.5 dynes/cm, using the default power
values and the average of the PERMX and PERMY values for each grid
block.

\subsection{JFUNCR -- Activate the Leverett J-function Saturation Table
Option}\label{jfuncr-activate-the-leverett-j-function-saturation-table-option}

\begin{longtable}[]{@{}
  >{\raggedright\arraybackslash}p{(\linewidth - 14\tabcolsep) * \real{0.1250}}
  >{\raggedright\arraybackslash}p{(\linewidth - 14\tabcolsep) * \real{0.1250}}
  >{\raggedright\arraybackslash}p{(\linewidth - 14\tabcolsep) * \real{0.1250}}
  >{\raggedright\arraybackslash}p{(\linewidth - 14\tabcolsep) * \real{0.1250}}
  >{\raggedright\arraybackslash}p{(\linewidth - 14\tabcolsep) * \real{0.1250}}
  >{\raggedright\arraybackslash}p{(\linewidth - 14\tabcolsep) * \real{0.1250}}
  >{\raggedright\arraybackslash}p{(\linewidth - 14\tabcolsep) * \real{0.1250}}
  >{\raggedright\arraybackslash}p{(\linewidth - 14\tabcolsep) * \real{0.1250}}@{}}
\toprule\noalign{}
\begin{minipage}[b]{\linewidth}\raggedright
RUNSPEC
\end{minipage} & \begin{minipage}[b]{\linewidth}\raggedright
GRID
\end{minipage} & \begin{minipage}[b]{\linewidth}\raggedright
EDIT
\end{minipage} & \begin{minipage}[b]{\linewidth}\raggedright
PROPS
\end{minipage} & \begin{minipage}[b]{\linewidth}\raggedright
REGIONS
\end{minipage} & \begin{minipage}[b]{\linewidth}\raggedright
SOLUTION
\end{minipage} & \begin{minipage}[b]{\linewidth}\raggedright
SUMMARY
\end{minipage} & \begin{minipage}[b]{\linewidth}\raggedright
SCHEDULE
\end{minipage} \\
\midrule\noalign{}
\endhead
\bottomrule\noalign{}
\endlastfoot
\end{longtable}

\subsubsection{Description}\label{description-265}

JFUNCR keyword activates Leverett-J-Function {[}Leverett, M. C.;
``Capillary Behaviour in Porous Solids'', Trans. AIME (1941) 142,
152-168.{]} Saturation Table option which is a commonly used technique
to normalize capillary pressure base on laboratory measured core plugs
porosity and permeability values and the resulting capillary pressure
data. This keyword is an extension of the JFUNC keyword in the GRID
section that uses the parameters on the JFUNC keyword combined with a
cell's porosity and permeability to perform the scaling globally. In
comparison, the JFUNCR allows for the J-Function parameters to be
declared per saturation table number, resulting in greater flexibility.

The keyword should only be used if end-point scaling is switched on
using the ENDSCALE keyword in the RUNSPEC section.

This keyword is not supported by OPM Flow but would change the results
if supported so the simulation will be stopped.

\begin{longtable}[]{@{}
  >{\raggedright\arraybackslash}p{(\linewidth - 6\tabcolsep) * \real{0.2500}}
  >{\raggedright\arraybackslash}p{(\linewidth - 6\tabcolsep) * \real{0.2500}}
  >{\raggedright\arraybackslash}p{(\linewidth - 6\tabcolsep) * \real{0.2500}}
  >{\raggedright\arraybackslash}p{(\linewidth - 6\tabcolsep) * \real{0.2500}}@{}}
\toprule\noalign{}
\begin{minipage}[b]{\linewidth}\raggedright
No.
\end{minipage} & \begin{minipage}[b]{\linewidth}\raggedright
Name
\end{minipage} & \begin{minipage}[b]{\linewidth}\raggedright
Description
\end{minipage} & \begin{minipage}[b]{\linewidth}\raggedright
Default
\end{minipage} \\
\midrule\noalign{}
\endhead
\bottomrule\noalign{}
\endlastfoot
Field & Metric & Laboratory & \\
1 & JFOPT & A character string that defines which capillary data sets
the J-Function option should be applied to, based on the following
options: & BOTH \\
2 & OWSTEN & A positive real number that defines oil-water surface
tension used to de-normalized J-Function data entered in the PROPS
section. & None \\
dynes/cm & dynes/cm & dynes/cm & \\
3 & OGSTEN & A positive real number that defines oil-gas surface tension
used to de-normalized J-Function data entered in the PROPS section. &
None \\
dynes/cm & dynes/cm & dynes/cm & \\
4 & ALPHA & A positive real value that defines an alternative power
value for the porosity term in the J-Function equation, that is instead
of\(\sqrt{\frac{k}{\mathrm{φ}}}\) use
\(\frac{{k}^{0.5}}{{\mathrm{φ}}^{\mathrm{α}}}\)instead in the
transformation. & 0.5 \\
5 & BETA & A positive real number that defines an alternative power
value for the permeability term in the J-Function equation, that is
instead of \(\sqrt{\frac{k}{\mathrm{φ}}}\)use
\(\frac{{k}^{\mathrm{β}}}{{\mathrm{φ}}^{0.5}}\)instead in the
transformation. & 0.5 \\
6 & PERM & PERM is a character string that sets the permeability array
to be used in the transform, based on the following options: & XY \\
Notes: & & & \\
\end{longtable}

\emph{Table 6.54: JFUNCR Keyword Description}

Just like the relative permeability data capillary pressure data are
measured on core plugs with varying quality and perhaps from different
reservoirs. It is therefore necessary to determine averaged data, before
employing the data in engineering calculations. This is commonly done by
using the Leverett J-function {[}Leverett, M. C.; ``Capillary Behaviour
in Porous Solids'', Trans. AIME (1941) 142, 152-168.{]}, which is
defined as:

\begin{equation}\phantomsection\label{eq-6-8}{
J ({S}_{w}) = \frac{{P}_{c,\mathit{res}}({S}_{w}) \sqrt{\frac{k}{φ}}}{σ}
}\end{equation}

Where:

J (Sw) = dimensionless function of water saturation

Pc (Sw) = capillary pressure (kPa)

k = permeability, (m2)

φ = porosity (fraction)

σ = interfacial tension (mN/m)

Θ = contact angle

Sometimes the equation is stated with the cos θ term included, that is:

\begin{equation}\phantomsection\label{eq-6-9}{
J ({S}_{w}) = \frac{{P}_{c,\mathit{res}}({S}_{w}) \sqrt{\frac{k}{φ}}}{σ cos Θ}
}\end{equation}

Since the above function is just a normalizing function, then units are
not important, as long as when we de-normalize the average curve we use
the same unit set. Secondly, if all the capillary pressure data has been
converted to reservoir conditions, we can actually ignore the
denominator as it is a constant, and we can therefore just use:

\begin{equation}\phantomsection\label{eq-6-10}{
J ({S}_{w}) = {P}_{c,\mathit{res}}({S}_{w}) \sqrt{\frac{k}{φ}}
}\end{equation}

However, in the simulator it is necessary to use the formal definition
as outlined in equation (6.8). In addition to the standard the equation
the keyword allows for de-normalizing the curve to use alternative power
functions instead of the standard 0.5 used in equation (6.8), that is:

\begin{equation}\phantomsection\label{eq-6-11}{
J ({S}_{w}) = \frac{{P}_{c,\mathit{res}}({S}_{w}) \left(\frac{{k}^{\mathrm{β}}}{{φ}^{\mathrm{α}}}\right)}{σ}
}\end{equation}

Where:

J (Sw) = dimensionless function of water saturation

Pc (Sw) = capillary pressure (kPa)

k = permeability, (m2)

φ = porosity (fraction)

σ = interfacial tension (mN/m)

Θ = contact angle

α = porosity power value

β = permeability value

The JFUNC keyword allows the data entered as capillary pressure in the
saturation tables, for example, by using the SGFN and SWFN keywords in
the PROPS section to be treated as J-functions instead, and to
de-normalize these curves for each active cell in the model using the
options and values defined with the JFUNC keyword combined with a cells
porosity and permeability values.

\begin{tcolorbox}[enhanced jigsaw, title=\textcolor{quarto-callout-note-color}{\faInfo}\hspace{0.5em}{Note}, coltitle=black, breakable, toprule=.15mm, colbacktitle=quarto-callout-note-color!10!white, leftrule=.75mm, opacitybacktitle=0.6, colback=white, bottomtitle=1mm, toptitle=1mm, titlerule=0mm, left=2mm, colframe=quarto-callout-note-color-frame, arc=.35mm, rightrule=.15mm, bottomrule=.15mm, opacityback=0]

If either the JFUNC or JFUNCR keywords are used to activate J-Function
scaling then the ENDSCALE keyword in the RUNSPEC section must also be
present in the input deck, in order for the dimensionless J-function
values entered on the SWFN, SGFN or the SWOF, SGOF, SLGOF keywords to be
re-scaled to capillary pressure data. Note if the ENDSCALE keyword is
absent, then like the commercial simulator, J-Function scaling is not
performed, and the values entered on the SWFN, SGFN or the SWOF, SGOF,
SLGOF keywords are used as entered.

\end{tcolorbox}

See also the JFUNC keyword in the GRID section that uses the parameters
on the JFUNC keyword combined with a cell's porosity and permeability to
perform the scaling globally.

\subsubsection{Example}\label{example-175}

The example below assumes NTSFUN is equal to five on the TABDIMS keyword
in the RUNSPEC section.

\begin{verbatim}
--
--       DEFINE LEVERETT J-FUNCTION PARAMETERS BY SATURATION TABLES
--       JFUN    OILWAT  GASOIL  PORO    PERM    PERM
--       OPTN    SDENS   SDEN    ALPHA   BETA    OPTN
JFUNCR
         WATER   22.5    1*      0.5     0.5     XY                            /
         WATER   22.5    1*      0.5     0.5     XY                            /
         WATER   22.5    1*      0.5     0.5     XY                            /
         WATER   22.5    1*      0.5     0.5     XY                            /
         WATER   22.5    1*      0.5     0.5     XY                            /
\end{verbatim}

Here the oil-water capillary pressure data entered on the SWFN keyword
in the PROPS section are treated as J-Functions, and that the J-Function
should be de-normalized using an oil-water surface density of 22.5
dynes/cm, using the default power values and the average of the PERMX
and PERMY values for each grid block, for all five tables. Note that
since all the JFUNCR parameters are the same for all saturation tables
then the JFUNC keyword could be used instead in this instance.

\subsection{LGRCOPY -- Activate Local Grid Refinement
Inheritance}\label{lgrcopy-activate-local-grid-refinement-inheritance-1}

The LGRCOPY keyword actives Local Grid Refinement (``LGR'') Inheritance
option that allows the LGR to inherit the properties of the global or
host cell containing a LGR grid block at the start of the GRID section,
as oppose to the normal process of applying this transform at the end of
the GRID section. LGRCOPY can be used in the RUNSPEC, GRID and EDIT
sections. If used in the RUNSPEC section then the option is applied to
all LGRs defined in the input file, whereas if used in the GRID and EDIT
sections the keyword must be placed inside a LGR definition section,
that is between a CARFIN (Cartesian LGR grid) or RADIN/RADIN4 (radial
LGR grid) and the ENDFIN keyword. In the latter case inheritance is
applied on an individual LGR basis.

See LGRCOPY -- Activate Local Grid Refinement Inheritance in the RUNSPEC
section for a full description.

\subsection{LINKPERM -- Assign Cell Permeabilities to Cell
Faces}\label{linkperm-assign-cell-permeabilities-to-cell-faces}

\begin{longtable}[]{@{}
  >{\raggedright\arraybackslash}p{(\linewidth - 14\tabcolsep) * \real{0.1250}}
  >{\raggedright\arraybackslash}p{(\linewidth - 14\tabcolsep) * \real{0.1250}}
  >{\raggedright\arraybackslash}p{(\linewidth - 14\tabcolsep) * \real{0.1250}}
  >{\raggedright\arraybackslash}p{(\linewidth - 14\tabcolsep) * \real{0.1250}}
  >{\raggedright\arraybackslash}p{(\linewidth - 14\tabcolsep) * \real{0.1250}}
  >{\raggedright\arraybackslash}p{(\linewidth - 14\tabcolsep) * \real{0.1250}}
  >{\raggedright\arraybackslash}p{(\linewidth - 14\tabcolsep) * \real{0.1250}}
  >{\raggedright\arraybackslash}p{(\linewidth - 14\tabcolsep) * \real{0.1250}}@{}}
\toprule\noalign{}
\begin{minipage}[b]{\linewidth}\raggedright
RUNSPEC
\end{minipage} & \begin{minipage}[b]{\linewidth}\raggedright
GRID
\end{minipage} & \begin{minipage}[b]{\linewidth}\raggedright
EDIT
\end{minipage} & \begin{minipage}[b]{\linewidth}\raggedright
PROPS
\end{minipage} & \begin{minipage}[b]{\linewidth}\raggedright
REGIONS
\end{minipage} & \begin{minipage}[b]{\linewidth}\raggedright
SOLUTION
\end{minipage} & \begin{minipage}[b]{\linewidth}\raggedright
SUMMARY
\end{minipage} & \begin{minipage}[b]{\linewidth}\raggedright
SCHEDULE
\end{minipage} \\
\midrule\noalign{}
\endhead
\bottomrule\noalign{}
\endlastfoot
\end{longtable}

\subsubsection{Description}\label{description-266}

The LINKPERM keyword assigns the grid cell permeabilities entered via
the PERMX, PERMY and PERMZ keywords to a cell face (I±, J±, or K±) and
results in the simulator using these values directly in the calculating
the transmissibility between grid blocks. This is different to the
conventional way of entering permeability data that consists of entering
the cell centered permeability and the simulator calculating a weighted
average transmissibility based on the cell centered permeability of the
up-stream and down-stream grid blocks.

This keyword is not supported by OPM Flow but would change the results
if supported so the simulation will be stopped.

\subsection{LTOSIGMA -- Dual Porosity Viscous Displacement Sigma
Parameters}\label{ltosigma-dual-porosity-viscous-displacement-sigma-parameters}

\begin{longtable}[]{@{}
  >{\raggedright\arraybackslash}p{(\linewidth - 14\tabcolsep) * \real{0.1250}}
  >{\raggedright\arraybackslash}p{(\linewidth - 14\tabcolsep) * \real{0.1250}}
  >{\raggedright\arraybackslash}p{(\linewidth - 14\tabcolsep) * \real{0.1250}}
  >{\raggedright\arraybackslash}p{(\linewidth - 14\tabcolsep) * \real{0.1250}}
  >{\raggedright\arraybackslash}p{(\linewidth - 14\tabcolsep) * \real{0.1250}}
  >{\raggedright\arraybackslash}p{(\linewidth - 14\tabcolsep) * \real{0.1250}}
  >{\raggedright\arraybackslash}p{(\linewidth - 14\tabcolsep) * \real{0.1250}}
  >{\raggedright\arraybackslash}p{(\linewidth - 14\tabcolsep) * \real{0.1250}}@{}}
\toprule\noalign{}
\begin{minipage}[b]{\linewidth}\raggedright
RUNSPEC
\end{minipage} & \begin{minipage}[b]{\linewidth}\raggedright
GRID
\end{minipage} & \begin{minipage}[b]{\linewidth}\raggedright
EDIT
\end{minipage} & \begin{minipage}[b]{\linewidth}\raggedright
PROPS
\end{minipage} & \begin{minipage}[b]{\linewidth}\raggedright
REGIONS
\end{minipage} & \begin{minipage}[b]{\linewidth}\raggedright
SOLUTION
\end{minipage} & \begin{minipage}[b]{\linewidth}\raggedright
SUMMARY
\end{minipage} & \begin{minipage}[b]{\linewidth}\raggedright
SCHEDULE
\end{minipage} \\
\midrule\noalign{}
\endhead
\bottomrule\noalign{}
\endlastfoot
\end{longtable}

\subsubsection{Description}\label{description-267}

The LTOSIGMA keyword defines parameters to calculate the sigma factor in
conjunction with the data entered via the LX, LY and LZ keywords in the
GRID section, for when the VISCD keyword has been used in the RUNSPEC
section to activate the Dual Porosity Viscous Displacement option. In
addition, either the DUALPORO or DUALPERM keyword should be entered in
the RUNSPEC section to activate the dual porosity or dual permeability
models.

This keyword is not supported by OPM Flow but would change the results
if supported so the simulation will be stopped.

\subsection{LX -- Dual Porosity Viscous Displacement X Direction Matrix
Size for All
Cells}\label{lx-dual-porosity-viscous-displacement-x-direction-matrix-size-for-all-cells}

\begin{longtable}[]{@{}
  >{\raggedright\arraybackslash}p{(\linewidth - 14\tabcolsep) * \real{0.1250}}
  >{\raggedright\arraybackslash}p{(\linewidth - 14\tabcolsep) * \real{0.1250}}
  >{\raggedright\arraybackslash}p{(\linewidth - 14\tabcolsep) * \real{0.1250}}
  >{\raggedright\arraybackslash}p{(\linewidth - 14\tabcolsep) * \real{0.1250}}
  >{\raggedright\arraybackslash}p{(\linewidth - 14\tabcolsep) * \real{0.1250}}
  >{\raggedright\arraybackslash}p{(\linewidth - 14\tabcolsep) * \real{0.1250}}
  >{\raggedright\arraybackslash}p{(\linewidth - 14\tabcolsep) * \real{0.1250}}
  >{\raggedright\arraybackslash}p{(\linewidth - 14\tabcolsep) * \real{0.1250}}@{}}
\toprule\noalign{}
\begin{minipage}[b]{\linewidth}\raggedright
RUNSPEC
\end{minipage} & \begin{minipage}[b]{\linewidth}\raggedright
GRID
\end{minipage} & \begin{minipage}[b]{\linewidth}\raggedright
EDIT
\end{minipage} & \begin{minipage}[b]{\linewidth}\raggedright
PROPS
\end{minipage} & \begin{minipage}[b]{\linewidth}\raggedright
REGIONS
\end{minipage} & \begin{minipage}[b]{\linewidth}\raggedright
SOLUTION
\end{minipage} & \begin{minipage}[b]{\linewidth}\raggedright
SUMMARY
\end{minipage} & \begin{minipage}[b]{\linewidth}\raggedright
SCHEDULE
\end{minipage} \\
\midrule\noalign{}
\endhead
\bottomrule\noalign{}
\endlastfoot
\end{longtable}

\subsubsection{Description}\label{description-268}

The LX keyword defines the size of ``representative'' matrix grid blocks
in the X direction via an array in dual porosity and dual permeability
runs, for when the VISCD keyword has been used in the RUNSPEC section to
activate the dual porosity viscous displacement option. In addition,
either the DUALPORO or DUALPERM keyword should be entered in the RUNSPEC
section to activate the dual porosity or dual permeability models. The
VISCD option is used to model the viscous displacement of fluids from
the matrix by the fracture pressure gradient, for when the fracture
system has a more moderate permeability, and flow to and from the matrix
caused by the fracture pressure gradient acts as an additional
production mechanism.

This keyword is not supported by OPM Flow but would change the results
if supported so the simulation will be stopped.

\begin{longtable}[]{@{}
  >{\raggedright\arraybackslash}p{(\linewidth - 6\tabcolsep) * \real{0.2500}}
  >{\raggedright\arraybackslash}p{(\linewidth - 6\tabcolsep) * \real{0.2500}}
  >{\raggedright\arraybackslash}p{(\linewidth - 6\tabcolsep) * \real{0.2500}}
  >{\raggedright\arraybackslash}p{(\linewidth - 6\tabcolsep) * \real{0.2500}}@{}}
\toprule\noalign{}
\begin{minipage}[b]{\linewidth}\raggedright
No.
\end{minipage} & \begin{minipage}[b]{\linewidth}\raggedright
Name
\end{minipage} & \begin{minipage}[b]{\linewidth}\raggedright
Description
\end{minipage} & \begin{minipage}[b]{\linewidth}\raggedright
Default
\end{minipage} \\
\midrule\noalign{}
\endhead
\bottomrule\noalign{}
\endlastfoot
Field & Metric & Laboratory & \\
1 & LX & LX is an array of real numbers describing the
``representative'' cell size in the X direction for each cell in the
model. Repeat counts may be used, for example 10*100.0. & 0 \\
feet & m & cm & \\
Notes: & & & \\
\end{longtable}

\emph{Table 6.55: LX Keyword Description}

If the VISCD keyword has been used to activate the Dual Porosity Viscous
Displacement option and LX has not been specified then LX is set to zero
in the calculation of the viscous displacement term. See also the LY, LZ
and LTOSIGMA keywords in the GRID section.

\subsubsection{Example}\label{example-176}

\begin{verbatim}
--
--       DEFINE INPUT BOX FOR EDITING INPUT ARRAYS
--
--       ---------- BOX ---------
--       I1  I2   J1  J2   K1  K2
BOX
         10  10   1   6    1   1                           / DEFINE BOX AREA
--
--       DEFINE DUAL POROSITY VISCOUS DISPLACEMENT X DIRECTION MATRIX SIZE
--
LX
         6*10.0                                            /
--
--       DEFINE END OF INPUT BOX EDITING OF INPUT ARRAYS
--
ENDBOX
\end{verbatim}

The above example defines a subset of the grid and the size of the
``representative'' matrix cells in the X direction to 10.0 ft.; after
which the ENDBOX keyword resets the input to be the full grid.

\subsection{LXFIN -- Define Logarithmic LGR Grid Block Spacing in the
X-Direction}\label{lxfin-define-logarithmic-lgr-grid-block-spacing-in-the-x-direction}

\begin{longtable}[]{@{}
  >{\raggedright\arraybackslash}p{(\linewidth - 14\tabcolsep) * \real{0.1250}}
  >{\raggedright\arraybackslash}p{(\linewidth - 14\tabcolsep) * \real{0.1250}}
  >{\raggedright\arraybackslash}p{(\linewidth - 14\tabcolsep) * \real{0.1250}}
  >{\raggedright\arraybackslash}p{(\linewidth - 14\tabcolsep) * \real{0.1250}}
  >{\raggedright\arraybackslash}p{(\linewidth - 14\tabcolsep) * \real{0.1250}}
  >{\raggedright\arraybackslash}p{(\linewidth - 14\tabcolsep) * \real{0.1250}}
  >{\raggedright\arraybackslash}p{(\linewidth - 14\tabcolsep) * \real{0.1250}}
  >{\raggedright\arraybackslash}p{(\linewidth - 14\tabcolsep) * \real{0.1250}}@{}}
\toprule\noalign{}
\begin{minipage}[b]{\linewidth}\raggedright
RUNSPEC
\end{minipage} & \begin{minipage}[b]{\linewidth}\raggedright
GRID
\end{minipage} & \begin{minipage}[b]{\linewidth}\raggedright
EDIT
\end{minipage} & \begin{minipage}[b]{\linewidth}\raggedright
PROPS
\end{minipage} & \begin{minipage}[b]{\linewidth}\raggedright
REGIONS
\end{minipage} & \begin{minipage}[b]{\linewidth}\raggedright
SOLUTION
\end{minipage} & \begin{minipage}[b]{\linewidth}\raggedright
SUMMARY
\end{minipage} & \begin{minipage}[b]{\linewidth}\raggedright
SCHEDULE
\end{minipage} \\
\midrule\noalign{}
\endhead
\bottomrule\noalign{}
\endlastfoot
\end{longtable}

\subsubsection{Description}\label{description-269}

The LXFIN keyword defines the parameters for automatically generating a
Local Grid Refinement (``LGR'') grid in the X direction based on
logarithmic block spacing, for when the LGR option has been activated by
the LGR keyword in the RUNSPEC section. LXFIN should be placed in
between the CARFIN and ENDFIN keywords in the GRID section.

This keyword is not supported by OPM Flow but would change the results
if supported so the simulation will be stopped.

\subsection{LY -- Dual Porosity Viscous Displacement Y Direction Matrix
Size for All
Cells}\label{ly-dual-porosity-viscous-displacement-y-direction-matrix-size-for-all-cells}

\begin{longtable}[]{@{}
  >{\raggedright\arraybackslash}p{(\linewidth - 14\tabcolsep) * \real{0.1250}}
  >{\raggedright\arraybackslash}p{(\linewidth - 14\tabcolsep) * \real{0.1250}}
  >{\raggedright\arraybackslash}p{(\linewidth - 14\tabcolsep) * \real{0.1250}}
  >{\raggedright\arraybackslash}p{(\linewidth - 14\tabcolsep) * \real{0.1250}}
  >{\raggedright\arraybackslash}p{(\linewidth - 14\tabcolsep) * \real{0.1250}}
  >{\raggedright\arraybackslash}p{(\linewidth - 14\tabcolsep) * \real{0.1250}}
  >{\raggedright\arraybackslash}p{(\linewidth - 14\tabcolsep) * \real{0.1250}}
  >{\raggedright\arraybackslash}p{(\linewidth - 14\tabcolsep) * \real{0.1250}}@{}}
\toprule\noalign{}
\begin{minipage}[b]{\linewidth}\raggedright
RUNSPEC
\end{minipage} & \begin{minipage}[b]{\linewidth}\raggedright
GRID
\end{minipage} & \begin{minipage}[b]{\linewidth}\raggedright
EDIT
\end{minipage} & \begin{minipage}[b]{\linewidth}\raggedright
PROPS
\end{minipage} & \begin{minipage}[b]{\linewidth}\raggedright
REGIONS
\end{minipage} & \begin{minipage}[b]{\linewidth}\raggedright
SOLUTION
\end{minipage} & \begin{minipage}[b]{\linewidth}\raggedright
SUMMARY
\end{minipage} & \begin{minipage}[b]{\linewidth}\raggedright
SCHEDULE
\end{minipage} \\
\midrule\noalign{}
\endhead
\bottomrule\noalign{}
\endlastfoot
\end{longtable}

\subsubsection{Description}\label{description-270}

The LY keyword defines the size of ``representative'' matrix grid blocks
in the Y direction via an array in dual porosity and dual permeability
runs, for when the VISCD keyword has been used in the RUNSPEC section to
activate the dual porosity viscous displacement option. In addition,
either the DUALPORO or DUALPERM keyword should be entered in the RUNSPEC
section to activate the dual porosity or dual permeability models. The
VISCD option is used to model the viscous displacement of fluids from
the matrix by the fracture pressure gradient, for when the fracture
system has a more moderate permeability, and flow to and from the matrix
caused by the fracture pressure gradient acts as an additional
production mechanism.

This keyword is not supported by OPM Flow but would change the results
if supported so the simulation will be stopped.

\begin{longtable}[]{@{}
  >{\raggedright\arraybackslash}p{(\linewidth - 6\tabcolsep) * \real{0.2500}}
  >{\raggedright\arraybackslash}p{(\linewidth - 6\tabcolsep) * \real{0.2500}}
  >{\raggedright\arraybackslash}p{(\linewidth - 6\tabcolsep) * \real{0.2500}}
  >{\raggedright\arraybackslash}p{(\linewidth - 6\tabcolsep) * \real{0.2500}}@{}}
\toprule\noalign{}
\begin{minipage}[b]{\linewidth}\raggedright
No.
\end{minipage} & \begin{minipage}[b]{\linewidth}\raggedright
Name
\end{minipage} & \begin{minipage}[b]{\linewidth}\raggedright
Description
\end{minipage} & \begin{minipage}[b]{\linewidth}\raggedright
Default
\end{minipage} \\
\midrule\noalign{}
\endhead
\bottomrule\noalign{}
\endlastfoot
Field & Metric & Laboratory & \\
1 & LY & LY is an array of real numbers describing the
``representative'' cell size in the Y direction for each cell in the
model. Repeat counts may be used, for example 10*100.0. & 0 \\
feet & m & cm & \\
Notes: & & & \\
\end{longtable}

\emph{Table 6.56: LY Keyword Description}

If the VISCD keyword has been used to activate the Dual Porosity Viscous
Displacement option and LY has not been specified then LY is set to zero
in the calculation of the viscous displacement term.

See also the LX, LZ and LTOSIGMA keywords in the GRID section.

\subsubsection{Example}\label{example-177}

\begin{verbatim}
--
--       DEFINE INPUT BOX FOR EDITING INPUT ARRAYS
--
--       ---------- BOX ---------
--       I1  I2   J1  J2   K1  K2
BOX
         10  10   1   6    1   1                           / DEFINE BOX AREA
--
--       DEFINE DUAL POROSITY VISCOUS DISPLACEMENT Y DIRECTION MATRIX SIZE
--
LY
         6*15.0                                            /
--
--       DEFINE END OF INPUT BOX EDITING OF INPUT ARRAYS
--
ENDBOX
\end{verbatim}

The above example defines a subset of the grid and the size of the
``representative'' matrix cells in the Y direction to 15.0 ft.; after
which the ENDBOX keyword resets the input to be the full grid.

\subsection{LYFIN -- Define Logarithmic LGR Grid Block Spacing in the
Y-Direction}\label{lyfin-define-logarithmic-lgr-grid-block-spacing-in-the-y-direction}

\begin{longtable}[]{@{}
  >{\raggedright\arraybackslash}p{(\linewidth - 14\tabcolsep) * \real{0.1250}}
  >{\raggedright\arraybackslash}p{(\linewidth - 14\tabcolsep) * \real{0.1250}}
  >{\raggedright\arraybackslash}p{(\linewidth - 14\tabcolsep) * \real{0.1250}}
  >{\raggedright\arraybackslash}p{(\linewidth - 14\tabcolsep) * \real{0.1250}}
  >{\raggedright\arraybackslash}p{(\linewidth - 14\tabcolsep) * \real{0.1250}}
  >{\raggedright\arraybackslash}p{(\linewidth - 14\tabcolsep) * \real{0.1250}}
  >{\raggedright\arraybackslash}p{(\linewidth - 14\tabcolsep) * \real{0.1250}}
  >{\raggedright\arraybackslash}p{(\linewidth - 14\tabcolsep) * \real{0.1250}}@{}}
\toprule\noalign{}
\begin{minipage}[b]{\linewidth}\raggedright
RUNSPEC
\end{minipage} & \begin{minipage}[b]{\linewidth}\raggedright
GRID
\end{minipage} & \begin{minipage}[b]{\linewidth}\raggedright
EDIT
\end{minipage} & \begin{minipage}[b]{\linewidth}\raggedright
PROPS
\end{minipage} & \begin{minipage}[b]{\linewidth}\raggedright
REGIONS
\end{minipage} & \begin{minipage}[b]{\linewidth}\raggedright
SOLUTION
\end{minipage} & \begin{minipage}[b]{\linewidth}\raggedright
SUMMARY
\end{minipage} & \begin{minipage}[b]{\linewidth}\raggedright
SCHEDULE
\end{minipage} \\
\midrule\noalign{}
\endhead
\bottomrule\noalign{}
\endlastfoot
\end{longtable}

\subsubsection{Description}\label{description-271}

The LYFIN keyword defines the parameters for automatically generating a
Local Grid Refinement (``LGR'') grid in the Y direction based on
logarithmic block spacing, for when the LGR option has been activated by
the LGR keyword in the RUNSPEC section. LYFIN should be placed in
between the CARFIN and ENDFIN keywords in the GRID section.

This keyword is not supported by OPM Flow but would change the results
if supported so the simulation will be stopped.

\subsection{LZ -- Dual Porosity Viscous Displacement Z Direction Matrix
Size for All
Cells}\label{lz-dual-porosity-viscous-displacement-z-direction-matrix-size-for-all-cells}

\begin{longtable}[]{@{}
  >{\raggedright\arraybackslash}p{(\linewidth - 14\tabcolsep) * \real{0.1250}}
  >{\raggedright\arraybackslash}p{(\linewidth - 14\tabcolsep) * \real{0.1250}}
  >{\raggedright\arraybackslash}p{(\linewidth - 14\tabcolsep) * \real{0.1250}}
  >{\raggedright\arraybackslash}p{(\linewidth - 14\tabcolsep) * \real{0.1250}}
  >{\raggedright\arraybackslash}p{(\linewidth - 14\tabcolsep) * \real{0.1250}}
  >{\raggedright\arraybackslash}p{(\linewidth - 14\tabcolsep) * \real{0.1250}}
  >{\raggedright\arraybackslash}p{(\linewidth - 14\tabcolsep) * \real{0.1250}}
  >{\raggedright\arraybackslash}p{(\linewidth - 14\tabcolsep) * \real{0.1250}}@{}}
\toprule\noalign{}
\begin{minipage}[b]{\linewidth}\raggedright
RUNSPEC
\end{minipage} & \begin{minipage}[b]{\linewidth}\raggedright
GRID
\end{minipage} & \begin{minipage}[b]{\linewidth}\raggedright
EDIT
\end{minipage} & \begin{minipage}[b]{\linewidth}\raggedright
PROPS
\end{minipage} & \begin{minipage}[b]{\linewidth}\raggedright
REGIONS
\end{minipage} & \begin{minipage}[b]{\linewidth}\raggedright
SOLUTION
\end{minipage} & \begin{minipage}[b]{\linewidth}\raggedright
SUMMARY
\end{minipage} & \begin{minipage}[b]{\linewidth}\raggedright
SCHEDULE
\end{minipage} \\
\midrule\noalign{}
\endhead
\bottomrule\noalign{}
\endlastfoot
\end{longtable}

\subsubsection{Description}\label{description-272}

The LZ keyword defines the size of ``representative'' matrix grid blocks
in the Z direction via an array in dual porosity and dual permeability
runs, for when the VISCD keyword has been used in the RUNSPEC section to
activate the dual porosity viscous displacement option. In addition,
either the DUALPORO or DUALPERM keyword should be entered in the RUNSPEC
section to activate the dual porosity or dual permeability models. The
VISCD option is used to model the viscous displacement of fluids from
the matrix by the fracture pressure gradient, for when the fracture
system has a more moderate permeability, and flow to and from the matrix
caused by the fracture pressure gradient acts as an additional
production mechanism.

This keyword is not supported by OPM Flow but would change the results
if supported so the simulation will be stopped.

\begin{longtable}[]{@{}
  >{\raggedright\arraybackslash}p{(\linewidth - 6\tabcolsep) * \real{0.2500}}
  >{\raggedright\arraybackslash}p{(\linewidth - 6\tabcolsep) * \real{0.2500}}
  >{\raggedright\arraybackslash}p{(\linewidth - 6\tabcolsep) * \real{0.2500}}
  >{\raggedright\arraybackslash}p{(\linewidth - 6\tabcolsep) * \real{0.2500}}@{}}
\toprule\noalign{}
\begin{minipage}[b]{\linewidth}\raggedright
No.
\end{minipage} & \begin{minipage}[b]{\linewidth}\raggedright
Name
\end{minipage} & \begin{minipage}[b]{\linewidth}\raggedright
Description
\end{minipage} & \begin{minipage}[b]{\linewidth}\raggedright
Default
\end{minipage} \\
\midrule\noalign{}
\endhead
\bottomrule\noalign{}
\endlastfoot
Field & Metric & Laboratory & \\
1 & LZ & LZ is an array of real numbers describing the
``representative'' cell size in the Z direction for each cell in the
model. Repeat counts may be used, for example 10*100.0. & 0 \\
feet & m & cm & \\
Notes: & & & \\
\end{longtable}

\emph{Table 6.57: LZ Keyword Description}

If the VISCD keyword has been used to activate the Dual Porosity Viscous
Displacement option and LZ has not been specified then LZ is set to zero
in the calculation of the viscous displacement term.

See also the LX, LY and LTOSIGMA keywords in the GRID section.

\subsubsection{Example}\label{example-178}

\begin{verbatim}
--
--       DEFINE INPUT BOX FOR EDITING INPUT ARRAYS
--
--       ---------- BOX ---------
--       I1  I2   J1  J2   K1  K2
BOX
         10  10   1   6    1   1                           / DEFINE BOX AREA
--
--       DEFINE DUAL POROSITY VISCOUS DISPLACEMENT Z DIRECTION MATRIX SIZE
--
LZ
         6*3.0                                             /
--
--       DEFINE END OF INPUT BOX EDITING OF INPUT ARRAYS
--
ENDBOX
\end{verbatim}

The example defines a subset of the grid and the size of the
``representative'' matrix cells in the Y direction to 15.0 ft.; after
which the ENDBOX keyword resets the input to be the full grid.

\subsection{LZFIN -- Define Logarithmic LGR Grid Block Spacing in the
Z-Direction}\label{lzfin-define-logarithmic-lgr-grid-block-spacing-in-the-z-direction}

\begin{longtable}[]{@{}
  >{\raggedright\arraybackslash}p{(\linewidth - 14\tabcolsep) * \real{0.1250}}
  >{\raggedright\arraybackslash}p{(\linewidth - 14\tabcolsep) * \real{0.1250}}
  >{\raggedright\arraybackslash}p{(\linewidth - 14\tabcolsep) * \real{0.1250}}
  >{\raggedright\arraybackslash}p{(\linewidth - 14\tabcolsep) * \real{0.1250}}
  >{\raggedright\arraybackslash}p{(\linewidth - 14\tabcolsep) * \real{0.1250}}
  >{\raggedright\arraybackslash}p{(\linewidth - 14\tabcolsep) * \real{0.1250}}
  >{\raggedright\arraybackslash}p{(\linewidth - 14\tabcolsep) * \real{0.1250}}
  >{\raggedright\arraybackslash}p{(\linewidth - 14\tabcolsep) * \real{0.1250}}@{}}
\toprule\noalign{}
\begin{minipage}[b]{\linewidth}\raggedright
RUNSPEC
\end{minipage} & \begin{minipage}[b]{\linewidth}\raggedright
GRID
\end{minipage} & \begin{minipage}[b]{\linewidth}\raggedright
EDIT
\end{minipage} & \begin{minipage}[b]{\linewidth}\raggedright
PROPS
\end{minipage} & \begin{minipage}[b]{\linewidth}\raggedright
REGIONS
\end{minipage} & \begin{minipage}[b]{\linewidth}\raggedright
SOLUTION
\end{minipage} & \begin{minipage}[b]{\linewidth}\raggedright
SUMMARY
\end{minipage} & \begin{minipage}[b]{\linewidth}\raggedright
SCHEDULE
\end{minipage} \\
\midrule\noalign{}
\endhead
\bottomrule\noalign{}
\endlastfoot
\end{longtable}

\subsubsection{Description}\label{description-273}

The LZFIN keyword defines the parameters for automatically generating a
Local Grid Refinement (``LGR'') grid in the Z direction based on
logarithmic block spacing, for when the LGR option has been activated by
the LGR keyword in the RUNSPEC section. LZFIN should be placed in
between the CARFIN and ENDFIN keywords in the GRID section.

This keyword is not supported by OPM Flow but would change the results
if supported so the simulation will be stopped.

\subsection{MAPAXES -- Define the Map Origin Input
Data}\label{mapaxes-define-the-map-origin-input-data}

\begin{longtable}[]{@{}
  >{\raggedright\arraybackslash}p{(\linewidth - 14\tabcolsep) * \real{0.1250}}
  >{\raggedright\arraybackslash}p{(\linewidth - 14\tabcolsep) * \real{0.1250}}
  >{\raggedright\arraybackslash}p{(\linewidth - 14\tabcolsep) * \real{0.1250}}
  >{\raggedright\arraybackslash}p{(\linewidth - 14\tabcolsep) * \real{0.1250}}
  >{\raggedright\arraybackslash}p{(\linewidth - 14\tabcolsep) * \real{0.1250}}
  >{\raggedright\arraybackslash}p{(\linewidth - 14\tabcolsep) * \real{0.1250}}
  >{\raggedright\arraybackslash}p{(\linewidth - 14\tabcolsep) * \real{0.1250}}
  >{\raggedright\arraybackslash}p{(\linewidth - 14\tabcolsep) * \real{0.1250}}@{}}
\toprule\noalign{}
\begin{minipage}[b]{\linewidth}\raggedright
RUNSPEC
\end{minipage} & \begin{minipage}[b]{\linewidth}\raggedright
GRID
\end{minipage} & \begin{minipage}[b]{\linewidth}\raggedright
EDIT
\end{minipage} & \begin{minipage}[b]{\linewidth}\raggedright
PROPS
\end{minipage} & \begin{minipage}[b]{\linewidth}\raggedright
REGIONS
\end{minipage} & \begin{minipage}[b]{\linewidth}\raggedright
SOLUTION
\end{minipage} & \begin{minipage}[b]{\linewidth}\raggedright
SUMMARY
\end{minipage} & \begin{minipage}[b]{\linewidth}\raggedright
SCHEDULE
\end{minipage} \\
\midrule\noalign{}
\endhead
\bottomrule\noalign{}
\endlastfoot
\end{longtable}

\subsubsection{Description}\label{description-274}

MAPAXES specifies the origin of the map used to create the grid. It is
usually output by pre-processing software when exporting the grid
geometry. The data is not used by OPM Flow intrinsically, but is merely
written to the output EGRID file, as specified by the GRIDFILE keyword,
for the use of post-processing software like OPM ResInsight.

\begin{longtable}[]{@{}
  >{\raggedright\arraybackslash}p{(\linewidth - 6\tabcolsep) * \real{0.2500}}
  >{\raggedright\arraybackslash}p{(\linewidth - 6\tabcolsep) * \real{0.2500}}
  >{\raggedright\arraybackslash}p{(\linewidth - 6\tabcolsep) * \real{0.2500}}
  >{\raggedright\arraybackslash}p{(\linewidth - 6\tabcolsep) * \real{0.2500}}@{}}
\toprule\noalign{}
\begin{minipage}[b]{\linewidth}\raggedright
No.
\end{minipage} & \begin{minipage}[b]{\linewidth}\raggedright
Name
\end{minipage} & \begin{minipage}[b]{\linewidth}\raggedright
Description
\end{minipage} & \begin{minipage}[b]{\linewidth}\raggedright
Default
\end{minipage} \\
\midrule\noalign{}
\endhead
\bottomrule\noalign{}
\endlastfoot
Field (feet) & Metric (metres) & Laboratory (metres) & \\
1 & X1 & X1 is a real number that defines the x co-ordinate of a point
on the y-axis. & None \\
2 & Y1 & Y1 is a real number that defines the y co-ordinate of a point
on the y-axis. & None \\
3 & X2 & X2 is a real number that defines the x co-ordinate of the
origin. & None \\
4 & Y2 & Y2 is a real number that defines the y co-ordinate of the
origin. & None \\
5 & X3 & X3 is a real number that defines the x co-ordinate of a point
on the x-axis. & None \\
6 & Y3 & Y3 is a real number that defines the y co-ordinate of a point
on the x-axis. & None \\
Notes: & & & \\
\end{longtable}

\emph{Table 6.58: MAPAXES Keyword Description}

\subsubsection{Example}\label{example-179}

\begin{verbatim}
--
--        ---------------- MAPAXES -----------------
--        X1      Y1      X2      Y2      X3      Y3
MAPAXES
          0.0     100.0   0.0     0.0     100.0   0.0   /
\end{verbatim}

The above example defines the map axes to be exported to the grid file
for use by post-processing software.

\subsection{MAPUNITS -- Define the Map Axes
Units}\label{mapunits-define-the-map-axes-units}

\begin{longtable}[]{@{}
  >{\raggedright\arraybackslash}p{(\linewidth - 14\tabcolsep) * \real{0.1250}}
  >{\raggedright\arraybackslash}p{(\linewidth - 14\tabcolsep) * \real{0.1250}}
  >{\raggedright\arraybackslash}p{(\linewidth - 14\tabcolsep) * \real{0.1250}}
  >{\raggedright\arraybackslash}p{(\linewidth - 14\tabcolsep) * \real{0.1250}}
  >{\raggedright\arraybackslash}p{(\linewidth - 14\tabcolsep) * \real{0.1250}}
  >{\raggedright\arraybackslash}p{(\linewidth - 14\tabcolsep) * \real{0.1250}}
  >{\raggedright\arraybackslash}p{(\linewidth - 14\tabcolsep) * \real{0.1250}}
  >{\raggedright\arraybackslash}p{(\linewidth - 14\tabcolsep) * \real{0.1250}}@{}}
\toprule\noalign{}
\begin{minipage}[b]{\linewidth}\raggedright
RUNSPEC
\end{minipage} & \begin{minipage}[b]{\linewidth}\raggedright
GRID
\end{minipage} & \begin{minipage}[b]{\linewidth}\raggedright
EDIT
\end{minipage} & \begin{minipage}[b]{\linewidth}\raggedright
PROPS
\end{minipage} & \begin{minipage}[b]{\linewidth}\raggedright
REGIONS
\end{minipage} & \begin{minipage}[b]{\linewidth}\raggedright
SOLUTION
\end{minipage} & \begin{minipage}[b]{\linewidth}\raggedright
SUMMARY
\end{minipage} & \begin{minipage}[b]{\linewidth}\raggedright
SCHEDULE
\end{minipage} \\
\midrule\noalign{}
\endhead
\bottomrule\noalign{}
\endlastfoot
\end{longtable}

\subsubsection{Description}\label{description-275}

The MAPUNITS keyword defines the units of the coordinates stated on the
MAPAXES keyword. It is usually output by pre-processing software when
exporting the grid geometry. The data is not used by OPM Flow
intrinsically, but is merely written to the output EGRID file, as
specified by the GRIDFILE keyword, for the use of post-processing
software like OPM ResInsight.

This keyword is not supported by OPM Flow but has no effect on the
results so it will be ignored.

\begin{longtable}[]{@{}
  >{\raggedright\arraybackslash}p{(\linewidth - 6\tabcolsep) * \real{0.2500}}
  >{\raggedright\arraybackslash}p{(\linewidth - 6\tabcolsep) * \real{0.2500}}
  >{\raggedright\arraybackslash}p{(\linewidth - 6\tabcolsep) * \real{0.2500}}
  >{\raggedright\arraybackslash}p{(\linewidth - 6\tabcolsep) * \real{0.2500}}@{}}
\toprule\noalign{}
\begin{minipage}[b]{\linewidth}\raggedright
No.
\end{minipage} & \begin{minipage}[b]{\linewidth}\raggedright
Name
\end{minipage} & \begin{minipage}[b]{\linewidth}\raggedright
Description
\end{minipage} & \begin{minipage}[b]{\linewidth}\raggedright
Default
\end{minipage} \\
\midrule\noalign{}
\endhead
\bottomrule\noalign{}
\endlastfoot
1 & MAPUNITS & A character string that defines the units of the
coordinates stated on the MAPAXES keyword, and should be set to: &
METRES \\
Notes: & & & \\
\end{longtable}

\emph{Table 6.59: MAPUNITS Keyword Description}

\subsubsection{Example}\label{example-180}

\begin{verbatim}
--
--       SET THE MAP UNITS FOR THE MAPAXES KEYWORD
MAPUNITS
         METRES                                                                /
\end{verbatim}

The above example specifies the units on the MAPAXES to be the default
METRES.

\subsection{MAXVALUE -- Sets a Maximum Value for an Array
Element}\label{maxvalue-sets-a-maximum-value-for-an-array-element}

\begin{longtable}[]{@{}
  >{\raggedright\arraybackslash}p{(\linewidth - 14\tabcolsep) * \real{0.1250}}
  >{\raggedright\arraybackslash}p{(\linewidth - 14\tabcolsep) * \real{0.1250}}
  >{\raggedright\arraybackslash}p{(\linewidth - 14\tabcolsep) * \real{0.1250}}
  >{\raggedright\arraybackslash}p{(\linewidth - 14\tabcolsep) * \real{0.1250}}
  >{\raggedright\arraybackslash}p{(\linewidth - 14\tabcolsep) * \real{0.1250}}
  >{\raggedright\arraybackslash}p{(\linewidth - 14\tabcolsep) * \real{0.1250}}
  >{\raggedright\arraybackslash}p{(\linewidth - 14\tabcolsep) * \real{0.1250}}
  >{\raggedright\arraybackslash}p{(\linewidth - 14\tabcolsep) * \real{0.1250}}@{}}
\toprule\noalign{}
\begin{minipage}[b]{\linewidth}\raggedright
RUNSPEC
\end{minipage} & \begin{minipage}[b]{\linewidth}\raggedright
GRID
\end{minipage} & \begin{minipage}[b]{\linewidth}\raggedright
EDIT
\end{minipage} & \begin{minipage}[b]{\linewidth}\raggedright
PROPS
\end{minipage} & \begin{minipage}[b]{\linewidth}\raggedright
REGIONS
\end{minipage} & \begin{minipage}[b]{\linewidth}\raggedright
SOLUTION
\end{minipage} & \begin{minipage}[b]{\linewidth}\raggedright
SUMMARY
\end{minipage} & \begin{minipage}[b]{\linewidth}\raggedright
SCHEDULE
\end{minipage} \\
\midrule\noalign{}
\endhead
\bottomrule\noalign{}
\endlastfoot
\end{longtable}

\subsubsection{Description}\label{description-276}

The MAXVALUE keyword sets a maximum value for the specified array or
part of an array. The constant can be an integer or real value depending
on the array type; however, the arrays that can be operated on are
dependent on which section the MAXVALUE keyword is being applied in.

\begin{longtable}[]{@{}
  >{\raggedright\arraybackslash}p{(\linewidth - 6\tabcolsep) * \real{0.2500}}
  >{\raggedright\arraybackslash}p{(\linewidth - 6\tabcolsep) * \real{0.2500}}
  >{\raggedright\arraybackslash}p{(\linewidth - 6\tabcolsep) * \real{0.2500}}
  >{\raggedright\arraybackslash}p{(\linewidth - 6\tabcolsep) * \real{0.2500}}@{}}
\toprule\noalign{}
\begin{minipage}[b]{\linewidth}\raggedright
No.
\end{minipage} & \begin{minipage}[b]{\linewidth}\raggedright
Name
\end{minipage} & \begin{minipage}[b]{\linewidth}\raggedright
Description
\end{minipage} & \begin{minipage}[b]{\linewidth}\raggedright
Default
\end{minipage} \\
\midrule\noalign{}
\endhead
\bottomrule\noalign{}
\endlastfoot
1 & ARRAY & A character string of up to eight characters in length that
defines the keyword identifying the array to be modified. & None \\
2 & CONSTANT & CONSTANT is a positive integer or positive real value
that an ARRAY element will be reset to if an element in the defined
input BOX, as defined by items (3) to (8), is greater than CONSTANT.
CONSTANT has in the same units as the ARRAY property. & None \\
3 & I1 & The lower bound of the array in the I-direction to be modified
must be greater than or equal to one and less than or equal to I2 and
NX. & 1 \\
4 & I2 & The upper bound of the array in the I-direction to be modified
must be greater than or equal to I1 and less than or equal to NX & NX \\
5 & J1 & The lower bound of the array in the J-direction to be modified
must be greater than or equal to one and less than or equal to J2 and
NY. & 1 \\
6 & J2 & The upper bound of the array in the J-direction to be modified
must be greater than or equal to J1 and less than or equal to NY. &
NY \\
7 & K1 & The lower bound of the array in the K-direction to be modified
must be greater than or equal to one and less than or equal to K2 and
NZ. & 1 \\
8 & K2 & The upper bound of the array in the K-direction to be modified
must be greater than or equal to K1 and less than or equal to NZ. &
NZ \\
Notes: & & & \\
\end{longtable}

\emph{Table 6.60: MAXVALUE Keyword Description}

Examples of the arrays most commonly operated on in each section are
given in Table 6.61. Cells colored red indicate arrays that are not
supported by OPM Flow operations.

\begin{longtable}[]{@{}lllllll@{}}
\toprule\noalign{}
MAXVALUE Keyword and Variable Options by Section & & & & & & \\
\midrule\noalign{}
\endhead
\bottomrule\noalign{}
\endlastfoot
GRID & EDIT & PROPS & REGIONS & SOLUTION & SUMMARY & SCHEDULE \\
DX & DEPTH & SWL & & & & \\
DY & PORV & SWCR & & & & \\
DZ & TRANX & SWU & & & & \\
PERMX & TRANY & SGL & & & & \\
PERMY & TRANZ & SGCR & & & & \\
PERMZ & DIFFX & SGU & & & & \\
MULTX & DIFFY & KRW & & & & \\
MULTY & DIFFZ & KRO & & & & \\
MULTZ & TRANR & KRG & & & & \\
DR & TRANTHT & PCG & & & & \\
DTHETA & DIFFR & PCW & & & & \\
PERMR & DIFFTHT & & & & & \\
PERMTHT & & & & & & \\
DZNET & & & & & & \\
PORO & & & & & & \\
NTG & & & & & & \\
FLUXNUM & & & & & & \\
MULTNUM & & & & & & \\
MPFANUM & & & & & & \\
DIFFX & & & & & & \\
DIFFY & & & & & & \\
DIFFZ & & & & & & \\
DIFFR & & & & & & \\
DIFFTHT & & & & & & \\
\end{longtable}

\emph{Table 6.61: MAXVALUE Keyword Applicable Arrays by Section}

\subsubsection{Example}\label{example-181}

\begin{verbatim}
--
--       ARRAY     CONSTANT       ---------- BOX ---------
--                                I1  I2   J1  J2   K1  K2
MAXVALUE
         PERMX     1.0E2          1*  1*   1*  1*   1*  1* / MAX VALUE FOR PERMX
         PERMY     1.0E2          1*  1*   1*  1*   1*  1* / MAX VALUE FOR PERMY
         PERMZ     1.0E1          1*  1*   1*  1*   1*  1* / MAX VALUE FOR PERMZ
/
\end{verbatim}

The above example resets the maximum values for the PERMX, PERMY and
PERMZ, arrays to 100.0, 100.0 and 10.0, respectively, for all cells.

\subsection{MESSAGE -- Output User
Message}\label{message-output-user-message-2}

The MESSAGE keyword outputs a user message to the terminal, as well as
to the print (\emph{.PRT) and debug (}.DBG) files. Note this is
different to the MESSAGES keyword, that defines OPM Flows message print
limits and stop limits generated by the simulator.

See MESSAGE -- Output User Message in the GLOBAL section for a full
description.

\subsection{MESSAGES -- Define Message Print Limits and Stop
Limits}\label{messages-define-message-print-limits-and-stop-limits-2}

The MESSAGES keyword defines the print and stops levels for various
messages. The ``print limits'' set the maximum number of messages that
will be printed, after which no more messages will be printed and the
``stop limits'' terminate the run when these limits are exceeded. There
are six levels of message that increase in severity from informative all
the way to programming errors, as outlined in Table 4.5.

See MESSAGES -- Define Message Print Limits and Stop Limits in the
GLOBAL section for a full description.

\subsection{MINNNCT -- Set a Minimum Non-Neighbor Connection
Transmissibility}\label{minnnct-set-a-minimum-non-neighbor-connection-transmissibility}

\begin{longtable}[]{@{}
  >{\raggedright\arraybackslash}p{(\linewidth - 14\tabcolsep) * \real{0.1250}}
  >{\raggedright\arraybackslash}p{(\linewidth - 14\tabcolsep) * \real{0.1250}}
  >{\raggedright\arraybackslash}p{(\linewidth - 14\tabcolsep) * \real{0.1250}}
  >{\raggedright\arraybackslash}p{(\linewidth - 14\tabcolsep) * \real{0.1250}}
  >{\raggedright\arraybackslash}p{(\linewidth - 14\tabcolsep) * \real{0.1250}}
  >{\raggedright\arraybackslash}p{(\linewidth - 14\tabcolsep) * \real{0.1250}}
  >{\raggedright\arraybackslash}p{(\linewidth - 14\tabcolsep) * \real{0.1250}}
  >{\raggedright\arraybackslash}p{(\linewidth - 14\tabcolsep) * \real{0.1250}}@{}}
\toprule\noalign{}
\begin{minipage}[b]{\linewidth}\raggedright
RUNSPEC
\end{minipage} & \begin{minipage}[b]{\linewidth}\raggedright
GRID
\end{minipage} & \begin{minipage}[b]{\linewidth}\raggedright
EDIT
\end{minipage} & \begin{minipage}[b]{\linewidth}\raggedright
PROPS
\end{minipage} & \begin{minipage}[b]{\linewidth}\raggedright
REGIONS
\end{minipage} & \begin{minipage}[b]{\linewidth}\raggedright
SOLUTION
\end{minipage} & \begin{minipage}[b]{\linewidth}\raggedright
SUMMARY
\end{minipage} & \begin{minipage}[b]{\linewidth}\raggedright
SCHEDULE
\end{minipage} \\
\midrule\noalign{}
\endhead
\bottomrule\noalign{}
\endlastfoot
\end{longtable}

\subsubsection{Description}\label{description-277}

The MINNNCT keyword defines a minimum non-neighbor connection
transmissibility below which the non-neighbor connection is deleted. The
keyword allows for three minimum values, one for the transmissibility,
one for the diffusivity and one for the thermal transmissibility. If the
keyword is absent from the input deck then no minimum cut-off is
applied.

This keyword is not supported by OPM Flow but would change the results
if supported so the simulation will be stopped.

\subsection{MINPORV -- Set a Minimum Grid Block Pore Volume Threshold
for All
Cells}\label{minporv-set-a-minimum-grid-block-pore-volume-threshold-for-all-cells}

\begin{longtable}[]{@{}
  >{\raggedright\arraybackslash}p{(\linewidth - 14\tabcolsep) * \real{0.1250}}
  >{\raggedright\arraybackslash}p{(\linewidth - 14\tabcolsep) * \real{0.1250}}
  >{\raggedright\arraybackslash}p{(\linewidth - 14\tabcolsep) * \real{0.1250}}
  >{\raggedright\arraybackslash}p{(\linewidth - 14\tabcolsep) * \real{0.1250}}
  >{\raggedright\arraybackslash}p{(\linewidth - 14\tabcolsep) * \real{0.1250}}
  >{\raggedright\arraybackslash}p{(\linewidth - 14\tabcolsep) * \real{0.1250}}
  >{\raggedright\arraybackslash}p{(\linewidth - 14\tabcolsep) * \real{0.1250}}
  >{\raggedright\arraybackslash}p{(\linewidth - 14\tabcolsep) * \real{0.1250}}@{}}
\toprule\noalign{}
\begin{minipage}[b]{\linewidth}\raggedright
RUNSPEC
\end{minipage} & \begin{minipage}[b]{\linewidth}\raggedright
GRID
\end{minipage} & \begin{minipage}[b]{\linewidth}\raggedright
EDIT
\end{minipage} & \begin{minipage}[b]{\linewidth}\raggedright
PROPS
\end{minipage} & \begin{minipage}[b]{\linewidth}\raggedright
REGIONS
\end{minipage} & \begin{minipage}[b]{\linewidth}\raggedright
SOLUTION
\end{minipage} & \begin{minipage}[b]{\linewidth}\raggedright
SUMMARY
\end{minipage} & \begin{minipage}[b]{\linewidth}\raggedright
SCHEDULE
\end{minipage} \\
\midrule\noalign{}
\endhead
\bottomrule\noalign{}
\endlastfoot
\end{longtable}

\subsubsection{Description}\label{description-278}

MINPORV defines a minimum threshold pore volume that makes all grid
blocks whose pore volume is below this value inactive in the model
(inactive cells are not used in OPM Flow calculations). Note this
keyword is different to the MINPVV keyword in the GRID section, that
sets a minimum threshold pore volume for individual cells in the model.

This keyword is an alias for the MINPV keyword in the GRID section, and
thus provides the same functionality.

\begin{longtable}[]{@{}
  >{\raggedright\arraybackslash}p{(\linewidth - 6\tabcolsep) * \real{0.2500}}
  >{\raggedright\arraybackslash}p{(\linewidth - 6\tabcolsep) * \real{0.2500}}
  >{\raggedright\arraybackslash}p{(\linewidth - 6\tabcolsep) * \real{0.2500}}
  >{\raggedright\arraybackslash}p{(\linewidth - 6\tabcolsep) * \real{0.2500}}@{}}
\toprule\noalign{}
\begin{minipage}[b]{\linewidth}\raggedright
No.
\end{minipage} & \begin{minipage}[b]{\linewidth}\raggedright
Name
\end{minipage} & \begin{minipage}[b]{\linewidth}\raggedright
Description
\end{minipage} & \begin{minipage}[b]{\linewidth}\raggedright
Default
\end{minipage} \\
\midrule\noalign{}
\endhead
\bottomrule\noalign{}
\endlastfoot
Field & Metric & Laboratory & \\
1 & MINPORV & MINPORV is a real positive number that defines the minimum
pore volume for a cell to be active in the model. & Defined \\
rb 1.0e-6 & rm3 1.0e-6 & rcc 1.0e-6 & \\
Notes: & & & \\
\end{longtable}

\emph{Table 6.62: MINPORV Keyword Description}

The MINPORV, MINPV, and MINPVV keywords only apply their minimum
threshold pore volume values to active cells. Thus, cells that have been
made inactive via setting their ACTNUM values to zero, remain inactive,
even if their pore volume exceeds the values set by the MINPORV, MINPV,
and MINPVV keywords.

Secondly, although the MINPORV keyword allows one to set a minimum
threshold pore volume below the default value, this is not recommended,
as cells with small pore volumes can cause significant numerical
convergence errors. Thus, in practice, values greater than the default
values are normally applied to eliminate cells that have relatively
small pore volumes. In addition, the simulator reports the total pore
volume and the number of cells made inactive, as well as the pore volume
reduction, when the keyword is invoked; allowing one to run some
sensitivities to the minimum pore volume value.

See also the PINCH keyword for the treatment of inactive grid cells and
pinch-outs.

\subsubsection{Example}\label{example-182}

\begin{verbatim}
--
--       MINIMUM PORE VOLUME FOR ACTIVE CELLS
--
MINPORV
         500.0                                                                  /
\end{verbatim}

The above example defines 500 rb (or m3) as the minimum pore volume for
a cell to be active in the model.

\subsection{MINPV -- Set a Minimum Grid Block Pore Volume Threshold for
All
Cells}\label{minpv-set-a-minimum-grid-block-pore-volume-threshold-for-all-cells}

\begin{longtable}[]{@{}
  >{\raggedright\arraybackslash}p{(\linewidth - 14\tabcolsep) * \real{0.1250}}
  >{\raggedright\arraybackslash}p{(\linewidth - 14\tabcolsep) * \real{0.1250}}
  >{\raggedright\arraybackslash}p{(\linewidth - 14\tabcolsep) * \real{0.1250}}
  >{\raggedright\arraybackslash}p{(\linewidth - 14\tabcolsep) * \real{0.1250}}
  >{\raggedright\arraybackslash}p{(\linewidth - 14\tabcolsep) * \real{0.1250}}
  >{\raggedright\arraybackslash}p{(\linewidth - 14\tabcolsep) * \real{0.1250}}
  >{\raggedright\arraybackslash}p{(\linewidth - 14\tabcolsep) * \real{0.1250}}
  >{\raggedright\arraybackslash}p{(\linewidth - 14\tabcolsep) * \real{0.1250}}@{}}
\toprule\noalign{}
\begin{minipage}[b]{\linewidth}\raggedright
RUNSPEC
\end{minipage} & \begin{minipage}[b]{\linewidth}\raggedright
GRID
\end{minipage} & \begin{minipage}[b]{\linewidth}\raggedright
EDIT
\end{minipage} & \begin{minipage}[b]{\linewidth}\raggedright
PROPS
\end{minipage} & \begin{minipage}[b]{\linewidth}\raggedright
REGIONS
\end{minipage} & \begin{minipage}[b]{\linewidth}\raggedright
SOLUTION
\end{minipage} & \begin{minipage}[b]{\linewidth}\raggedright
SUMMARY
\end{minipage} & \begin{minipage}[b]{\linewidth}\raggedright
SCHEDULE
\end{minipage} \\
\midrule\noalign{}
\endhead
\bottomrule\noalign{}
\endlastfoot
\end{longtable}

\subsubsection{Description}\label{description-279}

MINPV defines a minimum threshold pore volume that makes all grid blocks
whose pore volume is below this value inactive in the model (inactive
cells are not used in OPM Flow calculations). Note this keyword is
different to the MINPVV keyword in the GRID section, that sets a minimum
threshold pore volume for individual cells in the model.

Note that the MINPORV keyword is an alias for the MINPV keyword in the
GRID section, and thus provides the same functionality.

\begin{longtable}[]{@{}
  >{\raggedright\arraybackslash}p{(\linewidth - 6\tabcolsep) * \real{0.2500}}
  >{\raggedright\arraybackslash}p{(\linewidth - 6\tabcolsep) * \real{0.2500}}
  >{\raggedright\arraybackslash}p{(\linewidth - 6\tabcolsep) * \real{0.2500}}
  >{\raggedright\arraybackslash}p{(\linewidth - 6\tabcolsep) * \real{0.2500}}@{}}
\toprule\noalign{}
\begin{minipage}[b]{\linewidth}\raggedright
No.
\end{minipage} & \begin{minipage}[b]{\linewidth}\raggedright
Name
\end{minipage} & \begin{minipage}[b]{\linewidth}\raggedright
Description
\end{minipage} & \begin{minipage}[b]{\linewidth}\raggedright
Default
\end{minipage} \\
\midrule\noalign{}
\endhead
\bottomrule\noalign{}
\endlastfoot
Field & Metric & Laboratory & \\
1 & MINPV & MINPV is a real positive number that defines the minimum
pore volume for a cell to be active in the model. & Defined \\
rb 1.0e-6 & rm3 1.0e-6 & rcc 1.0e-6 & \\
Notes: & & & \\
\end{longtable}

\emph{Table 6.63: MINPV Keyword Description}

The MINPV, MINPORV, and MINPVV keywords only apply their minimum
threshold pore volume values to active cells. Thus, cells that have been
made inactive via setting their ACTNUM values to zero, remain inactive,
even if their pore volume exceeds the values set by the MINPV, MINPORV,
and MINPVV keywords.

Secondly, although the MINPV keyword allows one to set a minimum
threshold pore volume below the default value, this is not recommended,
as cells with small pore volumes can cause significant numerical
convergence errors. Thus, in practice, values greater than the default
values are normally applied to eliminate cells that have relatively
small pore volumes. In addition, the simulator reports the total pore
volume and the number of cells made inactive, as well as the pore volume
reduction, when the keyword is invoked; allowing one to run some
sensitivities to the minimum pore volume value.

See also the PINCH keyword for the treatment of inactive grid cells and
pinch-outs.

\subsubsection{Example}\label{example-183}

\begin{verbatim}
--
--       MINIMUM PORE VOLUME FOR ACTIVE CELLS
--
MINPV
         500.0                                                                  /
\end{verbatim}

The above example defines 500 rb (or m3) as the minimum pore volume for
a cell to be active in the model.

\subsection{MINPVV -- Set a Minimum Grid Block Pore Volume Threshold for
Individual
Cells}\label{minpvv-set-a-minimum-grid-block-pore-volume-threshold-for-individual-cells}

\begin{longtable}[]{@{}
  >{\raggedright\arraybackslash}p{(\linewidth - 14\tabcolsep) * \real{0.1250}}
  >{\raggedright\arraybackslash}p{(\linewidth - 14\tabcolsep) * \real{0.1250}}
  >{\raggedright\arraybackslash}p{(\linewidth - 14\tabcolsep) * \real{0.1250}}
  >{\raggedright\arraybackslash}p{(\linewidth - 14\tabcolsep) * \real{0.1250}}
  >{\raggedright\arraybackslash}p{(\linewidth - 14\tabcolsep) * \real{0.1250}}
  >{\raggedright\arraybackslash}p{(\linewidth - 14\tabcolsep) * \real{0.1250}}
  >{\raggedright\arraybackslash}p{(\linewidth - 14\tabcolsep) * \real{0.1250}}
  >{\raggedright\arraybackslash}p{(\linewidth - 14\tabcolsep) * \real{0.1250}}@{}}
\toprule\noalign{}
\begin{minipage}[b]{\linewidth}\raggedright
RUNSPEC
\end{minipage} & \begin{minipage}[b]{\linewidth}\raggedright
GRID
\end{minipage} & \begin{minipage}[b]{\linewidth}\raggedright
EDIT
\end{minipage} & \begin{minipage}[b]{\linewidth}\raggedright
PROPS
\end{minipage} & \begin{minipage}[b]{\linewidth}\raggedright
REGIONS
\end{minipage} & \begin{minipage}[b]{\linewidth}\raggedright
SOLUTION
\end{minipage} & \begin{minipage}[b]{\linewidth}\raggedright
SUMMARY
\end{minipage} & \begin{minipage}[b]{\linewidth}\raggedright
SCHEDULE
\end{minipage} \\
\midrule\noalign{}
\endhead
\bottomrule\noalign{}
\endlastfoot
\end{longtable}

\subsubsection{Description}\label{description-280}

MINPVV is an array that defines the minimum threshold pore volume for
each cell, that makes grid blocks whose pore volume is below this value
inactive in the model (inactive cells are not used in OPM Flow
calculations).

Note this keyword is different to the MINPV and MINPORV keywords in the
GRID section, that set a constant minimum threshold pore volume for all
cells in the model.

\begin{longtable}[]{@{}
  >{\raggedright\arraybackslash}p{(\linewidth - 6\tabcolsep) * \real{0.2500}}
  >{\raggedright\arraybackslash}p{(\linewidth - 6\tabcolsep) * \real{0.2500}}
  >{\raggedright\arraybackslash}p{(\linewidth - 6\tabcolsep) * \real{0.2500}}
  >{\raggedright\arraybackslash}p{(\linewidth - 6\tabcolsep) * \real{0.2500}}@{}}
\toprule\noalign{}
\begin{minipage}[b]{\linewidth}\raggedright
No.
\end{minipage} & \begin{minipage}[b]{\linewidth}\raggedright
Name
\end{minipage} & \begin{minipage}[b]{\linewidth}\raggedright
Description
\end{minipage} & \begin{minipage}[b]{\linewidth}\raggedright
Default
\end{minipage} \\
\midrule\noalign{}
\endhead
\bottomrule\noalign{}
\endlastfoot
Field & Metric & Laboratory & \\
1 & MINPVV & MINPVV is an array of real positive numbers that defines
the minimum pore volumes for each cell in the model in order for the
cells to be active. & Defined \\
rb 1.0e-6 & rm3 1.0e-6 & rcc 1.0e-6 & \\
Notes: & & & \\
\end{longtable}

\emph{Table 6.64: MINPVV Keyword Description}

The MINPVV, MINPV, and MINPORV keywords only apply their minimum
threshold pore volume values to active cells. Thus, cells that have been
made inactive via setting their ACTNUM values to zero, remain inactive,
even if their pore volume exceeds the values set by the MINPVV, MINPV,
and MINPORV keywords.

Secondly, although the MINPVV keyword allows one to set a minimum
threshold pore volume below the default value, this is not recommended,
as cells with small pore volumes can cause significant numerical
convergence errors. Thus, in practice, values greater than the default
values are normally applied to eliminate cells that have relatively
small pore volumes. In addition, the simulator reports the total pore
volume and the number of cells made inactive, as well as the pore volume
reduction, when the keyword is invoked; allowing one to run some
sensitivities to the minimum pore volume value.

See also the PINCH keyword for the treatment of inactive grid cells and
pinch-outs.

\subsubsection{Example}\label{example-184}

The example below shows how to define 500 rb (or m3) as the minimum pore
volume for all cells in layer 19 to be active in the model, and 750 rb
(or m3) as the minimum pore volume for all cells in layer 20, by using
the BOX keyword to set the portion of the grid of interest.

\begin{verbatim}
--
--       DEFINE A BOX GRID FOR THE BOTTOM TWO LAYERS OF A 100 X 100 X 20 MODEL
--
--       ---------- BOX ---------
--       I1  I2   J1  J2   K1  K2
BOX
         1*  1*   1*  1*   19  20 / SELECT THE BOTTOM LAYERS
--
--       MINIMUM PORE VOLUME FOR INDIVIDUAL CELLS TO BE ACTIVE
--
MINPVV
         10000*500.0   10000*750.0                                                                          /
--
--       RESET THE INPUT BOX TO BE THE FULL MODEL
--
ENDBOX

\end{verbatim}

Although this will work in the commercial simulators, it does not
currently work in OPM Flow, that is one cannot use the MINPVV keyword in
conjunction with the BOX keyword, as shown in the aforementioned
example.

Instead one can use:

\begin{verbatim}
--
--       MINIMUM PORE VOLUME FOR INDIVIDUAL CELLS TO BE ACTIVE
--
MINPVV
         10000*        10000*     10000*   10000*   10000*
         10000*        10000*     10000*   10000*   10000*
         10000*        10000*     10000*   10000*   10000*
         10000*        10000*     10000*
         10000*500.0   10000*750.0                                             /
\end{verbatim}

To accomplish the same thing, where the 10000* instructs the simulator
to use the default value of 1.0 x 10-6 for 10,000 cells, which in this
case is one layer.

\begin{verbatim}
\end{verbatim}

\subsection{MINVALUE -- Set a Minimum Value for an Array
Element}\label{minvalue-set-a-minimum-value-for-an-array-element}

\begin{longtable}[]{@{}
  >{\raggedright\arraybackslash}p{(\linewidth - 14\tabcolsep) * \real{0.1250}}
  >{\raggedright\arraybackslash}p{(\linewidth - 14\tabcolsep) * \real{0.1250}}
  >{\raggedright\arraybackslash}p{(\linewidth - 14\tabcolsep) * \real{0.1250}}
  >{\raggedright\arraybackslash}p{(\linewidth - 14\tabcolsep) * \real{0.1250}}
  >{\raggedright\arraybackslash}p{(\linewidth - 14\tabcolsep) * \real{0.1250}}
  >{\raggedright\arraybackslash}p{(\linewidth - 14\tabcolsep) * \real{0.1250}}
  >{\raggedright\arraybackslash}p{(\linewidth - 14\tabcolsep) * \real{0.1250}}
  >{\raggedright\arraybackslash}p{(\linewidth - 14\tabcolsep) * \real{0.1250}}@{}}
\toprule\noalign{}
\begin{minipage}[b]{\linewidth}\raggedright
RUNSPEC
\end{minipage} & \begin{minipage}[b]{\linewidth}\raggedright
GRID
\end{minipage} & \begin{minipage}[b]{\linewidth}\raggedright
EDIT
\end{minipage} & \begin{minipage}[b]{\linewidth}\raggedright
PROPS
\end{minipage} & \begin{minipage}[b]{\linewidth}\raggedright
REGIONS
\end{minipage} & \begin{minipage}[b]{\linewidth}\raggedright
SOLUTION
\end{minipage} & \begin{minipage}[b]{\linewidth}\raggedright
SUMMARY
\end{minipage} & \begin{minipage}[b]{\linewidth}\raggedright
SCHEDULE
\end{minipage} \\
\midrule\noalign{}
\endhead
\bottomrule\noalign{}
\endlastfoot
\end{longtable}

\subsubsection{Description}\label{description-281}

The MINVALUE keyword sets a minimum value for the specified array or
part of an array. The constant can be an integer or real value depending
on the array type; however, the arrays that can be operated on are
dependent on which section the MINVALUE keyword is being applied in.

\begin{longtable}[]{@{}
  >{\raggedright\arraybackslash}p{(\linewidth - 6\tabcolsep) * \real{0.2500}}
  >{\raggedright\arraybackslash}p{(\linewidth - 6\tabcolsep) * \real{0.2500}}
  >{\raggedright\arraybackslash}p{(\linewidth - 6\tabcolsep) * \real{0.2500}}
  >{\raggedright\arraybackslash}p{(\linewidth - 6\tabcolsep) * \real{0.2500}}@{}}
\toprule\noalign{}
\begin{minipage}[b]{\linewidth}\raggedright
No.
\end{minipage} & \begin{minipage}[b]{\linewidth}\raggedright
Name
\end{minipage} & \begin{minipage}[b]{\linewidth}\raggedright
Description
\end{minipage} & \begin{minipage}[b]{\linewidth}\raggedright
Default
\end{minipage} \\
\midrule\noalign{}
\endhead
\bottomrule\noalign{}
\endlastfoot
1 & ARRAY & A character string of up to eight characters in length that
defines the keyword identifying the array to be modified. & None \\
2 & CONSTANT & CONSTANT is a positive integer or positive real value
that an ARRAY element will be reset to if an element in the defined
input BOX, as defined by items (3) to (8), is less than CONSTANT.
CONSTANT has in the same units as the ARRAY property. & None \\
3 & I1 & A positive integer that defines the lower bound of the array in
the I-direction to be modified must be greater than or equal to one and
less than or equal to I2 and NX. & 1 \\
4 & I2 & A positive integer that defines the upper bound of the array in
the I-direction to be modified must be greater than or equal to I1 and
less than or equal to NX & NX \\
5 & J1 & A positive integer that defines the lower bound of the array in
the J-direction to be modified must be greater than or equal to one and
less than or equal to J2 and NY. & 1 \\
6 & J2 & A positive integer that defines the upper bound of the array in
the J-direction to be modified must be greater than or equal to J1 and
less than or equal to NY. & NY \\
7 & K1 & A positive integer that defines the lower bound of the array in
the K-direction to be modified must be greater than or equal to one and
less than or equal to K2 and NZ. & 1 \\
8 & K2 & A positive integer that defines the upper bound of the array in
the K-direction to be modified must be greater than or equal to K1 and
less than or equal to NZ. & NZ \\
Notes: & & & \\
\end{longtable}

\emph{Table 6.65: MINVALUE Keyword Description}

Examples of the arrays most commonly operated on in each section are
given in Table 6.66. Cells colored red indicate arrays that are not
supported by OPM Flow operations.

\begin{longtable}[]{@{}lllllll@{}}
\toprule\noalign{}
EQUALS Keyword and Variable Options by Section & & & & & & \\
\midrule\noalign{}
\endhead
\bottomrule\noalign{}
\endlastfoot
GRID & EDIT & PROPS & REGIONS & SOLUTION & SUMMARY & SCHEDULE \\
DX & DEPTH & SWL & & & & \\
DY & PORV & SWCR & & & & \\
DZ & TRANX & SWU & & & & \\
PERMX & TRANY & SGL & & & & \\
PERMY & TRANZ & SGCR & & & & \\
PERMZ & DIFFX & SGU & & & & \\
MULTX & DIFFY & KRW & & & & \\
MULTY & DIFFZ & KRO & & & & \\
MULTZ & TRANR & KRG & & & & \\
DR & TRANTHT & PCG & & & & \\
DTHETA & DIFFR & PCW & & & & \\
PERMR & DIFFTHT & & & & & \\
PERMTHT & & & & & & \\
DZNET & & & & & & \\
PORO & & & & & & \\
NTG & & & & & & \\
FLUXNUM & & & & & & \\
MULTNUM & & & & & & \\
MPFANUM & & & & & & \\
DIFFX & & & & & & \\
DIFFY & & & & & & \\
DIFFZ & & & & & & \\
DIFFR & & & & & & \\
DIFFTHT & & & & & & \\
\end{longtable}

\emph{Table 6.66: MINVALUE Keyword Applicable Arrays by Section}

\subsubsection{Example}\label{example-185}

\begin{verbatim}
--
--       ARRAY    CONSTANT        ---------- BOX ---------
--                                I1  I2   J1  J2   K1  K2
MINVALUE
         PERMX    1.0E0           1*  1*   1*  1*   1*  1* / MINIMUM PERMX
         PERMY    1.0E0           1*  1*   1*  1*   1*  1* / MINIMUM PERMY
         PERMZ    1.0E-1          1*  1*   1*  1*   1*  1* / MINIMUM PERMZ
/
\end{verbatim}

The above example resets the minimum values for the PERMX, PERMY and
PERMZ, arrays to 1.0, 1.0 and 0.1, respectively, for all cells.

\subsection{MPFANUM -- Define Multi-Point Flux Discretization
Regions}\label{mpfanum-define-multi-point-flux-discretization-regions}

\begin{longtable}[]{@{}
  >{\raggedright\arraybackslash}p{(\linewidth - 14\tabcolsep) * \real{0.1250}}
  >{\raggedright\arraybackslash}p{(\linewidth - 14\tabcolsep) * \real{0.1250}}
  >{\raggedright\arraybackslash}p{(\linewidth - 14\tabcolsep) * \real{0.1250}}
  >{\raggedright\arraybackslash}p{(\linewidth - 14\tabcolsep) * \real{0.1250}}
  >{\raggedright\arraybackslash}p{(\linewidth - 14\tabcolsep) * \real{0.1250}}
  >{\raggedright\arraybackslash}p{(\linewidth - 14\tabcolsep) * \real{0.1250}}
  >{\raggedright\arraybackslash}p{(\linewidth - 14\tabcolsep) * \real{0.1250}}
  >{\raggedright\arraybackslash}p{(\linewidth - 14\tabcolsep) * \real{0.1250}}@{}}
\toprule\noalign{}
\begin{minipage}[b]{\linewidth}\raggedright
RUNSPEC
\end{minipage} & \begin{minipage}[b]{\linewidth}\raggedright
GRID
\end{minipage} & \begin{minipage}[b]{\linewidth}\raggedright
EDIT
\end{minipage} & \begin{minipage}[b]{\linewidth}\raggedright
PROPS
\end{minipage} & \begin{minipage}[b]{\linewidth}\raggedright
REGIONS
\end{minipage} & \begin{minipage}[b]{\linewidth}\raggedright
SOLUTION
\end{minipage} & \begin{minipage}[b]{\linewidth}\raggedright
SUMMARY
\end{minipage} & \begin{minipage}[b]{\linewidth}\raggedright
SCHEDULE
\end{minipage} \\
\midrule\noalign{}
\endhead
\bottomrule\noalign{}
\endlastfoot
\end{longtable}

\subsubsection{Description}\label{description-282}

The MPFANUM keyword defines regions in the model where the multi-point
flux discretization should be applied.

This keyword is not supported by OPM Flow but would change the results
if supported so the simulation will be stopped.

\subsection{MPFNNC -- Define Multi-Point Flux Non-Neighbor
Connections}\label{mpfnnc-define-multi-point-flux-non-neighbor-connections}

\begin{longtable}[]{@{}
  >{\raggedright\arraybackslash}p{(\linewidth - 14\tabcolsep) * \real{0.1250}}
  >{\raggedright\arraybackslash}p{(\linewidth - 14\tabcolsep) * \real{0.1250}}
  >{\raggedright\arraybackslash}p{(\linewidth - 14\tabcolsep) * \real{0.1250}}
  >{\raggedright\arraybackslash}p{(\linewidth - 14\tabcolsep) * \real{0.1250}}
  >{\raggedright\arraybackslash}p{(\linewidth - 14\tabcolsep) * \real{0.1250}}
  >{\raggedright\arraybackslash}p{(\linewidth - 14\tabcolsep) * \real{0.1250}}
  >{\raggedright\arraybackslash}p{(\linewidth - 14\tabcolsep) * \real{0.1250}}
  >{\raggedright\arraybackslash}p{(\linewidth - 14\tabcolsep) * \real{0.1250}}@{}}
\toprule\noalign{}
\begin{minipage}[b]{\linewidth}\raggedright
RUNSPEC
\end{minipage} & \begin{minipage}[b]{\linewidth}\raggedright
GRID
\end{minipage} & \begin{minipage}[b]{\linewidth}\raggedright
EDIT
\end{minipage} & \begin{minipage}[b]{\linewidth}\raggedright
PROPS
\end{minipage} & \begin{minipage}[b]{\linewidth}\raggedright
REGIONS
\end{minipage} & \begin{minipage}[b]{\linewidth}\raggedright
SOLUTION
\end{minipage} & \begin{minipage}[b]{\linewidth}\raggedright
SUMMARY
\end{minipage} & \begin{minipage}[b]{\linewidth}\raggedright
SCHEDULE
\end{minipage} \\
\midrule\noalign{}
\endhead
\bottomrule\noalign{}
\endlastfoot
\end{longtable}

\subsubsection{Description}\label{description-283}

The MPFNNC keyword defines multi-point flux non-neighbor connections
explicitly.

This keyword is not supported by OPM Flow but would change the results
if supported so the simulation will be stopped.

\subsection{MULTFLT -- Multiply the Transmissibility of a Defined Fault
by a
Constant}\label{multflt-multiply-the-transmissibility-of-a-defined-fault-by-a-constant}

\begin{longtable}[]{@{}
  >{\raggedright\arraybackslash}p{(\linewidth - 14\tabcolsep) * \real{0.1250}}
  >{\raggedright\arraybackslash}p{(\linewidth - 14\tabcolsep) * \real{0.1250}}
  >{\raggedright\arraybackslash}p{(\linewidth - 14\tabcolsep) * \real{0.1250}}
  >{\raggedright\arraybackslash}p{(\linewidth - 14\tabcolsep) * \real{0.1250}}
  >{\raggedright\arraybackslash}p{(\linewidth - 14\tabcolsep) * \real{0.1250}}
  >{\raggedright\arraybackslash}p{(\linewidth - 14\tabcolsep) * \real{0.1250}}
  >{\raggedright\arraybackslash}p{(\linewidth - 14\tabcolsep) * \real{0.1250}}
  >{\raggedright\arraybackslash}p{(\linewidth - 14\tabcolsep) * \real{0.1250}}@{}}
\toprule\noalign{}
\begin{minipage}[b]{\linewidth}\raggedright
RUNSPEC
\end{minipage} & \begin{minipage}[b]{\linewidth}\raggedright
GRID
\end{minipage} & \begin{minipage}[b]{\linewidth}\raggedright
EDIT
\end{minipage} & \begin{minipage}[b]{\linewidth}\raggedright
PROPS
\end{minipage} & \begin{minipage}[b]{\linewidth}\raggedright
REGIONS
\end{minipage} & \begin{minipage}[b]{\linewidth}\raggedright
SOLUTION
\end{minipage} & \begin{minipage}[b]{\linewidth}\raggedright
SUMMARY
\end{minipage} & \begin{minipage}[b]{\linewidth}\raggedright
SCHEDULE
\end{minipage} \\
\midrule\noalign{}
\endhead
\bottomrule\noalign{}
\endlastfoot
\end{longtable}

\subsubsection{Description}\label{description-284}

MULTFLT enables the transmissibilities across defined faults, as
declared by the FAULTS keyword, to be modified. They keyword allows for
the re-scaling of the existing fault transmissibilities calculated by
OPM Flow, for example setting a fault to be completely sealing by
setting the multiplier to zero.

\begin{longtable}[]{@{}
  >{\raggedright\arraybackslash}p{(\linewidth - 6\tabcolsep) * \real{0.2500}}
  >{\raggedright\arraybackslash}p{(\linewidth - 6\tabcolsep) * \real{0.2500}}
  >{\raggedright\arraybackslash}p{(\linewidth - 6\tabcolsep) * \real{0.2500}}
  >{\raggedright\arraybackslash}p{(\linewidth - 6\tabcolsep) * \real{0.2500}}@{}}
\toprule\noalign{}
\begin{minipage}[b]{\linewidth}\raggedright
No.
\end{minipage} & \begin{minipage}[b]{\linewidth}\raggedright
Name
\end{minipage} & \begin{minipage}[b]{\linewidth}\raggedright
Description
\end{minipage} & \begin{minipage}[b]{\linewidth}\raggedright
Default
\end{minipage} \\
\midrule\noalign{}
\endhead
\bottomrule\noalign{}
\endlastfoot
1 & FLTNAME & FLTNAME is a character string enclosed in quotes with a
maximum length of eight characters, that defines the name of the fault
that FLTMULT will be applied to. FLTNAME must have previously been
defined using the FAULTS keyword in GRID section & None \\
2 & FLT-TRS & A positive real number that sets the transmissible
multiplier to be applied to the FLTNAME transmissibilities positive real
number that sets the transmissible multiplier to be applied to the
FLTNAME transmissibilities. & 1.0 \\
3 & FLT-DIF & A positive real number that sets the diffusivity
multiplier to be applied to the FLTNAME diffusivities. This option
should only be used if the Diffusion option has been made activate by
the DIFFUSE keyword in the RUNSPEC section. OPM Flow does not support
this Diffusion option. & 1.0 \\
Notes: & & & \\
\end{longtable}

\emph{Table 6.67: MULTFLT Keyword Description}

\subsubsection{Example}\label{example-186}

\begin{verbatim}
--
--       MODIFY THE TRANSMISSIBILITES ACROSS DEFINED FAULTS
--
--       FAULT            TRANS           DIFUSS
--       NAME             MULTIPLIER      MULTIPLIER
MULTFLT
         'FAULT01'        0.0                                 / FAULT MULTIPLIERS
         'FAULT02'        0.0                                 / FAULT MULTIPLIERS
         'FAULT03'        0.0                                 / FAULT MULTIPLIERS
/
\end{verbatim}

The above example sets the fault transmissibility multiplier for defined
faults named FAULT01, FAULT02, and FAULT03 to zero making the faults
sealing in the model.

\subsection{MULTIPLY -- Multiply a Specified Array by a
Constant}\label{multiply-multiply-a-specified-array-by-a-constant}

\begin{longtable}[]{@{}
  >{\raggedright\arraybackslash}p{(\linewidth - 14\tabcolsep) * \real{0.1250}}
  >{\raggedright\arraybackslash}p{(\linewidth - 14\tabcolsep) * \real{0.1250}}
  >{\raggedright\arraybackslash}p{(\linewidth - 14\tabcolsep) * \real{0.1250}}
  >{\raggedright\arraybackslash}p{(\linewidth - 14\tabcolsep) * \real{0.1250}}
  >{\raggedright\arraybackslash}p{(\linewidth - 14\tabcolsep) * \real{0.1250}}
  >{\raggedright\arraybackslash}p{(\linewidth - 14\tabcolsep) * \real{0.1250}}
  >{\raggedright\arraybackslash}p{(\linewidth - 14\tabcolsep) * \real{0.1250}}
  >{\raggedright\arraybackslash}p{(\linewidth - 14\tabcolsep) * \real{0.1250}}@{}}
\toprule\noalign{}
\begin{minipage}[b]{\linewidth}\raggedright
RUNSPEC
\end{minipage} & \begin{minipage}[b]{\linewidth}\raggedright
GRID
\end{minipage} & \begin{minipage}[b]{\linewidth}\raggedright
EDIT
\end{minipage} & \begin{minipage}[b]{\linewidth}\raggedright
PROPS
\end{minipage} & \begin{minipage}[b]{\linewidth}\raggedright
REGIONS
\end{minipage} & \begin{minipage}[b]{\linewidth}\raggedright
SOLUTION
\end{minipage} & \begin{minipage}[b]{\linewidth}\raggedright
SUMMARY
\end{minipage} & \begin{minipage}[b]{\linewidth}\raggedright
SCHEDULE
\end{minipage} \\
\midrule\noalign{}
\endhead
\bottomrule\noalign{}
\endlastfoot
\end{longtable}

\subsubsection{Description}\label{description-285}

The MULTIPLY keyword multiplies a specified array or part of an array by
a constant. The constant can be an integer or real value depending on
the array type; however, the arrays that can be operated on are
dependent on which section the MULTIPLY keyword is being applied in.

\begin{longtable}[]{@{}
  >{\raggedright\arraybackslash}p{(\linewidth - 6\tabcolsep) * \real{0.2500}}
  >{\raggedright\arraybackslash}p{(\linewidth - 6\tabcolsep) * \real{0.2500}}
  >{\raggedright\arraybackslash}p{(\linewidth - 6\tabcolsep) * \real{0.2500}}
  >{\raggedright\arraybackslash}p{(\linewidth - 6\tabcolsep) * \real{0.2500}}@{}}
\toprule\noalign{}
\begin{minipage}[b]{\linewidth}\raggedright
No.
\end{minipage} & \begin{minipage}[b]{\linewidth}\raggedright
Name
\end{minipage} & \begin{minipage}[b]{\linewidth}\raggedright
Description
\end{minipage} & \begin{minipage}[b]{\linewidth}\raggedright
Default
\end{minipage} \\
\midrule\noalign{}
\endhead
\bottomrule\noalign{}
\endlastfoot
1 & ARRAY & A character string of up to eight characters in length that
defines the keyword identifying the property to be modified. & None \\
2 & CONSTANT & An integer or real value to multiply the ARRAY by in the
same units as the ARRAY property. & None \\
3 & I1 & A positive integer that defines the lower bound of the array in
the I-direction to be modified must be greater than or equal to one and
less than or equal to I2 and NX. & 1 \\
4 & I2 & A positive integer that defines the upper bound of the array in
the I-direction to be modified must be greater than or equal to II and
less than or equal to NX & NX \\
5 & J1 & A positive integer that defines the lower bound of the array in
the J-direction to be modified must be greater than or equal to one and
less than or equal to J2 and NY. & 1 \\
6 & J2 & A positive integer that defines the upper bound of the array in
the J-direction to be modified must be greater than or equal to JI and
less than or equal to NY. & NY \\
7 & K1 & A positive integer that defines the lower bound of the array in
the K-direction to be modified must be greater than or equal to one and
less than or equal to K2 and NZ. & 1 \\
8 & K2 & A positive integer that defines the upper bound of the array in
the K-direction to be modified must be greater than or equal to KI and
less than or equal to NZ. & NZ \\
Notes: & & & \\
\end{longtable}

\emph{Table 6.68: MULTIPLY Keyword Description}

Examples of the arrays most commonly operated on in each section are
given in Table 6.69. Cells colored red indicate arrays that are not
supported by OPM Flow operations.

\begin{longtable}[]{@{}lllllll@{}}
\toprule\noalign{}
MULTIPLY Keyword and Variable Options by Section & & & & & & \\
\midrule\noalign{}
\endhead
\bottomrule\noalign{}
\endlastfoot
GRID & EDIT & PROPS & REGIONS & SOLUTION & SUMMARY & SCHEDULE \\
DX & DEPTH & SWL & ENDNUM & PRESSURE & & \\
DY & PORV & SWCR & EQLNUM & SWAT & & \\
DZ & TRANX & SWU & FIPNUM & SGAS & & \\
PERMX & TRANY & SGL & IMBNUM & RV & & \\
PERMY & TRANZ & SGCR & MISCNUM & RS & & \\
PERMZ & DIFFX & SGU & PVTNUM & TBLK & & \\
MULTX & DIFFY & KRW & ROCKNUM & GI & & \\
MULTY & DIFFZ & KRO & SATNUM & OILAPI & & \\
MULTZ & TRANR & KRG & WH2NUM & SALT & & \\
DR & TRANTHT & PCG & & GASCONC & & \\
DTHETA & DIFFR & PCW & & SOLVCONC & & \\
PERMR & DIFFTHT & & & SOLVFRAC & & \\
PERMTHT & & & & SFOAM & & \\
DZNET & & & & SPOLY & & \\
PORO & & & & & & \\
NTG & & & & & & \\
FLUXNUM & & & & & & \\
MULTNUM & & & & & & \\
MPFANUM & & & & & & \\
DIFFX & & & & & & \\
DIFFY & & & & & & \\
DIFFZ & & & & & & \\
DIFFR & & & & & & \\
DIFFTHT & & & & & & \\
\end{longtable}

\emph{Table 6.69: MULTIPLY Keyword Applicable Arrays by Section}

\subsubsection{Example}\label{example-187}

\begin{verbatim}
--
--       ARRAY       CONSTANT     ---------- BOX ---------
--                                I1  I2   J1  J2   K1  K2
MULTIPLY
         PERMZ       0.50000      1*  1*   1*  1*   1*  1* / PERMZ * 0.5
/
\end{verbatim}

The above example multiples the PERMZ property array by 0.5 throughout
the model.

\subsection{MULTIREG -- Multiply an Array by a Constant based on a
Region
Number}\label{multireg-multiply-an-array-by-a-constant-based-on-a-region-number}

\begin{longtable}[]{@{}
  >{\raggedright\arraybackslash}p{(\linewidth - 14\tabcolsep) * \real{0.1250}}
  >{\raggedright\arraybackslash}p{(\linewidth - 14\tabcolsep) * \real{0.1250}}
  >{\raggedright\arraybackslash}p{(\linewidth - 14\tabcolsep) * \real{0.1250}}
  >{\raggedright\arraybackslash}p{(\linewidth - 14\tabcolsep) * \real{0.1250}}
  >{\raggedright\arraybackslash}p{(\linewidth - 14\tabcolsep) * \real{0.1250}}
  >{\raggedright\arraybackslash}p{(\linewidth - 14\tabcolsep) * \real{0.1250}}
  >{\raggedright\arraybackslash}p{(\linewidth - 14\tabcolsep) * \real{0.1250}}
  >{\raggedright\arraybackslash}p{(\linewidth - 14\tabcolsep) * \real{0.1250}}@{}}
\toprule\noalign{}
\begin{minipage}[b]{\linewidth}\raggedright
RUNSPEC
\end{minipage} & \begin{minipage}[b]{\linewidth}\raggedright
GRID
\end{minipage} & \begin{minipage}[b]{\linewidth}\raggedright
EDIT
\end{minipage} & \begin{minipage}[b]{\linewidth}\raggedright
PROPS
\end{minipage} & \begin{minipage}[b]{\linewidth}\raggedright
REGIONS
\end{minipage} & \begin{minipage}[b]{\linewidth}\raggedright
SOLUTION
\end{minipage} & \begin{minipage}[b]{\linewidth}\raggedright
SUMMARY
\end{minipage} & \begin{minipage}[b]{\linewidth}\raggedright
SCHEDULE
\end{minipage} \\
\midrule\noalign{}
\endhead
\bottomrule\noalign{}
\endlastfoot
\end{longtable}

\subsubsection{Description}\label{description-286}

The MULTIREG keyword multiplies an array or part of an array by a
constant for cells with a specific region number. The region number
array can be FLUXNUM, MULTNUM or OPERNUM and these arrays must be
defined and be available before the MULTIREG keyword is read by the
simulator. The constant can be an integer or real value depending on the
array type; however, the arrays that can be operated on are dependent on
which section the MULTIREG keyword is being applied in.

\begin{longtable}[]{@{}
  >{\raggedright\arraybackslash}p{(\linewidth - 6\tabcolsep) * \real{0.2500}}
  >{\raggedright\arraybackslash}p{(\linewidth - 6\tabcolsep) * \real{0.2500}}
  >{\raggedright\arraybackslash}p{(\linewidth - 6\tabcolsep) * \real{0.2500}}
  >{\raggedright\arraybackslash}p{(\linewidth - 6\tabcolsep) * \real{0.2500}}@{}}
\toprule\noalign{}
\begin{minipage}[b]{\linewidth}\raggedright
No.
\end{minipage} & \begin{minipage}[b]{\linewidth}\raggedright
Name
\end{minipage} & \begin{minipage}[b]{\linewidth}\raggedright
Description
\end{minipage} & \begin{minipage}[b]{\linewidth}\raggedright
Default
\end{minipage} \\
\midrule\noalign{}
\endhead
\bottomrule\noalign{}
\endlastfoot
1 & ARRAY & A character string of up to eight characters in length that
defines the keyword identifying the array to be modified. & None \\
2 & CONSTANT & An integer or real value to multiply the ARRAY by in the
same units as the ARRAY property for a given REGION. & 0 \\
3 & REGION NUMBER & REGION NUMBER is a positive integer representing the
region for which the CONSTANT in (2) should be applied. & None \\
4 & REGION ARRAY & The REGION ARRAY to use for applying the CONSTANT in
(2) based on the REGION NUMBER in (3). REGION ARRAY can have the
following values: & M \\
Notes: & & & \\
\end{longtable}

\emph{Table 6.70: MULTIREG Keyword Description}

Examples of the arrays most commonly operated on in each section are
given in Table 6.71. Cells colored red indicate arrays that are not
supported by OPM Flow operations.

\begin{longtable}[]{@{}lllllll@{}}
\toprule\noalign{}
MULTREG Keyword and Variable Options by Section & & & & & & \\
\midrule\noalign{}
\endhead
\bottomrule\noalign{}
\endlastfoot
GRID & EDIT & PROPS & REGIONS & SOLUTION & SUMMARY & SCHEDULE \\
DX & DEPTH & SWL & ENDNUM & PRESSURE & & \\
DY & PORV & SWCR & EQLNUM & SWAT & & \\
DZ & TRANX & SWU & FIPNUM & SGAS & & \\
PERMX & TRANY & SGL & IMBNUM & RV & & \\
PERMY & TRANZ & SGCR & MISCNUM & RS & & \\
PERMZ & DIFFX & SGU & PVTNUM & TBLK & & \\
MULTX & DIFFY & KRW & ROCKNUM & GI & & \\
MULTY & DIFFZ & KRO & SATNUM & OILAPI & & \\
MULTZ & TRANR & KRG & WH2NUM & SALT & & \\
DR & TRANTHT & PCG & & GASCONC & & \\
DTHETA & DIFFR & PCW & & SOLVCONC & & \\
PERMR & DIFFTHT & & & SOLVFRAC & & \\
PERMTHT & & & & SFOAM & & \\
DZNET & & & & SPOLY & & \\
PORO & & & & & & \\
NTG & & & & & & \\
FLUXNUM & & & & & & \\
MULTNUM & & & & & & \\
MPFANUM & & & & & & \\
DIFFX & & & & & & \\
DIFFY & & & & & & \\
DIFFZ & & & & & & \\
DIFFR & & & & & & \\
DIFFTHT & & & & & & \\
\end{longtable}

\emph{Table 6.71: MULTIREG Keyword Applicable Arrays by Section}

\subsubsection{Example}\label{example-188}

\begin{verbatim}
--
-- FIRST DEFINE THE PROPERTY ARRAYS AND MULTNUM ARRAYS FOR 10 X 10 X 20 MODEL
--
--  ARRAY     CONSTANT     ---------- BOX ---------
--                          I1  I2   J1  J2   K1  K2
EQUALS
    PORO       0.2000       1*  1*   1*  1*   1*  1*  / PORO  TO 0.20 IN MODEL
    PERMX      100.00       1*  1*   1*  1*   1*  1*  / PERMX TO 0.10 IN MODEL
    MULTNUM    1            1*  1*   1*  1*   1*  1*  / MULTNUM IN MODEL
    MULTNUM    2            1*  5    1   5    6   6   / MULTNUM IN MODEL
    MULTNUM    3            1*  1*   1*  1*   10  10  / MULTNUM IN MODEL
/
--
-- NOW RESET PORO AND PERMX BASED ON THE MULTNUM REGION NUMBER
--
-- MULTIPLY AN ARRAY BY A CONSTANT BASED ON A REGION NUMBER
--
--    ARRAY     CONSTANT  REGION   REGION ARRAY
--              VALUE     NUMBER    M / F / O
MULTIREG
      PORO      1.050     1         M                /
      PORO      1.100     2         M                /
      PORO      0.950     3         M                /
      PERMX      1.25     1         M                /
      PERMX      1.30     2         M                /
      PERMX      0.90     3         M                /
/
\end{verbatim}

The example first defines the PORO and PERMX property arrays for the
model and then sets the MULTNUM array to 1 for all cells in the model,
after which selected areas of model are assigned various MULTNUM integer
values. The MULTIREG can then be invoked to multiply the PORO and PERMX
arrays by a constant for the various MULTNUM regions.

\subsection{MULTNUM -- Define the Multiple Transmissibility
Regions}\label{multnum-define-the-multiple-transmissibility-regions}

\begin{longtable}[]{@{}
  >{\raggedright\arraybackslash}p{(\linewidth - 14\tabcolsep) * \real{0.1250}}
  >{\raggedright\arraybackslash}p{(\linewidth - 14\tabcolsep) * \real{0.1250}}
  >{\raggedright\arraybackslash}p{(\linewidth - 14\tabcolsep) * \real{0.1250}}
  >{\raggedright\arraybackslash}p{(\linewidth - 14\tabcolsep) * \real{0.1250}}
  >{\raggedright\arraybackslash}p{(\linewidth - 14\tabcolsep) * \real{0.1250}}
  >{\raggedright\arraybackslash}p{(\linewidth - 14\tabcolsep) * \real{0.1250}}
  >{\raggedright\arraybackslash}p{(\linewidth - 14\tabcolsep) * \real{0.1250}}
  >{\raggedright\arraybackslash}p{(\linewidth - 14\tabcolsep) * \real{0.1250}}@{}}
\toprule\noalign{}
\begin{minipage}[b]{\linewidth}\raggedright
RUNSPEC
\end{minipage} & \begin{minipage}[b]{\linewidth}\raggedright
GRID
\end{minipage} & \begin{minipage}[b]{\linewidth}\raggedright
EDIT
\end{minipage} & \begin{minipage}[b]{\linewidth}\raggedright
PROPS
\end{minipage} & \begin{minipage}[b]{\linewidth}\raggedright
REGIONS
\end{minipage} & \begin{minipage}[b]{\linewidth}\raggedright
SOLUTION
\end{minipage} & \begin{minipage}[b]{\linewidth}\raggedright
SUMMARY
\end{minipage} & \begin{minipage}[b]{\linewidth}\raggedright
SCHEDULE
\end{minipage} \\
\midrule\noalign{}
\endhead
\bottomrule\noalign{}
\endlastfoot
\end{longtable}

\subsubsection{Description}\label{description-287}

The MULTNUM keyword defines the inter-region transmissibility region
numbers for each grid block, as such there must be one entry for each
cell in the model. The array can be used with the EQUALREG, ADDREG,
COPYREG, MULTIREG, MULTREGP and MULTREGT keywords in calculating various
grid properties in the GRID section.

\begin{longtable}[]{@{}
  >{\raggedright\arraybackslash}p{(\linewidth - 6\tabcolsep) * \real{0.2500}}
  >{\raggedright\arraybackslash}p{(\linewidth - 6\tabcolsep) * \real{0.2500}}
  >{\raggedright\arraybackslash}p{(\linewidth - 6\tabcolsep) * \real{0.2500}}
  >{\raggedright\arraybackslash}p{(\linewidth - 6\tabcolsep) * \real{0.2500}}@{}}
\toprule\noalign{}
\begin{minipage}[b]{\linewidth}\raggedright
No.
\end{minipage} & \begin{minipage}[b]{\linewidth}\raggedright
Name
\end{minipage} & \begin{minipage}[b]{\linewidth}\raggedright
Description
\end{minipage} & \begin{minipage}[b]{\linewidth}\raggedright
Default
\end{minipage} \\
\midrule\noalign{}
\endhead
\bottomrule\noalign{}
\endlastfoot
1 & MULTNUM & MULTNUM defines an array of positive integers assigning a
grid cell to a particular inter-region transmissibility region. The
maximum number of MULTNUM regions is set by the NRMULT variable on the
GRIDOPTS keyword in the RUNSPEC section. & 1 \\
Notes: & & & \\
\end{longtable}

\emph{Table 6.72: MULTNUM Keyword Description}

\subsubsection{Examples}\label{examples-13}

The example below sets three MULTNUM regions for a 4 x 5 x 2 model.

\begin{verbatim}
--
--       DEFINE MULTNUM REGIONS FOR ALL CELLS
--
MULTNUM
         2  2  1  1  2  2  1  1  1  1  1  1  1  1  1  1  1  1  1  1
         3  3  1  1  3  3  1  1  1  1  1  1  1  1  1  1  1  1  1  1
/
\end{verbatim}

Alternatively the EQUALS keyword could be employed to accomplish the
same task, that is:

\begin{verbatim}
--
--       ARRAY       CONSTANT     ---------- BOX ---------
--                                I1  I2   J1  J2   K1  K2
EQUALS
         MULTNUM     1            1*  1*   1*  1*   1*  1* / SET REGION 1
         MULTNUM     2            1   2    1   2    1   1  / SET REGION 2
         MULTNUM     3            1   2    1   2    2   2  / SET REGION 3
/
\end{verbatim}

One can then increase PERMX by 25\% in region three only.

\begin{verbatim}
--
--       MULTIPLY AN ARRAY BY A CONSTANT BASED ON A REGION NUMBER
--
--       ARRAY     CONSTANT  REGION   REGION ARRAY
--                 VALUE     NUMBER    M / F / O
MULTIREG
         PERMX     1.25     3         M                     /
/
\end{verbatim}

\subsection{MULTPV -- Multiply Cell Pore Volumes by a
Constant}\label{multpv-multiply-cell-pore-volumes-by-a-constant}

\begin{longtable}[]{@{}
  >{\raggedright\arraybackslash}p{(\linewidth - 14\tabcolsep) * \real{0.1250}}
  >{\raggedright\arraybackslash}p{(\linewidth - 14\tabcolsep) * \real{0.1250}}
  >{\raggedright\arraybackslash}p{(\linewidth - 14\tabcolsep) * \real{0.1250}}
  >{\raggedright\arraybackslash}p{(\linewidth - 14\tabcolsep) * \real{0.1250}}
  >{\raggedright\arraybackslash}p{(\linewidth - 14\tabcolsep) * \real{0.1250}}
  >{\raggedright\arraybackslash}p{(\linewidth - 14\tabcolsep) * \real{0.1250}}
  >{\raggedright\arraybackslash}p{(\linewidth - 14\tabcolsep) * \real{0.1250}}
  >{\raggedright\arraybackslash}p{(\linewidth - 14\tabcolsep) * \real{0.1250}}@{}}
\toprule\noalign{}
\begin{minipage}[b]{\linewidth}\raggedright
RUNSPEC
\end{minipage} & \begin{minipage}[b]{\linewidth}\raggedright
GRID
\end{minipage} & \begin{minipage}[b]{\linewidth}\raggedright
EDIT
\end{minipage} & \begin{minipage}[b]{\linewidth}\raggedright
PROPS
\end{minipage} & \begin{minipage}[b]{\linewidth}\raggedright
REGIONS
\end{minipage} & \begin{minipage}[b]{\linewidth}\raggedright
SOLUTION
\end{minipage} & \begin{minipage}[b]{\linewidth}\raggedright
SUMMARY
\end{minipage} & \begin{minipage}[b]{\linewidth}\raggedright
SCHEDULE
\end{minipage} \\
\midrule\noalign{}
\endhead
\bottomrule\noalign{}
\endlastfoot
\end{longtable}

\subsubsection{Description}\label{description-288}

MULTPV multiples the pore volumes of a cell by a real positive constant
for all the cells in the model via an array. An alternative to defining
the complete array is to use the BOX keyword to define an area of the
grid and then use the MULTPV keyword to set the multipliers just for the
area defined by the BOX keyword (see the example).

The keyword can be used for all grid types, except for the Radial Grid
geometry.

\begin{longtable}[]{@{}
  >{\raggedright\arraybackslash}p{(\linewidth - 6\tabcolsep) * \real{0.2500}}
  >{\raggedright\arraybackslash}p{(\linewidth - 6\tabcolsep) * \real{0.2500}}
  >{\raggedright\arraybackslash}p{(\linewidth - 6\tabcolsep) * \real{0.2500}}
  >{\raggedright\arraybackslash}p{(\linewidth - 6\tabcolsep) * \real{0.2500}}@{}}
\toprule\noalign{}
\begin{minipage}[b]{\linewidth}\raggedright
No.
\end{minipage} & \begin{minipage}[b]{\linewidth}\raggedright
Name
\end{minipage} & \begin{minipage}[b]{\linewidth}\raggedright
Description
\end{minipage} & \begin{minipage}[b]{\linewidth}\raggedright
Default
\end{minipage} \\
\midrule\noalign{}
\endhead
\bottomrule\noalign{}
\endlastfoot
1 & MULTPV & MULTPV is an array of real positive numbers assigning the
pore volume multipliers for each cell in the model. Repeat counts may be
used, for example 20*100.0. & 1.0 \\
Notes: & & & \\
\end{longtable}

\emph{Table 6.73: MULTPV Keyword Description}

See also the MULTREGP keyword for scaling the cell pore volumes by
region numbers.

\subsubsection{Example}\label{example-189}

\begin{verbatim}
--
--       DEFINE INPUT BOX FOR EDITING INPUT ARRAYS
--
--       ---------- BOX ---------
--       I1  I2   J1  J2   K1  K2
BOX
         10  10   1   6    1   3                           / DEFINE BOX AREA
--
--       SET PORE VOLUME MULTIPLIERS
--
MULTPV
         18*0.0500                                         /
--
--       DEFINE END OF INPUT BOX EDITING OF INPUT ARRAYS
--
ENDBOX
\end{verbatim}

The above example defines a 0.05 scaling multiplier for the 18 cells
defined by the preceding BOX statement. The ENDBOX keyword resets the
input box to the full grid.

\subsection{MULTR- -- Multiply Cell Transmissibility in the -R
Direction}\label{multr--multiply-cell-transmissibility-in-the--r-direction}

\begin{longtable}[]{@{}
  >{\raggedright\arraybackslash}p{(\linewidth - 14\tabcolsep) * \real{0.1250}}
  >{\raggedright\arraybackslash}p{(\linewidth - 14\tabcolsep) * \real{0.1250}}
  >{\raggedright\arraybackslash}p{(\linewidth - 14\tabcolsep) * \real{0.1250}}
  >{\raggedright\arraybackslash}p{(\linewidth - 14\tabcolsep) * \real{0.1250}}
  >{\raggedright\arraybackslash}p{(\linewidth - 14\tabcolsep) * \real{0.1250}}
  >{\raggedright\arraybackslash}p{(\linewidth - 14\tabcolsep) * \real{0.1250}}
  >{\raggedright\arraybackslash}p{(\linewidth - 14\tabcolsep) * \real{0.1250}}
  >{\raggedright\arraybackslash}p{(\linewidth - 14\tabcolsep) * \real{0.1250}}@{}}
\toprule\noalign{}
\begin{minipage}[b]{\linewidth}\raggedright
RUNSPEC
\end{minipage} & \begin{minipage}[b]{\linewidth}\raggedright
GRID
\end{minipage} & \begin{minipage}[b]{\linewidth}\raggedright
EDIT
\end{minipage} & \begin{minipage}[b]{\linewidth}\raggedright
PROPS
\end{minipage} & \begin{minipage}[b]{\linewidth}\raggedright
REGIONS
\end{minipage} & \begin{minipage}[b]{\linewidth}\raggedright
SOLUTION
\end{minipage} & \begin{minipage}[b]{\linewidth}\raggedright
SUMMARY
\end{minipage} & \begin{minipage}[b]{\linewidth}\raggedright
SCHEDULE
\end{minipage} \\
\midrule\noalign{}
\endhead
\bottomrule\noalign{}
\endlastfoot
\end{longtable}

\subsubsection{Description}\label{description-289}

MULTR- multiples the transmissibility between two cell faces in the -R
direction for all the cells in the model via an array, that is the
keyword sets the transmissibility multiplier of block (I, J, K) between
the cells (I-1, J, K) and (I, J, K). An alternative to defining the
complete array is to use the BOX keyword to define an area of the grid
and then use the MULTR- keyword to set the multipliers just for the area
defined by the BOX keyword (see the example).

The keyword should only be used with radial and spider grids, as
declared by the RADIAL or SPIDER keywords in the RUNSPEC section.

\begin{longtable}[]{@{}
  >{\raggedright\arraybackslash}p{(\linewidth - 6\tabcolsep) * \real{0.2500}}
  >{\raggedright\arraybackslash}p{(\linewidth - 6\tabcolsep) * \real{0.2500}}
  >{\raggedright\arraybackslash}p{(\linewidth - 6\tabcolsep) * \real{0.2500}}
  >{\raggedright\arraybackslash}p{(\linewidth - 6\tabcolsep) * \real{0.2500}}@{}}
\toprule\noalign{}
\begin{minipage}[b]{\linewidth}\raggedright
No.
\end{minipage} & \begin{minipage}[b]{\linewidth}\raggedright
Name
\end{minipage} & \begin{minipage}[b]{\linewidth}\raggedright
Description
\end{minipage} & \begin{minipage}[b]{\linewidth}\raggedright
Default
\end{minipage} \\
\midrule\noalign{}
\endhead
\bottomrule\noalign{}
\endlastfoot
1 & MULTR- & MULTR- is an array of real positive numbers assigning the
transmissibility multipliers in the -R direction to each cell face in
the model. Repeat counts may be used, for example 20*100.0. & 1.0 \\
Notes: & & & \\
\end{longtable}

\emph{Table 6.75: MULTR- Keyword Description}

Note that OPM Flow does not require the GRIDOPTS(TRANMULT) parameter in
the RUNSPEC section to be set to YES, in order to use this and other
negative directional dependent multiplier keywords in the input deck.
Whereas, the commercial simulator will terminate with an error if the
keyword is present, and the GRIDOPTS(TRANMULT) parameter has not been
set to YES.

See also the MULTR, MULTTHT, MULTTHT-, MULTZ and MULTZ- keywords for
scaling transmissible between cells in the R direction.

\subsubsection{Example}\label{example-190}

\begin{verbatim}
--
--       DEFINE INPUT BOX FOR EDITING INPUT ARRAYS
--
--       ---------- BOX ---------
--       I1  I2   J1  J2   K1  K2
BOX
         10  10   1   6    1   1                           / DEFINE BOX AREA
--
--       SET MULTR- TRANSMISSIBILITY MULTIPLIERS
--
MULTR-
         6*0.500                                           /
--
--       DEFINE END OF INPUT BOX EDITING OF INPUT ARRAYS
--
ENDBOX
\end{verbatim}

The above example defines a 0.5 scaling multiplier for the six cells
defined by the preceding BOX statement. The ENDBOX keyword resets the
input box to the full grid.

\subsection{MULTR -- Multiply Cell Transmissibility in the +R
Direction}\label{multr-multiply-cell-transmissibility-in-the-r-direction}

\begin{longtable}[]{@{}
  >{\raggedright\arraybackslash}p{(\linewidth - 14\tabcolsep) * \real{0.1250}}
  >{\raggedright\arraybackslash}p{(\linewidth - 14\tabcolsep) * \real{0.1250}}
  >{\raggedright\arraybackslash}p{(\linewidth - 14\tabcolsep) * \real{0.1250}}
  >{\raggedright\arraybackslash}p{(\linewidth - 14\tabcolsep) * \real{0.1250}}
  >{\raggedright\arraybackslash}p{(\linewidth - 14\tabcolsep) * \real{0.1250}}
  >{\raggedright\arraybackslash}p{(\linewidth - 14\tabcolsep) * \real{0.1250}}
  >{\raggedright\arraybackslash}p{(\linewidth - 14\tabcolsep) * \real{0.1250}}
  >{\raggedright\arraybackslash}p{(\linewidth - 14\tabcolsep) * \real{0.1250}}@{}}
\toprule\noalign{}
\begin{minipage}[b]{\linewidth}\raggedright
RUNSPEC
\end{minipage} & \begin{minipage}[b]{\linewidth}\raggedright
GRID
\end{minipage} & \begin{minipage}[b]{\linewidth}\raggedright
EDIT
\end{minipage} & \begin{minipage}[b]{\linewidth}\raggedright
PROPS
\end{minipage} & \begin{minipage}[b]{\linewidth}\raggedright
REGIONS
\end{minipage} & \begin{minipage}[b]{\linewidth}\raggedright
SOLUTION
\end{minipage} & \begin{minipage}[b]{\linewidth}\raggedright
SUMMARY
\end{minipage} & \begin{minipage}[b]{\linewidth}\raggedright
SCHEDULE
\end{minipage} \\
\midrule\noalign{}
\endhead
\bottomrule\noalign{}
\endlastfoot
\end{longtable}

\subsubsection{Description}\label{description-290}

MULTR multiples the transmissibility between two cell faces in the +R
direction for all the cells in the model via an array, that is the
keyword sets the transmissibility multiplier of block (I, J, K) between
the cells (I, J, K) and (I+I, J, K). An alternative to defining the
complete array is to use the BOX keyword to define an area of the grid
and then use the MULTR keyword to set the multipliers just for the area
defined by the BOX keyword (see the example).

The keyword should only be used with radial and spider grids, as
declared by the RADIAL or SPIDER keywords in the RUNSPEC section.

\begin{longtable}[]{@{}
  >{\raggedright\arraybackslash}p{(\linewidth - 6\tabcolsep) * \real{0.2500}}
  >{\raggedright\arraybackslash}p{(\linewidth - 6\tabcolsep) * \real{0.2500}}
  >{\raggedright\arraybackslash}p{(\linewidth - 6\tabcolsep) * \real{0.2500}}
  >{\raggedright\arraybackslash}p{(\linewidth - 6\tabcolsep) * \real{0.2500}}@{}}
\toprule\noalign{}
\begin{minipage}[b]{\linewidth}\raggedright
No.
\end{minipage} & \begin{minipage}[b]{\linewidth}\raggedright
Name
\end{minipage} & \begin{minipage}[b]{\linewidth}\raggedright
Description
\end{minipage} & \begin{minipage}[b]{\linewidth}\raggedright
Default
\end{minipage} \\
\midrule\noalign{}
\endhead
\bottomrule\noalign{}
\endlastfoot
1 & MULTR+ & MULTR+ is an array of real positive numbers assigning the
transmissibility multipliers in the +R direction to each cell face in
the model. Repeat counts may be used, for example 20*100.0. & 1.0 \\
Notes: & & & \\
\end{longtable}

\emph{Table 6.74: MULTR Keyword Description}

See also the MULTR-, MULTTHT, MULTTHT-, MULTZ and MULTZ- keywords for
scaling transmissible between between grid cells in the R direction.

\subsubsection{Example}\label{example-191}

\begin{verbatim}
--
--       DEFINE INPUT BOX FOR EDITING INPUT ARRAYS
--
--       ---------- BOX ---------
--       I1  I2   J1  J2   K1  K2
BOX
         10  10   1   6    1   3                           / DEFINE BOX AREA
--
--       SET MULTR+ TRANSMISSIBILITY MULTIPLIERS
--
MULTR
         18*0.300                                          /
--
--       DEFINE END OF INPUT BOX EDITING OF INPUT ARRAYS
--
ENDBOX
\end{verbatim}

The above example defines a 0.3 scaling multiplier for the 18 cells
defined by the preceding BOX statement. The ENDBOX keyword resets the
input box to the full grid.

\subsection{MULTREGD -- Multiply Diffusivities Between
Regions}\label{multregd-multiply-diffusivities-between-regions}

\begin{longtable}[]{@{}
  >{\raggedright\arraybackslash}p{(\linewidth - 14\tabcolsep) * \real{0.1250}}
  >{\raggedright\arraybackslash}p{(\linewidth - 14\tabcolsep) * \real{0.1250}}
  >{\raggedright\arraybackslash}p{(\linewidth - 14\tabcolsep) * \real{0.1250}}
  >{\raggedright\arraybackslash}p{(\linewidth - 14\tabcolsep) * \real{0.1250}}
  >{\raggedright\arraybackslash}p{(\linewidth - 14\tabcolsep) * \real{0.1250}}
  >{\raggedright\arraybackslash}p{(\linewidth - 14\tabcolsep) * \real{0.1250}}
  >{\raggedright\arraybackslash}p{(\linewidth - 14\tabcolsep) * \real{0.1250}}
  >{\raggedright\arraybackslash}p{(\linewidth - 14\tabcolsep) * \real{0.1250}}@{}}
\toprule\noalign{}
\begin{minipage}[b]{\linewidth}\raggedright
RUNSPEC
\end{minipage} & \begin{minipage}[b]{\linewidth}\raggedright
GRID
\end{minipage} & \begin{minipage}[b]{\linewidth}\raggedright
EDIT
\end{minipage} & \begin{minipage}[b]{\linewidth}\raggedright
PROPS
\end{minipage} & \begin{minipage}[b]{\linewidth}\raggedright
REGIONS
\end{minipage} & \begin{minipage}[b]{\linewidth}\raggedright
SOLUTION
\end{minipage} & \begin{minipage}[b]{\linewidth}\raggedright
SUMMARY
\end{minipage} & \begin{minipage}[b]{\linewidth}\raggedright
SCHEDULE
\end{minipage} \\
\midrule\noalign{}
\endhead
\bottomrule\noalign{}
\endlastfoot
\end{longtable}

\subsubsection{Description}\label{description-291}

The MULTREGD keyword multiplies the diffusivity between two regions by a
constant. The region number array can be FLUXNUM, MULTNUM or OPERNUM and
these arrays must be defined and be available before the MULTREGD
keyword is read by the simulator. The constant should be a real number.

This keyword is not supported by OPM Flow but would change the results
if supported so the simulation will be stopped.

\begin{longtable}[]{@{}
  >{\raggedright\arraybackslash}p{(\linewidth - 6\tabcolsep) * \real{0.2500}}
  >{\raggedright\arraybackslash}p{(\linewidth - 6\tabcolsep) * \real{0.2500}}
  >{\raggedright\arraybackslash}p{(\linewidth - 6\tabcolsep) * \real{0.2500}}
  >{\raggedright\arraybackslash}p{(\linewidth - 6\tabcolsep) * \real{0.2500}}@{}}
\toprule\noalign{}
\begin{minipage}[b]{\linewidth}\raggedright
No.
\end{minipage} & \begin{minipage}[b]{\linewidth}\raggedright
Name
\end{minipage} & \begin{minipage}[b]{\linewidth}\raggedright
Description
\end{minipage} & \begin{minipage}[b]{\linewidth}\raggedright
Default
\end{minipage} \\
\midrule\noalign{}
\endhead
\bottomrule\noalign{}
\endlastfoot
1 & REGION1 & A positive integer value that defines the from REGION
number for which the CONSTANT in (3) should be applied. & None \\
2 & REGION2 & A positive integer value that defines the to REGION number
for which the CONSTANT in (3) should be applied. & None \\
3 & CONSTANT & A real value to multiply the diffusivity between REGION1
and REGION2. & 1 \\
4 & DIR & A character string that defines the direction to apply the
diffusivity multiplier between the two regions, should be set to one of
the following X, Y, Z, XY, XZ, YZ, or XYZ. & XYZ \\
5 & TYPE & A character string that defines the type of connections the
diffusivity multiplier should be applied to, should be one of the
following: & ALL \\
6 & REGION ARRAY & The REGION ARRAY to use for applying the CONSTANT in
(3) based on the regions REGION1 and REGION2 in (1 and 2). REGION ARRAY
can have the following values: & M \\
Notes: & & & \\
\end{longtable}

\emph{Table 6.76: MULTREGD Keyword Description}

\subsubsection{Example}\label{example-192}

\begin{verbatim}
--
--    MULTIPLY DIFFUSIVITIES BETWEEN RESERVOIRS
--
--    REGION   REGION   DIFFS   DIREC   NNC    REGION ARRAY
--    FROM     TO       MULT    OPT     OPTS   M / F / O
MULTREGD
      1*       1*       1.05    1*     'ALL'   M           / ALL REGIONS
/
\end{verbatim}

The above example multiplies the diffusivities between all the MULTNUM
regions by 1.05 in all directions and for all connections types.

\subsection{MULTREGH -- Multiply Thermal Conductivities Between
Regions}\label{multregh-multiply-thermal-conductivities-between-regions}

\begin{longtable}[]{@{}
  >{\raggedright\arraybackslash}p{(\linewidth - 14\tabcolsep) * \real{0.1250}}
  >{\raggedright\arraybackslash}p{(\linewidth - 14\tabcolsep) * \real{0.1250}}
  >{\raggedright\arraybackslash}p{(\linewidth - 14\tabcolsep) * \real{0.1250}}
  >{\raggedright\arraybackslash}p{(\linewidth - 14\tabcolsep) * \real{0.1250}}
  >{\raggedright\arraybackslash}p{(\linewidth - 14\tabcolsep) * \real{0.1250}}
  >{\raggedright\arraybackslash}p{(\linewidth - 14\tabcolsep) * \real{0.1250}}
  >{\raggedright\arraybackslash}p{(\linewidth - 14\tabcolsep) * \real{0.1250}}
  >{\raggedright\arraybackslash}p{(\linewidth - 14\tabcolsep) * \real{0.1250}}@{}}
\toprule\noalign{}
\begin{minipage}[b]{\linewidth}\raggedright
RUNSPEC
\end{minipage} & \begin{minipage}[b]{\linewidth}\raggedright
GRID
\end{minipage} & \begin{minipage}[b]{\linewidth}\raggedright
EDIT
\end{minipage} & \begin{minipage}[b]{\linewidth}\raggedright
PROPS
\end{minipage} & \begin{minipage}[b]{\linewidth}\raggedright
REGIONS
\end{minipage} & \begin{minipage}[b]{\linewidth}\raggedright
SOLUTION
\end{minipage} & \begin{minipage}[b]{\linewidth}\raggedright
SUMMARY
\end{minipage} & \begin{minipage}[b]{\linewidth}\raggedright
SCHEDULE
\end{minipage} \\
\midrule\noalign{}
\endhead
\bottomrule\noalign{}
\endlastfoot
\end{longtable}

\subsubsection{Description}\label{description-292}

The MULTREGH keyword multiplies the thermal conductivity between two
regions by a constant. The region number array can be FLUXNUM, MULTNUM
or OPERNUM and these arrays must be defined and be available before the
MULTREGT keyword is read by the simulator. The constant should be a real
number.

This keyword is not supported by OPM Flow but would change the results
if supported so the simulation will be stopped.

\begin{longtable}[]{@{}
  >{\raggedright\arraybackslash}p{(\linewidth - 6\tabcolsep) * \real{0.2500}}
  >{\raggedright\arraybackslash}p{(\linewidth - 6\tabcolsep) * \real{0.2500}}
  >{\raggedright\arraybackslash}p{(\linewidth - 6\tabcolsep) * \real{0.2500}}
  >{\raggedright\arraybackslash}p{(\linewidth - 6\tabcolsep) * \real{0.2500}}@{}}
\toprule\noalign{}
\begin{minipage}[b]{\linewidth}\raggedright
No.
\end{minipage} & \begin{minipage}[b]{\linewidth}\raggedright
Name
\end{minipage} & \begin{minipage}[b]{\linewidth}\raggedright
Description
\end{minipage} & \begin{minipage}[b]{\linewidth}\raggedright
Default
\end{minipage} \\
\midrule\noalign{}
\endhead
\bottomrule\noalign{}
\endlastfoot
1 & REGION1 & A positive integer value that defines the from REGION
number for which the CONSTANT in (3) should be applied. & None \\
2 & REGION2 & A positive integer value that defines the to REGION number
for which the CONSTANT in (3) should be applied. & None \\
3 & CONSTANT & A real value to multiply the thermal conductivity between
REGION1 and REGION2. & 1 \\
4 & DIR & A character string that defines the direction to apply the
thermal conductivity multiplier between the two regions, should be set
to one of the following X, Y, Z, XY, XZ, YZ, or XYZ. & XYZ \\
5 & TYPE & A character string that defines the type of connections the
thermal conductivity multiplier should be applied to, should be one of
the following: & ALL \\
6 & REGION ARRAY & The REGION ARRAY to use for applying the CONSTANT in
(3) based on the regions REGION1 and REGION2 in (1 and 2). REGION ARRAY
can have the following values: & M \\
Notes: & & & \\
\end{longtable}

\emph{Table 6.77: MULTREGH Keyword Description}

\subsubsection{Example}\label{example-193}

\begin{verbatim}
--
--    MULTIPLY THERMAL CONDUCTIVITIES BETWEEN RESERVOIRS
--
--    REGION   REGION   CONDS   DIREC   NNC    REGION ARRAY
--    FROM     TO       MULT    OPT     OPTS   M / F / O
MULTREGH
      1*       1*       1.05    1*     'ALL'   M           / ALL REGIONS
/
\end{verbatim}

The above example multiplies the thermal conductivities between all the
MULTNUM regions by 1.05 in all directions and for all connections types.

\subsection{MULTREGP -- Multiply Pore Volumes Based On Region
Number}\label{multregp-multiply-pore-volumes-based-on-region-number}

\begin{longtable}[]{@{}
  >{\raggedright\arraybackslash}p{(\linewidth - 14\tabcolsep) * \real{0.1250}}
  >{\raggedright\arraybackslash}p{(\linewidth - 14\tabcolsep) * \real{0.1250}}
  >{\raggedright\arraybackslash}p{(\linewidth - 14\tabcolsep) * \real{0.1250}}
  >{\raggedright\arraybackslash}p{(\linewidth - 14\tabcolsep) * \real{0.1250}}
  >{\raggedright\arraybackslash}p{(\linewidth - 14\tabcolsep) * \real{0.1250}}
  >{\raggedright\arraybackslash}p{(\linewidth - 14\tabcolsep) * \real{0.1250}}
  >{\raggedright\arraybackslash}p{(\linewidth - 14\tabcolsep) * \real{0.1250}}
  >{\raggedright\arraybackslash}p{(\linewidth - 14\tabcolsep) * \real{0.1250}}@{}}
\toprule\noalign{}
\begin{minipage}[b]{\linewidth}\raggedright
RUNSPEC
\end{minipage} & \begin{minipage}[b]{\linewidth}\raggedright
GRID
\end{minipage} & \begin{minipage}[b]{\linewidth}\raggedright
EDIT
\end{minipage} & \begin{minipage}[b]{\linewidth}\raggedright
PROPS
\end{minipage} & \begin{minipage}[b]{\linewidth}\raggedright
REGIONS
\end{minipage} & \begin{minipage}[b]{\linewidth}\raggedright
SOLUTION
\end{minipage} & \begin{minipage}[b]{\linewidth}\raggedright
SUMMARY
\end{minipage} & \begin{minipage}[b]{\linewidth}\raggedright
SCHEDULE
\end{minipage} \\
\midrule\noalign{}
\endhead
\bottomrule\noalign{}
\endlastfoot
\end{longtable}

\subsubsection{Description}\label{description-293}

The MULTREGP keyword multiplies the pore volume of a cell by a constant
for all cells with a specific region number. The region number array can
be FLUXNUM, MULTNUM or OPERNUM and these arrays must be defined and be
available before the MULTREGP keyword is read by the simulator. The
constant should be a real number.

\begin{longtable}[]{@{}
  >{\raggedright\arraybackslash}p{(\linewidth - 6\tabcolsep) * \real{0.2500}}
  >{\raggedright\arraybackslash}p{(\linewidth - 6\tabcolsep) * \real{0.2500}}
  >{\raggedright\arraybackslash}p{(\linewidth - 6\tabcolsep) * \real{0.2500}}
  >{\raggedright\arraybackslash}p{(\linewidth - 6\tabcolsep) * \real{0.2500}}@{}}
\toprule\noalign{}
\begin{minipage}[b]{\linewidth}\raggedright
No.
\end{minipage} & \begin{minipage}[b]{\linewidth}\raggedright
Name
\end{minipage} & \begin{minipage}[b]{\linewidth}\raggedright
Description
\end{minipage} & \begin{minipage}[b]{\linewidth}\raggedright
Default
\end{minipage} \\
\midrule\noalign{}
\endhead
\bottomrule\noalign{}
\endlastfoot
1 & REGION & REGION is a positive integer representing the region for
which the CONSTANT in (2) should be applied. & None \\
2 & CONSTANT & A real value to multiply the pore volume by for a given
REGION. & 1 \\
3 & REGION ARRAY & The REGION ARRAY to use for applying the CONSTANT in
(2) based on the REGION in (1). ARRAY can have the following values: &
M \\
Notes: & & & \\
\end{longtable}

\emph{Table 6.78: MULTREGP Keyword Description}

\subsubsection{Example}\label{example-194}

\begin{verbatim}
--
-- RESET PORE VOLUME FOR DIFFERENT REGIONS
--
--    REGION    PORV          REGION ARRAY
--    NUMBER    MULT          M / F / O
MULTREGP
         1       1.0456573    M             /  Fault Block 1
         2       0            M             /  Fault Block 2
         3       0.9756715    M             /  Fault Block 3
         4       0            M             /  Inactive Blocks
/
\end{verbatim}

The above example re-scales the pore volumes for MULTNUM regions one and
three and makes regions two and four inactive by setting their pore
volumes to zero.

\subsection{MULTREGT -- Multiply Transmissibilities Between
Regions}\label{multregt-multiply-transmissibilities-between-regions}

\begin{longtable}[]{@{}
  >{\raggedright\arraybackslash}p{(\linewidth - 14\tabcolsep) * \real{0.1250}}
  >{\raggedright\arraybackslash}p{(\linewidth - 14\tabcolsep) * \real{0.1250}}
  >{\raggedright\arraybackslash}p{(\linewidth - 14\tabcolsep) * \real{0.1250}}
  >{\raggedright\arraybackslash}p{(\linewidth - 14\tabcolsep) * \real{0.1250}}
  >{\raggedright\arraybackslash}p{(\linewidth - 14\tabcolsep) * \real{0.1250}}
  >{\raggedright\arraybackslash}p{(\linewidth - 14\tabcolsep) * \real{0.1250}}
  >{\raggedright\arraybackslash}p{(\linewidth - 14\tabcolsep) * \real{0.1250}}
  >{\raggedright\arraybackslash}p{(\linewidth - 14\tabcolsep) * \real{0.1250}}@{}}
\toprule\noalign{}
\begin{minipage}[b]{\linewidth}\raggedright
RUNSPEC
\end{minipage} & \begin{minipage}[b]{\linewidth}\raggedright
GRID
\end{minipage} & \begin{minipage}[b]{\linewidth}\raggedright
EDIT
\end{minipage} & \begin{minipage}[b]{\linewidth}\raggedright
PROPS
\end{minipage} & \begin{minipage}[b]{\linewidth}\raggedright
REGIONS
\end{minipage} & \begin{minipage}[b]{\linewidth}\raggedright
SOLUTION
\end{minipage} & \begin{minipage}[b]{\linewidth}\raggedright
SUMMARY
\end{minipage} & \begin{minipage}[b]{\linewidth}\raggedright
SCHEDULE
\end{minipage} \\
\midrule\noalign{}
\endhead
\bottomrule\noalign{}
\endlastfoot
\end{longtable}

\subsubsection{Description}\label{description-294}

The MULTREGT keyword multiplies the transmissibility between two regions
by a constant. The region number array can be FLUXNUM, MULTNUM or
OPERNUM and these arrays must be defined and be available before the
MULTREGT keyword is read by the simulator. The constant should be a real
number.

\begin{longtable}[]{@{}
  >{\raggedright\arraybackslash}p{(\linewidth - 6\tabcolsep) * \real{0.2500}}
  >{\raggedright\arraybackslash}p{(\linewidth - 6\tabcolsep) * \real{0.2500}}
  >{\raggedright\arraybackslash}p{(\linewidth - 6\tabcolsep) * \real{0.2500}}
  >{\raggedright\arraybackslash}p{(\linewidth - 6\tabcolsep) * \real{0.2500}}@{}}
\toprule\noalign{}
\begin{minipage}[b]{\linewidth}\raggedright
No.
\end{minipage} & \begin{minipage}[b]{\linewidth}\raggedright
Name
\end{minipage} & \begin{minipage}[b]{\linewidth}\raggedright
Description
\end{minipage} & \begin{minipage}[b]{\linewidth}\raggedright
Default
\end{minipage} \\
\midrule\noalign{}
\endhead
\bottomrule\noalign{}
\endlastfoot
1 & REGION1 & A positive integer value that defines the from REGION
number for which the CONSTANT in (3) should be applied. & None \\
2 & REGION2 & A positive integer value that defines the to REGION number
for which the CONSTANT in (3) should be applied. & None \\
3 & CONSTANT & A real positive value to multiply the transmissibility
between REGION1 and REGION2. & 1 \\
4 & DIR & A character string that defines the direction to apply the
transmissibility multiplier between the two regions, should be set to
one of the following X, Y, Z, XY, XZ, YZ, or XYZ. & XYZ \\
5 & TYPE & A character string that defines the type of connections the
transmissibility multiplier should be applied to, should be one of the
following: & ALL \\
6 & REGION ARRAY & A single character that defines the REGION ARRAY that
is used to specify the regions identified by REGION1 and REGION2. REGION
ARRAY can have the following values: & M \\
Notes: & & & \\
\end{longtable}

\emph{Table 6.79: MULTREGT Keyword Description}

\begin{tcolorbox}[enhanced jigsaw, title=\textcolor{quarto-callout-note-color}{\faInfo}\hspace{0.5em}{Note}, coltitle=black, breakable, toprule=.15mm, colbacktitle=quarto-callout-note-color!10!white, leftrule=.75mm, opacitybacktitle=0.6, colback=white, bottomtitle=1mm, toptitle=1mm, titlerule=0mm, left=2mm, colframe=quarto-callout-note-color-frame, arc=.35mm, rightrule=.15mm, bottomrule=.15mm, opacityback=0]

Note if the MULTREGT keyword is used in the EDIT section, OPM Flow will
always apply the changes irrespective, of if the TRANX, TRANY and TRANZ
transmissibility arrays have been entered or not in the EDIT section.
This behavior is different to the commercial simulator that only applies
the keyword if the transmissibility arrays have been entered in the EDIT
section.

\end{tcolorbox}

\subsubsection{Example}\label{example-195}

\begin{verbatim}
--
--    SET TRANSMISSIBILITES ACROSS DIFFERENT RESERVOIRS TO ZERO
--
--    REGION   REGION   TRANS   DIREC   NNC    REGION ARRAY
--    FROM     TO       MULT    OPT     OPTS   M / F / O
MULTREGT
      1*       1*       0.0     1*     'ALL'   M           / ALL REGIONS SEALED
/
\end{verbatim}

The above example isolates all regions from one another by setting the
transmissibility for the MULTNUM regions to zero in all directions and
for all connections types.

\subsection{MULTTHT- -- Multiply Cell Transmissibility in the -Theta
Direction}\label{multtht--multiply-cell-transmissibility-in-the--theta-direction}

\begin{longtable}[]{@{}
  >{\raggedright\arraybackslash}p{(\linewidth - 14\tabcolsep) * \real{0.1250}}
  >{\raggedright\arraybackslash}p{(\linewidth - 14\tabcolsep) * \real{0.1250}}
  >{\raggedright\arraybackslash}p{(\linewidth - 14\tabcolsep) * \real{0.1250}}
  >{\raggedright\arraybackslash}p{(\linewidth - 14\tabcolsep) * \real{0.1250}}
  >{\raggedright\arraybackslash}p{(\linewidth - 14\tabcolsep) * \real{0.1250}}
  >{\raggedright\arraybackslash}p{(\linewidth - 14\tabcolsep) * \real{0.1250}}
  >{\raggedright\arraybackslash}p{(\linewidth - 14\tabcolsep) * \real{0.1250}}
  >{\raggedright\arraybackslash}p{(\linewidth - 14\tabcolsep) * \real{0.1250}}@{}}
\toprule\noalign{}
\begin{minipage}[b]{\linewidth}\raggedright
RUNSPEC
\end{minipage} & \begin{minipage}[b]{\linewidth}\raggedright
GRID
\end{minipage} & \begin{minipage}[b]{\linewidth}\raggedright
EDIT
\end{minipage} & \begin{minipage}[b]{\linewidth}\raggedright
PROPS
\end{minipage} & \begin{minipage}[b]{\linewidth}\raggedright
REGIONS
\end{minipage} & \begin{minipage}[b]{\linewidth}\raggedright
SOLUTION
\end{minipage} & \begin{minipage}[b]{\linewidth}\raggedright
SUMMARY
\end{minipage} & \begin{minipage}[b]{\linewidth}\raggedright
SCHEDULE
\end{minipage} \\
\midrule\noalign{}
\endhead
\bottomrule\noalign{}
\endlastfoot
\end{longtable}

\subsubsection{Description}\label{description-295}

MULTTHT- multiples the transmissibility between two cell faces in the
-Theta direction for all the cells in the model via an array, that is
the keyword sets the transmissibility multiplier of block (I, J, K)
between the cells (I, J-1, K) and (I, J, K). An alternative to defining
the complete array is to use the BOX keyword to define an area of the
grid and then use the MULTTHT- keyword to set the multipliers just for
the area defined by the BOX keyword (see the example).

The keyword should only be used with radial and spider grids, as
declared by the RADIAL or SPIDER keywords in the RUNSPEC section.

\begin{longtable}[]{@{}
  >{\raggedright\arraybackslash}p{(\linewidth - 6\tabcolsep) * \real{0.2500}}
  >{\raggedright\arraybackslash}p{(\linewidth - 6\tabcolsep) * \real{0.2500}}
  >{\raggedright\arraybackslash}p{(\linewidth - 6\tabcolsep) * \real{0.2500}}
  >{\raggedright\arraybackslash}p{(\linewidth - 6\tabcolsep) * \real{0.2500}}@{}}
\toprule\noalign{}
\begin{minipage}[b]{\linewidth}\raggedright
No.
\end{minipage} & \begin{minipage}[b]{\linewidth}\raggedright
Name
\end{minipage} & \begin{minipage}[b]{\linewidth}\raggedright
Description
\end{minipage} & \begin{minipage}[b]{\linewidth}\raggedright
Default
\end{minipage} \\
\midrule\noalign{}
\endhead
\bottomrule\noalign{}
\endlastfoot
1 & MULTTHT- & MULTTHT- is an array of real positive numbers assigning
the transmissibility multipliers in the -Theta direction to each cell
face in the model. Repeat counts may be used, for example 20*100.0. &
1.0 \\
Notes: & & & \\
\end{longtable}

\emph{Table 6.81: MULTTHT- Keyword Description}

Note that OPM Flow does not require the GRIDOPTS(TRANMULT) parameter in
the RUNSPEC section to be set to YES, in order to use this and other
negative directional dependent multiplier keywords in the input deck.
Whereas, the commercial simulator will terminate with an error if the
keyword is present, and the GRIDOPTS(TRANMULT) parameter has not been
set to YES.

See also the MULTTHT, MULTR, MULTR-, MULTZ and MULTZ- keywords for
scaling transmissible between l grid cells in the Theta direction.

\subsubsection{Example}\label{example-196}

\begin{verbatim}
--
--       DEFINE INPUT BOX FOR EDITING INPUT ARRAYS
--
--       ---------- BOX ---------
--       I1  I2   J1  J2   K1  K2
BOX
         10  10   1   6    1   1                           / DEFINE BOX AREA
--
--       SET MULTTHT- TRANSMISSIBILITY MULTIPLIERS
--
MULTTHT-
         6*0.500                                           /
--
--       DEFINE END OF INPUT BOX EDITING OF INPUT ARRAYS
--
ENDBOX
\end{verbatim}

The above example defines a 0.5 scaling multiplier for the six cells
defined by the preceding BOX statement. The ENDBOX keyword resets the
input box to the full grid.

\subsection{MULTTHT -- Multiply Cell Transmissibility in the +Theta
Direction}\label{multtht-multiply-cell-transmissibility-in-the-theta-direction}

\begin{longtable}[]{@{}
  >{\raggedright\arraybackslash}p{(\linewidth - 14\tabcolsep) * \real{0.1250}}
  >{\raggedright\arraybackslash}p{(\linewidth - 14\tabcolsep) * \real{0.1250}}
  >{\raggedright\arraybackslash}p{(\linewidth - 14\tabcolsep) * \real{0.1250}}
  >{\raggedright\arraybackslash}p{(\linewidth - 14\tabcolsep) * \real{0.1250}}
  >{\raggedright\arraybackslash}p{(\linewidth - 14\tabcolsep) * \real{0.1250}}
  >{\raggedright\arraybackslash}p{(\linewidth - 14\tabcolsep) * \real{0.1250}}
  >{\raggedright\arraybackslash}p{(\linewidth - 14\tabcolsep) * \real{0.1250}}
  >{\raggedright\arraybackslash}p{(\linewidth - 14\tabcolsep) * \real{0.1250}}@{}}
\toprule\noalign{}
\begin{minipage}[b]{\linewidth}\raggedright
RUNSPEC
\end{minipage} & \begin{minipage}[b]{\linewidth}\raggedright
GRID
\end{minipage} & \begin{minipage}[b]{\linewidth}\raggedright
EDIT
\end{minipage} & \begin{minipage}[b]{\linewidth}\raggedright
PROPS
\end{minipage} & \begin{minipage}[b]{\linewidth}\raggedright
REGIONS
\end{minipage} & \begin{minipage}[b]{\linewidth}\raggedright
SOLUTION
\end{minipage} & \begin{minipage}[b]{\linewidth}\raggedright
SUMMARY
\end{minipage} & \begin{minipage}[b]{\linewidth}\raggedright
SCHEDULE
\end{minipage} \\
\midrule\noalign{}
\endhead
\bottomrule\noalign{}
\endlastfoot
\end{longtable}

\subsubsection{Description}\label{description-296}

MULTTHT multiples the transmissibility between two cell faces in the
+Theta direction for all the cells in the model via an array, that is
the keyword sets the transmissibility multiplier of block (I, J, K)
between the cells (I, J, K) and (I, J+1, K). An alternative to defining
the complete array is to use the BOX keyword to define an area of the
grid and then use the MULTTHT keyword to set the multipliers just for
the area defined by the BOX keyword (see the example).

The keyword should only be used with radial and spider grids, as
declared by the RADIAL or SPIDER keywords in the RUNSPEC section.

\begin{longtable}[]{@{}
  >{\raggedright\arraybackslash}p{(\linewidth - 6\tabcolsep) * \real{0.2500}}
  >{\raggedright\arraybackslash}p{(\linewidth - 6\tabcolsep) * \real{0.2500}}
  >{\raggedright\arraybackslash}p{(\linewidth - 6\tabcolsep) * \real{0.2500}}
  >{\raggedright\arraybackslash}p{(\linewidth - 6\tabcolsep) * \real{0.2500}}@{}}
\toprule\noalign{}
\begin{minipage}[b]{\linewidth}\raggedright
No.
\end{minipage} & \begin{minipage}[b]{\linewidth}\raggedright
Name
\end{minipage} & \begin{minipage}[b]{\linewidth}\raggedright
Description
\end{minipage} & \begin{minipage}[b]{\linewidth}\raggedright
Default
\end{minipage} \\
\midrule\noalign{}
\endhead
\bottomrule\noalign{}
\endlastfoot
1 & MULTTHT+ & MULTTHT+ is an array of real positive numbers assigning
the transmissibility multipliers in the +Theta direction to each cell
face in the model. Repeat counts may be used, for example 20*100.0. &
1.0 \\
Notes: & & & \\
\end{longtable}

\emph{Table 6.80: MULTTHT Keyword Description}

See also the MULTTHT-, MULTR, MULTR-, MULTZ and MULTZ- keywords for
scaling transmissible between l grid cells in the Theta direction.

\subsubsection{Example}\label{example-197}

\begin{verbatim}
--
--       DEFINE INPUT BOX FOR EDITING INPUT ARRAYS
--
--       ---------- BOX ---------
--       I1  I2   J1  J2   K1  K2
BOX
         10  10   1   6    1   3                           / DEFINE BOX AREA
--
--       SET MULTTHT+ TRANSMISSIBILITY MULTIPLIERS
--
MULTTHT
         18*0.300                                          /
--
--       DEFINE END OF INPUT BOX EDITING OF INPUT ARRAYS
--
ENDBOX
\end{verbatim}

The above example defines a 0.3 scaling multiplier for the 18 cells
defined by the preceding BOX statement. The ENDBOX keyword resets the
input box to the full grid.

\subsection{MULTX- -- Multiply Cell Transmissibility in the -X
Direction}\label{multx--multiply-cell-transmissibility-in-the--x-direction}

\begin{longtable}[]{@{}
  >{\raggedright\arraybackslash}p{(\linewidth - 14\tabcolsep) * \real{0.1250}}
  >{\raggedright\arraybackslash}p{(\linewidth - 14\tabcolsep) * \real{0.1250}}
  >{\raggedright\arraybackslash}p{(\linewidth - 14\tabcolsep) * \real{0.1250}}
  >{\raggedright\arraybackslash}p{(\linewidth - 14\tabcolsep) * \real{0.1250}}
  >{\raggedright\arraybackslash}p{(\linewidth - 14\tabcolsep) * \real{0.1250}}
  >{\raggedright\arraybackslash}p{(\linewidth - 14\tabcolsep) * \real{0.1250}}
  >{\raggedright\arraybackslash}p{(\linewidth - 14\tabcolsep) * \real{0.1250}}
  >{\raggedright\arraybackslash}p{(\linewidth - 14\tabcolsep) * \real{0.1250}}@{}}
\toprule\noalign{}
\begin{minipage}[b]{\linewidth}\raggedright
RUNSPEC
\end{minipage} & \begin{minipage}[b]{\linewidth}\raggedright
GRID
\end{minipage} & \begin{minipage}[b]{\linewidth}\raggedright
EDIT
\end{minipage} & \begin{minipage}[b]{\linewidth}\raggedright
PROPS
\end{minipage} & \begin{minipage}[b]{\linewidth}\raggedright
REGIONS
\end{minipage} & \begin{minipage}[b]{\linewidth}\raggedright
SOLUTION
\end{minipage} & \begin{minipage}[b]{\linewidth}\raggedright
SUMMARY
\end{minipage} & \begin{minipage}[b]{\linewidth}\raggedright
SCHEDULE
\end{minipage} \\
\midrule\noalign{}
\endhead
\bottomrule\noalign{}
\endlastfoot
\end{longtable}

\subsubsection{Description}\label{description-297}

MULTX- multiples the transmissibility between two cell faces in the -X
direction for all the cells in the model via an array, that is the
keyword sets the transmissibility multiplier of block (I, J, K) between
the cells (I-1, J, K) and (I, J, K). An alternative to defining the
complete array is to use the BOX keyword to define an area of the grid
and then use the MULTX- keyword to set the multipliers just for the area
defined by the BOX keyword (see the example).

The keyword can be used for all grid types, except for the Radial Grid
geometry.

\begin{longtable}[]{@{}
  >{\raggedright\arraybackslash}p{(\linewidth - 6\tabcolsep) * \real{0.2500}}
  >{\raggedright\arraybackslash}p{(\linewidth - 6\tabcolsep) * \real{0.2500}}
  >{\raggedright\arraybackslash}p{(\linewidth - 6\tabcolsep) * \real{0.2500}}
  >{\raggedright\arraybackslash}p{(\linewidth - 6\tabcolsep) * \real{0.2500}}@{}}
\toprule\noalign{}
\begin{minipage}[b]{\linewidth}\raggedright
No.
\end{minipage} & \begin{minipage}[b]{\linewidth}\raggedright
Name
\end{minipage} & \begin{minipage}[b]{\linewidth}\raggedright
Description
\end{minipage} & \begin{minipage}[b]{\linewidth}\raggedright
Default
\end{minipage} \\
\midrule\noalign{}
\endhead
\bottomrule\noalign{}
\endlastfoot
1 & MULTX- & MULTX- is an array of real positive numbers assigning the
transmissibility multipliers in the -X direction to each cell face in
the model. Repeat counts may be used, for example 20*100.0. & 1.0 \\
Notes: & & & \\
\end{longtable}

\emph{Table 6.83: MULTX- Keyword Description}

Note that OPM Flow does not require the GRIDOPTS(TRANMULT) parameter in
the RUNSPEC section to be set to YES, in order to use this and other
negative directional dependent multiplier keywords in the input deck.
Whereas, the commercial simulator will terminate with an error if the
keyword is present, and the GRIDOPTS(TRANMULT) parameter has not been
set to YES.

See also the MULTX, MULTY, MULTY-, MULTZ and MULTZ- keywords for scaling
transmissible between grid cells.

\subsubsection{Example}\label{example-198}

\begin{verbatim}
--
--       DEFINE INPUT BOX FOR EDITING INPUT ARRAYS
--
--       ---------- BOX ---------
--       I1  I2   J1  J2   K1  K2
BOX
         10  10   1   6    1   1                           / DEFINE BOX AREA
--
--       SET MULTX- TRANSMISSIBILITY MULTIPLIERS
--
MULTX-
         6*0.500                                           /
--
--       DEFINE END OF INPUT BOX EDITING OF INPUT ARRAYS
--
ENDBOX
\end{verbatim}

The above example defines a 0.5 scaling multiplier for the six cells
defined by the preceding BOX statement. The ENDBOX keyword resets the
input box to the full grid.

\subsection{MULTX -- Multiply Cell Transmissibility in the +X
Direction}\label{multx-multiply-cell-transmissibility-in-the-x-direction}

\begin{longtable}[]{@{}
  >{\raggedright\arraybackslash}p{(\linewidth - 14\tabcolsep) * \real{0.1250}}
  >{\raggedright\arraybackslash}p{(\linewidth - 14\tabcolsep) * \real{0.1250}}
  >{\raggedright\arraybackslash}p{(\linewidth - 14\tabcolsep) * \real{0.1250}}
  >{\raggedright\arraybackslash}p{(\linewidth - 14\tabcolsep) * \real{0.1250}}
  >{\raggedright\arraybackslash}p{(\linewidth - 14\tabcolsep) * \real{0.1250}}
  >{\raggedright\arraybackslash}p{(\linewidth - 14\tabcolsep) * \real{0.1250}}
  >{\raggedright\arraybackslash}p{(\linewidth - 14\tabcolsep) * \real{0.1250}}
  >{\raggedright\arraybackslash}p{(\linewidth - 14\tabcolsep) * \real{0.1250}}@{}}
\toprule\noalign{}
\begin{minipage}[b]{\linewidth}\raggedright
RUNSPEC
\end{minipage} & \begin{minipage}[b]{\linewidth}\raggedright
GRID
\end{minipage} & \begin{minipage}[b]{\linewidth}\raggedright
EDIT
\end{minipage} & \begin{minipage}[b]{\linewidth}\raggedright
PROPS
\end{minipage} & \begin{minipage}[b]{\linewidth}\raggedright
REGIONS
\end{minipage} & \begin{minipage}[b]{\linewidth}\raggedright
SOLUTION
\end{minipage} & \begin{minipage}[b]{\linewidth}\raggedright
SUMMARY
\end{minipage} & \begin{minipage}[b]{\linewidth}\raggedright
SCHEDULE
\end{minipage} \\
\midrule\noalign{}
\endhead
\bottomrule\noalign{}
\endlastfoot
\end{longtable}

\subsubsection{Description}\label{description-298}

MULTX multiples the transmissibility between two cell faces in the +X
direction for all the cells in the model via an array, that is the
keyword sets the transmissibility multiplier of block (I, J, K) between
the cells (I, J, K) and (I+I, J, K).

An alternative to defining the complete array is to use the BOX keyword
to define an area of the grid and then use the MULTX keyword to set the
multipliers just for the area defined by the BOX keyword (see the
example).

The keyword can be used for all grid types, except for the Radial Grid
geometry.

\begin{longtable}[]{@{}
  >{\raggedright\arraybackslash}p{(\linewidth - 6\tabcolsep) * \real{0.2500}}
  >{\raggedright\arraybackslash}p{(\linewidth - 6\tabcolsep) * \real{0.2500}}
  >{\raggedright\arraybackslash}p{(\linewidth - 6\tabcolsep) * \real{0.2500}}
  >{\raggedright\arraybackslash}p{(\linewidth - 6\tabcolsep) * \real{0.2500}}@{}}
\toprule\noalign{}
\begin{minipage}[b]{\linewidth}\raggedright
No.
\end{minipage} & \begin{minipage}[b]{\linewidth}\raggedright
Name
\end{minipage} & \begin{minipage}[b]{\linewidth}\raggedright
Description
\end{minipage} & \begin{minipage}[b]{\linewidth}\raggedright
Default
\end{minipage} \\
\midrule\noalign{}
\endhead
\bottomrule\noalign{}
\endlastfoot
1 & MULTX+ & MULTX+ is an array of real positive numbers assigning the
transmissibility multipliers in the +X direction to each cell face in
the model. Repeat counts may be used, for example 20*100.0. & 1.0 \\
Notes: & & & \\
\end{longtable}

\emph{Table 6.82: MULTX Keyword Description}

See also the MULTX-, MULTY, MULTY-, MULTZ and MULTZ- keywords for
scaling transmissible between grid cells.

\subsubsection{Example}\label{example-199}

\begin{verbatim}
--
--       DEFINE INPUT BOX FOR EDITING INPUT ARRAYS
--
--       ---------- BOX ---------
--       I1  I2   J1  J2   K1  K2
BOX
         10  10   1   6    1   3                           / DEFINE BOX AREA
--
--       SET MULTX+ TRANSMISSIBILITY MULTIPLIERS
--
MULTX
         18*0.300                                          /
--
--       DEFINE END OF INPUT BOX EDITING OF INPUT ARRAYS
--
ENDBOX
\end{verbatim}

The above example defines a 0.3 scaling multiplier for the 18 cells
defined by the preceding BOX statement. The ENDBOX keyword resets the
input box to the full grid.

\subsection{MULTY- -- Multiply Cell Transmissibility in the -Y
Direction}\label{multy--multiply-cell-transmissibility-in-the--y-direction}

\begin{longtable}[]{@{}
  >{\raggedright\arraybackslash}p{(\linewidth - 14\tabcolsep) * \real{0.1250}}
  >{\raggedright\arraybackslash}p{(\linewidth - 14\tabcolsep) * \real{0.1250}}
  >{\raggedright\arraybackslash}p{(\linewidth - 14\tabcolsep) * \real{0.1250}}
  >{\raggedright\arraybackslash}p{(\linewidth - 14\tabcolsep) * \real{0.1250}}
  >{\raggedright\arraybackslash}p{(\linewidth - 14\tabcolsep) * \real{0.1250}}
  >{\raggedright\arraybackslash}p{(\linewidth - 14\tabcolsep) * \real{0.1250}}
  >{\raggedright\arraybackslash}p{(\linewidth - 14\tabcolsep) * \real{0.1250}}
  >{\raggedright\arraybackslash}p{(\linewidth - 14\tabcolsep) * \real{0.1250}}@{}}
\toprule\noalign{}
\begin{minipage}[b]{\linewidth}\raggedright
RUNSPEC
\end{minipage} & \begin{minipage}[b]{\linewidth}\raggedright
GRID
\end{minipage} & \begin{minipage}[b]{\linewidth}\raggedright
EDIT
\end{minipage} & \begin{minipage}[b]{\linewidth}\raggedright
PROPS
\end{minipage} & \begin{minipage}[b]{\linewidth}\raggedright
REGIONS
\end{minipage} & \begin{minipage}[b]{\linewidth}\raggedright
SOLUTION
\end{minipage} & \begin{minipage}[b]{\linewidth}\raggedright
SUMMARY
\end{minipage} & \begin{minipage}[b]{\linewidth}\raggedright
SCHEDULE
\end{minipage} \\
\midrule\noalign{}
\endhead
\bottomrule\noalign{}
\endlastfoot
\end{longtable}

\subsubsection{Description}\label{description-299}

MULTY- multiples the transmissibility between two cell faces in the -Y
direction for all the cells in the model via an array, that is the
keyword sets the transmissibility multiplier of block (I, J, K) between
the cells (I, J-1, K) and (I, J, K). An alternative to defining the
complete array is to use the BOX keyword to define an area of the grid
and then use the MULTY- keyword to set the multipliers just for the area
defined by the BOX keyword (see the example).

The keyword can be used for all grid types, except for the Radial Grid
geometry.

\begin{longtable}[]{@{}
  >{\raggedright\arraybackslash}p{(\linewidth - 6\tabcolsep) * \real{0.2500}}
  >{\raggedright\arraybackslash}p{(\linewidth - 6\tabcolsep) * \real{0.2500}}
  >{\raggedright\arraybackslash}p{(\linewidth - 6\tabcolsep) * \real{0.2500}}
  >{\raggedright\arraybackslash}p{(\linewidth - 6\tabcolsep) * \real{0.2500}}@{}}
\toprule\noalign{}
\begin{minipage}[b]{\linewidth}\raggedright
No.
\end{minipage} & \begin{minipage}[b]{\linewidth}\raggedright
Name
\end{minipage} & \begin{minipage}[b]{\linewidth}\raggedright
Description
\end{minipage} & \begin{minipage}[b]{\linewidth}\raggedright
Default
\end{minipage} \\
\midrule\noalign{}
\endhead
\bottomrule\noalign{}
\endlastfoot
1 & MULTY- & MULTY- is an array of real positive numbers assigning the
transmissibility multipliers in the -Y direction to each cell face in
the model. Repeat counts may be used, for example 20*100.0. & 1.0 \\
Notes: & & & \\
\end{longtable}

\emph{Table 6.85: MULTY- Keyword Description}

Note that OPM Flow does not require the GRIDOPTS(TRANMULT) parameter in
the RUNSPEC section to be set to YES, in order to use this and other
negative directional dependent multiplier keywords in the input deck.
Whereas, the commercial simulator will terminate with an error if the
keyword is present, and the GRIDOPTS(TRANMULT) parameter has not been
set to YES.

See also the MULTY, MULTX, MULTX-, MULTZ and MULTZ- keywords for scaling
transmissible between grid cells.

\subsubsection{Example}\label{example-200}

\begin{verbatim}
--
--       DEFINE INPUT BOX FOR EDITING INPUT ARRAYS
--
--       ---------- BOX ---------
--       I1  I2   J1  J2   K1  K2
BOX
         10  10   1   6    1   1                           / DEFINE BOX AREA
--
--       SET MULTY- TRANSMISSIBILITY MULTIPLIERS
--
MULTY-
         6*0.500                                           /
--
--       DEFINE END OF INPUT BOX EDITING OF INPUT ARRAYS
--
ENDBOX
\end{verbatim}

The above example defines a 0.5 scaling multiplier for the six cells
defined by the preceding BOX statement. The ENDBOX keyword resets the
input box to the full grid.

\subsection{MULTY -- Multiply Cell Transmissibility in the +Y
Direction}\label{multy-multiply-cell-transmissibility-in-the-y-direction}

\begin{longtable}[]{@{}
  >{\raggedright\arraybackslash}p{(\linewidth - 14\tabcolsep) * \real{0.1250}}
  >{\raggedright\arraybackslash}p{(\linewidth - 14\tabcolsep) * \real{0.1250}}
  >{\raggedright\arraybackslash}p{(\linewidth - 14\tabcolsep) * \real{0.1250}}
  >{\raggedright\arraybackslash}p{(\linewidth - 14\tabcolsep) * \real{0.1250}}
  >{\raggedright\arraybackslash}p{(\linewidth - 14\tabcolsep) * \real{0.1250}}
  >{\raggedright\arraybackslash}p{(\linewidth - 14\tabcolsep) * \real{0.1250}}
  >{\raggedright\arraybackslash}p{(\linewidth - 14\tabcolsep) * \real{0.1250}}
  >{\raggedright\arraybackslash}p{(\linewidth - 14\tabcolsep) * \real{0.1250}}@{}}
\toprule\noalign{}
\begin{minipage}[b]{\linewidth}\raggedright
RUNSPEC
\end{minipage} & \begin{minipage}[b]{\linewidth}\raggedright
GRID
\end{minipage} & \begin{minipage}[b]{\linewidth}\raggedright
EDIT
\end{minipage} & \begin{minipage}[b]{\linewidth}\raggedright
PROPS
\end{minipage} & \begin{minipage}[b]{\linewidth}\raggedright
REGIONS
\end{minipage} & \begin{minipage}[b]{\linewidth}\raggedright
SOLUTION
\end{minipage} & \begin{minipage}[b]{\linewidth}\raggedright
SUMMARY
\end{minipage} & \begin{minipage}[b]{\linewidth}\raggedright
SCHEDULE
\end{minipage} \\
\midrule\noalign{}
\endhead
\bottomrule\noalign{}
\endlastfoot
\end{longtable}

\subsubsection{Description}\label{description-300}

MULTY multiples the transmissibility between two cell faces in the +Y
direction for all the cells in the model via an array, that is the
keyword sets the transmissibility multiplier of block (I, J, K) between
the cells (I, J, K) and (I, J+1, K).

An alternative to defining the complete array is to use the BOX keyword
to define an area of the grid and then use the MULTY keyword to set the
multipliers just for the area defined by the BOX keyword (see the
example).

The keyword can be used for all grid types, except for the Radial Grid
geometry.

\begin{longtable}[]{@{}
  >{\raggedright\arraybackslash}p{(\linewidth - 6\tabcolsep) * \real{0.2500}}
  >{\raggedright\arraybackslash}p{(\linewidth - 6\tabcolsep) * \real{0.2500}}
  >{\raggedright\arraybackslash}p{(\linewidth - 6\tabcolsep) * \real{0.2500}}
  >{\raggedright\arraybackslash}p{(\linewidth - 6\tabcolsep) * \real{0.2500}}@{}}
\toprule\noalign{}
\begin{minipage}[b]{\linewidth}\raggedright
No.
\end{minipage} & \begin{minipage}[b]{\linewidth}\raggedright
Name
\end{minipage} & \begin{minipage}[b]{\linewidth}\raggedright
Description
\end{minipage} & \begin{minipage}[b]{\linewidth}\raggedright
Default
\end{minipage} \\
\midrule\noalign{}
\endhead
\bottomrule\noalign{}
\endlastfoot
1 & MULTY+ & MULTY+ is an array of real positive numbers assigning the
transmissibility multipliers in the +Y direction to each cell face in
the model. Repeat counts may be used, for example 20*100.0. & 1.0 \\
Notes: & & & \\
\end{longtable}

\emph{Table 6.84: MULTY Keyword Description}

See also the MULTY-, MULTX, MULTX-, MULTZ and MULTZ- keywords for
scaling transmissible between grid cells.

\subsubsection{Example}\label{example-201}

\begin{verbatim}
--
--       DEFINE INPUT BOX FOR EDITING INPUT ARRAYS
--
--       ---------- BOX ---------
--       I1  I2   J1  J2   K1  K2
BOX
         10  10   1   6    1   3                           / DEFINE BOX AREA
--
--       SET MULTY+ TRANSMISSIBILITY MULTIPLIERS
--
MULTY
         18*0.300                                          /
--
--       DEFINE END OF INPUT BOX EDITING OF INPUT ARRAYS
--
ENDBOX
\end{verbatim}

The above example defines a 0.3 scaling multiplier for the 18 cells
defined by the preceding BOX statement. The ENDBOX keyword resets the
input box to the full grid.

\subsection{MULTZ- -- Multiply Cell Transmissibility in the -Z
Direction}\label{multz--multiply-cell-transmissibility-in-the--z-direction}

\begin{longtable}[]{@{}
  >{\raggedright\arraybackslash}p{(\linewidth - 14\tabcolsep) * \real{0.1250}}
  >{\raggedright\arraybackslash}p{(\linewidth - 14\tabcolsep) * \real{0.1250}}
  >{\raggedright\arraybackslash}p{(\linewidth - 14\tabcolsep) * \real{0.1250}}
  >{\raggedright\arraybackslash}p{(\linewidth - 14\tabcolsep) * \real{0.1250}}
  >{\raggedright\arraybackslash}p{(\linewidth - 14\tabcolsep) * \real{0.1250}}
  >{\raggedright\arraybackslash}p{(\linewidth - 14\tabcolsep) * \real{0.1250}}
  >{\raggedright\arraybackslash}p{(\linewidth - 14\tabcolsep) * \real{0.1250}}
  >{\raggedright\arraybackslash}p{(\linewidth - 14\tabcolsep) * \real{0.1250}}@{}}
\toprule\noalign{}
\begin{minipage}[b]{\linewidth}\raggedright
RUNSPEC
\end{minipage} & \begin{minipage}[b]{\linewidth}\raggedright
GRID
\end{minipage} & \begin{minipage}[b]{\linewidth}\raggedright
EDIT
\end{minipage} & \begin{minipage}[b]{\linewidth}\raggedright
PROPS
\end{minipage} & \begin{minipage}[b]{\linewidth}\raggedright
REGIONS
\end{minipage} & \begin{minipage}[b]{\linewidth}\raggedright
SOLUTION
\end{minipage} & \begin{minipage}[b]{\linewidth}\raggedright
SUMMARY
\end{minipage} & \begin{minipage}[b]{\linewidth}\raggedright
SCHEDULE
\end{minipage} \\
\midrule\noalign{}
\endhead
\bottomrule\noalign{}
\endlastfoot
\end{longtable}

\subsubsection{Description}\label{description-301}

MULTZ- multiples the transmissibility between two cell faces in the -Z
direction for all the cells in the model via an array, that is the
keyword sets the transmissibility multiplier of block (I, J, K) between
the cells (1, J, K-1) and (I, J, K). An alternative to defining the
complete array is to use the BOX keyword to define an area of the grid
and then use the MULTZ- keyword to set the multipliers just for the area
defined by the BOX keyword (see the example).

The keyword can be used for all grid types, except for the Radial Grid
geometry.

\begin{longtable}[]{@{}
  >{\raggedright\arraybackslash}p{(\linewidth - 6\tabcolsep) * \real{0.2500}}
  >{\raggedright\arraybackslash}p{(\linewidth - 6\tabcolsep) * \real{0.2500}}
  >{\raggedright\arraybackslash}p{(\linewidth - 6\tabcolsep) * \real{0.2500}}
  >{\raggedright\arraybackslash}p{(\linewidth - 6\tabcolsep) * \real{0.2500}}@{}}
\toprule\noalign{}
\begin{minipage}[b]{\linewidth}\raggedright
No.
\end{minipage} & \begin{minipage}[b]{\linewidth}\raggedright
Name
\end{minipage} & \begin{minipage}[b]{\linewidth}\raggedright
Description
\end{minipage} & \begin{minipage}[b]{\linewidth}\raggedright
Default
\end{minipage} \\
\midrule\noalign{}
\endhead
\bottomrule\noalign{}
\endlastfoot
1 & MULTZ- & MULTZ- is an array of real positive numbers assigning the
transmissibility multipliers in the -X direction to each cell face in
the model. Repeat counts may be used, for example 20*100.0. & 1.0 \\
Notes: & & & \\
\end{longtable}

\emph{Table 6.87: MULTZ- Keyword Description}

Note that OPM Flow does not require the GRIDOPTS(TRANMULT) parameter in
the RUNSPEC section to be set to YES, in order to use this and other
negative directional dependent multiplier keywords in the input deck.
Whereas, the commercial simulator will terminate with an error if the
keyword is present, and the GRIDOPTS(TRANMULT) parameter has not been
set to YES.

See also the MULTZ, MULTX, MULTX-, MULTY and MULTY- keywords for scaling
transmissible between grid cells.

\subsubsection{Example}\label{example-202}

\begin{verbatim}
--
--       DEFINE INPUT BOX FOR EDITING INPUT ARRAYS
--
--       ---------- BOX ---------
--       I1  I2   J1  J2   K1  K2
BOX
         10  10   1   6    1   1                           / DEFINE BOX AREA
--
--       SET MULTZ- TRANSMISSIBILITY MULTIPLIERS
--
MULTZ-
         6*0.500                                           /
--
--       DEFINE END OF INPUT BOX EDITING OF INPUT ARRAYS
--
ENDBOX
\end{verbatim}

The above example defines a 0.5 scaling multiplier for the six cells
defined by the preceding BOX statement. The ENDBOX keyword resets the
input box to the full grid.

\subsection{MULTZ -- Multiply Cell Transmissibility in the +Z
Direction}\label{multz-multiply-cell-transmissibility-in-the-z-direction}

\begin{longtable}[]{@{}
  >{\raggedright\arraybackslash}p{(\linewidth - 14\tabcolsep) * \real{0.1250}}
  >{\raggedright\arraybackslash}p{(\linewidth - 14\tabcolsep) * \real{0.1250}}
  >{\raggedright\arraybackslash}p{(\linewidth - 14\tabcolsep) * \real{0.1250}}
  >{\raggedright\arraybackslash}p{(\linewidth - 14\tabcolsep) * \real{0.1250}}
  >{\raggedright\arraybackslash}p{(\linewidth - 14\tabcolsep) * \real{0.1250}}
  >{\raggedright\arraybackslash}p{(\linewidth - 14\tabcolsep) * \real{0.1250}}
  >{\raggedright\arraybackslash}p{(\linewidth - 14\tabcolsep) * \real{0.1250}}
  >{\raggedright\arraybackslash}p{(\linewidth - 14\tabcolsep) * \real{0.1250}}@{}}
\toprule\noalign{}
\begin{minipage}[b]{\linewidth}\raggedright
RUNSPEC
\end{minipage} & \begin{minipage}[b]{\linewidth}\raggedright
GRID
\end{minipage} & \begin{minipage}[b]{\linewidth}\raggedright
EDIT
\end{minipage} & \begin{minipage}[b]{\linewidth}\raggedright
PROPS
\end{minipage} & \begin{minipage}[b]{\linewidth}\raggedright
REGIONS
\end{minipage} & \begin{minipage}[b]{\linewidth}\raggedright
SOLUTION
\end{minipage} & \begin{minipage}[b]{\linewidth}\raggedright
SUMMARY
\end{minipage} & \begin{minipage}[b]{\linewidth}\raggedright
SCHEDULE
\end{minipage} \\
\midrule\noalign{}
\endhead
\bottomrule\noalign{}
\endlastfoot
\end{longtable}

\subsubsection{Description}\label{description-302}

MULTZ multiples the transmissibility between two cell faces in the +Z
direction for all the cells in the model via an array, that is the
keyword sets the transmissibility multiplier of block (I, J, K) between
the cells (I, J, K) and (I, J, K+1).

An alternative to defining the complete array is to use the BOX keyword
to define an area of the grid and then use the MULTZ keyword to set the
multipliers just for the area defined by the BOX keyword (see the
example).

The keyword can be used for all grid types, except for the Radial Grid
geometry.

\begin{longtable}[]{@{}
  >{\raggedright\arraybackslash}p{(\linewidth - 6\tabcolsep) * \real{0.2500}}
  >{\raggedright\arraybackslash}p{(\linewidth - 6\tabcolsep) * \real{0.2500}}
  >{\raggedright\arraybackslash}p{(\linewidth - 6\tabcolsep) * \real{0.2500}}
  >{\raggedright\arraybackslash}p{(\linewidth - 6\tabcolsep) * \real{0.2500}}@{}}
\toprule\noalign{}
\begin{minipage}[b]{\linewidth}\raggedright
No.
\end{minipage} & \begin{minipage}[b]{\linewidth}\raggedright
Name
\end{minipage} & \begin{minipage}[b]{\linewidth}\raggedright
Description
\end{minipage} & \begin{minipage}[b]{\linewidth}\raggedright
Default
\end{minipage} \\
\midrule\noalign{}
\endhead
\bottomrule\noalign{}
\endlastfoot
1 & MULTZ+ & MULTZ+ is an array of real positive numbers assigning the
transmissibility multipliers in the +Z direction to each cell face in
the model. Repeat counts may be used, for example 20*100.0. & 1.0 \\
Notes: & & & \\
\end{longtable}

\emph{Table 6.86: MULTZ Keyword Description}

See also the MULTZ-, MULTX, MULTX-, MULTY and MULTY- keywords for
scaling transmissible between grid cells.

\subsubsection{Example}\label{example-203}

\begin{verbatim}
--
--       DEFINE INPUT BOX FOR EDITING INPUT ARRAYS
--
--       ---------- BOX ---------
--       I1  I2   J1  J2   K1  K2
BOX
         10  10   1   18   1   1                           / DEFINE BOX AREA
--
--       SET MULTZ+ TRANSMISSIBILITY MULTIPLIERS
--
MULTZ
         18*0.300                                          /
--
--       DEFINE END OF INPUT BOX EDITING OF INPUT ARRAYS
--
ENDBOX
\end{verbatim}

The above example defines a 0.3 scaling multiplier for the 18 cells
defined by the preceding BOX statement. The ENDBOX keyword resets the
input box to the full grid.

\subsection{NEWTRAN -- Activate Irregular Corner-Point Grid
Transmissibilities}\label{newtran-activate-irregular-corner-point-grid-transmissibilities}

\begin{longtable}[]{@{}
  >{\raggedright\arraybackslash}p{(\linewidth - 14\tabcolsep) * \real{0.1250}}
  >{\raggedright\arraybackslash}p{(\linewidth - 14\tabcolsep) * \real{0.1250}}
  >{\raggedright\arraybackslash}p{(\linewidth - 14\tabcolsep) * \real{0.1250}}
  >{\raggedright\arraybackslash}p{(\linewidth - 14\tabcolsep) * \real{0.1250}}
  >{\raggedright\arraybackslash}p{(\linewidth - 14\tabcolsep) * \real{0.1250}}
  >{\raggedright\arraybackslash}p{(\linewidth - 14\tabcolsep) * \real{0.1250}}
  >{\raggedright\arraybackslash}p{(\linewidth - 14\tabcolsep) * \real{0.1250}}
  >{\raggedright\arraybackslash}p{(\linewidth - 14\tabcolsep) * \real{0.1250}}@{}}
\toprule\noalign{}
\begin{minipage}[b]{\linewidth}\raggedright
RUNSPEC
\end{minipage} & \begin{minipage}[b]{\linewidth}\raggedright
GRID
\end{minipage} & \begin{minipage}[b]{\linewidth}\raggedright
EDIT
\end{minipage} & \begin{minipage}[b]{\linewidth}\raggedright
PROPS
\end{minipage} & \begin{minipage}[b]{\linewidth}\raggedright
REGIONS
\end{minipage} & \begin{minipage}[b]{\linewidth}\raggedright
SOLUTION
\end{minipage} & \begin{minipage}[b]{\linewidth}\raggedright
SUMMARY
\end{minipage} & \begin{minipage}[b]{\linewidth}\raggedright
SCHEDULE
\end{minipage} \\
\midrule\noalign{}
\endhead
\bottomrule\noalign{}
\endlastfoot
\end{longtable}

\subsubsection{Description}\label{description-303}

This keyword switches on Irregular Corner-Point Grid geometry
transmissibility calculation, which is the default option for this type
of grid. Grids defined with the COORD and ZCORN keywords will always
invoke this option by default.

For Cartesian Regular Grids defined by the DX, DY, and DZ series of
keywords the block center geometry transmissibility calculations should
be activated via the OLDTRAN keyword. Again this is automatically
invoked if this type of grid is being employed.

There is no data required for this keyword and there is no terminating
``/'' for this keyword.

\subsubsection{Example}\label{example-204}

\begin{verbatim}
--
--       ACTIVATE IRREGULAR CORNER-POINT GRID TRANSMISSIBILITIES
--
NEWTRAN
\end{verbatim}

The above example manually activates Irregular Corner-Point Grid
transmissibility calculations.

\subsection{NINENUM -- Define the Nine-Point Discretization
Region}\label{ninenum-define-the-nine-point-discretization-region}

\begin{longtable}[]{@{}
  >{\raggedright\arraybackslash}p{(\linewidth - 14\tabcolsep) * \real{0.1250}}
  >{\raggedright\arraybackslash}p{(\linewidth - 14\tabcolsep) * \real{0.1250}}
  >{\raggedright\arraybackslash}p{(\linewidth - 14\tabcolsep) * \real{0.1250}}
  >{\raggedright\arraybackslash}p{(\linewidth - 14\tabcolsep) * \real{0.1250}}
  >{\raggedright\arraybackslash}p{(\linewidth - 14\tabcolsep) * \real{0.1250}}
  >{\raggedright\arraybackslash}p{(\linewidth - 14\tabcolsep) * \real{0.1250}}
  >{\raggedright\arraybackslash}p{(\linewidth - 14\tabcolsep) * \real{0.1250}}
  >{\raggedright\arraybackslash}p{(\linewidth - 14\tabcolsep) * \real{0.1250}}@{}}
\toprule\noalign{}
\begin{minipage}[b]{\linewidth}\raggedright
RUNSPEC
\end{minipage} & \begin{minipage}[b]{\linewidth}\raggedright
GRID
\end{minipage} & \begin{minipage}[b]{\linewidth}\raggedright
EDIT
\end{minipage} & \begin{minipage}[b]{\linewidth}\raggedright
PROPS
\end{minipage} & \begin{minipage}[b]{\linewidth}\raggedright
REGIONS
\end{minipage} & \begin{minipage}[b]{\linewidth}\raggedright
SOLUTION
\end{minipage} & \begin{minipage}[b]{\linewidth}\raggedright
SUMMARY
\end{minipage} & \begin{minipage}[b]{\linewidth}\raggedright
SCHEDULE
\end{minipage} \\
\midrule\noalign{}
\endhead
\bottomrule\noalign{}
\endlastfoot
\end{longtable}

\subsubsection{Description}\label{description-304}

The NINENUM keyword defines areas in the grid that should use the
Nine-Point Discretization formulation by setting a grid block's NINENUM
value to one, or zero for the conventional standard five-point
discretization formulation, for when the Nine-Point Discretization
formulation has been activated by the NINEPOIN keyword in the RUNSPEC
section. There should be a NINENUM value for each grid block in the
model. Note that if the if the NINEPOIN keyword in the RUNSPEC section
has been invoked and the NINENUM keyword has not been used in the input
deck, then all the grid will use the nine-point scheme.

This keyword is not supported by OPM Flow but would change the results
if supported so the simulation will be stopped.

\begin{longtable}[]{@{}
  >{\raggedright\arraybackslash}p{(\linewidth - 6\tabcolsep) * \real{0.2500}}
  >{\raggedright\arraybackslash}p{(\linewidth - 6\tabcolsep) * \real{0.2500}}
  >{\raggedright\arraybackslash}p{(\linewidth - 6\tabcolsep) * \real{0.2500}}
  >{\raggedright\arraybackslash}p{(\linewidth - 6\tabcolsep) * \real{0.2500}}@{}}
\toprule\noalign{}
\begin{minipage}[b]{\linewidth}\raggedright
No.
\end{minipage} & \begin{minipage}[b]{\linewidth}\raggedright
Name
\end{minipage} & \begin{minipage}[b]{\linewidth}\raggedright
Description
\end{minipage} & \begin{minipage}[b]{\linewidth}\raggedright
Default
\end{minipage} \\
\midrule\noalign{}
\endhead
\bottomrule\noalign{}
\endlastfoot
1 & NINENUM & NINENUM defines an integer array of zeros and ones
assigning a grid cell to a particular discretization region, a value of
zero for five-point or a value of one for nine-point discretization.
Note that the default value of one implies a cell is included in the
Nine-Point Discretization region; thus, if a cell is to use the
conventional standard five-point finite difference discretization
formulation, then NINENUM must be explicitly set to zero. & 1 \\
Notes: & & & \\
\end{longtable}

\emph{Table 6.88: NINENUM Keyword Description}

The NINENUM keyword cannot be used in models with Local Grid Refinements
(``LGR'') to set different discretization regions within the model, that
is if LGRs are present in the model either all the grid uses nine-point
discretization, if NINEPOIN is present in the RUNSPEC section, or
five-point if NINEPOIN is absent.

\subsubsection{Example}\label{example-205}

The example below sets a portion of the model to us the Nine-Point
Discretization formulation.

\begin{verbatim}
--
--       DEFINE NINE-POINT DISCRETIZATION REGION FOR ALL CELLS
--
--       ARRAY       CONSTANT     ---------- BOX ---------
--                                I1  I2   J1  J2   K1  K2
EQUALS
         NINENUM’    0            1*  1*   1*  1*   1*  1* / FIVE-POINT
         NINENUM’    1            1*  1*   1*  1*   1   5  / NINE-POINT
/
\end{verbatim}

Here the first line sets all the grid to us the five-point
discretization formulation, all values set to zero, and then the second
line sets all the cells in the layers one to five to use the nine-point
discretization formulation.

\subsection{NMATOPTS -- Define the Discretized Matrix Dual Porosity
Parameters}\label{nmatopts-define-the-discretized-matrix-dual-porosity-parameters}

\begin{longtable}[]{@{}
  >{\raggedright\arraybackslash}p{(\linewidth - 14\tabcolsep) * \real{0.1250}}
  >{\raggedright\arraybackslash}p{(\linewidth - 14\tabcolsep) * \real{0.1250}}
  >{\raggedright\arraybackslash}p{(\linewidth - 14\tabcolsep) * \real{0.1250}}
  >{\raggedright\arraybackslash}p{(\linewidth - 14\tabcolsep) * \real{0.1250}}
  >{\raggedright\arraybackslash}p{(\linewidth - 14\tabcolsep) * \real{0.1250}}
  >{\raggedright\arraybackslash}p{(\linewidth - 14\tabcolsep) * \real{0.1250}}
  >{\raggedright\arraybackslash}p{(\linewidth - 14\tabcolsep) * \real{0.1250}}
  >{\raggedright\arraybackslash}p{(\linewidth - 14\tabcolsep) * \real{0.1250}}@{}}
\toprule\noalign{}
\begin{minipage}[b]{\linewidth}\raggedright
RUNSPEC
\end{minipage} & \begin{minipage}[b]{\linewidth}\raggedright
GRID
\end{minipage} & \begin{minipage}[b]{\linewidth}\raggedright
EDIT
\end{minipage} & \begin{minipage}[b]{\linewidth}\raggedright
PROPS
\end{minipage} & \begin{minipage}[b]{\linewidth}\raggedright
REGIONS
\end{minipage} & \begin{minipage}[b]{\linewidth}\raggedright
SOLUTION
\end{minipage} & \begin{minipage}[b]{\linewidth}\raggedright
SUMMARY
\end{minipage} & \begin{minipage}[b]{\linewidth}\raggedright
SCHEDULE
\end{minipage} \\
\midrule\noalign{}
\endhead
\bottomrule\noalign{}
\endlastfoot
\end{longtable}

\subsubsection{Description}\label{description-305}

The NMATOPTS keyword defines the Discretized Matrix Dual Porosity
parameters for when the Discretized Matrix Dual Porosity option has been
activated by NMATRIX keyword in the RUNSPEC section. The option allows
the matrix grid blocks to be subdivided into smaller cells for more
accurate flow calculations, in particular the modeling of transient flow
within the matrix grid blocks.

This keyword is not supported by OPM Flow but would change the results
if supported so the simulation will be stopped.

\subsection{NNC -- Define Non-Neighbor Connections Between Cells
Manually}\label{nnc-define-non-neighbor-connections-between-cells-manually}

\begin{longtable}[]{@{}
  >{\raggedright\arraybackslash}p{(\linewidth - 14\tabcolsep) * \real{0.1250}}
  >{\raggedright\arraybackslash}p{(\linewidth - 14\tabcolsep) * \real{0.1250}}
  >{\raggedright\arraybackslash}p{(\linewidth - 14\tabcolsep) * \real{0.1250}}
  >{\raggedright\arraybackslash}p{(\linewidth - 14\tabcolsep) * \real{0.1250}}
  >{\raggedright\arraybackslash}p{(\linewidth - 14\tabcolsep) * \real{0.1250}}
  >{\raggedright\arraybackslash}p{(\linewidth - 14\tabcolsep) * \real{0.1250}}
  >{\raggedright\arraybackslash}p{(\linewidth - 14\tabcolsep) * \real{0.1250}}
  >{\raggedright\arraybackslash}p{(\linewidth - 14\tabcolsep) * \real{0.1250}}@{}}
\toprule\noalign{}
\begin{minipage}[b]{\linewidth}\raggedright
RUNSPEC
\end{minipage} & \begin{minipage}[b]{\linewidth}\raggedright
GRID
\end{minipage} & \begin{minipage}[b]{\linewidth}\raggedright
EDIT
\end{minipage} & \begin{minipage}[b]{\linewidth}\raggedright
PROPS
\end{minipage} & \begin{minipage}[b]{\linewidth}\raggedright
REGIONS
\end{minipage} & \begin{minipage}[b]{\linewidth}\raggedright
SOLUTION
\end{minipage} & \begin{minipage}[b]{\linewidth}\raggedright
SUMMARY
\end{minipage} & \begin{minipage}[b]{\linewidth}\raggedright
SCHEDULE
\end{minipage} \\
\midrule\noalign{}
\endhead
\bottomrule\noalign{}
\endlastfoot
\end{longtable}

\subsubsection{Description}\label{description-306}

NNC enables Non-Neighbor Connections (``NNC'') to be manually defined.
This keyword is normally generated by static modeling software as
opposed to be manually entered in the OPM Flow input deck due to the
verbosity and complexity of calculating the required parameters for this
keyword.

\begin{longtable}[]{@{}
  >{\raggedright\arraybackslash}p{(\linewidth - 6\tabcolsep) * \real{0.2500}}
  >{\raggedright\arraybackslash}p{(\linewidth - 6\tabcolsep) * \real{0.2500}}
  >{\raggedright\arraybackslash}p{(\linewidth - 6\tabcolsep) * \real{0.2500}}
  >{\raggedright\arraybackslash}p{(\linewidth - 6\tabcolsep) * \real{0.2500}}@{}}
\toprule\noalign{}
\begin{minipage}[b]{\linewidth}\raggedright
No.
\end{minipage} & \begin{minipage}[b]{\linewidth}\raggedright
Name
\end{minipage} & \begin{minipage}[b]{\linewidth}\raggedright
Description
\end{minipage} & \begin{minipage}[b]{\linewidth}\raggedright
Default
\end{minipage} \\
\midrule\noalign{}
\endhead
\bottomrule\noalign{}
\endlastfoot
Field & Metric & Laboratory & \\
1 & I1 & A positive integer that defines the first grid block in the
I-direction in a non-neighbor connection, must be greater than or equal
to one and less than or equal to NX on the DIMENS keyword in the RUNSPEC
section. & None \\
2 & J1 & A positive integer that defines the first grid block in the
J-direction in a non-neighbor connection, must be greater than or equal
to one and less than or equal to NY on the DIMENS keyword in the RUNSPEC
section. & None \\
3 & K1 & A positive integer that defines the first grid block in the
K-direction in a non-neighbor connection, must be greater than or equal
to one and less than or equal to NZ on the DIMENS keyword in the RUNSPEC
section. & None \\
4 & I2 & A positive integer that defines the second grid block in the
I-direction in a non-neighbor connection, must be greater than or equal
to one and less than or equal to NX on the DIMENS keyword in the RUNSPEC
section. & None \\
5 & J2 & A positive integer that defines the second grid block in the
J-direction in a non-neighbor connection, must be greater than or equal
to one and less than or equal to NY on the DIMENS keyword in the RUNSPEC
section. & None \\
6 & K2 & A positive integer that defines the second grid block in the
K-direction in a non-neighbor connection, must be greater than or equal
to one and less than or equal to NZ on the DIMENS keyword in the RUNSPEC
section. & None \\
7 & TRANSNNC & TRANSNNC is a positive real number greater than or equal
to zero that defines the transmissibility between the first grid block
(I1, J1, K1) and the second grid block (I2, J2, K2). The default value
of zero sets the transmissibility between the two cells to zero. &
0.0 \\
cP.rb/day/psia & cP.rm3/day/bars & cP.rcc/hr/atm & \\
8 & ISATNUM1 & ISATNUM1 is a positive integer defining which saturation
table number (relative permeability table) to be used for flow from the
first grid block to the second grid block. The default value of zero
means the existing saturation table allocated to the upstream cell
(I1,J1,K1). & 0 \\
9 & ISATNUM2 & ISATNUM2 is a positive integer defining which saturation
table number (relative permeability table) to be used for flow from the
second grid block to the first grid block. The default value of zero
means the existing saturation table allocated to the downstream cell
(I2,J2,K2). & 0 \\
10 & IPRSNUM1 & IPRSNUM1 is a positive integer defining which pressure
table number (PVT table) to be used for flow from the first grid block
to the second grid block. The default value of zero means the existing
PVT table allocated to the upstream cell (I1,J1,K1). & 0 \\
11 & IPRSNUM2 & IPRSNUM2 is a positive integer defining which pressure
table number (PVT table) to be used for flow from the second grid block
to the first grid block. The default value of zero means the existing
PVT table allocated to the downstream cell (I2,J2,K2). & 0 \\
12 & FACE1 & FACE1 is a character string that defines the face
associated with flow from the first grid block to the second grid block,
where FACE1 can have values of: X+, X-, Y+, Y-, Z+, or Z-. & None \\
13 & FACE2 & FACE2 is a character string that defines the face
associated with flow from the second grid block to the first grid block,
where FACE2 can have values of: X+, X-, Y+, Y-, Z+, or Z-. & None \\
14 & DIFFNNC & DIFFNNC is a positive real number that defines the
diffusivity between the first grid block (I1, J1, K1) and the second
grid block (I2, J2, K2). & 0.0 \\
feet & meters & cm & \\
15 & DISPNNC & DISPNNC s a positive real number that defines the
dispersion coefficient
\(\frac{1}{(\mathit{Area}x\mathit{Porosity})}\)between the first grid
block (I1, J1, K1) and the second grid block (I2, J2, K2), used with the
DISPERSE option. & 0.0 \\
ft-2 & m-2 & cm-2 & \\
16 & AREANNC & AREANNC is a positive real number that defines the area
associated with the connection between the first grid block (I1, J1, K1)
and the second grid block (I2, J2, K2). & None \\
ft2 & m2 & cm2 & \\
17 & PERMNNC & AREANNC is a positive real number that defines the
permeability associated with the connection between the first grid block
(I1, J1, K1) and the second grid block (I2, J2, K2). This used by the
non-Darcy option. & None \\
mD & mD & mD & \\
Notes: & & & \\
\end{longtable}

\emph{Table 6.89: NNC Keyword Description}

Note that although items (8) to (17) for this keyword are not available
in OPM Flow, even if they were, it is strongly recommended that these
items are defaulted if the data is being entered manually, as opposed to
being generated by a pre-processing software.

Care should be taken that cells in different PVTNUM regions (see the
PVTNUM keyword in the REGIONS section) are not connected, since the
fluid properties are associated with a cell. If for example, a rbbl or a
rm3 of oil flows from PVTNUM region 1 to PVTNUM region 2, then the oil
properties of that oil will change from the PVT 1 data set to the PVT
data set 2. This will result in material balance errors, that may or may
not cause numerical issues.

If the transmissibility across a fault needs to be modified see the
FAULTS and MULTFLT keywords in the GRID section for an alternative and
less complicated method to modifying fault transmissibilities.
Transmissibility between reservoir regions can be modified by using
MULTREGT keyword, provided MULTNUM has been used to define the
inter-region transmissibility region numbers for each grid block.
Finally, the MULTX, MULTY and MULTZ series of keywords can be used to
modify transmissibility between various cells. All the aforementioned
keywords are described in the GRID section.

\subsubsection{Example}\label{example-206}

\begin{verbatim}
--
--       MANUALLY DEFINE NON-NEIGHBOR CONNECTIONS
--
--       ------------ BOX -----------   -- TRANSNNC --
--       I1   J1   K1    I2   J2   K2
NNC
         1    1    1     1    1    2        0.2500     / SET NNC FOR FAULT
         1    1    2     1    1    3        0.2500     / SET NNC FOR FAULT
         1    1    3     1    1    4        0.2500     / SET NNC FOR FAULT
/
\end{verbatim}

The above example defines the transmissibility between cells (1, 1, 1)
and (1, 1, 2), (1, 1, 2) and (1, 1, 3) and finally between (1, 1, 3) and
(1, 1, 4) to be 0.2500.

\subsection{NODPPM -- Deactivate Fracture Porosity-Permeability
Calculation}\label{nodppm-deactivate-fracture-porosity-permeability-calculation-1}

The NODPPM keyword deactivates the default behavior of multiplying the
fracture porosity by the fracture permeability to calculate the
effective fracture permeability in dual porosity and dual permeability
runs. Either the DUALPORO or DUALPERM keywords in the RUNSPEC section
must be declared in the input file in order to use this keyword. If the
default calculation is switched off by this keyword, then the effective
fracture permeability is taken to be those entered for the fracture
using the PERMX, PERMY and PERMZ keywords in the GRID section. If the
keyword is absent from the input deck, then the entered PERMX, PERMY and
PERMZ arrays for the fractures are multiplied by fracture PORO array
values in order to obtain the effective fracture permeability.

See NODPPM -- Deactivate Fracture Porosity-Permeability Calculation in
the RUNSPEC section for a full description.

\subsection{NOECHO -- Deactivate Echoing of User Input Files to the
Print
File}\label{noecho-deactivate-echoing-of-user-input-files-to-the-print-file-2}

Turns off echoing of all the input files to the print file. Note by
default echoing of the inputs files is active. but can subsequently be
switched off by the NOECHO activation keyword.

See NOECHO -- Deactivate Echoing of User Input Files to the Print File
in the GLOBAL section for a full description.

\subsection{NOGGF -- Deactivate Output of Grid Geometry
File}\label{noggf-deactivate-output-of-grid-geometry-file}

\begin{longtable}[]{@{}
  >{\raggedright\arraybackslash}p{(\linewidth - 14\tabcolsep) * \real{0.1250}}
  >{\raggedright\arraybackslash}p{(\linewidth - 14\tabcolsep) * \real{0.1250}}
  >{\raggedright\arraybackslash}p{(\linewidth - 14\tabcolsep) * \real{0.1250}}
  >{\raggedright\arraybackslash}p{(\linewidth - 14\tabcolsep) * \real{0.1250}}
  >{\raggedright\arraybackslash}p{(\linewidth - 14\tabcolsep) * \real{0.1250}}
  >{\raggedright\arraybackslash}p{(\linewidth - 14\tabcolsep) * \real{0.1250}}
  >{\raggedright\arraybackslash}p{(\linewidth - 14\tabcolsep) * \real{0.1250}}
  >{\raggedright\arraybackslash}p{(\linewidth - 14\tabcolsep) * \real{0.1250}}@{}}
\toprule\noalign{}
\begin{minipage}[b]{\linewidth}\raggedright
RUNSPEC
\end{minipage} & \begin{minipage}[b]{\linewidth}\raggedright
GRID
\end{minipage} & \begin{minipage}[b]{\linewidth}\raggedright
EDIT
\end{minipage} & \begin{minipage}[b]{\linewidth}\raggedright
PROPS
\end{minipage} & \begin{minipage}[b]{\linewidth}\raggedright
REGIONS
\end{minipage} & \begin{minipage}[b]{\linewidth}\raggedright
SOLUTION
\end{minipage} & \begin{minipage}[b]{\linewidth}\raggedright
SUMMARY
\end{minipage} & \begin{minipage}[b]{\linewidth}\raggedright
SCHEDULE
\end{minipage} \\
\midrule\noalign{}
\endhead
\bottomrule\noalign{}
\endlastfoot
\end{longtable}

\subsubsection{Description}\label{description-307}

This keyword deactivates the output of a standard GRID or extended GRID
file, as well as the extensible EGRID file for post-processing
applications.

This keyword is not supported by OPM Flow but would change the results
if supported so the simulation will be stopped.

There is no data required for this keyword and there is no terminating
``/'' for this keyword.

\subsubsection{Example}\label{example-207}

\begin{verbatim}
--
--       DEACTIVATE GRID GEOMETRY OUTPUT
--
NOGGF
\end{verbatim}

The above example switches off the default behavior of writing out the
grid geometry files.

\subsection{NOWARN -- Deactivate Warning
Messages}\label{nowarn-deactivate-warning-messages-2}

Turns off warning messages to be printed to the print file; note that
this keyword is deactivated by default and can subsequently be switched
off by the WARN activation keyword. The warning messages may be turned
on and off using keywords WARN and NOWARN.

See NOWARN -- Deactivate Warning Messages in the GLOBAL section for a
full description.

\subsection{NTG -- Define the Net-to-Gross Ratio for All the
Cells}\label{ntg-define-the-net-to-gross-ratio-for-all-the-cells}

\begin{longtable}[]{@{}
  >{\raggedright\arraybackslash}p{(\linewidth - 14\tabcolsep) * \real{0.1250}}
  >{\raggedright\arraybackslash}p{(\linewidth - 14\tabcolsep) * \real{0.1250}}
  >{\raggedright\arraybackslash}p{(\linewidth - 14\tabcolsep) * \real{0.1250}}
  >{\raggedright\arraybackslash}p{(\linewidth - 14\tabcolsep) * \real{0.1250}}
  >{\raggedright\arraybackslash}p{(\linewidth - 14\tabcolsep) * \real{0.1250}}
  >{\raggedright\arraybackslash}p{(\linewidth - 14\tabcolsep) * \real{0.1250}}
  >{\raggedright\arraybackslash}p{(\linewidth - 14\tabcolsep) * \real{0.1250}}
  >{\raggedright\arraybackslash}p{(\linewidth - 14\tabcolsep) * \real{0.1250}}@{}}
\toprule\noalign{}
\begin{minipage}[b]{\linewidth}\raggedright
RUNSPEC
\end{minipage} & \begin{minipage}[b]{\linewidth}\raggedright
GRID
\end{minipage} & \begin{minipage}[b]{\linewidth}\raggedright
EDIT
\end{minipage} & \begin{minipage}[b]{\linewidth}\raggedright
PROPS
\end{minipage} & \begin{minipage}[b]{\linewidth}\raggedright
REGIONS
\end{minipage} & \begin{minipage}[b]{\linewidth}\raggedright
SOLUTION
\end{minipage} & \begin{minipage}[b]{\linewidth}\raggedright
SUMMARY
\end{minipage} & \begin{minipage}[b]{\linewidth}\raggedright
SCHEDULE
\end{minipage} \\
\midrule\noalign{}
\endhead
\bottomrule\noalign{}
\endlastfoot
\end{longtable}

\subsubsection{Description}\label{description-308}

NTG defines the Net-to-Gross Ratio (``NTG'') for all the cells in the
model via an array. The keyword can be used with all grid types.

\begin{longtable}[]{@{}
  >{\raggedright\arraybackslash}p{(\linewidth - 6\tabcolsep) * \real{0.2500}}
  >{\raggedright\arraybackslash}p{(\linewidth - 6\tabcolsep) * \real{0.2500}}
  >{\raggedright\arraybackslash}p{(\linewidth - 6\tabcolsep) * \real{0.2500}}
  >{\raggedright\arraybackslash}p{(\linewidth - 6\tabcolsep) * \real{0.2500}}@{}}
\toprule\noalign{}
\begin{minipage}[b]{\linewidth}\raggedright
No.
\end{minipage} & \begin{minipage}[b]{\linewidth}\raggedright
Name
\end{minipage} & \begin{minipage}[b]{\linewidth}\raggedright
Description
\end{minipage} & \begin{minipage}[b]{\linewidth}\raggedright
Default
\end{minipage} \\
\midrule\noalign{}
\endhead
\bottomrule\noalign{}
\endlastfoot
Field & Metric & Laboratory & \\
1 & NTG & NTG is an array of real numbers greater than or equal to zero
and less than or equal to one, that are assigned the net-to-gross ratio
values for each cell in the model. The number of entries should
correspond to the NX x NY x NZ parameters on the DIMENS keyword. Repeat
counts may be used, for example 200*0.850. & None \\
dimensionless & dimensionless & dimensionless & \\
Notes: & & & \\
\end{longtable}

\emph{Table 6.90: NTG Keyword Description}

See also the PORO, PERMX, PERMY and PERMZ keywords to fully define a
grid's properties.

\subsubsection{Example}\label{example-208}

\begin{verbatim}
--
--        DEFINE GRID BLOCK NTG DATA FOR ALL CELLS (BASED ON NX x NY x NZ = 300)
--
NTG
          100*1.000   100*0.850   100*0.500                                     /
\end{verbatim}

The above example defines a constant NTG of 1.00 for the first 100
cells, then 0.85 for the second 100 hundred cells, and finally 0.500 for
the last 100 cell, for the 300 cells in the model as defined by the
DIMENS keyword in the RUNSPEC section.

\subsection{NXFIN -- Define the Number of LGR Grid Blocks in the
X-Direction}\label{nxfin-define-the-number-of-lgr-grid-blocks-in-the-x-direction}

\begin{longtable}[]{@{}
  >{\raggedright\arraybackslash}p{(\linewidth - 14\tabcolsep) * \real{0.1250}}
  >{\raggedright\arraybackslash}p{(\linewidth - 14\tabcolsep) * \real{0.1250}}
  >{\raggedright\arraybackslash}p{(\linewidth - 14\tabcolsep) * \real{0.1250}}
  >{\raggedright\arraybackslash}p{(\linewidth - 14\tabcolsep) * \real{0.1250}}
  >{\raggedright\arraybackslash}p{(\linewidth - 14\tabcolsep) * \real{0.1250}}
  >{\raggedright\arraybackslash}p{(\linewidth - 14\tabcolsep) * \real{0.1250}}
  >{\raggedright\arraybackslash}p{(\linewidth - 14\tabcolsep) * \real{0.1250}}
  >{\raggedright\arraybackslash}p{(\linewidth - 14\tabcolsep) * \real{0.1250}}@{}}
\toprule\noalign{}
\begin{minipage}[b]{\linewidth}\raggedright
RUNSPEC
\end{minipage} & \begin{minipage}[b]{\linewidth}\raggedright
GRID
\end{minipage} & \begin{minipage}[b]{\linewidth}\raggedright
EDIT
\end{minipage} & \begin{minipage}[b]{\linewidth}\raggedright
PROPS
\end{minipage} & \begin{minipage}[b]{\linewidth}\raggedright
REGIONS
\end{minipage} & \begin{minipage}[b]{\linewidth}\raggedright
SOLUTION
\end{minipage} & \begin{minipage}[b]{\linewidth}\raggedright
SUMMARY
\end{minipage} & \begin{minipage}[b]{\linewidth}\raggedright
SCHEDULE
\end{minipage} \\
\midrule\noalign{}
\endhead
\bottomrule\noalign{}
\endlastfoot
\end{longtable}

\subsubsection{Description}\label{description-309}

NXFIN defines the number of Local Grid Refinement (``LGR'') cells within
a global or host cell in the x-direction via a vector, as opposed to
defining the size for each cell for a Cartesian LGR Grid. The LGR
keyword in the RUNSPEC section should be activated to indicate an LGR is
being used, and the keyword NXFIN should be placed in between the CARFIN
and ENDFIN keywords in the GRID section.

This keyword is not supported by OPM Flow but would change the results
if supported so the simulation will be stopped.

\begin{longtable}[]{@{}
  >{\raggedright\arraybackslash}p{(\linewidth - 6\tabcolsep) * \real{0.2500}}
  >{\raggedright\arraybackslash}p{(\linewidth - 6\tabcolsep) * \real{0.2500}}
  >{\raggedright\arraybackslash}p{(\linewidth - 6\tabcolsep) * \real{0.2500}}
  >{\raggedright\arraybackslash}p{(\linewidth - 6\tabcolsep) * \real{0.2500}}@{}}
\toprule\noalign{}
\begin{minipage}[b]{\linewidth}\raggedright
No.
\end{minipage} & \begin{minipage}[b]{\linewidth}\raggedright
Name
\end{minipage} & \begin{minipage}[b]{\linewidth}\raggedright
Description
\end{minipage} & \begin{minipage}[b]{\linewidth}\raggedright
Default
\end{minipage} \\
\midrule\noalign{}
\endhead
\bottomrule\noalign{}
\endlastfoot
Field & Metric & Laboratory & \\
1 & NXFIN & NXFIN is a vector of integer numbers describing the number
of LGR cells within each defined global or host grid block in the
x-direction in a Cartesian LGR grid. Repeat counts may be used, for
example 2*2.0. & None \\
dimensionless & dimensionless & dimensionless & \\
Notes: & & & \\
\end{longtable}

\emph{Table 6.91: NXFIN Keyword Description}

See also the CARFIN, ENDFIN, NYFIN, and NZFIN keywords in the GRID
section to fully define a Cartesian LGR grid model.

\subsubsection{Example}\label{example-209}

\begin{verbatim}
--
--       CARFIN LGR GRID COMMANDS
--
--       LGR        ----- HOST GRID ------   -- CARFIN GRID --  MAX     HOST
--       NAME       I1  I2  J1  J2  K1  K2     NX    NY    NZ   WELLS   NAME
CARFIN
         LGR-OP01   24  25  87  87   1  50      8     1    50     1     GLOBAL /
--
--       DEFINE LGR GRID BLOCKS IN THE X-DIRECTION
--
NXFIN
         4  4                                                                  /

ENDFIN
\end{verbatim}

The above example splits the global cells (24-25,87, 1-50) into four and
four LGR grid blocks in the x-direction, and since the HXFIN keyword has
not been supplied, then the host cells will split into equal
proportions.

\subsection{NYFIN -- Define the Number of LGR Grid Blocks in the
Y-Direction}\label{nyfin-define-the-number-of-lgr-grid-blocks-in-the-y-direction}

\begin{longtable}[]{@{}
  >{\raggedright\arraybackslash}p{(\linewidth - 14\tabcolsep) * \real{0.1250}}
  >{\raggedright\arraybackslash}p{(\linewidth - 14\tabcolsep) * \real{0.1250}}
  >{\raggedright\arraybackslash}p{(\linewidth - 14\tabcolsep) * \real{0.1250}}
  >{\raggedright\arraybackslash}p{(\linewidth - 14\tabcolsep) * \real{0.1250}}
  >{\raggedright\arraybackslash}p{(\linewidth - 14\tabcolsep) * \real{0.1250}}
  >{\raggedright\arraybackslash}p{(\linewidth - 14\tabcolsep) * \real{0.1250}}
  >{\raggedright\arraybackslash}p{(\linewidth - 14\tabcolsep) * \real{0.1250}}
  >{\raggedright\arraybackslash}p{(\linewidth - 14\tabcolsep) * \real{0.1250}}@{}}
\toprule\noalign{}
\begin{minipage}[b]{\linewidth}\raggedright
RUNSPEC
\end{minipage} & \begin{minipage}[b]{\linewidth}\raggedright
GRID
\end{minipage} & \begin{minipage}[b]{\linewidth}\raggedright
EDIT
\end{minipage} & \begin{minipage}[b]{\linewidth}\raggedright
PROPS
\end{minipage} & \begin{minipage}[b]{\linewidth}\raggedright
REGIONS
\end{minipage} & \begin{minipage}[b]{\linewidth}\raggedright
SOLUTION
\end{minipage} & \begin{minipage}[b]{\linewidth}\raggedright
SUMMARY
\end{minipage} & \begin{minipage}[b]{\linewidth}\raggedright
SCHEDULE
\end{minipage} \\
\midrule\noalign{}
\endhead
\bottomrule\noalign{}
\endlastfoot
\end{longtable}

\subsubsection{Description}\label{description-310}

NYFIN defines the number of Local Grid Refinement (``LGR'') cells within
a global or host cell in the x-direction via a vector, as opposed to
defining the size for each cell for a Cartesian LGR Grid. The LGR
keyword in the RUNSPEC section should be activated to indicate an LGR is
being used, and the keyword NYFIN should be placed in between the CARFIN
and ENDFIN keywords in the GRID section.

This keyword is not supported by OPM Flow but would change the results
if supported so the simulation will be stopped.

\begin{longtable}[]{@{}
  >{\raggedright\arraybackslash}p{(\linewidth - 6\tabcolsep) * \real{0.2500}}
  >{\raggedright\arraybackslash}p{(\linewidth - 6\tabcolsep) * \real{0.2500}}
  >{\raggedright\arraybackslash}p{(\linewidth - 6\tabcolsep) * \real{0.2500}}
  >{\raggedright\arraybackslash}p{(\linewidth - 6\tabcolsep) * \real{0.2500}}@{}}
\toprule\noalign{}
\begin{minipage}[b]{\linewidth}\raggedright
No.
\end{minipage} & \begin{minipage}[b]{\linewidth}\raggedright
Name
\end{minipage} & \begin{minipage}[b]{\linewidth}\raggedright
Description
\end{minipage} & \begin{minipage}[b]{\linewidth}\raggedright
Default
\end{minipage} \\
\midrule\noalign{}
\endhead
\bottomrule\noalign{}
\endlastfoot
Field & Metric & Laboratory & \\
1 & NYFIN & NYFIN is a vector of integer numbers describing the number
of LGR cells within each defined global or host grid block in the
y-direction in a Cartesian LGR grid. Repeat counts may be used, for
example 2*2.0. & None \\
dimensionless & dimensionless & dimensionless & \\
Notes: & & & \\
\end{longtable}

\emph{Table 6.92: NYFIN Keyword Description}

See also the CARFIN, ENDFIN, NXFIN, and NZFIN keywords in the GRID
section to fully define a Cartesian LGR grid model.

\subsubsection{Example}\label{example-210}

\begin{verbatim}
--
--       CARFIN LGR GRID COMMANDS
--
--       LGR        ----- HOST GRID ------   -- CARFIN GRID --  MAX     HOST
--       NAME       I1  I2  J1  J2  K1  K2     NX    NY    NZ   WELLS   NAME
CARFIN
         LGR-OP01   24  24  86  87   1  50      1     8    50     1     GLOBAL /
--
--       DEFINE LGR GRID BLOCKS IN THE Y-DIRECTION
--
NYFIN
         4  4                                                                  /

ENDFIN
\end{verbatim}

The above example splits the global cells (24, 86-87,1-50) into four and
four LGR grid blocks in the y-direction and since the HYFIN keyword has
not been supplied, then the host cells will split into equal
proportions.

\subsection{NZFIN -- Define the Number of LGR Grid Blocks in the
Z-Direction}\label{nzfin-define-the-number-of-lgr-grid-blocks-in-the-z-direction}

\begin{longtable}[]{@{}
  >{\raggedright\arraybackslash}p{(\linewidth - 14\tabcolsep) * \real{0.1250}}
  >{\raggedright\arraybackslash}p{(\linewidth - 14\tabcolsep) * \real{0.1250}}
  >{\raggedright\arraybackslash}p{(\linewidth - 14\tabcolsep) * \real{0.1250}}
  >{\raggedright\arraybackslash}p{(\linewidth - 14\tabcolsep) * \real{0.1250}}
  >{\raggedright\arraybackslash}p{(\linewidth - 14\tabcolsep) * \real{0.1250}}
  >{\raggedright\arraybackslash}p{(\linewidth - 14\tabcolsep) * \real{0.1250}}
  >{\raggedright\arraybackslash}p{(\linewidth - 14\tabcolsep) * \real{0.1250}}
  >{\raggedright\arraybackslash}p{(\linewidth - 14\tabcolsep) * \real{0.1250}}@{}}
\toprule\noalign{}
\begin{minipage}[b]{\linewidth}\raggedright
RUNSPEC
\end{minipage} & \begin{minipage}[b]{\linewidth}\raggedright
GRID
\end{minipage} & \begin{minipage}[b]{\linewidth}\raggedright
EDIT
\end{minipage} & \begin{minipage}[b]{\linewidth}\raggedright
PROPS
\end{minipage} & \begin{minipage}[b]{\linewidth}\raggedright
REGIONS
\end{minipage} & \begin{minipage}[b]{\linewidth}\raggedright
SOLUTION
\end{minipage} & \begin{minipage}[b]{\linewidth}\raggedright
SUMMARY
\end{minipage} & \begin{minipage}[b]{\linewidth}\raggedright
SCHEDULE
\end{minipage} \\
\midrule\noalign{}
\endhead
\bottomrule\noalign{}
\endlastfoot
\end{longtable}

\subsubsection{Description}\label{description-311}

NZFIN defines the number of Local Grid Refinement (``LGR'') cells within
a global or host cell in the z-direction via a vector, as opposed to
defining the size for each cell for a Cartesian LGR Grid. The LGR
keyword in the RUNSPEC section should be activated to indicate an LGR is
being used, and the keyword NXFIN should be placed in between the CARFIN
and ENDFIN keywords in the GRID section.

This keyword is not supported by OPM Flow but would change the results
if supported so the simulation will be stopped.

\begin{longtable}[]{@{}
  >{\raggedright\arraybackslash}p{(\linewidth - 6\tabcolsep) * \real{0.2500}}
  >{\raggedright\arraybackslash}p{(\linewidth - 6\tabcolsep) * \real{0.2500}}
  >{\raggedright\arraybackslash}p{(\linewidth - 6\tabcolsep) * \real{0.2500}}
  >{\raggedright\arraybackslash}p{(\linewidth - 6\tabcolsep) * \real{0.2500}}@{}}
\toprule\noalign{}
\begin{minipage}[b]{\linewidth}\raggedright
No.
\end{minipage} & \begin{minipage}[b]{\linewidth}\raggedright
Name
\end{minipage} & \begin{minipage}[b]{\linewidth}\raggedright
Description
\end{minipage} & \begin{minipage}[b]{\linewidth}\raggedright
Default
\end{minipage} \\
\midrule\noalign{}
\endhead
\bottomrule\noalign{}
\endlastfoot
Field & Metric & Laboratory & \\
1 & NZFIN & NZFIN is a vector of integer numbers describing the number
of LGR cells within each defined global or host grid block in the
x-direction in a Cartesian LGR grid. Repeat counts may be used, for
example 2*2.0. & None \\
dimensionless & dimensionless & dimensionless & \\
Notes: & & & \\
\end{longtable}

\emph{Table 6.93: NZFIN Keyword Description}

See also the CARFIN, ENDFIN, NXFIN, and NYFIN keywords in the GRID
section to fully define a Cartesian LGR grid model.

\subsubsection{Example}\label{example-211}

\begin{verbatim}
--
--       CARFIN LGR GRID COMMANDS
--
--       LGR        ----- HOST GRID ------   -- CARFIN GRID --  MAX     HOST
--       NAME       I1  I2  J1  J2  K1  K2     NX    NY    NZ   WELLS   NAME
CARFIN
         LGR-OP01   24  24  86  86   1  50      8     1    100    1     GLOBAL /
--
--       DEFINE LGR GRID BLOCKS IN THE Z-DIRECTION
--
NZFIN
         50*2                                                                  /

ENDFIN
\end{verbatim}

The above example splits the global cells (24, 86, 1-50) into two LGR
grid blocks per host cell in the z-direction, and since the HZFIN
keyword has not been supplied, then the host cells will split into equal
proportions.

\subsection{OLDTRAN -- Activate Cartesian Regular Grid
Transmissibilities}\label{oldtran-activate-cartesian-regular-grid-transmissibilities}

\begin{longtable}[]{@{}
  >{\raggedright\arraybackslash}p{(\linewidth - 14\tabcolsep) * \real{0.1250}}
  >{\raggedright\arraybackslash}p{(\linewidth - 14\tabcolsep) * \real{0.1250}}
  >{\raggedright\arraybackslash}p{(\linewidth - 14\tabcolsep) * \real{0.1250}}
  >{\raggedright\arraybackslash}p{(\linewidth - 14\tabcolsep) * \real{0.1250}}
  >{\raggedright\arraybackslash}p{(\linewidth - 14\tabcolsep) * \real{0.1250}}
  >{\raggedright\arraybackslash}p{(\linewidth - 14\tabcolsep) * \real{0.1250}}
  >{\raggedright\arraybackslash}p{(\linewidth - 14\tabcolsep) * \real{0.1250}}
  >{\raggedright\arraybackslash}p{(\linewidth - 14\tabcolsep) * \real{0.1250}}@{}}
\toprule\noalign{}
\begin{minipage}[b]{\linewidth}\raggedright
RUNSPEC
\end{minipage} & \begin{minipage}[b]{\linewidth}\raggedright
GRID
\end{minipage} & \begin{minipage}[b]{\linewidth}\raggedright
EDIT
\end{minipage} & \begin{minipage}[b]{\linewidth}\raggedright
PROPS
\end{minipage} & \begin{minipage}[b]{\linewidth}\raggedright
REGIONS
\end{minipage} & \begin{minipage}[b]{\linewidth}\raggedright
SOLUTION
\end{minipage} & \begin{minipage}[b]{\linewidth}\raggedright
SUMMARY
\end{minipage} & \begin{minipage}[b]{\linewidth}\raggedright
SCHEDULE
\end{minipage} \\
\midrule\noalign{}
\endhead
\bottomrule\noalign{}
\endlastfoot
\end{longtable}

\subsubsection{Description}\label{description-312}

This keyword switches on Cartesian Regular Grids geometry
transmissibility calculation (or block centered transmissibility
calculations), which is the default option for this type of grid. Grids
defined by the DX, DY, and DZ series of keywords will always invoke this
option by default.

For Irregular Corner-Point Grids defined by the COORD and ZCORN keywords
Irregular Corner-Point Grid geometry transmissibility calculations
should be activated via the NEWTRAN keyword. Again this is automatically
invoked if this type of grid is being employed.

There is no data required for this keyword and there is no terminating
``/'' for this keyword.

\subsubsection{Example}\label{example-212}

\begin{verbatim}
--
--       ACTIVATE CARTESIAN REGULAR GRID TRANSMISSIBILITIES
--
OLDTRAN
\end{verbatim}

The above example manually activates Cartesian Regular Grid
transmissibility calculations.

\subsection{OLDTRANR -- Activate Radial Regular Grid
Transmissibilities}\label{oldtranr-activate-radial-regular-grid-transmissibilities}

\begin{longtable}[]{@{}
  >{\raggedright\arraybackslash}p{(\linewidth - 14\tabcolsep) * \real{0.1250}}
  >{\raggedright\arraybackslash}p{(\linewidth - 14\tabcolsep) * \real{0.1250}}
  >{\raggedright\arraybackslash}p{(\linewidth - 14\tabcolsep) * \real{0.1250}}
  >{\raggedright\arraybackslash}p{(\linewidth - 14\tabcolsep) * \real{0.1250}}
  >{\raggedright\arraybackslash}p{(\linewidth - 14\tabcolsep) * \real{0.1250}}
  >{\raggedright\arraybackslash}p{(\linewidth - 14\tabcolsep) * \real{0.1250}}
  >{\raggedright\arraybackslash}p{(\linewidth - 14\tabcolsep) * \real{0.1250}}
  >{\raggedright\arraybackslash}p{(\linewidth - 14\tabcolsep) * \real{0.1250}}@{}}
\toprule\noalign{}
\begin{minipage}[b]{\linewidth}\raggedright
RUNSPEC
\end{minipage} & \begin{minipage}[b]{\linewidth}\raggedright
GRID
\end{minipage} & \begin{minipage}[b]{\linewidth}\raggedright
EDIT
\end{minipage} & \begin{minipage}[b]{\linewidth}\raggedright
PROPS
\end{minipage} & \begin{minipage}[b]{\linewidth}\raggedright
REGIONS
\end{minipage} & \begin{minipage}[b]{\linewidth}\raggedright
SOLUTION
\end{minipage} & \begin{minipage}[b]{\linewidth}\raggedright
SUMMARY
\end{minipage} & \begin{minipage}[b]{\linewidth}\raggedright
SCHEDULE
\end{minipage} \\
\midrule\noalign{}
\endhead
\bottomrule\noalign{}
\endlastfoot
\end{longtable}

\subsubsection{Description}\label{description-313}

This keyword switches on Radial Regular Grids geometry transmissibility
calculation (or block centered transmissibility calculations), which is
the default option for this type of grid. Grids defined by the DR,
DTHETA, and DZ series of keywords will always invoke this option by
default.

For Irregular Corner-Point Grids defined by the COORD and ZCORN keywords
Irregular Corner-Point Grid geometry transmissibility calculations
should be activated via the NEWTRAN keyword. Again this is automatically
invoked if this type of grid is being employed.

There is no data required for this keyword and there is no terminating
``/'' for this keyword.

\subsubsection{Example}\label{example-213}

\begin{verbatim}
--
--       ACTIVATE RADIAL REGULAR GRID TRANSMISSIBILITIES
--
OLDTRANR
\end{verbatim}

The above example manually activates Radial Regular Grid
transmissibility calculations.

\subsection{OPERATE -- Define Mathematical Operations on
Arrays}\label{operate-define-mathematical-operations-on-arrays}

\begin{longtable}[]{@{}
  >{\raggedright\arraybackslash}p{(\linewidth - 14\tabcolsep) * \real{0.1250}}
  >{\raggedright\arraybackslash}p{(\linewidth - 14\tabcolsep) * \real{0.1250}}
  >{\raggedright\arraybackslash}p{(\linewidth - 14\tabcolsep) * \real{0.1250}}
  >{\raggedright\arraybackslash}p{(\linewidth - 14\tabcolsep) * \real{0.1250}}
  >{\raggedright\arraybackslash}p{(\linewidth - 14\tabcolsep) * \real{0.1250}}
  >{\raggedright\arraybackslash}p{(\linewidth - 14\tabcolsep) * \real{0.1250}}
  >{\raggedright\arraybackslash}p{(\linewidth - 14\tabcolsep) * \real{0.1250}}
  >{\raggedright\arraybackslash}p{(\linewidth - 14\tabcolsep) * \real{0.1250}}@{}}
\toprule\noalign{}
\begin{minipage}[b]{\linewidth}\raggedright
RUNSPEC
\end{minipage} & \begin{minipage}[b]{\linewidth}\raggedright
GRID
\end{minipage} & \begin{minipage}[b]{\linewidth}\raggedright
EDIT
\end{minipage} & \begin{minipage}[b]{\linewidth}\raggedright
PROPS
\end{minipage} & \begin{minipage}[b]{\linewidth}\raggedright
REGIONS
\end{minipage} & \begin{minipage}[b]{\linewidth}\raggedright
SOLUTION
\end{minipage} & \begin{minipage}[b]{\linewidth}\raggedright
SUMMARY
\end{minipage} & \begin{minipage}[b]{\linewidth}\raggedright
SCHEDULE
\end{minipage} \\
\midrule\noalign{}
\endhead
\bottomrule\noalign{}
\endlastfoot
\end{longtable}

\subsubsection{Description}\label{description-314}

The OPERATE keyword performs a mathematical operation on a specified
array or part of an array, optionally using another specified array as
input to the operation. The keyword allows for various mathematical
functions and their associated variables to be defined and applied to
the specified array. Input constants can be integer or real valued
depending on the array type; however, the arrays that can be operated on
are dependent on which section the OPERATE keyword is being applied in.

\begin{longtable}[]{@{}
  >{\raggedright\arraybackslash}p{(\linewidth - 6\tabcolsep) * \real{0.2500}}
  >{\raggedright\arraybackslash}p{(\linewidth - 6\tabcolsep) * \real{0.2500}}
  >{\raggedright\arraybackslash}p{(\linewidth - 6\tabcolsep) * \real{0.2500}}
  >{\raggedright\arraybackslash}p{(\linewidth - 6\tabcolsep) * \real{0.2500}}@{}}
\toprule\noalign{}
\begin{minipage}[b]{\linewidth}\raggedright
No.
\end{minipage} & \begin{minipage}[b]{\linewidth}\raggedright
Name
\end{minipage} & \begin{minipage}[b]{\linewidth}\raggedright
Description
\end{minipage} & \begin{minipage}[b]{\linewidth}\raggedright
Default
\end{minipage} \\
\midrule\noalign{}
\endhead
\bottomrule\noalign{}
\endlastfoot
1 & Y & A character string of up to eight characters in length that
defines the keyword identifying the array to be modified. & None \\
2 & I1 & A positive integer that defines the lower bound of the array in
the I-direction to be modified must be greater than or equal to one and
less than or equal to I2 and NX. & 1 \\
3 & I2 & A positive integer that defines the upper bound of the array in
the I-direction to be modified must be greater than or equal to I1 and
less than or equal to NX & NX \\
4 & J1 & A positive integer that defines the lower bound of the array in
the J-direction to be modified must be greater than or equal to one and
less than or equal to J2 and NY. & 1 \\
5 & J2 & A positive integer that defines the upper bound of the array in
the J-direction to be modified must be greater than or equal to J1 and
less than or equal to NY. & NY \\
6 & K1 & A positive integer that defines the lower bound of the array in
the K-direction to be modified must be greater than or equal to one and
less than or equal to K2 and NZ. & 1 \\
7 & K2 & A positive integer that defines the upper bound of the array in
the K-direction to be modified must be greater than or equal to K1 and
less than or equal to NZ. & NZ \\
8 & EQUATION & A defined character string of up to eight characters that
defines the mathematical function to be applied, using the X array and
the ALPHA and BETA constants declared on this keyword. EQUATION should
be set to one of the following character strings: & None \\
`MULTA' -\(Y = αX + β\) `POLY' -\(Y =Y + α{X}^{β}\) `SLOG'
-\(Y = {10}^{α + βX}\) `LOG10' -\(Y = log(X)\) `LOGE' -\(Y = ln(X)\)
`INV' -\(Y = \frac{1}{X}\) `MULTX' -\(Y = αX\) & `ADDX' -\(Y = α + X\)
`COPY' -\(Y = X\) `MAXLIM' -\(Y = \mathit{min}(α,X)\) `MINLIM'
-\(Y = \mathit{max}(α,X)\) `MULTP' -\(Y = α{X}^{β}\) `ABS'
-\(Y = \left\|(X)\left\|\) `MULTIPLY' -\(Y = \mathit{XY}\) & & \\
9 & X & A character string of up to eight characters in length that
defines the keyword identifying the array to be used as an input
parameter. & None \\
10 & ALPHA & An integer or real value that is the α variable in the
EQUATION function. & None \\
11 & BETA & An integer or real value that is the β variable in the
EQUATION function. & None \\
Notes: & & & \\
\end{longtable}

\emph{Table 6.94: OPERATE Keyword Description}

Examples of the arrays most commonly operated on in each section are
given in Table 6.95 . Cells colored red indicate arrays that are not
supported by OPM Flow operations.

\begin{longtable}[]{@{}lllllll@{}}
\toprule\noalign{}
OPERATE Keyword and Variable Options by Section & & & & & & \\
\midrule\noalign{}
\endhead
\bottomrule\noalign{}
\endlastfoot
GRID & EDIT & PROPS & REGIONS & SOLUTION & SUMMARY & SCHEDULE \\
DX & DEPTH & SWL & ENDNUM & PRESSURE & & \\
DY & PORV & SWCR & EQLNUM & SWAT & & \\
DZ & TRANX & SWU & FIPNUM & SGAS & & \\
PERMX & TRANY & SGL & IMBNUM & RV & & \\
PERMY & TRANZ & SGCR & MISCNUM & RS & & \\
PERMZ & DIFFX & SGU & PVTNUM & TBLK & & \\
MULTX & DIFFY & KRW & ROCKNUM & GI & & \\
MULTY & DIFFZ & KRO & SATNUM & OILAPI & & \\
MULTZ & TRANR & KRG & WH2NUM & SALT & & \\
DR & TRANTHT & PCG & & GASCONC & & \\
DTHETA & DIFFR & PCW & & SOLVCONC & & \\
PERMR & DIFFTHT & & & SOLVFRAC & & \\
PERMTHT & & & & SFOAM & & \\
DZNET & & & & SPOLY & & \\
PORO & & & & & & \\
NTG & & & & & & \\
FLUXNUM & & & & & & \\
MULTNUM & & & & & & \\
MPFANUM & & & & & & \\
DIFFX & & & & & & \\
DIFFY & & & & & & \\
DIFFZ & & & & & & \\
DIFFR & & & & & & \\
DIFFTHT & & & & & & \\
\end{longtable}

\emph{Table 6.95: OPERATE Keyword Applicable Arrays by Section}

Double precision and integer work arrays may also be used to hold
temporary values for use in more complex operations. Note that these
work arrays are reset to zero at the start of each section of the data
file. The names of the double precision work arrays should be of the
form `WORKn', where n is a positive integer that defines the work array
index, for example `WORK1', `WORK2', and so on. Similarly, the names of
integer work arrays should be of the form `IWORKn', for example
`IWORK1', `IWORK2', and so on. OPM Flow currently only supports double
precision work arrays, and does not currently support using work arrays
with ``global'' arrays (MULTZ\emph{). Note that the permeability arrays
(PERM}) are treated as ``global'' arrays in the GRID section only. In
the commercial simulator the maximum number of double precision and
integer work arrays are specified using the REGDIMS keyword. Whereas,
OPM Flow does not impose a limit on the maximum number of work arrays.

Note that care should be exercised when performing operations on integer
array data as all transforms are performed using floating point
arithmetic operations. In addition, operations on any of the
transmissibility arrays (TRANX, TRANX-, TRANY, TRANY-, TRANZ, and
TRANZ-) may result in untended consequences as these arrays have zero
values on the boundary of the grid. In this use OPM ResInsight to verify
and visually inspect the results.

\subsubsection{Example}\label{example-214}

The first example uses the MULTP function combined with the Net-to-Gross
(NTG) array to re-scale the MULTX, MULTY and MULTZ arrays to reduce the
transmissibility in three separate reservoirs based on the reservoir
quality (NTG).

\begin{verbatim}
--
--       MATHEMATICAL OPERATIONS ON ARRAYS BY CELL
--
--       OUTPUT  ---------- BOX ---------  OPERATION  INPUT   ALPHA   BETA
--       ARRAY   I1  I2   J1  J2   K1  K2  ---------  ARRAY   -----   ----
OPERATE
         MULTX   1*  1*   1*  1*     1 32  ‘MULTP’    NTG     1.00    0.75 / RES1
         MULTY   1*  1*   1*  1*     1 32  ‘MULTP’    NTG     1.00    0.75 / RES1
         MULTZ   1*  1*   1*  1*     1 32  ‘MULTP’    NTG     1.00    0.75 / RES1

         MULTX   1*  1*   1*  1*    34 64  ‘MULTP’    NTG     1.00    0.85 / RES2
         MULTY   1*  1*   1*  1*    34 64  ‘MULTP’    NTG     1.00    0.85 / RES2
         MULTZ   1*  1*   1*  1*    34 64  ‘MULTP’    NTG     1.00    0.85 / RES2

         MULTX   1*  1*   1*  1*    67 96  ‘MULTP’    NTG     1.00    0.50 / RES3
         MULTY   1*  1*   1*  1*    67 96  ‘MULTP’    NTG     1.00    0.50 / RES3
         MULTZ   1*  1*   1*  1*    67 96  ‘MULTP’    NTG     1.00    0.50 / RES3
/
\end{verbatim}

The next example shows how to set the maximum gas saturation (SGU) based
on the minimum (lowest) water saturation (SWL) when using the End-Point
Scaling option.

\begin{verbatim}
--
--       MATHEMATICAL OPERATIONS ON ARRAYS
--
--       OUTPUT  ---------- BOX ---------  OPERATION  INPUT   ALPHA   BETA
--       ARRAY   I1  I2   J1  J2   K1  K2  ---------  ARRAY   -----   ----
OPERATE
         SGU     1*  1*   1*  1*   1*  1*  'MULTA'     SWL    -1.0    1.0      /
/
\end{verbatim}

The above example sets the maximum gas saturation to be one minus the
minimum water saturation.

\subsection{OPERATER -- Define Mathematical Operations on Arrays by
Region}\label{operater-define-mathematical-operations-on-arrays-by-region}

\begin{longtable}[]{@{}
  >{\raggedright\arraybackslash}p{(\linewidth - 14\tabcolsep) * \real{0.1250}}
  >{\raggedright\arraybackslash}p{(\linewidth - 14\tabcolsep) * \real{0.1250}}
  >{\raggedright\arraybackslash}p{(\linewidth - 14\tabcolsep) * \real{0.1250}}
  >{\raggedright\arraybackslash}p{(\linewidth - 14\tabcolsep) * \real{0.1250}}
  >{\raggedright\arraybackslash}p{(\linewidth - 14\tabcolsep) * \real{0.1250}}
  >{\raggedright\arraybackslash}p{(\linewidth - 14\tabcolsep) * \real{0.1250}}
  >{\raggedright\arraybackslash}p{(\linewidth - 14\tabcolsep) * \real{0.1250}}
  >{\raggedright\arraybackslash}p{(\linewidth - 14\tabcolsep) * \real{0.1250}}@{}}
\toprule\noalign{}
\begin{minipage}[b]{\linewidth}\raggedright
RUNSPEC
\end{minipage} & \begin{minipage}[b]{\linewidth}\raggedright
GRID
\end{minipage} & \begin{minipage}[b]{\linewidth}\raggedright
EDIT
\end{minipage} & \begin{minipage}[b]{\linewidth}\raggedright
PROPS
\end{minipage} & \begin{minipage}[b]{\linewidth}\raggedright
REGIONS
\end{minipage} & \begin{minipage}[b]{\linewidth}\raggedright
SOLUTION
\end{minipage} & \begin{minipage}[b]{\linewidth}\raggedright
SUMMARY
\end{minipage} & \begin{minipage}[b]{\linewidth}\raggedright
SCHEDULE
\end{minipage} \\
\midrule\noalign{}
\endhead
\bottomrule\noalign{}
\endlastfoot
\end{longtable}

\subsubsection{Description}\label{description-315}

The OPERATER keyword is similar to the OPERATE keyword, except it
applies the mathematical operation on specific regions, whereas, OPERATE
applies the operations on a cell by cell basis. Here the OPERATER
keyword performs a mathematical operation on a specified property array,
optionally using another property array as input to the function. The
keyword allows for various mathematical functions and their associated
variables to be defined and applied to the specified region data. Input
constants can be integer or real valued depending on the array type;
however, the arrays that can be operated on are dependent on which
section the OPERATER keyword is being applied in.

\begin{longtable}[]{@{}
  >{\raggedright\arraybackslash}p{(\linewidth - 6\tabcolsep) * \real{0.2500}}
  >{\raggedright\arraybackslash}p{(\linewidth - 6\tabcolsep) * \real{0.2500}}
  >{\raggedright\arraybackslash}p{(\linewidth - 6\tabcolsep) * \real{0.2500}}
  >{\raggedright\arraybackslash}p{(\linewidth - 6\tabcolsep) * \real{0.2500}}@{}}
\toprule\noalign{}
\begin{minipage}[b]{\linewidth}\raggedright
No.
\end{minipage} & \begin{minipage}[b]{\linewidth}\raggedright
Name
\end{minipage} & \begin{minipage}[b]{\linewidth}\raggedright
Description
\end{minipage} & \begin{minipage}[b]{\linewidth}\raggedright
Default
\end{minipage} \\
\midrule\noalign{}
\endhead
\bottomrule\noalign{}
\endlastfoot
1 & Y & A character string of up to eight characters in length that
defines the keyword identifying the array to be modified. & None \\
2 & REGION & REGION is a positive integer representing the region for
which the EQUATION should be applied. The default is to use the region
number from the OPERNUM keyword; however this can be reset to another
region array via the ARRAY item on this keyword, provided the array
exists at the time the keyword is declared in the input deck. Note also
the OPERNUM keyword must precede the use of the OPERATER keyword. & 0 \\
3 & EQUATION & A defined character string of up to eight characters that
defines the mathematical function to be applied, using the X array and
the ALPHA and BETA constants declared on this keyword. EQUATION should
be set to one of the following character strings: & None \\
`MULTA' -\(Y = αX + β\) `POLY' -\(Y =Y + α{X}^{β}\) `SLOG'
-\(Y = {10}^{α + βX}\) `LOG10' -\(Y = log(X)\) `LOGE' -\(Y = ln(X)\)
`INV' -\(Y = \frac{1}{X}\) `MULTX' -\(Y = αX\) & `ADDX' -\(Y = α + X\)
`COPY' -\(Y = X\) `MAXLIM' -\(Y = \mathit{min}(α,X)\) `MINLIM'
-\(Y = \mathit{max}(α,X)\) `MULTP' -\(Y = α{X}^{β}\) `ABS'
-\(Y = \left\|(X)\left\|\) `MULTIPLY' -\(Y = \mathit{XY}\) & & \\
4 & X & A character string of up to eight characters in length that
defines the keyword identifying the array to be used as an input
parameter. & None \\
5 & ALPHA & An integer or real value that is the α variable in the
EQUATION function. & None \\
6 & BETA & An integer or real value that is the β variable in the
EQUATION function. & None \\
7 & ARRAY & The name of the array for which the REGION variable
references. This can be any standard region array as declared in the
REGION section (FIPNUM, PVTNUM, etc.), provided the array exists at the
time the OPERATER keyword is invoked. In addition, the MULTNUM, FLUXNUM
and OPERNUM may be used. Only the default value of OPERNUM is supported
by OPM Flow. & OPERNUM \\
Notes: & & & \\
\end{longtable}

\emph{Table 6.96: OPERATER Keyword Description}

Examples of the arrays most commonly operated on in each section are
given in Table 6.97. Cells colored red indicate arrays that are not
supported by OPM Flow operations.

\begin{longtable}[]{@{}lllllll@{}}
\toprule\noalign{}
OPERATER Keyword and Variable Options by Section & & & & & & \\
\midrule\noalign{}
\endhead
\bottomrule\noalign{}
\endlastfoot
GRID & EDIT & PROPS & REGIONS & SOLUTION & SUMMARY & SCHEDULE \\
DX & DEPTH & SWL & ENDNUM & PRESSURE & & \\
DY & PORV & SWCR & EQLNUM & SWAT & & \\
DZ & TRANX & SWU & FIPNUM & SGAS & & \\
PERMX & TRANY & SGL & IMBNUM & RV & & \\
PERMY & TRANZ & SGCR & MISCNUM & RS & & \\
PERMZ & DIFFX & SGU & PVTNUM & TBLK & & \\
MULTX & DIFFY & KRW & ROCKNUM & GI & & \\
MULTY & DIFFZ & KRO & SATNUM & OILAPI & & \\
MULTZ & TRANR & KRG & WH2NUM & SALT & & \\
DR & TRANTHT & PCG & & GASCONC & & \\
DTHETA & DIFFR & PCW & & SOLVCONC & & \\
PERMR & DIFFTHT & & & SOLVFRAC & & \\
PERMTHT & & & & SFOAM & & \\
DZNET & & & & SPOLY & & \\
PORO & & & & & & \\
NTG & & & & & & \\
FLUXNUM & & & & & & \\
MULTNUM & & & & & & \\
MPFANUM & & & & & & \\
DIFFX & & & & & & \\
DIFFY & & & & & & \\
DIFFZ & & & & & & \\
DIFFR & & & & & & \\
DIFFTHT & & & & & & \\
\end{longtable}

\emph{Table 6.97: OPERATER Keyword Applicable Arrays by Section}

Double precision and integer work arrays may also be used to hold
temporary values for use in more complex operations. Note that these
work arrays are reset to zero at the start of each section of the data
file. The names of the double precision work arrays should be of the
form `WORKn', where n is a positive integer that defines the work array
index, for example `WORK1', `WORK2', and so on. Similarly, the names of
integer work arrays should be of the form `IWORKn', for example
`IWORK1', `IWORK2', and so on. OPM Flow currently only supports double
precision work arrays, and does not currently support using work arrays
with ``global'' arrays (MULTZ\emph{). Note that the permeability arrays
(PERM}) are treated as ``global'' arrays in the GRID section only. In
the commercial simulator the maximum number of double precision and
integer work arrays are specified using the REGDIMS keyword. Whereas,
OPM Flow does not impose a limit on the maximum number of work arrays.

Note that care should be exercised when performing operations on integer
array data as all transforms are performed using floating point
arithmetic operations. In addition, operations on any of the
transmissibility arrays (TRANX, TRANX-, TRANY, TRANY-, TRANZ, and
TRANZ-) may result in unintended consequences as these arrays have zero
values on the boundary of the grid. In this use OPM ResInsight to verify
and visually inspect the results.

\begin{tcolorbox}[enhanced jigsaw, title=\textcolor{quarto-callout-note-color}{\faInfo}\hspace{0.5em}{Note}, coltitle=black, breakable, toprule=.15mm, colbacktitle=quarto-callout-note-color!10!white, leftrule=.75mm, opacitybacktitle=0.6, colback=white, bottomtitle=1mm, toptitle=1mm, titlerule=0mm, left=2mm, colframe=quarto-callout-note-color-frame, arc=.35mm, rightrule=.15mm, bottomrule=.15mm, opacityback=0]

The OPERATER and OPERATE keywords allow for great flexibility in
generating or modifying the simulator's input arrays. In addition, OPM
Flow also has a Python facility to manipulate and calculate data that
offers even greater flexibility, but note that this feature is not
compatible with the commercial simulator. Finally, OPM ResInsight, the
post-processing plotting software, has both Octave and Python scripting
facilities that enable both data generation and visual inspection of the
results. The resulting calculated arrays can then be exported from OPM
ResInsight and ``included'' back into OPM Flow, thus maintaining
compatibility with the commercial simulator.

\end{tcolorbox}

\subsubsection{Example}\label{example-215}

The first example uses the MULTP function combined with the Net-to-Gross
(NTG) array to re-scale the MULTX, MULTY and MULTZ arrays to reduce the
transmissibility in three separate reservoirs based on the reservoir
quality (NTG). This keyword sequence should be in the GRID section.

\begin{verbatim}
--
--       MATHEMATICAL OPERATIONS ON ARRAYS BY REGION
--
--       OUTPUT   REGN   OPERATION   SOURCE   ALPHA   BETA    REGN
--       ARRAY    NUM    TYPE        ARRAY    CONST   CONST   ARRAY
OPERATER
         MULTX    1     'MULTP'      NTG      1.00    0.75        / RES1
         MULTY    1     'MULTP'      NTG      1.00    0.75        / RES1
         MULTZ    1     'MULTP'      NTG      1.00    0.75        / RES1

         MULTX    2     'MULTP'      NTG      1.00    0.85        / RES2
         MULTY    2     'MULTP'      NTG      1.00    0.85        / RES2
         MULTZ    2     'MULTP'      NTG      1.00    0.85        / RES2

         MULTX    3     'MULTP'      NTG      1.00    0.50        / RES3
         MULTY    3     'MULTP'      NTG      1.00    0.50        / RES3
         MULTZ    3     'MULTP'      NTG      1.00    0.50        / RES3
/
\end{verbatim}

Notice that the ARRAY variable has been defaulted, resulting in OPERNUM
being the regional array for the REGION variable.

The next example shows how to set the maximum gas saturation (SGU) based
on the minimum (lowest) water saturation (SWL) when using the End-Point
Scaling option, in the PROPS section.

\begin{verbatim}
--
--       MATHEMATICAL OPERATIONS ON ARRAYS
--
--       OUTPUT   REGN   OPERATION   SOURCE   ALPHA   BETA    REGN
--       ARRAY    NUM    TYPE        ARRAY    CONST   CONST   ARRAY
OPERATER
         SGU      1     'MULTA'      SWL      -1.0    1.0          /
         SGU      2     'MULTA'      SWL      -1.0    1.0          /
         SGU      3     'MULTA'      SWL      -1.0    1.0          /
/
\end{verbatim}

The above example sets the maximum gas saturation to be one minus the
minimum water saturation for regions one to three.

The final example shows how to reset the FIPNUM array when the exported
array from the earth model does not correspond to the simulator's
desired numbering scheme.

\begin{verbatim}
--
--       MATHEMATICAL OPERATIONS ON ARRAYS BY REGION
--
--       RESET FIPNUM BASED ON MULTNUM AND OPERNUM
--
--       DESTIN   REGN   OPERATION   SOURCE   ALPHA   BETA  INPUT  SEGNUM   EQUIL
--       ARRAY    NUM    TYPE        ARRAY    CONST   CONST ARRAY  NUMBER  NUMBER
OPERATER
         FIPNUM    26    'MULTA'   'MULTNUM'  0.00     1          /  26       1
         FIPNUM    44    'MULTA'   'MULTNUM'  0.00     2          /  44       2
         FIPNUM    62    'MULTA'   'MULTNUM'  0.00     3          /  62       3
         FIPNUM    98    'MULTA'   'MULTNUM'  0.00     4          /  98       4
         FIPNUM   116    'MULTA'   'MULTNUM'  0.00     5          / 116       5
         FIPNUM   134    'MULTA'   'MULTNUM'  0.00     6          / 134       6
         FIPNUM    46    'MULTA'   'MULTNUM'  0.00     7          /  46       7
         FIPNUM    64    'MULTA'   'MULTNUM'  0.00     8          /  64       8
         FIPNUM    82    'MULTA'   'MULTNUM'  0.00     9          /  82       9
         FIPNUM   226    'MULTA'   'MULTNUM'  0.00    10          / 226      10
         FIPNUM   262    'MULTA'   'MULTNUM'  0.00    11          / 262      11
         FIPNUM   280    'MULTA'   'MULTNUM'  0.00    12          / 280      12
         FIPNUM   298    'MULTA'   'MULTNUM'  0.00    13          / 298      13
         FIPNUM    33    'MULTA'   'MULTNUM'  0.00    14          /  33      14
         FIPNUM    51    'MULTA'   'MULTNUM'  0.00    15          /  51      15
         FIPNUM   105    'MULTA'   'MULTNUM'  0.00    16          / 105      16
         FIPNUM   159    'MULTA'   'MULTNUM'  0.00    17          / 159      17
         FIPNUM   195    'MULTA'   'MULTNUM'  0.00    18          / 195      18
         FIPNUM   267    'MULTA'   'MULTNUM'  0.00    19          / 267      19
         FIPNUM   303    'MULTA'   'MULTNUM'  0.00    20          / 303      20
         FIPNUM   321    'MULTA'   'MULTNUM'  0.00    21          / 321      21
         FIPNUM   339    'MULTA'   'MULTNUM'  0.00    22          / 339      22
         FIPNUM    54    'MULTA'   'MULTNUM'  0.00    23          /  54      23
         FIPNUM    72    'MULTA'   'MULTNUM'  0.00    24          /  72      24
         FIPNUM   108    'MULTA'   'MULTNUM'  0.00    25          / 108      25
         FIPNUM   144    'MULTA'   'MULTNUM'  0.00    26          / 144      26
         FIPNUM   270    'MULTA'   'MULTNUM'  0.00    27          / 270      27
/
\end{verbatim}

Note that operation can only be done in the REGION section as FIPNUM is
only available for use in this section and that the ARRAY variable has
been defaulted, resulting in OPERNUM being the regional array for the
REGION variable.

\subsection{OPERNUM -- Define Regions for Mathematical Operations on
Arrays}\label{opernum-define-regions-for-mathematical-operations-on-arrays}

\begin{longtable}[]{@{}
  >{\raggedright\arraybackslash}p{(\linewidth - 14\tabcolsep) * \real{0.1250}}
  >{\raggedright\arraybackslash}p{(\linewidth - 14\tabcolsep) * \real{0.1250}}
  >{\raggedright\arraybackslash}p{(\linewidth - 14\tabcolsep) * \real{0.1250}}
  >{\raggedright\arraybackslash}p{(\linewidth - 14\tabcolsep) * \real{0.1250}}
  >{\raggedright\arraybackslash}p{(\linewidth - 14\tabcolsep) * \real{0.1250}}
  >{\raggedright\arraybackslash}p{(\linewidth - 14\tabcolsep) * \real{0.1250}}
  >{\raggedright\arraybackslash}p{(\linewidth - 14\tabcolsep) * \real{0.1250}}
  >{\raggedright\arraybackslash}p{(\linewidth - 14\tabcolsep) * \real{0.1250}}@{}}
\toprule\noalign{}
\begin{minipage}[b]{\linewidth}\raggedright
RUNSPEC
\end{minipage} & \begin{minipage}[b]{\linewidth}\raggedright
GRID
\end{minipage} & \begin{minipage}[b]{\linewidth}\raggedright
EDIT
\end{minipage} & \begin{minipage}[b]{\linewidth}\raggedright
PROPS
\end{minipage} & \begin{minipage}[b]{\linewidth}\raggedright
REGIONS
\end{minipage} & \begin{minipage}[b]{\linewidth}\raggedright
SOLUTION
\end{minipage} & \begin{minipage}[b]{\linewidth}\raggedright
SUMMARY
\end{minipage} & \begin{minipage}[b]{\linewidth}\raggedright
SCHEDULE
\end{minipage} \\
\midrule\noalign{}
\endhead
\bottomrule\noalign{}
\endlastfoot
\end{longtable}

\subsubsection{Description}\label{description-316}

This keyword defines the OPERATER region numbers for each grid block.
The OPERNUM keyword defines the region numbers for each grid block, as
such there must be one entry for each cell in the model. The array can
also be used with the EQUALREG, ADDREG, COPYREG, MULTIREG, MULTREGP and
MULTREGT keywords, as well as the OPERATER keyword in calculating
various grid properties in the GRID and REGION section.

\begin{longtable}[]{@{}
  >{\raggedright\arraybackslash}p{(\linewidth - 6\tabcolsep) * \real{0.2500}}
  >{\raggedright\arraybackslash}p{(\linewidth - 6\tabcolsep) * \real{0.2500}}
  >{\raggedright\arraybackslash}p{(\linewidth - 6\tabcolsep) * \real{0.2500}}
  >{\raggedright\arraybackslash}p{(\linewidth - 6\tabcolsep) * \real{0.2500}}@{}}
\toprule\noalign{}
\begin{minipage}[b]{\linewidth}\raggedright
No.
\end{minipage} & \begin{minipage}[b]{\linewidth}\raggedright
Name
\end{minipage} & \begin{minipage}[b]{\linewidth}\raggedright
Description
\end{minipage} & \begin{minipage}[b]{\linewidth}\raggedright
Default
\end{minipage} \\
\midrule\noalign{}
\endhead
\bottomrule\noalign{}
\endlastfoot
1 & OPERNUM & OPERNUM defines an array of positive integers greater than
or equal to one that assigns a grid cell to a particular OPERNUM region.
The maximum number of OPERNUM regions is set by the NOPREG variable on
the REGDIMS keyword in the RUNSPEC section. Note that the default value
of zero implies that the calculations requested by the OPERATER keyword
will not be performed. & 0 \\
Notes: & & & \\
\end{longtable}

\emph{Table 6.98: OPERNUM Keyword Description}

\subsubsection{Examples}\label{examples-14}

The example below sets three OPERNUM regions for a 4 x 5 x 2 model.

\begin{verbatim}
--
--       DEFINE OPERNUM REGIONS FOR ALL CELLS
--
OPERNUM
         2  2  1  1  2  2  1  1  1  1  1  1  1  1  1  1  1  1  1  1
         3  3  1  1  3  3  1  1  1  1  1  1  1  1  1  1  1  1  1  1
/
\end{verbatim}

Alternatively the EQUALS keyword could be employed to accomplish the
same task, that is:

\begin{verbatim}
--
--       ARRAY       CONSTANT     ---------- BOX ---------
--                                I1  I2   J1  J2   K1  K2
EQUALS
         OPERNUM     1            1*  1*   1*  1*   1*  1* / SET REGION 1
         OPERNUM     2            1   2    1   2    1   1  / SET REGION 2
         OPERNUM     3            1   2    1   2    2   2  / SET REGION 3
/
\end{verbatim}

One can then increase PERMX by 25\% in region three only.

\begin{verbatim}
--
--       MULTIPLY AN ARRAY BY A CONSTANT BASED ON A REGION NUMBER
--
--       ARRAY     CONSTANT  REGION   REGION ARRAY
--                 VALUE     NUMBER    M / F / O
MULTIREG
         PERMX     1.25      3         O                   /
/
\end{verbatim}

\subsection{OUTRAD -- Define the Outer Radius of a Radial
Grid}\label{outrad-define-the-outer-radius-of-a-radial-grid}

\begin{longtable}[]{@{}
  >{\raggedright\arraybackslash}p{(\linewidth - 14\tabcolsep) * \real{0.1250}}
  >{\raggedright\arraybackslash}p{(\linewidth - 14\tabcolsep) * \real{0.1250}}
  >{\raggedright\arraybackslash}p{(\linewidth - 14\tabcolsep) * \real{0.1250}}
  >{\raggedright\arraybackslash}p{(\linewidth - 14\tabcolsep) * \real{0.1250}}
  >{\raggedright\arraybackslash}p{(\linewidth - 14\tabcolsep) * \real{0.1250}}
  >{\raggedright\arraybackslash}p{(\linewidth - 14\tabcolsep) * \real{0.1250}}
  >{\raggedright\arraybackslash}p{(\linewidth - 14\tabcolsep) * \real{0.1250}}
  >{\raggedright\arraybackslash}p{(\linewidth - 14\tabcolsep) * \real{0.1250}}@{}}
\toprule\noalign{}
\begin{minipage}[b]{\linewidth}\raggedright
RUNSPEC
\end{minipage} & \begin{minipage}[b]{\linewidth}\raggedright
GRID
\end{minipage} & \begin{minipage}[b]{\linewidth}\raggedright
EDIT
\end{minipage} & \begin{minipage}[b]{\linewidth}\raggedright
PROPS
\end{minipage} & \begin{minipage}[b]{\linewidth}\raggedright
REGIONS
\end{minipage} & \begin{minipage}[b]{\linewidth}\raggedright
SOLUTION
\end{minipage} & \begin{minipage}[b]{\linewidth}\raggedright
SUMMARY
\end{minipage} & \begin{minipage}[b]{\linewidth}\raggedright
SCHEDULE
\end{minipage} \\
\midrule\noalign{}
\endhead
\bottomrule\noalign{}
\endlastfoot
\end{longtable}

\subsubsection{Description}\label{description-317}

OUTRAD defines the OUTER radius of the reservoir model for a radial or
spider grid geometry. The RADIAL or SPIDER keyword in the RUNSPEC should
be activated to indicate that radial or spider geometry is being used.

This keyword is not supported by OPM Flow but would change the results
if supported so the simulation will be stopped.

\begin{longtable}[]{@{}
  >{\raggedright\arraybackslash}p{(\linewidth - 6\tabcolsep) * \real{0.2500}}
  >{\raggedright\arraybackslash}p{(\linewidth - 6\tabcolsep) * \real{0.2500}}
  >{\raggedright\arraybackslash}p{(\linewidth - 6\tabcolsep) * \real{0.2500}}
  >{\raggedright\arraybackslash}p{(\linewidth - 6\tabcolsep) * \real{0.2500}}@{}}
\toprule\noalign{}
\begin{minipage}[b]{\linewidth}\raggedright
No.
\end{minipage} & \begin{minipage}[b]{\linewidth}\raggedright
Name
\end{minipage} & \begin{minipage}[b]{\linewidth}\raggedright
Description
\end{minipage} & \begin{minipage}[b]{\linewidth}\raggedright
Default
\end{minipage} \\
\midrule\noalign{}
\endhead
\bottomrule\noalign{}
\endlastfoot
Field & Metric & Laboratory & \\
1 & OUTRAD & A single real positive number greater than INRAD defining
the outer radius of a radial grid. & None \\
feet & m & cm & \\
Notes: & & & \\
\end{longtable}

\emph{Table 6.99: OUTRAD Keyword Description}

Note that the SPIDER keyword activates OPM Flow's radial grid geometry
option for the model. This keyword will create a spiderweb-shaped grid
based on a corner-point grid using the standard radial grid keywords:
INRAD, DRV, DTHETAV, DZ/DZV etc. in the GRID the section. A spider grid
can be viewed in 3D in OPM ResInsight unlike radial grids that cannot be
viewed in the software. To overcome this, the simulator now converts
radial grids to Irregular Corner-Point Grids and adjusts the model's
pore volume to reflect radial coordinates; thus, overcoming the display
limitation.

The keyword allows for an alternative method of entering the size of the
R direction grid cells instead of entering the data using the DR or DRV
keywords in the GRID section. Given the internal radius set by the INRAD
keyword, the external radius set by the OUTRAD keyword and the number of
grid cells in the R direction set by the NX variable on the DIMENS
keyword in the RUNSPEC section, the R direction cells sizes are computed
automatically on a geometric spacing, as defined by:

\begin{equation}\phantomsection\label{eq-6-12}{
\frac{{R}_{i}}{{R}_{i-1}} = \left(\frac{\mathit{OUTRAD}}{{R}_{{i}_{j}-1}}\right){ }^{\frac{1}{\left(\mathit{NX} - {i}_{j} + 1\right)}}
}\end{equation}

or

\begin{equation}\phantomsection\label{eq-6-13}{
{R}_{i} = \left({R}_{{i}_{j}-1}\right)\left(\frac{\mathit{OUTRAD}}{{R}_{{i}_{j}-1}}\right){ }^{\frac{\left(i - {i}_{j} +1 \right)}{\left(\mathit{NX} - {i}_{j} + 1\right)}}
}\end{equation}

and the DR value for the ith cell, that is the value that can also be
manually entered on the DR keyword in the GRID section, is given by:

\begin{equation}\phantomsection\label{eq-6-14}{
{\mathit{DR}}_{i} = {R}_{i} - {R}_{i - 1}
}\end{equation}

Where:

DRi = DR value for the ith cell

Ri = current total radius for the i radii.

Ri-1 = total radius for the i -- 1 radii.

NX (NR) = number of radial grid cells excluding the inner radius

OUTRAD = the outer radius of the radial grid, the value includes the
inner radius.

For example, given an inner radius set to 0.25, an outer radius of 2,050
and the number of cells in the R direction set to ten, then Table 6.100
shows the grid size calculations.

\begin{longtable}[]{@{}llll@{}}
\toprule\noalign{}
OUTRAD Radial Grid Example & & & \\
\midrule\noalign{}
\endhead
\bottomrule\noalign{}
\endlastfoot
INRAD & 0.25 & & \\
OUTRAD & 2050.0 & & \\
NX & 10 & & \\
NX & Ri & DR & Ratio \\
0 & 0.250 & 0.250 & \\
1 & 0.616 & 0.366 & 1.463 \\
2 & 1.516 & 0.900 & 2.463 \\
3 & 3.733 & 2.217 & 2.463 \\
4 & 9.193 & 5.460 & 2.463 \\
5 & 22.638 & 13.445 & 2.463 \\
6 & 55.748 & 33.109 & 2.463 \\
7 & 137.281 & 81.533 & 2.463 \\
8 & 338.058 & 200.777 & 2.463 \\
9 & 832.477 & 494.420 & 2.463 \\
10 & 2050.000 & 1217.523 & 2.463 \\
Total & & 2050.000 & \\
\end{longtable}

\emph{Table 6.100: OUTRAD Radial Grid Example}

See also the DR, DRV, DTHETA, DTHETAV and TOPS keywords to fully define
a Radial or Spider Grid.

\subsubsection{Example}\label{example-216}

\begin{verbatim}
--
--       INNER RADIUS OF FIRST GRID BLOCK IN THE RADIAL DIRECTION
--
INRAD
         0.25                                                                  /
--
--       OUTER RADIUS OF FIRST GRID BLOCK IN THE RADIAL DIRECTION
--
OUTRAD
         2050.0                                                                /
\end{verbatim}

The above example defines the inner radius to be 0.25 and the outer
radius to be 2,050 feet. Note that the outer radius includes the inner
radius.

\subsection{PARAOPTS -- Define Parallel Run
Options}\label{paraopts-define-parallel-run-options}

\begin{longtable}[]{@{}
  >{\raggedright\arraybackslash}p{(\linewidth - 14\tabcolsep) * \real{0.1250}}
  >{\raggedright\arraybackslash}p{(\linewidth - 14\tabcolsep) * \real{0.1250}}
  >{\raggedright\arraybackslash}p{(\linewidth - 14\tabcolsep) * \real{0.1250}}
  >{\raggedright\arraybackslash}p{(\linewidth - 14\tabcolsep) * \real{0.1250}}
  >{\raggedright\arraybackslash}p{(\linewidth - 14\tabcolsep) * \real{0.1250}}
  >{\raggedright\arraybackslash}p{(\linewidth - 14\tabcolsep) * \real{0.1250}}
  >{\raggedright\arraybackslash}p{(\linewidth - 14\tabcolsep) * \real{0.1250}}
  >{\raggedright\arraybackslash}p{(\linewidth - 14\tabcolsep) * \real{0.1250}}@{}}
\toprule\noalign{}
\begin{minipage}[b]{\linewidth}\raggedright
RUNSPEC
\end{minipage} & \begin{minipage}[b]{\linewidth}\raggedright
GRID
\end{minipage} & \begin{minipage}[b]{\linewidth}\raggedright
EDIT
\end{minipage} & \begin{minipage}[b]{\linewidth}\raggedright
PROPS
\end{minipage} & \begin{minipage}[b]{\linewidth}\raggedright
REGIONS
\end{minipage} & \begin{minipage}[b]{\linewidth}\raggedright
SOLUTION
\end{minipage} & \begin{minipage}[b]{\linewidth}\raggedright
SUMMARY
\end{minipage} & \begin{minipage}[b]{\linewidth}\raggedright
SCHEDULE
\end{minipage} \\
\midrule\noalign{}
\endhead
\bottomrule\noalign{}
\endlastfoot
\end{longtable}

\subsubsection{Description}\label{description-318}

The PARAOPTS keyword defines various options for parallel runs, for when
the Parallel option has been invoked by the PARALLEL keyword in the
RUNSPEC section.

This keyword is not supported by OPM Flow but has no effect on the
results so it will be ignored.

See section 2.2 Running OPM Flow 2023-04 From The Command Line on how to
run OPM Flow in parallel mode.

\subsection{PEBI -- Activate and Defined PEBI Grid
Options}\label{pebi-activate-and-defined-pebi-grid-options}

\begin{longtable}[]{@{}
  >{\raggedright\arraybackslash}p{(\linewidth - 14\tabcolsep) * \real{0.1250}}
  >{\raggedright\arraybackslash}p{(\linewidth - 14\tabcolsep) * \real{0.1250}}
  >{\raggedright\arraybackslash}p{(\linewidth - 14\tabcolsep) * \real{0.1250}}
  >{\raggedright\arraybackslash}p{(\linewidth - 14\tabcolsep) * \real{0.1250}}
  >{\raggedright\arraybackslash}p{(\linewidth - 14\tabcolsep) * \real{0.1250}}
  >{\raggedright\arraybackslash}p{(\linewidth - 14\tabcolsep) * \real{0.1250}}
  >{\raggedright\arraybackslash}p{(\linewidth - 14\tabcolsep) * \real{0.1250}}
  >{\raggedright\arraybackslash}p{(\linewidth - 14\tabcolsep) * \real{0.1250}}@{}}
\toprule\noalign{}
\begin{minipage}[b]{\linewidth}\raggedright
RUNSPEC
\end{minipage} & \begin{minipage}[b]{\linewidth}\raggedright
GRID
\end{minipage} & \begin{minipage}[b]{\linewidth}\raggedright
EDIT
\end{minipage} & \begin{minipage}[b]{\linewidth}\raggedright
PROPS
\end{minipage} & \begin{minipage}[b]{\linewidth}\raggedright
REGIONS
\end{minipage} & \begin{minipage}[b]{\linewidth}\raggedright
SOLUTION
\end{minipage} & \begin{minipage}[b]{\linewidth}\raggedright
SUMMARY
\end{minipage} & \begin{minipage}[b]{\linewidth}\raggedright
SCHEDULE
\end{minipage} \\
\midrule\noalign{}
\endhead
\bottomrule\noalign{}
\endlastfoot
\end{longtable}

\subsubsection{Description}\label{description-319}

PEBI activates the unstructured Perpendicular Bisector (``PEBI'')
{[}Heinemann, Z.E. and Brand, C.W. 1988. Gridding Techniques in
Reservoir Simulation. Proc., First Intl. Forum on Reservoir Simulation,
Alpbach, Austria, 339.{]} and {[}Heinemann, Z.E., Brand, C.W., Munka, M.
et al.~1991. Modeling Reservoir Geometry With Irregular Grids. SPE Res
Eng 6 (2): 225--232. SPE-18412-PA. http://dx.doi.org/10.2118/18412-PA{]}
loading of grid data generated by an external pre-processing program for
generating simulation grids.

This keyword is not supported by OPM Flow but would change the results
if supported so the simulation will be stopped.

\begin{longtable}[]{@{}
  >{\raggedright\arraybackslash}p{(\linewidth - 6\tabcolsep) * \real{0.2500}}
  >{\raggedright\arraybackslash}p{(\linewidth - 6\tabcolsep) * \real{0.2500}}
  >{\raggedright\arraybackslash}p{(\linewidth - 6\tabcolsep) * \real{0.2500}}
  >{\raggedright\arraybackslash}p{(\linewidth - 6\tabcolsep) * \real{0.2500}}@{}}
\toprule\noalign{}
\begin{minipage}[b]{\linewidth}\raggedright
No.
\end{minipage} & \begin{minipage}[b]{\linewidth}\raggedright
Name
\end{minipage} & \begin{minipage}[b]{\linewidth}\raggedright
Description
\end{minipage} & \begin{minipage}[b]{\linewidth}\raggedright
Default
\end{minipage} \\
\midrule\noalign{}
\endhead
\bottomrule\noalign{}
\endlastfoot
1 & OPTION1 & A defined character string that activates or deactivates
the checking of negative transmissibility values. OPTION1 should be set
to YES to check for negative values, or NO switches off this option. &
NO \\
2 & OPTION2 & A defined character string that activates or deactivates
the calculation of pore volumes and transmissibilities. OPTION2 should
be set to YES if the pore volumes and transmissibilities are provided,
or NO for the values to calculated by the simulator. & NO \\
\end{longtable}

\emph{Table 6.101: PEBI Keyword Description}

\subsubsection{Example}\label{example-217}

\begin{verbatim}
–-
--      OPTION1  OPTION2
--      CHECK    CALCULATE
PEBI
        NO       YES                                                           /
\end{verbatim}

The above example switches off the negative transmissibility check and
requests that the simulator calculates pore volumes and
transmissibilities as they are not provided by the input data.

\subsection{PERMAVE -- Define Average Transmissibility
Coefficients}\label{permave-define-average-transmissibility-coefficients}

\begin{longtable}[]{@{}
  >{\raggedright\arraybackslash}p{(\linewidth - 14\tabcolsep) * \real{0.1250}}
  >{\raggedright\arraybackslash}p{(\linewidth - 14\tabcolsep) * \real{0.1250}}
  >{\raggedright\arraybackslash}p{(\linewidth - 14\tabcolsep) * \real{0.1250}}
  >{\raggedright\arraybackslash}p{(\linewidth - 14\tabcolsep) * \real{0.1250}}
  >{\raggedright\arraybackslash}p{(\linewidth - 14\tabcolsep) * \real{0.1250}}
  >{\raggedright\arraybackslash}p{(\linewidth - 14\tabcolsep) * \real{0.1250}}
  >{\raggedright\arraybackslash}p{(\linewidth - 14\tabcolsep) * \real{0.1250}}
  >{\raggedright\arraybackslash}p{(\linewidth - 14\tabcolsep) * \real{0.1250}}@{}}
\toprule\noalign{}
\begin{minipage}[b]{\linewidth}\raggedright
RUNSPEC
\end{minipage} & \begin{minipage}[b]{\linewidth}\raggedright
GRID
\end{minipage} & \begin{minipage}[b]{\linewidth}\raggedright
EDIT
\end{minipage} & \begin{minipage}[b]{\linewidth}\raggedright
PROPS
\end{minipage} & \begin{minipage}[b]{\linewidth}\raggedright
REGIONS
\end{minipage} & \begin{minipage}[b]{\linewidth}\raggedright
SOLUTION
\end{minipage} & \begin{minipage}[b]{\linewidth}\raggedright
SUMMARY
\end{minipage} & \begin{minipage}[b]{\linewidth}\raggedright
SCHEDULE
\end{minipage} \\
\midrule\noalign{}
\endhead
\bottomrule\noalign{}
\endlastfoot
\end{longtable}

\subsubsection{Description}\label{description-320}

The PERMAVE keyword defines the three directional exponent coefficients
used to average the grid block permeabilities between two neighboring
cells when calculating the transmissibility between the two blocks. The
keyword can be used to change from the default weighted harmonic
averaging (coefficient set equal to -1), to geometric (coefficient equal
to zero), or to arithmetic averaging (coefficient equal to 1). The three
coefficients represent the averaging in the x-, y- and z-directions.

This keyword is not supported by OPM Flow but would change the results
if supported so the simulation will be stopped.

\subsection{PERMJFUN -- Define Leverett J-Function Permeability for All
Cells}\label{permjfun-define-leverett-j-function-permeability-for-all-cells}

\begin{longtable}[]{@{}
  >{\raggedright\arraybackslash}p{(\linewidth - 14\tabcolsep) * \real{0.1250}}
  >{\raggedright\arraybackslash}p{(\linewidth - 14\tabcolsep) * \real{0.1250}}
  >{\raggedright\arraybackslash}p{(\linewidth - 14\tabcolsep) * \real{0.1250}}
  >{\raggedright\arraybackslash}p{(\linewidth - 14\tabcolsep) * \real{0.1250}}
  >{\raggedright\arraybackslash}p{(\linewidth - 14\tabcolsep) * \real{0.1250}}
  >{\raggedright\arraybackslash}p{(\linewidth - 14\tabcolsep) * \real{0.1250}}
  >{\raggedright\arraybackslash}p{(\linewidth - 14\tabcolsep) * \real{0.1250}}
  >{\raggedright\arraybackslash}p{(\linewidth - 14\tabcolsep) * \real{0.1250}}@{}}
\toprule\noalign{}
\begin{minipage}[b]{\linewidth}\raggedright
RUNSPEC
\end{minipage} & \begin{minipage}[b]{\linewidth}\raggedright
GRID
\end{minipage} & \begin{minipage}[b]{\linewidth}\raggedright
EDIT
\end{minipage} & \begin{minipage}[b]{\linewidth}\raggedright
PROPS
\end{minipage} & \begin{minipage}[b]{\linewidth}\raggedright
REGIONS
\end{minipage} & \begin{minipage}[b]{\linewidth}\raggedright
SOLUTION
\end{minipage} & \begin{minipage}[b]{\linewidth}\raggedright
SUMMARY
\end{minipage} & \begin{minipage}[b]{\linewidth}\raggedright
SCHEDULE
\end{minipage} \\
\midrule\noalign{}
\endhead
\bottomrule\noalign{}
\endlastfoot
\end{longtable}

\subsubsection{Description}\label{description-321}

PERMJFUN defines the permeability to be used in de-normalizing the
Leverett J-Functions {[}Leverett, M. C.; ``Capillary Behaviour in Porous
Solids'', Trans. AIME (1941) 142, 152-168.{]} for when the PERM variable
on the JFUNC or the JFUNCR keyword in the GRID section has been set to
``U'', as oppose to using PERMX, PERMY, PERMZ arrays etc.

This keyword is not supported by OPM Flow but would change the results
if supported so the simulation will be stopped.

\begin{longtable}[]{@{}
  >{\raggedright\arraybackslash}p{(\linewidth - 6\tabcolsep) * \real{0.2500}}
  >{\raggedright\arraybackslash}p{(\linewidth - 6\tabcolsep) * \real{0.2500}}
  >{\raggedright\arraybackslash}p{(\linewidth - 6\tabcolsep) * \real{0.2500}}
  >{\raggedright\arraybackslash}p{(\linewidth - 6\tabcolsep) * \real{0.2500}}@{}}
\toprule\noalign{}
\begin{minipage}[b]{\linewidth}\raggedright
No.
\end{minipage} & \begin{minipage}[b]{\linewidth}\raggedright
Name
\end{minipage} & \begin{minipage}[b]{\linewidth}\raggedright
Description
\end{minipage} & \begin{minipage}[b]{\linewidth}\raggedright
Default
\end{minipage} \\
\midrule\noalign{}
\endhead
\bottomrule\noalign{}
\endlastfoot
Field & Metric & Laboratory & \\
1 & PERMJFUN & PERMJFUN is an array of real positive numbers assigning
the permeability to be used in de-normalizing the Leverett J-Function to
each cell in the model. Repeat counts may be used, for example 20*100.0.
& None \\
mD & mD & mD & \\
Notes: & & & \\
\end{longtable}

\emph{Table 6.102: PERMJFUN Keyword Description}

For grid blocks that have not been assigned a PERMJFUN value the default
directional permeability will be used, that is the average of PERMX and
PERMY.

See also the PERMX, PERMY and PERMZ keywords to fully define the
permeability for the model.

\subsubsection{Example}\label{example-218}

\begin{verbatim}
--
--       DEFINE GRID BLOCK PERMJFUN FOR ALL CELLS (BASED ON NX x NY x NZ = 300)
--
PERMJFUN
         100*500.0   100*50.0   100*200.0                                      /
\end{verbatim}

The above example defines the PERMJFUN to be 500.0, 50.0, and 200.0 for
the first, second and third layers in the model for all 300 cells, as
defined by the DIMENS keyword in the RUNSPEC section.

\subsection{PERMR -- Define the Permeability for Each Cell in the R
Direction}\label{permr-define-the-permeability-for-each-cell-in-the-r-direction}

\begin{longtable}[]{@{}
  >{\raggedright\arraybackslash}p{(\linewidth - 14\tabcolsep) * \real{0.1250}}
  >{\raggedright\arraybackslash}p{(\linewidth - 14\tabcolsep) * \real{0.1250}}
  >{\raggedright\arraybackslash}p{(\linewidth - 14\tabcolsep) * \real{0.1250}}
  >{\raggedright\arraybackslash}p{(\linewidth - 14\tabcolsep) * \real{0.1250}}
  >{\raggedright\arraybackslash}p{(\linewidth - 14\tabcolsep) * \real{0.1250}}
  >{\raggedright\arraybackslash}p{(\linewidth - 14\tabcolsep) * \real{0.1250}}
  >{\raggedright\arraybackslash}p{(\linewidth - 14\tabcolsep) * \real{0.1250}}
  >{\raggedright\arraybackslash}p{(\linewidth - 14\tabcolsep) * \real{0.1250}}@{}}
\toprule\noalign{}
\begin{minipage}[b]{\linewidth}\raggedright
RUNSPEC
\end{minipage} & \begin{minipage}[b]{\linewidth}\raggedright
GRID
\end{minipage} & \begin{minipage}[b]{\linewidth}\raggedright
EDIT
\end{minipage} & \begin{minipage}[b]{\linewidth}\raggedright
PROPS
\end{minipage} & \begin{minipage}[b]{\linewidth}\raggedright
REGIONS
\end{minipage} & \begin{minipage}[b]{\linewidth}\raggedright
SOLUTION
\end{minipage} & \begin{minipage}[b]{\linewidth}\raggedright
SUMMARY
\end{minipage} & \begin{minipage}[b]{\linewidth}\raggedright
SCHEDULE
\end{minipage} \\
\midrule\noalign{}
\endhead
\bottomrule\noalign{}
\endlastfoot
\end{longtable}

\subsubsection{Description}\label{description-322}

PERMR sets the permeability for each cell in the R direction in a radial
geometry grid. The RADIAL or SPIDER keywords in the RUNSPEC should be
activated to indicate that radial or spider geometry is being used.

\begin{longtable}[]{@{}
  >{\raggedright\arraybackslash}p{(\linewidth - 6\tabcolsep) * \real{0.2500}}
  >{\raggedright\arraybackslash}p{(\linewidth - 6\tabcolsep) * \real{0.2500}}
  >{\raggedright\arraybackslash}p{(\linewidth - 6\tabcolsep) * \real{0.2500}}
  >{\raggedright\arraybackslash}p{(\linewidth - 6\tabcolsep) * \real{0.2500}}@{}}
\toprule\noalign{}
\begin{minipage}[b]{\linewidth}\raggedright
No.
\end{minipage} & \begin{minipage}[b]{\linewidth}\raggedright
Name
\end{minipage} & \begin{minipage}[b]{\linewidth}\raggedright
Description
\end{minipage} & \begin{minipage}[b]{\linewidth}\raggedright
Default
\end{minipage} \\
\midrule\noalign{}
\endhead
\bottomrule\noalign{}
\endlastfoot
Field & Metric & Laboratory & \\
1 & PERMR & PERMR is an array of real positive numbers assigning the
permeability in the R direction to each cell in the model. This
equivalent to PERMX in a Cartesian grid. Repeat counts may be used, for
example 20*100.0. & None \\
mD & mD & mD & \\
Notes: & & & \\
\end{longtable}

\emph{Table 6.103: PERMR Keyword Description}

Note that the SPIDER keyword activates OPM Flow's radial grid geometry
option for the model. This keyword will create a spiderweb-shaped grid
based on a corner-point grid using the standard radial grid keywords:
INRAD, DRV, DTHETAV, DZ/DZV etc. in the GRID the section. A spider grid
can be viewed in 3D in OPM ResInsight unlike radial grids that cannot be
viewed in the software. To overcome this, the simulator now converts
radial grids to Irregular Corner-Point Grids and adjusts the model's
pore volume to reflect radial coordinates; thus, overcoming the display
limitation.

See also the PERMTHT and PERMZ keywords to fully define the permeability
for a radial or spider grid model.

\subsubsection{Example}\label{example-219}

\begin{verbatim}
--
--       DEFINE GRID BLOCK PERMR DATA FOR ALL CELLS (BASED ON NR x NY x NZ = 300)
--
PERMR
         100*500.0   100*50.0   100*200.0                                       /
\end{verbatim}

The above example defines the PERMR to be 500.0, 50.0, and 200.0 for the
first, second and third layers in the model for all 300 cells, as
defined by the DIMENS keyword (10, 10, 3) in the RUNSPEC section.

\subsection{PERMTHT -- Define the Permeability for Each Cell in the
THETA
Direction}\label{permtht-define-the-permeability-for-each-cell-in-the-theta-direction}

\begin{longtable}[]{@{}
  >{\raggedright\arraybackslash}p{(\linewidth - 14\tabcolsep) * \real{0.1250}}
  >{\raggedright\arraybackslash}p{(\linewidth - 14\tabcolsep) * \real{0.1250}}
  >{\raggedright\arraybackslash}p{(\linewidth - 14\tabcolsep) * \real{0.1250}}
  >{\raggedright\arraybackslash}p{(\linewidth - 14\tabcolsep) * \real{0.1250}}
  >{\raggedright\arraybackslash}p{(\linewidth - 14\tabcolsep) * \real{0.1250}}
  >{\raggedright\arraybackslash}p{(\linewidth - 14\tabcolsep) * \real{0.1250}}
  >{\raggedright\arraybackslash}p{(\linewidth - 14\tabcolsep) * \real{0.1250}}
  >{\raggedright\arraybackslash}p{(\linewidth - 14\tabcolsep) * \real{0.1250}}@{}}
\toprule\noalign{}
\begin{minipage}[b]{\linewidth}\raggedright
RUNSPEC
\end{minipage} & \begin{minipage}[b]{\linewidth}\raggedright
GRID
\end{minipage} & \begin{minipage}[b]{\linewidth}\raggedright
EDIT
\end{minipage} & \begin{minipage}[b]{\linewidth}\raggedright
PROPS
\end{minipage} & \begin{minipage}[b]{\linewidth}\raggedright
REGIONS
\end{minipage} & \begin{minipage}[b]{\linewidth}\raggedright
SOLUTION
\end{minipage} & \begin{minipage}[b]{\linewidth}\raggedright
SUMMARY
\end{minipage} & \begin{minipage}[b]{\linewidth}\raggedright
SCHEDULE
\end{minipage} \\
\midrule\noalign{}
\endhead
\bottomrule\noalign{}
\endlastfoot
\end{longtable}

\subsubsection{Description}\label{description-323}

PERMTHT sets the permeability for each cell in the THETA direction in a
radial geometry grid. The RADIAL or SPIDER keyword in the RUNSPEC should
be activated to indicate that radial or spider geometry is being used.

\begin{longtable}[]{@{}
  >{\raggedright\arraybackslash}p{(\linewidth - 6\tabcolsep) * \real{0.2500}}
  >{\raggedright\arraybackslash}p{(\linewidth - 6\tabcolsep) * \real{0.2500}}
  >{\raggedright\arraybackslash}p{(\linewidth - 6\tabcolsep) * \real{0.2500}}
  >{\raggedright\arraybackslash}p{(\linewidth - 6\tabcolsep) * \real{0.2500}}@{}}
\toprule\noalign{}
\begin{minipage}[b]{\linewidth}\raggedright
No.
\end{minipage} & \begin{minipage}[b]{\linewidth}\raggedright
Name
\end{minipage} & \begin{minipage}[b]{\linewidth}\raggedright
Description
\end{minipage} & \begin{minipage}[b]{\linewidth}\raggedright
Default
\end{minipage} \\
\midrule\noalign{}
\endhead
\bottomrule\noalign{}
\endlastfoot
Field & Metric & Laboratory & \\
1 & PERMTHT & PERMTHT is an array of real positive numbers assigning the
permeability in the THETA direction to each cell in the model. This
equivalent to PERMY in a Cartesian grid. Repeat counts may be used, for
example 20*100.0. & None \\
mD & mD & mD & \\
Notes: & & & \\
\end{longtable}

\emph{Table 6.104: PERMTHT Keyword Description}

Note that the SPIDER keyword activates OPM Flow's radial grid geometry
option for the model. This keyword will create a spiderweb-shaped grid
based on a corner-point grid using the standard radial grid keywords:
INRAD, DRV, DTHETAV, DZ/DZV etc. in the GRID the section. A spider grid
can be viewed in 3D in OPM ResInsight unlike radial grids that cannot be
viewed in the software. To overcome this, the simulator now converts
radial grids to Irregular Corner-Point Grids and adjusts the model's
pore volume to reflect radial coordinates; thus, overcoming the display
limitation.

See also the PERMR and PERMZ keywords to fully define the permeability
for a radial or spider grid model.

\subsubsection{Example}\label{example-220}

\begin{verbatim}
--
--       DEFINE GRID BLOCK PERMTHT DATA FOR ALL CELLS (NX x NY x NZ = 300)
--
PERMTHT
         100*500.0   100*50.0   100*200.0                                       /
\end{verbatim}

The above example defines the PERMTHT to be 500.0, 50.0, and 200.0 for
the first, second and third layers in the model for all 300 cells, as
defined by the DIMENS keyword (10, 10, 3) in the RUNSPEC section.

\subsection{PERMX -- Define the Permeability in the X Direction for All
the
Cells}\label{permx-define-the-permeability-in-the-x-direction-for-all-the-cells}

\begin{longtable}[]{@{}
  >{\raggedright\arraybackslash}p{(\linewidth - 14\tabcolsep) * \real{0.1250}}
  >{\raggedright\arraybackslash}p{(\linewidth - 14\tabcolsep) * \real{0.1250}}
  >{\raggedright\arraybackslash}p{(\linewidth - 14\tabcolsep) * \real{0.1250}}
  >{\raggedright\arraybackslash}p{(\linewidth - 14\tabcolsep) * \real{0.1250}}
  >{\raggedright\arraybackslash}p{(\linewidth - 14\tabcolsep) * \real{0.1250}}
  >{\raggedright\arraybackslash}p{(\linewidth - 14\tabcolsep) * \real{0.1250}}
  >{\raggedright\arraybackslash}p{(\linewidth - 14\tabcolsep) * \real{0.1250}}
  >{\raggedright\arraybackslash}p{(\linewidth - 14\tabcolsep) * \real{0.1250}}@{}}
\toprule\noalign{}
\begin{minipage}[b]{\linewidth}\raggedright
RUNSPEC
\end{minipage} & \begin{minipage}[b]{\linewidth}\raggedright
GRID
\end{minipage} & \begin{minipage}[b]{\linewidth}\raggedright
EDIT
\end{minipage} & \begin{minipage}[b]{\linewidth}\raggedright
PROPS
\end{minipage} & \begin{minipage}[b]{\linewidth}\raggedright
REGIONS
\end{minipage} & \begin{minipage}[b]{\linewidth}\raggedright
SOLUTION
\end{minipage} & \begin{minipage}[b]{\linewidth}\raggedright
SUMMARY
\end{minipage} & \begin{minipage}[b]{\linewidth}\raggedright
SCHEDULE
\end{minipage} \\
\midrule\noalign{}
\endhead
\bottomrule\noalign{}
\endlastfoot
\end{longtable}

\subsubsection{Description}\label{description-324}

PERMX defines the permeability in the X direction for all the cells in
the model via an array. The keyword can be used for all grid types,
except for the Radial Grid geometry.

\begin{longtable}[]{@{}
  >{\raggedright\arraybackslash}p{(\linewidth - 6\tabcolsep) * \real{0.2500}}
  >{\raggedright\arraybackslash}p{(\linewidth - 6\tabcolsep) * \real{0.2500}}
  >{\raggedright\arraybackslash}p{(\linewidth - 6\tabcolsep) * \real{0.2500}}
  >{\raggedright\arraybackslash}p{(\linewidth - 6\tabcolsep) * \real{0.2500}}@{}}
\toprule\noalign{}
\begin{minipage}[b]{\linewidth}\raggedright
No.
\end{minipage} & \begin{minipage}[b]{\linewidth}\raggedright
Name
\end{minipage} & \begin{minipage}[b]{\linewidth}\raggedright
Description
\end{minipage} & \begin{minipage}[b]{\linewidth}\raggedright
Default
\end{minipage} \\
\midrule\noalign{}
\endhead
\bottomrule\noalign{}
\endlastfoot
Field & Metric & Laboratory & \\
1 & PERMX & PERMX is an array of real positive numbers assigning the
permeability in the X direction to each cell in the model. Repeat counts
may be used, for example 20*100.0. & None \\
mD & mD & mD & \\
Notes: & & & \\
\end{longtable}

\emph{Table 6.105: PERMX Keyword Description}

See also the PERMY and PERMZ keywords to fully define the permeability
for the model.

\subsubsection{Example}\label{example-221}

\begin{verbatim}
--
--       DEFINE GRID BLOCK PERMX DATA FOR ALL CELLS (BASED ON NX x NY x NZ = 300)
--
PERMX
         100*500.0   100*50.0   100*200.0                                       /
\end{verbatim}

The above example defines the PERMX to be 500.0, 50.0, and 200.0 for the
first, second and third layers in the model for all 300 cells, as
defined by the DIMENS keyword in the RUNSPEC section.

\subsection{PERMY -- Define the Permeability in the Y Direction for All
the
Cells}\label{permy-define-the-permeability-in-the-y-direction-for-all-the-cells}

\begin{longtable}[]{@{}
  >{\raggedright\arraybackslash}p{(\linewidth - 14\tabcolsep) * \real{0.1250}}
  >{\raggedright\arraybackslash}p{(\linewidth - 14\tabcolsep) * \real{0.1250}}
  >{\raggedright\arraybackslash}p{(\linewidth - 14\tabcolsep) * \real{0.1250}}
  >{\raggedright\arraybackslash}p{(\linewidth - 14\tabcolsep) * \real{0.1250}}
  >{\raggedright\arraybackslash}p{(\linewidth - 14\tabcolsep) * \real{0.1250}}
  >{\raggedright\arraybackslash}p{(\linewidth - 14\tabcolsep) * \real{0.1250}}
  >{\raggedright\arraybackslash}p{(\linewidth - 14\tabcolsep) * \real{0.1250}}
  >{\raggedright\arraybackslash}p{(\linewidth - 14\tabcolsep) * \real{0.1250}}@{}}
\toprule\noalign{}
\begin{minipage}[b]{\linewidth}\raggedright
RUNSPEC
\end{minipage} & \begin{minipage}[b]{\linewidth}\raggedright
GRID
\end{minipage} & \begin{minipage}[b]{\linewidth}\raggedright
EDIT
\end{minipage} & \begin{minipage}[b]{\linewidth}\raggedright
PROPS
\end{minipage} & \begin{minipage}[b]{\linewidth}\raggedright
REGIONS
\end{minipage} & \begin{minipage}[b]{\linewidth}\raggedright
SOLUTION
\end{minipage} & \begin{minipage}[b]{\linewidth}\raggedright
SUMMARY
\end{minipage} & \begin{minipage}[b]{\linewidth}\raggedright
SCHEDULE
\end{minipage} \\
\midrule\noalign{}
\endhead
\bottomrule\noalign{}
\endlastfoot
\end{longtable}

\subsubsection{Description}\label{description-325}

PERMY defines the permeability in the Y direction for all the cells in
the model via an array. The keyword can be used for all grid types,
except for the Radial Grid geometry.

\begin{longtable}[]{@{}
  >{\raggedright\arraybackslash}p{(\linewidth - 6\tabcolsep) * \real{0.2500}}
  >{\raggedright\arraybackslash}p{(\linewidth - 6\tabcolsep) * \real{0.2500}}
  >{\raggedright\arraybackslash}p{(\linewidth - 6\tabcolsep) * \real{0.2500}}
  >{\raggedright\arraybackslash}p{(\linewidth - 6\tabcolsep) * \real{0.2500}}@{}}
\toprule\noalign{}
\begin{minipage}[b]{\linewidth}\raggedright
No.
\end{minipage} & \begin{minipage}[b]{\linewidth}\raggedright
Name
\end{minipage} & \begin{minipage}[b]{\linewidth}\raggedright
Description
\end{minipage} & \begin{minipage}[b]{\linewidth}\raggedright
Default
\end{minipage} \\
\midrule\noalign{}
\endhead
\bottomrule\noalign{}
\endlastfoot
Field & Metric & Laboratory & \\
1 & PERMY & PERMY is an array of real positive numbers assigning the
permeability in the Y direction to each cell in the model. Repeat counts
may be used, for example 20*100.0. & None \\
mD & mD & mD & \\
Notes: & & & \\
\end{longtable}

\emph{Table 6.106: PERMY Keyword Description}

See also the PERMX and PERMZ keywords to fully define the permeability
for the model.

\subsubsection{Example}\label{example-222}

\begin{verbatim}
--
--       DEFINE GRID BLOCK PERMY DATA FOR ALL CELLS (BASED ON NX x NY x NZ = 300)
--
PERMY
         100*500.0   100*50.0   100*200.0                                       /
\end{verbatim}

The above example defines the PERMY to be 500.0, 50.0, and 200.0 for the
first, second and third layers in the model for all 300 cells, as
defined by the DIMENS keyword in the RUNSPEC section.

\subsection{PERMZ -- Define the Permeability in the Z Direction for All
the
Cells}\label{permz-define-the-permeability-in-the-z-direction-for-all-the-cells}

\begin{longtable}[]{@{}
  >{\raggedright\arraybackslash}p{(\linewidth - 14\tabcolsep) * \real{0.1250}}
  >{\raggedright\arraybackslash}p{(\linewidth - 14\tabcolsep) * \real{0.1250}}
  >{\raggedright\arraybackslash}p{(\linewidth - 14\tabcolsep) * \real{0.1250}}
  >{\raggedright\arraybackslash}p{(\linewidth - 14\tabcolsep) * \real{0.1250}}
  >{\raggedright\arraybackslash}p{(\linewidth - 14\tabcolsep) * \real{0.1250}}
  >{\raggedright\arraybackslash}p{(\linewidth - 14\tabcolsep) * \real{0.1250}}
  >{\raggedright\arraybackslash}p{(\linewidth - 14\tabcolsep) * \real{0.1250}}
  >{\raggedright\arraybackslash}p{(\linewidth - 14\tabcolsep) * \real{0.1250}}@{}}
\toprule\noalign{}
\begin{minipage}[b]{\linewidth}\raggedright
RUNSPEC
\end{minipage} & \begin{minipage}[b]{\linewidth}\raggedright
GRID
\end{minipage} & \begin{minipage}[b]{\linewidth}\raggedright
EDIT
\end{minipage} & \begin{minipage}[b]{\linewidth}\raggedright
PROPS
\end{minipage} & \begin{minipage}[b]{\linewidth}\raggedright
REGIONS
\end{minipage} & \begin{minipage}[b]{\linewidth}\raggedright
SOLUTION
\end{minipage} & \begin{minipage}[b]{\linewidth}\raggedright
SUMMARY
\end{minipage} & \begin{minipage}[b]{\linewidth}\raggedright
SCHEDULE
\end{minipage} \\
\midrule\noalign{}
\endhead
\bottomrule\noalign{}
\endlastfoot
\end{longtable}

\subsubsection{Description}\label{description-326}

PERMZ defines the permeability in the Z direction for all the cells in
the model via an array. The keyword can be used for all grid types,
except for the Radial Grid geometry.

\begin{longtable}[]{@{}
  >{\raggedright\arraybackslash}p{(\linewidth - 6\tabcolsep) * \real{0.2500}}
  >{\raggedright\arraybackslash}p{(\linewidth - 6\tabcolsep) * \real{0.2500}}
  >{\raggedright\arraybackslash}p{(\linewidth - 6\tabcolsep) * \real{0.2500}}
  >{\raggedright\arraybackslash}p{(\linewidth - 6\tabcolsep) * \real{0.2500}}@{}}
\toprule\noalign{}
\begin{minipage}[b]{\linewidth}\raggedright
No.
\end{minipage} & \begin{minipage}[b]{\linewidth}\raggedright
Name
\end{minipage} & \begin{minipage}[b]{\linewidth}\raggedright
Description
\end{minipage} & \begin{minipage}[b]{\linewidth}\raggedright
Default
\end{minipage} \\
\midrule\noalign{}
\endhead
\bottomrule\noalign{}
\endlastfoot
Field & Metric & Laboratory & \\
1 & PERMZ & PERMZ is an array of real positive numbers assigning the
permeability in the Z direction to each cell in the model. Repeat counts
may be used, for example 200*50.0. & None \\
mD & mD & mD & \\
Notes: & & & \\
\end{longtable}

\emph{Table 6.107: PERMZ Keyword Description}

\begin{tcolorbox}[enhanced jigsaw, title=\textcolor{quarto-callout-note-color}{\faInfo}\hspace{0.5em}{Note}, coltitle=black, breakable, toprule=.15mm, colbacktitle=quarto-callout-note-color!10!white, leftrule=.75mm, opacitybacktitle=0.6, colback=white, bottomtitle=1mm, toptitle=1mm, titlerule=0mm, left=2mm, colframe=quarto-callout-note-color-frame, arc=.35mm, rightrule=.15mm, bottomrule=.15mm, opacityback=0]

Although PERMX and PERMY are commonly set to be equal, PERMZ is
typically not equal to either PERMX or PERMY. Normally PERMZ is set as a
fraction of PERMX with typical values ranging from 0.1 to 0.5 times
PERMX.

\end{tcolorbox}

See also the PERMX and PERMY keywords to fully define the permeability
for the model.

\subsubsection{Example}\label{example-223}

The example below defines the PERMZ to be 50.0, 5.0, and 20.0 for the
first, second and third layers in the model for all 300 cells, as
defined by the DIMENS keyword in the RUNSPEC section.

\begin{verbatim}
--
--       DEFINE GRID BLOCK PERMZ DATA FOR ALL CELLS (BASED ON NX x NY x NZ = 300)
--
PERMZ
         100*50.0    100*5.0   100*20.0                                         /
\end{verbatim}

The next example sets PERMX to be 500.0, 50.0, and 200.0 for the first,
second and third layers in the model for all 300 cells, as defined by
the DIMENS keyword in the RUNSPEC section. It then copies the PERMX
values to the PERMY and PERMZ arrays, and finally multiplies PERMZ by
0.1 times to get the final values for PERMZ.

\begin{verbatim}
--
--       DEFINE GRID BLOCK PERMX DATA FOR ALL CELLS
--
PERMX
         100*500.0   100*50.0   100*200.0                                      /
--
--       SOURCE      DESTIN.      ---------- BOX ---------
--                                I1  I2   J1  J2   K1  K2
COPY
         PERMX       PERMY        1*  1*   1*  1*   1*  1* / CREATE PERMY
         PERMX       PERMZ        1*  1*   1*  1*   1*  1* / CREATE PERMZ
/
--
--       ARRAY       CONSTANT     ---------- BOX ---------
--                                I1  I2   J1  J2   K1  K2
MULTIPLY
         PERMZ       0.10000      1*  1*   1*  1*   1*  1* / PERMZ * 0.1
/
\end{verbatim}

The above sequence of keywords is quite common in input decks, that is
copying the PERMX data to the PERMY and PERMZ arrays and then adjusting
the PERMY and PERMZ arrays as required using the MULTIPLY keyword.

\subsection{PETGRID -- Load a Generic Simulation Grid
File}\label{petgrid-load-a-generic-simulation-grid-file}

\begin{longtable}[]{@{}
  >{\raggedright\arraybackslash}p{(\linewidth - 14\tabcolsep) * \real{0.1250}}
  >{\raggedright\arraybackslash}p{(\linewidth - 14\tabcolsep) * \real{0.1250}}
  >{\raggedright\arraybackslash}p{(\linewidth - 14\tabcolsep) * \real{0.1250}}
  >{\raggedright\arraybackslash}p{(\linewidth - 14\tabcolsep) * \real{0.1250}}
  >{\raggedright\arraybackslash}p{(\linewidth - 14\tabcolsep) * \real{0.1250}}
  >{\raggedright\arraybackslash}p{(\linewidth - 14\tabcolsep) * \real{0.1250}}
  >{\raggedright\arraybackslash}p{(\linewidth - 14\tabcolsep) * \real{0.1250}}
  >{\raggedright\arraybackslash}p{(\linewidth - 14\tabcolsep) * \real{0.1250}}@{}}
\toprule\noalign{}
\begin{minipage}[b]{\linewidth}\raggedright
RUNSPEC
\end{minipage} & \begin{minipage}[b]{\linewidth}\raggedright
GRID
\end{minipage} & \begin{minipage}[b]{\linewidth}\raggedright
EDIT
\end{minipage} & \begin{minipage}[b]{\linewidth}\raggedright
PROPS
\end{minipage} & \begin{minipage}[b]{\linewidth}\raggedright
REGIONS
\end{minipage} & \begin{minipage}[b]{\linewidth}\raggedright
SOLUTION
\end{minipage} & \begin{minipage}[b]{\linewidth}\raggedright
SUMMARY
\end{minipage} & \begin{minipage}[b]{\linewidth}\raggedright
SCHEDULE
\end{minipage} \\
\midrule\noalign{}
\endhead
\bottomrule\noalign{}
\endlastfoot
\end{longtable}

\subsubsection{Description}\label{description-327}

The PETGRID keyword instructs the simulator to load a Generic Simulation
Grid (*.GSG) file that contains grid geometry data.

This keyword is not supported by OPM Flow but would change the results
if supported so the simulation will be stopped.

\subsection{PINCH -- Define Pinch-Out Layer
Options}\label{pinch-define-pinch-out-layer-options}

\begin{longtable}[]{@{}
  >{\raggedright\arraybackslash}p{(\linewidth - 14\tabcolsep) * \real{0.1250}}
  >{\raggedright\arraybackslash}p{(\linewidth - 14\tabcolsep) * \real{0.1250}}
  >{\raggedright\arraybackslash}p{(\linewidth - 14\tabcolsep) * \real{0.1250}}
  >{\raggedright\arraybackslash}p{(\linewidth - 14\tabcolsep) * \real{0.1250}}
  >{\raggedright\arraybackslash}p{(\linewidth - 14\tabcolsep) * \real{0.1250}}
  >{\raggedright\arraybackslash}p{(\linewidth - 14\tabcolsep) * \real{0.1250}}
  >{\raggedright\arraybackslash}p{(\linewidth - 14\tabcolsep) * \real{0.1250}}
  >{\raggedright\arraybackslash}p{(\linewidth - 14\tabcolsep) * \real{0.1250}}@{}}
\toprule\noalign{}
\begin{minipage}[b]{\linewidth}\raggedright
RUNSPEC
\end{minipage} & \begin{minipage}[b]{\linewidth}\raggedright
GRID
\end{minipage} & \begin{minipage}[b]{\linewidth}\raggedright
EDIT
\end{minipage} & \begin{minipage}[b]{\linewidth}\raggedright
PROPS
\end{minipage} & \begin{minipage}[b]{\linewidth}\raggedright
REGIONS
\end{minipage} & \begin{minipage}[b]{\linewidth}\raggedright
SOLUTION
\end{minipage} & \begin{minipage}[b]{\linewidth}\raggedright
SUMMARY
\end{minipage} & \begin{minipage}[b]{\linewidth}\raggedright
SCHEDULE
\end{minipage} \\
\midrule\noalign{}
\endhead
\bottomrule\noalign{}
\endlastfoot
\end{longtable}

\subsubsection{Description}\label{description-328}

The PINCH keyword defines the parameters used to control the generation
of Non-Neighbor Connections (``NNCs'') in the vertical (K) direction due
to layers pinching out. This keyword is applied to all layers in the
model as opposed to the PINCHREG keyword that offers more flexibility by
applying the pinch-out controls to various regions in the model defined
by the PINCHNUM keyword.

OPM Flow will automatically generate connections between non neighbor
cells in the vertical direction based on the parameters on this keyword.

\begin{longtable}[]{@{}
  >{\raggedright\arraybackslash}p{(\linewidth - 6\tabcolsep) * \real{0.2500}}
  >{\raggedright\arraybackslash}p{(\linewidth - 6\tabcolsep) * \real{0.2500}}
  >{\raggedright\arraybackslash}p{(\linewidth - 6\tabcolsep) * \real{0.2500}}
  >{\raggedright\arraybackslash}p{(\linewidth - 6\tabcolsep) * \real{0.2500}}@{}}
\toprule\noalign{}
\begin{minipage}[b]{\linewidth}\raggedright
No.
\end{minipage} & \begin{minipage}[b]{\linewidth}\raggedright
Name
\end{minipage} & \begin{minipage}[b]{\linewidth}\raggedright
Description
\end{minipage} & \begin{minipage}[b]{\linewidth}\raggedright
Default
\end{minipage} \\
\midrule\noalign{}
\endhead
\bottomrule\noalign{}
\endlastfoot
Field & Metric & Laboratory & \\
1 & PINCHTHK & A real number defining the pinch-out threshold thickness
for any cell. NNCs are generated across inactive cells having a vertical
thickness less than PINCHTHK. & Defined \\
ft. 0.001 & m 0.001 & cm 0.001 & \\
2 & PINCHOPT & A character string controlling the generation of
pinch-outs when the MINPV keyword has been used to deactivate cells with
small pore volumes. PINCHOPT can either be set to: & GAP \\
3 & PINCHGAP & A real number defining the maximum ``empty'' thickness
allowed between grid blocks in adjacent grid layers for a non-zero
transmissibility to exist between them. & Defined \\
ft. 1.0E20 & m 1.0E20 & cm 1.0E20 & \\
4 & PINCHCAL & A character string controlling the calculation of the
pinch-out transmissibilities. PINCHCAL can either be set to: & TOPBOT \\
5 & PINCHMUL & A character string controlling the calculation of the
pinch-out transmissibilities when adjustments have been made by the
MULTZ keyword. PINCHMUL can either be set to: Note if PINCHCAL has been
set equal to ALL then PINCHMUL is reset to TOP, irrespective of the
entered value for PINCHMUL. & TOP \\
Notes: & & & \\
\end{longtable}

\emph{Table 6.108: PINCH Keyword Description}

\subsubsection{Examples}\label{examples-15}

The first example below will create NNCs between the cells above and
below any cell having vertical thickness less than 0.01 in either feet
or metres.

\begin{verbatim}
--
--       SET PINCH-OUT PARAMETERS FOR CALCULATING PINCH-OUT PROPERTIES
--
PINCH
--       THRESHOLD   GAP      EMPTY   TRANS    MULTZ
--       THICKNESS   NO GAP   GAP     CALC     CALC
         0.01        1*       1*      1*       1*                              /
\end{verbatim}

For the second example, the MINPV keyword is used to set the minimum
pore volume to 500 m3 (metric units) and then the PINCH keyword is
invoked with PINCHGAP set equal to GAP, as follows:

\begin{verbatim}
--
--       MINIMUM PORE VOLUME FOR ACTIVE CELLS
--
MINPV
         500.0                                                                 /
--
--       SET PINCH-OUT CRITERIA FOR THE MODEL
--
PINCH
--       THRESHOLD   GAP      EMPTY   TRANS    MULTZ
--       THICKNESS   NO GAP   GAP     CALC     CALC
         0.1         GAP       1*     1*       1*                              /
\end{verbatim}

In the above example the MINPV keyword will deactivate all cells with
pore volumes less than 500 m3. These deactivated cells are inactive in
the model and therefore are not included in the flow calculations;
however, by default they will result in no-flow barriers but may not be
thin enough for PINCH to create NNCs across them. By setting PINCHGAP
equal to GAP on the PINCH keyword (the default setting), then OPM Flow
generates NNCs across the cells that have been deactivated by the MINPV
keyword. However, in this case there may be grid blocks in the model
with a pore volume greater than MINPV but a thickness less than the
pinch-out threshold. These cells will not be deactivated by the PINCH
keyword.

\subsection{PINCHNUM -- Define Pinch-Out Regions for the PINCHREG
Keyword}\label{pinchnum-define-pinch-out-regions-for-the-pinchreg-keyword}

\begin{longtable}[]{@{}
  >{\raggedright\arraybackslash}p{(\linewidth - 14\tabcolsep) * \real{0.1250}}
  >{\raggedright\arraybackslash}p{(\linewidth - 14\tabcolsep) * \real{0.1250}}
  >{\raggedright\arraybackslash}p{(\linewidth - 14\tabcolsep) * \real{0.1250}}
  >{\raggedright\arraybackslash}p{(\linewidth - 14\tabcolsep) * \real{0.1250}}
  >{\raggedright\arraybackslash}p{(\linewidth - 14\tabcolsep) * \real{0.1250}}
  >{\raggedright\arraybackslash}p{(\linewidth - 14\tabcolsep) * \real{0.1250}}
  >{\raggedright\arraybackslash}p{(\linewidth - 14\tabcolsep) * \real{0.1250}}
  >{\raggedright\arraybackslash}p{(\linewidth - 14\tabcolsep) * \real{0.1250}}@{}}
\toprule\noalign{}
\begin{minipage}[b]{\linewidth}\raggedright
RUNSPEC
\end{minipage} & \begin{minipage}[b]{\linewidth}\raggedright
GRID
\end{minipage} & \begin{minipage}[b]{\linewidth}\raggedright
EDIT
\end{minipage} & \begin{minipage}[b]{\linewidth}\raggedright
PROPS
\end{minipage} & \begin{minipage}[b]{\linewidth}\raggedright
REGIONS
\end{minipage} & \begin{minipage}[b]{\linewidth}\raggedright
SOLUTION
\end{minipage} & \begin{minipage}[b]{\linewidth}\raggedright
SUMMARY
\end{minipage} & \begin{minipage}[b]{\linewidth}\raggedright
SCHEDULE
\end{minipage} \\
\midrule\noalign{}
\endhead
\bottomrule\noalign{}
\endlastfoot
\end{longtable}

\subsubsection{Description}\label{description-329}

The PINCHNUM keyword defines the pinch-out region numbers for each grid
block, as such there must be one entry for each cell in the model. The
array is used with the PINCHREG keyword to set the pinch-out options and
threshold thickness for each region.

This keyword is not supported by OPM Flow but would change the results
if supported so the simulation will be stopped.

\begin{longtable}[]{@{}
  >{\raggedright\arraybackslash}p{(\linewidth - 6\tabcolsep) * \real{0.2500}}
  >{\raggedright\arraybackslash}p{(\linewidth - 6\tabcolsep) * \real{0.2500}}
  >{\raggedright\arraybackslash}p{(\linewidth - 6\tabcolsep) * \real{0.2500}}
  >{\raggedright\arraybackslash}p{(\linewidth - 6\tabcolsep) * \real{0.2500}}@{}}
\toprule\noalign{}
\begin{minipage}[b]{\linewidth}\raggedright
No.
\end{minipage} & \begin{minipage}[b]{\linewidth}\raggedright
Name
\end{minipage} & \begin{minipage}[b]{\linewidth}\raggedright
Description
\end{minipage} & \begin{minipage}[b]{\linewidth}\raggedright
Default
\end{minipage} \\
\midrule\noalign{}
\endhead
\bottomrule\noalign{}
\endlastfoot
1 & PINCHNUM & PINCHNUM defines an array of positive integers assigning
a grid cell to a particular PINCHNUM region. The maximum number of
PINCHNUM regions is set by the NRPINC variable on the GRIDOPTS keyword
in the RUNSPEC section. & 1 \\
Notes: & & & \\
\end{longtable}

\emph{Table 6.109: PINCHNUM Keyword Description}

\subsubsection{Example}\label{example-224}

The example below sets defines three PINCHNUM regions for various layers
in a model based on the model's layering.

\begin{verbatim}
--
--       ARRAY       CONSTANT     ---------- BOX ---------
--                                I1  I2   J1  J2   K1  K2
EQUALS
         MULTNUM     1            1*  1*   1*  1*   1*  1* / SET REGION 1
         MULTNUM     2            1   2    1   2    10  50 / SET REGION 2
         MULTNUM     3            1   2    1   2    51 100 / SET REGION 3
/
\end{verbatim}

One can then set the pinch-out criteria for each region.

\begin{verbatim}
--
--       SET PINCH-OUT CRITERA VIA THE PINCHNUM REGION
--
PINCHREG
--       THRESHOLD   GAP      EMPTY   TRANS
--       THICKNESS   NO GAP   GAP     CALC
         0.1         1*       1*      1*                  / PINCHNUM 01
         1.0         1*       10      1*                  / PINCHNUM 02
         1.0         NOGAP    20      1*                  / PINCHNUM 03
\end{verbatim}

The above example sets the default pinch-out criteria for grid blocks
defined as region one via the PINCHNUM array and different criteria for
regions two and three.

\subsection{PINCHOUT -- Define Pinch-Out Layers Option
(Fixed)}\label{pinchout-define-pinch-out-layers-option-fixed}

\begin{longtable}[]{@{}
  >{\raggedright\arraybackslash}p{(\linewidth - 14\tabcolsep) * \real{0.1250}}
  >{\raggedright\arraybackslash}p{(\linewidth - 14\tabcolsep) * \real{0.1250}}
  >{\raggedright\arraybackslash}p{(\linewidth - 14\tabcolsep) * \real{0.1250}}
  >{\raggedright\arraybackslash}p{(\linewidth - 14\tabcolsep) * \real{0.1250}}
  >{\raggedright\arraybackslash}p{(\linewidth - 14\tabcolsep) * \real{0.1250}}
  >{\raggedright\arraybackslash}p{(\linewidth - 14\tabcolsep) * \real{0.1250}}
  >{\raggedright\arraybackslash}p{(\linewidth - 14\tabcolsep) * \real{0.1250}}
  >{\raggedright\arraybackslash}p{(\linewidth - 14\tabcolsep) * \real{0.1250}}@{}}
\toprule\noalign{}
\begin{minipage}[b]{\linewidth}\raggedright
RUNSPEC
\end{minipage} & \begin{minipage}[b]{\linewidth}\raggedright
GRID
\end{minipage} & \begin{minipage}[b]{\linewidth}\raggedright
EDIT
\end{minipage} & \begin{minipage}[b]{\linewidth}\raggedright
PROPS
\end{minipage} & \begin{minipage}[b]{\linewidth}\raggedright
REGIONS
\end{minipage} & \begin{minipage}[b]{\linewidth}\raggedright
SOLUTION
\end{minipage} & \begin{minipage}[b]{\linewidth}\raggedright
SUMMARY
\end{minipage} & \begin{minipage}[b]{\linewidth}\raggedright
SCHEDULE
\end{minipage} \\
\midrule\noalign{}
\endhead
\bottomrule\noalign{}
\endlastfoot
\end{longtable}

\subsubsection{Description}\label{description-330}

The PINCHOUT keyword activates the generation of Non-Neighbor
Connections (``NNCs'') in the vertical (K) direction due to layers
pinching out, using a constant threshold thickness of 0.001 for all unit
systems. See also the PINCH keyword in the GRID section that allows for
specifying the threshold thickness and other parameters on a layer
basis, and the PINCHREG keyword that applies the pinch-out controls to
various regions in the model defined by the PINCHNUM keyword.

There is no data required for this keyword and there is no terminating
``/'' for this keyword.

This keyword is not supported by OPM Flow but would change the results
if supported so the simulation will be stopped.

\subsubsection{Example}\label{example-225}

The example will create NNCs between the cells above and below any cell
having vertical thickness less than 0.001 in either feet or metres.

\begin{verbatim}
--
--       SET PINCH-OUT CRITERA WITH CONSTANT THRESHOLD THICKNESS OF 0.001
--
PINCHOUT
\end{verbatim}

\subsection{PINCHREG -- Define Pinch-Out Region
Options}\label{pinchreg-define-pinch-out-region-options}

\begin{longtable}[]{@{}
  >{\raggedright\arraybackslash}p{(\linewidth - 14\tabcolsep) * \real{0.1250}}
  >{\raggedright\arraybackslash}p{(\linewidth - 14\tabcolsep) * \real{0.1250}}
  >{\raggedright\arraybackslash}p{(\linewidth - 14\tabcolsep) * \real{0.1250}}
  >{\raggedright\arraybackslash}p{(\linewidth - 14\tabcolsep) * \real{0.1250}}
  >{\raggedright\arraybackslash}p{(\linewidth - 14\tabcolsep) * \real{0.1250}}
  >{\raggedright\arraybackslash}p{(\linewidth - 14\tabcolsep) * \real{0.1250}}
  >{\raggedright\arraybackslash}p{(\linewidth - 14\tabcolsep) * \real{0.1250}}
  >{\raggedright\arraybackslash}p{(\linewidth - 14\tabcolsep) * \real{0.1250}}@{}}
\toprule\noalign{}
\begin{minipage}[b]{\linewidth}\raggedright
RUNSPEC
\end{minipage} & \begin{minipage}[b]{\linewidth}\raggedright
GRID
\end{minipage} & \begin{minipage}[b]{\linewidth}\raggedright
EDIT
\end{minipage} & \begin{minipage}[b]{\linewidth}\raggedright
PROPS
\end{minipage} & \begin{minipage}[b]{\linewidth}\raggedright
REGIONS
\end{minipage} & \begin{minipage}[b]{\linewidth}\raggedright
SOLUTION
\end{minipage} & \begin{minipage}[b]{\linewidth}\raggedright
SUMMARY
\end{minipage} & \begin{minipage}[b]{\linewidth}\raggedright
SCHEDULE
\end{minipage} \\
\midrule\noalign{}
\endhead
\bottomrule\noalign{}
\endlastfoot
\end{longtable}

\subsubsection{Description}\label{description-331}

The PINCHREG keyword defines the parameters used to control the
generation of Non-Neighbor Connections (``NNCs'') in the vertical (K)
direction due to layers pinching out in combination with the PINCHNUM
keyword. This allows different regions in the model to use different
criteria in controlling the how pinch-outs are generated. The keyword
should contain NRPINC records defining the criteria for each pinch-out
region defined with the PINCHNUM keyword. NRPINC is the maximum number
of PINCHNUM regions defined via the GRIDOPTS keyword in the RUNSPEC
section.

An alternative method to set the pinch-out criteria is to use the PINCH
keyword, that applies the criteria to the whole model.

OPM Flow will automatically generate connections between non neighbor
cells in the vertical direction based on the parameters on this keyword.

\begin{longtable}[]{@{}
  >{\raggedright\arraybackslash}p{(\linewidth - 6\tabcolsep) * \real{0.2500}}
  >{\raggedright\arraybackslash}p{(\linewidth - 6\tabcolsep) * \real{0.2500}}
  >{\raggedright\arraybackslash}p{(\linewidth - 6\tabcolsep) * \real{0.2500}}
  >{\raggedright\arraybackslash}p{(\linewidth - 6\tabcolsep) * \real{0.2500}}@{}}
\toprule\noalign{}
\begin{minipage}[b]{\linewidth}\raggedright
No.
\end{minipage} & \begin{minipage}[b]{\linewidth}\raggedright
Name
\end{minipage} & \begin{minipage}[b]{\linewidth}\raggedright
Description
\end{minipage} & \begin{minipage}[b]{\linewidth}\raggedright
Default
\end{minipage} \\
\midrule\noalign{}
\endhead
\bottomrule\noalign{}
\endlastfoot
Field & Metric & Laboratory & \\
1 & PINCHTHK & A real number defining the pinch-out threshold thickness
for any cell. NNCs are generated across inactive cells having a vertical
thickness less than PINCHTHK. & Defined \\
ft. 0.001 & m 0.001 & cm 0.001 & \\
2 & PINCHOPT & A character string controlling the generation of
pinch-outs when the MINPV keyword has been used to deactivate cells with
small pore volumes. PINCHOPT can either be set to: & GAP \\
3 & PINCHGAP & A real number defining the maximum ``empty'' thickness
allowed between grid blocks in adjacent grid layers for a non-zero
transmissibility to exist between them. & Defined \\
ft. 1.0E20 & m 1.0E20 & cm 1.0E20 & \\
4 & PINCHCAL & A character string controlling the calculation of the
pinch-out transmissibilities. PINCHCAL can either be set to: & TOPBOT \\
5 & PINCHMUL & A character string controlling the calculation of the
pinch-out transmissibilities when adjustments have been made by the
MULTZ keyword. PINCHMUL can either be set to: Note if PINCHCAL has been
set equal to ALL then PINCHMUL is reset to TOP, irrespective of the
entered value for PINCHMUL. & TOP \\
Notes: & & & \\
\end{longtable}

\emph{Table 6.110: PINCHREG Keyword Description}

\subsubsection{Example}\label{example-226}

\begin{verbatim}
--
--       SET PINCH-OUT CRITERA VIA THE PINCHNUM REGION
--
PINCHREG
--       THRESHOLD   GAP      EMPTY   TRANS
--       THICKNESS   NO GAP   GAP     CALC
         0.1         1*       1*      1*                  / PINCHNUM 01
         1.0         1*       10      1*                  / PINCHNUM 02
         1.0         NOGAP    20      1*                  / PINCHNUM 03
\end{verbatim}

The above example sets the default pinch-out criteria for grid blocks
defined as region one via the PINCHNUM array and different values for
regions two and three.

\subsection{PINCHXY -- Define Pinch-Out Areal
Options}\label{pinchxy-define-pinch-out-areal-options}

\begin{longtable}[]{@{}
  >{\raggedright\arraybackslash}p{(\linewidth - 14\tabcolsep) * \real{0.1250}}
  >{\raggedright\arraybackslash}p{(\linewidth - 14\tabcolsep) * \real{0.1250}}
  >{\raggedright\arraybackslash}p{(\linewidth - 14\tabcolsep) * \real{0.1250}}
  >{\raggedright\arraybackslash}p{(\linewidth - 14\tabcolsep) * \real{0.1250}}
  >{\raggedright\arraybackslash}p{(\linewidth - 14\tabcolsep) * \real{0.1250}}
  >{\raggedright\arraybackslash}p{(\linewidth - 14\tabcolsep) * \real{0.1250}}
  >{\raggedright\arraybackslash}p{(\linewidth - 14\tabcolsep) * \real{0.1250}}
  >{\raggedright\arraybackslash}p{(\linewidth - 14\tabcolsep) * \real{0.1250}}@{}}
\toprule\noalign{}
\begin{minipage}[b]{\linewidth}\raggedright
RUNSPEC
\end{minipage} & \begin{minipage}[b]{\linewidth}\raggedright
GRID
\end{minipage} & \begin{minipage}[b]{\linewidth}\raggedright
EDIT
\end{minipage} & \begin{minipage}[b]{\linewidth}\raggedright
PROPS
\end{minipage} & \begin{minipage}[b]{\linewidth}\raggedright
REGIONS
\end{minipage} & \begin{minipage}[b]{\linewidth}\raggedright
SOLUTION
\end{minipage} & \begin{minipage}[b]{\linewidth}\raggedright
SUMMARY
\end{minipage} & \begin{minipage}[b]{\linewidth}\raggedright
SCHEDULE
\end{minipage} \\
\midrule\noalign{}
\endhead
\bottomrule\noalign{}
\endlastfoot
\end{longtable}

\subsubsection{Description}\label{description-332}

The PINCHXY keyword defines the x-direction and y-direction threshold
thickness used to control the generation of Non-Neighbor Connections
(``NNCs'') in the x- and y- directions for missing cells in the areal
plane.

This keyword is not supported by OPM Flow but would change the results
if supported so the simulation will be stopped.

\begin{longtable}[]{@{}
  >{\raggedright\arraybackslash}p{(\linewidth - 6\tabcolsep) * \real{0.2500}}
  >{\raggedright\arraybackslash}p{(\linewidth - 6\tabcolsep) * \real{0.2500}}
  >{\raggedright\arraybackslash}p{(\linewidth - 6\tabcolsep) * \real{0.2500}}
  >{\raggedright\arraybackslash}p{(\linewidth - 6\tabcolsep) * \real{0.2500}}@{}}
\toprule\noalign{}
\begin{minipage}[b]{\linewidth}\raggedright
No.
\end{minipage} & \begin{minipage}[b]{\linewidth}\raggedright
Name
\end{minipage} & \begin{minipage}[b]{\linewidth}\raggedright
Description
\end{minipage} & \begin{minipage}[b]{\linewidth}\raggedright
Default
\end{minipage} \\
\midrule\noalign{}
\endhead
\bottomrule\noalign{}
\endlastfoot
Field & Metric & Laboratory & \\
1 & PINCHTHX & A real number defining the pinch-out threshold width for
any cell in the x-direction. NNCs are generated across inactive cells
having a width less than PINCHTHX in the x-direction. & Defined \\
ft. 0.001 & m 0.001 & cm 0.001 & \\
2 & PINCHTHY & A real number defining the pinch-out threshold width for
any cell in the y-direction. NNCs are generated across inactive cells
having a width less than PINCHTHY in the y-direction. & Defined \\
ft. 0.001 & m 0.001 & cm 0.001 & \\
Notes: & & & \\
\end{longtable}

\emph{Table 6.111: PINCHXY Keyword Description}

\subsubsection{Example}\label{example-227}

The example below will create NNCs between the cells in the areal plane
having cell widths less than 0.01 in either feet or metres in both the
x- and y-directions.

\begin{verbatim}
--
--       SET PINCH-OUT PARAMETERS FOR AREAL PLANE
--
PINCHXY
--       X-DIRC      Y-DIRC
--       THRESHOLD   THRESHOLD
--
         0.01        0.01                                              /
\end{verbatim}

\subsection{PORO -- Define the Porosity Values for All the
Cells}\label{poro-define-the-porosity-values-for-all-the-cells}

\begin{longtable}[]{@{}
  >{\raggedright\arraybackslash}p{(\linewidth - 14\tabcolsep) * \real{0.1250}}
  >{\raggedright\arraybackslash}p{(\linewidth - 14\tabcolsep) * \real{0.1250}}
  >{\raggedright\arraybackslash}p{(\linewidth - 14\tabcolsep) * \real{0.1250}}
  >{\raggedright\arraybackslash}p{(\linewidth - 14\tabcolsep) * \real{0.1250}}
  >{\raggedright\arraybackslash}p{(\linewidth - 14\tabcolsep) * \real{0.1250}}
  >{\raggedright\arraybackslash}p{(\linewidth - 14\tabcolsep) * \real{0.1250}}
  >{\raggedright\arraybackslash}p{(\linewidth - 14\tabcolsep) * \real{0.1250}}
  >{\raggedright\arraybackslash}p{(\linewidth - 14\tabcolsep) * \real{0.1250}}@{}}
\toprule\noalign{}
\begin{minipage}[b]{\linewidth}\raggedright
RUNSPEC
\end{minipage} & \begin{minipage}[b]{\linewidth}\raggedright
GRID
\end{minipage} & \begin{minipage}[b]{\linewidth}\raggedright
EDIT
\end{minipage} & \begin{minipage}[b]{\linewidth}\raggedright
PROPS
\end{minipage} & \begin{minipage}[b]{\linewidth}\raggedright
REGIONS
\end{minipage} & \begin{minipage}[b]{\linewidth}\raggedright
SOLUTION
\end{minipage} & \begin{minipage}[b]{\linewidth}\raggedright
SUMMARY
\end{minipage} & \begin{minipage}[b]{\linewidth}\raggedright
SCHEDULE
\end{minipage} \\
\midrule\noalign{}
\endhead
\bottomrule\noalign{}
\endlastfoot
\end{longtable}

\subsubsection{Description}\label{description-333}

PORO defines the porosity for all the cells in the model via an array.
The keyword can be used with all grid types.

\begin{longtable}[]{@{}
  >{\raggedright\arraybackslash}p{(\linewidth - 6\tabcolsep) * \real{0.2500}}
  >{\raggedright\arraybackslash}p{(\linewidth - 6\tabcolsep) * \real{0.2500}}
  >{\raggedright\arraybackslash}p{(\linewidth - 6\tabcolsep) * \real{0.2500}}
  >{\raggedright\arraybackslash}p{(\linewidth - 6\tabcolsep) * \real{0.2500}}@{}}
\toprule\noalign{}
\begin{minipage}[b]{\linewidth}\raggedright
No.
\end{minipage} & \begin{minipage}[b]{\linewidth}\raggedright
Name
\end{minipage} & \begin{minipage}[b]{\linewidth}\raggedright
Description
\end{minipage} & \begin{minipage}[b]{\linewidth}\raggedright
Default
\end{minipage} \\
\midrule\noalign{}
\endhead
\bottomrule\noalign{}
\endlastfoot
Field & Metric & Laboratory & \\
1 & PORO & PORO is an array of real positive numbers that are greater
than or equal to zero and less than or equal to one that are the
porosity values for each cell in the model. Repeat counts may be used,
for example 3000*0.15 & None \\
dimensionless & dimensionless & dimensionless & \\
Notes: & & & \\
\end{longtable}

\emph{Table 6.112: PORO Keyword Description}

See also the NTG, PERMX, PERMY and PERMX keywords to fully define a
grid's properties

\subsubsection{Example}\label{example-228}

\begin{verbatim}
--
--    DEFINE GRID BLOCK POROSITY DATA FOR ALL CELLS (BASED ON NX x NY x NZ = 300)
--
PORO
      300*0.300                                                                 /
\end{verbatim}

The above example defines a constant porosity of 0.300 to all 300 cells
in the model as defined by the DIMENS keyword in the RUNSPEC section.

\subsection{PYEND -- End the Definition of a PYINPUT
Section}\label{pyend-end-the-definition-of-a-pyinput-section}

\begin{longtable}[]{@{}
  >{\raggedright\arraybackslash}p{(\linewidth - 14\tabcolsep) * \real{0.1250}}
  >{\raggedright\arraybackslash}p{(\linewidth - 14\tabcolsep) * \real{0.1250}}
  >{\raggedright\arraybackslash}p{(\linewidth - 14\tabcolsep) * \real{0.1250}}
  >{\raggedright\arraybackslash}p{(\linewidth - 14\tabcolsep) * \real{0.1250}}
  >{\raggedright\arraybackslash}p{(\linewidth - 14\tabcolsep) * \real{0.1250}}
  >{\raggedright\arraybackslash}p{(\linewidth - 14\tabcolsep) * \real{0.1250}}
  >{\raggedright\arraybackslash}p{(\linewidth - 14\tabcolsep) * \real{0.1250}}
  >{\raggedright\arraybackslash}p{(\linewidth - 14\tabcolsep) * \real{0.1250}}@{}}
\toprule\noalign{}
\begin{minipage}[b]{\linewidth}\raggedright
RUNSPEC
\end{minipage} & \begin{minipage}[b]{\linewidth}\raggedright
GRID
\end{minipage} & \begin{minipage}[b]{\linewidth}\raggedright
EDIT
\end{minipage} & \begin{minipage}[b]{\linewidth}\raggedright
PROPS
\end{minipage} & \begin{minipage}[b]{\linewidth}\raggedright
REGIONS
\end{minipage} & \begin{minipage}[b]{\linewidth}\raggedright
SOLUTION
\end{minipage} & \begin{minipage}[b]{\linewidth}\raggedright
SUMMARY
\end{minipage} & \begin{minipage}[b]{\linewidth}\raggedright
SCHEDULE
\end{minipage} \\
\midrule\noalign{}
\endhead
\bottomrule\noalign{}
\endlastfoot
\end{longtable}

\subsubsection{Description}\label{description-334}

The PYINPUT and PYEND keywords are a part of OPM Flow's Python scripting
facility that processes standard Python commands that can be used to
manipulate and define the simulators input parameters during processing
of the input deck. The main purpose of the facility is to script the
construction of the various keywords.

PYINPUT declares the start of a PYINPUT Definition Section on a single
separate line, which is then followed by various standard Python
commands, with one command per line. A PYINPUT Definition Section is
terminated by a PYEND keyword (this keyword) on a separate single line.

There is no data required for this keyword and there is no terminating
``/'' for this keyword.

Although this keyword is read by OPM Flow and the script processing has
been implemented, one should use caution when using this facility as it
may result in OPM Flow aborting. This is because the PYINPUT facility
allows the user to implement complex functionality and the
implementation is new for the 2020-04 release. Users should therefore
use caution when using this facility.

\begin{tcolorbox}[enhanced jigsaw, title=\textcolor{quarto-callout-note-color}{\faInfo}\hspace{0.5em}{Note}, coltitle=black, breakable, toprule=.15mm, colbacktitle=quarto-callout-note-color!10!white, leftrule=.75mm, opacitybacktitle=0.6, colback=white, bottomtitle=1mm, toptitle=1mm, titlerule=0mm, left=2mm, colframe=quarto-callout-note-color-frame, arc=.35mm, rightrule=.15mm, bottomrule=.15mm, opacityback=0]

This is an OPM Flow specific keyword for the simulator's scripting
facility using the standard Python interpreter, as such it gives more
flexibility than the commercial simulator's data editing keywords (ADD,
EQUAL, MULTIPLY, etc.), although OPM Flow also supports these keywords
as well. The PYINPUT facility should be considered experimental as
details of the OPM Flow - Python interface might change for future
releases. In particular, the current implementation is quite minimal;
however, future releases are expected to add more entry points in the
simulator's deck class which can be used to manipulate the input deck as
the data is loaded. As a user you are encouraged to come with wishes in
this regard. The PYINPUT facility is very powerful and allows for any
piece of Python code to be included and run, including potentially
malicious code. The important point is to scrutinize the Python code in
between PYINPUT and PYEND in a deck you receive from other parties.

\end{tcolorbox}

See also the PYACTION keyword in the SCHEDULE section which is also part
of OPM Flow's Python scripting facility, that loads a standard Python
script file that can be used to define a series of conditions and
actions as the simulation proceeds through time.

\subsubsection{Example}\label{example-229}

The example shows how to construct the DX variable in the GRID section
and to add the resulting DX array as part of the input deck.

\begin{verbatim}
--
--       START OF PYINPUT SECTION
--
PYINPUT
#
# Import Numpy Model
#
import numpy as np
#
# Define DX and Get the Input Decks Unit Systems
#
dx = np.array([100.0, 100.0, 100.0, 100.0])
active_unit_system  = context.deck.active_unit_system()
default_unit_system = context.deck.default_unit_system()
#
# Set DX in the Input Deck
#
kw = context.DeckKeyword( context.parser['DX'], dx, active_unit_system,
default_unit_system )
context.deck.add(kw)

PYEND
\end{verbatim}

The active Parser objects are accessible as context.parser and the
active Deck object is available as context.deck.

\subsection{PYINPUT -- Define the Start of a PYINPUT
Section}\label{pyinput-define-the-start-of-a-pyinput-section}

\begin{longtable}[]{@{}
  >{\raggedright\arraybackslash}p{(\linewidth - 14\tabcolsep) * \real{0.1250}}
  >{\raggedright\arraybackslash}p{(\linewidth - 14\tabcolsep) * \real{0.1250}}
  >{\raggedright\arraybackslash}p{(\linewidth - 14\tabcolsep) * \real{0.1250}}
  >{\raggedright\arraybackslash}p{(\linewidth - 14\tabcolsep) * \real{0.1250}}
  >{\raggedright\arraybackslash}p{(\linewidth - 14\tabcolsep) * \real{0.1250}}
  >{\raggedright\arraybackslash}p{(\linewidth - 14\tabcolsep) * \real{0.1250}}
  >{\raggedright\arraybackslash}p{(\linewidth - 14\tabcolsep) * \real{0.1250}}
  >{\raggedright\arraybackslash}p{(\linewidth - 14\tabcolsep) * \real{0.1250}}@{}}
\toprule\noalign{}
\begin{minipage}[b]{\linewidth}\raggedright
RUNSPEC
\end{minipage} & \begin{minipage}[b]{\linewidth}\raggedright
GRID
\end{minipage} & \begin{minipage}[b]{\linewidth}\raggedright
EDIT
\end{minipage} & \begin{minipage}[b]{\linewidth}\raggedright
PROPS
\end{minipage} & \begin{minipage}[b]{\linewidth}\raggedright
REGIONS
\end{minipage} & \begin{minipage}[b]{\linewidth}\raggedright
SOLUTION
\end{minipage} & \begin{minipage}[b]{\linewidth}\raggedright
SUMMARY
\end{minipage} & \begin{minipage}[b]{\linewidth}\raggedright
SCHEDULE
\end{minipage} \\
\midrule\noalign{}
\endhead
\bottomrule\noalign{}
\endlastfoot
\end{longtable}

\subsubsection{Description}\label{description-335}

The PYINPUT and PYEND keywords are a part of OPM Flow's Python scripting
facility that processes standard Python commands that can be used to
manipulate and define the simulators input parameters during processing
of the input deck. The main purpose of the facility is to script the
construction of the various keywords used by the simulator.

PYINPUT declares the start of a PYINPUT Definition Section on a single
separate line, which is then followed by various standard Python
commands, with one command per line. A PYINPUT Definition Section is
terminated by a PYEND keyword on a separate single line.

There is no data required for this keyword and there is no terminating
``/'' for this keyword.

Although this keyword is read by OPM Flow and the script processing has
been implemented, one should use caution when using this facility as it
may result in OPM Flow aborting. This is because the PYINPUT facility
allows the user to implement complex functionality and the
implementation is new for the 2020-04 release. Users should therefore
use caution when using this facility.

\begin{tcolorbox}[enhanced jigsaw, title=\textcolor{quarto-callout-note-color}{\faInfo}\hspace{0.5em}{Note}, coltitle=black, breakable, toprule=.15mm, colbacktitle=quarto-callout-note-color!10!white, leftrule=.75mm, opacitybacktitle=0.6, colback=white, bottomtitle=1mm, toptitle=1mm, titlerule=0mm, left=2mm, colframe=quarto-callout-note-color-frame, arc=.35mm, rightrule=.15mm, bottomrule=.15mm, opacityback=0]

This is an OPM Flow specific keyword for the simulator's scripting
facility using the standard Python interpreter, as such it gives more
flexibility than the commercial simulator's data editing keywords (ADD,
EQUALS, MULTIPLY, etc.), although OPM Flow also supports these keywords
as well. The PYINPUT facility should be considered experimental as
details of the OPM Flow - Python interface might change for future
releases. In particular, the current implementation is quite minimal;
however, future releases are expected to add more entry points in the
simulator's deck class which can be used to manipulate the input deck as
the data is loaded. As a user you are encouraged to come with wishes in
this regard. The PYINPUT facility is very powerful and allows for any
piece of Python code to be included and run, including potentially
malicious code. The important point is to scrutinize the Python code in
between PYINPUT and PYEND in a deck you receive from other parties.

\end{tcolorbox}

\begin{longtable}[]{@{}
  >{\raggedright\arraybackslash}p{(\linewidth - 6\tabcolsep) * \real{0.2500}}
  >{\raggedright\arraybackslash}p{(\linewidth - 6\tabcolsep) * \real{0.2500}}
  >{\raggedright\arraybackslash}p{(\linewidth - 6\tabcolsep) * \real{0.2500}}
  >{\raggedright\arraybackslash}p{(\linewidth - 6\tabcolsep) * \real{0.2500}}@{}}
\toprule\noalign{}
\begin{minipage}[b]{\linewidth}\raggedright
No.
\end{minipage} & \begin{minipage}[b]{\linewidth}\raggedright
Name
\end{minipage} & \begin{minipage}[b]{\linewidth}\raggedright
Description
\end{minipage} & \begin{minipage}[b]{\linewidth}\raggedright
Default
\end{minipage} \\
\midrule\noalign{}
\endhead
\bottomrule\noalign{}
\endlastfoot
PYINPUT & PYINPUT declares the start of a PYINPUT Definition Section.
This is then followed by any number Python commands. & Not Applicable
& \\
1-1 & PYTHON & A series of standard Python commands with one line per
command. The active Parser objects are accessible as context.parser and
the active Deck object is available as context.deck. & \\
PYEND & PYEND declares the end of a PYINPUT Definition Section. The
Python code between PYINPUT and PYEND is read and executed, and the
simulator thenreturns to reading the normal simulation input deck. &
& \\
Notes: & & & \\
\end{longtable}

\emph{Table 6.113: PYINPUT Keyword Description}

The PYINPUT/PYEND set of keywords is a result of combining two
programming languages, the interactive Python interpreter and OPM Flow's
source code language C++. When combining two languages one extends and
embeds one into the other. When extending Python with C++ the
functionality implemented in C++ is made available to Python
applications, when embedding Python in C++ one can call Python functions
from within C++. The PYINPUT/PYEND set of keywords is based on embedding
a Python interpreter in the C++ OPM Flow simulator, but the Python code
actually runs as part of the PYINPUT keyword is based on wrapping C++
objects in Python, that is extending Python.

The Python code in between the PYINPUT and PYEND keywords are imported
during processing of the input deck and as such this implies that basic
Python syntax checking is performed during reading the Python script.

See also the PYACTION keyword in the SCHEDULE section which is also part
of OPM Flow's Python scripting facility, that loads a standard Python
script file that can be used to define a series of conditions and
actions as the simulation proceeds through time.

\subsubsection{Example}\label{example-230}

The example shows how to construct the DX variable in the GRID section
and to add the resulting DX array as part of the input deck.

\begin{verbatim}
--
--       START OF PYINPUT SECTION
--
PYINPUT
#
# Import Numpy Model
#
import numpy as np
#
# Define DX and Get the Input Decks Unit Systems
#
dx = np.array([100.0, 100.0, 100.0, 100.0])
active_unit_system  = context.deck.active_unit_system()
default_unit_system = context.deck.default_unit_system()
#
# Set DX in the Input Deck
#
kw = context.DeckKeyword( context.parser['DX'], dx, active_unit_system,
default_unit_system )
context.deck.add(kw)

PYEND
\end{verbatim}

The active Parser objects is accessible as context.parser and the active
Deck object is available as context.deck.

\subsection{QMOBIL -- Activate or Deactivate LGR End-Point Mobility
Correction}\label{qmobil-activate-or-deactivate-lgr-end-point-mobility-correction}

\begin{longtable}[]{@{}
  >{\raggedright\arraybackslash}p{(\linewidth - 14\tabcolsep) * \real{0.1250}}
  >{\raggedright\arraybackslash}p{(\linewidth - 14\tabcolsep) * \real{0.1250}}
  >{\raggedright\arraybackslash}p{(\linewidth - 14\tabcolsep) * \real{0.1250}}
  >{\raggedright\arraybackslash}p{(\linewidth - 14\tabcolsep) * \real{0.1250}}
  >{\raggedright\arraybackslash}p{(\linewidth - 14\tabcolsep) * \real{0.1250}}
  >{\raggedright\arraybackslash}p{(\linewidth - 14\tabcolsep) * \real{0.1250}}
  >{\raggedright\arraybackslash}p{(\linewidth - 14\tabcolsep) * \real{0.1250}}
  >{\raggedright\arraybackslash}p{(\linewidth - 14\tabcolsep) * \real{0.1250}}@{}}
\toprule\noalign{}
\begin{minipage}[b]{\linewidth}\raggedright
RUNSPEC
\end{minipage} & \begin{minipage}[b]{\linewidth}\raggedright
GRID
\end{minipage} & \begin{minipage}[b]{\linewidth}\raggedright
EDIT
\end{minipage} & \begin{minipage}[b]{\linewidth}\raggedright
PROPS
\end{minipage} & \begin{minipage}[b]{\linewidth}\raggedright
REGIONS
\end{minipage} & \begin{minipage}[b]{\linewidth}\raggedright
SOLUTION
\end{minipage} & \begin{minipage}[b]{\linewidth}\raggedright
SUMMARY
\end{minipage} & \begin{minipage}[b]{\linewidth}\raggedright
SCHEDULE
\end{minipage} \\
\midrule\noalign{}
\endhead
\bottomrule\noalign{}
\endlastfoot
\end{longtable}

\subsubsection{Description}\label{description-336}

The QMOBIL keyword activates or deactivates the end-point mobility
correction for Local Grid Refinements (``LGR''), for when LGRs have been
activated for the input deck using the LGR keyword in the RUNSPEC
section. QMOBIL should be placed in between the LGR definition keywords
CARFIN, or RADIN (or RAFDIN4) and the ENDFIN keyword in the GRID
section.

This keyword is not supported by OPM Flow but would change the results
if supported so the simulation will be stopped.

\subsection{RADFIN -- Define a Radial Local Grid Refinement with One
Column}\label{radfin-define-a-radial-local-grid-refinement-with-one-column}

\begin{longtable}[]{@{}
  >{\raggedright\arraybackslash}p{(\linewidth - 14\tabcolsep) * \real{0.1250}}
  >{\raggedright\arraybackslash}p{(\linewidth - 14\tabcolsep) * \real{0.1250}}
  >{\raggedright\arraybackslash}p{(\linewidth - 14\tabcolsep) * \real{0.1250}}
  >{\raggedright\arraybackslash}p{(\linewidth - 14\tabcolsep) * \real{0.1250}}
  >{\raggedright\arraybackslash}p{(\linewidth - 14\tabcolsep) * \real{0.1250}}
  >{\raggedright\arraybackslash}p{(\linewidth - 14\tabcolsep) * \real{0.1250}}
  >{\raggedright\arraybackslash}p{(\linewidth - 14\tabcolsep) * \real{0.1250}}
  >{\raggedright\arraybackslash}p{(\linewidth - 14\tabcolsep) * \real{0.1250}}@{}}
\toprule\noalign{}
\begin{minipage}[b]{\linewidth}\raggedright
RUNSPEC
\end{minipage} & \begin{minipage}[b]{\linewidth}\raggedright
GRID
\end{minipage} & \begin{minipage}[b]{\linewidth}\raggedright
EDIT
\end{minipage} & \begin{minipage}[b]{\linewidth}\raggedright
PROPS
\end{minipage} & \begin{minipage}[b]{\linewidth}\raggedright
REGIONS
\end{minipage} & \begin{minipage}[b]{\linewidth}\raggedright
SOLUTION
\end{minipage} & \begin{minipage}[b]{\linewidth}\raggedright
SUMMARY
\end{minipage} & \begin{minipage}[b]{\linewidth}\raggedright
SCHEDULE
\end{minipage} \\
\midrule\noalign{}
\endhead
\bottomrule\noalign{}
\endlastfoot
\end{longtable}

\subsubsection{Description}\label{description-337}

This keyword defines a radial local grid refinement using one columns
Local grid refinement is currently not supported by OPM Flow.

This keyword is not supported by OPM Flow but would change the results
if supported so the simulation will be stopped.

\subsection{RADFIN4 -- Define a Radial Local Grid Refinement with Four
Columns}\label{radfin4-define-a-radial-local-grid-refinement-with-four-columns}

\begin{longtable}[]{@{}
  >{\raggedright\arraybackslash}p{(\linewidth - 14\tabcolsep) * \real{0.1250}}
  >{\raggedright\arraybackslash}p{(\linewidth - 14\tabcolsep) * \real{0.1250}}
  >{\raggedright\arraybackslash}p{(\linewidth - 14\tabcolsep) * \real{0.1250}}
  >{\raggedright\arraybackslash}p{(\linewidth - 14\tabcolsep) * \real{0.1250}}
  >{\raggedright\arraybackslash}p{(\linewidth - 14\tabcolsep) * \real{0.1250}}
  >{\raggedright\arraybackslash}p{(\linewidth - 14\tabcolsep) * \real{0.1250}}
  >{\raggedright\arraybackslash}p{(\linewidth - 14\tabcolsep) * \real{0.1250}}
  >{\raggedright\arraybackslash}p{(\linewidth - 14\tabcolsep) * \real{0.1250}}@{}}
\toprule\noalign{}
\begin{minipage}[b]{\linewidth}\raggedright
RUNSPEC
\end{minipage} & \begin{minipage}[b]{\linewidth}\raggedright
GRID
\end{minipage} & \begin{minipage}[b]{\linewidth}\raggedright
EDIT
\end{minipage} & \begin{minipage}[b]{\linewidth}\raggedright
PROPS
\end{minipage} & \begin{minipage}[b]{\linewidth}\raggedright
REGIONS
\end{minipage} & \begin{minipage}[b]{\linewidth}\raggedright
SOLUTION
\end{minipage} & \begin{minipage}[b]{\linewidth}\raggedright
SUMMARY
\end{minipage} & \begin{minipage}[b]{\linewidth}\raggedright
SCHEDULE
\end{minipage} \\
\midrule\noalign{}
\endhead
\bottomrule\noalign{}
\endlastfoot
\end{longtable}

\subsubsection{Description}\label{description-338}

This keyword defines a radial local grid refinement using four columns.
Local grid refinement is currently not supported by OPM Flow.

This keyword is not supported by OPM Flow but would change the results
if supported so the simulation will be stopped.

\subsection{REFINE -- Start the Definition of a Local Grid
Refinement}\label{refine-start-the-definition-of-a-local-grid-refinement}

\begin{longtable}[]{@{}
  >{\raggedright\arraybackslash}p{(\linewidth - 14\tabcolsep) * \real{0.1250}}
  >{\raggedright\arraybackslash}p{(\linewidth - 14\tabcolsep) * \real{0.1250}}
  >{\raggedright\arraybackslash}p{(\linewidth - 14\tabcolsep) * \real{0.1250}}
  >{\raggedright\arraybackslash}p{(\linewidth - 14\tabcolsep) * \real{0.1250}}
  >{\raggedright\arraybackslash}p{(\linewidth - 14\tabcolsep) * \real{0.1250}}
  >{\raggedright\arraybackslash}p{(\linewidth - 14\tabcolsep) * \real{0.1250}}
  >{\raggedright\arraybackslash}p{(\linewidth - 14\tabcolsep) * \real{0.1250}}
  >{\raggedright\arraybackslash}p{(\linewidth - 14\tabcolsep) * \real{0.1250}}@{}}
\toprule\noalign{}
\begin{minipage}[b]{\linewidth}\raggedright
RUNSPEC
\end{minipage} & \begin{minipage}[b]{\linewidth}\raggedright
GRID
\end{minipage} & \begin{minipage}[b]{\linewidth}\raggedright
EDIT
\end{minipage} & \begin{minipage}[b]{\linewidth}\raggedright
PROPS
\end{minipage} & \begin{minipage}[b]{\linewidth}\raggedright
REGIONS
\end{minipage} & \begin{minipage}[b]{\linewidth}\raggedright
SOLUTION
\end{minipage} & \begin{minipage}[b]{\linewidth}\raggedright
SUMMARY
\end{minipage} & \begin{minipage}[b]{\linewidth}\raggedright
SCHEDULE
\end{minipage} \\
\midrule\noalign{}
\endhead
\bottomrule\noalign{}
\endlastfoot
\end{longtable}

\subsubsection{Description}\label{description-339}

The REFINE keyword defines the start of a Cartesian or radial Local Grid
Refinement (``LGR'') definition that sets the properties of the selected
LGR. The keyword is then followed by the property keywords associated
with the section where the keyword is being invoked. For example, if the
REFINE keyword is used in the GRID section then most of the keywords in
that section can be used to set the grid properties for the LGR.

The ENDFIN keyword is used to terminate the LGR definition.

There is no data required for this keyword.

This keyword is not supported by OPM Flow but would change the results
if supported so the simulation will be stopped.

\subsection{RESVNUM -- Define Reservoir Coordinate Data
Set}\label{resvnum-define-reservoir-coordinate-data-set}

\begin{longtable}[]{@{}
  >{\raggedright\arraybackslash}p{(\linewidth - 14\tabcolsep) * \real{0.1250}}
  >{\raggedright\arraybackslash}p{(\linewidth - 14\tabcolsep) * \real{0.1250}}
  >{\raggedright\arraybackslash}p{(\linewidth - 14\tabcolsep) * \real{0.1250}}
  >{\raggedright\arraybackslash}p{(\linewidth - 14\tabcolsep) * \real{0.1250}}
  >{\raggedright\arraybackslash}p{(\linewidth - 14\tabcolsep) * \real{0.1250}}
  >{\raggedright\arraybackslash}p{(\linewidth - 14\tabcolsep) * \real{0.1250}}
  >{\raggedright\arraybackslash}p{(\linewidth - 14\tabcolsep) * \real{0.1250}}
  >{\raggedright\arraybackslash}p{(\linewidth - 14\tabcolsep) * \real{0.1250}}@{}}
\toprule\noalign{}
\begin{minipage}[b]{\linewidth}\raggedright
RUNSPEC
\end{minipage} & \begin{minipage}[b]{\linewidth}\raggedright
GRID
\end{minipage} & \begin{minipage}[b]{\linewidth}\raggedright
EDIT
\end{minipage} & \begin{minipage}[b]{\linewidth}\raggedright
PROPS
\end{minipage} & \begin{minipage}[b]{\linewidth}\raggedright
REGIONS
\end{minipage} & \begin{minipage}[b]{\linewidth}\raggedright
SOLUTION
\end{minipage} & \begin{minipage}[b]{\linewidth}\raggedright
SUMMARY
\end{minipage} & \begin{minipage}[b]{\linewidth}\raggedright
SCHEDULE
\end{minipage} \\
\midrule\noalign{}
\endhead
\bottomrule\noalign{}
\endlastfoot
\end{longtable}

\subsubsection{Description}\label{description-340}

The RESVNUM keyword is used to define the start of a reservoir
coordinate data set and stipulates the reservoir number for the data
set. The keyword is used in conjunction with the COORD keyword in the
GRID section, that specifies a set of coordinate lines or pillars for a
reservoir grid via an array. Note that the COORD keyword should
immediately follow the RESVNUM keyword.

Although the keyword is processed by OPM Flow, the keyword is
effectively ignored as only the default value of one is supported.

\begin{longtable}[]{@{}
  >{\raggedright\arraybackslash}p{(\linewidth - 6\tabcolsep) * \real{0.2500}}
  >{\raggedright\arraybackslash}p{(\linewidth - 6\tabcolsep) * \real{0.2500}}
  >{\raggedright\arraybackslash}p{(\linewidth - 6\tabcolsep) * \real{0.2500}}
  >{\raggedright\arraybackslash}p{(\linewidth - 6\tabcolsep) * \real{0.2500}}@{}}
\toprule\noalign{}
\begin{minipage}[b]{\linewidth}\raggedright
No.
\end{minipage} & \begin{minipage}[b]{\linewidth}\raggedright
Name
\end{minipage} & \begin{minipage}[b]{\linewidth}\raggedright
Description
\end{minipage} & \begin{minipage}[b]{\linewidth}\raggedright
Default
\end{minipage} \\
\midrule\noalign{}
\endhead
\bottomrule\noalign{}
\endlastfoot
1 & RESVNUM & A positive integer values that defines the reservoir
coordinate data set, or the independent reservoir, for which the
subsequent COORD data is to be associated with. RESVNUM should be less
than or equal to NUMRES on the NUMRES keyword in the RUNSPEC section.
OPM Flow currently only accepts a single data set, that is the default
value of one. & 1 \\
Notes: & & & \\
\end{longtable}

\emph{Table 6.114: RESVNUM Keyword Description}

See the NUMRES keyword in the RUNSPEC section that defines the number of
reservoir grids (COORD data sets) that the simulator should process.

The facility is useful to combine two separate reservoir grids into one
model in the simulator.

\subsubsection{Example}\label{example-231}

\begin{verbatim}
--
--       NUMRES
--       NUMBER
RESVNUM
         1                                                                     /
--
-- SPECIFY VERTICAL COORDINATE LINES FOR A REGULAR 3 x 2 GRID
--(DX = 100 and DY = 200)
--
-- X1     Y1    Z1      X2     Y2    Z2
-- ---    ---   ----    ---    ---   ----
COORD
     0      0   1000      0      0   5000
   100      0   1000    100      0   5000
   200      0   1000    200      0   5000
   300      0   1000    300      0   5000
     0    200   1000      0    200   5000
   100    200   1000    100    200   5000
   200    200   1000    200    200   5000
   300    200   1000    300    200   5000
     0    400   1000      0    400   5000
   100    400   1000    100    400   5000
   200    400   1000    200    400   5000
   300    400   1000    300    400   5000
/
--
--       NUMRES
--       NUMBER
RESVNUM
         2                                                                     /
--
-- SPECIFY VERTICAL COORDINATE LINES FOR A REGULAR 3 x 2 GRID
--(DX = 100 and DY = 200)
--
-- X1     Y1    Z1      X2     Y2    Z2
-- ---    ---   ----    ---    ---   ----
COORD
     0      0   1000      0      0   5000
   100      0   1000    100      0   5000
   200      0   1000    200      0   5000
   300      0   1000    300      0   5000
     0    200   1000      0    200   5000
   100    200   1000    100    200   5000
   200    200   1000    200    200   5000
   300    200   1000    300    200   5000
     0    400   1000      0    400   5000
   100    400   1000    100    400   5000
   200    400   1000    200    400   5000
   300    400   1000    300    400   5000
/

\end{verbatim}

\subsection{ROCKFRAC -- Define the Rock Volume to Bulk Volume Fraction
for All the
Cells}\label{rockfrac-define-the-rock-volume-to-bulk-volume-fraction-for-all-the-cells}

\begin{longtable}[]{@{}
  >{\raggedright\arraybackslash}p{(\linewidth - 14\tabcolsep) * \real{0.1250}}
  >{\raggedright\arraybackslash}p{(\linewidth - 14\tabcolsep) * \real{0.1250}}
  >{\raggedright\arraybackslash}p{(\linewidth - 14\tabcolsep) * \real{0.1250}}
  >{\raggedright\arraybackslash}p{(\linewidth - 14\tabcolsep) * \real{0.1250}}
  >{\raggedright\arraybackslash}p{(\linewidth - 14\tabcolsep) * \real{0.1250}}
  >{\raggedright\arraybackslash}p{(\linewidth - 14\tabcolsep) * \real{0.1250}}
  >{\raggedright\arraybackslash}p{(\linewidth - 14\tabcolsep) * \real{0.1250}}
  >{\raggedright\arraybackslash}p{(\linewidth - 14\tabcolsep) * \real{0.1250}}@{}}
\toprule\noalign{}
\begin{minipage}[b]{\linewidth}\raggedright
RUNSPEC
\end{minipage} & \begin{minipage}[b]{\linewidth}\raggedright
GRID
\end{minipage} & \begin{minipage}[b]{\linewidth}\raggedright
EDIT
\end{minipage} & \begin{minipage}[b]{\linewidth}\raggedright
PROPS
\end{minipage} & \begin{minipage}[b]{\linewidth}\raggedright
REGIONS
\end{minipage} & \begin{minipage}[b]{\linewidth}\raggedright
SOLUTION
\end{minipage} & \begin{minipage}[b]{\linewidth}\raggedright
SUMMARY
\end{minipage} & \begin{minipage}[b]{\linewidth}\raggedright
SCHEDULE
\end{minipage} \\
\midrule\noalign{}
\endhead
\bottomrule\noalign{}
\endlastfoot
\end{longtable}

\subsubsection{Description}\label{description-341}

ROCKFRAC defines the rock volume to bulk volume fraction for all the
cells, The keyword can be used with all grid types. Rock volume of a
grid block is calculated by multiply a cell's bulk volume by it's
ROCKFRAC volume. A cell's rock volume is used in the Coal option to
calculate the adsorbed gas in the rock (coal), as well as the Thermal
and Temp options to calculate the energy is stored in the rock.

This keyword is not supported by OPM Flow but would change the results
if supported so the simulation will be stopped.

\begin{longtable}[]{@{}
  >{\raggedright\arraybackslash}p{(\linewidth - 6\tabcolsep) * \real{0.2500}}
  >{\raggedright\arraybackslash}p{(\linewidth - 6\tabcolsep) * \real{0.2500}}
  >{\raggedright\arraybackslash}p{(\linewidth - 6\tabcolsep) * \real{0.2500}}
  >{\raggedright\arraybackslash}p{(\linewidth - 6\tabcolsep) * \real{0.2500}}@{}}
\toprule\noalign{}
\begin{minipage}[b]{\linewidth}\raggedright
No.
\end{minipage} & \begin{minipage}[b]{\linewidth}\raggedright
Name
\end{minipage} & \begin{minipage}[b]{\linewidth}\raggedright
Description
\end{minipage} & \begin{minipage}[b]{\linewidth}\raggedright
Default
\end{minipage} \\
\midrule\noalign{}
\endhead
\bottomrule\noalign{}
\endlastfoot
Field & Metric & Laboratory & \\
1 & ROCKFRAC & ROCKFRAC is an array of real numbers greater than or
equal to zero and less than or equal to one, that are assigned the rock
volume to bulk volume fraction values for each cell in the model. The
number of entries should correspond to the NX x NY x NZ parameters on
the DIMENS keyword. Repeat counts may be used, for example 200*0.850. &
None \\
dimensionless & dimensionless & dimensionless & \\
Notes: & & & \\
\end{longtable}

\emph{Table 6.115: NTG Keyword Description}

See also the PORO, PERMX, PERMY, PERMZ and NTG keywords to fully define
a grid's properties.

\subsubsection{Example}\label{example-232}

\begin{verbatim}
--
--       DEFINE GRID ROCKFRAC DATA FOR ALL CELLS (BASED ON NX x NY x NZ = 300)
--
ROCKFRAC
         100*1.000   100*0.850   100*0.500                                     /
\end{verbatim}

The above example defines a constant ROCKFRAC of 1.00 for the first 100
cells, then 0.85 for the second 100 hundred cells, and finally 0.500 for
the last 100 cell, for the 300 cells in the model as defined by the
DIMENS keyword in the RUNSPEC section.

\begin{verbatim}
\end{verbatim}

\subsection{RPTGRID -- Define GRID Section
Reporting}\label{rptgrid-define-grid-section-reporting}

\begin{longtable}[]{@{}
  >{\raggedright\arraybackslash}p{(\linewidth - 14\tabcolsep) * \real{0.1250}}
  >{\raggedright\arraybackslash}p{(\linewidth - 14\tabcolsep) * \real{0.1250}}
  >{\raggedright\arraybackslash}p{(\linewidth - 14\tabcolsep) * \real{0.1250}}
  >{\raggedright\arraybackslash}p{(\linewidth - 14\tabcolsep) * \real{0.1250}}
  >{\raggedright\arraybackslash}p{(\linewidth - 14\tabcolsep) * \real{0.1250}}
  >{\raggedright\arraybackslash}p{(\linewidth - 14\tabcolsep) * \real{0.1250}}
  >{\raggedright\arraybackslash}p{(\linewidth - 14\tabcolsep) * \real{0.1250}}
  >{\raggedright\arraybackslash}p{(\linewidth - 14\tabcolsep) * \real{0.1250}}@{}}
\toprule\noalign{}
\begin{minipage}[b]{\linewidth}\raggedright
RUNSPEC
\end{minipage} & \begin{minipage}[b]{\linewidth}\raggedright
GRID
\end{minipage} & \begin{minipage}[b]{\linewidth}\raggedright
EDIT
\end{minipage} & \begin{minipage}[b]{\linewidth}\raggedright
PROPS
\end{minipage} & \begin{minipage}[b]{\linewidth}\raggedright
REGIONS
\end{minipage} & \begin{minipage}[b]{\linewidth}\raggedright
SOLUTION
\end{minipage} & \begin{minipage}[b]{\linewidth}\raggedright
SUMMARY
\end{minipage} & \begin{minipage}[b]{\linewidth}\raggedright
SCHEDULE
\end{minipage} \\
\midrule\noalign{}
\endhead
\bottomrule\noalign{}
\endlastfoot
\end{longtable}

\subsubsection{Description}\label{description-342}

This keyword defines the data in the GRID section that is to be printed
to the output print file in human readable format. The keyword has two
distinct forms, the first of which consists of the keyword followed by a
series of integers on the next line indicating the data to be printed
(see the first example). This is the original formal in the commercial
simulator and was subsequently superseded by the second format. The
second format consists of the keyword followed by a series of character
strings that indicate the data to be printed. In most cases the
character string is the keyword used to load the data in the OPM Flow
input deck, for example PORO for the porosity array. Its is anticipated
that OPM Flow will eventually support the functionality of the second
format only, the first format although recognized will be completely
ignored.

This keyword is not supported by OPM Flow but has no effect on the
results so it will be ignored.

\begin{longtable}[]{@{}llll@{}}
\toprule\noalign{}
No. & Name & Description & Default \\
\midrule\noalign{}
\endhead
\bottomrule\noalign{}
\endlastfoot
1 & ALLNCC & Print all the non-neighbor connections. & N/A \\
2 & COORD & Print the coordinate lines. & N/A \\
3 & COORDYS & Print the coordinate systems. & N/A \\
4 & DEPTH & Print grid cells center depths. & N/A \\
\ldots. & & \ldots. & N/A \\
Notes: & & & \\
\end{longtable}

\emph{Table 6.116: RPTGRID Keyword Description}

\begin{tcolorbox}[enhanced jigsaw, title=\textcolor{quarto-callout-note-color}{\faInfo}\hspace{0.5em}{Note}, coltitle=black, breakable, toprule=.15mm, colbacktitle=quarto-callout-note-color!10!white, leftrule=.75mm, opacitybacktitle=0.6, colback=white, bottomtitle=1mm, toptitle=1mm, titlerule=0mm, left=2mm, colframe=quarto-callout-note-color-frame, arc=.35mm, rightrule=.15mm, bottomrule=.15mm, opacityback=0]

This keyword has the potential to produce very large print files that
some text editors may have difficulty loading, coupled with the fact
that reviewing the data in this format is very cumbersome. A more
efficient solution is to load the *.INIT file into OPM ResInsight to
view the data graphically, this also has the benefit of being able to
filter the grid based on I, J, K ranges and grid properties.

\end{tcolorbox}

\subsubsection{Examples}\label{examples-16}

The first example shows the original format of this keyword; although
the keyword and format are recognized by OPM Flow, the format is ignored
and is unlikely to be implemented in in the simulator.

\begin{verbatim}
--
--       DEFINE GRID SECTION REPORT OPTION (ORIGINAL FORMAT)
–
RPTGRID
         1        2*0      1        3*1                                        /
\end{verbatim}

The next example shows the second format of the keyword which may be
supported in a future release of OPM Flow.

\begin{verbatim}
--       DEFINE GRID SECTION REPORT OPTIONS
--
RPTGRID
         DX       DY       DZ       DEPTH    PORO     PERMX                    /
\end{verbatim}

\subsection{RPTGRIDL -- Define GRID Section Reporting for
LGRs}\label{rptgridl-define-grid-section-reporting-for-lgrs}

\begin{longtable}[]{@{}
  >{\raggedright\arraybackslash}p{(\linewidth - 14\tabcolsep) * \real{0.1250}}
  >{\raggedright\arraybackslash}p{(\linewidth - 14\tabcolsep) * \real{0.1250}}
  >{\raggedright\arraybackslash}p{(\linewidth - 14\tabcolsep) * \real{0.1250}}
  >{\raggedright\arraybackslash}p{(\linewidth - 14\tabcolsep) * \real{0.1250}}
  >{\raggedright\arraybackslash}p{(\linewidth - 14\tabcolsep) * \real{0.1250}}
  >{\raggedright\arraybackslash}p{(\linewidth - 14\tabcolsep) * \real{0.1250}}
  >{\raggedright\arraybackslash}p{(\linewidth - 14\tabcolsep) * \real{0.1250}}
  >{\raggedright\arraybackslash}p{(\linewidth - 14\tabcolsep) * \real{0.1250}}@{}}
\toprule\noalign{}
\begin{minipage}[b]{\linewidth}\raggedright
RUNSPEC
\end{minipage} & \begin{minipage}[b]{\linewidth}\raggedright
GRID
\end{minipage} & \begin{minipage}[b]{\linewidth}\raggedright
EDIT
\end{minipage} & \begin{minipage}[b]{\linewidth}\raggedright
PROPS
\end{minipage} & \begin{minipage}[b]{\linewidth}\raggedright
REGIONS
\end{minipage} & \begin{minipage}[b]{\linewidth}\raggedright
SOLUTION
\end{minipage} & \begin{minipage}[b]{\linewidth}\raggedright
SUMMARY
\end{minipage} & \begin{minipage}[b]{\linewidth}\raggedright
SCHEDULE
\end{minipage} \\
\midrule\noalign{}
\endhead
\bottomrule\noalign{}
\endlastfoot
\end{longtable}

\subsubsection{Description}\label{description-343}

This keyword defines the data in the GRID section that is to be printed
to the output print file in human readable format for Local Grid
Refinements (``LGRs''), for when LGRs have been activated for the input
deck using the LGR keyword in the RUNSPEC section.

The keyword has two distinct forms, the first of which consists of the
keyword followed by a series of integers on the next line indicating the
data to be printed (see the first example). This is the original formal
in the commercial simulator and was subsequently superseded by the
second format. The second format consists of the keyword followed by a
series of character strings that indicate the data to be printed. In
most cases the character string is the keyword used to load the data in
the OPM Flow input deck, for example PORO for the porosity array. Its is
anticipated that OPM Flow will eventually support the functionality of
the second format only, the first format although recognized will be
completely ignored.

This keyword is not supported by OPM Flow but has no effect on the
results so it will be ignored.

\begin{longtable}[]{@{}
  >{\raggedright\arraybackslash}p{(\linewidth - 6\tabcolsep) * \real{0.2500}}
  >{\raggedright\arraybackslash}p{(\linewidth - 6\tabcolsep) * \real{0.2500}}
  >{\raggedright\arraybackslash}p{(\linewidth - 6\tabcolsep) * \real{0.2500}}
  >{\raggedright\arraybackslash}p{(\linewidth - 6\tabcolsep) * \real{0.2500}}@{}}
\toprule\noalign{}
\begin{minipage}[b]{\linewidth}\raggedright
No.
\end{minipage} & \begin{minipage}[b]{\linewidth}\raggedright
Name
\end{minipage} & \begin{minipage}[b]{\linewidth}\raggedright
Description
\end{minipage} & \begin{minipage}[b]{\linewidth}\raggedright
Default
\end{minipage} \\
\midrule\noalign{}
\endhead
\bottomrule\noalign{}
\endlastfoot
1 & ALLNCC & Print all the non-neighbor connections. & N/A \\
2 & COORD & Print the coordinate lines. & N/A \\
3 & COORDYS & Print the coordinate systems. & N/A \\
4 & DEPTH & Print grid cells center depths. & N/A \\
\ldots. & & \ldots. & N/A \\
24 & ALLNNC & ALLNNC is a defined positive integer that specifies the
type of Non-Neighbor Connections (``NNC'') to be printed, and should be
set to one of the follow: & N/A \\
\ldots. & & \ldots. & N/A \\
57 & EXTHOST & EXTHOSTS outputs host cells for Perpendicular Bisector
(``PEBI'') {[}Heinemann, Z.E. and Brand, C.W. 1988. Gridding Techniques
in Reservoir Simulation. Proc., First Intl. Forum on Reservoir
Simulation, Alpbach, Austria, 339.{]} and {[}Heinemann, Z.E., Brand,
C.W., Munka, M. et al.~1991. Modeling Reservoir Geometry With Irregular
Grids. SPE Res Eng 6 (2): 225--232. SPE-18412-PA.
http://dx.doi.org/10.2118/18412-PA{]} LGRs. & \\
\ldots. & & \ldots. & N/A \\
Notes: & & & \\
\end{longtable}

\emph{Table 6.117: RPTGRIDL Keyword Description}

\begin{tcolorbox}[enhanced jigsaw, title=\textcolor{quarto-callout-note-color}{\faInfo}\hspace{0.5em}{Note}, coltitle=black, breakable, toprule=.15mm, colbacktitle=quarto-callout-note-color!10!white, leftrule=.75mm, opacitybacktitle=0.6, colback=white, bottomtitle=1mm, toptitle=1mm, titlerule=0mm, left=2mm, colframe=quarto-callout-note-color-frame, arc=.35mm, rightrule=.15mm, bottomrule=.15mm, opacityback=0]

This keyword has the potential to produce very large print files that
some text editors may have difficulty loading, coupled with the fact
that reviewing the data in this format is very cumbersome. A more
efficient solution is to load the *.INIT file into OPM ResInsight to
view the data graphically, this also has the benefit of being able to
filter the grid based on I, J, K ranges and grid properties.

\end{tcolorbox}

\subsubsection{Examples}\label{examples-17}

The first example shows the original format of this keyword; although
the keyword and format are recognized by OPM Flow, the format is ignored
and is unlikely to be implemented in in the simulator.

\begin{verbatim}
--
--       DEFINE LGR GRID SECTION REPORT OPTION (ORIGINAL FORMAT)
–
RPTGRIDL
         1        2*0      1        3*1                                        /
\end{verbatim}

The next example shows the second format of the keyword which may be
supported in a future release of OPM Flow.

\begin{verbatim}
--       DEFINE LGR GRID SECTION REPORT OPTIONS
--
RPTGRIDL
         DX       DY       DZ       DEPTH    PORO     PERMX                    /
\end{verbatim}

\subsection{RPTINIT -- Define Output to the INIT
File}\label{rptinit-define-output-to-the-init-file}

\begin{longtable}[]{@{}
  >{\raggedright\arraybackslash}p{(\linewidth - 14\tabcolsep) * \real{0.1250}}
  >{\raggedright\arraybackslash}p{(\linewidth - 14\tabcolsep) * \real{0.1250}}
  >{\raggedright\arraybackslash}p{(\linewidth - 14\tabcolsep) * \real{0.1250}}
  >{\raggedright\arraybackslash}p{(\linewidth - 14\tabcolsep) * \real{0.1250}}
  >{\raggedright\arraybackslash}p{(\linewidth - 14\tabcolsep) * \real{0.1250}}
  >{\raggedright\arraybackslash}p{(\linewidth - 14\tabcolsep) * \real{0.1250}}
  >{\raggedright\arraybackslash}p{(\linewidth - 14\tabcolsep) * \real{0.1250}}
  >{\raggedright\arraybackslash}p{(\linewidth - 14\tabcolsep) * \real{0.1250}}@{}}
\toprule\noalign{}
\begin{minipage}[b]{\linewidth}\raggedright
RUNSPEC
\end{minipage} & \begin{minipage}[b]{\linewidth}\raggedright
GRID
\end{minipage} & \begin{minipage}[b]{\linewidth}\raggedright
EDIT
\end{minipage} & \begin{minipage}[b]{\linewidth}\raggedright
PROPS
\end{minipage} & \begin{minipage}[b]{\linewidth}\raggedright
REGIONS
\end{minipage} & \begin{minipage}[b]{\linewidth}\raggedright
SOLUTION
\end{minipage} & \begin{minipage}[b]{\linewidth}\raggedright
SUMMARY
\end{minipage} & \begin{minipage}[b]{\linewidth}\raggedright
SCHEDULE
\end{minipage} \\
\midrule\noalign{}
\endhead
\bottomrule\noalign{}
\endlastfoot
\end{longtable}

\subsubsection{Description}\label{description-344}

This keyword defines the data in the GRID and EDIT sections that is to
be written out to the INIT file (\emph{.INIT or }.FINIT). The format
consists of the keyword followed by a series of character strings that
indicate the data to be written. In most cases the character string is
the keyword used to load the data into the OPM Flow input deck, for
example PORO for the porosity array in the GRID section. In addition,
values either read or calculated by the simulator in the EDIT section
can also be written to the INIT file. Again the keyword or property name
is used as the mnemonic for the character string, for example the PORV,
TRANX keywords etc. If the RPTINIT keyword is not used in the input deck
then a default set of data array are written to the file, in this case
the actual data written is dependent on the model's configuration and
the options being used.

This keyword is not supported by OPM Flow but has no effect on the
results so it will be ignored.

\subsection{RPTISOL -- Activate Isolated Reservoir Number
Reporting}\label{rptisol-activate-isolated-reservoir-number-reporting}

\begin{longtable}[]{@{}
  >{\raggedright\arraybackslash}p{(\linewidth - 14\tabcolsep) * \real{0.1250}}
  >{\raggedright\arraybackslash}p{(\linewidth - 14\tabcolsep) * \real{0.1250}}
  >{\raggedright\arraybackslash}p{(\linewidth - 14\tabcolsep) * \real{0.1250}}
  >{\raggedright\arraybackslash}p{(\linewidth - 14\tabcolsep) * \real{0.1250}}
  >{\raggedright\arraybackslash}p{(\linewidth - 14\tabcolsep) * \real{0.1250}}
  >{\raggedright\arraybackslash}p{(\linewidth - 14\tabcolsep) * \real{0.1250}}
  >{\raggedright\arraybackslash}p{(\linewidth - 14\tabcolsep) * \real{0.1250}}
  >{\raggedright\arraybackslash}p{(\linewidth - 14\tabcolsep) * \real{0.1250}}@{}}
\toprule\noalign{}
\begin{minipage}[b]{\linewidth}\raggedright
RUNSPEC
\end{minipage} & \begin{minipage}[b]{\linewidth}\raggedright
GRID
\end{minipage} & \begin{minipage}[b]{\linewidth}\raggedright
EDIT
\end{minipage} & \begin{minipage}[b]{\linewidth}\raggedright
PROPS
\end{minipage} & \begin{minipage}[b]{\linewidth}\raggedright
REGIONS
\end{minipage} & \begin{minipage}[b]{\linewidth}\raggedright
SOLUTION
\end{minipage} & \begin{minipage}[b]{\linewidth}\raggedright
SUMMARY
\end{minipage} & \begin{minipage}[b]{\linewidth}\raggedright
SCHEDULE
\end{minipage} \\
\midrule\noalign{}
\endhead
\bottomrule\noalign{}
\endlastfoot
\end{longtable}

\subsubsection{Description}\label{description-345}

The RPTISOL keyword activates the isolated reservoir report that
generates an array of isolated region numbers that is printed in the
debug file (*.DBG). The main purpose of this facility is to use the
generated array as input to the ISOLNUM keyword in the GRID section in
conjunction with the Independent Reservoir Regions option. If the model
can be divided into isolated reservoirs then the individual reservoirs
may be solved independently, resulting in increased computational
efficient, compared with solving the model as a whole.

This keyword is not supported by OPM Flow but has no effect on the
results so it will be ignored.

There is no data required for this keyword and there is no terminating
``/'' for this keyword.

\subsubsection{Example}\label{example-233}

\begin{verbatim}
--
--       ACTIVATE ISOLATED RESERVOIR NUMBER REPORTING
--
RPTISOL

\end{verbatim}

The above example activates the isolated reservoir report that generates
an array of isolated region numbers to the debug file (*.DBG).

\subsection{SIGMA -- Dual Porosity Matrix to Fracture Sigma (All
Cells)}\label{sigma-dual-porosity-matrix-to-fracture-sigma-all-cells}

\begin{longtable}[]{@{}
  >{\raggedright\arraybackslash}p{(\linewidth - 14\tabcolsep) * \real{0.1250}}
  >{\raggedright\arraybackslash}p{(\linewidth - 14\tabcolsep) * \real{0.1250}}
  >{\raggedright\arraybackslash}p{(\linewidth - 14\tabcolsep) * \real{0.1250}}
  >{\raggedright\arraybackslash}p{(\linewidth - 14\tabcolsep) * \real{0.1250}}
  >{\raggedright\arraybackslash}p{(\linewidth - 14\tabcolsep) * \real{0.1250}}
  >{\raggedright\arraybackslash}p{(\linewidth - 14\tabcolsep) * \real{0.1250}}
  >{\raggedright\arraybackslash}p{(\linewidth - 14\tabcolsep) * \real{0.1250}}
  >{\raggedright\arraybackslash}p{(\linewidth - 14\tabcolsep) * \real{0.1250}}@{}}
\toprule\noalign{}
\begin{minipage}[b]{\linewidth}\raggedright
RUNSPEC
\end{minipage} & \begin{minipage}[b]{\linewidth}\raggedright
GRID
\end{minipage} & \begin{minipage}[b]{\linewidth}\raggedright
EDIT
\end{minipage} & \begin{minipage}[b]{\linewidth}\raggedright
PROPS
\end{minipage} & \begin{minipage}[b]{\linewidth}\raggedright
REGIONS
\end{minipage} & \begin{minipage}[b]{\linewidth}\raggedright
SOLUTION
\end{minipage} & \begin{minipage}[b]{\linewidth}\raggedright
SUMMARY
\end{minipage} & \begin{minipage}[b]{\linewidth}\raggedright
SCHEDULE
\end{minipage} \\
\midrule\noalign{}
\endhead
\bottomrule\noalign{}
\endlastfoot
\end{longtable}

\subsubsection{Description}\label{description-346}

The SIGMA keyword defines the dual porosity matrix to fracture
transmissibility multiplier, sigma, that is applied to all cells, for
when the Dual Porosity model has been activated by either the DUALPORO
or the DUALPERM keywords in the RUNSPEC section. Sigma (σ) takes into
account the matrix-fracture interface area per unit volume and was
defined by Kazemi et al {[}Kazemi, H., Merrill JR., L. S., Porterfield,
K. L., and Zeman, P. R. ``Numerical Simulation of Water-Oil Flow in
Naturally Fractured Reservoirs,'' paper SPE 5719, Society of Petroleum
Engineers Journal (1976) 16, No.~6, 317-326.{]} to be:

\begin{equation}\phantomsection\label{eq-6-15}{
\mathrm{σ} = 4\left(\frac{1}{{{l}_{x}}^{2}} +\frac{1}{{{l}_{y}}^{2}} +\frac{1}{{{l}_{z}}^{2}} \right)
}\end{equation}

Where lx, ly and lz are not the grid block dimensions in the model in
the respective directions, but the dimensions of the blocks of the
matrix material. In practice, σ is used as a tuning parameter in dual
porosity runs to match reservoir and well performance.

See also the SIGMAV keyword in the GRID section that supplies the sigma
values on an individual cells basis.

This keyword is not supported by OPM Flow but would change the results
if supported so the simulation will be stopped.

\subsection{SIGMAGD -- Dual Porosity Matrix to Fracture Sigma for
Gravity Drainage (All
Cells)}\label{sigmagd-dual-porosity-matrix-to-fracture-sigma-for-gravity-drainage-all-cells}

\begin{longtable}[]{@{}
  >{\raggedright\arraybackslash}p{(\linewidth - 14\tabcolsep) * \real{0.1250}}
  >{\raggedright\arraybackslash}p{(\linewidth - 14\tabcolsep) * \real{0.1250}}
  >{\raggedright\arraybackslash}p{(\linewidth - 14\tabcolsep) * \real{0.1250}}
  >{\raggedright\arraybackslash}p{(\linewidth - 14\tabcolsep) * \real{0.1250}}
  >{\raggedright\arraybackslash}p{(\linewidth - 14\tabcolsep) * \real{0.1250}}
  >{\raggedright\arraybackslash}p{(\linewidth - 14\tabcolsep) * \real{0.1250}}
  >{\raggedright\arraybackslash}p{(\linewidth - 14\tabcolsep) * \real{0.1250}}
  >{\raggedright\arraybackslash}p{(\linewidth - 14\tabcolsep) * \real{0.1250}}@{}}
\toprule\noalign{}
\begin{minipage}[b]{\linewidth}\raggedright
RUNSPEC
\end{minipage} & \begin{minipage}[b]{\linewidth}\raggedright
GRID
\end{minipage} & \begin{minipage}[b]{\linewidth}\raggedright
EDIT
\end{minipage} & \begin{minipage}[b]{\linewidth}\raggedright
PROPS
\end{minipage} & \begin{minipage}[b]{\linewidth}\raggedright
REGIONS
\end{minipage} & \begin{minipage}[b]{\linewidth}\raggedright
SOLUTION
\end{minipage} & \begin{minipage}[b]{\linewidth}\raggedright
SUMMARY
\end{minipage} & \begin{minipage}[b]{\linewidth}\raggedright
SCHEDULE
\end{minipage} \\
\midrule\noalign{}
\endhead
\bottomrule\noalign{}
\endlastfoot
\end{longtable}

\subsubsection{Description}\label{description-347}

The SIGMAGD keyword defines the dual porosity matrix to fracture
transmissibility multiplier, sigma, that is applied to all cells, for
when the Dual Porosity model has been activated by either the DUALPORO
or the DUALPERM keywords in the RUNSPEC section. In addition, the GRAVDR
keyword in the RUNSPEC section should be used to enable the Gravity
Drainage model for the run. Sigma (σ) takes into account the
matrix-fracture interface area per unit volume and was defined by Kazemi
et al {[}Kazemi, H., Merrill JR., L. S., Porterfield, K. L., and Zeman,
P. R. ``Numerical Simulation of Water-Oil Flow in Naturally Fractured
Reservoirs,'' paper SPE 5719, Society of Petroleum Engineers Journal
(1976) 16, No.~6, 317-326.{]} to be:

\begin{equation}\phantomsection\label{eq-6-16}{
\mathrm{σ} = 4\left(\frac{1}{{{l}_{x}}^{2}} +\frac{1}{{{l}_{y}}^{2}} +\frac{1}{{{l}_{z}}^{2}} \right)
}\end{equation}

Where lx, ly and lz are not the grid block dimensions in the model in
the respective directions, but the dimensions of the blocks of the
matrix material. In practice, σ is used as a tuning parameter in dual
porosity runs to match reservoir and well performance.

Note that SIGMAGD keyword data is used for areas being swept by gas and
the SIGMA keyword data is used when the area is being invaded by water.
See also the SIGMAGDV keyword in the GRID section that supplies the
sigma values on an individual cells basis

This keyword is not supported by OPM Flow but would change the results
if supported so the simulation will be stopped.

\subsection{SIGMAGDV -- Dual Porosity Matrix to Fracture Sigma Gravity
Drainage (Individual
Cells)}\label{sigmagdv-dual-porosity-matrix-to-fracture-sigma-gravity-drainage-individual-cells}

\begin{longtable}[]{@{}
  >{\raggedright\arraybackslash}p{(\linewidth - 14\tabcolsep) * \real{0.1250}}
  >{\raggedright\arraybackslash}p{(\linewidth - 14\tabcolsep) * \real{0.1250}}
  >{\raggedright\arraybackslash}p{(\linewidth - 14\tabcolsep) * \real{0.1250}}
  >{\raggedright\arraybackslash}p{(\linewidth - 14\tabcolsep) * \real{0.1250}}
  >{\raggedright\arraybackslash}p{(\linewidth - 14\tabcolsep) * \real{0.1250}}
  >{\raggedright\arraybackslash}p{(\linewidth - 14\tabcolsep) * \real{0.1250}}
  >{\raggedright\arraybackslash}p{(\linewidth - 14\tabcolsep) * \real{0.1250}}
  >{\raggedright\arraybackslash}p{(\linewidth - 14\tabcolsep) * \real{0.1250}}@{}}
\toprule\noalign{}
\begin{minipage}[b]{\linewidth}\raggedright
RUNSPEC
\end{minipage} & \begin{minipage}[b]{\linewidth}\raggedright
GRID
\end{minipage} & \begin{minipage}[b]{\linewidth}\raggedright
EDIT
\end{minipage} & \begin{minipage}[b]{\linewidth}\raggedright
PROPS
\end{minipage} & \begin{minipage}[b]{\linewidth}\raggedright
REGIONS
\end{minipage} & \begin{minipage}[b]{\linewidth}\raggedright
SOLUTION
\end{minipage} & \begin{minipage}[b]{\linewidth}\raggedright
SUMMARY
\end{minipage} & \begin{minipage}[b]{\linewidth}\raggedright
SCHEDULE
\end{minipage} \\
\midrule\noalign{}
\endhead
\bottomrule\noalign{}
\endlastfoot
\end{longtable}

\subsubsection{Description}\label{description-348}

The SIGMAGD keyword defines the dual porosity matrix to fracture
transmissibility multiplier, sigma, that is applied to individual cells,
for when the Dual Porosity model has been activated by either the
DUALPORO or the DUALPERM keywords in the RUNSPEC section. In addition,
the GRAVDR keyword in the RUNSPEC section should be used to enable the
Gravity Drainage model for the run. Sigma (σ) takes into account the
matrix-fracture interface area per unit volume and was defined by Kazemi
et al {[}Kazemi, H., Merrill JR., L. S., Porterfield, K. L., and Zeman,
P. R. ``Numerical Simulation of Water-Oil Flow in Naturally Fractured
Reservoirs,'' paper SPE 5719, Society of Petroleum Engineers Journal
(1976) 16, No.~6, 317-326.{]} to be:

\begin{equation}\phantomsection\label{eq-6-17}{
\mathrm{σ} = 4\left(\frac{1}{{{l}_{x}}^{2}} +\frac{1}{{{l}_{y}}^{2}} +\frac{1}{{{l}_{z}}^{2}} \right)
}\end{equation}

Where lx, ly and lz are not the grid block dimensions in the model in
the respective directions, but the dimensions of the blocks of the
matrix material. In practice, σ is used as a tuning parameter in dual
porosity runs to match reservoir and well performance.

Note that SIGMAGDV keyword data is used for areas being swept by gas and
the SIGMA keyword data is used when the area is being invaded by water.
See also the SIGMAGD keyword in the GRID section that supplies a
constant sigma value for all cells.

This keyword is not supported by OPM Flow but would change the results
if supported so the simulation will be stopped.

\subsection{SIGMAV -- Dual Porosity Matrix to Fracture Sigma (Individual
Cells)}\label{sigmav-dual-porosity-matrix-to-fracture-sigma-individual-cells}

\begin{longtable}[]{@{}
  >{\raggedright\arraybackslash}p{(\linewidth - 14\tabcolsep) * \real{0.1250}}
  >{\raggedright\arraybackslash}p{(\linewidth - 14\tabcolsep) * \real{0.1250}}
  >{\raggedright\arraybackslash}p{(\linewidth - 14\tabcolsep) * \real{0.1250}}
  >{\raggedright\arraybackslash}p{(\linewidth - 14\tabcolsep) * \real{0.1250}}
  >{\raggedright\arraybackslash}p{(\linewidth - 14\tabcolsep) * \real{0.1250}}
  >{\raggedright\arraybackslash}p{(\linewidth - 14\tabcolsep) * \real{0.1250}}
  >{\raggedright\arraybackslash}p{(\linewidth - 14\tabcolsep) * \real{0.1250}}
  >{\raggedright\arraybackslash}p{(\linewidth - 14\tabcolsep) * \real{0.1250}}@{}}
\toprule\noalign{}
\begin{minipage}[b]{\linewidth}\raggedright
RUNSPEC
\end{minipage} & \begin{minipage}[b]{\linewidth}\raggedright
GRID
\end{minipage} & \begin{minipage}[b]{\linewidth}\raggedright
EDIT
\end{minipage} & \begin{minipage}[b]{\linewidth}\raggedright
PROPS
\end{minipage} & \begin{minipage}[b]{\linewidth}\raggedright
REGIONS
\end{minipage} & \begin{minipage}[b]{\linewidth}\raggedright
SOLUTION
\end{minipage} & \begin{minipage}[b]{\linewidth}\raggedright
SUMMARY
\end{minipage} & \begin{minipage}[b]{\linewidth}\raggedright
SCHEDULE
\end{minipage} \\
\midrule\noalign{}
\endhead
\bottomrule\noalign{}
\endlastfoot
\end{longtable}

\subsubsection{Description}\label{description-349}

The SIGMAV keyword defines a dual porosity matrix to fracture
multiplier, sigma, that is applied to individual cells, for when the
Dual Porosity model has been invoked by either the DUALPORO or the
DUALPERM keywords in the RUNSPEC section. Sigma (σ) takes into account
the matrix-fracture interface area per unit volume and was defined by
Kazemi et al {[}Kazemi, H., Merrill JR., L. S., Porterfield, K. L., and
Zeman, P. R. ``Numerical Simulation of Water-Oil Flow in Naturally
Fractured Reservoirs,'' paper SPE 5719, Society of Petroleum Engineers
Journal (1976) 16, No.~6, 317-326.{]} to be:

\begin{equation}\phantomsection\label{eq-6-18}{
\mathrm{σ} = 4\left(\frac{1}{{{l}_{x}}^{2}} +\frac{1}{{{l}_{y}}^{2}} +\frac{1}{{{l}_{z}}^{2}} \right)
}\end{equation}

Where lx, ly and lz are not the grid block dimensions in the model in
the respective directions, but the dimensions of the blocks of the
matrix material. In practice, σ is used as a tuning parameter in dual
porosity runs to match reservoir and well performance.

See also the SIGMA keyword in the GRID section that supplies a constant
sigma to all cells.

This keyword is not supported by OPM Flow but would change the results
if supported so the simulation will be stopped.

\subsection{SKIP -- Activate Skipping of All Keywords and Input
Data}\label{skip-activate-skipping-of-all-keywords-and-input-data-2}

The SKIP keyword activates skipping of all keywords and input data until
the ENDSKIP keyword is encountered. All keywords between the SKIP and
ENDSKIP keywords are ignored.

See SKIP -- Activate Skipping of All Keywords and Input Data in the
GLOBAL section for a full description.

\subsection{SKIP100 -- Activate Skipping of Black-Oil Keywords and Input
Data}\label{skip100-activate-skipping-of-black-oil-keywords-and-input-data-1}

The SKIP100 keyword activates skipping of all keywords and input data by
the commercial black-oil simulator until the ENDSKIP keyword is
encountered. All keywords between the SKIP100 and ENDSKIP keywords are
ignored by the commercial black-oil simulator. The SKIP100 keyword is
ignored by the commercial compositional simulator.

See SKIP100 -- Activate Skipping of Keywords by Black-Oil Simulator in
the GLOBAL section for a full description.

\subsection{SKIP300 -- Activate Skipping of Keywords by Compositional
Simulator}\label{skip300-activate-skipping-of-keywords-by-compositional-simulator-2}

The SKIP300 keyword activates skipping of all keywords and input data by
the commercial compositional simulator until the ENDSKIP keyword is
encountered. All keywords between the SKIP300 and ENDSKIP keywords are
ignored by the commercial compositional simulator. The SKIP300 keyword
is ignored by the commercial black-oil simulator.

See Error: Reference source not found in the GLOBAL section for a full
description.

\subsection{SMULTX -- Multiply Cell Transmissibility in the +X Direction
(Auto-Refinement)}\label{smultx-multiply-cell-transmissibility-in-the-x-direction-auto-refinement}

\begin{longtable}[]{@{}
  >{\raggedright\arraybackslash}p{(\linewidth - 14\tabcolsep) * \real{0.1250}}
  >{\raggedright\arraybackslash}p{(\linewidth - 14\tabcolsep) * \real{0.1250}}
  >{\raggedright\arraybackslash}p{(\linewidth - 14\tabcolsep) * \real{0.1250}}
  >{\raggedright\arraybackslash}p{(\linewidth - 14\tabcolsep) * \real{0.1250}}
  >{\raggedright\arraybackslash}p{(\linewidth - 14\tabcolsep) * \real{0.1250}}
  >{\raggedright\arraybackslash}p{(\linewidth - 14\tabcolsep) * \real{0.1250}}
  >{\raggedright\arraybackslash}p{(\linewidth - 14\tabcolsep) * \real{0.1250}}
  >{\raggedright\arraybackslash}p{(\linewidth - 14\tabcolsep) * \real{0.1250}}@{}}
\toprule\noalign{}
\begin{minipage}[b]{\linewidth}\raggedright
RUNSPEC
\end{minipage} & \begin{minipage}[b]{\linewidth}\raggedright
GRID
\end{minipage} & \begin{minipage}[b]{\linewidth}\raggedright
EDIT
\end{minipage} & \begin{minipage}[b]{\linewidth}\raggedright
PROPS
\end{minipage} & \begin{minipage}[b]{\linewidth}\raggedright
REGIONS
\end{minipage} & \begin{minipage}[b]{\linewidth}\raggedright
SOLUTION
\end{minipage} & \begin{minipage}[b]{\linewidth}\raggedright
SUMMARY
\end{minipage} & \begin{minipage}[b]{\linewidth}\raggedright
SCHEDULE
\end{minipage} \\
\midrule\noalign{}
\endhead
\bottomrule\noalign{}
\endlastfoot
\end{longtable}

\subsubsection{Description}\label{description-350}

SMULTX multiples the transmissibility between two cell faces in the +X
direction between cells in a host base grid and the connecting
auto-refined grid cells, via an array, that is the keyword sets the
transmissibility multiplier of block (Ihost, Jhost, Khost) in the host
base grid, multiplies the transmissibility all the cells (Iauto, Jauto,
Kauto) and (I+Iauto, Jauto, Kauto) in the auto-refinement grid. The Auto
Refinement option must be enabled to use this keyword via the AUTOREF
keyword in the RUNSPEC section.

An alternative to defining the complete array is to use the BOX keyword
to define an area of the grid and then use the SMULTX keyword to set the
multipliers just for the area defined by the BOX keyword (see the
example).

This keyword is not supported by OPM Flow but would change the results
if supported so the simulation will be stopped.

\begin{longtable}[]{@{}
  >{\raggedright\arraybackslash}p{(\linewidth - 6\tabcolsep) * \real{0.2500}}
  >{\raggedright\arraybackslash}p{(\linewidth - 6\tabcolsep) * \real{0.2500}}
  >{\raggedright\arraybackslash}p{(\linewidth - 6\tabcolsep) * \real{0.2500}}
  >{\raggedright\arraybackslash}p{(\linewidth - 6\tabcolsep) * \real{0.2500}}@{}}
\toprule\noalign{}
\begin{minipage}[b]{\linewidth}\raggedright
No.
\end{minipage} & \begin{minipage}[b]{\linewidth}\raggedright
Name
\end{minipage} & \begin{minipage}[b]{\linewidth}\raggedright
Description
\end{minipage} & \begin{minipage}[b]{\linewidth}\raggedright
Default
\end{minipage} \\
\midrule\noalign{}
\endhead
\bottomrule\noalign{}
\endlastfoot
1 & SMULTX+ & SMULTX+ is an array of real positive numbers assigning the
transmissibility multipliers in the +X direction to each cell face in
the model. Repeat counts may be used, for example 20*100.0. & 1.0 \\
Notes: & & & \\
\end{longtable}

\emph{Table 6.118: SMULTX Keyword Description}

See also the MULTX, MULTY and MULTZ keywords for scaling transmissible
between grid cells.

\subsubsection{Example}\label{example-234}

\begin{verbatim}
--
--       DEFINE INPUT BOX FOR EDITING INPUT ARRAYS
--
--       ---------- BOX ---------
--       I1  I2   J1  J2   K1  K2
BOX
         10  10   1   6    1   3                           / DEFINE BOX AREA
--
--       SET SMULTX+ TRANSMISSIBILITY MULTIPLIERS
--
SMULTX
         18*0.300                                          /
--
--       DEFINE END OF INPUT BOX EDITING OF INPUT ARRAYS
--
ENDBOX
\end{verbatim}

The above example defines a 0.3 scaling multiplier for the 18 cells
defined by the preceding BOX statement. The ENDBOX keyword resets the
input box to the full grid.

\subsection{SMULTY -- Multiply Cell Transmissibility in the +Y Direction
(Auto-Refinement)}\label{smulty-multiply-cell-transmissibility-in-the-y-direction-auto-refinement}

\begin{longtable}[]{@{}
  >{\raggedright\arraybackslash}p{(\linewidth - 14\tabcolsep) * \real{0.1250}}
  >{\raggedright\arraybackslash}p{(\linewidth - 14\tabcolsep) * \real{0.1250}}
  >{\raggedright\arraybackslash}p{(\linewidth - 14\tabcolsep) * \real{0.1250}}
  >{\raggedright\arraybackslash}p{(\linewidth - 14\tabcolsep) * \real{0.1250}}
  >{\raggedright\arraybackslash}p{(\linewidth - 14\tabcolsep) * \real{0.1250}}
  >{\raggedright\arraybackslash}p{(\linewidth - 14\tabcolsep) * \real{0.1250}}
  >{\raggedright\arraybackslash}p{(\linewidth - 14\tabcolsep) * \real{0.1250}}
  >{\raggedright\arraybackslash}p{(\linewidth - 14\tabcolsep) * \real{0.1250}}@{}}
\toprule\noalign{}
\begin{minipage}[b]{\linewidth}\raggedright
RUNSPEC
\end{minipage} & \begin{minipage}[b]{\linewidth}\raggedright
GRID
\end{minipage} & \begin{minipage}[b]{\linewidth}\raggedright
EDIT
\end{minipage} & \begin{minipage}[b]{\linewidth}\raggedright
PROPS
\end{minipage} & \begin{minipage}[b]{\linewidth}\raggedright
REGIONS
\end{minipage} & \begin{minipage}[b]{\linewidth}\raggedright
SOLUTION
\end{minipage} & \begin{minipage}[b]{\linewidth}\raggedright
SUMMARY
\end{minipage} & \begin{minipage}[b]{\linewidth}\raggedright
SCHEDULE
\end{minipage} \\
\midrule\noalign{}
\endhead
\bottomrule\noalign{}
\endlastfoot
\end{longtable}

\subsubsection{Description}\label{description-351}

SMULTY multiples the transmissibility between two cell faces in the +Y
direction between cells in a host base grid and the connecting
auto-refined grid cells, via an array, that is the keyword sets the
transmissibility multiplier of block (Ihost, Jhost, Khost) in the host
base grid, multiplies the transmissibility all the cells (Iauto, Jauto,
Kauto) and (Iauto, J+1auto, Kauto) in the auto-refinement grid. The Auto
Refinement option must be enabled to use this keyword via the AUTOREF
keyword in the RUNSPEC section.

An alternative to defining the complete array is to use the BOX keyword
to define an area of the grid and then use the SMULTY keyword to set the
multipliers just for the area defined by the BOX keyword (see the
example).

This keyword is not supported by OPM Flow but would change the results
if supported so the simulation will be stopped.

\begin{longtable}[]{@{}
  >{\raggedright\arraybackslash}p{(\linewidth - 6\tabcolsep) * \real{0.2500}}
  >{\raggedright\arraybackslash}p{(\linewidth - 6\tabcolsep) * \real{0.2500}}
  >{\raggedright\arraybackslash}p{(\linewidth - 6\tabcolsep) * \real{0.2500}}
  >{\raggedright\arraybackslash}p{(\linewidth - 6\tabcolsep) * \real{0.2500}}@{}}
\toprule\noalign{}
\begin{minipage}[b]{\linewidth}\raggedright
No.
\end{minipage} & \begin{minipage}[b]{\linewidth}\raggedright
Name
\end{minipage} & \begin{minipage}[b]{\linewidth}\raggedright
Description
\end{minipage} & \begin{minipage}[b]{\linewidth}\raggedright
Default
\end{minipage} \\
\midrule\noalign{}
\endhead
\bottomrule\noalign{}
\endlastfoot
1 & SMULTY+ & SMULTY+ is an array of real positive numbers assigning the
transmissibility multipliers in the +X direction to each cell face in
the model. Repeat counts may be used, for example 20*100.0. & 1.0 \\
Notes: & & & \\
\end{longtable}

\emph{Table 6.119: SMULTY Keyword Description}

See also the MULTX, MULTY and MULTZ keywords for scaling transmissible
between grid cells.

\subsubsection{Example}\label{example-235}

\begin{verbatim}
--
--       DEFINE INPUT BOX FOR EDITING INPUT ARRAYS
--
--       ---------- BOX ---------
--       I1  I2   J1  J2   K1  K2
BOX
         10  10   1   6    1   3                           / DEFINE BOX AREA
--
--       SET SMULTY+ TRANSMISSIBILITY MULTIPLIERS
--
SMULTY
         18*0.300                                          /
--
--       DEFINE END OF INPUT BOX EDITING OF INPUT ARRAYS
--
ENDBOX
\end{verbatim}

The above example defines a 0.3 scaling multiplier for the 18 cells
defined by the preceding BOX statement. The ENDBOX keyword resets the
input box to the full grid.

\subsection{SMULTZ -- Multiply Cell Transmissibility in the +Z Direction
(Auto-Refinement)}\label{smultz-multiply-cell-transmissibility-in-the-z-direction-auto-refinement}

\begin{longtable}[]{@{}
  >{\raggedright\arraybackslash}p{(\linewidth - 14\tabcolsep) * \real{0.1250}}
  >{\raggedright\arraybackslash}p{(\linewidth - 14\tabcolsep) * \real{0.1250}}
  >{\raggedright\arraybackslash}p{(\linewidth - 14\tabcolsep) * \real{0.1250}}
  >{\raggedright\arraybackslash}p{(\linewidth - 14\tabcolsep) * \real{0.1250}}
  >{\raggedright\arraybackslash}p{(\linewidth - 14\tabcolsep) * \real{0.1250}}
  >{\raggedright\arraybackslash}p{(\linewidth - 14\tabcolsep) * \real{0.1250}}
  >{\raggedright\arraybackslash}p{(\linewidth - 14\tabcolsep) * \real{0.1250}}
  >{\raggedright\arraybackslash}p{(\linewidth - 14\tabcolsep) * \real{0.1250}}@{}}
\toprule\noalign{}
\begin{minipage}[b]{\linewidth}\raggedright
RUNSPEC
\end{minipage} & \begin{minipage}[b]{\linewidth}\raggedright
GRID
\end{minipage} & \begin{minipage}[b]{\linewidth}\raggedright
EDIT
\end{minipage} & \begin{minipage}[b]{\linewidth}\raggedright
PROPS
\end{minipage} & \begin{minipage}[b]{\linewidth}\raggedright
REGIONS
\end{minipage} & \begin{minipage}[b]{\linewidth}\raggedright
SOLUTION
\end{minipage} & \begin{minipage}[b]{\linewidth}\raggedright
SUMMARY
\end{minipage} & \begin{minipage}[b]{\linewidth}\raggedright
SCHEDULE
\end{minipage} \\
\midrule\noalign{}
\endhead
\bottomrule\noalign{}
\endlastfoot
\end{longtable}

\subsubsection{Description}\label{description-352}

SMULTZ multiples the transmissibility between two cell faces in the +Z
direction between cells in a host base grid and the connecting
auto-refined grid cells, via an array, that is the keyword sets the
transmissibility multiplier of block (Ihost, Jhost, Khost) in the host
base grid, multiplies the transmissibility all the cells (Iauto, Jauto,
Kauto) and (Iauto, Jauto, K+1auto) in the auto-refinement grid. The Auto
Refinement option must be enabled to use this keyword via the AUTOREF
keyword in the RUNSPEC section.

An alternative to defining the complete array is to use the BOX keyword
to define an area of the grid and then use the SMULTZ keyword to set the
multipliers just for the area defined by the BOX keyword (see the
example).

This keyword is not supported by OPM Flow but would change the results
if supported so the simulation will be stopped.

\begin{longtable}[]{@{}
  >{\raggedright\arraybackslash}p{(\linewidth - 6\tabcolsep) * \real{0.2500}}
  >{\raggedright\arraybackslash}p{(\linewidth - 6\tabcolsep) * \real{0.2500}}
  >{\raggedright\arraybackslash}p{(\linewidth - 6\tabcolsep) * \real{0.2500}}
  >{\raggedright\arraybackslash}p{(\linewidth - 6\tabcolsep) * \real{0.2500}}@{}}
\toprule\noalign{}
\begin{minipage}[b]{\linewidth}\raggedright
No.
\end{minipage} & \begin{minipage}[b]{\linewidth}\raggedright
Name
\end{minipage} & \begin{minipage}[b]{\linewidth}\raggedright
Description
\end{minipage} & \begin{minipage}[b]{\linewidth}\raggedright
Default
\end{minipage} \\
\midrule\noalign{}
\endhead
\bottomrule\noalign{}
\endlastfoot
1 & SMULTZ+ & SMULTZ+ is an array of real positive numbers assigning the
transmissibility multipliers in the +X direction to each cell face in
the model. Repeat counts may be used, for example 20*100.0. & 1.0 \\
Notes: & & & \\
\end{longtable}

\emph{Table 6.120: SMULTX Keyword Description}

See also the MULTX, MULTY and MULTZ keywords for scaling transmissible
between grid cells.

\subsubsection{Example}\label{example-236}

\begin{verbatim}
--
--       DEFINE INPUT BOX FOR EDITING INPUT ARRAYS
--
--       ---------- BOX ---------
--       I1  I2   J1  J2   K1  K2
BOX
         10  10   1   6    1   3                           / DEFINE BOX AREA
--
--       SET SMULTZ+ TRANSMISSIBILITY MULTIPLIERS
--
SMULTZ
         18*0.300                                          /
--
--       DEFINE END OF INPUT BOX EDITING OF INPUT ARRAYS
--
ENDBOX
\end{verbatim}

The above example defines a 0.3 scaling multiplier for the 18 cells
defined by the preceding BOX statement. The ENDBOX keyword resets the
input box to the full grid.

\subsection{SOLVDIRS -- Define Linear Solver Principal
Directions}\label{solvdirs-define-linear-solver-principal-directions}

\begin{longtable}[]{@{}
  >{\raggedright\arraybackslash}p{(\linewidth - 14\tabcolsep) * \real{0.1250}}
  >{\raggedright\arraybackslash}p{(\linewidth - 14\tabcolsep) * \real{0.1250}}
  >{\raggedright\arraybackslash}p{(\linewidth - 14\tabcolsep) * \real{0.1250}}
  >{\raggedright\arraybackslash}p{(\linewidth - 14\tabcolsep) * \real{0.1250}}
  >{\raggedright\arraybackslash}p{(\linewidth - 14\tabcolsep) * \real{0.1250}}
  >{\raggedright\arraybackslash}p{(\linewidth - 14\tabcolsep) * \real{0.1250}}
  >{\raggedright\arraybackslash}p{(\linewidth - 14\tabcolsep) * \real{0.1250}}
  >{\raggedright\arraybackslash}p{(\linewidth - 14\tabcolsep) * \real{0.1250}}@{}}
\toprule\noalign{}
\begin{minipage}[b]{\linewidth}\raggedright
RUNSPEC
\end{minipage} & \begin{minipage}[b]{\linewidth}\raggedright
GRID
\end{minipage} & \begin{minipage}[b]{\linewidth}\raggedright
EDIT
\end{minipage} & \begin{minipage}[b]{\linewidth}\raggedright
PROPS
\end{minipage} & \begin{minipage}[b]{\linewidth}\raggedright
REGIONS
\end{minipage} & \begin{minipage}[b]{\linewidth}\raggedright
SOLUTION
\end{minipage} & \begin{minipage}[b]{\linewidth}\raggedright
SUMMARY
\end{minipage} & \begin{minipage}[b]{\linewidth}\raggedright
SCHEDULE
\end{minipage} \\
\midrule\noalign{}
\endhead
\bottomrule\noalign{}
\endlastfoot
\end{longtable}

\subsubsection{Description}\label{description-353}

The SOLVDIRS keyword defines the linear solver principal directions,
which should be set to XY, XZ, YX, YZ, ZX, or ZY. The default direction
is based on the direction of the highest transmissibility and SOLVDIRS
allows for over writing the default direction for when linear
convergence of the equations are problematic.

This keyword is not supported by OPM Flow but has no effect on the
results so it will be ignored.

See section 2.2 Running OPM Flow 2023-04 From The Command Line on how to
invoke various numerical schemes via the OPM Flow command line
interface.

\subsection{SOLVNUM -- Define PEBI Grid Correspondence to Solver
Order}\label{solvnum-define-pebi-grid-correspondence-to-solver-order}

\begin{longtable}[]{@{}
  >{\raggedright\arraybackslash}p{(\linewidth - 14\tabcolsep) * \real{0.1250}}
  >{\raggedright\arraybackslash}p{(\linewidth - 14\tabcolsep) * \real{0.1250}}
  >{\raggedright\arraybackslash}p{(\linewidth - 14\tabcolsep) * \real{0.1250}}
  >{\raggedright\arraybackslash}p{(\linewidth - 14\tabcolsep) * \real{0.1250}}
  >{\raggedright\arraybackslash}p{(\linewidth - 14\tabcolsep) * \real{0.1250}}
  >{\raggedright\arraybackslash}p{(\linewidth - 14\tabcolsep) * \real{0.1250}}
  >{\raggedright\arraybackslash}p{(\linewidth - 14\tabcolsep) * \real{0.1250}}
  >{\raggedright\arraybackslash}p{(\linewidth - 14\tabcolsep) * \real{0.1250}}@{}}
\toprule\noalign{}
\begin{minipage}[b]{\linewidth}\raggedright
RUNSPEC
\end{minipage} & \begin{minipage}[b]{\linewidth}\raggedright
GRID
\end{minipage} & \begin{minipage}[b]{\linewidth}\raggedright
EDIT
\end{minipage} & \begin{minipage}[b]{\linewidth}\raggedright
PROPS
\end{minipage} & \begin{minipage}[b]{\linewidth}\raggedright
REGIONS
\end{minipage} & \begin{minipage}[b]{\linewidth}\raggedright
SOLUTION
\end{minipage} & \begin{minipage}[b]{\linewidth}\raggedright
SUMMARY
\end{minipage} & \begin{minipage}[b]{\linewidth}\raggedright
SCHEDULE
\end{minipage} \\
\midrule\noalign{}
\endhead
\bottomrule\noalign{}
\endlastfoot
\end{longtable}

\subsubsection{Description}\label{description-354}

The SOLVNUM defines the unstructured Perpendicular Bisector (``PEBI'')
{[}Heinemann, Z.E. and Brand, C.W. 1988. Gridding Techniques in
Reservoir Simulation. Proc., First Intl. Forum on Reservoir Simulation,
Alpbach, Austria, 339.{]} and {[}Heinemann, Z.E., Brand, C.W., Munka, M.
et al.~1991. Modeling Reservoir Geometry With Irregular Grids. SPE Res
Eng 6 (2): 225--232. SPE-18412-PA. http://dx.doi.org/10.2118/18412-PA{]}
grid correspondence to the nested factorization solver order, for when
the grid has been entered as a PEBI list. This keyword is generated by
an external pre-processing program for generating simulation grids.

This keyword is not supported by OPM Flow but would change the results
if supported so the simulation will be stopped.

\subsection{SPECGRID -- Define the Dimensions of a Corner-Point
Grid}\label{specgrid-define-the-dimensions-of-a-corner-point-grid}

\begin{longtable}[]{@{}
  >{\raggedright\arraybackslash}p{(\linewidth - 14\tabcolsep) * \real{0.1250}}
  >{\raggedright\arraybackslash}p{(\linewidth - 14\tabcolsep) * \real{0.1250}}
  >{\raggedright\arraybackslash}p{(\linewidth - 14\tabcolsep) * \real{0.1250}}
  >{\raggedright\arraybackslash}p{(\linewidth - 14\tabcolsep) * \real{0.1250}}
  >{\raggedright\arraybackslash}p{(\linewidth - 14\tabcolsep) * \real{0.1250}}
  >{\raggedright\arraybackslash}p{(\linewidth - 14\tabcolsep) * \real{0.1250}}
  >{\raggedright\arraybackslash}p{(\linewidth - 14\tabcolsep) * \real{0.1250}}
  >{\raggedright\arraybackslash}p{(\linewidth - 14\tabcolsep) * \real{0.1250}}@{}}
\toprule\noalign{}
\begin{minipage}[b]{\linewidth}\raggedright
RUNSPEC
\end{minipage} & \begin{minipage}[b]{\linewidth}\raggedright
GRID
\end{minipage} & \begin{minipage}[b]{\linewidth}\raggedright
EDIT
\end{minipage} & \begin{minipage}[b]{\linewidth}\raggedright
PROPS
\end{minipage} & \begin{minipage}[b]{\linewidth}\raggedright
REGIONS
\end{minipage} & \begin{minipage}[b]{\linewidth}\raggedright
SOLUTION
\end{minipage} & \begin{minipage}[b]{\linewidth}\raggedright
SUMMARY
\end{minipage} & \begin{minipage}[b]{\linewidth}\raggedright
SCHEDULE
\end{minipage} \\
\midrule\noalign{}
\endhead
\bottomrule\noalign{}
\endlastfoot
\end{longtable}

\subsubsection{Description}\label{description-355}

SPECGRID defines the dimensions of corner-point and radial grids in the
x, y, and z directions as well as the number of reservoirs, where each
reservoir has its own set of corner-point geometry data.

The keyword can only be used with Irregular Corner-Point Grids and
Radial Grids.

\begin{longtable}[]{@{}
  >{\raggedright\arraybackslash}p{(\linewidth - 6\tabcolsep) * \real{0.2500}}
  >{\raggedright\arraybackslash}p{(\linewidth - 6\tabcolsep) * \real{0.2500}}
  >{\raggedright\arraybackslash}p{(\linewidth - 6\tabcolsep) * \real{0.2500}}
  >{\raggedright\arraybackslash}p{(\linewidth - 6\tabcolsep) * \real{0.2500}}@{}}
\toprule\noalign{}
\begin{minipage}[b]{\linewidth}\raggedright
No.
\end{minipage} & \begin{minipage}[b]{\linewidth}\raggedright
Name
\end{minipage} & \begin{minipage}[b]{\linewidth}\raggedright
Description
\end{minipage} & \begin{minipage}[b]{\linewidth}\raggedright
Default
\end{minipage} \\
\midrule\noalign{}
\endhead
\bottomrule\noalign{}
\endlastfoot
1 & NDIVIX & A positive integer value that defines the number of cells
in the X or R direction & 1 \\
2 & NDIVIY & A positive integer value that defines the number of cells
in the Y or THETA direction & 1 \\
3 & NDIVZ & A positive integer value that defines the number of cells in
the Z direction & 1 \\
4 & NUMRES & A positive integer values that defines number of coordinate
data sets, or independent reservoirs in the model. OPM Flow currently
only accepts a single data set, that is the default value of one. & 1 \\
5 & TYPE & A character string set to either T of F that defines the type
of grid to be defined by subsequent keywords: Only the default option F
is supported by OPM Flow. & F \\
Notes: & & & \\
\end{longtable}

\emph{Table 6.121: SPECGRID Keyword Description}

See also the COORD, COORDSYS and ZCORN keywords to fully define an
Irregular Corner-Point Grid.

\subsubsection{Example}\label{example-237}

\begin{verbatim}
--
--       MAX     MAX     MAX     MAX     GRID
--       NDIVIX  NDIVIY  NDIVIZ  NUMRES  TYPE
SPECGRID
         46      112     22      1       F                                     /
\end{verbatim}

The above example defines a 46 x 112 x 22 grid with one set of irregular
corner-point data.

\subsection{THCGAS -- Define Gas Phase Thermal Conductivity for All
Cells}\label{thcgas-define-gas-phase-thermal-conductivity-for-all-cells}

\begin{longtable}[]{@{}
  >{\raggedright\arraybackslash}p{(\linewidth - 14\tabcolsep) * \real{0.1250}}
  >{\raggedright\arraybackslash}p{(\linewidth - 14\tabcolsep) * \real{0.1250}}
  >{\raggedright\arraybackslash}p{(\linewidth - 14\tabcolsep) * \real{0.1250}}
  >{\raggedright\arraybackslash}p{(\linewidth - 14\tabcolsep) * \real{0.1250}}
  >{\raggedright\arraybackslash}p{(\linewidth - 14\tabcolsep) * \real{0.1250}}
  >{\raggedright\arraybackslash}p{(\linewidth - 14\tabcolsep) * \real{0.1250}}
  >{\raggedright\arraybackslash}p{(\linewidth - 14\tabcolsep) * \real{0.1250}}
  >{\raggedright\arraybackslash}p{(\linewidth - 14\tabcolsep) * \real{0.1250}}@{}}
\toprule\noalign{}
\begin{minipage}[b]{\linewidth}\raggedright
RUNSPEC
\end{minipage} & \begin{minipage}[b]{\linewidth}\raggedright
GRID
\end{minipage} & \begin{minipage}[b]{\linewidth}\raggedright
EDIT
\end{minipage} & \begin{minipage}[b]{\linewidth}\raggedright
PROPS
\end{minipage} & \begin{minipage}[b]{\linewidth}\raggedright
REGIONS
\end{minipage} & \begin{minipage}[b]{\linewidth}\raggedright
SOLUTION
\end{minipage} & \begin{minipage}[b]{\linewidth}\raggedright
SUMMARY
\end{minipage} & \begin{minipage}[b]{\linewidth}\raggedright
SCHEDULE
\end{minipage} \\
\midrule\noalign{}
\endhead
\bottomrule\noalign{}
\endlastfoot
\end{longtable}

\subsubsection{Description}\label{description-356}

The THCGAS keyword defines the gas phase thermal conductivity for when
the thermal calculation is activated by the THERMAL keyword in the
RUNSPEC section, and should be used in conjunction with THCROCK keyword
in the GRID section.

This keyword can only be used if the thermal option has been activated
by the THERMAL keyword in the RUNSPEC section. Note this is different to
the commercial simulator that uses the TEMP keyword in the RUNSPEC
section to activate the black-oil thermal model.

\begin{longtable}[]{@{}
  >{\raggedright\arraybackslash}p{(\linewidth - 6\tabcolsep) * \real{0.2500}}
  >{\raggedright\arraybackslash}p{(\linewidth - 6\tabcolsep) * \real{0.2500}}
  >{\raggedright\arraybackslash}p{(\linewidth - 6\tabcolsep) * \real{0.2500}}
  >{\raggedright\arraybackslash}p{(\linewidth - 6\tabcolsep) * \real{0.2500}}@{}}
\toprule\noalign{}
\begin{minipage}[b]{\linewidth}\raggedright
No.
\end{minipage} & \begin{minipage}[b]{\linewidth}\raggedright
Name
\end{minipage} & \begin{minipage}[b]{\linewidth}\raggedright
Description
\end{minipage} & \begin{minipage}[b]{\linewidth}\raggedright
Default
\end{minipage} \\
\midrule\noalign{}
\endhead
\bottomrule\noalign{}
\endlastfoot
Field & Metric & Laboratory & \\
1 & THCGAS & THCGAS is an array of real positive numbers that define the
thermal conductivity of the gas phase in each grid block. Repeat counts
may be used, for example 3000*20.0 & None \\
Btu/ft/day/°R & kJ/m/day/K & J/cm/hr/K & \\
Notes: & & & \\
\end{longtable}

\emph{Table 6.122: THCGAS Keyword Description}

Note that there two ways to define the rock and in situ fluids thermal
conductivity:

\begin{verbatim}
- Either by using the THCONR keyword to define the combined rock and fluid conductivity, and optionally the THCONSF keyword in the GRID section, or
- by specifying the rock and fluid conductivities individually using the THCROCK, THCOIL, THCGAS, and THCWATER keywords in the GRID section.
\end{verbatim}

Hence, the THCROCK and THCONR keywords are mutually exclusive.

Here, the THCGAS keyword is used in conjunction with the other thermal
conductivity arrays to calculate the porosity weighted thermal
conductivity of a grid block using:

\begin{equation}\phantomsection\label{eq-6-19}{
\text{Average Thermal Conductivity}=\frac{\mathit{PORO}\times \left(\mathit{THCOIL}+\mathit{THCGAS}+\mathit{THCWATER}+\mathit{THCSOLID}\right)}{\text{        NUMBER OF PHASES IN THE MODEL}}\times \left(1-\mathit{PORO}\right)\times \mathit{THCROCK}
}\end{equation}

See also the THCOIL, and THCWATER, and THCROCK keywords in the GRID
section. The commercial compositional simulator's THCSOLID keyword is
not supported or required by OPM Flow.

\subsubsection{Example}\label{example-238}

\begin{verbatim}
--
--       DEFINE GRID BLOCK GAS PHASE THERMAL CONDUCTIVITY
–        FOR ALL CELLS (BASED ON NX x NY x NZ = 300)
--
THCGAS
         300*20.0                                                              /
\end{verbatim}

The above example defines the gas phase thermal conductivity of 20.0 for
each cell in the 300 grid block model as defined by the DIMENS keyword
in the RUNSPEC section.

\subsection{THCOIL -- Define Oil Phase Thermal Conductivity for All
Cells}\label{thcoil-define-oil-phase-thermal-conductivity-for-all-cells}

\begin{longtable}[]{@{}
  >{\raggedright\arraybackslash}p{(\linewidth - 14\tabcolsep) * \real{0.1250}}
  >{\raggedright\arraybackslash}p{(\linewidth - 14\tabcolsep) * \real{0.1250}}
  >{\raggedright\arraybackslash}p{(\linewidth - 14\tabcolsep) * \real{0.1250}}
  >{\raggedright\arraybackslash}p{(\linewidth - 14\tabcolsep) * \real{0.1250}}
  >{\raggedright\arraybackslash}p{(\linewidth - 14\tabcolsep) * \real{0.1250}}
  >{\raggedright\arraybackslash}p{(\linewidth - 14\tabcolsep) * \real{0.1250}}
  >{\raggedright\arraybackslash}p{(\linewidth - 14\tabcolsep) * \real{0.1250}}
  >{\raggedright\arraybackslash}p{(\linewidth - 14\tabcolsep) * \real{0.1250}}@{}}
\toprule\noalign{}
\begin{minipage}[b]{\linewidth}\raggedright
RUNSPEC
\end{minipage} & \begin{minipage}[b]{\linewidth}\raggedright
GRID
\end{minipage} & \begin{minipage}[b]{\linewidth}\raggedright
EDIT
\end{minipage} & \begin{minipage}[b]{\linewidth}\raggedright
PROPS
\end{minipage} & \begin{minipage}[b]{\linewidth}\raggedright
REGIONS
\end{minipage} & \begin{minipage}[b]{\linewidth}\raggedright
SOLUTION
\end{minipage} & \begin{minipage}[b]{\linewidth}\raggedright
SUMMARY
\end{minipage} & \begin{minipage}[b]{\linewidth}\raggedright
SCHEDULE
\end{minipage} \\
\midrule\noalign{}
\endhead
\bottomrule\noalign{}
\endlastfoot
\end{longtable}

\subsubsection{Description}\label{description-357}

The THCOIL keyword defines the oil phase thermal conductivity for when
the thermal calculation is activated by the THERMAL keyword in the
RUNSPEC section, and should be used in conjunction with THCROCK keyword
in the GRID section.

This keyword can only be used if the thermal option has been activated
by the THERMAL keyword in the RUNSPEC section. Note this is different to
the commercial simulator that uses the TEMP keyword in the RUNSPEC
section to activate the black-oil thermal model.

\begin{longtable}[]{@{}
  >{\raggedright\arraybackslash}p{(\linewidth - 6\tabcolsep) * \real{0.2500}}
  >{\raggedright\arraybackslash}p{(\linewidth - 6\tabcolsep) * \real{0.2500}}
  >{\raggedright\arraybackslash}p{(\linewidth - 6\tabcolsep) * \real{0.2500}}
  >{\raggedright\arraybackslash}p{(\linewidth - 6\tabcolsep) * \real{0.2500}}@{}}
\toprule\noalign{}
\begin{minipage}[b]{\linewidth}\raggedright
No.
\end{minipage} & \begin{minipage}[b]{\linewidth}\raggedright
Name
\end{minipage} & \begin{minipage}[b]{\linewidth}\raggedright
Description
\end{minipage} & \begin{minipage}[b]{\linewidth}\raggedright
Default
\end{minipage} \\
\midrule\noalign{}
\endhead
\bottomrule\noalign{}
\endlastfoot
Field & Metric & Laboratory & \\
1 & THCOIL & THCOIL is an array of real positive numbers that define the
thermal conductivity of the oil phase in each grid block. Repeat counts
may be used, for example 3000*20.0 & None \\
Btu/ft/day/°R & kJ/m/day/K & J/cm/hr/K & \\
Notes: & & & \\
\end{longtable}

\emph{Table 6.123: THCOIL Keyword Description}

Note that there two ways to define the rock and in situ fluids thermal
conductivity:

\begin{verbatim}
- Either by using the THCONR keyword to define the combined rock and fluid conductivity, and optionally the THCONSF keyword in the GRID section, or
- by specifying the rock and fluid conductivities individually using the THCROCK, THCOIL, THCGAS, and THCWATER keywords in the GRID section.
\end{verbatim}

Hence, the THCROCK and THCONR keywords are mutually exclusive.

Here, the THCOIL keyword is used in conjunction with the other thermal
conductivity arrays to calculate the porosity weighted thermal
conductivity of a grid block using:

\begin{equation}\phantomsection\label{eq-6-20}{
\text{Average Thermal Conductivity}=\frac{\mathit{PORO}\times \left(\mathit{THCOIL}+\mathit{THCGAS}+\mathit{THCWATER}+\mathit{THCSOLID}\right)}{\text{        NUMBER OF PHASES IN THE MODEL}}\times \left(1-\text{ PORO}\right)\times \text{THCROCK}
}\end{equation}

See also the THCGAS, and THCWATER, and THCROCK keywords in the GRID
section. The commercial compositional simulator's THCSOLID keyword is
not supported or required by OPM Flow.

\subsubsection{Example}\label{example-239}

\begin{verbatim}
--
--       DEFINE GRID BLOCK OIL PHASE THERMAL CONDUCTIVITY
–        FOR ALL CELLS (BASED ON NX x NY x NZ = 300)
--
THCOIL
         300*20.0                                                              /
\end{verbatim}

The above example defines the oil phase thermal conductivity of 20.0 for
each cell in the 300 grid block model, as defined by the DIMENS keyword
in the RUNSPEC section.

\subsection{THCONR -- Define Rock and Fluid Thermal Conductivity for All
Cells}\label{thconr-define-rock-and-fluid-thermal-conductivity-for-all-cells}

\begin{longtable}[]{@{}
  >{\raggedright\arraybackslash}p{(\linewidth - 14\tabcolsep) * \real{0.1250}}
  >{\raggedright\arraybackslash}p{(\linewidth - 14\tabcolsep) * \real{0.1250}}
  >{\raggedright\arraybackslash}p{(\linewidth - 14\tabcolsep) * \real{0.1250}}
  >{\raggedright\arraybackslash}p{(\linewidth - 14\tabcolsep) * \real{0.1250}}
  >{\raggedright\arraybackslash}p{(\linewidth - 14\tabcolsep) * \real{0.1250}}
  >{\raggedright\arraybackslash}p{(\linewidth - 14\tabcolsep) * \real{0.1250}}
  >{\raggedright\arraybackslash}p{(\linewidth - 14\tabcolsep) * \real{0.1250}}
  >{\raggedright\arraybackslash}p{(\linewidth - 14\tabcolsep) * \real{0.1250}}@{}}
\toprule\noalign{}
\begin{minipage}[b]{\linewidth}\raggedright
RUNSPEC
\end{minipage} & \begin{minipage}[b]{\linewidth}\raggedright
GRID
\end{minipage} & \begin{minipage}[b]{\linewidth}\raggedright
EDIT
\end{minipage} & \begin{minipage}[b]{\linewidth}\raggedright
PROPS
\end{minipage} & \begin{minipage}[b]{\linewidth}\raggedright
REGIONS
\end{minipage} & \begin{minipage}[b]{\linewidth}\raggedright
SOLUTION
\end{minipage} & \begin{minipage}[b]{\linewidth}\raggedright
SUMMARY
\end{minipage} & \begin{minipage}[b]{\linewidth}\raggedright
SCHEDULE
\end{minipage} \\
\midrule\noalign{}
\endhead
\bottomrule\noalign{}
\endlastfoot
\end{longtable}

\subsubsection{Description}\label{description-358}

The THCONR keyword defines the reservoir rock plus fluid thermal
conductivity for all cells for when the thermal calculation is activated
by the THERMAL keywords in the RUNSPEC section.

This keyword can only be used if the thermal option has been activated
by the THERMAL keyword in the RUNSPEC section. Note this is different to
the commercial simulator that uses the TEMP keyword in the RUNSPEC
section to activate the black-oil thermal model.

\begin{longtable}[]{@{}
  >{\raggedright\arraybackslash}p{(\linewidth - 6\tabcolsep) * \real{0.2500}}
  >{\raggedright\arraybackslash}p{(\linewidth - 6\tabcolsep) * \real{0.2500}}
  >{\raggedright\arraybackslash}p{(\linewidth - 6\tabcolsep) * \real{0.2500}}
  >{\raggedright\arraybackslash}p{(\linewidth - 6\tabcolsep) * \real{0.2500}}@{}}
\toprule\noalign{}
\begin{minipage}[b]{\linewidth}\raggedright
No.
\end{minipage} & \begin{minipage}[b]{\linewidth}\raggedright
Name
\end{minipage} & \begin{minipage}[b]{\linewidth}\raggedright
Description
\end{minipage} & \begin{minipage}[b]{\linewidth}\raggedright
Default
\end{minipage} \\
\midrule\noalign{}
\endhead
\bottomrule\noalign{}
\endlastfoot
Field & Metric & Laboratory & \\
1 & THCONR & THCONR is an array of real positive numbers that define the
combined rock and fluid conductivity of a grid block. Repeat counts may
be used, for example 3000*25.0 & None \\
Btu/ft/day/°R & kJ/m/day/K & J/cm/hr/K & \\
Notes: & & & \\
\end{longtable}

\emph{Table 6.124: THCONR Keyword Description}

Note that there two ways to define the rock and in situ fluids thermal
conductivity:

\begin{verbatim}
- Either by using the THCONR keyword to define the combined rock and fluid conductivity, and optionally the THCONSF keyword in the GRID section, or
- by specifying the rock and fluid conductivities individually using the THCROCK, THCOIL, THCGAS, and THCWATER keywords in the GRID section.
\end{verbatim}

Hence, the THCROCK and THCONR keywords are mutually exclusive.

\subsubsection{Example}\label{example-240}

\begin{verbatim}
--
--       DEFINE GRID BLOCK ROCK-FLUID THERMAL CONDUCTIVITY
–        FOR ALL CELLS (BASED ON NX x NY x NZ = 300)
--
THCONR
         300*25.0                                                              /
\end{verbatim}

The above example defines the combined rock and fluid thermal
conductivity of 25.0 for each cell in the 300 grid block model, as
defined by the DIMENS keyword in the RUNSPEC section.

\subsection{THCONSF -- Define Gas Saturation Dependent Thermal
Conductivity Scaling Factor for All
Cells}\label{thconsf-define-gas-saturation-dependent-thermal-conductivity-scaling-factor-for-all-cells}

\begin{longtable}[]{@{}
  >{\raggedright\arraybackslash}p{(\linewidth - 14\tabcolsep) * \real{0.1250}}
  >{\raggedright\arraybackslash}p{(\linewidth - 14\tabcolsep) * \real{0.1250}}
  >{\raggedright\arraybackslash}p{(\linewidth - 14\tabcolsep) * \real{0.1250}}
  >{\raggedright\arraybackslash}p{(\linewidth - 14\tabcolsep) * \real{0.1250}}
  >{\raggedright\arraybackslash}p{(\linewidth - 14\tabcolsep) * \real{0.1250}}
  >{\raggedright\arraybackslash}p{(\linewidth - 14\tabcolsep) * \real{0.1250}}
  >{\raggedright\arraybackslash}p{(\linewidth - 14\tabcolsep) * \real{0.1250}}
  >{\raggedright\arraybackslash}p{(\linewidth - 14\tabcolsep) * \real{0.1250}}@{}}
\toprule\noalign{}
\begin{minipage}[b]{\linewidth}\raggedright
RUNSPEC
\end{minipage} & \begin{minipage}[b]{\linewidth}\raggedright
GRID
\end{minipage} & \begin{minipage}[b]{\linewidth}\raggedright
EDIT
\end{minipage} & \begin{minipage}[b]{\linewidth}\raggedright
PROPS
\end{minipage} & \begin{minipage}[b]{\linewidth}\raggedright
REGIONS
\end{minipage} & \begin{minipage}[b]{\linewidth}\raggedright
SOLUTION
\end{minipage} & \begin{minipage}[b]{\linewidth}\raggedright
SUMMARY
\end{minipage} & \begin{minipage}[b]{\linewidth}\raggedright
SCHEDULE
\end{minipage} \\
\midrule\noalign{}
\endhead
\bottomrule\noalign{}
\endlastfoot
\end{longtable}

\subsubsection{Description}\label{description-359}

The THCONSF keyword defines a gas saturation dependent scaling factor to
the fluid and reservoir rock thermal conductivities entered via the
THCONR keyword in the GRID section, for when the thermal calculation is
activated by the THERMAL keyword in the RUNSPEC.

This keyword can only be used if the thermal option has been activated
by the THERMAL keyword in the RUNSPEC section. Note this is different to
the commercial simulator that uses the TEMP keyword in the RUNSPEC
section to activate the black-oil thermal model.

\begin{longtable}[]{@{}
  >{\raggedright\arraybackslash}p{(\linewidth - 6\tabcolsep) * \real{0.2500}}
  >{\raggedright\arraybackslash}p{(\linewidth - 6\tabcolsep) * \real{0.2500}}
  >{\raggedright\arraybackslash}p{(\linewidth - 6\tabcolsep) * \real{0.2500}}
  >{\raggedright\arraybackslash}p{(\linewidth - 6\tabcolsep) * \real{0.2500}}@{}}
\toprule\noalign{}
\begin{minipage}[b]{\linewidth}\raggedright
No.
\end{minipage} & \begin{minipage}[b]{\linewidth}\raggedright
Name
\end{minipage} & \begin{minipage}[b]{\linewidth}\raggedright
Description
\end{minipage} & \begin{minipage}[b]{\linewidth}\raggedright
Default
\end{minipage} \\
\midrule\noalign{}
\endhead
\bottomrule\noalign{}
\endlastfoot
Field & Metric & Laboratory & \\
1 & THCONSF & THCONSF is an array of real positive numbers, greater than
zero and less than or equal to one, that define the gas saturation
dependent scaling factor that is applied to the THCONR data, entered via
the THCONR keyword, to adjust the thermal conductivity of the reservoir
cells in each grid block. Repeat counts may be used, for example
3000*0.15 & None \\
dimensionless & dimensionless & dimensionless & \\
Notes: & & & \\
\end{longtable}

\emph{Table 6.125: THCROCK Keyword Description}

Note that there two ways to define the rock and in situ fluids thermal
conductivity:

\begin{verbatim}
- Either by using the THCONR keyword to define the combined rock and fluid conductivity, and optionally the THCONSF keyword in the GRID section, or
- by specifying the rock and fluid conductivities individually using the THCROCK, THCOIL, THCGAS, and THCWATER keywords in the GRID section.
\end{verbatim}

Hence, the THCROCK and THCONR keywords are mutually exclusive.

Here, the THCONSF keyword defines a scaling factor which is a function
of the gas saturation that scales a cells total thermal conductivity
(reservoir fluids plus reservoir rock) entered via the THCONR keyword in
the GRID section. This combination of keywords, THCONSF and THCONR
implies that the oil and water phase thermal conductivities are
saturation independent with respect to the liquid phase, and that only
the gas saturation influences a cell's thermal conductivity as entered
via the THCONR keyword.

Thus, THCONSF scales the THCONR values via a multiplier Ω, by:

\begin{equation}\phantomsection\label{eq-6-21}{
{\mathrm{Ω}}_{i,j,k} = {\left(1-\text{THCONSF x Gas Saturation}\right)}_{i,j,k}
}\end{equation}

See also the THCGAS, THCOIL, THCWATER and THROCK keywords in the GRID
section, for an alternative way to enter the thermal conductivity
properties. However, the THCONSF keyword cannot be used with the THCGAS,
THCOIL, THCWATER and THCROCK keywords. Secondly, the commercial
compositional simulator's THCSOLID keyword is not supported or required
by OPM Flow.

\subsubsection{Example}\label{example-241}

\begin{verbatim}
--
--       DEFINE GRID SGAS DEPENDENT SCALING FACTOR FOR THE THCONR ARRAY                            --       FOR ALL CELLS (BASED ON NX x NY x NZ = 300)
--       (OPM FLOW THERMAL OPTION ONLY)
--
THCONSF
         300*0.12                                                              /
\end{verbatim}

The above example defines the gas saturation thermal conductivity
scaling factor to be applied to the THCONR to be 0.12 for all 300 cells
in the model, as defined by the DIMENS keyword in the RUNSPEC section.

\subsection{THCROCK -- Define Reservoir Rock Thermal Conductivity for
All
Cells}\label{thcrock-define-reservoir-rock-thermal-conductivity-for-all-cells}

\begin{longtable}[]{@{}
  >{\raggedright\arraybackslash}p{(\linewidth - 14\tabcolsep) * \real{0.1250}}
  >{\raggedright\arraybackslash}p{(\linewidth - 14\tabcolsep) * \real{0.1250}}
  >{\raggedright\arraybackslash}p{(\linewidth - 14\tabcolsep) * \real{0.1250}}
  >{\raggedright\arraybackslash}p{(\linewidth - 14\tabcolsep) * \real{0.1250}}
  >{\raggedright\arraybackslash}p{(\linewidth - 14\tabcolsep) * \real{0.1250}}
  >{\raggedright\arraybackslash}p{(\linewidth - 14\tabcolsep) * \real{0.1250}}
  >{\raggedright\arraybackslash}p{(\linewidth - 14\tabcolsep) * \real{0.1250}}
  >{\raggedright\arraybackslash}p{(\linewidth - 14\tabcolsep) * \real{0.1250}}@{}}
\toprule\noalign{}
\begin{minipage}[b]{\linewidth}\raggedright
RUNSPEC
\end{minipage} & \begin{minipage}[b]{\linewidth}\raggedright
GRID
\end{minipage} & \begin{minipage}[b]{\linewidth}\raggedright
EDIT
\end{minipage} & \begin{minipage}[b]{\linewidth}\raggedright
PROPS
\end{minipage} & \begin{minipage}[b]{\linewidth}\raggedright
REGIONS
\end{minipage} & \begin{minipage}[b]{\linewidth}\raggedright
SOLUTION
\end{minipage} & \begin{minipage}[b]{\linewidth}\raggedright
SUMMARY
\end{minipage} & \begin{minipage}[b]{\linewidth}\raggedright
SCHEDULE
\end{minipage} \\
\midrule\noalign{}
\endhead
\bottomrule\noalign{}
\endlastfoot
\end{longtable}

\subsubsection{Description}\label{description-360}

The THCROCK keyword defines the reservoir rock thermal conductivity for
when the thermal calculation is activated by the THERMAL keyword in the
RUNSPEC section.

This keyword can only be used if the thermal option has been activated
by the THERMAL keyword in the RUNSPEC section. Note this is different to
the commercial simulator that uses the TEMP keyword in the RUNSPEC
section to activate the black-oil thermal model.

\begin{longtable}[]{@{}
  >{\raggedright\arraybackslash}p{(\linewidth - 6\tabcolsep) * \real{0.2500}}
  >{\raggedright\arraybackslash}p{(\linewidth - 6\tabcolsep) * \real{0.2500}}
  >{\raggedright\arraybackslash}p{(\linewidth - 6\tabcolsep) * \real{0.2500}}
  >{\raggedright\arraybackslash}p{(\linewidth - 6\tabcolsep) * \real{0.2500}}@{}}
\toprule\noalign{}
\begin{minipage}[b]{\linewidth}\raggedright
No.
\end{minipage} & \begin{minipage}[b]{\linewidth}\raggedright
Name
\end{minipage} & \begin{minipage}[b]{\linewidth}\raggedright
Description
\end{minipage} & \begin{minipage}[b]{\linewidth}\raggedright
Default
\end{minipage} \\
\midrule\noalign{}
\endhead
\bottomrule\noalign{}
\endlastfoot
Field & Metric & Laboratory & \\
1 & THCROCK & THCROCK is an array of real positive numbers that define
the thermal conductivity of the reservoir rock in each grid block.
Repeat counts may be used, for example 3000*20.0 & None \\
Btu/ft/day/°R & kJ/m/day/K & J/cm/hr/K & \\
Notes: & & & \\
\end{longtable}

\emph{Table 6.126: THCROCK Keyword Description}

Note that there two ways to define the rock and in situ fluids thermal
conductivity:

\begin{verbatim}
- Either by using the THCONR keyword to define the combined rock and fluid conductivity, and optionally the THCONSF keyword in the GRID section, or
- by specifying the rock and fluid conductivities individually using the THCROCK, THCOIL, THCGAS, and THCWATER keywords in the GRID section.
\end{verbatim}

Hence, the THCROCK and THCONR keywords are mutually exclusive.

Here, the THCROCK keyword is used in conjunction with the other thermal
conductivity arrays to calculate the porosity weighted thermal
conductivity of a grid block using:

\begin{equation}\phantomsection\label{eq-6-22}{
\text{Average Thermal Conductivity}=\frac{\mathit{PORO}\times \left(\mathit{THCOIL}+\mathit{THCGAS}+\mathit{THCWATER}+\mathit{THCSOLID}\right)}{\text{        NUMBER OF PHASES IN THE MODEL}}\times \left(1-\mathit{PORO}\right)\times \mathit{THCROCK}
}\end{equation}

See also the THCOIL, and THCGAS, and THCWATER keywords in the GRID
section. The commercial compositional simulator's THCSOLID keyword is
not supported or required by OPM Flow.

\subsubsection{Example}\label{example-242}

\begin{verbatim}
--
--       DEFINE GRID BLOCK RESERVOIR ROCK THERMAL CONDUCTIVITY
–        FOR ALL CELLS (BASED ON NX x NY x NZ = 300)
--
THCROCK
         300*20.0                                                              /
\end{verbatim}

The above example defines the reservoir rock thermal conductivity of
20.0 for each cell in the 300 grid block model, as defined by the DIMENS
keyword in the RUNSPEC section.

\subsection{THCWATER -- Define Water Phase Thermal Conductivity for All
Cells}\label{thcwater-define-water-phase-thermal-conductivity-for-all-cells}

\begin{longtable}[]{@{}
  >{\raggedright\arraybackslash}p{(\linewidth - 14\tabcolsep) * \real{0.1250}}
  >{\raggedright\arraybackslash}p{(\linewidth - 14\tabcolsep) * \real{0.1250}}
  >{\raggedright\arraybackslash}p{(\linewidth - 14\tabcolsep) * \real{0.1250}}
  >{\raggedright\arraybackslash}p{(\linewidth - 14\tabcolsep) * \real{0.1250}}
  >{\raggedright\arraybackslash}p{(\linewidth - 14\tabcolsep) * \real{0.1250}}
  >{\raggedright\arraybackslash}p{(\linewidth - 14\tabcolsep) * \real{0.1250}}
  >{\raggedright\arraybackslash}p{(\linewidth - 14\tabcolsep) * \real{0.1250}}
  >{\raggedright\arraybackslash}p{(\linewidth - 14\tabcolsep) * \real{0.1250}}@{}}
\toprule\noalign{}
\begin{minipage}[b]{\linewidth}\raggedright
RUNSPEC
\end{minipage} & \begin{minipage}[b]{\linewidth}\raggedright
GRID
\end{minipage} & \begin{minipage}[b]{\linewidth}\raggedright
EDIT
\end{minipage} & \begin{minipage}[b]{\linewidth}\raggedright
PROPS
\end{minipage} & \begin{minipage}[b]{\linewidth}\raggedright
REGIONS
\end{minipage} & \begin{minipage}[b]{\linewidth}\raggedright
SOLUTION
\end{minipage} & \begin{minipage}[b]{\linewidth}\raggedright
SUMMARY
\end{minipage} & \begin{minipage}[b]{\linewidth}\raggedright
SCHEDULE
\end{minipage} \\
\midrule\noalign{}
\endhead
\bottomrule\noalign{}
\endlastfoot
\end{longtable}

\subsubsection{Description}\label{description-361}

The THCWATER keyword defines the water phase thermal conductivity for
when the thermal calculation is activated by the THERMAL keyword in the
RUNSPEC section, and should be used in conjunction with THCROCK keyword
in the GRID section.

This keyword can only be used if the thermal option has been activated
by the THERMAL keyword in the RUNSPEC section. Note this is different to
the commercial simulator that uses the TEMP keyword in the RUNSPEC
section to activate the black-oil thermal model.

\begin{longtable}[]{@{}
  >{\raggedright\arraybackslash}p{(\linewidth - 6\tabcolsep) * \real{0.2500}}
  >{\raggedright\arraybackslash}p{(\linewidth - 6\tabcolsep) * \real{0.2500}}
  >{\raggedright\arraybackslash}p{(\linewidth - 6\tabcolsep) * \real{0.2500}}
  >{\raggedright\arraybackslash}p{(\linewidth - 6\tabcolsep) * \real{0.2500}}@{}}
\toprule\noalign{}
\begin{minipage}[b]{\linewidth}\raggedright
No.
\end{minipage} & \begin{minipage}[b]{\linewidth}\raggedright
Name
\end{minipage} & \begin{minipage}[b]{\linewidth}\raggedright
Description
\end{minipage} & \begin{minipage}[b]{\linewidth}\raggedright
Default
\end{minipage} \\
\midrule\noalign{}
\endhead
\bottomrule\noalign{}
\endlastfoot
Field & Metric & Laboratory & \\
1 & THCWATER & THCWATER is an array of real positive numbers that define
the thermal conductivity of the water phase in each grid block. Repeat
counts may be used, for example 3000*20.0 & None \\
Btu/ft/day/°R & kJ/m/day/K & J/cm/hr/K & \\
Notes: & & & \\
\end{longtable}

\emph{Table 6.127: THCWATER Keyword Description}

Note that there two ways to define the rock and in situ fluids thermal
conductivity:

\begin{verbatim}
- Either by using the THCONR keyword to define the combined rock and fluid conductivity, and optionally the THCONSF keyword in the GRID section, or
- by specifying the rock and fluid conductivities individually using the THCROCK, THCOIL, THCGAS, and THCWATER keywords in the GRID section.
\end{verbatim}

Hence, the THCROCK and THCONR keywords are mutually exclusive.

Here, the THCWATER keyword is used in conjunction with the other thermal
conductivity arrays to calculate the porosity weighted thermal
conductivity of a grid block using:

\begin{equation}\phantomsection\label{eq-6-23}{
\text{Average Thermal Conductivity}=\frac{\mathit{PORO}\times \left(\mathit{THCOIL}+\mathit{THCGAS}+\mathit{THCWATER}+\mathit{THCSOLID}\right)}{\text{        NUMBER OF PHASES IN THE MODEL}}\times \left(1-\mathit{PORO}\right)\times \mathit{THCROCK}
}\end{equation}

See also the THCGAS, and THCOIL, and THCROCK keywords in the GRID
section. The commercial compositional simulator's THCSOLID keyword is
not supported or required by OPM Flow

\subsubsection{Example}\label{example-243}

\begin{verbatim}
--
--       DEFINE GRID BLOCK WATER PHASE THERMAL CONDUCTIVITY
–        FOR ALL CELLS (BASED ON NX x NY x NZ = 300)
--
THCWATER
         300*2O.0                                                              /
\end{verbatim}

The above example defines the water phase thermal conductivity of 20.0
for each cell in the 300 grid block model, as defined by the DIMENS
keyword in the RUNSPEC section.

\subsection{THPRESFT -- Define Fault Threshold
Pressures}\label{thpresft-define-fault-threshold-pressures}

\begin{longtable}[]{@{}
  >{\raggedright\arraybackslash}p{(\linewidth - 14\tabcolsep) * \real{0.1250}}
  >{\raggedright\arraybackslash}p{(\linewidth - 14\tabcolsep) * \real{0.1250}}
  >{\raggedright\arraybackslash}p{(\linewidth - 14\tabcolsep) * \real{0.1250}}
  >{\raggedright\arraybackslash}p{(\linewidth - 14\tabcolsep) * \real{0.1250}}
  >{\raggedright\arraybackslash}p{(\linewidth - 14\tabcolsep) * \real{0.1250}}
  >{\raggedright\arraybackslash}p{(\linewidth - 14\tabcolsep) * \real{0.1250}}
  >{\raggedright\arraybackslash}p{(\linewidth - 14\tabcolsep) * \real{0.1250}}
  >{\raggedright\arraybackslash}p{(\linewidth - 14\tabcolsep) * \real{0.1250}}@{}}
\toprule\noalign{}
\begin{minipage}[b]{\linewidth}\raggedright
RUNSPEC
\end{minipage} & \begin{minipage}[b]{\linewidth}\raggedright
GRID
\end{minipage} & \begin{minipage}[b]{\linewidth}\raggedright
EDIT
\end{minipage} & \begin{minipage}[b]{\linewidth}\raggedright
PROPS
\end{minipage} & \begin{minipage}[b]{\linewidth}\raggedright
REGIONS
\end{minipage} & \begin{minipage}[b]{\linewidth}\raggedright
SOLUTION
\end{minipage} & \begin{minipage}[b]{\linewidth}\raggedright
SUMMARY
\end{minipage} & \begin{minipage}[b]{\linewidth}\raggedright
SCHEDULE
\end{minipage} \\
\midrule\noalign{}
\endhead
\bottomrule\noalign{}
\endlastfoot
\end{longtable}

\subsubsection{Description}\label{description-362}

The THPRESFT keyword defines a fault threshold pressures that prevents
fluid flow from occurring across the fault plane until the threshold
pressure is exceeded, for when the threshold pressure option has been
activated via the THRPRES variable on the EQLOPTS keyword in the RUNSPEC
section.

Each row entry in the THPRESFT keyword defines a fault threshold
pressure.

\begin{longtable}[]{@{}
  >{\raggedright\arraybackslash}p{(\linewidth - 6\tabcolsep) * \real{0.2500}}
  >{\raggedright\arraybackslash}p{(\linewidth - 6\tabcolsep) * \real{0.2500}}
  >{\raggedright\arraybackslash}p{(\linewidth - 6\tabcolsep) * \real{0.2500}}
  >{\raggedright\arraybackslash}p{(\linewidth - 6\tabcolsep) * \real{0.2500}}@{}}
\toprule\noalign{}
\begin{minipage}[b]{\linewidth}\raggedright
No.
\end{minipage} & \begin{minipage}[b]{\linewidth}\raggedright
Name
\end{minipage} & \begin{minipage}[b]{\linewidth}\raggedright
Description
\end{minipage} & \begin{minipage}[b]{\linewidth}\raggedright
Default
\end{minipage} \\
\midrule\noalign{}
\endhead
\bottomrule\noalign{}
\endlastfoot
Field & Metric & Laboratory & \\
1 & FLTNAME & FLTNAME is a character string enclosed in quotes with a
maximum length of eight characters, that defines the name of the fault.
FLTNAME must have been previously defined using the FAULTS keyword in
the GROD section, otherwise an error will occur. & None \\
2 & PRESS & PRESS is a single positive real value that defines the
threshold pressure for the fault (FLTNAME). If PRESS is defaulted then
the simulator will set the threshold pressure to zero, that is the fault
is open to flow along the fault plane. & 0 \\
psia & barsa & atma & \\
Notes: & & & \\
\end{longtable}

\emph{Table 6.128: THPRESFT Keyword Description}

See also the FAULTS keyword in the GRID section that is used to created
named faults. Normally this fault data is automatically generated by
pre-processing software that generates the static model.

\subsubsection{Example}\label{example-244}

The example below defines two fault traces, `M\_WEST' and `BC' fault
having threshold pressures of 1000.0 and 2000 psis respectively.

\begin{verbatim}
--
--       DEFINE FAULTS IN THE GRID GEOMETRY
--
--       FAULT       ------------ FAULT TRACE -------------
--       NAME         I1   I2    J1   J2    K1   K2    FACE
FAULTS
         'M_WEST'      5    5     3    3    1    22   'X'   /
         'M_WEST'      5    5     4    4    1    22   'X'   /
         'M_WEST'      5    5     5    5    1    22   'X'   /
……………………………………………..

         'BC'         43   43     8    8    1    22   'Y'   /
         'BC'         42   42     9    9    1    22   'X'   /
         'BC'         44   44     8    8    1    22   'Y'   /
……………………………………………..
        /
--
--       DEFINE FAULT THRESHOLD PRESSURES
--
--       FAULT        THRESHOLD
--       NAME         PRESSURE
THPRESFT
         'M_WEST'     1000.0                                /
         'BC'         1200.0                                /
        /
\end{verbatim}

\subsection{TOPS -- Define the Depth at the Center of the Top Face for
Each
Cell}\label{tops-define-the-depth-at-the-center-of-the-top-face-for-each-cell}

\begin{longtable}[]{@{}
  >{\raggedright\arraybackslash}p{(\linewidth - 14\tabcolsep) * \real{0.1250}}
  >{\raggedright\arraybackslash}p{(\linewidth - 14\tabcolsep) * \real{0.1250}}
  >{\raggedright\arraybackslash}p{(\linewidth - 14\tabcolsep) * \real{0.1250}}
  >{\raggedright\arraybackslash}p{(\linewidth - 14\tabcolsep) * \real{0.1250}}
  >{\raggedright\arraybackslash}p{(\linewidth - 14\tabcolsep) * \real{0.1250}}
  >{\raggedright\arraybackslash}p{(\linewidth - 14\tabcolsep) * \real{0.1250}}
  >{\raggedright\arraybackslash}p{(\linewidth - 14\tabcolsep) * \real{0.1250}}
  >{\raggedright\arraybackslash}p{(\linewidth - 14\tabcolsep) * \real{0.1250}}@{}}
\toprule\noalign{}
\begin{minipage}[b]{\linewidth}\raggedright
RUNSPEC
\end{minipage} & \begin{minipage}[b]{\linewidth}\raggedright
GRID
\end{minipage} & \begin{minipage}[b]{\linewidth}\raggedright
EDIT
\end{minipage} & \begin{minipage}[b]{\linewidth}\raggedright
PROPS
\end{minipage} & \begin{minipage}[b]{\linewidth}\raggedright
REGIONS
\end{minipage} & \begin{minipage}[b]{\linewidth}\raggedright
SOLUTION
\end{minipage} & \begin{minipage}[b]{\linewidth}\raggedright
SUMMARY
\end{minipage} & \begin{minipage}[b]{\linewidth}\raggedright
SCHEDULE
\end{minipage} \\
\midrule\noalign{}
\endhead
\bottomrule\noalign{}
\endlastfoot
\end{longtable}

\subsubsection{Description}\label{description-363}

TOPS defines the depth of the top face of each cell in the model.

It can only be used with the Cartesian Regular Grid or Radial Grid
models.

\begin{longtable}[]{@{}
  >{\raggedright\arraybackslash}p{(\linewidth - 6\tabcolsep) * \real{0.2500}}
  >{\raggedright\arraybackslash}p{(\linewidth - 6\tabcolsep) * \real{0.2500}}
  >{\raggedright\arraybackslash}p{(\linewidth - 6\tabcolsep) * \real{0.2500}}
  >{\raggedright\arraybackslash}p{(\linewidth - 6\tabcolsep) * \real{0.2500}}@{}}
\toprule\noalign{}
\begin{minipage}[b]{\linewidth}\raggedright
No.
\end{minipage} & \begin{minipage}[b]{\linewidth}\raggedright
Name
\end{minipage} & \begin{minipage}[b]{\linewidth}\raggedright
Description
\end{minipage} & \begin{minipage}[b]{\linewidth}\raggedright
Default
\end{minipage} \\
\midrule\noalign{}
\endhead
\bottomrule\noalign{}
\endlastfoot
Field & Metric & Laboratory & \\
1 & TOPS & TOPS is an array of real numbers defining the depth at the
top face of each cell in the model. One can either just enter the TOPS
for the first layer only based on NX x NY entries and OPM Flow will
calculate the remaining TOPS based on either DZ or DZV. Alternatively NX
x NY x NZ TOPS may be entered for each cell in the model. See the DIMENS
keyword in the RUNSPEC section for the definition of NX, NY and NZ.
Repeat counts may be used, for example 10*5201.0. & None \\
feet & m & cm & \\
Notes: & & & \\
\end{longtable}

\emph{Table 6.129: TOPS Keyword Description}

See also the DEPTHS keyword to define the structural depth for the
cells.

\subsubsection{Examples}\label{examples-18}

The example below defines the TOPS of the cells for each cell for NX =
5, NY = 5 and NZ = 3 model, as well as the X and Y direction cells
sizes.

\begin{verbatim}
--
--       DEFINE GRID BLOCK TOPS FOR THE TOP LAYER (NX=5, NY=5, and NZ=3)
--
TOPS
         25*3100  25*3105  25*3110                                             /                                                                                 --
--       DEFINE GRID BLOCK X DIRECTION CELL SIZE (BASED ON NX = 5)
--
DXV
         5*100                                                                 /                                                                                 --
--       DEFINE GRID BLOCK X DIRECTION CELL SIZE (BASED ON NY = 5)
--
DYV
         5*100                                                                 /
\end{verbatim}

A second example is shown on the following page.

This example defines the same grid as before but with the TOPS keyword
only defining the top layer and DZV keyword defining the cells
thickness.

\begin{verbatim}
--
--       DEFINE GRID BLOCK TOPS FOR THE TOP LAYER (NX = 5, NY = 5, NZ = 3)
--
TOPS
         25*3100                                                               /
--
--       DEFINE GRID BLOCK Z DIRECTION CELL SIZE (BASED ON NZ = 3)
--
DZV
         3*5.0                                                                 /
--
--       DEFINE GRID BLOCK X DIRECTION CELL SIZE (BASED ON NX = 5)
--
DXV
         5*100                                                                 /                                                                                 --
--       DEFINE GRID BLOCK Y DIRECTION CELL SIZE (BASED ON NY = 5)
--
DYV
         5*100                                                                 /
\end{verbatim}

\subsection{TRANGL -- Define Non-Neighbor Connections Between Global and
LGR Cells
Manually}\label{trangl-define-non-neighbor-connections-between-global-and-lgr-cells-manually}

\begin{longtable}[]{@{}
  >{\raggedright\arraybackslash}p{(\linewidth - 14\tabcolsep) * \real{0.1250}}
  >{\raggedright\arraybackslash}p{(\linewidth - 14\tabcolsep) * \real{0.1250}}
  >{\raggedright\arraybackslash}p{(\linewidth - 14\tabcolsep) * \real{0.1250}}
  >{\raggedright\arraybackslash}p{(\linewidth - 14\tabcolsep) * \real{0.1250}}
  >{\raggedright\arraybackslash}p{(\linewidth - 14\tabcolsep) * \real{0.1250}}
  >{\raggedright\arraybackslash}p{(\linewidth - 14\tabcolsep) * \real{0.1250}}
  >{\raggedright\arraybackslash}p{(\linewidth - 14\tabcolsep) * \real{0.1250}}
  >{\raggedright\arraybackslash}p{(\linewidth - 14\tabcolsep) * \real{0.1250}}@{}}
\toprule\noalign{}
\begin{minipage}[b]{\linewidth}\raggedright
RUNSPEC
\end{minipage} & \begin{minipage}[b]{\linewidth}\raggedright
GRID
\end{minipage} & \begin{minipage}[b]{\linewidth}\raggedright
EDIT
\end{minipage} & \begin{minipage}[b]{\linewidth}\raggedright
PROPS
\end{minipage} & \begin{minipage}[b]{\linewidth}\raggedright
REGIONS
\end{minipage} & \begin{minipage}[b]{\linewidth}\raggedright
SOLUTION
\end{minipage} & \begin{minipage}[b]{\linewidth}\raggedright
SUMMARY
\end{minipage} & \begin{minipage}[b]{\linewidth}\raggedright
SCHEDULE
\end{minipage} \\
\midrule\noalign{}
\endhead
\bottomrule\noalign{}
\endlastfoot
\end{longtable}

\subsubsection{Description}\label{description-364}

TRANGL enables Non-Neighbor Connections (``NNC'') between the global
cells and the Local Grid Refinement (``LGR'') cells to be manually
specified, as oppose to the simulator calculating the
transmissibilities. The LGR keyword in the RUNSPEC section should be
utilized to define the presence of LGRs in the model and to define
various LGR dimension parameters.

This keyword is not supported by OPM Flow but would change the results
if supported so the simulation will be stopped.

\begin{longtable}[]{@{}
  >{\raggedright\arraybackslash}p{(\linewidth - 6\tabcolsep) * \real{0.2500}}
  >{\raggedright\arraybackslash}p{(\linewidth - 6\tabcolsep) * \real{0.2500}}
  >{\raggedright\arraybackslash}p{(\linewidth - 6\tabcolsep) * \real{0.2500}}
  >{\raggedright\arraybackslash}p{(\linewidth - 6\tabcolsep) * \real{0.2500}}@{}}
\toprule\noalign{}
\begin{minipage}[b]{\linewidth}\raggedright
No.
\end{minipage} & \begin{minipage}[b]{\linewidth}\raggedright
Name
\end{minipage} & \begin{minipage}[b]{\linewidth}\raggedright
Description
\end{minipage} & \begin{minipage}[b]{\linewidth}\raggedright
Default
\end{minipage} \\
\midrule\noalign{}
\endhead
\bottomrule\noalign{}
\endlastfoot
Field & Metric & Laboratory & \\
1 & I1 & A positive integer that defines the LGR grid block in the
I-direction in a non-neighbor connection, must be greater than or equal
to one and less than or equal to NX on the CARFIN keyword in the GRID
section. & None \\
2 & J1 & A positive integer that defines the LGR grid block in the
J-direction in a non-neighbor connection, must be greater than or equal
to one and less than or equal to NY on the CARFIN keyword in the GRID
section. & None \\
3 & K1 & A positive integer that defines the LGR grid block in the
K-direction in a non-neighbor connection, must be greater than or equal
to one and less than or equal to NZ on the CARFIN keyword in the GRID
section. & None \\
4 & I2 & A positive integer that defines the GLOBAL grid block in the
I-direction in a non-neighbor connection, must be greater than or equal
to one and less than or equal to NX on the DIMENS keyword in the RUNSPEC
section. & None \\
5 & J2 & A positive integer that defines the GLOBAL grid block in the
J-direction in a non-neighbor connection, must be greater than or equal
to one and less than or equal to NY on the DIMENS keyword in the RUNSPEC
section. & None \\
6 & K2 & A positive integer that defines the GLOBAL grid block in the
K-direction in a non-neighbor connection, must be greater than or equal
to one and less than or equal to NZ on the DIMENS keyword in the RUNSPEC
section. & None \\
7 & TRANSNNC & TRANSNNC is a positive real number greater than or equal
to zero that defines the transmissibility between the GLOBAL grid block
(I1, J1, K1) and the LGR grid block (I2, J2, K2). The default value of
zero sets the transmissibility between the two cells to zero. & 0.0 \\
cP.rb/day/psia & cP.rm3/day/bars & cP.rcc/hr/atm & \\
Notes: & & & \\
\end{longtable}

\emph{Table 6.130: TRANGL Keyword Description}

\subsubsection{Example}\label{example-245}

\begin{verbatim}
--
--       MANUALLY DEFINE LGR-GLOBAL GRID NON-NEIGHBOR CONNECTIONS
--
--       ----LGR-----   ---GLOBAL----   -- TRANSNCC --
--       I1   J1   K1    I2   J2   K2
TRANGL
         1    1    1     1    1    2        0.2500     /
         1    1    2     1    1    3        0.2500     /
         1    1    3     1    1    4        0.2500     /
/
\end{verbatim}

The above example defines the transmissibility between LGR cell (1, 1,
1) and global cell (1, 1, 2), LGR cell (1, 1, 2) and global cell (1, 1,
3) and finally between LGR cell (1, 1, 3) and global cell (1, 1, 4) to
be 0.2500.

\subsection{USEFLUX -- Activate Flux Boundary Model and Define Flux
File}\label{useflux-activate-flux-boundary-model-and-define-flux-file}

\begin{longtable}[]{@{}
  >{\raggedright\arraybackslash}p{(\linewidth - 14\tabcolsep) * \real{0.1250}}
  >{\raggedright\arraybackslash}p{(\linewidth - 14\tabcolsep) * \real{0.1250}}
  >{\raggedright\arraybackslash}p{(\linewidth - 14\tabcolsep) * \real{0.1250}}
  >{\raggedright\arraybackslash}p{(\linewidth - 14\tabcolsep) * \real{0.1250}}
  >{\raggedright\arraybackslash}p{(\linewidth - 14\tabcolsep) * \real{0.1250}}
  >{\raggedright\arraybackslash}p{(\linewidth - 14\tabcolsep) * \real{0.1250}}
  >{\raggedright\arraybackslash}p{(\linewidth - 14\tabcolsep) * \real{0.1250}}
  >{\raggedright\arraybackslash}p{(\linewidth - 14\tabcolsep) * \real{0.1250}}@{}}
\toprule\noalign{}
\begin{minipage}[b]{\linewidth}\raggedright
RUNSPEC
\end{minipage} & \begin{minipage}[b]{\linewidth}\raggedright
GRID
\end{minipage} & \begin{minipage}[b]{\linewidth}\raggedright
EDIT
\end{minipage} & \begin{minipage}[b]{\linewidth}\raggedright
PROPS
\end{minipage} & \begin{minipage}[b]{\linewidth}\raggedright
REGIONS
\end{minipage} & \begin{minipage}[b]{\linewidth}\raggedright
SOLUTION
\end{minipage} & \begin{minipage}[b]{\linewidth}\raggedright
SUMMARY
\end{minipage} & \begin{minipage}[b]{\linewidth}\raggedright
SCHEDULE
\end{minipage} \\
\midrule\noalign{}
\endhead
\bottomrule\noalign{}
\endlastfoot
\end{longtable}

\subsubsection{Description}\label{description-365}

The USEFLUX keyword activates the Flux Boundary model and defines the
name of the FLUX file. Only grid blocks that have been declared by the
FLUXREG keyword in the GRID section to be in an active flux region, are
active for the run.

This keyword is not supported by OPM Flow but would change the results
if supported so the simulation will be stopped.

\subsection{USENOFLO -- Activate Flux Boundary Model Without a Flux
File}\label{usenoflo-activate-flux-boundary-model-without-a-flux-file}

\begin{longtable}[]{@{}
  >{\raggedright\arraybackslash}p{(\linewidth - 14\tabcolsep) * \real{0.1250}}
  >{\raggedright\arraybackslash}p{(\linewidth - 14\tabcolsep) * \real{0.1250}}
  >{\raggedright\arraybackslash}p{(\linewidth - 14\tabcolsep) * \real{0.1250}}
  >{\raggedright\arraybackslash}p{(\linewidth - 14\tabcolsep) * \real{0.1250}}
  >{\raggedright\arraybackslash}p{(\linewidth - 14\tabcolsep) * \real{0.1250}}
  >{\raggedright\arraybackslash}p{(\linewidth - 14\tabcolsep) * \real{0.1250}}
  >{\raggedright\arraybackslash}p{(\linewidth - 14\tabcolsep) * \real{0.1250}}
  >{\raggedright\arraybackslash}p{(\linewidth - 14\tabcolsep) * \real{0.1250}}@{}}
\toprule\noalign{}
\begin{minipage}[b]{\linewidth}\raggedright
RUNSPEC
\end{minipage} & \begin{minipage}[b]{\linewidth}\raggedright
GRID
\end{minipage} & \begin{minipage}[b]{\linewidth}\raggedright
EDIT
\end{minipage} & \begin{minipage}[b]{\linewidth}\raggedright
PROPS
\end{minipage} & \begin{minipage}[b]{\linewidth}\raggedright
REGIONS
\end{minipage} & \begin{minipage}[b]{\linewidth}\raggedright
SOLUTION
\end{minipage} & \begin{minipage}[b]{\linewidth}\raggedright
SUMMARY
\end{minipage} & \begin{minipage}[b]{\linewidth}\raggedright
SCHEDULE
\end{minipage} \\
\midrule\noalign{}
\endhead
\bottomrule\noalign{}
\endlastfoot
\end{longtable}

\subsubsection{Description}\label{description-366}

The USENOFLUX keyword activates the Flux Boundary model without a FLUX
file. The USEFLUX keyword should still be in the input deck, but in this
case the FLUX filename is ignored. The option is useful when the no-flow
boundary condition is a reasonable assumption and avoids the pre-cursor
run used to generate the FLUX file via the DUMPFLUX keyword in the GRID
section. Only grid blocks that have been declared by the FLUXREG keyword
in the GRID section to be in an active flux region, are active for the
run.

There is no data required for this keyword and there is no terminating
``/'' for this keyword.

This keyword is not supported by OPM Flow but would change the results
if supported so the simulation will be stopped.

This keyword is not supported by OPM Flow but would change the results
if supported so the simulation will be stopped.

\subsubsection{Example}\label{example-246}

\begin{verbatim}
--
--       ACTIVATE FLUX BOUNDARY MODEL WITHOUT A FLUX FILE
--
USEFLUX
/

USENOFLO
\end{verbatim}

The above example activates the Flux Boundary model without a FLUX file.

\subsection{VEDEBUG -- Vertical Equilibrium Debug Data
Output}\label{vedebug-vertical-equilibrium-debug-data-output}

\begin{longtable}[]{@{}
  >{\raggedright\arraybackslash}p{(\linewidth - 14\tabcolsep) * \real{0.1250}}
  >{\raggedright\arraybackslash}p{(\linewidth - 14\tabcolsep) * \real{0.1250}}
  >{\raggedright\arraybackslash}p{(\linewidth - 14\tabcolsep) * \real{0.1250}}
  >{\raggedright\arraybackslash}p{(\linewidth - 14\tabcolsep) * \real{0.1250}}
  >{\raggedright\arraybackslash}p{(\linewidth - 14\tabcolsep) * \real{0.1250}}
  >{\raggedright\arraybackslash}p{(\linewidth - 14\tabcolsep) * \real{0.1250}}
  >{\raggedright\arraybackslash}p{(\linewidth - 14\tabcolsep) * \real{0.1250}}
  >{\raggedright\arraybackslash}p{(\linewidth - 14\tabcolsep) * \real{0.1250}}@{}}
\toprule\noalign{}
\begin{minipage}[b]{\linewidth}\raggedright
RUNSPEC
\end{minipage} & \begin{minipage}[b]{\linewidth}\raggedright
GRID
\end{minipage} & \begin{minipage}[b]{\linewidth}\raggedright
EDIT
\end{minipage} & \begin{minipage}[b]{\linewidth}\raggedright
PROPS
\end{minipage} & \begin{minipage}[b]{\linewidth}\raggedright
REGIONS
\end{minipage} & \begin{minipage}[b]{\linewidth}\raggedright
SOLUTION
\end{minipage} & \begin{minipage}[b]{\linewidth}\raggedright
SUMMARY
\end{minipage} & \begin{minipage}[b]{\linewidth}\raggedright
SCHEDULE
\end{minipage} \\
\midrule\noalign{}
\endhead
\bottomrule\noalign{}
\endlastfoot
\end{longtable}

\subsubsection{Description}\label{description-367}

This keyword defines the debug Vertical Equilibrium (``VE'') data to be
written to the debug file (*.DBG), for when the VE model has been
activated by the VE keyword in the RUNSPEC section.

This keyword is not supported by OPM Flow but would change the results
if supported so the simulation will be stopped.

\subsection{VEFIN -- Activate Vertical Equilibrium Model
(LGR)}\label{vefin-activate-vertical-equilibrium-model-lgr}

\begin{longtable}[]{@{}
  >{\raggedright\arraybackslash}p{(\linewidth - 14\tabcolsep) * \real{0.1250}}
  >{\raggedright\arraybackslash}p{(\linewidth - 14\tabcolsep) * \real{0.1250}}
  >{\raggedright\arraybackslash}p{(\linewidth - 14\tabcolsep) * \real{0.1250}}
  >{\raggedright\arraybackslash}p{(\linewidth - 14\tabcolsep) * \real{0.1250}}
  >{\raggedright\arraybackslash}p{(\linewidth - 14\tabcolsep) * \real{0.1250}}
  >{\raggedright\arraybackslash}p{(\linewidth - 14\tabcolsep) * \real{0.1250}}
  >{\raggedright\arraybackslash}p{(\linewidth - 14\tabcolsep) * \real{0.1250}}
  >{\raggedright\arraybackslash}p{(\linewidth - 14\tabcolsep) * \real{0.1250}}@{}}
\toprule\noalign{}
\begin{minipage}[b]{\linewidth}\raggedright
RUNSPEC
\end{minipage} & \begin{minipage}[b]{\linewidth}\raggedright
GRID
\end{minipage} & \begin{minipage}[b]{\linewidth}\raggedright
EDIT
\end{minipage} & \begin{minipage}[b]{\linewidth}\raggedright
PROPS
\end{minipage} & \begin{minipage}[b]{\linewidth}\raggedright
REGIONS
\end{minipage} & \begin{minipage}[b]{\linewidth}\raggedright
SOLUTION
\end{minipage} & \begin{minipage}[b]{\linewidth}\raggedright
SUMMARY
\end{minipage} & \begin{minipage}[b]{\linewidth}\raggedright
SCHEDULE
\end{minipage} \\
\midrule\noalign{}
\endhead
\bottomrule\noalign{}
\endlastfoot
\end{longtable}

\subsubsection{Description}\label{description-368}

If the VE keyword in the RUNSPEC section has been used to activate the
Vertical Equilibrium (``VE'') model for the global grid, then the VEFIN
keyword may used to set various options for the Local Grid Refinements
(``LGR''). The LGR keyword in the RUNSPEC section should be activated to
indicate the presence of LGRs and the keyword VEFIN should be placed in
between the CARFIN and ENDFIN keywords in the GRID section.

This keyword is not supported by OPM Flow but would change the results
if supported so the simulation will be stopped.

\subsection{WARN -- Activate Warning
Messages}\label{warn-activate-warning-messages-2}

Turns on warning messages to be printed to the print file (*.PRT); note
that this keyword is activated by default and can subsequently be
switched off by the NOWARN activation keyword. The warning messages may
be turned on and off using keywords WARN and NOWARN. OPM Flow always
prints error messages.

See WARN -- Activate Warning Messages in the GLOBAL section for a full
description.

\subsection{ZCORN -- Define the Depth of Each Corner-Point of a Grid
Block}\label{zcorn-define-the-depth-of-each-corner-point-of-a-grid-block}

\begin{longtable}[]{@{}
  >{\raggedright\arraybackslash}p{(\linewidth - 14\tabcolsep) * \real{0.1250}}
  >{\raggedright\arraybackslash}p{(\linewidth - 14\tabcolsep) * \real{0.1250}}
  >{\raggedright\arraybackslash}p{(\linewidth - 14\tabcolsep) * \real{0.1250}}
  >{\raggedright\arraybackslash}p{(\linewidth - 14\tabcolsep) * \real{0.1250}}
  >{\raggedright\arraybackslash}p{(\linewidth - 14\tabcolsep) * \real{0.1250}}
  >{\raggedright\arraybackslash}p{(\linewidth - 14\tabcolsep) * \real{0.1250}}
  >{\raggedright\arraybackslash}p{(\linewidth - 14\tabcolsep) * \real{0.1250}}
  >{\raggedright\arraybackslash}p{(\linewidth - 14\tabcolsep) * \real{0.1250}}@{}}
\toprule\noalign{}
\begin{minipage}[b]{\linewidth}\raggedright
RUNSPEC
\end{minipage} & \begin{minipage}[b]{\linewidth}\raggedright
GRID
\end{minipage} & \begin{minipage}[b]{\linewidth}\raggedright
EDIT
\end{minipage} & \begin{minipage}[b]{\linewidth}\raggedright
PROPS
\end{minipage} & \begin{minipage}[b]{\linewidth}\raggedright
REGIONS
\end{minipage} & \begin{minipage}[b]{\linewidth}\raggedright
SOLUTION
\end{minipage} & \begin{minipage}[b]{\linewidth}\raggedright
SUMMARY
\end{minipage} & \begin{minipage}[b]{\linewidth}\raggedright
SCHEDULE
\end{minipage} \\
\midrule\noalign{}
\endhead
\bottomrule\noalign{}
\endlastfoot
\end{longtable}

\subsubsection{Description}\label{description-369}

ZCORN defines the depth of each corner point of a grid block on the
pillars defining the reservoir grid. A total of 8 x NX x NY x NZ values
are needed to fully define all the depths in the model. The depths
specifying the top of the first layer are entered first with one point
for each pillar for each grid block. The points are entered with the X
axis cycling fastest. Next come the depths of the bottom of the first
layer. The top of layer two follows etc.

The keyword can be only used be used with Irregular Corner-Point Grids.

\begin{longtable}[]{@{}
  >{\raggedright\arraybackslash}p{(\linewidth - 6\tabcolsep) * \real{0.2500}}
  >{\raggedright\arraybackslash}p{(\linewidth - 6\tabcolsep) * \real{0.2500}}
  >{\raggedright\arraybackslash}p{(\linewidth - 6\tabcolsep) * \real{0.2500}}
  >{\raggedright\arraybackslash}p{(\linewidth - 6\tabcolsep) * \real{0.2500}}@{}}
\toprule\noalign{}
\begin{minipage}[b]{\linewidth}\raggedright
No.
\end{minipage} & \begin{minipage}[b]{\linewidth}\raggedright
Name
\end{minipage} & \begin{minipage}[b]{\linewidth}\raggedright
Description
\end{minipage} & \begin{minipage}[b]{\linewidth}\raggedright
Default
\end{minipage} \\
\midrule\noalign{}
\endhead
\bottomrule\noalign{}
\endlastfoot
Field & Metric & Laboratory & \\
1 & ZCORN & An array of depths with 8 depths for each cell, for a total
of 8 x Nx x NY x NZ entries & None \\
feet & metres & cm & \\
Notes: & & & \\
\end{longtable}

\emph{Table 6.131: ZCORN Keyword Description}

See also the SPECGRID, COORD and COORDSYS keywords to fully define an
Irregular Corner-Point Grid.

\subsubsection{Example}\label{example-247}

\begin{verbatim}
--
--       SPECIFY CORNER-POINT DEPTHS FOR A 3 x 2 x 2 GRID,
--       WITH CONSTANT SLOPE IN THE X AND Y DIRECTIONS
--       SUCH THAT ALL CORNER POINTS OF NEIGHBOURING BLOCKS ALIGN
ZCORN
--
--                    top of layer 1
--
         1450   1500   1500   1550   1550   1600
         1500   1550   1550   1600   1600   1650
         1500   1550   1550   1600   1600   1650
         1550   1600   1600   1650   1650   1700
--
--                     bottom of layer 1
--
         1460   1510   1510   1560   1560   1610
         1510   1560   1560   1610   1610   1660
         1510   1560   1560   1610   1610   1660
         1560   1610   1610   1660   1660   1710
--
--                     top of layer 2
--
         1460   1510   1510   1560   1560   1610
         1510   1560   1560   1610   1610   1660
         1510   1560   1560   1610   1610   1660
         1560   1610   1610   1660   1660   1710
--
--                     bottom of layer 2
--
         1470   1520   1520   1570   1570   1620
         1520   1570   1570   1620   1620   1670
         1520   1570   1570   1620   1620   1670
         1570   1620   1620   1670   1670   1720
/
\end{verbatim}

The above example defines depths of the vertical coordinate lines for a
regular 3 by 2 by 2 grid with a constant slope in the x and y directions
such that all the corner points of neighboring blocks are aligned.

\bookmarksetup{startatroot}

\chapter{EDIT SECTION}\label{edit-section}

\section{Introduction}\label{introduction-5}

This section allows users to define changes to be applied after OPM Flow
has processed the data in the GRID section, that is the resulting pore
volume (PORV) and transmissibility arrays (TRANX, TRANY, and TRANZ). The
entered primary static arrays (PORO, PERMX, etc.) in the GRID section
are no longer available and all modifications are applied to the pore
volume and transmissibility arrays.

Historically the intention of this section was allow for the editing of
the processed data; however, the features available in this section
have, through time, migrated to the GRID section. For example the array
operator keywords like ADD, COPY, MULTIPLY, etc. are available in the
GRID section and thus enable editing of the primary static arrays (PORO,
PERMX, etc.). This increased capability in the GRID section has
therefore made the EDIT section somewhat redundant.

Nevertheless the section is sometimes used by users to incorporate
history matching parameter changes and by static earth modeling software
packages to import the pore volumes and transmissibilities calculated in
the static model directly into the numerical model via the EDIT section.
Although the latter work flow is not very common.

\section{Keyword Definitions}\label{keyword-definitions-7}

\section{Data Requirements}\label{data-requirements-3}

As the primary purpose of this section is to modify the simulator's
calculated pore volumes and transmissibilities, then the properties used
to define these arrays must have been fully defined in the GRID section.
The arrays available for modification in the EDIT section are listed in
Table 7.1 together with the associated GRID arrays used to generate the
EDIT property array.

\begin{longtable}[]{@{}
  >{\raggedright\arraybackslash}p{(\linewidth - 6\tabcolsep) * \real{0.2500}}
  >{\raggedright\arraybackslash}p{(\linewidth - 6\tabcolsep) * \real{0.2500}}
  >{\raggedright\arraybackslash}p{(\linewidth - 6\tabcolsep) * \real{0.2500}}
  >{\raggedright\arraybackslash}p{(\linewidth - 6\tabcolsep) * \real{0.2500}}@{}}
\toprule\noalign{}
\begin{minipage}[b]{\linewidth}\raggedright
Cartesian and Irregular Corner-Point Grids Keywords
\end{minipage} & \begin{minipage}[b]{\linewidth}\raggedright
Radial Grid Keywords
\end{minipage} & \begin{minipage}[b]{\linewidth}\raggedright
\end{minipage} & \begin{minipage}[b]{\linewidth}\raggedright
\end{minipage} \\
\midrule\noalign{}
\endhead
\bottomrule\noalign{}
\endlastfoot
GRID & EDIT & GRID & EDIT \\
TOPS & DEPTH & TOPS & DEPTH \\
DX & PORV & DR & PORV \\
DY & DTHETA & & \\
DZ & DZ & & \\
DZNET & DZNET & & \\
PORO & PORO & & \\
NTG & NTG & & \\
PERMX & TRANX & PERMR & TRANR \\
MULTX & MULTR & & \\
PERMY & TRANY & PERMTHT & TRANTHT \\
MULTY & MULTTHT & & \\
PERMZ & TRANZ & PERMZ & TRANZ \\
MULTZ & MULTZ & & \\
Notes: & & & \\
\end{longtable}

\emph{Table 7.1: EDIT Section Arrays Available for Modification}

An example pore volume array (PORV property) from the Volve {[}The Volve
Data was approved for data sharing in 2018 by the initiative of the last
Operating company, Equinor and approved by the license partners
ExxonMobil E\&P Norway AS and Bayerngas Norge AS in the end of 2017.{]}
field is shown in Figure 7.1 and Figure 7.2 illustrates the model's
transmissibility in the x-direction (TRANX).

\pandocbounded{\includegraphics[keepaspectratio]{chapters/images/Frame10_0f44802fffa6.png}}
\pandocbounded{\includegraphics[keepaspectratio]{chapters/images/Frame11_0f44802fffa6.png}}
\pandocbounded{\includegraphics[keepaspectratio]{chapters/images/Image9_0f44802fffa6.png}}

\pandocbounded{\includegraphics[keepaspectratio]{chapters/images/Frame12_d4fccc8fc349.png}}
\pandocbounded{\includegraphics[keepaspectratio]{chapters/images/Image10_d4fccc8fc349.png}}

\section{Keyword Definitions}\label{keyword-definitions-8}

\subsection{ADD -- Add a Constant to a Specified
Array}\label{add-add-a-constant-to-a-specified-array-1}

The ADD keyword adds a constant to a specified array or part of an
array. The constant can be real or integer depending on the array type;
however, the arrays that can be operated on is dependent on which
section the ADD keyword is being applied.

See ADD -- Add a Constant to a Specified Array in the GRID section for a
full description.

\subsection{ADDREG -- Add a Constant to an Array based on a Region
Number}\label{addreg-add-a-constant-to-an-array-based-on-a-region-number-1}

The ADDREG keyword adds a constant to a specified array or part of an
array based on cells with a specific region number. The region number
array can be FLUXNUM, MULTNUM or OPERNUM and these arrays must be
defined and be available before the ADDREG keyword is read by the
simulator. The constant can be real or integer depending on the property
array type; however, the property arrays that can be operated on is
dependent on which section the ADDREG keyword is being used.

See ADDREG -- Add a Constant to an Array based on a Region Number in the
GRID section for a full description.

\subsection{BOUNDARY -- Define a Boundary Box for
Printing}\label{boundary-define-a-boundary-box-for-printing-1}

The BOUNDARY keyword defines a rectangular grid for printing various
arrays to the output print file (*.PRT); thus, avoiding printing all the
elements in the selected array.

See BOUNDARY -- Define a Boundary Box for Printing in the GRID section
for a full description.

\subsection{BOX -- Define a Range of Grid Blocks to Enter Property
Data}\label{box-define-a-range-of-grid-blocks-to-enter-property-data-1}

BOX defines a range of grid blocks for which subsequent data is assigned
for all the cells in the defined BOX. Note that the BOX grid is reset by
the keyword ENDBOX by resetting the current defined BOX to be the whole
grid. The keyword can be used for any array and for all grid types.

See BOX - Define a Range of Grid Blocks to Enter Property Data in the
GRID section for a full description.

\subsection{COLUMNS -- Define Input File Column
Margins}\label{columns-define-input-file-column-margins-3}

The COLUMNS keyword defines the input file column margins; characters
outside the margins are ignored by the input parser.

See COLUMNS -- Define Input File Column Margins in the GLOBAL section
for a full description.

\subsection{COPY -- Copy Array Data to Another
Array}\label{copy-copy-array-data-to-another-array-1}

The COPY keyword copies an array (or part of an array) to another array
or part of an array. The arrays can be real or integer depending on the
array type; however, the arrays that can be operated on is dependent on
which section the COPY keyword is being used.

See COPY -- Copy Array Data to Another Array in the GRID section for a
full description.

\subsection{COPYREG -- Copy an Array to Another Array based on a Region
Number}\label{copyreg-copy-an-array-to-another-array-based-on-a-region-number-1}

The COPYREG keyword copies a specified array or part of an array based
on cells with a specific region number to another array. The region
number array can be FLUXNUM, MULTNUM or OPERNUM and these arrays must be
defined and be available before the COPYREG keyword is read by the
simulator. The property arrays can be real or integer depending on the
property array type; however, the property arrays that can be operated
on is dependent on which section the COPYREG keyword is being used.

See COPYREG -- Copy an Array to Another Array based on a Region Number
in the GRID section for a full description.

\subsection{DEBUG -- Define the Debug Data to be Printed to
File}\label{debug-define-the-debug-data-to-be-printed-to-file-3}

This keyword defines the debug data to be written to the debug file
(*.DBG). This keyword is not supported by OPM Flow but has no effect on
the results so it will be ignored.

See DEBUG -- Define the Debug Data to be Printed to File in the GLOBAL
section for a full description.

\subsection{DEPTH -- Edits the Depth at the Center of Each
Cell}\label{depth-edits-the-depth-at-the-center-of-each-cell}

\begin{longtable}[]{@{}
  >{\raggedright\arraybackslash}p{(\linewidth - 14\tabcolsep) * \real{0.1250}}
  >{\raggedright\arraybackslash}p{(\linewidth - 14\tabcolsep) * \real{0.1250}}
  >{\raggedright\arraybackslash}p{(\linewidth - 14\tabcolsep) * \real{0.1250}}
  >{\raggedright\arraybackslash}p{(\linewidth - 14\tabcolsep) * \real{0.1250}}
  >{\raggedright\arraybackslash}p{(\linewidth - 14\tabcolsep) * \real{0.1250}}
  >{\raggedright\arraybackslash}p{(\linewidth - 14\tabcolsep) * \real{0.1250}}
  >{\raggedright\arraybackslash}p{(\linewidth - 14\tabcolsep) * \real{0.1250}}
  >{\raggedright\arraybackslash}p{(\linewidth - 14\tabcolsep) * \real{0.1250}}@{}}
\toprule\noalign{}
\begin{minipage}[b]{\linewidth}\raggedright
RUNSPEC
\end{minipage} & \begin{minipage}[b]{\linewidth}\raggedright
GRID
\end{minipage} & \begin{minipage}[b]{\linewidth}\raggedright
EDIT
\end{minipage} & \begin{minipage}[b]{\linewidth}\raggedright
PROPS
\end{minipage} & \begin{minipage}[b]{\linewidth}\raggedright
REGIONS
\end{minipage} & \begin{minipage}[b]{\linewidth}\raggedright
SOLUTION
\end{minipage} & \begin{minipage}[b]{\linewidth}\raggedright
SUMMARY
\end{minipage} & \begin{minipage}[b]{\linewidth}\raggedright
SCHEDULE
\end{minipage} \\
\midrule\noalign{}
\endhead
\bottomrule\noalign{}
\endlastfoot
\end{longtable}

\subsubsection{Description}\label{description-370}

The DEPTH keywords modifies the depth at the center of selected cells in
the model. The cells DEPTH are calculated by OPM Flow at the end of the
GRID section and this keyword allows the user to adjust the calculated
depths in the EDIT section. The area to be modified can be defined via
the various grid selection keywords, ADD, BOX, EQUALS, etc., and areas
that are not selected remain unchanged.

\begin{longtable}[]{@{}
  >{\raggedright\arraybackslash}p{(\linewidth - 6\tabcolsep) * \real{0.2500}}
  >{\raggedright\arraybackslash}p{(\linewidth - 6\tabcolsep) * \real{0.2500}}
  >{\raggedright\arraybackslash}p{(\linewidth - 6\tabcolsep) * \real{0.2500}}
  >{\raggedright\arraybackslash}p{(\linewidth - 6\tabcolsep) * \real{0.2500}}@{}}
\toprule\noalign{}
\begin{minipage}[b]{\linewidth}\raggedright
No.
\end{minipage} & \begin{minipage}[b]{\linewidth}\raggedright
Name
\end{minipage} & \begin{minipage}[b]{\linewidth}\raggedright
Description
\end{minipage} & \begin{minipage}[b]{\linewidth}\raggedright
Default
\end{minipage} \\
\midrule\noalign{}
\endhead
\bottomrule\noalign{}
\endlastfoot
Field & Metric & Laboratory & \\
1 & DEPTH & DEPTH is an array of real numbers defining the depth at the
center of each cell in the model. Only the values in the currently
defined input BOX needed be entered. Repeat counts may be used, for
example 30*5201.0. & None \\
feet & m & cm & \\
Notes: & & & \\
\end{longtable}

\emph{Table 7.2: Depth Keyword Description}

See also the TOPS keyword to define the top structural depth for the
cells.

\subsubsection{Examples}\label{examples-19}

The example below modifies the DEPTH of the cells for a selection of 10
cells from an NX = 10, NY = 11 and NZ = 20 model.

\begin{verbatim}
--
--       DEFINE INPUT BOX FOR EDITING INPUT ARRAYS
--
--       ---------- BOX ---------
--       I1  I2   J1  J2   K1  K2
BOX
          1  10   11  11   20  20 / SET BOX AREA TO BE MODIFIED
/
--
--       SET GRID BLOCK CENTER DEPTH FOR THE GRID BLOCKS
--
DEPTH    10*3500.0                                                              /
--
--       DEFINE END OF INPUT BOX EDITING OF INPUT ARRAYS
--
ENDBOX
\end{verbatim}

Alternatively the EQUALS keyword can be used to perform the same edit.

\begin{verbatim}
--
--       ARRAY       CONSTANT     ---------- BOX ---------
--                                I1  I2   J1  J2   K1  K2
EQUALS
         DEPTH       3500.0        1  10   11  11   20  20 / RESET DEPTH
/
\end{verbatim}

\subsection{DIFFR -- Define Grid Block Radial Direction Diffusivity
Values}\label{diffr-define-grid-block-radial-direction-diffusivity-values}

\begin{longtable}[]{@{}
  >{\raggedright\arraybackslash}p{(\linewidth - 14\tabcolsep) * \real{0.1250}}
  >{\raggedright\arraybackslash}p{(\linewidth - 14\tabcolsep) * \real{0.1250}}
  >{\raggedright\arraybackslash}p{(\linewidth - 14\tabcolsep) * \real{0.1250}}
  >{\raggedright\arraybackslash}p{(\linewidth - 14\tabcolsep) * \real{0.1250}}
  >{\raggedright\arraybackslash}p{(\linewidth - 14\tabcolsep) * \real{0.1250}}
  >{\raggedright\arraybackslash}p{(\linewidth - 14\tabcolsep) * \real{0.1250}}
  >{\raggedright\arraybackslash}p{(\linewidth - 14\tabcolsep) * \real{0.1250}}
  >{\raggedright\arraybackslash}p{(\linewidth - 14\tabcolsep) * \real{0.1250}}@{}}
\toprule\noalign{}
\begin{minipage}[b]{\linewidth}\raggedright
RUNSPEC
\end{minipage} & \begin{minipage}[b]{\linewidth}\raggedright
GRID
\end{minipage} & \begin{minipage}[b]{\linewidth}\raggedright
EDIT
\end{minipage} & \begin{minipage}[b]{\linewidth}\raggedright
PROPS
\end{minipage} & \begin{minipage}[b]{\linewidth}\raggedright
REGIONS
\end{minipage} & \begin{minipage}[b]{\linewidth}\raggedright
SOLUTION
\end{minipage} & \begin{minipage}[b]{\linewidth}\raggedright
SUMMARY
\end{minipage} & \begin{minipage}[b]{\linewidth}\raggedright
SCHEDULE
\end{minipage} \\
\midrule\noalign{}
\endhead
\bottomrule\noalign{}
\endlastfoot
\end{longtable}

\subsubsection{Description}\label{description-371}

The DIFFR keyword defines the radial direction diffusivity values for
cells in the current input box for when the Diffusivity option has been
activated by the DIFFUSE keyword in the RUNSPEC section.

This keyword is not supported by OPM Flow but would change the results
if supported so the simulation will be stopped.

\subsection{DIFFTHT -- Define Grid Block Theta Direction Diffusivity
Values}\label{difftht-define-grid-block-theta-direction-diffusivity-values}

\begin{longtable}[]{@{}
  >{\raggedright\arraybackslash}p{(\linewidth - 14\tabcolsep) * \real{0.1250}}
  >{\raggedright\arraybackslash}p{(\linewidth - 14\tabcolsep) * \real{0.1250}}
  >{\raggedright\arraybackslash}p{(\linewidth - 14\tabcolsep) * \real{0.1250}}
  >{\raggedright\arraybackslash}p{(\linewidth - 14\tabcolsep) * \real{0.1250}}
  >{\raggedright\arraybackslash}p{(\linewidth - 14\tabcolsep) * \real{0.1250}}
  >{\raggedright\arraybackslash}p{(\linewidth - 14\tabcolsep) * \real{0.1250}}
  >{\raggedright\arraybackslash}p{(\linewidth - 14\tabcolsep) * \real{0.1250}}
  >{\raggedright\arraybackslash}p{(\linewidth - 14\tabcolsep) * \real{0.1250}}@{}}
\toprule\noalign{}
\begin{minipage}[b]{\linewidth}\raggedright
RUNSPEC
\end{minipage} & \begin{minipage}[b]{\linewidth}\raggedright
GRID
\end{minipage} & \begin{minipage}[b]{\linewidth}\raggedright
EDIT
\end{minipage} & \begin{minipage}[b]{\linewidth}\raggedright
PROPS
\end{minipage} & \begin{minipage}[b]{\linewidth}\raggedright
REGIONS
\end{minipage} & \begin{minipage}[b]{\linewidth}\raggedright
SOLUTION
\end{minipage} & \begin{minipage}[b]{\linewidth}\raggedright
SUMMARY
\end{minipage} & \begin{minipage}[b]{\linewidth}\raggedright
SCHEDULE
\end{minipage} \\
\midrule\noalign{}
\endhead
\bottomrule\noalign{}
\endlastfoot
\end{longtable}

\subsubsection{Description}\label{description-372}

The DIFFTHT keyword defines the theta direction diffusivity values for
cells in the current input box for when the Diffusivity option has been
activated by the DIFFUSE keyword in the RUNSPEC section.

This keyword is not supported by OPM Flow but would change the results
if supported so the simulation will be stopped.

\subsection{DIFFX -- Define Grid Block X-Direction Diffusivity
Values}\label{diffx-define-grid-block-x-direction-diffusivity-values}

\begin{longtable}[]{@{}
  >{\raggedright\arraybackslash}p{(\linewidth - 14\tabcolsep) * \real{0.1250}}
  >{\raggedright\arraybackslash}p{(\linewidth - 14\tabcolsep) * \real{0.1250}}
  >{\raggedright\arraybackslash}p{(\linewidth - 14\tabcolsep) * \real{0.1250}}
  >{\raggedright\arraybackslash}p{(\linewidth - 14\tabcolsep) * \real{0.1250}}
  >{\raggedright\arraybackslash}p{(\linewidth - 14\tabcolsep) * \real{0.1250}}
  >{\raggedright\arraybackslash}p{(\linewidth - 14\tabcolsep) * \real{0.1250}}
  >{\raggedright\arraybackslash}p{(\linewidth - 14\tabcolsep) * \real{0.1250}}
  >{\raggedright\arraybackslash}p{(\linewidth - 14\tabcolsep) * \real{0.1250}}@{}}
\toprule\noalign{}
\begin{minipage}[b]{\linewidth}\raggedright
RUNSPEC
\end{minipage} & \begin{minipage}[b]{\linewidth}\raggedright
GRID
\end{minipage} & \begin{minipage}[b]{\linewidth}\raggedright
EDIT
\end{minipage} & \begin{minipage}[b]{\linewidth}\raggedright
PROPS
\end{minipage} & \begin{minipage}[b]{\linewidth}\raggedright
REGIONS
\end{minipage} & \begin{minipage}[b]{\linewidth}\raggedright
SOLUTION
\end{minipage} & \begin{minipage}[b]{\linewidth}\raggedright
SUMMARY
\end{minipage} & \begin{minipage}[b]{\linewidth}\raggedright
SCHEDULE
\end{minipage} \\
\midrule\noalign{}
\endhead
\bottomrule\noalign{}
\endlastfoot
\end{longtable}

\subsubsection{Description}\label{description-373}

The DIFFX keyword defines the x-direction diffusivity values for cells
in the current input box for when the Diffusivity option has been
activated by the DIFFUSE keyword in the RUNSPEC section.

This keyword is not supported by OPM Flow but would change the results
if supported so the simulation will be stopped.

\subsection{DIFFY -- Define Grid Block Y-Direction Diffusivity
Values}\label{diffy-define-grid-block-y-direction-diffusivity-values}

\begin{longtable}[]{@{}
  >{\raggedright\arraybackslash}p{(\linewidth - 14\tabcolsep) * \real{0.1250}}
  >{\raggedright\arraybackslash}p{(\linewidth - 14\tabcolsep) * \real{0.1250}}
  >{\raggedright\arraybackslash}p{(\linewidth - 14\tabcolsep) * \real{0.1250}}
  >{\raggedright\arraybackslash}p{(\linewidth - 14\tabcolsep) * \real{0.1250}}
  >{\raggedright\arraybackslash}p{(\linewidth - 14\tabcolsep) * \real{0.1250}}
  >{\raggedright\arraybackslash}p{(\linewidth - 14\tabcolsep) * \real{0.1250}}
  >{\raggedright\arraybackslash}p{(\linewidth - 14\tabcolsep) * \real{0.1250}}
  >{\raggedright\arraybackslash}p{(\linewidth - 14\tabcolsep) * \real{0.1250}}@{}}
\toprule\noalign{}
\begin{minipage}[b]{\linewidth}\raggedright
RUNSPEC
\end{minipage} & \begin{minipage}[b]{\linewidth}\raggedright
GRID
\end{minipage} & \begin{minipage}[b]{\linewidth}\raggedright
EDIT
\end{minipage} & \begin{minipage}[b]{\linewidth}\raggedright
PROPS
\end{minipage} & \begin{minipage}[b]{\linewidth}\raggedright
REGIONS
\end{minipage} & \begin{minipage}[b]{\linewidth}\raggedright
SOLUTION
\end{minipage} & \begin{minipage}[b]{\linewidth}\raggedright
SUMMARY
\end{minipage} & \begin{minipage}[b]{\linewidth}\raggedright
SCHEDULE
\end{minipage} \\
\midrule\noalign{}
\endhead
\bottomrule\noalign{}
\endlastfoot
\end{longtable}

\subsubsection{Description}\label{description-374}

The DIFFY keyword defines they y-direction diffusivity values for cells
in the current input box for when the Diffusivity option has been
activated by the DIFFUSE keyword in the RUNSPEC section.

This keyword is not supported by OPM Flow but would change the results
if supported so the simulation will be stopped.

\subsection{DIFFZ -- Define Grid Block Z-Direction Diffusivity
Values}\label{diffz-define-grid-block-z-direction-diffusivity-values}

\begin{longtable}[]{@{}
  >{\raggedright\arraybackslash}p{(\linewidth - 14\tabcolsep) * \real{0.1250}}
  >{\raggedright\arraybackslash}p{(\linewidth - 14\tabcolsep) * \real{0.1250}}
  >{\raggedright\arraybackslash}p{(\linewidth - 14\tabcolsep) * \real{0.1250}}
  >{\raggedright\arraybackslash}p{(\linewidth - 14\tabcolsep) * \real{0.1250}}
  >{\raggedright\arraybackslash}p{(\linewidth - 14\tabcolsep) * \real{0.1250}}
  >{\raggedright\arraybackslash}p{(\linewidth - 14\tabcolsep) * \real{0.1250}}
  >{\raggedright\arraybackslash}p{(\linewidth - 14\tabcolsep) * \real{0.1250}}
  >{\raggedright\arraybackslash}p{(\linewidth - 14\tabcolsep) * \real{0.1250}}@{}}
\toprule\noalign{}
\begin{minipage}[b]{\linewidth}\raggedright
RUNSPEC
\end{minipage} & \begin{minipage}[b]{\linewidth}\raggedright
GRID
\end{minipage} & \begin{minipage}[b]{\linewidth}\raggedright
EDIT
\end{minipage} & \begin{minipage}[b]{\linewidth}\raggedright
PROPS
\end{minipage} & \begin{minipage}[b]{\linewidth}\raggedright
REGIONS
\end{minipage} & \begin{minipage}[b]{\linewidth}\raggedright
SOLUTION
\end{minipage} & \begin{minipage}[b]{\linewidth}\raggedright
SUMMARY
\end{minipage} & \begin{minipage}[b]{\linewidth}\raggedright
SCHEDULE
\end{minipage} \\
\midrule\noalign{}
\endhead
\bottomrule\noalign{}
\endlastfoot
\end{longtable}

\subsubsection{Description}\label{description-375}

The DIFFZ keyword defines the z-direction diffusivity values for cells
in the current input box for when the Diffusivity option has been
activated by the DIFFUSE keyword in the RUNSPEC section.

This keyword is not supported by OPM Flow but would change the results
if supported so the simulation will be stopped.

\subsection{ECHO -- Activate Echoing of User Input Files to the Print
File}\label{echo-activate-echoing-of-user-input-files-to-the-print-file-3}

Turns on echoing of all the input files to the print file; note that
this keyword is activated by default and can subsequently be switched
off by the NOECHO activation keyword.

See ECHO -- Activate Echoing of User Input Files to the Print File in
the GLOBAL section for a full description.

\subsection{EDIT -- Define the Start of the EDIT Section of
Keywords}\label{edit-define-the-start-of-the-edit-section-of-keywords}

\begin{longtable}[]{@{}
  >{\raggedright\arraybackslash}p{(\linewidth - 14\tabcolsep) * \real{0.1250}}
  >{\raggedright\arraybackslash}p{(\linewidth - 14\tabcolsep) * \real{0.1250}}
  >{\raggedright\arraybackslash}p{(\linewidth - 14\tabcolsep) * \real{0.1250}}
  >{\raggedright\arraybackslash}p{(\linewidth - 14\tabcolsep) * \real{0.1250}}
  >{\raggedright\arraybackslash}p{(\linewidth - 14\tabcolsep) * \real{0.1250}}
  >{\raggedright\arraybackslash}p{(\linewidth - 14\tabcolsep) * \real{0.1250}}
  >{\raggedright\arraybackslash}p{(\linewidth - 14\tabcolsep) * \real{0.1250}}
  >{\raggedright\arraybackslash}p{(\linewidth - 14\tabcolsep) * \real{0.1250}}@{}}
\toprule\noalign{}
\begin{minipage}[b]{\linewidth}\raggedright
RUNSPEC
\end{minipage} & \begin{minipage}[b]{\linewidth}\raggedright
GRID
\end{minipage} & \begin{minipage}[b]{\linewidth}\raggedright
EDIT
\end{minipage} & \begin{minipage}[b]{\linewidth}\raggedright
PROPS
\end{minipage} & \begin{minipage}[b]{\linewidth}\raggedright
REGIONS
\end{minipage} & \begin{minipage}[b]{\linewidth}\raggedright
SOLUTION
\end{minipage} & \begin{minipage}[b]{\linewidth}\raggedright
SUMMARY
\end{minipage} & \begin{minipage}[b]{\linewidth}\raggedright
SCHEDULE
\end{minipage} \\
\midrule\noalign{}
\endhead
\bottomrule\noalign{}
\endlastfoot
\end{longtable}

\subsubsection{Description}\label{description-376}

The EDIT activation keyword marks the end of the GRID section and the
start of the EDIT section that enables modifications to the OPM Flow
calculated properties derived from the data entered in the GRID section,
for example grid block pore volumes via the PORV array and the
transmissibilities via the TRANX, TRANY and TRANZ family of keywords.

There is no data required for this keyword.

\subsubsection{Example}\label{example-248}

\begin{verbatim}
-- ==============================================================================
--
-- EDIT SECTION
--
-- ==============================================================================
EDIT
\end{verbatim}

The above example marks the end of the GRID section and the start of the
EDIT section in the OPM Flow data input file.

\subsection{EDITNNC -- Scale Non-Neighbor Connections Between Cells
Manually}\label{editnnc-scale-non-neighbor-connections-between-cells-manually}

\begin{longtable}[]{@{}
  >{\raggedright\arraybackslash}p{(\linewidth - 14\tabcolsep) * \real{0.1250}}
  >{\raggedright\arraybackslash}p{(\linewidth - 14\tabcolsep) * \real{0.1250}}
  >{\raggedright\arraybackslash}p{(\linewidth - 14\tabcolsep) * \real{0.1250}}
  >{\raggedright\arraybackslash}p{(\linewidth - 14\tabcolsep) * \real{0.1250}}
  >{\raggedright\arraybackslash}p{(\linewidth - 14\tabcolsep) * \real{0.1250}}
  >{\raggedright\arraybackslash}p{(\linewidth - 14\tabcolsep) * \real{0.1250}}
  >{\raggedright\arraybackslash}p{(\linewidth - 14\tabcolsep) * \real{0.1250}}
  >{\raggedright\arraybackslash}p{(\linewidth - 14\tabcolsep) * \real{0.1250}}@{}}
\toprule\noalign{}
\begin{minipage}[b]{\linewidth}\raggedright
RUNSPEC
\end{minipage} & \begin{minipage}[b]{\linewidth}\raggedright
GRID
\end{minipage} & \begin{minipage}[b]{\linewidth}\raggedright
EDIT
\end{minipage} & \begin{minipage}[b]{\linewidth}\raggedright
PROPS
\end{minipage} & \begin{minipage}[b]{\linewidth}\raggedright
REGIONS
\end{minipage} & \begin{minipage}[b]{\linewidth}\raggedright
SOLUTION
\end{minipage} & \begin{minipage}[b]{\linewidth}\raggedright
SUMMARY
\end{minipage} & \begin{minipage}[b]{\linewidth}\raggedright
SCHEDULE
\end{minipage} \\
\midrule\noalign{}
\endhead
\bottomrule\noalign{}
\endlastfoot
\end{longtable}

\subsubsection{Description}\label{description-377}

EDITNNC enables Non-Neighbor Connections (``NNC''), entered via the NNC
keyword or calculated by the simulator, to be multiplied (re-scaled) by
a constant. For example, if the existing transmissibility between
non-neighbor connections is Told and the multiplier is C, then the
resulting transmissibility, Tnew, will be
\({T}_{\mathit{new}} =C x {T}_{\mathit{old}}\). Only previously defined
NNC's entered via the NNC keyword or calculated by the simulator can be
edited, otherwise a warning message will be printed.

\begin{longtable}[]{@{}
  >{\raggedright\arraybackslash}p{(\linewidth - 6\tabcolsep) * \real{0.2500}}
  >{\raggedright\arraybackslash}p{(\linewidth - 6\tabcolsep) * \real{0.2500}}
  >{\raggedright\arraybackslash}p{(\linewidth - 6\tabcolsep) * \real{0.2500}}
  >{\raggedright\arraybackslash}p{(\linewidth - 6\tabcolsep) * \real{0.2500}}@{}}
\toprule\noalign{}
\begin{minipage}[b]{\linewidth}\raggedright
No.
\end{minipage} & \begin{minipage}[b]{\linewidth}\raggedright
Name
\end{minipage} & \begin{minipage}[b]{\linewidth}\raggedright
Description
\end{minipage} & \begin{minipage}[b]{\linewidth}\raggedright
Default
\end{minipage} \\
\midrule\noalign{}
\endhead
\bottomrule\noalign{}
\endlastfoot
Field & Metric & Laboratory & \\
1 & I1 & A positive integer that defines the first grid block in the
I-direction in a non-neighbor connection, must be greater than or equal
to one and less than or equal to NX on the DIMENS keyword in the RUNSPEC
section. & None \\
2 & J1 & A positive integer that defines the first grid block in the
J-direction in a non-neighbor connection, must be greater than or equal
to one and less than or equal to NY on the DIMENS keyword in the RUNSPEC
section. & None \\
3 & K1 & A positive integer that defines the first grid block in the
K-direction in a non-neighbor connection, must be greater than or equal
to one and less than or equal to NZ on the DIMENS keyword in the RUNSPEC
section. & None \\
4 & I2 & A positive integer that defines the second grid block in the
I-direction in a non-neighbor connection, must be greater than or equal
to one and less than or equal to NX on the DIMENS keyword in the RUNSPEC
section. & None \\
5 & J2 & A positive integer that defines the second grid block in the
J-direction in a non-neighbor connection, must be greater than or equal
to one and less than or equal to NY on the DIMENS keyword in the RUNSPEC
section. & None \\
6 & K2 & A positive integer that defines the second grid block in the
K-direction in a non-neighbor connection, must be greater than or equal
to one and less than or equal to NZ on the DIMENS keyword in the RUNSPEC
section. & None \\
7 & TRANSMUL & TRANSMUL is a positive real number greater than or equal
to zero that defines a constant that scales the transmissibility between
the first grid block (I1, J1, K1) and the second grid block (I2, J2,
K2). The default value of one means no scaling will be applied. & 1 \\
dimensionless & dimensionless & dimensionless & \\
8 & ISATNUM1 & The default value of zero means the existing saturation
table allocated to the upstream cell (I1,J1,K1). & 0 \\
9 & ISATNUM2 & ISATNUM2 is a positive integer defining which saturation
table number (relative permeability table) to be used for flow from the
second grid block to the first grid block. The default value of zero
means the existing saturation table allocated to the downstream cell
(I2,J2,K2). & 0 \\
10 & IPRSNUM1 & IPRSNUM1 is a positive integer defining which pressure
table number (PVT table) to be used for flow from the first grid block
to the second grid block. The default value of zero means the existing
PVT table allocated to the upstream cell (I1,J1,K1). & 0 \\
11 & IPRSNUM2 & IPRSNUM2 is a positive integer defining which pressure
table number (PVT table) to be used for flow from the second grid block
to the first grid block. The default value of zero means the existing
PVT table allocated to the downstream cell (I2,J2,K2). & 0 \\
12 & FACE1 & FACE1 is a character string that defines the face
associated with flow from the first grid block to the second grid block,
where FACE1 can have values of: X+, X-, Y+, Y-, Z+, or Z-. FACE1 is used
with the commercial simulator's Vertical Equilibrium option which is not
supported by OPM Flow. & None \\
13 & FACE2 & FACE2 is a character string that defines the face
associated with flow from the second grid block to the first grid block,
where FACE2 can have values of: X+, X-, Y+, Y-, Z+, or Z-. FACE2 is used
with the commercial simulator's Vertical Equilibrium option which is not
supported by OPM Flow. & None \\
14 & DIFFNNC & DIFFNNC is a positive real number greater than or equal
to zero that scales the diffusivity between the first grid block (I1,
J1, K1) and the second grid block (I2, J2, K2). & 0.0 \\
dimensionless & dimensionless & dimensionless & \\
Notes: & & & \\
\end{longtable}

\emph{Table 7.3: EDITNNC Keyword Description}

Note that although items (8) to (14) for this keyword are not available
in OPM Flow, even if they were, it is strongly recommended that these
items are defaulted if the data is being entered manually, as opposed to
being generated by pre-processing software.

Care should be taken that cells in different PVTNUM regions (see the
PVTNUM keyword in the REGIONS section) are not connected, since the
fluid properties are associated with a cell. If for example, a rbbl or a
rm3 of oil flows from PVTNUM region 1 to PVTNUM region 2, then the oil
properties of that oil will change from the PVT 1 data set to the PVT
data set 2. This will result in material balance errors, that may or may
not cause numerical issues.

If the transmissibility across a fault needs to be modified see the
FAULTS and MULTFLT keywords in the GRID section for an alternative and
less complicated method to modifying fault transmissibilities.
Transmissibility between reservoir regions can be modified by using
MULTREGT keyword, provided MULTNUM has been used to define the
inter-region transmissibility region numbers for each grid block.
Finally, the MULTX, MULTY and MULTZ series of keywords can be used to
modify transmissibility between various cells. All the aforementioned
keywords are described in the GRID section.

See also the EDITNNCR keyword in the EDIT section that resets an
existing NNC to a user defined value.

\subsubsection{Example}\label{example-249}

\begin{verbatim}
--
--       MANUALLY RESCALE NON-NEIGHBOR CONNECTIONS
--
--       ------------ BOX -----------   -- TRANSMUL --
--       I1   J1   K1    I2   J2   K2
EDITNCC
         1    1    1     1    1    2        0.2000     / SET NNC FOR FAULT
         1    1    2     1    1    3        0.2000     / SET NNC FOR FAULT
         1    1    3     1    1    4        0.2000     / SET NNC FOR FAULT
/
\end{verbatim}

The above example multiplies the transmissibility between cells (1, 1,
1) and (1, 1, 2), (1, 1, 2) and (1, 1, 3) and finally between (1, 1, 3)
and (1, 1, 4) by 0.200.

\subsection{EDITNNCR -- Reset Non-Neighbor Connections Between Cells
Manually}\label{editnncr-reset-non-neighbor-connections-between-cells-manually}

\begin{longtable}[]{@{}
  >{\raggedright\arraybackslash}p{(\linewidth - 14\tabcolsep) * \real{0.1250}}
  >{\raggedright\arraybackslash}p{(\linewidth - 14\tabcolsep) * \real{0.1250}}
  >{\raggedright\arraybackslash}p{(\linewidth - 14\tabcolsep) * \real{0.1250}}
  >{\raggedright\arraybackslash}p{(\linewidth - 14\tabcolsep) * \real{0.1250}}
  >{\raggedright\arraybackslash}p{(\linewidth - 14\tabcolsep) * \real{0.1250}}
  >{\raggedright\arraybackslash}p{(\linewidth - 14\tabcolsep) * \real{0.1250}}
  >{\raggedright\arraybackslash}p{(\linewidth - 14\tabcolsep) * \real{0.1250}}
  >{\raggedright\arraybackslash}p{(\linewidth - 14\tabcolsep) * \real{0.1250}}@{}}
\toprule\noalign{}
\begin{minipage}[b]{\linewidth}\raggedright
RUNSPEC
\end{minipage} & \begin{minipage}[b]{\linewidth}\raggedright
GRID
\end{minipage} & \begin{minipage}[b]{\linewidth}\raggedright
EDIT
\end{minipage} & \begin{minipage}[b]{\linewidth}\raggedright
PROPS
\end{minipage} & \begin{minipage}[b]{\linewidth}\raggedright
REGIONS
\end{minipage} & \begin{minipage}[b]{\linewidth}\raggedright
SOLUTION
\end{minipage} & \begin{minipage}[b]{\linewidth}\raggedright
SUMMARY
\end{minipage} & \begin{minipage}[b]{\linewidth}\raggedright
SCHEDULE
\end{minipage} \\
\midrule\noalign{}
\endhead
\bottomrule\noalign{}
\endlastfoot
\end{longtable}

\subsubsection{Description}\label{description-378}

EDITNNCR enables Non-Neighbor Connections (``NNC''), entered via the NNC
keyword or calculated by the simulator, to be reset to a user defined
value. Only previously defined NNC's entered via the NNC keyword or
calculated by the simulator can be edited, otherwise a warning message
will be printed. See also the EDITNNC keyword in the EDIT section that
scales an existing NNC.

\begin{longtable}[]{@{}
  >{\raggedright\arraybackslash}p{(\linewidth - 6\tabcolsep) * \real{0.2500}}
  >{\raggedright\arraybackslash}p{(\linewidth - 6\tabcolsep) * \real{0.2500}}
  >{\raggedright\arraybackslash}p{(\linewidth - 6\tabcolsep) * \real{0.2500}}
  >{\raggedright\arraybackslash}p{(\linewidth - 6\tabcolsep) * \real{0.2500}}@{}}
\toprule\noalign{}
\begin{minipage}[b]{\linewidth}\raggedright
No.
\end{minipage} & \begin{minipage}[b]{\linewidth}\raggedright
Name
\end{minipage} & \begin{minipage}[b]{\linewidth}\raggedright
Description
\end{minipage} & \begin{minipage}[b]{\linewidth}\raggedright
Default
\end{minipage} \\
\midrule\noalign{}
\endhead
\bottomrule\noalign{}
\endlastfoot
Field & Metric & Laboratory & \\
1 & I1 & A positive integer that defines the first grid block in the
I-direction in a non-neighbor connection, must be greater than or equal
to one and less than or equal to NX on the DIMENS keyword in the RUNSPEC
section. & None \\
2 & J1 & A positive integer that defines the first grid block in the
J-direction in a non-neighbor connection, must be greater than or equal
to one and less than or equal to NY on the DIMENS keyword in the RUNSPEC
section. & None \\
3 & K1 & A positive integer that defines the first grid block in the
K-direction in a non-neighbor connection, must be greater than or equal
to one and less than or equal to NZ on the DIMENS keyword in the RUNSPEC
section. & None \\
4 & I2 & A positive integer that defines the second grid block in the
I-direction in a non-neighbor connection, must be greater than or equal
to one and less than or equal to NX on the DIMENS keyword in the RUNSPEC
section. & None \\
5 & J2 & A positive integer that defines the second grid block in the
J-direction in a non-neighbor connection, must be greater than or equal
to one and less than or equal to NY on the DIMENS keyword in the RUNSPEC
section. & None \\
6 & K2 & A positive integer that defines the second grid block in the
K-direction in a non-neighbor connection, must be greater than or equal
to one and less than or equal to NZ on the DIMENS keyword in the RUNSPEC
section. & None \\
7 & TRANSNNC & TRANSNNC is a positive real number greater than or equal
to zero that defines the transmissibility between the first grid block
(I1, J1, K1) and the second grid block (I2, J2, K2). This value cannot
be defaulted and must be defined. & None \\
cP.rb/day/psia & cP.rm3/day/bars & cP.rcc/hr/atm & \\
8 & ISATNUM1 & ISATNUM1 is a positive integer defining which saturation
table number (relative permeability table) to be used for flow from the
first grid block to the second grid block. The default value of zero
means the existing saturation table allocated to the upstream cell
(I1,J1,K1). & 0 \\
9 & ISATNUM2 & ISATNUM2 is a positive integer defining which saturation
table number (relative permeability table) to be used for flow from the
second grid block to the first grid block. The default value of zero
means the existing saturation table allocated to the downstream cell
(I2,J2,K2). & 0 \\
10 & IPRSNUM1 & IPRSNUM1 is a positive integer defining which pressure
table number (PVT table) to be used for flow from the first grid block
to the second grid block. The default value of zero means the existing
PVT table allocated to the upstream cell (I1,J1,K1). & 0 \\
11 & IPRSNUM2 & IPRSNUM2 is a positive integer defining which pressure
table number (PVT table) to be used for flow from the second grid block
to the first grid block. The default value of zero means the existing
PVT table allocated to the downstream cell (I2,J2,K2). & 0 \\
12 & FACE1 & FACE1 is a character string that defines the face
associated with flow from the first grid block to the second grid block,
where FACE1 can have vales of: X+, X-, Y+, Y-, Z+, or Z-. & None \\
13 & FACE2 & FACE2 is a character string that defines the face
associated with flow from the second grid block to the first grid block,
where FACE2 can have vales of: X+, X-, Y+, Y-, Z+, or Z-. & None \\
14 & DIFFNNC & DIFFNNC is a positive real number greater than or equal
to zero that scales the diffusivity between the first grid block (I1,
J1, K1) and the second grid block (I2, J2, K2). The default value is the
value calculated in the GRID section. & 1* \\
feet & meters & cm & \\
Notes: & & & \\
\end{longtable}

\emph{Table 7.4: EDITNNCR Keyword Description}

Note that although items (8) to (14) for this keyword are not available
in OPM Flow, even if they were, it is strongly recommended that these
items are defaulted if the data is being entered manually, as opposed to
being generated by pre-processing software.

If the transmissibility across a fault needs to be modified see the
FAULTS and MULTFLT keywords in the GRID section for an alternative and
less complicated method to modifying fault transmissibilities.
Transmissibility between reservoir regions can be modified by using
MULTREGT keyword, provided MULTNUM has been used to define the
inter-region transmissibility region numbers for each grid block.
Finally, the MULTX, MULTY and MULTZ series of keywords can be used to
modify transmissibility between various cells. All the aforementioned
keywords are described in the GRID section.

\subsubsection{Example}\label{example-250}

\begin{verbatim}
--
--       MANUALLY RESET NON-NEIGHBOR CONNECTIONS
--
--       ------------ BOX -----------   -- TRANSNNC --
--       I1   J1   K1    I2   J2   K2
EDITNCCR
         1    1    1     1    1    2        0.2500    / RESET NNC TRANS FOR FAULT
         1    1    2     1    1    3        0.2500    / RESET NNC TRANS FAULT
         1    1    3     1    1    4        0.2500    / RESET NNC TRANS FAULT
/
\end{verbatim}

The above example res-sets the transmissibility between cells (1, 1, 1)
and (1, 1, 2), (1, 1, 2) and (1, 1, 3) and (1, 1, 3) and (1, 1, 4) to be
0.2500.

\subsection{END -- Define the End of the Input
File}\label{end-define-the-end-of-the-input-file-3}

This keyword marks the end of the input file and can occur in any
section. Any keywords and data after the END keyword are ignored.

See END -- Define the End of the Input File in the GLOBAL section for a
full description.

\subsection{ENDBOX -- Define the End of the BOX Defined
Grid}\label{endbox-define-the-end-of-the-box-defined-grid-1}

This keyword marks the end of a previously defined BOX sub-grid as
defined by a previously entered BOX keyword. The keyword resets the
input grid to be the full grid as defined by the NX, NY, and NZ
variables on the DIMENS keyword in the RUNSPEC section.

See ENDBOX -- Define the End of the BOX Defined Grid in the GRID section
for a full description.

\subsection{ENDFIN -- End the Definition of a Local Grid
Refinement}\label{endfin-end-the-definition-of-a-local-grid-refinement-1}

NDFIN defines the end of a Cartesian or radial local grid refinement
(``LGR'') definition and a LGR property definition data set.

See ENDFIN -- End the Definition of a Local Grid Refinement in the GRID
section for a full description.

\subsection{ENDINC -- Define the End of an Include
File}\label{endinc-define-the-end-of-an-include-file-3}

This keyword marks the end of an include file specified on the INCLUDE
keyword. When the ENDINC keyword is encountered in the INCLUDE file,
input data is read from the next keyword in the current file. Any
keywords and data after the ENDINC keyword in the INCLUDE file are
ignored.

See ENDINC -- Define the End of an Include File in the GLOBAL section
for a full description.

\subsection{ENDSKIP -- Deactivate Skipping of Keywords and Input
Data}\label{endskip-deactivate-skipping-of-keywords-and-input-data-3}

The ENDSKIP keyword deactivates the skipping of keywords that was
activated by the SKIP, SKIP100, or SKIP300 keywords. Each SKIP keyword
should be paired with an ENDSKIP keyword.

See ENDSKIP -- Deactivate Skipping of Keywords and Input Data in the
GLOBAL section for a full description.

\subsection{EQUALREG -- Sets an Array to a Constant by Region
Number}\label{equalreg-sets-an-array-to-a-constant-by-region-number-1}

The EQUALREG keyword sets a specified array to a constant for cells with
a specific region number. The region number array can be FLUXNUM,
MULTNUM or OPERNUM and these arrays must be defined and be available
before the EQUALREG keyword is read by the simulator. The constant can
be real or integer depending on the property array type; however, the
property arrays that can be operated on is dependent on which section
the EQUALREG keyword is being used.

See EQUALREG -- Sets an Array to a Constant by Region Number in the GRID
section for a full description.

\subsection{EQUALS -- Sets a Specified Array to a
Constant}\label{equals-sets-a-specified-array-to-a-constant-1}

The EQUALS keyword sets a specified array or part of an array to a
constant. The constant can be real or integer depending on the array
type; however, the arrays that can be operated on is dependent on which
section the EQUALS keyword is being used.

See EQUALS -- Sets a Specified Array to a Constant in the GRID section
for a full description.

\subsection{EXTRAPMS -- Activate Extrapolation Warning
Messages}\label{extrapms-activate-extrapolation-warning-messages-3}

The EXTRAPMS keyword activates extrapolation warning messages for when
OPM Flow extrapolates the PVT or VFP tables. Frequent extrapolation
warning messages should be investigated and resolved as this would
indicate possible incorrect data and may result in the simulator
extrapolating to unrealistic values.

See EXTRAPMS -- Activate Extrapolation Warning Messages in the GLOBAL
section for a full description.

\subsection{FILEUNIT -- Activate Unit Consistency
Verification}\label{fileunit-activate-unit-consistency-verification}

The FILEUNIT keyword defines the units of the data set, and is used to
verify that the units in the input deck and any associated include files
are consistent. The keyword does not provide for the conversion between
different sets of units.

See FILEUNIT -- Activate Unit Consistency Checking in the GRID section
for a full description.

\subsection{FORMFEED -- Defined the Print File Form-Feed
Character}\label{formfeed-defined-the-print-file-form-feed-character-3}

The FORMFEED keyword defines the form-feed character, or carriage
control character, for the output print (\emph{.PRT) run summary (}.RSM)
files. The keyword should be place at the very top of the input file.

See FORMFEED -- Defined the Print File Form-Feed Character in the GLOBAL
section for a dull description.

\subsection{GETDATA -- Load and Assign Data Array from INIT or RESTART
File}\label{getdata-load-and-assign-data-array-from-init-or-restart-file}

The GETDATA keyword loads a data array from a previously generated INIT
or RESTART file and assigns the loaded array to either same array in the
run or another array name.

See GETDATA -- Load and Assign Data Array from INIT or RESTART Files in
the GRID section for a full description

\subsection{HMMULT -- History Match Grid Transmissibility \& Pore Volume
Gradient Cumulative
Multipliers}\label{hmmult-history-match-grid-transmissibility-pore-volume-gradient-cumulative-multipliers}

\begin{longtable}[]{@{}
  >{\raggedright\arraybackslash}p{(\linewidth - 14\tabcolsep) * \real{0.1250}}
  >{\raggedright\arraybackslash}p{(\linewidth - 14\tabcolsep) * \real{0.1250}}
  >{\raggedright\arraybackslash}p{(\linewidth - 14\tabcolsep) * \real{0.1250}}
  >{\raggedright\arraybackslash}p{(\linewidth - 14\tabcolsep) * \real{0.1250}}
  >{\raggedright\arraybackslash}p{(\linewidth - 14\tabcolsep) * \real{0.1250}}
  >{\raggedright\arraybackslash}p{(\linewidth - 14\tabcolsep) * \real{0.1250}}
  >{\raggedright\arraybackslash}p{(\linewidth - 14\tabcolsep) * \real{0.1250}}
  >{\raggedright\arraybackslash}p{(\linewidth - 14\tabcolsep) * \real{0.1250}}@{}}
\toprule\noalign{}
\begin{minipage}[b]{\linewidth}\raggedright
RUNSPEC
\end{minipage} & \begin{minipage}[b]{\linewidth}\raggedright
GRID
\end{minipage} & \begin{minipage}[b]{\linewidth}\raggedright
EDIT
\end{minipage} & \begin{minipage}[b]{\linewidth}\raggedright
PROPS
\end{minipage} & \begin{minipage}[b]{\linewidth}\raggedright
REGIONS
\end{minipage} & \begin{minipage}[b]{\linewidth}\raggedright
SOLUTION
\end{minipage} & \begin{minipage}[b]{\linewidth}\raggedright
SUMMARY
\end{minipage} & \begin{minipage}[b]{\linewidth}\raggedright
SCHEDULE
\end{minipage} \\
\midrule\noalign{}
\endhead
\bottomrule\noalign{}
\endlastfoot
\end{longtable}

\subsubsection{Description}\label{description-379}

The HMMULT series of keywords defines the history match gradient
cumulative permeability multipliers, for when the History Match Gradient
option has been activated by the HMDIMS keyword in the RUNSPEC section.
The keyword consists of the first six characters of ``HMMULT'' followed
by a one or two character string shown in Table 7.5, that determines the
transmissibility direction, for example, HMMULTX.

This keyword is not supported by OPM Flow but has no effect on the
results so it will be ignored.

\begin{longtable}[]{@{}lllll@{}}
\toprule\noalign{}
Mnemonic & Cartesian Grid & Radial Grid & & \\
\midrule\noalign{}
\endhead
\bottomrule\noalign{}
\endlastfoot
Grid Keyword & HMMULT Keyword & Grid Keyword & HMMULT Keyword & \\
X/R & MULTX & HMMULTX & MULTR & HMMULTR \\
XY & & HMMULTXY & & \\
Y/HT & MULTY & hMMULTY & MULTTHT & HMMULTTH \\
z & MULTZ & HMMULTZ & MULTZ & HMMULTZ \\
PV & MULTPV & HMMULTPV & MULTPV & HMMULTPV \\
\end{longtable}

\emph{Table 7.5: HMMULT Keyword List}

See also the HMMLT keyword in the GRID section.

\subsection{IMPORT -- Import Grid File Data at the Current
Position}\label{import-import-grid-file-data-at-the-current-position-1}

The IMPORT keyword informs the simulator to continue reading input data
from the specified IMPORT file. When the end of the IMPORT file is
reached, input data is read from the next keyword in the current file.
Normally IMPORT files are generated by grid pre-processing software and
the keyword allows for both formatted and unformatted (binary) files to
be loaded.

See IMPORT -- Import Grid File Data at the Current Position in the GRID
section for a full description.

\subsection{INCLUDE -- Load Another Data File at the Current
Position}\label{include-load-another-data-file-at-the-current-position-3}

The INCLUDE keyword informs OPM Flow to continue reading input data from
the specified INCLUDE file. When the end of the INCLUDE file is reached,
or the ENDINC keyword is encountered in the included file, input data is
read from the next keyword in the current file.

See INCLUDE -- Load Another Data File at the Current Position in the
GLOBAL section for a full description.

\subsection{LGRCOPY -- Activate Local Grid Refinement
Inheritance}\label{lgrcopy-activate-local-grid-refinement-inheritance-2}

The LGRCOPY keyword actives Local Grid Refinement (``LGR'') Inheritance
option that allows the LGR to inherit the properties of the global or
host cell containing a LGR grid block at the start of the GRID section,
as oppose to the normal process of applying this transform at the end of
the GRID section. LGRCOPY can be used in the RUNSPEC, GRID and EDIT
sections. If used in the RUNSPEC section then the option is applied to
all LGRs defined in the input file, whereas if used in the GRID and EDIT
sections the keyword must be placed inside a LGR definition section,
that is between a CARFIN (Cartesian LGR grid) or RADIN/RADIN4 (radial
LGR grid) and the ENDFIN keyword. In the latter case inheritance is
applied on an individual LGR basis.

See LGRCOPY -- Activate Local Grid Refinement Inheritance in the RUNSPEC
section for a full description.

\subsection{MAXVALUE -- Sets a Maximum Value for an Array
Element}\label{maxvalue-sets-a-maximum-value-for-an-array-element-1}

The MAXVALUE keyword sets a maximum value for the specified array or
part of an array. The constant can be real or integer depending on the
array type; however, the arrays that can be operated on is dependent on
which section the MAXVALUE keyword is being used.

See MAXVALUE -- Sets a Maximum Value for an Array Element in the GRID
section for a full description.

\subsection{MESSAGE -- Output User
Message}\label{message-output-user-message-3}

The MESSAGE keyword outputs a user message to the terminal, as well as
to the print (\emph{.PRT) and debug (}.DBG) files. Note this is
different to the MESSAGES keyword, that defines OPM Flows message print
limits and stop limits generated by the simulator.

See MESSAGE -- Output User Message in the GLOBAL section for a full
description.

\subsection{MESSAGES -- Define Message Print Limits and Stop
Limits}\label{messages-define-message-print-limits-and-stop-limits-3}

The MESSAGES keyword defines the print and stops levels for various
messages. The ``print limits'' set the maximum number of messages that
will be printed, after which no more messages will be printed and the
``stop limits'' terminate the run when these limits are exceeded. There
are six levels of message that increase in severity from informative all
the way to programming errors, as outlined in Table 4.5.

See MESSAGES -- Define Message Print Limits and Stop Limits in the
GLOBAL section for a full description.

\subsection{MINVALUE -- Set a Minimum Value for an Array
Element}\label{minvalue-set-a-minimum-value-for-an-array-element-1}

The MINVALUE keyword sets a minimum value for the specified array or
part of an array. The constant can be real or integer depending on the
array type; however, the arrays that can be operated on is dependent on
which section the MINVALUE keyword is being used.

See MINVALUE -- Set a Minimum Value for an Array Element in the GRID
section for a full description.

\subsection{MULTFLT -- Multiply the Transmissibility of a Defined Fault
by a
Constant}\label{multflt-multiply-the-transmissibility-of-a-defined-fault-by-a-constant-1}

MULTFLT enables the transmissibilities across defined faults, as
declared by the FAULTS keyword, to be modified. They keyword allows for
the re-scaling of the existing fault transmissibilities calculated by
OPM Flow, for example setting a fault to be completely sealing by
setting the multiplier to zero..

See MULTFLT -- Multiply the Transmissibility of a Defined Fault by a
Constant in the GRID section for a full description.

\subsection{MULTIPLY -- Multiply a Specified Array by a
Constant}\label{multiply-multiply-a-specified-array-by-a-constant-1}

The MULTIPLY keyword multiplies a specified array or part of an array by
a constant. The constant can be real or integer depending on the array
type; however, the arrays that can be operated on is dependent on which
section the keyword is being used.

See MULTIPLY -- Multiply a Specified Array by a Constant in the GRID
section for a full description.

\subsection{MULTIREG -- Multiply an Array by a Constant based on a
Region
Number}\label{multireg-multiply-an-array-by-a-constant-based-on-a-region-number-1}

The MULTIREG keyword multiplies an array or part of an array by a
constant for cells with a specific region number. The region number
array can be FLUXNUM, MULTNUM or OPERNUM and these arrays must be
defined and be available before the MULTIREG keyword is read by the
simulator. The constant can be real or integer depending on the property
array type; however, the property arrays that can be operated on is
dependent on which section the MULTIREG keyword is being used.

See MULTIREG -- Multiply an Array by a Constant based on a Region Number
in the GRID section for a full description.

\subsection{MULTPV -- Multiply Cell Pore Volumes by a
Constant}\label{multpv-multiply-cell-pore-volumes-by-a-constant-1}

MULTPV multiples the pore volumes of a cell by a real positive constant
for all the cells in the model via an array. An alternative to defining
the complete array is to use the BOX keyword to define an area of the
grid and then use the MULTPV keyword to set the multipliers just for the
area defined by the BOX keyword (see the example).

See MULTPV -- Multiply Cell Pore Volumes by a Constantin the GRID
section for a full description.

\subsection{MULTR- -- Multiply Cell Transmissibility in the -R
Direction}\label{multr--multiply-cell-transmissibility-in-the--r-direction-1}

MULTR- multiples the transmissibility between two cell faces in the -X
direction for all the cells in the model via an array, that is the
keyword sets the transmissibility multiplier of block (I, J, K) between
the cells (I-1, J, K) and (I, J, K). The keyword should only be used
with radial grids, as declared by the RADIAL keyword in the RUNSPEC
section.

See MULTR- - Multiply Cell Transmissibility in the -R Direction in the
GRID section for a full description.

\subsection{MULTR -- Multiply Cell Transmissibility in the +R
Direction}\label{multr-multiply-cell-transmissibility-in-the-r-direction-1}

MULTR multiples the transmissibility between two cell faces in the +R
direction for all the cells in the model via an array, that is the
keyword sets the transmissibility multiplier of block (I, J, K) between
the cells (I, J, K) and (I+I, J, K). The keyword should only be used
with radial grids, as declared by the RADIAL keyword in the RUNSPEC
section.

See MULTR - Multiply Cell Transmissibility in the +R Direction in the
GRID section for a full description.

\subsection{MULTREGD -- Multiply Diffusivities Between
Regions}\label{multregd-multiply-diffusivities-between-regions-1}

The MULTREGT keyword multiplies the diffusivity between two regions by a
constant. The region number array can be FLUXNUM, MULTNUM or OPERNUM and
these arrays must be defined and be available before the MULTREGT
keyword is read by the simulator. The constant should be a real number.

See MULTREGD -- Multiply Diffusivities Between Regions in the GRID
section for a full description.

\subsection{MULTREGH -- Multiply Thermal Conductivities Between
Regions}\label{multregh-multiply-thermal-conductivities-between-regions-1}

The MULTREGH keyword multiplies the thermal conductivity between two
regions by a constant. The region number array can be FLUXNUM, MULTNUM
or OPERNUM and these arrays must be defined and be available before the
MULTREGT keyword is read by the simulator. The constant should be a real
number.

See MULTREGH -- Multiply Thermal Conductivities Between Regions in the
GRID section for a full description.

\subsection{MULTREGP -- Multiply Pore Volumes Based On Region
Number}\label{multregp-multiply-pore-volumes-based-on-region-number-1}

The MULTREGP keyword multiplies the pore volume of a cell by a constant
for all cells with a specific region number. The region number array can
be FLUXNUM, MULTNUM or OPERNUM and these arrays must be defined and be
available before the MULTREGP keyword is read by the simulator. The
constant should be a real number.

See MULTREGP -- Multiply Pore Volumes Based On Region Number in the GRID
section for a full description.

\subsection{MULTREGT -- Multiply Transmissibilities Between
Regions}\label{multregt-multiply-transmissibilities-between-regions-1}

The MULTREGT keyword multiplies the transmissibility between two regions
by a constant. The region number array can be FLUXNUM, MULTNUM or
OPERNUM and these arrays must be defined and be available before the
MULTREGT keyword is read by the simulator. The constant should be a real
number.

See MULTREGT -- Multiply Transmissibilities Between Regionsin the GRID
section for a full description.

\subsection{MULTTHT- -- Multiply Cell Transmissibility in the -Theta
Direction}\label{multtht--multiply-cell-transmissibility-in-the--theta-direction-1}

MULTTHT- multiples the transmissibility between two cell faces in the
-Theta direction for all the cells in the model via an array, that is
the keyword sets the transmissibility multiplier of block (I, J, K)
between the cells (I, J-1, K) and (I, J, K). The keyword should only be
used with radial grids, as declared by the RADIAL keyword in the RUNSPEC
section.

See MULTTHT- - Multiply Cell Transmissibility in the -Theta Direction in
the GRID section for a full description.

\subsection{MULTTHT -- Multiply Cell Transmissibility in the +Theta
Direction}\label{multtht-multiply-cell-transmissibility-in-the-theta-direction-1}

MULTTHT multiples the transmissibility between two cell faces in the
+Theta direction for all the cells in the model via an array, that is
the keyword sets the transmissibility multiplier of block (I, J, K)
between the cells (I, J, K) and (I, J+1, K). The keyword should only be
used with radial grids, as declared by the RADIAL keyword in the RUNSPEC
section.

See MULTTHT - Multiply Cell Transmissibility in the +Theta Direction in
the GRID section for a full description.

\subsection{MULTX- -- Multiply Cell Transmissibility in the -X
Direction}\label{multx--multiply-cell-transmissibility-in-the--x-direction-1}

MULTX- multiples the transmissibility between two cell faces in the -X
direction for all the cells in the model via an array, that is the
keyword sets the transmissibility multiplier of block (I, J, K) between
the cells (I-1, J, K) and (I, J, K).

See MULTX- - Multiply Cell Transmissibility in the -X Direction in the
GRID section for a full description.

\subsection{MULTX -- Multiply Cell Transmissibility in the +X
Direction}\label{multx-multiply-cell-transmissibility-in-the-x-direction-1}

MULTX multiples the transmissibility between two cell faces in the +X
direction for all the cells in the model via an array, that is the
keyword sets the transmissibility multiplier of block (I, J, K) between
the cells (I, J, K) and (I+I, J, K).

See MULTX - Multiply Cell Transmissibility in the +X Direction in the
GRID section for a full description.

\subsection{MULTY- -- Multiply Cell Transmissibility in the -Y
Direction}\label{multy--multiply-cell-transmissibility-in-the--y-direction-1}

MULTY- multiples the transmissibility between two cell faces in the -Y
direction for all the cells in the model via an array, that is the
keyword sets the transmissibility multiplier of block (I, J, K) between
the cells (I, J-1, K) and (I, J, K).

See MULTY- - Multiply Cell Transmissibility in the -Y Direction in the
GRID section for a full description.

\subsection{MULTY -- Multiply Cell Transmissibility in the +Y
Direction}\label{multy-multiply-cell-transmissibility-in-the-y-direction-1}

MULTY multiples the transmissibility between two cell faces in the +Y
direction for all the cells in the model via an array, that is the
keyword sets the transmissibility multiplier of block (I, J, K) between
the cells (I, J, K) and (I, J+1, K).

See MULTY - Multiply Cell Transmissibility in the +Y Direction in the
GRID section for a full description.

\subsection{MULTZ- -- Multiply Cell Transmissibility in the -Z
Direction}\label{multz--multiply-cell-transmissibility-in-the--z-direction-1}

MULTZ- multiples the transmissibility between two cell faces in the -Z
direction for all the cells in the model via an array, that is the
keyword sets the transmissibility multiplier of block (I, J, K) between
the cells (1, J, K-1) and (I, J, K).

See MULTZ- - Multiply Cell Transmissibility in the -Z Direction in the
GRID section for a full description.

\subsection{MULTZ -- Multiply Cell Transmissibility in the +Z
Direction}\label{multz-multiply-cell-transmissibility-in-the-z-direction-1}

MULTZ multiples the transmissibility between two cell faces in the +Z
direction for all the cells in the model via an array, that is the
keyword sets the transmissibility multiplier of block (I, J, K) between
the cells (I, J, K) and (I, J, K+1).

See MULTZ - Multiply Cell Transmissibility in the +Z Direction in the
GRID section for a full description.

\subsection{NOECHO -- Deactivate Echoing of User Input Files to the
Print
File}\label{noecho-deactivate-echoing-of-user-input-files-to-the-print-file-3}

Turns off echoing of all the input files to the print file. Note by
default echoing of the inputs files is active. but can subsequently be
switched off by the NOECHO activation keyword.

See NOECHO -- Deactivate Echoing of User Input Files to the Print File
in the GLOBAL section for a full description.

\subsection{NOWARN -- Deactivate Warning
Messages}\label{nowarn-deactivate-warning-messages-3}

Turns off warning messages to be printed to the print file; note that
this keyword is deactivated by default and can subsequently be switched
off by the WARN activation keyword. The warning messages may be turned
on and off using keywords WARN and NOWARN.

See NOWARN -- Deactivate Warning Messages in the GLOBAL section for a
full description.

\subsection{OPERATE -- Define Mathematical Operations on
Arrays}\label{operate-define-mathematical-operations-on-arrays-1}

This keyword, OPERATE, defines mathematical operations on property
arrays (NTG, PORO etc.) and optionally using another property array as
input to the function. The keyword allows for various mathematical
functions and their associated variables to be defined and applied to
the selected array data. Input constants can be real or integer
depending on the array type; however, the arrays that can be operated on
is dependent on which section the keyword is being used.

See OPERATE -- Define Mathematical Operations on Arrays in the GRID
section for a full description.

\subsection{OPERATER -- Define Mathematical Operations on Arrays by
Region}\label{operater-define-mathematical-operations-on-arrays-by-region-1}

This keyword defines the mathematical operations on arrays for specific
regions in the commercial simulator and is currently not supported by
OPM Flow. However, similar functionality is provided by the ADD and
MULTIPLY keywords.

See OPERATER -- Define Mathematical Operations on Arrays by Region in
the GRID section for a full description.

\subsection{PORV -- Define the Pore Volumes for All the
Cells}\label{porv-define-the-pore-volumes-for-all-the-cells}

\begin{longtable}[]{@{}
  >{\raggedright\arraybackslash}p{(\linewidth - 14\tabcolsep) * \real{0.1250}}
  >{\raggedright\arraybackslash}p{(\linewidth - 14\tabcolsep) * \real{0.1250}}
  >{\raggedright\arraybackslash}p{(\linewidth - 14\tabcolsep) * \real{0.1250}}
  >{\raggedright\arraybackslash}p{(\linewidth - 14\tabcolsep) * \real{0.1250}}
  >{\raggedright\arraybackslash}p{(\linewidth - 14\tabcolsep) * \real{0.1250}}
  >{\raggedright\arraybackslash}p{(\linewidth - 14\tabcolsep) * \real{0.1250}}
  >{\raggedright\arraybackslash}p{(\linewidth - 14\tabcolsep) * \real{0.1250}}
  >{\raggedright\arraybackslash}p{(\linewidth - 14\tabcolsep) * \real{0.1250}}@{}}
\toprule\noalign{}
\begin{minipage}[b]{\linewidth}\raggedright
RUNSPEC
\end{minipage} & \begin{minipage}[b]{\linewidth}\raggedright
GRID
\end{minipage} & \begin{minipage}[b]{\linewidth}\raggedright
EDIT
\end{minipage} & \begin{minipage}[b]{\linewidth}\raggedright
PROPS
\end{minipage} & \begin{minipage}[b]{\linewidth}\raggedright
REGIONS
\end{minipage} & \begin{minipage}[b]{\linewidth}\raggedright
SOLUTION
\end{minipage} & \begin{minipage}[b]{\linewidth}\raggedright
SUMMARY
\end{minipage} & \begin{minipage}[b]{\linewidth}\raggedright
SCHEDULE
\end{minipage} \\
\midrule\noalign{}
\endhead
\bottomrule\noalign{}
\endlastfoot
\end{longtable}

\subsubsection{Description}\label{description-380}

PORV defines the pore volumes for all the cells in the model via an
array. The keyword can be used for all grid types, except for the Radial
Grid geometry. The keyword effectively overwrites previously entered and
calculated data. The area to be modified can be defined via the various
grid selection keywords, ADD, BOX, EQUALS, etc., and areas that are not
selected remain unchanged.

\begin{longtable}[]{@{}
  >{\raggedright\arraybackslash}p{(\linewidth - 6\tabcolsep) * \real{0.2500}}
  >{\raggedright\arraybackslash}p{(\linewidth - 6\tabcolsep) * \real{0.2500}}
  >{\raggedright\arraybackslash}p{(\linewidth - 6\tabcolsep) * \real{0.2500}}
  >{\raggedright\arraybackslash}p{(\linewidth - 6\tabcolsep) * \real{0.2500}}@{}}
\toprule\noalign{}
\begin{minipage}[b]{\linewidth}\raggedright
No.
\end{minipage} & \begin{minipage}[b]{\linewidth}\raggedright
Name
\end{minipage} & \begin{minipage}[b]{\linewidth}\raggedright
Description
\end{minipage} & \begin{minipage}[b]{\linewidth}\raggedright
Default
\end{minipage} \\
\midrule\noalign{}
\endhead
\bottomrule\noalign{}
\endlastfoot
Field & Metric & Laboratory & \\
1 & PORV & PORV is an array of real positive numbers assigning a pore
volume to each cell in the model. Only the values in the currently
defined input BOX needed be entered. Repeat counts may be used, for
example 20*100.0. & None \\
rb & rm3 & rcc & \\
Notes: & & & \\
\end{longtable}

\emph{Table 7.6: PORV Keyword Description}

\subsubsection{Example}\label{example-251}

\begin{verbatim}
--
--       DEFINE INPUT BOX FOR EDITING INPUT ARRAYS
--
--       ---------- BOX ---------
--       I1  I2   J1  J2   K1  K2
BOX
         1* 100   1* 100   20  20                          / DEFINE BOX AREA
--
--       SET PORV FOR THE GRID BLOCKS
--
PORV
         1000*0.00                                         /
--
--       DEFINE END OF INPUT BOX EDITING OF INPUT ARRAYS
--
ENDBOX
\end{verbatim}

Here the BOX statement is used to define the input grid for the PORV
keyword, which overwrites the pore volume previously calculated with
pore volume values of zero, resulting in a no-flow boundary in that part
of the field between layers 19 and 21, since layer 20 is deactivated.
The ENDBOX keyword resets the input box to the full grid.

\subsection{PYEND -- End the Definition of a PYINPUT
Section}\label{pyend-end-the-definition-of-a-pyinput-section-1}

The PYINPUT and PYEND keywords are a part of OPM Flow's Python scripting
facility that processes standard Python commands that can be used to
manipulate and define the simulators input parameters during processing
of the input deck. The main purpose of the facility is to script the
construction of the various keywords. A PYINPUT Definition Section is
terminated by a PYEND keyword on a separate single line.

See PYEND -- End the Definition of a PYINPUT Section in the GRID section
for a full description.

\subsection{PYINPUT -- Define the Start of a PYINPUT
Section}\label{pyinput-define-the-start-of-a-pyinput-section-1}

The PYINPUT and PYEND keywords are a part of OPM Flow's Python scripting
facility that processes standard Python commands that can be used to
manipulate and define the simulators input parameters during processing
of the input deck. The main purpose of the facility is to script the
construction of the various keywords used by the simulator. PYINPUT
declares the start of a PYINPUT Definition Section on a single separate
line, which is then followed by various standard Python commands, with
one command per line.

See PYINPUT -- Define the Start of a PYINPUT Section in the GRID section
for a full description.

\subsection{REFINE -- Start the Definition of a Local Grid
Refinement}\label{refine-start-the-definition-of-a-local-grid-refinement-1}

The REFINE keyword defines the start of a Cartesian or radial local grid
refinement (``LGR'') definition that sets the properties of the selected
LGR. The keyword is then followed by the property keywords associated
with the section where the keyword is being invoked. For example, if the
REFINE keyword is used in the GRID section then most of the keywords in
that section can be used to set the grid properties for the LGR.

See REFINE -- Start the Definition of a Local Grid Refinement in the
GRID section for a full description.

\subsection{SKIP -- Activate Skipping of All Keywords and Input
Data}\label{skip-activate-skipping-of-all-keywords-and-input-data-3}

The SKIP keyword activates skipping of all keywords and input data until
the ENDSKIP keyword is encountered. All keywords between the SKIP and
ENDSKIP keywords are ignored.

See SKIP -- Activate Skipping of All Keywords and Input Data in the
GLOBAL section for a full description.

\subsection{SKIP100 -- Activate Skipping of Black-Oil Keywords and Input
Data}\label{skip100-activate-skipping-of-black-oil-keywords-and-input-data-2}

The SKIP100 keyword activates skipping of all keywords and input data by
the commercial black-oil simulator until the ENDSKIP keyword is
encountered. All keywords between the SKIP100 and ENDSKIP keywords are
ignored by the commercial black-oil simulator. The SKIP100 keyword is
ignored by the commercial compositional simulator.

See SKIP100 -- Activate Skipping of Keywords by Black-Oil Simulator in
the GLOBAL section for a full description.

\subsection{SKIP300 -- Activate Skipping of Keywords by Compositional
Simulator}\label{skip300-activate-skipping-of-keywords-by-compositional-simulator-3}

The SKIP300 keyword activates skipping of all keywords and input data by
the commercial compositional simulator until the ENDSKIP keyword is
encountered. All keywords between the SKIP300 and ENDSKIP keywords are
ignored by the commercial compositional simulator. The SKIP300 keyword
is ignored by the commercial black-oil simulator.

See Error: Reference source not found in the GLOBAL section for a full
description.

\subsection{TRANR -- Define the Transmissibility in the +R Direction for
All the
Cells}\label{tranr-define-the-transmissibility-in-the-r-direction-for-all-the-cells}

\begin{longtable}[]{@{}
  >{\raggedright\arraybackslash}p{(\linewidth - 14\tabcolsep) * \real{0.1250}}
  >{\raggedright\arraybackslash}p{(\linewidth - 14\tabcolsep) * \real{0.1250}}
  >{\raggedright\arraybackslash}p{(\linewidth - 14\tabcolsep) * \real{0.1250}}
  >{\raggedright\arraybackslash}p{(\linewidth - 14\tabcolsep) * \real{0.1250}}
  >{\raggedright\arraybackslash}p{(\linewidth - 14\tabcolsep) * \real{0.1250}}
  >{\raggedright\arraybackslash}p{(\linewidth - 14\tabcolsep) * \real{0.1250}}
  >{\raggedright\arraybackslash}p{(\linewidth - 14\tabcolsep) * \real{0.1250}}
  >{\raggedright\arraybackslash}p{(\linewidth - 14\tabcolsep) * \real{0.1250}}@{}}
\toprule\noalign{}
\begin{minipage}[b]{\linewidth}\raggedright
RUNSPEC
\end{minipage} & \begin{minipage}[b]{\linewidth}\raggedright
GRID
\end{minipage} & \begin{minipage}[b]{\linewidth}\raggedright
EDIT
\end{minipage} & \begin{minipage}[b]{\linewidth}\raggedright
PROPS
\end{minipage} & \begin{minipage}[b]{\linewidth}\raggedright
REGIONS
\end{minipage} & \begin{minipage}[b]{\linewidth}\raggedright
SOLUTION
\end{minipage} & \begin{minipage}[b]{\linewidth}\raggedright
SUMMARY
\end{minipage} & \begin{minipage}[b]{\linewidth}\raggedright
SCHEDULE
\end{minipage} \\
\midrule\noalign{}
\endhead
\bottomrule\noalign{}
\endlastfoot
\end{longtable}

\subsubsection{Description}\label{description-381}

TRANR defines the transmissibility in the +R direction for all the cells
in the model via an array. The keyword can only be used with Radial Grid
geometry grids. The keyword effectively overwrites previously entered
and calculated data. The transmissibility overwritten is the +R face
transmissibility of each grid block, that is for cell (I, J, K) the
transmissibility between cells (I, J, K) and (I+1, J, K).

This keyword is not supported by OPM Flow but would change the results
if supported so the simulation will be stopped.

\begin{longtable}[]{@{}
  >{\raggedright\arraybackslash}p{(\linewidth - 6\tabcolsep) * \real{0.2500}}
  >{\raggedright\arraybackslash}p{(\linewidth - 6\tabcolsep) * \real{0.2500}}
  >{\raggedright\arraybackslash}p{(\linewidth - 6\tabcolsep) * \real{0.2500}}
  >{\raggedright\arraybackslash}p{(\linewidth - 6\tabcolsep) * \real{0.2500}}@{}}
\toprule\noalign{}
\begin{minipage}[b]{\linewidth}\raggedright
No.
\end{minipage} & \begin{minipage}[b]{\linewidth}\raggedright
Name
\end{minipage} & \begin{minipage}[b]{\linewidth}\raggedright
Description
\end{minipage} & \begin{minipage}[b]{\linewidth}\raggedright
Default
\end{minipage} \\
\midrule\noalign{}
\endhead
\bottomrule\noalign{}
\endlastfoot
Field & Metric & Laboratory & \\
1 & TRANR & TRANR is an array of real positive numbers assigning the
transmissibility in the R direction to each cell in the model. Repeat
counts may be used, for example 20*100.0. & None \\
cP.rb/day/psia & cP.rm3/day/bars & cP.rcc/hr/atm & \\
Notes: & & & \\
\end{longtable}

\emph{Table 7.7: TRANR Keyword Description}

See also the TRANTHT and TRANYZ keywords to modify the
transmissibilities in the other directions.

\subsubsection{Example}\label{example-252}

\begin{verbatim}
--
--       DEFINE INPUT BOX FOR EDITING INPUT ARRAYS
--
--       ---------- BOX ---------
--       I1  I2   J1  J2   K1  K2
BOX
         1   1    10  10   1  120                         / DEFINE BOX AREA
--
--       SET TRANR+ TRANSMISSIBILITY
--
TRANR
         120*0.00                                         /
--
--       DEFINE END OF INPUT BOX EDITING OF INPUT ARRAYS
--
ENDBOX
\end{verbatim}

Here the BOX statement is used to define the input grid for the TRANR
keyword, which overwrites the transmissibility previously calculated
with transmissibility values of zero, resulting in a no-flow boundary in
that part of the grid. The ENDBOX keyword resets the input box to the
full grid.

\subsection{TRANTHT -- Define the Transmissibility in the +Theta
Direction for All the
Cells}\label{trantht-define-the-transmissibility-in-the-theta-direction-for-all-the-cells}

\begin{longtable}[]{@{}
  >{\raggedright\arraybackslash}p{(\linewidth - 14\tabcolsep) * \real{0.1250}}
  >{\raggedright\arraybackslash}p{(\linewidth - 14\tabcolsep) * \real{0.1250}}
  >{\raggedright\arraybackslash}p{(\linewidth - 14\tabcolsep) * \real{0.1250}}
  >{\raggedright\arraybackslash}p{(\linewidth - 14\tabcolsep) * \real{0.1250}}
  >{\raggedright\arraybackslash}p{(\linewidth - 14\tabcolsep) * \real{0.1250}}
  >{\raggedright\arraybackslash}p{(\linewidth - 14\tabcolsep) * \real{0.1250}}
  >{\raggedright\arraybackslash}p{(\linewidth - 14\tabcolsep) * \real{0.1250}}
  >{\raggedright\arraybackslash}p{(\linewidth - 14\tabcolsep) * \real{0.1250}}@{}}
\toprule\noalign{}
\begin{minipage}[b]{\linewidth}\raggedright
RUNSPEC
\end{minipage} & \begin{minipage}[b]{\linewidth}\raggedright
GRID
\end{minipage} & \begin{minipage}[b]{\linewidth}\raggedright
EDIT
\end{minipage} & \begin{minipage}[b]{\linewidth}\raggedright
PROPS
\end{minipage} & \begin{minipage}[b]{\linewidth}\raggedright
REGIONS
\end{minipage} & \begin{minipage}[b]{\linewidth}\raggedright
SOLUTION
\end{minipage} & \begin{minipage}[b]{\linewidth}\raggedright
SUMMARY
\end{minipage} & \begin{minipage}[b]{\linewidth}\raggedright
SCHEDULE
\end{minipage} \\
\midrule\noalign{}
\endhead
\bottomrule\noalign{}
\endlastfoot
\end{longtable}

\subsubsection{Description}\label{description-382}

TRANTHT defines the transmissibility in the +Theta direction for all the
cells in the model via an array. The keyword can only be used with
Radial Grid geometry grids. The keyword effectively overwrites
previously entered and calculated data. The transmissibility overwritten
is the +Theta face transmissibility of each grid block, that is for cell
(I, J, K) the transmissibility between cells (I, J, K) and (I, J+1, K).

This keyword is not supported by OPM Flow but would change the results
if supported so the simulation will be stopped.

\begin{longtable}[]{@{}
  >{\raggedright\arraybackslash}p{(\linewidth - 6\tabcolsep) * \real{0.2500}}
  >{\raggedright\arraybackslash}p{(\linewidth - 6\tabcolsep) * \real{0.2500}}
  >{\raggedright\arraybackslash}p{(\linewidth - 6\tabcolsep) * \real{0.2500}}
  >{\raggedright\arraybackslash}p{(\linewidth - 6\tabcolsep) * \real{0.2500}}@{}}
\toprule\noalign{}
\begin{minipage}[b]{\linewidth}\raggedright
No.
\end{minipage} & \begin{minipage}[b]{\linewidth}\raggedright
Name
\end{minipage} & \begin{minipage}[b]{\linewidth}\raggedright
Description
\end{minipage} & \begin{minipage}[b]{\linewidth}\raggedright
Default
\end{minipage} \\
\midrule\noalign{}
\endhead
\bottomrule\noalign{}
\endlastfoot
Field & Metric & Laboratory & \\
1 & TRANTHT & TRANTHT is an array of real positive numbers assigning the
transmissibility in the +Theta direction to each cell in the model.
Repeat counts may be used, for example 20*100.0. & None \\
cP.rb/day/psia & cP.rm3/day/bars & cP.rcc/hr/atm & \\
Notes: & & & \\
\end{longtable}

\emph{Table 7.8: TRANR Keyword Description}

See also the TRANR and TRANYZ keywords to modify the transmissibilities
in the other directions.

\subsubsection{Example}\label{example-253}

\begin{verbatim}
--
--       DEFINE INPUT BOX FOR EDITING INPUT ARRAYS
--
--       ---------- BOX ---------
--       I1  I2   J1  J2   K1  K2
BOX
         10  10   1   6    1   3                           / DEFINE BOX AREA
--
--       SET TRANTHT TRANSMISSIBILITY
--
TRANTHT
         18*0.00                                           /
--
--       DEFINE END OF INPUT BOX EDITING OF INPUT ARRAYS
--
ENDBOX
\end{verbatim}

Here the BOX statement is used to define the input grid for the TRANTHT
keyword, which overwrites the transmissibility previously calculated
with transmissibility values of zero, resulting in a no-flow boundary in
that part of the grid. The ENDBOX keyword resets the input box to the
full grid.

\subsection{TRANX -- Define the Transmissibility in the X Direction for
All the
Cells}\label{tranx-define-the-transmissibility-in-the-x-direction-for-all-the-cells}

\begin{longtable}[]{@{}
  >{\raggedright\arraybackslash}p{(\linewidth - 14\tabcolsep) * \real{0.1250}}
  >{\raggedright\arraybackslash}p{(\linewidth - 14\tabcolsep) * \real{0.1250}}
  >{\raggedright\arraybackslash}p{(\linewidth - 14\tabcolsep) * \real{0.1250}}
  >{\raggedright\arraybackslash}p{(\linewidth - 14\tabcolsep) * \real{0.1250}}
  >{\raggedright\arraybackslash}p{(\linewidth - 14\tabcolsep) * \real{0.1250}}
  >{\raggedright\arraybackslash}p{(\linewidth - 14\tabcolsep) * \real{0.1250}}
  >{\raggedright\arraybackslash}p{(\linewidth - 14\tabcolsep) * \real{0.1250}}
  >{\raggedright\arraybackslash}p{(\linewidth - 14\tabcolsep) * \real{0.1250}}@{}}
\toprule\noalign{}
\begin{minipage}[b]{\linewidth}\raggedright
RUNSPEC
\end{minipage} & \begin{minipage}[b]{\linewidth}\raggedright
GRID
\end{minipage} & \begin{minipage}[b]{\linewidth}\raggedright
EDIT
\end{minipage} & \begin{minipage}[b]{\linewidth}\raggedright
PROPS
\end{minipage} & \begin{minipage}[b]{\linewidth}\raggedright
REGIONS
\end{minipage} & \begin{minipage}[b]{\linewidth}\raggedright
SOLUTION
\end{minipage} & \begin{minipage}[b]{\linewidth}\raggedright
SUMMARY
\end{minipage} & \begin{minipage}[b]{\linewidth}\raggedright
SCHEDULE
\end{minipage} \\
\midrule\noalign{}
\endhead
\bottomrule\noalign{}
\endlastfoot
\end{longtable}

\subsubsection{Description}\label{description-383}

TRANX defines the transmissibility in the X direction for all the cells
in the model via an array. The keyword can be used for all grid types,
except for the Radial Grid geometry. The keyword effectively overwrites
previously entered and calculated data. The transmissibility overwritten
is the +X face transmissibility of each grid block, that is for cell (I,
J, K) the transmissibility between cells (I, J, K) and (I+1, J, K).

\begin{longtable}[]{@{}
  >{\raggedright\arraybackslash}p{(\linewidth - 6\tabcolsep) * \real{0.2500}}
  >{\raggedright\arraybackslash}p{(\linewidth - 6\tabcolsep) * \real{0.2500}}
  >{\raggedright\arraybackslash}p{(\linewidth - 6\tabcolsep) * \real{0.2500}}
  >{\raggedright\arraybackslash}p{(\linewidth - 6\tabcolsep) * \real{0.2500}}@{}}
\toprule\noalign{}
\begin{minipage}[b]{\linewidth}\raggedright
No.
\end{minipage} & \begin{minipage}[b]{\linewidth}\raggedright
Name
\end{minipage} & \begin{minipage}[b]{\linewidth}\raggedright
Description
\end{minipage} & \begin{minipage}[b]{\linewidth}\raggedright
Default
\end{minipage} \\
\midrule\noalign{}
\endhead
\bottomrule\noalign{}
\endlastfoot
Field & Metric & Laboratory & \\
1 & TRANX & TRANX is an array of real positive numbers assigning the
transmissibility in the X direction to each cell in the model. Repeat
counts may be used, for example 20*100.0. & None \\
cP.rb/day/psia & cP.rm3/day/bars & cP.rcc/hr/atm & \\
Notes: & & & \\
\end{longtable}

\emph{Table 7.9: TRANX Keyword Description}

See also the TRANY and TRANYZ keywords to modify the transmissibilities
in the other directions.

\subsubsection{Example}\label{example-254}

\begin{verbatim}
--
--       DEFINE INPUT BOX FOR EDITING INPUT ARRAYS
--
--       ---------- BOX ---------
--       I1  I2   J1  J2   K1  K2
BOX
         1   1    10  10   1  120                         / DEFINE BOX AREA
--
--       SET TRANX+ TRANSMISSIBILITY
--
TRANX
         120*0.00                                         /
--
--       DEFINE END OF INPUT BOX EDITING OF INPUT ARRAYS
--
ENDBOX
\end{verbatim}

Here the BOX statement is used to define the input grid for the TRANX
keyword, which overwrites the transmissibility previously calculated
with transmissibility values of zero, resulting in a no-flow boundary in
that part of the field. The ENDBOX keyword resets the input box to the
full grid.

\subsection{TRANY -- Define the Transmissibility in the Y Direction for
All the
Cells}\label{trany-define-the-transmissibility-in-the-y-direction-for-all-the-cells}

\begin{longtable}[]{@{}
  >{\raggedright\arraybackslash}p{(\linewidth - 14\tabcolsep) * \real{0.1250}}
  >{\raggedright\arraybackslash}p{(\linewidth - 14\tabcolsep) * \real{0.1250}}
  >{\raggedright\arraybackslash}p{(\linewidth - 14\tabcolsep) * \real{0.1250}}
  >{\raggedright\arraybackslash}p{(\linewidth - 14\tabcolsep) * \real{0.1250}}
  >{\raggedright\arraybackslash}p{(\linewidth - 14\tabcolsep) * \real{0.1250}}
  >{\raggedright\arraybackslash}p{(\linewidth - 14\tabcolsep) * \real{0.1250}}
  >{\raggedright\arraybackslash}p{(\linewidth - 14\tabcolsep) * \real{0.1250}}
  >{\raggedright\arraybackslash}p{(\linewidth - 14\tabcolsep) * \real{0.1250}}@{}}
\toprule\noalign{}
\begin{minipage}[b]{\linewidth}\raggedright
RUNSPEC
\end{minipage} & \begin{minipage}[b]{\linewidth}\raggedright
GRID
\end{minipage} & \begin{minipage}[b]{\linewidth}\raggedright
EDIT
\end{minipage} & \begin{minipage}[b]{\linewidth}\raggedright
PROPS
\end{minipage} & \begin{minipage}[b]{\linewidth}\raggedright
REGIONS
\end{minipage} & \begin{minipage}[b]{\linewidth}\raggedright
SOLUTION
\end{minipage} & \begin{minipage}[b]{\linewidth}\raggedright
SUMMARY
\end{minipage} & \begin{minipage}[b]{\linewidth}\raggedright
SCHEDULE
\end{minipage} \\
\midrule\noalign{}
\endhead
\bottomrule\noalign{}
\endlastfoot
\end{longtable}

\subsubsection{Description}\label{description-384}

TRANY defines the transmissibility in the Y direction for all the cells
in the model via an array. The keyword can be used for all grid types,
except for the Radial Grid geometry. The keyword effectively overwrites
previously entered and calculated data. The transmissibility overwritten
is the +Y face transmissibility of each grid block, that is for cell (I,
J, K) the transmissibility between cells (I, J, K) and (I, J+1, K).

\begin{longtable}[]{@{}
  >{\raggedright\arraybackslash}p{(\linewidth - 6\tabcolsep) * \real{0.2500}}
  >{\raggedright\arraybackslash}p{(\linewidth - 6\tabcolsep) * \real{0.2500}}
  >{\raggedright\arraybackslash}p{(\linewidth - 6\tabcolsep) * \real{0.2500}}
  >{\raggedright\arraybackslash}p{(\linewidth - 6\tabcolsep) * \real{0.2500}}@{}}
\toprule\noalign{}
\begin{minipage}[b]{\linewidth}\raggedright
No.
\end{minipage} & \begin{minipage}[b]{\linewidth}\raggedright
Name
\end{minipage} & \begin{minipage}[b]{\linewidth}\raggedright
Description
\end{minipage} & \begin{minipage}[b]{\linewidth}\raggedright
Default
\end{minipage} \\
\midrule\noalign{}
\endhead
\bottomrule\noalign{}
\endlastfoot
Field & Metric & Laboratory & \\
1 & TRANY & TRANY is an array of real positive numbers assigning the
transmissibility in the Y direction to each cell in the model. Repeat
counts may be used, for example 20*100.0. & None \\
cP.rb/day/psia & cP.rm3/day/bars & cP.rcc/hr/atm & \\
Notes: & & & \\
\end{longtable}

\emph{Table 7.10: TRANY Keyword Description}

See also the TRANX and TRANZ keywords to modify the transmissibilities
in the other directions.

\subsubsection{Example}\label{example-255}

\begin{verbatim}
--
--       DEFINE INPUT BOX FOR EDITING INPUT ARRAYS
--
--       ---------- BOX ---------
--       I1  I2   J1  J2   K1  K2
BOX
         1   1    10  10   1  120                         / DEFINE BOX AREA
--
--       SET TRANY+ TRANSMISSIBILITY
--
TRANY
         120*0.00                                         /
--
--       DEFINE END OF INPUT BOX EDITING OF INPUT ARRAYS
--
ENDBOX
\end{verbatim}

Here the BOX statement is used to define the input grid for the TRANY
keyword, which overwrites the transmissibility previously calculated
with transmissibility values of zero, resulting in a no-flow boundary in
that part of the field. The ENDBOX keyword resets the input box to the
full grid.

\subsection{TRANZ -- Define the Transmissibility in the Z Direction for
All the
Cells}\label{tranz-define-the-transmissibility-in-the-z-direction-for-all-the-cells}

\begin{longtable}[]{@{}
  >{\raggedright\arraybackslash}p{(\linewidth - 14\tabcolsep) * \real{0.1250}}
  >{\raggedright\arraybackslash}p{(\linewidth - 14\tabcolsep) * \real{0.1250}}
  >{\raggedright\arraybackslash}p{(\linewidth - 14\tabcolsep) * \real{0.1250}}
  >{\raggedright\arraybackslash}p{(\linewidth - 14\tabcolsep) * \real{0.1250}}
  >{\raggedright\arraybackslash}p{(\linewidth - 14\tabcolsep) * \real{0.1250}}
  >{\raggedright\arraybackslash}p{(\linewidth - 14\tabcolsep) * \real{0.1250}}
  >{\raggedright\arraybackslash}p{(\linewidth - 14\tabcolsep) * \real{0.1250}}
  >{\raggedright\arraybackslash}p{(\linewidth - 14\tabcolsep) * \real{0.1250}}@{}}
\toprule\noalign{}
\begin{minipage}[b]{\linewidth}\raggedright
RUNSPEC
\end{minipage} & \begin{minipage}[b]{\linewidth}\raggedright
GRID
\end{minipage} & \begin{minipage}[b]{\linewidth}\raggedright
EDIT
\end{minipage} & \begin{minipage}[b]{\linewidth}\raggedright
PROPS
\end{minipage} & \begin{minipage}[b]{\linewidth}\raggedright
REGIONS
\end{minipage} & \begin{minipage}[b]{\linewidth}\raggedright
SOLUTION
\end{minipage} & \begin{minipage}[b]{\linewidth}\raggedright
SUMMARY
\end{minipage} & \begin{minipage}[b]{\linewidth}\raggedright
SCHEDULE
\end{minipage} \\
\midrule\noalign{}
\endhead
\bottomrule\noalign{}
\endlastfoot
\end{longtable}

\subsubsection{Description}\label{description-385}

TRANX defines the transmissibility in the z direction for all the cells
in the model via an array. The keyword can be used for all grid types,
except for the Radial Grid geometry. The keyword effectively overwrites
previously entered and calculated data. The transmissibility overwritten
is the +Z face transmissibility of each grid block, that is for cell (I,
J, K) the transmissibility between cells (I, J, K) and (I, J, K+1).

\begin{longtable}[]{@{}
  >{\raggedright\arraybackslash}p{(\linewidth - 6\tabcolsep) * \real{0.2500}}
  >{\raggedright\arraybackslash}p{(\linewidth - 6\tabcolsep) * \real{0.2500}}
  >{\raggedright\arraybackslash}p{(\linewidth - 6\tabcolsep) * \real{0.2500}}
  >{\raggedright\arraybackslash}p{(\linewidth - 6\tabcolsep) * \real{0.2500}}@{}}
\toprule\noalign{}
\begin{minipage}[b]{\linewidth}\raggedright
No.
\end{minipage} & \begin{minipage}[b]{\linewidth}\raggedright
Name
\end{minipage} & \begin{minipage}[b]{\linewidth}\raggedright
Description
\end{minipage} & \begin{minipage}[b]{\linewidth}\raggedright
Default
\end{minipage} \\
\midrule\noalign{}
\endhead
\bottomrule\noalign{}
\endlastfoot
Field & Metric & Laboratory & \\
1 & TRANZ & TRANZ is an array of real positive numbers assigning the
transmissibility in the Z direction to each cell in the model. Repeat
counts may be used, for example 20*100.0. & None \\
cP.rb/day/psia & cP.rm3/day/bars & cP.rcc/hr/atm & \\
Notes: & & & \\
\end{longtable}

\emph{Table 7.11: TRANZ Keyword Description}

See also the TRANX and TRANY keywords to modify the transmissibilities
in the other directions.

\subsubsection{Example}\label{example-256}

\begin{verbatim}
--
--       DEFINE INPUT BOX FOR EDITING INPUT ARRAYS
--
--       ---------- BOX ---------
--       I1  I2   J1  J2   K1  K2
BOX
         1* 100   1* 100   20  20                          / DEFINE BOX AREA
--
--       SET TRANZ+ TRANSMISSIBILITY
--
TRANZ
         1000*0.00                                         /
--
--       DEFINE END OF INPUT BOX EDITING OF INPUT ARRAYS
--
ENDBOX
\end{verbatim}

Here the BOX statement is used to define the input grid for the TRANZ
keyword, which overwrites the transmissibility previously calculated
with transmissibility values of zero, resulting in a no-flow boundary in
that part of the field between layers 20 and 21. The ENDBOX keyword
resets the input box to the full grid.

\subsection{WARN -- Activate Warning
Messages}\label{warn-activate-warning-messages-3}

Turns on warning messages to be printed to the print file (*.PRT); note
that this keyword is activated by default and can subsequently be
switched off by the NOWARN activation keyword. The warning messages may
be turned on and off using keywords WARN and NOWARN. OPM Flow always
prints error messages.

See WARN -- Activate Warning Messages in the GLOBAL section for a full
description.

\bookmarksetup{startatroot}

\chapter{PROPS SECTION}\label{props-section-1}

\section{Introduction}\label{introduction-6}

The PROPS section is the section that contains the fluid property
keywords used to define the PVT behavior of the fluids in the model and
therefore the data is fluid type dependent. For example, if the oil
phase has been activated in the RUNSPEC section via the OIL keyword then
one of the oil PVT keywords needs to be defined in this section. The
other main type of data required for the PROPS section is the saturation
tables that govern how the various fluids flow in the model relative to
the other phases. Again, the saturation tables are fluid type dependent.

Various other fluid and rock property data are also entered in this
section including rock compressibility via the ROCK keyword for the
standard rock compressibility model, and rather unusually, the
Carter-Tracy influence function entered via the AQUTAB keyword.

\section{Keyword Definitions}\label{keyword-definitions-9}

\section{Data Requirements}\label{data-requirements-4}

For all phases activated in the model a complete description of the PVT
behavior (PVT tables) and fluid flow behavior (saturations tables) is
required. As the data is fluid type dependent, the next two sections
attempt to define the appropriate keywords for the various activated
fluid types. The first section outlines the fluid PVT property
requirements and the second section summarizes the keywords associated
with fluid flow property data, namely the relative permeability and
capillary pressure data.

\subsection{Fluid Properties}\label{fluid-properties}

\pandocbounded{\includegraphics[keepaspectratio]{chapters/images/Frame33_3c9982fcb963.png}}
\pandocbounded{\includegraphics[keepaspectratio]{chapters/images/graphics1_3c9982fcb963.png}}

A phase is a portion of a system which is (1) homogeneous in
composition, (2) bounded by a physical surface, and (3) mechanically
separable from other phases which may be present. The most important
phases which occur in petroleum production are liquid phases (oil,
distillate, condensate, and water) and a vapor or gas phase (natural
gas) component. Separators are used to separate the gas, water, and
hydrocarbon liquid phases produced from wells. Solid phases also occur
in petroleum production, mainly as:

\begin{enumerate}
\def\labelenumi{(\arabic{enumi})}
\item
  paraffin or other deposits within the producing formation, in the
  tubing, and in the surface equipment;
\item
  gas hydrates which may freeze and clog gas-flow lines and equipment;
  and
\item
  sand or reservoir rock {[}Bradley, H. B.: ``Petroleum Engineering
  Handbook'', Society of Petroleum Engineers, Richardson, Texas,
  USA.{]}.
\end{enumerate}

The phase behavior of petroleum systems is represented by a
Pressure-Volume-Temperature (``PVT'') phase diagram, as illustrated in
Figure 8.1, which shows a typical phase diagram after Clark {[}Clark,
Norman, ``It Pays to Know Your Petroleum'', World Oil, March and April,
1953.{]}.

Several important definitions should be noted from Figure 8.1 that
describes how a fluid behaves with respect to the composition, pressure,
and temperature, namely:

\begin{itemize}
\tightlist
\item
  Bubble Point: Pressure at which the first bubble of gas is formed.
\item
  Dew Point: Pressure at which the last drop of liquid turns to gas.
\item
  Critical Point: Point at which all properties of liquid and gas are
  identical, the point where the dew point and bubble point curves join.
\item
  Cricondentherm: Temperature above which two phases cannot co-exist
  regardless of pressure.
\item
  Cricondenbar: Pressure above which two phases cannot co-exist
  regardless of the temperature
\end{itemize}

Multi-component hydrocarbon systems are most accurately defined by their
compositions, which are usually expressed as mole fractions of the
components methane, ethane, propane, isobutane, n-butane, iso-pentane,
n-pentane, hexanes, and heptanes and heavier (heptanes plus). The
heptanes-plus fraction is usually defined by its average molecular
weight, its specific gravity, and/or the fractions which boil off in
selected temperature ranges. In addition to hydrocarbons, many naturally
occurring petroleum systems also contain hydrogen sulphide, nitrogen,
carbon dioxide, and helium in quantities which are sufficiently large to
appreciably affect the phase behavior of the system (see GPA Standard
2145 {[}Gas Processors Association, GPA, ``Standard 2145 - Physical
Constants of Paraffin Hydrocarbons and Other Components of Natural
gas'', 1986.{]} and Standing (1952) {[}Standing, M. B.:'' Volumetric and
Phase Behaviour of Oil Field Hydrocarbon Systems'', Renihold Publishing
Corp., New York City (1952).{]} for additional information).

\subsubsection{Gas Fluid Properties}\label{gas-fluid-properties}

The Ideal Gas Law is based on the theory of gases, a mathematical
equation called an Equation-of-State that can be derived to express the
relationship existing between pressure, volume, and temperature, for a
given quantity of gas. This relationship is called the Ideal Gas Law,
and is expresses mathematically as:

\begin{equation}\phantomsection\label{eq-8-1}{
\mathit{PV} = \mathit{nRT}
}\end{equation}

where

P = pressure (psia or Pa),

T = temperature (oR = 459.67 + oF or K = 273.15 + oC),

V = volume of the gas (ft3 or m3),

n = number of moles of gas (lb-mole or g-mol), and

R = universal gas constant which has the value of approximately

10.732 psi.ft3/lb-mole/oR or 8.314 Pa.m3/g-mol/K.

However, the above equation of state is not applicable for Real Gases
observed in petroleum reservoirs. In dealing with gases at a very low
pressure, the ideal gas relationship (equation of state) is a convenient
and accurate, within 2 - 3\% at atmospheric conditions, while for
pressures and temperatures found in petroleum reservoirs the physical
properties calculated can lead to errors in excess of 500\%. This is
because real gases (natural gases) do not behave as an ideal gas.
Basically the magnitude of deviations of real gases from the conditions
of the ideal gas law increases with increasing pressure and temperature,
and varies widely with the composition of the gas. The reason for this
is that the perfect gas law was derived under the assumption that the
volume of molecules is insignificant and no molecular attraction or
repulsion takes place. This is not the case for real gases.

In order to express a more exact relationship between the variables P,
V, and T, a correction factor called the gas compressibility factor (the
gas deviation factor, or simply the Z-factor) must be introduced into
the ideal gas law, that is:

\begin{equation}\phantomsection\label{eq-8-2}{
\mathit{PV} = \mathit{ZnRT}
}\end{equation}

where

P = pressure (psia or Pa),

T = temperature (oR = 459.67 + oF or K = 273.15 + oC),

V = volume of the gas (ft3 or m3),

n = number of moles of gas (lb-mole or g-mol),

R = universal gas constant which has the value of approximately

10.732 psi.ft3/lb-mole/oR or 8.314 Pa.m3/g-mol/K, and

Z = compressibility factor for the real gas.

For an ideal gas, the gas compressibility factor is equal to one, while
for a real gas, the Z-factor is greater or less than one depending on
the pressure, temperature, and composition of the gas. The Z-factor can
be calculated from the composition of the gas or from various
correlations including Standing {[}Standing, M.B. and Katz, D.L.:
``Density of Natural Gases,'' Trans., AIME (1942) 146, 140-44.{]} and
Dranchuk, P. M., Purvis, R. A., and Robinson {[}Dranchuk, P. M., Purvis,
R. A., and Robinson, D. B., ``Computer Calculations of Natural Gas
Compressibility Factors Using the Standing and Katz Correlation'',
Institute of Petroleum Technical Series, No.~IP 74-008, 1974.{]}.

Gas Density (ρg): Depends on composition, temperature, and pressure at
which the gas is measured, it is defined as the ratio of mass (m) to
volume (V), that is:

\begin{equation}\phantomsection\label{eq-8-3}{
{ρ}_{g} = \frac{m}{V} = \frac{P \mathit{MW}}{\mathit{ZRT}}
}\end{equation}

where \hyperref[REF_HEADING_KEYWORD_MW]{MW} is the molecular weight of
the gas and ρg has units of lbm/ft3 and kg/m3 for field and metric
units, respectively in equation (8.3). Gas density is often reported as
Relative Density, that is relative to air, as per 0.65 (air = 1).

Gas Formation Volume Factor (E): Is used to relate the volume of gas, as
measured at reservoir conditions, to the volume of gas as measured at
standard conditions (60 oF and 14.7 psia, or 15 oC and 101.325 kPa).
This gas property is then defined as the actual volume occupied by a
certain amount of gas at a specified pressure and temperature, divided
by the same amount of gas at standard conditions. Thus the gas formation
volume factor can be expressed as:

\begin{equation}\phantomsection\label{eq-8-4}{
E = \frac{{V}_{\mathit{sc}}}{{V}_{i}}
}\end{equation}

and substituting equation (8.3) into (8.4) we can derive the gas
formation volume factor for a given reservoir pressure and temperature:

\begin{equation}\phantomsection\label{eq-8-5}{
E = \left(\frac{{P}_{i}}{{P}_{\mathit{sc}}}\right) \left(\frac{{T}_{\mathit{sc}}}{{T}_{i}}\right) \left(\frac{1}{{Z}_{i}}\right)
}\end{equation}

\begin{tcolorbox}[enhanced jigsaw, title=\textcolor{quarto-callout-note-color}{\faInfo}\hspace{0.5em}{Note}, coltitle=black, breakable, toprule=.15mm, colbacktitle=quarto-callout-note-color!10!white, leftrule=.75mm, opacitybacktitle=0.6, colback=white, bottomtitle=1mm, toptitle=1mm, titlerule=0mm, left=2mm, colframe=quarto-callout-note-color-frame, arc=.35mm, rightrule=.15mm, bottomrule=.15mm, opacityback=0]

In oilfield units when we are dealing with oil reservoirs the gas
formation volume factor is more commonly denoted as Bg and has the units
of rb/scf. Whereas for gas reservoirs, the gas formation volume factor
is Eg and has the units scf/rcf and is often described as the Gas
Expansion Factor or GEF.

\end{tcolorbox}

Gas Isothermal Compressibility: A knowledge of the variability of a
fluid compressibility with pressure and temperature is essential in
performing many reservoir engineering calculations. For a liquid phase,
the compressibility is small and is usually assumed to be constant.
Whereas for a gas phase, the compressibility is neither small nor
constant. By definition the isothermal gas compressibility, cg, of a
substance is the change in volume per unit volume, for a unit change in
pressure, i.e.

\begin{equation}\phantomsection\label{eq-8-6}{
{c}_{g} = -\frac{1}{V} {\left(\frac{∂V}{∂P}\right)}_{T}
}\end{equation}

Most simulators calculate the isothermal gas compressibility based on
the entered PVT data, so it is not necessary to enter this variable
directly. If needed the correlation of Meehan, D. N., and Lyons, W. K.
{[}Meehan, D. N., and Lyons, W. K., ``Programmable Calculations for Gas
Compressibility'', Oil and Gas Journal, Oct.~8, 1979.{]} may be used to
estimate the values.

Gas Viscosity (μg): The gas viscosity of a fluid is a measure of the
internal fluid friction (resistance) to flow. If the friction between
layers of the fluid is small, i.e., low viscosity, an applied shearing
force will result in a large velocity gradient. As viscosity increases,
each fluid layer exerts a larger frictional drag on the adjacent layers
and velocity gradient increases. The viscosity of a fluid is generally
defined as the ratio of the shear force per unit area to the local
velocity gradient. Note that gas viscosity is not normally measured in
the laboratory because it can be estimated precisely from empirical
correlations. Gas viscosity increases with increasing temperature and
pressure. Usually it is not measured but is obtained from the Carr,
Kobayashi and Burrows {[}Carr, N.L., Kobayshi, R., and Burrows, D.B.;
``Viscosity of Natural Cases Under Pressure'', Trans. A.I.M.E. (1954),
201 pp.~264-272.{]} correlation, which include corrections for H2S, CO2,
and N2. For sour gases, this correlation is preferred to the Lee,
Gonzalez and Eakin {[}Lee, A.L., Gonzalez M.H. and Eakin B.E.: ``The
Viscosity of Natural Gases'', J.Pet.Tech. (1966) 18, pp.~997-1000.{]}
formulation (which does not account for H2S, CO2 and N2). Typically, gas
viscosity is in the range of 0.015 to 0.03 cp.

Condensate-Gas Ratio (``CGR''): Defines the volume of hydrocarbon liquid
that will dissolve in a unit volume of gas at a certain pressure and
temperature and has units of stb/MMscf and sm3/sm3 for field and metric
units, respectively. Field units values range from 1 or 2 stb/MMscf for
dry gas reservoirs, up to 100 stb/MMscf for wet gas reservoirs. The
empirical correlation of Meethan Vogel {[}Meehan, D. N., and Vogel, E.
L, HP-41 Reservoir Engineering Manual, PennWell Books, 1982.{]} can be
used to estimate CGR (Rv).

Water Gas Ratio (``WGR''): Gas in the reservoir is usually saturated
with water vapor, but as the gas moves up the wellbore, the temperature
and pressure decrease, and water condenses out of the gas. Often, water
that is produced at the surface is not formation water but is water of
condensation. Confirmation of this is a low water-gas ratio (1 -- 5
bbl/MMscf) and low water salinity (less than 100 ppm NaCl). The water
vapor content of a gas depends on pressure, temperature, and composition
of the gas. As pressure decreases, the water vapor content increases.
Conversely, as temperature decreases, the water vapor content decreases.
Units are stb/Mscf (although often quoted as stb/MMscf) and sm3/sm3 for
field and metric units, respectively. Bubacek's {[}Bukacek, R.F.:
``Equilibrium Moisture Content of Natural Gases,'' Bull., Inst. of Gas
Technology Bulletin (1955).{]} empirical correlation for sweet natural
gas (methane) may be used for gases without H2S to calculate WGR or
McKetta and Wehe's {[}McKetta, J.J. and Wehe, A.H.: ``Use This Chart for
Water Vapor Content of Natural Gases,'' Petroleum Refiner (Aug.~1958)
153-54.{]} correlation that takes into account water salinity and gas
gravity. Finally Robinson et al {[}Robinson, J. N., et al., ``Estimation
of the Water Content of Sour Natural Gases'', Paper Number SPE 6098,
51st Annual Fall Technical Conference, SPE, New Orleans, Oct.~3--6,
1976.{]} and Wichert and Wichert's {[}Wichert, G. C. and Wichert, E.,
``Chart Estimates Water Content of Sour Natural Gas'', O\&GJ, March 29,
1993, pp.~61-64.{]} correlations may be used for sour natural gases
(those containing H2S).

\subsubsection{Oil Fluid Properties}\label{oil-fluid-properties}

For oil and condensate fluids we have a similar set of properties;
however, oil properties are nearly always measured experimentally in the
laboratory, as although there are many correlations available, they
generally have low predictability unless tuned with the laboratory
measured data. Thus, the main oil properties are:

Oil Density (ρo): Like gas density, oil density depends on composition,
temperature, and pressure at which the fluid is measured, it is defined
as the ratio of mass (m) to volume (V), that is:

\begin{equation}\phantomsection\label{eq-8-7}{
{ρ}_{o} = \frac{m}{V}
}\end{equation}

with units of lbm/ft3, kg/m3, and g/cm3. Note that density is quoted in
various units even within a given unit system. In oilfield (English)
units it is quite common to quote density as a gradient that is psi/ft,
this is because fluid gradients are a common form of in-situ measurement
that can be easily related to density by:

\begin{equation}\phantomsection\label{eq-8-8}{
{ρ}_{o} = 144 x \mathit{fluid} \mathit{gradient} (\text{psi/ft})
}\end{equation}

For example, for an oil density measured as 0.338 psi/ft, the actual
in-situ density would simply be:

\begin{equation}\phantomsection\label{eq-8-9}{
{ρ}_{o} = 144 x 0.338 = 48.7 \frac{\mathit{lbm}}{{\mathit{ft}}^{3}}
}\end{equation}

Oil density can range from as low as about 0.270 psi/ft (0.62 g/cc) in
the Malay Basin (offshore Malaysia) to as high as 0.429 psi/ft (0.99
g/cc) in the heavy oil deposits in Cold Lake (Canada). However, most
common oil densities range from around 0.78 - 0.85 g/cc (0.338 - 0.369
psi/ft). When using oil field units most engineers drop the m in the lbm
units term, and therefore quote density as lb/ft3 (or lb/cu. ft.), as
lbm is implied.

Oil Relative Density (γo) defines the ratio of the density of a known
material to the density of reference material, at standard conditions of
pressure and temperature. Standard conditions vary throughout the world,
but for oil field units one normally uses14.7 psia and 60 oF, while for
SI units some areas use 101.325 kPa and 15 oC . The reference material
for liquids is always water, and thus the Oil Relative Density is
defined as:

\begin{equation}\phantomsection\label{eq-8-10}{
{γ}_{o} = \frac{{ρ}_{o}}{{ρ}_{w}}
}\end{equation}

Since the density of water at 14.7 psia and 60 oF is 62.4 lb/ft3, then
we have:

\begin{equation}\phantomsection\label{eq-8-11}{
{γ}_{o} = \frac{{ρ}_{o}}{62.4}
}\end{equation}

For SI units, with density measured in kg/m3, relative density is:

\begin{equation}\phantomsection\label{eq-8-12}{
{γ}_{o} = \frac{{ρ}_{o}}{1000 }
}\end{equation}

Commonly relative density is quoted with reference to the reference
phase, for example, γo = 0.780 (water = 1).

The American Petroleum Institute (``API'') classifies oils based on an
API gravity (γAPI), or degrees API (oAPI), the relationship between
relative density (γo) of oil and API gravity (γAPI) is given by:

\begin{equation}\phantomsection\label{eq-8-13}{
{γ}_{\mathit{API}} = \frac{141.5}{{γ}_{o}} - 131.5
}\end{equation}

where

γo = oil gravity (water = 1.0),

ρo = oil density (lb/ft3 or kg/m3), and

γAPI = API gravity (oAPI).

API density is commonly reported at the well site.

Bubble Point (Saturation Pressure): The bubble point pressure, or
saturation pressure (Pb or Psat), is defined as the pressure at which
the first bubbles of gas appear from a liquid as the pressure declines.
The correlations given in the oil formation volume section (later on),
also have the equivalent bubble point pressure correlations. It should
be noted that the comments addressing the application of these
correlations in the oil formation volume factor section, are equally
applicable here.

Gas-Oil Ratio (``GOR''): The GOR, (Rs) is defined as the number of
standard cubic meters of gas that will dissolve in one stock tank meter
cubed of crude oil at a certain pressure and temperature. The solubility
of a natural gas in a crude oil is a strong function of the pressure,
the temperature, the API gravity, and gas gravity. For a particular gas
and crude oil to exist at a constant temperature, the solubility
increases with pressure until the saturation pressure is reached. At the
saturation pressure (bubble point pressure) all the available gases are
dissolved in the oil, and the gas solubility reaches its maximum value.
Rather than measuring the amount of gas that will dissolve in given
stock tank crude oil as the pressure is increased, it is customary to
determine the amount of gas that will come out of the sample of
reservoir crude oil.

Oil Formation Volume Factor (Bo): The oil formation factor, Bo, is
defined as the ratio of the volume of oil (plus the gas in solution) at
the prevailing reservoir pressure and temperature, to the volume of oil
at standard conditions. Oil formation volume factor is generally greater
than or equal to one. Mathematically we can express the oil formation
volume factor as:

\begin{equation}\phantomsection\label{eq-8-14}{
{B}_{o} = \frac{{V}_{\mathit{rc}}}{{V}_{\mathit{sc}}}
}\end{equation}

where

Bo = oil formation volume factor (rb/stb or m3/Sm3),

Vrc = volume of oil at reservoir conditions (rb or m3), and

Vsc = volume of oil at standard conditions (stb or Sm3).

Values of oil FVF at reservoir temperature and various reservoir
pressures can be obtained from a standard PVT analysis of a reservoir
fluid sample. However, if this data is unavailable, the
geologist/engineer must then resort to empirical correlations, such as
Standing {[}Standing, M. B.: ``A Pressure-Volume-Temperature Correlation
for Mixtures of California Oils and Gases'', Drill, and Prod. Prac., API
(1947).{]}, Vasquez and Beggs {[}Vasquez M. and Beggs H. D.:
``Correlations for Fluid Physical Property Prediction'', J. Pet. Tech.
(June 1980).{]}, or Glasø {[}Glasø, O.: ``Generalized
Pressure-Volume-Temperature Correlations'', J.Pet.Tech. (May 1980).{]}.
Generally Bo ranges from a low of 1.05 for heavy oils to a high of 2.5
rb/stb for volatile oils. Care should be exercised when applying any of
the aforementioned correlations, as there is significant variation
between the various correlations. In general, a given correlation may be
more appropriate for a geological basin based on comparing actual
measured PVT data from other fields with the correlation predicted
results.

Oil Isothermal Compressibility (co): The oil isothermal compressibility
is defined as the rate of change in volume with pressure increase per
unit volume of liquid, all variables other than pressure being constant.
Thus, by definition, isothermal compressibility of a substance is
defined mathematically by the following expression:

\begin{equation}\phantomsection\label{eq-8-15}{
{c}_{o} = - {\frac{1}{V} \left(\frac{∂V}{∂P}\right)}_{T}
}\end{equation}

where

co = oil isothermal compressibility (1/kPa),

V = oil volume (m3), and

P = pressure (kPa).

In terms of the oil formation volume factor, we can derive:

\begin{equation}\phantomsection\label{eq-8-16}{
{c}_{o} = - \frac{1}{{B}_{o}} \frac{{\mathit{dB}}_{o}}{\mathit{dP}}
}\end{equation}

where

co = oil compressibility (1/kPa), and

Bo = oil formation volume factor (m3/Sm3).

Generally, the isothermal compressibility is determined from a
laboratory PVT study. In cases where a study has not been performed, the
correlations of Vasquez and Beggs {[}Vasquez M. and Beggs H. D.:
``Correlations for Fluid Physical Property Prediction'', J. Pet. Tech.
(June 1980).{]} or Ramey {[}Ramey, H. J.: ``Rapid Methods for Estimating
Reservoir Compressibilities'', J.Pet.Tech. (April 1964).{]} may be used.

Oil Viscosity (μo): Crude oil viscosity is an important physical
property that controls and influences the flow of oil through porous
media and pipes. The viscosity, in general, is defined as the internal
resistance of the fluid to flow, and oil viscosity is a strong function
of the temperature, pressure, oil gravity, gas gravity, and gas
solubility. Oil viscosity can be classified into three main groups,
based on the nature or state of the crude with respect to pressure,
these are:

\begin{longtable}[]{@{}
  >{\raggedright\arraybackslash}p{(\linewidth - 2\tabcolsep) * \real{0.5000}}
  >{\raggedright\arraybackslash}p{(\linewidth - 2\tabcolsep) * \real{0.5000}}@{}}
\toprule\noalign{}
\begin{minipage}[b]{\linewidth}\raggedright
\end{minipage} & \begin{minipage}[b]{\linewidth}\raggedright
The dead oil viscosity is defined as the viscosity of the oil at
atmospheric pressure and temperature conditions (no gas in solution).
\end{minipage} \\
\midrule\noalign{}
\endhead
\bottomrule\noalign{}
\endlastfoot
& The saturated (bubble point) oil viscosity is defined as the viscosity
of oil at the bubble point pressure and reservoir temperature. \\
& The undersaturated oil viscosity is defined as the viscosity of the
oil at a pressure above the bubble point and at reservoir
temperature. \\
\end{longtable}

Whenever possible, oil viscosity should be determined in the laboratory
at reservoir conditions, and is normally included in a standard PVT
analysis of the crude. However, if no data is available, the
correlations of Beggs and Robinson {[}Beggs, H. D. and Robinson J. R.:
``Estimating the Viscosity of Crude Oil Systems'', J. Pet. Tech. (Sept.,
1975).{]} and Andrade-Guzman-Reynolds {[}Andrade, E. N. da C., Nature,
125 (1930), 309.{]} may be employed.

\subsubsection{Water Fluid Properties}\label{water-fluid-properties}

Water fluid properties are similar to the oil fluid properties, and
include:

Water Density (ρw): Water density has the same definition as oil
density, equation (8.7), with units of lbm/ft3, kg/m3, and g/cm3. Fresh
water density at 14.7 psia and 60 oF is 62.4 lb/ft3, and for SI units at
101. 3 kPa and 15 oC, the density is 1000 kg/m3. Note that density is
quoted in various units even within a given unit system. In oilfield
(English) units it is quite common to quote density as a gradient that
is psi/ft, when laboratory data or actual water samples are unavailable.
The density of formation water at reservoir conditions can be estimated
roughly (usually to within +/−10\%) from correlations such as McCain's
{[}McCain, W.D. Jr.: McCain, W.D. Jr.~1990. The Properties of Petroleum
Fluids, second edition. Tulsa, Oklahoma: PennWell Books.{]} and {[}Cain
Jr., W.D. 1991. Reservoir-Fluid Property Correlations-State of the Art
(includes associated papers 23583 and 23594 ). SPE Res Eng 6 (2):
266-272. SPE-18571-PA. http://dx.doi.org/10.2118/18571-PA.{]}
correlations.

Water Formation Volume Factor (Bw): The water formation volume factor
has the same definition as for the oil formation volume factor. Normally
this property is estimated from correlations based on salt content, for
example the Numbere et al {[}Numbere, D., Brigham, W. E., and Standing,
M. B., Correlations for Physical Properties of Petroleum Reservoir
Brines, Petroleum Research Institute, Stanford University, November,
1977).{]} correlation.

Water Viscosity (μw): Again, the definition is the same as for oil and
is normally estimated from correlations based on salt content, for
example by the Meehan {[}Meehan, D. N., ``Estimating Water Viscosity at
Reservoir Conditions'', Petroleum Engineer, July 1980.{]}.

\subsection{Classification of Reservoirs by Fluid
Type}\label{classification-of-reservoirs-by-fluid-type}

Reservoir classification by fluid type includes the following reservoir
fluids:

\begin{itemize}
\tightlist
\item
  Low Shrinkage Oil Reservoirs.
\item
  High Shrinkage Oil Reservoirs.
\item
  Volatile Oil Reservoirs.
\item
  Dry Gas Reservoirs.
\item
  Wet Gas Reservoirs.
\item
  Retrograde Condensate Gas Reservoirs.
\end{itemize}

The following sections briefly outline the various reservoir types based
on their fluid properties.

\subsubsection{Low Shrinkage Oil
Reservoirs}\label{low-shrinkage-oil-reservoirs}

The phase diagram for a typical low shrinkage oil is shown in Figure 8.2
(after Clark {[}Clark, Norman, ``It Pays to Know Your Petroleum'', World
Oil, March and April, 1953.{]}), in this case the two-phase region
covers a wide range of pressure and temperature (the critical
temperature is above the reservoir temperature). The line marked 1 and 2
shows the effect of pressure reduction, with constant reservoir
temperature, as the reservoir is produced. The dashed line shows the
conditions encountered as the fluid travels up the tubing string and
into the surface facilities.

\pandocbounded{\includegraphics[keepaspectratio]{chapters/images/Frame34_703d61fd0177.png}}
\pandocbounded{\includegraphics[keepaspectratio]{chapters/images/graphics2_703d61fd0177.png}}

If the reservoir pressure and temperature is at point 2 on the phase
diagram, then the oil is at bubble point pressure and is said to be
saturated, i.e., the oil contains the maximum amount of dissolved gas at
the given conditions. However, if the reservoir conditions are at point
1 on the phase diagram, then the oil is said to be undersaturated. That
is the oil is able to dissolve more gas at the current reservoir
pressure and temperature.

The term low shrinkage oil is derived from the fact that oil shrinkage
is small, i.e., with an oil formation volume factors less than 1.5. This
oil can be described as having a broad based phase envelope, high
percentage of liquid, high proportion of heavier hydrocarbons, GOR's
less than 500 scf/stb or 100 Sm3/m3, oil gravity 30 oAPI or heavier, and
the stock tank liquid being black or a deep color. If the oil has a very
low GOR and therefore has an oil formation volumes factor close to one,
as observed in heavy oil reservoirs, then it is common to model this
type of reservoir neglecting the gas phase. That is this type of
reservoir fluid is typically modeled as a two-phase oil-water system
using a black-oil formulation, which results in greater computational
efficient due to having only two phases as opposed to the three phases
in an oil-gas-water system.

\subsubsection{High Shrinkage Oil
Reservoirs}\label{high-shrinkage-oil-reservoirs}

A high shrinkage oil's phase diagram is shown in Figure 8.3.

\pandocbounded{\includegraphics[keepaspectratio]{chapters/images/Frame35_1085786835ce.png}}
\pandocbounded{\includegraphics[keepaspectratio]{chapters/images/graphics3_1085786835ce.png}}

This type of reservoir fluid is characterized by not so broad a phase
envelope, fewer heavier hydrocarbons, deep colored stock tank fluid, oil
gravity lighter than 50 oAPI, GOR less than 8000 scf/stb or 1500 Sm3/m3,
and oil formation volume factors greater than 1.5 Sm3/m3.

\subsubsection{Dry Gas Reservoirs}\label{dry-gas-reservoirs}

Natural gas that occurs in the absence of condensate or liquid
hydrocarbons, or gas that had condensable hydrocarbons removed, is
called dry gas. It is primarily methane with some intermediates. The
hydrocarbon mixture is solely gas in the reservoir and there is no
liquid (condensate surface liquid) formed either in the reservoir or at
surface.

The phase diagram for a dry gas reservoir is illustrated in Figure 8.4.

\pandocbounded{\includegraphics[keepaspectratio]{chapters/images/Frame36_03ba828996ce.png}}
\pandocbounded{\includegraphics[keepaspectratio]{chapters/images/graphics4_03ba828996ce.png}}

Notice that the separator conditions lie outside the phase envelope in
the gas region. The term dry indicates that the gas does not contain
heavier hydrocarbons to form liquids at the surface conditions. Dry gas
typically has a GOR greater than 100,000 scf/stb or 18,000 Sm3/m3.

These type of reservoirs are typically modeled as two-phase gas-water
systems in black-oil formulations, which results in greater
computational efficient due to having only two phases as oppose to three
phases in an oil-gas-water system. Also it is not uncommon to model dry
gas reservoirs using analytical models, for example material balance, as
often the reservoirs behave like ``tanks'' either with or without
aquifer influx.

\subsubsection{Wet Gas Reservoirs}\label{wet-gas-reservoirs}

Natural gas that contains significant heavy hydrocarbons such as
propane, butane, and other liquid hydrocarbons is known as wet gas or
rich gas. The general rule of thumb is if the gas contains less methane
(typically less than 85\% methane) and more ethane, and other more
complex hydrocarbons, it is labelled as wet gas.

As can be seen from the figure below, wet gas exists in the gas phase in
the reservoir throughout depletion of the reservoir pressure. However,
the separator conditions lie in the two-phase region, resulting in
liquid production (condensate).

\pandocbounded{\includegraphics[keepaspectratio]{chapters/images/Frame37_4031757f087d.png}}
\pandocbounded{\includegraphics[keepaspectratio]{chapters/images/graphics5_4031757f087d.png}}

Wet gas normally has a GOR less than 100,000 scf/stb or 18,000 Sm3/m3,
with the condensate having a gravity greater than 50 oAPI.

\subsubsection{Retrograde Condensate Gas
Reservoirs}\label{retrograde-condensate-gas-reservoirs}

If the reservoir temperature lies between the critical point and the
cricondentherm, a retrograde reservoir exists. As the pressure declines
to point 2, which is the dew point, liquids begin to form in the
reservoir. As the pressure further declines to point 3, the liquid yield
increases. Further reduction in the pressure causes the liquid to
vaporize back into the gas phase.

\pandocbounded{\includegraphics[keepaspectratio]{chapters/images/Frame38_fcdacd8d01d2.png}}
\pandocbounded{\includegraphics[keepaspectratio]{chapters/images/graphics6_fcdacd8d01d2.png}}

The liquid drop out in the reservoir can cause a well's productivity to
be reduced, due to the system moving from a two-phase relative
permeability effect, to a three phase effect. This is normally diagnosed
by the well having a high skin factor.

This type of fluid has more lighter hydrocarbons than high shrinkage
oils, and fewer heavier hydrocarbons than high shrinkage oils. The
gravity can be as high as 60 oAPI, with GOR up to 70,000 scf/stb or
13,000 Sm3/m3. Generally the stock tank liquid is water-white, or
slightly colored.

\subsubsection{Reservoir Classification}\label{reservoir-classification}

The following tables outline some general parameters used to classify
reservoir fluids. Table 8.1 outlines some general producing
characteristics used to define reservoir fluids.

\begin{longtable}[]{@{}
  >{\raggedright\arraybackslash}p{(\linewidth - 8\tabcolsep) * \real{0.2000}}
  >{\raggedright\arraybackslash}p{(\linewidth - 8\tabcolsep) * \real{0.2000}}
  >{\raggedright\arraybackslash}p{(\linewidth - 8\tabcolsep) * \real{0.2000}}
  >{\raggedright\arraybackslash}p{(\linewidth - 8\tabcolsep) * \real{0.2000}}
  >{\raggedright\arraybackslash}p{(\linewidth - 8\tabcolsep) * \real{0.2000}}@{}}
\toprule\noalign{}
\begin{minipage}[b]{\linewidth}\raggedright
Item
\end{minipage} & \begin{minipage}[b]{\linewidth}\raggedright
Crude Oil
\end{minipage} & \begin{minipage}[b]{\linewidth}\raggedright
Volatile Oil
\end{minipage} & \begin{minipage}[b]{\linewidth}\raggedright
Gas Condensate
\end{minipage} & \begin{minipage}[b]{\linewidth}\raggedright
Gas
\end{minipage} \\
\midrule\noalign{}
\endhead
\bottomrule\noalign{}
\endlastfoot
Initial Separator Gas-Liquid Ratio (scf/stb) & \textless{} 2,000 & 2,000
-- 6,000 & 6,000 -- 100,000 Generally 7,000 -- 15,000 & \textgreater{}
100,000 Usually below 300,000 \\
Color of Produced Liquid & Black to Light Green & Dark Straw Colored &
Colorless to Light Straw Colored & Colorless \\
Usual API Gravity Range of Produced Fluid & 10 -- 45+ & 40 -- 50 & 45 --
65 & -- \\
Composition of Produced Fluid & C7+: usually \textgreater{} 40\% & C7+:
10 -- 40\% & C7+: 2 -- 10\% & Primarily Methane \\
Oil Formation Volume Factor (rb/stb) & \textless{} 2.0 frequently
\textless{} 1.5 & \textgreater{} 2.0 often \textgreater{} 2.5 & & \\
\end{longtable}

\emph{Table 8.1: General Reservoir Fluid Classification}

In terms of fluid composition reservoir fluids can be indicated by the
values shown in Table 8.2 below.

\begin{longtable}[]{@{}lllll@{}}
\toprule\noalign{}
Item & Gas & Gas Condensate & Volatile Oil & Crude Oil \\
\midrule\noalign{}
\endhead
\bottomrule\noalign{}
\endlastfoot
Methane (mole\%) & 96 & 87 & 64 & 49 \\
Ethane (mole \%) & 3 & 4 & 8 & 3 \\
Heptanes Plus (mole \%) & 1 & 3 & 15 & 42 \\
Solution GOR (scf/stb) & 105,000 & 18,200 & 2,000 & 625 \\
Stock Tank Oil oAPI & 68 & 61 & 50 & 34 \\
\end{longtable}

\emph{Table 8.2: Reservoir Fluids Composition Indicators}

In addition reservoir oil is further divided into Light, Medium, Heavy,
Extra Heavy, and Bituman (Wilmon {[}Wilmon, G. J. : ``Economic Outlook
for Extra-heavy Oils, Natural Bitumens and Shale Oils'', World Petroleum
Congress (1987).{]}) as shown in Table 8.3 below:

\begin{longtable}[]{@{}lll@{}}
\toprule\noalign{}
temperature & Stock Tank Oil oAPI & In Situ Viscosity (cp) \\
\midrule\noalign{}
\endhead
\bottomrule\noalign{}
\endlastfoot
Light & \textgreater{} 31.1 & \\
Medium & 22.3 -- 31.1 & \\
Heavy & 10.0 -- 22.3 & 100 -- 10,000 \\
Extra Heavy & \textless{} 10.0 & \textless{} 10,000 \\
Bitumen & \textless{} 10.0 & \textgreater{} 10,000 \\
\end{longtable}

\emph{Table 8.3: Oil Sub-Classification (after Wilmon)}

\subsection{Fluid Property Tables}\label{fluid-property-tables}

Table 8.4 outlines the oil, gas, and water fluid types that can be
active in the model, together with the related RUNSPEC section keywords
that activate the phases, versus the PVT keywords that can be used to
define the PVT behavior. The table also includes several water types,
the standard water phase used in most simulations models, the Brine
water phase used in the Brine Tracking model to track the flow of brine
through the simulation grid and the effect of brine on reservoir
performance, and finally OPM Flow's Vaporized Water phase that is used
in the simulator's Salt Precipitation model. The latter water phase is
not available in the commercial simulator.

\begin{longtable}[]{@{}l@{}}
\toprule\noalign{}
Fluid Property Keywords versus Oil, Gas \& Water Fluid Types \\
\midrule\noalign{}
\endhead
\bottomrule\noalign{}
\endlastfoot
\end{longtable}

\begin{longtable}[]{@{}
  >{\raggedright\arraybackslash}p{(\linewidth - 16\tabcolsep) * \real{0.1111}}
  >{\raggedright\arraybackslash}p{(\linewidth - 16\tabcolsep) * \real{0.1111}}
  >{\raggedright\arraybackslash}p{(\linewidth - 16\tabcolsep) * \real{0.1111}}
  >{\raggedright\arraybackslash}p{(\linewidth - 16\tabcolsep) * \real{0.1111}}
  >{\raggedright\arraybackslash}p{(\linewidth - 16\tabcolsep) * \real{0.1111}}
  >{\raggedright\arraybackslash}p{(\linewidth - 16\tabcolsep) * \real{0.1111}}
  >{\raggedright\arraybackslash}p{(\linewidth - 16\tabcolsep) * \real{0.1111}}
  >{\raggedright\arraybackslash}p{(\linewidth - 16\tabcolsep) * \real{0.1111}}
  >{\raggedright\arraybackslash}p{(\linewidth - 16\tabcolsep) * \real{0.1111}}@{}}
\toprule\noalign{}
\begin{minipage}[b]{\linewidth}\raggedright
Item
\end{minipage} & \begin{minipage}[b]{\linewidth}\raggedright
Oil
\end{minipage} & \begin{minipage}[b]{\linewidth}\raggedright
Gas
\end{minipage} & \begin{minipage}[b]{\linewidth}\raggedright
Water
\end{minipage} & \begin{minipage}[b]{\linewidth}\raggedright
\end{minipage} & \begin{minipage}[b]{\linewidth}\raggedright
\end{minipage} & \begin{minipage}[b]{\linewidth}\raggedright
\end{minipage} & \begin{minipage}[b]{\linewidth}\raggedright
\end{minipage} & \begin{minipage}[b]{\linewidth}\raggedright
\end{minipage} \\
\midrule\noalign{}
\endhead
\bottomrule\noalign{}
\endlastfoot
Fluid Type & Dead Oil & Live Oil & Dry Gas & Wet Gas & Dry / Wet Gas
with Vaporized Water & CO2 (EOR) & Water & Brine \\
RUNSPEC Keywords & OIL & OIL & GAS & GAS & GAS & GAS & WATER & WATER \\
& DISGAS & & VAPOIL & None / VAPOIL & & & & \\
& & & & VAPWAT5 & & & & \\
& & & & WATER & & & & \\
& & & & BRINE & & & BRINE & \\
Pressure Dependent PVT & PVCDO & PVCO & PVDG & PVTG & PVTWSALT & PVTSOL
& PVTW & PVTWSALT \\
PVDO & PVTO & PVZG & & PVTGW5 & PVDS & & & \\
& & & & PVTGWO5 & & & & \\
Miscellaneous & RSCONST & & RVCONST & & RWGSALT5 & & & \\
RSCONSTT & & RVCONSTT & & & & & & \\
Surface Density & DENSITY & BDENSITY & SDENSITY & DENSITY & BDENSITY & &
& \\
GRAVITY & & & GRAVITY & & & & & \\
\end{longtable}

\begin{longtable}[]{@{}l@{}}
\toprule\noalign{}
Notes: \\
\midrule\noalign{}
\endhead
\bottomrule\noalign{}
\endlastfoot
\end{longtable}

\emph{Table 8.4: Fluid Property Keywords versus Oil, Gas and Water Fluid
Type}

For the Dead Oil phases the RSCONST and RSCONSTT keywords are used to
set a constant gas-oil ratio (Rs). In this case the Rs is independent of
the reservoir pressure and Rs is also negligible, as in for example
heavy oil type fluids.

Similarly for the Dry Gas phase, where the RVCONST and RVCONSTT keywords
are used to set a condensate-gas ratio (Rv) which is independent of the
reservoir pressure and is also negligible, as in for example dry gas
type fluids.

For the Vaporized Water phase and model, both dry and wet gas can be
incorporated. Note that the PVTGW, PVTGWO, RWGSALT, and VAPWAT keywords
are OPM Flow specific keywords.

In addition for the Brine phase, then either the SALT or SALTVD keywords
in the SOLUTION section should be used to define the initial
equilibration salt concentration for the model.

CO2 can either be used as Enhanced Oil Recovery (``EOR'') fluid by
injecting the CO2 into an oil reservoir, or injected for CO2 storage.
For the former, CO2 is declared via the GAS keyword in the RUNSPEC
section and the PVT data is entered via the standard gas fluid
properties for the hydrocarbon gas and either the PVDS or the PVTSOL
keywords are used to describe the interaction of the in situ oil and the
injected CO2.

In addition to the above the ROCK keyword should be used to define the
rock compressibility.

Similarly, Table 8.5 outlines the fluid property data keywords for the
CO2 storage, foam, polymer, solvent, and MICP phases. Note that for
these phases multiple keywords can be used to define the desired
property behavior.

\begin{longtable}[]{@{}
  >{\raggedright\arraybackslash}p{(\linewidth - 10\tabcolsep) * \real{0.1667}}
  >{\raggedright\arraybackslash}p{(\linewidth - 10\tabcolsep) * \real{0.1667}}
  >{\raggedright\arraybackslash}p{(\linewidth - 10\tabcolsep) * \real{0.1667}}
  >{\raggedright\arraybackslash}p{(\linewidth - 10\tabcolsep) * \real{0.1667}}
  >{\raggedright\arraybackslash}p{(\linewidth - 10\tabcolsep) * \real{0.1667}}
  >{\raggedright\arraybackslash}p{(\linewidth - 10\tabcolsep) * \real{0.1667}}@{}}
\toprule\noalign{}
\begin{minipage}[b]{\linewidth}\raggedright
Fluid Property Keywords versus CO2 (Storage), Foam, Polymer, Solvent,
and MICP Fluid Types
\end{minipage} & \begin{minipage}[b]{\linewidth}\raggedright
\end{minipage} & \begin{minipage}[b]{\linewidth}\raggedright
\end{minipage} & \begin{minipage}[b]{\linewidth}\raggedright
\end{minipage} & \begin{minipage}[b]{\linewidth}\raggedright
\end{minipage} & \begin{minipage}[b]{\linewidth}\raggedright
\end{minipage} \\
\midrule\noalign{}
\endhead
\bottomrule\noalign{}
\endlastfoot
Fluid Type & CO2 (Storage) & Foam & Polymer & Solvent & MICP7 \\
RUNSPEC Keywords & CO2STORE & FOAM & POLYMER & SOLVENT & MICP \\
GAS & & & & WATER & \\
WATER & & & & & \\
Pressure Dependent PVT & Not required, calculated via correlations & & &
PVDS & PVTW \\
& & PVTSOL6 & & & \\
Surface Density & & & & SDENSITY & DENSITY \\
Miscellaneous & SALINITY5 & FOAMADS & PLMIXPAR & &
\hyperref[REF_HEADING_KEYWORD_BIOFPARA]{BIOFPARA} \\
& FOAMDCYO & PLYADS & & & \\
& FOAMDCYW & PLYADSS & & & \\
& FOAMFCN & PLYATEMP & & & \\
& FOAMFRM & PLYCAMAX & & & \\
& FOAMFSC & PLYDHFLF & & & \\
& FOAMFSO & PLYESAL & & & \\
& FOAMFST & PLYKRRF & & & \\
& FOAMFSW & PLYMAX & & & \\
& FOAMMOB & PLYMWINJ8 & & & \\
& FOAMMOBP & PLYRMDEN & & & \\
& FOAMMOBS & PLYROCK & & & \\
& FOAMOPTS & PLYSHEAR & & & \\
& FOAMROCK & PLYSHLOG & & & \\
& & PLYTRRF & & & \\
& & PLYTRRFA & & & \\
& & PLYVISC & & & \\
& & PLYVISCS & & & \\
& & PLYVISCT & & & \\
& & PLYVMH8 & & & \\
& & PLYVSCST & & & \\
& & SKPRPOLY8 & & & \\
Miscellaneous & & & SKPRWAT8 & & \\
Notes: & & & & & \\
\end{longtable}

\emph{Table 8.5: Fluid Property Keywords versus CO2 (Storage), Foam,
Polymer, Solvent, and MICP Fluid Types}

As mentioned previously, CO2 can either be used as an EOR fluid by
injecting the gas into an oil reservoir, or injected for CO2 storage.
For the latter, CO2 is declared via the GAS and CO2STORE keywords in the
RUNSPEC section, and, except for the SALINITY keyword, all the remaining
PVT data is calculated automatically via correlations by the simulator.
A full description of the underlying PVT models is described by Sandve
at al.~{[}Sandve, T. H., Gasda, S. E., Rasmussen, A., and Rustad, A. B.
Convective dissolution in field scale CO2 storage simulation using the
OPM Flow simulator. Submitted to TCCS 11 -- Trondheim Conference on CO2
Capture, Transport and Storage Trondheim, Norway -- June 21-23, 2021.{]}

Typical live oil and dry gas PVT data from the Volve {[}The Volve Data
was approved for data sharing in 2018 by the initiative of the last
Operating company, Equinor and approved by the license partners
ExxonMobil E\&P Norway AS and Bayerngas Norge AS in the end of 2017.{]}
field is shown in Figure 8.7 and Figure 8.8, respectively. To fully
define the live oil, dry gas, and water PVT properties for Volve the
DENSITY, PVDG, PVTO, and PVTW keywords are employed.

\pandocbounded{\includegraphics[keepaspectratio]{chapters/images/Frame13_64d35fcd8d22.png}}
\pandocbounded{\includegraphics[keepaspectratio]{chapters/images/Image11_64d35fcd8d22.png}}

\pandocbounded{\includegraphics[keepaspectratio]{chapters/images/Frame14_95472fdbfcec.png}}
\pandocbounded{\includegraphics[keepaspectratio]{chapters/images/Image12_95472fdbfcec.png}}

\subsection{Saturation Tables (Relative Permeability and Capillary
Pressure
Tables)}\label{saturation-tables-relative-permeability-and-capillary-pressure-tables}

Saturation tables contain the relative permeability and capillary
pressure data as a function of fluid saturation and are used to both
initialize the model and to describe multi-phase flow in the reservoir.
Multiple saturation tables can be entered and allocated to various areas
in the model, based on rock typing. Alternatively, a limited number of
saturation tables may be entered and allocated by region and combined
with end-point scaling option to enable a more robust reservoir rock
characterization.

\begin{tcolorbox}[enhanced jigsaw, title=\textcolor{quarto-callout-note-color}{\faInfo}\hspace{0.5em}{Note}, coltitle=black, breakable, toprule=.15mm, colbacktitle=quarto-callout-note-color!10!white, leftrule=.75mm, opacitybacktitle=0.6, colback=white, bottomtitle=1mm, toptitle=1mm, titlerule=0mm, left=2mm, colframe=quarto-callout-note-color-frame, arc=.35mm, rightrule=.15mm, bottomrule=.15mm, opacityback=0]

Confusion abounds with the terms: critical water saturation, connate
water saturation, and irreducible water saturation definitions, as the
terms are extensively used interchangeably to define the maximum water
saturation at which the water phase will remain immobile. For reference:
Connate Water Saturation: Water trapped in the interstices of the
sediments at the time of deposition, as opposed to water that migrated
into the formations after deposition
{[}https://wiki.seg.org/wiki/Dictionary:Connate\_water{]}. Normally the
nomenclature of Swc is used to describe this type of water saturation.
Irreducible Water Saturation: The fraction of the pore volume occupied
by water in a reservoir at maximum hydrocarbon saturation. It represents
water that has not been displaced by hydrocarbons because it is trapped
by adhering to rock surfaces, trapped in small pore spaces and narrow
interstices, etc. Irreducible water saturation is an equilibrium
situation. It differs from residual water saturation, the value measured
by core analysis, because of filtrate invasion and the gas expansion
that occurs when a core is removed from the bottom of the hole to the
surface. Also called immovable water
{[}https://wiki.seg.org/wiki/Dictionary:Irreducible\_water\_saturation{]}.
Here the nomenclature for irreducible water saturation is Swirr.
Critical Water Saturation: Defines the largest water saturation for
which the water relative permeability is zero.

\end{tcolorbox}

In addition to the above definitions, the simulator uses the term SWL,
which is lowest water saturation in a relative permeability table. Thus,
SWL can represent either the connate water saturation or the irreducible
water saturation in the relative permeability tables and functions.

Secondly, the term irreducible water saturation is replaced by term
critical water saturation (SWCR) in the relative permeability tables and
function, as both terms are equivalent with regards to the mobility of
the water phase. This terminology is also consistent with the critical
values of the oil and gas phases, which are dependent on the displacing
phase.

Thirdly, it is not uncommon for SWL to be set equal to the critical
water saturation; however, care should be taken in this instance if
end-point scaling is being used to ensure that the cell scaled relative
permeability curves are as one might expect.

A typical set of oil-water relative permeability curves is shown in
Figure 8.9 indicating the oil end-point data (KRO, KRORW, and (1 --
SOWCR)) and the water end-point data (KRWR, KRW, SWL, and SWCR).

\pandocbounded{\includegraphics[keepaspectratio]{chapters/images/Frame6_ef371648a600.png}}
\pandocbounded{\includegraphics[keepaspectratio]{chapters/images/Image7_ef371648a600.png}}

The associated oil-water end-point definitions are outlined in the
following table:

\begin{longtable}[]{@{}
  >{\raggedright\arraybackslash}p{(\linewidth - 4\tabcolsep) * \real{0.3333}}
  >{\raggedright\arraybackslash}p{(\linewidth - 4\tabcolsep) * \real{0.3333}}
  >{\raggedright\arraybackslash}p{(\linewidth - 4\tabcolsep) * \real{0.3333}}@{}}
\toprule\noalign{}
\begin{minipage}[b]{\linewidth}\raggedright
Type
\end{minipage} & \begin{minipage}[b]{\linewidth}\raggedright
End-Point Keyword
\end{minipage} & \begin{minipage}[b]{\linewidth}\raggedright
Oil-Water End-Point Definitions
\end{minipage} \\
\midrule\noalign{}
\endhead
\bottomrule\noalign{}
\endlastfoot
Saturation & SOWCR & Critical oil-in-water saturation, that is the
largest oil saturation for which the oil relative permeability is zero
in an oil-water system. \\
SWL & Connate water saturation, that is the lowest water saturation in a
water saturation function table. & \\
SWCR & Critical water saturation, that is the largest water saturation
for which the water relative permeability is zero. & \\
SWU & Maximum water saturation in a water saturation table. This is
commonly set equal to one, unless there is a residual oil saturation
below the defined OWC. & \\
Relative Permeability & KRO & Relative permeability of oil at the
maximum oil saturation. \\
KRORW & Relative permeability of oil at the critical water saturation
& \\
KRW & Relative permeability of water at the maximum water saturation
(normally the maximum water saturation is one). & \\
KRWR & Relative permeability of water at the residual oil saturation (or
at the residual gas saturation in a gas-water run). & \\
Capillary Pressure & SWLPC & Capillary pressure connate water
saturation, that is the lowest water saturation in a water saturation
function table to be scaled independently from the SWL relative
permeability end-point value. \\
\end{longtable}

\emph{Table 8.6: Oil-Water Relative Permeability End-Point Data
Definitions}

\pandocbounded{\includegraphics[keepaspectratio]{chapters/images/Frame5_be79ab2b1e5d.png}}
\pandocbounded{\includegraphics[keepaspectratio]{chapters/images/Image6_be79ab2b1e5d.png}}

Similarly for gas-oil systems, Figure 8.10 illustrates a typical gas-oil
relative permeability set of curves indicating the oil end-point data
(KRORG and (1 -- SOGCR)) and the gas end-point data (KRGR, KRG, and
SGCR).

The gas-oil end-point definitions are outlined in the following table:

\begin{longtable}[]{@{}
  >{\raggedright\arraybackslash}p{(\linewidth - 4\tabcolsep) * \real{0.3333}}
  >{\raggedright\arraybackslash}p{(\linewidth - 4\tabcolsep) * \real{0.3333}}
  >{\raggedright\arraybackslash}p{(\linewidth - 4\tabcolsep) * \real{0.3333}}@{}}
\toprule\noalign{}
\begin{minipage}[b]{\linewidth}\raggedright
Type
\end{minipage} & \begin{minipage}[b]{\linewidth}\raggedright
End-Point Keyword
\end{minipage} & \begin{minipage}[b]{\linewidth}\raggedright
Gas-Oil End-Point Definitions
\end{minipage} \\
\midrule\noalign{}
\endhead
\bottomrule\noalign{}
\endlastfoot
Saturation & SOGCR & Critical oil-in-gas saturation, that is the largest
oil saturation for which the oil relative permeability is zero in an
oil-gas-connate-water system. \\
SGL & Connate gas saturation, that is the lowest gas saturation in a gas
saturation function table. & \\
SGCR & Critical gas saturation, that is the largest gas saturation for
which the gas relative permeability is zero. & \\
SGU & Maximum gas saturation in a gas saturation table. This is
generally set equal to one minus the connate water saturation. & \\
Relative Permeability & KRG & Relative permeability of gas at the
maximum gas saturation. \\
KRGR & Relative permeability of gas at the residual oil saturation (or
the critical water saturation in a gas-water run). & \\
KRORG & Relative permeability of oil at the critical gas saturation.
& \\
Capillary Pressure & SGLPC & Capillary pressure connate gas saturation,
that is the smallest gas saturation in a gas saturation function table
to be scaled independently from the SGL relative permeability end-point
value. \\
\end{longtable}

\emph{Table 8.7: Gas-Oil Relative Permeability End-Point Data
Definitions}

Finally, for two phase gas-water systems, Figure 8.11 illustrates a
typical set of gas-water relative permeability curves indicating the gas
end-point data (KRG, KRGR, and (1 -- SGWCR)) and the water end-point
data (KRW, KRWR, SWL, and SWGCR).

\pandocbounded{\includegraphics[keepaspectratio]{chapters/images/Frame32_19b99162c724.png}}
\pandocbounded{\includegraphics[keepaspectratio]{chapters/images/Image28_19b99162c724.png}}

The gas-oil end-point definitions are outlined in the following table:

\begin{longtable}[]{@{}
  >{\raggedright\arraybackslash}p{(\linewidth - 4\tabcolsep) * \real{0.3333}}
  >{\raggedright\arraybackslash}p{(\linewidth - 4\tabcolsep) * \real{0.3333}}
  >{\raggedright\arraybackslash}p{(\linewidth - 4\tabcolsep) * \real{0.3333}}@{}}
\toprule\noalign{}
\begin{minipage}[b]{\linewidth}\raggedright
Type
\end{minipage} & \begin{minipage}[b]{\linewidth}\raggedright
End-Point Keyword
\end{minipage} & \begin{minipage}[b]{\linewidth}\raggedright
Gas-Water End-Point Definitions
\end{minipage} \\
\midrule\noalign{}
\endhead
\bottomrule\noalign{}
\endlastfoot
Saturation & SGL & Connate gas saturation, that is the lowest gas
saturation in a gas saturation function table. \\
SGU & Maximum gas saturation in a gas saturation table. & \\
SGWCR & Critical gas-in-water saturation, that is the largest gas
saturation for which the gas relative permeability is zero in a
gas-water system. & \\
SWL & Connate water saturation, that is the lowest water saturation in a
water saturation function table. & \\
SWGCR & Critical water saturation, that is the largest water saturation
for which the water relative permeability is zero in a gas-water system.
& \\
SWU & Maximum water saturation in a water saturation table. & \\
Relative Permeability & KRG & Relative permeability of gas at the
maximum gas saturation. \\
KRGR & Relative permeability of gas at the residual oil saturation (or
the critical water saturation in a gas-water run). & \\
KRWR & Relative permeability of water at the residual oil saturation (or
the residual gas saturation in a gas-water run). & \\
Capillary Pressure & SWLPC & Capillary pressure connate water
saturation, that is the smallest water saturation in a water saturation
function table to be scaled independently from the SWL relative
permeability end-point value. \\
\end{longtable}

\emph{Table 8.8: Gas-Water Relative Permeability End-Point Data
Definitions}

End-point scaling is activated in the RUNSPEC section with the ENDSCALE
keyword and the data used to apply end-point scaling is entered in the
PROPS section using the end-point keywords defined in Table 8.6, Table
8.7, and Table 8.8 to define each grid block's end-point data. There are
also direction dependent versions of these keywords for when directional
end-point scaling has been activated. For example for critical water
saturation, SWCR is used with non-directional end-point scaling and the
SWCRX, SWCRX-, SWCRY, SWCRY-, SWCRZ, and SWCRZ- series of keywords is
used for when directional end-point scaling has been activated. In
addition, there is also the facility to incorporate end-point scaling
based on the drainage and/or imbibition process which again can be
either non-directional or directional.

\begin{tcolorbox}[enhanced jigsaw, title=\textcolor{quarto-callout-note-color}{\faInfo}\hspace{0.5em}{Note}, coltitle=black, breakable, toprule=.15mm, colbacktitle=quarto-callout-note-color!10!white, leftrule=.75mm, opacitybacktitle=0.6, colback=white, bottomtitle=1mm, toptitle=1mm, titlerule=0mm, left=2mm, colframe=quarto-callout-note-color-frame, arc=.35mm, rightrule=.15mm, bottomrule=.15mm, opacityback=0]

If the hysteresis model option has been activated on the SATOPTS keyword
in the RUNSPEC section, then the standard end-point scaling arrays, for
example KRORW, can be used to scale the drainage tables, and the
equivalent imbibition arrays suffixed with the letter I, for example
IKRORW, can be used to scale the imbibition tables. However, if the
hysteresis model option has not been activated, then only the standard
keywords can be used to define the end-point scaling parameters, that is
KRORW should be used and not IKRORW, even though an imbibition process
with the wetting phase (normally water) increasing is usually being
modelled.

\end{tcolorbox}

Saturation functions can be entered via several keywords consisting of
two format types as depicted in the following table:

\begin{longtable}[]{@{}llllllll@{}}
\toprule\noalign{}
Format Type One & Format Type Two & & & & & & \\
\midrule\noalign{}
\endhead
\bottomrule\noalign{}
\endlastfoot
Keyword & Oil & Gas & Water & Keyword & Oil & Gas & Water \\
SGOF & Pcog & & SGFN3 & & Pcog & & \\
SLGOF & Pcog & & SGWFN & & Pcgw & & \\
SWOF & Pcwo & & Pcwo & SOF24 & No Pc & & \\
& & & & SOF35 & No Pc & & \\
& & & & SOF32D & No Pc & & \\
& & & & SWFN & & & Pcwo \\
OPM Flow LET6 & Format Type Three7 & & & & & & \\
Keyword & Oil & Gas & Water & Keyword & Oil & Gas & Water \\
SGOFLET & Pcog & & GSF & & Pcgw & & \\
SWGFLET & & Pcgw & WSF & & & No Pc & \\
SWOFLET & Pcwo & & Pcwo & & & & \\
Notes: & & & & & & & \\
\end{longtable}

\emph{Table 8.9: Saturation Table Formats and Phases}

Note that only one format type can be used in each run to define the
required saturation functions, that is one cannot combine the keywords
from the different format types in the same input deck.

Figure 8.12 shows an unscaled oil-water relative permeability curve from
the Volve model, and Figure 8.13 displays the corresponding scaled
relative permeability curves after applying the end-point scaling
parameters. Similarly for the gas-oil curves, Figure 8.14 illustrates
the unscaled curves and Figure 8.15 the scaled curves. The plots
indicate that the end-point scaling has been applied correctly.

\begin{longtable}[]{@{}
  >{\raggedright\arraybackslash}p{(\linewidth - 2\tabcolsep) * \real{0.5000}}
  >{\raggedright\arraybackslash}p{(\linewidth - 2\tabcolsep) * \real{0.5000}}@{}}
\toprule\noalign{}
\begin{minipage}[b]{\linewidth}\raggedright
\pandocbounded{\includegraphics[keepaspectratio]{chapters/images/Frame39_3d61873adcad.png}}
\pandocbounded{\includegraphics[keepaspectratio]{chapters/images/Image29_3d61873adcad.png}}
\end{minipage} & \begin{minipage}[b]{\linewidth}\raggedright
\pandocbounded{\includegraphics[keepaspectratio]{chapters/images/Frame49_73f6532732f0.png}}
\pandocbounded{\includegraphics[keepaspectratio]{chapters/images/Image30_73f6532732f0.png}}
\end{minipage} \\
\midrule\noalign{}
\endhead
\bottomrule\noalign{}
\endlastfoot
\pandocbounded{\includegraphics[keepaspectratio]{chapters/images/Frame50_010d8a902c9f.png}}
\pandocbounded{\includegraphics[keepaspectratio]{chapters/images/Image31_010d8a902c9f.png}}
&
\pandocbounded{\includegraphics[keepaspectratio]{chapters/images/Frame51_26d150ab4573.png}}
\pandocbounded{\includegraphics[keepaspectratio]{chapters/images/Image32_26d150ab4573.png}} \\
\end{longtable}

\subsection{Saturation Table Generation - Corey
Curves}\label{saturation-table-generation---corey-curves}

As mentioned in the 2.4.1 Data Quality and Assurance section, Corey
relative permeability curves are often used to ``smooth'' the intended
relative permeability curves to be used in the simulator in order to
improve convergence.

Corey {[}Corey, A. T. : ``The Interrelation Between gas and Oil Relative
Permeabilities'', Production Mon., 19. 38. (1954).{]} combined the work
of Purcell {[}Purcell, W. R., ``Capillary Pressures- Their Measurement
Using Mercury and the Calculation of Permeability Therefrom'',
Transactions AIME, 186, 39 (1949).{]} and Burdine {[}Burdine, N. T.,
``Relative Permeability Calculations from Pore Size Distribution Data'',
, Transactions AIME, 198, 71 (1953).{]} that was widely accepted for its
simplicity. His original equations were developed for the drainage cycle
in water-wet sandstones, but have also been used in carbonate
formations. Corey's original water-oil equations were as follows:

\begin{equation}\phantomsection\label{eq-8-17}{
{k}_{\mathit{ro}} ({S}_{w}) = {\left[\frac{1 - {S}_{w }}{1 - {S}_{\mathit{wcr} }}\right]}^{{n}_{o}}
}\end{equation}

\begin{equation}\phantomsection\label{eq-8-18}{
{k}_{\mathit{rw}} ({S}_{w}) = {\left[\frac{{S}_{w} - {S}_{\mathit{wcr} }}{1 - {S}_{\mathit{wcr} }}\right]}^{{n}_{w}}
}\end{equation}

where

kro(Sw) = relative permeability to oil,

krw(Sw) = relative permeability to water,

no = Corey oil exponent, set to four in the original paper,

nw = Corey water exponent, set to four in the original paper,

Sw = water saturation, and

Swcr = critical water saturation.

The denominator in equations (8.17) and (8.18) scales the water
saturation to the mobile water phase. There are several forms of these
equations, with the most common normalizing the saturation over the
mobile hydrocarbon phase, as depicted in equations (8.19) and (8.20).

\begin{equation}\phantomsection\label{eq-8-19}{
{k}_{\mathit{ro}} ({S}_{w}) = {k}_{\mathit{row}}{\left[\frac{1 - {S}_{w} - {S}_{\mathit{orw}}}{1 - {S}_{\mathit{orw}} - {S}_{\mathit{wcr}}}\right]}^{{n}_{o}}
}\end{equation}

\begin{equation}\phantomsection\label{eq-8-20}{
{k}_{\mathit{rw}} ({S}_{w}) = {k}_{\mathit{rww}}{\left[\frac{{S}_{w} - {S}_{\mathit{wcr}}}{1 - {S}_{\mathit{orw}} - {S}_{\mathit{wcr}}}\right]}^{{n}_{w}}
}\end{equation}

where

krow = maximum oil relative permeability at Swcr,

krw = maximum water relative permeability st Sorw,

Swcr = critical water saturation, and

Sorw = residual oil saturation under a water flood (SOWCR).

Similar equations exist for gas-oil and water-oil systems. As mentioned
above, the denominator in the equations, normalizes the saturation to
the mobile phase, as a consequence, it is still necessary to extend the
resulting Corey water curve to 100\% water saturation in order to
correctly model the water leg. In terms of the values for the various
Corey exponents, Table 8.10 offers some guidelines based on the rock's
wettability.

\begin{longtable}[]{@{}llllll@{}}
\toprule\noalign{}
Corey Guide (Water-Oil) & Corey Guide (Gas-Oil) & & & & \\
\midrule\noalign{}
\endhead
\bottomrule\noalign{}
\endlastfoot
Wettability & Kro & Krw & Wettability & Kro & Krg \\
Water wet & 2 - 4 & 5 - 8 & Water wet & & \\
Intermediate & 3 - 6 & 3 - 5 & Intermediate & & \\
Oil wet & 6 - 8 & 2 - 3 & Oil wet & 4 - 8 & 2 - 4 \\
\end{longtable}

\emph{Table 8.10: Corey Exponent Values}

\pandocbounded{\includegraphics[keepaspectratio]{chapters/images/Frame159_37e5b952fe69.png}}
\pandocbounded{\includegraphics[keepaspectratio]{chapters/images/Image148_37e5b952fe69.png}}
\pandocbounded{\includegraphics[keepaspectratio]{chapters/images/Frame160_13696c6e2167.png}}
\pandocbounded{\includegraphics[keepaspectratio]{chapters/images/Image150_13696c6e2167.png}}
\pandocbounded{\includegraphics[keepaspectratio]{chapters/images/Frame161_b043d8da4ab1.png}}
\pandocbounded{\includegraphics[keepaspectratio]{chapters/images/Image151_b043d8da4ab1.png}}

As can be seen in Figure 8.16 to Figure 8.18, the Corey exponents less
than one result in a concave curve, and values great than one result in
a convex curve, where as a value of one gives a straight line.

Unfortunately, neither OPM Flow nor the commercial simulator, unlike
some other simulators, support the direct entry of Corey type curves,
and thus Corey curves have to be generated outside of the simulator.
Fortunately, OPM Flow does support the more advanced and flexible LET
family of models instead, which are described in the next section.

\subsection{Saturation Table Generation - LET
Functions}\label{saturation-table-generation---let-functions}

OPM Flow has implemented the LET family of saturation functions models,
based on the work of Lomeland et al.~{[}Lomeland F., Ebeltoft E. and
Thomas W.H., 2005. A New Versatile Relative Permeability Correlation.
Paper SCA2005-32 presented at the International Symposium of the Society
of Core Analysts held in Toronto, Canada, 21-25 August, 2005.{]},
{[}Lomeland F. and Ebeltoft E., 2008. A New Versatile Capillary Pressure
Correlation. Paper SCA2008-08 presented at the International Symposium
of the Society of Core Analysts held in Abu Dhabi, UAE, 29 Oct.~-- 2
Nov., 2008.{]} {[}Lomeland F., Hasanov B., Ebeltoft E. and Berge M.,
2012. A Versatile Representation of Up-scaled Relative Permeability for
Field Applications. Paper SPE 154487-MS presented at the EAGE Annual
Conference \& Exhibition incorporating SPE Europec held in Copenhagen,
Denmark, 4-7 June 2012.{]} and {[}Lomeland F., 2018.Overview Of The Let
Family Of Versatile Correlations For Flow Functions. Paper SCA2018-056
presented at the International Symposium of the Society of Core Analysts
held in Trondheim, Norway, 27-30 August 2018.{]} via the SGOFLET and
SWOFLET keywords in this section. The keywords are used as replacements
for the SGOF and SWOF keywords for three-phase oil-gas-water systems,
and the SGWFN keyword for gas-water systems. Note there are two versions
of the LET functions, LET168 for two-phase flowing conditions and LETx
{[}Lomeland F. and Ebeltoft E., 2013. Versatile Three-phase Correlations
for Relative Permeability and Capillary Pressure. Paper SCA2013-034
presented at the International Symposium of the Society of Core Analysts
held in Napa Valley, California, USA, 16-19 September, 2013.{]} for
three-phase flowing conditions. All three keywords implement the LET
version for two phase systems.

The functions are dependent on the drainage and imbibition cycle of the
wetting phase as well as drainage and inhibition cycle number, since a
reservoir may undergo several flooding events. To account for this the
system defines the flooding event using the three saturations: Sw, So,
and Sg together with the state of the three saturations during the
flooding event. The saturation state can be increasing, decreasing, or
constant for a given flooding event cycle number (n). Thus, Sw(D),
So(I), Sg(C)1, or DIC1, means the water phase is decreasing, the oil
phase is increasing, and the gas phase is constant for the primary or
first cycle (n equals one). This is the case when oil is migrating into
the reservoir rock and displacing the initial water contained with the
reservoir. The various flooding processes are outlined in Table 8.11

\begin{longtable}[]{@{}
  >{\raggedright\arraybackslash}p{(\linewidth - 8\tabcolsep) * \real{0.2000}}
  >{\raggedright\arraybackslash}p{(\linewidth - 8\tabcolsep) * \real{0.2000}}
  >{\raggedright\arraybackslash}p{(\linewidth - 8\tabcolsep) * \real{0.2000}}
  >{\raggedright\arraybackslash}p{(\linewidth - 8\tabcolsep) * \real{0.2000}}
  >{\raggedright\arraybackslash}p{(\linewidth - 8\tabcolsep) * \real{0.2000}}@{}}
\toprule\noalign{}
\begin{minipage}[b]{\linewidth}\raggedright
Modeling of flooding processes
\end{minipage} & \begin{minipage}[b]{\linewidth}\raggedright
\end{minipage} & \begin{minipage}[b]{\linewidth}\raggedright
\end{minipage} & \begin{minipage}[b]{\linewidth}\raggedright
\end{minipage} & \begin{minipage}[b]{\linewidth}\raggedright
\end{minipage} \\
\midrule\noalign{}
\endhead
\bottomrule\noalign{}
\endlastfoot
Oil Field & Gas Field & Comment & & \\
Oil-Water Oil Leg & Oil-Gas Oil-Leg & Gas-Water Gas Cap & Oil-Gas Gas
Cap & Nomenclature LET(Sw, So, Sg, n) \\
DIC1 & CDI2 & DCI1 & CDI2 & Primary starting point. \\
IDC2 & CID3 & ICD2 & CID3 & Second event \\
DIC3 & CID4 & DCI3 & CID4 & Third event \\
\end{longtable}

\emph{Table 8.11: LET Modeling of Flooding Processes}

Thus, DIC1 is the well-known primary (first inflow number) drainage for
an oil/water system, and IDC2 is the well-known (second general inflow
number) imbibition for an oil/water system.

The following section describe the various flooding scenarios and the
standard LET equations used in the literature.

\subsubsection{Oil Displacing Water: Primary Drainage Cycle
(DIC1)}\label{oil-displacing-water-primary-drainage-cycle-dic1}

For the Primary Drainage, assuming water is the wetting phase, then,
Sw(D), So(I), Sg(C)1, or DIC1, means the water phase is decreasing, the
oil phase is increasing and the gas phase is constant for the primary or
first cycle (n equals one). This is the case for when oil is migrating
into the reservoir rock and displacing the initial water contained
within the reservoir. Under these circumstance the LET normalized water
saturation and relative permeability functions are defined as:

\begin{equation}\phantomsection\label{eq-8-21}{
{K}_{\mathit{row}} = \frac{{K}_{\mathit{rot}}{\left(1 - {S}_{wp}\right)}^{{L}_{o}}}{{\left(1 - {S}_{wp}\right)}^{{L}_{o}} + {E}_{o}{{S}_{wp}}^{{T}_{o}}}
}\end{equation}

\begin{equation}\phantomsection\label{eq-8-22}{
{K}_{\mathit{rw}} = \frac{{K}_{\mathit{rwt}}{S}_{wp}^{{L}_{w}}}{{S}_{wp}^{{L}_{w}} + {E}_{w}{\left(1 - {S}_{wp}\right)}^{{T}_{w}}}
}\end{equation}

\begin{equation}\phantomsection\label{eq-8-23}{
{S}_{wp} = \frac{{S}_{w} - {S}_{\mathit{wirr}}}{1 - {S}_{\mathit{wirr}}}
}\end{equation}

where

\({E}_{α}\) = empirical E parameter for the α phase,

\({K}_{\mathit{row}}\) = oil relative permeability with respect to water
saturation,

\({K}_{\mathit{rot}}\) = base or top oil relative permeability, used to
scale the maximum oil

permeability,

\({K}_{\mathit{rw}}\) = water relative permeability,

\({K}_{\mathit{rwt}}\) = base or top water relative permeability, used
to scale the maximum water

permeability,

\({L}_{α}\) = empirical L parameter for the α phase,

\({T}_{α}\) = empirical T parameter for the α phase,

\({S}_{w}\) = water saturation,

\({S}_{\mathit{wirr}}\) = irreducible water saturation, and

\({S}_{wp}\) = normalize water saturation.

The LET empirical parameters E, L, and T influence the shape of the
relative permeability curves, as depicted in Figure 8.19

\pandocbounded{\includegraphics[keepaspectratio]{chapters/images/Frame7_c68996d8b89d.png}}
\pandocbounded{\includegraphics[keepaspectratio]{chapters/images/Image152_c68996d8b89d.png}}

For the primary drainage or first cycle, DIC1, with the water phase
decreasing, the oil phase increasing and the gas phase constant, the LET
capillary pressure functions are:

\begin{equation}\phantomsection\label{eq-8-24}{
Y = \frac{{P}_{\mathit{cow}} - {P}_{\mathit{ct}}}{{P}_{\mathit{cir}} - {P}_{\mathit{ct}}} a\mathit{nd} Y = F\left({S}_{wp}\right) \mathit{where} {S}_{\mathit{wirr}}⩽{S}_{w}⩽1.0
}\end{equation}

\begin{equation}\phantomsection\label{eq-8-25}{
{P}_{\mathit{cow}}({S}_{w}={S}_{\mathit{wirr}}) = {P}_{\mathit{cir}} \text{and} {P}_{\mathit{cow}}({S}_{w}=1.0) = {P}_{\mathit{cit}}
}\end{equation}

\begin{equation}\phantomsection\label{eq-8-26}{
F = \frac{{\left(1 - {S}_{wp}\right)}^{L}}{{\left(1 - {S}_{wp}\right)}^{L} + E{S}_{wp}^{T}} \text{and} {S}_{wp} = \frac{{S}_{w} - {S}_{\mathit{wirr}}}{1 - {S}_{\mathit{wirr}}}
}\end{equation}

where

\(E\) = empirical E parameter,

\(L\) = empirical L parameter,

\({P}_{\mathit{cir}}\) = capillary pressure at irreducible water
saturation,

\({P}_{\mathit{ct}}\) = capillary threshold / entry pressure,

\(T\) = empirical T parameter,

\({S}_{w}\) = water saturation,

\({S}_{\mathit{wirr}}\) = irreducible water saturation, and

\({S}_{wp}\) = normalize water saturation.

The capillary pressure LET functions with either L or T equal to one are
called a semi-simple LET functions, and LET functions with both L and T
equal to one are called a simple LET function. The simple LET function
is equal to the Pc correlation of Honapour et al.~{[}Honarpour M.M.,
Djabbarah N.F. and Kralik J.G., 2004. Expert-Based Methodology for
Primary Drainage Capillary Pressure Measurements and Modeling. Paper
SPE-88709 presented at the 11th Abu Dhabi International Petroleum
Exhibition and Conference held in Abu Dhabi, UAE, 10-13 Oct., 2004.{]}
except for the arrangement of the empirical coefficients.

The L parameter describes the lower part of the curve, whereas the T
parameter describes the top part of the curve in a similar to how the L
parameter describes the lower part of the curve. Finally, the E
parameter describes the position of the slope of the curve, that is the
elevation. An E value of one is neutral, and the position of the slope
is governed by the L and T parameters in this case. Whereas, increasing
the value of the E moves the slope towards the high end of the curve,
and decreasing the value of the E moves the slope towards the lower end
of the curve. Lombard et al.169 noted that experience using the LET
correlation indicate that L, E, and T values should be greater than or
equal to 0.01, with no upper limit.

Figure 8.20 shows a typical LET derived drainage oil-water capillary
pressure curve, and Figure 8.21 shows a typical imbibition curves, the
later is discussed in the following section.

\pandocbounded{\includegraphics[keepaspectratio]{chapters/images/Frame157_c26f0918227b.png}}
\pandocbounded{\includegraphics[keepaspectratio]{chapters/images/Image153_c26f0918227b.png}}
\pandocbounded{\includegraphics[keepaspectratio]{chapters/images/Frame158_a9ece1e92dcf.png}}
\pandocbounded{\includegraphics[keepaspectratio]{chapters/images/Image154_a9ece1e92dcf.png}}

OPM has implemented a more restrictive form of the general LET
functions, more akin to the original paper168, as depicted in Figure
8.20. Table 8.12 relates the general LET equations presented previously
with the parameters on the SWOFLET keyword.

\begin{longtable}[]{@{}
  >{\raggedright\arraybackslash}p{(\linewidth - 10\tabcolsep) * \real{0.1667}}
  >{\raggedright\arraybackslash}p{(\linewidth - 10\tabcolsep) * \real{0.1667}}
  >{\raggedright\arraybackslash}p{(\linewidth - 10\tabcolsep) * \real{0.1667}}
  >{\raggedright\arraybackslash}p{(\linewidth - 10\tabcolsep) * \real{0.1667}}
  >{\raggedright\arraybackslash}p{(\linewidth - 10\tabcolsep) * \real{0.1667}}
  >{\raggedright\arraybackslash}p{(\linewidth - 10\tabcolsep) * \real{0.1667}}@{}}
\toprule\noalign{}
\begin{minipage}[b]{\linewidth}\raggedright
LET Function (DIC1) versus SWOFLET Keyword
\end{minipage} & \begin{minipage}[b]{\linewidth}\raggedright
\end{minipage} & \begin{minipage}[b]{\linewidth}\raggedright
\end{minipage} & \begin{minipage}[b]{\linewidth}\raggedright
\end{minipage} & \begin{minipage}[b]{\linewidth}\raggedright
\end{minipage} & \begin{minipage}[b]{\linewidth}\raggedright
\end{minipage} \\
\midrule\noalign{}
\endhead
\bottomrule\noalign{}
\endlastfoot
SWOFLET & LET Function & SWOFLET & LET Function & SWOFLET & LET
Function \\
SWL & \({S}_{\mathit{wirr}}\) & SORW & & PCIR &
\({P}_{\mathit{cir}}\) \\
SWCR & \({S}_{\mathit{wirr}}\) & SOWCR & & PCIT &
\({P}_{\mathit{ct}}\) \\
LWAT & \({L}_{w}\) & LOIL & \({L}_{o}\) & LPC & \(L\) \\
EWAT & \({E}_{w}\) & EOIL & \({E}_{o}\) & EPC & \(E\) \\
TWAT & \({T}_{w}\) & TOIL & \({T}_{o}\) & TPC & \(T\) \\
KRTWAT & \({K}_{\mathit{rwt}}\) & KRTOIL & \({K}_{\mathit{rot}}\) & & \\
Notes: & & & & & \\
\end{longtable}

\emph{Table 8.12: LET Functions (DIC1) versus SWOFLET Keyword}

\subsubsection{Water Displacing Oil: Primary Imbibition Cycle
(IDC2)}\label{water-displacing-oil-primary-imbibition-cycle-idc2}

Similarly for the oil producing phase, where the water phase is
increasing either by replacing the displaced oil or via water injection,
that is Sw(I), So(D), Sg(C), for the second flooding event, IDC2. Under
this scenario the LET equations are:

\begin{equation}\phantomsection\label{eq-8-27}{
{K}_{\mathit{row}} = \frac{{K}_{\mathit{rot}}{\left(1 - {S}_{\mathit{wn}}\right)}^{{L}_{o}}}{{\left(1 - {S}_{\mathit{wn}}\right)}^{{L}_{o}} + {E}_{o}{{S}_{\mathit{wn}}}^{{T}_{o}}}
}\end{equation}

\begin{equation}\phantomsection\label{eq-8-28}{
{K}_{\mathit{rw}} = \frac{{K}_{\mathit{rwr}}{S}_{\mathit{wn}}^{{L}_{w}}}{{S}_{\mathit{wn}}^{{L}_{w}} + {E}_{w}{\left(1 - {S}_{\mathit{wn}}\right)}^{{T}_{w}}}
}\end{equation}

\begin{equation}\phantomsection\label{eq-8-29}{
{S}_{\mathit{wn}} = \frac{{S}_{w} - {S}_{\mathit{wirr}}}{1 - {S}_{\mathit{orw}} - {S}_{\mathit{wirr}}}
}\end{equation}

where

\({E}_{α}\) = empirical E parameter for the α phase,

\({K}_{\mathit{row}}\) = oil relative permeability with respect to water
saturation,

\({K}_{\mathit{rot}}\) = base or top oil relative permeability, used to
scale the maximum oil

permeability,

\({K}_{\mathit{rw}}\) = water relative permeability,

\({K}_{\mathit{rwr}}\) = base or top water relative permeability, used
to scale the maximum water

permeability,

\({L}_{α}\) = empirical L parameter for the α phase,

\({T}_{α}\) = empirical T parameter for the α phase,

\({S}_{\mathit{orw}}\) = residual oil saturation under water
displacement,

\({S}_{w}\) = water saturation,

\({S}_{\mathit{wirr}}\) = irreducible water saturation, and

\({S}_{\mathit{wn}}\) = normalize water saturation.

For the primary imbibition or second cycle, IDC2, with the water phase
increasing, the oil phase decreasing and the gas phase constant, the LET
capillary pressure functions are:

\begin{equation}\phantomsection\label{eq-8-30}{
\begin{matrix}{P}_{\mathit{cow}}=({P}_{\mathit{cir}} - {P}_{\mathit{ct}}){G}_{n} + ({P}_{\mathit{cor}} - {P}_{\mathit{ct}}){F}_{n} + {P}_{\mathit{ct}} \\  \\ \mathit{where} {S}_{\mathit{wir}}\leq {S}_{w}\leq (1-{S}_{\mathit{orw}}) \text{and} {P}_{\mathit{cow}}({S}_{w}={S}_{\mathit{wzo}}) = 0\text{ => }{E}_{\mathit{sn}}\end{matrix}
}\end{equation}

\begin{equation}\phantomsection\label{eq-8-31}{
{G}_{n} = \frac{{\left(1 - {S}_{\mathit{wn}}\right)}^{{L}_{\mathit{sn}}}}{{\left(1 - {S}_{\mathit{wn}}\right)}^{{L}_{\mathit{sn}}} + {E}_{\mathit{sn}}{S}_{\mathit{wn}}^{{T}_{\mathit{sn}}}} {F}_{n} = \frac{{S}_{\mathit{wn}}^{{L}_{\mathit{fn}}}}{{S}_{\mathit{wn}}^{{L}_{\mathit{fn}}} + {E}_{\mathit{fn}}{\left(1 - {S}_{\mathit{wn}}\right)}^{{T}_{\mathit{fn}}}}
}\end{equation}

\begin{equation}\phantomsection\label{eq-8-32}{
{S}_{\mathit{wn}} = \frac{{S}_{w} - {S}_{\mathit{wirr}}}{1 - {S}_{\mathit{orw}} - {S}_{\mathit{wirr}}}
}\end{equation}

where

\({E}_{\mathit{fn}}\) = forced imbibition empirical E parameter,

\({E}_{\mathit{sn}}\) = spontaneous imbibition empirical E parameter,

\({L}_{\mathit{fn}}\) = forced imbibition empirical L parameter,

\({L}_{\mathit{sn}}\) = spontaneous imbibition empirical L parameter,

\({P}_{\mathit{cir}}\) = capillary pressure at irreducible water
saturation,

\({P}_{\mathit{ct}}\) = capillary threshold / entry pressure,

\({T}_{\mathit{fn}}\) = forced imbibition empirical T parameter,

\({T}_{\mathit{sn}}\) = spontaneous imbibition empirical T parameter,

\({S}_{\mathit{orw}}\) = residual oil saturation under water
displacement,

\({S}_{w}\) = water saturation,

\({S}_{\mathit{wirr}}\) = irreducible water saturation, and

\({S}_{\mathit{wn}}\) = normalize water saturation.

Spontaneous imbibition refers to the wetting phase displacing the
non-wetting phase, resulting in capillary pressure values greater than
zero. For example, in a water wet oil reservoir under water drive, the
water will easily imbibe into the pore space and displace the
non-wetting oil phase. Forced imbibition is when the non-wetting phase
displaces the wetting phase, and results in a negative capillary
pressure values necessary to overcome the capillary force holding the
wetting phase in-place. Thus, for a water wet oil reservoir, with the
oil zone expanding into the water zone, an additional force is required
to displace the water.

Table 8.13 relates the general LET equations presented previously with
the parameters on the SWOFLET keyword.

\begin{longtable}[]{@{}
  >{\raggedright\arraybackslash}p{(\linewidth - 10\tabcolsep) * \real{0.1667}}
  >{\raggedright\arraybackslash}p{(\linewidth - 10\tabcolsep) * \real{0.1667}}
  >{\raggedright\arraybackslash}p{(\linewidth - 10\tabcolsep) * \real{0.1667}}
  >{\raggedright\arraybackslash}p{(\linewidth - 10\tabcolsep) * \real{0.1667}}
  >{\raggedright\arraybackslash}p{(\linewidth - 10\tabcolsep) * \real{0.1667}}
  >{\raggedright\arraybackslash}p{(\linewidth - 10\tabcolsep) * \real{0.1667}}@{}}
\toprule\noalign{}
\begin{minipage}[b]{\linewidth}\raggedright
LET Function (IDC2) versus SWOFLET Keyword
\end{minipage} & \begin{minipage}[b]{\linewidth}\raggedright
\end{minipage} & \begin{minipage}[b]{\linewidth}\raggedright
\end{minipage} & \begin{minipage}[b]{\linewidth}\raggedright
\end{minipage} & \begin{minipage}[b]{\linewidth}\raggedright
\end{minipage} & \begin{minipage}[b]{\linewidth}\raggedright
\end{minipage} \\
\midrule\noalign{}
\endhead
\bottomrule\noalign{}
\endlastfoot
SWOFLET & LET Function & SWOFLET & LET Function & SWOFLET & LET
Function \\
SWL & \({S}_{\mathit{wirr}}\) & SORW & \({S}_{\mathit{orw}}\) & PCIR &
\({P}_{\mathit{cir}}\) \\
SWCR & \({S}_{\mathit{wirr}}\) & SOWCR & \({S}_{\mathit{orw}}\) & PCIT &
\({P}_{\mathit{ct}}\) \\
LWAT & \({L}_{w}\) & LOIL & \({L}_{o}\) & LPC & \({L}_{\mathit{sn}}\) \\
EWAT & \({E}_{w}\) & EOIL & \({E}_{o}\) & EPC & \({E}_{\mathit{sn}}\) \\
TWAT & \({T}_{w}\) & TOIL & \({T}_{o}\) & TPC & \({T}_{\mathit{sn}}\) \\
KRTWAT & \({K}_{\mathit{rwr}}\) & KRTOIL & \({K}_{\mathit{rot}}\) & & \\
Notes: & & & & & \\
\end{longtable}

\emph{Table 8.13: LET Functions (IDC2) versus SWOFLET Keyword}

\subsubsection{Gas Displacing Oil: Primary Imbibition
(CDI2)}\label{gas-displacing-oil-primary-imbibition-cdi2}

For the oil producing phase, where the gas phase is increasing either by
replacing the displaced oil or via gas injection, CDI2, the LET
equations are:

\begin{equation}\phantomsection\label{eq-8-33}{
{K}_{\mathit{rog}} = \frac{{K}_{\mathit{rot}}{\left(1 - {S}_{\mathit{gn}}\right)}^{{L}_{o}}}{{\left(1 - {S}_{\mathit{gn}}\right)}^{{L}_{o}} + {E}_{o}{{S}_{\mathit{gn}}}^{{T}_{o}}}
}\end{equation}

\begin{equation}\phantomsection\label{eq-8-34}{
{K}_{\mathit{rg}} = \frac{{K}_{\mathit{rgt}}{S}_{\mathit{gn}}^{{L}_{g}}}{{S}_{\mathit{gn}}^{{L}_{g}} + {E}_{g}{\left(1 - {S}_{\mathit{gn}}\right)}^{{T}_{g}}}
}\end{equation}

\begin{equation}\phantomsection\label{eq-8-35}{
{S}_{\mathit{gn}} = \frac{{S}_{g}}{1 - {S}_{\mathit{org}} - {S}_{\mathit{wirr}}}
}\end{equation}

where

\({E}_{α}\) = empirical E parameter for the α phase,

\({K}_{\mathit{rog}}\) = oil relative permeability with respect to gas
saturation,

\({K}_{\mathit{rot}}\) = base or top oil relative permeability, used to
scale the maximum oil

permeability,

\({K}_{\mathit{rg}}\) = gas relative permeability,

\({K}_{\mathit{rgt}}\) = base or top gas relative permeability, used to
scale the maximum gas

permeability,

\({L}_{α}\) = empirical L parameter for the α phase,

\({T}_{α}\) = empirical T parameter for the α phase,

\({S}_{g}\) = gas saturation,

\({S}_{\mathit{gn}}\) = normalize gas saturation,

\({S}_{\mathit{org}}\) = residual oil saturation under gas displacement,
and

\({S}_{\mathit{wirr}}\) = irreducible water saturation.

Similarly for the capillary pressure CDI2, with the water phase
constant, the oil phase decreasing and the gas phase increasing, the LET
capillary pressure functions are:

\begin{equation}\phantomsection\label{eq-8-36}{
\begin{matrix}{P}_{\mathit{cog}}=({P}_{\mathit{cir}} - {P}_{\mathit{ct}}){G}_{n} + ({P}_{\mathit{cor}} - {P}_{\mathit{ct}}){F}_{n} + {P}_{\mathit{ct}} \\  \\ \mathit{where}{S}_{\mathit{wir}}\leq {S}_{w}\leq (1-{S}_{\mathit{org}}-{S}_{\mathit{wirr}}) \text{and} {P}_{\mathit{cog}}({S}_{w}={S}_{\mathit{wzo}}) = 0\text{ => }{E}_{\mathit{sn}}\end{matrix}
}\end{equation}

\begin{equation}\phantomsection\label{eq-8-37}{
{G}_{n} = \frac{{\left(1 - {S}_{\mathit{gn}}\right)}^{{L}_{\mathit{sn}}}}{{\left(1 - {S}_{\mathit{gn}}\right)}^{{L}_{\mathit{sn}}} + {E}_{\mathit{sn}}{S}_{\mathit{gn}}^{{T}_{\mathit{sn}}}} {F}_{n} = \frac{{S}_{\mathit{gn}}^{{L}_{\mathit{fn}}}}{{S}_{\mathit{gn}}^{{L}_{\mathit{fn}}} + {E}_{\mathit{fn}}{\left(1 - {S}_{\mathit{gn}}\right)}^{{T}_{\mathit{fn}}}}
}\end{equation}

\begin{equation}\phantomsection\label{eq-8-38}{
{S}_{\mathit{gn}} = \frac{{S}_{g}}{1 - {S}_{\mathit{org}} - {S}_{\mathit{wirr}}}
}\end{equation}

where

\({E}_{\mathit{fn}}\) = forced imbibition empirical E parameter,

\({E}_{\mathit{sn}}\) = spontaneous imbibition empirical E parameter,

\({L}_{\mathit{fn}}\) = forced imbibition empirical L parameter,

\({L}_{\mathit{sn}}\) = spontaneous imbibition empirical L parameter,

\({P}_{\mathit{cir}}\) = capillary pressure at irreducible water
saturation,

\({P}_{\mathit{ct}}\) = capillary threshold/entry pressure,

\({T}_{\mathit{fn}}\) = forced imbibition empirical T parameter,

\({T}_{\mathit{sn}}\) = spontaneous imbibition empirical T parameter,

\({S}_{g}\) = gas saturation,

\({S}_{\mathit{gn}}\) = normalize gas saturation,

\({S}_{\mathit{org}}\) = residual oil saturation under gas displacement,
and

\({S}_{\mathit{wirr}}\) = irreducible water saturation.

Table 8.14 relates the general LET equations presented previously with
the parameters on the SGOFLET keyword.

\begin{longtable}[]{@{}
  >{\raggedright\arraybackslash}p{(\linewidth - 10\tabcolsep) * \real{0.1667}}
  >{\raggedright\arraybackslash}p{(\linewidth - 10\tabcolsep) * \real{0.1667}}
  >{\raggedright\arraybackslash}p{(\linewidth - 10\tabcolsep) * \real{0.1667}}
  >{\raggedright\arraybackslash}p{(\linewidth - 10\tabcolsep) * \real{0.1667}}
  >{\raggedright\arraybackslash}p{(\linewidth - 10\tabcolsep) * \real{0.1667}}
  >{\raggedright\arraybackslash}p{(\linewidth - 10\tabcolsep) * \real{0.1667}}@{}}
\toprule\noalign{}
\begin{minipage}[b]{\linewidth}\raggedright
LET Function (CDI2) versus SGOFLET Keyword
\end{minipage} & \begin{minipage}[b]{\linewidth}\raggedright
\end{minipage} & \begin{minipage}[b]{\linewidth}\raggedright
\end{minipage} & \begin{minipage}[b]{\linewidth}\raggedright
\end{minipage} & \begin{minipage}[b]{\linewidth}\raggedright
\end{minipage} & \begin{minipage}[b]{\linewidth}\raggedright
\end{minipage} \\
\midrule\noalign{}
\endhead
\bottomrule\noalign{}
\endlastfoot
SGOFLET & LET Function & SGOFLET & LET Function & SGOFLET & LET
Function \\
SGL & & SORG & \({S}_{\mathit{org}}\) & PCIR & \({P}_{\mathit{cir}}\) \\
SGCR & & SOGCR & \({S}_{\mathit{org}}\) & PCIT &
\({P}_{\mathit{ct}}\) \\
LGAS & \({L}_{g}\) & LOIL & \({L}_{o}\) & LPC & \({L}_{\mathit{sn}}\) \\
EGAS & \({E}_{g}\) & EOIL & \({E}_{o}\) & EPC & \({E}_{\mathit{sn}}\) \\
TGAS & \({T}_{g}\) & TOIL & \({T}_{o}\) & TPC & \({T}_{\mathit{sn}}\) \\
KRTGAS & \({K}_{\mathit{rgt}}\) & KRTOIL & \({K}_{\mathit{rot}}\) & & \\
Notes: & & & & & \\
\end{longtable}

\emph{Table 8.14: LET Functions (CDI2) versus SGOFLET Keyword}

\subsubsection{Gas Displacing Water: Primary Drainage Cycle
(DCI1)}\label{gas-displacing-water-primary-drainage-cycle-dci1}

For the Primary Drainage, assuming water is the wetting phase, then,
Sw(D), So(C), Sg(I)1, or DCI1, means the water phase is decreasing, the
oil phase is constant, and the gas phase is increasing for the primary
or first cycle (n equals one). This is the case for when gas is
migrating into the reservoir rock and displacing the initial water
contained within the reservoir. Under these circumstance the LET
normalized water saturation and relative permeability functions are
defined as:

\begin{equation}\phantomsection\label{eq-8-39}{
{K}_{\mathit{rgw}} = \frac{{K}_{\mathit{rgt}}{\left(1 - {S}_{wp}\right)}^{{L}_{g}}}{{\left(1 - {S}_{wp}\right)}^{{L}_{g}} + {E}_{g}{{S}_{wp}}^{{T}_{g}}}
}\end{equation}

\begin{equation}\phantomsection\label{eq-8-40}{
{K}_{\mathit{rw}} = \frac{{K}_{\mathit{rwt}}{S}_{wp}^{{L}_{w}}}{{S}_{wp}^{{L}_{w}} + {E}_{w}{\left(1 - {S}_{wp}\right)}^{{T}_{w}}}
}\end{equation}

\begin{equation}\phantomsection\label{eq-8-41}{
{S}_{wp} = \frac{{S}_{w} - {S}_{\mathit{wirr}}}{1 - {S}_{\mathit{wirr}}}
}\end{equation}

where

\({E}_{α}\) = empirical E parameter for the α phase,

\({K}_{\mathit{rgw}}\) = gas relative permeability with respect to water
saturation,

\({K}_{\mathit{rgt}}\) = base or top gas relative permeability, used to
scale the maximum gas

permeability,

\({K}_{\mathit{rw}}\) = water relative permeability,

\({K}_{\mathit{rwt}}\) = base or top water relative permeability, used
to scale the maximum water

permeability,

\({L}_{α}\) = empirical L parameter for the α phase,

\({T}_{α}\) = empirical T parameter for the α phase,

\({S}_{w}\) = water saturation,

\({S}_{\mathit{wirr}}\) = irreducible water saturation, and

\({S}_{wp}\) = normalize water saturation.

For the primary drainage or first cycle, DCI1, with the water phase
decreasing, the oil phase constant and the gas phase increasing, the LET
capillary pressure functions are:

\begin{equation}\phantomsection\label{eq-8-42}{
Y = \frac{{P}_{\mathit{cgw}} - {P}_{\mathit{ct}}}{{P}_{\mathit{cir}} - {P}_{\mathit{ct}}} a\mathit{nd} Y = F\left({S}_{wp}\right) \mathit{where} {S}_{\mathit{wirr}}⩽{S}_{w}⩽1.0
}\end{equation}

\begin{equation}\phantomsection\label{eq-8-43}{
{P}_{\mathit{cgw}}({S}_{w}={S}_{\mathit{wirr}}) = {P}_{\mathit{cir}} \text{and} {P}_{\mathit{cgw}}({S}_{w}=1.0) = {P}_{\mathit{cit}}
}\end{equation}

\begin{equation}\phantomsection\label{eq-8-44}{
F = \frac{{\left(1 - {S}_{wp}\right)}^{L}}{{\left(1 - {S}_{wp}\right)}^{L} + E{S}_{wp}^{T}} \text{and} {S}_{wp} = \frac{{S}_{w} - {S}_{\mathit{wirr}}}{1 - {S}_{\mathit{wirr}}}
}\end{equation}

where

\(E\) = empirical E parameter,

\(L\) = empirical L parameter,

\({P}_{\mathit{cir}}\) = capillary pressure at irreducible water
saturation,

\({P}_{\mathit{ct}}\) = capillary threshold/entry pressure,

\(T\) = empirical T parameter,

\({S}_{w}\) = water saturation,

\({S}_{\mathit{wirr}}\) = irreducible water saturation, and

\({S}_{wp}\) = normalize water saturation.

Table 8.15 relates the general LET equations presented previously with
the parameters on the \href{__RefHeading___Toc548127_947687768}{SGWFLET}
keyword.

\begin{longtable}[]{@{}
  >{\raggedright\arraybackslash}p{(\linewidth - 10\tabcolsep) * \real{0.1667}}
  >{\raggedright\arraybackslash}p{(\linewidth - 10\tabcolsep) * \real{0.1667}}
  >{\raggedright\arraybackslash}p{(\linewidth - 10\tabcolsep) * \real{0.1667}}
  >{\raggedright\arraybackslash}p{(\linewidth - 10\tabcolsep) * \real{0.1667}}
  >{\raggedright\arraybackslash}p{(\linewidth - 10\tabcolsep) * \real{0.1667}}
  >{\raggedright\arraybackslash}p{(\linewidth - 10\tabcolsep) * \real{0.1667}}@{}}
\toprule\noalign{}
\begin{minipage}[b]{\linewidth}\raggedright
LET Function (DCI1) versus SGWFLET Keyword
\end{minipage} & \begin{minipage}[b]{\linewidth}\raggedright
\end{minipage} & \begin{minipage}[b]{\linewidth}\raggedright
\end{minipage} & \begin{minipage}[b]{\linewidth}\raggedright
\end{minipage} & \begin{minipage}[b]{\linewidth}\raggedright
\end{minipage} & \begin{minipage}[b]{\linewidth}\raggedright
\end{minipage} \\
\midrule\noalign{}
\endhead
\bottomrule\noalign{}
\endlastfoot
SGWFLET & LET Function & SGWFLET & LET Function & SGWFLET & LET
Function \\
SGL & & SWL & \({S}_{\mathit{wirr}}\) & PCIR & \({P}_{\mathit{cir}}\) \\
SGCR & & SWCR & \({S}_{\mathit{wirr}}\) & PCIT &
\({P}_{\mathit{ct}}\) \\
LGAS & \({L}_{g}\) & LWAT & \({L}_{w}\) & LPC & \(L\) \\
EGAS & \({E}_{g}\) & EWAT & \({E}_{w}\) & EPC & \(E\) \\
TGAS & \({T}_{g}\) & TWAT & \({T}_{w}\) & TPC & \(T\) \\
KRTGAS & \({K}_{\mathit{rgt}}\) & KRTWAT & \({K}_{\mathit{rwr}}\) & & \\
Notes: & & & & & \\
\end{longtable}

\emph{Table 8.15: LET Functions (DCI1) versus SGWFLET Keyword}

\begin{tcolorbox}[enhanced jigsaw, title=\textcolor{quarto-callout-note-color}{\faInfo}\hspace{0.5em}{Note}, coltitle=black, breakable, toprule=.15mm, colbacktitle=quarto-callout-note-color!10!white, leftrule=.75mm, opacitybacktitle=0.6, colback=white, bottomtitle=1mm, toptitle=1mm, titlerule=0mm, left=2mm, colframe=quarto-callout-note-color-frame, arc=.35mm, rightrule=.15mm, bottomrule=.15mm, opacityback=0]

The SGWFLET keyword has not been implemented, as of this current
release.

\end{tcolorbox}

\subsubsection{Water Displacing Gas: Primary Imbibition Cycle
(ICD2)}\label{water-displacing-gas-primary-imbibition-cycle-icd2}

Similarly for the gas producing phase, where the water phase is
increasing either by replacing the displaced gas or via water influx
from the aquifer, that is Sw(I), So(C), Sg(D), for the second flooding
event, IDC2. Under this scenario the LET equations are:

\begin{equation}\phantomsection\label{eq-8-45}{
{K}_{\mathit{rgw}} = \frac{{K}_{\mathit{rgt}}{\left(1 - {S}_{\mathit{wn}}\right)}^{{L}_{o}}}{{\left(1 - {S}_{\mathit{wn}}\right)}^{{L}_{o}} + {E}_{o}{{S}_{\mathit{wn}}}^{{T}_{o}}}
}\end{equation}

\begin{equation}\phantomsection\label{eq-8-46}{
{K}_{\mathit{rw}} = \frac{{K}_{\mathit{rwr}}{S}_{\mathit{wn}}^{{L}_{w}}}{{S}_{\mathit{wn}}^{{L}_{w}} + {E}_{w}{\left(1 - {S}_{\mathit{wn}}\right)}^{{T}_{w}}}
}\end{equation}

\begin{equation}\phantomsection\label{eq-8-47}{
{S}_{\mathit{wn}} = \frac{{S}_{w} - {S}_{\mathit{wirr}}}{1 - {S}_{\mathit{grw}} - {S}_{\mathit{wirr}}}
}\end{equation}

where

\({E}_{α}\) = empirical E parameter for the α phase,

\({K}_{\mathit{rgw}}\) = gas relative permeability with respect to water
saturation,

\({K}_{\mathit{rot}}\) = base or top gas relative permeability, used to
scale the maximum gas

permeability,

\({K}_{\mathit{rw}}\) = water relative permeability,

\({K}_{\mathit{rwr}}\) = base or top water relative permeability, used
to scale the maximum water

permeability,

\({L}_{α}\) = empirical L parameter for the α phase,

\({T}_{α}\) = empirical T parameter for the α phase,

\({S}_{\mathit{ogw}}\) = residual gas saturation under water
displacement,

\({S}_{w}\) = water saturation,

\({S}_{\mathit{wirr}}\) = irreducible water saturation, and

\({S}_{\mathit{wn}}\) = normalize water saturation.

For the primary imbibition or second cycle, IDC2, with the water phase
increasing, the oil phase decreasing and the gas phase constant, the LET
capillary pressure functions are:

\begin{equation}\phantomsection\label{eq-8-48}{
\begin{matrix}{P}_{\mathit{cgw}}=({P}_{\mathit{cir}} - {P}_{\mathit{ct}}){G}_{n} + ({P}_{\mathit{cgr}} - {P}_{\mathit{ct}}){F}_{n} + {P}_{\mathit{ct}} \\  \\ \mathit{where} {S}_{\mathit{wir}}\leq {S}_{w}\leq (1-{S}_{\mathit{gorw}}) \text{and} {P}_{\mathit{cgw}}({S}_{w}={S}_{\mathit{wzo}}) = 0\text{ => }{E}_{\mathit{sn}}\end{matrix}
}\end{equation}

\begin{equation}\phantomsection\label{eq-8-49}{
{G}_{n} = \frac{{\left(1 - {S}_{\mathit{wn}}\right)}^{{L}_{\mathit{sn}}}}{{\left(1 - {S}_{\mathit{wn}}\right)}^{{L}_{\mathit{sn}}} + {E}_{\mathit{sn}}{S}_{\mathit{wn}}^{{T}_{\mathit{sn}}}} {F}_{n} = \frac{{S}_{\mathit{wn}}^{{L}_{\mathit{fn}}}}{{S}_{\mathit{wn}}^{{L}_{\mathit{fn}}} + {E}_{\mathit{fn}}{\left(1 - {S}_{\mathit{wn}}\right)}^{{T}_{\mathit{fn}}}}
}\end{equation}

\begin{equation}\phantomsection\label{eq-8-50}{
{S}_{\mathit{wn}} = \frac{{S}_{w} - {S}_{\mathit{wirr}}}{1 - {S}_{\mathit{grw}} - {S}_{\mathit{wirr}}}
}\end{equation}

where

\({E}_{\mathit{fn}}\) = forced imbibition empirical E parameter,

\({E}_{\mathit{sn}}\) = spontaneous imbibition empirical E parameter,

\({L}_{\mathit{fn}}\) = forced imbibition empirical L parameter,

\({L}_{\mathit{sn}}\) = spontaneous imbibition empirical L parameter,

\({P}_{\mathit{cir}}\) = capillary pressure at irreducible water
saturation,

\({P}_{\mathit{ct}}\) = capillary threshold / entry pressure,

\({T}_{\mathit{fn}}\) = forced imbibition empirical T parameter,

\({T}_{\mathit{sn}}\) = spontaneous imbibition empirical T parameter,

\({S}_{\mathit{ogw}}\) = residual gas saturation under water
displacement,

\({S}_{w}\) = water saturation,

\({S}_{\mathit{wirr}}\) = irreducible water saturation, and

\({S}_{\mathit{wn}}\) = normalize water saturation.

Table 8.16 relates the general LET equations presented previously with
the parameters on the \href{__RefHeading___Toc548127_947687768}{SGWFLET}
keyword.

\begin{longtable}[]{@{}
  >{\raggedright\arraybackslash}p{(\linewidth - 10\tabcolsep) * \real{0.1667}}
  >{\raggedright\arraybackslash}p{(\linewidth - 10\tabcolsep) * \real{0.1667}}
  >{\raggedright\arraybackslash}p{(\linewidth - 10\tabcolsep) * \real{0.1667}}
  >{\raggedright\arraybackslash}p{(\linewidth - 10\tabcolsep) * \real{0.1667}}
  >{\raggedright\arraybackslash}p{(\linewidth - 10\tabcolsep) * \real{0.1667}}
  >{\raggedright\arraybackslash}p{(\linewidth - 10\tabcolsep) * \real{0.1667}}@{}}
\toprule\noalign{}
\begin{minipage}[b]{\linewidth}\raggedright
LET Function (DCI1) versus SGWFLET Keyword
\end{minipage} & \begin{minipage}[b]{\linewidth}\raggedright
\end{minipage} & \begin{minipage}[b]{\linewidth}\raggedright
\end{minipage} & \begin{minipage}[b]{\linewidth}\raggedright
\end{minipage} & \begin{minipage}[b]{\linewidth}\raggedright
\end{minipage} & \begin{minipage}[b]{\linewidth}\raggedright
\end{minipage} \\
\midrule\noalign{}
\endhead
\bottomrule\noalign{}
\endlastfoot
SGWFLET & LET Function & SGWFLET & LET Function & SGWFLET & LET
Function \\
SGL & & SWL & \({S}_{\mathit{wirr}}\) & PCIR & \({P}_{\mathit{cir}}\) \\
SGCR & & SWCR & \({S}_{\mathit{wirr}}\) & PCIT &
\({P}_{\mathit{ct}}\) \\
LGAS & \({L}_{g}\) & LWAT & \({L}_{w}\) & LPC & \({L}_{\mathit{sn}}\) \\
EGAS & \({E}_{g}\) & EWAT & \({E}_{w}\) & EPC & \({E}_{\mathit{sn}}\) \\
TGAS & \({T}_{g}\) & TWAT & \({T}_{w}\) & TPC & \({T}_{\mathit{sn}}\) \\
KRTGAS & \({K}_{\mathit{rgt}}\) & KRTWAT & \({K}_{\mathit{rwr}}\) & & \\
Notes: & & & & & \\
\end{longtable}

\emph{Table 8.16: LET Functions (DCI1) versus SGWFLET Keyword}

\begin{tcolorbox}[enhanced jigsaw, title=\textcolor{quarto-callout-note-color}{\faInfo}\hspace{0.5em}{Note}, coltitle=black, breakable, toprule=.15mm, colbacktitle=quarto-callout-note-color!10!white, leftrule=.75mm, opacitybacktitle=0.6, colback=white, bottomtitle=1mm, toptitle=1mm, titlerule=0mm, left=2mm, colframe=quarto-callout-note-color-frame, arc=.35mm, rightrule=.15mm, bottomrule=.15mm, opacityback=0]

The SGWFLET keyword has not been implemented, as of this current
release.

\end{tcolorbox}

\subsection{Saturation Table
Allocation}\label{saturation-table-allocation}

If none of the relative permeability assignment options have been
activated on the SATOPTS keyword in the RUNSPEC section then saturation
tables may be alloacted to each cell using the SATNUM keyword in the
REGIONS section.

If the directional relative permeability option has been activated by
the DIRECT parameter on the SATOPTS keyword then direction dependent
relative permeability tables may be allocated to each cell using the
KRNUMX, KRNUMY, and KRNUMZ series of keywords in the REGIONS section. If
the irreversible directional relative permeability option has also been
activated by the DIRECT and IRREVERS parameters on the SATOPTS keyword
then relative permeability tables may be allocated to each face of each
cell using the KRNUMX, KRNUMX-, KRNUMY, KRNUMY-, KRNUMZ, and KRNUMZ-
series of keywords.

If the hysteresis option has been activated by the HYSTER parameter on
the SATOPTS keyword then the ``drainage'' and imbibition saturation
function tables may be allocated to each cell using the SATNUM and
IMBNUM keywords respectively in the REGIONS section.

If both the hystersis and directional relative permeability options have
been activated then the ``drainage'' relative permeabilities can be
allocated using the KRNUMX, KRNUMY, and KRNUMZ series of keywords, and
the imbibition relative permeabilities can be allocated using the
IMBNUMX, IMBNUMY, and IMBNUMZ series of keywords.

If the hystersis, directional, and irreversible relative permeability
options have been activated then the ``drainage'' relative
permeabilities can be allocated using the KRNUMX, KRNUMX-, KRNUMY,
KRNUMY-, KRNUMZ, and KRNUMZ- series of keywords, and the imbibition
relative permeabilities can be allocated using the IMBNUMX, IMBNUMX-,
IMBNUMY, IMBNUMY-, IMBNUMZ, and IMBNUMZ- series of keywords.

The above discussion assumes that the model has a Cartesian grid;
however, if the model has a radial grid then the equivalent radial
saturation table allocation keywords should be used. Thus, the
equivalent KRNUM keywords would be KRNUMR, KRNUMR-, KRNUMT, KRNUMT-,
KRNUMZ, and KRNUMZ-, and the equivalent IMBNUM keywords would be
IMBNUMR, IMBNUMR-, IMBNUMT, IMBNUMT-, IMBNUMZ, and IMBNUMZ-. Note that
the KRNUMZ, KRNUMZ-, IMBNUMZ, and IMBNUMZ- keywords are the same for
both Cartesian and radial grids.

\section{Keyword Definitions}\label{keyword-definitions-10}

\subsection{ACF -- Define Acentric
Factors}\label{acf-define-acentric-factors}

\begin{longtable}[]{@{}
  >{\raggedright\arraybackslash}p{(\linewidth - 14\tabcolsep) * \real{0.1250}}
  >{\raggedright\arraybackslash}p{(\linewidth - 14\tabcolsep) * \real{0.1250}}
  >{\raggedright\arraybackslash}p{(\linewidth - 14\tabcolsep) * \real{0.1250}}
  >{\raggedright\arraybackslash}p{(\linewidth - 14\tabcolsep) * \real{0.1250}}
  >{\raggedright\arraybackslash}p{(\linewidth - 14\tabcolsep) * \real{0.1250}}
  >{\raggedright\arraybackslash}p{(\linewidth - 14\tabcolsep) * \real{0.1250}}
  >{\raggedright\arraybackslash}p{(\linewidth - 14\tabcolsep) * \real{0.1250}}
  >{\raggedright\arraybackslash}p{(\linewidth - 14\tabcolsep) * \real{0.1250}}@{}}
\toprule\noalign{}
\begin{minipage}[b]{\linewidth}\raggedright
RUNSPEC
\end{minipage} & \begin{minipage}[b]{\linewidth}\raggedright
GRID
\end{minipage} & \begin{minipage}[b]{\linewidth}\raggedright
EDIT
\end{minipage} & \begin{minipage}[b]{\linewidth}\raggedright
PROPS
\end{minipage} & \begin{minipage}[b]{\linewidth}\raggedright
REGIONS
\end{minipage} & \begin{minipage}[b]{\linewidth}\raggedright
SOLUTION
\end{minipage} & \begin{minipage}[b]{\linewidth}\raggedright
SUMMARY
\end{minipage} & \begin{minipage}[b]{\linewidth}\raggedright
SCHEDULE
\end{minipage} \\
\midrule\noalign{}
\endhead
\bottomrule\noalign{}
\endlastfoot
\end{longtable}

\subsubsection{Description}\label{description-386}

The \hyperref[REF_HEADING_KEYWORD_ACF]{ACF} keyword defines the acentric
factors for each of the compositional components active in the model and
for each equation of state. The keyword should only be used if the
compositional mode has been requested using the COMPS keyword in the
RUNSPEC section.

OPM Flow does not currently support the general compositional modeling
formulation.

This keyword is not supported by OPM Flow but it will be parsed and its
data ignored.

\begin{longtable}[]{@{}
  >{\raggedright\arraybackslash}p{(\linewidth - 6\tabcolsep) * \real{0.2500}}
  >{\raggedright\arraybackslash}p{(\linewidth - 6\tabcolsep) * \real{0.2500}}
  >{\raggedright\arraybackslash}p{(\linewidth - 6\tabcolsep) * \real{0.2500}}
  >{\raggedright\arraybackslash}p{(\linewidth - 6\tabcolsep) * \real{0.2500}}@{}}
\toprule\noalign{}
\begin{minipage}[b]{\linewidth}\raggedright
No.
\end{minipage} & \begin{minipage}[b]{\linewidth}\raggedright
Name
\end{minipage} & \begin{minipage}[b]{\linewidth}\raggedright
Description
\end{minipage} & \begin{minipage}[b]{\linewidth}\raggedright
Default
\end{minipage} \\
\midrule\noalign{}
\endhead
\bottomrule\noalign{}
\endlastfoot
1 & \hyperref[REF_HEADING_KEYWORD_ACF]{ACF} & A series of real numbers
that define the acentric factors for each of the compositional
components active in the model. & None \\
Notes: & & & \\
\end{longtable}

\emph{Table 8.3.1.1: ACF Keyword Description}

\subsubsection{Examples}\label{examples-20}

The following example defines the acentric factors for each component in
a single three-component equation of state model.

\begin{verbatim}
--
-- Acentric Factors
--
ACF
  0.0108  0.2273  0.3434 /
\end{verbatim}

The following example defines the acentric factors for each component in
two three-component equation of state models.

\begin{verbatim}
--
-- Acentric Factors
--
ACF
  0.0108  0.2273  0.3434 /
  0.0110  0.2281  0.3428 /
\end{verbatim}

\subsection{ACTCO2S -- Activity Model for
CO2}\label{actco2s-activity-model-for-co2}

\begin{longtable}[]{@{}
  >{\raggedright\arraybackslash}p{(\linewidth - 14\tabcolsep) * \real{0.1250}}
  >{\raggedright\arraybackslash}p{(\linewidth - 14\tabcolsep) * \real{0.1250}}
  >{\raggedright\arraybackslash}p{(\linewidth - 14\tabcolsep) * \real{0.1250}}
  >{\raggedright\arraybackslash}p{(\linewidth - 14\tabcolsep) * \real{0.1250}}
  >{\raggedright\arraybackslash}p{(\linewidth - 14\tabcolsep) * \real{0.1250}}
  >{\raggedright\arraybackslash}p{(\linewidth - 14\tabcolsep) * \real{0.1250}}
  >{\raggedright\arraybackslash}p{(\linewidth - 14\tabcolsep) * \real{0.1250}}
  >{\raggedright\arraybackslash}p{(\linewidth - 14\tabcolsep) * \real{0.1250}}@{}}
\toprule\noalign{}
\begin{minipage}[b]{\linewidth}\raggedright
RUNSPEC
\end{minipage} & \begin{minipage}[b]{\linewidth}\raggedright
GRID
\end{minipage} & \begin{minipage}[b]{\linewidth}\raggedright
EDIT
\end{minipage} & \begin{minipage}[b]{\linewidth}\raggedright
PROPS
\end{minipage} & \begin{minipage}[b]{\linewidth}\raggedright
REGIONS
\end{minipage} & \begin{minipage}[b]{\linewidth}\raggedright
SOLUTION
\end{minipage} & \begin{minipage}[b]{\linewidth}\raggedright
SUMMARY
\end{minipage} & \begin{minipage}[b]{\linewidth}\raggedright
SCHEDULE
\end{minipage} \\
\midrule\noalign{}
\endhead
\bottomrule\noalign{}
\endlastfoot
\end{longtable}

\subsubsection{Description}\label{description-387}

The keyword \hyperref[REF_HEADING_KEYWORD_ACTCO2S]{ACTCO2S} specifies
the activity model for salting-out effects when calculating mutual
solubility in the CO2 storage module which is activated by the CO2STORE
keyword in the RUNSPEC section.

\begin{longtable}[]{@{}
  >{\raggedright\arraybackslash}p{(\linewidth - 6\tabcolsep) * \real{0.2500}}
  >{\raggedright\arraybackslash}p{(\linewidth - 6\tabcolsep) * \real{0.2500}}
  >{\raggedright\arraybackslash}p{(\linewidth - 6\tabcolsep) * \real{0.2500}}
  >{\raggedright\arraybackslash}p{(\linewidth - 6\tabcolsep) * \real{0.2500}}@{}}
\toprule\noalign{}
\begin{minipage}[b]{\linewidth}\raggedright
No.
\end{minipage} & \begin{minipage}[b]{\linewidth}\raggedright
Name
\end{minipage} & \begin{minipage}[b]{\linewidth}\raggedright
Description
\end{minipage} & \begin{minipage}[b]{\linewidth}\raggedright
Default
\end{minipage} \\
\midrule\noalign{}
\endhead
\bottomrule\noalign{}
\endlastfoot
1 & ACTMODEL & A positive integer value selecting the activity model for
the salting-out effect. The choices are: & 3 \\
Notes: & & & \\
\end{longtable}

\emph{Table 8.3.1.1: \hyperref[REF_HEADING_KEYWORD_ACTCO2S]{ACTCO2S}
Keyword Description}

\subsubsection{Example}\label{example-257}

The following example activates activity model number 1.

\begin{verbatim}
--
--       ACTIVITY MODEL FOR SALTING-OUT EFFECTS IN CO2
--
ACTCO2S
         1                                                 /
\end{verbatim}

\subsection{ADD -- Add a Constant to a Specified
Array}\label{add-add-a-constant-to-a-specified-array-2}

The ADD keyword adds a constant to a specified array or part of an
array. The constant can be real or integer depending on the array type;
however, the arrays that can be operated on is dependent on which
section the ADD keyword is being applied.

See ADD -- Add a Constant to a Specified Array in the GRID section for a
full description.

\subsection{ADDREG -- Add a Constant to an Array based on a Region
Number}\label{addreg-add-a-constant-to-an-array-based-on-a-region-number-2}

The ADDREG keyword adds a constant to a specified array or part of an
array based on cells with a specific region number. The region number
array can be FLUXNUM, MULTNUM or OPERNUM and these arrays must be
defined and be available before the ADDREG keyword is read by the
simulator. The constant can be real or integer depending on the property
array type; however, the property arrays that can be operated on is
dependent on which section the ADDREG keyword is being used.

See ADDREG -- Add a Constant to an Array based on a Region Number in the
GRID section for a full description.

\subsection{ADSALNOD -- Salt Concentration Based on SATNUM
Array}\label{adsalnod-salt-concentration-based-on-satnum-array}

\begin{longtable}[]{@{}
  >{\raggedright\arraybackslash}p{(\linewidth - 14\tabcolsep) * \real{0.1250}}
  >{\raggedright\arraybackslash}p{(\linewidth - 14\tabcolsep) * \real{0.1250}}
  >{\raggedright\arraybackslash}p{(\linewidth - 14\tabcolsep) * \real{0.1250}}
  >{\raggedright\arraybackslash}p{(\linewidth - 14\tabcolsep) * \real{0.1250}}
  >{\raggedright\arraybackslash}p{(\linewidth - 14\tabcolsep) * \real{0.1250}}
  >{\raggedright\arraybackslash}p{(\linewidth - 14\tabcolsep) * \real{0.1250}}
  >{\raggedright\arraybackslash}p{(\linewidth - 14\tabcolsep) * \real{0.1250}}
  >{\raggedright\arraybackslash}p{(\linewidth - 14\tabcolsep) * \real{0.1250}}@{}}
\toprule\noalign{}
\begin{minipage}[b]{\linewidth}\raggedright
RUNSPEC
\end{minipage} & \begin{minipage}[b]{\linewidth}\raggedright
GRID
\end{minipage} & \begin{minipage}[b]{\linewidth}\raggedright
EDIT
\end{minipage} & \begin{minipage}[b]{\linewidth}\raggedright
PROPS
\end{minipage} & \begin{minipage}[b]{\linewidth}\raggedright
REGIONS
\end{minipage} & \begin{minipage}[b]{\linewidth}\raggedright
SOLUTION
\end{minipage} & \begin{minipage}[b]{\linewidth}\raggedright
SUMMARY
\end{minipage} & \begin{minipage}[b]{\linewidth}\raggedright
SCHEDULE
\end{minipage} \\
\midrule\noalign{}
\endhead
\bottomrule\noalign{}
\endlastfoot
\end{longtable}

\subsubsection{Description}\label{description-388}

ADSALNOD defines the salt concentration value based on a cells SATNUM
number. The ADSALNOD property is used in the calculation of a polymer
viscosity when the polymer and the salt options has been activated by
the POLYMER and BRINE keywords in the RUNSPEC section. In the RUNSPEC
section the number of SATNUM functions is declared by the NTSFUN
variable on the TABDIMS keyword and allocated to individual cells by the
SATNUM property array in the REGIONS section. NSSFUN on the TABDIMS
keyword in the RUNSPEC section defines the maximum number of rows (or
saturation values) in the relative permeability saturation tables and
also sets the maximum number of entries for each ADSALNOD data set. The
number of values for each data set must correspond to the number of
polymer solution adsorption entries on the PLYADSS keyword. For example,
if there are three sets of relative permeability tables and four values
on the PLYADSS keyword, then three ADSALNOD data sets with four values
of salt concentrations need to be entered.

The salt concentrations within each data set should be positive and
monotonically increasing and each ADSALNOD data set is delimited by
``/'' including the last data set.

\begin{longtable}[]{@{}
  >{\raggedright\arraybackslash}p{(\linewidth - 10\tabcolsep) * \real{0.1667}}
  >{\raggedright\arraybackslash}p{(\linewidth - 10\tabcolsep) * \real{0.1667}}
  >{\raggedright\arraybackslash}p{(\linewidth - 10\tabcolsep) * \real{0.1667}}
  >{\raggedright\arraybackslash}p{(\linewidth - 10\tabcolsep) * \real{0.1667}}
  >{\raggedright\arraybackslash}p{(\linewidth - 10\tabcolsep) * \real{0.1667}}
  >{\raggedright\arraybackslash}p{(\linewidth - 10\tabcolsep) * \real{0.1667}}@{}}
\toprule\noalign{}
\begin{minipage}[b]{\linewidth}\raggedright
No.
\end{minipage} & \begin{minipage}[b]{\linewidth}\raggedright
Name
\end{minipage} & \begin{minipage}[b]{\linewidth}\raggedright
Description
\end{minipage} & \begin{minipage}[b]{\linewidth}\raggedright
Default
\end{minipage} & \begin{minipage}[b]{\linewidth}\raggedright
\end{minipage} & \begin{minipage}[b]{\linewidth}\raggedright
\end{minipage} \\
\midrule\noalign{}
\endhead
\bottomrule\noalign{}
\endlastfoot
1 & SALTCON & Field & Metric & Laboratory & None \\
A real positive columnar vector that sets the salt concentrations for
the given relative permeability saturation tables. & & & & & \\
lb/stb & kg/sm3 & gm/scc & & & \\
Notes: & & & & & \\
\end{longtable}

\emph{Table 8.17: ADSALNOD Keyword Description}

An alternative manner of entering the salt concentrations is by
utilizing the PVTNUM region array by using the SALTNODE keyword in the
PROPS section.

This keyword is not supported by OPM Flow but would change the results
if supported so the simulation will be stopped.

\subsubsection{Example}\label{example-258}

Given three sets of relative permeability tables and four values on the
PLYADSS keyword, then the data salt concentration should be entered as
follows:

\begin{verbatim}
--
--    SETS SALT CONCENTRATION FOR POLYMER SOLUTION ADSORPTION
--    VIA SATNUM ARRAY ALLOCATION
--
--    SALT
--
ADSALNOD
      1.0
      5.0
      10.5
      25.0      / SATNUM TABLE NO. 01
      1.0
      3.0
      7.5
      15.0      / SATNUM TABLE NO. 02
      1.0
      7.5
      20.5
      35.0      / SATNUM TABLE NO. 03
\end{verbatim}

See also the SALTNODE keyword.

\subsection{ADSORP -- Define Generalized Langmuir Adsorption
Function}\label{adsorp-define-generalized-langmuir-adsorption-function}

\begin{longtable}[]{@{}
  >{\raggedright\arraybackslash}p{(\linewidth - 14\tabcolsep) * \real{0.1250}}
  >{\raggedright\arraybackslash}p{(\linewidth - 14\tabcolsep) * \real{0.1250}}
  >{\raggedright\arraybackslash}p{(\linewidth - 14\tabcolsep) * \real{0.1250}}
  >{\raggedright\arraybackslash}p{(\linewidth - 14\tabcolsep) * \real{0.1250}}
  >{\raggedright\arraybackslash}p{(\linewidth - 14\tabcolsep) * \real{0.1250}}
  >{\raggedright\arraybackslash}p{(\linewidth - 14\tabcolsep) * \real{0.1250}}
  >{\raggedright\arraybackslash}p{(\linewidth - 14\tabcolsep) * \real{0.1250}}
  >{\raggedright\arraybackslash}p{(\linewidth - 14\tabcolsep) * \real{0.1250}}@{}}
\toprule\noalign{}
\begin{minipage}[b]{\linewidth}\raggedright
RUNSPEC
\end{minipage} & \begin{minipage}[b]{\linewidth}\raggedright
GRID
\end{minipage} & \begin{minipage}[b]{\linewidth}\raggedright
EDIT
\end{minipage} & \begin{minipage}[b]{\linewidth}\raggedright
PROPS
\end{minipage} & \begin{minipage}[b]{\linewidth}\raggedright
REGIONS
\end{minipage} & \begin{minipage}[b]{\linewidth}\raggedright
SOLUTION
\end{minipage} & \begin{minipage}[b]{\linewidth}\raggedright
SUMMARY
\end{minipage} & \begin{minipage}[b]{\linewidth}\raggedright
SCHEDULE
\end{minipage} \\
\midrule\noalign{}
\endhead
\bottomrule\noalign{}
\endlastfoot
\end{longtable}

\subsubsection{Description}\label{description-389}

The ADSORP keyword defines the parameters for the generalized Langmuir
Adsorption {[}Langmuir, Irving (June 1918). ``The Adsorption of Gases on
Plane Surface of Glass, Mica and Platinum''. The Research Laboratory of
the General Electric Company. 40 (9): 1361--1402.
doi:10.1021/ja02242a004{]} function for when the polymer, surfact,
alkaline, foam and tracers phases have been activated in the RUNSPEC
section by the POLYMER, SURFACT, ALKALINE, FOAM and TRACER keywords.

This keyword is not supported by OPM Flow but would change the results
if supported so the simulation will be stopped.

\subsection{ALKADS -- Define Alkaline Adsorption
Functions}\label{alkads-define-alkaline-adsorption-functions}

\begin{longtable}[]{@{}
  >{\raggedright\arraybackslash}p{(\linewidth - 14\tabcolsep) * \real{0.1250}}
  >{\raggedright\arraybackslash}p{(\linewidth - 14\tabcolsep) * \real{0.1250}}
  >{\raggedright\arraybackslash}p{(\linewidth - 14\tabcolsep) * \real{0.1250}}
  >{\raggedright\arraybackslash}p{(\linewidth - 14\tabcolsep) * \real{0.1250}}
  >{\raggedright\arraybackslash}p{(\linewidth - 14\tabcolsep) * \real{0.1250}}
  >{\raggedright\arraybackslash}p{(\linewidth - 14\tabcolsep) * \real{0.1250}}
  >{\raggedright\arraybackslash}p{(\linewidth - 14\tabcolsep) * \real{0.1250}}
  >{\raggedright\arraybackslash}p{(\linewidth - 14\tabcolsep) * \real{0.1250}}@{}}
\toprule\noalign{}
\begin{minipage}[b]{\linewidth}\raggedright
RUNSPEC
\end{minipage} & \begin{minipage}[b]{\linewidth}\raggedright
GRID
\end{minipage} & \begin{minipage}[b]{\linewidth}\raggedright
EDIT
\end{minipage} & \begin{minipage}[b]{\linewidth}\raggedright
PROPS
\end{minipage} & \begin{minipage}[b]{\linewidth}\raggedright
REGIONS
\end{minipage} & \begin{minipage}[b]{\linewidth}\raggedright
SOLUTION
\end{minipage} & \begin{minipage}[b]{\linewidth}\raggedright
SUMMARY
\end{minipage} & \begin{minipage}[b]{\linewidth}\raggedright
SCHEDULE
\end{minipage} \\
\midrule\noalign{}
\endhead
\bottomrule\noalign{}
\endlastfoot
\end{longtable}

\subsubsection{Description}\label{description-390}

ALKADS defines the alkaline adsorption functions for when the alkaline
model has been activated via the ALKALINE keyword in the RUNSPEC
section.

This keyword is not supported by OPM Flow but would change the results
if supported so the simulation will be stopped.

\subsection{ALKROCK -- Define Rock Alkaline
Properties}\label{alkrock-define-rock-alkaline-properties}

\begin{longtable}[]{@{}
  >{\raggedright\arraybackslash}p{(\linewidth - 14\tabcolsep) * \real{0.1250}}
  >{\raggedright\arraybackslash}p{(\linewidth - 14\tabcolsep) * \real{0.1250}}
  >{\raggedright\arraybackslash}p{(\linewidth - 14\tabcolsep) * \real{0.1250}}
  >{\raggedright\arraybackslash}p{(\linewidth - 14\tabcolsep) * \real{0.1250}}
  >{\raggedright\arraybackslash}p{(\linewidth - 14\tabcolsep) * \real{0.1250}}
  >{\raggedright\arraybackslash}p{(\linewidth - 14\tabcolsep) * \real{0.1250}}
  >{\raggedright\arraybackslash}p{(\linewidth - 14\tabcolsep) * \real{0.1250}}
  >{\raggedright\arraybackslash}p{(\linewidth - 14\tabcolsep) * \real{0.1250}}@{}}
\toprule\noalign{}
\begin{minipage}[b]{\linewidth}\raggedright
RUNSPEC
\end{minipage} & \begin{minipage}[b]{\linewidth}\raggedright
GRID
\end{minipage} & \begin{minipage}[b]{\linewidth}\raggedright
EDIT
\end{minipage} & \begin{minipage}[b]{\linewidth}\raggedright
PROPS
\end{minipage} & \begin{minipage}[b]{\linewidth}\raggedright
REGIONS
\end{minipage} & \begin{minipage}[b]{\linewidth}\raggedright
SOLUTION
\end{minipage} & \begin{minipage}[b]{\linewidth}\raggedright
SUMMARY
\end{minipage} & \begin{minipage}[b]{\linewidth}\raggedright
SCHEDULE
\end{minipage} \\
\midrule\noalign{}
\endhead
\bottomrule\noalign{}
\endlastfoot
\end{longtable}

\subsubsection{Description}\label{description-391}

The ALKROCK keyword defines the rock alkaline properties for when the
alkaline model has been activated via the ALKALINE keyword in the
RUNSPEC section.

This keyword is not supported by OPM Flow but would change the results
if supported so the simulation will be stopped.

\subsection{ALPOLADS -- Polymer Adsorption versus Alkaline Concentration
Multipliers}\label{alpolads-polymer-adsorption-versus-alkaline-concentration-multipliers}

\begin{longtable}[]{@{}
  >{\raggedright\arraybackslash}p{(\linewidth - 14\tabcolsep) * \real{0.1250}}
  >{\raggedright\arraybackslash}p{(\linewidth - 14\tabcolsep) * \real{0.1250}}
  >{\raggedright\arraybackslash}p{(\linewidth - 14\tabcolsep) * \real{0.1250}}
  >{\raggedright\arraybackslash}p{(\linewidth - 14\tabcolsep) * \real{0.1250}}
  >{\raggedright\arraybackslash}p{(\linewidth - 14\tabcolsep) * \real{0.1250}}
  >{\raggedright\arraybackslash}p{(\linewidth - 14\tabcolsep) * \real{0.1250}}
  >{\raggedright\arraybackslash}p{(\linewidth - 14\tabcolsep) * \real{0.1250}}
  >{\raggedright\arraybackslash}p{(\linewidth - 14\tabcolsep) * \real{0.1250}}@{}}
\toprule\noalign{}
\begin{minipage}[b]{\linewidth}\raggedright
RUNSPEC
\end{minipage} & \begin{minipage}[b]{\linewidth}\raggedright
GRID
\end{minipage} & \begin{minipage}[b]{\linewidth}\raggedright
EDIT
\end{minipage} & \begin{minipage}[b]{\linewidth}\raggedright
PROPS
\end{minipage} & \begin{minipage}[b]{\linewidth}\raggedright
REGIONS
\end{minipage} & \begin{minipage}[b]{\linewidth}\raggedright
SOLUTION
\end{minipage} & \begin{minipage}[b]{\linewidth}\raggedright
SUMMARY
\end{minipage} & \begin{minipage}[b]{\linewidth}\raggedright
SCHEDULE
\end{minipage} \\
\midrule\noalign{}
\endhead
\bottomrule\noalign{}
\endlastfoot
\end{longtable}

\subsubsection{Description}\label{description-392}

ALPOLDS defines the polymer adsorption versus alkaline concentration
multipliers for when the alkaline model has been activated via the
ALKALINE keyword in the RUNSPEC section.

This keyword is not supported by OPM Flow but would change the results
if supported so the simulation will be stopped.

\subsection{ALSURFAD -- Surfactant Adsorption versus Alkaline
Concentration
Multipliers}\label{alsurfad-surfactant-adsorption-versus-alkaline-concentration-multipliers}

\begin{longtable}[]{@{}
  >{\raggedright\arraybackslash}p{(\linewidth - 14\tabcolsep) * \real{0.1250}}
  >{\raggedright\arraybackslash}p{(\linewidth - 14\tabcolsep) * \real{0.1250}}
  >{\raggedright\arraybackslash}p{(\linewidth - 14\tabcolsep) * \real{0.1250}}
  >{\raggedright\arraybackslash}p{(\linewidth - 14\tabcolsep) * \real{0.1250}}
  >{\raggedright\arraybackslash}p{(\linewidth - 14\tabcolsep) * \real{0.1250}}
  >{\raggedright\arraybackslash}p{(\linewidth - 14\tabcolsep) * \real{0.1250}}
  >{\raggedright\arraybackslash}p{(\linewidth - 14\tabcolsep) * \real{0.1250}}
  >{\raggedright\arraybackslash}p{(\linewidth - 14\tabcolsep) * \real{0.1250}}@{}}
\toprule\noalign{}
\begin{minipage}[b]{\linewidth}\raggedright
RUNSPEC
\end{minipage} & \begin{minipage}[b]{\linewidth}\raggedright
GRID
\end{minipage} & \begin{minipage}[b]{\linewidth}\raggedright
EDIT
\end{minipage} & \begin{minipage}[b]{\linewidth}\raggedright
PROPS
\end{minipage} & \begin{minipage}[b]{\linewidth}\raggedright
REGIONS
\end{minipage} & \begin{minipage}[b]{\linewidth}\raggedright
SOLUTION
\end{minipage} & \begin{minipage}[b]{\linewidth}\raggedright
SUMMARY
\end{minipage} & \begin{minipage}[b]{\linewidth}\raggedright
SCHEDULE
\end{minipage} \\
\midrule\noalign{}
\endhead
\bottomrule\noalign{}
\endlastfoot
\end{longtable}

\subsubsection{Description}\label{description-393}

ALSURAD defines the surfactant adsorption versus alkaline concentration
multipliers for when the alkaline model has been activated via the
ALKALINE keyword in the RUNSPEC section.

This keyword is not supported by OPM Flow but would change the results
if supported so the simulation will be stopped.

\subsection{ALSURFST -- Water-Oil Surface Tension versus Alkaline
Concentration
Multipliers}\label{alsurfst-water-oil-surface-tension-versus-alkaline-concentration-multipliers}

\begin{longtable}[]{@{}
  >{\raggedright\arraybackslash}p{(\linewidth - 14\tabcolsep) * \real{0.1250}}
  >{\raggedright\arraybackslash}p{(\linewidth - 14\tabcolsep) * \real{0.1250}}
  >{\raggedright\arraybackslash}p{(\linewidth - 14\tabcolsep) * \real{0.1250}}
  >{\raggedright\arraybackslash}p{(\linewidth - 14\tabcolsep) * \real{0.1250}}
  >{\raggedright\arraybackslash}p{(\linewidth - 14\tabcolsep) * \real{0.1250}}
  >{\raggedright\arraybackslash}p{(\linewidth - 14\tabcolsep) * \real{0.1250}}
  >{\raggedright\arraybackslash}p{(\linewidth - 14\tabcolsep) * \real{0.1250}}
  >{\raggedright\arraybackslash}p{(\linewidth - 14\tabcolsep) * \real{0.1250}}@{}}
\toprule\noalign{}
\begin{minipage}[b]{\linewidth}\raggedright
RUNSPEC
\end{minipage} & \begin{minipage}[b]{\linewidth}\raggedright
GRID
\end{minipage} & \begin{minipage}[b]{\linewidth}\raggedright
EDIT
\end{minipage} & \begin{minipage}[b]{\linewidth}\raggedright
PROPS
\end{minipage} & \begin{minipage}[b]{\linewidth}\raggedright
REGIONS
\end{minipage} & \begin{minipage}[b]{\linewidth}\raggedright
SOLUTION
\end{minipage} & \begin{minipage}[b]{\linewidth}\raggedright
SUMMARY
\end{minipage} & \begin{minipage}[b]{\linewidth}\raggedright
SCHEDULE
\end{minipage} \\
\midrule\noalign{}
\endhead
\bottomrule\noalign{}
\endlastfoot
\end{longtable}

\subsubsection{Description}\label{description-394}

The ALSURFST keyword defines the water-oil surface tension versus
alkaline concentration multipliers for when the alkaline model has been
activated via the ALKALINE keyword in the RUNSPEC section.

This keyword is not supported by OPM Flow but would change the results
if supported so the simulation will be stopped.

\subsection{APIGROUP -- Define API Tracking Number of Grouped Oil PVT
Tables}\label{apigroup-define-api-tracking-number-of-grouped-oil-pvt-tables}

\begin{longtable}[]{@{}
  >{\raggedright\arraybackslash}p{(\linewidth - 14\tabcolsep) * \real{0.1250}}
  >{\raggedright\arraybackslash}p{(\linewidth - 14\tabcolsep) * \real{0.1250}}
  >{\raggedright\arraybackslash}p{(\linewidth - 14\tabcolsep) * \real{0.1250}}
  >{\raggedright\arraybackslash}p{(\linewidth - 14\tabcolsep) * \real{0.1250}}
  >{\raggedright\arraybackslash}p{(\linewidth - 14\tabcolsep) * \real{0.1250}}
  >{\raggedright\arraybackslash}p{(\linewidth - 14\tabcolsep) * \real{0.1250}}
  >{\raggedright\arraybackslash}p{(\linewidth - 14\tabcolsep) * \real{0.1250}}
  >{\raggedright\arraybackslash}p{(\linewidth - 14\tabcolsep) * \real{0.1250}}@{}}
\toprule\noalign{}
\begin{minipage}[b]{\linewidth}\raggedright
RUNSPEC
\end{minipage} & \begin{minipage}[b]{\linewidth}\raggedright
GRID
\end{minipage} & \begin{minipage}[b]{\linewidth}\raggedright
EDIT
\end{minipage} & \begin{minipage}[b]{\linewidth}\raggedright
PROPS
\end{minipage} & \begin{minipage}[b]{\linewidth}\raggedright
REGIONS
\end{minipage} & \begin{minipage}[b]{\linewidth}\raggedright
SOLUTION
\end{minipage} & \begin{minipage}[b]{\linewidth}\raggedright
SUMMARY
\end{minipage} & \begin{minipage}[b]{\linewidth}\raggedright
SCHEDULE
\end{minipage} \\
\midrule\noalign{}
\endhead
\bottomrule\noalign{}
\endlastfoot
\end{longtable}

\subsubsection{Description}\label{description-395}

The APIGROUP keyword defines the maximum number of groups of oil PVT
tables when the API tracking option has been activated via the API
keyword in the RUNSPEC section.

This keyword is not supported by OPM Flow but would change the results
if supported so the simulation will be stopped.

\subsection{AQUCT -- Define Carter-Tracy Analytical
Aquifers}\label{aquct-define-carter-tracy-analytical-aquifers-1}

The AQUCT keyword defines a Carter Tracy aquifer, the properties of the
aquifer, including the Carter-Tracy aquifer influence function
associated with the aquifer, defined by the AQUTAB keyword in the PROPS
section.

See AQUCT -- Define Carter-Tracy Analytical Aquifers in the GRID section
for a full description

\subsection{AQUTAB -- Define Carter-Tracy Aquifer Influence
Functions}\label{aqutab-define-carter-tracy-aquifer-influence-functions}

\begin{longtable}[]{@{}
  >{\raggedright\arraybackslash}p{(\linewidth - 14\tabcolsep) * \real{0.1250}}
  >{\raggedright\arraybackslash}p{(\linewidth - 14\tabcolsep) * \real{0.1250}}
  >{\raggedright\arraybackslash}p{(\linewidth - 14\tabcolsep) * \real{0.1250}}
  >{\raggedright\arraybackslash}p{(\linewidth - 14\tabcolsep) * \real{0.1250}}
  >{\raggedright\arraybackslash}p{(\linewidth - 14\tabcolsep) * \real{0.1250}}
  >{\raggedright\arraybackslash}p{(\linewidth - 14\tabcolsep) * \real{0.1250}}
  >{\raggedright\arraybackslash}p{(\linewidth - 14\tabcolsep) * \real{0.1250}}
  >{\raggedright\arraybackslash}p{(\linewidth - 14\tabcolsep) * \real{0.1250}}@{}}
\toprule\noalign{}
\begin{minipage}[b]{\linewidth}\raggedright
RUNSPEC
\end{minipage} & \begin{minipage}[b]{\linewidth}\raggedright
GRID
\end{minipage} & \begin{minipage}[b]{\linewidth}\raggedright
EDIT
\end{minipage} & \begin{minipage}[b]{\linewidth}\raggedright
PROPS
\end{minipage} & \begin{minipage}[b]{\linewidth}\raggedright
REGIONS
\end{minipage} & \begin{minipage}[b]{\linewidth}\raggedright
SOLUTION
\end{minipage} & \begin{minipage}[b]{\linewidth}\raggedright
SUMMARY
\end{minipage} & \begin{minipage}[b]{\linewidth}\raggedright
SCHEDULE
\end{minipage} \\
\midrule\noalign{}
\endhead
\bottomrule\noalign{}
\endlastfoot
\end{longtable}

\subsubsection{Description}\label{description-396}

The AQUTAB keyword defines additional Carter-Tracy {[}Carter, R. D., and
Tracy, G. W. ``An Improved Method for Calculating Water Influx.''
Transactions of AIME, Vol. 219 (1060), pp 415-417.{]} aquifer functions
to be used in the model. Carter-Tracy representation of the aquifer
influx is via a qw term in the non-linear aquifer influence function
Q(t). It allows the water influx from the aquifer to be represented in
the simulator by assuming that there is a constant water influx rate
over finite time periods. It is derived from the superposition methods
of van Everdingen and Hurst {[}Van Everdingen, A. F., and Hurst, W.
``The Application of the Laplace Transform to Flow Problems in
Reservoirs.'' Transactions of AIME, Vol. 186 (1949), pp.~305-324.{]},
whose superposition methods are not suitable for implementation in
reservoir simulation software, although they are very useful in
interpreting aquifer response. The storage requirements and calculation
complexity of handling the resulting superposition formulae can be
largely eliminated by use of the Carter-Tracy approximate water influx
method.

\begin{longtable}[]{@{}llll@{}}
\toprule\noalign{}
No. & Name & Description & Default \\
\midrule\noalign{}
\endhead
\bottomrule\noalign{}
\endlastfoot
Field & Metric & Laboratory & \\
1 & TD & Dimensionless Time & None \\
dimensionless & dimensionless & dimensionless & \\
2 & PD & Dimensionless Pressure & None \\
dimensionless & dimensionless & dimensionless & \\
Notes: & & & \\
\end{longtable}

\emph{Table 8.18: AQUTAB Keyword Description}

\begin{tcolorbox}[enhanced jigsaw, title=\textcolor{quarto-callout-note-color}{\faInfo}\hspace{0.5em}{Note}, coltitle=black, breakable, toprule=.15mm, colbacktitle=quarto-callout-note-color!10!white, leftrule=.75mm, opacitybacktitle=0.6, colback=white, bottomtitle=1mm, toptitle=1mm, titlerule=0mm, left=2mm, colframe=quarto-callout-note-color-frame, arc=.35mm, rightrule=.15mm, bottomrule=.15mm, opacityback=0]

OPM Flow includes the infinite acting Carter-Tracy aquifer influence
table as a default for table number one; thus data entered on this
keyword starts from table number two.

\end{tcolorbox}

Table 8.19 to Table 8.22 outline the standard finite tables derived by
van Everdingen and Hurst that are used for the Carter-Tracy analytical
aquifers and are taken from Bradely {[}Bradley Howard B., et al.,
Petroleum Engineering Handbook, Society of Petroleum Engineers (June
1989), ISBN:9781555630102.{]} table 38-6 on page 38-12. In the tables rD
is defined as the ratio of the aquifer external radius divided by
hydrocarbon radius, that is:\({r}_{D} = \frac{{r}_{e}}{{r}_{o}}\).

\begin{longtable}[]{@{}
  >{\raggedright\arraybackslash}p{(\linewidth - 16\tabcolsep) * \real{0.1111}}
  >{\raggedright\arraybackslash}p{(\linewidth - 16\tabcolsep) * \real{0.1111}}
  >{\raggedright\arraybackslash}p{(\linewidth - 16\tabcolsep) * \real{0.1111}}
  >{\raggedright\arraybackslash}p{(\linewidth - 16\tabcolsep) * \real{0.1111}}
  >{\raggedright\arraybackslash}p{(\linewidth - 16\tabcolsep) * \real{0.1111}}
  >{\raggedright\arraybackslash}p{(\linewidth - 16\tabcolsep) * \real{0.1111}}
  >{\raggedright\arraybackslash}p{(\linewidth - 16\tabcolsep) * \real{0.1111}}
  >{\raggedright\arraybackslash}p{(\linewidth - 16\tabcolsep) * \real{0.1111}}
  >{\raggedright\arraybackslash}p{(\linewidth - 16\tabcolsep) * \real{0.1111}}@{}}
\toprule\noalign{}
\begin{minipage}[b]{\linewidth}\raggedright
Carter-Tracy Aquifer Influence Functions
\end{minipage} & \begin{minipage}[b]{\linewidth}\raggedright
\end{minipage} & \begin{minipage}[b]{\linewidth}\raggedright
\end{minipage} & \begin{minipage}[b]{\linewidth}\raggedright
\end{minipage} & \begin{minipage}[b]{\linewidth}\raggedright
\end{minipage} & \begin{minipage}[b]{\linewidth}\raggedright
\end{minipage} & \begin{minipage}[b]{\linewidth}\raggedright
\end{minipage} & \begin{minipage}[b]{\linewidth}\raggedright
\end{minipage} & \begin{minipage}[b]{\linewidth}\raggedright
\end{minipage} \\
\midrule\noalign{}
\endhead
\bottomrule\noalign{}
\endlastfoot
No. & rD = 1.5 Dimensionless & rD = 2.0 Dimensionless & rD = 2.5
Dimensionless & rD = 3.0 Dimensionless & & & & \\
tD & pD & tD & pD & tD & pD & tD & pD & \\
1 & 0.0600 & 0.2510 & 0.2200 & 0.4430 & 0.4000 & 0.5650 & 0.5200 &
0.6270 \\
2 & 0.0800 & 0.2880 & 0.2400 & 0.4590 & 0.4200 & 0.5760 & 0.5400 &
0.6360 \\
3 & 0.1000 & 0.3220 & 0.2600 & 0.4760 & 0.4400 & 0.5870 & 0.5600 &
0.6450 \\
4 & 0.1200 & 0.3550 & 0.2800 & 0.4920 & 0.4600 & 0.5980 & 0.6000 &
0.6620 \\
5 & 0.1400 & 0.3870 & 0.3000 & 0.5070 & 0.4800 & 0.6080 & 0.6500 &
0.6830 \\
6 & 0.1600 & 0.4200 & 0.3200 & 0.5220 & 0.5000 & 0.6180 & 0.7000 &
0.7030 \\
7 & 0.1800 & 0.4520 & 0.3400 & 0.5360 & 0.5200 & 0.6280 & 0.7500 &
0.7210 \\
8 & 0.2000 & 0.4840 & 0.3600 & 0.5510 & 0.5400 & 0.6380 & 0.8000 &
0.7400 \\
9 & 0.2200 & 0.5160 & 0.3800 & 0.5650 & 0.5600 & 0.6470 & 0.8500 &
0.7580 \\
10 & 0.2400 & 0.5480 & 0.4000 & 0.5790 & 0.5800 & 0.6570 & 0.9000 &
0.7760 \\
11 & 0.2600 & 0.5800 & 0.4200 & 0.5930 & 0.6000 & 0.6660 & 0.9500 &
0.7910 \\
12 & 0.2800 & 0.6120 & 0.4400 & 0.6070 & 0.6500 & 0.6880 & 1.0000 &
0.8060 \\
13 & 0.3000 & 0.6440 & 0.4600 & 0.6210 & 0.7000 & 0.7100 & 1.2000 &
0.8650 \\
14 & 0.3500 & 0.7240 & 0.4800 & 0.6340 & 0.7500 & 0.7310 & 1.4000 &
0.9200 \\
15 & 0.4000 & 0.8040 & 0.5000 & 0.6480 & 0.8000 & 0.7520 & 1.6000 &
0.9730 \\
16 & 0.4500 & 0.8840 & 0.6000 & 0.7150 & 0.8500 & 0.7720 & 2.0000 &
1.0760 \\
17 & 0.5000 & 0.9640 & 0.7000 & 0.7820 & 0.9000 & 0.7920 & 3.0000 &
1.3280 \\
18 & 0.5500 & 1.0440 & 0.8000 & 0.8490 & 0.9500 & 0.8120 & 4.0000 &
1.5780 \\
19 & 0.6000 & 1.1240 & 0.9000 & 0.9150 & 1.0000 & 0.8320 & 5.0000 &
1.8280 \\
20 & & & 1.0000 & 0.9820 & 2.0000 & 1.2150 & & \\
21 & & & 2.0000 & 1.6490 & 3.0000 & 1.5960 & & \\
22 & & & 3.0000 & 2.3160 & 4.0000 & 1.9770 & & \\
23 & & & 5.0000 & 3.6490 & 5.0000 & 2.3580 & & \\
\end{longtable}

\emph{Table 8.19: Carter-Tracy Aquifer Influence Functions (RD =1.5,
2.0, 2.5 and 3.0)}

\begin{longtable}[]{@{}
  >{\raggedright\arraybackslash}p{(\linewidth - 16\tabcolsep) * \real{0.1111}}
  >{\raggedright\arraybackslash}p{(\linewidth - 16\tabcolsep) * \real{0.1111}}
  >{\raggedright\arraybackslash}p{(\linewidth - 16\tabcolsep) * \real{0.1111}}
  >{\raggedright\arraybackslash}p{(\linewidth - 16\tabcolsep) * \real{0.1111}}
  >{\raggedright\arraybackslash}p{(\linewidth - 16\tabcolsep) * \real{0.1111}}
  >{\raggedright\arraybackslash}p{(\linewidth - 16\tabcolsep) * \real{0.1111}}
  >{\raggedright\arraybackslash}p{(\linewidth - 16\tabcolsep) * \real{0.1111}}
  >{\raggedright\arraybackslash}p{(\linewidth - 16\tabcolsep) * \real{0.1111}}
  >{\raggedright\arraybackslash}p{(\linewidth - 16\tabcolsep) * \real{0.1111}}@{}}
\toprule\noalign{}
\begin{minipage}[b]{\linewidth}\raggedright
Carter-Tracy Aquifer Influence Functions
\end{minipage} & \begin{minipage}[b]{\linewidth}\raggedright
\end{minipage} & \begin{minipage}[b]{\linewidth}\raggedright
\end{minipage} & \begin{minipage}[b]{\linewidth}\raggedright
\end{minipage} & \begin{minipage}[b]{\linewidth}\raggedright
\end{minipage} & \begin{minipage}[b]{\linewidth}\raggedright
\end{minipage} & \begin{minipage}[b]{\linewidth}\raggedright
\end{minipage} & \begin{minipage}[b]{\linewidth}\raggedright
\end{minipage} & \begin{minipage}[b]{\linewidth}\raggedright
\end{minipage} \\
\midrule\noalign{}
\endhead
\bottomrule\noalign{}
\endlastfoot
No. & rD = 3.5 Dimensionless & rD = 4.0 Dimensionless & rD = 4.5
Dimensionless & rD = 5.0 Dimensionless & & & & \\
tD & pD & tD & pD & tD & pD & tD & pD & \\
1 & 1.0000 & 0.8020 & 1.5000 & 0.9270 & 2.0000 & 1.0230 & 3.0000 &
1.1670 \\
2 & 1.1000 & 0.8300 & 1.6000 & 0.9480 & 2.1000 & 1.0400 & 3.1000 &
1.1800 \\
3 & 1.2000 & 0.8570 & 1.7000 & 0.9680 & 2.2000 & 1.0560 & 3.2000 &
1.1920 \\
4 & 1.3000 & 0.8820 & 1.8000 & 0.9880 & 2.3000 & 1.0720 & 3.3000 &
1.2040 \\
5 & 1.4000 & 0.9060 & 1.9000 & 1.0070 & 2.4000 & 1.0870 & 3.4000 &
1.2150 \\
6 & 1.5000 & 0.9290 & 2.0000 & 1.0250 & 2.5000 & 1.1020 & 3.5000 &
1.2270 \\
7 & 1.6000 & 0.9510 & 2.2000 & 1.0590 & 2.6000 & 1.1160 & 3.6000 &
1.2380 \\
8 & 1.7000 & 0.9730 & 2.4000 & 1.0920 & 2.7000 & 1.1300 & 3.7000 &
1.2490 \\
9 & 1.8000 & 0.9940 & 2.6000 & 1.1230 & 2.8000 & 1.1440 & 3.8000 &
1.2590 \\
10 & 1.9000 & 1.0140 & 2.8000 & 1.1540 & 2.9000 & 1.1580 & 3.9000 &
1.2700 \\
11 & 2.0000 & 1.0340 & 3.0000 & 1.1840 & 3.0000 & 1.1710 & 4.0000 &
1.2810 \\
12 & 2.2500 & 1.0830 & 3.5000 & 1.2550 & 3.2000 & 1.1970 & 4.2000 &
1.3010 \\
13 & 2.5000 & 1.1300 & 4.0000 & 1.3240 & 3.4000 & 1.2220 & 4.4000 &
1.3210 \\
14 & 2.7500 & 1.1760 & 4.5000 & 1.3920 & 3.6000 & 1.2460 & 4.6000 &
1.3400 \\
15 & 3.0000 & 1.2210 & 5.0000 & 1.4600 & 3.8000 & 1.2690 & 4.8000 &
1.3600 \\
16 & 4.0000 & 1.4010 & 5.5000 & 1.5270 & 4.0000 & 1.2920 & 5.0000 &
1.3780 \\
17 & 5.0000 & 1.5790 & 6.0000 & 1.5940 & 4.5000 & 1.3490 & 5.5000 &
1.4240 \\
18 & 6.0000 & 1.7570 & 6.5000 & 1.6600 & 5.0000 & 1.4030 & 6.0000 &
1.4690 \\
19 & & & 7.0000 & 1.7270 & 5.5000 & 1.4570 & 6.5000 & 1.5130 \\
20 & & & 8.0000 & 1.8610 & 6.0000 & 1.5100 & 7.0000 & 1.5560 \\
21 & & & 9.0000 & 1.9940 & 7.0000 & 1.6150 & 7.5000 & 1.5980 \\
22 & & & 10.0000 & 2.1270 & 8.0000 & 1.7190 & 8.0000 & 1.6410 \\
23 & & & & & 9.0000 & 1.8230 & 9.0000 & 1.7250 \\
24 & & & & & 10.0000 & 1.9270 & 10.0000 & 1.8080 \\
25 & & & & & 11.0000 & 2.0310 & 11.0000 & 1.8920 \\
26 & & & & & 12.0000 & 2.1350 & 12.0000 & 1.9750 \\
27 & & & & & 13.0000 & 2.2390 & 13.0000 & 2.0590 \\
28 & & & & & 14.0000 & 2.3430 & 14.0000 & 2.1420 \\
29 & & & & & 15.0000 & 2.4470 & 15.0000 & 2.2250 \\
\end{longtable}

\emph{Table 8.20: Carter-Tracy Aquifer Influence Functions (RD =3.5,
4.0, 4.5 and 5.0)}

\begin{longtable}[]{@{}
  >{\raggedright\arraybackslash}p{(\linewidth - 16\tabcolsep) * \real{0.1111}}
  >{\raggedright\arraybackslash}p{(\linewidth - 16\tabcolsep) * \real{0.1111}}
  >{\raggedright\arraybackslash}p{(\linewidth - 16\tabcolsep) * \real{0.1111}}
  >{\raggedright\arraybackslash}p{(\linewidth - 16\tabcolsep) * \real{0.1111}}
  >{\raggedright\arraybackslash}p{(\linewidth - 16\tabcolsep) * \real{0.1111}}
  >{\raggedright\arraybackslash}p{(\linewidth - 16\tabcolsep) * \real{0.1111}}
  >{\raggedright\arraybackslash}p{(\linewidth - 16\tabcolsep) * \real{0.1111}}
  >{\raggedright\arraybackslash}p{(\linewidth - 16\tabcolsep) * \real{0.1111}}
  >{\raggedright\arraybackslash}p{(\linewidth - 16\tabcolsep) * \real{0.1111}}@{}}
\toprule\noalign{}
\begin{minipage}[b]{\linewidth}\raggedright
Carter-Tracy Aquifer Influence Functions
\end{minipage} & \begin{minipage}[b]{\linewidth}\raggedright
\end{minipage} & \begin{minipage}[b]{\linewidth}\raggedright
\end{minipage} & \begin{minipage}[b]{\linewidth}\raggedright
\end{minipage} & \begin{minipage}[b]{\linewidth}\raggedright
\end{minipage} & \begin{minipage}[b]{\linewidth}\raggedright
\end{minipage} & \begin{minipage}[b]{\linewidth}\raggedright
\end{minipage} & \begin{minipage}[b]{\linewidth}\raggedright
\end{minipage} & \begin{minipage}[b]{\linewidth}\raggedright
\end{minipage} \\
\midrule\noalign{}
\endhead
\bottomrule\noalign{}
\endlastfoot
No. & rD = 6.0 Dimensionless & rD = 7.0 Dimensionless & rD = 8.0
Dimensionless & rD = 9.0 Dimensionless & & & & \\
tD & pD & tD & pD & tD & pD & tD & pD & \\
1 & 4.0000 & 1.2750 & 6.0000 & 1.4360 & 8.0000 & 1.5560 & 10.0000 &
1.6510 \\
2 & 4.5000 & 1.3220 & 6.5000 & 1.4700 & 8.5000 & 1.5820 & 10.5000 &
1.6730 \\
3 & 5.0000 & 1.3640 & 7.0000 & 1.5010 & 9.0000 & 1.6070 & 11.0000 &
1.6930 \\
4 & 5.5000 & 1.4040 & 7.5000 & 1.5310 & 9.5000 & 1.6310 & 11.5000 &
1.7130 \\
5 & 6.0000 & 1.4410 & 8.0000 & 1.5590 & 10.0000 & 1.6530 & 12.0000 &
1.7320 \\
6 & 6.5000 & 1.4770 & 8.5000 & 1.5860 & 10.5000 & 1.6750 & 12.5000 &
1.7500 \\
7 & 7.0000 & 1.5110 & 9.0000 & 1.6130 & 11.0000 & 1.6970 & 13.0000 &
1.7680 \\
8 & 7.5000 & 1.5440 & 9.5000 & 1.6380 & 11.5000 & 1.7170 & 13.5000 &
1.7860 \\
9 & 8.0000 & 1.5760 & 10.0000 & 1.6630 & 12.0000 & 1.7370 & 14.0000 &
1.8030 \\
10 & 8.5000 & 1.6070 & 11.0000 & 1.7110 & 12.5000 & 1.7570 & 14.5000 &
1.8190 \\
11 & 9.0000 & 1.6380 & 12.0000 & 1.7570 & 13.0000 & 1.7760 & 15.0000 &
1.8350 \\
12 & 9.5000 & 1.6680 & 13.0000 & 1.8010 & 13.5000 & 1.7950 & 15.5000 &
1.8510 \\
13 & 10.0000 & 1.6980 & 14.0000 & 1.8450 & 14.0000 & 1.8130 & 16.0000 &
1.8670 \\
14 & 11.0000 & 1.7570 & 15.0000 & 1.8880 & 14.5000 & 1.8310 & 17.0000 &
1.8970 \\
15 & 12.0000 & 1.8150 & 16.0000 & 1.9310 & 15.0000 & 1.8490 & 18.0000 &
1.9260 \\
16 & 13.0000 & 1.8730 & 17.0000 & 1.9740 & 17.0000 & 1.9190 & 19.0000 &
1.9550 \\
17 & 14.0000 & 1.9310 & 18.0000 & 2.0160 & 19.0000 & 1.9860 & 20.0000 &
1.9830 \\
18 & 15.0000 & 1.9880 & 19.0000 & 2.0580 & 21.0000 & 2.0510 & 22.0000 &
2.0370 \\
19 & 16.0000 & 2.0450 & 20.0000 & 2.1000 & 23.0000 & 2.1160 & 24.0000 &
2.0900 \\
20 & 17.0000 & 2.1030 & 22.0000 & 2.1840 & 25.0000 & 2.1800 & 26.0000 &
2.1420 \\
21 & 18.0000 & 2.1600 & 24.0000 & 2.2670 & 30.0000 & 2.3400 & 28.0000 &
2.1930 \\
22 & 19.0000 & 2.2170 & 26.0000 & 2.3510 & 35.0000 & 2.4990 & 30.0000 &
2.2440 \\
23 & 20.0000 & 2.2740 & 28.0000 & 2.4340 & 40.0000 & 2.6580 & 34.0000 &
2.3450 \\
24 & 25.0000 & 2.5600 & 30.0000 & 2.5170 & 45.0000 & 2.8170 & 38.0000 &
2.4460 \\
25 & 30.0000 & 2.8460 & & & & & 40.0000 & 2.4960 \\
26 & & & & & & & 45.0000 & 2.6210 \\
27 & & & & & & & 50.0000 & 2.7460 \\
\end{longtable}

\emph{Table 8.21: Carter-Tracy Aquifer Influence Functions (RD =6.0,
7.0, 8.0 and 9.0)}

\begin{longtable}[]{@{}lllll@{}}
\toprule\noalign{}
Carter-Tracy Aquifer Influence Functions & & & & \\
\midrule\noalign{}
\endhead
\bottomrule\noalign{}
\endlastfoot
No. & rD = 10.0 Dimensionless & Finite Linear Dimensionless & & \\
tD & pD & tD & pD & \\
1 & 12.0000 & 1.7320 & 0.0050 & 0.0798 \\
2 & 12.5000 & 1.7500 & 0.0100 & 0.1130 \\
3 & 13.0000 & 1.7680 & 0.0200 & 0.1596 \\
4 & 13.5000 & 1.7840 & 0.0600 & 0.2764 \\
5 & 14.0000 & 1.8010 & 0.0800 & 0.3192 \\
6 & 14.5000 & 1.8170 & 0.1000 & 0.3568 \\
7 & 15.0000 & 1.8320 & 0.1200 & 0.3909 \\
8 & 15.5000 & 1.8470 & 0.1600 & 0.4515 \\
9 & 16.0000 & 1.8620 & 0.2000 & 0.5052 \\
10 & 17.0000 & 1.8900 & 0.2400 & 0.5544 \\
11 & 18.0000 & 1.9170 & 0.3000 & 0.6228 \\
12 & 19.0000 & 1.9430 & 0.4000 & 0.7294 \\
13 & 20.0000 & 1.9680 & 0.6000 & 0.9328 \\
14 & 22.0000 & 2.0170 & 0.8000 & 1.1333 \\
15 & 24.0000 & 2.0630 & 1.0000 & 1.3333 \\
16 & 26.0000 & 2.1080 & & \\
17 & 28.0000 & 2.1510 & & \\
18 & 30.0000 & 2.1940 & & \\
19 & 32.0000 & 2.2360 & & \\
20 & 34.0000 & 2.2780 & & \\
21 & 36.0000 & 2.3190 & & \\
22 & 38.0000 & 2.3600 & & \\
23 & 40.0000 & 2.4010 & & \\
24 & 50.0000 & 2.6040 & & \\
25 & 60.0000 & 2.8060 & & \\
26 & 70.0000 & 3.0080 & & \\
\end{longtable}

\emph{Table 8.22: Carter-Tracy Aquifer Influence Functions (RD =10 and
Finite Linear)}

For the finite linear Carter-Tracy influence function in Table 8.22 set
the inner radius of the aquifer to the length of linear aquifer and the
angle of influence to\$ =
\frac{360\times \mathit{Width}}{(2\times π\times \mathit{Length})}\$on
the AQUCT keyword in the grid section.

For reference, Table 8.23 outlines the content terminal rate case for an
infinite aquifer derived by van Everdingen and Hurst, that is the
default table number one used for the Carter-Tracy analytical aquifers,
as implemented in OPM Flow.

\begin{longtable}[]{@{}
  >{\raggedright\arraybackslash}p{(\linewidth - 10\tabcolsep) * \real{0.1667}}
  >{\raggedright\arraybackslash}p{(\linewidth - 10\tabcolsep) * \real{0.1667}}
  >{\raggedright\arraybackslash}p{(\linewidth - 10\tabcolsep) * \real{0.1667}}
  >{\raggedright\arraybackslash}p{(\linewidth - 10\tabcolsep) * \real{0.1667}}
  >{\raggedright\arraybackslash}p{(\linewidth - 10\tabcolsep) * \real{0.1667}}
  >{\raggedright\arraybackslash}p{(\linewidth - 10\tabcolsep) * \real{0.1667}}@{}}
\toprule\noalign{}
\begin{minipage}[b]{\linewidth}\raggedright
Carter-Tracy Infinite Radial Aquifer Influence Function (Default)
\end{minipage} & \begin{minipage}[b]{\linewidth}\raggedright
\end{minipage} & \begin{minipage}[b]{\linewidth}\raggedright
\end{minipage} & \begin{minipage}[b]{\linewidth}\raggedright
\end{minipage} & \begin{minipage}[b]{\linewidth}\raggedright
\end{minipage} & \begin{minipage}[b]{\linewidth}\raggedright
\end{minipage} \\
\midrule\noalign{}
\endhead
\bottomrule\noalign{}
\endlastfoot
No. & Infinite Dimensionless & No. & Infinite Dimensionless & & \\
tD & pD & t & pD & & \\
1 & 1.0 x 10-2 & 0,1120 & 27 & 20,00 & 1.9600 \\
2 & 5.0 x 10-2 & 0,2290 & 28 & 25,00 & 2.0670 \\
3 & 1.0 x 10-1 & 0,3150 & 29 & 30,00 & 2.1470 \\
4 & 1.5 x 10-1 & 0,3760 & 30 & 40,00 & 2.2820 \\
5 & 2.0 x 10-1 & 0,4240 & 31 & 50,00 & 2.3880 \\
6 & 2.5 x 10-1 & 0,4690 & 32 & 60,00 & 2.4760 \\
7 & 3.0 x 10-1 & 0,5030 & 33 & 70,00 & 2.5500 \\
8 & 4.0 x 10-1 & 0,5640 & 34 & 80,00 & 2.6150 \\
9 & 5.0 x 10-1 & 0,6160 & 35 & 90,00 & 2.6720 \\
10 & 6.0 x 10-1 & 0,6590 & 36 & 100,00 & 2.7230 \\
11 & 7.0 x 10-1 & 0,7020 & 37 & 150,00 & 2.9210 \\
12 & 8.0 x 10-1 & 0,7350 & 38 & 200,00 & 3.0640 \\
13 & 9.0 x 10-1 & 0,7720 & 39 & 250,00 & 3.1730 \\
14 & 1.0 & 0,8020 & 40 & 300,00 & 3.2630 \\
15 & 1.5 & 0,9270 & 41 & 400,00 & 3.4060 \\
16 & 2.0 & 1,0200 & 42 & 500,00 & 3.5160 \\
17 & 2.5 & 1,1010 & 43 & 600,00 & 3.6080 \\
18 & 3.0 & 1,1690 & 44 & 700,00 & 3.6840 \\
19 & 4.0 & 1,2750 & 45 & 800,00 & 3.7500 \\
20 & 5.0 & 1,3620 & 46 & 900,00 & 3.8090 \\
21 & 6.0 & 1,4360 & 47 & 1000,00 & 3.8600 \\
22 & 7.0 & 1,5000 & 48 & 2000,00 & 4.1801 \\
23 & 8.0 & 1,5560 & 49 & 5000,00 & 4.6350 \\
24 & 9.0 & 1,6040 & 50 & 8000,00 & 4.8732 \\
25 & 10.0 & 1,6510 & 51 & 10000,00 & 4.9815 \\
26 & 15.0 & 1,8290 & & & \\
\end{longtable}

\emph{Table 8.23: Carter-Tracy Infinite Radial Aquifer Influence
Function (Default)}

For an overview of analytical aquifers see Dake {[}Dake, L.P.
Fundamentals of Reservoir Engineering, Amsterdam, The Netherlands,
Elsevier Science BV (1978) Chapter 9.{]}.

\subsubsection{Example}\label{example-259}

\begin{verbatim}
--
--       CARTER-TRACY AQUIFER INFLUENCE TABLES
--       (STARTS FROM TABLE NO. 2, AS DEFAULT IS TABLE NO. 1)
--
AQUTAB
--
--         TD          PD
--       -------    ---------
          0.06        0.251
          0.08        0.288
          0.10        0.322
          0.12        0.355
          0.14        0.387
          0.16        0.420
          0.18        0.452
          0.20        0.484
          0.22        0.516
          0.24        0.548
          0.26        0.580
          0.28        0.612
          0.30        0.644
          0.35        0.724
          0.40        0.804
          0.45        0.884
          0.50        0.964
          0.55        1.044
          0.60        1.124                                / RD=1.5 TABLE NO. 02
--
--         TD          PD
--       -------    ---------
          0.22        0.443
          0.24        0.459
          0.26        0.476
          0.28        0.492
          0.30        0.507
          0.32        0.522
          0.34        0.536
          0.36        0.551
          0.38        0.565
          0.40        0.579
          0.42        0.593
          0.44        0.607
          0.46        0.621
          0.48        0.634
          0.50        0.648
          0.6         0.715
          0.7         0.782
          0.8         0.849
          0.9         0.915
          1.0         0.982
          2.0         1.649
          3.0         2.316
          5.0         3.649                                / RD=2.0 TABLE NO. 03
\end{verbatim}

The above example defines tables two and three Carter-Tracy aquifer
influence tables.

\subsection{BDENSITY -- Define the Surface Brine Density for the
Fluid}\label{bdensity-define-the-surface-brine-density-for-the-fluid}

\begin{longtable}[]{@{}
  >{\raggedright\arraybackslash}p{(\linewidth - 14\tabcolsep) * \real{0.1250}}
  >{\raggedright\arraybackslash}p{(\linewidth - 14\tabcolsep) * \real{0.1250}}
  >{\raggedright\arraybackslash}p{(\linewidth - 14\tabcolsep) * \real{0.1250}}
  >{\raggedright\arraybackslash}p{(\linewidth - 14\tabcolsep) * \real{0.1250}}
  >{\raggedright\arraybackslash}p{(\linewidth - 14\tabcolsep) * \real{0.1250}}
  >{\raggedright\arraybackslash}p{(\linewidth - 14\tabcolsep) * \real{0.1250}}
  >{\raggedright\arraybackslash}p{(\linewidth - 14\tabcolsep) * \real{0.1250}}
  >{\raggedright\arraybackslash}p{(\linewidth - 14\tabcolsep) * \real{0.1250}}@{}}
\toprule\noalign{}
\begin{minipage}[b]{\linewidth}\raggedright
RUNSPEC
\end{minipage} & \begin{minipage}[b]{\linewidth}\raggedright
GRID
\end{minipage} & \begin{minipage}[b]{\linewidth}\raggedright
EDIT
\end{minipage} & \begin{minipage}[b]{\linewidth}\raggedright
PROPS
\end{minipage} & \begin{minipage}[b]{\linewidth}\raggedright
REGIONS
\end{minipage} & \begin{minipage}[b]{\linewidth}\raggedright
SOLUTION
\end{minipage} & \begin{minipage}[b]{\linewidth}\raggedright
SUMMARY
\end{minipage} & \begin{minipage}[b]{\linewidth}\raggedright
SCHEDULE
\end{minipage} \\
\midrule\noalign{}
\endhead
\bottomrule\noalign{}
\endlastfoot
\end{longtable}

\subsubsection{Description}\label{description-397}

BDENSITY defines the brine surface density for when the brine phase has
been activated in the model by the BRINE keyword in the RUNSPEC section.
The number of BDENSITY vector data sets is defined by the NTPVT
parameter on the TABDIMS keyword in the RUNSPEC section. Each record
consists of a maximum of NPPVT values, as declared on the TABDIMS
keyword in the RUNSPEC section, with each value representing a brine
surface density.

This keyword is not supported by OPM Flow but would change the results
if supported so the simulation will be stopped.

The keyword is used in conjunction with the PVTWSALT keyword in the
PROPS section, with each brine density value matching with the salt
concentration values in column 1 of each table in the PVTWSALT keyword.
Note that the BDENSITY keyword is optional, and if absent from the input
file, then the brine surface densities will be set to the water density
values declared via the DENSITY keyword in the PROPS section. In this
case there is no variation in brine surface density with respect to salt
concentration.

\begin{longtable}[]{@{}
  >{\raggedright\arraybackslash}p{(\linewidth - 10\tabcolsep) * \real{0.1667}}
  >{\raggedright\arraybackslash}p{(\linewidth - 10\tabcolsep) * \real{0.1667}}
  >{\raggedright\arraybackslash}p{(\linewidth - 10\tabcolsep) * \real{0.1667}}
  >{\raggedright\arraybackslash}p{(\linewidth - 10\tabcolsep) * \real{0.1667}}
  >{\raggedright\arraybackslash}p{(\linewidth - 10\tabcolsep) * \real{0.1667}}
  >{\raggedright\arraybackslash}p{(\linewidth - 10\tabcolsep) * \real{0.1667}}@{}}
\toprule\noalign{}
\begin{minipage}[b]{\linewidth}\raggedright
No.
\end{minipage} & \begin{minipage}[b]{\linewidth}\raggedright
Name
\end{minipage} & \begin{minipage}[b]{\linewidth}\raggedright
Description
\end{minipage} & \begin{minipage}[b]{\linewidth}\raggedright
Default
\end{minipage} & \begin{minipage}[b]{\linewidth}\raggedright
\end{minipage} & \begin{minipage}[b]{\linewidth}\raggedright
\end{minipage} \\
\midrule\noalign{}
\endhead
\bottomrule\noalign{}
\endlastfoot
1 & WATDEN & Field & Metric & Laboratory & None \\
WATDEN is a real monotonically increasing positive row vector that
defines the brine density at surface conditions for the given salt
concentrations on the corresponding PVTWSALT keyword in the PROPS
section. There should be one row element for each salt concentration
columnar element (SALTCON) on the PVTWSALT keyword. & & & & & \\
lb/ft3 & kg/m & gm/cc & & & \\
Notes: & & & & & \\
\end{longtable}

\emph{Table 8.24: BDENSITY Keyword Description}

\begin{tcolorbox}[enhanced jigsaw, title=\textcolor{quarto-callout-note-color}{\faInfo}\hspace{0.5em}{Note}, coltitle=black, breakable, toprule=.15mm, colbacktitle=quarto-callout-note-color!10!white, leftrule=.75mm, opacitybacktitle=0.6, colback=white, bottomtitle=1mm, toptitle=1mm, titlerule=0mm, left=2mm, colframe=quarto-callout-note-color-frame, arc=.35mm, rightrule=.15mm, bottomrule=.15mm, opacityback=0]

In OPM Flow the tracer equations are solved decoupled from the reservoir
equations at the end of a time step. For each tracer an implicit system
is solved, however, the tracer equations are linear, resulting in
converge in two iterations. However, the Brine phase is solved fully
implicitly and is fully coupled with the other flow equations. This is
different to the commercial simulator, where the tracer equations are
solved explicitly after the flow equations have converged at the end of
a time step. This can lead to numerical instabilities if there are large
variations in brine densities.

\end{tcolorbox}

\subsubsection{Example}\label{example-260}

The following shows the BDENSITY and PVTWSALT keywords for when NTPVT on
the TABDIMS keyword in the RUNSPEC section is set equal to two and NPPVT
is set to greater than four on the TABDIMS keyword.

\begin{verbatim}
--
--       BRINE WATER DENSITY DATA FOR PVTWSALT KEYWORD
--
--       SALTCON  SALTCON  SALTCON  SALTCON  SALTCON                                          --       DENSITY  DENSITY  DENSITY  DENSITY  DENSITY
--       -------  -------  -------  -------  -------
BDENSITY
         62.20    63.50    64.75    65.90                / FOR PVTWSALT TABLE 1
         64.00    65.50    67.00                         / FOR PVTWSALT TABLE 2           --
--       WATER SALT PVT TABLE
--
PVTWSALT
--       REF PRES  REF SALT
--       PSIA      LB/STB
--       --------  --------
         4500.0    0.000                                 / TABLE NO. REF. DATA
--
--       SALTCONC  BW         CW        VISC     VISC
--       LB/STB    RB/STB     1/PSIA    CPOISE   GRAD
--       --------  --------   -------   ------   ------
            0.0    1.020      2.7E-6    0.370    0.0
            2.0    1.010      2.7E-6    0.370    0.0
            4.0    1.000      2.7E-6    0.370    0.0
           10.0    0.950      2.7E-6    0.370    0.0     / TABLE NO. 01 SALT DATA
--
--       REF PRES  REF SALT
--       PSIA      LB/STB
--       --------  --------
         4000.0    0.000                                 / TABLE NO. 02 REF. DATA
--
--       SALTCONC  BW         CW        VISC     VISC
--       LB/STB    RB/STB     1/PSIA    CPOISE   GRAD
--       --------  --------   -------   ------   ------
            0.0    1.005      2.5E-6    0.320    0.0
            6.0    0.985      2.5E-6    0.320    0.0
           12.0    0.930      2.5E-6    0.320    0.0     / TABLE NO. 02 SALT DATA

\end{verbatim}

\subsection{BGGI -- Define Gi Gas Formation Volume Factor Pressure
Tables}\label{bggi-define-gi-gas-formation-volume-factor-pressure-tables}

\begin{longtable}[]{@{}
  >{\raggedright\arraybackslash}p{(\linewidth - 14\tabcolsep) * \real{0.1250}}
  >{\raggedright\arraybackslash}p{(\linewidth - 14\tabcolsep) * \real{0.1250}}
  >{\raggedright\arraybackslash}p{(\linewidth - 14\tabcolsep) * \real{0.1250}}
  >{\raggedright\arraybackslash}p{(\linewidth - 14\tabcolsep) * \real{0.1250}}
  >{\raggedright\arraybackslash}p{(\linewidth - 14\tabcolsep) * \real{0.1250}}
  >{\raggedright\arraybackslash}p{(\linewidth - 14\tabcolsep) * \real{0.1250}}
  >{\raggedright\arraybackslash}p{(\linewidth - 14\tabcolsep) * \real{0.1250}}
  >{\raggedright\arraybackslash}p{(\linewidth - 14\tabcolsep) * \real{0.1250}}@{}}
\toprule\noalign{}
\begin{minipage}[b]{\linewidth}\raggedright
RUNSPEC
\end{minipage} & \begin{minipage}[b]{\linewidth}\raggedright
GRID
\end{minipage} & \begin{minipage}[b]{\linewidth}\raggedright
EDIT
\end{minipage} & \begin{minipage}[b]{\linewidth}\raggedright
PROPS
\end{minipage} & \begin{minipage}[b]{\linewidth}\raggedright
REGIONS
\end{minipage} & \begin{minipage}[b]{\linewidth}\raggedright
SOLUTION
\end{minipage} & \begin{minipage}[b]{\linewidth}\raggedright
SUMMARY
\end{minipage} & \begin{minipage}[b]{\linewidth}\raggedright
SCHEDULE
\end{minipage} \\
\midrule\noalign{}
\endhead
\bottomrule\noalign{}
\endlastfoot
\end{longtable}

\subsubsection{Description}\label{description-398}

The BGGI keyword defines Gi gas formation volume factor as a function of
Gi and pressure for when the Gi option has been invoked via the GIMODEL
keyword in the RUNSPEC section.

This keyword is not supported by OPM Flow but would change the results
if supported so the simulation will be stopped.

The accuracy of gas condensate and volatile oil modeling using a
black-oil reservoir simulator is now firmly accepted in the industry;
both depletion and gas cycling above the saturation point can be modeled
and yield an acceptable match compared with compositional simulators.
The main deficiency with the existing black-oil formulation is the
treatment of gas injection below the saturation point, where the
compositional effect of the stripping of liquid components is inversely
proportional to their molecular weights, is not considered. This is
because the black-oil model assumes that the saturated hydrocarbon fluid
properties are only functions of pressure. Thus, when dry gas is
injected into a condensate reservoir below the saturation pressure the
injected gas continues to re-vaporize liquid at a rate governed only by
the cell pressures. Compositional modeling indicate that this not the
case.

The GI option attempts to overcome the limitation of the standard
black-oil approach by extending the black-oil model using the method of
Cook et al.~{[}Cook, R. E., Jacoby, R. H., and Ramesh, A. B. ``A
Beta-type Reservoir Simulator for Approximating Compositional Effects
During Gas Injection'' paper SPE 4272, Society of Petroleum Engineers
Journal (1974) 14, No.~5, 471-481.{]} to take into account the fluid
property changes occurring during gas injection. This is done by
extending the fluid property treatment so that the saturated fluid
properties depend on pressure (as per the standard black-oil
formulation) and also on an additional parameter which characterizes the
compositional changes in the reservoir liquid and vapor phases at
constant pressure.

See also the PVTG and GINODE keywords in the PROPS section.

\subsection{BIC -- Binary Interaction
Coefficients}\label{bic-binary-interaction-coefficients}

\begin{longtable}[]{@{}
  >{\raggedright\arraybackslash}p{(\linewidth - 14\tabcolsep) * \real{0.1250}}
  >{\raggedright\arraybackslash}p{(\linewidth - 14\tabcolsep) * \real{0.1250}}
  >{\raggedright\arraybackslash}p{(\linewidth - 14\tabcolsep) * \real{0.1250}}
  >{\raggedright\arraybackslash}p{(\linewidth - 14\tabcolsep) * \real{0.1250}}
  >{\raggedright\arraybackslash}p{(\linewidth - 14\tabcolsep) * \real{0.1250}}
  >{\raggedright\arraybackslash}p{(\linewidth - 14\tabcolsep) * \real{0.1250}}
  >{\raggedright\arraybackslash}p{(\linewidth - 14\tabcolsep) * \real{0.1250}}
  >{\raggedright\arraybackslash}p{(\linewidth - 14\tabcolsep) * \real{0.1250}}@{}}
\toprule\noalign{}
\begin{minipage}[b]{\linewidth}\raggedright
RUNSPEC
\end{minipage} & \begin{minipage}[b]{\linewidth}\raggedright
GRID
\end{minipage} & \begin{minipage}[b]{\linewidth}\raggedright
EDIT
\end{minipage} & \begin{minipage}[b]{\linewidth}\raggedright
PROPS
\end{minipage} & \begin{minipage}[b]{\linewidth}\raggedright
REGIONS
\end{minipage} & \begin{minipage}[b]{\linewidth}\raggedright
SOLUTION
\end{minipage} & \begin{minipage}[b]{\linewidth}\raggedright
SUMMARY
\end{minipage} & \begin{minipage}[b]{\linewidth}\raggedright
SCHEDULE
\end{minipage} \\
\midrule\noalign{}
\endhead
\bottomrule\noalign{}
\endlastfoot
\end{longtable}

\subsubsection{Description}\label{description-399}

The \hyperref[REF_HEADING_KEYWORD_BIC]{BIC} keyword defines the binary
interaction coefficients for each pair of compositional components
active in the model and for each equation of state. The keyword should
only be used if the compositional mode has been requested using the
COMPS keyword in the RUNSPEC section.

OPM Flow does not currently support the general compositional modeling
formulation.

This keyword is not supported by OPM Flow but it will be parsed and its
data ignored.

\begin{longtable}[]{@{}
  >{\raggedright\arraybackslash}p{(\linewidth - 6\tabcolsep) * \real{0.2500}}
  >{\raggedright\arraybackslash}p{(\linewidth - 6\tabcolsep) * \real{0.2500}}
  >{\raggedright\arraybackslash}p{(\linewidth - 6\tabcolsep) * \real{0.2500}}
  >{\raggedright\arraybackslash}p{(\linewidth - 6\tabcolsep) * \real{0.2500}}@{}}
\toprule\noalign{}
\begin{minipage}[b]{\linewidth}\raggedright
No.
\end{minipage} & \begin{minipage}[b]{\linewidth}\raggedright
Name
\end{minipage} & \begin{minipage}[b]{\linewidth}\raggedright
Description
\end{minipage} & \begin{minipage}[b]{\linewidth}\raggedright
Default
\end{minipage} \\
\midrule\noalign{}
\endhead
\bottomrule\noalign{}
\endlastfoot
1 & \hyperref[REF_HEADING_KEYWORD_BIC]{BIC} & A series of real numbers
that define the binary interaction coefficients \({k}_{\mathit{ij}}\)
for each pair (i, j) of compositional components active in the model.
The matrix \({k}_{\mathit{ij}}\) is symmetrical with zeroes on the
diagonal so only the portion of the matrix below the diagonal is
required (i.e., i = \{2, \ldots{} COMPS\} and j \textless{} i), where
COMPS is specified by the COMPS keyword in the RUNSPEC section. The
values are ordered with j cycling fastest (i.e.~\({k}_{21}\),
\({k}_{31}\), \({k}_{32}\), \({k}_{41}\), \({k}_{42}\), \({k}_{43}\),
\ldots). & 0 \\
Notes: & & & \\
\end{longtable}

\emph{Table 8.3.17.1: \hyperref[REF_HEADING_KEYWORD_BIC]{BIC} Keyword
Description}

\subsubsection{Examples}\label{examples-21}

The following example defines the binary interaction coefficients for
each pair of compositional components in a single three-component
equation of state model.

\begin{verbatim}
--
-- Binary Interaction Coefficients
--
BIC
  0.1209
  0.1310  0.0339 /
\end{verbatim}

The following example defines the binary interaction coefficients for
each component in two three-component equation of state models.

\begin{verbatim}
--
-- Binary Interaction Coefficients
--
BIC
  0.1209
  0.1310  0.0339 /
  0.1209
  0.1310  0.0339 /
\end{verbatim}

\subsection{BIOFPARA -- Define Properties for Models with
Biofilms}\label{biofpara-define-properties-for-models-with-biofilms}

\begin{longtable}[]{@{}
  >{\raggedright\arraybackslash}p{(\linewidth - 14\tabcolsep) * \real{0.1250}}
  >{\raggedright\arraybackslash}p{(\linewidth - 14\tabcolsep) * \real{0.1250}}
  >{\raggedright\arraybackslash}p{(\linewidth - 14\tabcolsep) * \real{0.1250}}
  >{\raggedright\arraybackslash}p{(\linewidth - 14\tabcolsep) * \real{0.1250}}
  >{\raggedright\arraybackslash}p{(\linewidth - 14\tabcolsep) * \real{0.1250}}
  >{\raggedright\arraybackslash}p{(\linewidth - 14\tabcolsep) * \real{0.1250}}
  >{\raggedright\arraybackslash}p{(\linewidth - 14\tabcolsep) * \real{0.1250}}
  >{\raggedright\arraybackslash}p{(\linewidth - 14\tabcolsep) * \real{0.1250}}@{}}
\toprule\noalign{}
\begin{minipage}[b]{\linewidth}\raggedright
RUNSPEC
\end{minipage} & \begin{minipage}[b]{\linewidth}\raggedright
GRID
\end{minipage} & \begin{minipage}[b]{\linewidth}\raggedright
EDIT
\end{minipage} & \begin{minipage}[b]{\linewidth}\raggedright
PROPS
\end{minipage} & \begin{minipage}[b]{\linewidth}\raggedright
REGIONS
\end{minipage} & \begin{minipage}[b]{\linewidth}\raggedright
SOLUTION
\end{minipage} & \begin{minipage}[b]{\linewidth}\raggedright
SUMMARY
\end{minipage} & \begin{minipage}[b]{\linewidth}\raggedright
SCHEDULE
\end{minipage} \\
\midrule\noalign{}
\endhead
\bottomrule\noalign{}
\endlastfoot
\end{longtable}

\subsubsection{Description}\label{description-400}

The \hyperref[REF_HEADING_KEYWORD_BIOFPARA]{BIOFPARA} keyword defines
parameters for models including biofilms. For the parameters, biomass
means both suspended microbes in the water phase (labeled as microbial)
and biofilm. Currently the two available models with biofilm effects are
the \hyperref[REF_HEADING_KEYWORD_BIOFILM]{BIOFILM} and the MICP model.
See Landa-Marbán et al {[}Landa-Marbán, D., Tveit, S., Kumar, K., Gasda,
S.E., 2021. Practical approaches to study microbially induced calcite
precipitation at the field scale. Int. J. Greenh. Gas Control 106,
103256. https://doi.org/10.1016/j.ijggc.2021.103256.{]} and
{[}Landa-Marbán, D., Kumar, K., Tveit, S., Gasda, S.E., 2021. Numerical
studies of CO2 leakage remediation by micp-based plugging technology.
In: Røkke, N.A. and Knuutila, H.K. (Eds) Short Papers from the 11th
International Trondheim CCS conference, ISBN: 978-82-536-1714-5,
284-290.{]} for further information on the MICP model parameters, which
are also used in the \hyperref[REF_HEADING_KEYWORD_BIOFILM]{BIOFILM}
model.

\begin{longtable}[]{@{}
  >{\raggedright\arraybackslash}p{(\linewidth - 6\tabcolsep) * \real{0.2500}}
  >{\raggedright\arraybackslash}p{(\linewidth - 6\tabcolsep) * \real{0.2500}}
  >{\raggedright\arraybackslash}p{(\linewidth - 6\tabcolsep) * \real{0.2500}}
  >{\raggedright\arraybackslash}p{(\linewidth - 6\tabcolsep) * \real{0.2500}}@{}}
\toprule\noalign{}
\begin{minipage}[b]{\linewidth}\raggedright
No.
\end{minipage} & \begin{minipage}[b]{\linewidth}\raggedright
Name
\end{minipage} & \begin{minipage}[b]{\linewidth}\raggedright
Description
\end{minipage} & \begin{minipage}[b]{\linewidth}\raggedright
Default
\end{minipage} \\
\midrule\noalign{}
\endhead
\bottomrule\noalign{}
\endlastfoot
Field & Metric & Laboratory & \\
1 & DENSBIOF & A real positive value that defines the density of the
biofilm. & None \\
lb/rft3 & kg/rm3 & gm/rcc & \\
2 & DEATRATE & A real value that defines the biomass death rate. &
None \\
1/days & 1/days & 1/hours & \\
3 & GROWRATE & A real positive value that defines the maximum specific
biomass growth rate. & None \\
1/days & 1/days & 1/hours & \\
4 & HALFGROW & A real positive value that defines the half-velocity
coefficient in the Monod equation for the biomass growth rate. & None \\
lb/rft3 & kg/rm3 & gm/rcc & \\
5 & YIELGROW & A real value that defines the yield biomass growth
coefficient. & None \\
dimensionless & dimensionless & dimensionless & \\
6 & OXYGSUBS & A real value that defines the mass ratio of oxygen
consumed to substrate used for biomass growth. & None \\
dimensionless & dimensionless & dimensionless & \\
7 & ATTARATE & A real positive value that defines the microbial
attachment rate. & None \\
1/days & 1/days & 1/hours & \\
8 & DETARATE & A real positive value that defines the biofilm detachment
rate. & None \\
1/days & 1/days & 1/hours & \\
9 & DETAEXPO & A real value that defines the exponent in the norm for
the water velocity in the detachment term (this value was set to 0.58 in
the
\href{https://www.sciencedirect.com/science/article/pii/S1750583621000086}{MICP
publication}). & None \\
dimensionless & dimensionless & dimensionless & \\
10 & UREARATE & A real positive value that defines the maximum rate of
urea utilization. & None \\
1/days & 1/days & 1/hours & \\
11 & HALFUREA & A real positive value that defines the half-velocity
coefficient in the Monod equation for the urea utilization rate. &
None \\
lb/rft3 & kg/rm3 & gm/rcc & \\
12 & DENSCALC & A real positive value that defines the calcite density.
& None \\
lb/rft3 & kg/rm3 & gm/rcc & \\
13 & YIELURCA & A real value that defines the yield coefficient in the
calcite precipitation term (units of produced calcite over units of urea
utilization). & None \\
dimensionless & dimensionless & dimensionless & \\
Notes: & & & \\
\end{longtable}

\emph{Table 8.87: BIOFPARA Keyword Description}

\begin{tcolorbox}[enhanced jigsaw, title=\textcolor{quarto-callout-note-color}{\faInfo}\hspace{0.5em}{Note}, coltitle=black, breakable, toprule=.15mm, colbacktitle=quarto-callout-note-color!10!white, leftrule=.75mm, opacitybacktitle=0.6, colback=white, bottomtitle=1mm, toptitle=1mm, titlerule=0mm, left=2mm, colframe=quarto-callout-note-color-frame, arc=.35mm, rightrule=.15mm, bottomrule=.15mm, opacityback=0]

This is an OPM Flow specific keyword.

\end{tcolorbox}

\subsubsection{Example}\label{example-261}

The example below is based on metric units, with NTSFUN equal to two on
the TABDIMS keyword.

\begin{verbatim}
--
--       DEFINE BIOMASS PARAMETERS FOR THE MICP MODEL
--
--       DENS    DEAT    GROW    HALF    YIEL    OXYG
--       BIOF    RATE    RATE    GROW    GROW    SUBS
--       ------  ------  ------  ------  ------  ------  ------
--       ATTA    DETA    DETA    UREA    HALF    DENS    YIEL
--       RATE    RATE    EXPO    RATE    UREA    CALC    URCA
--       ------  ------  ------  ------  ------  ------  ------
BIOFPARA
         35      0.0275  3.6     2E-05   0.5     0.5
         0.0735  2.6E-13 0.58    1391    21.3    2710    1.67  /
         10      0.0275  3.6     2E-05   0.5     0.5
         0.0735  1.1E-05 0.58    1391    21.3    2710    1.67  /
\end{verbatim}

\subsection{BOGI -- Define Gi Oil Formation Volume Factor Pressure
Tables}\label{bogi-define-gi-oil-formation-volume-factor-pressure-tables}

\begin{longtable}[]{@{}
  >{\raggedright\arraybackslash}p{(\linewidth - 14\tabcolsep) * \real{0.1250}}
  >{\raggedright\arraybackslash}p{(\linewidth - 14\tabcolsep) * \real{0.1250}}
  >{\raggedright\arraybackslash}p{(\linewidth - 14\tabcolsep) * \real{0.1250}}
  >{\raggedright\arraybackslash}p{(\linewidth - 14\tabcolsep) * \real{0.1250}}
  >{\raggedright\arraybackslash}p{(\linewidth - 14\tabcolsep) * \real{0.1250}}
  >{\raggedright\arraybackslash}p{(\linewidth - 14\tabcolsep) * \real{0.1250}}
  >{\raggedright\arraybackslash}p{(\linewidth - 14\tabcolsep) * \real{0.1250}}
  >{\raggedright\arraybackslash}p{(\linewidth - 14\tabcolsep) * \real{0.1250}}@{}}
\toprule\noalign{}
\begin{minipage}[b]{\linewidth}\raggedright
RUNSPEC
\end{minipage} & \begin{minipage}[b]{\linewidth}\raggedright
GRID
\end{minipage} & \begin{minipage}[b]{\linewidth}\raggedright
EDIT
\end{minipage} & \begin{minipage}[b]{\linewidth}\raggedright
PROPS
\end{minipage} & \begin{minipage}[b]{\linewidth}\raggedright
REGIONS
\end{minipage} & \begin{minipage}[b]{\linewidth}\raggedright
SOLUTION
\end{minipage} & \begin{minipage}[b]{\linewidth}\raggedright
SUMMARY
\end{minipage} & \begin{minipage}[b]{\linewidth}\raggedright
SCHEDULE
\end{minipage} \\
\midrule\noalign{}
\endhead
\bottomrule\noalign{}
\endlastfoot
\end{longtable}

\subsubsection{Description}\label{description-401}

The BOGI keyword defines Gi oil formation volume factor as a function of
Gi and pressure for when the Gi option has been invoked via the GIMODEL
keyword in the RUNSPEC section.

This keyword is not supported by OPM Flow but would change the results
if supported so the simulation will be stopped.

The accuracy of gas condensate and volatile oil modeling using a
black-oil reservoir simulator is now firmly accepted in the industry;
both depletion and gas cycling above the saturation point can be modeled
and yield an acceptable match compared with compositional simulators.
The main deficiency with the existing black-oil formulation is the
treatment of gas injection below the saturation point, where the
compositional effect of the stripping of liquid components is inversely
proportional to their molecular weights, is not considered. This is
because the black-oil model assumes that the saturated hydrocarbon fluid
properties are only functions of pressure. Thus, when dry gas is
injected into a condensate reservoir below the saturation pressure the
injected gas continues to re-vaporize liquid at a rate governed only by
the cell pressures. Compositional modeling indicate that this not the
case.

The GI option attempts to overcome the limitation of the standard
black-oil approach by extending the black-oil model using the method of
Cook et al.~{[}Cook, R. E., Jacoby, R. H., and Ramesh, A. B. ``A
Beta-type Reservoir Simulator for Approximating Compositional Effects
During Gas Injection'' paper SPE 4272, Society of Petroleum Engineers
Journal (1974) 14, No.~5, 471-481.{]} to take into account the fluid
property changes occurring during gas injection. This is done by
extending the fluid property treatment so that the saturated fluid
properties depend on pressure (as per the standard black-oil
formulation) and also on an additional parameter which characterizes the
compositional changes in the reservoir liquid and vapor phases at
constant pressure.

See also the PVTG and GINODE keywords in the PROPS section.

\subsection{BOX -- Define a Range of Grid Blocks to Enter Property
Data}\label{box-define-a-range-of-grid-blocks-to-enter-property-data-2}

BOX defines a range of grid blocks for which subsequent data is assigned
for all the cells in the defined BOX. Note that the BOX grid is reset by
the keyword ENDBOX by resetting the current defined BOX to be the whole
grid. The keyword can be used for any array and for all grid types.

See BOX - Define a Range of Grid Blocks to Enter Property Data in the
GRID section for a full description.

\subsection{CNAMES -- Define Compositional Component
Names}\label{cnames-define-compositional-component-names}

\begin{longtable}[]{@{}
  >{\raggedright\arraybackslash}p{(\linewidth - 14\tabcolsep) * \real{0.1250}}
  >{\raggedright\arraybackslash}p{(\linewidth - 14\tabcolsep) * \real{0.1250}}
  >{\raggedright\arraybackslash}p{(\linewidth - 14\tabcolsep) * \real{0.1250}}
  >{\raggedright\arraybackslash}p{(\linewidth - 14\tabcolsep) * \real{0.1250}}
  >{\raggedright\arraybackslash}p{(\linewidth - 14\tabcolsep) * \real{0.1250}}
  >{\raggedright\arraybackslash}p{(\linewidth - 14\tabcolsep) * \real{0.1250}}
  >{\raggedright\arraybackslash}p{(\linewidth - 14\tabcolsep) * \real{0.1250}}
  >{\raggedright\arraybackslash}p{(\linewidth - 14\tabcolsep) * \real{0.1250}}@{}}
\toprule\noalign{}
\begin{minipage}[b]{\linewidth}\raggedright
RUNSPEC
\end{minipage} & \begin{minipage}[b]{\linewidth}\raggedright
GRID
\end{minipage} & \begin{minipage}[b]{\linewidth}\raggedright
EDIT
\end{minipage} & \begin{minipage}[b]{\linewidth}\raggedright
PROPS
\end{minipage} & \begin{minipage}[b]{\linewidth}\raggedright
REGIONS
\end{minipage} & \begin{minipage}[b]{\linewidth}\raggedright
SOLUTION
\end{minipage} & \begin{minipage}[b]{\linewidth}\raggedright
SUMMARY
\end{minipage} & \begin{minipage}[b]{\linewidth}\raggedright
SCHEDULE
\end{minipage} \\
\midrule\noalign{}
\endhead
\bottomrule\noalign{}
\endlastfoot
\end{longtable}

\subsubsection{Description}\label{description-402}

The CNAMES keyword defines the names for each of the compositional
components active in the model. The keyword should only be used if the
compositional mode has been requested using the COMPS keyword in the
RUNSPEC section.

OPM Flow does not currently support the general compositional modeling
formulation.

\begin{tcolorbox}[enhanced jigsaw, title=\textcolor{quarto-callout-note-color}{\faInfo}\hspace{0.5em}{Note}, coltitle=black, breakable, toprule=.15mm, colbacktitle=quarto-callout-note-color!10!white, leftrule=.75mm, opacitybacktitle=0.6, colback=white, bottomtitle=1mm, toptitle=1mm, titlerule=0mm, left=2mm, colframe=quarto-callout-note-color-frame, arc=.35mm, rightrule=.15mm, bottomrule=.15mm, opacityback=0]

This keyword is only supported by OPM Flow when the two phase gas-water
CO2 storage model has been activated using the CO2STORE keyword and
either the GASWAT or the GAS and WATER keywords in the RUNSPEC section.
Only the component names ``H2O'', ``CO2'' and ``NACL'' (water, CO2 and
salt respectively) are recognized when the CNAMES keyword is used with
the CO2STORE keyword; any other component names are ignored.

\end{tcolorbox}

\begin{longtable}[]{@{}
  >{\raggedright\arraybackslash}p{(\linewidth - 6\tabcolsep) * \real{0.2500}}
  >{\raggedright\arraybackslash}p{(\linewidth - 6\tabcolsep) * \real{0.2500}}
  >{\raggedright\arraybackslash}p{(\linewidth - 6\tabcolsep) * \real{0.2500}}
  >{\raggedright\arraybackslash}p{(\linewidth - 6\tabcolsep) * \real{0.2500}}@{}}
\toprule\noalign{}
\begin{minipage}[b]{\linewidth}\raggedright
No.
\end{minipage} & \begin{minipage}[b]{\linewidth}\raggedright
Name
\end{minipage} & \begin{minipage}[b]{\linewidth}\raggedright
Description
\end{minipage} & \begin{minipage}[b]{\linewidth}\raggedright
Default
\end{minipage} \\
\midrule\noalign{}
\endhead
\bottomrule\noalign{}
\endlastfoot
1 & CNAMES & A series of character strings of up to eight characters in
length that define the names for each of the compositional components
active in the model. & None \\
Notes: & & & \\
\end{longtable}

\emph{Table 8.25: CNAMES Keyword Description}

\subsubsection{Example}\label{example-262}

The following example defines how to confirm a three component
formulation, together with defining the names of the compositional
components, to be used with the CO2STORE and GASWAT options.

\begin{verbatim}
--
--       CONFIRM NUMBER OF COMPOSITIONAL COMPONENTS (OPM FLOW KEYWORD)
--
NCOMPS
                3                                                                      /
--
--       DEFINE COMPOSITIONAL COMPONENTS NAMES (OPM FLOW KEYWORD)
--
CNAMES
         'H2O'
         'CO2'
         'NACL'                                                                 /
\end{verbatim}

\subsection{COALADS -- Define Gas and Solvent Relative Adsorption
Tables}\label{coalads-define-gas-and-solvent-relative-adsorption-tables}

\begin{longtable}[]{@{}
  >{\raggedright\arraybackslash}p{(\linewidth - 14\tabcolsep) * \real{0.1250}}
  >{\raggedright\arraybackslash}p{(\linewidth - 14\tabcolsep) * \real{0.1250}}
  >{\raggedright\arraybackslash}p{(\linewidth - 14\tabcolsep) * \real{0.1250}}
  >{\raggedright\arraybackslash}p{(\linewidth - 14\tabcolsep) * \real{0.1250}}
  >{\raggedright\arraybackslash}p{(\linewidth - 14\tabcolsep) * \real{0.1250}}
  >{\raggedright\arraybackslash}p{(\linewidth - 14\tabcolsep) * \real{0.1250}}
  >{\raggedright\arraybackslash}p{(\linewidth - 14\tabcolsep) * \real{0.1250}}
  >{\raggedright\arraybackslash}p{(\linewidth - 14\tabcolsep) * \real{0.1250}}@{}}
\toprule\noalign{}
\begin{minipage}[b]{\linewidth}\raggedright
RUNSPEC
\end{minipage} & \begin{minipage}[b]{\linewidth}\raggedright
GRID
\end{minipage} & \begin{minipage}[b]{\linewidth}\raggedright
EDIT
\end{minipage} & \begin{minipage}[b]{\linewidth}\raggedright
PROPS
\end{minipage} & \begin{minipage}[b]{\linewidth}\raggedright
REGIONS
\end{minipage} & \begin{minipage}[b]{\linewidth}\raggedright
SOLUTION
\end{minipage} & \begin{minipage}[b]{\linewidth}\raggedright
SUMMARY
\end{minipage} & \begin{minipage}[b]{\linewidth}\raggedright
SCHEDULE
\end{minipage} \\
\midrule\noalign{}
\endhead
\bottomrule\noalign{}
\endlastfoot
\end{longtable}

\subsubsection{Description}\label{description-403}

The COALADS keyword defines the gas and solvent relative adsorption
tables for when the coal phase has been activated via the COAL keyword
in the RUNSPEC section.

This keyword is not supported by OPM Flow but would change the results
if supported so the simulation will be stopped.

\subsection{COALPP -- Define Gas and Solvent Partial Pressure Adsorption
Tables}\label{coalpp-define-gas-and-solvent-partial-pressure-adsorption-tables}

\begin{longtable}[]{@{}
  >{\raggedright\arraybackslash}p{(\linewidth - 14\tabcolsep) * \real{0.1250}}
  >{\raggedright\arraybackslash}p{(\linewidth - 14\tabcolsep) * \real{0.1250}}
  >{\raggedright\arraybackslash}p{(\linewidth - 14\tabcolsep) * \real{0.1250}}
  >{\raggedright\arraybackslash}p{(\linewidth - 14\tabcolsep) * \real{0.1250}}
  >{\raggedright\arraybackslash}p{(\linewidth - 14\tabcolsep) * \real{0.1250}}
  >{\raggedright\arraybackslash}p{(\linewidth - 14\tabcolsep) * \real{0.1250}}
  >{\raggedright\arraybackslash}p{(\linewidth - 14\tabcolsep) * \real{0.1250}}
  >{\raggedright\arraybackslash}p{(\linewidth - 14\tabcolsep) * \real{0.1250}}@{}}
\toprule\noalign{}
\begin{minipage}[b]{\linewidth}\raggedright
RUNSPEC
\end{minipage} & \begin{minipage}[b]{\linewidth}\raggedright
GRID
\end{minipage} & \begin{minipage}[b]{\linewidth}\raggedright
EDIT
\end{minipage} & \begin{minipage}[b]{\linewidth}\raggedright
PROPS
\end{minipage} & \begin{minipage}[b]{\linewidth}\raggedright
REGIONS
\end{minipage} & \begin{minipage}[b]{\linewidth}\raggedright
SOLUTION
\end{minipage} & \begin{minipage}[b]{\linewidth}\raggedright
SUMMARY
\end{minipage} & \begin{minipage}[b]{\linewidth}\raggedright
SCHEDULE
\end{minipage} \\
\midrule\noalign{}
\endhead
\bottomrule\noalign{}
\endlastfoot
\end{longtable}

\subsubsection{Description}\label{description-404}

The COALPP keyword defines the gas and solvent partial pressure
adsorption tables for when the coal phase has been activated via the
COAL keyword in the RUNSPEC section.

This keyword is not supported by OPM Flow but would change the results
if supported so the simulation will be stopped.

\subsection{COLUMNS -- Define Input File Column
Margins}\label{columns-define-input-file-column-margins-4}

The COLUMNS keyword defines the input file column margins; characters
outside the margins are ignored by the input parser.

See COLUMNS -- Define Input File Column Margins in the GLOBAL section
for a full description.

\subsection{COPY -- Copy Array Data to Another
Array}\label{copy-copy-array-data-to-another-array-2}

The COPY keyword copies an array (or part of an array) to another array
or part of an array. The arrays can be real or integer depending on the
array type; however, the arrays that can be operated on is dependent on
which section the COPY keyword is being used.

See COPY -- Copy Array Data to Another Array in the GRID section for a
full description.

\subsection{COPYBOX -- Copy Array Data Defined by a
Box}\label{copybox-copy-array-data-defined-by-a-box-1}

The COPYBOX keyword copies an array (or part of an array) to part of the
same array. The array can be real or integer depending on the array
type; however, the arrays that can be operated on is dependent on which
section the COPYBOX keyword is being used.

See COPYBOX -- Copy Array Data Defined by a Box in the GRID section for
a full description.

\subsection{COPYREG -- Copy an Array to Another Array based on a Region
Number}\label{copyreg-copy-an-array-to-another-array-based-on-a-region-number-2}

The COPYREG keyword copies a specified array or part of an array based
on cells with a specific region number to another array. The region
number array can be FLUXNUM, MULTNUM or OPERNUM and these arrays must be
defined and be available before the COPYREG keyword is read by the
simulator. The property arrays can be real or integer depending on the
property array type; however, the property arrays that can be operated
on is dependent on which section the COPYREG keyword is being used.

See COPYREG -- Copy an Array to Another Array based on a Region Number
in the GRID section for a full description.

\subsection{DEBUG -- Define the Debug Data to be Printed to
File}\label{debug-define-the-debug-data-to-be-printed-to-file-4}

This keyword defines the debug data to be written to the debug file
(*.DBG). This keyword is not supported by OPM Flow but has no effect on
the results so it will be ignored.

See DEBUG -- Define the Debug Data to be Printed to File in the GLOBAL
section for a full description.

\subsection{DENAQA -- Specify Ezrokhi Coefficients for Aqueous
Density}\label{denaqa-specify-ezrokhi-coefficients-for-aqueous-density}

\begin{longtable}[]{@{}
  >{\raggedright\arraybackslash}p{(\linewidth - 14\tabcolsep) * \real{0.1250}}
  >{\raggedright\arraybackslash}p{(\linewidth - 14\tabcolsep) * \real{0.1250}}
  >{\raggedright\arraybackslash}p{(\linewidth - 14\tabcolsep) * \real{0.1250}}
  >{\raggedright\arraybackslash}p{(\linewidth - 14\tabcolsep) * \real{0.1250}}
  >{\raggedright\arraybackslash}p{(\linewidth - 14\tabcolsep) * \real{0.1250}}
  >{\raggedright\arraybackslash}p{(\linewidth - 14\tabcolsep) * \real{0.1250}}
  >{\raggedright\arraybackslash}p{(\linewidth - 14\tabcolsep) * \real{0.1250}}
  >{\raggedright\arraybackslash}p{(\linewidth - 14\tabcolsep) * \real{0.1250}}@{}}
\toprule\noalign{}
\begin{minipage}[b]{\linewidth}\raggedright
RUNSPEC
\end{minipage} & \begin{minipage}[b]{\linewidth}\raggedright
GRID
\end{minipage} & \begin{minipage}[b]{\linewidth}\raggedright
EDIT
\end{minipage} & \begin{minipage}[b]{\linewidth}\raggedright
PROPS
\end{minipage} & \begin{minipage}[b]{\linewidth}\raggedright
REGIONS
\end{minipage} & \begin{minipage}[b]{\linewidth}\raggedright
SOLUTION
\end{minipage} & \begin{minipage}[b]{\linewidth}\raggedright
SUMMARY
\end{minipage} & \begin{minipage}[b]{\linewidth}\raggedright
SCHEDULE
\end{minipage} \\
\midrule\noalign{}
\endhead
\bottomrule\noalign{}
\endlastfoot
\end{longtable}

\subsubsection{Description}\label{description-405}

The \hyperref[REF_HEADING_KEYWORD_DENAQA_8_3]{DENAQA} keyword specifies
the three Ezrokhi coefficients for each compositional component and for
each equation of state that are used to calculate the aqueous phase
density. The keyword should only be used if the compositional mode has
been requested using the COMPS keyword in the RUNSPEC section.

OPM Flow does not currently support the general compositional modeling
formulation.

This keyword is only supported by OPM Flow when the two phase gas-water
CO2 storage model has been activated using the CO2STORE keyword and
either the GASWAT or the GAS and WATER keywords in the RUNSPEC section.
The component names ``H2O'', ``CO2'' and ``NACL'' (water, CO2 and salt
respectively) must be specified using the CNAMES keyword.

\begin{longtable}[]{@{}
  >{\raggedright\arraybackslash}p{(\linewidth - 6\tabcolsep) * \real{0.2500}}
  >{\raggedright\arraybackslash}p{(\linewidth - 6\tabcolsep) * \real{0.2500}}
  >{\raggedright\arraybackslash}p{(\linewidth - 6\tabcolsep) * \real{0.2500}}
  >{\raggedright\arraybackslash}p{(\linewidth - 6\tabcolsep) * \real{0.2500}}@{}}
\toprule\noalign{}
\begin{minipage}[b]{\linewidth}\raggedright
No.
\end{minipage} & \begin{minipage}[b]{\linewidth}\raggedright
Name
\end{minipage} & \begin{minipage}[b]{\linewidth}\raggedright
Description
\end{minipage} & \begin{minipage}[b]{\linewidth}\raggedright
Default
\end{minipage} \\
\midrule\noalign{}
\endhead
\bottomrule\noalign{}
\endlastfoot
Field & Metric & Laboratory & \\
1 & COEFFS & A series of real numbers that define the three Ezrokhi
coeffients for each of the compositional components active in the model.
& None \\
dimensionless & dimensionless & dimensionless & \\
Notes: & & & \\
\end{longtable}

\emph{Table 8.3.28.1: DENAQA Keyword Description}

\subsubsection{Examples}\label{examples-22}

The following example defines the Ezrokhi coefficients for each
component in a single three-component equation of state model.

\begin{verbatim}
--
-- Ezrokhi Coefficients for Aqueous Density Calculation
--
DENAQA
--COEFF0   COEFF1   COEFF2
  2.0E-6   -5.0E-7  2.0E-9
  2.0E-6   -5.0E-7  2.0E-9
  2.0E-6   -5.0E-7  2.0E-9   /
\end{verbatim}

The following example defines the Ezrokhi coefficients for each
component in two three-component equation of state models.

\begin{verbatim}
--
-- Ezrokhi Coefficients for Aqueous Density Calculation
--
DENAQA
--COEFF0   COEFF1   COEFF2
  2.0E-6   -5.0E-7  2.0E-9
  2.0E-6   -5.0E-7  2.0E-9
  2.0E-6   -5.0E-7  2.0E-9   /
--
  2.0E-6   -5.0E-7  2.0E-9
  2.0E-6   -5.0E-7  2.0E-9
  2.0E-6   -5.0E-7  2.0E-9   /
\end{verbatim}

\subsection{DENSITY -- Define the Surface Oil, Water Gas Densities for
the
Fluids}\label{density-define-the-surface-oil-water-gas-densities-for-the-fluids}

\begin{longtable}[]{@{}
  >{\raggedright\arraybackslash}p{(\linewidth - 14\tabcolsep) * \real{0.1250}}
  >{\raggedright\arraybackslash}p{(\linewidth - 14\tabcolsep) * \real{0.1250}}
  >{\raggedright\arraybackslash}p{(\linewidth - 14\tabcolsep) * \real{0.1250}}
  >{\raggedright\arraybackslash}p{(\linewidth - 14\tabcolsep) * \real{0.1250}}
  >{\raggedright\arraybackslash}p{(\linewidth - 14\tabcolsep) * \real{0.1250}}
  >{\raggedright\arraybackslash}p{(\linewidth - 14\tabcolsep) * \real{0.1250}}
  >{\raggedright\arraybackslash}p{(\linewidth - 14\tabcolsep) * \real{0.1250}}
  >{\raggedright\arraybackslash}p{(\linewidth - 14\tabcolsep) * \real{0.1250}}@{}}
\toprule\noalign{}
\begin{minipage}[b]{\linewidth}\raggedright
RUNSPEC
\end{minipage} & \begin{minipage}[b]{\linewidth}\raggedright
GRID
\end{minipage} & \begin{minipage}[b]{\linewidth}\raggedright
EDIT
\end{minipage} & \begin{minipage}[b]{\linewidth}\raggedright
PROPS
\end{minipage} & \begin{minipage}[b]{\linewidth}\raggedright
REGIONS
\end{minipage} & \begin{minipage}[b]{\linewidth}\raggedright
SOLUTION
\end{minipage} & \begin{minipage}[b]{\linewidth}\raggedright
SUMMARY
\end{minipage} & \begin{minipage}[b]{\linewidth}\raggedright
SCHEDULE
\end{minipage} \\
\midrule\noalign{}
\endhead
\bottomrule\noalign{}
\endlastfoot
\end{longtable}

\subsubsection{Description}\label{description-406}

DENSITY defines the oil, water and gas surface densities for the fluids
for various regions in the model. The number of DENSITY vector data sets
is defined by the NTPVT parameter on the TABDIMS keyword in the RUNSPEC
section and the allocation of the DENSITY data sets to different grid
blocks in the model is done via the PVTNUM keyword in the REGIONS
section. One data set consists of one record or line which is terminated
by a ``/''. The surface density or gravity must be entered using either
the DENSITY or GRAVITY keywords irrespective of which phases are active
in the model.

\begin{longtable}[]{@{}
  >{\raggedright\arraybackslash}p{(\linewidth - 6\tabcolsep) * \real{0.2500}}
  >{\raggedright\arraybackslash}p{(\linewidth - 6\tabcolsep) * \real{0.2500}}
  >{\raggedright\arraybackslash}p{(\linewidth - 6\tabcolsep) * \real{0.2500}}
  >{\raggedright\arraybackslash}p{(\linewidth - 6\tabcolsep) * \real{0.2500}}@{}}
\toprule\noalign{}
\begin{minipage}[b]{\linewidth}\raggedright
No.
\end{minipage} & \begin{minipage}[b]{\linewidth}\raggedright
Name
\end{minipage} & \begin{minipage}[b]{\linewidth}\raggedright
Description
\end{minipage} & \begin{minipage}[b]{\linewidth}\raggedright
Default
\end{minipage} \\
\midrule\noalign{}
\endhead
\bottomrule\noalign{}
\endlastfoot
Field & Metric & Laboratory & \\
1 & OILDEN & OILDEN is a real number defining the density of the oil
phase at surface conditions. & Defined \\
lb/ft3 37.457 & kg/m3 600 & gm/cc 0.6 & \\
2 & WATDEN & WATDEN is a real number defining the density of the water
phase at surface conditions. & Defined \\
lb/ft3 62.366 & kg/m3 999.014 & gm/cc 0.999014 & \\
3 & GASDEN & GASDEN is a real number defining the density of the gas
phase at surface conditions. & Defined \\
lb/ft3 0.062428 & kg/m3 1.000 & gm/cc 0.001 & \\
Notes: & & & \\
\end{longtable}

\emph{Table 8.26: DENSITY Keyword Description}

Note that provided the first table has been entered, subsequent tables
may be defaulted, in this case the prior table is copied to the current
table. See the third example for an illustration on how to use this
feature.

According to the SPE SI standard {[}The SI Metric System of Units and
SPE Metric Standard, Adopted for Use as a Voluntary Standard by the SPE
Board of Directors, June 1983, Society of Petroleum Engineers.{]},
Relative Density (γ) replaces Specific Gravity as the term used to
define the ratio of the density of a known material to the density of
reference material, at standard conditions of pressure and temperature.
Standard conditions vary throughout the world, but for oil field units
one normally uses14.7 psia and 60 oF, whereas for SI units some areas
use 101.325 kPa and 15 oC.

See also the GRAVITY keyword in the PROPS section, that can be used to
enter the relative density values instead of the density numbers.

\subsubsection{Examples}\label{examples-23}

The following shows the DENSITY keyword for when NTPVT on the TABDIMS
keyword in the RUNSPEC section is set to one.

\begin{verbatim}
--
--       OIL      WAT        GAS
--       DENSITY  DENSITY    DENSITY
--       -------  -------    -------
DENSITY
         39.0     62.37      0.04520                       / PVT DATA REGION 1
\end{verbatim}

The next example shows the DENSITY keyword for when NTPVT on the TABDIMS
keyword in the RUNSPEC section is set to three.

\begin{verbatim}
--
--       OIL      WAT        GAS
--       DENSITY  DENSITY    DENSITY
--       -------  -------    -------
DENSITY
         38.0     62.30      0.04500                       / PVT DATA REGION 1
         39.0     62.37      0.04520                       / PVT DATA REGION 2
         40.0     62.40      0.04800                       / PVT DATA REGION 3
\end{verbatim}

The third, and final, example shows the DENSITY keyword for when NTPVT
on the TABDIMS keyword in the RUNSPEC section is set to four. Here table
two defaults to table one, and table four defaults to table three.

\begin{verbatim}
--
--       OIL      WAT        GAS
--       DENSITY  DENSITY    DENSITY
--       -------  -------    -------
DENSITY
         38.0     62.30      0.04500                       / PVT DATA REGION 1
                                                           / PVT DATA REGION 2
         39.0     62.37      0.04520                       / PVT DATA REGION 3
                                                           / PVT DATA REGION 3
\end{verbatim}

Again, note that there is no terminating ``/'' for this keyword.

\subsection{DEPTHTAB -- River Time and Depth
Tables}\label{depthtab-river-time-and-depth-tables}

\begin{longtable}[]{@{}
  >{\raggedright\arraybackslash}p{(\linewidth - 14\tabcolsep) * \real{0.1250}}
  >{\raggedright\arraybackslash}p{(\linewidth - 14\tabcolsep) * \real{0.1250}}
  >{\raggedright\arraybackslash}p{(\linewidth - 14\tabcolsep) * \real{0.1250}}
  >{\raggedright\arraybackslash}p{(\linewidth - 14\tabcolsep) * \real{0.1250}}
  >{\raggedright\arraybackslash}p{(\linewidth - 14\tabcolsep) * \real{0.1250}}
  >{\raggedright\arraybackslash}p{(\linewidth - 14\tabcolsep) * \real{0.1250}}
  >{\raggedright\arraybackslash}p{(\linewidth - 14\tabcolsep) * \real{0.1250}}
  >{\raggedright\arraybackslash}p{(\linewidth - 14\tabcolsep) * \real{0.1250}}@{}}
\toprule\noalign{}
\begin{minipage}[b]{\linewidth}\raggedright
RUNSPEC
\end{minipage} & \begin{minipage}[b]{\linewidth}\raggedright
GRID
\end{minipage} & \begin{minipage}[b]{\linewidth}\raggedright
EDIT
\end{minipage} & \begin{minipage}[b]{\linewidth}\raggedright
PROPS
\end{minipage} & \begin{minipage}[b]{\linewidth}\raggedright
REGIONS
\end{minipage} & \begin{minipage}[b]{\linewidth}\raggedright
SOLUTION
\end{minipage} & \begin{minipage}[b]{\linewidth}\raggedright
SUMMARY
\end{minipage} & \begin{minipage}[b]{\linewidth}\raggedright
SCHEDULE
\end{minipage} \\
\midrule\noalign{}
\endhead
\bottomrule\noalign{}
\endlastfoot
\end{longtable}

\subsubsection{Description}\label{description-407}

This keyword, DEPTHTAB, defines the river time and depth tables, for
when the River option has been activated via the RIVRDIMS keyword in the
RUNSPEC section.

This keyword is not supported by OPM Flow but would change the results
if supported so the simulation will be stopped.

\subsection{DIAGDISP -- Activate Alternate Form of Tracer
Dispersion}\label{diagdisp-activate-alternate-form-of-tracer-dispersion}

\begin{longtable}[]{@{}
  >{\raggedright\arraybackslash}p{(\linewidth - 14\tabcolsep) * \real{0.1250}}
  >{\raggedright\arraybackslash}p{(\linewidth - 14\tabcolsep) * \real{0.1250}}
  >{\raggedright\arraybackslash}p{(\linewidth - 14\tabcolsep) * \real{0.1250}}
  >{\raggedright\arraybackslash}p{(\linewidth - 14\tabcolsep) * \real{0.1250}}
  >{\raggedright\arraybackslash}p{(\linewidth - 14\tabcolsep) * \real{0.1250}}
  >{\raggedright\arraybackslash}p{(\linewidth - 14\tabcolsep) * \real{0.1250}}
  >{\raggedright\arraybackslash}p{(\linewidth - 14\tabcolsep) * \real{0.1250}}
  >{\raggedright\arraybackslash}p{(\linewidth - 14\tabcolsep) * \real{0.1250}}@{}}
\toprule\noalign{}
\begin{minipage}[b]{\linewidth}\raggedright
RUNSPEC
\end{minipage} & \begin{minipage}[b]{\linewidth}\raggedright
GRID
\end{minipage} & \begin{minipage}[b]{\linewidth}\raggedright
EDIT
\end{minipage} & \begin{minipage}[b]{\linewidth}\raggedright
PROPS
\end{minipage} & \begin{minipage}[b]{\linewidth}\raggedright
REGIONS
\end{minipage} & \begin{minipage}[b]{\linewidth}\raggedright
SOLUTION
\end{minipage} & \begin{minipage}[b]{\linewidth}\raggedright
SUMMARY
\end{minipage} & \begin{minipage}[b]{\linewidth}\raggedright
SCHEDULE
\end{minipage} \\
\midrule\noalign{}
\endhead
\bottomrule\noalign{}
\endlastfoot
\end{longtable}

\subsubsection{Description}\label{description-408}

This keyword, DIAGDISP, activates the alternate form of tracer
dispersion matrix for when the Tracer facility has been activated by the
TRACERS keyword in the RUNSPEC section.

This keyword is not supported by OPM Flow but would change the results
if supported so the simulation will be stopped.

\subsection{DIFFAGAS -- Define PVT Region Gas Component Diffusion
Coefficients (Mass Fraction
Formulation)}\label{diffagas-define-pvt-region-gas-component-diffusion-coefficients-mass-fraction-formulation}

\begin{longtable}[]{@{}
  >{\raggedright\arraybackslash}p{(\linewidth - 14\tabcolsep) * \real{0.1250}}
  >{\raggedright\arraybackslash}p{(\linewidth - 14\tabcolsep) * \real{0.1250}}
  >{\raggedright\arraybackslash}p{(\linewidth - 14\tabcolsep) * \real{0.1250}}
  >{\raggedright\arraybackslash}p{(\linewidth - 14\tabcolsep) * \real{0.1250}}
  >{\raggedright\arraybackslash}p{(\linewidth - 14\tabcolsep) * \real{0.1250}}
  >{\raggedright\arraybackslash}p{(\linewidth - 14\tabcolsep) * \real{0.1250}}
  >{\raggedright\arraybackslash}p{(\linewidth - 14\tabcolsep) * \real{0.1250}}
  >{\raggedright\arraybackslash}p{(\linewidth - 14\tabcolsep) * \real{0.1250}}@{}}
\toprule\noalign{}
\begin{minipage}[b]{\linewidth}\raggedright
RUNSPEC
\end{minipage} & \begin{minipage}[b]{\linewidth}\raggedright
GRID
\end{minipage} & \begin{minipage}[b]{\linewidth}\raggedright
EDIT
\end{minipage} & \begin{minipage}[b]{\linewidth}\raggedright
PROPS
\end{minipage} & \begin{minipage}[b]{\linewidth}\raggedright
REGIONS
\end{minipage} & \begin{minipage}[b]{\linewidth}\raggedright
SOLUTION
\end{minipage} & \begin{minipage}[b]{\linewidth}\raggedright
SUMMARY
\end{minipage} & \begin{minipage}[b]{\linewidth}\raggedright
SCHEDULE
\end{minipage} \\
\midrule\noalign{}
\endhead
\bottomrule\noalign{}
\endlastfoot
\end{longtable}

\subsubsection{Description}\label{description-409}

The \hyperref[REF_HEADING_KEYWORD_DIFFAGAS_8_3]{DIFFAGAS} keyword
defines the gas diffusion coefficients assuming a mass fraction
formulation for each compositional component in the model and for each
PVT region, for when the Molecular Diffusion option has been activated
by the DIFFUSE keyword in the RUNSPEC section.

This keyword is optional as OPM Flow will automatically calculate the
coefficients assuming the mole fraction formulation, as described by
Sandve et al.~{[}Tor Harald Sandve, Sarah E. Gasda, Atgeirr Rasmussen,
and Alf Birger Rustad. Convective dissolution in field scale CO2 storage
simulation using the OPM Flow simulator. Submitted to TCCS 11 --
Trondheim Conference on CO2 Capture, Transport and Storage Trondheim,
Norway -- June 21-23, 2021.{]}, if the
\hyperref[REF_HEADING_KEYWORD_DIFFAGAS_8_3]{DIFFAGAS} and DIFFCGAS
keywords are absent from the input deck. The keyword thus allows one to
overwrite the automatically calculated values.

The keyword should only be used if the CO2STORE or
\hyperref[REF_HEADING_KEYWORD_H2STORE]{H2STORE} keywords and either the
GASWAT or the GAS and WATER keywords in the RUNSPEC section, have also
been activated for the gas-water two component model.

See also the DIFFCGAS and DIFFCWAT keywords that assume the standard
mole fraction formulation for diffusion rather than the mass fraction
formulation assumed by the
\hyperref[REF_HEADING_KEYWORD_DIFFAGAS_8_3]{DIFFAGAS} and
\hyperref[REF_HEADING_KEYWORD_DIFFAWAT_8_3]{DIFFAWAT} keywords. The
\hyperref[REF_HEADING_KEYWORD_DIFFAGAS_8_3]{DIFFAGAS} and
\hyperref[REF_HEADING_KEYWORD_DIFFAWAT_8_3]{DIFFAWAT} keywords cannot be
used in combination with the DIFFCGAS and DIFFCWAT keywords.

\begin{tcolorbox}[enhanced jigsaw, title=\textcolor{quarto-callout-note-color}{\faInfo}\hspace{0.5em}{Note}, coltitle=black, breakable, toprule=.15mm, colbacktitle=quarto-callout-note-color!10!white, leftrule=.75mm, opacitybacktitle=0.6, colback=white, bottomtitle=1mm, toptitle=1mm, titlerule=0mm, left=2mm, colframe=quarto-callout-note-color-frame, arc=.35mm, rightrule=.15mm, bottomrule=.15mm, opacityback=0]

This is an OPM Flow keyword used with OPM Flow's CO2STORE or
\hyperref[REF_HEADING_KEYWORD_H2STORE]{H2STORE} and GASWAT keywords in
the RUNSPEC section, and should not be confused with the version of the
\hyperref[REF_HEADING_KEYWORD_DIFFAGAS_8_3]{DIFFAGAS} keyword used in
the commercial compositional simulator that defines activity corrected
diffusion coefficients.

\end{tcolorbox}

\begin{longtable}[]{@{}
  >{\raggedright\arraybackslash}p{(\linewidth - 6\tabcolsep) * \real{0.2500}}
  >{\raggedright\arraybackslash}p{(\linewidth - 6\tabcolsep) * \real{0.2500}}
  >{\raggedright\arraybackslash}p{(\linewidth - 6\tabcolsep) * \real{0.2500}}
  >{\raggedright\arraybackslash}p{(\linewidth - 6\tabcolsep) * \real{0.2500}}@{}}
\toprule\noalign{}
\begin{minipage}[b]{\linewidth}\raggedright
No.
\end{minipage} & \begin{minipage}[b]{\linewidth}\raggedright
Name
\end{minipage} & \begin{minipage}[b]{\linewidth}\raggedright
Description
\end{minipage} & \begin{minipage}[b]{\linewidth}\raggedright
Default
\end{minipage} \\
\midrule\noalign{}
\endhead
\bottomrule\noalign{}
\endlastfoot
Field & Metric & Laboratory & \\
1 & CO2DIFF & A real positive number that declares the CO2 or H2 in gas
diffusion coefficient in the given PVT region. & None \\
ft2/day & m2/day & cm2/hour & \\
2 & WATDIFF & A real positive number that specifies the water in gas
diffusion coefficient in the given PVT region. & None \\
ft2/day & m2/day & cm2/hour & \\
Notes: & & & \\
\end{longtable}

\emph{Table 8.3.31.1: DIFFAGAS Keyword Description}

\begin{tcolorbox}[enhanced jigsaw, title=\textcolor{quarto-callout-note-color}{\faInfo}\hspace{0.5em}{Note}, coltitle=black, breakable, toprule=.15mm, colbacktitle=quarto-callout-note-color!10!white, leftrule=.75mm, opacitybacktitle=0.6, colback=white, bottomtitle=1mm, toptitle=1mm, titlerule=0mm, left=2mm, colframe=quarto-callout-note-color-frame, arc=.35mm, rightrule=.15mm, bottomrule=.15mm, opacityback=0]

The option has been tested in combination with the CO2STORE keyword, but
not for the general case at this point.

\end{tcolorbox}

See also the DIFFUSE keyword in the RUNSPEC section to activate the
Molecular Diffusion option and the
\hyperref[REF_HEADING_KEYWORD_DIFFAWAT_8_3]{DIFFAWAT} keyword in the
PROPS section that defines the water diffusion coefficients for each
compositional component in the model and for each PVT region. Finally,
for gas-oil systems the DIFFC keyword in the PROPS section should be
used.

Normally diffusion coefficients are measured in laboratory units, that
is cm2/s, for ease of use, Table 8.3.31.2 outlines the conversion
factors for converting the laboratory measured diffusion coefficients to
those used by the simulator.

\begin{longtable}[]{@{}lll@{}}
\toprule\noalign{}
Diffusivity Conversion Factors & & \\
\midrule\noalign{}
\endhead
\bottomrule\noalign{}
\endlastfoot
Laboratory Measured Units & Conversion Factor & Simulator Units \\
1 cm2/s & 92.9979 ft2/day & Field \\
8.64 m2/day & Metric & \\
3600 cm2/hour & Laboratory & \\
\end{longtable}

\emph{Table 8.3.31.2: Diffusivity Conversion Factors}

\subsubsection{Example}\label{example-263}

The example below is based on field units, with NTPVT equal to three on
the TABDIMS keyword.

\begin{verbatim}
--
--       PVT REGION GAS COMPONENT DIFFUSION COEFFICIENTS (OPM FLOW KEYWORD)
--
DIFFAGAS
--       CO2 IN   WAT IN
--       GAS DF   GAS DF
--       -------  -------
         0.160    0.150                / PVT REGION NO. 01
         0.165    0.155                / PVT REGION NO. 02
                                       / PVT REGION NO. 03
\end{verbatim}

Here the third PVT region has no values for the two component diffusion
coefficients, and therefore the simulator will use correlations to
define the diffusivity coefficients for this PVT region.

\subsection{DIFFAWAT -- Define PVT Region Water Component Diffusion
Coefficients (Mass Fraction
Formulation)}\label{diffawat-define-pvt-region-water-component-diffusion-coefficients-mass-fraction-formulation}

\begin{longtable}[]{@{}
  >{\raggedright\arraybackslash}p{(\linewidth - 14\tabcolsep) * \real{0.1250}}
  >{\raggedright\arraybackslash}p{(\linewidth - 14\tabcolsep) * \real{0.1250}}
  >{\raggedright\arraybackslash}p{(\linewidth - 14\tabcolsep) * \real{0.1250}}
  >{\raggedright\arraybackslash}p{(\linewidth - 14\tabcolsep) * \real{0.1250}}
  >{\raggedright\arraybackslash}p{(\linewidth - 14\tabcolsep) * \real{0.1250}}
  >{\raggedright\arraybackslash}p{(\linewidth - 14\tabcolsep) * \real{0.1250}}
  >{\raggedright\arraybackslash}p{(\linewidth - 14\tabcolsep) * \real{0.1250}}
  >{\raggedright\arraybackslash}p{(\linewidth - 14\tabcolsep) * \real{0.1250}}@{}}
\toprule\noalign{}
\begin{minipage}[b]{\linewidth}\raggedright
RUNSPEC
\end{minipage} & \begin{minipage}[b]{\linewidth}\raggedright
GRID
\end{minipage} & \begin{minipage}[b]{\linewidth}\raggedright
EDIT
\end{minipage} & \begin{minipage}[b]{\linewidth}\raggedright
PROPS
\end{minipage} & \begin{minipage}[b]{\linewidth}\raggedright
REGIONS
\end{minipage} & \begin{minipage}[b]{\linewidth}\raggedright
SOLUTION
\end{minipage} & \begin{minipage}[b]{\linewidth}\raggedright
SUMMARY
\end{minipage} & \begin{minipage}[b]{\linewidth}\raggedright
SCHEDULE
\end{minipage} \\
\midrule\noalign{}
\endhead
\bottomrule\noalign{}
\endlastfoot
\end{longtable}

\subsubsection{Description}\label{description-410}

The \hyperref[REF_HEADING_KEYWORD_DIFFAWAT_8_3]{DIFFAWAT} keyword
defines the water diffusion coefficients assuming a mass fraction
formulation for each compositional component in the model and for each
PVT region, for when the Molecular Diffusion option has been activated
by the DIFFUSE keyword in the RUNSPEC section.

This keyword is optional as OPM Flow will automatically calculate the
coefficients assuming the mole fraction formulation, as described by
Sandve et al.~{[}Tor Harald Sandve, Sarah E. Gasda, Atgeirr Rasmussen,
and Alf Birger Rustad. Convective dissolution in field scale CO2 storage
simulation using the OPM Flow simulator. Submitted to TCCS 11 --
Trondheim Conference on CO2 Capture, Transport and Storage Trondheim,
Norway -- June 21-23, 2021.{]}, if the
\hyperref[REF_HEADING_KEYWORD_DIFFAWAT_8_3]{DIFFAWAT} and DIFFCWAT
keywords are absent from the input deck. The keyword thus allows one to
overwrite the automatically calculated values.

The keyword should only be used if the CO2STORE or
\hyperref[REF_HEADING_KEYWORD_H2STORE]{H2STORE} keywords and either the
GASWAT or the GAS and WATER keywords in the RUNSPEC section, have also
been activated for the gas-water two component model.

See also the DIFFCGAS and DIFFCWAT keywords that assume the standard
mole fraction formulation for diffusion rather than the mass fraction
formulation assumed by the
\hyperref[REF_HEADING_KEYWORD_DIFFAGAS_8_3]{DIFFAGAS} and
\hyperref[REF_HEADING_KEYWORD_DIFFAWAT_8_3]{DIFFAWAT} keywords. The
\hyperref[REF_HEADING_KEYWORD_DIFFAGAS_8_3]{DIFFAGAS} and
\hyperref[REF_HEADING_KEYWORD_DIFFAWAT_8_3]{DIFFAWAT} keywords cannot be
used in combination with the DIFFCGAS and DIFFCWAT keywords.

\begin{tcolorbox}[enhanced jigsaw, title=\textcolor{quarto-callout-note-color}{\faInfo}\hspace{0.5em}{Note}, coltitle=black, breakable, toprule=.15mm, colbacktitle=quarto-callout-note-color!10!white, leftrule=.75mm, opacitybacktitle=0.6, colback=white, bottomtitle=1mm, toptitle=1mm, titlerule=0mm, left=2mm, colframe=quarto-callout-note-color-frame, arc=.35mm, rightrule=.15mm, bottomrule=.15mm, opacityback=0]

This is an OPM Flow specific keyword used with OPM Flow's CO2STORE or
\hyperref[REF_HEADING_KEYWORD_H2STORE]{H2STORE} and GASWAT keywords in
the RUNSPEC section.

\end{tcolorbox}

\begin{longtable}[]{@{}
  >{\raggedright\arraybackslash}p{(\linewidth - 6\tabcolsep) * \real{0.2500}}
  >{\raggedright\arraybackslash}p{(\linewidth - 6\tabcolsep) * \real{0.2500}}
  >{\raggedright\arraybackslash}p{(\linewidth - 6\tabcolsep) * \real{0.2500}}
  >{\raggedright\arraybackslash}p{(\linewidth - 6\tabcolsep) * \real{0.2500}}@{}}
\toprule\noalign{}
\begin{minipage}[b]{\linewidth}\raggedright
No.
\end{minipage} & \begin{minipage}[b]{\linewidth}\raggedright
Name
\end{minipage} & \begin{minipage}[b]{\linewidth}\raggedright
Description
\end{minipage} & \begin{minipage}[b]{\linewidth}\raggedright
Default
\end{minipage} \\
\midrule\noalign{}
\endhead
\bottomrule\noalign{}
\endlastfoot
Field & Metric & Laboratory & \\
1 & CO2DIFF & A real positive number that declares the CO2 or H2 in
water diffusion coefficient in the given PVT region. & None \\
ft2/day & m2/day & cm2/hour & \\
2 & WATDIFF & A real positive number that specifies the water in water
diffusion coefficient in the given PVT region. & None \\
ft2/day & m2/day & cm2/hour & \\
Notes: & & & \\
\end{longtable}

\emph{Table 8.3.32.1: DIFFAWAT Keyword Description}

\begin{tcolorbox}[enhanced jigsaw, title=\textcolor{quarto-callout-note-color}{\faInfo}\hspace{0.5em}{Note}, coltitle=black, breakable, toprule=.15mm, colbacktitle=quarto-callout-note-color!10!white, leftrule=.75mm, opacitybacktitle=0.6, colback=white, bottomtitle=1mm, toptitle=1mm, titlerule=0mm, left=2mm, colframe=quarto-callout-note-color-frame, arc=.35mm, rightrule=.15mm, bottomrule=.15mm, opacityback=0]

The option has been tested in combination with the CO2STORE keyword, but
not for the general case at this point.

\end{tcolorbox}

See also the DIFFUSE keyword in the RUNSPEC section to activate the
Molecular Diffusion option and the
\hyperref[REF_HEADING_KEYWORD_DIFFAGAS_8_3]{DIFFAGAS} keyword in the
PROPS section that defines the gas diffusion coefficients for each
compositional component in the model and for each PVT region. Finally,
for gas-oil systems the DIFFC keyword in the PROPS section should be
used.

Normally diffusion coefficients are measured in laboratory units, that
is cm2/s, for ease of use, Table 8.3.32.2 outlines the conversion
factors for converting the laboratory measured diffusion coefficients to
those used by the simulator.

\begin{longtable}[]{@{}lll@{}}
\toprule\noalign{}
Diffusivity Conversion Factors & & \\
\midrule\noalign{}
\endhead
\bottomrule\noalign{}
\endlastfoot
Laboratory Measured Units & Conversion Factor & Simulator Units \\
1 cm2/s & 92.9979 ft2/day & Field \\
8.64 m2/day & Metric & \\
3600 cm2/hour & Laboratory & \\
\end{longtable}

\emph{Table 8.3.32.2: Diffusivity Conversion Factors}

\subsubsection{Example}\label{example-264}

The example below is based on field units, with NTPVT equal to three on
the TABDIMS keyword.

\begin{verbatim}
--
--       PVT REGION WATER COMPONENT DIFFUSION COEFFICIENTS (OPM FLOW KEYWORD)
--
DIFFAWAT
--       CO2 IN   WAT IN
--       WAT DF   WAT DF
--       -------  -------
         0.160    0.150                / PVT REGION NO. 01
         0.165    0.155                / PVT REGION NO. 02
                                       / PVT REGION NO. 03
\end{verbatim}

Here the third PVT region has no values for the two component diffusion
coefficients, and therefore the simulator will use correlations to
define the diffusivity coefficients for this PVT region.

\subsection{DIFFC -- Define PVT Region Molecular Diffusion
Tables}\label{diffc-define-pvt-region-molecular-diffusion-tables}

\begin{longtable}[]{@{}
  >{\raggedright\arraybackslash}p{(\linewidth - 14\tabcolsep) * \real{0.1250}}
  >{\raggedright\arraybackslash}p{(\linewidth - 14\tabcolsep) * \real{0.1250}}
  >{\raggedright\arraybackslash}p{(\linewidth - 14\tabcolsep) * \real{0.1250}}
  >{\raggedright\arraybackslash}p{(\linewidth - 14\tabcolsep) * \real{0.1250}}
  >{\raggedright\arraybackslash}p{(\linewidth - 14\tabcolsep) * \real{0.1250}}
  >{\raggedright\arraybackslash}p{(\linewidth - 14\tabcolsep) * \real{0.1250}}
  >{\raggedright\arraybackslash}p{(\linewidth - 14\tabcolsep) * \real{0.1250}}
  >{\raggedright\arraybackslash}p{(\linewidth - 14\tabcolsep) * \real{0.1250}}@{}}
\toprule\noalign{}
\begin{minipage}[b]{\linewidth}\raggedright
RUNSPEC
\end{minipage} & \begin{minipage}[b]{\linewidth}\raggedright
GRID
\end{minipage} & \begin{minipage}[b]{\linewidth}\raggedright
EDIT
\end{minipage} & \begin{minipage}[b]{\linewidth}\raggedright
PROPS
\end{minipage} & \begin{minipage}[b]{\linewidth}\raggedright
REGIONS
\end{minipage} & \begin{minipage}[b]{\linewidth}\raggedright
SOLUTION
\end{minipage} & \begin{minipage}[b]{\linewidth}\raggedright
SUMMARY
\end{minipage} & \begin{minipage}[b]{\linewidth}\raggedright
SCHEDULE
\end{minipage} \\
\midrule\noalign{}
\endhead
\bottomrule\noalign{}
\endlastfoot
\end{longtable}

\subsubsection{Description}\label{description-411}

The DIFFC keyword defines the molecular weight of the fluids and
diffusion coefficients between phases for each PVT region, for when the
Molecular Diffusion option has been activated by the DIFFUSE keyword in
the RUNSPEC section. This keyword is optional as OPM Flow will
automatically calculate the coefficients, as described by Sandve et
al.~{[}Tor Harald Sandve1, Sarah E. Gasda, Atgeirr Rasmussen, and Alf
Birger Rustad. Convective dissolution in field scale CO2 storage
simulation using the OPM Flow simulator. Submitted to TCCS 11 --
Trondheim Conference on CO2 Capture, Transport and Storage Trondheim,
Norway -- June 21-23, 2021.{]}, if the DIFFC keyword is absent from the
input deck. The keyword thus allows one to overwrite the automatically
calculated values.

\begin{longtable}[]{@{}
  >{\raggedright\arraybackslash}p{(\linewidth - 6\tabcolsep) * \real{0.2500}}
  >{\raggedright\arraybackslash}p{(\linewidth - 6\tabcolsep) * \real{0.2500}}
  >{\raggedright\arraybackslash}p{(\linewidth - 6\tabcolsep) * \real{0.2500}}
  >{\raggedright\arraybackslash}p{(\linewidth - 6\tabcolsep) * \real{0.2500}}@{}}
\toprule\noalign{}
\begin{minipage}[b]{\linewidth}\raggedright
No.
\end{minipage} & \begin{minipage}[b]{\linewidth}\raggedright
Name
\end{minipage} & \begin{minipage}[b]{\linewidth}\raggedright
Description
\end{minipage} & \begin{minipage}[b]{\linewidth}\raggedright
Default
\end{minipage} \\
\midrule\noalign{}
\endhead
\bottomrule\noalign{}
\endlastfoot
Field & Metric & Laboratory & \\
1 & OILMW & OILMW is a real positive number that specifies the molecular
weight of the oil in the given PVT region. & None \\
lb/lb-M & kg/kg-M & gm/gm-M & \\
2 & GASMW & GASMW is a real positive number that defines the molecular
weight of the gas in the given PVT region. & None \\
lb/lb-M & kg/kg-M & gm/gm-M & \\
3 & GASGASDF & A real positive number that defines the gas in gas
diffusion coefficient in the given PVT region. & None \\
ft2/day & m2/day & cm2/hour & \\
4 & OILGASDF & A real positive number that declares the oil in gas
diffusion coefficient in the given PVT region. & None \\
ft2/day & m2/day & cm2/hour & \\
5 & GASOILDF & A real positive number that specifies the gas in oil
diffusion coefficient in the given PVT region. & None \\
ft2/day & m2/day & cm2/hour & \\
6 & OILOILDF & A real positive number that defines the oil in oil
diffusion coefficient in the given PVT region. & None \\
ft2/day & m2/day & cm2/hour & \\
7 & GASOILCD & A real positive number that defines the gas in oil cross
phase diffusion coefficient in the given PVT region. This parameter is
ignored by OPM Flow and should be defaulted or set equal to zero. &
0.0 \\
ft2/day & m2/day & cm2/hour & \\
8 & OILOILCD & A real positive number that defines the oil in oil cross
phase diffusion coefficient in the given PVT region. This parameter is
ignored by OPM Flow and should be defaulted or set equal to zero. &
0.0 \\
ft2/day & m2/day & cm2/hour & \\
Notes: & & & \\
\end{longtable}

\emph{Table 8.27: DIFFC Keyword Description}

\begin{tcolorbox}[enhanced jigsaw, title=\textcolor{quarto-callout-note-color}{\faInfo}\hspace{0.5em}{Note}, coltitle=black, breakable, toprule=.15mm, colbacktitle=quarto-callout-note-color!10!white, leftrule=.75mm, opacitybacktitle=0.6, colback=white, bottomtitle=1mm, toptitle=1mm, titlerule=0mm, left=2mm, colframe=quarto-callout-note-color-frame, arc=.35mm, rightrule=.15mm, bottomrule=.15mm, opacityback=0]

The option has been tested in combination with the CO2STORE keyword, but
not for the general case at this point.

\end{tcolorbox}

See also the DIFFUSE keyword in the RUNSPEC section to activate the
Molecular Diffusion option.

\subsubsection{Example}\label{example-265}

The example below is based on field units, with NTPVT equal to three on
the TABDIMS keyword.

\begin{verbatim}
--
--       PVT REGION MOLECULAR DIFFUSION TABLES
--
DIFFC
--       OIL     GAS    GAS IN   OIL IN   GAS IN   OIL IN   GAS IN   OIL IN
--       MW      MW     GAS DF   GAS DF   OIL DF   OIL DF   OIL CD   OIL CD
--       ------  -----  -------  -------  -------  -------  -------  ------
         103.20  1.120  1.35E-6  1.05E-7  4.50E-7  1.05E-8                 /TAB-1
         102.00  1.130  1.25E-6  1.25E-7  4.80E-7  1.05E-8                 /TAB-2
         100.00  1.250  1.22E-6                                            /TAB-3
\end{verbatim}

Here the third PVT region has no values for the various oil related
diffusion coefficients.

\subsection{DIFFCGAS -- Define PVT Region Gas Component Diffusion
Coefficients}\label{diffcgas-define-pvt-region-gas-component-diffusion-coefficients}

\begin{longtable}[]{@{}
  >{\raggedright\arraybackslash}p{(\linewidth - 14\tabcolsep) * \real{0.1250}}
  >{\raggedright\arraybackslash}p{(\linewidth - 14\tabcolsep) * \real{0.1250}}
  >{\raggedright\arraybackslash}p{(\linewidth - 14\tabcolsep) * \real{0.1250}}
  >{\raggedright\arraybackslash}p{(\linewidth - 14\tabcolsep) * \real{0.1250}}
  >{\raggedright\arraybackslash}p{(\linewidth - 14\tabcolsep) * \real{0.1250}}
  >{\raggedright\arraybackslash}p{(\linewidth - 14\tabcolsep) * \real{0.1250}}
  >{\raggedright\arraybackslash}p{(\linewidth - 14\tabcolsep) * \real{0.1250}}
  >{\raggedright\arraybackslash}p{(\linewidth - 14\tabcolsep) * \real{0.1250}}@{}}
\toprule\noalign{}
\begin{minipage}[b]{\linewidth}\raggedright
RUNSPEC
\end{minipage} & \begin{minipage}[b]{\linewidth}\raggedright
GRID
\end{minipage} & \begin{minipage}[b]{\linewidth}\raggedright
EDIT
\end{minipage} & \begin{minipage}[b]{\linewidth}\raggedright
PROPS
\end{minipage} & \begin{minipage}[b]{\linewidth}\raggedright
REGIONS
\end{minipage} & \begin{minipage}[b]{\linewidth}\raggedright
SOLUTION
\end{minipage} & \begin{minipage}[b]{\linewidth}\raggedright
SUMMARY
\end{minipage} & \begin{minipage}[b]{\linewidth}\raggedright
SCHEDULE
\end{minipage} \\
\midrule\noalign{}
\endhead
\bottomrule\noalign{}
\endlastfoot
\end{longtable}

\subsubsection{Description}\label{description-412}

The DIFFCGAS keyword defines the gas diffusion coefficients assuming the
standard mole fraction formulation for each compositional component in
the model and for each PVT region, for when the Molecular Diffusion
option has been activated by the DIFFUSE keyword in the RUNSPEC section.

This keyword is optional as OPM Flow will automatically calculate the
coefficients assuming the mole fraction formulation, as described by
Sandve et al.~{[}Tor Harald Sandve, Sarah E. Gasda, Atgeirr Rasmussen,
and Alf Birger Rustad. Convective dissolution in field scale CO2 storage
simulation using the OPM Flow simulator. Submitted to TCCS 11 --
Trondheim Conference on CO2 Capture, Transport and Storage Trondheim,
Norway -- June 21-23, 2021.{]}, if the
\hyperref[REF_HEADING_KEYWORD_DIFFAGAS_8_3]{DIFFAGAS} and DIFFCGAS
keywords are absent from the input deck. The keyword thus allows one to
overwrite the automatically calculated values.

The keyword should only be used if the CO2STORE or
\hyperref[REF_HEADING_KEYWORD_H2STORE]{H2STORE} keyword and either the
GASWAT or the GAS and WATER keywords in the RUNSPEC section, have also
been activated for the gas-water two component model.

See also the \hyperref[REF_HEADING_KEYWORD_DIFFAGAS_8_3]{DIFFAGAS} and
\hyperref[REF_HEADING_KEYWORD_DIFFAWAT_8_3]{DIFFAWAT} keywords that
assume a mass fraction formulation for diffusion rather than the
standard mole fraction formulation assumed by the DIFFCGAS and DIFFCWAT
keywords. The \hyperref[REF_HEADING_KEYWORD_DIFFAGAS_8_3]{DIFFAGAS} and
\hyperref[REF_HEADING_KEYWORD_DIFFAWAT_8_3]{DIFFAWAT} keywords cannot be
used in combination with the DIFFCGAS and DIFFCWAT keywords.

\begin{tcolorbox}[enhanced jigsaw, title=\textcolor{quarto-callout-note-color}{\faInfo}\hspace{0.5em}{Note}, coltitle=black, breakable, toprule=.15mm, colbacktitle=quarto-callout-note-color!10!white, leftrule=.75mm, opacitybacktitle=0.6, colback=white, bottomtitle=1mm, toptitle=1mm, titlerule=0mm, left=2mm, colframe=quarto-callout-note-color-frame, arc=.35mm, rightrule=.15mm, bottomrule=.15mm, opacityback=0]

This is an OPM Flow keyword used with OPM Flow's CO2STORE or
\hyperref[REF_HEADING_KEYWORD_H2STORE]{H2STORE} and GASWAT keywords in
the RUNSPEC section, and should not be confused with the more general
version of the DIFFCGAS keyword used in the commercial compositional
simulator.

\end{tcolorbox}

\begin{longtable}[]{@{}
  >{\raggedright\arraybackslash}p{(\linewidth - 6\tabcolsep) * \real{0.2500}}
  >{\raggedright\arraybackslash}p{(\linewidth - 6\tabcolsep) * \real{0.2500}}
  >{\raggedright\arraybackslash}p{(\linewidth - 6\tabcolsep) * \real{0.2500}}
  >{\raggedright\arraybackslash}p{(\linewidth - 6\tabcolsep) * \real{0.2500}}@{}}
\toprule\noalign{}
\begin{minipage}[b]{\linewidth}\raggedright
No.
\end{minipage} & \begin{minipage}[b]{\linewidth}\raggedright
Name
\end{minipage} & \begin{minipage}[b]{\linewidth}\raggedright
Description
\end{minipage} & \begin{minipage}[b]{\linewidth}\raggedright
Default
\end{minipage} \\
\midrule\noalign{}
\endhead
\bottomrule\noalign{}
\endlastfoot
Field & Metric & Laboratory & \\
1 & CO2DIFF & A real positive number that declares the CO2 or H2 in gas
diffusion coefficient in the given PVT region. & None \\
ft2/day & m2/day & cm2/hour & \\
2 & WATDIFF & A real positive number that specifies the water in gas
diffusion coefficient in the given PVT region. & None \\
ft2/day & m2/day & cm2/hour & \\
Notes: & & & \\
\end{longtable}

\emph{Table 8.28: DIFFCGAS Keyword Description}

\begin{tcolorbox}[enhanced jigsaw, title=\textcolor{quarto-callout-note-color}{\faInfo}\hspace{0.5em}{Note}, coltitle=black, breakable, toprule=.15mm, colbacktitle=quarto-callout-note-color!10!white, leftrule=.75mm, opacitybacktitle=0.6, colback=white, bottomtitle=1mm, toptitle=1mm, titlerule=0mm, left=2mm, colframe=quarto-callout-note-color-frame, arc=.35mm, rightrule=.15mm, bottomrule=.15mm, opacityback=0]

The option has been tested in combination with the CO2STORE keyword, but
not for the general case at this point.

\end{tcolorbox}

See also the DIFFUSE keyword in the RUNSPEC section to activate the
Molecular Diffusion option and the DIFFCWAT keyword in the PROPS section
that defines the water diffusion coefficients for each compositional
component in the model and for each PVT region. Finally, for gas-oil
systems the DIFFC keyword in the PROPS section should be used.

Normally diffusion coefficients are measured in laboratory units, that
is cm2/s, for ease of use, Table 8.29 outlines the conversion factors
for converting the laboratory measured diffusion coefficients to those
used by the simulator.

\begin{longtable}[]{@{}lll@{}}
\toprule\noalign{}
Diffusivity Conversion Factors & & \\
\midrule\noalign{}
\endhead
\bottomrule\noalign{}
\endlastfoot
Laboratory Measured Units & Conversion Factor & Simulator Units \\
1 cm2/s & 92.9979 ft2/day & Field \\
8.64 m2/day & Metric & \\
3600 cm2/hour & Laboratory & \\
\end{longtable}

\emph{Table 8.29: Diffusivity Conversion Factors}

\subsubsection{Example}\label{example-266}

The example below is based on field units, with NTPVT equal to three on
the TABDIMS keyword.

\begin{verbatim}
--
--       PVT REGION GAS COMPONENT DIFFUSION COEFFICIENTS (OPM FLOW KEYWORD)
--
DIFFCGAS
--       CO2 IN   WAT IN
--       GAS DF   GAS DF
--       -------  -------
         0.160    0.150                / PVT REGION NO. 01
         0.165    0.155                / PVT REGION NO. 02
                                       / PVT REGION NO. 03
\end{verbatim}

Here the third PVT region has no values for the two component diffusion
coefficients, and therefore the simulator will use correlations to
define the diffusivity coefficients for this PVT region.

\subsection{DIFFCOAL -- Define Coal Bed Methane Gas Diffusion
Data}\label{diffcoal-define-coal-bed-methane-gas-diffusion-data}

\begin{longtable}[]{@{}
  >{\raggedright\arraybackslash}p{(\linewidth - 14\tabcolsep) * \real{0.1250}}
  >{\raggedright\arraybackslash}p{(\linewidth - 14\tabcolsep) * \real{0.1250}}
  >{\raggedright\arraybackslash}p{(\linewidth - 14\tabcolsep) * \real{0.1250}}
  >{\raggedright\arraybackslash}p{(\linewidth - 14\tabcolsep) * \real{0.1250}}
  >{\raggedright\arraybackslash}p{(\linewidth - 14\tabcolsep) * \real{0.1250}}
  >{\raggedright\arraybackslash}p{(\linewidth - 14\tabcolsep) * \real{0.1250}}
  >{\raggedright\arraybackslash}p{(\linewidth - 14\tabcolsep) * \real{0.1250}}
  >{\raggedright\arraybackslash}p{(\linewidth - 14\tabcolsep) * \real{0.1250}}@{}}
\toprule\noalign{}
\begin{minipage}[b]{\linewidth}\raggedright
RUNSPEC
\end{minipage} & \begin{minipage}[b]{\linewidth}\raggedright
GRID
\end{minipage} & \begin{minipage}[b]{\linewidth}\raggedright
EDIT
\end{minipage} & \begin{minipage}[b]{\linewidth}\raggedright
PROPS
\end{minipage} & \begin{minipage}[b]{\linewidth}\raggedright
REGIONS
\end{minipage} & \begin{minipage}[b]{\linewidth}\raggedright
SOLUTION
\end{minipage} & \begin{minipage}[b]{\linewidth}\raggedright
SUMMARY
\end{minipage} & \begin{minipage}[b]{\linewidth}\raggedright
SCHEDULE
\end{minipage} \\
\midrule\noalign{}
\endhead
\bottomrule\noalign{}
\endlastfoot
\end{longtable}

\subsubsection{Description}\label{description-413}

The DIFF keyword defines the coal bed methane diffusion data for when
the Coal Bed Methane option has been activated via the COAL keyword in
the RUNSPEC section.

This keyword is not supported by OPM Flow but would change the results
if supported so the simulation will be stopped.

\subsection{DIFFCWAT -- Define PVT Region Water Component Diffusion
Coefficients}\label{diffcwat-define-pvt-region-water-component-diffusion-coefficients}

\begin{longtable}[]{@{}
  >{\raggedright\arraybackslash}p{(\linewidth - 14\tabcolsep) * \real{0.1250}}
  >{\raggedright\arraybackslash}p{(\linewidth - 14\tabcolsep) * \real{0.1250}}
  >{\raggedright\arraybackslash}p{(\linewidth - 14\tabcolsep) * \real{0.1250}}
  >{\raggedright\arraybackslash}p{(\linewidth - 14\tabcolsep) * \real{0.1250}}
  >{\raggedright\arraybackslash}p{(\linewidth - 14\tabcolsep) * \real{0.1250}}
  >{\raggedright\arraybackslash}p{(\linewidth - 14\tabcolsep) * \real{0.1250}}
  >{\raggedright\arraybackslash}p{(\linewidth - 14\tabcolsep) * \real{0.1250}}
  >{\raggedright\arraybackslash}p{(\linewidth - 14\tabcolsep) * \real{0.1250}}@{}}
\toprule\noalign{}
\begin{minipage}[b]{\linewidth}\raggedright
RUNSPEC
\end{minipage} & \begin{minipage}[b]{\linewidth}\raggedright
GRID
\end{minipage} & \begin{minipage}[b]{\linewidth}\raggedright
EDIT
\end{minipage} & \begin{minipage}[b]{\linewidth}\raggedright
PROPS
\end{minipage} & \begin{minipage}[b]{\linewidth}\raggedright
REGIONS
\end{minipage} & \begin{minipage}[b]{\linewidth}\raggedright
SOLUTION
\end{minipage} & \begin{minipage}[b]{\linewidth}\raggedright
SUMMARY
\end{minipage} & \begin{minipage}[b]{\linewidth}\raggedright
SCHEDULE
\end{minipage} \\
\midrule\noalign{}
\endhead
\bottomrule\noalign{}
\endlastfoot
\end{longtable}

\subsubsection{Description}\label{description-414}

The DIFFCWAT keyword defines the water diffusion coefficients assuming
the standard mole fraction formulation for each compositional component
in the model and for each PVT region, for when the Molecular Diffusion
option has been activated by the DIFFUSE keyword in the RUNSPEC section.

This keyword is optional as OPM Flow will automatically calculate the
coefficients assuming the mole fraction formulation, as described by
Sandve et al.~{[}Tor Harald Sandve, Sarah E. Gasda, Atgeirr Rasmussen,
and Alf Birger Rustad. Convective dissolution in field scale CO2 storage
simulation using the OPM Flow simulator. Submitted to TCCS 11 --
Trondheim Conference on CO2 Capture, Transport and Storage Trondheim,
Norway -- June 21-23, 2021.{]}, if the
\hyperref[REF_HEADING_KEYWORD_DIFFAWAT_8_3]{DIFFAWAT} and DIFFCWAT
keywords are absent from the input deck. The keyword thus allows one to
overwrite the automatically calculated values.

The keyword should only be used if the CO2STORE or
\hyperref[REF_HEADING_KEYWORD_H2STORE]{H2STORE} keyword and either the
GASWAT or the GAS and WATER keywords in the RUNSPEC section, have also
been activated for the gas-water two component model.

See also the \hyperref[REF_HEADING_KEYWORD_DIFFAGAS_8_3]{DIFFAGAS} and
\hyperref[REF_HEADING_KEYWORD_DIFFAWAT_8_3]{DIFFAWAT} keywords that
assume a mass fraction formulation for diffusion rather than the
standard mole fraction formulation assumed by the DIFFCGAS and DIFFCWAT
keywords. The \hyperref[REF_HEADING_KEYWORD_DIFFAGAS_8_3]{DIFFAGAS} and
\hyperref[REF_HEADING_KEYWORD_DIFFAWAT_8_3]{DIFFAWAT} keywords cannot be
used in combination with the DIFFCGAS and DIFFCWAT keywords.

\begin{tcolorbox}[enhanced jigsaw, title=\textcolor{quarto-callout-note-color}{\faInfo}\hspace{0.5em}{Note}, coltitle=black, breakable, toprule=.15mm, colbacktitle=quarto-callout-note-color!10!white, leftrule=.75mm, opacitybacktitle=0.6, colback=white, bottomtitle=1mm, toptitle=1mm, titlerule=0mm, left=2mm, colframe=quarto-callout-note-color-frame, arc=.35mm, rightrule=.15mm, bottomrule=.15mm, opacityback=0]

This is an OPM Flow keyword used with OPM Flow's CO2STORE or
\hyperref[REF_HEADING_KEYWORD_H2STORE]{H2STORE} and GASWAT keywords in
the RUNSPEC section, and should not be confused with the more general
version of the DIFFCWAT keyword used in the commercial compositional
simulator.

\end{tcolorbox}

\begin{longtable}[]{@{}
  >{\raggedright\arraybackslash}p{(\linewidth - 6\tabcolsep) * \real{0.2500}}
  >{\raggedright\arraybackslash}p{(\linewidth - 6\tabcolsep) * \real{0.2500}}
  >{\raggedright\arraybackslash}p{(\linewidth - 6\tabcolsep) * \real{0.2500}}
  >{\raggedright\arraybackslash}p{(\linewidth - 6\tabcolsep) * \real{0.2500}}@{}}
\toprule\noalign{}
\begin{minipage}[b]{\linewidth}\raggedright
No.
\end{minipage} & \begin{minipage}[b]{\linewidth}\raggedright
Name
\end{minipage} & \begin{minipage}[b]{\linewidth}\raggedright
Description
\end{minipage} & \begin{minipage}[b]{\linewidth}\raggedright
Default
\end{minipage} \\
\midrule\noalign{}
\endhead
\bottomrule\noalign{}
\endlastfoot
Field & Metric & Laboratory & \\
1 & CO2DIFF & A real positive number that declares the CO2 or H2 in
water diffusion coefficient in the given PVT region. & None \\
ft2/day & m2/day & cm2/hour & \\
2 & WATDIFF & A real positive number that specifies the water in water
diffusion coefficient in the given PVT region. & None \\
ft2/day & m2/day & cm2/hour & \\
Notes: & & & \\
\end{longtable}

\emph{Table 8.30: DIFFCWAT Keyword Description}

\begin{tcolorbox}[enhanced jigsaw, title=\textcolor{quarto-callout-note-color}{\faInfo}\hspace{0.5em}{Note}, coltitle=black, breakable, toprule=.15mm, colbacktitle=quarto-callout-note-color!10!white, leftrule=.75mm, opacitybacktitle=0.6, colback=white, bottomtitle=1mm, toptitle=1mm, titlerule=0mm, left=2mm, colframe=quarto-callout-note-color-frame, arc=.35mm, rightrule=.15mm, bottomrule=.15mm, opacityback=0]

The option has been tested in combination with the CO2STORE keyword, but
not for the general case at this point.

\end{tcolorbox}

See also the DIFFUSE keyword in the RUNSPEC section to activate the
Molecular Diffusion option and the DIFFCGAS keyword in the PROPS section
that defines the gas diffusion coefficients for each compositional
component in the model and for each PVT region. Finally, for gas-oil
systems the DIFFC keyword in the PROPS section should be used.

Normally diffusion coefficients are measured in laboratory units, that
is cm2/s, for ease of use Table 8.31 outlines the conversion factors for
converting the laboratory measured diffusion coefficients to those used
by the simulator.

\begin{longtable}[]{@{}lll@{}}
\toprule\noalign{}
Diffusivity Conversion Factors & & \\
\midrule\noalign{}
\endhead
\bottomrule\noalign{}
\endlastfoot
Laboratory Measured Units & Conversion Factor & Simulator Units \\
1 cm2/s & 92.9979 ft2/day & Field \\
8.64 m2/day & Metric & \\
3600 cm2/hour & Laboratory & \\
\end{longtable}

\emph{Table 8.31: Diffusivity Conversion Factors}

\subsubsection{Example}\label{example-267}

The example below is based on field units, with NTPVT equal to three on
the TABDIMS keyword.

\begin{verbatim}
--
--       PVT REGION WATER COMPONENT DIFFUSION COEFFICIENTS (OPM FLOW KEYWORD)
--
DIFFCWAT
--       CO2 IN   WAT IN
--       WAT DF   WAT DF
--       -------  -------
         0.160    0.150                / PVT REGION NO. 01
         0.165    0.155                / PVT REGION NO. 02
                                       / PVT REGION NO. 03
\end{verbatim}

Here the third PVT region has no values for the two component diffusion
coefficients, and therefore the simulator will use correlations to
define the diffusivity coefficients for this PVT region.

\subsection{DIFFDP -- Activate Dual Porosity Molecular Diffusion for
Matrix-Fracture Flow
Only}\label{diffdp-activate-dual-porosity-molecular-diffusion-for-matrix-fracture-flow-only}

\begin{longtable}[]{@{}
  >{\raggedright\arraybackslash}p{(\linewidth - 14\tabcolsep) * \real{0.1250}}
  >{\raggedright\arraybackslash}p{(\linewidth - 14\tabcolsep) * \real{0.1250}}
  >{\raggedright\arraybackslash}p{(\linewidth - 14\tabcolsep) * \real{0.1250}}
  >{\raggedright\arraybackslash}p{(\linewidth - 14\tabcolsep) * \real{0.1250}}
  >{\raggedright\arraybackslash}p{(\linewidth - 14\tabcolsep) * \real{0.1250}}
  >{\raggedright\arraybackslash}p{(\linewidth - 14\tabcolsep) * \real{0.1250}}
  >{\raggedright\arraybackslash}p{(\linewidth - 14\tabcolsep) * \real{0.1250}}
  >{\raggedright\arraybackslash}p{(\linewidth - 14\tabcolsep) * \real{0.1250}}@{}}
\toprule\noalign{}
\begin{minipage}[b]{\linewidth}\raggedright
RUNSPEC
\end{minipage} & \begin{minipage}[b]{\linewidth}\raggedright
GRID
\end{minipage} & \begin{minipage}[b]{\linewidth}\raggedright
EDIT
\end{minipage} & \begin{minipage}[b]{\linewidth}\raggedright
PROPS
\end{minipage} & \begin{minipage}[b]{\linewidth}\raggedright
REGIONS
\end{minipage} & \begin{minipage}[b]{\linewidth}\raggedright
SOLUTION
\end{minipage} & \begin{minipage}[b]{\linewidth}\raggedright
SUMMARY
\end{minipage} & \begin{minipage}[b]{\linewidth}\raggedright
SCHEDULE
\end{minipage} \\
\midrule\noalign{}
\endhead
\bottomrule\noalign{}
\endlastfoot
\end{longtable}

\subsubsection{Description}\label{description-415}

This keyword, DIFFDP, activates the dual porosity molecular diffusion
for matrix-fracture flow only option for when the Dual Porosity option
has be activated by either the DUALPORO or DUALPERM keywords, and the
Diffusivity option has been activated by the DIFFUSE keywords; three
keywords are in the RUNSPEC section.

This keyword is not supported by OPM Flow but would change the results
if supported so the simulation will be stopped.

\subsection{DIFFMICP -- Define PVT Region MICP Diffusion Coefficients
(Mass Concentration
Formulation)}\label{diffmicp-define-pvt-region-micp-diffusion-coefficients-mass-concentration-formulation}

\begin{longtable}[]{@{}
  >{\raggedright\arraybackslash}p{(\linewidth - 14\tabcolsep) * \real{0.1250}}
  >{\raggedright\arraybackslash}p{(\linewidth - 14\tabcolsep) * \real{0.1250}}
  >{\raggedright\arraybackslash}p{(\linewidth - 14\tabcolsep) * \real{0.1250}}
  >{\raggedright\arraybackslash}p{(\linewidth - 14\tabcolsep) * \real{0.1250}}
  >{\raggedright\arraybackslash}p{(\linewidth - 14\tabcolsep) * \real{0.1250}}
  >{\raggedright\arraybackslash}p{(\linewidth - 14\tabcolsep) * \real{0.1250}}
  >{\raggedright\arraybackslash}p{(\linewidth - 14\tabcolsep) * \real{0.1250}}
  >{\raggedright\arraybackslash}p{(\linewidth - 14\tabcolsep) * \real{0.1250}}@{}}
\toprule\noalign{}
\begin{minipage}[b]{\linewidth}\raggedright
RUNSPEC
\end{minipage} & \begin{minipage}[b]{\linewidth}\raggedright
GRID
\end{minipage} & \begin{minipage}[b]{\linewidth}\raggedright
EDIT
\end{minipage} & \begin{minipage}[b]{\linewidth}\raggedright
PROPS
\end{minipage} & \begin{minipage}[b]{\linewidth}\raggedright
REGIONS
\end{minipage} & \begin{minipage}[b]{\linewidth}\raggedright
SOLUTION
\end{minipage} & \begin{minipage}[b]{\linewidth}\raggedright
SUMMARY
\end{minipage} & \begin{minipage}[b]{\linewidth}\raggedright
SCHEDULE
\end{minipage} \\
\midrule\noalign{}
\endhead
\bottomrule\noalign{}
\endlastfoot
\end{longtable}

\subsubsection{Description}\label{description-416}

The \hyperref[REF_HEADING_KEYWORD_DIFFMICP_8_3]{DIFFMICP} keyword
defines the diffusion coefficients assuming the mass concentration
formulation for each component dissolved in water and for each PVT
region, for when the molecular diffusion option has been activated by
the DIFFUSE keyword in the RUNSPEC section. The keyword should only be
used if either the MICP or
\hyperref[REF_HEADING_KEYWORD_BIOFILM]{BIOFILM} model has been activated
in the RUNSPEC section.

\begin{tcolorbox}[enhanced jigsaw, title=\textcolor{quarto-callout-note-color}{\faInfo}\hspace{0.5em}{Note}, coltitle=black, breakable, toprule=.15mm, colbacktitle=quarto-callout-note-color!10!white, leftrule=.75mm, opacitybacktitle=0.6, colback=white, bottomtitle=1mm, toptitle=1mm, titlerule=0mm, left=2mm, colframe=quarto-callout-note-color-frame, arc=.35mm, rightrule=.15mm, bottomrule=.15mm, opacityback=0]

This is an OPM Flow specific keyword.

\end{tcolorbox}

\begin{longtable}[]{@{}
  >{\raggedright\arraybackslash}p{(\linewidth - 6\tabcolsep) * \real{0.2500}}
  >{\raggedright\arraybackslash}p{(\linewidth - 6\tabcolsep) * \real{0.2500}}
  >{\raggedright\arraybackslash}p{(\linewidth - 6\tabcolsep) * \real{0.2500}}
  >{\raggedright\arraybackslash}p{(\linewidth - 6\tabcolsep) * \real{0.2500}}@{}}
\toprule\noalign{}
\begin{minipage}[b]{\linewidth}\raggedright
No.
\end{minipage} & \begin{minipage}[b]{\linewidth}\raggedright
Name
\end{minipage} & \begin{minipage}[b]{\linewidth}\raggedright
Description
\end{minipage} & \begin{minipage}[b]{\linewidth}\raggedright
Default
\end{minipage} \\
\midrule\noalign{}
\endhead
\bottomrule\noalign{}
\endlastfoot
Field & Metric & Laboratory & \\
1 & MICRDIFF & A real positive number that declares the microbial
concentration in water diffusion coefficient in the given PVT region. &
None \\
ft2/day & m2/day & cm2/hour & \\
2 & OXYGDIFF & A real positive number that declares the oxygen
concentration in water diffusion coefficient in the given PVT region. &
None \\
ft2/day & m2/day & cm2/hour & \\
3 & UREADIFF & A real positive number that declares the urea
concentration in water diffusion coefficient in the given PVT region. &
None \\
ft2/day & m2/day & cm2/hour & \\
Notes: & & & \\
\end{longtable}

\emph{Table 8.28: DIFFMICP Keyword Description}

\subsubsection{Example}\label{example-268}

The example below is based on metric units, with NTPVT equal to one on
the TABDIMS keyword.

\begin{verbatim}
--
--       PVT REGION MICP COMPONENT DIFFUSION COEFFICIENTS (OPM FLOW KEYWORD)
--
DIFFMICP
--       MICR IN   OXYG IN   UREA IN
--       WAT DF    WAT DF    WAT DF
--       -------   -------   -------
         2E-04     2E-04     1E-04                / PVT REGION NO. 01
\end{verbatim}

\subsection{DIFFMMF -- Define Diffusivity Multipliers for
Matrix-Fractures}\label{diffmmf-define-diffusivity-multipliers-for-matrix-fractures}

\begin{longtable}[]{@{}
  >{\raggedright\arraybackslash}p{(\linewidth - 14\tabcolsep) * \real{0.1250}}
  >{\raggedright\arraybackslash}p{(\linewidth - 14\tabcolsep) * \real{0.1250}}
  >{\raggedright\arraybackslash}p{(\linewidth - 14\tabcolsep) * \real{0.1250}}
  >{\raggedright\arraybackslash}p{(\linewidth - 14\tabcolsep) * \real{0.1250}}
  >{\raggedright\arraybackslash}p{(\linewidth - 14\tabcolsep) * \real{0.1250}}
  >{\raggedright\arraybackslash}p{(\linewidth - 14\tabcolsep) * \real{0.1250}}
  >{\raggedright\arraybackslash}p{(\linewidth - 14\tabcolsep) * \real{0.1250}}
  >{\raggedright\arraybackslash}p{(\linewidth - 14\tabcolsep) * \real{0.1250}}@{}}
\toprule\noalign{}
\begin{minipage}[b]{\linewidth}\raggedright
RUNSPEC
\end{minipage} & \begin{minipage}[b]{\linewidth}\raggedright
GRID
\end{minipage} & \begin{minipage}[b]{\linewidth}\raggedright
EDIT
\end{minipage} & \begin{minipage}[b]{\linewidth}\raggedright
PROPS
\end{minipage} & \begin{minipage}[b]{\linewidth}\raggedright
REGIONS
\end{minipage} & \begin{minipage}[b]{\linewidth}\raggedright
SOLUTION
\end{minipage} & \begin{minipage}[b]{\linewidth}\raggedright
SUMMARY
\end{minipage} & \begin{minipage}[b]{\linewidth}\raggedright
SCHEDULE
\end{minipage} \\
\midrule\noalign{}
\endhead
\bottomrule\noalign{}
\endlastfoot
\end{longtable}

\subsubsection{Description}\label{description-417}

This keyword, DIFFMMF, defines the diffusivity multipliers for
matrix-fractures for when the Dual Porosity option has be activated by
either the DUALPORO or DUALPERM keywords, or the Coal Bed Methane option
is selected by the COAL keyword, and the Diffusivity option has been
activated by the DIFFUSE keywords; all four keywords are in the RUNSPEC
section.

This keyword is not supported by OPM Flow but would change the results
if supported so the simulation will be stopped.

\subsection{DISPERSE -- Define Dispersion
Tables}\label{disperse-define-dispersion-tables}

\begin{longtable}[]{@{}
  >{\raggedright\arraybackslash}p{(\linewidth - 14\tabcolsep) * \real{0.1250}}
  >{\raggedright\arraybackslash}p{(\linewidth - 14\tabcolsep) * \real{0.1250}}
  >{\raggedright\arraybackslash}p{(\linewidth - 14\tabcolsep) * \real{0.1250}}
  >{\raggedright\arraybackslash}p{(\linewidth - 14\tabcolsep) * \real{0.1250}}
  >{\raggedright\arraybackslash}p{(\linewidth - 14\tabcolsep) * \real{0.1250}}
  >{\raggedright\arraybackslash}p{(\linewidth - 14\tabcolsep) * \real{0.1250}}
  >{\raggedright\arraybackslash}p{(\linewidth - 14\tabcolsep) * \real{0.1250}}
  >{\raggedright\arraybackslash}p{(\linewidth - 14\tabcolsep) * \real{0.1250}}@{}}
\toprule\noalign{}
\begin{minipage}[b]{\linewidth}\raggedright
RUNSPEC
\end{minipage} & \begin{minipage}[b]{\linewidth}\raggedright
GRID
\end{minipage} & \begin{minipage}[b]{\linewidth}\raggedright
EDIT
\end{minipage} & \begin{minipage}[b]{\linewidth}\raggedright
PROPS
\end{minipage} & \begin{minipage}[b]{\linewidth}\raggedright
REGIONS
\end{minipage} & \begin{minipage}[b]{\linewidth}\raggedright
SOLUTION
\end{minipage} & \begin{minipage}[b]{\linewidth}\raggedright
SUMMARY
\end{minipage} & \begin{minipage}[b]{\linewidth}\raggedright
SCHEDULE
\end{minipage} \\
\midrule\noalign{}
\endhead
\bottomrule\noalign{}
\endlastfoot
\end{longtable}

\subsubsection{Description}\label{description-418}

This keyword, DISPERSE, defines the dispersion tables for when the
Dispersion option has been activated via declaring the dimensions of the
DISPERSE tables using the DISPDIMS keyword and activating the Tracer
option via the TRACERS keyword.

This keyword is not supported by OPM Flow but would change the results
if supported so the simulation will be stopped.

\subsection{DPKRMOD -- Modify Matrix Oil Relative Permeability
Data}\label{dpkrmod-modify-matrix-oil-relative-permeability-data}

\begin{longtable}[]{@{}
  >{\raggedright\arraybackslash}p{(\linewidth - 14\tabcolsep) * \real{0.1250}}
  >{\raggedright\arraybackslash}p{(\linewidth - 14\tabcolsep) * \real{0.1250}}
  >{\raggedright\arraybackslash}p{(\linewidth - 14\tabcolsep) * \real{0.1250}}
  >{\raggedright\arraybackslash}p{(\linewidth - 14\tabcolsep) * \real{0.1250}}
  >{\raggedright\arraybackslash}p{(\linewidth - 14\tabcolsep) * \real{0.1250}}
  >{\raggedright\arraybackslash}p{(\linewidth - 14\tabcolsep) * \real{0.1250}}
  >{\raggedright\arraybackslash}p{(\linewidth - 14\tabcolsep) * \real{0.1250}}
  >{\raggedright\arraybackslash}p{(\linewidth - 14\tabcolsep) * \real{0.1250}}@{}}
\toprule\noalign{}
\begin{minipage}[b]{\linewidth}\raggedright
RUNSPEC
\end{minipage} & \begin{minipage}[b]{\linewidth}\raggedright
GRID
\end{minipage} & \begin{minipage}[b]{\linewidth}\raggedright
EDIT
\end{minipage} & \begin{minipage}[b]{\linewidth}\raggedright
PROPS
\end{minipage} & \begin{minipage}[b]{\linewidth}\raggedright
REGIONS
\end{minipage} & \begin{minipage}[b]{\linewidth}\raggedright
SOLUTION
\end{minipage} & \begin{minipage}[b]{\linewidth}\raggedright
SUMMARY
\end{minipage} & \begin{minipage}[b]{\linewidth}\raggedright
SCHEDULE
\end{minipage} \\
\midrule\noalign{}
\endhead
\bottomrule\noalign{}
\endlastfoot
\end{longtable}

\subsubsection{Description}\label{description-419}

The DPKRMOD keyword can be used to modify the matrix oil relative
permeability data (oil-water, oil-gas) and the scaling of the fracture
to matrix relative permeabilities, for dual porosity runs for when a
Dual Porosity model has been activated by either the DUALPORO or
DUALPERM keywords in the RUNSPEC section.

This keyword is not supported by OPM Flow but would change the results
if supported so the simulation will be stopped.

\subsection{DSPDEINT -- Activate Brine Tracer Dispersion Interpolation
by Water
Density}\label{dspdeint-activate-brine-tracer-dispersion-interpolation-by-water-density}

\begin{longtable}[]{@{}
  >{\raggedright\arraybackslash}p{(\linewidth - 14\tabcolsep) * \real{0.1250}}
  >{\raggedright\arraybackslash}p{(\linewidth - 14\tabcolsep) * \real{0.1250}}
  >{\raggedright\arraybackslash}p{(\linewidth - 14\tabcolsep) * \real{0.1250}}
  >{\raggedright\arraybackslash}p{(\linewidth - 14\tabcolsep) * \real{0.1250}}
  >{\raggedright\arraybackslash}p{(\linewidth - 14\tabcolsep) * \real{0.1250}}
  >{\raggedright\arraybackslash}p{(\linewidth - 14\tabcolsep) * \real{0.1250}}
  >{\raggedright\arraybackslash}p{(\linewidth - 14\tabcolsep) * \real{0.1250}}
  >{\raggedright\arraybackslash}p{(\linewidth - 14\tabcolsep) * \real{0.1250}}@{}}
\toprule\noalign{}
\begin{minipage}[b]{\linewidth}\raggedright
RUNSPEC
\end{minipage} & \begin{minipage}[b]{\linewidth}\raggedright
GRID
\end{minipage} & \begin{minipage}[b]{\linewidth}\raggedright
EDIT
\end{minipage} & \begin{minipage}[b]{\linewidth}\raggedright
PROPS
\end{minipage} & \begin{minipage}[b]{\linewidth}\raggedright
REGIONS
\end{minipage} & \begin{minipage}[b]{\linewidth}\raggedright
SOLUTION
\end{minipage} & \begin{minipage}[b]{\linewidth}\raggedright
SUMMARY
\end{minipage} & \begin{minipage}[b]{\linewidth}\raggedright
SCHEDULE
\end{minipage} \\
\midrule\noalign{}
\endhead
\bottomrule\noalign{}
\endlastfoot
\end{longtable}

\subsubsection{Description}\label{description-420}

This keyword, DSPDEINT, activates the brine tracer dispersion
interpolation by water density option for when the Brine phase is
activated in the model by the BRINE keyword in the RUNSPEC section and
the DISPERSE keyword in the PROPS section is in the input file. They
keyword cause the lookup and interpolation of the DISPERSE tracer
concentration to water density, that is the tracer concentration data on
the DISPERSE keyword has been replaced by the water density data.

This keyword is not supported by OPM Flow but would change the results
if supported so the simulation will be stopped.

\subsection{ECHO -- Activate Echoing of User Input Files to the Print
File}\label{echo-activate-echoing-of-user-input-files-to-the-print-file-4}

Turns on echoing of all the input files to the print file; note that
this keyword is activated by default and can subsequently be switched
off by the NOECHO activation keyword.

See ECHO -- Activate Echoing of User Input Files to the Print File in
the GLOBAL section for a full description.

\subsection{EHYSTR -- Define Hysteresis Model and
Parameters}\label{ehystr-define-hysteresis-model-and-parameters}

\begin{longtable}[]{@{}
  >{\raggedright\arraybackslash}p{(\linewidth - 14\tabcolsep) * \real{0.1250}}
  >{\raggedright\arraybackslash}p{(\linewidth - 14\tabcolsep) * \real{0.1250}}
  >{\raggedright\arraybackslash}p{(\linewidth - 14\tabcolsep) * \real{0.1250}}
  >{\raggedright\arraybackslash}p{(\linewidth - 14\tabcolsep) * \real{0.1250}}
  >{\raggedright\arraybackslash}p{(\linewidth - 14\tabcolsep) * \real{0.1250}}
  >{\raggedright\arraybackslash}p{(\linewidth - 14\tabcolsep) * \real{0.1250}}
  >{\raggedright\arraybackslash}p{(\linewidth - 14\tabcolsep) * \real{0.1250}}
  >{\raggedright\arraybackslash}p{(\linewidth - 14\tabcolsep) * \real{0.1250}}@{}}
\toprule\noalign{}
\begin{minipage}[b]{\linewidth}\raggedright
RUNSPEC
\end{minipage} & \begin{minipage}[b]{\linewidth}\raggedright
GRID
\end{minipage} & \begin{minipage}[b]{\linewidth}\raggedright
EDIT
\end{minipage} & \begin{minipage}[b]{\linewidth}\raggedright
PROPS
\end{minipage} & \begin{minipage}[b]{\linewidth}\raggedright
REGIONS
\end{minipage} & \begin{minipage}[b]{\linewidth}\raggedright
SOLUTION
\end{minipage} & \begin{minipage}[b]{\linewidth}\raggedright
SUMMARY
\end{minipage} & \begin{minipage}[b]{\linewidth}\raggedright
SCHEDULE
\end{minipage} \\
\midrule\noalign{}
\endhead
\bottomrule\noalign{}
\endlastfoot
\end{longtable}

\subsubsection{Description}\label{description-421}

The EHYSTR keyword defines the hysteresis model and associated
parameters when the hysteresis option has been activated by the HYSTER
variable on the SATOPTS keyword in the RUNSPEC section. Both the Carlson
{[}Carlson, F. M. ``Simulation of Relative Permeability Hysteresis to
the Non-Wetting Phase,'' paper SPE 10157, presented at the SPE Annual
Technical Conference \& Exhibition, San Antonio, Texas, USA (October
5-7, 1981).{]} and Killough {[}Killough, J. E. ``Reservoir Simulation
with History-dependent Saturation Functions,'' paper SPE 5106, Society
of Petroleum Engineers Journal (1976) 16, No.~1, 37-48.{]} models are
available.

\begin{longtable}[]{@{}
  >{\raggedright\arraybackslash}p{(\linewidth - 6\tabcolsep) * \real{0.2500}}
  >{\raggedright\arraybackslash}p{(\linewidth - 6\tabcolsep) * \real{0.2500}}
  >{\raggedright\arraybackslash}p{(\linewidth - 6\tabcolsep) * \real{0.2500}}
  >{\raggedright\arraybackslash}p{(\linewidth - 6\tabcolsep) * \real{0.2500}}@{}}
\toprule\noalign{}
\begin{minipage}[b]{\linewidth}\raggedright
No.
\end{minipage} & \begin{minipage}[b]{\linewidth}\raggedright
Name
\end{minipage} & \begin{minipage}[b]{\linewidth}\raggedright
Description
\end{minipage} & \begin{minipage}[b]{\linewidth}\raggedright
Default
\end{minipage} \\
\midrule\noalign{}
\endhead
\bottomrule\noalign{}
\endlastfoot
1 & HYSTRCP & HYSTRCP is a positive real value that defines the Killough
curvature parameter for capillary pressure hysteresis model. The value
should range from 0.05 to 0.10. This item is not supported by OPM Flow
and should be defaulted (1*). & 0.1 \\
2 & HYSTMOD & An integer value that determines the relative permeability
hysteresis model to be used depending on the phase and the wettability
of the system. HYSTMOD should be set to one of the following values: &
0 \\
Water Wet Hysteresis Models & & & \\
HYSTMOD & Non-Wetting Phases & Wetting Phase & \\
-1 & Equilibration option for equilibrating the model with the SATNUM
(drainage curves) and running the model with imbibition curves (IMBNUM).
This option implies no hysteresis. & & \\
0 & Carlson Hysteresis Model & SATNUM & \\
1 & Carlson Hysteresis Model & IMBNUM & \\
2 & Killough Hysteresis Model & SATNUM & \\
3 & Killough Hysteresis Model & IMBNUM & \\
4 & Killough Hysteresis Model & Killough Hysteresis Model & \\
Oil Wet to Water Wet Water Wet & & & \\
5 & Carlson Non- Wetting Modeling for Gas and Water & SATNUM & \\
6 & Killough Non- Wetting Modeling for Gas and Water & SATNUM & \\
7 & Killough Non- Wetting Modeling for Gas and Water & Killough Non-
Wetting Modeling for the Wetting Oil Phase & \\
Only the HYSTMOD options not colored orange are currently support by OPM
Flow. & & & \\
3 & HYSTREL & HYSTREL is a positive real number that defines the
Killough's wetting phase relative permeability curvature parameter. This
parameter is only applicable if HYSMOD is set to either 4 or 7. This
item is not supported by OPM Flow and should be defaulted (1*). & 1.0 \\
4 & HYSTSGR & HYSTSGR is a positive real number that sets a scaling
parameter for the trapped non-wetting phase saturation in the Killough
model. This item is not supported by OPM Flow and should be defaulted
(1*). & 0.1 \\
5 & HYSTOPT & A defined character string that determines if the
hysteresis model should be activated for relative permeability,
capillary pressure, or both, and should be set to one of the following:
& BOTH \\
6 & HYSTSCAN & A defined character string that determines the shape of
Killough capillary pressure scanning curves when secondary reversal
occurs, that is for a drainage, imbibition, drainage cycle. Only the
RETR option is supported by OPM Flow. & RETR \\
7 & HYSTMOB & A defined character string that determines how to apply
the mobility control correction invoked by the MOBILE variable on the
EQLOPTS keyword in the RUNSPEC section. HYSTMOB should be set to one of
the following: The Mobility Control option is not supported in OPM Flow
so this parameter has no effect and will be ignored. This item should be
defaulted. & DRAIN \\
8 & HYSTWET & A defined character string that sets the wetting phase
between oil and gas in three phase systems and should be set to one of
the following: Note for all the above cases the gas relative
permeability is always modelled as a non-wetting phase. Note oil is
always the wetting phase to gas in two phase systems. This item is not
supported by OPM Flow and should be defaulted (1*). & OIL \\
9 & HYBAKOIL & Baker model one or model two relative permeability oil
phase hysteresis option used in the commercial compositional simulator.
This item is not supported by OPM Flow and should be defaulted (1*). &
NO \\
10 & HYBAKGAS & Baker model one or model two relative permeability gas
phase hysteresis option used in the commercial compositional simulator.
This item is not supported by OPM Flow and should be defaulted (1*). &
NO \\
11 & HYBAKWAT & Baker model one or model two relative permeability water
phase hysteresis option used in the commercial compositional simulator.
This item is not supported by OPM Flow and should be defaulted (1*). &
NO \\
12 & HYTHRESH & Killough's hysteresis threshold saturation value used in
the commercial compositional simulator. This item is not supported by
OPM Flow and should be defaulted (1*). & 0.0 \\
13 & HYSWETRP & Killough's hysteresis wetting phase modification used in
the commercial black-oil simulator. This item is not supported by OPM
Flow and should be defaulted (1*). & 0 \\
14 & HYSPCSCL & An integer value that determines if capillary pressure
scaling should be enabled to correct the construction of the scanning
capillary pressure curve. This option is currently disabled by default
but this is likely to change as the correction is believed to be a bug
fix. For more information see
\href{https://github.com/OPM/opm-simulators/issues/6383\#issuecomment-3329954105}{\#6383}.
This is an OPM Flow specific item that is not supported by the
commercial simulator. & 0 \\
Notes: & & & \\
\end{longtable}

\emph{Table 8.32: EHYSTR Keyword Description}

\subsubsection{Example}\label{example-269}

\begin{verbatim}
--
--       HYSTERESIS MODEL AND PARAMETERS
--
--       PC-CUR  MODEL   RELPERM TRAPPED OPTION  SHAPE    MOBILIT  WET
--       HYSTRCP HYSTMOD HYSTREL HYSTSGR HYSTOPT HYSTSCAN HYSTMOB  HYSTWET
EHYSTR
         0.1     0       0.1     1*      KR      1*       1*       1*          /
\end{verbatim}

The above example defines the hysteresis model and parameters used in
the Norne model. Here the default value is used for the Killough
curvature parameter for capillary pressure hysteresis mode, the Carlson
hysteresis model is used for the non-wetting phase and SATNUM for the
wetting phase, 0.1 is used for Killough's wetting phase relative
permeability curvature parameter (this parameter is ignored because the
Carlson model has been selected), the default values for the trapped
non-wetting phase saturation in the Killough mode (again, this parameter
is ignored because the Carlson model has been selected, and the
hysteresis modeling is only applied to relative permeability curves).

\subsection{EHYSTRR -- Define Hysteresis Model and Parameters via
SATNUM}\label{ehystrr-define-hysteresis-model-and-parameters-via-satnum}

\begin{longtable}[]{@{}
  >{\raggedright\arraybackslash}p{(\linewidth - 14\tabcolsep) * \real{0.1250}}
  >{\raggedright\arraybackslash}p{(\linewidth - 14\tabcolsep) * \real{0.1250}}
  >{\raggedright\arraybackslash}p{(\linewidth - 14\tabcolsep) * \real{0.1250}}
  >{\raggedright\arraybackslash}p{(\linewidth - 14\tabcolsep) * \real{0.1250}}
  >{\raggedright\arraybackslash}p{(\linewidth - 14\tabcolsep) * \real{0.1250}}
  >{\raggedright\arraybackslash}p{(\linewidth - 14\tabcolsep) * \real{0.1250}}
  >{\raggedright\arraybackslash}p{(\linewidth - 14\tabcolsep) * \real{0.1250}}
  >{\raggedright\arraybackslash}p{(\linewidth - 14\tabcolsep) * \real{0.1250}}@{}}
\toprule\noalign{}
\begin{minipage}[b]{\linewidth}\raggedright
RUNSPEC
\end{minipage} & \begin{minipage}[b]{\linewidth}\raggedright
GRID
\end{minipage} & \begin{minipage}[b]{\linewidth}\raggedright
EDIT
\end{minipage} & \begin{minipage}[b]{\linewidth}\raggedright
PROPS
\end{minipage} & \begin{minipage}[b]{\linewidth}\raggedright
REGIONS
\end{minipage} & \begin{minipage}[b]{\linewidth}\raggedright
SOLUTION
\end{minipage} & \begin{minipage}[b]{\linewidth}\raggedright
SUMMARY
\end{minipage} & \begin{minipage}[b]{\linewidth}\raggedright
SCHEDULE
\end{minipage} \\
\midrule\noalign{}
\endhead
\bottomrule\noalign{}
\endlastfoot
\end{longtable}

\subsubsection{Description}\label{description-422}

The EHYSTRR keyword defines the hysteresis model and associated
parameters via the drainage SATNUM allocation region array, for when the
hysteresis option has been activated by the HYSTER variable on the
SATOPTS keyword in the RUNSPEC section. Only the Killough {[}Killough,
J. E. ``Reservoir Simulation with History-dependent Saturation
Functions,'' paper SPE 5106, Society of Petroleum Engineers Journal
(1976) 16, No.~1, 37-48.{]} model is available for this keyword and the
keyword is optional.

This keyword is not supported by OPM Flow but would change the results
if supported so the simulation will be stopped. See also the EHYSTR
keyword in the RUNSPEC for an alternative keyword to enter the
hysteresis model and associated parameters that is supported by OPM Flow

\begin{longtable}[]{@{}
  >{\raggedright\arraybackslash}p{(\linewidth - 6\tabcolsep) * \real{0.2500}}
  >{\raggedright\arraybackslash}p{(\linewidth - 6\tabcolsep) * \real{0.2500}}
  >{\raggedright\arraybackslash}p{(\linewidth - 6\tabcolsep) * \real{0.2500}}
  >{\raggedright\arraybackslash}p{(\linewidth - 6\tabcolsep) * \real{0.2500}}@{}}
\toprule\noalign{}
\begin{minipage}[b]{\linewidth}\raggedright
No.
\end{minipage} & \begin{minipage}[b]{\linewidth}\raggedright
Name
\end{minipage} & \begin{minipage}[b]{\linewidth}\raggedright
Description
\end{minipage} & \begin{minipage}[b]{\linewidth}\raggedright
Default
\end{minipage} \\
\midrule\noalign{}
\endhead
\bottomrule\noalign{}
\endlastfoot
1 & HYSTRCP & HYSTRCP is a positive real value that defines the Killough
curvature parameter for capillary pressure hysteresis model. The value
should range from 0.05 to 0.10. & 0.1 \\
2 & HYSTREL & HYSTREL is a positive real number that defines the
Killough's wetting phase relative permeability curvature parameter. This
parameter is ignored if HYSMOD on the EHYSTR keyword is not set to 4. &
1.0 \\
3 & HYSTSGR & HYSTSGR is a positive real number that sets a scaling
parameter for the trapped non-wetting phase saturation in the Killough
model. & 0.1 \\
Notes: & & & \\
\end{longtable}

\emph{Table 8.33: EHYSTRR Keyword Description}

\subsubsection{Example}\label{example-270}

\begin{verbatim}
--
--       HYSTERESIS MODEL AND PARAMETERS VIA SATNUM
--
--       PC-CUR  RELPERM TRAPPED
--       HYSTRCP HYSTREL HYSTSGR
EHYSTRR
         0.04    1.0     1*                                 / SATNUM REGION 1
         0.06    1.0     1*                                 / SATNUM REGION 2
         0.08    1.0     1*                                 / SATNUM REGION 3
         0.10    1.0     1*                                 / SATNUM REGION 4
         0.10    1.0     1*                                 / SATNUM REGION 5
\end{verbatim}

The above example defines the hysteresis model and parameters for when
NTSFUN equals five on the TABDIMS keyword in the RUNSPEC section, that
is for five SATNUM regions.

\subsection{END -- Define the End of the Input
File}\label{end-define-the-end-of-the-input-file-4}

This keyword marks the end of the input file and can occur in any
section. Any keywords and data after the END keyword are ignored.

See END -- Define the End of the Input File in the GLOBAL section for a
full description.

\subsection{ENDBOX -- Define the End of the BOX Defined
Grid}\label{endbox-define-the-end-of-the-box-defined-grid-2}

This keyword marks the end of a previously defined BOX sub-grid as
defined by a previously entered BOX keyword. The keyword resets the
input grid to be the full grid as defined by the NX, NY, and NZ
variables on the DIMENS keyword in the RUNSPEC section.

See ENDBOX -- Define the End of the BOX Defined Grid in the GRID section
for a full description.

\subsection{ENDFIN -- End the Definition of a Local Grid
Refinement}\label{endfin-end-the-definition-of-a-local-grid-refinement-2}

ENDFIN defines the end of a Cartesian or radial local grid refinement
(``LGR'') definition and a LGR property definition data set.

See ENDFIN -- End the Definition of a Local Grid Refinement in the GRID
section for a full description.

\subsection{ENDINC -- Define the End of an Include
File}\label{endinc-define-the-end-of-an-include-file-4}

This keyword marks the end of an include file specified on the INCLUDE
keyword. When the ENDINC keyword is encountered in the INCLUDE file,
input data is read from the next keyword in the current file. Any
keywords and data after the ENDINC keyword in the INCLUDE file are
ignored.

See ENDINC -- Define the End of an Include File in the GLOBAL section
for a full description.

\subsection{ENDSKIP -- Deactivate Skipping of Keywords and Input
Data}\label{endskip-deactivate-skipping-of-keywords-and-input-data-4}

The ENDSKIP keyword deactivates the skipping of keywords that was
activated by the SKIP, SKIP100, or SKIP300 keywords. Each SKIP keyword
should be paired with an ENDSKIP keyword.

See ENDSKIP -- Deactivate Skipping of Keywords and Input Data in the
GLOBAL section for a full description.

\subsection{ENKRVD -- Define Relative Permeability End-Points versus
Depth
Functions}\label{enkrvd-define-relative-permeability-end-points-versus-depth-functions}

\begin{longtable}[]{@{}
  >{\raggedright\arraybackslash}p{(\linewidth - 14\tabcolsep) * \real{0.1250}}
  >{\raggedright\arraybackslash}p{(\linewidth - 14\tabcolsep) * \real{0.1250}}
  >{\raggedright\arraybackslash}p{(\linewidth - 14\tabcolsep) * \real{0.1250}}
  >{\raggedright\arraybackslash}p{(\linewidth - 14\tabcolsep) * \real{0.1250}}
  >{\raggedright\arraybackslash}p{(\linewidth - 14\tabcolsep) * \real{0.1250}}
  >{\raggedright\arraybackslash}p{(\linewidth - 14\tabcolsep) * \real{0.1250}}
  >{\raggedright\arraybackslash}p{(\linewidth - 14\tabcolsep) * \real{0.1250}}
  >{\raggedright\arraybackslash}p{(\linewidth - 14\tabcolsep) * \real{0.1250}}@{}}
\toprule\noalign{}
\begin{minipage}[b]{\linewidth}\raggedright
RUNSPEC
\end{minipage} & \begin{minipage}[b]{\linewidth}\raggedright
GRID
\end{minipage} & \begin{minipage}[b]{\linewidth}\raggedright
EDIT
\end{minipage} & \begin{minipage}[b]{\linewidth}\raggedright
PROPS
\end{minipage} & \begin{minipage}[b]{\linewidth}\raggedright
REGIONS
\end{minipage} & \begin{minipage}[b]{\linewidth}\raggedright
SOLUTION
\end{minipage} & \begin{minipage}[b]{\linewidth}\raggedright
SUMMARY
\end{minipage} & \begin{minipage}[b]{\linewidth}\raggedright
SCHEDULE
\end{minipage} \\
\midrule\noalign{}
\endhead
\bottomrule\noalign{}
\endlastfoot
\end{longtable}

\subsubsection{Description}\label{description-423}

This keyword defines the maximum oil, gas, and water relative
permeability values versus depth for the three phases and for when the
end-point scaling option has been invoked by the ENDSCALE keyword in the
RUNSPEC section. This functionality is not supported in OPM Flow.

This keyword is not supported by OPM Flow but would change the results
if supported so the simulation will be stopped.

\subsection{ENPCVD -- Define Maximum Capillary Pressure versus Depth
Functions}\label{enpcvd-define-maximum-capillary-pressure-versus-depth-functions}

\begin{longtable}[]{@{}
  >{\raggedright\arraybackslash}p{(\linewidth - 14\tabcolsep) * \real{0.1250}}
  >{\raggedright\arraybackslash}p{(\linewidth - 14\tabcolsep) * \real{0.1250}}
  >{\raggedright\arraybackslash}p{(\linewidth - 14\tabcolsep) * \real{0.1250}}
  >{\raggedright\arraybackslash}p{(\linewidth - 14\tabcolsep) * \real{0.1250}}
  >{\raggedright\arraybackslash}p{(\linewidth - 14\tabcolsep) * \real{0.1250}}
  >{\raggedright\arraybackslash}p{(\linewidth - 14\tabcolsep) * \real{0.1250}}
  >{\raggedright\arraybackslash}p{(\linewidth - 14\tabcolsep) * \real{0.1250}}
  >{\raggedright\arraybackslash}p{(\linewidth - 14\tabcolsep) * \real{0.1250}}@{}}
\toprule\noalign{}
\begin{minipage}[b]{\linewidth}\raggedright
RUNSPEC
\end{minipage} & \begin{minipage}[b]{\linewidth}\raggedright
GRID
\end{minipage} & \begin{minipage}[b]{\linewidth}\raggedright
EDIT
\end{minipage} & \begin{minipage}[b]{\linewidth}\raggedright
PROPS
\end{minipage} & \begin{minipage}[b]{\linewidth}\raggedright
REGIONS
\end{minipage} & \begin{minipage}[b]{\linewidth}\raggedright
SOLUTION
\end{minipage} & \begin{minipage}[b]{\linewidth}\raggedright
SUMMARY
\end{minipage} & \begin{minipage}[b]{\linewidth}\raggedright
SCHEDULE
\end{minipage} \\
\midrule\noalign{}
\endhead
\bottomrule\noalign{}
\endlastfoot
\end{longtable}

\subsubsection{Description}\label{description-424}

This keyword defines the maximum gas-oil and water-oil capillary
pressure values versus depth for when the end-point scaling option has
been invoked by the ENDSCALE keyword in the RUNSPEC section

This keyword is not supported by OPM Flow but would change the results
if supported so the simulation will be stopped.

\subsection{ENPTVD -- Define Relative Permeability Saturation End-Points
versus
Depth}\label{enptvd-define-relative-permeability-saturation-end-points-versus-depth}

\begin{longtable}[]{@{}
  >{\raggedright\arraybackslash}p{(\linewidth - 14\tabcolsep) * \real{0.1250}}
  >{\raggedright\arraybackslash}p{(\linewidth - 14\tabcolsep) * \real{0.1250}}
  >{\raggedright\arraybackslash}p{(\linewidth - 14\tabcolsep) * \real{0.1250}}
  >{\raggedright\arraybackslash}p{(\linewidth - 14\tabcolsep) * \real{0.1250}}
  >{\raggedright\arraybackslash}p{(\linewidth - 14\tabcolsep) * \real{0.1250}}
  >{\raggedright\arraybackslash}p{(\linewidth - 14\tabcolsep) * \real{0.1250}}
  >{\raggedright\arraybackslash}p{(\linewidth - 14\tabcolsep) * \real{0.1250}}
  >{\raggedright\arraybackslash}p{(\linewidth - 14\tabcolsep) * \real{0.1250}}@{}}
\toprule\noalign{}
\begin{minipage}[b]{\linewidth}\raggedright
RUNSPEC
\end{minipage} & \begin{minipage}[b]{\linewidth}\raggedright
GRID
\end{minipage} & \begin{minipage}[b]{\linewidth}\raggedright
EDIT
\end{minipage} & \begin{minipage}[b]{\linewidth}\raggedright
PROPS
\end{minipage} & \begin{minipage}[b]{\linewidth}\raggedright
REGIONS
\end{minipage} & \begin{minipage}[b]{\linewidth}\raggedright
SOLUTION
\end{minipage} & \begin{minipage}[b]{\linewidth}\raggedright
SUMMARY
\end{minipage} & \begin{minipage}[b]{\linewidth}\raggedright
SCHEDULE
\end{minipage} \\
\midrule\noalign{}
\endhead
\bottomrule\noalign{}
\endlastfoot
\end{longtable}

\subsubsection{Description}\label{description-425}

This keyword defines the variation of the relative permeability
saturation end-points (SWL, SWCR, etc.) for all three phases versus
depth for when the end-point scaling option has been invoked by the
ENDSCALE keyword in the RUNSPEC section. This functionality is not
supported in OPM Flow.

This keyword is not supported by OPM Flow but would change the results
if supported so the simulation will be stopped.

\subsection{ENSPCVD -- Define Capillary Pressure End-Points versus
Depth}\label{enspcvd-define-capillary-pressure-end-points-versus-depth}

\begin{longtable}[]{@{}
  >{\raggedright\arraybackslash}p{(\linewidth - 14\tabcolsep) * \real{0.1250}}
  >{\raggedright\arraybackslash}p{(\linewidth - 14\tabcolsep) * \real{0.1250}}
  >{\raggedright\arraybackslash}p{(\linewidth - 14\tabcolsep) * \real{0.1250}}
  >{\raggedright\arraybackslash}p{(\linewidth - 14\tabcolsep) * \real{0.1250}}
  >{\raggedright\arraybackslash}p{(\linewidth - 14\tabcolsep) * \real{0.1250}}
  >{\raggedright\arraybackslash}p{(\linewidth - 14\tabcolsep) * \real{0.1250}}
  >{\raggedright\arraybackslash}p{(\linewidth - 14\tabcolsep) * \real{0.1250}}
  >{\raggedright\arraybackslash}p{(\linewidth - 14\tabcolsep) * \real{0.1250}}@{}}
\toprule\noalign{}
\begin{minipage}[b]{\linewidth}\raggedright
RUNSPEC
\end{minipage} & \begin{minipage}[b]{\linewidth}\raggedright
GRID
\end{minipage} & \begin{minipage}[b]{\linewidth}\raggedright
EDIT
\end{minipage} & \begin{minipage}[b]{\linewidth}\raggedright
PROPS
\end{minipage} & \begin{minipage}[b]{\linewidth}\raggedright
REGIONS
\end{minipage} & \begin{minipage}[b]{\linewidth}\raggedright
SOLUTION
\end{minipage} & \begin{minipage}[b]{\linewidth}\raggedright
SUMMARY
\end{minipage} & \begin{minipage}[b]{\linewidth}\raggedright
SCHEDULE
\end{minipage} \\
\midrule\noalign{}
\endhead
\bottomrule\noalign{}
\endlastfoot
\end{longtable}

\subsubsection{Description}\label{description-426}

This keyword defines the variation of the capillary pressure saturation
end-points, connate gas (SGL) and connate water (SWL), versus depth for
when the end-point scaling option has been invoked by the ENDSCALE
keyword in the RUNSPEC section. This functionality is not supported in
OPM Flow.

This keyword is not supported by OPM Flow but would change the results
if supported so the simulation will be stopped.

\subsection{EOS -- Specify Equation of
State}\label{eos-specify-equation-of-state-1}

\begin{longtable}[]{@{}
  >{\raggedright\arraybackslash}p{(\linewidth - 14\tabcolsep) * \real{0.1250}}
  >{\raggedright\arraybackslash}p{(\linewidth - 14\tabcolsep) * \real{0.1250}}
  >{\raggedright\arraybackslash}p{(\linewidth - 14\tabcolsep) * \real{0.1250}}
  >{\raggedright\arraybackslash}p{(\linewidth - 14\tabcolsep) * \real{0.1250}}
  >{\raggedright\arraybackslash}p{(\linewidth - 14\tabcolsep) * \real{0.1250}}
  >{\raggedright\arraybackslash}p{(\linewidth - 14\tabcolsep) * \real{0.1250}}
  >{\raggedright\arraybackslash}p{(\linewidth - 14\tabcolsep) * \real{0.1250}}
  >{\raggedright\arraybackslash}p{(\linewidth - 14\tabcolsep) * \real{0.1250}}@{}}
\toprule\noalign{}
\begin{minipage}[b]{\linewidth}\raggedright
RUNSPEC
\end{minipage} & \begin{minipage}[b]{\linewidth}\raggedright
GRID
\end{minipage} & \begin{minipage}[b]{\linewidth}\raggedright
EDIT
\end{minipage} & \begin{minipage}[b]{\linewidth}\raggedright
PROPS
\end{minipage} & \begin{minipage}[b]{\linewidth}\raggedright
REGIONS
\end{minipage} & \begin{minipage}[b]{\linewidth}\raggedright
SOLUTION
\end{minipage} & \begin{minipage}[b]{\linewidth}\raggedright
SUMMARY
\end{minipage} & \begin{minipage}[b]{\linewidth}\raggedright
SCHEDULE
\end{minipage} \\
\midrule\noalign{}
\endhead
\bottomrule\noalign{}
\endlastfoot
\end{longtable}

\subsubsection{Description}\label{description-427}

The \hyperref[REF_HEADING_KEYWORD_EOS_5_3]{EOS} keyword specifies the
Equation Of State (\hyperref[REF_HEADING_KEYWORD_EOS_5_3]{EOS}) to be
used for each \hyperref[REF_HEADING_KEYWORD_EOS_5_3]{EOS} region. The
keyword should only be used if the compositional mode has been requested
using the COMPS keyword in the RUNSPEC section.

OPM Flow does not currently support the general compositional modeling
formulation.

This keyword is not supported by OPM Flow but it will be parsed and its
data ignored.

\begin{longtable}[]{@{}
  >{\raggedright\arraybackslash}p{(\linewidth - 6\tabcolsep) * \real{0.2500}}
  >{\raggedright\arraybackslash}p{(\linewidth - 6\tabcolsep) * \real{0.2500}}
  >{\raggedright\arraybackslash}p{(\linewidth - 6\tabcolsep) * \real{0.2500}}
  >{\raggedright\arraybackslash}p{(\linewidth - 6\tabcolsep) * \real{0.2500}}@{}}
\toprule\noalign{}
\begin{minipage}[b]{\linewidth}\raggedright
No.
\end{minipage} & \begin{minipage}[b]{\linewidth}\raggedright
Name
\end{minipage} & \begin{minipage}[b]{\linewidth}\raggedright
Description
\end{minipage} & \begin{minipage}[b]{\linewidth}\raggedright
Default
\end{minipage} \\
\midrule\noalign{}
\endhead
\bottomrule\noalign{}
\endlastfoot
1 & EQUATION & A defined character string that specifies the Equation of
State to be used, and should be set to one of the following: & PR \\
Notes: & & & \\
\end{longtable}

\emph{Table 8.3.51.1: \hyperref[REF_HEADING_KEYWORD_EOS_5_3]{EOS}
Keyword Description}

\subsubsection{Example}\label{example-271}

The following example defines Peng-Robinson as the Equation of State for
a model with a single \hyperref[REF_HEADING_KEYWORD_EOS_5_3]{EOS}
region.

\begin{verbatim}
--
--       SPECIFY EQUATION OF STATE
--
EOS
         PR                                     /
\end{verbatim}

The following example defines Peng-Robinson as the Equation of State for
\hyperref[REF_HEADING_KEYWORD_EOS_5_3]{EOS} region one, and
Soave-Redlich-Kwong as the Equation of State for
\hyperref[REF_HEADING_KEYWORD_EOS_5_3]{EOS} region two in a model with
two \hyperref[REF_HEADING_KEYWORD_EOS_5_3]{EOS} regions.

\begin{verbatim}
--
--       SPECIFY EQUATION OF STATE
--
EOS
         PR                                     /
         RK                                     /
\end{verbatim}

\subsection{EPSDBGS -- Write End-Point Debug Data to the DEBUG File
(Multiple)}\label{epsdbgs-write-end-point-debug-data-to-the-debug-file-multiple}

\begin{longtable}[]{@{}
  >{\raggedright\arraybackslash}p{(\linewidth - 14\tabcolsep) * \real{0.1250}}
  >{\raggedright\arraybackslash}p{(\linewidth - 14\tabcolsep) * \real{0.1250}}
  >{\raggedright\arraybackslash}p{(\linewidth - 14\tabcolsep) * \real{0.1250}}
  >{\raggedright\arraybackslash}p{(\linewidth - 14\tabcolsep) * \real{0.1250}}
  >{\raggedright\arraybackslash}p{(\linewidth - 14\tabcolsep) * \real{0.1250}}
  >{\raggedright\arraybackslash}p{(\linewidth - 14\tabcolsep) * \real{0.1250}}
  >{\raggedright\arraybackslash}p{(\linewidth - 14\tabcolsep) * \real{0.1250}}
  >{\raggedright\arraybackslash}p{(\linewidth - 14\tabcolsep) * \real{0.1250}}@{}}
\toprule\noalign{}
\begin{minipage}[b]{\linewidth}\raggedright
RUNSPEC
\end{minipage} & \begin{minipage}[b]{\linewidth}\raggedright
GRID
\end{minipage} & \begin{minipage}[b]{\linewidth}\raggedright
EDIT
\end{minipage} & \begin{minipage}[b]{\linewidth}\raggedright
PROPS
\end{minipage} & \begin{minipage}[b]{\linewidth}\raggedright
REGIONS
\end{minipage} & \begin{minipage}[b]{\linewidth}\raggedright
SOLUTION
\end{minipage} & \begin{minipage}[b]{\linewidth}\raggedright
SUMMARY
\end{minipage} & \begin{minipage}[b]{\linewidth}\raggedright
SCHEDULE
\end{minipage} \\
\midrule\noalign{}
\endhead
\bottomrule\noalign{}
\endlastfoot
\end{longtable}

\subsubsection{Description}\label{description-428}

This keyword, EPSDBGS, defines the end-point debug data for multiple
grid blocks that should be written to the DEBUG file (*.DBG) for when
the End-Point Scaling option has been activated by the ENDSCALE keyword
in the RONSPEC section.

This keyword is not supported by OPM Flow but has no effect on the
results so it will be ignored.

\subsection{EPSDEBUG -- Write End-Point Debug Data to the DEBUG File
(Individual)}\label{epsdebug-write-end-point-debug-data-to-the-debug-file-individual}

\begin{longtable}[]{@{}
  >{\raggedright\arraybackslash}p{(\linewidth - 14\tabcolsep) * \real{0.1250}}
  >{\raggedright\arraybackslash}p{(\linewidth - 14\tabcolsep) * \real{0.1250}}
  >{\raggedright\arraybackslash}p{(\linewidth - 14\tabcolsep) * \real{0.1250}}
  >{\raggedright\arraybackslash}p{(\linewidth - 14\tabcolsep) * \real{0.1250}}
  >{\raggedright\arraybackslash}p{(\linewidth - 14\tabcolsep) * \real{0.1250}}
  >{\raggedright\arraybackslash}p{(\linewidth - 14\tabcolsep) * \real{0.1250}}
  >{\raggedright\arraybackslash}p{(\linewidth - 14\tabcolsep) * \real{0.1250}}
  >{\raggedright\arraybackslash}p{(\linewidth - 14\tabcolsep) * \real{0.1250}}@{}}
\toprule\noalign{}
\begin{minipage}[b]{\linewidth}\raggedright
RUNSPEC
\end{minipage} & \begin{minipage}[b]{\linewidth}\raggedright
GRID
\end{minipage} & \begin{minipage}[b]{\linewidth}\raggedright
EDIT
\end{minipage} & \begin{minipage}[b]{\linewidth}\raggedright
PROPS
\end{minipage} & \begin{minipage}[b]{\linewidth}\raggedright
REGIONS
\end{minipage} & \begin{minipage}[b]{\linewidth}\raggedright
SOLUTION
\end{minipage} & \begin{minipage}[b]{\linewidth}\raggedright
SUMMARY
\end{minipage} & \begin{minipage}[b]{\linewidth}\raggedright
SCHEDULE
\end{minipage} \\
\midrule\noalign{}
\endhead
\bottomrule\noalign{}
\endlastfoot
\end{longtable}

\subsubsection{Description}\label{description-429}

This keyword, EPSDEBUG, defines the end-point debug data for individual
grid blocks that should be written to the Debug file (*.DBG) for when
the End-Point Scaling option has been activated by the ENDSCALE keyword
in the RUNSPEC section.

This keyword is not supported by OPM Flow but has no effect on the
results so it will be ignored.

\subsection{EQUALREG -- Sets an Array to a Constant by Region
Number}\label{equalreg-sets-an-array-to-a-constant-by-region-number-2}

The EQUALREG keyword sets a specified array to a constant for cells with
a specific region number. The region number array can be FLUXNUM,
MULTNUM or OPERNUM and these arrays must be defined and be available
before the EQUALREG keyword is read by the simulator. The constant can
be real or integer depending on the property array type; however, the
property arrays that can be operated on is dependent on which section
the EQUALREG keyword is being used.

See EQUALREG -- Sets an Array to a Constant by Region Number in the GRID
section for a full description.

\subsection{EQUALS -- Sets a Specified Array to a
Constant}\label{equals-sets-a-specified-array-to-a-constant-2}

The EQUALS keyword sets a specified array or part of an array to a
constant. The constant can be real or integer depending on the array
type; however, the arrays that can be operated on is dependent on which
section the EQUALS keyword is being used.

See EQUALS -- Sets a Specified Array to a Constant in the GRID section
for a full description.

\subsection{ESSNODE -- Define Salt Concentration Data for Water-Oil
Surface
Tension}\label{essnode-define-salt-concentration-data-for-water-oil-surface-tension}

\begin{longtable}[]{@{}
  >{\raggedright\arraybackslash}p{(\linewidth - 14\tabcolsep) * \real{0.1250}}
  >{\raggedright\arraybackslash}p{(\linewidth - 14\tabcolsep) * \real{0.1250}}
  >{\raggedright\arraybackslash}p{(\linewidth - 14\tabcolsep) * \real{0.1250}}
  >{\raggedright\arraybackslash}p{(\linewidth - 14\tabcolsep) * \real{0.1250}}
  >{\raggedright\arraybackslash}p{(\linewidth - 14\tabcolsep) * \real{0.1250}}
  >{\raggedright\arraybackslash}p{(\linewidth - 14\tabcolsep) * \real{0.1250}}
  >{\raggedright\arraybackslash}p{(\linewidth - 14\tabcolsep) * \real{0.1250}}
  >{\raggedright\arraybackslash}p{(\linewidth - 14\tabcolsep) * \real{0.1250}}@{}}
\toprule\noalign{}
\begin{minipage}[b]{\linewidth}\raggedright
RUNSPEC
\end{minipage} & \begin{minipage}[b]{\linewidth}\raggedright
GRID
\end{minipage} & \begin{minipage}[b]{\linewidth}\raggedright
EDIT
\end{minipage} & \begin{minipage}[b]{\linewidth}\raggedright
PROPS
\end{minipage} & \begin{minipage}[b]{\linewidth}\raggedright
REGIONS
\end{minipage} & \begin{minipage}[b]{\linewidth}\raggedright
SOLUTION
\end{minipage} & \begin{minipage}[b]{\linewidth}\raggedright
SUMMARY
\end{minipage} & \begin{minipage}[b]{\linewidth}\raggedright
SCHEDULE
\end{minipage} \\
\midrule\noalign{}
\endhead
\bottomrule\noalign{}
\endlastfoot
\end{longtable}

\subsubsection{Description}\label{description-430}

This keyword, ESSNODE, defines the salt concentration data that is used
in calculating the water-oil surface tension for when the Brine option
has been activated by the BRINE keyword in the RUNSPEC section.

This keyword is not supported by OPM Flow but would change the results
if supported so the simulation will be stopped.

\subsection{EXTRAPMS -- Activate Extrapolation Warning
Messages}\label{extrapms-activate-extrapolation-warning-messages-4}

The EXTRAPMS keyword activates extrapolation warning messages for when
OPM Flow extrapolates the PVT or VFP tables. Frequent extrapolation
warning messages should be investigated and resolved as this would
indicate possible incorrect data and may result in the simulator
extrapolating to unrealistic values.

See EXTRAPMS -- Activate Extrapolation Warning Messages in the GLOBAL
section for a full description.

\subsection{FHERCHBL -- Define Herschel-Bulkley Data versus Polymer
Concentration}\label{fherchbl-define-herschel-bulkley-data-versus-polymer-concentration}

\begin{longtable}[]{@{}
  >{\raggedright\arraybackslash}p{(\linewidth - 14\tabcolsep) * \real{0.1250}}
  >{\raggedright\arraybackslash}p{(\linewidth - 14\tabcolsep) * \real{0.1250}}
  >{\raggedright\arraybackslash}p{(\linewidth - 14\tabcolsep) * \real{0.1250}}
  >{\raggedright\arraybackslash}p{(\linewidth - 14\tabcolsep) * \real{0.1250}}
  >{\raggedright\arraybackslash}p{(\linewidth - 14\tabcolsep) * \real{0.1250}}
  >{\raggedright\arraybackslash}p{(\linewidth - 14\tabcolsep) * \real{0.1250}}
  >{\raggedright\arraybackslash}p{(\linewidth - 14\tabcolsep) * \real{0.1250}}
  >{\raggedright\arraybackslash}p{(\linewidth - 14\tabcolsep) * \real{0.1250}}@{}}
\toprule\noalign{}
\begin{minipage}[b]{\linewidth}\raggedright
RUNSPEC
\end{minipage} & \begin{minipage}[b]{\linewidth}\raggedright
GRID
\end{minipage} & \begin{minipage}[b]{\linewidth}\raggedright
EDIT
\end{minipage} & \begin{minipage}[b]{\linewidth}\raggedright
PROPS
\end{minipage} & \begin{minipage}[b]{\linewidth}\raggedright
REGIONS
\end{minipage} & \begin{minipage}[b]{\linewidth}\raggedright
SOLUTION
\end{minipage} & \begin{minipage}[b]{\linewidth}\raggedright
SUMMARY
\end{minipage} & \begin{minipage}[b]{\linewidth}\raggedright
SCHEDULE
\end{minipage} \\
\midrule\noalign{}
\endhead
\bottomrule\noalign{}
\endlastfoot
\end{longtable}

\subsubsection{Description}\label{description-431}

The FHERCHBL keyword defines Herschel-Bulkley rheological property data
for Non-Newtonian fluids versus polymer concentration, for when the
Polymer option has been invoked via the POLYMER keyword in the RUNSPEC
section and Non-Newtonian Fluid phase has been declared active by the
NNEWTF keyword, also in the RUNSPEC section.

This keyword is not supported by OPM Flow but would change the results
if supported so the simulation will be stopped.

\subsection{FILEUNIT -- Activate Unit Consistency
Verification}\label{fileunit-activate-unit-consistency-verification-1}

The FILEUNIT keyword defines the units of the data set, and is used to
verify that the units in the input deck and any associated include files
are consistent. The keyword does not provide for the conversion between
different sets of units.

See FILEUNIT -- Activate Unit Consistency Checking in the GRID section
for a full description.

\subsection{FILLEPS -- Activate Saturation End-Point Export to the INIT
File}\label{filleps-activate-saturation-end-point-export-to-the-init-file}

\begin{longtable}[]{@{}
  >{\raggedright\arraybackslash}p{(\linewidth - 14\tabcolsep) * \real{0.1250}}
  >{\raggedright\arraybackslash}p{(\linewidth - 14\tabcolsep) * \real{0.1250}}
  >{\raggedright\arraybackslash}p{(\linewidth - 14\tabcolsep) * \real{0.1250}}
  >{\raggedright\arraybackslash}p{(\linewidth - 14\tabcolsep) * \real{0.1250}}
  >{\raggedright\arraybackslash}p{(\linewidth - 14\tabcolsep) * \real{0.1250}}
  >{\raggedright\arraybackslash}p{(\linewidth - 14\tabcolsep) * \real{0.1250}}
  >{\raggedright\arraybackslash}p{(\linewidth - 14\tabcolsep) * \real{0.1250}}
  >{\raggedright\arraybackslash}p{(\linewidth - 14\tabcolsep) * \real{0.1250}}@{}}
\toprule\noalign{}
\begin{minipage}[b]{\linewidth}\raggedright
RUNSPEC
\end{minipage} & \begin{minipage}[b]{\linewidth}\raggedright
GRID
\end{minipage} & \begin{minipage}[b]{\linewidth}\raggedright
EDIT
\end{minipage} & \begin{minipage}[b]{\linewidth}\raggedright
PROPS
\end{minipage} & \begin{minipage}[b]{\linewidth}\raggedright
REGIONS
\end{minipage} & \begin{minipage}[b]{\linewidth}\raggedright
SOLUTION
\end{minipage} & \begin{minipage}[b]{\linewidth}\raggedright
SUMMARY
\end{minipage} & \begin{minipage}[b]{\linewidth}\raggedright
SCHEDULE
\end{minipage} \\
\midrule\noalign{}
\endhead
\bottomrule\noalign{}
\endlastfoot
\end{longtable}

\subsubsection{Description}\label{description-432}

This keyword switches on the export of the saturation end-point data
(SWL, SWCR, SOWCR array etc.) to the *.INIT file so that the data can be
viewed in post-processing software like OPM ResInsight.

There is no data required for this keyword.

\subsubsection{Example}\label{example-272}

\begin{verbatim}
--
--       ACTIVATE SATURATION END-POINT EXPORT TO THE INIT FILE
--
FILLEPS
\end{verbatim}

The above example switches on the export of the end-point saturation
data to the *.INIT file.

\subsection{FOAMADS -- Define Foam Rock Adsorption
Tables}\label{foamads-define-foam-rock-adsorption-tables}

\begin{longtable}[]{@{}
  >{\raggedright\arraybackslash}p{(\linewidth - 14\tabcolsep) * \real{0.1250}}
  >{\raggedright\arraybackslash}p{(\linewidth - 14\tabcolsep) * \real{0.1250}}
  >{\raggedright\arraybackslash}p{(\linewidth - 14\tabcolsep) * \real{0.1250}}
  >{\raggedright\arraybackslash}p{(\linewidth - 14\tabcolsep) * \real{0.1250}}
  >{\raggedright\arraybackslash}p{(\linewidth - 14\tabcolsep) * \real{0.1250}}
  >{\raggedright\arraybackslash}p{(\linewidth - 14\tabcolsep) * \real{0.1250}}
  >{\raggedright\arraybackslash}p{(\linewidth - 14\tabcolsep) * \real{0.1250}}
  >{\raggedright\arraybackslash}p{(\linewidth - 14\tabcolsep) * \real{0.1250}}@{}}
\toprule\noalign{}
\begin{minipage}[b]{\linewidth}\raggedright
RUNSPEC
\end{minipage} & \begin{minipage}[b]{\linewidth}\raggedright
GRID
\end{minipage} & \begin{minipage}[b]{\linewidth}\raggedright
EDIT
\end{minipage} & \begin{minipage}[b]{\linewidth}\raggedright
PROPS
\end{minipage} & \begin{minipage}[b]{\linewidth}\raggedright
REGIONS
\end{minipage} & \begin{minipage}[b]{\linewidth}\raggedright
SOLUTION
\end{minipage} & \begin{minipage}[b]{\linewidth}\raggedright
SUMMARY
\end{minipage} & \begin{minipage}[b]{\linewidth}\raggedright
SCHEDULE
\end{minipage} \\
\midrule\noalign{}
\endhead
\bottomrule\noalign{}
\endlastfoot
\end{longtable}

\subsubsection{Description}\label{description-433}

The FOAMADS keyword defines the foam rock adsorption tables for when the
Foam option has been activated by the FOAM keyword in the RUNSPEC
section.

\begin{longtable}[]{@{}
  >{\raggedright\arraybackslash}p{(\linewidth - 6\tabcolsep) * \real{0.2500}}
  >{\raggedright\arraybackslash}p{(\linewidth - 6\tabcolsep) * \real{0.2500}}
  >{\raggedright\arraybackslash}p{(\linewidth - 6\tabcolsep) * \real{0.2500}}
  >{\raggedright\arraybackslash}p{(\linewidth - 6\tabcolsep) * \real{0.2500}}@{}}
\toprule\noalign{}
\begin{minipage}[b]{\linewidth}\raggedright
No.
\end{minipage} & \begin{minipage}[b]{\linewidth}\raggedright
Name
\end{minipage} & \begin{minipage}[b]{\linewidth}\raggedright
Description
\end{minipage} & \begin{minipage}[b]{\linewidth}\raggedright
Default
\end{minipage} \\
\midrule\noalign{}
\endhead
\bottomrule\noalign{}
\endlastfoot
Field & Metric & Laboratory & \\
1 & FOAMCON & A columnar vector of real monotonically increasing down
the column values that defines the foam concentration in the solution
surrounding the rock. The first entry should be zero to define a no foam
concentration data set. Units are dependent on the transport phase
specified via the FOAMOPT1 variable on the FOAMOPTS keywod in the PROPS
section. & None \\
Gas: lb/Mscf Water: lb/stb & Gas: kg/sm3 Water: kg/sm3 & Gas: gm/scc
Water: gm/scc & \\
2 & FOAMRATI & A columnar vector of real increasing down the column
values that defines the mass of adsorbed foam per unit mass of rock for
a given FOAMCON. The first table data set entry should be zero to define
a no foam concentration data set. & None \\
lb/lb & kg/kg & gm/gm & \\
Notes: & & & \\
\end{longtable}

\emph{Table 8.3.58.1: FOAMADS Keyword Description}

\subsubsection{Example}\label{example-273}

\begin{verbatim}
--
--       FOAM ROCK ADSORPTION TABLE
--
FOAMADS
--       FOAM       FOAM
--       FOAMCON    FOAMRATI
--       -------    --------
             0.0     0.00000
             2.0     0.00003
             4.0     0.00005
             6.0     0.00007
             8.0     0.00009
            10.0     0.00011
            12.0     0.00012
            14.0     0.00015                               / TABLE NO. 01
--       FOAM       FOAM
--       FOAMCON    FOAMRATI
--       -------    --------
             0.0     0.00000
             3.0     0.00004
             5.0     0.00006
             7.0     0.00008
             8.0     0.00009
            10.0     0.00011                               / TABLE NO. 02
\end{verbatim}

The above example defines two foam rock adsorption tables assuming
NTSFUN equals two and NSSFUN is greater than or equal to eight on the
TABDIMS keyword in the RUNSPEC section.

\subsection{FOAMDCYO -- Define Foam Decay versus Oil Saturation
Tables}\label{foamdcyo-define-foam-decay-versus-oil-saturation-tables}

\begin{longtable}[]{@{}
  >{\raggedright\arraybackslash}p{(\linewidth - 14\tabcolsep) * \real{0.1250}}
  >{\raggedright\arraybackslash}p{(\linewidth - 14\tabcolsep) * \real{0.1250}}
  >{\raggedright\arraybackslash}p{(\linewidth - 14\tabcolsep) * \real{0.1250}}
  >{\raggedright\arraybackslash}p{(\linewidth - 14\tabcolsep) * \real{0.1250}}
  >{\raggedright\arraybackslash}p{(\linewidth - 14\tabcolsep) * \real{0.1250}}
  >{\raggedright\arraybackslash}p{(\linewidth - 14\tabcolsep) * \real{0.1250}}
  >{\raggedright\arraybackslash}p{(\linewidth - 14\tabcolsep) * \real{0.1250}}
  >{\raggedright\arraybackslash}p{(\linewidth - 14\tabcolsep) * \real{0.1250}}@{}}
\toprule\noalign{}
\begin{minipage}[b]{\linewidth}\raggedright
RUNSPEC
\end{minipage} & \begin{minipage}[b]{\linewidth}\raggedright
GRID
\end{minipage} & \begin{minipage}[b]{\linewidth}\raggedright
EDIT
\end{minipage} & \begin{minipage}[b]{\linewidth}\raggedright
PROPS
\end{minipage} & \begin{minipage}[b]{\linewidth}\raggedright
REGIONS
\end{minipage} & \begin{minipage}[b]{\linewidth}\raggedright
SOLUTION
\end{minipage} & \begin{minipage}[b]{\linewidth}\raggedright
SUMMARY
\end{minipage} & \begin{minipage}[b]{\linewidth}\raggedright
SCHEDULE
\end{minipage} \\
\midrule\noalign{}
\endhead
\bottomrule\noalign{}
\endlastfoot
\end{longtable}

\subsubsection{Description}\label{description-434}

The FOAMDCYO keyword defines the foam decay half-life versus oil
saturation for when the Foam option has been activated by the FOAM
keyword in the RUNSPEC section.

This keyword is not supported by OPM Flow but would change the results
if supported so the simulation will be stopped.

\subsection{FOAMDCYW -- Define Foam Decay versus Water Saturation
Tables}\label{foamdcyw-define-foam-decay-versus-water-saturation-tables}

\begin{longtable}[]{@{}
  >{\raggedright\arraybackslash}p{(\linewidth - 14\tabcolsep) * \real{0.1250}}
  >{\raggedright\arraybackslash}p{(\linewidth - 14\tabcolsep) * \real{0.1250}}
  >{\raggedright\arraybackslash}p{(\linewidth - 14\tabcolsep) * \real{0.1250}}
  >{\raggedright\arraybackslash}p{(\linewidth - 14\tabcolsep) * \real{0.1250}}
  >{\raggedright\arraybackslash}p{(\linewidth - 14\tabcolsep) * \real{0.1250}}
  >{\raggedright\arraybackslash}p{(\linewidth - 14\tabcolsep) * \real{0.1250}}
  >{\raggedright\arraybackslash}p{(\linewidth - 14\tabcolsep) * \real{0.1250}}
  >{\raggedright\arraybackslash}p{(\linewidth - 14\tabcolsep) * \real{0.1250}}@{}}
\toprule\noalign{}
\begin{minipage}[b]{\linewidth}\raggedright
RUNSPEC
\end{minipage} & \begin{minipage}[b]{\linewidth}\raggedright
GRID
\end{minipage} & \begin{minipage}[b]{\linewidth}\raggedright
EDIT
\end{minipage} & \begin{minipage}[b]{\linewidth}\raggedright
PROPS
\end{minipage} & \begin{minipage}[b]{\linewidth}\raggedright
REGIONS
\end{minipage} & \begin{minipage}[b]{\linewidth}\raggedright
SOLUTION
\end{minipage} & \begin{minipage}[b]{\linewidth}\raggedright
SUMMARY
\end{minipage} & \begin{minipage}[b]{\linewidth}\raggedright
SCHEDULE
\end{minipage} \\
\midrule\noalign{}
\endhead
\bottomrule\noalign{}
\endlastfoot
\end{longtable}

\subsubsection{Description}\label{description-435}

The FOAMDCYW keyword defines the foam decay half-life versus water
saturation for when the Foam option has been activated by the FOAM
keyword in the RUNSPEC section.

This keyword is not supported by OPM Flow but would change the results
if supported so the simulation will be stopped.

\subsection{FOAMFCN -- Define Foam Gas Mobility Reduction versus
Capillary
Number}\label{foamfcn-define-foam-gas-mobility-reduction-versus-capillary-number}

\begin{longtable}[]{@{}
  >{\raggedright\arraybackslash}p{(\linewidth - 14\tabcolsep) * \real{0.1250}}
  >{\raggedright\arraybackslash}p{(\linewidth - 14\tabcolsep) * \real{0.1250}}
  >{\raggedright\arraybackslash}p{(\linewidth - 14\tabcolsep) * \real{0.1250}}
  >{\raggedright\arraybackslash}p{(\linewidth - 14\tabcolsep) * \real{0.1250}}
  >{\raggedright\arraybackslash}p{(\linewidth - 14\tabcolsep) * \real{0.1250}}
  >{\raggedright\arraybackslash}p{(\linewidth - 14\tabcolsep) * \real{0.1250}}
  >{\raggedright\arraybackslash}p{(\linewidth - 14\tabcolsep) * \real{0.1250}}
  >{\raggedright\arraybackslash}p{(\linewidth - 14\tabcolsep) * \real{0.1250}}@{}}
\toprule\noalign{}
\begin{minipage}[b]{\linewidth}\raggedright
RUNSPEC
\end{minipage} & \begin{minipage}[b]{\linewidth}\raggedright
GRID
\end{minipage} & \begin{minipage}[b]{\linewidth}\raggedright
EDIT
\end{minipage} & \begin{minipage}[b]{\linewidth}\raggedright
PROPS
\end{minipage} & \begin{minipage}[b]{\linewidth}\raggedright
REGIONS
\end{minipage} & \begin{minipage}[b]{\linewidth}\raggedright
SOLUTION
\end{minipage} & \begin{minipage}[b]{\linewidth}\raggedright
SUMMARY
\end{minipage} & \begin{minipage}[b]{\linewidth}\raggedright
SCHEDULE
\end{minipage} \\
\midrule\noalign{}
\endhead
\bottomrule\noalign{}
\endlastfoot
\end{longtable}

\subsubsection{Description}\label{description-436}

The FOAMFCN keyword defines the reduction in gas mobility versus
capillary number, for when the Foam option has been activated by the
FOAM keyword in the RUNSPEC section, and the FOAMOPT2 parameter on the
FOAMOPTS keyword in the PROPS section has been set to the character
string FUNC.

This keyword is not supported by OPM Flow but would change the results
if supported so the simulation will be stopped.

\subsection{FOAMFRM -- Define Foam Gas Mobility Reduction versus
Reference
Mobility}\label{foamfrm-define-foam-gas-mobility-reduction-versus-reference-mobility}

\begin{longtable}[]{@{}
  >{\raggedright\arraybackslash}p{(\linewidth - 14\tabcolsep) * \real{0.1250}}
  >{\raggedright\arraybackslash}p{(\linewidth - 14\tabcolsep) * \real{0.1250}}
  >{\raggedright\arraybackslash}p{(\linewidth - 14\tabcolsep) * \real{0.1250}}
  >{\raggedright\arraybackslash}p{(\linewidth - 14\tabcolsep) * \real{0.1250}}
  >{\raggedright\arraybackslash}p{(\linewidth - 14\tabcolsep) * \real{0.1250}}
  >{\raggedright\arraybackslash}p{(\linewidth - 14\tabcolsep) * \real{0.1250}}
  >{\raggedright\arraybackslash}p{(\linewidth - 14\tabcolsep) * \real{0.1250}}
  >{\raggedright\arraybackslash}p{(\linewidth - 14\tabcolsep) * \real{0.1250}}@{}}
\toprule\noalign{}
\begin{minipage}[b]{\linewidth}\raggedright
RUNSPEC
\end{minipage} & \begin{minipage}[b]{\linewidth}\raggedright
GRID
\end{minipage} & \begin{minipage}[b]{\linewidth}\raggedright
EDIT
\end{minipage} & \begin{minipage}[b]{\linewidth}\raggedright
PROPS
\end{minipage} & \begin{minipage}[b]{\linewidth}\raggedright
REGIONS
\end{minipage} & \begin{minipage}[b]{\linewidth}\raggedright
SOLUTION
\end{minipage} & \begin{minipage}[b]{\linewidth}\raggedright
SUMMARY
\end{minipage} & \begin{minipage}[b]{\linewidth}\raggedright
SCHEDULE
\end{minipage} \\
\midrule\noalign{}
\endhead
\bottomrule\noalign{}
\endlastfoot
\end{longtable}

\subsubsection{Description}\label{description-437}

The FOAMFRM keyword defines the reduction in gas mobility versus the
reference mobility reduction factor, for when the Foam option has been
activated by the FOAM keyword in the RUNSPEC section, and the FOAMOPT2
parameter on the FOAMOPTS keyword in the PROPS section has been set to
the character string FUNC.

This keyword is not supported by OPM Flow but would change the results
if supported so the simulation will be stopped.

\subsection{FOAMFSC -- Define Foam Gas Mobility versus Surfactant
Concentration
Functions}\label{foamfsc-define-foam-gas-mobility-versus-surfactant-concentration-functions}

\begin{longtable}[]{@{}
  >{\raggedright\arraybackslash}p{(\linewidth - 14\tabcolsep) * \real{0.1250}}
  >{\raggedright\arraybackslash}p{(\linewidth - 14\tabcolsep) * \real{0.1250}}
  >{\raggedright\arraybackslash}p{(\linewidth - 14\tabcolsep) * \real{0.1250}}
  >{\raggedright\arraybackslash}p{(\linewidth - 14\tabcolsep) * \real{0.1250}}
  >{\raggedright\arraybackslash}p{(\linewidth - 14\tabcolsep) * \real{0.1250}}
  >{\raggedright\arraybackslash}p{(\linewidth - 14\tabcolsep) * \real{0.1250}}
  >{\raggedright\arraybackslash}p{(\linewidth - 14\tabcolsep) * \real{0.1250}}
  >{\raggedright\arraybackslash}p{(\linewidth - 14\tabcolsep) * \real{0.1250}}@{}}
\toprule\noalign{}
\begin{minipage}[b]{\linewidth}\raggedright
RUNSPEC
\end{minipage} & \begin{minipage}[b]{\linewidth}\raggedright
GRID
\end{minipage} & \begin{minipage}[b]{\linewidth}\raggedright
EDIT
\end{minipage} & \begin{minipage}[b]{\linewidth}\raggedright
PROPS
\end{minipage} & \begin{minipage}[b]{\linewidth}\raggedright
REGIONS
\end{minipage} & \begin{minipage}[b]{\linewidth}\raggedright
SOLUTION
\end{minipage} & \begin{minipage}[b]{\linewidth}\raggedright
SUMMARY
\end{minipage} & \begin{minipage}[b]{\linewidth}\raggedright
SCHEDULE
\end{minipage} \\
\midrule\noalign{}
\endhead
\bottomrule\noalign{}
\endlastfoot
\end{longtable}

\subsubsection{Description}\label{description-438}

The FOAMFSC keyword defines the reduction in gas mobility as a function
of the foam surfactant concentration within a grid block. The Foam
option must be activated by the FOAM keyword in the RUNSPEC section in
order to use this keyword. In addition, the FOAMOPT2 parameter on the
FOAMOPTS keyword in the PROPS section must be set to the character
string FUNC, in order to activate the functional form of the gas
mobility reduction calculations.

\begin{longtable}[]{@{}
  >{\raggedright\arraybackslash}p{(\linewidth - 6\tabcolsep) * \real{0.2500}}
  >{\raggedright\arraybackslash}p{(\linewidth - 6\tabcolsep) * \real{0.2500}}
  >{\raggedright\arraybackslash}p{(\linewidth - 6\tabcolsep) * \real{0.2500}}
  >{\raggedright\arraybackslash}p{(\linewidth - 6\tabcolsep) * \real{0.2500}}@{}}
\toprule\noalign{}
\begin{minipage}[b]{\linewidth}\raggedright
No.
\end{minipage} & \begin{minipage}[b]{\linewidth}\raggedright
Name
\end{minipage} & \begin{minipage}[b]{\linewidth}\raggedright
Description
\end{minipage} & \begin{minipage}[b]{\linewidth}\raggedright
Default
\end{minipage} \\
\midrule\noalign{}
\endhead
\bottomrule\noalign{}
\endlastfoot
Field & Metric & Laboratory & \\
1 & FOAMCON & A real positive value that defines the foam surfactant
concentration at which foam modeling becomes active in the model and a
strong foam is formed. FOAMCON cannot be defaulted and must be specified
for the first table. Subsequent tables can be defaulted and will in this
case use the previous tables' entries as the default value. & None \\
lb/stb & kg/sm3 & gm/scc & \\
2 & FOAMEXP & A real positive value that defines an exponent that
determines the gradient in the change of the reduction in gas mobility
due to foam (es in equation (8.51). Note if es is less than one then the
slope of Fs in equation (8.51) will be infinite at Cs equal to zero. In
this case, small surfactant concentrations have a significant effect on
the mobility, especially if the reference concentration Csr is also
small. If this is the case use MINSURF on this keyword to set a minimum
surfactant concentration to avoid small-scale numerical errors from
affecting the simulation. & Defined \\
dimensionless 1.0 & dimensionless 1.0 & dimensionless 1.0 & \\
3 & MINSURF & MINSURF is a real positive value that defines the minimum
surfactant concentration for which the reduction in gas mobility will be
calculated. The default value of 1 x 10-20 implies that there is no
minimum & Defined \\
lb/stb 1 x 10-20 & kg/sm3 1 x 10-20 & gm/scc 1 x 10-20 & \\
4 & MINSWAT & MINSWAT is a real positive value less than 1.0 that sets
the minimum water saturation below which foam has no effect. The default
value of 1 x 10-6 implies that there is no minimum. Note that this
parameter is only used in the commercial simulator's compositional
simulator and is therefore not used by OPM Flow or the commercial
simulators black-oil simulator. & Defined \\
dimensionless 1 x 10-6 & dimensionless 1 x 10-6 & dimensionless 1 x 10-6
& \\
Notes: & & & \\
\end{longtable}

\emph{Table 8.35: FOAMFSC Keyword Description}

The gas mobility reduction as a function of surfactant concentration is
of the form:

\begin{equation}\phantomsection\label{eq-8-51}{
{F}_{s} = {\left(\frac{{C}_{s}}{{C}_{s}^{r}}\right)}^{{e}_{s}}
}\end{equation}

Where:

Fs = the resulting gas mobility reduction factor as a function of
surfactant

concentration,

Cs = surfactant concentration,

Crs = reference surfactant concentration, that is Cs \textless{} Csr
defines a weak foam and

Cs \textgreater{} Csr defines a strong foam (FOAMCON), and

es = exponent that determines the gradient in the change of the
reduction in

gas mobility due to foam (FAOAMEXP).

The functional form of the reduction in gas mobility factor (Mrf) is:

\begin{equation}\phantomsection\label{eq-8-52}{
{M}_{\mathit{rf}} = \frac{1}{1 + \left({M}_{r} \times  {F}_{s} \times  {F}_{w} \times  {F}_{o} \times  {F}_{c}\right)}
}\end{equation}

Where:

Mr = the reference mobility reduction factor, see the FOAMFRM keyword in
the

PROPS section,

Fs = gas mobility reduction factor as a function of surfactant
concentration, see

the FOAMFSC keyword in the PROPS section,

Fw = gas mobility reduction factor as a function of water saturation,
see

the FOAMFSW keyword in the PROPS section,

Fo = gas mobility reduction factor as a function of oil saturation, see

the FOAMFSO keyword in the PROPS section, and

Fc = gas mobility reduction factor as a function of capillary number,
see

the FOAMFCN keyword in the PROPS section.

See also the FOAM keyword in the RUNSPEC section, the FOAMADS, FOAMMOB,
FOAMOPTS and FOAMROCK keywords in the PROPS section.

\subsubsection{Example}\label{example-274}

\begin{verbatim}
--
--       FOAM GAS MOBILITY VERSUS SURFACTANT CONCENTRATION FUNCTIONS
--
FOAMFSC
--       FOAM       FOAM      FOAM      FOAM
--       FOAMCON    FOAMEXP   MINSURF   MINSWAT
--       -------    -------   -------   -------
           0.001      1.010                            / TABLE NO. 01
           0.002      1.000                            / TABLE NO. 02
                                                       / TABLE NO. 03 (DEFAULTED)
           0.001      0.850  1.0E-10                   / TABLE NO. 04
           0.002      1.030                            / TABLE NO. 05
           0.002      1.000                            / TABLE NO. 06
\end{verbatim}

Here, NTSFUN equals six on the TABDIMS keyword in the RUNSPEC section
and therefore six entries are required for the FOAMFSC keyword. Table
number three is completed defaulted and will therefore use all the
properties from the previous table, that is table number two.

\subsection{FOAMFSO -- Define Foam Gas Mobility Reduction versus Oil
Saturation}\label{foamfso-define-foam-gas-mobility-reduction-versus-oil-saturation}

\begin{longtable}[]{@{}
  >{\raggedright\arraybackslash}p{(\linewidth - 14\tabcolsep) * \real{0.1250}}
  >{\raggedright\arraybackslash}p{(\linewidth - 14\tabcolsep) * \real{0.1250}}
  >{\raggedright\arraybackslash}p{(\linewidth - 14\tabcolsep) * \real{0.1250}}
  >{\raggedright\arraybackslash}p{(\linewidth - 14\tabcolsep) * \real{0.1250}}
  >{\raggedright\arraybackslash}p{(\linewidth - 14\tabcolsep) * \real{0.1250}}
  >{\raggedright\arraybackslash}p{(\linewidth - 14\tabcolsep) * \real{0.1250}}
  >{\raggedright\arraybackslash}p{(\linewidth - 14\tabcolsep) * \real{0.1250}}
  >{\raggedright\arraybackslash}p{(\linewidth - 14\tabcolsep) * \real{0.1250}}@{}}
\toprule\noalign{}
\begin{minipage}[b]{\linewidth}\raggedright
RUNSPEC
\end{minipage} & \begin{minipage}[b]{\linewidth}\raggedright
GRID
\end{minipage} & \begin{minipage}[b]{\linewidth}\raggedright
EDIT
\end{minipage} & \begin{minipage}[b]{\linewidth}\raggedright
PROPS
\end{minipage} & \begin{minipage}[b]{\linewidth}\raggedright
REGIONS
\end{minipage} & \begin{minipage}[b]{\linewidth}\raggedright
SOLUTION
\end{minipage} & \begin{minipage}[b]{\linewidth}\raggedright
SUMMARY
\end{minipage} & \begin{minipage}[b]{\linewidth}\raggedright
SCHEDULE
\end{minipage} \\
\midrule\noalign{}
\endhead
\bottomrule\noalign{}
\endlastfoot
\end{longtable}

\subsubsection{Description}\label{description-439}

The FOAMFSO keyword defines the reduction in gas mobility versus oil
saturation, for when the Foam option has been activated by the FOAM
keyword in the RUNSPEC section, and the FOAMOPT2 parameter on the
FOAMOPTS keyword in the PROPS section has been set to the character
string FUNC.

This keyword is not supported by OPM Flow but would change the results
if supported so the simulation will be stopped.

\subsection{FOAMFST -- Define Foam Gas-Water Surface Tension versus
Surfactant
Concentration}\label{foamfst-define-foam-gas-water-surface-tension-versus-surfactant-concentration}

\begin{longtable}[]{@{}
  >{\raggedright\arraybackslash}p{(\linewidth - 14\tabcolsep) * \real{0.1250}}
  >{\raggedright\arraybackslash}p{(\linewidth - 14\tabcolsep) * \real{0.1250}}
  >{\raggedright\arraybackslash}p{(\linewidth - 14\tabcolsep) * \real{0.1250}}
  >{\raggedright\arraybackslash}p{(\linewidth - 14\tabcolsep) * \real{0.1250}}
  >{\raggedright\arraybackslash}p{(\linewidth - 14\tabcolsep) * \real{0.1250}}
  >{\raggedright\arraybackslash}p{(\linewidth - 14\tabcolsep) * \real{0.1250}}
  >{\raggedright\arraybackslash}p{(\linewidth - 14\tabcolsep) * \real{0.1250}}
  >{\raggedright\arraybackslash}p{(\linewidth - 14\tabcolsep) * \real{0.1250}}@{}}
\toprule\noalign{}
\begin{minipage}[b]{\linewidth}\raggedright
RUNSPEC
\end{minipage} & \begin{minipage}[b]{\linewidth}\raggedright
GRID
\end{minipage} & \begin{minipage}[b]{\linewidth}\raggedright
EDIT
\end{minipage} & \begin{minipage}[b]{\linewidth}\raggedright
PROPS
\end{minipage} & \begin{minipage}[b]{\linewidth}\raggedright
REGIONS
\end{minipage} & \begin{minipage}[b]{\linewidth}\raggedright
SOLUTION
\end{minipage} & \begin{minipage}[b]{\linewidth}\raggedright
SUMMARY
\end{minipage} & \begin{minipage}[b]{\linewidth}\raggedright
SCHEDULE
\end{minipage} \\
\midrule\noalign{}
\endhead
\bottomrule\noalign{}
\endlastfoot
\end{longtable}

\subsubsection{Description}\label{description-440}

The FOAMFST keyword defines the gas-water surface tension versus the
foam surfactant concentration, for when the Foam option has been
activated by the FOAM keyword in the RUNSPEC section, and the FOAMOPT2
parameter on the FOAMOPTS keyword in the PROPS section has been set to
the character string FUNC.

This keyword is not supported by OPM Flow but would change the results
if supported so the simulation will be stopped.

\subsection{FOAMFSW -- Define Foam Gas Mobility Reduction versus Water
Saturation}\label{foamfsw-define-foam-gas-mobility-reduction-versus-water-saturation}

\begin{longtable}[]{@{}
  >{\raggedright\arraybackslash}p{(\linewidth - 14\tabcolsep) * \real{0.1250}}
  >{\raggedright\arraybackslash}p{(\linewidth - 14\tabcolsep) * \real{0.1250}}
  >{\raggedright\arraybackslash}p{(\linewidth - 14\tabcolsep) * \real{0.1250}}
  >{\raggedright\arraybackslash}p{(\linewidth - 14\tabcolsep) * \real{0.1250}}
  >{\raggedright\arraybackslash}p{(\linewidth - 14\tabcolsep) * \real{0.1250}}
  >{\raggedright\arraybackslash}p{(\linewidth - 14\tabcolsep) * \real{0.1250}}
  >{\raggedright\arraybackslash}p{(\linewidth - 14\tabcolsep) * \real{0.1250}}
  >{\raggedright\arraybackslash}p{(\linewidth - 14\tabcolsep) * \real{0.1250}}@{}}
\toprule\noalign{}
\begin{minipage}[b]{\linewidth}\raggedright
RUNSPEC
\end{minipage} & \begin{minipage}[b]{\linewidth}\raggedright
GRID
\end{minipage} & \begin{minipage}[b]{\linewidth}\raggedright
EDIT
\end{minipage} & \begin{minipage}[b]{\linewidth}\raggedright
PROPS
\end{minipage} & \begin{minipage}[b]{\linewidth}\raggedright
REGIONS
\end{minipage} & \begin{minipage}[b]{\linewidth}\raggedright
SOLUTION
\end{minipage} & \begin{minipage}[b]{\linewidth}\raggedright
SUMMARY
\end{minipage} & \begin{minipage}[b]{\linewidth}\raggedright
SCHEDULE
\end{minipage} \\
\midrule\noalign{}
\endhead
\bottomrule\noalign{}
\endlastfoot
\end{longtable}

\subsubsection{Description}\label{description-441}

The FOAMFRM keyword defines the reduction in gas mobility versus water
saturation, for when the Foam option has been activated by the FOAM
keyword in the RUNSPEC section, and the FOAMOPT2 parameter on the
FOAMOPTS keyword in the PROPS section has been set to the character
string FUNC.

This keyword is not supported by OPM Flow but would change the results
if supported so the simulation will be stopped.

\subsection{FOAMMOB -- Define Foam Gas Mobility versus Foam
Concentration
Tables}\label{foammob-define-foam-gas-mobility-versus-foam-concentration-tables}

\begin{longtable}[]{@{}
  >{\raggedright\arraybackslash}p{(\linewidth - 14\tabcolsep) * \real{0.1250}}
  >{\raggedright\arraybackslash}p{(\linewidth - 14\tabcolsep) * \real{0.1250}}
  >{\raggedright\arraybackslash}p{(\linewidth - 14\tabcolsep) * \real{0.1250}}
  >{\raggedright\arraybackslash}p{(\linewidth - 14\tabcolsep) * \real{0.1250}}
  >{\raggedright\arraybackslash}p{(\linewidth - 14\tabcolsep) * \real{0.1250}}
  >{\raggedright\arraybackslash}p{(\linewidth - 14\tabcolsep) * \real{0.1250}}
  >{\raggedright\arraybackslash}p{(\linewidth - 14\tabcolsep) * \real{0.1250}}
  >{\raggedright\arraybackslash}p{(\linewidth - 14\tabcolsep) * \real{0.1250}}@{}}
\toprule\noalign{}
\begin{minipage}[b]{\linewidth}\raggedright
RUNSPEC
\end{minipage} & \begin{minipage}[b]{\linewidth}\raggedright
GRID
\end{minipage} & \begin{minipage}[b]{\linewidth}\raggedright
EDIT
\end{minipage} & \begin{minipage}[b]{\linewidth}\raggedright
PROPS
\end{minipage} & \begin{minipage}[b]{\linewidth}\raggedright
REGIONS
\end{minipage} & \begin{minipage}[b]{\linewidth}\raggedright
SOLUTION
\end{minipage} & \begin{minipage}[b]{\linewidth}\raggedright
SUMMARY
\end{minipage} & \begin{minipage}[b]{\linewidth}\raggedright
SCHEDULE
\end{minipage} \\
\midrule\noalign{}
\endhead
\bottomrule\noalign{}
\endlastfoot
\end{longtable}

\subsubsection{Description}\label{description-442}

The FOAMMOB keyword defines the reduction in gas mobility as a function
of the foam concentration within a grid block. The Foam option must be
activated by the FOAM keyword in the RUNSPEC section in order to use
this keyword. In addition, this keyword must be supplied if the foam
model is activated.

\begin{longtable}[]{@{}
  >{\raggedright\arraybackslash}p{(\linewidth - 6\tabcolsep) * \real{0.2500}}
  >{\raggedright\arraybackslash}p{(\linewidth - 6\tabcolsep) * \real{0.2500}}
  >{\raggedright\arraybackslash}p{(\linewidth - 6\tabcolsep) * \real{0.2500}}
  >{\raggedright\arraybackslash}p{(\linewidth - 6\tabcolsep) * \real{0.2500}}@{}}
\toprule\noalign{}
\begin{minipage}[b]{\linewidth}\raggedright
No.
\end{minipage} & \begin{minipage}[b]{\linewidth}\raggedright
Name
\end{minipage} & \begin{minipage}[b]{\linewidth}\raggedright
Description
\end{minipage} & \begin{minipage}[b]{\linewidth}\raggedright
Default
\end{minipage} \\
\midrule\noalign{}
\endhead
\bottomrule\noalign{}
\endlastfoot
Field & Metric & Laboratory & \\
1 & FOAMCON & A columnar vector of real monotonically increasing down
the column values that defines the foam concentration for the
corresponding gas mobility reduction factor (FOAMRATI). The first entry
should be zero to define a no foam concentration data set. Units are
dependent on the transport phase specified via the FOAMOPT1 variable on
the FOAMOPTS keyword in the PROPS section. FOAMOPT1 should be set to
either GAS or WATER. & None \\
Gas: lb/Mscf Water: lb/stb & Gas: kg/sm3 Water: kg/sm3 & Gas: gm/scc
Water: gm/scc & \\
2 & FOAMRATI & A columnar vector of real decreasing down the column
values that defines the corresponding gas mobility reduction factor for
a given FOAMCON. The first table data set entry should be one to define
a no foam concentration data set. Each FOAMCON/FOAMRATI data set should
be terminated by a ``/'' & None \\
dimensionless & dimensionless & dimensionless & \\
Notes: & & & \\
\end{longtable}

\emph{Table 8.36: FOAMMOB Keyword Description}

See also the FOAM keyword in the RUNSPEC section, the FOAMADS, FOAMOPTS
and FOAMROCK keywords in the PROPS section.

\subsubsection{Example}\label{example-275}

\begin{verbatim}
--
--       FOAM GAS MOBILITY VERSUS FOAM CONCENTRATION TABLES
--
FOAMMOB
--       FOAM       FOAM
--       FOAMCON    FOAMRATI
--       -------    --------
           0.000     1.00000
           0.005     0.50000
           0.010     0.20000
           0.015     0.10000
           0.020     0.07500
           0.025     0.07000
           0.030     0.06500
           0.035     0.06500                               / TABLE NO. 01
--       FOAM       FOAM
--       FOAMCON    FOAMRATI
--       -------    --------
           0.000     1.00000
           0.010     0.50000
           0.015     0.25000
           0.020     0.07500
           0.025     0.07000
           0.030     0.07000                               / TABLE NO. 02
\end{verbatim}

Given NTPVT equals two and NPPVT is greater and or equal to eight on the
TABDIMS keyword in the RUNSPEC section, the example defines the foam gas
mobility versus foam concentration tables for two tables.

There is no terminating ``/'' for this keyword.

\subsection{FOAMMOBP -- Define Foam Mobility Reduction versus Oil
Pressure}\label{foammobp-define-foam-mobility-reduction-versus-oil-pressure}

\begin{longtable}[]{@{}
  >{\raggedright\arraybackslash}p{(\linewidth - 14\tabcolsep) * \real{0.1250}}
  >{\raggedright\arraybackslash}p{(\linewidth - 14\tabcolsep) * \real{0.1250}}
  >{\raggedright\arraybackslash}p{(\linewidth - 14\tabcolsep) * \real{0.1250}}
  >{\raggedright\arraybackslash}p{(\linewidth - 14\tabcolsep) * \real{0.1250}}
  >{\raggedright\arraybackslash}p{(\linewidth - 14\tabcolsep) * \real{0.1250}}
  >{\raggedright\arraybackslash}p{(\linewidth - 14\tabcolsep) * \real{0.1250}}
  >{\raggedright\arraybackslash}p{(\linewidth - 14\tabcolsep) * \real{0.1250}}
  >{\raggedright\arraybackslash}p{(\linewidth - 14\tabcolsep) * \real{0.1250}}@{}}
\toprule\noalign{}
\begin{minipage}[b]{\linewidth}\raggedright
RUNSPEC
\end{minipage} & \begin{minipage}[b]{\linewidth}\raggedright
GRID
\end{minipage} & \begin{minipage}[b]{\linewidth}\raggedright
EDIT
\end{minipage} & \begin{minipage}[b]{\linewidth}\raggedright
PROPS
\end{minipage} & \begin{minipage}[b]{\linewidth}\raggedright
REGIONS
\end{minipage} & \begin{minipage}[b]{\linewidth}\raggedright
SOLUTION
\end{minipage} & \begin{minipage}[b]{\linewidth}\raggedright
SUMMARY
\end{minipage} & \begin{minipage}[b]{\linewidth}\raggedright
SCHEDULE
\end{minipage} \\
\midrule\noalign{}
\endhead
\bottomrule\noalign{}
\endlastfoot
\end{longtable}

\subsubsection{Description}\label{description-443}

The FOAMMOBP keyword defines the reduction in foam mobility reduction
versus oil pressure, for when the Foam option has been activated by the
FOAM keyword in the RUNSPEC.

This keyword is not supported by OPM Flow but would change the results
if supported so the simulation will be stopped.

\subsection{FOAMMOBS -- Define Foam Mobility Reduction versus
Shear}\label{foammobs-define-foam-mobility-reduction-versus-shear}

\begin{longtable}[]{@{}
  >{\raggedright\arraybackslash}p{(\linewidth - 14\tabcolsep) * \real{0.1250}}
  >{\raggedright\arraybackslash}p{(\linewidth - 14\tabcolsep) * \real{0.1250}}
  >{\raggedright\arraybackslash}p{(\linewidth - 14\tabcolsep) * \real{0.1250}}
  >{\raggedright\arraybackslash}p{(\linewidth - 14\tabcolsep) * \real{0.1250}}
  >{\raggedright\arraybackslash}p{(\linewidth - 14\tabcolsep) * \real{0.1250}}
  >{\raggedright\arraybackslash}p{(\linewidth - 14\tabcolsep) * \real{0.1250}}
  >{\raggedright\arraybackslash}p{(\linewidth - 14\tabcolsep) * \real{0.1250}}
  >{\raggedright\arraybackslash}p{(\linewidth - 14\tabcolsep) * \real{0.1250}}@{}}
\toprule\noalign{}
\begin{minipage}[b]{\linewidth}\raggedright
RUNSPEC
\end{minipage} & \begin{minipage}[b]{\linewidth}\raggedright
GRID
\end{minipage} & \begin{minipage}[b]{\linewidth}\raggedright
EDIT
\end{minipage} & \begin{minipage}[b]{\linewidth}\raggedright
PROPS
\end{minipage} & \begin{minipage}[b]{\linewidth}\raggedright
REGIONS
\end{minipage} & \begin{minipage}[b]{\linewidth}\raggedright
SOLUTION
\end{minipage} & \begin{minipage}[b]{\linewidth}\raggedright
SUMMARY
\end{minipage} & \begin{minipage}[b]{\linewidth}\raggedright
SCHEDULE
\end{minipage} \\
\midrule\noalign{}
\endhead
\bottomrule\noalign{}
\endlastfoot
\end{longtable}

\subsubsection{Description}\label{description-444}

The FOAMMOBS keyword defines the reduction in foam mobility reduction
versus shear, for when the Foam option has been activated by the FOAM
keyword in the RUNSPEC.

This keyword is not supported by OPM Flow but would change the results
if supported so the simulation will be stopped.

\subsection{FOAMOPTS -- Define Foam Model
Options}\label{foamopts-define-foam-model-options}

\begin{longtable}[]{@{}
  >{\raggedright\arraybackslash}p{(\linewidth - 14\tabcolsep) * \real{0.1250}}
  >{\raggedright\arraybackslash}p{(\linewidth - 14\tabcolsep) * \real{0.1250}}
  >{\raggedright\arraybackslash}p{(\linewidth - 14\tabcolsep) * \real{0.1250}}
  >{\raggedright\arraybackslash}p{(\linewidth - 14\tabcolsep) * \real{0.1250}}
  >{\raggedright\arraybackslash}p{(\linewidth - 14\tabcolsep) * \real{0.1250}}
  >{\raggedright\arraybackslash}p{(\linewidth - 14\tabcolsep) * \real{0.1250}}
  >{\raggedright\arraybackslash}p{(\linewidth - 14\tabcolsep) * \real{0.1250}}
  >{\raggedright\arraybackslash}p{(\linewidth - 14\tabcolsep) * \real{0.1250}}@{}}
\toprule\noalign{}
\begin{minipage}[b]{\linewidth}\raggedright
RUNSPEC
\end{minipage} & \begin{minipage}[b]{\linewidth}\raggedright
GRID
\end{minipage} & \begin{minipage}[b]{\linewidth}\raggedright
EDIT
\end{minipage} & \begin{minipage}[b]{\linewidth}\raggedright
PROPS
\end{minipage} & \begin{minipage}[b]{\linewidth}\raggedright
REGIONS
\end{minipage} & \begin{minipage}[b]{\linewidth}\raggedright
SOLUTION
\end{minipage} & \begin{minipage}[b]{\linewidth}\raggedright
SUMMARY
\end{minipage} & \begin{minipage}[b]{\linewidth}\raggedright
SCHEDULE
\end{minipage} \\
\midrule\noalign{}
\endhead
\bottomrule\noalign{}
\endlastfoot
\end{longtable}

\subsubsection{Description}\label{description-445}

The FOAMOPTS keyword defines the transport phase for the foam (gas,
water or solvent) and how gas mobility reduction should be calculated
for when the Foam option has been activated by the FOAM keyword in the
RUNSPEC section.

\begin{longtable}[]{@{}
  >{\raggedright\arraybackslash}p{(\linewidth - 6\tabcolsep) * \real{0.2500}}
  >{\raggedright\arraybackslash}p{(\linewidth - 6\tabcolsep) * \real{0.2500}}
  >{\raggedright\arraybackslash}p{(\linewidth - 6\tabcolsep) * \real{0.2500}}
  >{\raggedright\arraybackslash}p{(\linewidth - 6\tabcolsep) * \real{0.2500}}@{}}
\toprule\noalign{}
\begin{minipage}[b]{\linewidth}\raggedright
No.
\end{minipage} & \begin{minipage}[b]{\linewidth}\raggedright
Name
\end{minipage} & \begin{minipage}[b]{\linewidth}\raggedright
Description
\end{minipage} & \begin{minipage}[b]{\linewidth}\raggedright
Default
\end{minipage} \\
\midrule\noalign{}
\endhead
\bottomrule\noalign{}
\endlastfoot
1 & FOAMOPT1 & A defined character string that specifies the transport
phase for the foam, and should be set to one of the following: & GAS \\
2 & FOAMOPT2 & A defined character string that specifies the method to
be used to calculate the reduction in gas mobility, and should be set to
one of the following: Only the default value of TAB is currently
supported by OPM Flow. & TAB \\
Notes: & & & \\
\end{longtable}

\emph{Table 8.3.70.1: FOAMOPTS Keyword Description}

\subsubsection{Example}\label{example-276}

\begin{verbatim}
--
--       FOAMOPT1  FOAMOPT2
--       PHASE     MOBILITY
--       --------  --------
FOAMOPTS
         GAS       TAB                                     / FOAM MODEL OPTIONS
\end{verbatim}

The above example defines the transport phase is to be gas and the gas
mobility reduction is to use a table as defined by the FOAMMOB keyword
as a function of foam concentration, the FOAMMOBS keyword as a function
of shear, or as a function of pressure using the FOAMMOBP keyword.

\subsection{FOAMROCK -- Define Foam Rock
Properties}\label{foamrock-define-foam-rock-properties}

\begin{longtable}[]{@{}
  >{\raggedright\arraybackslash}p{(\linewidth - 14\tabcolsep) * \real{0.1250}}
  >{\raggedright\arraybackslash}p{(\linewidth - 14\tabcolsep) * \real{0.1250}}
  >{\raggedright\arraybackslash}p{(\linewidth - 14\tabcolsep) * \real{0.1250}}
  >{\raggedright\arraybackslash}p{(\linewidth - 14\tabcolsep) * \real{0.1250}}
  >{\raggedright\arraybackslash}p{(\linewidth - 14\tabcolsep) * \real{0.1250}}
  >{\raggedright\arraybackslash}p{(\linewidth - 14\tabcolsep) * \real{0.1250}}
  >{\raggedright\arraybackslash}p{(\linewidth - 14\tabcolsep) * \real{0.1250}}
  >{\raggedright\arraybackslash}p{(\linewidth - 14\tabcolsep) * \real{0.1250}}@{}}
\toprule\noalign{}
\begin{minipage}[b]{\linewidth}\raggedright
RUNSPEC
\end{minipage} & \begin{minipage}[b]{\linewidth}\raggedright
GRID
\end{minipage} & \begin{minipage}[b]{\linewidth}\raggedright
EDIT
\end{minipage} & \begin{minipage}[b]{\linewidth}\raggedright
PROPS
\end{minipage} & \begin{minipage}[b]{\linewidth}\raggedright
REGIONS
\end{minipage} & \begin{minipage}[b]{\linewidth}\raggedright
SOLUTION
\end{minipage} & \begin{minipage}[b]{\linewidth}\raggedright
SUMMARY
\end{minipage} & \begin{minipage}[b]{\linewidth}\raggedright
SCHEDULE
\end{minipage} \\
\midrule\noalign{}
\endhead
\bottomrule\noalign{}
\endlastfoot
\end{longtable}

\subsubsection{Description}\label{description-446}

The FOAMROCK keyword defines the foam rock properties for when the Foam
option has been activated by the FOAM keyword in the RUNSPEC section.

The keyword is recognized by the input deck parser and simulator support
is available in the experimental ``ebos'' simulator.

\begin{longtable}[]{@{}
  >{\raggedright\arraybackslash}p{(\linewidth - 6\tabcolsep) * \real{0.2500}}
  >{\raggedright\arraybackslash}p{(\linewidth - 6\tabcolsep) * \real{0.2500}}
  >{\raggedright\arraybackslash}p{(\linewidth - 6\tabcolsep) * \real{0.2500}}
  >{\raggedright\arraybackslash}p{(\linewidth - 6\tabcolsep) * \real{0.2500}}@{}}
\toprule\noalign{}
\begin{minipage}[b]{\linewidth}\raggedright
No.
\end{minipage} & \begin{minipage}[b]{\linewidth}\raggedright
Name
\end{minipage} & \begin{minipage}[b]{\linewidth}\raggedright
Description
\end{minipage} & \begin{minipage}[b]{\linewidth}\raggedright
Default
\end{minipage} \\
\midrule\noalign{}
\endhead
\bottomrule\noalign{}
\endlastfoot
Field & Metric & Laboratory & \\
1 & ADINDX & A positive integer of 1 or 2 that defines foam desorption
option, as per: Only the default value of 1 is supported by OPM Flow. &
Defined \\
dimensionless 1 & dimensionless 1 & dimensionless 1 & \\
2 & DENSITY & A real value that defines the rock in situ density, that
is at reservoir conditions. & None \\
lb/rb & kg/rm3 & gm/rcc & \\
Notes: & & & \\
\end{longtable}

\emph{Table 8.38: FOAMROCK Keyword Description}

\begin{tcolorbox}[enhanced jigsaw, title=\textcolor{quarto-callout-note-color}{\faInfo}\hspace{0.5em}{Note}, coltitle=black, breakable, toprule=.15mm, colbacktitle=quarto-callout-note-color!10!white, leftrule=.75mm, opacitybacktitle=0.6, colback=white, bottomtitle=1mm, toptitle=1mm, titlerule=0mm, left=2mm, colframe=quarto-callout-note-color-frame, arc=.35mm, rightrule=.15mm, bottomrule=.15mm, opacityback=0]

In the commercial simulator if the POLYMER and SURFACT phases have been
activated in conjunction with the FOAM phase then the mass density of
rock will be set by the PLYROCK, SURFROCK, or the FOAMROCK keywords
depending on the order entered in the run deck. This is not the case for
OPM Flow. OPM Flow's FOAM phase is a standalone implementation and
cannot be used in conjunction with the either the POLYMER or SURFACT
phases.

\end{tcolorbox}

\subsubsection{Example}\label{example-277}

\begin{verbatim}
--
--       FOAM-ROCK PROPERTIES
--
FOAMROCK
--       DESORP   INSITU
--       OPTN     DENSITY
--       ------   -------
           1      1800.0                                   / TABLE NO. 01
           2      1980.0                                   / TABLE NO. 02
           1      2005.0                                   / TABLE NO. 03
\end{verbatim}

The above example defines three foam-rock tables, based on the NTSFUN
variable on the TABDIMS keyword in the RUNSPEC section being equal to
three.

There is no terminating ``/'' for this keyword.

\subsection{FORMFEED -- Defined the Print File Form-Feed
Character}\label{formfeed-defined-the-print-file-form-feed-character-4}

The FORMFEED keyword defines the form-feed character, or carriage
control character, for the output print (\emph{.PRT) run summary (}.RSM)
files. The keyword should be place at the very top of the input file.

See FORMFEED -- Defined the Print File Form-Feed Character in the GLOBAL
section for a dull description.

\subsection{GASDENT -- Define Gas Density Temperature
Coefficients}\label{gasdent-define-gas-density-temperature-coefficients}

\begin{longtable}[]{@{}
  >{\raggedright\arraybackslash}p{(\linewidth - 14\tabcolsep) * \real{0.1250}}
  >{\raggedright\arraybackslash}p{(\linewidth - 14\tabcolsep) * \real{0.1250}}
  >{\raggedright\arraybackslash}p{(\linewidth - 14\tabcolsep) * \real{0.1250}}
  >{\raggedright\arraybackslash}p{(\linewidth - 14\tabcolsep) * \real{0.1250}}
  >{\raggedright\arraybackslash}p{(\linewidth - 14\tabcolsep) * \real{0.1250}}
  >{\raggedright\arraybackslash}p{(\linewidth - 14\tabcolsep) * \real{0.1250}}
  >{\raggedright\arraybackslash}p{(\linewidth - 14\tabcolsep) * \real{0.1250}}
  >{\raggedright\arraybackslash}p{(\linewidth - 14\tabcolsep) * \real{0.1250}}@{}}
\toprule\noalign{}
\begin{minipage}[b]{\linewidth}\raggedright
RUNSPEC
\end{minipage} & \begin{minipage}[b]{\linewidth}\raggedright
GRID
\end{minipage} & \begin{minipage}[b]{\linewidth}\raggedright
EDIT
\end{minipage} & \begin{minipage}[b]{\linewidth}\raggedright
PROPS
\end{minipage} & \begin{minipage}[b]{\linewidth}\raggedright
REGIONS
\end{minipage} & \begin{minipage}[b]{\linewidth}\raggedright
SOLUTION
\end{minipage} & \begin{minipage}[b]{\linewidth}\raggedright
SUMMARY
\end{minipage} & \begin{minipage}[b]{\linewidth}\raggedright
SCHEDULE
\end{minipage} \\
\midrule\noalign{}
\endhead
\bottomrule\noalign{}
\endlastfoot
\end{longtable}

\subsubsection{Description}\label{description-447}

GASDENT defines the gas density as a function of temperature
coefficients for when OPM Flow's thermal option has been activated by
the THERMAL keyword in the RUNSPEC. Note this is an OPM Flow keyword
used with OPM Flow's black-oil thermal model that is not available in
the commercial simulator's black-oil thermal formulation.

This keyword can only be used if OPM Flow's thermal option has been
activated by the THERMAL keyword in the RUNSPEC section. Note this is
different to the commercial simulator that uses the TEMP keyword in the
RUNSPEC section to activate the black-oil thermal model.

\begin{longtable}[]{@{}
  >{\raggedright\arraybackslash}p{(\linewidth - 6\tabcolsep) * \real{0.2500}}
  >{\raggedright\arraybackslash}p{(\linewidth - 6\tabcolsep) * \real{0.2500}}
  >{\raggedright\arraybackslash}p{(\linewidth - 6\tabcolsep) * \real{0.2500}}
  >{\raggedright\arraybackslash}p{(\linewidth - 6\tabcolsep) * \real{0.2500}}@{}}
\toprule\noalign{}
\begin{minipage}[b]{\linewidth}\raggedright
No.
\end{minipage} & \begin{minipage}[b]{\linewidth}\raggedright
Name
\end{minipage} & \begin{minipage}[b]{\linewidth}\raggedright
Description
\end{minipage} & \begin{minipage}[b]{\linewidth}\raggedright
Default
\end{minipage} \\
\midrule\noalign{}
\endhead
\bottomrule\noalign{}
\endlastfoot
Field & Metric & Laboratory & \\
1 & TEMP & TEMP is a real positive value greater than zero that defines
the absolute reference temperature used with TEXP1 and TEXP2 to estimate
the change in gas density with respect to temperature. & Defined \\
oR 527.67 & oK 293.15 & oK 293.15 & \\
2 & TEXP1 & TEXP1 is a real positive value greater than zero that
defines the gas thermal expansion coefficient of the first order. &
Defined \\
1/oR 1.67 x 10-4 & 1/oK 3.0 x 10-4 & 1/oK 3.0 x 10-4 & \\
3 & TEXP2 & TEXP2 is a real positive value greater than zero that
defines the gas thermal expansion coefficient of the second order. &
Defined \\
1/oR2 9.26 x 10-7 & 1/oK2 3.0 x 10-6 & 1/oK2 3.0 x 10-6 & \\
Notes: & & & \\
\end{longtable}

\emph{Table 8.3.73.1: GASDENT Keyword Description}

The gas density at a given pressure and temperature is calculated from
its value at surface conditions and the gas expansion factor (the
reciprocal of the gas formation volume factor) as shown in the following
equation:

\begin{equation}\phantomsection\label{eq-8-3-73-1}{
{ρ}_{g}(p,T)={ρ}_{g}({p}_{s},{T}_{s}){b}_{g}(p,T)
}\end{equation}

Where the temperature dependence of the gas expansion factor relative to
its value at the reference temperature is calculated as shown in the
following equation:

\begin{equation}\phantomsection\label{eq-8-3-73-2}{
{b}_{g}(p,T)=\frac{{b}_{g}(p,{T}_{\mathit{ref}})}{1+{c}_{1}(T-{T}_{\mathit{ref}})+{c}_{2}{(T-{T}_{\mathit{ref}})}^{2}}
}\end{equation}

Where:

\({ρ}_{g}\) = gas density

\({b}_{g}\) = gas expansion factor

\(p\) = pressure

\(T\) = temperature

\({c}_{1},{c}_{2}\) = thermal expansion coefficients to first and second
order

\(s\) = subscript indicating surface conditions

\(\mathit{ref}\) = subscript indicating reference conditions

\subsubsection{Example}\label{example-278}

The following example shows the GASDENT keyword using the default
values, for when the thermal option has been activated by the THERMAL
keyword in the RUNSPEC section and for when NTPVT on the TABDIMS keyword
in the RUNSPEC section is set to two.

\begin{verbatim}
--
--       GAS DENSITY TEMPERATURE COEFFICIENTS (OPM FLOW THERMAL KEYWORD)
--
--       GAS        DENSITY   DENSITY
--       TEMP       COEFF1    COEFF2
--       --------   -------   -------
GASDENT
         1*         1*        1*                           / TABLE NO. 01
         1*         1*        1*                           / TABLE NO. 02
\end{verbatim}

There is no terminating ``/'' for this keyword.

\subsection{GASJT -- Define Gas Joule-Thomson
Coefficient}\label{gasjt-define-gas-joule-thomson-coefficient}

\begin{longtable}[]{@{}
  >{\raggedright\arraybackslash}p{(\linewidth - 14\tabcolsep) * \real{0.1250}}
  >{\raggedright\arraybackslash}p{(\linewidth - 14\tabcolsep) * \real{0.1250}}
  >{\raggedright\arraybackslash}p{(\linewidth - 14\tabcolsep) * \real{0.1250}}
  >{\raggedright\arraybackslash}p{(\linewidth - 14\tabcolsep) * \real{0.1250}}
  >{\raggedright\arraybackslash}p{(\linewidth - 14\tabcolsep) * \real{0.1250}}
  >{\raggedright\arraybackslash}p{(\linewidth - 14\tabcolsep) * \real{0.1250}}
  >{\raggedright\arraybackslash}p{(\linewidth - 14\tabcolsep) * \real{0.1250}}
  >{\raggedright\arraybackslash}p{(\linewidth - 14\tabcolsep) * \real{0.1250}}@{}}
\toprule\noalign{}
\begin{minipage}[b]{\linewidth}\raggedright
RUNSPEC
\end{minipage} & \begin{minipage}[b]{\linewidth}\raggedright
GRID
\end{minipage} & \begin{minipage}[b]{\linewidth}\raggedright
EDIT
\end{minipage} & \begin{minipage}[b]{\linewidth}\raggedright
PROPS
\end{minipage} & \begin{minipage}[b]{\linewidth}\raggedright
REGIONS
\end{minipage} & \begin{minipage}[b]{\linewidth}\raggedright
SOLUTION
\end{minipage} & \begin{minipage}[b]{\linewidth}\raggedright
SUMMARY
\end{minipage} & \begin{minipage}[b]{\linewidth}\raggedright
SCHEDULE
\end{minipage} \\
\midrule\noalign{}
\endhead
\bottomrule\noalign{}
\endlastfoot
\end{longtable}

\subsubsection{Description}\label{description-448}

GASJT activates the gas Joule-Thomson {[}Natural Gas Engineering
(McGraw-Hill chemical engineering series), Donald L. Katz, Robert l Lee,
McGraw-Hill Education, 1990 (ISBN 0071007776, 9780071007771).{]} effect
in temperature calculations, and defines the gas Joule-Thomson
Coefficient (``JTC'') at a given reference pressure, for when OPM Flow's
thermal option has been activated by the THERMAL keyword in the RUNSPEC.

\begin{tcolorbox}[enhanced jigsaw, title=\textcolor{quarto-callout-note-color}{\faInfo}\hspace{0.5em}{Note}, coltitle=black, breakable, toprule=.15mm, colbacktitle=quarto-callout-note-color!10!white, leftrule=.75mm, opacitybacktitle=0.6, colback=white, bottomtitle=1mm, toptitle=1mm, titlerule=0mm, left=2mm, colframe=quarto-callout-note-color-frame, arc=.35mm, rightrule=.15mm, bottomrule=.15mm, opacityback=0]

This is an OPM Flow keyword used with OPM Flow's black-oil thermal
model, that is not available in the commercial simulator's black-oil
thermal formulation.

\end{tcolorbox}

This keyword can only be used if OPM Flow's thermal option has been
activated by the THERMAL keyword in the RUNSPEC section. Note this is
different to the commercial simulator that uses the TEMP keyword in the
RUNSPEC section to activate the black-oil thermal model, and does not
include the Joule-Thomson effect in temperature calculations.

\begin{longtable}[]{@{}
  >{\raggedright\arraybackslash}p{(\linewidth - 6\tabcolsep) * \real{0.2500}}
  >{\raggedright\arraybackslash}p{(\linewidth - 6\tabcolsep) * \real{0.2500}}
  >{\raggedright\arraybackslash}p{(\linewidth - 6\tabcolsep) * \real{0.2500}}
  >{\raggedright\arraybackslash}p{(\linewidth - 6\tabcolsep) * \real{0.2500}}@{}}
\toprule\noalign{}
\begin{minipage}[b]{\linewidth}\raggedright
No.
\end{minipage} & \begin{minipage}[b]{\linewidth}\raggedright
Name
\end{minipage} & \begin{minipage}[b]{\linewidth}\raggedright
Description
\end{minipage} & \begin{minipage}[b]{\linewidth}\raggedright
Default
\end{minipage} \\
\midrule\noalign{}
\endhead
\bottomrule\noalign{}
\endlastfoot
Field & Metric & Laboratory & \\
1 & PRESS & A real positive value that defines the reference pressure
for the corresponding Joule-Thomson Coefficient, GASJTC. & None \\
psia & barsa & atma & \\
2 & GASJTC & GASJTC is a real positive or negative value that defines
the gas phase Joule-Thomson Coefficient. If the value is defaulted (1*)
or set to 0, then GASJTC is internally calculated using the thermal gas
density data on the GASDENT keyword in the PROPS section. If a non-zero
value is specified, then the GASJTC is assumed to be constant and equal
to that value. & 0 \\
oF/psia & oC/barsa & oC/atma & \\
Notes: & & & \\
\end{longtable}

\emph{Table 8.40: GASJT Keyword Description}

The Joule--Thomson effect is when a real gas, as oppose to an ideal gas,
expands, resulting in the temperature of the gas dropping. During
passage of a gas through a choke, the internal energy is transferred to
kinetic energy with a corresponding reduction in temperature as the
velocity increases. The effect for natural gas is approximately 7 oF for
every 100 psi pressure reduction, or 0.5 oC per barsa
{[}https://petrowiki.spe.org/Glossary:Joule-thompson\_effect.{]}, is
valid for ``normal'' pressures and temperatures at the surface.

Thermodynamically, the Joule--Thomson coefficient is defined as the
isenthalpic {[}An isenthalpic process or isoenthalpic process, is a
process that proceeds without any change in enthalpy, H; or specific
enthalpy, h.{]} change in temperature in a fluid caused by a unitary
pressure drop, as shown in the following equation:

\begin{equation}\phantomsection\label{eq-8-53}{
\mathrm{η} = \left(\frac{\partial T}{\partial P}\right)
}\end{equation}

One can also express enthalpy changes in terms of pressure, temperature
and volume changes:

\begin{equation}\phantomsection\label{eq-8-54}{
\mathrm{η} = \frac{{\mathit{RT}}^{2}}{{\mathit{Pc}}_{p}}{\left(\frac{\partial Z}{\partial T}\right)}_{P}
}\end{equation}

Where:

\(η\) = Joule--Thomson coefficient,

\({C}_{g}\) = specific heat at constant pressure,

\(P\) = pressure,

\(R\) = gas constant,

\(T\) = temperature, and

\(Z\) = gas compressibility factor.

\subsubsection{Example}\label{example-279}

The following example shows the GASJT keyword for when the thermal
option has been activated by the THERMAL keyword in the RUNSPEC section,
and for when NTPVT on the TABDIMS keyword in the RUNSPEC section is set
equal to two.

\begin{verbatim}
--
--       GAS JOULE-THOMSON COEFFICIENT (OPM FLOW EXTENSION KEYWORD)
--
--       REF        GAS
--       PRESS      JTC
--       --------   -------
GASJT
         20.0       1*                  / TABLE NO. 01
         20.0       0.50                / TABLE NO. 02
\end{verbatim}

Here the first entry is defaulted, and the simulator will therefore
calculate the gas JTC internally using the data on the GASDENT keyword
in the PROPS section.

There is no terminating ``/'' for this keyword.

\subsection{GASVISCT -- Define Gas Viscosity versus Temperature
Functions}\label{gasvisct-define-gas-viscosity-versus-temperature-functions}

\begin{longtable}[]{@{}
  >{\raggedright\arraybackslash}p{(\linewidth - 14\tabcolsep) * \real{0.1250}}
  >{\raggedright\arraybackslash}p{(\linewidth - 14\tabcolsep) * \real{0.1250}}
  >{\raggedright\arraybackslash}p{(\linewidth - 14\tabcolsep) * \real{0.1250}}
  >{\raggedright\arraybackslash}p{(\linewidth - 14\tabcolsep) * \real{0.1250}}
  >{\raggedright\arraybackslash}p{(\linewidth - 14\tabcolsep) * \real{0.1250}}
  >{\raggedright\arraybackslash}p{(\linewidth - 14\tabcolsep) * \real{0.1250}}
  >{\raggedright\arraybackslash}p{(\linewidth - 14\tabcolsep) * \real{0.1250}}
  >{\raggedright\arraybackslash}p{(\linewidth - 14\tabcolsep) * \real{0.1250}}@{}}
\toprule\noalign{}
\begin{minipage}[b]{\linewidth}\raggedright
RUNSPEC
\end{minipage} & \begin{minipage}[b]{\linewidth}\raggedright
GRID
\end{minipage} & \begin{minipage}[b]{\linewidth}\raggedright
EDIT
\end{minipage} & \begin{minipage}[b]{\linewidth}\raggedright
PROPS
\end{minipage} & \begin{minipage}[b]{\linewidth}\raggedright
REGIONS
\end{minipage} & \begin{minipage}[b]{\linewidth}\raggedright
SOLUTION
\end{minipage} & \begin{minipage}[b]{\linewidth}\raggedright
SUMMARY
\end{minipage} & \begin{minipage}[b]{\linewidth}\raggedright
SCHEDULE
\end{minipage} \\
\midrule\noalign{}
\endhead
\bottomrule\noalign{}
\endlastfoot
\end{longtable}

\subsubsection{Description}\label{description-449}

GASVISCT defines the gas viscosity as a function of temperature for when
OPM Flow's thermal option has been activated by the THERMAL keyword in
the RUNSPEC. The reference pressure for this table is given by the
VISCREF keyword in the PROPS section. Note this is an OPM Flow keyword
used with OPM Flow's black-oil thermal model that is not available in
the commercial simulator's black-oil thermal formulation. However, the
keyword and similar functionality is available in the commercial
compositional simulator.

This keyword can only be used if OPM Flow's thermal option has been
activated by the THERMAL keyword in the RUNSPEC section. Note this is
different to the commercial simulator that uses the TEMP keyword in the
RUNSPEC section to activate the black-oil thermal model.

\begin{longtable}[]{@{}
  >{\raggedright\arraybackslash}p{(\linewidth - 6\tabcolsep) * \real{0.2500}}
  >{\raggedright\arraybackslash}p{(\linewidth - 6\tabcolsep) * \real{0.2500}}
  >{\raggedright\arraybackslash}p{(\linewidth - 6\tabcolsep) * \real{0.2500}}
  >{\raggedright\arraybackslash}p{(\linewidth - 6\tabcolsep) * \real{0.2500}}@{}}
\toprule\noalign{}
\begin{minipage}[b]{\linewidth}\raggedright
No.
\end{minipage} & \begin{minipage}[b]{\linewidth}\raggedright
Name
\end{minipage} & \begin{minipage}[b]{\linewidth}\raggedright
Description
\end{minipage} & \begin{minipage}[b]{\linewidth}\raggedright
Default
\end{minipage} \\
\midrule\noalign{}
\endhead
\bottomrule\noalign{}
\endlastfoot
Field & Metric & Laboratory & \\
1 & TEMP & A columnar vector of real monotonically increasing down the
column values that defines the temperature values. & None \\
oF & oC & oC & \\
2 & VIS & A columnar vector of real increasing down the column values
that defines the gas viscosity for the corresponding temperature values
(TEMP). VIS should be given at the reference pressure defined by the
PRESS variable on the VISCREF keyword. & None \\
cP & cP & cP & \\
Notes: & & & \\
\end{longtable}

\emph{Table 8.41: GASVISCT Keyword Description}

\subsubsection{Example}\label{example-280}

The following example shows the GASVISCT keyword for when the thermal
option has been activated by the THERMAL keyword in the RUNSPEC section
and for when NTPVT on the TABDIMS keyword in the RUNSPEC section is set
equal to one.

\begin{verbatim}
--
--       GAS VISCOSITY VERSUS TEMPERATURE TABLES (OPM FLOW EXTENSION KEYWORD)
--
--       GAS        GAS
--       TEMP       VISC
--       --------   -------
GASVISCT
           100.0    0.0500
           110.0    0.0550
           120.0    0.0580
           150.0    0.0620
           165.0    0.0625                                 / TABLE NO. 01
\end{verbatim}

There is no terminating ``/'' for this keyword.

\subsection{GETDATA -- Load and Assign Data Array from INIT or RESTART
File}\label{getdata-load-and-assign-data-array-from-init-or-restart-file-1}

The GETDATA keyword loads a data array from a previously generated INIT
or RESTART file and assigns the loaded array to either same array in the
run or another array name.

See GETDATA -- Load and Assign Data Array from INIT or RESTART Files in
the GRID section for a full description.

\subsection{GIALL -- Define Gi Values and PVT Properties versus
Pressure}\label{giall-define-gi-values-and-pvt-properties-versus-pressure}

\begin{longtable}[]{@{}
  >{\raggedright\arraybackslash}p{(\linewidth - 14\tabcolsep) * \real{0.1250}}
  >{\raggedright\arraybackslash}p{(\linewidth - 14\tabcolsep) * \real{0.1250}}
  >{\raggedright\arraybackslash}p{(\linewidth - 14\tabcolsep) * \real{0.1250}}
  >{\raggedright\arraybackslash}p{(\linewidth - 14\tabcolsep) * \real{0.1250}}
  >{\raggedright\arraybackslash}p{(\linewidth - 14\tabcolsep) * \real{0.1250}}
  >{\raggedright\arraybackslash}p{(\linewidth - 14\tabcolsep) * \real{0.1250}}
  >{\raggedright\arraybackslash}p{(\linewidth - 14\tabcolsep) * \real{0.1250}}
  >{\raggedright\arraybackslash}p{(\linewidth - 14\tabcolsep) * \real{0.1250}}@{}}
\toprule\noalign{}
\begin{minipage}[b]{\linewidth}\raggedright
RUNSPEC
\end{minipage} & \begin{minipage}[b]{\linewidth}\raggedright
GRID
\end{minipage} & \begin{minipage}[b]{\linewidth}\raggedright
EDIT
\end{minipage} & \begin{minipage}[b]{\linewidth}\raggedright
PROPS
\end{minipage} & \begin{minipage}[b]{\linewidth}\raggedright
REGIONS
\end{minipage} & \begin{minipage}[b]{\linewidth}\raggedright
SOLUTION
\end{minipage} & \begin{minipage}[b]{\linewidth}\raggedright
SUMMARY
\end{minipage} & \begin{minipage}[b]{\linewidth}\raggedright
SCHEDULE
\end{minipage} \\
\midrule\noalign{}
\endhead
\bottomrule\noalign{}
\endlastfoot
\end{longtable}

\subsubsection{Description}\label{description-450}

The GIALL keyword defines the GI values and the associated RVGI, RSGI,
BGGI and BOGI values as a function of pressure, for when the GI Pseudo
Compositional option has been activated in the model via the GIMODEL
keyword in the RUNSPEC section.

This keyword is not supported by OPM Flow but would change the results
if supported so the simulation will be stopped.

\subsection{GINODE -- Define Gi Node
Values}\label{ginode-define-gi-node-values}

\begin{longtable}[]{@{}
  >{\raggedright\arraybackslash}p{(\linewidth - 14\tabcolsep) * \real{0.1250}}
  >{\raggedright\arraybackslash}p{(\linewidth - 14\tabcolsep) * \real{0.1250}}
  >{\raggedright\arraybackslash}p{(\linewidth - 14\tabcolsep) * \real{0.1250}}
  >{\raggedright\arraybackslash}p{(\linewidth - 14\tabcolsep) * \real{0.1250}}
  >{\raggedright\arraybackslash}p{(\linewidth - 14\tabcolsep) * \real{0.1250}}
  >{\raggedright\arraybackslash}p{(\linewidth - 14\tabcolsep) * \real{0.1250}}
  >{\raggedright\arraybackslash}p{(\linewidth - 14\tabcolsep) * \real{0.1250}}
  >{\raggedright\arraybackslash}p{(\linewidth - 14\tabcolsep) * \real{0.1250}}@{}}
\toprule\noalign{}
\begin{minipage}[b]{\linewidth}\raggedright
RUNSPEC
\end{minipage} & \begin{minipage}[b]{\linewidth}\raggedright
GRID
\end{minipage} & \begin{minipage}[b]{\linewidth}\raggedright
EDIT
\end{minipage} & \begin{minipage}[b]{\linewidth}\raggedright
PROPS
\end{minipage} & \begin{minipage}[b]{\linewidth}\raggedright
REGIONS
\end{minipage} & \begin{minipage}[b]{\linewidth}\raggedright
SOLUTION
\end{minipage} & \begin{minipage}[b]{\linewidth}\raggedright
SUMMARY
\end{minipage} & \begin{minipage}[b]{\linewidth}\raggedright
SCHEDULE
\end{minipage} \\
\midrule\noalign{}
\endhead
\bottomrule\noalign{}
\endlastfoot
\end{longtable}

\subsubsection{Description}\label{description-451}

The GINODE keyword defines the Gi node values used when the GIMODEL
keyword in the RUNSPEC section has been used to activate the GI Pseudo
Compositional option for the run. The keyword is used in conjunction
with the RSGI, RVGI, BGGI and BOGI keywords in the PROPS section to
describe the fluid properties for the GI Pseudo Compositional option.

This keyword is not supported by OPM Flow but would change the results
if supported so the simulation will be stopped.

\subsection{GRAVCONS -- Re-Define Gravity
Constant}\label{gravcons-re-define-gravity-constant}

\begin{longtable}[]{@{}
  >{\raggedright\arraybackslash}p{(\linewidth - 14\tabcolsep) * \real{0.1250}}
  >{\raggedright\arraybackslash}p{(\linewidth - 14\tabcolsep) * \real{0.1250}}
  >{\raggedright\arraybackslash}p{(\linewidth - 14\tabcolsep) * \real{0.1250}}
  >{\raggedright\arraybackslash}p{(\linewidth - 14\tabcolsep) * \real{0.1250}}
  >{\raggedright\arraybackslash}p{(\linewidth - 14\tabcolsep) * \real{0.1250}}
  >{\raggedright\arraybackslash}p{(\linewidth - 14\tabcolsep) * \real{0.1250}}
  >{\raggedright\arraybackslash}p{(\linewidth - 14\tabcolsep) * \real{0.1250}}
  >{\raggedright\arraybackslash}p{(\linewidth - 14\tabcolsep) * \real{0.1250}}@{}}
\toprule\noalign{}
\begin{minipage}[b]{\linewidth}\raggedright
RUNSPEC
\end{minipage} & \begin{minipage}[b]{\linewidth}\raggedright
GRID
\end{minipage} & \begin{minipage}[b]{\linewidth}\raggedright
EDIT
\end{minipage} & \begin{minipage}[b]{\linewidth}\raggedright
PROPS
\end{minipage} & \begin{minipage}[b]{\linewidth}\raggedright
REGIONS
\end{minipage} & \begin{minipage}[b]{\linewidth}\raggedright
SOLUTION
\end{minipage} & \begin{minipage}[b]{\linewidth}\raggedright
SUMMARY
\end{minipage} & \begin{minipage}[b]{\linewidth}\raggedright
SCHEDULE
\end{minipage} \\
\midrule\noalign{}
\endhead
\bottomrule\noalign{}
\endlastfoot
\end{longtable}

\subsubsection{Description}\label{description-452}

The GRAVCONS keyword re-defines the gravity constant used in various
calculations from the default value used by the simulator. Normally this
keyword should not be used.

This keyword is not supported by OPM Flow but would change the results
if supported so the simulation will be stopped.

\begin{longtable}[]{@{}
  >{\raggedright\arraybackslash}p{(\linewidth - 6\tabcolsep) * \real{0.2500}}
  >{\raggedright\arraybackslash}p{(\linewidth - 6\tabcolsep) * \real{0.2500}}
  >{\raggedright\arraybackslash}p{(\linewidth - 6\tabcolsep) * \real{0.2500}}
  >{\raggedright\arraybackslash}p{(\linewidth - 6\tabcolsep) * \real{0.2500}}@{}}
\toprule\noalign{}
\begin{minipage}[b]{\linewidth}\raggedright
No.
\end{minipage} & \begin{minipage}[b]{\linewidth}\raggedright
Name
\end{minipage} & \begin{minipage}[b]{\linewidth}\raggedright
Description
\end{minipage} & \begin{minipage}[b]{\linewidth}\raggedright
Default
\end{minipage} \\
\midrule\noalign{}
\endhead
\bottomrule\noalign{}
\endlastfoot
Field & Metric & Laboratory & \\
1 & GRAVCONS & GRAVCONS is a positive real number number that defines
the gravity constant used in various calculations. & Defined \\
ft2psi/lb 0.00694 & m2bars/kg 0.0000981 & cm2atm/gm 0.000968 & \\
Notes: & & & \\
\end{longtable}

\emph{Table 8.42: GRAVCONS Keyword Description}

\subsubsection{Example}\label{example-281}

\begin{verbatim}
--
--       RE-DEFINE GRAVITY CONSTANT
--
GRAVITY
         0.0000980665                                      /
\end{verbatim}

The above example re-defines the gravity constant to be 0.0000980665
ft2psi/lb from the default value of 0.00694 ft2psi/lb.

\subsection{GRAVITY -- Define the Surface Oil, Water Gas Gravities for
the
Fluids}\label{gravity-define-the-surface-oil-water-gas-gravities-for-the-fluids}

\begin{longtable}[]{@{}
  >{\raggedright\arraybackslash}p{(\linewidth - 14\tabcolsep) * \real{0.1250}}
  >{\raggedright\arraybackslash}p{(\linewidth - 14\tabcolsep) * \real{0.1250}}
  >{\raggedright\arraybackslash}p{(\linewidth - 14\tabcolsep) * \real{0.1250}}
  >{\raggedright\arraybackslash}p{(\linewidth - 14\tabcolsep) * \real{0.1250}}
  >{\raggedright\arraybackslash}p{(\linewidth - 14\tabcolsep) * \real{0.1250}}
  >{\raggedright\arraybackslash}p{(\linewidth - 14\tabcolsep) * \real{0.1250}}
  >{\raggedright\arraybackslash}p{(\linewidth - 14\tabcolsep) * \real{0.1250}}
  >{\raggedright\arraybackslash}p{(\linewidth - 14\tabcolsep) * \real{0.1250}}@{}}
\toprule\noalign{}
\begin{minipage}[b]{\linewidth}\raggedright
RUNSPEC
\end{minipage} & \begin{minipage}[b]{\linewidth}\raggedright
GRID
\end{minipage} & \begin{minipage}[b]{\linewidth}\raggedright
EDIT
\end{minipage} & \begin{minipage}[b]{\linewidth}\raggedright
PROPS
\end{minipage} & \begin{minipage}[b]{\linewidth}\raggedright
REGIONS
\end{minipage} & \begin{minipage}[b]{\linewidth}\raggedright
SOLUTION
\end{minipage} & \begin{minipage}[b]{\linewidth}\raggedright
SUMMARY
\end{minipage} & \begin{minipage}[b]{\linewidth}\raggedright
SCHEDULE
\end{minipage} \\
\midrule\noalign{}
\endhead
\bottomrule\noalign{}
\endlastfoot
\end{longtable}

\subsubsection{Description}\label{description-453}

GRAVITY defines the oil API gravity and water and gas surface specific
gravities for the fluids for various regions in the model. The number of
GRAVITY vector data sets is defined by the NTPVT parameter on the
TABDIMS keyword in the RUNSPEC section and the allocation of the GRAVITY
data sets to different grid blocks in the model is done via the PVTNUM
keyword in the REGION section. One data set consists of one record or
line which is terminated by a ``/''.

This surface density or gravity must be entered using either the DENSITY
or GRAVITY keywords irrespective of which phases are active in the
model.

\begin{longtable}[]{@{}
  >{\raggedright\arraybackslash}p{(\linewidth - 6\tabcolsep) * \real{0.2500}}
  >{\raggedright\arraybackslash}p{(\linewidth - 6\tabcolsep) * \real{0.2500}}
  >{\raggedright\arraybackslash}p{(\linewidth - 6\tabcolsep) * \real{0.2500}}
  >{\raggedright\arraybackslash}p{(\linewidth - 6\tabcolsep) * \real{0.2500}}@{}}
\toprule\noalign{}
\begin{minipage}[b]{\linewidth}\raggedright
No.
\end{minipage} & \begin{minipage}[b]{\linewidth}\raggedright
Name
\end{minipage} & \begin{minipage}[b]{\linewidth}\raggedright
Description
\end{minipage} & \begin{minipage}[b]{\linewidth}\raggedright
Default
\end{minipage} \\
\midrule\noalign{}
\endhead
\bottomrule\noalign{}
\endlastfoot
Field & Metric & Laboratory & \\
1 & OILAPI & OILAPI is a real number defining the API gravity of the oil
phase at surface conditions. The American Petroleum Institute (``API'')
classifies oils based on an API gravity (γAPI), or degrees API (oAPI),
the relationship between relative density (γo) of oil and API gravity
(γAPI) is given by:
\({\mathrm{γ}}_{\mathit{API}} = \frac{141.5}{{\mathrm{γ}}_{o}} - 131.5\)
& None \\
oAPI & oAPI & oAPI & \\
2 & WATGRAV & WATGRAV is a real number defining the specific gravity of
the water phase relative to pure water at surface conditions. &
Defined \\
(water =1.0) 0.7773 & (water =1.0) 0.7773 & (water =1.0) 0.7773 & \\
3 & GASGRAV & GASGRAV is a real number defining the specific gravity of
the gas phase relative to air at surface conditions. & Defined \\
(air =1.0) 1.000 & (air =1.0) 1.000 & (air =1.0) 1.000 & \\
Notes: & & & \\
\end{longtable}

\emph{Table 8.43: GRAVITY Keyword Description}

According to the SPE SI standard {[}The SI Metric System of Units and
SPE Metric Standard, Adopted for Use as a Voluntary Standard by the SPE
Board of Directors, June 1983, Society of Petroleum Engineers.{]},
Relative Density (γ) replaces Specific Gravity as the term used to
define the ratio of the density of a known material to the density of
reference material, at standard conditions of pressure and temperature.
Standard conditions vary throughout the world, but for oil field units
one normally uses14.7 psia and 60 oF, while for SI units some areas use
101.325 kPa and 15 oC.

Note that provided the first table has been entered, subsequent tables
may be defaulted, in this case the prior table is copied to the current
table. See the third example for an illustration on how to use this
feature.

See also the DENSITY keyword in the PROPS section.

\subsubsection{Examples}\label{examples-24}

The following shows the GRAVITY keyword for when NTPVT on the TABDIMS
keyword in the RUNSPEC section is set to one.

\begin{verbatim}
--
--       OIL      WAT        GAS
--       GRAVITY  GRAVITY    GRAVITY
--       -------  -------    -------
GRAVITY
          39.0     1.012      0.650                / GRAVITY PVT DATA REGION 1
\end{verbatim}

The next example shows the GRAVITY keyword for when NTPVT on the TABDIMS
keyword in the RUNSPEC section is set to three.

\begin{verbatim}
--
--       OIL      WAT        GAS
--       GRAVITY  GRAVITY    GRAVITY
--       -------  -------    -------
GRAVITY
          37.0     1.012      0.650                / GRAVITY PVT DATA REGION 1
          38.0     1.012      0.646                / GRAVITY PVT DATA REGION 2
          39.0     1.012      0.640                / GRAVITY PVT DATA REGION 3
\end{verbatim}

The third and final example shows the GRAVITY keyword for when NTPVT on
the TABDIMS keyword in the RUNSPEC section is set to four. Here table
two defaults to table one, and table four defaults to table three.

\begin{verbatim}
--
--       OIL      WAT        GAS
--       GRAVITY  GRAVITY    GRAVITY
--       -------  -------    -------
GRAVITY
          37.0     1.012      0.650                / GRAVITY PVT DATA REGION 1
                                                   / GRAVITY PVT DATA REGION 2
          38.0     1.012      0.646                / GRAVITY PVT DATA REGION 3
                                                   / GRAVITY PVT DATA REGION 4
\end{verbatim}

Again, note that there is no terminating ``/'' for this keyword.

\subsection{GSF -- Gas Saturation Function Tables (Gas-Water
Systems)}\label{gsf-gas-saturation-function-tables-gas-water-systems}

\begin{longtable}[]{@{}
  >{\raggedright\arraybackslash}p{(\linewidth - 14\tabcolsep) * \real{0.1250}}
  >{\raggedright\arraybackslash}p{(\linewidth - 14\tabcolsep) * \real{0.1250}}
  >{\raggedright\arraybackslash}p{(\linewidth - 14\tabcolsep) * \real{0.1250}}
  >{\raggedright\arraybackslash}p{(\linewidth - 14\tabcolsep) * \real{0.1250}}
  >{\raggedright\arraybackslash}p{(\linewidth - 14\tabcolsep) * \real{0.1250}}
  >{\raggedright\arraybackslash}p{(\linewidth - 14\tabcolsep) * \real{0.1250}}
  >{\raggedright\arraybackslash}p{(\linewidth - 14\tabcolsep) * \real{0.1250}}
  >{\raggedright\arraybackslash}p{(\linewidth - 14\tabcolsep) * \real{0.1250}}@{}}
\toprule\noalign{}
\begin{minipage}[b]{\linewidth}\raggedright
RUNSPEC
\end{minipage} & \begin{minipage}[b]{\linewidth}\raggedright
GRID
\end{minipage} & \begin{minipage}[b]{\linewidth}\raggedright
EDIT
\end{minipage} & \begin{minipage}[b]{\linewidth}\raggedright
PROPS
\end{minipage} & \begin{minipage}[b]{\linewidth}\raggedright
REGIONS
\end{minipage} & \begin{minipage}[b]{\linewidth}\raggedright
SOLUTION
\end{minipage} & \begin{minipage}[b]{\linewidth}\raggedright
SUMMARY
\end{minipage} & \begin{minipage}[b]{\linewidth}\raggedright
SCHEDULE
\end{minipage} \\
\midrule\noalign{}
\endhead
\bottomrule\noalign{}
\endlastfoot
\end{longtable}

\subsubsection{Description}\label{description-454}

The GSF keyword defines the gas relative permeability and gas-water
capillary pressure data versus gas saturation tables for when only gas
and water are present in the input deck. This keyword should only be
used if the gas and water phases are present in the run, and can
therefore also be used with the CO2STORE and
\hyperref[REF_HEADING_KEYWORD_H2STORE]{H2STORE} models. In addition, the
keyword must be used in conjunction with the WSF keyword in the PROPS
section, that defines the water relative permeability versus water
saturation for gas-water systems.

\begin{tcolorbox}[enhanced jigsaw, title=\textcolor{quarto-callout-note-color}{\faInfo}\hspace{0.5em}{Note}, coltitle=black, breakable, toprule=.15mm, colbacktitle=quarto-callout-note-color!10!white, leftrule=.75mm, opacitybacktitle=0.6, colback=white, bottomtitle=1mm, toptitle=1mm, titlerule=0mm, left=2mm, colframe=quarto-callout-note-color-frame, arc=.35mm, rightrule=.15mm, bottomrule=.15mm, opacityback=0]

GSF is a compositional keyword in the commercial compositional
simulator, and will therefore cause an error in the commercial black-oil
simulator. Currently, both the GSF and WSF keywords can only be used
with the CO2STORE and \hyperref[REF_HEADING_KEYWORD_H2STORE]{H2STORE}
models.

\end{tcolorbox}

\begin{longtable}[]{@{}
  >{\raggedright\arraybackslash}p{(\linewidth - 6\tabcolsep) * \real{0.2500}}
  >{\raggedright\arraybackslash}p{(\linewidth - 6\tabcolsep) * \real{0.2500}}
  >{\raggedright\arraybackslash}p{(\linewidth - 6\tabcolsep) * \real{0.2500}}
  >{\raggedright\arraybackslash}p{(\linewidth - 6\tabcolsep) * \real{0.2500}}@{}}
\toprule\noalign{}
\begin{minipage}[b]{\linewidth}\raggedright
No.
\end{minipage} & \begin{minipage}[b]{\linewidth}\raggedright
Name
\end{minipage} & \begin{minipage}[b]{\linewidth}\raggedright
Description
\end{minipage} & \begin{minipage}[b]{\linewidth}\raggedright
Default
\end{minipage} \\
\midrule\noalign{}
\endhead
\bottomrule\noalign{}
\endlastfoot
Field & Metric & Laboratory & \\
1 & SGAS & A columnar vector of real values monotonically increasing
down the column starting from zero and terminating at one minus the
connate water saturation, that defines the gas saturation. & None \\
dimensionless & dimensionless & dimensionless & \\
2 & KRG & A columnar vector of real values that are either equal or
increasing down the column and that are greater than or equal to zero
and less than or equal to one that defines the gas relative
permeability. Note that the first entry in the column must be zero. &
None \\
dimensionless & dimensionless & dimensionless & \\
3 & PCWG & A columnar vector of real values that are either equal or
increasing down the column that defines the gas-water capillary
pressure. & None \\
psia & bars & atm & \\
Notes: & & & \\
\end{longtable}

\emph{Table 8.44: GSF Keyword Description}

See also the WSF - Water Saturation Tables versus Water Saturation
(Gas-Water and CO2STORE Systems) keyword in the PROPS section, that
defines the water saturation as a function of water saturation for when
only the gas and water phases are present in the model.

\subsubsection{Example}\label{example-282}

\begin{verbatim}
--
--       GAS RELATIVE PERMEABILITY TABLES (OPM FLOW KEYWORD)
--
GSF
--       SGAS       KRG        PCGW
--       FRAC                  PSIA
--       --------  --------   -------
             0.00    0.0000      0.0
             0.05    0.0000      0.0
             0.10    0.0000      0.0
             0.15    0.0000      0.0
             0.20    0.0002      0.0
             0.25    0.0010      0.0
             0.30    0.0062      0.0
             0.35    0.0140      0.0
             0.40    0.0273      0.0
             0.45    0.0450      0.0
             0.50    0.0707      0.0
             0.55    0.1020      0.0
             0.60    0.1412      0.0
             0.65    0.1870      0.0
             0.70    0.2412      0.0
             0.77    0.3288      0.0
             0.82    0.4000      0.0
             0.85    0.4450      0.0                       / TABLE NO. 01
--       SGAS       KRG        PCGW
--       FRAC                  PSIA
--       --------  --------   -------
             0.00    0.0000      0.0
             0.05    0.0000      0.0
             0.10    0.0000      0.0
             0.15    0.0000      0.0
             0.20    0.0002      0.0
             0.25    0.0010      0.0
             0.30    0.0062      0.0
             0.35    0.0140      0.0
             0.40    0.0273      0.0
             0.45    0.0450      0.0
             0.50    0.0707      0.0
             0.55    0.1020      0.0
             0.60    0.1412      0.0
             0.65    0.1870      0.0
             0.70    0.2412      0.0
             0.77    0.3288      0.0
             0.82    0.4000      0.0
             0.85    0.4450      0.0                       / TABLE NO. 02
\end{verbatim}

The example defines two GSF tables for when gas and water are present in
the input deck. In the tables the gas-water capillary pressure data has
been set to zero.

\subsection{HA -- History Match End-Point Gradient Additive
Modifier}\label{ha-history-match-end-point-gradient-additive-modifier}

\begin{longtable}[]{@{}
  >{\raggedright\arraybackslash}p{(\linewidth - 14\tabcolsep) * \real{0.1250}}
  >{\raggedright\arraybackslash}p{(\linewidth - 14\tabcolsep) * \real{0.1250}}
  >{\raggedright\arraybackslash}p{(\linewidth - 14\tabcolsep) * \real{0.1250}}
  >{\raggedright\arraybackslash}p{(\linewidth - 14\tabcolsep) * \real{0.1250}}
  >{\raggedright\arraybackslash}p{(\linewidth - 14\tabcolsep) * \real{0.1250}}
  >{\raggedright\arraybackslash}p{(\linewidth - 14\tabcolsep) * \real{0.1250}}
  >{\raggedright\arraybackslash}p{(\linewidth - 14\tabcolsep) * \real{0.1250}}
  >{\raggedright\arraybackslash}p{(\linewidth - 14\tabcolsep) * \real{0.1250}}@{}}
\toprule\noalign{}
\begin{minipage}[b]{\linewidth}\raggedright
RUNSPEC
\end{minipage} & \begin{minipage}[b]{\linewidth}\raggedright
GRID
\end{minipage} & \begin{minipage}[b]{\linewidth}\raggedright
EDIT
\end{minipage} & \begin{minipage}[b]{\linewidth}\raggedright
PROPS
\end{minipage} & \begin{minipage}[b]{\linewidth}\raggedright
REGIONS
\end{minipage} & \begin{minipage}[b]{\linewidth}\raggedright
SOLUTION
\end{minipage} & \begin{minipage}[b]{\linewidth}\raggedright
SUMMARY
\end{minipage} & \begin{minipage}[b]{\linewidth}\raggedright
SCHEDULE
\end{minipage} \\
\midrule\noalign{}
\endhead
\bottomrule\noalign{}
\endlastfoot
\end{longtable}

\subsubsection{Description}\label{description-455}

The HA series of keywords defines the history match end-point gradient
parameters used to set the additive cumulative end point data, for when
the History Match Gradient option has been activated by the HMDIMS
keyword in the RUNSPEC section. In addition, the End-Point Scaling
option must also be activated by the ENDSCALE keyword which is also in
the RUNSPEC section. The keyword consists of the first two characters of
``HA'' followed by the end-point keyword shown in Table 8.45, for
example, HASWL.

This keyword is not supported by OPM Flow but has no effect on the
results so it will be ignored.

\begin{longtable}[]{@{}
  >{\raggedright\arraybackslash}p{(\linewidth - 4\tabcolsep) * \real{0.3333}}
  >{\raggedright\arraybackslash}p{(\linewidth - 4\tabcolsep) * \real{0.3333}}
  >{\raggedright\arraybackslash}p{(\linewidth - 4\tabcolsep) * \real{0.3333}}@{}}
\toprule\noalign{}
\begin{minipage}[b]{\linewidth}\raggedright
Type
\end{minipage} & \begin{minipage}[b]{\linewidth}\raggedright
End-Point Keyword
\end{minipage} & \begin{minipage}[b]{\linewidth}\raggedright
Oil-Water End-Point Definitions
\end{minipage} \\
\midrule\noalign{}
\endhead
\bottomrule\noalign{}
\endlastfoot
Saturation & SWL & Connate water saturation, that is the smallest water
saturation in a water saturation function table. \\
SWCR & Critical water saturation, that is the largest water saturation
for which the water relative permeability is zero. & \\
SOWCR & Critical oil-in-water saturation, that is the largest oil
saturation for which the oil relative permeability is zero in an
oil-water system. & \\
Relative Permeability & KRW & Relative permeability of water at the
maximum water saturation (normally the maximum water saturation is
one). \\
KRO & Relative permeability of oil at the maximum oil saturation. & \\
KRWR & Relative permeability of water at the residual oil saturation or
the residual gas saturation in a gas-water run. & \\
KRORW & Relative permeability of oil at the critical water saturation.
& \\
Capillary Pressure & SWLPC & Capillary pressure connate water
saturation, that is the smallest water saturation in a water saturation
function table. \\
Type & End-Point Keyword & Gas-Oil End-Point Definitions \\
Saturation & SGL & Connate gas saturation, that is the smallest gas
saturation in a gas saturation function table. \\
SGCR & Critical gas saturation, that is the largest gas saturation for
which the gas relative permeability is zero. & \\
SOGCR & Critical oil-in-gas saturation, that is the largest oil
saturation for which the oil relative permeability is zero in an
oil-gas-connate water system. & \\
Relative Permeability & KRG & Relative permeability of gas at the
maximum gas saturation. \\
KRGR & Relative permeability of gas at the residual oil saturation or
the critical water saturation in a gas-water run. & \\
KRORG & Relative permeability of oil at the critical gas saturation.
& \\
Capillary Pressure & SGLPC & Capillary pressure connate gas saturation,
that is the smallest gas saturation in a gas saturation function
table. \\
\end{longtable}

\emph{Table 8.45: HA Keyword List}

See also the HMPROPS keyword in the PROPS section that allows the use of
the ADD, BOX, EQUALS, COPY, MINVALUE, and MAXVALUE keywords to be used
with the HA and HM series of keywords.

\subsection{HDISP -- Define Tracer Mechanical Dispersivity
Parameters}\label{hdisp-define-tracer-mechanical-dispersivity-parameters}

\begin{longtable}[]{@{}
  >{\raggedright\arraybackslash}p{(\linewidth - 14\tabcolsep) * \real{0.1250}}
  >{\raggedright\arraybackslash}p{(\linewidth - 14\tabcolsep) * \real{0.1250}}
  >{\raggedright\arraybackslash}p{(\linewidth - 14\tabcolsep) * \real{0.1250}}
  >{\raggedright\arraybackslash}p{(\linewidth - 14\tabcolsep) * \real{0.1250}}
  >{\raggedright\arraybackslash}p{(\linewidth - 14\tabcolsep) * \real{0.1250}}
  >{\raggedright\arraybackslash}p{(\linewidth - 14\tabcolsep) * \real{0.1250}}
  >{\raggedright\arraybackslash}p{(\linewidth - 14\tabcolsep) * \real{0.1250}}
  >{\raggedright\arraybackslash}p{(\linewidth - 14\tabcolsep) * \real{0.1250}}@{}}
\toprule\noalign{}
\begin{minipage}[b]{\linewidth}\raggedright
RUNSPEC
\end{minipage} & \begin{minipage}[b]{\linewidth}\raggedright
GRID
\end{minipage} & \begin{minipage}[b]{\linewidth}\raggedright
EDIT
\end{minipage} & \begin{minipage}[b]{\linewidth}\raggedright
PROPS
\end{minipage} & \begin{minipage}[b]{\linewidth}\raggedright
REGIONS
\end{minipage} & \begin{minipage}[b]{\linewidth}\raggedright
SOLUTION
\end{minipage} & \begin{minipage}[b]{\linewidth}\raggedright
SUMMARY
\end{minipage} & \begin{minipage}[b]{\linewidth}\raggedright
SCHEDULE
\end{minipage} \\
\midrule\noalign{}
\endhead
\bottomrule\noalign{}
\endlastfoot
\end{longtable}

\subsubsection{Description}\label{description-456}

The HDISP keyword is combined with three character tracer name,
specified by the TRACER keyword in the PROPS section, to define the
tracer's mechanical dispersivity parameters.

This keyword is not supported by OPM Flow but would change the results
if supported so the simulation will be stopped.

\subsection{HM -- History Match End-Point Gradient Multiplicative
Modifier}\label{hm-history-match-end-point-gradient-multiplicative-modifier}

\begin{longtable}[]{@{}
  >{\raggedright\arraybackslash}p{(\linewidth - 14\tabcolsep) * \real{0.1250}}
  >{\raggedright\arraybackslash}p{(\linewidth - 14\tabcolsep) * \real{0.1250}}
  >{\raggedright\arraybackslash}p{(\linewidth - 14\tabcolsep) * \real{0.1250}}
  >{\raggedright\arraybackslash}p{(\linewidth - 14\tabcolsep) * \real{0.1250}}
  >{\raggedright\arraybackslash}p{(\linewidth - 14\tabcolsep) * \real{0.1250}}
  >{\raggedright\arraybackslash}p{(\linewidth - 14\tabcolsep) * \real{0.1250}}
  >{\raggedright\arraybackslash}p{(\linewidth - 14\tabcolsep) * \real{0.1250}}
  >{\raggedright\arraybackslash}p{(\linewidth - 14\tabcolsep) * \real{0.1250}}@{}}
\toprule\noalign{}
\begin{minipage}[b]{\linewidth}\raggedright
RUNSPEC
\end{minipage} & \begin{minipage}[b]{\linewidth}\raggedright
GRID
\end{minipage} & \begin{minipage}[b]{\linewidth}\raggedright
EDIT
\end{minipage} & \begin{minipage}[b]{\linewidth}\raggedright
PROPS
\end{minipage} & \begin{minipage}[b]{\linewidth}\raggedright
REGIONS
\end{minipage} & \begin{minipage}[b]{\linewidth}\raggedright
SOLUTION
\end{minipage} & \begin{minipage}[b]{\linewidth}\raggedright
SUMMARY
\end{minipage} & \begin{minipage}[b]{\linewidth}\raggedright
SCHEDULE
\end{minipage} \\
\midrule\noalign{}
\endhead
\bottomrule\noalign{}
\endlastfoot
\end{longtable}

\subsubsection{Description}\label{description-457}

The HM series of keywords defines the history match end-point gradient
parameters used to set the multiplicative cumulative end point data, for
when the History Match Gradient option has been activated by the HMDIMS
keyword in the RUNSPEC section. In addition, the End-Point Scaling
option must also be activated by the ENDSCALE keyword which is also in
the RUNSPEC section. The keyword consists of the first two characters of
``HM'' followed by the end-point keyword shown in Table 8.46, for
example, HMSWL.

This keyword is not supported by OPM Flow but has no effect on the
results so it will be ignored.

\begin{longtable}[]{@{}
  >{\raggedright\arraybackslash}p{(\linewidth - 4\tabcolsep) * \real{0.3333}}
  >{\raggedright\arraybackslash}p{(\linewidth - 4\tabcolsep) * \real{0.3333}}
  >{\raggedright\arraybackslash}p{(\linewidth - 4\tabcolsep) * \real{0.3333}}@{}}
\toprule\noalign{}
\begin{minipage}[b]{\linewidth}\raggedright
Type
\end{minipage} & \begin{minipage}[b]{\linewidth}\raggedright
End-Point Keyword
\end{minipage} & \begin{minipage}[b]{\linewidth}\raggedright
Oil-Water End-Point Definitions
\end{minipage} \\
\midrule\noalign{}
\endhead
\bottomrule\noalign{}
\endlastfoot
Saturation & SWL & Connate water saturation, that is the smallest water
saturation in a water saturation function table. \\
SWCR & Critical water saturation, that is the largest water saturation
for which the water relative permeability is zero. & \\
SOWCR & Critical oil-in-water saturation, that is the largest oil
saturation for which the oil relative permeability is zero in an
oil-water system. & \\
Relative Permeability & KRW & Relative permeability of water at the
maximum water saturation (normally the maximum water saturation is
one). \\
KRO & Relative permeability of oil at the maximum oil saturation. & \\
KRWR & Relative permeability of water at the residual oil saturation or
the residual gas saturation in a gas-water run. & \\
KRORW & Relative permeability of oil at the critical water saturation.
& \\
Capillary Pressure & SWLPC & Capillary pressure connate water
saturation, that is the smallest water saturation in a water saturation
function table. \\
Type & End-Point Keyword & Gas-Oil End-Point Definitions \\
Saturation & SGL & Connate gas saturation, that is the smallest gas
saturation in a gas saturation function table. \\
SGCR & Critical gas saturation, that is the largest gas saturation for
which the gas relative permeability is zero. & \\
SOGCR & Critical oil-in-gas saturation, that is the largest oil
saturation for which the oil relative permeability is zero in an
oil-gas-connate water system. & \\
Relative Permeability & KRG & Relative permeability of gas at the
maximum gas saturation. \\
KRGR & Relative permeability of gas at the residual oil saturation or
the critical water saturation in a gas-water run. & \\
KRORG & Relative permeability of oil at the critical gas saturation.
& \\
Capillary Pressure & SGLPC & Capillary pressure connate gas saturation,
that is the smallest gas saturation in a gas saturation function
table. \\
\end{longtable}

\emph{Table 8.46: HM Keyword List}

See also the HMPROPS keyword in the PROPS section that allows the use of
the ADD, BOX, EQUALS, COPY, MINVALUE, and MAXVALUE keywords to be used
with the HA and HM series of keywords.

\subsection{HMMROCK -- History Match Rock Compressibility Gradient
Cumulative
Multipliers}\label{hmmrock-history-match-rock-compressibility-gradient-cumulative-multipliers}

\begin{longtable}[]{@{}
  >{\raggedright\arraybackslash}p{(\linewidth - 14\tabcolsep) * \real{0.1250}}
  >{\raggedright\arraybackslash}p{(\linewidth - 14\tabcolsep) * \real{0.1250}}
  >{\raggedright\arraybackslash}p{(\linewidth - 14\tabcolsep) * \real{0.1250}}
  >{\raggedright\arraybackslash}p{(\linewidth - 14\tabcolsep) * \real{0.1250}}
  >{\raggedright\arraybackslash}p{(\linewidth - 14\tabcolsep) * \real{0.1250}}
  >{\raggedright\arraybackslash}p{(\linewidth - 14\tabcolsep) * \real{0.1250}}
  >{\raggedright\arraybackslash}p{(\linewidth - 14\tabcolsep) * \real{0.1250}}
  >{\raggedright\arraybackslash}p{(\linewidth - 14\tabcolsep) * \real{0.1250}}@{}}
\toprule\noalign{}
\begin{minipage}[b]{\linewidth}\raggedright
RUNSPEC
\end{minipage} & \begin{minipage}[b]{\linewidth}\raggedright
GRID
\end{minipage} & \begin{minipage}[b]{\linewidth}\raggedright
EDIT
\end{minipage} & \begin{minipage}[b]{\linewidth}\raggedright
PROPS
\end{minipage} & \begin{minipage}[b]{\linewidth}\raggedright
REGIONS
\end{minipage} & \begin{minipage}[b]{\linewidth}\raggedright
SOLUTION
\end{minipage} & \begin{minipage}[b]{\linewidth}\raggedright
SUMMARY
\end{minipage} & \begin{minipage}[b]{\linewidth}\raggedright
SCHEDULE
\end{minipage} \\
\midrule\noalign{}
\endhead
\bottomrule\noalign{}
\endlastfoot
\end{longtable}

\subsubsection{Description}\label{description-458}

HMMROCK defines the rock compressibility gradient cumulative multipliers
to be applied to the rock compressibility as defined by the ROCK keyword
in the PROPS section, for when the History Match Gradient option has
been activated by the HMDIMS keyword in the RUNSPEC section. The
constant should be a real number.

The allocation of the ROCK tables to different grid blocks in the model
is done via the PVTNUM or the SATNUM keywords in the REGION section.

This keyword is not supported by OPM Flow but has no effect on the
results so it will be ignored.

\subsection{HMMROCKT -- History Match Rock Compaction Gradient
Cumulative
Multipliers}\label{hmmrockt-history-match-rock-compaction-gradient-cumulative-multipliers}

\begin{longtable}[]{@{}
  >{\raggedright\arraybackslash}p{(\linewidth - 14\tabcolsep) * \real{0.1250}}
  >{\raggedright\arraybackslash}p{(\linewidth - 14\tabcolsep) * \real{0.1250}}
  >{\raggedright\arraybackslash}p{(\linewidth - 14\tabcolsep) * \real{0.1250}}
  >{\raggedright\arraybackslash}p{(\linewidth - 14\tabcolsep) * \real{0.1250}}
  >{\raggedright\arraybackslash}p{(\linewidth - 14\tabcolsep) * \real{0.1250}}
  >{\raggedright\arraybackslash}p{(\linewidth - 14\tabcolsep) * \real{0.1250}}
  >{\raggedright\arraybackslash}p{(\linewidth - 14\tabcolsep) * \real{0.1250}}
  >{\raggedright\arraybackslash}p{(\linewidth - 14\tabcolsep) * \real{0.1250}}@{}}
\toprule\noalign{}
\begin{minipage}[b]{\linewidth}\raggedright
RUNSPEC
\end{minipage} & \begin{minipage}[b]{\linewidth}\raggedright
GRID
\end{minipage} & \begin{minipage}[b]{\linewidth}\raggedright
EDIT
\end{minipage} & \begin{minipage}[b]{\linewidth}\raggedright
PROPS
\end{minipage} & \begin{minipage}[b]{\linewidth}\raggedright
REGIONS
\end{minipage} & \begin{minipage}[b]{\linewidth}\raggedright
SOLUTION
\end{minipage} & \begin{minipage}[b]{\linewidth}\raggedright
SUMMARY
\end{minipage} & \begin{minipage}[b]{\linewidth}\raggedright
SCHEDULE
\end{minipage} \\
\midrule\noalign{}
\endhead
\bottomrule\noalign{}
\endlastfoot
\end{longtable}

\subsubsection{Description}\label{description-459}

HMMROCKT defines the rock compaction gradient cumulative multipliers to
be applied to the compaction data entered by the ROCTAB or ROCKTABH
keywords in the PRROPS section, for when the History Match Gradient
option has been activated by the HMDIMS keyword in the RUNSPEC section

The ROCKTAB keyword defines the rock compaction attributes to be applied
for when the rock compaction option has been invoked by the ROCKCOMP
keyword in the RUNSPEC section. ROCKTAB defines pore volume and
transmissibility multipliers versus pressure that are used in the
compaction calculations. If the RKTRMDIR has been activated in the PROPS
section, then the transmissibility multiplier is directional dependent

This keyword should only be used if compaction option has been enabled.

This keyword is not supported by OPM Flow but has no effect on the
results so it will be ignored.

\subsection{HMPROPS -- History Match End-Point Section
Start}\label{hmprops-history-match-end-point-section-start}

\begin{longtable}[]{@{}
  >{\raggedright\arraybackslash}p{(\linewidth - 14\tabcolsep) * \real{0.1250}}
  >{\raggedright\arraybackslash}p{(\linewidth - 14\tabcolsep) * \real{0.1250}}
  >{\raggedright\arraybackslash}p{(\linewidth - 14\tabcolsep) * \real{0.1250}}
  >{\raggedright\arraybackslash}p{(\linewidth - 14\tabcolsep) * \real{0.1250}}
  >{\raggedright\arraybackslash}p{(\linewidth - 14\tabcolsep) * \real{0.1250}}
  >{\raggedright\arraybackslash}p{(\linewidth - 14\tabcolsep) * \real{0.1250}}
  >{\raggedright\arraybackslash}p{(\linewidth - 14\tabcolsep) * \real{0.1250}}
  >{\raggedright\arraybackslash}p{(\linewidth - 14\tabcolsep) * \real{0.1250}}@{}}
\toprule\noalign{}
\begin{minipage}[b]{\linewidth}\raggedright
RUNSPEC
\end{minipage} & \begin{minipage}[b]{\linewidth}\raggedright
GRID
\end{minipage} & \begin{minipage}[b]{\linewidth}\raggedright
EDIT
\end{minipage} & \begin{minipage}[b]{\linewidth}\raggedright
PROPS
\end{minipage} & \begin{minipage}[b]{\linewidth}\raggedright
REGIONS
\end{minipage} & \begin{minipage}[b]{\linewidth}\raggedright
SOLUTION
\end{minipage} & \begin{minipage}[b]{\linewidth}\raggedright
SUMMARY
\end{minipage} & \begin{minipage}[b]{\linewidth}\raggedright
SCHEDULE
\end{minipage} \\
\midrule\noalign{}
\endhead
\bottomrule\noalign{}
\endlastfoot
\end{longtable}

\subsubsection{Description}\label{description-460}

HMPROPS defines the start of a history match end-points section, for
when the History Match Gradient option has been activated by the HMDIMS
keyword in the RUNSPEC section. In addition, the End-Point Scaling
option must also be activated by the ENDSCALE keyword which is also in
the RUNSPEC section. The keyword allows for the BOX, EQUALS, COPY,
MINVALUE, MAXVALUE and ADD keywords to be used with the HA and HM series
of keywords that reference the end-point scaling arrays, that is: HMKRG,
HMKRGR, HMKRO, HMKRORG, HMKRORW, HMKRW, HMKRWR, HMPCW, HMPCG, HMSGCR,
HMSOWCR, HMSOGCR, HMSWCR, and HMSWL keywords.

This keyword is not supported by OPM Flow but would change the results
if supported so the simulation will be stopped.

\subsection{HMROCK -- History Match Rock Compressibility Gradient
Parameters}\label{hmrock-history-match-rock-compressibility-gradient-parameters}

\begin{longtable}[]{@{}
  >{\raggedright\arraybackslash}p{(\linewidth - 14\tabcolsep) * \real{0.1250}}
  >{\raggedright\arraybackslash}p{(\linewidth - 14\tabcolsep) * \real{0.1250}}
  >{\raggedright\arraybackslash}p{(\linewidth - 14\tabcolsep) * \real{0.1250}}
  >{\raggedright\arraybackslash}p{(\linewidth - 14\tabcolsep) * \real{0.1250}}
  >{\raggedright\arraybackslash}p{(\linewidth - 14\tabcolsep) * \real{0.1250}}
  >{\raggedright\arraybackslash}p{(\linewidth - 14\tabcolsep) * \real{0.1250}}
  >{\raggedright\arraybackslash}p{(\linewidth - 14\tabcolsep) * \real{0.1250}}
  >{\raggedright\arraybackslash}p{(\linewidth - 14\tabcolsep) * \real{0.1250}}@{}}
\toprule\noalign{}
\begin{minipage}[b]{\linewidth}\raggedright
RUNSPEC
\end{minipage} & \begin{minipage}[b]{\linewidth}\raggedright
GRID
\end{minipage} & \begin{minipage}[b]{\linewidth}\raggedright
EDIT
\end{minipage} & \begin{minipage}[b]{\linewidth}\raggedright
PROPS
\end{minipage} & \begin{minipage}[b]{\linewidth}\raggedright
REGIONS
\end{minipage} & \begin{minipage}[b]{\linewidth}\raggedright
SOLUTION
\end{minipage} & \begin{minipage}[b]{\linewidth}\raggedright
SUMMARY
\end{minipage} & \begin{minipage}[b]{\linewidth}\raggedright
SCHEDULE
\end{minipage} \\
\midrule\noalign{}
\endhead
\bottomrule\noalign{}
\endlastfoot
\end{longtable}

\subsubsection{Description}\label{description-461}

The HMROCK keyword defines the history match rock compressibility
gradient parameters for when the History Match Gradient option has been
activated by the HMDIMS keyword in the RUNSPEC section

See also the HMDIMS keyword in the RUNSPEC section that specifies the
dimensions for the gradient option, including the maximum number of rock
gradient parameters that can be used with the History Match Gradient
option.

This keyword is not supported by OPM Flow but would change the results
if supported so the simulation will be stopped.

\subsection{HMROCKT -- History Match Rock Compaction Gradient
Parameters}\label{hmrockt-history-match-rock-compaction-gradient-parameters}

\begin{longtable}[]{@{}
  >{\raggedright\arraybackslash}p{(\linewidth - 14\tabcolsep) * \real{0.1250}}
  >{\raggedright\arraybackslash}p{(\linewidth - 14\tabcolsep) * \real{0.1250}}
  >{\raggedright\arraybackslash}p{(\linewidth - 14\tabcolsep) * \real{0.1250}}
  >{\raggedright\arraybackslash}p{(\linewidth - 14\tabcolsep) * \real{0.1250}}
  >{\raggedright\arraybackslash}p{(\linewidth - 14\tabcolsep) * \real{0.1250}}
  >{\raggedright\arraybackslash}p{(\linewidth - 14\tabcolsep) * \real{0.1250}}
  >{\raggedright\arraybackslash}p{(\linewidth - 14\tabcolsep) * \real{0.1250}}
  >{\raggedright\arraybackslash}p{(\linewidth - 14\tabcolsep) * \real{0.1250}}@{}}
\toprule\noalign{}
\begin{minipage}[b]{\linewidth}\raggedright
RUNSPEC
\end{minipage} & \begin{minipage}[b]{\linewidth}\raggedright
GRID
\end{minipage} & \begin{minipage}[b]{\linewidth}\raggedright
EDIT
\end{minipage} & \begin{minipage}[b]{\linewidth}\raggedright
PROPS
\end{minipage} & \begin{minipage}[b]{\linewidth}\raggedright
REGIONS
\end{minipage} & \begin{minipage}[b]{\linewidth}\raggedright
SOLUTION
\end{minipage} & \begin{minipage}[b]{\linewidth}\raggedright
SUMMARY
\end{minipage} & \begin{minipage}[b]{\linewidth}\raggedright
SCHEDULE
\end{minipage} \\
\midrule\noalign{}
\endhead
\bottomrule\noalign{}
\endlastfoot
\end{longtable}

\subsubsection{Description}\label{description-462}

The HMROCKT keyword defines the history match rock compaction gradient
parameters for when the History Match Gradient option has been activated
by the HMDIMS keyword in the RUNSPEC section, and the history match rock
compaction data has been entered via the HMMROCKT, ROCKTAB and ROCKTABH
keywords in the PROPS section.

See also the HMDIMS keyword in the RUNSPEC section that specifies the
dimensions for the gradient option, including the maximum number of rock
gradient parameters that can be used with the History Match Gradient
option.

This keyword is not supported by OPM Flow but would change the results
if supported so the simulation will be stopped.

\subsection{HMRREF -- History Match Rock Table Reference Pressure
Values}\label{hmrref-history-match-rock-table-reference-pressure-values}

\begin{longtable}[]{@{}
  >{\raggedright\arraybackslash}p{(\linewidth - 14\tabcolsep) * \real{0.1250}}
  >{\raggedright\arraybackslash}p{(\linewidth - 14\tabcolsep) * \real{0.1250}}
  >{\raggedright\arraybackslash}p{(\linewidth - 14\tabcolsep) * \real{0.1250}}
  >{\raggedright\arraybackslash}p{(\linewidth - 14\tabcolsep) * \real{0.1250}}
  >{\raggedright\arraybackslash}p{(\linewidth - 14\tabcolsep) * \real{0.1250}}
  >{\raggedright\arraybackslash}p{(\linewidth - 14\tabcolsep) * \real{0.1250}}
  >{\raggedright\arraybackslash}p{(\linewidth - 14\tabcolsep) * \real{0.1250}}
  >{\raggedright\arraybackslash}p{(\linewidth - 14\tabcolsep) * \real{0.1250}}@{}}
\toprule\noalign{}
\begin{minipage}[b]{\linewidth}\raggedright
RUNSPEC
\end{minipage} & \begin{minipage}[b]{\linewidth}\raggedright
GRID
\end{minipage} & \begin{minipage}[b]{\linewidth}\raggedright
EDIT
\end{minipage} & \begin{minipage}[b]{\linewidth}\raggedright
PROPS
\end{minipage} & \begin{minipage}[b]{\linewidth}\raggedright
REGIONS
\end{minipage} & \begin{minipage}[b]{\linewidth}\raggedright
SOLUTION
\end{minipage} & \begin{minipage}[b]{\linewidth}\raggedright
SUMMARY
\end{minipage} & \begin{minipage}[b]{\linewidth}\raggedright
SCHEDULE
\end{minipage} \\
\midrule\noalign{}
\endhead
\bottomrule\noalign{}
\endlastfoot
\end{longtable}

\subsubsection{Description}\label{description-463}

The HMRREF keyword defines the history match rock compaction reference
pressure gradient values to be used in conjunction with HMMROCKT,
ROCKTAB and ROCKTABH keywords in the PROPS section, for when the History
Match Gradient option has been activated by the HMDIMS keyword in the
RUNSPEC section. The history match rock compaction data is entered via
the HMMROCKT, ROCKTAB and ROCKTABH keywords in the PROPS section.

See also the HMDIMS keyword in the RUNSPEC section that specifies the
dimensions for the gradient option, including the maximum number of rock
gradient parameters that can be used with the History Match Gradient
option.

This keyword is not supported by OPM Flow but would change the results
if supported so the simulation will be stopped.

\subsection{HWKRO -- End-Point Scaling of Grid Cell Kro(Swl) (High
Salinity and Water
Wet)}\label{hwkro-end-point-scaling-of-grid-cell-kroswl-high-salinity-and-water-wet}

\begin{longtable}[]{@{}
  >{\raggedright\arraybackslash}p{(\linewidth - 14\tabcolsep) * \real{0.1250}}
  >{\raggedright\arraybackslash}p{(\linewidth - 14\tabcolsep) * \real{0.1250}}
  >{\raggedright\arraybackslash}p{(\linewidth - 14\tabcolsep) * \real{0.1250}}
  >{\raggedright\arraybackslash}p{(\linewidth - 14\tabcolsep) * \real{0.1250}}
  >{\raggedright\arraybackslash}p{(\linewidth - 14\tabcolsep) * \real{0.1250}}
  >{\raggedright\arraybackslash}p{(\linewidth - 14\tabcolsep) * \real{0.1250}}
  >{\raggedright\arraybackslash}p{(\linewidth - 14\tabcolsep) * \real{0.1250}}
  >{\raggedright\arraybackslash}p{(\linewidth - 14\tabcolsep) * \real{0.1250}}@{}}
\toprule\noalign{}
\begin{minipage}[b]{\linewidth}\raggedright
RUNSPEC
\end{minipage} & \begin{minipage}[b]{\linewidth}\raggedright
GRID
\end{minipage} & \begin{minipage}[b]{\linewidth}\raggedright
EDIT
\end{minipage} & \begin{minipage}[b]{\linewidth}\raggedright
PROPS
\end{minipage} & \begin{minipage}[b]{\linewidth}\raggedright
REGIONS
\end{minipage} & \begin{minipage}[b]{\linewidth}\raggedright
SOLUTION
\end{minipage} & \begin{minipage}[b]{\linewidth}\raggedright
SUMMARY
\end{minipage} & \begin{minipage}[b]{\linewidth}\raggedright
SCHEDULE
\end{minipage} \\
\midrule\noalign{}
\endhead
\bottomrule\noalign{}
\endlastfoot
\end{longtable}

\subsubsection{Description}\label{description-464}

HWKRO defines the scaling parameter for the oil relative permeability
value at the connate water saturation (SWL), for all the cells in the
model via an array, for when the Low Salt and Surfactant Wettability
options have been selected. The data is used to scale the oil relative
permeability in the high salinity water wet oil relative permeability
saturation tables. The ENDSCALE keyword in the RUNSPEC section should be
activated to enable end-point scaling and the use of this keyword. In
addition the Low Salt option should be activated by the LOWSALT keyword
in the RUNSPEC section and the Surfactant Wettability option activated
by the SURFACT or SURFACTW keywords, which are also in the RUNSPEC
section.

This keyword is not supported by OPM Flow but would change the results
if supported so the simulation will be stopped.

\subsection{HWKRORG -- End-Point Scaling of Grid Cell Kro(Sgcr) (High
Salinity and Water
Wet)}\label{hwkrorg-end-point-scaling-of-grid-cell-krosgcr-high-salinity-and-water-wet}

\begin{longtable}[]{@{}
  >{\raggedright\arraybackslash}p{(\linewidth - 14\tabcolsep) * \real{0.1250}}
  >{\raggedright\arraybackslash}p{(\linewidth - 14\tabcolsep) * \real{0.1250}}
  >{\raggedright\arraybackslash}p{(\linewidth - 14\tabcolsep) * \real{0.1250}}
  >{\raggedright\arraybackslash}p{(\linewidth - 14\tabcolsep) * \real{0.1250}}
  >{\raggedright\arraybackslash}p{(\linewidth - 14\tabcolsep) * \real{0.1250}}
  >{\raggedright\arraybackslash}p{(\linewidth - 14\tabcolsep) * \real{0.1250}}
  >{\raggedright\arraybackslash}p{(\linewidth - 14\tabcolsep) * \real{0.1250}}
  >{\raggedright\arraybackslash}p{(\linewidth - 14\tabcolsep) * \real{0.1250}}@{}}
\toprule\noalign{}
\begin{minipage}[b]{\linewidth}\raggedright
RUNSPEC
\end{minipage} & \begin{minipage}[b]{\linewidth}\raggedright
GRID
\end{minipage} & \begin{minipage}[b]{\linewidth}\raggedright
EDIT
\end{minipage} & \begin{minipage}[b]{\linewidth}\raggedright
PROPS
\end{minipage} & \begin{minipage}[b]{\linewidth}\raggedright
REGIONS
\end{minipage} & \begin{minipage}[b]{\linewidth}\raggedright
SOLUTION
\end{minipage} & \begin{minipage}[b]{\linewidth}\raggedright
SUMMARY
\end{minipage} & \begin{minipage}[b]{\linewidth}\raggedright
SCHEDULE
\end{minipage} \\
\midrule\noalign{}
\endhead
\bottomrule\noalign{}
\endlastfoot
\end{longtable}

\subsubsection{Description}\label{description-465}

HWKRORG defines the scaling parameter for the relative permeability of
oil at the critical gas saturation (SGCR), for all the cells in the
model via an array, for when the Low Salt and Surfactant Wettability
options have been selected. The data is used to scale the oil relative
permeability in the high salinity water wet oil relative permeability
saturation tables. The ENDSCALE keyword in the RUNSPEC section should be
activated to enable end-point scaling and the use of this keyword. In
addition the Low Salt option should be activated by the LOWSALT keyword
in the RUNSPEC section and the Surfactant Wettability option activated
by the SURFACT or SURFACTW keywords, which are also in the RUNSPEC
section.

This keyword is not supported by OPM Flow but would change the results
if supported so the simulation will be stopped.

\subsection{HWKRORW -- End-Point Scaling of Grid Cell Kro(Swcr) (High
Salinity and Water
Wet)}\label{hwkrorw-end-point-scaling-of-grid-cell-kroswcr-high-salinity-and-water-wet}

\begin{longtable}[]{@{}
  >{\raggedright\arraybackslash}p{(\linewidth - 14\tabcolsep) * \real{0.1250}}
  >{\raggedright\arraybackslash}p{(\linewidth - 14\tabcolsep) * \real{0.1250}}
  >{\raggedright\arraybackslash}p{(\linewidth - 14\tabcolsep) * \real{0.1250}}
  >{\raggedright\arraybackslash}p{(\linewidth - 14\tabcolsep) * \real{0.1250}}
  >{\raggedright\arraybackslash}p{(\linewidth - 14\tabcolsep) * \real{0.1250}}
  >{\raggedright\arraybackslash}p{(\linewidth - 14\tabcolsep) * \real{0.1250}}
  >{\raggedright\arraybackslash}p{(\linewidth - 14\tabcolsep) * \real{0.1250}}
  >{\raggedright\arraybackslash}p{(\linewidth - 14\tabcolsep) * \real{0.1250}}@{}}
\toprule\noalign{}
\begin{minipage}[b]{\linewidth}\raggedright
RUNSPEC
\end{minipage} & \begin{minipage}[b]{\linewidth}\raggedright
GRID
\end{minipage} & \begin{minipage}[b]{\linewidth}\raggedright
EDIT
\end{minipage} & \begin{minipage}[b]{\linewidth}\raggedright
PROPS
\end{minipage} & \begin{minipage}[b]{\linewidth}\raggedright
REGIONS
\end{minipage} & \begin{minipage}[b]{\linewidth}\raggedright
SOLUTION
\end{minipage} & \begin{minipage}[b]{\linewidth}\raggedright
SUMMARY
\end{minipage} & \begin{minipage}[b]{\linewidth}\raggedright
SCHEDULE
\end{minipage} \\
\midrule\noalign{}
\endhead
\bottomrule\noalign{}
\endlastfoot
\end{longtable}

\subsubsection{Description}\label{description-466}

HWKRORW defines the scaling parameter for the relative permeability of
oil at the critical water saturation (SWCR), for all the cells in the
model via an array, for when the Low Salt and Surfactant Wettability
options have been selected. The data is used to scale the oil relative
permeability in the high salinity water wet oil relative permeability
saturation tables. The ENDSCALE keyword in the RUNSPEC section should be
activated to enable end-point scaling and the use of this keyword. In
addition the Low Salt option should be activated by the LOWSALT keyword
in the RUNSPEC section and the Surfactant Wettability option activated
by the SURFACT or SURFACTW keywords, which are also in the RUNSPEC
section.

This keyword is not supported by OPM Flow but would change the results
if supported so the simulation will be stopped.

\subsection{HWKRW -- End-Point Scaling of Grid Cell Krw(Sw =1.0) (High
Salinity and Water
Wet)}\label{hwkrw-end-point-scaling-of-grid-cell-krwsw-1.0-high-salinity-and-water-wet}

\begin{longtable}[]{@{}
  >{\raggedright\arraybackslash}p{(\linewidth - 14\tabcolsep) * \real{0.1250}}
  >{\raggedright\arraybackslash}p{(\linewidth - 14\tabcolsep) * \real{0.1250}}
  >{\raggedright\arraybackslash}p{(\linewidth - 14\tabcolsep) * \real{0.1250}}
  >{\raggedright\arraybackslash}p{(\linewidth - 14\tabcolsep) * \real{0.1250}}
  >{\raggedright\arraybackslash}p{(\linewidth - 14\tabcolsep) * \real{0.1250}}
  >{\raggedright\arraybackslash}p{(\linewidth - 14\tabcolsep) * \real{0.1250}}
  >{\raggedright\arraybackslash}p{(\linewidth - 14\tabcolsep) * \real{0.1250}}
  >{\raggedright\arraybackslash}p{(\linewidth - 14\tabcolsep) * \real{0.1250}}@{}}
\toprule\noalign{}
\begin{minipage}[b]{\linewidth}\raggedright
RUNSPEC
\end{minipage} & \begin{minipage}[b]{\linewidth}\raggedright
GRID
\end{minipage} & \begin{minipage}[b]{\linewidth}\raggedright
EDIT
\end{minipage} & \begin{minipage}[b]{\linewidth}\raggedright
PROPS
\end{minipage} & \begin{minipage}[b]{\linewidth}\raggedright
REGIONS
\end{minipage} & \begin{minipage}[b]{\linewidth}\raggedright
SOLUTION
\end{minipage} & \begin{minipage}[b]{\linewidth}\raggedright
SUMMARY
\end{minipage} & \begin{minipage}[b]{\linewidth}\raggedright
SCHEDULE
\end{minipage} \\
\midrule\noalign{}
\endhead
\bottomrule\noalign{}
\endlastfoot
\end{longtable}

\subsubsection{Description}\label{description-467}

HWKRW defines the scaling parameter at the maximum oil relative
permeability value (SWU), that is for Sw = 1.0, for all the cells in the
model via an array, for when the Low Salt and Surfactant Wettability
options have been selected. The data is used to scale the water relative
permeability in the high salinity water wet water relative permeability
saturation tables. The ENDSCALE keyword in the RUNSPEC section should be
activated to enable end-point scaling and the use of this keyword. In
addition the Low Salt option should be activated by the LOWSALT keyword
in the RUNSPEC section and the Surfactant Wettability option activated
by the SURFACT or SURFACTW keywords, which are also in the RUNSPEC
section.

This keyword is not supported by OPM Flow but would change the results
if supported so the simulation will be stopped.

\subsection{HWKRWR -- End-Point Scaling of Grid Cell KRWR(Sw=1.0) (High
Salinity and Water
Wet)}\label{hwkrwr-end-point-scaling-of-grid-cell-krwrsw1.0-high-salinity-and-water-wet}

\begin{longtable}[]{@{}
  >{\raggedright\arraybackslash}p{(\linewidth - 14\tabcolsep) * \real{0.1250}}
  >{\raggedright\arraybackslash}p{(\linewidth - 14\tabcolsep) * \real{0.1250}}
  >{\raggedright\arraybackslash}p{(\linewidth - 14\tabcolsep) * \real{0.1250}}
  >{\raggedright\arraybackslash}p{(\linewidth - 14\tabcolsep) * \real{0.1250}}
  >{\raggedright\arraybackslash}p{(\linewidth - 14\tabcolsep) * \real{0.1250}}
  >{\raggedright\arraybackslash}p{(\linewidth - 14\tabcolsep) * \real{0.1250}}
  >{\raggedright\arraybackslash}p{(\linewidth - 14\tabcolsep) * \real{0.1250}}
  >{\raggedright\arraybackslash}p{(\linewidth - 14\tabcolsep) * \real{0.1250}}@{}}
\toprule\noalign{}
\begin{minipage}[b]{\linewidth}\raggedright
RUNSPEC
\end{minipage} & \begin{minipage}[b]{\linewidth}\raggedright
GRID
\end{minipage} & \begin{minipage}[b]{\linewidth}\raggedright
EDIT
\end{minipage} & \begin{minipage}[b]{\linewidth}\raggedright
PROPS
\end{minipage} & \begin{minipage}[b]{\linewidth}\raggedright
REGIONS
\end{minipage} & \begin{minipage}[b]{\linewidth}\raggedright
SOLUTION
\end{minipage} & \begin{minipage}[b]{\linewidth}\raggedright
SUMMARY
\end{minipage} & \begin{minipage}[b]{\linewidth}\raggedright
SCHEDULE
\end{minipage} \\
\midrule\noalign{}
\endhead
\bottomrule\noalign{}
\endlastfoot
\end{longtable}

\subsubsection{Description}\label{description-468}

HWKRWR defines the scaling parameter at the maximum oil relative
permeability value (SWU), that is for Sw = 1.0, for all the cells in the
model via an array, for when the Low Salt and Surfactant Wettability
options have been selected. The data is used to scale the water relative
permeability in the high salinity water wet water relative permeability
saturation tables. The ENDSCALE keyword in the RUNSPEC section should be
activated to enable end-point scaling and the use of this keyword. In
addition the Low Salt option should be activated by the LOWSALT keyword
in the RUNSPEC section and the Surfactant Wettability option activated
by the SURFACT or SURFACTW keywords, which are also in the RUNSPEC
section.

This keyword is not supported by OPM Flow but would change the results
if supported so the simulation will be stopped.

\subsection{HWPCW -- End-Point Scaling of Grid Cell Water Capillary
Pressure (High Salinity and Water
Wet)}\label{hwpcw-end-point-scaling-of-grid-cell-water-capillary-pressure-high-salinity-and-water-wet}

\begin{longtable}[]{@{}
  >{\raggedright\arraybackslash}p{(\linewidth - 14\tabcolsep) * \real{0.1250}}
  >{\raggedright\arraybackslash}p{(\linewidth - 14\tabcolsep) * \real{0.1250}}
  >{\raggedright\arraybackslash}p{(\linewidth - 14\tabcolsep) * \real{0.1250}}
  >{\raggedright\arraybackslash}p{(\linewidth - 14\tabcolsep) * \real{0.1250}}
  >{\raggedright\arraybackslash}p{(\linewidth - 14\tabcolsep) * \real{0.1250}}
  >{\raggedright\arraybackslash}p{(\linewidth - 14\tabcolsep) * \real{0.1250}}
  >{\raggedright\arraybackslash}p{(\linewidth - 14\tabcolsep) * \real{0.1250}}
  >{\raggedright\arraybackslash}p{(\linewidth - 14\tabcolsep) * \real{0.1250}}@{}}
\toprule\noalign{}
\begin{minipage}[b]{\linewidth}\raggedright
RUNSPEC
\end{minipage} & \begin{minipage}[b]{\linewidth}\raggedright
GRID
\end{minipage} & \begin{minipage}[b]{\linewidth}\raggedright
EDIT
\end{minipage} & \begin{minipage}[b]{\linewidth}\raggedright
PROPS
\end{minipage} & \begin{minipage}[b]{\linewidth}\raggedright
REGIONS
\end{minipage} & \begin{minipage}[b]{\linewidth}\raggedright
SOLUTION
\end{minipage} & \begin{minipage}[b]{\linewidth}\raggedright
SUMMARY
\end{minipage} & \begin{minipage}[b]{\linewidth}\raggedright
SCHEDULE
\end{minipage} \\
\midrule\noalign{}
\endhead
\bottomrule\noalign{}
\endlastfoot
\end{longtable}

\subsubsection{Description}\label{description-469}

HWPCW defines the maximum water-oil pressure values for all the cells in
the model via an array, for when the Low Salt and Surfactant Wettability
options have been selected. The ENDSCALE keyword in the RUNSPEC section
should be activated to enable end-point scaling and the use of this
keyword. In addition the Low Salt option should be activated by the
LOWSALT keyword in the RUNSPEC section and the Surfactant Wettability
option activated by the SURFACT or SURFACTW keywords, which are also in
the RUNSPEC section. The keyword re-scales the oil-water capillary
pressure in the high salinity water wet capillary saturation tables from
a cell's assigned saturation function by the grid block's HWPCW value.

The capillary pressure for a grid block is scaled by:

\begin{equation}\phantomsection\label{eq-8-55}{
{P}_{c} = {P}_{{c}_{\mathit{TABLE}}}\left(\frac{\mathit{HWPCW}}{{P}_{{c}_{\mathit{TABLE}-\mathit{MAX}}}}\right)
}\end{equation}

Where:

Pc = the resulting high salinity water wet water capillary pressure for
a grid cell.

HWPCW = the maximum capillary pressure from the HWPCW array for a given
cell.

\({P}_{{c}_{\mathit{TABLE}}}\) = the capillary pressure in the high
salinity water wet water capillary pressure

table allocated to the grid block.

\({P}_{{c}_{\mathit{TABLE}-\mathit{MAX}}}\) = the maximum capillary
pressure in the high salinity water wet water

capillary pressure table allocated to the grid block (that is at the
connate

water saturation).

This keyword is not supported by OPM Flow but would change the results
if supported so the simulation will be stopped.

\subsection{HWSOGCR -- End-Point Scaling Grid Cell SOGCR (High Salinity
and Water
Wet)}\label{hwsogcr-end-point-scaling-grid-cell-sogcr-high-salinity-and-water-wet}

\begin{longtable}[]{@{}
  >{\raggedright\arraybackslash}p{(\linewidth - 14\tabcolsep) * \real{0.1250}}
  >{\raggedright\arraybackslash}p{(\linewidth - 14\tabcolsep) * \real{0.1250}}
  >{\raggedright\arraybackslash}p{(\linewidth - 14\tabcolsep) * \real{0.1250}}
  >{\raggedright\arraybackslash}p{(\linewidth - 14\tabcolsep) * \real{0.1250}}
  >{\raggedright\arraybackslash}p{(\linewidth - 14\tabcolsep) * \real{0.1250}}
  >{\raggedright\arraybackslash}p{(\linewidth - 14\tabcolsep) * \real{0.1250}}
  >{\raggedright\arraybackslash}p{(\linewidth - 14\tabcolsep) * \real{0.1250}}
  >{\raggedright\arraybackslash}p{(\linewidth - 14\tabcolsep) * \real{0.1250}}@{}}
\toprule\noalign{}
\begin{minipage}[b]{\linewidth}\raggedright
RUNSPEC
\end{minipage} & \begin{minipage}[b]{\linewidth}\raggedright
GRID
\end{minipage} & \begin{minipage}[b]{\linewidth}\raggedright
EDIT
\end{minipage} & \begin{minipage}[b]{\linewidth}\raggedright
PROPS
\end{minipage} & \begin{minipage}[b]{\linewidth}\raggedright
REGIONS
\end{minipage} & \begin{minipage}[b]{\linewidth}\raggedright
SOLUTION
\end{minipage} & \begin{minipage}[b]{\linewidth}\raggedright
SUMMARY
\end{minipage} & \begin{minipage}[b]{\linewidth}\raggedright
SCHEDULE
\end{minipage} \\
\midrule\noalign{}
\endhead
\bottomrule\noalign{}
\endlastfoot
\end{longtable}

\subsubsection{Description}\label{description-470}

HWSOGCR defines the critical oil saturation with respect to gas
(``SOGCR''), for all the cells in the model via an array, for when the
Low Salt and Surfactant Wettability options have been selected. The data
is used to scale the oil saturation in the high salinity water wet
oil-gas relative permeability saturation tables. The ENDSCALE keyword in
the RUNSPEC section should be activated to enable end-point scaling and
the use of this keyword. In addition, the Low Salt option should be
activated by the LOWSALT keyword in the RUNSPEC section and the
Surfactant Wettability option activated by the SURFACT or SURFACTW
keywords, which are also in the RUNSPEC section.

This keyword is not supported by OPM Flow but would change the results
if supported so the simulation will be stopped.

\subsection{HWSOWCR -- End-Point Scaling Grid Cell SOWCR (High Salinity
and Water
Wet)}\label{hwsowcr-end-point-scaling-grid-cell-sowcr-high-salinity-and-water-wet}

\begin{longtable}[]{@{}
  >{\raggedright\arraybackslash}p{(\linewidth - 14\tabcolsep) * \real{0.1250}}
  >{\raggedright\arraybackslash}p{(\linewidth - 14\tabcolsep) * \real{0.1250}}
  >{\raggedright\arraybackslash}p{(\linewidth - 14\tabcolsep) * \real{0.1250}}
  >{\raggedright\arraybackslash}p{(\linewidth - 14\tabcolsep) * \real{0.1250}}
  >{\raggedright\arraybackslash}p{(\linewidth - 14\tabcolsep) * \real{0.1250}}
  >{\raggedright\arraybackslash}p{(\linewidth - 14\tabcolsep) * \real{0.1250}}
  >{\raggedright\arraybackslash}p{(\linewidth - 14\tabcolsep) * \real{0.1250}}
  >{\raggedright\arraybackslash}p{(\linewidth - 14\tabcolsep) * \real{0.1250}}@{}}
\toprule\noalign{}
\begin{minipage}[b]{\linewidth}\raggedright
RUNSPEC
\end{minipage} & \begin{minipage}[b]{\linewidth}\raggedright
GRID
\end{minipage} & \begin{minipage}[b]{\linewidth}\raggedright
EDIT
\end{minipage} & \begin{minipage}[b]{\linewidth}\raggedright
PROPS
\end{minipage} & \begin{minipage}[b]{\linewidth}\raggedright
REGIONS
\end{minipage} & \begin{minipage}[b]{\linewidth}\raggedright
SOLUTION
\end{minipage} & \begin{minipage}[b]{\linewidth}\raggedright
SUMMARY
\end{minipage} & \begin{minipage}[b]{\linewidth}\raggedright
SCHEDULE
\end{minipage} \\
\midrule\noalign{}
\endhead
\bottomrule\noalign{}
\endlastfoot
\end{longtable}

\subsubsection{Description}\label{description-471}

HWSOWCR defines the critical oil saturation with respect to water
(``SOWCR''), for all the cells in the model via an array, for when the
Low Salt and Surfactant Wettability options have been selected. The data
is used to scale the oil saturation in the high salinity water wet
water-oil relative permeability saturation tables. The ENDSCALE keyword
in the RUNSPEC section should be activated to enable end-point scaling
and the use of this keyword. In addition, the Low Salt option should be
activated by the LOWSALT keyword in the RUNSPEC section and the
Surfactant Wettability option activated by the SURFACT or SURFACTW
keywords, which are also in the RUNSPEC section.

This keyword is not supported by OPM Flow but would change the results
if supported so the simulation will be stopped.

\subsection{HWSWCR -- End-Point Scaling Grid Cell SWCR (High Salinity
and Water
Wet)}\label{hwswcr-end-point-scaling-grid-cell-swcr-high-salinity-and-water-wet}

\begin{longtable}[]{@{}
  >{\raggedright\arraybackslash}p{(\linewidth - 14\tabcolsep) * \real{0.1250}}
  >{\raggedright\arraybackslash}p{(\linewidth - 14\tabcolsep) * \real{0.1250}}
  >{\raggedright\arraybackslash}p{(\linewidth - 14\tabcolsep) * \real{0.1250}}
  >{\raggedright\arraybackslash}p{(\linewidth - 14\tabcolsep) * \real{0.1250}}
  >{\raggedright\arraybackslash}p{(\linewidth - 14\tabcolsep) * \real{0.1250}}
  >{\raggedright\arraybackslash}p{(\linewidth - 14\tabcolsep) * \real{0.1250}}
  >{\raggedright\arraybackslash}p{(\linewidth - 14\tabcolsep) * \real{0.1250}}
  >{\raggedright\arraybackslash}p{(\linewidth - 14\tabcolsep) * \real{0.1250}}@{}}
\toprule\noalign{}
\begin{minipage}[b]{\linewidth}\raggedright
RUNSPEC
\end{minipage} & \begin{minipage}[b]{\linewidth}\raggedright
GRID
\end{minipage} & \begin{minipage}[b]{\linewidth}\raggedright
EDIT
\end{minipage} & \begin{minipage}[b]{\linewidth}\raggedright
PROPS
\end{minipage} & \begin{minipage}[b]{\linewidth}\raggedright
REGIONS
\end{minipage} & \begin{minipage}[b]{\linewidth}\raggedright
SOLUTION
\end{minipage} & \begin{minipage}[b]{\linewidth}\raggedright
SUMMARY
\end{minipage} & \begin{minipage}[b]{\linewidth}\raggedright
SCHEDULE
\end{minipage} \\
\midrule\noalign{}
\endhead
\bottomrule\noalign{}
\endlastfoot
\end{longtable}

\subsubsection{Description}\label{description-472}

HWSWCR defines the critical water saturation (``SWCR''), for all the
cells in the model via an array, for when the Low Salt and Surfactant
Wettability options have been selected. The data is used to scale the
water saturation in the high salinity water wet water-oil relative
permeability saturation tables. The ENDSCALE keyword in the RUNSPEC
section should be activated to enable end-point scaling and the use of
this keyword. In addition, the Low Salt option should be activated by
the LOWSALT keyword in the RUNSPEC section and the Surfactant
Wettability option activated by the SURFACT or SURFACTW keywords, which
are also in the RUNSPEC section.

This keyword is not supported by OPM Flow but would change the results
if supported so the simulation will be stopped.

\subsection{HWSWL -- End-Point Scaling Grid Cell SWL (High Salinity and
Water
Wet)}\label{hwswl-end-point-scaling-grid-cell-swl-high-salinity-and-water-wet}

\begin{longtable}[]{@{}
  >{\raggedright\arraybackslash}p{(\linewidth - 14\tabcolsep) * \real{0.1250}}
  >{\raggedright\arraybackslash}p{(\linewidth - 14\tabcolsep) * \real{0.1250}}
  >{\raggedright\arraybackslash}p{(\linewidth - 14\tabcolsep) * \real{0.1250}}
  >{\raggedright\arraybackslash}p{(\linewidth - 14\tabcolsep) * \real{0.1250}}
  >{\raggedright\arraybackslash}p{(\linewidth - 14\tabcolsep) * \real{0.1250}}
  >{\raggedright\arraybackslash}p{(\linewidth - 14\tabcolsep) * \real{0.1250}}
  >{\raggedright\arraybackslash}p{(\linewidth - 14\tabcolsep) * \real{0.1250}}
  >{\raggedright\arraybackslash}p{(\linewidth - 14\tabcolsep) * \real{0.1250}}@{}}
\toprule\noalign{}
\begin{minipage}[b]{\linewidth}\raggedright
RUNSPEC
\end{minipage} & \begin{minipage}[b]{\linewidth}\raggedright
GRID
\end{minipage} & \begin{minipage}[b]{\linewidth}\raggedright
EDIT
\end{minipage} & \begin{minipage}[b]{\linewidth}\raggedright
PROPS
\end{minipage} & \begin{minipage}[b]{\linewidth}\raggedright
REGIONS
\end{minipage} & \begin{minipage}[b]{\linewidth}\raggedright
SOLUTION
\end{minipage} & \begin{minipage}[b]{\linewidth}\raggedright
SUMMARY
\end{minipage} & \begin{minipage}[b]{\linewidth}\raggedright
SCHEDULE
\end{minipage} \\
\midrule\noalign{}
\endhead
\bottomrule\noalign{}
\endlastfoot
\end{longtable}

\subsubsection{Description}\label{description-473}

HWSWL defines the connate water saturation (``SWL''), for all the cells
in the model via an array, for when the Low Salt and Surfactant
Wettability options have been selected. The data is used to scale the
water saturation in the high salinity water wet water-oil relative
permeability saturation tables, as well as the associated capillary
pressure tables. The ENDSCALE keyword in the RUNSPEC section should be
activated to enable end-point scaling and the use of this keyword. In
addition, the Low Salt option should be activated by the LOWSALT keyword
in the RUNSPEC section and the Surfactant Wettability option activated
by the SURFACT or SURFACTW keywords, which are also in the RUNSPEC
section.

This keyword is not supported by OPM Flow but would change the results
if supported so the simulation will be stopped.

\subsection{HWSWLPC -- End-Point Scaling Grid Cell SWLPC (High Salinity
and Water
Wet)}\label{hwswlpc-end-point-scaling-grid-cell-swlpc-high-salinity-and-water-wet}

\begin{longtable}[]{@{}
  >{\raggedright\arraybackslash}p{(\linewidth - 14\tabcolsep) * \real{0.1250}}
  >{\raggedright\arraybackslash}p{(\linewidth - 14\tabcolsep) * \real{0.1250}}
  >{\raggedright\arraybackslash}p{(\linewidth - 14\tabcolsep) * \real{0.1250}}
  >{\raggedright\arraybackslash}p{(\linewidth - 14\tabcolsep) * \real{0.1250}}
  >{\raggedright\arraybackslash}p{(\linewidth - 14\tabcolsep) * \real{0.1250}}
  >{\raggedright\arraybackslash}p{(\linewidth - 14\tabcolsep) * \real{0.1250}}
  >{\raggedright\arraybackslash}p{(\linewidth - 14\tabcolsep) * \real{0.1250}}
  >{\raggedright\arraybackslash}p{(\linewidth - 14\tabcolsep) * \real{0.1250}}@{}}
\toprule\noalign{}
\begin{minipage}[b]{\linewidth}\raggedright
RUNSPEC
\end{minipage} & \begin{minipage}[b]{\linewidth}\raggedright
GRID
\end{minipage} & \begin{minipage}[b]{\linewidth}\raggedright
EDIT
\end{minipage} & \begin{minipage}[b]{\linewidth}\raggedright
PROPS
\end{minipage} & \begin{minipage}[b]{\linewidth}\raggedright
REGIONS
\end{minipage} & \begin{minipage}[b]{\linewidth}\raggedright
SOLUTION
\end{minipage} & \begin{minipage}[b]{\linewidth}\raggedright
SUMMARY
\end{minipage} & \begin{minipage}[b]{\linewidth}\raggedright
SCHEDULE
\end{minipage} \\
\midrule\noalign{}
\endhead
\bottomrule\noalign{}
\endlastfoot
\end{longtable}

\subsubsection{Description}\label{description-474}

HWSWLPC defines the capillary pressure connate water saturation
(``SWLPC''), for all the cells in the model via an array, for when the
Low Salt and Surfactant Wettability options have been selected. The data
is used to scale the water saturation in the high salinity water wet
water-oil capillary pressure tables. The ENDSCALE keyword in the RUNSPEC
section should be activated to enable end-point scaling and the use of
this keyword. In addition, the Low Salt option should be activated by
the LOWSALT keyword in the RUNSPEC section and the Surfactant
Wettability option activated by the SURFACT or SURFACTW keywords, which
are also in the RUNSPEC section.

Note the keyword only applies the scaling to the capillary pressures
tables, unlike the HWSWL keyword that applies the scaling to both the
capillary pressure and relative permeability tables.

This keyword is not supported by OPM Flow but would change the results
if supported so the simulation will be stopped.

\subsection{HWSWU -- End-Point Scaling Grid Cell SWU (High Salinity and
Water
Wet)}\label{hwswu-end-point-scaling-grid-cell-swu-high-salinity-and-water-wet}

\begin{longtable}[]{@{}
  >{\raggedright\arraybackslash}p{(\linewidth - 14\tabcolsep) * \real{0.1250}}
  >{\raggedright\arraybackslash}p{(\linewidth - 14\tabcolsep) * \real{0.1250}}
  >{\raggedright\arraybackslash}p{(\linewidth - 14\tabcolsep) * \real{0.1250}}
  >{\raggedright\arraybackslash}p{(\linewidth - 14\tabcolsep) * \real{0.1250}}
  >{\raggedright\arraybackslash}p{(\linewidth - 14\tabcolsep) * \real{0.1250}}
  >{\raggedright\arraybackslash}p{(\linewidth - 14\tabcolsep) * \real{0.1250}}
  >{\raggedright\arraybackslash}p{(\linewidth - 14\tabcolsep) * \real{0.1250}}
  >{\raggedright\arraybackslash}p{(\linewidth - 14\tabcolsep) * \real{0.1250}}@{}}
\toprule\noalign{}
\begin{minipage}[b]{\linewidth}\raggedright
RUNSPEC
\end{minipage} & \begin{minipage}[b]{\linewidth}\raggedright
GRID
\end{minipage} & \begin{minipage}[b]{\linewidth}\raggedright
EDIT
\end{minipage} & \begin{minipage}[b]{\linewidth}\raggedright
PROPS
\end{minipage} & \begin{minipage}[b]{\linewidth}\raggedright
REGIONS
\end{minipage} & \begin{minipage}[b]{\linewidth}\raggedright
SOLUTION
\end{minipage} & \begin{minipage}[b]{\linewidth}\raggedright
SUMMARY
\end{minipage} & \begin{minipage}[b]{\linewidth}\raggedright
SCHEDULE
\end{minipage} \\
\midrule\noalign{}
\endhead
\bottomrule\noalign{}
\endlastfoot
\end{longtable}

\subsubsection{Description}\label{description-475}

HWSWU defines the maximum water saturation(``SWU''), for all the cells
in the model via an array, for when the Low Salt and Surfactant
Wettability options have been selected. The data is used to scale the
water saturation in the high salinity water wet water-oil relative
permeability saturation tables, as well as the associated capillary
pressure tables. The ENDSCALE keyword in the RUNSPEC section should be
activated to enable end-point scaling and the use of this keyword. In
addition, the Low Salt option should be activated by the LOWSALT keyword
in the RUNSPEC section and the Surfactant Wettability option activated
by the SURFACT or SURFACTW keywords, which are also in the RUNSPEC
section.

This keyword is not supported by OPM Flow but would change the results
if supported so the simulation will be stopped.

\subsection{HYDRHEAD -- Define Hydraulic Head Output Reference
Data}\label{hydrhead-define-hydraulic-head-output-reference-data}

\begin{longtable}[]{@{}
  >{\raggedright\arraybackslash}p{(\linewidth - 14\tabcolsep) * \real{0.1250}}
  >{\raggedright\arraybackslash}p{(\linewidth - 14\tabcolsep) * \real{0.1250}}
  >{\raggedright\arraybackslash}p{(\linewidth - 14\tabcolsep) * \real{0.1250}}
  >{\raggedright\arraybackslash}p{(\linewidth - 14\tabcolsep) * \real{0.1250}}
  >{\raggedright\arraybackslash}p{(\linewidth - 14\tabcolsep) * \real{0.1250}}
  >{\raggedright\arraybackslash}p{(\linewidth - 14\tabcolsep) * \real{0.1250}}
  >{\raggedright\arraybackslash}p{(\linewidth - 14\tabcolsep) * \real{0.1250}}
  >{\raggedright\arraybackslash}p{(\linewidth - 14\tabcolsep) * \real{0.1250}}@{}}
\toprule\noalign{}
\begin{minipage}[b]{\linewidth}\raggedright
RUNSPEC
\end{minipage} & \begin{minipage}[b]{\linewidth}\raggedright
GRID
\end{minipage} & \begin{minipage}[b]{\linewidth}\raggedright
EDIT
\end{minipage} & \begin{minipage}[b]{\linewidth}\raggedright
PROPS
\end{minipage} & \begin{minipage}[b]{\linewidth}\raggedright
REGIONS
\end{minipage} & \begin{minipage}[b]{\linewidth}\raggedright
SOLUTION
\end{minipage} & \begin{minipage}[b]{\linewidth}\raggedright
SUMMARY
\end{minipage} & \begin{minipage}[b]{\linewidth}\raggedright
SCHEDULE
\end{minipage} \\
\midrule\noalign{}
\endhead
\bottomrule\noalign{}
\endlastfoot
\end{longtable}

\subsubsection{Description}\label{description-476}

The HYDRHEAD keyword defines the hydraulic head reference data for when
the hydraulic head information is requested to be written out via one on
the SUMMARY keywords (BHD, BHDF, etc.) in the SUMMARY section, or to the
RESTART file via the HYDH or HYDHFW variables on the RESTART keyword

This keyword is not supported by OPM Flow but would change the results
if supported so the simulation will be stopped.

\subsection{HYMOBGDR -- Activate Carlson and Killough Alternative
Drainage Hysteresis
Option}\label{hymobgdr-activate-carlson-and-killough-alternative-drainage-hysteresis-option}

\begin{longtable}[]{@{}
  >{\raggedright\arraybackslash}p{(\linewidth - 14\tabcolsep) * \real{0.1250}}
  >{\raggedright\arraybackslash}p{(\linewidth - 14\tabcolsep) * \real{0.1250}}
  >{\raggedright\arraybackslash}p{(\linewidth - 14\tabcolsep) * \real{0.1250}}
  >{\raggedright\arraybackslash}p{(\linewidth - 14\tabcolsep) * \real{0.1250}}
  >{\raggedright\arraybackslash}p{(\linewidth - 14\tabcolsep) * \real{0.1250}}
  >{\raggedright\arraybackslash}p{(\linewidth - 14\tabcolsep) * \real{0.1250}}
  >{\raggedright\arraybackslash}p{(\linewidth - 14\tabcolsep) * \real{0.1250}}
  >{\raggedright\arraybackslash}p{(\linewidth - 14\tabcolsep) * \real{0.1250}}@{}}
\toprule\noalign{}
\begin{minipage}[b]{\linewidth}\raggedright
RUNSPEC
\end{minipage} & \begin{minipage}[b]{\linewidth}\raggedright
GRID
\end{minipage} & \begin{minipage}[b]{\linewidth}\raggedright
EDIT
\end{minipage} & \begin{minipage}[b]{\linewidth}\raggedright
PROPS
\end{minipage} & \begin{minipage}[b]{\linewidth}\raggedright
REGIONS
\end{minipage} & \begin{minipage}[b]{\linewidth}\raggedright
SOLUTION
\end{minipage} & \begin{minipage}[b]{\linewidth}\raggedright
SUMMARY
\end{minipage} & \begin{minipage}[b]{\linewidth}\raggedright
SCHEDULE
\end{minipage} \\
\midrule\noalign{}
\endhead
\bottomrule\noalign{}
\endlastfoot
\end{longtable}

\subsubsection{Description}\label{description-477}

This keyword, HYMOBGDR, activates the Carlson and Killough alternative
secondary drainage hysteresis option for when the hysteresis option has
been activated by the HYSTER variable on the SATOPTS keyword in the
RUNSPEC section, and either the Carlson {[}Carlson, F. M. ``Simulation
of Relative Permeability Hysteresis to the Non-Wetting Phase,'' paper
SPE 10157, presented at the SPE Annual Technical Conference \&
Exhibition, San Antonio, Texas, USA (October 5-7, 1981).{]} or Killough
{[}Killough, J. E. ``Reservoir Simulation with History-dependent
Saturation Functions,'' paper SPE 5106, Society of Petroleum Engineers
Journal (1976) 16, No.~1, 37-48.{]} models have been selected via the
EHYSTR keyword in the PROPS section. Due to numerical accuracy, the gas
saturation may fall below the critical gas saturation (SGCR), that is
the largest gas saturation for which the gas relative permeability is
zero, and gas would therefore be immobile until the gas saturation
increases above SGCR. This option overcomes this situation by letting
the gas become mobile once it starts increasing, effectively setting the
SGCR to the current gas saturation.

There is no data required for this keyword and there is no terminating
``/'' for this keyword.

This keyword is not supported by OPM Flow but would change the results
if supported so the simulation will be stopped.

\subsubsection{Example}\label{example-283}

\begin{verbatim}
--
--       ACTIVATE CARLSON AND KILLOUGH ALTERNATIVE DRAINAGE HYSTERESIS OPTION
--
HYMOBGDR
\end{verbatim}

\subsection{HYSTCHCK -- Activate Hysteresis Imbibition and Drainage
End-Point
Validation}\label{hystchck-activate-hysteresis-imbibition-and-drainage-end-point-validation}

\begin{longtable}[]{@{}
  >{\raggedright\arraybackslash}p{(\linewidth - 14\tabcolsep) * \real{0.1250}}
  >{\raggedright\arraybackslash}p{(\linewidth - 14\tabcolsep) * \real{0.1250}}
  >{\raggedright\arraybackslash}p{(\linewidth - 14\tabcolsep) * \real{0.1250}}
  >{\raggedright\arraybackslash}p{(\linewidth - 14\tabcolsep) * \real{0.1250}}
  >{\raggedright\arraybackslash}p{(\linewidth - 14\tabcolsep) * \real{0.1250}}
  >{\raggedright\arraybackslash}p{(\linewidth - 14\tabcolsep) * \real{0.1250}}
  >{\raggedright\arraybackslash}p{(\linewidth - 14\tabcolsep) * \real{0.1250}}
  >{\raggedright\arraybackslash}p{(\linewidth - 14\tabcolsep) * \real{0.1250}}@{}}
\toprule\noalign{}
\begin{minipage}[b]{\linewidth}\raggedright
RUNSPEC
\end{minipage} & \begin{minipage}[b]{\linewidth}\raggedright
GRID
\end{minipage} & \begin{minipage}[b]{\linewidth}\raggedright
EDIT
\end{minipage} & \begin{minipage}[b]{\linewidth}\raggedright
PROPS
\end{minipage} & \begin{minipage}[b]{\linewidth}\raggedright
REGIONS
\end{minipage} & \begin{minipage}[b]{\linewidth}\raggedright
SOLUTION
\end{minipage} & \begin{minipage}[b]{\linewidth}\raggedright
SUMMARY
\end{minipage} & \begin{minipage}[b]{\linewidth}\raggedright
SCHEDULE
\end{minipage} \\
\midrule\noalign{}
\endhead
\bottomrule\noalign{}
\endlastfoot
\end{longtable}

\subsubsection{Description}\label{description-478}

The HYSTCHCK keyword activate the hysteresis imbibition and drainage
end-point check to validate that the two sets of end-points are
consistent, for when the Hysteresis option has been activated by the
HYSTER variable on the SATOPTS keyword in the RUNSPEC section, and the
ENDSCALE keyword in the RUNSPEC section has been activated to enable
end-point scaling.

This keyword is not supported by OPM Flow but would change the results
if supported so the simulation will be stopped.

\subsection{IKRG -- End-Point Scaling of Grid Cell Krg(Sgu)
(Imbibition)}\label{ikrg-end-point-scaling-of-grid-cell-krgsgu-imbibition}

\begin{longtable}[]{@{}
  >{\raggedright\arraybackslash}p{(\linewidth - 14\tabcolsep) * \real{0.1250}}
  >{\raggedright\arraybackslash}p{(\linewidth - 14\tabcolsep) * \real{0.1250}}
  >{\raggedright\arraybackslash}p{(\linewidth - 14\tabcolsep) * \real{0.1250}}
  >{\raggedright\arraybackslash}p{(\linewidth - 14\tabcolsep) * \real{0.1250}}
  >{\raggedright\arraybackslash}p{(\linewidth - 14\tabcolsep) * \real{0.1250}}
  >{\raggedright\arraybackslash}p{(\linewidth - 14\tabcolsep) * \real{0.1250}}
  >{\raggedright\arraybackslash}p{(\linewidth - 14\tabcolsep) * \real{0.1250}}
  >{\raggedright\arraybackslash}p{(\linewidth - 14\tabcolsep) * \real{0.1250}}@{}}
\toprule\noalign{}
\begin{minipage}[b]{\linewidth}\raggedright
RUNSPEC
\end{minipage} & \begin{minipage}[b]{\linewidth}\raggedright
GRID
\end{minipage} & \begin{minipage}[b]{\linewidth}\raggedright
EDIT
\end{minipage} & \begin{minipage}[b]{\linewidth}\raggedright
PROPS
\end{minipage} & \begin{minipage}[b]{\linewidth}\raggedright
REGIONS
\end{minipage} & \begin{minipage}[b]{\linewidth}\raggedright
SOLUTION
\end{minipage} & \begin{minipage}[b]{\linewidth}\raggedright
SUMMARY
\end{minipage} & \begin{minipage}[b]{\linewidth}\raggedright
SCHEDULE
\end{minipage} \\
\midrule\noalign{}
\endhead
\bottomrule\noalign{}
\endlastfoot
\end{longtable}

\subsubsection{Description}\label{description-479}

IKRG defines the imbibition scaling parameter at the maximum gas
relative permeability value (ISGU), normally ISGU is equal to 1.0 - Swc,
for all the cells in the model via an array. The ENDSCALE keyword in the
RUNSPEC section should be activated to enable end-point scaling and the
use of this keyword. In addition, the HYSTER option on the SATOPTS
keyword in the RUNSPEC section has to be activated to invoke the
Hysteresis option. The SCALCERS keyword in the PROPS section defines the
options used in the re-scaling process, the options are two point
scaling and three point scaling.

\begin{longtable}[]{@{}
  >{\raggedright\arraybackslash}p{(\linewidth - 6\tabcolsep) * \real{0.2500}}
  >{\raggedright\arraybackslash}p{(\linewidth - 6\tabcolsep) * \real{0.2500}}
  >{\raggedright\arraybackslash}p{(\linewidth - 6\tabcolsep) * \real{0.2500}}
  >{\raggedright\arraybackslash}p{(\linewidth - 6\tabcolsep) * \real{0.2500}}@{}}
\toprule\noalign{}
\begin{minipage}[b]{\linewidth}\raggedright
No.
\end{minipage} & \begin{minipage}[b]{\linewidth}\raggedright
Name
\end{minipage} & \begin{minipage}[b]{\linewidth}\raggedright
Description
\end{minipage} & \begin{minipage}[b]{\linewidth}\raggedright
Default
\end{minipage} \\
\midrule\noalign{}
\endhead
\bottomrule\noalign{}
\endlastfoot
Field & Metric & Laboratory & \\
1 & IKRG & IKRG is an array of positive real numbers which are greater
than zero and less than or equal to 1.0, that are the assigned scaling
IKRG imbibition values for each cell in the model. Repeat counts may be
used, for example 50*0.400. dimensionless & Taken from cell allocated
relative permeability table. \\
dimensionless & dimensionless & dimensionless & \\
Notes: & & & \\
\end{longtable}

\emph{Table 8.47: IKRG Keyword Description}

\begin{verbatim}
\end{verbatim}

For the two point scaling option and for the IKRGR gas relative
permeability array NOT present in the input deck the krg value for a
grid block is scaled by:

\begin{equation}\phantomsection\label{eq-8-56}{
{k}_{\mathit{rg}} = {k}_{{\mathit{rg}}_{ \mathit{TABLE}}}\left(\frac{\mathit{IKRG}}{{k}_{{\mathit{rg}}_{ \mathit{TABLE}-\mathit{MAX}}}}\right)
}\end{equation}

Where:

krg = the resulting krg value for a grid cell.

IKRG = the scaling gas relative permeability value from the IKRG array
for a given

cell.

\({k}_{{\mathit{rg}}_{ \mathit{TABLE}}}\) = the gas relative
permeability from a grid block's gas-oil table at the grid

blocks gas saturation.

\({k}_{{\mathit{rg}}_{ \mathit{TABLE}-\mathit{MAX}}}\) = the maximum gas
relative permeability from a grid block's gas-oil table, that

is at the connate water saturation (Swc).

If the IKRGR keyword is present in the input deck then the scaling
matches the imbibition relative permeability at the critical saturation
of the displacing phase.

If three point scaling option has been selected via the SCALECRS keyword
in the PROPS section the critical displacing phase is defined as:

\begin{longtable}[]{@{}lll@{}}
\toprule\noalign{}
No & Phases Present & Critical Saturation \\
\midrule\noalign{}
\endhead
\bottomrule\noalign{}
\endlastfoot
1 & Gas-Oil & S critical = 1.0 -- ISOGCR - ISWL \\
2 & Gas-Oil-Water & S critical = 1.0 -- ISOGCR - ISWL \\
3 & Gas-Water & S critical = 1.0 -- ISWCR \\
\end{longtable}

\emph{Table 8.48: Critical Displacement Relationships}

End-point scaling allows the entered relative permeability functions to
be re-scaled based on the saturation end-points allocated to each cell
by the ISWL, ISWCR, ISWU, ISGL, ISGCR, ISGU, ISOWCR, and ISOGCR
saturation grid arrays for the saturation end-points, In addition
end-point scaling may be directional dependent in which case the
directional dependent versions of the aforementioned arrays should be
used, that is ISWUX, ISWUY and ISWUZ instead of ISWU, There is also the
facility to make the directional end-point scaling reversible or
non-reversible and if the non-reversible option is selected the
non-reversible versions of the aforementioned arrays should be used,
that is ISWUX, ISWUX-, ISWUY, ISWUY-, ISWUZ and ISWUZ-, instead of the
ISWU keyword.

End-point scaling also allows the entered relative permeability
functions to be scale on the relative permeability values using the
IKRG, IKRGR, IKRO, IKRORG, IKRORW, IKRW and IKRWR relative permeability
grid cell arrays for the relative permeability end-point data. In
addition end-point scaling may be directional dependent in which case
the directional dependent versions of the aforementioned arrays should
be used, that is IKRGX, IKRGY and IKRGZ instead of IKRG, There is also
the facility to make the directional end-point scaling reversible or
non-reversible and if the non-reversible option is selected the
non-reversible versions of the aforementioned arrays should be used,
that is IKRGX, IKRGX-, IKRGY, IKRGY-,I KRGZ and IKRGZ-, instead of the
IKRG keyword.

\subsubsection{Example}\label{example-284}

The example below defines an input box for the whole grid and for layers
one to three, for layer one IKRG is set equal to 0.550, for layer two
IKRG equals 0.575, and for layer three IKRG equals 0.600.

\begin{verbatim}
--
--       DEFINE INPUT BOX FOR EDITING INPUT ARRAYS (NX=100, NY=100)
--
--       ---------- BOX ---------
--       I1  I2   J1  J2   K1  K2
BOX
         1*  1*   1*  1*   1   3                           / DEFINE BOX AREA
--
--       SET IKRG VALUES FOR THREE LAYERS IN THE MODEL
--
IKRG
         1000*0.550  1000*0.575  1000*0.600                /
--
--       DEFINE END OF INPUT BOX EDITING OF INPUT ARRAYS
--
ENDBOX
\end{verbatim}

\subsection{IKRGR -- End-Point Scaling of Grid Cell Krgr(1-Sogcr)
(Imbibition)}\label{ikrgr-end-point-scaling-of-grid-cell-krgr1-sogcr-imbibition}

\begin{longtable}[]{@{}
  >{\raggedright\arraybackslash}p{(\linewidth - 14\tabcolsep) * \real{0.1250}}
  >{\raggedright\arraybackslash}p{(\linewidth - 14\tabcolsep) * \real{0.1250}}
  >{\raggedright\arraybackslash}p{(\linewidth - 14\tabcolsep) * \real{0.1250}}
  >{\raggedright\arraybackslash}p{(\linewidth - 14\tabcolsep) * \real{0.1250}}
  >{\raggedright\arraybackslash}p{(\linewidth - 14\tabcolsep) * \real{0.1250}}
  >{\raggedright\arraybackslash}p{(\linewidth - 14\tabcolsep) * \real{0.1250}}
  >{\raggedright\arraybackslash}p{(\linewidth - 14\tabcolsep) * \real{0.1250}}
  >{\raggedright\arraybackslash}p{(\linewidth - 14\tabcolsep) * \real{0.1250}}@{}}
\toprule\noalign{}
\begin{minipage}[b]{\linewidth}\raggedright
RUNSPEC
\end{minipage} & \begin{minipage}[b]{\linewidth}\raggedright
GRID
\end{minipage} & \begin{minipage}[b]{\linewidth}\raggedright
EDIT
\end{minipage} & \begin{minipage}[b]{\linewidth}\raggedright
PROPS
\end{minipage} & \begin{minipage}[b]{\linewidth}\raggedright
REGIONS
\end{minipage} & \begin{minipage}[b]{\linewidth}\raggedright
SOLUTION
\end{minipage} & \begin{minipage}[b]{\linewidth}\raggedright
SUMMARY
\end{minipage} & \begin{minipage}[b]{\linewidth}\raggedright
SCHEDULE
\end{minipage} \\
\midrule\noalign{}
\endhead
\bottomrule\noalign{}
\endlastfoot
\end{longtable}

\subsubsection{Description}\label{description-480}

IKRGR defines the imbibition scaling parameter at the relative
permeability of gas at residual oil saturation (1 -- ISOGCR), or
critical water saturation in a gas-water run (Swc), for all the cells in
the model via an array. The ENDSCALE keyword in the RUNSPEC section
should be activated to enable end-point scaling and the use of this
keyword. In addition, the HYSTER option on the SATOPTS keyword in the
RUNSPEC section has to be activated to invoke the Hysteresis option. The
SCALCERS keyword in the PROPS section defines the options used in the
re-scaling process, the options are two point scaling and three point
scaling.

\begin{longtable}[]{@{}
  >{\raggedright\arraybackslash}p{(\linewidth - 6\tabcolsep) * \real{0.2500}}
  >{\raggedright\arraybackslash}p{(\linewidth - 6\tabcolsep) * \real{0.2500}}
  >{\raggedright\arraybackslash}p{(\linewidth - 6\tabcolsep) * \real{0.2500}}
  >{\raggedright\arraybackslash}p{(\linewidth - 6\tabcolsep) * \real{0.2500}}@{}}
\toprule\noalign{}
\begin{minipage}[b]{\linewidth}\raggedright
No.
\end{minipage} & \begin{minipage}[b]{\linewidth}\raggedright
Name
\end{minipage} & \begin{minipage}[b]{\linewidth}\raggedright
Description
\end{minipage} & \begin{minipage}[b]{\linewidth}\raggedright
Default
\end{minipage} \\
\midrule\noalign{}
\endhead
\bottomrule\noalign{}
\endlastfoot
Field & Metric & Laboratory & \\
1 & IKRGR & IKRGR is an array of positive real numbers which are greater
than zero and less than or equal to 1.0, that are the assigned
imbibition scaling IKRGR values for each cell in the model. In addition,
for a given grid block IKGRGT should be less than IKRG. Repeat counts
may be used, for example 50*0.400. & Taken from cell allocated relative
permeability table. \\
dimensionless & dimensionless & dimensionless & \\
Notes: & & & \\
\end{longtable}

\emph{Table 8.49: IKRGR Keyword Description}

\begin{verbatim}
\end{verbatim}

When the IKRGR keyword is present in the input deck then the scaling
matches the imbibition relative permeability at the critical saturation
of the displacing phase (oil or water).

If three point scaling option has been selected via the SCALECRS keyword
in the PROPS section the critical displacing phase is defined as:

\begin{longtable}[]{@{}lll@{}}
\toprule\noalign{}
No & Phases Present & Critical Saturation \\
\midrule\noalign{}
\endhead
\bottomrule\noalign{}
\endlastfoot
1 & Gas-Oil & S critical = 1.0 -- ISOGCR - ISWL \\
2 & Gas-Oil-Water & S critical = 1.0 -- ISOGCR - ISWL \\
3 & Gas-Water & S critical = 1.0 -- ISWCR \\
\end{longtable}

\emph{Table 8.50: Critical Displacement Relationships}

End-point scaling allows the entered relative permeability functions to
be re-scaled based on the saturation end-points allocated to each cell
by the ISWL, ISWCR, ISWU, ISGL, ISGCR, ISGU, ISOWCR, and ISOGCR
saturation grid arrays for the saturation end-points, In addition
end-point scaling may be directional dependent in which case the
directional dependent versions of the aforementioned arrays should be
used, that is ISWUX, ISWUY and ISWUZ instead of ISWU, There is also the
facility to make the directional end-point scaling reversible or
non-reversible and if the non-reversible option is selected the
non-reversible versions of the aforementioned arrays should be used,
that is ISWUX, ISWUX-, ISWUY, ISWUY-, ISWUZ and ISWUZ-, instead of the
ISWU keyword.

End-point scaling also allows the entered relative permeability
functions to be scale on the relative permeability values using the
IKRG, IKRGR, IKRO, IKRORG, IKRORW, IKRW and IKRWR relative permeability
grid cell arrays for the relative permeability end-point data. In
addition end-point scaling may be directional dependent in which case
the directional dependent versions of the aforementioned arrays should
be used, that is IKRGX, IKRGY and IKRGZ instead of IKRG, There is also
the facility to make the directional end-point scaling reversible or
non-reversible and if the non-reversible option is selected the
non-reversible versions of the aforementioned arrays should be used,
that is IKRGX, IKRGX-, IKRGY, IKRGY-,I KRGZ and IKRGZ-, instead of the
IKRG keyword.

\subsubsection{Examples}\label{examples-25}

The first example defines an input box for the whole grid and for layers
one to three, for layer one IKRRG is set equal to 0.500, for layer two
IKRGR equals 0.570, and for layer three IKRGR equals 0.580.

\begin{verbatim}
--
--       DEFINE INPUT BOX FOR EDITING INPUT ARRAYS (NX=100, NY=100)
--
--       ---------- BOX ---------
--       I1  I2   J1  J2   K1  K2
BOX
         1*  1*   1*  1*   1   3                           / DEFINE BOX AREA
--
--       SET IKRGR VALUES FOR THREE LAYERS IN THE MODEL
--
IKRGR
         10000*0.500  10000*0.570  10000*0.580                /
--
--       DEFINE END OF INPUT BOX EDITING OF INPUT ARRAYS
--
ENDBOX
\end{verbatim}

The next example does exactly the same thing using the EQUALS keyword
instead.

\begin{verbatim}
--
--       ARRAY       CONSTANT     ---------- BOX ---------
--                                I1  I2   J1  J2   K1  K2
EQUALS
         IKRGR’      0.5000       1*  1*   1*  1*    1   1  / IKRGR FOR LAYER 1
         IKRGR’      0.5700       1*  1*   1*  1*    2   2  / IKRGR FOR LAYER 2
         IKRGR’      0.5800       1*  1*   1*  1*    3   3  / IKRGR FOR LAYER 3
/
\end{verbatim}

\subsection{IKRO -- End-Point Scaling of Grid Cell Kro(Swl)
(Imbibition)}\label{ikro-end-point-scaling-of-grid-cell-kroswl-imbibition}

\begin{longtable}[]{@{}
  >{\raggedright\arraybackslash}p{(\linewidth - 14\tabcolsep) * \real{0.1250}}
  >{\raggedright\arraybackslash}p{(\linewidth - 14\tabcolsep) * \real{0.1250}}
  >{\raggedright\arraybackslash}p{(\linewidth - 14\tabcolsep) * \real{0.1250}}
  >{\raggedright\arraybackslash}p{(\linewidth - 14\tabcolsep) * \real{0.1250}}
  >{\raggedright\arraybackslash}p{(\linewidth - 14\tabcolsep) * \real{0.1250}}
  >{\raggedright\arraybackslash}p{(\linewidth - 14\tabcolsep) * \real{0.1250}}
  >{\raggedright\arraybackslash}p{(\linewidth - 14\tabcolsep) * \real{0.1250}}
  >{\raggedright\arraybackslash}p{(\linewidth - 14\tabcolsep) * \real{0.1250}}@{}}
\toprule\noalign{}
\begin{minipage}[b]{\linewidth}\raggedright
RUNSPEC
\end{minipage} & \begin{minipage}[b]{\linewidth}\raggedright
GRID
\end{minipage} & \begin{minipage}[b]{\linewidth}\raggedright
EDIT
\end{minipage} & \begin{minipage}[b]{\linewidth}\raggedright
PROPS
\end{minipage} & \begin{minipage}[b]{\linewidth}\raggedright
REGIONS
\end{minipage} & \begin{minipage}[b]{\linewidth}\raggedright
SOLUTION
\end{minipage} & \begin{minipage}[b]{\linewidth}\raggedright
SUMMARY
\end{minipage} & \begin{minipage}[b]{\linewidth}\raggedright
SCHEDULE
\end{minipage} \\
\midrule\noalign{}
\endhead
\bottomrule\noalign{}
\endlastfoot
\end{longtable}

\subsubsection{Description}\label{description-481}

IKRO defines the scaling parameter for the imbibition oil relative
permeability value at the connate water saturation (ISWL), for all the
cells in the model via an array. The ENDSCALE keyword in the RUNSPEC
section should be activated to enable end-point scaling and the use of
this keyword. In addition, the HYSTER option on the SATOPTS keyword in
the RUNSPEC section has to be activated to invoke the Hysteresis option.
The SCALCERS keyword in the PROPS section defines the options used in
the re-scaling process, the options are two point scaling and three
point scaling.

\begin{longtable}[]{@{}
  >{\raggedright\arraybackslash}p{(\linewidth - 6\tabcolsep) * \real{0.2500}}
  >{\raggedright\arraybackslash}p{(\linewidth - 6\tabcolsep) * \real{0.2500}}
  >{\raggedright\arraybackslash}p{(\linewidth - 6\tabcolsep) * \real{0.2500}}
  >{\raggedright\arraybackslash}p{(\linewidth - 6\tabcolsep) * \real{0.2500}}@{}}
\toprule\noalign{}
\begin{minipage}[b]{\linewidth}\raggedright
No.
\end{minipage} & \begin{minipage}[b]{\linewidth}\raggedright
Name
\end{minipage} & \begin{minipage}[b]{\linewidth}\raggedright
Description
\end{minipage} & \begin{minipage}[b]{\linewidth}\raggedright
Default
\end{minipage} \\
\midrule\noalign{}
\endhead
\bottomrule\noalign{}
\endlastfoot
Field & Metric & Laboratory & \\
1 & IKRO & IKRO is an array of positive real numbers which are greater
than zero and less than or equal to 1.0, that are the assigned
imbibition scaling IKRO values for each cell in the model. Repeat counts
may be used, for example 50*0.500. & Taken from cell allocated relative
permeability table. \\
dimensionless & dimensionless & dimensionless & \\
Notes: & & & \\
\end{longtable}

\emph{Table 8.51: IKRO Keyword Description}

For the two point scaling option and for the IKRORW or IKRORG oil
imbibition relative permeability arrays NOT being present in the input
deck the kro value for a grid block is scaled by:

\begin{equation}\phantomsection\label{eq-8-57}{
{k}_{\mathit{ro}} = {k}_{{\mathit{ro}}_{ \mathit{TABLE}}}\left(\frac{\mathit{IKRO}}{{k}_{{\mathit{ro}}_{ \mathit{TABLE}-\mathit{MAX}}}}\right)
}\end{equation}

Where:

\(\mathit{kro}\) = the resulting kro value for a grid cell.

\(\mathit{IKRO}\) = the scaling oil relative permeability value from the
IKRO array for a given

cell.

\({k}_{{\mathit{ro}}_{ \mathit{TABLE}}}\) = the oil relative
permeability from a grid block's oil relative permeability

table at the grid blocks oil saturation.

\({k}_{{\mathit{ro}}_{ \mathit{TABLE}-\mathit{MAX}}}\) = the maximum oil
relative permeability from a grid block's oil relative table,

that is at the critical water saturation (Swcr).

If the IKRORW or IKRORG keywords are present in the input deck then the
scaling matches the relative permeability at the critical saturation of
the displacing phase.

If three point scaling option has been selected via the SCALECRS keyword
in the PROPS section the critical displacing phase is defined as:

\begin{longtable}[]{@{}lll@{}}
\toprule\noalign{}
No & Keywords Present & Critical Saturation \\
\midrule\noalign{}
\endhead
\bottomrule\noalign{}
\endlastfoot
1 & KRORW & S critical = 1.0 -- SWCR - SGL \\
2 & KRORG & S critical = 1.0 -- SGCR - SWL \\
\end{longtable}

\emph{Table 8.52: Critical Displacement Relationships}

End-point scaling allows the entered relative permeability functions to
be re-scaled based on the saturation end-points allocated to each cell
by the ISWL, ISWCR, ISWU, ISGL, ISGCR, ISGU, ISOWCR, and ISOGCR
saturation grid arrays for the saturation end-points, In addition
end-point scaling may be directional dependent in which case the
directional dependent versions of the aforementioned arrays should be
used, that is ISWUX, ISWUY and ISWUZ instead of ISWU, There is also the
facility to make the directional end-point scaling reversible or
non-reversible and if the non-reversible option is selected the
non-reversible versions of the aforementioned arrays should be used,
that is ISWUX, ISWUX-, ISWUY, ISWUY-, ISWUZ and ISWUZ-, instead of the
ISWU keyword.

End-point scaling also allows the entered relative permeability
functions to be scale on the relative permeability values using the
IKRG, IKRGR, IKRO, IKRORG, IKRORW, IKRW and IKRWR relative permeability
grid cell arrays for the relative permeability end-point data. In
addition end-point scaling may be directional dependent in which case
the directional dependent versions of the aforementioned arrays should
be used, that is IKRGX, IKRGY and IKRGZ instead of IKRG, There is also
the facility to make the directional end-point scaling reversible or
non-reversible and if the non-reversible option is selected the
non-reversible versions of the aforementioned arrays should be used,
that is IKRGX, IKRGX-, IKRGY, IKRGY-,I KRGZ and IKRGZ-, instead of the
IKRG keyword.

\subsubsection{Example}\label{example-285}

The example below defines an input box for the whole grid and for layers
one to three, for layer one IKRO is set equal to 0.850, for layer two
IKRO equals 0.875, and for layer three IKRO equals 0.900.

\begin{verbatim}
--
--       ARRAY       CONSTANT     ---------- BOX ---------
--                                I1  I2   J1  J2   K1  K2
EQUALS
         IKRO        0.8500       1*  1*   1*  1*    1   1  / IKRO FOR LAYER 1
         IKRO        0.8750       1*  1*   1*  1*    2   2  / IKRO FOR LAYER 2
         IKRO        0.9000       1*  1*   1*  1*    3   3  / IKRO FOR LAYER 3
/
\end{verbatim}

\subsection{IKRORG -- End-Point Scaling of Grid Cell Kro(Sgcr)
(Imbibition)}\label{ikrorg-end-point-scaling-of-grid-cell-krosgcr-imbibition}

\begin{longtable}[]{@{}
  >{\raggedright\arraybackslash}p{(\linewidth - 14\tabcolsep) * \real{0.1250}}
  >{\raggedright\arraybackslash}p{(\linewidth - 14\tabcolsep) * \real{0.1250}}
  >{\raggedright\arraybackslash}p{(\linewidth - 14\tabcolsep) * \real{0.1250}}
  >{\raggedright\arraybackslash}p{(\linewidth - 14\tabcolsep) * \real{0.1250}}
  >{\raggedright\arraybackslash}p{(\linewidth - 14\tabcolsep) * \real{0.1250}}
  >{\raggedright\arraybackslash}p{(\linewidth - 14\tabcolsep) * \real{0.1250}}
  >{\raggedright\arraybackslash}p{(\linewidth - 14\tabcolsep) * \real{0.1250}}
  >{\raggedright\arraybackslash}p{(\linewidth - 14\tabcolsep) * \real{0.1250}}@{}}
\toprule\noalign{}
\begin{minipage}[b]{\linewidth}\raggedright
RUNSPEC
\end{minipage} & \begin{minipage}[b]{\linewidth}\raggedright
GRID
\end{minipage} & \begin{minipage}[b]{\linewidth}\raggedright
EDIT
\end{minipage} & \begin{minipage}[b]{\linewidth}\raggedright
PROPS
\end{minipage} & \begin{minipage}[b]{\linewidth}\raggedright
REGIONS
\end{minipage} & \begin{minipage}[b]{\linewidth}\raggedright
SOLUTION
\end{minipage} & \begin{minipage}[b]{\linewidth}\raggedright
SUMMARY
\end{minipage} & \begin{minipage}[b]{\linewidth}\raggedright
SCHEDULE
\end{minipage} \\
\midrule\noalign{}
\endhead
\bottomrule\noalign{}
\endlastfoot
\end{longtable}

\subsubsection{Description}\label{description-482}

IKRORG defines the scaling parameter for the imbibition relative
permeability of oil at the critical gas saturation (ISGCR), for all the
cells in the model via an array. The ENDSCALE keyword in the RUNSPEC
section should be activated to enable end-point scaling and the use of
this keyword. In addition, the HYSTER option on the SATOPTS keyword in
the RUNSPEC section has to be activated to invoke the Hysteresis option.
The SCALCERS keyword in the PROPS section defines the options used in
the re-scaling process, the options are two point scaling and three
point scaling.

\begin{longtable}[]{@{}
  >{\raggedright\arraybackslash}p{(\linewidth - 6\tabcolsep) * \real{0.2500}}
  >{\raggedright\arraybackslash}p{(\linewidth - 6\tabcolsep) * \real{0.2500}}
  >{\raggedright\arraybackslash}p{(\linewidth - 6\tabcolsep) * \real{0.2500}}
  >{\raggedright\arraybackslash}p{(\linewidth - 6\tabcolsep) * \real{0.2500}}@{}}
\toprule\noalign{}
\begin{minipage}[b]{\linewidth}\raggedright
No.
\end{minipage} & \begin{minipage}[b]{\linewidth}\raggedright
Name
\end{minipage} & \begin{minipage}[b]{\linewidth}\raggedright
Description
\end{minipage} & \begin{minipage}[b]{\linewidth}\raggedright
Default
\end{minipage} \\
\midrule\noalign{}
\endhead
\bottomrule\noalign{}
\endlastfoot
Field & Metric & Laboratory & \\
1 & IKRORG & IKRORG is an array of positive real numbers which are
greater than zero and less than or equal to 1.0, that are the assigned
imbibition scaling IKRORG values for each cell in the model. Repeat
counts may be used, for example 50*0.850. & Taken from cell allocated
relative permeability table. \\
dimensionless & dimensionless & dimensionless & \\
Notes: & & & \\
\end{longtable}

\emph{Table 8.53: IKRORG Keyword Description}

When the IKRORG keyword is present in the input deck then the scaling
matches the imbibition relative permeability at the critical saturation
of the displacing phase.

If three point scaling option has been selected via the SCALECRS keyword
in the PROPS section the critical displacing phase is defined as:

\begin{longtable}[]{@{}lll@{}}
\toprule\noalign{}
No & Keywords Present & Critical Saturation \\
\midrule\noalign{}
\endhead
\bottomrule\noalign{}
\endlastfoot
1 & IKRORW & S critical = 1.0 -- ISWCR - ISGL \\
2 & IKRORG & S critical = 1.0 -- ISGCR - SWL \\
\end{longtable}

\emph{Table 8.54: Critical Displacement Relationships}

End-point scaling allows the entered relative permeability functions to
be re-scaled based on the saturation end-points allocated to each cell
by the ISWL, ISWCR, ISWU, ISGL, ISGCR, ISGU, ISOWCR, and ISOGCR
saturation grid arrays for the saturation end-points, In addition
end-point scaling may be directional dependent in which case the
directional dependent versions of the aforementioned arrays should be
used, that is ISWUX, ISWUY and ISWUZ instead of ISWU, There is also the
facility to make the directional end-point scaling reversible or
non-reversible and if the non-reversible option is selected the
non-reversible versions of the aforementioned arrays should be used,
that is ISWUX, ISWUX-, ISWUY, ISWUY-, ISWUZ and ISWUZ-, instead of the
ISWU keyword.

End-point scaling also allows the entered relative permeability
functions to be scale on the relative permeability values using the
IKRG, IKRGR, IKRO, IKRORG, IKRORW, IKRW and IKRWR relative permeability
grid cell arrays for the relative permeability end-point data. In
addition end-point scaling may be directional dependent in which case
the directional dependent versions of the aforementioned arrays should
be used, that is IKRGX, IKRGY and IKRGZ instead of IKRG, There is also
the facility to make the directional end-point scaling reversible or
non-reversible and if the non-reversible option is selected the
non-reversible versions of the aforementioned arrays should be used,
that is IKRGX, IKRGX-, IKRGY, IKRGY-,I KRGZ and IKRGZ-, instead of the
IKRG keyword.

\subsubsection{Examples}\label{examples-26}

The first example defines an input box for the whole grid and for layers
one to three, for layer one IKRORG is set equal to 0.755, for layer two
IKRORG equals 0.775, and for layer three IKRORG equals 0.800.

\begin{verbatim}
--
--       DEFINE INPUT BOX FOR EDITING INPUT ARRAYS (NX=100, NY=100)
--
--       ---------- BOX ---------
--       I1  I2   J1  J2   K1  K2
BOX
         1*  1*   1*  1*   1   3                           / DEFINE BOX AREA
--
--       SET IKRORG VALUES FOR THREE LAYERS IN THE MODEL
--
IKRORG
         10000*0.755  10000*0.775  10000*0.800                /
--
--       DEFINE END OF INPUT BOX EDITING OF INPUT ARRAYS
--
ENDBOX
\end{verbatim}

The next example does exactly the same thing using the EQUALS keyword
instead.

\begin{verbatim}
--
--       ARRAY       CONSTANT     ---------- BOX ---------
--                                I1  I2   J1  J2   K1  K2
EQUALS
         IKRORG      0.7550       1*  1*   1*  1*    1   1  / IKRORG FOR LAYER 1
         IKRORG      0.7750       1*  1*   1*  1*    2   2  / IKRORG FOR LAYER 2
         IKRORG      0.8000       1*  1*   1*  1*    3   3  / IKRORG FOR LAYER 3
/
\end{verbatim}

\subsection{IKRORW -- End-Point Scaling of Grid Cell Kro(Swcr)
(Imbibition)}\label{ikrorw-end-point-scaling-of-grid-cell-kroswcr-imbibition}

\begin{longtable}[]{@{}
  >{\raggedright\arraybackslash}p{(\linewidth - 14\tabcolsep) * \real{0.1250}}
  >{\raggedright\arraybackslash}p{(\linewidth - 14\tabcolsep) * \real{0.1250}}
  >{\raggedright\arraybackslash}p{(\linewidth - 14\tabcolsep) * \real{0.1250}}
  >{\raggedright\arraybackslash}p{(\linewidth - 14\tabcolsep) * \real{0.1250}}
  >{\raggedright\arraybackslash}p{(\linewidth - 14\tabcolsep) * \real{0.1250}}
  >{\raggedright\arraybackslash}p{(\linewidth - 14\tabcolsep) * \real{0.1250}}
  >{\raggedright\arraybackslash}p{(\linewidth - 14\tabcolsep) * \real{0.1250}}
  >{\raggedright\arraybackslash}p{(\linewidth - 14\tabcolsep) * \real{0.1250}}@{}}
\toprule\noalign{}
\begin{minipage}[b]{\linewidth}\raggedright
RUNSPEC
\end{minipage} & \begin{minipage}[b]{\linewidth}\raggedright
GRID
\end{minipage} & \begin{minipage}[b]{\linewidth}\raggedright
EDIT
\end{minipage} & \begin{minipage}[b]{\linewidth}\raggedright
PROPS
\end{minipage} & \begin{minipage}[b]{\linewidth}\raggedright
REGIONS
\end{minipage} & \begin{minipage}[b]{\linewidth}\raggedright
SOLUTION
\end{minipage} & \begin{minipage}[b]{\linewidth}\raggedright
SUMMARY
\end{minipage} & \begin{minipage}[b]{\linewidth}\raggedright
SCHEDULE
\end{minipage} \\
\midrule\noalign{}
\endhead
\bottomrule\noalign{}
\endlastfoot
\end{longtable}

\subsubsection{Description}\label{description-483}

IKRORW defines the scaling parameter for the imbibition relative
permeability of oil at the critical water saturation (ISWCR), for all
the cells in the model via an array. The ENDSCALE keyword in the RUNSPEC
section should be activated to enable end-point scaling and the use of
this keyword. In addition, the HYSTER option on the SATOPTS keyword in
the RUNSPEC section has to be activated to invoke the Hysteresis option.
The SCALECRS keyword in the PROPS section defines the options used in
the re-scaling process, the options are two point scaling and three
point scaling.

\begin{longtable}[]{@{}
  >{\raggedright\arraybackslash}p{(\linewidth - 6\tabcolsep) * \real{0.2500}}
  >{\raggedright\arraybackslash}p{(\linewidth - 6\tabcolsep) * \real{0.2500}}
  >{\raggedright\arraybackslash}p{(\linewidth - 6\tabcolsep) * \real{0.2500}}
  >{\raggedright\arraybackslash}p{(\linewidth - 6\tabcolsep) * \real{0.2500}}@{}}
\toprule\noalign{}
\begin{minipage}[b]{\linewidth}\raggedright
No.
\end{minipage} & \begin{minipage}[b]{\linewidth}\raggedright
Name
\end{minipage} & \begin{minipage}[b]{\linewidth}\raggedright
Description
\end{minipage} & \begin{minipage}[b]{\linewidth}\raggedright
Default
\end{minipage} \\
\midrule\noalign{}
\endhead
\bottomrule\noalign{}
\endlastfoot
Field & Metric & Laboratory & \\
1 & IKRORW & IKRORW is an array of positive real numbers which are
greater than zero and less than or equal to 1.0, that are the assigned
imbibition scaling IKRORW values for each cell in the model. Repeat
counts may be used, for example 50*0.850 & Taken from cell allocated
relative permeability table. \\
dimensionless & dimensionless & dimensionless & \\
Notes: & & & \\
\end{longtable}

\emph{Table 8.55: IKRORW Keyword Description}

When the IKRORW keyword is present in the input deck then the scaling
matches the imbibition relative permeability at the critical saturation
of the displacing phase.

If three point scaling option has been selected via the SCALECRS keyword
in the PROPS section the critical displacing phase is defined as:

\begin{longtable}[]{@{}lll@{}}
\toprule\noalign{}
No & Keywords Present & Critical Saturation \\
\midrule\noalign{}
\endhead
\bottomrule\noalign{}
\endlastfoot
1 & IKRORW & S critical = 1.0 -- ISWCR - ISGL \\
2 & IKRORG & S critical = 1.0 -- ISGCR - ISWL \\
\end{longtable}

\emph{Table 8.56: Critical Displacement Relationships}

End-point scaling allows the entered relative permeability functions to
be re-scaled based on the saturation end-points allocated to each cell
by the ISWL, ISWCR, ISWU, ISGL, ISGCR, ISGU, ISOWCR, and ISOGCR
saturation grid arrays for the saturation end-points, In addition
end-point scaling may be directional dependent in which case the
directional dependent versions of the aforementioned arrays should be
used, that is ISWUX, ISWUY and ISWUZ instead of ISWU, There is also the
facility to make the directional end-point scaling reversible or
non-reversible and if the non-reversible option is selected the
non-reversible versions of the aforementioned arrays should be used,
that is ISWUX, ISWUX-, ISWUY, ISWUY-, ISWUZ and ISWUZ-, instead of the
ISWU keyword.

End-point scaling also allows the entered relative permeability
functions to be scale on the relative permeability values using the
IKRG, IKRGR, IKRO, IKRORG, IKRORW, IKRW and IKRWR relative permeability
grid cell arrays for the relative permeability end-point data. In
addition end-point scaling may be directional dependent in which case
the directional dependent versions of the aforementioned arrays should
be used, that is IKRGX, IKRGY and IKRGZ instead of IKRG, There is also
the facility to make the directional end-point scaling reversible or
non-reversible and if the non-reversible option is selected the
non-reversible versions of the aforementioned arrays should be used,
that is IKRGX, IKRGX-, IKRGY, IKRGY-,I KRGZ and IKRGZ-, instead of the
IKRG keyword.

\subsubsection{Examples}\label{examples-27}

The first example defines an input box for the whole grid and for layers
one to three, for layer one IKRORW is set equal to 0.755, for layer two
IKRORW equals 0.775, and for layer three IKRORW equals 0.800.

\begin{verbatim}
--
--       DEFINE INPUT BOX FOR EDITING INPUT ARRAYS (NX=100, NY=100)
--
--       ---------- BOX ---------
--       I1  I2   J1  J2   K1  K2
BOX
         1*  1*   1*  1*   1   3                           / DEFINE BOX AREA
--
--       SET IKRORW VALUES FOR THREE LAYERS IN THE MODEL
--
IKRORW
         10000*0.755  10000*0.775  10000*0.800                /
--
--       DEFINE END OF INPUT BOX EDITING OF INPUT ARRAYS
--
ENDBOX
\end{verbatim}

The next example does exactly the same thing using the EQUALS keyword
instead.

\begin{verbatim}
--
--       ARRAY       CONSTANT     ---------- BOX ---------
--                                I1  I2   J1  J2   K1  K2
EQUALS
         IKRORW      0.7550       1*  1*   1*  1*    1   1  / IKRORW FOR LAYER 1
         IKRORW      0.7750       1*  1*   1*  1*    2   2  / IKRORW FOR LAYER 2
         IKRORW      0.8000       1*  1*   1*  1*    3   3  / IKRORW FOR LAYER 3
/
\end{verbatim}

\subsection{IKRW -- End-Point Scaling of Grid Cell Krw(Sw =1.0)
(Imbibition)}\label{ikrw-end-point-scaling-of-grid-cell-krwsw-1.0-imbibition}

\begin{longtable}[]{@{}
  >{\raggedright\arraybackslash}p{(\linewidth - 14\tabcolsep) * \real{0.1250}}
  >{\raggedright\arraybackslash}p{(\linewidth - 14\tabcolsep) * \real{0.1250}}
  >{\raggedright\arraybackslash}p{(\linewidth - 14\tabcolsep) * \real{0.1250}}
  >{\raggedright\arraybackslash}p{(\linewidth - 14\tabcolsep) * \real{0.1250}}
  >{\raggedright\arraybackslash}p{(\linewidth - 14\tabcolsep) * \real{0.1250}}
  >{\raggedright\arraybackslash}p{(\linewidth - 14\tabcolsep) * \real{0.1250}}
  >{\raggedright\arraybackslash}p{(\linewidth - 14\tabcolsep) * \real{0.1250}}
  >{\raggedright\arraybackslash}p{(\linewidth - 14\tabcolsep) * \real{0.1250}}@{}}
\toprule\noalign{}
\begin{minipage}[b]{\linewidth}\raggedright
RUNSPEC
\end{minipage} & \begin{minipage}[b]{\linewidth}\raggedright
GRID
\end{minipage} & \begin{minipage}[b]{\linewidth}\raggedright
EDIT
\end{minipage} & \begin{minipage}[b]{\linewidth}\raggedright
PROPS
\end{minipage} & \begin{minipage}[b]{\linewidth}\raggedright
REGIONS
\end{minipage} & \begin{minipage}[b]{\linewidth}\raggedright
SOLUTION
\end{minipage} & \begin{minipage}[b]{\linewidth}\raggedright
SUMMARY
\end{minipage} & \begin{minipage}[b]{\linewidth}\raggedright
SCHEDULE
\end{minipage} \\
\midrule\noalign{}
\endhead
\bottomrule\noalign{}
\endlastfoot
\end{longtable}

\subsubsection{Description}\label{description-484}

IKRW defines the scaling parameter at the maximum imbibition water
relative permeability value (ISWU), that is for Sw = 1.0, for all the
cells in the model via an array. The ENDSCALE keyword in the RUNSPEC
section should be activated to enable end-point scaling and the use of
this keyword. In addition, the HYSTER option on the SATOPTS keyword in
the RUNSPEC section has to be activated to invoke the Hysteresis option.
The SCALCERS keyword in the PROPS section defines the options used in
the re-scaling process, the options are two point scaling and three
point scaling.

\begin{longtable}[]{@{}
  >{\raggedright\arraybackslash}p{(\linewidth - 6\tabcolsep) * \real{0.2500}}
  >{\raggedright\arraybackslash}p{(\linewidth - 6\tabcolsep) * \real{0.2500}}
  >{\raggedright\arraybackslash}p{(\linewidth - 6\tabcolsep) * \real{0.2500}}
  >{\raggedright\arraybackslash}p{(\linewidth - 6\tabcolsep) * \real{0.2500}}@{}}
\toprule\noalign{}
\begin{minipage}[b]{\linewidth}\raggedright
No.
\end{minipage} & \begin{minipage}[b]{\linewidth}\raggedright
Name
\end{minipage} & \begin{minipage}[b]{\linewidth}\raggedright
Description
\end{minipage} & \begin{minipage}[b]{\linewidth}\raggedright
Default
\end{minipage} \\
\midrule\noalign{}
\endhead
\bottomrule\noalign{}
\endlastfoot
Field & Metric & Laboratory & \\
1 & IKRW & IKRW is an array of positive real numbers which are greater
than zero and less than or equal to 1.0, that are the assigned
imbibition scaling IKRW values for each cell in the model. Repeat counts
may be used, for example 50*1.000. & Taken from cell allocated relative
permeability table. \\
dimensionless & dimensionless & dimensionless & \\
Notes: & & & \\
\end{longtable}

\emph{Table 8.57: IKRW Keyword Description}

\begin{verbatim}
\end{verbatim}

For the two point scaling option and for the IKRWR water relative
permeability array NOT present in the input deck the krw value for a
grid block is scaled by:

\begin{equation}\phantomsection\label{eq-8-58}{
{k}_{\mathit{rw}} = {k}_{{\mathit{rw}}_{ \mathit{TABLE}}}\left(\frac{\mathit{IKRW}}{{k}_{{\mathit{rw}}_{ \mathit{TABLE}-\mathit{MAX}}}}\right)
}\end{equation}

Where:

\(\mathit{krw}\) = the resulting IKRW value for a grid cell.

\(\mathit{IKRW}\) = the scaling water relative permeability value from
the IKRW array for a

given cell.

\({k}_{{\mathit{rw}}_{ \mathit{TABLE}}}\) = the water relative
permeability from a grid block's oil relative permeability

table at the grid blocks water saturation.

\({k}_{{\mathit{rw}}_{ \mathit{TABLE}-\mathit{MAX}}}\) = the maximum
water relative permeability from a grid block's water

relative table, that is at the maximum water saturation.

If the IKRWR keyword is present in the input deck then the scaling
matches the imbibition relative permeability at the critical saturation
of the displacing phase.

If three point scaling option has been selected via the SCALECRS keyword
in the PROPS section the critical displacing phase is defined as:

\begin{longtable}[]{@{}lll@{}}
\toprule\noalign{}
No & Phases Present & Critical Saturation \\
\midrule\noalign{}
\endhead
\bottomrule\noalign{}
\endlastfoot
1 & Gas-Oil & S critical = 1.0 -- ISOWCR - ISGL \\
2 & Gas-Oil-Water & S critical = 1.0 -- ISOWCR - ISGL \\
3 & Gas-Water & S critical = 1.0 -- ISGCR \\
\end{longtable}

\emph{Table 8.58: Critical Displacement Relationships}

End-point scaling allows the entered relative permeability functions to
be re-scaled based on the saturation end-points allocated to each cell
by the ISWL, ISWCR, ISWU, ISGL, ISGCR, ISGU, ISOWCR, and ISOGCR
saturation grid arrays for the saturation end-points, In addition
end-point scaling may be directional dependent in which case the
directional dependent versions of the aforementioned arrays should be
used, that is ISWUX, ISWUY and ISWUZ instead of ISWU, There is also the
facility to make the directional end-point scaling reversible or
non-reversible and if the non-reversible option is selected the
non-reversible versions of the aforementioned arrays should be used,
that is ISWUX, ISWUX-, ISWUY, ISWUY-, ISWUZ and ISWUZ-, instead of the
ISWU keyword.

End-point scaling also allows the entered relative permeability
functions to be scale on the relative permeability values using the
IKRG, IKRGR, IKRO, IKRORG, IKRORW, IKRW and IKRWR relative permeability
grid cell arrays for the relative permeability end-point data. In
addition end-point scaling may be directional dependent in which case
the directional dependent versions of the aforementioned arrays should
be used, that is IKRGX, IKRGY and IKRGZ instead of IKRG, There is also
the facility to make the directional end-point scaling reversible or
non-reversible and if the non-reversible option is selected the
non-reversible versions of the aforementioned arrays should be used,
that is IKRGX, IKRGX-, IKRGY, IKRGY-,I KRGZ and IKRGZ-, instead of the
IKRG keyword.

\subsubsection{Example}\label{example-286}

The example below defines an input box for the whole grid and for layers
one to three, for layer one IKRW is set equal to 0.850, for layer two
IKRW equals 0.875, and for layer three IKRW equals 0.900.

\begin{verbatim}
--
--       DEFINE INPUT BOX FOR EDITING INPUT ARRAYS (NX=100, NY=100)
--
--       ---------- BOX ---------
--       I1  I2   J1  J2   K1  K2
BOX
         1*  1*   1*  1*   1   3                           / DEFINE BOX AREA
--
--       SET IKRW VALUES FOR THREE LAYERS IN THE MODEL
--
IKRW
         10000*0.850  10000*0.875  10000*0.900                /
--
--       DEFINE END OF INPUT BOX EDITING OF INPUT ARRAYS
--
ENDBOX
\end{verbatim}

\subsection{IKRWR -- End-Point Scaling of Grid Cell KRWR(Sowcr)
(Imbibition)}\label{ikrwr-end-point-scaling-of-grid-cell-krwrsowcr-imbibition}

\begin{longtable}[]{@{}
  >{\raggedright\arraybackslash}p{(\linewidth - 14\tabcolsep) * \real{0.1250}}
  >{\raggedright\arraybackslash}p{(\linewidth - 14\tabcolsep) * \real{0.1250}}
  >{\raggedright\arraybackslash}p{(\linewidth - 14\tabcolsep) * \real{0.1250}}
  >{\raggedright\arraybackslash}p{(\linewidth - 14\tabcolsep) * \real{0.1250}}
  >{\raggedright\arraybackslash}p{(\linewidth - 14\tabcolsep) * \real{0.1250}}
  >{\raggedright\arraybackslash}p{(\linewidth - 14\tabcolsep) * \real{0.1250}}
  >{\raggedright\arraybackslash}p{(\linewidth - 14\tabcolsep) * \real{0.1250}}
  >{\raggedright\arraybackslash}p{(\linewidth - 14\tabcolsep) * \real{0.1250}}@{}}
\toprule\noalign{}
\begin{minipage}[b]{\linewidth}\raggedright
RUNSPEC
\end{minipage} & \begin{minipage}[b]{\linewidth}\raggedright
GRID
\end{minipage} & \begin{minipage}[b]{\linewidth}\raggedright
EDIT
\end{minipage} & \begin{minipage}[b]{\linewidth}\raggedright
PROPS
\end{minipage} & \begin{minipage}[b]{\linewidth}\raggedright
REGIONS
\end{minipage} & \begin{minipage}[b]{\linewidth}\raggedright
SOLUTION
\end{minipage} & \begin{minipage}[b]{\linewidth}\raggedright
SUMMARY
\end{minipage} & \begin{minipage}[b]{\linewidth}\raggedright
SCHEDULE
\end{minipage} \\
\midrule\noalign{}
\endhead
\bottomrule\noalign{}
\endlastfoot
\end{longtable}

\subsubsection{Description}\label{description-485}

IKRWR defines the scaling parameter at the imbibition critical oil to
water saturation value (ISOWCR), for the imbibition water relative
permeability curve, for all the cells in the model via an array. The
ENDSCALE keyword in the RUNSPEC section should be activated to enable
end-point scaling and the use of this keyword. In addition, the HYSTER
option on the SATOPTS keyword in the RUNSPEC section has to be activated
to invoke the Hysteresis option. The SCALCERS keyword in the PROPS
section defines the options used in the re-scaling process, the options
are two point scaling and three point scaling.

\begin{longtable}[]{@{}
  >{\raggedright\arraybackslash}p{(\linewidth - 6\tabcolsep) * \real{0.2500}}
  >{\raggedright\arraybackslash}p{(\linewidth - 6\tabcolsep) * \real{0.2500}}
  >{\raggedright\arraybackslash}p{(\linewidth - 6\tabcolsep) * \real{0.2500}}
  >{\raggedright\arraybackslash}p{(\linewidth - 6\tabcolsep) * \real{0.2500}}@{}}
\toprule\noalign{}
\begin{minipage}[b]{\linewidth}\raggedright
No.
\end{minipage} & \begin{minipage}[b]{\linewidth}\raggedright
Name
\end{minipage} & \begin{minipage}[b]{\linewidth}\raggedright
Description
\end{minipage} & \begin{minipage}[b]{\linewidth}\raggedright
Default
\end{minipage} \\
\midrule\noalign{}
\endhead
\bottomrule\noalign{}
\endlastfoot
Field & Metric & Laboratory & \\
1 & IKRWR & IKRWR is an array of positive real numbers which are greater
than zero and less than or equal to 1.0, that are the assigned
imbibition scaling IKRWR values for each cell in the model. Repeat
counts may be used, for example 50*1.000. & Taken from cell allocated
relative permeability table. \\
dimensionless & dimensionless & dimensionless & \\
Notes: & & & \\
\end{longtable}

\emph{Table 8.59: IKRWR Keyword Description}

\begin{verbatim}
\end{verbatim}

When the IKRWR keyword is present in the input deck then the scaling
matches the imbibition relative permeability at the critical saturation
of the displacing phase.

If three point scaling option has been selected via the SCALECRS keyword
in the PROPS section the critical displacing phase is defined as:

\begin{longtable}[]{@{}lll@{}}
\toprule\noalign{}
No & Phases Present & Critical Saturation \\
\midrule\noalign{}
\endhead
\bottomrule\noalign{}
\endlastfoot
1 & Gas-Oil & S critical = 1.0 -- ISOWCR - ISGL \\
2 & Gas-Oil-Water & S critical = 1.0 -- ISOWCR - ISGL \\
3 & Gas-Water & S critical = 1.0 -- ISGCR \\
\end{longtable}

\emph{Table 8.60: Critical Displacement Relationships}

End-point scaling allows the entered relative permeability functions to
be re-scaled based on the saturation end-points allocated to each cell
by the ISWL, ISWCR, ISWU, ISGL, ISGCR, ISGU, ISOWCR, and ISOGCR
saturation grid arrays for the saturation end-points, In addition
end-point scaling may be directional dependent in which case the
directional dependent versions of the aforementioned arrays should be
used, that is ISWUX, ISWUY and ISWUZ instead of ISWU, There is also the
facility to make the directional end-point scaling reversible or
non-reversible and if the non-reversible option is selected the
non-reversible versions of the aforementioned arrays should be used,
that is ISWUX, ISWUX-, ISWUY, ISWUY-, ISWUZ and ISWUZ-, instead of the
ISWU keyword.

End-point scaling also allows the entered relative permeability
functions to be scale on the relative permeability values using the
IKRG, IKRGR, IKRO, IKRORG, IKRORW, IKRW and IKRWR relative permeability
grid cell arrays for the relative permeability end-point data. In
addition end-point scaling may be directional dependent in which case
the directional dependent versions of the aforementioned arrays should
be used, that is IKRGX, IKRGY and IKRGZ instead of IKRG, There is also
the facility to make the directional end-point scaling reversible or
non-reversible and if the non-reversible option is selected the
non-reversible versions of the aforementioned arrays should be used,
that is IKRGX, IKRGX-, IKRGY, IKRGY-,I KRGZ and IKRGZ-, instead of the
IKRG keyword.

\subsubsection{Examples}\label{examples-28}

The first example defines an input box for the whole grid and for layers
one to three, for layer one IKRWR is set equal to 0.755, for layer two
IKRWR equals 0.775, and for layer three IKRWR equals 0.800.

\begin{verbatim}
--
--       DEFINE INPUT BOX FOR EDITING INPUT ARRAYS (NX=100, NY=100)
--
--       ---------- BOX ---------
--       I1  I2   J1  J2   K1  K2
BOX
         1*  1*   1*  1*   1   3                           / DEFINE BOX AREA
--
--       SET IKRWR VALUES FOR THREE LAYERS IN THE MODEL
--
IKRWR
         10000*0.755  10000*0.775  10000*0.800                /
--
--       DEFINE END OF INPUT BOX EDITING OF INPUT ARRAYS
--
ENDBOX
\end{verbatim}

The next example does exactly the same thing using the EQUALS keyword
instead.

\begin{verbatim}
--
--       ARRAY       CONSTANT     ---------- BOX ---------
--                                I1  I2   J1  J2   K1  K2
EQUALS
         IKRWR       0.7550       1*  1*   1*  1*    1   1  / IKRWR FOR LAYER 1
         IKRWR       0.7750       1*  1*   1*  1*    2   2  / IKRWR FOR LAYER 2
         IKRWR       0.8000       1*  1*   1*  1*    3   3  / IKRWR FOR LAYER 3
/
\end{verbatim}

\subsection{IMKRVD -- Imbibition Relative Permeability End-Points versus
Depth
Functions}\label{imkrvd-imbibition-relative-permeability-end-points-versus-depth-functions}

\begin{longtable}[]{@{}
  >{\raggedright\arraybackslash}p{(\linewidth - 14\tabcolsep) * \real{0.1250}}
  >{\raggedright\arraybackslash}p{(\linewidth - 14\tabcolsep) * \real{0.1250}}
  >{\raggedright\arraybackslash}p{(\linewidth - 14\tabcolsep) * \real{0.1250}}
  >{\raggedright\arraybackslash}p{(\linewidth - 14\tabcolsep) * \real{0.1250}}
  >{\raggedright\arraybackslash}p{(\linewidth - 14\tabcolsep) * \real{0.1250}}
  >{\raggedright\arraybackslash}p{(\linewidth - 14\tabcolsep) * \real{0.1250}}
  >{\raggedright\arraybackslash}p{(\linewidth - 14\tabcolsep) * \real{0.1250}}
  >{\raggedright\arraybackslash}p{(\linewidth - 14\tabcolsep) * \real{0.1250}}@{}}
\toprule\noalign{}
\begin{minipage}[b]{\linewidth}\raggedright
RUNSPEC
\end{minipage} & \begin{minipage}[b]{\linewidth}\raggedright
GRID
\end{minipage} & \begin{minipage}[b]{\linewidth}\raggedright
EDIT
\end{minipage} & \begin{minipage}[b]{\linewidth}\raggedright
PROPS
\end{minipage} & \begin{minipage}[b]{\linewidth}\raggedright
REGIONS
\end{minipage} & \begin{minipage}[b]{\linewidth}\raggedright
SOLUTION
\end{minipage} & \begin{minipage}[b]{\linewidth}\raggedright
SUMMARY
\end{minipage} & \begin{minipage}[b]{\linewidth}\raggedright
SCHEDULE
\end{minipage} \\
\midrule\noalign{}
\endhead
\bottomrule\noalign{}
\endlastfoot
\end{longtable}

\subsubsection{Description}\label{description-486}

This keyword defines the maximum imbibition oil, gas, and water relative
permeability versus depth for the three phases. The ENDSCALE keyword in
the RUNSPEC section should be activated to enable end-point scaling and
the use of this keyword. In addition, the HYSTER option on the SATOPTS
keyword in the RUNSPEC section has to be activated to invoke the
Hysteresis option.

This keyword is not supported by OPM Flow but would change the results
if supported so the simulation will be stopped.

\subsection{IMPCVD -- Imbibition Maximum Capillary Pressure versus Depth
Functions}\label{impcvd-imbibition-maximum-capillary-pressure-versus-depth-functions}

\begin{longtable}[]{@{}
  >{\raggedright\arraybackslash}p{(\linewidth - 14\tabcolsep) * \real{0.1250}}
  >{\raggedright\arraybackslash}p{(\linewidth - 14\tabcolsep) * \real{0.1250}}
  >{\raggedright\arraybackslash}p{(\linewidth - 14\tabcolsep) * \real{0.1250}}
  >{\raggedright\arraybackslash}p{(\linewidth - 14\tabcolsep) * \real{0.1250}}
  >{\raggedright\arraybackslash}p{(\linewidth - 14\tabcolsep) * \real{0.1250}}
  >{\raggedright\arraybackslash}p{(\linewidth - 14\tabcolsep) * \real{0.1250}}
  >{\raggedright\arraybackslash}p{(\linewidth - 14\tabcolsep) * \real{0.1250}}
  >{\raggedright\arraybackslash}p{(\linewidth - 14\tabcolsep) * \real{0.1250}}@{}}
\toprule\noalign{}
\begin{minipage}[b]{\linewidth}\raggedright
RUNSPEC
\end{minipage} & \begin{minipage}[b]{\linewidth}\raggedright
GRID
\end{minipage} & \begin{minipage}[b]{\linewidth}\raggedright
EDIT
\end{minipage} & \begin{minipage}[b]{\linewidth}\raggedright
PROPS
\end{minipage} & \begin{minipage}[b]{\linewidth}\raggedright
REGIONS
\end{minipage} & \begin{minipage}[b]{\linewidth}\raggedright
SOLUTION
\end{minipage} & \begin{minipage}[b]{\linewidth}\raggedright
SUMMARY
\end{minipage} & \begin{minipage}[b]{\linewidth}\raggedright
SCHEDULE
\end{minipage} \\
\midrule\noalign{}
\endhead
\bottomrule\noalign{}
\endlastfoot
\end{longtable}

\subsubsection{Description}\label{description-487}

This keyword defines the maximum imbibition gas-oil and water-oil
capillary pressure values versus depth for when the end-point scaling
option has been invoked by the ENDSCALE keyword in the RUNSPEC section
and the HYSTER option on the SATOPTS keyword in the RUNSPEC section has
been activated to invoke the Hysteresis option.

This keyword is not supported by OPM Flow but would change the results
if supported so the simulation will be stopped.

\subsection{IMPORT -- Import Grid File Data at the Current
Position}\label{import-import-grid-file-data-at-the-current-position-2}

The IMPORT keyword informs the simulator to continue reading input data
from the specified IMPORT file. When the end of the IMPORT file is
reached, input data is read from the next keyword in the current file.
Normally IMPORT files are generated by grid pre-processing software and
the keyword allows for both formatted and unformatted (binary) files to
be loaded.

See IMPORT -- Import Grid File Data at the Current Position in the GRID
section for a full description.

\subsection{IMPTVD -- Imbibition Relative Permeability Saturation
End-Points versus
Depth}\label{imptvd-imbibition-relative-permeability-saturation-end-points-versus-depth}

\begin{longtable}[]{@{}
  >{\raggedright\arraybackslash}p{(\linewidth - 14\tabcolsep) * \real{0.1250}}
  >{\raggedright\arraybackslash}p{(\linewidth - 14\tabcolsep) * \real{0.1250}}
  >{\raggedright\arraybackslash}p{(\linewidth - 14\tabcolsep) * \real{0.1250}}
  >{\raggedright\arraybackslash}p{(\linewidth - 14\tabcolsep) * \real{0.1250}}
  >{\raggedright\arraybackslash}p{(\linewidth - 14\tabcolsep) * \real{0.1250}}
  >{\raggedright\arraybackslash}p{(\linewidth - 14\tabcolsep) * \real{0.1250}}
  >{\raggedright\arraybackslash}p{(\linewidth - 14\tabcolsep) * \real{0.1250}}
  >{\raggedright\arraybackslash}p{(\linewidth - 14\tabcolsep) * \real{0.1250}}@{}}
\toprule\noalign{}
\begin{minipage}[b]{\linewidth}\raggedright
RUNSPEC
\end{minipage} & \begin{minipage}[b]{\linewidth}\raggedright
GRID
\end{minipage} & \begin{minipage}[b]{\linewidth}\raggedright
EDIT
\end{minipage} & \begin{minipage}[b]{\linewidth}\raggedright
PROPS
\end{minipage} & \begin{minipage}[b]{\linewidth}\raggedright
REGIONS
\end{minipage} & \begin{minipage}[b]{\linewidth}\raggedright
SOLUTION
\end{minipage} & \begin{minipage}[b]{\linewidth}\raggedright
SUMMARY
\end{minipage} & \begin{minipage}[b]{\linewidth}\raggedright
SCHEDULE
\end{minipage} \\
\midrule\noalign{}
\endhead
\bottomrule\noalign{}
\endlastfoot
\end{longtable}

\subsubsection{Description}\label{description-488}

This keyword defines the variation of the imbibition relative
permeability saturation end-points (SWL, SWCR, etc.) for all three
phases versus depth., for when the end-point scaling option has been
invoked by the ENDSCALE keyword in the RUNSPEC section, and the HYSTER
option on the SATOPTS keyword in the RUNSPEC section has been activated
to invoke the Hysteresis option. This functionality is not supported in
OPM Flow.

This keyword is not supported by OPM Flow but would change the results
if supported so the simulation will be stopped.

\subsection{IMSPCVD -- Imbibition Capillary Pressure Connate Saturations
versus
Depth}\label{imspcvd-imbibition-capillary-pressure-connate-saturations-versus-depth}

\begin{longtable}[]{@{}
  >{\raggedright\arraybackslash}p{(\linewidth - 14\tabcolsep) * \real{0.1250}}
  >{\raggedright\arraybackslash}p{(\linewidth - 14\tabcolsep) * \real{0.1250}}
  >{\raggedright\arraybackslash}p{(\linewidth - 14\tabcolsep) * \real{0.1250}}
  >{\raggedright\arraybackslash}p{(\linewidth - 14\tabcolsep) * \real{0.1250}}
  >{\raggedright\arraybackslash}p{(\linewidth - 14\tabcolsep) * \real{0.1250}}
  >{\raggedright\arraybackslash}p{(\linewidth - 14\tabcolsep) * \real{0.1250}}
  >{\raggedright\arraybackslash}p{(\linewidth - 14\tabcolsep) * \real{0.1250}}
  >{\raggedright\arraybackslash}p{(\linewidth - 14\tabcolsep) * \real{0.1250}}@{}}
\toprule\noalign{}
\begin{minipage}[b]{\linewidth}\raggedright
RUNSPEC
\end{minipage} & \begin{minipage}[b]{\linewidth}\raggedright
GRID
\end{minipage} & \begin{minipage}[b]{\linewidth}\raggedright
EDIT
\end{minipage} & \begin{minipage}[b]{\linewidth}\raggedright
PROPS
\end{minipage} & \begin{minipage}[b]{\linewidth}\raggedright
REGIONS
\end{minipage} & \begin{minipage}[b]{\linewidth}\raggedright
SOLUTION
\end{minipage} & \begin{minipage}[b]{\linewidth}\raggedright
SUMMARY
\end{minipage} & \begin{minipage}[b]{\linewidth}\raggedright
SCHEDULE
\end{minipage} \\
\midrule\noalign{}
\endhead
\bottomrule\noalign{}
\endlastfoot
\end{longtable}

\subsubsection{Description}\label{description-489}

This keyword defines the imbibition capillary pressure gas and water
connate saturations values versus depth for when the end-point scaling
option has been invoked by the ENDSCALE keyword in the RUNSPEC section,
and the HYSTER option on the SATOPTS keyword in the RUNSPEC section has
been activated to invoke the Hysteresis option. This functionality is
not supported in OPM Flow.

This keyword is not supported by OPM Flow but would change the results
if supported so the simulation will be stopped.

\subsection{INCLUDE -- Load Another Data File at the Current
Position}\label{include-load-another-data-file-at-the-current-position-4}

The INCLUDE keyword informs OPM Flow to continue reading input data from
the specified INCLUDE file. When the end of the INCLUDE file is reached,
or the ENDINC keyword is encountered in the included file, input data is
read from the next keyword in the current file.

See INCLUDE -- Load Another Data File at the Current Position in the
GLOBAL section for a full description.

\subsection{INTPC -- Activate Dual Porosity Integrated Capillary
Pressure
Option}\label{intpc-activate-dual-porosity-integrated-capillary-pressure-option}

\begin{longtable}[]{@{}
  >{\raggedright\arraybackslash}p{(\linewidth - 14\tabcolsep) * \real{0.1250}}
  >{\raggedright\arraybackslash}p{(\linewidth - 14\tabcolsep) * \real{0.1250}}
  >{\raggedright\arraybackslash}p{(\linewidth - 14\tabcolsep) * \real{0.1250}}
  >{\raggedright\arraybackslash}p{(\linewidth - 14\tabcolsep) * \real{0.1250}}
  >{\raggedright\arraybackslash}p{(\linewidth - 14\tabcolsep) * \real{0.1250}}
  >{\raggedright\arraybackslash}p{(\linewidth - 14\tabcolsep) * \real{0.1250}}
  >{\raggedright\arraybackslash}p{(\linewidth - 14\tabcolsep) * \real{0.1250}}
  >{\raggedright\arraybackslash}p{(\linewidth - 14\tabcolsep) * \real{0.1250}}@{}}
\toprule\noalign{}
\begin{minipage}[b]{\linewidth}\raggedright
RUNSPEC
\end{minipage} & \begin{minipage}[b]{\linewidth}\raggedright
GRID
\end{minipage} & \begin{minipage}[b]{\linewidth}\raggedright
EDIT
\end{minipage} & \begin{minipage}[b]{\linewidth}\raggedright
PROPS
\end{minipage} & \begin{minipage}[b]{\linewidth}\raggedright
REGIONS
\end{minipage} & \begin{minipage}[b]{\linewidth}\raggedright
SOLUTION
\end{minipage} & \begin{minipage}[b]{\linewidth}\raggedright
SUMMARY
\end{minipage} & \begin{minipage}[b]{\linewidth}\raggedright
SCHEDULE
\end{minipage} \\
\midrule\noalign{}
\endhead
\bottomrule\noalign{}
\endlastfoot
\end{longtable}

\subsubsection{Description}\label{description-490}

The INTPC keyword activates the integrated capillary pressure option for
the oil, gas or both phases, for when a Dual Porosity model has been
activated by either the DUALPORO or DUALPERM keywords in the RUNSPEC
section. In addition, the keyword can only be used if the Gravity
Drainage option has been specified by either the GRAVDR or GRAVDRM in
the RUNSPEC section. Basically, activating this feature results in the
simulator adjusting the capillary pressure curves by integrating the
matrix capillary pressure curves over the matrix block height to
calculate the average saturation.

This keyword is not supported by OPM Flow but would change the results
if supported so the simulation will be stopped.

\subsection{IONXROCK -- Define Ion Exchange Constant by Saturation Table
Regions}\label{ionxrock-define-ion-exchange-constant-by-saturation-table-regions}

\begin{longtable}[]{@{}
  >{\raggedright\arraybackslash}p{(\linewidth - 14\tabcolsep) * \real{0.1250}}
  >{\raggedright\arraybackslash}p{(\linewidth - 14\tabcolsep) * \real{0.1250}}
  >{\raggedright\arraybackslash}p{(\linewidth - 14\tabcolsep) * \real{0.1250}}
  >{\raggedright\arraybackslash}p{(\linewidth - 14\tabcolsep) * \real{0.1250}}
  >{\raggedright\arraybackslash}p{(\linewidth - 14\tabcolsep) * \real{0.1250}}
  >{\raggedright\arraybackslash}p{(\linewidth - 14\tabcolsep) * \real{0.1250}}
  >{\raggedright\arraybackslash}p{(\linewidth - 14\tabcolsep) * \real{0.1250}}
  >{\raggedright\arraybackslash}p{(\linewidth - 14\tabcolsep) * \real{0.1250}}@{}}
\toprule\noalign{}
\begin{minipage}[b]{\linewidth}\raggedright
RUNSPEC
\end{minipage} & \begin{minipage}[b]{\linewidth}\raggedright
GRID
\end{minipage} & \begin{minipage}[b]{\linewidth}\raggedright
EDIT
\end{minipage} & \begin{minipage}[b]{\linewidth}\raggedright
PROPS
\end{minipage} & \begin{minipage}[b]{\linewidth}\raggedright
REGIONS
\end{minipage} & \begin{minipage}[b]{\linewidth}\raggedright
SOLUTION
\end{minipage} & \begin{minipage}[b]{\linewidth}\raggedright
SUMMARY
\end{minipage} & \begin{minipage}[b]{\linewidth}\raggedright
SCHEDULE
\end{minipage} \\
\midrule\noalign{}
\endhead
\bottomrule\noalign{}
\endlastfoot
\end{longtable}

\subsubsection{Description}\label{description-491}

The IONXROCK keyword activates ion exchange and defines the ion exchange
constant by saturation table regions, for when the brine phase has been
activated by the BRINE keyword and the Multi-Component Brine model, that
allows for the water phase to have multiple water salinities, has been
activated by the ECLMC keyword. Both keywords are in the RUNSPEC
section.

This keyword is not supported by OPM Flow but would change the results
if supported so the simulation will be stopped.

\subsection{IONXSURF -- Define Surfactant Ion Exchange Constant by
Saturation Table
Regions}\label{ionxsurf-define-surfactant-ion-exchange-constant-by-saturation-table-regions}

\begin{longtable}[]{@{}
  >{\raggedright\arraybackslash}p{(\linewidth - 14\tabcolsep) * \real{0.1250}}
  >{\raggedright\arraybackslash}p{(\linewidth - 14\tabcolsep) * \real{0.1250}}
  >{\raggedright\arraybackslash}p{(\linewidth - 14\tabcolsep) * \real{0.1250}}
  >{\raggedright\arraybackslash}p{(\linewidth - 14\tabcolsep) * \real{0.1250}}
  >{\raggedright\arraybackslash}p{(\linewidth - 14\tabcolsep) * \real{0.1250}}
  >{\raggedright\arraybackslash}p{(\linewidth - 14\tabcolsep) * \real{0.1250}}
  >{\raggedright\arraybackslash}p{(\linewidth - 14\tabcolsep) * \real{0.1250}}
  >{\raggedright\arraybackslash}p{(\linewidth - 14\tabcolsep) * \real{0.1250}}@{}}
\toprule\noalign{}
\begin{minipage}[b]{\linewidth}\raggedright
RUNSPEC
\end{minipage} & \begin{minipage}[b]{\linewidth}\raggedright
GRID
\end{minipage} & \begin{minipage}[b]{\linewidth}\raggedright
EDIT
\end{minipage} & \begin{minipage}[b]{\linewidth}\raggedright
PROPS
\end{minipage} & \begin{minipage}[b]{\linewidth}\raggedright
REGIONS
\end{minipage} & \begin{minipage}[b]{\linewidth}\raggedright
SOLUTION
\end{minipage} & \begin{minipage}[b]{\linewidth}\raggedright
SUMMARY
\end{minipage} & \begin{minipage}[b]{\linewidth}\raggedright
SCHEDULE
\end{minipage} \\
\midrule\noalign{}
\endhead
\bottomrule\noalign{}
\endlastfoot
\end{longtable}

\subsubsection{Description}\label{description-492}

The IONXROCK keyword activates ion exchange on surfactant micellae
{[}Particle of colloidal dimensions that exists in equilibrium with the
molecules or ions in solution from which it is formed. A micella or
micelle (plural micellae or micelles, respectively) is an aggregate (or
supramolecular assembly) of surfactant molecules dispersed in a liquid
colloid. A typical micella in aqueous solution forms an aggregate with
the hydrophilic ``head'' regions in contact with surrounding solvent,
sequestering the hydrophobic single-tail regions in the micella centre
(https://en.wikipedia.org/wiki/Micelle).{]} and defines the ion exchange
constant by surfactant equivalent molecular weight for saturation table
regions, for when the brine and surfactant phases has been activated by
the BRINE and SURFACT keywords, and the Multi-Component Brine model,
that allows for the water phase to have multiple water salinities, has
been activated by the ECLMC keyword. All three keywords are in the
RUNSPEC section.

This keyword is not supported by OPM Flow but would change the results
if supported so the simulation will be stopped.

\subsection{IPCG -- End-Point Scaling of Grid Cell Gas Capillary
Pressure
(Imbibition)}\label{ipcg-end-point-scaling-of-grid-cell-gas-capillary-pressure-imbibition}

\begin{longtable}[]{@{}
  >{\raggedright\arraybackslash}p{(\linewidth - 14\tabcolsep) * \real{0.1250}}
  >{\raggedright\arraybackslash}p{(\linewidth - 14\tabcolsep) * \real{0.1250}}
  >{\raggedright\arraybackslash}p{(\linewidth - 14\tabcolsep) * \real{0.1250}}
  >{\raggedright\arraybackslash}p{(\linewidth - 14\tabcolsep) * \real{0.1250}}
  >{\raggedright\arraybackslash}p{(\linewidth - 14\tabcolsep) * \real{0.1250}}
  >{\raggedright\arraybackslash}p{(\linewidth - 14\tabcolsep) * \real{0.1250}}
  >{\raggedright\arraybackslash}p{(\linewidth - 14\tabcolsep) * \real{0.1250}}
  >{\raggedright\arraybackslash}p{(\linewidth - 14\tabcolsep) * \real{0.1250}}@{}}
\toprule\noalign{}
\begin{minipage}[b]{\linewidth}\raggedright
RUNSPEC
\end{minipage} & \begin{minipage}[b]{\linewidth}\raggedright
GRID
\end{minipage} & \begin{minipage}[b]{\linewidth}\raggedright
EDIT
\end{minipage} & \begin{minipage}[b]{\linewidth}\raggedright
PROPS
\end{minipage} & \begin{minipage}[b]{\linewidth}\raggedright
REGIONS
\end{minipage} & \begin{minipage}[b]{\linewidth}\raggedright
SOLUTION
\end{minipage} & \begin{minipage}[b]{\linewidth}\raggedright
SUMMARY
\end{minipage} & \begin{minipage}[b]{\linewidth}\raggedright
SCHEDULE
\end{minipage} \\
\midrule\noalign{}
\endhead
\bottomrule\noalign{}
\endlastfoot
\end{longtable}

\subsubsection{Description}\label{description-493}

IPCG defines the maximum imbibition gas-oil capillary pressure values
for all the cells in the model via an array. The ENDSCALE keyword in the
RUNSPEC section should be activated to enable end-point scaling and the
use of this keyword. In addition, the HYSTER option on the SATOPTS
keyword in the RUNSPEC section has to be activated to invoke the
hysteresis option. The keyword can be used with all grid types.

\begin{longtable}[]{@{}
  >{\raggedright\arraybackslash}p{(\linewidth - 6\tabcolsep) * \real{0.2500}}
  >{\raggedright\arraybackslash}p{(\linewidth - 6\tabcolsep) * \real{0.2500}}
  >{\raggedright\arraybackslash}p{(\linewidth - 6\tabcolsep) * \real{0.2500}}
  >{\raggedright\arraybackslash}p{(\linewidth - 6\tabcolsep) * \real{0.2500}}@{}}
\toprule\noalign{}
\begin{minipage}[b]{\linewidth}\raggedright
No.
\end{minipage} & \begin{minipage}[b]{\linewidth}\raggedright
Name
\end{minipage} & \begin{minipage}[b]{\linewidth}\raggedright
Description
\end{minipage} & \begin{minipage}[b]{\linewidth}\raggedright
Default
\end{minipage} \\
\midrule\noalign{}
\endhead
\bottomrule\noalign{}
\endlastfoot
Field & Metric & Laboratory & \\
1 & IPCG & IPCG is an array of positive real numbers assigning the
maximum imbibition gas capillary pressure values for each cell in the
model. Repeat counts may be used, for example 30*100.0. & None \\
psia & bars & atm & \\
Notes: & & & \\
\end{longtable}

\emph{Table 8.61: IPCG Keyword Description}

The capillary pressure for a grid block is scaled by:

\begin{equation}\phantomsection\label{eq-8-59}{
{P}_{c} = {P}_{{c}_{\mathit{TABLE}}}\left(\frac{\mathit{IPCG}}{{P}_{{c}_{\mathit{TABLE}-\mathit{MAX}}}}\right)
}\end{equation}

Where:

\(\mathit{Pc}\) = the resulting imbibition gas-oil capillary pressure
for a grid cell.

\(\mathit{IPCG}\) = the maximum capillary pressure from the IPCG array
for a given cell.

\({P}_{{c}_{\mathit{TABLE}}}\) = the capillary pressure in the
imbibition capillary pressure table

allocated to the grid block.

\({P}_{{c}_{\mathit{TABLE}-\mathit{MAX}}}\) = the maximum capillary
pressure in the imbibition capillary pressure table

allocated to the grid block at\({S}_{g}=1-{S}_{\mathit{wco}}\).

See also the PCG keyword for the equivalent drainage functionality.

\subsubsection{Example}\label{example-287}

\begin{verbatim}
--
--        DEFINE GRID BLOCK IPCG DATA FOR ALL CELLS (BASED ON NX x NY x NZ = 300)
--
IPCG
          100*50.0  100*75.0  100*125.0                                         /
\end{verbatim}

The above example defines the IPCG for 300 cells in the model as defined
by the DIMENS keyword in the RUNSPEC section.

\subsection{IPCW -- End-Point Scaling of Grid Cell Water Capillary
Pressure
(Imbibition)}\label{ipcw-end-point-scaling-of-grid-cell-water-capillary-pressure-imbibition}

\begin{longtable}[]{@{}
  >{\raggedright\arraybackslash}p{(\linewidth - 14\tabcolsep) * \real{0.1250}}
  >{\raggedright\arraybackslash}p{(\linewidth - 14\tabcolsep) * \real{0.1250}}
  >{\raggedright\arraybackslash}p{(\linewidth - 14\tabcolsep) * \real{0.1250}}
  >{\raggedright\arraybackslash}p{(\linewidth - 14\tabcolsep) * \real{0.1250}}
  >{\raggedright\arraybackslash}p{(\linewidth - 14\tabcolsep) * \real{0.1250}}
  >{\raggedright\arraybackslash}p{(\linewidth - 14\tabcolsep) * \real{0.1250}}
  >{\raggedright\arraybackslash}p{(\linewidth - 14\tabcolsep) * \real{0.1250}}
  >{\raggedright\arraybackslash}p{(\linewidth - 14\tabcolsep) * \real{0.1250}}@{}}
\toprule\noalign{}
\begin{minipage}[b]{\linewidth}\raggedright
RUNSPEC
\end{minipage} & \begin{minipage}[b]{\linewidth}\raggedright
GRID
\end{minipage} & \begin{minipage}[b]{\linewidth}\raggedright
EDIT
\end{minipage} & \begin{minipage}[b]{\linewidth}\raggedright
PROPS
\end{minipage} & \begin{minipage}[b]{\linewidth}\raggedright
REGIONS
\end{minipage} & \begin{minipage}[b]{\linewidth}\raggedright
SOLUTION
\end{minipage} & \begin{minipage}[b]{\linewidth}\raggedright
SUMMARY
\end{minipage} & \begin{minipage}[b]{\linewidth}\raggedright
SCHEDULE
\end{minipage} \\
\midrule\noalign{}
\endhead
\bottomrule\noalign{}
\endlastfoot
\end{longtable}

\subsubsection{Description}\label{description-494}

IPCW defines the maximum imbibition water-oil or water-gas capillary
pressure values for all the cells in the model via an array. The
ENDSCALE keyword in the RUNSPEC section should be activated to enable
end-point scaling and the use of this keyword. In addition, the HYSTER
option on the SATOPTS keyword in the RUNSPEC section has to be activated
to invoke the hysteresis option. The keyword can be used with all grid
types.

\begin{longtable}[]{@{}
  >{\raggedright\arraybackslash}p{(\linewidth - 6\tabcolsep) * \real{0.2500}}
  >{\raggedright\arraybackslash}p{(\linewidth - 6\tabcolsep) * \real{0.2500}}
  >{\raggedright\arraybackslash}p{(\linewidth - 6\tabcolsep) * \real{0.2500}}
  >{\raggedright\arraybackslash}p{(\linewidth - 6\tabcolsep) * \real{0.2500}}@{}}
\toprule\noalign{}
\begin{minipage}[b]{\linewidth}\raggedright
No.
\end{minipage} & \begin{minipage}[b]{\linewidth}\raggedright
Name
\end{minipage} & \begin{minipage}[b]{\linewidth}\raggedright
Description
\end{minipage} & \begin{minipage}[b]{\linewidth}\raggedright
Default
\end{minipage} \\
\midrule\noalign{}
\endhead
\bottomrule\noalign{}
\endlastfoot
Field & Metric & Laboratory & \\
1 & IPCW & IPCW is an array of positive real numbers assigning the
maximum imbibition water capillary pressure values for each cell in the
model. Repeat counts may be used, for example 30*100.0. & None \\
psia & bars & atm & \\
Notes: & & & \\
\end{longtable}

\emph{Table 8.62: IPCW Keyword Description}

The capillary pressure for a grid block is scaled by:

\begin{equation}\phantomsection\label{eq-8-60}{
{P}_{c} = {P}_{{c}_{\mathit{TABLE}}}\left(\frac{\mathit{IPCW}}{{P}_{{c}_{\mathit{TABLE}-\mathit{MAX}}}}\right)
}\end{equation}

Where:

\(\mathit{Pc}\) = the resulting imbibition water capillary pressure for
a grid cell.

\(\mathit{IPCW}\) = the maximum capillary pressure from the IPCW array
for a given cell.

\({P}_{{c}_{\mathit{TABLE}}}\) = the capillary pressure in the
imbibition capillary pressure table

allocated to the grid block.

\({P}_{{c}_{\mathit{TABLE}-\mathit{MAX}}}\) = the maximum capillary
pressure in the imbibition capillary pressure table

allocated to the grid block (that is at the connate water saturation).

See also the PCW keyword for the equivalent drainage functionality.

\subsubsection{Example}\label{example-288}

\begin{verbatim}
--
--        DEFINE GRID BLOCK IPCW DATA FOR ALL CELLS (BASED ON NX x NY x NZ = 300)
--
IPCW
          100*50.0  100*75.0   100*125.0                                        /
\end{verbatim}

The above example defines the IPCW for 300 cells in the model as defined
by the DIMENS keyword in the RUNSPEC section.

\subsection{ISGCR -- End-Point Scaling of Grid Cell Critical Gas
Saturation
(Imbibition)}\label{isgcr-end-point-scaling-of-grid-cell-critical-gas-saturation-imbibition}

\begin{longtable}[]{@{}
  >{\raggedright\arraybackslash}p{(\linewidth - 14\tabcolsep) * \real{0.1250}}
  >{\raggedright\arraybackslash}p{(\linewidth - 14\tabcolsep) * \real{0.1250}}
  >{\raggedright\arraybackslash}p{(\linewidth - 14\tabcolsep) * \real{0.1250}}
  >{\raggedright\arraybackslash}p{(\linewidth - 14\tabcolsep) * \real{0.1250}}
  >{\raggedright\arraybackslash}p{(\linewidth - 14\tabcolsep) * \real{0.1250}}
  >{\raggedright\arraybackslash}p{(\linewidth - 14\tabcolsep) * \real{0.1250}}
  >{\raggedright\arraybackslash}p{(\linewidth - 14\tabcolsep) * \real{0.1250}}
  >{\raggedright\arraybackslash}p{(\linewidth - 14\tabcolsep) * \real{0.1250}}@{}}
\toprule\noalign{}
\begin{minipage}[b]{\linewidth}\raggedright
RUNSPEC
\end{minipage} & \begin{minipage}[b]{\linewidth}\raggedright
GRID
\end{minipage} & \begin{minipage}[b]{\linewidth}\raggedright
EDIT
\end{minipage} & \begin{minipage}[b]{\linewidth}\raggedright
PROPS
\end{minipage} & \begin{minipage}[b]{\linewidth}\raggedright
REGIONS
\end{minipage} & \begin{minipage}[b]{\linewidth}\raggedright
SOLUTION
\end{minipage} & \begin{minipage}[b]{\linewidth}\raggedright
SUMMARY
\end{minipage} & \begin{minipage}[b]{\linewidth}\raggedright
SCHEDULE
\end{minipage} \\
\midrule\noalign{}
\endhead
\bottomrule\noalign{}
\endlastfoot
\end{longtable}

\subsubsection{Description}\label{description-495}

ISGCR defines the imbibition critical gas saturation for all the cells
in the model via an array when the end-point scaling option has been
invoked via the ENDSCALE keyword in the RUNSPEC section and the
hysteresis model option has been activated on the SATOPTS keyword in the
RUNSPEC section. The critical gas saturation is defined as the maximum
gas saturation for which the gas relative permeability is zero in a
two-phase relative permeability table.

The keyword can be used with all grid types.

\begin{longtable}[]{@{}
  >{\raggedright\arraybackslash}p{(\linewidth - 6\tabcolsep) * \real{0.2500}}
  >{\raggedright\arraybackslash}p{(\linewidth - 6\tabcolsep) * \real{0.2500}}
  >{\raggedright\arraybackslash}p{(\linewidth - 6\tabcolsep) * \real{0.2500}}
  >{\raggedright\arraybackslash}p{(\linewidth - 6\tabcolsep) * \real{0.2500}}@{}}
\toprule\noalign{}
\begin{minipage}[b]{\linewidth}\raggedright
No.
\end{minipage} & \begin{minipage}[b]{\linewidth}\raggedright
Name
\end{minipage} & \begin{minipage}[b]{\linewidth}\raggedright
Description
\end{minipage} & \begin{minipage}[b]{\linewidth}\raggedright
Default
\end{minipage} \\
\midrule\noalign{}
\endhead
\bottomrule\noalign{}
\endlastfoot
Field & Metric & Laboratory & \\
1 & ISGCR & ISGCR is an array of real numbers assigning the critical gas
saturation values to each cell in the model. The number of entries
should correspond to the NX x NY x NZ parameters on the DIMENS keyword.
Repeat counts may be used, for example 30*0.03 & Taken from cell
allocated relative permeability table. \\
dimensionless & dimensionless & dimensionless & \\
Notes: & & & \\
\end{longtable}

\emph{Table 8.63: ISGCR Keyword Description}

End-point scaling allows the entered relative permeability functions to
be re-scaled based on the saturation end-points allocated to each cell
by the ISWL, ISWCR, ISWU, ISGL, ISGCR, ISGU, ISOWCR, and ISOGCR
saturation grid arrays for the saturation end-points, and the IKRG,
IKRORG, IKRORW and IKRW relative permeability grid cell arrays for the
relative permeability end-point data. In addition end-point scaling may
be directional dependent in which case the directional dependent
versions of the aforementioned arrays should be used, that is ISGCRX,
ISGCRY and ISGCRZ instead of ISGCR. There is also the facility to make
the directional end-point scaling reversible or non-reversible and if
the non-reversible option is selected the non-reversible versions of the
aforementioned arrays should be used, that is ISGCRX, ISGCRX-, ISGCRY,
ISGCRY-, ISGCRZ and ISGCRZ-, instead of the ISGCR keyword.

\subsubsection{Example}\label{example-289}

\begin{verbatim}
--
--       DEFINE GRID BLOCK END-POINT ISGCR DATA FOR CELLS (NX x NY x NZ = 300)
ISGCR
         300*0.050                                                             /
\end{verbatim}

The above example defines a constant critical gas saturation of 0.05 to
all 300 cells in the model as defined by the DIMENS keyword in the
RUNSPEC section.

\subsection{ISGL -- End-Point Scaling of Grid Cell Connate Gas
Saturation
(Imbibition)}\label{isgl-end-point-scaling-of-grid-cell-connate-gas-saturation-imbibition}

\begin{longtable}[]{@{}
  >{\raggedright\arraybackslash}p{(\linewidth - 14\tabcolsep) * \real{0.1250}}
  >{\raggedright\arraybackslash}p{(\linewidth - 14\tabcolsep) * \real{0.1250}}
  >{\raggedright\arraybackslash}p{(\linewidth - 14\tabcolsep) * \real{0.1250}}
  >{\raggedright\arraybackslash}p{(\linewidth - 14\tabcolsep) * \real{0.1250}}
  >{\raggedright\arraybackslash}p{(\linewidth - 14\tabcolsep) * \real{0.1250}}
  >{\raggedright\arraybackslash}p{(\linewidth - 14\tabcolsep) * \real{0.1250}}
  >{\raggedright\arraybackslash}p{(\linewidth - 14\tabcolsep) * \real{0.1250}}
  >{\raggedright\arraybackslash}p{(\linewidth - 14\tabcolsep) * \real{0.1250}}@{}}
\toprule\noalign{}
\begin{minipage}[b]{\linewidth}\raggedright
RUNSPEC
\end{minipage} & \begin{minipage}[b]{\linewidth}\raggedright
GRID
\end{minipage} & \begin{minipage}[b]{\linewidth}\raggedright
EDIT
\end{minipage} & \begin{minipage}[b]{\linewidth}\raggedright
PROPS
\end{minipage} & \begin{minipage}[b]{\linewidth}\raggedright
REGIONS
\end{minipage} & \begin{minipage}[b]{\linewidth}\raggedright
SOLUTION
\end{minipage} & \begin{minipage}[b]{\linewidth}\raggedright
SUMMARY
\end{minipage} & \begin{minipage}[b]{\linewidth}\raggedright
SCHEDULE
\end{minipage} \\
\midrule\noalign{}
\endhead
\bottomrule\noalign{}
\endlastfoot
\end{longtable}

\subsubsection{Description}\label{description-496}

ISGL defines the imbibition connate gas saturation for all the cells in
the model via an array when the end-point scaling option has been
invoked via the ENDSCALE keyword in the RUNSPEC section and the
hysteresis model option has been activated on the SATOPTS keyword in the
RUNSPEC section. The connate gas saturation is defined as the minimum
gas saturation in a two-phase gas relative permeability table.

The keyword can be used with all grid types.

\begin{longtable}[]{@{}
  >{\raggedright\arraybackslash}p{(\linewidth - 6\tabcolsep) * \real{0.2500}}
  >{\raggedright\arraybackslash}p{(\linewidth - 6\tabcolsep) * \real{0.2500}}
  >{\raggedright\arraybackslash}p{(\linewidth - 6\tabcolsep) * \real{0.2500}}
  >{\raggedright\arraybackslash}p{(\linewidth - 6\tabcolsep) * \real{0.2500}}@{}}
\toprule\noalign{}
\begin{minipage}[b]{\linewidth}\raggedright
No.
\end{minipage} & \begin{minipage}[b]{\linewidth}\raggedright
Name
\end{minipage} & \begin{minipage}[b]{\linewidth}\raggedright
Description
\end{minipage} & \begin{minipage}[b]{\linewidth}\raggedright
Default
\end{minipage} \\
\midrule\noalign{}
\endhead
\bottomrule\noalign{}
\endlastfoot
Field & Metric & Laboratory & \\
1 & ISGL & ISGL is an array of real numbers assigning the connate gas
saturation values to each cell in the model. The number of entries
should correspond to the NX x NY x NZ parameters on the DIMENS keyword.
Repeat counts may be used, for example 30*0.03 dimensionless & Taken
from cell allocated relative permeability table. \\
dimensionless & dimensionless & dimensionless & \\
Notes: & & & \\
\end{longtable}

\emph{Table 8.64: ISGL Keyword Description}

End-point scaling allows the entered relative permeability functions to
be re-scaled based on the saturation end-points allocated to each cell
by the ISWL, ISWCR, ISWU, ISGL, ISGCR, ISGU, ISOWCR, and ISOGCR
saturation grid arrays for the saturation end-points, and the IKRG,
IKRORG, IKRORW and IKRW relative permeability grid cell arrays for the
relative permeability end-point data. In addition end-point scaling may
be directional dependent in which case the directional dependent
versions of the aforementioned arrays should be used, that is ISGLX,
ISGLY and ISGLZ instead of ISGL. There is also the facility to make the
directional end-point scaling reversible or non-reversible and if the
non-reversible option is selected the non-reversible versions of the
aforementioned arrays should be used, that is ISGLX, ISGLX-, ISGLY,
ISGLY-, ISGLZ and ISGLZ-, instead of the ISGL keyword.

\subsubsection{Example}\label{example-290}

\begin{verbatim}
--
--       DEFINE GRID BLOCK END-POINT ISGL DATA FOR ALL CELLS (NX x NY x NZ = 300)
ISGL
         300*0.030                                                              /
\end{verbatim}

The above example defines a constant connate gas saturation of 0.03 to
all 300 cells in the model as defined by the DIMENS keyword in the
RUNSPEC section.

\subsection{ISGLPC -- End-Point Scaling of Grid Cell Capillary Pressure
Connate Gas Saturation
(Imbibition)}\label{isglpc-end-point-scaling-of-grid-cell-capillary-pressure-connate-gas-saturation-imbibition}

\begin{longtable}[]{@{}
  >{\raggedright\arraybackslash}p{(\linewidth - 14\tabcolsep) * \real{0.1250}}
  >{\raggedright\arraybackslash}p{(\linewidth - 14\tabcolsep) * \real{0.1250}}
  >{\raggedright\arraybackslash}p{(\linewidth - 14\tabcolsep) * \real{0.1250}}
  >{\raggedright\arraybackslash}p{(\linewidth - 14\tabcolsep) * \real{0.1250}}
  >{\raggedright\arraybackslash}p{(\linewidth - 14\tabcolsep) * \real{0.1250}}
  >{\raggedright\arraybackslash}p{(\linewidth - 14\tabcolsep) * \real{0.1250}}
  >{\raggedright\arraybackslash}p{(\linewidth - 14\tabcolsep) * \real{0.1250}}
  >{\raggedright\arraybackslash}p{(\linewidth - 14\tabcolsep) * \real{0.1250}}@{}}
\toprule\noalign{}
\begin{minipage}[b]{\linewidth}\raggedright
RUNSPEC
\end{minipage} & \begin{minipage}[b]{\linewidth}\raggedright
GRID
\end{minipage} & \begin{minipage}[b]{\linewidth}\raggedright
EDIT
\end{minipage} & \begin{minipage}[b]{\linewidth}\raggedright
PROPS
\end{minipage} & \begin{minipage}[b]{\linewidth}\raggedright
REGIONS
\end{minipage} & \begin{minipage}[b]{\linewidth}\raggedright
SOLUTION
\end{minipage} & \begin{minipage}[b]{\linewidth}\raggedright
SUMMARY
\end{minipage} & \begin{minipage}[b]{\linewidth}\raggedright
SCHEDULE
\end{minipage} \\
\midrule\noalign{}
\endhead
\bottomrule\noalign{}
\endlastfoot
\end{longtable}

\subsubsection{Description}\label{description-497}

ISGLPC defines the imbibition connate gas saturation for all the cells
in the model via an array when the end-point scaling option has been
invoked via the ENDSCALE keyword in the RUNSPEC section and the
hysteresis model option has been activated on the SATOPTS keyword in the
RUNSPEC section. The connate gas saturation is defined as the minimum
gas saturation in a two-phase gas relative permeability table. The
keyword only applies the scaling to the imbibition capillary pressures
tables, unlike the ISGL keyword that applies the scaling to both the
capillary pressure and relative permeability tables. The keyword can be
used with all grid types.

\begin{longtable}[]{@{}
  >{\raggedright\arraybackslash}p{(\linewidth - 6\tabcolsep) * \real{0.2500}}
  >{\raggedright\arraybackslash}p{(\linewidth - 6\tabcolsep) * \real{0.2500}}
  >{\raggedright\arraybackslash}p{(\linewidth - 6\tabcolsep) * \real{0.2500}}
  >{\raggedright\arraybackslash}p{(\linewidth - 6\tabcolsep) * \real{0.2500}}@{}}
\toprule\noalign{}
\begin{minipage}[b]{\linewidth}\raggedright
No.
\end{minipage} & \begin{minipage}[b]{\linewidth}\raggedright
Name
\end{minipage} & \begin{minipage}[b]{\linewidth}\raggedright
Description
\end{minipage} & \begin{minipage}[b]{\linewidth}\raggedright
Default
\end{minipage} \\
\midrule\noalign{}
\endhead
\bottomrule\noalign{}
\endlastfoot
Field & Metric & Laboratory & \\
1 & ISGLPC & ISGLPC is an array of real numbers assigning the connate
gas saturation values to each cell in the model. The number of entries
should correspond to the NX x NY x NZ parameters on the DIMENS keyword.
If ISGLPC is omitted from the input deck the values will be defaulted to
those on the ISGL series of keywords. If the ISGL series of keywords are
missing from the input deck then the values are taken from the cell
allocated capillary pressure table. Repeat counts may be used, for
example 30*0.03 & Taken from SGL or from the cell allocated capillary
pressure table. \\
dimensionless & dimensionless & dimensionless & \\
Notes: & & & \\
\end{longtable}

\emph{Table 8.65: ISGLPC Keyword Description}

End-point scaling allows the entered relative permeability functions to
be re-scaled based on the saturation end-points allocated to each cell
by the ISWL, ISWCR, ISWU, ISGL, ISGCR, ISGU, ISOWCR, and ISOGCR
saturation grid arrays for the saturation end-points, and the IKRG,
IKRORG, IKRORW and IKRW relative permeability grid cell arrays for the
relative permeability end-point data. In addition end-point scaling may
be directional dependent in which case the directional dependent
versions of the aforementioned arrays should be used, that is ISGLX,
ISGLY and ISGLZ instead of ISGL or ISGLPC. There is also the facility to
make the directional end-point scaling reversible or non-reversible and
if the non-reversible option is selected the non-reversible versions of
the aforementioned arrays should be used, that is ISGLX, ISGLX-, ISGLY,
ISGLY-, ISGLZ and ISGLZ-, instead of the ISGL or ISGLPC keywords.

Missing Some Functionality - Use with Caution.

\subsubsection{Example}\label{example-291}

\begin{verbatim}
--
--       DEFINE GRID BLOCK END-POINT ISGLPC DATA FOR ALL CELLS
–       (NX x NY x NZ = 300)
ISGLPC
         300*0.030                                                           /
\end{verbatim}

The above example defines a constant connate gas saturation of 0.03 to
all 300 cells in the model as defined by the DIMENS keyword in the
RUNSPEC section.

\subsection{ISGU -- End-Point Scaling of Grid Cell Maximum Gas
Saturation
(Imbibition)}\label{isgu-end-point-scaling-of-grid-cell-maximum-gas-saturation-imbibition}

\begin{longtable}[]{@{}
  >{\raggedright\arraybackslash}p{(\linewidth - 14\tabcolsep) * \real{0.1250}}
  >{\raggedright\arraybackslash}p{(\linewidth - 14\tabcolsep) * \real{0.1250}}
  >{\raggedright\arraybackslash}p{(\linewidth - 14\tabcolsep) * \real{0.1250}}
  >{\raggedright\arraybackslash}p{(\linewidth - 14\tabcolsep) * \real{0.1250}}
  >{\raggedright\arraybackslash}p{(\linewidth - 14\tabcolsep) * \real{0.1250}}
  >{\raggedright\arraybackslash}p{(\linewidth - 14\tabcolsep) * \real{0.1250}}
  >{\raggedright\arraybackslash}p{(\linewidth - 14\tabcolsep) * \real{0.1250}}
  >{\raggedright\arraybackslash}p{(\linewidth - 14\tabcolsep) * \real{0.1250}}@{}}
\toprule\noalign{}
\begin{minipage}[b]{\linewidth}\raggedright
RUNSPEC
\end{minipage} & \begin{minipage}[b]{\linewidth}\raggedright
GRID
\end{minipage} & \begin{minipage}[b]{\linewidth}\raggedright
EDIT
\end{minipage} & \begin{minipage}[b]{\linewidth}\raggedright
PROPS
\end{minipage} & \begin{minipage}[b]{\linewidth}\raggedright
REGIONS
\end{minipage} & \begin{minipage}[b]{\linewidth}\raggedright
SOLUTION
\end{minipage} & \begin{minipage}[b]{\linewidth}\raggedright
SUMMARY
\end{minipage} & \begin{minipage}[b]{\linewidth}\raggedright
SCHEDULE
\end{minipage} \\
\midrule\noalign{}
\endhead
\bottomrule\noalign{}
\endlastfoot
\end{longtable}

\subsubsection{Description}\label{description-498}

ISGU defines the imbibition maximum gas saturation for all the cells in
the model via an array when the end-point scaling option has been
invoked via the ENDSCALE keyword in the RUNSPEC section and the
hysteresis model option has been activated on the SATOPTS keyword in the
RUNSPEC section. The maximum gas saturation is defined as the maximum
gas saturation in a two-phase gas relative permeability table. The
keyword can be used with all grid types.

\begin{longtable}[]{@{}
  >{\raggedright\arraybackslash}p{(\linewidth - 6\tabcolsep) * \real{0.2500}}
  >{\raggedright\arraybackslash}p{(\linewidth - 6\tabcolsep) * \real{0.2500}}
  >{\raggedright\arraybackslash}p{(\linewidth - 6\tabcolsep) * \real{0.2500}}
  >{\raggedright\arraybackslash}p{(\linewidth - 6\tabcolsep) * \real{0.2500}}@{}}
\toprule\noalign{}
\begin{minipage}[b]{\linewidth}\raggedright
No.
\end{minipage} & \begin{minipage}[b]{\linewidth}\raggedright
Name
\end{minipage} & \begin{minipage}[b]{\linewidth}\raggedright
Description
\end{minipage} & \begin{minipage}[b]{\linewidth}\raggedright
Default
\end{minipage} \\
\midrule\noalign{}
\endhead
\bottomrule\noalign{}
\endlastfoot
Field & Metric & Laboratory & \\
1 & ISGU & ISGU is an array of real numbers assigning the maximum gas
saturation values to each cell in the model. The number of entries
should correspond to the NX x NY x NZ parameters on the DIMENS keyword.
Repeat counts may be used, for example 30*0.70 & Taken from cell
allocated relative permeability table. \\
dimensionless & dimensionless & dimensionless & \\
Notes: & & & \\
\end{longtable}

\emph{Table 8.66: ISGU Keyword Description}

End-point scaling allows the entered relative permeability functions to
be re-scaled based on the saturation end-points allocated to each cell
by the ISWL, ISWCR, ISWU, ISGL, ISGCR, ISGU, ISOWCR, and ISOGCR
saturation grid arrays for the saturation end-points, and the IKRG,
IKRORG, IKRORW and IKRW relative permeability grid cell arrays for the
relative permeability end-point data. In addition end-point scaling may
be directional dependent in which case the directional dependent
versions of the aforementioned arrays should be used, that is ISGUX,
ISGUY and ISGUZ instead of ISGU. There is also the facility to make the
directional end-point scaling reversible or non-reversible and if the
non-reversible option is selected the non-reversible versions of the
aforementioned arrays should be used, that is ISGUX, ISGUX-, ISGUY,
ISGUY-, ISGUZ and ISGUZ-, instead of the ISGU keyword.

\subsubsection{Example}\label{example-292}

\begin{verbatim}
--
--       DEFINE GRID BLOCK END-POINT ISGU DATA FOR ALL CELLS (NX x NY x NZ = 300)
ISGU
         300*0.700                                                             /
\end{verbatim}

The above example defines a constant connate gas saturation of 0.70 to
all 300 cells in the model as defined by the DIMENS keyword in the
RUNSPEC section.

\subsection{ISOGCR -- End-Point Scaling of Grid Cell Critical Oil
Saturation to Gas
(Imbibition)}\label{isogcr-end-point-scaling-of-grid-cell-critical-oil-saturation-to-gas-imbibition}

\begin{longtable}[]{@{}
  >{\raggedright\arraybackslash}p{(\linewidth - 14\tabcolsep) * \real{0.1250}}
  >{\raggedright\arraybackslash}p{(\linewidth - 14\tabcolsep) * \real{0.1250}}
  >{\raggedright\arraybackslash}p{(\linewidth - 14\tabcolsep) * \real{0.1250}}
  >{\raggedright\arraybackslash}p{(\linewidth - 14\tabcolsep) * \real{0.1250}}
  >{\raggedright\arraybackslash}p{(\linewidth - 14\tabcolsep) * \real{0.1250}}
  >{\raggedright\arraybackslash}p{(\linewidth - 14\tabcolsep) * \real{0.1250}}
  >{\raggedright\arraybackslash}p{(\linewidth - 14\tabcolsep) * \real{0.1250}}
  >{\raggedright\arraybackslash}p{(\linewidth - 14\tabcolsep) * \real{0.1250}}@{}}
\toprule\noalign{}
\begin{minipage}[b]{\linewidth}\raggedright
RUNSPEC
\end{minipage} & \begin{minipage}[b]{\linewidth}\raggedright
GRID
\end{minipage} & \begin{minipage}[b]{\linewidth}\raggedright
EDIT
\end{minipage} & \begin{minipage}[b]{\linewidth}\raggedright
PROPS
\end{minipage} & \begin{minipage}[b]{\linewidth}\raggedright
REGIONS
\end{minipage} & \begin{minipage}[b]{\linewidth}\raggedright
SOLUTION
\end{minipage} & \begin{minipage}[b]{\linewidth}\raggedright
SUMMARY
\end{minipage} & \begin{minipage}[b]{\linewidth}\raggedright
SCHEDULE
\end{minipage} \\
\midrule\noalign{}
\endhead
\bottomrule\noalign{}
\endlastfoot
\end{longtable}

\subsubsection{Description}\label{description-499}

ISOGCR defines the imbibition critical oil saturation for all the cells
in the model via an array when the end-point scaling option has been
invoked via the ENDSCALE in the RUNSPEC section and the hysteresis model
option has been activated on the SATOPTS keyword in the RUNSPEC section.
The critical oil saturation with respect to gas is defined as the
maximum oil saturation for which the oil relative permeability is zero
in a two-phase gas-oil relative permeability table. The keyword can be
used with all grid types.

\begin{longtable}[]{@{}
  >{\raggedright\arraybackslash}p{(\linewidth - 6\tabcolsep) * \real{0.2500}}
  >{\raggedright\arraybackslash}p{(\linewidth - 6\tabcolsep) * \real{0.2500}}
  >{\raggedright\arraybackslash}p{(\linewidth - 6\tabcolsep) * \real{0.2500}}
  >{\raggedright\arraybackslash}p{(\linewidth - 6\tabcolsep) * \real{0.2500}}@{}}
\toprule\noalign{}
\begin{minipage}[b]{\linewidth}\raggedright
No.
\end{minipage} & \begin{minipage}[b]{\linewidth}\raggedright
Name
\end{minipage} & \begin{minipage}[b]{\linewidth}\raggedright
Description
\end{minipage} & \begin{minipage}[b]{\linewidth}\raggedright
Default
\end{minipage} \\
\midrule\noalign{}
\endhead
\bottomrule\noalign{}
\endlastfoot
Field & Metric & Laboratory & \\
1 & ISOGCR & ISOGCR is an array of real numbers assigning the critical
oil saturation with respect to gas values to each cell in the model. The
number of entries should correspond to the NX x NY x NZ parameters on
the DIMENS keyword. Repeat counts may be used, for example 30*0.30 &
Taken from cell allocated relative permeability table. \\
dimensionless & dimensionless & dimensionless & \\
Notes: & & & \\
\end{longtable}

\emph{Table 8.67: ISOGCR Keyword Description}

End-point scaling allows the entered relative permeability functions to
be re-scaled based on the saturation end-points allocated to each cell
by the ISWL, ISWCR, ISWU, ISGL, ISGCR, ISGU, ISOWCR, and ISOGCR
saturation grid arrays for the saturation end-points, and the IKRG,
IKRORG, IKRORW and IKRW\textgreater{} relative permeability grid cell
arrays for the relative permeability end-point data. In addition
end-point scaling may be directional dependent in which case the
directional dependent versions of the aforementioned arrays should be
used, that is ISOGCRX, ISOGCRY and ISOGCRZ instead of ISOGCR. There is
also the facility to make the directional end-point scaling reversible
or non-reversible and if the non-reversible option is selected the
non-reversible versions of the aforementioned arrays should be used,
that is ISOGCRX, ISOGCRX-, ISOGCRY, ISOGCRY-, ISOGCRZ and ISOGCRZ-,
instead of the ISOGCR keyword.

\subsubsection{Example}\label{example-293}

\begin{verbatim}
--
--       DEFINE GRID BLOCK END-POINT ISOGCR DATA FOR ALL CELLS
--       (NX x NY x NZ = 300)
--
ISOGCR
         300*0.200                                                              /
\end{verbatim}

The above example defines a constant critical gas saturation of 0.20 to
all 300 cells in the model as defined by the DIMENS keyword in the
RUNSPEC section.

\subsection{ISOWCR -- End-Point Scaling of Grid Cell Critical Oil
Saturation to Water
(Imbib.)}\label{isowcr-end-point-scaling-of-grid-cell-critical-oil-saturation-to-water-imbib.}

\begin{longtable}[]{@{}
  >{\raggedright\arraybackslash}p{(\linewidth - 14\tabcolsep) * \real{0.1250}}
  >{\raggedright\arraybackslash}p{(\linewidth - 14\tabcolsep) * \real{0.1250}}
  >{\raggedright\arraybackslash}p{(\linewidth - 14\tabcolsep) * \real{0.1250}}
  >{\raggedright\arraybackslash}p{(\linewidth - 14\tabcolsep) * \real{0.1250}}
  >{\raggedright\arraybackslash}p{(\linewidth - 14\tabcolsep) * \real{0.1250}}
  >{\raggedright\arraybackslash}p{(\linewidth - 14\tabcolsep) * \real{0.1250}}
  >{\raggedright\arraybackslash}p{(\linewidth - 14\tabcolsep) * \real{0.1250}}
  >{\raggedright\arraybackslash}p{(\linewidth - 14\tabcolsep) * \real{0.1250}}@{}}
\toprule\noalign{}
\begin{minipage}[b]{\linewidth}\raggedright
RUNSPEC
\end{minipage} & \begin{minipage}[b]{\linewidth}\raggedright
GRID
\end{minipage} & \begin{minipage}[b]{\linewidth}\raggedright
EDIT
\end{minipage} & \begin{minipage}[b]{\linewidth}\raggedright
PROPS
\end{minipage} & \begin{minipage}[b]{\linewidth}\raggedright
REGIONS
\end{minipage} & \begin{minipage}[b]{\linewidth}\raggedright
SOLUTION
\end{minipage} & \begin{minipage}[b]{\linewidth}\raggedright
SUMMARY
\end{minipage} & \begin{minipage}[b]{\linewidth}\raggedright
SCHEDULE
\end{minipage} \\
\midrule\noalign{}
\endhead
\bottomrule\noalign{}
\endlastfoot
\end{longtable}

\subsubsection{Description}\label{description-500}

ISOWCR defines the imbibition critical oil saturation for all the cells
in the model via an array when the end-point scaling option has been
invoked via the ENDSCALE in the RUNSPEC section and the hysteresis model
option has been activated on the SATOPTS keyword in the RUNSPEC section.
The critical oil saturation with respect to water is defined as the
maximum oil saturation for which the oil relative permeability is zero
in a two-phase oil-water relative permeability table. The keyword can be
used with all grid types.

\begin{longtable}[]{@{}
  >{\raggedright\arraybackslash}p{(\linewidth - 6\tabcolsep) * \real{0.2500}}
  >{\raggedright\arraybackslash}p{(\linewidth - 6\tabcolsep) * \real{0.2500}}
  >{\raggedright\arraybackslash}p{(\linewidth - 6\tabcolsep) * \real{0.2500}}
  >{\raggedright\arraybackslash}p{(\linewidth - 6\tabcolsep) * \real{0.2500}}@{}}
\toprule\noalign{}
\begin{minipage}[b]{\linewidth}\raggedright
No.
\end{minipage} & \begin{minipage}[b]{\linewidth}\raggedright
Name
\end{minipage} & \begin{minipage}[b]{\linewidth}\raggedright
Description
\end{minipage} & \begin{minipage}[b]{\linewidth}\raggedright
Default
\end{minipage} \\
\midrule\noalign{}
\endhead
\bottomrule\noalign{}
\endlastfoot
Field & Metric & Laboratory & \\
1 & ISOWCR & ISOWCR is an array of real numbers assigning the critical
oil saturation with respect to water values to each cell in the model.
The number of entries should correspond to the NX x NY x NZ parameters
on the DIMENS keyword. Repeat counts may be used, for example 30*0.30 &
Taken from cell allocated relative permeability table. \\
dimensionless & dimensionless & dimensionless & \\
Notes: & & & \\
\end{longtable}

\emph{Table 8.68: ISOWCR Keyword Description}

End-point scaling allows the entered relative permeability functions to
be re-scaled based on the saturation end-points allocated to each cell
by the ISWL, ISWCR, ISWU, ISGL, ISGCR, ISGU, ISOWCR, and ISOGCR
saturation grid arrays for the saturation end-points, and the IKRG,
IKRORG, IKRORW and IKRW relative permeability grid cell arrays for the
relative permeability end-point data. In addition end-point scaling may
be directional dependent in which case the directional dependent
versions of the aforementioned arrays should be used, that is ISOWCRX,
ISOWCRY and ISOWCRZ instead of ISOWCR. There is also the facility to
make the directional end-point scaling reversible or non-reversible and
if the non-reversible option is selected the non-reversible versions of
the aforementioned arrays should be used, that is ISOWCRX, ISOWCRX-,
ISOWCRY, ISOWCRY-, ISOWCRZ and ISOWCRZ-, instead of the ISOWCR keyword.

\subsubsection{Example}\label{example-294}

\begin{verbatim}
--
--       DEFINE GRID BLOCK END-POINT ISOWCR DATA FOR ALL CELLS
--       (NX x NY x NZ = 300)
--
ISOWCR
         300*0.200                                                              /
\end{verbatim}

The above example defines a constant critical gas saturation of 0.20 to
all 300 cells in the model as defined by the DIMENS keyword in the
RUNSPEC section

\subsection{ISWCR -- End-Point Scaling of Grid Cell Critical Water
Saturation
(Imbibition)}\label{iswcr-end-point-scaling-of-grid-cell-critical-water-saturation-imbibition}

\begin{longtable}[]{@{}
  >{\raggedright\arraybackslash}p{(\linewidth - 14\tabcolsep) * \real{0.1250}}
  >{\raggedright\arraybackslash}p{(\linewidth - 14\tabcolsep) * \real{0.1250}}
  >{\raggedright\arraybackslash}p{(\linewidth - 14\tabcolsep) * \real{0.1250}}
  >{\raggedright\arraybackslash}p{(\linewidth - 14\tabcolsep) * \real{0.1250}}
  >{\raggedright\arraybackslash}p{(\linewidth - 14\tabcolsep) * \real{0.1250}}
  >{\raggedright\arraybackslash}p{(\linewidth - 14\tabcolsep) * \real{0.1250}}
  >{\raggedright\arraybackslash}p{(\linewidth - 14\tabcolsep) * \real{0.1250}}
  >{\raggedright\arraybackslash}p{(\linewidth - 14\tabcolsep) * \real{0.1250}}@{}}
\toprule\noalign{}
\begin{minipage}[b]{\linewidth}\raggedright
RUNSPEC
\end{minipage} & \begin{minipage}[b]{\linewidth}\raggedright
GRID
\end{minipage} & \begin{minipage}[b]{\linewidth}\raggedright
EDIT
\end{minipage} & \begin{minipage}[b]{\linewidth}\raggedright
PROPS
\end{minipage} & \begin{minipage}[b]{\linewidth}\raggedright
REGIONS
\end{minipage} & \begin{minipage}[b]{\linewidth}\raggedright
SOLUTION
\end{minipage} & \begin{minipage}[b]{\linewidth}\raggedright
SUMMARY
\end{minipage} & \begin{minipage}[b]{\linewidth}\raggedright
SCHEDULE
\end{minipage} \\
\midrule\noalign{}
\endhead
\bottomrule\noalign{}
\endlastfoot
\end{longtable}

\subsubsection{Description}\label{description-501}

ISWCR defines the imbibition critical water saturation for all the cells
in the model via an array when the end-point scaling option has been
invoked via the ENDSCALE in the RUNSPEC section and the hysteresis model
option has been activated on the SATOPTS keyword in the RUNSPEC section.
The critical water saturation is defined as the maximum water saturation
for which the water relative permeability is zero in a two-phase
relative permeability table. The keyword can be used with all grid
types.

\begin{longtable}[]{@{}
  >{\raggedright\arraybackslash}p{(\linewidth - 6\tabcolsep) * \real{0.2500}}
  >{\raggedright\arraybackslash}p{(\linewidth - 6\tabcolsep) * \real{0.2500}}
  >{\raggedright\arraybackslash}p{(\linewidth - 6\tabcolsep) * \real{0.2500}}
  >{\raggedright\arraybackslash}p{(\linewidth - 6\tabcolsep) * \real{0.2500}}@{}}
\toprule\noalign{}
\begin{minipage}[b]{\linewidth}\raggedright
No.
\end{minipage} & \begin{minipage}[b]{\linewidth}\raggedright
Name
\end{minipage} & \begin{minipage}[b]{\linewidth}\raggedright
Description
\end{minipage} & \begin{minipage}[b]{\linewidth}\raggedright
Default
\end{minipage} \\
\midrule\noalign{}
\endhead
\bottomrule\noalign{}
\endlastfoot
Field & Metric & Laboratory & \\
1 & ISWCR & ISWCR is an array of real numbers assigning the critical
water saturation values to each cell in the model. The number of entries
should correspond to the NX x NY x NZ parameters on the DIMENS keyword.
Repeat counts may be used, for example 30*0.20 & Taken from cell
allocated relative permeability table. \\
dimensionless & dimensionless & dimensionless & \\
Notes: & & & \\
\end{longtable}

\emph{Table 8.69: ISWCR Keyword Description}

End-point scaling allows the entered relative permeability functions to
be re-scaled based on the saturation end-points allocated to each cell
by the ISWL, ISWCR, ISWU, ISGL, ISGCR, ISGU, ISOWCR, and ISOGCR
saturation grid arrays for the saturation end-points, and the IKRG,
IKRORG, IKRORW and IKRW relative permeability grid cell arrays for the
relative permeability end-point data. In addition end-point scaling may
be directional dependent in which case the directional dependent
versions of the aforementioned arrays should be used, that is ISWCRX,
ISWCRY and ISWCRZ instead of ISWCR. There is also the facility to make
the directional end-point scaling reversible or non-reversible and if
the non-reversible option is selected the non-reversible versions of the
aforementioned arrays should be used, that is ISWCRX, ISWCRX-, ISWCRY,
ISWCRY-, ISWCRZ and ISWCRZ-, instead of the ISWCR keyword.

\subsubsection{Example}\label{example-295}

\begin{verbatim}
--
--       DEFINE GRID BLOCK END-POINT ISWCR DATA FOR ALL CELLS
--       (NX x NY x NZ = 300)
--
ISWCR
  300*0.200                                                                     /
\end{verbatim}

The above example defines a constant critical water saturation of 0.20
to all 300 cells in the model as defined by the DIMENS keyword in the
RUNSPEC section.

\subsection{ISWL -- End-Point Scaling of Grid Cell Connate Water
Saturation
(Imbibition)}\label{iswl-end-point-scaling-of-grid-cell-connate-water-saturation-imbibition}

\begin{longtable}[]{@{}
  >{\raggedright\arraybackslash}p{(\linewidth - 14\tabcolsep) * \real{0.1250}}
  >{\raggedright\arraybackslash}p{(\linewidth - 14\tabcolsep) * \real{0.1250}}
  >{\raggedright\arraybackslash}p{(\linewidth - 14\tabcolsep) * \real{0.1250}}
  >{\raggedright\arraybackslash}p{(\linewidth - 14\tabcolsep) * \real{0.1250}}
  >{\raggedright\arraybackslash}p{(\linewidth - 14\tabcolsep) * \real{0.1250}}
  >{\raggedright\arraybackslash}p{(\linewidth - 14\tabcolsep) * \real{0.1250}}
  >{\raggedright\arraybackslash}p{(\linewidth - 14\tabcolsep) * \real{0.1250}}
  >{\raggedright\arraybackslash}p{(\linewidth - 14\tabcolsep) * \real{0.1250}}@{}}
\toprule\noalign{}
\begin{minipage}[b]{\linewidth}\raggedright
RUNSPEC
\end{minipage} & \begin{minipage}[b]{\linewidth}\raggedright
GRID
\end{minipage} & \begin{minipage}[b]{\linewidth}\raggedright
EDIT
\end{minipage} & \begin{minipage}[b]{\linewidth}\raggedright
PROPS
\end{minipage} & \begin{minipage}[b]{\linewidth}\raggedright
REGIONS
\end{minipage} & \begin{minipage}[b]{\linewidth}\raggedright
SOLUTION
\end{minipage} & \begin{minipage}[b]{\linewidth}\raggedright
SUMMARY
\end{minipage} & \begin{minipage}[b]{\linewidth}\raggedright
SCHEDULE
\end{minipage} \\
\midrule\noalign{}
\endhead
\bottomrule\noalign{}
\endlastfoot
\end{longtable}

\subsubsection{Description}\label{description-502}

ISWL defines the imbibition connate water saturation for all the cells
in the model via an array when the end-point scaling option has been
invoked via the ENDSCALE in the RUNSPEC section and the hysteresis model
option has been activated on the SATOPTS keyword in the RUNSPEC section.
The connate water saturation is defined as the minimum water saturation
in a two-phase water relative permeability table.

The keyword can be used with all grid types.

\begin{longtable}[]{@{}
  >{\raggedright\arraybackslash}p{(\linewidth - 6\tabcolsep) * \real{0.2500}}
  >{\raggedright\arraybackslash}p{(\linewidth - 6\tabcolsep) * \real{0.2500}}
  >{\raggedright\arraybackslash}p{(\linewidth - 6\tabcolsep) * \real{0.2500}}
  >{\raggedright\arraybackslash}p{(\linewidth - 6\tabcolsep) * \real{0.2500}}@{}}
\toprule\noalign{}
\begin{minipage}[b]{\linewidth}\raggedright
No.
\end{minipage} & \begin{minipage}[b]{\linewidth}\raggedright
Name
\end{minipage} & \begin{minipage}[b]{\linewidth}\raggedright
Description
\end{minipage} & \begin{minipage}[b]{\linewidth}\raggedright
Default
\end{minipage} \\
\midrule\noalign{}
\endhead
\bottomrule\noalign{}
\endlastfoot
Field & Metric & Laboratory & \\
1 & ISWL & ISWL is an array of real numbers assigning the connate water
saturation values to each cell in the model. The number of entries
should correspond to the NX x NY x NZ parameters on the DIMENS keyword.
Repeat counts may be used, for example 30*0.15 & Taken from cell
allocated relative permeability table. \\
dimensionless & dimensionless & dimensionless & \\
Notes: & & & \\
\end{longtable}

\emph{Table 8.70: ISWL Keyword Description}

End-point scaling allows the entered relative permeability functions to
be re-scaled based on the saturation end-points allocated to each cell
by the ISWL, ISWCR, ISWU, ISGL, ISGCR, ISGU, ISOWCR, and ISOGCR
saturation grid arrays for the saturation end-points, and the IKRG,
IKRORG, IKRORW and IKRW relative permeability grid cell arrays for the
relative permeability end-point data. In addition end-point scaling may
be directional dependent in which case the directional dependent
versions of the aforementioned arrays should be used, that is ISWLX,
ISWLY and ISWLZ instead of ISWL. There is also the facility to make the
directional end-point scaling reversible or non-reversible and if the
non-reversible option is selected the non-reversible versions of the
aforementioned arrays should be used, that is ISWLX, ISWLX-, ISWLY,
ISWLY-, ISWLZ and ISWLZ-, instead of the ISWL keyword.

\subsubsection{Example}\label{example-296}

\begin{verbatim}
--
--       DEFINE GRID BLOCK END-POINT ISWL DATA FOR ALL CELLS (NX x NY x NZ = 300)
--
ISWL
         300*0.150                                                              /
\end{verbatim}

The above example defines a constant connate water saturation of 0.15 to
all 300 cells in the model as defined by the DIMENS keyword in the
RUNSPEC section.

\subsection{ISWLPC -- End-Point Scaling Grid Cell Capillary Pressure
Connate Water Saturations
(Imbibition)}\label{iswlpc-end-point-scaling-grid-cell-capillary-pressure-connate-water-saturations-imbibition}

\begin{longtable}[]{@{}
  >{\raggedright\arraybackslash}p{(\linewidth - 14\tabcolsep) * \real{0.1250}}
  >{\raggedright\arraybackslash}p{(\linewidth - 14\tabcolsep) * \real{0.1250}}
  >{\raggedright\arraybackslash}p{(\linewidth - 14\tabcolsep) * \real{0.1250}}
  >{\raggedright\arraybackslash}p{(\linewidth - 14\tabcolsep) * \real{0.1250}}
  >{\raggedright\arraybackslash}p{(\linewidth - 14\tabcolsep) * \real{0.1250}}
  >{\raggedright\arraybackslash}p{(\linewidth - 14\tabcolsep) * \real{0.1250}}
  >{\raggedright\arraybackslash}p{(\linewidth - 14\tabcolsep) * \real{0.1250}}
  >{\raggedright\arraybackslash}p{(\linewidth - 14\tabcolsep) * \real{0.1250}}@{}}
\toprule\noalign{}
\begin{minipage}[b]{\linewidth}\raggedright
RUNSPEC
\end{minipage} & \begin{minipage}[b]{\linewidth}\raggedright
GRID
\end{minipage} & \begin{minipage}[b]{\linewidth}\raggedright
EDIT
\end{minipage} & \begin{minipage}[b]{\linewidth}\raggedright
PROPS
\end{minipage} & \begin{minipage}[b]{\linewidth}\raggedright
REGIONS
\end{minipage} & \begin{minipage}[b]{\linewidth}\raggedright
SOLUTION
\end{minipage} & \begin{minipage}[b]{\linewidth}\raggedright
SUMMARY
\end{minipage} & \begin{minipage}[b]{\linewidth}\raggedright
SCHEDULE
\end{minipage} \\
\midrule\noalign{}
\endhead
\bottomrule\noalign{}
\endlastfoot
\end{longtable}

\subsubsection{Description}\label{description-503}

ISWLPC defines the imbibition connate water saturation for all the cells
in the model via an array when the end-point scaling option has been
invoked via the ENDSCALE in the RUNSPEC section and the hysteresis model
option has been activated on the SATOPTS keyword in the RUNSPEC section.
The connate water saturation is defined as the minimum water saturation
in a two-phase water relative permeability table. The keyword only
applies the scaling to the imbibition capillary pressures tables, unlike
the ISWL keyword that applies the scaling to both the capillary pressure
and relative permeability tables. The keyword can be used with all grid
types.

\begin{longtable}[]{@{}
  >{\raggedright\arraybackslash}p{(\linewidth - 6\tabcolsep) * \real{0.2500}}
  >{\raggedright\arraybackslash}p{(\linewidth - 6\tabcolsep) * \real{0.2500}}
  >{\raggedright\arraybackslash}p{(\linewidth - 6\tabcolsep) * \real{0.2500}}
  >{\raggedright\arraybackslash}p{(\linewidth - 6\tabcolsep) * \real{0.2500}}@{}}
\toprule\noalign{}
\begin{minipage}[b]{\linewidth}\raggedright
No.
\end{minipage} & \begin{minipage}[b]{\linewidth}\raggedright
Name
\end{minipage} & \begin{minipage}[b]{\linewidth}\raggedright
Description
\end{minipage} & \begin{minipage}[b]{\linewidth}\raggedright
Default
\end{minipage} \\
\midrule\noalign{}
\endhead
\bottomrule\noalign{}
\endlastfoot
Field & Metric & Laboratory & \\
1 & ISWLPC & ISWLPC is an array of real numbers assigning the connate
water saturation values to each cell in the model. The number of entries
should correspond to the NX x NY x NZ parameters on the DIMENS keyword.
If ISWLPC is omitted from the input deck the values will be defaulted to
those on the ISGL series of keywords. If the ISWL series of keywords are
missing from the input deck then the values are taken from the cell
allocated capillary pressure table. Repeat counts may be used, for
example 30*0.15 & Taken from SWL or from the cell allocated capillary
pressure table. \\
dimensionless & dimensionless & dimensionless & \\
Notes: & & & \\
\end{longtable}

\emph{Table 8.71: ISWLPC Keyword Description}

End-point scaling allows the entered relative permeability functions to
be re-scaled based on the saturation end-points allocated to each cell
by the ISWL, ISWCR, ISWU, ISGL, ISGCR, ISGU, ISOWCR, and ISOGCR
saturation grid arrays for the saturation end-points, and the IKRG,
IKRORG, IKRORW and IKRW relative permeability grid cell arrays for the
relative permeability end-point data. In addition end-point scaling may
be directional dependent in which case the directional dependent
versions of the aforementioned arrays should be used, that is ISWLX,
ISWLY and ISWLZ instead of ISWL or ISWLPC. There is also the facility to
make the directional end-point scaling reversible or non-reversible and
if the non-reversible option is selected the non-reversible versions of
the aforementioned arrays should be used, that is ISWLX, ISWLX-, ISWLY,
ISWLY-, ISWLZ and ISWLZ-, instead of the ISWL or ISWLPC keywords.

Missing Some Functionality - Use with Caution.

\subsubsection{Example}\label{example-297}

\begin{verbatim}
--
--      DEFINE GRID BLOCK END-POINT ISWLPC DATA FOR ALL CELLS (
--      NX x NY x NZ = 300)
--
ISWLPC
         300*0.150                                                             /
\end{verbatim}

The above example defines a constant connate water saturation of 0.15 to
all 300 cells in the model as defined by the DIMENS keyword in the
RUNSPEC section.

\subsection{ISWU -- End-Point Scaling of Grid Cell Maximum Water
Saturation
(Imbibition)}\label{iswu-end-point-scaling-of-grid-cell-maximum-water-saturation-imbibition}

\begin{longtable}[]{@{}
  >{\raggedright\arraybackslash}p{(\linewidth - 14\tabcolsep) * \real{0.1250}}
  >{\raggedright\arraybackslash}p{(\linewidth - 14\tabcolsep) * \real{0.1250}}
  >{\raggedright\arraybackslash}p{(\linewidth - 14\tabcolsep) * \real{0.1250}}
  >{\raggedright\arraybackslash}p{(\linewidth - 14\tabcolsep) * \real{0.1250}}
  >{\raggedright\arraybackslash}p{(\linewidth - 14\tabcolsep) * \real{0.1250}}
  >{\raggedright\arraybackslash}p{(\linewidth - 14\tabcolsep) * \real{0.1250}}
  >{\raggedright\arraybackslash}p{(\linewidth - 14\tabcolsep) * \real{0.1250}}
  >{\raggedright\arraybackslash}p{(\linewidth - 14\tabcolsep) * \real{0.1250}}@{}}
\toprule\noalign{}
\begin{minipage}[b]{\linewidth}\raggedright
RUNSPEC
\end{minipage} & \begin{minipage}[b]{\linewidth}\raggedright
GRID
\end{minipage} & \begin{minipage}[b]{\linewidth}\raggedright
EDIT
\end{minipage} & \begin{minipage}[b]{\linewidth}\raggedright
PROPS
\end{minipage} & \begin{minipage}[b]{\linewidth}\raggedright
REGIONS
\end{minipage} & \begin{minipage}[b]{\linewidth}\raggedright
SOLUTION
\end{minipage} & \begin{minipage}[b]{\linewidth}\raggedright
SUMMARY
\end{minipage} & \begin{minipage}[b]{\linewidth}\raggedright
SCHEDULE
\end{minipage} \\
\midrule\noalign{}
\endhead
\bottomrule\noalign{}
\endlastfoot
\end{longtable}

\subsubsection{Description}\label{description-504}

ISWU defines the imbibition maximum water saturation for all the cells
in the model via an array when the end-point scaling option has been
invoked via the ENDSCALE in the RUNSPEC section and the hysteresis model
option has been activated on the SATOPTS keyword in the RUNSPEC section.
The maximum water saturation is defined as the maximum water saturation
in a two-phase water relative permeability table. The keyword can be
used with all grid types.

\begin{longtable}[]{@{}
  >{\raggedright\arraybackslash}p{(\linewidth - 6\tabcolsep) * \real{0.2500}}
  >{\raggedright\arraybackslash}p{(\linewidth - 6\tabcolsep) * \real{0.2500}}
  >{\raggedright\arraybackslash}p{(\linewidth - 6\tabcolsep) * \real{0.2500}}
  >{\raggedright\arraybackslash}p{(\linewidth - 6\tabcolsep) * \real{0.2500}}@{}}
\toprule\noalign{}
\begin{minipage}[b]{\linewidth}\raggedright
No.
\end{minipage} & \begin{minipage}[b]{\linewidth}\raggedright
Name
\end{minipage} & \begin{minipage}[b]{\linewidth}\raggedright
Description
\end{minipage} & \begin{minipage}[b]{\linewidth}\raggedright
Default
\end{minipage} \\
\midrule\noalign{}
\endhead
\bottomrule\noalign{}
\endlastfoot
Field & Metric & Laboratory & \\
1 & ISWU & ISWU is an array of real numbers assigning the maximum water
saturation values to each cell in the model. The number of entries
should correspond to the NX x NY x NZ parameters on the DIMENS keyword.
Repeat counts may be used, for example 30*0.70 & Taken from cell
allocated relative permeability table. \\
dimensionless & dimensionless & dimensionless & \\
Notes: & & & \\
\end{longtable}

\emph{Table 8.72: ISWU Keyword Description}

End-point scaling allows the entered relative permeability functions to
be re-scaled based on the saturation end-points allocated to each cell
by the ISWL, ISWCR, ISWU, ISGL, ISGCR, ISGU, ISOWCR, and ISOGCR
saturation grid arrays for the saturation end-points, and the IKRG,
IKRORG, IKRORW and IKRW relative permeability grid cell arrays for the
relative permeability end-point data. In addition end-point scaling may
be directional dependent in which case the directional dependent
versions of the aforementioned arrays should be used, that is ISWUX,
ISWUY and ISWUZ instead of ISWU. There is also the facility to make the
directional end-point scaling reversible or non-reversible and if the
non-reversible option is selected the non-reversible versions of the
aforementioned arrays should be used, that is ISWUX, ISWUX-, ISWUY,
ISWUY-, ISWUZ and ISWUZ-, instead of the ISWU keyword.

\subsubsection{Example}\label{example-298}

\begin{verbatim}
--
--       DEFINE GRID BLOCK END-POINT ISWU DATA FOR ALL CELLS (NX x NY x NZ = 300)
--
ISWU
         300*0.700                                                             /
\end{verbatim}

The above example defines a constant connate gas saturation of 0.70 to
all 300 cells in the model as defined by the DIMENS keyword in the
RUNSPEC section.

\subsection{KRG -- End-Point Scaling of Grid Cell Krg(Sgu)
(Drainage)}\label{krg-end-point-scaling-of-grid-cell-krgsgu-drainage}

\begin{longtable}[]{@{}
  >{\raggedright\arraybackslash}p{(\linewidth - 14\tabcolsep) * \real{0.1250}}
  >{\raggedright\arraybackslash}p{(\linewidth - 14\tabcolsep) * \real{0.1250}}
  >{\raggedright\arraybackslash}p{(\linewidth - 14\tabcolsep) * \real{0.1250}}
  >{\raggedright\arraybackslash}p{(\linewidth - 14\tabcolsep) * \real{0.1250}}
  >{\raggedright\arraybackslash}p{(\linewidth - 14\tabcolsep) * \real{0.1250}}
  >{\raggedright\arraybackslash}p{(\linewidth - 14\tabcolsep) * \real{0.1250}}
  >{\raggedright\arraybackslash}p{(\linewidth - 14\tabcolsep) * \real{0.1250}}
  >{\raggedright\arraybackslash}p{(\linewidth - 14\tabcolsep) * \real{0.1250}}@{}}
\toprule\noalign{}
\begin{minipage}[b]{\linewidth}\raggedright
RUNSPEC
\end{minipage} & \begin{minipage}[b]{\linewidth}\raggedright
GRID
\end{minipage} & \begin{minipage}[b]{\linewidth}\raggedright
EDIT
\end{minipage} & \begin{minipage}[b]{\linewidth}\raggedright
PROPS
\end{minipage} & \begin{minipage}[b]{\linewidth}\raggedright
REGIONS
\end{minipage} & \begin{minipage}[b]{\linewidth}\raggedright
SOLUTION
\end{minipage} & \begin{minipage}[b]{\linewidth}\raggedright
SUMMARY
\end{minipage} & \begin{minipage}[b]{\linewidth}\raggedright
SCHEDULE
\end{minipage} \\
\midrule\noalign{}
\endhead
\bottomrule\noalign{}
\endlastfoot
\end{longtable}

\subsubsection{Description}\label{description-505}

KRG defines the scaling parameter at the maximum drainage gas relative
permeability value (SGU), normally SGU is equal to 1.0 - Swc, for all
the cells in the model via an array. The ENDSCALE keyword in the RUNSPEC
section should be activated to enable end-point scaling and the use of
this keyword. The SCALCERS keyword in the PROPS section defines the
options used in the re-scaling process, the options are two point
scaling and three point scaling.

\begin{longtable}[]{@{}
  >{\raggedright\arraybackslash}p{(\linewidth - 6\tabcolsep) * \real{0.2500}}
  >{\raggedright\arraybackslash}p{(\linewidth - 6\tabcolsep) * \real{0.2500}}
  >{\raggedright\arraybackslash}p{(\linewidth - 6\tabcolsep) * \real{0.2500}}
  >{\raggedright\arraybackslash}p{(\linewidth - 6\tabcolsep) * \real{0.2500}}@{}}
\toprule\noalign{}
\begin{minipage}[b]{\linewidth}\raggedright
No.
\end{minipage} & \begin{minipage}[b]{\linewidth}\raggedright
Name
\end{minipage} & \begin{minipage}[b]{\linewidth}\raggedright
Description
\end{minipage} & \begin{minipage}[b]{\linewidth}\raggedright
Default
\end{minipage} \\
\midrule\noalign{}
\endhead
\bottomrule\noalign{}
\endlastfoot
Field & Metric & Laboratory & \\
1 & KRG & KRG is an array of positive real numbers which are greater
than zero and less than or equal to 1.0, that are the assigned scaling
KRG values for each cell in the model. Repeat counts may be used, for
example 50*0.400. dimensionless & Taken from cell allocated relative
permeability table. \\
dimensionless & dimensionless & dimensionless & \\
Notes: & & & \\
\end{longtable}

\emph{Table 8.73: KRG Keyword Description}

\begin{verbatim}
\end{verbatim}

For the two point scaling option and for the KRGR gas relative
permeability array NOT present in the input deck the krg value for a
grid block is scaled by:

\begin{equation}\phantomsection\label{eq-8-61}{
{k}_{\mathit{rg}} = {k}_{{\mathit{rg}}_{ \mathit{TABLE}}}\left(\frac{\mathit{KRG}}{{k}_{{\mathit{rg}}_{ \mathit{TABLE}-\mathit{MAX}}}}\right)
}\end{equation}

Where:

\(\mathit{krg}\) = the resulting krg value for a grid cell.

\(\mathit{KRG}\) = the scaling gas relative permeability value from the
KRG array for a given

cell.

\({k}_{{\mathit{rg}}_{ \mathit{TABLE}}}\) = the gas relative
permeability from a grid block's gas-oil table at the grid

blocks gas saturation.

\({k}_{{\mathit{rg}}_{ \mathit{TABLE}-\mathit{MAX}}}\) = the maximum gas
relative permeability from a grid block's gas-oil table, that

is at the connate water saturation (Swc).

If the KRGR keyword is present in the input deck then the scaling
matches the relative permeability at the critical saturation of the
displacing phase.

If three point scaling option has been selected via the SCALECRS keyword
in the PROPS section the critical displacing phase is defined as:

\begin{longtable}[]{@{}lll@{}}
\toprule\noalign{}
No & Phases Present & Critical Saturation \\
\midrule\noalign{}
\endhead
\bottomrule\noalign{}
\endlastfoot
1 & Gas-Oil & S critical = 1.0 -- SOGCR - SWL \\
2 & Gas-Oil-Water & S critical = 1.0 -- SOGCR - SWL \\
3 & Gas-Water & S critical = 1.0 -- SWCR \\
\end{longtable}

\emph{Table 8.74: Critical Displacement Relationships}

End-point scaling allows the entered relative permeability functions to
be re-scaled based on the saturation end-points allocated to each cell
by the SWL, SWCR, SWU, SGL, SGCR, SGU, SOWCR, and SOGCR saturation grid
arrays for the saturation end-points, In addition end-point scaling may
be directional dependent in which case the directional dependent
versions of the aforementioned arrays should be used, that is SWUX, SWUY
and SWUZ instead of SWU, There is also the facility to make the
directional end-point scaling reversible or non-reversible and if the
non-reversible option is selected the non-reversible versions of the
aforementioned arrays should be used, that is SWUX, SWUX-, SWUY, SWUY-,
SWUZ and SWUZ-, instead of the SWU keyword.

End-point scaling also allows the entered relative permeability
functions to be scale on the relative permeability values using the KRG,
KRGR, KRO, KRORG, KRORW, KRW and KRWR relative permeability grid cell
arrays for the relative permeability end-point data. In addition
end-point scaling may be directional dependent in which case the
directional dependent versions of the aforementioned arrays should be
used, that is KRGX, KRGY and KRGZ instead of KRG, There is also the
facility to make the directional end-point scaling reversible or
non-reversible and if the non-reversible option is selected the
non-reversible versions of the aforementioned arrays should be used,
that is KRGX, KRGX-, KRGY, KRGY-, KRGZ and KRGZ-, instead of the KRG
keyword.

If the hysteresis model option has been activated on the SATOPTS keyword
in the RUNSPEC section, then the equivalent imbibition arrays prefixed
with the letter I, for example IKRG, can be used to scale KRG for the
relative permeability imbibition tables.

\subsubsection{Examples}\label{examples-29}

The first example defines an input box for the whole grid and for layers
one to three, for layer one KRG is set equal to 0.555, for layer two KRG
equals 0.575, and for layer three KRG equals 0.600.

\begin{verbatim}
--
--       DEFINE INPUT BOX FOR EDITING INPUT ARRAYS (NX=100, NY=100)
--
--       ---------- BOX ---------
--       I1  I2   J1  J2   K1  K2
BOX
         1*  1*   1*  1*   1   3                           / DEFINE BOX AREA
--
--       SET KRG VALUES FOR THREE LAYERS IN THE MODEL
--
KRG
         10000*0.555  10000*0.575  10000*0.600                /
--
--       DEFINE END OF INPUT BOX EDITING OF INPUT ARRAYS
--
ENDBOX

\end{verbatim}

The next example does exactly the same thing using the EQUALS keyword
instead.

\begin{verbatim}
--
--       ARRAY       CONSTANT     ---------- BOX ---------
--                               I1  I2   J1  J2   K1  K2
EQUALS
         KRG         0.5550       1*  1*   1*  1*   1   1  / KRG FOR LAYER 1
         KRG         0.5750       1*  1*   1*  1*   2   2  / KRG FOR LAYER 2
         KRG         0.6000       1*  1*   1*  1*   3   3  / KRG FOR LAYER 3
/
\end{verbatim}

\subsection{KRGR -- End-Point Scaling of Grid Cell Krgr(1-Sogcr)
(Drainage)}\label{krgr-end-point-scaling-of-grid-cell-krgr1-sogcr-drainage}

\begin{longtable}[]{@{}
  >{\raggedright\arraybackslash}p{(\linewidth - 14\tabcolsep) * \real{0.1250}}
  >{\raggedright\arraybackslash}p{(\linewidth - 14\tabcolsep) * \real{0.1250}}
  >{\raggedright\arraybackslash}p{(\linewidth - 14\tabcolsep) * \real{0.1250}}
  >{\raggedright\arraybackslash}p{(\linewidth - 14\tabcolsep) * \real{0.1250}}
  >{\raggedright\arraybackslash}p{(\linewidth - 14\tabcolsep) * \real{0.1250}}
  >{\raggedright\arraybackslash}p{(\linewidth - 14\tabcolsep) * \real{0.1250}}
  >{\raggedright\arraybackslash}p{(\linewidth - 14\tabcolsep) * \real{0.1250}}
  >{\raggedright\arraybackslash}p{(\linewidth - 14\tabcolsep) * \real{0.1250}}@{}}
\toprule\noalign{}
\begin{minipage}[b]{\linewidth}\raggedright
RUNSPEC
\end{minipage} & \begin{minipage}[b]{\linewidth}\raggedright
GRID
\end{minipage} & \begin{minipage}[b]{\linewidth}\raggedright
EDIT
\end{minipage} & \begin{minipage}[b]{\linewidth}\raggedright
PROPS
\end{minipage} & \begin{minipage}[b]{\linewidth}\raggedright
REGIONS
\end{minipage} & \begin{minipage}[b]{\linewidth}\raggedright
SOLUTION
\end{minipage} & \begin{minipage}[b]{\linewidth}\raggedright
SUMMARY
\end{minipage} & \begin{minipage}[b]{\linewidth}\raggedright
SCHEDULE
\end{minipage} \\
\midrule\noalign{}
\endhead
\bottomrule\noalign{}
\endlastfoot
\end{longtable}

\subsubsection{Description}\label{description-506}

KRGR defines the scaling parameter at the relative permeability of gas
at residual oil saturation (1 -- SOGCR), or critical water saturation in
a gas-water run (Swc), for all the cells in the model via an array. The
ENDSCALE keyword in the RUNSPEC section should be activated to enable
end-point scaling and the use of this keyword. The SCALCERS keyword in
the PROPS section defines the options used in the re-scaling process,
the options are two point scaling and three point scaling.

\begin{longtable}[]{@{}
  >{\raggedright\arraybackslash}p{(\linewidth - 6\tabcolsep) * \real{0.2500}}
  >{\raggedright\arraybackslash}p{(\linewidth - 6\tabcolsep) * \real{0.2500}}
  >{\raggedright\arraybackslash}p{(\linewidth - 6\tabcolsep) * \real{0.2500}}
  >{\raggedright\arraybackslash}p{(\linewidth - 6\tabcolsep) * \real{0.2500}}@{}}
\toprule\noalign{}
\begin{minipage}[b]{\linewidth}\raggedright
No.
\end{minipage} & \begin{minipage}[b]{\linewidth}\raggedright
Name
\end{minipage} & \begin{minipage}[b]{\linewidth}\raggedright
Description
\end{minipage} & \begin{minipage}[b]{\linewidth}\raggedright
Default
\end{minipage} \\
\midrule\noalign{}
\endhead
\bottomrule\noalign{}
\endlastfoot
Field & Metric & Laboratory & \\
1 & KRGR & KRGR is an array of positive real numbers which are greater
than zero and less than or equal to 1.0, that are the assigned scaling
KRGR values for each cell in the model. In addition, for a given grid
block KGRG should be less than KRG. Repeat counts may be used, for
example 50*0.400. & Taken from cell allocated relative permeability
table. \\
dimensionless & dimensionless & dimensionless & \\
Notes: & & & \\
\end{longtable}

\emph{Table 8.75: KRGR Keyword Description}

\begin{verbatim}
\end{verbatim}

When the KRGR keyword is present in the input deck then the scaling
matches the relative permeability at the critical saturation of the
displacing phase (oil or water).

If three point scaling option has been selected via the SCALECRS keyword
in the PROPS section the critical displacing phase is defined as:

\begin{longtable}[]{@{}lll@{}}
\toprule\noalign{}
No & Phases Present & Critical Saturation \\
\midrule\noalign{}
\endhead
\bottomrule\noalign{}
\endlastfoot
1 & Gas-Oil & S critical = 1.0 -- SOGCR - SWL \\
2 & Gas-Oil-Water & S critical = 1.0 -- SOGCR - SWL \\
3 & Gas-Water & S critical = 1.0 -- SWCR \\
\end{longtable}

\emph{Table 8.76: Critical Displacement Relationships}

End-point scaling allows the entered relative permeability functions to
be re-scaled based on the saturation end-points allocated to each cell
by the SWL, SWCR, SWU, SGL, SGCR, SGU, SOWCR, and SOGCR saturation grid
arrays for the saturation end-points, In addition end-point scaling may
be directional dependent in which case the directional dependent
versions of the aforementioned arrays should be used, that is SWUX, SWUY
and SWUZ instead of SWU, There is also the facility to make the
directional end-point scaling reversible or non-reversible and if the
non-reversible option is selected the non-reversible versions of the
aforementioned arrays should be used, that is SWUX, SWUX-, SWUY, SWUY-,
SWUZ and SWUZ-, instead of the SWU keyword.

End-point scaling also allows the entered relative permeability
functions to be scale on the relative permeability values using the KRG,
KRGR, KRO, KRORG, KRORW, KRW and KRWR relative permeability grid cell
arrays for the relative permeability end-point data. In addition
end-point scaling may be directional dependent in which case the
directional dependent versions of the aforementioned arrays should be
used, that is KRGRX, KRGRY and KRGRZ instead of KRGR, There is also the
facility to make the directional end-point scaling reversible or
non-reversible and if the non-reversible option is selected the
non-reversible versions of the aforementioned arrays should be used,
that is KRGRX, KRGRX-, KRGRY, KRGRY-, KRGRZ and KRGRZ-, instead of the
KRGR keyword.

If the hysteresis model option has been activated on the SATOPTS keyword
in the RUNSPEC section, then the equivalent imbibition arrays suffixed
with the letter I, for example IKRGR, can be used to define the KRGR for
the relative permeability imbibition tables.

\subsubsection{Examples}\label{examples-30}

The first example defines an input box for the whole grid and for layers
one to three, for layer one KRGR is set equal to 0.500, for layer two
KRGR equals 0.570, and for layer three KRGR equals 0.580.

\begin{verbatim}
--
--       DEFINE INPUT BOX FOR EDITING INPUT ARRAYS (NX=100, NY=100)
--
--       ---------- BOX ---------
--       I1  I2   J1  J2   K1  K2
BOX
         1*  1*   1*  1*   1   3                           / DEFINE BOX AREA
--
--       SET KRGR VALUES FOR THREE LAYERS IN THE MODEL
--
KRGR
         10000*0.500  10000*0.570  10000*0.580                /
--
--       DEFINE END OF INPUT BOX EDITING OF INPUT ARRAYS
--
ENDBOX
\end{verbatim}

The next example does exactly the same thing using the EQUALS keyword
instead.

\begin{verbatim}
--
--       ARRAY       CONSTANT     ---------- BOX ---------
--                                I1  I2   J1  J2   K1  K2
EQUALS
         KRGR        0.5000       1*  1*   1*  1*    1   1  / KRGR FOR LAYER 1
         KRGR        0.5700       1*  1*   1*  1*    2   2  / KRGR FOR LAYER 2
         KRGR        0.5800       1*  1*   1*  1*    3   3  / KRGR FOR LAYER 3
/
\end{verbatim}

\subsection{KRO -- End-Point Scaling of Grid Cell Kro(Swl)
(Drainage)}\label{kro-end-point-scaling-of-grid-cell-kroswl-drainage}

\begin{longtable}[]{@{}
  >{\raggedright\arraybackslash}p{(\linewidth - 14\tabcolsep) * \real{0.1250}}
  >{\raggedright\arraybackslash}p{(\linewidth - 14\tabcolsep) * \real{0.1250}}
  >{\raggedright\arraybackslash}p{(\linewidth - 14\tabcolsep) * \real{0.1250}}
  >{\raggedright\arraybackslash}p{(\linewidth - 14\tabcolsep) * \real{0.1250}}
  >{\raggedright\arraybackslash}p{(\linewidth - 14\tabcolsep) * \real{0.1250}}
  >{\raggedright\arraybackslash}p{(\linewidth - 14\tabcolsep) * \real{0.1250}}
  >{\raggedright\arraybackslash}p{(\linewidth - 14\tabcolsep) * \real{0.1250}}
  >{\raggedright\arraybackslash}p{(\linewidth - 14\tabcolsep) * \real{0.1250}}@{}}
\toprule\noalign{}
\begin{minipage}[b]{\linewidth}\raggedright
RUNSPEC
\end{minipage} & \begin{minipage}[b]{\linewidth}\raggedright
GRID
\end{minipage} & \begin{minipage}[b]{\linewidth}\raggedright
EDIT
\end{minipage} & \begin{minipage}[b]{\linewidth}\raggedright
PROPS
\end{minipage} & \begin{minipage}[b]{\linewidth}\raggedright
REGIONS
\end{minipage} & \begin{minipage}[b]{\linewidth}\raggedright
SOLUTION
\end{minipage} & \begin{minipage}[b]{\linewidth}\raggedright
SUMMARY
\end{minipage} & \begin{minipage}[b]{\linewidth}\raggedright
SCHEDULE
\end{minipage} \\
\midrule\noalign{}
\endhead
\bottomrule\noalign{}
\endlastfoot
\end{longtable}

\subsubsection{Description}\label{description-507}

KRO defines the scaling parameter for the drainage oil relative
permeability value at the connate water saturation (SWL), for all the
cells in the model via an array. The ENDSCALE keyword in the RUNSPEC
section should be activated to enable end-point scaling and the use of
this keyword. The SCALCERS keyword in the PROPS section defines the
options used in the re-scaling process, the options are two point
scaling and three point scaling.

\begin{longtable}[]{@{}
  >{\raggedright\arraybackslash}p{(\linewidth - 6\tabcolsep) * \real{0.2500}}
  >{\raggedright\arraybackslash}p{(\linewidth - 6\tabcolsep) * \real{0.2500}}
  >{\raggedright\arraybackslash}p{(\linewidth - 6\tabcolsep) * \real{0.2500}}
  >{\raggedright\arraybackslash}p{(\linewidth - 6\tabcolsep) * \real{0.2500}}@{}}
\toprule\noalign{}
\begin{minipage}[b]{\linewidth}\raggedright
No.
\end{minipage} & \begin{minipage}[b]{\linewidth}\raggedright
Name
\end{minipage} & \begin{minipage}[b]{\linewidth}\raggedright
Description
\end{minipage} & \begin{minipage}[b]{\linewidth}\raggedright
Default
\end{minipage} \\
\midrule\noalign{}
\endhead
\bottomrule\noalign{}
\endlastfoot
Field & Metric & Laboratory & \\
1 & KRO & KRO is an array of positive real numbers which are greater
than zero and less than or equal to 1.0, that are the assigned scaling
KRO values for each cell in the model. Repeat counts may be used, for
example 50*0.500. & Taken from cell allocated relative permeability
table. \\
dimensionless & dimensionless & dimensionless & \\
Notes: & & & \\
\end{longtable}

\emph{Table 8.77: KRO Keyword Description}

For the two point scaling option and for the KRORW or KRORG oil relative
permeability arrays NOT being present in the input deck the kro value
for a grid block is scaled by:

\begin{equation}\phantomsection\label{eq-8-62}{
{k}_{\mathit{ro}} = {k}_{{\mathit{ro}}_{ \mathit{TABLE}}}\left(\frac{\mathit{KRO}}{{k}_{{\mathit{ro}}_{ \mathit{TABLE}-\mathit{MAX}}}}\right)
}\end{equation}

Where:

\(\mathit{kro}\) = the resulting kro value for a grid cell.

\(\mathit{KRO}\) = the scaling oil relative permeability value from the
KRO array for a given

cell.

\({k}_{{\mathit{ro}}_{ \mathit{TABLE}}}\) = the oil relative
permeability from a grid block's oil relative permeability

table at the grid blocks oil saturation.

\({k}_{{\mathit{ro}}_{ \mathit{TABLE}-\mathit{MAX}}}\) = the maximum oil
relative permeability from a grid block's oil relative table,

that is at the critical water saturation (Swcr).

If the KRORW or KRORG keywords are present in the input deck then the
scaling matches the relative permeability at the critical saturation of
the displacing phase.

If three point scaling option has been selected via the SCALECRS keyword
in the PROPS section the critical displacing phase is defined as:

\begin{longtable}[]{@{}lll@{}}
\toprule\noalign{}
No & Keywords Present & Critical Saturation \\
\midrule\noalign{}
\endhead
\bottomrule\noalign{}
\endlastfoot
1 & KRORW & S critical = 1.0 -- SWCR - SGL \\
2 & KRORG & S critical = 1.0 -- SGCR - SWL \\
\end{longtable}

\emph{Table 8.78: Critical Displacement Relationships}

End-point scaling allows the entered relative permeability functions to
be re-scaled based on the saturation end-points allocated to each cell
by the SWL, SWCR, SWU, SGL, SGCR, SGU, SOWCR, and SOGCR saturation grid
arrays for the saturation end-points, In addition end-point scaling may
be directional dependent in which case the directional dependent
versions of the aforementioned arrays should be used, that is SWUX, SWUY
and SWUZ instead of SWU, There is also the facility to make the
directional end-point scaling reversible or non-reversible and if the
non-reversible option is selected the non-reversible versions of the
aforementioned arrays should be used, that is SWUX, SWUX-, SWUY, SWUY-,
SWUZ and SWUZ-, instead of the SWU keyword.

End-point scaling also allows the entered relative permeability
functions to be scale on the relative permeability values using the KRG,
KRGR, KRO, KRORG, KRORW, KRW and KRWR relative permeability grid cell
arrays for the relative permeability end-point data. In addition
end-point scaling may be directional dependent in which case the
directional dependent versions of the aforementioned arrays should be
used, that is KROX, KROY and KROZ instead of KRO, There is also the
facility to make the directional end-point scaling reversible or
non-reversible and if the non-reversible option is selected the
non-reversible versions of the aforementioned arrays should be used,
that is KROX, KROX-, KROY, KROY-, KROZ and KROZ-, instead of the KRO
keyword.

If the hysteresis model option has been activated on the SATOPTS keyword
in the RUNSPEC section, then the equivalent imbibition arrays suffixed
with the letter I, for example IKRO, can be used to define the KRO for
the relative permeability imbibition tables.

\subsubsection{Examples}\label{examples-31}

The first example defines an input box for the whole grid and for layers
one to three, for layer one KRO is set equal to 0.855, for layer two KRO
equals 0.875, and for layer three KRO equals 0.900.

\begin{verbatim}
--
--       DEFINE INPUT BOX FOR EDITING INPUT ARRAYS (NX=100, NY=100)
--
--       ---------- BOX ---------
--       I1  I2   J1  J2   K1  K2
BOX
         1*  1*   1*  1*   1   3                           / DEFINE BOX AREA
--
--       SET KRO VALUES FOR THREE LAYERS IN THE MODEL
--
KRO
         10000*0.855  10000*0.875  10000*0.900                /
--
--       DEFINE END OF INPUT BOX EDITING OF INPUT ARRAYS
--
ENDBOX
\end{verbatim}

The next example does exactly the same thing using the EQUALS keyword
instead.

\begin{verbatim}
--
--       ARRAY       CONSTANT     ---------- BOX ---------
--                                I1  I2   J1  J2   K1  K2
EQUALS
         KRO         0.8550       1*  1*   1*  1*    1   1  / KRO FOR LAYER 1
         KOG         0.8750       1*  1*   1*  1*    2   2  / KRO FOR LAYER 2
         KRO         0.9000       1*  1*   1*  1*    3   3  / KRO FOR LAYER 3
/
\end{verbatim}

\subsection{KRORG -- End-Point Scaling of Grid Cell Kro(Sgcr)
(Drainage)}\label{krorg-end-point-scaling-of-grid-cell-krosgcr-drainage}

\begin{longtable}[]{@{}
  >{\raggedright\arraybackslash}p{(\linewidth - 14\tabcolsep) * \real{0.1250}}
  >{\raggedright\arraybackslash}p{(\linewidth - 14\tabcolsep) * \real{0.1250}}
  >{\raggedright\arraybackslash}p{(\linewidth - 14\tabcolsep) * \real{0.1250}}
  >{\raggedright\arraybackslash}p{(\linewidth - 14\tabcolsep) * \real{0.1250}}
  >{\raggedright\arraybackslash}p{(\linewidth - 14\tabcolsep) * \real{0.1250}}
  >{\raggedright\arraybackslash}p{(\linewidth - 14\tabcolsep) * \real{0.1250}}
  >{\raggedright\arraybackslash}p{(\linewidth - 14\tabcolsep) * \real{0.1250}}
  >{\raggedright\arraybackslash}p{(\linewidth - 14\tabcolsep) * \real{0.1250}}@{}}
\toprule\noalign{}
\begin{minipage}[b]{\linewidth}\raggedright
RUNSPEC
\end{minipage} & \begin{minipage}[b]{\linewidth}\raggedright
GRID
\end{minipage} & \begin{minipage}[b]{\linewidth}\raggedright
EDIT
\end{minipage} & \begin{minipage}[b]{\linewidth}\raggedright
PROPS
\end{minipage} & \begin{minipage}[b]{\linewidth}\raggedright
REGIONS
\end{minipage} & \begin{minipage}[b]{\linewidth}\raggedright
SOLUTION
\end{minipage} & \begin{minipage}[b]{\linewidth}\raggedright
SUMMARY
\end{minipage} & \begin{minipage}[b]{\linewidth}\raggedright
SCHEDULE
\end{minipage} \\
\midrule\noalign{}
\endhead
\bottomrule\noalign{}
\endlastfoot
\end{longtable}

\subsubsection{Description}\label{description-508}

KRORG defines the scaling parameter for the drainage relative
permeability of oil at the critical gas saturation (SGCR), for all the
cells in the model via an array. The ENDSCALE keyword in the RUNSPEC
section should be activated to enable end-point scaling and the use of
this keyword. The SCALCERS keyword in the PROPS section defines the
options used in the re-scaling process, the options are two point
scaling and three point scaling.

\begin{longtable}[]{@{}
  >{\raggedright\arraybackslash}p{(\linewidth - 6\tabcolsep) * \real{0.2500}}
  >{\raggedright\arraybackslash}p{(\linewidth - 6\tabcolsep) * \real{0.2500}}
  >{\raggedright\arraybackslash}p{(\linewidth - 6\tabcolsep) * \real{0.2500}}
  >{\raggedright\arraybackslash}p{(\linewidth - 6\tabcolsep) * \real{0.2500}}@{}}
\toprule\noalign{}
\begin{minipage}[b]{\linewidth}\raggedright
No.
\end{minipage} & \begin{minipage}[b]{\linewidth}\raggedright
Name
\end{minipage} & \begin{minipage}[b]{\linewidth}\raggedright
Description
\end{minipage} & \begin{minipage}[b]{\linewidth}\raggedright
Default
\end{minipage} \\
\midrule\noalign{}
\endhead
\bottomrule\noalign{}
\endlastfoot
Field & Metric & Laboratory & \\
1 & KRORG & KRORG is an array of positive real numbers which are greater
than zero and less than or equal to 1.0, that are the assigned scaling
KRORG values for each cell in the model. Repeat counts may be used, for
example 50*0.850. & Taken from cell allocated relative permeability
table. \\
dimensionless & dimensionless & dimensionless & \\
Notes: & & & \\
\end{longtable}

\emph{Table 8.79: KRORG Keyword Description}

When the KRORG keyword is present in the input deck then the scaling
matches the relative permeability at the critical saturation of the
displacing phase.

If three point scaling option has been selected via the SCALECRS keyword
in the PROPS section the critical displacing phase is defined as:

\begin{longtable}[]{@{}lll@{}}
\toprule\noalign{}
No & Keywords Present & Critical Saturation \\
\midrule\noalign{}
\endhead
\bottomrule\noalign{}
\endlastfoot
1 & KRORW & S critical = 1.0 -- SWCR - SGL \\
2 & KRORG & S critical = 1.0 -- SGCR - SWL \\
\end{longtable}

\emph{Table 8.80: Critical Displacement Relationships}

End-point scaling allows the entered relative permeability functions to
be re-scaled based on the saturation end-points allocated to each cell
by the SWL, SWCR, SWU, SGL, SGCR, SGU, SOWCR, and SOGCR saturation grid
arrays for the saturation end-points, In addition end-point scaling may
be directional dependent in which case the directional dependent
versions of the aforementioned arrays should be used, that is SWUX, SWUY
and SWUZ instead of SWU, There is also the facility to make the
directional end-point scaling reversible or non-reversible and if the
non-reversible option is selected the non-reversible versions of the
aforementioned arrays should be used, that is SWUX, SWUX-, SWUY, SWUY-,
SWUZ and SWUZ-, instead of the SWU keyword.

End-point scaling also allows the entered relative permeability
functions to be scale on the relative permeability values using the KRG,
KRGR, KRORG, KRORGRG, KRORGRW, KRW and KRWR relative permeability grid
cell arrays for the relative permeability end-point data. In addition
end-point scaling may be directional dependent in which case the
directional dependent versions of the aforementioned arrays should be
used, that is KRORGX, KRORGY and KRORGZ instead of KRORG, There is also
the facility to make the directional end-point scaling reversible or
non-reversible and if the non-reversible option is selected the
non-reversible versions of the aforementioned arrays should be used,
that is KRORGX, KRORGX-, KRORGY, KRORGY-, KRORGZ and KRORGZ-, instead of
the KRORG keyword.

If the hysteresis model option has been activated on the SATOPTS keyword
in the RUNSPEC section, then the equivalent imbibition arrays suffixed
with the letter I, for example IKRORG, can be used to define the KRORG
for the relative permeability imbibition tables.

\subsubsection{Examples}\label{examples-32}

The first example defines an input box for the whole grid and for layers
one to three, for layer one KRORG is set equal to 0.755, for layer two
KRORG equals 0.775, and for layer three KRORG equals 0.800.

\begin{verbatim}
--
--       DEFINE INPUT BOX FOR EDITING INPUT ARRAYS (NX=100, NY=100)
--
--       ---------- BOX ---------
--       I1  I2   J1  J2   K1  K2
BOX
         1*  1*   1*  1*   1   3                           / DEFINE BOX AREA
--
--       SET KRORG VALUES FOR THREE LAYERS IN THE MODEL
--
KRORG
         10000*0.755  10000*0.775  10000*0.800                /
--
--       DEFINE END OF INPUT BOX EDITING OF INPUT ARRAYS
--
ENDBOX
\end{verbatim}

The next example does exactly the same thing using the EQUALS keyword
instead.

\begin{verbatim}
--
--       ARRAY       CONSTANT     ---------- BOX ---------
--                                I1  I2   J1  J2   K1  K2
EQUALS
         KRORG       0.7550       1*  1*   1*  1*    1   1  / KRORG FOR LAYER 1
         KRORG       0.7750       1*  1*   1*  1*    2   2  / KRORG FOR LAYER 2
         KRORG       0.8000       1*  1*   1*  1*    3   3  / KRORG FOR LAYER 3
/
\end{verbatim}

\subsection{KRORW -- End-Point Scaling of Grid Cell Kro(Swcr)
(Drainage)}\label{krorw-end-point-scaling-of-grid-cell-kroswcr-drainage}

\begin{longtable}[]{@{}
  >{\raggedright\arraybackslash}p{(\linewidth - 14\tabcolsep) * \real{0.1250}}
  >{\raggedright\arraybackslash}p{(\linewidth - 14\tabcolsep) * \real{0.1250}}
  >{\raggedright\arraybackslash}p{(\linewidth - 14\tabcolsep) * \real{0.1250}}
  >{\raggedright\arraybackslash}p{(\linewidth - 14\tabcolsep) * \real{0.1250}}
  >{\raggedright\arraybackslash}p{(\linewidth - 14\tabcolsep) * \real{0.1250}}
  >{\raggedright\arraybackslash}p{(\linewidth - 14\tabcolsep) * \real{0.1250}}
  >{\raggedright\arraybackslash}p{(\linewidth - 14\tabcolsep) * \real{0.1250}}
  >{\raggedright\arraybackslash}p{(\linewidth - 14\tabcolsep) * \real{0.1250}}@{}}
\toprule\noalign{}
\begin{minipage}[b]{\linewidth}\raggedright
RUNSPEC
\end{minipage} & \begin{minipage}[b]{\linewidth}\raggedright
GRID
\end{minipage} & \begin{minipage}[b]{\linewidth}\raggedright
EDIT
\end{minipage} & \begin{minipage}[b]{\linewidth}\raggedright
PROPS
\end{minipage} & \begin{minipage}[b]{\linewidth}\raggedright
REGIONS
\end{minipage} & \begin{minipage}[b]{\linewidth}\raggedright
SOLUTION
\end{minipage} & \begin{minipage}[b]{\linewidth}\raggedright
SUMMARY
\end{minipage} & \begin{minipage}[b]{\linewidth}\raggedright
SCHEDULE
\end{minipage} \\
\midrule\noalign{}
\endhead
\bottomrule\noalign{}
\endlastfoot
\end{longtable}

\subsubsection{Description}\label{description-509}

KRORW defines the scaling parameter for the drainage relative
permeability of oil at the critical water saturation (SWCR), for all the
cells in the model via an array. The ENDSCALE keyword in the RUNSPEC
section should be activated to enable end-point scaling and the use of
this keyword. The SCALECRS keyword in the PROPS section defines the
options used in the re-scaling process, the options are two point
scaling and three point scaling.

\begin{longtable}[]{@{}
  >{\raggedright\arraybackslash}p{(\linewidth - 6\tabcolsep) * \real{0.2500}}
  >{\raggedright\arraybackslash}p{(\linewidth - 6\tabcolsep) * \real{0.2500}}
  >{\raggedright\arraybackslash}p{(\linewidth - 6\tabcolsep) * \real{0.2500}}
  >{\raggedright\arraybackslash}p{(\linewidth - 6\tabcolsep) * \real{0.2500}}@{}}
\toprule\noalign{}
\begin{minipage}[b]{\linewidth}\raggedright
No.
\end{minipage} & \begin{minipage}[b]{\linewidth}\raggedright
Name
\end{minipage} & \begin{minipage}[b]{\linewidth}\raggedright
Description
\end{minipage} & \begin{minipage}[b]{\linewidth}\raggedright
Default
\end{minipage} \\
\midrule\noalign{}
\endhead
\bottomrule\noalign{}
\endlastfoot
Field & Metric & Laboratory & \\
1 & KRORW & KRORW is an array of positive real numbers which are greater
than zero and less than or equal to 1.0, that are the assigned scaling
KRORW values for each cell in the model. Repeat counts may be used, for
example 50*0.850 & Taken from cell allocated relative permeability
table. \\
dimensionless & dimensionless & dimensionless & \\
Notes: & & & \\
\end{longtable}

\emph{Table 8.81: KRORW Keyword Description}

When the KRORW keyword is present in the input deck then the scaling
matches the relative permeability at the critical saturation of the
displacing phase.

If three point scaling option has been selected via the SCALECRS keyword
in the PROPS section the critical displacing phase is defined as:

\begin{longtable}[]{@{}lll@{}}
\toprule\noalign{}
No & Keywords Present & Critical Saturation \\
\midrule\noalign{}
\endhead
\bottomrule\noalign{}
\endlastfoot
1 & KRORW & S critical = 1.0 -- SWCR - SGL \\
2 & KRORG & S critical = 1.0 -- SGCR - SWL \\
\end{longtable}

\emph{Table 8.82: Critical Displacement Relationships}

End-point scaling allows the entered relative permeability functions to
be re-scaled based on the saturation end-points allocated to each cell
by the SWL, SWCR, SWU, SGL, SGCR, SGU, SOWCR, and SOGCR saturation grid
arrays for the saturation end-points, In addition end-point scaling may
be directional dependent in which case the directional dependent
versions of the aforementioned arrays should be used, that is SWUX, SWUY
and SWUZ instead of SWU, There is also the facility to make the
directional end-point scaling reversible or non-reversible and if the
non-reversible option is selected the non-reversible versions of the
aforementioned arrays should be used, that is SWUX, SWUX-, SWUY, SWUY-,
SWUZ and SWUZ-, instead of the SWU keyword.

End-point scaling also allows the entered relative permeability
functions to be scale on the relative permeability values using the KRG,
KRGR, KRORW, KRORWRG, KRORWRW, KRW and KRWR relative permeability grid
cell arrays for the relative permeability end-point data. In addition
end-point scaling may be directional dependent in which case the
directional dependent versions of the aforementioned arrays should be
used, that is KRORWX, KRORWY and KRORWZ instead of KRORW, There is also
the facility to make the directional end-point scaling reversible or
non-reversible and if the non-reversible option is selected the
non-reversible versions of the aforementioned arrays should be used,
that is KRORWX, KRORWX-, KRORWY, KRORWY-, KRORWZ and KRORWZ-, instead of
the KRORW keyword.

If the hysteresis model option has been activated on the SATOPTS keyword
in the RUNSPEC section, then the equivalent imbibition arrays suffixed
with the letter I, for example IKRORW, can be used to define the KRORW
for the relative permeability imbibition tables.

\subsubsection{Examples}\label{examples-33}

The first example defines an input box for the whole grid and for layers
one to three, for layer one KRORW is set equal to 0.755, for layer two
KRORW equals 0.775, and for layer three KRORW equals 0.800.

\begin{verbatim}
--
--       DEFINE INPUT BOX FOR EDITING INPUT ARRAYS (NX=100, NY=100)
--
--       ---------- BOX ---------
--       I1  I2   J1  J2   K1  K2
BOX
         1*  1*   1*  1*   1   3                           / DEFINE BOX AREA
--
--       SET KRORW VALUES FOR THREE LAYERS IN THE MODEL
--
KRORW
         10000*0.755  10000*0.775  10000*0.800               /
--
--       DEFINE END OF INPUT BOX EDITING OF INPUT ARRAYS
--
ENDBOX
\end{verbatim}

The next example does exactly the same thing using the EQUALS keyword
instead.

\begin{verbatim}
--
--       ARRAY       CONSTANT     ---------- BOX ---------
--                                I1  I2   J1  J2   K1  K2
EQUALS
         KRORW       0.7550       1*  1*   1*  1*    1   1  / KRORW FOR LAYER 1
         KRORW       0.7750       1*  1*   1*  1*    2   2  / KRORW FOR LAYER 2
         KRORW       0.8000       1*  1*   1*  1*    3   3  / KRORW FOR LAYER 3
/
\end{verbatim}

\subsection{KRW -- End-Point Scaling of Grid Cell Krw(Sw =1.0)
(Drainage)}\label{krw-end-point-scaling-of-grid-cell-krwsw-1.0-drainage}

\begin{longtable}[]{@{}
  >{\raggedright\arraybackslash}p{(\linewidth - 14\tabcolsep) * \real{0.1250}}
  >{\raggedright\arraybackslash}p{(\linewidth - 14\tabcolsep) * \real{0.1250}}
  >{\raggedright\arraybackslash}p{(\linewidth - 14\tabcolsep) * \real{0.1250}}
  >{\raggedright\arraybackslash}p{(\linewidth - 14\tabcolsep) * \real{0.1250}}
  >{\raggedright\arraybackslash}p{(\linewidth - 14\tabcolsep) * \real{0.1250}}
  >{\raggedright\arraybackslash}p{(\linewidth - 14\tabcolsep) * \real{0.1250}}
  >{\raggedright\arraybackslash}p{(\linewidth - 14\tabcolsep) * \real{0.1250}}
  >{\raggedright\arraybackslash}p{(\linewidth - 14\tabcolsep) * \real{0.1250}}@{}}
\toprule\noalign{}
\begin{minipage}[b]{\linewidth}\raggedright
RUNSPEC
\end{minipage} & \begin{minipage}[b]{\linewidth}\raggedright
GRID
\end{minipage} & \begin{minipage}[b]{\linewidth}\raggedright
EDIT
\end{minipage} & \begin{minipage}[b]{\linewidth}\raggedright
PROPS
\end{minipage} & \begin{minipage}[b]{\linewidth}\raggedright
REGIONS
\end{minipage} & \begin{minipage}[b]{\linewidth}\raggedright
SOLUTION
\end{minipage} & \begin{minipage}[b]{\linewidth}\raggedright
SUMMARY
\end{minipage} & \begin{minipage}[b]{\linewidth}\raggedright
SCHEDULE
\end{minipage} \\
\midrule\noalign{}
\endhead
\bottomrule\noalign{}
\endlastfoot
\end{longtable}

\subsubsection{Description}\label{description-510}

KRW defines the scaling parameter at the maximum drainage water relative
permeability value (SWU), that is for Sw = 1.0, for all the cells in the
model via an array. The ENDSCALE keyword in the RUNSPEC section should
be activated to enable end-point scaling and the use of this keyword.
The SCALCERS keyword in the PROPS section defines the options used in
the re-scaling process, the options are two point scaling and three
point scaling.

\begin{longtable}[]{@{}
  >{\raggedright\arraybackslash}p{(\linewidth - 6\tabcolsep) * \real{0.2500}}
  >{\raggedright\arraybackslash}p{(\linewidth - 6\tabcolsep) * \real{0.2500}}
  >{\raggedright\arraybackslash}p{(\linewidth - 6\tabcolsep) * \real{0.2500}}
  >{\raggedright\arraybackslash}p{(\linewidth - 6\tabcolsep) * \real{0.2500}}@{}}
\toprule\noalign{}
\begin{minipage}[b]{\linewidth}\raggedright
No.
\end{minipage} & \begin{minipage}[b]{\linewidth}\raggedright
Name
\end{minipage} & \begin{minipage}[b]{\linewidth}\raggedright
Description
\end{minipage} & \begin{minipage}[b]{\linewidth}\raggedright
Default
\end{minipage} \\
\midrule\noalign{}
\endhead
\bottomrule\noalign{}
\endlastfoot
Field & Metric & Laboratory & \\
1 & KRW & KRW is an array of positive real numbers which are greater
than zero and less than or equal to 1.0, that are the assigned scaling
KRW values for each cell in the model. Repeat counts may be used, for
example 50*1.000. & Taken from cell allocated relative permeability
table. \\
dimensionless & dimensionless & dimensionless & \\
Notes: & & & \\
\end{longtable}

\emph{Table 8.83: KRW Keyword Description}

\begin{verbatim}
\end{verbatim}

For the two point scaling option and for the KRWR water relative
permeability array NOT present in the input deck the krw value for a
grid block is scaled by:

\begin{equation}\phantomsection\label{eq-8-63}{
{k}_{\mathit{rw}} = {k}_{{\mathit{rw}}_{ \mathit{TABLE}}}\left(\frac{\mathit{KRW}}{{k}_{{\mathit{rw}}_{ \mathit{TABLE}-\mathit{MAX}}}}\right)
}\end{equation}

Where:

\(\mathit{krw}\) = the resulting KRW value for a grid cell.

\(\mathit{KRW}\) = the scaling water relative permeability value from
the KRW array for a

given cell.

\({k}_{{\mathit{rw}}_{ \mathit{TABLE}}}\) = the water relative
permeability from a grid block's oil relative permeability

table at the grid blocks water saturation.

\({k}_{{\mathit{rw}}_{ \mathit{TABLE}-\mathit{MAX}}}\) = the maximum
water relative permeability from a grid block's water

relative table, that is at the maximum water saturation.

If the KRWR keyword is present in the input deck then the scaling
matches the relative permeability at the critical saturation of the
displacing phase.

If three point scaling option has been selected via the SCALECRS keyword
in the PROPS section the critical displacing phase is defined as:

\begin{longtable}[]{@{}lll@{}}
\toprule\noalign{}
No & Phases Present & Critical Saturation \\
\midrule\noalign{}
\endhead
\bottomrule\noalign{}
\endlastfoot
1 & Gas-Oil & S critical = 1.0 -- SOWCR - SGL \\
2 & Gas-Oil-Water & S critical = 1.0 -- SOWCR - SGL \\
3 & Gas-Water & S critical = 1.0 -- SGCR \\
\end{longtable}

\emph{Table 8.84: Critical Displacement Relationships}

End-point scaling allows the entered relative permeability functions to
be re-scaled based on the saturation end-points allocated to each cell
by the SWL, SWCR, SWU, SGL, SGCR, SGU, SOWCR, and SOGCR saturation grid
arrays for the saturation end-points, In addition end-point scaling may
be directional dependent in which case the directional dependent
versions of the aforementioned arrays should be used, that is SWUX, SWUY
and SWUZ instead of SWU, There is also the facility to make the
directional end-point scaling reversible or non-reversible and if the
non-reversible option is selected the non-reversible versions of the
aforementioned arrays should be used, that is SWUX, SWUX-, SWUY, SWUY-,
SWUZ and SWUZ-, instead of the SWU keyword.

End-point scaling also allows the entered relative permeability
functions to be scale on the relative permeability values using the KRG,
KRGR, KRW, KRORG, KRORW, KRW and KRWR relative permeability grid cell
arrays for the relative permeability end-point data. In addition
end-point scaling may be directional dependent in which case the
directional dependent versions of the aforementioned arrays should be
used, that is KRWX, KRWY and KRWZ instead of KRW, There is also the
facility to make the directional end-point scaling reversible or
non-reversible and if the non-reversible option is selected the
non-reversible versions of the aforementioned arrays should be used,
that is KRWX, KRWX-, KRWY, KRWY-, KRWZ and KRWZ-, instead of the KRW
keyword.

If the hysteresis model option has been activated on the SATOPTS keyword
in the RUNSPEC section, then the equivalent imbibition arrays suffixed
with the letter I, for example IKRW, can be used to define the KRW for
the relative permeability imbibition tables.

\subsubsection{Examples}\label{examples-34}

The first example defines an input box for the whole grid and for layers
one to three, for layer one KRW is set equal to 0.855, for layer two KRW
equals 0.875, and for layer three KRW equals 0.900.

\begin{verbatim}
--
--       DEFINE INPUT BOX FOR EDITING INPUT ARRAYS (NX=100, NY=100)
--
--       ---------- BOX ---------
--       I1  I2   J1  J2   K1  K2
BOX
         1*  1*   1*  1*   1   3                           / DEFINE BOX AREA
--
--       SET KRW VALUES FOR THREE LAYERS IN THE MODEL
--
KRW
         10000*0.855  10000*0.875  10000*0.900                /
--
--       DEFINE END OF INPUT BOX EDITING OF INPUT ARRAYS
--
ENDBOX
\end{verbatim}

The next example does exactly the same thing using the EQUALS keyword
instead.

\begin{verbatim}
--       ARRAY       CONSTANT     ---------- BOX ---------
--                                I1  I2   J1  J2   K1  K2
EQUALS
         KRW         0.8550       1*  1*   1*  1*    1   1  / KRW FOR LAYER 1
         KRW         0.8750       1*  1*   1*  1*    2   2  / KRW FOR LAYER 2
         KRW         0.9000       1*  1*   1*  1*    3   3  / KRW FOR LAYER 3
/
\end{verbatim}

\subsection{KRWR -- End-Point Scaling of Grid Cell KRWR(Sowcr)
(Drainage)}\label{krwr-end-point-scaling-of-grid-cell-krwrsowcr-drainage}

\begin{longtable}[]{@{}
  >{\raggedright\arraybackslash}p{(\linewidth - 14\tabcolsep) * \real{0.1250}}
  >{\raggedright\arraybackslash}p{(\linewidth - 14\tabcolsep) * \real{0.1250}}
  >{\raggedright\arraybackslash}p{(\linewidth - 14\tabcolsep) * \real{0.1250}}
  >{\raggedright\arraybackslash}p{(\linewidth - 14\tabcolsep) * \real{0.1250}}
  >{\raggedright\arraybackslash}p{(\linewidth - 14\tabcolsep) * \real{0.1250}}
  >{\raggedright\arraybackslash}p{(\linewidth - 14\tabcolsep) * \real{0.1250}}
  >{\raggedright\arraybackslash}p{(\linewidth - 14\tabcolsep) * \real{0.1250}}
  >{\raggedright\arraybackslash}p{(\linewidth - 14\tabcolsep) * \real{0.1250}}@{}}
\toprule\noalign{}
\begin{minipage}[b]{\linewidth}\raggedright
RUNSPEC
\end{minipage} & \begin{minipage}[b]{\linewidth}\raggedright
GRID
\end{minipage} & \begin{minipage}[b]{\linewidth}\raggedright
EDIT
\end{minipage} & \begin{minipage}[b]{\linewidth}\raggedright
PROPS
\end{minipage} & \begin{minipage}[b]{\linewidth}\raggedright
REGIONS
\end{minipage} & \begin{minipage}[b]{\linewidth}\raggedright
SOLUTION
\end{minipage} & \begin{minipage}[b]{\linewidth}\raggedright
SUMMARY
\end{minipage} & \begin{minipage}[b]{\linewidth}\raggedright
SCHEDULE
\end{minipage} \\
\midrule\noalign{}
\endhead
\bottomrule\noalign{}
\endlastfoot
\end{longtable}

\subsubsection{Description}\label{description-511}

KRWR defines the scaling parameter at the drainage critical oil to water
saturation value (SOWCR), for the drainage water relative permeability
curve, for all the cells in the model via an array. The ENDSCALE keyword
in the RUNSPEC section should be activated to enable end-point scaling
and the use of this keyword. The SCALCERS keyword in the PROPS section
defines the options used in the re-scaling process, the options are two
point scaling and three point scaling.

\begin{longtable}[]{@{}
  >{\raggedright\arraybackslash}p{(\linewidth - 6\tabcolsep) * \real{0.2500}}
  >{\raggedright\arraybackslash}p{(\linewidth - 6\tabcolsep) * \real{0.2500}}
  >{\raggedright\arraybackslash}p{(\linewidth - 6\tabcolsep) * \real{0.2500}}
  >{\raggedright\arraybackslash}p{(\linewidth - 6\tabcolsep) * \real{0.2500}}@{}}
\toprule\noalign{}
\begin{minipage}[b]{\linewidth}\raggedright
No.
\end{minipage} & \begin{minipage}[b]{\linewidth}\raggedright
Name
\end{minipage} & \begin{minipage}[b]{\linewidth}\raggedright
Description
\end{minipage} & \begin{minipage}[b]{\linewidth}\raggedright
Default
\end{minipage} \\
\midrule\noalign{}
\endhead
\bottomrule\noalign{}
\endlastfoot
Field & Metric & Laboratory & \\
1 & KRWR & KRWR is an array of positive real numbers which are greater
than zero and less than or equal to 1.0, that are the assigned scaling
KRWR values for each cell in the model. Repeat counts may be used, for
example 50*1.000. & Taken from cell allocated relative permeability
table. \\
dimensionless & dimensionless & dimensionless & \\
Notes: & & & \\
\end{longtable}

\emph{Table 8.85: KRWR Keyword Description}

\begin{verbatim}
\end{verbatim}

When the KRWR keyword is present in the input deck then the scaling
matches the relative permeability at the critical saturation of the
displacing phase.

If three point scaling option has been selected via the SCALECRS keyword
in the PROPS section the critical displacing phase is defined as:

\begin{longtable}[]{@{}lll@{}}
\toprule\noalign{}
No & Phases Present & Critical Saturation \\
\midrule\noalign{}
\endhead
\bottomrule\noalign{}
\endlastfoot
1 & Gas-Oil & S critical = 1.0 -- SOWCR - SGL \\
2 & Gas-Oil-Water & S critical = 1.0 -- SOWCR - SGL \\
3 & Gas-Water & S critical = 1.0 -- SGCR \\
\end{longtable}

\emph{Table 8.86: Critical Displacement Relationships}

End-point scaling allows the entered relative permeability functions to
be re-scaled based on the saturation end-points allocated to each cell
by the SWL, SWCR, SWU, SGL, SGCR, SGU, SOWCR, and SOGCR saturation grid
arrays for the saturation end-points, In addition end-point scaling may
be directional dependent in which case the directional dependent
versions of the aforementioned arrays should be used, that is SWUX, SWUY
and SWUZ instead of SWU, There is also the facility to make the
directional end-point scaling reversible or non-reversible and if the
non-reversible option is selected the non-reversible versions of the
aforementioned arrays should be used, that is SWUX, SWUX-, SWUY, SWUY-,
SWUZ and SWUZ-, instead of the SWU keyword.

End-point scaling also allows the entered relative permeability
functions to be scale on the relative permeability values using the KRG,
KRGR, KRWR, KRORG, KRORW, KRWR and KRWRR relative permeability grid cell
arrays for the relative permeability end-point data. In addition
end-point scaling may be directional dependent in which case the
directional dependent versions of the aforementioned arrays should be
used, that is KRWRX, KRWRY and KRWRZ instead of KRWR, There is also the
facility to make the directional end-point scaling reversible or
non-reversible and if the non-reversible option is selected the
non-reversible versions of the aforementioned arrays should be used,
that is KRWRX, KRWRX-, KRWRY, KRWRY-, KRWRZ and KRWRZ-, instead of the
KRWR keyword.

If the hysteresis model option has been activated on the SATOPTS keyword
in the RUNSPEC section, then the equivalent imbibition arrays suffixed
with the letter I, for example IKRWR, can be used to define the KRWR for
the relative permeability imbibition tables.

\subsubsection{Examples}\label{examples-35}

The first example defines an input box for the whole grid and for layers
one to three, for layer one KRWR is set equal to 0.755, for layer two
KRWR equals 0.775, and for layer three KRWR equals 0.800.

\begin{verbatim}
--
--       DEFINE INPUT BOX FOR EDITING INPUT ARRAYS (NX=100, NY=100)
--
--       ---------- BOX ---------
--       I1  I2   J1  J2   K1  K2
BOX
         1*  1*   1*  1*   1   3                           / DEFINE BOX AREA
--
--       SET KRWR VALUES FOR THREE LAYERS IN THE MODEL
--
KRWR
         10000*0.755  10000*0.775  10000*0.800                /
--
--       DEFINE END OF INPUT BOX EDITING OF INPUT ARRAYS
--
ENDBOX
\end{verbatim}

The next example does exactly the same thing using the EQUALS keyword
instead.

\begin{verbatim}
--
--       ARRAY       CONSTANT     ---------- BOX ---------
--                                I1  I2   J1  J2   K1  K2
EQUALS
         KRWR        0.7550       1*  1*   1*  1*    1   1  / KRWR FOR LAYER 1
         KRWR        0.7750       1*  1*   1*  1*    2   2  / KRWR FOR LAYER 2
         KRWR        0.8000       1*  1*   1*  1*    3   3  / KRWR FOR LAYER 3
/
\end{verbatim}

\subsection{LANGMPL -- Define Langmuir Pressure Grid Cell
Multiplier}\label{langmpl-define-langmuir-pressure-grid-cell-multiplier}

\begin{longtable}[]{@{}
  >{\raggedright\arraybackslash}p{(\linewidth - 14\tabcolsep) * \real{0.1250}}
  >{\raggedright\arraybackslash}p{(\linewidth - 14\tabcolsep) * \real{0.1250}}
  >{\raggedright\arraybackslash}p{(\linewidth - 14\tabcolsep) * \real{0.1250}}
  >{\raggedright\arraybackslash}p{(\linewidth - 14\tabcolsep) * \real{0.1250}}
  >{\raggedright\arraybackslash}p{(\linewidth - 14\tabcolsep) * \real{0.1250}}
  >{\raggedright\arraybackslash}p{(\linewidth - 14\tabcolsep) * \real{0.1250}}
  >{\raggedright\arraybackslash}p{(\linewidth - 14\tabcolsep) * \real{0.1250}}
  >{\raggedright\arraybackslash}p{(\linewidth - 14\tabcolsep) * \real{0.1250}}@{}}
\toprule\noalign{}
\begin{minipage}[b]{\linewidth}\raggedright
RUNSPEC
\end{minipage} & \begin{minipage}[b]{\linewidth}\raggedright
GRID
\end{minipage} & \begin{minipage}[b]{\linewidth}\raggedright
EDIT
\end{minipage} & \begin{minipage}[b]{\linewidth}\raggedright
PROPS
\end{minipage} & \begin{minipage}[b]{\linewidth}\raggedright
REGIONS
\end{minipage} & \begin{minipage}[b]{\linewidth}\raggedright
SOLUTION
\end{minipage} & \begin{minipage}[b]{\linewidth}\raggedright
SUMMARY
\end{minipage} & \begin{minipage}[b]{\linewidth}\raggedright
SCHEDULE
\end{minipage} \\
\midrule\noalign{}
\endhead
\bottomrule\noalign{}
\endlastfoot
\end{longtable}

\subsubsection{Description}\label{description-512}

This keyword, LANGMPL, defines the coal bed methane Langmuir Adsorption
{[}Langmuir, Irving (June 1918). ``The Adsorption of Gases on Plane
Surface of Glass, Mica and Platinum''. The Research Laboratory of the
General Electric Company. 40 (9): 1361--1402. doi:10.1021/ja02242a004{]}
pressure multiplier for each grid block, for when the Coal Bed Methane
option has been activated via the COAL keyword in the RUNSPEC section.
The keyword applies the multiplier to the pressure values in a cell's
allocated Langmuir table when calculating a cell's adsorption capacity.
See the LANGMUIR keyword in the PROPS section for specifying the
Langmuir tables for the model.

Note that if he Dual Porosity model has been activated by either the
DUALPORO or the DUALPERM keywords in the RUNSPEC section, then LANGMPL
applies to only the matrix grid block.

This keyword is not supported by OPM Flow but would change the results
if supported so the simulation will be stopped.

\subsection{LANGMUIR -- Langmuir Adsorption Isotherm
Tables}\label{langmuir-langmuir-adsorption-isotherm-tables}

\begin{longtable}[]{@{}
  >{\raggedright\arraybackslash}p{(\linewidth - 14\tabcolsep) * \real{0.1250}}
  >{\raggedright\arraybackslash}p{(\linewidth - 14\tabcolsep) * \real{0.1250}}
  >{\raggedright\arraybackslash}p{(\linewidth - 14\tabcolsep) * \real{0.1250}}
  >{\raggedright\arraybackslash}p{(\linewidth - 14\tabcolsep) * \real{0.1250}}
  >{\raggedright\arraybackslash}p{(\linewidth - 14\tabcolsep) * \real{0.1250}}
  >{\raggedright\arraybackslash}p{(\linewidth - 14\tabcolsep) * \real{0.1250}}
  >{\raggedright\arraybackslash}p{(\linewidth - 14\tabcolsep) * \real{0.1250}}
  >{\raggedright\arraybackslash}p{(\linewidth - 14\tabcolsep) * \real{0.1250}}@{}}
\toprule\noalign{}
\begin{minipage}[b]{\linewidth}\raggedright
RUNSPEC
\end{minipage} & \begin{minipage}[b]{\linewidth}\raggedright
GRID
\end{minipage} & \begin{minipage}[b]{\linewidth}\raggedright
EDIT
\end{minipage} & \begin{minipage}[b]{\linewidth}\raggedright
PROPS
\end{minipage} & \begin{minipage}[b]{\linewidth}\raggedright
REGIONS
\end{minipage} & \begin{minipage}[b]{\linewidth}\raggedright
SOLUTION
\end{minipage} & \begin{minipage}[b]{\linewidth}\raggedright
SUMMARY
\end{minipage} & \begin{minipage}[b]{\linewidth}\raggedright
SCHEDULE
\end{minipage} \\
\midrule\noalign{}
\endhead
\bottomrule\noalign{}
\endlastfoot
\end{longtable}

\subsubsection{Description}\label{description-513}

The LANGMUIR keyword defines the coal bed methane Langmuir Adsorption
Isotherms {[}Langmuir, Irving (June 1918). ``The Adsorption of Gases on
Plane Surface of Glass, Mica and Platinum''. The Research Laboratory of
the General Electric Company. 40 (9): 1361--1402.
doi:10.1021/ja02242a004{]} tables, for when the Coal Bed Methane option
has been activated via the COAL keyword in the RUNSPEC section. See the
COALNUM keyword in the GRID section for allocating the Langmuir tables
to the grid blocks and also the LANGMPL keyword in the PROPS section for
re-scaling the pressure values in the tables that are allocated to a
cell.

This keyword is not supported by OPM Flow but would change the results
if supported so the simulation will be stopped.

\subsection{LANGSOLV -- Langmuir Adsorption Isotherm Solvent
Tables}\label{langsolv-langmuir-adsorption-isotherm-solvent-tables}

\begin{longtable}[]{@{}
  >{\raggedright\arraybackslash}p{(\linewidth - 14\tabcolsep) * \real{0.1250}}
  >{\raggedright\arraybackslash}p{(\linewidth - 14\tabcolsep) * \real{0.1250}}
  >{\raggedright\arraybackslash}p{(\linewidth - 14\tabcolsep) * \real{0.1250}}
  >{\raggedright\arraybackslash}p{(\linewidth - 14\tabcolsep) * \real{0.1250}}
  >{\raggedright\arraybackslash}p{(\linewidth - 14\tabcolsep) * \real{0.1250}}
  >{\raggedright\arraybackslash}p{(\linewidth - 14\tabcolsep) * \real{0.1250}}
  >{\raggedright\arraybackslash}p{(\linewidth - 14\tabcolsep) * \real{0.1250}}
  >{\raggedright\arraybackslash}p{(\linewidth - 14\tabcolsep) * \real{0.1250}}@{}}
\toprule\noalign{}
\begin{minipage}[b]{\linewidth}\raggedright
RUNSPEC
\end{minipage} & \begin{minipage}[b]{\linewidth}\raggedright
GRID
\end{minipage} & \begin{minipage}[b]{\linewidth}\raggedright
EDIT
\end{minipage} & \begin{minipage}[b]{\linewidth}\raggedright
PROPS
\end{minipage} & \begin{minipage}[b]{\linewidth}\raggedright
REGIONS
\end{minipage} & \begin{minipage}[b]{\linewidth}\raggedright
SOLUTION
\end{minipage} & \begin{minipage}[b]{\linewidth}\raggedright
SUMMARY
\end{minipage} & \begin{minipage}[b]{\linewidth}\raggedright
SCHEDULE
\end{minipage} \\
\midrule\noalign{}
\endhead
\bottomrule\noalign{}
\endlastfoot
\end{longtable}

\subsubsection{Description}\label{description-514}

The LANGMUIR keyword defines the coal bed methane Langmuir Adsorption
Isotherms {[}Langmuir, Irving (June 1918). ``The Adsorption of Gases on
Plane Surface of Glass, Mica and Platinum''. The Research Laboratory of
the General Electric Company. 40 (9): 1361--1402.
doi:10.1021/ja02242a004{]} Solvent tables, for when the Coal Bed Methane
option has been activated via the COAL keyword and the Solvent phase has
been declared by the SOLVENT keyword in the RUNSPEC section. See the
COALNUM keyword in the GRID section for allocating the Langmuir solvent
tables to the grid blocks, and also the LANGMUIR keyword in the PROPS
section for defining the Langmuir Adsorption Isotherm tables. Keywords
COALADS and COALPP, also in the PROPS section, are used to specify the
relative adsorption data in runs containing the solvent phase.

This keyword is not supported by OPM Flow but would change the results
if supported so the simulation will be stopped.

\subsection{LCUNIT -- Define Linear Combination Rate and Volume
Units}\label{lcunit-define-linear-combination-rate-and-volume-units}

\begin{longtable}[]{@{}
  >{\raggedright\arraybackslash}p{(\linewidth - 14\tabcolsep) * \real{0.1250}}
  >{\raggedright\arraybackslash}p{(\linewidth - 14\tabcolsep) * \real{0.1250}}
  >{\raggedright\arraybackslash}p{(\linewidth - 14\tabcolsep) * \real{0.1250}}
  >{\raggedright\arraybackslash}p{(\linewidth - 14\tabcolsep) * \real{0.1250}}
  >{\raggedright\arraybackslash}p{(\linewidth - 14\tabcolsep) * \real{0.1250}}
  >{\raggedright\arraybackslash}p{(\linewidth - 14\tabcolsep) * \real{0.1250}}
  >{\raggedright\arraybackslash}p{(\linewidth - 14\tabcolsep) * \real{0.1250}}
  >{\raggedright\arraybackslash}p{(\linewidth - 14\tabcolsep) * \real{0.1250}}@{}}
\toprule\noalign{}
\begin{minipage}[b]{\linewidth}\raggedright
RUNSPEC
\end{minipage} & \begin{minipage}[b]{\linewidth}\raggedright
GRID
\end{minipage} & \begin{minipage}[b]{\linewidth}\raggedright
EDIT
\end{minipage} & \begin{minipage}[b]{\linewidth}\raggedright
PROPS
\end{minipage} & \begin{minipage}[b]{\linewidth}\raggedright
REGIONS
\end{minipage} & \begin{minipage}[b]{\linewidth}\raggedright
SOLUTION
\end{minipage} & \begin{minipage}[b]{\linewidth}\raggedright
SUMMARY
\end{minipage} & \begin{minipage}[b]{\linewidth}\raggedright
SCHEDULE
\end{minipage} \\
\midrule\noalign{}
\endhead
\bottomrule\noalign{}
\endlastfoot
\end{longtable}

\subsubsection{Description}\label{description-515}

The LCUNIT keyword defines the units for the Linear Combination facility
which allows for a linear combination of oil, gas and water rates and
volumes to be used as combination targets and constraints in controlling
group and well production and injection data. See also the LINCOM in the
SCHEDULE section that defines the actual linear combination equation.

This keyword is not supported by OPM Flow but would change the results
if supported so the simulation will be stopped.

\subsection{LKRO -- End-Point Scaling of Grid Cell Kro(Swl) (Low
Salinity and Oil
Wet)}\label{lkro-end-point-scaling-of-grid-cell-kroswl-low-salinity-and-oil-wet}

\begin{longtable}[]{@{}
  >{\raggedright\arraybackslash}p{(\linewidth - 14\tabcolsep) * \real{0.1250}}
  >{\raggedright\arraybackslash}p{(\linewidth - 14\tabcolsep) * \real{0.1250}}
  >{\raggedright\arraybackslash}p{(\linewidth - 14\tabcolsep) * \real{0.1250}}
  >{\raggedright\arraybackslash}p{(\linewidth - 14\tabcolsep) * \real{0.1250}}
  >{\raggedright\arraybackslash}p{(\linewidth - 14\tabcolsep) * \real{0.1250}}
  >{\raggedright\arraybackslash}p{(\linewidth - 14\tabcolsep) * \real{0.1250}}
  >{\raggedright\arraybackslash}p{(\linewidth - 14\tabcolsep) * \real{0.1250}}
  >{\raggedright\arraybackslash}p{(\linewidth - 14\tabcolsep) * \real{0.1250}}@{}}
\toprule\noalign{}
\begin{minipage}[b]{\linewidth}\raggedright
RUNSPEC
\end{minipage} & \begin{minipage}[b]{\linewidth}\raggedright
GRID
\end{minipage} & \begin{minipage}[b]{\linewidth}\raggedright
EDIT
\end{minipage} & \begin{minipage}[b]{\linewidth}\raggedright
PROPS
\end{minipage} & \begin{minipage}[b]{\linewidth}\raggedright
REGIONS
\end{minipage} & \begin{minipage}[b]{\linewidth}\raggedright
SOLUTION
\end{minipage} & \begin{minipage}[b]{\linewidth}\raggedright
SUMMARY
\end{minipage} & \begin{minipage}[b]{\linewidth}\raggedright
SCHEDULE
\end{minipage} \\
\midrule\noalign{}
\endhead
\bottomrule\noalign{}
\endlastfoot
\end{longtable}

\subsubsection{Description}\label{description-516}

LKRO defines the scaling parameter for the oil relative permeability
value at the connate water saturation (SWL), for all the cells in the
model via an array, for when the Low Salt option and the End-point
Scaling options has been activated by the LOWSALT and the ENDSCALE
keywords in the RUNSPEC section. The data is used to scale the oil
relative permeability in the low salinity oil wet oil relative
permeability saturation tables.

This keyword is not supported by OPM Flow but would change the results
if supported so the simulation will be stopped.

\subsection{LKRORG -- End-Point Scaling of Grid Cell Kro(Sgcr) (Low
Salinity and Oil
Wet)}\label{lkrorg-end-point-scaling-of-grid-cell-krosgcr-low-salinity-and-oil-wet}

\begin{longtable}[]{@{}
  >{\raggedright\arraybackslash}p{(\linewidth - 14\tabcolsep) * \real{0.1250}}
  >{\raggedright\arraybackslash}p{(\linewidth - 14\tabcolsep) * \real{0.1250}}
  >{\raggedright\arraybackslash}p{(\linewidth - 14\tabcolsep) * \real{0.1250}}
  >{\raggedright\arraybackslash}p{(\linewidth - 14\tabcolsep) * \real{0.1250}}
  >{\raggedright\arraybackslash}p{(\linewidth - 14\tabcolsep) * \real{0.1250}}
  >{\raggedright\arraybackslash}p{(\linewidth - 14\tabcolsep) * \real{0.1250}}
  >{\raggedright\arraybackslash}p{(\linewidth - 14\tabcolsep) * \real{0.1250}}
  >{\raggedright\arraybackslash}p{(\linewidth - 14\tabcolsep) * \real{0.1250}}@{}}
\toprule\noalign{}
\begin{minipage}[b]{\linewidth}\raggedright
RUNSPEC
\end{minipage} & \begin{minipage}[b]{\linewidth}\raggedright
GRID
\end{minipage} & \begin{minipage}[b]{\linewidth}\raggedright
EDIT
\end{minipage} & \begin{minipage}[b]{\linewidth}\raggedright
PROPS
\end{minipage} & \begin{minipage}[b]{\linewidth}\raggedright
REGIONS
\end{minipage} & \begin{minipage}[b]{\linewidth}\raggedright
SOLUTION
\end{minipage} & \begin{minipage}[b]{\linewidth}\raggedright
SUMMARY
\end{minipage} & \begin{minipage}[b]{\linewidth}\raggedright
SCHEDULE
\end{minipage} \\
\midrule\noalign{}
\endhead
\bottomrule\noalign{}
\endlastfoot
\end{longtable}

\subsubsection{Description}\label{description-517}

LKRORG defines the scaling parameter for the relative permeability of
oil at the critical gas saturation (SGCR), for all the cells in the
model via an array, for when the Low Salt option and the End-point
Scaling options has been activated by the LOWSALT and the ENDSCALE
keywords in the RUNSPEC section. The data is used to scale the oil
relative permeability in the low salinity oil wet oil relative
permeability saturation tables.

This keyword is not supported by OPM Flow but would change the results
if supported so the simulation will be stopped.

\subsection{LKRORW -- End-Point Scaling of Grid Cell Kro(Swcr) (Low
Salinity and Oil
Wet)}\label{lkrorw-end-point-scaling-of-grid-cell-kroswcr-low-salinity-and-oil-wet}

\begin{longtable}[]{@{}
  >{\raggedright\arraybackslash}p{(\linewidth - 14\tabcolsep) * \real{0.1250}}
  >{\raggedright\arraybackslash}p{(\linewidth - 14\tabcolsep) * \real{0.1250}}
  >{\raggedright\arraybackslash}p{(\linewidth - 14\tabcolsep) * \real{0.1250}}
  >{\raggedright\arraybackslash}p{(\linewidth - 14\tabcolsep) * \real{0.1250}}
  >{\raggedright\arraybackslash}p{(\linewidth - 14\tabcolsep) * \real{0.1250}}
  >{\raggedright\arraybackslash}p{(\linewidth - 14\tabcolsep) * \real{0.1250}}
  >{\raggedright\arraybackslash}p{(\linewidth - 14\tabcolsep) * \real{0.1250}}
  >{\raggedright\arraybackslash}p{(\linewidth - 14\tabcolsep) * \real{0.1250}}@{}}
\toprule\noalign{}
\begin{minipage}[b]{\linewidth}\raggedright
RUNSPEC
\end{minipage} & \begin{minipage}[b]{\linewidth}\raggedright
GRID
\end{minipage} & \begin{minipage}[b]{\linewidth}\raggedright
EDIT
\end{minipage} & \begin{minipage}[b]{\linewidth}\raggedright
PROPS
\end{minipage} & \begin{minipage}[b]{\linewidth}\raggedright
REGIONS
\end{minipage} & \begin{minipage}[b]{\linewidth}\raggedright
SOLUTION
\end{minipage} & \begin{minipage}[b]{\linewidth}\raggedright
SUMMARY
\end{minipage} & \begin{minipage}[b]{\linewidth}\raggedright
SCHEDULE
\end{minipage} \\
\midrule\noalign{}
\endhead
\bottomrule\noalign{}
\endlastfoot
\end{longtable}

\subsubsection{Description}\label{description-518}

LKRORW defines the scaling parameter for the relative permeability of
oil at the critical water saturation (SWCR), for all the cells in the
model via an array, for when the Low Salt option and the End-point
Scaling options has been activated by the LOWSALT and the ENDSCALE
keywords in the RUNSPEC section. The data is used to scale the oil
relative permeability in the low salinity oil wet oil relative
permeability saturation tables.

This keyword is not supported by OPM Flow but would change the results
if supported so the simulation will be stopped.

\subsection{LKRW -- End-Point Scaling of Grid Cell Krw(Sw =1.0) (Low
Salinity and Oil
Wet)}\label{lkrw-end-point-scaling-of-grid-cell-krwsw-1.0-low-salinity-and-oil-wet}

\begin{longtable}[]{@{}
  >{\raggedright\arraybackslash}p{(\linewidth - 14\tabcolsep) * \real{0.1250}}
  >{\raggedright\arraybackslash}p{(\linewidth - 14\tabcolsep) * \real{0.1250}}
  >{\raggedright\arraybackslash}p{(\linewidth - 14\tabcolsep) * \real{0.1250}}
  >{\raggedright\arraybackslash}p{(\linewidth - 14\tabcolsep) * \real{0.1250}}
  >{\raggedright\arraybackslash}p{(\linewidth - 14\tabcolsep) * \real{0.1250}}
  >{\raggedright\arraybackslash}p{(\linewidth - 14\tabcolsep) * \real{0.1250}}
  >{\raggedright\arraybackslash}p{(\linewidth - 14\tabcolsep) * \real{0.1250}}
  >{\raggedright\arraybackslash}p{(\linewidth - 14\tabcolsep) * \real{0.1250}}@{}}
\toprule\noalign{}
\begin{minipage}[b]{\linewidth}\raggedright
RUNSPEC
\end{minipage} & \begin{minipage}[b]{\linewidth}\raggedright
GRID
\end{minipage} & \begin{minipage}[b]{\linewidth}\raggedright
EDIT
\end{minipage} & \begin{minipage}[b]{\linewidth}\raggedright
PROPS
\end{minipage} & \begin{minipage}[b]{\linewidth}\raggedright
REGIONS
\end{minipage} & \begin{minipage}[b]{\linewidth}\raggedright
SOLUTION
\end{minipage} & \begin{minipage}[b]{\linewidth}\raggedright
SUMMARY
\end{minipage} & \begin{minipage}[b]{\linewidth}\raggedright
SCHEDULE
\end{minipage} \\
\midrule\noalign{}
\endhead
\bottomrule\noalign{}
\endlastfoot
\end{longtable}

\subsubsection{Description}\label{description-519}

LKRW defines the scaling parameter at the maximum oil relative
permeability value (SWU), that is for Sw = 1.0, for all the cells in the
model via an array, for when the Low Salt option and the End-point
Scaling options has been activated by the LOWSALT and the ENDSCALE
keywords in the RUNSPEC section. The data is used to scale the water
relative permeability in the low salinity oil wet water relative
permeability saturation tables.

This keyword is not supported by OPM Flow but would change the results
if supported so the simulation will be stopped.

\subsection{LKRWR -- End-Point Scaling of Grid Cell KRWR(Sowcr) (Low
Salinity and Oil
Wet)}\label{lkrwr-end-point-scaling-of-grid-cell-krwrsowcr-low-salinity-and-oil-wet}

\begin{longtable}[]{@{}
  >{\raggedright\arraybackslash}p{(\linewidth - 14\tabcolsep) * \real{0.1250}}
  >{\raggedright\arraybackslash}p{(\linewidth - 14\tabcolsep) * \real{0.1250}}
  >{\raggedright\arraybackslash}p{(\linewidth - 14\tabcolsep) * \real{0.1250}}
  >{\raggedright\arraybackslash}p{(\linewidth - 14\tabcolsep) * \real{0.1250}}
  >{\raggedright\arraybackslash}p{(\linewidth - 14\tabcolsep) * \real{0.1250}}
  >{\raggedright\arraybackslash}p{(\linewidth - 14\tabcolsep) * \real{0.1250}}
  >{\raggedright\arraybackslash}p{(\linewidth - 14\tabcolsep) * \real{0.1250}}
  >{\raggedright\arraybackslash}p{(\linewidth - 14\tabcolsep) * \real{0.1250}}@{}}
\toprule\noalign{}
\begin{minipage}[b]{\linewidth}\raggedright
RUNSPEC
\end{minipage} & \begin{minipage}[b]{\linewidth}\raggedright
GRID
\end{minipage} & \begin{minipage}[b]{\linewidth}\raggedright
EDIT
\end{minipage} & \begin{minipage}[b]{\linewidth}\raggedright
PROPS
\end{minipage} & \begin{minipage}[b]{\linewidth}\raggedright
REGIONS
\end{minipage} & \begin{minipage}[b]{\linewidth}\raggedright
SOLUTION
\end{minipage} & \begin{minipage}[b]{\linewidth}\raggedright
SUMMARY
\end{minipage} & \begin{minipage}[b]{\linewidth}\raggedright
SCHEDULE
\end{minipage} \\
\midrule\noalign{}
\endhead
\bottomrule\noalign{}
\endlastfoot
\end{longtable}

\subsubsection{Description}\label{description-520}

LKRWR defines the scaling parameter at the maximum oil relative
permeability value (SWU), that is for Sw = 1.0, for all the cells in the
model via an array, for when the Low Salt option and the End-point
Scaling options has been activated by the LOWSALT and the ENDSCALE
keywords in the RUNSPEC section. The data is used to scale the water
relative permeability in the low salinity oil wet water relative
permeability saturation tables.

This keyword is not supported by OPM Flow but would change the results
if supported so the simulation will be stopped.

\subsection{LPCW -- End-Point Scaling of Grid Cell Water Capillary
Pressure (Low Salinity and Oil
Wet)}\label{lpcw-end-point-scaling-of-grid-cell-water-capillary-pressure-low-salinity-and-oil-wet}

\begin{longtable}[]{@{}
  >{\raggedright\arraybackslash}p{(\linewidth - 14\tabcolsep) * \real{0.1250}}
  >{\raggedright\arraybackslash}p{(\linewidth - 14\tabcolsep) * \real{0.1250}}
  >{\raggedright\arraybackslash}p{(\linewidth - 14\tabcolsep) * \real{0.1250}}
  >{\raggedright\arraybackslash}p{(\linewidth - 14\tabcolsep) * \real{0.1250}}
  >{\raggedright\arraybackslash}p{(\linewidth - 14\tabcolsep) * \real{0.1250}}
  >{\raggedright\arraybackslash}p{(\linewidth - 14\tabcolsep) * \real{0.1250}}
  >{\raggedright\arraybackslash}p{(\linewidth - 14\tabcolsep) * \real{0.1250}}
  >{\raggedright\arraybackslash}p{(\linewidth - 14\tabcolsep) * \real{0.1250}}@{}}
\toprule\noalign{}
\begin{minipage}[b]{\linewidth}\raggedright
RUNSPEC
\end{minipage} & \begin{minipage}[b]{\linewidth}\raggedright
GRID
\end{minipage} & \begin{minipage}[b]{\linewidth}\raggedright
EDIT
\end{minipage} & \begin{minipage}[b]{\linewidth}\raggedright
PROPS
\end{minipage} & \begin{minipage}[b]{\linewidth}\raggedright
REGIONS
\end{minipage} & \begin{minipage}[b]{\linewidth}\raggedright
SOLUTION
\end{minipage} & \begin{minipage}[b]{\linewidth}\raggedright
SUMMARY
\end{minipage} & \begin{minipage}[b]{\linewidth}\raggedright
SCHEDULE
\end{minipage} \\
\midrule\noalign{}
\endhead
\bottomrule\noalign{}
\endlastfoot
\end{longtable}

\subsubsection{Description}\label{description-521}

LPCW defines the maximum oil-water pressure values for all the cells in
the model via an array, for when the Low Salt option and the End-point
Scaling options has been activated by the LOWSALT and the ENDSCALE
keywords in the RUNSPEC section. The keyword re-scales the oil-water
capillary pressure in the low salinity oil wet capillary saturation
tables from a cell's assigned saturation function by the grid block's
LPCW value.

The capillary pressure for a grid block is scaled by:

\begin{equation}\phantomsection\label{eq-8-64}{
{P}_{c} = {P}_{{c}_{\mathit{TABLE}}}\left(\frac{\mathit{LPCW}}{{P}_{{c}_{\mathit{TABLE}-\mathit{MAX}}}}\right)
}\end{equation}

Where:

\({P}_{c}\) = the resulting low salt oil wet water capillary pressure
for a grid cell.

\(\mathit{LPCW}\) = the maximum capillary pressure from the LPCW array
for a given cell.

\({P}_{{c}_{\mathit{TABLE}}}\) = the capillary pressure in the low salt
oil wet capillary pressure table

allocated to the grid block.

\({P}_{{c}_{\mathit{TABLE}-\mathit{MAX}}}\) = the maximum capillary
pressure in the low salt oil capillary pressure table

allocated to the grid block (that is at the connate water saturation).

This keyword is not supported by OPM Flow but would change the results
if supported so the simulation will be stopped.

\subsection{LSALTFNC -- Define Low Salt Weighting Factors versus Salt
Concentration
Functions}\label{lsaltfnc-define-low-salt-weighting-factors-versus-salt-concentration-functions}

\begin{longtable}[]{@{}
  >{\raggedright\arraybackslash}p{(\linewidth - 14\tabcolsep) * \real{0.1250}}
  >{\raggedright\arraybackslash}p{(\linewidth - 14\tabcolsep) * \real{0.1250}}
  >{\raggedright\arraybackslash}p{(\linewidth - 14\tabcolsep) * \real{0.1250}}
  >{\raggedright\arraybackslash}p{(\linewidth - 14\tabcolsep) * \real{0.1250}}
  >{\raggedright\arraybackslash}p{(\linewidth - 14\tabcolsep) * \real{0.1250}}
  >{\raggedright\arraybackslash}p{(\linewidth - 14\tabcolsep) * \real{0.1250}}
  >{\raggedright\arraybackslash}p{(\linewidth - 14\tabcolsep) * \real{0.1250}}
  >{\raggedright\arraybackslash}p{(\linewidth - 14\tabcolsep) * \real{0.1250}}@{}}
\toprule\noalign{}
\begin{minipage}[b]{\linewidth}\raggedright
RUNSPEC
\end{minipage} & \begin{minipage}[b]{\linewidth}\raggedright
GRID
\end{minipage} & \begin{minipage}[b]{\linewidth}\raggedright
EDIT
\end{minipage} & \begin{minipage}[b]{\linewidth}\raggedright
PROPS
\end{minipage} & \begin{minipage}[b]{\linewidth}\raggedright
REGIONS
\end{minipage} & \begin{minipage}[b]{\linewidth}\raggedright
SOLUTION
\end{minipage} & \begin{minipage}[b]{\linewidth}\raggedright
SUMMARY
\end{minipage} & \begin{minipage}[b]{\linewidth}\raggedright
SCHEDULE
\end{minipage} \\
\midrule\noalign{}
\endhead
\bottomrule\noalign{}
\endlastfoot
\end{longtable}

\subsubsection{Description}\label{description-522}

The LSALTFNC keyword defines the low salt weighting factors versus salt
concentration functions for when the Low Salt option has been activated
by the LOWSALT keyword in the RUNSPEC section. The tables are used to
modify the oil and water relative permeability saturation end-points, as
well as the water-oil capillary pressure end-points, for different salt
concentrations within a grid cell.

This keyword is not supported by OPM Flow but would change the results
if supported so the simulation will be stopped.

\subsection{LSOGCR -- End-Point Scaling Grid Cell SOGCR (Low Salinity
and Oil
Wet)}\label{lsogcr-end-point-scaling-grid-cell-sogcr-low-salinity-and-oil-wet}

\begin{longtable}[]{@{}
  >{\raggedright\arraybackslash}p{(\linewidth - 14\tabcolsep) * \real{0.1250}}
  >{\raggedright\arraybackslash}p{(\linewidth - 14\tabcolsep) * \real{0.1250}}
  >{\raggedright\arraybackslash}p{(\linewidth - 14\tabcolsep) * \real{0.1250}}
  >{\raggedright\arraybackslash}p{(\linewidth - 14\tabcolsep) * \real{0.1250}}
  >{\raggedright\arraybackslash}p{(\linewidth - 14\tabcolsep) * \real{0.1250}}
  >{\raggedright\arraybackslash}p{(\linewidth - 14\tabcolsep) * \real{0.1250}}
  >{\raggedright\arraybackslash}p{(\linewidth - 14\tabcolsep) * \real{0.1250}}
  >{\raggedright\arraybackslash}p{(\linewidth - 14\tabcolsep) * \real{0.1250}}@{}}
\toprule\noalign{}
\begin{minipage}[b]{\linewidth}\raggedright
RUNSPEC
\end{minipage} & \begin{minipage}[b]{\linewidth}\raggedright
GRID
\end{minipage} & \begin{minipage}[b]{\linewidth}\raggedright
EDIT
\end{minipage} & \begin{minipage}[b]{\linewidth}\raggedright
PROPS
\end{minipage} & \begin{minipage}[b]{\linewidth}\raggedright
REGIONS
\end{minipage} & \begin{minipage}[b]{\linewidth}\raggedright
SOLUTION
\end{minipage} & \begin{minipage}[b]{\linewidth}\raggedright
SUMMARY
\end{minipage} & \begin{minipage}[b]{\linewidth}\raggedright
SCHEDULE
\end{minipage} \\
\midrule\noalign{}
\endhead
\bottomrule\noalign{}
\endlastfoot
\end{longtable}

\subsubsection{Description}\label{description-523}

LSOGCR defines the critical oil saturation with respect to gas
(``SOGCR''), for all the cells in the model via an array, for when the
Low Salt option has been selected. The data is used to scale the oil
saturation in the low salinity oil wet oil-gas relative permeability
saturation tables. The ENDSCALE keyword in the RUNSPEC section should be
activated to enable end-point scaling and the use of this keyword. In
addition, the Low Salt option should be activated by the LOWSALT keyword
in the RUNSPEC section.

This keyword is not supported by OPM Flow but would change the results
if supported so the simulation will be stopped.

\subsection{LSOWCR -- End-Point Scaling Grid Cell SOWCR (Low Salinity
and Oil
Wet)}\label{lsowcr-end-point-scaling-grid-cell-sowcr-low-salinity-and-oil-wet}

\begin{longtable}[]{@{}
  >{\raggedright\arraybackslash}p{(\linewidth - 14\tabcolsep) * \real{0.1250}}
  >{\raggedright\arraybackslash}p{(\linewidth - 14\tabcolsep) * \real{0.1250}}
  >{\raggedright\arraybackslash}p{(\linewidth - 14\tabcolsep) * \real{0.1250}}
  >{\raggedright\arraybackslash}p{(\linewidth - 14\tabcolsep) * \real{0.1250}}
  >{\raggedright\arraybackslash}p{(\linewidth - 14\tabcolsep) * \real{0.1250}}
  >{\raggedright\arraybackslash}p{(\linewidth - 14\tabcolsep) * \real{0.1250}}
  >{\raggedright\arraybackslash}p{(\linewidth - 14\tabcolsep) * \real{0.1250}}
  >{\raggedright\arraybackslash}p{(\linewidth - 14\tabcolsep) * \real{0.1250}}@{}}
\toprule\noalign{}
\begin{minipage}[b]{\linewidth}\raggedright
RUNSPEC
\end{minipage} & \begin{minipage}[b]{\linewidth}\raggedright
GRID
\end{minipage} & \begin{minipage}[b]{\linewidth}\raggedright
EDIT
\end{minipage} & \begin{minipage}[b]{\linewidth}\raggedright
PROPS
\end{minipage} & \begin{minipage}[b]{\linewidth}\raggedright
REGIONS
\end{minipage} & \begin{minipage}[b]{\linewidth}\raggedright
SOLUTION
\end{minipage} & \begin{minipage}[b]{\linewidth}\raggedright
SUMMARY
\end{minipage} & \begin{minipage}[b]{\linewidth}\raggedright
SCHEDULE
\end{minipage} \\
\midrule\noalign{}
\endhead
\bottomrule\noalign{}
\endlastfoot
\end{longtable}

\subsubsection{Description}\label{description-524}

LSOWCR defines the critical oil saturation with respect to water
(``SOWCR''), for all the cells in the model via an array, for when the
Low Salt option has been selected. The data is used to scale the oil
saturation in the low salinity oil wet water-oil relative permeability
saturation tables. The ENDSCALE keyword in the RUNSPEC section should be
activated to enable end-point scaling and the use of this keyword. In
addition, the Low Salt option should be activated by the LOWSALT keyword
in the RUNSPEC section.

This keyword is not supported by OPM Flow but would change the results
if supported so the simulation will be stopped.

\subsection{LSWCR -- End-Point Scaling Grid Cell SWCR (Low Salinity and
Oil
Wet)}\label{lswcr-end-point-scaling-grid-cell-swcr-low-salinity-and-oil-wet}

\begin{longtable}[]{@{}
  >{\raggedright\arraybackslash}p{(\linewidth - 14\tabcolsep) * \real{0.1250}}
  >{\raggedright\arraybackslash}p{(\linewidth - 14\tabcolsep) * \real{0.1250}}
  >{\raggedright\arraybackslash}p{(\linewidth - 14\tabcolsep) * \real{0.1250}}
  >{\raggedright\arraybackslash}p{(\linewidth - 14\tabcolsep) * \real{0.1250}}
  >{\raggedright\arraybackslash}p{(\linewidth - 14\tabcolsep) * \real{0.1250}}
  >{\raggedright\arraybackslash}p{(\linewidth - 14\tabcolsep) * \real{0.1250}}
  >{\raggedright\arraybackslash}p{(\linewidth - 14\tabcolsep) * \real{0.1250}}
  >{\raggedright\arraybackslash}p{(\linewidth - 14\tabcolsep) * \real{0.1250}}@{}}
\toprule\noalign{}
\begin{minipage}[b]{\linewidth}\raggedright
RUNSPEC
\end{minipage} & \begin{minipage}[b]{\linewidth}\raggedright
GRID
\end{minipage} & \begin{minipage}[b]{\linewidth}\raggedright
EDIT
\end{minipage} & \begin{minipage}[b]{\linewidth}\raggedright
PROPS
\end{minipage} & \begin{minipage}[b]{\linewidth}\raggedright
REGIONS
\end{minipage} & \begin{minipage}[b]{\linewidth}\raggedright
SOLUTION
\end{minipage} & \begin{minipage}[b]{\linewidth}\raggedright
SUMMARY
\end{minipage} & \begin{minipage}[b]{\linewidth}\raggedright
SCHEDULE
\end{minipage} \\
\midrule\noalign{}
\endhead
\bottomrule\noalign{}
\endlastfoot
\end{longtable}

\subsubsection{Description}\label{description-525}

LSWCR defines the critical water saturation (``SWCR''), for all the
cells in the model via an array, for when the Low Salt option has been
selected. The data is used to scale the water saturation in the low
salinity oil wet water-oil relative permeability saturation tables. The
ENDSCALE keyword in the RUNSPEC section should be activated to enable
end-point scaling and the use of this keyword. In addition, the Low Salt
option should be activated by the LOWSALT keyword in the RUNSPEC
section.

This keyword is not supported by OPM Flow but would change the results
if supported so the simulation will be stopped.

\subsection{LSWL -- End-Point Scaling Grid Cell SWL (Low Salinity and
Oil
Wet)}\label{lswl-end-point-scaling-grid-cell-swl-low-salinity-and-oil-wet}

\begin{longtable}[]{@{}
  >{\raggedright\arraybackslash}p{(\linewidth - 14\tabcolsep) * \real{0.1250}}
  >{\raggedright\arraybackslash}p{(\linewidth - 14\tabcolsep) * \real{0.1250}}
  >{\raggedright\arraybackslash}p{(\linewidth - 14\tabcolsep) * \real{0.1250}}
  >{\raggedright\arraybackslash}p{(\linewidth - 14\tabcolsep) * \real{0.1250}}
  >{\raggedright\arraybackslash}p{(\linewidth - 14\tabcolsep) * \real{0.1250}}
  >{\raggedright\arraybackslash}p{(\linewidth - 14\tabcolsep) * \real{0.1250}}
  >{\raggedright\arraybackslash}p{(\linewidth - 14\tabcolsep) * \real{0.1250}}
  >{\raggedright\arraybackslash}p{(\linewidth - 14\tabcolsep) * \real{0.1250}}@{}}
\toprule\noalign{}
\begin{minipage}[b]{\linewidth}\raggedright
RUNSPEC
\end{minipage} & \begin{minipage}[b]{\linewidth}\raggedright
GRID
\end{minipage} & \begin{minipage}[b]{\linewidth}\raggedright
EDIT
\end{minipage} & \begin{minipage}[b]{\linewidth}\raggedright
PROPS
\end{minipage} & \begin{minipage}[b]{\linewidth}\raggedright
REGIONS
\end{minipage} & \begin{minipage}[b]{\linewidth}\raggedright
SOLUTION
\end{minipage} & \begin{minipage}[b]{\linewidth}\raggedright
SUMMARY
\end{minipage} & \begin{minipage}[b]{\linewidth}\raggedright
SCHEDULE
\end{minipage} \\
\midrule\noalign{}
\endhead
\bottomrule\noalign{}
\endlastfoot
\end{longtable}

\subsubsection{Description}\label{description-526}

LSWL defines the connate water saturation (``SWL''), for all the cells
in the model via an array, for when the Low Salt option has been
selected. The data is used to scale the water saturation in the low
salinity oil wet water-oil relative permeability saturation tables, as
well as the associated capillary pressure tables. The ENDSCALE keyword
in the RUNSPEC section should be activated to enable end-point scaling
and the use of this keyword. In addition, the Low Salt option should be
activated by the LOWSALT keyword in the RUNSPEC section.

This keyword is not supported by OPM Flow but would change the results
if supported so the simulation will be stopped.

\subsection{LSWLPC -- End-Point Scaling Grid Cell SWLPC (Low Salinity
and Oil
Wet)}\label{lswlpc-end-point-scaling-grid-cell-swlpc-low-salinity-and-oil-wet}

\begin{longtable}[]{@{}
  >{\raggedright\arraybackslash}p{(\linewidth - 14\tabcolsep) * \real{0.1250}}
  >{\raggedright\arraybackslash}p{(\linewidth - 14\tabcolsep) * \real{0.1250}}
  >{\raggedright\arraybackslash}p{(\linewidth - 14\tabcolsep) * \real{0.1250}}
  >{\raggedright\arraybackslash}p{(\linewidth - 14\tabcolsep) * \real{0.1250}}
  >{\raggedright\arraybackslash}p{(\linewidth - 14\tabcolsep) * \real{0.1250}}
  >{\raggedright\arraybackslash}p{(\linewidth - 14\tabcolsep) * \real{0.1250}}
  >{\raggedright\arraybackslash}p{(\linewidth - 14\tabcolsep) * \real{0.1250}}
  >{\raggedright\arraybackslash}p{(\linewidth - 14\tabcolsep) * \real{0.1250}}@{}}
\toprule\noalign{}
\begin{minipage}[b]{\linewidth}\raggedright
RUNSPEC
\end{minipage} & \begin{minipage}[b]{\linewidth}\raggedright
GRID
\end{minipage} & \begin{minipage}[b]{\linewidth}\raggedright
EDIT
\end{minipage} & \begin{minipage}[b]{\linewidth}\raggedright
PROPS
\end{minipage} & \begin{minipage}[b]{\linewidth}\raggedright
REGIONS
\end{minipage} & \begin{minipage}[b]{\linewidth}\raggedright
SOLUTION
\end{minipage} & \begin{minipage}[b]{\linewidth}\raggedright
SUMMARY
\end{minipage} & \begin{minipage}[b]{\linewidth}\raggedright
SCHEDULE
\end{minipage} \\
\midrule\noalign{}
\endhead
\bottomrule\noalign{}
\endlastfoot
\end{longtable}

\subsubsection{Description}\label{description-527}

LSWLPC defines the capillary pressure connate water saturation
(``SWLPC''), for all the cells in the model via an array, for when the
Low Salt option has been selected. The data is used to scale the water
saturation in the low salinity oil wet water-oil capillary pressure
tables. The ENDSCALE keyword in the RUNSPEC section should be activated
to enable end-point scaling and the use of this keyword. In addition,
the Low Salt option should be activated by the LOWSALT keyword in the
RUNSPEC section.

Note the keyword only applies the scaling to the capillary pressures
tables, unlike the LSWL keyword that applies the scaling to both the
capillary pressure and relative permeability tables.

This keyword is not supported by OPM Flow but would change the results
if supported so the simulation will be stopped.

\subsection{LSWU -- End-Point Scaling Grid Cell SWU (Low Salinity and
Oil
Wet)}\label{lswu-end-point-scaling-grid-cell-swu-low-salinity-and-oil-wet}

\begin{longtable}[]{@{}
  >{\raggedright\arraybackslash}p{(\linewidth - 14\tabcolsep) * \real{0.1250}}
  >{\raggedright\arraybackslash}p{(\linewidth - 14\tabcolsep) * \real{0.1250}}
  >{\raggedright\arraybackslash}p{(\linewidth - 14\tabcolsep) * \real{0.1250}}
  >{\raggedright\arraybackslash}p{(\linewidth - 14\tabcolsep) * \real{0.1250}}
  >{\raggedright\arraybackslash}p{(\linewidth - 14\tabcolsep) * \real{0.1250}}
  >{\raggedright\arraybackslash}p{(\linewidth - 14\tabcolsep) * \real{0.1250}}
  >{\raggedright\arraybackslash}p{(\linewidth - 14\tabcolsep) * \real{0.1250}}
  >{\raggedright\arraybackslash}p{(\linewidth - 14\tabcolsep) * \real{0.1250}}@{}}
\toprule\noalign{}
\begin{minipage}[b]{\linewidth}\raggedright
RUNSPEC
\end{minipage} & \begin{minipage}[b]{\linewidth}\raggedright
GRID
\end{minipage} & \begin{minipage}[b]{\linewidth}\raggedright
EDIT
\end{minipage} & \begin{minipage}[b]{\linewidth}\raggedright
PROPS
\end{minipage} & \begin{minipage}[b]{\linewidth}\raggedright
REGIONS
\end{minipage} & \begin{minipage}[b]{\linewidth}\raggedright
SOLUTION
\end{minipage} & \begin{minipage}[b]{\linewidth}\raggedright
SUMMARY
\end{minipage} & \begin{minipage}[b]{\linewidth}\raggedright
SCHEDULE
\end{minipage} \\
\midrule\noalign{}
\endhead
\bottomrule\noalign{}
\endlastfoot
\end{longtable}

\subsubsection{Description}\label{description-528}

LSWU defines the maximum water saturation(``SWU''), for all the cells in
the model via an array, for when the Low Salt option has been selected.
The data is used to scale the water saturation in the low salinity oil
wet water-oil relative permeability saturation tables, as well as the
associated capillary pressure tables. The ENDSCALE keyword in the
RUNSPEC section should be activated to enable end-point scaling and the
use of this keyword. In addition, the Low Salt option should be
activated by the LOWSALT keyword in the RUNSPEC section.

This keyword is not supported by OPM Flow but would change the results
if supported so the simulation will be stopped.

\subsection{LWKRO -- End-Point Scaling of Grid Cell Kro(Swl) (Low
Salinity and Water
Wet)}\label{lwkro-end-point-scaling-of-grid-cell-kroswl-low-salinity-and-water-wet}

\begin{longtable}[]{@{}
  >{\raggedright\arraybackslash}p{(\linewidth - 14\tabcolsep) * \real{0.1250}}
  >{\raggedright\arraybackslash}p{(\linewidth - 14\tabcolsep) * \real{0.1250}}
  >{\raggedright\arraybackslash}p{(\linewidth - 14\tabcolsep) * \real{0.1250}}
  >{\raggedright\arraybackslash}p{(\linewidth - 14\tabcolsep) * \real{0.1250}}
  >{\raggedright\arraybackslash}p{(\linewidth - 14\tabcolsep) * \real{0.1250}}
  >{\raggedright\arraybackslash}p{(\linewidth - 14\tabcolsep) * \real{0.1250}}
  >{\raggedright\arraybackslash}p{(\linewidth - 14\tabcolsep) * \real{0.1250}}
  >{\raggedright\arraybackslash}p{(\linewidth - 14\tabcolsep) * \real{0.1250}}@{}}
\toprule\noalign{}
\begin{minipage}[b]{\linewidth}\raggedright
RUNSPEC
\end{minipage} & \begin{minipage}[b]{\linewidth}\raggedright
GRID
\end{minipage} & \begin{minipage}[b]{\linewidth}\raggedright
EDIT
\end{minipage} & \begin{minipage}[b]{\linewidth}\raggedright
PROPS
\end{minipage} & \begin{minipage}[b]{\linewidth}\raggedright
REGIONS
\end{minipage} & \begin{minipage}[b]{\linewidth}\raggedright
SOLUTION
\end{minipage} & \begin{minipage}[b]{\linewidth}\raggedright
SUMMARY
\end{minipage} & \begin{minipage}[b]{\linewidth}\raggedright
SCHEDULE
\end{minipage} \\
\midrule\noalign{}
\endhead
\bottomrule\noalign{}
\endlastfoot
\end{longtable}

\subsubsection{Description}\label{description-529}

LWKRO defines the scaling parameter for the oil relative permeability
value at the connate water saturation (SWL), for all the cells in the
model via an array, for when the Low Salt and Surfactant Wettability
options have been selected. The data is used to scale the oil relative
permeability in the low salinity water wet oil relative permeability
saturation tables. The ENDSCALE keyword in the RUNSPEC section should be
activated to enable end-point scaling and the use of this keyword. In
addition the Low Salt option should be activated by the LOWSALT keyword
in the RUNSPEC section and the Surfactant Wettability option activated
by the SURFACT or SURFACTW keywords, which are also in the RUNSPEC
section.

This keyword is not supported by OPM Flow but would change the results
if supported so the simulation will be stopped.

\subsection{LWKRORG -- End-Point Scaling of Grid Cell Kro(Sgcr) (Low
Salinity and Water
Wet)}\label{lwkrorg-end-point-scaling-of-grid-cell-krosgcr-low-salinity-and-water-wet}

\begin{longtable}[]{@{}
  >{\raggedright\arraybackslash}p{(\linewidth - 14\tabcolsep) * \real{0.1250}}
  >{\raggedright\arraybackslash}p{(\linewidth - 14\tabcolsep) * \real{0.1250}}
  >{\raggedright\arraybackslash}p{(\linewidth - 14\tabcolsep) * \real{0.1250}}
  >{\raggedright\arraybackslash}p{(\linewidth - 14\tabcolsep) * \real{0.1250}}
  >{\raggedright\arraybackslash}p{(\linewidth - 14\tabcolsep) * \real{0.1250}}
  >{\raggedright\arraybackslash}p{(\linewidth - 14\tabcolsep) * \real{0.1250}}
  >{\raggedright\arraybackslash}p{(\linewidth - 14\tabcolsep) * \real{0.1250}}
  >{\raggedright\arraybackslash}p{(\linewidth - 14\tabcolsep) * \real{0.1250}}@{}}
\toprule\noalign{}
\begin{minipage}[b]{\linewidth}\raggedright
RUNSPEC
\end{minipage} & \begin{minipage}[b]{\linewidth}\raggedright
GRID
\end{minipage} & \begin{minipage}[b]{\linewidth}\raggedright
EDIT
\end{minipage} & \begin{minipage}[b]{\linewidth}\raggedright
PROPS
\end{minipage} & \begin{minipage}[b]{\linewidth}\raggedright
REGIONS
\end{minipage} & \begin{minipage}[b]{\linewidth}\raggedright
SOLUTION
\end{minipage} & \begin{minipage}[b]{\linewidth}\raggedright
SUMMARY
\end{minipage} & \begin{minipage}[b]{\linewidth}\raggedright
SCHEDULE
\end{minipage} \\
\midrule\noalign{}
\endhead
\bottomrule\noalign{}
\endlastfoot
\end{longtable}

\subsubsection{Description}\label{description-530}

LWKRORG defines the scaling parameter for the relative permeability of
oil at the critical gas saturation (SGCR), for all the cells in the
model via an array, for when the Low Salt and Surfactant Wettability
options have been selected. The data is used to scale the oil relative
permeability in the low salinity water wet oil relative permeability
saturation tables. The ENDSCALE keyword in the RUNSPEC section should be
activated to enable end-point scaling and the use of this keyword. In
addition the Low Salt option should be activated by the LOWSALT keyword
in the RUNSPEC section and the Surfactant Wettability option activated
by the SURFACT or SURFACTW keywords, which are also in the RUNSPEC
section.

This keyword is not supported by OPM Flow but would change the results
if supported so the simulation will be stopped.

\subsection{LWKRORW -- End-Point Scaling of Grid Cell Kro(Swcr) (Low
Salinity and Water
Wet)}\label{lwkrorw-end-point-scaling-of-grid-cell-kroswcr-low-salinity-and-water-wet}

\begin{longtable}[]{@{}
  >{\raggedright\arraybackslash}p{(\linewidth - 14\tabcolsep) * \real{0.1250}}
  >{\raggedright\arraybackslash}p{(\linewidth - 14\tabcolsep) * \real{0.1250}}
  >{\raggedright\arraybackslash}p{(\linewidth - 14\tabcolsep) * \real{0.1250}}
  >{\raggedright\arraybackslash}p{(\linewidth - 14\tabcolsep) * \real{0.1250}}
  >{\raggedright\arraybackslash}p{(\linewidth - 14\tabcolsep) * \real{0.1250}}
  >{\raggedright\arraybackslash}p{(\linewidth - 14\tabcolsep) * \real{0.1250}}
  >{\raggedright\arraybackslash}p{(\linewidth - 14\tabcolsep) * \real{0.1250}}
  >{\raggedright\arraybackslash}p{(\linewidth - 14\tabcolsep) * \real{0.1250}}@{}}
\toprule\noalign{}
\begin{minipage}[b]{\linewidth}\raggedright
RUNSPEC
\end{minipage} & \begin{minipage}[b]{\linewidth}\raggedright
GRID
\end{minipage} & \begin{minipage}[b]{\linewidth}\raggedright
EDIT
\end{minipage} & \begin{minipage}[b]{\linewidth}\raggedright
PROPS
\end{minipage} & \begin{minipage}[b]{\linewidth}\raggedright
REGIONS
\end{minipage} & \begin{minipage}[b]{\linewidth}\raggedright
SOLUTION
\end{minipage} & \begin{minipage}[b]{\linewidth}\raggedright
SUMMARY
\end{minipage} & \begin{minipage}[b]{\linewidth}\raggedright
SCHEDULE
\end{minipage} \\
\midrule\noalign{}
\endhead
\bottomrule\noalign{}
\endlastfoot
\end{longtable}

\subsubsection{Description}\label{description-531}

LWKRORW defines the scaling parameter for the relative permeability of
oil at the critical water saturation (SWCR), for all the cells in the
model via an array, for when the Low Salt and Surfactant Wettability
options have been selected. The data is used to scale the oil relative
permeability in the low salinity water wet oil relative permeability
saturation tables. The ENDSCALE keyword in the RUNSPEC section should be
activated to enable end-point scaling and the use of this keyword. In
addition the Low Salt option should be activated by the LOWSALT keyword
in the RUNSPEC section and the Surfactant Wettability option activated
by the SURFACT or SURFACTW keywords, which are also in the RUNSPEC
section.

This keyword is not supported by OPM Flow but would change the results
if supported so the simulation will be stopped.

\subsection{LWKRW -- End-Point Scaling of Grid Cell Krw(Sw =1.0) (Low
Salinity and Water
Wet)}\label{lwkrw-end-point-scaling-of-grid-cell-krwsw-1.0-low-salinity-and-water-wet}

\begin{longtable}[]{@{}
  >{\raggedright\arraybackslash}p{(\linewidth - 14\tabcolsep) * \real{0.1250}}
  >{\raggedright\arraybackslash}p{(\linewidth - 14\tabcolsep) * \real{0.1250}}
  >{\raggedright\arraybackslash}p{(\linewidth - 14\tabcolsep) * \real{0.1250}}
  >{\raggedright\arraybackslash}p{(\linewidth - 14\tabcolsep) * \real{0.1250}}
  >{\raggedright\arraybackslash}p{(\linewidth - 14\tabcolsep) * \real{0.1250}}
  >{\raggedright\arraybackslash}p{(\linewidth - 14\tabcolsep) * \real{0.1250}}
  >{\raggedright\arraybackslash}p{(\linewidth - 14\tabcolsep) * \real{0.1250}}
  >{\raggedright\arraybackslash}p{(\linewidth - 14\tabcolsep) * \real{0.1250}}@{}}
\toprule\noalign{}
\begin{minipage}[b]{\linewidth}\raggedright
RUNSPEC
\end{minipage} & \begin{minipage}[b]{\linewidth}\raggedright
GRID
\end{minipage} & \begin{minipage}[b]{\linewidth}\raggedright
EDIT
\end{minipage} & \begin{minipage}[b]{\linewidth}\raggedright
PROPS
\end{minipage} & \begin{minipage}[b]{\linewidth}\raggedright
REGIONS
\end{minipage} & \begin{minipage}[b]{\linewidth}\raggedright
SOLUTION
\end{minipage} & \begin{minipage}[b]{\linewidth}\raggedright
SUMMARY
\end{minipage} & \begin{minipage}[b]{\linewidth}\raggedright
SCHEDULE
\end{minipage} \\
\midrule\noalign{}
\endhead
\bottomrule\noalign{}
\endlastfoot
\end{longtable}

\subsubsection{Description}\label{description-532}

LWKRW defines the scaling parameter at the maximum water relative
permeability value (SWU), that is for Sw = 1.0, for all the cells in the
model via an array, for when the Low Salt and Surfactant Wettability
options have been selected. The data is used to scale the water relative
permeability in the low salinity water wet water relative permeability
saturation tables. The ENDSCALE keyword in the RUNSPEC section should be
activated to enable end-point scaling and the use of this keyword. In
addition the Low Salt option should be activated by the LOWSALT keyword
in the RUNSPEC section and the Surfactant Wettability option activated
by the SURFACT or SURFACTW keywords, which are also in the RUNSPEC
section.

This keyword is not supported by OPM Flow but would change the results
if supported so the simulation will be stopped.

\subsection{LWKRWR -- End-Point Scaling of Grid Cell KRWR(Sw=1.0) (Low
Salinity and Water
Wet)}\label{lwkrwr-end-point-scaling-of-grid-cell-krwrsw1.0-low-salinity-and-water-wet}

\begin{longtable}[]{@{}
  >{\raggedright\arraybackslash}p{(\linewidth - 14\tabcolsep) * \real{0.1250}}
  >{\raggedright\arraybackslash}p{(\linewidth - 14\tabcolsep) * \real{0.1250}}
  >{\raggedright\arraybackslash}p{(\linewidth - 14\tabcolsep) * \real{0.1250}}
  >{\raggedright\arraybackslash}p{(\linewidth - 14\tabcolsep) * \real{0.1250}}
  >{\raggedright\arraybackslash}p{(\linewidth - 14\tabcolsep) * \real{0.1250}}
  >{\raggedright\arraybackslash}p{(\linewidth - 14\tabcolsep) * \real{0.1250}}
  >{\raggedright\arraybackslash}p{(\linewidth - 14\tabcolsep) * \real{0.1250}}
  >{\raggedright\arraybackslash}p{(\linewidth - 14\tabcolsep) * \real{0.1250}}@{}}
\toprule\noalign{}
\begin{minipage}[b]{\linewidth}\raggedright
RUNSPEC
\end{minipage} & \begin{minipage}[b]{\linewidth}\raggedright
GRID
\end{minipage} & \begin{minipage}[b]{\linewidth}\raggedright
EDIT
\end{minipage} & \begin{minipage}[b]{\linewidth}\raggedright
PROPS
\end{minipage} & \begin{minipage}[b]{\linewidth}\raggedright
REGIONS
\end{minipage} & \begin{minipage}[b]{\linewidth}\raggedright
SOLUTION
\end{minipage} & \begin{minipage}[b]{\linewidth}\raggedright
SUMMARY
\end{minipage} & \begin{minipage}[b]{\linewidth}\raggedright
SCHEDULE
\end{minipage} \\
\midrule\noalign{}
\endhead
\bottomrule\noalign{}
\endlastfoot
\end{longtable}

\subsubsection{Description}\label{description-533}

LWKRWR defines the scaling parameter at the critical oil to water
saturation value (SOWCR), for the water relative permeability curve, for
all the cells in the model via an array, and for when the Low Salt and
Surfactant Wettability options have been selected. The data is used to
scale the water relative permeability in the low salinity water wet
water relative permeability saturation tables. The ENDSCALE keyword in
the RUNSPEC section should be activated to enable end-point scaling and
the use of this keyword. In addition the Low Salt option should be
activated by the LOWSALT keyword in the RUNSPEC section and the
Surfactant Wettability option activated by the SURFACT or SURFACTW
keywords, which are also in the RUNSPEC section.

This keyword is not supported by OPM Flow but would change the results
if supported so the simulation will be stopped.

\subsection{LWPCW -- End-Point Scaling of Grid Cell Water Capillary
Pressure (Low Salinity and Water
Wet)}\label{lwpcw-end-point-scaling-of-grid-cell-water-capillary-pressure-low-salinity-and-water-wet}

\begin{longtable}[]{@{}
  >{\raggedright\arraybackslash}p{(\linewidth - 14\tabcolsep) * \real{0.1250}}
  >{\raggedright\arraybackslash}p{(\linewidth - 14\tabcolsep) * \real{0.1250}}
  >{\raggedright\arraybackslash}p{(\linewidth - 14\tabcolsep) * \real{0.1250}}
  >{\raggedright\arraybackslash}p{(\linewidth - 14\tabcolsep) * \real{0.1250}}
  >{\raggedright\arraybackslash}p{(\linewidth - 14\tabcolsep) * \real{0.1250}}
  >{\raggedright\arraybackslash}p{(\linewidth - 14\tabcolsep) * \real{0.1250}}
  >{\raggedright\arraybackslash}p{(\linewidth - 14\tabcolsep) * \real{0.1250}}
  >{\raggedright\arraybackslash}p{(\linewidth - 14\tabcolsep) * \real{0.1250}}@{}}
\toprule\noalign{}
\begin{minipage}[b]{\linewidth}\raggedright
RUNSPEC
\end{minipage} & \begin{minipage}[b]{\linewidth}\raggedright
GRID
\end{minipage} & \begin{minipage}[b]{\linewidth}\raggedright
EDIT
\end{minipage} & \begin{minipage}[b]{\linewidth}\raggedright
PROPS
\end{minipage} & \begin{minipage}[b]{\linewidth}\raggedright
REGIONS
\end{minipage} & \begin{minipage}[b]{\linewidth}\raggedright
SOLUTION
\end{minipage} & \begin{minipage}[b]{\linewidth}\raggedright
SUMMARY
\end{minipage} & \begin{minipage}[b]{\linewidth}\raggedright
SCHEDULE
\end{minipage} \\
\midrule\noalign{}
\endhead
\bottomrule\noalign{}
\endlastfoot
\end{longtable}

\subsubsection{Description}\label{description-534}

LWPCW defines the maximum water-oil pressure values for all the cells in
the model via an array, for when the Low Salt and Surfactant Wettability
options have been selected. The ENDSCALE keyword in the RUNSPEC section
should be activated to enable end-point scaling and the use of this
keyword. In addition the Low Salt option should be activated by the
LOWSALT keyword in the RUNSPEC section and the Surfactant Wettability
option activated by the SURFACT or SURFACTW keywords, which are also in
the RUNSPEC section. The keyword re-scales the oil-water capillary
pressure in the low salinity water wet capillary saturation tables from
the cell's assigned saturation function by the grid block's LWPCW value.

The capillary pressure for a grid block is scaled by:

\begin{equation}\phantomsection\label{eq-8-65}{
{P}_{c} = {P}_{{c}_{\mathit{TABLE}}}\left(\frac{\mathit{HWPCW}}{{P}_{{c}_{\mathit{TABLE}-\mathit{MAX}}}}\right)
}\end{equation}

Where:

\({P}_{c}\) = the resulting high salinity water wet water capillary
pressure for a grid cell.

\(\mathit{HWPCW}\) = the maximum capillary pressure from the HWPCW array
for a given cell.

\({P}_{{c}_{\mathit{TABLE}}}\) = the capillary pressure in the high
salinity water wet water capillary pressure

table allocated to the grid block.

\({P}_{{c}_{\mathit{TABLE}-\mathit{MAX}}}\) = the maximum capillary
pressure in the high salinity water wet water

capillary pressure table allocated to the grid block (that is at the
connate

water saturation).

This keyword is not supported by OPM Flow but would change the results
if supported so the simulation will be stopped.

\subsection{LWSOGCR -- End-Point Scaling Grid Cell SOGCR (Low Salinity
and Water
Wet)}\label{lwsogcr-end-point-scaling-grid-cell-sogcr-low-salinity-and-water-wet}

\begin{longtable}[]{@{}
  >{\raggedright\arraybackslash}p{(\linewidth - 14\tabcolsep) * \real{0.1250}}
  >{\raggedright\arraybackslash}p{(\linewidth - 14\tabcolsep) * \real{0.1250}}
  >{\raggedright\arraybackslash}p{(\linewidth - 14\tabcolsep) * \real{0.1250}}
  >{\raggedright\arraybackslash}p{(\linewidth - 14\tabcolsep) * \real{0.1250}}
  >{\raggedright\arraybackslash}p{(\linewidth - 14\tabcolsep) * \real{0.1250}}
  >{\raggedright\arraybackslash}p{(\linewidth - 14\tabcolsep) * \real{0.1250}}
  >{\raggedright\arraybackslash}p{(\linewidth - 14\tabcolsep) * \real{0.1250}}
  >{\raggedright\arraybackslash}p{(\linewidth - 14\tabcolsep) * \real{0.1250}}@{}}
\toprule\noalign{}
\begin{minipage}[b]{\linewidth}\raggedright
RUNSPEC
\end{minipage} & \begin{minipage}[b]{\linewidth}\raggedright
GRID
\end{minipage} & \begin{minipage}[b]{\linewidth}\raggedright
EDIT
\end{minipage} & \begin{minipage}[b]{\linewidth}\raggedright
PROPS
\end{minipage} & \begin{minipage}[b]{\linewidth}\raggedright
REGIONS
\end{minipage} & \begin{minipage}[b]{\linewidth}\raggedright
SOLUTION
\end{minipage} & \begin{minipage}[b]{\linewidth}\raggedright
SUMMARY
\end{minipage} & \begin{minipage}[b]{\linewidth}\raggedright
SCHEDULE
\end{minipage} \\
\midrule\noalign{}
\endhead
\bottomrule\noalign{}
\endlastfoot
\end{longtable}

\subsubsection{Description}\label{description-535}

LWSOGCR defines the critical oil saturation with respect to gas
(``SOGCR''), for all the cells in the model via an array, for when the
Low Salt and Surfactant Wettability options have been selected. The data
is used to scale the oil saturation in the low salinity water wet
oil-gas relative permeability saturation tables. The ENDSCALE keyword in
the RUNSPEC section should be activated to enable end-point scaling and
the use of this keyword. In addition, the Low Salt option should be
activated by the LOWSALT keyword in the RUNSPEC section and the
Surfactant Wettability option activated by the SURFACT or SURFACTW
keywords, which are also in the RUNSPEC section.

This keyword is not supported by OPM Flow but would change the results
if supported so the simulation will be stopped.

\subsection{LWSOWCR -- End-Point Scaling Grid Cell SOWCR (Low Salinity
and Water
Wet)}\label{lwsowcr-end-point-scaling-grid-cell-sowcr-low-salinity-and-water-wet}

\begin{longtable}[]{@{}
  >{\raggedright\arraybackslash}p{(\linewidth - 14\tabcolsep) * \real{0.1250}}
  >{\raggedright\arraybackslash}p{(\linewidth - 14\tabcolsep) * \real{0.1250}}
  >{\raggedright\arraybackslash}p{(\linewidth - 14\tabcolsep) * \real{0.1250}}
  >{\raggedright\arraybackslash}p{(\linewidth - 14\tabcolsep) * \real{0.1250}}
  >{\raggedright\arraybackslash}p{(\linewidth - 14\tabcolsep) * \real{0.1250}}
  >{\raggedright\arraybackslash}p{(\linewidth - 14\tabcolsep) * \real{0.1250}}
  >{\raggedright\arraybackslash}p{(\linewidth - 14\tabcolsep) * \real{0.1250}}
  >{\raggedright\arraybackslash}p{(\linewidth - 14\tabcolsep) * \real{0.1250}}@{}}
\toprule\noalign{}
\begin{minipage}[b]{\linewidth}\raggedright
RUNSPEC
\end{minipage} & \begin{minipage}[b]{\linewidth}\raggedright
GRID
\end{minipage} & \begin{minipage}[b]{\linewidth}\raggedright
EDIT
\end{minipage} & \begin{minipage}[b]{\linewidth}\raggedright
PROPS
\end{minipage} & \begin{minipage}[b]{\linewidth}\raggedright
REGIONS
\end{minipage} & \begin{minipage}[b]{\linewidth}\raggedright
SOLUTION
\end{minipage} & \begin{minipage}[b]{\linewidth}\raggedright
SUMMARY
\end{minipage} & \begin{minipage}[b]{\linewidth}\raggedright
SCHEDULE
\end{minipage} \\
\midrule\noalign{}
\endhead
\bottomrule\noalign{}
\endlastfoot
\end{longtable}

\subsubsection{Description}\label{description-536}

LWSOWCR defines the critical oil saturation with respect to water
(``SOWCR''), for all the cells in the model via an array, for when the
Low Salt and Surfactant Wettability options have been selected. The data
is used to scale the oil saturation in the low salinity water wet
water-oil relative permeability saturation tables. The ENDSCALE keyword
in the RUNSPEC section should be activated to enable end-point scaling
and the use of this keyword. In addition, the Low Salt option should be
activated by the LOWSALT keyword in the RUNSPEC section and the
Surfactant Wettability option activated by the SURFACT or SURFACTW
keywords, which are also in the RUNSPEC section.

This keyword is not supported by OPM Flow but would change the results
if supported so the simulation will be stopped.

\subsection{LWSWCR -- End-Point Scaling Grid Cell SWCR (Low Salinity and
Water
Wet)}\label{lwswcr-end-point-scaling-grid-cell-swcr-low-salinity-and-water-wet}

\begin{longtable}[]{@{}
  >{\raggedright\arraybackslash}p{(\linewidth - 14\tabcolsep) * \real{0.1250}}
  >{\raggedright\arraybackslash}p{(\linewidth - 14\tabcolsep) * \real{0.1250}}
  >{\raggedright\arraybackslash}p{(\linewidth - 14\tabcolsep) * \real{0.1250}}
  >{\raggedright\arraybackslash}p{(\linewidth - 14\tabcolsep) * \real{0.1250}}
  >{\raggedright\arraybackslash}p{(\linewidth - 14\tabcolsep) * \real{0.1250}}
  >{\raggedright\arraybackslash}p{(\linewidth - 14\tabcolsep) * \real{0.1250}}
  >{\raggedright\arraybackslash}p{(\linewidth - 14\tabcolsep) * \real{0.1250}}
  >{\raggedright\arraybackslash}p{(\linewidth - 14\tabcolsep) * \real{0.1250}}@{}}
\toprule\noalign{}
\begin{minipage}[b]{\linewidth}\raggedright
RUNSPEC
\end{minipage} & \begin{minipage}[b]{\linewidth}\raggedright
GRID
\end{minipage} & \begin{minipage}[b]{\linewidth}\raggedright
EDIT
\end{minipage} & \begin{minipage}[b]{\linewidth}\raggedright
PROPS
\end{minipage} & \begin{minipage}[b]{\linewidth}\raggedright
REGIONS
\end{minipage} & \begin{minipage}[b]{\linewidth}\raggedright
SOLUTION
\end{minipage} & \begin{minipage}[b]{\linewidth}\raggedright
SUMMARY
\end{minipage} & \begin{minipage}[b]{\linewidth}\raggedright
SCHEDULE
\end{minipage} \\
\midrule\noalign{}
\endhead
\bottomrule\noalign{}
\endlastfoot
\end{longtable}

\subsubsection{Description}\label{description-537}

LWSWCR defines the critical water saturation (``SWCR''), for all the
cells in the model via an array, for when the Low Salt and Surfactant
Wettability options have been selected. The data is used to scale the
water saturation in the low salinity water wet water-oil relative
permeability saturation tables. The ENDSCALE keyword in the RUNSPEC
section should be activated to enable end-point scaling and the use of
this keyword. In addition, the Low Salt option should be activated by
the LOWSALT keyword in the RUNSPEC section and the Surfactant
Wettability option activated by the SURFACT or SURFACTW keywords, which
are also in the RUNSPEC section.

This keyword is not supported by OPM Flow but would change the results
if supported so the simulation will be stopped.

\subsection{LWSWL -- End-Point Scaling Grid Cell SWL (Low Salinity and
Water
Wet)}\label{lwswl-end-point-scaling-grid-cell-swl-low-salinity-and-water-wet}

\begin{longtable}[]{@{}
  >{\raggedright\arraybackslash}p{(\linewidth - 14\tabcolsep) * \real{0.1250}}
  >{\raggedright\arraybackslash}p{(\linewidth - 14\tabcolsep) * \real{0.1250}}
  >{\raggedright\arraybackslash}p{(\linewidth - 14\tabcolsep) * \real{0.1250}}
  >{\raggedright\arraybackslash}p{(\linewidth - 14\tabcolsep) * \real{0.1250}}
  >{\raggedright\arraybackslash}p{(\linewidth - 14\tabcolsep) * \real{0.1250}}
  >{\raggedright\arraybackslash}p{(\linewidth - 14\tabcolsep) * \real{0.1250}}
  >{\raggedright\arraybackslash}p{(\linewidth - 14\tabcolsep) * \real{0.1250}}
  >{\raggedright\arraybackslash}p{(\linewidth - 14\tabcolsep) * \real{0.1250}}@{}}
\toprule\noalign{}
\begin{minipage}[b]{\linewidth}\raggedright
RUNSPEC
\end{minipage} & \begin{minipage}[b]{\linewidth}\raggedright
GRID
\end{minipage} & \begin{minipage}[b]{\linewidth}\raggedright
EDIT
\end{minipage} & \begin{minipage}[b]{\linewidth}\raggedright
PROPS
\end{minipage} & \begin{minipage}[b]{\linewidth}\raggedright
REGIONS
\end{minipage} & \begin{minipage}[b]{\linewidth}\raggedright
SOLUTION
\end{minipage} & \begin{minipage}[b]{\linewidth}\raggedright
SUMMARY
\end{minipage} & \begin{minipage}[b]{\linewidth}\raggedright
SCHEDULE
\end{minipage} \\
\midrule\noalign{}
\endhead
\bottomrule\noalign{}
\endlastfoot
\end{longtable}

\subsubsection{Description}\label{description-538}

LWSWL defines the connate water saturation (``SWL''), for all the cells
in the model via an array, for when the Low Salt and Surfactant
Wettability options have been selected. The data is used to scale the
water saturation in the low salinity water wet water-oil relative
permeability saturation tables, as well as the associated capillary
pressure tables. The ENDSCALE keyword in the RUNSPEC section should be
activated to enable end-point scaling and the use of this keyword. In
addition, the Low Salt option should be activated by the LOWSALT keyword
in the RUNSPEC section and the Surfactant Wettability option activated
by the SURFACT or SURFACTW keywords, which are also in the RUNSPEC
section.

This keyword is not supported by OPM Flow but would change the results
if supported so the simulation will be stopped.

\subsection{LWSWLPC -- End-Point Scaling Grid Cell SWLPC (Low Salinity
and Water
Wet)}\label{lwswlpc-end-point-scaling-grid-cell-swlpc-low-salinity-and-water-wet}

\begin{longtable}[]{@{}
  >{\raggedright\arraybackslash}p{(\linewidth - 14\tabcolsep) * \real{0.1250}}
  >{\raggedright\arraybackslash}p{(\linewidth - 14\tabcolsep) * \real{0.1250}}
  >{\raggedright\arraybackslash}p{(\linewidth - 14\tabcolsep) * \real{0.1250}}
  >{\raggedright\arraybackslash}p{(\linewidth - 14\tabcolsep) * \real{0.1250}}
  >{\raggedright\arraybackslash}p{(\linewidth - 14\tabcolsep) * \real{0.1250}}
  >{\raggedright\arraybackslash}p{(\linewidth - 14\tabcolsep) * \real{0.1250}}
  >{\raggedright\arraybackslash}p{(\linewidth - 14\tabcolsep) * \real{0.1250}}
  >{\raggedright\arraybackslash}p{(\linewidth - 14\tabcolsep) * \real{0.1250}}@{}}
\toprule\noalign{}
\begin{minipage}[b]{\linewidth}\raggedright
RUNSPEC
\end{minipage} & \begin{minipage}[b]{\linewidth}\raggedright
GRID
\end{minipage} & \begin{minipage}[b]{\linewidth}\raggedright
EDIT
\end{minipage} & \begin{minipage}[b]{\linewidth}\raggedright
PROPS
\end{minipage} & \begin{minipage}[b]{\linewidth}\raggedright
REGIONS
\end{minipage} & \begin{minipage}[b]{\linewidth}\raggedright
SOLUTION
\end{minipage} & \begin{minipage}[b]{\linewidth}\raggedright
SUMMARY
\end{minipage} & \begin{minipage}[b]{\linewidth}\raggedright
SCHEDULE
\end{minipage} \\
\midrule\noalign{}
\endhead
\bottomrule\noalign{}
\endlastfoot
\end{longtable}

\subsubsection{Description}\label{description-539}

LWSWLPC defines the capillary pressure connate water saturation
(``SWLPC''), for all the cells in the model via an array, for when the
Low Salt and Surfactant Wettability options have been selected. The data
is used to scale the water saturation in the low salinity water wet
water-oil capillary pressure tables. The ENDSCALE keyword in the RUNSPEC
section should be activated to enable end-point scaling and the use of
this keyword. In addition, the Low Salt option should be activated by
the LOWSALT keyword in the RUNSPEC section and the Surfactant
Wettability option activated by the SURFACT or SURFACTW keywords, which
are also in the RUNSPEC section.

Note the keyword only applies the scaling to the capillary pressures
tables, unlike the LWSWL keyword that applies the scaling to both the
capillary pressure and relative permeability tables.

This keyword is not supported by OPM Flow but would change the results
if supported so the simulation will be stopped.

\subsection{LWSWU -- End-Point Scaling Grid Cell SWU (Low Salinity and
Water
Wet)}\label{lwswu-end-point-scaling-grid-cell-swu-low-salinity-and-water-wet}

\begin{longtable}[]{@{}
  >{\raggedright\arraybackslash}p{(\linewidth - 14\tabcolsep) * \real{0.1250}}
  >{\raggedright\arraybackslash}p{(\linewidth - 14\tabcolsep) * \real{0.1250}}
  >{\raggedright\arraybackslash}p{(\linewidth - 14\tabcolsep) * \real{0.1250}}
  >{\raggedright\arraybackslash}p{(\linewidth - 14\tabcolsep) * \real{0.1250}}
  >{\raggedright\arraybackslash}p{(\linewidth - 14\tabcolsep) * \real{0.1250}}
  >{\raggedright\arraybackslash}p{(\linewidth - 14\tabcolsep) * \real{0.1250}}
  >{\raggedright\arraybackslash}p{(\linewidth - 14\tabcolsep) * \real{0.1250}}
  >{\raggedright\arraybackslash}p{(\linewidth - 14\tabcolsep) * \real{0.1250}}@{}}
\toprule\noalign{}
\begin{minipage}[b]{\linewidth}\raggedright
RUNSPEC
\end{minipage} & \begin{minipage}[b]{\linewidth}\raggedright
GRID
\end{minipage} & \begin{minipage}[b]{\linewidth}\raggedright
EDIT
\end{minipage} & \begin{minipage}[b]{\linewidth}\raggedright
PROPS
\end{minipage} & \begin{minipage}[b]{\linewidth}\raggedright
REGIONS
\end{minipage} & \begin{minipage}[b]{\linewidth}\raggedright
SOLUTION
\end{minipage} & \begin{minipage}[b]{\linewidth}\raggedright
SUMMARY
\end{minipage} & \begin{minipage}[b]{\linewidth}\raggedright
SCHEDULE
\end{minipage} \\
\midrule\noalign{}
\endhead
\bottomrule\noalign{}
\endlastfoot
\end{longtable}

\subsubsection{Description}\label{description-540}

LWSWU defines the maximum water saturation(``SWU''), for all the cells
in the model via an array, for when the Low Salt and Surfactant
Wettability options have been selected. The data is used to scale the
water saturation in the low salinity water wet water-oil relative
permeability saturation tables, as well as the associated capillary
pressure tables. The ENDSCALE keyword in the RUNSPEC section should be
activated to enable end-point scaling and the use of this keyword. In
addition, the Low Salt option should be activated by the LOWSALT keyword
in the RUNSPEC section and the Surfactant Wettability option activated
by the SURFACT or SURFACTW keywords, which are also in the RUNSPEC
section.

This keyword is not supported by OPM Flow but would change the results
if supported so the simulation will be stopped.

\subsection{MASSFLOW -- Define River Mass Flow versus Time
Tables}\label{massflow-define-river-mass-flow-versus-time-tables}

\begin{longtable}[]{@{}
  >{\raggedright\arraybackslash}p{(\linewidth - 14\tabcolsep) * \real{0.1250}}
  >{\raggedright\arraybackslash}p{(\linewidth - 14\tabcolsep) * \real{0.1250}}
  >{\raggedright\arraybackslash}p{(\linewidth - 14\tabcolsep) * \real{0.1250}}
  >{\raggedright\arraybackslash}p{(\linewidth - 14\tabcolsep) * \real{0.1250}}
  >{\raggedright\arraybackslash}p{(\linewidth - 14\tabcolsep) * \real{0.1250}}
  >{\raggedright\arraybackslash}p{(\linewidth - 14\tabcolsep) * \real{0.1250}}
  >{\raggedright\arraybackslash}p{(\linewidth - 14\tabcolsep) * \real{0.1250}}
  >{\raggedright\arraybackslash}p{(\linewidth - 14\tabcolsep) * \real{0.1250}}@{}}
\toprule\noalign{}
\begin{minipage}[b]{\linewidth}\raggedright
RUNSPEC
\end{minipage} & \begin{minipage}[b]{\linewidth}\raggedright
GRID
\end{minipage} & \begin{minipage}[b]{\linewidth}\raggedright
EDIT
\end{minipage} & \begin{minipage}[b]{\linewidth}\raggedright
PROPS
\end{minipage} & \begin{minipage}[b]{\linewidth}\raggedright
REGIONS
\end{minipage} & \begin{minipage}[b]{\linewidth}\raggedright
SOLUTION
\end{minipage} & \begin{minipage}[b]{\linewidth}\raggedright
SUMMARY
\end{minipage} & \begin{minipage}[b]{\linewidth}\raggedright
SCHEDULE
\end{minipage} \\
\midrule\noalign{}
\endhead
\bottomrule\noalign{}
\endlastfoot
\end{longtable}

\subsubsection{Description}\label{description-541}

The MASSFLOW keyword defines the upstream river mass flow versus time
tables for rivers, for when the River option has been activated via the
RIVRDIMS keyword in the RUNSPEC section.

This keyword is not supported by OPM Flow but would change the results
if supported so the simulation will be stopped.

\subsection{MAXVALUE -- Sets a Maximum Value for an Array
Element}\label{maxvalue-sets-a-maximum-value-for-an-array-element-2}

The MAXVALUE keyword sets a maximum value for the specified array or
part of an array. The constant can be real or integer depending on the
array type; however, the arrays that can be operated on is dependent on
which section the MAXVALUE keyword is being used.

See MAXVALUE -- Sets a Maximum Value for an Array Element in the GRID
section for a full description.

\subsection{MESSAGE -- Output User
Message}\label{message-output-user-message-4}

The MESSAGE keyword outputs a user message to the terminal, as well as
to the print (\emph{.PRT) and debug (}.DBG) files. Note this is
different to the MESSAGES keyword, that defines OPM Flows message print
limits and stop limits generated by the simulator.

See MESSAGE -- Output User Message in the GLOBAL section for a full
description.

\subsection{MESSAGES -- Define Message Print Limits and Stop
Limits}\label{messages-define-message-print-limits-and-stop-limits-4}

The MESSAGES keyword defines the print and stops levels for various
messages. The ``print limits'' set the maximum number of messages that
will be printed, after which no more messages will be printed and the
``stop limits'' terminate the run when these limits are exceeded. There
are six levels of message that increase in severity from informative all
the way to programming errors, as outlined in Table 4.5.

See MESSAGES -- Define Message Print Limits and Stop Limits in the
GLOBAL section for a full description.

\subsection{MICPPARA -- Define Microbially Induced Calcite Precipitation
Parameters
(Retired)}\label{micppara-define-microbially-induced-calcite-precipitation-parameters-retired}

\begin{longtable}[]{@{}
  >{\raggedright\arraybackslash}p{(\linewidth - 14\tabcolsep) * \real{0.1250}}
  >{\raggedright\arraybackslash}p{(\linewidth - 14\tabcolsep) * \real{0.1250}}
  >{\raggedright\arraybackslash}p{(\linewidth - 14\tabcolsep) * \real{0.1250}}
  >{\raggedright\arraybackslash}p{(\linewidth - 14\tabcolsep) * \real{0.1250}}
  >{\raggedright\arraybackslash}p{(\linewidth - 14\tabcolsep) * \real{0.1250}}
  >{\raggedright\arraybackslash}p{(\linewidth - 14\tabcolsep) * \real{0.1250}}
  >{\raggedright\arraybackslash}p{(\linewidth - 14\tabcolsep) * \real{0.1250}}
  >{\raggedright\arraybackslash}p{(\linewidth - 14\tabcolsep) * \real{0.1250}}@{}}
\toprule\noalign{}
\begin{minipage}[b]{\linewidth}\raggedright
RUNSPEC
\end{minipage} & \begin{minipage}[b]{\linewidth}\raggedright
GRID
\end{minipage} & \begin{minipage}[b]{\linewidth}\raggedright
EDIT
\end{minipage} & \begin{minipage}[b]{\linewidth}\raggedright
PROPS
\end{minipage} & \begin{minipage}[b]{\linewidth}\raggedright
REGIONS
\end{minipage} & \begin{minipage}[b]{\linewidth}\raggedright
SOLUTION
\end{minipage} & \begin{minipage}[b]{\linewidth}\raggedright
SUMMARY
\end{minipage} & \begin{minipage}[b]{\linewidth}\raggedright
SCHEDULE
\end{minipage} \\
\midrule\noalign{}
\endhead
\bottomrule\noalign{}
\endlastfoot
\end{longtable}

\subsubsection{Description}\label{description-542}

The MICPPARA keyword used to define the model parameters for when the
MICP model had been activated via the MICP keyword in the RUNSPEC
section.

The parameters in this keyword are now defined by the
\hyperref[REF_HEADING_KEYWORD_BIOFPARA]{BIOFPARA} and PERMFACT keywords.

\subsection{MINVALUE -- Set a Minimum Value for an Array
Element}\label{minvalue-set-a-minimum-value-for-an-array-element-2}

The MINVALUE keyword sets a minimum value for the specified array or
part of an array. The constant can be real or integer depending on the
array type; however, the arrays that can be operated on is dependent on
which section the MINVALUE keyword is being used.

See MINVALUE -- Set a Minimum Value for an Array Element in the GRID
section for a full description.

\subsection{MISC -- Define Solvent Miscibility-Immiscibility Transform
Functions}\label{misc-define-solvent-miscibility-immiscibility-transform-functions}

\begin{longtable}[]{@{}
  >{\raggedright\arraybackslash}p{(\linewidth - 14\tabcolsep) * \real{0.1250}}
  >{\raggedright\arraybackslash}p{(\linewidth - 14\tabcolsep) * \real{0.1250}}
  >{\raggedright\arraybackslash}p{(\linewidth - 14\tabcolsep) * \real{0.1250}}
  >{\raggedright\arraybackslash}p{(\linewidth - 14\tabcolsep) * \real{0.1250}}
  >{\raggedright\arraybackslash}p{(\linewidth - 14\tabcolsep) * \real{0.1250}}
  >{\raggedright\arraybackslash}p{(\linewidth - 14\tabcolsep) * \real{0.1250}}
  >{\raggedright\arraybackslash}p{(\linewidth - 14\tabcolsep) * \real{0.1250}}
  >{\raggedright\arraybackslash}p{(\linewidth - 14\tabcolsep) * \real{0.1250}}@{}}
\toprule\noalign{}
\begin{minipage}[b]{\linewidth}\raggedright
RUNSPEC
\end{minipage} & \begin{minipage}[b]{\linewidth}\raggedright
GRID
\end{minipage} & \begin{minipage}[b]{\linewidth}\raggedright
EDIT
\end{minipage} & \begin{minipage}[b]{\linewidth}\raggedright
PROPS
\end{minipage} & \begin{minipage}[b]{\linewidth}\raggedright
REGIONS
\end{minipage} & \begin{minipage}[b]{\linewidth}\raggedright
SOLUTION
\end{minipage} & \begin{minipage}[b]{\linewidth}\raggedright
SUMMARY
\end{minipage} & \begin{minipage}[b]{\linewidth}\raggedright
SCHEDULE
\end{minipage} \\
\midrule\noalign{}
\endhead
\bottomrule\noalign{}
\endlastfoot
\end{longtable}

\subsubsection{Description}\label{description-543}

MISC defines the transformation between the miscible and immiscible
relative permeability models, for when the MISCIBLE and SOLVENT keywords
in the RUNSPEC section have been activated. The keyword can only be used
with the MISCIBLE option and for when the oil, water, gas and solvent
phases are active in the model.

\begin{longtable}[]{@{}
  >{\raggedright\arraybackslash}p{(\linewidth - 6\tabcolsep) * \real{0.2500}}
  >{\raggedright\arraybackslash}p{(\linewidth - 6\tabcolsep) * \real{0.2500}}
  >{\raggedright\arraybackslash}p{(\linewidth - 6\tabcolsep) * \real{0.2500}}
  >{\raggedright\arraybackslash}p{(\linewidth - 6\tabcolsep) * \real{0.2500}}@{}}
\toprule\noalign{}
\begin{minipage}[b]{\linewidth}\raggedright
No.
\end{minipage} & \begin{minipage}[b]{\linewidth}\raggedright
Name
\end{minipage} & \begin{minipage}[b]{\linewidth}\raggedright
Description
\end{minipage} & \begin{minipage}[b]{\linewidth}\raggedright
Default
\end{minipage} \\
\midrule\noalign{}
\endhead
\bottomrule\noalign{}
\endlastfoot
Field & Metric & Laboratory & \\
1 & SSOL & A columnar vector of real monotonically increasing down the
column values starting from zero and terminating at one, that defines
the solvent fraction with respect to the solvent and gas saturation, and
is defined by: \(\frac{{S}_{s}}{\left({S}_{g} + {S}_{s}\right)}\) Where
Sg is the gas saturation and Ss is the solvent saturation. Note that the
first entry in the columnar vector should be zero and the last entry
should be one to fully define the solvent fraction range. & None \\
dimensionless & dimensionless & dimensionless & \\
2 & MISC & A columnar vector of real equal or increasing down the column
values that are greater than or equal to zero and less then one, that
define the corresponding miscibility for the corresponding solvent
fraction SSOL. The first entry in the columnar vector should be zero and
the last entry should be one to fully define the miscible-immiscible
relationship. & None \\
dimensionless & dimensionless & dimensionless & \\
Notes: & & & \\
\end{longtable}

\emph{Table 8.88: MISC Keyword Description}

\subsubsection{Example}\label{example-299}

\begin{verbatim}
--
--       SOLVENT MISCIBILITY-IMMISCIBLITY TRANSFORM TABLE
--
MISC
--       SSOL       MISC
--       FRAC       FRAC
--       -------    --------
          0.0000     0.0000
          0.2000     0.2500
          0.5000     0.7500
          1.0000     1.0000                                / TABLE NO. 01

--       SSOL       MISC
--       FRAC       FRAC
--       -------    --------
          0.0000     0.0000
          0.3000     0.2500
          0.6000     1.0000
          1.0000     1.0000                                / TABLE NO. 02
\end{verbatim}

The above example defines two solvent miscible-immiscible transform
tables assuming NTMISC equals two and NSMISC is greater than or equal to
four on the MISCIBLE keyword in the RUNSPEC section.

\subsection{MLANG -- Define Langmuir Maximum Gas Concentration for All
Grid
Cells}\label{mlang-define-langmuir-maximum-gas-concentration-for-all-grid-cells}

\begin{longtable}[]{@{}
  >{\raggedright\arraybackslash}p{(\linewidth - 14\tabcolsep) * \real{0.1250}}
  >{\raggedright\arraybackslash}p{(\linewidth - 14\tabcolsep) * \real{0.1250}}
  >{\raggedright\arraybackslash}p{(\linewidth - 14\tabcolsep) * \real{0.1250}}
  >{\raggedright\arraybackslash}p{(\linewidth - 14\tabcolsep) * \real{0.1250}}
  >{\raggedright\arraybackslash}p{(\linewidth - 14\tabcolsep) * \real{0.1250}}
  >{\raggedright\arraybackslash}p{(\linewidth - 14\tabcolsep) * \real{0.1250}}
  >{\raggedright\arraybackslash}p{(\linewidth - 14\tabcolsep) * \real{0.1250}}
  >{\raggedright\arraybackslash}p{(\linewidth - 14\tabcolsep) * \real{0.1250}}@{}}
\toprule\noalign{}
\begin{minipage}[b]{\linewidth}\raggedright
RUNSPEC
\end{minipage} & \begin{minipage}[b]{\linewidth}\raggedright
GRID
\end{minipage} & \begin{minipage}[b]{\linewidth}\raggedright
EDIT
\end{minipage} & \begin{minipage}[b]{\linewidth}\raggedright
PROPS
\end{minipage} & \begin{minipage}[b]{\linewidth}\raggedright
REGIONS
\end{minipage} & \begin{minipage}[b]{\linewidth}\raggedright
SOLUTION
\end{minipage} & \begin{minipage}[b]{\linewidth}\raggedright
SUMMARY
\end{minipage} & \begin{minipage}[b]{\linewidth}\raggedright
SCHEDULE
\end{minipage} \\
\midrule\noalign{}
\endhead
\bottomrule\noalign{}
\endlastfoot
\end{longtable}

\subsubsection{Description}\label{description-544}

This keyword, MLANG, defines the coal bed methane Langmuir Adsorption
{[}Langmuir, Irving (June 1918). ``The Adsorption of Gases on Plane
Surface of Glass, Mica and Platinum''. The Research Laboratory of the
General Electric Company. 40 (9): 1361--1402. doi:10.1021/ja02242a004{]}
maximum gas concentration for each grid cell used to scale the Langmuir
isotherm table allocated to the cell, for when the Coal Bed Methane
option has been activated via the COAL keyword in the RUNSPEC section.
See the LANGMUIR keyword in the PROPS section for specifying the
Langmuir tables for the model.

Note that if he Dual Porosity model has been activated by either the
DUALPORO or the DUALPERM keywords in the RUNSPEC section, then MLANG
applies to only the matrix grid block.

This keyword is not supported by OPM Flow but would change the results
if supported so the simulation will be stopped.

\subsection{MLANGSLV -- Define Langmuir Maximum Solvent Concentration
for All Grid
Cells}\label{mlangslv-define-langmuir-maximum-solvent-concentration-for-all-grid-cells}

\begin{longtable}[]{@{}
  >{\raggedright\arraybackslash}p{(\linewidth - 14\tabcolsep) * \real{0.1250}}
  >{\raggedright\arraybackslash}p{(\linewidth - 14\tabcolsep) * \real{0.1250}}
  >{\raggedright\arraybackslash}p{(\linewidth - 14\tabcolsep) * \real{0.1250}}
  >{\raggedright\arraybackslash}p{(\linewidth - 14\tabcolsep) * \real{0.1250}}
  >{\raggedright\arraybackslash}p{(\linewidth - 14\tabcolsep) * \real{0.1250}}
  >{\raggedright\arraybackslash}p{(\linewidth - 14\tabcolsep) * \real{0.1250}}
  >{\raggedright\arraybackslash}p{(\linewidth - 14\tabcolsep) * \real{0.1250}}
  >{\raggedright\arraybackslash}p{(\linewidth - 14\tabcolsep) * \real{0.1250}}@{}}
\toprule\noalign{}
\begin{minipage}[b]{\linewidth}\raggedright
RUNSPEC
\end{minipage} & \begin{minipage}[b]{\linewidth}\raggedright
GRID
\end{minipage} & \begin{minipage}[b]{\linewidth}\raggedright
EDIT
\end{minipage} & \begin{minipage}[b]{\linewidth}\raggedright
PROPS
\end{minipage} & \begin{minipage}[b]{\linewidth}\raggedright
REGIONS
\end{minipage} & \begin{minipage}[b]{\linewidth}\raggedright
SOLUTION
\end{minipage} & \begin{minipage}[b]{\linewidth}\raggedright
SUMMARY
\end{minipage} & \begin{minipage}[b]{\linewidth}\raggedright
SCHEDULE
\end{minipage} \\
\midrule\noalign{}
\endhead
\bottomrule\noalign{}
\endlastfoot
\end{longtable}

\subsubsection{Description}\label{description-545}

This keyword, MLANGSLV, defines the coal bed methane Langmuir Adsorption
{[}Langmuir, Irving (June 1918). ``The Adsorption of Gases on Plane
Surface of Glass, Mica and Platinum''. The Research Laboratory of the
General Electric Company. 40 (9): 1361--1402. doi:10.1021/ja02242a004{]}
maximum solvent concentration for each grid cell used to scale the
Langmuir isotherm solvent table allocated to the cell, for when the Coal
Bed Methane option has been activated via the COAL keyword in the
RUNSPEC section. In addition, the Solvent phase must have been declared
by the SOLVENT keyword in the RUNSPEC section. See the COALNUM keyword
in the GRID section for allocating the Langmuir solvent tables to the
grid blocks, and also the LANGMUIR keyword in the PROPS section for
defining the Langmuir Adsorption Isotherm tables. Keywords COALADS and
COALPP, also in the PROPS section, are used to specify the relative
adsorption data in runs containing the solvent phase.

Note that if he Dual Porosity model has been activated by either the
DUALPORO or the DUALPERM keywords in the RUNSPEC section, then MLANGSLV
applies to only the matrix grid block.

This keyword is not supported by OPM Flow but would change the results
if supported so the simulation will be stopped.

\subsection{MSFN -- Miscible Normalized Relative Permeability
Tables}\label{msfn-miscible-normalized-relative-permeability-tables}

\begin{longtable}[]{@{}
  >{\raggedright\arraybackslash}p{(\linewidth - 14\tabcolsep) * \real{0.1250}}
  >{\raggedright\arraybackslash}p{(\linewidth - 14\tabcolsep) * \real{0.1250}}
  >{\raggedright\arraybackslash}p{(\linewidth - 14\tabcolsep) * \real{0.1250}}
  >{\raggedright\arraybackslash}p{(\linewidth - 14\tabcolsep) * \real{0.1250}}
  >{\raggedright\arraybackslash}p{(\linewidth - 14\tabcolsep) * \real{0.1250}}
  >{\raggedright\arraybackslash}p{(\linewidth - 14\tabcolsep) * \real{0.1250}}
  >{\raggedright\arraybackslash}p{(\linewidth - 14\tabcolsep) * \real{0.1250}}
  >{\raggedright\arraybackslash}p{(\linewidth - 14\tabcolsep) * \real{0.1250}}@{}}
\toprule\noalign{}
\begin{minipage}[b]{\linewidth}\raggedright
RUNSPEC
\end{minipage} & \begin{minipage}[b]{\linewidth}\raggedright
GRID
\end{minipage} & \begin{minipage}[b]{\linewidth}\raggedright
EDIT
\end{minipage} & \begin{minipage}[b]{\linewidth}\raggedright
PROPS
\end{minipage} & \begin{minipage}[b]{\linewidth}\raggedright
REGIONS
\end{minipage} & \begin{minipage}[b]{\linewidth}\raggedright
SOLUTION
\end{minipage} & \begin{minipage}[b]{\linewidth}\raggedright
SUMMARY
\end{minipage} & \begin{minipage}[b]{\linewidth}\raggedright
SCHEDULE
\end{minipage} \\
\midrule\noalign{}
\endhead
\bottomrule\noalign{}
\endlastfoot
\end{longtable}

\subsubsection{Description}\label{description-546}

The MSFN keyword defines the miscible normalized relative permeability
tables for when the MISCIBLE and or SOLVENT options have been activated
in the RUNSPEC section using the respective keyword. The MISCIBLE
keyword invokes a three component formulation (oil, water and solvent
gas or an oil, water and solvent oil). Whereas the SOLVENT keyword
results in a four component model (oil, water and gas plus a solvent).
This keyword should only be used if the MISCIBLE and or SOLVENT options
have been activated.

\begin{longtable}[]{@{}
  >{\raggedright\arraybackslash}p{(\linewidth - 6\tabcolsep) * \real{0.2500}}
  >{\raggedright\arraybackslash}p{(\linewidth - 6\tabcolsep) * \real{0.2500}}
  >{\raggedright\arraybackslash}p{(\linewidth - 6\tabcolsep) * \real{0.2500}}
  >{\raggedright\arraybackslash}p{(\linewidth - 6\tabcolsep) * \real{0.2500}}@{}}
\toprule\noalign{}
\begin{minipage}[b]{\linewidth}\raggedright
No.
\end{minipage} & \begin{minipage}[b]{\linewidth}\raggedright
Name
\end{minipage} & \begin{minipage}[b]{\linewidth}\raggedright
Description
\end{minipage} & \begin{minipage}[b]{\linewidth}\raggedright
Default
\end{minipage} \\
\midrule\noalign{}
\endhead
\bottomrule\noalign{}
\endlastfoot
1 & SGAS & A columnar vector of real monotonically increasing down the
column values starting from zero and terminating at one, that defines
the gas plus solvent saturation. & None \\
2 & KRSG & A columnar vector of real values that are either equal or
increasing down the column and that are greater than or equal to zero
and less than or equal to one that defines the gas plus solvent relative
permeability multiplier. & None \\
3 & KRO & A columnar vector of real values that are either equal or
increasing down the column and that are greater than or equal to zero
and less than or equal to one that defines the oil relative permeability
multiplier. & None \\
Notes: & & & \\
\end{longtable}

\emph{Table 8.89: MSFN Keyword Description}

\subsubsection{Example}\label{example-300}

\begin{verbatim}
--
--       MISCIBLE NORMALIZED RELATIVE PERMEABILITY TABLES
--
MSFN
--       SGAS       KRSG         KRO
--       FRAC
--       --------  --------     --------
          0.0000    0.0000       1.0000
          1.0000    1.0000       0.0000                    / TABLE NO. 01

--       SGAS       KRSG         KRO
--       FRAC
--       --------  --------     --------
          0.0000    0.0000       1.0000
          0.2000    0.2000       0.8000
          0.4000    0.3000       0.7000
          0.6000    0.4000       0.6000
          0.8000    0.5000       0.4000
          1.0000    1.0000       0.0000                    / TABLE NO. 02
\end{verbatim}

The above example defines two MSN tables for use the MISCIBLE and
SOLVENT options.

\subsection{MULTIPLY -- Multiply a Specified Array by a
Constant}\label{multiply-multiply-a-specified-array-by-a-constant-2}

The MULTIPLY keyword multiplies a specified array or part of an array by
a constant. The constant can be real or integer depending on the array
type; however, the arrays that can be operated on is dependent on which
section the keyword is being used.

See MULTIPLY -- Multiply a Specified Array by a Constant in the GRID
section for a full description.

\subsection{MULTIREG -- Multiply an Array by a Constant based on a
Region
Number}\label{multireg-multiply-an-array-by-a-constant-based-on-a-region-number-2}

The MULTIREG keyword multiplies an array or part of an array by a
constant for cells with a specific region number. The region number
array can be FLUXNUM, MULTNUM or OPERNUM and these arrays must be
defined and be available before the MULTIREG keyword is read by the
simulator. The constant can be real or integer depending on the property
array type; however, the property arrays that can be operated on is
dependent on which section the MULTIREG keyword is being used.

See MULTIREG -- Multiply an Array by a Constant based on a Region Number
in the GRID section for a full description.

\subsection{MW -- Molecular Weights}\label{mw-molecular-weights}

\begin{longtable}[]{@{}
  >{\raggedright\arraybackslash}p{(\linewidth - 14\tabcolsep) * \real{0.1250}}
  >{\raggedright\arraybackslash}p{(\linewidth - 14\tabcolsep) * \real{0.1250}}
  >{\raggedright\arraybackslash}p{(\linewidth - 14\tabcolsep) * \real{0.1250}}
  >{\raggedright\arraybackslash}p{(\linewidth - 14\tabcolsep) * \real{0.1250}}
  >{\raggedright\arraybackslash}p{(\linewidth - 14\tabcolsep) * \real{0.1250}}
  >{\raggedright\arraybackslash}p{(\linewidth - 14\tabcolsep) * \real{0.1250}}
  >{\raggedright\arraybackslash}p{(\linewidth - 14\tabcolsep) * \real{0.1250}}
  >{\raggedright\arraybackslash}p{(\linewidth - 14\tabcolsep) * \real{0.1250}}@{}}
\toprule\noalign{}
\begin{minipage}[b]{\linewidth}\raggedright
RUNSPEC
\end{minipage} & \begin{minipage}[b]{\linewidth}\raggedright
GRID
\end{minipage} & \begin{minipage}[b]{\linewidth}\raggedright
EDIT
\end{minipage} & \begin{minipage}[b]{\linewidth}\raggedright
PROPS
\end{minipage} & \begin{minipage}[b]{\linewidth}\raggedright
REGIONS
\end{minipage} & \begin{minipage}[b]{\linewidth}\raggedright
SOLUTION
\end{minipage} & \begin{minipage}[b]{\linewidth}\raggedright
SUMMARY
\end{minipage} & \begin{minipage}[b]{\linewidth}\raggedright
SCHEDULE
\end{minipage} \\
\midrule\noalign{}
\endhead
\bottomrule\noalign{}
\endlastfoot
\end{longtable}

\subsubsection{Description}\label{description-547}

The \hyperref[REF_HEADING_KEYWORD_MW]{MW} keyword defines the molecular
weights for each of the compositional components active in the model and
for each equation of state. The keyword should only be used if the
compositional mode has been requested using the COMPS keyword in the
RUNSPEC section.

OPM Flow does not currently support the general compositional modeling
formulation.

This keyword is not supported by OPM Flow but it will be parsed and its
data ignored.

\begin{longtable}[]{@{}
  >{\raggedright\arraybackslash}p{(\linewidth - 6\tabcolsep) * \real{0.2500}}
  >{\raggedright\arraybackslash}p{(\linewidth - 6\tabcolsep) * \real{0.2500}}
  >{\raggedright\arraybackslash}p{(\linewidth - 6\tabcolsep) * \real{0.2500}}
  >{\raggedright\arraybackslash}p{(\linewidth - 6\tabcolsep) * \real{0.2500}}@{}}
\toprule\noalign{}
\begin{minipage}[b]{\linewidth}\raggedright
No.
\end{minipage} & \begin{minipage}[b]{\linewidth}\raggedright
Name
\end{minipage} & \begin{minipage}[b]{\linewidth}\raggedright
Description
\end{minipage} & \begin{minipage}[b]{\linewidth}\raggedright
Default
\end{minipage} \\
\midrule\noalign{}
\endhead
\bottomrule\noalign{}
\endlastfoot
Field & Metric & Laboratory & \\
1 & \hyperref[REF_HEADING_KEYWORD_MW]{MW} & A series of real numbers
that define the molecular weights for each of the compositional
components active in the model. & None \\
lb/lb-M & kg/kg-M & gm/gm-M & \\
Notes: & & & \\
\end{longtable}

\emph{Table 8.3.184.1: \hyperref[REF_HEADING_KEYWORD_MW]{MW} Keyword
Description}

\subsubsection{Examples}\label{examples-36}

The following example defines the molecular weights for each component
in a single three-component equation of state model.

\begin{verbatim}
--
-- Molecular Weights
--
MW
  16.0425  44.009   18.01528 /
\end{verbatim}

The following example defines the molecular weights for each component
in two three-component equation of state models.

\begin{verbatim}
--
-- Molecular Weights
--
MW
  16.0425  44.009   18.01528 /
  16.0425  44.009   18.01528 /
\end{verbatim}

\subsection{NCOMPS -- Confirm Number of Compositional
Components}\label{ncomps-confirm-number-of-compositional-components}

\begin{longtable}[]{@{}
  >{\raggedright\arraybackslash}p{(\linewidth - 14\tabcolsep) * \real{0.1250}}
  >{\raggedright\arraybackslash}p{(\linewidth - 14\tabcolsep) * \real{0.1250}}
  >{\raggedright\arraybackslash}p{(\linewidth - 14\tabcolsep) * \real{0.1250}}
  >{\raggedright\arraybackslash}p{(\linewidth - 14\tabcolsep) * \real{0.1250}}
  >{\raggedright\arraybackslash}p{(\linewidth - 14\tabcolsep) * \real{0.1250}}
  >{\raggedright\arraybackslash}p{(\linewidth - 14\tabcolsep) * \real{0.1250}}
  >{\raggedright\arraybackslash}p{(\linewidth - 14\tabcolsep) * \real{0.1250}}
  >{\raggedright\arraybackslash}p{(\linewidth - 14\tabcolsep) * \real{0.1250}}@{}}
\toprule\noalign{}
\begin{minipage}[b]{\linewidth}\raggedright
RUNSPEC
\end{minipage} & \begin{minipage}[b]{\linewidth}\raggedright
GRID
\end{minipage} & \begin{minipage}[b]{\linewidth}\raggedright
EDIT
\end{minipage} & \begin{minipage}[b]{\linewidth}\raggedright
PROPS
\end{minipage} & \begin{minipage}[b]{\linewidth}\raggedright
REGIONS
\end{minipage} & \begin{minipage}[b]{\linewidth}\raggedright
SOLUTION
\end{minipage} & \begin{minipage}[b]{\linewidth}\raggedright
SUMMARY
\end{minipage} & \begin{minipage}[b]{\linewidth}\raggedright
SCHEDULE
\end{minipage} \\
\midrule\noalign{}
\endhead
\bottomrule\noalign{}
\endlastfoot
\end{longtable}

\subsubsection{Description}\label{description-548}

This keyword confirms the maximum number of compositional components in
the model, as such it should have the same number of components as that
declared via the COMPS keyword in the RUNSPEC section. The keyword
should only be used if the CO2STORE keyword and either the GASWAT or the
GAS and WATER keywords in the RUNSPEC section, have also be activated
for the gas-water two component model.

\begin{tcolorbox}[enhanced jigsaw, title=\textcolor{quarto-callout-note-color}{\faInfo}\hspace{0.5em}{Note}, coltitle=black, breakable, toprule=.15mm, colbacktitle=quarto-callout-note-color!10!white, leftrule=.75mm, opacitybacktitle=0.6, colback=white, bottomtitle=1mm, toptitle=1mm, titlerule=0mm, left=2mm, colframe=quarto-callout-note-color-frame, arc=.35mm, rightrule=.15mm, bottomrule=.15mm, opacityback=0]

This is an OPM Flow keyword used with OPM Flow's CO2STORE and GASWAT
keywords in the RUNSPEC section, and should not be confused with the
more general version of the NCOMPS keyword used in the commercial
compositional simulator. Secondly, although OPM Flow parses the keyword,
the simulator currently ignores the data for this keyword.

\end{tcolorbox}

\begin{longtable}[]{@{}
  >{\raggedright\arraybackslash}p{(\linewidth - 6\tabcolsep) * \real{0.2500}}
  >{\raggedright\arraybackslash}p{(\linewidth - 6\tabcolsep) * \real{0.2500}}
  >{\raggedright\arraybackslash}p{(\linewidth - 6\tabcolsep) * \real{0.2500}}
  >{\raggedright\arraybackslash}p{(\linewidth - 6\tabcolsep) * \real{0.2500}}@{}}
\toprule\noalign{}
\begin{minipage}[b]{\linewidth}\raggedright
No.
\end{minipage} & \begin{minipage}[b]{\linewidth}\raggedright
Name
\end{minipage} & \begin{minipage}[b]{\linewidth}\raggedright
Description
\end{minipage} & \begin{minipage}[b]{\linewidth}\raggedright
Default
\end{minipage} \\
\midrule\noalign{}
\endhead
\bottomrule\noalign{}
\endlastfoot
1 & NCOMPS & A positive integer defining the maximum number of
compositional components in the model. Secondly the number of components
must be the same as that enter via the COMPS keyword in the RUNSPEC
section. Only the default value of two is currently supported by OPM
Flow. & 2 \\
Notes: & & & \\
\end{longtable}

\emph{Table 8.90: NCOMPS Keyword Description}

\subsubsection{Example}\label{example-301}

The following example defines how to confirm a two component formulation
to be used with the CO2STORE and GASWAT options.

\begin{verbatim}
--
--       CONFIRM NUMBER OF COMPOSITIONAL COMPONENTS (OPM FLOW KEYWORD)
--
NCOMPS
         2                                                                     /
\end{verbatim}

\subsection{NOECHO -- Deactivate Echoing of User Input Files to the
Print
File}\label{noecho-deactivate-echoing-of-user-input-files-to-the-print-file-4}

Turns off echoing of all the input files to the print file. Note by
default echoing of the inputs files is active. but can subsequently be
switched off by the NOECHO activation keyword.

See NOECHO -- Deactivate Echoing of User Input Files to the Print File
in the GLOBAL section for a full description.

\subsection{NOWARN -- Deactivate Warning
Messages}\label{nowarn-deactivate-warning-messages-4}

Turns off warning messages to be printed to the print file; note that
this keyword is deactivated by default and can subsequently be switched
off by the WARN activation keyword. The warning messages may be turned
on and off using keywords WARN and NOWARN.

See NOWARN -- Deactivate Warning Messages in the GLOBAL section for a
full description.

\subsection{NOWARNEP -- Deactivate End-Point Scaling Warning
Messages}\label{nowarnep-deactivate-end-point-scaling-warning-messages}

\begin{longtable}[]{@{}
  >{\raggedright\arraybackslash}p{(\linewidth - 14\tabcolsep) * \real{0.1250}}
  >{\raggedright\arraybackslash}p{(\linewidth - 14\tabcolsep) * \real{0.1250}}
  >{\raggedright\arraybackslash}p{(\linewidth - 14\tabcolsep) * \real{0.1250}}
  >{\raggedright\arraybackslash}p{(\linewidth - 14\tabcolsep) * \real{0.1250}}
  >{\raggedright\arraybackslash}p{(\linewidth - 14\tabcolsep) * \real{0.1250}}
  >{\raggedright\arraybackslash}p{(\linewidth - 14\tabcolsep) * \real{0.1250}}
  >{\raggedright\arraybackslash}p{(\linewidth - 14\tabcolsep) * \real{0.1250}}
  >{\raggedright\arraybackslash}p{(\linewidth - 14\tabcolsep) * \real{0.1250}}@{}}
\toprule\noalign{}
\begin{minipage}[b]{\linewidth}\raggedright
RUNSPEC
\end{minipage} & \begin{minipage}[b]{\linewidth}\raggedright
GRID
\end{minipage} & \begin{minipage}[b]{\linewidth}\raggedright
EDIT
\end{minipage} & \begin{minipage}[b]{\linewidth}\raggedright
PROPS
\end{minipage} & \begin{minipage}[b]{\linewidth}\raggedright
REGIONS
\end{minipage} & \begin{minipage}[b]{\linewidth}\raggedright
SOLUTION
\end{minipage} & \begin{minipage}[b]{\linewidth}\raggedright
SUMMARY
\end{minipage} & \begin{minipage}[b]{\linewidth}\raggedright
SCHEDULE
\end{minipage} \\
\midrule\noalign{}
\endhead
\bottomrule\noalign{}
\endlastfoot
\end{longtable}

\subsubsection{Description}\label{description-549}

The NOWARNEP keyword deactivates the writing out of warning messages
associated with checking the consistency of saturation table end-points;
however error messages are still reported by the simulator.

Hence, OPM Flow ignores this keyword but it is documented here for
completeness.

There is no data required for this keyword and there is no terminating
``/'' for this keyword.

\subsubsection{Example}\label{example-302}

\begin{verbatim}
--
--       DEACTIVATE END-POINT SCALING WARNING MESSAGES
--
NOWARNEP
\end{verbatim}

The above example switches off the writing out of warning messages
associated with checking the consistency of saturation table end-points;

\subsection{OILDENT -- Define Oil Density Temperature
Coefficients}\label{oildent-define-oil-density-temperature-coefficients}

\begin{longtable}[]{@{}
  >{\raggedright\arraybackslash}p{(\linewidth - 14\tabcolsep) * \real{0.1250}}
  >{\raggedright\arraybackslash}p{(\linewidth - 14\tabcolsep) * \real{0.1250}}
  >{\raggedright\arraybackslash}p{(\linewidth - 14\tabcolsep) * \real{0.1250}}
  >{\raggedright\arraybackslash}p{(\linewidth - 14\tabcolsep) * \real{0.1250}}
  >{\raggedright\arraybackslash}p{(\linewidth - 14\tabcolsep) * \real{0.1250}}
  >{\raggedright\arraybackslash}p{(\linewidth - 14\tabcolsep) * \real{0.1250}}
  >{\raggedright\arraybackslash}p{(\linewidth - 14\tabcolsep) * \real{0.1250}}
  >{\raggedright\arraybackslash}p{(\linewidth - 14\tabcolsep) * \real{0.1250}}@{}}
\toprule\noalign{}
\begin{minipage}[b]{\linewidth}\raggedright
RUNSPEC
\end{minipage} & \begin{minipage}[b]{\linewidth}\raggedright
GRID
\end{minipage} & \begin{minipage}[b]{\linewidth}\raggedright
EDIT
\end{minipage} & \begin{minipage}[b]{\linewidth}\raggedright
PROPS
\end{minipage} & \begin{minipage}[b]{\linewidth}\raggedright
REGIONS
\end{minipage} & \begin{minipage}[b]{\linewidth}\raggedright
SOLUTION
\end{minipage} & \begin{minipage}[b]{\linewidth}\raggedright
SUMMARY
\end{minipage} & \begin{minipage}[b]{\linewidth}\raggedright
SCHEDULE
\end{minipage} \\
\midrule\noalign{}
\endhead
\bottomrule\noalign{}
\endlastfoot
\end{longtable}

\subsubsection{Description}\label{description-550}

OILDENT defines the oil density as a function of temperature
coefficients for when OPM Flow's thermal option has been activated by
the THERMAL keyword in the RUNSPEC. Note this is an OPM Flow keyword
used with OPM Flow's black-oil thermal model that is not available in
the commercial simulator's black-oil thermal formulation.

This keyword can only be used if OPM Flow's thermal option has been
activated by the THERMAL keyword in the RUNSPEC section. Note this is
different to the commercial simulator that uses the TEMP keyword in the
RUNSPEC section to activate the black-oil thermal model.

\begin{longtable}[]{@{}
  >{\raggedright\arraybackslash}p{(\linewidth - 6\tabcolsep) * \real{0.2500}}
  >{\raggedright\arraybackslash}p{(\linewidth - 6\tabcolsep) * \real{0.2500}}
  >{\raggedright\arraybackslash}p{(\linewidth - 6\tabcolsep) * \real{0.2500}}
  >{\raggedright\arraybackslash}p{(\linewidth - 6\tabcolsep) * \real{0.2500}}@{}}
\toprule\noalign{}
\begin{minipage}[b]{\linewidth}\raggedright
No.
\end{minipage} & \begin{minipage}[b]{\linewidth}\raggedright
Name
\end{minipage} & \begin{minipage}[b]{\linewidth}\raggedright
Description
\end{minipage} & \begin{minipage}[b]{\linewidth}\raggedright
Default
\end{minipage} \\
\midrule\noalign{}
\endhead
\bottomrule\noalign{}
\endlastfoot
Field & Metric & Laboratory & \\
1 & TEMP & TEMP is a real positive value greater than zero that defines
the absolute reference temperature used with TEXP1 and TEXP2 to estimate
the change in oil density with respect to temperature. & Defined \\
oR 527.67 & K 293.15 & K 293.15 & \\
2 & TEXP1 & TEXP1 is a real positive value greater than zero that
defines the oil thermal expansion coefficient of the first order. &
Defined \\
1/oR 1.67 x 10-4 & 1/K 3.0 x 10-4 & 1/K 3.0 x 10-4 & \\
3 & TEXP2 & TEXP2 is a real positive value greater than zero that
defines the oil thermal expansion coefficient of the second order. &
Defined \\
1/oR2 9.26 x 10-7 & 1/K2 3.0 x 10-6 & 1/K2 3.0 x 10-6 & \\
Notes: & & & \\
\end{longtable}

\emph{Table 8.3.186.1: OILDENT Keyword Description}

The oil density at a given pressure and temperature is calculated from
its value at surface conditions and the oil shrinkage factor (the
reciprocal of the oil formation volume factor) as shown in the following
equation:

\begin{equation}\phantomsection\label{eq-8-3-186-1}{
{ρ}_{o}(p,T)={ρ}_{o}({p}_{s},{T}_{s}){b}_{o}(p,T)
}\end{equation}

Where the temperature dependence of the oil shrinkage factor relative to
its value at the reference temperature is calculated as shown in the
following equation:

\begin{equation}\phantomsection\label{eq-8-3-186-2}{
{b}_{o}(p,T)=\frac{{b}_{o}(p,{T}_{\mathit{ref}})}{1+{c}_{1}(T-{T}_{\mathit{ref}})+{c}_{2}{(T-{T}_{\mathit{ref}})}^{2}}
}\end{equation}

Where:

\({ρ}_{o}\) = oil density

\({b}_{o}\) = oil shrinkage factor

\(p\) = pressure

\(T\) = temperature

\({c}_{1},{c}_{2}\) = thermal expansion coefficients to first and second
order

\(s\) = subscript indicating surface conditions

\(\mathit{ref}\) = subscript indicating reference conditions

\subsubsection{Example}\label{example-303}

The following example shows the OILDENT keyword using the default
values, for when the thermal option has been activated by the THERMAL
keyword in the RUNSPEC section and for when NTPVT on the TABDIMS keyword
in the RUNSPEC section is set to two.

\begin{verbatim}
--
--       OIL DENSITY TEMPERATURE COEFFICIENTS (OPM FLOW THERMAL KEYWORD)
--
--       OIL        DENSITY   DENSITY
--       TEMP       COEFF1    COEFF2
--       --------   -------   -------
OILDENT
         1*         1*        1*                           / TABLE NO. 01
         1*         1*        1*                           / TABLE NO. 02
\end{verbatim}

There is no terminating ``/'' for this keyword.

\subsection{OILJT -- Define Oil Joule-Thomson
Coefficient}\label{oiljt-define-oil-joule-thomson-coefficient}

\begin{longtable}[]{@{}
  >{\raggedright\arraybackslash}p{(\linewidth - 14\tabcolsep) * \real{0.1250}}
  >{\raggedright\arraybackslash}p{(\linewidth - 14\tabcolsep) * \real{0.1250}}
  >{\raggedright\arraybackslash}p{(\linewidth - 14\tabcolsep) * \real{0.1250}}
  >{\raggedright\arraybackslash}p{(\linewidth - 14\tabcolsep) * \real{0.1250}}
  >{\raggedright\arraybackslash}p{(\linewidth - 14\tabcolsep) * \real{0.1250}}
  >{\raggedright\arraybackslash}p{(\linewidth - 14\tabcolsep) * \real{0.1250}}
  >{\raggedright\arraybackslash}p{(\linewidth - 14\tabcolsep) * \real{0.1250}}
  >{\raggedright\arraybackslash}p{(\linewidth - 14\tabcolsep) * \real{0.1250}}@{}}
\toprule\noalign{}
\begin{minipage}[b]{\linewidth}\raggedright
RUNSPEC
\end{minipage} & \begin{minipage}[b]{\linewidth}\raggedright
GRID
\end{minipage} & \begin{minipage}[b]{\linewidth}\raggedright
EDIT
\end{minipage} & \begin{minipage}[b]{\linewidth}\raggedright
PROPS
\end{minipage} & \begin{minipage}[b]{\linewidth}\raggedright
REGIONS
\end{minipage} & \begin{minipage}[b]{\linewidth}\raggedright
SOLUTION
\end{minipage} & \begin{minipage}[b]{\linewidth}\raggedright
SUMMARY
\end{minipage} & \begin{minipage}[b]{\linewidth}\raggedright
SCHEDULE
\end{minipage} \\
\midrule\noalign{}
\endhead
\bottomrule\noalign{}
\endlastfoot
\end{longtable}

\subsubsection{Description}\label{description-551}

OILJT activates the oil Joule-Thomson effect {[}The Joule--Thomson
coefficient is defined as the change in temperature with respect to an
increase in pressure at constant enthalpy.{]} in temperature
calculations, and defines the oil Joule-Thomson Coefficient (``JTC'') at
a given reference pressure, for when OPM Flow's thermal option has been
activated by the THERMAL keyword in the RUNSPEC.

\begin{tcolorbox}[enhanced jigsaw, title=\textcolor{quarto-callout-note-color}{\faInfo}\hspace{0.5em}{Note}, coltitle=black, breakable, toprule=.15mm, colbacktitle=quarto-callout-note-color!10!white, leftrule=.75mm, opacitybacktitle=0.6, colback=white, bottomtitle=1mm, toptitle=1mm, titlerule=0mm, left=2mm, colframe=quarto-callout-note-color-frame, arc=.35mm, rightrule=.15mm, bottomrule=.15mm, opacityback=0]

This is an OPM Flow keyword used with OPM Flow's black-oil thermal
model, that is not available in the commercial simulator's black-oil
thermal formulation.

\end{tcolorbox}

This keyword can only be used if OPM Flow's thermal option has been
activated by the THERMAL keyword in the RUNSPEC section. Note this is
different to the commercial simulator that uses the TEMP keyword in the
RUNSPEC section to activate the black-oil thermal model, and does not
include the Joule-Thomson effect in temperature calculations.

\begin{longtable}[]{@{}
  >{\raggedright\arraybackslash}p{(\linewidth - 6\tabcolsep) * \real{0.2500}}
  >{\raggedright\arraybackslash}p{(\linewidth - 6\tabcolsep) * \real{0.2500}}
  >{\raggedright\arraybackslash}p{(\linewidth - 6\tabcolsep) * \real{0.2500}}
  >{\raggedright\arraybackslash}p{(\linewidth - 6\tabcolsep) * \real{0.2500}}@{}}
\toprule\noalign{}
\begin{minipage}[b]{\linewidth}\raggedright
No.
\end{minipage} & \begin{minipage}[b]{\linewidth}\raggedright
Name
\end{minipage} & \begin{minipage}[b]{\linewidth}\raggedright
Description
\end{minipage} & \begin{minipage}[b]{\linewidth}\raggedright
Default
\end{minipage} \\
\midrule\noalign{}
\endhead
\bottomrule\noalign{}
\endlastfoot
Field & Metric & Laboratory & \\
1 & PRESS & A real positive value that defines the reference pressure
for the corresponding Joule-Thomson Coefficient, OILJTC. & None \\
psia & barsa & atma & \\
2 & OILJTC & OILJTC is a real positive or negative value that defines
the oil phase Joule-Thomson Coefficient. If the value is defaulted (1*)
or set to 0, then OILJTC is internally calculated using the thermal oil
density data on the OILDENT keyword in the PROPS section. If a non-zero
value is specified, then the OILJTC is assumed to be constant and equal
to that value. & 0 \\
oF/psia & oC/barsa & oC/atma & \\
Notes: & & & \\
\end{longtable}

\emph{Table 8.92: OILJT Keyword Description}

The Joule--Thomson effect is when a real gas, as oppose to an ideal gas,
expands, resulting in the temperature of the gas dropping {[}Natural Gas
Engineering (McGraw-Hill chemical engineering series), Donald L. Katz,
Robert l Lee, McGraw-Hill Education, 1990 (ISBN 0071007776,
9780071007771).{]}. For liquids the effect is the opposite, that is the
internal energy is transferred to kinetic energy with a corresponding
increase in temperature as velocity increases.

Thermodynamically, the Joule--Thomson coefficient is defined as the
isenthalpic {[}An isenthalpic process or isoenthalpic process, is a
process that proceeds without any change in enthalpy, H; or specific
enthalpy, h.{]} change in temperature in a fluid caused by a unitary
pressure drop, as shown in the following equation:

\begin{equation}\phantomsection\label{eq-8-66}{
\mathrm{η} = \left(\frac{\partial T}{\partial P}\right)
}\end{equation}

Which can also express as {[}Pippard, A.B.: Elements of Classical
Thermodynamics: For Advanced Students of Physics. Cambridge University
Press, Cambridge, UK (1957){]}:

\begin{equation}\phantomsection\label{eq-8-67}{
\mathrm{η} = \left(T\mathrm{α} - 1\right)\frac{1}{(\mathrm{ρ}{C}_{p})} - {\left(\frac{g}{{C}_{p}}\frac{\mathit{dp}}{\mathit{dz}}\right)}^{-1}
}\end{equation}

Setting the gravity term, \(g\), to zero we have:

\begin{equation}\phantomsection\label{eq-8-68}{
\mathrm{η} = \left(T\mathrm{α} - 1\right)\frac{1}{(\mathrm{ρ}{C}_{b})}
}\end{equation}

Where:

\(η\) = Joule--Thomson coefficient (oC/Pa),

\(α\) = thermal expansivity at constant pressure (1/oC),

\({C}_{p}\) = specific heat at constant pressure (J/kg oC),

\(g\) = gravitational acceleration (m/s2)

\(P\) = pressure (Pa),

\(T\) = temperature (oC), and

\(z\) = height (m).

\subsubsection{Example}\label{example-304}

The following example shows the OILJT keyword for when the thermal
option has been activated by the THERMAL keyword in the RUNSPEC section,
and for when NTPVT on the TABDIMS keyword in the RUNSPEC section is set
equal to two.

\begin{verbatim}
--
--       OIL JOULE-THOMSON COEFFICIENT (OPM FLOW EXTENSION KEYWORD)
--
--       REF        OIL
--       PRESS      JTC
--       --------   -------
OILJT
         20.0       1*                  / TABLE NO. 01
         20.0       -0.20               / TABLE NO. 02
\end{verbatim}

Here the first entry is defaulted, and the simulator will therefore
calculate the oil JTC internally using the data on the OILDENT keyword
in the PROPS section.

There is no terminating ``/'' for this keyword.

\subsection{OILVISCT -- Define Oil Viscosity versus Temperature
Functions}\label{oilvisct-define-oil-viscosity-versus-temperature-functions}

\begin{longtable}[]{@{}
  >{\raggedright\arraybackslash}p{(\linewidth - 14\tabcolsep) * \real{0.1250}}
  >{\raggedright\arraybackslash}p{(\linewidth - 14\tabcolsep) * \real{0.1250}}
  >{\raggedright\arraybackslash}p{(\linewidth - 14\tabcolsep) * \real{0.1250}}
  >{\raggedright\arraybackslash}p{(\linewidth - 14\tabcolsep) * \real{0.1250}}
  >{\raggedright\arraybackslash}p{(\linewidth - 14\tabcolsep) * \real{0.1250}}
  >{\raggedright\arraybackslash}p{(\linewidth - 14\tabcolsep) * \real{0.1250}}
  >{\raggedright\arraybackslash}p{(\linewidth - 14\tabcolsep) * \real{0.1250}}
  >{\raggedright\arraybackslash}p{(\linewidth - 14\tabcolsep) * \real{0.1250}}@{}}
\toprule\noalign{}
\begin{minipage}[b]{\linewidth}\raggedright
RUNSPEC
\end{minipage} & \begin{minipage}[b]{\linewidth}\raggedright
GRID
\end{minipage} & \begin{minipage}[b]{\linewidth}\raggedright
EDIT
\end{minipage} & \begin{minipage}[b]{\linewidth}\raggedright
PROPS
\end{minipage} & \begin{minipage}[b]{\linewidth}\raggedright
REGIONS
\end{minipage} & \begin{minipage}[b]{\linewidth}\raggedright
SOLUTION
\end{minipage} & \begin{minipage}[b]{\linewidth}\raggedright
SUMMARY
\end{minipage} & \begin{minipage}[b]{\linewidth}\raggedright
SCHEDULE
\end{minipage} \\
\midrule\noalign{}
\endhead
\bottomrule\noalign{}
\endlastfoot
\end{longtable}

\subsubsection{Description}\label{description-552}

OILVISCT defines the oil viscosity as a function of temperature for when
OPM Flow's thermal option has been activated by the THERMAL keyword in
the RUNSPEC section. The reference pressure and solution gas-oil ratio
of the oil for this table is given by the VISCREF keyword in the PROPS
section. Note this is an OPM Flow keyword used with OPM Flow's black-oil
thermal model that is not available in the commercial simulator's
black-oil thermal formulation.

This keyword can only be used if OPM Flow's thermal option has been
activated by the THERMAL keyword in the RUNSPEC section. Note this is
different to the commercial simulator that uses the TEMP keyword in the
RUNSPEC section to activate the black-oil thermal model.

\begin{longtable}[]{@{}
  >{\raggedright\arraybackslash}p{(\linewidth - 6\tabcolsep) * \real{0.2500}}
  >{\raggedright\arraybackslash}p{(\linewidth - 6\tabcolsep) * \real{0.2500}}
  >{\raggedright\arraybackslash}p{(\linewidth - 6\tabcolsep) * \real{0.2500}}
  >{\raggedright\arraybackslash}p{(\linewidth - 6\tabcolsep) * \real{0.2500}}@{}}
\toprule\noalign{}
\begin{minipage}[b]{\linewidth}\raggedright
No.
\end{minipage} & \begin{minipage}[b]{\linewidth}\raggedright
Name
\end{minipage} & \begin{minipage}[b]{\linewidth}\raggedright
Description
\end{minipage} & \begin{minipage}[b]{\linewidth}\raggedright
Default
\end{minipage} \\
\midrule\noalign{}
\endhead
\bottomrule\noalign{}
\endlastfoot
Field & Metric & Laboratory & \\
1 & TEMP & A columnar vector of real monotonically increasing down the
column values that defines the temperature values. & None \\
oF & oC & oC & \\
2 & VIS & A columnar vector of real increasing down the column values
that defines the oil viscosity for the corresponding temperature values
(TEMP). VIS should be given at the reference pressure and solution
gas-oil ratio as defined by PRESS and RS variables on the VISCREF
keyword. & None \\
cP & cP & cP & \\
Notes: & & & \\
\end{longtable}

\emph{Table 8.93: OILVISCT Keyword Description}

There is no terminating ``/'' for this keyword.

\subsubsection{Example}\label{example-305}

The following example shows the OILVISCT keyword for when the thermal
option has been activated by the THERMAL keyword in the RUNSPEC section
and for when NTPVT on the TABDIMS keyword in the RUNSPEC section is set
equal to one.

\begin{verbatim}
--
--       OIL VISCOSITY VERSUS TEMPERATURE TABLES (OPM FLOW EXTENSION KEYWORD)
--
--       OIL        OIL
--       TEMP       VISC
--       --------   -------
OILVISCT
           100.0    0.600
           110.0    0.650
           120.0    0.680
           150.0    0.720
           165.0    0.725                                  / TABLE NO. 01
\end{verbatim}

\subsection{OPERATE -- Define Mathematical Operations on
Arrays}\label{operate-define-mathematical-operations-on-arrays-2}

This keyword, OPERATE, defines mathematical operations on property
arrays (NTG, PORO etc.) and optionally using another property array as
input to the function. The keyword allows for various mathematical
functions and their associated variables to be defined and applied to
the selected array data. Input constants can be real or integer
depending on the array type; however, the arrays that can be operated on
is dependent on which section the keyword is being used.

See OPERATE -- Define Mathematical Operations on Arrays in the GRID
section for a full description.

\subsection{OPERATER -- Define Mathematical Operations on Arrays by
Region}\label{operater-define-mathematical-operations-on-arrays-by-region-2}

This keyword defines the mathematical operations on arrays for specific
regions in the commercial simulator and is currently not supported by
OPM Flow. However, similar functionality is provided by the ADD and
MULTIPLY keywords.

See OPERATER -- Define Mathematical Operations on Arrays by Region in
the GRID section for a full description.

\subsection{OVERBURD -- Define Rock Overburden Pressure versus Depth
Tables}\label{overburd-define-rock-overburden-pressure-versus-depth-tables}

\begin{longtable}[]{@{}
  >{\raggedright\arraybackslash}p{(\linewidth - 14\tabcolsep) * \real{0.1250}}
  >{\raggedright\arraybackslash}p{(\linewidth - 14\tabcolsep) * \real{0.1250}}
  >{\raggedright\arraybackslash}p{(\linewidth - 14\tabcolsep) * \real{0.1250}}
  >{\raggedright\arraybackslash}p{(\linewidth - 14\tabcolsep) * \real{0.1250}}
  >{\raggedright\arraybackslash}p{(\linewidth - 14\tabcolsep) * \real{0.1250}}
  >{\raggedright\arraybackslash}p{(\linewidth - 14\tabcolsep) * \real{0.1250}}
  >{\raggedright\arraybackslash}p{(\linewidth - 14\tabcolsep) * \real{0.1250}}
  >{\raggedright\arraybackslash}p{(\linewidth - 14\tabcolsep) * \real{0.1250}}@{}}
\toprule\noalign{}
\begin{minipage}[b]{\linewidth}\raggedright
RUNSPEC
\end{minipage} & \begin{minipage}[b]{\linewidth}\raggedright
GRID
\end{minipage} & \begin{minipage}[b]{\linewidth}\raggedright
EDIT
\end{minipage} & \begin{minipage}[b]{\linewidth}\raggedright
PROPS
\end{minipage} & \begin{minipage}[b]{\linewidth}\raggedright
REGIONS
\end{minipage} & \begin{minipage}[b]{\linewidth}\raggedright
SOLUTION
\end{minipage} & \begin{minipage}[b]{\linewidth}\raggedright
SUMMARY
\end{minipage} & \begin{minipage}[b]{\linewidth}\raggedright
SCHEDULE
\end{minipage} \\
\midrule\noalign{}
\endhead
\bottomrule\noalign{}
\endlastfoot
\end{longtable}

\subsubsection{Description}\label{description-553}

The OVERBURD keyword defines the overburden pressures versus depth
relationship to be applied for when the rock compaction option has been
invoked by the ROCKCOMP keyword in the RUNSPEC section.

The rock compaction pore volume and transmissibility multipliers,
entered via the ROCKTAB, ROCK2D and ROCK2DTR keywords, are applied to
the pore pressure, unless the OVERBURD keyword is included in the input
deck. When the OVERBURD keyword is present the multipliers are applied
to the effective pore volume pressure, that is
\({P}_{(\mathit{effective})} = {P}_{(\mathit{Pressure})} - {P}_{(\mathit{overburden})}\).
If the keyword is not present in the input deck then the overburden
pressure is set to zero.

This keyword should only be used if compaction option has been enabled.

\begin{longtable}[]{@{}
  >{\raggedright\arraybackslash}p{(\linewidth - 6\tabcolsep) * \real{0.2500}}
  >{\raggedright\arraybackslash}p{(\linewidth - 6\tabcolsep) * \real{0.2500}}
  >{\raggedright\arraybackslash}p{(\linewidth - 6\tabcolsep) * \real{0.2500}}
  >{\raggedright\arraybackslash}p{(\linewidth - 6\tabcolsep) * \real{0.2500}}@{}}
\toprule\noalign{}
\begin{minipage}[b]{\linewidth}\raggedright
No.
\end{minipage} & \begin{minipage}[b]{\linewidth}\raggedright
Name
\end{minipage} & \begin{minipage}[b]{\linewidth}\raggedright
Description
\end{minipage} & \begin{minipage}[b]{\linewidth}\raggedright
Default
\end{minipage} \\
\midrule\noalign{}
\endhead
\bottomrule\noalign{}
\endlastfoot
Field & Metric & Laboratory & \\
1 & DEPTH & A columnar vector of real monotonically increasing down the
column values that defines the depth for corresponding overburden
pressure parameter PRESS. & None \\
feet & m & cm & \\
2 & PRESS & A columnar vector of real monotonically increasing down the
column values that defines the corresponding overburden pressure for the
given DEPTH. & None \\
psia & bars & atm & \\
Notes: & & & \\
\end{longtable}

\emph{Table 8.94: OVERBURD Keyword Description}

See also the ROCKTAB, ROCK2D, ROCK2DTR, and ROCKWNOD keywords in the
PROPS section.

\subsubsection{Examples}\label{examples-37}

The example below defines three overburden tables, assuming NTROCC is
equal to three on the ROCKCOMP keyword and NPPVT is greater than or
equal to four on the TABDIMS keyword.

\begin{verbatim}
--
--       OVERBURDEN PRESSURE VERSUS DEPTH TABLES
--
OVERBURD
--       DEPTH    OVERBURDEN
--       FEET     PRESSURE
--       ------   ----------
         1000.0    300.000
         2000.0    600.000
         3000.0    900.000
         4000.0   1200.000                                 / TABLE N0. 01
--       DEPTH    OVERBURDEN
--       FEET     PRESSURE
--       ------   ----------
         1000.0    200.000
         2000.0    400.000
         3000.0    800.000
         4000.0   1000.000                                 / TABLE N0. 02
--       DEPTH    OVERBURDEN
--       FEET     PRESSURE
--       ------   ----------
         1000.0    400.000
         2000.0    800.000
         3000.0   1100.000
         4000.0   1500.000                                 / TABLE N0. 03
\end{verbatim}

Note that there must be exactly NTROCC tables entered for this keyword,
otherwise an error will occur.

\begin{verbatim}
--
--       ROCK COMPACTION TABLES
--
ROCKTAB
--       PRESS    PORV     TX(YZ)
--                MULT     MULT
--       ------   ------   ------
         1000.0   0.9600   0.9650
         1500.0   0.9800   0.9850
         3000.0   0.9900   0.9950
         4500.0   1.0000   1.0000
         4750.0   1.0100   1.0100                          / TABLE NO. 01
--       PRESS    PORV     TX(YZ)
--                MULT     MULT
--       ------   ------   ------
         1000.0   0.9600   0.9650
         1500.0   0.9800   0.9850
         3000.0   0.9900   0.9950
         4500.0   1.0000   1.0000
         4750.0   1.0100   1.0100                          / TABLE NO. 02
--       PRESS    PORV     TX(YZ)
--                MULT     MULT
--       ------   ------   ------
         1000.0   0.9600   0.9650
         2000.0   0.9800   0.9850
         3000.0   0.9900   0.9950
         4000.0   1.0100   1.0100                          / TABLE NO. 03
\end{verbatim}

Here ROCKTAB tables one and two are identical.

\subsection{PCFACT -- Capillary Pressure Multiplication Factor as a
Function of Porosity
Change}\label{pcfact-capillary-pressure-multiplication-factor-as-a-function-of-porosity-change}

\begin{longtable}[]{@{}
  >{\raggedright\arraybackslash}p{(\linewidth - 14\tabcolsep) * \real{0.1250}}
  >{\raggedright\arraybackslash}p{(\linewidth - 14\tabcolsep) * \real{0.1250}}
  >{\raggedright\arraybackslash}p{(\linewidth - 14\tabcolsep) * \real{0.1250}}
  >{\raggedright\arraybackslash}p{(\linewidth - 14\tabcolsep) * \real{0.1250}}
  >{\raggedright\arraybackslash}p{(\linewidth - 14\tabcolsep) * \real{0.1250}}
  >{\raggedright\arraybackslash}p{(\linewidth - 14\tabcolsep) * \real{0.1250}}
  >{\raggedright\arraybackslash}p{(\linewidth - 14\tabcolsep) * \real{0.1250}}
  >{\raggedright\arraybackslash}p{(\linewidth - 14\tabcolsep) * \real{0.1250}}@{}}
\toprule\noalign{}
\begin{minipage}[b]{\linewidth}\raggedright
RUNSPEC
\end{minipage} & \begin{minipage}[b]{\linewidth}\raggedright
GRID
\end{minipage} & \begin{minipage}[b]{\linewidth}\raggedright
EDIT
\end{minipage} & \begin{minipage}[b]{\linewidth}\raggedright
PROPS
\end{minipage} & \begin{minipage}[b]{\linewidth}\raggedright
REGIONS
\end{minipage} & \begin{minipage}[b]{\linewidth}\raggedright
SOLUTION
\end{minipage} & \begin{minipage}[b]{\linewidth}\raggedright
SUMMARY
\end{minipage} & \begin{minipage}[b]{\linewidth}\raggedright
SCHEDULE
\end{minipage} \\
\midrule\noalign{}
\endhead
\bottomrule\noalign{}
\endlastfoot
\end{longtable}

\subsubsection{Description}\label{description-554}

\hyperref[REF_HEADING_KEYWORD_PCFACT_8_3]{PCFACT} defines the capillary
pressure multiplication factor due to a change in porosity. Currently
the keyword is used in conjunction with OPM Flow's Salt Precipitation
model, in which the pore space is reduced due to salt precipitating in
the pore space, causing a reduction in porosity and an associated
increase in capillary pressure.

\begin{tcolorbox}[enhanced jigsaw, title=\textcolor{quarto-callout-note-color}{\faInfo}\hspace{0.5em}{Note}, coltitle=black, breakable, toprule=.15mm, colbacktitle=quarto-callout-note-color!10!white, leftrule=.75mm, opacitybacktitle=0.6, colback=white, bottomtitle=1mm, toptitle=1mm, titlerule=0mm, left=2mm, colframe=quarto-callout-note-color-frame, arc=.35mm, rightrule=.15mm, bottomrule=.15mm, opacityback=0]

This is an OPM Flow specific keyword for the simulator's Salt
Precipitation model that is activated by the BRINE and PRECSALT keywords
and declaring that vaporized water is present in the run via the VAPWAT
keyword. All three keywords are in the RUNSPEC section.

\end{tcolorbox}

\begin{longtable}[]{@{}
  >{\raggedright\arraybackslash}p{(\linewidth - 6\tabcolsep) * \real{0.2500}}
  >{\raggedright\arraybackslash}p{(\linewidth - 6\tabcolsep) * \real{0.2500}}
  >{\raggedright\arraybackslash}p{(\linewidth - 6\tabcolsep) * \real{0.2500}}
  >{\raggedright\arraybackslash}p{(\linewidth - 6\tabcolsep) * \real{0.2500}}@{}}
\toprule\noalign{}
\begin{minipage}[b]{\linewidth}\raggedright
No.
\end{minipage} & \begin{minipage}[b]{\linewidth}\raggedright
Name
\end{minipage} & \begin{minipage}[b]{\linewidth}\raggedright
Description
\end{minipage} & \begin{minipage}[b]{\linewidth}\raggedright
Default
\end{minipage} \\
\midrule\noalign{}
\endhead
\bottomrule\noalign{}
\endlastfoot
Field & Metric & Laboratory & \\
1 & POROFAC & A real monotonically increasing positive columnar vector
that defines the porosity factor (\(\frac{ϕ}{{ϕ}_{0}}\)) for the
corresponding PCFAC vector. In the simulator's Salt Precipitation model,
the maximum value of \(ϕ\) is \({ϕ}_{0}\), implying a maximum value of
one for POROFAC. & None \\
dimensionless & dimensionless & dimensionless & \\
2 & PCFAC & A real positive monotonically decreasing columnar vector
that defines the capillary pressure (\({p}_{c}\)) multiplier associated
with POROFAC and used to scale a grid block's capillary pressure due to
the reduction in pore volume caused by salt precipitation. Where:
\(\begin{matrix}\mathit{PCFAC} = m\left(\frac{ϕ}{{ϕ}_{0}}\right) \\ \mathit{with} {p}_{c} = m\left(\frac{ϕ}{{ϕ}_{0}}\right){p}_{c0}\end{matrix}\)
& None \\
dimensionless & dimensionless & dimensionless & \\
Notes: & & & \\
\end{longtable}

\emph{Table 8.3.199.1: PCFACT Keyword Description}

The porosity reduction is a function of the volume fraction of salt
(\({s}_{s}\)) precipitated out of the vaporized water phase, that is:

\begin{equation}\phantomsection\label{eq-8-3-199-1}{
ϕ = (1 - {s}_{s}){ϕ}_{0}
}\end{equation}

The capillary pressure factor data can be estimated from the
permeability and porosity factors defined by the PERMFACT keyword in the
PROPS section using the Leverett J-function, for example:

\begin{equation}\phantomsection\label{eq-8-3-199-2}{
\frac{{p}_{c}}{{p}_{c0}} = {\left(\frac{ϕ}{{ϕ}_{0}}\frac{{k}_{0}}{k}\right)}^{\frac{1}{2}}
}\end{equation}

Where:

\({ϕ}_{0}\), \(ϕ\) = the initial and scaled porosity,

\({k}_{0}\), \(k\) = the initial and scaled permeability, and

\({p}_{c0}\), \({p}_{c}\) = the initial and scaled capillary pressure.

\subsubsection{Example}\label{example-306}

The example below defines two
\hyperref[REF_HEADING_KEYWORD_PCFACT_8_3]{PCFACT} tables assuming NTSFUN
equals two and NSSFUN is greater than or equal to five on the TABDIMS
keyword in the RUNSPEC section.

\begin{verbatim}
--
--       CAPILLARY PRESSURE FACTOR INCREASE DUE TO SALT PRECIPITATION
--       (OPM FLOW KEYWORD)
--
PCFACT
--       PORO       PC
--       FACTOR     FACTOR
--       -------    --------
         0.00        2.0000
         0.25        2.0000
         0.50        1.4142
         0.75        1.1547
         1.00        1.0000                    / TABLE NO. 01
--       -------    --------
         0.00        2.0000
         0.25        2.0000
         0.50        1.4142
         0.75        1.1547
         1.00        1.0000                    / TABLE NO. 02
\end{verbatim}

Note that the capillary pressure factor has been capped for POROFAC less
than 0.25.

\subsection{PCG -- End-Point Scaling of Grid Cell Maximum Gas Capillary
Pressure
(Drainage)}\label{pcg-end-point-scaling-of-grid-cell-maximum-gas-capillary-pressure-drainage}

\begin{longtable}[]{@{}
  >{\raggedright\arraybackslash}p{(\linewidth - 14\tabcolsep) * \real{0.1250}}
  >{\raggedright\arraybackslash}p{(\linewidth - 14\tabcolsep) * \real{0.1250}}
  >{\raggedright\arraybackslash}p{(\linewidth - 14\tabcolsep) * \real{0.1250}}
  >{\raggedright\arraybackslash}p{(\linewidth - 14\tabcolsep) * \real{0.1250}}
  >{\raggedright\arraybackslash}p{(\linewidth - 14\tabcolsep) * \real{0.1250}}
  >{\raggedright\arraybackslash}p{(\linewidth - 14\tabcolsep) * \real{0.1250}}
  >{\raggedright\arraybackslash}p{(\linewidth - 14\tabcolsep) * \real{0.1250}}
  >{\raggedright\arraybackslash}p{(\linewidth - 14\tabcolsep) * \real{0.1250}}@{}}
\toprule\noalign{}
\begin{minipage}[b]{\linewidth}\raggedright
RUNSPEC
\end{minipage} & \begin{minipage}[b]{\linewidth}\raggedright
GRID
\end{minipage} & \begin{minipage}[b]{\linewidth}\raggedright
EDIT
\end{minipage} & \begin{minipage}[b]{\linewidth}\raggedright
PROPS
\end{minipage} & \begin{minipage}[b]{\linewidth}\raggedright
REGIONS
\end{minipage} & \begin{minipage}[b]{\linewidth}\raggedright
SOLUTION
\end{minipage} & \begin{minipage}[b]{\linewidth}\raggedright
SUMMARY
\end{minipage} & \begin{minipage}[b]{\linewidth}\raggedright
SCHEDULE
\end{minipage} \\
\midrule\noalign{}
\endhead
\bottomrule\noalign{}
\endlastfoot
\end{longtable}

\subsubsection{Description}\label{description-555}

PCG defines the maximum drainage gas-oil capillary pressure values for
all the cells in the model via an array. The ENDSCALE keyword in the
RUNSPEC section should be activated to enable end-point scaling and the
use of this keyword. The keyword can be used with all grid types.

This keyword is not supported by OPM Flow but would change the results
if supported so the simulation will be stopped. See also the IPCG
keyword for the equivalent imbibition functionality.

\begin{longtable}[]{@{}
  >{\raggedright\arraybackslash}p{(\linewidth - 6\tabcolsep) * \real{0.2500}}
  >{\raggedright\arraybackslash}p{(\linewidth - 6\tabcolsep) * \real{0.2500}}
  >{\raggedright\arraybackslash}p{(\linewidth - 6\tabcolsep) * \real{0.2500}}
  >{\raggedright\arraybackslash}p{(\linewidth - 6\tabcolsep) * \real{0.2500}}@{}}
\toprule\noalign{}
\begin{minipage}[b]{\linewidth}\raggedright
No.
\end{minipage} & \begin{minipage}[b]{\linewidth}\raggedright
Name
\end{minipage} & \begin{minipage}[b]{\linewidth}\raggedright
Description
\end{minipage} & \begin{minipage}[b]{\linewidth}\raggedright
Default
\end{minipage} \\
\midrule\noalign{}
\endhead
\bottomrule\noalign{}
\endlastfoot
Field & Metric & Laboratory & \\
1 & PCG & PCG is an array of positive real numbers assigning the maximum
drainage gas-oil capillary pressure values for each cell in the model.
Repeat counts may be used, for example 30*100.0. & None \\
psia & bars & atm & \\
Notes: & & & \\
\end{longtable}

\emph{Table 8.95: PCG Keyword Description}

The capillary pressure for a grid block is scaled by:

\begin{equation}\phantomsection\label{eq-8-69}{
{P}_{c} = {P}_{{c}_{\mathit{TABLE}}}\left(\frac{\mathit{PCG}}{{P}_{{c}_{\mathit{TABLE}-\mathit{MAX}}}}\right)
}\end{equation}

Where:

\({P}_{c}\) = the resulting drainage gas-oil capillary pressure for a
grid cell.

\(\mathit{PCG}\) = the maximum capillary pressure from the PCG array for
a given cell.

\({P}_{{c}_{\mathit{TABLE}}}\) = the capillary pressure in the drainage
capillary pressure table

allocated to the grid block.

\({P}_{{c}_{\mathit{TABLE}-\mathit{MAX}}}\) = the maximum capillary
pressure in the drainage capillary pressure table

allocated to the grid block at \({S}_{g}=1-{S}_{\mathit{wco}}\).

\subsubsection{Example}\label{example-307}

\begin{verbatim}
--
--       DEFINE GRID BLOCK PCG DATA FOR ALL CELLS (BASED ON NX x NY x NZ = 300)
--
PCG
         100*50.0  100*75.0  100*125.0                                         /
\end{verbatim}

The above example defines the PCW for 300 cells in the model as defined
by the DIMENS keyword in the RUNSPEC section.

\subsection{PCG32D -- Gas-Oil Capillary Pressure versus Oil and Water
Saturation
Tables}\label{pcg32d-gas-oil-capillary-pressure-versus-oil-and-water-saturation-tables}

\begin{longtable}[]{@{}
  >{\raggedright\arraybackslash}p{(\linewidth - 14\tabcolsep) * \real{0.1250}}
  >{\raggedright\arraybackslash}p{(\linewidth - 14\tabcolsep) * \real{0.1250}}
  >{\raggedright\arraybackslash}p{(\linewidth - 14\tabcolsep) * \real{0.1250}}
  >{\raggedright\arraybackslash}p{(\linewidth - 14\tabcolsep) * \real{0.1250}}
  >{\raggedright\arraybackslash}p{(\linewidth - 14\tabcolsep) * \real{0.1250}}
  >{\raggedright\arraybackslash}p{(\linewidth - 14\tabcolsep) * \real{0.1250}}
  >{\raggedright\arraybackslash}p{(\linewidth - 14\tabcolsep) * \real{0.1250}}
  >{\raggedright\arraybackslash}p{(\linewidth - 14\tabcolsep) * \real{0.1250}}@{}}
\toprule\noalign{}
\begin{minipage}[b]{\linewidth}\raggedright
RUNSPEC
\end{minipage} & \begin{minipage}[b]{\linewidth}\raggedright
GRID
\end{minipage} & \begin{minipage}[b]{\linewidth}\raggedright
EDIT
\end{minipage} & \begin{minipage}[b]{\linewidth}\raggedright
PROPS
\end{minipage} & \begin{minipage}[b]{\linewidth}\raggedright
REGIONS
\end{minipage} & \begin{minipage}[b]{\linewidth}\raggedright
SOLUTION
\end{minipage} & \begin{minipage}[b]{\linewidth}\raggedright
SUMMARY
\end{minipage} & \begin{minipage}[b]{\linewidth}\raggedright
SCHEDULE
\end{minipage} \\
\midrule\noalign{}
\endhead
\bottomrule\noalign{}
\endlastfoot
\end{longtable}

\subsubsection{Description}\label{description-556}

This keyword, PCG32D, enables the gas-oil capillary pressure data to be
entered as a function of both oil and water saturations. The keyword
should be used in conjunction with the SGF32D keyword in the PROPS
section. See also the PCW32D keyword in the PROPS section that provides
similar functionality for the water-oil capillary pressure data.

This keyword is not supported by OPM Flow but would change the results
if supported so the simulation will be stopped.

\subsection{PCRIT -- Critical Pressures}\label{pcrit-critical-pressures}

\begin{longtable}[]{@{}
  >{\raggedright\arraybackslash}p{(\linewidth - 14\tabcolsep) * \real{0.1250}}
  >{\raggedright\arraybackslash}p{(\linewidth - 14\tabcolsep) * \real{0.1250}}
  >{\raggedright\arraybackslash}p{(\linewidth - 14\tabcolsep) * \real{0.1250}}
  >{\raggedright\arraybackslash}p{(\linewidth - 14\tabcolsep) * \real{0.1250}}
  >{\raggedright\arraybackslash}p{(\linewidth - 14\tabcolsep) * \real{0.1250}}
  >{\raggedright\arraybackslash}p{(\linewidth - 14\tabcolsep) * \real{0.1250}}
  >{\raggedright\arraybackslash}p{(\linewidth - 14\tabcolsep) * \real{0.1250}}
  >{\raggedright\arraybackslash}p{(\linewidth - 14\tabcolsep) * \real{0.1250}}@{}}
\toprule\noalign{}
\begin{minipage}[b]{\linewidth}\raggedright
RUNSPEC
\end{minipage} & \begin{minipage}[b]{\linewidth}\raggedright
GRID
\end{minipage} & \begin{minipage}[b]{\linewidth}\raggedright
EDIT
\end{minipage} & \begin{minipage}[b]{\linewidth}\raggedright
PROPS
\end{minipage} & \begin{minipage}[b]{\linewidth}\raggedright
REGIONS
\end{minipage} & \begin{minipage}[b]{\linewidth}\raggedright
SOLUTION
\end{minipage} & \begin{minipage}[b]{\linewidth}\raggedright
SUMMARY
\end{minipage} & \begin{minipage}[b]{\linewidth}\raggedright
SCHEDULE
\end{minipage} \\
\midrule\noalign{}
\endhead
\bottomrule\noalign{}
\endlastfoot
\end{longtable}

\subsubsection{Description}\label{description-557}

The \hyperref[REF_HEADING_KEYWORD_PCRIT]{PCRIT} keyword defines the
critical pressures for each of the compositional components active in
the model and for each equation of state. The keyword should only be
used if the compositional mode has been requested using the COMPS
keyword in the RUNSPEC section.

OPM Flow does not currently support the general compositional modeling
formulation.

This keyword is not supported by OPM Flow but it will be parsed and its
data ignored.

\begin{longtable}[]{@{}
  >{\raggedright\arraybackslash}p{(\linewidth - 6\tabcolsep) * \real{0.2500}}
  >{\raggedright\arraybackslash}p{(\linewidth - 6\tabcolsep) * \real{0.2500}}
  >{\raggedright\arraybackslash}p{(\linewidth - 6\tabcolsep) * \real{0.2500}}
  >{\raggedright\arraybackslash}p{(\linewidth - 6\tabcolsep) * \real{0.2500}}@{}}
\toprule\noalign{}
\begin{minipage}[b]{\linewidth}\raggedright
No.
\end{minipage} & \begin{minipage}[b]{\linewidth}\raggedright
Name
\end{minipage} & \begin{minipage}[b]{\linewidth}\raggedright
Description
\end{minipage} & \begin{minipage}[b]{\linewidth}\raggedright
Default
\end{minipage} \\
\midrule\noalign{}
\endhead
\bottomrule\noalign{}
\endlastfoot
Field & Metric & Laboratory & \\
1 & \hyperref[REF_HEADING_KEYWORD_PCRIT]{PCRIT} & A series of real
numbers that define the critical pressures for each of the compositional
components active in the model. & None \\
psia & barsa & atma & \\
Notes: & & & \\
\end{longtable}

\emph{Table 8.3.196.1: \hyperref[REF_HEADING_KEYWORD_PCRIT]{PCRIT}
Keyword Description}

\subsubsection{Examples}\label{examples-38}

The following example defines the critical pressures for each component
in a single three-component equation of state model.

\begin{verbatim}
--
-- Critical Pressures
--
PCRIT
  1070.132 667.0576 616.5844 /
\end{verbatim}

The following example defines the critical pressures for each component
in two three-component equation of state models.

\begin{verbatim}
--
-- Critical Pressures
--
PCRIT
  1070.132 667.0576 616.5844 /
  1070.132 667.0576 616.5844 /
\end{verbatim}

\subsection{PCW -- End-Point Scaling of Grid Cell Water Capillary
Pressure
(Drainage)}\label{pcw-end-point-scaling-of-grid-cell-water-capillary-pressure-drainage}

\begin{longtable}[]{@{}
  >{\raggedright\arraybackslash}p{(\linewidth - 14\tabcolsep) * \real{0.1250}}
  >{\raggedright\arraybackslash}p{(\linewidth - 14\tabcolsep) * \real{0.1250}}
  >{\raggedright\arraybackslash}p{(\linewidth - 14\tabcolsep) * \real{0.1250}}
  >{\raggedright\arraybackslash}p{(\linewidth - 14\tabcolsep) * \real{0.1250}}
  >{\raggedright\arraybackslash}p{(\linewidth - 14\tabcolsep) * \real{0.1250}}
  >{\raggedright\arraybackslash}p{(\linewidth - 14\tabcolsep) * \real{0.1250}}
  >{\raggedright\arraybackslash}p{(\linewidth - 14\tabcolsep) * \real{0.1250}}
  >{\raggedright\arraybackslash}p{(\linewidth - 14\tabcolsep) * \real{0.1250}}@{}}
\toprule\noalign{}
\begin{minipage}[b]{\linewidth}\raggedright
RUNSPEC
\end{minipage} & \begin{minipage}[b]{\linewidth}\raggedright
GRID
\end{minipage} & \begin{minipage}[b]{\linewidth}\raggedright
EDIT
\end{minipage} & \begin{minipage}[b]{\linewidth}\raggedright
PROPS
\end{minipage} & \begin{minipage}[b]{\linewidth}\raggedright
REGIONS
\end{minipage} & \begin{minipage}[b]{\linewidth}\raggedright
SOLUTION
\end{minipage} & \begin{minipage}[b]{\linewidth}\raggedright
SUMMARY
\end{minipage} & \begin{minipage}[b]{\linewidth}\raggedright
SCHEDULE
\end{minipage} \\
\midrule\noalign{}
\endhead
\bottomrule\noalign{}
\endlastfoot
\end{longtable}

\subsubsection{Description}\label{description-558}

PCW defines the maximum drainage water-oil or water-gas capillary
pressure values for all the cells in the model via an array. The
ENDSCALE keyword in the RUNSPEC section should be activated to enable
end-point scaling and the use of this keyword. The keyword can be used
with all grid types.

This keyword is not supported by OPM Flow but would change the results
if supported so the simulation will be stopped.

\begin{longtable}[]{@{}
  >{\raggedright\arraybackslash}p{(\linewidth - 6\tabcolsep) * \real{0.2500}}
  >{\raggedright\arraybackslash}p{(\linewidth - 6\tabcolsep) * \real{0.2500}}
  >{\raggedright\arraybackslash}p{(\linewidth - 6\tabcolsep) * \real{0.2500}}
  >{\raggedright\arraybackslash}p{(\linewidth - 6\tabcolsep) * \real{0.2500}}@{}}
\toprule\noalign{}
\begin{minipage}[b]{\linewidth}\raggedright
No.
\end{minipage} & \begin{minipage}[b]{\linewidth}\raggedright
Name
\end{minipage} & \begin{minipage}[b]{\linewidth}\raggedright
Description
\end{minipage} & \begin{minipage}[b]{\linewidth}\raggedright
Default
\end{minipage} \\
\midrule\noalign{}
\endhead
\bottomrule\noalign{}
\endlastfoot
Field & Metric & Laboratory & \\
1 & PCW & PCW is an array of positive real numbers assigning the maximum
drainage water capillary pressure values for each cell in the model.
Repeat counts may be used, for example 30*100.0. & None \\
psia & bars & atm & \\
Notes: & & & \\
\end{longtable}

\emph{Table 8.96: PCW Keyword Description}

See also the IPCW keyword for the equivalent imbibition functionality.

The capillary pressure for a grid block is scaled by:

\begin{equation}\phantomsection\label{eq-8-70}{
{P}_{c} = {P}_{{c}_{\mathit{TABLE}}}\left(\frac{\mathit{PCW}}{{P}_{{c}_{\mathit{TABLE}-\mathit{MAX}}}}\right)
}\end{equation}

Where:

\({P}_{c}\) = the resulting drainage water capillary pressure for a grid
cell.

\(\mathit{PCW}\) = the maximum capillary pressure from the PCW array for
a given cell.

\({P}_{{c}_{\mathit{TABLE}}}\) = the capillary pressure in the drainage
capillary pressure table allocated

to the grid block.

\({P}_{{c}_{\mathit{TABLE}-\mathit{MAX}}}\) = the maximum capillary
pressure in the drainage capillary pressure table

allocated to the grid block (that is at the connate water saturation).

\subsubsection{Example}\label{example-308}

\begin{verbatim}
--
--       DEFINE GRID BLOCK PCW DATA FOR ALL CELLS (BASED ON NX x NY x NZ = 300)
--
PCW
         100*50.0  100*75.0  100*125.0                                        /
\end{verbatim}

The above example defines the PCW for 300 cells in the model as defined
by the DIMENS keyword in the RUNSPEC section.

\subsection{PCW32D -- Water-Oil Capillary Pressure versus Oil and Gas
Saturation
Tables}\label{pcw32d-water-oil-capillary-pressure-versus-oil-and-gas-saturation-tables}

\begin{longtable}[]{@{}
  >{\raggedright\arraybackslash}p{(\linewidth - 14\tabcolsep) * \real{0.1250}}
  >{\raggedright\arraybackslash}p{(\linewidth - 14\tabcolsep) * \real{0.1250}}
  >{\raggedright\arraybackslash}p{(\linewidth - 14\tabcolsep) * \real{0.1250}}
  >{\raggedright\arraybackslash}p{(\linewidth - 14\tabcolsep) * \real{0.1250}}
  >{\raggedright\arraybackslash}p{(\linewidth - 14\tabcolsep) * \real{0.1250}}
  >{\raggedright\arraybackslash}p{(\linewidth - 14\tabcolsep) * \real{0.1250}}
  >{\raggedright\arraybackslash}p{(\linewidth - 14\tabcolsep) * \real{0.1250}}
  >{\raggedright\arraybackslash}p{(\linewidth - 14\tabcolsep) * \real{0.1250}}@{}}
\toprule\noalign{}
\begin{minipage}[b]{\linewidth}\raggedright
RUNSPEC
\end{minipage} & \begin{minipage}[b]{\linewidth}\raggedright
GRID
\end{minipage} & \begin{minipage}[b]{\linewidth}\raggedright
EDIT
\end{minipage} & \begin{minipage}[b]{\linewidth}\raggedright
PROPS
\end{minipage} & \begin{minipage}[b]{\linewidth}\raggedright
REGIONS
\end{minipage} & \begin{minipage}[b]{\linewidth}\raggedright
SOLUTION
\end{minipage} & \begin{minipage}[b]{\linewidth}\raggedright
SUMMARY
\end{minipage} & \begin{minipage}[b]{\linewidth}\raggedright
SCHEDULE
\end{minipage} \\
\midrule\noalign{}
\endhead
\bottomrule\noalign{}
\endlastfoot
\end{longtable}

\subsubsection{Description}\label{description-559}

This keyword, PCW32D, enables the water-oil capillary pressure data to
be entered as a function of both oil and gas saturations. The keyword
should be used in conjunction with the SWF32D keyword in the PROPS
section. See also the PCG32D keyword in the PROPS section that provides
similar functionality for the gas-oil capillary pressure data.

This keyword is not supported by OPM Flow but would change the results
if supported so the simulation will be stopped.

\subsection{PECOEFS -- Define Petro-Elastic Model
Coefficients}\label{pecoefs-define-petro-elastic-model-coefficients}

\begin{longtable}[]{@{}
  >{\raggedright\arraybackslash}p{(\linewidth - 14\tabcolsep) * \real{0.1250}}
  >{\raggedright\arraybackslash}p{(\linewidth - 14\tabcolsep) * \real{0.1250}}
  >{\raggedright\arraybackslash}p{(\linewidth - 14\tabcolsep) * \real{0.1250}}
  >{\raggedright\arraybackslash}p{(\linewidth - 14\tabcolsep) * \real{0.1250}}
  >{\raggedright\arraybackslash}p{(\linewidth - 14\tabcolsep) * \real{0.1250}}
  >{\raggedright\arraybackslash}p{(\linewidth - 14\tabcolsep) * \real{0.1250}}
  >{\raggedright\arraybackslash}p{(\linewidth - 14\tabcolsep) * \real{0.1250}}
  >{\raggedright\arraybackslash}p{(\linewidth - 14\tabcolsep) * \real{0.1250}}@{}}
\toprule\noalign{}
\begin{minipage}[b]{\linewidth}\raggedright
RUNSPEC
\end{minipage} & \begin{minipage}[b]{\linewidth}\raggedright
GRID
\end{minipage} & \begin{minipage}[b]{\linewidth}\raggedright
EDIT
\end{minipage} & \begin{minipage}[b]{\linewidth}\raggedright
PROPS
\end{minipage} & \begin{minipage}[b]{\linewidth}\raggedright
REGIONS
\end{minipage} & \begin{minipage}[b]{\linewidth}\raggedright
SOLUTION
\end{minipage} & \begin{minipage}[b]{\linewidth}\raggedright
SUMMARY
\end{minipage} & \begin{minipage}[b]{\linewidth}\raggedright
SCHEDULE
\end{minipage} \\
\midrule\noalign{}
\endhead
\bottomrule\noalign{}
\endlastfoot
\end{longtable}

\subsubsection{Description}\label{description-560}

The PECOEFS keyword defines the Petro-Elastic model coefficients.

This keyword is not supported by OPM Flow but would change the results
if supported so the simulation will be stopped.

\subsection{PEGTAB -- Petro-Elastic Pressure Shear Modulus
Table}\label{pegtab-petro-elastic-pressure-shear-modulus-table}

\begin{longtable}[]{@{}
  >{\raggedright\arraybackslash}p{(\linewidth - 14\tabcolsep) * \real{0.1250}}
  >{\raggedright\arraybackslash}p{(\linewidth - 14\tabcolsep) * \real{0.1250}}
  >{\raggedright\arraybackslash}p{(\linewidth - 14\tabcolsep) * \real{0.1250}}
  >{\raggedright\arraybackslash}p{(\linewidth - 14\tabcolsep) * \real{0.1250}}
  >{\raggedright\arraybackslash}p{(\linewidth - 14\tabcolsep) * \real{0.1250}}
  >{\raggedright\arraybackslash}p{(\linewidth - 14\tabcolsep) * \real{0.1250}}
  >{\raggedright\arraybackslash}p{(\linewidth - 14\tabcolsep) * \real{0.1250}}
  >{\raggedright\arraybackslash}p{(\linewidth - 14\tabcolsep) * \real{0.1250}}@{}}
\toprule\noalign{}
\begin{minipage}[b]{\linewidth}\raggedright
RUNSPEC
\end{minipage} & \begin{minipage}[b]{\linewidth}\raggedright
GRID
\end{minipage} & \begin{minipage}[b]{\linewidth}\raggedright
EDIT
\end{minipage} & \begin{minipage}[b]{\linewidth}\raggedright
PROPS
\end{minipage} & \begin{minipage}[b]{\linewidth}\raggedright
REGIONS
\end{minipage} & \begin{minipage}[b]{\linewidth}\raggedright
SOLUTION
\end{minipage} & \begin{minipage}[b]{\linewidth}\raggedright
SUMMARY
\end{minipage} & \begin{minipage}[b]{\linewidth}\raggedright
SCHEDULE
\end{minipage} \\
\midrule\noalign{}
\endhead
\bottomrule\noalign{}
\endlastfoot
\end{longtable}

\subsubsection{Description}\label{description-561}

The PEGTAB series of keywords define a series of coefficients of a
polynomial function used in the calculation of the shear modulus in the
petro-elastic model. The series of keywords consist of: PEGTAB0 PEGTAB1,
PEGTAB2, PEGTAB3, PEGTAB4, PEGTAB5, PEGTAB6, and PEGTAB7.

This series of keywords are ignored by OPM Flow and have no effect on
the simulation.

\subsection{PEKTAB -- Petro-Elastic Pressure Bulk Modulus
Table}\label{pektab-petro-elastic-pressure-bulk-modulus-table}

\begin{longtable}[]{@{}
  >{\raggedright\arraybackslash}p{(\linewidth - 14\tabcolsep) * \real{0.1250}}
  >{\raggedright\arraybackslash}p{(\linewidth - 14\tabcolsep) * \real{0.1250}}
  >{\raggedright\arraybackslash}p{(\linewidth - 14\tabcolsep) * \real{0.1250}}
  >{\raggedright\arraybackslash}p{(\linewidth - 14\tabcolsep) * \real{0.1250}}
  >{\raggedright\arraybackslash}p{(\linewidth - 14\tabcolsep) * \real{0.1250}}
  >{\raggedright\arraybackslash}p{(\linewidth - 14\tabcolsep) * \real{0.1250}}
  >{\raggedright\arraybackslash}p{(\linewidth - 14\tabcolsep) * \real{0.1250}}
  >{\raggedright\arraybackslash}p{(\linewidth - 14\tabcolsep) * \real{0.1250}}@{}}
\toprule\noalign{}
\begin{minipage}[b]{\linewidth}\raggedright
RUNSPEC
\end{minipage} & \begin{minipage}[b]{\linewidth}\raggedright
GRID
\end{minipage} & \begin{minipage}[b]{\linewidth}\raggedright
EDIT
\end{minipage} & \begin{minipage}[b]{\linewidth}\raggedright
PROPS
\end{minipage} & \begin{minipage}[b]{\linewidth}\raggedright
REGIONS
\end{minipage} & \begin{minipage}[b]{\linewidth}\raggedright
SOLUTION
\end{minipage} & \begin{minipage}[b]{\linewidth}\raggedright
SUMMARY
\end{minipage} & \begin{minipage}[b]{\linewidth}\raggedright
SCHEDULE
\end{minipage} \\
\midrule\noalign{}
\endhead
\bottomrule\noalign{}
\endlastfoot
\end{longtable}

\subsubsection{Description}\label{description-562}

The PEKTAB series of keywords define a series of coefficients of a
polynomial function used in the calculation of the bulk modulus in the
petro-elastic model. The series of keywords consist of: PEKTAB0 PEKTAB1,
PEKTAB2, PEKTAB3, PEKTAB4, PEKTAB5, PEKTAB6, and PEKTAB7.

This series of keywords are ignored by OPM Flow and have no effect on
the simulation.

\subsection{PERMFACT -- Permeability Multiplication Factor as a Function
of Porosity
Change}\label{permfact-permeability-multiplication-factor-as-a-function-of-porosity-change}

\begin{longtable}[]{@{}
  >{\raggedright\arraybackslash}p{(\linewidth - 14\tabcolsep) * \real{0.1250}}
  >{\raggedright\arraybackslash}p{(\linewidth - 14\tabcolsep) * \real{0.1250}}
  >{\raggedright\arraybackslash}p{(\linewidth - 14\tabcolsep) * \real{0.1250}}
  >{\raggedright\arraybackslash}p{(\linewidth - 14\tabcolsep) * \real{0.1250}}
  >{\raggedright\arraybackslash}p{(\linewidth - 14\tabcolsep) * \real{0.1250}}
  >{\raggedright\arraybackslash}p{(\linewidth - 14\tabcolsep) * \real{0.1250}}
  >{\raggedright\arraybackslash}p{(\linewidth - 14\tabcolsep) * \real{0.1250}}
  >{\raggedright\arraybackslash}p{(\linewidth - 14\tabcolsep) * \real{0.1250}}@{}}
\toprule\noalign{}
\begin{minipage}[b]{\linewidth}\raggedright
RUNSPEC
\end{minipage} & \begin{minipage}[b]{\linewidth}\raggedright
GRID
\end{minipage} & \begin{minipage}[b]{\linewidth}\raggedright
EDIT
\end{minipage} & \begin{minipage}[b]{\linewidth}\raggedright
PROPS
\end{minipage} & \begin{minipage}[b]{\linewidth}\raggedright
REGIONS
\end{minipage} & \begin{minipage}[b]{\linewidth}\raggedright
SOLUTION
\end{minipage} & \begin{minipage}[b]{\linewidth}\raggedright
SUMMARY
\end{minipage} & \begin{minipage}[b]{\linewidth}\raggedright
SCHEDULE
\end{minipage} \\
\midrule\noalign{}
\endhead
\bottomrule\noalign{}
\endlastfoot
\end{longtable}

\subsubsection{Description}\label{description-563}

PERMFACT defines the permeability multiplication factor due to a change
in porosity. The keyword is used in conjunction with OPM Flow's Salt
Precipitation model, in which the pore space is reduced due to salt
precipitating in the pore space, causing a reduction in porosity and an
associated reduction in permeability. This keyword is also used for the
\hyperref[REF_HEADING_KEYWORD_BIOFILM]{BIOFILM} and MICP model, where
the pore space is reduced due to biofilm formation, and for the MICP
model, also due to calcite precipitation.

\begin{tcolorbox}[enhanced jigsaw, title=\textcolor{quarto-callout-note-color}{\faInfo}\hspace{0.5em}{Note}, coltitle=black, breakable, toprule=.15mm, colbacktitle=quarto-callout-note-color!10!white, leftrule=.75mm, opacitybacktitle=0.6, colback=white, bottomtitle=1mm, toptitle=1mm, titlerule=0mm, left=2mm, colframe=quarto-callout-note-color-frame, arc=.35mm, rightrule=.15mm, bottomrule=.15mm, opacityback=0]

This is an OPM Flow specific keyword.

\end{tcolorbox}

\begin{longtable}[]{@{}
  >{\raggedright\arraybackslash}p{(\linewidth - 6\tabcolsep) * \real{0.2500}}
  >{\raggedright\arraybackslash}p{(\linewidth - 6\tabcolsep) * \real{0.2500}}
  >{\raggedright\arraybackslash}p{(\linewidth - 6\tabcolsep) * \real{0.2500}}
  >{\raggedright\arraybackslash}p{(\linewidth - 6\tabcolsep) * \real{0.2500}}@{}}
\toprule\noalign{}
\begin{minipage}[b]{\linewidth}\raggedright
No.
\end{minipage} & \begin{minipage}[b]{\linewidth}\raggedright
Name
\end{minipage} & \begin{minipage}[b]{\linewidth}\raggedright
Description
\end{minipage} & \begin{minipage}[b]{\linewidth}\raggedright
Default
\end{minipage} \\
\midrule\noalign{}
\endhead
\bottomrule\noalign{}
\endlastfoot
Field & Metric & Laboratory & \\
1 & POROFAC & A real monotonically increasing positive columnar vector
that defines the porosity (\(\frac{ϕ}{{ϕ}_{0}}\)) factor for the
corresponding PERMFAC vector. In the simulator's Salt Precipitation
model, the maximum value of \(ϕ\) is \({ϕ}_{0}\), implying a maximum
value of one for POROFAC. & None \\
dimensionless & dimensionless & dimensionless & \\
2 & PERMFAC & A real positive monotonically increasing columnar vector
that defines the permeability (k) multiplier associated with POROFAC and
used to scale a grid block's permeability due to the reduction in pore
volume caused by salt precipitation. Where:
\(\begin{matrix}\mathit{PERMFAC} = m\left(ϕ\right) \\ \mathit{with} k = m\left(ϕ\right){k}_{0}\end{matrix}\)
& None \\
dimensionless & dimensionless & dimensionless & \\
Notes: & & & \\
\end{longtable}

\emph{Table 8.97: PERMFACT Keyword Description}

The porosity reduction can be written as a function proportional to the
initial porosity, e.g., considering the volume fraction of salt
(\({s}_{s}\)) precipitated out of the vaporized water phase, that is:

\begin{equation}\phantomsection\label{eq-8-71}{
\mathrm{ϕ} = (1 - {s}_{s}){\mathrm{ϕ}}_{o}
}\end{equation}

The porosity and associated permeability factor data can be calculated
using a permeability-porosity relationship, for example:

\begin{equation}\phantomsection\label{eq-8-72}{
\frac{k}{{k}_{o}} = {\left(\frac{\mathrm{ϕ}-{\mathrm{ϕ}}_{c}}{{\mathrm{ϕ}}_{o}-{\mathrm{ϕ}}_{c}}\right)}^{\mathrm{γ}}
}\end{equation}

Where:

ko and\({ϕ}_{o}\) = the initial permeability and porosity,

k and\(ϕ\) = the actual permeability and porosity that changes due to
the model dynamics (e.g., precipitation or dissolution of salt, biofilm
formation, calcite precipitation),

\({ϕ}_{c}\) = the residual porosity at which permeability is zero, and

\(γ\) = a positive exponent.

See also Kozeny-Carment (extended) {[}J. Kozeny, ``Ueber kapillare
Leitung des Wassers im Boden.'' Sitzungsber Akad. Wiss., Wien, 136(2a):
271-306, 1927.{]}, {[}P.C. Carman, ``Fluid flow through granular beds.''
Transactions, Institution of Chemical Engineers, London, 15: 150-166,
1937.{]} and {[}P.C. Carman, ``Flow of gases through porous media.''
Butterworths, London, 1956.{]} and Verma-Pruess {[}Verma, A., \& Pruess,
K. (1988). Thermohydrological conditions and silica redistribution near
high-level nuclear wastes emplaced in saturated geological formations.
Journal of Geophysical Research,93, 1159--1173.{]} for additional
functional forms that can be used to derive the tabulated data that can
be entered via the PERMFACT keyword.

\subsubsection{Example}\label{example-309}

The example below defines two PERMFACT tables assuming NTSFUN equals two
and NSSFUN is greater than or equal to five on the TABDIMS keyword in
the RUNSPEC section.

\begin{verbatim}
--
--       PERMEABILITY FACTOR AS A FUNCTION OF POROSITY REDUCTION
--       (OPM FLOW KEYWORD)
--
PERMFACT
--       PORO       PERM
--       FACTOR     FACTOR
--       -------    --------
         0.00        0.0000
         0.25        0.0625
         0.50        0.2500
         0.75        0.5625
         1.00        1.0000                    / TABLE NO. 01
--       -------    --------
         0.00        0.0000
         0.25        0.0625
         0.50        0.2500
         0.75        0.5625
         1.00        1.0000                    / TABLE NO. 02
\end{verbatim}

Both tables use equation (8.72) with Φc equal to zero and ɣ equal to
2.0. Note that PERMFACT changes the permeability in all three
dimensions, that is the PERMX, PERMY and PERMZ arrays are all modified.

\subsection{PLMIXPAR -- Define the Polymer Todd-Longstaff Mixing
Parameters}\label{plmixpar-define-the-polymer-todd-longstaff-mixing-parameters}

\begin{longtable}[]{@{}
  >{\raggedright\arraybackslash}p{(\linewidth - 14\tabcolsep) * \real{0.1250}}
  >{\raggedright\arraybackslash}p{(\linewidth - 14\tabcolsep) * \real{0.1250}}
  >{\raggedright\arraybackslash}p{(\linewidth - 14\tabcolsep) * \real{0.1250}}
  >{\raggedright\arraybackslash}p{(\linewidth - 14\tabcolsep) * \real{0.1250}}
  >{\raggedright\arraybackslash}p{(\linewidth - 14\tabcolsep) * \real{0.1250}}
  >{\raggedright\arraybackslash}p{(\linewidth - 14\tabcolsep) * \real{0.1250}}
  >{\raggedright\arraybackslash}p{(\linewidth - 14\tabcolsep) * \real{0.1250}}
  >{\raggedright\arraybackslash}p{(\linewidth - 14\tabcolsep) * \real{0.1250}}@{}}
\toprule\noalign{}
\begin{minipage}[b]{\linewidth}\raggedright
RUNSPEC
\end{minipage} & \begin{minipage}[b]{\linewidth}\raggedright
GRID
\end{minipage} & \begin{minipage}[b]{\linewidth}\raggedright
EDIT
\end{minipage} & \begin{minipage}[b]{\linewidth}\raggedright
PROPS
\end{minipage} & \begin{minipage}[b]{\linewidth}\raggedright
REGIONS
\end{minipage} & \begin{minipage}[b]{\linewidth}\raggedright
SOLUTION
\end{minipage} & \begin{minipage}[b]{\linewidth}\raggedright
SUMMARY
\end{minipage} & \begin{minipage}[b]{\linewidth}\raggedright
SCHEDULE
\end{minipage} \\
\midrule\noalign{}
\endhead
\bottomrule\noalign{}
\endlastfoot
\end{longtable}

\subsubsection{Description}\label{description-564}

The PLMIXPAR keyword defines the Todd-Longstaff {[}Todd, M. and
Longstaff, W. ``The Development, Testing and Application of a Numerical
Simulator for Predicting Miscible Flood Performance,'' paper SPE 3484,
Journal of Canadian Petroleum Technology (1972) 24, No.~7, 874-882.{]}
mixing parameters for when the polymer option has been activated by the
POLYMER keyword in the RUNSPEC section. This keyword must be present in
the input deck if the POLYMER keyword has been activated.

Note that this keyword is used only for the polymer option, if the
MISCIBLE keyword in the RUNSPEC section has been invoked then in
addition the TLMIXPAR keyword is also required to define the
Todd-Longstaff mixing parameters for the MISCIBLE option.

\begin{longtable}[]{@{}
  >{\raggedright\arraybackslash}p{(\linewidth - 6\tabcolsep) * \real{0.2500}}
  >{\raggedright\arraybackslash}p{(\linewidth - 6\tabcolsep) * \real{0.2500}}
  >{\raggedright\arraybackslash}p{(\linewidth - 6\tabcolsep) * \real{0.2500}}
  >{\raggedright\arraybackslash}p{(\linewidth - 6\tabcolsep) * \real{0.2500}}@{}}
\toprule\noalign{}
\begin{minipage}[b]{\linewidth}\raggedright
No.
\end{minipage} & \begin{minipage}[b]{\linewidth}\raggedright
Name
\end{minipage} & \begin{minipage}[b]{\linewidth}\raggedright
Description
\end{minipage} & \begin{minipage}[b]{\linewidth}\raggedright
Default
\end{minipage} \\
\midrule\noalign{}
\endhead
\bottomrule\noalign{}
\endlastfoot
Field & Metric & Laboratory & \\
1 & PLMVIS & A real positive value that is greater than or equal to zero
and less than or equal to one, that defines the viscosity Todd-Longstaff
mixing parameter for each polymer region. & None \\
dimensionless & dimensionless & dimensionless & \\
Notes: & & & \\
\end{longtable}

\emph{Table 8.98: PLMIXPAR Keyword Description}

\subsubsection{Example}\label{example-310}

\begin{verbatim}
--
--       POLYMER TODD-LONGSTAFF MIXING PARAMETERS
--
PLMIXPAR
--       PLM
--       VISCOS
--       -------
         0.3500                                            / TABLE NO. 01
         0.2500                                            / TABLE NO. 02
         0.6500                                            / TABLE NO. 03
\end{verbatim}

The above example defines three polymer Todd-Longstaff mixing parameter
data sets, based on the NPLMIX variable on the REGDIMS keyword in the
RUNSPEC section being equal to three.

\subsection{PLYADS -- Define Polymer Rock Adsorption
Tables}\label{plyads-define-polymer-rock-adsorption-tables}

\begin{longtable}[]{@{}
  >{\raggedright\arraybackslash}p{(\linewidth - 14\tabcolsep) * \real{0.1250}}
  >{\raggedright\arraybackslash}p{(\linewidth - 14\tabcolsep) * \real{0.1250}}
  >{\raggedright\arraybackslash}p{(\linewidth - 14\tabcolsep) * \real{0.1250}}
  >{\raggedright\arraybackslash}p{(\linewidth - 14\tabcolsep) * \real{0.1250}}
  >{\raggedright\arraybackslash}p{(\linewidth - 14\tabcolsep) * \real{0.1250}}
  >{\raggedright\arraybackslash}p{(\linewidth - 14\tabcolsep) * \real{0.1250}}
  >{\raggedright\arraybackslash}p{(\linewidth - 14\tabcolsep) * \real{0.1250}}
  >{\raggedright\arraybackslash}p{(\linewidth - 14\tabcolsep) * \real{0.1250}}@{}}
\toprule\noalign{}
\begin{minipage}[b]{\linewidth}\raggedright
RUNSPEC
\end{minipage} & \begin{minipage}[b]{\linewidth}\raggedright
GRID
\end{minipage} & \begin{minipage}[b]{\linewidth}\raggedright
EDIT
\end{minipage} & \begin{minipage}[b]{\linewidth}\raggedright
PROPS
\end{minipage} & \begin{minipage}[b]{\linewidth}\raggedright
REGIONS
\end{minipage} & \begin{minipage}[b]{\linewidth}\raggedright
SOLUTION
\end{minipage} & \begin{minipage}[b]{\linewidth}\raggedright
SUMMARY
\end{minipage} & \begin{minipage}[b]{\linewidth}\raggedright
SCHEDULE
\end{minipage} \\
\midrule\noalign{}
\endhead
\bottomrule\noalign{}
\endlastfoot
\end{longtable}

\subsubsection{Description}\label{description-565}

The PLYADS keyword defines the rock polymer adsorption tables for when
the polymer option has been activated by the POLYMER keyword in the
RUNSPEC section. Alternatively, the functions can be entered via the
PLYADSS keyword in the PROPS section for when salt sensitivity is to be
considered.

\begin{longtable}[]{@{}
  >{\raggedright\arraybackslash}p{(\linewidth - 6\tabcolsep) * \real{0.2500}}
  >{\raggedright\arraybackslash}p{(\linewidth - 6\tabcolsep) * \real{0.2500}}
  >{\raggedright\arraybackslash}p{(\linewidth - 6\tabcolsep) * \real{0.2500}}
  >{\raggedright\arraybackslash}p{(\linewidth - 6\tabcolsep) * \real{0.2500}}@{}}
\toprule\noalign{}
\begin{minipage}[b]{\linewidth}\raggedright
No.
\end{minipage} & \begin{minipage}[b]{\linewidth}\raggedright
Name
\end{minipage} & \begin{minipage}[b]{\linewidth}\raggedright
Description
\end{minipage} & \begin{minipage}[b]{\linewidth}\raggedright
Default
\end{minipage} \\
\midrule\noalign{}
\endhead
\bottomrule\noalign{}
\endlastfoot
Field & Metric & Laboratory & \\
1 & POLCON & A columnar vector of real monotonically increasing down the
column values that defines the polymer concentration in the solution
surrounding the rock. The first entry should be zero to define a no
polymer concentration. & None \\
lb/stb & kg/sm3 & gm/scc & \\
2 & POLRATIO & A columnar vector of real increasing down the column
values that defines the mass of adsorbed polymer per unit mass of rock.
The first entry should be zero to define a zero ratio of polymer
concentration. & None \\
lb/lb & kg/kg & gm/gm & \\
Notes: & & & \\
\end{longtable}

\emph{Table 8.99: PLYADS Keyword Description}

See also the PLYADSS keyword in the PROPS section to also define rock
polymer adsorption tables when the polymer concentration is a function
of salinity.

\subsubsection{Example}\label{example-311}

\begin{verbatim}
--
--       POLYMER ROCK ADSORPTION TABLE
--
PLYADS
--       POLYMER    POLYMER
--       POLCON     POLRATIO
--       -------    --------
             0.0     0.00000
             2.0     0.00003
             4.0     0.00005
             6.0     0.00007
             8.0     0.00009
            10.0     0.00011
            12.0     0.00012
            14.0     0.00015                               / TABLE NO. 01
--       POLYMER    POLYMER
--       POLCON     POLRATIO
--       -------    --------
             0.0     0.00000
             3.0     0.00004
             5.0     0.00006
             7.0     0.00008
             8.0     0.00009
            10.0     0.00011                               / TABLE NO. 02
\end{verbatim}

The above example defines two polymer rock adsorption tables assuming
NTSFUN equals two and NSSFUN is greater than or equal to eight on the
TABDIMS keyword in the RUNSPEC section.

There is no terminating ``/'' for this keyword.

\subsection{PLYADSS -- Define Polymer Rock Adsorption with Salt
Dependence
Tables}\label{plyadss-define-polymer-rock-adsorption-with-salt-dependence-tables}

\begin{longtable}[]{@{}
  >{\raggedright\arraybackslash}p{(\linewidth - 14\tabcolsep) * \real{0.1250}}
  >{\raggedright\arraybackslash}p{(\linewidth - 14\tabcolsep) * \real{0.1250}}
  >{\raggedright\arraybackslash}p{(\linewidth - 14\tabcolsep) * \real{0.1250}}
  >{\raggedright\arraybackslash}p{(\linewidth - 14\tabcolsep) * \real{0.1250}}
  >{\raggedright\arraybackslash}p{(\linewidth - 14\tabcolsep) * \real{0.1250}}
  >{\raggedright\arraybackslash}p{(\linewidth - 14\tabcolsep) * \real{0.1250}}
  >{\raggedright\arraybackslash}p{(\linewidth - 14\tabcolsep) * \real{0.1250}}
  >{\raggedright\arraybackslash}p{(\linewidth - 14\tabcolsep) * \real{0.1250}}@{}}
\toprule\noalign{}
\begin{minipage}[b]{\linewidth}\raggedright
RUNSPEC
\end{minipage} & \begin{minipage}[b]{\linewidth}\raggedright
GRID
\end{minipage} & \begin{minipage}[b]{\linewidth}\raggedright
EDIT
\end{minipage} & \begin{minipage}[b]{\linewidth}\raggedright
PROPS
\end{minipage} & \begin{minipage}[b]{\linewidth}\raggedright
REGIONS
\end{minipage} & \begin{minipage}[b]{\linewidth}\raggedright
SOLUTION
\end{minipage} & \begin{minipage}[b]{\linewidth}\raggedright
SUMMARY
\end{minipage} & \begin{minipage}[b]{\linewidth}\raggedright
SCHEDULE
\end{minipage} \\
\midrule\noalign{}
\endhead
\bottomrule\noalign{}
\endlastfoot
\end{longtable}

\subsubsection{Description}\label{description-566}

The PLYADSS keyword defines the rock polymer adsorption tables for when
the polymer and the salt options has been activated by the POLYMER and
BRINE keywords in the RUNSPEC section.

Note that the BRINE option is not currently supported by OPM Flow;
however, the polymer rock adsorption functions without salt dependence
may be entered via the PLYADS keyword in the PROPS section, for when
salt sensitivity is not to be considered.

This keyword is not supported by OPM Flow but would change the results
if supported so the simulation will be stopped.

\begin{longtable}[]{@{}
  >{\raggedright\arraybackslash}p{(\linewidth - 6\tabcolsep) * \real{0.2500}}
  >{\raggedright\arraybackslash}p{(\linewidth - 6\tabcolsep) * \real{0.2500}}
  >{\raggedright\arraybackslash}p{(\linewidth - 6\tabcolsep) * \real{0.2500}}
  >{\raggedright\arraybackslash}p{(\linewidth - 6\tabcolsep) * \real{0.2500}}@{}}
\toprule\noalign{}
\begin{minipage}[b]{\linewidth}\raggedright
No.
\end{minipage} & \begin{minipage}[b]{\linewidth}\raggedright
Name
\end{minipage} & \begin{minipage}[b]{\linewidth}\raggedright
Description
\end{minipage} & \begin{minipage}[b]{\linewidth}\raggedright
Default
\end{minipage} \\
\midrule\noalign{}
\endhead
\bottomrule\noalign{}
\endlastfoot
Field & Metric & Laboratory & \\
1 & POLCON & A columnar vector of real monotonically increasing down the
column values that defines the polymer concentration in the solution
surrounding the rock. The first entry should be zero to define a no
polymer and no salt concentration data set. POLCON should only be given
for the first entry of the POLCON/POLRATIO set and skipped until another
POLCON/POLRATIO table is entered. & None \\
lb/stb & kg/sm3 & gm/scc & \\
2 & POLRATIO & A columnar vector of real increasing down the column
values that defines the mass of adsorbed polymer per unit mass of rock
of the saturated concentration of polymer adsorbed by the rock for a
given POLCON and the salt concentration given by SALTCON on the ADSALNOD
keyword in the PROPS section. The first table data set entry should be
zero to define a no polymer and no salt concentration data set.
Subsequent POLRATIO values define the POLCON/POLRATIO combinations for a
given salt concentration as listed (and in the same order) by the
SALTCON variable on the ADSALNOD keyword in the PROPS section. Each
POLCON/POLRATIO/SALT data sets should be terminated by a ``/'' & None \\
lb/lb & kg/kg & gm/gm & \\
Notes: & & & \\
\end{longtable}

\emph{Table 8.100: PLYADSS Keyword Description}

See also the PLYADS keyword in the PROPS section to also define rock
polymer adsorption tables when the polymer concentration is not a
function of salinity.

\subsubsection{Example}\label{example-312}

\begin{verbatim}
--
--       SETS SALT CONCENTRATION FOR POLYMER SOLUTION ADSORPTION
--       VIA SATNUM ARRAY ALLOCATION
--
--       SALT
--
ADSALNOD
         1.0
         5.0
         10.5
         25.0      / SATNUM TABLE NO. 01
--
--       POLYMER ROCK ADSORPTION WITH SALT DEPENDENCY TABLE
--
PLYADSS
--       POLYMER    POLYMER
--       POLCON     POLRATIO
--       -------    --------
             0.0     0.00000
                     0.00000
                     0.00000
                     0.00000                               / TABLE NO. 01
--       POLYMER    POLYMER
--       POLCON     POLRATIO
--       -------    --------
             1.0     0.00002
                     0.00003
                     0.00004
                     0.00005                               / TABLE NO. 02

--       POLYMER    POLYMER
--       POLCON     POLRATIO
--       -------    --------
             2.0     0.00003
                     0.00004
                     0.00005
                     0.00006                               / TABLE NO. 03
--       POLYMER    POLYMER
--       POLCON     POLRATIO
--       -------    --------
             3.0     0.00004
                     0.00005
                     0.00006
                     0.00007                               / TABLE NO. 04
\end{verbatim}

The above example defines four polymer rock adsorption tables for four
salt concentration on the ADSALNOD keyword, assuming NTSFUN equals one
and NSSFUN is greater than or equal to four on the TABDIMS keyword in
the RUNSPEC section.

There is no terminating ``/'' for this keyword.

\subsection{PLYATEMP -- Define Polymer Adsorption Table
Temperature}\label{plyatemp-define-polymer-adsorption-table-temperature}

\begin{longtable}[]{@{}
  >{\raggedright\arraybackslash}p{(\linewidth - 14\tabcolsep) * \real{0.1250}}
  >{\raggedright\arraybackslash}p{(\linewidth - 14\tabcolsep) * \real{0.1250}}
  >{\raggedright\arraybackslash}p{(\linewidth - 14\tabcolsep) * \real{0.1250}}
  >{\raggedright\arraybackslash}p{(\linewidth - 14\tabcolsep) * \real{0.1250}}
  >{\raggedright\arraybackslash}p{(\linewidth - 14\tabcolsep) * \real{0.1250}}
  >{\raggedright\arraybackslash}p{(\linewidth - 14\tabcolsep) * \real{0.1250}}
  >{\raggedright\arraybackslash}p{(\linewidth - 14\tabcolsep) * \real{0.1250}}
  >{\raggedright\arraybackslash}p{(\linewidth - 14\tabcolsep) * \real{0.1250}}@{}}
\toprule\noalign{}
\begin{minipage}[b]{\linewidth}\raggedright
RUNSPEC
\end{minipage} & \begin{minipage}[b]{\linewidth}\raggedright
GRID
\end{minipage} & \begin{minipage}[b]{\linewidth}\raggedright
EDIT
\end{minipage} & \begin{minipage}[b]{\linewidth}\raggedright
PROPS
\end{minipage} & \begin{minipage}[b]{\linewidth}\raggedright
REGIONS
\end{minipage} & \begin{minipage}[b]{\linewidth}\raggedright
SOLUTION
\end{minipage} & \begin{minipage}[b]{\linewidth}\raggedright
SUMMARY
\end{minipage} & \begin{minipage}[b]{\linewidth}\raggedright
SCHEDULE
\end{minipage} \\
\midrule\noalign{}
\endhead
\bottomrule\noalign{}
\endlastfoot
\end{longtable}

\subsubsection{Description}\label{description-567}

This keyword defines the polymer adsorption temperature for subsequent
polymer adsorption tables entered via the PLYADS and PLYADSS keywords in
the PROPS section. The Polymer option must have been activated by the
POLYMER keyword in the RUNSPEC section and the Thermal option invoked by
the THERMAL keyword, also in the RUNSPEC section.

This keyword is not supported by OPM Flow but would change the results
if supported so the simulation will be stopped.

\begin{longtable}[]{@{}
  >{\raggedright\arraybackslash}p{(\linewidth - 6\tabcolsep) * \real{0.2500}}
  >{\raggedright\arraybackslash}p{(\linewidth - 6\tabcolsep) * \real{0.2500}}
  >{\raggedright\arraybackslash}p{(\linewidth - 6\tabcolsep) * \real{0.2500}}
  >{\raggedright\arraybackslash}p{(\linewidth - 6\tabcolsep) * \real{0.2500}}@{}}
\toprule\noalign{}
\begin{minipage}[b]{\linewidth}\raggedright
No.
\end{minipage} & \begin{minipage}[b]{\linewidth}\raggedright
Name
\end{minipage} & \begin{minipage}[b]{\linewidth}\raggedright
Description
\end{minipage} & \begin{minipage}[b]{\linewidth}\raggedright
Default
\end{minipage} \\
\midrule\noalign{}
\endhead
\bottomrule\noalign{}
\endlastfoot
Field & Metric & Laboratory & \\
1 & PLYATEMP & Single real positive value that defines polymer
adsorption temperature for subsequent polymer adsorption tables. &
None \\
oF & oC & oC & \\
Notes: & & & \\
\end{longtable}

\emph{Table 8.101: PLYATEMP Keyword Description}

\subsubsection{Example}\label{example-313}

The example shows how to enter the polymer adsorption data using the
PLYADS keyword for two different temperatures.

\begin{verbatim}
--
--       RESERVOIR
--       TEMPERATURE
--       -----------
PLYATEMP
         60.0                                             / TEMPERATURE
--
--       POLYMER ROCK ADSORPTION TABLE
--
PLYADS
--       POLYMER    POLYMER
--       POLCON     POLRATIO
--       -------    --------
             0.0     0.00000
             2.0     0.00003
             4.0     0.00005
             6.0     0.00007
             8.0     0.00009
            10.0     0.00011
            12.0     0.00012
            14.0     0.00015                               / TABLE NO. 01

--
--       POLYMER    POLYMER
--       POLCON     POLRATIO
--       -------    --------
             0.0     0.00000
             3.0     0.00004
             5.0     0.00006
             7.0     0.00008
             8.0     0.00009
            10.0     0.00011                               / TABLE NO. 02
--
--       RESERVOIR
--       TEMPERATURE
--       -----------
PLYATEMP
         120.0                                            / TEMPERATURE
--
--       POLYMER ROCK ADSORPTION TABLE
--
PLYADS
--       POLYMER    POLYMER
--       POLCON     POLRATIO
--       -------    --------
             0.0     0.00000
             2.0     0.00003
             4.0     0.00005
             6.0     0.00007
             8.0     0.00009
            10.0     0.00011
            12.0     0.00012
            14.0     0.00015                               / TABLE NO. 01
--
--       POLYMER    POLYMER
--       POLCON     POLRATIO
--       -------    --------
             0.0     0.00000
             3.0     0.00004
             5.0     0.00006
             7.0     0.00008
             8.0     0.00009
            10.0     0.00011                               / TABLE NO. 02
\end{verbatim}

Here the first PLYATEMP keyword defines the temperature to be 60 oF for
the subsequent two polymer rock adsorption tables, assuming NTSFUN
equals four and NSSFUN is greater than or equal to eight on the TABDIMS
keyword in the RUNSPEC section. The next PLYATEMP keyword defines the
temperature to be 120 oF for the subsequent two polymer rock adsorption
tables.

\subsection{PLYCAMAX -- Define Polymer-Rock Maximum Adsorption by
Cell}\label{plycamax-define-polymer-rock-maximum-adsorption-by-cell}

\begin{longtable}[]{@{}
  >{\raggedright\arraybackslash}p{(\linewidth - 14\tabcolsep) * \real{0.1250}}
  >{\raggedright\arraybackslash}p{(\linewidth - 14\tabcolsep) * \real{0.1250}}
  >{\raggedright\arraybackslash}p{(\linewidth - 14\tabcolsep) * \real{0.1250}}
  >{\raggedright\arraybackslash}p{(\linewidth - 14\tabcolsep) * \real{0.1250}}
  >{\raggedright\arraybackslash}p{(\linewidth - 14\tabcolsep) * \real{0.1250}}
  >{\raggedright\arraybackslash}p{(\linewidth - 14\tabcolsep) * \real{0.1250}}
  >{\raggedright\arraybackslash}p{(\linewidth - 14\tabcolsep) * \real{0.1250}}
  >{\raggedright\arraybackslash}p{(\linewidth - 14\tabcolsep) * \real{0.1250}}@{}}
\toprule\noalign{}
\begin{minipage}[b]{\linewidth}\raggedright
RUNSPEC
\end{minipage} & \begin{minipage}[b]{\linewidth}\raggedright
GRID
\end{minipage} & \begin{minipage}[b]{\linewidth}\raggedright
EDIT
\end{minipage} & \begin{minipage}[b]{\linewidth}\raggedright
PROPS
\end{minipage} & \begin{minipage}[b]{\linewidth}\raggedright
REGIONS
\end{minipage} & \begin{minipage}[b]{\linewidth}\raggedright
SOLUTION
\end{minipage} & \begin{minipage}[b]{\linewidth}\raggedright
SUMMARY
\end{minipage} & \begin{minipage}[b]{\linewidth}\raggedright
SCHEDULE
\end{minipage} \\
\midrule\noalign{}
\endhead
\bottomrule\noalign{}
\endlastfoot
\end{longtable}

\subsubsection{Description}\label{description-568}

The PLYCAMAX keyword defines the maximum polymer-rock adsorption value
used in the calculation of the resistance factor for the water phase by
individual grid block, for when the Polymer option has been activated by
the POLYMER keyword in the RUNSPEC section. See the POLMAX parameter on
the PLYROCK keyword in the PROPS section for setting the property for
the whole grid.

This keyword is not supported by OPM Flow but would change the results
if supported so the simulation will be stopped.

\subsection{PLYDHFLF -- Define Polymer Thermal Degradation Half-Life
Tables}\label{plydhflf-define-polymer-thermal-degradation-half-life-tables}

\begin{longtable}[]{@{}
  >{\raggedright\arraybackslash}p{(\linewidth - 14\tabcolsep) * \real{0.1250}}
  >{\raggedright\arraybackslash}p{(\linewidth - 14\tabcolsep) * \real{0.1250}}
  >{\raggedright\arraybackslash}p{(\linewidth - 14\tabcolsep) * \real{0.1250}}
  >{\raggedright\arraybackslash}p{(\linewidth - 14\tabcolsep) * \real{0.1250}}
  >{\raggedright\arraybackslash}p{(\linewidth - 14\tabcolsep) * \real{0.1250}}
  >{\raggedright\arraybackslash}p{(\linewidth - 14\tabcolsep) * \real{0.1250}}
  >{\raggedright\arraybackslash}p{(\linewidth - 14\tabcolsep) * \real{0.1250}}
  >{\raggedright\arraybackslash}p{(\linewidth - 14\tabcolsep) * \real{0.1250}}@{}}
\toprule\noalign{}
\begin{minipage}[b]{\linewidth}\raggedright
RUNSPEC
\end{minipage} & \begin{minipage}[b]{\linewidth}\raggedright
GRID
\end{minipage} & \begin{minipage}[b]{\linewidth}\raggedright
EDIT
\end{minipage} & \begin{minipage}[b]{\linewidth}\raggedright
PROPS
\end{minipage} & \begin{minipage}[b]{\linewidth}\raggedright
REGIONS
\end{minipage} & \begin{minipage}[b]{\linewidth}\raggedright
SOLUTION
\end{minipage} & \begin{minipage}[b]{\linewidth}\raggedright
SUMMARY
\end{minipage} & \begin{minipage}[b]{\linewidth}\raggedright
SCHEDULE
\end{minipage} \\
\midrule\noalign{}
\endhead
\bottomrule\noalign{}
\endlastfoot
\end{longtable}

\subsubsection{Description}\label{description-569}

The PLYDHFLF keyword defines the polymer thermal degradation half-life
with respect to temperature functions for when the polymer option has
been activated by the POLYMER keyword in the RUNSPEC section.

This keyword is not supported by OPM Flow but would change the results
if supported so the simulation will be stopped.

\begin{longtable}[]{@{}
  >{\raggedright\arraybackslash}p{(\linewidth - 6\tabcolsep) * \real{0.2500}}
  >{\raggedright\arraybackslash}p{(\linewidth - 6\tabcolsep) * \real{0.2500}}
  >{\raggedright\arraybackslash}p{(\linewidth - 6\tabcolsep) * \real{0.2500}}
  >{\raggedright\arraybackslash}p{(\linewidth - 6\tabcolsep) * \real{0.2500}}@{}}
\toprule\noalign{}
\begin{minipage}[b]{\linewidth}\raggedright
No.
\end{minipage} & \begin{minipage}[b]{\linewidth}\raggedright
Name
\end{minipage} & \begin{minipage}[b]{\linewidth}\raggedright
Description
\end{minipage} & \begin{minipage}[b]{\linewidth}\raggedright
Default
\end{minipage} \\
\midrule\noalign{}
\endhead
\bottomrule\noalign{}
\endlastfoot
Field & Metric & Laboratory & \\
1 & TEMP & A columnar vector of real monotonically increasing down the
column values that defines the polymer thermal degradation temperature.
& None \\
oF & oC & oC & \\
2 & POLHFLF & A columnar vector of real values that defines the
corresponding polymer half-life. & None \\
days & days & hours & \\
Notes: & & & \\
\end{longtable}

\emph{Table 8.102: PLYDHFLF Keyword Description}

\subsubsection{Example}\label{example-314}

\begin{verbatim}
--
--       POLYMER THERMAL DEGRADATION HALF-LIFE TABLE
--
PLYDHFLF
--       POLYMER    POLYMER
--       TEMP       HALF-LIFE
--       -------    ---------
             0.0     365.000
            40.0     200.000
            80.0     150.000
           120.0     100.000                               / TABLE NO. 01
--       POLYMER    POLYMER
--       TEMP       HALF-LIFE
--       -------    --------
             0.0     365.000
            50.0     175.000
            75.0     140.000
           100.0     120.000
           125.0      90.000
           150.0      85.000                               / TABLE NO. 02
\end{verbatim}

The example defines two polymer thermal degradation half-life tables,
based on the NTPVT variable on the TABDIMS keyword in the RUNSPEC
section being equal to two and NPPVT variable on the same keyword being
greater than or equal to six.

\subsection{PLYESAL -- Define Polymer Effective Salinity
Coefficient}\label{plyesal-define-polymer-effective-salinity-coefficient}

\begin{longtable}[]{@{}
  >{\raggedright\arraybackslash}p{(\linewidth - 14\tabcolsep) * \real{0.1250}}
  >{\raggedright\arraybackslash}p{(\linewidth - 14\tabcolsep) * \real{0.1250}}
  >{\raggedright\arraybackslash}p{(\linewidth - 14\tabcolsep) * \real{0.1250}}
  >{\raggedright\arraybackslash}p{(\linewidth - 14\tabcolsep) * \real{0.1250}}
  >{\raggedright\arraybackslash}p{(\linewidth - 14\tabcolsep) * \real{0.1250}}
  >{\raggedright\arraybackslash}p{(\linewidth - 14\tabcolsep) * \real{0.1250}}
  >{\raggedright\arraybackslash}p{(\linewidth - 14\tabcolsep) * \real{0.1250}}
  >{\raggedright\arraybackslash}p{(\linewidth - 14\tabcolsep) * \real{0.1250}}@{}}
\toprule\noalign{}
\begin{minipage}[b]{\linewidth}\raggedright
RUNSPEC
\end{minipage} & \begin{minipage}[b]{\linewidth}\raggedright
GRID
\end{minipage} & \begin{minipage}[b]{\linewidth}\raggedright
EDIT
\end{minipage} & \begin{minipage}[b]{\linewidth}\raggedright
PROPS
\end{minipage} & \begin{minipage}[b]{\linewidth}\raggedright
REGIONS
\end{minipage} & \begin{minipage}[b]{\linewidth}\raggedright
SOLUTION
\end{minipage} & \begin{minipage}[b]{\linewidth}\raggedright
SUMMARY
\end{minipage} & \begin{minipage}[b]{\linewidth}\raggedright
SCHEDULE
\end{minipage} \\
\midrule\noalign{}
\endhead
\bottomrule\noalign{}
\endlastfoot
\end{longtable}

\subsubsection{Description}\label{description-570}

This keyword, PLYESAL, defines the polymer effective salinity
coefficient as well as enabling the effective salinity calculation for
polymer adsorption. The keyword should only be used if the BRINE keyword
has been declared to activate the brine phase, the ECLMC keyword to
enable the Multi-Component Brine model, and the POLYMER keyword has been
used to activate the polymer phase. All three keywords are in the
RUNSPEC section.

This keyword is not supported by OPM Flow but would change the results
if supported so the simulation will be stopped.

\subsection{PLYKRRF -- Define Polymer Rock Permeability Reduction by
Cell}\label{plykrrf-define-polymer-rock-permeability-reduction-by-cell}

\begin{longtable}[]{@{}
  >{\raggedright\arraybackslash}p{(\linewidth - 14\tabcolsep) * \real{0.1250}}
  >{\raggedright\arraybackslash}p{(\linewidth - 14\tabcolsep) * \real{0.1250}}
  >{\raggedright\arraybackslash}p{(\linewidth - 14\tabcolsep) * \real{0.1250}}
  >{\raggedright\arraybackslash}p{(\linewidth - 14\tabcolsep) * \real{0.1250}}
  >{\raggedright\arraybackslash}p{(\linewidth - 14\tabcolsep) * \real{0.1250}}
  >{\raggedright\arraybackslash}p{(\linewidth - 14\tabcolsep) * \real{0.1250}}
  >{\raggedright\arraybackslash}p{(\linewidth - 14\tabcolsep) * \real{0.1250}}
  >{\raggedright\arraybackslash}p{(\linewidth - 14\tabcolsep) * \real{0.1250}}@{}}
\toprule\noalign{}
\begin{minipage}[b]{\linewidth}\raggedright
RUNSPEC
\end{minipage} & \begin{minipage}[b]{\linewidth}\raggedright
GRID
\end{minipage} & \begin{minipage}[b]{\linewidth}\raggedright
EDIT
\end{minipage} & \begin{minipage}[b]{\linewidth}\raggedright
PROPS
\end{minipage} & \begin{minipage}[b]{\linewidth}\raggedright
REGIONS
\end{minipage} & \begin{minipage}[b]{\linewidth}\raggedright
SOLUTION
\end{minipage} & \begin{minipage}[b]{\linewidth}\raggedright
SUMMARY
\end{minipage} & \begin{minipage}[b]{\linewidth}\raggedright
SCHEDULE
\end{minipage} \\
\midrule\noalign{}
\endhead
\bottomrule\noalign{}
\endlastfoot
\end{longtable}

\subsubsection{Description}\label{description-571}

The PLYKRRF keyword defines the polymer rock permeability reduction
factor to the water phase by individual cell, for when the Polymer
option has been activated by the POLYMER keyword in the RUNSPEC section.
PLYKRRF should consist of an array of real positive values that are
greater than or equal to one. See the PERMFAC parameter on the PLYROCK
keyword in the PROPS section for setting the property for the whole
grid.

This keyword is not supported by OPM Flow but would change the results
if supported so the simulation will be stopped.

\subsection{PLYMAX -- Define Polymer-Salt Viscosity Mixing
Concentrations}\label{plymax-define-polymer-salt-viscosity-mixing-concentrations}

\begin{longtable}[]{@{}
  >{\raggedright\arraybackslash}p{(\linewidth - 14\tabcolsep) * \real{0.1250}}
  >{\raggedright\arraybackslash}p{(\linewidth - 14\tabcolsep) * \real{0.1250}}
  >{\raggedright\arraybackslash}p{(\linewidth - 14\tabcolsep) * \real{0.1250}}
  >{\raggedright\arraybackslash}p{(\linewidth - 14\tabcolsep) * \real{0.1250}}
  >{\raggedright\arraybackslash}p{(\linewidth - 14\tabcolsep) * \real{0.1250}}
  >{\raggedright\arraybackslash}p{(\linewidth - 14\tabcolsep) * \real{0.1250}}
  >{\raggedright\arraybackslash}p{(\linewidth - 14\tabcolsep) * \real{0.1250}}
  >{\raggedright\arraybackslash}p{(\linewidth - 14\tabcolsep) * \real{0.1250}}@{}}
\toprule\noalign{}
\begin{minipage}[b]{\linewidth}\raggedright
RUNSPEC
\end{minipage} & \begin{minipage}[b]{\linewidth}\raggedright
GRID
\end{minipage} & \begin{minipage}[b]{\linewidth}\raggedright
EDIT
\end{minipage} & \begin{minipage}[b]{\linewidth}\raggedright
PROPS
\end{minipage} & \begin{minipage}[b]{\linewidth}\raggedright
REGIONS
\end{minipage} & \begin{minipage}[b]{\linewidth}\raggedright
SOLUTION
\end{minipage} & \begin{minipage}[b]{\linewidth}\raggedright
SUMMARY
\end{minipage} & \begin{minipage}[b]{\linewidth}\raggedright
SCHEDULE
\end{minipage} \\
\midrule\noalign{}
\endhead
\bottomrule\noalign{}
\endlastfoot
\end{longtable}

\subsubsection{Description}\label{description-572}

The PLYMAX keyword defines maximum polymer and salt concentrations that
are to be used in the mixing parameter calculation of the fluid
component viscosities, for when the Polymer option has been activated by
the POLYMER keyword in the RUNSPEC section.

Note that If the BRINE option has not be activated by the BRINE keyword
in the RUNSPEC section, then the salt concentrations in the second
column are ignored.

\begin{longtable}[]{@{}
  >{\raggedright\arraybackslash}p{(\linewidth - 6\tabcolsep) * \real{0.2500}}
  >{\raggedright\arraybackslash}p{(\linewidth - 6\tabcolsep) * \real{0.2500}}
  >{\raggedright\arraybackslash}p{(\linewidth - 6\tabcolsep) * \real{0.2500}}
  >{\raggedright\arraybackslash}p{(\linewidth - 6\tabcolsep) * \real{0.2500}}@{}}
\toprule\noalign{}
\begin{minipage}[b]{\linewidth}\raggedright
No.
\end{minipage} & \begin{minipage}[b]{\linewidth}\raggedright
Name
\end{minipage} & \begin{minipage}[b]{\linewidth}\raggedright
Description
\end{minipage} & \begin{minipage}[b]{\linewidth}\raggedright
Default
\end{minipage} \\
\midrule\noalign{}
\endhead
\bottomrule\noalign{}
\endlastfoot
Field & Metric & Laboratory & \\
1 & POLCON & A real value that defines the polymer concentration in the
solution which is used to calculate maximum polymer fluid component
viscosity. & None \\
lb/stb & kg/sm3 & gm/scc & \\
2 & SALTCON & A real value that defines the salt concentration in the
solution which is used to calculate maximum polymer fluid component
viscosity. Note that If the BRINE option has not been activated by the
BRINE keyword in the RUNSPEC section, then this variable is ignored;
however, there should still be dummy entries in this case. This variable
is ignored as the BRINE and POLYMER combination is not implemented in
OPM Flow. & None \\
lb/stb & kg/sm3 & gm/scc & \\
Notes: & & & \\
\end{longtable}

\emph{Table 8.103: PLYMAX Keyword Description}

\begin{tcolorbox}[enhanced jigsaw, title=\textcolor{quarto-callout-note-color}{\faInfo}\hspace{0.5em}{Note}, coltitle=black, breakable, toprule=.15mm, colbacktitle=quarto-callout-note-color!10!white, leftrule=.75mm, opacitybacktitle=0.6, colback=white, bottomtitle=1mm, toptitle=1mm, titlerule=0mm, left=2mm, colframe=quarto-callout-note-color-frame, arc=.35mm, rightrule=.15mm, bottomrule=.15mm, opacityback=0]

Currently, combining the BRINE and POLYMER models is not implemented in
OPM Flow, and therefore SALTCON parameter on the PLYMAX keyword is
ignored.

\end{tcolorbox}

\subsubsection{Example}\label{example-315}

\begin{verbatim}
--
--       POLYMER-SALT VISCOSITY MIXING CONCENTRATIONS
--
PLYMAX
--       POLYMER    SALT
--       POLCON     SALTCON
--       -------    --------
         0.0100      0.0500                                / TABLE NO. 01
         0.0075      0.0400                                / TABLE NO. 02
         0.0050      0.0300                                / TABLE NO. 03
\end{verbatim}

The above example defines three polymer-salt viscosity mixing
concentrations, based on the NPLMIX variable on the REGDIMS keyword in
the RUNSPEC section being equal to three.

\subsection{PLYMWINJ -- Polymer Molecular Weight Model Throughput and
Velocity
Table}\label{plymwinj-polymer-molecular-weight-model-throughput-and-velocity-table}

\begin{longtable}[]{@{}
  >{\raggedright\arraybackslash}p{(\linewidth - 14\tabcolsep) * \real{0.1250}}
  >{\raggedright\arraybackslash}p{(\linewidth - 14\tabcolsep) * \real{0.1250}}
  >{\raggedright\arraybackslash}p{(\linewidth - 14\tabcolsep) * \real{0.1250}}
  >{\raggedright\arraybackslash}p{(\linewidth - 14\tabcolsep) * \real{0.1250}}
  >{\raggedright\arraybackslash}p{(\linewidth - 14\tabcolsep) * \real{0.1250}}
  >{\raggedright\arraybackslash}p{(\linewidth - 14\tabcolsep) * \real{0.1250}}
  >{\raggedright\arraybackslash}p{(\linewidth - 14\tabcolsep) * \real{0.1250}}
  >{\raggedright\arraybackslash}p{(\linewidth - 14\tabcolsep) * \real{0.1250}}@{}}
\toprule\noalign{}
\begin{minipage}[b]{\linewidth}\raggedright
RUNSPEC
\end{minipage} & \begin{minipage}[b]{\linewidth}\raggedright
GRID
\end{minipage} & \begin{minipage}[b]{\linewidth}\raggedright
EDIT
\end{minipage} & \begin{minipage}[b]{\linewidth}\raggedright
PROPS
\end{minipage} & \begin{minipage}[b]{\linewidth}\raggedright
REGIONS
\end{minipage} & \begin{minipage}[b]{\linewidth}\raggedright
SOLUTION
\end{minipage} & \begin{minipage}[b]{\linewidth}\raggedright
SUMMARY
\end{minipage} & \begin{minipage}[b]{\linewidth}\raggedright
SCHEDULE
\end{minipage} \\
\midrule\noalign{}
\endhead
\bottomrule\noalign{}
\endlastfoot
\end{longtable}

\subsubsection{Description}\label{description-573}

This keyword, PLYMWINJ, describes the relationship of the injected
polymer molecular weight as a function of polymer throughput and polymer
velocity, for the simulator's Polymer Molecular Weight Transport option,
that uses the polymer molecular weight in calculating the polymer
viscosity. The table is a two dimensional table that relates the polymer
throughput values and velocity values to derive the resulting molecular
weight of the injected polymer, which is then used via the PLYVHM
keyword in the PROPS section, to derive the polymer molecular weight
scaled viscosity.

This keyword should only be used if the POLYMER and POLYMW keywords in
the RUNSPEC section are also activated.

\begin{tcolorbox}[enhanced jigsaw, title=\textcolor{quarto-callout-note-color}{\faInfo}\hspace{0.5em}{Note}, coltitle=black, breakable, toprule=.15mm, colbacktitle=quarto-callout-note-color!10!white, leftrule=.75mm, opacitybacktitle=0.6, colback=white, bottomtitle=1mm, toptitle=1mm, titlerule=0mm, left=2mm, colframe=quarto-callout-note-color-frame, arc=.35mm, rightrule=.15mm, bottomrule=.15mm, opacityback=0]

This is an OPM Flow specific keyword that employs an alternative polymer
flood model based on a Polymer Molecular Weight Transport equation, that
is not available in the commercial simulator. The model has been tested
using metric units; however, using either field or laboratory units with
the option should be considered experimental.

\end{tcolorbox}

\begin{longtable}[]{@{}
  >{\raggedright\arraybackslash}p{(\linewidth - 6\tabcolsep) * \real{0.2500}}
  >{\raggedright\arraybackslash}p{(\linewidth - 6\tabcolsep) * \real{0.2500}}
  >{\raggedright\arraybackslash}p{(\linewidth - 6\tabcolsep) * \real{0.2500}}
  >{\raggedright\arraybackslash}p{(\linewidth - 6\tabcolsep) * \real{0.2500}}@{}}
\toprule\noalign{}
\begin{minipage}[b]{\linewidth}\raggedright
No.
\end{minipage} & \begin{minipage}[b]{\linewidth}\raggedright
Name
\end{minipage} & \begin{minipage}[b]{\linewidth}\raggedright
Description
\end{minipage} & \begin{minipage}[b]{\linewidth}\raggedright
Default
\end{minipage} \\
\midrule\noalign{}
\endhead
\bottomrule\noalign{}
\endlastfoot
Field & Metric & Laboratory & \\
1-1 & PLYMWNUM & A positive integer value greater than zero and less
than or equal to the NTPMWINJ variable, as defined on the PINTDIMS
keyword in the RUNSPEC section, that defines the PLYMWINJ Polymer
Molecular Weight Model throughput and velocity table number. & None \\
2-1 & THRUPUT & A real positive monotonically increasing vector, that
defines the polymer throughput values. The first entry should be zero to
define a no throughput data set, and each vector record should be on a
separate line terminated by a ``/''. & None \\
feet3/feet2 & m3/m2 & cm3/cm2 & \\
3-1 & VELOCITY & A real positive monotonically increasing vector, that
defines the polymer velocity values. The first entry should be zero to
define a no velocity data set, and each vector record should be on a
separate line terminated by a ``/''. & None \\
feet/day & m/day & cm/hour & \\
4-1 & POLYMW & A series of real positive vectors representing the
polymer molecular weights for all combinations of the polymer throughput
values (THRUPUT) and velocity values (VELOCITY), organized as a series
of vectors POLYMW(THRUPUT, VELOCITY). Thus, the first vector represents
the molecular weights of the first THRUPUT value and each entry in the
vector is the corresponding molecular weight of the associated VELOCITY
vector. Each vector should be on a separate line and should be
terminated by a ``/''. Thus, if THRUPUT has three entries and VELOCITY
has four, then there should be three vectors, with each vector
containing four elements representing molecular weight values, as a
function of THRUPUT and VELOCITY. & None \\
lb/lb-M & kg/kg-M & gm/gm-M & \\
Notes: & & & \\
\end{longtable}

\emph{Table 8.104: PLYMWINJ Keyword Description}

Unlike other PROPS section table keywords, that enable multiple tables
following the keyword to be entered, the PLYMWINJ keyword requires that
the keyword itself must be repeated for each table, as is shown in the
example on the following page.

See also the WPMITAB keyword in the SCHEDULE section, that assigns the
PLYMWINJ tables to the water injection wells, as well as the SKPRWAT,
SKPRPOLY, and PLYVMH keywords, in the PROPS section, that are the
additional property keywords required for the Polymer Molecular Weight
Transport option.

The WSKPTAB keyword in the SCHEDULE section may be used to assign the
SKPRPOLY and SKPRWAT tables to water injections wells, that enable the
calculation of the wellbore skin pressure based on the fluids being
injected.

Note that the standard polymer property data keywords: PLYROCK, PLYADS,
PLYMAX, etc., are still required to fully describe the polymer fluid.

\subsubsection{Example}\label{example-316}

Given NTPMWINJ equals two on the PINTDIMS keyword in the RUNSPEC
section, then two PLYWINJ tables are required to be entered:

\begin{verbatim}
--
--       POLYMER MOLECULAR WEIGHT MODEL THROUGHPUT AND VELOCITY TABLE
--       (OPM FLOW PROPS KEYWORD)
--
PLYMWINJ
         1                             / TABLE NUMBER
--
--       THROUGHPUT VALUES
--
         0.0     200.0   400.0         /
--
--       VELOCITY VALUES
--
         0.0     50.0    80.0    100.0 /
--
--       POLYMW VALUES
--
         20.0    19.0    18.0    16.0  / POLYMW(OUTPUT=1, VELOCITY=1 TO N)
         20.0    16.0    14.0    12.0  / POLYMW(OUTPUT=2, VELOCITY=1 TO N)
         20.0    12.0     8.0     4.0  / POLYMW(OUTPUT=3, VELOCITY=1 TO N)
/
PLYMWINJ
         2                             / TABLE NUMBER
--
--       THROUGHPUT VALUES
--
         0.0     200.0   400.0         /
--
--       VELOCITY VALUES
--
         0.0     50.0    70.0    100.0 /
--
--       POLYMW VALUES
--
         20.0    19.0    18.0    16.0  / POLYMW(OUTPUT=1, VELOCITY=1 TO N)
         20.0    16.0    14.0    12.0  / POLYMW(OUTPUT=2, VELOCITY=1 TO N)
         20.0    12.0     8.0     4.0  / POLYMW(OUTPUT=3, VELOCITY=1 TO N)
/
\end{verbatim}

As mentioned previously, the PLYMWINJ keyword requires that the keyword
itself must be repeated for each table, as is shown in the above
example.

\subsection{PLYRMDEN -- Define Polymer Model In Situ Rock
Density}\label{plyrmden-define-polymer-model-in-situ-rock-density}

\begin{longtable}[]{@{}
  >{\raggedright\arraybackslash}p{(\linewidth - 14\tabcolsep) * \real{0.1250}}
  >{\raggedright\arraybackslash}p{(\linewidth - 14\tabcolsep) * \real{0.1250}}
  >{\raggedright\arraybackslash}p{(\linewidth - 14\tabcolsep) * \real{0.1250}}
  >{\raggedright\arraybackslash}p{(\linewidth - 14\tabcolsep) * \real{0.1250}}
  >{\raggedright\arraybackslash}p{(\linewidth - 14\tabcolsep) * \real{0.1250}}
  >{\raggedright\arraybackslash}p{(\linewidth - 14\tabcolsep) * \real{0.1250}}
  >{\raggedright\arraybackslash}p{(\linewidth - 14\tabcolsep) * \real{0.1250}}
  >{\raggedright\arraybackslash}p{(\linewidth - 14\tabcolsep) * \real{0.1250}}@{}}
\toprule\noalign{}
\begin{minipage}[b]{\linewidth}\raggedright
RUNSPEC
\end{minipage} & \begin{minipage}[b]{\linewidth}\raggedright
GRID
\end{minipage} & \begin{minipage}[b]{\linewidth}\raggedright
EDIT
\end{minipage} & \begin{minipage}[b]{\linewidth}\raggedright
PROPS
\end{minipage} & \begin{minipage}[b]{\linewidth}\raggedright
REGIONS
\end{minipage} & \begin{minipage}[b]{\linewidth}\raggedright
SOLUTION
\end{minipage} & \begin{minipage}[b]{\linewidth}\raggedright
SUMMARY
\end{minipage} & \begin{minipage}[b]{\linewidth}\raggedright
SCHEDULE
\end{minipage} \\
\midrule\noalign{}
\endhead
\bottomrule\noalign{}
\endlastfoot
\end{longtable}

\subsubsection{Description}\label{description-574}

The PLYRMDEN keyword defines the in situ rock density at reservoir
conditions by individual cell, for when the Polymer option has been
activated by the POLYMER keyword in the RUNSPEC section. PLYRMDEN should
consist of an array of real positive values. See the DENSITY parameter
on the PLYROCK keyword in the PROPS section for setting the property for
the whole grid.

\subsection{PLYROCK -- Define Polymer-Rock
Properties}\label{plyrock-define-polymer-rock-properties}

\begin{longtable}[]{@{}
  >{\raggedright\arraybackslash}p{(\linewidth - 14\tabcolsep) * \real{0.1250}}
  >{\raggedright\arraybackslash}p{(\linewidth - 14\tabcolsep) * \real{0.1250}}
  >{\raggedright\arraybackslash}p{(\linewidth - 14\tabcolsep) * \real{0.1250}}
  >{\raggedright\arraybackslash}p{(\linewidth - 14\tabcolsep) * \real{0.1250}}
  >{\raggedright\arraybackslash}p{(\linewidth - 14\tabcolsep) * \real{0.1250}}
  >{\raggedright\arraybackslash}p{(\linewidth - 14\tabcolsep) * \real{0.1250}}
  >{\raggedright\arraybackslash}p{(\linewidth - 14\tabcolsep) * \real{0.1250}}
  >{\raggedright\arraybackslash}p{(\linewidth - 14\tabcolsep) * \real{0.1250}}@{}}
\toprule\noalign{}
\begin{minipage}[b]{\linewidth}\raggedright
RUNSPEC
\end{minipage} & \begin{minipage}[b]{\linewidth}\raggedright
GRID
\end{minipage} & \begin{minipage}[b]{\linewidth}\raggedright
EDIT
\end{minipage} & \begin{minipage}[b]{\linewidth}\raggedright
PROPS
\end{minipage} & \begin{minipage}[b]{\linewidth}\raggedright
REGIONS
\end{minipage} & \begin{minipage}[b]{\linewidth}\raggedright
SOLUTION
\end{minipage} & \begin{minipage}[b]{\linewidth}\raggedright
SUMMARY
\end{minipage} & \begin{minipage}[b]{\linewidth}\raggedright
SCHEDULE
\end{minipage} \\
\midrule\noalign{}
\endhead
\bottomrule\noalign{}
\endlastfoot
\end{longtable}

\subsubsection{Description}\label{description-575}

The PLYROCK keyword defines rock properties for when the Polymer option
has been activated by the POLYMER keyword in the RUNSPEC section.

\begin{longtable}[]{@{}
  >{\raggedright\arraybackslash}p{(\linewidth - 6\tabcolsep) * \real{0.2500}}
  >{\raggedright\arraybackslash}p{(\linewidth - 6\tabcolsep) * \real{0.2500}}
  >{\raggedright\arraybackslash}p{(\linewidth - 6\tabcolsep) * \real{0.2500}}
  >{\raggedright\arraybackslash}p{(\linewidth - 6\tabcolsep) * \real{0.2500}}@{}}
\toprule\noalign{}
\begin{minipage}[b]{\linewidth}\raggedright
No.
\end{minipage} & \begin{minipage}[b]{\linewidth}\raggedright
Name
\end{minipage} & \begin{minipage}[b]{\linewidth}\raggedright
Description
\end{minipage} & \begin{minipage}[b]{\linewidth}\raggedright
Default
\end{minipage} \\
\midrule\noalign{}
\endhead
\bottomrule\noalign{}
\endlastfoot
Field & Metric & Laboratory & \\
1 & PSPACE & A real positive value that is greater than or equal to zero
and less the maximum water saturation and less than one, that defines
available pore space for this rock type. & None \\
dimensionless & dimensionless & dimensionless & \\
2 & PERMFAC & A real positive value that is greater than or equal to one
that defines decrease in the rock permeability to the water phase when
the maximum amount of polymer has been adsorbed. & None \\
dimensionless & dimensionless & dimensionless & \\
3 & DENSITY & A real value that defines the rock in-situ density, that
is at reservoir conditions. & None \\
lb/rb & kg/rm3 & gm/rcc & \\
4 & ADINDX & A positive integer of 1 or 2 that defines the polymer
desorption option. & Defined \\
dimensionless 1 & dimensionless 1 & dimensionless 1 & \\
5 & POLMAX & A real positive non-zero value that defines the maximum
polymer adsorption to be used in the calculation of the resistance
factor for the water phase. & None \\
lb/lb & kg/kg & gm/gm & \\
Notes: & & & \\
\end{longtable}

\emph{Table 8.105: PLYROCK Keyword Description}

\subsubsection{Example}\label{example-317}

\begin{verbatim}
--
--       POLYMER-ROCK PROPERTIES
--
PLYROCK
--       PORE      PERM      INSITU    DESORP  MAX
--       SPACE     FACTOR    DENSITY   OPTN    POLY
--       ------   --------   -------   ------  -------
         0.1200    1.7500    1800.0      1     0.00012     / TABLE NO. 01
         0.1300    1.8500    1980.0      2     0.00015     / TABLE NO. 02
         0.1500    1.9500    2005.0      1     0.00014     / TABLE NO. 03
\end{verbatim}

The above example defines three polymer-rock tables, based on the NTSFUN
variable on the TABDIMS keyword in the RUNSPEC section being equal to
three.

There is no terminating ``/'' for this keyword.

\subsection{PLYSHEAR -- Activate and Define Polymer Shearing
Parameters}\label{plyshear-activate-and-define-polymer-shearing-parameters}

\begin{longtable}[]{@{}
  >{\raggedright\arraybackslash}p{(\linewidth - 14\tabcolsep) * \real{0.1250}}
  >{\raggedright\arraybackslash}p{(\linewidth - 14\tabcolsep) * \real{0.1250}}
  >{\raggedright\arraybackslash}p{(\linewidth - 14\tabcolsep) * \real{0.1250}}
  >{\raggedright\arraybackslash}p{(\linewidth - 14\tabcolsep) * \real{0.1250}}
  >{\raggedright\arraybackslash}p{(\linewidth - 14\tabcolsep) * \real{0.1250}}
  >{\raggedright\arraybackslash}p{(\linewidth - 14\tabcolsep) * \real{0.1250}}
  >{\raggedright\arraybackslash}p{(\linewidth - 14\tabcolsep) * \real{0.1250}}
  >{\raggedright\arraybackslash}p{(\linewidth - 14\tabcolsep) * \real{0.1250}}@{}}
\toprule\noalign{}
\begin{minipage}[b]{\linewidth}\raggedright
RUNSPEC
\end{minipage} & \begin{minipage}[b]{\linewidth}\raggedright
GRID
\end{minipage} & \begin{minipage}[b]{\linewidth}\raggedright
EDIT
\end{minipage} & \begin{minipage}[b]{\linewidth}\raggedright
PROPS
\end{minipage} & \begin{minipage}[b]{\linewidth}\raggedright
REGIONS
\end{minipage} & \begin{minipage}[b]{\linewidth}\raggedright
SOLUTION
\end{minipage} & \begin{minipage}[b]{\linewidth}\raggedright
SUMMARY
\end{minipage} & \begin{minipage}[b]{\linewidth}\raggedright
SCHEDULE
\end{minipage} \\
\midrule\noalign{}
\endhead
\bottomrule\noalign{}
\endlastfoot
\end{longtable}

\subsubsection{Description}\label{description-576}

The PLYSHEAR keyword activates and defines the polymer shear
thinning-thickening option for when the polymer option has been
activated by the POLYMER keyword in the RUNSPEC section.

This keyword is not supported by OPM Flow but would change the results
if supported so the simulation will be stopped.

\begin{longtable}[]{@{}
  >{\raggedright\arraybackslash}p{(\linewidth - 6\tabcolsep) * \real{0.2500}}
  >{\raggedright\arraybackslash}p{(\linewidth - 6\tabcolsep) * \real{0.2500}}
  >{\raggedright\arraybackslash}p{(\linewidth - 6\tabcolsep) * \real{0.2500}}
  >{\raggedright\arraybackslash}p{(\linewidth - 6\tabcolsep) * \real{0.2500}}@{}}
\toprule\noalign{}
\begin{minipage}[b]{\linewidth}\raggedright
No.
\end{minipage} & \begin{minipage}[b]{\linewidth}\raggedright
Name
\end{minipage} & \begin{minipage}[b]{\linewidth}\raggedright
Description
\end{minipage} & \begin{minipage}[b]{\linewidth}\raggedright
Default
\end{minipage} \\
\midrule\noalign{}
\endhead
\bottomrule\noalign{}
\endlastfoot
Field & Metric & Laboratory & \\
1 & VELOCITY & A columnar vector of real monotonically increasing down
the column values that defines the water-polymer flow velocity. The
VELOCITY value for the first row in the table should be zero. & None \\
feet/day & m/day & cm/hour & \\
2 & VISFAC & A columnar vector of real values that defines a factor that
scales the effective water and polymer viscosities for when shear
thinning-thickening of the polymer occurs. Normally VISFAC value for the
first row in the table should be one. & None \\
dimensionless & dimensionless & dimensionless & \\
Notes: & & & \\
\end{longtable}

\emph{Table 8.106: PLYSHEAR Keyword Description}

\subsubsection{Example}\label{example-318}

\begin{verbatim}
--
--       ACTIVATE AND DEFINE POLYMER SHEARING PARAMETERS
--
PLYSHEAR
--       WAT-POLY    VISCOSITY
--       VELOCITY    FACTOR
--       --------    ---------
             0.0       1.000
             1.0       0.900
             3.0       0.800
             6.0       0.700                               / TABLE NO. 01
--       WAT-POLY    VISCOSITY
--       VELOCITY    FACTOR
--       --------    ---------
             0.0       1.000
             1.0       0.900
             2.0       0.800
             4.0       0.750
             6.0       0.700
             8.0       0.650                               / TABLE NO. 02
\end{verbatim}

The above example activates the polymer shear thinning-thickening option
and defines two polymer shear thinning-thickening tables, based on the
NTPVT variable on the TABDIMS keyword in the RUNSPEC section being equal
to two and NPPVT variable on the same keyword being greater than or
equal to six.

\subsection{PLYSHLOG -- Activate and Define the Polymer Shearing
Logarithmic
Parameters}\label{plyshlog-activate-and-define-the-polymer-shearing-logarithmic-parameters}

\begin{longtable}[]{@{}
  >{\raggedright\arraybackslash}p{(\linewidth - 14\tabcolsep) * \real{0.1250}}
  >{\raggedright\arraybackslash}p{(\linewidth - 14\tabcolsep) * \real{0.1250}}
  >{\raggedright\arraybackslash}p{(\linewidth - 14\tabcolsep) * \real{0.1250}}
  >{\raggedright\arraybackslash}p{(\linewidth - 14\tabcolsep) * \real{0.1250}}
  >{\raggedright\arraybackslash}p{(\linewidth - 14\tabcolsep) * \real{0.1250}}
  >{\raggedright\arraybackslash}p{(\linewidth - 14\tabcolsep) * \real{0.1250}}
  >{\raggedright\arraybackslash}p{(\linewidth - 14\tabcolsep) * \real{0.1250}}
  >{\raggedright\arraybackslash}p{(\linewidth - 14\tabcolsep) * \real{0.1250}}@{}}
\toprule\noalign{}
\begin{minipage}[b]{\linewidth}\raggedright
RUNSPEC
\end{minipage} & \begin{minipage}[b]{\linewidth}\raggedright
GRID
\end{minipage} & \begin{minipage}[b]{\linewidth}\raggedright
EDIT
\end{minipage} & \begin{minipage}[b]{\linewidth}\raggedright
PROPS
\end{minipage} & \begin{minipage}[b]{\linewidth}\raggedright
REGIONS
\end{minipage} & \begin{minipage}[b]{\linewidth}\raggedright
SOLUTION
\end{minipage} & \begin{minipage}[b]{\linewidth}\raggedright
SUMMARY
\end{minipage} & \begin{minipage}[b]{\linewidth}\raggedright
SCHEDULE
\end{minipage} \\
\midrule\noalign{}
\endhead
\bottomrule\noalign{}
\endlastfoot
\end{longtable}

\subsubsection{Description}\label{description-577}

This keyword activates and defines the parameters for the
logarithm-based polymer shear thinning/thickening option.

\begin{longtable}[]{@{}
  >{\raggedright\arraybackslash}p{(\linewidth - 6\tabcolsep) * \real{0.2500}}
  >{\raggedright\arraybackslash}p{(\linewidth - 6\tabcolsep) * \real{0.2500}}
  >{\raggedright\arraybackslash}p{(\linewidth - 6\tabcolsep) * \real{0.2500}}
  >{\raggedright\arraybackslash}p{(\linewidth - 6\tabcolsep) * \real{0.2500}}@{}}
\toprule\noalign{}
\begin{minipage}[b]{\linewidth}\raggedright
No.
\end{minipage} & \begin{minipage}[b]{\linewidth}\raggedright
Name
\end{minipage} & \begin{minipage}[b]{\linewidth}\raggedright
Description
\end{minipage} & \begin{minipage}[b]{\linewidth}\raggedright
Default
\end{minipage} \\
\midrule\noalign{}
\endhead
\bottomrule\noalign{}
\endlastfoot
Field & Metric & Laboratory & \\
1-1 & POLCON & A real positive value that defines the reference polymer
concentration for the VELOCITY and VISFAC data for this keyword. &
None \\
lb/stb & kg/sm3 & gm/scc & \\
1-2 & SALTCON & A real positive value that defines the reference salt
concentration for the VELOCITY and VISFAC data for this keyword. Note
that If the BRINE option has not been activated by the BRINE keyword in
the RUNSPEC section, then this variable is ignored. This variable is
ignored as the BRINE option is not implemented in OPM Flow. & None \\
lb/stb & kg/sm3 & gm/scc & \\
1-3 & TEMP & A real positive value defines the reference polymer
temperature for the VELOCITY and VISFAC data for this keyword. Note that
If the TEMP option has not been activated by the TEMP keyword in the
RUNSPEC section, then this variable is ignored. This variable is ignored
as the TEMP and POLYMER options combination is not implemented in OPM
Flow. & None \\
oF & oC & oC & \\
1-4 & / & Record terminated by a ``/'' & Not Applicable \\
2-1 & VELOCITY & A columnar vector of real monotonically increasing down
the column values that defines the water-polymer flow velocity for the
reference conditions of POLCON, SALTCON and TEMP. The VELOCITY value for
the first row in the table should be a very small value that is greater
than zero and less than 1 x 10-6. & None \\
feet/day & m/day & cm/hour & \\
2-2 & VISFAC & A columnar vector of real positive values that define the
dimensionless shear effect multiplier for the given VELOCITY entry for
the reference conditions of POLCON, SALTCON and TEMP. Normally VISFAC
value for the first row in the table should be one. & None \\
dimensionless & dimensionless & dimensionless & \\
1-4 & / & Record terminated by a ``/'' & Not Applicable \\
Notes: & & & \\
\end{longtable}

\emph{Table 8.107: PLYSHLOG Keyword Description}

See the PLYSHEAR keyword for the alternative polymer shear
thinning/thickening option that is also implemented in OPM Flow.

\subsubsection{Example}\label{example-319}

The following example show how to enter two PLYSHLOG tables given that
the NTPVT variable on the TABDIMS keyword in the RUNSPEC section is set
equal to two.

\begin{verbatim}
--
--       POLYMER SHEARING LOGARITHMIC PARAMETERS
--
PLYSHLOG
--       REF         REF        REF
--       POLCON      SALTCON    TEMP
--       --------    -------    ----
         0.5
/
--
--       VELOCITY    VISFAC
--       --------    -------
         0.0000001    1.00
         0.000001     1.10
         0.0001       1.30
         0.001        1.47
         0.01         1.67
         0.1          2.00
         1.0          2.20
         10.0         2.30
         100.0        2.40
         1000.0       2.40
                                                           / TABLE NO. 01
--       REF         REF        REF
--       POLCON      SALTCON    TEMP
--       --------    -------    ----
         0.5
/


--
--       VELOCITY    VISFAC
--       --------    -------
         0.0000001    1.00
         0.000001     1.10
         0.0001       1.35
         0.001        1.57
         0.01         1.87
         0.1          2.20
         1.0          2.40
         10.0         2.60
         100.0        2.65
         1000.0       2.65
                                                           / TABLE NO. 02
\end{verbatim}

The example activates the polymer logarithmic shear thinning-thickening
option and defines two polymer shear thinning-thickening tables, based
on the NTPVT variable on the TABDIMS keyword in the RUNSPEC section
being equal to two and NPPVT variable on the same keyword being greater
than or equal to ten.

\subsection{PLYTRRF -- Define Polymer Rock Permeability Reduction versus
Temperature}\label{plytrrf-define-polymer-rock-permeability-reduction-versus-temperature}

\begin{longtable}[]{@{}
  >{\raggedright\arraybackslash}p{(\linewidth - 14\tabcolsep) * \real{0.1250}}
  >{\raggedright\arraybackslash}p{(\linewidth - 14\tabcolsep) * \real{0.1250}}
  >{\raggedright\arraybackslash}p{(\linewidth - 14\tabcolsep) * \real{0.1250}}
  >{\raggedright\arraybackslash}p{(\linewidth - 14\tabcolsep) * \real{0.1250}}
  >{\raggedright\arraybackslash}p{(\linewidth - 14\tabcolsep) * \real{0.1250}}
  >{\raggedright\arraybackslash}p{(\linewidth - 14\tabcolsep) * \real{0.1250}}
  >{\raggedright\arraybackslash}p{(\linewidth - 14\tabcolsep) * \real{0.1250}}
  >{\raggedright\arraybackslash}p{(\linewidth - 14\tabcolsep) * \real{0.1250}}@{}}
\toprule\noalign{}
\begin{minipage}[b]{\linewidth}\raggedright
RUNSPEC
\end{minipage} & \begin{minipage}[b]{\linewidth}\raggedright
GRID
\end{minipage} & \begin{minipage}[b]{\linewidth}\raggedright
EDIT
\end{minipage} & \begin{minipage}[b]{\linewidth}\raggedright
PROPS
\end{minipage} & \begin{minipage}[b]{\linewidth}\raggedright
REGIONS
\end{minipage} & \begin{minipage}[b]{\linewidth}\raggedright
SOLUTION
\end{minipage} & \begin{minipage}[b]{\linewidth}\raggedright
SUMMARY
\end{minipage} & \begin{minipage}[b]{\linewidth}\raggedright
SCHEDULE
\end{minipage} \\
\midrule\noalign{}
\endhead
\bottomrule\noalign{}
\endlastfoot
\end{longtable}

\subsubsection{Description}\label{description-578}

The PLYTRRF keyword defines the polymer rock permeability reduction
factor to the water phase as a function of temperature, for when the
Polymer option has been activated by the POLYMER keyword in the RUNSPEC
section. See the PLYTRRF keyword for the options on how this data is
used in the polymer model and the PERMFAC parameter on the PLYROCK
keyword for setting the property for the whole grid for a constant
temperature. Both keywords are in the PROPS section.

This keyword is not supported by OPM Flow but would change the results
if supported so the simulation will be stopped.

\subsection{PLYTRRFA -- Define Polymer Rock Permeability Reduction
versus Temperature
Option}\label{plytrrfa-define-polymer-rock-permeability-reduction-versus-temperature-option}

\begin{longtable}[]{@{}
  >{\raggedright\arraybackslash}p{(\linewidth - 14\tabcolsep) * \real{0.1250}}
  >{\raggedright\arraybackslash}p{(\linewidth - 14\tabcolsep) * \real{0.1250}}
  >{\raggedright\arraybackslash}p{(\linewidth - 14\tabcolsep) * \real{0.1250}}
  >{\raggedright\arraybackslash}p{(\linewidth - 14\tabcolsep) * \real{0.1250}}
  >{\raggedright\arraybackslash}p{(\linewidth - 14\tabcolsep) * \real{0.1250}}
  >{\raggedright\arraybackslash}p{(\linewidth - 14\tabcolsep) * \real{0.1250}}
  >{\raggedright\arraybackslash}p{(\linewidth - 14\tabcolsep) * \real{0.1250}}
  >{\raggedright\arraybackslash}p{(\linewidth - 14\tabcolsep) * \real{0.1250}}@{}}
\toprule\noalign{}
\begin{minipage}[b]{\linewidth}\raggedright
RUNSPEC
\end{minipage} & \begin{minipage}[b]{\linewidth}\raggedright
GRID
\end{minipage} & \begin{minipage}[b]{\linewidth}\raggedright
EDIT
\end{minipage} & \begin{minipage}[b]{\linewidth}\raggedright
PROPS
\end{minipage} & \begin{minipage}[b]{\linewidth}\raggedright
REGIONS
\end{minipage} & \begin{minipage}[b]{\linewidth}\raggedright
SOLUTION
\end{minipage} & \begin{minipage}[b]{\linewidth}\raggedright
SUMMARY
\end{minipage} & \begin{minipage}[b]{\linewidth}\raggedright
SCHEDULE
\end{minipage} \\
\midrule\noalign{}
\endhead
\bottomrule\noalign{}
\endlastfoot
\end{longtable}

\subsubsection{Description}\label{description-579}

The PLYTRRFA keyword defines the how the polymer rock permeability
reduction factor to the water phase as a function of temperature data,
entered via the PLYTRRA keyword in the PROPS section, should be used.
This keyword should only be used if the Polymer option has been
activated by the POLYMER keyword in the RUNSPEC section. See the PERMFAC
parameter on the PLYROCK keyword in the PROPS section for setting the
property for the whole grid for a constant temperature.

This keyword is not supported by OPM Flow but would change the results
if supported so the simulation will be stopped.

\subsection{PLYVISC -- Define Polymer Viscosity Scaling
Factors}\label{plyvisc-define-polymer-viscosity-scaling-factors}

\begin{longtable}[]{@{}
  >{\raggedright\arraybackslash}p{(\linewidth - 14\tabcolsep) * \real{0.1250}}
  >{\raggedright\arraybackslash}p{(\linewidth - 14\tabcolsep) * \real{0.1250}}
  >{\raggedright\arraybackslash}p{(\linewidth - 14\tabcolsep) * \real{0.1250}}
  >{\raggedright\arraybackslash}p{(\linewidth - 14\tabcolsep) * \real{0.1250}}
  >{\raggedright\arraybackslash}p{(\linewidth - 14\tabcolsep) * \real{0.1250}}
  >{\raggedright\arraybackslash}p{(\linewidth - 14\tabcolsep) * \real{0.1250}}
  >{\raggedright\arraybackslash}p{(\linewidth - 14\tabcolsep) * \real{0.1250}}
  >{\raggedright\arraybackslash}p{(\linewidth - 14\tabcolsep) * \real{0.1250}}@{}}
\toprule\noalign{}
\begin{minipage}[b]{\linewidth}\raggedright
RUNSPEC
\end{minipage} & \begin{minipage}[b]{\linewidth}\raggedright
GRID
\end{minipage} & \begin{minipage}[b]{\linewidth}\raggedright
EDIT
\end{minipage} & \begin{minipage}[b]{\linewidth}\raggedright
PROPS
\end{minipage} & \begin{minipage}[b]{\linewidth}\raggedright
REGIONS
\end{minipage} & \begin{minipage}[b]{\linewidth}\raggedright
SOLUTION
\end{minipage} & \begin{minipage}[b]{\linewidth}\raggedright
SUMMARY
\end{minipage} & \begin{minipage}[b]{\linewidth}\raggedright
SCHEDULE
\end{minipage} \\
\midrule\noalign{}
\endhead
\bottomrule\noalign{}
\endlastfoot
\end{longtable}

\subsubsection{Description}\label{description-580}

PLYSVISC defines the polymer viscosity scaling factors used to determine
the relationship of pure water viscosity with respect to increasing
polymer concentration within a grid block. The polymer option must be
activated by the POLYMER keyword in the RUNSPEC section in order to use
this keyword.

The BRINE option in the RUNSPEC should be deactivated if this keyword is
to be used.

\begin{longtable}[]{@{}
  >{\raggedright\arraybackslash}p{(\linewidth - 6\tabcolsep) * \real{0.2500}}
  >{\raggedright\arraybackslash}p{(\linewidth - 6\tabcolsep) * \real{0.2500}}
  >{\raggedright\arraybackslash}p{(\linewidth - 6\tabcolsep) * \real{0.2500}}
  >{\raggedright\arraybackslash}p{(\linewidth - 6\tabcolsep) * \real{0.2500}}@{}}
\toprule\noalign{}
\begin{minipage}[b]{\linewidth}\raggedright
No.
\end{minipage} & \begin{minipage}[b]{\linewidth}\raggedright
Name
\end{minipage} & \begin{minipage}[b]{\linewidth}\raggedright
Description
\end{minipage} & \begin{minipage}[b]{\linewidth}\raggedright
Default
\end{minipage} \\
\midrule\noalign{}
\endhead
\bottomrule\noalign{}
\endlastfoot
Field & Metric & Laboratory & \\
1 & POLCON & A columnar vector of real monotonically increasing down the
column values that defines the polymer concentration in the solution
surrounding the rock. The first entry should be zero to define a no
polymer concentration. & None \\
lb/stb & kg/sm3 & gm/scc & \\
2 & VISFAC & A columnar vector of real increasing or equal values that
defines a factor that scales the effective viscosity of the solution for
the given POLCON entry. Normally VISFAC value for the first row in the
table should be one. & None \\
dimensionless & dimensionless & dimensionless & \\
Notes: & & & \\
\end{longtable}

\emph{Table 8.108: PLYVISC Keyword Description}

\subsubsection{Example}\label{example-320}

\begin{verbatim}
--
--       POLYMER VISCOSITY SCALING FACTOR TABLES
--
PLYVISC
--       POLYMER     VISCOSITY
--       POLCON      VISFAC
--       --------    ---------
          0.0000       1.000
          0.0002      10.000
          0.0004      20.000
          0.0008      40.000                               / TABLE NO. 01

--       POLYMER     VISCOSITY
--       POLCON      VISFAC
--       --------    ---------
          0.0000       1.000
          0.0003      10.000
          0.0005      20.000
          0.0007      40.000
          0.0009      45.000
          0.0011      55.000                               / TABLE NO. 02
\end{verbatim}

The example defines two polymer viscosity scaling factor tables, based
on the NTPVT variable on the TABDIMS keyword in the RUNSPEC section
being equal to two and NPPVT variable on the same keyword being greater
than or equal to six.

\subsection{PLYVISCS -- Define Polymer-Salt Viscosity Scaling
Factors}\label{plyviscs-define-polymer-salt-viscosity-scaling-factors}

\begin{longtable}[]{@{}
  >{\raggedright\arraybackslash}p{(\linewidth - 14\tabcolsep) * \real{0.1250}}
  >{\raggedright\arraybackslash}p{(\linewidth - 14\tabcolsep) * \real{0.1250}}
  >{\raggedright\arraybackslash}p{(\linewidth - 14\tabcolsep) * \real{0.1250}}
  >{\raggedright\arraybackslash}p{(\linewidth - 14\tabcolsep) * \real{0.1250}}
  >{\raggedright\arraybackslash}p{(\linewidth - 14\tabcolsep) * \real{0.1250}}
  >{\raggedright\arraybackslash}p{(\linewidth - 14\tabcolsep) * \real{0.1250}}
  >{\raggedright\arraybackslash}p{(\linewidth - 14\tabcolsep) * \real{0.1250}}
  >{\raggedright\arraybackslash}p{(\linewidth - 14\tabcolsep) * \real{0.1250}}@{}}
\toprule\noalign{}
\begin{minipage}[b]{\linewidth}\raggedright
RUNSPEC
\end{minipage} & \begin{minipage}[b]{\linewidth}\raggedright
GRID
\end{minipage} & \begin{minipage}[b]{\linewidth}\raggedright
EDIT
\end{minipage} & \begin{minipage}[b]{\linewidth}\raggedright
PROPS
\end{minipage} & \begin{minipage}[b]{\linewidth}\raggedright
REGIONS
\end{minipage} & \begin{minipage}[b]{\linewidth}\raggedright
SOLUTION
\end{minipage} & \begin{minipage}[b]{\linewidth}\raggedright
SUMMARY
\end{minipage} & \begin{minipage}[b]{\linewidth}\raggedright
SCHEDULE
\end{minipage} \\
\midrule\noalign{}
\endhead
\bottomrule\noalign{}
\endlastfoot
\end{longtable}

\subsubsection{Description}\label{description-581}

PLYSVISCS defines the polymer-salt viscosity scaling factor tables
applied to pure water that are used to determine the viscosity of a
polymer-salt mixture with respect to increasing polymer concentration
within a grid block. The polymer option must be activated by the POLYMER
keyword, as well as the brine phase declared by the BRINE keyword in the
RUNSPEC section in order to use this keyword. However the ECLM keyword
in the RUNSPEC must not be used with this keyword.

See also the PLYVSCST keyword in the PROPS section to enter polymer
viscosity scaling factor data that is dependent on both salt and
reservoir temperature.

This keyword is not supported by OPM Flow but would change the results
if supported so the simulation will be stopped.

\subsection{PLYVISCT -- Define Polymer-Temperature Viscosity Scaling
Factors}\label{plyvisct-define-polymer-temperature-viscosity-scaling-factors}

\begin{longtable}[]{@{}
  >{\raggedright\arraybackslash}p{(\linewidth - 14\tabcolsep) * \real{0.1250}}
  >{\raggedright\arraybackslash}p{(\linewidth - 14\tabcolsep) * \real{0.1250}}
  >{\raggedright\arraybackslash}p{(\linewidth - 14\tabcolsep) * \real{0.1250}}
  >{\raggedright\arraybackslash}p{(\linewidth - 14\tabcolsep) * \real{0.1250}}
  >{\raggedright\arraybackslash}p{(\linewidth - 14\tabcolsep) * \real{0.1250}}
  >{\raggedright\arraybackslash}p{(\linewidth - 14\tabcolsep) * \real{0.1250}}
  >{\raggedright\arraybackslash}p{(\linewidth - 14\tabcolsep) * \real{0.1250}}
  >{\raggedright\arraybackslash}p{(\linewidth - 14\tabcolsep) * \real{0.1250}}@{}}
\toprule\noalign{}
\begin{minipage}[b]{\linewidth}\raggedright
RUNSPEC
\end{minipage} & \begin{minipage}[b]{\linewidth}\raggedright
GRID
\end{minipage} & \begin{minipage}[b]{\linewidth}\raggedright
EDIT
\end{minipage} & \begin{minipage}[b]{\linewidth}\raggedright
PROPS
\end{minipage} & \begin{minipage}[b]{\linewidth}\raggedright
REGIONS
\end{minipage} & \begin{minipage}[b]{\linewidth}\raggedright
SOLUTION
\end{minipage} & \begin{minipage}[b]{\linewidth}\raggedright
SUMMARY
\end{minipage} & \begin{minipage}[b]{\linewidth}\raggedright
SCHEDULE
\end{minipage} \\
\midrule\noalign{}
\endhead
\bottomrule\noalign{}
\endlastfoot
\end{longtable}

\subsubsection{Description}\label{description-582}

PLYSVISCT defines the polymer-temperature viscosity scaling factor
tables applied to pure water that are used to determine the viscosity of
the polymer at a given temperature with respect to increasing polymer
concentration within a grid block. Both the polymer option must be
activated by the POLYMER keyword and the temperature option invoked by
the TEMP keyword in the RUNSPEC section in order to use this keyword.
However the BRINE keyword in the RUNSPEC must not be used with this
keyword, that is the salt sensitivity options should be deactivated.

See also the PLYVSCST keyword in the PROPS section to enter polymer
viscosity scaling factor data that is dependent on both salt and
reservoir temperature.

This keyword is not supported by OPM Flow but would change the results
if supported so the simulation will be stopped.

\subsection{PLYVMH -- Polymer Molecular Weight Model Polymer Viscosity
Constants}\label{plyvmh-polymer-molecular-weight-model-polymer-viscosity-constants}

\begin{longtable}[]{@{}
  >{\raggedright\arraybackslash}p{(\linewidth - 14\tabcolsep) * \real{0.1250}}
  >{\raggedright\arraybackslash}p{(\linewidth - 14\tabcolsep) * \real{0.1250}}
  >{\raggedright\arraybackslash}p{(\linewidth - 14\tabcolsep) * \real{0.1250}}
  >{\raggedright\arraybackslash}p{(\linewidth - 14\tabcolsep) * \real{0.1250}}
  >{\raggedright\arraybackslash}p{(\linewidth - 14\tabcolsep) * \real{0.1250}}
  >{\raggedright\arraybackslash}p{(\linewidth - 14\tabcolsep) * \real{0.1250}}
  >{\raggedright\arraybackslash}p{(\linewidth - 14\tabcolsep) * \real{0.1250}}
  >{\raggedright\arraybackslash}p{(\linewidth - 14\tabcolsep) * \real{0.1250}}@{}}
\toprule\noalign{}
\begin{minipage}[b]{\linewidth}\raggedright
RUNSPEC
\end{minipage} & \begin{minipage}[b]{\linewidth}\raggedright
GRID
\end{minipage} & \begin{minipage}[b]{\linewidth}\raggedright
EDIT
\end{minipage} & \begin{minipage}[b]{\linewidth}\raggedright
PROPS
\end{minipage} & \begin{minipage}[b]{\linewidth}\raggedright
REGIONS
\end{minipage} & \begin{minipage}[b]{\linewidth}\raggedright
SOLUTION
\end{minipage} & \begin{minipage}[b]{\linewidth}\raggedright
SUMMARY
\end{minipage} & \begin{minipage}[b]{\linewidth}\raggedright
SCHEDULE
\end{minipage} \\
\midrule\noalign{}
\endhead
\bottomrule\noalign{}
\endlastfoot
\end{longtable}

\subsubsection{Description}\label{description-583}

This keyword, PLYVMH, defines the constants used to calculate viscosity
of the polymer solution as a function of the polymer molecular weight
and the polymer concentration, for the simulator's Polymer Molecular
Weight Transport option, that uses the polymer molecular weight in
calculating the polymer viscosity. The keyword consists of a series of
row vectors, which each vector having four elements, that define the
constants used in calculating the polymer viscosity

This keyword should only be used if the POLYMER and POLYMW keywords in
the RUNSPEC section are also activated.

\begin{tcolorbox}[enhanced jigsaw, title=\textcolor{quarto-callout-note-color}{\faInfo}\hspace{0.5em}{Note}, coltitle=black, breakable, toprule=.15mm, colbacktitle=quarto-callout-note-color!10!white, leftrule=.75mm, opacitybacktitle=0.6, colback=white, bottomtitle=1mm, toptitle=1mm, titlerule=0mm, left=2mm, colframe=quarto-callout-note-color-frame, arc=.35mm, rightrule=.15mm, bottomrule=.15mm, opacityback=0]

This is an OPM Flow specific keyword that employs an alternative polymer
flood model based on a Polymer Molecular Weight Transport equation, that
is not available in the commercial simulator. The model has been tested
using metric units; however, using either field or laboratory units with
the option should be considered experimental.

\end{tcolorbox}

\begin{longtable}[]{@{}
  >{\raggedright\arraybackslash}p{(\linewidth - 6\tabcolsep) * \real{0.2500}}
  >{\raggedright\arraybackslash}p{(\linewidth - 6\tabcolsep) * \real{0.2500}}
  >{\raggedright\arraybackslash}p{(\linewidth - 6\tabcolsep) * \real{0.2500}}
  >{\raggedright\arraybackslash}p{(\linewidth - 6\tabcolsep) * \real{0.2500}}@{}}
\toprule\noalign{}
\begin{minipage}[b]{\linewidth}\raggedright
No.
\end{minipage} & \begin{minipage}[b]{\linewidth}\raggedright
Name
\end{minipage} & \begin{minipage}[b]{\linewidth}\raggedright
Description
\end{minipage} & \begin{minipage}[b]{\linewidth}\raggedright
Default
\end{minipage} \\
\midrule\noalign{}
\endhead
\bottomrule\noalign{}
\endlastfoot
Field & Metric & Laboratory & \\
1 ‍ & MHK & The Mark-Houwink K polymer specific constant in the
Mark-Houwink equation, see equation (8.77). & None \\
& ml/g & & \\
2 & MHA & The Mark-Houwink exponent parameter a, in the Mark-Houwink
equation, see equation (8.77). & None \\
dimensionless & dimensionless & dimensionless & \\
3 & GAMMA & The ɣ intrinsic water-polymer viscosity constant in the
Huggins equation, see equation (8.79). & None \\
4 & KAPPA & The κ intrinsic water-polymer viscosity constant in the
Huggins equation, see equation (8.79). & None \\
Notes: & & & \\
\end{longtable}

\emph{Table 8.109: PLYVMH Keyword Description}

The high molecular weight of polymers greatly increase the viscosity of
the injected water in which they are dissolved, thus adjusting the
mobility ration of the displacing phase. The increase in viscosity is
caused by strong internal friction between the randomly coiled and
swollen macro molecules and the surrounding water molecules. And is
dependent on both the nature of the polymer and the injected water
acting as the solvent. There are several formulations of viscosity
associated with polymer rheology, namely:

Relative viscosity is defined as:

\begin{equation}\phantomsection\label{eq-8-73}{
{\mathrm{η}}_{\mathit{rel}} = \frac{\mathrm{η}}{{\mathrm{η}}_{s}}
}\end{equation}

where:

\({η}_{\mathit{rel}}\) = is the relative viscosity,

\(η\) = is the viscosity of the solution, and

\({η}_{s}\) = is the viscosity of the solvent (injected water).

Specific viscosity is defined as:

\begin{equation}\phantomsection\label{eq-8-74}{
\begin{matrix}{\mathrm{η}}_{\text{sp}} = \frac{\left(\mathrm{η}-{\mathrm{η}}_{s}\right)}{{\mathrm{η}}_{s}} \\   = \frac{\mathrm{η}}{{\mathrm{η}}_{s}} - 1 \\   = {\mathrm{η}}_{\mathit{rel}} - 1\end{matrix}
}\end{equation}

where:

\({η}_{\text{sp}}\) = is the specific water-polymer viscosity, and

\({η}_{s}\) = is the viscosity of the solvent (injected water).

Reduced Specific Viscosity is given by:

\begin{equation}\phantomsection\label{eq-8-75}{
\begin{matrix}{\mathrm{η}}_{\text{red}} = \frac{{\mathrm{η}}_{\text{sp}}}{{C}_{p}} \\   = \frac{\left({\mathrm{η}}_{\mathit{rel}} - 1\right)}{{C}_{p}}\end{matrix}
}\end{equation}

where:

\({η}_{\text{red}}\) = is the reduced specific water-polymer viscosity,
and

\({C}_{p}\) = is the polymer concentration.

Finally, the Intrinsic Velocity, which is defined as a measure for the
internal friction in polymer solutions at the limit of zero polymer
concentration. Thus, this quantity describes the effect of completely
separated polymer chains on the solution viscosity, and is defined as:

\begin{equation}\phantomsection\label{eq-8-76}{
\begin{matrix}\left[\mathrm{η}\right] = \underset{{C}_{p}\rightarrow 0}{lim}\frac{{\mathrm{η}}_{\text{sp}}}{{C}_{p}} \\   = \underset{{C}_{p}\rightarrow 0}{lim}\frac{\mathrm{η}-{\mathrm{η}}_{0}}{{\mathrm{η}}_{0}{C}_{p}}\end{matrix}
}\end{equation}

where:

\(\left[η\right]\) = is the intrinsic water-polymer viscosity and
describes the increase in

viscosity arising from an individual polymer chain and is a measure of
the polymers' thickening power, and

\({η}_{o}\) = zero-shear viscosity.

For any given solvent pair, the intrinsic viscosity increases as the
molecular weight of the polymer increases; here, the Mark-Houwink
equation, also known as the Mark-Houwink-Staudinger equation {[}H. Mark,
in R. Saenger, Der feste Koerper, Hirzel, Leipzig, 1938.{]}, {[}R.
Houwink , J. Prakt. Chem., Vol. 157, Issue 1-3, p.~15 (1940).{]}, and
{[}H. Staudinger, Die Hochmolekulare Organischen Verbindungen, Julius
Springer, Berlin 1932.{]} is used to calculate {[}η{]}, that is:

\begin{equation}\phantomsection\label{eq-8-77}{
\left[\mathrm{η}\right] = K ⋅ {{M}_{w}}^{a}
}\end{equation}

where:

\(K\) = is the polymer specific constant in the Mark-Houwink equation,
the MHK

parameter in Table 8.109,

\({M}_{w}\) = is the polymer molecular weight, and

\(a\) = the exponent constant in the Mark-Houwink equation, the MHA
parameter

in Table 8.109.

The Mark-Houwink parameters can be determined from a double logarithmic
plot of intrinsic viscosity versus molecular weight which yields
straight lines, that is:

\begin{equation}\phantomsection\label{eq-8-78}{
ln\left([\mathrm{η}]\right) = ln(K) + a \times  ln({M}_{w})
}\end{equation}

OPM Flow uses a form of the Huggins {[}Huggins, M. L. 1942. The
viscosity of dilute solutions of long-chain molecules. IV. Dependence on
concentration.» Journal of the American Chemical Society 64 (11):
2716-2718.{]} equation to calculate the polymer apparent viscosity, as
shown below:

\begin{equation}\phantomsection\label{eq-8-79}{
\begin{matrix} \frac{{\mathrm{η}}_{0}}{{\mathrm{η}}_{s}} = 1 + \mathrm{γ}\left(X + \mathrm{κ}{X}^{2}\right) \\ \text{where} X = \left[\mathrm{η}\right]{C}_{p}\end{matrix}
}\end{equation}

where:

\(γ\) = a user defined constant, GAMMA in Table 8.109, and

\(κ\) = a user defined constant, KAPPA in Table 8.109.

Which can be used to calculate the zero-shear viscosity, η0, based on
the quadratic function in equation (8.79) time polymer concentration(Cp)
and polymer intrinsic viscosity described in equation (8.77).

In terms of the keyword units, given the intrinsic viscosity in ml/g,
polymer concentration in kg/m3, and the molecular weight as kg/kg-M,
then equation (8.77) becomes:

\begin{equation}\phantomsection\label{eq-8-80}{
\left[\mathrm{η}\right] = K{\left({M}_{w} ⋅ 1.0\times {10}^{-3}\right)}^{a}\cdot 1.0\times {10}^{-3}
}\end{equation}

and equation (8.79) becomes:

\begin{equation}\phantomsection\label{eq-8-81}{
\begin{matrix} \frac{{\mathrm{η}}_{0}}{{\mathrm{η}}_{s}} = 1 + \mathrm{γ}\left(X + \mathrm{κ}{X}^{2}\right) \\ \text{where} X = 1.0\times {10}^{-6}\left[\mathrm{η}\right]{C}_{p}\end{matrix}
}\end{equation}

Note that the model does not account for non-Newtonian flow; the
apparent viscosity is simply set equal to the zero-shear viscosity, and
that the model only considers the full mixing between the polymer and
water.

See also the PLYMWINJ keyword in the PROPS section, that describes the
relationship of the injected polymer molecular weight as a function of
polymer throughput and polymer velocity. Note that the standard polymer
property data keywords: PLYROCK, PLYADS, PLYMAX, etc., are still
required to fully describe the polymer fluid.

\subsubsection{Example}\label{example-321}

Given NPLYVMH equals two on the PINTDIMS keyword in the RUNSPEC section,
then:

\begin{verbatim}
--
--       POLYMER MOLECULAR WEIGHT MODEL POLYMER VISCOSITY CONSTANTS
--       (OPM FLOW PROPS KEYWORD)
--
--       MHK     MHA     VISC    VISC
--       CONST   EXPON   GAMMA   KAPPA
PLYVMH
         0.02    0.50    0.40    0.60                      /
         0.03    0.51    0.42    0.63                      /
/
\end{verbatim}

Two sets of data should be entered with the PLYVMH keyword, as shown
above.

\subsection{PLYVSCST -- Define Polymer-Salt-Temperature Viscosity
Scaling
Factors}\label{plyvscst-define-polymer-salt-temperature-viscosity-scaling-factors}

\begin{longtable}[]{@{}
  >{\raggedright\arraybackslash}p{(\linewidth - 14\tabcolsep) * \real{0.1250}}
  >{\raggedright\arraybackslash}p{(\linewidth - 14\tabcolsep) * \real{0.1250}}
  >{\raggedright\arraybackslash}p{(\linewidth - 14\tabcolsep) * \real{0.1250}}
  >{\raggedright\arraybackslash}p{(\linewidth - 14\tabcolsep) * \real{0.1250}}
  >{\raggedright\arraybackslash}p{(\linewidth - 14\tabcolsep) * \real{0.1250}}
  >{\raggedright\arraybackslash}p{(\linewidth - 14\tabcolsep) * \real{0.1250}}
  >{\raggedright\arraybackslash}p{(\linewidth - 14\tabcolsep) * \real{0.1250}}
  >{\raggedright\arraybackslash}p{(\linewidth - 14\tabcolsep) * \real{0.1250}}@{}}
\toprule\noalign{}
\begin{minipage}[b]{\linewidth}\raggedright
RUNSPEC
\end{minipage} & \begin{minipage}[b]{\linewidth}\raggedright
GRID
\end{minipage} & \begin{minipage}[b]{\linewidth}\raggedright
EDIT
\end{minipage} & \begin{minipage}[b]{\linewidth}\raggedright
PROPS
\end{minipage} & \begin{minipage}[b]{\linewidth}\raggedright
REGIONS
\end{minipage} & \begin{minipage}[b]{\linewidth}\raggedright
SOLUTION
\end{minipage} & \begin{minipage}[b]{\linewidth}\raggedright
SUMMARY
\end{minipage} & \begin{minipage}[b]{\linewidth}\raggedright
SCHEDULE
\end{minipage} \\
\midrule\noalign{}
\endhead
\bottomrule\noalign{}
\endlastfoot
\end{longtable}

\subsubsection{Description}\label{description-584}

PLYVSCST defines the polymer-salt-temperature viscosity scaling factor
tables applied to pure water that are used to determine the viscosity of
the polymer at a given salt concentration and for a given temperature,
with respect to increasing polymer concentration within a grid block.
Both the polymer option must be activated by the POLYMER keyword and the
temperature option invoked by the TEMP keyword in the RUNSPEC section in
order to use this keyword. In addition, the BRINE keyword in the RUNSPEC
must also be invoked. The keyword is used in conjunction with the
SALTNODE keyword to define the various salt concentrations and the
TEMPNODE keyword to define the various reservoir temperatures. Both
keywords are in the PROPS section.

See also the PLYVISCS keyword in the PROPS section to enter polymer
viscosity scaling factor data that is dependent just salt concentration
and the PLYVISCT keyword in the PROPS section to enter polymer viscosity
scaling factor data that is dependent just on reservoir temperature.

This keyword is not supported by OPM Flow but would change the results
if supported so the simulation will be stopped.

\subsection{PMAX -- Maximum and Minimum Pressure for Total
Compressibility
Check}\label{pmax-maximum-and-minimum-pressure-for-total-compressibility-check}

\begin{longtable}[]{@{}
  >{\raggedright\arraybackslash}p{(\linewidth - 14\tabcolsep) * \real{0.1250}}
  >{\raggedright\arraybackslash}p{(\linewidth - 14\tabcolsep) * \real{0.1250}}
  >{\raggedright\arraybackslash}p{(\linewidth - 14\tabcolsep) * \real{0.1250}}
  >{\raggedright\arraybackslash}p{(\linewidth - 14\tabcolsep) * \real{0.1250}}
  >{\raggedright\arraybackslash}p{(\linewidth - 14\tabcolsep) * \real{0.1250}}
  >{\raggedright\arraybackslash}p{(\linewidth - 14\tabcolsep) * \real{0.1250}}
  >{\raggedright\arraybackslash}p{(\linewidth - 14\tabcolsep) * \real{0.1250}}
  >{\raggedright\arraybackslash}p{(\linewidth - 14\tabcolsep) * \real{0.1250}}@{}}
\toprule\noalign{}
\begin{minipage}[b]{\linewidth}\raggedright
RUNSPEC
\end{minipage} & \begin{minipage}[b]{\linewidth}\raggedright
GRID
\end{minipage} & \begin{minipage}[b]{\linewidth}\raggedright
EDIT
\end{minipage} & \begin{minipage}[b]{\linewidth}\raggedright
PROPS
\end{minipage} & \begin{minipage}[b]{\linewidth}\raggedright
REGIONS
\end{minipage} & \begin{minipage}[b]{\linewidth}\raggedright
SOLUTION
\end{minipage} & \begin{minipage}[b]{\linewidth}\raggedright
SUMMARY
\end{minipage} & \begin{minipage}[b]{\linewidth}\raggedright
SCHEDULE
\end{minipage} \\
\midrule\noalign{}
\endhead
\bottomrule\noalign{}
\endlastfoot
\end{longtable}

\subsubsection{Description}\label{description-585}

The PMAX keyword defines the maximum and minimum pressures expected to
be encountered during the run. The data is used to perform the PVT total
compressibility check that ensures that the total compressibility of a
mixture of oil-gas, for when the gas-oil ratio is increasing for an oil,
or the condensate gas ratio is increasing for a gas condensate, is
positive respect to pressure. The total compressibility check is used to
ensure that the entered oil and gas PVT data is consistent. If the check
fails for given oil-gas mixture at a given pressure, resulting in a
negative total compressibility, then this will result in numerical
instabilities in the run causing this simulator difficulties in
converging to a solution.

This keyword is not supported by OPM Flow but would change the results
if supported so the simulation will be stopped.

\subsection{PMISC -- Define Miscibility versus Pressure
Tables}\label{pmisc-define-miscibility-versus-pressure-tables}

\begin{longtable}[]{@{}
  >{\raggedright\arraybackslash}p{(\linewidth - 14\tabcolsep) * \real{0.1250}}
  >{\raggedright\arraybackslash}p{(\linewidth - 14\tabcolsep) * \real{0.1250}}
  >{\raggedright\arraybackslash}p{(\linewidth - 14\tabcolsep) * \real{0.1250}}
  >{\raggedright\arraybackslash}p{(\linewidth - 14\tabcolsep) * \real{0.1250}}
  >{\raggedright\arraybackslash}p{(\linewidth - 14\tabcolsep) * \real{0.1250}}
  >{\raggedright\arraybackslash}p{(\linewidth - 14\tabcolsep) * \real{0.1250}}
  >{\raggedright\arraybackslash}p{(\linewidth - 14\tabcolsep) * \real{0.1250}}
  >{\raggedright\arraybackslash}p{(\linewidth - 14\tabcolsep) * \real{0.1250}}@{}}
\toprule\noalign{}
\begin{minipage}[b]{\linewidth}\raggedright
RUNSPEC
\end{minipage} & \begin{minipage}[b]{\linewidth}\raggedright
GRID
\end{minipage} & \begin{minipage}[b]{\linewidth}\raggedright
EDIT
\end{minipage} & \begin{minipage}[b]{\linewidth}\raggedright
PROPS
\end{minipage} & \begin{minipage}[b]{\linewidth}\raggedright
REGIONS
\end{minipage} & \begin{minipage}[b]{\linewidth}\raggedright
SOLUTION
\end{minipage} & \begin{minipage}[b]{\linewidth}\raggedright
SUMMARY
\end{minipage} & \begin{minipage}[b]{\linewidth}\raggedright
SCHEDULE
\end{minipage} \\
\midrule\noalign{}
\endhead
\bottomrule\noalign{}
\endlastfoot
\end{longtable}

\subsubsection{Description}\label{description-586}

PMISC defines the transition between immiscible and miscible
displacement as a function of oil pressure tables, for when the MISCIBLE
keyword in the RUNSPEC section has be activated. If this keyword is
absent from the input deck and MISCIBLE keyword in the RUNSPEC keyword
has been activated, then miscibility is independent of the oil phase
pressure.

\begin{longtable}[]{@{}
  >{\raggedright\arraybackslash}p{(\linewidth - 6\tabcolsep) * \real{0.2500}}
  >{\raggedright\arraybackslash}p{(\linewidth - 6\tabcolsep) * \real{0.2500}}
  >{\raggedright\arraybackslash}p{(\linewidth - 6\tabcolsep) * \real{0.2500}}
  >{\raggedright\arraybackslash}p{(\linewidth - 6\tabcolsep) * \real{0.2500}}@{}}
\toprule\noalign{}
\begin{minipage}[b]{\linewidth}\raggedright
No.
\end{minipage} & \begin{minipage}[b]{\linewidth}\raggedright
Name
\end{minipage} & \begin{minipage}[b]{\linewidth}\raggedright
Description
\end{minipage} & \begin{minipage}[b]{\linewidth}\raggedright
Default
\end{minipage} \\
\midrule\noalign{}
\endhead
\bottomrule\noalign{}
\endlastfoot
Field & Metric & Laboratory & \\
1 & PRESS & A columnar vector of real monotonically increasing down the
column values that defines the oil phase pressure. & None \\
psia & barsa & atma & \\
2 & MISC & A columnar vector of real equal or increasing down the column
values that defines the corresponding miscibility factor. MISC is a
scaling that should lie be zero and one, where zero means no miscibility
and one means full miscibility. & None \\
dimensionless & dimensionless & dimensionless & \\
Notes: & & & \\
\end{longtable}

\emph{Table 8.110: PMISC Keyword Description}

\subsubsection{Example}\label{example-322}

\begin{verbatim}
--
--       MISCIBILITY VERSUS PRESSURE TABLES
--
PMISC
--       OIL        MISCIBILE
--       PRESS      FACTOR
--       -------    ---------
          1000.0       0.000
          2000.0       0.250
          3000.0       1.000
          4000.0       1.000                               / TABLE NO. 01
--       OIL        MISCIBILE
--       PRESS      FACTOR
--       -------    ---------
          1500.0       0.000
          2000.0       0.000
          2500.0       0.250
          3000.0       0.350
          3500.0       1.000
          4000.0       1.000                               / TABLE NO. 02
\end{verbatim}

The above example defines two miscibility versus pressure tables
assuming NTMISC equals two and NSMISC is greater than or equal to six on
the MISCIBLE keyword in the RUNSPEC section.

\subsection{PPCWMAX -- Define SWATINIT Calculated Capillary Pressure
Constraints}\label{ppcwmax-define-swatinit-calculated-capillary-pressure-constraints}

\begin{longtable}[]{@{}
  >{\raggedright\arraybackslash}p{(\linewidth - 14\tabcolsep) * \real{0.1250}}
  >{\raggedright\arraybackslash}p{(\linewidth - 14\tabcolsep) * \real{0.1250}}
  >{\raggedright\arraybackslash}p{(\linewidth - 14\tabcolsep) * \real{0.1250}}
  >{\raggedright\arraybackslash}p{(\linewidth - 14\tabcolsep) * \real{0.1250}}
  >{\raggedright\arraybackslash}p{(\linewidth - 14\tabcolsep) * \real{0.1250}}
  >{\raggedright\arraybackslash}p{(\linewidth - 14\tabcolsep) * \real{0.1250}}
  >{\raggedright\arraybackslash}p{(\linewidth - 14\tabcolsep) * \real{0.1250}}
  >{\raggedright\arraybackslash}p{(\linewidth - 14\tabcolsep) * \real{0.1250}}@{}}
\toprule\noalign{}
\begin{minipage}[b]{\linewidth}\raggedright
RUNSPEC
\end{minipage} & \begin{minipage}[b]{\linewidth}\raggedright
GRID
\end{minipage} & \begin{minipage}[b]{\linewidth}\raggedright
EDIT
\end{minipage} & \begin{minipage}[b]{\linewidth}\raggedright
PROPS
\end{minipage} & \begin{minipage}[b]{\linewidth}\raggedright
REGIONS
\end{minipage} & \begin{minipage}[b]{\linewidth}\raggedright
SOLUTION
\end{minipage} & \begin{minipage}[b]{\linewidth}\raggedright
SUMMARY
\end{minipage} & \begin{minipage}[b]{\linewidth}\raggedright
SCHEDULE
\end{minipage} \\
\midrule\noalign{}
\endhead
\bottomrule\noalign{}
\endlastfoot
\end{longtable}

\subsubsection{Description}\label{description-587}

The PPCWMAX keyword defines the maximum capillary pressure allowed when
scaling the capillary pressure tables to match the inputted SWATINIT
array. This is primary used for when the SWATINIT array has values of
water saturation above the connate water saturation significantly
outside than capillary pressure transition zone, that is high on the
structure. In this case OPM Flow may generate large values for the
capillary pressure which may result in numerical converge problems. This
keyword sets the maximum allowable calculated capillary pressure and how
the water saturation should be treated when the limit is exceeded.

\begin{longtable}[]{@{}
  >{\raggedright\arraybackslash}p{(\linewidth - 6\tabcolsep) * \real{0.2500}}
  >{\raggedright\arraybackslash}p{(\linewidth - 6\tabcolsep) * \real{0.2500}}
  >{\raggedright\arraybackslash}p{(\linewidth - 6\tabcolsep) * \real{0.2500}}
  >{\raggedright\arraybackslash}p{(\linewidth - 6\tabcolsep) * \real{0.2500}}@{}}
\toprule\noalign{}
\begin{minipage}[b]{\linewidth}\raggedright
No.
\end{minipage} & \begin{minipage}[b]{\linewidth}\raggedright
Name
\end{minipage} & \begin{minipage}[b]{\linewidth}\raggedright
Description
\end{minipage} & \begin{minipage}[b]{\linewidth}\raggedright
Default
\end{minipage} \\
\midrule\noalign{}
\endhead
\bottomrule\noalign{}
\endlastfoot
Field & Metric & Laboratory & \\
1 & PCWO & A columnar vector of real values that defines the maximum
allowable capillary pressure for each SATNUM region. The default value
of infinity means there is no limit applied. & Infinity \\
psia & barsa & atma & \\
2 & OPTN & A columnar vector of character strings that should be set to:
& No \\
Notes: & & & \\
\end{longtable}

\emph{Table 8.111: PPCWMAX Keyword Description}

\begin{tcolorbox}[enhanced jigsaw, title=\textcolor{quarto-callout-note-color}{\faInfo}\hspace{0.5em}{Note}, coltitle=black, breakable, toprule=.15mm, colbacktitle=quarto-callout-note-color!10!white, leftrule=.75mm, opacitybacktitle=0.6, colback=white, bottomtitle=1mm, toptitle=1mm, titlerule=0mm, left=2mm, colframe=quarto-callout-note-color-frame, arc=.35mm, rightrule=.15mm, bottomrule=.15mm, opacityback=0]

Using this keyword to limit the re-scaled grid block capillary pressure
values will effect the fluids in-place when the simulator has to
re-calculate values due to the capillary pressure limit being exceeded.
In addition, the high grid block capillary pressures may be indicative
of an inconsistency between the tabular SATNUM capillary pressure values
and the provided SWATINIT array water saturations. This inconsistency
may be a result of the SWATINIT array being derived using a saturation
height function, as is customary in static modeling software, and the
numerical models tabulated capillary pressure. Rather than resetting the
maximum calculated capillary pressure using the PPCWMAX keyword, it may
be more appropriate to investigate the reason for the high capillary
pressures values first, prior to applying the keyword.

\end{tcolorbox}

\subsubsection{Example}\label{example-323}

\begin{verbatim}
--
--       SET MAXIMUM PC FOR SWATINIT INITIALIZATION
--       MAX       MATCH
--       PC        SWATINIT
--       --------  --------
PPCWMAX
         100.0     YES                                     / TABLE NO 01
         125.0     YES                                     / TABLE NO 02
         135.0     YES                                     / TABLE NO 03
\end{verbatim}

The above example sets the maximum capillary pressure for three
saturation regions to 100, 125 and 135 with SWATINIT reset to the
connate water saturation for when the capillary pressure limit is
exceeded.

\subsection{PROPS -- Define the Start of the PROPS Section of
Keywords}\label{props-define-the-start-of-the-props-section-of-keywords}

\begin{longtable}[]{@{}
  >{\raggedright\arraybackslash}p{(\linewidth - 14\tabcolsep) * \real{0.1250}}
  >{\raggedright\arraybackslash}p{(\linewidth - 14\tabcolsep) * \real{0.1250}}
  >{\raggedright\arraybackslash}p{(\linewidth - 14\tabcolsep) * \real{0.1250}}
  >{\raggedright\arraybackslash}p{(\linewidth - 14\tabcolsep) * \real{0.1250}}
  >{\raggedright\arraybackslash}p{(\linewidth - 14\tabcolsep) * \real{0.1250}}
  >{\raggedright\arraybackslash}p{(\linewidth - 14\tabcolsep) * \real{0.1250}}
  >{\raggedright\arraybackslash}p{(\linewidth - 14\tabcolsep) * \real{0.1250}}
  >{\raggedright\arraybackslash}p{(\linewidth - 14\tabcolsep) * \real{0.1250}}@{}}
\toprule\noalign{}
\begin{minipage}[b]{\linewidth}\raggedright
RUNSPEC
\end{minipage} & \begin{minipage}[b]{\linewidth}\raggedright
GRID
\end{minipage} & \begin{minipage}[b]{\linewidth}\raggedright
EDIT
\end{minipage} & \begin{minipage}[b]{\linewidth}\raggedright
PROPS
\end{minipage} & \begin{minipage}[b]{\linewidth}\raggedright
REGIONS
\end{minipage} & \begin{minipage}[b]{\linewidth}\raggedright
SOLUTION
\end{minipage} & \begin{minipage}[b]{\linewidth}\raggedright
SUMMARY
\end{minipage} & \begin{minipage}[b]{\linewidth}\raggedright
SCHEDULE
\end{minipage} \\
\midrule\noalign{}
\endhead
\bottomrule\noalign{}
\endlastfoot
\end{longtable}

\subsubsection{Description}\label{description-588}

The PROPS activation keyword marks the end of the EDIT section and the
start of the PROPS section that defines the key fluid and rock property
data for the simulator

There is no data required for this keyword.

\subsubsection{Example}\label{example-324}

\begin{verbatim}
-- ==============================================================================
--
-- PROPS SECTION
--
-- ==============================================================================
PROPS
\end{verbatim}

The above example marks the end of the EDIT section and the start of the
PROPS section in the OPM Flow data input file.

\subsection{PVCDO -- Oil PVT Properties for Dead Oil (Constant
Compressibility)}\label{pvcdo-oil-pvt-properties-for-dead-oil-constant-compressibility}

\begin{longtable}[]{@{}
  >{\raggedright\arraybackslash}p{(\linewidth - 14\tabcolsep) * \real{0.1250}}
  >{\raggedright\arraybackslash}p{(\linewidth - 14\tabcolsep) * \real{0.1250}}
  >{\raggedright\arraybackslash}p{(\linewidth - 14\tabcolsep) * \real{0.1250}}
  >{\raggedright\arraybackslash}p{(\linewidth - 14\tabcolsep) * \real{0.1250}}
  >{\raggedright\arraybackslash}p{(\linewidth - 14\tabcolsep) * \real{0.1250}}
  >{\raggedright\arraybackslash}p{(\linewidth - 14\tabcolsep) * \real{0.1250}}
  >{\raggedright\arraybackslash}p{(\linewidth - 14\tabcolsep) * \real{0.1250}}
  >{\raggedright\arraybackslash}p{(\linewidth - 14\tabcolsep) * \real{0.1250}}@{}}
\toprule\noalign{}
\begin{minipage}[b]{\linewidth}\raggedright
RUNSPEC
\end{minipage} & \begin{minipage}[b]{\linewidth}\raggedright
GRID
\end{minipage} & \begin{minipage}[b]{\linewidth}\raggedright
EDIT
\end{minipage} & \begin{minipage}[b]{\linewidth}\raggedright
PROPS
\end{minipage} & \begin{minipage}[b]{\linewidth}\raggedright
REGIONS
\end{minipage} & \begin{minipage}[b]{\linewidth}\raggedright
SOLUTION
\end{minipage} & \begin{minipage}[b]{\linewidth}\raggedright
SUMMARY
\end{minipage} & \begin{minipage}[b]{\linewidth}\raggedright
SCHEDULE
\end{minipage} \\
\midrule\noalign{}
\endhead
\bottomrule\noalign{}
\endlastfoot
\end{longtable}

\subsubsection{Description}\label{description-589}

PVCDO defines the oil PVT properties for dead oil {[}``Dead'' oil is oil
that it contains no dissolved gas or a relatively thick oil or residue
that has lost its volatile components.{]} with constant compressibility.
If the oil has a constant and uniform dissolved gas concentration,
Gas-Oil Ratio (``GOR''), and if the reservoir pressure never drops below
the saturation pressure (bubble point pressure), then the model can be
run more efficiently by omitting the GAS and DISGAS keywords from the
RUNSPEC section, treating the oil as a dead oil, and defining a constant
Rs (GOR) value with keyword RSCONST or RSCONSTT in the PROPS section.
This results in the model being run as a dead oil problem with no active
gas phase. However, OPM Flow takes into account the constant Rs in the
calculations and reporting.

\begin{longtable}[]{@{}
  >{\raggedright\arraybackslash}p{(\linewidth - 6\tabcolsep) * \real{0.2500}}
  >{\raggedright\arraybackslash}p{(\linewidth - 6\tabcolsep) * \real{0.2500}}
  >{\raggedright\arraybackslash}p{(\linewidth - 6\tabcolsep) * \real{0.2500}}
  >{\raggedright\arraybackslash}p{(\linewidth - 6\tabcolsep) * \real{0.2500}}@{}}
\toprule\noalign{}
\begin{minipage}[b]{\linewidth}\raggedright
No.
\end{minipage} & \begin{minipage}[b]{\linewidth}\raggedright
Name
\end{minipage} & \begin{minipage}[b]{\linewidth}\raggedright
Description
\end{minipage} & \begin{minipage}[b]{\linewidth}\raggedright
Default
\end{minipage} \\
\midrule\noalign{}
\endhead
\bottomrule\noalign{}
\endlastfoot
Field & Metric & Laboratory & \\
1 & PRESS & PRESS is a real positive value defining the oil reference
pressure for the other parameters for this data set. & None \\
psia & barsa & atma & \\
2 & OFVF & OFVF is a real positive value defining the oil formation
volume factor (Bo) at the reference pressure. & None \\
rb/stb & rm3/sm3 & rcc/scc & \\
3 & OCOMP & OCOMP is a real positive value defining the oil
compressibility (Co) at the oil reference pressure and is defined as:
\({C}_{o} = -\frac{1}{{B}_{o}}\left(\frac{{\mathit{dB}}_{o}}{\mathit{dP}}\right)\)
& None \\
1/psia & 1/barsa & 1/atma & \\
4 & OVISC & OVISC is a real positive value defining the oil viscosity
(µo) at the oil reference pressure. & None \\
cP & cP & cP & \\
5 & OVISCOMP & OVISCOMP is a real positive value defining the oil
viscosibility (µoc) at the oil reference pressure and is defined as:
\({μ}_{\mathit{oc}} = \frac{1}{{μ}_{o}}\left(\frac{d{μ}_{o}}{\mathit{dP}}\right)\)
& None \\
1/psia & 1/barsa & 1/atma & \\
Notes: & & & \\
\end{longtable}

\emph{Table 8.3.225.1: PVCDO Keyword Description}

See also the RSCONST and RSCONSTT keywords to define the constant Rs for
dead oil and PVDO as an alternative keyword to enter the dead oil
properties.

\subsubsection{Example}\label{example-325}

\begin{verbatim}
--
--       OIL PVT TABLE FOR DEAD WITH CONSTANT COMPRESSIBILITY
--
PVCDO
--       REF PRES  BO         CO        VISC     VISC
--       PSIA      RB/STB     1/PSIA    CPOISE   GRAD
--       --------  --------   -------   ------   ------
          3840.0    1.080     1.5E-6     1.750    0.0       / TABLE NO. 01
          3840.0    1.100     1.5E-6     1.050    0.0       / TABLE NO. 02
          3840.0    1.120     1.6E-6     0.950    0.0       / TABLE NO. 03
          3840.0    1.140     1.7E-6     0.850    0.0       / TABLE NO. 04
          3840.0    1.160     1.7E-6     0.800    0.0       / TABLE NO. 05
\end{verbatim}

The above example defines five dead oil PVT tables with constant
compressibility and viscosity, and assumes that NTPVT equals five on the
TABDIMS keyword in the RUNSPEC section.

There is no terminating ``/'' for this keyword.

\subsection{PVCO -- Oil PVT Properties for Live
Oil}\label{pvco-oil-pvt-properties-for-live-oil}

\begin{longtable}[]{@{}
  >{\raggedright\arraybackslash}p{(\linewidth - 14\tabcolsep) * \real{0.1250}}
  >{\raggedright\arraybackslash}p{(\linewidth - 14\tabcolsep) * \real{0.1250}}
  >{\raggedright\arraybackslash}p{(\linewidth - 14\tabcolsep) * \real{0.1250}}
  >{\raggedright\arraybackslash}p{(\linewidth - 14\tabcolsep) * \real{0.1250}}
  >{\raggedright\arraybackslash}p{(\linewidth - 14\tabcolsep) * \real{0.1250}}
  >{\raggedright\arraybackslash}p{(\linewidth - 14\tabcolsep) * \real{0.1250}}
  >{\raggedright\arraybackslash}p{(\linewidth - 14\tabcolsep) * \real{0.1250}}
  >{\raggedright\arraybackslash}p{(\linewidth - 14\tabcolsep) * \real{0.1250}}@{}}
\toprule\noalign{}
\begin{minipage}[b]{\linewidth}\raggedright
RUNSPEC
\end{minipage} & \begin{minipage}[b]{\linewidth}\raggedright
GRID
\end{minipage} & \begin{minipage}[b]{\linewidth}\raggedright
EDIT
\end{minipage} & \begin{minipage}[b]{\linewidth}\raggedright
PROPS
\end{minipage} & \begin{minipage}[b]{\linewidth}\raggedright
REGIONS
\end{minipage} & \begin{minipage}[b]{\linewidth}\raggedright
SOLUTION
\end{minipage} & \begin{minipage}[b]{\linewidth}\raggedright
SUMMARY
\end{minipage} & \begin{minipage}[b]{\linewidth}\raggedright
SCHEDULE
\end{minipage} \\
\midrule\noalign{}
\endhead
\bottomrule\noalign{}
\endlastfoot
\end{longtable}

\subsubsection{Description}\label{description-590}

PVCO defines the oil PVT properties for live {[}``Live'' oil is oil that
contains gas in solution, which is normally the case for most
conventional oil reservoirs. However, for oil reservoirs classified as
heavy oil reservoirs, the in situ dissolved gas may be negligible and
oil would then be classified as gas-free oil which is commonly referred
to as ``dead'' oil.{]} and the keyword should only be used if the there
is both oil and gas phases in the model. This keyword should be used
when the DISGAS keyword has be declared in the RUNSPEC section
indicating that dissolved gas (more commonly referred to as solution
gas) is present in the oil. The keyword may be used for oil-water and
oil-water-gas input decks. This is an alternative keyword to the PVTO
keyword in the PROPS section that also enables entering live oil PVT
data. Here, the PVCO keyword assumes that for the undersaturated oil
with a given Gas-Oil Ratio (``GOR'' or ``Rs''), the oil compressibility
is independent of the pressure. Hence, is not necessary to enter the
undersaturated oil formation volume factor versus pressure data.
Similarly, the viscosity of the same type of oil is assumed to have a
pressure independent ``viscosibility'' derivative, and therefore it is
not necessary to enter undersaturated viscosity versus pressure data.

This keyword is not supported by OPM Flow but would change the results
if supported so the simulation will be stopped; however, the PVTO
keyword in the PROPS section may be used to enter live oil PVT data
instead.

\begin{longtable}[]{@{}
  >{\raggedright\arraybackslash}p{(\linewidth - 6\tabcolsep) * \real{0.2500}}
  >{\raggedright\arraybackslash}p{(\linewidth - 6\tabcolsep) * \real{0.2500}}
  >{\raggedright\arraybackslash}p{(\linewidth - 6\tabcolsep) * \real{0.2500}}
  >{\raggedright\arraybackslash}p{(\linewidth - 6\tabcolsep) * \real{0.2500}}@{}}
\toprule\noalign{}
\begin{minipage}[b]{\linewidth}\raggedright
No.
\end{minipage} & \begin{minipage}[b]{\linewidth}\raggedright
Name
\end{minipage} & \begin{minipage}[b]{\linewidth}\raggedright
Description
\end{minipage} & \begin{minipage}[b]{\linewidth}\raggedright
Default
\end{minipage} \\
\midrule\noalign{}
\endhead
\bottomrule\noalign{}
\endlastfoot
Field & Metric & Laboratory & \\
1 & PRESS & PRESS is a real columnar vector of real monotonically
increasing down the column values that defines the oil phase saturation
pressure (bubble-point pressure), that defines the RS, oil formation
volume factor and the oil viscosity at PRESS. & None \\
psia & barsa & atma & \\
2 & RS & RS is a real monotonically increasing down the column values
that defines the saturated gas-oil ratio (``GOR'') or Rs, for the given
value of PRESS. & 1* \\
Mscf/stb & sm3/sm3 & scc/scc & \\
3 & OFVF & OFVF is a real positive value defining the oil saturated
formation volume factor (Bo) at the saturation pressure PRESS. & None \\
rb/stb & rm3/sm3 & rcc/scc & \\
4 & OVISC & OVISC is a real positive value defining the oil viscosity
(µo) at the oil saturated reference pressure, PRESS. & 1* \\
cP & cP & cP & \\
5 & OCOMP & OCOMP is a real positive value defining the oil
compressibility (Co) at the saturated oil reference pressure and is
defined as:
\({C}_{o} = -\frac{1}{{B}_{o}}\left(\frac{{\mathit{dB}}_{o}}{\mathit{dP}}\right)\)
& 1* \\
1/psia & 1/barsa & 1/atma & \\
6 & OVISCOMP & OVISCOMP is a real positive value defining the oil
viiscosibility (µoc) at the saturated oil reference pressure with the
given RS, where (µoc) is defined as:
\({\mathrm{μ}}_{\mathit{oc}} = -\frac{1}{{\mathrm{μ}}_{o}}\left(\frac{d{\mathrm{μ}}_{o}}{\mathit{dP}}\right)\)
& 1* \\
1/psia & 1/barsa & 1/atma & \\
Notes: & & & \\
\end{longtable}

\emph{Table 8.113: PVCO Keyword Description}

\subsubsection{Example}\label{example-326}

\begin{verbatim}
--
--       OIL PVT TABLE FOR LIVE OIL
--
PVCO
--       PSAT      RS        BO       VISC    OIL      OIL
--       PSIA      MSCF/STB  RB/STB   CPOISE  COMPRES  VISCOS
--       --------  --------  -------  ------  -------  ------
            14.7    0.0010   1.05340  1.7230  3.0E-5   1*
           500.0    0.0890   1.08890  1.1670  1*       1*
          1000.0    0.2060   1.13850  0.8570  1*       1*
          1500.0    0.3360   1.19640  0.6840  1*       1*
          2000.0    0.4750   1.26110  0.5750  1*       1*
          2500.0    0.6220   1.33160  0.5000  1*       1*
          3000.0    0.7750   1.40740  0.4450  1*       1*
          3500.0    0.9330   1.48790  0.4020  1*       1*
          4000.0    1.0960   1.57280  0.3680  1*       1*
          4258.0    1.1800   1.61760  0.3530  1*       1*
          4500.0    1.2630   1.66190  0.3400  1*       1*
          5000.0    1.4340   1.75480  0.3170  1*       1*
          5500.0    1.6060   1.85020  0.2980  1*       1*    / TABLE NO. 01
--
--       PSAT      RS        BO       VISC    OIL      OIL
--       PSIA      MSCF/STB  RB/STB   CPOISE  COMPRES  VISCOS
--       --------  --------  -------  ------  -------  ------
            14.7    0.0010   1.05340  1.7230  3.0E-5   1*
           500.0    0.0890   1.08890  1.1670  1*       1*
          1000.0    0.2060   1.13850  0.8570  1*       1*
          1500.0    0.3360   1.19640  0.6840  1*       1*
          2000.0    0.4750   1.26110  0.5750  1*       1*
          2500.0    0.6220   1.33160  0.5000  1*       1*
          3000.0    0.7750   1.40740  0.4450  1*       1*
          3500.0    0.9330   1.48790  0.4020  1*       1*
          4000.0    1.0960   1.57280  0.3680  1*       1*
          4258.0    1.1800   1.61760  0.3530  1*       1*
          4500.0    1.2630   1.66190  0.3400  1*       1*
          5000.0    1.4340   1.75480  0.3170  1*       1*
          5500.0    1.6060   1.85020  0.2980  1*       1*    / TABLE NO. 02
\end{verbatim}

The example defines two live oil PVT tables with constant
compressibility above the saturation pressure, and assumes that NTPVT
equals two on the TABDIMS keyword in the RUNSPEC section.

\subsection{PVDG -- Gas PVT Properties for Dry
Gas}\label{pvdg-gas-pvt-properties-for-dry-gas}

\begin{longtable}[]{@{}
  >{\raggedright\arraybackslash}p{(\linewidth - 14\tabcolsep) * \real{0.1250}}
  >{\raggedright\arraybackslash}p{(\linewidth - 14\tabcolsep) * \real{0.1250}}
  >{\raggedright\arraybackslash}p{(\linewidth - 14\tabcolsep) * \real{0.1250}}
  >{\raggedright\arraybackslash}p{(\linewidth - 14\tabcolsep) * \real{0.1250}}
  >{\raggedright\arraybackslash}p{(\linewidth - 14\tabcolsep) * \real{0.1250}}
  >{\raggedright\arraybackslash}p{(\linewidth - 14\tabcolsep) * \real{0.1250}}
  >{\raggedright\arraybackslash}p{(\linewidth - 14\tabcolsep) * \real{0.1250}}
  >{\raggedright\arraybackslash}p{(\linewidth - 14\tabcolsep) * \real{0.1250}}@{}}
\toprule\noalign{}
\begin{minipage}[b]{\linewidth}\raggedright
RUNSPEC
\end{minipage} & \begin{minipage}[b]{\linewidth}\raggedright
GRID
\end{minipage} & \begin{minipage}[b]{\linewidth}\raggedright
EDIT
\end{minipage} & \begin{minipage}[b]{\linewidth}\raggedright
PROPS
\end{minipage} & \begin{minipage}[b]{\linewidth}\raggedright
REGIONS
\end{minipage} & \begin{minipage}[b]{\linewidth}\raggedright
SOLUTION
\end{minipage} & \begin{minipage}[b]{\linewidth}\raggedright
SUMMARY
\end{minipage} & \begin{minipage}[b]{\linewidth}\raggedright
SCHEDULE
\end{minipage} \\
\midrule\noalign{}
\endhead
\bottomrule\noalign{}
\endlastfoot
\end{longtable}

\subsubsection{Description}\label{description-591}

PVDG defines the gas PVT properties for dry gas {[}Natural gas that
occurs in the absence of condensate or liquid hydrocarbons, or gas that
had condensable hydrocarbons removed, is called dry gas. It is primarily
methane with some intermediates. The hydrocarbon mixture is solely gas
in the reservoir and there is no liquid (condensate surface liquid)
formed either in the reservoir or at surface. The term dry indicates
that the gas does not contain heavier hydrocarbons to form liquids at
the surface conditions. Dry gas typically has GOR's greater than 100,000
scf/stb or 18,000 Sm3/m3.{]}. If the gas has a constant and uniform
vaporized oil concentration, Condensate-Gas Ratio (``CGR''), and if the
reservoir pressure never drops below the saturation pressure (dew point
pressure), then the model can be run more efficiently by omitting the
OIL and VAPOIL keywords from the RUNSPEC section, treating the gas as a
dry gas, and defining a constant Rv (CGR) value with keyword RVCONST or
RVCONSTT in the PROPS section. This results in the model being run with
as a dry gas problem with no active oil (condensate) phase. However, OPM
Flow takes into account the constant Rv in the calculations and
reporting.

\begin{longtable}[]{@{}
  >{\raggedright\arraybackslash}p{(\linewidth - 6\tabcolsep) * \real{0.2500}}
  >{\raggedright\arraybackslash}p{(\linewidth - 6\tabcolsep) * \real{0.2500}}
  >{\raggedright\arraybackslash}p{(\linewidth - 6\tabcolsep) * \real{0.2500}}
  >{\raggedright\arraybackslash}p{(\linewidth - 6\tabcolsep) * \real{0.2500}}@{}}
\toprule\noalign{}
\begin{minipage}[b]{\linewidth}\raggedright
No.
\end{minipage} & \begin{minipage}[b]{\linewidth}\raggedright
Name
\end{minipage} & \begin{minipage}[b]{\linewidth}\raggedright
Description
\end{minipage} & \begin{minipage}[b]{\linewidth}\raggedright
Default
\end{minipage} \\
\midrule\noalign{}
\endhead
\bottomrule\noalign{}
\endlastfoot
Field & Metric & Laboratory & \\
1 & PRESS & A columnar vector of real monotonically increasing down the
column values that defines the gas phase pressure. & None \\
psia & barsa & atma & \\
2 & GFVF & A columnar vector of real decreasing down the column values
that defines the corresponding gas phase formation volume factor. &
None \\
rb/Mscf & rm3/sm3 & rcc/scc & \\
3 & GVISC & A columnar vector of real increasing down the column values
that defines the corresponding gas phase viscosity. & None \\
cP & cP & cP & \\
Notes: & & & \\
\end{longtable}

\emph{Table 8.114: PVDG Keyword Description}

Note that provided the first table has been entered, subsequent tables
may be defaulted, in this case the prior table is copied to the current
table. See the second example for an illustration on how to use this
feature.

See also the RVCONST and RVCONSTT keywords to define the constant Rv for
dry gas.

\subsubsection{Example}\label{example-327}

The first example below defines two dry gas PVT tables assuming NTPVT
equals two and NPPVT is greater than or equal to 22 on the TABDIMS
keyword in the RUNSPEC section.

\begin{verbatim}
--
--       GAS PVT TABLE FOR DRY GAS
--
PVDG
--       PRES     BG          VISC
--       PSIA     RB/MSCF     CPOISE
--       ------   --------    ------
           14.7   197.8092    0.0129
           50.0    65.9364    0.0130
          100.0    31.6495    0.0130
          230.0    13.8813    0.0131
          460.0     6.8210    0.0132
          690.0     4.4703    0.0135
          920.0     3.2968    0.0138
         1150.0     2.6113    0.0141
         1380.0     2.1560    0.0145
         1610.0     1.8316    0.0150
         1840.0     1.5952    0.0155
         2070.0     1.4129    0.0161
         2300.0     1.2700    0.0167
         2372.0     1.2305    0.0169
         2530.0     1.1551    0.0174
         2760.0     1.0621    0.0181
         2990.0     0.9841    0.0189
         3220.0     0.9190    0.0196
         3450.0     0.8638    0.0204
         4500.0     0.6910    0.0242
         6000.0     0.5616    0.0293                       / TABLE N0. 01
--
--       PRES     BG          VISC
--       PSIA     RB/MSCF     CPOISE
--       ------   --------    ------
           14.7   265.0126    0.0133
           50.0    66.2531    0.0133
          100.0    33.1266    0.0133
          230.0    14.4552    0.0134
          460.0     7.0357    0.0136
          690.0     4.6493    0.0138
          920.0     3.4417    0.0140
         1150.0     2.7227    0.0144
         1380.0     2.2522    0.0147
         1610.0     1.9158    0.0151
         1840.0     1.6702    0.0156
         2070.0     1.4805    0.0162
         2300.0     1.3317    0.0167
         2372.0     1.2927    0.0169
         2530.0     1.2119    0.0173
         2760.0     1.1135    0.0180
         2990.0     1.0325    0.0187
         3220.0     0.9637    0.0194
         3450.0     0.9055    0.0201
         4500.0     0.7228    0.0236
         6000.0     0.5837    0.0285                       / TABLE N0. 02
\end{verbatim}

The second example defines four dry gas PVT tables assuming NTPVT equals
four and NPPVT is greater than or equal to 22 on the TABDIMS keyword in
the RUNSPEC section. Here table two defaults to table one, and table
four defaults to table three.

\begin{verbatim}
--
--       GAS PVT TABLE FOR DRY GAS
--
PVDG
--       PRES     BG          VISC
--       PSIA     RB/MSCF     CPOISE
--       ------   --------    ------
           14.7   197.8092    0.0129
           50.0    65.9364    0.0130
          100.0    31.6495    0.0130
          230.0    13.8813    0.0131
          460.0     6.8210    0.0132
          690.0     4.4703    0.0135
          920.0     3.2968    0.0138
         1150.0     2.6113    0.0141
         1380.0     2.1560    0.0145
         1610.0     1.8316    0.0150
         1840.0     1.5952    0.0155
         2070.0     1.4129    0.0161
         2300.0     1.2700    0.0167
         2372.0     1.2305    0.0169
         2530.0     1.1551    0.0174
         2760.0     1.0621    0.0181
         2990.0     0.9841    0.0189
         3220.0     0.9190    0.0196
         3450.0     0.8638    0.0204
         4500.0     0.6910    0.0242
         6000.0     0.5616    0.0293                       / TABLE N0. 01
                                                           / TABLE N0. 02
--
--       PRES     BG          VISC
--       PSIA     RB/MSCF     CPOISE
--       ------   --------    ------
           14.7   265.0126    0.0133
           50.0    66.2531    0.0133
          100.0    33.1266    0.0133
          230.0    14.4552    0.0134
          460.0     7.0357    0.0136
          690.0     4.6493    0.0138
          920.0     3.4417    0.0140
         1150.0     2.7227    0.0144
         1380.0     2.2522    0.0147
         1610.0     1.9158    0.0151
         1840.0     1.6702    0.0156
         2070.0     1.4805    0.0162
         2300.0     1.3317    0.0167
         2372.0     1.2927    0.0169
         2530.0     1.2119    0.0173
         2760.0     1.1135    0.0180
         2990.0     1.0325    0.0187
         3220.0     0.9637    0.0194
         3450.0     0.9055    0.0201
         4500.0     0.7228    0.0236
         6000.0     0.5837    0.0285                       / TABLE N0. 03
                                                           / TABLE N0. 04
\end{verbatim}

There is no terminating ``/'' for this keyword.

\subsection{PVDO -- Oil PVT Properties for Dead
Oil}\label{pvdo-oil-pvt-properties-for-dead-oil}

\begin{longtable}[]{@{}
  >{\raggedright\arraybackslash}p{(\linewidth - 14\tabcolsep) * \real{0.1250}}
  >{\raggedright\arraybackslash}p{(\linewidth - 14\tabcolsep) * \real{0.1250}}
  >{\raggedright\arraybackslash}p{(\linewidth - 14\tabcolsep) * \real{0.1250}}
  >{\raggedright\arraybackslash}p{(\linewidth - 14\tabcolsep) * \real{0.1250}}
  >{\raggedright\arraybackslash}p{(\linewidth - 14\tabcolsep) * \real{0.1250}}
  >{\raggedright\arraybackslash}p{(\linewidth - 14\tabcolsep) * \real{0.1250}}
  >{\raggedright\arraybackslash}p{(\linewidth - 14\tabcolsep) * \real{0.1250}}
  >{\raggedright\arraybackslash}p{(\linewidth - 14\tabcolsep) * \real{0.1250}}@{}}
\toprule\noalign{}
\begin{minipage}[b]{\linewidth}\raggedright
RUNSPEC
\end{minipage} & \begin{minipage}[b]{\linewidth}\raggedright
GRID
\end{minipage} & \begin{minipage}[b]{\linewidth}\raggedright
EDIT
\end{minipage} & \begin{minipage}[b]{\linewidth}\raggedright
PROPS
\end{minipage} & \begin{minipage}[b]{\linewidth}\raggedright
REGIONS
\end{minipage} & \begin{minipage}[b]{\linewidth}\raggedright
SOLUTION
\end{minipage} & \begin{minipage}[b]{\linewidth}\raggedright
SUMMARY
\end{minipage} & \begin{minipage}[b]{\linewidth}\raggedright
SCHEDULE
\end{minipage} \\
\midrule\noalign{}
\endhead
\bottomrule\noalign{}
\endlastfoot
\end{longtable}

\subsubsection{Description}\label{description-592}

PVDO defines the oil PVT properties for dead oil {[}``Dead'' oil is oil
that it contains no dissolved gas or a relatively thick oil or residue
that has lost its volatile components.{]}. If the oil has a constant and
uniform dissolved gas concentration, Gas-Oil Ratio (``GOR''), and if the
reservoir pressure never drops below the saturation pressure (bubble
point pressure), then the model can be run more efficiently by omitting
the GAS and DISGAS keywords from the RUNSPEC section, treating the oil
as a dead oil, and defining a constant Rs (GOR) value with keyword
RSCONST or RSCONSTT in the PROPS section. This results in the model
being run as a dead oil problem with no active gas phase. However, OPM
Flow takes into account the constant Rs in the calculations and
reporting.

\begin{longtable}[]{@{}
  >{\raggedright\arraybackslash}p{(\linewidth - 6\tabcolsep) * \real{0.2500}}
  >{\raggedright\arraybackslash}p{(\linewidth - 6\tabcolsep) * \real{0.2500}}
  >{\raggedright\arraybackslash}p{(\linewidth - 6\tabcolsep) * \real{0.2500}}
  >{\raggedright\arraybackslash}p{(\linewidth - 6\tabcolsep) * \real{0.2500}}@{}}
\toprule\noalign{}
\begin{minipage}[b]{\linewidth}\raggedright
No.
\end{minipage} & \begin{minipage}[b]{\linewidth}\raggedright
Name
\end{minipage} & \begin{minipage}[b]{\linewidth}\raggedright
Description
\end{minipage} & \begin{minipage}[b]{\linewidth}\raggedright
Default
\end{minipage} \\
\midrule\noalign{}
\endhead
\bottomrule\noalign{}
\endlastfoot
Field & Metric & Laboratory & \\
1 & PRESS & A columnar vector of real monotonically increasing down the
column values that defines the oil phase pressure. & None \\
psia & barsa & atma & \\
2 & OFVF & A columnar vector of real decreasing down the column values
that defines the corresponding oil phase formation volume factor. &
None \\
rb/stb & rm3/sm3 & rcc/scc & \\
3 & OVISC & A columnar vector of real increasing down the column values
that defines the corresponding oil phase viscosity. & None \\
cP & cP & cP & \\
Notes: & & & \\
\end{longtable}

\emph{Table 8.115: PVDO Keyword Description}

Note that provided the first table has been entered, subsequent tables
may be defaulted, in this case the prior table is copied to the current
table. See the second example for an illustration on how to use this
feature.

See also the RSCONST and RSCONSTT keywords to define the constant Rs for
dead oil and PVCDO as an alternative keyword to enter the dead oil
properties.

\subsubsection{Example}\label{example-328}

The example below defines two dead oil PVT tables with variable
viscosity and compressibility with respect to pressure, and assumes that
NTPVT equals two and NPPVT is greater than or equal to eight on the
TABDIMS keyword in the RUNSPEC section.

\begin{verbatim}
--
--      OIL PVT TABLE FOR DEAD OIL
--
PVDO
--      PSAT      BO        VISC
--      PSIA      RB/STB    CPOISE
--      --------  -------   ------
           400    1.0102     1.16
          1200    1.0040     1.164
          2000    0.9960     1.167
          2800    0.9880     1.172
          3600    0.9802     1.177
          4400    0.9724     1.181
          5200    0.9646     1.185
          5600    0.9607     1.19                           / TABLE NO. 01
--      --------  -------   ------
           800    1.0255     1.14
          1600    1.0172     1.14
          2400    1.0091     1.14
          3200    1.0011     1.14
          4000    0.9931     1.14
          4800    0.9852     1.14
          5600    0.9774     1.14                           / TABLE NO. 02
\end{verbatim}

The second example defines four dead oil PVT tables with variable
viscosity and compressibility with respect to pressure, and assumes that
NTPVT equals four and NPPVT is greater than or equal to eight on the
TABDIMS keyword in the RUNSPEC section. Here table two defaults to table
one, and table four defaults to table three.

\begin{verbatim}
--
--      OIL PVT TABLE FOR DEAD OIL
--
PVDO
--      PSAT      BO        VISC
--      PSIA      RB/STB    CPOISE
--      --------  -------   ------
           400    1.0102     1.16
          1200    1.0040     1.164
          2000    0.9960     1.167
          2800    0.9880     1.172
          3600    0.9802     1.177
          4400    0.9724     1.181
          5200    0.9646     1.185
          5600    0.9607     1.19                           / TABLE NO. 01
                                                            / TABLE NO. 02
           800    1.0255     1.14
          1600    1.0172     1.14
          2400    1.0091     1.14
          3200    1.0011     1.14
          4000    0.9931     1.14
          4800    0.9852     1.14
          5600    0.9774     1.14                           / TABLE NO. 03
                                                            / TABLE NO. 04
\end{verbatim}

There is no terminating ``/'' for this keyword.

\subsection{PVDS -- Solvent PVT Properties for the Solvent
Model}\label{pvds-solvent-pvt-properties-for-the-solvent-model}

\begin{longtable}[]{@{}
  >{\raggedright\arraybackslash}p{(\linewidth - 14\tabcolsep) * \real{0.1250}}
  >{\raggedright\arraybackslash}p{(\linewidth - 14\tabcolsep) * \real{0.1250}}
  >{\raggedright\arraybackslash}p{(\linewidth - 14\tabcolsep) * \real{0.1250}}
  >{\raggedright\arraybackslash}p{(\linewidth - 14\tabcolsep) * \real{0.1250}}
  >{\raggedright\arraybackslash}p{(\linewidth - 14\tabcolsep) * \real{0.1250}}
  >{\raggedright\arraybackslash}p{(\linewidth - 14\tabcolsep) * \real{0.1250}}
  >{\raggedright\arraybackslash}p{(\linewidth - 14\tabcolsep) * \real{0.1250}}
  >{\raggedright\arraybackslash}p{(\linewidth - 14\tabcolsep) * \real{0.1250}}@{}}
\toprule\noalign{}
\begin{minipage}[b]{\linewidth}\raggedright
RUNSPEC
\end{minipage} & \begin{minipage}[b]{\linewidth}\raggedright
GRID
\end{minipage} & \begin{minipage}[b]{\linewidth}\raggedright
EDIT
\end{minipage} & \begin{minipage}[b]{\linewidth}\raggedright
PROPS
\end{minipage} & \begin{minipage}[b]{\linewidth}\raggedright
REGIONS
\end{minipage} & \begin{minipage}[b]{\linewidth}\raggedright
SOLUTION
\end{minipage} & \begin{minipage}[b]{\linewidth}\raggedright
SUMMARY
\end{minipage} & \begin{minipage}[b]{\linewidth}\raggedright
SCHEDULE
\end{minipage} \\
\midrule\noalign{}
\endhead
\bottomrule\noalign{}
\endlastfoot
\end{longtable}

\subsubsection{Description}\label{description-593}

PVDS defines the solvent PVT properties for use with SOLVENT option. The
solvent is treated as an additional dry gas phase within the model. This
keyword should only be used if the SOLVENT model has been invoked in the
RUNSPEC section.

\begin{longtable}[]{@{}
  >{\raggedright\arraybackslash}p{(\linewidth - 6\tabcolsep) * \real{0.2500}}
  >{\raggedright\arraybackslash}p{(\linewidth - 6\tabcolsep) * \real{0.2500}}
  >{\raggedright\arraybackslash}p{(\linewidth - 6\tabcolsep) * \real{0.2500}}
  >{\raggedright\arraybackslash}p{(\linewidth - 6\tabcolsep) * \real{0.2500}}@{}}
\toprule\noalign{}
\begin{minipage}[b]{\linewidth}\raggedright
No.
\end{minipage} & \begin{minipage}[b]{\linewidth}\raggedright
Name
\end{minipage} & \begin{minipage}[b]{\linewidth}\raggedright
Description
\end{minipage} & \begin{minipage}[b]{\linewidth}\raggedright
Default
\end{minipage} \\
\midrule\noalign{}
\endhead
\bottomrule\noalign{}
\endlastfoot
Field & Metric & Laboratory & \\
1 & PRESS & A columnar vector of real monotonically increasing down the
column values that defines the solvent phase pressure. & None \\
psia & barsa & atma & \\
2 & GFVF & A columnar vector of real decreasing down the column values
that defines the corresponding solvent phase formation volume factor. &
None \\
rb/Mscf & rm3/sm3 & rcc/scc & \\
3 & GVISC & A columnar vector of real increasing down the column values
that defines the corresponding solvent phase viscosity. & None \\
cP & cP & cP & \\
Notes: & & & \\
\end{longtable}

\emph{Table 8.116: PVDS Keyword Description}

\subsubsection{Example}\label{example-329}

\begin{verbatim}
--
--       GAS SOLVENT PVT TABLE
--
PVDS
--       PRES     BG          VISC
--       PSIA     RB/MSCF     CPOISE
--       ------   --------    ------
          700.0     4.4703    0.0135
          920.0     3.2968    0.0138
         1150.0     2.6113    0.0141
         1380.0     2.1560    0.0145
         1610.0     1.8316    0.0150
         1840.0     1.5952    0.0155
         2070.0     1.4129    0.0161
         2300.0     1.2700    0.0167
         2372.0     1.2305    0.0169
         2530.0     1.1551    0.0174
         2760.0     1.0621    0.0181
         2990.0     0.9841    0.0189
         3220.0     0.9190    0.0196
         3450.0     0.8638    0.0204
         4500.0     0.6910    0.0242
         6000.0     0.5616    0.0293                       / TABLE N0. 01
--
--       PRES     BG          VISC
--       PSIA     RB/MSCF     CPOISE
--       ------   --------    ------
          700.0     4.6493    0.0138
          920.0     3.4417    0.0140
         1150.0     2.7227    0.0144
         1380.0     2.2522    0.0147
         1610.0     1.9158    0.0151
         1840.0     1.6702    0.0156
         2070.0     1.4805    0.0162
         2300.0     1.3317    0.0167
         2372.0     1.2927    0.0169
         2530.0     1.2119    0.0173
         2760.0     1.1135    0.0180
         2990.0     1.0325    0.0187
         3220.0     0.9637    0.0194
         3450.0     0.9055    0.0201
         4500.0     0.7228    0.0236
         6000.0     0.5837    0.0285                       / TABLE N0. 02
\end{verbatim}

The above example defines two solvent PVT tables assuming NTPVT equals
two and NPPVT is greater than or equal to 16 on the TABDIMS keyword in
the RUNSPEC section.

There is no terminating ``/'' for this keyword.

\subsection{PVTG -- Gas PVT Properties for Wet Gas with Vaporized
Oil}\label{pvtg-gas-pvt-properties-for-wet-gas-with-vaporized-oil}

\begin{longtable}[]{@{}
  >{\raggedright\arraybackslash}p{(\linewidth - 14\tabcolsep) * \real{0.1250}}
  >{\raggedright\arraybackslash}p{(\linewidth - 14\tabcolsep) * \real{0.1250}}
  >{\raggedright\arraybackslash}p{(\linewidth - 14\tabcolsep) * \real{0.1250}}
  >{\raggedright\arraybackslash}p{(\linewidth - 14\tabcolsep) * \real{0.1250}}
  >{\raggedright\arraybackslash}p{(\linewidth - 14\tabcolsep) * \real{0.1250}}
  >{\raggedright\arraybackslash}p{(\linewidth - 14\tabcolsep) * \real{0.1250}}
  >{\raggedright\arraybackslash}p{(\linewidth - 14\tabcolsep) * \real{0.1250}}
  >{\raggedright\arraybackslash}p{(\linewidth - 14\tabcolsep) * \real{0.1250}}@{}}
\toprule\noalign{}
\begin{minipage}[b]{\linewidth}\raggedright
RUNSPEC
\end{minipage} & \begin{minipage}[b]{\linewidth}\raggedright
GRID
\end{minipage} & \begin{minipage}[b]{\linewidth}\raggedright
EDIT
\end{minipage} & \begin{minipage}[b]{\linewidth}\raggedright
PROPS
\end{minipage} & \begin{minipage}[b]{\linewidth}\raggedright
REGIONS
\end{minipage} & \begin{minipage}[b]{\linewidth}\raggedright
SOLUTION
\end{minipage} & \begin{minipage}[b]{\linewidth}\raggedright
SUMMARY
\end{minipage} & \begin{minipage}[b]{\linewidth}\raggedright
SCHEDULE
\end{minipage} \\
\midrule\noalign{}
\endhead
\bottomrule\noalign{}
\endlastfoot
\end{longtable}

\subsubsection{Description}\label{description-594}

PVTG defines the gas PVT properties for wet gas {[}Natural gas that
contains significant heavy hydrocarbons such as propane, butane and
other liquid hydrocarbons is known as wet gas or rich gas. The general
rule of thumb is if the gas contains less methane (typically less than
85\% methane) and more ethane, and other more complex hydrocarbons, it
is labeled as wet gas. Wet gas normally has GOR's less than 100,000
scf/stb or 18,000 Sm3/m3, with the condensate having a gravity greater
than 50 oAPI.{]}. This keyword should be used when the VAPOIL keyword
has be declared in the RUNSPEC section indicating that that vaporized
oil (more commonly referred to as condensate) is present in the wet gas
phase. The keyword may be used for gas-water and oil-water-gas input
decks that contain the oil and gas phases.

\begin{longtable}[]{@{}
  >{\raggedright\arraybackslash}p{(\linewidth - 8\tabcolsep) * \real{0.2000}}
  >{\raggedright\arraybackslash}p{(\linewidth - 8\tabcolsep) * \real{0.2000}}
  >{\raggedright\arraybackslash}p{(\linewidth - 8\tabcolsep) * \real{0.2000}}
  >{\raggedright\arraybackslash}p{(\linewidth - 8\tabcolsep) * \real{0.2000}}
  >{\raggedright\arraybackslash}p{(\linewidth - 8\tabcolsep) * \real{0.2000}}@{}}
\toprule\noalign{}
\begin{minipage}[b]{\linewidth}\raggedright
No.
\end{minipage} & \begin{minipage}[b]{\linewidth}\raggedright
Name
\end{minipage} & \begin{minipage}[b]{\linewidth}\raggedright
Description
\end{minipage} & \begin{minipage}[b]{\linewidth}\raggedright
Default
\end{minipage} & \begin{minipage}[b]{\linewidth}\raggedright
\end{minipage} \\
\midrule\noalign{}
\endhead
\bottomrule\noalign{}
\endlastfoot
Field & Metric & Laboratory & & \\
1 & PRESS & A real monotonically increasing down the column vector that
defines the gas phase pressure, associated with the saturated
condensate-gas ratio (``CGR'') or Rv, the gas formation volume factor
and the gas viscosity for the corresponding pressure for the stated
saturated RVS. For a given PRESS the variability of the gas formation
volume factor and the gas viscosity with respect to the under-saturated
Rv is optionally included as a sub table under RVU, FVFU and VISU
columns, that is it is not necessary to repeat PRESS for each sub table
entry. However, each sub table must be terminated by a ``/''. The under
saturated Rv entries are optional, except for perhaps the last PRESS
entry to define the PVT properties above the initial saturation
pressure. & None & \\
psia & barsa & atma & & \\
2 & RVS & RVU & A columnar vector of real positive number for both the
saturated (RVS) and under saturated (RVU) Rv sub table entries. The RVS
entry on the main table is the saturated CGR at the pressure indicated
by PRESS and may be increasing or decreasing in value as PRESS varies.
Subsequent under-saturated Rv for a sub table at the given PRESS, as
defined by RVU, are monotonically decreasing for entries in a given sub
table. & None \\
stb/Mscf & sm3/sm3 & scc/scc & & \\
3 & FVFS & FVFU & A columnar vector of real decreasing down the column
values that defines the corresponding gas phase formation volume factor
for a given pressure (PRESS) and for a given Rv (either RVS or RVU). &
None \\
rb/Mscf & rm3/sm3 & rcc/scc & & \\
4 & VISS & VISU & VISS a columnar vector of real increasing down the
column values that defines the corresponding gas phase viscosity for a
given pressure (PRESS) and for a given RVS. VISU a columnar vector of
real decreasing from VISS down the column values that defines the
corresponding gas phase viscosity for a given pressure (PRESS) and for a
given RVU. & None \\
cP & cP & cP & & \\
Notes: & & & & \\
\end{longtable}

\emph{Table 8.117: PVTG Keyword Description}

Note that provided the first table has been entered, subsequent tables
may be defaulted, in this case the prior table is copied to the current
table. See the second example for an illustration on how to use this
feature.

\begin{tcolorbox}[enhanced jigsaw, title=\textcolor{quarto-callout-note-color}{\faInfo}\hspace{0.5em}{Note}, coltitle=black, breakable, toprule=.15mm, colbacktitle=quarto-callout-note-color!10!white, leftrule=.75mm, opacitybacktitle=0.6, colback=white, bottomtitle=1mm, toptitle=1mm, titlerule=0mm, left=2mm, colframe=quarto-callout-note-color-frame, arc=.35mm, rightrule=.15mm, bottomrule=.15mm, opacityback=0]

If the VAPWAT keyword in the RUNSPEC section is also present in the
input deck, then the PVTG keyword in the PROPS section should be used to
define the gas properties as function of pressure and RV, assuming
water-saturated gas. Also, in this case, the PVTGW keyword, also in the
PROPS section, should also be in the input deck. In this case, PVTGW
defines the gas properties as function of pressure and RVW, assuming
oil-saturated gas.

\end{tcolorbox}

\subsubsection{Example}\label{example-330}

The first example defines two wet gas PVT tables assuming NTPVT equals
two, NPPVT is greater than or equal to eight, and NRPVT greater than or
equal to two on the TABDIMS keyword in the RUNSPEC section.

\begin{verbatim}
--
--       GAS PVT TABLE FOR WET GAS WITH VAPORIZED OIL
--
PVTG
--       PRES       RV           BG           VISC
--       BARSA    SM^3/SM^3     RM^3/SM^3       CPOISE
--      ------    ---------    -------       ------
           20     0.000132      0.042340      0.01344
                  0             0.042310      0.01389      /
           40     0.000124      0.020460      0.01420
                  0             0.020430      0.01450      /
           60     0.000126      0.013280      0.01526
                  0             0.013250      0.01532      /
           80     0.000135      0.009770      0.01660
                  0             0.009730      0.01634      /
          100     0.000149      0.007730      0.01818
                  0             0.007690      0.01752      /
          120     0.000163      0.006426      0.01994
                  0             0.006405      0.01883      /
          140     0.000191      0.005541      0.02181
                  0             0.005553      0.02021      /
          160     0.000225      0.004919      0.02370
                  0             0.004952      0.02163      /
                                                           / TABLE NO. 1
--       PRES       RV           BG           VISC
--       BARSA    SM^3/SM^3     RM^3/SM^3       CPOISE
--      ------    ---------    -------       ------
           20     0.000132      0.042340      0.01344      /
           40     0.000124      0.020460      0.01420      /
           60     0.000126      0.013280      0.01526      /
           80     0.000135      0.009770      0.01660      /
          100     0.000149      0.007730      0.01818      /
          120     0.000163      0.006426      0.01994      /
          140     0.000191      0.005541      0.02181      /
          160     0.000225      0.004919      0.02370
                  0             0.004952      0.02163      /
                                                           / TABLE NO. 2
\end{verbatim}

The second example defines four wet gas PVT tables assuming NTPVT equals
four, NPPVT is greater than or equal to eight, and NRPVT greater than or
equal to two on the TABDIMS keyword in the RUNSPEC section. Here table
two defaults to table one, and table four defaults to table three.

\begin{verbatim}
--
--       GAS PVT TABLE FOR WET GAS WITH VAPORIZED OIL
--
PVTG
--       PRES       RV           BG           VISC
--       BARSA    SM^3/SM^3     RM^3/SM^3       CPOISE
--      ------    ---------    -------       ------
           20     0.000132      0.042340      0.01344
                  0             0.042310      0.01389      /
           40     0.000124      0.020460      0.01420
                  0             0.020430      0.01450      /
           60     0.000126      0.013280      0.01526
                  0             0.013250      0.01532      /
           80     0.000135      0.009770      0.01660
                  0             0.009730      0.01634      /
          100     0.000149      0.007730      0.01818
                  0             0.007690      0.01752      /
          120     0.000163      0.006426      0.01994
                  0             0.006405      0.01883      /
          140     0.000191      0.005541      0.02181
                  0             0.005553      0.02021      /
          160     0.000225      0.004919      0.02370
                  0             0.004952      0.02163      /
                                                           / TABLE NO. 1
                                                           / TABLE NO. 2
--       PRES       RV           BG           VISC
--       BARSA    SM^3/SM^3     RM^3/SM^3       CPOISE
--      ------    ---------    -------       ------
           20     0.000132      0.042340      0.01344      /
           40     0.000124      0.020460      0.01420      /
           60     0.000126      0.013280      0.01526      /
           80     0.000135      0.009770      0.01660      /
          100     0.000149      0.007730      0.01818      /
          120     0.000163      0.006426      0.01994      /
          140     0.000191      0.005541      0.02181      /
          160     0.000225      0.004919      0.02370
                  0             0.004952      0.02163      /
                                                           / TABLE NO. 3
                                                           / TABLE NO. 4
\end{verbatim}

Notice that in both examples there is no terminating ``/'' for this
keyword only for a table and a sub table.

\subsection{PVTGW -- Gas PVT Properties for Dry Gas with Vaporized
Water}\label{pvtgw-gas-pvt-properties-for-dry-gas-with-vaporized-water}

\begin{longtable}[]{@{}
  >{\raggedright\arraybackslash}p{(\linewidth - 14\tabcolsep) * \real{0.1250}}
  >{\raggedright\arraybackslash}p{(\linewidth - 14\tabcolsep) * \real{0.1250}}
  >{\raggedright\arraybackslash}p{(\linewidth - 14\tabcolsep) * \real{0.1250}}
  >{\raggedright\arraybackslash}p{(\linewidth - 14\tabcolsep) * \real{0.1250}}
  >{\raggedright\arraybackslash}p{(\linewidth - 14\tabcolsep) * \real{0.1250}}
  >{\raggedright\arraybackslash}p{(\linewidth - 14\tabcolsep) * \real{0.1250}}
  >{\raggedright\arraybackslash}p{(\linewidth - 14\tabcolsep) * \real{0.1250}}
  >{\raggedright\arraybackslash}p{(\linewidth - 14\tabcolsep) * \real{0.1250}}@{}}
\toprule\noalign{}
\begin{minipage}[b]{\linewidth}\raggedright
RUNSPEC
\end{minipage} & \begin{minipage}[b]{\linewidth}\raggedright
GRID
\end{minipage} & \begin{minipage}[b]{\linewidth}\raggedright
EDIT
\end{minipage} & \begin{minipage}[b]{\linewidth}\raggedright
PROPS
\end{minipage} & \begin{minipage}[b]{\linewidth}\raggedright
REGIONS
\end{minipage} & \begin{minipage}[b]{\linewidth}\raggedright
SOLUTION
\end{minipage} & \begin{minipage}[b]{\linewidth}\raggedright
SUMMARY
\end{minipage} & \begin{minipage}[b]{\linewidth}\raggedright
SCHEDULE
\end{minipage} \\
\midrule\noalign{}
\endhead
\bottomrule\noalign{}
\endlastfoot
\end{longtable}

\subsubsection{Description}\label{description-595}

PVTGW defines the gas PVT properties for dry gas {[}Natural gas that
occurs in the absence of condensate or liquid hydrocarbons, or gas that
had condensable hydrocarbons removed, is called dry gas. It is primarily
methane with some intermediates. The hydrocarbon mixture is solely gas
in the reservoir and there is no liquid (condensate surface liquid)
formed either in the reservoir or at surface. The term dry indicates
that the gas does not contain heavier hydrocarbons to form liquids at
the surface conditions. Dry gas typically has GOR's greater than 100,000
scf/stb or 18,000 Sm3/m3.{]} with vaporized water. This keyword should
be used when the VAPWAT keyword has be declared in the RUNSPEC section
indicating that that vaporized water is present in the dry gas phase.
The keyword may be used for gas-water and oil-water-gas input decks that
contain the dry gas and vaporized water phases.

\begin{tcolorbox}[enhanced jigsaw, title=\textcolor{quarto-callout-note-color}{\faInfo}\hspace{0.5em}{Note}, coltitle=black, breakable, toprule=.15mm, colbacktitle=quarto-callout-note-color!10!white, leftrule=.75mm, opacitybacktitle=0.6, colback=white, bottomtitle=1mm, toptitle=1mm, titlerule=0mm, left=2mm, colframe=quarto-callout-note-color-frame, arc=.35mm, rightrule=.15mm, bottomrule=.15mm, opacityback=0]

This is an OPM Flow specific keyword for the simulator's Water
Vaporization Model that is activated by declaring that vaporized water
is present in the run.

\end{tcolorbox}

\begin{longtable}[]{@{}
  >{\raggedright\arraybackslash}p{(\linewidth - 8\tabcolsep) * \real{0.2000}}
  >{\raggedright\arraybackslash}p{(\linewidth - 8\tabcolsep) * \real{0.2000}}
  >{\raggedright\arraybackslash}p{(\linewidth - 8\tabcolsep) * \real{0.2000}}
  >{\raggedright\arraybackslash}p{(\linewidth - 8\tabcolsep) * \real{0.2000}}
  >{\raggedright\arraybackslash}p{(\linewidth - 8\tabcolsep) * \real{0.2000}}@{}}
\toprule\noalign{}
\begin{minipage}[b]{\linewidth}\raggedright
No.
\end{minipage} & \begin{minipage}[b]{\linewidth}\raggedright
Name
\end{minipage} & \begin{minipage}[b]{\linewidth}\raggedright
Description
\end{minipage} & \begin{minipage}[b]{\linewidth}\raggedright
Default
\end{minipage} & \begin{minipage}[b]{\linewidth}\raggedright
\end{minipage} \\
\midrule\noalign{}
\endhead
\bottomrule\noalign{}
\endlastfoot
Field & Metric & Laboratory & & \\
1 & PRESS & A real monotonically increasing down the column vector that
defines the gas phase pressure, associated with the corresponding
saturated water-gas ratio (``WGR'') or Rw, the gas formation volume
factor, and the gas viscosity for the stated saturated RWS. For a given
PRESS the variability of the gas formation volume factor and the gas
viscosity with respect to the under-saturated Rw is optionally included
as a sub table under RWU, FVFU and VISU columns, that is it is not
necessary to repeat PRESS for each sub table entry. However, each sub
table must be terminated by a ``/''. The under saturated Rw entries are
optional, except for perhaps the last PRESS entry to define the PVT
properties above the initial saturation pressure. & None & \\
psia & barsa & atma & & \\
2 & RWS & RWU & A columnar vector of real positive numbers for both the
saturated (RWS) and under saturated (RWU) Rw sub table entries. The RWS
entry on the main table is the saturated WGR at the pressure indicated
by PRESS and may be increasing or decreasing in value as PRESS varies.
Subsequent under-saturated Rw for a sub table at the given PRESS, as
defined by RWU, are monotonically decreasing for entries in a given sub
table. & None \\
stb/Mscf & sm3/sm3 & scc/scc & & \\
3 & FVFS & FVFU & A columnar vector of real decreasing down the column
values that defines the corresponding gas phase formation volume factor
for a given pressure (PRESS) and for a given Rw (either RWS or RWU). &
None \\
rb/Mscf & rm3/sm3 & rcc/scc & & \\
4 & VISS & VISU & VISS a columnar vector of real increasing down the
column values that defines the corresponding gas phase viscosity for a
given pressure (PRESS) and for a given RWS. VISU a columnar vector of
real decreasing from VISS down the column values that defines the
corresponding gas phase viscosity for a given pressure (PRESS) and for a
given RWU. & None \\
cP & cP & cP & & \\
Notes: & & & & \\
\end{longtable}

\emph{Table 8.118: PVTGW Keyword Description}

Note that provided the first table has been entered, subsequent tables
may be defaulted, in this case the prior table is copied to the current
table.

\begin{tcolorbox}[enhanced jigsaw, title=\textcolor{quarto-callout-note-color}{\faInfo}\hspace{0.5em}{Note}, coltitle=black, breakable, toprule=.15mm, colbacktitle=quarto-callout-note-color!10!white, leftrule=.75mm, opacitybacktitle=0.6, colback=white, bottomtitle=1mm, toptitle=1mm, titlerule=0mm, left=2mm, colframe=quarto-callout-note-color-frame, arc=.35mm, rightrule=.15mm, bottomrule=.15mm, opacityback=0]

If both the VAPWAT and VAPOIL keywords have been declared in the RUNSPEC
section indicating that both vaporized water and vaporized oil are
present in the wet gas, then the PVTGW keyword should be used along with
the PVTG keyword in the PROPS section to fully define the wet gas PVT
properties. The PVTGW keyword should be used to define the gas
properties as a function of pressure and water-gas ratio (RVW), assuming
oil-saturated gas. The PVTG keyword should be used to define the gas
properties as a function of pressure and oil-gas ratio (RV), assuming
water-saturated gas. Alternatively, the PVTGWO keyword in the PROPS
section may be used instead of the PVTGW and PVTG keywords to fully
define the wet gas PVT properties.

\end{tcolorbox}

See also the PVTG keyword in the PROPS section that defines the wet gas
PVT for when vaporized oil is present in the gas phase. Alternatively,
the PVTGWO keyword in the PROPS section may be utilized instead of PVTG
and PVTGW to fully define the wet gas PVT properties, for when both
vaporized oil and water are present in the gas phase.

\subsubsection{Example}\label{example-331}

\begin{verbatim}
--
--       GAS PVT TABLE FOR DRY GAS WITH VAPORIZED WATER (OPM FLOW KEYWORD)
--
PVTGW
--       PRES       RW           BG           VISC
--       PSIA     SM^3/SM^3     RM^3/SM^3       CPOISE
--      ------    ---------    -------       ------
          300     0.000479      0.042340      0.01344
                  0             0.042310      0.01389      /
          600     0.000469      0.020460      0.01420
                  0             0.020430      0.01450      /
          900     0.000403      0.013280      0.01526
                  0             0.013250      0.01532      /
         1200     0.000354      0.009770      0.01660
                  0             0.009730      0.01634      /
         1500     0.000272      0.007730      0.01818
                  0             0.007690      0.01752      /
         1800     0.000225      0.006426      0.01994
                  0             0.006405      0.01883      /
         2100     0.000191      0.005541      0.02181
                  0             0.005553      0.02021      /
         2400     0.000163      0.004919      0.02370
                  0             0.004952      0.02163      /
                                                           / TABLE NO. 1
--       PRES       RW           BG           VISC
--       PSIA     SM^3/SM^3     RM^3/SM^3       CPOISE
--      ------    ---------    -------       ------
          300     0.000479      0.042340      0.01344      /
          600     0.000469      0.020460      0.01420      /
          900     0.000403      0.013280      0.01526      /
         1200     0.000354      0.009770      0.01660      /
         1500     0.000272      0.007730      0.01818      /
         1800     0.000225      0.006426      0.01994      /
         2100     0.000191      0.005541      0.02181      /
         2400     0.000163      0.004919      0.02370      /
                                                           / TABLE NO. 2
\end{verbatim}

The above example defines two dry gas PVT tables assuming NTPVT equals
two and NPPVT is greater than or equal to eight on the TABDIMS keyword
in the RUNSPEC section.

Notice that there is no terminating ``/'' for this keyword only for a
table and a sub table.

\subsection{PVTGWO -- Gas PVT Properties for Wet Gas with Vaporized
Water and
Oil}\label{pvtgwo-gas-pvt-properties-for-wet-gas-with-vaporized-water-and-oil}

\begin{longtable}[]{@{}
  >{\raggedright\arraybackslash}p{(\linewidth - 14\tabcolsep) * \real{0.1250}}
  >{\raggedright\arraybackslash}p{(\linewidth - 14\tabcolsep) * \real{0.1250}}
  >{\raggedright\arraybackslash}p{(\linewidth - 14\tabcolsep) * \real{0.1250}}
  >{\raggedright\arraybackslash}p{(\linewidth - 14\tabcolsep) * \real{0.1250}}
  >{\raggedright\arraybackslash}p{(\linewidth - 14\tabcolsep) * \real{0.1250}}
  >{\raggedright\arraybackslash}p{(\linewidth - 14\tabcolsep) * \real{0.1250}}
  >{\raggedright\arraybackslash}p{(\linewidth - 14\tabcolsep) * \real{0.1250}}
  >{\raggedright\arraybackslash}p{(\linewidth - 14\tabcolsep) * \real{0.1250}}@{}}
\toprule\noalign{}
\begin{minipage}[b]{\linewidth}\raggedright
RUNSPEC
\end{minipage} & \begin{minipage}[b]{\linewidth}\raggedright
GRID
\end{minipage} & \begin{minipage}[b]{\linewidth}\raggedright
EDIT
\end{minipage} & \begin{minipage}[b]{\linewidth}\raggedright
PROPS
\end{minipage} & \begin{minipage}[b]{\linewidth}\raggedright
REGIONS
\end{minipage} & \begin{minipage}[b]{\linewidth}\raggedright
SOLUTION
\end{minipage} & \begin{minipage}[b]{\linewidth}\raggedright
SUMMARY
\end{minipage} & \begin{minipage}[b]{\linewidth}\raggedright
SCHEDULE
\end{minipage} \\
\midrule\noalign{}
\endhead
\bottomrule\noalign{}
\endlastfoot
\end{longtable}

\subsubsection{Description}\label{description-596}

PVTGWO defines the gas PVT properties for wet gas {[}Natural gas that
contains significant heavy hydrocarbons such as propane, butane and
other liquid hydrocarbons is known as wet gas or rich gas. The general
rule of thumb is if the gas contains less methane (typically less than
85\% methane) and more ethane, and other more complex hydrocarbons, it
is labeled as wet gas. Wet gas normally has GOR's less than 100,000
scf/stb or 18,000 Sm3/m3, with the condensate having a gravity greater
than 50 oAPI.{]} with vaporized water and oil. This keyword should be
used when the VAPOIL and VAPWAT keywords have been declared in the
RUNSPEC section indicating that vaporized oil and water are present in
the wet gas phase. The keyword may be used for oil-water-gas input decks
that contain the wet gas with vaporized oil and water phases.

This keyword is not supported by OPM Flow but would change the results
if supported so the simulation will be stopped.

\begin{tcolorbox}[enhanced jigsaw, title=\textcolor{quarto-callout-note-color}{\faInfo}\hspace{0.5em}{Note}, coltitle=black, breakable, toprule=.15mm, colbacktitle=quarto-callout-note-color!10!white, leftrule=.75mm, opacitybacktitle=0.6, colback=white, bottomtitle=1mm, toptitle=1mm, titlerule=0mm, left=2mm, colframe=quarto-callout-note-color-frame, arc=.35mm, rightrule=.15mm, bottomrule=.15mm, opacityback=0]

This is an OPM Flow specific keyword for the simulator's Water
Vaporization Model that is activated by declaring that vaporized water
is present in the run.

\end{tcolorbox}

\begin{longtable}[]{@{}
  >{\raggedright\arraybackslash}p{(\linewidth - 8\tabcolsep) * \real{0.2000}}
  >{\raggedright\arraybackslash}p{(\linewidth - 8\tabcolsep) * \real{0.2000}}
  >{\raggedright\arraybackslash}p{(\linewidth - 8\tabcolsep) * \real{0.2000}}
  >{\raggedright\arraybackslash}p{(\linewidth - 8\tabcolsep) * \real{0.2000}}
  >{\raggedright\arraybackslash}p{(\linewidth - 8\tabcolsep) * \real{0.2000}}@{}}
\toprule\noalign{}
\begin{minipage}[b]{\linewidth}\raggedright
No.
\end{minipage} & \begin{minipage}[b]{\linewidth}\raggedright
Name
\end{minipage} & \begin{minipage}[b]{\linewidth}\raggedright
Description
\end{minipage} & \begin{minipage}[b]{\linewidth}\raggedright
Default
\end{minipage} & \begin{minipage}[b]{\linewidth}\raggedright
\end{minipage} \\
\midrule\noalign{}
\endhead
\bottomrule\noalign{}
\endlastfoot
Field & Metric & Laboratory & & \\
1 & PRESS & A real monotonically increasing down the column vector that
defines the gas phase pressure, associated with the saturated water-gas
ratio (``WGR'') or Rw, the saturated condensate-gas ratio (``CGR'') or
Rv, the gas formation volume factor, and the gas viscosity for the
corresponding pressure for the stated saturated RWS. For a given PRESS
the variability of the gas formation volume factor and the gas viscosity
with respect to the under-saturated Rw and Rv is optionally included as
a sub table under RWU, RVU, FVFU and VISU columns, that is it is not
necessary to repeat PRESS for each sub table entry. However, each sub
table must be terminated by a ``/''. The under saturated Rw and Rv
entries are optional, except for perhaps the last PRESS entry to define
the PVT properties above the initial saturation pressure. & None & \\
psia & barsa & atma & & \\
2 & RWS & RWU & A columnar vector of real positive numbers for both the
saturated (RWS) and under saturated (RWU) Rw sub table entries. The RWS
entry on the main table is the saturated WGR at the pressure indicated
by PRESS and may be increasing or decreasing in value as PRESS varies.
Subsequent under-saturated Rw for a sub table at the given PRESS, as
defined by RWU, are monotonically decreasing for entries in a given sub
table. & None \\
stb/Mscf & sm3/sm3 & rcc/scc & & \\
3 & RVS & RVU & A columnar vector of real positive numbers for both the
saturated (RVS) and under saturated (RVU) Rv sub table entries. The RVS
entry on the main table is the saturated CGR at the pressure indicated
by PRESS and may be increasing or decreasing in value as PRESS varies.
Subsequent under-saturated Rv for a sub table at the given PRESS, as
defined by RVU, are monotonically decreasing for entries in a given sub
table. & None \\
stb/Mscf & sm3/sm3 & rcc/scc & & \\
4 & FVFS & FVFU & A columnar vector of real decreasing down the column
values that defines the corresponding gas phase formation volume factor
for a given pressure (PRESS) and for a given Rw (either RWS or RWU) and
Rv (either RVS or RVU). & None \\
rb/Mscf & rm3/sm3 & rcc/scc & & \\
5 & VISS & VISU & VISS a columnar vector of real increasing down the
column values that defines the corresponding gas phase viscosity for a
given pressure (PRESS) and for a given RWS and RVS. VISU a columnar
vector of real decreasing from VISS down the column values that defines
the corresponding gas phase viscosity for a given pressure (PRESS) and
for a given RWU and RVU. & None \\
cP & cP & cP & & \\
Notes: & & & & \\
\end{longtable}

\emph{Table 8.3.232.1: PVTGWO Keyword Description}

See also the PVTG keyword in the PROPS section that defines the wet gas
PVT properties when vaporized oil is present in the gas phase, and the
PVTGW keyword in the PROPS section that defines the dry gas PVT
properties when vaporized water is present in the gas phase. As an
alternative to using the PVTGWO keyword, the PVTGW and PVTG keywords may
be combined to fully define the wet gas PVT properties if both vaporized
water and vaporized oil are present in the gas phase.

\subsubsection{Example}\label{example-332}

\begin{verbatim}
--
--       GAS PVT TABLE FOR WET GAS WITH VAPORIZED WATER & OIL (OPM FLOW KEYWORD)
--
PVTGWO
--       PRES     RW         RV         BG        VISC
--       PSIA   STB/MSCF   STB/MSCF   RB/MSCF    CPOISE
--      ------  --------   ---------  -------    ------
          300   0.000479   0.000132   0.042340   0.01344
                0          0          0.042310   0.01389   /
          600   0.000469   0.000124   0.020460   0.01420
                0          0          0.020430   0.01450   /
          900   0.000403   0.000126   0.013280   0.01526
                0          0          0.013250   0.01532   /
         1200   0.000354   0.000135   0.009770   0.01660
                0          0          0.009730   0.01634   /
         1500   0.000272   0.000149   0.007730   0.01818
                0          0          0.007690   0.01752   /
         1800   0.000225   0.000163   0.006426   0.01994
                0          0          0.006405   0.01883   /
         2100   0.000191   0.000191   0.005541   0.02181
                0          0          0.005553   0.02021   /
         2400   0.000163   0.000225   0.004919   0.02370
                0          0          0.004952   0.02163   /
                                                           / TABLE NO. 1
--       PRES     RW         RV        BG        VISC
--       PSIA   STB/MSCF   STB/MSCF   RB/MSCF    CPOISE
--      ------  --------   ---------  -------    ------
          300   0.000479   0.000132   0.042340   0.01344   /
          600   0.000469   0.000124   0.020460   0.01420   /
          900   0.000403   0.000126   0.013280   0.01526   /
         1200   0.000354   0.000135   0.009770   0.01660   /
         1500   0.000272   0.000149   0.007730   0.01818   /
         1800   0.000225   0.000163   0.006426   0.01994   /
         2100   0.000191   0.000191   0.005541   0.02181   /
         2400   0.000163   0.000225   0.004919   0.02370   /
                                                           / TABLE NO. 2
\end{verbatim}

The above example defines two wet PVT tables assuming NTPVT equals two
and NPPVT is greater than or equal to eight on the TABDIMS keyword in
the RUNSPEC section.

Notice that there is no terminating ``/'' for this keyword only for a
table and a sub table.

\subsection{PVTO -- Oil PVT Properties for Live
Oil}\label{pvto-oil-pvt-properties-for-live-oil}

\begin{longtable}[]{@{}
  >{\raggedright\arraybackslash}p{(\linewidth - 14\tabcolsep) * \real{0.1250}}
  >{\raggedright\arraybackslash}p{(\linewidth - 14\tabcolsep) * \real{0.1250}}
  >{\raggedright\arraybackslash}p{(\linewidth - 14\tabcolsep) * \real{0.1250}}
  >{\raggedright\arraybackslash}p{(\linewidth - 14\tabcolsep) * \real{0.1250}}
  >{\raggedright\arraybackslash}p{(\linewidth - 14\tabcolsep) * \real{0.1250}}
  >{\raggedright\arraybackslash}p{(\linewidth - 14\tabcolsep) * \real{0.1250}}
  >{\raggedright\arraybackslash}p{(\linewidth - 14\tabcolsep) * \real{0.1250}}
  >{\raggedright\arraybackslash}p{(\linewidth - 14\tabcolsep) * \real{0.1250}}@{}}
\toprule\noalign{}
\begin{minipage}[b]{\linewidth}\raggedright
RUNSPEC
\end{minipage} & \begin{minipage}[b]{\linewidth}\raggedright
GRID
\end{minipage} & \begin{minipage}[b]{\linewidth}\raggedright
EDIT
\end{minipage} & \begin{minipage}[b]{\linewidth}\raggedright
PROPS
\end{minipage} & \begin{minipage}[b]{\linewidth}\raggedright
REGIONS
\end{minipage} & \begin{minipage}[b]{\linewidth}\raggedright
SOLUTION
\end{minipage} & \begin{minipage}[b]{\linewidth}\raggedright
SUMMARY
\end{minipage} & \begin{minipage}[b]{\linewidth}\raggedright
SCHEDULE
\end{minipage} \\
\midrule\noalign{}
\endhead
\bottomrule\noalign{}
\endlastfoot
\end{longtable}

\subsubsection{Description}\label{description-597}

PVTO defines the oil PVT properties for live oil {[}``Live'' oil is oil
that contains gas in solution, which is normally the case for most
conventional oil reservoirs. However, for oil reservoirs classified as
heavy oil reservoirs, the in situ dissolved gas may be negligible and
oil would then be classified as gas-free oil which is commonly referred
to as ``dead'' oil.{]} and the keyword should only be used if the there
is both oil and gas phases in the model. This keyword should be used
when the DISGAS keyword has be declared in the RUNSPEC section
indicating that dissolved gas (more commonly referred to as solution
gas) is present in the oil. The keyword may be used for oil-water and
oil-water-gas input decks.

\begin{longtable}[]{@{}
  >{\raggedright\arraybackslash}p{(\linewidth - 8\tabcolsep) * \real{0.2000}}
  >{\raggedright\arraybackslash}p{(\linewidth - 8\tabcolsep) * \real{0.2000}}
  >{\raggedright\arraybackslash}p{(\linewidth - 8\tabcolsep) * \real{0.2000}}
  >{\raggedright\arraybackslash}p{(\linewidth - 8\tabcolsep) * \real{0.2000}}
  >{\raggedright\arraybackslash}p{(\linewidth - 8\tabcolsep) * \real{0.2000}}@{}}
\toprule\noalign{}
\begin{minipage}[b]{\linewidth}\raggedright
No.
\end{minipage} & \begin{minipage}[b]{\linewidth}\raggedright
Name
\end{minipage} & \begin{minipage}[b]{\linewidth}\raggedright
Description
\end{minipage} & \begin{minipage}[b]{\linewidth}\raggedright
Default
\end{minipage} & \begin{minipage}[b]{\linewidth}\raggedright
\end{minipage} \\
\midrule\noalign{}
\endhead
\bottomrule\noalign{}
\endlastfoot
Field & Metric & Laboratory & & \\
1 & RS & A real monotonically increasing down the column values that
defines the saturated gas-oil ratio (``GOR'') or Rs, that defines the
oil formation volume factor and the oil viscosity for the tabulated
corresponding pressure for stated saturated RS. For a given RS the
variability of the oil formation volume factor and the oil viscosity
with respect to the saturated RS and pressure is optionally included as
a sub table under PRSU, FVFU and VISU columns, that is it is not
necessary to repeat RS for each sub table entry. However, each sub table
must be terminated by a ``/''. The under-saturated PRSU entries are
optional, except for perhaps the last RS entry to define the PVT
properties above the initial saturation pressure. If there are no
following under-saturated PRSU entries then the RS entry row should be
terminated by a ``/'', if there are under-saturated PRSU entries then
the last PRSU entry row should be terminated by a ``/''. & None & \\
Mscf/stb & sm3/sm3 & scc/scc & & \\
2 & PRSS & PRSU & PRSS is a real columnar vector of real monotonically
increasing down the column values that defines the oil phase saturation
pressure (bubble-point pressure), that defines the oil formation volume
factor and the oil viscosity for the corresponding PRSS pressure for a
given saturated RS. PRSU is a real columnar vector of real monotonically
increasing down the column values that defines the oil phase
under-saturated pressure that defines the oil formation volume factor
and the oil viscosity for the corresponding PRSU pressure for a given
saturated RS. Note that PRSU should be greater than PRSS. & None \\
psia & barsa & atma & & \\
3 & FVFS & FVFU & FVFS is a columnar vector of real increasing down the
column values that defines the corresponding oil phase saturated
formation volume factor for a given pressure (PRSS) and for a given RS.
FVFU is a columnar vector of real decreasing down the column values that
defines the corresponding oil phase under-saturated formation volume
factor for a given pressure (PRSU) and for a given RS. & None \\
rb/stb & rm3/sm3 & rcc/scc & & \\
4 & VISS & VISU & VISS a columnar vector of real increasing down the
column values that defines the corresponding oil phase saturated
viscosity for a given pressure (PRSS) and for a given RS. If this is the
only entry for a given RS and PRSS then the record should be terminate
by a ``/''. VISU a columnar vector of real decreasing from VISS down the
column values that defines the corresponding oil phase under-saturated
viscosity for a given pressure (PRSU) and for a given RS. If this is the
only entry for a given RS and PRSU then the record should be terminate
by a ``/''. & None \\
cP & cP & cP & & \\
Notes: & & & & \\
\end{longtable}

\emph{Table 8.120: PVTO Keyword Description}

Note that provided the first table has been entered, subsequent tables
may be defaulted, in this case the prior table is copied to the current
table. See the second example for an illustration on how to use this
feature.

\subsubsection{Example}\label{example-333}

The first example defines live oil PVT tables assuming NTPVT equals two,
NPPVT is greater than or equal to two, and NRPVT is greater than or
equal to 18 on the TABDIMS keyword in the RUNSPEC section.

\begin{verbatim}
--
--       OIL PVT TABLE FOR LIVE OIL
--
PVTO
--       RS        PSAT       BO        VISC
--       MSCF/STB  PSIA       RB/STB    CPOISE
--       --------  --------   -------   ------
          0.0010      14.7    1.05340   1.7230  /
          0.0890     500.0    1.08890   1.1670  /
          0.2060    1000.0    1.13850   0.8570  /
          0.3360    1500.0    1.19640   0.6840  /
          0.4750    2000.0    1.26110   0.5750  /
          0.5480    2250.0    1.29570   0.5340  /
          0.6220    2500.0    1.33160   0.5000  /
          0.6980    2750.0    1.36890   0.4700  /
          0.7750    3000.0    1.40740   0.4450  /
          0.8530    3250.0    1.44710   0.4220  /
          0.9330    3500.0    1.48790   0.4020  /
          1.0140    3750.0    1.52980   0.3840  /
          1.0960    4000.0    1.57280   0.3680  /
          1.1800    4258.0    1.61760   0.3530  /
          1.2630    4500.0    1.66190   0.3400  /
          1.3480    4750.0    1.70780   0.3280  /
          1.4340    5000.0    1.75480   0.3170  /
          1.6060    5500.0    1.85020   0.2980
                    6242.0    1.83040   0.3186  /
                                                / TABLE NO. 1
--
--
--       RS        PSAT       BO        VISC
--       MSCF/STB  PSIA       RB/STB    CPOISE
--       --------  --------   -------   ------
          0.0010      14.7    1.05340   1.7230  /
          0.0390     250.0    1.06830   1.4220  /
          0.0890     500.0    1.08890   1.1670  /
          0.1460     750.0    1.11250   0.9850  /
          0.2060    1000.0    1.13850   0.8570  /
          0.2700    1250.0    1.16660   0.7590  /
          0.3360    1500.0    1.19640   0.6840  /
          0.4050    1750.0    1.22800   0.6240  /
          0.4750    2000.0    1.26110   0.5750  /
          0.5480    2250.0    1.29570   0.5340  /
          0.6220    2500.0    1.33160   0.5000  /
          0.6980    2750.0    1.36890   0.4700  /
          0.7750    3000.0    1.40740   0.4450  /
          0.8530    3250.0    1.44710   0.4220  /
          0.9330    3500.0    1.48790   0.4020  /
          1.0140    3750.0    1.52980   0.3840  /
          1.0960    4000.0    1.57280   0.3680  /
          1.1800    4258.0    1.61760   0.3530  /
          1.2630    4500.0    1.66190   0.3400  /
          1.3480    4750.0    1.70780   0.3280  /
          1.4340    5000.0    1.75480   0.3170  /
          1.6060    5500.0    1.85020   0.2980
                    6242.0    1.83040   0.3186  /
                                                / TABLE NO. 2
\end{verbatim}

Notice that there must be at least two entries for the last Rs value to
enable the simulator to interpolate over the undersaturated pressure
region.

The second example defines live oil PVT tables assuming NTPVT equals
four, NPPVT is greater than or equal to two, and NRPVT is greater than
or equal to 13 on the TABDIMS keyword in the RUNSPEC section. Here,
tables two to four all default to table number one.

\begin{verbatim}
--
--       OIL PVT TABLE FOR LIVE OIL
--
PVTO
--       RS        PSAT       BO        VISC
--       MSCF/STB  PSIA       RB/STB    CPOISE
--       --------  --------   -------   ------
          0.0010      14.7    1.05340   1.7230  /
          0.0890     500.0    1.08890   1.1670  /
          0.2060    1000.0    1.13850   0.8570  /
          0.3360    1500.0    1.19640   0.6840  /
          0.4750    2000.0    1.26110   0.5750  /
          0.5480    2250.0    1.29570   0.5340  /
          0.6220    2500.0    1.33160   0.5000  /
          0.7750    3000.0    1.40740   0.4450  /
          0.9330    3500.0    1.48790   0.4020  /
          1.0960    4000.0    1.57280   0.3680  /
          1.2630    4500.0    1.66190   0.3400  /
          1.4340    5000.0    1.75480   0.3170  /
          1.6060    5500.0    1.85020   0.2980
                    6242.0    1.83040   0.3186  /
                                                / TABLE NO. 1
                                                / TABLE NO. 2
                                                / TABLE NO. 3
                                                / TABLE NO. 4
\end{verbatim}

Again, note that there is no terminating ``/'' for this keyword only for
a table and a sub table.

\subsection{PVTSOL -- Oil PVT Properties for Live Oil versus CO2 Mass
Fraction}\label{pvtsol-oil-pvt-properties-for-live-oil-versus-co2-mass-fraction}

\begin{longtable}[]{@{}
  >{\raggedright\arraybackslash}p{(\linewidth - 14\tabcolsep) * \real{0.1250}}
  >{\raggedright\arraybackslash}p{(\linewidth - 14\tabcolsep) * \real{0.1250}}
  >{\raggedright\arraybackslash}p{(\linewidth - 14\tabcolsep) * \real{0.1250}}
  >{\raggedright\arraybackslash}p{(\linewidth - 14\tabcolsep) * \real{0.1250}}
  >{\raggedright\arraybackslash}p{(\linewidth - 14\tabcolsep) * \real{0.1250}}
  >{\raggedright\arraybackslash}p{(\linewidth - 14\tabcolsep) * \real{0.1250}}
  >{\raggedright\arraybackslash}p{(\linewidth - 14\tabcolsep) * \real{0.1250}}
  >{\raggedright\arraybackslash}p{(\linewidth - 14\tabcolsep) * \real{0.1250}}@{}}
\toprule\noalign{}
\begin{minipage}[b]{\linewidth}\raggedright
RUNSPEC
\end{minipage} & \begin{minipage}[b]{\linewidth}\raggedright
GRID
\end{minipage} & \begin{minipage}[b]{\linewidth}\raggedright
EDIT
\end{minipage} & \begin{minipage}[b]{\linewidth}\raggedright
PROPS
\end{minipage} & \begin{minipage}[b]{\linewidth}\raggedright
REGIONS
\end{minipage} & \begin{minipage}[b]{\linewidth}\raggedright
SOLUTION
\end{minipage} & \begin{minipage}[b]{\linewidth}\raggedright
SUMMARY
\end{minipage} & \begin{minipage}[b]{\linewidth}\raggedright
SCHEDULE
\end{minipage} \\
\midrule\noalign{}
\endhead
\bottomrule\noalign{}
\endlastfoot
\end{longtable}

\subsubsection{Description}\label{description-598}

PVTSOL defines the live oil PVT properties as a function of CO2 mass
fraction. The keyword automatically invokes the simulator's CO2 Dynamic
EOR Model {[}T. H. Sandve, O. Sævareid and I. Aavatsmark: ``Improved
Extended Blackoil Formulation -- for CO2 EOR Simulations.'' in ECMOR
XVII -- The 17th European Conference on the -- Mathematics of Oil
Recovery, September 2020.{]}, that uses a fourth component to model the
injected CO2., for use in evaluating CO2 Enhanced Oil Recovery (``EOR'')
projects. Normally CO2 EOR projects are evaluated via compositional
simulators to account for the mass transfer of the various components
and phases. Unfortunately, compositional models are computationally
expensive compared to the black-oil approach, which for field studies is
challenging, especially if an ensemble approach is being used to capture
the uncertainties. Previous extended black-oil formulations often poorly
represent the PVT properties of the oil-CO2 mixtures, resulting in poor
agreement with the compositional formulation.

To overcome the limitations of the standard four component black-oil
formulation, OPM Flow uses an improved extended black-oil formulation,
the CO2 Dynamic EOR Model, with the black-oil properties dependent on
the fraction of CO2 in the cell, as described by the PVTSOL keyword.
This approach models the oil-CO2 mixture more accurately and thus give
results closer to the compositional simulator. Data for the keyword
should normally be derived from laboratory or numerical slim-tube
experiments based on one-dimensional compositional Equation Of State
(``EOS'') simulations.

Note, if this keyword is absent from the input deck then the CO2
Standard EOR Model is used instead.

\begin{longtable}[]{@{}
  >{\raggedright\arraybackslash}p{(\linewidth - 8\tabcolsep) * \real{0.2000}}
  >{\raggedright\arraybackslash}p{(\linewidth - 8\tabcolsep) * \real{0.2000}}
  >{\raggedright\arraybackslash}p{(\linewidth - 8\tabcolsep) * \real{0.2000}}
  >{\raggedright\arraybackslash}p{(\linewidth - 8\tabcolsep) * \real{0.2000}}
  >{\raggedright\arraybackslash}p{(\linewidth - 8\tabcolsep) * \real{0.2000}}@{}}
\toprule\noalign{}
\begin{minipage}[b]{\linewidth}\raggedright
No.
\end{minipage} & \begin{minipage}[b]{\linewidth}\raggedright
Name
\end{minipage} & \begin{minipage}[b]{\linewidth}\raggedright
Description
\end{minipage} & \begin{minipage}[b]{\linewidth}\raggedright
Default
\end{minipage} & \begin{minipage}[b]{\linewidth}\raggedright
\end{minipage} \\
\midrule\noalign{}
\endhead
\bottomrule\noalign{}
\endlastfoot
Field & Metric & Laboratory & & \\
1 & CO2 & A real monotonically increasing down the column values that
stipulates the CO2 mass fraction, that defines the oil and gas
properties, formation volume factor, viscosity etc., for the tabulated
corresponding pressure for the stated CO2 mass fraction. CO2 should be
greater than or equal to zero and less than or equal to one. Note it is
not necessary to repeat the value of CO2 for the pressure column.
However, for a given CO2 mass fraction, the last pressure entry (PRESS)
of the sub table should be terminated by a ``/''. & None & \\
dimensionless & dimensionless & dimensionless & & \\
2 & & PRESS & PRESS is a real columnar vector of real monotonically
increasing down the column values that defines the oil phase saturation
pressure (bubble-point pressure), that defines the oil formation volume
factor and the oil viscosity for the corresponding PRESS pressure for a
given saturated value of CO2. & None \\
psia & barsa & atma & & \\
3 & & OFVF & OFVF is a columnar vector of real increasing down the
column values that defines the corresponding oil phase saturated
formation volume factor for a given pressure (PRESS) and for a given
saturated value of CO2. & None \\
rb/stb & rm3/sm3 & rcc/scc & & \\
4 & & GFVF & GFVF is a columnar vector of real decreasing down the
column values that defines the corresponding gas phase saturated
formation volume factor for a given pressure (PRESS) and for a given
saturated value of CO2. & None \\
rb/Mscf & rm3/sm3 & rcc/scc & & \\
5 & & RS & RS is a real monotonically increasing down the column values
that defines the saturated gas-oil ratio (``GOR'') or Rs, for the given
value of PRESS and for a given saturated value of CO2. & None \\
Mscf/stb & sm3/sm3 & scc/scc & & \\
6 & & RV & RV is a real monotonically increasing down the column values
that defines the saturated condensate-gas ratio (``CGR'') or Rv, for the
given value of PRESS and for a given saturated value of CO2. & None \\
stb/Mscf & sm3/sm3 & scc/scc & & \\
7 & & XVOL & XVOL is a real positive value greater than or equal to zero
and less than or equal to one, that stipulates the volumetric fraction
of CO2 in the oil phase, that is:
\(\mathit{XVOL} = \frac{{\mathit{Volume}}_{\mathit{oil}}({\mathit{CO}}_{2})}{{\mathit{Volume}}_{\mathit{oil}}({\mathit{CO}}_{2}) + {\mathit{Volume}}_{\mathit{oil}}(\mathit{Gas}) + {\mathit{Volume}}_{\mathit{oil}}(\mathit{Oil})}\)
where: Volume oil (CO2) is the surface volume of CO2 in the oil phase,
Volume oil (Gas) is the surface volume of gas in the oil phase, and
Volume oil (Oil) is the surface volume of oil in the oil phase. &
None \\
dimensionless & dimensionless & dimensionless & & \\
8 & & YVOL & YVOL is a real positive value greater than or equal to zero
and less than or equal to one, that defines the volumetric fraction of
CO2 in the gas phase, that is:
\(\mathit{YVOL} = \frac{{\mathit{Volume}}_{\mathit{gas}}({\mathit{CO}}_{2})}{{\mathit{Volume}}_{\mathit{gas}}({\mathit{CO}}_{2}) + {\mathit{Volume}}_{\mathit{gas}}(\mathit{Gas}) + {\mathit{Volume}}_{\mathit{gas}}(\mathit{Oil})}\)
where: Volume gas (CO2) is the surface volume of CO2 in the gas phase,
Volume gas (Gas) is the surface volume of gas in the gas phase, and
Volume gas (Oil) is the surface volume of oil in the gas phase. &
None \\
dimensionless & dimensionless & dimensionless & & \\
9 & & OVISC & OVISC is a columnar vector of real decreasing down the
column values that defines the corresponding oil phase saturated
viscosity for a given pressure (PRESS) and for a given saturated value
of CO2. & None \\
cP & cP & cP & & \\
10 & & GVISC & GVISC is a columnar vector of real increasing down the
column values that defines the corresponding gas phase saturated
viscosity for a given pressure (PRESS) and for a given saturated value
of CO2. & None \\
cP & cP & cP & & \\
Notes: & & & & \\
\end{longtable}

\emph{Table 8.121: PVTSOL Keyword Description}

See also the PVDS keyword in the PROPS section that can also be used to
model CO2 injection, using the SOLVENT model (CO2 Standard EOR), this is
the conventional approach, and does not take into account black-oil
properties being dependent on the fraction of CO2.

\begin{tcolorbox}[enhanced jigsaw, title=\textcolor{quarto-callout-note-color}{\faInfo}\hspace{0.5em}{Note}, coltitle=black, breakable, toprule=.15mm, colbacktitle=quarto-callout-note-color!10!white, leftrule=.75mm, opacitybacktitle=0.6, colback=white, bottomtitle=1mm, toptitle=1mm, titlerule=0mm, left=2mm, colframe=quarto-callout-note-color-frame, arc=.35mm, rightrule=.15mm, bottomrule=.15mm, opacityback=0]

In the CO2 Dynamic EOR Model the oil properties within a cell are
dependent on the CO2 saturation in the cell. However, the model is not
restricted to just CO2 utilization, as long as the appropriate
oil-injection-gas dependent properties are entered in the model via the
PVTSOL keyword. Thus, N2 or other miscible gases may be used instead of
CO2, provided due care is taken in generating the required data for the
PVTSOL keyword. Another example would be a gas-condensate re-cycling
project where the produced gas is stripped of the liquids and then
re-injected back into the reservoir in order to maintain the reservoir
pressure. Here, the CO2 fraction in the PVTSOL keyword would be replaced
by the lean injected gas fraction. Again the other properties would be
derived from laboratory or numerical slim-tube experiments based on
one-dimensional compositional
\hyperref[REF_HEADING_KEYWORD_EOS_5_3]{EOS} simulations.

\end{tcolorbox}

As an aside, the scope and accuracy of gas condensate modeling using
black-oil reservoir simulators is well established in the industry; both
depletion and gas cycling above the dew point can be modeled and yield
an adequate match with the results from multi-component compositional
simulators. The main inadequacy is with the treatment of gas injection
below the dew point where the primary compositional effect, the
stripping of liquid components is inversely proportional to their
molecular weights, is completely ignored. The standard black-oil model
assumes that the saturated hydrocarbon fluid properties are functions of
pressure only and disregards any compositional dependence in the
saturated fluid PVT properties. Thus, when dry gas is injected into a
condensate below its dew point the gas continues to re-vaporize liquid
at a rate governed only by the ambient pressure. The vapor saturates
over a zone whose thickness is of the order of one grid block; in
particular all the liquid in the vicinity of the injectors evaporates
rapidly. Results obtained with fully compositional simulation models
suggest that liquid saturation profiles would vary more slowly with
increasing distance from the gas injectors. This compositional effect
can thus be modeled via the CO2 Dynamic EOR Model, although there are no
test cases at this stage.

\subsubsection{Example}\label{example-334}

\begin{verbatim}
--
--       OIL PVT PROPERTIES FOR LIVE OIL VERSUS CO2 MASS FRACTION
--
PVTSOL
-- CO2    PSAT     BO        BG        RS        RV          XVOL    YVOL    OIL VISC    GAS VISC
-- MFRAC  PSIA     RB/STB  RB/MSCF     MSCF/STB  STB/MSCF                    CPOISE      CPOISE
-- -----  -------  ------  ----------  --------  ----------  ------  ------  ----------  ----------
   0.000    14.50  1.0000  1.7810e+02  0.0000    0.0000e+00  0.0000  0.0000  2.0340e-01  1.3016e-02
           725.18  1.0655  6.0855e+00  0.1302    1.4432e-04  0.0000  0.0000  2.0340e-01  1.3016e-02
           870.22  1.0806  5.0340e+00  0.1591    1.9594e-04  0.0000  0.0000  1.9973e-01  1.3222e-02
          1015.26  1.0961  4.2814e+00  0.1886    2.8389e-04  0.0000  0.0000  1.9609e-01  1.3452e-02
          1160.30  1.1121  3.7174e+00  0.2189    4.1914e-04  0.0000  0.0000  1.9246e-01  1.3708e-02
          1305.33  1.1285  3.2799e+00  0.2499    6.0779e-04  0.0000  0.0000  1.8886e-01  1.3994e-02
          1450.37  1.1454  2.9314e+00  0.2819    8.5105e-04  0.0000  0.0000  1.8527e-01  1.4309e-02
          1595.41  1.1629  2.6479e+00  0.3148    1.1490e-03  0.0000  0.0000  1.8171e-01  1.4657e-02
          1740.45  1.1809  2.4133e+00  0.3487    1.5030e-03  0.0000  0.0000  1.7819e-01  1.5037e-02
          1885.49  1.1996  2.2163e+00  0.3838    1.9155e-03  0.0000  0.0000  1.7469e-01  1.5450e-02
          2030.52  1.2189  2.0490e+00  0.4200    2.3899e-03  0.0000  0.0000  1.7122e-01  1.5897e-02
          2175.56  1.2389  1.9055e+00  0.4574    2.9302e-03  0.0000  0.0000  1.6779e-01  1.6377e-02
          2320.60  1.2597  1.7812e+00  0.4962    3.5412e-03  0.0000  0.0000  1.6440e-01  1.6890e-02
          2465.64  1.2812  1.6728e+00  0.5364    4.2273e-03  0.0000  0.0000  1.6104e-01  1.7435e-02
          2610.67  1.3037  1.5775e+00  0.5782    4.9942e-03  0.0000  0.0000  1.5772e-01  1.8011e-02
          2755.71  1.3270  1.4933e+00  0.6216    5.8471e-03  0.0000  0.0000  1.5443e-01  1.8617e-02
          2900.75  1.3514  1.4185e+00  0.6668    6.7918e-03  0.0000  0.0000  1.5117e-01  1.9253e-02
          3045.79  1.3768  1.3517e+00  0.7139    7.8347e-03  0.0000  0.0000  1.4795e-01  1.9917e-02
          3137.32  1.3937  1.3132e+00  0.7450    8.5549e-03  0.0000  0.0000  1.4592e-01  2.0353e-02
          4061.05  1.3705  1.3132e+00  0.7450    8.5549e-03  0.0000  0.0000  1.5939e-01  2.0353e-02 /
   0.010    14.50  1.0000  1.7810e+02  0.0000    0.0000e+00  0.0464  0.0381  2.0301e-01  1.3203e-02
           725.18  1.0670  6.0041e+00  0.1341    1.4845e-04  0.0464  0.0381  2.0301e-01  1.3203e-02
           …………….………….….…….…….…….…….……….………….…….…….…….………….…….……….……………….……….…...…….…….….…….…………...
          3135.04  1.4082  1.2950e+00  0.7755    8.9113e-03  0.0298  0.0277  1.4422e-01  2.0743e-02
          4061.05  1.3840  1.2950e+00  0.7755    8.9113e-03  0.0298  0.0277  1.5772e-01  2.0743e-02 /
   0.100    14.50  1.0000  1.7810e+02  0.0000    0.0000e+00  0.3426  0.3023  1.9984e-01  1.4740e-02
           725.18  1.0795  5.3329e+00  0.1666    1.8243e-04  0.3426  0.3023  1.9984e-01  1.4740e-02
           …………….………….….…….…….…….…….……….………….…….…….…….………….…….……….……………….……….…...…….…….….…….…………...
          3116.25  1.5281  1.1444e+00  1.0268    1.1849e-02  0.2492  0.2431  1.3023e-01  2.3961e-02
          4061.05  1.4955  1.1444e+00  1.0268    1.1849e-02  0.2492  0.2431  1.4400e-01  2.3961e-02 /
   0.200    14.50  1.0000  1.7810e+02  0.0000    0.0000e+00  0.5303  0.4898  1.9699e-01  1.5838e-02
           725.18  1.0909  4.8402e+00  0.1964    2.1919e-04  0.5303  0.4898  1.9699e-01  1.5838e-02
           …………….………….….…….…….…….…….……….………….…….…….…….………….…….……….……………….……….…...…….…….….…….…………...
          4061.05  1.6501  1.0017e+00  1.3720    1.6728e-02  0.4220  0.4240  1.2963e-01  2.7757e-02 /
   0.300    14.50  1.0000  1.7810e+02  0.0000    0.0000e+00  0.6486  0.6159  1.9471e-01  1.6588e-02
           725.18  1.1003  4.4992e+00  0.2210    2.5440e-04  0.6486  0.6159  1.9471e-01  1.6588e-02
           …………….………….….…….…….…….…….……….………….…….…….…….………….…….……….……………….……….…...…….…….….…….…………...
          4061.05  1.8529  8.8189e-01  1.8200    2.4541e-02  0.5491  0.5597  1.1636e-01  3.2195e-02 /
   0.400    14.50  1.0000  1.7810e+02  0.0000    0.0000e+00  0.7303  0.7059  1.9292e-01  1.7131e-02
           725.18  1.1079  4.2506e+00  0.2413    2.8832e-04  0.7303  0.7059  1.9292e-01  1.7131e-02
           …………….………….….…….…….…….…….……….………….…….…….…….………….…….……….……………….……….…...…….…….….…….…………...
          4061.05  2.1298  7.8515e-01  2.4251    3.7664e-02  0.6468  0.6616  1.0441e-01  3.7755e-02 /
   0.500    14.50  1.0000  1.7810e+02  0.0000    0.0000e+00  0.7905  0.7733  1.9156e-01  1.7543e-02
           725.18  1.1140  4.0628e+00  0.2581    3.2217e-04  0.7905  0.7733  1.9156e-01  1.7543e-02
           …………….………….….…….…….…….…….……….………….…….…….…….………….…….……….……………….……….…...…….…….….…….…………...
          4061.05  2.5282  7.1557e-01  3.2879    6.0741e-02  0.7245  0.7379  9.3885e-02  4.5038e-02 /
   0.600    14.50  1.0000  1.7810e+02  0.0000    0.0000e+00  0.8372  0.8259  1.9061e-01  1.7868e-02
           725.18  1.1188  3.9168e+00  0.2720    3.5813e-04  0.8372  0.8259  1.9061e-01  1.7868e-02
           …………….………….….…….…….…….…….……….………….…….…….…….………….…….……….……………….……….…...…….…….….…….…………...
          4061.05  3.1455  6.8307e-01  4.6154    1.0176e-01  0.7880  0.7951  8.4761e-02  5.4583e-02 /
   0.700    14.50  1.0000  1.7810e+02  0.0000    0.0000e+00  0.8755  0.8686  1.9018e-01  1.8132e-02
           725.18  1.1223  3.8017e+00  0.2833    4.0084e-04  0.8755  0.8686  1.9018e-01  1.8132e-02
           …………….………….….…….…….…….…….……….………….…….…….…….………….…….……….……………….……….…...…….…….….…….…………...
          4061.05  3.3568  6.6803e-01  5.4784    1.3282e-01  0.9106  0.9124  8.9321e-02  6.2163e-02 /
   0.800    14.50  1.0000  1.7810e+02  0.0000    0.0000e+00  0.9520  0.9591  1.8807e-01  1.8720e-02
           725.18  1.1335  3.5254e+00  0.3153    5.0228e-04  0.9520  0.9591  1.8807e-01  1.8720e-02
           …………….………….….…….…….…….…….……….………….…….…….…….………….…….……….……………….……….…...…….…….….…….…………...
          4061.05  2.7406  5.5115e-01  3.5079    8.5588e-02  0.8828  0.8905  8.1404e-02  7.0425e-02 /
   0.900    14.50  1.0000  1.7810e+02  0.0000    0.0000e+00  0.9846  0.9855  1.9003e-01  1.8852e-02
           725.18  1.1293  3.5076e+00  0.3007    4.6299e-04  0.9846  0.9855  1.9003e-01  1.8852e-02
           …………….………….….…….…….…….…….……….………….…….…….…….………….…….……….……………….……….…...…….…….….…….…………...
          4061.05  1.9854  5.0045e-01  2.0535    3.9219e-02  0.9400  0.9448  1.0088e-01  6.5831e-02 /
   1.000    14.50  1.0000  1.7810e+02  0.0000    0.0000e+00  1.0000  1.0000  1.7716e-01  1.8887e-02
           725.18  1.1418  3.5421e+00  0.2487    1.7810e-10  1.0000  1.0000  1.7716e-01  1.8887e-02
           …………….………….….…….…….…….…….……….………….…….…….…….………….…….……….……………….……….…...…….…….….…….…………...
          4061.05  1.2145  4.3773e-01  0.6633    4.2745e-09  1.0000  1.0000  1.4288e-01  6.2211e-02 /
/
\end{verbatim}

The above example defines one oil PVT Properties for Live Oil versus CO2
Mass Fraction table, and assumes that NTPVT equals one the TABDIMS
keyword in the RUNSPEC section.

\subsection{PVTW -- Define Water Fluid Properties for Various
Regions}\label{pvtw-define-water-fluid-properties-for-various-regions}

\begin{longtable}[]{@{}
  >{\raggedright\arraybackslash}p{(\linewidth - 14\tabcolsep) * \real{0.1250}}
  >{\raggedright\arraybackslash}p{(\linewidth - 14\tabcolsep) * \real{0.1250}}
  >{\raggedright\arraybackslash}p{(\linewidth - 14\tabcolsep) * \real{0.1250}}
  >{\raggedright\arraybackslash}p{(\linewidth - 14\tabcolsep) * \real{0.1250}}
  >{\raggedright\arraybackslash}p{(\linewidth - 14\tabcolsep) * \real{0.1250}}
  >{\raggedright\arraybackslash}p{(\linewidth - 14\tabcolsep) * \real{0.1250}}
  >{\raggedright\arraybackslash}p{(\linewidth - 14\tabcolsep) * \real{0.1250}}
  >{\raggedright\arraybackslash}p{(\linewidth - 14\tabcolsep) * \real{0.1250}}@{}}
\toprule\noalign{}
\begin{minipage}[b]{\linewidth}\raggedright
RUNSPEC
\end{minipage} & \begin{minipage}[b]{\linewidth}\raggedright
GRID
\end{minipage} & \begin{minipage}[b]{\linewidth}\raggedright
EDIT
\end{minipage} & \begin{minipage}[b]{\linewidth}\raggedright
PROPS
\end{minipage} & \begin{minipage}[b]{\linewidth}\raggedright
REGIONS
\end{minipage} & \begin{minipage}[b]{\linewidth}\raggedright
SOLUTION
\end{minipage} & \begin{minipage}[b]{\linewidth}\raggedright
SUMMARY
\end{minipage} & \begin{minipage}[b]{\linewidth}\raggedright
SCHEDULE
\end{minipage} \\
\midrule\noalign{}
\endhead
\bottomrule\noalign{}
\endlastfoot
\end{longtable}

\subsubsection{Description}\label{description-599}

PVTW defines the water properties for various regions in the model. The
number of PVTW vector data sets is defined by the NTPVT parameter on the
TABDIMS keyword in the RUNSPEC section and the allocation of the PVTW
tables to different grid blocks in the model is done via the PVTNUM
keyword in the REGIONS section. One data set consists of one record or
line which is terminated by a ``/''. If the water phase is active in the
model, which is normally the case, then this keyword must be defined in
the OPM Flow input deck.

\begin{longtable}[]{@{}
  >{\raggedright\arraybackslash}p{(\linewidth - 6\tabcolsep) * \real{0.2500}}
  >{\raggedright\arraybackslash}p{(\linewidth - 6\tabcolsep) * \real{0.2500}}
  >{\raggedright\arraybackslash}p{(\linewidth - 6\tabcolsep) * \real{0.2500}}
  >{\raggedright\arraybackslash}p{(\linewidth - 6\tabcolsep) * \real{0.2500}}@{}}
\toprule\noalign{}
\begin{minipage}[b]{\linewidth}\raggedright
No.
\end{minipage} & \begin{minipage}[b]{\linewidth}\raggedright
Name
\end{minipage} & \begin{minipage}[b]{\linewidth}\raggedright
Description
\end{minipage} & \begin{minipage}[b]{\linewidth}\raggedright
Default
\end{minipage} \\
\midrule\noalign{}
\endhead
\bottomrule\noalign{}
\endlastfoot
Field & Metric & Laboratory & \\
1 & PRES & PRES is a real number defining the water reference pressure
(P) for the other parameters for this data set. & None \\
psia & barsa & atma & \\
2 & WFVF & WFVF is a real number defining the water formation volume
factor (Bw) at the water reference pressure. & Defined \\
rb/stb 1.0 & rm3/sm3 1.0 & rcc/scc 1.0 & \\
3 & WCOMP & WCOMP is a real number defining the water compressibility
(Cw) at the water reference pressure and is defined as:
\({C}_{w} = -\frac{1}{{B}_{w}}\left(\frac{{\mathit{dB}}_{w}}{\mathit{dP}}\right)\)
& Defined \\
1/psia 0.00004 & 1/barsa 0.00004 & 1/atma 0.00004 & \\
4 & WVISC & WVISC is a real number defining the water viscosity (µw) at
the water reference pressure. & Defined \\
cP 0.50 & cP 0.50 & cP 0.50 & \\
5 & WVISCOMP & WVISCOMP is a real number defining the water
viscosibility (µwc) at the water reference pressure, µwc(Pref) and is
defined as:
\({\mathrm{μ}}_{\mathit{wc}} = -\frac{1}{{\mathrm{μ}}_{w}}\left(\frac{d{\mathrm{μ}}_{w}}{\mathit{dP}}\right)\)
& Defined \\
1/psia 0.0 & 1/barsa 0.0 & 1/atma 0.0 & \\
Notes: & & & \\
\end{longtable}

\emph{Table 8.122: PVTW Keyword Description}

Note that provided the first table has been entered, subsequent tables
may be defaulted, in this case the prior table is copied to the current
table. See the third example for an illustration on how to use this
feature.

\subsubsection{Examples}\label{examples-39}

The following shows the PVTW keyword for when NTPVT on the TABDIMS
keyword in the RUNSPEC section is set to one.

\begin{verbatim}
--
--       WATER PVT TABLE
--
PVTW
--       REF PRES  BW         CW        VISC     VISC
--       PSIA      RB/STB     1/PSIA    CPOISE   GRAD
--       --------  --------   -------   ------   ------
         4840.0    1.019     2.7E-6     0.370    1*         / TABLE NO. 01
\end{verbatim}

The next example shows the PVTW keyword for when NTPVT on the TABDIMS
keyword in the RUNSPEC section is set to three.

\begin{verbatim}
--
--       WATER PVT TABLE
--
PVTW
--       REF PRES  BW         CW        VISC     VISC
--       PSIA      RB/STB     1/PSIA    CPOISE   GRAD
--       --------  --------   -------   ------   ------
         4640.0    1.008     2.5E-6     0.350    1*         / TABLE NO. 01
         4840.0    1.019     2.7E-6     0.370    1*         / TABLE NO. 02
         4940.0    1.030     2.8E-6     0.390    1*         / TABLE NO. 03
\end{verbatim}

The above example defines three water PVT tables and assumes that NTPVT
equals three on the TABDIMS keyword in the RUNSPEC section.

The third, and final example, shows the PVTW keyword for when NTPVT on
the TABDIMS keyword in the RUNSPEC section is set to four. Here table
two defaults to table one, and table four defaults to table three

\begin{verbatim}
--
--       WATER PVT TABLE
--
PVTW
--       REF PRES  BW         CW        VISC     VISC
--       PSIA      RB/STB     1/PSIA    CPOISE   GRAD
--       --------  --------   -------   ------   ------
         4640.0    1.008     2.5E-6     0.350    1*         / TABLE NO. 01
                                                            / TABLE NO. 02
         4940.0    1.030     2.8E-6     0.390    1*         / TABLE NO. 03
                                                            / TABLE NO. 04
\end{verbatim}

Note that there is no terminating ``/'' for this keyword.

\subsection{PVTWSALT -- Define Brine Water Fluid Properties for Various
Regions}\label{pvtwsalt-define-brine-water-fluid-properties-for-various-regions}

\begin{longtable}[]{@{}
  >{\raggedright\arraybackslash}p{(\linewidth - 14\tabcolsep) * \real{0.1250}}
  >{\raggedright\arraybackslash}p{(\linewidth - 14\tabcolsep) * \real{0.1250}}
  >{\raggedright\arraybackslash}p{(\linewidth - 14\tabcolsep) * \real{0.1250}}
  >{\raggedright\arraybackslash}p{(\linewidth - 14\tabcolsep) * \real{0.1250}}
  >{\raggedright\arraybackslash}p{(\linewidth - 14\tabcolsep) * \real{0.1250}}
  >{\raggedright\arraybackslash}p{(\linewidth - 14\tabcolsep) * \real{0.1250}}
  >{\raggedright\arraybackslash}p{(\linewidth - 14\tabcolsep) * \real{0.1250}}
  >{\raggedright\arraybackslash}p{(\linewidth - 14\tabcolsep) * \real{0.1250}}@{}}
\toprule\noalign{}
\begin{minipage}[b]{\linewidth}\raggedright
RUNSPEC
\end{minipage} & \begin{minipage}[b]{\linewidth}\raggedright
GRID
\end{minipage} & \begin{minipage}[b]{\linewidth}\raggedright
EDIT
\end{minipage} & \begin{minipage}[b]{\linewidth}\raggedright
PROPS
\end{minipage} & \begin{minipage}[b]{\linewidth}\raggedright
REGIONS
\end{minipage} & \begin{minipage}[b]{\linewidth}\raggedright
SOLUTION
\end{minipage} & \begin{minipage}[b]{\linewidth}\raggedright
SUMMARY
\end{minipage} & \begin{minipage}[b]{\linewidth}\raggedright
SCHEDULE
\end{minipage} \\
\midrule\noalign{}
\endhead
\bottomrule\noalign{}
\endlastfoot
\end{longtable}

\subsubsection{Description}\label{description-600}

PVTWSALT defines the brine water properties for various regions in the
model, for when the brine phase has been activated by the BRINE keyword
in the RUNSPEC section. In this case PVTWSALT is used instead of PVTW in
the input file. However, if the ECLMC keyword has been entered in the
RUNSPEC section to invoke the Multi-Component Brine model, the PVTW
keyword should be used instead of PVTWSALT, as with this combination the
salinity effect on the density is ignored.

The number of PVTWSALT table data sets is defined by the NTPVT parameter
on the TABDIMS keyword in the RUNSPEC section and the allocation of the
PVTWSALT tables to different grid blocks in the model is done via the
PVTNUM keyword in the REGION section.

\begin{longtable}[]{@{}
  >{\raggedright\arraybackslash}p{(\linewidth - 6\tabcolsep) * \real{0.2500}}
  >{\raggedright\arraybackslash}p{(\linewidth - 6\tabcolsep) * \real{0.2500}}
  >{\raggedright\arraybackslash}p{(\linewidth - 6\tabcolsep) * \real{0.2500}}
  >{\raggedright\arraybackslash}p{(\linewidth - 6\tabcolsep) * \real{0.2500}}@{}}
\toprule\noalign{}
\begin{minipage}[b]{\linewidth}\raggedright
No.
\end{minipage} & \begin{minipage}[b]{\linewidth}\raggedright
Name
\end{minipage} & \begin{minipage}[b]{\linewidth}\raggedright
Description
\end{minipage} & \begin{minipage}[b]{\linewidth}\raggedright
Default
\end{minipage} \\
\midrule\noalign{}
\endhead
\bottomrule\noalign{}
\endlastfoot
Field & Metric & Laboratory & \\
1-1 & PRESS & Single real positive value that defines the reference
pressure for the data in the following records (Pref). PRESS should be
approximately equal to the average reservoir pressures in the model. The
simulator uses the previous time step values to forecast the current
time step water properties by linear interpolation. If PRESS is not
representative of the average reservoir pressures in the model then the
linear interpolation might result in nonphysical values of the water
saturation and water viscosity. & None \\
psia & barsa & atma & \\
1-2 & SALTSURF & A real value that defines the reference salt
concentration in the solution in the surface stock tank water (Cs, ref)
If defaulted SALTSURF is taken as the minimum salt concentration entered
in the SALTCON columnar vector in the second record for this keyword.
This should be in most cases be zero. & Defined \\
1-1 & & & \\
lb/stb & kg/sm3 & gm/scc & \\
1-3 & / & Record terminated by a ``/'' & Not Applicable \\
2-1 & SALTCON & A columnar vector of real monotonically increasing
values that defines the salt concentration in the solution (Cs). &
None \\
lb/stb & kg/sm3 & gm/scc & \\
2-2 & WFVF & WFVF is a real columnar vector defining the water formation
volume factor (Bw) at the reference pressure PRESS, for the
corresponding salt concentration SALTCON. & None \\
2-1 & & & \\
rb/stb & rm3/sm3 & rcc/scc & \\
2-3 & WCOMP & WCOMP is a real columnar vector defining the water
compressibility (Cw) at the water reference pressure PRESS, for the
corresponding salt concentration SALTCON. The water compressibility is
defined as:
\({C}_{w} = -\frac{1}{{B}_{w}}\left(\frac{{\mathit{dB}}_{w}}{\mathit{dP}}\right)\)
& None \\
1/psia & 1/bars & 1/atma & \\
2.4 & WVISC & WVISC is a real columnar vector defining the water
viscosity (µw) at the water reference pressure PRESS, for the
corresponding salt concentration SALTCON. & None \\
cP & cP & cP & \\
2.5 & WVISCOMP & WVISCOMP is a real columnar vector defining the water
viscosibility (µwc) at the water reference pressure PRESS, for the
corresponding salt concentration SALTCON. The water viscosibility is
defined as:
\({\mathrm{μ}}_{\mathit{wc}} = -\frac{1}{{\mathrm{μ}}_{w}}\left(\frac{d{\mathrm{μ}}_{w}}{\mathit{dP}}\right)\)
& None \\
1/psia & 1/barsa & 1/atma & \\
2-6 & / & Table and record terminated by a ``/'' & Not Applicable \\
Notes: & & & \\
\end{longtable}

\emph{Table 8.123: PVTWSALT Keyword Description}

As mentioned above, the simulator first calculates the water properties
as functions of the salt concentration at the previous time step by
linear interpolation in salt concentration for water compressibility
(Cw), water viscosibility (μwc),
\(\frac{1}{{B}_{w}}\)and\(\frac{1}{{B}_{w} {μ}_{w}}\). It then
calculates the values of\({B}_{w}\)and\({B}_{w} {μ}_{w}\)at the current
time step using the current pressure P, using the following equations:

\begin{equation}\phantomsection\label{eq-8-82}{
{B}_{w}(P,{C}_{s}) = \frac{{B}_{w}({P}_{\mathit{ref}},{C}_{s,\mathit{ref}})}{1 + {C}_{w}(P-{P}_{\mathit{ref}}) + \frac{{\left({C}_{w}(P-{P}_{\mathit{ref}})\right)}^{2}}{2}}
}\end{equation}

and

\begin{equation}\phantomsection\label{eq-8-83}{
{B}_{w}(P,{C}_{s}) {\mathrm{μ}}_{w}(P,{C}_{s}) = \frac{{B}_{w}({P}_{\mathit{ref}},{C}_{s,\mathit{ref}}) {\mathrm{μ}}_{w}({P}_{\mathit{ref}},{C}_{s,\mathit{ref}})}{1 + ({C}_{w}- {\mathrm{μ}}_{\mathit{wc}})(P-{P}_{\mathit{ref}}) + \frac{{\left(({C}_{w}- {\mathrm{μ}}_{\mathit{wc}})(P-{P}_{\mathit{ref}})\right)}^{2}}{2}}
}\end{equation}

See also the BDENSITY keyword in the PROPS section that defines the
brine surface densities for the salt concentrations declared on the
PVTWSALT keyword. Note that if the BDENSITY keyword is absent from the
input file then the brine surface densities will be set to the water
density values declared via the DENSITY keyword in the PROPS section. In
this case there is no variation in brine surface density with respect to
salt concentration.

\subsubsection{Example}\label{example-335}

The following shows the PVTWSALT keyword for when NTPVT on the TABDIMS
keyword in the RUNSPEC section is set equal to two and NPPVT is set to
greater than four on the TABDIMS keyword.

\begin{verbatim}
--
--       WATER SALT PVT TABLE
--
PVTWSALT
--       REF PRES  REF SALT
--       PSIA      LB/STB
--       --------  --------
         4500.0    0.000                                 / TABLE NO. REF. DATA
--
--       SALTCON BW         CW        VISC     VISC
--       LB/STB    RB/STB     1/PSIA    CPOISE   GRAD
--       --------  --------   -------   ------   ------
            0.0    1.020      2.7E-6    0.370    0.0
            2.0    1.010      2.7E-6    0.370    0.0
            4.0    1.000      2.7E-6    0.370    0.0
           10.0    0.950      2.7E-6    0.370    0.0     / TABLE NO. 01 SALT DATA
--
--       REF PRES  REF SALT
--       PSIA      LB/STB
--       --------  --------
         4000.0    0.000                                 / TABLE NO. 02 REF. DATA
--
--       SALTCON  BW         CW        VISC     VISC
--       LB/STB    RB/STB     1/PSIA    CPOISE   GRAD
--       --------  --------   -------   ------   ------
            0.0    1.005      2.5E-6    0.320    0.0
            3.0    1.000      2.5E-6    0.320    0.0
            6.0    0.985      2.5E-6    0.320    0.0
           12.0    0.930      2.5E-6    0.320    0.0     / TABLE NO. 02 SALT DATA

\end{verbatim}

Note that each table is terminated by a ``/'' and there is no ``/''
terminator for the keyword.

\begin{verbatim}
\end{verbatim}

\subsection{PVZG -- Gas PVT Properties for Dry Gas
(Z-Factor)}\label{pvzg-gas-pvt-properties-for-dry-gas-z-factor}

\begin{longtable}[]{@{}
  >{\raggedright\arraybackslash}p{(\linewidth - 14\tabcolsep) * \real{0.1250}}
  >{\raggedright\arraybackslash}p{(\linewidth - 14\tabcolsep) * \real{0.1250}}
  >{\raggedright\arraybackslash}p{(\linewidth - 14\tabcolsep) * \real{0.1250}}
  >{\raggedright\arraybackslash}p{(\linewidth - 14\tabcolsep) * \real{0.1250}}
  >{\raggedright\arraybackslash}p{(\linewidth - 14\tabcolsep) * \real{0.1250}}
  >{\raggedright\arraybackslash}p{(\linewidth - 14\tabcolsep) * \real{0.1250}}
  >{\raggedright\arraybackslash}p{(\linewidth - 14\tabcolsep) * \real{0.1250}}
  >{\raggedright\arraybackslash}p{(\linewidth - 14\tabcolsep) * \real{0.1250}}@{}}
\toprule\noalign{}
\begin{minipage}[b]{\linewidth}\raggedright
RUNSPEC
\end{minipage} & \begin{minipage}[b]{\linewidth}\raggedright
GRID
\end{minipage} & \begin{minipage}[b]{\linewidth}\raggedright
EDIT
\end{minipage} & \begin{minipage}[b]{\linewidth}\raggedright
PROPS
\end{minipage} & \begin{minipage}[b]{\linewidth}\raggedright
REGIONS
\end{minipage} & \begin{minipage}[b]{\linewidth}\raggedright
SOLUTION
\end{minipage} & \begin{minipage}[b]{\linewidth}\raggedright
SUMMARY
\end{minipage} & \begin{minipage}[b]{\linewidth}\raggedright
SCHEDULE
\end{minipage} \\
\midrule\noalign{}
\endhead
\bottomrule\noalign{}
\endlastfoot
\end{longtable}

\subsubsection{Description}\label{description-601}

PVZG defines the gas PVT properties for dry gas {[}Natural gas that
occurs in the absence of condensate or liquid hydrocarbons, or gas that
had condensable hydrocarbons removed, is called dry gas. It is primarily
methane with some intermediates. The hydrocarbon mixture is solely gas
in the reservoir and there is no liquid (condensate surface liquid)
formed either in the reservoir or at surface. The term dry indicates
that the gas does not contain heavier hydrocarbons to form liquids at
the surface conditions. Dry gas typically has GOR's greater than 100,000
scf/stb or 18,000 Sm3/m3.{]} via the gas compressibility factor
(z-factor), instead of the gas formation volume factor. If the gas has a
constant and uniform vaporized oil concentration, Condensate-Gas Ratio
(``CGR''), and if the reservoir pressure never drops below the
saturation pressure (dew point pressure), then the model can be run more
efficiently by omitting the OIL and VAPOIL keywords from the RUNSPEC
section, treating the gas as a dry gas, and defining a constant Rv (CGR)
value with keyword RVCONST or RVCONSTT in the PROPS section. This
results in the model being run with as a dry gas problem with no active
oil (condensate) phase. However, OPM Flow takes into account the
constant Rv in the calculations and reporting.

This keyword is not supported by OPM Flow but would change the results
if supported so the simulation will be stopped.

\begin{longtable}[]{@{}
  >{\raggedright\arraybackslash}p{(\linewidth - 6\tabcolsep) * \real{0.2500}}
  >{\raggedright\arraybackslash}p{(\linewidth - 6\tabcolsep) * \real{0.2500}}
  >{\raggedright\arraybackslash}p{(\linewidth - 6\tabcolsep) * \real{0.2500}}
  >{\raggedright\arraybackslash}p{(\linewidth - 6\tabcolsep) * \real{0.2500}}@{}}
\toprule\noalign{}
\begin{minipage}[b]{\linewidth}\raggedright
No.
\end{minipage} & \begin{minipage}[b]{\linewidth}\raggedright
Name
\end{minipage} & \begin{minipage}[b]{\linewidth}\raggedright
Description
\end{minipage} & \begin{minipage}[b]{\linewidth}\raggedright
Default
\end{minipage} \\
\midrule\noalign{}
\endhead
\bottomrule\noalign{}
\endlastfoot
Field & Metric & Laboratory & \\
1-1 & RTEMP & Single real positive value that defines the reservoir
temperature for the data in the following records. & \\
oF & oC & oC & \\
1-2 & / & Record terminated by a ``/'' & Not Applicable \\
2-1 & PRESS & A columnar vector of real monotonically increasing down
the column values that defines the gas phase pressure. & None \\
psia & barsa & atma & \\
2-2 & GZFAC & A columnar vector of real values that defines the
corresponding gas phase z-factor at the given pressure, PRESS. & None \\
2-1 & & & \\
dimensionless & dimensionless & dimensionless & \\
2-3 & GVISC & A columnar vector of real increasing down the column
values that defines the corresponding gas phase viscosity. & None \\
cP & cP & cP & \\
Notes: & & & \\
\end{longtable}

\emph{Table 8.124: PVZG Keyword Description}

See also the RVCONST and RVCONSTT keywords to define the constant Rv for
dry gas.

The ideal gas law provides a relationship between the pressure, the
temperature and the specific volume of an ideal gas (pure component).
This relationship is modified by use of a compressibility factor, Z
{[}Standing, M. B.:'' Volumetric and Phase Behaviour of Oil Field
Hydrocarbon Systems'', Renihold Publishing Corp., New York City
(1952).{]}, to account for deviations, from ideal, to the behavior of
real gases. The PVT relation for a real gas can be defined by:

\begin{equation}\phantomsection\label{eq-8-84}{
\mathit{PV} = \mathit{ZnRT}
}\end{equation}

As the gas formation volume factor is used to relate the volume of gas,
as measured at reservoir conditions, to the volume of gas as measured at
standard conditions (60 oF and 14.7 psia, or 15 oC and 101.325 kPa).
This gas property is then defined as the actual volume occupied by a
certain amount of gas at a specified pressure and temperature, divided
by the same amount of gas at standard conditions. Thus, using the above
equation one can obtain the gas volumes at reservoir and standard
conditions, i.e.

\begin{equation}\phantomsection\label{eq-8-85}{
{V}_{\mathit{sc}} = \frac{{Z}_{\mathit{sc}}{\mathit{nRT}}_{\mathit{sc}}}{{P}_{\mathit{sc}}}
}\end{equation}

\begin{equation}\phantomsection\label{eq-8-86}{
{V}_{i} = \frac{{Z}_{i}n{\mathit{RT}}_{i}}{{P}_{i}}
}\end{equation}

Thus the gas formation volume factor can be expressed as:

\begin{equation}\phantomsection\label{eq-8-87}{
E = \frac{{V}_{\mathit{sc}}}{{V}_{i}}
}\end{equation}

And substituting equation (8.85) and (8.86) into (8.87) we obtain

\begin{equation}\phantomsection\label{eq-8-88}{
E = \left(\frac{{P}_{i}}{{P}_{\mathit{sc}}}\right) \left(\frac{{T}_{\mathit{sc}}}{{T}_{i}}\right) \left(\frac{1}{{Z}_{i}}\right)
}\end{equation}

Incorporating standard pressure and temperature values gives in SI
units:

\begin{equation}\phantomsection\label{eq-8-89}{
E = \left(\frac{{P}_{i}}{101.325}\right) \left(\frac{273.15 + 15}{{T}_{i}}\right) \left(\frac{1}{{Z}_{i}}\right) = 2.84 \left(\frac{{P}_{i}}{{Z}_{i}{T}_{i}}\right)
}\end{equation}

or in field units:

\begin{equation}\phantomsection\label{eq-8-90}{
E = \left(\frac{{P}_{i}}{14.7}\right) \left(\frac{460 + 60}{{T}_{i}}\right) \left(\frac{1}{{Z}_{i}}\right) = 35.37 \left(\frac{{P}_{i}}{{Z}_{i}{T}_{i}}\right)
}\end{equation}

Where,

\(E\) = gas formation volume factor (scf /rcft or Sm3/m3)

\(P\) = pressure (psia or kPa)

\(\mathit{Psc}\) = pressure standard conditions (psia or kPa )

\(\mathit{Pi}\) = initial reservoir pressure (psia or kPa)

\(V\) = volume (ft3 or m3)

\(T\) = absolute temperature (oR or K)

\(\mathit{Tsc}\) = temperature at standard conditions (oR or K)

\(\mathit{Ti}\) = initial reservoir temperature (oR or K)

\(R\) = gas constant (10.73 or 8.314)

RTEMP on this keyword is the \(\mathit{Ti}\)in the above equations
and\({P}_{i}\) is PRESS columnar vector.

\subsubsection{Example}\label{example-336}

\begin{verbatim}
---
--       GAS PVT TABLE USING GAS Z-FACTOR
--
PVZG
--       RESERVOIR TEMPERATURE FOR Z TO BG CONVERSION
--
         180.0                                             /
--
--        PRES       ZG          VISC
--        PSIA     DIMLESS      CPOISE
--       ------    ---------    -------
           14.7    0.998970     0.0130
          250.0    0.976260     0.0131
          500.0    0.954790     0.0134
          750.0    0.932050     0.0137
         1000.0    0.912990     0.0142
         1250.0    0.896320     0.0147
         1500.0    0.881610     0.0152
         1750.0    0.870830     0.0159
         2000.0    0.863130     0.0166
         2250.0    0.858920     0.0173
         2500.0    0.857800     0.0181
         2750.0    0.860430     0.0189
         3000.0    0.866440     0.0197
         3250.0    0.874980     0.0206
         3500.0    0.885470     0.0214
         3750.0    0.898350     0.0223
         4000.0    1.025120     0.0277                     / TABLE NO 01              --
--       GAS PVT TABLE USING GAS Z-FACTOR
--
PVZG
--       RESERVOIR TEMPERATURE FOR Z TO BG CONVERSION
--
         180.0                                             /
--
--        PRES       ZG          VISC
--        PSIA     DIMLESS      CPOISE
--       ------    ---------    -------
           14.7    0.998970     0.0130
          250.0    0.976260     0.0131
          500.0    0.954790     0.0134
          750.0    0.932050     0.0137
         1000.0    0.912990     0.0142
         1250.0    0.896320     0.0147
         1500.0    0.881610     0.0152
         1750.0    0.870830     0.0159
         2000.0    0.863130     0.0166
         2250.0    0.858920     0.0173
         2500.0    0.857800     0.0181
         2750.0    0.860430     0.0189
         3000.0    0.866440     0.0197
         3250.0    0.874980     0.0206
         3500.0    0.885470     0.0214
         3750.0    0.898350     0.0223
         4000.0    1.025120     0.0277                     / TABLE NO 01
\end{verbatim}

The above example defines two dry PVZG tables assuming NTPVT equals two
and NPPVT is greater than or equal to 17 on the TABDIMS keyword in the
RUNSPEC section. There is no terminating ``/'' for this keyword.

\subsection{PYEND -- End the Definition of a PYINPUT
Section}\label{pyend-end-the-definition-of-a-pyinput-section-2}

The PYINPUT and PYEND keywords are a part of OPM Flow's Python scripting
facility that processes standard Python commands that can be used to
manipulate and define the simulators input parameters during processing
of the input deck. The main purpose of the facility is to script the
construction of the various keywords. A PYINPUT Definition Section is
terminated by a PYEND keyword on a separate single line.

See PYEND -- End the Definition of a PYINPUT Section in the GRID section
for a full description.

\subsection{PYINPUT -- Define the Start of a PYINPUT
Section}\label{pyinput-define-the-start-of-a-pyinput-section-2}

The PYINPUT and PYEND keywords are a part of OPM Flow's Python scripting
facility that processes standard Python commands that can be used to
manipulate and define the simulators input parameters during processing
of the input deck. The main purpose of the facility is to script the
construction of the various keywords used by the simulator. PYINPUT
declares the start of a PYINPUT Definition Section on a single separate
line, which is then followed by various standard Python commands, with
one command per line.

See PYINPUT -- Define the Start of a PYINPUT Section in the GRID section
for a full description.

\subsection{QHRATING -- Define River Mass Flow versus Depth
Tables}\label{qhrating-define-river-mass-flow-versus-depth-tables}

\begin{longtable}[]{@{}
  >{\raggedright\arraybackslash}p{(\linewidth - 14\tabcolsep) * \real{0.1250}}
  >{\raggedright\arraybackslash}p{(\linewidth - 14\tabcolsep) * \real{0.1250}}
  >{\raggedright\arraybackslash}p{(\linewidth - 14\tabcolsep) * \real{0.1250}}
  >{\raggedright\arraybackslash}p{(\linewidth - 14\tabcolsep) * \real{0.1250}}
  >{\raggedright\arraybackslash}p{(\linewidth - 14\tabcolsep) * \real{0.1250}}
  >{\raggedright\arraybackslash}p{(\linewidth - 14\tabcolsep) * \real{0.1250}}
  >{\raggedright\arraybackslash}p{(\linewidth - 14\tabcolsep) * \real{0.1250}}
  >{\raggedright\arraybackslash}p{(\linewidth - 14\tabcolsep) * \real{0.1250}}@{}}
\toprule\noalign{}
\begin{minipage}[b]{\linewidth}\raggedright
RUNSPEC
\end{minipage} & \begin{minipage}[b]{\linewidth}\raggedright
GRID
\end{minipage} & \begin{minipage}[b]{\linewidth}\raggedright
EDIT
\end{minipage} & \begin{minipage}[b]{\linewidth}\raggedright
PROPS
\end{minipage} & \begin{minipage}[b]{\linewidth}\raggedright
REGIONS
\end{minipage} & \begin{minipage}[b]{\linewidth}\raggedright
SOLUTION
\end{minipage} & \begin{minipage}[b]{\linewidth}\raggedright
SUMMARY
\end{minipage} & \begin{minipage}[b]{\linewidth}\raggedright
SCHEDULE
\end{minipage} \\
\midrule\noalign{}
\endhead
\bottomrule\noalign{}
\endlastfoot
\end{longtable}

\subsubsection{Description}\label{description-602}

The QHRATING keyword defines a river's mass flow rate versus depth
tables, for when the River option has been activated via the RIVRDIMS
keyword in the RUNSPEC section.

This keyword is not supported by OPM Flow but would change the results
if supported so the simulation will be stopped.

\subsection{REFINE -- Start the Definition of a Local Grid
Refinement}\label{refine-start-the-definition-of-a-local-grid-refinement-2}

The REFINE keyword defines the start of a Cartesian or radial local grid
refinement (``LGR'') definition that sets the properties of the selected
LGR. The keyword is then followed by the property keywords associated
with the section where the keyword is being invoked. For example, if the
REFINE keyword is used in the GRID section then most of the keywords in
that section can be used to set the grid properties for the LGR.

See REFINE -- Start the Definition of a Local Grid Refinement in the
GRID section for a full description.

\subsection{RIVRXSEC -- Define River Cross-Section versus Depth
Parameters}\label{rivrxsec-define-river-cross-section-versus-depth-parameters}

\begin{longtable}[]{@{}
  >{\raggedright\arraybackslash}p{(\linewidth - 14\tabcolsep) * \real{0.1250}}
  >{\raggedright\arraybackslash}p{(\linewidth - 14\tabcolsep) * \real{0.1250}}
  >{\raggedright\arraybackslash}p{(\linewidth - 14\tabcolsep) * \real{0.1250}}
  >{\raggedright\arraybackslash}p{(\linewidth - 14\tabcolsep) * \real{0.1250}}
  >{\raggedright\arraybackslash}p{(\linewidth - 14\tabcolsep) * \real{0.1250}}
  >{\raggedright\arraybackslash}p{(\linewidth - 14\tabcolsep) * \real{0.1250}}
  >{\raggedright\arraybackslash}p{(\linewidth - 14\tabcolsep) * \real{0.1250}}
  >{\raggedright\arraybackslash}p{(\linewidth - 14\tabcolsep) * \real{0.1250}}@{}}
\toprule\noalign{}
\begin{minipage}[b]{\linewidth}\raggedright
RUNSPEC
\end{minipage} & \begin{minipage}[b]{\linewidth}\raggedright
GRID
\end{minipage} & \begin{minipage}[b]{\linewidth}\raggedright
EDIT
\end{minipage} & \begin{minipage}[b]{\linewidth}\raggedright
PROPS
\end{minipage} & \begin{minipage}[b]{\linewidth}\raggedright
REGIONS
\end{minipage} & \begin{minipage}[b]{\linewidth}\raggedright
SOLUTION
\end{minipage} & \begin{minipage}[b]{\linewidth}\raggedright
SUMMARY
\end{minipage} & \begin{minipage}[b]{\linewidth}\raggedright
SCHEDULE
\end{minipage} \\
\midrule\noalign{}
\endhead
\bottomrule\noalign{}
\endlastfoot
\end{longtable}

\subsubsection{Description}\label{description-603}

The RIVRXSEC keyword defines a river's cross-sectional area and
perimeter versus depth parameters. The River option must be activated
via the RIVRDIMS keyword in the RUNSPEC section in order to use this
keyword.

This keyword is not supported by OPM Flow but would change the results
if supported so the simulation will be stopped.

\subsection{RKTRMDIR -- Activate ROCKTAB Keyword Directional
Transmissibility
Multipliers}\label{rktrmdir-activate-rocktab-keyword-directional-transmissibility-multipliers}

\begin{longtable}[]{@{}
  >{\raggedright\arraybackslash}p{(\linewidth - 14\tabcolsep) * \real{0.1250}}
  >{\raggedright\arraybackslash}p{(\linewidth - 14\tabcolsep) * \real{0.1250}}
  >{\raggedright\arraybackslash}p{(\linewidth - 14\tabcolsep) * \real{0.1250}}
  >{\raggedright\arraybackslash}p{(\linewidth - 14\tabcolsep) * \real{0.1250}}
  >{\raggedright\arraybackslash}p{(\linewidth - 14\tabcolsep) * \real{0.1250}}
  >{\raggedright\arraybackslash}p{(\linewidth - 14\tabcolsep) * \real{0.1250}}
  >{\raggedright\arraybackslash}p{(\linewidth - 14\tabcolsep) * \real{0.1250}}
  >{\raggedright\arraybackslash}p{(\linewidth - 14\tabcolsep) * \real{0.1250}}@{}}
\toprule\noalign{}
\begin{minipage}[b]{\linewidth}\raggedright
RUNSPEC
\end{minipage} & \begin{minipage}[b]{\linewidth}\raggedright
GRID
\end{minipage} & \begin{minipage}[b]{\linewidth}\raggedright
EDIT
\end{minipage} & \begin{minipage}[b]{\linewidth}\raggedright
PROPS
\end{minipage} & \begin{minipage}[b]{\linewidth}\raggedright
REGIONS
\end{minipage} & \begin{minipage}[b]{\linewidth}\raggedright
SOLUTION
\end{minipage} & \begin{minipage}[b]{\linewidth}\raggedright
SUMMARY
\end{minipage} & \begin{minipage}[b]{\linewidth}\raggedright
SCHEDULE
\end{minipage} \\
\midrule\noalign{}
\endhead
\bottomrule\noalign{}
\endlastfoot
\end{longtable}

\subsubsection{Description}\label{description-604}

This keyword activates the directional transmissibility multipliers for
the ROCKTAB keyword. This results in two additional columns being
inputted on the ROCKTAB keyword. This feature is currently not supported
in OPM Flow.

This keyword is not supported by OPM Flow but would change the results
if supported so the simulation will be stopped.

\subsection{ROCK -- Define the Rock Compressibility for Various
Regions}\label{rock-define-the-rock-compressibility-for-various-regions}

\begin{longtable}[]{@{}
  >{\raggedright\arraybackslash}p{(\linewidth - 14\tabcolsep) * \real{0.1250}}
  >{\raggedright\arraybackslash}p{(\linewidth - 14\tabcolsep) * \real{0.1250}}
  >{\raggedright\arraybackslash}p{(\linewidth - 14\tabcolsep) * \real{0.1250}}
  >{\raggedright\arraybackslash}p{(\linewidth - 14\tabcolsep) * \real{0.1250}}
  >{\raggedright\arraybackslash}p{(\linewidth - 14\tabcolsep) * \real{0.1250}}
  >{\raggedright\arraybackslash}p{(\linewidth - 14\tabcolsep) * \real{0.1250}}
  >{\raggedright\arraybackslash}p{(\linewidth - 14\tabcolsep) * \real{0.1250}}
  >{\raggedright\arraybackslash}p{(\linewidth - 14\tabcolsep) * \real{0.1250}}@{}}
\toprule\noalign{}
\begin{minipage}[b]{\linewidth}\raggedright
RUNSPEC
\end{minipage} & \begin{minipage}[b]{\linewidth}\raggedright
GRID
\end{minipage} & \begin{minipage}[b]{\linewidth}\raggedright
EDIT
\end{minipage} & \begin{minipage}[b]{\linewidth}\raggedright
PROPS
\end{minipage} & \begin{minipage}[b]{\linewidth}\raggedright
REGIONS
\end{minipage} & \begin{minipage}[b]{\linewidth}\raggedright
SOLUTION
\end{minipage} & \begin{minipage}[b]{\linewidth}\raggedright
SUMMARY
\end{minipage} & \begin{minipage}[b]{\linewidth}\raggedright
SCHEDULE
\end{minipage} \\
\midrule\noalign{}
\endhead
\bottomrule\noalign{}
\endlastfoot
\end{longtable}

\subsubsection{Description}\label{description-605}

ROCK defines the rock compressibility for various regions in the model.
The number of ROCK vector data sets is defined by the NTPVT parameter on
the TABDIMS keyword in the RUNSPEC section and the allocation of the
ROCK tables to different grid blocks in the model is done via the PVTNUM
keyword in the REGIONS section. One data set consists of one record or
line which is terminated by a ``/''.

This keyword must be defined in the OPM Flow input deck.

\begin{longtable}[]{@{}
  >{\raggedright\arraybackslash}p{(\linewidth - 6\tabcolsep) * \real{0.2500}}
  >{\raggedright\arraybackslash}p{(\linewidth - 6\tabcolsep) * \real{0.2500}}
  >{\raggedright\arraybackslash}p{(\linewidth - 6\tabcolsep) * \real{0.2500}}
  >{\raggedright\arraybackslash}p{(\linewidth - 6\tabcolsep) * \real{0.2500}}@{}}
\toprule\noalign{}
\begin{minipage}[b]{\linewidth}\raggedright
No.
\end{minipage} & \begin{minipage}[b]{\linewidth}\raggedright
Name
\end{minipage} & \begin{minipage}[b]{\linewidth}\raggedright
Description
\end{minipage} & \begin{minipage}[b]{\linewidth}\raggedright
Default
\end{minipage} \\
\midrule\noalign{}
\endhead
\bottomrule\noalign{}
\endlastfoot
Field & Metric & Laboratory & \\
1 & PRESS & PRESS is a real number defining the rock reference pressure
for the other parameters for this data set. & Default \\
psia 14.7 & barsa 1.0132 & atma 1.0 & \\
2 & RCOMP & RCOMP is a real number defining the rock compressibility
(cf) at the rock reference pressure and is defined as:
\({c}_{f} = -\frac{1}{V}\left(\frac{\mathit{dV}}{\mathit{dP}}\right)\) &
Defined \\
1/psia 0.0 & 1/barsa 0.0 & 1/atma 0.0 & \\
Notes: Note, however, that if ROCKOPTS(ROCKOPT3) parameter has been used
to set the allocation of the ROCK data via the ROCKNUM array, then the
number of ROCK vectors should correspond to the value entered on
TABDIMS(NTROCC) in the RUNSPEC section. Similarly, if the
ROCKOPTS(ROCKOPT3) has been used to set the assignment of the ROCK data
via the SATNUM array, then the number of vectors should correspond to
the value entered via the TABDIMS(NTSFUN) parameter, since the tables
will be allocated via the SATNUM array. & & & \\
\end{longtable}

\emph{Table 8.125: ROCK Keyword Description}

The simulator adjusts the pore volume based on the reference pressure
(PRESS), that is:

\begin{equation}\phantomsection\label{eq-8-91}{
V({P}_{i}) = V({P}_{r})\left(1 + {c}_{f}\left({P}_{i} - {P}_{r}\right) + \frac{{\left({c}_{f}({P}_{i} - {P}_{r})\right)}^{2}}{2}\right)
}\end{equation}

where:

\({c}_{f}\) = rock compressibility (RCOMP) at the reference pressure
(PRESS),

\({P}_{i}\) = initial grid cell pressure,

\({P}_{r}\) = reference pressure (PRESS),

\(V({P}_{i})\) = pore volume at initial conditions, and

\(V({P}_{r})\) = pore volume at at the reference pressure.

\begin{tcolorbox}[enhanced jigsaw, title=\textcolor{quarto-callout-note-color}{\faInfo}\hspace{0.5em}{Note}, coltitle=black, breakable, toprule=.15mm, colbacktitle=quarto-callout-note-color!10!white, leftrule=.75mm, opacitybacktitle=0.6, colback=white, bottomtitle=1mm, toptitle=1mm, titlerule=0mm, left=2mm, colframe=quarto-callout-note-color-frame, arc=.35mm, rightrule=.15mm, bottomrule=.15mm, opacityback=0]

If the Rock Compaction option has been activated via the ROCKCOMP
keyword in the RUNSPEC section, then the ROCKTAB keyword in the PROPS
section should be used instead of ROCK keyword.

\end{tcolorbox}

See also the ROCKOPTS and ROCKTAB keywords in the PROPS section.

\subsubsection{Examples}\label{examples-40}

The following shows the ROCK keyword for when NTPVT on the TABDIMS
keyword in the RUNSPEC section is set to one.

\begin{verbatim}
--
--       ROCK COMPRESSIBILITY
--
--       REFERENCE PRESSURE IS TAKEN FROM THE HCPV WEIGHTED RESERVOIR PRESSURE
--       AS THE PORV IS ALREADY AT RESERVOIR CONDITIONS (OPM FLOW USES THE
--       REFERENCE PRESSURE) TO CONVERT THE GIVEN PORV TO RESERVOIR CONDITIONS
--       USING THE DATA ON THE ROCK KEYWORD)
--
--       REF PRES  CF
--       PSIA      1/PSIA
--       --------  --------
ROCK
         3966.9    5.0E-06                                 / ROCK COMPRESSIBILITY
\end{verbatim}

The next example shows the ROCK keyword for when NTPVT on the TABDIMS
keyword in the RUNSPEC section is set to three.

\begin{verbatim}
--
--       ROCK COMPRESSIBILITY
--
--       REFERENCE PRESSURE IS TAKEN FROM THE HCPV WEIGHTED RESERVOIR PRESSURE
--       AS THE PORV IS ALREADY AT RESERVOIR CONDITIONS (OPM FLOW USES THE
--       REFERENCE PRESSURE) TO CONVERT THE GIVEN PORV TO RESERVOIR CONDITIONS
--       USING THE DATA ON THE ROCK KEYWORD)
--
--       REF PRES  CF
--       PSIA      1/PSIA
--       --------  --------
ROCK
         3566.9    5.0E-06                        / ROCK COMPRESSIBILITY REGION 1
         3966.9    5.5E-06                        / ROCK COMPRESSIBILITY REGION 2
         4566.9    6.0E-06                        / ROCK COMPRESSIBILITY REGION 3
\end{verbatim}

The above example defines three ROCK tables and assumes that NTPVT
equals three on the TABDIMS keyword in the RUNSPEC section.

There is no terminating ``/'' for this keyword, and thus the number
entries must match the value entered via the TABDIMS(NTPVT),
TABDIMS(NTROCC), or TABDIMS(NTSFUN) parameters, depending on the option
selected via the ROCKOPTS keyword.

\subsection{ROCK2D -- Pore Volume Compaction versus Pressure and Sw
Tables}\label{rock2d-pore-volume-compaction-versus-pressure-and-sw-tables}

\begin{longtable}[]{@{}
  >{\raggedright\arraybackslash}p{(\linewidth - 14\tabcolsep) * \real{0.1250}}
  >{\raggedright\arraybackslash}p{(\linewidth - 14\tabcolsep) * \real{0.1250}}
  >{\raggedright\arraybackslash}p{(\linewidth - 14\tabcolsep) * \real{0.1250}}
  >{\raggedright\arraybackslash}p{(\linewidth - 14\tabcolsep) * \real{0.1250}}
  >{\raggedright\arraybackslash}p{(\linewidth - 14\tabcolsep) * \real{0.1250}}
  >{\raggedright\arraybackslash}p{(\linewidth - 14\tabcolsep) * \real{0.1250}}
  >{\raggedright\arraybackslash}p{(\linewidth - 14\tabcolsep) * \real{0.1250}}
  >{\raggedright\arraybackslash}p{(\linewidth - 14\tabcolsep) * \real{0.1250}}@{}}
\toprule\noalign{}
\begin{minipage}[b]{\linewidth}\raggedright
RUNSPEC
\end{minipage} & \begin{minipage}[b]{\linewidth}\raggedright
GRID
\end{minipage} & \begin{minipage}[b]{\linewidth}\raggedright
EDIT
\end{minipage} & \begin{minipage}[b]{\linewidth}\raggedright
PROPS
\end{minipage} & \begin{minipage}[b]{\linewidth}\raggedright
REGIONS
\end{minipage} & \begin{minipage}[b]{\linewidth}\raggedright
SOLUTION
\end{minipage} & \begin{minipage}[b]{\linewidth}\raggedright
SUMMARY
\end{minipage} & \begin{minipage}[b]{\linewidth}\raggedright
SCHEDULE
\end{minipage} \\
\midrule\noalign{}
\endhead
\bottomrule\noalign{}
\endlastfoot
\end{longtable}

\subsubsection{Description}\label{description-606}

The ROCK2D keyword defines rock compressibility pore volume multipliers
as a function of pressure and water saturation (``Sw'') for when the
rock compaction option has been invoked by the ROCKCOMP keyword in the
RUNSPEC section. The pressure values are defined on this keyword and the
water saturations are declared on the associated ROCKWNOD keyword in the
PROPS section

The rock compaction pore volume and transmissibility multipliers,
entered via the ROCKTAB, ROCK2D and ROCK2DTR keywords, are applied to
the pore pressure, unless the OVERBURD keyword in PROPS section is
included in the input deck. When the OVERBURD keyword is present the
multipliers are applied to the effective pore volume pressure, that
is\({P}_{(\mathit{effective})} = {P}_{(\mathit{Pressure})} - {P}_{(\mathit{overburden})}\).
If the keyword is not present in the input deck then the overburden
pressure is set to zero.

This keyword should only be used if compaction option has been enabled.

\begin{longtable}[]{@{}
  >{\raggedright\arraybackslash}p{(\linewidth - 6\tabcolsep) * \real{0.2500}}
  >{\raggedright\arraybackslash}p{(\linewidth - 6\tabcolsep) * \real{0.2500}}
  >{\raggedright\arraybackslash}p{(\linewidth - 6\tabcolsep) * \real{0.2500}}
  >{\raggedright\arraybackslash}p{(\linewidth - 6\tabcolsep) * \real{0.2500}}@{}}
\toprule\noalign{}
\begin{minipage}[b]{\linewidth}\raggedright
No.
\end{minipage} & \begin{minipage}[b]{\linewidth}\raggedright
Name
\end{minipage} & \begin{minipage}[b]{\linewidth}\raggedright
Description
\end{minipage} & \begin{minipage}[b]{\linewidth}\raggedright
Default
\end{minipage} \\
\midrule\noalign{}
\endhead
\bottomrule\noalign{}
\endlastfoot
Field & Metric & Laboratory & \\
1 & PRESS & A columnar vector of real monotonically increasing down the
column values that defines the corresponding overburden pressure for the
subsequent MULT columnar vector. & None \\
psia & bars & atm & \\
2 & MULT & A columnar vector of real equal or decreasing down the column
values that are less than or equal to one, that defines the rock
compressibility pore volume multipliers corresponding to PRESS and for
each water saturation entry in the ROCKWNOD keyword. & None \\
dimensionless & dimensionless & dimensionless & \\
Notes: & & & \\
\end{longtable}

\emph{Table 8.126: ROCK2D Keyword Description}

See also the OVERBURD, ROCKTAB, ROCK2DTR, and ROCKWNOD keywords in the
PROPS section.

\subsubsection{Example}\label{example-337}

The following example defines two pore volume compaction tables,
assuming NTROCC is equal to two on the ROCKCOMP keyword and NSSFUN is
greater than or equal to four on the TABDIMS keyword.

\begin{verbatim}
--
--       ROCK COMPACTION VERSUS PRESSURE AND SW TABLES
--
ROCK2D
--       PRESS    PORV          FIRST ROCK2D TABLE DATA
--       PSIA     MULTIPLER
--       ------   ----------
            0.0      0.850
                     0.850
                     0.850
                     0.085                                / P-SW SET TABLE NO. 01
--       PRESS    PORV
--       PSIA     MULTIPLER
--       ------   ----------
         1000.0      0.900
                     0.900
                     0.900
                     0.900                                / P-SW SET TABLE NO. 01
--       PRESS    PORV
--       PSIA     MULTIPLER
--       ------   ----------
         2500.0      0.950
                     0.950
                     0.950
                     0.950                                / P-SW SET TABLE NO. 01
--       PRESS    PORV
--       PSIA     MULTIPLER
--       ------   ----------
         5000.0      1.000
                     1.000
                     1.000
                     1.000                                / P-SW SET TABLE NO. 01
--
--       PRESS    PORV          SECOND ROCK2D TABLE DATA
--       PSIA     MULTIPLER
--       ------   ----------
            0.0      0.800
                     0.800
                     0.800
                     0.800                                / P-SW SET TABLE NO. 02
--       PRESS    PORV
--       PSIA     MULTIPLER
--       ------   ----------
         1000.0      0.880
                     0.880
                     0.880
                     0.880                                / P-SW SET TABLE NO. 02
--       PRESS    PORV
--       PSIA     MULTIPLER
--       ------   ----------
         2500.0      0.950
                     0.950
                     0.950
                     0.950                                / P-SW SET TABLE NO. 02
--       PRESS    PORV
--       PSIA     MULTIPLER
--       ------   ----------
         5000.0      1.000
                     1.000
                     1.000
                     1.000                                / P-SW SET TABLE NO. 02
\end{verbatim}

Note that there must be exactly NTROCC tables entered for this keyword,
otherwise an error will occur.

\subsection{ROCK2DTR -- Transmissibility Compaction versus Pressure and
Sw
Tables}\label{rock2dtr-transmissibility-compaction-versus-pressure-and-sw-tables}

\begin{longtable}[]{@{}
  >{\raggedright\arraybackslash}p{(\linewidth - 14\tabcolsep) * \real{0.1250}}
  >{\raggedright\arraybackslash}p{(\linewidth - 14\tabcolsep) * \real{0.1250}}
  >{\raggedright\arraybackslash}p{(\linewidth - 14\tabcolsep) * \real{0.1250}}
  >{\raggedright\arraybackslash}p{(\linewidth - 14\tabcolsep) * \real{0.1250}}
  >{\raggedright\arraybackslash}p{(\linewidth - 14\tabcolsep) * \real{0.1250}}
  >{\raggedright\arraybackslash}p{(\linewidth - 14\tabcolsep) * \real{0.1250}}
  >{\raggedright\arraybackslash}p{(\linewidth - 14\tabcolsep) * \real{0.1250}}
  >{\raggedright\arraybackslash}p{(\linewidth - 14\tabcolsep) * \real{0.1250}}@{}}
\toprule\noalign{}
\begin{minipage}[b]{\linewidth}\raggedright
RUNSPEC
\end{minipage} & \begin{minipage}[b]{\linewidth}\raggedright
GRID
\end{minipage} & \begin{minipage}[b]{\linewidth}\raggedright
EDIT
\end{minipage} & \begin{minipage}[b]{\linewidth}\raggedright
PROPS
\end{minipage} & \begin{minipage}[b]{\linewidth}\raggedright
REGIONS
\end{minipage} & \begin{minipage}[b]{\linewidth}\raggedright
SOLUTION
\end{minipage} & \begin{minipage}[b]{\linewidth}\raggedright
SUMMARY
\end{minipage} & \begin{minipage}[b]{\linewidth}\raggedright
SCHEDULE
\end{minipage} \\
\midrule\noalign{}
\endhead
\bottomrule\noalign{}
\endlastfoot
\end{longtable}

\subsubsection{Description}\label{description-607}

The ROCK2DTR keyword defines rock compressibility transmissibility
multipliers as a function of pressure and water saturation (``Sw'') for
when the rock compaction option has been invoked by the ROCKCOMP keyword
in the RUNSPEC section. The pressure values are defined on this keyword
and the water saturations are declared on the associated ROCKWNOD
keyword in the PROPS section

The rock compaction pore volume and transmissibility multipliers,
entered via the ROCKTAB, ROCK2D and ROCK2DTR keywords, are applied to
the pore pressure, unless the OVERBURD keyword in PROPS section is
included in the input deck. When the OVERBURD keyword is present the
multipliers are applied to the effective pore volume pressure, that
is\({P}_{(\mathit{effective})} = {P}_{(\mathit{Pressure})} - {P}_{(\mathit{overburden})}\).
If the keyword is not present in the input deck then the overburden
pressure is set to zero.

This keyword should only be used if compaction option has been enabled.

\begin{longtable}[]{@{}
  >{\raggedright\arraybackslash}p{(\linewidth - 6\tabcolsep) * \real{0.2500}}
  >{\raggedright\arraybackslash}p{(\linewidth - 6\tabcolsep) * \real{0.2500}}
  >{\raggedright\arraybackslash}p{(\linewidth - 6\tabcolsep) * \real{0.2500}}
  >{\raggedright\arraybackslash}p{(\linewidth - 6\tabcolsep) * \real{0.2500}}@{}}
\toprule\noalign{}
\begin{minipage}[b]{\linewidth}\raggedright
No.
\end{minipage} & \begin{minipage}[b]{\linewidth}\raggedright
Name
\end{minipage} & \begin{minipage}[b]{\linewidth}\raggedright
Description
\end{minipage} & \begin{minipage}[b]{\linewidth}\raggedright
Default
\end{minipage} \\
\midrule\noalign{}
\endhead
\bottomrule\noalign{}
\endlastfoot
Field & Metric & Laboratory & \\
1 & PRESS & A columnar vector of real monotonically increasing down the
column values that defines the corresponding overburden pressure for the
subsequent MULT columnar vector. & None \\
psia & bars & atm & \\
2 & MULT & A columnar vector of real equal or decreasing down the column
values that are less than or equal to one, that defines the rock
compressibility transmissibility multipliers corresponding to PRESS and
for each water saturation entry in the ROCKWNOD keyword. & None \\
dimensionless & dimensionless & dimensionless & \\
Notes: & & & \\
\end{longtable}

\emph{Table 8.127: ROCK2DTR Keyword Description}

See also the OVERBURD, ROCKTAB, ROCK2D, and ROCKWNOD keywords in the
PROPS section.

\subsubsection{Example}\label{example-338}

The following example defines two rock compressibility transmissibility
compaction tables, assuming NTROCC is equal to two on the ROCKCOMP
keyword and NSSFUN is greater than or equal to four on the TABDIMS
keyword.

\begin{verbatim}
--
--       TRANSMISSIBILITY COMPACTION VERSUS PRESSURE AND SW TABLES
--
ROCK2DTR
--       PRESS    TRAN          FIRST ROCK2DTR TABLE DATA
--       PSIA     MULTIPLER
--       ------   ----------
            0.0      0.850
                     0.850
                     0.850
                     0.085                                / P-SW SET TABLE NO. 01
--       PRESS    TRAN
--       PSIA     MULTIPLER
--       ------   ----------
         1000.0      0.900
                     0.900
                     0.900
                     0.900                                / P-SW SET TABLE NO. 01
--       PRESS    TRAN
--       PSIA     MULTIPLER
--       ------   ----------
         2500.0      0.950
                     0.950
                     0.950
                     0.950                                / P-SW SET TABLE NO. 01
--       PRESS    TRAN
--       PSIA     MULTIPLER
--       ------   ----------
         5000.0      1.000
                     1.000
                     1.000
                     1.000                                / P-SW SET TABLE NO. 01
--
--       PRESS    TRAN          SECOND ROCK2DTR TABLE DATA
--       PSIA     MULTIPLER
--       ------   ----------
            0.0      0.800
                     0.800
                     0.800
                     0.800                                / P-SW SET TABLE NO. 02
--       PRESS    TRAN
--       PSIA     MULTIPLER
--       ------   ----------
         1000.0      0.880
                     0.880
                     0.880
                     0.880                                / P-SW SET TABLE NO. 02
--       PRESS    TRAN
--       PSIA     MULTIPLER
--       ------   ----------
         2500.0      0.950
                     0.950
                     0.950
                     0.950                                / P-SW SET TABLE NO. 02
--       PRESS    TRAN
--       PSIA     MULTIPLER
--       ------   ----------
         5000.0      1.000
                     1.000
                     1.000
                     1.000                                / P-SW SET TABLE NO. 02
\end{verbatim}

Note that there must be exactly NTROCC tables entered for this keyword,
otherwise an error will occur.

\subsection{ROCKOPTS -- Define Rock Compaction and Compressibility
Options}\label{rockopts-define-rock-compaction-and-compressibility-options}

\begin{longtable}[]{@{}
  >{\raggedright\arraybackslash}p{(\linewidth - 14\tabcolsep) * \real{0.1250}}
  >{\raggedright\arraybackslash}p{(\linewidth - 14\tabcolsep) * \real{0.1250}}
  >{\raggedright\arraybackslash}p{(\linewidth - 14\tabcolsep) * \real{0.1250}}
  >{\raggedright\arraybackslash}p{(\linewidth - 14\tabcolsep) * \real{0.1250}}
  >{\raggedright\arraybackslash}p{(\linewidth - 14\tabcolsep) * \real{0.1250}}
  >{\raggedright\arraybackslash}p{(\linewidth - 14\tabcolsep) * \real{0.1250}}
  >{\raggedright\arraybackslash}p{(\linewidth - 14\tabcolsep) * \real{0.1250}}
  >{\raggedright\arraybackslash}p{(\linewidth - 14\tabcolsep) * \real{0.1250}}@{}}
\toprule\noalign{}
\begin{minipage}[b]{\linewidth}\raggedright
RUNSPEC
\end{minipage} & \begin{minipage}[b]{\linewidth}\raggedright
GRID
\end{minipage} & \begin{minipage}[b]{\linewidth}\raggedright
EDIT
\end{minipage} & \begin{minipage}[b]{\linewidth}\raggedright
PROPS
\end{minipage} & \begin{minipage}[b]{\linewidth}\raggedright
REGIONS
\end{minipage} & \begin{minipage}[b]{\linewidth}\raggedright
SOLUTION
\end{minipage} & \begin{minipage}[b]{\linewidth}\raggedright
SUMMARY
\end{minipage} & \begin{minipage}[b]{\linewidth}\raggedright
SCHEDULE
\end{minipage} \\
\midrule\noalign{}
\endhead
\bottomrule\noalign{}
\endlastfoot
\end{longtable}

\subsubsection{Description}\label{description-608}

The ROCKOPTS keyword defines various options with respect to rock
compaction and rock compressibility.

\begin{longtable}[]{@{}
  >{\raggedright\arraybackslash}p{(\linewidth - 6\tabcolsep) * \real{0.2500}}
  >{\raggedright\arraybackslash}p{(\linewidth - 6\tabcolsep) * \real{0.2500}}
  >{\raggedright\arraybackslash}p{(\linewidth - 6\tabcolsep) * \real{0.2500}}
  >{\raggedright\arraybackslash}p{(\linewidth - 6\tabcolsep) * \real{0.2500}}@{}}
\toprule\noalign{}
\begin{minipage}[b]{\linewidth}\raggedright
No.
\end{minipage} & \begin{minipage}[b]{\linewidth}\raggedright
Name
\end{minipage} & \begin{minipage}[b]{\linewidth}\raggedright
Description
\end{minipage} & \begin{minipage}[b]{\linewidth}\raggedright
Default
\end{minipage} \\
\midrule\noalign{}
\endhead
\bottomrule\noalign{}
\endlastfoot
1 & ROCKOPT1 & A defined character string that specifies how the
overburden pressures supplied by the OVERBURD keyword are applied to the
tabulated pressures in the ROCKTAB keywords: ROCKOPT1 should be set to
PRESSURE if the OVERBURD is not used in the input deck. Only the default
value of PRESSURE is supported. & PRESSURE \\
2 & ROCKOPT2 & A defined character string that sets the reference
pressure option: Note that STORE option should not be used with the
OVERBURD keywords as the OVERBURD data will be overwritten. & NOSTORE \\
3 & ROCKOPT3 & A defined character string that specifies which region
array should be used to allocate the various ROCK and ROCKTAB property
tables in the model: Only the PVTNUM and ROCKNUM options are currently
supported. & PVTNUM \\
4 & ROCKOPT4 & A defined character string that sets the initial
conditions for the HYSTER and BOBERG options: This parameter is ignored
by OPM Flow as the ROCKCOMP(ROCKOPT) options of HYSTER and BOBERG are
not supported by the simulator. & DEFLATION \\
Notes: & & & \\
\end{longtable}

\emph{Table 8.128: ROCKOPTS Keyword Description}

\subsubsection{Example}\label{example-339}

\begin{verbatim}
--
--       ROCKOPT1  ROCKOPT2   ROCKOPT3  ROCKOPT4
--       PRS/STRE  NO/STORE   ARRAY
--       --------  --------   --------  --------
ROCKOPTS
         PRESSURE  NOSTORE    PVTNUM                       / ROCK COMP OPTIONS
\end{verbatim}

The above example defines the default values for the ROCKOPTS keyword.

\subsection{ROCKPAMA -- Define Coal Palmer-Mansorri Rock Model
Parameters}\label{rockpama-define-coal-palmer-mansorri-rock-model-parameters}

\begin{longtable}[]{@{}
  >{\raggedright\arraybackslash}p{(\linewidth - 14\tabcolsep) * \real{0.1250}}
  >{\raggedright\arraybackslash}p{(\linewidth - 14\tabcolsep) * \real{0.1250}}
  >{\raggedright\arraybackslash}p{(\linewidth - 14\tabcolsep) * \real{0.1250}}
  >{\raggedright\arraybackslash}p{(\linewidth - 14\tabcolsep) * \real{0.1250}}
  >{\raggedright\arraybackslash}p{(\linewidth - 14\tabcolsep) * \real{0.1250}}
  >{\raggedright\arraybackslash}p{(\linewidth - 14\tabcolsep) * \real{0.1250}}
  >{\raggedright\arraybackslash}p{(\linewidth - 14\tabcolsep) * \real{0.1250}}
  >{\raggedright\arraybackslash}p{(\linewidth - 14\tabcolsep) * \real{0.1250}}@{}}
\toprule\noalign{}
\begin{minipage}[b]{\linewidth}\raggedright
RUNSPEC
\end{minipage} & \begin{minipage}[b]{\linewidth}\raggedright
GRID
\end{minipage} & \begin{minipage}[b]{\linewidth}\raggedright
EDIT
\end{minipage} & \begin{minipage}[b]{\linewidth}\raggedright
PROPS
\end{minipage} & \begin{minipage}[b]{\linewidth}\raggedright
REGIONS
\end{minipage} & \begin{minipage}[b]{\linewidth}\raggedright
SOLUTION
\end{minipage} & \begin{minipage}[b]{\linewidth}\raggedright
SUMMARY
\end{minipage} & \begin{minipage}[b]{\linewidth}\raggedright
SCHEDULE
\end{minipage} \\
\midrule\noalign{}
\endhead
\bottomrule\noalign{}
\endlastfoot
\end{longtable}

\subsubsection{Description}\label{description-609}

ROCKPAMA defines the Palmer-Mansoori {[}Palmer, I. and Mansoori, J.
``How Permeability Depends on Stress and Pore Pressure in Coalbeds: A
New Model,'' paper SPE 52607, SPE Reservoir Evaluation \& Engineering
(1998) 1, No.~6, 539-544.{]} and {[}Clarkson, C.R., Pan, Z., Palmer, I.
and Harpalani, S. ``Predicting Sorption-Induced Strain and Permeability
Increase With Depletion for Coalbed-Methane Reservoirs'', SPE 114778-PA,
SPE Journal (2010) 15, No.~1, 152--159.{]} parameters used for this rock
model, for when the Coal Bed Methane option has been activated via the
COAL keyword, and PALM-MAN has been declared for the ROCKOPT variable on
the ROCKCOMP keyword; both keywords are in the RUNSPEC section. The
Palmer-Mansoori rock model is used to calculate the impact on pore
volume and permeability due to rock compaction.

This keyword is not supported by OPM Flow but would change the results
if supported so the simulation will be stopped.

\subsection{ROCKTAB -- Rock Compaction
Tables}\label{rocktab-rock-compaction-tables}

\begin{longtable}[]{@{}
  >{\raggedright\arraybackslash}p{(\linewidth - 14\tabcolsep) * \real{0.1250}}
  >{\raggedright\arraybackslash}p{(\linewidth - 14\tabcolsep) * \real{0.1250}}
  >{\raggedright\arraybackslash}p{(\linewidth - 14\tabcolsep) * \real{0.1250}}
  >{\raggedright\arraybackslash}p{(\linewidth - 14\tabcolsep) * \real{0.1250}}
  >{\raggedright\arraybackslash}p{(\linewidth - 14\tabcolsep) * \real{0.1250}}
  >{\raggedright\arraybackslash}p{(\linewidth - 14\tabcolsep) * \real{0.1250}}
  >{\raggedright\arraybackslash}p{(\linewidth - 14\tabcolsep) * \real{0.1250}}
  >{\raggedright\arraybackslash}p{(\linewidth - 14\tabcolsep) * \real{0.1250}}@{}}
\toprule\noalign{}
\begin{minipage}[b]{\linewidth}\raggedright
RUNSPEC
\end{minipage} & \begin{minipage}[b]{\linewidth}\raggedright
GRID
\end{minipage} & \begin{minipage}[b]{\linewidth}\raggedright
EDIT
\end{minipage} & \begin{minipage}[b]{\linewidth}\raggedright
PROPS
\end{minipage} & \begin{minipage}[b]{\linewidth}\raggedright
REGIONS
\end{minipage} & \begin{minipage}[b]{\linewidth}\raggedright
SOLUTION
\end{minipage} & \begin{minipage}[b]{\linewidth}\raggedright
SUMMARY
\end{minipage} & \begin{minipage}[b]{\linewidth}\raggedright
SCHEDULE
\end{minipage} \\
\midrule\noalign{}
\endhead
\bottomrule\noalign{}
\endlastfoot
\end{longtable}

\subsubsection{Description}\label{description-610}

The ROCKTAB keyword defines the rock compaction attributes to be applied
when the rock compaction option has been invoked by the ROCKCOMP keyword
in the RUNSPEC section. ROCKTAB defines pore volume and transmissibility
multipliers versus pressure that are used in the compaction
calculations. If the RKTRMDIR has been activated in the PROPS section,
then the transmissibility multiplier is directional dependent and two
additional columns are used to define the y and z direction
transmissibility multipliers.

\begin{longtable}[]{@{}
  >{\raggedright\arraybackslash}p{(\linewidth - 6\tabcolsep) * \real{0.2500}}
  >{\raggedright\arraybackslash}p{(\linewidth - 6\tabcolsep) * \real{0.2500}}
  >{\raggedright\arraybackslash}p{(\linewidth - 6\tabcolsep) * \real{0.2500}}
  >{\raggedright\arraybackslash}p{(\linewidth - 6\tabcolsep) * \real{0.2500}}@{}}
\toprule\noalign{}
\begin{minipage}[b]{\linewidth}\raggedright
No.
\end{minipage} & \begin{minipage}[b]{\linewidth}\raggedright
Name
\end{minipage} & \begin{minipage}[b]{\linewidth}\raggedright
Description
\end{minipage} & \begin{minipage}[b]{\linewidth}\raggedright
Default
\end{minipage} \\
\midrule\noalign{}
\endhead
\bottomrule\noalign{}
\endlastfoot
Field & Metric & Laboratory & \\
1 & PRESS & If the ROCKOPT1 variable has been set to PRESSURE on the
ROCKOPTS keyword in the PROPS section, then PRESS should be a columnar
vector of real monotonically increasing down the column values, that
define the reference pressure for which the other parameters correspond
to. If ROCKOPT1 has been set to STRESS, then PRESS should be a columnar
vector of real monotonically decreasing down the column values. &
None \\
psia & bars & atm & \\
2 & PORV & A columnar vector of real positive values that are either
equal or increasing down the column that define the rock pore volume
multiplier for a given PRESS. & None \\
dimensionless & dimensionless & dimensionless & \\
3 & TRANS & If the RKTRMDIR is absent from the input deck, then TRANS is
a columnar vector of real positive values that are either equal or
increasing down the column that define the x, y, and z directional
transmissibility multipliers for the corresponding PRESS. If the
RKTRMDIR is present in the input deck, then TRANS is a columnar vector
of real positive values that are either equal or increasing down the
column that define only the x directional transmissibility multipliers
for the corresponding PRESS. & None \\
dimensionless & dimensionless & dimensionless & \\
4 & TRANSY & If the RKTRMDIR is absent from the input deck, then TRANSY
is ignored. If the RKTRMDIR is present in the input deck, then TRANSY is
a columnar vector of real positive values that are either equal or
increasing down the column that define only the y directional
transmissibility multipliers for the corresponding PRESS. & None \\
dimensionless & dimensionless & dimensionless & \\
5 & TRANSZ & If the RKTRMDIR is absent from the input deck, then TRANSZ
is ignored. If the RKTRMDIR is present in the input deck, then TRANSZ is
a columnar vector of real positive values that are either equal or
increasing down the column that define only the z directional
transmissibility multipliers for the corresponding PRESS. & None \\
dimensionless & dimensionless & dimensionless & \\
Notes: & & & \\
\end{longtable}

\emph{Table 8.129: ROCKTAB Keyword Description}

\subsubsection{Examples}\label{examples-41}

The example below defines two rock compaction tables, assuming NTROCC is
equal to two on the ROCKCOMP keyword and NPPVT is greater than or equal
to five on the TABDIMS keyword and that the RKTRMDIR keyword is present
in the input deck.

\begin{verbatim}
--
--       ROCK COMPACTION TABLES
--
ROCKTAB
--       PRESS    PORV     TX(YZ)   TY       TZ
--                MULT     MULT     MULT     MULT
--       ------   ------   ------   ------   ------
         1000.0   0.9600   0.9650   0.9650   0.9650
         1500.0   0.9800   0.9850   0.9850   0.9500
         3000.0   0.9900   0.9950   0.9950   0.9950
         4500.0   1.0000   1.0000   1.0000   1.0000
         4750.0   1.0100   1.0100   1.0100   1.0100        / TABLE NO. 01
--       PRESS    PORV     TX(YZ)   TY       TZ
--                MULT     MULT     MULT     MULT
--       ------   ------   ------   ------   ------
         1000.0   0.9600   0.9650   0.9650   0.9650
         1500.0   0.9800   0.9850   0.9850   0.9500
         3000.0   0.9900   0.9950   0.9950   0.9950
         4500.0   1.0000   1.0000   1.0000   1.0000
         4750.0   1.0100   1.0100   1.0100   1.0100        / TABLE NO. 02
\end{verbatim}

As the x, y and z directional transmissibility multipliers are identical
in the above example, we could eliminate the RKTRMDIR keyword from the
input deck and enter the data in the three column format, as shown on
the next page.

\begin{verbatim}
--
--       ROCK COMPACTION TABLES
--
ROCKTAB
--       PRESS    PORV     TX(YZ)
--                MULT     MULT
--       ------   ------   ------
         1000.0   0.9600   0.9650
         1500.0   0.9800   0.9850
         3000.0   0.9900   0.9950
         4500.0   1.0000   1.0000
         4750.0   1.0100   1.0100                          / TABLE NO. 01
--       PRESS    PORV     TX(YZ)
--                MULT     MULT
--       ------   ------   ------
         1000.0   0.9600   0.9650
         1500.0   0.9800   0.9850
         3000.0   0.9900   0.9950
         4500.0   1.0000   1.0000
         4750.0   1.0100   1.0100                          / TABLE NO. 02
\end{verbatim}

The net result of the two examples in this case is identical.

\subsection{ROCKTABH -- Rock Compaction Hysteresis
Tables}\label{rocktabh-rock-compaction-hysteresis-tables}

\begin{longtable}[]{@{}
  >{\raggedright\arraybackslash}p{(\linewidth - 14\tabcolsep) * \real{0.1250}}
  >{\raggedright\arraybackslash}p{(\linewidth - 14\tabcolsep) * \real{0.1250}}
  >{\raggedright\arraybackslash}p{(\linewidth - 14\tabcolsep) * \real{0.1250}}
  >{\raggedright\arraybackslash}p{(\linewidth - 14\tabcolsep) * \real{0.1250}}
  >{\raggedright\arraybackslash}p{(\linewidth - 14\tabcolsep) * \real{0.1250}}
  >{\raggedright\arraybackslash}p{(\linewidth - 14\tabcolsep) * \real{0.1250}}
  >{\raggedright\arraybackslash}p{(\linewidth - 14\tabcolsep) * \real{0.1250}}
  >{\raggedright\arraybackslash}p{(\linewidth - 14\tabcolsep) * \real{0.1250}}@{}}
\toprule\noalign{}
\begin{minipage}[b]{\linewidth}\raggedright
RUNSPEC
\end{minipage} & \begin{minipage}[b]{\linewidth}\raggedright
GRID
\end{minipage} & \begin{minipage}[b]{\linewidth}\raggedright
EDIT
\end{minipage} & \begin{minipage}[b]{\linewidth}\raggedright
PROPS
\end{minipage} & \begin{minipage}[b]{\linewidth}\raggedright
REGIONS
\end{minipage} & \begin{minipage}[b]{\linewidth}\raggedright
SOLUTION
\end{minipage} & \begin{minipage}[b]{\linewidth}\raggedright
SUMMARY
\end{minipage} & \begin{minipage}[b]{\linewidth}\raggedright
SCHEDULE
\end{minipage} \\
\midrule\noalign{}
\endhead
\bottomrule\noalign{}
\endlastfoot
\end{longtable}

\subsubsection{Description}\label{description-611}

The ROCKTABH keyword defines the rock compaction hysteresis attributes
to be applied for when the rock compaction option has been invoked by
the ROCKCOMP keyword in the RUNSPEC section. ROCKTABH defines pore
volume and transmissibility multipliers versus pressure that are used in
the compaction calculations. If the RKTRMDIR has been activated in the
PROPS section, then the transmissibility multiplier is directional
dependent and two additional columns are used to define the y and z
direction transmissibility multipliers. The keyword should only be used
if the Rock Compaction Hysteresis option has been activated by either
setting the ROCKOPT parameter on the ROCKCOMP keyword to HYSTER or
BOBERG.

This keyword is not supported by OPM Flow but would change the results
if supported so the simulation will be stopped.

\begin{longtable}[]{@{}
  >{\raggedright\arraybackslash}p{(\linewidth - 6\tabcolsep) * \real{0.2500}}
  >{\raggedright\arraybackslash}p{(\linewidth - 6\tabcolsep) * \real{0.2500}}
  >{\raggedright\arraybackslash}p{(\linewidth - 6\tabcolsep) * \real{0.2500}}
  >{\raggedright\arraybackslash}p{(\linewidth - 6\tabcolsep) * \real{0.2500}}@{}}
\toprule\noalign{}
\begin{minipage}[b]{\linewidth}\raggedright
No.
\end{minipage} & \begin{minipage}[b]{\linewidth}\raggedright
Name
\end{minipage} & \begin{minipage}[b]{\linewidth}\raggedright
Description
\end{minipage} & \begin{minipage}[b]{\linewidth}\raggedright
Default
\end{minipage} \\
\midrule\noalign{}
\endhead
\bottomrule\noalign{}
\endlastfoot
Field & Metric & Laboratory & \\
1 & PRESS & If the ROCKOPT1 variable has been set to PRESSURE on the
ROCKOPTS keyword in the PROPS section, then PRESS should be a columnar
vector of real monotonically increasing down the column values, that
define the reference pressure for which the other parameters correspond
to. If ROCKOPT1 has been set to STRESS, then PRESS should be a columnar
vector of real monotonically decreasing down the column values. &
None \\
psia & bars & atm & \\
2 & PORV & A columnar vector of real positive values that are either
equal or increasing down the column that define the rock pore volume
multiplier for a given PRESS. & None \\
dimensionless & dimensionless & dimensionless & \\
3 & TRANS & If the RKTRMDIR is absent from the input deck, then TRANS is
a columnar vector of real positive values that are either equal or
increasing down the column that define the x, y, and z directional
transmissibility multipliers for the corresponding PRESS. If the
RKTRMDIR is present in the input deck, then TRANS is a columnar vector
of real positive values that are either equal or increasing down the
column that define only the x directional transmissibility multipliers
for the corresponding PRESS. & None \\
dimensionless & dimensionless & dimensionless & \\
4 & TRANSY & If the RKTRMDIR is absent from the input deck, then TRANSY
is ignored. If the RKTRMDIR is present in the input deck, then TRANSY is
a columnar vector of real positive values that are either equal or
increasing down the column that define only the y directional
transmissibility multipliers for the corresponding PRESS. & None \\
dimensionless & dimensionless & dimensionless & \\
5 & TRANSZ & If the RKTRMDIR is absent from the input deck, then TRANSZ
is ignored. If the RKTRMDIR is present in the input deck, then TRANSZ is
a columnar vector of real positive values that are either equal or
increasing down the column that define only the z directional
transmissibility multipliers for the corresponding PRESS. & None \\
dimensionless & dimensionless & dimensionless & \\
Notes: & & & \\
\end{longtable}

\emph{Table 8.130: ROCKTABH Keyword Description}

Each data set consists of columnar vectors of pore volume and
transmissibility multipliers versus pressure that specify the elastic
contraction and expansion and of the reservoir rock. The deflation curve
is derived from the first data elements on each elastic curve. If the
ROCKOPT parameter on the ROCKCOMP keyword has been set to HYSTER, then
the dilation curves are extrapolated to infinite pressure, that is the
curves are unbounded. However, if ROCKCOMP is set to BOBERG the last
points of each elastic curve are used as the dilation curves.

\subsubsection{Example}\label{example-340}

The example below defines two rock compaction tables, assuming NTROCC is
equal to two on the ROCKCOMP keyword and NPPVT is greater than or equal
to four on the TABDIMS keyword and that the RKTRMDIR keyword is not
present in the input deck.

\begin{verbatim}
--
--       ROCK COMPACTION HYSTERESIS TABLES
--
ROCKTABH
--       PRESS    PORV     TX(YZ)   TY       TZ
--                MULT     MULT     MULT     MULT
--       ------   ------   ------   ------   ------
         1500.0   0.9600   0.9800
         2500.0   0.9700   0.9850
         3500.0   0.9800   0.9900
         4500.0   0.9900   0.9950                          / NPPVT = 1
         2500.0   0.9900   0.9900
         3500.0   0.9950   0.9950
         4750.0   0.9980   0.9980                          / NPPVT = 2
         3500.0   1.0000   1.0000
         5500.0   1.0100   1.0100                          / NPPVT = 3
         4500.0   1.0100   1.0100
         5750.0   1.0200   1.0200                          / NPPVT = 4

                                                           / TABLE NO. 01       --       PRESS    PORV     TX(YZ)   TY       TZ
--                MULT     MULT     MULT     MULT
--       ------   ------   ------   ------   ------
         1500.0   0.9400   0.9700
         2750.0   0.9400   0.9700                          / NPPVT = 1
         2250.0   0.9800   0.9900
         3250.0   0.9800   0.9900                          / NPPVT = 2
         3000.0   1.0000   1.0000
         4250.0   1.0000   1.0000                          / NPPVT = 3
         4550.0   1.0200   1.0100
         5750.0   1.0200   1.0100                          / NPPVT = 4
                                                           / TABLE NO. 02
\end{verbatim}

Here the deflation curve is define for table number one is:

\begin{verbatim}
1500.0   0.9600   0.9800
2500.0   0.9900   0.9900
3500.0   1.0000   1.0000
4500.0   1.0100   1.0100
\end{verbatim}

and for table number 2:

\begin{verbatim}
1500.0   0.9400   0.9700
2250.0   0.9800   0.9900
3250.0   1.0000   1.0000
4250.0   1.0200   1.0100
\end{verbatim}

And the dilation curve is define for table number one is:

\begin{verbatim}
4500.0   0.9900   0.9950
4750.0   0.9980   0.9980
5500.0   1.0100   1.0100
5750.0   1.0200   1.0200
\end{verbatim}

and for table number 2:

\begin{verbatim}
2250.0   0.9400   0.9700
3250.0   0.9800   0.9900
4250.0   1.0000   1.0000
5250.0   1.0200   1.0100
\end{verbatim}

\subsection{ROCKTABW -- Rock Compaction Tables (Water
Induced)}\label{rocktabw-rock-compaction-tables-water-induced}

\begin{longtable}[]{@{}
  >{\raggedright\arraybackslash}p{(\linewidth - 14\tabcolsep) * \real{0.1250}}
  >{\raggedright\arraybackslash}p{(\linewidth - 14\tabcolsep) * \real{0.1250}}
  >{\raggedright\arraybackslash}p{(\linewidth - 14\tabcolsep) * \real{0.1250}}
  >{\raggedright\arraybackslash}p{(\linewidth - 14\tabcolsep) * \real{0.1250}}
  >{\raggedright\arraybackslash}p{(\linewidth - 14\tabcolsep) * \real{0.1250}}
  >{\raggedright\arraybackslash}p{(\linewidth - 14\tabcolsep) * \real{0.1250}}
  >{\raggedright\arraybackslash}p{(\linewidth - 14\tabcolsep) * \real{0.1250}}
  >{\raggedright\arraybackslash}p{(\linewidth - 14\tabcolsep) * \real{0.1250}}@{}}
\toprule\noalign{}
\begin{minipage}[b]{\linewidth}\raggedright
RUNSPEC
\end{minipage} & \begin{minipage}[b]{\linewidth}\raggedright
GRID
\end{minipage} & \begin{minipage}[b]{\linewidth}\raggedright
EDIT
\end{minipage} & \begin{minipage}[b]{\linewidth}\raggedright
PROPS
\end{minipage} & \begin{minipage}[b]{\linewidth}\raggedright
REGIONS
\end{minipage} & \begin{minipage}[b]{\linewidth}\raggedright
SOLUTION
\end{minipage} & \begin{minipage}[b]{\linewidth}\raggedright
SUMMARY
\end{minipage} & \begin{minipage}[b]{\linewidth}\raggedright
SCHEDULE
\end{minipage} \\
\midrule\noalign{}
\endhead
\bottomrule\noalign{}
\endlastfoot
\end{longtable}

\subsubsection{Description}\label{description-612}

The ROCKTABW keyword defines the rock compaction tables induced by
increasing water saturation within a grid cell due to water invasion,
for when the rock compaction option has been invoked by the ROCKCOMP
keyword in the RUNSPEC section. ROCKTABW defines pore volume and
transmissibility multipliers versus water saturation that are used in
the compaction calculations. The keyword should be used together with
the ROCK, ROCKTAB or ROCKTABH keywords that specify the pore volume and
transmissibility multipliers as functions of pressure. Alternatively the
ROCKWNOD, ROCK2D and ROCK2DTR keywords can be used to enter two
dimensional tables of the data. All keywords are in the PROPS section.

This keyword is not supported by OPM Flow but would change the results
if supported so the simulation will be stopped.

\subsection{ROCKTHSG -- Rock Compaction Hysteresis Tables (Dual
Porosity)}\label{rockthsg-rock-compaction-hysteresis-tables-dual-porosity}

\begin{longtable}[]{@{}
  >{\raggedright\arraybackslash}p{(\linewidth - 14\tabcolsep) * \real{0.1250}}
  >{\raggedright\arraybackslash}p{(\linewidth - 14\tabcolsep) * \real{0.1250}}
  >{\raggedright\arraybackslash}p{(\linewidth - 14\tabcolsep) * \real{0.1250}}
  >{\raggedright\arraybackslash}p{(\linewidth - 14\tabcolsep) * \real{0.1250}}
  >{\raggedright\arraybackslash}p{(\linewidth - 14\tabcolsep) * \real{0.1250}}
  >{\raggedright\arraybackslash}p{(\linewidth - 14\tabcolsep) * \real{0.1250}}
  >{\raggedright\arraybackslash}p{(\linewidth - 14\tabcolsep) * \real{0.1250}}
  >{\raggedright\arraybackslash}p{(\linewidth - 14\tabcolsep) * \real{0.1250}}@{}}
\toprule\noalign{}
\begin{minipage}[b]{\linewidth}\raggedright
RUNSPEC
\end{minipage} & \begin{minipage}[b]{\linewidth}\raggedright
GRID
\end{minipage} & \begin{minipage}[b]{\linewidth}\raggedright
EDIT
\end{minipage} & \begin{minipage}[b]{\linewidth}\raggedright
PROPS
\end{minipage} & \begin{minipage}[b]{\linewidth}\raggedright
REGIONS
\end{minipage} & \begin{minipage}[b]{\linewidth}\raggedright
SOLUTION
\end{minipage} & \begin{minipage}[b]{\linewidth}\raggedright
SUMMARY
\end{minipage} & \begin{minipage}[b]{\linewidth}\raggedright
SCHEDULE
\end{minipage} \\
\midrule\noalign{}
\endhead
\bottomrule\noalign{}
\endlastfoot
\end{longtable}

\subsubsection{Description}\label{description-613}

The ROCKTHSG keyword defines the rock compaction hysteresis attributes
to be applied for when the rock compaction option has been invoked by
the ROCKCOMP keyword in the RUNSPEC section and the either the Dual
Permeability or Dual Porosity models are activated by the DUALPERM and
DUALPORO keywords in the RUNSPEC section. ROCKTHSG specifies sigma
multipliers versus pressure that are used in the dual porosity rock
compaction calculations. The keyword should only be used if the Rock
Compaction Hysteresis option has been activated by either setting the
ROCKOPT parameter on the ROCKCOMP keyword to one of the available
options.

Each data set consists of columnar vectors of sigma multipliers versus
pressure that specify the elastic contraction and expansion and of the
reservoir rock. The deflation curve is derived from the first data
elements on each elastic curve. If the ROCKOPT parameter on the ROCKCOMP
keyword has been set to HYSTER, then the dilation curves are
extrapolated to infinite pressure, that is the curves are unbounded.
However, if ROCKCOMP is set to BOBERG the last points of each elastic
curve are used as the dilation curves.

This keyword is not supported by OPM Flow but would change the results
if supported so the simulation will be stopped.

\subsection{ROCKTSIG -- Rock Compaction Tables (Dual
Porosity)}\label{rocktsig-rock-compaction-tables-dual-porosity}

\begin{longtable}[]{@{}
  >{\raggedright\arraybackslash}p{(\linewidth - 14\tabcolsep) * \real{0.1250}}
  >{\raggedright\arraybackslash}p{(\linewidth - 14\tabcolsep) * \real{0.1250}}
  >{\raggedright\arraybackslash}p{(\linewidth - 14\tabcolsep) * \real{0.1250}}
  >{\raggedright\arraybackslash}p{(\linewidth - 14\tabcolsep) * \real{0.1250}}
  >{\raggedright\arraybackslash}p{(\linewidth - 14\tabcolsep) * \real{0.1250}}
  >{\raggedright\arraybackslash}p{(\linewidth - 14\tabcolsep) * \real{0.1250}}
  >{\raggedright\arraybackslash}p{(\linewidth - 14\tabcolsep) * \real{0.1250}}
  >{\raggedright\arraybackslash}p{(\linewidth - 14\tabcolsep) * \real{0.1250}}@{}}
\toprule\noalign{}
\begin{minipage}[b]{\linewidth}\raggedright
RUNSPEC
\end{minipage} & \begin{minipage}[b]{\linewidth}\raggedright
GRID
\end{minipage} & \begin{minipage}[b]{\linewidth}\raggedright
EDIT
\end{minipage} & \begin{minipage}[b]{\linewidth}\raggedright
PROPS
\end{minipage} & \begin{minipage}[b]{\linewidth}\raggedright
REGIONS
\end{minipage} & \begin{minipage}[b]{\linewidth}\raggedright
SOLUTION
\end{minipage} & \begin{minipage}[b]{\linewidth}\raggedright
SUMMARY
\end{minipage} & \begin{minipage}[b]{\linewidth}\raggedright
SCHEDULE
\end{minipage} \\
\midrule\noalign{}
\endhead
\bottomrule\noalign{}
\endlastfoot
\end{longtable}

\subsubsection{Description}\label{description-614}

The ROCKTSIG keyword defines the rock compaction attributes to be
applied for when the rock compaction option has been invoked by the
ROCKCOMP keyword in the RUNSPEC section, and the either the Dual
Permeability or Dual Porosity models are activated by the DUALPERM and
DUALPORO keywords in the RUNSPEC section. ROCKTSIG specifies sigma
multipliers versus pressure that are used in the dual porosity rock
compaction calculations. The keyword should only be used if the Rock
Compaction Hysteresis option has been activated by either setting the
ROCKOPT parameter on the ROCKCOMP keyword to one of the available
options.

This keyword is not supported by OPM Flow but would change the results
if supported so the simulation will be stopped.

\subsection{ROCKWNOD -- Water Saturation Values for Compaction
Pressure-Sw
Tables}\label{rockwnod-water-saturation-values-for-compaction-pressure-sw-tables}

\begin{longtable}[]{@{}
  >{\raggedright\arraybackslash}p{(\linewidth - 14\tabcolsep) * \real{0.1250}}
  >{\raggedright\arraybackslash}p{(\linewidth - 14\tabcolsep) * \real{0.1250}}
  >{\raggedright\arraybackslash}p{(\linewidth - 14\tabcolsep) * \real{0.1250}}
  >{\raggedright\arraybackslash}p{(\linewidth - 14\tabcolsep) * \real{0.1250}}
  >{\raggedright\arraybackslash}p{(\linewidth - 14\tabcolsep) * \real{0.1250}}
  >{\raggedright\arraybackslash}p{(\linewidth - 14\tabcolsep) * \real{0.1250}}
  >{\raggedright\arraybackslash}p{(\linewidth - 14\tabcolsep) * \real{0.1250}}
  >{\raggedright\arraybackslash}p{(\linewidth - 14\tabcolsep) * \real{0.1250}}@{}}
\toprule\noalign{}
\begin{minipage}[b]{\linewidth}\raggedright
RUNSPEC
\end{minipage} & \begin{minipage}[b]{\linewidth}\raggedright
GRID
\end{minipage} & \begin{minipage}[b]{\linewidth}\raggedright
EDIT
\end{minipage} & \begin{minipage}[b]{\linewidth}\raggedright
PROPS
\end{minipage} & \begin{minipage}[b]{\linewidth}\raggedright
REGIONS
\end{minipage} & \begin{minipage}[b]{\linewidth}\raggedright
SOLUTION
\end{minipage} & \begin{minipage}[b]{\linewidth}\raggedright
SUMMARY
\end{minipage} & \begin{minipage}[b]{\linewidth}\raggedright
SCHEDULE
\end{minipage} \\
\midrule\noalign{}
\endhead
\bottomrule\noalign{}
\endlastfoot
\end{longtable}

\subsubsection{Description}\label{description-615}

The ROCK2D and the ROCK2DTR keywords in the PROPS section define rock
compressibility pore volume and transmissibility multipliers as a
function of pressure and water saturation (``Sw''), for when the rock
compaction option has been invoked by the ROCKCOMP keyword in the
RUNSPEC section. The pressure values are defined on ROCK2D and the
ROCK2DTR keywords together with the multipliers. This keyword ROCKWNOD,
defines the water saturations that are used in conjunction with the
ROCK2D and the ROCK2DTR keywords.

This keyword should only be used if compaction option has been enabled.

\begin{longtable}[]{@{}
  >{\raggedright\arraybackslash}p{(\linewidth - 6\tabcolsep) * \real{0.2500}}
  >{\raggedright\arraybackslash}p{(\linewidth - 6\tabcolsep) * \real{0.2500}}
  >{\raggedright\arraybackslash}p{(\linewidth - 6\tabcolsep) * \real{0.2500}}
  >{\raggedright\arraybackslash}p{(\linewidth - 6\tabcolsep) * \real{0.2500}}@{}}
\toprule\noalign{}
\begin{minipage}[b]{\linewidth}\raggedright
No.
\end{minipage} & \begin{minipage}[b]{\linewidth}\raggedright
Name
\end{minipage} & \begin{minipage}[b]{\linewidth}\raggedright
Description
\end{minipage} & \begin{minipage}[b]{\linewidth}\raggedright
Default
\end{minipage} \\
\midrule\noalign{}
\endhead
\bottomrule\noalign{}
\endlastfoot
Field & Metric & Laboratory & \\
1 & SWAT & A columnar vector of real monotonically increasing down the
column values that defines the water saturations to be associated with
the data on the ROCK2D and the ROCKTR keywords. & None \\
dimensionless & dimensionless & dimensionless & \\
Notes: & & & \\
\end{longtable}

\emph{Table 8.131: ROCKWNOD Keyword Description}

See also the OVERBURD, ROCKTAB, ROCK2D and ROCK2DTR keywords in the
PROPS section.

\subsubsection{Example}\label{example-341}

The following example defines two ROCKWNOD tables for the pore volume
and transmissibility compaction tables, assuming NTROCC is equal to two
on the ROCKCOMP keyword and NSSFUN is greater than or equal to four on
the TABDIMS keyword.

\begin{verbatim}
--
--       WATER SATURATION VALUES FOR COMPACTION PRESSURE-SW TABLES
--
ROCKWNOD
--       COMPACT
--       SWAT
--       ------
          0.000
          0.200
          0.400
          1.000                                           / P-SW SET TABLE NO. 01
--       COMPACT
--       SWAT
--       ------
          0.000
          0.250
          0.750
          1.000                                           / P-SW SET TABLE NO. 02
\end{verbatim}

Note that there must be exactly NTROCC tables entered for this keyword,
otherwise an error will occur.

\subsection{RPTPROPS -- Define PROPS Section
Reporting}\label{rptprops-define-props-section-reporting}

\begin{longtable}[]{@{}
  >{\raggedright\arraybackslash}p{(\linewidth - 14\tabcolsep) * \real{0.1250}}
  >{\raggedright\arraybackslash}p{(\linewidth - 14\tabcolsep) * \real{0.1250}}
  >{\raggedright\arraybackslash}p{(\linewidth - 14\tabcolsep) * \real{0.1250}}
  >{\raggedright\arraybackslash}p{(\linewidth - 14\tabcolsep) * \real{0.1250}}
  >{\raggedright\arraybackslash}p{(\linewidth - 14\tabcolsep) * \real{0.1250}}
  >{\raggedright\arraybackslash}p{(\linewidth - 14\tabcolsep) * \real{0.1250}}
  >{\raggedright\arraybackslash}p{(\linewidth - 14\tabcolsep) * \real{0.1250}}
  >{\raggedright\arraybackslash}p{(\linewidth - 14\tabcolsep) * \real{0.1250}}@{}}
\toprule\noalign{}
\begin{minipage}[b]{\linewidth}\raggedright
RUNSPEC
\end{minipage} & \begin{minipage}[b]{\linewidth}\raggedright
GRID
\end{minipage} & \begin{minipage}[b]{\linewidth}\raggedright
EDIT
\end{minipage} & \begin{minipage}[b]{\linewidth}\raggedright
PROPS
\end{minipage} & \begin{minipage}[b]{\linewidth}\raggedright
REGIONS
\end{minipage} & \begin{minipage}[b]{\linewidth}\raggedright
SOLUTION
\end{minipage} & \begin{minipage}[b]{\linewidth}\raggedright
SUMMARY
\end{minipage} & \begin{minipage}[b]{\linewidth}\raggedright
SCHEDULE
\end{minipage} \\
\midrule\noalign{}
\endhead
\bottomrule\noalign{}
\endlastfoot
\end{longtable}

\subsubsection{Description}\label{description-616}

This keyword defines the data in the PROPS section that is to be printed
to the output print file in human readable format. The keyword has two
distinct forms, the first of which consists of the keyword followed by a
series of integers on the next line indicating the data to be printed
(see the first example). This is the original formal in the commercial
simulator and was subsequently superseded by the second format. The
second format consists of the keyword followed by a series of character
strings that indicate the data to be printed. In most cases the
character string is the keyword used to load the data in the OPM Flow
input deck, for example PVDG for the dry gas PVT tables. Its is
anticipated that OPM Flow will eventually support the functionality of
the second format only, the first format although recognized will be
completely ignored.

This keyword is not supported by OPM Flow but has no effect on the
results so it will be ignored.

\begin{longtable}[]{@{}
  >{\raggedright\arraybackslash}p{(\linewidth - 6\tabcolsep) * \real{0.2500}}
  >{\raggedright\arraybackslash}p{(\linewidth - 6\tabcolsep) * \real{0.2500}}
  >{\raggedright\arraybackslash}p{(\linewidth - 6\tabcolsep) * \real{0.2500}}
  >{\raggedright\arraybackslash}p{(\linewidth - 6\tabcolsep) * \real{0.2500}}@{}}
\toprule\noalign{}
\begin{minipage}[b]{\linewidth}\raggedright
No.
\end{minipage} & \begin{minipage}[b]{\linewidth}\raggedright
Name
\end{minipage} & \begin{minipage}[b]{\linewidth}\raggedright
Description
\end{minipage} & \begin{minipage}[b]{\linewidth}\raggedright
Default
\end{minipage} \\
\midrule\noalign{}
\endhead
\bottomrule\noalign{}
\endlastfoot
1 & PVDG & Print dry gas PVT tables & N/A \\
2 & PVTG & Print wet gas PVT tables & N/A \\
3 & SGFN & Print gas relative permeability saturation function tables. &
N/A \\
4 & SGL & Print connate gas saturation array. & N/A \\
\ldots. & & \ldots. & N/A \\
Notes: & & & \\
\end{longtable}

\emph{Table 8.132: RPTPROPS Keyword Description}

\begin{tcolorbox}[enhanced jigsaw, title=\textcolor{quarto-callout-note-color}{\faInfo}\hspace{0.5em}{Note}, coltitle=black, breakable, toprule=.15mm, colbacktitle=quarto-callout-note-color!10!white, leftrule=.75mm, opacitybacktitle=0.6, colback=white, bottomtitle=1mm, toptitle=1mm, titlerule=0mm, left=2mm, colframe=quarto-callout-note-color-frame, arc=.35mm, rightrule=.15mm, bottomrule=.15mm, opacityback=0]

Except for tabular like data, PVDG etc., this keyword has the potential
to produce very large print files that some text editors may have
difficulty loading. A more efficient solution for array type data is to
load the *.INIT file into OPM ResInsight to view the data graphically,
this also has the benefit of being able to filter the grid based on I,
J, K ranges and grid properties.

\end{tcolorbox}

\subsubsection{Examples}\label{examples-42}

The first example shows the original format of this keyword; although
the keyword and format are recognized by OPM Flow, the format is ignored
and is unlikely to be implemented in in the simulator.

\begin{verbatim}
--
--       DEFINE PROPS SECTION REPORT OPTION (ORIGINAL FORMAT)
–
RPTPROPS
         1        2*0      1        3*1                                        /
\end{verbatim}

The next example shows the second format of the keyword which may be
supported in a future release of OPM Flow.

\begin{verbatim}
--
--       DEFINE PROPS SECTION REPORT OPTIONS
--
RPTPROPS
         PVDO     SOF2     SGFN     SWFN                                       /
\end{verbatim}

\subsection{RSCONST -- Define Constant GOR (Rs) for All Dead Oil PVT
Fluids}\label{rsconst-define-constant-gor-rs-for-all-dead-oil-pvt-fluids}

\begin{longtable}[]{@{}
  >{\raggedright\arraybackslash}p{(\linewidth - 14\tabcolsep) * \real{0.1250}}
  >{\raggedright\arraybackslash}p{(\linewidth - 14\tabcolsep) * \real{0.1250}}
  >{\raggedright\arraybackslash}p{(\linewidth - 14\tabcolsep) * \real{0.1250}}
  >{\raggedright\arraybackslash}p{(\linewidth - 14\tabcolsep) * \real{0.1250}}
  >{\raggedright\arraybackslash}p{(\linewidth - 14\tabcolsep) * \real{0.1250}}
  >{\raggedright\arraybackslash}p{(\linewidth - 14\tabcolsep) * \real{0.1250}}
  >{\raggedright\arraybackslash}p{(\linewidth - 14\tabcolsep) * \real{0.1250}}
  >{\raggedright\arraybackslash}p{(\linewidth - 14\tabcolsep) * \real{0.1250}}@{}}
\toprule\noalign{}
\begin{minipage}[b]{\linewidth}\raggedright
RUNSPEC
\end{minipage} & \begin{minipage}[b]{\linewidth}\raggedright
GRID
\end{minipage} & \begin{minipage}[b]{\linewidth}\raggedright
EDIT
\end{minipage} & \begin{minipage}[b]{\linewidth}\raggedright
PROPS
\end{minipage} & \begin{minipage}[b]{\linewidth}\raggedright
REGIONS
\end{minipage} & \begin{minipage}[b]{\linewidth}\raggedright
SOLUTION
\end{minipage} & \begin{minipage}[b]{\linewidth}\raggedright
SUMMARY
\end{minipage} & \begin{minipage}[b]{\linewidth}\raggedright
SCHEDULE
\end{minipage} \\
\midrule\noalign{}
\endhead
\bottomrule\noalign{}
\endlastfoot
\end{longtable}

\subsubsection{Description}\label{description-617}

RSCONST defines a constant Gas-Oil Ratio (``GOR''), for all dead oil
{[}``Dead'' oil is oil that it contains no dissolved gas or a relatively
thick oil or residue that has lost its volatile components.{]} PVT
fluids. If the oil has a constant and uniform dissolved gas
concentration, GOR, and if the reservoir pressure never drops below the
saturation pressure (bubble point pressure), then the model can be run
more efficiently by omitting the GAS and DISGAS keywords from the
RUNSPEC section, treating the oil as a dead oil, and defining a constant
Rs (GOR) value with keywords RSCONST or RSCONSTT in the PROPS section.
This results in the model being run as a dead oil problem with no active
gas phase. However, OPM Flow takes into account the constant Rs in the
calculations and reporting.

This keyword is not supported by OPM Flow but would change the results
if supported so the simulation will be stopped.

\begin{longtable}[]{@{}
  >{\raggedright\arraybackslash}p{(\linewidth - 6\tabcolsep) * \real{0.2500}}
  >{\raggedright\arraybackslash}p{(\linewidth - 6\tabcolsep) * \real{0.2500}}
  >{\raggedright\arraybackslash}p{(\linewidth - 6\tabcolsep) * \real{0.2500}}
  >{\raggedright\arraybackslash}p{(\linewidth - 6\tabcolsep) * \real{0.2500}}@{}}
\toprule\noalign{}
\begin{minipage}[b]{\linewidth}\raggedright
No.
\end{minipage} & \begin{minipage}[b]{\linewidth}\raggedright
Name
\end{minipage} & \begin{minipage}[b]{\linewidth}\raggedright
Description
\end{minipage} & \begin{minipage}[b]{\linewidth}\raggedright
Default
\end{minipage} \\
\midrule\noalign{}
\endhead
\bottomrule\noalign{}
\endlastfoot
Field & Metric & Laboratory & \\
1 & RS & A real positive value that defines the dead oil GOR for all oil
PVT tables in the model & None \\
Mscf/stb & sm3/sm3 & scc/scc & \\
2 & PRESS & A real positive value that defines that saturation pressure
(bubble point pressure) for all the oil PVT tables in the model. &
None \\
psia & barsa & atma & \\
Notes: & & & \\
\end{longtable}

\emph{Table 8.133: RSCONST Keyword Description}

See also the RSCONSTT keyword to define a different constant Rs to the
various dead oil PVT tables and the PVDO and PVCDO keywords to enter the
dead oil properties. All of the aforementioned keywords are in the PROPS
section.

\subsubsection{Example}\label{example-342}

The example sets the dead oil GOR to 5 scf/stb and the bubble point
pressure to 14.7 psia.

\begin{verbatim}
--
--       DEAD OIL PVT CONSTANT GOR AND SATURATION PRESSURE
--
RSCONST
--       RS        PSAT
--       MSCF/STB  PSIA
--       --------  ------
          0.0050    14.7                                   /
\end{verbatim}

\subsection{RSCONSTT -- Define Constant GOR (Rs) for Each Dead Oil PVT
Fluid}\label{rsconstt-define-constant-gor-rs-for-each-dead-oil-pvt-fluid}

\begin{longtable}[]{@{}
  >{\raggedright\arraybackslash}p{(\linewidth - 14\tabcolsep) * \real{0.1250}}
  >{\raggedright\arraybackslash}p{(\linewidth - 14\tabcolsep) * \real{0.1250}}
  >{\raggedright\arraybackslash}p{(\linewidth - 14\tabcolsep) * \real{0.1250}}
  >{\raggedright\arraybackslash}p{(\linewidth - 14\tabcolsep) * \real{0.1250}}
  >{\raggedright\arraybackslash}p{(\linewidth - 14\tabcolsep) * \real{0.1250}}
  >{\raggedright\arraybackslash}p{(\linewidth - 14\tabcolsep) * \real{0.1250}}
  >{\raggedright\arraybackslash}p{(\linewidth - 14\tabcolsep) * \real{0.1250}}
  >{\raggedright\arraybackslash}p{(\linewidth - 14\tabcolsep) * \real{0.1250}}@{}}
\toprule\noalign{}
\begin{minipage}[b]{\linewidth}\raggedright
RUNSPEC
\end{minipage} & \begin{minipage}[b]{\linewidth}\raggedright
GRID
\end{minipage} & \begin{minipage}[b]{\linewidth}\raggedright
EDIT
\end{minipage} & \begin{minipage}[b]{\linewidth}\raggedright
PROPS
\end{minipage} & \begin{minipage}[b]{\linewidth}\raggedright
REGIONS
\end{minipage} & \begin{minipage}[b]{\linewidth}\raggedright
SOLUTION
\end{minipage} & \begin{minipage}[b]{\linewidth}\raggedright
SUMMARY
\end{minipage} & \begin{minipage}[b]{\linewidth}\raggedright
SCHEDULE
\end{minipage} \\
\midrule\noalign{}
\endhead
\bottomrule\noalign{}
\endlastfoot
\end{longtable}

\subsubsection{Description}\label{description-618}

RSCONSTT defines a constant Gas-Oil Ratio (``GOR''), for each dead oil
{[}``Dead'' oil is oil that it contains no dissolved gas or a relatively
thick oil or residue that has lost its volatile components.{]} PVT fluid
in the model. If the oil has a constant and uniform dissolved gas
concentration, GOR, and if the reservoir pressure never drops below the
saturation pressure (bubble point pressure), then the model can be run
more efficiently by omitting the GAS and DISGAS keywords from the
RUNSPEC section, treating the oil as a dead oil, and defining a constant
Rs (GOR) value with keywords RSCONST or RSCONSTT in the PROPS section.
This results in the model being run as a dead oil problem with no active
gas phase. However, OPM Flow takes into account the constant Rs in the
calculations and reporting.

This keyword is not supported by OPM Flow but would change the results
if supported so the simulation will be stopped.

\begin{longtable}[]{@{}
  >{\raggedright\arraybackslash}p{(\linewidth - 6\tabcolsep) * \real{0.2500}}
  >{\raggedright\arraybackslash}p{(\linewidth - 6\tabcolsep) * \real{0.2500}}
  >{\raggedright\arraybackslash}p{(\linewidth - 6\tabcolsep) * \real{0.2500}}
  >{\raggedright\arraybackslash}p{(\linewidth - 6\tabcolsep) * \real{0.2500}}@{}}
\toprule\noalign{}
\begin{minipage}[b]{\linewidth}\raggedright
No.
\end{minipage} & \begin{minipage}[b]{\linewidth}\raggedright
Name
\end{minipage} & \begin{minipage}[b]{\linewidth}\raggedright
Description
\end{minipage} & \begin{minipage}[b]{\linewidth}\raggedright
Default
\end{minipage} \\
\midrule\noalign{}
\endhead
\bottomrule\noalign{}
\endlastfoot
Field & Metric & Laboratory & \\
1 & RS & A real positive columnar vector that defines the dead oil GOR
for each oil PVT table in the model & None \\
Mscf/stb & sm3/sm3 & scc/scc & \\
2 & PRESS & A real positive columnar vector that defines the saturation
pressure (bubble point pressure) for each the oil PVT table in the
model. & None \\
psia & barsa & atma & \\
Notes: & & & \\
\end{longtable}

\emph{Table 8.134: RSCONSTT Keyword Description}

See also the RSCONST keyword to define a constant Rs to all the various
dead oil PVT tables and the PVDO and PVCDO keywords to enter the dead
oil properties. All of the aforementioned keywords are in the PROPS
section.

\subsubsection{Example}\label{example-343}

The example sets the dead oil GOR to 5, 6.5 and 8.0 scf/stb for PVT
tables one, two and three, respectively and the bubble point pressure to
14.7 psia for all three tables.

\begin{verbatim}
--
--       DEAD OIL PVT CONSTANT GOR AND SATURATION PRESSURE
--
RSCONSTT
--       RS        PSAT
--       MSCF/STB  PSIA
--       --------  ------
          0.0050     14.7                                  / TABLE NO. 01
          0.0065     14.7                                  / TABLE NO. 02
          0.0080     14.7                                  / TABLE NO. 03
\end{verbatim}

\subsection{RSGI -- Define Gas-Oil Ratio versus Pressure and Gi
Tables}\label{rsgi-define-gas-oil-ratio-versus-pressure-and-gi-tables}

\begin{longtable}[]{@{}
  >{\raggedright\arraybackslash}p{(\linewidth - 14\tabcolsep) * \real{0.1250}}
  >{\raggedright\arraybackslash}p{(\linewidth - 14\tabcolsep) * \real{0.1250}}
  >{\raggedright\arraybackslash}p{(\linewidth - 14\tabcolsep) * \real{0.1250}}
  >{\raggedright\arraybackslash}p{(\linewidth - 14\tabcolsep) * \real{0.1250}}
  >{\raggedright\arraybackslash}p{(\linewidth - 14\tabcolsep) * \real{0.1250}}
  >{\raggedright\arraybackslash}p{(\linewidth - 14\tabcolsep) * \real{0.1250}}
  >{\raggedright\arraybackslash}p{(\linewidth - 14\tabcolsep) * \real{0.1250}}
  >{\raggedright\arraybackslash}p{(\linewidth - 14\tabcolsep) * \real{0.1250}}@{}}
\toprule\noalign{}
\begin{minipage}[b]{\linewidth}\raggedright
RUNSPEC
\end{minipage} & \begin{minipage}[b]{\linewidth}\raggedright
GRID
\end{minipage} & \begin{minipage}[b]{\linewidth}\raggedright
EDIT
\end{minipage} & \begin{minipage}[b]{\linewidth}\raggedright
PROPS
\end{minipage} & \begin{minipage}[b]{\linewidth}\raggedright
REGIONS
\end{minipage} & \begin{minipage}[b]{\linewidth}\raggedright
SOLUTION
\end{minipage} & \begin{minipage}[b]{\linewidth}\raggedright
SUMMARY
\end{minipage} & \begin{minipage}[b]{\linewidth}\raggedright
SCHEDULE
\end{minipage} \\
\midrule\noalign{}
\endhead
\bottomrule\noalign{}
\endlastfoot
\end{longtable}

\subsubsection{Description}\label{description-619}

The RSGI keyword specifies the saturated oil Gas-Oil Ratio (``GOR'')
factors used to specify the variation of the maximum possible GOR of oil
with respect to pressure and Gi values, for when the GIMODEL keyword in
the RUNSPEC section has been used to activate the GI Pseudo
Compositional option for the run. See also the GINODE, RSGI, RVGI, BGGI
and BOGI keywords in the PROPS section to describe the fluid properties
for the GI Pseudo Compositional option.

This keyword is not supported by OPM Flow but would change the results
if supported so the simulation will be stopped.

\subsection{RTEMP -- Define the Initial Reservoir Temperature for the
Model}\label{rtemp-define-the-initial-reservoir-temperature-for-the-model}

\begin{longtable}[]{@{}
  >{\raggedright\arraybackslash}p{(\linewidth - 14\tabcolsep) * \real{0.1250}}
  >{\raggedright\arraybackslash}p{(\linewidth - 14\tabcolsep) * \real{0.1250}}
  >{\raggedright\arraybackslash}p{(\linewidth - 14\tabcolsep) * \real{0.1250}}
  >{\raggedright\arraybackslash}p{(\linewidth - 14\tabcolsep) * \real{0.1250}}
  >{\raggedright\arraybackslash}p{(\linewidth - 14\tabcolsep) * \real{0.1250}}
  >{\raggedright\arraybackslash}p{(\linewidth - 14\tabcolsep) * \real{0.1250}}
  >{\raggedright\arraybackslash}p{(\linewidth - 14\tabcolsep) * \real{0.1250}}
  >{\raggedright\arraybackslash}p{(\linewidth - 14\tabcolsep) * \real{0.1250}}@{}}
\toprule\noalign{}
\begin{minipage}[b]{\linewidth}\raggedright
RUNSPEC
\end{minipage} & \begin{minipage}[b]{\linewidth}\raggedright
GRID
\end{minipage} & \begin{minipage}[b]{\linewidth}\raggedright
EDIT
\end{minipage} & \begin{minipage}[b]{\linewidth}\raggedright
PROPS
\end{minipage} & \begin{minipage}[b]{\linewidth}\raggedright
REGIONS
\end{minipage} & \begin{minipage}[b]{\linewidth}\raggedright
SOLUTION
\end{minipage} & \begin{minipage}[b]{\linewidth}\raggedright
SUMMARY
\end{minipage} & \begin{minipage}[b]{\linewidth}\raggedright
SCHEDULE
\end{minipage} \\
\midrule\noalign{}
\endhead
\bottomrule\noalign{}
\endlastfoot
\end{longtable}

\subsubsection{Description}\label{description-620}

This keyword defines the initial reservoir temperature for the model.
Note that the RTEMP keyword is an alias for RTEMPA, and that both
keywords are supported by OPM Flow, in both the PROPS and SOLUTION
sections, but are treated as being mutually exclusive.

The initial reservoir temperature must be defined when OPM Flow's
thermal option has been activated by the THERMAL keyword in the RUNSPEC
section. Note this is different to the commercial simulator that uses
the TEMP keyword in the RUNSPEC section to activate the black-oil
temperature model, and the THERMAL keyword to activate the compositional
thermal model.

The initial reservoir temperature should be defined when OPM Flow's CO2
or H2 storage option has been activated by the CO2STORE or
\hyperref[REF_HEADING_KEYWORD_H2STORE]{H2STORE} keyword in the RUNSPEC
section.

\begin{longtable}[]{@{}
  >{\raggedright\arraybackslash}p{(\linewidth - 6\tabcolsep) * \real{0.2500}}
  >{\raggedright\arraybackslash}p{(\linewidth - 6\tabcolsep) * \real{0.2500}}
  >{\raggedright\arraybackslash}p{(\linewidth - 6\tabcolsep) * \real{0.2500}}
  >{\raggedright\arraybackslash}p{(\linewidth - 6\tabcolsep) * \real{0.2500}}@{}}
\toprule\noalign{}
\begin{minipage}[b]{\linewidth}\raggedright
No.
\end{minipage} & \begin{minipage}[b]{\linewidth}\raggedright
Name
\end{minipage} & \begin{minipage}[b]{\linewidth}\raggedright
Description
\end{minipage} & \begin{minipage}[b]{\linewidth}\raggedright
Default
\end{minipage} \\
\midrule\noalign{}
\endhead
\bottomrule\noalign{}
\endlastfoot
Field & Metric & Laboratory & \\
1 & RTEMP & Single real positive value that defines the reservoir
temperature for the model. & None \\
oF & oC & oC & \\
Notes: & & & \\
\end{longtable}

\emph{Table 8.3.259.1: RTEMP Keyword Description}

See also the RTEMPVD keyword in SOLUTION section to define the reservoir
temperature as a function of depth.

\subsubsection{Example}\label{example-344}

\begin{verbatim}
--
--       RESERVOIR
--       TEMPERATURE
--       -----------
RTEMP
         190.0                                            / RESERVOIR TEMPERATURE
\end{verbatim}

The above example defines the reservoir temperature to be 190 oF.

\subsection{RTEMPA -- Define the Initial Reservoir Temperature for the
Model}\label{rtempa-define-the-initial-reservoir-temperature-for-the-model}

\begin{longtable}[]{@{}
  >{\raggedright\arraybackslash}p{(\linewidth - 14\tabcolsep) * \real{0.1250}}
  >{\raggedright\arraybackslash}p{(\linewidth - 14\tabcolsep) * \real{0.1250}}
  >{\raggedright\arraybackslash}p{(\linewidth - 14\tabcolsep) * \real{0.1250}}
  >{\raggedright\arraybackslash}p{(\linewidth - 14\tabcolsep) * \real{0.1250}}
  >{\raggedright\arraybackslash}p{(\linewidth - 14\tabcolsep) * \real{0.1250}}
  >{\raggedright\arraybackslash}p{(\linewidth - 14\tabcolsep) * \real{0.1250}}
  >{\raggedright\arraybackslash}p{(\linewidth - 14\tabcolsep) * \real{0.1250}}
  >{\raggedright\arraybackslash}p{(\linewidth - 14\tabcolsep) * \real{0.1250}}@{}}
\toprule\noalign{}
\begin{minipage}[b]{\linewidth}\raggedright
RUNSPEC
\end{minipage} & \begin{minipage}[b]{\linewidth}\raggedright
GRID
\end{minipage} & \begin{minipage}[b]{\linewidth}\raggedright
EDIT
\end{minipage} & \begin{minipage}[b]{\linewidth}\raggedright
PROPS
\end{minipage} & \begin{minipage}[b]{\linewidth}\raggedright
REGIONS
\end{minipage} & \begin{minipage}[b]{\linewidth}\raggedright
SOLUTION
\end{minipage} & \begin{minipage}[b]{\linewidth}\raggedright
SUMMARY
\end{minipage} & \begin{minipage}[b]{\linewidth}\raggedright
SCHEDULE
\end{minipage} \\
\midrule\noalign{}
\endhead
\bottomrule\noalign{}
\endlastfoot
\end{longtable}

\subsubsection{Description}\label{description-621}

This keyword defines the initial reservoir temperature for the model.
Note that the RTEMPA keyword is an alias for RTEMP, and that both
keywords are supported by OPM Flow, in both the PROPS and SOLUTION
sections, but are treated as being mutually exclusive.

The initial reservoir temperature must be defined when OPM Flow's
thermal option has been activated by the THERMAL keyword in the RUNSPEC
section. Note this is different to the commercial simulator that uses
the TEMP keyword in the RUNSPEC section to activate the black-oil
temperature model, and the THERMAL keyword to activate the compositional
thermal model.

The initial reservoir temperature should be defined when OPM Flow's CO2
or H2 storage option has been activated by the CO2STORE or
\hyperref[REF_HEADING_KEYWORD_H2STORE]{H2STORE} keyword in the RUNSPEC
section.

\begin{longtable}[]{@{}
  >{\raggedright\arraybackslash}p{(\linewidth - 6\tabcolsep) * \real{0.2500}}
  >{\raggedright\arraybackslash}p{(\linewidth - 6\tabcolsep) * \real{0.2500}}
  >{\raggedright\arraybackslash}p{(\linewidth - 6\tabcolsep) * \real{0.2500}}
  >{\raggedright\arraybackslash}p{(\linewidth - 6\tabcolsep) * \real{0.2500}}@{}}
\toprule\noalign{}
\begin{minipage}[b]{\linewidth}\raggedright
No.
\end{minipage} & \begin{minipage}[b]{\linewidth}\raggedright
Name
\end{minipage} & \begin{minipage}[b]{\linewidth}\raggedright
Description
\end{minipage} & \begin{minipage}[b]{\linewidth}\raggedright
Default
\end{minipage} \\
\midrule\noalign{}
\endhead
\bottomrule\noalign{}
\endlastfoot
Field & Metric & Laboratory & \\
1 & RTEMPA & Single real positive value that define the reservoir
temperature for the model. & None \\
oF & oC & oC & \\
Notes: & & & \\
\end{longtable}

\emph{Table 8.136: RTEMPA Keyword Description}

See also the RTEMPVD keyword in SOLUTION section to define the reservoir
temperature as a function of depth.

\subsubsection{Example}\label{example-345}

\begin{verbatim}
--
--       RESERVOIR
--       TEMPERATURE
--       -----------
RTEMPA
         190.0                                            / RESERVOIR TEMPERATURE
\end{verbatim}

The above example defines the reservoir temperature to be 190 oF.

\subsection{RTEMPVD -- Define the Initial Reservoir Temperature versus
Depth
Tables}\label{rtempvd-define-the-initial-reservoir-temperature-versus-depth-tables}

\begin{longtable}[]{@{}
  >{\raggedright\arraybackslash}p{(\linewidth - 14\tabcolsep) * \real{0.1250}}
  >{\raggedright\arraybackslash}p{(\linewidth - 14\tabcolsep) * \real{0.1250}}
  >{\raggedright\arraybackslash}p{(\linewidth - 14\tabcolsep) * \real{0.1250}}
  >{\raggedright\arraybackslash}p{(\linewidth - 14\tabcolsep) * \real{0.1250}}
  >{\raggedright\arraybackslash}p{(\linewidth - 14\tabcolsep) * \real{0.1250}}
  >{\raggedright\arraybackslash}p{(\linewidth - 14\tabcolsep) * \real{0.1250}}
  >{\raggedright\arraybackslash}p{(\linewidth - 14\tabcolsep) * \real{0.1250}}
  >{\raggedright\arraybackslash}p{(\linewidth - 14\tabcolsep) * \real{0.1250}}@{}}
\toprule\noalign{}
\begin{minipage}[b]{\linewidth}\raggedright
RUNSPEC
\end{minipage} & \begin{minipage}[b]{\linewidth}\raggedright
GRID
\end{minipage} & \begin{minipage}[b]{\linewidth}\raggedright
EDIT
\end{minipage} & \begin{minipage}[b]{\linewidth}\raggedright
PROPS
\end{minipage} & \begin{minipage}[b]{\linewidth}\raggedright
REGIONS
\end{minipage} & \begin{minipage}[b]{\linewidth}\raggedright
SOLUTION
\end{minipage} & \begin{minipage}[b]{\linewidth}\raggedright
SUMMARY
\end{minipage} & \begin{minipage}[b]{\linewidth}\raggedright
SCHEDULE
\end{minipage} \\
\midrule\noalign{}
\endhead
\bottomrule\noalign{}
\endlastfoot
\end{longtable}

\subsubsection{Description}\label{description-622}

This keyword defines the initial reservoir temperature versus depth
tables for each equilibration region. Note that the RTEMPVD keyword is
an alias for TEMPVD, and that both keywords are supported by OPM Flow,
in both the PROPS and SOLUTION sections, but are treated as being
mutually exclusive.

See Error: Reference source not found in the SOLUTION section, as the
keyword is documented in the SOLUTION section, the same as the
commercial simulator, but it can also be used in the PROPS section by
OPM Flow.

\subsection{RVCONST -- Define Constant CGR (Rv) for All Dry Gas PVT
Fluids}\label{rvconst-define-constant-cgr-rv-for-all-dry-gas-pvt-fluids}

\begin{longtable}[]{@{}
  >{\raggedright\arraybackslash}p{(\linewidth - 14\tabcolsep) * \real{0.1250}}
  >{\raggedright\arraybackslash}p{(\linewidth - 14\tabcolsep) * \real{0.1250}}
  >{\raggedright\arraybackslash}p{(\linewidth - 14\tabcolsep) * \real{0.1250}}
  >{\raggedright\arraybackslash}p{(\linewidth - 14\tabcolsep) * \real{0.1250}}
  >{\raggedright\arraybackslash}p{(\linewidth - 14\tabcolsep) * \real{0.1250}}
  >{\raggedright\arraybackslash}p{(\linewidth - 14\tabcolsep) * \real{0.1250}}
  >{\raggedright\arraybackslash}p{(\linewidth - 14\tabcolsep) * \real{0.1250}}
  >{\raggedright\arraybackslash}p{(\linewidth - 14\tabcolsep) * \real{0.1250}}@{}}
\toprule\noalign{}
\begin{minipage}[b]{\linewidth}\raggedright
RUNSPEC
\end{minipage} & \begin{minipage}[b]{\linewidth}\raggedright
GRID
\end{minipage} & \begin{minipage}[b]{\linewidth}\raggedright
EDIT
\end{minipage} & \begin{minipage}[b]{\linewidth}\raggedright
PROPS
\end{minipage} & \begin{minipage}[b]{\linewidth}\raggedright
REGIONS
\end{minipage} & \begin{minipage}[b]{\linewidth}\raggedright
SOLUTION
\end{minipage} & \begin{minipage}[b]{\linewidth}\raggedright
SUMMARY
\end{minipage} & \begin{minipage}[b]{\linewidth}\raggedright
SCHEDULE
\end{minipage} \\
\midrule\noalign{}
\endhead
\bottomrule\noalign{}
\endlastfoot
\end{longtable}

\subsubsection{Description}\label{description-623}

RVCONST defines a constant Condensate-Gas Ratio (``CGR'' or Rv), for all
dry gas {[}Natural gas that occurs in the absence of condensate or
liquid hydrocarbons, or gas that had condensable hydrocarbons removed,
is called dry gas. It is primarily methane with some intermediates. The
hydrocarbon mixture is solely gas in the reservoir and there is no
liquid (condensate surface liquid) formed either in the reservoir or at
surface. The term dry indicates that the gas does not contain heavier
hydrocarbons to form liquids at the surface conditions. Dry gas
typically has GOR's greater than 100,000 scf/stb or 18,000 Sm3/m3.{]}
PVT fluids. If the gas has a constant and uniform dissolved condensate
concentration, and if the reservoir pressure never drops below the
saturation pressure (dew point pressure), then the model can be run more
efficiently by omitting the OIL and VAPOIL keywords from the RUNSPEC
section, treating the gas as a dry gas, and defining a constant Rv (CGR)
value with keywords RVCONST or RVCONSTT in the PROPS section. This
results in the model being run as a dry gas problem with no active oil
phase. However, OPM Flow takes into account the constant Rv in the
calculations and reporting.

This keyword is not supported by OPM Flow but would change the results
if supported so the simulation will be stopped.

\begin{longtable}[]{@{}
  >{\raggedright\arraybackslash}p{(\linewidth - 6\tabcolsep) * \real{0.2500}}
  >{\raggedright\arraybackslash}p{(\linewidth - 6\tabcolsep) * \real{0.2500}}
  >{\raggedright\arraybackslash}p{(\linewidth - 6\tabcolsep) * \real{0.2500}}
  >{\raggedright\arraybackslash}p{(\linewidth - 6\tabcolsep) * \real{0.2500}}@{}}
\toprule\noalign{}
\begin{minipage}[b]{\linewidth}\raggedright
No.
\end{minipage} & \begin{minipage}[b]{\linewidth}\raggedright
Name
\end{minipage} & \begin{minipage}[b]{\linewidth}\raggedright
Description
\end{minipage} & \begin{minipage}[b]{\linewidth}\raggedright
Default
\end{minipage} \\
\midrule\noalign{}
\endhead
\bottomrule\noalign{}
\endlastfoot
Field & Metric & Laboratory & \\
1 & RV & A real positive value that defines the dry gas CGR for all dry
gas PVT tables in the model & None \\
stb/Mscf & sm3/sm3 & scc/scc & \\
2 & PRESS & A real positive value that defines that saturation pressure
(dew point pressure) for all the dry gas PVT tables in the model. &
0.0 \\
psia & barsa & atma & \\
Notes: & & & \\
\end{longtable}

\emph{Table 8.137: RVCONST Keyword Description}

See also the RVCONSTT keyword to define a different constant Rv to the
various dry gas PVT tables and the PVDG keyword to enter the dry gas
properties. All of the aforementioned keywords are in the PROPS section.

\subsubsection{Example}\label{example-346}

The example sets the dry gas CGR to 5 stb/MMscf and the bubble point
pressure to 14.7 psia.

\begin{verbatim}
--
--       DRY GAS PVT CONSTANT GCR AND SATURATION PRESSURE
--
RVCONST
--       RV        PSAT
--       STB/MSCF  PSIA
--       --------  ------
          0.0050    14.7                                   /
\end{verbatim}

\subsection{RVCONSTT -- Define Constant CGR (Rv) for Each Dry Gas PVT
Fluid}\label{rvconstt-define-constant-cgr-rv-for-each-dry-gas-pvt-fluid}

\begin{longtable}[]{@{}
  >{\raggedright\arraybackslash}p{(\linewidth - 14\tabcolsep) * \real{0.1250}}
  >{\raggedright\arraybackslash}p{(\linewidth - 14\tabcolsep) * \real{0.1250}}
  >{\raggedright\arraybackslash}p{(\linewidth - 14\tabcolsep) * \real{0.1250}}
  >{\raggedright\arraybackslash}p{(\linewidth - 14\tabcolsep) * \real{0.1250}}
  >{\raggedright\arraybackslash}p{(\linewidth - 14\tabcolsep) * \real{0.1250}}
  >{\raggedright\arraybackslash}p{(\linewidth - 14\tabcolsep) * \real{0.1250}}
  >{\raggedright\arraybackslash}p{(\linewidth - 14\tabcolsep) * \real{0.1250}}
  >{\raggedright\arraybackslash}p{(\linewidth - 14\tabcolsep) * \real{0.1250}}@{}}
\toprule\noalign{}
\begin{minipage}[b]{\linewidth}\raggedright
RUNSPEC
\end{minipage} & \begin{minipage}[b]{\linewidth}\raggedright
GRID
\end{minipage} & \begin{minipage}[b]{\linewidth}\raggedright
EDIT
\end{minipage} & \begin{minipage}[b]{\linewidth}\raggedright
PROPS
\end{minipage} & \begin{minipage}[b]{\linewidth}\raggedright
REGIONS
\end{minipage} & \begin{minipage}[b]{\linewidth}\raggedright
SOLUTION
\end{minipage} & \begin{minipage}[b]{\linewidth}\raggedright
SUMMARY
\end{minipage} & \begin{minipage}[b]{\linewidth}\raggedright
SCHEDULE
\end{minipage} \\
\midrule\noalign{}
\endhead
\bottomrule\noalign{}
\endlastfoot
\end{longtable}

\subsubsection{Description}\label{description-624}

RVCONSTT defines a constant Condensate-Gas Ratio (``CGR'' or Rv), for
each dry gas {[}Natural gas that occurs in the absence of condensate or
liquid hydrocarbons, or gas that had condensable hydrocarbons removed,
is called dry gas. It is primarily methane with some intermediates. The
hydrocarbon mixture is solely gas in the reservoir and there is no
liquid (condensate surface liquid) formed either in the reservoir or at
surface. The term dry indicates that the gas does not contain heavier
hydrocarbons to form liquids at the surface conditions. Dry gas
typically has GOR's greater than 100,000 scf/stb or 18,000 Sm3/m3.{]}
PVT fluid. If the gas has a constant and uniform dissolved condensate
concentration, and if the reservoir pressure never drops below the
saturation pressure (dew point pressure), then the model can be run more
efficiently by omitting the OIL and VAPGAS keywords from the RUNSPEC
section, treating the gas as a dry gas, and defining a constant Rv (CGR)
value with keywords RVCONST or RVCONSTT in the PROPS section. This
results in the model being run as a dry gas problem with no active oil
phase. However, OPM Flow takes into account the constant Rv in the
calculations and reporting.

This keyword is ignored by OPM Flow but is documented here for
completeness.

\begin{longtable}[]{@{}
  >{\raggedright\arraybackslash}p{(\linewidth - 6\tabcolsep) * \real{0.2500}}
  >{\raggedright\arraybackslash}p{(\linewidth - 6\tabcolsep) * \real{0.2500}}
  >{\raggedright\arraybackslash}p{(\linewidth - 6\tabcolsep) * \real{0.2500}}
  >{\raggedright\arraybackslash}p{(\linewidth - 6\tabcolsep) * \real{0.2500}}@{}}
\toprule\noalign{}
\begin{minipage}[b]{\linewidth}\raggedright
No.
\end{minipage} & \begin{minipage}[b]{\linewidth}\raggedright
Name
\end{minipage} & \begin{minipage}[b]{\linewidth}\raggedright
Description
\end{minipage} & \begin{minipage}[b]{\linewidth}\raggedright
Default
\end{minipage} \\
\midrule\noalign{}
\endhead
\bottomrule\noalign{}
\endlastfoot
Field & Metric & Laboratory & \\
1 & RS & A real positive value that defines the dry gas CGR for each dry
gas PVT table in the model & None \\
stb/Mscf & sm3/sm3 & scc/scc & \\
2 & PRESS & A real positive value that defines that saturation pressure
(dew point pressure) for each dry gas PVT table in the model. & 0.0 \\
psia & barsa & atma & \\
Notes: & & & \\
\end{longtable}

\emph{Table 8.138: RVCONSTT Keyword Description}

See also the RVCONST keyword to define a constant Rv to all the various
dry gas PVT tables and the PVDG keyword to enter the dry gas properties.
All of the aforementioned keywords are in the PROPS section.

\subsubsection{Example}\label{example-347}

The example sets the dry gas CGR to 5, 6.5 and 8.0 stb/MMscf for PVT
tables one, two and three, respectively and the bubble point pressure to
14.7 psia for all three tables.

\begin{verbatim}
--
--       DRY GAS PVT CONSTANT GCR AND SATURATION PRESSURE
--
RVCONSTT
--       RV        PSAT
--       STB/MSCF  PSIA
--       --------  ------
          0.0050     14.7                                  / TABLE NO. 01
          0.0065     14.7                                  / TABLE NO. 02
          0.0080     14.7                                  / TABLE NO. 03
\end{verbatim}

\subsection{RVGI -- Define Condensate-Gas Ratio versus Pressure and Gi
Tables}\label{rvgi-define-condensate-gas-ratio-versus-pressure-and-gi-tables}

\begin{longtable}[]{@{}
  >{\raggedright\arraybackslash}p{(\linewidth - 14\tabcolsep) * \real{0.1250}}
  >{\raggedright\arraybackslash}p{(\linewidth - 14\tabcolsep) * \real{0.1250}}
  >{\raggedright\arraybackslash}p{(\linewidth - 14\tabcolsep) * \real{0.1250}}
  >{\raggedright\arraybackslash}p{(\linewidth - 14\tabcolsep) * \real{0.1250}}
  >{\raggedright\arraybackslash}p{(\linewidth - 14\tabcolsep) * \real{0.1250}}
  >{\raggedright\arraybackslash}p{(\linewidth - 14\tabcolsep) * \real{0.1250}}
  >{\raggedright\arraybackslash}p{(\linewidth - 14\tabcolsep) * \real{0.1250}}
  >{\raggedright\arraybackslash}p{(\linewidth - 14\tabcolsep) * \real{0.1250}}@{}}
\toprule\noalign{}
\begin{minipage}[b]{\linewidth}\raggedright
RUNSPEC
\end{minipage} & \begin{minipage}[b]{\linewidth}\raggedright
GRID
\end{minipage} & \begin{minipage}[b]{\linewidth}\raggedright
EDIT
\end{minipage} & \begin{minipage}[b]{\linewidth}\raggedright
PROPS
\end{minipage} & \begin{minipage}[b]{\linewidth}\raggedright
REGIONS
\end{minipage} & \begin{minipage}[b]{\linewidth}\raggedright
SOLUTION
\end{minipage} & \begin{minipage}[b]{\linewidth}\raggedright
SUMMARY
\end{minipage} & \begin{minipage}[b]{\linewidth}\raggedright
SCHEDULE
\end{minipage} \\
\midrule\noalign{}
\endhead
\bottomrule\noalign{}
\endlastfoot
\end{longtable}

\subsubsection{Description}\label{description-625}

The RVGI keyword specifies the saturated gas Condensate-Gas Ratio
(``CGR'') factors used to specify the variation of the maximum possible
CGR of gas with respect to pressure and Gi values, for when the GIMODEL
keyword in the RUNSPEC section has been used to activate the GI Pseudo
Compositional option for the run. See also the GINODE, RSGI, RVGI, BGGI
and BOGI keywords in the PROPS section to describe the fluid properties
for the GI Pseudo Compositional option.

This keyword is not supported by OPM Flow but would change the results
if supported so the simulation will be stopped.

\subsection{RWGSALT -- Water Vaporization versus Pressure and Salt
Concentration}\label{rwgsalt-water-vaporization-versus-pressure-and-salt-concentration}

\begin{longtable}[]{@{}
  >{\raggedright\arraybackslash}p{(\linewidth - 14\tabcolsep) * \real{0.1250}}
  >{\raggedright\arraybackslash}p{(\linewidth - 14\tabcolsep) * \real{0.1250}}
  >{\raggedright\arraybackslash}p{(\linewidth - 14\tabcolsep) * \real{0.1250}}
  >{\raggedright\arraybackslash}p{(\linewidth - 14\tabcolsep) * \real{0.1250}}
  >{\raggedright\arraybackslash}p{(\linewidth - 14\tabcolsep) * \real{0.1250}}
  >{\raggedright\arraybackslash}p{(\linewidth - 14\tabcolsep) * \real{0.1250}}
  >{\raggedright\arraybackslash}p{(\linewidth - 14\tabcolsep) * \real{0.1250}}
  >{\raggedright\arraybackslash}p{(\linewidth - 14\tabcolsep) * \real{0.1250}}@{}}
\toprule\noalign{}
\begin{minipage}[b]{\linewidth}\raggedright
RUNSPEC
\end{minipage} & \begin{minipage}[b]{\linewidth}\raggedright
GRID
\end{minipage} & \begin{minipage}[b]{\linewidth}\raggedright
EDIT
\end{minipage} & \begin{minipage}[b]{\linewidth}\raggedright
PROPS
\end{minipage} & \begin{minipage}[b]{\linewidth}\raggedright
REGIONS
\end{minipage} & \begin{minipage}[b]{\linewidth}\raggedright
SOLUTION
\end{minipage} & \begin{minipage}[b]{\linewidth}\raggedright
SUMMARY
\end{minipage} & \begin{minipage}[b]{\linewidth}\raggedright
SCHEDULE
\end{minipage} \\
\midrule\noalign{}
\endhead
\bottomrule\noalign{}
\endlastfoot
\end{longtable}

\subsubsection{Description}\label{description-626}

RWGSALT defines the relationship of water vaporization versus pressure
and salt concentration. This keyword should be used when the VAPWAT
keyword has be declared in the RUNSPEC section indicating that vaporized
water is present in the gas phase. In addition, if the Salt
Precipitation model has been activated via the BRINE and PRECSALT
keywords, also in the RUNSPEC section, then this keyword must be
present. The keyword may be used for gas-water and oil-water-gas input
decks that contain the either dry or wet gas and vaporized water phases.

\begin{tcolorbox}[enhanced jigsaw, title=\textcolor{quarto-callout-note-color}{\faInfo}\hspace{0.5em}{Note}, coltitle=black, breakable, toprule=.15mm, colbacktitle=quarto-callout-note-color!10!white, leftrule=.75mm, opacitybacktitle=0.6, colback=white, bottomtitle=1mm, toptitle=1mm, titlerule=0mm, left=2mm, colframe=quarto-callout-note-color-frame, arc=.35mm, rightrule=.15mm, bottomrule=.15mm, opacityback=0]

This is an OPM Flow specific keyword for the simulator's Water
Vaporization and Salt Precipitation Models, note that these are
extensions to the simulator's standard Brine model.

\end{tcolorbox}

\begin{longtable}[]{@{}
  >{\raggedright\arraybackslash}p{(\linewidth - 6\tabcolsep) * \real{0.2500}}
  >{\raggedright\arraybackslash}p{(\linewidth - 6\tabcolsep) * \real{0.2500}}
  >{\raggedright\arraybackslash}p{(\linewidth - 6\tabcolsep) * \real{0.2500}}
  >{\raggedright\arraybackslash}p{(\linewidth - 6\tabcolsep) * \real{0.2500}}@{}}
\toprule\noalign{}
\begin{minipage}[b]{\linewidth}\raggedright
No.
\end{minipage} & \begin{minipage}[b]{\linewidth}\raggedright
Name
\end{minipage} & \begin{minipage}[b]{\linewidth}\raggedright
Description
\end{minipage} & \begin{minipage}[b]{\linewidth}\raggedright
Default
\end{minipage} \\
\midrule\noalign{}
\endhead
\bottomrule\noalign{}
\endlastfoot
Field & Metric & Laboratory & \\
1 & PRESS & A real monotonically increasing down the column values that
define the gas phase pressure, that together with salt concentration,
defines the vaporized water in gas ratio (``VWGR'') or Rw & None \\
psia & barsa & atma & \\
2 & SALTCON & A real monotonically increasing positive columnar vector
defining the salt concentration in water. & None \\
lb/stb & kg/sm3 & gm/scc & \\
3 & RW & A columnar vector of real positive number values defining the
vaporized water in gas ratio (Rw) that for a given PRESS and SALTCON. &
None \\
stb/Mscf & sm3/sm3 & scc/scc & \\
Notes: & & & \\
\end{longtable}

\emph{Table 8.139: RWGSALT Keyword Description}

Since the water component is in both the water and the gas phases,
RWGSALT controls the amount of water component evaporated from the water
phase into the gas phase, which is a function of both the water phase
salinity and a grid cells pressure. The keyword should be used in
conjunction with the dry gas PVTGW keyword or the wet gas PVTGWO
keyword, both of which are in the PROPS section.

\subsubsection{Example}\label{example-348}

The example defines two RWGSALT tables assuming NTPVT equals two and
NPPVT is greater than or equal to eight on the TABDIMS keyword in the
RUNSPEC section.

\begin{verbatim}
--
--       WATER VAPORIZATION TABLE FOR BRINE  (OPM FLOW KEYWORD)
--
RWGSALT
--       PRES     SALTCONC    RVW
--       PSIA     LB/STB      STB/MSCF
--       ------  ----------   ---------
          300     0           0.000132
                  0.5         0.000132
                  1           0.000132                     /
          600     0           0.000132
                  0.5         0.000132
                  1           0.000132                     /
          900     0           0.000132
                  0.5         0.000132
                  1           0.000132                     /
         1200     0           0.000132
                  0.5         0.000132
                  1           0.000132                     /
         1500     0           0.000132
                  0.5         0.000132
                  1           0.000132                     /
         1800     0           0.000132
                  0.5         0.000132
                  1           0.000132                     /
         2100     0           0.000132
                  0.5         0.000132
                  1           0.000132                     /
         2400     0           0.000132
                  0.5         0.000132
                  1           0.000132                     /
                                                           / Table NO. 1
--       PRES     SALTCONC    RW
--       PSIA     LB/STB      STB/MSCF
--       ------  ----------   ---------
         300      0           0.000132
                  0.5         0.000132
                  1           0.000132                     /
         600      0           0.000132
                  0.5         0.000132
                  1           0.000132                     /
         900      0           0.000132
                  0.5         0.000132
                  1           0.000132                     /
         1200     0           0.000132
                  0.5         0.000132
                  1           0.000132                     /
         1500     0           0.000132
                  0.5         0.000132
                  1           0.000132                     /
         1800     0           0.000132
                  0.5         0.000132
                  1           0.000132                     /
         2100     0           0.000132
                  0.5         0.000132
                  1           0.000132                     /
         2400     0           0.000132
                  0.5         0.000132
                  1           0.000132                     /
                                                           / Table NO. 2
\end{verbatim}

\subsection{SALINITY -- Define the Reservoir Salinity for All
Cells}\label{salinity-define-the-reservoir-salinity-for-all-cells}

\begin{longtable}[]{@{}
  >{\raggedright\arraybackslash}p{(\linewidth - 14\tabcolsep) * \real{0.1250}}
  >{\raggedright\arraybackslash}p{(\linewidth - 14\tabcolsep) * \real{0.1250}}
  >{\raggedright\arraybackslash}p{(\linewidth - 14\tabcolsep) * \real{0.1250}}
  >{\raggedright\arraybackslash}p{(\linewidth - 14\tabcolsep) * \real{0.1250}}
  >{\raggedright\arraybackslash}p{(\linewidth - 14\tabcolsep) * \real{0.1250}}
  >{\raggedright\arraybackslash}p{(\linewidth - 14\tabcolsep) * \real{0.1250}}
  >{\raggedright\arraybackslash}p{(\linewidth - 14\tabcolsep) * \real{0.1250}}
  >{\raggedright\arraybackslash}p{(\linewidth - 14\tabcolsep) * \real{0.1250}}@{}}
\toprule\noalign{}
\begin{minipage}[b]{\linewidth}\raggedright
RUNSPEC
\end{minipage} & \begin{minipage}[b]{\linewidth}\raggedright
GRID
\end{minipage} & \begin{minipage}[b]{\linewidth}\raggedright
EDIT
\end{minipage} & \begin{minipage}[b]{\linewidth}\raggedright
PROPS
\end{minipage} & \begin{minipage}[b]{\linewidth}\raggedright
REGIONS
\end{minipage} & \begin{minipage}[b]{\linewidth}\raggedright
SOLUTION
\end{minipage} & \begin{minipage}[b]{\linewidth}\raggedright
SUMMARY
\end{minipage} & \begin{minipage}[b]{\linewidth}\raggedright
SCHEDULE
\end{minipage} \\
\midrule\noalign{}
\endhead
\bottomrule\noalign{}
\endlastfoot
\end{longtable}

\subsubsection{Description}\label{description-627}

The SALINITY keyword defines a uniform salinity for all cells in the
model. The keyword should only be used with OPM Flow's CO2-Brine model
which is activated via the CO2STORE keyword in the RUNSPEC section. This
keyword is a compositional keyword in the commercial simulator but has
been implemented in OPM Flow's black-oil CO2-Brine model.

\begin{longtable}[]{@{}
  >{\raggedright\arraybackslash}p{(\linewidth - 10\tabcolsep) * \real{0.1667}}
  >{\raggedright\arraybackslash}p{(\linewidth - 10\tabcolsep) * \real{0.1667}}
  >{\raggedright\arraybackslash}p{(\linewidth - 10\tabcolsep) * \real{0.1667}}
  >{\raggedright\arraybackslash}p{(\linewidth - 10\tabcolsep) * \real{0.1667}}
  >{\raggedright\arraybackslash}p{(\linewidth - 10\tabcolsep) * \real{0.1667}}
  >{\raggedright\arraybackslash}p{(\linewidth - 10\tabcolsep) * \real{0.1667}}@{}}
\toprule\noalign{}
\begin{minipage}[b]{\linewidth}\raggedright
No.
\end{minipage} & \begin{minipage}[b]{\linewidth}\raggedright
Name
\end{minipage} & \begin{minipage}[b]{\linewidth}\raggedright
Description
\end{minipage} & \begin{minipage}[b]{\linewidth}\raggedright
Default
\end{minipage} & \begin{minipage}[b]{\linewidth}\raggedright
\end{minipage} & \begin{minipage}[b]{\linewidth}\raggedright
\end{minipage} \\
\midrule\noalign{}
\endhead
\bottomrule\noalign{}
\endlastfoot
& & Field & Metric & Laboratory & 0 \\
1 & SALINITY & A real positive value that defines the salinity for all
grid blocks in the model for when the CO2-Brine model has been
activated. Note that the units for salinity are molality, that is
gm-M/Kg, and therefore the units are defined as given below with the
10-3 prefix. & & & \\
10-3 x lb-M/lb & 10-3 x kg-M/kg & 10-3 x gm-M/gm & & & \\
Notes: & & & & & \\
\end{longtable}

\emph{Table 8.140: SALINITY Keyword Description}

See also the CO2STORE keyword in the RUNSPEC section.

\subsubsection{Example}\label{example-349}

The example sets the salt salinity for all cells in the model to 0.001
lb-M/Ib.

\begin{verbatim}
--
--       SET SALINITY FOR ALL CELLS (OPM FLOW KEYWORD)
--
SALINITY
         1.0                                               /
\end{verbatim}

Note that units for salinity are to the 10-3, that is a value of 0.001
lb-M/lb should be entered as 1.0 (10-3 x lb-M/lb), as per the example.

\subsection{SALTMF -- Define the Salt Liquid-Phase Mole Fraction for All
Cells}\label{saltmf-define-the-salt-liquid-phase-mole-fraction-for-all-cells}

\begin{longtable}[]{@{}
  >{\raggedright\arraybackslash}p{(\linewidth - 14\tabcolsep) * \real{0.1250}}
  >{\raggedright\arraybackslash}p{(\linewidth - 14\tabcolsep) * \real{0.1250}}
  >{\raggedright\arraybackslash}p{(\linewidth - 14\tabcolsep) * \real{0.1250}}
  >{\raggedright\arraybackslash}p{(\linewidth - 14\tabcolsep) * \real{0.1250}}
  >{\raggedright\arraybackslash}p{(\linewidth - 14\tabcolsep) * \real{0.1250}}
  >{\raggedright\arraybackslash}p{(\linewidth - 14\tabcolsep) * \real{0.1250}}
  >{\raggedright\arraybackslash}p{(\linewidth - 14\tabcolsep) * \real{0.1250}}
  >{\raggedright\arraybackslash}p{(\linewidth - 14\tabcolsep) * \real{0.1250}}@{}}
\toprule\noalign{}
\begin{minipage}[b]{\linewidth}\raggedright
RUNSPEC
\end{minipage} & \begin{minipage}[b]{\linewidth}\raggedright
GRID
\end{minipage} & \begin{minipage}[b]{\linewidth}\raggedright
EDIT
\end{minipage} & \begin{minipage}[b]{\linewidth}\raggedright
PROPS
\end{minipage} & \begin{minipage}[b]{\linewidth}\raggedright
REGIONS
\end{minipage} & \begin{minipage}[b]{\linewidth}\raggedright
SOLUTION
\end{minipage} & \begin{minipage}[b]{\linewidth}\raggedright
SUMMARY
\end{minipage} & \begin{minipage}[b]{\linewidth}\raggedright
SCHEDULE
\end{minipage} \\
\midrule\noalign{}
\endhead
\bottomrule\noalign{}
\endlastfoot
\end{longtable}

\subsubsection{Description}\label{description-628}

The SALTMF keyword defines a uniform salt liquid-phase mole fraction for
all cells in the model. The keyword should only be used with OPM Flow's
CO2-Brine model which is activated via the CO2STORE keyword in the
RUNSPEC section.

This is an OPM Flow specific keyword.

\begin{longtable}[]{@{}
  >{\raggedright\arraybackslash}p{(\linewidth - 10\tabcolsep) * \real{0.1667}}
  >{\raggedright\arraybackslash}p{(\linewidth - 10\tabcolsep) * \real{0.1667}}
  >{\raggedright\arraybackslash}p{(\linewidth - 10\tabcolsep) * \real{0.1667}}
  >{\raggedright\arraybackslash}p{(\linewidth - 10\tabcolsep) * \real{0.1667}}
  >{\raggedright\arraybackslash}p{(\linewidth - 10\tabcolsep) * \real{0.1667}}
  >{\raggedright\arraybackslash}p{(\linewidth - 10\tabcolsep) * \real{0.1667}}@{}}
\toprule\noalign{}
\begin{minipage}[b]{\linewidth}\raggedright
No.
\end{minipage} & \begin{minipage}[b]{\linewidth}\raggedright
Name
\end{minipage} & \begin{minipage}[b]{\linewidth}\raggedright
Description
\end{minipage} & \begin{minipage}[b]{\linewidth}\raggedright
Default
\end{minipage} & \begin{minipage}[b]{\linewidth}\raggedright
\end{minipage} & \begin{minipage}[b]{\linewidth}\raggedright
\end{minipage} \\
\midrule\noalign{}
\endhead
\bottomrule\noalign{}
\endlastfoot
& & Field & Metric & Laboratory & 0 \\
1 & SALTMF & A real positive value that defines the salt liquid-phase
mole fraction for all grid blocks in the model for when the CO2-Brine
model has been activated. & & & \\
mole fraction & mole fraction & mole fraction & & & \\
Notes: & & & & & \\
\end{longtable}

\emph{Table 8.3.277.1: SALTMF Keyword Description}

See also the CO2STORE keyword in the RUNSPEC section.

\subsubsection{Example}\label{example-350}

The example sets the salt liquid-phase mole fraction for all cells in
the model to 0.018.

\begin{verbatim}
--
--       SET SALT LIQUID-PHASE MOLE FRACTION FOR ALL CELLS (OPM FLOW KEYWORD)
--
SALTMF
         0.018                                             /
\end{verbatim}

\subsection{SALTNODE -- Salt Concentration Based PVTNUM
Array}\label{saltnode-salt-concentration-based-pvtnum-array}

\begin{longtable}[]{@{}
  >{\raggedright\arraybackslash}p{(\linewidth - 14\tabcolsep) * \real{0.1250}}
  >{\raggedright\arraybackslash}p{(\linewidth - 14\tabcolsep) * \real{0.1250}}
  >{\raggedright\arraybackslash}p{(\linewidth - 14\tabcolsep) * \real{0.1250}}
  >{\raggedright\arraybackslash}p{(\linewidth - 14\tabcolsep) * \real{0.1250}}
  >{\raggedright\arraybackslash}p{(\linewidth - 14\tabcolsep) * \real{0.1250}}
  >{\raggedright\arraybackslash}p{(\linewidth - 14\tabcolsep) * \real{0.1250}}
  >{\raggedright\arraybackslash}p{(\linewidth - 14\tabcolsep) * \real{0.1250}}
  >{\raggedright\arraybackslash}p{(\linewidth - 14\tabcolsep) * \real{0.1250}}@{}}
\toprule\noalign{}
\begin{minipage}[b]{\linewidth}\raggedright
RUNSPEC
\end{minipage} & \begin{minipage}[b]{\linewidth}\raggedright
GRID
\end{minipage} & \begin{minipage}[b]{\linewidth}\raggedright
EDIT
\end{minipage} & \begin{minipage}[b]{\linewidth}\raggedright
PROPS
\end{minipage} & \begin{minipage}[b]{\linewidth}\raggedright
REGIONS
\end{minipage} & \begin{minipage}[b]{\linewidth}\raggedright
SOLUTION
\end{minipage} & \begin{minipage}[b]{\linewidth}\raggedright
SUMMARY
\end{minipage} & \begin{minipage}[b]{\linewidth}\raggedright
SCHEDULE
\end{minipage} \\
\midrule\noalign{}
\endhead
\bottomrule\noalign{}
\endlastfoot
\end{longtable}

\subsubsection{Description}\label{description-629}

SALTNODE defines the salt concentration value based on a cells PVTNUM
number. The SALTNODE property is used in the calculation of a polymer
viscosity when the polymer and the salt options has been activated by
the POLYMER and BRINE keywords in the RUNSPEC section. In the RUNSPEC
section the number of PVTNUM functions is declared by NTPVT variable on
the TABDIMS keyword and allocated to individual cells by the PVTNUM
property array in the REGIONS section. NPPVT on the TABDIMS keyword in
the RUNSPEC section defines the maximum number of rows (or pressure
values) in the PVT tables and also sets the maximum number of entries
for each SALNODE data set. The number of values for each data set must
correspond to the number of polymer solution adsorption entries on the
PLYADSS keyword. For example if there are three sets of PVT tables and
four values on the PLYADSS keyword, then three SALTNODE data sets with
four values of salt concentrations need to be entered.

This keyword is not supported by OPM Flow but would change the results
if supported so the simulation will be stopped.

\begin{longtable}[]{@{}
  >{\raggedright\arraybackslash}p{(\linewidth - 10\tabcolsep) * \real{0.1667}}
  >{\raggedright\arraybackslash}p{(\linewidth - 10\tabcolsep) * \real{0.1667}}
  >{\raggedright\arraybackslash}p{(\linewidth - 10\tabcolsep) * \real{0.1667}}
  >{\raggedright\arraybackslash}p{(\linewidth - 10\tabcolsep) * \real{0.1667}}
  >{\raggedright\arraybackslash}p{(\linewidth - 10\tabcolsep) * \real{0.1667}}
  >{\raggedright\arraybackslash}p{(\linewidth - 10\tabcolsep) * \real{0.1667}}@{}}
\toprule\noalign{}
\begin{minipage}[b]{\linewidth}\raggedright
No.
\end{minipage} & \begin{minipage}[b]{\linewidth}\raggedright
Name
\end{minipage} & \begin{minipage}[b]{\linewidth}\raggedright
Description
\end{minipage} & \begin{minipage}[b]{\linewidth}\raggedright
Default
\end{minipage} & \begin{minipage}[b]{\linewidth}\raggedright
\end{minipage} & \begin{minipage}[b]{\linewidth}\raggedright
\end{minipage} \\
\midrule\noalign{}
\endhead
\bottomrule\noalign{}
\endlastfoot
& & Field & Metric & Laboratory & \\
1 & SALTNODE & A real monotonically increasing positive columnar vector
defining the salt concentration for a given PVTNUM table. & None & & \\
lb/stb & kg/sm3 & gm/scc & & & \\
Notes: & & & & & \\
\end{longtable}

\emph{Table 8.141: SALTNODE Keyword Description}

An alternative manner of entering the salt concentrations is by
utilizing the PVTNUM region array by using the ADSALNOD keyword in the
PROPS section.

\subsubsection{Example}\label{example-351}

Given three sets of relative permeability tables and four values on the
PLYADSS keyword and two SALNODE data sets with four values of salt
concentrations then the data should be entered as follows:

\begin{verbatim}
--
--    SETS SALT CONCENTRATION FOR POLYMER SOLUTION ADSORPTION
--    VIA PVTNUM ARRAY ALLOCATION
--
--    SALT
--
SALTNODE
      1.0
      5.0
      10.5
      25.0      / PVTNUM TABLE NO. 01
      1.0
      3.0
      7.5
      15.0      / PVTNUM TABLE NO. 02
\end{verbatim}

See also the ADSALNOD keyword.

\subsection{SALTSOL -- Define the Salt Solubility Limit by
Region}\label{saltsol-define-the-salt-solubility-limit-by-region}

\begin{longtable}[]{@{}
  >{\raggedright\arraybackslash}p{(\linewidth - 14\tabcolsep) * \real{0.1250}}
  >{\raggedright\arraybackslash}p{(\linewidth - 14\tabcolsep) * \real{0.1250}}
  >{\raggedright\arraybackslash}p{(\linewidth - 14\tabcolsep) * \real{0.1250}}
  >{\raggedright\arraybackslash}p{(\linewidth - 14\tabcolsep) * \real{0.1250}}
  >{\raggedright\arraybackslash}p{(\linewidth - 14\tabcolsep) * \real{0.1250}}
  >{\raggedright\arraybackslash}p{(\linewidth - 14\tabcolsep) * \real{0.1250}}
  >{\raggedright\arraybackslash}p{(\linewidth - 14\tabcolsep) * \real{0.1250}}
  >{\raggedright\arraybackslash}p{(\linewidth - 14\tabcolsep) * \real{0.1250}}@{}}
\toprule\noalign{}
\begin{minipage}[b]{\linewidth}\raggedright
RUNSPEC
\end{minipage} & \begin{minipage}[b]{\linewidth}\raggedright
GRID
\end{minipage} & \begin{minipage}[b]{\linewidth}\raggedright
EDIT
\end{minipage} & \begin{minipage}[b]{\linewidth}\raggedright
PROPS
\end{minipage} & \begin{minipage}[b]{\linewidth}\raggedright
REGIONS
\end{minipage} & \begin{minipage}[b]{\linewidth}\raggedright
SOLUTION
\end{minipage} & \begin{minipage}[b]{\linewidth}\raggedright
SUMMARY
\end{minipage} & \begin{minipage}[b]{\linewidth}\raggedright
SCHEDULE
\end{minipage} \\
\midrule\noalign{}
\endhead
\bottomrule\noalign{}
\endlastfoot
\end{longtable}

\subsubsection{Description}\label{description-630}

SALTSOL defines a grid block's maximum salt solubility for each PVTNUM
region. The keyword should only be used with OPM Flow's Salt
Precipitation model which is activated via the PRECSALT keyword in the
RUNSPEC section.

\begin{tcolorbox}[enhanced jigsaw, title=\textcolor{quarto-callout-note-color}{\faInfo}\hspace{0.5em}{Note}, coltitle=black, breakable, toprule=.15mm, colbacktitle=quarto-callout-note-color!10!white, leftrule=.75mm, opacitybacktitle=0.6, colback=white, bottomtitle=1mm, toptitle=1mm, titlerule=0mm, left=2mm, colframe=quarto-callout-note-color-frame, arc=.35mm, rightrule=.15mm, bottomrule=.15mm, opacityback=0]

This is an OPM Flow specific keyword for the simulator's Salt
Precipitation model that is activated by the PRECSALT keyword and
declaring that vaporized water is present in the run via the VAPWAT in
the RUNSPEC section.

\end{tcolorbox}

\begin{longtable}[]{@{}
  >{\raggedright\arraybackslash}p{(\linewidth - 6\tabcolsep) * \real{0.2500}}
  >{\raggedright\arraybackslash}p{(\linewidth - 6\tabcolsep) * \real{0.2500}}
  >{\raggedright\arraybackslash}p{(\linewidth - 6\tabcolsep) * \real{0.2500}}
  >{\raggedright\arraybackslash}p{(\linewidth - 6\tabcolsep) * \real{0.2500}}@{}}
\toprule\noalign{}
\begin{minipage}[b]{\linewidth}\raggedright
No.
\end{minipage} & \begin{minipage}[b]{\linewidth}\raggedright
Name
\end{minipage} & \begin{minipage}[b]{\linewidth}\raggedright
Description
\end{minipage} & \begin{minipage}[b]{\linewidth}\raggedright
Default
\end{minipage} \\
\midrule\noalign{}
\endhead
\bottomrule\noalign{}
\endlastfoot
Field & Metric & Laboratory & \\
1 & SALTSOL & A real positive value that defines the maximum salt
solubility for all grid blocks in a PVTNUM region. & None \\
lb/stb & kg/sm3 & gm/scc & \\
‍2 ‍ & SALTDEN & SALTDEN is a real number defining the density of salt at
surface conditions. & Defined \\
lb/ft3 135.469 & kg/m3 2170 & gm/cc 2.170 & \\
Notes: & & & \\
\end{longtable}

\emph{Table 8.142: SALTSOL Keyword Description}

See also the PRECSALT and VAPWAT keywords in the RUNSPEC section and the
PVTGW and PVTGWO keywords in the PROPS section.

\subsubsection{Example}\label{example-352}

The first example sets the maximum salt solubility for all cells in the
model to 134.6 lb/stb, assuming that there is only one PVT region, that
is NTPVT is equal to one on the TABDIMS keyword in the RUNSPEC section.

\begin{verbatim}
--
--       SET SALT SOLUBILITY LIMIT FOR EACH REGION (OPM FLOW KEYWORD)
--
SALTSOL
--       MAX       SALT
--       SALTSOL   DENSITY
         134.6     1*                                             /
\end{verbatim}

The 134.6 lb/stb, (380 kg/sm3 or 0.384 gm/scc for metric and laboratory
units, respectively) is based on the solubility of NACL at 212 oF (100
oC) and should be used with care.

The next example shows how to set the maximum salt solubility for when
NTPVT is equal to three on the TABDIMS keyword in the RUNSPEC section.

\begin{verbatim}
--
--       SET SALT SOLUBILITY LIMIT FOR EACH REGION (OPM FLOW KEYWORD)
--
SALTSOL
         134.6                                             / PVT REGION NO. 1
         124.0                                             / PVT REGION NO. 2
                                                           / PVT REGION NO. 3
\end{verbatim}

Here the last entry, which is for region number three, is defaulted, and
results in region's three maximum salt solubility to take the previous
value, in this case 124.0 lb/stb.

\subsection{SCALECRS -- Define End-Point Scaling
Option}\label{scalecrs-define-end-point-scaling-option}

\begin{longtable}[]{@{}
  >{\raggedright\arraybackslash}p{(\linewidth - 14\tabcolsep) * \real{0.1250}}
  >{\raggedright\arraybackslash}p{(\linewidth - 14\tabcolsep) * \real{0.1250}}
  >{\raggedright\arraybackslash}p{(\linewidth - 14\tabcolsep) * \real{0.1250}}
  >{\raggedright\arraybackslash}p{(\linewidth - 14\tabcolsep) * \real{0.1250}}
  >{\raggedright\arraybackslash}p{(\linewidth - 14\tabcolsep) * \real{0.1250}}
  >{\raggedright\arraybackslash}p{(\linewidth - 14\tabcolsep) * \real{0.1250}}
  >{\raggedright\arraybackslash}p{(\linewidth - 14\tabcolsep) * \real{0.1250}}
  >{\raggedright\arraybackslash}p{(\linewidth - 14\tabcolsep) * \real{0.1250}}@{}}
\toprule\noalign{}
\begin{minipage}[b]{\linewidth}\raggedright
RUNSPEC
\end{minipage} & \begin{minipage}[b]{\linewidth}\raggedright
GRID
\end{minipage} & \begin{minipage}[b]{\linewidth}\raggedright
EDIT
\end{minipage} & \begin{minipage}[b]{\linewidth}\raggedright
PROPS
\end{minipage} & \begin{minipage}[b]{\linewidth}\raggedright
REGIONS
\end{minipage} & \begin{minipage}[b]{\linewidth}\raggedright
SOLUTION
\end{minipage} & \begin{minipage}[b]{\linewidth}\raggedright
SUMMARY
\end{minipage} & \begin{minipage}[b]{\linewidth}\raggedright
SCHEDULE
\end{minipage} \\
\midrule\noalign{}
\endhead
\bottomrule\noalign{}
\endlastfoot
\end{longtable}

\subsubsection{Description}\label{description-631}

The SCALECRS keyword sets the end-point scaling option to be either
two-point or three-point scaling, for when the End-Point Scaling option
has been invoked by the ENDSCALE keyword in the RUNSPEC section. This
determines which end-points on the relative permeability curves are used
for scaling based on the supplied end-point arrays (SGCR, SWCR, etc.).

\begin{longtable}[]{@{}
  >{\raggedright\arraybackslash}p{(\linewidth - 6\tabcolsep) * \real{0.2500}}
  >{\raggedright\arraybackslash}p{(\linewidth - 6\tabcolsep) * \real{0.2500}}
  >{\raggedright\arraybackslash}p{(\linewidth - 6\tabcolsep) * \real{0.2500}}
  >{\raggedright\arraybackslash}p{(\linewidth - 6\tabcolsep) * \real{0.2500}}@{}}
\toprule\noalign{}
\begin{minipage}[b]{\linewidth}\raggedright
No.
\end{minipage} & \begin{minipage}[b]{\linewidth}\raggedright
Name
\end{minipage} & \begin{minipage}[b]{\linewidth}\raggedright
Description
\end{minipage} & \begin{minipage}[b]{\linewidth}\raggedright
Default
\end{minipage} \\
\midrule\noalign{}
\endhead
\bottomrule\noalign{}
\endlastfoot
1 & SCALEOPT & SCALEOPT is a character string that sets the end-point
scaling option and should be set to either NO or YES: & NO \\
Notes: & & & \\
\end{longtable}

\emph{Table 8.143: SCALECRS Keyword Description}

The end-point scaled for each option and the arrays used in the
end-point scaling are summarized in the following table:

\begin{longtable}[]{@{}
  >{\raggedright\arraybackslash}p{(\linewidth - 10\tabcolsep) * \real{0.1667}}
  >{\raggedright\arraybackslash}p{(\linewidth - 10\tabcolsep) * \real{0.1667}}
  >{\raggedright\arraybackslash}p{(\linewidth - 10\tabcolsep) * \real{0.1667}}
  >{\raggedright\arraybackslash}p{(\linewidth - 10\tabcolsep) * \real{0.1667}}
  >{\raggedright\arraybackslash}p{(\linewidth - 10\tabcolsep) * \real{0.1667}}
  >{\raggedright\arraybackslash}p{(\linewidth - 10\tabcolsep) * \real{0.1667}}@{}}
\toprule\noalign{}
\begin{minipage}[b]{\linewidth}\raggedright
Option
\end{minipage} & \begin{minipage}[b]{\linewidth}\raggedright
Phases
\end{minipage} & \begin{minipage}[b]{\linewidth}\raggedright
Relative Permeability End-Point
\end{minipage} & \begin{minipage}[b]{\linewidth}\raggedright
Minimum Saturation End-Point
\end{minipage} & \begin{minipage}[b]{\linewidth}\raggedright
Middle Saturation End-Point
\end{minipage} & \begin{minipage}[b]{\linewidth}\raggedright
Maximum Saturation End-Point
\end{minipage} \\
\midrule\noalign{}
\endhead
\bottomrule\noalign{}
\endlastfoot
Two-Point & Water & KRW & SWCR & & SWU \\
Gas & KRG & SGCR & & SGU & \\
Oil-Water & KRORW & SOWCR & & (1.0 -- SWL - SGL) & \\
Oil-Gas & KRORG & SOGCR & & (1.0 -- SWL - SGL) & \\
Three-Point & Water & KRW & SWCR & (1.0 -- SOWCR - SGL) & SWU \\
Gas & KRG & SGCR & (1.0 - SOGCR-SWL) & SGU & \\
Oil-Water & KRORW & SOWCR & (1.0 -- SWCR - SGL) & (1.0 -- SWL - SGL)
& \\
Oil-Gas & KRORG & SOGCR & (1.0 -- SGCR - SGL) & (1.0 -- SWL - SGL) & \\
Two Phase Gas-Water Simulations & & & & & \\
Water & KRW & SWCR & (1.0 - SGCR) & SWU & \\
Gas & KRG & SGCR & (1.0 -SWCR) & SGU & \\
\end{longtable}

\emph{Table 8.144: End-Point Arrays Used in the End-Point Scaling
Options}

See also the TZONE keyword in the PROPS section that sets the transition
zone end-point scaling options for the oil, gas and water phases.

\subsubsection{Example}\label{example-353}

\begin{verbatim}
--
--       TWO-POINT END-POINT SCALING IS NO THREE POINT IS YES
--
--       SCALEOPT
--       ---------
SCALECRS
         YES                                               / SCALING OPTION
\end{verbatim}

The above example activates three-point end-point scaling of the
relative permeability curves.

\subsection{SCALELIM -- End-Point Scaling versus Depth Maximum Water
Saturation}\label{scalelim-end-point-scaling-versus-depth-maximum-water-saturation}

\begin{longtable}[]{@{}
  >{\raggedright\arraybackslash}p{(\linewidth - 14\tabcolsep) * \real{0.1250}}
  >{\raggedright\arraybackslash}p{(\linewidth - 14\tabcolsep) * \real{0.1250}}
  >{\raggedright\arraybackslash}p{(\linewidth - 14\tabcolsep) * \real{0.1250}}
  >{\raggedright\arraybackslash}p{(\linewidth - 14\tabcolsep) * \real{0.1250}}
  >{\raggedright\arraybackslash}p{(\linewidth - 14\tabcolsep) * \real{0.1250}}
  >{\raggedright\arraybackslash}p{(\linewidth - 14\tabcolsep) * \real{0.1250}}
  >{\raggedright\arraybackslash}p{(\linewidth - 14\tabcolsep) * \real{0.1250}}
  >{\raggedright\arraybackslash}p{(\linewidth - 14\tabcolsep) * \real{0.1250}}@{}}
\toprule\noalign{}
\begin{minipage}[b]{\linewidth}\raggedright
RUNSPEC
\end{minipage} & \begin{minipage}[b]{\linewidth}\raggedright
GRID
\end{minipage} & \begin{minipage}[b]{\linewidth}\raggedright
EDIT
\end{minipage} & \begin{minipage}[b]{\linewidth}\raggedright
PROPS
\end{minipage} & \begin{minipage}[b]{\linewidth}\raggedright
REGIONS
\end{minipage} & \begin{minipage}[b]{\linewidth}\raggedright
SOLUTION
\end{minipage} & \begin{minipage}[b]{\linewidth}\raggedright
SUMMARY
\end{minipage} & \begin{minipage}[b]{\linewidth}\raggedright
SCHEDULE
\end{minipage} \\
\midrule\noalign{}
\endhead
\bottomrule\noalign{}
\endlastfoot
\end{longtable}

\subsubsection{Description}\label{description-632}

This keyword defines the maximum water saturation allowed in a cell for
when the end-point versus depth tables are used in the End-Point Scaling
option to calculate the water saturation for a grid block. The End-Point
Scaling option must be invoked by the ENDSCALE keyword in the RUNSPEC
section to use this keyword, and the keyword may only be used in two
phase runs containing water, or if the Miscible Flood option has been
activated by the MISCIBLE keyword in the RUNSPEC section. This keyword
functionality is not supported in OPM Flow.

This keyword is not supported by OPM Flow but would change the results
if supported so the simulation will be stopped.

\subsection{SDENSITY -- Define the Miscible or Solvent Surface Gas
Density}\label{sdensity-define-the-miscible-or-solvent-surface-gas-density}

\begin{longtable}[]{@{}
  >{\raggedright\arraybackslash}p{(\linewidth - 14\tabcolsep) * \real{0.1250}}
  >{\raggedright\arraybackslash}p{(\linewidth - 14\tabcolsep) * \real{0.1250}}
  >{\raggedright\arraybackslash}p{(\linewidth - 14\tabcolsep) * \real{0.1250}}
  >{\raggedright\arraybackslash}p{(\linewidth - 14\tabcolsep) * \real{0.1250}}
  >{\raggedright\arraybackslash}p{(\linewidth - 14\tabcolsep) * \real{0.1250}}
  >{\raggedright\arraybackslash}p{(\linewidth - 14\tabcolsep) * \real{0.1250}}
  >{\raggedright\arraybackslash}p{(\linewidth - 14\tabcolsep) * \real{0.1250}}
  >{\raggedright\arraybackslash}p{(\linewidth - 14\tabcolsep) * \real{0.1250}}@{}}
\toprule\noalign{}
\begin{minipage}[b]{\linewidth}\raggedright
RUNSPEC
\end{minipage} & \begin{minipage}[b]{\linewidth}\raggedright
GRID
\end{minipage} & \begin{minipage}[b]{\linewidth}\raggedright
EDIT
\end{minipage} & \begin{minipage}[b]{\linewidth}\raggedright
PROPS
\end{minipage} & \begin{minipage}[b]{\linewidth}\raggedright
REGIONS
\end{minipage} & \begin{minipage}[b]{\linewidth}\raggedright
SOLUTION
\end{minipage} & \begin{minipage}[b]{\linewidth}\raggedright
SUMMARY
\end{minipage} & \begin{minipage}[b]{\linewidth}\raggedright
SCHEDULE
\end{minipage} \\
\midrule\noalign{}
\endhead
\bottomrule\noalign{}
\endlastfoot
\end{longtable}

\subsubsection{Description}\label{description-633}

The SDENSITY keyword defines density at surface conditions of either the
miscible injection gas for when the MISCIBLE option has been invoked in
the RUNSPEC section, or the solvent for when the SOLVENT option has been
invoked in the RUNSPEC section. This keyword must be invoked if either
the MISCIBLE or SOLVENT options have been activated in the RUNSPEC
section.

\begin{longtable}[]{@{}
  >{\raggedright\arraybackslash}p{(\linewidth - 6\tabcolsep) * \real{0.2500}}
  >{\raggedright\arraybackslash}p{(\linewidth - 6\tabcolsep) * \real{0.2500}}
  >{\raggedright\arraybackslash}p{(\linewidth - 6\tabcolsep) * \real{0.2500}}
  >{\raggedright\arraybackslash}p{(\linewidth - 6\tabcolsep) * \real{0.2500}}@{}}
\toprule\noalign{}
\begin{minipage}[b]{\linewidth}\raggedright
No.
\end{minipage} & \begin{minipage}[b]{\linewidth}\raggedright
Name
\end{minipage} & \begin{minipage}[b]{\linewidth}\raggedright
Description
\end{minipage} & \begin{minipage}[b]{\linewidth}\raggedright
Default
\end{minipage} \\
\midrule\noalign{}
\endhead
\bottomrule\noalign{}
\endlastfoot
Field & Metric & Laboratory & \\
1 & SDENSITY & SDENSITY is a real positive number defining the density
at surface conditions of either: & None \\
lb/ft3 & kg/m3 & gm/cc & \\
Notes: & & & \\
\end{longtable}

\emph{Table 8.145: SDENSITY Keyword Description}

In addition to this keyword, the surface density or gravity of the
in-place natural gas must be entered using either the DENSITY or GRAVITY
keywords.

\subsubsection{Examples}\label{examples-43}

The following shows the SDENSITY keyword for when NTPVT on the TABDIMS
keyword in the RUNSPEC section is set to one.

\begin{verbatim}
--
--       MIS-SOL
--       DENSITY
--       -------
SDENSITY
         0.04520                                           / MIS-SOL DENSITY
\end{verbatim}

The next example shows the SDENSITY keyword for when NTPVT on the
TABDIMS keyword in the RUNSPEC section is set to three.

\begin{verbatim}
--
--       MIS-SOL
--       DENSITY
--       -------
SDENSITY
         0.04520                                           / MIS-SOL DENSITY 1
         0.05520                                           / MIS-SOL DENSITY 2
         0.06420                                           / MIS-SOL DENSITY 3
\end{verbatim}

There is no terminating ``/'' for this keyword.

\subsection{SGCR -- End-Point Scaling Grid Cell Critical Gas
Saturations}\label{sgcr-end-point-scaling-grid-cell-critical-gas-saturations}

\begin{longtable}[]{@{}
  >{\raggedright\arraybackslash}p{(\linewidth - 14\tabcolsep) * \real{0.1250}}
  >{\raggedright\arraybackslash}p{(\linewidth - 14\tabcolsep) * \real{0.1250}}
  >{\raggedright\arraybackslash}p{(\linewidth - 14\tabcolsep) * \real{0.1250}}
  >{\raggedright\arraybackslash}p{(\linewidth - 14\tabcolsep) * \real{0.1250}}
  >{\raggedright\arraybackslash}p{(\linewidth - 14\tabcolsep) * \real{0.1250}}
  >{\raggedright\arraybackslash}p{(\linewidth - 14\tabcolsep) * \real{0.1250}}
  >{\raggedright\arraybackslash}p{(\linewidth - 14\tabcolsep) * \real{0.1250}}
  >{\raggedright\arraybackslash}p{(\linewidth - 14\tabcolsep) * \real{0.1250}}@{}}
\toprule\noalign{}
\begin{minipage}[b]{\linewidth}\raggedright
RUNSPEC
\end{minipage} & \begin{minipage}[b]{\linewidth}\raggedright
GRID
\end{minipage} & \begin{minipage}[b]{\linewidth}\raggedright
EDIT
\end{minipage} & \begin{minipage}[b]{\linewidth}\raggedright
PROPS
\end{minipage} & \begin{minipage}[b]{\linewidth}\raggedright
REGIONS
\end{minipage} & \begin{minipage}[b]{\linewidth}\raggedright
SOLUTION
\end{minipage} & \begin{minipage}[b]{\linewidth}\raggedright
SUMMARY
\end{minipage} & \begin{minipage}[b]{\linewidth}\raggedright
SCHEDULE
\end{minipage} \\
\midrule\noalign{}
\endhead
\bottomrule\noalign{}
\endlastfoot
\end{longtable}

\subsubsection{Description}\label{description-634}

SGCR defines the critical gas saturation for all the cells in the model
via an array when the end-point scaling option has been invoked via the
ENDSCALE keyword in the RUNSPEC section. The critical gas saturation is
defined as the maximum gas saturation for which the gas relative
permeability is zero in a two-phase relative permeability table.

The keyword can be used with all grid types.

\begin{longtable}[]{@{}
  >{\raggedright\arraybackslash}p{(\linewidth - 6\tabcolsep) * \real{0.2500}}
  >{\raggedright\arraybackslash}p{(\linewidth - 6\tabcolsep) * \real{0.2500}}
  >{\raggedright\arraybackslash}p{(\linewidth - 6\tabcolsep) * \real{0.2500}}
  >{\raggedright\arraybackslash}p{(\linewidth - 6\tabcolsep) * \real{0.2500}}@{}}
\toprule\noalign{}
\begin{minipage}[b]{\linewidth}\raggedright
No.
\end{minipage} & \begin{minipage}[b]{\linewidth}\raggedright
Name
\end{minipage} & \begin{minipage}[b]{\linewidth}\raggedright
Description
\end{minipage} & \begin{minipage}[b]{\linewidth}\raggedright
Default
\end{minipage} \\
\midrule\noalign{}
\endhead
\bottomrule\noalign{}
\endlastfoot
Field & Metric & Laboratory & \\
1 & SGCR & SGCR is an array of real numbers assigning the critical gas
saturation values to each cell in the model. The number of entries
should correspond to the NX x NY x NZ parameters on the DIMENS keyword.
Repeat counts may be used, for example 30*0.03 & Taken from cell
allocated relative permeability table. \\
dimensionless & dimensionless & dimensionless & \\
Notes: & & & \\
\end{longtable}

\emph{Table 8.146: SGCR Keyword Description}

End-point scaling allows the entered relative permeability functions to
be re-scaled based on the saturation end-points allocated to each cell
by the SWL, SWCR, SWU, SGL, SGCR, SGU, SOWCR, and SOGCR saturation grid
arrays for the saturation end-points, and the KRG, KRGR, KRO, KRORG,
KRORW, KRW and KRWR relative permeability grid cell arrays for the
relative permeability end-point data. In addition end-point scaling may
be directional dependent in which case the directional dependent
versions of the aforementioned arrays should be used, that is SGCRX,
SGCRY and SGCRZ instead of SGCR. There is also the facility to make the
directional end-point scaling reversible or non-reversible and if the
non-reversible option is selected the non-reversible versions of the
aforementioned arrays should be used, that is SGCRX, SGCRX-, SGCRY,
SGCRY-, SGCRZ and SGCRZ-, instead of the SGCR keyword.

\subsubsection{Example}\label{example-354}

\begin{verbatim}
--
--       DEFINE GRID BLOCK END-POINT SGCR DATA FOR ALL CELLS
--       (FOR NX x NY x NZ = 300)
--
SGCR
         300*0.050                                                              /
\end{verbatim}

The above example defines a constant critical gas saturation of 0.05 to
all 300 cells in the model as defined by the DIMENS keyword in the
RUNSPEC section.

\subsection{SGCWMIS -- Miscible Critical Gas versus Water Saturation
Functions}\label{sgcwmis-miscible-critical-gas-versus-water-saturation-functions}

\begin{longtable}[]{@{}
  >{\raggedright\arraybackslash}p{(\linewidth - 14\tabcolsep) * \real{0.1250}}
  >{\raggedright\arraybackslash}p{(\linewidth - 14\tabcolsep) * \real{0.1250}}
  >{\raggedright\arraybackslash}p{(\linewidth - 14\tabcolsep) * \real{0.1250}}
  >{\raggedright\arraybackslash}p{(\linewidth - 14\tabcolsep) * \real{0.1250}}
  >{\raggedright\arraybackslash}p{(\linewidth - 14\tabcolsep) * \real{0.1250}}
  >{\raggedright\arraybackslash}p{(\linewidth - 14\tabcolsep) * \real{0.1250}}
  >{\raggedright\arraybackslash}p{(\linewidth - 14\tabcolsep) * \real{0.1250}}
  >{\raggedright\arraybackslash}p{(\linewidth - 14\tabcolsep) * \real{0.1250}}@{}}
\toprule\noalign{}
\begin{minipage}[b]{\linewidth}\raggedright
RUNSPEC
\end{minipage} & \begin{minipage}[b]{\linewidth}\raggedright
GRID
\end{minipage} & \begin{minipage}[b]{\linewidth}\raggedright
EDIT
\end{minipage} & \begin{minipage}[b]{\linewidth}\raggedright
PROPS
\end{minipage} & \begin{minipage}[b]{\linewidth}\raggedright
REGIONS
\end{minipage} & \begin{minipage}[b]{\linewidth}\raggedright
SOLUTION
\end{minipage} & \begin{minipage}[b]{\linewidth}\raggedright
SUMMARY
\end{minipage} & \begin{minipage}[b]{\linewidth}\raggedright
SCHEDULE
\end{minipage} \\
\midrule\noalign{}
\endhead
\bottomrule\noalign{}
\endlastfoot
\end{longtable}

\subsubsection{Description}\label{description-635}

SGCWMIS defines the dependency between the miscible critical gas
saturation and the water saturation, for when the MISCIBLE keyword in
the RUNSPEC section has been activated. The keyword can only be used
with the MISCIBLE option and for when the oil, water and gas phases are
active in the model.

\begin{longtable}[]{@{}
  >{\raggedright\arraybackslash}p{(\linewidth - 6\tabcolsep) * \real{0.2500}}
  >{\raggedright\arraybackslash}p{(\linewidth - 6\tabcolsep) * \real{0.2500}}
  >{\raggedright\arraybackslash}p{(\linewidth - 6\tabcolsep) * \real{0.2500}}
  >{\raggedright\arraybackslash}p{(\linewidth - 6\tabcolsep) * \real{0.2500}}@{}}
\toprule\noalign{}
\begin{minipage}[b]{\linewidth}\raggedright
No.
\end{minipage} & \begin{minipage}[b]{\linewidth}\raggedright
Name
\end{minipage} & \begin{minipage}[b]{\linewidth}\raggedright
Description
\end{minipage} & \begin{minipage}[b]{\linewidth}\raggedright
Default
\end{minipage} \\
\midrule\noalign{}
\endhead
\bottomrule\noalign{}
\endlastfoot
Field & Metric & Laboratory & \\
1 & SWAT & A columnar vector of real monotonically increasing down the
column values starting from zero and terminating atone, that defines the
water saturation. & None \\
dimensionless & dimensionless & dimensionless & \\
2 & SGCMIS & A columnar vector of real equal or increasing down the
column values that are greater than or equal to zero and less than one,
that define the corresponding miscible gas critical gas saturation for
the corresponding water saturation SWAT. & None \\
dimensionless & dimensionless & dimensionless & \\
Notes: & & & \\
\end{longtable}

\emph{Table 8.147: SGCWMIS Keyword Description}

\subsubsection{Example}\label{example-355}

\begin{verbatim}
--
--       MISCIBLE CRITICAL GAS VERSUS WATER SATURATION TABLE
--
SGCWMIS
--       SWAT       SGCRMIS
--       FRAC       FRAC
--       -------    --------
          0.0000     0.0000
          0.2000     0.0300
          1.0000     0.0300                                / TABLE NO. 01
--       SWAT       SGCRMIS
--       FRAC       FRAC
--       -------    --------
          0.0000     0.0000
          0.3000     0.0500
          1.0000     0.0500                                / TABLE NO. 02
\end{verbatim}

The above example defines two miscible critical gas saturation versus
water saturation tables assuming NTMISC equals two and NSMISC is greater
than or equal to three on the MISCIBLE keyword in the RUNSPEC section.

\subsection{SGF32D -- Gas Saturation Tables versus Oil and Water
Saturations}\label{sgf32d-gas-saturation-tables-versus-oil-and-water-saturations}

\begin{longtable}[]{@{}
  >{\raggedright\arraybackslash}p{(\linewidth - 14\tabcolsep) * \real{0.1250}}
  >{\raggedright\arraybackslash}p{(\linewidth - 14\tabcolsep) * \real{0.1250}}
  >{\raggedright\arraybackslash}p{(\linewidth - 14\tabcolsep) * \real{0.1250}}
  >{\raggedright\arraybackslash}p{(\linewidth - 14\tabcolsep) * \real{0.1250}}
  >{\raggedright\arraybackslash}p{(\linewidth - 14\tabcolsep) * \real{0.1250}}
  >{\raggedright\arraybackslash}p{(\linewidth - 14\tabcolsep) * \real{0.1250}}
  >{\raggedright\arraybackslash}p{(\linewidth - 14\tabcolsep) * \real{0.1250}}
  >{\raggedright\arraybackslash}p{(\linewidth - 14\tabcolsep) * \real{0.1250}}@{}}
\toprule\noalign{}
\begin{minipage}[b]{\linewidth}\raggedright
RUNSPEC
\end{minipage} & \begin{minipage}[b]{\linewidth}\raggedright
GRID
\end{minipage} & \begin{minipage}[b]{\linewidth}\raggedright
EDIT
\end{minipage} & \begin{minipage}[b]{\linewidth}\raggedright
PROPS
\end{minipage} & \begin{minipage}[b]{\linewidth}\raggedright
REGIONS
\end{minipage} & \begin{minipage}[b]{\linewidth}\raggedright
SOLUTION
\end{minipage} & \begin{minipage}[b]{\linewidth}\raggedright
SUMMARY
\end{minipage} & \begin{minipage}[b]{\linewidth}\raggedright
SCHEDULE
\end{minipage} \\
\midrule\noalign{}
\endhead
\bottomrule\noalign{}
\endlastfoot
\end{longtable}

\subsubsection{Description}\label{description-636}

The SGF32D keyword defines the gas relative permeability as a function
of both oil and water saturations. This keyword should only be used if
the gas is present in the run.

This keyword is not supported by OPM Flow but would change the results
if supported so the simulation will be stopped.

See also the SWOF, SGOF, SLGOF series of keywords and the SWFN, SGFN,
SGWFN, SOF2, SOF3, SOF32D series of keywords to enter relative
permeability data.

\subsection{SGFN -- Gas Saturation Tables (Format Type
2)}\label{sgfn-gas-saturation-tables-format-type-2}

\begin{longtable}[]{@{}
  >{\raggedright\arraybackslash}p{(\linewidth - 14\tabcolsep) * \real{0.1250}}
  >{\raggedright\arraybackslash}p{(\linewidth - 14\tabcolsep) * \real{0.1250}}
  >{\raggedright\arraybackslash}p{(\linewidth - 14\tabcolsep) * \real{0.1250}}
  >{\raggedright\arraybackslash}p{(\linewidth - 14\tabcolsep) * \real{0.1250}}
  >{\raggedright\arraybackslash}p{(\linewidth - 14\tabcolsep) * \real{0.1250}}
  >{\raggedright\arraybackslash}p{(\linewidth - 14\tabcolsep) * \real{0.1250}}
  >{\raggedright\arraybackslash}p{(\linewidth - 14\tabcolsep) * \real{0.1250}}
  >{\raggedright\arraybackslash}p{(\linewidth - 14\tabcolsep) * \real{0.1250}}@{}}
\toprule\noalign{}
\begin{minipage}[b]{\linewidth}\raggedright
RUNSPEC
\end{minipage} & \begin{minipage}[b]{\linewidth}\raggedright
GRID
\end{minipage} & \begin{minipage}[b]{\linewidth}\raggedright
EDIT
\end{minipage} & \begin{minipage}[b]{\linewidth}\raggedright
PROPS
\end{minipage} & \begin{minipage}[b]{\linewidth}\raggedright
REGIONS
\end{minipage} & \begin{minipage}[b]{\linewidth}\raggedright
SOLUTION
\end{minipage} & \begin{minipage}[b]{\linewidth}\raggedright
SUMMARY
\end{minipage} & \begin{minipage}[b]{\linewidth}\raggedright
SCHEDULE
\end{minipage} \\
\midrule\noalign{}
\endhead
\bottomrule\noalign{}
\endlastfoot
\end{longtable}

\subsubsection{Description}\label{description-637}

The SGFN keyword defines the gas relative permeability and oil-gas
capillary pressure data versus gas saturation tables for when gas is
present in the input deck. This keyword should only be used if the gas
is present in the run.

\begin{longtable}[]{@{}
  >{\raggedright\arraybackslash}p{(\linewidth - 6\tabcolsep) * \real{0.2500}}
  >{\raggedright\arraybackslash}p{(\linewidth - 6\tabcolsep) * \real{0.2500}}
  >{\raggedright\arraybackslash}p{(\linewidth - 6\tabcolsep) * \real{0.2500}}
  >{\raggedright\arraybackslash}p{(\linewidth - 6\tabcolsep) * \real{0.2500}}@{}}
\toprule\noalign{}
\begin{minipage}[b]{\linewidth}\raggedright
No.
\end{minipage} & \begin{minipage}[b]{\linewidth}\raggedright
Name
\end{minipage} & \begin{minipage}[b]{\linewidth}\raggedright
Description
\end{minipage} & \begin{minipage}[b]{\linewidth}\raggedright
Default
\end{minipage} \\
\midrule\noalign{}
\endhead
\bottomrule\noalign{}
\endlastfoot
Field & Metric & Laboratory & \\
1 & SGAS & A columnar vector of real monotonically increasing down the
column values starting from zero and terminating at one, that defines
the gas saturation. & None \\
dimensionless & dimensionless & dimensionless & \\
2 & KRG & A columnar vector of real values that are either equal or
increasing down the column and that are greater than or equal to zero
and less than or equal to one that defines the gas relative
permeability. & None \\
dimensionless & dimensionless & dimensionless & \\
3 & PCOG & A columnar vector of real values that are either equal or
increasing down the column that defines the oil-gas capillary pressure.
& None \\
psia & bars & atm & \\
Notes: & & & \\
\end{longtable}

\emph{Table 8.148: SGFN Keyword Description}

See also the GSF - Gas Saturation Tables versus Gas Saturation
(Gas-Water and CO2STORE Systems) keyword in the PROPS section, that
defines the gas relative permeability and gas-water capillary pressure
data as a function of gas saturation, for when only the gas and water
phases are present in the model. Note that the GSF keyword must be used
in conjunction with the WSF - Water Saturation Tables versus Water
Saturation (Gas-Water and CO2STORE Systems) keyword, also in the PROPS
section. WSF defines the water saturation as a function of water
saturation for when only the gas and water phases are present in the
model.

\subsubsection{Example}\label{example-356}

\begin{verbatim}
--
--       GAS RELATIVE PERMEABILITY TABLES (SGFN)
--
SGFN
--       SGAS       KRG        PCGO
--       FRAC                  PSIA
--       --------  --------   -------
             0.00    0.0000      0.0
             0.05    0.0000      0.0
             0.10    0.0000      0.0
             0.15    0.0000      0.0
             0.20    0.0002      0.0
             0.25    0.0010      0.0
             0.30    0.0062      0.0
             0.35    0.0140      0.0
             0.40    0.0273      0.0
             0.45    0.0450      0.0
             0.50    0.0707      0.0
             0.55    0.1020      0.0
             0.60    0.1412      0.0
             0.65    0.1870      0.0
             0.70    0.2412      0.0
             0.77    0.3288      0.0
             0.82    0.4000      0.0
             0.85    0.4450      0.0                       / TABLE NO. 01
--       SGAS       KRG        PCGO
--       FRAC                  PSIA
--       --------  --------   -------
             0.00    0.0000      0.0
             0.05    0.0000      0.0
             0.10    0.0000      0.0
             0.15    0.0000      0.0
             0.20    0.0002      0.0
             0.25    0.0010      0.0
             0.30    0.0062      0.0
             0.35    0.0140      0.0
             0.40    0.0273      0.0
             0.45    0.0450      0.0
             0.50    0.0707      0.0
             0.55    0.1020      0.0
             0.60    0.1412      0.0
             0.65    0.1870      0.0
             0.70    0.2412      0.0
             0.77    0.3288      0.0
             0.82    0.4000      0.0
             0.85    0.4450      0.0                       / TABLE NO. 02
\end{verbatim}

The example defines two SGFN tables for when gas is present in the input
deck.

\subsection{SGL -- End-Point Scaling Grid Cell Connate Gas
Saturations}\label{sgl-end-point-scaling-grid-cell-connate-gas-saturations}

\begin{longtable}[]{@{}
  >{\raggedright\arraybackslash}p{(\linewidth - 14\tabcolsep) * \real{0.1250}}
  >{\raggedright\arraybackslash}p{(\linewidth - 14\tabcolsep) * \real{0.1250}}
  >{\raggedright\arraybackslash}p{(\linewidth - 14\tabcolsep) * \real{0.1250}}
  >{\raggedright\arraybackslash}p{(\linewidth - 14\tabcolsep) * \real{0.1250}}
  >{\raggedright\arraybackslash}p{(\linewidth - 14\tabcolsep) * \real{0.1250}}
  >{\raggedright\arraybackslash}p{(\linewidth - 14\tabcolsep) * \real{0.1250}}
  >{\raggedright\arraybackslash}p{(\linewidth - 14\tabcolsep) * \real{0.1250}}
  >{\raggedright\arraybackslash}p{(\linewidth - 14\tabcolsep) * \real{0.1250}}@{}}
\toprule\noalign{}
\begin{minipage}[b]{\linewidth}\raggedright
RUNSPEC
\end{minipage} & \begin{minipage}[b]{\linewidth}\raggedright
GRID
\end{minipage} & \begin{minipage}[b]{\linewidth}\raggedright
EDIT
\end{minipage} & \begin{minipage}[b]{\linewidth}\raggedright
PROPS
\end{minipage} & \begin{minipage}[b]{\linewidth}\raggedright
REGIONS
\end{minipage} & \begin{minipage}[b]{\linewidth}\raggedright
SOLUTION
\end{minipage} & \begin{minipage}[b]{\linewidth}\raggedright
SUMMARY
\end{minipage} & \begin{minipage}[b]{\linewidth}\raggedright
SCHEDULE
\end{minipage} \\
\midrule\noalign{}
\endhead
\bottomrule\noalign{}
\endlastfoot
\end{longtable}

\subsubsection{Description}\label{description-638}

SGL defines the connate gas saturation for all the cells in the model
via an array when the end-point scaling option has been invoked via the
ENDSCALE keyword in the RUNSPEC section. The connate gas saturation is
defined as the minimum gas saturation in a two-phase gas relative
permeability table.

The keyword can be used with all grid types.

\begin{longtable}[]{@{}
  >{\raggedright\arraybackslash}p{(\linewidth - 6\tabcolsep) * \real{0.2500}}
  >{\raggedright\arraybackslash}p{(\linewidth - 6\tabcolsep) * \real{0.2500}}
  >{\raggedright\arraybackslash}p{(\linewidth - 6\tabcolsep) * \real{0.2500}}
  >{\raggedright\arraybackslash}p{(\linewidth - 6\tabcolsep) * \real{0.2500}}@{}}
\toprule\noalign{}
\begin{minipage}[b]{\linewidth}\raggedright
No.
\end{minipage} & \begin{minipage}[b]{\linewidth}\raggedright
Name
\end{minipage} & \begin{minipage}[b]{\linewidth}\raggedright
Description
\end{minipage} & \begin{minipage}[b]{\linewidth}\raggedright
Default
\end{minipage} \\
\midrule\noalign{}
\endhead
\bottomrule\noalign{}
\endlastfoot
Field & Metric & Laboratory & \\
1 & SGL & SGL is an array of real numbers assigning the connate gas
saturation values to each cell in the model. The number of entries
should correspond to the NX x NY x NZ parameters on the DIMENS keyword.
Repeat counts may be used, for example 30*0.03 & Taken from cell
allocated relative permeability table. \\
dimensionless & dimensionless & dimensionless & \\
Notes: & & & \\
\end{longtable}

\emph{Table 8.149: SGL Keyword Description}

End-point scaling allows the entered relative permeability functions to
be re-scaled based on the saturation end-points allocated to each cell
by the SWL, SWCR, SWU, SGL, SGCR, SGU, SOWCR, and SOGCR saturation grid
arrays for the saturation end-points, and the KRG, KRORG, KRORW and KRW
relative permeability grid cell arrays for the relative permeability
end-point data. In addition end-point scaling may be directional
dependent in which case the directional dependent versions of the
aforementioned arrays should be used, that is SGLX, SGLY and SGLZ
instead of SGL. There is also the facility to make the directional
end-point scaling reversible or non-reversible and if the non-reversible
option is selected the non-reversible versions of the aforementioned
arrays should be used, that is SGLX, SGLX-, SGLY, SGLY-, SGLZ and SGLZ-,
instead of the SGL keyword.

\subsubsection{Example}\label{example-357}

\begin{verbatim}
--
--       DEFINE GRID BLOCK END-POINT SGL DATA FOR ALL CELLS
–       (FOR NX x NY x NZ = 300)
--
SGL
         300*0.030                                                              /
\end{verbatim}

The above example defines a constant connate gas saturation of 0.03 to
all 300 cells in the model as defined by the DIMENS keyword in the
RUNSPEC section.

\subsection{SGLPC -- End-Point Scaling Grid Cell Capillary Pressure
Connate Gas
Saturations}\label{sglpc-end-point-scaling-grid-cell-capillary-pressure-connate-gas-saturations}

\begin{longtable}[]{@{}
  >{\raggedright\arraybackslash}p{(\linewidth - 14\tabcolsep) * \real{0.1250}}
  >{\raggedright\arraybackslash}p{(\linewidth - 14\tabcolsep) * \real{0.1250}}
  >{\raggedright\arraybackslash}p{(\linewidth - 14\tabcolsep) * \real{0.1250}}
  >{\raggedright\arraybackslash}p{(\linewidth - 14\tabcolsep) * \real{0.1250}}
  >{\raggedright\arraybackslash}p{(\linewidth - 14\tabcolsep) * \real{0.1250}}
  >{\raggedright\arraybackslash}p{(\linewidth - 14\tabcolsep) * \real{0.1250}}
  >{\raggedright\arraybackslash}p{(\linewidth - 14\tabcolsep) * \real{0.1250}}
  >{\raggedright\arraybackslash}p{(\linewidth - 14\tabcolsep) * \real{0.1250}}@{}}
\toprule\noalign{}
\begin{minipage}[b]{\linewidth}\raggedright
RUNSPEC
\end{minipage} & \begin{minipage}[b]{\linewidth}\raggedright
GRID
\end{minipage} & \begin{minipage}[b]{\linewidth}\raggedright
EDIT
\end{minipage} & \begin{minipage}[b]{\linewidth}\raggedright
PROPS
\end{minipage} & \begin{minipage}[b]{\linewidth}\raggedright
REGIONS
\end{minipage} & \begin{minipage}[b]{\linewidth}\raggedright
SOLUTION
\end{minipage} & \begin{minipage}[b]{\linewidth}\raggedright
SUMMARY
\end{minipage} & \begin{minipage}[b]{\linewidth}\raggedright
SCHEDULE
\end{minipage} \\
\midrule\noalign{}
\endhead
\bottomrule\noalign{}
\endlastfoot
\end{longtable}

\subsubsection{Description}\label{description-639}

SGLPC defines the connate gas saturation for all the cells in the model
via an array when the end-point scaling option has been invoked via the
ENDSCALE keyword in the RUNSPEC section. The connate gas saturation is
defined as the minimum gas saturation in a two-phase gas relative
permeability table. The keyword only applies the scaling to the drainage
capillary pressure tables, unlike the SGL keyword that applies the
scaling to both the capillary pressure and relative permeability tables.
The keyword can be used with all grid types.

\begin{longtable}[]{@{}
  >{\raggedright\arraybackslash}p{(\linewidth - 6\tabcolsep) * \real{0.2500}}
  >{\raggedright\arraybackslash}p{(\linewidth - 6\tabcolsep) * \real{0.2500}}
  >{\raggedright\arraybackslash}p{(\linewidth - 6\tabcolsep) * \real{0.2500}}
  >{\raggedright\arraybackslash}p{(\linewidth - 6\tabcolsep) * \real{0.2500}}@{}}
\toprule\noalign{}
\begin{minipage}[b]{\linewidth}\raggedright
No.
\end{minipage} & \begin{minipage}[b]{\linewidth}\raggedright
Name
\end{minipage} & \begin{minipage}[b]{\linewidth}\raggedright
Description
\end{minipage} & \begin{minipage}[b]{\linewidth}\raggedright
Default
\end{minipage} \\
\midrule\noalign{}
\endhead
\bottomrule\noalign{}
\endlastfoot
Field & Metric & Laboratory & \\
1 & SGLPC & SGLPC is an array of real numbers assigning the connate gas
saturation values to each cell in the model. The number of entries
should correspond to the NX x NY x NZ parameters on the DIMENS keyword.
If SGLPC is omitted from the input deck the values will be defaulted to
those on the SGL series of keywords. If the SGL series of keywords are
missing from the input deck then the values are taken from the cell
allocated capillary pressure table. Repeat counts may be used, for
example 30*0.03 & Taken from SGL or from the cell allocated capillary
pressure table. \\
dimensionless & dimensionless & dimensionless & \\
Notes: & & & \\
\end{longtable}

\emph{Table 8.150: SGLPC Keyword Description}

End-point scaling allows the entered relative permeability functions to
be re-scaled based on the saturation end-points allocated to each cell
by the SWL, SWCR, SWU, SGL, SGCR, SGU, SOWCR, and SOGCR saturation grid
arrays for the saturation end-points, and the KRG, KRGR, KRO, KRORG,
KRORW, KRW and KRWR relative permeability grid cell arrays for the
relative permeability end-point data. In addition end-point scaling may
be directional dependent in which case the directional dependent
versions of the aforementioned arrays should be used, that is SGLX, SGLY
and SGLZ instead of SGL or SGLPC. There is also the facility to make the
directional end-point scaling reversible or non-reversible and if the
non-reversible option is selected the non-reversible versions of the
aforementioned arrays should be used, that is SGLX, SGLX-, SGLY, SGLY-,
SGLZ and SGLZ-, instead of the SGL or SGLPC keywords.

Missing Some Functionality - Use with Caution.

\subsubsection{Example}\label{example-358}

\begin{verbatim}
--
--       DEFINE GRID BLOCK END-POINT SGLPC DATA FOR ALL CELLS
–-       (FOR NX x NY x NZ = 300)
--
SGLPC
         300*0.030                                                              /
\end{verbatim}

The above example defines a constant connate gas saturation of 0.03 to
all 300 cells in the model as defined by the DIMENS keyword in the
RUNSPEC section.

\subsection{SGOF -- Gas-Oil Saturation Tables versus Gas (Format Type
1)}\label{sgof-gas-oil-saturation-tables-versus-gas-format-type-1}

\begin{longtable}[]{@{}
  >{\raggedright\arraybackslash}p{(\linewidth - 14\tabcolsep) * \real{0.1250}}
  >{\raggedright\arraybackslash}p{(\linewidth - 14\tabcolsep) * \real{0.1250}}
  >{\raggedright\arraybackslash}p{(\linewidth - 14\tabcolsep) * \real{0.1250}}
  >{\raggedright\arraybackslash}p{(\linewidth - 14\tabcolsep) * \real{0.1250}}
  >{\raggedright\arraybackslash}p{(\linewidth - 14\tabcolsep) * \real{0.1250}}
  >{\raggedright\arraybackslash}p{(\linewidth - 14\tabcolsep) * \real{0.1250}}
  >{\raggedright\arraybackslash}p{(\linewidth - 14\tabcolsep) * \real{0.1250}}
  >{\raggedright\arraybackslash}p{(\linewidth - 14\tabcolsep) * \real{0.1250}}@{}}
\toprule\noalign{}
\begin{minipage}[b]{\linewidth}\raggedright
RUNSPEC
\end{minipage} & \begin{minipage}[b]{\linewidth}\raggedright
GRID
\end{minipage} & \begin{minipage}[b]{\linewidth}\raggedright
EDIT
\end{minipage} & \begin{minipage}[b]{\linewidth}\raggedright
PROPS
\end{minipage} & \begin{minipage}[b]{\linewidth}\raggedright
REGIONS
\end{minipage} & \begin{minipage}[b]{\linewidth}\raggedright
SOLUTION
\end{minipage} & \begin{minipage}[b]{\linewidth}\raggedright
SUMMARY
\end{minipage} & \begin{minipage}[b]{\linewidth}\raggedright
SCHEDULE
\end{minipage} \\
\midrule\noalign{}
\endhead
\bottomrule\noalign{}
\endlastfoot
\end{longtable}

\subsubsection{Description}\label{description-640}

The SGOF keyword defines the oil and gas relative permeability and
oil-gas capillary versus gas saturation tables for when oil and gas are
present in the input deck. This keyword should only be used if both oil
and gas are present in the run.

\begin{longtable}[]{@{}
  >{\raggedright\arraybackslash}p{(\linewidth - 6\tabcolsep) * \real{0.2500}}
  >{\raggedright\arraybackslash}p{(\linewidth - 6\tabcolsep) * \real{0.2500}}
  >{\raggedright\arraybackslash}p{(\linewidth - 6\tabcolsep) * \real{0.2500}}
  >{\raggedright\arraybackslash}p{(\linewidth - 6\tabcolsep) * \real{0.2500}}@{}}
\toprule\noalign{}
\begin{minipage}[b]{\linewidth}\raggedright
No.
\end{minipage} & \begin{minipage}[b]{\linewidth}\raggedright
Name
\end{minipage} & \begin{minipage}[b]{\linewidth}\raggedright
Description
\end{minipage} & \begin{minipage}[b]{\linewidth}\raggedright
Default
\end{minipage} \\
\midrule\noalign{}
\endhead
\bottomrule\noalign{}
\endlastfoot
Field & Metric & Laboratory & \\
1 & SGAS & A columnar vector of real monotonically increasing down the
column values starting from zero and terminating at one, that defines
the gas saturation. & None \\
dimensionless & dimensionless & dimensionless & \\
2 & KRG & A columnar vector of real values that are either equal or
increasing down the column and that are greater than or equal to zero
and less than or equal to one that defines the gas relative
permeability. & None \\
dimensionless & dimensionless & dimensionless & \\
3 & KRO & A columnar vector of real values that are either equal or
decreasing down the column and that are greater than or equal to zero
and less than or equal to one that defines the oil relative permeability
with respect to gas and connate water saturation. When water is active
in the run, the first entry the column, that is at krog(Sg = 0), must be
the same as the first entry in the corresponding SWOF table, that is at
krow(So = 1 - Swco). The last value in the column should be zero. &
None \\
dimensionless & dimensionless & dimensionless & \\
4 & PCOG & A columnar vector of real values that are either equal or
increasing down the column that defines the oil-gas relative capillary
pressure. & None \\
psia & bars & atm & \\
Notes: & & & \\
\end{longtable}

\emph{Table 8.151: SGOF Keyword Description}

\subsubsection{Example}\label{example-359}

The following example is based on NTSFUN equals two on the TABDIMS
keyword in the RUNSPEC section.

\begin{verbatim}
--
--       GAS-OIL RELATIVE PERMEABILITY TABLES (SGOF)
SGOF
--       SG         KRG       KROG      PCOG
--       FRAC                           PSIA
--       -------   --------  -------   -------
         0.00000   0.000000  0.90000    0.0000
         0.03000   0.000000  0.82500    0.0000
         0.80000   0.900000  0.00000    0.0000             / TABLE No. 01
--       -------   --------  -------   -------
         0.00000   0.000000  0.90000    0.0000
         0.03000   0.000000  0.82500    0.0000
         0.04420   0.024200  0.80000    0.0000
         0.05850   0.048500  0.77500    0.0000
         0.07270   0.072700  0.75000    0.0000
         0.08700   0.097000  0.72500    0.0000
         0.10120   0.121200  0.70000    0.0000
         0.11550   0.145500  0.67500    0.0000
         0.12970   0.169700  0.65000    0.0000
         0.14390   0.193900  0.62500    0.0000
         0.15820   0.218200  0.60000    0.0000
         0.17240   0.242400  0.57500    0.0000
         0.18670   0.266700  0.55000    0.0000
         0.20090   0.290900  0.52500    0.0000
         0.21520   0.315200  0.50000    0.0000
         0.22940   0.339400  0.47500    0.0000
         0.24360   0.363600  0.45000    0.0000
         0.25790   0.387900  0.42500    0.0000
         0.27210   0.412100  0.40000    0.0000
         0.28640   0.436400  0.37500    0.0000
         0.30060   0.460600  0.35000    0.0000
         0.31480   0.484800  0.32500    0.0000
         0.32910   0.509100  0.30000    0.0000
         0.34330   0.533300  0.27500    0.0000
         0.35760   0.557600  0.25000    0.0000
         0.37180   0.581800  0.22500    0.0000
         0.38610   0.606100  0.20000    0.0000
         0.40030   0.630300  0.17500    0.0000
         0.41450   0.654500  0.15000    0.0000
         0.42880   0.678800  0.12500    0.0000
         0.44300   0.703000  0.10000    0.0000
         0.45730   0.727300  0.07500    0.0000
         0.47150   0.751500  0.05000    0.0000
         0.48580   0.775800  0.02500    0.0000
         0.50000   0.800000  0.00000    0.0000
         0.80000   0.900000  0.00000    0.0000             / TABLE No. 02
\end{verbatim}

The example defines two SGOF tables for use when oil, gas and water are
present in the run.

\subsection{SGOFLET -- Gas-Oil LET Relative Permeability
Functions}\label{sgoflet-gas-oil-let-relative-permeability-functions}

\begin{longtable}[]{@{}
  >{\raggedright\arraybackslash}p{(\linewidth - 14\tabcolsep) * \real{0.1250}}
  >{\raggedright\arraybackslash}p{(\linewidth - 14\tabcolsep) * \real{0.1250}}
  >{\raggedright\arraybackslash}p{(\linewidth - 14\tabcolsep) * \real{0.1250}}
  >{\raggedright\arraybackslash}p{(\linewidth - 14\tabcolsep) * \real{0.1250}}
  >{\raggedright\arraybackslash}p{(\linewidth - 14\tabcolsep) * \real{0.1250}}
  >{\raggedright\arraybackslash}p{(\linewidth - 14\tabcolsep) * \real{0.1250}}
  >{\raggedright\arraybackslash}p{(\linewidth - 14\tabcolsep) * \real{0.1250}}
  >{\raggedright\arraybackslash}p{(\linewidth - 14\tabcolsep) * \real{0.1250}}@{}}
\toprule\noalign{}
\begin{minipage}[b]{\linewidth}\raggedright
RUNSPEC
\end{minipage} & \begin{minipage}[b]{\linewidth}\raggedright
GRID
\end{minipage} & \begin{minipage}[b]{\linewidth}\raggedright
EDIT
\end{minipage} & \begin{minipage}[b]{\linewidth}\raggedright
PROPS
\end{minipage} & \begin{minipage}[b]{\linewidth}\raggedright
REGIONS
\end{minipage} & \begin{minipage}[b]{\linewidth}\raggedright
SOLUTION
\end{minipage} & \begin{minipage}[b]{\linewidth}\raggedright
SUMMARY
\end{minipage} & \begin{minipage}[b]{\linewidth}\raggedright
SCHEDULE
\end{minipage} \\
\midrule\noalign{}
\endhead
\bottomrule\noalign{}
\endlastfoot
\end{longtable}

\subsubsection{Description}\label{description-641}

SGOFLET defines the relative permeability and capillary pressure
parameters for the gas-oil LET family of models. Both the gas and oil
phases should be made active in the model via the GAS and OIL keywords
in the RUNSPEC section. See section 8.2.6Saturation Table Generation -
LET Functions and Lomeland et al.~{[}Lomeland F., Ebeltoft E. and Thomas
W.H., 2005. A New Versatile Relative Permeability Correlation. Paper
SCA2005-32 presented at the International Symposium of the Society of
Core Analysts held in Toronto, Canada, 21-25 August, 2005.{]},
{[}Lomeland F. and Ebeltoft E., 2008. A New Versatile Capillary Pressure
Correlation. Paper SCA2008-08 presented at the International Symposium
of the Society of Core Analysts held in Abu Dhabi, UAE, 29 Oct.~-- 2
Nov., 2008.{]} {[}Lomeland F., Hasanov B., Ebeltoft E. and Berge M.,
2012. A Versatile Representation of Up-scaled Relative Permeability for
Field Applications. Paper SPE 154487-MS presented at the EAGE Annual
Conference \& Exhibition incorporating SPE Europec held in Copenhagen,
Denmark, 4-7 June 2012.{]} and {[}Lomeland F., 2018.Overview Of The Let
Family Of Versatile Correlations For Flow Functions. Paper SCA2018-056
presented at the International Symposium of the Society of Core Analysts
held in Trondheim, Norway, 27-30 August 2018.{]} for further information
on the model.

The keyword is used as a replacement for the SGOF keyword for
three-phase oil-gas-water systems, and the LET series of keywords cannot
be combined with the standard set of relative permeability keywords.

\begin{tcolorbox}[enhanced jigsaw, title=\textcolor{quarto-callout-note-color}{\faInfo}\hspace{0.5em}{Note}, coltitle=black, breakable, toprule=.15mm, colbacktitle=quarto-callout-note-color!10!white, leftrule=.75mm, opacitybacktitle=0.6, colback=white, bottomtitle=1mm, toptitle=1mm, titlerule=0mm, left=2mm, colframe=quarto-callout-note-color-frame, arc=.35mm, rightrule=.15mm, bottomrule=.15mm, opacityback=0]

This is an OPM Flow specific keyword and will therefore cause an error
in the commercial simulator.

\end{tcolorbox}

\begin{longtable}[]{@{}
  >{\raggedright\arraybackslash}p{(\linewidth - 6\tabcolsep) * \real{0.2500}}
  >{\raggedright\arraybackslash}p{(\linewidth - 6\tabcolsep) * \real{0.2500}}
  >{\raggedright\arraybackslash}p{(\linewidth - 6\tabcolsep) * \real{0.2500}}
  >{\raggedright\arraybackslash}p{(\linewidth - 6\tabcolsep) * \real{0.2500}}@{}}
\toprule\noalign{}
\begin{minipage}[b]{\linewidth}\raggedright
No.
\end{minipage} & \begin{minipage}[b]{\linewidth}\raggedright
Name
\end{minipage} & \begin{minipage}[b]{\linewidth}\raggedright
Description
\end{minipage} & \begin{minipage}[b]{\linewidth}\raggedright
Default
\end{minipage} \\
\midrule\noalign{}
\endhead
\bottomrule\noalign{}
\endlastfoot
Field & Metric & Laboratory & \\
1 & SGL & SGL is a real positive number less than one that defines the
connate gas saturation, that is smallest gas saturation in the LET
function. & 0,0 \\
dimensionless & dimensionless & dimensionless & \\
2 & SGCR & SGCR is a real positive number greater than or equal to SGL
and less than one, that defines the critical gas saturation, that is the
largest gas saturation for which the gas relative permeability is zero.
& 0,0 \\
dimensionless & dimensionless & dimensionless & \\
3 & LGAS & LGAS is a real positive number that defines the LET Lower
empirical parameter Lg for the gas phase with the associated oil phase
in the LET relative permeability equation. & 1,0 \\
dimensionless & dimensionless & dimensionless & \\
4 & EGAS & EGAS is a real positive number that defines the LET Elevation
empirical parameter Eg for the gas phase with the associated oil phase
in the LET relative permeability equation. & 1,0 \\
dimensionless & dimensionless & dimensionless & \\
5 & TGAS & TGAS is a real positive number that defines the LET Top
empirical parameter Tg for the gas phase with the associated oil phase
in the LET relative permeability equation. & 1,0 \\
dimensionless & dimensionless & dimensionless & \\
6 & KRTGAS & KRTGAS is a real positive number less than one, that
defines the relative permeability of gas at the maximum gas saturation
Krgt in the LET relative permeability equation. & 1,0 \\
dimensionless & dimensionless & dimensionless & \\
7 & SORG & SORG is a real positive number less than one that defines the
residual oil saturation in a gas-oil system in the LET equations. &
0,0 \\
dimensionless & dimensionless & dimensionless & \\
8 & SOGCR & SOGCR is a real positive number less than one that defines
critical oil-in-gas saturation, that is the largest oil saturation for
which the oil relative permeability is zero in a gas-oil system. &
0,0 \\
dimensionless & dimensionless & dimensionless & \\
9 & LOIL & LOIL is a real positive number that defines the LET Lower
empirical parameter Lo for the oil phase with the associated gas phase
in the LET relative permeability equation. & 1,0 \\
dimensionless & dimensionless & dimensionless & \\
10 & EOIL & EOIL is a real positive number that defines the LET
Elevation empirical parameter Eo for the oil phase with the associated
gas phase in the LET relative permeability equation. & 1,0 \\
dimensionless & dimensionless & dimensionless & \\
11 & TOIL & TOIL is a real positive number that defines the LET Top
empirical parameter To for the oil phase with the associated gas phase
in the LET relative permeability equation. & 1,0 \\
dimensionless & dimensionless & dimensionless & \\
12 & KRTOIL & KRTOIL is a real positive number less than or equal to
one, that defines the relative permeability of oil at the residual oil
saturation, Krot in the LET relative permeability equation. & 1,0 \\
dimensionless & dimensionless & dimensionless & \\
13 & LPC & LPC is a real positive number that defines the gas-oil LET
Lower empirical parameter L in the LET capillary pressure equation. &
1,0 \\
dimensionless & dimensionless & dimensionless & \\
14 & EPC & EPC is a real positive number that defines the gas-oil LET
Elevation empirical parameter E in the LET capillary pressure equation.
& 1,0 \\
dimensionless & dimensionless & dimensionless & \\
15 & TPC & TPC is a real positive number that defines the gas-oil LET
Top empirical parameter T in the LET capillary pressure equation. &
1,0 \\
dimensionless & dimensionless & dimensionless & \\
16 & PCIR & PCIR is a real positive number that defines the gas-oil
capillary pressure at connate water saturation (SWL) in the LET
capillary pressure equations. & 0,0 \\
psia & bars & atm & \\
17 & PCIT & PCIT is a real positive number that defines the gas-oil
threshold capillary pressure at the maximum water saturation in the LET
capillary pressure equations. & 0,0 \\
psia & bars & atm & \\
Notes: & & & \\
\end{longtable}

\emph{Table 8.152: SGOFLET Keyword Description}

Note there a two versions of the LET functions, LET232 for two-phase
flowing conditions and LETx {[}Lomeland F. and Ebeltoft E., 2013.
Versatile Three-phase Correlations for Relative Permeability and
Capillary Pressure. Paper SCA2013-034 presented at the International
Symposium of the Society of Core Analysts held in Napa Valley,
California, USA, 16-19 September, 2013.{]} for three-phase flowing
conditions. This keyword implements the LET version for a gas-oil
system.

The functions are dependent on the drainage and imbibition cycle of the
wetting phase as well as drainage and inhibition cycle number, since a
reservoir may undergo several flooding events. To account for this the
system defines the flooding event using the three saturations: Sw, So,
and Sg together with the state of the three saturations during the
flooding event. The saturation state can be Increasing, Decreasing, or
Constant, for a given flooding event cycle number (n). Thus, Sw(D),
So(I), Sg(C)1or DIC1, means the water phase is decreasing, the oil phase
is increasing and the gas phase is constant for the primary or first
cycle (n equals one). This is case for when oil is migrating into the
reservoir rock and displacing the initial water contained with the
reservoir.

\begin{tcolorbox}[enhanced jigsaw, title=\textcolor{quarto-callout-note-color}{\faInfo}\hspace{0.5em}{Note}, coltitle=black, breakable, toprule=.15mm, colbacktitle=quarto-callout-note-color!10!white, leftrule=.75mm, opacitybacktitle=0.6, colback=white, bottomtitle=1mm, toptitle=1mm, titlerule=0mm, left=2mm, colframe=quarto-callout-note-color-frame, arc=.35mm, rightrule=.15mm, bottomrule=.15mm, opacityback=0]

All the LET parameters are dependent on the flooding event and flooding
cycle, and thus are expected vary as such. To be clear, the values of
SGCR, Lg, Lo etc. should be different for each flooding cycle.

\end{tcolorbox}

See also the SGWFLET -- Gas-Water LET Relative Permeability Functions,
and SWOFLET -- Water-Oil LET Relative Permeability Functions keywords in
this section.

\subsubsection{Example}\label{example-360}

The following example uses the SGOFLET keyword to define two relative
gas-oil relative permeability tables, based on NTSFUN equals two on the
TABDIMS keyword in the RUNSPEC section.

\begin{verbatim}
--
--       SGOFLET – GAS-OIL LET REL. PERMEABILITY FUNCTIONS (OPM FLOW KEYWORD)
--
SGOFLET
--       SGL       SGCR     L-GAS     E-GAS     T-GAS     KRT-GAS
--       SOR       SOGCR    L-OIL     E-OIL     T-OIL     KRT-OIL
--       L-PC      E-PC     T-PC      PCIR      PCIT
--       -------   ------   -------   -------   -------   -------
         0.00000   0.0000   1.00000   1.00000   1.00000   1.00000
         0.00000   0.0000   1.00000   1.00000   1.00000   1.00000
         1*        1*       1*        1*        1*                 / TABLE NO. 01

         0.00000   0.0000   1.00000   1.00000   1.00000   1.00000
         0.00000   0.0000   1.00000   1.00000   1.00000   1.00000
         1*        1*       1*        1*        1*                 / TABLE NO. 02
\end{verbatim}

Here the SGOFLET keyword parameters are all set to their default values.

\subsection{SGU -- End-Point Scaling Grid Cell Gas
Saturation}\label{sgu-end-point-scaling-grid-cell-gas-saturation}

\begin{longtable}[]{@{}
  >{\raggedright\arraybackslash}p{(\linewidth - 14\tabcolsep) * \real{0.1250}}
  >{\raggedright\arraybackslash}p{(\linewidth - 14\tabcolsep) * \real{0.1250}}
  >{\raggedright\arraybackslash}p{(\linewidth - 14\tabcolsep) * \real{0.1250}}
  >{\raggedright\arraybackslash}p{(\linewidth - 14\tabcolsep) * \real{0.1250}}
  >{\raggedright\arraybackslash}p{(\linewidth - 14\tabcolsep) * \real{0.1250}}
  >{\raggedright\arraybackslash}p{(\linewidth - 14\tabcolsep) * \real{0.1250}}
  >{\raggedright\arraybackslash}p{(\linewidth - 14\tabcolsep) * \real{0.1250}}
  >{\raggedright\arraybackslash}p{(\linewidth - 14\tabcolsep) * \real{0.1250}}@{}}
\toprule\noalign{}
\begin{minipage}[b]{\linewidth}\raggedright
RUNSPEC
\end{minipage} & \begin{minipage}[b]{\linewidth}\raggedright
GRID
\end{minipage} & \begin{minipage}[b]{\linewidth}\raggedright
EDIT
\end{minipage} & \begin{minipage}[b]{\linewidth}\raggedright
PROPS
\end{minipage} & \begin{minipage}[b]{\linewidth}\raggedright
REGIONS
\end{minipage} & \begin{minipage}[b]{\linewidth}\raggedright
SOLUTION
\end{minipage} & \begin{minipage}[b]{\linewidth}\raggedright
SUMMARY
\end{minipage} & \begin{minipage}[b]{\linewidth}\raggedright
SCHEDULE
\end{minipage} \\
\midrule\noalign{}
\endhead
\bottomrule\noalign{}
\endlastfoot
\end{longtable}

\subsubsection{Description}\label{description-642}

SGU defines the maximum gas saturation for all the cells in the model
via an array when the end-point scaling option has been invoked via the
ENDSCALE keyword in the RUNSPEC section. The maximum gas saturation is
defined as the maximum gas saturation in a two-phase gas relative
permeability table.

The keyword can be used with all grid types.

\begin{longtable}[]{@{}
  >{\raggedright\arraybackslash}p{(\linewidth - 6\tabcolsep) * \real{0.2500}}
  >{\raggedright\arraybackslash}p{(\linewidth - 6\tabcolsep) * \real{0.2500}}
  >{\raggedright\arraybackslash}p{(\linewidth - 6\tabcolsep) * \real{0.2500}}
  >{\raggedright\arraybackslash}p{(\linewidth - 6\tabcolsep) * \real{0.2500}}@{}}
\toprule\noalign{}
\begin{minipage}[b]{\linewidth}\raggedright
No.
\end{minipage} & \begin{minipage}[b]{\linewidth}\raggedright
Name
\end{minipage} & \begin{minipage}[b]{\linewidth}\raggedright
Description
\end{minipage} & \begin{minipage}[b]{\linewidth}\raggedright
Default
\end{minipage} \\
\midrule\noalign{}
\endhead
\bottomrule\noalign{}
\endlastfoot
Field & Metric & Laboratory & \\
1 & SGU & SGU is an array of real numbers assigning the maximum gas
saturation values to each cell in the model. The number of entries
should correspond to the NX x NY x NZ parameters on the DIMENS keyword.
Repeat counts may be used, for example 30*0.70 & Taken from cell
allocated relative permeability table. \\
dimensionless & dimensionless & dimensionless & \\
Notes: & & & \\
\end{longtable}

\emph{Table 8.153: SGU Keyword Description}

End-point scaling allows the entered relative permeability functions to
be re-scaled based on the saturation end-points allocated to each cell
by the SWL, SWCR, SWU, SGL, SGCR, SGU, SOWCR, and SOGCR saturation grid
arrays for the saturation end-points, and the KRG, KRORG, KRORW and KRW
relative permeability grid cell arrays for the relative permeability
end-point data. In addition end-point scaling may be directional
dependent in which case the directional dependent versions of the
aforementioned arrays should be used, that is SGUX, SGUY and SGUZ
instead of SGU. There is also the facility to make the directional
end-point scaling reversible or non-reversible and if the non-reversible
option is selected the non-reversible versions of the aforementioned
arrays should be used, that is SGUX, SGUX-, SGUY, SGUY-, SGUZ and SGUZ-,
instead of the SGU keyword.

\subsubsection{Example}\label{example-361}

\begin{verbatim}
--
--       DEFINE GRID BLOCK END-POINT SGU DATA FOR ALL CELLS
--       (FOR NX x NY x NZ = 300)
--
SGU
         300*0.700                                                         /
\end{verbatim}

The above example defines a constant connate gas saturation of 0.70 to
all 300 cells in the model as defined by the DIMENS keyword in the
RUNSPEC section.

\subsection{SGWFLET -- Gas-Water LET Relative Permeability
Functions}\label{sgwflet-gas-water-let-relative-permeability-functions}

\begin{longtable}[]{@{}
  >{\raggedright\arraybackslash}p{(\linewidth - 14\tabcolsep) * \real{0.1250}}
  >{\raggedright\arraybackslash}p{(\linewidth - 14\tabcolsep) * \real{0.1250}}
  >{\raggedright\arraybackslash}p{(\linewidth - 14\tabcolsep) * \real{0.1250}}
  >{\raggedright\arraybackslash}p{(\linewidth - 14\tabcolsep) * \real{0.1250}}
  >{\raggedright\arraybackslash}p{(\linewidth - 14\tabcolsep) * \real{0.1250}}
  >{\raggedright\arraybackslash}p{(\linewidth - 14\tabcolsep) * \real{0.1250}}
  >{\raggedright\arraybackslash}p{(\linewidth - 14\tabcolsep) * \real{0.1250}}
  >{\raggedright\arraybackslash}p{(\linewidth - 14\tabcolsep) * \real{0.1250}}@{}}
\toprule\noalign{}
\begin{minipage}[b]{\linewidth}\raggedright
RUNSPEC
\end{minipage} & \begin{minipage}[b]{\linewidth}\raggedright
GRID
\end{minipage} & \begin{minipage}[b]{\linewidth}\raggedright
EDIT
\end{minipage} & \begin{minipage}[b]{\linewidth}\raggedright
PROPS
\end{minipage} & \begin{minipage}[b]{\linewidth}\raggedright
REGIONS
\end{minipage} & \begin{minipage}[b]{\linewidth}\raggedright
SOLUTION
\end{minipage} & \begin{minipage}[b]{\linewidth}\raggedright
SUMMARY
\end{minipage} & \begin{minipage}[b]{\linewidth}\raggedright
SCHEDULE
\end{minipage} \\
\midrule\noalign{}
\endhead
\bottomrule\noalign{}
\endlastfoot
\end{longtable}

\subsubsection{Description}\label{description-643}

SWGFLET defines the relative permeability and capillary pressure
parameters for the water-gas LET family of models. Both the gas and
water phases should be made active in the model via the GAS and WATER
keywords in the RUNSPEC section. This keyword should only be used in
two-phase models containing the gas and water phases only, that is the
oil phase must be absent. See section 8.2.6Saturation Table Generation -
LET Functions and Lomeland et al.~{[}Lomeland F., Ebeltoft E. and Thomas
W.H., 2005. A New Versatile Relative Permeability Correlation. Paper
SCA2005-32 presented at the International Symposium of the Society of
Core Analysts held in Toronto, Canada, 21-25 August, 2005.{]},
{[}Lomeland F. and Ebeltoft E., 2008. A New Versatile Capillary Pressure
Correlation. Paper SCA2008-08 presented at the International Symposium
of the Society of Core Analysts held in Abu Dhabi, UAE, 29 Oct.~-- 2
Nov., 2008.{]} {[}Lomeland F., Hasanov B., Ebeltoft E. and Berge M.,
2012. A Versatile Representation of Up-scaled Relative Permeability for
Field Applications. Paper SPE 154487-MS presented at the EAGE Annual
Conference \& Exhibition incorporating SPE Europec held in Copenhagen,
Denmark, 4-7 June 2012.{]} and {[}Lomeland F., 2018.Overview Of The Let
Family Of Versatile Correlations For Flow Functions. Paper SCA2018-056
presented at the International Symposium of the Society of Core Analysts
held in Trondheim, Norway, 27-30 August 2018.{]} for further information
on the model.

The keyword is used as a replacement for the SGWFN keyword for two-phase
gas-water systems, and the LET series of keywords cannot be combined
with the standard set of relative permeability keywords.

\begin{tcolorbox}[enhanced jigsaw, title=\textcolor{quarto-callout-note-color}{\faInfo}\hspace{0.5em}{Note}, coltitle=black, breakable, toprule=.15mm, colbacktitle=quarto-callout-note-color!10!white, leftrule=.75mm, opacitybacktitle=0.6, colback=white, bottomtitle=1mm, toptitle=1mm, titlerule=0mm, left=2mm, colframe=quarto-callout-note-color-frame, arc=.35mm, rightrule=.15mm, bottomrule=.15mm, opacityback=0]

This is an OPM Flow specific keyword and will therefore cause an error
in the commercial simulator.

\end{tcolorbox}

\begin{longtable}[]{@{}
  >{\raggedright\arraybackslash}p{(\linewidth - 6\tabcolsep) * \real{0.2500}}
  >{\raggedright\arraybackslash}p{(\linewidth - 6\tabcolsep) * \real{0.2500}}
  >{\raggedright\arraybackslash}p{(\linewidth - 6\tabcolsep) * \real{0.2500}}
  >{\raggedright\arraybackslash}p{(\linewidth - 6\tabcolsep) * \real{0.2500}}@{}}
\toprule\noalign{}
\begin{minipage}[b]{\linewidth}\raggedright
No.
\end{minipage} & \begin{minipage}[b]{\linewidth}\raggedright
Name
\end{minipage} & \begin{minipage}[b]{\linewidth}\raggedright
Description
\end{minipage} & \begin{minipage}[b]{\linewidth}\raggedright
Default
\end{minipage} \\
\midrule\noalign{}
\endhead
\bottomrule\noalign{}
\endlastfoot
Field & Metric & Laboratory & \\
1 & SGL & SGL is a real positive number less than one that defines the
connate gas saturation, that is smallest gas saturation in the LET
function. & 0,0 \\
dimensionless & dimensionless & dimensionless & \\
2 & SGCR & SGCR is a real positive number greater than or equal to SGL
and less than one, that defines the critical gas saturation, that is the
largest gas saturation for which the gas relative permeability is zero.
& 0,0 \\
dimensionless & dimensionless & dimensionless & \\
3 & LGAS & LGAS is a real positive number that defines the LET Lower
empirical parameter Lg for the gas phase with the associated water phase
in the LET relative permeability equations. & 1,0 \\
dimensionless & dimensionless & dimensionless & \\
4 & EGAS & EGAS is a real positive number that defines the LET Elevation
empirical parameter Eg for the gas phase with the associated water phase
in the LET relative permeability equations. & 1,0 \\
dimensionless & dimensionless & dimensionless & \\
5 & TGAS & TGAS is a real positive number that defines the LET Top
empirical parameter Tg for the gas phase with the associated water phase
in the LET relative permeability equations. & 1,0 \\
dimensionless & dimensionless & dimensionless & \\
6 & KRTGAS & KRTGAS is a real positive number less than one, that
defines the relative permeability of gas at the maximum gas saturation
Krgt in the LET relative permeability equations. & 1,0 \\
dimensionless & dimensionless & dimensionless & \\
7 & SWL & SWL is a real positive number less than one, that defines the
connate water saturation, that is the smallest water saturation in the
LET function. & 0,0 \\
dimensionless & dimensionless & dimensionless & \\
8 & SWCR & SWCR is a real positive number greater than or equal to SWL
and less than one, that defines the critical water saturation, that is
the largest water saturation for which the water relative permeability
is zero, Swirr in equations. & 0,0 \\
dimensionless & dimensionless & dimensionless & \\
9 & LWAT & LWAT is a real positive number that defines the LET Lower
empirical parameter Lw for the water phase with the associated gas phase
in the LET relative permeability equations. & 1,0 \\
dimensionless & dimensionless & dimensionless & \\
10 & EWAT & EWAT is a real positive number that defines the LET
Elevation empirical parameter Ew for the water phase with the associated
gas phase in the LET relative permeability equations. & 1,0 \\
dimensionless & dimensionless & dimensionless & \\
11 & TWAT & TWAT is a real positive number that defines the LET Top
empirical parameter Tw for the water phase with the associated gas phase
in the LET relative permeability equations. & 1,0 \\
dimensionless & dimensionless & dimensionless & \\
12 & KRTWAT & KRTWAT is a real positive number less than one, that
defines the relative permeability of water at the maximum water
saturation (normally the maximum water saturation is one) Krwt in the
LET relative permeability equations. & 1,0 \\
dimensionless & dimensionless & dimensionless & \\
13 & LPC & LPC is a real positive number that defines the gas-water LET
Lower empirical parameter L in the LET capillary pressure equation. &
1,0 \\
dimensionless & dimensionless & dimensionless & \\
14 & EPC & EPC is a real positive number that defines the gas-water LET
Elevation empirical parameter E in the LET capillary pressure equation.
& 1,0 \\
dimensionless & dimensionless & dimensionless & \\
15 & TPC & TPC is a real positive number that defines the gas-water LET
empirical parameter T in the LET capillary pressure equation. & 1,0 \\
dimensionless & dimensionless & dimensionless & \\
16 & PCIR & PCIR is a real positive number that defines the gas-water
capillary pressure at connate water saturation (SWL) in the LET
capillary pressure equations. & 0,0 \\
psia & bars & atm & \\
17 & PCIT & PCIT is a real positive number that defines the gas-water
threshold capillary pressure at the maximum water saturation in the LET
capillary pressure equations. & 0,0 \\
psia & bars & atm & \\
Notes: & & & \\
\end{longtable}

\emph{Table 8.154: SGWFLET Keyword Description}

Note there a two versions of the LET functions, LET237 for two-phase
flowing conditions and LETx {[}Lomeland F. and Ebeltoft E., 2013.
Versatile Three-phase Correlations for Relative Permeability and
Capillary Pressure. Paper SCA2013-034 presented at the International
Symposium of the Society of Core Analysts held in Napa Valley,
California, USA, 16-19 September, 2013.{]} for three-phase flowing
conditions. This keyword implements the LET version for a gas-water
system.

The functions are dependent on the drainage and imbibition cycle of the
wetting phase as well as drainage and inhibition cycle number, since a
reservoir may undergo several flooding events. To account for this the
system defines the flooding event using the three saturations: Sw, So,
and Sg together with the state of the three saturations during the
flooding event. The saturation state can be Increasing, Decreasing, or
Constant, for a given flooding event cycle number (n). Thus, Sw(D),
So(C), Sg(I), or DCI1, means the water phase is decreasing, the oil
phase is constant and the gas phase is increasing for the primary or
first cycle (n equals one). This is case for when gas is migrating into
the reservoir rock and displacing the initial water contained with the
reservoir.

\begin{tcolorbox}[enhanced jigsaw, title=\textcolor{quarto-callout-note-color}{\faInfo}\hspace{0.5em}{Note}, coltitle=black, breakable, toprule=.15mm, colbacktitle=quarto-callout-note-color!10!white, leftrule=.75mm, opacitybacktitle=0.6, colback=white, bottomtitle=1mm, toptitle=1mm, titlerule=0mm, left=2mm, colframe=quarto-callout-note-color-frame, arc=.35mm, rightrule=.15mm, bottomrule=.15mm, opacityback=0]

All the LET parameters are dependent on the flooding event and flooding
cycle, and thus are expected vary as such. To be clear, the values of
SWCR, Lg, Lw etc. should be different for each flooding cycle.

\end{tcolorbox}

See also the SGOFLET-- Gas-Oil LET Relative Permeability Functions, and
SWOFLET -- Water-Oil LET Relative Permeability Functions keywords in
this section, that may be used for three-phase systems.

\subsubsection{Example}\label{example-362}

The following example uses the SWGFLET keyword to define two relative
gas-water relative permeability tables, based on NTSFUN equals two on
the TABDIMS keyword in the RUNSPEC section.

\begin{verbatim}

--
--       SGWFLET – GAS-WATER LET REL. PERMEABILITY FUNCTIONS (OPM FLOW KEYWORD)
--
SGWFLET
--       SGL       SGCR     L-GAS     E-GAS     T-GAS     KRT-GAS
--       SWL       SWCR     L-WAT     E-WAT     T-WAT     KRT-WAT
--       L-PC      E-PC     T-PC      PCIR      PCIT
--       -------   ------   -------   -------   -------   -------
         0.00000   0.0000   1.00000   1.00000   1.00000   1.00000
         0.00000   0.0000   1.00000   1.00000   1.00000   1.00000
         1*        1*       1*        1*        1*                 / TABLE NO. 01

         0.00000   0.0000   1.00000   1.00000   1.00000   1.00000
         0.00000   0.0000   1.00000   1.00000   1.00000   1.00000
         1*        1*       1*        1*        1*                 / TABLE NO. 02
\end{verbatim}

Here the SGWFLET keyword parameters are all set to their default values.

\subsection{SGWFN -- Gas-Water Saturation Tables (Format Type
2)}\label{sgwfn-gas-water-saturation-tables-format-type-2}

\begin{longtable}[]{@{}
  >{\raggedright\arraybackslash}p{(\linewidth - 14\tabcolsep) * \real{0.1250}}
  >{\raggedright\arraybackslash}p{(\linewidth - 14\tabcolsep) * \real{0.1250}}
  >{\raggedright\arraybackslash}p{(\linewidth - 14\tabcolsep) * \real{0.1250}}
  >{\raggedright\arraybackslash}p{(\linewidth - 14\tabcolsep) * \real{0.1250}}
  >{\raggedright\arraybackslash}p{(\linewidth - 14\tabcolsep) * \real{0.1250}}
  >{\raggedright\arraybackslash}p{(\linewidth - 14\tabcolsep) * \real{0.1250}}
  >{\raggedright\arraybackslash}p{(\linewidth - 14\tabcolsep) * \real{0.1250}}
  >{\raggedright\arraybackslash}p{(\linewidth - 14\tabcolsep) * \real{0.1250}}@{}}
\toprule\noalign{}
\begin{minipage}[b]{\linewidth}\raggedright
RUNSPEC
\end{minipage} & \begin{minipage}[b]{\linewidth}\raggedright
GRID
\end{minipage} & \begin{minipage}[b]{\linewidth}\raggedright
EDIT
\end{minipage} & \begin{minipage}[b]{\linewidth}\raggedright
PROPS
\end{minipage} & \begin{minipage}[b]{\linewidth}\raggedright
REGIONS
\end{minipage} & \begin{minipage}[b]{\linewidth}\raggedright
SOLUTION
\end{minipage} & \begin{minipage}[b]{\linewidth}\raggedright
SUMMARY
\end{minipage} & \begin{minipage}[b]{\linewidth}\raggedright
SCHEDULE
\end{minipage} \\
\midrule\noalign{}
\endhead
\bottomrule\noalign{}
\endlastfoot
\end{longtable}

\subsubsection{Description}\label{description-644}

The SGWFN keyword defines the gas and water relative permeability and
gas-water capillary pressure data versus gas saturation tables for when
gas and water are present in the input deck. This keyword should only be
used if gas and water are present in the run.

\begin{longtable}[]{@{}
  >{\raggedright\arraybackslash}p{(\linewidth - 6\tabcolsep) * \real{0.2500}}
  >{\raggedright\arraybackslash}p{(\linewidth - 6\tabcolsep) * \real{0.2500}}
  >{\raggedright\arraybackslash}p{(\linewidth - 6\tabcolsep) * \real{0.2500}}
  >{\raggedright\arraybackslash}p{(\linewidth - 6\tabcolsep) * \real{0.2500}}@{}}
\toprule\noalign{}
\begin{minipage}[b]{\linewidth}\raggedright
No.
\end{minipage} & \begin{minipage}[b]{\linewidth}\raggedright
Name
\end{minipage} & \begin{minipage}[b]{\linewidth}\raggedright
Description
\end{minipage} & \begin{minipage}[b]{\linewidth}\raggedright
Default
\end{minipage} \\
\midrule\noalign{}
\endhead
\bottomrule\noalign{}
\endlastfoot
Field & Metric & Laboratory & \\
1 & SGAS & A columnar vector of real monotonically increasing down the
column values starting from zero and terminating at one, that defines
the gas saturation. & None \\
dimensionless & dimensionless & dimensionless & \\
2 & KRG & A columnar vector of real values that are either equal or
increasing down the column and that are greater than or equal to zero
and less than or equal to one that defines the gas relative
permeability. Note that the first entry in the column must be zero. &
None \\
dimensionless & dimensionless & dimensionless & \\
3 & KRW & A columnar vector of real values that are either equal or
decreasing down the column and that are greater than or equal to zero
and less than or equal to one that defines the water relative
permeability with respect to gas saturation. The last value in the
column should be zero. & None \\
dimensionless & dimensionless & dimensionless & \\
4 & PCGW & A columnar vector of real values that are either equal or
increasing down the column that defines the gas-water relative capillary
pressure. & None \\
psia & bars & atm & \\
Notes: & & & \\
\end{longtable}

\emph{Table 8.155: SGWFN Keyword Description}

\subsubsection{Example}\label{example-363}

\begin{verbatim}
--
--       GAS-WATER RELATIVE PERMEABILITY TABLES (SGWFN)
SGWFN
--       SG         KRG       KRW       PCOW
--       FRAC                           PSIA
--       --------   --------  -------   -------
         0.000000   0.0000    0.9000    0.000000
         0.200000   0.0002    0.7664    0.000000
         0.699099   0.4973    0.0000    0.000000
         0.700000   1.0000    0.0000    0.000000           / TABLE NO. 01
--       --------   --------  -------   -------
         0.000000   0.0000    0.9000    0.000000
         0.200000   0.0002    0.7664    0.000000
         0.245309   0.0004    0.7443    0.000000
         0.261989   0.0010    0.6907    0.000000
         0.303091   0.0044    0.5671    0.000000
         0.368269   0.0191    0.3962    0.000000
         0.435026   0.0519    0.2528    0.000000
         0.486387   0.0940    0.1643    0.000000
         0.522283   0.1339    0.1137    0.000000
         0.550683   0.1725    0.0803    0.000000
         0.575342   0.2115    0.0559    0.000000
         0.599076   0.2542    0.0367    0.000000
         0.621294   0.2991    0.0223    0.000000
         0.642171   0.3458    0.0120    0.000000
         0.658984   0.3868    0.0061    0.000000
         0.671123   0.4183    0.0030    0.000000
         0.679268   0.4403    0.0015    0.000000
         0.684963   0.4562    0.0008    0.000000
         0.688893   0.4674    0.0004    0.000000
         0.692025   0.4765    0.0002    0.000000
         0.694641   0.4841    0.0001    0.000000
         0.696976   0.4910    0.0000    0.000000
         0.699099   0.4973    0.0000    0.000000
         0.700000   1.0000    0.0000    0.000000           / TABLE NO. 02
\end{verbatim}

The example defines two SGWFN tables for use when gas and water are
present in the run.

\subsection{SHRATE -- Activate Log-based Polymer Shearing and Define the
Shear Rate
Constant}\label{shrate-activate-log-based-polymer-shearing-and-define-the-shear-rate-constant}

\begin{longtable}[]{@{}
  >{\raggedright\arraybackslash}p{(\linewidth - 14\tabcolsep) * \real{0.1250}}
  >{\raggedright\arraybackslash}p{(\linewidth - 14\tabcolsep) * \real{0.1250}}
  >{\raggedright\arraybackslash}p{(\linewidth - 14\tabcolsep) * \real{0.1250}}
  >{\raggedright\arraybackslash}p{(\linewidth - 14\tabcolsep) * \real{0.1250}}
  >{\raggedright\arraybackslash}p{(\linewidth - 14\tabcolsep) * \real{0.1250}}
  >{\raggedright\arraybackslash}p{(\linewidth - 14\tabcolsep) * \real{0.1250}}
  >{\raggedright\arraybackslash}p{(\linewidth - 14\tabcolsep) * \real{0.1250}}
  >{\raggedright\arraybackslash}p{(\linewidth - 14\tabcolsep) * \real{0.1250}}@{}}
\toprule\noalign{}
\begin{minipage}[b]{\linewidth}\raggedright
RUNSPEC
\end{minipage} & \begin{minipage}[b]{\linewidth}\raggedright
GRID
\end{minipage} & \begin{minipage}[b]{\linewidth}\raggedright
EDIT
\end{minipage} & \begin{minipage}[b]{\linewidth}\raggedright
PROPS
\end{minipage} & \begin{minipage}[b]{\linewidth}\raggedright
REGIONS
\end{minipage} & \begin{minipage}[b]{\linewidth}\raggedright
SOLUTION
\end{minipage} & \begin{minipage}[b]{\linewidth}\raggedright
SUMMARY
\end{minipage} & \begin{minipage}[b]{\linewidth}\raggedright
SCHEDULE
\end{minipage} \\
\midrule\noalign{}
\endhead
\bottomrule\noalign{}
\endlastfoot
\end{longtable}

\subsubsection{Description}\label{description-645}

This keyword activates the logarithm-based polymer shear
thinning/thickening option and defines the shear rate constant. This
keyword can only be used in conjunction with the PLYSHLOG in the PROPS
section

\begin{longtable}[]{@{}
  >{\raggedright\arraybackslash}p{(\linewidth - 6\tabcolsep) * \real{0.2500}}
  >{\raggedright\arraybackslash}p{(\linewidth - 6\tabcolsep) * \real{0.2500}}
  >{\raggedright\arraybackslash}p{(\linewidth - 6\tabcolsep) * \real{0.2500}}
  >{\raggedright\arraybackslash}p{(\linewidth - 6\tabcolsep) * \real{0.2500}}@{}}
\toprule\noalign{}
\begin{minipage}[b]{\linewidth}\raggedright
No.
\end{minipage} & \begin{minipage}[b]{\linewidth}\raggedright
Name
\end{minipage} & \begin{minipage}[b]{\linewidth}\raggedright
Description
\end{minipage} & \begin{minipage}[b]{\linewidth}\raggedright
Default
\end{minipage} \\
\midrule\noalign{}
\endhead
\bottomrule\noalign{}
\endlastfoot
Field & Metric & Laboratory & \\
1 & SHRATE & A positive real value that defines the shear rate constant.
& 4.8 \\
dimensionless & dimensionless & dimensionless & \\
Notes: & & & \\
\end{longtable}

\emph{Table 8.3.283.1: SHRATE Keyword Description}

\subsubsection{Example}\label{example-364}

The following example activates the logarithm-based polymer shear
thinning/thickening option and defines the shear rate constants for a
run with two PVT regions.

\begin{verbatim}
--
--       ACTIVATE LOG-BASED POLYMER SHEAR THINNING-THICKENING OPTION
--       AND DEFINE THE SHEAR RATE CONSTANT
--
SHRATE
--       SHEAR RATE
--       CONSTANT
         4.8      /
         4.8      /
\end{verbatim}

\subsection{SKIP -- Activate Skipping of All Keywords and Input
Data}\label{skip-activate-skipping-of-all-keywords-and-input-data-4}

The SKIP keyword activates skipping of all keywords and input data until
the ENDSKIP keyword is encountered. All keywords between the SKIP and
ENDSKIP keywords are ignored.

See SKIP -- Activate Skipping of All Keywords and Input Data in the
GLOBAL section for a full description.

\subsection{SKIP100 -- Activate Skipping of Black-Oil Keywords and Input
Data}\label{skip100-activate-skipping-of-black-oil-keywords-and-input-data-3}

The SKIP100 keyword activates skipping of all keywords and input data by
the commercial black-oil simulator until the ENDSKIP keyword is
encountered. All keywords between the SKIP100 and ENDSKIP keywords are
ignored by the commercial black-oil simulator. The SKIP100 keyword is
ignored by the commercial compositional simulator.

See SKIP100 -- Activate Skipping of Keywords by Black-Oil Simulator in
the GLOBAL section for a full description.

\subsection{SKIP300 -- Activate Skipping of Keywords by Compositional
Simulator}\label{skip300-activate-skipping-of-keywords-by-compositional-simulator-4}

The SKIP300 keyword activates skipping of all keywords and input data by
the commercial compositional simulator until the ENDSKIP keyword is
encountered. All keywords between the SKIP300 and ENDSKIP keywords are
ignored by the commercial compositional simulator. The SKIP300 keyword
is ignored by the commercial black-oil simulator.

See Error: Reference source not found in the GLOBAL section for a full
description.

\subsection{SKPRPOLY -- Polymer Molecular Weight Model Polymer Injection
Skin
Table}\label{skprpoly-polymer-molecular-weight-model-polymer-injection-skin-table}

\begin{longtable}[]{@{}
  >{\raggedright\arraybackslash}p{(\linewidth - 14\tabcolsep) * \real{0.1250}}
  >{\raggedright\arraybackslash}p{(\linewidth - 14\tabcolsep) * \real{0.1250}}
  >{\raggedright\arraybackslash}p{(\linewidth - 14\tabcolsep) * \real{0.1250}}
  >{\raggedright\arraybackslash}p{(\linewidth - 14\tabcolsep) * \real{0.1250}}
  >{\raggedright\arraybackslash}p{(\linewidth - 14\tabcolsep) * \real{0.1250}}
  >{\raggedright\arraybackslash}p{(\linewidth - 14\tabcolsep) * \real{0.1250}}
  >{\raggedright\arraybackslash}p{(\linewidth - 14\tabcolsep) * \real{0.1250}}
  >{\raggedright\arraybackslash}p{(\linewidth - 14\tabcolsep) * \real{0.1250}}@{}}
\toprule\noalign{}
\begin{minipage}[b]{\linewidth}\raggedright
RUNSPEC
\end{minipage} & \begin{minipage}[b]{\linewidth}\raggedright
GRID
\end{minipage} & \begin{minipage}[b]{\linewidth}\raggedright
EDIT
\end{minipage} & \begin{minipage}[b]{\linewidth}\raggedright
PROPS
\end{minipage} & \begin{minipage}[b]{\linewidth}\raggedright
REGIONS
\end{minipage} & \begin{minipage}[b]{\linewidth}\raggedright
SOLUTION
\end{minipage} & \begin{minipage}[b]{\linewidth}\raggedright
SUMMARY
\end{minipage} & \begin{minipage}[b]{\linewidth}\raggedright
SCHEDULE
\end{minipage} \\
\midrule\noalign{}
\endhead
\bottomrule\noalign{}
\endlastfoot
\end{longtable}

\subsubsection{Description}\label{description-646}

This keyword, SKPRPOLY, describes the relationship of a water injection
well's injected polymer skin pressure as a function of polymer
throughput and water velocity, for the simulator's Polymer Molecular
Weight Transport option. The table is a two dimensional table that
relates the polymer throughput values and water velocity values to
derive the resulting wellbore skin pressure of the injected polymer,
which is then used to calculate the total wellbore skin pressure based
on the polymer concentration.

This keyword should only be used if the POLYMER and POLYMW keywords in
the RUNSPEC section are also activated.

\begin{tcolorbox}[enhanced jigsaw, title=\textcolor{quarto-callout-note-color}{\faInfo}\hspace{0.5em}{Note}, coltitle=black, breakable, toprule=.15mm, colbacktitle=quarto-callout-note-color!10!white, leftrule=.75mm, opacitybacktitle=0.6, colback=white, bottomtitle=1mm, toptitle=1mm, titlerule=0mm, left=2mm, colframe=quarto-callout-note-color-frame, arc=.35mm, rightrule=.15mm, bottomrule=.15mm, opacityback=0]

This is an OPM Flow specific keyword that employs an alternative polymer
flood model based on a Polymer Molecular Weight Transport equation, that
is not available in the commercial simulator. The model has been tested
using metric units; however, using either field or laboratory units with
the option should be considered experimental.

\end{tcolorbox}

\begin{longtable}[]{@{}
  >{\raggedright\arraybackslash}p{(\linewidth - 6\tabcolsep) * \real{0.2500}}
  >{\raggedright\arraybackslash}p{(\linewidth - 6\tabcolsep) * \real{0.2500}}
  >{\raggedright\arraybackslash}p{(\linewidth - 6\tabcolsep) * \real{0.2500}}
  >{\raggedright\arraybackslash}p{(\linewidth - 6\tabcolsep) * \real{0.2500}}@{}}
\toprule\noalign{}
\begin{minipage}[b]{\linewidth}\raggedright
No.
\end{minipage} & \begin{minipage}[b]{\linewidth}\raggedright
Name
\end{minipage} & \begin{minipage}[b]{\linewidth}\raggedright
Description
\end{minipage} & \begin{minipage}[b]{\linewidth}\raggedright
Default
\end{minipage} \\
\midrule\noalign{}
\endhead
\bottomrule\noalign{}
\endlastfoot
Field & Metric & Laboratory & \\
1-1 & SKPRPNUM & A positive integer value greater than zero and less
than or equal to the NTSKPOLY variable, as defined on the PINTDIMS
keyword in the RUNSPEC section, that defines the SKPRPOLY Polymer
Molecular Weight Model polymer injection skin pressure table number. &
None \\
1-2 & POLCON & A real positive value that the defines the reference
polymer concentration for the table. & None \\
lb/stb & kg/sm3 & gm/scc & \\
2-1 & THRUPUT & A real positive monotonically increasing vector, that
defines the polymer throughput values. The first entry should be zero to
define a no throughput data set, and each vector record should be on a
separate line terminated by a ``/''. & None \\
feet3/feet2 & m3/m2 & cm3/cm2 & \\
3-1 & VELOCITY & A real positive monotonically increasing vector, that
defines the polymer velocity values. The first entry should be zero to
define a no velocity data set, and each vector record should be on a
separate line terminated by a ``/''. & None \\
feet/day & m/day & cm/hour & \\
4-1 & PRESS & A series of real positive vectors representing the
wellbore skin pressure, for the given reference polymer concentration
(POLCON), for all combinations of the polymer throughput values
(THRUPUT) and velocity values (VELOCITY), organized as a series of
vectors PRESS(THRUPUT, VELOCITY). Thus, the first vector represents the
wellbore skin pressure of the first THRUPUT value and each entry in the
vector is the corresponding wellbore skin pressure of the associated
VELOCITY vector. Each vector should be on a separate line and should be
terminated by a ``/''. Thus, if THRUPUT has three entries and VELOCITY
has four, then there should be three vectors, with each vector
containing four elements representing wellbore skin pressure values, as
a function of THRUPUT and VELOCITY. & None \\
psia & barsa & atma & \\
Notes: & & & \\
\end{longtable}

\emph{Table 8.156: SKPRPOLY Keyword Description}

Unlike other PROPS section table keywords, that enable multiple tables
following the keyword to be entered, the SKPRPOLY keyword requires that
the keyword itself must be repeated for each table, as is shown in the
example on the following page.

The WSKPTAB keyword in the SCHEDULE section may be used to assign the
SKPRPOLY and SKPRWAT tables to water injections wells, that enable the
calculation of the wellbore skin pressure based on the fluids being
injected.

See also the SKPRWAT, PLYMWINJ, and PLYVMH keywords, in the PROPS
section, that are the additional property keywords required for the
Polymer Molecular Weight Transport option. In addition, see also the
WPMITAB keyword in the SCHEDULE section, that assigns the PLYMWINJ
tables to the water injection wells.

Note that the standard polymer property data keywords: PLYROCK, PLYADS,
PLYMAX, etc., are still required to fully describe the polymer fluid.

\subsubsection{Example}\label{example-365}

Given NTSKPOLY equals two on the PINTDIMS keyword in the RUNSPEC
section, then two SKPRPOLY tables are required to be entered:

\begin{verbatim}
--
--       POLYMER MOLECULAR WEIGHT MODEL POLYMER INJECTION SKIN TABLE
--       (OPM FLOW PROPS KEYWORD)
--
SKPRPOLY
--       TABLE   POLYMER REF
--       NO.     CONCENTRATION
         1       2.00                  /
--
--       THROUGHPUT VALUES
--
         0.0     200.0   400.0         /
--
--       VELOCITY VALUES
--
         0.0     50.0    80.0    100.0 /
--
--       PRESS SKIN VALUES
--
         0.0     10.0    20.0    40.0  / PRESS(THRUPUT=1, VELOCITY=1 TO N)
         0.0     50.0    100.0   200.0 / PRESS(THRUPUT=2, VELOCITY=1 TO N)
         0.0     100.0   200.0   400.0 / PRESS(THRUPUT=3, VELOCITY=1 TO N)
/
SKPRPOLY
--       TABLE   POLYMER REF
--       NO.     CONCENTRATION
         2       2.0                   /
--
--       THROUGHPUT VALUES
--
         0.0     200.0   400.0         /
--
--       VELOCITY VALUES
--
         0.0     30.0    50.0    100.0 /
--
--       PRESS SKIN VALUES
--
         0.0     10.0    20.0    40.0  / PRESS(THRUPUT=1, VELOCITY=1 TO N)
         0.0     50.0    100.0   200.0 / PRESS(THRUPUT=2, VELOCITY=1 TO N)
         0.0     100.0   200.0   400.0 / PRESS(THRUPUT=3, VELOCITY=1 TO N)
/
\end{verbatim}

As mentioned previously, the SKPRPOLY keyword requires that the keyword
itself must be repeated for each table, as is shown in the above
example.

\subsection{SKPRWAT -- Polymer Molecular Weight Model Water Injection
Skin
Table}\label{skprwat-polymer-molecular-weight-model-water-injection-skin-table}

\begin{longtable}[]{@{}
  >{\raggedright\arraybackslash}p{(\linewidth - 14\tabcolsep) * \real{0.1250}}
  >{\raggedright\arraybackslash}p{(\linewidth - 14\tabcolsep) * \real{0.1250}}
  >{\raggedright\arraybackslash}p{(\linewidth - 14\tabcolsep) * \real{0.1250}}
  >{\raggedright\arraybackslash}p{(\linewidth - 14\tabcolsep) * \real{0.1250}}
  >{\raggedright\arraybackslash}p{(\linewidth - 14\tabcolsep) * \real{0.1250}}
  >{\raggedright\arraybackslash}p{(\linewidth - 14\tabcolsep) * \real{0.1250}}
  >{\raggedright\arraybackslash}p{(\linewidth - 14\tabcolsep) * \real{0.1250}}
  >{\raggedright\arraybackslash}p{(\linewidth - 14\tabcolsep) * \real{0.1250}}@{}}
\toprule\noalign{}
\begin{minipage}[b]{\linewidth}\raggedright
RUNSPEC
\end{minipage} & \begin{minipage}[b]{\linewidth}\raggedright
GRID
\end{minipage} & \begin{minipage}[b]{\linewidth}\raggedright
EDIT
\end{minipage} & \begin{minipage}[b]{\linewidth}\raggedright
PROPS
\end{minipage} & \begin{minipage}[b]{\linewidth}\raggedright
REGIONS
\end{minipage} & \begin{minipage}[b]{\linewidth}\raggedright
SOLUTION
\end{minipage} & \begin{minipage}[b]{\linewidth}\raggedright
SUMMARY
\end{minipage} & \begin{minipage}[b]{\linewidth}\raggedright
SCHEDULE
\end{minipage} \\
\midrule\noalign{}
\endhead
\bottomrule\noalign{}
\endlastfoot
\end{longtable}

\subsubsection{Description}\label{description-647}

This keyword, SKPRWAT, describes the relationship of a water injection
well's injected water skin pressure as a function of water throughput
and water velocity, for the simulator's Polymer Molecular Weight
Transport option. The table is a two dimensional table that relates the
water throughput values and water velocity values to derive the
resulting wellbore skin pressure of the injected water, which is then
used to calculate the total wellbore skin pressure based on the polymer
concentration.

This keyword should only be used if the POLYMER and POLYMW keywords in
the RUNSPEC section are also activated.

\begin{tcolorbox}[enhanced jigsaw, title=\textcolor{quarto-callout-note-color}{\faInfo}\hspace{0.5em}{Note}, coltitle=black, breakable, toprule=.15mm, colbacktitle=quarto-callout-note-color!10!white, leftrule=.75mm, opacitybacktitle=0.6, colback=white, bottomtitle=1mm, toptitle=1mm, titlerule=0mm, left=2mm, colframe=quarto-callout-note-color-frame, arc=.35mm, rightrule=.15mm, bottomrule=.15mm, opacityback=0]

This is an OPM Flow specific keyword that employs an alternative polymer
flood model based on a Polymer Molecular Weight Transport equation, that
is not available in the commercial simulator. The model has been tested
using metric units; however, using either field or laboratory units with
the option should be considered experimental.

\end{tcolorbox}

\begin{longtable}[]{@{}
  >{\raggedright\arraybackslash}p{(\linewidth - 6\tabcolsep) * \real{0.2500}}
  >{\raggedright\arraybackslash}p{(\linewidth - 6\tabcolsep) * \real{0.2500}}
  >{\raggedright\arraybackslash}p{(\linewidth - 6\tabcolsep) * \real{0.2500}}
  >{\raggedright\arraybackslash}p{(\linewidth - 6\tabcolsep) * \real{0.2500}}@{}}
\toprule\noalign{}
\begin{minipage}[b]{\linewidth}\raggedright
No.
\end{minipage} & \begin{minipage}[b]{\linewidth}\raggedright
Name
\end{minipage} & \begin{minipage}[b]{\linewidth}\raggedright
Description
\end{minipage} & \begin{minipage}[b]{\linewidth}\raggedright
Default
\end{minipage} \\
\midrule\noalign{}
\endhead
\bottomrule\noalign{}
\endlastfoot
Field & Metric & Laboratory & \\
1-1 & SKPRWNUM & A positive integer value greater than zero and less
than or equal to the NTSKWAT variable, as defined on the PINTDIMS
keyword in the RUNSPEC section, that defines the SKPRWAT Polymer
Molecular Weight Model water injection skin pressure table number. &
None \\
2-1 & THRUPUT & A real positive monotonically increasing vector, that
defines the water throughput values. The first entry should be zero to
define a no throughput data set, and each vector record should be on a
separate line terminated by a ``/''. & None \\
feet3/feet2 & m3/m2 & cm3/cm2 & \\
3-1 & VELOCITY & A real positive monotonically increasing vector, that
defines the water velocity values. The first entry should be zero to
define a no velocity data set, and each vector record should be on a
separate line terminated by a ``/''. & None \\
feet/day & m/day & cm/hour & \\
4-1 & PRESS & A series of real positive vectors representing the
wellbore skin pressure, for all combinations of the water throughput
values (THRUPUT) and velocity values (VELOCITY), organized as a series
of vectors PRESS(THRUPUT, VELOCITY). Thus, the first vector represents
the wellbore skin pressure of the first THRUPUT value and each entry in
the vector is the corresponding wellbore skin pressure of the associated
VELOCITY vector. Each vector should be on a separate line and should be
terminated by a ``/''. Thus, if THRUPUT has three entries and VELOCITY
has four, then there should be three vectors, with each vector
containing four elements representing wellbore skin pressure values, as
a function of THRUPUT and VELOCITY. & None \\
psia & barsa & atma & \\
Notes: & & & \\
\end{longtable}

\emph{Table 8.157: SKPRWAT Keyword Description}

Unlike other PROPS section table keywords, that enable multiple tables
following the keyword to be entered, the SKPRWAT keyword requires that
the keyword itself must be repeated for each table, as is shown in the
example on the following page.

The WSKPTAB keyword in the SCHEDULE section may be used to assign the
SKPRPOLY and SKPRWAT tables to water injections wells, that enable the
calculation of the wellbore skin pressure based on the fluids being
injected.

See also the SKPRPOLY, PLYMWINJ, and PLYVMH keywords, in the PROPS
section, that are the additional property keywords required for the
Polymer Molecular Weight Transport option. In addition, see also the
WPMITAB keyword in the SCHEDULE section, that assigns the PLYMWINJ
tables to the water injection wells.

Note that the standard polymer property data keywords: PLYROCK, PLYADS,
PLYMAX, etc., are still required to fully describe the polymer fluid.

\subsubsection{Example}\label{example-366}

Given NTSKWAT equals two on the PINTDIMS keyword in the RUNSPEC section,
then two SKPRWAT tables are required to be entered:

\begin{verbatim}
--
--       POLYMER MOLECULAR WEIGHT MODEL WATER INJECTION SKIN TABLE
--       (OPM FLOW PROPS KEYWORD)
--
SKPRWAT
         1                             / TABLE NUMBER
--
--       THROUGHPUT VALUES
--
         0.0     200.0   400.0         /
--
--       VELOCITY VALUES
--
         0.0     50.0    70.0    100.0 /
--
--       PRESS SKIN VALUES
--
         0.0     2.0     4.0     8.0   / PRESS(THRUPUT=1, VELOCITY=1 TO N)
         0.0     20.0    40.0    80.0  / PRESS(THRUPUT=2, VELOCITY=1 TO N)
         0.0     50.0    100.0   200.0 / PRESS(THRUPUT=3, VELOCITY=1 TO N)
/
SKPRWAT
         2                             / TABLE NUMBER
--
--       THROUGHPUT VALUES
--
         0.0     200.0   400.0         /
--
--       VELOCITY VALUES
--
         0.0     30.0    50.0    100.0 /
--
--       PRESS SKIN VALUES
--
         0.0     2.0     4.0     8.0   / PRESS(THRUPUT=1, VELOCITY=1 TO N)
         0.0     20.0    40.0    80.0  / PRESS(THRUPUT=2, VELOCITY=1 TO N)
         0.0     50.0    100.0   200.0 / PRESS(THRUPUT=3, VELOCITY=1 TO N)
/
\end{verbatim}

As mentioned previously, the SKPRWAT keyword requires that the keyword
itself must be repeated for each table, as is shown in the above
example.

\subsection{SKRO -- End-Point Scaling of Grid Cell Kro(Swl)
(Surfactant)}\label{skro-end-point-scaling-of-grid-cell-kroswl-surfactant}

\begin{longtable}[]{@{}
  >{\raggedright\arraybackslash}p{(\linewidth - 14\tabcolsep) * \real{0.1250}}
  >{\raggedright\arraybackslash}p{(\linewidth - 14\tabcolsep) * \real{0.1250}}
  >{\raggedright\arraybackslash}p{(\linewidth - 14\tabcolsep) * \real{0.1250}}
  >{\raggedright\arraybackslash}p{(\linewidth - 14\tabcolsep) * \real{0.1250}}
  >{\raggedright\arraybackslash}p{(\linewidth - 14\tabcolsep) * \real{0.1250}}
  >{\raggedright\arraybackslash}p{(\linewidth - 14\tabcolsep) * \real{0.1250}}
  >{\raggedright\arraybackslash}p{(\linewidth - 14\tabcolsep) * \real{0.1250}}
  >{\raggedright\arraybackslash}p{(\linewidth - 14\tabcolsep) * \real{0.1250}}@{}}
\toprule\noalign{}
\begin{minipage}[b]{\linewidth}\raggedright
RUNSPEC
\end{minipage} & \begin{minipage}[b]{\linewidth}\raggedright
GRID
\end{minipage} & \begin{minipage}[b]{\linewidth}\raggedright
EDIT
\end{minipage} & \begin{minipage}[b]{\linewidth}\raggedright
PROPS
\end{minipage} & \begin{minipage}[b]{\linewidth}\raggedright
REGIONS
\end{minipage} & \begin{minipage}[b]{\linewidth}\raggedright
SOLUTION
\end{minipage} & \begin{minipage}[b]{\linewidth}\raggedright
SUMMARY
\end{minipage} & \begin{minipage}[b]{\linewidth}\raggedright
SCHEDULE
\end{minipage} \\
\midrule\noalign{}
\endhead
\bottomrule\noalign{}
\endlastfoot
\end{longtable}

\subsubsection{Description}\label{description-648}

SKRO defines the scaling parameter for the surfactant oil relative
permeability value at the connate water saturation (SWL), for all the
cells in the model via an array. The ENDSCALE keyword in the RUNSPEC
section should be activated to enable end-point scaling and the use of
this keyword. In addition, the Surfactant option must be enabled by
either the SURFST or SURFSTES keywords in the RUNSPEC section.

This keyword is not supported by OPM Flow but would change the results
if supported so the simulation will be stopped.

\begin{longtable}[]{@{}
  >{\raggedright\arraybackslash}p{(\linewidth - 6\tabcolsep) * \real{0.2500}}
  >{\raggedright\arraybackslash}p{(\linewidth - 6\tabcolsep) * \real{0.2500}}
  >{\raggedright\arraybackslash}p{(\linewidth - 6\tabcolsep) * \real{0.2500}}
  >{\raggedright\arraybackslash}p{(\linewidth - 6\tabcolsep) * \real{0.2500}}@{}}
\toprule\noalign{}
\begin{minipage}[b]{\linewidth}\raggedright
No.
\end{minipage} & \begin{minipage}[b]{\linewidth}\raggedright
Name
\end{minipage} & \begin{minipage}[b]{\linewidth}\raggedright
Description
\end{minipage} & \begin{minipage}[b]{\linewidth}\raggedright
Default
\end{minipage} \\
\midrule\noalign{}
\endhead
\bottomrule\noalign{}
\endlastfoot
Field & Metric & Laboratory & \\
1 & SKRO & SKRO is an array of positive real numbers which are greater
than zero and less than or equal to 1.0, that are the assigned scaling
SKRO values for each cell in the model. Repeat counts may be used, for
example 50*0.500. & Taken from cell allocated relative permeability
table. \\
dimensionless & dimensionless & dimensionless & \\
Notes: & & & \\
\end{longtable}

\emph{Table 8.158: SKRO Keyword Description}

End-point scaling allows the entered surfactant relative permeability
functions to be scale on the relative permeability values using the
SKRO, SKRORG, SKRORW, SKRW and SKRWR surfactant relative permeability
grid cell arrays for the relative permeability end-point data.

\subsubsection{Example}\label{example-367}

The example defines an input box for the whole grid and for layers one
to three, for layer one SKRO is set equal to 0.850, for layer two SKRO
equals 0.875, and for layer three SKRO equals 0.900.

\begin{verbatim}
--
--       DEFINE INPUT BOX FOR EDITING INPUT ARRAYS (NX=100, NY=100)
--
--       ---------- BOX ---------
--       I1  I2   J1  J2   K1  K2
BOX
         1*  1*   1*  1*   1   3                           / DEFINE BOX AREA
--
--       SET SKRO VALUES FOR THREE LAYERS IN THE MODEL
--
SKRO
         1000*0.855  1000*0.875  1000.0.900                /
--
--       DEFINE END OF INPUT BOX EDITING OF INPUT ARRAYS
--
ENDBOX
\end{verbatim}

\subsection{SKRORG -- End-Point Scaling of Grid Cell Kro(Sgcr)
(Surfactant)}\label{skrorg-end-point-scaling-of-grid-cell-krosgcr-surfactant}

\begin{longtable}[]{@{}
  >{\raggedright\arraybackslash}p{(\linewidth - 14\tabcolsep) * \real{0.1250}}
  >{\raggedright\arraybackslash}p{(\linewidth - 14\tabcolsep) * \real{0.1250}}
  >{\raggedright\arraybackslash}p{(\linewidth - 14\tabcolsep) * \real{0.1250}}
  >{\raggedright\arraybackslash}p{(\linewidth - 14\tabcolsep) * \real{0.1250}}
  >{\raggedright\arraybackslash}p{(\linewidth - 14\tabcolsep) * \real{0.1250}}
  >{\raggedright\arraybackslash}p{(\linewidth - 14\tabcolsep) * \real{0.1250}}
  >{\raggedright\arraybackslash}p{(\linewidth - 14\tabcolsep) * \real{0.1250}}
  >{\raggedright\arraybackslash}p{(\linewidth - 14\tabcolsep) * \real{0.1250}}@{}}
\toprule\noalign{}
\begin{minipage}[b]{\linewidth}\raggedright
RUNSPEC
\end{minipage} & \begin{minipage}[b]{\linewidth}\raggedright
GRID
\end{minipage} & \begin{minipage}[b]{\linewidth}\raggedright
EDIT
\end{minipage} & \begin{minipage}[b]{\linewidth}\raggedright
PROPS
\end{minipage} & \begin{minipage}[b]{\linewidth}\raggedright
REGIONS
\end{minipage} & \begin{minipage}[b]{\linewidth}\raggedright
SOLUTION
\end{minipage} & \begin{minipage}[b]{\linewidth}\raggedright
SUMMARY
\end{minipage} & \begin{minipage}[b]{\linewidth}\raggedright
SCHEDULE
\end{minipage} \\
\midrule\noalign{}
\endhead
\bottomrule\noalign{}
\endlastfoot
\end{longtable}

\subsubsection{Description}\label{description-649}

SKRORG defines the scaling parameter for the surfactant relative
permeability of oil at the critical gas saturation (SGCR), for all the
cells in the model via an array. The ENDSCALE keyword in the RUNSPEC
section should be activated to enable end-point scaling and the use of
this keyword. In addition, the Surfactant option must be enabled by
either the SURFST or SURFSTES keywords in the RUNSPEC section.

This keyword is not supported by OPM Flow but would change the results
if supported so the simulation will be stopped.

\begin{longtable}[]{@{}
  >{\raggedright\arraybackslash}p{(\linewidth - 6\tabcolsep) * \real{0.2500}}
  >{\raggedright\arraybackslash}p{(\linewidth - 6\tabcolsep) * \real{0.2500}}
  >{\raggedright\arraybackslash}p{(\linewidth - 6\tabcolsep) * \real{0.2500}}
  >{\raggedright\arraybackslash}p{(\linewidth - 6\tabcolsep) * \real{0.2500}}@{}}
\toprule\noalign{}
\begin{minipage}[b]{\linewidth}\raggedright
No.
\end{minipage} & \begin{minipage}[b]{\linewidth}\raggedright
Name
\end{minipage} & \begin{minipage}[b]{\linewidth}\raggedright
Description
\end{minipage} & \begin{minipage}[b]{\linewidth}\raggedright
Default
\end{minipage} \\
\midrule\noalign{}
\endhead
\bottomrule\noalign{}
\endlastfoot
Field & Metric & Laboratory & \\
1 & SKRORG & SKRORG is an array of positive real numbers which are
greater than zero and less than or equal to 1.0, that are the assigned
scaling SKRORG values for each cell in the model. Repeat counts may be
used, for example 50*0.850. & Taken from cell allocated relative
permeability table. \\
dimensionless & dimensionless & dimensionless & \\
Notes: & & & \\
\end{longtable}

\emph{Table 8.159: SKRORG Keyword Description}

End-point scaling allows the entered surfactant relative permeability
functions to be scale on the relative permeability values using the
SKRO, SKRORG, SKRORW, SKRW and SKRWR surfactant relative permeability
grid cell arrays for the relative permeability end-point data.

\subsubsection{Example}\label{example-368}

The example uses the EQUALS keyword to set layer one SKRORG equal to
0.750, layer two SKRORG equals 0.775, and layer three SKRORG equals
0.800.

\begin{verbatim}
--
--       ARRAY       CONSTANT     ---------- BOX ---------
--                                I1  I2   J1  J2   K1  K2
EQUALS
         SKRORG      0.7550       1*  1*   1*  1*    1   1  / SKRORG FOR LAYER 1
         SKRORG      0.7750       1*  1*   1*  1*    2   2  / SKRORG FOR LAYER 2
         SKRORG      0.8000       1*  1*   1*  1*    3   3  / SKRORG FOR LAYER 3
/
\end{verbatim}

\subsection{SKRORW -- End-Point Scaling of Grid Cell Kro(Swcr)
(Surfactant)}\label{skrorw-end-point-scaling-of-grid-cell-kroswcr-surfactant}

\begin{longtable}[]{@{}
  >{\raggedright\arraybackslash}p{(\linewidth - 14\tabcolsep) * \real{0.1250}}
  >{\raggedright\arraybackslash}p{(\linewidth - 14\tabcolsep) * \real{0.1250}}
  >{\raggedright\arraybackslash}p{(\linewidth - 14\tabcolsep) * \real{0.1250}}
  >{\raggedright\arraybackslash}p{(\linewidth - 14\tabcolsep) * \real{0.1250}}
  >{\raggedright\arraybackslash}p{(\linewidth - 14\tabcolsep) * \real{0.1250}}
  >{\raggedright\arraybackslash}p{(\linewidth - 14\tabcolsep) * \real{0.1250}}
  >{\raggedright\arraybackslash}p{(\linewidth - 14\tabcolsep) * \real{0.1250}}
  >{\raggedright\arraybackslash}p{(\linewidth - 14\tabcolsep) * \real{0.1250}}@{}}
\toprule\noalign{}
\begin{minipage}[b]{\linewidth}\raggedright
RUNSPEC
\end{minipage} & \begin{minipage}[b]{\linewidth}\raggedright
GRID
\end{minipage} & \begin{minipage}[b]{\linewidth}\raggedright
EDIT
\end{minipage} & \begin{minipage}[b]{\linewidth}\raggedright
PROPS
\end{minipage} & \begin{minipage}[b]{\linewidth}\raggedright
REGIONS
\end{minipage} & \begin{minipage}[b]{\linewidth}\raggedright
SOLUTION
\end{minipage} & \begin{minipage}[b]{\linewidth}\raggedright
SUMMARY
\end{minipage} & \begin{minipage}[b]{\linewidth}\raggedright
SCHEDULE
\end{minipage} \\
\midrule\noalign{}
\endhead
\bottomrule\noalign{}
\endlastfoot
\end{longtable}

\subsubsection{Description}\label{description-650}

SKRORW defines the scaling parameter for the surfactant relative
permeability of oil at the critical water saturation (SWCR), for all the
cells in the model via an array. The ENDSCALE keyword in the RUNSPEC
section should be activated to enable end-point scaling and the use of
this keyword. In addition, the Surfactant option must be enabled by
either the SURFST or SURFSTES keywords in the RUNSPEC section.

This keyword is not supported by OPM Flow but would change the results
if supported so the simulation will be stopped.

\begin{longtable}[]{@{}
  >{\raggedright\arraybackslash}p{(\linewidth - 6\tabcolsep) * \real{0.2500}}
  >{\raggedright\arraybackslash}p{(\linewidth - 6\tabcolsep) * \real{0.2500}}
  >{\raggedright\arraybackslash}p{(\linewidth - 6\tabcolsep) * \real{0.2500}}
  >{\raggedright\arraybackslash}p{(\linewidth - 6\tabcolsep) * \real{0.2500}}@{}}
\toprule\noalign{}
\begin{minipage}[b]{\linewidth}\raggedright
No.
\end{minipage} & \begin{minipage}[b]{\linewidth}\raggedright
Name
\end{minipage} & \begin{minipage}[b]{\linewidth}\raggedright
Description
\end{minipage} & \begin{minipage}[b]{\linewidth}\raggedright
Default
\end{minipage} \\
\midrule\noalign{}
\endhead
\bottomrule\noalign{}
\endlastfoot
Field & Metric & Laboratory & \\
1 & SKRORW & SKRORW is an array of positive real numbers which are
greater than zero and less than or equal to 1.0, that are the assigned
scaling SKRORW values for each cell in the model. Repeat counts may be
used, for example 50*0.850 & Taken from cell allocated relative
permeability table. \\
dimensionless & dimensionless & dimensionless & \\
Notes: & & & \\
\end{longtable}

\emph{Table 8.160: SKRORW Keyword Description}

End-point scaling allows the entered surfactant relative permeability
functions to be scale on the relative permeability values using the
SKRO, SKRORG, SKRORW, SKRW and SKRWR surfactant relative permeability
grid cell arrays for the relative permeability end-point data.

\subsubsection{Example}\label{example-369}

The example defines an input box for the whole grid and for layers one
to three, for layer one SKRORW is set equal to 0.750, for layer two
SKRORW equals 0.775, and for layer three SKRORW equals 0.800.

\begin{verbatim}
--
--       DEFINE INPUT BOX FOR EDITING INPUT ARRAYS (NX=100, NY=100)
--
--       ---------- BOX ---------
--       I1  I2   J1  J2   K1  K2
BOX
         1*  1*   1*  1*   1   3                           / DEFINE BOX AREA
--
--       SET SKRORW VALUES FOR THREE LAYERS IN THE MODEL
--
SKRORW
         1000*0.755  1000*0.775  1000.0.800                /
--
--       DEFINE END OF INPUT BOX EDITING OF INPUT ARRAYS
--
ENDBOX
\end{verbatim}

\subsection{SKRW -- End-Point Scaling of Grid Cell Krw(Sw =1.0)
(Surfactant)}\label{skrw-end-point-scaling-of-grid-cell-krwsw-1.0-surfactant}

\begin{longtable}[]{@{}
  >{\raggedright\arraybackslash}p{(\linewidth - 14\tabcolsep) * \real{0.1250}}
  >{\raggedright\arraybackslash}p{(\linewidth - 14\tabcolsep) * \real{0.1250}}
  >{\raggedright\arraybackslash}p{(\linewidth - 14\tabcolsep) * \real{0.1250}}
  >{\raggedright\arraybackslash}p{(\linewidth - 14\tabcolsep) * \real{0.1250}}
  >{\raggedright\arraybackslash}p{(\linewidth - 14\tabcolsep) * \real{0.1250}}
  >{\raggedright\arraybackslash}p{(\linewidth - 14\tabcolsep) * \real{0.1250}}
  >{\raggedright\arraybackslash}p{(\linewidth - 14\tabcolsep) * \real{0.1250}}
  >{\raggedright\arraybackslash}p{(\linewidth - 14\tabcolsep) * \real{0.1250}}@{}}
\toprule\noalign{}
\begin{minipage}[b]{\linewidth}\raggedright
RUNSPEC
\end{minipage} & \begin{minipage}[b]{\linewidth}\raggedright
GRID
\end{minipage} & \begin{minipage}[b]{\linewidth}\raggedright
EDIT
\end{minipage} & \begin{minipage}[b]{\linewidth}\raggedright
PROPS
\end{minipage} & \begin{minipage}[b]{\linewidth}\raggedright
REGIONS
\end{minipage} & \begin{minipage}[b]{\linewidth}\raggedright
SOLUTION
\end{minipage} & \begin{minipage}[b]{\linewidth}\raggedright
SUMMARY
\end{minipage} & \begin{minipage}[b]{\linewidth}\raggedright
SCHEDULE
\end{minipage} \\
\midrule\noalign{}
\endhead
\bottomrule\noalign{}
\endlastfoot
\end{longtable}

\subsubsection{Description}\label{description-651}

SKRW defines the scaling parameter at the maximum surfactant water
relative permeability value (SWU), that is for Sw = 1.0, for all the
cells in the model via an array. The ENDSCALE keyword in the RUNSPEC
section should be activated to enable end-point scaling and the use of
this keyword. In addition, the Surfactant option must be enabled by
either the SURFST or SURFSTES keywords in the RUNSPEC section.

This keyword is not supported by OPM Flow but would change the results
if supported so the simulation will be stopped.

\begin{longtable}[]{@{}
  >{\raggedright\arraybackslash}p{(\linewidth - 6\tabcolsep) * \real{0.2500}}
  >{\raggedright\arraybackslash}p{(\linewidth - 6\tabcolsep) * \real{0.2500}}
  >{\raggedright\arraybackslash}p{(\linewidth - 6\tabcolsep) * \real{0.2500}}
  >{\raggedright\arraybackslash}p{(\linewidth - 6\tabcolsep) * \real{0.2500}}@{}}
\toprule\noalign{}
\begin{minipage}[b]{\linewidth}\raggedright
No.
\end{minipage} & \begin{minipage}[b]{\linewidth}\raggedright
Name
\end{minipage} & \begin{minipage}[b]{\linewidth}\raggedright
Description
\end{minipage} & \begin{minipage}[b]{\linewidth}\raggedright
Default
\end{minipage} \\
\midrule\noalign{}
\endhead
\bottomrule\noalign{}
\endlastfoot
Field & Metric & Laboratory & \\
1 & SKRW & SKRW is an array of positive real numbers which are greater
than zero and less than or equal to 1.0, that are the assigned scaling
SKRW values for each cell in the model. Repeat counts may be used, for
example 50*1.000. & Taken from cell allocated relative permeability
table. \\
dimensionless & dimensionless & dimensionless & \\
Notes: & & & \\
\end{longtable}

\emph{Table 8.161: SKRW Keyword Description}

\begin{verbatim}
\end{verbatim}

End-point scaling allows the entered surfactant relative permeability
functions to be scale on the relative permeability values using the
SKRO, SKRORG, SKRORW, SKRW and SKRWR surfactant relative permeability
grid cell arrays for the relative permeability end-point data.

\subsubsection{Example}\label{example-370}

The example uses the EQUALS keyword to set SKRW for layer one equal to
0.850, layer two SKRW to 0.875, and layer three KRW to 0.900.

\begin{verbatim}
--
--       ARRAY       CONSTANT     ---------- BOX ---------
--                                I1  I2   J1  J2   K1  K2
EQUALS
         SKRW        0.8550       1*  1*   1*  1*    1   1  / SKRW FOR LAYER 1
         SKRW        0.8750       1*  1*   1*  1*    2   2  / SKRW FOR LAYER 2
         SKRW        0.9000       1*  1*   1*  1*    3   3  / SKRW FOR LAYER 3
/
\end{verbatim}

\subsection{SKRWR -- End-Point Scaling of Grid Cell KRWR(Sowcr)
(Surfactant)}\label{skrwr-end-point-scaling-of-grid-cell-krwrsowcr-surfactant}

\begin{longtable}[]{@{}
  >{\raggedright\arraybackslash}p{(\linewidth - 14\tabcolsep) * \real{0.1250}}
  >{\raggedright\arraybackslash}p{(\linewidth - 14\tabcolsep) * \real{0.1250}}
  >{\raggedright\arraybackslash}p{(\linewidth - 14\tabcolsep) * \real{0.1250}}
  >{\raggedright\arraybackslash}p{(\linewidth - 14\tabcolsep) * \real{0.1250}}
  >{\raggedright\arraybackslash}p{(\linewidth - 14\tabcolsep) * \real{0.1250}}
  >{\raggedright\arraybackslash}p{(\linewidth - 14\tabcolsep) * \real{0.1250}}
  >{\raggedright\arraybackslash}p{(\linewidth - 14\tabcolsep) * \real{0.1250}}
  >{\raggedright\arraybackslash}p{(\linewidth - 14\tabcolsep) * \real{0.1250}}@{}}
\toprule\noalign{}
\begin{minipage}[b]{\linewidth}\raggedright
RUNSPEC
\end{minipage} & \begin{minipage}[b]{\linewidth}\raggedright
GRID
\end{minipage} & \begin{minipage}[b]{\linewidth}\raggedright
EDIT
\end{minipage} & \begin{minipage}[b]{\linewidth}\raggedright
PROPS
\end{minipage} & \begin{minipage}[b]{\linewidth}\raggedright
REGIONS
\end{minipage} & \begin{minipage}[b]{\linewidth}\raggedright
SOLUTION
\end{minipage} & \begin{minipage}[b]{\linewidth}\raggedright
SUMMARY
\end{minipage} & \begin{minipage}[b]{\linewidth}\raggedright
SCHEDULE
\end{minipage} \\
\midrule\noalign{}
\endhead
\bottomrule\noalign{}
\endlastfoot
\end{longtable}

\subsubsection{Description}\label{description-652}

SKRWR defines the scaling parameter at the critical oil to water
saturation value (SOWCR), for the surfactant water relative permeability
curve, for all the cells in the model via an array. The ENDSCALE keyword
in the RUNSPEC section should be activated to enable end-point scaling
and the use of this keyword.In addition, the Surfactant option must be
enabled by either the SURFST or SURFSTES keywords in the RUNSPEC
section.

This keyword is not supported by OPM Flow but would change the results
if supported so the simulation will be stopped.

\begin{longtable}[]{@{}
  >{\raggedright\arraybackslash}p{(\linewidth - 6\tabcolsep) * \real{0.2500}}
  >{\raggedright\arraybackslash}p{(\linewidth - 6\tabcolsep) * \real{0.2500}}
  >{\raggedright\arraybackslash}p{(\linewidth - 6\tabcolsep) * \real{0.2500}}
  >{\raggedright\arraybackslash}p{(\linewidth - 6\tabcolsep) * \real{0.2500}}@{}}
\toprule\noalign{}
\begin{minipage}[b]{\linewidth}\raggedright
No.
\end{minipage} & \begin{minipage}[b]{\linewidth}\raggedright
Name
\end{minipage} & \begin{minipage}[b]{\linewidth}\raggedright
Description
\end{minipage} & \begin{minipage}[b]{\linewidth}\raggedright
Default
\end{minipage} \\
\midrule\noalign{}
\endhead
\bottomrule\noalign{}
\endlastfoot
Field & Metric & Laboratory & \\
1 & SKRWR & SKRWR is an array of positive real numbers which are greater
than zero and less than or equal to 1.0, that are the assigned scaling
SKRWR values for each cell in the model. Repeat counts may be used, for
example 50*1.000. & Taken from cell allocated relative permeability
table. \\
dimensionless & dimensionless & dimensionless & \\
Notes: & & & \\
\end{longtable}

\emph{Table 8.162: SKRWR Keyword Description}

\begin{verbatim}
\end{verbatim}

End-point scaling allows the entered surfactant relative permeability
functions to be scale on the relative permeability values using the
SKRO, SKRORG, SKRORW, SKRW and SKRWR surfactant relative permeability
grid cell arrays for the relative permeability end-point data.

\subsubsection{Example}\label{example-371}

The first example defines an input box for the whole grid and for layers
one to three, for layer one SKRWR is set equal to 0.750, for layer two
SKRWR equals 0.775, and for layer three SKRWR equals 0.800.

\begin{verbatim}
--
--       DEFINE INPUT BOX FOR EDITING INPUT ARRAYS (NX=100, NY=100)
--
--       ---------- BOX ---------
--       I1  I2   J1  J2   K1  K2
BOX
         1*  1*   1*  1*   1   3                           / DEFINE BOX AREA
--
--       SET SKRWR VALUES FOR THREE LAYERS IN THE MODEL
--
SKRWR
         1000*0.755  1000*0.775  1000.0.800                /
--
--       DEFINE END OF INPUT BOX EDITING OF INPUT ARRAYS
--
ENDBOX
\end{verbatim}

\subsection{SLGOF -- Gas-Oil Saturation Tables versus Gas (Format Type
1)}\label{slgof-gas-oil-saturation-tables-versus-gas-format-type-1}

\begin{longtable}[]{@{}
  >{\raggedright\arraybackslash}p{(\linewidth - 14\tabcolsep) * \real{0.1250}}
  >{\raggedright\arraybackslash}p{(\linewidth - 14\tabcolsep) * \real{0.1250}}
  >{\raggedright\arraybackslash}p{(\linewidth - 14\tabcolsep) * \real{0.1250}}
  >{\raggedright\arraybackslash}p{(\linewidth - 14\tabcolsep) * \real{0.1250}}
  >{\raggedright\arraybackslash}p{(\linewidth - 14\tabcolsep) * \real{0.1250}}
  >{\raggedright\arraybackslash}p{(\linewidth - 14\tabcolsep) * \real{0.1250}}
  >{\raggedright\arraybackslash}p{(\linewidth - 14\tabcolsep) * \real{0.1250}}
  >{\raggedright\arraybackslash}p{(\linewidth - 14\tabcolsep) * \real{0.1250}}@{}}
\toprule\noalign{}
\begin{minipage}[b]{\linewidth}\raggedright
RUNSPEC
\end{minipage} & \begin{minipage}[b]{\linewidth}\raggedright
GRID
\end{minipage} & \begin{minipage}[b]{\linewidth}\raggedright
EDIT
\end{minipage} & \begin{minipage}[b]{\linewidth}\raggedright
PROPS
\end{minipage} & \begin{minipage}[b]{\linewidth}\raggedright
REGIONS
\end{minipage} & \begin{minipage}[b]{\linewidth}\raggedright
SOLUTION
\end{minipage} & \begin{minipage}[b]{\linewidth}\raggedright
SUMMARY
\end{minipage} & \begin{minipage}[b]{\linewidth}\raggedright
SCHEDULE
\end{minipage} \\
\midrule\noalign{}
\endhead
\bottomrule\noalign{}
\endlastfoot
\end{longtable}

\subsubsection{Description}\label{description-653}

The SLGOF keyword defines the oil and gas relative permeability and
oil-gas capillary pressure versus liquid saturation tables for when oil
and gas are present in the input deck. This keyword should only be used
if both oil and gas are present in the run.

\begin{longtable}[]{@{}
  >{\raggedright\arraybackslash}p{(\linewidth - 6\tabcolsep) * \real{0.2500}}
  >{\raggedright\arraybackslash}p{(\linewidth - 6\tabcolsep) * \real{0.2500}}
  >{\raggedright\arraybackslash}p{(\linewidth - 6\tabcolsep) * \real{0.2500}}
  >{\raggedright\arraybackslash}p{(\linewidth - 6\tabcolsep) * \real{0.2500}}@{}}
\toprule\noalign{}
\begin{minipage}[b]{\linewidth}\raggedright
No.
\end{minipage} & \begin{minipage}[b]{\linewidth}\raggedright
Name
\end{minipage} & \begin{minipage}[b]{\linewidth}\raggedright
Description
\end{minipage} & \begin{minipage}[b]{\linewidth}\raggedright
Default
\end{minipage} \\
\midrule\noalign{}
\endhead
\bottomrule\noalign{}
\endlastfoot
Field & Metric & Laboratory & \\
1 & SLIQ & A columnar vector of real monotonically increasing down the
column values starting from zero and terminating at one, that defines
the liquid saturation, that is the connate water saturation (SWL) plus
the oil saturation. The first entry should correspond to residual
liquid, that is Swc + Sorg and the last entry should be 1.0 to
correspond to a gas saturation of zero. & None \\
dimensionless & dimensionless & dimensionless & \\
2 & KRG & A columnar vector of real values that are either equal or
decreasing down the column and that are greater than or equal to zero
and less than or equal to one that defines the gas relative
permeability.. & None \\
dimensionless & dimensionless & dimensionless & \\
3 & KRO & A columnar vector of real values that are either equal or
increasing down the column and that are greater than or equal to zero
and less than or equal to one that defines the oil relative permeability
with respect to gas and connate water saturation. When water is active
in the run, the last entry the column, that is at krog(Sg = 0), must be
the same as the first entry in the corresponding SWOF table, that is at
krow(So = 1 - Swco). The first value in the column should be zero. &
None \\
dimensionless & dimensionless & dimensionless & \\
4 & PCOG & A columnar vector of real values that are either equal or
decreasing down the column that defines the oil-gas relative capillary
pressure. & None \\
psia & bars & atm & \\
Notes: & & & \\
\end{longtable}

\emph{Table 8.163: SLGOF Keyword Description}

\subsubsection{Example}\label{example-372}

\begin{verbatim}
--
--       GAS-OIL RELATIVE PERMEABILITY TABLES (SLGOF)
SLGOF
--       SLIQ       KRG       KROG      PCOG
--       FRAC                           PSIA
--       -------   --------  -------   -------
         0.30060    0.55000  0.0000     0.0000
         0.31480    0.42500  0.2848     0.0000
         0.32910    0.35000  0.3091     0.0000
         0.34330    0.27500  0.4333     0.0000
         0.35760    0.25000  0.5576     0.0000
         0.37180    0.22500  0.5818     0.0000
         0.38610    0.20000  0.6061     0.0000
         0.40030    0.17500  0.6303     0.0000
         0.41450    0.15000  0.6545     0.0000
         0.42880    0.12500  0.6788     0.0000
         0.44300    0.10000  0.7030     0.0000
         0.45730    0.07500  0.7273     0.0000
         0.47150    0.05000  0.7515     0.0000
         0.48580    0.02500  0.7758     0.0000
         0.50000    0.00000  0.8000     0.0000
         0.80000    0.00000  0.9000     0.0000             / TABLE No. 01
--       -------   --------  -------   -------
         0.30060    0.55000  0.0000     0.0000
         0.31480    0.42500  0.2848     0.0000
         0.32910    0.35000  0.3091     0.0000
         0.34330    0.27500  0.4333     0.0000
         0.35760    0.25000  0.5576     0.0000
         0.37180    0.22500  0.5818     0.0000
         0.38610    0.20000  0.6061     0.0000
         0.40030    0.17500  0.6303     0.0000
         0.41450    0.15000  0.6545     0.0000
         0.42880    0.12500  0.6788     0.0000
         0.44300    0.10000  0.7030     0.0000
         0.45730    0.07500  0.7273     0.0000
         0.47150    0.05000  0.7515     0.0000
         0.48580    0.02500  0.7758     0.0000
         0.50000    0.00000  0.8000     0.0000
         0.80000    0.00000  0.9000     0.0000             / TABLE No. 02
\end{verbatim}

The example defines two SLGOF tables for use when oil, gas and water are
present in the run.

\subsection{SOCRS -- End-Point Scaling Grid Cell Miscible Critical Oil
Saturation with Respect to
Water}\label{socrs-end-point-scaling-grid-cell-miscible-critical-oil-saturation-with-respect-to-water}

\begin{longtable}[]{@{}
  >{\raggedright\arraybackslash}p{(\linewidth - 14\tabcolsep) * \real{0.1250}}
  >{\raggedright\arraybackslash}p{(\linewidth - 14\tabcolsep) * \real{0.1250}}
  >{\raggedright\arraybackslash}p{(\linewidth - 14\tabcolsep) * \real{0.1250}}
  >{\raggedright\arraybackslash}p{(\linewidth - 14\tabcolsep) * \real{0.1250}}
  >{\raggedright\arraybackslash}p{(\linewidth - 14\tabcolsep) * \real{0.1250}}
  >{\raggedright\arraybackslash}p{(\linewidth - 14\tabcolsep) * \real{0.1250}}
  >{\raggedright\arraybackslash}p{(\linewidth - 14\tabcolsep) * \real{0.1250}}
  >{\raggedright\arraybackslash}p{(\linewidth - 14\tabcolsep) * \real{0.1250}}@{}}
\toprule\noalign{}
\begin{minipage}[b]{\linewidth}\raggedright
RUNSPEC
\end{minipage} & \begin{minipage}[b]{\linewidth}\raggedright
GRID
\end{minipage} & \begin{minipage}[b]{\linewidth}\raggedright
EDIT
\end{minipage} & \begin{minipage}[b]{\linewidth}\raggedright
PROPS
\end{minipage} & \begin{minipage}[b]{\linewidth}\raggedright
REGIONS
\end{minipage} & \begin{minipage}[b]{\linewidth}\raggedright
SOLUTION
\end{minipage} & \begin{minipage}[b]{\linewidth}\raggedright
SUMMARY
\end{minipage} & \begin{minipage}[b]{\linewidth}\raggedright
SCHEDULE
\end{minipage} \\
\midrule\noalign{}
\endhead
\bottomrule\noalign{}
\endlastfoot
\end{longtable}

\subsubsection{Description}\label{description-654}

SOCRS defines the miscible critical oil saturation for all the cells in
the model via an array when the end-point scaling option has been
invoked via the ENDSCALE keyword in the RUNSPEC section. The critical
oil saturation with respect to water is defined as the maximum oil
saturation for which the oil relative permeability is zero in a
two-phase oil-water relative permeability table. The keyword is used
with the Surfactant model to re-scale the surfactant relative
permeability saturation tables allocated to a grid block by the by the
SURFNUM keyword in the REGIONS section.

This keyword is not supported by OPM Flow but would change the results
if supported so the simulation will be stopped.

\begin{longtable}[]{@{}
  >{\raggedright\arraybackslash}p{(\linewidth - 6\tabcolsep) * \real{0.2500}}
  >{\raggedright\arraybackslash}p{(\linewidth - 6\tabcolsep) * \real{0.2500}}
  >{\raggedright\arraybackslash}p{(\linewidth - 6\tabcolsep) * \real{0.2500}}
  >{\raggedright\arraybackslash}p{(\linewidth - 6\tabcolsep) * \real{0.2500}}@{}}
\toprule\noalign{}
\begin{minipage}[b]{\linewidth}\raggedright
No.
\end{minipage} & \begin{minipage}[b]{\linewidth}\raggedright
Name
\end{minipage} & \begin{minipage}[b]{\linewidth}\raggedright
Description
\end{minipage} & \begin{minipage}[b]{\linewidth}\raggedright
Default
\end{minipage} \\
\midrule\noalign{}
\endhead
\bottomrule\noalign{}
\endlastfoot
Field & Metric & Laboratory & \\
1 & SOCRS & SOCRS is an array of real numbers assigning the critical oil
saturation with respect to water values to each cell in the model. The
number of entries should correspond to the NX x NY x NZ parameters on
the DIMENS keyword. Repeat counts may be used, for example 30*0.30 &
Taken from cell allocated relative permeability table. \\
dimensionless & dimensionless & dimensionless & \\
Notes: & & & \\
\end{longtable}

\emph{Table 8.164: SOGCR Keyword Description}

End-point scaling allows the entered relative permeability functions to
be re-scaled based on the saturation end-points allocated to each cell
by the SWL, SWCR, SWU, SGL, SGCR, SGU, SOCRS, and SOGCR saturation grid
arrays for the saturation end-points, and the KRG, KRORG, KRORW and KRW
relative permeability grid cell arrays for the relative permeability
end-point data.

\subsubsection{Example}\label{example-373}

\begin{verbatim}
--
--       DEFINE GRID BLOCK END-POINT SOCRS DATA FOR ALL CELLS
--       (FOR NX x NY x NZ = 300)
--
SOCRS
         300*0.200                                                              /
\end{verbatim}

The above example defines a constant critical oil saturation of 0.20 to
all 300 cells in the model as defined by the DIMENS keyword in the
RUNSPEC section.

\subsection{SOF2 -- Oil Saturation Tables with Respect to Gas or Water
(Format Type
2)}\label{sof2-oil-saturation-tables-with-respect-to-gas-or-water-format-type-2}

\begin{longtable}[]{@{}
  >{\raggedright\arraybackslash}p{(\linewidth - 14\tabcolsep) * \real{0.1250}}
  >{\raggedright\arraybackslash}p{(\linewidth - 14\tabcolsep) * \real{0.1250}}
  >{\raggedright\arraybackslash}p{(\linewidth - 14\tabcolsep) * \real{0.1250}}
  >{\raggedright\arraybackslash}p{(\linewidth - 14\tabcolsep) * \real{0.1250}}
  >{\raggedright\arraybackslash}p{(\linewidth - 14\tabcolsep) * \real{0.1250}}
  >{\raggedright\arraybackslash}p{(\linewidth - 14\tabcolsep) * \real{0.1250}}
  >{\raggedright\arraybackslash}p{(\linewidth - 14\tabcolsep) * \real{0.1250}}
  >{\raggedright\arraybackslash}p{(\linewidth - 14\tabcolsep) * \real{0.1250}}@{}}
\toprule\noalign{}
\begin{minipage}[b]{\linewidth}\raggedright
RUNSPEC
\end{minipage} & \begin{minipage}[b]{\linewidth}\raggedright
GRID
\end{minipage} & \begin{minipage}[b]{\linewidth}\raggedright
EDIT
\end{minipage} & \begin{minipage}[b]{\linewidth}\raggedright
PROPS
\end{minipage} & \begin{minipage}[b]{\linewidth}\raggedright
REGIONS
\end{minipage} & \begin{minipage}[b]{\linewidth}\raggedright
SOLUTION
\end{minipage} & \begin{minipage}[b]{\linewidth}\raggedright
SUMMARY
\end{minipage} & \begin{minipage}[b]{\linewidth}\raggedright
SCHEDULE
\end{minipage} \\
\midrule\noalign{}
\endhead
\bottomrule\noalign{}
\endlastfoot
\end{longtable}

\subsubsection{Description}\label{description-655}

The SOF2 keyword defines the oil relative permeability versus oil
saturation tables for when oil and gas or oil and water are present in
the input deck. The keyword is also used to define the relative
permeability of the miscible hydrocarbon phase in SOLVENT runs. This
keyword should only be used if oil is present in the run.

\begin{longtable}[]{@{}
  >{\raggedright\arraybackslash}p{(\linewidth - 6\tabcolsep) * \real{0.2500}}
  >{\raggedright\arraybackslash}p{(\linewidth - 6\tabcolsep) * \real{0.2500}}
  >{\raggedright\arraybackslash}p{(\linewidth - 6\tabcolsep) * \real{0.2500}}
  >{\raggedright\arraybackslash}p{(\linewidth - 6\tabcolsep) * \real{0.2500}}@{}}
\toprule\noalign{}
\begin{minipage}[b]{\linewidth}\raggedright
No.
\end{minipage} & \begin{minipage}[b]{\linewidth}\raggedright
Name
\end{minipage} & \begin{minipage}[b]{\linewidth}\raggedright
Description
\end{minipage} & \begin{minipage}[b]{\linewidth}\raggedright
Default
\end{minipage} \\
\midrule\noalign{}
\endhead
\bottomrule\noalign{}
\endlastfoot
Field & Metric & Laboratory & \\
1 & SOIL & A columnar vector of real monotonically increasing down the
column values starting from zero and terminating at one, that defines
the oil or the hydrocarbon solvent saturation. For two phase runs the
oil saturation should be entered and for when the SOLVENT option has
been activated in the RUNSPEC section the total hydrocarbon phase
(including the solvent) should be entered, that is SOIL = So + Sg + Ss.
& None \\
dimensionless & dimensionless & dimensionless & \\
3 & KRO & A columnar vector of real values that are either equal or
increasing down the column and that are greater than or equal to zero
and less than or equal to one that defines the oil relative permeability
with respect to gas and connate water saturation. For two phase runs the
oil relative permeability should be entered and for when the SOLVENT
option has been activated in the RUNSPEC section the relative
permeability of the miscible hydrocarbon phase with respect to water.
The last value in the column should be zero. & None \\
dimensionless & dimensionless & dimensionless & \\
Notes: & & & \\
\end{longtable}

\emph{Table 8.165: SOF2 Keyword Description}

Note this keyword should only be used for when the SOLVENT keyword in
the RUNSPEC section has been activated. It should not be use for
two-phase oil-water runs.

\subsubsection{Example}\label{example-374}

\begin{verbatim}
--
--       OIL RELATIVE PERMEABILITY TABLES (SOF2)
--
SOF2
--       SOIL         KRO
--       FRAC         FRAC
--       --------     --------
            0.00      0.000000
            0.05      1.197e-5
            0.10      0.000191
            0.15      0.000969
            0.20      0.003065
            0.25      0.007483
            0.30      0.015517
            0.35      0.028747
            0.40      0.049041
            0.45      0.078555
            0.50      0.119730
            0.55      0.175297
            0.60      0.248272
            0.65      0.341961
            0.70      0.459956
            0.75      0.606134
            0.80      0.784664
            0.85      1.000000                             / TABLE NO. 01
--       --------     --------
            0.00      0.000000
            0.05      1.197e-5
            0.10      0.000191
            0.15      0.000969
            0.20      0.003065
            0.25      0.007483
            0.30      0.015517
            0.35      0.028747
            0.40      0.049041
            0.45      0.078555
            0.50      0.119730
            0.55      0.175297
            0.60      0.248272
            0.65      0.341961
            0.70      0.459956
            0.75      0.606134
            0.80      0.784664
            0.85      1.000000                             / TABLE NO. 02
\end{verbatim}

The example defines two SOF2 tables for when oil and gas or oil and
water are present in the input deck.

\subsection{SOF3 -- Oil Saturation Tables with Respect to Gas and Water
(Format Type
2)}\label{sof3-oil-saturation-tables-with-respect-to-gas-and-water-format-type-2}

\begin{longtable}[]{@{}
  >{\raggedright\arraybackslash}p{(\linewidth - 14\tabcolsep) * \real{0.1250}}
  >{\raggedright\arraybackslash}p{(\linewidth - 14\tabcolsep) * \real{0.1250}}
  >{\raggedright\arraybackslash}p{(\linewidth - 14\tabcolsep) * \real{0.1250}}
  >{\raggedright\arraybackslash}p{(\linewidth - 14\tabcolsep) * \real{0.1250}}
  >{\raggedright\arraybackslash}p{(\linewidth - 14\tabcolsep) * \real{0.1250}}
  >{\raggedright\arraybackslash}p{(\linewidth - 14\tabcolsep) * \real{0.1250}}
  >{\raggedright\arraybackslash}p{(\linewidth - 14\tabcolsep) * \real{0.1250}}
  >{\raggedright\arraybackslash}p{(\linewidth - 14\tabcolsep) * \real{0.1250}}@{}}
\toprule\noalign{}
\begin{minipage}[b]{\linewidth}\raggedright
RUNSPEC
\end{minipage} & \begin{minipage}[b]{\linewidth}\raggedright
GRID
\end{minipage} & \begin{minipage}[b]{\linewidth}\raggedright
EDIT
\end{minipage} & \begin{minipage}[b]{\linewidth}\raggedright
PROPS
\end{minipage} & \begin{minipage}[b]{\linewidth}\raggedright
REGIONS
\end{minipage} & \begin{minipage}[b]{\linewidth}\raggedright
SOLUTION
\end{minipage} & \begin{minipage}[b]{\linewidth}\raggedright
SUMMARY
\end{minipage} & \begin{minipage}[b]{\linewidth}\raggedright
SCHEDULE
\end{minipage} \\
\midrule\noalign{}
\endhead
\bottomrule\noalign{}
\endlastfoot
\end{longtable}

\subsubsection{Description}\label{description-656}

The SOF3 keyword defines the oil relative permeability versus oil
saturation tables for when oil, gas and water are present in the input
deck. The keyword should only be used if oil, gas and water are present
in the input deck.

\begin{longtable}[]{@{}
  >{\raggedright\arraybackslash}p{(\linewidth - 6\tabcolsep) * \real{0.2500}}
  >{\raggedright\arraybackslash}p{(\linewidth - 6\tabcolsep) * \real{0.2500}}
  >{\raggedright\arraybackslash}p{(\linewidth - 6\tabcolsep) * \real{0.2500}}
  >{\raggedright\arraybackslash}p{(\linewidth - 6\tabcolsep) * \real{0.2500}}@{}}
\toprule\noalign{}
\begin{minipage}[b]{\linewidth}\raggedright
No.
\end{minipage} & \begin{minipage}[b]{\linewidth}\raggedright
Name
\end{minipage} & \begin{minipage}[b]{\linewidth}\raggedright
Description
\end{minipage} & \begin{minipage}[b]{\linewidth}\raggedright
Default
\end{minipage} \\
\midrule\noalign{}
\endhead
\bottomrule\noalign{}
\endlastfoot
Field & Metric & Laboratory & \\
1 & SOIL & A columnar vector of real monotonically increasing down the
column values starting from zero and terminating at one, that defines
the oil or the hydrocarbon solvent saturation. The final entry should be
at the connate water saturation, that is 1- Swc. & None \\
dimensionless & dimensionless & dimensionless & \\
3 & KROW & A columnar vector of real values that are either equal or
increasing down the column and that are greater than or equal to zero
and less than or equal to one that defines the oil relative permeability
with respect to oil and water saturation. The first value in the column
should be zero. & None \\
dimensionless & dimensionless & dimensionless & \\
4 & KROG & A columnar vector of real values that are either equal or
increasing down the column and that are greater than or equal to zero
and less than or equal to one that defines the oil relative permeability
with respect to oil, gas and connate water saturation. The first value
in the column should be zero. & None \\
dimensionless & dimensionless & dimensionless & \\
Notes: & & & \\
\end{longtable}

\emph{Table 8.166: SOF3 Keyword Description}

\subsubsection{Example}\label{example-375}

\begin{verbatim}
--
--       OIL RELATIVE PERMEABILITY TABLES (SOF3)
--
SOF3
--       SOIL         KRO          KROG
--       FRAC         FRAC         FRAC
--       --------     --------    --------
            0.00      0.000000    0.00000
            0.05      1.197e-5    0.00000
            0.10      0.000191    0.00000
            0.15      0.000969    0.00000
            0.20      0.003065    0.00000
            0.25      0.007483    0.00000
            0.30      0.015517    0.05932
            0.35      0.028747    0.13158
            0.40      0.049041    0.21082
            0.45      0.078555    0.29960
            0.50      0.119730    0.40095
            0.55      0.175297    0.51818
            0.60      0.248272    0.65476
            0.65      0.341961    0.81420
            0.70      0.459956    1.00000
            0.75      0.606134    1.00000
            0.80      0.784664    1.00000
            0.85      1.000000    1.00000                  / TABLE NO. 1
--       --------     --------    --------
            0.00      0.000000    0.00000
            0.05      1.197e-5    0.00000
            0.10      0.000191    0.00000
            0.15      0.000969    0.00000
            0.20      0.003065    0.00000
            0.25      0.007483    0.00000
            0.30      0.015517    0.05932
            0.35      0.028747    0.13158
            0.40      0.049041    0.21082
            0.45      0.078555    0.29960
            0.50      0.119730    0.40095
            0.55      0.175297    0.51818
            0.60      0.248272    0.65476
            0.65      0.341961    0.81420
            0.70      0.459956    1.00000
            0.75      0.606134    1.00000
            0.80      0.784664    1.00000
            0.85      1.000000    1.00000                  / TABLE NO. 2
\end{verbatim}

The example defines two SOF3 tables for when oil, gas and water are
present in the input deck.

\subsection{SOF32D -- Oil Saturation Tables with Respect to Water and
Gas (Three
Phase)}\label{sof32d-oil-saturation-tables-with-respect-to-water-and-gas-three-phase}

\begin{longtable}[]{@{}
  >{\raggedright\arraybackslash}p{(\linewidth - 14\tabcolsep) * \real{0.1250}}
  >{\raggedright\arraybackslash}p{(\linewidth - 14\tabcolsep) * \real{0.1250}}
  >{\raggedright\arraybackslash}p{(\linewidth - 14\tabcolsep) * \real{0.1250}}
  >{\raggedright\arraybackslash}p{(\linewidth - 14\tabcolsep) * \real{0.1250}}
  >{\raggedright\arraybackslash}p{(\linewidth - 14\tabcolsep) * \real{0.1250}}
  >{\raggedright\arraybackslash}p{(\linewidth - 14\tabcolsep) * \real{0.1250}}
  >{\raggedright\arraybackslash}p{(\linewidth - 14\tabcolsep) * \real{0.1250}}
  >{\raggedright\arraybackslash}p{(\linewidth - 14\tabcolsep) * \real{0.1250}}@{}}
\toprule\noalign{}
\begin{minipage}[b]{\linewidth}\raggedright
RUNSPEC
\end{minipage} & \begin{minipage}[b]{\linewidth}\raggedright
GRID
\end{minipage} & \begin{minipage}[b]{\linewidth}\raggedright
EDIT
\end{minipage} & \begin{minipage}[b]{\linewidth}\raggedright
PROPS
\end{minipage} & \begin{minipage}[b]{\linewidth}\raggedright
REGIONS
\end{minipage} & \begin{minipage}[b]{\linewidth}\raggedright
SOLUTION
\end{minipage} & \begin{minipage}[b]{\linewidth}\raggedright
SUMMARY
\end{minipage} & \begin{minipage}[b]{\linewidth}\raggedright
SCHEDULE
\end{minipage} \\
\midrule\noalign{}
\endhead
\bottomrule\noalign{}
\endlastfoot
\end{longtable}

\subsubsection{Description}\label{description-657}

The SOF32D keyword defines the three phase oil relative permeability
versus water and gas saturation tables for when oil, gas and water are
present in the input deck. The keyword should only be used if oil, gas
and water are present in the input deck. Normally the simulator
calculates the three-phase oil relative permeabilities based on the
entered two phase tables of water-oil and gas-oil, combined with the
STONE1 and STONE2 keywords in the PROPS section that determine the
method used to generate the thee phase oil relative permeability curves.
SOF32D allows for the direct input of the three phase tables, as such
the STONE1 and STONE2 keywords should not be entered if SOF32D is used
in the input deck.

This keyword is not supported by OPM Flow but would change the results
if supported so the simulation will be stopped.

\subsection{SOGCR -- End-Point Scaling Grid Cell Critical Oil Saturation
with Respect to
Gas}\label{sogcr-end-point-scaling-grid-cell-critical-oil-saturation-with-respect-to-gas}

\begin{longtable}[]{@{}
  >{\raggedright\arraybackslash}p{(\linewidth - 14\tabcolsep) * \real{0.1250}}
  >{\raggedright\arraybackslash}p{(\linewidth - 14\tabcolsep) * \real{0.1250}}
  >{\raggedright\arraybackslash}p{(\linewidth - 14\tabcolsep) * \real{0.1250}}
  >{\raggedright\arraybackslash}p{(\linewidth - 14\tabcolsep) * \real{0.1250}}
  >{\raggedright\arraybackslash}p{(\linewidth - 14\tabcolsep) * \real{0.1250}}
  >{\raggedright\arraybackslash}p{(\linewidth - 14\tabcolsep) * \real{0.1250}}
  >{\raggedright\arraybackslash}p{(\linewidth - 14\tabcolsep) * \real{0.1250}}
  >{\raggedright\arraybackslash}p{(\linewidth - 14\tabcolsep) * \real{0.1250}}@{}}
\toprule\noalign{}
\begin{minipage}[b]{\linewidth}\raggedright
RUNSPEC
\end{minipage} & \begin{minipage}[b]{\linewidth}\raggedright
GRID
\end{minipage} & \begin{minipage}[b]{\linewidth}\raggedright
EDIT
\end{minipage} & \begin{minipage}[b]{\linewidth}\raggedright
PROPS
\end{minipage} & \begin{minipage}[b]{\linewidth}\raggedright
REGIONS
\end{minipage} & \begin{minipage}[b]{\linewidth}\raggedright
SOLUTION
\end{minipage} & \begin{minipage}[b]{\linewidth}\raggedright
SUMMARY
\end{minipage} & \begin{minipage}[b]{\linewidth}\raggedright
SCHEDULE
\end{minipage} \\
\midrule\noalign{}
\endhead
\bottomrule\noalign{}
\endlastfoot
\end{longtable}

\subsubsection{Description}\label{description-658}

SOGCR defines the critical oil saturation for all the cells in the model
via an array when the end-point scaling option has been invoked via the
ENDSCALE keyword in the RUNSPEC section. The critical oil saturation
with respect to gas is defined as the maximum oil saturation for which
the oil relative permeability is zero in a two-phase gas-oil relative
permeability table. The keyword can be used with all grid types.

\begin{longtable}[]{@{}
  >{\raggedright\arraybackslash}p{(\linewidth - 6\tabcolsep) * \real{0.2500}}
  >{\raggedright\arraybackslash}p{(\linewidth - 6\tabcolsep) * \real{0.2500}}
  >{\raggedright\arraybackslash}p{(\linewidth - 6\tabcolsep) * \real{0.2500}}
  >{\raggedright\arraybackslash}p{(\linewidth - 6\tabcolsep) * \real{0.2500}}@{}}
\toprule\noalign{}
\begin{minipage}[b]{\linewidth}\raggedright
No.
\end{minipage} & \begin{minipage}[b]{\linewidth}\raggedright
Name
\end{minipage} & \begin{minipage}[b]{\linewidth}\raggedright
Description
\end{minipage} & \begin{minipage}[b]{\linewidth}\raggedright
Default
\end{minipage} \\
\midrule\noalign{}
\endhead
\bottomrule\noalign{}
\endlastfoot
Field & Metric & Laboratory & \\
1 & SOGCR & SOGCR is an array of real numbers assigning the critical oil
saturation with respect to gas values to each cell in the model. The
number of entries should correspond to the NX x NY x NZ parameters on
the DIMENS keyword. Repeat counts may be used, for example 30*0.30
dimensionless & Taken from cell allocated relative permeability
table. \\
dimensionless & dimensionless & dimensionless & \\
Notes: & & & \\
\end{longtable}

\emph{Table 8.167: SOGCR Keyword Description}

End-point scaling allows the entered relative permeability functions to
be re-scaled based on the saturation end-points allocated to each cell
by the SWL, SWCR, SWU, SGL, SGCR, SGU, SOWCR, and SOGCR saturation grid
arrays for the saturation end-points, and the KRG, KRGR, KRO, KRORG,
KRORW, KRW and KRWR relative permeability grid cell arrays for the
relative permeability end-point data. In addition end-point scaling may
be directional dependent in which case the directional dependent
versions of the aforementioned arrays should be used, that is SOGCRX,
SOGCRY and SOGCRZ instead of SOGCR. There is also the facility to make
the directional end-point scaling reversible or non-reversible and if
the non-reversible option is selected the non-reversible versions of the
aforementioned arrays should be used, that is SOGCRX, SOGCRX-, SOGCRY,
SOGCRY-, SOGCRZ and SOGCRZ-, instead of the SOGCR keyword.

\subsubsection{Example}\label{example-376}

\begin{verbatim}
--
--       DEFINE GRID BLOCK END-POINT SOGCR DATA FOR ALL CELLS
--       (FOR NX x NY x NZ = 300)
--
SOGCR
         300*0.200                                                              /
\end{verbatim}

The above example defines a constant critical gas saturation of 0.20 to
all 300 cells in the model as defined by the DIMENS keyword in the
RUNSPEC section.

\subsection{SOMGAS -- STONE1 Model Minimum Oil Saturation versus Gas
Saturation}\label{somgas-stone1-model-minimum-oil-saturation-versus-gas-saturation}

\begin{longtable}[]{@{}
  >{\raggedright\arraybackslash}p{(\linewidth - 14\tabcolsep) * \real{0.1250}}
  >{\raggedright\arraybackslash}p{(\linewidth - 14\tabcolsep) * \real{0.1250}}
  >{\raggedright\arraybackslash}p{(\linewidth - 14\tabcolsep) * \real{0.1250}}
  >{\raggedright\arraybackslash}p{(\linewidth - 14\tabcolsep) * \real{0.1250}}
  >{\raggedright\arraybackslash}p{(\linewidth - 14\tabcolsep) * \real{0.1250}}
  >{\raggedright\arraybackslash}p{(\linewidth - 14\tabcolsep) * \real{0.1250}}
  >{\raggedright\arraybackslash}p{(\linewidth - 14\tabcolsep) * \real{0.1250}}
  >{\raggedright\arraybackslash}p{(\linewidth - 14\tabcolsep) * \real{0.1250}}@{}}
\toprule\noalign{}
\begin{minipage}[b]{\linewidth}\raggedright
RUNSPEC
\end{minipage} & \begin{minipage}[b]{\linewidth}\raggedright
GRID
\end{minipage} & \begin{minipage}[b]{\linewidth}\raggedright
EDIT
\end{minipage} & \begin{minipage}[b]{\linewidth}\raggedright
PROPS
\end{minipage} & \begin{minipage}[b]{\linewidth}\raggedright
REGIONS
\end{minipage} & \begin{minipage}[b]{\linewidth}\raggedright
SOLUTION
\end{minipage} & \begin{minipage}[b]{\linewidth}\raggedright
SUMMARY
\end{minipage} & \begin{minipage}[b]{\linewidth}\raggedright
SCHEDULE
\end{minipage} \\
\midrule\noalign{}
\endhead
\bottomrule\noalign{}
\endlastfoot
\end{longtable}

\subsubsection{Description}\label{description-659}

This keyword defines the minimum oil saturation as a function of gas
saturation for Stone's {[}Stone, H. L. ``Probability Model for
Estimating Three-Phase Relative Permeability,'' paper SPE 2116, Journal
of Canadian Petroleum Technology (1973) 22, No.~2, 214-218.{]} first
three phase oil relative permeability model as modified by Aziz and
Settari {[}Aziz, K. and Settari, A. Petroleum Reservoir Simulation,
London, UK, Applied Science Publishers (1979), page 398.{]}. If the
STONE1 and STONE2 keywords are not present in the input deck then the
default three phase oil relative permeability model is employed. The
SOMGAS and STONE1 keywords should only be used in three phase runs
containing the oil, gas and water phases. The keyword is optional.

This keyword is not supported by OPM Flow but would change the results
if supported so the simulation will be stopped.

\subsection{SOMWAT -- STONE1 Model Minimum Oil Saturation versus Water
Saturation}\label{somwat-stone1-model-minimum-oil-saturation-versus-water-saturation}

\begin{longtable}[]{@{}
  >{\raggedright\arraybackslash}p{(\linewidth - 14\tabcolsep) * \real{0.1250}}
  >{\raggedright\arraybackslash}p{(\linewidth - 14\tabcolsep) * \real{0.1250}}
  >{\raggedright\arraybackslash}p{(\linewidth - 14\tabcolsep) * \real{0.1250}}
  >{\raggedright\arraybackslash}p{(\linewidth - 14\tabcolsep) * \real{0.1250}}
  >{\raggedright\arraybackslash}p{(\linewidth - 14\tabcolsep) * \real{0.1250}}
  >{\raggedright\arraybackslash}p{(\linewidth - 14\tabcolsep) * \real{0.1250}}
  >{\raggedright\arraybackslash}p{(\linewidth - 14\tabcolsep) * \real{0.1250}}
  >{\raggedright\arraybackslash}p{(\linewidth - 14\tabcolsep) * \real{0.1250}}@{}}
\toprule\noalign{}
\begin{minipage}[b]{\linewidth}\raggedright
RUNSPEC
\end{minipage} & \begin{minipage}[b]{\linewidth}\raggedright
GRID
\end{minipage} & \begin{minipage}[b]{\linewidth}\raggedright
EDIT
\end{minipage} & \begin{minipage}[b]{\linewidth}\raggedright
PROPS
\end{minipage} & \begin{minipage}[b]{\linewidth}\raggedright
REGIONS
\end{minipage} & \begin{minipage}[b]{\linewidth}\raggedright
SOLUTION
\end{minipage} & \begin{minipage}[b]{\linewidth}\raggedright
SUMMARY
\end{minipage} & \begin{minipage}[b]{\linewidth}\raggedright
SCHEDULE
\end{minipage} \\
\midrule\noalign{}
\endhead
\bottomrule\noalign{}
\endlastfoot
\end{longtable}

\subsubsection{Description}\label{description-660}

This keyword defines the minimum oil saturation as a function of water
saturation for Stone's {[}Stone, H. L. ``Probability Model for
Estimating Three-Phase Relative Permeability,'' paper SPE 2116, Journal
of Canadian Petroleum Technology (1973) 22, No.~2, 214-218.{]} first
three phase oil relative permeability model as modified by Aziz and
Settari {[}Aziz, K. and Settari, A. Petroleum Reservoir Simulation,
London, UK, Applied Science Publishers (1979), page 398.{]}. If the
STONE1 and STONE2 keywords are not present in the input deck then the
default three phase oil relative permeability model is employed. The
SOMWAT and STONE1 keywords should only be used in three phase runs
containing the oil, gas and water phases. The keyword is optional.

This keyword is not supported by OPM Flow but would change the results
if supported so the simulation will be stopped.

\subsection{SORWMIS -- Miscible Residual Oil versus Water Saturation
Functions}\label{sorwmis-miscible-residual-oil-versus-water-saturation-functions}

\begin{longtable}[]{@{}
  >{\raggedright\arraybackslash}p{(\linewidth - 14\tabcolsep) * \real{0.1250}}
  >{\raggedright\arraybackslash}p{(\linewidth - 14\tabcolsep) * \real{0.1250}}
  >{\raggedright\arraybackslash}p{(\linewidth - 14\tabcolsep) * \real{0.1250}}
  >{\raggedright\arraybackslash}p{(\linewidth - 14\tabcolsep) * \real{0.1250}}
  >{\raggedright\arraybackslash}p{(\linewidth - 14\tabcolsep) * \real{0.1250}}
  >{\raggedright\arraybackslash}p{(\linewidth - 14\tabcolsep) * \real{0.1250}}
  >{\raggedright\arraybackslash}p{(\linewidth - 14\tabcolsep) * \real{0.1250}}
  >{\raggedright\arraybackslash}p{(\linewidth - 14\tabcolsep) * \real{0.1250}}@{}}
\toprule\noalign{}
\begin{minipage}[b]{\linewidth}\raggedright
RUNSPEC
\end{minipage} & \begin{minipage}[b]{\linewidth}\raggedright
GRID
\end{minipage} & \begin{minipage}[b]{\linewidth}\raggedright
EDIT
\end{minipage} & \begin{minipage}[b]{\linewidth}\raggedright
PROPS
\end{minipage} & \begin{minipage}[b]{\linewidth}\raggedright
REGIONS
\end{minipage} & \begin{minipage}[b]{\linewidth}\raggedright
SOLUTION
\end{minipage} & \begin{minipage}[b]{\linewidth}\raggedright
SUMMARY
\end{minipage} & \begin{minipage}[b]{\linewidth}\raggedright
SCHEDULE
\end{minipage} \\
\midrule\noalign{}
\endhead
\bottomrule\noalign{}
\endlastfoot
\end{longtable}

\subsubsection{Description}\label{description-661}

SORWMIS defines the dependency between the miscible residual oil
saturation and the water saturation, for when the MISCIBLE keyword in
the RUNSPEC section has been activated. The keyword can only be used
with the MISCIBLE option and for when the oil, water and gas phases are
active in the model.

\begin{longtable}[]{@{}
  >{\raggedright\arraybackslash}p{(\linewidth - 6\tabcolsep) * \real{0.2500}}
  >{\raggedright\arraybackslash}p{(\linewidth - 6\tabcolsep) * \real{0.2500}}
  >{\raggedright\arraybackslash}p{(\linewidth - 6\tabcolsep) * \real{0.2500}}
  >{\raggedright\arraybackslash}p{(\linewidth - 6\tabcolsep) * \real{0.2500}}@{}}
\toprule\noalign{}
\begin{minipage}[b]{\linewidth}\raggedright
No.
\end{minipage} & \begin{minipage}[b]{\linewidth}\raggedright
Name
\end{minipage} & \begin{minipage}[b]{\linewidth}\raggedright
Description
\end{minipage} & \begin{minipage}[b]{\linewidth}\raggedright
Default
\end{minipage} \\
\midrule\noalign{}
\endhead
\bottomrule\noalign{}
\endlastfoot
Field & Metric & Laboratory & \\
1 & SWAT & A columnar vector of real monotonically increasing down the
column values starting from zero and terminating at one, that defines
the water saturation. & None \\
dimensionless & dimensionless & dimensionless & \\
2 & SORMIS & A columnar vector of real equal or increasing down the
column values that are greater than or equal to zero and less than one,
that define the corresponding miscible residual oil saturation for the
corresponding water saturation SWAT. & None \\
dimensionless & dimensionless & dimensionless & \\
Notes: & & & \\
\end{longtable}

\emph{Table 8.168: SORWMIS Keyword Description}

\subsubsection{Example}\label{example-377}

\begin{verbatim}
--
--       MISCIBLE RESIDUAL OIL VERSUS WATER SATURATION TABLE
--
SORWMIS
--       SWAT       SORWMIS
--       FRAC       FRAC
--       -------    --------
          0.0000     0.0000
          0.2000     0.0000
          1.0000     0.0000                                / TABLE NO. 01
--       SWAT       SORWMIS
--       FRAC       FRAC
--       -------    --------
          0.0000     0.0000
          0.3000     0.1000
          0.7500     0.1500                                / TABLE NO. 02
\end{verbatim}

The above example defines two miscible residual oil versus water
saturation tables assuming NTMISC equals two and NSMISC is greater than
or equal to three on the MISCIBLE keyword in the RUNSPEC section.

\subsection{SOWCR -- End-Point Scaling Grid Cell Critical Oil Saturation
with Respect to
Water}\label{sowcr-end-point-scaling-grid-cell-critical-oil-saturation-with-respect-to-water}

\begin{longtable}[]{@{}
  >{\raggedright\arraybackslash}p{(\linewidth - 14\tabcolsep) * \real{0.1250}}
  >{\raggedright\arraybackslash}p{(\linewidth - 14\tabcolsep) * \real{0.1250}}
  >{\raggedright\arraybackslash}p{(\linewidth - 14\tabcolsep) * \real{0.1250}}
  >{\raggedright\arraybackslash}p{(\linewidth - 14\tabcolsep) * \real{0.1250}}
  >{\raggedright\arraybackslash}p{(\linewidth - 14\tabcolsep) * \real{0.1250}}
  >{\raggedright\arraybackslash}p{(\linewidth - 14\tabcolsep) * \real{0.1250}}
  >{\raggedright\arraybackslash}p{(\linewidth - 14\tabcolsep) * \real{0.1250}}
  >{\raggedright\arraybackslash}p{(\linewidth - 14\tabcolsep) * \real{0.1250}}@{}}
\toprule\noalign{}
\begin{minipage}[b]{\linewidth}\raggedright
RUNSPEC
\end{minipage} & \begin{minipage}[b]{\linewidth}\raggedright
GRID
\end{minipage} & \begin{minipage}[b]{\linewidth}\raggedright
EDIT
\end{minipage} & \begin{minipage}[b]{\linewidth}\raggedright
PROPS
\end{minipage} & \begin{minipage}[b]{\linewidth}\raggedright
REGIONS
\end{minipage} & \begin{minipage}[b]{\linewidth}\raggedright
SOLUTION
\end{minipage} & \begin{minipage}[b]{\linewidth}\raggedright
SUMMARY
\end{minipage} & \begin{minipage}[b]{\linewidth}\raggedright
SCHEDULE
\end{minipage} \\
\midrule\noalign{}
\endhead
\bottomrule\noalign{}
\endlastfoot
\end{longtable}

\subsubsection{Description}\label{description-662}

SOWCR defines the critical oil saturation for all the cells in the model
via an array when the end-point scaling option has been invoked via the
ENDSCALE keyword in the RUNSPEC section. The critical oil saturation
with respect to water is defined as the maximum oil saturation for which
the oil relative permeability is zero in a two-phase oil-water relative
permeability table.

The keyword can be used with all grid types.

\begin{longtable}[]{@{}
  >{\raggedright\arraybackslash}p{(\linewidth - 6\tabcolsep) * \real{0.2500}}
  >{\raggedright\arraybackslash}p{(\linewidth - 6\tabcolsep) * \real{0.2500}}
  >{\raggedright\arraybackslash}p{(\linewidth - 6\tabcolsep) * \real{0.2500}}
  >{\raggedright\arraybackslash}p{(\linewidth - 6\tabcolsep) * \real{0.2500}}@{}}
\toprule\noalign{}
\begin{minipage}[b]{\linewidth}\raggedright
No.
\end{minipage} & \begin{minipage}[b]{\linewidth}\raggedright
Name
\end{minipage} & \begin{minipage}[b]{\linewidth}\raggedright
Description
\end{minipage} & \begin{minipage}[b]{\linewidth}\raggedright
Default
\end{minipage} \\
\midrule\noalign{}
\endhead
\bottomrule\noalign{}
\endlastfoot
Field & Metric & Laboratory & \\
1 & SOWCR & SOWCR is an array of real numbers assigning the critical oil
saturation with respect to water values to each cell in the model. The
number of entries should correspond to the NX x NY x NZ parameters on
the DIMENS keyword. Repeat counts may be used, for example 30*0.30 &
Taken from cell allocated relative permeability table. \\
dimensionless & dimensionless & dimensionless & \\
Notes: & & & \\
\end{longtable}

\emph{Table 8.169: SOWCR Keyword Description}

End-point scaling allows the entered relative permeability functions to
be re-scaled based on the saturation end-points allocated to each cell
by the SWL, SWCR, SWU, SGL, SGCR, SGU, SOWCR, and SOGCR saturation grid
arrays for the saturation end-points, and the KRG, KRGR, KRO, KRORG,
KRORW, KRW and KRWR relative permeability grid cell arrays for the
relative permeability end-point data. In addition end-point scaling may
be directional dependent in which case the directional dependent
versions of the aforementioned arrays should be used, that is SOWCRX,
SOWCRY and SOWCRZ instead of SOWCR. There is also the facility to make
the directional end-point scaling reversible or non-reversible and if
the non-reversible option is selected the non-reversible versions of the
aforementioned arrays should be used, that is SOWCRX, SOWCRX-, SOWCRY,
SOWCRY-, SOWCRZ and SOWCRZ-, instead of the SOWCR keyword.

\subsubsection{Example}\label{example-378}

\begin{verbatim}
--
--       DEFINE GRID BLOCK END-POINT SOWCR DATA FOR ALL CELLS
–        (FOR NX x NY x NZ = 300)
--
SOWCR
         300*0.200                                                              /
\end{verbatim}

The above example defines a constant critical oil saturation of 0.20 to
all 300 cells in the model as defined by the DIMENS keyword in the
RUNSPEC section.

\subsection{SPECHEAT -- Define the Specific Heat of Oil, Water and
Gas}\label{specheat-define-the-specific-heat-of-oil-water-and-gas}

\begin{longtable}[]{@{}
  >{\raggedright\arraybackslash}p{(\linewidth - 14\tabcolsep) * \real{0.1250}}
  >{\raggedright\arraybackslash}p{(\linewidth - 14\tabcolsep) * \real{0.1250}}
  >{\raggedright\arraybackslash}p{(\linewidth - 14\tabcolsep) * \real{0.1250}}
  >{\raggedright\arraybackslash}p{(\linewidth - 14\tabcolsep) * \real{0.1250}}
  >{\raggedright\arraybackslash}p{(\linewidth - 14\tabcolsep) * \real{0.1250}}
  >{\raggedright\arraybackslash}p{(\linewidth - 14\tabcolsep) * \real{0.1250}}
  >{\raggedright\arraybackslash}p{(\linewidth - 14\tabcolsep) * \real{0.1250}}
  >{\raggedright\arraybackslash}p{(\linewidth - 14\tabcolsep) * \real{0.1250}}@{}}
\toprule\noalign{}
\begin{minipage}[b]{\linewidth}\raggedright
RUNSPEC
\end{minipage} & \begin{minipage}[b]{\linewidth}\raggedright
GRID
\end{minipage} & \begin{minipage}[b]{\linewidth}\raggedright
EDIT
\end{minipage} & \begin{minipage}[b]{\linewidth}\raggedright
PROPS
\end{minipage} & \begin{minipage}[b]{\linewidth}\raggedright
REGIONS
\end{minipage} & \begin{minipage}[b]{\linewidth}\raggedright
SOLUTION
\end{minipage} & \begin{minipage}[b]{\linewidth}\raggedright
SUMMARY
\end{minipage} & \begin{minipage}[b]{\linewidth}\raggedright
SCHEDULE
\end{minipage} \\
\midrule\noalign{}
\endhead
\bottomrule\noalign{}
\endlastfoot
\end{longtable}

\subsubsection{Description}\label{description-663}

SPECHEAT defines the specific heat of the oil, water and gas phases for
various PVT regions in the model for when the THERMAL option has been
activated in the RUNSPEC section. The number of SPECHEAT vector data
sets is defined by the NTPVT parameter on the TABDIMS keyword in the
RUNSPEC section and the allocation of the SPECHEAT data sets to
different grid blocks in the model is done via the PVTNUM keyword in the
REGIONS section.

This keyword can only be used if OPM Flow's thermal option has been
activated by the THERMAL keyword in the RUNSPEC section. Note this is
different to the commercial simulator that uses the TEMP keyword in the
RUNSPEC section to activate the black-oil thermal model.

\begin{longtable}[]{@{}
  >{\raggedright\arraybackslash}p{(\linewidth - 6\tabcolsep) * \real{0.2500}}
  >{\raggedright\arraybackslash}p{(\linewidth - 6\tabcolsep) * \real{0.2500}}
  >{\raggedright\arraybackslash}p{(\linewidth - 6\tabcolsep) * \real{0.2500}}
  >{\raggedright\arraybackslash}p{(\linewidth - 6\tabcolsep) * \real{0.2500}}@{}}
\toprule\noalign{}
\begin{minipage}[b]{\linewidth}\raggedright
No.
\end{minipage} & \begin{minipage}[b]{\linewidth}\raggedright
Name
\end{minipage} & \begin{minipage}[b]{\linewidth}\raggedright
Description
\end{minipage} & \begin{minipage}[b]{\linewidth}\raggedright
Default
\end{minipage} \\
\midrule\noalign{}
\endhead
\bottomrule\noalign{}
\endlastfoot
Field & Metric & Laboratory & \\
1 & TEMP & A columnar vector of real monotonically increasing down the
column values that define the temperature for the corresponding oil,
water and gas specific heat values. & None \\
oF & oC & oC & \\
2 & OILSHEAT & OILSHEAT is a columnar vector of positive real numbers
defining the specific heat of oil at the corresponding temperature,
TEMP. & None \\
Btu/lb/oR & kJ/kg/K & J/gm/K & \\
3 & WATSHEAT & WATSHEAT is a columnar vector of positive real numbers
defining the specific heat of water at the corresponding temperature,
TEMP. & None \\
Btu/lb/oR & kJ/kg/K & J/gm/K & \\
4 & GASSHEAT & GASHEAT is a columnar vector of positive real numbers
defining the specific heat of gas at the corresponding temperature,
TEMP. & None \\
Btu/lb/oR & kJ/kg/K & J/gm/K & \\
Notes: & & & \\
\end{longtable}

\emph{Table 8.170: SPECHEAT Keyword Description}

See also the SPECROCK keyword to define the reservoir rock specific
heat.

\subsubsection{Example}\label{example-379}

The example below defines three fluid phases specific heat versus
temperature tables assuming NTPVT equals three and NPPVT is greater than
or equal to two on the TABDIMS keyword in the RUNSPEC section.

\begin{verbatim}
--
--       SPECIFIC HEAT OF OIL, WATER AND GAS TABLE
--
SPECHEAT
--       TEMP       SPECHEAT   SPECHEAT   SPECHEAT
--                  OIL        WATER      GAS
--       -------    --------   --------   --------
           0.000     0.5000     1.5000     0.5000
         250.000     0.5000     1.5000     0.5000          / TABLE NO. 01
--       TEMP       SPECHEAT   SPECHEAT   SPECHEAT
--                  OIL        WATER      GAS
--       -------    --------   --------   --------
           0.000     0.5500     1.5000     0.5000
         260.000     0.5500     1.5000     0.5000          / TABLE NO. 02
--       TEMP       SPECHEAT   SPECHEAT   SPECHEAT
--                  OIL        WATER      GAS
--       -------    --------   --------   --------
           0.000     0.5500     1.5500     0.5000
         270.000     0.6000     1.5500     0.5000          / TABLE NO. 03
\end{verbatim}

There is no terminating ``/'' for this keyword.

\subsection{SPECROCK -- Define the Specific Heat of the Reservoir
Rock}\label{specrock-define-the-specific-heat-of-the-reservoir-rock}

\begin{longtable}[]{@{}
  >{\raggedright\arraybackslash}p{(\linewidth - 14\tabcolsep) * \real{0.1250}}
  >{\raggedright\arraybackslash}p{(\linewidth - 14\tabcolsep) * \real{0.1250}}
  >{\raggedright\arraybackslash}p{(\linewidth - 14\tabcolsep) * \real{0.1250}}
  >{\raggedright\arraybackslash}p{(\linewidth - 14\tabcolsep) * \real{0.1250}}
  >{\raggedright\arraybackslash}p{(\linewidth - 14\tabcolsep) * \real{0.1250}}
  >{\raggedright\arraybackslash}p{(\linewidth - 14\tabcolsep) * \real{0.1250}}
  >{\raggedright\arraybackslash}p{(\linewidth - 14\tabcolsep) * \real{0.1250}}
  >{\raggedright\arraybackslash}p{(\linewidth - 14\tabcolsep) * \real{0.1250}}@{}}
\toprule\noalign{}
\begin{minipage}[b]{\linewidth}\raggedright
RUNSPEC
\end{minipage} & \begin{minipage}[b]{\linewidth}\raggedright
GRID
\end{minipage} & \begin{minipage}[b]{\linewidth}\raggedright
EDIT
\end{minipage} & \begin{minipage}[b]{\linewidth}\raggedright
PROPS
\end{minipage} & \begin{minipage}[b]{\linewidth}\raggedright
REGIONS
\end{minipage} & \begin{minipage}[b]{\linewidth}\raggedright
SOLUTION
\end{minipage} & \begin{minipage}[b]{\linewidth}\raggedright
SUMMARY
\end{minipage} & \begin{minipage}[b]{\linewidth}\raggedright
SCHEDULE
\end{minipage} \\
\midrule\noalign{}
\endhead
\bottomrule\noalign{}
\endlastfoot
\end{longtable}

\subsubsection{Description}\label{description-664}

SPECROCK defines the specific heat of the reservoir rock for various PVT
regions in the model for when the THERMAL option has been activated in
the RUNSPEC section. The number of SPECROCK vector data sets is defined
by the NTSFUN parameter on the TABDIMS keyword in the RUNSPEC section
and the allocation of the SPECROCK data sets to different grid blocks in
the model is done via the SATNUM keyword in the REGIONS section.

This keyword can only be used if OPM's Flow's thermal option has been
activated by the THERMAL keyword in the RUNSPEC section. Note this is
different to the commercial simulator that uses the TEMP keyword in the
RUNSPEC section to activate the black-oil thermal model.

\begin{longtable}[]{@{}
  >{\raggedright\arraybackslash}p{(\linewidth - 6\tabcolsep) * \real{0.2500}}
  >{\raggedright\arraybackslash}p{(\linewidth - 6\tabcolsep) * \real{0.2500}}
  >{\raggedright\arraybackslash}p{(\linewidth - 6\tabcolsep) * \real{0.2500}}
  >{\raggedright\arraybackslash}p{(\linewidth - 6\tabcolsep) * \real{0.2500}}@{}}
\toprule\noalign{}
\begin{minipage}[b]{\linewidth}\raggedright
No.
\end{minipage} & \begin{minipage}[b]{\linewidth}\raggedright
Name
\end{minipage} & \begin{minipage}[b]{\linewidth}\raggedright
Description
\end{minipage} & \begin{minipage}[b]{\linewidth}\raggedright
Default
\end{minipage} \\
\midrule\noalign{}
\endhead
\bottomrule\noalign{}
\endlastfoot
Field & Metric & Laboratory & \\
1 & TEMP & A columnar vector of real monotonically increasing down the
column values that define the temperature for the corresponding rock
specific heat values. & None \\
oF & oC & oC & \\
2 & ROCKHEAT & ROCKHEAT is a columnar vector of positive real numbers
defining the specific heat of the rock at the corresponding temperature,
TEMP. & None \\
Btu/ft3/oR & kJ/m3/K & J/cc/K & \\
Notes: & & & \\
\end{longtable}

\emph{Table 8.171: SPECROCK Keyword Description}

See also the SPECHEAT keyword to define the specific heat relationships
for the oil, water and gas phases.

\subsubsection{Example}\label{example-380}

The example below defines three rock specific heat versus temperature
tables assuming NTSFUN equals three and NSSFUN is greater than or equal
to two on the TABDIMS keyword in the RUNSPEC section.

\begin{verbatim}
--
--       SPECIFIC HEAT OF ROCK
--
SPECROCK
--       TEMP       SPECHEAT
--                  ROCK
--       -------    --------
           0.000     20.000
         250.000     20.000                                / TABLE NO. 01
--       -------    --------
           0.000     21.000
         260.000     21.000                                / TABLE NO. 02
--       -------    --------
           0.000     23.000
         270.000     23.000                                / TABLE NO. 03
\end{verbatim}

There is no terminating ``/'' for this keyword.

\subsection{SSFN -- Solvent and Gas Relative Permeability
Tables}\label{ssfn-solvent-and-gas-relative-permeability-tables}

\begin{longtable}[]{@{}
  >{\raggedright\arraybackslash}p{(\linewidth - 14\tabcolsep) * \real{0.1250}}
  >{\raggedright\arraybackslash}p{(\linewidth - 14\tabcolsep) * \real{0.1250}}
  >{\raggedright\arraybackslash}p{(\linewidth - 14\tabcolsep) * \real{0.1250}}
  >{\raggedright\arraybackslash}p{(\linewidth - 14\tabcolsep) * \real{0.1250}}
  >{\raggedright\arraybackslash}p{(\linewidth - 14\tabcolsep) * \real{0.1250}}
  >{\raggedright\arraybackslash}p{(\linewidth - 14\tabcolsep) * \real{0.1250}}
  >{\raggedright\arraybackslash}p{(\linewidth - 14\tabcolsep) * \real{0.1250}}
  >{\raggedright\arraybackslash}p{(\linewidth - 14\tabcolsep) * \real{0.1250}}@{}}
\toprule\noalign{}
\begin{minipage}[b]{\linewidth}\raggedright
RUNSPEC
\end{minipage} & \begin{minipage}[b]{\linewidth}\raggedright
GRID
\end{minipage} & \begin{minipage}[b]{\linewidth}\raggedright
EDIT
\end{minipage} & \begin{minipage}[b]{\linewidth}\raggedright
PROPS
\end{minipage} & \begin{minipage}[b]{\linewidth}\raggedright
REGIONS
\end{minipage} & \begin{minipage}[b]{\linewidth}\raggedright
SOLUTION
\end{minipage} & \begin{minipage}[b]{\linewidth}\raggedright
SUMMARY
\end{minipage} & \begin{minipage}[b]{\linewidth}\raggedright
SCHEDULE
\end{minipage} \\
\midrule\noalign{}
\endhead
\bottomrule\noalign{}
\endlastfoot
\end{longtable}

\subsubsection{Description}\label{description-665}

The SSFN keyword defines the miscible normalized relative permeability
tables for when the SOLVENT option has been activated in the RUNSPEC
section using the respective keyword. The SOLVENT keyword results in a
four component model (oil, water and gas, plus a solvent). This keyword
should only be used if the SOLVENT option has been activated.

\begin{longtable}[]{@{}
  >{\raggedright\arraybackslash}p{(\linewidth - 6\tabcolsep) * \real{0.2500}}
  >{\raggedright\arraybackslash}p{(\linewidth - 6\tabcolsep) * \real{0.2500}}
  >{\raggedright\arraybackslash}p{(\linewidth - 6\tabcolsep) * \real{0.2500}}
  >{\raggedright\arraybackslash}p{(\linewidth - 6\tabcolsep) * \real{0.2500}}@{}}
\toprule\noalign{}
\begin{minipage}[b]{\linewidth}\raggedright
No.
\end{minipage} & \begin{minipage}[b]{\linewidth}\raggedright
Name
\end{minipage} & \begin{minipage}[b]{\linewidth}\raggedright
Description
\end{minipage} & \begin{minipage}[b]{\linewidth}\raggedright
Default
\end{minipage} \\
\midrule\noalign{}
\endhead
\bottomrule\noalign{}
\endlastfoot
Field & Metric & Laboratory & \\
1 & SGAS & A columnar vector of real monotonically increasing down the
column values starting from zero and terminating at one, that defines
the gas plus solvent saturation ration which is defined as either:
\(\frac{{S}_{g}}{\left({S}_{g} + {S}_{s}\right)}\)or
\(\frac{{S}_{s}}{\left({S}_{g} + {S}_{s}\right)}\) Where Sg is the gas
saturation and Ss is the solvent saturation. & None \\
dimensionless & dimensionless & dimensionless & \\
2 & KRGt & A columnar vector of real values that are either equal or
increasing down the column and that are greater than or equal to zero
and less than or equal to one that defines the gas relative
permeability. The resulting gas relative permeability is calculated
from:
\({k}_{\mathit{rg}} = {k}_{\mathit{rgt}}\left({S}_{g} + {S}_{s}\right){{k}_{\mathit{rg}}}^{t}\)
where krgt is the data in this column and krgt is the gas relative
permeability from the SGFN keyword. & None \\
dimensionless & dimensionless & dimensionless & \\
3 & KRSt & A columnar vector of real values that are either equal or
increasing down the column and that are greater than or equal to zero
and less than or equal to one that defines the solvent relative
permeability. The resulting solvent relative permeability is calculated
from:
\({k}_{\mathit{rs}} = {k}_{\mathit{rgt}}\left({S}_{g} + {S}_{s}\right){{k}_{\mathit{rs}}}^{t}\)
where krSt is the data in this column and krgt is the gas relative
permeability from the SGFN keyword. & None \\
dimensionless & dimensionless & dimensionless & \\
Notes: & & & \\
\end{longtable}

\emph{Table 8.172: SSFN Keyword Description}

\subsubsection{Example}\label{example-381}

\begin{verbatim}
--
--       SOLVENT RELATIVE PERMEABILITY TABLES
--
SSFN
--       SGAS       KRGT         KRST
--       FRAC
--       --------  --------     --------
          0.0000    0.0000       0.0000
          1.0000    1.0000       1.0000                    / TABLE NO. 01
--       --------  --------     --------
          0.0000    0.0000       0.0000
          0.2000    0.2000       0.3000
          0.4000    0.3000       0.5000
          0.6000    0.4000       0.7000
          0.8000    0.5000       0.7500
          1.0000    1.0000       1.0000                    / TABLE NO. 02
\end{verbatim}

The above example defines two SSFN tables for use with the SOLVENT
option.

\subsection{SSGCR -- End-Point Scaling Grid Cell Surfactant Critical Gas
Saturations}\label{ssgcr-end-point-scaling-grid-cell-surfactant-critical-gas-saturations}

\begin{longtable}[]{@{}
  >{\raggedright\arraybackslash}p{(\linewidth - 14\tabcolsep) * \real{0.1250}}
  >{\raggedright\arraybackslash}p{(\linewidth - 14\tabcolsep) * \real{0.1250}}
  >{\raggedright\arraybackslash}p{(\linewidth - 14\tabcolsep) * \real{0.1250}}
  >{\raggedright\arraybackslash}p{(\linewidth - 14\tabcolsep) * \real{0.1250}}
  >{\raggedright\arraybackslash}p{(\linewidth - 14\tabcolsep) * \real{0.1250}}
  >{\raggedright\arraybackslash}p{(\linewidth - 14\tabcolsep) * \real{0.1250}}
  >{\raggedright\arraybackslash}p{(\linewidth - 14\tabcolsep) * \real{0.1250}}
  >{\raggedright\arraybackslash}p{(\linewidth - 14\tabcolsep) * \real{0.1250}}@{}}
\toprule\noalign{}
\begin{minipage}[b]{\linewidth}\raggedright
RUNSPEC
\end{minipage} & \begin{minipage}[b]{\linewidth}\raggedright
GRID
\end{minipage} & \begin{minipage}[b]{\linewidth}\raggedright
EDIT
\end{minipage} & \begin{minipage}[b]{\linewidth}\raggedright
PROPS
\end{minipage} & \begin{minipage}[b]{\linewidth}\raggedright
REGIONS
\end{minipage} & \begin{minipage}[b]{\linewidth}\raggedright
SOLUTION
\end{minipage} & \begin{minipage}[b]{\linewidth}\raggedright
SUMMARY
\end{minipage} & \begin{minipage}[b]{\linewidth}\raggedright
SCHEDULE
\end{minipage} \\
\midrule\noalign{}
\endhead
\bottomrule\noalign{}
\endlastfoot
\end{longtable}

\subsubsection{Description}\label{description-666}

SSGCR defines the surfactant critical gas saturation for all the cells
in the model via an array when the end-point scaling option has been
invoked via the ENDSCALE keyword and the surfactant phase has been
activated by the SURFACT keyword in the RUNSPEC section. The critical
gas saturation is defined as the maximum gas saturation for which the
gas relative permeability is zero in a two-phase relative permeability
table. SSGCR is used to scale the surfactant oil relative permeability
to gas data.

This keyword is not supported by OPM Flow but would change the results
if supported so the simulation will be stopped.

\begin{longtable}[]{@{}
  >{\raggedright\arraybackslash}p{(\linewidth - 6\tabcolsep) * \real{0.2500}}
  >{\raggedright\arraybackslash}p{(\linewidth - 6\tabcolsep) * \real{0.2500}}
  >{\raggedright\arraybackslash}p{(\linewidth - 6\tabcolsep) * \real{0.2500}}
  >{\raggedright\arraybackslash}p{(\linewidth - 6\tabcolsep) * \real{0.2500}}@{}}
\toprule\noalign{}
\begin{minipage}[b]{\linewidth}\raggedright
No.
\end{minipage} & \begin{minipage}[b]{\linewidth}\raggedright
Name
\end{minipage} & \begin{minipage}[b]{\linewidth}\raggedright
Description
\end{minipage} & \begin{minipage}[b]{\linewidth}\raggedright
Default
\end{minipage} \\
\midrule\noalign{}
\endhead
\bottomrule\noalign{}
\endlastfoot
Field & Metric & Laboratory & \\
1 & SSGCR & SSGCR is an array of real numbers assigning the surfactant
critical gas saturation values to each cell in the model. The number of
entries should correspond to the NX x NY x NZ parameters on the DIMENS
keyword. Repeat counts may be used, for example 30*0.03 & Taken from
cell allocated relative permeability table. \\
dimensionless & dimensionless & dimensionless & \\
Notes: & & & \\
\end{longtable}

\emph{Table 8.173: SSGCR Keyword Description}

End-point scaling allows the entered surfactant relative permeability
functions to be re-scaled based on the saturation end-points allocated
to each cell by the SSWL, SSWCR, SSWU, SSGL, SSGCR, SSGU, SSOWCR, and
SSOGCR saturation grid arrays for the saturation end-points, and the
SKRG, SKROG, SKROW and SKRW relative permeability grid cell arrays for
the relative permeability end-point data.

\subsubsection{Example}\label{example-382}

\begin{verbatim}
--
--       DEFINE GRID BLOCK END-POINT SSGCR DATA FOR ALL CELLS
--       (FOR NX x NY x NZ = 300)
--
SSGCR
         300*0.050                                                              /
\end{verbatim}

The above example defines a constant surfactant critical oil saturation
with respect to gas of 0.05 to all 300 cells in the model as defined by
the DIMENS keyword in the RUNSPEC section.

\subsection{SSGL -- End-Point Scaling Grid Cell Surfactant Connate Gas
Saturations}\label{ssgl-end-point-scaling-grid-cell-surfactant-connate-gas-saturations}

\begin{longtable}[]{@{}
  >{\raggedright\arraybackslash}p{(\linewidth - 14\tabcolsep) * \real{0.1250}}
  >{\raggedright\arraybackslash}p{(\linewidth - 14\tabcolsep) * \real{0.1250}}
  >{\raggedright\arraybackslash}p{(\linewidth - 14\tabcolsep) * \real{0.1250}}
  >{\raggedright\arraybackslash}p{(\linewidth - 14\tabcolsep) * \real{0.1250}}
  >{\raggedright\arraybackslash}p{(\linewidth - 14\tabcolsep) * \real{0.1250}}
  >{\raggedright\arraybackslash}p{(\linewidth - 14\tabcolsep) * \real{0.1250}}
  >{\raggedright\arraybackslash}p{(\linewidth - 14\tabcolsep) * \real{0.1250}}
  >{\raggedright\arraybackslash}p{(\linewidth - 14\tabcolsep) * \real{0.1250}}@{}}
\toprule\noalign{}
\begin{minipage}[b]{\linewidth}\raggedright
RUNSPEC
\end{minipage} & \begin{minipage}[b]{\linewidth}\raggedright
GRID
\end{minipage} & \begin{minipage}[b]{\linewidth}\raggedright
EDIT
\end{minipage} & \begin{minipage}[b]{\linewidth}\raggedright
PROPS
\end{minipage} & \begin{minipage}[b]{\linewidth}\raggedright
REGIONS
\end{minipage} & \begin{minipage}[b]{\linewidth}\raggedright
SOLUTION
\end{minipage} & \begin{minipage}[b]{\linewidth}\raggedright
SUMMARY
\end{minipage} & \begin{minipage}[b]{\linewidth}\raggedright
SCHEDULE
\end{minipage} \\
\midrule\noalign{}
\endhead
\bottomrule\noalign{}
\endlastfoot
\end{longtable}

\subsubsection{Description}\label{description-667}

SSGL defines the surfactant connate gas saturation for all the cells in
the model via an array when the end-point scaling option has been
invoked via the ENDSCALE keyword and the surfactant phase has been
activated by the SURFACT keyword in the RUNSPEC section. The connate gas
saturation is defined as the minimum gas saturation in a two-phase gas
relative permeability table. SSGL is used to scale the surfactant oil
and water relative permeability data.

This keyword is not supported by OPM Flow but would change the results
if supported so the simulation will be stopped.

\begin{longtable}[]{@{}
  >{\raggedright\arraybackslash}p{(\linewidth - 6\tabcolsep) * \real{0.2500}}
  >{\raggedright\arraybackslash}p{(\linewidth - 6\tabcolsep) * \real{0.2500}}
  >{\raggedright\arraybackslash}p{(\linewidth - 6\tabcolsep) * \real{0.2500}}
  >{\raggedright\arraybackslash}p{(\linewidth - 6\tabcolsep) * \real{0.2500}}@{}}
\toprule\noalign{}
\begin{minipage}[b]{\linewidth}\raggedright
No.
\end{minipage} & \begin{minipage}[b]{\linewidth}\raggedright
Name
\end{minipage} & \begin{minipage}[b]{\linewidth}\raggedright
Description
\end{minipage} & \begin{minipage}[b]{\linewidth}\raggedright
Default
\end{minipage} \\
\midrule\noalign{}
\endhead
\bottomrule\noalign{}
\endlastfoot
Field & Metric & Laboratory & \\
1 & SSGL & SSGL is an array of real numbers assigning the surfactant
connate gas saturation values to each cell in the model. The number of
entries should correspond to the NX x NY x NZ parameters on the DIMENS
keyword. Repeat counts may be used, for example 30*0.03 & Taken from
cell allocated relative permeability table. \\
dimensionless & dimensionless & dimensionless & \\
Notes: & & & \\
\end{longtable}

\emph{Table 8.174: SSGL Keyword Description}

End-point scaling allows the entered surfactant relative permeability
functions to be re-scaled based on the saturation end-points allocated
to each cell by the SSWL, SSWCR, SSWU, SSGL, SSGCR, SSGU, SSOWCR, and
SSOGCR saturation grid arrays for the saturation end-points, and the
SKRG, SKROG, SKROW and SKRW relative permeability grid cell arrays for
the relative permeability end-point data.

\subsubsection{Example}\label{example-383}

\begin{verbatim}
--
--       DEFINE GRID BLOCK END-POINT SSGL DATA FOR ALL CELLS
--       (FOR NX x NY x NZ = 300)
--
SSGL
         300*0.030                                                              /
\end{verbatim}

The above example defines a constant surfactant connate gas saturation
of 0.03 to all 300 cells in the model as defined by the DIMENS keyword
in the RUNSPEC section.

\subsection{SSOGCR -- End-Point Scaling Grid Cell Surfactant Critical
Oil Saturation with Respect to
Gas}\label{ssogcr-end-point-scaling-grid-cell-surfactant-critical-oil-saturation-with-respect-to-gas}

\begin{longtable}[]{@{}
  >{\raggedright\arraybackslash}p{(\linewidth - 14\tabcolsep) * \real{0.1250}}
  >{\raggedright\arraybackslash}p{(\linewidth - 14\tabcolsep) * \real{0.1250}}
  >{\raggedright\arraybackslash}p{(\linewidth - 14\tabcolsep) * \real{0.1250}}
  >{\raggedright\arraybackslash}p{(\linewidth - 14\tabcolsep) * \real{0.1250}}
  >{\raggedright\arraybackslash}p{(\linewidth - 14\tabcolsep) * \real{0.1250}}
  >{\raggedright\arraybackslash}p{(\linewidth - 14\tabcolsep) * \real{0.1250}}
  >{\raggedright\arraybackslash}p{(\linewidth - 14\tabcolsep) * \real{0.1250}}
  >{\raggedright\arraybackslash}p{(\linewidth - 14\tabcolsep) * \real{0.1250}}@{}}
\toprule\noalign{}
\begin{minipage}[b]{\linewidth}\raggedright
RUNSPEC
\end{minipage} & \begin{minipage}[b]{\linewidth}\raggedright
GRID
\end{minipage} & \begin{minipage}[b]{\linewidth}\raggedright
EDIT
\end{minipage} & \begin{minipage}[b]{\linewidth}\raggedright
PROPS
\end{minipage} & \begin{minipage}[b]{\linewidth}\raggedright
REGIONS
\end{minipage} & \begin{minipage}[b]{\linewidth}\raggedright
SOLUTION
\end{minipage} & \begin{minipage}[b]{\linewidth}\raggedright
SUMMARY
\end{minipage} & \begin{minipage}[b]{\linewidth}\raggedright
SCHEDULE
\end{minipage} \\
\midrule\noalign{}
\endhead
\bottomrule\noalign{}
\endlastfoot
\end{longtable}

\subsubsection{Description}\label{description-668}

SSOGCR defines the surfactant critical oil saturation for all the cells
in the model via an array when the end-point scaling option has been
invoked via the ENDSCALE keyword and the surfactant phase has been
activated by the SURFACT keyword in the RUNSPEC section. The critical
oil saturation with respect to gas is defined as the maximum oil
saturation for which the oil relative permeability is zero in a
two-phase gas-oil relative permeability table. SSOGCR scales the
surfactant oil relative permeability to gas data.

This keyword is not supported by OPM Flow but would change the results
if supported so the simulation will be stopped.

\begin{longtable}[]{@{}
  >{\raggedright\arraybackslash}p{(\linewidth - 6\tabcolsep) * \real{0.2500}}
  >{\raggedright\arraybackslash}p{(\linewidth - 6\tabcolsep) * \real{0.2500}}
  >{\raggedright\arraybackslash}p{(\linewidth - 6\tabcolsep) * \real{0.2500}}
  >{\raggedright\arraybackslash}p{(\linewidth - 6\tabcolsep) * \real{0.2500}}@{}}
\toprule\noalign{}
\begin{minipage}[b]{\linewidth}\raggedright
No.
\end{minipage} & \begin{minipage}[b]{\linewidth}\raggedright
Name
\end{minipage} & \begin{minipage}[b]{\linewidth}\raggedright
Description
\end{minipage} & \begin{minipage}[b]{\linewidth}\raggedright
Default
\end{minipage} \\
\midrule\noalign{}
\endhead
\bottomrule\noalign{}
\endlastfoot
Field & Metric & Laboratory & \\
1 & SSOGCR & SSOGCR is an array of real numbers assigning the surfactant
critical oil saturation with respect to gas values to each cell in the
model. The number of entries should correspond to the NX x NY x NZ
parameters on the DIMENS keyword. Repeat counts may be used, for example
30*0.30 dimensionless & Taken from cell allocated relative permeability
table. \\
dimensionless & dimensionless & dimensionless & \\
Notes: & & & \\
\end{longtable}

\emph{Table 8.175: SSOGCR Keyword Description}

End-point scaling allows the entered surfactant relative permeability
functions to be re-scaled based on the saturation end-points allocated
to each cell by the SSWL, SSWCR, SSWU, SSGL, SSGCR, SSGU, SSOWCR, and
SSOGCR saturation grid arrays for the saturation end-points, and the
SKRG, SKROG, SKROW and SKRW relative permeability grid cell arrays for
the relative permeability end-point data.

\subsubsection{Example}\label{example-384}

\begin{verbatim}
--
--       DEFINE GRID BLOCK END-POINT SSOGCR DATA FOR ALL CELLS
--       (FOR NX x NY x NZ = 300)
--
SSOGCR
         300*0.200                                                              /
\end{verbatim}

The above example defines a surfactant constant critical oil saturation
with respect to gas of 0.20 to all 300 cells in the model as defined by
the DIMENS keyword in the RUNSPEC section.

\subsection{SSOWCR -- End-Point Scaling Grid Cell Surfactant Critical
Oil Saturation with Respect to
Water}\label{ssowcr-end-point-scaling-grid-cell-surfactant-critical-oil-saturation-with-respect-to-water}

\begin{longtable}[]{@{}
  >{\raggedright\arraybackslash}p{(\linewidth - 14\tabcolsep) * \real{0.1250}}
  >{\raggedright\arraybackslash}p{(\linewidth - 14\tabcolsep) * \real{0.1250}}
  >{\raggedright\arraybackslash}p{(\linewidth - 14\tabcolsep) * \real{0.1250}}
  >{\raggedright\arraybackslash}p{(\linewidth - 14\tabcolsep) * \real{0.1250}}
  >{\raggedright\arraybackslash}p{(\linewidth - 14\tabcolsep) * \real{0.1250}}
  >{\raggedright\arraybackslash}p{(\linewidth - 14\tabcolsep) * \real{0.1250}}
  >{\raggedright\arraybackslash}p{(\linewidth - 14\tabcolsep) * \real{0.1250}}
  >{\raggedright\arraybackslash}p{(\linewidth - 14\tabcolsep) * \real{0.1250}}@{}}
\toprule\noalign{}
\begin{minipage}[b]{\linewidth}\raggedright
RUNSPEC
\end{minipage} & \begin{minipage}[b]{\linewidth}\raggedright
GRID
\end{minipage} & \begin{minipage}[b]{\linewidth}\raggedright
EDIT
\end{minipage} & \begin{minipage}[b]{\linewidth}\raggedright
PROPS
\end{minipage} & \begin{minipage}[b]{\linewidth}\raggedright
REGIONS
\end{minipage} & \begin{minipage}[b]{\linewidth}\raggedright
SOLUTION
\end{minipage} & \begin{minipage}[b]{\linewidth}\raggedright
SUMMARY
\end{minipage} & \begin{minipage}[b]{\linewidth}\raggedright
SCHEDULE
\end{minipage} \\
\midrule\noalign{}
\endhead
\bottomrule\noalign{}
\endlastfoot
\end{longtable}

\subsubsection{Description}\label{description-669}

SSOWCR defines the surfactant critical oil saturation for all the cells
in the model via an array when the end-point scaling option has been
invoked via the ENDSCALE keyword and the surfactant phase has been
activated by the SURFACT keyword in the RUNSPEC section. The critical
oil saturation with respect to water is defined as the maximum oil
saturation for which the oil relative permeability is zero in a
two-phase oil-water relative permeability table. SSOWCR scales the
surfactant oil relative permeability to water data.

This keyword is not supported by OPM Flow but would change the results
if supported so the simulation will be stopped.

\begin{longtable}[]{@{}
  >{\raggedright\arraybackslash}p{(\linewidth - 6\tabcolsep) * \real{0.2500}}
  >{\raggedright\arraybackslash}p{(\linewidth - 6\tabcolsep) * \real{0.2500}}
  >{\raggedright\arraybackslash}p{(\linewidth - 6\tabcolsep) * \real{0.2500}}
  >{\raggedright\arraybackslash}p{(\linewidth - 6\tabcolsep) * \real{0.2500}}@{}}
\toprule\noalign{}
\begin{minipage}[b]{\linewidth}\raggedright
No.
\end{minipage} & \begin{minipage}[b]{\linewidth}\raggedright
Name
\end{minipage} & \begin{minipage}[b]{\linewidth}\raggedright
Description
\end{minipage} & \begin{minipage}[b]{\linewidth}\raggedright
Default
\end{minipage} \\
\midrule\noalign{}
\endhead
\bottomrule\noalign{}
\endlastfoot
Field & Metric & Laboratory & \\
1 & SSOWCR & SSOWCR is an array of real numbers assigning the surfactant
critical oil saturation with respect to water values to each cell in the
model. The number of entries should correspond to the NX x NY x NZ
parameters on the DIMENS keyword. Repeat counts may be used, for example
30*0.30 & Taken from cell allocated relative permeability table. \\
dimensionless & dimensionless & dimensionless & \\
Notes: & & & \\
\end{longtable}

\emph{Table 8.176: SSOWCR Keyword Description}

End-point scaling allows the entered surfactant relative permeability
functions to be re-scaled based on the saturation end-points allocated
to each cell by the SSWL, SSWCR, SSWU, SSGL, SSGCR, SSGU, SSOWCR, and
SSOGCR saturation grid arrays for the saturation end-points, and the
SKRG, SKROG, SKROW and SKRW relative permeability grid cell arrays for
the relative permeability end-point data.

\subsubsection{Example}\label{example-385}

\begin{verbatim}
--
--       DEFINE GRID BLOCK END-POINT SSOWCR DATA FOR ALL CELLS
--       (FOR NX x NY x NZ = 300)
--
SSOWCR
         300*0.200                                                              /
\end{verbatim}

The above example defines a constant critical oil saturation of 0.20 to
all 300 cells in the model as defined by the DIMENS keyword in the
RUNSPEC section.

\subsection{SSWCR -- End-Point Scaling Grid Cell Critical Water
Saturation}\label{sswcr-end-point-scaling-grid-cell-critical-water-saturation}

\begin{longtable}[]{@{}
  >{\raggedright\arraybackslash}p{(\linewidth - 14\tabcolsep) * \real{0.1250}}
  >{\raggedright\arraybackslash}p{(\linewidth - 14\tabcolsep) * \real{0.1250}}
  >{\raggedright\arraybackslash}p{(\linewidth - 14\tabcolsep) * \real{0.1250}}
  >{\raggedright\arraybackslash}p{(\linewidth - 14\tabcolsep) * \real{0.1250}}
  >{\raggedright\arraybackslash}p{(\linewidth - 14\tabcolsep) * \real{0.1250}}
  >{\raggedright\arraybackslash}p{(\linewidth - 14\tabcolsep) * \real{0.1250}}
  >{\raggedright\arraybackslash}p{(\linewidth - 14\tabcolsep) * \real{0.1250}}
  >{\raggedright\arraybackslash}p{(\linewidth - 14\tabcolsep) * \real{0.1250}}@{}}
\toprule\noalign{}
\begin{minipage}[b]{\linewidth}\raggedright
RUNSPEC
\end{minipage} & \begin{minipage}[b]{\linewidth}\raggedright
GRID
\end{minipage} & \begin{minipage}[b]{\linewidth}\raggedright
EDIT
\end{minipage} & \begin{minipage}[b]{\linewidth}\raggedright
PROPS
\end{minipage} & \begin{minipage}[b]{\linewidth}\raggedright
REGIONS
\end{minipage} & \begin{minipage}[b]{\linewidth}\raggedright
SOLUTION
\end{minipage} & \begin{minipage}[b]{\linewidth}\raggedright
SUMMARY
\end{minipage} & \begin{minipage}[b]{\linewidth}\raggedright
SCHEDULE
\end{minipage} \\
\midrule\noalign{}
\endhead
\bottomrule\noalign{}
\endlastfoot
\end{longtable}

\subsubsection{Description}\label{description-670}

SSWCR defines the surfactant critical water saturation for all the cells
in the model via an array when the end-point scaling option has been
invoked via the ENDSCALE keyword and the surfactant phase has been
activated by the SURFACT keyword in the RUNSPEC section. The critical
water saturation is defined as the maximum water saturation for which
the water relative permeability is zero in a two-phase relative
permeability table. SSWCR scales the surfactant water relative
permeability data.

This keyword is not supported by OPM Flow but would change the results
if supported so the simulation will be stopped.

\begin{longtable}[]{@{}
  >{\raggedright\arraybackslash}p{(\linewidth - 6\tabcolsep) * \real{0.2500}}
  >{\raggedright\arraybackslash}p{(\linewidth - 6\tabcolsep) * \real{0.2500}}
  >{\raggedright\arraybackslash}p{(\linewidth - 6\tabcolsep) * \real{0.2500}}
  >{\raggedright\arraybackslash}p{(\linewidth - 6\tabcolsep) * \real{0.2500}}@{}}
\toprule\noalign{}
\begin{minipage}[b]{\linewidth}\raggedright
No.
\end{minipage} & \begin{minipage}[b]{\linewidth}\raggedright
Name
\end{minipage} & \begin{minipage}[b]{\linewidth}\raggedright
Description
\end{minipage} & \begin{minipage}[b]{\linewidth}\raggedright
Default
\end{minipage} \\
\midrule\noalign{}
\endhead
\bottomrule\noalign{}
\endlastfoot
Field & Metric & Laboratory & \\
1 & SSWCR & SSWCR is an array of real numbers assigning the surfactant
critical water saturation values to each cell in the model. The number
of entries should correspond to the NX x NY x NZ parameters on the
DIMENS keyword. Repeat counts may be used, for example 30*0.20 & Taken
from cell allocated relative permeability table. \\
dimensionless & dimensionless & dimensionless & \\
Notes: & & & \\
\end{longtable}

\emph{Table 8.177: SSWCR Keyword Description}

End-point scaling allows the entered surfactant relative permeability
functions to be re-scaled based on the saturation end-points allocated
to each cell by the SSWL, SSWCR, SSWU, SSGL, SSGCR, SSGU, SSOWCR, and
SSOGCR saturation grid arrays for the saturation end-points, and the
SKRG, SKROG, SKROW and SKRW relative permeability grid cell arrays for
the relative permeability end-point data.

\subsubsection{Example}\label{example-386}

\begin{verbatim}
--
--       DEFINE GRID BLOCK END-POINT SSWCR DATA FOR ALL CELLS
--       (FOR NX x NY x NZ = 300)
--
SSWCR
         300*0.200                                                              /
\end{verbatim}

The above example defines a constant surfactant critical water
saturation of 0.20 to all 300 cells in the model as defined by the
DIMENS keyword in the RUNSPEC section.

\subsection{SSWL -- End-Point Scaling Grid Cell Surfactant Connate Water
Saturation}\label{sswl-end-point-scaling-grid-cell-surfactant-connate-water-saturation}

\begin{longtable}[]{@{}
  >{\raggedright\arraybackslash}p{(\linewidth - 14\tabcolsep) * \real{0.1250}}
  >{\raggedright\arraybackslash}p{(\linewidth - 14\tabcolsep) * \real{0.1250}}
  >{\raggedright\arraybackslash}p{(\linewidth - 14\tabcolsep) * \real{0.1250}}
  >{\raggedright\arraybackslash}p{(\linewidth - 14\tabcolsep) * \real{0.1250}}
  >{\raggedright\arraybackslash}p{(\linewidth - 14\tabcolsep) * \real{0.1250}}
  >{\raggedright\arraybackslash}p{(\linewidth - 14\tabcolsep) * \real{0.1250}}
  >{\raggedright\arraybackslash}p{(\linewidth - 14\tabcolsep) * \real{0.1250}}
  >{\raggedright\arraybackslash}p{(\linewidth - 14\tabcolsep) * \real{0.1250}}@{}}
\toprule\noalign{}
\begin{minipage}[b]{\linewidth}\raggedright
RUNSPEC
\end{minipage} & \begin{minipage}[b]{\linewidth}\raggedright
GRID
\end{minipage} & \begin{minipage}[b]{\linewidth}\raggedright
EDIT
\end{minipage} & \begin{minipage}[b]{\linewidth}\raggedright
PROPS
\end{minipage} & \begin{minipage}[b]{\linewidth}\raggedright
REGIONS
\end{minipage} & \begin{minipage}[b]{\linewidth}\raggedright
SOLUTION
\end{minipage} & \begin{minipage}[b]{\linewidth}\raggedright
SUMMARY
\end{minipage} & \begin{minipage}[b]{\linewidth}\raggedright
SCHEDULE
\end{minipage} \\
\midrule\noalign{}
\endhead
\bottomrule\noalign{}
\endlastfoot
\end{longtable}

\subsubsection{Description}\label{description-671}

SSWL defines the surfactant connate water saturation for all the cells
in the model via an array when the end-point scaling option has been
invoked via the ENDSCALE keyword and the surfactant phase has been
activated by the SURFACT keyword in the RUNSPEC section. The connate
water saturation is defined as the minimum water saturation in a
two-phase water relative permeability table. SSWL scales the surfactant
oil relative permeability to water and gas data.

This keyword is not supported by OPM Flow but would change the results
if supported so the simulation will be stopped.

\begin{longtable}[]{@{}
  >{\raggedright\arraybackslash}p{(\linewidth - 6\tabcolsep) * \real{0.2500}}
  >{\raggedright\arraybackslash}p{(\linewidth - 6\tabcolsep) * \real{0.2500}}
  >{\raggedright\arraybackslash}p{(\linewidth - 6\tabcolsep) * \real{0.2500}}
  >{\raggedright\arraybackslash}p{(\linewidth - 6\tabcolsep) * \real{0.2500}}@{}}
\toprule\noalign{}
\begin{minipage}[b]{\linewidth}\raggedright
No.
\end{minipage} & \begin{minipage}[b]{\linewidth}\raggedright
Name
\end{minipage} & \begin{minipage}[b]{\linewidth}\raggedright
Description
\end{minipage} & \begin{minipage}[b]{\linewidth}\raggedright
Default
\end{minipage} \\
\midrule\noalign{}
\endhead
\bottomrule\noalign{}
\endlastfoot
Field & Metric & Laboratory & \\
1 & SSWL & SSWL is an array of real numbers assigning the connate water
saturation values to each cell in the model. The number of entries
should correspond to the NX x NY x NZ parameters on the DIMENS keyword.
Repeat counts may be used, for example 30*0.15 & Taken from cell
allocated relative permeability table. \\
dimensionless & dimensionless & dimensionless & \\
Notes: & & & \\
\end{longtable}

\emph{Table 8.178: SSWL Keyword Description}

End-point scaling allows the entered surfactant relative permeability
functions to be re-scaled based on the saturation end-points allocated
to each cell by the SSWL, SSWCR, SSWU, SSGL, SSGCR, SSGU, SSOWCR, and
SSOGCR saturation grid arrays for the saturation end-points, and the
SKRG, SKROG, SKROW and SKRW relative permeability grid cell arrays for
the relative permeability end-point data.

\subsubsection{Example}\label{example-387}

\begin{verbatim}
--
--       DEFINE GRID BLOCK END-POINT SSWL DATA FOR ALL CELLS
--       (FOR NX x NY x NZ = 300)
--
SSWL
         300*0.150                                                              /
\end{verbatim}

The above example defines a constant surfactant connate water saturation
of 0.15 to all 300 cells in the model as defined by the DIMENS keyword
in the RUNSPEC section.

\subsection{SSWU -- End-Point Scaling Grid Cell Surfactant Maximum Water
Saturation}\label{sswu-end-point-scaling-grid-cell-surfactant-maximum-water-saturation}

\begin{longtable}[]{@{}
  >{\raggedright\arraybackslash}p{(\linewidth - 14\tabcolsep) * \real{0.1250}}
  >{\raggedright\arraybackslash}p{(\linewidth - 14\tabcolsep) * \real{0.1250}}
  >{\raggedright\arraybackslash}p{(\linewidth - 14\tabcolsep) * \real{0.1250}}
  >{\raggedright\arraybackslash}p{(\linewidth - 14\tabcolsep) * \real{0.1250}}
  >{\raggedright\arraybackslash}p{(\linewidth - 14\tabcolsep) * \real{0.1250}}
  >{\raggedright\arraybackslash}p{(\linewidth - 14\tabcolsep) * \real{0.1250}}
  >{\raggedright\arraybackslash}p{(\linewidth - 14\tabcolsep) * \real{0.1250}}
  >{\raggedright\arraybackslash}p{(\linewidth - 14\tabcolsep) * \real{0.1250}}@{}}
\toprule\noalign{}
\begin{minipage}[b]{\linewidth}\raggedright
RUNSPEC
\end{minipage} & \begin{minipage}[b]{\linewidth}\raggedright
GRID
\end{minipage} & \begin{minipage}[b]{\linewidth}\raggedright
EDIT
\end{minipage} & \begin{minipage}[b]{\linewidth}\raggedright
PROPS
\end{minipage} & \begin{minipage}[b]{\linewidth}\raggedright
REGIONS
\end{minipage} & \begin{minipage}[b]{\linewidth}\raggedright
SOLUTION
\end{minipage} & \begin{minipage}[b]{\linewidth}\raggedright
SUMMARY
\end{minipage} & \begin{minipage}[b]{\linewidth}\raggedright
SCHEDULE
\end{minipage} \\
\midrule\noalign{}
\endhead
\bottomrule\noalign{}
\endlastfoot
\end{longtable}

\subsubsection{Description}\label{description-672}

SSWU defines the surfactant maximum water saturation for all the cells
in the model via an array when the end-point scaling option has been
invoked via the ENDSCALE keyword and the surfactant phase has been
activated by the SURFACT keyword in the RUNSPEC section. The maximum
water saturation is defined as the maximum water saturation in a
two-phase water relative permeability table. SSWU scales the surfactant
water relative permeability data.

This keyword is not supported by OPM Flow but would change the results
if supported so the simulation will be stopped.

\begin{longtable}[]{@{}
  >{\raggedright\arraybackslash}p{(\linewidth - 6\tabcolsep) * \real{0.2500}}
  >{\raggedright\arraybackslash}p{(\linewidth - 6\tabcolsep) * \real{0.2500}}
  >{\raggedright\arraybackslash}p{(\linewidth - 6\tabcolsep) * \real{0.2500}}
  >{\raggedright\arraybackslash}p{(\linewidth - 6\tabcolsep) * \real{0.2500}}@{}}
\toprule\noalign{}
\begin{minipage}[b]{\linewidth}\raggedright
No.
\end{minipage} & \begin{minipage}[b]{\linewidth}\raggedright
Name
\end{minipage} & \begin{minipage}[b]{\linewidth}\raggedright
Description
\end{minipage} & \begin{minipage}[b]{\linewidth}\raggedright
Default
\end{minipage} \\
\midrule\noalign{}
\endhead
\bottomrule\noalign{}
\endlastfoot
Field & Metric & Laboratory & \\
1 & SSWU & SSWU is an array of real numbers assigning the surfactant
maximum water saturation values to each cell in the model. The number of
entries should correspond to the NX x NY x NZ parameters on the DIMENS
keyword. Repeat counts may be used, for example 30*0.70 & Taken from
cell allocated relative permeability table. \\
dimensionless & dimensionless & dimensionless & \\
Notes: & & & \\
\end{longtable}

\emph{Table 8.179: SSWU Keyword Description}

End-point scaling allows the entered surfactant relative permeability
functions to be re-scaled based on the saturation end-points allocated
to each cell by the SSWL, SSWCR, SSWU, SSGL, SSGCR, SSGU, SSOWCR, and
SSOGCR saturation grid arrays for the saturation end-points, and the
SKRG, SKROG, SKROW and SKRW relative permeability grid cell arrays for
the relative permeability end-point data.

\subsubsection{Example}\label{example-388}

\begin{verbatim}
--
--       DEFINE GRID BLOCK END-POINT SSWU DATA FOR ALL CELLS
--       (FOR NX x NY x NZ = 300)
--
SSWU
         300*0.700                                                              /
\end{verbatim}

The above example defines a constant surfactant connate gas saturation
of 0.70 to all 300 cells in the model as defined by the DIMENS keyword
in the RUNSPEC section.

\subsection{STOG -- Define Capillary Pressure Oil-Gas Surface Tension
versus
Pressure}\label{stog-define-capillary-pressure-oil-gas-surface-tension-versus-pressure}

\begin{longtable}[]{@{}
  >{\raggedright\arraybackslash}p{(\linewidth - 14\tabcolsep) * \real{0.1250}}
  >{\raggedright\arraybackslash}p{(\linewidth - 14\tabcolsep) * \real{0.1250}}
  >{\raggedright\arraybackslash}p{(\linewidth - 14\tabcolsep) * \real{0.1250}}
  >{\raggedright\arraybackslash}p{(\linewidth - 14\tabcolsep) * \real{0.1250}}
  >{\raggedright\arraybackslash}p{(\linewidth - 14\tabcolsep) * \real{0.1250}}
  >{\raggedright\arraybackslash}p{(\linewidth - 14\tabcolsep) * \real{0.1250}}
  >{\raggedright\arraybackslash}p{(\linewidth - 14\tabcolsep) * \real{0.1250}}
  >{\raggedright\arraybackslash}p{(\linewidth - 14\tabcolsep) * \real{0.1250}}@{}}
\toprule\noalign{}
\begin{minipage}[b]{\linewidth}\raggedright
RUNSPEC
\end{minipage} & \begin{minipage}[b]{\linewidth}\raggedright
GRID
\end{minipage} & \begin{minipage}[b]{\linewidth}\raggedright
EDIT
\end{minipage} & \begin{minipage}[b]{\linewidth}\raggedright
PROPS
\end{minipage} & \begin{minipage}[b]{\linewidth}\raggedright
REGIONS
\end{minipage} & \begin{minipage}[b]{\linewidth}\raggedright
SOLUTION
\end{minipage} & \begin{minipage}[b]{\linewidth}\raggedright
SUMMARY
\end{minipage} & \begin{minipage}[b]{\linewidth}\raggedright
SCHEDULE
\end{minipage} \\
\midrule\noalign{}
\endhead
\bottomrule\noalign{}
\endlastfoot
\end{longtable}

\subsubsection{Description}\label{description-673}

The STOG keyword defines capillary pressure oil-gas surface tension
versus pressure tables used in adjusting the pressure independent
capillary pressure vectors in the SGFN, SGOF or SLGOF saturation tables,
entered by their respective keywords in the PROPS section. The SATOPTS
keyword in the RUNSPEC section should state the SURFTENS character
string to activate the Capillary Pressure Surface Tension Pressure
Dependency option. If the STOG keyword is not supplied then no capillary
pressure surface tension pressure scaling will occur and the capillary
pressure values on the SGFN, SGOF or SLGOF saturation tables will be
used directly.

Capillary pressure surface tension pressure scaling can also be used
with multi-segment wells, but this facility has not been incorporated in
OPM Flow's multi-segment well implementation.

This keyword is not supported by OPM Flow but would change the results
if supported so the simulation will be stopped.

\subsection{STONE -- Activate Stone's Second Three Phase Oil Relative
Permeability Model (Alias for
STONE2)}\label{stone-activate-stones-second-three-phase-oil-relative-permeability-model-alias-for-stone2}

\begin{longtable}[]{@{}
  >{\raggedright\arraybackslash}p{(\linewidth - 14\tabcolsep) * \real{0.1250}}
  >{\raggedright\arraybackslash}p{(\linewidth - 14\tabcolsep) * \real{0.1250}}
  >{\raggedright\arraybackslash}p{(\linewidth - 14\tabcolsep) * \real{0.1250}}
  >{\raggedright\arraybackslash}p{(\linewidth - 14\tabcolsep) * \real{0.1250}}
  >{\raggedright\arraybackslash}p{(\linewidth - 14\tabcolsep) * \real{0.1250}}
  >{\raggedright\arraybackslash}p{(\linewidth - 14\tabcolsep) * \real{0.1250}}
  >{\raggedright\arraybackslash}p{(\linewidth - 14\tabcolsep) * \real{0.1250}}
  >{\raggedright\arraybackslash}p{(\linewidth - 14\tabcolsep) * \real{0.1250}}@{}}
\toprule\noalign{}
\begin{minipage}[b]{\linewidth}\raggedright
RUNSPEC
\end{minipage} & \begin{minipage}[b]{\linewidth}\raggedright
GRID
\end{minipage} & \begin{minipage}[b]{\linewidth}\raggedright
EDIT
\end{minipage} & \begin{minipage}[b]{\linewidth}\raggedright
PROPS
\end{minipage} & \begin{minipage}[b]{\linewidth}\raggedright
REGIONS
\end{minipage} & \begin{minipage}[b]{\linewidth}\raggedright
SOLUTION
\end{minipage} & \begin{minipage}[b]{\linewidth}\raggedright
SUMMARY
\end{minipage} & \begin{minipage}[b]{\linewidth}\raggedright
SCHEDULE
\end{minipage} \\
\midrule\noalign{}
\endhead
\bottomrule\noalign{}
\endlastfoot
\end{longtable}

\subsubsection{Description}\label{description-674}

This keyword is an alias for STONE2 keyword that activates Stone's
{[}Stone, H. L. ``Estimation of Three-Phase Relative Permeability and
Residual Oil Data,'' Journal of Canadian Petroleum Technology (1973) 12,
No.~4, 53-61.{]} second three phase oil relative permeability model as
modified by Aziz and Settari {[}Aziz, K. and Settari, A. Petroleum
Reservoir Simulation, London, UK, Applied Science Publishers (1979),
page 398.{]}. If the STONE, STONE1 and STONE2 keywords are not present
in the input deck then the default three phase oil relative permeability
model is employed. The STONE keyword should only be used in three phase
runs containing the oil, gas and water phases.

There is no data required for this keyword.

\subsubsection{Example}\label{example-389}

\begin{verbatim}
--
--       ACTIVATE STONE’S SECOND THREE PHASE RELATIVE PERMEABILITY MODEL
--
STONE
\end{verbatim}

The above example switches on the Modified Stone three phase relative
permeability model.

\subsection{STONE1 -- Activate Stone's First Three Phase Oil Relative
Permeability
Model}\label{stone1-activate-stones-first-three-phase-oil-relative-permeability-model}

\begin{longtable}[]{@{}
  >{\raggedright\arraybackslash}p{(\linewidth - 14\tabcolsep) * \real{0.1250}}
  >{\raggedright\arraybackslash}p{(\linewidth - 14\tabcolsep) * \real{0.1250}}
  >{\raggedright\arraybackslash}p{(\linewidth - 14\tabcolsep) * \real{0.1250}}
  >{\raggedright\arraybackslash}p{(\linewidth - 14\tabcolsep) * \real{0.1250}}
  >{\raggedright\arraybackslash}p{(\linewidth - 14\tabcolsep) * \real{0.1250}}
  >{\raggedright\arraybackslash}p{(\linewidth - 14\tabcolsep) * \real{0.1250}}
  >{\raggedright\arraybackslash}p{(\linewidth - 14\tabcolsep) * \real{0.1250}}
  >{\raggedright\arraybackslash}p{(\linewidth - 14\tabcolsep) * \real{0.1250}}@{}}
\toprule\noalign{}
\begin{minipage}[b]{\linewidth}\raggedright
RUNSPEC
\end{minipage} & \begin{minipage}[b]{\linewidth}\raggedright
GRID
\end{minipage} & \begin{minipage}[b]{\linewidth}\raggedright
EDIT
\end{minipage} & \begin{minipage}[b]{\linewidth}\raggedright
PROPS
\end{minipage} & \begin{minipage}[b]{\linewidth}\raggedright
REGIONS
\end{minipage} & \begin{minipage}[b]{\linewidth}\raggedright
SOLUTION
\end{minipage} & \begin{minipage}[b]{\linewidth}\raggedright
SUMMARY
\end{minipage} & \begin{minipage}[b]{\linewidth}\raggedright
SCHEDULE
\end{minipage} \\
\midrule\noalign{}
\endhead
\bottomrule\noalign{}
\endlastfoot
\end{longtable}

\subsubsection{Description}\label{description-675}

This keyword activates Stone's {[}Stone, H. L. ``Probability Model for
Estimating Three-Phase Relative Permeability,'' paper SPE 2116, Journal
of Canadian Petroleum Technology (1973) 22, No.~2, 214-218.{]} first
three phase oil relative permeability model as modified by Aziz and
Settari {[}Aziz, K. and Settari, A. Petroleum Reservoir Simulation,
London, UK, Applied Science Publishers (1979), page 398.{]}. If the
STONE1 and STONE2 keywords are not present in the input deck then the
default three phase oil relative permeability model is employed. The
STONE1 keyword should only be used in three phase runs containing the
oil, gas and water phases.

There is no data required for this keyword.

\subsubsection{Example}\label{example-390}

\begin{verbatim}
--
--       ACTIVATE STONE’S FIRST THREE PHASE RELATIVE PERMEABILITY MODEL
--
STONE1
\end{verbatim}

The above example switches on the Modified Stone three phase relative
permeability model.

\subsection{STONE1EX -- Define Stone's First Three Phase Oil Relative
Permeability
Parameter}\label{stone1ex-define-stones-first-three-phase-oil-relative-permeability-parameter}

\begin{longtable}[]{@{}
  >{\raggedright\arraybackslash}p{(\linewidth - 14\tabcolsep) * \real{0.1250}}
  >{\raggedright\arraybackslash}p{(\linewidth - 14\tabcolsep) * \real{0.1250}}
  >{\raggedright\arraybackslash}p{(\linewidth - 14\tabcolsep) * \real{0.1250}}
  >{\raggedright\arraybackslash}p{(\linewidth - 14\tabcolsep) * \real{0.1250}}
  >{\raggedright\arraybackslash}p{(\linewidth - 14\tabcolsep) * \real{0.1250}}
  >{\raggedright\arraybackslash}p{(\linewidth - 14\tabcolsep) * \real{0.1250}}
  >{\raggedright\arraybackslash}p{(\linewidth - 14\tabcolsep) * \real{0.1250}}
  >{\raggedright\arraybackslash}p{(\linewidth - 14\tabcolsep) * \real{0.1250}}@{}}
\toprule\noalign{}
\begin{minipage}[b]{\linewidth}\raggedright
RUNSPEC
\end{minipage} & \begin{minipage}[b]{\linewidth}\raggedright
GRID
\end{minipage} & \begin{minipage}[b]{\linewidth}\raggedright
EDIT
\end{minipage} & \begin{minipage}[b]{\linewidth}\raggedright
PROPS
\end{minipage} & \begin{minipage}[b]{\linewidth}\raggedright
REGIONS
\end{minipage} & \begin{minipage}[b]{\linewidth}\raggedright
SOLUTION
\end{minipage} & \begin{minipage}[b]{\linewidth}\raggedright
SUMMARY
\end{minipage} & \begin{minipage}[b]{\linewidth}\raggedright
SCHEDULE
\end{minipage} \\
\midrule\noalign{}
\endhead
\bottomrule\noalign{}
\endlastfoot
\end{longtable}

\subsubsection{Description}\label{description-676}

This keyword defines the exponent used in Stone's {[}Stone, H. L.
``Probability Model for Estimating Three-Phase Relative Permeability,''
paper SPE 2116, Journal of Canadian Petroleum Technology (1973) 22,
No.~2, 214-218.{]} first three phase oil relative permeability model as
modified by Aziz and Settari {[}Aziz, K. and Settari, A. Petroleum
Reservoir Simulation, London, UK, Applied Science Publishers (1979),
page 398.{]}. The STONE1EX keyword should only be used in three phase
runs containing the oil, gas and water phases and when the STONE1
keyword in the PROPS section has been used to activate Stone's first
three phase oil relative permeability model.

\begin{longtable}[]{@{}
  >{\raggedright\arraybackslash}p{(\linewidth - 6\tabcolsep) * \real{0.2500}}
  >{\raggedright\arraybackslash}p{(\linewidth - 6\tabcolsep) * \real{0.2500}}
  >{\raggedright\arraybackslash}p{(\linewidth - 6\tabcolsep) * \real{0.2500}}
  >{\raggedright\arraybackslash}p{(\linewidth - 6\tabcolsep) * \real{0.2500}}@{}}
\toprule\noalign{}
\begin{minipage}[b]{\linewidth}\raggedright
No.
\end{minipage} & \begin{minipage}[b]{\linewidth}\raggedright
Name
\end{minipage} & \begin{minipage}[b]{\linewidth}\raggedright
Description
\end{minipage} & \begin{minipage}[b]{\linewidth}\raggedright
Default
\end{minipage} \\
\midrule\noalign{}
\endhead
\bottomrule\noalign{}
\endlastfoot
Field & Metric & Laboratory & \\
1 & STONEPAR1 & A real positive value that defines the exponent to be
used in the Modified Stone first three phase oil relative permeability
model. & None \\
dimensionless & dimensionless & dimensionless & \\
Notes: & & & \\
\end{longtable}

\emph{Table 8.180: STONE1EX Keyword Description}

If the STONE, STONE1 and STONE2 keywords are not present in the input
deck then the default three phase oil relative permeability model is
employed.

\subsubsection{Example}\label{example-391}

Given NTSFUN equals five on the TABDIMS keyword in the RUNSPEC section,
then:

\begin{verbatim}
--
--       DEFINE STONE'S FIRST THREE PHASE RELATIVE PERMEABILITY MODEL PARAMETER
--
STONE1EX
         1.000                                        / SATURATION TABLE NO. 01
         1.000                                        / SATURATION TABLE NO. 02
         2.000                                        / SATURATION TABLE NO. 03
         1.000                                        / SATURATION TABLE NO. 04
         3.000                                        / SATURATION TABLE NO. 05
\end{verbatim}

Defines the exponents to be used in the Modified Stone first three phase
oil relative permeability model, for each of the five saturation tables.

\subsection{STONE2 -- Activate Stone's Second Three Phase Oil Relative
Permeability
Model}\label{stone2-activate-stones-second-three-phase-oil-relative-permeability-model}

\begin{longtable}[]{@{}
  >{\raggedright\arraybackslash}p{(\linewidth - 14\tabcolsep) * \real{0.1250}}
  >{\raggedright\arraybackslash}p{(\linewidth - 14\tabcolsep) * \real{0.1250}}
  >{\raggedright\arraybackslash}p{(\linewidth - 14\tabcolsep) * \real{0.1250}}
  >{\raggedright\arraybackslash}p{(\linewidth - 14\tabcolsep) * \real{0.1250}}
  >{\raggedright\arraybackslash}p{(\linewidth - 14\tabcolsep) * \real{0.1250}}
  >{\raggedright\arraybackslash}p{(\linewidth - 14\tabcolsep) * \real{0.1250}}
  >{\raggedright\arraybackslash}p{(\linewidth - 14\tabcolsep) * \real{0.1250}}
  >{\raggedright\arraybackslash}p{(\linewidth - 14\tabcolsep) * \real{0.1250}}@{}}
\toprule\noalign{}
\begin{minipage}[b]{\linewidth}\raggedright
RUNSPEC
\end{minipage} & \begin{minipage}[b]{\linewidth}\raggedright
GRID
\end{minipage} & \begin{minipage}[b]{\linewidth}\raggedright
EDIT
\end{minipage} & \begin{minipage}[b]{\linewidth}\raggedright
PROPS
\end{minipage} & \begin{minipage}[b]{\linewidth}\raggedright
REGIONS
\end{minipage} & \begin{minipage}[b]{\linewidth}\raggedright
SOLUTION
\end{minipage} & \begin{minipage}[b]{\linewidth}\raggedright
SUMMARY
\end{minipage} & \begin{minipage}[b]{\linewidth}\raggedright
SCHEDULE
\end{minipage} \\
\midrule\noalign{}
\endhead
\bottomrule\noalign{}
\endlastfoot
\end{longtable}

\subsubsection{Description}\label{description-677}

This keyword activates Stone's {[}Stone, H. L. ``Estimation of
Three-Phase Relative Permeability and Residual Oil Data,'' Journal of
Canadian Petroleum Technology (1973) 12, No.~4, 53-61.{]} second three
phase oil relative permeability model as modified by Aziz and Settari
{[}Aziz, K. and Settari, A. Petroleum Reservoir Simulation, London, UK,
Applied Science Publishers (1979), page 398.{]}. If the STONE, STONE1
and STONE2 keywords are not present in the input deck then the default
three phase oil relative permeability model is employed. The STONE2
keyword should only be used in three phase runs containing the oil, gas
and water phases.

There is no data required for this keyword.

\subsubsection{Example}\label{example-392}

\begin{verbatim}
--
--       ACTIVATE STONE’S SECOND THREE PHASE RELATIVE PERMEABILITY MODEL
--
STONE2
\end{verbatim}

The above example switches on the Modified Stone three phase relative
permeability model.

\subsection{STOW -- Define Capillary Pressure Oil-Water Surface Tension
versus
Pressure}\label{stow-define-capillary-pressure-oil-water-surface-tension-versus-pressure}

\begin{longtable}[]{@{}
  >{\raggedright\arraybackslash}p{(\linewidth - 14\tabcolsep) * \real{0.1250}}
  >{\raggedright\arraybackslash}p{(\linewidth - 14\tabcolsep) * \real{0.1250}}
  >{\raggedright\arraybackslash}p{(\linewidth - 14\tabcolsep) * \real{0.1250}}
  >{\raggedright\arraybackslash}p{(\linewidth - 14\tabcolsep) * \real{0.1250}}
  >{\raggedright\arraybackslash}p{(\linewidth - 14\tabcolsep) * \real{0.1250}}
  >{\raggedright\arraybackslash}p{(\linewidth - 14\tabcolsep) * \real{0.1250}}
  >{\raggedright\arraybackslash}p{(\linewidth - 14\tabcolsep) * \real{0.1250}}
  >{\raggedright\arraybackslash}p{(\linewidth - 14\tabcolsep) * \real{0.1250}}@{}}
\toprule\noalign{}
\begin{minipage}[b]{\linewidth}\raggedright
RUNSPEC
\end{minipage} & \begin{minipage}[b]{\linewidth}\raggedright
GRID
\end{minipage} & \begin{minipage}[b]{\linewidth}\raggedright
EDIT
\end{minipage} & \begin{minipage}[b]{\linewidth}\raggedright
PROPS
\end{minipage} & \begin{minipage}[b]{\linewidth}\raggedright
REGIONS
\end{minipage} & \begin{minipage}[b]{\linewidth}\raggedright
SOLUTION
\end{minipage} & \begin{minipage}[b]{\linewidth}\raggedright
SUMMARY
\end{minipage} & \begin{minipage}[b]{\linewidth}\raggedright
SCHEDULE
\end{minipage} \\
\midrule\noalign{}
\endhead
\bottomrule\noalign{}
\endlastfoot
\end{longtable}

\subsubsection{Description}\label{description-678}

The STOW keyword defines capillary pressure oil-water surface tension
versus pressure tables used in adjusting the pressure independent
capillary pressure vectors in the SWFN or SWOF saturation tables,
entered by their respective keywords in the PROPS section. The SATOPTS
keyword in the RUNSPEC section should state the SURFTENS character
string to activate the Capillary Pressure Surface Tension Pressure
Dependency option. If the STOW keyword is not supplied then no capillary
pressure surface tension pressure scaling will occur and the capillary
pressure values on the SWFN or SWOF saturation tables will be used
directly.

Capillary pressure surface tension pressure scaling can also be used
with multi-segment wells, but this facility has not been incorporated in
OPM Flow's multi-segment well implementation.

This keyword is not supported by OPM Flow but would change the results
if supported so the simulation will be stopped.

\subsection{STWG -- Define Capillary Pressure Water-Gas Surface Tension
versus
Pressure}\label{stwg-define-capillary-pressure-water-gas-surface-tension-versus-pressure}

\begin{longtable}[]{@{}
  >{\raggedright\arraybackslash}p{(\linewidth - 14\tabcolsep) * \real{0.1250}}
  >{\raggedright\arraybackslash}p{(\linewidth - 14\tabcolsep) * \real{0.1250}}
  >{\raggedright\arraybackslash}p{(\linewidth - 14\tabcolsep) * \real{0.1250}}
  >{\raggedright\arraybackslash}p{(\linewidth - 14\tabcolsep) * \real{0.1250}}
  >{\raggedright\arraybackslash}p{(\linewidth - 14\tabcolsep) * \real{0.1250}}
  >{\raggedright\arraybackslash}p{(\linewidth - 14\tabcolsep) * \real{0.1250}}
  >{\raggedright\arraybackslash}p{(\linewidth - 14\tabcolsep) * \real{0.1250}}
  >{\raggedright\arraybackslash}p{(\linewidth - 14\tabcolsep) * \real{0.1250}}@{}}
\toprule\noalign{}
\begin{minipage}[b]{\linewidth}\raggedright
RUNSPEC
\end{minipage} & \begin{minipage}[b]{\linewidth}\raggedright
GRID
\end{minipage} & \begin{minipage}[b]{\linewidth}\raggedright
EDIT
\end{minipage} & \begin{minipage}[b]{\linewidth}\raggedright
PROPS
\end{minipage} & \begin{minipage}[b]{\linewidth}\raggedright
REGIONS
\end{minipage} & \begin{minipage}[b]{\linewidth}\raggedright
SOLUTION
\end{minipage} & \begin{minipage}[b]{\linewidth}\raggedright
SUMMARY
\end{minipage} & \begin{minipage}[b]{\linewidth}\raggedright
SCHEDULE
\end{minipage} \\
\midrule\noalign{}
\endhead
\bottomrule\noalign{}
\endlastfoot
\end{longtable}

\subsubsection{Description}\label{description-679}

The STWG keyword defines capillary pressure water-gas surface tension
versus pressure tables for use with multi-segment wells. This facility
has not been incorporated in OPM Flow's multi-segment well
implementation. Note that STWG is not required for Capillary Pressure
Surface Tension Pressure Dependency option.

This keyword is not supported by OPM Flow but would change the results
if supported so the simulation will be stopped.

\subsection{SURFADDW -- Defined Surfactant Adsorbed Concentration versus
Wettability
Fraction}\label{surfaddw-defined-surfactant-adsorbed-concentration-versus-wettability-fraction}

\begin{longtable}[]{@{}
  >{\raggedright\arraybackslash}p{(\linewidth - 14\tabcolsep) * \real{0.1250}}
  >{\raggedright\arraybackslash}p{(\linewidth - 14\tabcolsep) * \real{0.1250}}
  >{\raggedright\arraybackslash}p{(\linewidth - 14\tabcolsep) * \real{0.1250}}
  >{\raggedright\arraybackslash}p{(\linewidth - 14\tabcolsep) * \real{0.1250}}
  >{\raggedright\arraybackslash}p{(\linewidth - 14\tabcolsep) * \real{0.1250}}
  >{\raggedright\arraybackslash}p{(\linewidth - 14\tabcolsep) * \real{0.1250}}
  >{\raggedright\arraybackslash}p{(\linewidth - 14\tabcolsep) * \real{0.1250}}
  >{\raggedright\arraybackslash}p{(\linewidth - 14\tabcolsep) * \real{0.1250}}@{}}
\toprule\noalign{}
\begin{minipage}[b]{\linewidth}\raggedright
RUNSPEC
\end{minipage} & \begin{minipage}[b]{\linewidth}\raggedright
GRID
\end{minipage} & \begin{minipage}[b]{\linewidth}\raggedright
EDIT
\end{minipage} & \begin{minipage}[b]{\linewidth}\raggedright
PROPS
\end{minipage} & \begin{minipage}[b]{\linewidth}\raggedright
REGIONS
\end{minipage} & \begin{minipage}[b]{\linewidth}\raggedright
SOLUTION
\end{minipage} & \begin{minipage}[b]{\linewidth}\raggedright
SUMMARY
\end{minipage} & \begin{minipage}[b]{\linewidth}\raggedright
SCHEDULE
\end{minipage} \\
\midrule\noalign{}
\endhead
\bottomrule\noalign{}
\endlastfoot
\end{longtable}

\subsubsection{Description}\label{description-680}

SURFADDW defines tables of surfactant adsorbed concentration versus
wettability fraction for when the SURFACTW keyword in the RUNSPEC
section as been declared to activate the surfactant phase with changing
wettability. The tables consists of columnar vectors of adsorbed
surfactant concentration versus a wettability fraction that indicates
the fraction of phase wettability. Here, a wettability fraction of zero
indicates a 100\% water wet rock resulting in the SURFWNUM allocated
saturation tables being used, and a value of one meaning 100\% oil wet
rock, with the SATNUM allocated saturations tables being employed. Both
the SURFWNUM and SATNUM keywords are in the REGIONS section.

This keyword is not supported by OPM Flow but would change the results
if supported so the simulation will be stopped.

\subsection{SURFADS -- Define Surfactant Rock Adsorption
Tables}\label{surfads-define-surfactant-rock-adsorption-tables}

\begin{longtable}[]{@{}
  >{\raggedright\arraybackslash}p{(\linewidth - 14\tabcolsep) * \real{0.1250}}
  >{\raggedright\arraybackslash}p{(\linewidth - 14\tabcolsep) * \real{0.1250}}
  >{\raggedright\arraybackslash}p{(\linewidth - 14\tabcolsep) * \real{0.1250}}
  >{\raggedright\arraybackslash}p{(\linewidth - 14\tabcolsep) * \real{0.1250}}
  >{\raggedright\arraybackslash}p{(\linewidth - 14\tabcolsep) * \real{0.1250}}
  >{\raggedright\arraybackslash}p{(\linewidth - 14\tabcolsep) * \real{0.1250}}
  >{\raggedright\arraybackslash}p{(\linewidth - 14\tabcolsep) * \real{0.1250}}
  >{\raggedright\arraybackslash}p{(\linewidth - 14\tabcolsep) * \real{0.1250}}@{}}
\toprule\noalign{}
\begin{minipage}[b]{\linewidth}\raggedright
RUNSPEC
\end{minipage} & \begin{minipage}[b]{\linewidth}\raggedright
GRID
\end{minipage} & \begin{minipage}[b]{\linewidth}\raggedright
EDIT
\end{minipage} & \begin{minipage}[b]{\linewidth}\raggedright
PROPS
\end{minipage} & \begin{minipage}[b]{\linewidth}\raggedright
REGIONS
\end{minipage} & \begin{minipage}[b]{\linewidth}\raggedright
SOLUTION
\end{minipage} & \begin{minipage}[b]{\linewidth}\raggedright
SUMMARY
\end{minipage} & \begin{minipage}[b]{\linewidth}\raggedright
SCHEDULE
\end{minipage} \\
\midrule\noalign{}
\endhead
\bottomrule\noalign{}
\endlastfoot
\end{longtable}

\subsubsection{Description}\label{description-681}

The SURFADS keyword defines the rock surfactant adsorption tables for
when the surfactant option has been activated by the SURFACT keyword in
the RUNSPEC section.

This keyword is not supported by OPM Flow but would change the results
if supported so the simulation will be stopped.

\begin{longtable}[]{@{}
  >{\raggedright\arraybackslash}p{(\linewidth - 6\tabcolsep) * \real{0.2500}}
  >{\raggedright\arraybackslash}p{(\linewidth - 6\tabcolsep) * \real{0.2500}}
  >{\raggedright\arraybackslash}p{(\linewidth - 6\tabcolsep) * \real{0.2500}}
  >{\raggedright\arraybackslash}p{(\linewidth - 6\tabcolsep) * \real{0.2500}}@{}}
\toprule\noalign{}
\begin{minipage}[b]{\linewidth}\raggedright
No.
\end{minipage} & \begin{minipage}[b]{\linewidth}\raggedright
Name
\end{minipage} & \begin{minipage}[b]{\linewidth}\raggedright
Description
\end{minipage} & \begin{minipage}[b]{\linewidth}\raggedright
Default
\end{minipage} \\
\midrule\noalign{}
\endhead
\bottomrule\noalign{}
\endlastfoot
Field & Metric & Laboratory & \\
1 & SURCON & A columnar vector of real monotonically increasing down the
column values that defines the surfactant concentration in the solution
surrounding the rock. The first entry should be zero to define a no
surfactant concentration. & None \\
lb/stb & kg/sm3 & gm/scc & \\
2 & SURRATIO & A columnar vector of real increasing down the column
values that defines the mass of adsorbed surfactant per unit mass of
rock of the saturated concentration of surfactant adsorbed by the rock.
The first entry should be zero to define a zero ratio of surfactant
concentration. & None \\
lb/lb & kg/kg & gm/gm & \\
Notes: & & & \\
\end{longtable}

\emph{Table 8.181: SURFADS Keyword Description}

See also the ADSORP keyword in the PROPS section that employs adsorption
functions, as oppose to adsorption tables, to define rock surfactant
adsorption behavior.

\subsubsection{Example}\label{example-393}

\begin{verbatim}
--
--       SURFACTANT ROCK ADSORPTION TABLE
--
SURFADS
--       SURF       SURF
--       SURCON     SURRATIO
--       -------    --------
             0.0     0.00000
             2.0     0.00003
             4.0     0.00005
             6.0     0.00007
             8.0     0.00009
            10.0     0.00011
            12.0     0.00012
            14.0     0.00015                               / TABLE NO. 01
--       SURF       SURF
--       SURCON     SURRATIO
--       -------    --------
             0.0     0.00000
             3.0     0.00004
             5.0     0.00006
             7.0     0.00008
             8.0     0.00009
            10.0     0.00011                               / TABLE NO. 02
\end{verbatim}

The above example defines two surfactant rock adsorption tables assuming
NTSFUN equals two and NSSFUN is greater than or equal to eight on the
TABDIMS keyword in the RUNSPEC section.

There is no terminating ``/'' for this keyword.

\subsection{SURFCAPD -- Capillary Number versus Miscibility
Tables}\label{surfcapd-capillary-number-versus-miscibility-tables}

\begin{longtable}[]{@{}
  >{\raggedright\arraybackslash}p{(\linewidth - 14\tabcolsep) * \real{0.1250}}
  >{\raggedright\arraybackslash}p{(\linewidth - 14\tabcolsep) * \real{0.1250}}
  >{\raggedright\arraybackslash}p{(\linewidth - 14\tabcolsep) * \real{0.1250}}
  >{\raggedright\arraybackslash}p{(\linewidth - 14\tabcolsep) * \real{0.1250}}
  >{\raggedright\arraybackslash}p{(\linewidth - 14\tabcolsep) * \real{0.1250}}
  >{\raggedright\arraybackslash}p{(\linewidth - 14\tabcolsep) * \real{0.1250}}
  >{\raggedright\arraybackslash}p{(\linewidth - 14\tabcolsep) * \real{0.1250}}
  >{\raggedright\arraybackslash}p{(\linewidth - 14\tabcolsep) * \real{0.1250}}@{}}
\toprule\noalign{}
\begin{minipage}[b]{\linewidth}\raggedright
RUNSPEC
\end{minipage} & \begin{minipage}[b]{\linewidth}\raggedright
GRID
\end{minipage} & \begin{minipage}[b]{\linewidth}\raggedright
EDIT
\end{minipage} & \begin{minipage}[b]{\linewidth}\raggedright
PROPS
\end{minipage} & \begin{minipage}[b]{\linewidth}\raggedright
REGIONS
\end{minipage} & \begin{minipage}[b]{\linewidth}\raggedright
SOLUTION
\end{minipage} & \begin{minipage}[b]{\linewidth}\raggedright
SUMMARY
\end{minipage} & \begin{minipage}[b]{\linewidth}\raggedright
SCHEDULE
\end{minipage} \\
\midrule\noalign{}
\endhead
\bottomrule\noalign{}
\endlastfoot
\end{longtable}

\subsubsection{Description}\label{description-682}

The SURFCAPD keyword defines the relationship between the log of the
capillary number and the level of miscibility, for when the Surfactant
option has been activated by the SURFACT keyword in the RUNSPEC section.
A value of zero for the level of miscibility means fully immiscible
conditions and consequently a value of one implies fully miscible
conditions.

This keyword is not supported by OPM Flow but would change the results
if supported so the simulation will be stopped.

\subsection{SURFESAL -- Define Surfactant Effective Salinity
Coefficient}\label{surfesal-define-surfactant-effective-salinity-coefficient}

\begin{longtable}[]{@{}
  >{\raggedright\arraybackslash}p{(\linewidth - 14\tabcolsep) * \real{0.1250}}
  >{\raggedright\arraybackslash}p{(\linewidth - 14\tabcolsep) * \real{0.1250}}
  >{\raggedright\arraybackslash}p{(\linewidth - 14\tabcolsep) * \real{0.1250}}
  >{\raggedright\arraybackslash}p{(\linewidth - 14\tabcolsep) * \real{0.1250}}
  >{\raggedright\arraybackslash}p{(\linewidth - 14\tabcolsep) * \real{0.1250}}
  >{\raggedright\arraybackslash}p{(\linewidth - 14\tabcolsep) * \real{0.1250}}
  >{\raggedright\arraybackslash}p{(\linewidth - 14\tabcolsep) * \real{0.1250}}
  >{\raggedright\arraybackslash}p{(\linewidth - 14\tabcolsep) * \real{0.1250}}@{}}
\toprule\noalign{}
\begin{minipage}[b]{\linewidth}\raggedright
RUNSPEC
\end{minipage} & \begin{minipage}[b]{\linewidth}\raggedright
GRID
\end{minipage} & \begin{minipage}[b]{\linewidth}\raggedright
EDIT
\end{minipage} & \begin{minipage}[b]{\linewidth}\raggedright
PROPS
\end{minipage} & \begin{minipage}[b]{\linewidth}\raggedright
REGIONS
\end{minipage} & \begin{minipage}[b]{\linewidth}\raggedright
SOLUTION
\end{minipage} & \begin{minipage}[b]{\linewidth}\raggedright
SUMMARY
\end{minipage} & \begin{minipage}[b]{\linewidth}\raggedright
SCHEDULE
\end{minipage} \\
\midrule\noalign{}
\endhead
\bottomrule\noalign{}
\endlastfoot
\end{longtable}

\subsubsection{Description}\label{description-683}

This keyword, SURFESAL, defines the surfactant effective salinity
coefficient as well as enabling the effective salinity calculation for
surfactant adsorption. The keyword should only be used if the BRINE
keyword has been declared to activate the brine phase, the ECLMC keyword
to enable the Multi-Component Brine model, and the SURFACT keyword has
been used to activate the surfact phase. All three keywords are in the
RUNSPEC section.

This keyword is not supported by OPM Flow but would change the results
if supported so the simulation will be stopped.

\subsection{SURFROCK -- Define Surfactant-Rock
Properties}\label{surfrock-define-surfactant-rock-properties}

\begin{longtable}[]{@{}
  >{\raggedright\arraybackslash}p{(\linewidth - 14\tabcolsep) * \real{0.1250}}
  >{\raggedright\arraybackslash}p{(\linewidth - 14\tabcolsep) * \real{0.1250}}
  >{\raggedright\arraybackslash}p{(\linewidth - 14\tabcolsep) * \real{0.1250}}
  >{\raggedright\arraybackslash}p{(\linewidth - 14\tabcolsep) * \real{0.1250}}
  >{\raggedright\arraybackslash}p{(\linewidth - 14\tabcolsep) * \real{0.1250}}
  >{\raggedright\arraybackslash}p{(\linewidth - 14\tabcolsep) * \real{0.1250}}
  >{\raggedright\arraybackslash}p{(\linewidth - 14\tabcolsep) * \real{0.1250}}
  >{\raggedright\arraybackslash}p{(\linewidth - 14\tabcolsep) * \real{0.1250}}@{}}
\toprule\noalign{}
\begin{minipage}[b]{\linewidth}\raggedright
RUNSPEC
\end{minipage} & \begin{minipage}[b]{\linewidth}\raggedright
GRID
\end{minipage} & \begin{minipage}[b]{\linewidth}\raggedright
EDIT
\end{minipage} & \begin{minipage}[b]{\linewidth}\raggedright
PROPS
\end{minipage} & \begin{minipage}[b]{\linewidth}\raggedright
REGIONS
\end{minipage} & \begin{minipage}[b]{\linewidth}\raggedright
SOLUTION
\end{minipage} & \begin{minipage}[b]{\linewidth}\raggedright
SUMMARY
\end{minipage} & \begin{minipage}[b]{\linewidth}\raggedright
SCHEDULE
\end{minipage} \\
\midrule\noalign{}
\endhead
\bottomrule\noalign{}
\endlastfoot
\end{longtable}

\subsubsection{Description}\label{description-684}

The SURFROCK keyword defines rock properties for when the Surfactant
option has been activated by the SURFACTANT keyword in the RUNSPEC
section.

This keyword is not supported by OPM Flow but would change the results
if supported so the simulation will be stopped.

\begin{longtable}[]{@{}
  >{\raggedright\arraybackslash}p{(\linewidth - 6\tabcolsep) * \real{0.2500}}
  >{\raggedright\arraybackslash}p{(\linewidth - 6\tabcolsep) * \real{0.2500}}
  >{\raggedright\arraybackslash}p{(\linewidth - 6\tabcolsep) * \real{0.2500}}
  >{\raggedright\arraybackslash}p{(\linewidth - 6\tabcolsep) * \real{0.2500}}@{}}
\toprule\noalign{}
\begin{minipage}[b]{\linewidth}\raggedright
No.
\end{minipage} & \begin{minipage}[b]{\linewidth}\raggedright
Name
\end{minipage} & \begin{minipage}[b]{\linewidth}\raggedright
Description
\end{minipage} & \begin{minipage}[b]{\linewidth}\raggedright
Default
\end{minipage} \\
\midrule\noalign{}
\endhead
\bottomrule\noalign{}
\endlastfoot
Field & Metric & Laboratory & \\
1 & ADINDX & A positive integer of 1 or 2 that defines the surfactant
desorption option. & Defined \\
dimensionless 1 & dimensionless 1 & dimensionless 1 & \\
2 & DENSITY & A real value that defines the rock in-situ density, that
is at reservoir conditions. & None \\
lb/rtb & kg/rm3 & gm/rcc & \\
Notes: & & & \\
\end{longtable}

\emph{Table 8.182: SURFROCK Keyword Description}

\subsubsection{Example}\label{example-394}

\begin{verbatim}
--
--       SURFACTANT-ROCK PROPERTIES
--
SURFROCK
--       DESORP   INSITU
--       OPTN     DENSITY
--       ------   -------
           1      1800.0                                   / TABLE NO. 01
           2      1980.0                                   / TABLE NO. 02
           1      2005.0                                   / TABLE NO. 03
\end{verbatim}

The above example defines three surfactant-rock tables, based on the
NTSFUN variable on the TABDIMS keyword in the RUNSPEC section being
equal to three.

There is no terminating ``/'' for this keyword.

\subsection{SURFST -- Surfactant Water-Oil Surface Tension versus
Surfactant
Concentration}\label{surfst-surfactant-water-oil-surface-tension-versus-surfactant-concentration}

\begin{longtable}[]{@{}
  >{\raggedright\arraybackslash}p{(\linewidth - 14\tabcolsep) * \real{0.1250}}
  >{\raggedright\arraybackslash}p{(\linewidth - 14\tabcolsep) * \real{0.1250}}
  >{\raggedright\arraybackslash}p{(\linewidth - 14\tabcolsep) * \real{0.1250}}
  >{\raggedright\arraybackslash}p{(\linewidth - 14\tabcolsep) * \real{0.1250}}
  >{\raggedright\arraybackslash}p{(\linewidth - 14\tabcolsep) * \real{0.1250}}
  >{\raggedright\arraybackslash}p{(\linewidth - 14\tabcolsep) * \real{0.1250}}
  >{\raggedright\arraybackslash}p{(\linewidth - 14\tabcolsep) * \real{0.1250}}
  >{\raggedright\arraybackslash}p{(\linewidth - 14\tabcolsep) * \real{0.1250}}@{}}
\toprule\noalign{}
\begin{minipage}[b]{\linewidth}\raggedright
RUNSPEC
\end{minipage} & \begin{minipage}[b]{\linewidth}\raggedright
GRID
\end{minipage} & \begin{minipage}[b]{\linewidth}\raggedright
EDIT
\end{minipage} & \begin{minipage}[b]{\linewidth}\raggedright
PROPS
\end{minipage} & \begin{minipage}[b]{\linewidth}\raggedright
REGIONS
\end{minipage} & \begin{minipage}[b]{\linewidth}\raggedright
SOLUTION
\end{minipage} & \begin{minipage}[b]{\linewidth}\raggedright
SUMMARY
\end{minipage} & \begin{minipage}[b]{\linewidth}\raggedright
SCHEDULE
\end{minipage} \\
\midrule\noalign{}
\endhead
\bottomrule\noalign{}
\endlastfoot
\end{longtable}

\subsubsection{Description}\label{description-685}

The SURFST keyword defines surfactant water-oil surface tension versus
surfactant concentration in the water phase tables, used in adjusting
the pressure independent capillary pressure vectors in the SWFN or SWOF
saturation tables, entered by their respective keywords in the PROPS
section. SURFST is also used to adjust the relative permeability curves
on the aforementioned tables via the capillary number. The Surfactant
option must have been activated by the SURFACTANT keyword in the RUNSPEC
section to use this keyword and either this keyword or the SURFSTES
keyword, also in the PROPS section, is obligatory in this case.

See also the SURFSTES that defines the surfactant water-oil surface
tension as a function of surfactant concentration in the water phase and
salt concentration or the effective salinity.

This keyword is not supported by OPM Flow but would change the results
if supported so the simulation will be stopped.

\subsection{SURFSTES -- Surfactant Water-Oil Surface Tension versus
Surfactant and Salt
Concentrations}\label{surfstes-surfactant-water-oil-surface-tension-versus-surfactant-and-salt-concentrations}

\begin{longtable}[]{@{}
  >{\raggedright\arraybackslash}p{(\linewidth - 14\tabcolsep) * \real{0.1250}}
  >{\raggedright\arraybackslash}p{(\linewidth - 14\tabcolsep) * \real{0.1250}}
  >{\raggedright\arraybackslash}p{(\linewidth - 14\tabcolsep) * \real{0.1250}}
  >{\raggedright\arraybackslash}p{(\linewidth - 14\tabcolsep) * \real{0.1250}}
  >{\raggedright\arraybackslash}p{(\linewidth - 14\tabcolsep) * \real{0.1250}}
  >{\raggedright\arraybackslash}p{(\linewidth - 14\tabcolsep) * \real{0.1250}}
  >{\raggedright\arraybackslash}p{(\linewidth - 14\tabcolsep) * \real{0.1250}}
  >{\raggedright\arraybackslash}p{(\linewidth - 14\tabcolsep) * \real{0.1250}}@{}}
\toprule\noalign{}
\begin{minipage}[b]{\linewidth}\raggedright
RUNSPEC
\end{minipage} & \begin{minipage}[b]{\linewidth}\raggedright
GRID
\end{minipage} & \begin{minipage}[b]{\linewidth}\raggedright
EDIT
\end{minipage} & \begin{minipage}[b]{\linewidth}\raggedright
PROPS
\end{minipage} & \begin{minipage}[b]{\linewidth}\raggedright
REGIONS
\end{minipage} & \begin{minipage}[b]{\linewidth}\raggedright
SOLUTION
\end{minipage} & \begin{minipage}[b]{\linewidth}\raggedright
SUMMARY
\end{minipage} & \begin{minipage}[b]{\linewidth}\raggedright
SCHEDULE
\end{minipage} \\
\midrule\noalign{}
\endhead
\bottomrule\noalign{}
\endlastfoot
\end{longtable}

\subsubsection{Description}\label{description-686}

The SURFSTES keyword defines surfactant water-oil surface tension versus
surfactant concentration in the water phase tables, used in adjusting
the pressure independent capillary pressure vectors in the SWFN or SWOF
saturation tables, entered by their respective keywords in the PROPS
section. SURFSTES is also used to adjust the relative permeability
curves on the aforementioned tables via the capillary number. The
Surfactant option must have been activated by the SURFACTANT keyword in
the RUNSPEC section to use this keyword and either this keyword or the
SURFST keyword, also in the PROPS section, is obligatory in this case.
In addition, the BRINE keyword in the RUNSPEC section must be activated
and the ESSNODE keyword in the PROPS section must be used to define the
salt concentration or the effective salinity.

See also the SURFSTS that defines the surfactant water-oil surface
tension as a function of surfactant concentration in the water phase
only.

This keyword is not supported by OPM Flow but would change the results
if supported so the simulation will be stopped.

\subsection{SURFVISC -- Surfactant Solution Viscosity versus
Concentration}\label{surfvisc-surfactant-solution-viscosity-versus-concentration}

\begin{longtable}[]{@{}
  >{\raggedright\arraybackslash}p{(\linewidth - 14\tabcolsep) * \real{0.1250}}
  >{\raggedright\arraybackslash}p{(\linewidth - 14\tabcolsep) * \real{0.1250}}
  >{\raggedright\arraybackslash}p{(\linewidth - 14\tabcolsep) * \real{0.1250}}
  >{\raggedright\arraybackslash}p{(\linewidth - 14\tabcolsep) * \real{0.1250}}
  >{\raggedright\arraybackslash}p{(\linewidth - 14\tabcolsep) * \real{0.1250}}
  >{\raggedright\arraybackslash}p{(\linewidth - 14\tabcolsep) * \real{0.1250}}
  >{\raggedright\arraybackslash}p{(\linewidth - 14\tabcolsep) * \real{0.1250}}
  >{\raggedright\arraybackslash}p{(\linewidth - 14\tabcolsep) * \real{0.1250}}@{}}
\toprule\noalign{}
\begin{minipage}[b]{\linewidth}\raggedright
RUNSPEC
\end{minipage} & \begin{minipage}[b]{\linewidth}\raggedright
GRID
\end{minipage} & \begin{minipage}[b]{\linewidth}\raggedright
EDIT
\end{minipage} & \begin{minipage}[b]{\linewidth}\raggedright
PROPS
\end{minipage} & \begin{minipage}[b]{\linewidth}\raggedright
REGIONS
\end{minipage} & \begin{minipage}[b]{\linewidth}\raggedright
SOLUTION
\end{minipage} & \begin{minipage}[b]{\linewidth}\raggedright
SUMMARY
\end{minipage} & \begin{minipage}[b]{\linewidth}\raggedright
SCHEDULE
\end{minipage} \\
\midrule\noalign{}
\endhead
\bottomrule\noalign{}
\endlastfoot
\end{longtable}

\subsubsection{Description}\label{description-687}

SURFSVISC defines the surfactant viscosity relationship of solution
water viscosity with respect to increasing surfactant concentration
within a grid block. The surfactant option must be activated by the
SURFACT keyword in the RUNSPEC section in order to use this keyword.

This keyword is not supported by OPM Flow but would change the results
if supported so the simulation will be stopped.

\begin{longtable}[]{@{}
  >{\raggedright\arraybackslash}p{(\linewidth - 6\tabcolsep) * \real{0.2500}}
  >{\raggedright\arraybackslash}p{(\linewidth - 6\tabcolsep) * \real{0.2500}}
  >{\raggedright\arraybackslash}p{(\linewidth - 6\tabcolsep) * \real{0.2500}}
  >{\raggedright\arraybackslash}p{(\linewidth - 6\tabcolsep) * \real{0.2500}}@{}}
\toprule\noalign{}
\begin{minipage}[b]{\linewidth}\raggedright
No.
\end{minipage} & \begin{minipage}[b]{\linewidth}\raggedright
Name
\end{minipage} & \begin{minipage}[b]{\linewidth}\raggedright
Description
\end{minipage} & \begin{minipage}[b]{\linewidth}\raggedright
Default
\end{minipage} \\
\midrule\noalign{}
\endhead
\bottomrule\noalign{}
\endlastfoot
Field & Metric & Laboratory & \\
1 & SURFCON & A columnar vector of real monotonically increasing down
the column values that defines the surfactant concentration in the
solution surrounding the rock. The first entry should be zero to define
a no surfactant concentration. & None \\
lb/stb & kg/sm3 & gm/scc & \\
2 & SURFVISC & A columnar vector of real positive values that defines
the solution water viscosity of the solution for the given SURFCON entry
at the reference pressure value, PRES, entered on the PVTW keyword in
the PROPS section. & None \\
cP & cP & cP & \\
Notes: & & & \\
\end{longtable}

\emph{Table 8.183: SURFVISC Keyword Description}

\subsubsection{Example}\label{example-395}

\begin{verbatim}
--
--       SURFACTANT SOLUTION WATER VISCOSITY TABLES
--
SURFVISC
--       SURF         VISCOSITY
--       SURFCON      SURFVISC
--       --------    -----------
          0.0000      0.3500
          0.0100      0.3900
          0.0200      0.4200
          0.0300      0.4300                               / TABLE NO. 01

--       SURF         VISCOSITY
--       SURFCON      VISFAC
--       --------    ---------
          0.0000       1.000
          0.0003      10.000
          0.0005      20.000
          0.0007      40.000
          0.0009      45.000
          0.0011      55.000                               / TABLE NO. 02
\end{verbatim}

The example defines two surfactant viscosity scaling factor tables,
based on the NTPVT variable on the TABDIMS keyword in the RUNSPEC
section being equal to two and NPPVT variable on the same keyword being
greater than or equal to six.

\subsection{SWATINIT -- Define the Initial Water Saturation Array for
Capillary Pressure
Scaling}\label{swatinit-define-the-initial-water-saturation-array-for-capillary-pressure-scaling}

\begin{longtable}[]{@{}
  >{\raggedright\arraybackslash}p{(\linewidth - 14\tabcolsep) * \real{0.1250}}
  >{\raggedright\arraybackslash}p{(\linewidth - 14\tabcolsep) * \real{0.1250}}
  >{\raggedright\arraybackslash}p{(\linewidth - 14\tabcolsep) * \real{0.1250}}
  >{\raggedright\arraybackslash}p{(\linewidth - 14\tabcolsep) * \real{0.1250}}
  >{\raggedright\arraybackslash}p{(\linewidth - 14\tabcolsep) * \real{0.1250}}
  >{\raggedright\arraybackslash}p{(\linewidth - 14\tabcolsep) * \real{0.1250}}
  >{\raggedright\arraybackslash}p{(\linewidth - 14\tabcolsep) * \real{0.1250}}
  >{\raggedright\arraybackslash}p{(\linewidth - 14\tabcolsep) * \real{0.1250}}@{}}
\toprule\noalign{}
\begin{minipage}[b]{\linewidth}\raggedright
RUNSPEC
\end{minipage} & \begin{minipage}[b]{\linewidth}\raggedright
GRID
\end{minipage} & \begin{minipage}[b]{\linewidth}\raggedright
EDIT
\end{minipage} & \begin{minipage}[b]{\linewidth}\raggedright
PROPS
\end{minipage} & \begin{minipage}[b]{\linewidth}\raggedright
REGIONS
\end{minipage} & \begin{minipage}[b]{\linewidth}\raggedright
SOLUTION
\end{minipage} & \begin{minipage}[b]{\linewidth}\raggedright
SUMMARY
\end{minipage} & \begin{minipage}[b]{\linewidth}\raggedright
SCHEDULE
\end{minipage} \\
\midrule\noalign{}
\endhead
\bottomrule\noalign{}
\endlastfoot
\end{longtable}

\subsubsection{Description}\label{description-688}

SWATINIT defines the initial water saturation for all the cells in the
model via an array. The keyword can be used with all grid types.
SWATINIT is used to initialize the model by setting each grid block's
initial water saturation (``Sw''). If the array is present in the input
deck, then OPM Flow will re-scale the water-oil capillary pressure
curves entered via the SWFN saturation functions in the PROPS section,
so that the resulting initialized Sw matches the values in the SWATINIT
array.

Normally the SWATINIT array is generated in the static earth model when
calculating the hydrocarbons in-place. volumes using Saturation Height
Functions (``SHF'') derived from capillary pressure functions. Static
earth models do not directly use capillary pressure in these type of
calculations as individual cell pressures are not required. There is
therefore some potential for inconsistencies to arise between the two
sets of formulations. This is normally manifested by extreme scaling in
the scaled capillary pressure values calculated by the simulator. If
this is the case then the PPCWMAX keyword can be used to set a maximum
scaled capillary pressure value. Note that as large values of scaled
capillary pressures can result in numerical issues, a more technically
sound approach would be to resolve these inconsistencies before
continuing with the model build.

\begin{longtable}[]{@{}
  >{\raggedright\arraybackslash}p{(\linewidth - 6\tabcolsep) * \real{0.2500}}
  >{\raggedright\arraybackslash}p{(\linewidth - 6\tabcolsep) * \real{0.2500}}
  >{\raggedright\arraybackslash}p{(\linewidth - 6\tabcolsep) * \real{0.2500}}
  >{\raggedright\arraybackslash}p{(\linewidth - 6\tabcolsep) * \real{0.2500}}@{}}
\toprule\noalign{}
\begin{minipage}[b]{\linewidth}\raggedright
No.
\end{minipage} & \begin{minipage}[b]{\linewidth}\raggedright
Name
\end{minipage} & \begin{minipage}[b]{\linewidth}\raggedright
Description
\end{minipage} & \begin{minipage}[b]{\linewidth}\raggedright
Default
\end{minipage} \\
\midrule\noalign{}
\endhead
\bottomrule\noalign{}
\endlastfoot
Field & Metric & Laboratory & \\
1 & SWATINIT & SWATINIT is an array of real positive numbers that are
greater than or equal to zero and less than or equal to one, that define
the initial water saturation values to each cell in the model. Repeat
counts may be used, for example 3000*0.15 & None \\
dimensionless & dimensionless & dimensionless & \\
Notes: & & & \\
\end{longtable}

\emph{Table 8.184: SWATINIT Keyword Description}

See also the PPCWMAX to control the re-scaling of the capillary pressure
entries on the SWFN saturation function keyword in the PROPS section.

\subsubsection{Example}\label{example-396}

\begin{verbatim}
--
--        DEFINE GRID BLOCK INITIAL SW DATA FOR ALL CELLS
--        (BASED ON NX x NY x NZ = 300)
--
SWATINIT
          300*0.300                                                            /
\end{verbatim}

The above example defines a constant initial water saturation of 0.300
to all 300 cells in the model as defined by the DIMENS keyword in the
RUNSPEC section.

\subsection{SWCR -- End-Point Scaling Grid Cell Critical Water
Saturation}\label{swcr-end-point-scaling-grid-cell-critical-water-saturation}

\begin{longtable}[]{@{}
  >{\raggedright\arraybackslash}p{(\linewidth - 14\tabcolsep) * \real{0.1250}}
  >{\raggedright\arraybackslash}p{(\linewidth - 14\tabcolsep) * \real{0.1250}}
  >{\raggedright\arraybackslash}p{(\linewidth - 14\tabcolsep) * \real{0.1250}}
  >{\raggedright\arraybackslash}p{(\linewidth - 14\tabcolsep) * \real{0.1250}}
  >{\raggedright\arraybackslash}p{(\linewidth - 14\tabcolsep) * \real{0.1250}}
  >{\raggedright\arraybackslash}p{(\linewidth - 14\tabcolsep) * \real{0.1250}}
  >{\raggedright\arraybackslash}p{(\linewidth - 14\tabcolsep) * \real{0.1250}}
  >{\raggedright\arraybackslash}p{(\linewidth - 14\tabcolsep) * \real{0.1250}}@{}}
\toprule\noalign{}
\begin{minipage}[b]{\linewidth}\raggedright
RUNSPEC
\end{minipage} & \begin{minipage}[b]{\linewidth}\raggedright
GRID
\end{minipage} & \begin{minipage}[b]{\linewidth}\raggedright
EDIT
\end{minipage} & \begin{minipage}[b]{\linewidth}\raggedright
PROPS
\end{minipage} & \begin{minipage}[b]{\linewidth}\raggedright
REGIONS
\end{minipage} & \begin{minipage}[b]{\linewidth}\raggedright
SOLUTION
\end{minipage} & \begin{minipage}[b]{\linewidth}\raggedright
SUMMARY
\end{minipage} & \begin{minipage}[b]{\linewidth}\raggedright
SCHEDULE
\end{minipage} \\
\midrule\noalign{}
\endhead
\bottomrule\noalign{}
\endlastfoot
\end{longtable}

\subsubsection{Description}\label{description-689}

SWCR defines the critical water saturation for all the cells in the
model via an array when the end-point scaling option has been invoked
via the ENDSCALE keyword in the RUNSPEC section. The critical water
saturation is defined as the maximum water saturation for which the
water relative permeability is zero in a two-phase relative permeability
table.

The keyword can be used with all grid types.

\begin{longtable}[]{@{}
  >{\raggedright\arraybackslash}p{(\linewidth - 6\tabcolsep) * \real{0.2500}}
  >{\raggedright\arraybackslash}p{(\linewidth - 6\tabcolsep) * \real{0.2500}}
  >{\raggedright\arraybackslash}p{(\linewidth - 6\tabcolsep) * \real{0.2500}}
  >{\raggedright\arraybackslash}p{(\linewidth - 6\tabcolsep) * \real{0.2500}}@{}}
\toprule\noalign{}
\begin{minipage}[b]{\linewidth}\raggedright
No.
\end{minipage} & \begin{minipage}[b]{\linewidth}\raggedright
Name
\end{minipage} & \begin{minipage}[b]{\linewidth}\raggedright
Description
\end{minipage} & \begin{minipage}[b]{\linewidth}\raggedright
Default
\end{minipage} \\
\midrule\noalign{}
\endhead
\bottomrule\noalign{}
\endlastfoot
Field & Metric & Laboratory & \\
1 & SWCR & SWCR is an array of real numbers assigning the critical water
saturation values to each cell in the model. The number of entries
should correspond to the NX x NY x NZ parameters on the DIMENS keyword.
Repeat counts may be used, for example 30*0.20 & Taken from cell
allocated relative permeability table. \\
dimensionless & dimensionless & dimensionless & \\
Notes: & & & \\
\end{longtable}

\emph{Table 8.185: SWCR Keyword Description}

End-point scaling allows the entered relative permeability functions to
be re-scaled based on the saturation end-points allocated to each cell
by the SWL, SWCR, SWU, SGL, SGCR, SGU, SOWCR, and SOGCR saturation grid
arrays for the saturation end-points, and the KRG, KRGR, KRO, KRORG,
KRORW, KRW and KRWR relative permeability grid cell arrays for the
relative permeability end-point data. In addition end-point scaling may
be directional dependent in which case the directional dependent
versions of the aforementioned arrays should be used, that is SWCRX,
SWCRY and SWCRZ instead of SWCR. There is also the facility to make the
directional end-point scaling reversible or non-reversible and if the
non-reversible option is selected the non-reversible versions of the
aforementioned arrays should be used, that is SWCRX, SWCRX-, SWCRY,
SWCRY-, SWCRZ and SWCRZ-, instead of the SWCR keyword.

\subsubsection{Example}\label{example-397}

\begin{verbatim}
--
--       DEFINE GRID BLOCK END-POINT SWCR DATA FOR ALL CELLS
--       (FOR NX x NY x NZ = 300)
--
SWCR
         300*0.200                                                              /
\end{verbatim}

The above example defines a constant critical water saturation of 0.20
to all 300 cells in the model as defined by the DIMENS keyword in the
RUNSPEC section.

\subsection{SWF32D -- Water Saturation Tables with Respect to Oil and
Gas (Three
Phase)}\label{swf32d-water-saturation-tables-with-respect-to-oil-and-gas-three-phase}

\begin{longtable}[]{@{}
  >{\raggedright\arraybackslash}p{(\linewidth - 14\tabcolsep) * \real{0.1250}}
  >{\raggedright\arraybackslash}p{(\linewidth - 14\tabcolsep) * \real{0.1250}}
  >{\raggedright\arraybackslash}p{(\linewidth - 14\tabcolsep) * \real{0.1250}}
  >{\raggedright\arraybackslash}p{(\linewidth - 14\tabcolsep) * \real{0.1250}}
  >{\raggedright\arraybackslash}p{(\linewidth - 14\tabcolsep) * \real{0.1250}}
  >{\raggedright\arraybackslash}p{(\linewidth - 14\tabcolsep) * \real{0.1250}}
  >{\raggedright\arraybackslash}p{(\linewidth - 14\tabcolsep) * \real{0.1250}}
  >{\raggedright\arraybackslash}p{(\linewidth - 14\tabcolsep) * \real{0.1250}}@{}}
\toprule\noalign{}
\begin{minipage}[b]{\linewidth}\raggedright
RUNSPEC
\end{minipage} & \begin{minipage}[b]{\linewidth}\raggedright
GRID
\end{minipage} & \begin{minipage}[b]{\linewidth}\raggedright
EDIT
\end{minipage} & \begin{minipage}[b]{\linewidth}\raggedright
PROPS
\end{minipage} & \begin{minipage}[b]{\linewidth}\raggedright
REGIONS
\end{minipage} & \begin{minipage}[b]{\linewidth}\raggedright
SOLUTION
\end{minipage} & \begin{minipage}[b]{\linewidth}\raggedright
SUMMARY
\end{minipage} & \begin{minipage}[b]{\linewidth}\raggedright
SCHEDULE
\end{minipage} \\
\midrule\noalign{}
\endhead
\bottomrule\noalign{}
\endlastfoot
\end{longtable}

\subsubsection{Description}\label{description-690}

The SWF32D keyword defines the three phase water relative permeability
versus oil and gas saturation tables for when oil, gas and water are
present in the input deck. The keyword should only be used if oil, gas
and water are present in the input deck. Normally the simulator
calculates the three-phase oil relative permeabilities based on the
entered two phase tables of water-oil and gas-oil, combined with the
STONE1 and STONE2 keywords in the PROPS section that determine the
method used to generate the thee phase water relative permeability
curves. SWF32D allows for the direct input of the three phase tables, as
such the STONE1 and STONE2 keywords should not be entered if SWF32D is
used in the input deck.

This keyword is not supported by OPM Flow but would change the results
if supported so the simulation will be stopped.

\subsection{SWFN -- Water Saturation Tables (Format Type
2)}\label{swfn-water-saturation-tables-format-type-2}

\begin{longtable}[]{@{}
  >{\raggedright\arraybackslash}p{(\linewidth - 14\tabcolsep) * \real{0.1250}}
  >{\raggedright\arraybackslash}p{(\linewidth - 14\tabcolsep) * \real{0.1250}}
  >{\raggedright\arraybackslash}p{(\linewidth - 14\tabcolsep) * \real{0.1250}}
  >{\raggedright\arraybackslash}p{(\linewidth - 14\tabcolsep) * \real{0.1250}}
  >{\raggedright\arraybackslash}p{(\linewidth - 14\tabcolsep) * \real{0.1250}}
  >{\raggedright\arraybackslash}p{(\linewidth - 14\tabcolsep) * \real{0.1250}}
  >{\raggedright\arraybackslash}p{(\linewidth - 14\tabcolsep) * \real{0.1250}}
  >{\raggedright\arraybackslash}p{(\linewidth - 14\tabcolsep) * \real{0.1250}}@{}}
\toprule\noalign{}
\begin{minipage}[b]{\linewidth}\raggedright
RUNSPEC
\end{minipage} & \begin{minipage}[b]{\linewidth}\raggedright
GRID
\end{minipage} & \begin{minipage}[b]{\linewidth}\raggedright
EDIT
\end{minipage} & \begin{minipage}[b]{\linewidth}\raggedright
PROPS
\end{minipage} & \begin{minipage}[b]{\linewidth}\raggedright
REGIONS
\end{minipage} & \begin{minipage}[b]{\linewidth}\raggedright
SOLUTION
\end{minipage} & \begin{minipage}[b]{\linewidth}\raggedright
SUMMARY
\end{minipage} & \begin{minipage}[b]{\linewidth}\raggedright
SCHEDULE
\end{minipage} \\
\midrule\noalign{}
\endhead
\bottomrule\noalign{}
\endlastfoot
\end{longtable}

\subsubsection{Description}\label{description-691}

The SWFN keyword defines the water relative permeability and water-oil
capillary pressure data versus water saturation tables for when water is
present in the input deck. This keyword should only be used if water is
present in the run.

\begin{longtable}[]{@{}
  >{\raggedright\arraybackslash}p{(\linewidth - 6\tabcolsep) * \real{0.2500}}
  >{\raggedright\arraybackslash}p{(\linewidth - 6\tabcolsep) * \real{0.2500}}
  >{\raggedright\arraybackslash}p{(\linewidth - 6\tabcolsep) * \real{0.2500}}
  >{\raggedright\arraybackslash}p{(\linewidth - 6\tabcolsep) * \real{0.2500}}@{}}
\toprule\noalign{}
\begin{minipage}[b]{\linewidth}\raggedright
No.
\end{minipage} & \begin{minipage}[b]{\linewidth}\raggedright
Name
\end{minipage} & \begin{minipage}[b]{\linewidth}\raggedright
Description
\end{minipage} & \begin{minipage}[b]{\linewidth}\raggedright
Default
\end{minipage} \\
\midrule\noalign{}
\endhead
\bottomrule\noalign{}
\endlastfoot
Field & Metric & Laboratory & \\
1 & SWAT & A columnar vector of real monotonically increasing down the
column values starting from zero and terminating at one, that defines
the water saturation. & None \\
dimensionless & dimensionless & dimensionless & \\
2 & KRW & A columnar vector of real values that are either equal or
increasing down the column and that are greater than or equal to zero
and less than or equal to one that defines the water relative
permeability with respect to gas saturation. The first value in the
column should be zero. & None \\
dimensionless & dimensionless & dimensionless & \\
3 & PCWO & A columnar vector of real values that are either equal or
increasing down the column that defines the water-oil relative capillary
pressure. If the SWATINIT keyword has been used to initialize the model
then columnar vector has to be strictly monotonically increasing. &
None \\
psia & bars & atm & \\
Notes: & & & \\
\end{longtable}

\emph{Table 8.3.332.1: SWFN Keyword Description}

\subsubsection{Example}\label{example-398}

\begin{verbatim}
--
--       WATER RELATIVE PERMEABILITY TABLES (SWFN)
--
SWFN
--       SWAT       KRW        PCOW
--       FRAC       FRAC       PSIA
--       --------   --------   -------
            0.15     0.00000     0.0
            0.20     6.25e-6     0.0
            0.25     0.00010     0.0
            0.30     0.00050     0.0
            0.35     0.00160     0.0
            0.40     0.00390     0.0
            0.45     0.00810     0.0
            0.50     0.01500     0.0
            0.55     0.02560     0.0
            0.60     0.04100     0.0
            0.65     0.06250     0.0
            0.70     0.09150     0.0
            0.75     0.12960     0.0
            0.80     0.17850     0.0
            0.85     0.24010     0.0
            0.90     0.31640     0.0
            0.95     0.40960     0.0
            1.00     0.52200     0.0                       / TABLE NO. 1
--       --------   --------   -------
            0.15     0.00000     0.0
            0.20     6.25e-6     0.0
            0.25     0.00010     0.0
            0.30     0.00050     0.0
            0.35     0.00160     0.0
            0.40     0.00390     0.0
            0.45     0.00810     0.0
            0.50     0.01500     0.0
            0.55     0.02560     0.0
            0.60     0.04100     0.0
            0.65     0.06250     0.0
            0.70     0.09150     0.0
            0.75     0.12960     0.0
            0.80     0.17850     0.0
            0.85     0.24010     0.0
            0.90     0.31640     0.0
            0.95     0.40960     0.0
            1.00     0.52200     0.0                       / TABLE NO. 2
\end{verbatim}

The example defines two SWFN tables for use when water is present in the
run.

\subsection{SWL -- End-Point Scaling Grid Cell Connate Water
Saturation}\label{swl-end-point-scaling-grid-cell-connate-water-saturation}

\begin{longtable}[]{@{}
  >{\raggedright\arraybackslash}p{(\linewidth - 14\tabcolsep) * \real{0.1250}}
  >{\raggedright\arraybackslash}p{(\linewidth - 14\tabcolsep) * \real{0.1250}}
  >{\raggedright\arraybackslash}p{(\linewidth - 14\tabcolsep) * \real{0.1250}}
  >{\raggedright\arraybackslash}p{(\linewidth - 14\tabcolsep) * \real{0.1250}}
  >{\raggedright\arraybackslash}p{(\linewidth - 14\tabcolsep) * \real{0.1250}}
  >{\raggedright\arraybackslash}p{(\linewidth - 14\tabcolsep) * \real{0.1250}}
  >{\raggedright\arraybackslash}p{(\linewidth - 14\tabcolsep) * \real{0.1250}}
  >{\raggedright\arraybackslash}p{(\linewidth - 14\tabcolsep) * \real{0.1250}}@{}}
\toprule\noalign{}
\begin{minipage}[b]{\linewidth}\raggedright
RUNSPEC
\end{minipage} & \begin{minipage}[b]{\linewidth}\raggedright
GRID
\end{minipage} & \begin{minipage}[b]{\linewidth}\raggedright
EDIT
\end{minipage} & \begin{minipage}[b]{\linewidth}\raggedright
PROPS
\end{minipage} & \begin{minipage}[b]{\linewidth}\raggedright
REGIONS
\end{minipage} & \begin{minipage}[b]{\linewidth}\raggedright
SOLUTION
\end{minipage} & \begin{minipage}[b]{\linewidth}\raggedright
SUMMARY
\end{minipage} & \begin{minipage}[b]{\linewidth}\raggedright
SCHEDULE
\end{minipage} \\
\midrule\noalign{}
\endhead
\bottomrule\noalign{}
\endlastfoot
\end{longtable}

\subsubsection{Description}\label{description-692}

SWL defines the connate water saturation for all the cells in the model
via an array when the end-point scaling option has been invoked via the
ENDSCALE keyword in the RUNSPEC section. The connate water saturation is
defined as the minimum water saturation in a two-phase water relative
permeability table.

The keyword can be used with all grid types.

\begin{longtable}[]{@{}
  >{\raggedright\arraybackslash}p{(\linewidth - 6\tabcolsep) * \real{0.2500}}
  >{\raggedright\arraybackslash}p{(\linewidth - 6\tabcolsep) * \real{0.2500}}
  >{\raggedright\arraybackslash}p{(\linewidth - 6\tabcolsep) * \real{0.2500}}
  >{\raggedright\arraybackslash}p{(\linewidth - 6\tabcolsep) * \real{0.2500}}@{}}
\toprule\noalign{}
\begin{minipage}[b]{\linewidth}\raggedright
No.
\end{minipage} & \begin{minipage}[b]{\linewidth}\raggedright
Name
\end{minipage} & \begin{minipage}[b]{\linewidth}\raggedright
Description
\end{minipage} & \begin{minipage}[b]{\linewidth}\raggedright
Default
\end{minipage} \\
\midrule\noalign{}
\endhead
\bottomrule\noalign{}
\endlastfoot
Field & Metric & Laboratory & \\
1 & SWL & SWL is an array of real numbers assigning the connate water
saturation values to each cell in the model. The number of entries
should correspond to the NX x NY x NZ parameters on the DIMENS keyword.
Repeat counts may be used, for example 30*0.15 & Taken from cell
allocated relative permeability table. \\
dimensionless & dimensionless & dimensionless & \\
Notes: & & & \\
\end{longtable}

\emph{Table 8.187: SWL Keyword Description}

End-point scaling allows the entered relative permeability functions to
be re-scaled based on the saturation end-points allocated to each cell
by the SWL, SWCR, SWU, SGL, SGCR, SGU, SOWCR, and SOGCR saturation grid
arrays for the saturation end-points, and the KRG, KRORG, KRORW and KRW
relative permeability grid cell arrays for the relative permeability
end-point data. In addition end-point scaling may be directional
dependent in which case the directional dependent versions of the
aforementioned arrays should be used, that is SWLX, SWLY and SWLZ
instead of SWL. There is also the facility to make the directional
end-point scaling reversible or non-reversible and if the non-reversible
option is selected the non-reversible versions of the aforementioned
arrays should be used, that is SWLX, SWLX-, SWLY, SWLY-, SWLZ and SWLZ-,
instead of the SWL keyword.

\subsubsection{Example}\label{example-399}

\begin{verbatim}
--
--       DEFINE GRID BLOCK END-POINT SWL DATA FOR ALL CELLS
--       (FOR NX x NY x NZ = 300)
--
SWL
         300*0.150                                                              /
\end{verbatim}

The above example defines a constant connate water saturation of 0.15 to
all 300 cells in the model as defined by the DIMENS keyword in the
RUNSPEC section.

\subsection{SWLPC -- End-Point Scaling Grid Cell Capillary Pressure
Connate Water
Saturations}\label{swlpc-end-point-scaling-grid-cell-capillary-pressure-connate-water-saturations}

\begin{longtable}[]{@{}
  >{\raggedright\arraybackslash}p{(\linewidth - 14\tabcolsep) * \real{0.1250}}
  >{\raggedright\arraybackslash}p{(\linewidth - 14\tabcolsep) * \real{0.1250}}
  >{\raggedright\arraybackslash}p{(\linewidth - 14\tabcolsep) * \real{0.1250}}
  >{\raggedright\arraybackslash}p{(\linewidth - 14\tabcolsep) * \real{0.1250}}
  >{\raggedright\arraybackslash}p{(\linewidth - 14\tabcolsep) * \real{0.1250}}
  >{\raggedright\arraybackslash}p{(\linewidth - 14\tabcolsep) * \real{0.1250}}
  >{\raggedright\arraybackslash}p{(\linewidth - 14\tabcolsep) * \real{0.1250}}
  >{\raggedright\arraybackslash}p{(\linewidth - 14\tabcolsep) * \real{0.1250}}@{}}
\toprule\noalign{}
\begin{minipage}[b]{\linewidth}\raggedright
RUNSPEC
\end{minipage} & \begin{minipage}[b]{\linewidth}\raggedright
GRID
\end{minipage} & \begin{minipage}[b]{\linewidth}\raggedright
EDIT
\end{minipage} & \begin{minipage}[b]{\linewidth}\raggedright
PROPS
\end{minipage} & \begin{minipage}[b]{\linewidth}\raggedright
REGIONS
\end{minipage} & \begin{minipage}[b]{\linewidth}\raggedright
SOLUTION
\end{minipage} & \begin{minipage}[b]{\linewidth}\raggedright
SUMMARY
\end{minipage} & \begin{minipage}[b]{\linewidth}\raggedright
SCHEDULE
\end{minipage} \\
\midrule\noalign{}
\endhead
\bottomrule\noalign{}
\endlastfoot
\end{longtable}

\subsubsection{Description}\label{description-693}

SWLPC defines the connate water saturation for all the cells in the
model via an array when the end-point scaling option has been invoked
via the ENDSCALE keyword in the RUNSPEC section. The connate water
saturation is defined as the minimum water saturation in a two-phase
water relative permeability table. The keyword only applies the scaling
to the drainage capillary pressures tables, unlike the SWL keyword that
applies the scaling to both the capillary pressure and relative
permeability tables. The keyword can be used with all grid types.

\begin{longtable}[]{@{}
  >{\raggedright\arraybackslash}p{(\linewidth - 6\tabcolsep) * \real{0.2500}}
  >{\raggedright\arraybackslash}p{(\linewidth - 6\tabcolsep) * \real{0.2500}}
  >{\raggedright\arraybackslash}p{(\linewidth - 6\tabcolsep) * \real{0.2500}}
  >{\raggedright\arraybackslash}p{(\linewidth - 6\tabcolsep) * \real{0.2500}}@{}}
\toprule\noalign{}
\begin{minipage}[b]{\linewidth}\raggedright
No.
\end{minipage} & \begin{minipage}[b]{\linewidth}\raggedright
Name
\end{minipage} & \begin{minipage}[b]{\linewidth}\raggedright
Description
\end{minipage} & \begin{minipage}[b]{\linewidth}\raggedright
Default
\end{minipage} \\
\midrule\noalign{}
\endhead
\bottomrule\noalign{}
\endlastfoot
Field & Metric & Laboratory & \\
1 & SWLPC & SWLPC is an array of real numbers assigning the connate
water saturation values to each cell in the model. The number of entries
should correspond to the NX x NY x NZ parameters on the DIMENS keyword.
If SWLPC is omitted from the input deck the values will be defaulted to
those on the SWL series of keywords. If the SWL series of keywords are
missing from the input deck then the values are taken from the cell
allocated capillary pressure table. Repeat counts may be used, for
example 30*0.03 & Taken from SGL or from the cell allocated capillary
pressure table. \\
dimensionless & dimensionless & dimensionless & \\
Notes: & & & \\
\end{longtable}

\emph{Table 8.188: SWLPC Keyword Description}

End-point scaling allows the entered relative permeability functions to
be re-scaled based on the saturation end-points allocated to each cell
by the SWL, SWCR, SWU, SGL, SGCR, SGU, SOWCR, and SOGCR saturation grid
arrays for the saturation end-points, and the KRG, KROG, KROW and KRW
relative permeability grid cell arrays for the relative permeability
end-point data. In addition end-point scaling may be directional
dependent in which case the directional dependent versions of the
aforementioned arrays should be used, that is SWLX, SWLY and SWLZ
instead of SWL or SWLPC. There is also the facility to make the
directional end-point scaling reversible or non-reversible and if the
non-reversible option is selected the non-reversible versions of the
aforementioned arrays should be used, that is SWLX, SWLX-, SWLY, SWLY-,
SWLZ and SWLZ-, instead of the SWL or SWLPC keywords.

Missing Some Functionality - Use with Caution.

\subsubsection{Example}\label{example-400}

\begin{verbatim}
--
--       DEFINE GRID BLOCK END-POINT SWLPC DATA FOR ALL CELLS
--       (FOR NX x NY x NZ = 300)
--
SWLPC
         300*0.150                                                              /
\end{verbatim}

The above example defines a constant connate water saturation of 0.15 to
all 300 cells in the model as defined by the DIMENS keyword in the
RUNSPEC section.

\subsection{SWOF -- Water-Oil Saturation Tables (Format Type
1)}\label{swof-water-oil-saturation-tables-format-type-1}

\begin{longtable}[]{@{}
  >{\raggedright\arraybackslash}p{(\linewidth - 14\tabcolsep) * \real{0.1250}}
  >{\raggedright\arraybackslash}p{(\linewidth - 14\tabcolsep) * \real{0.1250}}
  >{\raggedright\arraybackslash}p{(\linewidth - 14\tabcolsep) * \real{0.1250}}
  >{\raggedright\arraybackslash}p{(\linewidth - 14\tabcolsep) * \real{0.1250}}
  >{\raggedright\arraybackslash}p{(\linewidth - 14\tabcolsep) * \real{0.1250}}
  >{\raggedright\arraybackslash}p{(\linewidth - 14\tabcolsep) * \real{0.1250}}
  >{\raggedright\arraybackslash}p{(\linewidth - 14\tabcolsep) * \real{0.1250}}
  >{\raggedright\arraybackslash}p{(\linewidth - 14\tabcolsep) * \real{0.1250}}@{}}
\toprule\noalign{}
\begin{minipage}[b]{\linewidth}\raggedright
RUNSPEC
\end{minipage} & \begin{minipage}[b]{\linewidth}\raggedright
GRID
\end{minipage} & \begin{minipage}[b]{\linewidth}\raggedright
EDIT
\end{minipage} & \begin{minipage}[b]{\linewidth}\raggedright
PROPS
\end{minipage} & \begin{minipage}[b]{\linewidth}\raggedright
REGIONS
\end{minipage} & \begin{minipage}[b]{\linewidth}\raggedright
SOLUTION
\end{minipage} & \begin{minipage}[b]{\linewidth}\raggedright
SUMMARY
\end{minipage} & \begin{minipage}[b]{\linewidth}\raggedright
SCHEDULE
\end{minipage} \\
\midrule\noalign{}
\endhead
\bottomrule\noalign{}
\endlastfoot
\end{longtable}

\subsubsection{Description}\label{description-694}

The SWOF keyword defines the water and oil relative permeability and
water-oil capillary pressure data versus water saturation tables for
when water and oil are present in the input deck. This keyword should
only be used if water and oil are present in the run.

\begin{longtable}[]{@{}
  >{\raggedright\arraybackslash}p{(\linewidth - 6\tabcolsep) * \real{0.2500}}
  >{\raggedright\arraybackslash}p{(\linewidth - 6\tabcolsep) * \real{0.2500}}
  >{\raggedright\arraybackslash}p{(\linewidth - 6\tabcolsep) * \real{0.2500}}
  >{\raggedright\arraybackslash}p{(\linewidth - 6\tabcolsep) * \real{0.2500}}@{}}
\toprule\noalign{}
\begin{minipage}[b]{\linewidth}\raggedright
No.
\end{minipage} & \begin{minipage}[b]{\linewidth}\raggedright
Name
\end{minipage} & \begin{minipage}[b]{\linewidth}\raggedright
Description
\end{minipage} & \begin{minipage}[b]{\linewidth}\raggedright
Default
\end{minipage} \\
\midrule\noalign{}
\endhead
\bottomrule\noalign{}
\endlastfoot
Field & Metric & Laboratory & \\
1 & SWAT & A columnar vector of real monotonically increasing down the
column values starting from zero and terminating at one, that defines
the water saturation. The first entry is the connate water saturation
Swc and the last entry should be 1.0. & None \\
dimensionless & dimensionless & dimensionless & \\
2 & KRW & A columnar vector of real values that are either equal or
increasing down the column and that are greater than or equal to zero
and less than or equal to one that defines the water relative
permeability with respect to gas saturation. The first value in the
column should be zero. & None \\
dimensionless & dimensionless & dimensionless & \\
3 & KRO & A columnar vector of real values that are either equal or
decreasing down the column and that are greater than or equal to zero
and less than or equal to one that defines the oil relative permeability
with respect to oil and water saturation. When gas is active in the run,
the first entry the column, that is at krow(So = 1-Swc), must be the
same as the first entry in the corresponding SGOF or SLGOF table, that
is at krog(Sg = 0). The first value in the column should be one. &
None \\
dimensionless & dimensionless & dimensionless & \\
4 & PCWO & A columnar vector of real values that are either equal or
decreasing down the column that defines the water-oil relative capillary
pressure. If the SWATINIT keyword has been used to initialize the model
then columnar vector has to be strictly monotonically increasing. &
None \\
psia & bars & atm & \\
Notes: & & & \\
\end{longtable}

\emph{Table 8.189: SWOF Keyword Description}

\subsubsection{Example}\label{example-401}

The following example is based on NTSFUN equals two on the TABDIMS
keyword in the RUNSPEC section.

\begin{verbatim}
--
--       WATER-OIL RELATIVE PERMEABILITY TABLES (SWOF)
--
SWOF
--       SWAT       KRW       KROW      PCOW
--       FRAC                           PSIA
--       --------   --------  -------   -------
         0.200000   0.0000    0.9000    0.000000
         0.238616   0.0002    0.7664    0.000000
         0.245309   0.0004    0.7443    0.000000
         0.261989   0.0010    0.6907    0.000000
         0.303091   0.0044    0.5671    0.000000
         0.368269   0.0191    0.3962    0.000000
         0.435026   0.0519    0.2528    0.000000
         0.486387   0.0940    0.1643    0.000000
         0.550683   0.1725    0.0803    0.000000
         0.575342   0.2115    0.0559    0.000000
         0.599076   0.2542    0.0367    0.000000
         0.621294   0.2991    0.0223    0.000000
         0.642171   0.3458    0.0120    0.000000
         0.658984   0.3868    0.0061    0.000000
         0.671123   0.4183    0.0030    0.000000
         0.679268   0.4403    0.0015    0.000000
         0.684963   0.4562    0.0008    0.000000
         0.688893   0.4674    0.0004    0.000000
         0.692025   0.4765    0.0002    0.000000
         0.694641   0.4841    0.0001    0.000000
         0.696976   0.4910    0.0000    0.000000
         0.699099   0.4973    0.0000    0.000000
         0.700000   0.5000    0.0000    0.000000
         1.000000   0.9000    0.0000    0.000000           / TABLE NO. 01
--       --------   --------  -------   -------
         0.200000   0.0000    0.9000    0.000000
         0.238616   0.0002    0.7664    0.000000
         0.245309   0.0004    0.7443    0.000000
         0.261989   0.0010    0.6907    0.000000
         0.303091   0.0044    0.5671    0.000000
         0.368269   0.0191    0.3962    0.000000
         0.435026   0.0519    0.2528    0.000000
         0.486387   0.0940    0.1643    0.000000
         0.550683   0.1725    0.0803    0.000000
         0.575342   0.2115    0.0559    0.000000
         0.599076   0.2542    0.0367    0.000000
         0.621294   0.2991    0.0223    0.000000
         0.642171   0.3458    0.0120    0.000000
         0.658984   0.3868    0.0061    0.000000
         0.671123   0.4183    0.0030    0.000000
         0.679268   0.4403    0.0015    0.000000
         0.684963   0.4562    0.0008    0.000000
         0.688893   0.4674    0.0004    0.000000
         0.692025   0.4765    0.0002    0.000000
         0.694641   0.4841    0.0001    0.000000
         0.696976   0.4910    0.0000    0.000000
         0.699099   0.4973    0.0000    0.000000
         0.700000   0.5000    0.0000    0.000000
         1.000000   0.9000    0.0000    0.000000           / TABLE NO. 01
\end{verbatim}

The example defines two SWOF tables for use when water and oil are
present in the run. In the tables the water-oil capillary pressure data
has been set to zero.

\subsection{SWOFLET -- Water-Oil LET Relative Permeability
Functions}\label{swoflet-water-oil-let-relative-permeability-functions}

\begin{longtable}[]{@{}
  >{\raggedright\arraybackslash}p{(\linewidth - 14\tabcolsep) * \real{0.1250}}
  >{\raggedright\arraybackslash}p{(\linewidth - 14\tabcolsep) * \real{0.1250}}
  >{\raggedright\arraybackslash}p{(\linewidth - 14\tabcolsep) * \real{0.1250}}
  >{\raggedright\arraybackslash}p{(\linewidth - 14\tabcolsep) * \real{0.1250}}
  >{\raggedright\arraybackslash}p{(\linewidth - 14\tabcolsep) * \real{0.1250}}
  >{\raggedright\arraybackslash}p{(\linewidth - 14\tabcolsep) * \real{0.1250}}
  >{\raggedright\arraybackslash}p{(\linewidth - 14\tabcolsep) * \real{0.1250}}
  >{\raggedright\arraybackslash}p{(\linewidth - 14\tabcolsep) * \real{0.1250}}@{}}
\toprule\noalign{}
\begin{minipage}[b]{\linewidth}\raggedright
RUNSPEC
\end{minipage} & \begin{minipage}[b]{\linewidth}\raggedright
GRID
\end{minipage} & \begin{minipage}[b]{\linewidth}\raggedright
EDIT
\end{minipage} & \begin{minipage}[b]{\linewidth}\raggedright
PROPS
\end{minipage} & \begin{minipage}[b]{\linewidth}\raggedright
REGIONS
\end{minipage} & \begin{minipage}[b]{\linewidth}\raggedright
SOLUTION
\end{minipage} & \begin{minipage}[b]{\linewidth}\raggedright
SUMMARY
\end{minipage} & \begin{minipage}[b]{\linewidth}\raggedright
SCHEDULE
\end{minipage} \\
\midrule\noalign{}
\endhead
\bottomrule\noalign{}
\endlastfoot
\end{longtable}

\subsubsection{Description}\label{description-695}

SWOFLET defines the relative permeability and capillary pressure
parameters for the water-oil LET family of models. Both the oil and
water phases should be made active in the model via the OIL and WATER
keywords in the RUNSPEC section. See section 8.2.6Saturation Table
Generation - LET Functions and Lomeland et al.~{[}Lomeland F., Ebeltoft
E. and Thomas W.H., 2005. A New Versatile Relative Permeability
Correlation. Paper SCA2005-32 presented at the International Symposium
of the Society of Core Analysts held in Toronto, Canada, 21-25 August,
2005.{]}, {[}Lomeland F. and Ebeltoft E., 2008. A New Versatile
Capillary Pressure Correlation. Paper SCA2008-08 presented at the
International Symposium of the Society of Core Analysts held in Abu
Dhabi, UAE, 29 Oct.~-- 2 Nov., 2008.{]} {[}Lomeland F., Hasanov B.,
Ebeltoft E. and Berge M., 2012. A Versatile Representation of Up-scaled
Relative Permeability for Field Applications. Paper SPE 154487-MS
presented at the EAGE Annual Conference \& Exhibition incorporating SPE
Europec held in Copenhagen, Denmark, 4-7 June 2012.{]} and {[}Lomeland
F., 2018.Overview Of The Let Family Of Versatile Correlations For Flow
Functions. Paper SCA2018-056 presented at the International Symposium of
the Society of Core Analysts held in Trondheim, Norway, 27-30 August
2018.{]} for further information on the model.

The keyword is used as a replacement for the SWOF keyword for
three-phase oil-gas-water systems, and the LET series of keywords cannot
be combined with the standard set of relative permeability keywords.

\begin{tcolorbox}[enhanced jigsaw, title=\textcolor{quarto-callout-note-color}{\faInfo}\hspace{0.5em}{Note}, coltitle=black, breakable, toprule=.15mm, colbacktitle=quarto-callout-note-color!10!white, leftrule=.75mm, opacitybacktitle=0.6, colback=white, bottomtitle=1mm, toptitle=1mm, titlerule=0mm, left=2mm, colframe=quarto-callout-note-color-frame, arc=.35mm, rightrule=.15mm, bottomrule=.15mm, opacityback=0]

This is an OPM Flow specific keyword and will therefore cause an error
in the commercial simulator.

\end{tcolorbox}

\begin{longtable}[]{@{}
  >{\raggedright\arraybackslash}p{(\linewidth - 6\tabcolsep) * \real{0.2500}}
  >{\raggedright\arraybackslash}p{(\linewidth - 6\tabcolsep) * \real{0.2500}}
  >{\raggedright\arraybackslash}p{(\linewidth - 6\tabcolsep) * \real{0.2500}}
  >{\raggedright\arraybackslash}p{(\linewidth - 6\tabcolsep) * \real{0.2500}}@{}}
\toprule\noalign{}
\begin{minipage}[b]{\linewidth}\raggedright
No.
\end{minipage} & \begin{minipage}[b]{\linewidth}\raggedright
Name
\end{minipage} & \begin{minipage}[b]{\linewidth}\raggedright
Description
\end{minipage} & \begin{minipage}[b]{\linewidth}\raggedright
Default
\end{minipage} \\
\midrule\noalign{}
\endhead
\bottomrule\noalign{}
\endlastfoot
Field & Metric & Laboratory & \\
1 & SWL & SWL is a real positive number less than one, that defines the
connate water saturation, that is the smallest water saturation in the
LET function. & 0,0 \\
dimensionless & dimensionless & dimensionless & \\
2 & SWCR & SWCR is a real positive number greater than or equal to SWL
and less than one, that defines the critical water saturation, that is
the largest water saturation for which the water relative permeability
is zero, Swirr in equations. & 0,0 \\
dimensionless & dimensionless & dimensionless & \\
3 & LWAT & LWAT is a real positive number that defines the LET Lower
empirical parameter Lw for the water phase with the associated oil phase
in the LET relative permeability equations. & 1,0 \\
dimensionless & dimensionless & dimensionless & \\
4 & EWAT & EWAT is a real positive number that defines the LET Elevation
empirical parameter Ew for the water phase with the associated oil phase
in the LET relative permeability equations. & 1,0 \\
dimensionless & dimensionless & dimensionless & \\
5 & TWAT & TWAT is a real positive number that defines the LET Top
empirical parameter Tw for the water phase with the associated oil phase
in the LET relative permeability equations. & 1,0 \\
dimensionless & dimensionless & dimensionless & \\
6 & KRTWAT & KRTWAT is a real positive number less than one, that
defines the relative permeability of water at the maximum water
saturation (normally the maximum water saturation is one) Krwt in the
LET relative permeability equations. & 1,0 \\
dimensionless & dimensionless & dimensionless & \\
7 & SORW & SORW is a real positive number less than one that defines the
residual oil saturation in an oil-water system in the LET equation. &
0,0 \\
dimensionless & dimensionless & dimensionless & \\
8 & SOWCR & SOWCR is a real positive number less than one that defines
critical oil-in-water saturation, that is the largest oil saturation for
which the oil relative permeability is zero in an oil-water system. &
0,0 \\
dimensionless & dimensionless & dimensionless & \\
9 & LOIL & LOIL is a real positive number that defines the LET Lower
empirical parameter Lo for the oil phase with the associated water phase
in the LET relative permeability equationss & 1,0 \\
dimensionless & dimensionless & dimensionless & \\
10 & EOIL & EOIL is a real positive number that defines the LET
Elevation empirical parameter Eo for the oil phase with the associated
water phase in the LET relative permeability equations. & 1,0 \\
dimensionless & dimensionless & dimensionless & \\
11 & TOIL & TOIL is a real positive number that defines the LET Top
empirical parameter To for the oil phase with the associated water phase
in the LET relative permeability equations. & 1,0 \\
dimensionless & dimensionless & dimensionless & \\
12 & KRTOIL & KRTOIL is a real positive number less than or equal to
one, that defines the relative permeability of oil at the residual oil
saturation, Krot in the LET relative permeability equations. & 1,0 \\
dimensionless & dimensionless & dimensionless & \\
13 & LPC & LPC is a real positive number that defines the water-oil LET
Lower empirical parameter L in the LET capillary pressure equation. &
1,0 \\
dimensionless & dimensionless & dimensionless & \\
14 & EPC & EPC is a real positive number that defines the water-oil LET
Elevation empirical parameter E in the LET capillary pressure equation.
& 1,0 \\
dimensionless & dimensionless & dimensionless & \\
15 & TPC & TPC is a real positive number that defines the water-oil LET
empirical parameter T in the LET capillary pressure equations. & 1,0 \\
dimensionless & dimensionless & dimensionless & \\
16 & PCIR & PCIR is a real positive number that defines the water-oil
capillary pressure at connate water saturation (SWL) in the LET
capillary pressure equations. & 0,0 \\
psia & bars & atm & \\
17 & PCIT & PCIT is a real positive number that defines the water-oil
threshold capillary pressure at the maximum water saturation in the LET
capillary pressure equations. & 0,0 \\
psia & bars & atm & \\
Notes: & & & \\
\end{longtable}

\emph{Table 8.190: SWOFLET Keyword Description}

Note there a two versions of the LET functions, LET255 for two-phase
flowing conditions and LETx {[}Lomeland F. and Ebeltoft E., 2013.
Versatile Three-phase Correlations for Relative Permeability and
Capillary Pressure. Paper SCA2013-034 presented at the International
Symposium of the Society of Core Analysts held in Napa Valley,
California, USA, 16-19 September, 2013.{]} for three-phase flowing
conditions. This keyword implements the LET version for an oil-water
system.

The functions are dependent on the drainage and imbibition cycle of the
wetting phase as well as drainage and inhibition cycle number, since a
reservoir may undergo several flooding events. To account for this the
system defines the flooding event using the three saturations: Sw, So,
and Sg together with the state of the three saturations during the
flooding event. The saturation state can be Increasing, Decreasing, or
Constant, for a given flooding event cycle number (n). Thus, Sw(D),
So(I), Sg(C)1or DIC1, means the water phase is decreasing, the oil phase
is increasing and the gas phase is constant for the primary or first
cycle (n equals one). This is case for when oil is migrating into the
reservoir rock and displacing the initial water contained with the
reservoir.

\begin{tcolorbox}[enhanced jigsaw, title=\textcolor{quarto-callout-note-color}{\faInfo}\hspace{0.5em}{Note}, coltitle=black, breakable, toprule=.15mm, colbacktitle=quarto-callout-note-color!10!white, leftrule=.75mm, opacitybacktitle=0.6, colback=white, bottomtitle=1mm, toptitle=1mm, titlerule=0mm, left=2mm, colframe=quarto-callout-note-color-frame, arc=.35mm, rightrule=.15mm, bottomrule=.15mm, opacityback=0]

All the LET parameters are dependent on the flooding event and flooding
cycle, and thus are expected vary as such. To be clear, the values of
SWCR, Lo, Lw etc. should be different for each flooding cycle.

\end{tcolorbox}

See also the SGOFLET-- Gas-Oil LET Relative Permeability Functions, and
the SGWFLET -- Gas-Water LET Relative Permeability Functions keywords in
this section.

\subsubsection{Example}\label{example-402}

The following example uses the SWOFLET keyword to define two relative
oi-water relative permeability tables, based on NTSFUN equals two on the
TABDIMS keyword in the RUNSPEC section.

\begin{verbatim}

--
--       SWOFLET – WATER-OIL LET REL. PERMEABILITY FUNCTIONS (OPM FLOW KEYWORD)
--
SWOFLET
--       SWL       SWCR     L-WAT     E-WAT     T-WAT     KRT-WAT
--       SOR       SOWCR    L-OIL     E-OIL     T-OIL     KRT-OIL
--       L-PC      E-PC     T-PC      PCIR      PCIT
--       -------   ------   -------   -------   -------   -------
         0.00000   0.0000   1.00000   1.00000   1.00000   1.00000
         0.00000   0.0000   1.00000   1.00000   1.00000   1.00000
         1*        1*       1*        1*        1*                 / TABLE NO. 01

         0.00000   0.0000   1.00000   1.00000   1.00000   1.00000
         0.00000   0.0000   1.00000   1.00000   1.00000   1.00000
         1*        1*       1*        1*        1*                 / TABLE NO. 02
\end{verbatim}

Here the SWOFLET keyword parameters are all set to their default values.

\subsection{SWU -- End-Point Scaling Grid Cell Maximum Water
Saturation}\label{swu-end-point-scaling-grid-cell-maximum-water-saturation}

\begin{longtable}[]{@{}
  >{\raggedright\arraybackslash}p{(\linewidth - 14\tabcolsep) * \real{0.1250}}
  >{\raggedright\arraybackslash}p{(\linewidth - 14\tabcolsep) * \real{0.1250}}
  >{\raggedright\arraybackslash}p{(\linewidth - 14\tabcolsep) * \real{0.1250}}
  >{\raggedright\arraybackslash}p{(\linewidth - 14\tabcolsep) * \real{0.1250}}
  >{\raggedright\arraybackslash}p{(\linewidth - 14\tabcolsep) * \real{0.1250}}
  >{\raggedright\arraybackslash}p{(\linewidth - 14\tabcolsep) * \real{0.1250}}
  >{\raggedright\arraybackslash}p{(\linewidth - 14\tabcolsep) * \real{0.1250}}
  >{\raggedright\arraybackslash}p{(\linewidth - 14\tabcolsep) * \real{0.1250}}@{}}
\toprule\noalign{}
\begin{minipage}[b]{\linewidth}\raggedright
RUNSPEC
\end{minipage} & \begin{minipage}[b]{\linewidth}\raggedright
GRID
\end{minipage} & \begin{minipage}[b]{\linewidth}\raggedright
EDIT
\end{minipage} & \begin{minipage}[b]{\linewidth}\raggedright
PROPS
\end{minipage} & \begin{minipage}[b]{\linewidth}\raggedright
REGIONS
\end{minipage} & \begin{minipage}[b]{\linewidth}\raggedright
SOLUTION
\end{minipage} & \begin{minipage}[b]{\linewidth}\raggedright
SUMMARY
\end{minipage} & \begin{minipage}[b]{\linewidth}\raggedright
SCHEDULE
\end{minipage} \\
\midrule\noalign{}
\endhead
\bottomrule\noalign{}
\endlastfoot
\end{longtable}

\subsubsection{Description}\label{description-696}

SWU defines the maximum water saturation for all the cells in the model
via an array when the end-point scaling option has been invoked via the
ENDSCALE keyword in the RUNSPEC section. The maximum water saturation is
defined as the maximum water saturation in a two-phase water relative
permeability table.

The keyword can be used with all grid types.

\begin{longtable}[]{@{}
  >{\raggedright\arraybackslash}p{(\linewidth - 6\tabcolsep) * \real{0.2500}}
  >{\raggedright\arraybackslash}p{(\linewidth - 6\tabcolsep) * \real{0.2500}}
  >{\raggedright\arraybackslash}p{(\linewidth - 6\tabcolsep) * \real{0.2500}}
  >{\raggedright\arraybackslash}p{(\linewidth - 6\tabcolsep) * \real{0.2500}}@{}}
\toprule\noalign{}
\begin{minipage}[b]{\linewidth}\raggedright
No.
\end{minipage} & \begin{minipage}[b]{\linewidth}\raggedright
Name
\end{minipage} & \begin{minipage}[b]{\linewidth}\raggedright
Description
\end{minipage} & \begin{minipage}[b]{\linewidth}\raggedright
Default
\end{minipage} \\
\midrule\noalign{}
\endhead
\bottomrule\noalign{}
\endlastfoot
Field & Metric & Laboratory & \\
1 & SWU & SWU is an array of real numbers assigning the maximum water
saturation values to each cell in the model. The number of entries
should correspond to the NX x NY x NZ parameters on the DIMENS keyword.
Repeat counts may be used, for example 30*0.70 & Taken from cell
allocated relative permeability table. \\
dimensionless & dimensionless & dimensionless & \\
Notes: & & & \\
\end{longtable}

\emph{Table 8.191: SWU Keyword Description}

End-point scaling allows the entered relative permeability functions to
be re-scaled based on the saturation end-points allocated to each cell
by the SWL, SWCR, SWU, SGL, SGCR, SGU, SOWCR, and SOGCR saturation grid
arrays for the saturation end-points, and the KRG, KRORG, KRORW and KRW
relative permeability grid cell arrays for the relative permeability
end-point data. In addition end-point scaling may be directional
dependent in which case the directional dependent versions of the
aforementioned arrays should be used, that is SWUX, SWUY and SWUZ
instead of SWU. There is also the facility to make the directional
end-point scaling reversible or non-reversible and if the non-reversible
option is selected the non-reversible versions of the aforementioned
arrays should be used, that is SWUX, SWUX-, SWUY, SWUY-, SWUZ and SWUZ-,
instead of the SWU keyword.

\subsubsection{Example}\label{example-403}

\begin{verbatim}
--
--       DEFINE GRID BLOCK END-POINT SWU DATA FOR ALL CELLS
--       (FOR NX x NY x NZ = 300)
--
SWU
         300*0.700                                                              /
\end{verbatim}

The above example defines a constant connate gas saturation of 0.70 to
all 300 cells in the model as defined by the DIMENS keyword in the
RUNSPEC section.

\subsection{TCRIT -- Critical
Temperatures}\label{tcrit-critical-temperatures}

\begin{longtable}[]{@{}
  >{\raggedright\arraybackslash}p{(\linewidth - 14\tabcolsep) * \real{0.1250}}
  >{\raggedright\arraybackslash}p{(\linewidth - 14\tabcolsep) * \real{0.1250}}
  >{\raggedright\arraybackslash}p{(\linewidth - 14\tabcolsep) * \real{0.1250}}
  >{\raggedright\arraybackslash}p{(\linewidth - 14\tabcolsep) * \real{0.1250}}
  >{\raggedright\arraybackslash}p{(\linewidth - 14\tabcolsep) * \real{0.1250}}
  >{\raggedright\arraybackslash}p{(\linewidth - 14\tabcolsep) * \real{0.1250}}
  >{\raggedright\arraybackslash}p{(\linewidth - 14\tabcolsep) * \real{0.1250}}
  >{\raggedright\arraybackslash}p{(\linewidth - 14\tabcolsep) * \real{0.1250}}@{}}
\toprule\noalign{}
\begin{minipage}[b]{\linewidth}\raggedright
RUNSPEC
\end{minipage} & \begin{minipage}[b]{\linewidth}\raggedright
GRID
\end{minipage} & \begin{minipage}[b]{\linewidth}\raggedright
EDIT
\end{minipage} & \begin{minipage}[b]{\linewidth}\raggedright
PROPS
\end{minipage} & \begin{minipage}[b]{\linewidth}\raggedright
REGIONS
\end{minipage} & \begin{minipage}[b]{\linewidth}\raggedright
SOLUTION
\end{minipage} & \begin{minipage}[b]{\linewidth}\raggedright
SUMMARY
\end{minipage} & \begin{minipage}[b]{\linewidth}\raggedright
SCHEDULE
\end{minipage} \\
\midrule\noalign{}
\endhead
\bottomrule\noalign{}
\endlastfoot
\end{longtable}

\subsubsection{Description}\label{description-697}

The \hyperref[REF_HEADING_KEYWORD_TCRIT]{TCRIT} keyword defines the
critical temperatures for each of the compositional components active in
the model and for each equation of state. The keyword should only be
used if the compositional mode has been requested using the COMPS
keyword in the RUNSPEC section.

OPM Flow does not currently support the general compositional modeling
formulation.

This keyword is not supported by OPM Flow but it will be parsed and its
data ignored.

\begin{longtable}[]{@{}
  >{\raggedright\arraybackslash}p{(\linewidth - 6\tabcolsep) * \real{0.2500}}
  >{\raggedright\arraybackslash}p{(\linewidth - 6\tabcolsep) * \real{0.2500}}
  >{\raggedright\arraybackslash}p{(\linewidth - 6\tabcolsep) * \real{0.2500}}
  >{\raggedright\arraybackslash}p{(\linewidth - 6\tabcolsep) * \real{0.2500}}@{}}
\toprule\noalign{}
\begin{minipage}[b]{\linewidth}\raggedright
No.
\end{minipage} & \begin{minipage}[b]{\linewidth}\raggedright
Name
\end{minipage} & \begin{minipage}[b]{\linewidth}\raggedright
Description
\end{minipage} & \begin{minipage}[b]{\linewidth}\raggedright
Default
\end{minipage} \\
\midrule\noalign{}
\endhead
\bottomrule\noalign{}
\endlastfoot
Field & Metric & Laboratory & \\
1 & \hyperref[REF_HEADING_KEYWORD_TCRIT]{TCRIT} & A series of real
numbers that define the critical temperatures for each of the
compositional components active in the model. & None \\
°R & K & K & \\
Notes: & & & \\
\end{longtable}

\emph{Table 8.3.340.1: \hyperref[REF_HEADING_KEYWORD_TCRIT]{TCRIT}
Keyword Description}

\subsubsection{Examples}\label{examples-44}

The following example defines the critical temperatures for each
component in a single three-component equation of state model.

\begin{verbatim}
--
-- Critical Temperatures
--
TCRIT
  547.4412 343.0152 665.802 /
\end{verbatim}

The following example defines the critical temperatures for each
component in two three-component equation of state models.

\begin{verbatim}
--
-- Critical Temperatures
--
TCRIT
  547.4412 343.0152 665.802 /
  547.4412 343.0152 665.802 /
\end{verbatim}

\subsection{TEMPNODE -- Temperature Table for Polymer Solution
Viscosity}\label{tempnode-temperature-table-for-polymer-solution-viscosity}

\begin{longtable}[]{@{}
  >{\raggedright\arraybackslash}p{(\linewidth - 14\tabcolsep) * \real{0.1250}}
  >{\raggedright\arraybackslash}p{(\linewidth - 14\tabcolsep) * \real{0.1250}}
  >{\raggedright\arraybackslash}p{(\linewidth - 14\tabcolsep) * \real{0.1250}}
  >{\raggedright\arraybackslash}p{(\linewidth - 14\tabcolsep) * \real{0.1250}}
  >{\raggedright\arraybackslash}p{(\linewidth - 14\tabcolsep) * \real{0.1250}}
  >{\raggedright\arraybackslash}p{(\linewidth - 14\tabcolsep) * \real{0.1250}}
  >{\raggedright\arraybackslash}p{(\linewidth - 14\tabcolsep) * \real{0.1250}}
  >{\raggedright\arraybackslash}p{(\linewidth - 14\tabcolsep) * \real{0.1250}}@{}}
\toprule\noalign{}
\begin{minipage}[b]{\linewidth}\raggedright
RUNSPEC
\end{minipage} & \begin{minipage}[b]{\linewidth}\raggedright
GRID
\end{minipage} & \begin{minipage}[b]{\linewidth}\raggedright
EDIT
\end{minipage} & \begin{minipage}[b]{\linewidth}\raggedright
PROPS
\end{minipage} & \begin{minipage}[b]{\linewidth}\raggedright
REGIONS
\end{minipage} & \begin{minipage}[b]{\linewidth}\raggedright
SOLUTION
\end{minipage} & \begin{minipage}[b]{\linewidth}\raggedright
SUMMARY
\end{minipage} & \begin{minipage}[b]{\linewidth}\raggedright
SCHEDULE
\end{minipage} \\
\midrule\noalign{}
\endhead
\bottomrule\noalign{}
\endlastfoot
\end{longtable}

\subsubsection{Description}\label{description-698}

This keyword defines the reservoir temperature table used to calculate
the polymer solution viscosity when the temperature option has been
activated by the TEMP keyword in the RUNSPEC section in the commercial
simulator. Naturally, the polymer option must also be activated by the
POLYMER keyword in the RUNSPEC section in order to use this keyword.

This keyword is not supported by OPM Flow but would change the results
if supported so the simulation will be stopped.

\subsection{TEMPTVD -- Activate Temperature Flux Limited Transport
Option}\label{temptvd-activate-temperature-flux-limited-transport-option}

\begin{longtable}[]{@{}
  >{\raggedright\arraybackslash}p{(\linewidth - 14\tabcolsep) * \real{0.1250}}
  >{\raggedright\arraybackslash}p{(\linewidth - 14\tabcolsep) * \real{0.1250}}
  >{\raggedright\arraybackslash}p{(\linewidth - 14\tabcolsep) * \real{0.1250}}
  >{\raggedright\arraybackslash}p{(\linewidth - 14\tabcolsep) * \real{0.1250}}
  >{\raggedright\arraybackslash}p{(\linewidth - 14\tabcolsep) * \real{0.1250}}
  >{\raggedright\arraybackslash}p{(\linewidth - 14\tabcolsep) * \real{0.1250}}
  >{\raggedright\arraybackslash}p{(\linewidth - 14\tabcolsep) * \real{0.1250}}
  >{\raggedright\arraybackslash}p{(\linewidth - 14\tabcolsep) * \real{0.1250}}@{}}
\toprule\noalign{}
\begin{minipage}[b]{\linewidth}\raggedright
RUNSPEC
\end{minipage} & \begin{minipage}[b]{\linewidth}\raggedright
GRID
\end{minipage} & \begin{minipage}[b]{\linewidth}\raggedright
EDIT
\end{minipage} & \begin{minipage}[b]{\linewidth}\raggedright
PROPS
\end{minipage} & \begin{minipage}[b]{\linewidth}\raggedright
REGIONS
\end{minipage} & \begin{minipage}[b]{\linewidth}\raggedright
SOLUTION
\end{minipage} & \begin{minipage}[b]{\linewidth}\raggedright
SUMMARY
\end{minipage} & \begin{minipage}[b]{\linewidth}\raggedright
SCHEDULE
\end{minipage} \\
\midrule\noalign{}
\endhead
\bottomrule\noalign{}
\endlastfoot
\end{longtable}

\subsubsection{Description}\label{description-699}

TEMPTVD activates the Temperature Flux Limited Transport option in the
commercial simulator, to reduce numerical dispersion for when either the
TEMP or THERMAL keywords in the RUNSPEC section have been declared.

This keyword is not supported by OPM Flow but would change the results
if supported so the simulation will be stopped.

\subsection{TEMPVD -- Define the Initial Reservoir Temperature versus
Depth
Table}\label{tempvd-define-the-initial-reservoir-temperature-versus-depth-table}

\begin{longtable}[]{@{}
  >{\raggedright\arraybackslash}p{(\linewidth - 14\tabcolsep) * \real{0.1250}}
  >{\raggedright\arraybackslash}p{(\linewidth - 14\tabcolsep) * \real{0.1250}}
  >{\raggedright\arraybackslash}p{(\linewidth - 14\tabcolsep) * \real{0.1250}}
  >{\raggedright\arraybackslash}p{(\linewidth - 14\tabcolsep) * \real{0.1250}}
  >{\raggedright\arraybackslash}p{(\linewidth - 14\tabcolsep) * \real{0.1250}}
  >{\raggedright\arraybackslash}p{(\linewidth - 14\tabcolsep) * \real{0.1250}}
  >{\raggedright\arraybackslash}p{(\linewidth - 14\tabcolsep) * \real{0.1250}}
  >{\raggedright\arraybackslash}p{(\linewidth - 14\tabcolsep) * \real{0.1250}}@{}}
\toprule\noalign{}
\begin{minipage}[b]{\linewidth}\raggedright
RUNSPEC
\end{minipage} & \begin{minipage}[b]{\linewidth}\raggedright
GRID
\end{minipage} & \begin{minipage}[b]{\linewidth}\raggedright
EDIT
\end{minipage} & \begin{minipage}[b]{\linewidth}\raggedright
PROPS
\end{minipage} & \begin{minipage}[b]{\linewidth}\raggedright
REGIONS
\end{minipage} & \begin{minipage}[b]{\linewidth}\raggedright
SOLUTION
\end{minipage} & \begin{minipage}[b]{\linewidth}\raggedright
SUMMARY
\end{minipage} & \begin{minipage}[b]{\linewidth}\raggedright
SCHEDULE
\end{minipage} \\
\midrule\noalign{}
\endhead
\bottomrule\noalign{}
\endlastfoot
\end{longtable}

\subsubsection{Description}\label{description-700}

This keyword defines the initial reservoir temperature versus depth
tables for each equilibration region. Note that the TEMPVD keyword is an
alias for RTEMPVD, and that both keywords are supported by OPM Flow, in
both the PROPS and SOLUTION sections, but are treated as being mutually
exclusive.

See TEMPVD - Define the Initial Reservoir Temperature versus Depth
Tables in the SOLUTION section, as the keyword is documented in the
SOLUTION section in this instance, but it can also be used in the PROPS
section by OPM Flow.

\subsection{THCO2MIX -- Specify Thermal Mixing
Models}\label{thco2mix-specify-thermal-mixing-models}

\begin{longtable}[]{@{}
  >{\raggedright\arraybackslash}p{(\linewidth - 14\tabcolsep) * \real{0.1250}}
  >{\raggedright\arraybackslash}p{(\linewidth - 14\tabcolsep) * \real{0.1250}}
  >{\raggedright\arraybackslash}p{(\linewidth - 14\tabcolsep) * \real{0.1250}}
  >{\raggedright\arraybackslash}p{(\linewidth - 14\tabcolsep) * \real{0.1250}}
  >{\raggedright\arraybackslash}p{(\linewidth - 14\tabcolsep) * \real{0.1250}}
  >{\raggedright\arraybackslash}p{(\linewidth - 14\tabcolsep) * \real{0.1250}}
  >{\raggedright\arraybackslash}p{(\linewidth - 14\tabcolsep) * \real{0.1250}}
  >{\raggedright\arraybackslash}p{(\linewidth - 14\tabcolsep) * \real{0.1250}}@{}}
\toprule\noalign{}
\begin{minipage}[b]{\linewidth}\raggedright
RUNSPEC
\end{minipage} & \begin{minipage}[b]{\linewidth}\raggedright
GRID
\end{minipage} & \begin{minipage}[b]{\linewidth}\raggedright
EDIT
\end{minipage} & \begin{minipage}[b]{\linewidth}\raggedright
PROPS
\end{minipage} & \begin{minipage}[b]{\linewidth}\raggedright
REGIONS
\end{minipage} & \begin{minipage}[b]{\linewidth}\raggedright
SOLUTION
\end{minipage} & \begin{minipage}[b]{\linewidth}\raggedright
SUMMARY
\end{minipage} & \begin{minipage}[b]{\linewidth}\raggedright
SCHEDULE
\end{minipage} \\
\midrule\noalign{}
\endhead
\bottomrule\noalign{}
\endlastfoot
\end{longtable}

\subsubsection{Description}\label{description-701}

The \hyperref[REF_HEADING_KEYWORD_THCO2MIX_8_3]{THCO2MIX} keyword
specifies the thermal mixing models for salt in the water phase, CO2 in
the liquid phase and vaporized water in gas phase.

This is an OPM Flow specific keyword that should only be used if the
CO2STORE keyword has been specified in the RUNSPEC section.

\begin{longtable}[]{@{}
  >{\raggedright\arraybackslash}p{(\linewidth - 6\tabcolsep) * \real{0.2500}}
  >{\raggedright\arraybackslash}p{(\linewidth - 6\tabcolsep) * \real{0.2500}}
  >{\raggedright\arraybackslash}p{(\linewidth - 6\tabcolsep) * \real{0.2500}}
  >{\raggedright\arraybackslash}p{(\linewidth - 6\tabcolsep) * \real{0.2500}}@{}}
\toprule\noalign{}
\begin{minipage}[b]{\linewidth}\raggedright
No.
\end{minipage} & \begin{minipage}[b]{\linewidth}\raggedright
Name
\end{minipage} & \begin{minipage}[b]{\linewidth}\raggedright
Description
\end{minipage} & \begin{minipage}[b]{\linewidth}\raggedright
Default
\end{minipage} \\
\midrule\noalign{}
\endhead
\bottomrule\noalign{}
\endlastfoot
1 & SALTMOD & A defined character string that specifies the thermal
mixing model for salt in the liquid phase, and should be set to one of
the following: & MICHAELIDES \\
2 & LIQMOD & A defined character string that specifies the thermal
mixing model for CO2 in the liquid phase, and should be set to one of
the following: & DUANSUN \\
3 & GASMOD & A defined character string that specifies the thermal
mixing model for vaporized water in the gas phase, and should be set to
one of the following: & NONE \\
Notes: & & & \\
\end{longtable}

\emph{Table 8.3.343.1:
\hyperref[REF_HEADING_KEYWORD_THCO2MIX_8_3]{THCO2MIX} Keyword
Description}

\subsubsection{Example}\label{example-404}

The following example specifies the default thermal mixing models for
salt in the liquid phase, CO2 in the liquid phase, and vaporized water
in the gas phase.

\begin{verbatim}
--
--       SPECIFY THERMAL MIXING MODELS
--
--       SALT      LIQUID    GAS
--       --------  --------  --------
THCO2MIX
      MICHAELIDES  DUANSUN   NONE    /
\end{verbatim}

\subsection{TLMIXPAR -- Define the Miscible Todd-Longstaff Mixing
Parameters}\label{tlmixpar-define-the-miscible-todd-longstaff-mixing-parameters}

\begin{longtable}[]{@{}
  >{\raggedright\arraybackslash}p{(\linewidth - 14\tabcolsep) * \real{0.1250}}
  >{\raggedright\arraybackslash}p{(\linewidth - 14\tabcolsep) * \real{0.1250}}
  >{\raggedright\arraybackslash}p{(\linewidth - 14\tabcolsep) * \real{0.1250}}
  >{\raggedright\arraybackslash}p{(\linewidth - 14\tabcolsep) * \real{0.1250}}
  >{\raggedright\arraybackslash}p{(\linewidth - 14\tabcolsep) * \real{0.1250}}
  >{\raggedright\arraybackslash}p{(\linewidth - 14\tabcolsep) * \real{0.1250}}
  >{\raggedright\arraybackslash}p{(\linewidth - 14\tabcolsep) * \real{0.1250}}
  >{\raggedright\arraybackslash}p{(\linewidth - 14\tabcolsep) * \real{0.1250}}@{}}
\toprule\noalign{}
\begin{minipage}[b]{\linewidth}\raggedright
RUNSPEC
\end{minipage} & \begin{minipage}[b]{\linewidth}\raggedright
GRID
\end{minipage} & \begin{minipage}[b]{\linewidth}\raggedright
EDIT
\end{minipage} & \begin{minipage}[b]{\linewidth}\raggedright
PROPS
\end{minipage} & \begin{minipage}[b]{\linewidth}\raggedright
REGIONS
\end{minipage} & \begin{minipage}[b]{\linewidth}\raggedright
SOLUTION
\end{minipage} & \begin{minipage}[b]{\linewidth}\raggedright
SUMMARY
\end{minipage} & \begin{minipage}[b]{\linewidth}\raggedright
SCHEDULE
\end{minipage} \\
\midrule\noalign{}
\endhead
\bottomrule\noalign{}
\endlastfoot
\end{longtable}

\subsubsection{Description}\label{description-702}

The TLMIXPAR keyword defines the Todd-Longstaff {[}Todd, M. and
Longstaff, W. ``The Development, Testing and Application of a Numerical
Simulator for Predicting Miscible Flood Performance,'' paper SPE 3484,
Journal of Canadian Petroleum Technology (1972) 24, No.~7, 874-882.{]}
mixing parameters, for when either the miscible or solvent options have
been activated by the MISCIBLE or SOLVENT keywords in the RUNSPEC
section. This keyword must be present in the input deck if the MISCIBLE
or SOLVENT keywords have been activated.

Note that If the POLYMER option has been activated by the POLYMER
keyword in the RUNSPEC section, then this keyword is ignored and the
mixing parameters are taken from the PLMIXPAR keyword instead.

\begin{longtable}[]{@{}
  >{\raggedright\arraybackslash}p{(\linewidth - 6\tabcolsep) * \real{0.2500}}
  >{\raggedright\arraybackslash}p{(\linewidth - 6\tabcolsep) * \real{0.2500}}
  >{\raggedright\arraybackslash}p{(\linewidth - 6\tabcolsep) * \real{0.2500}}
  >{\raggedright\arraybackslash}p{(\linewidth - 6\tabcolsep) * \real{0.2500}}@{}}
\toprule\noalign{}
\begin{minipage}[b]{\linewidth}\raggedright
No.
\end{minipage} & \begin{minipage}[b]{\linewidth}\raggedright
Name
\end{minipage} & \begin{minipage}[b]{\linewidth}\raggedright
Description
\end{minipage} & \begin{minipage}[b]{\linewidth}\raggedright
Default
\end{minipage} \\
\midrule\noalign{}
\endhead
\bottomrule\noalign{}
\endlastfoot
Field & Metric & Laboratory & \\
1 & TLMVIS & A real positive value that is greater than or equal to zero
and less than or equal to one, that defines the viscosity Todd-Longstaff
mixing parameter for each miscibility region. & None \\
dimensionless & dimensionless & dimensionless & \\
2 & TLMDEN & A real positive value that is greater than or equal to zero
and less than or equal to one, that defines the density Todd-Longstaff
mixing parameter for each miscibility region. & The same value as
entered for TLMVIS \\
dimensionless & dimensionless & dimensionless & \\
Notes: & & & \\
\end{longtable}

\emph{Table 8.192: TLMIXPAR Keyword Description}

\subsubsection{Example}\label{example-405}

\begin{verbatim}
--
--       TODD-LONGSTAFF MIXING PARAMETERS
--
TLMIXPAR
--       TLM        TLM
--       VISCOS     DENSITY
--       -------    --------
         0.3500      0.3500                                / TABLE NO. 01
         0.2500      1*                                    / TABLE NO. 02
         0.6500      0.7500                                / TABLE NO. 03
\end{verbatim}

The above example defines three Todd-Longstaff mixing parameter data
sets, based on the NTMISC variable on the MISCIBLE keyword in the
RUNSPEC section being equal to three.

\subsection{TOLCRIT -- Define The Critical Saturation
Tolerance}\label{tolcrit-define-the-critical-saturation-tolerance}

\begin{longtable}[]{@{}
  >{\raggedright\arraybackslash}p{(\linewidth - 14\tabcolsep) * \real{0.1250}}
  >{\raggedright\arraybackslash}p{(\linewidth - 14\tabcolsep) * \real{0.1250}}
  >{\raggedright\arraybackslash}p{(\linewidth - 14\tabcolsep) * \real{0.1250}}
  >{\raggedright\arraybackslash}p{(\linewidth - 14\tabcolsep) * \real{0.1250}}
  >{\raggedright\arraybackslash}p{(\linewidth - 14\tabcolsep) * \real{0.1250}}
  >{\raggedright\arraybackslash}p{(\linewidth - 14\tabcolsep) * \real{0.1250}}
  >{\raggedright\arraybackslash}p{(\linewidth - 14\tabcolsep) * \real{0.1250}}
  >{\raggedright\arraybackslash}p{(\linewidth - 14\tabcolsep) * \real{0.1250}}@{}}
\toprule\noalign{}
\begin{minipage}[b]{\linewidth}\raggedright
RUNSPEC
\end{minipage} & \begin{minipage}[b]{\linewidth}\raggedright
GRID
\end{minipage} & \begin{minipage}[b]{\linewidth}\raggedright
EDIT
\end{minipage} & \begin{minipage}[b]{\linewidth}\raggedright
PROPS
\end{minipage} & \begin{minipage}[b]{\linewidth}\raggedright
REGIONS
\end{minipage} & \begin{minipage}[b]{\linewidth}\raggedright
SOLUTION
\end{minipage} & \begin{minipage}[b]{\linewidth}\raggedright
SUMMARY
\end{minipage} & \begin{minipage}[b]{\linewidth}\raggedright
SCHEDULE
\end{minipage} \\
\midrule\noalign{}
\endhead
\bottomrule\noalign{}
\endlastfoot
\end{longtable}

\subsubsection{Description}\label{description-703}

Critical fluid saturations are determined from the relative permeability
tables, that is the last saturation in a relative permeability table
where the relative permeability of a phase is set equal to zero. Since
floating-point numbers (as implemented in computers) are never exact,
one cannot compare floating point numbers for exact equality. Thus, this
keyword defines a value below which is considered equivalent to zero in
determining the critical saturation for a phase.

\begin{longtable}[]{@{}
  >{\raggedright\arraybackslash}p{(\linewidth - 6\tabcolsep) * \real{0.2500}}
  >{\raggedright\arraybackslash}p{(\linewidth - 6\tabcolsep) * \real{0.2500}}
  >{\raggedright\arraybackslash}p{(\linewidth - 6\tabcolsep) * \real{0.2500}}
  >{\raggedright\arraybackslash}p{(\linewidth - 6\tabcolsep) * \real{0.2500}}@{}}
\toprule\noalign{}
\begin{minipage}[b]{\linewidth}\raggedright
No.
\end{minipage} & \begin{minipage}[b]{\linewidth}\raggedright
Name
\end{minipage} & \begin{minipage}[b]{\linewidth}\raggedright
Description
\end{minipage} & \begin{minipage}[b]{\linewidth}\raggedright
Default
\end{minipage} \\
\midrule\noalign{}
\endhead
\bottomrule\noalign{}
\endlastfoot
Field & Metric & Laboratory & \\
1 & TOLCRIT & TOLCRIT is a real positive number greater than zero and
less than one that defines the critical saturation tolerance used to
determine the critical saturation of a fluid in the relative
permeability tables. The default value of 1 x 10-6 means that relative
permeabilty values less than this value will be treated as being equal
to zero. & 1 x 10-6 \\
dimensionless & dimensionless & dimensionless & \\
Notes: & & & \\
\end{longtable}

\emph{Table 8.193: TOLCRIT Keyword Description}

See also section 8.2.4 Saturation Tables (Relative Permeability and
Capillary Pressure Tables) for a description of the relative
permeability tables and the various end-point definitions, including
oil, water and gas critical saturations.

\subsubsection{Example}\label{example-406}

\begin{verbatim}
---
--       SET THE CRITICAL SATURATION TOLERANCE
--
TOLCRIT
         1.0E-6                                                                /
\end{verbatim}

The above example defines the critical saturation tolerance to be the
default value of 1 x 10-6.

\subsection{TPAMEPS -- Volumetric Strain versus Coal Gas Concentration
Tables}\label{tpameps-volumetric-strain-versus-coal-gas-concentration-tables}

\begin{longtable}[]{@{}
  >{\raggedright\arraybackslash}p{(\linewidth - 14\tabcolsep) * \real{0.1250}}
  >{\raggedright\arraybackslash}p{(\linewidth - 14\tabcolsep) * \real{0.1250}}
  >{\raggedright\arraybackslash}p{(\linewidth - 14\tabcolsep) * \real{0.1250}}
  >{\raggedright\arraybackslash}p{(\linewidth - 14\tabcolsep) * \real{0.1250}}
  >{\raggedright\arraybackslash}p{(\linewidth - 14\tabcolsep) * \real{0.1250}}
  >{\raggedright\arraybackslash}p{(\linewidth - 14\tabcolsep) * \real{0.1250}}
  >{\raggedright\arraybackslash}p{(\linewidth - 14\tabcolsep) * \real{0.1250}}
  >{\raggedright\arraybackslash}p{(\linewidth - 14\tabcolsep) * \real{0.1250}}@{}}
\toprule\noalign{}
\begin{minipage}[b]{\linewidth}\raggedright
RUNSPEC
\end{minipage} & \begin{minipage}[b]{\linewidth}\raggedright
GRID
\end{minipage} & \begin{minipage}[b]{\linewidth}\raggedright
EDIT
\end{minipage} & \begin{minipage}[b]{\linewidth}\raggedright
PROPS
\end{minipage} & \begin{minipage}[b]{\linewidth}\raggedright
REGIONS
\end{minipage} & \begin{minipage}[b]{\linewidth}\raggedright
SOLUTION
\end{minipage} & \begin{minipage}[b]{\linewidth}\raggedright
SUMMARY
\end{minipage} & \begin{minipage}[b]{\linewidth}\raggedright
SCHEDULE
\end{minipage} \\
\midrule\noalign{}
\endhead
\bottomrule\noalign{}
\endlastfoot
\end{longtable}

\subsubsection{Description}\label{description-704}

TPAMEPS defines the volumetric strain versus coal gas concentration
tables, for when the Coal Bed Methane option has been activated via the
COAL keyword, and PALM-MAN has been declared for the ROCKOPT variable on
the ROCKCOMP keyword; both keywords are in the RUNSPEC section. The
Palmer-Mansoorii {[}Palmer, I. and Mansoori, J. ``How Permeability
Depends on Stress and Pore Pressure in Coalbeds: A New Model,'' paper
SPE 52607, SPE Reservoir Evaluation \& Engineering (1998) 1, No.~6,
539-544.{]} and {[}Clarkson, C.R., Pan, Z., Palmer, I. and Harpalani, S.
``Predicting Sorption-Induced Strain and Permeability Increase With
Depletion for Coalbed-Methane Reservoirs'', SPE 114778-PA, SPE Journal
(2010) 15, No.~1, 152--159.{]} rock model is used to calculate the
impact on pore volume and permeability due to rock compaction.

See also the ROCKPARMA keyword in the PROPS section that defines the
Palmer-Mansoori parameters.

This keyword is not supported by OPM Flow but would change the results
if supported so the simulation will be stopped.

\subsection{TPAMEPSS -- Volumetric Strain versus Coal Solvent
Concentration
Tables}\label{tpamepss-volumetric-strain-versus-coal-solvent-concentration-tables}

\begin{longtable}[]{@{}
  >{\raggedright\arraybackslash}p{(\linewidth - 14\tabcolsep) * \real{0.1250}}
  >{\raggedright\arraybackslash}p{(\linewidth - 14\tabcolsep) * \real{0.1250}}
  >{\raggedright\arraybackslash}p{(\linewidth - 14\tabcolsep) * \real{0.1250}}
  >{\raggedright\arraybackslash}p{(\linewidth - 14\tabcolsep) * \real{0.1250}}
  >{\raggedright\arraybackslash}p{(\linewidth - 14\tabcolsep) * \real{0.1250}}
  >{\raggedright\arraybackslash}p{(\linewidth - 14\tabcolsep) * \real{0.1250}}
  >{\raggedright\arraybackslash}p{(\linewidth - 14\tabcolsep) * \real{0.1250}}
  >{\raggedright\arraybackslash}p{(\linewidth - 14\tabcolsep) * \real{0.1250}}@{}}
\toprule\noalign{}
\begin{minipage}[b]{\linewidth}\raggedright
RUNSPEC
\end{minipage} & \begin{minipage}[b]{\linewidth}\raggedright
GRID
\end{minipage} & \begin{minipage}[b]{\linewidth}\raggedright
EDIT
\end{minipage} & \begin{minipage}[b]{\linewidth}\raggedright
PROPS
\end{minipage} & \begin{minipage}[b]{\linewidth}\raggedright
REGIONS
\end{minipage} & \begin{minipage}[b]{\linewidth}\raggedright
SOLUTION
\end{minipage} & \begin{minipage}[b]{\linewidth}\raggedright
SUMMARY
\end{minipage} & \begin{minipage}[b]{\linewidth}\raggedright
SCHEDULE
\end{minipage} \\
\midrule\noalign{}
\endhead
\bottomrule\noalign{}
\endlastfoot
\end{longtable}

\subsubsection{Description}\label{description-705}

TPAMEPSS defines the volumetric strain versus coal solvent concentration
tables, for when the Coal Bed Methane option has been activated via the
COAL keyword, and PALM-MAN has been declared for the ROCKOPT variable on
the ROCKCOMP keyword; both keywords are in the RUNSPEC section. The
Palmer-Mansoorii {[}Palmer, I. and Mansoori, J. ``How Permeability
Depends on Stress and Pore Pressure in Coalbeds: A New Model,'' paper
SPE 52607, SPE Reservoir Evaluation \& Engineering (1998) 1, No.~6,
539-544.{]} and {[}Clarkson, C.R., Pan, Z., Palmer, I. and Harpalani, S.
``Predicting Sorption-Induced Strain and Permeability Increase With
Depletion for Coalbed-Methane Reservoirs'', SPE 114778-PA, SPE Journal
(2010) 15, No.~1, 152--159.{]} rock model is used to calculate the
impact on pore volume and permeability due to rock compaction.

See also the ROCKPARMA keyword in the PROPS section that defines the
Palmer-Mansoori parameters.

This keyword is not supported by OPM Flow but would change the results
if supported so the simulation will be stopped.

\subsection{TRACER -- Define Passive Tracer
Variables}\label{tracer-define-passive-tracer-variables}

\begin{longtable}[]{@{}
  >{\raggedright\arraybackslash}p{(\linewidth - 14\tabcolsep) * \real{0.1250}}
  >{\raggedright\arraybackslash}p{(\linewidth - 14\tabcolsep) * \real{0.1250}}
  >{\raggedright\arraybackslash}p{(\linewidth - 14\tabcolsep) * \real{0.1250}}
  >{\raggedright\arraybackslash}p{(\linewidth - 14\tabcolsep) * \real{0.1250}}
  >{\raggedright\arraybackslash}p{(\linewidth - 14\tabcolsep) * \real{0.1250}}
  >{\raggedright\arraybackslash}p{(\linewidth - 14\tabcolsep) * \real{0.1250}}
  >{\raggedright\arraybackslash}p{(\linewidth - 14\tabcolsep) * \real{0.1250}}
  >{\raggedright\arraybackslash}p{(\linewidth - 14\tabcolsep) * \real{0.1250}}@{}}
\toprule\noalign{}
\begin{minipage}[b]{\linewidth}\raggedright
RUNSPEC
\end{minipage} & \begin{minipage}[b]{\linewidth}\raggedright
GRID
\end{minipage} & \begin{minipage}[b]{\linewidth}\raggedright
EDIT
\end{minipage} & \begin{minipage}[b]{\linewidth}\raggedright
PROPS
\end{minipage} & \begin{minipage}[b]{\linewidth}\raggedright
REGIONS
\end{minipage} & \begin{minipage}[b]{\linewidth}\raggedright
SOLUTION
\end{minipage} & \begin{minipage}[b]{\linewidth}\raggedright
SUMMARY
\end{minipage} & \begin{minipage}[b]{\linewidth}\raggedright
SCHEDULE
\end{minipage} \\
\midrule\noalign{}
\endhead
\bottomrule\noalign{}
\endlastfoot
\end{longtable}

\subsubsection{Description}\label{description-706}

The TRACER keyword defines a series of passive tracers that are
associated with a phase (oil, water, or gas) in the model. The maximum
number of tracers for each phase are declared on the TRACERS keyword in
the RUNSPEC section.

\begin{longtable}[]{@{}
  >{\raggedright\arraybackslash}p{(\linewidth - 6\tabcolsep) * \real{0.2500}}
  >{\raggedright\arraybackslash}p{(\linewidth - 6\tabcolsep) * \real{0.2500}}
  >{\raggedright\arraybackslash}p{(\linewidth - 6\tabcolsep) * \real{0.2500}}
  >{\raggedright\arraybackslash}p{(\linewidth - 6\tabcolsep) * \real{0.2500}}@{}}
\toprule\noalign{}
\begin{minipage}[b]{\linewidth}\raggedright
No.
\end{minipage} & \begin{minipage}[b]{\linewidth}\raggedright
Name
\end{minipage} & \begin{minipage}[b]{\linewidth}\raggedright
Description
\end{minipage} & \begin{minipage}[b]{\linewidth}\raggedright
Default
\end{minipage} \\
\midrule\noalign{}
\endhead
\bottomrule\noalign{}
\endlastfoot
Field & Metric & Laboratory & \\
1 & NAME & A three letter character string defining the tracer's name.
NAME is used by the TNUM keyword in the REGIONS section, and unlike
other keywords, the TNUM keyword itself must be concatenated with the
phase and the name of the tracer defined by NAME. Similarly for the TVDP
keyword in the SOLUTION section, where the TVDP keyword itself must be
concatenated with the either letter F (for free) or S (for solution),
followed by the name of the tracer defined by NAME. Note it is best to
void names beginning with the letters F, S, and T as these names may
create naming issues in post-processing software. & None \\
2 & PHASE & A three letter character string that defines the tracer
given by NAME to a particular fluid phase. The character should be set
to OIL, WAT or GAS. & None \\
3 & UNITS & The units for the tracer. This can be any unit but should be
consistent with values entered via the TBLK and TVDP keywords in the
SOLUTION section. The default values are the same as the PHASE in the
model. & Same as the phases in the model \\
Liquid: stb Gas: Mscf & Liquid: sm3 Gas: sm3 & Liquid: scc Gas: scc & \\
4 & SOLPHASE & A three or four letter character string defining the
partitioned tracer's solution phase. The character string should be set
to OIL, WAT, GAS or MULT. Note that SOLPHASE only needs to be defined if
the partitioned tracer option has been activate with the PARTTRAC
keyword in the RUNSPEC section. & None \\
5 & KPNUM & The table number to be used with the partitioned tracers
defined by the PARTTRAC, TRACERKP and TRACERKM keywords. Note that KPNUM
only needs to be defined if the partitioned tracer option has been
activate with the PARTTRAC keyword in the RUNSPEC section. & None \\
6 & ADSPHASE & A three letter character string defining the phase used
for the adsorption calculation for when the MULT option has been for
SOLPHASE. The character string should be set to OIL, WAT, GAS or ALL.
Note that ADSPHASE only needs to be defined if the partitioned tracer
option has been activate with the PARTTRAC keyword in the RUNSPEC
section. & None \\
Notes: & & & \\
\end{longtable}

\emph{Table 8.194: TRACER Keyword Description}

See also the TNUM keyword in the REGION section that defines tracer
regions, and the TVDP keyword in the SOLUTION section that sets the
initial tracer saturation for all the cells as a function of depth.

\subsubsection{Example}\label{example-407}

\begin{verbatim}
--
--       DEFINE TRACER NAMES
--
--       TRACER   TRACER
--       NAME     PHASE
--       ------   ------
TRACER
        'IGS'     'GAS'                                    / GAS INJECTOR
        'DGS'     'GAS'                                    / DISOLVED GAS
        'IW1'     'WAT'                                    / WAT INJECTOR 1
        'Iw2'     'WAT'                                    / WAT INJECTOR 2
/
\end{verbatim}

The above example defines four passive tracers one for a gas injection
well, one for tracking the dissolved gas, and two to track the injected
water from two water injection wells.

\subsection{TRACERKM -- Multi-Partitioned Tracer Option K(P)
Tables}\label{tracerkm-multi-partitioned-tracer-option-kp-tables}

\begin{longtable}[]{@{}
  >{\raggedright\arraybackslash}p{(\linewidth - 14\tabcolsep) * \real{0.1250}}
  >{\raggedright\arraybackslash}p{(\linewidth - 14\tabcolsep) * \real{0.1250}}
  >{\raggedright\arraybackslash}p{(\linewidth - 14\tabcolsep) * \real{0.1250}}
  >{\raggedright\arraybackslash}p{(\linewidth - 14\tabcolsep) * \real{0.1250}}
  >{\raggedright\arraybackslash}p{(\linewidth - 14\tabcolsep) * \real{0.1250}}
  >{\raggedright\arraybackslash}p{(\linewidth - 14\tabcolsep) * \real{0.1250}}
  >{\raggedright\arraybackslash}p{(\linewidth - 14\tabcolsep) * \real{0.1250}}
  >{\raggedright\arraybackslash}p{(\linewidth - 14\tabcolsep) * \real{0.1250}}@{}}
\toprule\noalign{}
\begin{minipage}[b]{\linewidth}\raggedright
RUNSPEC
\end{minipage} & \begin{minipage}[b]{\linewidth}\raggedright
GRID
\end{minipage} & \begin{minipage}[b]{\linewidth}\raggedright
EDIT
\end{minipage} & \begin{minipage}[b]{\linewidth}\raggedright
PROPS
\end{minipage} & \begin{minipage}[b]{\linewidth}\raggedright
REGIONS
\end{minipage} & \begin{minipage}[b]{\linewidth}\raggedright
SOLUTION
\end{minipage} & \begin{minipage}[b]{\linewidth}\raggedright
SUMMARY
\end{minipage} & \begin{minipage}[b]{\linewidth}\raggedright
SCHEDULE
\end{minipage} \\
\midrule\noalign{}
\endhead
\bottomrule\noalign{}
\endlastfoot
\end{longtable}

\subsubsection{Description}\label{description-707}

This keyword, TRACERKM, defines the Multi-Partitioned Tracer option K(P)
tables, for when the Partitioned Tracer option has been activate with
the PARTTRAC keyword in the RUNSPEC section, and the SOLPHASE parameter
on the TRACER keyword in the PROPS section has been set to MULT to
activate the Multi-Partitioned Tracer option. Multi-partitioned tracers
can partition into any number of phases (oil, water, gas etc.) and have
adsorption, decay and diffusion parameters specific to each phase;
whereas the standard partitioned tracers only have a ``free'' and
``solution'' phases. For the TRACERKM keyword the K(P) tables relate the
ratio of the reference phase to the other phases versus pressure. So for
example, given a multi-partitioned tracer in oil, water and gas, with
the water phase acting as the reference phase, then TRACERKM would
consist of columnar vectors of:

\begin{equation}\phantomsection\label{eq-8-92}{
{K}_{\mathit{ow}}\left(P\right) = \frac{{C}_{\mathit{oil}}}{{C}_{\mathit{water}}}\text{    and    }{K}_{\mathit{gw}}\left(P\right) = \frac{{C}_{\mathit{gas}}}{{C}_{\mathit{water}}}
}\end{equation}

Where:

\(\mathit{Kow}(P)\) = multi-partitioned oil-water K(P)

\(\mathit{Kgw}(P)\) = multi-partitioned gas-water K(P)

\(\mathit{Coil}\) = oil concentration

\(\mathit{Cgas}\) = gas concentration

\(\mathit{Cwater}\) = water concentration

See also the TRACERKP keyword in the PROPS section that provides similar
data for standard partitioned tracers.

This keyword is not supported by OPM Flow but would change the results
if supported so the simulation will be stopped.

\subsection{TRACERKP -- Standard Partitioned Tracer Option K(P)
Tables}\label{tracerkp-standard-partitioned-tracer-option-kp-tables}

\begin{longtable}[]{@{}
  >{\raggedright\arraybackslash}p{(\linewidth - 14\tabcolsep) * \real{0.1250}}
  >{\raggedright\arraybackslash}p{(\linewidth - 14\tabcolsep) * \real{0.1250}}
  >{\raggedright\arraybackslash}p{(\linewidth - 14\tabcolsep) * \real{0.1250}}
  >{\raggedright\arraybackslash}p{(\linewidth - 14\tabcolsep) * \real{0.1250}}
  >{\raggedright\arraybackslash}p{(\linewidth - 14\tabcolsep) * \real{0.1250}}
  >{\raggedright\arraybackslash}p{(\linewidth - 14\tabcolsep) * \real{0.1250}}
  >{\raggedright\arraybackslash}p{(\linewidth - 14\tabcolsep) * \real{0.1250}}
  >{\raggedright\arraybackslash}p{(\linewidth - 14\tabcolsep) * \real{0.1250}}@{}}
\toprule\noalign{}
\begin{minipage}[b]{\linewidth}\raggedright
RUNSPEC
\end{minipage} & \begin{minipage}[b]{\linewidth}\raggedright
GRID
\end{minipage} & \begin{minipage}[b]{\linewidth}\raggedright
EDIT
\end{minipage} & \begin{minipage}[b]{\linewidth}\raggedright
PROPS
\end{minipage} & \begin{minipage}[b]{\linewidth}\raggedright
REGIONS
\end{minipage} & \begin{minipage}[b]{\linewidth}\raggedright
SOLUTION
\end{minipage} & \begin{minipage}[b]{\linewidth}\raggedright
SUMMARY
\end{minipage} & \begin{minipage}[b]{\linewidth}\raggedright
SCHEDULE
\end{minipage} \\
\midrule\noalign{}
\endhead
\bottomrule\noalign{}
\endlastfoot
\end{longtable}

\subsubsection{Description}\label{description-708}

This keyword, TRACERKP, defines the Standard Partitioned Tracer option
K(P) tables, for when the Partitioned Tracer option has been activate
with the PARTTRAC keyword in the RUNSPEC section. Standard partitioned
tracers only have a ``free'' and ``solution'' phases; whereas,
Multi-partitioned tracers can partition into any number of phases (oil,
water, gas etc.) and have adsorption, decay and diffusion parameters
specific to each phase. For the TRACERKP keyword the K(P) tables relate
the ratio of the reference phase (the ``free'' phase) to the solution
phase versus pressure. So for example, given a standard partitioned
tracer in oil and gas, with the oil phase acting as the reference phase,
then TRACERKP would consist of columnar vectors of:

\begin{equation}\phantomsection\label{eq-8-93}{
K\left(P\right) = \frac{{C}_{\mathit{gas}}}{{C}_{\mathit{oil}}}
}\end{equation}

Where:

\(K(P)\) = standard partitioned K(P)

\(\mathit{Coil}\) = oil concentration

\(\mathit{Cgas}\) = gas concentration

See also the TRACERKM keyword in the PROPS section that provides similar
data for tmulti-partitioned tracers.

This keyword is not supported by OPM Flow but would change the results
if supported so the simulation will be stopped.

\subsection{TRACITVD -- Activate and Define Tracer Implicit Flux Limited
Transport
Option}\label{tracitvd-activate-and-define-tracer-implicit-flux-limited-transport-option}

\begin{longtable}[]{@{}
  >{\raggedright\arraybackslash}p{(\linewidth - 14\tabcolsep) * \real{0.1250}}
  >{\raggedright\arraybackslash}p{(\linewidth - 14\tabcolsep) * \real{0.1250}}
  >{\raggedright\arraybackslash}p{(\linewidth - 14\tabcolsep) * \real{0.1250}}
  >{\raggedright\arraybackslash}p{(\linewidth - 14\tabcolsep) * \real{0.1250}}
  >{\raggedright\arraybackslash}p{(\linewidth - 14\tabcolsep) * \real{0.1250}}
  >{\raggedright\arraybackslash}p{(\linewidth - 14\tabcolsep) * \real{0.1250}}
  >{\raggedright\arraybackslash}p{(\linewidth - 14\tabcolsep) * \real{0.1250}}
  >{\raggedright\arraybackslash}p{(\linewidth - 14\tabcolsep) * \real{0.1250}}@{}}
\toprule\noalign{}
\begin{minipage}[b]{\linewidth}\raggedright
RUNSPEC
\end{minipage} & \begin{minipage}[b]{\linewidth}\raggedright
GRID
\end{minipage} & \begin{minipage}[b]{\linewidth}\raggedright
EDIT
\end{minipage} & \begin{minipage}[b]{\linewidth}\raggedright
PROPS
\end{minipage} & \begin{minipage}[b]{\linewidth}\raggedright
REGIONS
\end{minipage} & \begin{minipage}[b]{\linewidth}\raggedright
SOLUTION
\end{minipage} & \begin{minipage}[b]{\linewidth}\raggedright
SUMMARY
\end{minipage} & \begin{minipage}[b]{\linewidth}\raggedright
SCHEDULE
\end{minipage} \\
\midrule\noalign{}
\endhead
\bottomrule\noalign{}
\endlastfoot
\end{longtable}

\subsubsection{Description}\label{description-709}

TRACITVD activates the Tracer Implicit Flux Limited Transport option and
sets various parameters for this option. Basically the option is used to
control numerical dispersion for tracers. Both the TRACERS keyword in
the RUNSPEC section and the TRACER keyword in the PROPS section must be
declared to activate tracers and to define the tracers.

See also the TRACTVD keyword in the PROPS section activates the Tracer
Explicit Flux Limited Transport option.

This keyword is not supported by OPM Flow but would change the results
if supported so the simulation will be stopped.

\subsection{TRACTVD -- Activate Tracer Explicit Flux Limited Transport
Option}\label{tractvd-activate-tracer-explicit-flux-limited-transport-option}

\begin{longtable}[]{@{}
  >{\raggedright\arraybackslash}p{(\linewidth - 14\tabcolsep) * \real{0.1250}}
  >{\raggedright\arraybackslash}p{(\linewidth - 14\tabcolsep) * \real{0.1250}}
  >{\raggedright\arraybackslash}p{(\linewidth - 14\tabcolsep) * \real{0.1250}}
  >{\raggedright\arraybackslash}p{(\linewidth - 14\tabcolsep) * \real{0.1250}}
  >{\raggedright\arraybackslash}p{(\linewidth - 14\tabcolsep) * \real{0.1250}}
  >{\raggedright\arraybackslash}p{(\linewidth - 14\tabcolsep) * \real{0.1250}}
  >{\raggedright\arraybackslash}p{(\linewidth - 14\tabcolsep) * \real{0.1250}}
  >{\raggedright\arraybackslash}p{(\linewidth - 14\tabcolsep) * \real{0.1250}}@{}}
\toprule\noalign{}
\begin{minipage}[b]{\linewidth}\raggedright
RUNSPEC
\end{minipage} & \begin{minipage}[b]{\linewidth}\raggedright
GRID
\end{minipage} & \begin{minipage}[b]{\linewidth}\raggedright
EDIT
\end{minipage} & \begin{minipage}[b]{\linewidth}\raggedright
PROPS
\end{minipage} & \begin{minipage}[b]{\linewidth}\raggedright
REGIONS
\end{minipage} & \begin{minipage}[b]{\linewidth}\raggedright
SOLUTION
\end{minipage} & \begin{minipage}[b]{\linewidth}\raggedright
SUMMARY
\end{minipage} & \begin{minipage}[b]{\linewidth}\raggedright
SCHEDULE
\end{minipage} \\
\midrule\noalign{}
\endhead
\bottomrule\noalign{}
\endlastfoot
\end{longtable}

\subsubsection{Description}\label{description-710}

TRACTVD activates the Tracer Explicit Flux Limited Transport option.
Basically the option is used to control numerical dispersion for
tracers. Both the TRACERS keyword in the RUNSPEC section and the TRACER
keyword in the PROPS section must be declared to activate tracers and to
define the tracers.

See also the TRACITVD keyword in the PROPS section activates the Tracer
Implicit Flux Limited Transport option.

This keyword is not supported by OPM Flow but would change the results
if supported so the simulation will be stopped.

\subsection{TRADS -- Environmental Tracer Adsorption
Tables}\label{trads-environmental-tracer-adsorption-tables}

\begin{longtable}[]{@{}
  >{\raggedright\arraybackslash}p{(\linewidth - 14\tabcolsep) * \real{0.1250}}
  >{\raggedright\arraybackslash}p{(\linewidth - 14\tabcolsep) * \real{0.1250}}
  >{\raggedright\arraybackslash}p{(\linewidth - 14\tabcolsep) * \real{0.1250}}
  >{\raggedright\arraybackslash}p{(\linewidth - 14\tabcolsep) * \real{0.1250}}
  >{\raggedright\arraybackslash}p{(\linewidth - 14\tabcolsep) * \real{0.1250}}
  >{\raggedright\arraybackslash}p{(\linewidth - 14\tabcolsep) * \real{0.1250}}
  >{\raggedright\arraybackslash}p{(\linewidth - 14\tabcolsep) * \real{0.1250}}
  >{\raggedright\arraybackslash}p{(\linewidth - 14\tabcolsep) * \real{0.1250}}@{}}
\toprule\noalign{}
\begin{minipage}[b]{\linewidth}\raggedright
RUNSPEC
\end{minipage} & \begin{minipage}[b]{\linewidth}\raggedright
GRID
\end{minipage} & \begin{minipage}[b]{\linewidth}\raggedright
EDIT
\end{minipage} & \begin{minipage}[b]{\linewidth}\raggedright
PROPS
\end{minipage} & \begin{minipage}[b]{\linewidth}\raggedright
REGIONS
\end{minipage} & \begin{minipage}[b]{\linewidth}\raggedright
SOLUTION
\end{minipage} & \begin{minipage}[b]{\linewidth}\raggedright
SUMMARY
\end{minipage} & \begin{minipage}[b]{\linewidth}\raggedright
SCHEDULE
\end{minipage} \\
\midrule\noalign{}
\endhead
\bottomrule\noalign{}
\endlastfoot
\end{longtable}

\subsubsection{Description}\label{description-711}

This keyword, TRADS, specifies the environmental tracer adsorption
tables that describe how a tracer is absorbed by the surrounding rock,
for when the MXENVTR parameter has been set to greater than zero on the
TRACERS keyword in the RUNSPEC section to activate environmental
tracers. The keyword can only be used with environmental tracers.

Unlike other keywords, the TRADS keyword must be concatenated with the
three character name of the tracer declared by TRACER keyword in the
PROPS section.

This keyword is not supported by OPM Flow but would change the results
if supported so the simulation will be stopped.

\subsection{TRDCY -- Environmental Tracer Decay
Tables}\label{trdcy-environmental-tracer-decay-tables}

\begin{longtable}[]{@{}
  >{\raggedright\arraybackslash}p{(\linewidth - 14\tabcolsep) * \real{0.1250}}
  >{\raggedright\arraybackslash}p{(\linewidth - 14\tabcolsep) * \real{0.1250}}
  >{\raggedright\arraybackslash}p{(\linewidth - 14\tabcolsep) * \real{0.1250}}
  >{\raggedright\arraybackslash}p{(\linewidth - 14\tabcolsep) * \real{0.1250}}
  >{\raggedright\arraybackslash}p{(\linewidth - 14\tabcolsep) * \real{0.1250}}
  >{\raggedright\arraybackslash}p{(\linewidth - 14\tabcolsep) * \real{0.1250}}
  >{\raggedright\arraybackslash}p{(\linewidth - 14\tabcolsep) * \real{0.1250}}
  >{\raggedright\arraybackslash}p{(\linewidth - 14\tabcolsep) * \real{0.1250}}@{}}
\toprule\noalign{}
\begin{minipage}[b]{\linewidth}\raggedright
RUNSPEC
\end{minipage} & \begin{minipage}[b]{\linewidth}\raggedright
GRID
\end{minipage} & \begin{minipage}[b]{\linewidth}\raggedright
EDIT
\end{minipage} & \begin{minipage}[b]{\linewidth}\raggedright
PROPS
\end{minipage} & \begin{minipage}[b]{\linewidth}\raggedright
REGIONS
\end{minipage} & \begin{minipage}[b]{\linewidth}\raggedright
SOLUTION
\end{minipage} & \begin{minipage}[b]{\linewidth}\raggedright
SUMMARY
\end{minipage} & \begin{minipage}[b]{\linewidth}\raggedright
SCHEDULE
\end{minipage} \\
\midrule\noalign{}
\endhead
\bottomrule\noalign{}
\endlastfoot
\end{longtable}

\subsubsection{Description}\label{description-712}

This keyword, TRDCY, specifies the environmental tracer decay tables
that specifies the tracer decay half-life, for when the MXENVTR
parameter has been set to greater than zero on the TRACERS keyword in
the RUNSPEC section to activate environmental tracers. The keyword can
only be used with environmental tracers.

Unlike other keywords, the TRSCY keyword must be concatenated with the
three character name of the tracer declared by TRACER keyword in the
PROPS section.

This keyword is not supported by OPM Flow but would change the results
if supported so the simulation will be stopped.

\subsection{TRDIF -- Tracer Diffusion
Tables}\label{trdif-tracer-diffusion-tables}

\begin{longtable}[]{@{}
  >{\raggedright\arraybackslash}p{(\linewidth - 14\tabcolsep) * \real{0.1250}}
  >{\raggedright\arraybackslash}p{(\linewidth - 14\tabcolsep) * \real{0.1250}}
  >{\raggedright\arraybackslash}p{(\linewidth - 14\tabcolsep) * \real{0.1250}}
  >{\raggedright\arraybackslash}p{(\linewidth - 14\tabcolsep) * \real{0.1250}}
  >{\raggedright\arraybackslash}p{(\linewidth - 14\tabcolsep) * \real{0.1250}}
  >{\raggedright\arraybackslash}p{(\linewidth - 14\tabcolsep) * \real{0.1250}}
  >{\raggedright\arraybackslash}p{(\linewidth - 14\tabcolsep) * \real{0.1250}}
  >{\raggedright\arraybackslash}p{(\linewidth - 14\tabcolsep) * \real{0.1250}}@{}}
\toprule\noalign{}
\begin{minipage}[b]{\linewidth}\raggedright
RUNSPEC
\end{minipage} & \begin{minipage}[b]{\linewidth}\raggedright
GRID
\end{minipage} & \begin{minipage}[b]{\linewidth}\raggedright
EDIT
\end{minipage} & \begin{minipage}[b]{\linewidth}\raggedright
PROPS
\end{minipage} & \begin{minipage}[b]{\linewidth}\raggedright
REGIONS
\end{minipage} & \begin{minipage}[b]{\linewidth}\raggedright
SOLUTION
\end{minipage} & \begin{minipage}[b]{\linewidth}\raggedright
SUMMARY
\end{minipage} & \begin{minipage}[b]{\linewidth}\raggedright
SCHEDULE
\end{minipage} \\
\midrule\noalign{}
\endhead
\bottomrule\noalign{}
\endlastfoot
\end{longtable}

\subsubsection{Description}\label{description-713}

This keyword, TRDIF, specifies the tracer diffusion tables that specify
the diffusion coefficient for a tracer. The keyword can be used with
Environmental Tracers if the MXENVTR parameter has been set greater than
zero on the TRACERS keyword in the RUNSPEC section. When used with a
Standard Partitioned Tracer the diffusion coefficient applies to the
solution phase, whereas as for a Multi-Partitioned Tracer the diffusion
coefficient can be entered for each defined tracer phase. Unlike other
keywords, the TRADS keyword must be concatenated with the three
character name of the tracer declared by TRACER keyword in the PROPS
section.

This keyword is not supported by OPM Flow but would change the results
if supported so the simulation will be stopped.

\subsection{TRDIS -- Tracer Dispersion Table Number
Allocation}\label{trdis-tracer-dispersion-table-number-allocation}

\begin{longtable}[]{@{}
  >{\raggedright\arraybackslash}p{(\linewidth - 14\tabcolsep) * \real{0.1250}}
  >{\raggedright\arraybackslash}p{(\linewidth - 14\tabcolsep) * \real{0.1250}}
  >{\raggedright\arraybackslash}p{(\linewidth - 14\tabcolsep) * \real{0.1250}}
  >{\raggedright\arraybackslash}p{(\linewidth - 14\tabcolsep) * \real{0.1250}}
  >{\raggedright\arraybackslash}p{(\linewidth - 14\tabcolsep) * \real{0.1250}}
  >{\raggedright\arraybackslash}p{(\linewidth - 14\tabcolsep) * \real{0.1250}}
  >{\raggedright\arraybackslash}p{(\linewidth - 14\tabcolsep) * \real{0.1250}}
  >{\raggedright\arraybackslash}p{(\linewidth - 14\tabcolsep) * \real{0.1250}}@{}}
\toprule\noalign{}
\begin{minipage}[b]{\linewidth}\raggedright
RUNSPEC
\end{minipage} & \begin{minipage}[b]{\linewidth}\raggedright
GRID
\end{minipage} & \begin{minipage}[b]{\linewidth}\raggedright
EDIT
\end{minipage} & \begin{minipage}[b]{\linewidth}\raggedright
PROPS
\end{minipage} & \begin{minipage}[b]{\linewidth}\raggedright
REGIONS
\end{minipage} & \begin{minipage}[b]{\linewidth}\raggedright
SOLUTION
\end{minipage} & \begin{minipage}[b]{\linewidth}\raggedright
SUMMARY
\end{minipage} & \begin{minipage}[b]{\linewidth}\raggedright
SCHEDULE
\end{minipage} \\
\midrule\noalign{}
\endhead
\bottomrule\noalign{}
\endlastfoot
\end{longtable}

\subsubsection{Description}\label{description-714}

This keyword, TRDIS, specifies the tracer diffusion tables that should
be allocated to a tracer, the actual dispersion tables are specified by
the DISPERSE keyword in the PROPS section. The keyword can be used with
Environmental Tracers if the MXENVTR parameter has been set greater than
zero on the TRACERS keyword in the RUNSPEC section. The option does not
work with two-phase Standard Partitioned Tracers and Multi-Partitioned
Tracers. Unlike other keywords, the TRADS keyword must be concatenated
with the three character name of the tracer declared by TRACER keyword
in the PROPS section.

This keyword is not supported by OPM Flow but would change the results
if supported so the simulation will be stopped.

\subsection{TRNHD -- Activate Dispersion Non-Homogeneous Diffusion
Option}\label{trnhd-activate-dispersion-non-homogeneous-diffusion-option}

\begin{longtable}[]{@{}
  >{\raggedright\arraybackslash}p{(\linewidth - 14\tabcolsep) * \real{0.1250}}
  >{\raggedright\arraybackslash}p{(\linewidth - 14\tabcolsep) * \real{0.1250}}
  >{\raggedright\arraybackslash}p{(\linewidth - 14\tabcolsep) * \real{0.1250}}
  >{\raggedright\arraybackslash}p{(\linewidth - 14\tabcolsep) * \real{0.1250}}
  >{\raggedright\arraybackslash}p{(\linewidth - 14\tabcolsep) * \real{0.1250}}
  >{\raggedright\arraybackslash}p{(\linewidth - 14\tabcolsep) * \real{0.1250}}
  >{\raggedright\arraybackslash}p{(\linewidth - 14\tabcolsep) * \real{0.1250}}
  >{\raggedright\arraybackslash}p{(\linewidth - 14\tabcolsep) * \real{0.1250}}@{}}
\toprule\noalign{}
\begin{minipage}[b]{\linewidth}\raggedright
RUNSPEC
\end{minipage} & \begin{minipage}[b]{\linewidth}\raggedright
GRID
\end{minipage} & \begin{minipage}[b]{\linewidth}\raggedright
EDIT
\end{minipage} & \begin{minipage}[b]{\linewidth}\raggedright
PROPS
\end{minipage} & \begin{minipage}[b]{\linewidth}\raggedright
REGIONS
\end{minipage} & \begin{minipage}[b]{\linewidth}\raggedright
SOLUTION
\end{minipage} & \begin{minipage}[b]{\linewidth}\raggedright
SUMMARY
\end{minipage} & \begin{minipage}[b]{\linewidth}\raggedright
SCHEDULE
\end{minipage} \\
\midrule\noalign{}
\endhead
\bottomrule\noalign{}
\endlastfoot
\end{longtable}

\subsubsection{Description}\label{description-715}

The TRNHD keyword activates the Dispersion Non-Homogeneous Diffusion
option for when tracer dispersion is independent of velocity or tracer
concentration. Unlike other keywords, the TRNHD keyword must be
concatenated with the name of the tracer declared by TRACER keyword in
the PROPS section.

This keyword is not supported by OPM Flow but would change the results
if supported so the simulation will be stopped.

\subsection{TRROCK -- Environmental Tracer-Rock Property
Data}\label{trrock-environmental-tracer-rock-property-data}

\begin{longtable}[]{@{}
  >{\raggedright\arraybackslash}p{(\linewidth - 14\tabcolsep) * \real{0.1250}}
  >{\raggedright\arraybackslash}p{(\linewidth - 14\tabcolsep) * \real{0.1250}}
  >{\raggedright\arraybackslash}p{(\linewidth - 14\tabcolsep) * \real{0.1250}}
  >{\raggedright\arraybackslash}p{(\linewidth - 14\tabcolsep) * \real{0.1250}}
  >{\raggedright\arraybackslash}p{(\linewidth - 14\tabcolsep) * \real{0.1250}}
  >{\raggedright\arraybackslash}p{(\linewidth - 14\tabcolsep) * \real{0.1250}}
  >{\raggedright\arraybackslash}p{(\linewidth - 14\tabcolsep) * \real{0.1250}}
  >{\raggedright\arraybackslash}p{(\linewidth - 14\tabcolsep) * \real{0.1250}}@{}}
\toprule\noalign{}
\begin{minipage}[b]{\linewidth}\raggedright
RUNSPEC
\end{minipage} & \begin{minipage}[b]{\linewidth}\raggedright
GRID
\end{minipage} & \begin{minipage}[b]{\linewidth}\raggedright
EDIT
\end{minipage} & \begin{minipage}[b]{\linewidth}\raggedright
PROPS
\end{minipage} & \begin{minipage}[b]{\linewidth}\raggedright
REGIONS
\end{minipage} & \begin{minipage}[b]{\linewidth}\raggedright
SOLUTION
\end{minipage} & \begin{minipage}[b]{\linewidth}\raggedright
SUMMARY
\end{minipage} & \begin{minipage}[b]{\linewidth}\raggedright
SCHEDULE
\end{minipage} \\
\midrule\noalign{}
\endhead
\bottomrule\noalign{}
\endlastfoot
\end{longtable}

\subsubsection{Description}\label{description-716}

This keyword, TRROCK, specifies the environmental tracer rock data for
the tracer adsorption model, for when the MXENVTR parameter has been set
to greater than zero on the TRACERS keyword in the RUNSPEC section to
activate environmental tracers. The keyword can only be used with
environmental tracers.

This keyword is not supported by OPM Flow but would change the results
if supported so the simulation will be stopped.

\subsection{TZONE -- End-Point Scaling Transition Zone
Options}\label{tzone-end-point-scaling-transition-zone-options}

\begin{longtable}[]{@{}
  >{\raggedright\arraybackslash}p{(\linewidth - 14\tabcolsep) * \real{0.1250}}
  >{\raggedright\arraybackslash}p{(\linewidth - 14\tabcolsep) * \real{0.1250}}
  >{\raggedright\arraybackslash}p{(\linewidth - 14\tabcolsep) * \real{0.1250}}
  >{\raggedright\arraybackslash}p{(\linewidth - 14\tabcolsep) * \real{0.1250}}
  >{\raggedright\arraybackslash}p{(\linewidth - 14\tabcolsep) * \real{0.1250}}
  >{\raggedright\arraybackslash}p{(\linewidth - 14\tabcolsep) * \real{0.1250}}
  >{\raggedright\arraybackslash}p{(\linewidth - 14\tabcolsep) * \real{0.1250}}
  >{\raggedright\arraybackslash}p{(\linewidth - 14\tabcolsep) * \real{0.1250}}@{}}
\toprule\noalign{}
\begin{minipage}[b]{\linewidth}\raggedright
RUNSPEC
\end{minipage} & \begin{minipage}[b]{\linewidth}\raggedright
GRID
\end{minipage} & \begin{minipage}[b]{\linewidth}\raggedright
EDIT
\end{minipage} & \begin{minipage}[b]{\linewidth}\raggedright
PROPS
\end{minipage} & \begin{minipage}[b]{\linewidth}\raggedright
REGIONS
\end{minipage} & \begin{minipage}[b]{\linewidth}\raggedright
SOLUTION
\end{minipage} & \begin{minipage}[b]{\linewidth}\raggedright
SUMMARY
\end{minipage} & \begin{minipage}[b]{\linewidth}\raggedright
SCHEDULE
\end{minipage} \\
\midrule\noalign{}
\endhead
\bottomrule\noalign{}
\endlastfoot
\end{longtable}

\subsubsection{Description}\label{description-717}

The TZONE keyword sets the transition end-point scaling options for the
oil, water and gas phases, for when the End-Point Scaling option has
been activated by the ENDSCALE keyword in the RUNSPEC section. The
keyword determines if the phase critical saturation should or should not
be set to the initial immobile saturation in areas where the initial
saturation is below the entered critical saturation.

This keyword is not supported by OPM Flow but would change the results
if supported so the simulation will be stopped.

\begin{longtable}[]{@{}
  >{\raggedright\arraybackslash}p{(\linewidth - 6\tabcolsep) * \real{0.2500}}
  >{\raggedright\arraybackslash}p{(\linewidth - 6\tabcolsep) * \real{0.2500}}
  >{\raggedright\arraybackslash}p{(\linewidth - 6\tabcolsep) * \real{0.2500}}
  >{\raggedright\arraybackslash}p{(\linewidth - 6\tabcolsep) * \real{0.2500}}@{}}
\toprule\noalign{}
\begin{minipage}[b]{\linewidth}\raggedright
No.
\end{minipage} & \begin{minipage}[b]{\linewidth}\raggedright
Name
\end{minipage} & \begin{minipage}[b]{\linewidth}\raggedright
Description
\end{minipage} & \begin{minipage}[b]{\linewidth}\raggedright
Default
\end{minipage} \\
\midrule\noalign{}
\endhead
\bottomrule\noalign{}
\endlastfoot
1 & OILZONE & OILZONE is a single character that sets the oil phase
transition zone end-point scaling option and should be set to either T
or F: & F \\
2 & WATZONE & WATZONE is a single character that sets the water phase
transition zone end-point scaling option and should be set to either T
or F: & F \\
3 & GASZONE & GASZONE is a single character that sets the gas phase
transition zone end-point scaling option and should be set to either T
or F: & F \\
Notes: & & & \\
\end{longtable}

\emph{Table 8.195: TZONE Keyword Description}

See also the SCALECRS keyword in the PROPS section that sets the
end-point scaling option to be either two-point or three-point scaling.

\subsubsection{Example}\label{example-408}

\begin{verbatim}
--
--       END-POINT SCALING TRANSITION ZONE OPTIONS
--
--       OILZONE  WATZONE  GASZONE
--       -------  -------  -------
TZONE
         F        T        F                               / SCALING OPTION
\end{verbatim}

The above example results in the SWCR values being adjusted to the
initial immobile saturations.

\subsection{VCRIT -- Critical Volumes}\label{vcrit-critical-volumes}

\begin{longtable}[]{@{}
  >{\raggedright\arraybackslash}p{(\linewidth - 14\tabcolsep) * \real{0.1250}}
  >{\raggedright\arraybackslash}p{(\linewidth - 14\tabcolsep) * \real{0.1250}}
  >{\raggedright\arraybackslash}p{(\linewidth - 14\tabcolsep) * \real{0.1250}}
  >{\raggedright\arraybackslash}p{(\linewidth - 14\tabcolsep) * \real{0.1250}}
  >{\raggedright\arraybackslash}p{(\linewidth - 14\tabcolsep) * \real{0.1250}}
  >{\raggedright\arraybackslash}p{(\linewidth - 14\tabcolsep) * \real{0.1250}}
  >{\raggedright\arraybackslash}p{(\linewidth - 14\tabcolsep) * \real{0.1250}}
  >{\raggedright\arraybackslash}p{(\linewidth - 14\tabcolsep) * \real{0.1250}}@{}}
\toprule\noalign{}
\begin{minipage}[b]{\linewidth}\raggedright
RUNSPEC
\end{minipage} & \begin{minipage}[b]{\linewidth}\raggedright
GRID
\end{minipage} & \begin{minipage}[b]{\linewidth}\raggedright
EDIT
\end{minipage} & \begin{minipage}[b]{\linewidth}\raggedright
PROPS
\end{minipage} & \begin{minipage}[b]{\linewidth}\raggedright
REGIONS
\end{minipage} & \begin{minipage}[b]{\linewidth}\raggedright
SOLUTION
\end{minipage} & \begin{minipage}[b]{\linewidth}\raggedright
SUMMARY
\end{minipage} & \begin{minipage}[b]{\linewidth}\raggedright
SCHEDULE
\end{minipage} \\
\midrule\noalign{}
\endhead
\bottomrule\noalign{}
\endlastfoot
\end{longtable}

\subsubsection{Description}\label{description-718}

The \hyperref[REF_HEADING_KEYWORD_VCRIT]{VCRIT} keyword defines the
critical volumes for each of the compositional components active in the
model and for each equation of state. The keyword should only be used if
the compositional mode has been requested using the COMPS keyword in the
RUNSPEC section.

OPM Flow does not currently support the general compositional modeling
formulation.

This keyword is not supported by OPM Flow but it will be parsed and its
data ignored.

\begin{longtable}[]{@{}
  >{\raggedright\arraybackslash}p{(\linewidth - 6\tabcolsep) * \real{0.2500}}
  >{\raggedright\arraybackslash}p{(\linewidth - 6\tabcolsep) * \real{0.2500}}
  >{\raggedright\arraybackslash}p{(\linewidth - 6\tabcolsep) * \real{0.2500}}
  >{\raggedright\arraybackslash}p{(\linewidth - 6\tabcolsep) * \real{0.2500}}@{}}
\toprule\noalign{}
\begin{minipage}[b]{\linewidth}\raggedright
No.
\end{minipage} & \begin{minipage}[b]{\linewidth}\raggedright
Name
\end{minipage} & \begin{minipage}[b]{\linewidth}\raggedright
Description
\end{minipage} & \begin{minipage}[b]{\linewidth}\raggedright
Default
\end{minipage} \\
\midrule\noalign{}
\endhead
\bottomrule\noalign{}
\endlastfoot
Field & Metric & Laboratory & \\
1 & \hyperref[REF_HEADING_KEYWORD_VCRIT]{VCRIT} & A series of real
numbers that define the critical volumes for each of the compositional
components active in the model. & None \\
ft3/lb-M & m3/kg-M & m3/kg-M & \\
Notes: & & & \\
\end{longtable}

\emph{Table 8.3.359.1: VCRIT Keyword Description}

\subsubsection{Examples}\label{examples-45}

The following example defines the critical volumes for each component in
a single three-component equation of state model.

\begin{verbatim}
--
-- Critical Volumes
--
VCRIT
  1.505  1.585  3.250 /
\end{verbatim}

The following example defines the critical volumes for each component in
two three-component equation of state models.

\begin{verbatim}
--
-- Critical Volumes
--
VCRIT
  1.505  1.585  3.250 /
  1.505  1.585  3.250 /
\end{verbatim}

\subsection{VDFLOW -- Velocity Dependent Flow Coefficient for Grid Block
Flow
(Grid)}\label{vdflow-velocity-dependent-flow-coefficient-for-grid-block-flow-grid}

\begin{longtable}[]{@{}
  >{\raggedright\arraybackslash}p{(\linewidth - 14\tabcolsep) * \real{0.1250}}
  >{\raggedright\arraybackslash}p{(\linewidth - 14\tabcolsep) * \real{0.1250}}
  >{\raggedright\arraybackslash}p{(\linewidth - 14\tabcolsep) * \real{0.1250}}
  >{\raggedright\arraybackslash}p{(\linewidth - 14\tabcolsep) * \real{0.1250}}
  >{\raggedright\arraybackslash}p{(\linewidth - 14\tabcolsep) * \real{0.1250}}
  >{\raggedright\arraybackslash}p{(\linewidth - 14\tabcolsep) * \real{0.1250}}
  >{\raggedright\arraybackslash}p{(\linewidth - 14\tabcolsep) * \real{0.1250}}
  >{\raggedright\arraybackslash}p{(\linewidth - 14\tabcolsep) * \real{0.1250}}@{}}
\toprule\noalign{}
\begin{minipage}[b]{\linewidth}\raggedright
RUNSPEC
\end{minipage} & \begin{minipage}[b]{\linewidth}\raggedright
GRID
\end{minipage} & \begin{minipage}[b]{\linewidth}\raggedright
EDIT
\end{minipage} & \begin{minipage}[b]{\linewidth}\raggedright
PROPS
\end{minipage} & \begin{minipage}[b]{\linewidth}\raggedright
REGIONS
\end{minipage} & \begin{minipage}[b]{\linewidth}\raggedright
SOLUTION
\end{minipage} & \begin{minipage}[b]{\linewidth}\raggedright
SUMMARY
\end{minipage} & \begin{minipage}[b]{\linewidth}\raggedright
SCHEDULE
\end{minipage} \\
\midrule\noalign{}
\endhead
\bottomrule\noalign{}
\endlastfoot
\end{longtable}

\subsubsection{Description}\label{description-719}

VDFLOW activates non-Darcy flow between grid blocks and defines a
constant non-Darcy flow coefficient for the whole grid, the coefficient
only applies to the gas phase. The coefficient is normally derived from
well tests or calculated analytically based on the coefficient of
inertial resistance, usually known as β, in Forchheimer's flow equation,
{[}Geertsma, J., 1974. Estimating the Coefficient of Inertial Resistance
in Fluid Flow Through Porous Media. Soc.Pet.Eng.J., October:
445-450.{]}, {[}Gewers, C.W.W. and Nichol, L.R., 1969. Gas Turbulence
Factor in a Microvugular Carbonate. J.Can.Pet.Tech., April.{]} and
{[}Wong, S.W., 1970. Effects of Liquid Saturation on Turbulence Factors
for Gas Liquid Systems. J.Can.Pet.Tech., October{]}. Dake {[}Dake, L.P.
Fundamentals of Reservoir Engineering, Amsterdam, The Netherlands,
Elsevier Science BV (1978) Chapter 8.6, pages 252-257.{]}, in chapter
eight, reports a typical value of β to be 10.07 cm-1.

See also the VDFLOWR keyword in the PROPS section that allows the
non-Darcy coefficient to be entered for individual regions, and the
WDFAC and WDFACCOR keywords in the SCHEDULE section that assigns the
non-Darcy coefficient to well connections.

This keyword is not supported by OPM Flow but would change the results
if supported so the simulation will be stopped.

\subsection{VDFLOWR -- Velocity Dependent Flow Coefficient for Grid
Block Flow
(Region)}\label{vdflowr-velocity-dependent-flow-coefficient-for-grid-block-flow-region}

\begin{longtable}[]{@{}
  >{\raggedright\arraybackslash}p{(\linewidth - 14\tabcolsep) * \real{0.1250}}
  >{\raggedright\arraybackslash}p{(\linewidth - 14\tabcolsep) * \real{0.1250}}
  >{\raggedright\arraybackslash}p{(\linewidth - 14\tabcolsep) * \real{0.1250}}
  >{\raggedright\arraybackslash}p{(\linewidth - 14\tabcolsep) * \real{0.1250}}
  >{\raggedright\arraybackslash}p{(\linewidth - 14\tabcolsep) * \real{0.1250}}
  >{\raggedright\arraybackslash}p{(\linewidth - 14\tabcolsep) * \real{0.1250}}
  >{\raggedright\arraybackslash}p{(\linewidth - 14\tabcolsep) * \real{0.1250}}
  >{\raggedright\arraybackslash}p{(\linewidth - 14\tabcolsep) * \real{0.1250}}@{}}
\toprule\noalign{}
\begin{minipage}[b]{\linewidth}\raggedright
RUNSPEC
\end{minipage} & \begin{minipage}[b]{\linewidth}\raggedright
GRID
\end{minipage} & \begin{minipage}[b]{\linewidth}\raggedright
EDIT
\end{minipage} & \begin{minipage}[b]{\linewidth}\raggedright
PROPS
\end{minipage} & \begin{minipage}[b]{\linewidth}\raggedright
REGIONS
\end{minipage} & \begin{minipage}[b]{\linewidth}\raggedright
SOLUTION
\end{minipage} & \begin{minipage}[b]{\linewidth}\raggedright
SUMMARY
\end{minipage} & \begin{minipage}[b]{\linewidth}\raggedright
SCHEDULE
\end{minipage} \\
\midrule\noalign{}
\endhead
\bottomrule\noalign{}
\endlastfoot
\end{longtable}

\subsubsection{Description}\label{description-720}

VDFLOW activates non-Darcy flow between grid blocks and defines a
constant non-Darcy flow coefficient for individual regions allocated by
the SATNUM keyword in the REGIONS section. Note that the coefficient
only applies to the gas phase. The coefficient is normally derived from
well tests or calculated analytically based on the coefficient of
inertial resistance, usually known as β, in Forchheimer's flow equation,
{[}Geertsma, J., 1974. Estimating the Coefficient of Inertial Resistance
in Fluid Flow Through Porous Media. Soc.Pet.Eng.J., October:
445-450.{]}, {[}Gewers, C.W.W. and Nichol, L.R., 1969. Gas Turbulence
Factor in a Microvugular Carbonate. J.Can.Pet.Tech., April.{]} and
{[}Wong, S.W., 1970. Effects of Liquid Saturation on Turbulence Factors
for Gas Liquid Systems. J.Can.Pet.Tech., October{]}. Dake {[}Dake, L.P.
Fundamentals of Reservoir Engineering, Amsterdam, The Netherlands,
Elsevier Science BV (1978) Chapter 8.6, pages 252-257.{]}, in chapter
eight, reports a typical value of β to be 10.07 cm-1.

See also the VDFLOW keyword in the PROPS section that allows the
non-Darcy coefficient to be entered for the whole grid, and the WDFAC
and WDFACCOR keywords in the SCHEDULE section that assigns the non-Darcy
coefficient to well connections.

This keyword is not supported by OPM Flow but would change the results
if supported so the simulation will be stopped.

\subsection{VEFRAC -- Vertical Equilibrium Relative Permeability
Fraction
(Grid)}\label{vefrac-vertical-equilibrium-relative-permeability-fraction-grid}

\begin{longtable}[]{@{}
  >{\raggedright\arraybackslash}p{(\linewidth - 14\tabcolsep) * \real{0.1250}}
  >{\raggedright\arraybackslash}p{(\linewidth - 14\tabcolsep) * \real{0.1250}}
  >{\raggedright\arraybackslash}p{(\linewidth - 14\tabcolsep) * \real{0.1250}}
  >{\raggedright\arraybackslash}p{(\linewidth - 14\tabcolsep) * \real{0.1250}}
  >{\raggedright\arraybackslash}p{(\linewidth - 14\tabcolsep) * \real{0.1250}}
  >{\raggedright\arraybackslash}p{(\linewidth - 14\tabcolsep) * \real{0.1250}}
  >{\raggedright\arraybackslash}p{(\linewidth - 14\tabcolsep) * \real{0.1250}}
  >{\raggedright\arraybackslash}p{(\linewidth - 14\tabcolsep) * \real{0.1250}}@{}}
\toprule\noalign{}
\begin{minipage}[b]{\linewidth}\raggedright
RUNSPEC
\end{minipage} & \begin{minipage}[b]{\linewidth}\raggedright
GRID
\end{minipage} & \begin{minipage}[b]{\linewidth}\raggedright
EDIT
\end{minipage} & \begin{minipage}[b]{\linewidth}\raggedright
PROPS
\end{minipage} & \begin{minipage}[b]{\linewidth}\raggedright
REGIONS
\end{minipage} & \begin{minipage}[b]{\linewidth}\raggedright
SOLUTION
\end{minipage} & \begin{minipage}[b]{\linewidth}\raggedright
SUMMARY
\end{minipage} & \begin{minipage}[b]{\linewidth}\raggedright
SCHEDULE
\end{minipage} \\
\midrule\noalign{}
\endhead
\bottomrule\noalign{}
\endlastfoot
\end{longtable}

\subsubsection{Description}\label{description-721}

This keyword defines the Vertical Equilibrium (``VE'') relative
permeability weighting factor (α) used to calculate the VE relative
permeability curves to be used in the simulation, for when the VE option
has been activated by the VE keyword in the RUNSPEC section. If α = 1.0,
then the VE model calculated relative permeability curves will be used,
and if α = 0.0, then the curves entered via the SWOF, SGOF, SLGOF series
of keywords or the SWFN, SGFN, SGWFN, SOF2, SOF3, SOF32D series of
keywords, will be used. A value of α between zero and one will result in
weighted average relative permeability curves being employed, that is:

\begin{equation}\phantomsection\label{eq-8-94}{
{\mathit{VE}}_{(\mathit{average})} = \left(1.0 - \mathrm{α}\right) \times  \left({\mathit{SATNUM}}_{\mathit{curves}}\right) + \mathrm{α} \times  \left(\mathit{VE}{\mathit{Model}}_{\mathit{curves}}\right)
}\end{equation}

Note that VEFRAC sets α for the whole grid; whereas, the VEFRACV keyword
in the PROPS section assigns α on a cell by cell basis, See also the
VEFRACP and VEFRACPV keywords that apply the weighting factors to the
capillary pressure data.

This keyword is not supported by OPM Flow but would change the results
if supported so the simulation will be stopped.

\subsection{VEFRACP -- Vertical Equilibrium Capillary Pressure Fraction
(Grid)}\label{vefracp-vertical-equilibrium-capillary-pressure-fraction-grid}

\begin{longtable}[]{@{}
  >{\raggedright\arraybackslash}p{(\linewidth - 14\tabcolsep) * \real{0.1250}}
  >{\raggedright\arraybackslash}p{(\linewidth - 14\tabcolsep) * \real{0.1250}}
  >{\raggedright\arraybackslash}p{(\linewidth - 14\tabcolsep) * \real{0.1250}}
  >{\raggedright\arraybackslash}p{(\linewidth - 14\tabcolsep) * \real{0.1250}}
  >{\raggedright\arraybackslash}p{(\linewidth - 14\tabcolsep) * \real{0.1250}}
  >{\raggedright\arraybackslash}p{(\linewidth - 14\tabcolsep) * \real{0.1250}}
  >{\raggedright\arraybackslash}p{(\linewidth - 14\tabcolsep) * \real{0.1250}}
  >{\raggedright\arraybackslash}p{(\linewidth - 14\tabcolsep) * \real{0.1250}}@{}}
\toprule\noalign{}
\begin{minipage}[b]{\linewidth}\raggedright
RUNSPEC
\end{minipage} & \begin{minipage}[b]{\linewidth}\raggedright
GRID
\end{minipage} & \begin{minipage}[b]{\linewidth}\raggedright
EDIT
\end{minipage} & \begin{minipage}[b]{\linewidth}\raggedright
PROPS
\end{minipage} & \begin{minipage}[b]{\linewidth}\raggedright
REGIONS
\end{minipage} & \begin{minipage}[b]{\linewidth}\raggedright
SOLUTION
\end{minipage} & \begin{minipage}[b]{\linewidth}\raggedright
SUMMARY
\end{minipage} & \begin{minipage}[b]{\linewidth}\raggedright
SCHEDULE
\end{minipage} \\
\midrule\noalign{}
\endhead
\bottomrule\noalign{}
\endlastfoot
\end{longtable}

\subsubsection{Description}\label{description-722}

This keyword defines the Vertical Equilibrium (``VE'') capillary
pressure weighting factor (α) used to calculate the VE capillary
pressure curves to be used in the simulation, for when the VE option has
been activated by the VE keyword in the RUNSPEC section. If α = 1.0,
then the VE model calculated capillary pressure curves will be used, and
if α = 0.0, then the curves entered via the SWOF, SGOF, SLGOF series of
keywords or the SWFN, SGFN, SGWFN, SOF2, SOF3, SOF32D series of
keywords, will be used. A value of α between zero and one will result in
weighted average capillary pressure curves being employed, that is:

\begin{equation}\phantomsection\label{eq-8-95}{
{\mathit{VE}}_{(\mathit{average})} = \left(1.0 - \mathrm{α}\right) \times  \left({\mathit{SATNUM}}_{\mathit{curves}}\right) + \mathrm{α} \times  \left(\mathit{VE}{\mathit{Model}}_{\mathit{curves}}\right)
}\end{equation}

Note that VEFRACP sets α for the whole grid; whereas, the VEFRACPV
keyword in the PROPS section assigns α on a cell by cell basis, See also
the VEFRAC and VEFRACV keywords that apply the weighting factors to the
relative permeability data.

This keyword is not supported by OPM Flow but would change the results
if supported so the simulation will be stopped.

\subsection{VEFRACPV -- Vertical Equilibrium Capillary Pressure Fraction
(Cell)}\label{vefracpv-vertical-equilibrium-capillary-pressure-fraction-cell}

\begin{longtable}[]{@{}
  >{\raggedright\arraybackslash}p{(\linewidth - 14\tabcolsep) * \real{0.1250}}
  >{\raggedright\arraybackslash}p{(\linewidth - 14\tabcolsep) * \real{0.1250}}
  >{\raggedright\arraybackslash}p{(\linewidth - 14\tabcolsep) * \real{0.1250}}
  >{\raggedright\arraybackslash}p{(\linewidth - 14\tabcolsep) * \real{0.1250}}
  >{\raggedright\arraybackslash}p{(\linewidth - 14\tabcolsep) * \real{0.1250}}
  >{\raggedright\arraybackslash}p{(\linewidth - 14\tabcolsep) * \real{0.1250}}
  >{\raggedright\arraybackslash}p{(\linewidth - 14\tabcolsep) * \real{0.1250}}
  >{\raggedright\arraybackslash}p{(\linewidth - 14\tabcolsep) * \real{0.1250}}@{}}
\toprule\noalign{}
\begin{minipage}[b]{\linewidth}\raggedright
RUNSPEC
\end{minipage} & \begin{minipage}[b]{\linewidth}\raggedright
GRID
\end{minipage} & \begin{minipage}[b]{\linewidth}\raggedright
EDIT
\end{minipage} & \begin{minipage}[b]{\linewidth}\raggedright
PROPS
\end{minipage} & \begin{minipage}[b]{\linewidth}\raggedright
REGIONS
\end{minipage} & \begin{minipage}[b]{\linewidth}\raggedright
SOLUTION
\end{minipage} & \begin{minipage}[b]{\linewidth}\raggedright
SUMMARY
\end{minipage} & \begin{minipage}[b]{\linewidth}\raggedright
SCHEDULE
\end{minipage} \\
\midrule\noalign{}
\endhead
\bottomrule\noalign{}
\endlastfoot
\end{longtable}

\subsubsection{Description}\label{description-723}

This keyword defines the Vertical Equilibrium (``VE'') capillary
pressure weighting factor (α) used to calculate the VE capillary
pressure curves to be used in the simulation, for when the VE option has
been activated by the VE keyword in the RUNSPEC section. If α = 1.0,
then the VE model calculated capillary pressure curves will be used, and
if α = 0.0, then the curves entered via the SWOF, SGOF, SLGOF series of
keywords or the SWFN, SGFN, SGWFN, SOF2, SOF3, SOF32D series of
keywords, will be used. A value of α between zero and one will result in
weighted average capillary pressure curves being employed, that is:

\begin{equation}\phantomsection\label{eq-8-96}{
{\mathit{VE}}_{(\mathit{average})} = \left(1.0 - \mathrm{α}\right) \times  \left({\mathit{SATNUM}}_{\mathit{curves}}\right) + \mathrm{α} \times  \left(\mathit{VE}{\mathit{Model}}_{\mathit{curves}}\right)
}\end{equation}

Note that VEFRACPV sets α on a cell by cell basis; whereas, the VEFRACP
keyword in the PROPS section assigns α for the whole grid, See also the
VEFRAC and VEFRACV keywords that apply the weighting factors to the
relative permeability data.

This keyword is not supported by OPM Flow but would change the results
if supported so the simulation will be stopped.

\subsection{VEFRACV -- Vertical Equilibrium Relative Permeability
Fraction
(Cell)}\label{vefracv-vertical-equilibrium-relative-permeability-fraction-cell}

\begin{longtable}[]{@{}
  >{\raggedright\arraybackslash}p{(\linewidth - 14\tabcolsep) * \real{0.1250}}
  >{\raggedright\arraybackslash}p{(\linewidth - 14\tabcolsep) * \real{0.1250}}
  >{\raggedright\arraybackslash}p{(\linewidth - 14\tabcolsep) * \real{0.1250}}
  >{\raggedright\arraybackslash}p{(\linewidth - 14\tabcolsep) * \real{0.1250}}
  >{\raggedright\arraybackslash}p{(\linewidth - 14\tabcolsep) * \real{0.1250}}
  >{\raggedright\arraybackslash}p{(\linewidth - 14\tabcolsep) * \real{0.1250}}
  >{\raggedright\arraybackslash}p{(\linewidth - 14\tabcolsep) * \real{0.1250}}
  >{\raggedright\arraybackslash}p{(\linewidth - 14\tabcolsep) * \real{0.1250}}@{}}
\toprule\noalign{}
\begin{minipage}[b]{\linewidth}\raggedright
RUNSPEC
\end{minipage} & \begin{minipage}[b]{\linewidth}\raggedright
GRID
\end{minipage} & \begin{minipage}[b]{\linewidth}\raggedright
EDIT
\end{minipage} & \begin{minipage}[b]{\linewidth}\raggedright
PROPS
\end{minipage} & \begin{minipage}[b]{\linewidth}\raggedright
REGIONS
\end{minipage} & \begin{minipage}[b]{\linewidth}\raggedright
SOLUTION
\end{minipage} & \begin{minipage}[b]{\linewidth}\raggedright
SUMMARY
\end{minipage} & \begin{minipage}[b]{\linewidth}\raggedright
SCHEDULE
\end{minipage} \\
\midrule\noalign{}
\endhead
\bottomrule\noalign{}
\endlastfoot
\end{longtable}

\subsubsection{Description}\label{description-724}

This keyword defines the Vertical Equilibrium (``VE'') relative
permeability weighting factor (α) used to calculate the VE relative
permeability curves to be used in the simulation, for when the VE option
has been activated by the VE keyword in the RUNSPEC section. If α = 1.0,
then the VE model calculated relative permeability curves will be used,
and if α = 0.0, then the curves entered via the SWOF, SGOF, SLGOF series
of keywords or the SWFN, SGFN, SGWFN, SOF2, SOF3, SOF32D series of
keywords, will be used. A value of α between zero and one will result in
weighted average relative permeability curves being employed, that is:

\begin{equation}\phantomsection\label{eq-8-97}{
{\mathit{VE}}_{(\mathit{average})} = \left(1.0 - \mathrm{α}\right) \times  \left({\mathit{SATNUM}}_{\mathit{curves}}\right) + \mathrm{α} \times  \left(\mathit{VE}{\mathit{Model}}_{\mathit{curves}}\right)
}\end{equation}

Note that VEFRACVV sets α on a cell by cell basis; whereas, the VEFRAC
keyword in the PROPS section assigns α for the whole grid, See also the
VEFRACP and VEFRACPV keywords that apply the weighting factors to the
capillary pressure data.

This keyword is not supported by OPM Flow but would change the results
if supported so the simulation will be stopped.

\subsection{VISCAQA -- Specify Ezrokhi Coefficients for Aqueous
Viscosity}\label{viscaqa-specify-ezrokhi-coefficients-for-aqueous-viscosity}

\begin{longtable}[]{@{}
  >{\raggedright\arraybackslash}p{(\linewidth - 14\tabcolsep) * \real{0.1250}}
  >{\raggedright\arraybackslash}p{(\linewidth - 14\tabcolsep) * \real{0.1250}}
  >{\raggedright\arraybackslash}p{(\linewidth - 14\tabcolsep) * \real{0.1250}}
  >{\raggedright\arraybackslash}p{(\linewidth - 14\tabcolsep) * \real{0.1250}}
  >{\raggedright\arraybackslash}p{(\linewidth - 14\tabcolsep) * \real{0.1250}}
  >{\raggedright\arraybackslash}p{(\linewidth - 14\tabcolsep) * \real{0.1250}}
  >{\raggedright\arraybackslash}p{(\linewidth - 14\tabcolsep) * \real{0.1250}}
  >{\raggedright\arraybackslash}p{(\linewidth - 14\tabcolsep) * \real{0.1250}}@{}}
\toprule\noalign{}
\begin{minipage}[b]{\linewidth}\raggedright
RUNSPEC
\end{minipage} & \begin{minipage}[b]{\linewidth}\raggedright
GRID
\end{minipage} & \begin{minipage}[b]{\linewidth}\raggedright
EDIT
\end{minipage} & \begin{minipage}[b]{\linewidth}\raggedright
PROPS
\end{minipage} & \begin{minipage}[b]{\linewidth}\raggedright
REGIONS
\end{minipage} & \begin{minipage}[b]{\linewidth}\raggedright
SOLUTION
\end{minipage} & \begin{minipage}[b]{\linewidth}\raggedright
SUMMARY
\end{minipage} & \begin{minipage}[b]{\linewidth}\raggedright
SCHEDULE
\end{minipage} \\
\midrule\noalign{}
\endhead
\bottomrule\noalign{}
\endlastfoot
\end{longtable}

\subsubsection{Description}\label{description-725}

The \hyperref[REF_HEADING_KEYWORD_VISCAQA_8_3]{VISCAQA} keyword
specifies the three Ezrokhi coefficients for each compositional
component and for each equation of state that are used to calculate the
aqueous phase viscosity. The keyword should only be used if the
compositional mode has been requested using the COMPS keyword in the
RUNSPEC section.

OPM Flow does not currently support the general compositional modeling
formulation.

This keyword is only supported by OPM Flow when the two phase gas-water
CO2 storage model has been activated using the CO2STORE keyword and
either the GASWAT or the GAS and WATER keywords in the RUNSPEC section.
The component names ``H2O'', ``CO2'' and ``NACL'' (water, CO2 and salt
respectively) must be specified using the CNAMES keyword.

\begin{longtable}[]{@{}
  >{\raggedright\arraybackslash}p{(\linewidth - 6\tabcolsep) * \real{0.2500}}
  >{\raggedright\arraybackslash}p{(\linewidth - 6\tabcolsep) * \real{0.2500}}
  >{\raggedright\arraybackslash}p{(\linewidth - 6\tabcolsep) * \real{0.2500}}
  >{\raggedright\arraybackslash}p{(\linewidth - 6\tabcolsep) * \real{0.2500}}@{}}
\toprule\noalign{}
\begin{minipage}[b]{\linewidth}\raggedright
No.
\end{minipage} & \begin{minipage}[b]{\linewidth}\raggedright
Name
\end{minipage} & \begin{minipage}[b]{\linewidth}\raggedright
Description
\end{minipage} & \begin{minipage}[b]{\linewidth}\raggedright
Default
\end{minipage} \\
\midrule\noalign{}
\endhead
\bottomrule\noalign{}
\endlastfoot
Field & Metric & Laboratory & \\
1 & COEFFS & A series of real numbers that define the three Ezrokhi
coeffients for each of the compositional components active in the model.
& None \\
dimensionless & dimensionless & dimensionless & \\
Notes: & & & \\
\end{longtable}

\emph{Table 8.3.376.1: VISCAQA Keyword Description}

\subsubsection{Examples}\label{examples-46}

The following example defines the Ezrokhi coefficients for each
component in a single three-component equation of state model.

\begin{verbatim}
--
-- Ezrokhi Coefficients for Aqueous Viscosity Calculation
--
VISCAQA
--COEFF0   COEFF1   COEFF2
  2.0E-7   0.0      1.0E-9
  2.0E-7   0.0      1.0E-9
  2.0E-7   0.0      1.0E-9   /
\end{verbatim}

The following example defines the Ezrokhi coefficients for each
component in two three-component equation of state models.

\begin{verbatim}
--
-- Ezrokhi Coefficients for Aqueous Viscosity Calculation
--
VISCAQA
--COEFF0   COEFF1   COEFF2
  2.0E-7   0.0      1.0E-9
  2.0E-7   0.0      1.0E-9
  2.0E-7   0.0      1.0E-9   /
--
  2.0E-7   0.0      1.0E-9
  2.0E-7   0.0      1.0E-9
  2.0E-7   0.0      1.0E-9   /
\end{verbatim}

\subsection{VISCREF -- Define Viscosity-Temperature Reference
Conditions}\label{viscref-define-viscosity-temperature-reference-conditions}

\begin{longtable}[]{@{}
  >{\raggedright\arraybackslash}p{(\linewidth - 14\tabcolsep) * \real{0.1250}}
  >{\raggedright\arraybackslash}p{(\linewidth - 14\tabcolsep) * \real{0.1250}}
  >{\raggedright\arraybackslash}p{(\linewidth - 14\tabcolsep) * \real{0.1250}}
  >{\raggedright\arraybackslash}p{(\linewidth - 14\tabcolsep) * \real{0.1250}}
  >{\raggedright\arraybackslash}p{(\linewidth - 14\tabcolsep) * \real{0.1250}}
  >{\raggedright\arraybackslash}p{(\linewidth - 14\tabcolsep) * \real{0.1250}}
  >{\raggedright\arraybackslash}p{(\linewidth - 14\tabcolsep) * \real{0.1250}}
  >{\raggedright\arraybackslash}p{(\linewidth - 14\tabcolsep) * \real{0.1250}}@{}}
\toprule\noalign{}
\begin{minipage}[b]{\linewidth}\raggedright
RUNSPEC
\end{minipage} & \begin{minipage}[b]{\linewidth}\raggedright
GRID
\end{minipage} & \begin{minipage}[b]{\linewidth}\raggedright
EDIT
\end{minipage} & \begin{minipage}[b]{\linewidth}\raggedright
PROPS
\end{minipage} & \begin{minipage}[b]{\linewidth}\raggedright
REGIONS
\end{minipage} & \begin{minipage}[b]{\linewidth}\raggedright
SOLUTION
\end{minipage} & \begin{minipage}[b]{\linewidth}\raggedright
SUMMARY
\end{minipage} & \begin{minipage}[b]{\linewidth}\raggedright
SCHEDULE
\end{minipage} \\
\midrule\noalign{}
\endhead
\bottomrule\noalign{}
\endlastfoot
\end{longtable}

\subsubsection{Description}\label{description-726}

VISCREF defines the reference conditions for the viscosity-temperature
tables, GASVISCT, OILVISCT and WATVISCT, for when the thermal option has
been activated by THERMAL keyword in the RUNSPEC section. This keyword
can only be used if the thermal option has been activated by the THERMAL
keyword in the RUNSPEC section. Note this is different to the commercial
simulator that uses the TEMP keyword in the RUNSPEC section to activate
the black-oil thermal model.

\begin{longtable}[]{@{}
  >{\raggedright\arraybackslash}p{(\linewidth - 6\tabcolsep) * \real{0.2500}}
  >{\raggedright\arraybackslash}p{(\linewidth - 6\tabcolsep) * \real{0.2500}}
  >{\raggedright\arraybackslash}p{(\linewidth - 6\tabcolsep) * \real{0.2500}}
  >{\raggedright\arraybackslash}p{(\linewidth - 6\tabcolsep) * \real{0.2500}}@{}}
\toprule\noalign{}
\begin{minipage}[b]{\linewidth}\raggedright
No.
\end{minipage} & \begin{minipage}[b]{\linewidth}\raggedright
Name
\end{minipage} & \begin{minipage}[b]{\linewidth}\raggedright
Description
\end{minipage} & \begin{minipage}[b]{\linewidth}\raggedright
Default
\end{minipage} \\
\midrule\noalign{}
\endhead
\bottomrule\noalign{}
\endlastfoot
Field & Metric & Laboratory & \\
1 & PRES & PRES is a real positive number defining the reference
pressure for the viscosity and temperature tables & None \\
psia & barsa & atma & \\
2 & RS & RS is a real positive number defining the reference gas-oil
ratio for when the model contains gas dissolved as activated by the
DISGAS keyword in the RUNSPEC section & None \\
Mscf/stb & sm3/sm3 & scc/scc & \\
3 & API & API is a real number defining the oil API for when the API
tracking option has been invoked by the API keyword in the RUNSPEC
section. Note that OPM Flow does not support API tracking, and therefore
this variable is ignored. & None \\
oAPI & oAPI & oAPI & \\
Notes: & & & \\
\end{longtable}

\emph{Table 8.196: VISCREF Keyword Description}

OPM Flow currently does not support API tracking and therefore item (3)
of this keyword is ignored. See also the OILVISCT, GASVISCT and WATVISCT
keywords in the PROPS section.

\subsubsection{Example}\label{example-409}

The following example shows the VISCREF keyword for when the thermal
option has been activated by the THERMAL keyword in the RUNSPEC section
and for when NTPVT on the TABDIMS keyword in the RUNSPEC section is set
to five.

\begin{verbatim}
--
--       REF        REF       REF
--       PRESSURE   GOR       API
--       --------   -------   -------
VISCREF
          3000.0    0.500                                  / TABLE NO. 01
          3200.0    0.550                                  / TABLE NO. 02
          3300.0    0.580                                  / TABLE NO. 03
          3400.0    0.620                                  / TABLE NO. 04
          3500.0    0.625                                  / TABLE NO. 05
\end{verbatim}

There is no terminating ``/'' for this keyword.

\subsection{WAGHYSTR -- Define Water-Alternating-Gas Hysteresis
Parameters}\label{waghystr-define-water-alternating-gas-hysteresis-parameters}

\begin{longtable}[]{@{}
  >{\raggedright\arraybackslash}p{(\linewidth - 14\tabcolsep) * \real{0.1250}}
  >{\raggedright\arraybackslash}p{(\linewidth - 14\tabcolsep) * \real{0.1250}}
  >{\raggedright\arraybackslash}p{(\linewidth - 14\tabcolsep) * \real{0.1250}}
  >{\raggedright\arraybackslash}p{(\linewidth - 14\tabcolsep) * \real{0.1250}}
  >{\raggedright\arraybackslash}p{(\linewidth - 14\tabcolsep) * \real{0.1250}}
  >{\raggedright\arraybackslash}p{(\linewidth - 14\tabcolsep) * \real{0.1250}}
  >{\raggedright\arraybackslash}p{(\linewidth - 14\tabcolsep) * \real{0.1250}}
  >{\raggedright\arraybackslash}p{(\linewidth - 14\tabcolsep) * \real{0.1250}}@{}}
\toprule\noalign{}
\begin{minipage}[b]{\linewidth}\raggedright
RUNSPEC
\end{minipage} & \begin{minipage}[b]{\linewidth}\raggedright
GRID
\end{minipage} & \begin{minipage}[b]{\linewidth}\raggedright
EDIT
\end{minipage} & \begin{minipage}[b]{\linewidth}\raggedright
PROPS
\end{minipage} & \begin{minipage}[b]{\linewidth}\raggedright
REGIONS
\end{minipage} & \begin{minipage}[b]{\linewidth}\raggedright
SOLUTION
\end{minipage} & \begin{minipage}[b]{\linewidth}\raggedright
SUMMARY
\end{minipage} & \begin{minipage}[b]{\linewidth}\raggedright
SCHEDULE
\end{minipage} \\
\midrule\noalign{}
\endhead
\bottomrule\noalign{}
\endlastfoot
\end{longtable}

\subsubsection{Description}\label{description-727}

This keyword defines the parameters for the Water-Alternating-Gas
(``WAG'') hysteresis option, for when the hysteresis option has been
activated by the WAGHYSTR variable on the SATOPTS keyword in the RUNSPEC
section.

The WAG recovery mechanism is an Enhanced Oil Recovery (``EOR'') process
to optimize oil recovery by improving volumetric sweep efficiency. It
was originally proposed as a method to improve the sweep efficiency of
gas by using water to control the mobility ratio and to stabilize the
front (Caudle and Dyes, 1958 {[}Caudle, B. H., \& Dyes, A. B. (1958,
January 1). Improving Miscible Displacement by Gas-Water Injection.
Society of Petroleum Engineers.{]}; Christensen et al., 1998
{[}Christensen, J. R., Stenby, E. H., \& Skauge, A. (1998, January 1).
Review of WAG Field Experience. Society of Petroleum Engineers.
doi:10.2118/39883-MS.{]}; and Christensen et al., 2001 {[}Christensen,
J. R., Stenby, E. H., \& Skauge, A. (2001, April 1). Review of WAG Field
Experience. Society of Petroleum Engineers. doi:10.2118/71203-PA.{]}).
WAG injection can lead to improved oil recovery by combining better
mobility control and contacting upswept zones, and by leading to
improved microscopic displacement. Although initially the injected gas
was immiscible with respect to the oil (WAG Immiscible) the more common
process is WAG Miscible, with alternating different types of hydrocarbon
gases and non-hydrocarbon gases, such as N2 and CO2 Gases. WAG flooding
has been successfully applied to various fields worldwide.

Only the gas phase relative permeability WAG hysteresis model has been
implemented in the simulator.

\begin{longtable}[]{@{}
  >{\raggedright\arraybackslash}p{(\linewidth - 6\tabcolsep) * \real{0.2500}}
  >{\raggedright\arraybackslash}p{(\linewidth - 6\tabcolsep) * \real{0.2500}}
  >{\raggedright\arraybackslash}p{(\linewidth - 6\tabcolsep) * \real{0.2500}}
  >{\raggedright\arraybackslash}p{(\linewidth - 6\tabcolsep) * \real{0.2500}}@{}}
\toprule\noalign{}
\begin{minipage}[b]{\linewidth}\raggedright
No.
\end{minipage} & \begin{minipage}[b]{\linewidth}\raggedright
Name
\end{minipage} & \begin{minipage}[b]{\linewidth}\raggedright
Description
\end{minipage} & \begin{minipage}[b]{\linewidth}\raggedright
Default
\end{minipage} \\
\midrule\noalign{}
\endhead
\bottomrule\noalign{}
\endlastfoot
Field & Metric & Laboratory & \\
1 & LANDS\_PARAMETER & A real value greater than zero that defines
Land's parameter, \(C\). The Land's parameter controls how the trapped
gas saturation depends on the maximum gas saturation attained and the
critical gas saturation; and the shape of the imbibition curve.
\({s}_{\mathit{gtrap}}={s}_{\mathit{gcr}}+\frac{({s}_{\mathit{gm}}-{s}_{\mathit{gcr}})}{(1+C({s}_{\mathit{gm}}-{s}_{\mathit{gcr}}))}\)
where, \({s}_{\mathit{gtrap}}\) is the trapped gas saturation,
\({s}_{\mathit{gm}}\) is the maximum gas saturation attained, and
\({s}_{\mathit{gcr}}\) is the critical gas saturation. Values of the
Land's parameter that are too small give a trapped gas saturation close
to the maximum gas saturation attained. This results in an unphysical
steep relative permeability curve giving potential convergence problems.
& None \\
dimensionless & dimensionless & dimensionless & \\
2 & SECONDARY\_DRAINAGE\_REDUCTION & A real value greater than or equal
to zero that defines the secondary drainage reduction factor, alpha. As
alpha increases the reduction in gas mobility on secondary drainage
increases. & 0.0 \\
dimensionless & dimensionless & dimensionless & \\
3 & GAS\_MODEL & A defined character string that defines whether the gas
hysteresis model should be used, and should be set to one of the
following character strings: Only the YES option is currently supported
by the simulator. & YES \\
4 & RES\_OIL & A defined character string that defines whether the
residual oil model should be used, and should be set to one of the
following character strings: Only the NO option is currently supported
by the simulator. & YES \\
5 & WATER\_MODEL & A defined character string that defines whether the
water hysteresis model should be used, and should be set to one of the
following character strings: Only the NO option is currently supported
by the simulator. & YES \\
6 & IMB\_LINEAR\_FRACTION & A real value greater than zero that defines
the imbibtion curve linear fraction. This is the fraction of the curve
between Sgm and Sgtrap that uses a linear transformation. A non-zero
value for the linear fraction prevents the potential infinite gradient
in the imbibition curve when using the Carlson analytic model. & 0.1 \\
dimensionless & dimensionless & dimensionless & \\
7 & THREEPHASE\_SAT\_LIMIT & A real value between zero and one that
defines the three-phase model threshold saturation. When the water
saturation exceeds this threshold above the connate water saturation the
gas (non-wetting) phase hysteresis switches from the two-phase model to
the three-phase model. In the two-phase model a secondary drainage
process follows the imbibition curve. However, if the water saturation
exceeds the connate saturation by the given threshold, at the beginning
of the secondary drainage process a three-phase secondary drainage curve
is followed. This value also defines the minimum percentage change in
gas saturation to allow switching from drainage to imbibition curve and
vice-versa. This threshold allows better control of the numerical
sensitivity of the system, preventing it from being too unstable. &
0.001 \\
dimensionless & dimensionless & dimensionless & \\
8 & RES\_OIL\_MOD\_FRACTION & A real value between zero and one that
defines the residual oil modification fraction. This is the fraction of
the trapped gas saturation subtracted from the residual oil (SOM) in the
STONE 1 three-phase oil relative permeability model. This is not
supported and will be ignored by the simulator. & 1.0 \\
dimensionless & dimensionless & dimensionless & \\
Notes: & & & \\
\end{longtable}

\emph{Table 8.3.364.1: WAGHYSTR Keyword Description}

\subsubsection{Example}\label{example-410}

The following example defines the WAG hysteresis model parameters using
the WAGHYSTR keyword for a case with three saturation table regions

\begin{verbatim}
--
--       WAG HYSTERESIS PARAMETERS
--
-- LAND  ALPHA    GAS    RES    WATER  LINEAR
-- C     FACTOR   MODEL  OIL    MODEL  FRAC
WAGHYSTR
2.0      1.0      YES    YES    YES    0.2    /
2.0      1.0      YES    NO     /
2.0      /
\end{verbatim}

Here, saturation table region one uses the gas WAG hysteresis, residual
oil and water WAG hysteresis models, with a Land's parameter of 2.0, an
alpha factor of 1.0, and a imbibition curve linear factor of 0.2. The
gas and water WAG hysteresis models are used in region two but the
residual oil model is turned off, with a Land's parameter of 2.0, an
alpha factor of 1.0, and a default imbibition curve linear factor of
0.1. A Land's parameter of 2.0 has been specified for region three with
the remainder of the parameters defaulted.

Note that the above example uses some options that are not currently
supported by OPM Flow.

\subsection{WARN -- Activate Warning
Messages}\label{warn-activate-warning-messages-4}

Turns on warning messages to be printed to the print file (*.PRT); note
that this keyword is activated by default and can subsequently be
switched off by the NOWARN activation keyword. The warning messages may
be turned on and off using keywords WARN and NOWARN. OPM Flow always
prints error messages.

See WARN -- Activate Warning Messages in the GLOBAL section for a full
description.

\subsection{WATDENT -- Define Water Density Temperature
Coefficients}\label{watdent-define-water-density-temperature-coefficients}

\begin{longtable}[]{@{}
  >{\raggedright\arraybackslash}p{(\linewidth - 14\tabcolsep) * \real{0.1250}}
  >{\raggedright\arraybackslash}p{(\linewidth - 14\tabcolsep) * \real{0.1250}}
  >{\raggedright\arraybackslash}p{(\linewidth - 14\tabcolsep) * \real{0.1250}}
  >{\raggedright\arraybackslash}p{(\linewidth - 14\tabcolsep) * \real{0.1250}}
  >{\raggedright\arraybackslash}p{(\linewidth - 14\tabcolsep) * \real{0.1250}}
  >{\raggedright\arraybackslash}p{(\linewidth - 14\tabcolsep) * \real{0.1250}}
  >{\raggedright\arraybackslash}p{(\linewidth - 14\tabcolsep) * \real{0.1250}}
  >{\raggedright\arraybackslash}p{(\linewidth - 14\tabcolsep) * \real{0.1250}}@{}}
\toprule\noalign{}
\begin{minipage}[b]{\linewidth}\raggedright
RUNSPEC
\end{minipage} & \begin{minipage}[b]{\linewidth}\raggedright
GRID
\end{minipage} & \begin{minipage}[b]{\linewidth}\raggedright
EDIT
\end{minipage} & \begin{minipage}[b]{\linewidth}\raggedright
PROPS
\end{minipage} & \begin{minipage}[b]{\linewidth}\raggedright
REGIONS
\end{minipage} & \begin{minipage}[b]{\linewidth}\raggedright
SOLUTION
\end{minipage} & \begin{minipage}[b]{\linewidth}\raggedright
SUMMARY
\end{minipage} & \begin{minipage}[b]{\linewidth}\raggedright
SCHEDULE
\end{minipage} \\
\midrule\noalign{}
\endhead
\bottomrule\noalign{}
\endlastfoot
\end{longtable}

\subsubsection{Description}\label{description-728}

WATDENT defines the water density as a function of temperature
coefficients for when the thermal option has been activated by the
THERMAL keyword in the RUNSPEC. Note this is an OPM Flow keyword used
with OPM Flow's black-oil thermal model that is not available in the
commercial simulator's black-oil thermal formulation.

This keyword can only be used if OPM Flow's thermal option has been
activated by the THERMAL keyword in the RUNSPEC section. Note this is
different to the commercial simulator that uses the TEMP keyword in the
RUNSPEC section to activate the black-oil thermal model.

\begin{longtable}[]{@{}
  >{\raggedright\arraybackslash}p{(\linewidth - 6\tabcolsep) * \real{0.2500}}
  >{\raggedright\arraybackslash}p{(\linewidth - 6\tabcolsep) * \real{0.2500}}
  >{\raggedright\arraybackslash}p{(\linewidth - 6\tabcolsep) * \real{0.2500}}
  >{\raggedright\arraybackslash}p{(\linewidth - 6\tabcolsep) * \real{0.2500}}@{}}
\toprule\noalign{}
\begin{minipage}[b]{\linewidth}\raggedright
No.
\end{minipage} & \begin{minipage}[b]{\linewidth}\raggedright
Name
\end{minipage} & \begin{minipage}[b]{\linewidth}\raggedright
Description
\end{minipage} & \begin{minipage}[b]{\linewidth}\raggedright
Default
\end{minipage} \\
\midrule\noalign{}
\endhead
\bottomrule\noalign{}
\endlastfoot
Field & Metric & Laboratory & \\
1 & TEMP & TEMP is a real positive value greater than zero that defines
the absolute reference temperature used with TEXP1 and TEXP2 to estimate
the change in water density with respect to temperature. & Defined \\
oR 527.67 & K 293.15 & K 293.15 & \\
2 & TEXP1 & TEXP1 is a real positive value greater than zero that
defines the water thermal expansion coefficient of the first order. &
Defined \\
1/oR 1.67 x 10-4 & 1/K 3.0 x 10-4 & 1/K 3.0 x 10-4 & \\
3 & TEXP2 & TEXP2 is a real positive value greater than zero that
defines the water thermal expansion coefficient of the second order. &
Defined \\
1/oR2 9.26 x 10-7 & 1/K2 3.0 x 10-6 & 1/K2 3.0 x 10-6 & \\
Notes: & & & \\
\end{longtable}

\emph{Table 8.3.366.1: WATDENT Keyword Description}

The water density at a given pressure and temperature is calculated from
its value at surface conditions and the water shrinkage factor (the
reciprocal of the water formation volume factor) as shown in the
following equation:

\begin{equation}\phantomsection\label{eq-8-3-366-1}{
{ρ}_{w}(p,T)={ρ}_{w}({p}_{s},{T}_{s}){b}_{w}(p,T)
}\end{equation}

Where the temperature dependence of the water shrinkage factor relative
to its value at the reference temperature is calculated as shown in the
following equation:

\begin{equation}\phantomsection\label{eq-8-3-366-2}{
{b}_{w}(p,T)=\frac{{b}_{w}(p,{T}_{\mathit{ref}})}{1+{c}_{1}(T-{T}_{\mathit{ref}})+{c}_{2}{(T-{T}_{\mathit{ref}})}^{2}}
}\end{equation}

Where:

\({ρ}_{w}\) = water density

\({b}_{w}\) = water shrinkage factor

\(p\) = pressure

\(T\) = temperature

\({c}_{1},{c}_{2}\) = thermal expansion coefficients to first and second
order

\(s\) = subscript indicating surface conditions

\(\mathit{ref}\) = subscript indicating reference conditions

\subsubsection{Example}\label{example-411}

The following example shows the WATDENT keyword using the default
values, for when the thermal option has been activated by the THERMAL
keyword in the RUNSPEC section and for when NTPVT on the TABDIMS keyword
in the RUNSPEC section is set to two.

\begin{verbatim}
--
--       WATER DENSITY TEMPERATURE COEFFICIENTS (OPM FLOW EXTENSION KEYWORD)
--
--       WATER      DENSITY   DENSITY
--       TEMP       COEFF1    COEFF2
--       --------   -------   -------
WATDENT
         1*         1*        1*                           / TABLE NO. 01
         1*         1*        1*                           / TABLE NO. 02
\end{verbatim}

There is no terminating ``/'' for this keyword.

\subsection{WATJT -- Define Water Joule-Thomson
Coefficient}\label{watjt-define-water-joule-thomson-coefficient}

\begin{longtable}[]{@{}
  >{\raggedright\arraybackslash}p{(\linewidth - 14\tabcolsep) * \real{0.1250}}
  >{\raggedright\arraybackslash}p{(\linewidth - 14\tabcolsep) * \real{0.1250}}
  >{\raggedright\arraybackslash}p{(\linewidth - 14\tabcolsep) * \real{0.1250}}
  >{\raggedright\arraybackslash}p{(\linewidth - 14\tabcolsep) * \real{0.1250}}
  >{\raggedright\arraybackslash}p{(\linewidth - 14\tabcolsep) * \real{0.1250}}
  >{\raggedright\arraybackslash}p{(\linewidth - 14\tabcolsep) * \real{0.1250}}
  >{\raggedright\arraybackslash}p{(\linewidth - 14\tabcolsep) * \real{0.1250}}
  >{\raggedright\arraybackslash}p{(\linewidth - 14\tabcolsep) * \real{0.1250}}@{}}
\toprule\noalign{}
\begin{minipage}[b]{\linewidth}\raggedright
RUNSPEC
\end{minipage} & \begin{minipage}[b]{\linewidth}\raggedright
GRID
\end{minipage} & \begin{minipage}[b]{\linewidth}\raggedright
EDIT
\end{minipage} & \begin{minipage}[b]{\linewidth}\raggedright
PROPS
\end{minipage} & \begin{minipage}[b]{\linewidth}\raggedright
REGIONS
\end{minipage} & \begin{minipage}[b]{\linewidth}\raggedright
SOLUTION
\end{minipage} & \begin{minipage}[b]{\linewidth}\raggedright
SUMMARY
\end{minipage} & \begin{minipage}[b]{\linewidth}\raggedright
SCHEDULE
\end{minipage} \\
\midrule\noalign{}
\endhead
\bottomrule\noalign{}
\endlastfoot
\end{longtable}

\subsubsection{Description}\label{description-729}

WATJT activates the water Joule-Thomson effect {[}The Joule--Thomson
coefficient is defined as the change in temperature with respect to an
increase in pressure at constant enthalpy.{]} in temperature
calculations, and defines the water Joule-Thomson Coefficient (``JTC'')
at a given reference pressure, for when OPM Flow's thermal option has
been activated by the THERMAL keyword in the RUNSPEC.

\begin{tcolorbox}[enhanced jigsaw, title=\textcolor{quarto-callout-note-color}{\faInfo}\hspace{0.5em}{Note}, coltitle=black, breakable, toprule=.15mm, colbacktitle=quarto-callout-note-color!10!white, leftrule=.75mm, opacitybacktitle=0.6, colback=white, bottomtitle=1mm, toptitle=1mm, titlerule=0mm, left=2mm, colframe=quarto-callout-note-color-frame, arc=.35mm, rightrule=.15mm, bottomrule=.15mm, opacityback=0]

This is an OPM Flow keyword used with OPM Flow's black-oil thermal
model, that is not available in the commercial simulator's black-oil
thermal formulation.

\end{tcolorbox}

This keyword can only be used if OPM Flow's thermal option has been
activated by the THERMAL keyword in the RUNSPEC section. Note this is
different to the commercial simulator that uses the TEMP keyword in the
RUNSPEC section to activate the black-oil thermal model, and does not
include the Joule-Thomson effect in temperature calculations.

\begin{longtable}[]{@{}
  >{\raggedright\arraybackslash}p{(\linewidth - 6\tabcolsep) * \real{0.2500}}
  >{\raggedright\arraybackslash}p{(\linewidth - 6\tabcolsep) * \real{0.2500}}
  >{\raggedright\arraybackslash}p{(\linewidth - 6\tabcolsep) * \real{0.2500}}
  >{\raggedright\arraybackslash}p{(\linewidth - 6\tabcolsep) * \real{0.2500}}@{}}
\toprule\noalign{}
\begin{minipage}[b]{\linewidth}\raggedright
No.
\end{minipage} & \begin{minipage}[b]{\linewidth}\raggedright
Name
\end{minipage} & \begin{minipage}[b]{\linewidth}\raggedright
Description
\end{minipage} & \begin{minipage}[b]{\linewidth}\raggedright
Default
\end{minipage} \\
\midrule\noalign{}
\endhead
\bottomrule\noalign{}
\endlastfoot
Field & Metric & Laboratory & \\
1 & PRESS & A real positive value that defines the reference pressure
for the corresponding Joule-Thomson Coefficient, WATJTC. & None \\
psia & barsa & atma & \\
2 & WATJTC & WATJTC is a real positive or negative value that defines
the water phase Joule-Thomson Coefficient. If the value is defaulted
(1*) or set to 0, then WATJTC is internally calculated using thermal
water density data on the WATSDENT keyword in the PROPS section. If a
non-zero value is specified, then the WATJTC is assumed to be constant
and equal to that value. & 0 \\
oF/psia & oC/barsa & oC/atma & \\
Notes: & & & \\
\end{longtable}

\emph{Table 8.198: WATJT Keyword Description}

The Joule--Thomson effect is when a real gas, as oppose to an ideal gas,
expands, resulting in the temperature of the gas dropping {[}Natural Gas
Engineering (McGraw-Hill chemical engineering series), Donald L. Katz,
Robert l Lee, McGraw-Hill Education, 1990 (ISBN 0071007776,
9780071007771).{]}. For liquids the effect is the opposite, that is the
internal energy is transferred to kinetic energy with a corresponding
increase in temperature as velocity increases.

Thermodynamically, the Joule--Thomson coefficient is defined as the
isenthalpic {[}An isenthalpic process or isoenthalpic process, is a
process that proceeds without any change in enthalpy, H; or specific
enthalpy, h.{]} change in temperature in a fluid caused by a unitary
pressure drop, as shown in the following equation:

\begin{equation}\phantomsection\label{eq-8-98}{
\mathrm{η} = \left(\frac{\partial T}{\partial P}\right)
}\end{equation}

Which can also express as {[}Pippard, A.B.: Elements of Classical
Thermodynamics: For Advanced Students of Physics. Cambridge University
Press, Cambridge, UK (1957){]}:

\begin{equation}\phantomsection\label{eq-8-99}{
\mathrm{η} = \left(T\mathrm{α} - 1\right)\frac{1}{(\mathrm{ρ}{C}_{p})} - {\left(\frac{g}{{C}_{p}}\frac{\mathit{dp}}{\mathit{dz}}\right)}^{-1}
}\end{equation}

Setting the gravity term, \(g\), to zero we have:

\begin{equation}\phantomsection\label{eq-8-100}{
\mathrm{η} = \left(T\mathrm{α} - 1\right)\frac{1}{(\mathrm{ρ}{C}_{b})}
}\end{equation}

Where:

\(η\) = Joule--Thomson coefficient (oC/Pa),

\(α\) = thermal expansivity at constant pressure (1/oC),

\({C}_{p}\) = specific heat at constant pressure (J/kg oC),

\(g\) = gravitational acceleration (m/s2)

\(P\) = pressure (Pa),

\(T\) = temperature (oC), and

\(z\) = height (m).

\subsubsection{Example}\label{example-412}

The following example shows the WATJT keyword for when the thermal
option has been activated by the THERMAL keyword in the RUNSPEC section,
and for when NTPVT on the TABDIMS keyword in the RUNSPEC section is set
equal to two.

\begin{verbatim}
--
--       WATER JOULE-THOMSON COEFFICIENT (OPM FLOW EXTENSION KEYWORD)
--
--       REF        OIL
--       PRESS      JTC
--       --------   -------
WATJT
         20.0       1*                  / TABLE NO. 01
         20.0       -0.20               / TABLE NO. 02
\end{verbatim}

Here the first entry is defaulted, and the simulator will therefore
calculate the water JTC internally using the data on the WATDENT keyword
in the PROPS section.

There is no terminating ``/'' for this keyword.

\subsection{WATVISCT -- Define Water Viscosity versus Temperature
Functions}\label{watvisct-define-water-viscosity-versus-temperature-functions}

\begin{longtable}[]{@{}
  >{\raggedright\arraybackslash}p{(\linewidth - 14\tabcolsep) * \real{0.1250}}
  >{\raggedright\arraybackslash}p{(\linewidth - 14\tabcolsep) * \real{0.1250}}
  >{\raggedright\arraybackslash}p{(\linewidth - 14\tabcolsep) * \real{0.1250}}
  >{\raggedright\arraybackslash}p{(\linewidth - 14\tabcolsep) * \real{0.1250}}
  >{\raggedright\arraybackslash}p{(\linewidth - 14\tabcolsep) * \real{0.1250}}
  >{\raggedright\arraybackslash}p{(\linewidth - 14\tabcolsep) * \real{0.1250}}
  >{\raggedright\arraybackslash}p{(\linewidth - 14\tabcolsep) * \real{0.1250}}
  >{\raggedright\arraybackslash}p{(\linewidth - 14\tabcolsep) * \real{0.1250}}@{}}
\toprule\noalign{}
\begin{minipage}[b]{\linewidth}\raggedright
RUNSPEC
\end{minipage} & \begin{minipage}[b]{\linewidth}\raggedright
GRID
\end{minipage} & \begin{minipage}[b]{\linewidth}\raggedright
EDIT
\end{minipage} & \begin{minipage}[b]{\linewidth}\raggedright
PROPS
\end{minipage} & \begin{minipage}[b]{\linewidth}\raggedright
REGIONS
\end{minipage} & \begin{minipage}[b]{\linewidth}\raggedright
SOLUTION
\end{minipage} & \begin{minipage}[b]{\linewidth}\raggedright
SUMMARY
\end{minipage} & \begin{minipage}[b]{\linewidth}\raggedright
SCHEDULE
\end{minipage} \\
\midrule\noalign{}
\endhead
\bottomrule\noalign{}
\endlastfoot
\end{longtable}

\subsubsection{Description}\label{description-730}

WATVISCT defines the water viscosity as a function of temperature for
when thermal option has been activated by the THERMAL keywords in the
RUNSPEC. The reference pressure for this table is given by the VISCREF
keyword in the PROPS section.

This keyword can only be used if OPM Flow's thermal option has been
activated by the THERMAL keyword in the RUNSPEC section. Note this is
different to the commercial simulator that uses the TEMP keyword in the
RUNSPEC section to activate the black-oil thermal model.

\begin{longtable}[]{@{}
  >{\raggedright\arraybackslash}p{(\linewidth - 6\tabcolsep) * \real{0.2500}}
  >{\raggedright\arraybackslash}p{(\linewidth - 6\tabcolsep) * \real{0.2500}}
  >{\raggedright\arraybackslash}p{(\linewidth - 6\tabcolsep) * \real{0.2500}}
  >{\raggedright\arraybackslash}p{(\linewidth - 6\tabcolsep) * \real{0.2500}}@{}}
\toprule\noalign{}
\begin{minipage}[b]{\linewidth}\raggedright
No.
\end{minipage} & \begin{minipage}[b]{\linewidth}\raggedright
Name
\end{minipage} & \begin{minipage}[b]{\linewidth}\raggedright
Description
\end{minipage} & \begin{minipage}[b]{\linewidth}\raggedright
Default
\end{minipage} \\
\midrule\noalign{}
\endhead
\bottomrule\noalign{}
\endlastfoot
Field & Metric & Laboratory & \\
1 & TEMP & A columnar vector of real monotonically increasing down the
column values that defines the temperature values. & None \\
oF & oC & oC & \\
2 & VIS & A columnar vector of real decreasing down the column values
that defines the water viscosity for the corresponding temperature
values (TEMP). VIS should be given at the reference pressure defined by
the PRESS variable on the VISCREF keyword. & None \\
cP & cP & cP & \\
Notes: & & & \\
\end{longtable}

\emph{Table 8.199: WATVISCT Keyword Description}

\subsubsection{Example}\label{example-413}

The following example shows the WATVISCT keyword for when the thermal
option has been activated by the THERMAL keyword in the RUNSPEC section
and for when NTPVT on the TABDIMS keyword in the RUNSPEC section is set
equal to one.

\begin{verbatim}

--
--       WATER VISCOSITY VERSUS TEMPERATURE TABLES
--
--       WATER      WATER
--       TEMP       VISC
--       --------   -------
WATVISCT
           100.0    0.625
           110.0    0.620
           120.0    0.580
           150.0    0.550
           165.0    0.500                                  / TABLE NO. 01
\end{verbatim}

There is no terminating ``/'' for this keyword.

\subsection{WSF -- Water Saturation Function Tables (Gas-Water
Systems)}\label{wsf-water-saturation-function-tables-gas-water-systems}

\begin{longtable}[]{@{}
  >{\raggedright\arraybackslash}p{(\linewidth - 14\tabcolsep) * \real{0.1250}}
  >{\raggedright\arraybackslash}p{(\linewidth - 14\tabcolsep) * \real{0.1250}}
  >{\raggedright\arraybackslash}p{(\linewidth - 14\tabcolsep) * \real{0.1250}}
  >{\raggedright\arraybackslash}p{(\linewidth - 14\tabcolsep) * \real{0.1250}}
  >{\raggedright\arraybackslash}p{(\linewidth - 14\tabcolsep) * \real{0.1250}}
  >{\raggedright\arraybackslash}p{(\linewidth - 14\tabcolsep) * \real{0.1250}}
  >{\raggedright\arraybackslash}p{(\linewidth - 14\tabcolsep) * \real{0.1250}}
  >{\raggedright\arraybackslash}p{(\linewidth - 14\tabcolsep) * \real{0.1250}}@{}}
\toprule\noalign{}
\begin{minipage}[b]{\linewidth}\raggedright
RUNSPEC
\end{minipage} & \begin{minipage}[b]{\linewidth}\raggedright
GRID
\end{minipage} & \begin{minipage}[b]{\linewidth}\raggedright
EDIT
\end{minipage} & \begin{minipage}[b]{\linewidth}\raggedright
PROPS
\end{minipage} & \begin{minipage}[b]{\linewidth}\raggedright
REGIONS
\end{minipage} & \begin{minipage}[b]{\linewidth}\raggedright
SOLUTION
\end{minipage} & \begin{minipage}[b]{\linewidth}\raggedright
SUMMARY
\end{minipage} & \begin{minipage}[b]{\linewidth}\raggedright
SCHEDULE
\end{minipage} \\
\midrule\noalign{}
\endhead
\bottomrule\noalign{}
\endlastfoot
\end{longtable}

\subsubsection{Description}\label{description-731}

The WSF keyword defines the water relative permeability data versus
water saturation tables for when only gas and water are present in the
input deck. This keyword should only be used if the gas and water phases
are present in the run, and can therefore also be used with the CO2STORE
and \hyperref[REF_HEADING_KEYWORD_H2STORE]{H2STORE} models. In addition,
the keyword must be used in conjunction with the GSF keyword in the
PROPS section, that defines the gas relative permeability and gas-water
capillary pressure data versus gas saturation for gas-water systems.

\begin{tcolorbox}[enhanced jigsaw, title=\textcolor{quarto-callout-note-color}{\faInfo}\hspace{0.5em}{Note}, coltitle=black, breakable, toprule=.15mm, colbacktitle=quarto-callout-note-color!10!white, leftrule=.75mm, opacitybacktitle=0.6, colback=white, bottomtitle=1mm, toptitle=1mm, titlerule=0mm, left=2mm, colframe=quarto-callout-note-color-frame, arc=.35mm, rightrule=.15mm, bottomrule=.15mm, opacityback=0]

WSF is a compositional keyword in the commercial compositional
simulator, and will therefore cause an error in the commercial black-oil
simulator. Currently, both the GSF and WSF keywords can only be used
with the CO2STORE and \hyperref[REF_HEADING_KEYWORD_H2STORE]{H2STORE}
models.

\end{tcolorbox}

\begin{longtable}[]{@{}
  >{\raggedright\arraybackslash}p{(\linewidth - 6\tabcolsep) * \real{0.2500}}
  >{\raggedright\arraybackslash}p{(\linewidth - 6\tabcolsep) * \real{0.2500}}
  >{\raggedright\arraybackslash}p{(\linewidth - 6\tabcolsep) * \real{0.2500}}
  >{\raggedright\arraybackslash}p{(\linewidth - 6\tabcolsep) * \real{0.2500}}@{}}
\toprule\noalign{}
\begin{minipage}[b]{\linewidth}\raggedright
No.
\end{minipage} & \begin{minipage}[b]{\linewidth}\raggedright
Name
\end{minipage} & \begin{minipage}[b]{\linewidth}\raggedright
Description
\end{minipage} & \begin{minipage}[b]{\linewidth}\raggedright
Default
\end{minipage} \\
\midrule\noalign{}
\endhead
\bottomrule\noalign{}
\endlastfoot
Field & Metric & Laboratory & \\
1 & SWAT & A columnar vector of real values monotonically increasing
down the column starting from the connate water saturation and
terminating at one, that defines the water saturation. & None \\
dimensionless & dimensionless & dimensionless & \\
2 & KRW & A columnar vector of real values that are either equal or
increasing down the column and that are greater than or equal to zero
and less than or equal to one that defines the water relative
permeability with respect to water saturation. Note that the first entry
in the column must be zero. & None \\
dimensionless & dimensionless & dimensionless & \\
Notes: & & & \\
\end{longtable}

\emph{Table 8.200: WSF Keyword Description}

See also the GSF - Gas Saturation Tables versus Gas Saturation
(Gas-Water and CO2STORE Systems) keyword in the PROPS section, that
defines the gas relative permeability and gas-water capillary pressure
data as a function of gas saturation, for when only the gas and water
water phases are present in the model.

\subsubsection{Example}\label{example-414}

\begin{verbatim}
--
--       WATER RELATIVE PERMEABILITY TABLES (OPM FLOW KEYWORD)
--
WSF
--       SWAT       KRW
--       FRAC
--       --------   --------
         0.200000   0.0000
         0.238616   0.0002
         0.245309   0.0004
         0.261989   0.0010
         0.303091   0.0044
         0.368269   0.0191
         0.435026   0.0519
         0.486387   0.0940
         0.550683   0.1725
         0.575342   0.2115
         0.599076   0.2542
         0.621294   0.2991
         0.642171   0.3458
         0.658984   0.3868
         0.671123   0.4183
         0.679268   0.4403
         0.684963   0.4562
         0.688893   0.4674
         0.692025   0.4765
         0.694641   0.4841
         0.696976   0.4910
         0.699099   0.4973
         0.700000   0.5000
         1.000000   0.9000             / TABLE NO. 01
--       --------   --------
         0.200000   0.0000
         0.238616   0.0002
         0.245309   0.0004
         0.261989   0.0010
         0.303091   0.0044
         0.368269   0.0191
         0.435026   0.0519
         0.486387   0.0940
         0.550683   0.1725
         0.575342   0.2115
         0.599076   0.2542
         0.621294   0.2991
         0.642171   0.3458
         0.658984   0.3868
         0.671123   0.4183
         0.679268   0.4403
         0.684963   0.4562
         0.688893   0.4674
         0.692025   0.4765
         0.694641   0.4841
         0.696976   0.4910
         0.699099   0.4973
         0.700000   0.5000
         1.000000   0.9000             / TABLE NO. 01
\end{verbatim}

The example defines two WSF tables for use when only gas and water are
present in the run.

\subsection{ZMFVD -- Define Compositional Components versus
Depth}\label{zmfvd-define-compositional-components-versus-depth}

\begin{longtable}[]{@{}
  >{\raggedright\arraybackslash}p{(\linewidth - 14\tabcolsep) * \real{0.1250}}
  >{\raggedright\arraybackslash}p{(\linewidth - 14\tabcolsep) * \real{0.1250}}
  >{\raggedright\arraybackslash}p{(\linewidth - 14\tabcolsep) * \real{0.1250}}
  >{\raggedright\arraybackslash}p{(\linewidth - 14\tabcolsep) * \real{0.1250}}
  >{\raggedright\arraybackslash}p{(\linewidth - 14\tabcolsep) * \real{0.1250}}
  >{\raggedright\arraybackslash}p{(\linewidth - 14\tabcolsep) * \real{0.1250}}
  >{\raggedright\arraybackslash}p{(\linewidth - 14\tabcolsep) * \real{0.1250}}
  >{\raggedright\arraybackslash}p{(\linewidth - 14\tabcolsep) * \real{0.1250}}@{}}
\toprule\noalign{}
\begin{minipage}[b]{\linewidth}\raggedright
RUNSPEC
\end{minipage} & \begin{minipage}[b]{\linewidth}\raggedright
GRID
\end{minipage} & \begin{minipage}[b]{\linewidth}\raggedright
EDIT
\end{minipage} & \begin{minipage}[b]{\linewidth}\raggedright
PROPS
\end{minipage} & \begin{minipage}[b]{\linewidth}\raggedright
REGIONS
\end{minipage} & \begin{minipage}[b]{\linewidth}\raggedright
SOLUTION
\end{minipage} & \begin{minipage}[b]{\linewidth}\raggedright
SUMMARY
\end{minipage} & \begin{minipage}[b]{\linewidth}\raggedright
SCHEDULE
\end{minipage} \\
\midrule\noalign{}
\endhead
\bottomrule\noalign{}
\endlastfoot
\end{longtable}

\subsubsection{Description}\label{description-732}

The ZMFVD keyword defines the compositional component mole fractions,
for each component, as a function of depth, as such the keyword should
have the same number of component columnar vectors as that declared via
the COMPS keyword in the RUNSPEC section, and the NCOMPS keyword in the
PROPS section. The keyword should only be used if the CO2STORE and
GASWAT keywords in the RUNSPEC section have also be activated for the
gas-water two component model.

\begin{tcolorbox}[enhanced jigsaw, title=\textcolor{quarto-callout-note-color}{\faInfo}\hspace{0.5em}{Note}, coltitle=black, breakable, toprule=.15mm, colbacktitle=quarto-callout-note-color!10!white, leftrule=.75mm, opacitybacktitle=0.6, colback=white, bottomtitle=1mm, toptitle=1mm, titlerule=0mm, left=2mm, colframe=quarto-callout-note-color-frame, arc=.35mm, rightrule=.15mm, bottomrule=.15mm, opacityback=0]

This is an OPM Flow keyword used with OPM Flow's CO2STORE and GASWAT
keywords in the RUNSPEC section, and should not be confused with the
more general version of the ZMFVD keyword used in the commercial
compositional simulator.

\end{tcolorbox}

\begin{longtable}[]{@{}
  >{\raggedright\arraybackslash}p{(\linewidth - 6\tabcolsep) * \real{0.2500}}
  >{\raggedright\arraybackslash}p{(\linewidth - 6\tabcolsep) * \real{0.2500}}
  >{\raggedright\arraybackslash}p{(\linewidth - 6\tabcolsep) * \real{0.2500}}
  >{\raggedright\arraybackslash}p{(\linewidth - 6\tabcolsep) * \real{0.2500}}@{}}
\toprule\noalign{}
\begin{minipage}[b]{\linewidth}\raggedright
No.
\end{minipage} & \begin{minipage}[b]{\linewidth}\raggedright
Name
\end{minipage} & \begin{minipage}[b]{\linewidth}\raggedright
Description
\end{minipage} & \begin{minipage}[b]{\linewidth}\raggedright
Default
\end{minipage} \\
\midrule\noalign{}
\endhead
\bottomrule\noalign{}
\endlastfoot
Field & Metric & Laboratory & \\
1 & DEPTH & A columnar vector of real monotonically increasing down the
column values that defines the depth values for the corresponding
compositional component mole fractions. The default number of DEPTH
values is 20, as defined by the NDRXVD parameter on the EQLDIMS keyword
in the RUNSPEC section, and which may be used to reset the number of
DEPTH values. & None \\
feet & m & cm & \\
2 & ZCOMP & A series of columnar vectors, with each columnar vector
representing a compositional component mole fraction as a function of
DEPTH. In addition, the sum of the compositional component mole
fractions must sum to one for a given DEPTH value, otherwise an error
will occur. Secondly, if the composition is independent of depth, then
only one single row may be entered. Note that the number of columnar
vectors, should be the same as that entered via the NCOMPS keyword in
the PROPS section, and the COMPS keyword in the RUNSPEC section.
Finally, only the default value of two components are currently
supported by OPM Flow. & None \\
mole fraction & mole fraction & mole fraction & \\
Notes: & & & \\
\end{longtable}

\emph{Table 8.201: ZMFVD Keyword Description}

\subsubsection{Example}\label{example-415}

The following example defines how to confirm a two component
formulation, together with defining the names of the composition
components, as well as the compositional gradient, to be used with the
CO2STORE and GASWAT options.

\begin{verbatim}
--
--       CONFIRM NUMBER OF COMPOSITIONAL COMPONENTS (OPM FLOW KEYWORD)
--
NCOMPS
                2                                                                      /
--
--       DEFINE COMPOSITIONAL COMPONENTS NAMES (OPM FLOW KEYWORD)
--
CNAMES
         'CO2'
         'H2O'                                                                  /
--
--       COMPOSITIONAL COMPONENT MOLE FRACTIONS VS DEPTH (OPM FLOW KEYWORD)
--
--       DEPTH   CO2   H2O
ZMFVD
         2000    0.0   1.0
         2100    0.0   1.0
/
\end{verbatim}

\bookmarksetup{startatroot}

\chapter{REGIONS SECTION}\label{regions-section}

\section{Introduction}\label{introduction-7}

The REGIONS section defines how various properties in the PROPS and
SOLUTION sections are allocated to individual cells within the model, as
well as defining various fluid in-place reporting regions. This is
accomplished by assigning an integer value to each cell that represents
the data set of the property to be assigned to the grid block.

\section{Keyword Definitions}\label{keyword-definitions-11}

\section{Data Requirements}\label{data-requirements-5}

OPM Flow, like most numerical modeling software, uses a default value of
one for the various region arrays and thus if there is only one PVT data
set for example, then there is no need to define the region array
associated with allocating the PVT tables (via the PVTNUM keyword in the
REGIONS section), as all cells will be allocated PVT table number one.
However, if there are more than one PVT table entered in the PROPS
section and PVTNUM is not defined in the model, then PVT tables greater
than one will not be used and there will be no warning message
indicating the fact.

Table 9.1 summarizes the various property region data and their
associated keywords.

\begin{longtable}[]{@{}
  >{\raggedright\arraybackslash}p{(\linewidth - 6\tabcolsep) * \real{0.2500}}
  >{\raggedright\arraybackslash}p{(\linewidth - 6\tabcolsep) * \real{0.2500}}
  >{\raggedright\arraybackslash}p{(\linewidth - 6\tabcolsep) * \real{0.2500}}
  >{\raggedright\arraybackslash}p{(\linewidth - 6\tabcolsep) * \real{0.2500}}@{}}
\toprule\noalign{}
\begin{minipage}[b]{\linewidth}\raggedright
No.
\end{minipage} & \begin{minipage}[b]{\linewidth}\raggedright
REGIONS Keyword
\end{minipage} & \begin{minipage}[b]{\linewidth}\raggedright
Standard Property Regions
\end{minipage} & \begin{minipage}[b]{\linewidth}\raggedright
Property Section
\end{minipage} \\
\midrule\noalign{}
\endhead
\bottomrule\noalign{}
\endlastfoot
1 & EQLNUM & EQLNUM -- Define the Equilibration Region Numbers.
Equilibrium region allocation based on the EQUIL keyword in the SOLUTION
section. & SOLUTION \\
2 & FIPNUM & FIPNUM -- Define the Fluid In-Place Region Numbers. Fluid
In-Place reporting via the FIPNUM array that divides the model into
different fluid in-place reporting regions. & REGIONS \\
3 & PVTNUM & PVTNUM -- Define the PVT Regions. PVT table allocation of
the DENSITY, PVDG, PVDO, PVTG, PVTO, PVCO, PVTW, and ROCK tables. &
PROPS \\
4 & SATNUM & SATNUM -- Define the Saturation Table Region Numbers.
Saturation (relative permeability) table allocation of the SGFN, SWFN,
SOF2, SOF3, SGOF, and SWOF tables. This array is for the standard
drainage, as oppose to the imbibition curves, used in all models. There
are separate region arrays for alternative relative permeability curves.
& PROPS \\
& Function Specific Regions & & \\
5 & ENDNUM & ENDNUM -- Define the End-Point Scaling Depth Region
Numbers. ENPTVD and ENKRVD versus depth table allocation for when
ENDSCALE option has been activated in the RUNSPEC section. & PROPS \\
6 & FIP & FIP -- Define the Fluid In-Place Names and Region and Numbers.
Fluid In-Place reporting via the FIP array that divides the model into
different fluid in-place reporting regions. & REGIONS \\
7 & FIPOWG & FIPOWG -- Activate Oil, Gas, and Water FIP Zone Reporting.
Fluid In-Place reporting for the total oil, gas, and water in the model.
There is no region property array for this keyword, as the simulator
calculates the total phase volumes based on the contacts entered via the
EQUIL keyword in the SOLUTION section. & REGIONS \\
8 & HBNUM & HBNUM -- Define Herschel-Bulkley Region Numbers.
Herschel-Bulkley rheological property table region numbers. & PROPS \\
9 & HWSNUM & HWSNUM -- Define the Saturation Table Region Numbers (High
Salinity and Water Wet). High salinity water wet saturation table
allocation using the high salinity water wet saturation SWFN and SOFN
tables. This is an alias for the SURFWNUM property region which is also
not supported by OPM Flow. & PROPS \\
10 & IMBNUM & IMBNUM -- Define the Imbibition Saturation Table Region
Numbers. Imbibition saturation table allocation of the SWFN, SOF2, SOF3,
or SWOF imbibition tables. & PROPS \\
11 & IMBNUMMF & IMBNUMMF -- Define the Imbibition Saturation Table
Region Numbers (Matrix-Fracture). Imbibition saturation tables region
numbers for flow between the matrix and fracture blocks, for when the
Hysteresis option has been activated in Dual Porosity or Dual
Permeability runs. & PROPS \\
12 & KRNUM & KRNUM -- Define the Directional Saturation Table Region
Numbers. Direction dependent saturation tables region numbers for when
Directional Dependent Saturation Function option has been activated by
the DIRECT parameter on the SATOPTS. & PROPS \\
13 & KRNUMMF & KRNUMMF -- Define the Saturation Table Region Numbers
(Matrix-Fracture). Drainage saturation tables region numbers for flow
between the matrix and fracture blocks in a Dual Porosity or Dual
Permeability model. & PROPS \\
14 & LSLTWNUM & LSLTWNUM -- Define the Low Salt Water Wet Saturation
Table Region Numbers. Low salt water wet saturation tables region
numbers for when the Low Salinity option for the Brine model and the
Surfactant Wettability option have been activated for the model. &
PROPS \\
15 & LSNUM & LSNUM -- Define the Low Salt Oil Wet Saturation Table
Region Numbers. Low salt oil wet saturation tables region numbers for
the Low Salinity Brine model. & PROPS \\
16 & LWSLTNUM & LWSLTNUM -- Define the Low Salt Oil Wet Saturation Table
Region Numbers. Low salt oil wet saturation table region numbers for the
Low Salinity Brine model. & PROPS \\
17 & LWSNUM & LWSNUM -- Define the Low Salt Water Wet Saturation Table
Region Numbers. Low salt water wet saturation table region numbers for
when the Low Salinity option for the Brine model and the Surfactant
Wettability option have been activated for the model. & PROPS \\
18 & MISCNUM & MISCNUM -- Define the Miscibility Region Numbers.
Miscible regions based on the TLMIXPAR records when the MISCIBLE or
SOLVENT keywords have been activated in the RUNSPEC section. & PROPS \\
19 & OPERNUM & OPERNUM -- Define Regions for Mathematical Operations on
Arrays. OPERATER keyword region numbers that can also be used with the
EQUALREG, ADDREG, COPYREG, MULTIREG, MULTREGP, and MULTREGT keywords, as
well as the OPERATER keyword in calculating various grid properties in
the GRID and REGIONS section. & GRID EDIT PROPS REGIONS SOLUTION \\
20 & PENUM & PENUM -- Define the Petro-Elastic Region Numbers.
Petro-elastic region numbers for petro-elastic coefficients, bulk
modulus functions and shear modulus functions. & PROPS \\
21 & PLMIXNUM & PLMIXNUM -- Define the Polymer Region Numbers. Polymer
region numbers for the mixing tables, and the maximum polymer and salt
concentrations for the Polymer option. & PROPS \\
22 & RESIDNUM & RESIDNUM -- Define Vertical Equilibrium Residual Flow
Region Numbers. Vertical Equilibrium residual flow calculation
saturation tables or when region numbers for the Vertical Equilibrium
option. & PROPS \\
23 & ROCKNUM & ROCKNUM -- Define Rock Compaction Table Region Numbers.
Rock compaction table allocation for when the ROCKCOMP keyword as been
activated in the RUNSPEC section, that allocates the ROCKTAB series of
tables to a cell. & PROPS \\
24 & SURFNUM & SURFNUM -- Define the Surfactant Miscible Saturation
Table Region Numbers. Surfactant saturation (relative permeability)
tables allocation allocating the SWFN, SOF2, SOF3, or SWOF as miscible
tables. & PROPS \\
25 & SURFWNUM & SURFWNUM -- Define the Saturation Table Region Numbers
(High Salinity and Water Wet). High salinity water wet saturation table
allocation using the high salinity water wet saturation SWFN and SOFN
tables. & PROPS \\
26 & TNUM & TNUM -- Define Passive Tracer Concentration Regions.
Partitioned Tracer option region numbers for this option. & PROPS \\
27 & TRKPF & TRKPF -- Define Partitioned Tracer Regions. Defines the
regions associated with partition tracers and the partitioning tables
allocated to grid blocks for when the Partitioned Tracer option has been
enabled by the PARTTRAC keyword in the RUNSPEC section. & PROPS \\
28 & WH2NUM & WH2NUM -- Define WAG Hysteresis Saturation Table Region
Numbers (Two Phase). Two phase Water-Alternating-Gas (``WAG'')
hysteresis saturation table region numbers for when the hysteresis
option has been activated by the WAGHYSTR variable on the SATOPTS
keyword in the RUNSPEC. & PROPS \\
29 & WH3NUM & WH3NUM -- Define WAG Hysteresis Saturation Table Region
Numbers (Three Phase). Three phase WAG hysteresis saturation table
region numbers for when the hysteresis option has been activated by the
WAGHYSTR variable on the SATOPTS keyword in the RUNSPEC. & PROPS \\
Notes: & & & \\
\end{longtable}

\emph{Table 9.1: REGIONS Section Allocation Array Summary}

\begin{tcolorbox}[enhanced jigsaw, title=\textcolor{quarto-callout-note-color}{\faInfo}\hspace{0.5em}{Note}, coltitle=black, breakable, toprule=.15mm, colbacktitle=quarto-callout-note-color!10!white, leftrule=.75mm, opacitybacktitle=0.6, colback=white, bottomtitle=1mm, toptitle=1mm, titlerule=0mm, left=2mm, colframe=quarto-callout-note-color-frame, arc=.35mm, rightrule=.15mm, bottomrule=.15mm, opacityback=0]

Earth modeling software may export the region arrays to be loaded into
the simulator with values starting from zero instead of one. This may
result in an error message if these cells are active in the model, if
the cells are inactive then the zero values will be ignored by the
simulator.

\end{tcolorbox}

A complete property data set associated with a region keyword is
required for each region number allocated to the cells. For example, if
the fluid properties for the model are the same (for example, PVTO and
PVDG keyword data), but the rock compressibility is varying with depth
resulting in, say three different ROCK keyword records, then there has
to be three PVTNUM regions and three complete PVT data sets in order to
allocate the three ROCK records. This would mean that the PVTO and PVDG
records, in this instance, would have to be repeated three times to
match the three ROCK keyword records.

\pandocbounded{\includegraphics[keepaspectratio]{chapters/images/Frame16_20913f22c3f7.png}}
\pandocbounded{\includegraphics[keepaspectratio]{chapters/images/Image14_20913f22c3f7.png}}

Example SATNUM and EQLNUM arrays from the Volve {[}The Volve Data was
approved for data sharing in 2018 by the initiative of the last
Operating company, Equinor and approved by the license partners
ExxonMobil E\&P Norway AS and Bayerngas Norge AS in the end of 2017.{]}
field are displayed in Figure 9.1 and Figure 9.2, respectively.

\pandocbounded{\includegraphics[keepaspectratio]{chapters/images/Frame15_9c05c1ebe79a.png}}
\pandocbounded{\includegraphics[keepaspectratio]{chapters/images/Image13_9c05c1ebe79a.png}}

\section{Keyword Definitions}\label{keyword-definitions-12}

\subsection{ADD -- Add a Constant to a Specified
Array}\label{add-add-a-constant-to-a-specified-array-3}

The ADD keyword adds a constant to a specified array or part of an
array. The constant can be real or integer depending on the array type;
however, the arrays that can be operated on is dependent on which
section the ADD keyword is being applied.

See ADD -- Add a Constant to a Specified Array in the GRID section for a
full description.

\subsection{ADDREG -- Add a Constant to an Array based on a Region
Number}\label{addreg-add-a-constant-to-an-array-based-on-a-region-number-3}

The ADDREG keyword adds a constant to a specified array or part of an
array based on cells with a specific region number. The region number
array can be FLUXNUM, MULTNUM or OPERNUM and these arrays must be
defined and be available before the ADDREG keyword is read by the
simulator. The constant can be real or integer depending on the property
array type; however, the property arrays that can be operated on is
dependent on which section the ADDREG keyword is being used.

See ADDREG -- Add a Constant to an Array based on a Region Number in the
GRID section for a full description.

\subsection{BOUNDARY -- Define a Boundary Box for
Printing}\label{boundary-define-a-boundary-box-for-printing-2}

The BOUNDARY keyword defines a rectangular grid for printing various
arrays to the output print file (*.PRT); thus, avoiding printing all the
elements in the selected array.

See BOUNDARY -- Define a Boundary Box for Printing in the GRID section
for a full description.

\subsection{BOX -- Define a Range of Grid Blocks to Enter Property
Data}\label{box-define-a-range-of-grid-blocks-to-enter-property-data-3}

BOX defines a range of grid blocks for which subsequent data is assigned
for all the cells in the defined BOX. Note that the BOX grid is reset by
the keyword ENDBOX by resetting the current defined BOX to be the whole
grid. The keyword can be used for any array and for all grid types.

See BOX - Define a Range of Grid Blocks to Enter Property Data in the
GRID section for a full description.

\subsection{COLUMNS -- Define Input File Column
Margins}\label{columns-define-input-file-column-margins-5}

The COLUMNS keyword defines the input file column margins; characters
outside the margins are ignored by the input parser.

See COLUMNS -- Define Input File Column Margins in the GLOBAL section
for a full description.

\subsection{COPY -- Copy Array Data to Another
Array}\label{copy-copy-array-data-to-another-array-3}

The COPY keyword copies an array (or part of an array) to another array
or part of an array. The arrays can be real or integer depending on the
array type; however, the arrays that can be operated on is dependent on
which section the COPY keyword is being used.

See COPY -- Copy Array Data to Another Array in the GRID section for a
full description.

\subsection{COPYBOX -- Copy Array Data Defined by a
Box}\label{copybox-copy-array-data-defined-by-a-box-2}

The COPYBOX keyword copies an array (or part of an array) to part of the
same array. The array can be real or integer depending on the array
type; however, the arrays that can be operated on is dependent on which
section the COPYBOX keyword is being used.

See COPYBOX -- Copy Array Data Defined by a Box in the GRID section for
a full description.

\subsection{COPYREG -- Copy an Array to Another Array based on a Region
Number}\label{copyreg-copy-an-array-to-another-array-based-on-a-region-number-3}

The COPYREG keyword copies a specified array or part of an array based
on cells with a specific region number to another array. The region
number array can be FLUXNUM, MULTNUM or OPERNUM and these arrays must be
defined and be available before the COPYREG keyword is read by the
simulator. The property arrays can be real or integer depending on the
property array type; however, the property arrays that can be operated
on is dependent on which section the COPYREG keyword is being used.

See COPYREG -- Copy an Array to Another Array based on a Region Number
in the GRID section for a full description.

\subsection{DEBUG -- Define the Debug Data to be Printed to
File}\label{debug-define-the-debug-data-to-be-printed-to-file-5}

This keyword defines the debug data to be written to the debug file
(*.DBG). This keyword is not supported by OPM Flow but has no effect on
the results so it will be ignored.

See DEBUG -- Define the Debug Data to be Printed to File in the GLOBAL
section for a full description.

\subsection{ECHO -- Activate Echoing of User Input Files to the Print
File}\label{echo-activate-echoing-of-user-input-files-to-the-print-file-5}

Turns on echoing of all the input files to the print file; note that
this keyword is activated by default and can subsequently be switched
off by the NOECHO activation keyword.

See ECHO -- Activate Echoing of User Input Files to the Print File in
the GLOBAL section for a full description.

\subsection{END -- Define the End of the Input
File}\label{end-define-the-end-of-the-input-file-5}

This keyword marks the end of the input file and can occur in any
section. Any keywords and data after the END keyword are ignored.

See END -- Define the End of the Input File in the GLOBAL section for a
full description.

\subsection{ENDBOX -- Define the End of the BOX Defined
Grid}\label{endbox-define-the-end-of-the-box-defined-grid-3}

This keyword marks the end of a previously defined BOX sub-grid as
defined by a previously entered BOX keyword. The keyword resets the
input grid to be the full grid as defined by the NX, NY, and NZ
variables on the DIMENS keyword in the RUNSPEC section.

See ENDBOX -- Define the End of the BOX Defined Grid in the GRID section
for a full description.

\subsection{ENDFIN -- End the Definition of a Local Grid
Refinement}\label{endfin-end-the-definition-of-a-local-grid-refinement-3}

ENDFIN defines the end of a Cartesian or radial local grid refinement
(``LGR'') definition and a LGR property definition data set.

See ENDFIN -- End the Definition of a Local Grid Refinement in the GRID
section for a full description.

\subsection{ENDINC -- Define the End of an Include
File}\label{endinc-define-the-end-of-an-include-file-5}

This keyword marks the end of an include file specified on the INCLUDE
keyword. When the ENDINC keyword is encountered in the INCLUDE file,
input data is read from the next keyword in the current file. Any
keywords and data after the ENDINC keyword in the INCLUDE file are
ignored.

See ENDINC -- Define the End of an Include File in the GLOBAL section
for a full description.

\subsection{ENDNUM -- Define the End-Point Scaling Depth Region
Numbers}\label{endnum-define-the-end-point-scaling-depth-region-numbers}

\begin{longtable}[]{@{}
  >{\raggedright\arraybackslash}p{(\linewidth - 14\tabcolsep) * \real{0.1250}}
  >{\raggedright\arraybackslash}p{(\linewidth - 14\tabcolsep) * \real{0.1250}}
  >{\raggedright\arraybackslash}p{(\linewidth - 14\tabcolsep) * \real{0.1250}}
  >{\raggedright\arraybackslash}p{(\linewidth - 14\tabcolsep) * \real{0.1250}}
  >{\raggedright\arraybackslash}p{(\linewidth - 14\tabcolsep) * \real{0.1250}}
  >{\raggedright\arraybackslash}p{(\linewidth - 14\tabcolsep) * \real{0.1250}}
  >{\raggedright\arraybackslash}p{(\linewidth - 14\tabcolsep) * \real{0.1250}}
  >{\raggedright\arraybackslash}p{(\linewidth - 14\tabcolsep) * \real{0.1250}}@{}}
\toprule\noalign{}
\begin{minipage}[b]{\linewidth}\raggedright
RUNSPEC
\end{minipage} & \begin{minipage}[b]{\linewidth}\raggedright
GRID
\end{minipage} & \begin{minipage}[b]{\linewidth}\raggedright
EDIT
\end{minipage} & \begin{minipage}[b]{\linewidth}\raggedright
PROPS
\end{minipage} & \begin{minipage}[b]{\linewidth}\raggedright
REGIONS
\end{minipage} & \begin{minipage}[b]{\linewidth}\raggedright
SOLUTION
\end{minipage} & \begin{minipage}[b]{\linewidth}\raggedright
SUMMARY
\end{minipage} & \begin{minipage}[b]{\linewidth}\raggedright
SCHEDULE
\end{minipage} \\
\midrule\noalign{}
\endhead
\bottomrule\noalign{}
\endlastfoot
\end{longtable}

\subsubsection{Description}\label{description-733}

The ENDNUM keyword defines the end-point scaling depth table region
numbers for each grid block. The end-point scaling depth tables for
various regions are defined by the ENPVTD {[}This keyword is not
supported by OPM Flow but would change the results if supported so the
simulation will be stopped.{]} and the ENKRVD {[}This keyword is not
supported by OPM Flow but would change the results if supported so the
simulation will be stopped.{]} keywords in the PROPS section. In the
RUNSPEC section the NTENDP variable on the ENDSCALE keyword defines the
maximum number of depth tables.

\begin{longtable}[]{@{}
  >{\raggedright\arraybackslash}p{(\linewidth - 6\tabcolsep) * \real{0.2500}}
  >{\raggedright\arraybackslash}p{(\linewidth - 6\tabcolsep) * \real{0.2500}}
  >{\raggedright\arraybackslash}p{(\linewidth - 6\tabcolsep) * \real{0.2500}}
  >{\raggedright\arraybackslash}p{(\linewidth - 6\tabcolsep) * \real{0.2500}}@{}}
\toprule\noalign{}
\begin{minipage}[b]{\linewidth}\raggedright
No.
\end{minipage} & \begin{minipage}[b]{\linewidth}\raggedright
Name
\end{minipage} & \begin{minipage}[b]{\linewidth}\raggedright
Description
\end{minipage} & \begin{minipage}[b]{\linewidth}\raggedright
Default
\end{minipage} \\
\midrule\noalign{}
\endhead
\bottomrule\noalign{}
\endlastfoot
1 & ENDNUM & ENDNUM defines an array of positive integers assigning a
grid cell to a particular end-point scaling depth table region. The
maximum number of ENDNUM regions is set by the NTENDP variable on the
ENDSCALE keyword in the RUNSPEC section. & 1 \\
Notes: & & & \\
\end{longtable}

\emph{Table 9.2: ENDNUM Keyword Description}

This keyword is not supported by OPM Flow but would change the results
if supported so the simulation will be stopped.

\subsubsection{Examples}\label{examples-47}

The example below sets three ENDNUM regions for a 4 x 5 x 2 model.

\begin{verbatim}
--
--       DEFINE ENDNUM REGIONS FOR ALL CELLS
--
ENDNUM
         2  2  1  1  2  2  1  1  1  1  1  1  1  1  1  1  1  1  1  1
         3  3  1  1  3  3  1  1  1  1  1  1  1  1  1  1  1  1  1  1
/
\end{verbatim}

Alternatively the EQUALS keyword could be employed to accomplish the
same task, that is:

\begin{verbatim}
--
--       ARRAY       CONSTANT     ---------- BOX ---------
--                                I1  I2   J1  J2   K1  K2
EQUALS
         ENDNUM      1            1*  1*   1*  1*   1*  1* / SET REGION 1
         ENDNUM      2            1   2    1   2    1   1  / SET REGION 2
         ENDNUM      3            1   2    1   2    2   2  / SET REGION 3
/
\end{verbatim}

\subsection{ENDSKIP -- Deactivate Skipping of Keywords and Input
Data}\label{endskip-deactivate-skipping-of-keywords-and-input-data-5}

The ENDSKIP keyword deactivates the skipping of keywords that was
activated by the SKIP, SKIP100, or SKIP300 keywords. Each SKIP keyword
should be paired with an ENDSKIP keyword.

See ENDSKIP -- Deactivate Skipping of Keywords and Input Data in the
GLOBAL section for a full description.

\subsection{EOSNUM -- Define the Equation of State Region
Numbers}\label{eosnum-define-the-equation-of-state-region-numbers}

\begin{longtable}[]{@{}
  >{\raggedright\arraybackslash}p{(\linewidth - 14\tabcolsep) * \real{0.1250}}
  >{\raggedright\arraybackslash}p{(\linewidth - 14\tabcolsep) * \real{0.1250}}
  >{\raggedright\arraybackslash}p{(\linewidth - 14\tabcolsep) * \real{0.1250}}
  >{\raggedright\arraybackslash}p{(\linewidth - 14\tabcolsep) * \real{0.1250}}
  >{\raggedright\arraybackslash}p{(\linewidth - 14\tabcolsep) * \real{0.1250}}
  >{\raggedright\arraybackslash}p{(\linewidth - 14\tabcolsep) * \real{0.1250}}
  >{\raggedright\arraybackslash}p{(\linewidth - 14\tabcolsep) * \real{0.1250}}
  >{\raggedright\arraybackslash}p{(\linewidth - 14\tabcolsep) * \real{0.1250}}@{}}
\toprule\noalign{}
\begin{minipage}[b]{\linewidth}\raggedright
RUNSPEC
\end{minipage} & \begin{minipage}[b]{\linewidth}\raggedright
GRID
\end{minipage} & \begin{minipage}[b]{\linewidth}\raggedright
EDIT
\end{minipage} & \begin{minipage}[b]{\linewidth}\raggedright
PROPS
\end{minipage} & \begin{minipage}[b]{\linewidth}\raggedright
REGIONS
\end{minipage} & \begin{minipage}[b]{\linewidth}\raggedright
SOLUTION
\end{minipage} & \begin{minipage}[b]{\linewidth}\raggedright
SUMMARY
\end{minipage} & \begin{minipage}[b]{\linewidth}\raggedright
SCHEDULE
\end{minipage} \\
\midrule\noalign{}
\endhead
\bottomrule\noalign{}
\endlastfoot
\end{longtable}

\subsubsection{Description}\label{description-734}

The \hyperref[REF_HEADING_KEYWORD_EOSNUM_9_3]{EOSNUM} keyword defines
the Equation Of State (\hyperref[REF_HEADING_KEYWORD_EOS_5_3]{EOS})
region number for all the cells in the model via an array. The region
number specifies which Equation Of State is used in each grid block. The
keyword should only be used if the compositional mode has been requested
using the COMPS keyword in the RUNSPEC section.

OPM Flow does not currently support the general compositional modeling
formulation.

This keyword is not supported by OPM Flow but it will be parsed and its
data ignored.

\begin{longtable}[]{@{}
  >{\raggedright\arraybackslash}p{(\linewidth - 6\tabcolsep) * \real{0.2500}}
  >{\raggedright\arraybackslash}p{(\linewidth - 6\tabcolsep) * \real{0.2500}}
  >{\raggedright\arraybackslash}p{(\linewidth - 6\tabcolsep) * \real{0.2500}}
  >{\raggedright\arraybackslash}p{(\linewidth - 6\tabcolsep) * \real{0.2500}}@{}}
\toprule\noalign{}
\begin{minipage}[b]{\linewidth}\raggedright
No.
\end{minipage} & \begin{minipage}[b]{\linewidth}\raggedright
Name
\end{minipage} & \begin{minipage}[b]{\linewidth}\raggedright
Description
\end{minipage} & \begin{minipage}[b]{\linewidth}\raggedright
Default
\end{minipage} \\
\midrule\noalign{}
\endhead
\bottomrule\noalign{}
\endlastfoot
1 & \hyperref[REF_HEADING_KEYWORD_EOSNUM_9_3]{EOSNUM} &
\hyperref[REF_HEADING_KEYWORD_EOSNUM_9_3]{EOSNUM} is an array of
positive integers less than or equal to NMEOSR that define the
\hyperref[REF_HEADING_KEYWORD_EOS_5_3]{EOS} region number for each cell
in the model. The NMEOSR variable on the TABDIMS keyword in the RUNSPEC
section defines the number of
\hyperref[REF_HEADING_KEYWORD_EOS_5_3]{EOS} regions in the model. & 1 \\
Notes: & & & \\
\end{longtable}

\emph{Table 9.3.17.1: \hyperref[REF_HEADING_KEYWORD_EOSNUM_9_3]{EOSNUM}
Keyword Description}

\subsubsection{Examples}\label{examples-48}

The example below sets two
\hyperref[REF_HEADING_KEYWORD_EOS_5_3]{EOS}NUM regions for a 4 x 5 x 2
model.

\begin{verbatim}
--
--       DEFINE EOSNUM REGIONS FOR ALL CELLS
--
EOSNUM
         20*1
         20*2
/
\end{verbatim}

\subsection{EQLNUM -- Define the Equilibration Region
Numbers}\label{eqlnum-define-the-equilibration-region-numbers}

\begin{longtable}[]{@{}
  >{\raggedright\arraybackslash}p{(\linewidth - 14\tabcolsep) * \real{0.1250}}
  >{\raggedright\arraybackslash}p{(\linewidth - 14\tabcolsep) * \real{0.1250}}
  >{\raggedright\arraybackslash}p{(\linewidth - 14\tabcolsep) * \real{0.1250}}
  >{\raggedright\arraybackslash}p{(\linewidth - 14\tabcolsep) * \real{0.1250}}
  >{\raggedright\arraybackslash}p{(\linewidth - 14\tabcolsep) * \real{0.1250}}
  >{\raggedright\arraybackslash}p{(\linewidth - 14\tabcolsep) * \real{0.1250}}
  >{\raggedright\arraybackslash}p{(\linewidth - 14\tabcolsep) * \real{0.1250}}
  >{\raggedright\arraybackslash}p{(\linewidth - 14\tabcolsep) * \real{0.1250}}@{}}
\toprule\noalign{}
\begin{minipage}[b]{\linewidth}\raggedright
RUNSPEC
\end{minipage} & \begin{minipage}[b]{\linewidth}\raggedright
GRID
\end{minipage} & \begin{minipage}[b]{\linewidth}\raggedright
EDIT
\end{minipage} & \begin{minipage}[b]{\linewidth}\raggedright
PROPS
\end{minipage} & \begin{minipage}[b]{\linewidth}\raggedright
REGIONS
\end{minipage} & \begin{minipage}[b]{\linewidth}\raggedright
SOLUTION
\end{minipage} & \begin{minipage}[b]{\linewidth}\raggedright
SUMMARY
\end{minipage} & \begin{minipage}[b]{\linewidth}\raggedright
SCHEDULE
\end{minipage} \\
\midrule\noalign{}
\endhead
\bottomrule\noalign{}
\endlastfoot
\end{longtable}

\subsubsection{Description}\label{description-735}

The EQLNUM keyword defines the equilibration region numbers for each
grid block. The equilibration data for various regions are defined in
the SOLUTION section. For example, the EQUIL keyword in the SOLUTION
section defines the initial pressures and fluid contacts for each
equilibration region identified by the EQLNUM region array.

\begin{longtable}[]{@{}
  >{\raggedright\arraybackslash}p{(\linewidth - 6\tabcolsep) * \real{0.2500}}
  >{\raggedright\arraybackslash}p{(\linewidth - 6\tabcolsep) * \real{0.2500}}
  >{\raggedright\arraybackslash}p{(\linewidth - 6\tabcolsep) * \real{0.2500}}
  >{\raggedright\arraybackslash}p{(\linewidth - 6\tabcolsep) * \real{0.2500}}@{}}
\toprule\noalign{}
\begin{minipage}[b]{\linewidth}\raggedright
No.
\end{minipage} & \begin{minipage}[b]{\linewidth}\raggedright
Name
\end{minipage} & \begin{minipage}[b]{\linewidth}\raggedright
Description
\end{minipage} & \begin{minipage}[b]{\linewidth}\raggedright
Default
\end{minipage} \\
\midrule\noalign{}
\endhead
\bottomrule\noalign{}
\endlastfoot
1 & EQLNUM & EQLNUM defines an array of positive integers assigning a
grid cell to a particular equilibration region. The maximum number of
EQLNUM regions is set by the NTEQUL variable on the EQLDIMS keyword in
the RUNSPEC section. & 1 \\
Notes: & & & \\
\end{longtable}

\emph{Table 9.3: EQLNUM Keyword Description}

\subsubsection{Examples}\label{examples-49}

The example below sets three EQLNUM regions for a 4 x 5 x 2 model.

\begin{verbatim}
--
--       DEFINE EQLNUM REGIONS FOR ALL CELLS
--
EQLNUM
         2  2  1  1  2  2  1  1  1  1  1  1  1  1  1  1  1  1  1  1
         3  3  1  1  3  3  1  1  1  1  1  1  1  1  1  1  1  1  1  1
/
\end{verbatim}

Alternatively the EQUALS keyword could be employed to accomplish the
same task, that is:

\begin{verbatim}
--
--       ARRAY       CONSTANT     ---------- BOX ---------
--                                I1  I2   J1  J2   K1  K2
EQUALS
         EQLNUM’     1            1*  1*   1*  1*   1*  1* / SET REGION 1
         EQLNUM’     2            1   2    1   2    1   1  / SET REGION 2
         EQLNUM’     3            1   2    1   2    2   2  / SET REGION 3
/
\end{verbatim}

\subsection{EQUALREG -- Sets an Array to a Constant by Region
Number}\label{equalreg-sets-an-array-to-a-constant-by-region-number-3}

The EQUALREG keyword sets a specified array to a constant for cells with
a specific region number. The region number array can be FLUXNUM,
MULTNUM or OPERNUM and these arrays must be defined and be available
before the EQUALREG keyword is read by the simulator. The constant can
be real or integer depending on the property array type; however, the
property arrays that can be operated on is dependent on which section
the EQUALREG keyword is being used.

See EQUALREG -- Sets an Array to a Constant by Region Number in the GRID
section for a full description.

\subsection{EQUALS -- Sets a Specified Array to a
Constant}\label{equals-sets-a-specified-array-to-a-constant-3}

The EQUALS keyword sets a specified array or part of an array to a
constant. The constant can be real or integer depending on the array
type; however, the arrays that can be operated on is dependent on which
section the EQUALS keyword is being used.

See EQUALS -- Sets a Specified Array to a Constant in the GRID section
for a full description.

\subsection{EXTRAPMS -- Activate Extrapolation Warning
Messages}\label{extrapms-activate-extrapolation-warning-messages-5}

The EXTRAPMS keyword activates extrapolation warning messages for when
OPM Flow extrapolates the PVT or VFP tables. Frequent extrapolation
warning messages should be investigated and resolved as this would
indicate possible incorrect data and may result in the simulator
extrapolating to unrealistic values.

See EXTRAPMS -- Activate Extrapolation Warning Messages in the GLOBAL
section for a full description.

\subsection{FILEUNIT -- Activate Unit Consistency
Verification}\label{fileunit-activate-unit-consistency-verification-2}

The FILEUNIT keyword defines the units of the data set, and is used to
verify that the units in the input deck and any associated include files
are consistent. The keyword does not provide for the conversion between
different sets of units.

See FILEUNIT -- Activate Unit Consistency Checking in the GRID section
for a full description.

\subsection{FIP -- Define the Fluid In-Place Names and Region and
Numbers}\label{fip-define-the-fluid-in-place-names-and-region-and-numbers}

\begin{longtable}[]{@{}
  >{\raggedright\arraybackslash}p{(\linewidth - 14\tabcolsep) * \real{0.1250}}
  >{\raggedright\arraybackslash}p{(\linewidth - 14\tabcolsep) * \real{0.1250}}
  >{\raggedright\arraybackslash}p{(\linewidth - 14\tabcolsep) * \real{0.1250}}
  >{\raggedright\arraybackslash}p{(\linewidth - 14\tabcolsep) * \real{0.1250}}
  >{\raggedright\arraybackslash}p{(\linewidth - 14\tabcolsep) * \real{0.1250}}
  >{\raggedright\arraybackslash}p{(\linewidth - 14\tabcolsep) * \real{0.1250}}
  >{\raggedright\arraybackslash}p{(\linewidth - 14\tabcolsep) * \real{0.1250}}
  >{\raggedright\arraybackslash}p{(\linewidth - 14\tabcolsep) * \real{0.1250}}@{}}
\toprule\noalign{}
\begin{minipage}[b]{\linewidth}\raggedright
RUNSPEC
\end{minipage} & \begin{minipage}[b]{\linewidth}\raggedright
GRID
\end{minipage} & \begin{minipage}[b]{\linewidth}\raggedright
EDIT
\end{minipage} & \begin{minipage}[b]{\linewidth}\raggedright
PROPS
\end{minipage} & \begin{minipage}[b]{\linewidth}\raggedright
REGIONS
\end{minipage} & \begin{minipage}[b]{\linewidth}\raggedright
SOLUTION
\end{minipage} & \begin{minipage}[b]{\linewidth}\raggedright
SUMMARY
\end{minipage} & \begin{minipage}[b]{\linewidth}\raggedright
SCHEDULE
\end{minipage} \\
\midrule\noalign{}
\endhead
\bottomrule\noalign{}
\endlastfoot
\end{longtable}

\subsubsection{Description}\label{description-736}

The FIP keyword defines the fluid in-place name and the associated
region numbers for each grid block. The simulator can print out
summaries of the fluid in-place in each region, the current flow rates
between regions, and the cumulative flows between regions. This keyword
is not in the standard keyword format due to the fluid in-place name
being concatenated to the keyword FIP to fully define the keyword.

Note that the total number of FIP and FIPNUM regions must be defined by
the NMFIPR variable on the REGDIMS keyword, or the NTFIP variable on the
TABDIMS keyword. Both keywords are in the RUNSPEC section.

\begin{longtable}[]{@{}
  >{\raggedright\arraybackslash}p{(\linewidth - 6\tabcolsep) * \real{0.2500}}
  >{\raggedright\arraybackslash}p{(\linewidth - 6\tabcolsep) * \real{0.2500}}
  >{\raggedright\arraybackslash}p{(\linewidth - 6\tabcolsep) * \real{0.2500}}
  >{\raggedright\arraybackslash}p{(\linewidth - 6\tabcolsep) * \real{0.2500}}@{}}
\toprule\noalign{}
\begin{minipage}[b]{\linewidth}\raggedright
No.
\end{minipage} & \begin{minipage}[b]{\linewidth}\raggedright
Name
\end{minipage} & \begin{minipage}[b]{\linewidth}\raggedright
Description
\end{minipage} & \begin{minipage}[b]{\linewidth}\raggedright
Default
\end{minipage} \\
\midrule\noalign{}
\endhead
\bottomrule\noalign{}
\endlastfoot
1 & FIPNAME & A character string of up to eight characters, consisting
of FIP as the first three characters followed by up to a five letter
character string defining the fluid in-place's name. & None \\
2 & FIPNUM & FIPNUM defines an array of positive integers greater than
or equal to one, that assigns a grid cell to a particular fluid in-place
region named by FIPNAME. The maximum number of FIP and FIPNUM regions is
set by the REGDIMS(NMFIPR) or the TABDIMS(NTFIP) keywords(variables) in
the RUNSPEC section. If both REGDIMS(NMFIPR) and TABDIMS(NTFIP) have
been defined then the maximum of the two is used. & 1 \\
Notes: & & & \\
\end{longtable}

\emph{Table 9.4: FIP Keyword Description}

The keyword behaves the same as the FIPNUM keyword except the full name
of the keyword, including the concatenated characters, are used as the
property region name. For example, if we wish define a fluid in-place
region name called UNIT, then the keyword would be FIPUNIT.

\begin{tcolorbox}[enhanced jigsaw, title=\textcolor{quarto-callout-note-color}{\faInfo}\hspace{0.5em}{Note}, coltitle=black, breakable, toprule=.15mm, colbacktitle=quarto-callout-note-color!10!white, leftrule=.75mm, opacitybacktitle=0.6, colback=white, bottomtitle=1mm, toptitle=1mm, titlerule=0mm, left=2mm, colframe=quarto-callout-note-color-frame, arc=.35mm, rightrule=.15mm, bottomrule=.15mm, opacityback=0]

The commercial simulator prints out a fluid in-place report if the FIP
option on the RPTSCHED keyword is set equal to three, that is: FIP=3.
This option is currently not available in OPM Flow.

\end{tcolorbox}

The region property data for FIP arrays can be written to the SUMMARY
file, and the RSM file if requested, similar to the FIPNUM regions, with
some caveats:

\begin{verbatim}
- Only SUMMARY keywords for regions may be used, that is the SUMMARY variable name must begin with the letter R.
- The SUMMARY variable name must consist of a character string length of exactly five characters, if less than five characters, then the “_” character should be used to fill out the SUMMARY variable name. For example, instead of RPRUNIT, use RPR__UNI, or instead of ROIPUNIT, use ROIP_UNI.
- Only the first three characters of the FIP region name should be concatenated with the SUMMARY variable name. This means if the FIP region name is UNIT, the FIP keyword would be FIPUNIT; however, to access the regional pressures for FIPUNIT, one should use RPR__UNI, or to access the regional oil in-place one would use ROIP_UNI.
\end{verbatim}

See also the FIPOWG keyword in the REGIONS section that automatically
defines the fluid-in-place regions at the start of the run based the
gas, oil and water zones at the time the model was initialized.

\subsubsection{Example}\label{example-416}

The example below defines a region name of UNIT and sets three FIPUNIT
regions for a 4 x 5 x 2 model.

\begin{verbatim}
--
--       DEFINE FIPUNIT FIP REGIONS FOR ALL CELLS
--
FIPUNIT
         2  2  1  1  2  2  1  1  1  1  1  1  1  1  1  1  1  1  1  1
         3  3  1  1  3  3  1  1  1  1  1  1  1  1  1  1  1  1  1  1
/
\end{verbatim}

The second example is based on the Norne model that defines two FIP
regions based on geological layers and numerical layers using the EQUALS
keyword.

\begin{verbatim}
--       FIPGL BASED ON GEOLOGICAL LAYERS
--       ARRAY       CONSTANT     ---------- BOX ---------
--                                I1  I2   J1  J2   K1  K2
EQUALS
         FIPGL       1             1  46    1 112    1   3   /  Garn
         FIPGL       2             1  46    1 112    4   4   /  Not
         FIPGL       3             1  46    1 112    5   5   /  Ile 2.2
         FIPGL       4             1  46    1 112    6   8   /  Ile 2.1
         FIPGL       5             1  46    1 112    9  11   /  Ile 1
         FIPGL       6             1  46    1 112   12  12   /  Tofte 2.2
         FIPGL       7             1  46    1 112   13  15   /  Tofte 2.1
         FIPGL       8             1  46    1 112   16  18   /  Tofte 1
         FIPGL       9             1  46    1 112   19  22   /  Tilje
--
--       FIPNL BASED ON NUMERICAL LAYERS
--
         FIPNL       1             1  46    1 112    1   1   /  Garn 3
         FIPNL       2             1  46    1 112    2   2   /  Garn 2
         FIPNL       3             1  46    1 112    3   3   /  Garn 1w
         FIPNL       4             1  46    1 112    4   4   /  Not
         FIPNL       5             1  46    1 112    5   5   /  Ile 2.2
         FIPNL       6             1  46    1 112    6   6   /  Ile 2.1.3
         FIPNL       7             1  46    1 112    7   7   /  Ile 2.1.2
         FIPNL       8             1  46    1 112    8   8   /  Ile 2.1.1
         FIPNL       9             1  46    1 112    9   9   /  Ile 1.3
         FIPNL      10             1  46    1 112   10  10   /  Ile 1.2
         FIPNL      11             1  46    1 112   11  11   /  Ile 1.1
         FIPNL      12             1  46    1 112   12  12   /  Tofte 2.2
         FIPNL      13             1  46    1 112   13  13   /  Tofte 2.1.3
         FIPNL      14             1  46    1 112   14  14   /  Tofte 2.1.2
         FIPNL      15             1  46    1 112   15  15   /  Tofte 2.1.1
         FIPNL      16             1  46    1 112   16  16   /  Tofte 1.2.2
         FIPNL      17             1  46    1 112   17  17   /  Tofte 1.2.1
         FIPNL      18             1  46    1 112   18  18   /  Tofte 1.1
         FIPNL      19             1  46    1 112   19  19   /  Tilje 4
         FIPNL      20             1  46    1 112   20  20   /  Tilje 3
         FIPNL      21             1  46    1 112   21  21   /  Tilje 2
         FIPNL      22             1  46    1 112   22  22   /  Tilje 1
/
\end{verbatim}

Then in order to get the reservoir pressure for the regions and the
in-place oil volumes SUMMARY variables written to the SUMMARY file, one
would use the following SUMMARY variable names in the SUMMARY section

\begin{verbatim}
-- ==============================================================================
--
-- SUMMARY SECTION
--
-- ==============================================================================
SUMMARY
--       FIP REGION REPORTING
--
ROIP_GL
/
ROIP_NL
/
RPR__GL
/
RPR__NL
/
\end{verbatim}

Notice how the ``\_'' character has been used to extend the SUMMARY
variable name to five characters.

\subsection{FIPNUM -- Define the Fluid In-Place Region
Numbers}\label{fipnum-define-the-fluid-in-place-region-numbers}

\begin{longtable}[]{@{}
  >{\raggedright\arraybackslash}p{(\linewidth - 14\tabcolsep) * \real{0.1250}}
  >{\raggedright\arraybackslash}p{(\linewidth - 14\tabcolsep) * \real{0.1250}}
  >{\raggedright\arraybackslash}p{(\linewidth - 14\tabcolsep) * \real{0.1250}}
  >{\raggedright\arraybackslash}p{(\linewidth - 14\tabcolsep) * \real{0.1250}}
  >{\raggedright\arraybackslash}p{(\linewidth - 14\tabcolsep) * \real{0.1250}}
  >{\raggedright\arraybackslash}p{(\linewidth - 14\tabcolsep) * \real{0.1250}}
  >{\raggedright\arraybackslash}p{(\linewidth - 14\tabcolsep) * \real{0.1250}}
  >{\raggedright\arraybackslash}p{(\linewidth - 14\tabcolsep) * \real{0.1250}}@{}}
\toprule\noalign{}
\begin{minipage}[b]{\linewidth}\raggedright
RUNSPEC
\end{minipage} & \begin{minipage}[b]{\linewidth}\raggedright
GRID
\end{minipage} & \begin{minipage}[b]{\linewidth}\raggedright
EDIT
\end{minipage} & \begin{minipage}[b]{\linewidth}\raggedright
PROPS
\end{minipage} & \begin{minipage}[b]{\linewidth}\raggedright
REGIONS
\end{minipage} & \begin{minipage}[b]{\linewidth}\raggedright
SOLUTION
\end{minipage} & \begin{minipage}[b]{\linewidth}\raggedright
SUMMARY
\end{minipage} & \begin{minipage}[b]{\linewidth}\raggedright
SCHEDULE
\end{minipage} \\
\midrule\noalign{}
\endhead
\bottomrule\noalign{}
\endlastfoot
\end{longtable}

\subsubsection{Description}\label{description-737}

The FIPNUM keyword defines the fluid in-place region numbers for each
grid block. The simulator can print out summaries of the fluid in-place
in each region, the current flow rates between regions, and the
cumulative flows between regions.

Note that the total number of FIPNUM and FIP regions must be defined by
the NMFIPR variable on the REGDIMS keyword, or the NTFIP variable on the
TABDIMS keyword. Both keywords are in the RUNSPEC section.

\begin{longtable}[]{@{}
  >{\raggedright\arraybackslash}p{(\linewidth - 6\tabcolsep) * \real{0.2500}}
  >{\raggedright\arraybackslash}p{(\linewidth - 6\tabcolsep) * \real{0.2500}}
  >{\raggedright\arraybackslash}p{(\linewidth - 6\tabcolsep) * \real{0.2500}}
  >{\raggedright\arraybackslash}p{(\linewidth - 6\tabcolsep) * \real{0.2500}}@{}}
\toprule\noalign{}
\begin{minipage}[b]{\linewidth}\raggedright
No.
\end{minipage} & \begin{minipage}[b]{\linewidth}\raggedright
Name
\end{minipage} & \begin{minipage}[b]{\linewidth}\raggedright
Description
\end{minipage} & \begin{minipage}[b]{\linewidth}\raggedright
Default
\end{minipage} \\
\midrule\noalign{}
\endhead
\bottomrule\noalign{}
\endlastfoot
1 & FIPNUM & FIPNUM defines an array of positive integers greater than
or equal to one, that assigns a grid cell to a particular fluid in-place
region. The maximum number of FIPNUM and FIP regions is set by the
REGDIMS(NMFIPR) or the TABDIMS(NTFIP) keywords(variables) in the RUNSPEC
section. If both REGDIMS(NMFIPR) and TABDIMS(NTFIP) have been defined
then the maximum of the two is used. & 1 \\
Notes: & & & \\
\end{longtable}

\emph{Table 9.5: FIPNUM Keyword Description}

\begin{tcolorbox}[enhanced jigsaw, title=\textcolor{quarto-callout-note-color}{\faInfo}\hspace{0.5em}{Note}, coltitle=black, breakable, toprule=.15mm, colbacktitle=quarto-callout-note-color!10!white, leftrule=.75mm, opacitybacktitle=0.6, colback=white, bottomtitle=1mm, toptitle=1mm, titlerule=0mm, left=2mm, colframe=quarto-callout-note-color-frame, arc=.35mm, rightrule=.15mm, bottomrule=.15mm, opacityback=0]

In most simulation models the FIPNUM array is used to define various
regions in the model for fluid in-place reporting and to identify (or
report) the flow between the different regions. When calibrating a
model's in-place volumes it would be useful to use the FIPNUM array
combined with the MULTREGP keyword to accomplish this. However, the
FIPNUM array cannot be used in the GRID section. A work around is to:
The above work flow will ensure that both arrays and the reporting of
fluid in-place regions are consistent.

\end{tcolorbox}

\subsubsection{Examples}\label{examples-50}

The example below sets three FIPNUM regions for a 4 x 5 x 2 model.

\begin{verbatim}
--
--       DEFINE FIPNUM REGIONS FOR ALL CELLS
--
FIPNUM
         2  2  1  1  2  2  1  1  1  1  1  1  1  1  1  1  1  1  1  1
         3  3  1  1  3  3  1  1  1  1  1  1  1  1  1  1  1  1  1  1
/
\end{verbatim}

Alternatively the EQUALS keyword could be employed to accomplish the
same task, that is:

\begin{verbatim}
--
--       ARRAY       CONSTANT     ---------- BOX ---------
--                                I1  I2   J1  J2   K1  K2
EQUALS
         FIPNUM     1            1*  1*   1*  1*    1*  1* / SET REGION 1
         FIPNUM     2            1   2    1   2     1   1  / SET REGION 2
         FIPNUM     3            1   2    1   2     2   2  / SET REGION 3
/

\end{verbatim}

\subsection{FIPOWG -- Activate Oil, Gas, and Water FIP Zone
Reporting}\label{fipowg-activate-oil-gas-and-water-fip-zone-reporting}

\begin{longtable}[]{@{}
  >{\raggedright\arraybackslash}p{(\linewidth - 14\tabcolsep) * \real{0.1250}}
  >{\raggedright\arraybackslash}p{(\linewidth - 14\tabcolsep) * \real{0.1250}}
  >{\raggedright\arraybackslash}p{(\linewidth - 14\tabcolsep) * \real{0.1250}}
  >{\raggedright\arraybackslash}p{(\linewidth - 14\tabcolsep) * \real{0.1250}}
  >{\raggedright\arraybackslash}p{(\linewidth - 14\tabcolsep) * \real{0.1250}}
  >{\raggedright\arraybackslash}p{(\linewidth - 14\tabcolsep) * \real{0.1250}}
  >{\raggedright\arraybackslash}p{(\linewidth - 14\tabcolsep) * \real{0.1250}}
  >{\raggedright\arraybackslash}p{(\linewidth - 14\tabcolsep) * \real{0.1250}}@{}}
\toprule\noalign{}
\begin{minipage}[b]{\linewidth}\raggedright
RUNSPEC
\end{minipage} & \begin{minipage}[b]{\linewidth}\raggedright
GRID
\end{minipage} & \begin{minipage}[b]{\linewidth}\raggedright
EDIT
\end{minipage} & \begin{minipage}[b]{\linewidth}\raggedright
PROPS
\end{minipage} & \begin{minipage}[b]{\linewidth}\raggedright
REGIONS
\end{minipage} & \begin{minipage}[b]{\linewidth}\raggedright
SOLUTION
\end{minipage} & \begin{minipage}[b]{\linewidth}\raggedright
SUMMARY
\end{minipage} & \begin{minipage}[b]{\linewidth}\raggedright
SCHEDULE
\end{minipage} \\
\midrule\noalign{}
\endhead
\bottomrule\noalign{}
\endlastfoot
\end{longtable}

\subsubsection{Description}\label{description-738}

The FIPOWG keyword activates automatic fluid in-place reporting based on
the initial oil, gas and water zones defined by the initial
equilibration. The fluid contacts on the EQUIL keyword in the SOLUTION
section determine the reporting fluid category a grid cell belongs to.
For example all grid cells with depths above the gas-oil contact on the
EQUIL keyword will be assigned to the gas zone and reported accordingly.
Similarly, grid cells with depths between the gas-oil contact and the
water-oil contact will be assigned to the oil zone. And finally, grid
cells with depths below the oil-water contact will be assigned to the
water zone. The simulator can print out summaries of the fluid in-place
in each region, the current flow rates between regions, and the
cumulative flows between regions.

Note that the total number of FIP and FIPNUM regions must be defined by
the NMFIPR variable on the REGDIMS keyword in the RUNSPEC section.

There is no data required for this keyword.

This keyword is not supported by OPM Flow but has no effect on the
results so it will be ignored.

\subsubsection{Example}\label{example-417}

\begin{verbatim}
--
--       ACTIVATE OIL, GAS, AND WATER FIP ZONE REPORTING
--
FIPOWG
\end{verbatim}

The above example switches on automatic fluid in-place reporting based
on the initial oil, gas and water zones defined by the initial
equilibration.

\subsection{FORMFEED -- Defined the Print File Form-Feed
Character}\label{formfeed-defined-the-print-file-form-feed-character-5}

The FORMFEED keyword defines the form-feed character, or carriage
control character, for the output print (\emph{.PRT) run summary (}.RSM)
files. The keyword should be place at the very top of the input file.

See FORMFEED -- Defined the Print File Form-Feed Character in the GLOBAL
section for a dull description.

\subsection{GETDATA -- Load and Assign Data Array from INIT or RESTART
File}\label{getdata-load-and-assign-data-array-from-init-or-restart-file-2}

The GETDATA keyword loads a data array from a previously generated INIT
or RESTART file and assigns the loaded array to either same array in the
run or another array name.

See GETDATA -- Load and Assign Data Array from INIT or RESTART Files in
the GRID section for a full description.

\subsection{HA -- History Match End-Point Gradient Additive
Modifier}\label{ha-history-match-end-point-gradient-additive-modifier-1}

The HA series of keywords defines the history match end-point gradient
parameters used to set the additive cumulative end point data, for when
the History Match Gradient option has been activated by the HMDIMS
keyword in the RUNSPEC section.

See HA -- History Match End-Point Gradient Additive Modifier in the
PROPS section for a complete description.

\subsection{HBNUM -- Define Herschel-Bulkley Region
Numbers}\label{hbnum-define-herschel-bulkley-region-numbers}

\begin{longtable}[]{@{}
  >{\raggedright\arraybackslash}p{(\linewidth - 14\tabcolsep) * \real{0.1250}}
  >{\raggedright\arraybackslash}p{(\linewidth - 14\tabcolsep) * \real{0.1250}}
  >{\raggedright\arraybackslash}p{(\linewidth - 14\tabcolsep) * \real{0.1250}}
  >{\raggedright\arraybackslash}p{(\linewidth - 14\tabcolsep) * \real{0.1250}}
  >{\raggedright\arraybackslash}p{(\linewidth - 14\tabcolsep) * \real{0.1250}}
  >{\raggedright\arraybackslash}p{(\linewidth - 14\tabcolsep) * \real{0.1250}}
  >{\raggedright\arraybackslash}p{(\linewidth - 14\tabcolsep) * \real{0.1250}}
  >{\raggedright\arraybackslash}p{(\linewidth - 14\tabcolsep) * \real{0.1250}}@{}}
\toprule\noalign{}
\begin{minipage}[b]{\linewidth}\raggedright
RUNSPEC
\end{minipage} & \begin{minipage}[b]{\linewidth}\raggedright
GRID
\end{minipage} & \begin{minipage}[b]{\linewidth}\raggedright
EDIT
\end{minipage} & \begin{minipage}[b]{\linewidth}\raggedright
PROPS
\end{minipage} & \begin{minipage}[b]{\linewidth}\raggedright
REGIONS
\end{minipage} & \begin{minipage}[b]{\linewidth}\raggedright
SOLUTION
\end{minipage} & \begin{minipage}[b]{\linewidth}\raggedright
SUMMARY
\end{minipage} & \begin{minipage}[b]{\linewidth}\raggedright
SCHEDULE
\end{minipage} \\
\midrule\noalign{}
\endhead
\bottomrule\noalign{}
\endlastfoot
\end{longtable}

\subsubsection{Description}\label{description-739}

The HBNUM keyword defines the Herschel-Bulkley rheological property
table region numbers for each grid block, as such there must be one
entry for each cell in the model. The region number specifies which
table of Herschel-Bulkley rheological property data is assigned to a
grid block, for when the Polymer option has been invoked via the POLYMER
keyword in the RUNSPEC section and the Non-Newtonian Fluid phase has
been declared active by the NNEWTF keyword, also in the RUNSPEC section.
The FHERCHBL keyword in the PROPS section is used to specify the
Herschel-Bulkley rheological property table data.

This keyword is not supported by OPM Flow but would change the results
if supported so the simulation will be stopped.

\subsection{HM -- History Match Region Gradient
Parameters}\label{hm-history-match-region-gradient-parameters}

\begin{longtable}[]{@{}
  >{\raggedright\arraybackslash}p{(\linewidth - 14\tabcolsep) * \real{0.1250}}
  >{\raggedright\arraybackslash}p{(\linewidth - 14\tabcolsep) * \real{0.1250}}
  >{\raggedright\arraybackslash}p{(\linewidth - 14\tabcolsep) * \real{0.1250}}
  >{\raggedright\arraybackslash}p{(\linewidth - 14\tabcolsep) * \real{0.1250}}
  >{\raggedright\arraybackslash}p{(\linewidth - 14\tabcolsep) * \real{0.1250}}
  >{\raggedright\arraybackslash}p{(\linewidth - 14\tabcolsep) * \real{0.1250}}
  >{\raggedright\arraybackslash}p{(\linewidth - 14\tabcolsep) * \real{0.1250}}
  >{\raggedright\arraybackslash}p{(\linewidth - 14\tabcolsep) * \real{0.1250}}@{}}
\toprule\noalign{}
\begin{minipage}[b]{\linewidth}\raggedright
RUNSPEC
\end{minipage} & \begin{minipage}[b]{\linewidth}\raggedright
GRID
\end{minipage} & \begin{minipage}[b]{\linewidth}\raggedright
EDIT
\end{minipage} & \begin{minipage}[b]{\linewidth}\raggedright
PROPS
\end{minipage} & \begin{minipage}[b]{\linewidth}\raggedright
REGIONS
\end{minipage} & \begin{minipage}[b]{\linewidth}\raggedright
SOLUTION
\end{minipage} & \begin{minipage}[b]{\linewidth}\raggedright
SUMMARY
\end{minipage} & \begin{minipage}[b]{\linewidth}\raggedright
SCHEDULE
\end{minipage} \\
\midrule\noalign{}
\endhead
\bottomrule\noalign{}
\endlastfoot
\end{longtable}

\subsubsection{Description}\label{description-740}

The HM series of keywords in the REGION section defines the history
match gradient regions and sub-regions, for when the History Match
Gradient option has been activated by the HMDIMS keyword in the RUNSPEC
section.

This keyword is not supported by OPM Flow but has no effect on the
results so it will be ignored.

For grid properties, the region name (or region property array) is based
on the property arrays defined in Table 9.6.

\begin{longtable}[]{@{}
  >{\raggedright\arraybackslash}p{(\linewidth - 4\tabcolsep) * \real{0.3333}}
  >{\raggedright\arraybackslash}p{(\linewidth - 4\tabcolsep) * \real{0.3333}}
  >{\raggedright\arraybackslash}p{(\linewidth - 4\tabcolsep) * \real{0.3333}}@{}}
\toprule\noalign{}
\begin{minipage}[b]{\linewidth}\raggedright
Property Array
\end{minipage} & \begin{minipage}[b]{\linewidth}\raggedright
Region Name
\end{minipage} & \begin{minipage}[b]{\linewidth}\raggedright
Grid Property Data Desription
\end{minipage} \\
\midrule\noalign{}
\endhead
\bottomrule\noalign{}
\endlastfoot
PERMX & HMPERMX & Permeability multipliers in the x-direction. \\
PERMXY & HMPRMXY & Permeability multipliers in the x-direction and
y-direction \\
pERMY & HMPERMY & Permeability multipliers in the y-direction. \\
PERMZ & HMPERMZ & Permeability multipliers in the z-direction. \\
PORV & HMPORV & Pore volume multiplier \\
SIGMA & HMSIGMA & Dual porosity and/or dual permeability SIGMA
multiplier \\
TRANX & HMTRANX & Transmissibility multipliers in the x-direction. \\
TRANXY & HMTRNXY & Transmissibility multipliers in the x-direction and
y-direction \\
TRANY & HMTRANY & Transmissibility multipliers in the y-direction. \\
TRANZ & HMTRANZ & Transmissibility multipliers in the z-direction. \\
\end{longtable}

\emph{Table 9.6: HM Region Grid Gradient Parameter Keyword List}

In addition, if the End-Point Scaling option has been activated by the
ENDSCALE keyword in the RUNSPEC section, then the history match gradient
regions and sub-regions for the end-point data can be specified. In this
the keyword consists of the first two characters of ``HM'' followed by
the end-point keyword (Table 9.7), for example, HMSWL.

\begin{longtable}[]{@{}
  >{\raggedright\arraybackslash}p{(\linewidth - 4\tabcolsep) * \real{0.3333}}
  >{\raggedright\arraybackslash}p{(\linewidth - 4\tabcolsep) * \real{0.3333}}
  >{\raggedright\arraybackslash}p{(\linewidth - 4\tabcolsep) * \real{0.3333}}@{}}
\toprule\noalign{}
\begin{minipage}[b]{\linewidth}\raggedright
Type
\end{minipage} & \begin{minipage}[b]{\linewidth}\raggedright
End-Point Keyword
\end{minipage} & \begin{minipage}[b]{\linewidth}\raggedright
Oil-Water End-Point Definitions
\end{minipage} \\
\midrule\noalign{}
\endhead
\bottomrule\noalign{}
\endlastfoot
Saturation & HMSWL & Connate water saturation, that is the smallest
water saturation in a water saturation function table. \\
HMSWCR & Critical water saturation, that is the largest water saturation
for which the water relative permeability is zero. & \\
HMSOWCR & Critical oil-in-water saturation, that is the largest oil
saturation for which the oil relative permeability is zero in an
oil-water system. & \\
Relative Permeability & HMKRW & Relative permeability of water at the
maximum water saturation (normally the maximum water saturation is
one). \\
HMKRO & Relative permeability of oil at the maximum oil saturation. & \\
HMKRWR & Relative permeability of water at the residual oil saturation
or the residual gas saturation in a gas-water run. & \\
HMKRORW & Relative permeability of oil at the critical water saturation.
& \\
Capillary Pressure & HMSWLPC & Capillary pressure connate water
saturation, that is the smallest water saturation in a water saturation
function table. \\
Type & End-Point Keyword & Gas-Oil End-Point Definitions \\
Saturation & HMSGL & Connate gas saturation, that is the smallest gas
saturation in a gas saturation function table. \\
HMSGCR & Critical gas saturation, that is the largest gas saturation for
which the gas relative permeability is zero. & \\
HMSOGCR & Critical oil-in-gas saturation, that is the largest oil
saturation for which the oil relative permeability is zero in an
oil-gas-connate water system. & \\
Relative Permeability & HMKRG & Relative permeability of gas at the
maximum gas saturation. \\
HMKRGR & Relative permeability of gas at the residual oil saturation or
the critical water saturation in a gas-water run. & \\
HMKRORG & Relative permeability of oil at the critical gas saturation.
& \\
Capillary Pressure & HMSGLPC & Capillary pressure connate gas
saturation, that is the smallest gas saturation in a gas saturation
function table. \\
\end{longtable}

\emph{Table 9.7: HM Region End-Point Gradient Parameter Keyword List}

\subsection{HMPROPS -- History Match End-Point Section
Start}\label{hmprops-history-match-end-point-section-start-1}

HMPROPS defines the start of a history match end-points section, for
when the History Match Gradient option has been activated by the HMDIMS
keyword in the RUNSPEC section. In addition, the End-Point Scaling
option must also be activated by the ENDSCALE keyword which is also in
the RUNSPEC section. The keyword allows for the BOX, EQUALS, COPY,
MINVALUE, MAXVALUE and ADD keywords to be used with the HM series of
keywords that reference the end-point scaling arrays, that is: HMKRG,
HMKRGR, HMKRO, HMKRORG, HMKRORW, HMKRW, HMKRWR, HMPCW, HMPCG, HMSGCR,
HMSOWCR, HMSOGCR, HMSWCR, and HMSWL.

See HMPROPS -- History Match End-Point Section Start in the PROPS
section for a full description.

\subsection{HWSNUM -- Define the Saturation Table Region Numbers (High
Salinity and Water
Wet)}\label{hwsnum-define-the-saturation-table-region-numbers-high-salinity-and-water-wet}

\begin{longtable}[]{@{}
  >{\raggedright\arraybackslash}p{(\linewidth - 14\tabcolsep) * \real{0.1250}}
  >{\raggedright\arraybackslash}p{(\linewidth - 14\tabcolsep) * \real{0.1250}}
  >{\raggedright\arraybackslash}p{(\linewidth - 14\tabcolsep) * \real{0.1250}}
  >{\raggedright\arraybackslash}p{(\linewidth - 14\tabcolsep) * \real{0.1250}}
  >{\raggedright\arraybackslash}p{(\linewidth - 14\tabcolsep) * \real{0.1250}}
  >{\raggedright\arraybackslash}p{(\linewidth - 14\tabcolsep) * \real{0.1250}}
  >{\raggedright\arraybackslash}p{(\linewidth - 14\tabcolsep) * \real{0.1250}}
  >{\raggedright\arraybackslash}p{(\linewidth - 14\tabcolsep) * \real{0.1250}}@{}}
\toprule\noalign{}
\begin{minipage}[b]{\linewidth}\raggedright
RUNSPEC
\end{minipage} & \begin{minipage}[b]{\linewidth}\raggedright
GRID
\end{minipage} & \begin{minipage}[b]{\linewidth}\raggedright
EDIT
\end{minipage} & \begin{minipage}[b]{\linewidth}\raggedright
PROPS
\end{minipage} & \begin{minipage}[b]{\linewidth}\raggedright
REGIONS
\end{minipage} & \begin{minipage}[b]{\linewidth}\raggedright
SOLUTION
\end{minipage} & \begin{minipage}[b]{\linewidth}\raggedright
SUMMARY
\end{minipage} & \begin{minipage}[b]{\linewidth}\raggedright
SCHEDULE
\end{minipage} \\
\midrule\noalign{}
\endhead
\bottomrule\noalign{}
\endlastfoot
\end{longtable}

\subsubsection{Description}\label{description-741}

The HWSNUM keyword defines the saturation tables (relative permeability
and capillary pressure tables) region numbers for each grid block, as
such there must be one entry for each cell in the mode, for when the
Surfactant Wettability option has been selected. The keyword may also be
used with the Low Salt Brine option, in this case the water wet curves
are calculated as a function of the low and high water wet salinity
curves. The region number specifies which set of relative permeability
tables are used to calculate the relative permeability and capillary
pressure in a grid block. Note that the keyword is obligatory if the
SURFACTW keyword in the RUNSPEC section has been used to invoke the
Surfactant Wettability option.

This keyword is not supported by OPM Flow but would change the results
if supported so the simulation will be stopped.

\begin{longtable}[]{@{}
  >{\raggedright\arraybackslash}p{(\linewidth - 6\tabcolsep) * \real{0.2500}}
  >{\raggedright\arraybackslash}p{(\linewidth - 6\tabcolsep) * \real{0.2500}}
  >{\raggedright\arraybackslash}p{(\linewidth - 6\tabcolsep) * \real{0.2500}}
  >{\raggedright\arraybackslash}p{(\linewidth - 6\tabcolsep) * \real{0.2500}}@{}}
\toprule\noalign{}
\begin{minipage}[b]{\linewidth}\raggedright
No.
\end{minipage} & \begin{minipage}[b]{\linewidth}\raggedright
Name
\end{minipage} & \begin{minipage}[b]{\linewidth}\raggedright
Description
\end{minipage} & \begin{minipage}[b]{\linewidth}\raggedright
Default
\end{minipage} \\
\midrule\noalign{}
\endhead
\bottomrule\noalign{}
\endlastfoot
1 & HWSNUM & HWSNUM defines an array of positive integers assigning a
grid cell to a particular saturation table region. The maximum number of
HWSNUM regions is set by the NTSFUN variable on the TABDIMS keyword in
the RUNSPEC section. & 1 \\
Notes: & & & \\
\end{longtable}

\emph{Table 9.8: HWSNUM Keyword Description}

\subsubsection{Examples}\label{examples-51}

The example below sets three HWSNUM regions for a 4 x 5 x 2 model.

\begin{verbatim}
--
--       DEFINE HWSNUM REGIONS FOR ALL CELLS
--
HWSNUM
         2  2  1  1  2  2  1  1  1  1  1  1  1  1  1  1  1  1  1  1
         3  3  1  1  3  3  1  1  1  1  1  1  1  1  1  1  1  1  1  1
/
\end{verbatim}

Alternatively the EQUALS keyword could be employed to accomplish the
same task, that is:

\begin{verbatim}
--
--       ARRAY       CONSTANT     ---------- BOX ---------
--                                I1  I2   J1  J2   K1  K2
EQUALS
         HWSNUM      1            1*  1*   1*  1*   1*  1* / SET REGION 1
         HWSNUM      2            1   2    1   2    1   1  / SET REGION 2
         HWSNUM      3            1   2    1   2    2   2  / SET REGION 3
/
\end{verbatim}

\subsection{IMBNUM -- Define the Imbibition Saturation Table Region
Numbers}\label{imbnum-define-the-imbibition-saturation-table-region-numbers}

\begin{longtable}[]{@{}
  >{\raggedright\arraybackslash}p{(\linewidth - 14\tabcolsep) * \real{0.1250}}
  >{\raggedright\arraybackslash}p{(\linewidth - 14\tabcolsep) * \real{0.1250}}
  >{\raggedright\arraybackslash}p{(\linewidth - 14\tabcolsep) * \real{0.1250}}
  >{\raggedright\arraybackslash}p{(\linewidth - 14\tabcolsep) * \real{0.1250}}
  >{\raggedright\arraybackslash}p{(\linewidth - 14\tabcolsep) * \real{0.1250}}
  >{\raggedright\arraybackslash}p{(\linewidth - 14\tabcolsep) * \real{0.1250}}
  >{\raggedright\arraybackslash}p{(\linewidth - 14\tabcolsep) * \real{0.1250}}
  >{\raggedright\arraybackslash}p{(\linewidth - 14\tabcolsep) * \real{0.1250}}@{}}
\toprule\noalign{}
\begin{minipage}[b]{\linewidth}\raggedright
RUNSPEC
\end{minipage} & \begin{minipage}[b]{\linewidth}\raggedright
GRID
\end{minipage} & \begin{minipage}[b]{\linewidth}\raggedright
EDIT
\end{minipage} & \begin{minipage}[b]{\linewidth}\raggedright
PROPS
\end{minipage} & \begin{minipage}[b]{\linewidth}\raggedright
REGIONS
\end{minipage} & \begin{minipage}[b]{\linewidth}\raggedright
SOLUTION
\end{minipage} & \begin{minipage}[b]{\linewidth}\raggedright
SUMMARY
\end{minipage} & \begin{minipage}[b]{\linewidth}\raggedright
SCHEDULE
\end{minipage} \\
\midrule\noalign{}
\endhead
\bottomrule\noalign{}
\endlastfoot
\end{longtable}

\subsubsection{Description}\label{description-742}

The IMBNUM keyword defines the imbibition saturation tables (relative
permeability and capillary pressure tables) region numbers for each grid
block, as such there must be one entry for each cell in the model. The
region number specifies which set of relative permeability tables (SGFN,
SWFN, SOF2, SOF3, SOF32D, SGOF, SLGOF and SWOF) are used to calculate
the relative permeability and capillary pressure in a grid block.

\begin{longtable}[]{@{}
  >{\raggedright\arraybackslash}p{(\linewidth - 6\tabcolsep) * \real{0.2500}}
  >{\raggedright\arraybackslash}p{(\linewidth - 6\tabcolsep) * \real{0.2500}}
  >{\raggedright\arraybackslash}p{(\linewidth - 6\tabcolsep) * \real{0.2500}}
  >{\raggedright\arraybackslash}p{(\linewidth - 6\tabcolsep) * \real{0.2500}}@{}}
\toprule\noalign{}
\begin{minipage}[b]{\linewidth}\raggedright
No.
\end{minipage} & \begin{minipage}[b]{\linewidth}\raggedright
Name
\end{minipage} & \begin{minipage}[b]{\linewidth}\raggedright
Description
\end{minipage} & \begin{minipage}[b]{\linewidth}\raggedright
Default
\end{minipage} \\
\midrule\noalign{}
\endhead
\bottomrule\noalign{}
\endlastfoot
1 & IMBNUM & IMBNUM defines an array of positive integers assigning a
grid cell to a particular saturation table region. The maximum number of
IMBNUM regions is set by the NTSFUN variable on the TABDIMS keyword in
the RUNSPEC section. & 1 \\
Notes: & & & \\
\end{longtable}

\emph{Table 9.9: IMBNUM Keyword Description}

In addition, saturation table assignment may be directional dependent in
which case the directional dependent versions of the aforementioned
array should be used, that is IMBNUMX, IMBNUMY and IMBNUMZ instead of
IMBNUM. There is also the facility to make the directional end-point
scaling reversible or non-reversible and if the non-reversible option is
selected, the non-reversible versions of the aforementioned arrays
should be used, that is IMBNUMX, IMBNUMX-, IMBNUMY, IMBNUMY-, IMBNUMZ
and IMBNUMZ-, instead of the IMBNUM keyword. For reference, see Table
5.40, that lists the various keywords that may be used with directional
dependent relative permeability tables.

Note, currently IMBNUMX-, IMBNUMY-, and IMBNUMZ- are not supported by
OPM Flow.

\subsubsection{Example}\label{example-418}

The example below sets three IMBNUM regions for a 4 x 5 x 2 model using
the EQUALS keyword.

\begin{verbatim}
--
--       ARRAY       CONSTANT     ---------- BOX ---------
--                                I1  I2   J1  J2   K1  K2
EQUALS
         IMBNUM      1            1*  1*   1*  1*   1*  1* / SET REGION 1
         IMBNUM      2            1   2    1   2    1   1  / SET REGION 2
         IMBNUM      3            1   2    1   2    2   2  / SET REGION 3
/
\end{verbatim}

\subsection{IMBNUMMF -- Define the Imbibition Saturation Table Region
Numbers
(Matrix-Fracture)}\label{imbnummf-define-the-imbibition-saturation-table-region-numbers-matrix-fracture}

\begin{longtable}[]{@{}
  >{\raggedright\arraybackslash}p{(\linewidth - 14\tabcolsep) * \real{0.1250}}
  >{\raggedright\arraybackslash}p{(\linewidth - 14\tabcolsep) * \real{0.1250}}
  >{\raggedright\arraybackslash}p{(\linewidth - 14\tabcolsep) * \real{0.1250}}
  >{\raggedright\arraybackslash}p{(\linewidth - 14\tabcolsep) * \real{0.1250}}
  >{\raggedright\arraybackslash}p{(\linewidth - 14\tabcolsep) * \real{0.1250}}
  >{\raggedright\arraybackslash}p{(\linewidth - 14\tabcolsep) * \real{0.1250}}
  >{\raggedright\arraybackslash}p{(\linewidth - 14\tabcolsep) * \real{0.1250}}
  >{\raggedright\arraybackslash}p{(\linewidth - 14\tabcolsep) * \real{0.1250}}@{}}
\toprule\noalign{}
\begin{minipage}[b]{\linewidth}\raggedright
RUNSPEC
\end{minipage} & \begin{minipage}[b]{\linewidth}\raggedright
GRID
\end{minipage} & \begin{minipage}[b]{\linewidth}\raggedright
EDIT
\end{minipage} & \begin{minipage}[b]{\linewidth}\raggedright
PROPS
\end{minipage} & \begin{minipage}[b]{\linewidth}\raggedright
REGIONS
\end{minipage} & \begin{minipage}[b]{\linewidth}\raggedright
SOLUTION
\end{minipage} & \begin{minipage}[b]{\linewidth}\raggedright
SUMMARY
\end{minipage} & \begin{minipage}[b]{\linewidth}\raggedright
SCHEDULE
\end{minipage} \\
\midrule\noalign{}
\endhead
\bottomrule\noalign{}
\endlastfoot
\end{longtable}

\subsubsection{Description}\label{description-743}

The IMBNUMMF keyword defines the imbibition saturation tables (relative
permeability and capillary pressure tables) region numbers for flow
between the matrix and fracture blocks, for when the HYSTER option on
the SATOPTS keyword has been invoked to activate the Hysteresis option,
and the Dual Porosity or Dual Permeability models have been activated
via the DUALPORO or DUALPERM keywords. All keywords are in the RUNSPEC
section.

The region number specifies which set of relative permeability tables
(SGFN, SWFN, SOF2, SOF3, SOF32D, SGOF, SLGOF and SWOF) are used to
calculate the relative permeability and capillary pressure between the
matrix and fracture blocks. The keyword is optional and any cell not
assigned a value will use the assignment from the IMBNUM array.

This keyword is not supported by OPM Flow but would change the results
if supported so the simulation will be stopped.

\subsection{IMPORT -- Import Grid File Data at the Current
Position}\label{import-import-grid-file-data-at-the-current-position-3}

The IMPORT keyword informs the simulator to continue reading input data
from the specified IMPORT file. When the end of the IMPORT file is
reached, input data is read from the next keyword in the current file.
Normally IMPORT files are generated by grid pre-processing software and
the keyword allows for both formatted and unformatted (binary) files to
be loaded.

See IMPORT -- Import Grid File Data at the Current Position in the GRID
section for a full description.

\subsection{INCLUDE -- Load Another Data File at the Current
Position}\label{include-load-another-data-file-at-the-current-position-5}

The INCLUDE keyword informs OPM Flow to continue reading input data from
the specified INCLUDE file. When the end of the INCLUDE file is reached,
or the ENDINC keyword is encountered in the included file, input data is
read from the next keyword in the current file.

See INCLUDE -- Load Another Data File at the Current Position in the
GLOBAL section for a full description.

\subsection{KRNUM -- Define the Directional Saturation Table Region
Numbers}\label{krnum-define-the-directional-saturation-table-region-numbers}

\begin{longtable}[]{@{}
  >{\raggedright\arraybackslash}p{(\linewidth - 14\tabcolsep) * \real{0.1250}}
  >{\raggedright\arraybackslash}p{(\linewidth - 14\tabcolsep) * \real{0.1250}}
  >{\raggedright\arraybackslash}p{(\linewidth - 14\tabcolsep) * \real{0.1250}}
  >{\raggedright\arraybackslash}p{(\linewidth - 14\tabcolsep) * \real{0.1250}}
  >{\raggedright\arraybackslash}p{(\linewidth - 14\tabcolsep) * \real{0.1250}}
  >{\raggedright\arraybackslash}p{(\linewidth - 14\tabcolsep) * \real{0.1250}}
  >{\raggedright\arraybackslash}p{(\linewidth - 14\tabcolsep) * \real{0.1250}}
  >{\raggedright\arraybackslash}p{(\linewidth - 14\tabcolsep) * \real{0.1250}}@{}}
\toprule\noalign{}
\begin{minipage}[b]{\linewidth}\raggedright
RUNSPEC
\end{minipage} & \begin{minipage}[b]{\linewidth}\raggedright
GRID
\end{minipage} & \begin{minipage}[b]{\linewidth}\raggedright
EDIT
\end{minipage} & \begin{minipage}[b]{\linewidth}\raggedright
PROPS
\end{minipage} & \begin{minipage}[b]{\linewidth}\raggedright
REGIONS
\end{minipage} & \begin{minipage}[b]{\linewidth}\raggedright
SOLUTION
\end{minipage} & \begin{minipage}[b]{\linewidth}\raggedright
SUMMARY
\end{minipage} & \begin{minipage}[b]{\linewidth}\raggedright
SCHEDULE
\end{minipage} \\
\midrule\noalign{}
\endhead
\bottomrule\noalign{}
\endlastfoot
\end{longtable}

\subsubsection{Description}\label{description-744}

The KRNUM keyword defines the direction dependent saturation tables
(relative permeability and capillary pressure tables) region numbers for
each grid block face, as such there must be one entry for each cell in
the model. The region number specifies which set of relative
permeability tables (SGFN, SWFN, SOF2, SOF3, SOF32D, SGOF, SLGOF and
SWOF) are used to calculate the relative permeability and capillary
pressure in a grid block. The keyword should only be used if Directional
Dependent Saturation Function option has been activated by the DIRECT
parameter on the SATOPTS keyword in the RUNSPEC section. Otherwise the
standard none directional relative permeability curves should be
assigned by the SATNUM keyword in the REGIONS section.

This keyword is not in the standard keyword format due to the cell face
(X, X-, Y, Y-, Z, and Z- for Cartesian grids and R, R-, T, T-, Z, and Z-
for radial grids) being concatenated to the end of the keyword KRNUM to
fully define the keyword.

\begin{longtable}[]{@{}
  >{\raggedright\arraybackslash}p{(\linewidth - 6\tabcolsep) * \real{0.2500}}
  >{\raggedright\arraybackslash}p{(\linewidth - 6\tabcolsep) * \real{0.2500}}
  >{\raggedright\arraybackslash}p{(\linewidth - 6\tabcolsep) * \real{0.2500}}
  >{\raggedright\arraybackslash}p{(\linewidth - 6\tabcolsep) * \real{0.2500}}@{}}
\toprule\noalign{}
\begin{minipage}[b]{\linewidth}\raggedright
No.
\end{minipage} & \begin{minipage}[b]{\linewidth}\raggedright
Name
\end{minipage} & \begin{minipage}[b]{\linewidth}\raggedright
Description
\end{minipage} & \begin{minipage}[b]{\linewidth}\raggedright
Default
\end{minipage} \\
\midrule\noalign{}
\endhead
\bottomrule\noalign{}
\endlastfoot
1 & KRNUM & KRNUM defines an array of positive integers assigning a grid
cell to a particular saturation table region. The maximum number of
KRNUM regions is set by the NTSFUN variable on the TABDIMS keyword in
the RUNSPEC section. & 1 \\
Notes: & & & \\
\end{longtable}

\emph{Table 9.10: KRNUM Keyword Description}

If the Directional Dependent Saturation Function option has been
activated by the DIRECT parameter on the SATOPTS keyword in the RUNSPEC
section, then the KRNUMX, KRNUMY and KRNUMZ versions of the keyword
should be used for cartesian grids. Secondly, if the Non-Reversible
End-Point Scaling option has also been activated by the IRREVERS
parameter on the SATOPTS keyword in the RUNSPEC section, then the
non-reversible versions of the KRNUM should be used, that is KRNUMX,
KRNUMX-, KRNUMY, KRNUMY-, KRNUMZ and KRNUMZ-. For reference, see Table
5.40, that lists the various keywords that may be used with directional
dependent relative permeability tables.

The keywords KRNUMX-, KRNUMY-, and KRNUMZ- are not supported by OPM Flow
but would change the results if supported so the simulation will be
stopped.

\subsubsection{Example}\label{example-419}

The example below sets the directional saturation tables in all three
directions using the EQUALS keyword.

\begin{verbatim}
--
--       ARRAY       CONSTANT     ---------- BOX ---------
--                                I1  I2   J1  J2   K1  K2
EQUALS
         KRNUMX      1            1*  1*   1*  1*   1*  1* / SET X-DIR TABLES
         KRNUMY      2            1*  1*   1*  1*   1*  1* / SET Y-DIR TABLES
         KRNUMZ      3            1*  1*   1*  1*   1*  1* / SET Z-DIR TABLES
/
\end{verbatim}

\subsection{KRNUMMF -- Define the Saturation Table Region Numbers
(Matrix-Fracture)}\label{krnummf-define-the-saturation-table-region-numbers-matrix-fracture}

\begin{longtable}[]{@{}
  >{\raggedright\arraybackslash}p{(\linewidth - 14\tabcolsep) * \real{0.1250}}
  >{\raggedright\arraybackslash}p{(\linewidth - 14\tabcolsep) * \real{0.1250}}
  >{\raggedright\arraybackslash}p{(\linewidth - 14\tabcolsep) * \real{0.1250}}
  >{\raggedright\arraybackslash}p{(\linewidth - 14\tabcolsep) * \real{0.1250}}
  >{\raggedright\arraybackslash}p{(\linewidth - 14\tabcolsep) * \real{0.1250}}
  >{\raggedright\arraybackslash}p{(\linewidth - 14\tabcolsep) * \real{0.1250}}
  >{\raggedright\arraybackslash}p{(\linewidth - 14\tabcolsep) * \real{0.1250}}
  >{\raggedright\arraybackslash}p{(\linewidth - 14\tabcolsep) * \real{0.1250}}@{}}
\toprule\noalign{}
\begin{minipage}[b]{\linewidth}\raggedright
RUNSPEC
\end{minipage} & \begin{minipage}[b]{\linewidth}\raggedright
GRID
\end{minipage} & \begin{minipage}[b]{\linewidth}\raggedright
EDIT
\end{minipage} & \begin{minipage}[b]{\linewidth}\raggedright
PROPS
\end{minipage} & \begin{minipage}[b]{\linewidth}\raggedright
REGIONS
\end{minipage} & \begin{minipage}[b]{\linewidth}\raggedright
SOLUTION
\end{minipage} & \begin{minipage}[b]{\linewidth}\raggedright
SUMMARY
\end{minipage} & \begin{minipage}[b]{\linewidth}\raggedright
SCHEDULE
\end{minipage} \\
\midrule\noalign{}
\endhead
\bottomrule\noalign{}
\endlastfoot
\end{longtable}

\subsubsection{Description}\label{description-745}

The KRNUMMF keyword defines the drainage saturation tables (relative
permeability and capillary pressure tables) region numbers for flow
between the matrix and fracture blocks, for when the Dual Porosity or
Dual Permeability models have been activated via the DUALPORO or
DUALPERM keywords in the RUNSPEC section.

The region number specifies which set of relative permeability tables
(SGFN, SWFN, SOF2, SOF3, SOF32D, SGOF, SLGOF and SWOF) are used to
calculate the relative permeability and capillary pressure between the
matrix and fracture blocks. The keyword is optional and any cell not
assigned a value will use the assignment from the SATNUM array.

This keyword is not supported by OPM Flow but would change the results
if supported so the simulation will be stopped.

\subsection{LSLTWNUM -- Define the Low Salt Water Wet Saturation Table
Region
Numbers}\label{lsltwnum-define-the-low-salt-water-wet-saturation-table-region-numbers}

\begin{longtable}[]{@{}
  >{\raggedright\arraybackslash}p{(\linewidth - 14\tabcolsep) * \real{0.1250}}
  >{\raggedright\arraybackslash}p{(\linewidth - 14\tabcolsep) * \real{0.1250}}
  >{\raggedright\arraybackslash}p{(\linewidth - 14\tabcolsep) * \real{0.1250}}
  >{\raggedright\arraybackslash}p{(\linewidth - 14\tabcolsep) * \real{0.1250}}
  >{\raggedright\arraybackslash}p{(\linewidth - 14\tabcolsep) * \real{0.1250}}
  >{\raggedright\arraybackslash}p{(\linewidth - 14\tabcolsep) * \real{0.1250}}
  >{\raggedright\arraybackslash}p{(\linewidth - 14\tabcolsep) * \real{0.1250}}
  >{\raggedright\arraybackslash}p{(\linewidth - 14\tabcolsep) * \real{0.1250}}@{}}
\toprule\noalign{}
\begin{minipage}[b]{\linewidth}\raggedright
RUNSPEC
\end{minipage} & \begin{minipage}[b]{\linewidth}\raggedright
GRID
\end{minipage} & \begin{minipage}[b]{\linewidth}\raggedright
EDIT
\end{minipage} & \begin{minipage}[b]{\linewidth}\raggedright
PROPS
\end{minipage} & \begin{minipage}[b]{\linewidth}\raggedright
REGIONS
\end{minipage} & \begin{minipage}[b]{\linewidth}\raggedright
SOLUTION
\end{minipage} & \begin{minipage}[b]{\linewidth}\raggedright
SUMMARY
\end{minipage} & \begin{minipage}[b]{\linewidth}\raggedright
SCHEDULE
\end{minipage} \\
\midrule\noalign{}
\endhead
\bottomrule\noalign{}
\endlastfoot
\end{longtable}

\subsubsection{Description}\label{description-746}

The LSLTWNUM keyword defines the saturation tables (relative
permeability and capillary pressure tables) region numbers for each grid
block, as such there must be one entry for each cell in the model. The
region number specifies which set of relative permeability tables (SWFN,
SOF3 and related keywords) are used to calculate the relative
permeability and capillary pressure in a grid block. The keyword should
only be used if the Low Salinity option for the Brine model and the
Surfactant Wettability option have been activated by the LOWSALT and
SURFACTW keywords, respectively, in the RUNSPEC section.

The water wet curves are calculated as a weighted average of the low
salinity saturation tables (allocated by this keyword) and the high
salinity water wet saturation tables (allocated by the SURFWNUM keyword
in the REGIONS section), using the weights provided by the LSALTFNC
keyword in the PROPS section.

This keyword is not supported by OPM Flow but would change the results
if supported so the simulation will be stopped.

\begin{longtable}[]{@{}
  >{\raggedright\arraybackslash}p{(\linewidth - 6\tabcolsep) * \real{0.2500}}
  >{\raggedright\arraybackslash}p{(\linewidth - 6\tabcolsep) * \real{0.2500}}
  >{\raggedright\arraybackslash}p{(\linewidth - 6\tabcolsep) * \real{0.2500}}
  >{\raggedright\arraybackslash}p{(\linewidth - 6\tabcolsep) * \real{0.2500}}@{}}
\toprule\noalign{}
\begin{minipage}[b]{\linewidth}\raggedright
No.
\end{minipage} & \begin{minipage}[b]{\linewidth}\raggedright
Name
\end{minipage} & \begin{minipage}[b]{\linewidth}\raggedright
Description
\end{minipage} & \begin{minipage}[b]{\linewidth}\raggedright
Default
\end{minipage} \\
\midrule\noalign{}
\endhead
\bottomrule\noalign{}
\endlastfoot
1 & LSLTWNUM & LSLTWNUM defines an array of positive integers assigning
a grid cell to a particular saturation table region. The maximum number
of LSLTWNUM regions is set by the NTSFUN variable on the TABDIMS keyword
in the RUNSPEC section. & 1 \\
Notes: & & & \\
\end{longtable}

\emph{Table 9.11: LSLTWNUM Keyword Description}

\subsubsection{Example}\label{example-420}

The example below sets three LSLTWNUM regions for the model.

\begin{verbatim}
--
--       ARRAY       CONSTANT     ---------- BOX ---------
--                                I1  I2   J1  J2   K1  K2
EQUALS
         LSLTWNUM    1            1*  1*   1*  1*   1*  1* / SET REGION 1
         LSLTWNUM    2            1   2    1   2    1   1  / SET REGION 2
         LSLTWNUM    3            1   2    1   2    2   2  / SET REGION 3
/
\end{verbatim}

\subsection{LSNUM -- Define the Low Salt Oil Wet Saturation Table Region
Numbers}\label{lsnum-define-the-low-salt-oil-wet-saturation-table-region-numbers}

\begin{longtable}[]{@{}
  >{\raggedright\arraybackslash}p{(\linewidth - 14\tabcolsep) * \real{0.1250}}
  >{\raggedright\arraybackslash}p{(\linewidth - 14\tabcolsep) * \real{0.1250}}
  >{\raggedright\arraybackslash}p{(\linewidth - 14\tabcolsep) * \real{0.1250}}
  >{\raggedright\arraybackslash}p{(\linewidth - 14\tabcolsep) * \real{0.1250}}
  >{\raggedright\arraybackslash}p{(\linewidth - 14\tabcolsep) * \real{0.1250}}
  >{\raggedright\arraybackslash}p{(\linewidth - 14\tabcolsep) * \real{0.1250}}
  >{\raggedright\arraybackslash}p{(\linewidth - 14\tabcolsep) * \real{0.1250}}
  >{\raggedright\arraybackslash}p{(\linewidth - 14\tabcolsep) * \real{0.1250}}@{}}
\toprule\noalign{}
\begin{minipage}[b]{\linewidth}\raggedright
RUNSPEC
\end{minipage} & \begin{minipage}[b]{\linewidth}\raggedright
GRID
\end{minipage} & \begin{minipage}[b]{\linewidth}\raggedright
EDIT
\end{minipage} & \begin{minipage}[b]{\linewidth}\raggedright
PROPS
\end{minipage} & \begin{minipage}[b]{\linewidth}\raggedright
REGIONS
\end{minipage} & \begin{minipage}[b]{\linewidth}\raggedright
SOLUTION
\end{minipage} & \begin{minipage}[b]{\linewidth}\raggedright
SUMMARY
\end{minipage} & \begin{minipage}[b]{\linewidth}\raggedright
SCHEDULE
\end{minipage} \\
\midrule\noalign{}
\endhead
\bottomrule\noalign{}
\endlastfoot
\end{longtable}

\subsubsection{Description}\label{description-747}

The LSNUM keyword defines the saturation tables (relative permeability
and capillary pressure tables) region numbers for each grid block, as
such there must be one entry for each cell in the model. The region
number specifies which set of relative permeability tables (SWFN, SOF3
and related keywords) are used to calculate the relative permeability
and capillary pressure in a grid block. The keyword should only be used
if the Low Salinity option for the Brine model has been activated in the
RUNSPEC section using the LOWSALT keyword.

The oil wet curves are calculated as a weighted average of the low
salinity saturation tables (allocated by this keyword) and the high
salinity oil wet saturation tables (allocated by the SATNUM keyword),
using the weights provided by the LSALTFNC keyword in the PROPS section.

This keyword is not supported by OPM Flow but would change the results
if supported so the simulation will be stopped.

\begin{longtable}[]{@{}
  >{\raggedright\arraybackslash}p{(\linewidth - 6\tabcolsep) * \real{0.2500}}
  >{\raggedright\arraybackslash}p{(\linewidth - 6\tabcolsep) * \real{0.2500}}
  >{\raggedright\arraybackslash}p{(\linewidth - 6\tabcolsep) * \real{0.2500}}
  >{\raggedright\arraybackslash}p{(\linewidth - 6\tabcolsep) * \real{0.2500}}@{}}
\toprule\noalign{}
\begin{minipage}[b]{\linewidth}\raggedright
No.
\end{minipage} & \begin{minipage}[b]{\linewidth}\raggedright
Name
\end{minipage} & \begin{minipage}[b]{\linewidth}\raggedright
Description
\end{minipage} & \begin{minipage}[b]{\linewidth}\raggedright
Default
\end{minipage} \\
\midrule\noalign{}
\endhead
\bottomrule\noalign{}
\endlastfoot
1 & LSNUM & LSNUM defines an array of positive integers assigning a grid
cell to a particular saturation table region. The maximum number of
LSNUM regions is set by the NTSFUN variable on the TABDIMS keyword in
the RUNSPEC section. & 1 \\
Notes: & & & \\
\end{longtable}

\emph{Table 9.12: LSNUM Keyword Description}

If the Surfactant Wettability option have been activated by the SURFACTW
keyword, the LSNUM tables correspond to the immiscible low salinity
curves.

\subsubsection{Example}\label{example-421}

The example below sets three LSNUM regions for the model.

\begin{verbatim}
--
--       ARRAY       CONSTANT     ---------- BOX ---------
--                                I1  I2   J1  J2   K1  K2
EQUALS
         LSNUM       1            1*  1*   1*  1*   1*  1* / SET REGION 1
         LSNUM       2            1   2    1   2    1   1  / SET REGION 2
         LSNUM       3            1   2    1   2    2   2  / SET REGION 3
/
\end{verbatim}

\subsection{LWSLTNUM -- Define the Low Salt Oil Wet Saturation Table
Region
Numbers}\label{lwsltnum-define-the-low-salt-oil-wet-saturation-table-region-numbers}

\begin{longtable}[]{@{}
  >{\raggedright\arraybackslash}p{(\linewidth - 14\tabcolsep) * \real{0.1250}}
  >{\raggedright\arraybackslash}p{(\linewidth - 14\tabcolsep) * \real{0.1250}}
  >{\raggedright\arraybackslash}p{(\linewidth - 14\tabcolsep) * \real{0.1250}}
  >{\raggedright\arraybackslash}p{(\linewidth - 14\tabcolsep) * \real{0.1250}}
  >{\raggedright\arraybackslash}p{(\linewidth - 14\tabcolsep) * \real{0.1250}}
  >{\raggedright\arraybackslash}p{(\linewidth - 14\tabcolsep) * \real{0.1250}}
  >{\raggedright\arraybackslash}p{(\linewidth - 14\tabcolsep) * \real{0.1250}}
  >{\raggedright\arraybackslash}p{(\linewidth - 14\tabcolsep) * \real{0.1250}}@{}}
\toprule\noalign{}
\begin{minipage}[b]{\linewidth}\raggedright
RUNSPEC
\end{minipage} & \begin{minipage}[b]{\linewidth}\raggedright
GRID
\end{minipage} & \begin{minipage}[b]{\linewidth}\raggedright
EDIT
\end{minipage} & \begin{minipage}[b]{\linewidth}\raggedright
PROPS
\end{minipage} & \begin{minipage}[b]{\linewidth}\raggedright
REGIONS
\end{minipage} & \begin{minipage}[b]{\linewidth}\raggedright
SOLUTION
\end{minipage} & \begin{minipage}[b]{\linewidth}\raggedright
SUMMARY
\end{minipage} & \begin{minipage}[b]{\linewidth}\raggedright
SCHEDULE
\end{minipage} \\
\midrule\noalign{}
\endhead
\bottomrule\noalign{}
\endlastfoot
\end{longtable}

\subsubsection{Description}\label{description-748}

The LWSLTNUM keyword defines the saturation tables (relative
permeability and capillary pressure tables) region numbers for each grid
block, as such there must be one entry for each cell in the model. The
region number specifies which set of relative permeability tables (SWFN,
SOF3 and related keywords) are used to calculate the relative
permeability and capillary pressure in a grid block. The keyword should
only be used if the Low Salinity option for the Brine model has been
activated by the LOWSALT kwyword in the RUNSPEC section.

The oil wet curves are calculated as a weighted average of the low
salinity saturation tables (allocated by this keyword) and the high
salinity oil wet saturation tables (allocated by the SATNUM keyword in
the REGIONS section), using the weights provided by the LSALTFNC keyword
in the PROPS section.

This keyword is not supported by OPM Flow but would change the results
if supported so the simulation will be stopped.

\begin{longtable}[]{@{}
  >{\raggedright\arraybackslash}p{(\linewidth - 6\tabcolsep) * \real{0.2500}}
  >{\raggedright\arraybackslash}p{(\linewidth - 6\tabcolsep) * \real{0.2500}}
  >{\raggedright\arraybackslash}p{(\linewidth - 6\tabcolsep) * \real{0.2500}}
  >{\raggedright\arraybackslash}p{(\linewidth - 6\tabcolsep) * \real{0.2500}}@{}}
\toprule\noalign{}
\begin{minipage}[b]{\linewidth}\raggedright
No.
\end{minipage} & \begin{minipage}[b]{\linewidth}\raggedright
Name
\end{minipage} & \begin{minipage}[b]{\linewidth}\raggedright
Description
\end{minipage} & \begin{minipage}[b]{\linewidth}\raggedright
Default
\end{minipage} \\
\midrule\noalign{}
\endhead
\bottomrule\noalign{}
\endlastfoot
1 & LWSLTNUM & LWSLTNUM defines an array of positive integers assigning
a grid cell to a particular saturation table region. The maximum number
of LSLTWNUM regions is set by the NTSFUN variable on the TABDIMS keyword
in the RUNSPEC section. & 1 \\
Notes: & & & \\
\end{longtable}

\emph{Table 9.13: LWSLTNUM Keyword Description}

\subsubsection{Example}\label{example-422}

The example below sets three LWSLTNUM regions for the model.

\begin{verbatim}
--
--       ARRAY       CONSTANT     ---------- BOX ---------
--                                I1  I2   J1  J2   K1  K2
EQUALS
         LWSLTNUM    1            1*  1*   1*  1*   1*  1* / SET REGION 1
         LWSLTNUM    2            1   2    1   2    1   1  / SET REGION 2
         LWSLTNUM    3            1   2    1   2    2   2  / SET REGION 3
/
\end{verbatim}

\subsection{LWSNUM -- Define the Low Salt Water Wet Saturation Table
Region
Numbers}\label{lwsnum-define-the-low-salt-water-wet-saturation-table-region-numbers}

\begin{longtable}[]{@{}
  >{\raggedright\arraybackslash}p{(\linewidth - 14\tabcolsep) * \real{0.1250}}
  >{\raggedright\arraybackslash}p{(\linewidth - 14\tabcolsep) * \real{0.1250}}
  >{\raggedright\arraybackslash}p{(\linewidth - 14\tabcolsep) * \real{0.1250}}
  >{\raggedright\arraybackslash}p{(\linewidth - 14\tabcolsep) * \real{0.1250}}
  >{\raggedright\arraybackslash}p{(\linewidth - 14\tabcolsep) * \real{0.1250}}
  >{\raggedright\arraybackslash}p{(\linewidth - 14\tabcolsep) * \real{0.1250}}
  >{\raggedright\arraybackslash}p{(\linewidth - 14\tabcolsep) * \real{0.1250}}
  >{\raggedright\arraybackslash}p{(\linewidth - 14\tabcolsep) * \real{0.1250}}@{}}
\toprule\noalign{}
\begin{minipage}[b]{\linewidth}\raggedright
RUNSPEC
\end{minipage} & \begin{minipage}[b]{\linewidth}\raggedright
GRID
\end{minipage} & \begin{minipage}[b]{\linewidth}\raggedright
EDIT
\end{minipage} & \begin{minipage}[b]{\linewidth}\raggedright
PROPS
\end{minipage} & \begin{minipage}[b]{\linewidth}\raggedright
REGIONS
\end{minipage} & \begin{minipage}[b]{\linewidth}\raggedright
SOLUTION
\end{minipage} & \begin{minipage}[b]{\linewidth}\raggedright
SUMMARY
\end{minipage} & \begin{minipage}[b]{\linewidth}\raggedright
SCHEDULE
\end{minipage} \\
\midrule\noalign{}
\endhead
\bottomrule\noalign{}
\endlastfoot
\end{longtable}

\subsubsection{Description}\label{description-749}

The LWSNUM keyword defines the saturation tables (relative permeability
and capillary pressure tables) region numbers for each grid block, as
such there must be one entry for each cell in the model. The region
number specifies which set of relative permeability tables (SWFN, SOF3
and related keywords) are used to calculate the relative permeability
and capillary pressure in a grid block. The keyword should only be used
if the Low Salinity option for the Brine model and the Surfactant
Wettability option have been activated by the LOWSALT and SURFACTW
keywords, respectively, in the RUNSPEC section.

The water wet curves are calculated as a weighted average of the low
salinity saturation tables (allocated by this keyword) and the high
salinity water wet saturation tables (allocated by the HWSNUM keyword),
using the weights provided by the LSALTFNC keyword in the PROPS section.

This keyword is not supported by OPM Flow but would change the results
if supported so the simulation will be stopped.

\begin{longtable}[]{@{}
  >{\raggedright\arraybackslash}p{(\linewidth - 6\tabcolsep) * \real{0.2500}}
  >{\raggedright\arraybackslash}p{(\linewidth - 6\tabcolsep) * \real{0.2500}}
  >{\raggedright\arraybackslash}p{(\linewidth - 6\tabcolsep) * \real{0.2500}}
  >{\raggedright\arraybackslash}p{(\linewidth - 6\tabcolsep) * \real{0.2500}}@{}}
\toprule\noalign{}
\begin{minipage}[b]{\linewidth}\raggedright
No.
\end{minipage} & \begin{minipage}[b]{\linewidth}\raggedright
Name
\end{minipage} & \begin{minipage}[b]{\linewidth}\raggedright
Description
\end{minipage} & \begin{minipage}[b]{\linewidth}\raggedright
Default
\end{minipage} \\
\midrule\noalign{}
\endhead
\bottomrule\noalign{}
\endlastfoot
1 & LWSNUM & LWSNUM defines an array of positive integers assigning a
grid cell to a particular saturation table region. The maximum number of
LWSNUM regions is set by the NTSFUN variable on the TABDIMS keyword in
the RUNSPEC section. & 1 \\
Notes: & & & \\
\end{longtable}

\emph{Table 9.14: LWSNUM Keyword Description}

The HWSNUM allocated tables correspond to the immiscible high salinity
water wet curves.

\subsubsection{Example}\label{example-423}

The example below sets three LWSNUM regions for the model.

\begin{verbatim}
--
--       ARRAY       CONSTANT     ---------- BOX ---------
--                                I1  I2   J1  J2   K1  K2
EQUALS
         LWSNUM      1            1*  1*   1*  1*   1*  1* / SET REGION 1
         LWSNUM      2            1   2    1   2    1   1  / SET REGION 2
         LWSNUM      3            1   2    1   2    2   2  / SET REGION 3
/
\end{verbatim}

\subsection{MESSAGE -- Output User
Message}\label{message-output-user-message-5}

The MESSAGE keyword outputs a user message to the terminal, as well as
to the print (\emph{.PRT) and debug (}.DBG) files. Note this is
different to the MESSAGES keyword, that defines OPM Flows message print
limits and stop limits generated by the simulator.

See MESSAGE -- Output User Message in the GLOBAL section for a full
description.

\subsection{MESSAGES -- Define Message Print Limits and Stop
Limits}\label{messages-define-message-print-limits-and-stop-limits-5}

The MESSAGES keyword defines the print and stops levels for various
messages. The ``print limits'' set the maximum number of messages that
will be printed, after which no more messages will be printed and the
``stop limits'' terminate the run when these limits are exceeded. There
are six levels of message that increase in severity from informative all
the way to programming errors, as outlined in Table 4.5.

See MESSAGES -- Define Message Print Limits and Stop Limits in the
GLOBAL section for a full description.

\subsection{MISCNUM -- Define the Miscibility Region
Numbers}\label{miscnum-define-the-miscibility-region-numbers}

\begin{longtable}[]{@{}
  >{\raggedright\arraybackslash}p{(\linewidth - 14\tabcolsep) * \real{0.1250}}
  >{\raggedright\arraybackslash}p{(\linewidth - 14\tabcolsep) * \real{0.1250}}
  >{\raggedright\arraybackslash}p{(\linewidth - 14\tabcolsep) * \real{0.1250}}
  >{\raggedright\arraybackslash}p{(\linewidth - 14\tabcolsep) * \real{0.1250}}
  >{\raggedright\arraybackslash}p{(\linewidth - 14\tabcolsep) * \real{0.1250}}
  >{\raggedright\arraybackslash}p{(\linewidth - 14\tabcolsep) * \real{0.1250}}
  >{\raggedright\arraybackslash}p{(\linewidth - 14\tabcolsep) * \real{0.1250}}
  >{\raggedright\arraybackslash}p{(\linewidth - 14\tabcolsep) * \real{0.1250}}@{}}
\toprule\noalign{}
\begin{minipage}[b]{\linewidth}\raggedright
RUNSPEC
\end{minipage} & \begin{minipage}[b]{\linewidth}\raggedright
GRID
\end{minipage} & \begin{minipage}[b]{\linewidth}\raggedright
EDIT
\end{minipage} & \begin{minipage}[b]{\linewidth}\raggedright
PROPS
\end{minipage} & \begin{minipage}[b]{\linewidth}\raggedright
REGIONS
\end{minipage} & \begin{minipage}[b]{\linewidth}\raggedright
SOLUTION
\end{minipage} & \begin{minipage}[b]{\linewidth}\raggedright
SUMMARY
\end{minipage} & \begin{minipage}[b]{\linewidth}\raggedright
SCHEDULE
\end{minipage} \\
\midrule\noalign{}
\endhead
\bottomrule\noalign{}
\endlastfoot
\end{longtable}

\subsubsection{Description}\label{description-750}

The MISCNUM keyword defines the miscibility region number mixing tables
as defined by the TLMIXPAR keyword in the PROPS section, for when the
miscibility option has been activated by the MISCIBLE keyword in the
RUNSPEC section. MISCNUM also allocates miscible residual oil saturation
versus water saturation tables (SORWMIS keyword in the PROPS section)
used to calculate the relative permeability and PVT properties for a
grid cell.

Note that although this keyword can only be used when the miscibility
option is active, it is not necessary to use this keyword even if the
MISCIBLE keyword in the RUNSPEC has been activated as the default value
of one will be applied to all grid blocks. Secondly, a value of zero for
a grid cell results in immiscible fluids in that grid cell.

\begin{longtable}[]{@{}
  >{\raggedright\arraybackslash}p{(\linewidth - 6\tabcolsep) * \real{0.2500}}
  >{\raggedright\arraybackslash}p{(\linewidth - 6\tabcolsep) * \real{0.2500}}
  >{\raggedright\arraybackslash}p{(\linewidth - 6\tabcolsep) * \real{0.2500}}
  >{\raggedright\arraybackslash}p{(\linewidth - 6\tabcolsep) * \real{0.2500}}@{}}
\toprule\noalign{}
\begin{minipage}[b]{\linewidth}\raggedright
No.
\end{minipage} & \begin{minipage}[b]{\linewidth}\raggedright
Name
\end{minipage} & \begin{minipage}[b]{\linewidth}\raggedright
Description
\end{minipage} & \begin{minipage}[b]{\linewidth}\raggedright
Default
\end{minipage} \\
\midrule\noalign{}
\endhead
\bottomrule\noalign{}
\endlastfoot
1 & MISCNUM & MISCNUM defines an array of positive integers greater than
or equal to zero, that assign a grid cell to a particular table of
mixing parameters as defined by the TLMIXPAR and SORWMIS keywords. A
value of zero sets the fluids within a grid cell to be immiscible. The
maximum number of MISCNUM regions is set by the NTMIS variable on the
MISCIBLE keyword in the RUNSPEC section. & 1 \\
Notes: & & & \\
\end{longtable}

\emph{Table 9.15: MISCNUM Keyword Description}

See also the TLMIXPAR and SORWMIS keyword in the PROPS section.

\subsubsection{Example}\label{example-424}

The example below sets three MISCNUM regions in the model on a layer by
layer basis, using the EQUALS keyword.

\begin{verbatim}
--
--       ARRAY       CONSTANT     ---------- BOX ---------
--                                I1  I2   J1  J2   K1  K2
EQUALS
         MISCNUM      1            1*  1*   1*  1*   1   12 / SET REGION 1
         MISCNUM      2            1*  1*   1*  1*   13  55 / SET REGION 2
         MISCNUM      3            1*  1*   1*  1*   56 120 / SET REGION 3
/
\end{verbatim}

\subsection{MULTIPLY -- Multiply a Specified Array by a
Constant}\label{multiply-multiply-a-specified-array-by-a-constant-3}

The MULTIPLY keyword multiplies a specified array or part of an array by
a constant. The constant can be real or integer depending on the array
type; however, the arrays that can be operated on is dependent on which
section the keyword is being used.

See MULTIPLY -- Multiply a Specified Array by a Constant in the GRID
section for a full description.

\subsection{MULTIREG -- Multiply an Array by a Constant based on a
Region
Number}\label{multireg-multiply-an-array-by-a-constant-based-on-a-region-number-3}

The MULTIREG keyword multiplies an array or part of an array by a
constant for cells with a specific region number. The region number
array can be FLUXNUM, MULTNUM or OPERNUM and these arrays must be
defined and be available before the MULTIREG keyword is read by the
simulator. The constant can be real or integer depending on the property
array type; however, the property arrays that can be operated on is
dependent on which section the MULTIREG keyword is being used.

See MULTIREG -- Multiply an Array by a Constant based on a Region Number
in the GRID section for a full description.

\subsection{NOECHO -- Deactivate Echoing of User Input Files to the
Print
File}\label{noecho-deactivate-echoing-of-user-input-files-to-the-print-file-5}

Turns off echoing of all the input files to the print file. Note by
default echoing of the inputs files is active. but can subsequently be
switched off by the NOECHO activation keyword.

See NOECHO -- Deactivate Echoing of User Input Files to the Print File
in the GLOBAL section for a full description.

\subsection{NOWARN -- Deactivate Warning
Messages}\label{nowarn-deactivate-warning-messages-5}

Turns off warning messages to be printed to the print file; note that
this keyword is deactivated by default and can subsequently be switched
off by the WARN activation keyword. The warning messages may be turned
on and off using keywords WARN and NOWARN.

See NOWARN -- Deactivate Warning Messages in the GLOBAL section for a
full description.

\subsection{OPERATE -- Define Mathematical Operations on
Arrays}\label{operate-define-mathematical-operations-on-arrays-3}

This keyword, OPERATE, defines mathematical operations on property
arrays (NTG, PORO etc.) and optionally using another property array as
input to the function. The keyword allows for various mathematical
functions and their associated variables to be defined and applied to
the selected array data. Input constants can be real or integer
depending on the array type; however, the arrays that can be operated on
is dependent on which section the keyword is being used.

See OPERATE -- Define Mathematical Operations on Arrays in the GRID
section for a full description.

\subsection{OPERATER -- Define Mathematical Operations on Arrays by
Region}\label{operater-define-mathematical-operations-on-arrays-by-region-3}

This keyword defines the mathematical operations on arrays for specific
regions in the commercial simulator and is currently not supported by
OPM Flow. However, similar functionality is provided by the ADD and
MULTIPLY keywords.

See OPERATER -- Define Mathematical Operations on Arrays by Region in
the GRID section for a full description.

\subsection{OPERNUM -- Define Regions for Mathematical Operations on
Arrays}\label{opernum-define-regions-for-mathematical-operations-on-arrays-1}

This keyword defines the OPERATE region numbers for each grid block. The
OPERATE keyword defines mathematical operations on arrays in the
commercial simulator and is currently not supported by OPM Flow.
However, similar functionality is provided by the ADD and MULTIPLY
keywords combined with MULTNUM region array.

See OPERNUM -- Define Regions for Mathematical Operations on Arrays in
the GRID section for a full description.

\subsection{PENUM -- Define the Petro-Elastic Region
Numbers}\label{penum-define-the-petro-elastic-region-numbers}

\begin{longtable}[]{@{}
  >{\raggedright\arraybackslash}p{(\linewidth - 14\tabcolsep) * \real{0.1250}}
  >{\raggedright\arraybackslash}p{(\linewidth - 14\tabcolsep) * \real{0.1250}}
  >{\raggedright\arraybackslash}p{(\linewidth - 14\tabcolsep) * \real{0.1250}}
  >{\raggedright\arraybackslash}p{(\linewidth - 14\tabcolsep) * \real{0.1250}}
  >{\raggedright\arraybackslash}p{(\linewidth - 14\tabcolsep) * \real{0.1250}}
  >{\raggedright\arraybackslash}p{(\linewidth - 14\tabcolsep) * \real{0.1250}}
  >{\raggedright\arraybackslash}p{(\linewidth - 14\tabcolsep) * \real{0.1250}}
  >{\raggedright\arraybackslash}p{(\linewidth - 14\tabcolsep) * \real{0.1250}}@{}}
\toprule\noalign{}
\begin{minipage}[b]{\linewidth}\raggedright
RUNSPEC
\end{minipage} & \begin{minipage}[b]{\linewidth}\raggedright
GRID
\end{minipage} & \begin{minipage}[b]{\linewidth}\raggedright
EDIT
\end{minipage} & \begin{minipage}[b]{\linewidth}\raggedright
PROPS
\end{minipage} & \begin{minipage}[b]{\linewidth}\raggedright
REGIONS
\end{minipage} & \begin{minipage}[b]{\linewidth}\raggedright
SOLUTION
\end{minipage} & \begin{minipage}[b]{\linewidth}\raggedright
SUMMARY
\end{minipage} & \begin{minipage}[b]{\linewidth}\raggedright
SCHEDULE
\end{minipage} \\
\midrule\noalign{}
\endhead
\bottomrule\noalign{}
\endlastfoot
\end{longtable}

\subsubsection{Description}\label{description-751}

The PENNUM keyword defines the petro-elastic region number for each grid
block that is used to assign the of petro-elastic coefficients, bulk
modulus functions and shear modulus functions as defined by the PECOEFS,
PEKTAB and PEGTAB series of keywords in the PROPS section.

This keyword is not supported by OPM Flow but would change the results
if supported so the simulation will be stopped.

\subsection{PLMIXNUM -- Define the Polymer Region
Numbers}\label{plmixnum-define-the-polymer-region-numbers}

\begin{longtable}[]{@{}
  >{\raggedright\arraybackslash}p{(\linewidth - 14\tabcolsep) * \real{0.1250}}
  >{\raggedright\arraybackslash}p{(\linewidth - 14\tabcolsep) * \real{0.1250}}
  >{\raggedright\arraybackslash}p{(\linewidth - 14\tabcolsep) * \real{0.1250}}
  >{\raggedright\arraybackslash}p{(\linewidth - 14\tabcolsep) * \real{0.1250}}
  >{\raggedright\arraybackslash}p{(\linewidth - 14\tabcolsep) * \real{0.1250}}
  >{\raggedright\arraybackslash}p{(\linewidth - 14\tabcolsep) * \real{0.1250}}
  >{\raggedright\arraybackslash}p{(\linewidth - 14\tabcolsep) * \real{0.1250}}
  >{\raggedright\arraybackslash}p{(\linewidth - 14\tabcolsep) * \real{0.1250}}@{}}
\toprule\noalign{}
\begin{minipage}[b]{\linewidth}\raggedright
RUNSPEC
\end{minipage} & \begin{minipage}[b]{\linewidth}\raggedright
GRID
\end{minipage} & \begin{minipage}[b]{\linewidth}\raggedright
EDIT
\end{minipage} & \begin{minipage}[b]{\linewidth}\raggedright
PROPS
\end{minipage} & \begin{minipage}[b]{\linewidth}\raggedright
REGIONS
\end{minipage} & \begin{minipage}[b]{\linewidth}\raggedright
SOLUTION
\end{minipage} & \begin{minipage}[b]{\linewidth}\raggedright
SUMMARY
\end{minipage} & \begin{minipage}[b]{\linewidth}\raggedright
SCHEDULE
\end{minipage} \\
\midrule\noalign{}
\endhead
\bottomrule\noalign{}
\endlastfoot
\end{longtable}

\subsubsection{Description}\label{description-752}

The PLMIXNUM keyword defines the polymer region number for each grid
block that is used to assign the mixing tables as well as the maximum
polymer and salt concentrations, as defined by the PLMIXPAR and PLYMAX
keywords in the PROPS section, for when the polymer option has been
activated by the POLYMER keyword in the RUNSPEC section.

The maximum polymer concentration and the associated salt concentration
are declared on the PLYMAX keyword.

\begin{longtable}[]{@{}
  >{\raggedright\arraybackslash}p{(\linewidth - 6\tabcolsep) * \real{0.2500}}
  >{\raggedright\arraybackslash}p{(\linewidth - 6\tabcolsep) * \real{0.2500}}
  >{\raggedright\arraybackslash}p{(\linewidth - 6\tabcolsep) * \real{0.2500}}
  >{\raggedright\arraybackslash}p{(\linewidth - 6\tabcolsep) * \real{0.2500}}@{}}
\toprule\noalign{}
\begin{minipage}[b]{\linewidth}\raggedright
No.
\end{minipage} & \begin{minipage}[b]{\linewidth}\raggedright
Name
\end{minipage} & \begin{minipage}[b]{\linewidth}\raggedright
Description
\end{minipage} & \begin{minipage}[b]{\linewidth}\raggedright
Default
\end{minipage} \\
\midrule\noalign{}
\endhead
\bottomrule\noalign{}
\endlastfoot
1 & PLMIXNUM & PLMIXNUM defines an array of positive integers greater
than or equal to one, that assign a grid cell to a particular table of
mixing parameters as defined by the PLMIXPAR and PLYMAX keywords. The
maximum number of PLMIXNUM regions is set by the NPLMIX variable on the
REGDIMS keyword in the RUNSPEC section. & 1 \\
Notes: & & & \\
\end{longtable}

\emph{Table 9.16: PLMIXNUM Keyword Description}

See also the PLYADS, PLYADSS, PLYDHLF, PLYMAX, PLYROCK, PLYSHEAR,
PLYSHLOG and PLYVISC keywords in the PROPS section.

\subsubsection{Example}\label{example-425}

The example below sets three PLMIXNUM regions in the model on a layer by
layer basis, using the EQUALS keyword.

\begin{verbatim}
--
--       ARRAY       CONSTANT     ---------- BOX ---------
--                                I1  I2   J1  J2   K1  K2
EQUALS
         PLMIXNUM      1          1*  1*   1*  1*   1   12 / SET REGION 1
         PLMIXNUM      2          1*  1*   1*  1*   13  55 / SET REGION 2
         PLMIXNUM      3          1*  1*   1*  1*   56 120 / SET REGION 3
/
\end{verbatim}

\subsection{PVTNUM -- Define the PVT
Regions}\label{pvtnum-define-the-pvt-regions}

\begin{longtable}[]{@{}
  >{\raggedright\arraybackslash}p{(\linewidth - 14\tabcolsep) * \real{0.1250}}
  >{\raggedright\arraybackslash}p{(\linewidth - 14\tabcolsep) * \real{0.1250}}
  >{\raggedright\arraybackslash}p{(\linewidth - 14\tabcolsep) * \real{0.1250}}
  >{\raggedright\arraybackslash}p{(\linewidth - 14\tabcolsep) * \real{0.1250}}
  >{\raggedright\arraybackslash}p{(\linewidth - 14\tabcolsep) * \real{0.1250}}
  >{\raggedright\arraybackslash}p{(\linewidth - 14\tabcolsep) * \real{0.1250}}
  >{\raggedright\arraybackslash}p{(\linewidth - 14\tabcolsep) * \real{0.1250}}
  >{\raggedright\arraybackslash}p{(\linewidth - 14\tabcolsep) * \real{0.1250}}@{}}
\toprule\noalign{}
\begin{minipage}[b]{\linewidth}\raggedright
RUNSPEC
\end{minipage} & \begin{minipage}[b]{\linewidth}\raggedright
GRID
\end{minipage} & \begin{minipage}[b]{\linewidth}\raggedright
EDIT
\end{minipage} & \begin{minipage}[b]{\linewidth}\raggedright
PROPS
\end{minipage} & \begin{minipage}[b]{\linewidth}\raggedright
REGIONS
\end{minipage} & \begin{minipage}[b]{\linewidth}\raggedright
SOLUTION
\end{minipage} & \begin{minipage}[b]{\linewidth}\raggedright
SUMMARY
\end{minipage} & \begin{minipage}[b]{\linewidth}\raggedright
SCHEDULE
\end{minipage} \\
\midrule\noalign{}
\endhead
\bottomrule\noalign{}
\endlastfoot
\end{longtable}

\subsubsection{Description}\label{description-753}

The PVTNUM keyword defines the PVT region numbers for each grid block,
as such there must be one entry for each cell in the model. The region
number specifies which set of PVT tables (DENSITY, PVDG, PVDO, PVTG,
PVTO, PVCO, PVTW and ROCK) are used to calculate the PVT properties in a
grid block.

\begin{longtable}[]{@{}
  >{\raggedright\arraybackslash}p{(\linewidth - 6\tabcolsep) * \real{0.2500}}
  >{\raggedright\arraybackslash}p{(\linewidth - 6\tabcolsep) * \real{0.2500}}
  >{\raggedright\arraybackslash}p{(\linewidth - 6\tabcolsep) * \real{0.2500}}
  >{\raggedright\arraybackslash}p{(\linewidth - 6\tabcolsep) * \real{0.2500}}@{}}
\toprule\noalign{}
\begin{minipage}[b]{\linewidth}\raggedright
No.
\end{minipage} & \begin{minipage}[b]{\linewidth}\raggedright
Name
\end{minipage} & \begin{minipage}[b]{\linewidth}\raggedright
Description
\end{minipage} & \begin{minipage}[b]{\linewidth}\raggedright
Default
\end{minipage} \\
\midrule\noalign{}
\endhead
\bottomrule\noalign{}
\endlastfoot
1 & PVTNUM & PVTNUM defines an array of positive integers assigning a
grid cell to a particular PVT region. The maximum number of PVTNUM
regions is set by the NTPVT variable on the TABDIMS keyword in the
RUNSPEC section. & 1 \\
Notes: & & & \\
\end{longtable}

\emph{Table 9.17: PVTNUM Keyword Description}

\begin{tcolorbox}[enhanced jigsaw, title=\textcolor{quarto-callout-note-color}{\faInfo}\hspace{0.5em}{Note}, coltitle=black, breakable, toprule=.15mm, colbacktitle=quarto-callout-note-color!10!white, leftrule=.75mm, opacitybacktitle=0.6, colback=white, bottomtitle=1mm, toptitle=1mm, titlerule=0mm, left=2mm, colframe=quarto-callout-note-color-frame, arc=.35mm, rightrule=.15mm, bottomrule=.15mm, opacityback=0]

Care should be taken that cells in different PVTNUM regions are not in
communication, since the fluid properties are associated with a cell. If
for example, a rbbl or a rm3 of oil flows from PVTNUM region 1 to PVTNUM
region 2, then the oil properties of that oil will change from the PVT 1
data set to the PVT data set 2. This will result in material balance
errors, that may or may not cause numerical issues. To avoid this one
should use the MULTNUM (or FLUXNUM, or OPERNUM) array with the MULTREGT
array to ensure that the various PVTNUM regions are not in
communication.

\end{tcolorbox}

\subsubsection{Examples}\label{examples-52}

The example below sets three PVTNUM regions for a 4 x 5 x 2 model.

\begin{verbatim}
--
--       DEFINE PVTNUM REGION FOR ALL CELLS
--
PVTNUM
         2  2  1  1  2  2  1  1  1  1  1  1  1  1  1  1  1  1  1  1
         3  3  1  1  3  3  1  1  1  1  1  1  1  1  1  1  1  1  1  1
/
\end{verbatim}

Alternatively the EQUALS keyword could be employed to accomplish the
same task, that is:

\begin{verbatim}
--
--       ARRAY       CONSTANT     ---------- BOX ---------
--                                I1  I2   J1  J2   K1  K2
EQUALS
         PVTNUM      1            1*  1*   1*  1*   1*  1* / SET REGION 1
         PVTNUM      2            1   2    1   2    1   1  / SET REGION 2
         PVTNUM      3            1   2    1   2    2   2  / SET REGION 3
/
\end{verbatim}

There third example shows how to ensure the various PVT regions are
isolated. First of all define the MULTNUM array in the GRID section and
ensure all the regions are isolated.

\begin{verbatim}
-- ==============================================================================
--
-- GRID SECTION
--
-- ==============================================================================
GRID
--
--       ARRAY       CONSTANT     ---------- BOX ---------
--                                I1  I2   J1  J2   K1  K2
EQUALS
         MULTNUM     1            1*  1*   1*  1*   1*  1* / SET REGION 1
         MULTNUM     2            1   2    1   2    1   1  / SET REGION 2
         MULTNUM     3            1   2    1   2    2   2  / SET REGION 3
/
--
--       SET TRANSMISSIBILITES ACROSS DIFFERENT RESERVOIRS TO ZERO TO ISOLATE
--       RESERVOIRS
--
--       REGION   REGION   TRANS   DIREC   NNC    REGION ARRAY
--       FROM     TO       MULT    OPT     OPTS   M / F / O
MULTREGT
         1*       1*       0.0     1*     'ALL'   M        / ALL REGIONS SEALED
/
\end{verbatim}

Then in the REGIONS section copy the MULTNUM array to the PVTNUM array.

\begin{verbatim}
-- ==============================================================================
--
-- REGIONS SECTION
--
-- ==============================================================================
REGIONS

--
--       COPY AN ARRAY TO ANOTHER ARRAY BASED ON A REGION NUMBER
--
--       ARRAY     ARRAY     REGION   REGION ARRAY
--       FROM      TO        NUMBER    M / F / O
COPYREG
         MULTNUM   PVTNUM    1         M                   / COPY MULT TO PVT 1
         MULTNUM   PVTNUM    2         M                   / COPY MULT TO PVT 2
         MULTNUM   PVTNUM    3         M                   / COPY MULT TO PVT 3
/
\end{verbatim}

All the separate PVT regions are now isolated.

\subsection{PYEND -- End the Definition of a PYINPUT
Section}\label{pyend-end-the-definition-of-a-pyinput-section-3}

The PYINPUT and PYEND keywords are a part of OPM Flow's Python scripting
facility that processes standard Python commands that can be used to
manipulate and define the simulators input parameters during processing
of the input deck. The main purpose of the facility is to script the
construction of the various keywords. A PYINPUT Definition Section is
terminated by a PYEND keyword on a separate single line.

See PYEND -- End the Definition of a PYINPUT Section in the GRID section
for a full description.

\subsection{PYINPUT -- Define the Start of a PYINPUT
Section}\label{pyinput-define-the-start-of-a-pyinput-section-3}

The PYINPUT and PYEND keywords are a part of OPM Flow's Python scripting
facility that processes standard Python commands that can be used to
manipulate and define the simulators input parameters during processing
of the input deck. The main purpose of the facility is to script the
construction of the various keywords used by the simulator. PYINPUT
declares the start of a PYINPUT Definition Section on a single separate
line, which is then followed by various standard Python commands, with
one command per line.

See PYINPUT -- Define the Start of a PYINPUT Section in the GRID section
for a full description.

\subsection{REFINE -- Start the Definition of a Local Grid
Refinement}\label{refine-start-the-definition-of-a-local-grid-refinement-3}

The REFINE keyword defines the start of a Cartesian or radial local grid
refinement (``LGR'') definition that sets the properties of the selected
LGR. The keyword is then followed by the property keywords associated
with the section where the keyword is being invoked. For example, if the
REFINE keyword is used in the GRID section then most of the keywords in
that section can be used to set the grid properties for the LGR.

See REFINE -- Start the Definition of a Local Grid Refinement in the
GRID section for a full description.

\subsection{REGIONS -- Define the Start of the REGIONS Section of
Keywords}\label{regions-define-the-start-of-the-regions-section-of-keywords}

\begin{longtable}[]{@{}
  >{\raggedright\arraybackslash}p{(\linewidth - 14\tabcolsep) * \real{0.1250}}
  >{\raggedright\arraybackslash}p{(\linewidth - 14\tabcolsep) * \real{0.1250}}
  >{\raggedright\arraybackslash}p{(\linewidth - 14\tabcolsep) * \real{0.1250}}
  >{\raggedright\arraybackslash}p{(\linewidth - 14\tabcolsep) * \real{0.1250}}
  >{\raggedright\arraybackslash}p{(\linewidth - 14\tabcolsep) * \real{0.1250}}
  >{\raggedright\arraybackslash}p{(\linewidth - 14\tabcolsep) * \real{0.1250}}
  >{\raggedright\arraybackslash}p{(\linewidth - 14\tabcolsep) * \real{0.1250}}
  >{\raggedright\arraybackslash}p{(\linewidth - 14\tabcolsep) * \real{0.1250}}@{}}
\toprule\noalign{}
\begin{minipage}[b]{\linewidth}\raggedright
RUNSPEC
\end{minipage} & \begin{minipage}[b]{\linewidth}\raggedright
GRID
\end{minipage} & \begin{minipage}[b]{\linewidth}\raggedright
EDIT
\end{minipage} & \begin{minipage}[b]{\linewidth}\raggedright
PROPS
\end{minipage} & \begin{minipage}[b]{\linewidth}\raggedright
REGIONS
\end{minipage} & \begin{minipage}[b]{\linewidth}\raggedright
SOLUTION
\end{minipage} & \begin{minipage}[b]{\linewidth}\raggedright
SUMMARY
\end{minipage} & \begin{minipage}[b]{\linewidth}\raggedright
SCHEDULE
\end{minipage} \\
\midrule\noalign{}
\endhead
\bottomrule\noalign{}
\endlastfoot
\end{longtable}

\subsubsection{Description}\label{description-754}

The REGIONS activation keyword marks the end of the PROPS section and
the start of the REGIONS section that defines how various fluid and rock
property data defined in the PROPS section are allocated to the
individual cells in the model.

There is no data required for this keyword.

\subsubsection{Example}\label{example-426}

\begin{verbatim}
-- ==============================================================================
--
-- REGIONS SECTION
--
-- ==============================================================================
REGIONS
\end{verbatim}

The above example marks the end of the PROPS section and the start of
the REGIONS section in the OPM Flow data input file.

\subsection{RESIDNUM -- Define Vertical Equilibrium Residual Flow Region
Numbers}\label{residnum-define-vertical-equilibrium-residual-flow-region-numbers}

\begin{longtable}[]{@{}
  >{\raggedright\arraybackslash}p{(\linewidth - 14\tabcolsep) * \real{0.1250}}
  >{\raggedright\arraybackslash}p{(\linewidth - 14\tabcolsep) * \real{0.1250}}
  >{\raggedright\arraybackslash}p{(\linewidth - 14\tabcolsep) * \real{0.1250}}
  >{\raggedright\arraybackslash}p{(\linewidth - 14\tabcolsep) * \real{0.1250}}
  >{\raggedright\arraybackslash}p{(\linewidth - 14\tabcolsep) * \real{0.1250}}
  >{\raggedright\arraybackslash}p{(\linewidth - 14\tabcolsep) * \real{0.1250}}
  >{\raggedright\arraybackslash}p{(\linewidth - 14\tabcolsep) * \real{0.1250}}
  >{\raggedright\arraybackslash}p{(\linewidth - 14\tabcolsep) * \real{0.1250}}@{}}
\toprule\noalign{}
\begin{minipage}[b]{\linewidth}\raggedright
RUNSPEC
\end{minipage} & \begin{minipage}[b]{\linewidth}\raggedright
GRID
\end{minipage} & \begin{minipage}[b]{\linewidth}\raggedright
EDIT
\end{minipage} & \begin{minipage}[b]{\linewidth}\raggedright
PROPS
\end{minipage} & \begin{minipage}[b]{\linewidth}\raggedright
REGIONS
\end{minipage} & \begin{minipage}[b]{\linewidth}\raggedright
SOLUTION
\end{minipage} & \begin{minipage}[b]{\linewidth}\raggedright
SUMMARY
\end{minipage} & \begin{minipage}[b]{\linewidth}\raggedright
SCHEDULE
\end{minipage} \\
\midrule\noalign{}
\endhead
\bottomrule\noalign{}
\endlastfoot
\end{longtable}

\subsubsection{Description}\label{description-755}

The RESIDNUM keyword defines the Vertical Equilibrium (``VE'') residual
flow calculation saturation tables (relative permeability and capillary
pressure tables) region numbers for each grid block, as such there must
be one entry for each cell in the model. The region number specifies
which set of relative permeability tables (SGFN, SWFN, SOF2, SOF3,
SOF32D, SGOF, SLGOF and SWOF) are used to calculate the relative
permeability and capillary pressure in a grid block. The keyword should
only be used if the Vertical Equilibrium option has been invoked via the
VE keyword in the RUNSPEC section.

This keyword is not supported by OPM Flow but would change the results
if supported so the simulation will be stopped.

\begin{longtable}[]{@{}
  >{\raggedright\arraybackslash}p{(\linewidth - 6\tabcolsep) * \real{0.2500}}
  >{\raggedright\arraybackslash}p{(\linewidth - 6\tabcolsep) * \real{0.2500}}
  >{\raggedright\arraybackslash}p{(\linewidth - 6\tabcolsep) * \real{0.2500}}
  >{\raggedright\arraybackslash}p{(\linewidth - 6\tabcolsep) * \real{0.2500}}@{}}
\toprule\noalign{}
\begin{minipage}[b]{\linewidth}\raggedright
No.
\end{minipage} & \begin{minipage}[b]{\linewidth}\raggedright
Name
\end{minipage} & \begin{minipage}[b]{\linewidth}\raggedright
Description
\end{minipage} & \begin{minipage}[b]{\linewidth}\raggedright
Default
\end{minipage} \\
\midrule\noalign{}
\endhead
\bottomrule\noalign{}
\endlastfoot
1 & RESIDNUM & RESIDNUM defines an array of positive integers assigning
a grid cell to a particular saturation table region. The maximum number
of RESIDNUM regions is set by the NTSFUN variable on the TABDIMS keyword
in the RUNSPEC section. & 1 \\
Notes: & & & \\
\end{longtable}

\emph{Table 9.18: RESIDNUM Keyword Description}

Note that any capillary pressure data on the relative permeability
tables (SGFN, SWFN, SOF2, SOF3, SOF32D, SGOF, SLGOF and SWOF) will be
ignored by the VE option.

\subsubsection{Example}\label{example-427}

The example below sets three RESIDNUM regions for the model, by first
setting all values to one, then setting layers 2 to 10 to two, and
finally setting layers 30 to 50 to three.

\begin{verbatim}
--
--       ARRAY       CONSTANT     ---------- BOX ---------
--                                I1  I2   J1  J2   K1  K2
EQUALS
         SATNUM      1            1*  1*   1*  1*   1*  1* / SET REGION 1
         SATNUM      2            1   2    1   2    2   10 / SET REGION 2
         SATNUM      3            1   2    1   2    30  50 / SET REGION 3
/
\end{verbatim}

\subsection{ROCKNUM -- Define Rock Compaction Table Region
Numbers}\label{rocknum-define-rock-compaction-table-region-numbers}

\begin{longtable}[]{@{}
  >{\raggedright\arraybackslash}p{(\linewidth - 14\tabcolsep) * \real{0.1250}}
  >{\raggedright\arraybackslash}p{(\linewidth - 14\tabcolsep) * \real{0.1250}}
  >{\raggedright\arraybackslash}p{(\linewidth - 14\tabcolsep) * \real{0.1250}}
  >{\raggedright\arraybackslash}p{(\linewidth - 14\tabcolsep) * \real{0.1250}}
  >{\raggedright\arraybackslash}p{(\linewidth - 14\tabcolsep) * \real{0.1250}}
  >{\raggedright\arraybackslash}p{(\linewidth - 14\tabcolsep) * \real{0.1250}}
  >{\raggedright\arraybackslash}p{(\linewidth - 14\tabcolsep) * \real{0.1250}}
  >{\raggedright\arraybackslash}p{(\linewidth - 14\tabcolsep) * \real{0.1250}}@{}}
\toprule\noalign{}
\begin{minipage}[b]{\linewidth}\raggedright
RUNSPEC
\end{minipage} & \begin{minipage}[b]{\linewidth}\raggedright
GRID
\end{minipage} & \begin{minipage}[b]{\linewidth}\raggedright
EDIT
\end{minipage} & \begin{minipage}[b]{\linewidth}\raggedright
PROPS
\end{minipage} & \begin{minipage}[b]{\linewidth}\raggedright
REGIONS
\end{minipage} & \begin{minipage}[b]{\linewidth}\raggedright
SOLUTION
\end{minipage} & \begin{minipage}[b]{\linewidth}\raggedright
SUMMARY
\end{minipage} & \begin{minipage}[b]{\linewidth}\raggedright
SCHEDULE
\end{minipage} \\
\midrule\noalign{}
\endhead
\bottomrule\noalign{}
\endlastfoot
\end{longtable}

\subsubsection{Description}\label{description-756}

The ROCKNUM keyword defines the rock compaction table region numbers for
each grid block, as such there must be one entry for each cell in the
model. The region number specifies which set of rock compaction tables
defined by the ROCKTAB keyword are used to calculate the rock compaction
in a grid block.

\begin{longtable}[]{@{}
  >{\raggedright\arraybackslash}p{(\linewidth - 6\tabcolsep) * \real{0.2500}}
  >{\raggedright\arraybackslash}p{(\linewidth - 6\tabcolsep) * \real{0.2500}}
  >{\raggedright\arraybackslash}p{(\linewidth - 6\tabcolsep) * \real{0.2500}}
  >{\raggedright\arraybackslash}p{(\linewidth - 6\tabcolsep) * \real{0.2500}}@{}}
\toprule\noalign{}
\begin{minipage}[b]{\linewidth}\raggedright
No.
\end{minipage} & \begin{minipage}[b]{\linewidth}\raggedright
Name
\end{minipage} & \begin{minipage}[b]{\linewidth}\raggedright
Description
\end{minipage} & \begin{minipage}[b]{\linewidth}\raggedright
Default
\end{minipage} \\
\midrule\noalign{}
\endhead
\bottomrule\noalign{}
\endlastfoot
1 & ROCKNUM & ROCKNUM defines an array of positive integers assigning a
grid cell to a particular rock compaction table region. The maximum
number of ROCKNUM regions is set by the NTROCC variable on the ROCKCOMP
keyword in the RUNSPEC section. & 1 \\
Notes: & & & \\
\end{longtable}

\emph{Table 9.19: ROCKNUM Keyword Description}

\subsubsection{Examples}\label{examples-53}

The example below sets three ROCKNUM regions for a 4 x 5 x 2 model.

\begin{verbatim}
--
--       DEFINE ROCKNUM REGION FOR ALL CELLS
--
ROCKNUM
         2  2  1  1  2  2  1  1  1  1  1  1  1  1  1  1  1  1  1  1
         3  3  1  1  3  3  1  1  1  1  1  1  1  1  1  1  1  1  1  1
/
\end{verbatim}

Alternatively the EQUALS keyword could be employed to accomplish the
same task, that is:

\begin{verbatim}
--
--       ARRAY       CONSTANT     ---------- BOX ---------
--                                I1  I2   J1  J2   K1  K2
EQUALS
         ROCKNUM     1            1*  1*   1*  1*   1*  1* / SET REGION 1
         ROCKNUM     2            1   2    1   2    1   1  / SET REGION 2
         ROCKNUM     3            1   2    1   2    2   2  / SET REGION 3
/
\end{verbatim}

\subsection{RPTREGS -- Define REGIONS Section
Reporting}\label{rptregs-define-regions-section-reporting}

\begin{longtable}[]{@{}
  >{\raggedright\arraybackslash}p{(\linewidth - 14\tabcolsep) * \real{0.1250}}
  >{\raggedright\arraybackslash}p{(\linewidth - 14\tabcolsep) * \real{0.1250}}
  >{\raggedright\arraybackslash}p{(\linewidth - 14\tabcolsep) * \real{0.1250}}
  >{\raggedright\arraybackslash}p{(\linewidth - 14\tabcolsep) * \real{0.1250}}
  >{\raggedright\arraybackslash}p{(\linewidth - 14\tabcolsep) * \real{0.1250}}
  >{\raggedright\arraybackslash}p{(\linewidth - 14\tabcolsep) * \real{0.1250}}
  >{\raggedright\arraybackslash}p{(\linewidth - 14\tabcolsep) * \real{0.1250}}
  >{\raggedright\arraybackslash}p{(\linewidth - 14\tabcolsep) * \real{0.1250}}@{}}
\toprule\noalign{}
\begin{minipage}[b]{\linewidth}\raggedright
RUNSPEC
\end{minipage} & \begin{minipage}[b]{\linewidth}\raggedright
GRID
\end{minipage} & \begin{minipage}[b]{\linewidth}\raggedright
EDIT
\end{minipage} & \begin{minipage}[b]{\linewidth}\raggedright
PROPS
\end{minipage} & \begin{minipage}[b]{\linewidth}\raggedright
REGIONS
\end{minipage} & \begin{minipage}[b]{\linewidth}\raggedright
SOLUTION
\end{minipage} & \begin{minipage}[b]{\linewidth}\raggedright
SUMMARY
\end{minipage} & \begin{minipage}[b]{\linewidth}\raggedright
SCHEDULE
\end{minipage} \\
\midrule\noalign{}
\endhead
\bottomrule\noalign{}
\endlastfoot
\end{longtable}

\subsubsection{Description}\label{description-757}

This keyword defines the data in the REGIONS section that is to be
printed to the output print file in human readable format. The keyword
has two distinct forms, the first of which consists of the keyword
followed by a series of integers on the next line indicating the data to
be printed (see the first example). This is the original formal in the
commercial simulator and was subsequently superseded by the second
format. The second format consists of the keyword followed by a series
of character strings that indicate the data to be printed. In most cases
the character string is the keyword used to load the data in the OPM
Flow input deck, for example FIPNUM for the fluid in-place array. Its is
anticipated that OPM Flow will eventually support the functionality of
the second format only, the first format although recognized will be
completely ignored.

\begin{longtable}[]{@{}
  >{\raggedright\arraybackslash}p{(\linewidth - 6\tabcolsep) * \real{0.2500}}
  >{\raggedright\arraybackslash}p{(\linewidth - 6\tabcolsep) * \real{0.2500}}
  >{\raggedright\arraybackslash}p{(\linewidth - 6\tabcolsep) * \real{0.2500}}
  >{\raggedright\arraybackslash}p{(\linewidth - 6\tabcolsep) * \real{0.2500}}@{}}
\toprule\noalign{}
\begin{minipage}[b]{\linewidth}\raggedright
No.
\end{minipage} & \begin{minipage}[b]{\linewidth}\raggedright
Name
\end{minipage} & \begin{minipage}[b]{\linewidth}\raggedright
Description
\end{minipage} & \begin{minipage}[b]{\linewidth}\raggedright
Default
\end{minipage} \\
\midrule\noalign{}
\endhead
\bottomrule\noalign{}
\endlastfoot
1 & EQLNUM & Print the equilibration region array. & N/A \\
2 & FIPNUM & Print the fluid in-place array. & N/A \\
3 & PVTNUM & Print the PVT table assignment array. & N/A \\
4 & SATNUM & Print the saturation function (relative permeability)
assignment array. & N/A \\
\ldots. & & \ldots. & N/A \\
Notes: & & & \\
\end{longtable}

\emph{Table 9.20: RPTREGS Keyword Description}

This keyword is not supported by OPM Flow but has no effect on the
results so it will be ignored.

\begin{tcolorbox}[enhanced jigsaw, title=\textcolor{quarto-callout-note-color}{\faInfo}\hspace{0.5em}{Note}, coltitle=black, breakable, toprule=.15mm, colbacktitle=quarto-callout-note-color!10!white, leftrule=.75mm, opacitybacktitle=0.6, colback=white, bottomtitle=1mm, toptitle=1mm, titlerule=0mm, left=2mm, colframe=quarto-callout-note-color-frame, arc=.35mm, rightrule=.15mm, bottomrule=.15mm, opacityback=0]

This keyword has the potential to produce very large print files that
some text editors may have difficulty loading, coupled with the fact
that reviewing the data in this format is very cumbersome. A more
efficient solution is to load the *.INIT file into OPM ResInsight to
view the data graphically, this also has the benefit of being able to
filter the grid based on I, J, K ranges and grid properties.

\end{tcolorbox}

\subsubsection{Examples}\label{examples-54}

The first example shows the original format of this keyword; although
the keyword and format are recognized by OPM Flow, the format is ignored
and is unlikely to be implemented in in the simulator.

\begin{verbatim}
--
--       DEFINE REGIONS SECTION REPORT OPTION (ORIGINAL FORMAT)
–
RPTREGS
         1        2*0      1        3*1                                        /
\end{verbatim}

The next example shows the second format of the keyword which may be
supported in a future release of OPM Flow.

\begin{verbatim}
--       DEFINE REGIONS SECTION REPORT OPTIONS
--
RPTREGS
         FIPMUM   EQLNUM   PVTNUM   SATNUM                                     /
\end{verbatim}

\subsection{SATNUM -- Define the Saturation Table Region
Numbers}\label{satnum-define-the-saturation-table-region-numbers}

\begin{longtable}[]{@{}
  >{\raggedright\arraybackslash}p{(\linewidth - 14\tabcolsep) * \real{0.1250}}
  >{\raggedright\arraybackslash}p{(\linewidth - 14\tabcolsep) * \real{0.1250}}
  >{\raggedright\arraybackslash}p{(\linewidth - 14\tabcolsep) * \real{0.1250}}
  >{\raggedright\arraybackslash}p{(\linewidth - 14\tabcolsep) * \real{0.1250}}
  >{\raggedright\arraybackslash}p{(\linewidth - 14\tabcolsep) * \real{0.1250}}
  >{\raggedright\arraybackslash}p{(\linewidth - 14\tabcolsep) * \real{0.1250}}
  >{\raggedright\arraybackslash}p{(\linewidth - 14\tabcolsep) * \real{0.1250}}
  >{\raggedright\arraybackslash}p{(\linewidth - 14\tabcolsep) * \real{0.1250}}@{}}
\toprule\noalign{}
\begin{minipage}[b]{\linewidth}\raggedright
RUNSPEC
\end{minipage} & \begin{minipage}[b]{\linewidth}\raggedright
GRID
\end{minipage} & \begin{minipage}[b]{\linewidth}\raggedright
EDIT
\end{minipage} & \begin{minipage}[b]{\linewidth}\raggedright
PROPS
\end{minipage} & \begin{minipage}[b]{\linewidth}\raggedright
REGIONS
\end{minipage} & \begin{minipage}[b]{\linewidth}\raggedright
SOLUTION
\end{minipage} & \begin{minipage}[b]{\linewidth}\raggedright
SUMMARY
\end{minipage} & \begin{minipage}[b]{\linewidth}\raggedright
SCHEDULE
\end{minipage} \\
\midrule\noalign{}
\endhead
\bottomrule\noalign{}
\endlastfoot
\end{longtable}

\subsubsection{Description}\label{description-758}

The SATNUM keyword defines the saturation tables (relative permeability
and capillary pressure tables) region numbers for each grid block, as
such there must be one entry for each cell in the model. The region
number specifies which set of relative permeability tables (SGFN, SWFN,
SOF2, SOF3, SOF32D, SGOF, SLGOF and SWOF) are used to calculate the
relative permeability and capillary pressure in a grid block.

\begin{longtable}[]{@{}
  >{\raggedright\arraybackslash}p{(\linewidth - 6\tabcolsep) * \real{0.2500}}
  >{\raggedright\arraybackslash}p{(\linewidth - 6\tabcolsep) * \real{0.2500}}
  >{\raggedright\arraybackslash}p{(\linewidth - 6\tabcolsep) * \real{0.2500}}
  >{\raggedright\arraybackslash}p{(\linewidth - 6\tabcolsep) * \real{0.2500}}@{}}
\toprule\noalign{}
\begin{minipage}[b]{\linewidth}\raggedright
No.
\end{minipage} & \begin{minipage}[b]{\linewidth}\raggedright
Name
\end{minipage} & \begin{minipage}[b]{\linewidth}\raggedright
Description
\end{minipage} & \begin{minipage}[b]{\linewidth}\raggedright
Default
\end{minipage} \\
\midrule\noalign{}
\endhead
\bottomrule\noalign{}
\endlastfoot
1 & SATNUM & SATNUM defines an array of positive integers assigning a
grid cell to a particular saturation table region. The maximum number of
SATNUM regions is set by the NTSFUN variable on the TABDIMS keyword in
the RUNSPEC section. & 1 \\
Notes: & & & \\
\end{longtable}

\emph{Table 9.21: SATNUM Keyword Description}

\subsubsection{Examples}\label{examples-55}

The example below sets three SATNUM regions for a 4 x 5 x 2 model.

\begin{verbatim}
--
--       DEFINE SATNUM REGIONS FOR ALL CELLS
--
SATNUM
         2  2  1  1  2  2  1  1  1  1  1  1  1  1  1  1  1  1  1  1
         3  3  1  1  3  3  1  1  1  1  1  1  1  1  1  1  1  1  1  1
/
\end{verbatim}

Alternatively the EQUALS keyword could be employed to accomplish the
same task, that is:

\begin{verbatim}
--
--       ARRAY       CONSTANT     ---------- BOX ---------
--                                I1  I2   J1  J2   K1  K2
EQUALS
         SATNUM      1            1*  1*   1*  1*   1*  1* / SET REGION 1
         SATNUM      2            1   2    1   2    1   1  / SET REGION 2
         SATNUM      3            1   2    1   2    2   2  / SET REGION 3
/
\end{verbatim}

\subsection{SKIP -- Activate Skipping of All Keywords and Input
Data}\label{skip-activate-skipping-of-all-keywords-and-input-data-5}

The SKIP keyword activates skipping of all keywords and input data until
the ENDSKIP keyword is encountered. All keywords between the SKIP and
ENDSKIP keywords are ignored.

See SKIP -- Activate Skipping of All Keywords and Input Data in the
GLOBAL section for a full description.

\subsection{SKIP100 -- Activate Skipping of Black-Oil Keywords and Input
Data}\label{skip100-activate-skipping-of-black-oil-keywords-and-input-data-4}

The SKIP100 keyword activates skipping of all keywords and input data by
the commercial black-oil simulator until the ENDSKIP keyword is
encountered. All keywords between the SKIP100 and ENDSKIP keywords are
ignored by the commercial black-oil simulator. The SKIP100 keyword is
ignored by the commercial compositional simulator.

See SKIP100 -- Activate Skipping of Keywords by Black-Oil Simulator in
the GLOBAL section for a full description.

\subsection{SKIP300 -- Activate Skipping of Keywords by Compositional
Simulator}\label{skip300-activate-skipping-of-keywords-by-compositional-simulator-5}

The SKIP300 keyword activates skipping of all keywords and input data by
the commercial compositional simulator until the ENDSKIP keyword is
encountered. All keywords between the SKIP300 and ENDSKIP keywords are
ignored by the commercial compositional simulator. The SKIP300 keyword
is ignored by the commercial black-oil simulator.

See Error: Reference source not found in the GLOBAL section for a full
description.

\subsection{SURFNUM -- Define the Surfactant Miscible Saturation Table
Region
Numbers}\label{surfnum-define-the-surfactant-miscible-saturation-table-region-numbers}

\begin{longtable}[]{@{}
  >{\raggedright\arraybackslash}p{(\linewidth - 14\tabcolsep) * \real{0.1250}}
  >{\raggedright\arraybackslash}p{(\linewidth - 14\tabcolsep) * \real{0.1250}}
  >{\raggedright\arraybackslash}p{(\linewidth - 14\tabcolsep) * \real{0.1250}}
  >{\raggedright\arraybackslash}p{(\linewidth - 14\tabcolsep) * \real{0.1250}}
  >{\raggedright\arraybackslash}p{(\linewidth - 14\tabcolsep) * \real{0.1250}}
  >{\raggedright\arraybackslash}p{(\linewidth - 14\tabcolsep) * \real{0.1250}}
  >{\raggedright\arraybackslash}p{(\linewidth - 14\tabcolsep) * \real{0.1250}}
  >{\raggedright\arraybackslash}p{(\linewidth - 14\tabcolsep) * \real{0.1250}}@{}}
\toprule\noalign{}
\begin{minipage}[b]{\linewidth}\raggedright
RUNSPEC
\end{minipage} & \begin{minipage}[b]{\linewidth}\raggedright
GRID
\end{minipage} & \begin{minipage}[b]{\linewidth}\raggedright
EDIT
\end{minipage} & \begin{minipage}[b]{\linewidth}\raggedright
PROPS
\end{minipage} & \begin{minipage}[b]{\linewidth}\raggedright
REGIONS
\end{minipage} & \begin{minipage}[b]{\linewidth}\raggedright
SOLUTION
\end{minipage} & \begin{minipage}[b]{\linewidth}\raggedright
SUMMARY
\end{minipage} & \begin{minipage}[b]{\linewidth}\raggedright
SCHEDULE
\end{minipage} \\
\midrule\noalign{}
\endhead
\bottomrule\noalign{}
\endlastfoot
\end{longtable}

\subsubsection{Description}\label{description-759}

The SURFNUM keyword defines the saturation tables (relative permeability
and capillary pressure tables) region numbers for each grid block, as
such there must be one entry for each cell in the model. The region
number specifies which set of oil-water relative permeability tables
(SWFN, SOF2, SOF3, and SWOF) are used to calculate the relative
permeability and capillary pressure in a grid block. In this case the
SURFNUM allocated tables assume that oil and water are miscible, whereas
the SATNUM allocated tables are used to allocate the immiscible
saturation tables. To use this keyword the Surfactant option must have
been activated by the SURFACT keyword in the RUNSPEC section.

This keyword is not supported by OPM Flow but would change the results
if supported so the simulation will be stopped.

\begin{longtable}[]{@{}
  >{\raggedright\arraybackslash}p{(\linewidth - 6\tabcolsep) * \real{0.2500}}
  >{\raggedright\arraybackslash}p{(\linewidth - 6\tabcolsep) * \real{0.2500}}
  >{\raggedright\arraybackslash}p{(\linewidth - 6\tabcolsep) * \real{0.2500}}
  >{\raggedright\arraybackslash}p{(\linewidth - 6\tabcolsep) * \real{0.2500}}@{}}
\toprule\noalign{}
\begin{minipage}[b]{\linewidth}\raggedright
No.
\end{minipage} & \begin{minipage}[b]{\linewidth}\raggedright
Name
\end{minipage} & \begin{minipage}[b]{\linewidth}\raggedright
Description
\end{minipage} & \begin{minipage}[b]{\linewidth}\raggedright
Default
\end{minipage} \\
\midrule\noalign{}
\endhead
\bottomrule\noalign{}
\endlastfoot
1 & SURFNUM & SURFNUM defines an array of positive integers assigning a
grid cell to a particular saturation table region. The maximum number of
SURFNUM regions is set by the NTSFUN variable on the TABDIMS keyword in
the RUNSPEC section. & 1 \\
Notes: & & & \\
\end{longtable}

\emph{Table 9.22: SURFNUM Keyword Description}

\subsubsection{Example}\label{example-428}

The example below sets three SURFNUM for the model.

\begin{verbatim}
--
--       ARRAY       CONSTANT     ---------- BOX ---------
--                                I1  I2   J1  J2   K1  K2
EQUALS
         SURFNUM     1            1*  1*   1*  1*   1*  1* / SET REGION 1
         SURFNUM     2            1*  1*   1*  1*   1   1  / SET REGION 2
         SURFNUM     3            1*  1*   1*  1*   2   2  / SET REGION 3
/
\end{verbatim}

\subsection{SURFWNUM -- Define the Saturation Table Region Numbers (High
Salinity and Water
Wet)}\label{surfwnum-define-the-saturation-table-region-numbers-high-salinity-and-water-wet}

\begin{longtable}[]{@{}
  >{\raggedright\arraybackslash}p{(\linewidth - 14\tabcolsep) * \real{0.1250}}
  >{\raggedright\arraybackslash}p{(\linewidth - 14\tabcolsep) * \real{0.1250}}
  >{\raggedright\arraybackslash}p{(\linewidth - 14\tabcolsep) * \real{0.1250}}
  >{\raggedright\arraybackslash}p{(\linewidth - 14\tabcolsep) * \real{0.1250}}
  >{\raggedright\arraybackslash}p{(\linewidth - 14\tabcolsep) * \real{0.1250}}
  >{\raggedright\arraybackslash}p{(\linewidth - 14\tabcolsep) * \real{0.1250}}
  >{\raggedright\arraybackslash}p{(\linewidth - 14\tabcolsep) * \real{0.1250}}
  >{\raggedright\arraybackslash}p{(\linewidth - 14\tabcolsep) * \real{0.1250}}@{}}
\toprule\noalign{}
\begin{minipage}[b]{\linewidth}\raggedright
RUNSPEC
\end{minipage} & \begin{minipage}[b]{\linewidth}\raggedright
GRID
\end{minipage} & \begin{minipage}[b]{\linewidth}\raggedright
EDIT
\end{minipage} & \begin{minipage}[b]{\linewidth}\raggedright
PROPS
\end{minipage} & \begin{minipage}[b]{\linewidth}\raggedright
REGIONS
\end{minipage} & \begin{minipage}[b]{\linewidth}\raggedright
SOLUTION
\end{minipage} & \begin{minipage}[b]{\linewidth}\raggedright
SUMMARY
\end{minipage} & \begin{minipage}[b]{\linewidth}\raggedright
SCHEDULE
\end{minipage} \\
\midrule\noalign{}
\endhead
\bottomrule\noalign{}
\endlastfoot
\end{longtable}

\subsubsection{Description}\label{description-760}

The SURFWNUM keyword defines the saturation tables (relative
permeability and capillary pressure tables) region numbers for each grid
block, as such there must be one entry for each cell in the mode, for
when the Surfactant Wettability option has been selected. The keyword
may also be used with the Low Salt Brine option, in this case the water
wet curves are calculated as a function of the low and high water wet
salinity curves. The region number specifies which set of relative
permeability tables are used to calculate the relative permeability and
capillary pressure in a grid block. Note that the keyword is obligatory
if the SURFACTW keyword in the RUNSPEC section has been used to invoke
the Surfactant Wettability option.

This keyword is not supported by OPM Flow but would change the results
if supported so the simulation will be stopped.

\begin{longtable}[]{@{}
  >{\raggedright\arraybackslash}p{(\linewidth - 6\tabcolsep) * \real{0.2500}}
  >{\raggedright\arraybackslash}p{(\linewidth - 6\tabcolsep) * \real{0.2500}}
  >{\raggedright\arraybackslash}p{(\linewidth - 6\tabcolsep) * \real{0.2500}}
  >{\raggedright\arraybackslash}p{(\linewidth - 6\tabcolsep) * \real{0.2500}}@{}}
\toprule\noalign{}
\begin{minipage}[b]{\linewidth}\raggedright
No.
\end{minipage} & \begin{minipage}[b]{\linewidth}\raggedright
Name
\end{minipage} & \begin{minipage}[b]{\linewidth}\raggedright
Description
\end{minipage} & \begin{minipage}[b]{\linewidth}\raggedright
Default
\end{minipage} \\
\midrule\noalign{}
\endhead
\bottomrule\noalign{}
\endlastfoot
1 & SURFWNUM & SURFWNUM defines an array of positive integers assigning
a grid cell to a particular saturation table region. The maximum number
of SURFWNUM regions is set by the NTSFUN variable on the TABDIMS keyword
in the RUNSPEC section. & 1 \\
Notes: & & & \\
\end{longtable}

\emph{Table 9.23: SURFWNUM Keyword Description}

\subsubsection{Example}\label{example-429}

The example below sets three SURFWNUM regions for the model.

\begin{verbatim}
--
--       ARRAY       CONSTANT     ---------- BOX ---------
--                                I1  I2   J1  J2   K1  K2
EQUALS
         SURFWNUM    1            1*  1*   1*  1*   1*  1* / SET REGION 1
         SURFWNUM    2            1   2    1   2    1   1  / SET REGION 2
         SURFWNUM    3            1   2    1   2    2   2  / SET REGION 3
/
\end{verbatim}

\subsection{TNUM -- Define Passive Tracer Concentration
Regions}\label{tnum-define-passive-tracer-concentration-regions}

\begin{longtable}[]{@{}
  >{\raggedright\arraybackslash}p{(\linewidth - 14\tabcolsep) * \real{0.1250}}
  >{\raggedright\arraybackslash}p{(\linewidth - 14\tabcolsep) * \real{0.1250}}
  >{\raggedright\arraybackslash}p{(\linewidth - 14\tabcolsep) * \real{0.1250}}
  >{\raggedright\arraybackslash}p{(\linewidth - 14\tabcolsep) * \real{0.1250}}
  >{\raggedright\arraybackslash}p{(\linewidth - 14\tabcolsep) * \real{0.1250}}
  >{\raggedright\arraybackslash}p{(\linewidth - 14\tabcolsep) * \real{0.1250}}
  >{\raggedright\arraybackslash}p{(\linewidth - 14\tabcolsep) * \real{0.1250}}
  >{\raggedright\arraybackslash}p{(\linewidth - 14\tabcolsep) * \real{0.1250}}@{}}
\toprule\noalign{}
\begin{minipage}[b]{\linewidth}\raggedright
RUNSPEC
\end{minipage} & \begin{minipage}[b]{\linewidth}\raggedright
GRID
\end{minipage} & \begin{minipage}[b]{\linewidth}\raggedright
EDIT
\end{minipage} & \begin{minipage}[b]{\linewidth}\raggedright
PROPS
\end{minipage} & \begin{minipage}[b]{\linewidth}\raggedright
REGIONS
\end{minipage} & \begin{minipage}[b]{\linewidth}\raggedright
SOLUTION
\end{minipage} & \begin{minipage}[b]{\linewidth}\raggedright
SUMMARY
\end{minipage} & \begin{minipage}[b]{\linewidth}\raggedright
SCHEDULE
\end{minipage} \\
\midrule\noalign{}
\endhead
\bottomrule\noalign{}
\endlastfoot
\end{longtable}

\subsubsection{Description}\label{description-761}

The TNUM keyword defines the regions associated with the series of
tracers associated with a phase (oil, water, or gas) in the model. The
maximum number of tracers for each phase are declared on the TRACER
keyword in the RUNSPEC section. Unlike other keywords, the TNUM keyword
must be concatenated with the phase and the name of the tracer declared
by TRACER keyword in the PROPS section. Table 9.24 outlines the format
of the TNUM keyword name.

This keyword is not supported by OPM Flow but would change the results
if supported so the simulation will be stopped.

\begin{longtable}[]{@{}
  >{\raggedright\arraybackslash}p{(\linewidth - 6\tabcolsep) * \real{0.2500}}
  >{\raggedright\arraybackslash}p{(\linewidth - 6\tabcolsep) * \real{0.2500}}
  >{\raggedright\arraybackslash}p{(\linewidth - 6\tabcolsep) * \real{0.2500}}
  >{\raggedright\arraybackslash}p{(\linewidth - 6\tabcolsep) * \real{0.2500}}@{}}
\toprule\noalign{}
\begin{minipage}[b]{\linewidth}\raggedright
No.
\end{minipage} & \begin{minipage}[b]{\linewidth}\raggedright
Name
\end{minipage} & \begin{minipage}[b]{\linewidth}\raggedright
Description
\end{minipage} & \begin{minipage}[b]{\linewidth}\raggedright
Default
\end{minipage} \\
\midrule\noalign{}
\endhead
\bottomrule\noalign{}
\endlastfoot
1 & TNUM & A four letter character equal to TNUM that is the root
keyword name for this data set array. & None \\
2 & PHASE & A one letter character string that must be equal to F or S,
that is concatenated to TNUM. The letter F states that the tracer is for
the ``free'' phase, for example oil or water, as well as gas cap gas
(free gas). The letter S signifies that the tracer is a ``solution''
phase tracer, for example gas dissolved in oil (as activated by the
DISGAS keyword in the RUNSPEC section), or condensate (vaporized oil) in
the gas (as per the VAPOIL keyword in the RUNSPEC section). Note tracers
that are defined by the letter S to be in the ``solution'' phase, must
also be initialized by the ``free'' phase as well. & None \\
3 & NAME & A three letter character string defining the tracer's name,
which is concatenate to TNUM and PHASE to given the full name of the
keyword Note it is best to void names beginning with the letters F, S,
and T as these names may great naming issues in post-processing
software. & None \\
\end{longtable}

\emph{Table 9.24: TNUM Keyword Name Format}

Following the declaration of the full keyword name, TNUMPHASENAME, the
keyword is followed by the data as outlined below.

\begin{longtable}[]{@{}
  >{\raggedright\arraybackslash}p{(\linewidth - 6\tabcolsep) * \real{0.2500}}
  >{\raggedright\arraybackslash}p{(\linewidth - 6\tabcolsep) * \real{0.2500}}
  >{\raggedright\arraybackslash}p{(\linewidth - 6\tabcolsep) * \real{0.2500}}
  >{\raggedright\arraybackslash}p{(\linewidth - 6\tabcolsep) * \real{0.2500}}@{}}
\toprule\noalign{}
\begin{minipage}[b]{\linewidth}\raggedright
No.
\end{minipage} & \begin{minipage}[b]{\linewidth}\raggedright
Name
\end{minipage} & \begin{minipage}[b]{\linewidth}\raggedright
Description
\end{minipage} & \begin{minipage}[b]{\linewidth}\raggedright
Default
\end{minipage} \\
\midrule\noalign{}
\endhead
\bottomrule\noalign{}
\endlastfoot
1 & TNUMREG & TUNREG defines an array of positive integers assigning a
grid cell to a particular tracer table region. The maximum number of
TNUMREG regions is set by the NTTRVD variable on the EQLDIMS keyword in
the RUNSPEC section. & 1 \\
Notes: & & & \\
\end{longtable}

\emph{Table 9.25: TNUM Keyword Data Description}

See also the TRACER keyword in the PROPS section and the TBLK keyword in
the SOLUTION section.

\subsubsection{Example}\label{example-430}

First define four passive tracers one for a free gas, one for dissolved
gas, one for oil and one to track the water.

\begin{verbatim}
--
--       DEFINE TRACER NAMES
--
--       TRACER   TRACER
--       NAME     PHASE
--       ------   ------
TRACER
        'GCG'     'GAS'                                    / GAS CAP GAS
        'DGS'     'GAS'                                    / DISOLVED GAS
        'OIL'     'OIL'                                    / OIL
        'WAT'     'WAT'                                    / WAT
/
\end{verbatim}

Given a 100 x 100 x 5 grid with DISGAS activated in the RUNSPEC section,
then the following TNUM keywords define the various tracer regions given
that NTTRVD equals four on the EQLDIMS keyword in the RUNSPEC section.

\begin{verbatim}
--
--       DEFINE PASSIVE TRACER CONCENTRATION REGIONS
--
TNUMFGCG
         1000*1
         1000*2
         1000*2
         1000*2
         1000*2
/
TNUMSDGS
         1000*1
         1000*1
         1000*1
         1000*1
         1000*1
/
TNUMFOIL
         1000*3
         1000*3
         1000*3
         1000*3
         1000*3
/
TNUMFWAT
         1000*4
         1000*4
         1000*4
         1000*4
         1000*4
/
\end{verbatim}

The keyword name is derived from the TNUM keyword, plus either F or S,
plus the tracer name declared in the TRACER keyword. For example for the
gas cap (free gas) this would be TNUM+F+GAS to give the TNUMFGAS
keyword. And for the dissolved (solution) gas this would be TNUM+S+DGS
resulting in the TNUMSDGS keyword.

\subsection{TRKPF -- Define Partitioned Tracer
Regions}\label{trkpf-define-partitioned-tracer-regions}

\begin{longtable}[]{@{}
  >{\raggedright\arraybackslash}p{(\linewidth - 14\tabcolsep) * \real{0.1250}}
  >{\raggedright\arraybackslash}p{(\linewidth - 14\tabcolsep) * \real{0.1250}}
  >{\raggedright\arraybackslash}p{(\linewidth - 14\tabcolsep) * \real{0.1250}}
  >{\raggedright\arraybackslash}p{(\linewidth - 14\tabcolsep) * \real{0.1250}}
  >{\raggedright\arraybackslash}p{(\linewidth - 14\tabcolsep) * \real{0.1250}}
  >{\raggedright\arraybackslash}p{(\linewidth - 14\tabcolsep) * \real{0.1250}}
  >{\raggedright\arraybackslash}p{(\linewidth - 14\tabcolsep) * \real{0.1250}}
  >{\raggedright\arraybackslash}p{(\linewidth - 14\tabcolsep) * \real{0.1250}}@{}}
\toprule\noalign{}
\begin{minipage}[b]{\linewidth}\raggedright
RUNSPEC
\end{minipage} & \begin{minipage}[b]{\linewidth}\raggedright
GRID
\end{minipage} & \begin{minipage}[b]{\linewidth}\raggedright
EDIT
\end{minipage} & \begin{minipage}[b]{\linewidth}\raggedright
PROPS
\end{minipage} & \begin{minipage}[b]{\linewidth}\raggedright
REGIONS
\end{minipage} & \begin{minipage}[b]{\linewidth}\raggedright
SOLUTION
\end{minipage} & \begin{minipage}[b]{\linewidth}\raggedright
SUMMARY
\end{minipage} & \begin{minipage}[b]{\linewidth}\raggedright
SCHEDULE
\end{minipage} \\
\midrule\noalign{}
\endhead
\bottomrule\noalign{}
\endlastfoot
\end{longtable}

\subsubsection{Description}\label{description-762}

The TRKPF keyword defines the regions associated with the series of
partition tracers and the partitioning tables allocated to grid blocks
in the model, for when the Partitioned Tracer option has been enabled by
the PARTTRAC keyword in the RUNSPEC section. The maximum number of
tracers for each phase are declared on the TRACERS keyword in the
RUNSPEC section. Unlike other keywords, the TRKPF keyword must be
concatenated with the name of the tracer declared by TRACER keyword in
the PROPS section. Table 9.26 outlines the format of the TRKPF keyword
name.

This keyword is not supported by OPM Flow but would change the results
if supported so the simulation will be stopped.

\begin{longtable}[]{@{}
  >{\raggedright\arraybackslash}p{(\linewidth - 6\tabcolsep) * \real{0.2500}}
  >{\raggedright\arraybackslash}p{(\linewidth - 6\tabcolsep) * \real{0.2500}}
  >{\raggedright\arraybackslash}p{(\linewidth - 6\tabcolsep) * \real{0.2500}}
  >{\raggedright\arraybackslash}p{(\linewidth - 6\tabcolsep) * \real{0.2500}}@{}}
\toprule\noalign{}
\begin{minipage}[b]{\linewidth}\raggedright
No.
\end{minipage} & \begin{minipage}[b]{\linewidth}\raggedright
Name
\end{minipage} & \begin{minipage}[b]{\linewidth}\raggedright
Description
\end{minipage} & \begin{minipage}[b]{\linewidth}\raggedright
Default
\end{minipage} \\
\midrule\noalign{}
\endhead
\bottomrule\noalign{}
\endlastfoot
1 & TRKPF & A five letter character string equal to TRKPF that is the
root keyword name for this data set array. & None \\
2 & NAME & A three letter character string defining the tracer's name,
as declared by the TRACER keyword, which is concatenate to TRKPF to
given the full name of the keyword Note it is best to void names
beginning with the letters F, S, and T as these names may great naming
issues in post-processing software. & None \\
\end{longtable}

\emph{Table 9.26: TRKPF Keyword Name Format}

Following the declaration of the full keyword name, TRKPFNAME, the
keyword is followed by the data as outlined below.

\begin{longtable}[]{@{}
  >{\raggedright\arraybackslash}p{(\linewidth - 6\tabcolsep) * \real{0.2500}}
  >{\raggedright\arraybackslash}p{(\linewidth - 6\tabcolsep) * \real{0.2500}}
  >{\raggedright\arraybackslash}p{(\linewidth - 6\tabcolsep) * \real{0.2500}}
  >{\raggedright\arraybackslash}p{(\linewidth - 6\tabcolsep) * \real{0.2500}}@{}}
\toprule\noalign{}
\begin{minipage}[b]{\linewidth}\raggedright
No.
\end{minipage} & \begin{minipage}[b]{\linewidth}\raggedright
Name
\end{minipage} & \begin{minipage}[b]{\linewidth}\raggedright
Description
\end{minipage} & \begin{minipage}[b]{\linewidth}\raggedright
Default
\end{minipage} \\
\midrule\noalign{}
\endhead
\bottomrule\noalign{}
\endlastfoot
1 & TRKPFREG & TRKPFREG defines an array of positive integers assigning
a grid cell to a particular tracer table region. The maximum number of
TRKPFREG regions is set by the NTTRVD variable on the EQLDIMS keyword in
the RUNSPEC section. & 1 \\
Notes: & & & \\
\end{longtable}

\emph{Table 9.27: TRKPF Keyword Data Description}

See also the TRACER and TRACERKP keywords in the PROPS section and the
TBLK keyword in the SOLUTION section.

\subsubsection{Example}\label{example-431}

First define one mult-partitioned tracer for the water phase.

\begin{verbatim}
--
--       DEFINE TRACER NAMES
--
--       TRACER   TRACER   TRACER  PARTITION  NUM   ADSOR
--       NAME     PHASE    VOLUME  PHASE      K(P)  PHASE
--       ------   ------   ------  ---------  ----  -----
TRACER
        'WAT'     'WAT'    1*      MULT       2     ALL    / WAT
/
\end{verbatim}

Then for a given a 100 x 100 x 5 grid assign the partitioned tracer
regions and K(P) tables, based on two regions.

\begin{verbatim}
--
--       DEFINE PARTITIONED TRACER REGIONS
--
TRKPFWAT
         1000*1
         1000*1
         1000*2
         1000*2
         1000*2
/
\end{verbatim}

The keyword name is derived from the TRKPF keyword, plus the tracer name
declared in the TRACER keyword, in this case the keyword name is
TRKPFWAT.

\subsection{WARN -- Activate Warning
Messages}\label{warn-activate-warning-messages-5}

Turns on warning messages to be printed to the print file (*.PRT); note
that this keyword is activated by default and can subsequently be
switched off by the NOWARN activation keyword. The warning messages may
be turned on and off using keywords WARN and NOWARN. OPM Flow always
prints error messages.

See WARN -- Activate Warning Messages in the GLOBAL section for a full
description.

\subsection{WH2NUM -- Define WAG Hysteresis Saturation Table Region
Numbers (Two
Phase)}\label{wh2num-define-wag-hysteresis-saturation-table-region-numbers-two-phase}

\begin{longtable}[]{@{}
  >{\raggedright\arraybackslash}p{(\linewidth - 14\tabcolsep) * \real{0.1250}}
  >{\raggedright\arraybackslash}p{(\linewidth - 14\tabcolsep) * \real{0.1250}}
  >{\raggedright\arraybackslash}p{(\linewidth - 14\tabcolsep) * \real{0.1250}}
  >{\raggedright\arraybackslash}p{(\linewidth - 14\tabcolsep) * \real{0.1250}}
  >{\raggedright\arraybackslash}p{(\linewidth - 14\tabcolsep) * \real{0.1250}}
  >{\raggedright\arraybackslash}p{(\linewidth - 14\tabcolsep) * \real{0.1250}}
  >{\raggedright\arraybackslash}p{(\linewidth - 14\tabcolsep) * \real{0.1250}}
  >{\raggedright\arraybackslash}p{(\linewidth - 14\tabcolsep) * \real{0.1250}}@{}}
\toprule\noalign{}
\begin{minipage}[b]{\linewidth}\raggedright
RUNSPEC
\end{minipage} & \begin{minipage}[b]{\linewidth}\raggedright
GRID
\end{minipage} & \begin{minipage}[b]{\linewidth}\raggedright
EDIT
\end{minipage} & \begin{minipage}[b]{\linewidth}\raggedright
PROPS
\end{minipage} & \begin{minipage}[b]{\linewidth}\raggedright
REGIONS
\end{minipage} & \begin{minipage}[b]{\linewidth}\raggedright
SOLUTION
\end{minipage} & \begin{minipage}[b]{\linewidth}\raggedright
SUMMARY
\end{minipage} & \begin{minipage}[b]{\linewidth}\raggedright
SCHEDULE
\end{minipage} \\
\midrule\noalign{}
\endhead
\bottomrule\noalign{}
\endlastfoot
\end{longtable}

\subsubsection{Description}\label{description-763}

The WH2NUM keyword defines the two phase Water-Alternating-Gas (``WAG'')
hysteresis tables (relative permeability and capillary pressure tables)
region numbers for each grid block, for when the hysteresis option has
been activated by the WAGHYSTR variable on the SATOPTS keyword in the
RUNSPEC section. The region number specifies which set of relative
permeability tables (SGFN, SWFN, SOF2, SOF3, SOF32D, SGOF, SLGOF and
SWOF) are used to calculate the relative permeability and capillary
pressure in a grid block. Note that this keyword if the two phase water
relative permeabilities WAG option.

This keyword is not supported by OPM Flow but would change the results
if supported so the simulation will be stopped.

\begin{longtable}[]{@{}
  >{\raggedright\arraybackslash}p{(\linewidth - 6\tabcolsep) * \real{0.2500}}
  >{\raggedright\arraybackslash}p{(\linewidth - 6\tabcolsep) * \real{0.2500}}
  >{\raggedright\arraybackslash}p{(\linewidth - 6\tabcolsep) * \real{0.2500}}
  >{\raggedright\arraybackslash}p{(\linewidth - 6\tabcolsep) * \real{0.2500}}@{}}
\toprule\noalign{}
\begin{minipage}[b]{\linewidth}\raggedright
No.
\end{minipage} & \begin{minipage}[b]{\linewidth}\raggedright
Name
\end{minipage} & \begin{minipage}[b]{\linewidth}\raggedright
Description
\end{minipage} & \begin{minipage}[b]{\linewidth}\raggedright
Default
\end{minipage} \\
\midrule\noalign{}
\endhead
\bottomrule\noalign{}
\endlastfoot
1 & WH2NUM & WH2NUM defines an array of positive integers assigning a
grid cell to a particular saturation table region. The maximum number of
WH2NUM regions is set by the NTSFUN variable on the TABDIMS keyword in
the RUNSPEC section. & Taken from cell allocated SATNUM \\
Notes: & & & \\
\end{longtable}

\emph{Table 9.28: WH2NUM Keyword Description}

\subsubsection{Example}\label{example-432}

The example below sets three WH2NUM regions for a model.

\begin{verbatim}
--
--       ARRAY       CONSTANT     ---------- BOX ---------
--                                I1  I2   J1  J2   K1  K2
EQUALS
         WH2NUM      1            1*  1*   1*  1*   1*  1* / SET REGION 1
         WH2NUM      2            1   2    1   2    1   1  / SET REGION 2
         WH2NUM      3            1   2    1   2    2   2  / SET REGION 3
/
\end{verbatim}

\subsection{WH3NUM -- Define WAG Hysteresis Saturation Table Region
Numbers (Three
Phase)}\label{wh3num-define-wag-hysteresis-saturation-table-region-numbers-three-phase}

\begin{longtable}[]{@{}
  >{\raggedright\arraybackslash}p{(\linewidth - 14\tabcolsep) * \real{0.1250}}
  >{\raggedright\arraybackslash}p{(\linewidth - 14\tabcolsep) * \real{0.1250}}
  >{\raggedright\arraybackslash}p{(\linewidth - 14\tabcolsep) * \real{0.1250}}
  >{\raggedright\arraybackslash}p{(\linewidth - 14\tabcolsep) * \real{0.1250}}
  >{\raggedright\arraybackslash}p{(\linewidth - 14\tabcolsep) * \real{0.1250}}
  >{\raggedright\arraybackslash}p{(\linewidth - 14\tabcolsep) * \real{0.1250}}
  >{\raggedright\arraybackslash}p{(\linewidth - 14\tabcolsep) * \real{0.1250}}
  >{\raggedright\arraybackslash}p{(\linewidth - 14\tabcolsep) * \real{0.1250}}@{}}
\toprule\noalign{}
\begin{minipage}[b]{\linewidth}\raggedright
RUNSPEC
\end{minipage} & \begin{minipage}[b]{\linewidth}\raggedright
GRID
\end{minipage} & \begin{minipage}[b]{\linewidth}\raggedright
EDIT
\end{minipage} & \begin{minipage}[b]{\linewidth}\raggedright
PROPS
\end{minipage} & \begin{minipage}[b]{\linewidth}\raggedright
REGIONS
\end{minipage} & \begin{minipage}[b]{\linewidth}\raggedright
SOLUTION
\end{minipage} & \begin{minipage}[b]{\linewidth}\raggedright
SUMMARY
\end{minipage} & \begin{minipage}[b]{\linewidth}\raggedright
SCHEDULE
\end{minipage} \\
\midrule\noalign{}
\endhead
\bottomrule\noalign{}
\endlastfoot
\end{longtable}

\subsubsection{Description}\label{description-764}

The WH3NUM keyword defines the three phase Water-Alternating-Gas
(``WAG'') hysteresis tables (relative permeability and capillary
pressure tables) region numbers for each grid block, for when the
hysteresis option has been activated by the WAGHYSTR variable on the
SATOPTS keyword in the RUNSPEC section. The region number specifies
which set of relative permeability tables (SGFN, SWFN, SOF2, SOF3,
SOF32D, SGOF, SLGOF and SWOF) are used to calculate the relative
permeability and capillary pressure in a grid block. Note that this
keyword if the three phase water relative permeabilities WAG option.

This keyword is not supported by OPM Flow but would change the results
if supported so the simulation will be stopped.

\begin{longtable}[]{@{}
  >{\raggedright\arraybackslash}p{(\linewidth - 6\tabcolsep) * \real{0.2500}}
  >{\raggedright\arraybackslash}p{(\linewidth - 6\tabcolsep) * \real{0.2500}}
  >{\raggedright\arraybackslash}p{(\linewidth - 6\tabcolsep) * \real{0.2500}}
  >{\raggedright\arraybackslash}p{(\linewidth - 6\tabcolsep) * \real{0.2500}}@{}}
\toprule\noalign{}
\begin{minipage}[b]{\linewidth}\raggedright
No.
\end{minipage} & \begin{minipage}[b]{\linewidth}\raggedright
Name
\end{minipage} & \begin{minipage}[b]{\linewidth}\raggedright
Description
\end{minipage} & \begin{minipage}[b]{\linewidth}\raggedright
Default
\end{minipage} \\
\midrule\noalign{}
\endhead
\bottomrule\noalign{}
\endlastfoot
1 & WH3NUM & WH3NUM defines an array of positive integers assigning a
grid cell to a particular saturation table region. The maximum number of
WH3NUM regions is set by the NTSFUN variable on the TABDIMS keyword in
the RUNSPEC section. & Taken from cell allocated SATNUM \\
Notes: & & & \\
\end{longtable}

\emph{Table 9.29: WH3NUM Keyword Description}

\subsubsection{Example}\label{example-433}

The example below sets three WH3NUM regions for a model.

\begin{verbatim}
--
--       ARRAY       CONSTANT     ---------- BOX ---------
--                                I1  I2   J1  J2   K1  K2
EQUALS
         WH3NUM      1            1*  1*   1*  1*   1*  1* / SET REGION 1
         WH3NUM      2            1   2    1   2    1   1  / SET REGION 2
         WH3NUM      3            1   2    1   2    2   2  / SET REGION 3
/

\end{verbatim}

\bookmarksetup{startatroot}

\chapter{SOLUTION SECTION}\label{solution-section}

\section{Introduction}\label{introduction-8}

The SOLUTION section defines the initial conditions for the model, that
is the conditions at time zero, before any flow occurs in the reservoir.
As such, sufficient information from the previous sections is required
for the simulator to calculate the initial saturation and pressure
distributions through out the model, such that the fluids are in
gravity/hydrostatic equilibrium. This means that if the simulator is run
for a period of time without any production or injection then there
should be no saturation or pressure changes in the model. In fact this
``test'' of running the simulator for a few years without any production
or injection should always be performed to ensure that the model is
indeed in hydrostatic equilibrium.

There are basically three methods to initialize a model:

\begin{verbatim}
- Equilibrium Initialization:  This is the primary method to initialize a model, and is in fact, the only manner that should be used to set the model at equilibrium conditions, unless a restart run is being initialized.  Here, OPM Flow uses the previously entered PVT data and the fluid saturation tables (SWFN, SGFN, etc.), combined with the EQUIL keyword in the SOLUTION section to compute the capillary and fluid pressure gradients, and hence the fluid saturations and fluid densities in each cell. The EQUIL keyword basically sets the pressure at a datum depth together with fluid contacts for an equilibrium region, from which all the required grid block properties can be calculated.  Note that it is not uncommon for hydrostatic equilibrium not be be achieved initially, for example due to different PVT regions being in communication, aquifer properties incorrectly defined, etc. Which is why it is important to run the simulator for a few years without any production or injection to ensure that the model is indeed in hydrostatic equilibrium, and to rectify any issues.
- Enumeration Initialization:  This is the option where one defines the pressures for all the grid cells together with the cells initial gas and water saturations, the oil saturation in this case is back calculated by the simulator using ${S}_{o}=1.0 - {S}_{g} - {S}_{w}$. Thus, in three phase runs only two phase saturations need to be entered and for two phase runs only one saturation array needs to be provided. This type of initialization is rarely used in the industry due to the complexity of ensuring that the model is in gravity equilibrium; however, the facility has been used by some companies in the past, using in-house proprietary software, to generate suitable input to the simulator. Enumeration of course can also be used to initialize the model in non-equilibrium state, similar to the restart initialization described next.
- Restart Initialization:  This is a form of enumeration initialization but in this case the condition of gravity equilibrium is not applied; in other words the model is initialized at a point in time when after the model has produced or injected fluids. This type of initialization is very common, as it enables the run to “restart” from a previously run case.  For example, once a model has been history matched, all the prediction cases can be “restarted” from the end of the final history match run, as oppose to running the cases from the start of the history match to the end of the predication cases. In some fields, the reservoir may not be in hydrostatic equilibrium, for example, due to a pressure differential across the field caused by a re-charging aquifer.  In these rare cases the model is initialized using the hydrostatic equilibrium approach but then run until the present day conditions match the measured present day conditions by adjusting the aquifer influx. The actual history match/prediction runs are then restarted from the resulting present day restart file.
\end{verbatim}

Equilibrium Initialization uses the PVT data and saturation tables
entered in the PROPS section combined with data entered on the EQUIL
keyword in the SOLUTION section. As capillary forces are generally small
relative to viscous forces in field scale simulations, one may expect
the errors introduced by use of the drainage capillary pressure curves
to simulate an imbibition process to be small. Thus, it is common
practice in field scale simulations to use the capillary pressure
drainage curves in initializing a model, mainly as drainage curves are
routinely measured in the laboratory, whereas imbibition curve are not.
The exception of course is when hysteresis is applicable or when
simulating small scale systems. If the hysteresis option has been
activated via HYSTER variable on the SATOPTS keyword in the RUNSPEC
section, then the simulator will use the appropriate imbibition
capillary pressure data entered via the saturation table keywords (SWFN,
SGFN, etc.) in determining the fluid saturation.

The initial vertical distribution of fluids will be governed by the
fluid densities and capillary forces, with a gas zone, a gas-oil
transition zone, an oil zone, an oil-water transition zone, and finally
the water zone within a hydrostatic unit or region. Not all zones may be
present in a hydrostatic unit, for example in dry gas reservoirs there
will be only gas, gas-water transition, and water zones. This is
fundamental to construct given the density of the in situ phases.
However, there are rare exceptions to this principle. When there is an
exception, the more likely explanation is the various fluids are not in
the same hydrostatic unit.

In the methodology outline below, it has been assumed that the initial
oil pressure has been specified at some depth, zinit within the oil
zone. Note it is not necessary that the pressure is specified within the
oil zone, but it makes describing the analysis somewhat simpler. Thus we
have:

\begin{equation}\phantomsection\label{eq-10-1}{
p = {{p}^{\mathit{oil}}}_{\mathit{init}} \mathit{at} z = {z}_{\mathit{init}}
}\end{equation}

The oil pressure at any other point in the oil column can be calculated
as shown below:

\begin{equation}\phantomsection\label{eq-10-2}{
{p}_{\mathit{oil}} = {{p}^{\mathit{oil}}}_{\mathit{init}} + {γ}^{\mathit{oil}} \left( z - {z}_{\mathit{init}} \right)
}\end{equation}

where ɣoil is the oil pressure gradient and z is the given depth. We can
then calculate the oil pressure at the Water-Oil Contact (``WOC''), that
is poil woc from:

\begin{equation}\phantomsection\label{eq-10-3}{
{{p}^{\mathit{oil}}}_{\mathit{woc}} = {{p}^{\mathit{oil}}}_{\mathit{init}} + {γ}^{\mathit{oil}} \left( {z}_{\mathit{woc}} - {z}_{\mathit{init}} \right)
}\end{equation}

Employing the definition of water-oil capillary pressure, one can write
an expression for the water pressure at the water-oil contact:

\begin{equation}\phantomsection\label{eq-10-4}{
{{p}^{\mathit{water}}}_{\mathit{woc}} = {{p}^{\mathit{oil}}}_{\mathit{woc}} - {p}_{\mathit{cwoc}}
}\end{equation}

The water pressure at any other point in the reservoir can now be
related to the water pressure at the water-oil contact as shown below:

\begin{equation}\phantomsection\label{eq-10-5}{
{p}_{\mathit{water}} = {{p}^{\mathit{water}}}_{\mathit{woc}} + {γ}^{\mathit{water}} \left( z - {z}_{\mathit{woc}} \right)
}\end{equation}

where ɣwater is the water pressure gradient and z is the given depth.

Similarly, for gas we have the oil phase pressure at the Gas-Oil Contact
(``GOC'').

\begin{equation}\phantomsection\label{eq-10-6}{
{{p}^{\mathit{oil}}}_{\mathit{goc}} = {{p}^{\mathit{oil}}}_{\mathit{init}} + {γ}^{\mathit{oil}} \left( {z}_{\mathit{goc}} - {z}_{\mathit{init}} \right)
}\end{equation}

Again, we can use gas-oil capillary pressure to obtain the gas phase
pressure at the GOC, that is pgas goc :

\begin{equation}\phantomsection\label{eq-10-7}{
{{p}^{\mathit{gas}}}_{\mathit{goc}} = {{p}^{\mathit{oil}}}_{\mathit{goc}} + {p}_{\mathit{cgoc}}
}\end{equation}

and the gas phase pressure at any point in the reservoir is therefore
determined from:

\begin{equation}\phantomsection\label{eq-10-8}{
{p}_{\mathit{gas}} = {{p}^{\mathit{gas}}}_{\mathit{goc}} + {γ}^{\mathit{gas}} \left( z - {z}_{\mathit{goc}} \right)
}\end{equation}

where ɣgas is the gas pressure gradient.

Having calculated the phase pressures at all depths in the reservoir
(all grid block cells in the model), we can then calculate the capillary
pressures, as the difference between the oil and water pressures, and
the oil and gas pressures, depending on the phases present, that is:

\begin{equation}\phantomsection\label{eq-10-9}{
{p}_{\mathit{cwo}} = {p}_{\mathit{oil}} - {p}_{\mathit{water}}
}\end{equation}

\begin{equation}\phantomsection\label{eq-10-10}{
{p}_{\mathit{cgo}} = {p}_{\mathit{gas}} - {p}_{\mathit{oil}}
}\end{equation}

Under the assumption that the reservoir is initially at capillary
equilibrium, we now calculate the average water saturation for each grid
block in the computational grid on the basis of the capillary pressure
at the center of that grid block.

\begin{equation}\phantomsection\label{eq-10-11}{
{S}_{w} = {S}_{w} \left( {P}_{{\mathit{cow}}_{\mathit{mid}-\mathit{point} \mathit{of} \mathit{grid} \mathit{cell}} }\right)
}\end{equation}

The calculation of initial gas saturations is exactly analogous to the
above treatment of water saturations.

\begin{equation}\phantomsection\label{eq-10-12}{
{S}_{g} = {S}_{g} \left( {P}_{{\mathit{cog}}_{\mathit{mid}-\mathit{point} \mathit{of} \mathit{grid} \mathit{cell}} }\right)
}\end{equation}

Note also that fluid saturations in a grid block should sum to one, that
is, \({S}_{g} + {S}_{o} + {S}_{w} = 1.0\).

With grid blocks of sufficiently small size, the procedure described
above will yield accurate saturation distributions for the static
equilibrium situation. Unfortunately, in practice we are often forced to
used grid blocks that are much larger than we would like because of
computer time and storage considerations. One of the results of using
such large blocks is that the saturation distributions and hence the
oil- and gas-in-place volumes can sometimes be substantially in error.

This is easily demonstrated by imagining a case where the water-oil
contact lies infinitesimally above the center of a 20-foot thick grid
block. In this case, the procedure described above would cause the water
saturation to be calculated as 1.0, whereas in actual fact it should be
very nearly 0.5 (1.0 + Swc). Hence, we would make an error approaching
half the pore volume of the grid block in the oil and water in-place in
that block. To overcome this we can sub-divide each grid block in the
vertical direction, and calculate the saturations for the sub-blocks,
with the number of sub-blocks defined via the EQLOPT3 option on the
EQLNUM keyword. The simulator can then sum the sub-blocks to obtain the
average grid block saturation and thus obtain a more accurate estimate
of the volumes in-place. Note this is less of a concern these days due
to fact that most models now use 0.5 to 1.0 metres (one to three feet)
in the vertical resolution due to the improvement in computer
efficiency.

There are variations to the above general methodology depending on the
data entered, for example:

\begin{verbatim}
- If the SWATINIT array, that stipulates the initial water saturation for each cell, has been entered in the PROPS section then the capillary pressure data entered on the saturation tables are re-scaled to honor the SWATINIT array whenever possible, and the re-scaled capillary pressure data are used in the equilibrium initialization calculations. Thus, each cell will have its own scaling of the capillary pressure that needs to be taken into account during the simulation.
- For reservoirs with both oil and dissolved gas phases (OIL and DISGAS keywords in the RUNSPEC section), the dissolved gas ratio, Rs, or the saturation pressure Pb, may also be entered as a function of depth. See the RSVD and PBVD keywords in the SOLUTION section. If neither of these keywords have been entered then Rs is set to equal to the Rs value at the GOC. This also means that the datum depth on the EQUIL keyword in the SOLUTION section, must equal the GOC depth, in order for the saturated Rs to be calculated from the datum pressure and the associated PVT table.
- Similarly, for reservoirs with both wet gas and vaporized oil (GAS and VAPOIL keywords in the RUNSPEC section), the dissolved vaporized oil (more commonly referred to as condensate) gas ratio, Rv, or the dew-point pressure Pd, may also be entered as a function of depth. See the RVVD and PDVD keywords in the SOLUTION section. If neither of these keywords have been entered then Rv is set to equal to the Rv value at the GOC. Again, this also means that the datum depth must set to the GOC depth in order for the saturated Rv to be calculated from the datum pressure and the associated PVT table.
\end{verbatim}

In OPM Flow {[}Rasmussen, A.F., Sandve, T.H., Bao, K., Lauser, A., Hove,
J., Skaflestad, B., Klöfkorn, R., Blatt, M., Rustad, A.B., Sævareid, O.
et al.~{[}2019{]}. The Open Porous Media Flow Reservoir Simulator.
https://doi.org/10.1016/j.camwa.2020.05.014.{]} initial conditions are
calculated in a similar manner by using the pressure at a reference
depth, the depths of the water-oil and oil-gas contacts, and the
capillary pressures at those contact depths. It then solves the
following ordinary differential equations:

\begin{equation}\phantomsection\label{eq-10-13}{
\begin{matrix}\frac{{\mathit{dp}}_{\mathit{oil}}}{\mathit{dz}} = {ρ}_{\mathit{oil}}(z,{p}_{\mathit{oil}})g \\  \\ \frac{{\mathit{dp}}_{\mathit{gas}}}{\mathit{dz}} = {ρ}_{\mathit{gas}}(z,{p}_{\mathit{gas}})g \\  \\ \frac{{\mathit{dp}}_{\mathit{wat}}}{\mathit{dz}} ={ρ}_{\mathit{wat}}(z,{p}_{\mathit{wat}})g\end{matrix}
}\end{equation}

where

ρi = the oil, gas, water density,

z = the depth,

pi = the oil, gas, water pressure, and

g = the gravity term.

The equations are solved numerically using a 4th-order Runge-Kutta
{[}Runge, C. D. T. (1895), ``Über die numerische Auflösung von
Differentialgleichungen'', Mathematische Annalen, Springer, 46 (2):
167--178, doi:10.1007/BF01446807.{]} and {[}Kutta, M. (1901), ``Beitrag
zur näherungsweisen Integration totaler Differentialgleichungen'',
Zeitschrift für Mathematik und Physik, 46: 435--453.{]} method, with the
order of the phases decided by the location of the datum depth. First,
the simulator solves for the phase in whose zone the datum is located,
using the datum pressure to bind the pressure solution. Then the
simulator solves for phases corresponding to the zones above and below.
To tie the phase pressures, the simulator uses the already computed
phase pressure(s) evaluated at the zone contact depth and the input
capillary pressures.

If the initial water saturations have been specified using the SWATINIT
keyword then the same approach is used for the water saturation but the
capillary pressure is scaled to match the specified initial water
saturation. Each cell will thus have its own scaling of the capillary
pressure that needs to be taken into account during the simulations.
Initial fluid properties (GOR, formation volume factors, etc.) are
calculated from the cell equilibrium pressure and the entered fluid
property tables, as described earlier.

\section{Keyword Definitions}\label{keyword-definitions-13}

\section{Data Requirements}\label{data-requirements-6}

\subsection{Equilibrium
Initialization}\label{equilibrium-initialization}

For the standard Equilibrium Initialization the EQUIL keyword in the
SOLUTION section is used to initialize the model. The keyword contains
the datum depth and pressure, oil-water and gas-oil contacts, the
capillary pressure values at the fluid contacts, as well as various
equilibrium initialization options, including how the dissolve gas and
vaporized oil should be calculated, and if a refined in-place volume
calculation should be performed. There should be one record for each
equilibration region with the number of regions defined by the NTEQUL
parameter on the EQLDIMS keyword in the RUNSPEC section.

Note that if a fluid contact is not present in the reservoir, for
example in two phase runs containing only gas and water, the water
oil-water contact on the EQUIL keyword should be set to the gas-water
contact. Also, in cases where there is no gas cap (or free gas) then
gas-oil should be set to a value shallower than the top of the
reservoir, and for three phases cases where there is initially no oil
zone, as for a gas condensate field for example, the gas-oil contact
should be set to the same depth as water-oil contact. Table 10.1
outlines the minimum set of keywords required by OPM Flow in order for
the simulator to initialize the model For the standard Equilibrium
Initialization option.

\begin{longtable}[]{@{}
  >{\raggedright\arraybackslash}p{(\linewidth - 6\tabcolsep) * \real{0.2500}}
  >{\raggedright\arraybackslash}p{(\linewidth - 6\tabcolsep) * \real{0.2500}}
  >{\raggedright\arraybackslash}p{(\linewidth - 6\tabcolsep) * \real{0.2500}}
  >{\raggedright\arraybackslash}p{(\linewidth - 6\tabcolsep) * \real{0.2500}}@{}}
\toprule\noalign{}
\begin{minipage}[b]{\linewidth}\raggedright
No.
\end{minipage} & \begin{minipage}[b]{\linewidth}\raggedright
SOLUTION Keyword
\end{minipage} & \begin{minipage}[b]{\linewidth}\raggedright
Equilibrium Initialization Description
\end{minipage} & \begin{minipage}[b]{\linewidth}\raggedright
Type of Initialization
\end{minipage} \\
\midrule\noalign{}
\endhead
\bottomrule\noalign{}
\endlastfoot
1 & EQUIL & EQUIL -- Define the Equilibration Initialization Data - this
is the main keyword for the Equilibration initialization that sets the
datum depth, the pressure and fluid contacts for an equilibrium region.
& Equilibration \\
2 & PBVD & PBVD -- Equilibration Bubble-Point versus Depth Tables - used
for Live Black-Oil equilibration regions to define the saturation
pressure versus depth relationship. & Equilibration \\
3 & PDVD & PDVD -- Define Equilibration Dew-Point versus Depth Tables -
used for Wet Gas equilibration regions to define the saturation pressure
versus depth relationship. & Equilibration \\
4 & RSVD & RSVD -- Equilibration Dissolved Gas-Oil Ratio (Rs) versus
Depth Tables - used for Live Black-Oil equilibration regions to define
the gas-oil ratio versus depth relationship. & Equilibration \\
5 & RVVD & RVVD -- Equilibration Vaporized Oil-Gas Ratio (Rv) versus
Depth Tables - use for Wet Gas equilibration regions to declare the
condensate-gas ratio versus depth relationship. & Equilibration \\
6 & SALTVD & SALTVD -- Equilibration Salt Concentration versus Depth
Tables - this is a standard commercial simulator keyword used with the
Brine model, but may also be used with OPM Flow's Vaporized Water and
Salt Precipitation models. & Equilibration \\
7 & THPRES & THPRES - Define Equilibration Region Threshold Pressures &
Equilibration Enumeration \\
& & Vaporized Water Model (OPM Flow Specific Keywords) & \\
1 & RVWVD & RVWVD -- Equilibration Vaporized Water-Gas Ratio (Rvw)
versus Depth Tables & Equilibration \\
2 & SALTPVD & SALTPVD -- Initial Precipitated Salt Volume Fraction
versus Depth Tables & Equilibration \\
3 & SALTVD & SALTVD -- Equilibration Salt Concentration versus Depth
Tables - this is a standard commercial simulator keyword used with the
Brine model, but may also be used with OPM Flow's Vaporized Water and
Salt Precipitation models. & Equilibration \\
Notes: & & & \\
\end{longtable}

\emph{Table 10.1: Minimum Set of SOLUTION Equilibrium Keywords Required
by OPM Flow}

The minimum set of SOLUTION keywords is generally for three phases runs,
that is where the OIL, DISGAS, GAS, VAPOIL, and/or WATER keywords in the
RUNSPEC section have been invoked. If additional phases or options have
been invoked in the RUNSPEC section, then additional SOLUTION keywords
may be required in order to initialize the model. For example, if the
BRINE keyword in the RUNSPEC section has been used to activate the BRINE
model then the SALTVD keyword in the SOLUTION section may or may not be
used to set the initial salt concentration versus depth relationship.

\subsection{Enumeration
Initialization}\label{enumeration-initialization}

Enumeration Initialization is based on user supplying all the required
data for the simulator to determine the pressures, fluid saturations,
and fluid properties for each cell in the model. This is similar to
entering the grid property array data in the GRID section, in that all
the data is required for all the cells in the model.

To define the cell pressures for all the grid blocks, the PRESSURE
keyword in the SOLUTION section should be employed. There is also the
PRVD keyword in the SOLUTION section that can be used to define a
pressure versus depth relationship to enable the simulator to calculate
the cell pressures instead of entering the data via the PRESSURE
keyword; however, this keyword is not implemented in OPM Flow.

For the fluid saturations, it is only necessary to enter the saturations
sufficient to define the system, so for a two phase system it is only
necessary to define one phase, normally the water phase, as the other
phase will be calculated by the simulator. Similarly, for three phases
only two phases must be entered with the water phase always normally
entered, and the other phase dependent on the type of reservoir. The
data is entered using the SGAS, SOIL, and SWAT keywords in the SOLUTION
section.

Finally the initial fluid property data needs to be defined based on the
fluid phases active in the model. For oil reservoirs either the RS
keyword that defines a cell's initial gas-oil ratio, or the PBUB keyword
that defines the oil's initial bubble point pressure for each cell are
required. For gas reservoirs, either the RV keyword that defines
vaporized oil-gas ratio (or condensate gas-oil ratio), or the PDEW
keyword that defines the initial dew point pressure for each cell must
be entered. Table 10.2 outlines the minimum set of keywords required by
OPM Flow in order for the simulator to initialize the model. For the
Enumeration Initialization option - all the keywords are in the SOLUTION
section.

\begin{longtable}[]{@{}
  >{\raggedright\arraybackslash}p{(\linewidth - 6\tabcolsep) * \real{0.2500}}
  >{\raggedright\arraybackslash}p{(\linewidth - 6\tabcolsep) * \real{0.2500}}
  >{\raggedright\arraybackslash}p{(\linewidth - 6\tabcolsep) * \real{0.2500}}
  >{\raggedright\arraybackslash}p{(\linewidth - 6\tabcolsep) * \real{0.2500}}@{}}
\toprule\noalign{}
\begin{minipage}[b]{\linewidth}\raggedright
No.
\end{minipage} & \begin{minipage}[b]{\linewidth}\raggedright
SOLUTION Keyword
\end{minipage} & \begin{minipage}[b]{\linewidth}\raggedright
Enumeration Initialization Description
\end{minipage} & \begin{minipage}[b]{\linewidth}\raggedright
Type of Initialization
\end{minipage} \\
\midrule\noalign{}
\endhead
\bottomrule\noalign{}
\endlastfoot
1 & PRESSURE & PRESSURE -- Define the Initial Equilibration Pressures
for All Grid Blocks. This keyword must be present for all Enumeration
equilibration regions. & Enumeration \\
2 & PBUB & PBUB -- Define the Initial Equilibration Bubble-Point
Pressure for All Grid Blocks - used for Live Black-Oil equilibration
regions to define the saturation pressure for each grid cell. &
Enumeration \\
3 & PDEW & PDEW -- Define the Initial Equilibration Dew-Point Pressure
for All Grid Blocks - used for Wet Gas equilibration regions to set the
saturation pressure for each cell. & Enumeration \\
4 & RS & RS -- Define the Initial Equilibration GOR (Rs) for All Grid
Blocks - used for Live Black-Oil equilibration regions to define the
gas-oil ratio for each grid blocks. & Enumeration \\
5 & RV & RV -- Define the Initial Equilibration CGR (Rv) for All Grid
Blocks - use for Wet Gas equilibration regions to set the condensate-gas
ratio for each cell. & Enumeration \\
6 & SALT & SALT -- Define the Initial Equilibration Salt Concentration
for All Grid Blocks - used to set the initial salt concentration in the
Brine model and OPM Flow's Salt Precipitation model. & Enumeration \\
7 & SGAS & SGAS -- Define the Initial Equilibration Gas Saturation for
All Grid Blocks - if gas is present in the model then the keyword is
used to set the initial gas saturation. & Enumeration \\
8 & SOIL & SOIL -- Define the Initial Equilibration Oil Saturation for
All Grid Blocks - if oil is present in the model then the keyword is
used to set the initial oil saturation. & Enumeration \\
9 & SWAT & SWAT -- Define the Initial Equilibration Water Saturation for
All Grid Blocks - since water is normally always present in petroleum
reservoirs, this keyword is always used to define the initial water
saturation. For two-phase runs, oil-water or gas-water, then it is only
necessary to use this keyword, as the second phase saturations are
automatically calculated by the simulator. & Enumeration \\
10 & THPRES & THPRES - Define Equilibration Region Threshold Pressures &
Equilibration Enumeration \\
& & Microbial Induced Calcite Precipitation And Biofilm Model
Enumeration (OPM Flow Specific Keywords) & \\
1 & SBIOF & SBIOF - Define The Initial Equilibration Biofilm Volume
Fraction For All Grid Blocks & Enumeration \\
2 & SCALC & SCALC -- Define The Initial Equilibration Calcite Volume
Fraction For All Grid Blocks & Enumeration \\
3 & SMICR & SMICR -- Define The Initial Equilibration Microbial
Concentration For All Grid Blocks & Enumeration \\
4 & SOXYG & SOXYG - Define The Initial Equilibration Oxygen
Concentration For All Grid Blocks & Enumeration \\
5 & SUREA & SUREA - Define The Initial Equilibration Urea Concentration
For All Grid Blocks & Enumeration \\
& & Vaporized Water Model (OPM Flow Specific Keywords) & \\
1 & RVW & RVW -- Define the Initial Equilibration Vaporized Water in Gas
Ratio for All Grid Blocks - used in both the Vaporized Water model and
the Salt Precipitation model. & Enumeration \\
2 & SALT & SALT -- Define the Initial Equilibration Salt Concentration
for All Grid Blocks - this is a standard commercial simulator keyword
used with the Brine model, but may also be used with OPM Flow's
Vaporized Water and Salt Precipitation models. & Enumeration \\
3 & SALTP & SALTP -- Define the Initial Precipitated Salt Volume
Fraction for All Grid Blocks. & Enumeration \\
& & Dissolved Gas in Water Model (OPM Flow Specific Keywords) & \\
1 & \hyperref[REF_HEADING_KEYWORD_RSW_10_3]{RSW} &
\hyperref[0.0.0.RSW]{RSW -- Define the Initial Equilibrium Solution Gas
in Water Ratio for All Grid Blocks} - this is a commercial compositional
simulator keyword, but may also be used with OPM Flow's Dissolved Gas in
Water model. & Enumeration \\
\end{longtable}

\emph{Table 10.2: Minimum Set of SOLUTION Enumeration Keywords Required
by OPM Flow}

Again, the minimum set of SOLUTION keywords is for general three phases
runs, that is where the OIL, DISGAS, GAS, VAPOIL, and/or WATER keywords
in the RUNSPEC section have been invoked. If additional phases or
options have been invoked, then additional SOLUTION keywords may be
required in order to initialize the model.

\subsection{Restart Initialization}\label{restart-initialization}

Restart Initialization is performed by the simulator by loading a
previously generated RESTART file and equilibrating the model
accordingly. In nearly all cases the model will not satisfy the
condition of gravity equilibrium; in other words the model is
initialized at a point in time after the model has produced or injected
fluids.

To write a restart record to the RESTART file, either the RPTSOL or
RPTRST keyword in the SOLUTION section can be used to write the initial
restart record at time equal to zero. To write restart records at
various time points as the simulation progresses through time, then
either the RPTRST or RPTSCHED keywords may be used. Both the RPTSOL and
RPTSCHED keywords enable just the basic restart record to be written to
the RESTART file, which is generally sufficient in most cases. In
comparison with the dedicated RPTRST keyword, RPTRST provides additional
functionality on the data to be written to the RESTART file and is
therefore the preferred keyword to use. For example to request the
standard restart data be written out every month using the RPTRST
keyword, one would use the following:

\begin{verbatim}
--
--       RESTART CONTROL BASIC = 4 (YEARLY) 5 (MONTHLY)
--
RPTRST
         BASIC=5                                                               /
\end{verbatim}

Subsequently one of the monthly restart points can be used to restart
the run.

The keywords used to create restart points in various sections of the
input deck and to initialize the model from a RESTART file are shown in
Table 10.3.

\begin{longtable}[]{@{}
  >{\raggedright\arraybackslash}p{(\linewidth - 6\tabcolsep) * \real{0.2500}}
  >{\raggedright\arraybackslash}p{(\linewidth - 6\tabcolsep) * \real{0.2500}}
  >{\raggedright\arraybackslash}p{(\linewidth - 6\tabcolsep) * \real{0.2500}}
  >{\raggedright\arraybackslash}p{(\linewidth - 6\tabcolsep) * \real{0.2500}}@{}}
\toprule\noalign{}
\begin{minipage}[b]{\linewidth}\raggedright
No.
\end{minipage} & \begin{minipage}[b]{\linewidth}\raggedright
Keyword
\end{minipage} & \begin{minipage}[b]{\linewidth}\raggedright
Restart Creation Description
\end{minipage} & \begin{minipage}[b]{\linewidth}\raggedright
Keyword Section
\end{minipage} \\
\midrule\noalign{}
\endhead
\bottomrule\noalign{}
\endlastfoot
1 & RPTRST & RPTRST -- Define Data to be Written to the RESTART File &
SOLUTION \\
2 & RPTSOL & RPTSOL -- Define SOLUTION Section Reporting & \\
3 & RPTRST & RPTRST -- Define Data to be Written to the RESTART File &
SCHEDULE \\
4 & RPTSCHED & RPTSCHED -- Define SCHEDULE Section Reporting & \\
No. & Keyword & Restart Initialization Description & Keyword Section \\
5 & RESTART & RESTART -- Restart Run From an Existing Restart File -
this keyword is used to initialize the model from a previously run OPM
Flow run or from the commercial simulator. Note that due to the
complexities of the RESTART file, OPM Flow may not always be able to
restart from the commercial simulators RESTART file. & SOLUTION \\
Notes: & & & \\
\end{longtable}

\emph{Table 10.3: SOLUTION Restart Initialization Keywords Required by
OPM Flow}

In the input deck that is being employed to restart the run, all the
equilibration keywords (EQUIL, RSVD, etc.) or the enumeration
equilibration keywords (PRESSURE, SGAS, SOIL, SWAT, etc.) in the
SOLUTION section used to initialize the model should be deleted.
Secondly, it is good practice to insert the SKIPREST keyword at the very
beginning of the SCHEDULE section. The SKIPREST keyword causes the
simulator to only read in data it requires for restarting the run up to
the RESTART point. This is because certain keywords always need to be
present in a restart run in the SCHEDULE section, as the data is not
stored on the RESTART file, for example the VFP tables (VFPPROD and
VFPINJ keywords in the SCHEDULE section). The SKIPREST keyword
automatically processes the input deck and reads the required data.

\subsection{Section Keywords}\label{section-keywords}

A complete list of SOLUTION keywords in alphabetic order is shown in
Table 10.4 together with a generalized Topic column that classifies the
functionality of the keyword. Note that not all keywords and features
listed in Table 10.4 are implemented in OPM Flow. Cells not colored in
the No.~column indicate the keyword is supported by OPM Flow, cells
colored gray indicate that the keyword is not applicable, and finally,
cells colored in orange indicate keywords that are not currently
supported by OPM Flow.

\begin{longtable}[]{@{}
  >{\raggedright\arraybackslash}p{(\linewidth - 6\tabcolsep) * \real{0.2500}}
  >{\raggedright\arraybackslash}p{(\linewidth - 6\tabcolsep) * \real{0.2500}}
  >{\raggedright\arraybackslash}p{(\linewidth - 6\tabcolsep) * \real{0.2500}}
  >{\raggedright\arraybackslash}p{(\linewidth - 6\tabcolsep) * \real{0.2500}}@{}}
\toprule\noalign{}
\begin{minipage}[b]{\linewidth}\raggedright
No.
\end{minipage} & \begin{minipage}[b]{\linewidth}\raggedright
RUNSPEC Keyword
\end{minipage} & \begin{minipage}[b]{\linewidth}\raggedright
Description
\end{minipage} & \begin{minipage}[b]{\linewidth}\raggedright
Topic
\end{minipage} \\
\midrule\noalign{}
\endhead
\bottomrule\noalign{}
\endlastfoot
1 & ADD & ADD -- Add a Constant to a Specified Array. & Input \\
2 & ADDREG & ADDREG -- Add a Constant to an Array based on a Region
Number. & Input \\
3 & APIVD & APIVD - Equilibration Oil API Gravity versus Depth Tables. &
API \\
4 & AQANCONL & Error: Reference source not found. & Aquifer \\
5 & AQANNC & AQANNC -- Define Analytic Aquifer Non-Neighbor Connections.
& Aquifer \\
6 & AQANTRC & AQANTRC - Define Analytic Aquifer Initial Tracer
Concentrations. & Aquifer \\
7 & AQUALIST & AQUALIST -- Define An Analytic Aquifer Name to Aquifer
Numbers . & Aquifer \\
8 & AQUANCON & AQUANCON -- Define Analytical Connections to the Grid. &
Aquifer \\
9 & AQUCHGAS & AQUCHGAS -- Define Constant Pressure Gas Analytical
Aquifer Properties. & Aquifer \\
10 & AQUCHWAT & AQUCHWAT -- Define Constant Pressure Water Analytical
Aquifer Properties. & Aquifer \\
11 & AQUCON & AQUCON -- Define Numerical Aquifer Connections to the
Grid. & Aquifer \\
12 & AQUCT & AQUCT -- Define Carter-Tracy Analytical Aquifers. &
Aquifer \\
13 & AQUFET & AQUFET -- Define Fetkovich Analytical Aquifer and
Connections. & Aquifer \\
14 & AQUFETP & AQUFETP -- Define Fetkovich Analytical Aquifers. &
Aquifer \\
15 & AQUFLUX & AQUFLUX - Define Constant Flux Analytical Aquifer. &
Aquifer \\
16 & BC & BC -- Define Boundary Conditions. This is an OPM Flow specific
keyword defines the boundary conditions for the model, and can be used
to set boundary conditions for when external influx or efflux volumes
are influencing the reservoir pressure and production history. &
Equilibration \\
17 & BOUNDARY & BOUNDARY -- Define a Boundary Box for Printing. &
Output \\
18 & BOX & BOX - Define a Range of Grid Blocks to Enter Property Data. &
Input \\
19 & COLUMNS & COLUMNS -- Define Input File Column Margins. & Input \\
20 & COPY & COPY -- Copy Array Data to Another Array. & Input \\
21 & COPYREG & COPYREG -- Copy an Array to Another Array based on a
Region Number. & Input \\
22 & DATUM & DATUM -- Define the Datum Depth for the Model. & Output \\
23 & DATUMR & DATUMR -- Define Datum Depths for the FIPNUM Regions. &
Output \\
24 & DATUMRX & DATUMRX -- Define Datum Depths for the FIP Allocated
Regions. & Output \\
25 & DEBUG & DEBUG -- Define the Debug Data to be Printed to File. &
Output \\
26 & DYNAMICR & DYNAMICR -- Start of Dynamic Region Parameter
Definition. & Input \\
27 & ECHO & ECHO -- Activate Echoing of User Input Files to the Print
File. & Output \\
28 & END & END -- Define the End of the Input File. & Input \\
29 & ENDBOX & ENDBOX -- Define the End of the BOX Defined Grid. &
Input \\
30 & ENDDYN & ENDDYN-- End of Dynamic Region Parameter Definition. &
Input \\
31 & ENDFIN & ENDFIN -- End the Definition of a Local Grid Refinement. &
LGR \\
32 & ENDINC & ENDINC -- Define the End of an Include File. & Input \\
33 & ENDSKIP & ENDSKIP -- Deactivate Skipping of Keywords and Input
Data. & Input \\
34 & EQUALREG & EQUALREG -- Sets an Array to a Constant by Region
Number. & Input \\
35 & EQUALS & EQUALS -- Sets a Specified Array to a Constant. & Input \\
36 & EQUIL & EQUIL -- Define the Equilibration Initialization Data. &
Equilibration \\
37 & EXTRAPMS & EXTRAPMS -- Activate Extrapolation Warning Messages. &
Output \\
38 & FILEUNIT & FILEUNIT -- Activate Unit Consistency Checking. &
Input \\
39 & FORMFEED & FORMFEED -- Defined the Print File Form-Feed Character.
& Output \\
40 & GASCONC & GASCONC -- Define the Initial Equilibration Coal Gas
Concentration for All Grid Blocks. & CBM \\
41 & GASSATC & GASSATC -- Define the Initial Equilibration Saturated
Coal Gas Concentration for All Grid Blocks. & CBM \\
42 & GCVD & GCVD -- Define Equilibration Coal Gas Concentration versus
Depth Tables. & CBM \\
43 & GETDATA & GETDATA -- Load and Assign Data Array from INIT or
RESTART Files. & Input \\
44 & GETGLOB & GETGLOB -- Activate Loading of Global Grid Restart Data
Option. & Input \\
45 & GI & GI - Define the Initial Equilibration Gi Values for All Grid
Blocks. & Enumeration \\
46 & HMAQUCT & HMAQUCT -- History Match Carter-Tracy Aquifer Gradient
Parameters. & Gradient \\
47 & HMAQUFET & HMAQUFET -- History Match Fetkovich Aquifer Gradient
Parameters. & Gradient \\
48 & HMMLCTAQ & HMMLCTAQ -- History Match Carter-Tracy Aquifer Gradient
Multipliers. & Gradient \\
49 & HMMLFTAQ & HMMLFTAQ -- History Match Fetkovich Aquifer Gradient
Multipliers. & Gradient \\
50 & HMMLTWCN & HMMLTWCN -- History Match Well Connection and Skin
Multipliers. & Gradient \\
51 & HMWELCON & HMWELCON -- History Match Well Connection and Skin
Parameters. & Gradient \\
52 & IMPORT & IMPORT -- Import Grid File Data at the Current Position. &
Input \\
53 & INCLUDE & INCLUDE -- Load Another Data File at the Current
Position. & Input \\
54 & MESSAGE & MESSAGE -- Output User Message. The MESSAGE keyword
outputs a user message to the terminal, as well as to the print
(\emph{.PRT) and debug (}.DBG) files. & Output \\
55 & MESSAGES & MESSAGES -- Define Message Print Limits and Stop Limits.
The MESSAGES keyword defines the print and stops levels for various
messages. & Output \\
56 & MULTIPLY & MULTIPLY -- Multiply a Specified Array by a Constant. &
Input \\
57 & MULTIREG & MULTIREG -- Multiply an Array by a Constant based on a
Region Number. & Input \\
58 & NOECHO & NOECHO -- Deactivate Echoing of User Input Files to the
Print File. & Output \\
59 & NOHMD & NOHMD -- Deactivate History Match Gradient Derivative
Calculations. & Gradient \\
60 & NOHMO & NOHMO -- Deactivate History Match Gradient Derivative
Calculations (Alias). & Gradient \\
61 & NOWARN & NOWARN -- Deactivate Warning Messages. & Output \\
62 & OILAPI & OILAPI -- Define the Initial Equilibration Oil API for All
Grid Blocks. & API \\
63 & OPERATE & OPERATE -- Define Mathematical Operations on Arrays. &
Input \\
64 & OPERATER & OPERATER -- Define Mathematical Operations on Arrays by
Region. & Input \\
65 & OUTSOL & OUTSOL -- Define Data to be Written to the RESTART File
(Retired). & Output \\
66 & PBUB & PBUB -- Define the Initial Equilibration Bubble-Point
Pressure for All Grid Blocks. & Enumeration \\
67 & PBVD & PBVD -- Equilibration Bubble-Point versus Depth Tables. &
Equilibration \\
68 & PDEW & PDEW -- Define the Initial Equilibration Dew-Point Pressure
for All Grid Blocks. & Enumeration \\
69 & PDVD & PDVD -- Define Equilibration Dew-Point versus Depth Tables.
& Equilibration \\
70 & PRESSURE & PRESSURE -- Define the Initial Equilibration Pressures
for All Grid Blocks. & Enumeration \\
71 & PRVD & PRVD -- Define the Initial Equilibration Pressures versus
Depth. & Enumeration \\
72 & PYEND & PYEND -- End the Definition of a PYINPUT Section. This
keyword is an OPM Flow specific keyword for the simulator's scripting
facility using the standard Python interpreter. & Action \\
73 & PYINPUT & PYINPUT -- Define the Start of a PYINPUT Section. This
keyword is an OPM Flow specific keyword for the simulator's scripting
facility using the standard Python interpreter. & Action \\
74 & RAINFALL & RAINFALL -- Constant Flux Aquifer Rainfall Flux by
Month. & Aquifer \\
75 & RBEDCONT & RBEDCONT -- Define River Grid Block Contact Area versus
Depth. & Rivers \\
76 & REFINE & REFINE -- Start the Definition of a Local Grid Refinement.
& LGR \\
77 & RESTART & RESTART -- Restart Run From an Existing Restart File. &
Restart \\
78 & RIVERSYS & RIVERSYS - Define River System (Branch Structure and
Boundary Conditions). & Rivers \\
79 & RPTRST & RPTRST -- Define Data to be Written to the RESTART File. &
Output \\
80 & RPTSOL & RPTSOL -- Define SOLUTION Section Reporting. & Output \\
81 & RS & RS -- Define the Initial Equilibration GOR (Rs) for All Grid
Blocks. & Enumeration \\
82 & RSVD & RSVD -- Equilibration Dissolved Gas-Oil Ratio (Rs) versus
Depth Tables. & Equilibration \\
83 & \hyperref[REF_HEADING_KEYWORD_RSW_10_3]{RSW} &
\hyperref[0.0.0.RSW]{RSW -- Define the Initial Equilibrium Solution Gas
in Water Ratio for All Grid Blocks}. This is a commercial compositional
simulator keyword, but may also be used with OPM Flow's Dissolved Gas in
Water model. & Enumeration \\
84 & RTEMP & RTEMP - Define the Initial Reservoir Temperature for the
Model. & Thermal \\
85 & RTEMPA & RTEMPA - Define the Initial Reservoir Temperature for the
Model. & Thermal \\
86 & RTEMPVD & RTEMPVD - Define the Initial Reservoir Temperature versus
Depth Tables. & Thermal \\
87 & RV & RV -- Define the Initial Equilibration CGR (Rv) for All Grid
Blocks. & Enumeration \\
88 & RVVD & RVVD -- Equilibration Vaporized Oil-Gas Ratio (Rv) versus
Depth Tables. & Equilibration \\
89 & RVW & RVW -- Define the Initial Equilibration Vaporized Water in
Gas Ratio for All Grid Blocks. This is an OPM Flow specific keyword for
the simulator's Water Vaporization Model that is activated by the VAPWAT
in the RUNSPEC section. & Water Vaporization \\
90 & RVWVD & RVWVD -- Equilibration Vaporized Water-Gas Ratio (Rvw)
versus Depth Tables This is an OPM Flow specific keyword for the
simulator's Water Vaporization Model that is activated by the VAPWAT in
the RUNSPEC section. & Water Vaporization \\
91 & SALT & SALT -- Define the Initial Equilibration Salt Concentration
for All Grid Blocks. A standard commercial simulator keyword used with
the Brine Model, but may also be used with OPM Flow's Vaporized Water
and Salt Precipitation models. & Brine \\
92 & SALTP & SALTP -- Define the Initial Precipitated Salt Volume
Fraction for All Grid Blocks. This is an OPM Flow specific keyword for
the simulator's Salt Precipitation Model that is activated by the
PRECSALT keyword and declaring that vaporized water is present in the
run via the VAPWAT in the RUNSPEC section. & Water Vaporization \\
93 & SALTPVD & SALTPVD -- Initial Precipitated Salt Volume Fraction
versus Depth Tables. This is an OPM Flow specific keyword for the
simulator's Salt Precipitation Model that is activated by the PRECSALT
keyword and declaring that vaporized water is present in the run via the
VAPWAT in the RUNSPEC section. & Water Vaporization \\
94 & SALTREST & SALTREST -- Define the Restart Salt Concentration for
All Grid Blocks. & Brine \\
95 & SALTVD & SALTVD -- Equilibration Salt Concentration versus Depth
Tables. A standard commercial simulator keyword used with the Brine
model, but may also be used with OPM Flow's Vaporized Water and Salt
Precipitation models. & Brine \\
96 & SBIOF & SBIOF - Define The Initial Equilibration Biofilm Volume
Fraction For All Grid Blocks. & MICP/Biofilm \\
97 & SCALC & SCALC -- Define The Initial Equilibration Calcite Volume
Fraction For All Grid Blocks. & MICP \\
98 & SCVD & SCVD -- Define Equilibration Coal Solvent Concentration
versus Depth Tables. & CBM \\
99 & SFOAM & SFOAM -- Define the Initial Equilibration Foam
Concentration for All Grid Blocks. & Enumeration \\
100 & SGAS & SGAS -- Define the Initial Equilibration Gas Saturation for
All Grid Blocks. & Enumeration \\
101 & SKIP & SKIP -- Activate Skipping of All Keywords and Input Data. &
Input \\
102 & SKIP100 & SKIP100 -- Activate Skipping of Black-Oil Keywords and
Input Data. & Input \\
103 & SKIP300 & Error: Reference source not found. & Input \\
104 & SMICR & SMICR -- Define The Initial Equilibration Microbial
Concentration For All Grid Blocks. & MICP/Biofilm \\
105 & SOIL & SOIL -- Define the Initial Equilibration Oil Saturation for
All Grid Blocks. & Enumeration \\
106 & SOLUTION & SOLUTION - Define the Start of the SOLUTION Section of
Keywords. & Input \\
107 & SOLVCONC & SOLVCONC -- Define the Initial Coal Solvent
Concentration for All Grid Blocks. & CBM \\
108 & SOLVFRAC & SOLVFRAC -- Define the Initial Gas Solvent Fraction for
All Grid Blocks & Solvent \\
109 & SOXYG & SOXYG - Define The Initial Equilibration Oxygen
Concentration For All Grid Blocks. & MICP \\
110 & SPOLY & SPOLY -- Define the Initial Equilibration Polymer
Concentration for All Grid Blocks. & Polymer \\
111 & SPOLYMW & SPOLYMW -- Define The Initial Equilibration Polymer
Molecular Weights For All Grid Blocks. This is an OPM Flow specific
keyword for the simulator's Polymer Molecular Weight Transport option
that uses the polymer molecular weight in calculating the polymer
viscosity, as well as accounting for formation damage due to the water
and polymer injection, by adjusting the wellbore skin pressure. &
Polymer \\
112 & SSOL & SSOL -- Define the Initial Equilibration Solvent Saturation
for All Grid Blocks. & Solvent \\
113 & SUREA & SUREA - Define The Initial Equilibration Urea
Concentration For All Grid Blocks. & MICP \\
114 & SURF & SURF -- Define the Initial Equilibration Polymer
Concentration for All Grid Blocks. & Surfactant \\
115 & SWAT & SWAT -- Define the Initial Equilibration Water Saturation
for All Grid Blocks. & Enumeration \\
116 & TBLK & TBLK -- Define Tracer Initial Grid Block Concentrations. &
Tracers \\
117 & TEMPI & TEMPI -- Define the Initial Temperature Values for All
Cells. & Thermal \\
118 & TEMPVD & TEMPVD - Define the Initial Reservoir Temperature versus
Depth Tables. & Thermal \\
119 & THPRES & THPRES - Define Equilibration Region Threshold Pressures.
& Equilibration \\
120 & TVDP & TVDP -- Define the Initial Equilibration Tracer Saturation
versus Depth Functions. & Tracers \\
121 & VAPPARS & VAPPARS -- Oil Vaporization Parameters. & Phases \\
122 & VISDATES & VISDATES -- Define External Reservoir Geo-Mechanics
VISAGE Stress Dates. & Geo-Mechanics \\
123 & VISOPTS & VISOPTS -- Define External Reservoir Geo-Mechanics
VISAGE Options. & Geo-Mechanics \\
124 & WARN & WARN -- Activate Warning Messages. & Output \\
Notes: & & & \\
\end{longtable}

\emph{Table 10.4: Alphabetic List of SOLUTION Keywords}

Note that a number of keywords can be classified under several topics,
so the Topic column should only be used as a general reference guide.

\subsubsection{Example}\label{example-434}

The first example is taken from the Norne field, in which the model is
in SI units, and the NTEQUL variable on the EQLDIMS keyword in the
RUNSPEC section has been set to five for five equilibration regions in
the model.

\begin{verbatim}
--
--       DATUM   DATUM   OWC     PCOW   GOC    PCGO   RS   RV   N
--       DEPTH   PRESS   DEPTH   ----   DEPTH  ----   OPT  OPT  OPT
EQUIL
         2582.0  268.56  2692.0   0.0  2582.0  0.0    1    0    0 / C+D:Garn
         2500.0  263.41  2585.5   0.0  2500.0  0.0    2    0    0 / G:  Garn
         2582.0  269.46  2618.0   0.0  2582.0  0.0    3    0    0 / E:  Garn
         2200.0  236.92  2400.0   0.0  2200.0  0.0    4    0    0 / G:Ile-Tilje
         2585.0  268.77  2693.3   0.0  2585.0  0.0    5    0    0 / C+D+E:Ile-Ti
--
--       DEPTH    RS
--                sm3/sm3
--       ------   --------
RSVD
         2582.0    120.25
         2597.0    110.00
         2660.7    106.77
         2697.0    106.77                             / RS VS DEPTH EQUIL REGN 01
--       ------   --------
         2500.0     94.50
         2510.0     94.30
         2590.0     94.10                             / RS VS DEPTH EQUIL REGN 02
--       ------   --------
         2582.0    120.25
         2597.0    110.00
         2660.7    106.77
         2697.0    106.77                             / RS VS DEPTH EQUIL REGN 03
--       ------   --------
         2500.0     94.50
         2510.0     94.30
         2590.0     94.10                             / RS VS DEPTH EQUIL REGN 04
--       ------   --------
         2585.9    120.29
         2599.9    110.00
         2663.6    106.77
         2699.9    106.77                             / RS VS DEPTH EQUIL REGN 05
--
--       EQLNUM  EQLNUM  THPRES
--       FROM    TO      VALUE
THPRES
         1       2       0.588031                     / REGN 1 TO REGN 2
         2       1       0.588031                     / REGN 2 TO REGN 1
         1       3       0.787619                     / REGN 1 TO REGN 3
         3       1       0.787619                     / REGN 3 TO REGN 1
         1       4       7.000830                     / REGN 1 TO REGN 4
         4       1       7.000830                     / REGN 4 TO REGN 1
/
\end{verbatim}

Note that even though the datum depth are the gas-oil contact depths are
the same the RSVD keyword has been used to defined the RS (GOR) versus
depth relationship. Secondly, the THPRES keyword has been used to set a
threshold pressure between the various equilibration regions to prevent
flow between the regions until the threshold is exceeded.

The next example is taken from a field containing only dead {[}``Dead''
oil is oil that it contains no dissolved gas or a relatively thick oil
or residue that has lost its volatile components.{]} oil with one
equilibration region and Fetkovich {[}Fetkovich, M. J. ``A Simplified
Approach to Water Influx Calculations - Finite Aquifer Systems,''
Journal of Petroleum Technology, (1971) 23, No.~7, 814-828.{]}
analytical aquifer connected on the flanks of the field.

\begin{verbatim}
--
--       DATUM   DATUM   OWC     PCOW   GOC    PCGO   RS   RV   N
--       DEPTH   PRESS   DEPTH   ----   DEPTH  ----   OPT  OPT  OPT
EQUIL
         3707    1609    3707    0      1*      1*     0    0     0  /
--
--       DATUM
--       DEPTH
DATUM
         3500                                                      / DATUM DEPTH
--
--                      FETKOVICH AQUIFER DESCRIPTION
--
--      ID   DATUM   AQF    AQF    AQF     AQF    AQF   SALT
--      NUM  DEPTH   PRESS  VOLM   COMP    PI     PVT   CONC
--
AQUFETP
         1   1130.   1*    1.0E+12 3.0E-5  500E3  1     0.0        /
/
--
--                      ANALYTIC AQUIFER CONNECTION
--
--  ID     ---------- BOX ---------   CONNECT  AQF    AQF      ADJOIN
--  NUMBER I1  I2   J1  J2   K1  K2   FACE     INFLX  MULTI    CELLS
--
--                      ANALYTIC AQUIFER CONNECTION
--
--       ID     ---------- BOX ---------   CONNECT  AQF    AQF    ADJOIN
--       NUMBER I1  I2   J1  J2   K1  K2   FACE     INFLX  MULTI  CELLS
AQUANCON
         1     57  57    28  36   46  58    'I+'    1*     1*     'NO' /
         1      1   1    14  14    1  29    'J-'    1*     1*     'NO' / NORTHERN
         1      2   2    12  12    1  29    'J-'    1*     1*     'NO' /
         1      3   3    10  10    1  29    'J-'    1*     1*     'NO' /
         1      4   5     9   9    1  29    'J-'    1*     1*     'NO' /
         1      6   6     8   8    1  29    'J-'    1*     1*     'NO' /
         1     14  14     4   4    1  29    'J-'    1*     1*     'NO' /
         1     15  15     2   2    1  29    'J-'    1*     1*     'NO' /
         1     16  77     1   1    1  29    'J-'    1*     1*     'NO' /
         1      1   1    14 136    1  29    'I-'    1*     1*     'NO' / WESTERN
         1     13  13   137 141    1  29    'I-'    1*     1*     'NO' /
         1     12  12   142 143    1  29    'I-'    1*     1*     'NO' /
         1     11  11   144 144    1  29    'I-'    1*     1*     'NO' /
         1     10  10   145 145    1  29    'I-'    1*     1*     'NO' /
         1      9   9   146 146    1  29    'I-'    1*     1*     'NO' /
         1      8   8   147 147    1  29    'I-'    1*     1*     'NO' /
\end{verbatim}

Notice how the gas-oil contact has been defaulted as there is no gas in
the model only oil and water. Alternatively, one could have located the
gas-oil contact above the reservoir. Secondly the DATUM keyword has been
used to reset the datum depth for pressure reporting.

\section{Keyword Definitions}\label{keyword-definitions-14}

\subsection{ADD -- Add a Constant to a Specified
Array}\label{add-add-a-constant-to-a-specified-array-4}

The ADD keyword adds a constant to a specified array or part of an
array. The constant can be real or integer depending on the array type;
however, the arrays that can be operated on is dependent on which
section the ADD keyword is being applied.

See ADD -- Add a Constant to a Specified Array in the GRID section for a
full description.

\subsection{ADDREG -- Add a Constant to an Array based on a Region
Number}\label{addreg-add-a-constant-to-an-array-based-on-a-region-number-4}

The ADDREG keyword adds a constant to a specified array or part of an
array based on cells with a specific region number. The region number
array can be FLUXNUM, MULTNUM or OPERNUM and these arrays must be
defined and be available before the ADDREG keyword is read by the
simulator. The constant can be real or integer depending on the property
array type; however, the property arrays that can be operated on is
dependent on which section the ADDREG keyword is being used.

See ADDREG -- Add a Constant to an Array based on a Region Number in the
GRID section for a full description.

\subsection{APIVD -- Equilibration Oil API Gravity versus Depth
Tables}\label{apivd-equilibration-oil-api-gravity-versus-depth-tables}

\begin{longtable}[]{@{}
  >{\raggedright\arraybackslash}p{(\linewidth - 14\tabcolsep) * \real{0.1250}}
  >{\raggedright\arraybackslash}p{(\linewidth - 14\tabcolsep) * \real{0.1250}}
  >{\raggedright\arraybackslash}p{(\linewidth - 14\tabcolsep) * \real{0.1250}}
  >{\raggedright\arraybackslash}p{(\linewidth - 14\tabcolsep) * \real{0.1250}}
  >{\raggedright\arraybackslash}p{(\linewidth - 14\tabcolsep) * \real{0.1250}}
  >{\raggedright\arraybackslash}p{(\linewidth - 14\tabcolsep) * \real{0.1250}}
  >{\raggedright\arraybackslash}p{(\linewidth - 14\tabcolsep) * \real{0.1250}}
  >{\raggedright\arraybackslash}p{(\linewidth - 14\tabcolsep) * \real{0.1250}}@{}}
\toprule\noalign{}
\begin{minipage}[b]{\linewidth}\raggedright
RUNSPEC
\end{minipage} & \begin{minipage}[b]{\linewidth}\raggedright
GRID
\end{minipage} & \begin{minipage}[b]{\linewidth}\raggedright
EDIT
\end{minipage} & \begin{minipage}[b]{\linewidth}\raggedright
PROPS
\end{minipage} & \begin{minipage}[b]{\linewidth}\raggedright
REGIONS
\end{minipage} & \begin{minipage}[b]{\linewidth}\raggedright
SOLUTION
\end{minipage} & \begin{minipage}[b]{\linewidth}\raggedright
SUMMARY
\end{minipage} & \begin{minipage}[b]{\linewidth}\raggedright
SCHEDULE
\end{minipage} \\
\midrule\noalign{}
\endhead
\bottomrule\noalign{}
\endlastfoot
\end{longtable}

\subsubsection{Description}\label{description-765}

The APIVD keyword defines the oil API gravity versus depth tables for
each equilibration region when API Tracking as been activated by the API
keyword in the RUNSPEC section.

This keyword is not supported by OPM Flow but would change the results
if supported so the simulation will be stopped.

\begin{longtable}[]{@{}
  >{\raggedright\arraybackslash}p{(\linewidth - 6\tabcolsep) * \real{0.2500}}
  >{\raggedright\arraybackslash}p{(\linewidth - 6\tabcolsep) * \real{0.2500}}
  >{\raggedright\arraybackslash}p{(\linewidth - 6\tabcolsep) * \real{0.2500}}
  >{\raggedright\arraybackslash}p{(\linewidth - 6\tabcolsep) * \real{0.2500}}@{}}
\toprule\noalign{}
\begin{minipage}[b]{\linewidth}\raggedright
No.
\end{minipage} & \begin{minipage}[b]{\linewidth}\raggedright
Name
\end{minipage} & \begin{minipage}[b]{\linewidth}\raggedright
Description
\end{minipage} & \begin{minipage}[b]{\linewidth}\raggedright
Default
\end{minipage} \\
\midrule\noalign{}
\endhead
\bottomrule\noalign{}
\endlastfoot
Field & Metric & Laboratory & \\
1 & DEPTH & A columnar vector of real monotonically increasing down the
column values that defines the depth values for the corresponding API
gravity values, API. & None \\
feet & m & cm & \\
2 & API & A columnar vector of real values that defines the API gravity
at the corresponding DEPTH. The American Petroleum Institute (``API'')
classifies oils based on an API gravity (γAPI), or degrees API (oAPI),
the relationship between relative density (γo) of oil and API gravity
(γAPI) is given by:
\({\mathrm{γ}}_{\mathit{API}} = \frac{141.5}{{\mathrm{γ}}_{o}} - 131.5\)
& None \\
oAPI & oAPI & oAPI & \\
Notes: & & & \\
\end{longtable}

\emph{Table 10.5: APIVD Keyword Description}

\subsubsection{Example}\label{example-435}

Given NTEQUL equals three and NDRXVD is greater than or equal to two on
the EQLDIMS keyword in the RUNSPEC section, then the following example
defines the bubble-point versus depth functions.

\begin{verbatim}
--
--       DEPTH    API
--                GRAVITY
--       ------   --------
APIVD
         3000.0    41.10
         8000.0    41.10                             / API VS DEPTH EQUIL REGN 01
--       ------   --------
         3000.0    41.10
         8000.0    38.50                             / API VS DEPTH EQUIL REGN 02
--       ------   --------
         3000.0    41.10
         8000.0    38.50                             / API VS DEPTH EQUIL REGN 03
\end{verbatim}

Here three tables are entered; the first table has a constant API
gravity versus depth relationship for equilibration region number one
and the other two equilibration regions have the API gravity varying
with depth.

\subsection{AQANCONL -- Define Analytical Connections to a LGR
Grid}\label{aqanconl-define-analytical-connections-to-a-lgr-grid}

\begin{longtable}[]{@{}
  >{\raggedright\arraybackslash}p{(\linewidth - 14\tabcolsep) * \real{0.1250}}
  >{\raggedright\arraybackslash}p{(\linewidth - 14\tabcolsep) * \real{0.1250}}
  >{\raggedright\arraybackslash}p{(\linewidth - 14\tabcolsep) * \real{0.1250}}
  >{\raggedright\arraybackslash}p{(\linewidth - 14\tabcolsep) * \real{0.1250}}
  >{\raggedright\arraybackslash}p{(\linewidth - 14\tabcolsep) * \real{0.1250}}
  >{\raggedright\arraybackslash}p{(\linewidth - 14\tabcolsep) * \real{0.1250}}
  >{\raggedright\arraybackslash}p{(\linewidth - 14\tabcolsep) * \real{0.1250}}
  >{\raggedright\arraybackslash}p{(\linewidth - 14\tabcolsep) * \real{0.1250}}@{}}
\toprule\noalign{}
\begin{minipage}[b]{\linewidth}\raggedright
RUNSPEC
\end{minipage} & \begin{minipage}[b]{\linewidth}\raggedright
GRID
\end{minipage} & \begin{minipage}[b]{\linewidth}\raggedright
EDIT
\end{minipage} & \begin{minipage}[b]{\linewidth}\raggedright
PROPS
\end{minipage} & \begin{minipage}[b]{\linewidth}\raggedright
REGIONS
\end{minipage} & \begin{minipage}[b]{\linewidth}\raggedright
SOLUTION
\end{minipage} & \begin{minipage}[b]{\linewidth}\raggedright
SUMMARY
\end{minipage} & \begin{minipage}[b]{\linewidth}\raggedright
SCHEDULE
\end{minipage} \\
\midrule\noalign{}
\endhead
\bottomrule\noalign{}
\endlastfoot
\end{longtable}

\subsubsection{Description}\label{description-766}

The AQANCON keyword defines how analytical aquifers are connected to a
Local Grid Refinement (``LGR'') grid, this includes the Carter-Tracy,
Fetkovich and Constant Flux analytical aquifers, all of which are
implemented in OPM Flow. Carter-Tracy analytical aquifers are
characterized by the AQUCT keyword in the GRID section and Fetkovich
analytical aquifers are defined by either the AQUFET or AQUFETP keywords
in the SOLUTION section. Finally, the Constant Flux aquifer is defined
by the AQUFLUX keyword in SOLUTION section.

This keyword is not supported by OPM Flow but would change the results
if supported so the simulation will be stopped.

\begin{longtable}[]{@{}
  >{\raggedright\arraybackslash}p{(\linewidth - 6\tabcolsep) * \real{0.2500}}
  >{\raggedright\arraybackslash}p{(\linewidth - 6\tabcolsep) * \real{0.2500}}
  >{\raggedright\arraybackslash}p{(\linewidth - 6\tabcolsep) * \real{0.2500}}
  >{\raggedright\arraybackslash}p{(\linewidth - 6\tabcolsep) * \real{0.2500}}@{}}
\toprule\noalign{}
\begin{minipage}[b]{\linewidth}\raggedright
No.
\end{minipage} & \begin{minipage}[b]{\linewidth}\raggedright
Name
\end{minipage} & \begin{minipage}[b]{\linewidth}\raggedright
Description
\end{minipage} & \begin{minipage}[b]{\linewidth}\raggedright
Default
\end{minipage} \\
\midrule\noalign{}
\endhead
\bottomrule\noalign{}
\endlastfoot
Field & Metric & Laboratory & \\
1 & AQUID & AQUID is a positive integer greater than or equal to one and
less than the maximum number of analytical aquifers as defined by the
NANAQU variable on the AQUDIMS keyword in the RUNSPEC section, that
defines the aquifer to be connected to the grid. & None \\
2 & LGRNAME & A character string of up to eight characters in length
that defines the name of the LGR that will connect to an analytical
aquifer AQUID. The LGR must have been previously defined by the either
the CARFIN (Cartesian LGR grid) keyword, or the RADIN/RADIN4 (radial LGR
grid) keyword in the GRID section. & None \\
3 & I1 & A positive integer that defines the LGR's lower bound of the
cells in the I-direction to be connected to the aquifer, and must be
greater than or equal to one and less than or equal to I2 and NX on the
CARFIN keyword for Cartesian grids. & 1 \\
4 & I2 & A positive integer that defines the LGR's upper bound of the
cells in the I-direction to be connected to the aquifer, and must be
greater than or equal to I1 and less than or equal to NX on the CARFIN
keyword for Cartesian grids. & NX \\
5 & J1 & A positive integer that defines the LGR's lower bound of the
cells in the J-direction to be connected to the aquifer, and must be
greater than or equal to one and less than or equal to J2 and NY on the
CARFIN keyword for Cartesian grids. & 1 \\
6 & J2 & A positive integer that defines the LGR's upper bound of the
cells in the J-direction to be connected to the aquifer, and must be
greater than or equal to JI and less than or equal to NY on the CARFIN
keyword for Cartesian grids. & NY \\
7 & K1 & A positive integer that defines the LGR's lower bound of the
cells in the K-direction to be to be connected to the aquifer, and must
be greater than or equal to one and less than or equal to K2 and NZ on
the CARFIN keyword for Cartesian grids. & 1 \\
8 & K2 & A positive integer that defines the LGR's upper bound of the
cells in the K-direction to be connected to the aquifer, and must be
greater than or equal to KI and less than or equal to NZ on the CARFIN
keyword for Cartesian grids. & NZ \\
9 & AQUFACE & AQUFACE is a character string that sets the connection
``face'' of the cells declared by this record and should be set to one
of the following: & None \\
10 & AQUFLUX & AQUFLUX is a positive real value that sets the fraction
of the total influx between the aquifer and the defined cells declared
on this keyword. If defaulted the cell face for each cell is applied and
if a values is declared then this values is applied to all cells
declared by this record. & 1* \\
ft2 & m2 & cm2 & \\
11 & AQUCOEF & AQUCOEF is a real positive values that scales the
calculated connection between the aquifer and the cells declared on this
record. & 1.0 \\
dimensionless & dimensionless & dimensionless & \\
12 & AQUOPT & AQUOPT is a character string that sets the cell face
connection and should be set to one of the following: & NO \\
Notes: & & & \\
\end{longtable}

\emph{Table 10.6: AQANCONL Keyword Description}

\begin{tcolorbox}[enhanced jigsaw, title=\textcolor{quarto-callout-note-color}{\faInfo}\hspace{0.5em}{Note}, coltitle=black, breakable, toprule=.15mm, colbacktitle=quarto-callout-note-color!10!white, leftrule=.75mm, opacitybacktitle=0.6, colback=white, bottomtitle=1mm, toptitle=1mm, titlerule=0mm, left=2mm, colframe=quarto-callout-note-color-frame, arc=.35mm, rightrule=.15mm, bottomrule=.15mm, opacityback=0]

If the AQANCONL keyword has been utilized in the run deck then OPM Flow
will write the AQUIFERA array to the *.INIT file in order to visualize
the aquifer connections in OPM ResInsight. This is accomplished by
setting the AQUIFERA value to 2(AQUID-1) for cells connected to aquifer
AQUID. If a cell is connected to multiple analytical aquifers then
AQUIFERA is summed for all aquifers connected to a cell. Note that
connecting cells to multiple aquifers is best avoided.

\end{tcolorbox}

\subsubsection{Example}\label{example-436}

The following example defines aquifer number one connected to the J-
face of various cells in the LGROP001 LGR, and a second basal aquifer
connected to the K+ face.

\begin{verbatim}
--
--                      LGR ANALYTIC AQUIFER CONNECTION
--
--       ID   LGR       ---------- BOX ---------   CONNECT  AQF    AQF    ADJOIN
--       NUM  NAME      I1  I2   J1  J2   K1  K2   FACE     INFLX  MULTI  CELLS
AQANCONL
         1    LGROP001   1  10   1   15   1   10    J-      1*     1.0     'NO' /
         1    LGROP001   1  10   1   15   12  18    J-      1*     1.0     'NO' /
         2    LGROP001   1  10   1   15   20  20    K+      1*     1.0     'NO' /
/
\end{verbatim}

See the AQUCT keyword in the GRID section for a complete example on
defining and connecting a Carter-Tracy aquifer to a simulation grid.

\subsection{AQANNC -- Define Analytic Aquifer Non-Neighbor
Connections}\label{aqannc-define-analytic-aquifer-non-neighbor-connections}

\begin{longtable}[]{@{}
  >{\raggedright\arraybackslash}p{(\linewidth - 14\tabcolsep) * \real{0.1250}}
  >{\raggedright\arraybackslash}p{(\linewidth - 14\tabcolsep) * \real{0.1250}}
  >{\raggedright\arraybackslash}p{(\linewidth - 14\tabcolsep) * \real{0.1250}}
  >{\raggedright\arraybackslash}p{(\linewidth - 14\tabcolsep) * \real{0.1250}}
  >{\raggedright\arraybackslash}p{(\linewidth - 14\tabcolsep) * \real{0.1250}}
  >{\raggedright\arraybackslash}p{(\linewidth - 14\tabcolsep) * \real{0.1250}}
  >{\raggedright\arraybackslash}p{(\linewidth - 14\tabcolsep) * \real{0.1250}}
  >{\raggedright\arraybackslash}p{(\linewidth - 14\tabcolsep) * \real{0.1250}}@{}}
\toprule\noalign{}
\begin{minipage}[b]{\linewidth}\raggedright
RUNSPEC
\end{minipage} & \begin{minipage}[b]{\linewidth}\raggedright
GRID
\end{minipage} & \begin{minipage}[b]{\linewidth}\raggedright
EDIT
\end{minipage} & \begin{minipage}[b]{\linewidth}\raggedright
PROPS
\end{minipage} & \begin{minipage}[b]{\linewidth}\raggedright
REGIONS
\end{minipage} & \begin{minipage}[b]{\linewidth}\raggedright
SOLUTION
\end{minipage} & \begin{minipage}[b]{\linewidth}\raggedright
SUMMARY
\end{minipage} & \begin{minipage}[b]{\linewidth}\raggedright
SCHEDULE
\end{minipage} \\
\midrule\noalign{}
\endhead
\bottomrule\noalign{}
\endlastfoot
\end{longtable}

\subsubsection{Description}\label{description-767}

AQANNC defines the analytic aquifer non-neighbor connections.

This keyword is not supported by OPM Flow but would change the results
if supported so the simulation will be stopped.

\subsection{AQANTRC -- Define Analytic Aquifer Initial Tracer
Concentrations}\label{aqantrc-define-analytic-aquifer-initial-tracer-concentrations}

\begin{longtable}[]{@{}
  >{\raggedright\arraybackslash}p{(\linewidth - 14\tabcolsep) * \real{0.1250}}
  >{\raggedright\arraybackslash}p{(\linewidth - 14\tabcolsep) * \real{0.1250}}
  >{\raggedright\arraybackslash}p{(\linewidth - 14\tabcolsep) * \real{0.1250}}
  >{\raggedright\arraybackslash}p{(\linewidth - 14\tabcolsep) * \real{0.1250}}
  >{\raggedright\arraybackslash}p{(\linewidth - 14\tabcolsep) * \real{0.1250}}
  >{\raggedright\arraybackslash}p{(\linewidth - 14\tabcolsep) * \real{0.1250}}
  >{\raggedright\arraybackslash}p{(\linewidth - 14\tabcolsep) * \real{0.1250}}
  >{\raggedright\arraybackslash}p{(\linewidth - 14\tabcolsep) * \real{0.1250}}@{}}
\toprule\noalign{}
\begin{minipage}[b]{\linewidth}\raggedright
RUNSPEC
\end{minipage} & \begin{minipage}[b]{\linewidth}\raggedright
GRID
\end{minipage} & \begin{minipage}[b]{\linewidth}\raggedright
EDIT
\end{minipage} & \begin{minipage}[b]{\linewidth}\raggedright
PROPS
\end{minipage} & \begin{minipage}[b]{\linewidth}\raggedright
REGIONS
\end{minipage} & \begin{minipage}[b]{\linewidth}\raggedright
SOLUTION
\end{minipage} & \begin{minipage}[b]{\linewidth}\raggedright
SUMMARY
\end{minipage} & \begin{minipage}[b]{\linewidth}\raggedright
SCHEDULE
\end{minipage} \\
\midrule\noalign{}
\endhead
\bottomrule\noalign{}
\endlastfoot
\end{longtable}

\subsubsection{Description}\label{description-768}

The AQANTRC keyword defines the initial tracer concentrations for
analytic aquifers that have previously been defined by the AQCT keyword
in the GRID, PROPS, or SOLUTION sections for Carter-Tracy analytical
aquifers, or the AQFET and AQFETP keywords in the SOLUTION section for
Fetkovich analytical aquifers.

This keyword is not supported by OPM Flow but would change the results
if supported so the simulation will be stopped.

\subsection{AQUALIST -- Define An Analytic Aquifer Name to Aquifer
Numbers}\label{aqualist-define-an-analytic-aquifer-name-to-aquifer-numbers}

\begin{longtable}[]{@{}
  >{\raggedright\arraybackslash}p{(\linewidth - 14\tabcolsep) * \real{0.1250}}
  >{\raggedright\arraybackslash}p{(\linewidth - 14\tabcolsep) * \real{0.1250}}
  >{\raggedright\arraybackslash}p{(\linewidth - 14\tabcolsep) * \real{0.1250}}
  >{\raggedright\arraybackslash}p{(\linewidth - 14\tabcolsep) * \real{0.1250}}
  >{\raggedright\arraybackslash}p{(\linewidth - 14\tabcolsep) * \real{0.1250}}
  >{\raggedright\arraybackslash}p{(\linewidth - 14\tabcolsep) * \real{0.1250}}
  >{\raggedright\arraybackslash}p{(\linewidth - 14\tabcolsep) * \real{0.1250}}
  >{\raggedright\arraybackslash}p{(\linewidth - 14\tabcolsep) * \real{0.1250}}@{}}
\toprule\noalign{}
\begin{minipage}[b]{\linewidth}\raggedright
RUNSPEC
\end{minipage} & \begin{minipage}[b]{\linewidth}\raggedright
GRID
\end{minipage} & \begin{minipage}[b]{\linewidth}\raggedright
EDIT
\end{minipage} & \begin{minipage}[b]{\linewidth}\raggedright
PROPS
\end{minipage} & \begin{minipage}[b]{\linewidth}\raggedright
REGIONS
\end{minipage} & \begin{minipage}[b]{\linewidth}\raggedright
SOLUTION
\end{minipage} & \begin{minipage}[b]{\linewidth}\raggedright
SUMMARY
\end{minipage} & \begin{minipage}[b]{\linewidth}\raggedright
SCHEDULE
\end{minipage} \\
\midrule\noalign{}
\endhead
\bottomrule\noalign{}
\endlastfoot
\end{longtable}

\subsubsection{Description}\label{description-769}

The AQUALIST keyword defines named lists of analytic aquifers identified
by aquifer numbers for greater readability in the output.

This keyword is not supported by OPM Flow but would change the results
if supported so the simulation will be stopped.

\subsection{AQUANCON -- Define Analytical Connections to the
Grid}\label{aquancon-define-analytical-connections-to-the-grid-1}

The AQUANCON keyword defines how analytical aquifers are connected to
the simulation grid, this includes Carter-Tracy and Fetkovich analytical
aquifers, both of which are implemented in OPM Flow.

See AQUANCON -- Define Analytical Connections to the Grid in the GRID
section for a full description.

\subsection{AQUCHGAS -- Define Constant Pressure Gas Analytical Aquifer
Properties}\label{aquchgas-define-constant-pressure-gas-analytical-aquifer-properties}

\begin{longtable}[]{@{}
  >{\raggedright\arraybackslash}p{(\linewidth - 14\tabcolsep) * \real{0.1250}}
  >{\raggedright\arraybackslash}p{(\linewidth - 14\tabcolsep) * \real{0.1250}}
  >{\raggedright\arraybackslash}p{(\linewidth - 14\tabcolsep) * \real{0.1250}}
  >{\raggedright\arraybackslash}p{(\linewidth - 14\tabcolsep) * \real{0.1250}}
  >{\raggedright\arraybackslash}p{(\linewidth - 14\tabcolsep) * \real{0.1250}}
  >{\raggedright\arraybackslash}p{(\linewidth - 14\tabcolsep) * \real{0.1250}}
  >{\raggedright\arraybackslash}p{(\linewidth - 14\tabcolsep) * \real{0.1250}}
  >{\raggedright\arraybackslash}p{(\linewidth - 14\tabcolsep) * \real{0.1250}}@{}}
\toprule\noalign{}
\begin{minipage}[b]{\linewidth}\raggedright
RUNSPEC
\end{minipage} & \begin{minipage}[b]{\linewidth}\raggedright
GRID
\end{minipage} & \begin{minipage}[b]{\linewidth}\raggedright
EDIT
\end{minipage} & \begin{minipage}[b]{\linewidth}\raggedright
PROPS
\end{minipage} & \begin{minipage}[b]{\linewidth}\raggedright
REGIONS
\end{minipage} & \begin{minipage}[b]{\linewidth}\raggedright
SOLUTION
\end{minipage} & \begin{minipage}[b]{\linewidth}\raggedright
SUMMARY
\end{minipage} & \begin{minipage}[b]{\linewidth}\raggedright
SCHEDULE
\end{minipage} \\
\midrule\noalign{}
\endhead
\bottomrule\noalign{}
\endlastfoot
\end{longtable}

\subsubsection{Description}\label{description-770}

The AQUCHGAS keyword defines the properties of constant pressure gas
analytical aquifers.

This keyword is not supported by OPM Flow but would change the results
if supported so the simulation will be stopped.

\subsection{AQUCHWAT -- Define Constant Pressure Water Analytical
Aquifer
Properties}\label{aquchwat-define-constant-pressure-water-analytical-aquifer-properties}

\begin{longtable}[]{@{}
  >{\raggedright\arraybackslash}p{(\linewidth - 14\tabcolsep) * \real{0.1250}}
  >{\raggedright\arraybackslash}p{(\linewidth - 14\tabcolsep) * \real{0.1250}}
  >{\raggedright\arraybackslash}p{(\linewidth - 14\tabcolsep) * \real{0.1250}}
  >{\raggedright\arraybackslash}p{(\linewidth - 14\tabcolsep) * \real{0.1250}}
  >{\raggedright\arraybackslash}p{(\linewidth - 14\tabcolsep) * \real{0.1250}}
  >{\raggedright\arraybackslash}p{(\linewidth - 14\tabcolsep) * \real{0.1250}}
  >{\raggedright\arraybackslash}p{(\linewidth - 14\tabcolsep) * \real{0.1250}}
  >{\raggedright\arraybackslash}p{(\linewidth - 14\tabcolsep) * \real{0.1250}}@{}}
\toprule\noalign{}
\begin{minipage}[b]{\linewidth}\raggedright
RUNSPEC
\end{minipage} & \begin{minipage}[b]{\linewidth}\raggedright
GRID
\end{minipage} & \begin{minipage}[b]{\linewidth}\raggedright
EDIT
\end{minipage} & \begin{minipage}[b]{\linewidth}\raggedright
PROPS
\end{minipage} & \begin{minipage}[b]{\linewidth}\raggedright
REGIONS
\end{minipage} & \begin{minipage}[b]{\linewidth}\raggedright
SOLUTION
\end{minipage} & \begin{minipage}[b]{\linewidth}\raggedright
SUMMARY
\end{minipage} & \begin{minipage}[b]{\linewidth}\raggedright
SCHEDULE
\end{minipage} \\
\midrule\noalign{}
\endhead
\bottomrule\noalign{}
\endlastfoot
\end{longtable}

\subsubsection{Description}\label{description-771}

The AQUCHWAT keyword defines the properties of constant pressure water
analytical aquifers.

This keyword is not supported by OPM Flow but would change the results
if supported so the simulation will be stopped.

\subsection{AQUCON -- Define Numerical Aquifer Connections to the
Grid}\label{aqucon-define-numerical-aquifer-connections-to-the-grid-1}

AQUCON keyword defines how numerical aquifers are connected to the
simulation grid.

See AQUCON -- Define Numerical Aquifer Connections to the Grid in the
GRID section for a full description.

\subsection{AQUCT -- Define Carter-Tracy Analytical
Aquifers}\label{aquct-define-carter-tracy-analytical-aquifers-2}

The AQUCT keyword defines a Carter Tracy aquifer, the properties of the
aquifer, including the Carter-Tracy aquifer influence function
associated with the aquifer, defined by the AQUTAB keyword in the PROPS
section.

See AQUCT -- Define Carter-Tracy Analytical Aquifers in the GRID section
for a full description

\subsection{AQUFET -- Define Fetkovich Analytical Aquifer and
Connections}\label{aqufet-define-fetkovich-analytical-aquifer-and-connections}

\begin{longtable}[]{@{}
  >{\raggedright\arraybackslash}p{(\linewidth - 14\tabcolsep) * \real{0.1250}}
  >{\raggedright\arraybackslash}p{(\linewidth - 14\tabcolsep) * \real{0.1250}}
  >{\raggedright\arraybackslash}p{(\linewidth - 14\tabcolsep) * \real{0.1250}}
  >{\raggedright\arraybackslash}p{(\linewidth - 14\tabcolsep) * \real{0.1250}}
  >{\raggedright\arraybackslash}p{(\linewidth - 14\tabcolsep) * \real{0.1250}}
  >{\raggedright\arraybackslash}p{(\linewidth - 14\tabcolsep) * \real{0.1250}}
  >{\raggedright\arraybackslash}p{(\linewidth - 14\tabcolsep) * \real{0.1250}}
  >{\raggedright\arraybackslash}p{(\linewidth - 14\tabcolsep) * \real{0.1250}}@{}}
\toprule\noalign{}
\begin{minipage}[b]{\linewidth}\raggedright
RUNSPEC
\end{minipage} & \begin{minipage}[b]{\linewidth}\raggedright
GRID
\end{minipage} & \begin{minipage}[b]{\linewidth}\raggedright
EDIT
\end{minipage} & \begin{minipage}[b]{\linewidth}\raggedright
PROPS
\end{minipage} & \begin{minipage}[b]{\linewidth}\raggedright
REGIONS
\end{minipage} & \begin{minipage}[b]{\linewidth}\raggedright
SOLUTION
\end{minipage} & \begin{minipage}[b]{\linewidth}\raggedright
SUMMARY
\end{minipage} & \begin{minipage}[b]{\linewidth}\raggedright
SCHEDULE
\end{minipage} \\
\midrule\noalign{}
\endhead
\bottomrule\noalign{}
\endlastfoot
\end{longtable}

\subsubsection{Description}\label{description-772}

The AQUFET keyword defines Fetkovich {[}Fetkovich, M. J. ``A Simplified
Approach to Water Influx Calculations - Finite Aquifer Systems,''
Journal of Petroleum Technology, (1971) 23, No.~7, 814-828.{]}
analytical aquifers, the aquifer properties, together with the cell
connections to the aquifer. Each row entry in the AQUFETP keyword
defines one Fetkovich analytical aquifer and one cell face to be
connected to the aquifer.

This keyword is not supported by OPM Flow but would change the results
if supported so the simulation will be stopped; however, see the AQUFETP
keyword in the SOLUTION section and AQUANCON keyword in the GRID
section, on how to define and connect Fetkovich analytical aquifers.

\begin{longtable}[]{@{}
  >{\raggedright\arraybackslash}p{(\linewidth - 6\tabcolsep) * \real{0.2500}}
  >{\raggedright\arraybackslash}p{(\linewidth - 6\tabcolsep) * \real{0.2500}}
  >{\raggedright\arraybackslash}p{(\linewidth - 6\tabcolsep) * \real{0.2500}}
  >{\raggedright\arraybackslash}p{(\linewidth - 6\tabcolsep) * \real{0.2500}}@{}}
\toprule\noalign{}
\begin{minipage}[b]{\linewidth}\raggedright
No.
\end{minipage} & \begin{minipage}[b]{\linewidth}\raggedright
Name
\end{minipage} & \begin{minipage}[b]{\linewidth}\raggedright
Description
\end{minipage} & \begin{minipage}[b]{\linewidth}\raggedright
Default
\end{minipage} \\
\midrule\noalign{}
\endhead
\bottomrule\noalign{}
\endlastfoot
Field & Metric & Laboratory & \\
1 & DATUM & DATUM is a single positive value that defines the Fetkovich
reference datum depth for PRESS. & None \\
feet & m & cm & \\
2 & PRESS & PRESS is a single positive value that defines the aquifer
pressure at DATUM. If PRESS is defaulted then the simulator will set the
aquifer's initial reservoir pressure to be in equilibrium with the cells
the aquifer is contacted to. Defaulting this parameter will avoid
inconsistent equilibration pressures between the reservoir cells and the
aquifer. & 1* \\
psia & barsa & atma & \\
3 & PORV & & None \\
stb & sm3 & scc & \\
4 & COMP & COMP is a real number defining the total compressibility (Ct)
of the aquifer, that is the rock compressibility (Cf) plus the water
compressibility (Cw) at the aquifer datum pressure (DATUM) and is
defined as: \({C}_{t} = {C}_{f} + {C}_{w}\) & None \\
1/psia & 1/barsa & 1/atma & \\
5 & PI & A real positive number that defines the aquifer productivity
index based on the aquifer influx rate per unit pressure drop. & None \\
stb/d/psia & sm3/barsa & scc/hr/atma & \\
6 & PVTW & A positive integer that defines the aquifer's PVTW water
property table. & 1 \\
7 & I1 & A positive integer that defines the lower bound of the cells in
the I-direction to be connected to the aquifer and must be greater than
or equal to one and less than or equal to I2 and NX. & 1 \\
8 & I2 & A positive integer that defines the upper bound of the cells in
the I-direction to be connected to the aquifer and must be greater than
or equal to I1 and less than or equal to NX & NX \\
9 & J1 & A positive integer that defines the lower bound of the cells in
the J-direction to be connected to the aquifer and must be greater than
or equal to one and less than or equal to J2 and NY. & 1 \\
10 & J2 & A positive integer that defines the upper bound of the cells
in the J-direction to be connected to the aquifer and must be greater
than or equal to JI and less than or equal to NY. & NY \\
11 & K1 & A positive integer that defines the lower bound of the cells
in the K-direction to be to be connected to the aquifer and must be
greater than or equal to one and less than or equal to K2 and NZ. & 1 \\
12 & K2 & A positive integer that defines the upper bound of the cells
in the K-direction to be connected to the aquifer and must be greater
than or equal to KI and less than or equal to NZ. & NZ \\
13 & AQUFACE & AQUFACE is a character string that sets the connection
``face'' of the cells declared by this record and should be set to one
of the following: & None \\
14 & SALTCON & SALTCON is a real positive number that defines the
initial salt concentration in the aquifer. This variable is ignored by
OPM Flow. & 0.0 \\
lb/stb & kg/sm3 & gm/scc & \\
Notes: & & & \\
\end{longtable}

\emph{Table 10.7: AQUFET Keyword Description}

Note this keyword should only be used in equilibration and enumerated
initialize runs, that is it should be omitted from RESTART runs.

\begin{tcolorbox}[enhanced jigsaw, title=\textcolor{quarto-callout-note-color}{\faInfo}\hspace{0.5em}{Note}, coltitle=black, breakable, toprule=.15mm, colbacktitle=quarto-callout-note-color!10!white, leftrule=.75mm, opacitybacktitle=0.6, colback=white, bottomtitle=1mm, toptitle=1mm, titlerule=0mm, left=2mm, colframe=quarto-callout-note-color-frame, arc=.35mm, rightrule=.15mm, bottomrule=.15mm, opacityback=0]

If the model is unstable then this may be due to an aquifer not being in
equilibrium with the connecting reservoir blocks, for example the
aquifer is connected to only hydrocarbon reservoir cells. Try commenting
out the aquifer and see if this resolves the instabilities.

\end{tcolorbox}

\subsubsection{Example}\label{example-437}

Given the following grid and aquifer dimensions in the RUNSPEC section:

\begin{verbatim}
-- ==============================================================================
--
-- RUNSPEC SECTION
--
-- ==============================================================================
RUNSPEC   --
--       MAX     MAX     MAX
--       NDIVIX  NDIVIY  NDIVIZ
DIMENS
         20      1       5                                                    /

--       AQF     AQF     AQF     AQF     AQF     AQF    AQF    AQF
--       MXAQN   MXNAQC  NIFTBL  NRIFTB  NANAQU  NCAMAX MXNALI MXAAQL
AQUDIMS
         1*      1*      5       100     1       1*     1*     1*              /
\end{verbatim}

The Fetkovich Analytical aquifer is defined in the SOLUTION sections as:

\begin{verbatim}
-- ==============================================================================
--
-- SOLUTION SECTION
--
-- ==============================================================================
SOLUTION --
--                      FETKOVICH AQUIFER DESCRIPTION AND CONNECTIONS
--
--       DATUM  AQF    AQF    AQF     AQF  AQF -------- BOX ------- CONNECT SALT
--       DEPTH  PRESS  VOLM   COMP    PI   PVT I1  I2  J1  J2  K1  K2 FACE  CONC
--
AQUFET
         1130.  1*    1.0E+12 3.0E-5  500E3 1   1   1   1   1   1   1  'J-'  /
\end{verbatim}

Here one Fetkovich Analytical aquifer is connected to a single cell (1,
1, 1) at the J- face (or X- face) of the grid.

\subsection{AQUFETP -- Define Fetkovich Analytical
Aquifers}\label{aqufetp-define-fetkovich-analytical-aquifers}

\begin{longtable}[]{@{}
  >{\raggedright\arraybackslash}p{(\linewidth - 14\tabcolsep) * \real{0.1250}}
  >{\raggedright\arraybackslash}p{(\linewidth - 14\tabcolsep) * \real{0.1250}}
  >{\raggedright\arraybackslash}p{(\linewidth - 14\tabcolsep) * \real{0.1250}}
  >{\raggedright\arraybackslash}p{(\linewidth - 14\tabcolsep) * \real{0.1250}}
  >{\raggedright\arraybackslash}p{(\linewidth - 14\tabcolsep) * \real{0.1250}}
  >{\raggedright\arraybackslash}p{(\linewidth - 14\tabcolsep) * \real{0.1250}}
  >{\raggedright\arraybackslash}p{(\linewidth - 14\tabcolsep) * \real{0.1250}}
  >{\raggedright\arraybackslash}p{(\linewidth - 14\tabcolsep) * \real{0.1250}}@{}}
\toprule\noalign{}
\begin{minipage}[b]{\linewidth}\raggedright
RUNSPEC
\end{minipage} & \begin{minipage}[b]{\linewidth}\raggedright
GRID
\end{minipage} & \begin{minipage}[b]{\linewidth}\raggedright
EDIT
\end{minipage} & \begin{minipage}[b]{\linewidth}\raggedright
PROPS
\end{minipage} & \begin{minipage}[b]{\linewidth}\raggedright
REGIONS
\end{minipage} & \begin{minipage}[b]{\linewidth}\raggedright
SOLUTION
\end{minipage} & \begin{minipage}[b]{\linewidth}\raggedright
SUMMARY
\end{minipage} & \begin{minipage}[b]{\linewidth}\raggedright
SCHEDULE
\end{minipage} \\
\midrule\noalign{}
\endhead
\bottomrule\noalign{}
\endlastfoot
\end{longtable}

\subsubsection{Description}\label{description-773}

The AQUFETP keyword defines Fetkovich {[}Fetkovich, M. J. ``A Simplified
Approach to Water Influx Calculations - Finite Aquifer Systems,''
Journal of Petroleum Technology, (1971) 23, No.~7, 814-828.{]}
analytical aquifers and the aquifer properties. Each row entry in the
AQUFETP keyword defines one Fetkovich analytical aquifer. In order to
fully define this type of aquifer, the aquifer must be connected to the
reservoir using the AQUANCON keyword in the GRID or SOLUTION sections.

\begin{longtable}[]{@{}
  >{\raggedright\arraybackslash}p{(\linewidth - 6\tabcolsep) * \real{0.2500}}
  >{\raggedright\arraybackslash}p{(\linewidth - 6\tabcolsep) * \real{0.2500}}
  >{\raggedright\arraybackslash}p{(\linewidth - 6\tabcolsep) * \real{0.2500}}
  >{\raggedright\arraybackslash}p{(\linewidth - 6\tabcolsep) * \real{0.2500}}@{}}
\toprule\noalign{}
\begin{minipage}[b]{\linewidth}\raggedright
No.
\end{minipage} & \begin{minipage}[b]{\linewidth}\raggedright
Name
\end{minipage} & \begin{minipage}[b]{\linewidth}\raggedright
Description
\end{minipage} & \begin{minipage}[b]{\linewidth}\raggedright
Default
\end{minipage} \\
\midrule\noalign{}
\endhead
\bottomrule\noalign{}
\endlastfoot
Field & Metric & Laboratory & \\
1 & AQUID & A positive integer greater than or equal to one and less
than or equal to NANAQ on the AQUDIMS keyword in the RUNSPEC section,
that defines the Fetkovich aquifer number. & 1 \\
2 & DATUM & DATUM is a single positive value that defines the Fetkovich
reference datum depth for PRESS. & None \\
feet & m & cm & \\
3 & PRESS & PRESS is a single positive value that defines the aquifer
pressure at DATUM. If PRESS is defaulted then the simulator will set the
aquifer's initial reservoir pressure to be in equilibrium with the cells
the aquifer is contacted to. Defaulting this parameter will avoid
inconsistent equilibration pressures between the reservoir cells and the
aquifer. & 1* \\
psia & barsa & atma & \\
4 & PORV & A real positive value that defines the initial water volume
of the aquifer. & None \\
stb & sm3 & scc & \\
5 & COMP & COMP is a real number defining the total compressibility (Ct)
of the aquifer, that is the rock compressibility (Cf) plus the water
compressibility (Cw) at the aquifer datum pressure (DATUM) and is
defined as: \({C}_{t} = {C}_{f} + {C}_{w}\) & None \\
1/psia & 1/barsa & 1/atma & \\
6 & PI & A real positive number that defines the aquifer productivity
index based on the aquifer influx rate per unit pressure drop. & None \\
stb/d/psia & sm3/d/barsa & scc/hr/atma & \\
7 & PVTW & A positive integer that defines the aquifer's PVTW water
property table. & 1 \\
8 & SALTCON & SALTCON is a real positive number that defines the initial
salt concentration in the aquifer, for when the simulator's Brine Model
has been activated via the BRINE keyword in the RUNSPEC section. This
variable is ignored by OPM Flow. & 0.0 \\
lb/stb & kg/sm3 & gm/scc & \\
9 & TEMP & TEMP is a real positive number that defines the initial
temperature of the aquifer at DATUM for use with OPM Flow's thermal
option. The THERMAL keyword in the RUNSPEC section must be activated to
use this option. & 1* \\
oF & oC & oC & \\
Notes: & & & \\
\end{longtable}

\emph{Table 10.8: AQUFETP Keyword Description}

Note this keyword should only be used in equilibration and enumerated
initialize runs, that is it should be omitted from RESTART runs.

\begin{tcolorbox}[enhanced jigsaw, title=\textcolor{quarto-callout-note-color}{\faInfo}\hspace{0.5em}{Note}, coltitle=black, breakable, toprule=.15mm, colbacktitle=quarto-callout-note-color!10!white, leftrule=.75mm, opacitybacktitle=0.6, colback=white, bottomtitle=1mm, toptitle=1mm, titlerule=0mm, left=2mm, colframe=quarto-callout-note-color-frame, arc=.35mm, rightrule=.15mm, bottomrule=.15mm, opacityback=0]

If the model is unstable then this may be due to an aquifer not being in
equilibrium with the connecting reservoir blocks, for example if the
aquifer is connected to some hydrocarbon reservoir cells. Try commenting
out the aquifer and see if this resolves the instabilities, and if so
amend the aquifer connections accordingly.

\end{tcolorbox}

\subsubsection{Example}\label{example-438}

Given the following grid and aquifer dimensions in the RUNSPEC section:

\begin{verbatim}
-- ==============================================================================
--
-- RUNSPEC SECTION
--
-- ==============================================================================
RUNSPEC
--
--       MAX     MAX     MAX
--       NDIVIX  NDIVIY  NDIVIZ
DIMENS
         20      1       5                                                    /

--       AQF     AQF     AQF     AQF     AQF     AQF    AQF    AQF
--       MXAQN   MXNAQC  NIFTBL  NRIFTB  NANAQU  NCAMAX MXNALI MXAAQL
AQUDIMS
         1*      1*      5       100     1       1*     1*     1*              /
\end{verbatim}

The Fetkovich analytical aquifer is defined in the SOLUTION sections as:

\begin{verbatim}
-- ==============================================================================
--
-- SOLUTION SECTION
--
-- ==============================================================================
SOLUTION --
--                      FETKOVICH AQUIFER DESCRIPTION
--
--      ID   DATUM   AQF    AQF    AQF     AQF    AQF   SALT
--      NUM  DEPTH   PRESS  VOLM   COMP    PI     PVT   CONC
--
AQUFETP
         1   1130.   1*    1.0E+12 3.0E-5  500E3  1     0.0                    /
/
\end{verbatim}

And the connection of the aquifer is set in the GRID or the SOLUTION
sections as:

\begin{verbatim}
--
--                      ANALYTIC AQUIFER CONNECTION
--
--       ID     ---------- BOX ---------   CONNECT  AQF    AQF      ADJOIN
--       NUMBER I1  I2   J1  J2   K1  K2   FACE     INFLX  MULTI    CELLS
AQUANCON
         1      1   1    1   1    1   1     J-      1.0    1.0      'NO'        /
/
\end{verbatim}

Here one Fetkovich analytical aquifer is connected to a single cell (1,
1, 1) at the J- face (or Y- face) of the cell.

\subsection{AQUFLUX -- Define Constant Flux Analytical
Aquifer}\label{aquflux-define-constant-flux-analytical-aquifer}

\begin{longtable}[]{@{}
  >{\raggedright\arraybackslash}p{(\linewidth - 14\tabcolsep) * \real{0.1250}}
  >{\raggedright\arraybackslash}p{(\linewidth - 14\tabcolsep) * \real{0.1250}}
  >{\raggedright\arraybackslash}p{(\linewidth - 14\tabcolsep) * \real{0.1250}}
  >{\raggedright\arraybackslash}p{(\linewidth - 14\tabcolsep) * \real{0.1250}}
  >{\raggedright\arraybackslash}p{(\linewidth - 14\tabcolsep) * \real{0.1250}}
  >{\raggedright\arraybackslash}p{(\linewidth - 14\tabcolsep) * \real{0.1250}}
  >{\raggedright\arraybackslash}p{(\linewidth - 14\tabcolsep) * \real{0.1250}}
  >{\raggedright\arraybackslash}p{(\linewidth - 14\tabcolsep) * \real{0.1250}}@{}}
\toprule\noalign{}
\begin{minipage}[b]{\linewidth}\raggedright
RUNSPEC
\end{minipage} & \begin{minipage}[b]{\linewidth}\raggedright
GRID
\end{minipage} & \begin{minipage}[b]{\linewidth}\raggedright
EDIT
\end{minipage} & \begin{minipage}[b]{\linewidth}\raggedright
PROPS
\end{minipage} & \begin{minipage}[b]{\linewidth}\raggedright
REGIONS
\end{minipage} & \begin{minipage}[b]{\linewidth}\raggedright
SOLUTION
\end{minipage} & \begin{minipage}[b]{\linewidth}\raggedright
SUMMARY
\end{minipage} & \begin{minipage}[b]{\linewidth}\raggedright
SCHEDULE
\end{minipage} \\
\midrule\noalign{}
\endhead
\bottomrule\noalign{}
\endlastfoot
\end{longtable}

\subsubsection{Description}\label{description-774}

The AQUFLUX keyword defines the properties of Constant Flux Analytical
Aquifers, that allows for a constant water influx to the connected grid
blocks. This type of aquifer is connected to the model using either the
AQUANCON keyword for global cells, or the AQUANCONL keyword for cells
belonging to a Local Grid Refinement (``LGR''), both the aforementioned
keywords are in the GRID and SOLUTION sections. The keyword itself may
be utilized in both the SOLUTION and SCHEDULE sections, with subsequent
entries overwriting the previous entry.

\begin{longtable}[]{@{}
  >{\raggedright\arraybackslash}p{(\linewidth - 6\tabcolsep) * \real{0.2500}}
  >{\raggedright\arraybackslash}p{(\linewidth - 6\tabcolsep) * \real{0.2500}}
  >{\raggedright\arraybackslash}p{(\linewidth - 6\tabcolsep) * \real{0.2500}}
  >{\raggedright\arraybackslash}p{(\linewidth - 6\tabcolsep) * \real{0.2500}}@{}}
\toprule\noalign{}
\begin{minipage}[b]{\linewidth}\raggedright
No.
\end{minipage} & \begin{minipage}[b]{\linewidth}\raggedright
Name
\end{minipage} & \begin{minipage}[b]{\linewidth}\raggedright
Description
\end{minipage} & \begin{minipage}[b]{\linewidth}\raggedright
Default
\end{minipage} \\
\midrule\noalign{}
\endhead
\bottomrule\noalign{}
\endlastfoot
Field & Metric & Laboratory & \\
1 & AQUID & A positive integer greater than or equal to one and less
than or equal to NANAQ on the AQUDIMS keyword in the RUNSPEC section,
that defines the AQUFLUX aquifer number. & 1 \\
‍2 ‍ & AQFLUX & A real positive value the defines the aquifer water influx
rate, per unit area of the connected grid cell face. & None \\
stb/day/ft2 & sm3/day/m2 & scc/hour/cm2 & \\
3 & SALTCON & SALTCON is a real positive number that defines the initial
salt concentration in the aquifer, for when with simulator's Brine Model
has been activated via the BRINE keyword in the RUNSPEC section. This
variable is ignored by OPM Flow. & 0.0 \\
lb/stb & kg/sm3 & gm/scc & \\
4 & TEMP & TEMP is a real positive number that defines the initial
temperature of the aquifer for use with OPM Flow's thermal option. The
THERMAL keyword in the RUNSPEC section must be activated to use this
option. If the THERMAL keyword is absent from the input deck, then the
parameter is ignored. & 1* \\
oF & oC & oC & \\
5 & PRESS & PRESS is a single positive value that defines the aquifer
pressure for use with OPM Flow's thermal option. The THERMAL keyword in
the RUNSPEC section must be activated to use this option. If the THERMAL
keyword is absent from the input deck, then the parameter is ignored. &
1* \\
psia & barsa & atma & \\
Notes: & & & \\
\end{longtable}

\emph{Table 10.9: AQUFLUX Keyword Description}

The water flow rate into a connected grid cell for this type of aquifer,
is calculated from:

\begin{equation}\phantomsection\label{eq-10-14}{
{q}_{w} = \mathit{AQUFLUX}(\mathit{AQFLUX}) \times  {A}_{(i,j,k)} \times  \mathit{AQUANCON}(\mathit{AQUCOEF})
}\end{equation}

Where:

AQUFLUX(AQFLUX) = the AQFLUX parameter on the AQUFLUX keyword,

AQUANCON(AQUCOEF) = the AQUCOEF parameter on the AQUANCON keyword

Note that \({A}_{(i,j,k)}\)is calculated from the connected cell
geometry, and thus the AQUANCON(AQUFLUX) and the AQUANCONL(AQUFLUX)
parameters are ignored for this type of aquifer.

The AQUFLUX keyword should only be used in runs that have been
initialized by equilibration or with enumerated values, and should not
be used in RESTART runs.

See also the \hyperref[REF_HEADING_KEYWORD_BCCON]{BCCON} keyword in the
GRID section and the \hyperref[REF_HEADING_KEYWORD_BCPROP]{BCPROP}
keyword in the SCHEDULE section that can be used to specify boundary
conditions.

\subsubsection{Example}\label{example-439}

Given a total of five analytical aquifers as declared by the NANAQ
parameter on the AQUDIMS keyword being set to five, of which three are
Constant Flux Aquifers and two Carter-Tracy Aquifers, then:

\begin{verbatim}
--
--                      CONSTANT FLUX AQUIFER DESCRIPTION
--
--      ID   WATER    SALT   AQF    AQU
--      NUM  INFLUX   CONC   TEMP   PRES
--
AQUFLUX
         1   0.0004   1*     1*     1*                      /
         2   0.0005   1*     1*     1*                      /
         3   0.0003   1*     1*     1*                      /
/
--
--                      CARTER-TRACY AQUIFER DESCRIPTION
--
--      ID   DATUM   AQF    AQF    AQF    AQF     AQF   AQF   INFL   PVT  AQU
--      NUM  DEPTH   PRESS  PERM   PORO   RCOMP   RE    DZ    ANGLE  NUM  TAB
--
AQUCT
         4   2000.0  269    100.0  0.30   3.0e-5  330   10.0  360.0   1    2   /
         5   2000.0  269    100.0  0.30   3.0e-5  330   10.0  360.0   1    2   /
/
\end{verbatim}

Defines three Constant Flux Aquifers and three Carter-Tracy Aquifers,
and the connection of the aquifers are set in the GRID or SOLUTION
sections via:

\begin{verbatim}
--
--                      ANALYTIC AQUIFER CONNECTION
--
--       ID     ---------- BOX ---------   CONNECT  AQF    AQF      ADJOIN
--       NUMBER I1  I2   J1  J2   K1  K2   FACE     INFLX  MULTI    CELLS
AQUANCON
         1      1   198  1   14   1   10    J-      1*     1.0      'NO'        /
         2      1   198  1   14   12  22    J-      1*     1.0      'NO'        /
         3      1   198  1   14   23  54    J-      1*     1.0      'NO'        /
         4      1   198  1   14   61  70    J-      1*     1.0      'NO'        /
         5      1   198  15  172  71  71    K+      1*     1.0      'NO'        /
/

\end{verbatim}

\subsection{BC -- Define Boundary
Conditions}\label{bc-define-boundary-conditions}

\begin{longtable}[]{@{}
  >{\raggedright\arraybackslash}p{(\linewidth - 14\tabcolsep) * \real{0.1250}}
  >{\raggedright\arraybackslash}p{(\linewidth - 14\tabcolsep) * \real{0.1250}}
  >{\raggedright\arraybackslash}p{(\linewidth - 14\tabcolsep) * \real{0.1250}}
  >{\raggedright\arraybackslash}p{(\linewidth - 14\tabcolsep) * \real{0.1250}}
  >{\raggedright\arraybackslash}p{(\linewidth - 14\tabcolsep) * \real{0.1250}}
  >{\raggedright\arraybackslash}p{(\linewidth - 14\tabcolsep) * \real{0.1250}}
  >{\raggedright\arraybackslash}p{(\linewidth - 14\tabcolsep) * \real{0.1250}}
  >{\raggedright\arraybackslash}p{(\linewidth - 14\tabcolsep) * \real{0.1250}}@{}}
\toprule\noalign{}
\begin{minipage}[b]{\linewidth}\raggedright
RUNSPEC
\end{minipage} & \begin{minipage}[b]{\linewidth}\raggedright
GRID
\end{minipage} & \begin{minipage}[b]{\linewidth}\raggedright
EDIT
\end{minipage} & \begin{minipage}[b]{\linewidth}\raggedright
PROPS
\end{minipage} & \begin{minipage}[b]{\linewidth}\raggedright
REGIONS
\end{minipage} & \begin{minipage}[b]{\linewidth}\raggedright
SOLUTION
\end{minipage} & \begin{minipage}[b]{\linewidth}\raggedright
SUMMARY
\end{minipage} & \begin{minipage}[b]{\linewidth}\raggedright
SCHEDULE
\end{minipage} \\
\midrule\noalign{}
\endhead
\bottomrule\noalign{}
\endlastfoot
\end{longtable}

\subsubsection{Description}\label{description-775}

The BC keyword defines the boundary conditions for the model, and can be
used to set boundary conditions for when external influx or efflux
volumes are influencing the reservoir pressure and production history.
For example, when the average reservoir pressure remains constant
throughout the production period due to water influx, or gas migration
from an external source.

The BC keyword has been replaced by the
\hyperref[REF_HEADING_KEYWORD_BCCON]{BCCON} keyword in the GRID section
that defines the boundary condition connections and the
\hyperref[REF_HEADING_KEYWORD_BCPROP]{BCPROP} keyword in the SCHEDULE
section that defines the boundary condition properties.

\begin{tcolorbox}[enhanced jigsaw, title=\textcolor{quarto-callout-note-color}{\faInfo}\hspace{0.5em}{Note}, coltitle=black, breakable, toprule=.15mm, colbacktitle=quarto-callout-note-color!10!white, leftrule=.75mm, opacitybacktitle=0.6, colback=white, bottomtitle=1mm, toptitle=1mm, titlerule=0mm, left=2mm, colframe=quarto-callout-note-color-frame, arc=.35mm, rightrule=.15mm, bottomrule=.15mm, opacityback=0]

This is an OPM Flow specific keyword and will therefore cause an error
in the commercial simulator.

\end{tcolorbox}

\begin{longtable}[]{@{}
  >{\raggedright\arraybackslash}p{(\linewidth - 6\tabcolsep) * \real{0.2500}}
  >{\raggedright\arraybackslash}p{(\linewidth - 6\tabcolsep) * \real{0.2500}}
  >{\raggedright\arraybackslash}p{(\linewidth - 6\tabcolsep) * \real{0.2500}}
  >{\raggedright\arraybackslash}p{(\linewidth - 6\tabcolsep) * \real{0.2500}}@{}}
\toprule\noalign{}
\begin{minipage}[b]{\linewidth}\raggedright
No.
\end{minipage} & \begin{minipage}[b]{\linewidth}\raggedright
Name
\end{minipage} & \begin{minipage}[b]{\linewidth}\raggedright
Description
\end{minipage} & \begin{minipage}[b]{\linewidth}\raggedright
Default
\end{minipage} \\
\midrule\noalign{}
\endhead
\bottomrule\noalign{}
\endlastfoot
Field & Metric & Laboratory & \\
1 & I1 & A positive integer that defines the lower bound of the grid in
the I-direction for which the boundary conditions are to be applied,
must be greater than or equal 1 and less than or equal to I2 and NX. &
1 \\
2 & I2 & A positive integer that defines the upper bound of the grid in
the I-direction for which the boundary conditions are to be applied,
must be greater than or equal to II and less than or equal to NX & NX \\
3 & J1 & A positive integer that defines the lower bound of the grid in
the J-direction for which the boundary conditions are to be applied,
must be greater than or equal 1 and less than or equal to J2 and NY. &
1 \\
4 & J2 & A positive integer that defines the upper bound of the grid in
the J-direction for which the boundary conditions are to be applied,
must be greater than or equal to JI and less than or equal to NY. &
NY \\
5 & K1 & A positive integer that defines the lower bound of the grid in
the K-direction for which the boundary conditions are to be applied,
must be greater than or equal to one and less than or equal to K2 and
NZ. & 1 \\
6 & K2 & A positive integer that defines the upper bound of the grid in
the K-direction for which the boundary conditions are to be applied,
must be greater than or equal to KI and less than or equal to NZ. &
NZ \\
7 & DIRECT & A character string that defines the direction to apply the
boundary conditions, and should be set to one of the following X, Y, or
Z for the positive direction, or X-, Y- or Z- for the negative
direction. & None \\
8 & TYPE & A defined character string that defines the type of boundary
condition to be applied, and should be set to one of the following
character strings: Only the FREE and RATE options are currently
supported; however the next release will support the DIRICHLET option. &
None \\
9 & PHASE & A defined character string that sets fluid type used in the
boundary calculations, and should be set to one of the following
character strings: & None \\
10 & RATE & A real value that defines the constant surface oil, gas or
water rate to be injected or withdrawn at the boundary, for when TYPE
has been set to RATE. Note a negative value implies an injection rate,
whereas, a positive value indicates a withdrawal rate. & 0,0 \\
Liquid: stb/day Gas: Mscf/day & Liquid: sm3/day Gas: sm3/day & Liquid:
scc/hr Gas: scc/hr & \\
11 & PRESS & PRESS is a real positive value that defines the constant
pressure boundary condition. PRESS should only be entered if TYPE has
been set to DIRICHLET. If the pressure at the boundary is less than
PRESS, then the fluid type declared via PHASE will flow across the
boundary. The default value of 1* will use the simulator's calculated
value based on data entered via the EQUIL keyword in the SOLUTION
section. & 1* \\
psia & barsa & atma & \\
12 & TEMP & TEMP is a real positive number that defines the constant
temperature boundary condition. TEMP should only be entered if TYPE has
been set to DIRICHLET. The default value of 1* will use the simulator's
calculated value based on data entered via one of the following
reservoir temperature keywords: RTEMP, RTEMPA, RTEMPVD, TEMPI, or
TEMPVD, in the SOLUTION section. Note that all of the aforementioned
reservoir temperature keywords, except for TEMPI, may also be used in
the PROPS section as well. & 1* \\
oF & oC & oC & \\
Notes: & & & \\
\end{longtable}

\emph{Table 10.10: BC Keyword Description}

See also the AQUFLUX - Define Constant Flux Analytical Aquifer keyword
that is supported by OPM Flow in both the SOLUTION and SCHEDULE
sections, to define the aquifer influx with greater flexibility.

If the BC keyword is not present in the input deck, then the boundary
conditions for the model are set to be no flow, which is the normal
behavior in both OPM Flow and the commercial simulator.

\subsubsection{Examples}\label{examples-56}

The first example shows how to set a constant pressure boundary using
TYPE equal to FREE:

\begin{verbatim}
--
--       DEFINE MODEL BOUNDARY CONDITIONS (OPM FLOW SOLUTION GRID KEYWORD)
--
--       ---------- BOX ---------   BC    BC         BC     BC    BC     BC
--       I1  I2   J1  J2   K1  K2   DIRC  TYPE       PHASE  RATE  PRESS  TEMP
BC
         1   1    1   1*   1   1*   X-    FREE       1*     1*    1*     1*    /
         1   1*   1   1    1   1*   Y     FREE                                 /
/
\end{verbatim}

With this option it is only necessary to define the boundary cells and
all the other parameters (PHASE, RATE, PRESS, and TEMP) can be
defaulted, as they are ignored for when TYPE equals FREE.

The next example is based on NX, NY and NZ equal to 20, 1, 10
respectively, on the DIMENS keyword in the RUNSPEC section, and shows
how different boundary types can be assigned to different parts of the
model.

\begin{verbatim}
--
--       DEFINE MODEL BOUNDARY CONDITIONS (OPM FLOW SOLUTION GRID KEYWORD)
--
--       ---------- BOX ---------   BC    BC         BC     BC    BC     BC
--       I1  I2   J1  J2   K1  K2   DIRC  TYPE       PHASE  RATE  PRESS  TEMP
BC
         1   1    1   1    1   10   X-    RATE       GAS    1*    256.0  100.0 /
         20  20   1   1    1   10   X     FREE       4*                        /
/
\end{verbatim}

The last example shows how the DIRICHLET boundary condition option may
be used:

\begin{verbatim}
--
--       DEFINE MODEL BOUNDARY CONDITIONS (OPM FLOW SOLUTION GRID KEYWORD)
--
--       ---------- BOX ---------   BC    BC         BC     BC    BC     BC
--       I1  I2   J1  J2   K1  K2   DIRC  TYPE       PHASE  RATE  PRESS  TEMP
BC
         1   1    1   1*   1   1*   X-    DIRICHLET  WAT    1*    256.0  100.0 /
         1   1*   1   1    1   1*   Y     DIRICHLET  WAT    1*    1*     100.0 /
/
\end{verbatim}

Here, the first line sets both the pressure and temperature at the
boundary, and the second line defaults the pressure entry, so that the
simulator calculated initial boundary pressure will be used.

\subsection{BOUNDARY -- Define a Boundary Box for
Printing}\label{boundary-define-a-boundary-box-for-printing-3}

The BOUNDARY keyword defines a rectangular grid for printing various
arrays to the output print file (*.PRT); thus, avoiding printing all the
elements in the selected array.

See BOUNDARY -- Define a Boundary Box for Printing in the GRID section
for a full description.

\subsection{BOX -- Define a Range of Grid Blocks to Enter Property
Data}\label{box-define-a-range-of-grid-blocks-to-enter-property-data-4}

BOX defines a range of grid blocks for which subsequent data is assigned
for all the cells in the defined BOX. Note that the BOX grid is reset by
the keyword ENDBOX by resetting the current defined BOX to be the whole
grid. The keyword can be used for any array and for all grid types.

See BOX - Define a Range of Grid Blocks to Enter Property Data in the
GRID section for a full description.

\subsection{COLUMNS -- Define Input File Column
Margins}\label{columns-define-input-file-column-margins-6}

The COLUMNS keyword defines the input file column margins; characters
outside the margins are ignored by the input parser.

See COLUMNS -- Define Input File Column Margins in the GLOBAL section
for a full description.

\subsection{COPY -- Copy Array Data to Another
Array}\label{copy-copy-array-data-to-another-array-4}

The COPY keyword copies an array (or part of an array) to another array
or part of an array. The arrays can be real or integer depending on the
array type; however, the arrays that can be operated on is dependent on
which section the COPY keyword is being used.

See COPY -- Copy Array Data to Another Array in the GRID section for a
full description.

\subsection{COPYREG -- Copy an Array to Another Array based on a Region
Number}\label{copyreg-copy-an-array-to-another-array-based-on-a-region-number-4}

The COPYREG keyword copies a specified array or part of an array based
on cells with a specific region number to another array. The region
number array can be FLUXNUM, MULTNUM or OPERNUM and these arrays must be
defined and be available before the COPYREG keyword is read by the
simulator. The property arrays can be real or integer depending on the
property array type; however, the property arrays that can be operated
on is dependent on which section the COPYREG keyword is being used.

See COPYREG -- Copy an Array to Another Array based on a Region Number
in the GRID section for a full description.

\subsection{DATUM -- Define the Datum Depth for the
Model}\label{datum-define-the-datum-depth-for-the-model}

\begin{longtable}[]{@{}
  >{\raggedright\arraybackslash}p{(\linewidth - 14\tabcolsep) * \real{0.1250}}
  >{\raggedright\arraybackslash}p{(\linewidth - 14\tabcolsep) * \real{0.1250}}
  >{\raggedright\arraybackslash}p{(\linewidth - 14\tabcolsep) * \real{0.1250}}
  >{\raggedright\arraybackslash}p{(\linewidth - 14\tabcolsep) * \real{0.1250}}
  >{\raggedright\arraybackslash}p{(\linewidth - 14\tabcolsep) * \real{0.1250}}
  >{\raggedright\arraybackslash}p{(\linewidth - 14\tabcolsep) * \real{0.1250}}
  >{\raggedright\arraybackslash}p{(\linewidth - 14\tabcolsep) * \real{0.1250}}
  >{\raggedright\arraybackslash}p{(\linewidth - 14\tabcolsep) * \real{0.1250}}@{}}
\toprule\noalign{}
\begin{minipage}[b]{\linewidth}\raggedright
RUNSPEC
\end{minipage} & \begin{minipage}[b]{\linewidth}\raggedright
GRID
\end{minipage} & \begin{minipage}[b]{\linewidth}\raggedright
EDIT
\end{minipage} & \begin{minipage}[b]{\linewidth}\raggedright
PROPS
\end{minipage} & \begin{minipage}[b]{\linewidth}\raggedright
REGIONS
\end{minipage} & \begin{minipage}[b]{\linewidth}\raggedright
SOLUTION
\end{minipage} & \begin{minipage}[b]{\linewidth}\raggedright
SUMMARY
\end{minipage} & \begin{minipage}[b]{\linewidth}\raggedright
SCHEDULE
\end{minipage} \\
\midrule\noalign{}
\endhead
\bottomrule\noalign{}
\endlastfoot
\end{longtable}

\subsubsection{Description}\label{description-776}

The DATUM keyword defines the datum depth for the model. This allows for
all grid block potentials (depth corrected pressures) to be calculated
at a common depth.

\begin{longtable}[]{@{}
  >{\raggedright\arraybackslash}p{(\linewidth - 6\tabcolsep) * \real{0.2500}}
  >{\raggedright\arraybackslash}p{(\linewidth - 6\tabcolsep) * \real{0.2500}}
  >{\raggedright\arraybackslash}p{(\linewidth - 6\tabcolsep) * \real{0.2500}}
  >{\raggedright\arraybackslash}p{(\linewidth - 6\tabcolsep) * \real{0.2500}}@{}}
\toprule\noalign{}
\begin{minipage}[b]{\linewidth}\raggedright
No.
\end{minipage} & \begin{minipage}[b]{\linewidth}\raggedright
Name
\end{minipage} & \begin{minipage}[b]{\linewidth}\raggedright
Description
\end{minipage} & \begin{minipage}[b]{\linewidth}\raggedright
Default
\end{minipage} \\
\midrule\noalign{}
\endhead
\bottomrule\noalign{}
\endlastfoot
Field & Metric & Laboratory & \\
1 & DATUM & DATUM is a single positive value that defines the datum
depth for the model. & None \\
feet & m & cm & \\
Notes: & & & \\
\end{longtable}

\emph{Table 10.11: DATUM Keyword Description}

See also the DATUMR and DATUMRX keywords in the SOLUTION section that
also define the datum depth for the model.

\subsubsection{Example}\label{example-440}

\begin{verbatim}
--       DATUM
--       DEPTH
--       ------
DATUM
         5000.0                                      / DATUM DEPTH FOR REPORTING
\end{verbatim}

The above example defines the datum for the model to be 5000.0

\subsection{DATUMR -- Define Datum Depths for the FIPNUM
Regions}\label{datumr-define-datum-depths-for-the-fipnum-regions}

\begin{longtable}[]{@{}
  >{\raggedright\arraybackslash}p{(\linewidth - 14\tabcolsep) * \real{0.1250}}
  >{\raggedright\arraybackslash}p{(\linewidth - 14\tabcolsep) * \real{0.1250}}
  >{\raggedright\arraybackslash}p{(\linewidth - 14\tabcolsep) * \real{0.1250}}
  >{\raggedright\arraybackslash}p{(\linewidth - 14\tabcolsep) * \real{0.1250}}
  >{\raggedright\arraybackslash}p{(\linewidth - 14\tabcolsep) * \real{0.1250}}
  >{\raggedright\arraybackslash}p{(\linewidth - 14\tabcolsep) * \real{0.1250}}
  >{\raggedright\arraybackslash}p{(\linewidth - 14\tabcolsep) * \real{0.1250}}
  >{\raggedright\arraybackslash}p{(\linewidth - 14\tabcolsep) * \real{0.1250}}@{}}
\toprule\noalign{}
\begin{minipage}[b]{\linewidth}\raggedright
RUNSPEC
\end{minipage} & \begin{minipage}[b]{\linewidth}\raggedright
GRID
\end{minipage} & \begin{minipage}[b]{\linewidth}\raggedright
EDIT
\end{minipage} & \begin{minipage}[b]{\linewidth}\raggedright
PROPS
\end{minipage} & \begin{minipage}[b]{\linewidth}\raggedright
REGIONS
\end{minipage} & \begin{minipage}[b]{\linewidth}\raggedright
SOLUTION
\end{minipage} & \begin{minipage}[b]{\linewidth}\raggedright
SUMMARY
\end{minipage} & \begin{minipage}[b]{\linewidth}\raggedright
SCHEDULE
\end{minipage} \\
\midrule\noalign{}
\endhead
\bottomrule\noalign{}
\endlastfoot
\end{longtable}

\subsubsection{Description}\label{description-777}

The DATUMR keyword defines the datum depth for each fluid in-place
region (FIPNUM) declared in the model. This allows for all grid block
potentials (depth corrected pressures) to be calculated at a common
depth within a FIPNUM region.

\begin{longtable}[]{@{}
  >{\raggedright\arraybackslash}p{(\linewidth - 6\tabcolsep) * \real{0.2500}}
  >{\raggedright\arraybackslash}p{(\linewidth - 6\tabcolsep) * \real{0.2500}}
  >{\raggedright\arraybackslash}p{(\linewidth - 6\tabcolsep) * \real{0.2500}}
  >{\raggedright\arraybackslash}p{(\linewidth - 6\tabcolsep) * \real{0.2500}}@{}}
\toprule\noalign{}
\begin{minipage}[b]{\linewidth}\raggedright
No.
\end{minipage} & \begin{minipage}[b]{\linewidth}\raggedright
Name
\end{minipage} & \begin{minipage}[b]{\linewidth}\raggedright
Description
\end{minipage} & \begin{minipage}[b]{\linewidth}\raggedright
Default
\end{minipage} \\
\midrule\noalign{}
\endhead
\bottomrule\noalign{}
\endlastfoot
Field & Metric & Laboratory & \\
1 & DATUMR & DATUMR is a vector of positive values that defines the
datum depth for each fluid in-place region. & None \\
feet & m & cm & \\
Notes: & & & \\
\end{longtable}

\emph{Table 10.12: DATUMR Keyword Description}

See also the DATUM and DATUMRX keywords in the SOLUTION section that
also define the datum depth for the model.

\subsubsection{Example}\label{example-441}

\begin{verbatim}
--
--       DATUM
--       DEPTH
--       ------
DATUMR
         4800.0
         4900.0
         5000.0                                      / DATUM DEPTH FOR REPORTING
\end{verbatim}

The above example defines the datum depth for three FIPNUM regions, for
when NTFIP has been set equal to three on the REGDIMS keyword in the
RUNSPEC section.

\subsection{DATUMRX -- Define Datum Depths for the FIP Allocated
Regions}\label{datumrx-define-datum-depths-for-the-fip-allocated-regions}

\begin{longtable}[]{@{}
  >{\raggedright\arraybackslash}p{(\linewidth - 14\tabcolsep) * \real{0.1250}}
  >{\raggedright\arraybackslash}p{(\linewidth - 14\tabcolsep) * \real{0.1250}}
  >{\raggedright\arraybackslash}p{(\linewidth - 14\tabcolsep) * \real{0.1250}}
  >{\raggedright\arraybackslash}p{(\linewidth - 14\tabcolsep) * \real{0.1250}}
  >{\raggedright\arraybackslash}p{(\linewidth - 14\tabcolsep) * \real{0.1250}}
  >{\raggedright\arraybackslash}p{(\linewidth - 14\tabcolsep) * \real{0.1250}}
  >{\raggedright\arraybackslash}p{(\linewidth - 14\tabcolsep) * \real{0.1250}}
  >{\raggedright\arraybackslash}p{(\linewidth - 14\tabcolsep) * \real{0.1250}}@{}}
\toprule\noalign{}
\begin{minipage}[b]{\linewidth}\raggedright
RUNSPEC
\end{minipage} & \begin{minipage}[b]{\linewidth}\raggedright
GRID
\end{minipage} & \begin{minipage}[b]{\linewidth}\raggedright
EDIT
\end{minipage} & \begin{minipage}[b]{\linewidth}\raggedright
PROPS
\end{minipage} & \begin{minipage}[b]{\linewidth}\raggedright
REGIONS
\end{minipage} & \begin{minipage}[b]{\linewidth}\raggedright
SOLUTION
\end{minipage} & \begin{minipage}[b]{\linewidth}\raggedright
SUMMARY
\end{minipage} & \begin{minipage}[b]{\linewidth}\raggedright
SCHEDULE
\end{minipage} \\
\midrule\noalign{}
\endhead
\bottomrule\noalign{}
\endlastfoot
\end{longtable}

\subsubsection{Description}\label{description-778}

The DATUMRX keyword defines the datum depth for each fluid in-place
family region defined by the FIP keyword. This allows for all grid block
potentials (depth corrected pressures) to be calculated at a common
depth within a given FIP region. The FIP keyword in the REGIONS section
allows one to define additional sets of fluid in-place regions to the
standard FIPNUM keyword. For example, one could use FIPNUM to define the
reservoir layers as fluid in-place regions and the FIP keyword to define
the fluid in-place region for fault blocks.

\begin{longtable}[]{@{}
  >{\raggedright\arraybackslash}p{(\linewidth - 6\tabcolsep) * \real{0.2500}}
  >{\raggedright\arraybackslash}p{(\linewidth - 6\tabcolsep) * \real{0.2500}}
  >{\raggedright\arraybackslash}p{(\linewidth - 6\tabcolsep) * \real{0.2500}}
  >{\raggedright\arraybackslash}p{(\linewidth - 6\tabcolsep) * \real{0.2500}}@{}}
\toprule\noalign{}
\begin{minipage}[b]{\linewidth}\raggedright
No.
\end{minipage} & \begin{minipage}[b]{\linewidth}\raggedright
Name
\end{minipage} & \begin{minipage}[b]{\linewidth}\raggedright
Description
\end{minipage} & \begin{minipage}[b]{\linewidth}\raggedright
Default
\end{minipage} \\
\midrule\noalign{}
\endhead
\bottomrule\noalign{}
\endlastfoot
Field & Metric & Laboratory & \\
1 & FIPNAME & A character string of up to five characters in length that
defines the FIP family name for which the datum depth data is being
defined. The default value of 1* will set DATUMR to the standard FIPNUM
region numbers. & 1* \\
2 & DATUMR & DATUMR is a vector of positive values that defines the
datum depth for each fluid in-place family region. There must be one
entry for each region in the FIP family name. A maximum of NTFIP values,
as declared by the REGDIMS keyword in the RUNSPEC section, may be
entered for each FIPNAME entry. & None \\
feet & m & cm & \\
Notes: & & & \\
\end{longtable}

\emph{Table 10.13: DATUMR Keyword Description}

See also the FIP keyword in the REGIONS section to define FIP family
regions, and the DATUM and DATUMR keywords in the SOLUTION section that
also define the datum depth for the model.

\subsubsection{Example}\label{example-442}

\begin{verbatim}
--
--       FIP       DATUM
--       NAME      DEPTH
DATUMRX
         'FLTBL'   5000.0  5000.0  5000.0  5000.0    / DATUM DEPTH FOR REPORTING
         'LICBL'   5000.0  5050.0                    / DATUM DEPTH FOR REPORTING
/
\end{verbatim}

The above example defines the datum depth for two FIP families, FLTBL
and LICBL, with the datum set to a constant 5000.0 psia for FLTBL family
and different values for each of the regions in the LICBL family of
regions.

\subsection{DEBUG -- Define the Debug Data to be Printed to
File}\label{debug-define-the-debug-data-to-be-printed-to-file-6}

This keyword defines the debug data to be written to the debug file
(*.DBG). This keyword is not supported by OPM Flow but has no effect on
the results so it will be ignored.

See DEBUG -- Define the Debug Data to be Printed to File in the GLOBAL
section for a full description.

\subsection{DYNAMICR -- Start of Dynamic Region Parameter
Definition}\label{dynamicr-start-of-dynamic-region-parameter-definition}

\begin{longtable}[]{@{}
  >{\raggedright\arraybackslash}p{(\linewidth - 14\tabcolsep) * \real{0.1250}}
  >{\raggedright\arraybackslash}p{(\linewidth - 14\tabcolsep) * \real{0.1250}}
  >{\raggedright\arraybackslash}p{(\linewidth - 14\tabcolsep) * \real{0.1250}}
  >{\raggedright\arraybackslash}p{(\linewidth - 14\tabcolsep) * \real{0.1250}}
  >{\raggedright\arraybackslash}p{(\linewidth - 14\tabcolsep) * \real{0.1250}}
  >{\raggedright\arraybackslash}p{(\linewidth - 14\tabcolsep) * \real{0.1250}}
  >{\raggedright\arraybackslash}p{(\linewidth - 14\tabcolsep) * \real{0.1250}}
  >{\raggedright\arraybackslash}p{(\linewidth - 14\tabcolsep) * \real{0.1250}}@{}}
\toprule\noalign{}
\begin{minipage}[b]{\linewidth}\raggedright
RUNSPEC
\end{minipage} & \begin{minipage}[b]{\linewidth}\raggedright
GRID
\end{minipage} & \begin{minipage}[b]{\linewidth}\raggedright
EDIT
\end{minipage} & \begin{minipage}[b]{\linewidth}\raggedright
PROPS
\end{minipage} & \begin{minipage}[b]{\linewidth}\raggedright
REGIONS
\end{minipage} & \begin{minipage}[b]{\linewidth}\raggedright
SOLUTION
\end{minipage} & \begin{minipage}[b]{\linewidth}\raggedright
SUMMARY
\end{minipage} & \begin{minipage}[b]{\linewidth}\raggedright
SCHEDULE
\end{minipage} \\
\midrule\noalign{}
\endhead
\bottomrule\noalign{}
\endlastfoot
\end{longtable}

\subsubsection{Description}\label{description-779}

The DYNAMICR keyword marks the start of a Dynamic Region section and
defines the parameters used for Dynamic Regions that allows for property
and reporting regions to vary as the run progresses, based on the
parameters and logic defined by this keyword and section. A Dynamic
Region section is terminated by the ENDDYN keyword in the SOLUTION or
SCHEDULE sections.

This keyword is not supported by OPM Flow but would change the results
if supported so the simulation will be stopped.

\subsection{ECHO -- Activate Echoing of User Input Files to the Print
File}\label{echo-activate-echoing-of-user-input-files-to-the-print-file-6}

Turns on echoing of all the input files to the print file; note that
this keyword is activated by default and can subsequently be switched
off by the NOECHO activation keyword.

See ECHO -- Activate Echoing of User Input Files to the Print File in
the GLOBAL section for a full description.

\subsection{END -- Define the End of the Input
File}\label{end-define-the-end-of-the-input-file-6}

This keyword marks the end of the input file and can occur in any
section. Any keywords and data after the END keyword are ignored.

See END -- Define the End of the Input File in the GLOBAL section for a
full description.

\subsection{ENDBOX -- Define the End of the BOX Defined
Grid}\label{endbox-define-the-end-of-the-box-defined-grid-4}

This keyword marks the end of a previously defined BOX sub-grid as
defined by a previously entered BOX keyword. The keyword resets the
input grid to be the full grid as defined by the NX, NY, and NZ
variables on the DIMENS keyword in the RUNSPEC section.

See ENDBOX -- Define the End of the BOX Defined Grid in the GRID section
for a full description.

\subsection{ENDDYN -- End of Dynamic Region Parameter
Definition}\label{enddyn-end-of-dynamic-region-parameter-definition}

\begin{longtable}[]{@{}
  >{\raggedright\arraybackslash}p{(\linewidth - 14\tabcolsep) * \real{0.1250}}
  >{\raggedright\arraybackslash}p{(\linewidth - 14\tabcolsep) * \real{0.1250}}
  >{\raggedright\arraybackslash}p{(\linewidth - 14\tabcolsep) * \real{0.1250}}
  >{\raggedright\arraybackslash}p{(\linewidth - 14\tabcolsep) * \real{0.1250}}
  >{\raggedright\arraybackslash}p{(\linewidth - 14\tabcolsep) * \real{0.1250}}
  >{\raggedright\arraybackslash}p{(\linewidth - 14\tabcolsep) * \real{0.1250}}
  >{\raggedright\arraybackslash}p{(\linewidth - 14\tabcolsep) * \real{0.1250}}
  >{\raggedright\arraybackslash}p{(\linewidth - 14\tabcolsep) * \real{0.1250}}@{}}
\toprule\noalign{}
\begin{minipage}[b]{\linewidth}\raggedright
RUNSPEC
\end{minipage} & \begin{minipage}[b]{\linewidth}\raggedright
GRID
\end{minipage} & \begin{minipage}[b]{\linewidth}\raggedright
EDIT
\end{minipage} & \begin{minipage}[b]{\linewidth}\raggedright
PROPS
\end{minipage} & \begin{minipage}[b]{\linewidth}\raggedright
REGIONS
\end{minipage} & \begin{minipage}[b]{\linewidth}\raggedright
SOLUTION
\end{minipage} & \begin{minipage}[b]{\linewidth}\raggedright
SUMMARY
\end{minipage} & \begin{minipage}[b]{\linewidth}\raggedright
SCHEDULE
\end{minipage} \\
\midrule\noalign{}
\endhead
\bottomrule\noalign{}
\endlastfoot
\end{longtable}

\subsubsection{Description}\label{description-780}

The ENDDYN keyword marks the end of a Dynamic Region section that was
started with the DYNAMICR keyword in the SOLUTION or SCHEDULE sections.
Dynamic Regions allow for property and reporting regions to vary as the
run progresses, based on the parameters and logic defined within the
section.

This keyword is not supported by OPM Flow but would change the results
if supported so the simulation will be stopped.

\subsection{ENDFIN -- End the Definition of a Local Grid
Refinement}\label{endfin-end-the-definition-of-a-local-grid-refinement-4}

ENDFIN defines the end of a Cartesian or radial local grid refinement
(``LGR'') definition and a LGR property definition data set.

See ENDFIN -- End the Definition of a Local Grid Refinement in the GRID
section for a full description.

\subsection{ENDINC -- Define the End of an Include
File}\label{endinc-define-the-end-of-an-include-file-6}

This keyword marks the end of an include file specified on the INCLUDE
keyword. When the ENDINC keyword is encountered in the INCLUDE file,
input data is read from the next keyword in the current file. Any
keywords and data after the ENDINC keyword in the INCLUDE file are
ignored.

See ENDINC -- Define the End of an Include File in the GLOBAL section
for a full description.

\subsection{ENDSKIP -- Deactivate Skipping of Keywords and Input
Data}\label{endskip-deactivate-skipping-of-keywords-and-input-data-6}

The ENDSKIP keyword deactivates the skipping of keywords that was
activated by the SKIP, SKIP100, or SKIP300 keywords. Each SKIP keyword
should be paired with an ENDSKIP keyword.

See ENDSKIP -- Deactivate Skipping of Keywords and Input Data in the
GLOBAL section for a full description.

\subsection{EQUALREG -- Sets an Array to a Constant by Region
Number}\label{equalreg-sets-an-array-to-a-constant-by-region-number-4}

The EQUALREG keyword sets a specified array to a constant for cells with
a specific region number. The region number array can be FLUXNUM,
MULTNUM or OPERNUM and these arrays must be defined and be available
before the EQUALREG keyword is read by the simulator. The constant can
be real or integer depending on the property array type; however, the
property arrays that can be operated on is dependent on which section
the EQUALREG keyword is being used.

See EQUALREG -- Sets an Array to a Constant by Region Number in the GRID
section for a full description.

\subsection{EQUALS -- Sets a Specified Array to a
Constant}\label{equals-sets-a-specified-array-to-a-constant-4}

The EQUALS keyword sets a specified array or part of an array to a
constant. The constant can be real or integer depending on the array
type; however, the arrays that can be operated on is dependent on which
section the EQUALS keyword is being used.

See EQUALS -- Sets a Specified Array to a Constant in the GRID section
for a full description.

\subsection{EQUIL -- Define the Equilibration Initialization
Data}\label{equil-define-the-equilibration-initialization-data}

\begin{longtable}[]{@{}
  >{\raggedright\arraybackslash}p{(\linewidth - 14\tabcolsep) * \real{0.1250}}
  >{\raggedright\arraybackslash}p{(\linewidth - 14\tabcolsep) * \real{0.1250}}
  >{\raggedright\arraybackslash}p{(\linewidth - 14\tabcolsep) * \real{0.1250}}
  >{\raggedright\arraybackslash}p{(\linewidth - 14\tabcolsep) * \real{0.1250}}
  >{\raggedright\arraybackslash}p{(\linewidth - 14\tabcolsep) * \real{0.1250}}
  >{\raggedright\arraybackslash}p{(\linewidth - 14\tabcolsep) * \real{0.1250}}
  >{\raggedright\arraybackslash}p{(\linewidth - 14\tabcolsep) * \real{0.1250}}
  >{\raggedright\arraybackslash}p{(\linewidth - 14\tabcolsep) * \real{0.1250}}@{}}
\toprule\noalign{}
\begin{minipage}[b]{\linewidth}\raggedright
RUNSPEC
\end{minipage} & \begin{minipage}[b]{\linewidth}\raggedright
GRID
\end{minipage} & \begin{minipage}[b]{\linewidth}\raggedright
EDIT
\end{minipage} & \begin{minipage}[b]{\linewidth}\raggedright
PROPS
\end{minipage} & \begin{minipage}[b]{\linewidth}\raggedright
REGIONS
\end{minipage} & \begin{minipage}[b]{\linewidth}\raggedright
SOLUTION
\end{minipage} & \begin{minipage}[b]{\linewidth}\raggedright
SUMMARY
\end{minipage} & \begin{minipage}[b]{\linewidth}\raggedright
SCHEDULE
\end{minipage} \\
\midrule\noalign{}
\endhead
\bottomrule\noalign{}
\endlastfoot
\end{longtable}

\subsubsection{Description}\label{description-781}

This keyword defines the parameters used to initialize the model for
when equilibration is calculated by OPM Flow. This is the standard
methodology to initialize a model, the non-standard formulation of
entering the pressures and saturations for each grid cell is seldom
employed in the industry. The keyword can be used with all grid types.

\begin{longtable}[]{@{}
  >{\raggedright\arraybackslash}p{(\linewidth - 6\tabcolsep) * \real{0.2500}}
  >{\raggedright\arraybackslash}p{(\linewidth - 6\tabcolsep) * \real{0.2500}}
  >{\raggedright\arraybackslash}p{(\linewidth - 6\tabcolsep) * \real{0.2500}}
  >{\raggedright\arraybackslash}p{(\linewidth - 6\tabcolsep) * \real{0.2500}}@{}}
\toprule\noalign{}
\begin{minipage}[b]{\linewidth}\raggedright
No.
\end{minipage} & \begin{minipage}[b]{\linewidth}\raggedright
Name
\end{minipage} & \begin{minipage}[b]{\linewidth}\raggedright
Description
\end{minipage} & \begin{minipage}[b]{\linewidth}\raggedright
Default
\end{minipage} \\
\midrule\noalign{}
\endhead
\bottomrule\noalign{}
\endlastfoot
Field & Metric & Laboratory & \\
1 & DATUM & DATUM is a single positive value that defines the reference
datum depth for PRESS. & 0.0 \\
feet & m & cm & \\
2 & PRESS & PRESS is a single positive value that defines the pressure
at DATUM. If the DATUM depth lies above the GOC then PRESS is the
pressure with respect to the gas phase. If the DATUM depth is below OWC
then PRESS refers to the water phase pressure. Otherwise, PRESS refers
to the oil phase pressure. & 0.0 \\
psia & barsa & atma & \\
3 & WATCONT & & 0.0 \\
feet & m & cm & \\
4 & WATCAP & & 0.0 \\
psia & barsa & atma & \\
5 & GASCONT & & 0.0 \\
feet & m & cm & \\
6 & GASCAP & & 0.0 \\
psia & barsa & atma & \\
7 & EQLOPT1 & EQLOPT1 is an integer value that sets the initialization
option for when dissolved gas is present in the run, as activated by the
DISGAS keyword in the RUNSPEC section. EQLOPT1 is ignored if there is no
dissolved gas in the run. & 0 \\
dimensionless & dimensionless & dimensionless & \\
8 & EQLOPT2 & EQLOPT2 is an integer value that sets the initialization
option for when vaporized oil (condensate) is present in the run, as
activated by the VAPOIL keyword in the RUNSPEC section. EQLOPT2 is
ignored if there is no vaporized oil in the run. & 0 \\
dimensionless & dimensionless & dimensionless & \\
9 & EQLOPT3 & EQLOPT3 is an integer value that sets the initialization
accuracy options for the equilibration calculation. Increasing the value
of N increases the accuracy of the calculation, with the maximum value
of N being set to 20 by OPM Flow. Note this option should be used with
Irregular Corner-Point Grids. EQLOPT3 is ignored for Radial Grids. &
0 \\
dimensionless & dimensionless & dimensionless & \\
10 & EQLOPT4 & A positive integer value greater than or equal one and
less than or equal to three, that sets the initialization option in the
commercial compositional simulator. EQLOPT4 should be defaulted with 1*,
as it is not used by OPM Flow. & None \\
dimensionless & dimensionless & dimensionless & \\
11 & EQLOPT5 & A positive integer value that if set to one forces PRESS
to be used for the datum pressure in the commercial compositional
simulator. EQLOPT5 should be defaulted with either 1*, as it is not used
by OPM Flow. & None \\
dimensionless & dimensionless & dimensionless & \\
12 & EQLOPT6 & EQLOPT6 is an integer value that sets the initialization
option for when vaporized water is present in the run, as activated by
the VAPWAT keyword in the RUNSPEC section. Note this is an OPM Flow
specific parameter for use with simulator's Vaporized Water Model. Note
that the allocation of multiple RVWVD tables to each grid cell is
through the EQLNUM keyword and not the PVTNUM keyword. EQLOPT6 is
ignored if there is no vaporized water in the run. & 0 \\
dimensionless & dimensionless & dimensionless & \\
Notes: & & & \\
\end{longtable}

\emph{Table 10.14: EQUIL Keyword Description}

\begin{tcolorbox}[enhanced jigsaw, title=\textcolor{quarto-callout-note-color}{\faInfo}\hspace{0.5em}{Note}, coltitle=black, breakable, toprule=.15mm, colbacktitle=quarto-callout-note-color!10!white, leftrule=.75mm, opacitybacktitle=0.6, colback=white, bottomtitle=1mm, toptitle=1mm, titlerule=0mm, left=2mm, colframe=quarto-callout-note-color-frame, arc=.35mm, rightrule=.15mm, bottomrule=.15mm, opacityback=0]

A common method to initialize a model is by using the SWATINIT property
array to set the initial water saturation for each cell in the model.
This property is normally exported from a static model, where Saturation
Height Functions (``SHF'') have been used to describe the water
saturation profile with depth. In the dynamic model capillary pressure
functions are used to described the water profile versus depth. Note
that if the SWATINIT array has been used to initialize the model then
the fine grid block initialization via the EQLOPT3 variable, should not
normally be used, and should be defaulted or set equal to zero;
otherwise, the resulting water saturation will not strictly honor the
SWATINIT array.

\end{tcolorbox}

See also the PRESSURE, SGAS, SOIL and SWAT keywords in the SOLUTION
section to initialize the model using the non-standard formulation of
entering the pressures and saturations for each grid cell.

\subsubsection{Example}\label{example-443}

\begin{verbatim}
--
--       DATUM   DATUM   OWC     PCOW   GOC    PCGO   RS   RV   N    E300  RVW
--       DEPTH   PRESS   DEPTH   ----   DEPTH  ----   OPT  OPT  OPT  OPTS  OPT
EQUIL
         3650.0  1560.0  3712.0  0.00  1000.0  0.00   1    0    -5   2*    1*   /
         3650.0  1560.0  3741.0  0.00  1000.0  0.00   1    0    -5   2*    1*   /
         3650.0  1560.0  3741.0  0.00  1000.0  0.00   1    0    -5   2*    1*   /
\end{verbatim}

The above example defines three equilibration records for when NTEQUL
equals three on the EQLDIMS keyword in the RUNSPEC section. Here there
is no gas cap and the GOC has been set to a value above the reservoirs
(1000.0), and the value of EQLOPT3 (-5) has been explicitly stated.

\subsection{EXTRAPMS -- Activate Extrapolation Warning
Messages}\label{extrapms-activate-extrapolation-warning-messages-6}

The EXTRAPMS keyword activates extrapolation warning messages for when
OPM Flow extrapolates the PVT or VFP tables. Frequent extrapolation
warning messages should be investigated and resolved as this would
indicate possible incorrect data and may result in the simulator
extrapolating to unrealistic values.

See EXTRAPMS -- Activate Extrapolation Warning Messages in the GLOBAL
section for a full description.

\subsection{FILEUNIT -- Activate Unit Consistency
Verification}\label{fileunit-activate-unit-consistency-verification-3}

The FILEUNIT keyword defines the units of the data set, and is used to
verify that the units in the input deck and any associated include files
are consistent. The keyword does not provide for the conversion between
different sets of units.

See FILEUNIT -- Activate Unit Consistency Checking in the GRID section
for a full description.

\subsection{FORMFEED -- Defined the Print File Form-Feed
Character}\label{formfeed-defined-the-print-file-form-feed-character-6}

The FORMFEED keyword defines the form-feed character, or carriage
control character, for the output print (\emph{.PRT) run summary (}.RSM)
files. The keyword should be place at the very top of the input file.

See FORMFEED -- Defined the Print File Form-Feed Character in the GLOBAL
section for a dull description.

\subsection{GASCONC -- Define the Initial Equilibration Coal Gas
Concentration for All Grid
Blocks}\label{gasconc-define-the-initial-equilibration-coal-gas-concentration-for-all-grid-blocks}

\begin{longtable}[]{@{}
  >{\raggedright\arraybackslash}p{(\linewidth - 14\tabcolsep) * \real{0.1250}}
  >{\raggedright\arraybackslash}p{(\linewidth - 14\tabcolsep) * \real{0.1250}}
  >{\raggedright\arraybackslash}p{(\linewidth - 14\tabcolsep) * \real{0.1250}}
  >{\raggedright\arraybackslash}p{(\linewidth - 14\tabcolsep) * \real{0.1250}}
  >{\raggedright\arraybackslash}p{(\linewidth - 14\tabcolsep) * \real{0.1250}}
  >{\raggedright\arraybackslash}p{(\linewidth - 14\tabcolsep) * \real{0.1250}}
  >{\raggedright\arraybackslash}p{(\linewidth - 14\tabcolsep) * \real{0.1250}}
  >{\raggedright\arraybackslash}p{(\linewidth - 14\tabcolsep) * \real{0.1250}}@{}}
\toprule\noalign{}
\begin{minipage}[b]{\linewidth}\raggedright
RUNSPEC
\end{minipage} & \begin{minipage}[b]{\linewidth}\raggedright
GRID
\end{minipage} & \begin{minipage}[b]{\linewidth}\raggedright
EDIT
\end{minipage} & \begin{minipage}[b]{\linewidth}\raggedright
PROPS
\end{minipage} & \begin{minipage}[b]{\linewidth}\raggedright
REGIONS
\end{minipage} & \begin{minipage}[b]{\linewidth}\raggedright
SOLUTION
\end{minipage} & \begin{minipage}[b]{\linewidth}\raggedright
SUMMARY
\end{minipage} & \begin{minipage}[b]{\linewidth}\raggedright
SCHEDULE
\end{minipage} \\
\midrule\noalign{}
\endhead
\bottomrule\noalign{}
\endlastfoot
\end{longtable}

\subsubsection{Description}\label{description-782}

The GASCONC keyword defines the initial equilibration coal gas
concentration values for all matrix grid cells in the model and should
be used in conjunction with the GCVD keyword in the SOLUTION section, to
fully describe the initial state of the model. The keyword should only
be used if the coal phase has been activated in the model via the COAL
keyword in the RUNSPEC section. Note both GASCONC and GCVD are optional
as the simulator will calculate the coal gas concentration based on the
equilibrium concentration and the block pressure.

This keyword is not supported by OPM Flow but would change the results
if supported so the simulation will be stopped.

This is the non-standard method to initialize the model via enumeration
and is seldom employed in the industry. The standard methodology is for
OPM Flow to initialize a model using the parameters on the EQUIL keyword
combined with other keywords to fully describe the initial state of the
model. The keyword can be used with all grid types.

\begin{longtable}[]{@{}
  >{\raggedright\arraybackslash}p{(\linewidth - 6\tabcolsep) * \real{0.2500}}
  >{\raggedright\arraybackslash}p{(\linewidth - 6\tabcolsep) * \real{0.2500}}
  >{\raggedright\arraybackslash}p{(\linewidth - 6\tabcolsep) * \real{0.2500}}
  >{\raggedright\arraybackslash}p{(\linewidth - 6\tabcolsep) * \real{0.2500}}@{}}
\toprule\noalign{}
\begin{minipage}[b]{\linewidth}\raggedright
No.
\end{minipage} & \begin{minipage}[b]{\linewidth}\raggedright
Name
\end{minipage} & \begin{minipage}[b]{\linewidth}\raggedright
Description
\end{minipage} & \begin{minipage}[b]{\linewidth}\raggedright
Default
\end{minipage} \\
\midrule\noalign{}
\endhead
\bottomrule\noalign{}
\endlastfoot
Field & Metric & Laboratory & \\
1 & GASCONC & GASCONC is an array of real positive numbers that define
the initial equilibration coal gas concentration values to each matrix
cell in the model. Repeat counts may be used, for example 20*75.0. &
None \\
Mscf/ft3 & sm3/m3 & scc/cc & \\
Notes: & & & \\
\end{longtable}

\emph{Table 10.15: GASCONC Keyword Description}

See also the GCVD keyword in the SOLUTION section to fully define the
initial state of the model.

\subsubsection{Example}\label{example-444}

\begin{verbatim}
--
--       DEFINE INITIAL EQUILIBRATION COAL GAS CONCENTRATION FOR ALL CELLS
--       BASED ON NX = 100, NY = 100 AND NZ = 6
--
GASCONC
         1000*75.500    1000*65.500    1000*60.000                             /
\end{verbatim}

The above example defines the initial equilibration coal gas
concentration values to be 75.500 for all the matrix cells in the first
layer, 65.500 for all the cells in the second layer, and finally 60.000
for all the cells in the third layer.

\subsection{GASSATC -- Define the Initial Equilibration Saturated Coal
Gas Concentration for All Grid
Blocks}\label{gassatc-define-the-initial-equilibration-saturated-coal-gas-concentration-for-all-grid-blocks}

\begin{longtable}[]{@{}
  >{\raggedright\arraybackslash}p{(\linewidth - 14\tabcolsep) * \real{0.1250}}
  >{\raggedright\arraybackslash}p{(\linewidth - 14\tabcolsep) * \real{0.1250}}
  >{\raggedright\arraybackslash}p{(\linewidth - 14\tabcolsep) * \real{0.1250}}
  >{\raggedright\arraybackslash}p{(\linewidth - 14\tabcolsep) * \real{0.1250}}
  >{\raggedright\arraybackslash}p{(\linewidth - 14\tabcolsep) * \real{0.1250}}
  >{\raggedright\arraybackslash}p{(\linewidth - 14\tabcolsep) * \real{0.1250}}
  >{\raggedright\arraybackslash}p{(\linewidth - 14\tabcolsep) * \real{0.1250}}
  >{\raggedright\arraybackslash}p{(\linewidth - 14\tabcolsep) * \real{0.1250}}@{}}
\toprule\noalign{}
\begin{minipage}[b]{\linewidth}\raggedright
RUNSPEC
\end{minipage} & \begin{minipage}[b]{\linewidth}\raggedright
GRID
\end{minipage} & \begin{minipage}[b]{\linewidth}\raggedright
EDIT
\end{minipage} & \begin{minipage}[b]{\linewidth}\raggedright
PROPS
\end{minipage} & \begin{minipage}[b]{\linewidth}\raggedright
REGIONS
\end{minipage} & \begin{minipage}[b]{\linewidth}\raggedright
SOLUTION
\end{minipage} & \begin{minipage}[b]{\linewidth}\raggedright
SUMMARY
\end{minipage} & \begin{minipage}[b]{\linewidth}\raggedright
SCHEDULE
\end{minipage} \\
\midrule\noalign{}
\endhead
\bottomrule\noalign{}
\endlastfoot
\end{longtable}

\subsubsection{Description}\label{description-783}

The GASSATC keyword defines the initial equilibration saturated coal gas
concentration values for all grid cells in the model. The keyword should
only be used if the coal phase has been activated in the model via the
COAL keyword in the RUNSPEC section. The keyword is used to re-scale the
Langmuir isotherms entered via the LANGMUIR keyword in the PROPS
section, in conjunction with a matrix grid blocks initial reservoir
pressure. The keyword is optional, and if absent from the input file,
the matrix grid block Langmuir isotherm is left unscaled.

This keyword is not supported by OPM Flow but would change the results
if supported so the simulation will be stopped.

This is the non-standard method to initialize the model via enumeration
and is seldom employed in the industry. The standard methodology is for
OPM Flow to initialize a model using the parameters on the EQUIL keyword
combined with other keywords to fully describe the initial state of the
model. The keyword can be used with all grid types.

\begin{longtable}[]{@{}
  >{\raggedright\arraybackslash}p{(\linewidth - 6\tabcolsep) * \real{0.2500}}
  >{\raggedright\arraybackslash}p{(\linewidth - 6\tabcolsep) * \real{0.2500}}
  >{\raggedright\arraybackslash}p{(\linewidth - 6\tabcolsep) * \real{0.2500}}
  >{\raggedright\arraybackslash}p{(\linewidth - 6\tabcolsep) * \real{0.2500}}@{}}
\toprule\noalign{}
\begin{minipage}[b]{\linewidth}\raggedright
No.
\end{minipage} & \begin{minipage}[b]{\linewidth}\raggedright
Name
\end{minipage} & \begin{minipage}[b]{\linewidth}\raggedright
Description
\end{minipage} & \begin{minipage}[b]{\linewidth}\raggedright
Default
\end{minipage} \\
\midrule\noalign{}
\endhead
\bottomrule\noalign{}
\endlastfoot
Field & Metric & Laboratory & \\
1 & GASSATC & GASSATC is an array of real positive numbers that define
the initial equilibration saturated coal gas concentration values to
each cell in the model. Repeat counts may be used, for example 20*75.0.
& None \\
Mscf/ft3 & sm3/m3 & scc/cc & \\
Notes: & & & \\
\end{longtable}

\emph{Table 10.16: GASSATC Keyword Description}

See also the GASCONC and the GCVD keywords in the SOLUTION section to
fully define the initial state of the model.

\subsubsection{Example}\label{example-445}

\begin{verbatim}
--
--       DEFINE INITIAL EQUILIBRATION SAT COAL GAS CONCENTRATION ALL CELLS MODEL
--       BASED ON NX = 100, NY = 100 AND NZ = 6
--
GASSATC
         1000*75.500    1000*65.500    1000*60.000                             /
\end{verbatim}

The above example defines the initial equilibration saturated coal gas
concentration values to be 75.500 for all the matrix cells in the first
layer, 65.500 for all the cells in the second layer, and finally 60.000
for all the cells in the third layer.

\subsection{GCVD -- Define Equilibration Coal Gas Concentration versus
Depth
Tables}\label{gcvd-define-equilibration-coal-gas-concentration-versus-depth-tables}

\begin{longtable}[]{@{}
  >{\raggedright\arraybackslash}p{(\linewidth - 14\tabcolsep) * \real{0.1250}}
  >{\raggedright\arraybackslash}p{(\linewidth - 14\tabcolsep) * \real{0.1250}}
  >{\raggedright\arraybackslash}p{(\linewidth - 14\tabcolsep) * \real{0.1250}}
  >{\raggedright\arraybackslash}p{(\linewidth - 14\tabcolsep) * \real{0.1250}}
  >{\raggedright\arraybackslash}p{(\linewidth - 14\tabcolsep) * \real{0.1250}}
  >{\raggedright\arraybackslash}p{(\linewidth - 14\tabcolsep) * \real{0.1250}}
  >{\raggedright\arraybackslash}p{(\linewidth - 14\tabcolsep) * \real{0.1250}}
  >{\raggedright\arraybackslash}p{(\linewidth - 14\tabcolsep) * \real{0.1250}}@{}}
\toprule\noalign{}
\begin{minipage}[b]{\linewidth}\raggedright
RUNSPEC
\end{minipage} & \begin{minipage}[b]{\linewidth}\raggedright
GRID
\end{minipage} & \begin{minipage}[b]{\linewidth}\raggedright
EDIT
\end{minipage} & \begin{minipage}[b]{\linewidth}\raggedright
PROPS
\end{minipage} & \begin{minipage}[b]{\linewidth}\raggedright
REGIONS
\end{minipage} & \begin{minipage}[b]{\linewidth}\raggedright
SOLUTION
\end{minipage} & \begin{minipage}[b]{\linewidth}\raggedright
SUMMARY
\end{minipage} & \begin{minipage}[b]{\linewidth}\raggedright
SCHEDULE
\end{minipage} \\
\midrule\noalign{}
\endhead
\bottomrule\noalign{}
\endlastfoot
\end{longtable}

\subsubsection{Description}\label{description-784}

The GCVD keyword defines the initial coal gas concentration versus depth
tables for each equilibration region for when the coal phase has been
activated in the run via the COAL keyword in the RUNSPEC section. The
keyword may be used in conjunction with the GASCONC keyword in the
SOLUTION section, to fully describe the initial state of the model. Note
both GASCONC and GCVD are optional as the simulator will calculate the
coal gas concentration based on the equilibrium concentration and the
block pressure.

This keyword is not supported by OPM Flow but would change the results
if supported so the simulation will be stopped.

\begin{longtable}[]{@{}
  >{\raggedright\arraybackslash}p{(\linewidth - 6\tabcolsep) * \real{0.2500}}
  >{\raggedright\arraybackslash}p{(\linewidth - 6\tabcolsep) * \real{0.2500}}
  >{\raggedright\arraybackslash}p{(\linewidth - 6\tabcolsep) * \real{0.2500}}
  >{\raggedright\arraybackslash}p{(\linewidth - 6\tabcolsep) * \real{0.2500}}@{}}
\toprule\noalign{}
\begin{minipage}[b]{\linewidth}\raggedright
No.
\end{minipage} & \begin{minipage}[b]{\linewidth}\raggedright
Name
\end{minipage} & \begin{minipage}[b]{\linewidth}\raggedright
Description
\end{minipage} & \begin{minipage}[b]{\linewidth}\raggedright
Default
\end{minipage} \\
\midrule\noalign{}
\endhead
\bottomrule\noalign{}
\endlastfoot
Field & Metric & Laboratory & \\
1 & DEPTH & A columnar vector of real monotonically increasing down the
column values that defines the depth values for the corresponding coal
gas concentration, GCVALS. & None \\
feet & m & cm & \\
2 & GCVALS & A columnar vector of real values that defines the coal gas
concentration values at the corresponding DEPTH. & None \\
Mscf/ft3 & sm3/m3 & scc/cc & \\
Notes: & & & \\
\end{longtable}

\emph{Table 10.17: GCVD Keyword Description}

See also the GASCONC and GASSATC keywords in the SOLUTION section.

\subsubsection{Example}\label{example-446}

Given NTEQUL equals three and NDRXVD is greater than or equal to two on
the EQLDIMS keyword in the RUNSPEC section, then the following example
defines the coal gas concentration versus depth functions.

\begin{verbatim}
--
--       DEPTH    GC
--                MSCF/FT
--       ------   --------
GCVD
          100.0   75.5000
         1000.0   75.5000                            / GC VS DEPTH EQUIL REGN 01
--       ------   --------
          100.0   65.5000
         1000.0   65.5000                            / GC VS DEPTH EQUIL REGN 02
--       ------   --------
          100.0   60.0000
         1000.0   60.0000                            / GC VS DEPTH EQUIL REGN 03
\end{verbatim}

Here three tables are entered with a constant coal gas concentration
versus depth relationship for each equilibration region.

\subsection{GETDATA -- Load and Assign Data Array from INIT or RESTART
File}\label{getdata-load-and-assign-data-array-from-init-or-restart-file-3}

The GETDATA keyword loads a data array from a previously generated INIT
or RESTART file and assigns the loaded array to either same array in the
run or another array name.

See GETDATA -- Load and Assign Data Array from INIT or RESTART Files in
the GRID section for a full description

\subsection{GETGLOB -- Activate Loading of Global Grid Restart Data
Option}\label{getglob-activate-loading-of-global-grid-restart-data-option}

\begin{longtable}[]{@{}
  >{\raggedright\arraybackslash}p{(\linewidth - 14\tabcolsep) * \real{0.1250}}
  >{\raggedright\arraybackslash}p{(\linewidth - 14\tabcolsep) * \real{0.1250}}
  >{\raggedright\arraybackslash}p{(\linewidth - 14\tabcolsep) * \real{0.1250}}
  >{\raggedright\arraybackslash}p{(\linewidth - 14\tabcolsep) * \real{0.1250}}
  >{\raggedright\arraybackslash}p{(\linewidth - 14\tabcolsep) * \real{0.1250}}
  >{\raggedright\arraybackslash}p{(\linewidth - 14\tabcolsep) * \real{0.1250}}
  >{\raggedright\arraybackslash}p{(\linewidth - 14\tabcolsep) * \real{0.1250}}
  >{\raggedright\arraybackslash}p{(\linewidth - 14\tabcolsep) * \real{0.1250}}@{}}
\toprule\noalign{}
\begin{minipage}[b]{\linewidth}\raggedright
RUNSPEC
\end{minipage} & \begin{minipage}[b]{\linewidth}\raggedright
GRID
\end{minipage} & \begin{minipage}[b]{\linewidth}\raggedright
EDIT
\end{minipage} & \begin{minipage}[b]{\linewidth}\raggedright
PROPS
\end{minipage} & \begin{minipage}[b]{\linewidth}\raggedright
REGIONS
\end{minipage} & \begin{minipage}[b]{\linewidth}\raggedright
SOLUTION
\end{minipage} & \begin{minipage}[b]{\linewidth}\raggedright
SUMMARY
\end{minipage} & \begin{minipage}[b]{\linewidth}\raggedright
SCHEDULE
\end{minipage} \\
\midrule\noalign{}
\endhead
\bottomrule\noalign{}
\endlastfoot
\end{longtable}

\subsubsection{Description}\label{description-785}

This keyword, GETGLOB, switches on the global grid read option for when
the run is restarting from a RESTART file. Only the global grid will be
loaded in the subsequent RESTART keyword and any Local Grid Refinements
(``LGR'') on the RESTART file will be ignored.

There is no data required for this keyword.

This keyword is not supported by OPM Flow but would change the results
if supported so the simulation will be stopped.

\subsubsection{Example}\label{example-447}

\begin{verbatim}
--
--       ACTIVATE LOADING OF GLOBAL GRID RESTART DATA OPTION
--
GETGLOB
\end{verbatim}

The above example switches on the option to only load the global grid
from the RESTART file.

\subsection{GI -- Define the Initial Equilibration Gi Values for All
Grid
Blocks}\label{gi-define-the-initial-equilibration-gi-values-for-all-grid-blocks}

\begin{longtable}[]{@{}
  >{\raggedright\arraybackslash}p{(\linewidth - 14\tabcolsep) * \real{0.1250}}
  >{\raggedright\arraybackslash}p{(\linewidth - 14\tabcolsep) * \real{0.1250}}
  >{\raggedright\arraybackslash}p{(\linewidth - 14\tabcolsep) * \real{0.1250}}
  >{\raggedright\arraybackslash}p{(\linewidth - 14\tabcolsep) * \real{0.1250}}
  >{\raggedright\arraybackslash}p{(\linewidth - 14\tabcolsep) * \real{0.1250}}
  >{\raggedright\arraybackslash}p{(\linewidth - 14\tabcolsep) * \real{0.1250}}
  >{\raggedright\arraybackslash}p{(\linewidth - 14\tabcolsep) * \real{0.1250}}
  >{\raggedright\arraybackslash}p{(\linewidth - 14\tabcolsep) * \real{0.1250}}@{}}
\toprule\noalign{}
\begin{minipage}[b]{\linewidth}\raggedright
RUNSPEC
\end{minipage} & \begin{minipage}[b]{\linewidth}\raggedright
GRID
\end{minipage} & \begin{minipage}[b]{\linewidth}\raggedright
EDIT
\end{minipage} & \begin{minipage}[b]{\linewidth}\raggedright
PROPS
\end{minipage} & \begin{minipage}[b]{\linewidth}\raggedright
REGIONS
\end{minipage} & \begin{minipage}[b]{\linewidth}\raggedright
SOLUTION
\end{minipage} & \begin{minipage}[b]{\linewidth}\raggedright
SUMMARY
\end{minipage} & \begin{minipage}[b]{\linewidth}\raggedright
SCHEDULE
\end{minipage} \\
\midrule\noalign{}
\endhead
\bottomrule\noalign{}
\endlastfoot
\end{longtable}

\subsubsection{Description}\label{description-786}

The GI keyword defines the initial equilibration GI values for all grid
cells in the model and should be used in conjunction with the other
enumeration equilibration keywords; PBUB, PDEW, PRESSURE, RS, RV, SOIL
and SWAT keywords etc., to fully describe the initial state of the
model. The keyword should only be used if the GI Pseudo Compositional
option has been activated in the model via the GIMODEL keyword in the
RUNSPEC section.

This is the non-standard method to initialize the model via enumeration
and is seldom employed in the industry. The standard methodology is for
OPM Flow to initialize a model using the parameters on the EQUIL keyword
combined with other keywords to fully describe the initial state of the
model. The keyword can be used with all grid types.

This keyword is not supported by OPM Flow but would change the results
if supported so the simulation will be stopped.

See also the GIALL keyword in the PROPS section that sets the GI values
as a function of pressure, as well as setting the corresponding RVGI,
RSGI, BGGI and BOGI values at the same time.

\subsection{HMAQUCT -- History Match Carter-Tracy Aquifer Gradient
Parameters}\label{hmaquct-history-match-carter-tracy-aquifer-gradient-parameters}

\begin{longtable}[]{@{}
  >{\raggedright\arraybackslash}p{(\linewidth - 14\tabcolsep) * \real{0.1250}}
  >{\raggedright\arraybackslash}p{(\linewidth - 14\tabcolsep) * \real{0.1250}}
  >{\raggedright\arraybackslash}p{(\linewidth - 14\tabcolsep) * \real{0.1250}}
  >{\raggedright\arraybackslash}p{(\linewidth - 14\tabcolsep) * \real{0.1250}}
  >{\raggedright\arraybackslash}p{(\linewidth - 14\tabcolsep) * \real{0.1250}}
  >{\raggedright\arraybackslash}p{(\linewidth - 14\tabcolsep) * \real{0.1250}}
  >{\raggedright\arraybackslash}p{(\linewidth - 14\tabcolsep) * \real{0.1250}}
  >{\raggedright\arraybackslash}p{(\linewidth - 14\tabcolsep) * \real{0.1250}}@{}}
\toprule\noalign{}
\begin{minipage}[b]{\linewidth}\raggedright
RUNSPEC
\end{minipage} & \begin{minipage}[b]{\linewidth}\raggedright
GRID
\end{minipage} & \begin{minipage}[b]{\linewidth}\raggedright
EDIT
\end{minipage} & \begin{minipage}[b]{\linewidth}\raggedright
PROPS
\end{minipage} & \begin{minipage}[b]{\linewidth}\raggedright
REGIONS
\end{minipage} & \begin{minipage}[b]{\linewidth}\raggedright
SOLUTION
\end{minipage} & \begin{minipage}[b]{\linewidth}\raggedright
SUMMARY
\end{minipage} & \begin{minipage}[b]{\linewidth}\raggedright
SCHEDULE
\end{minipage} \\
\midrule\noalign{}
\endhead
\bottomrule\noalign{}
\endlastfoot
\end{longtable}

\subsubsection{Description}\label{description-787}

The HMAQUCT keyword defines the history match analytical Carter-Tracy
aquifer gradient parameters for when the History Match Gradient option
has been activated by the HMDIMS keyword in the RUNSPEC section, and
analytical Carter-Tracy aquifers have been specified in the model via
the AQUCT and connected to the grid using the AQUANCON or AQANCONL
keywords. All keywords are in the SOLUTION section.

See also the HMDIMS keyword in the RUNSPEC section that specifies the
dimensions for the gradient option, including the maximum number of
aquifers that can be used with the History Match Gradient option.

This keyword is not supported by OPM Flow but would change the results
if supported so the simulation will be stopped.

\subsection{HMAQUFET -- History Match Fetkovich Aquifer Gradient
Parameters}\label{hmaqufet-history-match-fetkovich-aquifer-gradient-parameters}

\begin{longtable}[]{@{}
  >{\raggedright\arraybackslash}p{(\linewidth - 14\tabcolsep) * \real{0.1250}}
  >{\raggedright\arraybackslash}p{(\linewidth - 14\tabcolsep) * \real{0.1250}}
  >{\raggedright\arraybackslash}p{(\linewidth - 14\tabcolsep) * \real{0.1250}}
  >{\raggedright\arraybackslash}p{(\linewidth - 14\tabcolsep) * \real{0.1250}}
  >{\raggedright\arraybackslash}p{(\linewidth - 14\tabcolsep) * \real{0.1250}}
  >{\raggedright\arraybackslash}p{(\linewidth - 14\tabcolsep) * \real{0.1250}}
  >{\raggedright\arraybackslash}p{(\linewidth - 14\tabcolsep) * \real{0.1250}}
  >{\raggedright\arraybackslash}p{(\linewidth - 14\tabcolsep) * \real{0.1250}}@{}}
\toprule\noalign{}
\begin{minipage}[b]{\linewidth}\raggedright
RUNSPEC
\end{minipage} & \begin{minipage}[b]{\linewidth}\raggedright
GRID
\end{minipage} & \begin{minipage}[b]{\linewidth}\raggedright
EDIT
\end{minipage} & \begin{minipage}[b]{\linewidth}\raggedright
PROPS
\end{minipage} & \begin{minipage}[b]{\linewidth}\raggedright
REGIONS
\end{minipage} & \begin{minipage}[b]{\linewidth}\raggedright
SOLUTION
\end{minipage} & \begin{minipage}[b]{\linewidth}\raggedright
SUMMARY
\end{minipage} & \begin{minipage}[b]{\linewidth}\raggedright
SCHEDULE
\end{minipage} \\
\midrule\noalign{}
\endhead
\bottomrule\noalign{}
\endlastfoot
\end{longtable}

\subsubsection{Description}\label{description-788}

The HMAQUFET keyword defines the history match analytical Fetkovich
aquifer gradient parameters for when the History Match Gradient option
has been activated by the HMDIMS keyword in the RUNSPEC section, and
analytical Fetkovich aquifers have been specified in the model via the
AQUFET and/or the AQUFETP keywords and connected to the grid using
AQUANCON or AQANCONL keywords. All keywords are in the SOLUTION section.

See also the HMDIMS keyword in the RUNSPEC section that specifies the
dimensions for the gradient option, including the maximum number of
aquifers that can be used with the History Match Gradient option.

This keyword is not supported by OPM Flow but would change the results
if supported so the simulation will be stopped.

\subsection{HMMLCTAQ -- History Match Carter-Tracy Aquifer Gradient
Multipliers}\label{hmmlctaq-history-match-carter-tracy-aquifer-gradient-multipliers}

\begin{longtable}[]{@{}
  >{\raggedright\arraybackslash}p{(\linewidth - 14\tabcolsep) * \real{0.1250}}
  >{\raggedright\arraybackslash}p{(\linewidth - 14\tabcolsep) * \real{0.1250}}
  >{\raggedright\arraybackslash}p{(\linewidth - 14\tabcolsep) * \real{0.1250}}
  >{\raggedright\arraybackslash}p{(\linewidth - 14\tabcolsep) * \real{0.1250}}
  >{\raggedright\arraybackslash}p{(\linewidth - 14\tabcolsep) * \real{0.1250}}
  >{\raggedright\arraybackslash}p{(\linewidth - 14\tabcolsep) * \real{0.1250}}
  >{\raggedright\arraybackslash}p{(\linewidth - 14\tabcolsep) * \real{0.1250}}
  >{\raggedright\arraybackslash}p{(\linewidth - 14\tabcolsep) * \real{0.1250}}@{}}
\toprule\noalign{}
\begin{minipage}[b]{\linewidth}\raggedright
RUNSPEC
\end{minipage} & \begin{minipage}[b]{\linewidth}\raggedright
GRID
\end{minipage} & \begin{minipage}[b]{\linewidth}\raggedright
EDIT
\end{minipage} & \begin{minipage}[b]{\linewidth}\raggedright
PROPS
\end{minipage} & \begin{minipage}[b]{\linewidth}\raggedright
REGIONS
\end{minipage} & \begin{minipage}[b]{\linewidth}\raggedright
SOLUTION
\end{minipage} & \begin{minipage}[b]{\linewidth}\raggedright
SUMMARY
\end{minipage} & \begin{minipage}[b]{\linewidth}\raggedright
SCHEDULE
\end{minipage} \\
\midrule\noalign{}
\endhead
\bottomrule\noalign{}
\endlastfoot
\end{longtable}

\subsubsection{Description}\label{description-789}

The HMMLCTAQ keyword defines the history match analytical Carter-Tracy
aquifer gradient multipliers for when the History Match Gradient option
has been activated by the HMDIMS keyword in the RUNSPEC section, and
analytical Carter-Tracy aquifers have been specified in the model via
the AQUCT and connected to the grid using the AQUANCON or AQANCONL
keywords. All keywords are in the SOLUTION section.

Multipliers can be declared for the Carter-Tracy aquifer permeability,
aquifer angle of influence and the aquifer depth.

This keyword is not supported by OPM Flow but would change the results
if supported so the simulation will be stopped.

\subsection{HMMLFTAQ -- History Match Fetkovich Aquifer Gradient
Multipliers}\label{hmmlftaq-history-match-fetkovich-aquifer-gradient-multipliers}

\begin{longtable}[]{@{}
  >{\raggedright\arraybackslash}p{(\linewidth - 14\tabcolsep) * \real{0.1250}}
  >{\raggedright\arraybackslash}p{(\linewidth - 14\tabcolsep) * \real{0.1250}}
  >{\raggedright\arraybackslash}p{(\linewidth - 14\tabcolsep) * \real{0.1250}}
  >{\raggedright\arraybackslash}p{(\linewidth - 14\tabcolsep) * \real{0.1250}}
  >{\raggedright\arraybackslash}p{(\linewidth - 14\tabcolsep) * \real{0.1250}}
  >{\raggedright\arraybackslash}p{(\linewidth - 14\tabcolsep) * \real{0.1250}}
  >{\raggedright\arraybackslash}p{(\linewidth - 14\tabcolsep) * \real{0.1250}}
  >{\raggedright\arraybackslash}p{(\linewidth - 14\tabcolsep) * \real{0.1250}}@{}}
\toprule\noalign{}
\begin{minipage}[b]{\linewidth}\raggedright
RUNSPEC
\end{minipage} & \begin{minipage}[b]{\linewidth}\raggedright
GRID
\end{minipage} & \begin{minipage}[b]{\linewidth}\raggedright
EDIT
\end{minipage} & \begin{minipage}[b]{\linewidth}\raggedright
PROPS
\end{minipage} & \begin{minipage}[b]{\linewidth}\raggedright
REGIONS
\end{minipage} & \begin{minipage}[b]{\linewidth}\raggedright
SOLUTION
\end{minipage} & \begin{minipage}[b]{\linewidth}\raggedright
SUMMARY
\end{minipage} & \begin{minipage}[b]{\linewidth}\raggedright
SCHEDULE
\end{minipage} \\
\midrule\noalign{}
\endhead
\bottomrule\noalign{}
\endlastfoot
\end{longtable}

\subsubsection{Description}\label{description-790}

The HMAQUFET keyword defines the history match analytical Fetkovich
aquifer gradient multipliers for when the History Match Gradient option
has been activated by the HMDIMS keyword in the RUNSPEC section, and
analytical Fetkovich aquifers have been specified in the model via the
AQUFET and/or the AQUFETP keywords and connected to the grid using
AQUANCON or AQANCONL keywords. All keywords are in the SOLUTION section.

Multipliers can be declared for the Fetkovich aquifer water volume,
aquifer permeability, and the aquifer depth.

This keyword is not supported by OPM Flow but would change the results
if supported so the simulation will be stopped.

\subsection{HMMLTWCN -- History Match Well Connection and Skin
Multipliers}\label{hmmltwcn-history-match-well-connection-and-skin-multipliers}

\begin{longtable}[]{@{}
  >{\raggedright\arraybackslash}p{(\linewidth - 14\tabcolsep) * \real{0.1250}}
  >{\raggedright\arraybackslash}p{(\linewidth - 14\tabcolsep) * \real{0.1250}}
  >{\raggedright\arraybackslash}p{(\linewidth - 14\tabcolsep) * \real{0.1250}}
  >{\raggedright\arraybackslash}p{(\linewidth - 14\tabcolsep) * \real{0.1250}}
  >{\raggedright\arraybackslash}p{(\linewidth - 14\tabcolsep) * \real{0.1250}}
  >{\raggedright\arraybackslash}p{(\linewidth - 14\tabcolsep) * \real{0.1250}}
  >{\raggedright\arraybackslash}p{(\linewidth - 14\tabcolsep) * \real{0.1250}}
  >{\raggedright\arraybackslash}p{(\linewidth - 14\tabcolsep) * \real{0.1250}}@{}}
\toprule\noalign{}
\begin{minipage}[b]{\linewidth}\raggedright
RUNSPEC
\end{minipage} & \begin{minipage}[b]{\linewidth}\raggedright
GRID
\end{minipage} & \begin{minipage}[b]{\linewidth}\raggedright
EDIT
\end{minipage} & \begin{minipage}[b]{\linewidth}\raggedright
PROPS
\end{minipage} & \begin{minipage}[b]{\linewidth}\raggedright
REGIONS
\end{minipage} & \begin{minipage}[b]{\linewidth}\raggedright
SOLUTION
\end{minipage} & \begin{minipage}[b]{\linewidth}\raggedright
SUMMARY
\end{minipage} & \begin{minipage}[b]{\linewidth}\raggedright
SCHEDULE
\end{minipage} \\
\midrule\noalign{}
\endhead
\bottomrule\noalign{}
\endlastfoot
\end{longtable}

\subsubsection{Description}\label{description-791}

This keyword, HMMLTWCN, defines the history match gradient multipliers
for well connection factors and connection skins, for when the History
Match Gradient option has been activated by the HMDIMS keyword in the
RUNSPEC section. Wells must be specified using the WELPSECS keyword in
the SCHEDULE section and their connections defined by the COMPDAT and/or
COMPDATL keywords, also in the SCHEDULE section

This keyword is not supported by OPM Flow but would change the results
if supported so the simulation will be stopped.

\subsection{HMWELCON -- History Match Well Connection and Skin
Parameters}\label{hmwelcon-history-match-well-connection-and-skin-parameters}

\begin{longtable}[]{@{}
  >{\raggedright\arraybackslash}p{(\linewidth - 14\tabcolsep) * \real{0.1250}}
  >{\raggedright\arraybackslash}p{(\linewidth - 14\tabcolsep) * \real{0.1250}}
  >{\raggedright\arraybackslash}p{(\linewidth - 14\tabcolsep) * \real{0.1250}}
  >{\raggedright\arraybackslash}p{(\linewidth - 14\tabcolsep) * \real{0.1250}}
  >{\raggedright\arraybackslash}p{(\linewidth - 14\tabcolsep) * \real{0.1250}}
  >{\raggedright\arraybackslash}p{(\linewidth - 14\tabcolsep) * \real{0.1250}}
  >{\raggedright\arraybackslash}p{(\linewidth - 14\tabcolsep) * \real{0.1250}}
  >{\raggedright\arraybackslash}p{(\linewidth - 14\tabcolsep) * \real{0.1250}}@{}}
\toprule\noalign{}
\begin{minipage}[b]{\linewidth}\raggedright
RUNSPEC
\end{minipage} & \begin{minipage}[b]{\linewidth}\raggedright
GRID
\end{minipage} & \begin{minipage}[b]{\linewidth}\raggedright
EDIT
\end{minipage} & \begin{minipage}[b]{\linewidth}\raggedright
PROPS
\end{minipage} & \begin{minipage}[b]{\linewidth}\raggedright
REGIONS
\end{minipage} & \begin{minipage}[b]{\linewidth}\raggedright
SOLUTION
\end{minipage} & \begin{minipage}[b]{\linewidth}\raggedright
SUMMARY
\end{minipage} & \begin{minipage}[b]{\linewidth}\raggedright
SCHEDULE
\end{minipage} \\
\midrule\noalign{}
\endhead
\bottomrule\noalign{}
\endlastfoot
\end{longtable}

\subsubsection{Description}\label{description-792}

This keyword, HMWELCON, defines the history match gradient parameters
for well connection factors and connection skins, for when the History
Match Gradient option has been activated by the HMDIMS keyword in the
RUNSPEC section. Wells must be specified using the WELPSECS keyword in
the SCHEDULE section and their connections defined by the COMPDAT and/or
COMPDATL keywords, also in the SCHEDULE section

This keyword is not supported by OPM Flow but would change the results
if supported so the simulation will be stopped.

\subsection{IMPORT -- Import Grid File Data at the Current
Position}\label{import-import-grid-file-data-at-the-current-position-4}

The IMPORT keyword informs the simulator to continue reading input data
from the specified IMPORT file. When the end of the IMPORT file is
reached, input data is read from the next keyword in the current file.
Normally IMPORT files are generated by grid pre-processing software and
the keyword allows for both formatted and unformatted (binary) files to
be loaded.

See IMPORT -- Import Grid File Data at the Current Position in the GRID
section for a full description.

\subsection{INCLUDE -- Load Another Data File at the Current
Position}\label{include-load-another-data-file-at-the-current-position-6}

The INCLUDE keyword informs OPM Flow to continue reading input data from
the specified INCLUDE file. When the end of the INCLUDE file is reached,
or the ENDINC keyword is encountered in the included file, input data is
read from the next keyword in the current file.

See INCLUDE -- Load Another Data File at the Current Position in the
GLOBAL section for a full description.

\subsection{MESSAGE -- Output User
Message}\label{message-output-user-message-6}

The MESSAGE keyword outputs a user message to the terminal, as well as
to the print (\emph{.PRT) and debug (}.DBG) files. Note this is
different to the MESSAGES keyword, that defines OPM Flows message print
limits and stop limits generated by the simulator.

See MESSAGE -- Output User Message in the GLOBAL section for a full
description.

\subsection{MESSAGES -- Define Message Print Limits and Stop
Limits}\label{messages-define-message-print-limits-and-stop-limits-6}

The MESSAGES keyword defines the print and stops levels for various
messages. The ``print limits'' set the maximum number of messages that
will be printed, after which no more messages will be printed and the
``stop limits'' terminate the run when these limits are exceeded. There
are six levels of message that increase in severity from informative all
the way to programming errors, as outlined in Table 4.5.

See MESSAGES -- Define Message Print Limits and Stop Limits in the
GLOBAL section for a full description.

\subsection{MULTIPLY -- Multiply a Specified Array by a
Constant}\label{multiply-multiply-a-specified-array-by-a-constant-4}

The MULTIPLY keyword multiplies a specified array or part of an array by
a constant. The constant can be real or integer depending on the array
type; however, the arrays that can be operated on is dependent on which
section the keyword is being used.

See MULTIPLY -- Multiply a Specified Array by a Constant in the GRID
section for a full description.

\subsection{MULTIREG -- Multiply an Array by a Constant based on a
Region
Number}\label{multireg-multiply-an-array-by-a-constant-based-on-a-region-number-4}

The MULTIREG keyword multiplies an array or part of an array by a
constant for cells with a specific region number. The region number
array can be FLUXNUM, MULTNUM or OPERNUM and these arrays must be
defined and be available before the MULTIREG keyword is read by the
simulator. The constant can be real or integer depending on the property
array type; however, the property arrays that can be operated on is
dependent on which section the MULTIREG keyword is being used.

See MULTIREG -- Multiply an Array by a Constant based on a Region Number
in the GRID section for a full description.

\subsection{NOECHO -- Deactivate Echoing of User Input Files to the
Print
File}\label{noecho-deactivate-echoing-of-user-input-files-to-the-print-file-6}

Turns off echoing of all the input files to the print file. Note by
default echoing of the inputs files is active. but can subsequently be
switched off by the NOECHO activation keyword.

See NOECHO -- Deactivate Echoing of User Input Files to the Print File
in the GLOBAL section for a full description.

\subsection{NOHMD -- Deactivate History Match Gradient Derivative
Calculations}\label{nohmd-deactivate-history-match-gradient-derivative-calculations}

\begin{longtable}[]{@{}
  >{\raggedright\arraybackslash}p{(\linewidth - 14\tabcolsep) * \real{0.1250}}
  >{\raggedright\arraybackslash}p{(\linewidth - 14\tabcolsep) * \real{0.1250}}
  >{\raggedright\arraybackslash}p{(\linewidth - 14\tabcolsep) * \real{0.1250}}
  >{\raggedright\arraybackslash}p{(\linewidth - 14\tabcolsep) * \real{0.1250}}
  >{\raggedright\arraybackslash}p{(\linewidth - 14\tabcolsep) * \real{0.1250}}
  >{\raggedright\arraybackslash}p{(\linewidth - 14\tabcolsep) * \real{0.1250}}
  >{\raggedright\arraybackslash}p{(\linewidth - 14\tabcolsep) * \real{0.1250}}
  >{\raggedright\arraybackslash}p{(\linewidth - 14\tabcolsep) * \real{0.1250}}@{}}
\toprule\noalign{}
\begin{minipage}[b]{\linewidth}\raggedright
RUNSPEC
\end{minipage} & \begin{minipage}[b]{\linewidth}\raggedright
GRID
\end{minipage} & \begin{minipage}[b]{\linewidth}\raggedright
EDIT
\end{minipage} & \begin{minipage}[b]{\linewidth}\raggedright
PROPS
\end{minipage} & \begin{minipage}[b]{\linewidth}\raggedright
REGIONS
\end{minipage} & \begin{minipage}[b]{\linewidth}\raggedright
SOLUTION
\end{minipage} & \begin{minipage}[b]{\linewidth}\raggedright
SUMMARY
\end{minipage} & \begin{minipage}[b]{\linewidth}\raggedright
SCHEDULE
\end{minipage} \\
\midrule\noalign{}
\endhead
\bottomrule\noalign{}
\endlastfoot
\end{longtable}

\subsubsection{Description}\label{description-793}

The NOHMD deactivates various history match gradient derivative
calculations for when the History Match Gradient option has been
activated by the HMDIMS keyword in the RUNSPEC section. The keyword
consists of a series of character strings that define which derivative
should be switch off based on the keyword that requested the derivatives
to be calculated, for example HMFAULTS keyword in the GRID section. If
an empty list is entered then all the gradient derivative calculations
previously requested are switch off. The keyword is useful for changing
from history matching runs to predication cases, as the prediction cases
will be more computationally efficient without the burden of the
gradient derivative calculations.

This keyword is not supported by OPM Flow but would change the results
if supported so the simulation will be stopped.

\subsection{NOHMO -- Deactivate History Match Gradient Derivative
Calculations
(Alias)}\label{nohmo-deactivate-history-match-gradient-derivative-calculations-alias}

\begin{longtable}[]{@{}
  >{\raggedright\arraybackslash}p{(\linewidth - 14\tabcolsep) * \real{0.1250}}
  >{\raggedright\arraybackslash}p{(\linewidth - 14\tabcolsep) * \real{0.1250}}
  >{\raggedright\arraybackslash}p{(\linewidth - 14\tabcolsep) * \real{0.1250}}
  >{\raggedright\arraybackslash}p{(\linewidth - 14\tabcolsep) * \real{0.1250}}
  >{\raggedright\arraybackslash}p{(\linewidth - 14\tabcolsep) * \real{0.1250}}
  >{\raggedright\arraybackslash}p{(\linewidth - 14\tabcolsep) * \real{0.1250}}
  >{\raggedright\arraybackslash}p{(\linewidth - 14\tabcolsep) * \real{0.1250}}
  >{\raggedright\arraybackslash}p{(\linewidth - 14\tabcolsep) * \real{0.1250}}@{}}
\toprule\noalign{}
\begin{minipage}[b]{\linewidth}\raggedright
RUNSPEC
\end{minipage} & \begin{minipage}[b]{\linewidth}\raggedright
GRID
\end{minipage} & \begin{minipage}[b]{\linewidth}\raggedright
EDIT
\end{minipage} & \begin{minipage}[b]{\linewidth}\raggedright
PROPS
\end{minipage} & \begin{minipage}[b]{\linewidth}\raggedright
REGIONS
\end{minipage} & \begin{minipage}[b]{\linewidth}\raggedright
SOLUTION
\end{minipage} & \begin{minipage}[b]{\linewidth}\raggedright
SUMMARY
\end{minipage} & \begin{minipage}[b]{\linewidth}\raggedright
SCHEDULE
\end{minipage} \\
\midrule\noalign{}
\endhead
\bottomrule\noalign{}
\endlastfoot
\end{longtable}

\subsubsection{Description}\label{description-794}

The NOHMO deactivates various history match gradient derivative
calculations for when the History Match Gradient option has been
activated by the HMDIMS keyword in the RUNSPEC section. The keyword
consists of a series of character strings that define which derivative
should be switch off based on the keyword that requested the derivatives
to be calculated, for example HMFAULTS keyword in the GRID section. If
an empty list is entered then all the gradient derivative calculations
previously requested are switch off. The keyword is useful for changing
from history matching runs to predication cases, as the prediction cases
will be more computationally efficient without the burden of the
gradient derivative calculations.

The keyword is an alias for the NOHMD keyword in the SOLUTION section.

This keyword is not supported by OPM Flow but would change the results
if supported so the simulation will be stopped.

\subsection{NOWARN -- Deactivate Warning
Messages}\label{nowarn-deactivate-warning-messages-6}

Turns off warning messages to be printed to the print file; note that
this keyword is deactivated by default and can subsequently be switched
off by the WARN activation keyword. The warning messages may be turned
on and off using keywords WARN and NOWARN.

See NOWARN -- Deactivate Warning Messages in the GLOBAL section for a
full description.

\subsection{OILAPI -- Define the Initial Equilibration Oil API for All
Grid
Blocks}\label{oilapi-define-the-initial-equilibration-oil-api-for-all-grid-blocks}

\begin{longtable}[]{@{}
  >{\raggedright\arraybackslash}p{(\linewidth - 14\tabcolsep) * \real{0.1250}}
  >{\raggedright\arraybackslash}p{(\linewidth - 14\tabcolsep) * \real{0.1250}}
  >{\raggedright\arraybackslash}p{(\linewidth - 14\tabcolsep) * \real{0.1250}}
  >{\raggedright\arraybackslash}p{(\linewidth - 14\tabcolsep) * \real{0.1250}}
  >{\raggedright\arraybackslash}p{(\linewidth - 14\tabcolsep) * \real{0.1250}}
  >{\raggedright\arraybackslash}p{(\linewidth - 14\tabcolsep) * \real{0.1250}}
  >{\raggedright\arraybackslash}p{(\linewidth - 14\tabcolsep) * \real{0.1250}}
  >{\raggedright\arraybackslash}p{(\linewidth - 14\tabcolsep) * \real{0.1250}}@{}}
\toprule\noalign{}
\begin{minipage}[b]{\linewidth}\raggedright
RUNSPEC
\end{minipage} & \begin{minipage}[b]{\linewidth}\raggedright
GRID
\end{minipage} & \begin{minipage}[b]{\linewidth}\raggedright
EDIT
\end{minipage} & \begin{minipage}[b]{\linewidth}\raggedright
PROPS
\end{minipage} & \begin{minipage}[b]{\linewidth}\raggedright
REGIONS
\end{minipage} & \begin{minipage}[b]{\linewidth}\raggedright
SOLUTION
\end{minipage} & \begin{minipage}[b]{\linewidth}\raggedright
SUMMARY
\end{minipage} & \begin{minipage}[b]{\linewidth}\raggedright
SCHEDULE
\end{minipage} \\
\midrule\noalign{}
\endhead
\bottomrule\noalign{}
\endlastfoot
\end{longtable}

\subsubsection{Description}\label{description-795}

The OILAPI keyword defines the initial equilibration oil API gravity
pressures for all grid cells in the model, for when the Oil API Tracking
option as been invoked by the API keyword in the RUNSPEC section. The
keyword should be used in conjunction with the PBUB, PDEW, PRESSURE, RS,
RV, SGAS, SOIL and SWAT keywords etc., to fully describe the initial
state of the model.

The keyword is used by the Enumeration Initialization method to
initialize the model, as opposed to the Equilibration Initialization
method that utilizes the EQUIL keyword in the SOLUTION section. This is
the non-standard formulation to initialize the model and is seldom
employed in the industry. The standard methodology is for OPM Flow to
initialize a model using the parameters on the EQUIL keyword combined
with other keywords to fully describe the initial state of the model.
The keyword can be used with all grid types.

\begin{longtable}[]{@{}
  >{\raggedright\arraybackslash}p{(\linewidth - 6\tabcolsep) * \real{0.2500}}
  >{\raggedright\arraybackslash}p{(\linewidth - 6\tabcolsep) * \real{0.2500}}
  >{\raggedright\arraybackslash}p{(\linewidth - 6\tabcolsep) * \real{0.2500}}
  >{\raggedright\arraybackslash}p{(\linewidth - 6\tabcolsep) * \real{0.2500}}@{}}
\toprule\noalign{}
\begin{minipage}[b]{\linewidth}\raggedright
No.
\end{minipage} & \begin{minipage}[b]{\linewidth}\raggedright
Name
\end{minipage} & \begin{minipage}[b]{\linewidth}\raggedright
Description
\end{minipage} & \begin{minipage}[b]{\linewidth}\raggedright
Default
\end{minipage} \\
\midrule\noalign{}
\endhead
\bottomrule\noalign{}
\endlastfoot
Field & Metric & Laboratory & \\
1 & OILAPI & OILAPI is an array of real positive numbers assigning the
initial equilibration oil API gravity to each cell in the model. The
American Petroleum Institute (``API'') classifies oils based on an API
gravity (γAPI), or degrees API (oAPI), the relationship between relative
density (γo) of oil and API gravity (γAPI) is given by:
\({\mathrm{γ}}_{\mathit{API}} = \frac{141.5}{{\mathrm{γ}}_{o}} - 131.5\)
Repeat counts may be used, for example 20*38.5 & None \\
oAPI & oAPI & oAPI & \\
Notes: & & & \\
\end{longtable}

\emph{Table 10.18: OILAPI Keyword Description}

See also the PBUB, PDEW, PRESSURE, RS, RV, SGAS, SOIL and SWAT keywords
to fully define the initial state of the model.

\subsubsection{Example}\label{example-448}

\begin{verbatim}
--
--       DEFINE INITIAL EQUILIBRATION OIL API FOR ALL CELLS IN THE MODEL
--       BASED ON NX = 100, NY = 100 AND NZ = 3
--
OILAPI
         1000*40.2     1000*39.5     1000*38.2                                 /
\end{verbatim}

The above example defines the initial equilibration oil API gravity to
be 40.2 for all the cells in the first layer, 39.5 for all the cells in
the second layer, and finally 38.2 for all the cells in the third layer.

\subsection{OPERATE -- Define Mathematical Operations on
Arrays}\label{operate-define-mathematical-operations-on-arrays-4}

This keyword, OPERATE, defines mathematical operations on property
arrays (NTG, PORO etc.) and optionally using another property array as
input to the function. The keyword allows for various mathematical
functions and their associated variables to be defined and applied to
the selected array data. Input constants can be real or integer
depending on the array type; however, the arrays that can be operated on
is dependent on which section the keyword is being used.

See OPERATE -- Define Mathematical Operations on Arrays in the GRID
section for a full description.

\subsection{OPERATER -- Define Mathematical Operations on Arrays by
Region}\label{operater-define-mathematical-operations-on-arrays-by-region-4}

This keyword defines the mathematical operations on arrays for specific
regions in the commercial simulator and is currently not supported by
OPM Flow. However, similar functionality is provided by the ADD and
MULTIPLY keywords.

See OPERATER -- Define Mathematical Operations on Arrays by Region in
the GRID section for a full description.

\subsection{OUTSOL -- Define Data to be Written to the RESTART File
(Retired)}\label{outsol-define-data-to-be-written-to-the-restart-file-retired}

\begin{longtable}[]{@{}
  >{\raggedright\arraybackslash}p{(\linewidth - 14\tabcolsep) * \real{0.1250}}
  >{\raggedright\arraybackslash}p{(\linewidth - 14\tabcolsep) * \real{0.1250}}
  >{\raggedright\arraybackslash}p{(\linewidth - 14\tabcolsep) * \real{0.1250}}
  >{\raggedright\arraybackslash}p{(\linewidth - 14\tabcolsep) * \real{0.1250}}
  >{\raggedright\arraybackslash}p{(\linewidth - 14\tabcolsep) * \real{0.1250}}
  >{\raggedright\arraybackslash}p{(\linewidth - 14\tabcolsep) * \real{0.1250}}
  >{\raggedright\arraybackslash}p{(\linewidth - 14\tabcolsep) * \real{0.1250}}
  >{\raggedright\arraybackslash}p{(\linewidth - 14\tabcolsep) * \real{0.1250}}@{}}
\toprule\noalign{}
\begin{minipage}[b]{\linewidth}\raggedright
RUNSPEC
\end{minipage} & \begin{minipage}[b]{\linewidth}\raggedright
GRID
\end{minipage} & \begin{minipage}[b]{\linewidth}\raggedright
EDIT
\end{minipage} & \begin{minipage}[b]{\linewidth}\raggedright
PROPS
\end{minipage} & \begin{minipage}[b]{\linewidth}\raggedright
REGIONS
\end{minipage} & \begin{minipage}[b]{\linewidth}\raggedright
SOLUTION
\end{minipage} & \begin{minipage}[b]{\linewidth}\raggedright
SUMMARY
\end{minipage} & \begin{minipage}[b]{\linewidth}\raggedright
SCHEDULE
\end{minipage} \\
\midrule\noalign{}
\endhead
\bottomrule\noalign{}
\endlastfoot
\end{longtable}

\subsubsection{Description}\label{description-796}

This keyword defines the data and frequency of the data to be written to
the RESTART file at each requested restart point. The keyword has been
replaced by the RPTRST keyword in the SOLUTION and SCHEDULE sections and
is therefore considered retired.

This keyword is not supported by OPM Flow but has no effect on the
results so it will be ignored.

\subsection{PBUB -- Define the Initial Equilibration Bubble-Point
Pressure for All Grid
Blocks}\label{pbub-define-the-initial-equilibration-bubble-point-pressure-for-all-grid-blocks}

\begin{longtable}[]{@{}
  >{\raggedright\arraybackslash}p{(\linewidth - 14\tabcolsep) * \real{0.1250}}
  >{\raggedright\arraybackslash}p{(\linewidth - 14\tabcolsep) * \real{0.1250}}
  >{\raggedright\arraybackslash}p{(\linewidth - 14\tabcolsep) * \real{0.1250}}
  >{\raggedright\arraybackslash}p{(\linewidth - 14\tabcolsep) * \real{0.1250}}
  >{\raggedright\arraybackslash}p{(\linewidth - 14\tabcolsep) * \real{0.1250}}
  >{\raggedright\arraybackslash}p{(\linewidth - 14\tabcolsep) * \real{0.1250}}
  >{\raggedright\arraybackslash}p{(\linewidth - 14\tabcolsep) * \real{0.1250}}
  >{\raggedright\arraybackslash}p{(\linewidth - 14\tabcolsep) * \real{0.1250}}@{}}
\toprule\noalign{}
\begin{minipage}[b]{\linewidth}\raggedright
RUNSPEC
\end{minipage} & \begin{minipage}[b]{\linewidth}\raggedright
GRID
\end{minipage} & \begin{minipage}[b]{\linewidth}\raggedright
EDIT
\end{minipage} & \begin{minipage}[b]{\linewidth}\raggedright
PROPS
\end{minipage} & \begin{minipage}[b]{\linewidth}\raggedright
REGIONS
\end{minipage} & \begin{minipage}[b]{\linewidth}\raggedright
SOLUTION
\end{minipage} & \begin{minipage}[b]{\linewidth}\raggedright
SUMMARY
\end{minipage} & \begin{minipage}[b]{\linewidth}\raggedright
SCHEDULE
\end{minipage} \\
\midrule\noalign{}
\endhead
\bottomrule\noalign{}
\endlastfoot
\end{longtable}

\subsubsection{Description}\label{description-797}

The PBUB keyword defines the initial equilibration buble-point
saturation pressures values for all grid cells in the model and should
be used in conjunction with the PDEW, PRESSURE, RS, RV, SGAS, SOIL and
SWAT keywords etc., to fully describe the initial state of the model.
The keyword should only be used if dissolved gas has been activated in
the model via the DISGAS keyword in the RUNSPEC section.

This is the non-standard method to initialize the model via enumeration
and is seldom employed in the industry. The standard methodology is for
OPM Flow to initialize a model using the parameters on the EQUIL keyword
combined with other keywords to fully describe the initial state of the
model. The keyword can be used with all grid types.

This keyword is not supported by OPM Flow but would change the results
if supported so the simulation will be stopped.

\begin{longtable}[]{@{}
  >{\raggedright\arraybackslash}p{(\linewidth - 6\tabcolsep) * \real{0.2500}}
  >{\raggedright\arraybackslash}p{(\linewidth - 6\tabcolsep) * \real{0.2500}}
  >{\raggedright\arraybackslash}p{(\linewidth - 6\tabcolsep) * \real{0.2500}}
  >{\raggedright\arraybackslash}p{(\linewidth - 6\tabcolsep) * \real{0.2500}}@{}}
\toprule\noalign{}
\begin{minipage}[b]{\linewidth}\raggedright
No.
\end{minipage} & \begin{minipage}[b]{\linewidth}\raggedright
Name
\end{minipage} & \begin{minipage}[b]{\linewidth}\raggedright
Description
\end{minipage} & \begin{minipage}[b]{\linewidth}\raggedright
Default
\end{minipage} \\
\midrule\noalign{}
\endhead
\bottomrule\noalign{}
\endlastfoot
Field & Metric & Laboratory & \\
1 & PBUB & PBUB is an array of real positive numbers assigning the
initial equilibration bubble-point saturation pressure values to each
cell in the model. Repeat counts may be used, for example 20*3500.0 &
None \\
psia & barsa & atma & \\
Notes: & & & \\
\end{longtable}

\emph{Table 10.19: PBUB Keyword Description}

See also the PBVD, PDEW, PRESSURE, RV, SGAS, SOIL and SWAT keywords to
fully define the initial state of the model.

\subsubsection{Example}\label{example-449}

\begin{verbatim}
--
--       DEFINE INITIAL EQUILIBRATION PSAT VALUES FOR ALL CELLS IN THE MODEL
--       BASED ON NX = 100, NY = 100 AND NZ = 3
--
PBUB
         1000*3500.0    1000*3525.0    1000*0.3535.0                           /
\end{verbatim}

The above example defines the initial equilibration bubble-point
saturation pressure values to be 3500.0 for all the cells in the first
layer, 3525.0 for all the cells in the second layer, and finally 3535.0
for all the cells in the third layer.

\subsection{PBVD -- Equilibration Bubble-Point versus Depth
Tables}\label{pbvd-equilibration-bubble-point-versus-depth-tables}

\begin{longtable}[]{@{}
  >{\raggedright\arraybackslash}p{(\linewidth - 14\tabcolsep) * \real{0.1250}}
  >{\raggedright\arraybackslash}p{(\linewidth - 14\tabcolsep) * \real{0.1250}}
  >{\raggedright\arraybackslash}p{(\linewidth - 14\tabcolsep) * \real{0.1250}}
  >{\raggedright\arraybackslash}p{(\linewidth - 14\tabcolsep) * \real{0.1250}}
  >{\raggedright\arraybackslash}p{(\linewidth - 14\tabcolsep) * \real{0.1250}}
  >{\raggedright\arraybackslash}p{(\linewidth - 14\tabcolsep) * \real{0.1250}}
  >{\raggedright\arraybackslash}p{(\linewidth - 14\tabcolsep) * \real{0.1250}}
  >{\raggedright\arraybackslash}p{(\linewidth - 14\tabcolsep) * \real{0.1250}}@{}}
\toprule\noalign{}
\begin{minipage}[b]{\linewidth}\raggedright
RUNSPEC
\end{minipage} & \begin{minipage}[b]{\linewidth}\raggedright
GRID
\end{minipage} & \begin{minipage}[b]{\linewidth}\raggedright
EDIT
\end{minipage} & \begin{minipage}[b]{\linewidth}\raggedright
PROPS
\end{minipage} & \begin{minipage}[b]{\linewidth}\raggedright
REGIONS
\end{minipage} & \begin{minipage}[b]{\linewidth}\raggedright
SOLUTION
\end{minipage} & \begin{minipage}[b]{\linewidth}\raggedright
SUMMARY
\end{minipage} & \begin{minipage}[b]{\linewidth}\raggedright
SCHEDULE
\end{minipage} \\
\midrule\noalign{}
\endhead
\bottomrule\noalign{}
\endlastfoot
\end{longtable}

\subsubsection{Description}\label{description-798}

The PBVD keyword defines the bubble-point pressure versus depth tables
for each equilibration region that should be used when there is
dissolved gas in the model (DISGAS has been activated in the RUNSPEC
section) and the EQLOPT1 variable has been set to a positive integer on
the EQUIL keyword in the SOLUTION section.

\begin{longtable}[]{@{}
  >{\raggedright\arraybackslash}p{(\linewidth - 6\tabcolsep) * \real{0.2500}}
  >{\raggedright\arraybackslash}p{(\linewidth - 6\tabcolsep) * \real{0.2500}}
  >{\raggedright\arraybackslash}p{(\linewidth - 6\tabcolsep) * \real{0.2500}}
  >{\raggedright\arraybackslash}p{(\linewidth - 6\tabcolsep) * \real{0.2500}}@{}}
\toprule\noalign{}
\begin{minipage}[b]{\linewidth}\raggedright
No.
\end{minipage} & \begin{minipage}[b]{\linewidth}\raggedright
Name
\end{minipage} & \begin{minipage}[b]{\linewidth}\raggedright
Description
\end{minipage} & \begin{minipage}[b]{\linewidth}\raggedright
Default
\end{minipage} \\
\midrule\noalign{}
\endhead
\bottomrule\noalign{}
\endlastfoot
Field & Metric & Laboratory & \\
1 & DEPTH & A columnar vector of real monotonically increasing down the
column values that defines the depth values for the corresponding
bubble-point values, PBVALS. & None \\
feet & m & cm & \\
2 & PBVALS & A columnar vector of real values that defines the oil
bubble-point values at the corresponding DEPTH. & None \\
psia & barsa & atma & \\
Notes: & & & \\
\end{longtable}

\emph{Table 10.20: PBVD Keyword Description}

Alternatively, the dissolved gas-oil ratio versus depth tables may be
entered using the RSVD keyword in the SOLUTION section instead of this
keyword. See also the RSVD and EQUIL keywords in the SOLUTION section.

\subsubsection{Example}\label{example-450}

Given NTEQUL equals three and NDRXVD is greater than or equal to two on
the EQLDIMS keyword in the RUNSPEC section, then the following example
defines the bubble-point versus depth functions.

\begin{verbatim}
--
--       DEPTH    PSAT
--                PRESS
--       ------   ------
PBVD
         3000.0   3000.0
         8000.0   3025.0                            / PSAT VS DEPTH EQUIL REGN 01
--       ------   ------
         3000.0   3100.0
         8000.0   3125.0                            / PSAT VS DEPTH EQUIL REGN 02
--       ------   ------
         3000.0   3200.0
         8000.0   3225.0                            / PSAT VS DEPTH EQUIL REGN 03
\end{verbatim}

Here three tables are entered and each table is terminated by a ``/''
and there is no keyword terminating ``/''.

\subsection{PDEW -- Define the Initial Equilibration Dew-Point Pressure
for All Grid
Blocks}\label{pdew-define-the-initial-equilibration-dew-point-pressure-for-all-grid-blocks}

\begin{longtable}[]{@{}
  >{\raggedright\arraybackslash}p{(\linewidth - 14\tabcolsep) * \real{0.1250}}
  >{\raggedright\arraybackslash}p{(\linewidth - 14\tabcolsep) * \real{0.1250}}
  >{\raggedright\arraybackslash}p{(\linewidth - 14\tabcolsep) * \real{0.1250}}
  >{\raggedright\arraybackslash}p{(\linewidth - 14\tabcolsep) * \real{0.1250}}
  >{\raggedright\arraybackslash}p{(\linewidth - 14\tabcolsep) * \real{0.1250}}
  >{\raggedright\arraybackslash}p{(\linewidth - 14\tabcolsep) * \real{0.1250}}
  >{\raggedright\arraybackslash}p{(\linewidth - 14\tabcolsep) * \real{0.1250}}
  >{\raggedright\arraybackslash}p{(\linewidth - 14\tabcolsep) * \real{0.1250}}@{}}
\toprule\noalign{}
\begin{minipage}[b]{\linewidth}\raggedright
RUNSPEC
\end{minipage} & \begin{minipage}[b]{\linewidth}\raggedright
GRID
\end{minipage} & \begin{minipage}[b]{\linewidth}\raggedright
EDIT
\end{minipage} & \begin{minipage}[b]{\linewidth}\raggedright
PROPS
\end{minipage} & \begin{minipage}[b]{\linewidth}\raggedright
REGIONS
\end{minipage} & \begin{minipage}[b]{\linewidth}\raggedright
SOLUTION
\end{minipage} & \begin{minipage}[b]{\linewidth}\raggedright
SUMMARY
\end{minipage} & \begin{minipage}[b]{\linewidth}\raggedright
SCHEDULE
\end{minipage} \\
\midrule\noalign{}
\endhead
\bottomrule\noalign{}
\endlastfoot
\end{longtable}

\subsubsection{Description}\label{description-799}

The PDEW keyword defines the initial equilibration dew-point pressure
values for all grid cells in the model and should be used in conjunction
with the PBUB, PRESSURE, RS, RV, SGAS, SOIL and SWAT keywords etc., to
fully describe the initial state of the model. The keyword should only
be used if vaporized oil been activated in the model via the VAPOIL
keyword in the RUNSPEC section.

This is the non-standard method to initialize the model via enumeration
and is seldom employed in the industry. The standard methodology is for
OPM Flow to initialize a model using the parameters on the EQUIL keyword
combined with other keywords to fully describe the initial state of the
model. The keyword can be used with all grid types.

This keyword is not supported by OPM Flow but would change the results
if supported so the simulation will be stopped.

\begin{longtable}[]{@{}
  >{\raggedright\arraybackslash}p{(\linewidth - 6\tabcolsep) * \real{0.2500}}
  >{\raggedright\arraybackslash}p{(\linewidth - 6\tabcolsep) * \real{0.2500}}
  >{\raggedright\arraybackslash}p{(\linewidth - 6\tabcolsep) * \real{0.2500}}
  >{\raggedright\arraybackslash}p{(\linewidth - 6\tabcolsep) * \real{0.2500}}@{}}
\toprule\noalign{}
\begin{minipage}[b]{\linewidth}\raggedright
No.
\end{minipage} & \begin{minipage}[b]{\linewidth}\raggedright
Name
\end{minipage} & \begin{minipage}[b]{\linewidth}\raggedright
Description
\end{minipage} & \begin{minipage}[b]{\linewidth}\raggedright
Default
\end{minipage} \\
\midrule\noalign{}
\endhead
\bottomrule\noalign{}
\endlastfoot
Field & Metric & Laboratory & \\
1 & PDEW & PDEW is an array of real positive numbers assigning the
initial equilibration dew-point pressure values to each cell in the
model. Repeat counts may be used, for example 20*3525.0 & None \\
psia & barsa & atma & \\
Notes: & & & \\
\end{longtable}

\emph{Table 10.21: PDEW Keyword Description}

See also the PBUB, PRESSURE, RS, RV, SGAS, SOIL and SWAT keywords to
fully define the initial state of the model.

\subsubsection{Example}\label{example-451}

\begin{verbatim}
--
--       DEFINE INITIAL EQUILIBRATION PSAT VALUES FOR ALL CELLS IN THE MODEL
--       BASED ON NX = 100, NY = 100 AND NZ = 3
--
PDEW
         1000*3500.0    1000*3525.0    1000*0.3535.0                           /
\end{verbatim}

The above example defines the initial equilibration dew-point saturation
pressure values to be 3500.0 for all the cells in the first layer,
3525.0 for all the cells in the second layer, and finally 3535.0 for all
the cells in the third layer.

\subsection{PDVD -- Define Equilibration Dew-Point versus Depth
Tables}\label{pdvd-define-equilibration-dew-point-versus-depth-tables}

\begin{longtable}[]{@{}
  >{\raggedright\arraybackslash}p{(\linewidth - 14\tabcolsep) * \real{0.1250}}
  >{\raggedright\arraybackslash}p{(\linewidth - 14\tabcolsep) * \real{0.1250}}
  >{\raggedright\arraybackslash}p{(\linewidth - 14\tabcolsep) * \real{0.1250}}
  >{\raggedright\arraybackslash}p{(\linewidth - 14\tabcolsep) * \real{0.1250}}
  >{\raggedright\arraybackslash}p{(\linewidth - 14\tabcolsep) * \real{0.1250}}
  >{\raggedright\arraybackslash}p{(\linewidth - 14\tabcolsep) * \real{0.1250}}
  >{\raggedright\arraybackslash}p{(\linewidth - 14\tabcolsep) * \real{0.1250}}
  >{\raggedright\arraybackslash}p{(\linewidth - 14\tabcolsep) * \real{0.1250}}@{}}
\toprule\noalign{}
\begin{minipage}[b]{\linewidth}\raggedright
RUNSPEC
\end{minipage} & \begin{minipage}[b]{\linewidth}\raggedright
GRID
\end{minipage} & \begin{minipage}[b]{\linewidth}\raggedright
EDIT
\end{minipage} & \begin{minipage}[b]{\linewidth}\raggedright
PROPS
\end{minipage} & \begin{minipage}[b]{\linewidth}\raggedright
REGIONS
\end{minipage} & \begin{minipage}[b]{\linewidth}\raggedright
SOLUTION
\end{minipage} & \begin{minipage}[b]{\linewidth}\raggedright
SUMMARY
\end{minipage} & \begin{minipage}[b]{\linewidth}\raggedright
SCHEDULE
\end{minipage} \\
\midrule\noalign{}
\endhead
\bottomrule\noalign{}
\endlastfoot
\end{longtable}

\subsubsection{Description}\label{description-800}

The PDVD keyword defines the dew-point pressure versus depth tables for
each equilibration region that should be used when there is vaporized
oil in the model (VAPOIL has been activated in the RUNSPEC section) and
the EQLOPT2 variable has been set to a positive integer on the EQUIL
keyword in the SOLUTION section.

\begin{longtable}[]{@{}
  >{\raggedright\arraybackslash}p{(\linewidth - 6\tabcolsep) * \real{0.2500}}
  >{\raggedright\arraybackslash}p{(\linewidth - 6\tabcolsep) * \real{0.2500}}
  >{\raggedright\arraybackslash}p{(\linewidth - 6\tabcolsep) * \real{0.2500}}
  >{\raggedright\arraybackslash}p{(\linewidth - 6\tabcolsep) * \real{0.2500}}@{}}
\toprule\noalign{}
\begin{minipage}[b]{\linewidth}\raggedright
No.
\end{minipage} & \begin{minipage}[b]{\linewidth}\raggedright
Name
\end{minipage} & \begin{minipage}[b]{\linewidth}\raggedright
Description
\end{minipage} & \begin{minipage}[b]{\linewidth}\raggedright
Default
\end{minipage} \\
\midrule\noalign{}
\endhead
\bottomrule\noalign{}
\endlastfoot
Field & Metric & Laboratory & \\
1 & DEPTH & A columnar vector of real monotonically increasing down the
column values that defines the depth values for the corresponding
dew-point values, PDVALS. & None \\
feet & m & cm & \\
2 & PDVALS & A columnar vector of real values that defines the gas
dew-point values at the corresponding DEPTH. & None \\
psia & barsa & atma & \\
Notes: & & & \\
\end{longtable}

\emph{Table 10.22: PDVD Keyword Description}

Alternatively, the vaporized oil-gas ratio (condensate-gas ratio) versus
depth tables may be entered using the RVVD keyword in the SOLUTION
section instead of this keyword.

See also the RVVD and EQUIL keywords in the SOLUTION section.

\subsubsection{Example}\label{example-452}

Given NTEQUL equals three and NDRXVD is greater than or equal to two on
the EQLDIMS keyword in the RUNSPEC section, then the following example
defines the bubble-point versus depth functions.

\begin{verbatim}
--
--       DEPTH    PSAT
--                PRESS
--       ------   ------
PDVD
         3000.0   2000.0
         8000.0   2025.0                            / PSAT VS DEPTH EQUIL REGN 01
--       ------   ------
         3000.0   2100.0
         8000.0   2125.0                            / PSAT VS DEPTH EQUIL REGN 02
--       ------   ------
         3000.0   2200.0
         8000.0   2225.0                            / PSAT VS DEPTH EQUIL REGN 03
\end{verbatim}

Here three tables are entered and each table is terminated by a ``/''
and there is no keyword terminating ``/''.

\subsection{PRESSURE -- Define the Initial Equilibration Pressures for
All Grid
Blocks}\label{pressure-define-the-initial-equilibration-pressures-for-all-grid-blocks}

\begin{longtable}[]{@{}
  >{\raggedright\arraybackslash}p{(\linewidth - 14\tabcolsep) * \real{0.1250}}
  >{\raggedright\arraybackslash}p{(\linewidth - 14\tabcolsep) * \real{0.1250}}
  >{\raggedright\arraybackslash}p{(\linewidth - 14\tabcolsep) * \real{0.1250}}
  >{\raggedright\arraybackslash}p{(\linewidth - 14\tabcolsep) * \real{0.1250}}
  >{\raggedright\arraybackslash}p{(\linewidth - 14\tabcolsep) * \real{0.1250}}
  >{\raggedright\arraybackslash}p{(\linewidth - 14\tabcolsep) * \real{0.1250}}
  >{\raggedright\arraybackslash}p{(\linewidth - 14\tabcolsep) * \real{0.1250}}
  >{\raggedright\arraybackslash}p{(\linewidth - 14\tabcolsep) * \real{0.1250}}@{}}
\toprule\noalign{}
\begin{minipage}[b]{\linewidth}\raggedright
RUNSPEC
\end{minipage} & \begin{minipage}[b]{\linewidth}\raggedright
GRID
\end{minipage} & \begin{minipage}[b]{\linewidth}\raggedright
EDIT
\end{minipage} & \begin{minipage}[b]{\linewidth}\raggedright
PROPS
\end{minipage} & \begin{minipage}[b]{\linewidth}\raggedright
REGIONS
\end{minipage} & \begin{minipage}[b]{\linewidth}\raggedright
SOLUTION
\end{minipage} & \begin{minipage}[b]{\linewidth}\raggedright
SUMMARY
\end{minipage} & \begin{minipage}[b]{\linewidth}\raggedright
SCHEDULE
\end{minipage} \\
\midrule\noalign{}
\endhead
\bottomrule\noalign{}
\endlastfoot
\end{longtable}

\subsubsection{Description}\label{description-801}

The PRESSURE keyword defines the initial equilibration pressures for all
grid cells in the model and should be used in conjunction with the PBUB,
PDEW, RS, RV, SGAS, SOIL and SWAT keywords etc., to fully describe the
initial state of the model.

The keyword is used by the Enumeration Initialization method to
initialize the model, as opposed to the Equilibration Initialization
method that utilizes the EQUIL keyword in the SOLUTION section. This is
the non-standard formulation to initialize the model and is seldom
employed in the industry. The standard methodology is for OPM Flow to
initialize a model using the parameters on the EQUIL keyword combined
with other keywords to fully describe the initial state of the model.
The keyword can be used with all grid types.

\begin{longtable}[]{@{}
  >{\raggedright\arraybackslash}p{(\linewidth - 6\tabcolsep) * \real{0.2500}}
  >{\raggedright\arraybackslash}p{(\linewidth - 6\tabcolsep) * \real{0.2500}}
  >{\raggedright\arraybackslash}p{(\linewidth - 6\tabcolsep) * \real{0.2500}}
  >{\raggedright\arraybackslash}p{(\linewidth - 6\tabcolsep) * \real{0.2500}}@{}}
\toprule\noalign{}
\begin{minipage}[b]{\linewidth}\raggedright
No.
\end{minipage} & \begin{minipage}[b]{\linewidth}\raggedright
Name
\end{minipage} & \begin{minipage}[b]{\linewidth}\raggedright
Description
\end{minipage} & \begin{minipage}[b]{\linewidth}\raggedright
Default
\end{minipage} \\
\midrule\noalign{}
\endhead
\bottomrule\noalign{}
\endlastfoot
Field & Metric & Laboratory & \\
1 & PRESSURE & PRESSURE is an array of real positive numbers assigning
the initial equilibration pressures to each cell in the model. Repeat
counts may be used, for example 20*4200.0. & None \\
psia & barsa & atma & \\
Notes: & & & \\
\end{longtable}

\emph{Table 10.23: PRESSURE Keyword Description}

See also the PBUB, PDEW, RS, RV, SGAS, SOIL and SWAT keywords to fully
define the initial state of the model.

\subsubsection{Example}\label{example-453}

\begin{verbatim}
--
--       DEFINE INITIAL EQUILIBRATION PRESSURES FOR ALL CELLS IN THE MODEL
--       BASED ON NX = 100, NY = 100 AND NZ = 3
--
PRESSURE
         1000*4500.0   1000*4510.0   1000*4520.0                               /
\end{verbatim}

The above example defines the initial equilibration pressures to be
4500.0 for all the cells in the first layer, 4510.0 for all the cells in
the second layer, and finally 4520.0 for all the cells in the third
layer.

\subsection{PRVD -- Define the Initial Equilibration Pressures versus
Depth}\label{prvd-define-the-initial-equilibration-pressures-versus-depth}

\begin{longtable}[]{@{}
  >{\raggedright\arraybackslash}p{(\linewidth - 14\tabcolsep) * \real{0.1250}}
  >{\raggedright\arraybackslash}p{(\linewidth - 14\tabcolsep) * \real{0.1250}}
  >{\raggedright\arraybackslash}p{(\linewidth - 14\tabcolsep) * \real{0.1250}}
  >{\raggedright\arraybackslash}p{(\linewidth - 14\tabcolsep) * \real{0.1250}}
  >{\raggedright\arraybackslash}p{(\linewidth - 14\tabcolsep) * \real{0.1250}}
  >{\raggedright\arraybackslash}p{(\linewidth - 14\tabcolsep) * \real{0.1250}}
  >{\raggedright\arraybackslash}p{(\linewidth - 14\tabcolsep) * \real{0.1250}}
  >{\raggedright\arraybackslash}p{(\linewidth - 14\tabcolsep) * \real{0.1250}}@{}}
\toprule\noalign{}
\begin{minipage}[b]{\linewidth}\raggedright
RUNSPEC
\end{minipage} & \begin{minipage}[b]{\linewidth}\raggedright
GRID
\end{minipage} & \begin{minipage}[b]{\linewidth}\raggedright
EDIT
\end{minipage} & \begin{minipage}[b]{\linewidth}\raggedright
PROPS
\end{minipage} & \begin{minipage}[b]{\linewidth}\raggedright
REGIONS
\end{minipage} & \begin{minipage}[b]{\linewidth}\raggedright
SOLUTION
\end{minipage} & \begin{minipage}[b]{\linewidth}\raggedright
SUMMARY
\end{minipage} & \begin{minipage}[b]{\linewidth}\raggedright
SCHEDULE
\end{minipage} \\
\midrule\noalign{}
\endhead
\bottomrule\noalign{}
\endlastfoot
\end{longtable}

\subsubsection{Description}\label{description-802}

The PRVD keyword defines the initial reservoir pressure versus depth and
should be used in conjunction with the PBUB, PDEW, RS, RV, SGAS, SOIL
and SWAT keywords etc., to fully describe the initial state of the
model. PRVD is an alternative to the PRESSURE keyword in the SOLUTION
section, that defines the initial equilibration pressures for all grid
cells in the model

The keyword is used by the Enumeration Initialization method to
initialize the model, as opposed to the Equilibration Initialization
method that utilizes the EQUIL keyword in the SOLUTION section. This is
the non-standard formulation to initialize the model and is seldom
employed in the industry. The standard methodology is for OPM Flow to
initialize a model using the parameters on the EQUIL keyword combined
with other keywords to fully describe the initial state of the model.
The keyword can be used with all grid types.

This keyword is not supported by OPM Flow but would change the results
if supported so the simulation will be stopped.

\begin{longtable}[]{@{}
  >{\raggedright\arraybackslash}p{(\linewidth - 6\tabcolsep) * \real{0.2500}}
  >{\raggedright\arraybackslash}p{(\linewidth - 6\tabcolsep) * \real{0.2500}}
  >{\raggedright\arraybackslash}p{(\linewidth - 6\tabcolsep) * \real{0.2500}}
  >{\raggedright\arraybackslash}p{(\linewidth - 6\tabcolsep) * \real{0.2500}}@{}}
\toprule\noalign{}
\begin{minipage}[b]{\linewidth}\raggedright
No.
\end{minipage} & \begin{minipage}[b]{\linewidth}\raggedright
Name
\end{minipage} & \begin{minipage}[b]{\linewidth}\raggedright
Description
\end{minipage} & \begin{minipage}[b]{\linewidth}\raggedright
Default
\end{minipage} \\
\midrule\noalign{}
\endhead
\bottomrule\noalign{}
\endlastfoot
Field & Metric & Laboratory & \\
1 & DEPTH & A columnar vector of real monotonically increasing down the
column values that defines the depth values for the corresponding
reservoir oil pressures values, PRESSURE. & None \\
feet & m & cm & \\
2 & PRESSURE & A columnar vector of real values that defines the initial
equilibration oil pressure values at the corresponding DEPTH. & None \\
psia & barsa & atma & \\
Notes: & & & \\
\end{longtable}

\emph{Table 10.24: PRVD Keyword Description}

See also the PBUB, PDEW, RS, RV, SGAS, SOIL and SWAT keywords to fully
define the initial state of the model.

\subsubsection{Example}\label{example-454}

Given NTEQUL equals three and NDRXVD is greater than or equal to five on
the EQLDIMS keyword in the RUNSPEC section, then the following example
defines the initial oil reservoir pressure versus depth

\begin{verbatim}
--
--       DEPTH    INIT
--                PRESS
--       ------   ------
PRVD
         3000.0   3000.0
         4000.0   3345.0
         5000.0   3690.0
         7000.0   4700.0
         7200.0   4769.0                            / POIL VS DEPTH EQUIL REGN 01
--       ------   ------
         3000.0   3100.0
         4000.0   3445.0
         5000.0   3790.0
         7000.0   4700.0
         7200.0   4769.0                            / POIL VS DEPTH EQUIL REGN 02
--       ------   ------
         3000.0   3150.0
         4000.0   3495.0
         5000.0   3840.0
         7000.0   4700.0
         7200.0   4769.0                            / POIL VS DEPTH EQUIL REGN 03
\end{verbatim}

Here three tables are entered and each table is terminated by a ``/''
and there is no keyword terminating ``/''.

\subsection{PYEND -- End the Definition of a PYINPUT
Section}\label{pyend-end-the-definition-of-a-pyinput-section-4}

The PYINPUT and PYEND keywords are a part of OPM Flow's Python scripting
facility that processes standard Python commands that can be used to
manipulate and define the simulators input parameters during processing
of the input deck. The main purpose of the facility is to script the
construction of the various keywords. A PYINPUT Definition Section is
terminated by a PYEND keyword on a separate single line.

See PYEND -- End the Definition of a PYINPUT Section in the GRID section
for a full description.

\subsection{PYINPUT -- Define the Start of a PYINPUT
Section}\label{pyinput-define-the-start-of-a-pyinput-section-4}

The PYINPUT and PYEND keywords are a part of OPM Flow's Python scripting
facility that processes standard Python commands that can be used to
manipulate and define the simulators input parameters during processing
of the input deck. The main purpose of the facility is to script the
construction of the various keywords used by the simulator. PYINPUT
declares the start of a PYINPUT Definition Section on a single separate
line, which is then followed by various standard Python commands, with
one command per line.

See PYINPUT -- Define the Start of a PYINPUT Section in the GRID section
for a full description.

\subsection{RAINFALL -- Constant Flux Aquifer Rainfall Flux by
Month}\label{rainfall-constant-flux-aquifer-rainfall-flux-by-month}

\begin{longtable}[]{@{}
  >{\raggedright\arraybackslash}p{(\linewidth - 14\tabcolsep) * \real{0.1250}}
  >{\raggedright\arraybackslash}p{(\linewidth - 14\tabcolsep) * \real{0.1250}}
  >{\raggedright\arraybackslash}p{(\linewidth - 14\tabcolsep) * \real{0.1250}}
  >{\raggedright\arraybackslash}p{(\linewidth - 14\tabcolsep) * \real{0.1250}}
  >{\raggedright\arraybackslash}p{(\linewidth - 14\tabcolsep) * \real{0.1250}}
  >{\raggedright\arraybackslash}p{(\linewidth - 14\tabcolsep) * \real{0.1250}}
  >{\raggedright\arraybackslash}p{(\linewidth - 14\tabcolsep) * \real{0.1250}}
  >{\raggedright\arraybackslash}p{(\linewidth - 14\tabcolsep) * \real{0.1250}}@{}}
\toprule\noalign{}
\begin{minipage}[b]{\linewidth}\raggedright
RUNSPEC
\end{minipage} & \begin{minipage}[b]{\linewidth}\raggedright
GRID
\end{minipage} & \begin{minipage}[b]{\linewidth}\raggedright
EDIT
\end{minipage} & \begin{minipage}[b]{\linewidth}\raggedright
PROPS
\end{minipage} & \begin{minipage}[b]{\linewidth}\raggedright
REGIONS
\end{minipage} & \begin{minipage}[b]{\linewidth}\raggedright
SOLUTION
\end{minipage} & \begin{minipage}[b]{\linewidth}\raggedright
SUMMARY
\end{minipage} & \begin{minipage}[b]{\linewidth}\raggedright
SCHEDULE
\end{minipage} \\
\midrule\noalign{}
\endhead
\bottomrule\noalign{}
\endlastfoot
\end{longtable}

\subsubsection{Description}\label{description-803}

This RAINFALL keyword defines the month by month rainfall flux for
constant flux aquifers.

This keyword is not supported by OPM Flow but would change the results
if supported so the simulation will be stopped.

\subsection{RBEDCONT -- Define River Grid Block Contact Area versus
Depth}\label{rbedcont-define-river-grid-block-contact-area-versus-depth}

\begin{longtable}[]{@{}
  >{\raggedright\arraybackslash}p{(\linewidth - 14\tabcolsep) * \real{0.1250}}
  >{\raggedright\arraybackslash}p{(\linewidth - 14\tabcolsep) * \real{0.1250}}
  >{\raggedright\arraybackslash}p{(\linewidth - 14\tabcolsep) * \real{0.1250}}
  >{\raggedright\arraybackslash}p{(\linewidth - 14\tabcolsep) * \real{0.1250}}
  >{\raggedright\arraybackslash}p{(\linewidth - 14\tabcolsep) * \real{0.1250}}
  >{\raggedright\arraybackslash}p{(\linewidth - 14\tabcolsep) * \real{0.1250}}
  >{\raggedright\arraybackslash}p{(\linewidth - 14\tabcolsep) * \real{0.1250}}
  >{\raggedright\arraybackslash}p{(\linewidth - 14\tabcolsep) * \real{0.1250}}@{}}
\toprule\noalign{}
\begin{minipage}[b]{\linewidth}\raggedright
RUNSPEC
\end{minipage} & \begin{minipage}[b]{\linewidth}\raggedright
GRID
\end{minipage} & \begin{minipage}[b]{\linewidth}\raggedright
EDIT
\end{minipage} & \begin{minipage}[b]{\linewidth}\raggedright
PROPS
\end{minipage} & \begin{minipage}[b]{\linewidth}\raggedright
REGIONS
\end{minipage} & \begin{minipage}[b]{\linewidth}\raggedright
SOLUTION
\end{minipage} & \begin{minipage}[b]{\linewidth}\raggedright
SUMMARY
\end{minipage} & \begin{minipage}[b]{\linewidth}\raggedright
SCHEDULE
\end{minipage} \\
\midrule\noalign{}
\endhead
\bottomrule\noalign{}
\endlastfoot
\end{longtable}

\subsubsection{Description}\label{description-804}

The RBEDCONT keyword defines the river grid block contact area versus
depth tables, for when the River option has been activated via the
RIVRDIMS keyword in the RUNSPEC section.

This keyword is not supported by OPM Flow but would change the results
if supported so the simulation will be stopped.

\subsection{REFINE -- Start the Definition of a Local Grid
Refinement}\label{refine-start-the-definition-of-a-local-grid-refinement-4}

The REFINE keyword defines the start of a Cartesian or radial local grid
refinement (``LGR'') definition that sets the properties of the selected
LGR. The keyword is then followed by the property keywords associated
with the section where the keyword is being invoked. For example, if the
REFINE keyword is used in the GRID section then most of the keywords in
that section can be used to set the grid properties for the LGR.

See REFINE -- Start the Definition of a Local Grid Refinement in the
GRID section for a full description.

\subsection{RESTART -- Restart Run From an Existing Restart
File}\label{restart-restart-run-from-an-existing-restart-file}

\begin{longtable}[]{@{}
  >{\raggedright\arraybackslash}p{(\linewidth - 14\tabcolsep) * \real{0.1250}}
  >{\raggedright\arraybackslash}p{(\linewidth - 14\tabcolsep) * \real{0.1250}}
  >{\raggedright\arraybackslash}p{(\linewidth - 14\tabcolsep) * \real{0.1250}}
  >{\raggedright\arraybackslash}p{(\linewidth - 14\tabcolsep) * \real{0.1250}}
  >{\raggedright\arraybackslash}p{(\linewidth - 14\tabcolsep) * \real{0.1250}}
  >{\raggedright\arraybackslash}p{(\linewidth - 14\tabcolsep) * \real{0.1250}}
  >{\raggedright\arraybackslash}p{(\linewidth - 14\tabcolsep) * \real{0.1250}}
  >{\raggedright\arraybackslash}p{(\linewidth - 14\tabcolsep) * \real{0.1250}}@{}}
\toprule\noalign{}
\begin{minipage}[b]{\linewidth}\raggedright
RUNSPEC
\end{minipage} & \begin{minipage}[b]{\linewidth}\raggedright
GRID
\end{minipage} & \begin{minipage}[b]{\linewidth}\raggedright
EDIT
\end{minipage} & \begin{minipage}[b]{\linewidth}\raggedright
PROPS
\end{minipage} & \begin{minipage}[b]{\linewidth}\raggedright
REGIONS
\end{minipage} & \begin{minipage}[b]{\linewidth}\raggedright
SOLUTION
\end{minipage} & \begin{minipage}[b]{\linewidth}\raggedright
SUMMARY
\end{minipage} & \begin{minipage}[b]{\linewidth}\raggedright
SCHEDULE
\end{minipage} \\
\midrule\noalign{}
\endhead
\bottomrule\noalign{}
\endlastfoot
\end{longtable}

\subsubsection{Description}\label{description-805}

The RESTART keyword defines the parameters to restart the simulation
from a previous run that has written a RESTART file out to disk. Only
restarting from RESTART files is permitted by OPM Flow; restarting from
SAVE files is not implemented.

\begin{longtable}[]{@{}
  >{\raggedright\arraybackslash}p{(\linewidth - 6\tabcolsep) * \real{0.2500}}
  >{\raggedright\arraybackslash}p{(\linewidth - 6\tabcolsep) * \real{0.2500}}
  >{\raggedright\arraybackslash}p{(\linewidth - 6\tabcolsep) * \real{0.2500}}
  >{\raggedright\arraybackslash}p{(\linewidth - 6\tabcolsep) * \real{0.2500}}@{}}
\toprule\noalign{}
\begin{minipage}[b]{\linewidth}\raggedright
No.
\end{minipage} & \begin{minipage}[b]{\linewidth}\raggedright
Name
\end{minipage} & \begin{minipage}[b]{\linewidth}\raggedright
Description
\end{minipage} & \begin{minipage}[b]{\linewidth}\raggedright
Default
\end{minipage} \\
\midrule\noalign{}
\endhead
\bottomrule\noalign{}
\endlastfoot
1 & RSNAME & The RSNAME variable is a character string that defines the
root name of the RESTART file to be read into the current input deck. &
None \\
2 & RSNUM & A positive integer that defines the restart point on the
RESTART file to be read and to be used to initialize the model. When OPM
Flow writes a restart point a message is printed to the *.PRT file
indicating the time step the restart was written out. & None \\
3 & RSTYPE & A defined character sting set to SAVE to read the restart
data from the SAVE file, otherwise defaulted to 1* to read the data from
the RESTART file. The SAVE file option is not supported by OPM Flow and
should be defaulted with 1*. & 1* \\
4 & RSFORMAT & A defined character string that defines the format of the
SAVE file to be read if RSTYPE has been set to SAVE, and should be set
to one of the following: If the variable RSFORMAT omitted then the
default is for binary file input. This option is not supported by OPM
Flow and should be defaulted with 1*. & U \\
Notes: & & & \\
\end{longtable}

\emph{Table 10.25: RESTART Keyword Description}

The most direct way to start a restart run is to:

\begin{verbatim}
- Copy the existing data file that created the RESTART file and give it a new name. For example if the RESTART file is from a case named NOR-OPM-A01.DATA, then the copied data file could be named NOR-OPM-A01-R1.DATA.
- Edit the copied data file (NOR-OPM-A01-R1.DATA) and delete all equilibration keywords (EQUIL, RSVD, etc.) or the enumeration equilibration keywords (PRESSURE, SGAS, SOIL. SWAT, etc.) in the SOLUTION section used to initialize the model.
- In the SOLUTION section of NOR-OPM-A01-R1.DATA file insert the RESTART keyword, using NOR-OPM-A01 as RSNAME and the required RSNUM value for the time step to restart from.
- In the SCHEDULE section of NOR-OPM-A01-R1.DATA file insert the SKIPREST keyword at the very beginning of the SCHEDULE section. The SKIPREST keyword causes the simulator to only read in data it requires for restarting the run up to the RESTART point (RSNUM). Note that certain keywords always need to be present in a restart run in the SCHEDULE section as the data is not stored on the RESTART file, for example the VFP tables (VFPPROD and VFPINJ keywords). The SKIPREST keyword automatically processes the input deck and reads the required data.
- In the SCHEDULE section of NOR-OPM-A01-R1.DATA file after the RESTART point make any required changes, save the file and run the NOR-OPM-A01-R1.DATA with OPM Flow.
\end{verbatim}

See also RPTRST, RPTSCHED and SKIPREST keywords.

\subsubsection{Example}\label{example-455}

The example below defines a restart from the previously run NOR-OPM-A01
case at time step number 40.

\begin{verbatim}
-- ==============================================================================
--
-- SOLUTION SECTION
--
-- ==============================================================================
SOLUTION
--
--       FLEXIBLE RESTART FROM PREVIOUS SIMULATION RUN
--
--       FILE                   RESTART   RESTART   FILE
--       NAME                   NUMBER    TYPE      FORMAT
RESTART
         'NOR-OPM-A01'          40        1*        1*     /
\end{verbatim}

In addition in the SCHEDULE section the SKIPREST keyword should be used
to correctly read in the schedule data up to the RESTART point.

\begin{verbatim}
-- ==============================================================================
--
-- SCHEDULE SECTION
--
-- ==============================================================================
SCHEDULE
--
--       ACTIVATE SKIPREST OPTION TO AVOID MODIFYING SCHEDULE SECTION
--
SKIPREST
\end{verbatim}

Note is advisable to place the SKIPREST keyword at the very beginning of
the SCHEDULE section.

\subsection{RIVERSYS -- Define River System (Branch Structure and
Boundary
Conditions)}\label{riversys-define-river-system-branch-structure-and-boundary-conditions}

\begin{longtable}[]{@{}
  >{\raggedright\arraybackslash}p{(\linewidth - 14\tabcolsep) * \real{0.1250}}
  >{\raggedright\arraybackslash}p{(\linewidth - 14\tabcolsep) * \real{0.1250}}
  >{\raggedright\arraybackslash}p{(\linewidth - 14\tabcolsep) * \real{0.1250}}
  >{\raggedright\arraybackslash}p{(\linewidth - 14\tabcolsep) * \real{0.1250}}
  >{\raggedright\arraybackslash}p{(\linewidth - 14\tabcolsep) * \real{0.1250}}
  >{\raggedright\arraybackslash}p{(\linewidth - 14\tabcolsep) * \real{0.1250}}
  >{\raggedright\arraybackslash}p{(\linewidth - 14\tabcolsep) * \real{0.1250}}
  >{\raggedright\arraybackslash}p{(\linewidth - 14\tabcolsep) * \real{0.1250}}@{}}
\toprule\noalign{}
\begin{minipage}[b]{\linewidth}\raggedright
RUNSPEC
\end{minipage} & \begin{minipage}[b]{\linewidth}\raggedright
GRID
\end{minipage} & \begin{minipage}[b]{\linewidth}\raggedright
EDIT
\end{minipage} & \begin{minipage}[b]{\linewidth}\raggedright
PROPS
\end{minipage} & \begin{minipage}[b]{\linewidth}\raggedright
REGIONS
\end{minipage} & \begin{minipage}[b]{\linewidth}\raggedright
SOLUTION
\end{minipage} & \begin{minipage}[b]{\linewidth}\raggedright
SUMMARY
\end{minipage} & \begin{minipage}[b]{\linewidth}\raggedright
SCHEDULE
\end{minipage} \\
\midrule\noalign{}
\endhead
\bottomrule\noalign{}
\endlastfoot
\end{longtable}

\subsubsection{Description}\label{description-806}

RIVERSYS defines a river system by specifying the branch structure of
the river together with the branch's associated boundary conditions, for
when the River option has been activated via the RIVRDIMS keyword in the
RUNSPEC section.

This keyword is not supported by OPM Flow but would change the results
if supported so the simulation will be stopped.

\subsection{RPTRST -- Define Data to be Written to the RESTART
File}\label{rptrst-define-data-to-be-written-to-the-restart-file}

\begin{longtable}[]{@{}
  >{\raggedright\arraybackslash}p{(\linewidth - 14\tabcolsep) * \real{0.1250}}
  >{\raggedright\arraybackslash}p{(\linewidth - 14\tabcolsep) * \real{0.1250}}
  >{\raggedright\arraybackslash}p{(\linewidth - 14\tabcolsep) * \real{0.1250}}
  >{\raggedright\arraybackslash}p{(\linewidth - 14\tabcolsep) * \real{0.1250}}
  >{\raggedright\arraybackslash}p{(\linewidth - 14\tabcolsep) * \real{0.1250}}
  >{\raggedright\arraybackslash}p{(\linewidth - 14\tabcolsep) * \real{0.1250}}
  >{\raggedright\arraybackslash}p{(\linewidth - 14\tabcolsep) * \real{0.1250}}
  >{\raggedright\arraybackslash}p{(\linewidth - 14\tabcolsep) * \real{0.1250}}@{}}
\toprule\noalign{}
\begin{minipage}[b]{\linewidth}\raggedright
RUNSPEC
\end{minipage} & \begin{minipage}[b]{\linewidth}\raggedright
GRID
\end{minipage} & \begin{minipage}[b]{\linewidth}\raggedright
EDIT
\end{minipage} & \begin{minipage}[b]{\linewidth}\raggedright
PROPS
\end{minipage} & \begin{minipage}[b]{\linewidth}\raggedright
REGIONS
\end{minipage} & \begin{minipage}[b]{\linewidth}\raggedright
SOLUTION
\end{minipage} & \begin{minipage}[b]{\linewidth}\raggedright
SUMMARY
\end{minipage} & \begin{minipage}[b]{\linewidth}\raggedright
SCHEDULE
\end{minipage} \\
\midrule\noalign{}
\endhead
\bottomrule\noalign{}
\endlastfoot
\end{longtable}

\subsubsection{Description}\label{description-807}

This keyword defines the data to be written to the RESTART file and the
frequency at which restart points will be created. In addition to the
solution data arrays required to restart a run, the user may request
extra data to be written to the restart file for visualization in OPM
ResInsight.

The RPTRST keyword should be followed by a series of character strings
that indicate the data to be written. In most cases the character string
is the keyword used to load the data in the OPM Flow input deck, for
example KRG for the gas relative permeability of each grid block at the
requested times. It is anticipated that OPM Flow will support additional
functionality as development progresses.

\begin{longtable}[]{@{}
  >{\raggedright\arraybackslash}p{(\linewidth - 6\tabcolsep) * \real{0.2500}}
  >{\raggedright\arraybackslash}p{(\linewidth - 6\tabcolsep) * \real{0.2500}}
  >{\raggedright\arraybackslash}p{(\linewidth - 6\tabcolsep) * \real{0.2500}}
  >{\raggedright\arraybackslash}p{(\linewidth - 6\tabcolsep) * \real{0.2500}}@{}}
\toprule\noalign{}
\begin{minipage}[b]{\linewidth}\raggedright
No.
\end{minipage} & \begin{minipage}[b]{\linewidth}\raggedright
Mnemonic
\end{minipage} & \begin{minipage}[b]{\linewidth}\raggedright
Description
\end{minipage} & \begin{minipage}[b]{\linewidth}\raggedright
Array Name
\end{minipage} \\
\midrule\noalign{}
\endhead
\bottomrule\noalign{}
\endlastfoot
1 & ALLPROPS & An alias for the BG, BO, BW, DEN, KRG, KRO, KRW, VGAS,
VOIL and VWAT mnemonics combined that writes all the associated
properties. & \\
2 & BASIC & The BASIC mnemonic defines the frequency at which the
restart data and the additional requested data is written to the RESTART
file. The mnemonic is assigned a value, using the form BASIC=OPTION,
where OPTION is an integer variable set to one of the following: & \\
3 & BG & Gas formation volume factor. & 1OVERBG \\
4 & BO & Oil formation volume factor. & 1OVERBO \\
5 & BW & Water formation volume factor. & 1OVERBW \\
6 & CONV & Identify the worst n cells causing convergence problems based
on the residuals of each equation (CNV\_XXX), using the form CONV=n.~In
addition, the number of Newton iterations required by each cell in order
to satisfy the solution change convergence criteria at the last time
step (CONV\_NEW). & CNV\_GAS CNV\_OIL CNV\_WAT CONV\_NEW \\
7 & DEN & Fluid phase densities at reservoir conditions. & GAS\_DEN
OIL\_DEN WAT\_DEN \\
8 & FIP & Volumes in place for each phase at surface conditions. &
FIPGAS FIPOIL FIPWAT \\
9 & FLOWS & Inter-block volumetric flows at surface conditions for each
phase (GAS, OIL, WAT) in the I+, J+ and K+ directions. The Array Name
column shows the arrays for just the GAS phase, the other phases will
have similar arrays. & FLOGASI+ FLOGASJ+ FLOGASK+ \\
10 & FLOWS- & Inter-block volumetric flows at surface conditions for
each phase (GAS, OIL, WAT) in the I-, J- and K- directions. The Array
Name column shows the arrays for just the GAS phase, the other phases
will have similar arrays. & FLOGASI- FLOGASJ- FLOGASK- \\
11 & FLORES & Inter-block volumetric flows at reservoir conditions for
each phase (GAS, OIL, WAT) in the I+, J+ and K+ directions. The Array
Name column shows the arrays for just the GAS phase, the other phases
will have similar arrays. & FLRGASI+ FLRGASJ+ FLRGASK+ \\
12 & FLORES- & Inter-block volumetric flows at reservoir conditions for
each phase (GAS, OIL, WAT) in the I-, J- and K- directions. The Array
Name column shows the arrays for just the GAS phase, the other phases
will have similar arrays. & FLRGASI- FLRGASJ- FLRGASK- \\
13 & FREQ & The FREQ mnemonic defines the frequency at which restart
data is written when BASIC=3, 4 or 5, using the form FREQ=n. & \\
14 & KRG & Gas relative permeability at the grid block's gas saturation.
& GASKR \\
15 & KRO & Oil relative permeability at the grid block's oil saturation.
& OILKR \\
16 & KRW & Water relative permeability at the grid block's water
saturation. & WATKR \\
17 & PBPD & Bubble point and dew point pressures. & PBUB PDEW \\
18 & PCGW & Gas-water capillary pressure (gas-water simulations). &
PCGW \\
19 & PCOG & Gas-oil capillary pressure. & PCOG \\
20 & PCOW & Oil-water capillary pressure. & PCOW \\
21 & RESIDUAL & Component residuals after the last Newton iteration. &
RES\_GAS RES\_OIL RES\_WAT \\
22 & RFIP & Volumes in place for each phase at reservoir conditions. &
RFIPGAS RFIPOIL RFIPWAT \\
23 & ROCKC & Rock compaction: Overburden pressure, minimum oil pressure,
maximum water saturation, pore volume multiplier and transmissibility
multiplier. & PRES\_OVB PRESROCC SWMAX PORV\_RC TMULT\_RC \\
24 & RSSAT & Saturated dissolved gas-oil ratio. & RSSAT \\
25 & RSWSAT & Saturated dissolved gas-water ratio. This is an OPM Flow
specific mnemonic that requires the command line option
--enable-opm-rst-file=true is specified. & RSWSAT \\
26 & RVSAT & Saturated vaporized oil-gas ratio. & RVSAT \\
27 & RVWSAT & Saturated vaporized water-gas ratio. This is an OPM Flow
specific mnemonic that requires the command line option
--enable-opm-rst-file=true is specified. & RVWSAT \\
28 & SFIP & Volumes in place for each phase at surface/separator
conditions. & SFIPGAS SFIPOIL SFIPWAT \\
29 & VGAS & Gas viscosity at reservoir conditions. & GAS\_VISC \\
30 & VISC & Fluid phase viscosities at reservoir conditions. & GAS\_VISC
OIL\_VISC WAT\_VISC \\
31 & VOIL & Oil viscosity at reservoir conditions. & OIL\_VISC \\
32 & VWAT & Water viscosity at reservoir conditions. & WAT\_VISC \\
Notes: & & & \\
\end{longtable}

\emph{Table 10.26: RPTRST Keyword Description}

Note, older versions of the commercial simulator used integer values to
control the output to the RESTART file, this type of format is not
supported by OPM Flow and thus the integers should be converted to the
mnemonic format as outlined in Table 10.26.

Note that OPM Flow automatically writes out all the data required to
make a restart run as outlined in the table below:

\begin{longtable}[]{@{}
  >{\raggedright\arraybackslash}p{(\linewidth - 6\tabcolsep) * \real{0.2500}}
  >{\raggedright\arraybackslash}p{(\linewidth - 6\tabcolsep) * \real{0.2500}}
  >{\raggedright\arraybackslash}p{(\linewidth - 6\tabcolsep) * \real{0.2500}}
  >{\raggedright\arraybackslash}p{(\linewidth - 6\tabcolsep) * \real{0.2500}}@{}}
\toprule\noalign{}
\begin{minipage}[b]{\linewidth}\raggedright
No.
\end{minipage} & \begin{minipage}[b]{\linewidth}\raggedright
Mnemonic
\end{minipage} & \begin{minipage}[b]{\linewidth}\raggedright
Description
\end{minipage} & \begin{minipage}[b]{\linewidth}\raggedright
Array Name
\end{minipage} \\
\midrule\noalign{}
\endhead
\bottomrule\noalign{}
\endlastfoot
1 & BIOFILM & Biofilm volume fraction used when the
\hyperref[REF_HEADING_KEYWORD_BIOFILM]{BIOFILM} or MICP option has been
activated in the RUNSPEC section. & BIOFILM \\
2 & CALCITE & Calcite volume fraction used when the MICP option has been
activated in the RUNSPEC section. & CALCITE \\
3 & FOAM & Foam concentration used when the FOAM option has been
activated in the RUNSPEC section. & FOAM \\
4 & MICROBES & Microbe concentration used when the
\hyperref[REF_HEADING_KEYWORD_BIOFILM]{BIOFILM} or MICP option has been
activated in the RUNSPEC section. & MICROBES \\
5 & MULTPV & Pore volume multiplier output when MULTPV keyword is
present in the input deck. & MULTPV \\
6 & OXYGEN & Oxygen concentration used when the MICP option has been
activated in the RUNSPEC section. & OXYGEN \\
7 & PERMFACT & Permeability multiplication factor as a function of
porosity change used when the PRECSALT option has been activated in the
RUNSPEC section & PERMFACT \\
8 & POLYMER & Polymer concentration used when the POLYMER option has
been activated in the RUNSPEC section & POLYMER \\
9 & PPCW & Modified oil-water capillary pressure end-points used when
the SWATINIT keyword has been included in the input deck. & PPCW \\
10 & PRESSURE & Pressure data for each grid block. & PRESSURE \\
11 & RS & Dissolved gas-oil ratio for each grid block. & RS \\
12 & \hyperref[REF_HEADING_KEYWORD_RSW_10_3]{RSW} & Dissolved gas-water
ratio for each grid block for use when the DISGASW option has been
activated in the RUNSPEC section. &
\hyperref[REF_HEADING_KEYWORD_RSW_10_3]{RSW} \\
13 & RSWSOL & Dissolved solvent-water ratio for each grid block for use
when the SOLVENT and DISGASW options have been activated in the RUNSPEC
section. & RSWSOL \\
14 & RV & Vaporized oil-gas ratio for each grid block. & RV \\
15 & RVW & Vaporized water-gas ratio for each grid block for use when
the VAPWAT option has been activated in the RUNSPEC section. & RVW \\
16 & SALT & Salt concentration for each grid block for use when the
BRINE option has been activated in the RUNSPEC section. & SALT \\
17 & SALTP & Precipitated salt volume fraction for each grid block for
use when the PRECSALT option has been activated in the RUNSPEC section.
& SALTP \\
18 & SGAS & Gas saturation for each grid block. & SGAS \\
19 & SGMAX & Maximum gas saturation reached used when gas-oil
non-wetting phase relative permeability hysteresis has been activated. &
SGMAX \\
20 & SHMAX & Maximum wetting phase saturation reached used when gas-oil
wetting phase relative permeability hysteresis has been activated. &
SHMAX \\
21 & SOIL & Oil saturation each grid block. & SOIL \\
22 & SOMAX & Maximum oil saturation reached used when the VAPPARS
keyword has been included in the input deck, or when oil-water
non-wetting phase relative permeability hysteresis has been activated. &
SOMAX \\
23 & SOMIN & Minimum oil saturation reached used when gas-oil capillary
pressure hysteresis has been activated. & SOMIN \\
24 & SSOLVENT & Solvent saturation for each grid block used when the
SOLVENT option has been activated in the RUNSPEC section. & SSOLVENT \\
25 & SWAT & Water saturation for each grid block. & SWAT \\
26 & SWHY1 & Minimum water saturation reached used when oil-water
capillary pressure hysteresis has been activated. & SWHY1 \\
27 & SWMAX & Maximum water saturation reached used when oil-water
wetting phase relative permeability hysteresis has been activated. &
SWMAX \\
28 & TEMP & Temperature of each grid block used when the THERMAL option
has been activated in the RUNSPEC section. OPM Flow also supports the
output of temperatures for non-thermal simulations when the TEMP
mnemonic is specified. & TEMP \\
29 & UREA & Urea concentration used when the MICP option has been
activated in the RUNSPEC section. & UREA \\
Notes: & & & \\
\end{longtable}

\emph{Table 10.27: Data Sets Automatically Written to the RESTART File}

\begin{tcolorbox}[enhanced jigsaw, title=\textcolor{quarto-callout-note-color}{\faInfo}\hspace{0.5em}{Note}, coltitle=black, breakable, toprule=.15mm, colbacktitle=quarto-callout-note-color!10!white, leftrule=.75mm, opacitybacktitle=0.6, colback=white, bottomtitle=1mm, toptitle=1mm, titlerule=0mm, left=2mm, colframe=quarto-callout-note-color-frame, arc=.35mm, rightrule=.15mm, bottomrule=.15mm, opacityback=0]

Currently, OPM Flow distinguishes between those arrays which are
required for restarting a run and those which are ``merely'' for
visualization and analysis. However, this classification is imperfect
and leads to potentially exporting arrays that are incompatible with the
commercial simulator, such as the TEMP array in non-thermal simulation
runs. If one is willing to forego the ability to restart an OPM Flow
simulation run, using the commercial simulator, e.g., simulating the
historic period using OPM Flow and the prediction period using the
commercial simulator, then additional arrays are available, including
the PBUB and PDEW vectors generated by the PBPD option, by using the
following command line option: --enable-opm-rst-file=true

\end{tcolorbox}

If the --enable-opm-rst-file=true command line option has been specified
then the following OPM Flow specific data are also automatically output
for visualization only as outlined in the table below:

\begin{longtable}[]{@{}
  >{\raggedright\arraybackslash}p{(\linewidth - 6\tabcolsep) * \real{0.2500}}
  >{\raggedright\arraybackslash}p{(\linewidth - 6\tabcolsep) * \real{0.2500}}
  >{\raggedright\arraybackslash}p{(\linewidth - 6\tabcolsep) * \real{0.2500}}
  >{\raggedright\arraybackslash}p{(\linewidth - 6\tabcolsep) * \real{0.2500}}@{}}
\toprule\noalign{}
\begin{minipage}[b]{\linewidth}\raggedright
No.
\end{minipage} & \begin{minipage}[b]{\linewidth}\raggedright
Mnemonic
\end{minipage} & \begin{minipage}[b]{\linewidth}\raggedright
Description
\end{minipage} & \begin{minipage}[b]{\linewidth}\raggedright
Array Name
\end{minipage} \\
\midrule\noalign{}
\endhead
\bottomrule\noalign{}
\endlastfoot
1 & XMFCO2 & CO2 liquid-phase mole fractions written when the CO2STORE
option has been activated in the RUNSPEC section. & XMFCO2 \\
2 & XMFH2 & H2 liquid-phase mole fractions written when the
\hyperref[REF_HEADING_KEYWORD_H2STORE]{H2STORE} option has been
activated in the RUNSPEC section. & XMFH2 \\
3 & YMFWAT & Water gas-phase mole fractions written when either the
CO2STORE or \hyperref[REF_HEADING_KEYWORD_H2STORE]{H2STORE} option has
been activated in the RUNSPEC section. & YMFWAT \\
Notes: & & & \\
\end{longtable}

\emph{Table 10.27: Data Sets Automatically Written to the RESTART File
for Visualization Only}

\subsubsection{Examples}\label{examples-57}

The first example request that the standard restart data be written out
every month.

\begin{verbatim}
--
--       RESTART CONTROL BASIC = 4 (YEARLY) 5 (MONTHLY)
--
RPTRST
         BASIC=5                                                               /
\end{verbatim}

The next example requests that the standard restart data be written at
every report time step until this switch is reset and all the restarts
are kept. In addition to the standard the data the gas, oil and water
relative permeability data will also be written out at each report time
step.

\begin{verbatim}
--
--       RESTART CONTROL BASIC = 4 (YEARLY) 5 (MONTHLY)
--
RPTRST
         BASIC=2  KRG   KRO   KRW                                             /
\end{verbatim}

\subsection{RPTSOL -- Define SOLUTION Section
Reporting}\label{rptsol-define-solution-section-reporting}

\begin{longtable}[]{@{}
  >{\raggedright\arraybackslash}p{(\linewidth - 14\tabcolsep) * \real{0.1250}}
  >{\raggedright\arraybackslash}p{(\linewidth - 14\tabcolsep) * \real{0.1250}}
  >{\raggedright\arraybackslash}p{(\linewidth - 14\tabcolsep) * \real{0.1250}}
  >{\raggedright\arraybackslash}p{(\linewidth - 14\tabcolsep) * \real{0.1250}}
  >{\raggedright\arraybackslash}p{(\linewidth - 14\tabcolsep) * \real{0.1250}}
  >{\raggedright\arraybackslash}p{(\linewidth - 14\tabcolsep) * \real{0.1250}}
  >{\raggedright\arraybackslash}p{(\linewidth - 14\tabcolsep) * \real{0.1250}}
  >{\raggedright\arraybackslash}p{(\linewidth - 14\tabcolsep) * \real{0.1250}}@{}}
\toprule\noalign{}
\begin{minipage}[b]{\linewidth}\raggedright
RUNSPEC
\end{minipage} & \begin{minipage}[b]{\linewidth}\raggedright
GRID
\end{minipage} & \begin{minipage}[b]{\linewidth}\raggedright
EDIT
\end{minipage} & \begin{minipage}[b]{\linewidth}\raggedright
PROPS
\end{minipage} & \begin{minipage}[b]{\linewidth}\raggedright
REGIONS
\end{minipage} & \begin{minipage}[b]{\linewidth}\raggedright
SOLUTION
\end{minipage} & \begin{minipage}[b]{\linewidth}\raggedright
SUMMARY
\end{minipage} & \begin{minipage}[b]{\linewidth}\raggedright
SCHEDULE
\end{minipage} \\
\midrule\noalign{}
\endhead
\bottomrule\noalign{}
\endlastfoot
\end{longtable}

\subsubsection{Description}\label{description-808}

This keyword defines the data in the SOLUTION section that is to be
printed to the output print file in human readable format. The keyword
has two distinct forms, the first of which consists of the keyword
followed by a series of integers on the next line indicating the data to
be printed (see the first example). This is the original formal in the
commercial simulator and was subsequently superseded by the second
format. The second format consists of the keyword followed by a series
of character strings that indicate the data to be printed. In most cases
the character string is the keyword used to load the data in the OPM
Flow input deck, for example PVDG for the dry gas PVT tables. Its is
anticipated that OPM Flow will eventually support the functionality of
the second format only, the first format although recognized will be
completely ignored.

OPM Flow provides only limited supported for this keyword and will
ignore the unsupported options because they have no effect on the
results.

\begin{longtable}[]{@{}
  >{\raggedright\arraybackslash}p{(\linewidth - 6\tabcolsep) * \real{0.2500}}
  >{\raggedright\arraybackslash}p{(\linewidth - 6\tabcolsep) * \real{0.2500}}
  >{\raggedright\arraybackslash}p{(\linewidth - 6\tabcolsep) * \real{0.2500}}
  >{\raggedright\arraybackslash}p{(\linewidth - 6\tabcolsep) * \real{0.2500}}@{}}
\toprule\noalign{}
\begin{minipage}[b]{\linewidth}\raggedright
No.
\end{minipage} & \begin{minipage}[b]{\linewidth}\raggedright
Name
\end{minipage} & \begin{minipage}[b]{\linewidth}\raggedright
Description
\end{minipage} & \begin{minipage}[b]{\linewidth}\raggedright
Default
\end{minipage} \\
\midrule\noalign{}
\endhead
\bottomrule\noalign{}
\endlastfoot
1 & DENO & Print the oil reservoir density array & N/A \\
2 & EQUIL & Print the equilibration report. & N/A \\
3 & FIP & Print the fluid in-place report. The parameter is assigned a
value, OPTION, using the form FIP=OPTION, where OPTION is an integer
variable set to: & FIP=2 \\
4 & FIPRESV & Print the reservoir volumes in-place report. & N/A \\
5 & WELSPECS & WELSPECS switches on reporting of the well connections,
wells and groups at each report time step. There are numerous reports
associated with this option. Unlike the other reporting parameters that
produce a report for each reporting time step, the WELSPECS report
option only produces a report if an associated keyword has been
activated at the current reporting time step. For example, if the
reporting time steps are January, February, and March 2020, and the
RPTSCHED WELSPECS option is activated in January, with wells OP01 and
OP02 being declared via the WELSPECS and COMPDAT keywords, then a report
will be printed for January for these two wells. If there are no further
well activations until March, with well OP03 being declared, then there
will be no report for February, and only well OP03 will reported at the
March reporting time step. & N/A \\
Notes: & & & \\
\end{longtable}

\emph{Table 10.3.80.1: RPTSOL Keyword Description}

\begin{tcolorbox}[enhanced jigsaw, title=\textcolor{quarto-callout-note-color}{\faInfo}\hspace{0.5em}{Note}, coltitle=black, breakable, toprule=.15mm, colbacktitle=quarto-callout-note-color!10!white, leftrule=.75mm, opacitybacktitle=0.6, colback=white, bottomtitle=1mm, toptitle=1mm, titlerule=0mm, left=2mm, colframe=quarto-callout-note-color-frame, arc=.35mm, rightrule=.15mm, bottomrule=.15mm, opacityback=0]

Except for non-array like data, FIP etc., this keyword has the potential
to produce very large print files that some text editors may have
difficulty loading. A more efficient solution for array type data is to
load the \emph{.INIT and }.RESTART files into OPM ResInsight to view the
data graphically, this also has the benefit of being able to filter the
grid based on I, J, K ranges and grid properties.

\end{tcolorbox}

\subsubsection{Examples}\label{examples-58}

The first example shows the original format of this keyword; although
the keyword and format are recognized by OPM Flow, the format is ignored
and is unlikely to be implemented in in the simulator.

\begin{verbatim}
--
--       DEFINE SOLUTION SECTION REPORT OPTION (ORIGINAL FORMAT)
–
RPTSOL
         1        2*0      1        3*1                                        /
\end{verbatim}

The next example shows the second format of the keyword which may be
supported in a future release of OPM Flow.

\begin{verbatim}
--
--       DEFINE SOLUTION SECTION REPORT OPTIONS
--
RPTSOL
         FIP=2    FIPRESV  RESTART=3                                           /

\end{verbatim}

\subsection{RS -- Define the Initial Equilibration GOR (Rs) for All Grid
Blocks}\label{rs-define-the-initial-equilibration-gor-rs-for-all-grid-blocks}

\begin{longtable}[]{@{}
  >{\raggedright\arraybackslash}p{(\linewidth - 14\tabcolsep) * \real{0.1250}}
  >{\raggedright\arraybackslash}p{(\linewidth - 14\tabcolsep) * \real{0.1250}}
  >{\raggedright\arraybackslash}p{(\linewidth - 14\tabcolsep) * \real{0.1250}}
  >{\raggedright\arraybackslash}p{(\linewidth - 14\tabcolsep) * \real{0.1250}}
  >{\raggedright\arraybackslash}p{(\linewidth - 14\tabcolsep) * \real{0.1250}}
  >{\raggedright\arraybackslash}p{(\linewidth - 14\tabcolsep) * \real{0.1250}}
  >{\raggedright\arraybackslash}p{(\linewidth - 14\tabcolsep) * \real{0.1250}}
  >{\raggedright\arraybackslash}p{(\linewidth - 14\tabcolsep) * \real{0.1250}}@{}}
\toprule\noalign{}
\begin{minipage}[b]{\linewidth}\raggedright
RUNSPEC
\end{minipage} & \begin{minipage}[b]{\linewidth}\raggedright
GRID
\end{minipage} & \begin{minipage}[b]{\linewidth}\raggedright
EDIT
\end{minipage} & \begin{minipage}[b]{\linewidth}\raggedright
PROPS
\end{minipage} & \begin{minipage}[b]{\linewidth}\raggedright
REGIONS
\end{minipage} & \begin{minipage}[b]{\linewidth}\raggedright
SOLUTION
\end{minipage} & \begin{minipage}[b]{\linewidth}\raggedright
SUMMARY
\end{minipage} & \begin{minipage}[b]{\linewidth}\raggedright
SCHEDULE
\end{minipage} \\
\midrule\noalign{}
\endhead
\bottomrule\noalign{}
\endlastfoot
\end{longtable}

\subsubsection{Description}\label{description-809}

The RS keyword defines the initial equilibration gas-oil ratio values
for all grid cells in the model and should be used in conjunction with
the PBUB, PDEW, PRESSURE, RV, SGAS, SOIL and SWAT keywords etc., to
fully describe the initial state of the model. The keyword should only
be used if dissolved gas has been activated in the model via the DISGAS
keyword in the RUNSPEC section.

This is the non-standard method to initialize the model via enumeration
and is seldom employed in the industry. The standard methodology is for
OPM Flow to initialize a model using the parameters on the EQUIL keyword
combined with other keywords to fully describe the initial state of the
model. The keyword can be used with all grid types.

\begin{longtable}[]{@{}
  >{\raggedright\arraybackslash}p{(\linewidth - 6\tabcolsep) * \real{0.2500}}
  >{\raggedright\arraybackslash}p{(\linewidth - 6\tabcolsep) * \real{0.2500}}
  >{\raggedright\arraybackslash}p{(\linewidth - 6\tabcolsep) * \real{0.2500}}
  >{\raggedright\arraybackslash}p{(\linewidth - 6\tabcolsep) * \real{0.2500}}@{}}
\toprule\noalign{}
\begin{minipage}[b]{\linewidth}\raggedright
No.
\end{minipage} & \begin{minipage}[b]{\linewidth}\raggedright
Name
\end{minipage} & \begin{minipage}[b]{\linewidth}\raggedright
Description
\end{minipage} & \begin{minipage}[b]{\linewidth}\raggedright
Default
\end{minipage} \\
\midrule\noalign{}
\endhead
\bottomrule\noalign{}
\endlastfoot
Field & Metric & Laboratory & \\
1 & RS & RS is an array of real positive numbers assigning the initial
equilibration gas-oil ratio values to each cell in the model. Repeat
counts may be used, for example 20*1.30. & None \\
Mscf/stb & sm3/sm3 & scc/scc & \\
Notes: & & & \\
\end{longtable}

\emph{Table 10.29: RS Keyword Description}

See also the PBUB, PDEW, PRESSURE, RV, SGAS, SOIL and SWAT keywords to
fully define the initial state of the model.

\subsubsection{Example}\label{example-456}

\begin{verbatim}
--
--       DEFINE INITIAL EQUILIBRATION GOR VALUES FOR ALL CELLS IN THE MODEL
--       BASED ON NX = 100, NY = 100 AND NZ = 3
--
RS
         1000*1.3500   1000*1.3010   1000*1.3000                               /
\end{verbatim}

The above example defines the initial equilibration GOR values to be
1.3500 for all the cells in the first layer, 1.3010 for all the cells in
the second layer, and finally 1.3000 for all the cells in the third
layer.

\subsection{RSVD -- Equilibration Dissolved Gas-Oil Ratio (Rs) versus
Depth
Tables}\label{rsvd-equilibration-dissolved-gas-oil-ratio-rs-versus-depth-tables}

\begin{longtable}[]{@{}
  >{\raggedright\arraybackslash}p{(\linewidth - 14\tabcolsep) * \real{0.1250}}
  >{\raggedright\arraybackslash}p{(\linewidth - 14\tabcolsep) * \real{0.1250}}
  >{\raggedright\arraybackslash}p{(\linewidth - 14\tabcolsep) * \real{0.1250}}
  >{\raggedright\arraybackslash}p{(\linewidth - 14\tabcolsep) * \real{0.1250}}
  >{\raggedright\arraybackslash}p{(\linewidth - 14\tabcolsep) * \real{0.1250}}
  >{\raggedright\arraybackslash}p{(\linewidth - 14\tabcolsep) * \real{0.1250}}
  >{\raggedright\arraybackslash}p{(\linewidth - 14\tabcolsep) * \real{0.1250}}
  >{\raggedright\arraybackslash}p{(\linewidth - 14\tabcolsep) * \real{0.1250}}@{}}
\toprule\noalign{}
\begin{minipage}[b]{\linewidth}\raggedright
RUNSPEC
\end{minipage} & \begin{minipage}[b]{\linewidth}\raggedright
GRID
\end{minipage} & \begin{minipage}[b]{\linewidth}\raggedright
EDIT
\end{minipage} & \begin{minipage}[b]{\linewidth}\raggedright
PROPS
\end{minipage} & \begin{minipage}[b]{\linewidth}\raggedright
REGIONS
\end{minipage} & \begin{minipage}[b]{\linewidth}\raggedright
SOLUTION
\end{minipage} & \begin{minipage}[b]{\linewidth}\raggedright
SUMMARY
\end{minipage} & \begin{minipage}[b]{\linewidth}\raggedright
SCHEDULE
\end{minipage} \\
\midrule\noalign{}
\endhead
\bottomrule\noalign{}
\endlastfoot
\end{longtable}

\subsubsection{Description}\label{description-810}

The RSVD keyword defines the dissolved gas-oil ratio (Rs) versus depth
tables for each equilibration region that should be used when there is
dissolved gas in the model (DISGAS has been activated in the RUNSPEC
section) and the EQLOPT1 variable has been set to a positive integer on
the EQUIL keyword in the SOLUTION section.

\begin{longtable}[]{@{}
  >{\raggedright\arraybackslash}p{(\linewidth - 6\tabcolsep) * \real{0.2500}}
  >{\raggedright\arraybackslash}p{(\linewidth - 6\tabcolsep) * \real{0.2500}}
  >{\raggedright\arraybackslash}p{(\linewidth - 6\tabcolsep) * \real{0.2500}}
  >{\raggedright\arraybackslash}p{(\linewidth - 6\tabcolsep) * \real{0.2500}}@{}}
\toprule\noalign{}
\begin{minipage}[b]{\linewidth}\raggedright
No.
\end{minipage} & \begin{minipage}[b]{\linewidth}\raggedright
Name
\end{minipage} & \begin{minipage}[b]{\linewidth}\raggedright
Description
\end{minipage} & \begin{minipage}[b]{\linewidth}\raggedright
Default
\end{minipage} \\
\midrule\noalign{}
\endhead
\bottomrule\noalign{}
\endlastfoot
Field & Metric & Laboratory & \\
1 & DEPTH & A columnar vector of real monotonically increasing down the
column values that defines the depth values for the corresponding
dissolve gas-oil ratio values, RS. & None \\
feet & m & cm & \\
2 & RS & A columnar vector of real values that defines the dissolved
gas-oil ratio values at the corresponding DEPTH. & None \\
Mscf/stb & sm3/sm3 & scc/scc & \\
Notes: & & & \\
\end{longtable}

\emph{Table 10.30: RSVD Keyword Description}

Alternatively, the oil bubble-point pressure versus depth tables may be
entered using the PBVD keyword in the SOLUTION section instead of this
keyword.

See also the PBVD and EQUIL keywords in the SOLUTION section.

\subsubsection{Example}\label{example-457}

Given NTEQUL equals three and NDRXVD is greater than or equal to two on
the EQLDIMS keyword in the RUNSPEC section, then the following example
defines the bubble-point versus depth functions.

\begin{verbatim}
--
--       DEPTH    RS
--                MSCF/STB
--       ------   --------
RSVD
         3000.0    1.400
         8000.0    1.400                             / RS VS DEPTH EQUIL REGN 01
--       ------   --------
         3000.0    1.400
         8000.0    1.400                             / RS VS DEPTH EQUIL REGN 02
--       ------   --------
         3000.0    1.400
         8000.0    1.400                             / RS VS DEPTH EQUIL REGN 03
\end{verbatim}

Here three tables are entered with a constant GOR versus depth
relationship.

\subsection{RSW -- Define the Initial Equilibrium Solution Gas in Water
Ratio for All Grid
Blocks}\label{rsw-define-the-initial-equilibrium-solution-gas-in-water-ratio-for-all-grid-blocks}

\begin{longtable}[]{@{}
  >{\raggedright\arraybackslash}p{(\linewidth - 14\tabcolsep) * \real{0.1250}}
  >{\raggedright\arraybackslash}p{(\linewidth - 14\tabcolsep) * \real{0.1250}}
  >{\raggedright\arraybackslash}p{(\linewidth - 14\tabcolsep) * \real{0.1250}}
  >{\raggedright\arraybackslash}p{(\linewidth - 14\tabcolsep) * \real{0.1250}}
  >{\raggedright\arraybackslash}p{(\linewidth - 14\tabcolsep) * \real{0.1250}}
  >{\raggedright\arraybackslash}p{(\linewidth - 14\tabcolsep) * \real{0.1250}}
  >{\raggedright\arraybackslash}p{(\linewidth - 14\tabcolsep) * \real{0.1250}}
  >{\raggedright\arraybackslash}p{(\linewidth - 14\tabcolsep) * \real{0.1250}}@{}}
\toprule\noalign{}
\begin{minipage}[b]{\linewidth}\raggedright
RUNSPEC
\end{minipage} & \begin{minipage}[b]{\linewidth}\raggedright
GRID
\end{minipage} & \begin{minipage}[b]{\linewidth}\raggedright
EDIT
\end{minipage} & \begin{minipage}[b]{\linewidth}\raggedright
PROPS
\end{minipage} & \begin{minipage}[b]{\linewidth}\raggedright
REGIONS
\end{minipage} & \begin{minipage}[b]{\linewidth}\raggedright
SOLUTION
\end{minipage} & \begin{minipage}[b]{\linewidth}\raggedright
SUMMARY
\end{minipage} & \begin{minipage}[b]{\linewidth}\raggedright
SCHEDULE
\end{minipage} \\
\midrule\noalign{}
\endhead
\bottomrule\noalign{}
\endlastfoot
\end{longtable}

\subsubsection{Description}\label{description-811}

The \hyperref[REF_HEADING_KEYWORD_RSW_10_3]{RSW} keyword defines the
initial equilibration solution gas in water ratio values for all grid
cells in the model and should be used in conjunction with the PBUB,
PDEW, PRESSURE, RS, SGAS, SOIL and SWAT keywords etc., to fully describe
the initial state of the model. The keyword should only be used if both
gas and water phases have been activated in the model via the GAS and
WATER keywords, and the DISGASW keyword is also present activating OPM
Flow's Dissolved Gas in Water Model. All the aforementioned keywords are
in the RUNSPEC section.

This is the non-standard method to initialize the model via enumeration
and is seldom employed in the industry. The standard methodology is for
OPM Flow to initialize a model using the parameters on the EQUIL keyword
combined with other keywords to fully describe the initial state of the
model. The keyword can be used with all grid types.

\begin{longtable}[]{@{}
  >{\raggedright\arraybackslash}p{(\linewidth - 6\tabcolsep) * \real{0.2500}}
  >{\raggedright\arraybackslash}p{(\linewidth - 6\tabcolsep) * \real{0.2500}}
  >{\raggedright\arraybackslash}p{(\linewidth - 6\tabcolsep) * \real{0.2500}}
  >{\raggedright\arraybackslash}p{(\linewidth - 6\tabcolsep) * \real{0.2500}}@{}}
\toprule\noalign{}
\begin{minipage}[b]{\linewidth}\raggedright
No.
\end{minipage} & \begin{minipage}[b]{\linewidth}\raggedright
Name
\end{minipage} & \begin{minipage}[b]{\linewidth}\raggedright
Description
\end{minipage} & \begin{minipage}[b]{\linewidth}\raggedright
Default
\end{minipage} \\
\midrule\noalign{}
\endhead
\bottomrule\noalign{}
\endlastfoot
Field & Metric & Laboratory & \\
1 & RSW & RSW is an array of positive real numbers assigning the initial
equilibration solution gas-water ratio values to each cell in the model.
Repeat counts may be used, for example 20*1.30. & None \\
Mscf/stb & sm3/sm3 & scc/scc & \\
Notes: & & & \\
\end{longtable}

\emph{Table 10.34: RSW Keyword Description}

See also the PBUB, PDEW, PRESSURE, RS, SGAS, SOIL and SWAT keywords to
fully define the initial state of the model.

\begin{tcolorbox}[enhanced jigsaw, title=\textcolor{quarto-callout-note-color}{\faInfo}\hspace{0.5em}{Note}, coltitle=black, breakable, toprule=.15mm, colbacktitle=quarto-callout-note-color!10!white, leftrule=.75mm, opacitybacktitle=0.6, colback=white, bottomtitle=1mm, toptitle=1mm, titlerule=0mm, left=2mm, colframe=quarto-callout-note-color-frame, arc=.35mm, rightrule=.15mm, bottomrule=.15mm, opacityback=0]

This is an OPM Flow specific keyword for the simulator's Dissolved Gas
in Water Model that is activated by declaring that dissolved gas in
water is present in the run using the DISGASW keyword in the RUNSPEC
section. Use the command line option --enable-opm-rst-file=true to
output the \hyperref[REF_HEADING_KEYWORD_RSW_10_3]{RSW} data to the
RESTART file.

\end{tcolorbox}

\subsubsection{Example}\label{example-458}

\begin{verbatim}
--
--       INITIAL EQUILIBRATION SOLUTION GAS IN WATER RATIO VALUES FOR ALL CELLS
--       BASED ON NX = 100, NY = 100 AND NZ = 3
--
RSW
         1000*0.0000   1000* 0.0000   1000*1.3000                               /
\end{verbatim}

The above example defines the initial equilibration solution gas-water
values to be 0.000 for all the cells in the first and second layers and
1.3000 for all the cells in the third layer.

\subsection{RTEMP -- Define the Initial Reservoir Temperature for the
Model}\label{rtemp-define-the-initial-reservoir-temperature-for-the-model-1}

This keyword defines the initial reservoir temperature for the model.
Note that the RTEMP keyword is an alias for RTEMPA, and that both
keywords are supported by OPM Flow, in both the PROPS and SOLUTION
sections, but are treated as being mutually exclusive.

See RTEMP - Define the Initial Reservoir Temperature for the Model in
the PROPS section for a full description.

\subsection{RTEMPA -- Define the Initial Reservoir Temperature for the
Model}\label{rtempa-define-the-initial-reservoir-temperature-for-the-model-1}

This keyword defines the initial reservoir temperature for the model.
Note that the RTEMPA keyword is an alias for RTEMP, and that both
keywords are supported by OPM Flow, in both the PROPS and SOLUTION
sections, but are treated as being mutually exclusive.

See RTEMPA - Define the Initial Reservoir Temperature for the Model in
the PROPS section for a full description.

\subsection{RTEMPVD -- Define the Initial Reservoir Temperature versus
Depth
Tables}\label{rtempvd-define-the-initial-reservoir-temperature-versus-depth-tables-1}

\begin{longtable}[]{@{}
  >{\raggedright\arraybackslash}p{(\linewidth - 14\tabcolsep) * \real{0.1250}}
  >{\raggedright\arraybackslash}p{(\linewidth - 14\tabcolsep) * \real{0.1250}}
  >{\raggedright\arraybackslash}p{(\linewidth - 14\tabcolsep) * \real{0.1250}}
  >{\raggedright\arraybackslash}p{(\linewidth - 14\tabcolsep) * \real{0.1250}}
  >{\raggedright\arraybackslash}p{(\linewidth - 14\tabcolsep) * \real{0.1250}}
  >{\raggedright\arraybackslash}p{(\linewidth - 14\tabcolsep) * \real{0.1250}}
  >{\raggedright\arraybackslash}p{(\linewidth - 14\tabcolsep) * \real{0.1250}}
  >{\raggedright\arraybackslash}p{(\linewidth - 14\tabcolsep) * \real{0.1250}}@{}}
\toprule\noalign{}
\begin{minipage}[b]{\linewidth}\raggedright
RUNSPEC
\end{minipage} & \begin{minipage}[b]{\linewidth}\raggedright
GRID
\end{minipage} & \begin{minipage}[b]{\linewidth}\raggedright
EDIT
\end{minipage} & \begin{minipage}[b]{\linewidth}\raggedright
PROPS
\end{minipage} & \begin{minipage}[b]{\linewidth}\raggedright
REGIONS
\end{minipage} & \begin{minipage}[b]{\linewidth}\raggedright
SOLUTION
\end{minipage} & \begin{minipage}[b]{\linewidth}\raggedright
SUMMARY
\end{minipage} & \begin{minipage}[b]{\linewidth}\raggedright
SCHEDULE
\end{minipage} \\
\midrule\noalign{}
\endhead
\bottomrule\noalign{}
\endlastfoot
\end{longtable}

\subsubsection{Description}\label{description-812}

This keyword defines the initial reservoir temperature versus depth
tables for each equilibration region. Note that the RTEMPVD keyword is
an alias for TEMPVD, and that both keywords are supported by OPM Flow,
in both the PROPS and SOLUTION sections, but are treated as being
mutually exclusive.

The initial reservoir temperature must be defined when OPM Flow's
thermal option has been activated by the THERMAL keyword in the RUNSPEC
section. Note this is different to the commercial simulator that uses
the TEMP keyword in the RUNSPEC section to activate the black-oil
temperature model, and the THERMAL keyword to activate the compositional
thermal model.

The initial reservoir temperature should be defined when OPM Flow's CO2
or H2 storage option has been activated by the CO2STORE or
\hyperref[REF_HEADING_KEYWORD_H2STORE]{H2STORE} keyword in the RUNSPEC
section.

\begin{longtable}[]{@{}
  >{\raggedright\arraybackslash}p{(\linewidth - 6\tabcolsep) * \real{0.2500}}
  >{\raggedright\arraybackslash}p{(\linewidth - 6\tabcolsep) * \real{0.2500}}
  >{\raggedright\arraybackslash}p{(\linewidth - 6\tabcolsep) * \real{0.2500}}
  >{\raggedright\arraybackslash}p{(\linewidth - 6\tabcolsep) * \real{0.2500}}@{}}
\toprule\noalign{}
\begin{minipage}[b]{\linewidth}\raggedright
No.
\end{minipage} & \begin{minipage}[b]{\linewidth}\raggedright
Name
\end{minipage} & \begin{minipage}[b]{\linewidth}\raggedright
Description
\end{minipage} & \begin{minipage}[b]{\linewidth}\raggedright
Default
\end{minipage} \\
\midrule\noalign{}
\endhead
\bottomrule\noalign{}
\endlastfoot
Field & Metric & Laboratory & \\
1 & DEPTH & A columnar vector of real monotonically increasing down the
column values that defines the depth for corresponding reservoir
temperature parameter TEMP. & None \\
feet & m & cm & \\
2 & TEMP & A columnar vector of real monotonically increasing down the
column values that defines the corresponding reservoir temperature for
the given depth. & None \\
oF & oC & oC & \\
Notes: & & & \\
\end{longtable}

\emph{Table 10.31: RTEMPVD Keyword Description}

See also the RTEMP keyword in the PROPS section for an alternative way
to define a uniform initial reservoir temperature.

\begin{tcolorbox}[enhanced jigsaw, title=\textcolor{quarto-callout-note-color}{\faInfo}\hspace{0.5em}{Note}, coltitle=black, breakable, toprule=.15mm, colbacktitle=quarto-callout-note-color!10!white, leftrule=.75mm, opacitybacktitle=0.6, colback=white, bottomtitle=1mm, toptitle=1mm, titlerule=0mm, left=2mm, colframe=quarto-callout-note-color-frame, arc=.35mm, rightrule=.15mm, bottomrule=.15mm, opacityback=0]

The keyword is documented here in the SOLUTION section, the same as the
commercial simulator, but it can also be used in the PROPS section by
OPM Flow.

\end{tcolorbox}

\subsubsection{Example}\label{example-459}

\begin{verbatim}
--
--       INITIAL RESERVOIR TEMPERATURE VERSUS DEPTH TABLE
--
RTEMPVD
--       DEPTH    TEMPERATURE
--       FEET     DEG F
--       ------   ----------
         1000.0     90.000
         2000.0    100.000
         3000.0    130.000
         4000.0    160.000                                 / TABLE N0. 01
--       ------   ----------
         1000.0     90.000
         2000.0    100.000
         3000.0    130.000
         4000.0    160.000                                 / TABLE N0. 02
--       ------   ----------
         1000.0     90.000
         2000.0    100.000
         3000.0    130.000
         4000.0    160.000                                 / TABLE N0. 03
\end{verbatim}

The above example defines three identical reservoir depth versus
temperature tables for the three NTEQUIL regions defined on the EQLDIMS
keyword in the RUNSPEC section.

\subsection{RV -- Define the Initial Equilibration CGR (Rv) for All Grid
Blocks}\label{rv-define-the-initial-equilibration-cgr-rv-for-all-grid-blocks}

\begin{longtable}[]{@{}
  >{\raggedright\arraybackslash}p{(\linewidth - 14\tabcolsep) * \real{0.1250}}
  >{\raggedright\arraybackslash}p{(\linewidth - 14\tabcolsep) * \real{0.1250}}
  >{\raggedright\arraybackslash}p{(\linewidth - 14\tabcolsep) * \real{0.1250}}
  >{\raggedright\arraybackslash}p{(\linewidth - 14\tabcolsep) * \real{0.1250}}
  >{\raggedright\arraybackslash}p{(\linewidth - 14\tabcolsep) * \real{0.1250}}
  >{\raggedright\arraybackslash}p{(\linewidth - 14\tabcolsep) * \real{0.1250}}
  >{\raggedright\arraybackslash}p{(\linewidth - 14\tabcolsep) * \real{0.1250}}
  >{\raggedright\arraybackslash}p{(\linewidth - 14\tabcolsep) * \real{0.1250}}@{}}
\toprule\noalign{}
\begin{minipage}[b]{\linewidth}\raggedright
RUNSPEC
\end{minipage} & \begin{minipage}[b]{\linewidth}\raggedright
GRID
\end{minipage} & \begin{minipage}[b]{\linewidth}\raggedright
EDIT
\end{minipage} & \begin{minipage}[b]{\linewidth}\raggedright
PROPS
\end{minipage} & \begin{minipage}[b]{\linewidth}\raggedright
REGIONS
\end{minipage} & \begin{minipage}[b]{\linewidth}\raggedright
SOLUTION
\end{minipage} & \begin{minipage}[b]{\linewidth}\raggedright
SUMMARY
\end{minipage} & \begin{minipage}[b]{\linewidth}\raggedright
SCHEDULE
\end{minipage} \\
\midrule\noalign{}
\endhead
\bottomrule\noalign{}
\endlastfoot
\end{longtable}

\subsubsection{Description}\label{description-813}

The RV keyword defines the initial equilibration vaporized oil-gas ratio
values for all grid cells in the model and should be used in conjunction
with the PBUB, PDEW, PRESSURE, RS, SGAS, SOIL and SWAT keywords etc., to
fully describe the initial state of the model. The keyword should only
be used if vaporized oil been activated in the model via the VAPOIL
keyword in the RUNSPEC section.

This is the non-standard method to initialize the model via enumeration
and is seldom employed in the industry. The standard methodology is for
OPM Flow to initialize a model using the parameters on the EQUIL keyword
combined with other keywords to fully describe the initial state of the
model. The keyword can be used with all grid types.

\begin{longtable}[]{@{}
  >{\raggedright\arraybackslash}p{(\linewidth - 6\tabcolsep) * \real{0.2500}}
  >{\raggedright\arraybackslash}p{(\linewidth - 6\tabcolsep) * \real{0.2500}}
  >{\raggedright\arraybackslash}p{(\linewidth - 6\tabcolsep) * \real{0.2500}}
  >{\raggedright\arraybackslash}p{(\linewidth - 6\tabcolsep) * \real{0.2500}}@{}}
\toprule\noalign{}
\begin{minipage}[b]{\linewidth}\raggedright
No.
\end{minipage} & \begin{minipage}[b]{\linewidth}\raggedright
Name
\end{minipage} & \begin{minipage}[b]{\linewidth}\raggedright
Description
\end{minipage} & \begin{minipage}[b]{\linewidth}\raggedright
Default
\end{minipage} \\
\midrule\noalign{}
\endhead
\bottomrule\noalign{}
\endlastfoot
Field & Metric & Laboratory & \\
1 & RV & RV is an array of real positive numbers assigning the initial
equilibration vaporized oil-gas ratio values to each cell in the model.
Repeat counts may be used, for example 20*0.00720 & None \\
stb/Mscf & sm3/sm3 & scc/scc & \\
Notes: & & & \\
\end{longtable}

\emph{Table 10.32: RV Keyword Description}

See also the PBUB, PDEW, PRESSURE, RS, SGAS, SOIL and SWAT keywords to
fully define the initial state of the model.

\subsubsection{Example}\label{example-460}

\begin{verbatim}
--
--       DEFINE INITIAL EQUILIBRATION CGR VALUES FOR ALL CELLS IN THE MODEL
--       BASED ON NX = 100, NY = 100 AND NZ = 3
--
RV
         1000*0.00720   1000*0.00725   1000*0.00730                            /
\end{verbatim}

The above example defines the initial equilibration GOR values to be
0.00720 for all the cells in the first layer, 0.00725 for all the cells
in the second layer, and finally 0.00730 for all the cells in the third
layer.

\subsection{RVVD -- Equilibration Vaporized Oil-Gas Ratio (Rv) versus
Depth
Tables}\label{rvvd-equilibration-vaporized-oil-gas-ratio-rv-versus-depth-tables}

\begin{longtable}[]{@{}
  >{\raggedright\arraybackslash}p{(\linewidth - 14\tabcolsep) * \real{0.1250}}
  >{\raggedright\arraybackslash}p{(\linewidth - 14\tabcolsep) * \real{0.1250}}
  >{\raggedright\arraybackslash}p{(\linewidth - 14\tabcolsep) * \real{0.1250}}
  >{\raggedright\arraybackslash}p{(\linewidth - 14\tabcolsep) * \real{0.1250}}
  >{\raggedright\arraybackslash}p{(\linewidth - 14\tabcolsep) * \real{0.1250}}
  >{\raggedright\arraybackslash}p{(\linewidth - 14\tabcolsep) * \real{0.1250}}
  >{\raggedright\arraybackslash}p{(\linewidth - 14\tabcolsep) * \real{0.1250}}
  >{\raggedright\arraybackslash}p{(\linewidth - 14\tabcolsep) * \real{0.1250}}@{}}
\toprule\noalign{}
\begin{minipage}[b]{\linewidth}\raggedright
RUNSPEC
\end{minipage} & \begin{minipage}[b]{\linewidth}\raggedright
GRID
\end{minipage} & \begin{minipage}[b]{\linewidth}\raggedright
EDIT
\end{minipage} & \begin{minipage}[b]{\linewidth}\raggedright
PROPS
\end{minipage} & \begin{minipage}[b]{\linewidth}\raggedright
REGIONS
\end{minipage} & \begin{minipage}[b]{\linewidth}\raggedright
SOLUTION
\end{minipage} & \begin{minipage}[b]{\linewidth}\raggedright
SUMMARY
\end{minipage} & \begin{minipage}[b]{\linewidth}\raggedright
SCHEDULE
\end{minipage} \\
\midrule\noalign{}
\endhead
\bottomrule\noalign{}
\endlastfoot
\end{longtable}

\subsubsection{Description}\label{description-814}

The RVVD keyword defines the vaporized oil-gas ratio (Rv) versus depth
tables for each equilibration region that should be used when there is
vaporized oil in the model (VAPOIL has been activated in the RUNSPEC
section) and the EQLOPT2 variable has been set to a positive integer on
the EQUIL keyword in the SOLUTION section.

\begin{longtable}[]{@{}
  >{\raggedright\arraybackslash}p{(\linewidth - 6\tabcolsep) * \real{0.2500}}
  >{\raggedright\arraybackslash}p{(\linewidth - 6\tabcolsep) * \real{0.2500}}
  >{\raggedright\arraybackslash}p{(\linewidth - 6\tabcolsep) * \real{0.2500}}
  >{\raggedright\arraybackslash}p{(\linewidth - 6\tabcolsep) * \real{0.2500}}@{}}
\toprule\noalign{}
\begin{minipage}[b]{\linewidth}\raggedright
No.
\end{minipage} & \begin{minipage}[b]{\linewidth}\raggedright
Name
\end{minipage} & \begin{minipage}[b]{\linewidth}\raggedright
Description
\end{minipage} & \begin{minipage}[b]{\linewidth}\raggedright
Default
\end{minipage} \\
\midrule\noalign{}
\endhead
\bottomrule\noalign{}
\endlastfoot
Field & Metric & Laboratory & \\
1 & DEPTH & A columnar vector of real monotonically increasing down the
column values that defines the depth values for the corresponding
vaporized oil-gas ratio values, RV. & None \\
feet & m & cm & \\
2 & RV & A columnar vector of real values that defines the vaporized
oil-gas ratio values at the corresponding DEPTH. & None \\
stb/Mscf & sm3/sm3 & scc/scc & \\
Notes: & & & \\
\end{longtable}

\emph{Table 10.33: RVVD Keyword Description}

Alternatively, the gas dew-point pressure versus depth tables may be
entered using the PDVD keyword in the SOLUTION section instead of this
keyword.

See also the PDVD and EQUIL keywords in the SOLUTION section.

\subsubsection{Example}\label{example-461}

Given NTEQUL equals three and NDRXVD is greater than or equal to two on
the EQLDIMS keyword in the RUNSPEC section, then the following example
defines the dew-point versus depth functions.

\begin{verbatim}
--
--       DEPTH    RV
--                STB/MSCF
--       ------   --------
RVVD
         3000.0   0.00725
         8000.0   0.00725                            / RV VS DEPTH EQUIL REGN 01
--       ------   --------
         3000.0   0.00730
         8000.0   0.00730                            / RV VS DEPTH EQUIL REGN 02
--       ------   --------
         3000.0   0.00750
         8000.0   0.00750                            / RV VS DEPTH EQUIL REGN 03
\end{verbatim}

Here three tables are entered with a constant CGR versus depth
relationship for each equilibration region.

\subsection{RVW -- Define the Initial Equilibration Vaporized Water in
Gas Ratio for All Grid
Blocks}\label{rvw-define-the-initial-equilibration-vaporized-water-in-gas-ratio-for-all-grid-blocks}

\begin{longtable}[]{@{}
  >{\raggedright\arraybackslash}p{(\linewidth - 14\tabcolsep) * \real{0.1250}}
  >{\raggedright\arraybackslash}p{(\linewidth - 14\tabcolsep) * \real{0.1250}}
  >{\raggedright\arraybackslash}p{(\linewidth - 14\tabcolsep) * \real{0.1250}}
  >{\raggedright\arraybackslash}p{(\linewidth - 14\tabcolsep) * \real{0.1250}}
  >{\raggedright\arraybackslash}p{(\linewidth - 14\tabcolsep) * \real{0.1250}}
  >{\raggedright\arraybackslash}p{(\linewidth - 14\tabcolsep) * \real{0.1250}}
  >{\raggedright\arraybackslash}p{(\linewidth - 14\tabcolsep) * \real{0.1250}}
  >{\raggedright\arraybackslash}p{(\linewidth - 14\tabcolsep) * \real{0.1250}}@{}}
\toprule\noalign{}
\begin{minipage}[b]{\linewidth}\raggedright
RUNSPEC
\end{minipage} & \begin{minipage}[b]{\linewidth}\raggedright
GRID
\end{minipage} & \begin{minipage}[b]{\linewidth}\raggedright
EDIT
\end{minipage} & \begin{minipage}[b]{\linewidth}\raggedright
PROPS
\end{minipage} & \begin{minipage}[b]{\linewidth}\raggedright
REGIONS
\end{minipage} & \begin{minipage}[b]{\linewidth}\raggedright
SOLUTION
\end{minipage} & \begin{minipage}[b]{\linewidth}\raggedright
SUMMARY
\end{minipage} & \begin{minipage}[b]{\linewidth}\raggedright
SCHEDULE
\end{minipage} \\
\midrule\noalign{}
\endhead
\bottomrule\noalign{}
\endlastfoot
\end{longtable}

\subsubsection{Description}\label{description-815}

The RVW keyword defines the initial equilibration vaporized water in gas
ratio values for all grid cells in the model and should be used in
conjunction with the PBUB, PDEW, PRESSURE, RV, SGAS, SOIL and SWAT
keywords etc., to fully describe the initial state of the model. The
keyword should only be used if both gas and water phases have been
activated in the model via the GAS and WATER keywords, and the VAPWAT is
also present activating OPM Flow's Vaporized Water Model. All the
aforementioned keywords are in the RUNSPEC section.

This is the non-standard method to initialize the model via enumeration
and is seldom employed in the industry. The standard methodology is for
OPM Flow to initialize a model using the parameters on the EQUIL keyword
combined with other keywords to fully describe the initial state of the
model. The keyword can be used with all grid types.

\begin{longtable}[]{@{}
  >{\raggedright\arraybackslash}p{(\linewidth - 6\tabcolsep) * \real{0.2500}}
  >{\raggedright\arraybackslash}p{(\linewidth - 6\tabcolsep) * \real{0.2500}}
  >{\raggedright\arraybackslash}p{(\linewidth - 6\tabcolsep) * \real{0.2500}}
  >{\raggedright\arraybackslash}p{(\linewidth - 6\tabcolsep) * \real{0.2500}}@{}}
\toprule\noalign{}
\begin{minipage}[b]{\linewidth}\raggedright
No.
\end{minipage} & \begin{minipage}[b]{\linewidth}\raggedright
Name
\end{minipage} & \begin{minipage}[b]{\linewidth}\raggedright
Description
\end{minipage} & \begin{minipage}[b]{\linewidth}\raggedright
Default
\end{minipage} \\
\midrule\noalign{}
\endhead
\bottomrule\noalign{}
\endlastfoot
Field & Metric & Laboratory & \\
1 & RVW & RVW is an array of real positive numbers assigning the initial
equilibration gas-vaporized water ratio values to each cell in the
model. Repeat counts may be used, for example 20*1.30. & None \\
stb/Mscf & sm3/sm3 & rcc/scc & \\
Notes: & & & \\
\end{longtable}

\emph{Table 10.34: RVW Keyword Description}

See also the PBUB, PDEW, PRESSURE, RV, SGAS, SOIL and SWAT keywords to
fully define the initial state of the model.

\begin{tcolorbox}[enhanced jigsaw, title=\textcolor{quarto-callout-note-color}{\faInfo}\hspace{0.5em}{Note}, coltitle=black, breakable, toprule=.15mm, colbacktitle=quarto-callout-note-color!10!white, leftrule=.75mm, opacitybacktitle=0.6, colback=white, bottomtitle=1mm, toptitle=1mm, titlerule=0mm, left=2mm, colframe=quarto-callout-note-color-frame, arc=.35mm, rightrule=.15mm, bottomrule=.15mm, opacityback=0]

This is an OPM Flow specific keyword for the simulator's Water
Vaporization Model that is activated by declaring that vaporized water
is present in the run using the VAPWAT keyword in the RUNSPEC section.
Use the command line option --enable-opm-rst-file=true to output the RVW
data to the RESTART file.

\end{tcolorbox}

\subsubsection{Example}\label{example-462}

\begin{verbatim}
--
--       INITIAL EQUILIBRATION WATER VAPOR IN GAS RATIO VALUES FOR ALL CELLS
--       BASED ON NX = 100, NY = 100 AND NZ = 3
--
RVW
         1000*0.0000   1000* 0.0000   1000*1.3000                               /
\end{verbatim}

The above example defines the initial equilibration gas-vaporized water
values to be 0.000 for all the cells in the first and second layers and
1.3000 for all the cells in the third layer.

\subsection{RVWVD -- Equilibration Vaporized Water-Gas Ratio (Rvw)
versus Depth
Tables}\label{rvwvd-equilibration-vaporized-water-gas-ratio-rvw-versus-depth-tables}

\begin{longtable}[]{@{}
  >{\raggedright\arraybackslash}p{(\linewidth - 14\tabcolsep) * \real{0.1250}}
  >{\raggedright\arraybackslash}p{(\linewidth - 14\tabcolsep) * \real{0.1250}}
  >{\raggedright\arraybackslash}p{(\linewidth - 14\tabcolsep) * \real{0.1250}}
  >{\raggedright\arraybackslash}p{(\linewidth - 14\tabcolsep) * \real{0.1250}}
  >{\raggedright\arraybackslash}p{(\linewidth - 14\tabcolsep) * \real{0.1250}}
  >{\raggedright\arraybackslash}p{(\linewidth - 14\tabcolsep) * \real{0.1250}}
  >{\raggedright\arraybackslash}p{(\linewidth - 14\tabcolsep) * \real{0.1250}}
  >{\raggedright\arraybackslash}p{(\linewidth - 14\tabcolsep) * \real{0.1250}}@{}}
\toprule\noalign{}
\begin{minipage}[b]{\linewidth}\raggedright
RUNSPEC
\end{minipage} & \begin{minipage}[b]{\linewidth}\raggedright
GRID
\end{minipage} & \begin{minipage}[b]{\linewidth}\raggedright
EDIT
\end{minipage} & \begin{minipage}[b]{\linewidth}\raggedright
PROPS
\end{minipage} & \begin{minipage}[b]{\linewidth}\raggedright
REGIONS
\end{minipage} & \begin{minipage}[b]{\linewidth}\raggedright
SOLUTION
\end{minipage} & \begin{minipage}[b]{\linewidth}\raggedright
SUMMARY
\end{minipage} & \begin{minipage}[b]{\linewidth}\raggedright
SCHEDULE
\end{minipage} \\
\midrule\noalign{}
\endhead
\bottomrule\noalign{}
\endlastfoot
\end{longtable}

\subsubsection{Description}\label{description-816}

The RVWVD keyword defines the vaporized water-gas ratio (Rvw) versus
depth tables for each equilibration region that should be used when
there is vaporized water in the model and the EQLOPT6 variable has been
set to a positive integer on the EQUIL keyword in the SOLUTION section.

The keyword should only be used if both gas and water phases haves been
activated in the model via the GAS and WATER keywords, and the VAPWAT is
also present activating OPM Flow's Vaporized Water Model. All the
aforementioned keywords are in the RUNSPEC section.

\begin{longtable}[]{@{}
  >{\raggedright\arraybackslash}p{(\linewidth - 6\tabcolsep) * \real{0.2500}}
  >{\raggedright\arraybackslash}p{(\linewidth - 6\tabcolsep) * \real{0.2500}}
  >{\raggedright\arraybackslash}p{(\linewidth - 6\tabcolsep) * \real{0.2500}}
  >{\raggedright\arraybackslash}p{(\linewidth - 6\tabcolsep) * \real{0.2500}}@{}}
\toprule\noalign{}
\begin{minipage}[b]{\linewidth}\raggedright
No.
\end{minipage} & \begin{minipage}[b]{\linewidth}\raggedright
Name
\end{minipage} & \begin{minipage}[b]{\linewidth}\raggedright
Description
\end{minipage} & \begin{minipage}[b]{\linewidth}\raggedright
Default
\end{minipage} \\
\midrule\noalign{}
\endhead
\bottomrule\noalign{}
\endlastfoot
Field & Metric & Laboratory & \\
1 & DEPTH & A columnar vector of real monotonically increasing down the
column values that defines the depth values for the corresponding
vaporized oil-gas ratio values, RVW & None \\
feet & m & cm & \\
2 & RVW & A columnar vector of real values that defines the vaporized
water-gas ratio values, values at the corresponding DEPTH. & None \\
stb/Mscf & sm3/sm3 & scc/scc & \\
Notes: & & & \\
\end{longtable}

\emph{Table 10.35: RVWVD Keyword Description}

Alternatively, the vaporized water-gas ratio for each cell may be set
via the RVW keyword in the SOLUTION section, if the non-standard method
to initialize the model via enumeration is being employed.

See also the EQUIL keywords in the SOLUTION section.

\begin{tcolorbox}[enhanced jigsaw, title=\textcolor{quarto-callout-note-color}{\faInfo}\hspace{0.5em}{Note}, coltitle=black, breakable, toprule=.15mm, colbacktitle=quarto-callout-note-color!10!white, leftrule=.75mm, opacitybacktitle=0.6, colback=white, bottomtitle=1mm, toptitle=1mm, titlerule=0mm, left=2mm, colframe=quarto-callout-note-color-frame, arc=.35mm, rightrule=.15mm, bottomrule=.15mm, opacityback=0]

This is an OPM Flow specific keyword for the simulator's Water
Vaporization Model that is activated by declaring that vaporized water
is present in the run using the VAPWAT keyword in the RUNSPEC section.
Use the command line option --enable-opm-rst-file=true to output the RVW
data to the RESTART file.

\end{tcolorbox}

\subsubsection{Example}\label{example-463}

Given NTEQUL equals three and NDRXVD is greater than or equal to two on
the EQLDIMS keyword in the RUNSPEC section, then the following example
defines the vaporized water-gas ratio versus depth functions.

\begin{verbatim}
--
--       DEPTH    RVW
--                STB/MSCF
--       ------   --------
RVWVD
         3000.0   0.00000
         8000.0   0.00000                            / RVW VS DEPTH EQUIL REGN 01
--       ------   --------
         3000.0   0.00000
         8000.0   0.00000                            / RVW VS DEPTH EQUIL REGN 02
--       ------   --------
         3000.0   0.00100
         8000.0   0.00100                            / RVW VS DEPTH EQUIL REGN 03
\end{verbatim}

The example shows three tables for three regions with constant RVW
versus depth relationships for each equilibration region, with the first
two tables having a zero vaporized water-gas ratio and the last region
having a constant 0.001 stb/Mscf versus depth relationship.

\subsection{SALT -- Define the Initial Equilibration Salt Concentration
for All Grid
Blocks}\label{salt-define-the-initial-equilibration-salt-concentration-for-all-grid-blocks}

\begin{longtable}[]{@{}
  >{\raggedright\arraybackslash}p{(\linewidth - 14\tabcolsep) * \real{0.1250}}
  >{\raggedright\arraybackslash}p{(\linewidth - 14\tabcolsep) * \real{0.1250}}
  >{\raggedright\arraybackslash}p{(\linewidth - 14\tabcolsep) * \real{0.1250}}
  >{\raggedright\arraybackslash}p{(\linewidth - 14\tabcolsep) * \real{0.1250}}
  >{\raggedright\arraybackslash}p{(\linewidth - 14\tabcolsep) * \real{0.1250}}
  >{\raggedright\arraybackslash}p{(\linewidth - 14\tabcolsep) * \real{0.1250}}
  >{\raggedright\arraybackslash}p{(\linewidth - 14\tabcolsep) * \real{0.1250}}
  >{\raggedright\arraybackslash}p{(\linewidth - 14\tabcolsep) * \real{0.1250}}@{}}
\toprule\noalign{}
\begin{minipage}[b]{\linewidth}\raggedright
RUNSPEC
\end{minipage} & \begin{minipage}[b]{\linewidth}\raggedright
GRID
\end{minipage} & \begin{minipage}[b]{\linewidth}\raggedright
EDIT
\end{minipage} & \begin{minipage}[b]{\linewidth}\raggedright
PROPS
\end{minipage} & \begin{minipage}[b]{\linewidth}\raggedright
REGIONS
\end{minipage} & \begin{minipage}[b]{\linewidth}\raggedright
SOLUTION
\end{minipage} & \begin{minipage}[b]{\linewidth}\raggedright
SUMMARY
\end{minipage} & \begin{minipage}[b]{\linewidth}\raggedright
SCHEDULE
\end{minipage} \\
\midrule\noalign{}
\endhead
\bottomrule\noalign{}
\endlastfoot
\end{longtable}

\subsubsection{Description}\label{description-817}

The SALT keyword defines the initial equilibration salt concentration
values for all grid cells in the model and should be used in conjunction
with the PBUB, PDEW, PRESSURE, RS, RV, SGAS, SGAS and SWAT keywords
etc., to fully describe the initial state of the model. The keyword
should only be used if the salt (brine) phase has been activated in the
model via the BRINE keyword in the RUNSPEC section.

This is the non-standard method to initialize the model via enumeration
and is seldom employed in the industry. The standard methodology is for
OPM Flow to initialize a model using the parameters on the EQUIL keyword
combined with other keywords to fully describe the initial state of the
model. The keyword can be used with all grid types.

\begin{longtable}[]{@{}
  >{\raggedright\arraybackslash}p{(\linewidth - 6\tabcolsep) * \real{0.2500}}
  >{\raggedright\arraybackslash}p{(\linewidth - 6\tabcolsep) * \real{0.2500}}
  >{\raggedright\arraybackslash}p{(\linewidth - 6\tabcolsep) * \real{0.2500}}
  >{\raggedright\arraybackslash}p{(\linewidth - 6\tabcolsep) * \real{0.2500}}@{}}
\toprule\noalign{}
\begin{minipage}[b]{\linewidth}\raggedright
No.
\end{minipage} & \begin{minipage}[b]{\linewidth}\raggedright
Name
\end{minipage} & \begin{minipage}[b]{\linewidth}\raggedright
Description
\end{minipage} & \begin{minipage}[b]{\linewidth}\raggedright
Default
\end{minipage} \\
\midrule\noalign{}
\endhead
\bottomrule\noalign{}
\endlastfoot
Field & Metric & Laboratory & \\
1 & SALT & SALT is an array of real positive numbers that are greater
than or equal to zero assigning the initial equilibration salt
concentration values to each cell in the model. Repeat counts may be
used, for example 20*15.0. & None \\
lb/stb & kg/sm3 & gm/scc & \\
Notes: & & & \\
\end{longtable}

\emph{Table 10.36: SALT Keyword Description}

See also the PBUB, PDEW, PRESSURE, RS, RV, SGAS, SOIL and SWAT keywords
to fully define the initial state of the model.

\subsubsection{Example}\label{example-464}

\begin{verbatim}
--
--       DEFINE INITIAL EQUILIBRATION SALT CONCENTRATION FOR ALL CELLS
--       BASED ON NX = 100, NY = 100 AND NZ = 3
--
SALT
         10000*0.0000    10000*0.0000    10000*15.000                             /
\end{verbatim}

The above example defines the initial equilibration salt concentration
values to be 0.0000 for all the cells in the first and second layers and
finally 15.000 for all the cells in the third layer.

\subsection{SALTP -- Define the Initial Precipitated Salt Volume
Fraction for All Grid
Blocks}\label{saltp-define-the-initial-precipitated-salt-volume-fraction-for-all-grid-blocks}

\begin{longtable}[]{@{}
  >{\raggedright\arraybackslash}p{(\linewidth - 14\tabcolsep) * \real{0.1250}}
  >{\raggedright\arraybackslash}p{(\linewidth - 14\tabcolsep) * \real{0.1250}}
  >{\raggedright\arraybackslash}p{(\linewidth - 14\tabcolsep) * \real{0.1250}}
  >{\raggedright\arraybackslash}p{(\linewidth - 14\tabcolsep) * \real{0.1250}}
  >{\raggedright\arraybackslash}p{(\linewidth - 14\tabcolsep) * \real{0.1250}}
  >{\raggedright\arraybackslash}p{(\linewidth - 14\tabcolsep) * \real{0.1250}}
  >{\raggedright\arraybackslash}p{(\linewidth - 14\tabcolsep) * \real{0.1250}}
  >{\raggedright\arraybackslash}p{(\linewidth - 14\tabcolsep) * \real{0.1250}}@{}}
\toprule\noalign{}
\begin{minipage}[b]{\linewidth}\raggedright
RUNSPEC
\end{minipage} & \begin{minipage}[b]{\linewidth}\raggedright
GRID
\end{minipage} & \begin{minipage}[b]{\linewidth}\raggedright
EDIT
\end{minipage} & \begin{minipage}[b]{\linewidth}\raggedright
PROPS
\end{minipage} & \begin{minipage}[b]{\linewidth}\raggedright
REGIONS
\end{minipage} & \begin{minipage}[b]{\linewidth}\raggedright
SOLUTION
\end{minipage} & \begin{minipage}[b]{\linewidth}\raggedright
SUMMARY
\end{minipage} & \begin{minipage}[b]{\linewidth}\raggedright
SCHEDULE
\end{minipage} \\
\midrule\noalign{}
\endhead
\bottomrule\noalign{}
\endlastfoot
\end{longtable}

\subsubsection{Description}\label{description-818}

The SALTP keyword defines the initial equilibration precipitated salt
volume fraction values for all grid cells in the model and should be
used in conjunction with the PBUB, PDEW, PRESSURE, RS, RV, SGAS, SOIL
and SWAT keywords etc., to fully describe the initial state of the
model. The keyword should only be used if the salt (brine) phase has
been activated in the model via the BRINE keyword, and the PRECSALT
keyword to activate OPM Flow's Salt Precipitation Model. Both keywords
are in the RUNSPEC section.

This is the non-standard method to initialize the model via enumeration
and is seldom employed in the industry. The standard methodology is for
OPM Flow to initialize a model using the parameters on the EQUIL keyword
combined with other keywords to fully describe the initial state of the
model. The keyword can be used with all grid types.

\begin{longtable}[]{@{}
  >{\raggedright\arraybackslash}p{(\linewidth - 6\tabcolsep) * \real{0.2500}}
  >{\raggedright\arraybackslash}p{(\linewidth - 6\tabcolsep) * \real{0.2500}}
  >{\raggedright\arraybackslash}p{(\linewidth - 6\tabcolsep) * \real{0.2500}}
  >{\raggedright\arraybackslash}p{(\linewidth - 6\tabcolsep) * \real{0.2500}}@{}}
\toprule\noalign{}
\begin{minipage}[b]{\linewidth}\raggedright
No.
\end{minipage} & \begin{minipage}[b]{\linewidth}\raggedright
Name
\end{minipage} & \begin{minipage}[b]{\linewidth}\raggedright
Description
\end{minipage} & \begin{minipage}[b]{\linewidth}\raggedright
Default
\end{minipage} \\
\midrule\noalign{}
\endhead
\bottomrule\noalign{}
\endlastfoot
Field & Metric & Laboratory & \\
1 & SALTP & SALTP is an array of real positive numbers that are greater
than or equal to zero and less than or equal to one, that define the
initial equilibration salt volume fraction values to each cell in the
model. Repeat counts may be used, for example 20*0.15. & None \\
dimensionless & dimensionless & dimensionless & \\
Notes: & & & \\
\end{longtable}

\emph{Table 10.37: SALTP Keyword Description}

See also the PBUB, PDEW, PRESSURE, RS, RV, SGAS, SOIL and SWAT keywords
to fully define the initial state of the model.

\begin{tcolorbox}[enhanced jigsaw, title=\textcolor{quarto-callout-note-color}{\faInfo}\hspace{0.5em}{Note}, coltitle=black, breakable, toprule=.15mm, colbacktitle=quarto-callout-note-color!10!white, leftrule=.75mm, opacitybacktitle=0.6, colback=white, bottomtitle=1mm, toptitle=1mm, titlerule=0mm, left=2mm, colframe=quarto-callout-note-color-frame, arc=.35mm, rightrule=.15mm, bottomrule=.15mm, opacityback=0]

This is an OPM Flow specific keyword for the simulator's Salt
Precipitation Model that is activated by the PRECSALT keyword and
declaring that vaporized water is present in the run via the VAPWAT in
the RUNSPEC section. This keyword defines the initial precipitated salt
volume fraction contained within the pore space. See SALT in the
SOLUTION section that defines the initial salt concentration within the
water phase.

\end{tcolorbox}

\subsubsection{Example}\label{example-465}

The example activates the standard Brine Tracking model using the BRINE
keyword, OPM Flow's Salt Precipitation model using the PRECSALT keyword,
and OPM Flow's vaporized water phase with the VAPWAT keyword; all three
keywords are in the RUNSPEC section.

\begin{verbatim}
-- ==============================================================================
--
-- RUNSPEC SECTION
--
-- ==============================================================================
RUNSPEC
--
--       MAX     MAX     RSVD    TVDP    TVDP
--       EQLNUM  DEPTH   NODES   TABLE   NODES
EQLDIMS
         3       1*      20      1*      1*                                    /
--
--       ACTIVATE STANDARD BRINE MODEL
--
BRINE
--
--       ACTIVATE THE OPM FLOW SALT PRECIPITATION MODEL (OPM FLOW KEYWORD)
--
PRECSALT
--
--       VAPORIZED WATER IN WET GAS IS PRESENT IN THE RUN (OPM FLOW KEYWORD)
--
VAPWAT
\end{verbatim}

Then in the SOLUTION section the SALTP keyword would be of the form:

\begin{verbatim}
-- ==============================================================================
--
-- SOLUTION SECTION
--
-- ==============================================================================
SOLUTION
--
--       DEFINE INITIAL PRECIPITATED SALT VOLUME FRACTION FOR ALL CELLS
--       BASED ON NX = 100, NY = 100 AND NZ = 3
--
SALTP
         1000*0.0000    1000*0.0000    1000*0.100                             /
\end{verbatim}

Here the initial equilibration precipitated salt volume fraction values
are set to 0.0000 for all the cells in the first and second layers and
finally 0.1000 for all the cells in the third layer.

\subsection{SALTPVD -- Initial Precipitated Salt Volume Fraction versus
Depth
Tables}\label{saltpvd-initial-precipitated-salt-volume-fraction-versus-depth-tables}

\begin{longtable}[]{@{}
  >{\raggedright\arraybackslash}p{(\linewidth - 14\tabcolsep) * \real{0.1250}}
  >{\raggedright\arraybackslash}p{(\linewidth - 14\tabcolsep) * \real{0.1250}}
  >{\raggedright\arraybackslash}p{(\linewidth - 14\tabcolsep) * \real{0.1250}}
  >{\raggedright\arraybackslash}p{(\linewidth - 14\tabcolsep) * \real{0.1250}}
  >{\raggedright\arraybackslash}p{(\linewidth - 14\tabcolsep) * \real{0.1250}}
  >{\raggedright\arraybackslash}p{(\linewidth - 14\tabcolsep) * \real{0.1250}}
  >{\raggedright\arraybackslash}p{(\linewidth - 14\tabcolsep) * \real{0.1250}}
  >{\raggedright\arraybackslash}p{(\linewidth - 14\tabcolsep) * \real{0.1250}}@{}}
\toprule\noalign{}
\begin{minipage}[b]{\linewidth}\raggedright
RUNSPEC
\end{minipage} & \begin{minipage}[b]{\linewidth}\raggedright
GRID
\end{minipage} & \begin{minipage}[b]{\linewidth}\raggedright
EDIT
\end{minipage} & \begin{minipage}[b]{\linewidth}\raggedright
PROPS
\end{minipage} & \begin{minipage}[b]{\linewidth}\raggedright
REGIONS
\end{minipage} & \begin{minipage}[b]{\linewidth}\raggedright
SOLUTION
\end{minipage} & \begin{minipage}[b]{\linewidth}\raggedright
SUMMARY
\end{minipage} & \begin{minipage}[b]{\linewidth}\raggedright
SCHEDULE
\end{minipage} \\
\midrule\noalign{}
\endhead
\bottomrule\noalign{}
\endlastfoot
\end{longtable}

\subsubsection{Description}\label{description-819}

The SALTPVD keyword defines the initial precipitated salt volume
fraction versus depth tables for each equilibration region for when OPM
Flow's Salt Precipitation Model has been activated in the input deck via
the PRECSALT keyword in the RUNSPEC section. The keyword defines the
initial deposited salt as a volume fraction (Ss), that is solid salt
saturation.

\begin{tcolorbox}[enhanced jigsaw, title=\textcolor{quarto-callout-note-color}{\faInfo}\hspace{0.5em}{Note}, coltitle=black, breakable, toprule=.15mm, colbacktitle=quarto-callout-note-color!10!white, leftrule=.75mm, opacitybacktitle=0.6, colback=white, bottomtitle=1mm, toptitle=1mm, titlerule=0mm, left=2mm, colframe=quarto-callout-note-color-frame, arc=.35mm, rightrule=.15mm, bottomrule=.15mm, opacityback=0]

This is an OPM Flow specific keyword for the simulator's Salt
Precipitation Model that is activated by the PRECSALT keyword and
declaring that vaporized water is present in the run via the VAPWAT in
the RUNSPEC section. This is the initial precipitated salt volume
fraction contained within the pore space, see SALTVD in the SOLUTION
section that defines the initial salt concentration within the water
phase.

\end{tcolorbox}

\begin{longtable}[]{@{}
  >{\raggedright\arraybackslash}p{(\linewidth - 6\tabcolsep) * \real{0.2500}}
  >{\raggedright\arraybackslash}p{(\linewidth - 6\tabcolsep) * \real{0.2500}}
  >{\raggedright\arraybackslash}p{(\linewidth - 6\tabcolsep) * \real{0.2500}}
  >{\raggedright\arraybackslash}p{(\linewidth - 6\tabcolsep) * \real{0.2500}}@{}}
\toprule\noalign{}
\begin{minipage}[b]{\linewidth}\raggedright
No.
\end{minipage} & \begin{minipage}[b]{\linewidth}\raggedright
Name
\end{minipage} & \begin{minipage}[b]{\linewidth}\raggedright
Description
\end{minipage} & \begin{minipage}[b]{\linewidth}\raggedright
Default
\end{minipage} \\
\midrule\noalign{}
\endhead
\bottomrule\noalign{}
\endlastfoot
Field & Metric & Laboratory & \\
1 & DEPTH & A columnar vector of real monotonically increasing down the
column values that defines the depth for corresponding salt volume
fraction SALTPSAT. & None \\
feet & m & cm & \\
2 & SALTPSAT & A columnar vector of real values that defines the
corresponding volume fraction of precipitated salt for the given depth.
Note only the standard Brine Model is supported and therefore there
should be only one columnar vector of SALTPSAT. & None \\
dimensionless & dimensionless & dimensionless & \\
Notes: & & & \\
\end{longtable}

\emph{Table 10.38: SALTPVD Keyword Description}

Precipitated salt can be viewed in OPM ResInsight via the SALTP dynamic
grid variable that represents the precipitated salt saturation.

\subsubsection{Example}\label{example-466}

The example activates the standard Brine Tracking model using the BRINE
keyword, OPM Flow's Salt Precipitation model using the PRECSALT keyword,
and OPM Flow's vaporized water phase with the VAPWAT keyword; all three
keywords are in the RUNSPEC section. The example also sets the number of
equilibrium regions to three (NTEQUL set to three on the EQLDIMS keyword
also in the RUNSPEC), that is:

\begin{verbatim}
-- ==============================================================================
--
-- RUNSPEC SECTION
--
-- ==============================================================================
RUNSPEC
--
--       MAX     MAX     RSVD    TVDP    TVDP
--       EQLNUM  DEPTH   NODES   TABLE   NODES
EQLDIMS
         3       1*      20      1*      1*                                    /
--
--       ACTIVATE STANDARD BRINE MODEL
--
BRINE
--
--       ACTIVATE THE OPM FLOW SALT PRECIPITATION MODEL (OPM FLOW KEYWORD)
--
PRECSALT
--
--       VAPORIZED WATER IN WET GAS IS PRESENT IN THE RUN (OPM FLOW KEYWORD)
--
VAPWAT
\end{verbatim}

Then in the SOLUTION section the SALTPVD keyword would be of the form:

\begin{verbatim}
-- ==============================================================================
--
-- SOLUTION SECTION
--
-- ==============================================================================
SOLUTION
--
--       PRECIPITATED SALT VOLUME FRACTION VERSUS DEPTH (OPM FLOW KEYWORD)
--
--       DEPTH    SALTPSAT
--       ------   --------
SALTPVD
         3000.0    0.000
         8000.0    0.000                                   / EQUIL REGN 01
--       ------   --------
         3000.0    0.000
         8000.0    0.000                                   / EQUIL REGN 02
--       ------   --------
         3000.0    0.000
         8000.0    0.000                                   / EQUIL REGN 03
\end{verbatim}

Here the initial precipitated salt volume fraction has been set to zero
for all three equilibration regions.

\begin{verbatim}
\end{verbatim}

\subsection{SALTREST -- Define the Restart Salt Concentration for All
Grid
Blocks}\label{saltrest-define-the-restart-salt-concentration-for-all-grid-blocks}

\begin{longtable}[]{@{}
  >{\raggedright\arraybackslash}p{(\linewidth - 14\tabcolsep) * \real{0.1250}}
  >{\raggedright\arraybackslash}p{(\linewidth - 14\tabcolsep) * \real{0.1250}}
  >{\raggedright\arraybackslash}p{(\linewidth - 14\tabcolsep) * \real{0.1250}}
  >{\raggedright\arraybackslash}p{(\linewidth - 14\tabcolsep) * \real{0.1250}}
  >{\raggedright\arraybackslash}p{(\linewidth - 14\tabcolsep) * \real{0.1250}}
  >{\raggedright\arraybackslash}p{(\linewidth - 14\tabcolsep) * \real{0.1250}}
  >{\raggedright\arraybackslash}p{(\linewidth - 14\tabcolsep) * \real{0.1250}}
  >{\raggedright\arraybackslash}p{(\linewidth - 14\tabcolsep) * \real{0.1250}}@{}}
\toprule\noalign{}
\begin{minipage}[b]{\linewidth}\raggedright
RUNSPEC
\end{minipage} & \begin{minipage}[b]{\linewidth}\raggedright
GRID
\end{minipage} & \begin{minipage}[b]{\linewidth}\raggedright
EDIT
\end{minipage} & \begin{minipage}[b]{\linewidth}\raggedright
PROPS
\end{minipage} & \begin{minipage}[b]{\linewidth}\raggedright
REGIONS
\end{minipage} & \begin{minipage}[b]{\linewidth}\raggedright
SOLUTION
\end{minipage} & \begin{minipage}[b]{\linewidth}\raggedright
SUMMARY
\end{minipage} & \begin{minipage}[b]{\linewidth}\raggedright
SCHEDULE
\end{minipage} \\
\midrule\noalign{}
\endhead
\bottomrule\noalign{}
\endlastfoot
\end{longtable}

\subsubsection{Description}\label{description-820}

The SALTREST keyword defines restart salt concentration values for all
grid cells in the model and should be used in runs that are using the
RESTART facility, where the initial run has not used the Low Salt or
Brine options. This allows for initial runs that have used the standard
water PVT properties via the PVTW keyword in the PROPS section, to be
restarted with salt dependent water properties. The keyword should only
be used if the salt (brine) phase has been activated in the current
restart run (not the initial run) via the BRINE keyword in the RUNSPEC
section.

This keyword is not supported by OPM Flow but would change the results
if supported so the simulation will be stopped.

\begin{longtable}[]{@{}
  >{\raggedright\arraybackslash}p{(\linewidth - 6\tabcolsep) * \real{0.2500}}
  >{\raggedright\arraybackslash}p{(\linewidth - 6\tabcolsep) * \real{0.2500}}
  >{\raggedright\arraybackslash}p{(\linewidth - 6\tabcolsep) * \real{0.2500}}
  >{\raggedright\arraybackslash}p{(\linewidth - 6\tabcolsep) * \real{0.2500}}@{}}
\toprule\noalign{}
\begin{minipage}[b]{\linewidth}\raggedright
No.
\end{minipage} & \begin{minipage}[b]{\linewidth}\raggedright
Name
\end{minipage} & \begin{minipage}[b]{\linewidth}\raggedright
Description
\end{minipage} & \begin{minipage}[b]{\linewidth}\raggedright
Default
\end{minipage} \\
\midrule\noalign{}
\endhead
\bottomrule\noalign{}
\endlastfoot
Field & Metric & Laboratory & \\
1 & SALTREST & SALTREST is an array of real positive numbers that are
greater than or equal to zero assigning the restart salt concentration
values to each cell in the model. Repeat counts may be used, for example
20*15.0. & None \\
lb/stb & kg/sm3 & gm/scc & \\
Notes: & & & \\
\end{longtable}

\emph{Table 10.39: SALTREST Keyword Description}

See also the PVTWSALT keyword in the PROPS section and the RESTART
keyword in the SOLUTION section.

\subsubsection{Example}\label{example-467}

\begin{verbatim}
--
--       DEFINE RESTART SALTREST VALUES FOR ALL CELLS IN THE MODEL
--       BASED ON NX = 100, NY = 100 AND NZ = 3
--
SALTREST
         1000*0.0000    1000*0.0000    1000*15.000                             /
\end{verbatim}

The above example defines the restart salt concentration values to be
0.0000 for all the cells in the first and second layers and finally
15.000 for all the cells in the third layer.

\subsection{SALTVD -- Equilibration Salt Concentration versus Depth
Tables}\label{saltvd-equilibration-salt-concentration-versus-depth-tables}

\begin{longtable}[]{@{}
  >{\raggedright\arraybackslash}p{(\linewidth - 14\tabcolsep) * \real{0.1250}}
  >{\raggedright\arraybackslash}p{(\linewidth - 14\tabcolsep) * \real{0.1250}}
  >{\raggedright\arraybackslash}p{(\linewidth - 14\tabcolsep) * \real{0.1250}}
  >{\raggedright\arraybackslash}p{(\linewidth - 14\tabcolsep) * \real{0.1250}}
  >{\raggedright\arraybackslash}p{(\linewidth - 14\tabcolsep) * \real{0.1250}}
  >{\raggedright\arraybackslash}p{(\linewidth - 14\tabcolsep) * \real{0.1250}}
  >{\raggedright\arraybackslash}p{(\linewidth - 14\tabcolsep) * \real{0.1250}}
  >{\raggedright\arraybackslash}p{(\linewidth - 14\tabcolsep) * \real{0.1250}}@{}}
\toprule\noalign{}
\begin{minipage}[b]{\linewidth}\raggedright
RUNSPEC
\end{minipage} & \begin{minipage}[b]{\linewidth}\raggedright
GRID
\end{minipage} & \begin{minipage}[b]{\linewidth}\raggedright
EDIT
\end{minipage} & \begin{minipage}[b]{\linewidth}\raggedright
PROPS
\end{minipage} & \begin{minipage}[b]{\linewidth}\raggedright
REGIONS
\end{minipage} & \begin{minipage}[b]{\linewidth}\raggedright
SOLUTION
\end{minipage} & \begin{minipage}[b]{\linewidth}\raggedright
SUMMARY
\end{minipage} & \begin{minipage}[b]{\linewidth}\raggedright
SCHEDULE
\end{minipage} \\
\midrule\noalign{}
\endhead
\bottomrule\noalign{}
\endlastfoot
\end{longtable}

\subsubsection{Description}\label{description-821}

The SALTVD keyword defines the initial salt concentration versus depth
tables for each equilibration region for when the salt (brine) phase has
been activated in the model via the BRINE keyword in the RUNSPEC
section, and the EQLOPT1 variable has been set to a positive integer on
the EQUIL keyword in the SOLUTION section. Secondly, the keyword should
also be used to set the initial salt concentration versus depth if OPM
Flow's PRECSALT keyword in the RUNSPEC section has been used to activate
the simulators Salt Precipitation model.

\begin{longtable}[]{@{}
  >{\raggedright\arraybackslash}p{(\linewidth - 6\tabcolsep) * \real{0.2500}}
  >{\raggedright\arraybackslash}p{(\linewidth - 6\tabcolsep) * \real{0.2500}}
  >{\raggedright\arraybackslash}p{(\linewidth - 6\tabcolsep) * \real{0.2500}}
  >{\raggedright\arraybackslash}p{(\linewidth - 6\tabcolsep) * \real{0.2500}}@{}}
\toprule\noalign{}
\begin{minipage}[b]{\linewidth}\raggedright
No.
\end{minipage} & \begin{minipage}[b]{\linewidth}\raggedright
Name
\end{minipage} & \begin{minipage}[b]{\linewidth}\raggedright
Description
\end{minipage} & \begin{minipage}[b]{\linewidth}\raggedright
Default
\end{minipage} \\
\midrule\noalign{}
\endhead
\bottomrule\noalign{}
\endlastfoot
Field & Metric & Laboratory & \\
1 & DEPTH & A columnar vector of real monotonically increasing down the
column values that defines the depth for corresponding salt
concentrations SALTCON. & None \\
feet & m & cm & \\
2 & SALTCON & A columnar vector of real monotonically increasing down
the column values that defines the corresponding salt concentration
within the water phase for the given depth. There should be one columnar
vector for each type of salt. For the standard Brine Model there is only
one salt type and therefore there should be only one columnar vector of
SALTCON. However, if the BRINE keyword has been invoked with the ECLMC
keyword in the RUNSPEC section, then there should one columnar SALTCON
vector for each declared salt type. It is recommended to provide initial
salt concentrations less then or equal to values provided by SALTSOL
keyword in the PROPS section. & None \\
lb/stb & kg/sm3 & gm/scc & \\
Notes: & & & \\
\end{longtable}

\emph{Table 10.40: SALTVD Keyword Description}

\begin{tcolorbox}[enhanced jigsaw, title=\textcolor{quarto-callout-note-color}{\faInfo}\hspace{0.5em}{Note}, coltitle=black, breakable, toprule=.15mm, colbacktitle=quarto-callout-note-color!10!white, leftrule=.75mm, opacitybacktitle=0.6, colback=white, bottomtitle=1mm, toptitle=1mm, titlerule=0mm, left=2mm, colframe=quarto-callout-note-color-frame, arc=.35mm, rightrule=.15mm, bottomrule=.15mm, opacityback=0]

This is the initial salt concentration contained within the water phase,
see the SALTPVD keyword in the SOLUTION section that defines the initial
salt volume fraction that has been precipitated into the pore space.

\end{tcolorbox}

\subsubsection{Examples}\label{examples-59}

The first example activates the standard Brine Tracking model using the
BRINE keyword in the RUNSPEC section and sets the number of equilibrium
regions to three (NTEQUIL set to 3 on the EQLDIMS keyword also in the
RUNSPEC), that is:

\begin{verbatim}
-- ==============================================================================
--
-- RUNSPEC SECTION
--
-- ==============================================================================
RUNSPEC
--
--       MAX     MAX     RSVD    TVDP    TVDP
--       EQLNUM  DEPTH   NODES   TABLE   NODES
EQLDIMS
         3       1*      20      1*      1*                                    /
--
--       ACTIVATE STANDARD BRINE MODEL
--
BRINE
\end{verbatim}

Then in the SOLUTION section the SALTVD keyword would be of the form:

\begin{verbatim}
-- ==============================================================================
--
-- SOLUTION SECTION
--
-- ==============================================================================
SOLUTION
--
--       DEPTH    SALT-1    SALT-2    SALT-3    SALT-4
--                SALTCON   SALTCON   SALTCON   SALTCON
--       ------   -------   -------   -------   -------                                                         SALTVD
         3000.0    1.200
         8000.0    1.200                                / EQUIL REGN 01
--       ------   --------
         3000.0    1.300
         8000.0    1.300                                / EQUIL REGN 02
--       ------   --------
         3000.0    1.400
         8000.0    1.400                                / EQUIL REGN 03

\end{verbatim}

The next example shows how the SALTVD keyword is entered when both the
ECLMC and BRINE keywords have activated the Multi-Component Brine model
in the RUNSPEC section, that is:

\begin{verbatim}
-- ==============================================================================
--
-- RUNSPEC SECTION
--
-- ==============================================================================
RUNSPEC
--
--       MAX     MAX     RSVD    TVDP    TVDP
--       EQLNUM  DEPTH   NODES   TABLE   NODES
EQLDIMS
         3       1*      20      1*      1*                                    /
--
--
--       ACTIVATE MULTI-COMPONENT BRINE MODEL
--
ECLMC
--
--       DEFINE WATER PHASE MULTI-COMPONENT BRINE COMPONENTS
--
--       SALT1   SALT2   SALT3   SALT4   SALT5
BRINE
         NACL    CACL    MGC03                                                 /
\end{verbatim}

The above example activates the Multi-Component Brine model with three
different water salinities for three equilibrium regions. In this case
the resulting SALTVD keyword would be of the form:

\begin{verbatim}
-- ==============================================================================
--
-- SOLUTION SECTION
--
-- ==============================================================================
SOLUTION
--
--       DEPTH    SALT-01   SALT-02   SALT-03   SALT-04
--                CONCENTR  CONCENTR  CONCENTR  CONCENTR
--       ------   --------  --------  --------  --------                                                         SALTVD
         3000.0    1.200    0.540     0.020
         7000.0    1.200    0.640     0.040             / EQUIL REGN 01
--       ------   --------  --------  --------
         3000.0    1.300    0.440     0.020
         8000.0    1.300    0.540     0.040             / EQUIL REGN 02
--       ------   --------  --------  --------
         5000.0    1.400    0.640     0.002
         8000.0    1.400    0.640     0.002             / EQUIL REGN 03
\end{verbatim}

In this case there are three data sets, on one for each equilibrium
region and three SALTCON columnar vectors, one for each salt type (NACL,
CACL and MGC03) declared via the BRINE keyword in the RUNSPEC section.

Note that the Multi-Component Brine model is not available in OPM Flow.

\begin{verbatim}
\end{verbatim}

\subsection{SBIOF -- Define The Initial Equilibration Biofilm Volume
Fraction For All Grid
Blocks}\label{sbiof-define-the-initial-equilibration-biofilm-volume-fraction-for-all-grid-blocks}

\begin{longtable}[]{@{}
  >{\raggedright\arraybackslash}p{(\linewidth - 14\tabcolsep) * \real{0.1250}}
  >{\raggedright\arraybackslash}p{(\linewidth - 14\tabcolsep) * \real{0.1250}}
  >{\raggedright\arraybackslash}p{(\linewidth - 14\tabcolsep) * \real{0.1250}}
  >{\raggedright\arraybackslash}p{(\linewidth - 14\tabcolsep) * \real{0.1250}}
  >{\raggedright\arraybackslash}p{(\linewidth - 14\tabcolsep) * \real{0.1250}}
  >{\raggedright\arraybackslash}p{(\linewidth - 14\tabcolsep) * \real{0.1250}}
  >{\raggedright\arraybackslash}p{(\linewidth - 14\tabcolsep) * \real{0.1250}}
  >{\raggedright\arraybackslash}p{(\linewidth - 14\tabcolsep) * \real{0.1250}}@{}}
\toprule\noalign{}
\begin{minipage}[b]{\linewidth}\raggedright
RUNSPEC
\end{minipage} & \begin{minipage}[b]{\linewidth}\raggedright
GRID
\end{minipage} & \begin{minipage}[b]{\linewidth}\raggedright
EDIT
\end{minipage} & \begin{minipage}[b]{\linewidth}\raggedright
PROPS
\end{minipage} & \begin{minipage}[b]{\linewidth}\raggedright
REGIONS
\end{minipage} & \begin{minipage}[b]{\linewidth}\raggedright
SOLUTION
\end{minipage} & \begin{minipage}[b]{\linewidth}\raggedright
SUMMARY
\end{minipage} & \begin{minipage}[b]{\linewidth}\raggedright
SCHEDULE
\end{minipage} \\
\midrule\noalign{}
\endhead
\bottomrule\noalign{}
\endlastfoot
\end{longtable}

\subsubsection{Description}\label{description-822}

The SBIOF keyword defines the initial equilibration biofilm volume
fraction for all grid cells in the model. The keyword should only be
used if either the \hyperref[REF_HEADING_KEYWORD_BIOFILM]{BIOFILM} or
MICP model has been activated in the RUNSPEC section.

\begin{tcolorbox}[enhanced jigsaw, title=\textcolor{quarto-callout-note-color}{\faInfo}\hspace{0.5em}{Note}, coltitle=black, breakable, toprule=.15mm, colbacktitle=quarto-callout-note-color!10!white, leftrule=.75mm, opacitybacktitle=0.6, colback=white, bottomtitle=1mm, toptitle=1mm, titlerule=0mm, left=2mm, colframe=quarto-callout-note-color-frame, arc=.35mm, rightrule=.15mm, bottomrule=.15mm, opacityback=0]

This is an OPM Flow specific keyword.

\end{tcolorbox}

\begin{longtable}[]{@{}
  >{\raggedright\arraybackslash}p{(\linewidth - 6\tabcolsep) * \real{0.2500}}
  >{\raggedright\arraybackslash}p{(\linewidth - 6\tabcolsep) * \real{0.2500}}
  >{\raggedright\arraybackslash}p{(\linewidth - 6\tabcolsep) * \real{0.2500}}
  >{\raggedright\arraybackslash}p{(\linewidth - 6\tabcolsep) * \real{0.2500}}@{}}
\toprule\noalign{}
\begin{minipage}[b]{\linewidth}\raggedright
No.
\end{minipage} & \begin{minipage}[b]{\linewidth}\raggedright
Name
\end{minipage} & \begin{minipage}[b]{\linewidth}\raggedright
Description
\end{minipage} & \begin{minipage}[b]{\linewidth}\raggedright
Default
\end{minipage} \\
\midrule\noalign{}
\endhead
\bottomrule\noalign{}
\endlastfoot
Field & Metric & Laboratory & \\
1 & SBIOF & SBIOF is an array of real numbers that are greater than or
equal to zero and less than or equal to one assigning the initial
equilibration biofilm volume fraction values to each cell in the model.
Repeat counts may be used, for example 20*0.0010. & None \\
dimensionless & dimensionless & dimensionless & \\
Notes: & & & \\
\end{longtable}

\emph{Table 10.41: SBIOF Keyword Description}

For both \hyperref[REF_HEADING_KEYWORD_BIOFILM]{BIOFILM} and MICP
models, see also the SMICR keyword, and for the MICP model, the SCALC,
SOXYG, and SUREA keywords to define the initial state of the model.

\subsubsection{Example}\label{example-468}

\begin{verbatim}
--
--       DEFINE INITIAL EQUILIBRATION BIOFILM VOLUME FRACTION FOR ALL CELLS
--       BASED ON NX = 100, NY = 100 AND NZ = 3
--
SBIOF
         10000*0.0000 10000*0.0000 10000*0.0010                                   /
\end{verbatim}

The above example defines the initial equilibration biofilm volume
fraction values to be 0.0000 for all the cells in the first and second
layers and finally 0.0010 for all the cells in the third layer.

\subsection{SCALC -- Define The Initial Equilibration Calcite Volume
Fraction For All Grid
Blocks}\label{scalc-define-the-initial-equilibration-calcite-volume-fraction-for-all-grid-blocks}

\begin{longtable}[]{@{}
  >{\raggedright\arraybackslash}p{(\linewidth - 14\tabcolsep) * \real{0.1250}}
  >{\raggedright\arraybackslash}p{(\linewidth - 14\tabcolsep) * \real{0.1250}}
  >{\raggedright\arraybackslash}p{(\linewidth - 14\tabcolsep) * \real{0.1250}}
  >{\raggedright\arraybackslash}p{(\linewidth - 14\tabcolsep) * \real{0.1250}}
  >{\raggedright\arraybackslash}p{(\linewidth - 14\tabcolsep) * \real{0.1250}}
  >{\raggedright\arraybackslash}p{(\linewidth - 14\tabcolsep) * \real{0.1250}}
  >{\raggedright\arraybackslash}p{(\linewidth - 14\tabcolsep) * \real{0.1250}}
  >{\raggedright\arraybackslash}p{(\linewidth - 14\tabcolsep) * \real{0.1250}}@{}}
\toprule\noalign{}
\begin{minipage}[b]{\linewidth}\raggedright
RUNSPEC
\end{minipage} & \begin{minipage}[b]{\linewidth}\raggedright
GRID
\end{minipage} & \begin{minipage}[b]{\linewidth}\raggedright
EDIT
\end{minipage} & \begin{minipage}[b]{\linewidth}\raggedright
PROPS
\end{minipage} & \begin{minipage}[b]{\linewidth}\raggedright
REGIONS
\end{minipage} & \begin{minipage}[b]{\linewidth}\raggedright
SOLUTION
\end{minipage} & \begin{minipage}[b]{\linewidth}\raggedright
SUMMARY
\end{minipage} & \begin{minipage}[b]{\linewidth}\raggedright
SCHEDULE
\end{minipage} \\
\midrule\noalign{}
\endhead
\bottomrule\noalign{}
\endlastfoot
\end{longtable}

\subsubsection{Description}\label{description-823}

The SCALC keyword defines the initial equilibration calcite volume
fraction for all grid cells in the model. The keyword should only be
used if the MICP model has been activated in the RUNSPEC section.

\begin{tcolorbox}[enhanced jigsaw, title=\textcolor{quarto-callout-note-color}{\faInfo}\hspace{0.5em}{Note}, coltitle=black, breakable, toprule=.15mm, colbacktitle=quarto-callout-note-color!10!white, leftrule=.75mm, opacitybacktitle=0.6, colback=white, bottomtitle=1mm, toptitle=1mm, titlerule=0mm, left=2mm, colframe=quarto-callout-note-color-frame, arc=.35mm, rightrule=.15mm, bottomrule=.15mm, opacityback=0]

This is an OPM Flow specific keyword.

\end{tcolorbox}

\begin{longtable}[]{@{}
  >{\raggedright\arraybackslash}p{(\linewidth - 6\tabcolsep) * \real{0.2500}}
  >{\raggedright\arraybackslash}p{(\linewidth - 6\tabcolsep) * \real{0.2500}}
  >{\raggedright\arraybackslash}p{(\linewidth - 6\tabcolsep) * \real{0.2500}}
  >{\raggedright\arraybackslash}p{(\linewidth - 6\tabcolsep) * \real{0.2500}}@{}}
\toprule\noalign{}
\begin{minipage}[b]{\linewidth}\raggedright
No.
\end{minipage} & \begin{minipage}[b]{\linewidth}\raggedright
Name
\end{minipage} & \begin{minipage}[b]{\linewidth}\raggedright
Description
\end{minipage} & \begin{minipage}[b]{\linewidth}\raggedright
Default
\end{minipage} \\
\midrule\noalign{}
\endhead
\bottomrule\noalign{}
\endlastfoot
Field & Metric & Laboratory & \\
1 & SCALC & SCALC is an array of real numbers that are greater than or
equal to zero and less than or equal to one assigning the initial
equilibration calcite volume fraction values to each cell in the model.
Repeat counts may be used, for example 20*0.0010. & None \\
dimensionless & dimensionless & dimensionless & \\
Notes: & & & \\
\end{longtable}

\emph{Table 10.42: SCALC Keyword Description}

See also the SBIOF, SMICR, SOXYG and SUREA keywords to define the
initial state of the model.

\subsubsection{Example}\label{example-469}

\begin{verbatim}
--
--       DEFINE INITIAL EQUILIBRATION CALCITE VOLUME FRACTION FOR ALL CELLS
--       BASED ON NX = 100, NY = 100 AND NZ = 3
--
SCALC
         1000*0.0000 1000*0.0000 1000*0.0010                                   /
\end{verbatim}

The above example defines the initial equilibration calcite volume
fraction values to be 0.0000 for all the cells in the first and second
layers and finally 0.0010 for all the cells in the third layer.

\subsection{SCVD -- Define Equilibration Coal Solvent Concentration
versus Depth
Tables}\label{scvd-define-equilibration-coal-solvent-concentration-versus-depth-tables}

\begin{longtable}[]{@{}
  >{\raggedright\arraybackslash}p{(\linewidth - 14\tabcolsep) * \real{0.1250}}
  >{\raggedright\arraybackslash}p{(\linewidth - 14\tabcolsep) * \real{0.1250}}
  >{\raggedright\arraybackslash}p{(\linewidth - 14\tabcolsep) * \real{0.1250}}
  >{\raggedright\arraybackslash}p{(\linewidth - 14\tabcolsep) * \real{0.1250}}
  >{\raggedright\arraybackslash}p{(\linewidth - 14\tabcolsep) * \real{0.1250}}
  >{\raggedright\arraybackslash}p{(\linewidth - 14\tabcolsep) * \real{0.1250}}
  >{\raggedright\arraybackslash}p{(\linewidth - 14\tabcolsep) * \real{0.1250}}
  >{\raggedright\arraybackslash}p{(\linewidth - 14\tabcolsep) * \real{0.1250}}@{}}
\toprule\noalign{}
\begin{minipage}[b]{\linewidth}\raggedright
RUNSPEC
\end{minipage} & \begin{minipage}[b]{\linewidth}\raggedright
GRID
\end{minipage} & \begin{minipage}[b]{\linewidth}\raggedright
EDIT
\end{minipage} & \begin{minipage}[b]{\linewidth}\raggedright
PROPS
\end{minipage} & \begin{minipage}[b]{\linewidth}\raggedright
REGIONS
\end{minipage} & \begin{minipage}[b]{\linewidth}\raggedright
SOLUTION
\end{minipage} & \begin{minipage}[b]{\linewidth}\raggedright
SUMMARY
\end{minipage} & \begin{minipage}[b]{\linewidth}\raggedright
SCHEDULE
\end{minipage} \\
\midrule\noalign{}
\endhead
\bottomrule\noalign{}
\endlastfoot
\end{longtable}

\subsubsection{Description}\label{description-824}

The SCVD keyword defines the initial coal solvent concentration versus
depth tables for each equilibration region for when the coal phase has
been activated in the run via the COAL keyword in the RUNSPEC section.
The keyword may be used in conjunction with the SOLVCONC keyword in the
SOLUTION section, to fully describe the initial state of the model. Note
both SOLVCONC and SCVD are optional as the simulator will calculate the
coal gas concentration based on the equilibrium concentration and the
block pressure.

This keyword is not supported by OPM Flow but would change the results
if supported so the simulation will be stopped.

\begin{longtable}[]{@{}
  >{\raggedright\arraybackslash}p{(\linewidth - 6\tabcolsep) * \real{0.2500}}
  >{\raggedright\arraybackslash}p{(\linewidth - 6\tabcolsep) * \real{0.2500}}
  >{\raggedright\arraybackslash}p{(\linewidth - 6\tabcolsep) * \real{0.2500}}
  >{\raggedright\arraybackslash}p{(\linewidth - 6\tabcolsep) * \real{0.2500}}@{}}
\toprule\noalign{}
\begin{minipage}[b]{\linewidth}\raggedright
No.
\end{minipage} & \begin{minipage}[b]{\linewidth}\raggedright
Name
\end{minipage} & \begin{minipage}[b]{\linewidth}\raggedright
Description
\end{minipage} & \begin{minipage}[b]{\linewidth}\raggedright
Default
\end{minipage} \\
\midrule\noalign{}
\endhead
\bottomrule\noalign{}
\endlastfoot
Field & Metric & Laboratory & \\
1 & DEPTH & A columnar vector of real monotonically increasing down the
column values that defines the depth values for the corresponding coal
solvent concentration, SCVALS. & None \\
feet & m & cm & \\
2 & SCVALS & A columnar vector of real values that defines the coal
solvent concentration values at the corresponding DEPTH. & None \\
Mscf/ft3 & sm3/m3 & scc/cc & \\
Notes: & & & \\
\end{longtable}

\emph{Table 10.43: SCVD Keyword Description}

See also the SOLVCONC, GCVD, GASCONC and GASSATC keywords in the
SOLUTION section.

\subsubsection{Example}\label{example-470}

Given NTEQUL equals three and NDRXVD is greater than or equal to two on
the EQLDIMS keyword in the RUNSPEC section, then the following example
defines the coal solvent concentration versus depth functions.

\begin{verbatim}
--
--       DEPTH    SOLVC
--                MSCF/FT
--       ------   --------
SCVD
          100.0   75.5000
         1000.0   75.5000                            / SC VS DEPTH EQUIL REGN 01
--       ------   --------
          100.0   65.5000
         1000.0   65.5000                            / SC VS DEPTH EQUIL REGN 02
--       ------   --------
          100.0   60.0000
         1000.0   60.0000                            / SC VS DEPTH EQUIL REGN 03
\end{verbatim}

Here three tables are entered with a constant coal solvent concentration
versus depth relationship for each equilibration region

\subsection{SFOAM -- Define the Initial Equilibration Foam Concentration
for All Grid
Blocks}\label{sfoam-define-the-initial-equilibration-foam-concentration-for-all-grid-blocks}

\begin{longtable}[]{@{}
  >{\raggedright\arraybackslash}p{(\linewidth - 14\tabcolsep) * \real{0.1250}}
  >{\raggedright\arraybackslash}p{(\linewidth - 14\tabcolsep) * \real{0.1250}}
  >{\raggedright\arraybackslash}p{(\linewidth - 14\tabcolsep) * \real{0.1250}}
  >{\raggedright\arraybackslash}p{(\linewidth - 14\tabcolsep) * \real{0.1250}}
  >{\raggedright\arraybackslash}p{(\linewidth - 14\tabcolsep) * \real{0.1250}}
  >{\raggedright\arraybackslash}p{(\linewidth - 14\tabcolsep) * \real{0.1250}}
  >{\raggedright\arraybackslash}p{(\linewidth - 14\tabcolsep) * \real{0.1250}}
  >{\raggedright\arraybackslash}p{(\linewidth - 14\tabcolsep) * \real{0.1250}}@{}}
\toprule\noalign{}
\begin{minipage}[b]{\linewidth}\raggedright
RUNSPEC
\end{minipage} & \begin{minipage}[b]{\linewidth}\raggedright
GRID
\end{minipage} & \begin{minipage}[b]{\linewidth}\raggedright
EDIT
\end{minipage} & \begin{minipage}[b]{\linewidth}\raggedright
PROPS
\end{minipage} & \begin{minipage}[b]{\linewidth}\raggedright
REGIONS
\end{minipage} & \begin{minipage}[b]{\linewidth}\raggedright
SOLUTION
\end{minipage} & \begin{minipage}[b]{\linewidth}\raggedright
SUMMARY
\end{minipage} & \begin{minipage}[b]{\linewidth}\raggedright
SCHEDULE
\end{minipage} \\
\midrule\noalign{}
\endhead
\bottomrule\noalign{}
\endlastfoot
\end{longtable}

\subsubsection{Description}\label{description-825}

The SFOAM keyword defines the initial equilibration foam concentration
values for all grid cells in the model and should be used in conjunction
with the PBUB, PDEW, PRESSURE, RS, RV, SGAS, SGAS and SWAT keywords
etc., to fully describe the initial state of the model. The keyword
should only be used if the foam phase has been activated in the model
via the FOAM keyword in the RUNSPEC section.

This keyword is not supported by OPM Flow but would change the results
if supported so the simulation will be stopped.

This is the non-standard method to initialize the model via enumeration
and is seldom employed in the industry. The standard methodology is for
OPM Flow to initialize a model using the parameters on the EQUIL keyword
combined with other keywords to fully describe the initial state of the
model. The keyword can be used with all grid types.

\begin{longtable}[]{@{}
  >{\raggedright\arraybackslash}p{(\linewidth - 6\tabcolsep) * \real{0.2500}}
  >{\raggedright\arraybackslash}p{(\linewidth - 6\tabcolsep) * \real{0.2500}}
  >{\raggedright\arraybackslash}p{(\linewidth - 6\tabcolsep) * \real{0.2500}}
  >{\raggedright\arraybackslash}p{(\linewidth - 6\tabcolsep) * \real{0.2500}}@{}}
\toprule\noalign{}
\begin{minipage}[b]{\linewidth}\raggedright
No.
\end{minipage} & \begin{minipage}[b]{\linewidth}\raggedright
Name
\end{minipage} & \begin{minipage}[b]{\linewidth}\raggedright
Description
\end{minipage} & \begin{minipage}[b]{\linewidth}\raggedright
Default
\end{minipage} \\
\midrule\noalign{}
\endhead
\bottomrule\noalign{}
\endlastfoot
Field & Metric & Laboratory & \\
1 & SFOAM & SFOAM is an array of real positive numbers that are greater
than or equal to zero assigning the initial equilibration foam
concentration values to each cell in the model. Units are dependent on
the transport phase specified via the FOAMOPT1 variable on the FOAMOPTS
keyword in the PROPS section. FOAMOPT1 should be set to either GAS or
WATER. Repeat counts may be used, for example 20*0.5 & None \\
Gas: lb/Mscf Water: lb/stb & Gas: kg/sm3 Water: kg/sm3 & Gas: gm/scc
Water: gm/scc & \\
Notes: & & & \\
\end{longtable}

\emph{Table 10.44: SFOAM Keyword Description}

See also the PBUB, PDEW, PRESSURE, RS, RV, SGAS, SOIL and SWAT keywords
to fully define the initial state of the model.

\subsubsection{Example}\label{example-471}

\begin{verbatim}
--
--       DEFINE INITIAL EQUILIBRATION FOAM VALUES FOR ALL CELLS IN THE MODEL
--       BASED ON NX = 100, NY = 100 AND NZ = 3
--
SFOAM
         1000*0.0000    1000*0.0000    1000*0.500                             /
\end{verbatim}

The above example defines the initial equilibration foam concentration
values to be 0.0000 for all the cells in the first and second layers and
finally 0.500 for all the cells in the third layer.

\subsection{SGAS -- Define the Initial Equilibration Gas Saturation for
All Grid
Blocks}\label{sgas-define-the-initial-equilibration-gas-saturation-for-all-grid-blocks}

\begin{longtable}[]{@{}
  >{\raggedright\arraybackslash}p{(\linewidth - 14\tabcolsep) * \real{0.1250}}
  >{\raggedright\arraybackslash}p{(\linewidth - 14\tabcolsep) * \real{0.1250}}
  >{\raggedright\arraybackslash}p{(\linewidth - 14\tabcolsep) * \real{0.1250}}
  >{\raggedright\arraybackslash}p{(\linewidth - 14\tabcolsep) * \real{0.1250}}
  >{\raggedright\arraybackslash}p{(\linewidth - 14\tabcolsep) * \real{0.1250}}
  >{\raggedright\arraybackslash}p{(\linewidth - 14\tabcolsep) * \real{0.1250}}
  >{\raggedright\arraybackslash}p{(\linewidth - 14\tabcolsep) * \real{0.1250}}
  >{\raggedright\arraybackslash}p{(\linewidth - 14\tabcolsep) * \real{0.1250}}@{}}
\toprule\noalign{}
\begin{minipage}[b]{\linewidth}\raggedright
RUNSPEC
\end{minipage} & \begin{minipage}[b]{\linewidth}\raggedright
GRID
\end{minipage} & \begin{minipage}[b]{\linewidth}\raggedright
EDIT
\end{minipage} & \begin{minipage}[b]{\linewidth}\raggedright
PROPS
\end{minipage} & \begin{minipage}[b]{\linewidth}\raggedright
REGIONS
\end{minipage} & \begin{minipage}[b]{\linewidth}\raggedright
SOLUTION
\end{minipage} & \begin{minipage}[b]{\linewidth}\raggedright
SUMMARY
\end{minipage} & \begin{minipage}[b]{\linewidth}\raggedright
SCHEDULE
\end{minipage} \\
\midrule\noalign{}
\endhead
\bottomrule\noalign{}
\endlastfoot
\end{longtable}

\subsubsection{Description}\label{description-826}

The SGAS keyword defines the initial equilibration gas saturation values
for all grid cells in the model and should be used in conjunction with
the PBUB, PDEW, PRESSURE, RS, RV, SOIL and SWAT keywords etc., to fully
describe the initial state of the model. The keyword should only be used
if the gas phase has been activated in the model via the GAS keyword in
the RUNSPEC section.

This is the non-standard method to initialize the model via enumeration
and is seldom employed in the industry. The standard methodology is for
OPM Flow to initialize a model using the parameters on the EQUIL keyword
combined with other keywords to fully describe the initial state of the
model. The keyword can be used with all grid types.

\begin{longtable}[]{@{}
  >{\raggedright\arraybackslash}p{(\linewidth - 6\tabcolsep) * \real{0.2500}}
  >{\raggedright\arraybackslash}p{(\linewidth - 6\tabcolsep) * \real{0.2500}}
  >{\raggedright\arraybackslash}p{(\linewidth - 6\tabcolsep) * \real{0.2500}}
  >{\raggedright\arraybackslash}p{(\linewidth - 6\tabcolsep) * \real{0.2500}}@{}}
\toprule\noalign{}
\begin{minipage}[b]{\linewidth}\raggedright
No.
\end{minipage} & \begin{minipage}[b]{\linewidth}\raggedright
Name
\end{minipage} & \begin{minipage}[b]{\linewidth}\raggedright
Description
\end{minipage} & \begin{minipage}[b]{\linewidth}\raggedright
Default
\end{minipage} \\
\midrule\noalign{}
\endhead
\bottomrule\noalign{}
\endlastfoot
Field & Metric & Laboratory & \\
1 & SGAS & SGAS is an array of real positive numbers that are greater
than or equal to zero and less than or equal to one assigning the
initial equilibration gas saturation values to each cell in the model.
Repeat counts may be used, for example 20*0.600. & None \\
dimensionless & dimensionless & dimensionless & \\
Notes: & & & \\
\end{longtable}

\emph{Table 10.45: SGAS Keyword Description}

Note for two phase runs it is only necessary to enter one saturation
array of one of the phases present in the run (SGAS, SOIL, or SWAT), as
the simulator will calculate the other phases by difference. Similarly
for three phase runs it is only necessary to enter the array data for
two of the phases, as the third saturation will again be calculated by
the simulator.

See also the PBUB, PDEW, PRESSURE, RS, RV, SOIL and SWAT keywords to
fully define the initial state of the model.

\subsubsection{Example}\label{example-472}

\begin{verbatim}
--
--       DEFINE INITIAL EQUILIBRATION GAS SAT VALUES FOR ALL CELLS IN THE MODEL
--       BASED ON NX = 100, NY = 100 AND NZ = 3
--
SGAS
         1000*0.7000    1000*0.6500    1000*0.6000                             /
\end{verbatim}

The above example defines the initial equilibration gas saturation
values to be 0.7000 for all the cells in the first layer, 0.6500 for all
the cells in the second layer, and finally 0.6000 for all the cells in
the third layer.

\subsection{SKIP -- Activate Skipping of All Keywords and Input
Data}\label{skip-activate-skipping-of-all-keywords-and-input-data-6}

The SKIP keyword activates skipping of all keywords and input data until
the ENDSKIP keyword is encountered. All keywords between the SKIP and
ENDSKIP keywords are ignored.

See SKIP -- Activate Skipping of All Keywords and Input Data in the
GLOBAL section for a full description.

\subsection{SKIP100 -- Activate Skipping of Black-Oil Keywords and Input
Data}\label{skip100-activate-skipping-of-black-oil-keywords-and-input-data-5}

The SKIP100 keyword activates skipping of all keywords and input data by
the commercial black-oil simulator until the ENDSKIP keyword is
encountered. All keywords between the SKIP100 and ENDSKIP keywords are
ignored by the commercial black-oil simulator. The SKIP100 keyword is
ignored by the commercial compositional simulator.

See SKIP100 -- Activate Skipping of Keywords by Black-Oil Simulator in
the GLOBAL section for a full description.

\subsection{SKIP300 -- Activate Skipping of Keywords by Compositional
Simulator}\label{skip300-activate-skipping-of-keywords-by-compositional-simulator-6}

The SKIP300 keyword activates skipping of all keywords and input data by
the commercial compositional simulator until the ENDSKIP keyword is
encountered. All keywords between the SKIP300 and ENDSKIP keywords are
ignored by the commercial compositional simulator. The SKIP300 keyword
is ignored by the commercial black-oil simulator.

See Error: Reference source not found in the GLOBAL section for a full
description.

\subsection{SMICR -- Define The Initial Equilibration Microbial
Concentration For All Grid
Blocks}\label{smicr-define-the-initial-equilibration-microbial-concentration-for-all-grid-blocks}

\begin{longtable}[]{@{}
  >{\raggedright\arraybackslash}p{(\linewidth - 14\tabcolsep) * \real{0.1250}}
  >{\raggedright\arraybackslash}p{(\linewidth - 14\tabcolsep) * \real{0.1250}}
  >{\raggedright\arraybackslash}p{(\linewidth - 14\tabcolsep) * \real{0.1250}}
  >{\raggedright\arraybackslash}p{(\linewidth - 14\tabcolsep) * \real{0.1250}}
  >{\raggedright\arraybackslash}p{(\linewidth - 14\tabcolsep) * \real{0.1250}}
  >{\raggedright\arraybackslash}p{(\linewidth - 14\tabcolsep) * \real{0.1250}}
  >{\raggedright\arraybackslash}p{(\linewidth - 14\tabcolsep) * \real{0.1250}}
  >{\raggedright\arraybackslash}p{(\linewidth - 14\tabcolsep) * \real{0.1250}}@{}}
\toprule\noalign{}
\begin{minipage}[b]{\linewidth}\raggedright
RUNSPEC
\end{minipage} & \begin{minipage}[b]{\linewidth}\raggedright
GRID
\end{minipage} & \begin{minipage}[b]{\linewidth}\raggedright
EDIT
\end{minipage} & \begin{minipage}[b]{\linewidth}\raggedright
PROPS
\end{minipage} & \begin{minipage}[b]{\linewidth}\raggedright
REGIONS
\end{minipage} & \begin{minipage}[b]{\linewidth}\raggedright
SOLUTION
\end{minipage} & \begin{minipage}[b]{\linewidth}\raggedright
SUMMARY
\end{minipage} & \begin{minipage}[b]{\linewidth}\raggedright
SCHEDULE
\end{minipage} \\
\midrule\noalign{}
\endhead
\bottomrule\noalign{}
\endlastfoot
\end{longtable}

\subsubsection{Description}\label{description-827}

The SMICR keyword defines the initial equilibration microbial
concentration values for all grid cells in the model. The keyword should
only be used if either the
\hyperref[REF_HEADING_KEYWORD_BIOFILM]{BIOFILM} or MICP model has been
activated in the RUNSPEC section.

\begin{tcolorbox}[enhanced jigsaw, title=\textcolor{quarto-callout-note-color}{\faInfo}\hspace{0.5em}{Note}, coltitle=black, breakable, toprule=.15mm, colbacktitle=quarto-callout-note-color!10!white, leftrule=.75mm, opacitybacktitle=0.6, colback=white, bottomtitle=1mm, toptitle=1mm, titlerule=0mm, left=2mm, colframe=quarto-callout-note-color-frame, arc=.35mm, rightrule=.15mm, bottomrule=.15mm, opacityback=0]

This is an OPM Flow specific keyword.

\end{tcolorbox}

\begin{longtable}[]{@{}
  >{\raggedright\arraybackslash}p{(\linewidth - 6\tabcolsep) * \real{0.2500}}
  >{\raggedright\arraybackslash}p{(\linewidth - 6\tabcolsep) * \real{0.2500}}
  >{\raggedright\arraybackslash}p{(\linewidth - 6\tabcolsep) * \real{0.2500}}
  >{\raggedright\arraybackslash}p{(\linewidth - 6\tabcolsep) * \real{0.2500}}@{}}
\toprule\noalign{}
\begin{minipage}[b]{\linewidth}\raggedright
No.
\end{minipage} & \begin{minipage}[b]{\linewidth}\raggedright
Name
\end{minipage} & \begin{minipage}[b]{\linewidth}\raggedright
Description
\end{minipage} & \begin{minipage}[b]{\linewidth}\raggedright
Default
\end{minipage} \\
\midrule\noalign{}
\endhead
\bottomrule\noalign{}
\endlastfoot
Field & Metric & Laboratory & \\
1 & SMICR & SMICR is an array of real numbers that are greater than or
equal to zero assigning the initial equilibration microbial
concentration values to each cell in the model. Repeat counts may be
used, for example 20*0.1500. & None \\
lb/stb & kg/sm3 & gm/scc & \\
Notes: & & & \\
\end{longtable}

\emph{Table 10.46: SMICR Keyword Description}

For both \hyperref[REF_HEADING_KEYWORD_BIOFILM]{BIOFILM} and MICP
models, see also the SBIOF keyword, and for the MICP model, the SCALC,
SOXYG, and SUREA keywords to define the initial state of the model.

\subsubsection{Example}\label{example-473}

\begin{verbatim}
--
--       DEFINE INITIAL EQUILIBRATION MICROBIAL CONCENTRATION FOR ALL CELLS
--       BASED ON NX = 100, NY = 100 AND NZ = 3
--
SMICR
         1000*0.0000 1000*0.0000 1000*0.1500                                   /
\end{verbatim}

The above example defines the initial equilibration microbial
concentration values to be 0.0000 for all the cells in the first and
second layers and finally 0.1500 for all the cells in the third layer.

\subsection{SOIL -- Define the Initial Equilibration Oil Saturation for
All Grid
Blocks}\label{soil-define-the-initial-equilibration-oil-saturation-for-all-grid-blocks}

\begin{longtable}[]{@{}
  >{\raggedright\arraybackslash}p{(\linewidth - 14\tabcolsep) * \real{0.1250}}
  >{\raggedright\arraybackslash}p{(\linewidth - 14\tabcolsep) * \real{0.1250}}
  >{\raggedright\arraybackslash}p{(\linewidth - 14\tabcolsep) * \real{0.1250}}
  >{\raggedright\arraybackslash}p{(\linewidth - 14\tabcolsep) * \real{0.1250}}
  >{\raggedright\arraybackslash}p{(\linewidth - 14\tabcolsep) * \real{0.1250}}
  >{\raggedright\arraybackslash}p{(\linewidth - 14\tabcolsep) * \real{0.1250}}
  >{\raggedright\arraybackslash}p{(\linewidth - 14\tabcolsep) * \real{0.1250}}
  >{\raggedright\arraybackslash}p{(\linewidth - 14\tabcolsep) * \real{0.1250}}@{}}
\toprule\noalign{}
\begin{minipage}[b]{\linewidth}\raggedright
RUNSPEC
\end{minipage} & \begin{minipage}[b]{\linewidth}\raggedright
GRID
\end{minipage} & \begin{minipage}[b]{\linewidth}\raggedright
EDIT
\end{minipage} & \begin{minipage}[b]{\linewidth}\raggedright
PROPS
\end{minipage} & \begin{minipage}[b]{\linewidth}\raggedright
REGIONS
\end{minipage} & \begin{minipage}[b]{\linewidth}\raggedright
SOLUTION
\end{minipage} & \begin{minipage}[b]{\linewidth}\raggedright
SUMMARY
\end{minipage} & \begin{minipage}[b]{\linewidth}\raggedright
SCHEDULE
\end{minipage} \\
\midrule\noalign{}
\endhead
\bottomrule\noalign{}
\endlastfoot
\end{longtable}

\subsubsection{Description}\label{description-828}

The SOIL keyword defines the initial equilibration oil saturation values
for all grid cells in the model and should be used in conjunction with
the PBUB, PDEW, PRESSURE, RS, RV, SGAS and SWAT keywords etc., to fully
describe the initial state of the model. The keyword should only be used
if the oil phase has been activated in the model via the OIL keyword in
the RUNSPEC section.

This is the non-standard method to initialize the model via enumeration
and is seldom employed in the industry. The standard methodology is for
OPM Flow to initialize a model using the parameters on the EQUIL keyword
combined with other keywords to fully describe the initial state of the
model. The keyword can be used with all grid types.

\begin{longtable}[]{@{}
  >{\raggedright\arraybackslash}p{(\linewidth - 6\tabcolsep) * \real{0.2500}}
  >{\raggedright\arraybackslash}p{(\linewidth - 6\tabcolsep) * \real{0.2500}}
  >{\raggedright\arraybackslash}p{(\linewidth - 6\tabcolsep) * \real{0.2500}}
  >{\raggedright\arraybackslash}p{(\linewidth - 6\tabcolsep) * \real{0.2500}}@{}}
\toprule\noalign{}
\begin{minipage}[b]{\linewidth}\raggedright
No.
\end{minipage} & \begin{minipage}[b]{\linewidth}\raggedright
Name
\end{minipage} & \begin{minipage}[b]{\linewidth}\raggedright
Description
\end{minipage} & \begin{minipage}[b]{\linewidth}\raggedright
Default
\end{minipage} \\
\midrule\noalign{}
\endhead
\bottomrule\noalign{}
\endlastfoot
Field & Metric & Laboratory & \\
1 & SOIL & SOIL is an array of real positive numbers that are greater
than or equal to zero and less than or equal to one assigning the
initial equilibration oil saturation values to each cell in the model.
Repeat counts may be used, for example 20*0.600. & None \\
dimensionless & dimensionless & dimensionless & \\
Notes: & & & \\
\end{longtable}

\emph{Table 10.47: SOIL Keyword Description}

Note for two phase runs it is only necessary to enter one saturation
array of one of the phases present in the run (SGAS, SOIL, or SWAT), as
the simulator will calculate the other phases by difference. Similarly
for three phase runs it is only necessary to enter the array data for
two of the phases, as the third saturation will again be calculated by
the simulator.

See also the PBUB, PDEW, PRESSURE, RS, RV, SGAS and SWAT keywords to
fully define the initial state of the model.

\subsubsection{Example}\label{example-474}

\begin{verbatim}
--
--       DEFINE INITIAL EQUILIBRATION OIL SAT VALUES FOR ALL CELLS IN THE MODEL
--       BASED ON NX = 100, NY = 100 AND NZ = 3
--
SOIL
         1000*0.7000    1000*0.6500    1000*0.6000                             /
\end{verbatim}

The above example defines the initial equilibration oil saturation
values to be 0.7000 for all the cells in the first layer, 0.6500 for all
the cells in the second layer, and finally 0.6000 for all the cells in
the third layer.

\subsection{SOLUTION -- Define the Start of the SOLUTION Section of
Keywords}\label{solution-define-the-start-of-the-solution-section-of-keywords}

\begin{longtable}[]{@{}
  >{\raggedright\arraybackslash}p{(\linewidth - 14\tabcolsep) * \real{0.1250}}
  >{\raggedright\arraybackslash}p{(\linewidth - 14\tabcolsep) * \real{0.1250}}
  >{\raggedright\arraybackslash}p{(\linewidth - 14\tabcolsep) * \real{0.1250}}
  >{\raggedright\arraybackslash}p{(\linewidth - 14\tabcolsep) * \real{0.1250}}
  >{\raggedright\arraybackslash}p{(\linewidth - 14\tabcolsep) * \real{0.1250}}
  >{\raggedright\arraybackslash}p{(\linewidth - 14\tabcolsep) * \real{0.1250}}
  >{\raggedright\arraybackslash}p{(\linewidth - 14\tabcolsep) * \real{0.1250}}
  >{\raggedright\arraybackslash}p{(\linewidth - 14\tabcolsep) * \real{0.1250}}@{}}
\toprule\noalign{}
\begin{minipage}[b]{\linewidth}\raggedright
RUNSPEC
\end{minipage} & \begin{minipage}[b]{\linewidth}\raggedright
GRID
\end{minipage} & \begin{minipage}[b]{\linewidth}\raggedright
EDIT
\end{minipage} & \begin{minipage}[b]{\linewidth}\raggedright
PROPS
\end{minipage} & \begin{minipage}[b]{\linewidth}\raggedright
REGIONS
\end{minipage} & \begin{minipage}[b]{\linewidth}\raggedright
SOLUTION
\end{minipage} & \begin{minipage}[b]{\linewidth}\raggedright
SUMMARY
\end{minipage} & \begin{minipage}[b]{\linewidth}\raggedright
SCHEDULE
\end{minipage} \\
\midrule\noalign{}
\endhead
\bottomrule\noalign{}
\endlastfoot
\end{longtable}

\subsubsection{Description}\label{description-829}

The SOLUTION activation keyword marks the end of the REGIONS section and
the start of the SOLUTION section that defines the parameters used to
initialize the model, by:

\begin{itemize}
\tightlist
\item
  defining fluid contacts and pressures, or
\item
  defining pressures and fluid saturations for all cells in the model,
  or
\item
  by restarting from a previously OPM Flow completed run.
\end{itemize}

There is no data required for this keyword.

\subsubsection{Example}\label{example-475}

\begin{verbatim}
-- ==============================================================================
--
-- SOLUTION SECTION
--
-- ==============================================================================
SOLUTION
\end{verbatim}

The above example marks the end of the REGIONS section and the start of
the SOLUTION section in the OPM Flow data input file.

\subsection{SOLVCONC -- Define the Initial Coal Solvent Concentration
for All Grid
Blocks}\label{solvconc-define-the-initial-coal-solvent-concentration-for-all-grid-blocks}

\begin{longtable}[]{@{}
  >{\raggedright\arraybackslash}p{(\linewidth - 14\tabcolsep) * \real{0.1250}}
  >{\raggedright\arraybackslash}p{(\linewidth - 14\tabcolsep) * \real{0.1250}}
  >{\raggedright\arraybackslash}p{(\linewidth - 14\tabcolsep) * \real{0.1250}}
  >{\raggedright\arraybackslash}p{(\linewidth - 14\tabcolsep) * \real{0.1250}}
  >{\raggedright\arraybackslash}p{(\linewidth - 14\tabcolsep) * \real{0.1250}}
  >{\raggedright\arraybackslash}p{(\linewidth - 14\tabcolsep) * \real{0.1250}}
  >{\raggedright\arraybackslash}p{(\linewidth - 14\tabcolsep) * \real{0.1250}}
  >{\raggedright\arraybackslash}p{(\linewidth - 14\tabcolsep) * \real{0.1250}}@{}}
\toprule\noalign{}
\begin{minipage}[b]{\linewidth}\raggedright
RUNSPEC
\end{minipage} & \begin{minipage}[b]{\linewidth}\raggedright
GRID
\end{minipage} & \begin{minipage}[b]{\linewidth}\raggedright
EDIT
\end{minipage} & \begin{minipage}[b]{\linewidth}\raggedright
PROPS
\end{minipage} & \begin{minipage}[b]{\linewidth}\raggedright
REGIONS
\end{minipage} & \begin{minipage}[b]{\linewidth}\raggedright
SOLUTION
\end{minipage} & \begin{minipage}[b]{\linewidth}\raggedright
SUMMARY
\end{minipage} & \begin{minipage}[b]{\linewidth}\raggedright
SCHEDULE
\end{minipage} \\
\midrule\noalign{}
\endhead
\bottomrule\noalign{}
\endlastfoot
\end{longtable}

\subsubsection{Description}\label{description-830}

The SOLVCONC keyword defines the initial coal solvent concentration
values for all matrix grid cells in the model and should be used in
conjunction with the SCVD keyword in the SOLUTION section, to fully
describe the initial state of the model. The keyword should only be used
if the coal phase has been activated in the model via the COAL keyword
in the RUNSPEC section. Note both SOLVCONC and SCVD are optional as the
simulator will calculate the coal solvent concentration based on the
equilibrium concentration and the block pressure.

This keyword is not supported by OPM Flow but would change the results
if supported so the simulation will be stopped.

This is the non-standard method to initialize the model via enumeration
and is seldom employed in the industry. The standard methodology is for
OPM Flow to initialize a model using the parameters on the EQUIL keyword
combined with other keywords to fully describe the initial state of the
model. The keyword can be used with all grid types.

\begin{longtable}[]{@{}
  >{\raggedright\arraybackslash}p{(\linewidth - 6\tabcolsep) * \real{0.2500}}
  >{\raggedright\arraybackslash}p{(\linewidth - 6\tabcolsep) * \real{0.2500}}
  >{\raggedright\arraybackslash}p{(\linewidth - 6\tabcolsep) * \real{0.2500}}
  >{\raggedright\arraybackslash}p{(\linewidth - 6\tabcolsep) * \real{0.2500}}@{}}
\toprule\noalign{}
\begin{minipage}[b]{\linewidth}\raggedright
No.
\end{minipage} & \begin{minipage}[b]{\linewidth}\raggedright
Name
\end{minipage} & \begin{minipage}[b]{\linewidth}\raggedright
Description
\end{minipage} & \begin{minipage}[b]{\linewidth}\raggedright
Default
\end{minipage} \\
\midrule\noalign{}
\endhead
\bottomrule\noalign{}
\endlastfoot
Field & Metric & Laboratory & \\
1 & SOLVCONC & SOLVCONC is an array of real positive numbers that define
the initial equilibration coal solvent concentration values to each
matrix cell in the model. Repeat counts may be used, for example
20*75.0. & None \\
Mscf/ft3 & sm3/m3 & scc/cc & \\
Notes: & & & \\
\end{longtable}

\emph{Table 10.48: SOLVCONC Keyword Description}

See also the SCVD keyword in the SOLUTION section to fully define the
initial state of the model.

\subsubsection{Example}\label{example-476}

\begin{verbatim}
--
--       DEFINE INITIAL EQUILIBRATION COAL SOLVENT CONCENTRATION FOR ALL CELLS
--       BASED ON NX = 100, NY = 100 AND NZ = 6
--
SOLVCONC
         1000*75.500    1000*65.500    1000*60.000                             /
\end{verbatim}

The above example defines the initial coal solvent concentration values
to be 75.500 for all the matrix cells in the first layer, 65.500 for all
the cells in the second layer, and finally 60.000 for all the cells in
the third layer.

\subsection{SOLVFRAC -- Define the Initial Gas Solvent Fraction for All
Grid
Blocks}\label{solvfrac-define-the-initial-gas-solvent-fraction-for-all-grid-blocks}

\begin{longtable}[]{@{}
  >{\raggedright\arraybackslash}p{(\linewidth - 14\tabcolsep) * \real{0.1250}}
  >{\raggedright\arraybackslash}p{(\linewidth - 14\tabcolsep) * \real{0.1250}}
  >{\raggedright\arraybackslash}p{(\linewidth - 14\tabcolsep) * \real{0.1250}}
  >{\raggedright\arraybackslash}p{(\linewidth - 14\tabcolsep) * \real{0.1250}}
  >{\raggedright\arraybackslash}p{(\linewidth - 14\tabcolsep) * \real{0.1250}}
  >{\raggedright\arraybackslash}p{(\linewidth - 14\tabcolsep) * \real{0.1250}}
  >{\raggedright\arraybackslash}p{(\linewidth - 14\tabcolsep) * \real{0.1250}}
  >{\raggedright\arraybackslash}p{(\linewidth - 14\tabcolsep) * \real{0.1250}}@{}}
\toprule\noalign{}
\begin{minipage}[b]{\linewidth}\raggedright
RUNSPEC
\end{minipage} & \begin{minipage}[b]{\linewidth}\raggedright
GRID
\end{minipage} & \begin{minipage}[b]{\linewidth}\raggedright
EDIT
\end{minipage} & \begin{minipage}[b]{\linewidth}\raggedright
PROPS
\end{minipage} & \begin{minipage}[b]{\linewidth}\raggedright
REGIONS
\end{minipage} & \begin{minipage}[b]{\linewidth}\raggedright
SOLUTION
\end{minipage} & \begin{minipage}[b]{\linewidth}\raggedright
SUMMARY
\end{minipage} & \begin{minipage}[b]{\linewidth}\raggedright
SCHEDULE
\end{minipage} \\
\midrule\noalign{}
\endhead
\bottomrule\noalign{}
\endlastfoot
\end{longtable}

\subsubsection{Description}\label{description-831}

The SOLVFRAC keyword defines the initial solvent faction within the gas
phase values for all matrix grid cells in the model. The keyword should
only be used if the coal phase has been activated in the model via the
COAL keyword in the RUNSPEC section.

This keyword is not supported by OPM Flow but would change the results
if supported so the simulation will be stopped.

SOLVFRAC is used with the standard equilibration method to initialize
the model via the EQUIL keyword in the RUNSPEC section, as oppose to the
non-standard enumeration method.

\begin{longtable}[]{@{}
  >{\raggedright\arraybackslash}p{(\linewidth - 6\tabcolsep) * \real{0.2500}}
  >{\raggedright\arraybackslash}p{(\linewidth - 6\tabcolsep) * \real{0.2500}}
  >{\raggedright\arraybackslash}p{(\linewidth - 6\tabcolsep) * \real{0.2500}}
  >{\raggedright\arraybackslash}p{(\linewidth - 6\tabcolsep) * \real{0.2500}}@{}}
\toprule\noalign{}
\begin{minipage}[b]{\linewidth}\raggedright
No.
\end{minipage} & \begin{minipage}[b]{\linewidth}\raggedright
Name
\end{minipage} & \begin{minipage}[b]{\linewidth}\raggedright
Description
\end{minipage} & \begin{minipage}[b]{\linewidth}\raggedright
Default
\end{minipage} \\
\midrule\noalign{}
\endhead
\bottomrule\noalign{}
\endlastfoot
Field & Metric & Laboratory & \\
1 & SOLVFRAC & SOLVFRAC is an array of real positive numbers that define
the initial solvent fraction within the gas phase values for each matrix
cell in the model. Repeat counts may be used, for example 20*0.075. &
None \\
dimensionless & dimensionless & dimensionless & \\
Notes: & & & \\
\end{longtable}

\emph{Table 10.49: SOLVFRAC Keyword Description}

See also the EQUIL keyword in the SOLUTION section to fully define the
initial state of the model.

\subsubsection{Example}\label{example-477}

\begin{verbatim}
--
--       DEFINE INITIAL EQUILIBRATION GAS SOLVENT FRACTION FOR ALL CELLS
--       BASED ON NX = 100, NY = 100 AND NZ = 6
--
SOLVFRAC
         1000*0.0250    1000*0.0350    1000*0.0500                             /
\end{verbatim}

The above example defines the initial gas solvent fraction values to be
0.250 for all the matrix cells in the first layer, 0.0350 for all the
cells in the second layer, and finally 0.0500 for all the cells in the
third layer.

\subsection{SOXYG -- Define The Initial Equilibration Oxygen
Concentration For All Grid
Blocks}\label{soxyg-define-the-initial-equilibration-oxygen-concentration-for-all-grid-blocks}

\begin{longtable}[]{@{}
  >{\raggedright\arraybackslash}p{(\linewidth - 14\tabcolsep) * \real{0.1250}}
  >{\raggedright\arraybackslash}p{(\linewidth - 14\tabcolsep) * \real{0.1250}}
  >{\raggedright\arraybackslash}p{(\linewidth - 14\tabcolsep) * \real{0.1250}}
  >{\raggedright\arraybackslash}p{(\linewidth - 14\tabcolsep) * \real{0.1250}}
  >{\raggedright\arraybackslash}p{(\linewidth - 14\tabcolsep) * \real{0.1250}}
  >{\raggedright\arraybackslash}p{(\linewidth - 14\tabcolsep) * \real{0.1250}}
  >{\raggedright\arraybackslash}p{(\linewidth - 14\tabcolsep) * \real{0.1250}}
  >{\raggedright\arraybackslash}p{(\linewidth - 14\tabcolsep) * \real{0.1250}}@{}}
\toprule\noalign{}
\begin{minipage}[b]{\linewidth}\raggedright
RUNSPEC
\end{minipage} & \begin{minipage}[b]{\linewidth}\raggedright
GRID
\end{minipage} & \begin{minipage}[b]{\linewidth}\raggedright
EDIT
\end{minipage} & \begin{minipage}[b]{\linewidth}\raggedright
PROPS
\end{minipage} & \begin{minipage}[b]{\linewidth}\raggedright
REGIONS
\end{minipage} & \begin{minipage}[b]{\linewidth}\raggedright
SOLUTION
\end{minipage} & \begin{minipage}[b]{\linewidth}\raggedright
SUMMARY
\end{minipage} & \begin{minipage}[b]{\linewidth}\raggedright
SCHEDULE
\end{minipage} \\
\midrule\noalign{}
\endhead
\bottomrule\noalign{}
\endlastfoot
\end{longtable}

\subsubsection{Description}\label{description-832}

The SOXYG keyword defines the initial equilibration oxygen concentration
values for all grid cells in the model. The keyword should only be used
if the MICP model has been activated in the RUNSPEC section.

\begin{tcolorbox}[enhanced jigsaw, title=\textcolor{quarto-callout-note-color}{\faInfo}\hspace{0.5em}{Note}, coltitle=black, breakable, toprule=.15mm, colbacktitle=quarto-callout-note-color!10!white, leftrule=.75mm, opacitybacktitle=0.6, colback=white, bottomtitle=1mm, toptitle=1mm, titlerule=0mm, left=2mm, colframe=quarto-callout-note-color-frame, arc=.35mm, rightrule=.15mm, bottomrule=.15mm, opacityback=0]

This is an OPM Flow specific keyword.

\end{tcolorbox}

\begin{longtable}[]{@{}
  >{\raggedright\arraybackslash}p{(\linewidth - 6\tabcolsep) * \real{0.2500}}
  >{\raggedright\arraybackslash}p{(\linewidth - 6\tabcolsep) * \real{0.2500}}
  >{\raggedright\arraybackslash}p{(\linewidth - 6\tabcolsep) * \real{0.2500}}
  >{\raggedright\arraybackslash}p{(\linewidth - 6\tabcolsep) * \real{0.2500}}@{}}
\toprule\noalign{}
\begin{minipage}[b]{\linewidth}\raggedright
No.
\end{minipage} & \begin{minipage}[b]{\linewidth}\raggedright
Name
\end{minipage} & \begin{minipage}[b]{\linewidth}\raggedright
Description
\end{minipage} & \begin{minipage}[b]{\linewidth}\raggedright
Default
\end{minipage} \\
\midrule\noalign{}
\endhead
\bottomrule\noalign{}
\endlastfoot
Field & Metric & Laboratory & \\
1 & SOXYG & SOXYG is an array of real numbers that are greater than or
equal to zero assigning the initial equilibration oxygen concentration
values to each cell in the model. Repeat counts may be used, for example
20*0.1500 & None \\
lb/stb & kg/sm3 & gm/scc & \\
Notes: & & & \\
\end{longtable}

\emph{Table 10.50: SOXYG Keyword Description}

See also the SBIOF, SCALC, SMICR, and SUREA keywords to define the
initial state of the model.

\subsubsection{Example}\label{example-478}

\begin{verbatim}
--
--       DEFINE INITIAL EQUILIBRATION OXYGEN CONCENTRATION FOR ALL CELLS
--       BASED ON NX = 100, NY = 100 AND NZ = 3
--
SOXYG
         1000*0.0000 1000*0.0000 1000*0.1500                                   /
\end{verbatim}

The above example defines the initial equilibration oxygen concentration
values to be 0.0000 for all the cells in the first and second layers and
finally 0.1500 for all the cells in the third layer.

\subsection{SPOLY -- Define the Initial Equilibration Polymer
Concentration for All Grid
Blocks}\label{spoly-define-the-initial-equilibration-polymer-concentration-for-all-grid-blocks}

\begin{longtable}[]{@{}
  >{\raggedright\arraybackslash}p{(\linewidth - 14\tabcolsep) * \real{0.1250}}
  >{\raggedright\arraybackslash}p{(\linewidth - 14\tabcolsep) * \real{0.1250}}
  >{\raggedright\arraybackslash}p{(\linewidth - 14\tabcolsep) * \real{0.1250}}
  >{\raggedright\arraybackslash}p{(\linewidth - 14\tabcolsep) * \real{0.1250}}
  >{\raggedright\arraybackslash}p{(\linewidth - 14\tabcolsep) * \real{0.1250}}
  >{\raggedright\arraybackslash}p{(\linewidth - 14\tabcolsep) * \real{0.1250}}
  >{\raggedright\arraybackslash}p{(\linewidth - 14\tabcolsep) * \real{0.1250}}
  >{\raggedright\arraybackslash}p{(\linewidth - 14\tabcolsep) * \real{0.1250}}@{}}
\toprule\noalign{}
\begin{minipage}[b]{\linewidth}\raggedright
RUNSPEC
\end{minipage} & \begin{minipage}[b]{\linewidth}\raggedright
GRID
\end{minipage} & \begin{minipage}[b]{\linewidth}\raggedright
EDIT
\end{minipage} & \begin{minipage}[b]{\linewidth}\raggedright
PROPS
\end{minipage} & \begin{minipage}[b]{\linewidth}\raggedright
REGIONS
\end{minipage} & \begin{minipage}[b]{\linewidth}\raggedright
SOLUTION
\end{minipage} & \begin{minipage}[b]{\linewidth}\raggedright
SUMMARY
\end{minipage} & \begin{minipage}[b]{\linewidth}\raggedright
SCHEDULE
\end{minipage} \\
\midrule\noalign{}
\endhead
\bottomrule\noalign{}
\endlastfoot
\end{longtable}

\subsubsection{Description}\label{description-833}

The SPOLY keyword defines the initial equilibration polymer
concentration values for all grid cells in the model and should be used
in conjunction with the PBUB, PDEW, PRESSURE, RS, RV, SGAS, SGAS and
SWAT keywords etc., to fully describe the initial state of the model.
The keyword should only be used if the polymer phase has been activated
in the model via the POLYMER keyword in the RUNSPEC section.

This is the non-standard method to initialize the model via enumeration
and is seldom employed in the industry. The standard methodology is for
OPM Flow to initialize a model using the parameters on the EQUIL keyword
combined with other keywords to fully describe the initial state of the
model. The keyword can be used with all grid types.

\begin{longtable}[]{@{}
  >{\raggedright\arraybackslash}p{(\linewidth - 6\tabcolsep) * \real{0.2500}}
  >{\raggedright\arraybackslash}p{(\linewidth - 6\tabcolsep) * \real{0.2500}}
  >{\raggedright\arraybackslash}p{(\linewidth - 6\tabcolsep) * \real{0.2500}}
  >{\raggedright\arraybackslash}p{(\linewidth - 6\tabcolsep) * \real{0.2500}}@{}}
\toprule\noalign{}
\begin{minipage}[b]{\linewidth}\raggedright
No.
\end{minipage} & \begin{minipage}[b]{\linewidth}\raggedright
Name
\end{minipage} & \begin{minipage}[b]{\linewidth}\raggedright
Description
\end{minipage} & \begin{minipage}[b]{\linewidth}\raggedright
Default
\end{minipage} \\
\midrule\noalign{}
\endhead
\bottomrule\noalign{}
\endlastfoot
Field & Metric & Laboratory & \\
1 & SPOLY & SPOLY is an array of real positive numbers that are greater
than or equal to zero assigning the initial equilibration polymer
concentration values to each cell in the model. Repeat counts may be
used, for example 20*25.0. & None \\
lb/stb & kg/sm3 & gm/scc & \\
Notes: & & & \\
\end{longtable}

\emph{Table 10.51: SPOLY Keyword Description}

See also the PBUB, PDEW, PRESSURE, RS, RV, SGAS, SOIL and SWAT keywords
to fully define the initial state of the model.

\subsubsection{Example}\label{example-479}

\begin{verbatim}
--
--       DEFINE INITIAL EQUILIBRATION POLYMER VALUES FOR ALL CELLS IN THE MODEL
--       BASED ON NX = 100, NY = 100 AND NZ = 3
--
SPOLY
         1000*0.0000    1000*0.0000    1000*15.000                             /
\end{verbatim}

The above example defines the initial equilibration polymer
concentration values to be 0.0000 for all the cells in the first and
second layers and finally 15.000 for all the cells in the third layer.

\subsection{SPOLYMW -- Define The Initial Equilibration Polymer
Molecular Weights For All Grid
Blocks}\label{spolymw-define-the-initial-equilibration-polymer-molecular-weights-for-all-grid-blocks}

\begin{longtable}[]{@{}
  >{\raggedright\arraybackslash}p{(\linewidth - 14\tabcolsep) * \real{0.1250}}
  >{\raggedright\arraybackslash}p{(\linewidth - 14\tabcolsep) * \real{0.1250}}
  >{\raggedright\arraybackslash}p{(\linewidth - 14\tabcolsep) * \real{0.1250}}
  >{\raggedright\arraybackslash}p{(\linewidth - 14\tabcolsep) * \real{0.1250}}
  >{\raggedright\arraybackslash}p{(\linewidth - 14\tabcolsep) * \real{0.1250}}
  >{\raggedright\arraybackslash}p{(\linewidth - 14\tabcolsep) * \real{0.1250}}
  >{\raggedright\arraybackslash}p{(\linewidth - 14\tabcolsep) * \real{0.1250}}
  >{\raggedright\arraybackslash}p{(\linewidth - 14\tabcolsep) * \real{0.1250}}@{}}
\toprule\noalign{}
\begin{minipage}[b]{\linewidth}\raggedright
RUNSPEC
\end{minipage} & \begin{minipage}[b]{\linewidth}\raggedright
GRID
\end{minipage} & \begin{minipage}[b]{\linewidth}\raggedright
EDIT
\end{minipage} & \begin{minipage}[b]{\linewidth}\raggedright
PROPS
\end{minipage} & \begin{minipage}[b]{\linewidth}\raggedright
REGIONS
\end{minipage} & \begin{minipage}[b]{\linewidth}\raggedright
SOLUTION
\end{minipage} & \begin{minipage}[b]{\linewidth}\raggedright
SUMMARY
\end{minipage} & \begin{minipage}[b]{\linewidth}\raggedright
SCHEDULE
\end{minipage} \\
\midrule\noalign{}
\endhead
\bottomrule\noalign{}
\endlastfoot
\end{longtable}

\subsubsection{Description}\label{description-834}

The SPOLYMW keyword defines the initial equilibration polymer molecular
weights for all grid cells in the model and should only be be used with
OPM Flow's Polymer Molecular Weight Transport option, together with the
other standard equilibration keywords, in order to fully describe the
initial state of the model.

This keyword should only be used if the POLYMER and POLYMW keywords in
the RUNSPEC section are also activated.

\begin{tcolorbox}[enhanced jigsaw, title=\textcolor{quarto-callout-note-color}{\faInfo}\hspace{0.5em}{Note}, coltitle=black, breakable, toprule=.15mm, colbacktitle=quarto-callout-note-color!10!white, leftrule=.75mm, opacitybacktitle=0.6, colback=white, bottomtitle=1mm, toptitle=1mm, titlerule=0mm, left=2mm, colframe=quarto-callout-note-color-frame, arc=.35mm, rightrule=.15mm, bottomrule=.15mm, opacityback=0]

This is an OPM Flow specific keyword that employs an alternative polymer
flood model based on a Polymer Molecular Weight Transport equation, that
is not available in the commercial simulator. The model has been tested
using metric units; however, using either field or laboratory units with
the option should be considered experimental.

\end{tcolorbox}

\begin{longtable}[]{@{}
  >{\raggedright\arraybackslash}p{(\linewidth - 6\tabcolsep) * \real{0.2500}}
  >{\raggedright\arraybackslash}p{(\linewidth - 6\tabcolsep) * \real{0.2500}}
  >{\raggedright\arraybackslash}p{(\linewidth - 6\tabcolsep) * \real{0.2500}}
  >{\raggedright\arraybackslash}p{(\linewidth - 6\tabcolsep) * \real{0.2500}}@{}}
\toprule\noalign{}
\begin{minipage}[b]{\linewidth}\raggedright
No.
\end{minipage} & \begin{minipage}[b]{\linewidth}\raggedright
Name
\end{minipage} & \begin{minipage}[b]{\linewidth}\raggedright
Description
\end{minipage} & \begin{minipage}[b]{\linewidth}\raggedright
Default
\end{minipage} \\
\midrule\noalign{}
\endhead
\bottomrule\noalign{}
\endlastfoot
Field & Metric & Laboratory & \\
1 & SPOLYMW & SPOLYMW is an array of real positive numbers that are
greater than or equal to zero assigning the initial equilibration
polymer molecular weights to each cell in the model. Repeat counts may
be used, for example 20*5.0 & 0,0 \\
lb/lb-M & kg/kg-M & gm/gm-M & \\
Notes: & & & \\
\end{longtable}

\emph{Table 10.52: SPOLYMW Keyword Description}

See also the PBUB, PDEW, PRESSURE, RS, RV, SGAS, SOIL and SWAT keywords
to fully define the initial state of the model.

\subsubsection{Example}\label{example-480}

\begin{verbatim}
--
--       INITIAL EQUILIBRATION POLYMER MOLECULAR WEIGHTS FOR ALL CELLS
--
--       ARRAY       CONSTANT     ---------- BOX ---------
--                                I1  I2   J1  J2   K1  K2
EQUALS
         SPOLYMW     0.0000       1*  1*   1*  1*    1   5 / LAYERS 1 TO 5
         SPOLYMW     5.0000       1*  1*   1*  1*    6   7 / LAYERS 6 TO 7
         SPOLYMW     0.0000       1*  1*   1*  1*    8  20 / LAYERS 8 TO 20
/
\end{verbatim}

The above example defines the initial equilibration polymer molecular
weights to be 0.0000 for all the cells, except for layers six to seven,
where the polymer molecular weight is set to five for these cells.

\subsection{SSOL -- Define the Initial Equilibration Solvent Saturation
for All Grid
Blocks}\label{ssol-define-the-initial-equilibration-solvent-saturation-for-all-grid-blocks}

\begin{longtable}[]{@{}
  >{\raggedright\arraybackslash}p{(\linewidth - 14\tabcolsep) * \real{0.1250}}
  >{\raggedright\arraybackslash}p{(\linewidth - 14\tabcolsep) * \real{0.1250}}
  >{\raggedright\arraybackslash}p{(\linewidth - 14\tabcolsep) * \real{0.1250}}
  >{\raggedright\arraybackslash}p{(\linewidth - 14\tabcolsep) * \real{0.1250}}
  >{\raggedright\arraybackslash}p{(\linewidth - 14\tabcolsep) * \real{0.1250}}
  >{\raggedright\arraybackslash}p{(\linewidth - 14\tabcolsep) * \real{0.1250}}
  >{\raggedright\arraybackslash}p{(\linewidth - 14\tabcolsep) * \real{0.1250}}
  >{\raggedright\arraybackslash}p{(\linewidth - 14\tabcolsep) * \real{0.1250}}@{}}
\toprule\noalign{}
\begin{minipage}[b]{\linewidth}\raggedright
RUNSPEC
\end{minipage} & \begin{minipage}[b]{\linewidth}\raggedright
GRID
\end{minipage} & \begin{minipage}[b]{\linewidth}\raggedright
EDIT
\end{minipage} & \begin{minipage}[b]{\linewidth}\raggedright
PROPS
\end{minipage} & \begin{minipage}[b]{\linewidth}\raggedright
REGIONS
\end{minipage} & \begin{minipage}[b]{\linewidth}\raggedright
SOLUTION
\end{minipage} & \begin{minipage}[b]{\linewidth}\raggedright
SUMMARY
\end{minipage} & \begin{minipage}[b]{\linewidth}\raggedright
SCHEDULE
\end{minipage} \\
\midrule\noalign{}
\endhead
\bottomrule\noalign{}
\endlastfoot
\end{longtable}

\subsubsection{Description}\label{description-835}

The SSOL keyword defines the initial equilibration solvent saturation
values for all grid cells in the model and should be used in conjunction
with the PBUB, PDEW, PRESSURE, RS, RV, SGAS, SOIL and SWAT keywords
etc., to fully describe the initial state of the model. The keyword
should only be used if the solvent phase has been activated in the model
via the SOLVENT keyword in the RUNSPEC section.

This is the non-standard method to initialize the model via enumeration
and is seldom employed in the industry. The standard methodology is for
OPM Flow to initialize a model using the parameters on the EQUIL keyword
combined with other keywords to fully describe the initial state of the
model. The keyword can be used with all grid types.

\begin{longtable}[]{@{}
  >{\raggedright\arraybackslash}p{(\linewidth - 6\tabcolsep) * \real{0.2500}}
  >{\raggedright\arraybackslash}p{(\linewidth - 6\tabcolsep) * \real{0.2500}}
  >{\raggedright\arraybackslash}p{(\linewidth - 6\tabcolsep) * \real{0.2500}}
  >{\raggedright\arraybackslash}p{(\linewidth - 6\tabcolsep) * \real{0.2500}}@{}}
\toprule\noalign{}
\begin{minipage}[b]{\linewidth}\raggedright
No.
\end{minipage} & \begin{minipage}[b]{\linewidth}\raggedright
Name
\end{minipage} & \begin{minipage}[b]{\linewidth}\raggedright
Description
\end{minipage} & \begin{minipage}[b]{\linewidth}\raggedright
Default
\end{minipage} \\
\midrule\noalign{}
\endhead
\bottomrule\noalign{}
\endlastfoot
Field & Metric & Laboratory & \\
1 & SSOL & SSOL is an array of real positive numbers that are greater
than or equal to zero and less than or equal to one assigning the
initial equilibration solvent saturation values to each cell in the
model. Repeat counts may be used, for example 20*0.000. & None \\
dimensionless & dimensionless & dimensionless & \\
Notes: & & & \\
\end{longtable}

\emph{Table 10.53: SSOL Keyword Description}

See also the PBUB, PDEW, PRESSURE, RS, RV, SGAS, SOIL, and SWAT keywords
to fully define the initial state of the model.

\subsubsection{Example}\label{example-481}

\begin{verbatim}
--
--       DEFINE INITIAL EQUILIBRATION GAS SAT VALUES FOR ALL CELLS IN THE MODEL
--       BASED ON NX = 100, NY = 100 AND NZ = 3
--
SSOL
         1000*0.0000    1000*0.0000    1000*0.0000                             /
\end{verbatim}

The above example defines the initial equilibration solvent saturation
values to be 0.0 for all the cells in the in the model.

\subsection{SUREA -- Define The Initial Equilibration Urea Concentration
For All Grid
Blocks}\label{surea-define-the-initial-equilibration-urea-concentration-for-all-grid-blocks}

\begin{longtable}[]{@{}
  >{\raggedright\arraybackslash}p{(\linewidth - 14\tabcolsep) * \real{0.1250}}
  >{\raggedright\arraybackslash}p{(\linewidth - 14\tabcolsep) * \real{0.1250}}
  >{\raggedright\arraybackslash}p{(\linewidth - 14\tabcolsep) * \real{0.1250}}
  >{\raggedright\arraybackslash}p{(\linewidth - 14\tabcolsep) * \real{0.1250}}
  >{\raggedright\arraybackslash}p{(\linewidth - 14\tabcolsep) * \real{0.1250}}
  >{\raggedright\arraybackslash}p{(\linewidth - 14\tabcolsep) * \real{0.1250}}
  >{\raggedright\arraybackslash}p{(\linewidth - 14\tabcolsep) * \real{0.1250}}
  >{\raggedright\arraybackslash}p{(\linewidth - 14\tabcolsep) * \real{0.1250}}@{}}
\toprule\noalign{}
\begin{minipage}[b]{\linewidth}\raggedright
RUNSPEC
\end{minipage} & \begin{minipage}[b]{\linewidth}\raggedright
GRID
\end{minipage} & \begin{minipage}[b]{\linewidth}\raggedright
EDIT
\end{minipage} & \begin{minipage}[b]{\linewidth}\raggedright
PROPS
\end{minipage} & \begin{minipage}[b]{\linewidth}\raggedright
REGIONS
\end{minipage} & \begin{minipage}[b]{\linewidth}\raggedright
SOLUTION
\end{minipage} & \begin{minipage}[b]{\linewidth}\raggedright
SUMMARY
\end{minipage} & \begin{minipage}[b]{\linewidth}\raggedright
SCHEDULE
\end{minipage} \\
\midrule\noalign{}
\endhead
\bottomrule\noalign{}
\endlastfoot
\end{longtable}

\subsubsection{Description}\label{description-836}

The SUREA keyword defines the initial equilibration urea concentration
values for all grid cells in the model. The keyword should only be used
if the MICP model has been activated in the RUNSPEC section.

\begin{tcolorbox}[enhanced jigsaw, title=\textcolor{quarto-callout-note-color}{\faInfo}\hspace{0.5em}{Note}, coltitle=black, breakable, toprule=.15mm, colbacktitle=quarto-callout-note-color!10!white, leftrule=.75mm, opacitybacktitle=0.6, colback=white, bottomtitle=1mm, toptitle=1mm, titlerule=0mm, left=2mm, colframe=quarto-callout-note-color-frame, arc=.35mm, rightrule=.15mm, bottomrule=.15mm, opacityback=0]

This is an OPM Flow specific keyword.

\end{tcolorbox}

\begin{longtable}[]{@{}
  >{\raggedright\arraybackslash}p{(\linewidth - 6\tabcolsep) * \real{0.2500}}
  >{\raggedright\arraybackslash}p{(\linewidth - 6\tabcolsep) * \real{0.2500}}
  >{\raggedright\arraybackslash}p{(\linewidth - 6\tabcolsep) * \real{0.2500}}
  >{\raggedright\arraybackslash}p{(\linewidth - 6\tabcolsep) * \real{0.2500}}@{}}
\toprule\noalign{}
\begin{minipage}[b]{\linewidth}\raggedright
No.
\end{minipage} & \begin{minipage}[b]{\linewidth}\raggedright
Name
\end{minipage} & \begin{minipage}[b]{\linewidth}\raggedright
Description
\end{minipage} & \begin{minipage}[b]{\linewidth}\raggedright
Default
\end{minipage} \\
\midrule\noalign{}
\endhead
\bottomrule\noalign{}
\endlastfoot
Field & Metric & Laboratory & \\
1 & SUREA & SUREA is an array of real numbers that are greater than or
equal to zero assigning the initial equilibration urea concentration
values to each cell in the model. Repeat counts may be used, for example
20*30.0. & None \\
lb/stb & kg/sm3 & gm/scc & \\
Notes: & & & \\
\end{longtable}

\emph{Table 10.54: SUREA Keyword Description}

See also the SBIOF, SCALC, SMICR, and SOXYG keywords to define the
initial state of the model.

\subsubsection{Example}\label{example-482}

\begin{verbatim}
--
--       DEFINE INITIAL EQUILIBRATION UREA CONCENTRATION FOR ALL CELLS
--       BASED ON NX = 100, NY = 100 AND NZ = 3
--
SUREA
         1000*0.0000 1000*0.0000 1000*20.0                                     /
\end{verbatim}

The above example defines the initial equilibration urea concentration
values to be 0.0000 for all the cells in the first and second layers and
finally 20.0 for all the cells in the third layer.

\subsection{SURF -- Define the Initial Equilibration Surfactant
Concentration for All Grid
Blocks}\label{surf-define-the-initial-equilibration-surfactant-concentration-for-all-grid-blocks}

\begin{longtable}[]{@{}
  >{\raggedright\arraybackslash}p{(\linewidth - 14\tabcolsep) * \real{0.1250}}
  >{\raggedright\arraybackslash}p{(\linewidth - 14\tabcolsep) * \real{0.1250}}
  >{\raggedright\arraybackslash}p{(\linewidth - 14\tabcolsep) * \real{0.1250}}
  >{\raggedright\arraybackslash}p{(\linewidth - 14\tabcolsep) * \real{0.1250}}
  >{\raggedright\arraybackslash}p{(\linewidth - 14\tabcolsep) * \real{0.1250}}
  >{\raggedright\arraybackslash}p{(\linewidth - 14\tabcolsep) * \real{0.1250}}
  >{\raggedright\arraybackslash}p{(\linewidth - 14\tabcolsep) * \real{0.1250}}
  >{\raggedright\arraybackslash}p{(\linewidth - 14\tabcolsep) * \real{0.1250}}@{}}
\toprule\noalign{}
\begin{minipage}[b]{\linewidth}\raggedright
RUNSPEC
\end{minipage} & \begin{minipage}[b]{\linewidth}\raggedright
GRID
\end{minipage} & \begin{minipage}[b]{\linewidth}\raggedright
EDIT
\end{minipage} & \begin{minipage}[b]{\linewidth}\raggedright
PROPS
\end{minipage} & \begin{minipage}[b]{\linewidth}\raggedright
REGIONS
\end{minipage} & \begin{minipage}[b]{\linewidth}\raggedright
SOLUTION
\end{minipage} & \begin{minipage}[b]{\linewidth}\raggedright
SUMMARY
\end{minipage} & \begin{minipage}[b]{\linewidth}\raggedright
SCHEDULE
\end{minipage} \\
\midrule\noalign{}
\endhead
\bottomrule\noalign{}
\endlastfoot
\end{longtable}

\subsubsection{Description}\label{description-837}

The SURF keyword defines the initial equilibration surfactant
concentration values for all grid cells in the model and should be used
in conjunction with the PBUB, PDEW, PRESSURE, RS, RV, SGAS, SGAS and
SWAT keywords etc., to fully describe the initial state of the model.
The keyword should only be used if the surfactant phase has been
activated in the model via the SURFACT keyword in the RUNSPEC section.

This keyword is not supported by OPM Flow but would change the results
if supported so the simulation will be stopped.

This is the non-standard method to initialize the model via enumeration
and is seldom employed in the industry. The standard methodology is for
OPM Flow to initialize a model using the parameters on the EQUIL keyword
combined with other keywords to fully describe the initial state of the
model. The keyword can be used with all grid types.

\begin{longtable}[]{@{}
  >{\raggedright\arraybackslash}p{(\linewidth - 6\tabcolsep) * \real{0.2500}}
  >{\raggedright\arraybackslash}p{(\linewidth - 6\tabcolsep) * \real{0.2500}}
  >{\raggedright\arraybackslash}p{(\linewidth - 6\tabcolsep) * \real{0.2500}}
  >{\raggedright\arraybackslash}p{(\linewidth - 6\tabcolsep) * \real{0.2500}}@{}}
\toprule\noalign{}
\begin{minipage}[b]{\linewidth}\raggedright
No.
\end{minipage} & \begin{minipage}[b]{\linewidth}\raggedright
Name
\end{minipage} & \begin{minipage}[b]{\linewidth}\raggedright
Description
\end{minipage} & \begin{minipage}[b]{\linewidth}\raggedright
Default
\end{minipage} \\
\midrule\noalign{}
\endhead
\bottomrule\noalign{}
\endlastfoot
Field & Metric & Laboratory & \\
1 & SURF & SURF is an array of real positive numbers that are greater
than or equal to zero assigning the initial equilibration surfactant
concentration values to each cell in the model. Repeat counts may be
used, for example 20*25.0. & None \\
lb/stb & kg/sm3 & gm/scc & \\
Notes: & & & \\
\end{longtable}

\emph{Table 10.55: SURF Keyword Description}

See also the PBUB, PDEW, PRESSURE, RS, RV, SGAS, SOIL and SWAT keywords
to fully define the initial state of the model.

\subsubsection{Example}\label{example-483}

\begin{verbatim}
--
--       DEFINE INITIAL EQUILIBRATION SURFACTANT VALUES FOR ALL CELLS
--       BASED ON NX = 100, NY = 100 AND NZ = 3
--
SURF
         1000*0.0000    1000*0.0000    1000*0.2500                             /
\end{verbatim}

The above example defines the initial equilibration surfactant
concentration values to be 0.0000 for all the cells in the first and
second layers and finally 0.2500 for all the cells in the third layer.

\subsection{SWAT -- Define the Initial Equilibration Water Saturation
for All Grid
Blocks}\label{swat-define-the-initial-equilibration-water-saturation-for-all-grid-blocks}

\begin{longtable}[]{@{}
  >{\raggedright\arraybackslash}p{(\linewidth - 14\tabcolsep) * \real{0.1250}}
  >{\raggedright\arraybackslash}p{(\linewidth - 14\tabcolsep) * \real{0.1250}}
  >{\raggedright\arraybackslash}p{(\linewidth - 14\tabcolsep) * \real{0.1250}}
  >{\raggedright\arraybackslash}p{(\linewidth - 14\tabcolsep) * \real{0.1250}}
  >{\raggedright\arraybackslash}p{(\linewidth - 14\tabcolsep) * \real{0.1250}}
  >{\raggedright\arraybackslash}p{(\linewidth - 14\tabcolsep) * \real{0.1250}}
  >{\raggedright\arraybackslash}p{(\linewidth - 14\tabcolsep) * \real{0.1250}}
  >{\raggedright\arraybackslash}p{(\linewidth - 14\tabcolsep) * \real{0.1250}}@{}}
\toprule\noalign{}
\begin{minipage}[b]{\linewidth}\raggedright
RUNSPEC
\end{minipage} & \begin{minipage}[b]{\linewidth}\raggedright
GRID
\end{minipage} & \begin{minipage}[b]{\linewidth}\raggedright
EDIT
\end{minipage} & \begin{minipage}[b]{\linewidth}\raggedright
PROPS
\end{minipage} & \begin{minipage}[b]{\linewidth}\raggedright
REGIONS
\end{minipage} & \begin{minipage}[b]{\linewidth}\raggedright
SOLUTION
\end{minipage} & \begin{minipage}[b]{\linewidth}\raggedright
SUMMARY
\end{minipage} & \begin{minipage}[b]{\linewidth}\raggedright
SCHEDULE
\end{minipage} \\
\midrule\noalign{}
\endhead
\bottomrule\noalign{}
\endlastfoot
\end{longtable}

\subsubsection{Description}\label{description-838}

The SWAT keyword defines the initial equilibration water saturation
values for all grid cells in the model and should be used in conjunction
with the PBUB, PDEW, PRESSURE, RS, RV, SGAS and SOIL keywords etc., to
fully describe the initial state of the model. The keyword should only
be used if the water phase has been activated in the model via the WATER
keyword in the RUNSPEC section.

This is the non-standard method to initialize the model via enumeration
and is seldom employed in the industry. The standard methodology is for
OPM Flow to initialize a model using the parameters on the EQUIL keyword
combined with other keywords to fully describe the initial state of the
model. The keyword can be used with all grid types.

\begin{longtable}[]{@{}
  >{\raggedright\arraybackslash}p{(\linewidth - 6\tabcolsep) * \real{0.2500}}
  >{\raggedright\arraybackslash}p{(\linewidth - 6\tabcolsep) * \real{0.2500}}
  >{\raggedright\arraybackslash}p{(\linewidth - 6\tabcolsep) * \real{0.2500}}
  >{\raggedright\arraybackslash}p{(\linewidth - 6\tabcolsep) * \real{0.2500}}@{}}
\toprule\noalign{}
\begin{minipage}[b]{\linewidth}\raggedright
No.
\end{minipage} & \begin{minipage}[b]{\linewidth}\raggedright
Name
\end{minipage} & \begin{minipage}[b]{\linewidth}\raggedright
Description
\end{minipage} & \begin{minipage}[b]{\linewidth}\raggedright
Default
\end{minipage} \\
\midrule\noalign{}
\endhead
\bottomrule\noalign{}
\endlastfoot
Field & Metric & Laboratory & \\
1 & SWAT & SWAT is an array of real positive numbers that are greater
than or equal to zero and less than or equal to one assigning the
initial equilibration water saturation values to each cell in the model.
Repeat counts may be used, for example 20*0.300. & None \\
dimensionless & dimensionless & dimensionless & \\
Notes: & & & \\
\end{longtable}

\emph{Table 10.56: SWAT Keyword Description}

Note for two phase runs it is only necessary to enter one saturation
array of one of the phases present in the run (SGAS, SOIL, or SWAT), as
the simulator will calculate the other phases by difference. Similarly
for three phase runs it is only necessary to enter the array data for
two of the phases, as the third saturation will again be calculated by
the simulator.

See also the PBUB, PDEW, PRESSURE, RS, RV, SGAS and SOIL keywords to
fully define the initial state of the model.

\subsubsection{Example}\label{example-484}

\begin{verbatim}
--
--       DEFINE INITIAL EQUILIBRATION WAT SAT VALUES FOR ALL CELLS IN THE MODEL
--       BASED ON NX = 100, NY = 100 AND NZ = 3
--
SWAT
         1000*0.2000    1000*0.2500    1000*0.4500                             /
\end{verbatim}

The above example defines the initial equilibration water saturation
values to be 0.2000 for all the cells in the first layer, 0.2500 for all
the cells in the second layer, and finally 0.4500 for all the cells in
the third layer.

\subsection{TBLK -- Define Tracer Initial Grid Block
Concentrations}\label{tblk-define-tracer-initial-grid-block-concentrations}

\begin{longtable}[]{@{}
  >{\raggedright\arraybackslash}p{(\linewidth - 14\tabcolsep) * \real{0.1250}}
  >{\raggedright\arraybackslash}p{(\linewidth - 14\tabcolsep) * \real{0.1250}}
  >{\raggedright\arraybackslash}p{(\linewidth - 14\tabcolsep) * \real{0.1250}}
  >{\raggedright\arraybackslash}p{(\linewidth - 14\tabcolsep) * \real{0.1250}}
  >{\raggedright\arraybackslash}p{(\linewidth - 14\tabcolsep) * \real{0.1250}}
  >{\raggedright\arraybackslash}p{(\linewidth - 14\tabcolsep) * \real{0.1250}}
  >{\raggedright\arraybackslash}p{(\linewidth - 14\tabcolsep) * \real{0.1250}}
  >{\raggedright\arraybackslash}p{(\linewidth - 14\tabcolsep) * \real{0.1250}}@{}}
\toprule\noalign{}
\begin{minipage}[b]{\linewidth}\raggedright
RUNSPEC
\end{minipage} & \begin{minipage}[b]{\linewidth}\raggedright
GRID
\end{minipage} & \begin{minipage}[b]{\linewidth}\raggedright
EDIT
\end{minipage} & \begin{minipage}[b]{\linewidth}\raggedright
PROPS
\end{minipage} & \begin{minipage}[b]{\linewidth}\raggedright
REGIONS
\end{minipage} & \begin{minipage}[b]{\linewidth}\raggedright
SOLUTION
\end{minipage} & \begin{minipage}[b]{\linewidth}\raggedright
SUMMARY
\end{minipage} & \begin{minipage}[b]{\linewidth}\raggedright
SCHEDULE
\end{minipage} \\
\midrule\noalign{}
\endhead
\bottomrule\noalign{}
\endlastfoot
\end{longtable}

\subsubsection{Description}\label{description-839}

TBLK keyword defines the initial tracer concentration for all or
selected cells in the model, for when the TRACERS keyword in the RUNSPEC
section has declared the maximum number of tracers for each phase, and
the TRACER keyword in the PROPS section has defined the tracer. This
keyword is not in the standard keyword format due to the tracer name
being concatenated to the keyword TBLK to fully define the tracer being
initialized.

\begin{longtable}[]{@{}
  >{\raggedright\arraybackslash}p{(\linewidth - 6\tabcolsep) * \real{0.2500}}
  >{\raggedright\arraybackslash}p{(\linewidth - 6\tabcolsep) * \real{0.2500}}
  >{\raggedright\arraybackslash}p{(\linewidth - 6\tabcolsep) * \real{0.2500}}
  >{\raggedright\arraybackslash}p{(\linewidth - 6\tabcolsep) * \real{0.2500}}@{}}
\toprule\noalign{}
\begin{minipage}[b]{\linewidth}\raggedright
No.
\end{minipage} & \begin{minipage}[b]{\linewidth}\raggedright
Name
\end{minipage} & \begin{minipage}[b]{\linewidth}\raggedright
Description
\end{minipage} & \begin{minipage}[b]{\linewidth}\raggedright
Default
\end{minipage} \\
\midrule\noalign{}
\endhead
\bottomrule\noalign{}
\endlastfoot
Field & Metric & Laboratory & \\
1 & NAME & A character string of up to eight characters, consisting of
TBLK as the first four characters followed by a four letter character
string defining the tracer's name. The fifth character should either be
the letter F or the letter S, that indicates the state of the tracer
either to be free (F) or in solution (S). For example, TBLKFIGS (free)
or TBLKSIGS (solution). The last three characters of NAME (the effective
tracer name) must also match an entry on the TRACER keyword's NAME
parameter, in the PROPS section. Note it is best to void names beginning
with the letters F, S, and T as these names may create naming issues in
post-processing software. & None \\
2 & TBLK & TBLK is an array of real numbers greater than or equal to
zero, that are assigned the tracer concentration values for each cell in
the model or the current input BOX. Repeat counts may be used, for
example 200*0.0. The units for the tracer, if required, are set on the
TRACER keyword in the PROPS section. This should be the same as the
PHASE in the model. & None \\
Liquid: TBLK/stb Gas: TBLK/Mscf & Liquid: TBLK/sm3 Gas: TBLK/sm3 &
Liquid: TBLK/scc Gas: TBLK/scc & \\
Notes: & & & \\
\end{longtable}

\emph{Table 10.57: TBLK Keyword Description}

See also the TRACERS keyword in the RUNSPEC section to declared the
maximum number of tracers for each phase, the TRACER keyword in the
PROPS section to define the tracer, and the WTRACER keyword in the
SCHEDULE section that defines the wells injecting the tracer.

\begin{tcolorbox}[enhanced jigsaw, title=\textcolor{quarto-callout-note-color}{\faInfo}\hspace{0.5em}{Note}, coltitle=black, breakable, toprule=.15mm, colbacktitle=quarto-callout-note-color!10!white, leftrule=.75mm, opacitybacktitle=0.6, colback=white, bottomtitle=1mm, toptitle=1mm, titlerule=0mm, left=2mm, colframe=quarto-callout-note-color-frame, arc=.35mm, rightrule=.15mm, bottomrule=.15mm, opacityback=0]

Currently, one cannot initialize tracers using the EQUALS keyword.
Instead use the array format, that is the keyword followed by the
required number of values, or the TVDP keyword in the SOLUTION section
to set the initial tracer concentrations as a function of depth.

\end{tcolorbox}

\subsubsection{Example}\label{example-485}

The following TRACERS keyword in the RUNSPEC section declares the number
of tracers in the model.

\begin{verbatim}
--
--       NUMBER AND TYPE OF TRACERS
--       NO OIL  NO WAT  NO GAS  NO ENV  DIFF    MAX    MIN    TRACER
--       TRACERS TRACERS TRACERS TRACERS CONTL   NONLIN NONLIN NONLIN
TRACERS
         0       0       1       0      'NODIFF' 1*     1*     1*              /
\end{verbatim}

And the TRACER keyword in the PROPS section declares the tracer name and
the phase for the tracer.

\begin{verbatim}
--
--       DEFINE TRACER NAMES
--
--       TRACER   TRACER
--       NAME     PHASE
--       ------   ------
TRACER
        'IGS'     'GAS'                                    / INJECTED GAS
/
\end{verbatim}

Finally, the TBLK keyword in the SOLUTION section sets the initial
tracer grid block concentrations in both the free and solution states.

\begin{verbatim}
--
--       INITIAL TRACER CONCENTRATIONS
--
TBLKFIGS
         1000*0.0                       / TRACER FIGS CONCENTRATIONS

TBLKSIGS
         1000*0.0                       / TRACER SIGS CONCENTRATIONS
\end{verbatim}

Here the initial concentrations are set to zero.

Then in the SCHEDULE section one can us the WTRACER keyword to define
the well injecting the tracer and the tracer concentration being
injected,.

\begin{verbatim}
--
--       DEFINE CONCENTRATION OF TRACERS IN THE INJECTION STREAMS,
--       INJECTION TRACER CONCENTRATIONS NOT DEFINED USING THE WTRACER
--       KEYWORD ARE ASSUMED TO BE ZERO.
--
-- WELL  NAME    TRACER  TRACER  TRACER
-- NAME  TRACER  VALUE   CUM     GROUP
WTRACER
'GI01'   'IGS'   1.0                                                            /
/
\end{verbatim}

In this case, well GI01 is a gas injection well injecting gas with a
tracer concentration of 1.0. The example shows how to track dry gas
injection in a gas condensate reservoir, although, the example can be
used for any type of gas injection.

\subsection{TEMPI -- Define the Initial Temperature Values for All
Cells}\label{tempi-define-the-initial-temperature-values-for-all-cells}

\begin{longtable}[]{@{}
  >{\raggedright\arraybackslash}p{(\linewidth - 14\tabcolsep) * \real{0.1250}}
  >{\raggedright\arraybackslash}p{(\linewidth - 14\tabcolsep) * \real{0.1250}}
  >{\raggedright\arraybackslash}p{(\linewidth - 14\tabcolsep) * \real{0.1250}}
  >{\raggedright\arraybackslash}p{(\linewidth - 14\tabcolsep) * \real{0.1250}}
  >{\raggedright\arraybackslash}p{(\linewidth - 14\tabcolsep) * \real{0.1250}}
  >{\raggedright\arraybackslash}p{(\linewidth - 14\tabcolsep) * \real{0.1250}}
  >{\raggedright\arraybackslash}p{(\linewidth - 14\tabcolsep) * \real{0.1250}}
  >{\raggedright\arraybackslash}p{(\linewidth - 14\tabcolsep) * \real{0.1250}}@{}}
\toprule\noalign{}
\begin{minipage}[b]{\linewidth}\raggedright
RUNSPEC
\end{minipage} & \begin{minipage}[b]{\linewidth}\raggedright
GRID
\end{minipage} & \begin{minipage}[b]{\linewidth}\raggedright
EDIT
\end{minipage} & \begin{minipage}[b]{\linewidth}\raggedright
PROPS
\end{minipage} & \begin{minipage}[b]{\linewidth}\raggedright
REGIONS
\end{minipage} & \begin{minipage}[b]{\linewidth}\raggedright
SOLUTION
\end{minipage} & \begin{minipage}[b]{\linewidth}\raggedright
SUMMARY
\end{minipage} & \begin{minipage}[b]{\linewidth}\raggedright
SCHEDULE
\end{minipage} \\
\midrule\noalign{}
\endhead
\bottomrule\noalign{}
\endlastfoot
\end{longtable}

\subsubsection{Description}\label{description-840}

This keyword defines the initial reservoir temperature for each cell in
the model. The keyword is used to explicitly define the initial
reservoir temperature via the Enumeration Initialization method rather
than defining a uniform initial temperature or defining temperature
versus depth tables.

The initial reservoir temperature must be defined when OPM Flow's
thermal option has been activated by the THERMAL keyword in the RUNSPEC
section. Note this is different to the commercial simulator that uses
the TEMP keyword in the RUNSPEC section to activate the black-oil
temperature model, and the THERMAL keyword to activate the compositional
thermal model.

The initial reservoir temperature should be defined when OPM Flow's CO2
or H2 storage option has been activated by the CO2STORE or
\hyperref[REF_HEADING_KEYWORD_H2STORE]{H2STORE} keyword in the RUNSPEC
section.

The keyword can be used with all grid types.

\begin{longtable}[]{@{}
  >{\raggedright\arraybackslash}p{(\linewidth - 6\tabcolsep) * \real{0.2500}}
  >{\raggedright\arraybackslash}p{(\linewidth - 6\tabcolsep) * \real{0.2500}}
  >{\raggedright\arraybackslash}p{(\linewidth - 6\tabcolsep) * \real{0.2500}}
  >{\raggedright\arraybackslash}p{(\linewidth - 6\tabcolsep) * \real{0.2500}}@{}}
\toprule\noalign{}
\begin{minipage}[b]{\linewidth}\raggedright
No.
\end{minipage} & \begin{minipage}[b]{\linewidth}\raggedright
Name
\end{minipage} & \begin{minipage}[b]{\linewidth}\raggedright
Description
\end{minipage} & \begin{minipage}[b]{\linewidth}\raggedright
Default
\end{minipage} \\
\midrule\noalign{}
\endhead
\bottomrule\noalign{}
\endlastfoot
Field & Metric & Laboratory & \\
1 & TEMPI & TEMPI is an array of real positive numbers assigning the
initial temperature to each cell in the model. Repeat counts may be
used, for example 20*100.0. & None \\
oF & oC & oC & \\
Notes: & & & \\
\end{longtable}

\emph{Table 10.58: TEMPI Keyword Description}

See also the RTEMP keyword in the PROPS section and the RTEMPVD keyword
in the SOLUTION section for alternative ways to initialize the model's
initial temperature.

\subsubsection{Example}\label{example-486}

\begin{verbatim}
--
--       DEFINE GRID BLOCK TEMPERATURE FOR ALL CELLS
–        (BASED ON NX x NY x NZ = 300)
--
TEMPI
         100*212.0   100*215.0   100*220.0                                     /
\end{verbatim}

The above example defines the initial temperature to be 212.0, 215.0,
and 220.0 oF for the first, second and third layers in the model for all
300 cells, as defined by the DIMENS keyword in the RUNSPEC section.

\subsection{TEMPVD -- Define the Initial Reservoir Temperature versus
Depth
Tables}\label{tempvd-define-the-initial-reservoir-temperature-versus-depth-tables}

\begin{longtable}[]{@{}
  >{\raggedright\arraybackslash}p{(\linewidth - 14\tabcolsep) * \real{0.1250}}
  >{\raggedright\arraybackslash}p{(\linewidth - 14\tabcolsep) * \real{0.1250}}
  >{\raggedright\arraybackslash}p{(\linewidth - 14\tabcolsep) * \real{0.1250}}
  >{\raggedright\arraybackslash}p{(\linewidth - 14\tabcolsep) * \real{0.1250}}
  >{\raggedright\arraybackslash}p{(\linewidth - 14\tabcolsep) * \real{0.1250}}
  >{\raggedright\arraybackslash}p{(\linewidth - 14\tabcolsep) * \real{0.1250}}
  >{\raggedright\arraybackslash}p{(\linewidth - 14\tabcolsep) * \real{0.1250}}
  >{\raggedright\arraybackslash}p{(\linewidth - 14\tabcolsep) * \real{0.1250}}@{}}
\toprule\noalign{}
\begin{minipage}[b]{\linewidth}\raggedright
RUNSPEC
\end{minipage} & \begin{minipage}[b]{\linewidth}\raggedright
GRID
\end{minipage} & \begin{minipage}[b]{\linewidth}\raggedright
EDIT
\end{minipage} & \begin{minipage}[b]{\linewidth}\raggedright
PROPS
\end{minipage} & \begin{minipage}[b]{\linewidth}\raggedright
REGIONS
\end{minipage} & \begin{minipage}[b]{\linewidth}\raggedright
SOLUTION
\end{minipage} & \begin{minipage}[b]{\linewidth}\raggedright
SUMMARY
\end{minipage} & \begin{minipage}[b]{\linewidth}\raggedright
SCHEDULE
\end{minipage} \\
\midrule\noalign{}
\endhead
\bottomrule\noalign{}
\endlastfoot
\end{longtable}

\subsubsection{Description}\label{description-841}

This keyword defines the initial reservoir temperature versus depth
tables for each equilibration region. Note that the TEMPVD keyword is an
alias for RTEMPVD, and that both keywords are supported by OPM Flow, in
both the PROPS and SOLUTION sections, but are treated as being mutually
exclusive.

The initial reservoir temperature must be defined when OPM Flow's
thermal option has been activated by the THERMAL keyword in the RUNSPEC
section. Note this is different to the commercial simulator that uses
the TEMP keyword in the RUNSPEC section to activate the black-oil
temperature model, and the THERMAL keyword to activate the compositional
thermal model.

The initial reservoir temperature should be defined when OPM Flow's CO2
or H2 storage option has been activated by the CO2STORE or
\hyperref[REF_HEADING_KEYWORD_H2STORE]{H2STORE} keyword in the RUNSPEC
section.

\begin{longtable}[]{@{}
  >{\raggedright\arraybackslash}p{(\linewidth - 6\tabcolsep) * \real{0.2500}}
  >{\raggedright\arraybackslash}p{(\linewidth - 6\tabcolsep) * \real{0.2500}}
  >{\raggedright\arraybackslash}p{(\linewidth - 6\tabcolsep) * \real{0.2500}}
  >{\raggedright\arraybackslash}p{(\linewidth - 6\tabcolsep) * \real{0.2500}}@{}}
\toprule\noalign{}
\begin{minipage}[b]{\linewidth}\raggedright
No.
\end{minipage} & \begin{minipage}[b]{\linewidth}\raggedright
Name
\end{minipage} & \begin{minipage}[b]{\linewidth}\raggedright
Description
\end{minipage} & \begin{minipage}[b]{\linewidth}\raggedright
Default
\end{minipage} \\
\midrule\noalign{}
\endhead
\bottomrule\noalign{}
\endlastfoot
Field & Metric & Laboratory & \\
1 & DEPTH & A columnar vector of real monotonically increasing down the
column values that defines the depth for corresponding reservoir
temperature parameter RTEMP. & None \\
feet & m & cm & \\
2 & RTEMP & A columnar vector of real monotonically increasing down the
column values that defines the corresponding reservoir temperature for
the given depth. & None \\
oF & oC & oC & \\
Notes: & & & \\
\end{longtable}

\emph{Table 10.59: TEMPVD Keyword Description}

See also the RTEMP keyword in the PROPS section for an alternative way
to define a uniform initial reservoir temperature.

\subsubsection{Example}\label{example-487}

\begin{verbatim}
--
--       INITIAL RESERVOIR TEMPERATURE VERSUS DEPTH TABLE
--
TEMPVD
--       DEPTH    TEMPERATURE
--       FEET     DEG F
--       ------   ----------
         1000.0     90.000
         2000.0    100.000
         3000.0    130.000
         4000.0    160.000                                 / TABLE N0. 01
--       ------   ----------
         1000.0     90.000
         2000.0    100.000
         3000.0    130.000
         4000.0    160.000                                 / TABLE N0. 02
--       ------   ----------
         1000.0     90.000
         2000.0    100.000
         3000.0    130.000
         4000.0    160.000                                 / TABLE N0. 03
\end{verbatim}

The above example defines three identical reservoir depth versus
temperature tables for the three NTEQUL regions defined on the EQLDIMS
keyword in the RUNSPEC section.

\subsection{THPRES -- Define Equilibration Region Threshold
Pressures}\label{thpres-define-equilibration-region-threshold-pressures}

\begin{longtable}[]{@{}
  >{\raggedright\arraybackslash}p{(\linewidth - 14\tabcolsep) * \real{0.1250}}
  >{\raggedright\arraybackslash}p{(\linewidth - 14\tabcolsep) * \real{0.1250}}
  >{\raggedright\arraybackslash}p{(\linewidth - 14\tabcolsep) * \real{0.1250}}
  >{\raggedright\arraybackslash}p{(\linewidth - 14\tabcolsep) * \real{0.1250}}
  >{\raggedright\arraybackslash}p{(\linewidth - 14\tabcolsep) * \real{0.1250}}
  >{\raggedright\arraybackslash}p{(\linewidth - 14\tabcolsep) * \real{0.1250}}
  >{\raggedright\arraybackslash}p{(\linewidth - 14\tabcolsep) * \real{0.1250}}
  >{\raggedright\arraybackslash}p{(\linewidth - 14\tabcolsep) * \real{0.1250}}@{}}
\toprule\noalign{}
\begin{minipage}[b]{\linewidth}\raggedright
RUNSPEC
\end{minipage} & \begin{minipage}[b]{\linewidth}\raggedright
GRID
\end{minipage} & \begin{minipage}[b]{\linewidth}\raggedright
EDIT
\end{minipage} & \begin{minipage}[b]{\linewidth}\raggedright
PROPS
\end{minipage} & \begin{minipage}[b]{\linewidth}\raggedright
REGIONS
\end{minipage} & \begin{minipage}[b]{\linewidth}\raggedright
SOLUTION
\end{minipage} & \begin{minipage}[b]{\linewidth}\raggedright
SUMMARY
\end{minipage} & \begin{minipage}[b]{\linewidth}\raggedright
SCHEDULE
\end{minipage} \\
\midrule\noalign{}
\endhead
\bottomrule\noalign{}
\endlastfoot
\end{longtable}

\subsubsection{Description}\label{description-842}

The THPRES defines the threshold pressure between various equilibration
regions that have been defined by the EQLNUM keyword in the REGIONS
section. The threshold pressure defines the potential difference between
two regions which must be exceeded before flow can occur between the two
regions. Once flow occurs the potential between the two regions is
reduced by the threshold pressure.

This option must be activated by THPRES variable on EQLOPTS keyword in
the RUNSPEC section in order to utilize this feature. Note that the
irreversible option, as defined by IRREVER variable on EQLOPTS keyword
in the RUNSPEC section, is not supported.

\begin{longtable}[]{@{}
  >{\raggedright\arraybackslash}p{(\linewidth - 6\tabcolsep) * \real{0.2500}}
  >{\raggedright\arraybackslash}p{(\linewidth - 6\tabcolsep) * \real{0.2500}}
  >{\raggedright\arraybackslash}p{(\linewidth - 6\tabcolsep) * \real{0.2500}}
  >{\raggedright\arraybackslash}p{(\linewidth - 6\tabcolsep) * \real{0.2500}}@{}}
\toprule\noalign{}
\begin{minipage}[b]{\linewidth}\raggedright
No.
\end{minipage} & \begin{minipage}[b]{\linewidth}\raggedright
Name
\end{minipage} & \begin{minipage}[b]{\linewidth}\raggedright
Description
\end{minipage} & \begin{minipage}[b]{\linewidth}\raggedright
Default
\end{minipage} \\
\midrule\noalign{}
\endhead
\bottomrule\noalign{}
\endlastfoot
Field & Metric & Laboratory & \\
1 & EQLNUM1 & EQLNUM1 is an a positive integer that is greater or equal
to one and less than or equal to NTEQUL on the EQLDIMS keyword in the
RUNSPEC section, that defines the ``from'' equilibration region number.
& None \\
dimensionless & dimensionless & dimensionless & \\
2 & EQLNUM2 & EQLNUM1 is an a positive integer that is greater or equal
to one and less than or equal to NTEQUL on the EQLDIMS keyword in the
RUNSPEC section, that defines the ``to'' equilibration region number. &
None \\
dimensionless & dimensionless & dimensionless & \\
3 & THPRES & THPRES defines the threshold pressure from EQLNUM1 to
EQLNUM2 and from EQLNUM2 to EQLNUM1. The default value of 1* sets the
threshold pressure to a value that initially prevents flow between the
two equilibration regions. Any subsequent production or injection in
either of the two equilibration regions will therefore result in flow
between the two regions. Thus, this default initially isolates the two
equilibration regions. If a equilibration region number pair has not
been explicitly defined by this keyword the THPRES is set to zero, for
no threshold pressure. & 1* \\
psia & barsa & atma & \\
Notes: & & & \\
\end{longtable}

\emph{Table 10.60: THPRES Keyword Description}

See also the MULTREGT keyword in the GRID section that uses the
transmissibility between the MULTNUM, FLUXNUM or OPERNUM region arrays
to control the flow between various regions within the model.

\begin{tcolorbox}[enhanced jigsaw, title=\textcolor{quarto-callout-note-color}{\faInfo}\hspace{0.5em}{Note}, coltitle=black, breakable, toprule=.15mm, colbacktitle=quarto-callout-note-color!10!white, leftrule=.75mm, opacitybacktitle=0.6, colback=white, bottomtitle=1mm, toptitle=1mm, titlerule=0mm, left=2mm, colframe=quarto-callout-note-color-frame, arc=.35mm, rightrule=.15mm, bottomrule=.15mm, opacityback=0]

Care should be taken that cells in different EQLNUM regions are not in
communication, as this will result in an unstable initial equilibration.

\end{tcolorbox}

\subsubsection{Examples}\label{examples-60}

Given NTEQUL is equal to six on the EQLDIMS keyword in the RUNSPEC
section,

\begin{verbatim}
--
--       EQLNUM  EQLNUM  THPRES
--       FROM    TO      VALUE
THPRES
         1       2       0.588031                          / REGN 1 TO REGN 2
         2       1       0.588031                          / REGN 2 TO REGN 1
         1       3       0.787619                          / REGN 1 TO REGN 3
         3       1       0.787619                          / REGN 3 TO REGN 1
         1       4       7.000830                          / REGN 1 TO REGN 4
         4       1       7.000830                          / REGN 4 TO REGN 1
/
\end{verbatim}

The above example defines the threshold pressures between equilibration
regions one and two, one and three and one and four. As the threshold
pressures between regions one and five and one and six (as well as other
combinations), have not been explicitly set in the example, the
threshold pressures for these combinations are set to zero.

However, as the irreversible option, as defined by IRREVER variable on
EQLOPTS keyword in the RUNSPEC section, is not supported, then example
can be simplified to:

\begin{verbatim}
--
--       EQLNUM  EQLNUM  THPRES
--       FROM    TO      VALUE
THPRES
         1       2       0.588031                          / REGN 1 AND REGN 2
         1       3       0.787619                          / REGN 1 AND REGN 3
         1       4       7.000830                          / REGN 1 AND REGN 4
/
\end{verbatim}

Again, as the threshold pressures between regions one and five and one
and six (as well as other combinations), have not been explicitly set in
the example, the threshold pressures for these combinations are set to
zero.

\subsection{TVDP -- Define the Initial Equilibration Tracer Saturation
versus Depth
Functions}\label{tvdp-define-the-initial-equilibration-tracer-saturation-versus-depth-functions}

\begin{longtable}[]{@{}
  >{\raggedright\arraybackslash}p{(\linewidth - 14\tabcolsep) * \real{0.1250}}
  >{\raggedright\arraybackslash}p{(\linewidth - 14\tabcolsep) * \real{0.1250}}
  >{\raggedright\arraybackslash}p{(\linewidth - 14\tabcolsep) * \real{0.1250}}
  >{\raggedright\arraybackslash}p{(\linewidth - 14\tabcolsep) * \real{0.1250}}
  >{\raggedright\arraybackslash}p{(\linewidth - 14\tabcolsep) * \real{0.1250}}
  >{\raggedright\arraybackslash}p{(\linewidth - 14\tabcolsep) * \real{0.1250}}
  >{\raggedright\arraybackslash}p{(\linewidth - 14\tabcolsep) * \real{0.1250}}
  >{\raggedright\arraybackslash}p{(\linewidth - 14\tabcolsep) * \real{0.1250}}@{}}
\toprule\noalign{}
\begin{minipage}[b]{\linewidth}\raggedright
RUNSPEC
\end{minipage} & \begin{minipage}[b]{\linewidth}\raggedright
GRID
\end{minipage} & \begin{minipage}[b]{\linewidth}\raggedright
EDIT
\end{minipage} & \begin{minipage}[b]{\linewidth}\raggedright
PROPS
\end{minipage} & \begin{minipage}[b]{\linewidth}\raggedright
REGIONS
\end{minipage} & \begin{minipage}[b]{\linewidth}\raggedright
SOLUTION
\end{minipage} & \begin{minipage}[b]{\linewidth}\raggedright
SUMMARY
\end{minipage} & \begin{minipage}[b]{\linewidth}\raggedright
SCHEDULE
\end{minipage} \\
\midrule\noalign{}
\endhead
\bottomrule\noalign{}
\endlastfoot
\end{longtable}

\subsubsection{Description}\label{description-843}

This keyword defines the initial equilibration tracer concentration
versus depth functions for each grid cell in the model, for when the
Tracer option has been enabled by the TRACERS keyword in the RUNSPEC
section. The maximum number of tracers for each phase are declared on
the TRACERS keyword in the RUNSPEC section. Unlike other keywords, the
TVDP keyword must be concatenated with the name of the tracer declared
by the TRACER keyword in the PROPS section as outlined in Table 10.61.

\begin{longtable}[]{@{}
  >{\raggedright\arraybackslash}p{(\linewidth - 6\tabcolsep) * \real{0.2500}}
  >{\raggedright\arraybackslash}p{(\linewidth - 6\tabcolsep) * \real{0.2500}}
  >{\raggedright\arraybackslash}p{(\linewidth - 6\tabcolsep) * \real{0.2500}}
  >{\raggedright\arraybackslash}p{(\linewidth - 6\tabcolsep) * \real{0.2500}}@{}}
\toprule\noalign{}
\begin{minipage}[b]{\linewidth}\raggedright
No.
\end{minipage} & \begin{minipage}[b]{\linewidth}\raggedright
Name
\end{minipage} & \begin{minipage}[b]{\linewidth}\raggedright
Description
\end{minipage} & \begin{minipage}[b]{\linewidth}\raggedright
Default
\end{minipage} \\
\midrule\noalign{}
\endhead
\bottomrule\noalign{}
\endlastfoot
1 & TVDP & A four letter character string equal to TVDP that is the root
keyword name for this data set function. & None \\
2 & STATE & A one letter character string defining the tracer's state,
which is concatenated to TVDP to give the partial name of the keyword.
STATE should be set to F (``free'') or S (``solution''), with the latter
related to a tracer that is dissolved in an associated phase. If the
tracer defined by NAME is associated with the OIL, GAS or WATER phases
then STATE should be set to F. If the tracer is associated with the
DISGAS or VAPOIL phases then STATE should be set to S. & None \\
3 & NAME & A three letter character string defining the tracer's name as
defined by the TRACER keyword, which is concatenated to TVDP and STATE
to given the full name of the keyword. For example, if the TRACER
keyword has been used to define a tracer named SEA in the WAT phase,
then the full keyword name would be TVDPFSEA. Whereas, if a tracer was
defined as IGS for the GAS phase then it could be TVDPFIGS and/or
TVDPSIGS depending if one was wishing to track the ``free'' or
``solution'' gas phase, or both. Note it is best to avoid names
beginning with the letters F, S, and T as these names may create naming
issues in post-processing software. & None \\
\end{longtable}

\emph{Table 10.61: TVDP Keyword Name Format}

Following the declaration of the full keyword name, TVDPSTATENAME, the
keyword is followed by the data as outlined below.

\begin{longtable}[]{@{}
  >{\raggedright\arraybackslash}p{(\linewidth - 6\tabcolsep) * \real{0.2500}}
  >{\raggedright\arraybackslash}p{(\linewidth - 6\tabcolsep) * \real{0.2500}}
  >{\raggedright\arraybackslash}p{(\linewidth - 6\tabcolsep) * \real{0.2500}}
  >{\raggedright\arraybackslash}p{(\linewidth - 6\tabcolsep) * \real{0.2500}}@{}}
\toprule\noalign{}
\begin{minipage}[b]{\linewidth}\raggedright
No.
\end{minipage} & \begin{minipage}[b]{\linewidth}\raggedright
Name
\end{minipage} & \begin{minipage}[b]{\linewidth}\raggedright
Description
\end{minipage} & \begin{minipage}[b]{\linewidth}\raggedright
Default
\end{minipage} \\
\midrule\noalign{}
\endhead
\bottomrule\noalign{}
\endlastfoot
Field & Metric & Laboratory & \\
1 & DEPTH & A columnar vector of real monotonically increasing down the
column values that defines the depth values for the corresponding
initial tracer saturations, TVDP & None \\
feet & m & cm & \\
2 & TVDPVAL & A columnar vector of real values, greater than or equal to
zero, that defines the initial tracer concentration values at the
corresponding DEPTH. If tracer units have been defined by the UNITS
parameter on the TRACER keyword in the PROPS section, then the units of
TVDPVAL are the ratio of UNITS divided by the TVDPVAL units given below.
For example, if UNITS was defined as kg, then for metric units TVDPVAL
units would be kg/sm3. & None \\
Liquid: stb Gas: Mscf & Liquid: sm3 Gas: sm3 & Liquid: scc Gas: scc & \\
Notes: & & & \\
\end{longtable}

\emph{Table 10.62: TVDP Keyword Description}

See also the TBLK keyword in the SOLUTION section that also sets the
initial tracer concentration for each grid block.

\subsubsection{Example}\label{example-488}

This example is taken from Norne model, in which there are seven tracers
related to the water phase.

\begin{verbatim}
-- ==============================================================================
--
-- RUNSPEC SECTION
--
-- ==============================================================================
RUNSPEC
--
--       MAX     MAX     RSVD    TVDP    TVDP
--       EQLNUM  DEPTH   NODES   TABLE   NODES
EQLDIMS
         9       1*      20      7       1*                                    /
-- ==============================================================================
--
-- PROPS SECTION
--
-- ==============================================================================
PROPS
--
--       DEFINE TRACER NAMES
--
--       TRACER   TRACER
--       NAME     PHASE
--       ------   ------
TRACER
         SEA      WAT                                      /
         HTO      WAT                                      /
         S36      WAT                                      /
         2FB      WAT                                      /
         4FB      WAT                                      /
         DFB      WAT                                      /
         TFB      WAT                                      /
/

-- ==============================================================================
--
-- SOLUTION SECTION
--
-- ==============================================================================
SOLUTION
--
--       INITIAL EQUILIBRATION TRACER SATURATION VERSUS DEPTH
--
--       DEPTH    TRACER
--                CONCENT
--       ------   --------
TVDPFSEA
         1000.0   0.0
         5000.0   0.0                       / TRACER FSEA CONCENTRATION VS DEPTH

TVDPFHTO
         1000.0   0.0
         5000.0   0.0                       / TRACER FHTO CONCENTRATION VS DEPTH

TVDPFS36
         1000.0   0.0
         5000.0   0.0                       / TRACER FS36 CONCENTRATION VS DEPTH

TVDPF2FB
         1000.0   0.0
         5000.0   0.0                       / TRACER F2FB CONCENTRATION VS DEPTH

TVDPF4FB
         1000.0   0.0
         5000.0   0.0                       / TRACER F4FB CONCENTRATION VS DEPTH

TVDPFDFB
         1000.0   0.0
         5000.0   0.0                       / TRACER FDFB CONCENTRATION VS DEPTH

TVDPFTFB
         1000.0   0.0
         5000.0   0.0                       / TRACER FTFB CONCENTRATION VS DEPTH
\end{verbatim}

Here we first define the number of tracers in the model via the EQLDIMS
keyword in the RUNSPEC section, then the actual tracers themselves in
the PROPS section using the TRACER keyword, and finally the initial
tracer concentrations are all set to zero via the TVDP keyword in the
PROPS section.

\subsection{VAPPARS -- Oil Vaporization
Parameters}\label{vappars-oil-vaporization-parameters}

\begin{longtable}[]{@{}
  >{\raggedright\arraybackslash}p{(\linewidth - 14\tabcolsep) * \real{0.1250}}
  >{\raggedright\arraybackslash}p{(\linewidth - 14\tabcolsep) * \real{0.1250}}
  >{\raggedright\arraybackslash}p{(\linewidth - 14\tabcolsep) * \real{0.1250}}
  >{\raggedright\arraybackslash}p{(\linewidth - 14\tabcolsep) * \real{0.1250}}
  >{\raggedright\arraybackslash}p{(\linewidth - 14\tabcolsep) * \real{0.1250}}
  >{\raggedright\arraybackslash}p{(\linewidth - 14\tabcolsep) * \real{0.1250}}
  >{\raggedright\arraybackslash}p{(\linewidth - 14\tabcolsep) * \real{0.1250}}
  >{\raggedright\arraybackslash}p{(\linewidth - 14\tabcolsep) * \real{0.1250}}@{}}
\toprule\noalign{}
\begin{minipage}[b]{\linewidth}\raggedright
RUNSPEC
\end{minipage} & \begin{minipage}[b]{\linewidth}\raggedright
GRID
\end{minipage} & \begin{minipage}[b]{\linewidth}\raggedright
EDIT
\end{minipage} & \begin{minipage}[b]{\linewidth}\raggedright
PROPS
\end{minipage} & \begin{minipage}[b]{\linewidth}\raggedright
REGIONS
\end{minipage} & \begin{minipage}[b]{\linewidth}\raggedright
SOLUTION
\end{minipage} & \begin{minipage}[b]{\linewidth}\raggedright
SUMMARY
\end{minipage} & \begin{minipage}[b]{\linewidth}\raggedright
SCHEDULE
\end{minipage} \\
\midrule\noalign{}
\endhead
\bottomrule\noalign{}
\endlastfoot
\end{longtable}

\subsubsection{Description}\label{description-844}

VAPPARS defines the rate of oil vaporization in the presence of
undersaturated gas and the rate at which the remaining oil gets
``heavier'' via the reduction in the solution gas-oil ratio (``Rs'').
This keyword should only be used if the OIL, GAS, DISGAS and VAPOIL
keywords in the RUNSPEC section have been invoked to allow oil, gas,
dissolved gas and vaporized oil to be present in the model.

\begin{longtable}[]{@{}
  >{\raggedright\arraybackslash}p{(\linewidth - 6\tabcolsep) * \real{0.2500}}
  >{\raggedright\arraybackslash}p{(\linewidth - 6\tabcolsep) * \real{0.2500}}
  >{\raggedright\arraybackslash}p{(\linewidth - 6\tabcolsep) * \real{0.2500}}
  >{\raggedright\arraybackslash}p{(\linewidth - 6\tabcolsep) * \real{0.2500}}@{}}
\toprule\noalign{}
\begin{minipage}[b]{\linewidth}\raggedright
No.
\end{minipage} & \begin{minipage}[b]{\linewidth}\raggedright
Name
\end{minipage} & \begin{minipage}[b]{\linewidth}\raggedright
Description
\end{minipage} & \begin{minipage}[b]{\linewidth}\raggedright
Default
\end{minipage} \\
\midrule\noalign{}
\endhead
\bottomrule\noalign{}
\endlastfoot
Field & Metric & Laboratory & \\
1 & VAPPAR1 & VAPPAR1 is a real positive dimensionless number that
defines the rate at which oil vaporizes into the available
undersaturated gas in a grid block. The default value of zero invokes
the standard black-oil formulation in which all oil vaporizes into the
available undersaturated phase in a grid cell. Increasing this parameter
decreases the rate of vaporization. Typical values for VAPPAR1 range
from zero to five. & 0 \\
dimensionless & dimensionless & dimensionless & \\
2 & VAPPAR2 & VAPPAR2 is a real positive dimensionless number that
defines the rate at which the Rs of the remaining oil in a grid cell
decreases The default value of zero invokes the standard black-oil
formulation in which the remaining oil's Rs does not change as the oil
vaporizes into the available undersaturated gas in a grid cell.
Increasing this parameter increases the difference between the remaining
oil and the vaporized oil Rs values. Typical values for VAPPAR2 are less
than one. & 0 \\
dimensionless & dimensionless & dimensionless & \\
Notes: & & & \\
\end{longtable}

\emph{Table 10.63: VAPPARS Keyword Description}

Note this keyword is normally used in history matching field performance
to control the availability of the vaporized oil phase.

See also the DRSDT and DRVDT keywords in the SCHEDULE section that
control the rate at which the solution gas-oil ratio and the vaporized
oil-gas ratio increase within a grid block, respectively.

\subsubsection{Examples}\label{examples-61}

The first example sets the black-oil default parameters

\begin{verbatim}
--
--       OIL VAPORIZATION PARAMETERS
--
--       OIL-VAP   RS-INCS
--       VAPPAR1   VAPPAR2
VAPPARS
         0         0                                       /
\end{verbatim}

And the second example decreases the rate at which the oil vaporizes
into the available undersaturated gas and increases the difference
between the grid block oil saturation Rs and the vaporized oil Rs within
a grid cell.

\begin{verbatim}
--
--       OIL VAPORIZATION PARAMETERS
--
--       OIL-VAP   RS-INCS
--       VAPPAR1   VAPPAR2
VAPPARS
         1.5       0.150                                   /
\end{verbatim}

Again, the keyword is normally used in history matching field
performance to control the availability of the vaporized oil phase.

\subsection{VISDATES -- Define External Reservoir Geo-Mechanics VISAGE
Stress
Dates}\label{visdates-define-external-reservoir-geo-mechanics-visage-stress-dates}

\begin{longtable}[]{@{}
  >{\raggedright\arraybackslash}p{(\linewidth - 14\tabcolsep) * \real{0.1250}}
  >{\raggedright\arraybackslash}p{(\linewidth - 14\tabcolsep) * \real{0.1250}}
  >{\raggedright\arraybackslash}p{(\linewidth - 14\tabcolsep) * \real{0.1250}}
  >{\raggedright\arraybackslash}p{(\linewidth - 14\tabcolsep) * \real{0.1250}}
  >{\raggedright\arraybackslash}p{(\linewidth - 14\tabcolsep) * \real{0.1250}}
  >{\raggedright\arraybackslash}p{(\linewidth - 14\tabcolsep) * \real{0.1250}}
  >{\raggedright\arraybackslash}p{(\linewidth - 14\tabcolsep) * \real{0.1250}}
  >{\raggedright\arraybackslash}p{(\linewidth - 14\tabcolsep) * \real{0.1250}}@{}}
\toprule\noalign{}
\begin{minipage}[b]{\linewidth}\raggedright
RUNSPEC
\end{minipage} & \begin{minipage}[b]{\linewidth}\raggedright
GRID
\end{minipage} & \begin{minipage}[b]{\linewidth}\raggedright
EDIT
\end{minipage} & \begin{minipage}[b]{\linewidth}\raggedright
PROPS
\end{minipage} & \begin{minipage}[b]{\linewidth}\raggedright
REGIONS
\end{minipage} & \begin{minipage}[b]{\linewidth}\raggedright
SOLUTION
\end{minipage} & \begin{minipage}[b]{\linewidth}\raggedright
SUMMARY
\end{minipage} & \begin{minipage}[b]{\linewidth}\raggedright
SCHEDULE
\end{minipage} \\
\midrule\noalign{}
\endhead
\bottomrule\noalign{}
\endlastfoot
\end{longtable}

\subsubsection{Description}\label{description-845}

The VISDATES keyword defines External Reservoir Geo-Mechanics VISAGE
option stress dates. The keyword should not be used in input decks as
the associated data is generated by an external program.

This keyword is not supported by OPM Flow but would change the results
if supported so the simulation will be stopped.

\subsection{VISOPTS -- Define External Reservoir Geo-Mechanics VISAGE
Options}\label{visopts-define-external-reservoir-geo-mechanics-visage-options}

\begin{longtable}[]{@{}
  >{\raggedright\arraybackslash}p{(\linewidth - 14\tabcolsep) * \real{0.1250}}
  >{\raggedright\arraybackslash}p{(\linewidth - 14\tabcolsep) * \real{0.1250}}
  >{\raggedright\arraybackslash}p{(\linewidth - 14\tabcolsep) * \real{0.1250}}
  >{\raggedright\arraybackslash}p{(\linewidth - 14\tabcolsep) * \real{0.1250}}
  >{\raggedright\arraybackslash}p{(\linewidth - 14\tabcolsep) * \real{0.1250}}
  >{\raggedright\arraybackslash}p{(\linewidth - 14\tabcolsep) * \real{0.1250}}
  >{\raggedright\arraybackslash}p{(\linewidth - 14\tabcolsep) * \real{0.1250}}
  >{\raggedright\arraybackslash}p{(\linewidth - 14\tabcolsep) * \real{0.1250}}@{}}
\toprule\noalign{}
\begin{minipage}[b]{\linewidth}\raggedright
RUNSPEC
\end{minipage} & \begin{minipage}[b]{\linewidth}\raggedright
GRID
\end{minipage} & \begin{minipage}[b]{\linewidth}\raggedright
EDIT
\end{minipage} & \begin{minipage}[b]{\linewidth}\raggedright
PROPS
\end{minipage} & \begin{minipage}[b]{\linewidth}\raggedright
REGIONS
\end{minipage} & \begin{minipage}[b]{\linewidth}\raggedright
SOLUTION
\end{minipage} & \begin{minipage}[b]{\linewidth}\raggedright
SUMMARY
\end{minipage} & \begin{minipage}[b]{\linewidth}\raggedright
SCHEDULE
\end{minipage} \\
\midrule\noalign{}
\endhead
\bottomrule\noalign{}
\endlastfoot
\end{longtable}

\subsubsection{Description}\label{description-846}

The VISDATES keyword defines External Reservoir Geo-Mechanics VISAGE
option modeling options. The keyword should not be used in input decks
as the associated data is generated by an external program.

This keyword is not supported by OPM Flow but would change the results
if supported so the simulation will be stopped.

\subsection{WARN -- Activate Warning
Messages}\label{warn-activate-warning-messages-6}

Turns on warning messages to be printed to the print file (*.PRT); note
that this keyword is activated by default and can subsequently be
switched off by the NOWARN activation keyword. The warning messages may
be turned on and off using keywords WARN and NOWARN. OPM Flow always
prints error messages.

See WARN -- Activate Warning Messages in the GLOBAL section for a full
description.

\bookmarksetup{startatroot}

\chapter{SUMMARY SECTION}\label{summary-section}

\section{Introduction}\label{introduction-9}

The SUMMARY section defines the variables to be written to the summary
files that are used to generate line graphs of properties such as oil
flow rate versus time, grid block pressure versus time, etc. Unlike the
other sections, the SUMMARY section has basically two types of keywords:
keywords that request a variable to be written out to the summary file
as described in section 11.2 Data Requirements, and the conventional
type of keyword described in section 11.3 Keyword Definitions that
perform an action. There are literally hundreds of keywords associated
with the former indicating the variety of data that can be written to
the summary file. Whereas, for the latter there are less than twenty
keywords.

Standard summary file output consists of three files as outlined in
Table 11.1, with the raw requested data stored in SUMMARY index and
SUMMARY data files, the SUMMARY RSM file is an ASCII file that is
generated from the other two files by the simulator at the end of the
run. There is also the option to output an OPM Flow specific Enhanced
SUMMARY output file.

\begin{longtable}[]{@{}
  >{\raggedright\arraybackslash}p{(\linewidth - 6\tabcolsep) * \real{0.2500}}
  >{\raggedright\arraybackslash}p{(\linewidth - 6\tabcolsep) * \real{0.2500}}
  >{\raggedright\arraybackslash}p{(\linewidth - 6\tabcolsep) * \real{0.2500}}
  >{\raggedright\arraybackslash}p{(\linewidth - 6\tabcolsep) * \real{0.2500}}@{}}
\toprule\noalign{}
\begin{minipage}[b]{\linewidth}\raggedright
File Type
\end{minipage} & \begin{minipage}[b]{\linewidth}\raggedright
Data Type
\end{minipage} & \begin{minipage}[b]{\linewidth}\raggedright
Description
\end{minipage} & \begin{minipage}[b]{\linewidth}\raggedright
OPM Flow Status
\end{minipage} \\
\midrule\noalign{}
\endhead
\bottomrule\noalign{}
\endlastfoot
ESMRY & Dynamic Vector Data & This is an OPM FLOW specific Enhanced
SUMMARY (\emph{.ESMRY) output file format that is optimized for fast
loading of selected vectors by post-processing applications. The }.ESMRY
file is re-written for every time step and does not have a performance
impact on the simulation. Load times are in the order of 30 times faster
compared with fetching selected vectors using the *.UNSMRY file format
in post processing software. The file type is compatible with OPM
ResInsight and the option is activated via the --enable-esmry=true
command line option (true is the default value). Note that the FMTOUT,
MULTOUT, and UNIFOUT keywords have no impact on this file type. & OPM
Flow Fully Supported. \\
RSM & Output File & The RSM file contains the variables requested to be
written to the SUMMARY file via the keywords described in the 11.2 Data
Requirements, in a tabulated ASCII output format. The format of the file
enables the data to be easily loaded into the LibreOffice Calc program
for further processing, as each vector represents one column. The report
is written at the end of the simulation run by parsing the SUMMARY Index
and SUMMARY Data files. & Fully Supported \\
SUMMARY Index & Dynamic Vector Index & The SUMMARY index file type
specifies and defines the format and data type written to the SUMMARY
Data file. & Fully Supported \\
SUMMARY Data & Dynamic Vector Data & The SUMMARY data file contains the
variables requested to be written to the file via the keywords described
in the 11.2 Data Requirements. The data consists of vectors that are
used to generate line graphs of properties such as oil flow rate versus
time, grid block pressure versus time, etc. The properties to be stored
on the SUMMARY file are written to the summary file at the end of each
successful time step. The data can be used to compare actual production
data with the simulation derived results in post processing software
like OPM ResInsight. & Fully Supported \\
Notes: & & & \\
\end{longtable}

\emph{Table 11.1: OPM Flow Summary Output File Summary}

As mentioned in Table 11.1 all files can be written out in either ASCII
or binary formats and in addition the SUMMARY file can be loaded or
written out in either unified or non-unified formats. The file type
(ASCII or binary) and file structure format (unified or non-unified
formats) is set via a set of keywords in the RUNSPEC section, as
described in Table 11.2 for easy of reference.

\begin{longtable}[]{@{}
  >{\raggedright\arraybackslash}p{(\linewidth - 6\tabcolsep) * \real{0.2500}}
  >{\raggedright\arraybackslash}p{(\linewidth - 6\tabcolsep) * \real{0.2500}}
  >{\raggedright\arraybackslash}p{(\linewidth - 6\tabcolsep) * \real{0.2500}}
  >{\raggedright\arraybackslash}p{(\linewidth - 6\tabcolsep) * \real{0.2500}}@{}}
\toprule\noalign{}
\begin{minipage}[b]{\linewidth}\raggedright
Process
\end{minipage} & \begin{minipage}[b]{\linewidth}\raggedright
RUNSPEC Keyword
\end{minipage} & \begin{minipage}[b]{\linewidth}\raggedright
Description
\end{minipage} & \begin{minipage}[b]{\linewidth}\raggedright
Files
\end{minipage} \\
\midrule\noalign{}
\endhead
\bottomrule\noalign{}
\endlastfoot
Input & FMTIN & The keyword defines the input files to be formatted as
ASCII, i.e., text files, as opposed to binary files. & \emph{.FSMSPEC
}.FUNSMRY \\
MULTIN & The keyword defines the input files to be non-unified multiple
files, as opposed to unified files. & \emph{.SMSPEC }.S0001 & \\
UNIFIN & UNIFIN defines the input files to be unified files, as opposed
to non-unified multiple files. & \emph{.SMSPEC }.UNSMRY & \\
Output & FMTOUT & The keyword sets all output files to be formatted as
ASCII, i.e., text files, as opposed to binary files. & \emph{.FSMSPEC
}.FUNSMRY \\
MULTOUT & MULTOUT defines the output files to be non-unified multiple
files, as opposed to unified files. & \emph{.SMSPEC }.S0001 & \\
UNIFOUT & UNIFOUT defines the output files to be unified files, as
opposed to non-unified multiple files. & \emph{.SMSPEC }.UNSMRY & \\
\end{longtable}

\emph{Table 11.2: RUNSPEC Input and Output File Format Keywords}

OPM Flow automatically generates the SUMMARY file names based on the
input file name and the output options selected via the keywords in the
RUNSPEC section as summarized in in Table 11.2. For example, starting
OPM Flow using the following command from the terminal:

\begin{verbatim}
flow CASENAME
\end{verbatim}

will start the simulator and run the case specified by CASENAME.DATA.
Here CASENAME is the ``root'' of the filename and DATA is the extension
of the filename. OPM Flow will generate the SUMMARY output files based
on the CASENAME with the extension based on the type (ASCII or binary)
and file structure format (unified or non-unified formats), as outlined
in Table 11.3.

\begin{longtable}[]{@{}
  >{\raggedright\arraybackslash}p{(\linewidth - 6\tabcolsep) * \real{0.2500}}
  >{\raggedright\arraybackslash}p{(\linewidth - 6\tabcolsep) * \real{0.2500}}
  >{\raggedright\arraybackslash}p{(\linewidth - 6\tabcolsep) * \real{0.2500}}
  >{\raggedright\arraybackslash}p{(\linewidth - 6\tabcolsep) * \real{0.2500}}@{}}
\toprule\noalign{}
\begin{minipage}[b]{\linewidth}\raggedright
File Type
\end{minipage} & \begin{minipage}[b]{\linewidth}\raggedright
Unformatted Extension
\end{minipage} & \begin{minipage}[b]{\linewidth}\raggedright
Description
\end{minipage} & \begin{minipage}[b]{\linewidth}\raggedright
Formatted Extension
\end{minipage} \\
\midrule\noalign{}
\endhead
\bottomrule\noalign{}
\endlastfoot
RSM & *.RSM & Output RSM file. & *.RSM \\
SUMMARY Index & *.SMSPEC & The SUMMARY index file for both unified and
non-unified formats. & *.FSMSPEC \\
SUMMARY Data Non-Unified & *.Snnnn & The SUMMARY data files containing
the variables requested to be written to the SUMMARY files. For example,
the SUMMARY non-unified binary files would be: CASENAME.S0001
CASENAME.S0002 CASENAME.S0003 etc. CASENAME.SMSPEC & *.Annnn \\
SUMMARY Data Unified & *.UNSMRY & The SUMMARY data file containing the
variables requested to be written to the SUMMARY file. For example, the
SUMMARY data unified binary files: CASENAME.UNSMRY CASENAME.SMSPEC &
*.FUNSMRY \\
Notes: & & & \\
\end{longtable}

\emph{Table 11.3: OPM Flow File Naming Conventions}

The default behavior is to write out the requested variables at each
time step. As this can lead to large files (or a large number of files
for non-unified output), especially for full field simulation models,
the RPTONLY - Activate the Report Time Steps Only Option for the SUMMARY
File keyword allows one to write out the data only at a report time step
instead.

The SUMMARY section is terminated by the SCHEDULE keyword.

\section{Keyword Definitions}\label{keyword-definitions-15}

\section{Data Requirements}\label{data-requirements-7}

To generate a summary vector, the appropriate keyword must be entered in
the SUMMARY section of the input data file; only variables explicitly
requested will be written to the SUMMARY files. Normally the data is
written to the summary file at the end of each successful time step, but
as mentioned previously, this can be changed by the RPTONLY keyword to
write out the data only at report time steps, thus reducing the overall
size of the SUMMARY files.

The following sections describe the summary variable mnemonic syntax
which defines the type of summary variable object (Aquifer, Field,
Group, Well, etc.), the variable format for a given variable object, and
the variable names for the various variable objects.

\subsection{Summary Variable Mnemonic
Syntax}\label{summary-variable-mnemonic-syntax}

As mentioned earlier, there are literally hundreds of variables that can
be written to the SUMMARY file, which can make the process of requesting
the data rather complex. Fortunately in most cases, but not always, the
variable names follow a four or five letter syntax that defines the
variable mnemonic used to describe the data to be written out. Table
11.2.1 outlines the general syntax employed in deriving the variable
mnemonic for some of the more commonly used summary variables.

\begin{longtable}[]{@{}
  >{\raggedright\arraybackslash}p{(\linewidth - 18\tabcolsep) * \real{0.1000}}
  >{\raggedright\arraybackslash}p{(\linewidth - 18\tabcolsep) * \real{0.1000}}
  >{\raggedright\arraybackslash}p{(\linewidth - 18\tabcolsep) * \real{0.1000}}
  >{\raggedright\arraybackslash}p{(\linewidth - 18\tabcolsep) * \real{0.1000}}
  >{\raggedright\arraybackslash}p{(\linewidth - 18\tabcolsep) * \real{0.1000}}
  >{\raggedright\arraybackslash}p{(\linewidth - 18\tabcolsep) * \real{0.1000}}
  >{\raggedright\arraybackslash}p{(\linewidth - 18\tabcolsep) * \real{0.1000}}
  >{\raggedright\arraybackslash}p{(\linewidth - 18\tabcolsep) * \real{0.1000}}
  >{\raggedright\arraybackslash}p{(\linewidth - 18\tabcolsep) * \real{0.1000}}
  >{\raggedright\arraybackslash}p{(\linewidth - 18\tabcolsep) * \real{0.1000}}@{}}
\toprule\noalign{}
\begin{minipage}[b]{\linewidth}\raggedright
Summary Variable Mnemonic Syntax
\end{minipage} & \begin{minipage}[b]{\linewidth}\raggedright
\end{minipage} & \begin{minipage}[b]{\linewidth}\raggedright
\end{minipage} & \begin{minipage}[b]{\linewidth}\raggedright
\end{minipage} & \begin{minipage}[b]{\linewidth}\raggedright
\end{minipage} & \begin{minipage}[b]{\linewidth}\raggedright
\end{minipage} & \begin{minipage}[b]{\linewidth}\raggedright
\end{minipage} & \begin{minipage}[b]{\linewidth}\raggedright
\end{minipage} & \begin{minipage}[b]{\linewidth}\raggedright
\end{minipage} & \begin{minipage}[b]{\linewidth}\raggedright
\end{minipage} \\
\midrule\noalign{}
\endhead
\bottomrule\noalign{}
\endlastfoot
First Character Object & Second Character Object Type & Third Character
Flow Type & Four Character Value Type & Fifth Character Special & & & &
& \\
A & Aquifers & A & Analytical Aquifer & Q & Influx & R & & & \\
L1 & Analytical Aquifer List & L & Liquid influx & T & & & & & \\
N & Numerical Aquifer & & & & & & & & \\
B & Block (Grid Cell) & O & Oil & P & Production & R & Rate & L &
Liquid \\
F & Field & W & Water & I & Injection & T & Total (cumulative) & H &
History \\
G & Group & G & Gas & F & Flow & & G & Gas & \\
R & Region & V & Volume at reservoir conditions & & & & & & \\
W & Well & L & Liquid & & & & & & \\
C & Well Connection & P2 & Pressure & & & & & & \\
S & Well Segment & T & Tracer & & & & & & \\
Notes: & & & & & & & & & \\
\end{longtable}

\emph{Table 11.2.1: Summary Variable Mnemonic Syntax}

So for example, FOPR would mean Field Oil Production Rate. Also if there
is historical data associated with a variable, then the same variable
mnemonic is used with the suffix H, that is for the historical field oil
production rate the mnemonic would be FOPRH. Table 11.2.2 shows some
examples of variable mnemonics that follow the general syntax.

\begin{longtable}[]{@{}ll@{}}
\toprule\noalign{}
Summary Variable Mnemonics Examples & \\
\midrule\noalign{}
\endhead
\bottomrule\noalign{}
\endlastfoot
Name & Description \\
BOFT & Block Oil Flow Total \\
COFR & Connection Oil Flow Rate \\
FOPR & Field Oil Production Rate \\
GVPR & Group Volume Production Rate \\
RTIT & Region Tracer Injection Total \\
WWPT & Well Water Production Total \\
\end{longtable}

\emph{Table 11.2.2: Summary Variable Mnemonics Examples}

And Table 11.2.3 shows some examples that do not follow the general
syntax. Where applicable, each of these keywords would also have an
equivalent mnemonic for historical input data available in history
matching runs, that is: FWCTH, WBHPH, WTHPH, WWCTH, etc.

\begin{longtable}[]{@{}
  >{\raggedright\arraybackslash}p{(\linewidth - 2\tabcolsep) * \real{0.5000}}
  >{\raggedright\arraybackslash}p{(\linewidth - 2\tabcolsep) * \real{0.5000}}@{}}
\toprule\noalign{}
\begin{minipage}[b]{\linewidth}\raggedright
Summary Variable Mnemonics Examples Not Following the General Syntax
\end{minipage} & \begin{minipage}[b]{\linewidth}\raggedright
\end{minipage} \\
\midrule\noalign{}
\endhead
\bottomrule\noalign{}
\endlastfoot
Name & Description \\
BAPI & Block API \\
BOSAT & Block Oil Saturation \\
BPR & Block Pressure \\
FAQR & Field Aquifer Influx Rate \\
FPR & Field Average Pressure \\
FWCT & Field Water Cut \\
WBHP & Well Bottom-Hole Pressure \\
WTHP & Well Tubing Head Pressure \\
WWCT & Well Water Cut \\
\end{longtable}

\emph{Table 11.2.3: Summary Variable Mnemonics Examples Not Following
the General Syntax}

\subsection{Summary Variable Format}\label{summary-variable-format}

In addition to the general mnemonic syntax, each object type (Field,
Group, Region, Well, etc.) has additional syntax governing which
specific object (group, well, etc.) should be written out to the SUMMARY
file as explained in Table 11.2.4.

\begin{longtable}[]{@{}
  >{\raggedright\arraybackslash}p{(\linewidth - 4\tabcolsep) * \real{0.3333}}
  >{\raggedright\arraybackslash}p{(\linewidth - 4\tabcolsep) * \real{0.3333}}
  >{\raggedright\arraybackslash}p{(\linewidth - 4\tabcolsep) * \real{0.3333}}@{}}
\toprule\noalign{}
\begin{minipage}[b]{\linewidth}\raggedright
Summary Variable Format
\end{minipage} & \begin{minipage}[b]{\linewidth}\raggedright
\end{minipage} & \begin{minipage}[b]{\linewidth}\raggedright
\end{minipage} \\
\midrule\noalign{}
\endhead
\bottomrule\noalign{}
\endlastfoot
Summary Object & Format & Example \\
AQUIFER & Analytical aquifers are identified by the first two characters
AA in the variable mnemonic and numerical aquifers are identified by the
first two characters of AN. Aquifer variables for a given aquifer can be
followed by a list of aquifer numbers and therefore a terminating ``/''
is required to end the list of aquifers. A blank list requests output
for all the aquifers. & For instance to obtain the aquifer influx rate
and cumulative influx for all analytical aquifers in the model, one
would use: AAQR / AAQT / or just for numerical aquifers one, two, and
three: ANQR 1 2 3 / ANQT 1 2 3 / \\
Analytical aquifer lists defined via the AQUALIST keyword in the
SOLUTION section are identified by the first two characters AL in the
variable mnemonic. Aquifer variables for aquifer lists can be followed
by any number of aquifer list names and therefore a terminating ``/'' is
required to end the list of aquifers. A blank list requests output for
all the aquifers. & For instance to obtain the total aquifer influx rate
and cumulative influx for all analytical aquifers in aquifer lists EAST
and WEST, one would use: ALQR `EAST' `WEST'/ ALQT `EAST' `WEST'/ Note
OPM Flow does not currently support Analytical Aquifer Lists. & \\
BLOCK (Grid Cells) & Block variables are followed by a list of cell (I,
J, K) indices, with each line terminated by a ``/'', and the list
terminated with another terminating ``/''. & To request the oil
saturation for a series of grid cells one would use the following: BOSAT
1 1 1 / 10 10 3 / 15 15 10 / / \\
FIELD & Field variables take no additional parameters and therefore do
not require a terminating ``/''. & For example: FOPR FOPT FPR Would
export the field oil production rate and total, plus the field average
reservoir pressure to the SUMMARY file. \\
GROUP & Group variables can be followed by a list of group names
enclosed in quotes and therefore a terminating ``/'' is required to end
the list of groups. A blank list requests output for all groups. & As an
example to get the oil production rates and totals for two groups
representing two platforms one would use GOPR `PLAT1' `PLAT2' / GOPT
`PLAT1' `PLAT2' / or GOPR / GOPT / if there are only two groups in the
model. \\
REGION & Region variables for a given region can be followed by a list
of region numbers and therefore a terminating ``/'' is required to end
the list of regions. A blank list requests output for all the regions. &
For instance to obtain the reservoir pressures and oil production totals
for all regions: RPR / ROPT / or just for regions one, two, and three:
RPR 1 2 3 / ROPT 1 2 3 / \\
Region variables that are based on flows between two regions must be
followed by a list of region number pairs, indicating the two regions.
Each region pair is terminated by a ``/''. In addition the variable
mnemonic should also be terminated by a terminating ``/''. & To obtain
the oil and water flow total (or cumulative) between regions one and
two, and also three and four, one would use: ROFT 1 2 / 3 4 / / RWFT 1 2
/ 3 4 / / & \\
WELL & Well variables can be followed by a list of well names enclosed
in quotes and therefore a terminating ``/'' is required to end the list
of wells. A blank list requests output for all wells. & As an example to
get the oil production rates and totals for all oil wells beginning with
the letters OP one would use WOPR `OP\emph{' / WOPT `OP}' / To output
the oil production rate for all wells, one may use: WOPR / Note, this
would write out the oil production rate for all well types: producers
and injectors, as well as wells declared as gas and water wells. \\
‍WELL COMPLETIONS & Well variables for well completion vectors {[}The
term well connection is used to describe individual connections from the
wellbore to the reservoir grid, as opposed to well completions. A well
completion is used to describe a set of connections, for example, a well
may consist of several completions with each completion consisting of
multiple connections.{]} generally consist of the well summary mnemonic
name suffixed with the letter ``L'', for example WOPRL, for the
completion oil production rate. The completion mnemonic is then followed
by a list of well names enclosed in quotes and the completion number,
each terminated by a ``/''. In addition the list is terminated by a
terminating ``/''. A blank list for the completions results in all
completions for a well being written out, and a blank list for the well
will result in all well completions being written to the SUMMARY file.
Care should be exercised when defaulting the list of wells and
completions as there is the potential to generate large volumes of data.
& For example, to get the completion oil production rates for well OP01
completions two, three, and four, and for completion two for well OP02
one would use: WOPRL `OP01' 2 / `OP01' 3 / `OP01' 4 / `OP02' 2 / / or to
get the oil production rates for each of the completions for well OP01:
WORPL `OP01' / / \\
WELL CONNECTIONS & Connection variables are followed by a list of well
names in quotes and connection I, J, K indices, each terminated by a
``/''. In addition the list is terminated by a terminating ``/''. A
blank list for the connections results in all connections for a well
being written out. And a blank list for the well will result in all well
connections being written to the SUMMARY file. Again, care should be
exercised when defaulting the list of wells and connections as there is
the potential to generate large volumes of data. & To request the
connection gas-oil ratio and water cut for wells OP01 and OP02 for
several connections, one would use: CGOR `OP01' 1 1 1 / `OP01' 1 1 2 /
`OP02' 10 10 1 / / CWCT `OP01' 1 1 1 / `OP01' 1 1 2 / `OP02' 10 10 1 /
/ \\
‍WELL CONNECTIONS (COMPLETIONS) & This type of data is the same as for
the Well Completion data, except that the Well Connection series of
mnemonics are used instead of the well mnemonics. Here the completion
vectors consist of the connection mnemonic name suffixed with the letter
``L'', for example COPRL, for the completion oil production rate. The
completion mnemonic is then followed by a list of well names enclosed in
quotes and the connection I, J, K indices, each terminated by a ``/''.
In addition, the list is terminated by a terminating ``/''. If any
connection has been defined as being within a completion via the
COMPLUMP keyword in the SCHEDULE section, then the completion data is
written out instead of the connection data. A blank list for the
completions results in all completions for a well being written out.
However, the well name cannot be defaulted, unlike the other well and
connection mnemonics. Care should be exercised when defaulting the list
of completions as there is the potential to generate large volumes of
data. & To request the gas-oil ratio and water cut for wells OP01 and
OP02 for several completions identified by connections that have been
lumped into those completions, one would use: CGORL `OP01' 1 1 1 /
`OP01' 1 1 2 / `OP02' 10 10 1 / / CWCTL `OP01' 1 1 1 / `OP01' 1 1 2 /
`OP02' 10 10 1 / / \\
WELL SEGMENTS & Well variables for multi-segment well segments can be
followed by a list of well names enclosed in quotes and the segment
number, each terminated by a ``/''. In addition, the list is terminated
by a terminating ``/''. A blank list for the segments results in all
segments for a well being written out. And a blank list for the well
will result in all well segments being written to the SUMMARY file. Care
should be exercised when defaulting the list of wells and segments as
there is the potential to generate large volumes of data. & For example,
to get the segment oil flow rates for well OP01 segments two, three, and
four and for segment two for well OP02 one would use: SOFR `OP01' 2 /
`OP01' 3 / `OP01' 4 / `OP02' 2 / / or to get all the segment oil flow
rates for all the segments for well OP01: SOFR `OP01' / / \\
\end{longtable}

\emph{Table 11.2.4: Summary Variable Format}

\subsection{Aquifer Variables}\label{aquifer-variables}

Table 11.2.5 outlines the aquifer summary variables based on the type of
aquifer: analytical, analytical list and numerical. Note that the
analytical list aquifer type is just an analytical aquifer ``set''
defined using the AQUALIST keyword in the SOLUTION section, that assigns
an analytic aquifer list name to a set of aquifer numbers for greater
readability in the output. This type of aquifer is not supported by OPM
Flow and neither are numerical aquifers.

\begin{longtable}[]{@{}
  >{\raggedright\arraybackslash}p{(\linewidth - 12\tabcolsep) * \real{0.1429}}
  >{\raggedright\arraybackslash}p{(\linewidth - 12\tabcolsep) * \real{0.1429}}
  >{\raggedright\arraybackslash}p{(\linewidth - 12\tabcolsep) * \real{0.1429}}
  >{\raggedright\arraybackslash}p{(\linewidth - 12\tabcolsep) * \real{0.1429}}
  >{\raggedright\arraybackslash}p{(\linewidth - 12\tabcolsep) * \real{0.1429}}
  >{\raggedright\arraybackslash}p{(\linewidth - 12\tabcolsep) * \real{0.1429}}
  >{\raggedright\arraybackslash}p{(\linewidth - 12\tabcolsep) * \real{0.1429}}@{}}
\toprule\noalign{}
\begin{minipage}[b]{\linewidth}\raggedright
Aquifer Summary Variables
\end{minipage} & \begin{minipage}[b]{\linewidth}\raggedright
\end{minipage} & \begin{minipage}[b]{\linewidth}\raggedright
\end{minipage} & \begin{minipage}[b]{\linewidth}\raggedright
\end{minipage} & \begin{minipage}[b]{\linewidth}\raggedright
\end{minipage} & \begin{minipage}[b]{\linewidth}\raggedright
\end{minipage} & \begin{minipage}[b]{\linewidth}\raggedright
\end{minipage} \\
\midrule\noalign{}
\endhead
\bottomrule\noalign{}
\endlastfoot
Variable & Root & Field & Analytical Aquifer & Analytical Aquifer List &
Numerical Aquifer & Comment \\
Aquifer Influx Rate (Water Aquifers) & QR & FAQR & AAQR & ALQR & ANQR
& \\
Aquifer Influx Total (Water Aquifers) & QT & FAQT & AAQT & ALQT & ANQT
& \\
Aquifer Influx Rate (Gas Aquifers) & QRG & FAQRG & AAQRG & ALQRG & & \\
Aquifer Influx Total (Gas Aquifers) & QTG & FAQTG & AAQTG & ALQTG & & \\
Aquifer Pressure & QP & & AAQP & & ANQP & Water pore volume weighted
average \\
Aquifer Carter-Tracy Dimensionless Pressure & QPD & & AAQPD & & & \\
Aquifer Carter-Tracy Dimensionless Time & QTD & & AAQTD & & & \\
Notes: & & & & & & \\
\end{longtable}

\emph{Table 11.2.5: Aquifer Summary Variables}

\subsubsection{Example}\label{example-489}

The following example requests the field and analytical aquifer
pressure, rate and cumulative water influxes for all analytical aquifers
to be written to the SUMMARY file.

\begin{verbatim}
-- ==============================================================================
--
-- SUMMARY SECTION
--
-- ==============================================================================
SUMMARY
--
--       ANALYTICAL AQUIFERS (FIELD)
--
FAQR
FAQT
--
--       ANALYTICAL AQUIFERS
--
AAQP
/
AAQR
/
AAQT
/
\end{verbatim}

\subsection{Field, Group, Well, Well Connection, and Completion
Variables}\label{field-group-well-well-connection-and-completion-variables}

The following table (Table 11.2.6) lists the field, group, well, well
connection, and well completion production, injection, pressure, and
productivity summary variables. Note that not all of these variables are
available in OPM Flow; however, the simulator will issue a warning
messages if this is indeed the case. It is anticipated that the number
of recognized summary variables will increase in future releases of OPM
Flow.

See also section 11.2.6 Field, Region, and Block Variables that outlines
the field, region, and block production and injection data summary
variables, and section 11.2.7 Field and Region Recovery Variables that
describes the summary variables for oil recovery efficiency and the oil
recovery mechanism.

\begin{longtable}[]{@{}
  >{\raggedright\arraybackslash}p{(\linewidth - 16\tabcolsep) * \real{0.1111}}
  >{\raggedright\arraybackslash}p{(\linewidth - 16\tabcolsep) * \real{0.1111}}
  >{\raggedright\arraybackslash}p{(\linewidth - 16\tabcolsep) * \real{0.1111}}
  >{\raggedright\arraybackslash}p{(\linewidth - 16\tabcolsep) * \real{0.1111}}
  >{\raggedright\arraybackslash}p{(\linewidth - 16\tabcolsep) * \real{0.1111}}
  >{\raggedright\arraybackslash}p{(\linewidth - 16\tabcolsep) * \real{0.1111}}
  >{\raggedright\arraybackslash}p{(\linewidth - 16\tabcolsep) * \real{0.1111}}
  >{\raggedright\arraybackslash}p{(\linewidth - 16\tabcolsep) * \real{0.1111}}
  >{\raggedright\arraybackslash}p{(\linewidth - 16\tabcolsep) * \real{0.1111}}@{}}
\toprule\noalign{}
\begin{minipage}[b]{\linewidth}\raggedright
Field, Group, Well, Well Connection, and Completion Summary Variables
\end{minipage} & \begin{minipage}[b]{\linewidth}\raggedright
\end{minipage} & \begin{minipage}[b]{\linewidth}\raggedright
\end{minipage} & \begin{minipage}[b]{\linewidth}\raggedright
\end{minipage} & \begin{minipage}[b]{\linewidth}\raggedright
\end{minipage} & \begin{minipage}[b]{\linewidth}\raggedright
\end{minipage} & \begin{minipage}[b]{\linewidth}\raggedright
\end{minipage} & \begin{minipage}[b]{\linewidth}\raggedright
\end{minipage} & \begin{minipage}[b]{\linewidth}\raggedright
\end{minipage} \\
\midrule\noalign{}
\endhead
\bottomrule\noalign{}
\endlastfoot
Type & Variable & Root & Field & Group & Well & Well Connection & Well
Completion & Comment \\
Flow & Artificial Lift Quantity & ALQ & & & WALQ & & & \\
Flow & Artificial Lift and VFP Table Number & MVFP & & & WMVFP & & &
Well VFP table number. \\
Flow & Efficiency Factor & EFF & & GEFF & WEFF & & & \\
Flow & Efficiency Factor Product & EFFG & & & WEFFG & & & Efficiency
factor of the well times the product of all the groups above the
well. \\
Flow & Energy Production Rate & EPR & FEPR & GEPR & WEPR & & & \\
Flow & Energy Production Total & EPT & FEPT & GEPT & WEPT & & & \\
Flow & Gas Consumption Rate & GCR & FGCR & GGCR & & & & \\
Flow & Gas Consumption Total & GCT & FGCT & GGCT & & & & \\
Flow & Gas Flow Rate & GFR & & & & CGFR & CGFRL & Positive for
production and negative for injection. \\
Flow & Gas Flow Rate (Free) & GFRF & & & & CGFRF & & \\
Flow & Gas Flow Rate (Solution) & GFRS & & & & CGFRS & & \\
Flow & Gas Flow Rate (Upstream) & GFRU & & & & CGFRU & & Sum of
connection gas flow rates upstream of, and including, this
connection. \\
Flow & Gas Import Rate & GIMR & FGIMR & GGIMR & & & & \\
Flow & Gas Import Total & GIMT & FGIMT & GGIMT & & & & \\
Flow & Gas Injection Guide Rate & GIGR & & GGIGR & WGIGR & & & \\
Flow & Gas Injection Rate & GIR & FGIR & GGIR & WGIR & CGIR & CGIRL & \\
Flow & Gas Injection Rate History & GIRH & FGIRH & GGIRH & WGIRH & &
& \\
Flow & Gas Injection Rate Target/Limit & GIRT & FGIRT & GGIRT & WGIRT &
& & \\
Flow & Gas Injection Total & GIT & FGIT & GGIT & WGIT & CGIT & CGITL
& \\
Flow & Gas Injection Total History & GITH & FGITH & GGITH & WGITH & &
& \\
Flow & Gas Lift Injection Rate & GLIR & FGLIR & GGLIR & WGLIR & & & Gas
lift gas with group and field total rates based on well efficiency. \\
Flow & Gas Lift Injection Total & GLIT & FGLIT & GGLIT & WGLIT & & & Gas
lift gas with group and field total rates based on well efficiency. \\
Flow & Gas Potential Injection Rate & GPI & FGPI & GGPI & WGPI & CGPI &
& Also WGIP for WGPI. \\
Flow & Gas Potential Injection Rate & GPI2 & FGPI2 & GGPI2 & WGPI2 & & &
Including shut and stopped wells. \\
Flow & Gas Potential Production Rate & GPP & FGPP & GGPP & WGPP & CGPP &
& \\
Flow & Gas Potential Production Rate & GPP2 & FGPP2 & GGPP2 & WGPP2 & &
& Including shut and stopped wells. \\
Flow & Gas Potential Production Rate (Free) & GPPF & FGPPF & GGPPF &
WGPPF & & & \\
Flow & Gas Potential Production Rate (Free) & GPPF2 & FGPPF2 & GGPPF2 &
WGPPF2 & & & Including shut and stopped wells. \\
Flow & Gas Potential Production Rate (Soln) & GPPS & FGPPS & GGPPS &
WGPPS & & & \\
Flow & Gas Potential Production Rate (Soln) & GPPS2 & FGPPS2 & GGPPS2 &
WGPPS2 & & & Including shut and stopped wells. \\
Flow & Gas Production Guide Rate & GPGR & & GGPGR & WGPGR & & & \\
Flow & Gas Production Rate & GPR & FGPR & GGPR & WGPR & CGPR & &
Produced reservoir gas only, gas lift gas excluded. \\
Flow & Gas Production Rate (Free) & GPRF & FGPRF & GGPRF & WGPRF & &
& \\
Flow & Gas Production Rate (Solution) & GPRS & FGPRS & GGPRS & WGPRS & &
& \\
Flow & Gas Production Rate History & GPRH & FGPRH & GGPRH & WGPRH & &
& \\
Flow & Gas Production Rate Target/Limit & GPRT & FGPRT & GGPRT & WGPRT &
& & \\
Flow & Gas Production Total & GPT & FGPT & GGPT & WGPT & CGPT & CGPTL
& \\
Flow & Gas Production Total (Free) & GPTF & FGPTF & GGPTF & WGPTF & &
& \\
Flow & Gas Production Total (Solution) & GPTS & FGPTS & GGPTS & WGPTS &
& & \\
Flow & Gas Production Total History & GPTH & FGPTH & GGPTH & WGPTH & &
& \\
Flow & Gas Sales Rate & GSR & FGSR & GGSR & & & & \\
Flow & Gas Sales Rate & SGR & FSGR & GSGR & & & & Same as GSR. \\
Flow & Gas Sales Total & GST & FGST & GGST & & & & \\
Flow & Gas Sales Total & SGT & FSGT & GSGT & & & & Same as GST. \\
Flow & Gas Voidage Injection Rate & GVIR & & & WGVIR & & & \\
Flow & Gas Voidage Production Rate & GVPR & & & WGVPR & & & \\
Flow & Liquid Flow Rate & LFR & & & & CLFR & CLFRL & Positive for
production and negative for injection. \\
Flow & Liquid Production Rate & LPR & FLPR & GLPR & WLPR & CLPR & & \\
Flow & Liquid Production Rate History & LPRH & FLPRH & GLPRH & WLPRH & &
& \\
Flow & Liquid Production Rate Target/Limit & LPRT & FLPRT & GLPRT &
WLPRT & & & \\
Flow & Liquid Production Total & LPT & FLPT & GLPT & WLPT & CLPT & CLPTL
& \\
Flow & Liquid Production Total History & LPTH & FLPTH & GLPTH & WLPTH &
& & \\
Flow & Oil Flow Rate & OFR & & & & COFR & COFRL & Positive for
production and negative for injection. \\
Flow & Oil Flow Rate (Free) & OFRF & & & & COFRF & & \\
Flow & Oil Flow Rate (Solution) & OFRS & & & & COFRS & & \\
Flow & Oil Flow Rate (Upstream) & OFRU & & & & COFRU & & Sum of
connection oil flow rates upstream of, and including this connection. \\
Flow & Oil Injection Guide Rate & OIGR & & GOIGR & WOIGR & & & \\
Flow & Oil Injection Rate & OIR & FOIR & GOIR & WOIR & COIR & COIRL & \\
Flow & Oil Injection Rate History & OIRH & FOIRH & GOIRH & WOIRH & &
& \\
Flow & Oil Injection Rate Target/Limit & OIRT & FOIRT & GOIRT & WOIRT &
& & \\
Flow & Oil Injection Total & OIT & FOIT & GOIT & WOIT & COIT & COITL
& \\
Flow & Oil Injection Total History & OITH & FOITH & GOITH & WOITH & &
& \\
Flow & Oil Potential Injection Rate & OPI & FOPI & GOPI & WOPI & & & \\
Flow & Oil Potential Injection Rate & OPI2 & FOPI2 & GOPI2 & WOPI2 & & &
Including shut and stopped wells. \\
Flow & Oil Potential Production Rate & OPP & FOPP & GOPP & WOPP & COPP &
& \\
Flow & Oil Potential Production Rate & OPP2 & FOPP2 & GOPP2 & WOPP2 & &
& Including shut and stopped wells. \\
Flow & Oil Production Guide Rate & OPGR & & GOPGR & WOPGR & & & \\
Flow & Oil Production Rate & OPR & FOPR & GOPR & WOPR & COPR & & \\
Flow & Oil Production Rate (Free) & OPRF & FOPRF & GOPRF & WOPRF & &
& \\
Flow & Oil Production Rate (Solution) & OPRS & FOPRS & GOPRS & WOPRS & &
& \\
Flow & Oil Production Rate History & OPRH & FOPRH & GOPRH & WOPRH & &
& \\
Flow & Oil Production Rate Target/Limit & OPRT & FOPRT & GOPRT & WOPRT &
& & \\
Flow & Oil Production Total & OPT & FOPT & GOPT & WOPT & & & \\
Flow & Oil Production Total (Free) & OPTF & FOPTF & GOPTF & WOPTF & &
& \\
Flow & Oil Production Total (Solution) & OPTS & FOPTS & GOPTS & WOPTS &
& & \\
Flow & Oil Production Total History & OPTH & FOPTH & GOPTH & WOPTH & &
& \\
Flow & Res. Vol. Flow Rate & VFR & & & & CVFR & CVFRL & Positive for
production and negative for injection. \\
Flow & Res. Vol. Flow Rate (Free) & VFRF & & & & CVFRF & & \\
Flow & Res. Vol. Injection Guide Rate & VIGR & & GVIGR & WVIGR & & & \\
Flow & Res. Vol. Injection Rate & VIR & FVIR & GVIR & WVIR & CVIR &
CVIRL & \\
Flow & Res. Vol. Injection Rate Target/Limit & VIRT & FVIRT & GVIRT &
WVIRT & & & \\
Flow & Res. Vol. Injection Total & VIT & FVIT & GVIT & WVIT & CVIT &
CVITL & \\
Flow & Res. Vol. Potential Injection Rate & VPI & FVPI & GVPI & WVPI & &
& \\
Flow & Res. Vol. Potential Injection Rate & VPI2 & FVPI2 & GVPI2 & WVPI2
& & & Including shut and stopped wells. \\
Flow & Res. Vol. Potential Production Rate & VPP & FVPP & GVPP & WVPP &
CVPP & & \\
Flow & Res. Vol. Potential Production Rate & VPP2 & FVPP2 & GVPP2 &
WVPP2 & & & Including shut and stopped wells. \\
Flow & Res. Vol. Production Guide Rate & VPGR & & GVPGR & WVPGR & & & \\
Flow & Res. Vol. Production Rate & VPR & FVPR & GVPR & WVPR & CVPR &
& \\
Flow & Res. Vol. Production Rate Target/Limit & VPRT & FVPRT & GVPRT &
WVPRT & & & \\
Flow & Res. Vol. Production Total & VPT & FVPT & GVPT & WVPT & CVPT &
& \\
Flow & Water Flow Rate & WFR & & & & CWFR & CWFRL & Positive for
production and negative for injection. \\
Flow & Water Flow Rate (Upstream) & WFRU & & & & CWFRU & & Sum of
connection water flow rates upstream of, and including, this
connection. \\
Flow & Water Injection Guide Rate & WIGR & & GWIGR & WWIGR & & & \\
Flow & Water Injection Rate & WIR & FWIR & GWIR & WWIR & CWIR & CWIRL
& \\
Flow & Water Injection Rate History & WIRH & FWIRH & GWIRH & WWIRH & &
& \\
Flow & Water Injection Rate Target/Limit & WIRT & FWIRT & GWIRT & WWIRT
& & & \\
Flow & Water Injection Total & WIT & FWIT & GWIT & WWIT & CWIT & CWITL
& \\
Flow & Water Injection Total History & WITH & FWITH & GWITH & WWITH & &
& \\
Flow & Water Potential Injection Rate & WPI & FWPI & GWPI & WWPI & & &
Also WWIP for WWPI. \\
Flow & Water Potential Injection Rate & WPI2 & FWPI2 & GWPI2 & WWPI2 & &
& Including shut and stopped wells. \\
Flow & Water Potential Production Rate & WPP & FWPP & GWPP & WWPP & CWPP
& & \\
Flow & Water Potential Production Rate & WPP2 & FWPP2 & GWPP2 & WWPP2 &
& & Including shut and stopped wells. \\
Flow & Water Produce/Injected Ratio & WPIR & FWPIR & GWPIR & WWPIR & & &
Reported as a percent. \\
Flow & Water Production Guide Rate & WPGR & & GWPGR & WWPGR & & & \\
Flow & Water Production Rate & WPR & FWPR & GWPR & WWPR & CWPR & & \\
Flow & Water Production Rate History & WPRH & FWPRH & GWPRH & WWPRH & &
& \\
Flow & Water Production Rate Target/Limit & WPRT & FWPRT & GWPRT & WWPRT
& & & \\
Flow & Water Production Total & WPT & FWPT & GWPT & WWPT & CWPT & CWPTL
& \\
Flow & Water Production Total History & WPTH & FWPTH & GWPTH & WWPTH & &
& \\
Flow & Water Voidage Injection Guide Rate & WVIR & & & WWVIR & & & \\
Flow & Water Voidage Injection Guide Total & WVIT & & & WWVIT & & & \\
Flow & Water Mass Injection Rate & AMIR & FAMIR & GAMIR & WAMIR & CAMIR
& CAMIRL & \\
Flow & Water Mass Injection Total & AMIT & FAMIT & GAMIT & WAMIT & CAMIT
& CAMITL & \\
Flow & Water Mass Production Rate & AMPR & FAMPR & GAMPR & WAMPR & CAMPR
& CAMPRL & \\
Flow & Water Mass Production Total & AMPT & FAMPT & GAMPT & WAMPT &
CAMPT & CAMPTL & \\
No.~Wells & Number of abandoned injection wells & MWIA & FMWIA & GMWIA &
& & & \\
No.~Wells & Number of abandoned production wells & MWPA & FMWPA & GMWPA
& & & & \\
No.~Wells & Number of drilling events in total & MWDT & FMWDT & GMWDT &
& & & \\
No.~Wells & Number of drilling events this timestep & MWDR & FMWDR &
GMWDR & & & & \\
No.~Wells & Number of injection wells currently flowing & MWIN & FMWIN &
GMWIN & & & & \\
No.~Wells & Number of injectors on group control & MWIG & FMWIG & GMWIG
& & & & \\
No.~Wells & Number of injectors on own reservoir volume rate limit
control & MWIV & FMWIV & GMWIV & & & & \\
No.~Wells & Number of injectors on own surface rate limit control & MWIS
& FMWIS & GMWIS & & & & \\
No.~Wells & Number of injectors on pressure control & MWIP & FMWIP &
GMWIP & & & & \\
No.~Wells & Number of producers controlled by own oil rate limit & MWPO
& FMWPO & GMWPO & & & & \\
No.~Wells & Number of producers on group control & MWPG & FMWPG & GMWPG
& & & & \\
No.~Wells & Number of producers on own reservoir volume rate limit
control & MWPV & FMWPV & GMWPV & & & & \\
No.~Wells & Number of producers on own surface rate limit control & MWPS
& FMWPS & GMWPS & & & & \\
No.~Wells & Number of producers on pressure control & MWPP & FMWPP &
GMWPP & & & & \\
No.~Wells & Number of producers using artificial lift (with ALQ
\textgreater{} 0.0) & MWPL & FMWPL & GMWPL & & & & \\
No.~Wells & Number of production wells currently flowing & MWPR & FMWPR
& GMWPR & & & & \\
No.~Wells & Number of unused injection wells & MWIU & FMWIU & GMWIU & &
& & \\
No.~Wells & Number of unused production wells & MWPU & FMWPU & GMWPU & &
& & \\
No.~Wells & Number of workover events in total & MWWT & FMWWT & GMWWT &
& & & \\
No.~Wells & Number of workover events this time step. & MWWO & FMWWO &
GMWWO & & & & \\
No.~Wells & The number of connections capable of flowing in the well &
MCON & & & WMCON & & & \\
No.~Wells & Total number of injection wells & MWIT & FMWIT & GMWIT & & &
& \\
No.~Wells & Total number of production wells & MWPT & FMWPT & GMWPT & &
& & \\
Pressure & Blocking Factor (GPP) & CDBF & & & & CCDBF & & \\
Pressure & Bottom-Hole Pressure & BHP & & & WBHP & & & \\
Pressure & Bottom-Hole Pressure (Five Point & BP5 & & & WBP5 & & & \\
Pressure & Bottom-Hole Pressure (Four Point) & BP4 & & & WBP4 & & & \\
Pressure & Bottom-Hole Pressure (Nine Point) & BP9 & & & WBP9 & & & \\
Pressure & Bottom-Hole Pressure (One Point) & BP & & & WBP & & & \\
Pressure & Bottom-Hole Pressure History & BHPH & & & WBHPH & & & \\
Pressure & Bottom-Hole Pressure Target/Limit & BHPT & & & WBHPT & & & \\
Pressure & Connection Pressure & PR & & & & CPR & CPRL & Completion
pressure is an average pressure for the completion. \\
Pressure & Connection D-Factor & DFAC & & & & CDFAC & & \\
Pressure & Connection Transmissibility Factor & TFAC & & & & CTFAC &
& \\
Pressure & Productivity (Gas Phase) & PIG & & & WPIG & & & \\
Pressure & Productivity (Liquid Phase) & PIL & & & WPIL & & & \\
Pressure & Productivity (Oil Phase) & PIO & & & WPIO & & & \\
Pressure & Productivity (Preferred Phase -- Five Point) & PI5 & & & WPI5
& & & \\
Pressure & Productivity (Preferred Phase -- Four Point) & PI4 & & & WPI4
& & & \\
Pressure & Productivity (Preferred Phase -- Nine Point) & PI9 & & & WPI9
& & & \\
Pressure & Productivity (Preferred Phase -- One Point) & PI1 & & & WPI1
& & & \\
Pressure & Productivity (Preferred Phase) & PI & & & WPI & CPI & & \\
Pressure & Productivity (Water Phase) & PIW & & & WPIW & & & \\
Pressure & Tubing Head Pressure & THP & & & WTHP & & & \\
Pressure & Tubing Head Pressure History & THPH & & & WTHPH & & & \\
Ratio & Gas-Liquid Ratio & GLR & FGLR & GGLR & WGLR & CGLR & CGLRL & \\
Ratio & Gas-Liquid Ratio (Bottom-Hole) & BGLR & & & WBGLR & & & \\
Ratio & Gas-Liquid Ratio History & GLRH & FGLRH & GGLRH & WGLRH & & & \\
Ratio & Gas-Oil Ratio & GOR & FGOR & GGOR & WGOR & CGOR & CGORL & \\
Ratio & Gas-Oil Ratio History & GORH & FGORH & GGORH & WGORH & & & \\
Ratio & Oil-Gas Ratio & OGR & FOGR & GOGR & WOGR & COGR & COGRL & \\
Ratio & Oil-Gas Ratio History & OGRH & FOGRH & GOGRH & WOGRH & & & \\
Ratio & Water -Gas Ratio & WGR & FWGR & GWGR & WWGR & CWGR & CWGRL & \\
Ratio & Water Cut & WCT & FWCT & GWCT & WWCT & CWCT & CWCTL & \\
Ratio & Water Cut History & WCTH & FWCTH & GWCTH & WWCTH & & & \\
Ratio & Water-Gas Ratio History & WGRH & FWGRH & GWGRH & WWGRH & & & \\
Notes: & & & & & & & & \\
\end{longtable}

\emph{Table 11.2.6: Field, Group, Well, Well Connection, and Completion
Summary Variables}

\subsection{Field, Group and Well Control Mode
Variables}\label{field-group-and-well-control-mode-variables}

In addition to the production, injection, pressure, and productivity
summary variables, there are also summary variables that output the
control mode at which the field, groups and wells are being controlled.
Table 11.2.7 summarizes the mnemonics for the field and groups together
with a description and the values written out to the SUMMARY file.

\begin{longtable}[]{@{}
  >{\raggedright\arraybackslash}p{(\linewidth - 12\tabcolsep) * \real{0.1429}}
  >{\raggedright\arraybackslash}p{(\linewidth - 12\tabcolsep) * \real{0.1429}}
  >{\raggedright\arraybackslash}p{(\linewidth - 12\tabcolsep) * \real{0.1429}}
  >{\raggedright\arraybackslash}p{(\linewidth - 12\tabcolsep) * \real{0.1429}}
  >{\raggedright\arraybackslash}p{(\linewidth - 12\tabcolsep) * \real{0.1429}}
  >{\raggedright\arraybackslash}p{(\linewidth - 12\tabcolsep) * \real{0.1429}}
  >{\raggedright\arraybackslash}p{(\linewidth - 12\tabcolsep) * \real{0.1429}}@{}}
\toprule\noalign{}
\begin{minipage}[b]{\linewidth}\raggedright
Field and Group Control Mode Reporting
\end{minipage} & \begin{minipage}[b]{\linewidth}\raggedright
\end{minipage} & \begin{minipage}[b]{\linewidth}\raggedright
\end{minipage} & \begin{minipage}[b]{\linewidth}\raggedright
\end{minipage} & \begin{minipage}[b]{\linewidth}\raggedright
\end{minipage} & \begin{minipage}[b]{\linewidth}\raggedright
\end{minipage} & \begin{minipage}[b]{\linewidth}\raggedright
\end{minipage} \\
\midrule\noalign{}
\endhead
\bottomrule\noalign{}
\endlastfoot
Object & Field & Group & Field & Group & Field & Group \\
Mnemonic & FMCTP & GMCTP & FMCTW & GMCTW & FMCTG & GMCTG \\
Description & Production Group. & Water Injection Group. & Gas Injection
Group. & & & \\
Values & Value & Mnemonic & Value & Mnemonic & Value & Mnemonic \\
0.0 & NONE & 0.0 & NONE & 0.0 & NONE & \\
1.0 & ORAT & 1.0 & RATE & 1.0 & RATE & \\
2.0 & WRAT & 2.0 & RESV & 2.0 & RESV & \\
3.0 & GRAT & 3.0 & REIN & 3.0 & REIN & \\
4.0 & LRAT & 4.0 & VREP & 4.0 & VREP & \\
5.0 & RESV & 5.0 & & 5.0 & & \\
6.0 & PRBL & 6.0 & & 6.0 & & \\
7.0 & ENERGY & 7.0 & & 7.0 & & \\
8.0 & WGRA & 8.0 & WGRA & 8.0 & WGRA & \\
9.0 & CRAT & 9.0 & Availability & 9.0 & Availability & \\
10.0 & PBGS & 10.0 & & 10.0 & & \\
11.0 & PBWS & 11.0 & & 11.0 & & \\
12.0 & CRAT & 12.0 & & 12.0 & & \\
Notes: & & & & & & \\
\end{longtable}

\emph{Table 11.2.7: Field and Group Control Mode Reporting}

Similarly, Table 11.2.8 shows the control mode reporting variables for
wells.

\begin{longtable}[]{@{}
  >{\raggedright\arraybackslash}p{(\linewidth - 12\tabcolsep) * \real{0.1429}}
  >{\raggedright\arraybackslash}p{(\linewidth - 12\tabcolsep) * \real{0.1429}}
  >{\raggedright\arraybackslash}p{(\linewidth - 12\tabcolsep) * \real{0.1429}}
  >{\raggedright\arraybackslash}p{(\linewidth - 12\tabcolsep) * \real{0.1429}}
  >{\raggedright\arraybackslash}p{(\linewidth - 12\tabcolsep) * \real{0.1429}}
  >{\raggedright\arraybackslash}p{(\linewidth - 12\tabcolsep) * \real{0.1429}}
  >{\raggedright\arraybackslash}p{(\linewidth - 12\tabcolsep) * \real{0.1429}}@{}}
\toprule\noalign{}
\begin{minipage}[b]{\linewidth}\raggedright
Well Control Mode Reporting
\end{minipage} & \begin{minipage}[b]{\linewidth}\raggedright
\end{minipage} & \begin{minipage}[b]{\linewidth}\raggedright
\end{minipage} & \begin{minipage}[b]{\linewidth}\raggedright
\end{minipage} & \begin{minipage}[b]{\linewidth}\raggedright
\end{minipage} & \begin{minipage}[b]{\linewidth}\raggedright
\end{minipage} & \begin{minipage}[b]{\linewidth}\raggedright
\end{minipage} \\
\midrule\noalign{}
\endhead
\bottomrule\noalign{}
\endlastfoot
Object & Well & Well & & & & \\
Mnemonic & WSTAT & WMCTL & & & & \\
Description & Well Status indicator. & Well Mode of Control indicator. &
& & & \\
Values & Value & Mnemonic & Description & Value & Mnemonic &
Description \\
1.0 & PROD & Well is a producing well. & -1.0 & & Well is under group
control. & \\
2.0 & INJ & Well is an injection well. & 0.0 & SHUT or STOP & Well is
shut-in or stopped. & \\
3.0 & SHUT & Well is shut-in except those under priority control. & 1.0
& ORAT & Well is on oil rate control. & \\
4.0 & STOP & Well is stopped except those under priority control. & 2.0
& WRAT & Well is on water rate control. & \\
5.0 & PSHUT & Well is shut-in under priority control. & 3.0 & GRAT &
Well is on gas rate control. & \\
6.0 & PSTOP & Well is stopped under priority control & 4.0 & LRAT & Well
is on liquid rate control. & \\
& 5.0 & RESV & Well is on reservoir voidage rate control. & & & \\
6.0 & THP & Well is on THP control. & & & & \\
7.0 & BHP & Well is on BHP control. & & & & \\
9.0 & CRAT & Not used & & & & \\
11.0 & GOR Penalty & & & & & \\
12.0 & Drawdown & Well is under drawdown control. & & & & \\
13.0 & NGL & Not used & & & & \\
30.0 & Availability & & & & & \\
31.0 & REIN & & & & & \\
32.0 & TMRA & & & & & \\
33.0 & WGRA & & & & & \\
34.0 & CVAL & & & & & \\
41.0 - 51.0 & & & & & & \\
Notes: & & & & & & \\
\end{longtable}

\emph{Table 11.2.8: Well Control Mode Reporting}

\subsection{Field, Region, and Block
Variables}\label{field-region-and-block-variables}

The next set of summary data is based on the block, field and region
summary variables and the variables are tabulated in Table 11.2.6. Note
when the data is for a block variable, like for example the oil
saturation, this will be the actual value for the a grid cell (block),
whereas, for the field and region items this will be an average value.

See also section 11.2.4 Field, Group, Well, Well Connection, and
Completion Variables that outlines the field production data summary
variables, and section 11.2.7 Field and Region Recovery Variables that
describes the summary variables for oil recovery efficiency and the oil
recovery mechanism.

\begin{longtable}[]{@{}
  >{\raggedright\arraybackslash}p{(\linewidth - 12\tabcolsep) * \real{0.1429}}
  >{\raggedright\arraybackslash}p{(\linewidth - 12\tabcolsep) * \real{0.1429}}
  >{\raggedright\arraybackslash}p{(\linewidth - 12\tabcolsep) * \real{0.1429}}
  >{\raggedright\arraybackslash}p{(\linewidth - 12\tabcolsep) * \real{0.1429}}
  >{\raggedright\arraybackslash}p{(\linewidth - 12\tabcolsep) * \real{0.1429}}
  >{\raggedright\arraybackslash}p{(\linewidth - 12\tabcolsep) * \real{0.1429}}
  >{\raggedright\arraybackslash}p{(\linewidth - 12\tabcolsep) * \real{0.1429}}@{}}
\toprule\noalign{}
\begin{minipage}[b]{\linewidth}\raggedright
Block, Field, and Region Summary Variables
\end{minipage} & \begin{minipage}[b]{\linewidth}\raggedright
\end{minipage} & \begin{minipage}[b]{\linewidth}\raggedright
\end{minipage} & \begin{minipage}[b]{\linewidth}\raggedright
\end{minipage} & \begin{minipage}[b]{\linewidth}\raggedright
\end{minipage} & \begin{minipage}[b]{\linewidth}\raggedright
\end{minipage} & \begin{minipage}[b]{\linewidth}\raggedright
\end{minipage} \\
\midrule\noalign{}
\endhead
\bottomrule\noalign{}
\endlastfoot
Type & Variable & Root & Field & Region & Block & Comment \\
Flow & Gas Flow (Inter-Block in I+ Direction) & FLOGI & & & BFLOGI & \\
Flow & Gas Flow (Inter-Block in I- Direction) & FLOGI- & & & BFLOGI-
& \\
Flow & Gas Flow (Inter-Block in J+ Direction) & FLOGJ & & & BFLOGJ & \\
Flow & Gas Flow (Inter-Block in J- Direction) & FLOGJ- & & & BFLOGJ-
& \\
Flow & Gas Flow (Inter-Block in K+ Direction) & FLOGK & & & BFLOGK & \\
Flow & Gas Flow (Inter-Block in K- Direction) & FLOGK- & & & BFLOGK-
& \\
Flow-Reg & Gas Flow Rate Inter-Region & GFR & & RGFR & & OPM Flow
implementation. \\
Flow-Reg & Gas Flow Rate Inter-Region (Positive Contribution) & GFR+ & &
RGFR+ & & \\
Flow-Reg & Gas Flow Rate Inter-Region (Negative Contribution) & GFR- & &
RGFR- & & \\
Flow-Reg & Gas Flow Total Inter-Region (Liquid and Gas Phases) & GFT & &
RGFT & & \\
Flow-Reg & Gas Flow Total Inter-Region (+VE Liquid and Gas Phases) &
GFT+ & & RGFT+ & & OPM Flow Implementation \\
Flow-Reg & Gas Flow Total Inter-Region (-VE Liquid and Gas Phases) &
GFT- & & RGFT- & & \\
Flow-Reg & Gas Flow Total Inter-Region (Gas Phase) & GFTG & & RGFTG &
& \\
Flow-Reg & Gas Flow Total Inter-Region (Liquid Phase) & GFTL & & RGFTL &
& \\
Flow & Gas Injection Rate & GIR & & RGIR & & \\
Flow & Gas Injection Total & GIT & & RGIT & & \\
Flow & Gas Potential & PPG & FPPG & RPPG & BPPG & \\
Flow & Gas Production Rate & GPR & & RGPR & & \\
Flow & Gas Production Rate (Free) & GPRF & & RGPRF & & \\
Flow & Gas Production Rate (Solution) & GPRS & & RGPRS & & \\
Flow & Gas Production Total & GPT & & RGPT & & Liquid and gas phases \\
Flow & Gas Production Total (Free) & GPTF & & RGPTF & & \\
Flow & Gas Production Total (Net From Region) & GP & & RGP & & Minus
injected gas. \\
Flow & Gas Production Total (Solution) & GPTS & & RGPTS & & \\
Flow & Gas Velocity (Inter-Block in I+ Direction) & VELGI & & & BVELGI
& \\
Flow & Gas Velocity (Inter-Block in I- Direction) & VELGI- & & & BVELGI-
& \\
Flow & Gas Velocity (Inter-Block in J+ Direction) & VELGJ & & & BVELGJ
& \\
Flow & Gas Velocity (Inter-Block in J- Direction) & VELGJ- & & & BVELGJ-
& \\
Flow & Gas Velocity (Inter-Block in K+ Direction) & VELGK & & & BVELGK
& \\
Flow & Gas Velocity (Inter-Block in K- Direction) & VELGK- & & & BVELGK-
& \\
Flow & Oil Flow (Inter-Block in I+ Direction) & FLOOI & & & BFLOOI & \\
Flow & Oil Flow (Inter-Block in I- Direction) & FLOOI- & & & BFLOOI-
& \\
Flow & Oil Flow (Inter-Block in J+ Direction) & FLOOJ & & & BFLOOJ & \\
Flow & Oil Flow (Inter-Block in J- Direction) & FLOOJ- & & & BFLOOJ-
& \\
Flow & Oil Flow (Inter-Block in K+ Direction) & FLOOK & & & BFLOOK & \\
Flow & Oil Flow (Inter-Block in K- Direction) & FLOOK- & & & BFLOOK-
& \\
Flow-Reg & Oil Flow Rate Inter-Region & OFR & & ROFR & & OPM Flow
implementation. \\
Flow-Reg & Oil Flow Rate Inter-Region (Positive Contribution) & OFR+ & &
ROFR+ & & \\
Flow-Reg & Oil Flow Rate Inter-Region (Negative Contribution) & OFR- & &
ROFR- & & \\
Flow-Reg & Oil Flow Total Inter-Region (Liquid and Gas Phases) & OFT & &
ROFT & & \\
Flow-Reg & Oil Flow Total Inter-Region (+VE Liquid and Gas Phases) &
OFT+ & & ROFT+ & & OPM Flow implementation. \\
Flow-Reg & Oil Flow Total Inter-Region (-VE Liquid and Gas Phases) &
OFT- & & ROFT- & & \\
Flow-Reg & Oil Flow Total Inter-Region (Gas Phase) & OFTG & & ROFTG &
& \\
Flow-Reg & Oil Flow Total Inter-Region (Liquid Phase) & OFTL & & ROFTL &
& \\
Flow & Oil Injection Rate & OIR & & ROIR & & \\
Flow & Oil Injection Total & OIT & & ROIT & & \\
Flow & Oil Potential & PPO & FPPO & RPPO & BPPO & \\
Flow & Oil Production Rate & OPR & & ROPR & & \\
Flow & Oil Production Total & OPT & & ROPT & & \\
Flow & Oil Production Total (Net From Region) & OP & & ROP & & \\
Flow & Oil Velocity (Inter-Block in I+ Direction) & VELOI & & & BVELOI
& \\
Flow & Oil Velocity (Inter-Block in I- Direction) & VELOI- & & & BVELOI-
& \\
Flow & Oil Velocity (Inter-Block in J+ Direction) & VELOJ & & & BVELOJ
& \\
Flow & Oil Velocity (Inter-Block in J- Direction) & VELOJ- & & & BVELOJ-
& \\
Flow & Oil Velocity (Inter-Block in K+ Direction) & VELOK & & & BVELOK
& \\
Flow & Oil Velocity (Inter-Block in K- Direction) & VELOK- & & & BVELOK-
& \\
Flow & Water Flow (Inter-Block in I+ Direction) & FLOWI & & & BFLOWI
& \\
Flow & Water Flow (Inter-Block in I- Direction) & FLOWI- & & & BFLOWI-
& \\
Flow & Water Flow (Inter-Block in J+ Direction) & FLOWJ & & & BFLOWJ
& \\
Flow & Water Flow (Inter-Block in J- Direction) & FLOWJ- & & & BFLOWJ-
& \\
Flow & Water Flow (Inter-Block in K+ Direction) & FLOWK & & & BFLOWK
& \\
Flow & Water Flow (Inter-Block in K- Direction) & FLOWK- & & & BFLOWK-
& \\
Flow-Reg & Water Flow Rate Inter-Region & WFR & & RWFR & & OPM Flow
implementation. \\
Flow-Reg & Water Flow Rate Inter-Region & WFR+ & & RWFR+ & & \\
Flow-Reg & Water Flow Rate Inter-Region & WFR- & & RWFR- & & \\
Flow-Reg & Water Flow Total Inter-Region & WFT & & RWFT & & \\
Flow-Reg & Water Flow Total Inter-Region & WFT+ & & RWFT+ & & OPM Flow
implementation. \\
Flow-Reg & Water Flow Total Inter-Region & WFT- & & RWFT- & & \\
Flow & Water Injection Rate & WIR & & RWIR & & \\
Flow & Water Injection Total & WIT & & RWIT & & \\
Flow & Water Potential & PPW & FPPW & RPPW & BPPW & \\
Flow & Water Production Rate & WPR & & RWPR & & \\
Flow & Water Production Total & WPT & & RWPT & & \\
Flow & Water Production Total (Net From Region) & WP & & RWP & & \\
Flow & Water Velocity (Inter-Block in I+ Direction) & VELWI & & & BVELWI
& \\
Flow & Water Velocity (Inter-Block in I- Direction) & VELWI- & & &
BVELWI- & \\
Flow & Water Velocity (Inter-Block in J+ Direction) & VELWJ & & & BVELWJ
& \\
Flow & Water Velocity (Inter-Block in J- Direction) & VELWJ- & & &
BVELWJ- & \\
Flow & Water Velocity (Inter-Block in K+ Direction) & VELWK & & & BVELWK
& \\
Flow & Water Velocity (Inter-Block in K- Direction) & VELWK- & & &
BVELWK- & \\
Mass & Water In-place & AMIP & FAMIP & RAMIP & BAMIP & \\
Pressure & Gas Phase Pressure & GPR & & & BGPR & \\
Pressure & Oil Phase Pressure & PR & FPR & RPR & BPR & Field and Region
HCPV weighted. \\
BPRESSUR & & & & & & \\
Pressure & Oil Phase Pressure & PRH & FPRH & RPRH & & Field and Region
HCPV weighted. \\
Pressure & Oil Phase Pressure (PV Weighted) & PRP & FPRP & RPRP & &
Field and Region PV weighted. \\
Pressure & Water Phase Pressure & WPR & & & BWPR & \\
Property & Block Hydrocarbon Phase State & STATE & & & BSTATE & Gas (1),
Gas \& Oil (2), and Oil (3). \\
Property & Bubble Point Pressure & PBUB & & & BPBUB & \\
Property & Capillary Pressure (Gas-Oil) & GPC & & & BGPC & \\
Property & Capillary Pressure (Water-Oil) & WPC & & & BWPC & \\
Property & Dew Point & PDEW & & & BPDEW & \\
Property & Gas P/Z & PRGZ & FPRGZ & RPRGZ & & \\
Property & Gas Reservoir Density & GDEN & FGDEN & RGDEN & BGDEN & \\
Property & Gas Reservoir Density & DENG & & & BDENG & Same as BGDEN. \\
Property & Gas Saturation & GSAT & FGSAT & RGSAT & BGSAT & Also
BSGAS. \\
Property & Gas Saturation (Drainage to Imbibition) & GSHY & & & BGSHY &
Leaving gas saturation for gas capillary pressure hysteresis. \\
Property & Gas Saturation (Dynamically Trapped) & GTPD & & & BGTPD & WAG
hysteresis only. \\
Property & Gas Saturation (Trapped Critical) & GSTRP & & & BGSTRP & For
gas capillary pressure hysteresis. \\
Property & Gas Saturation (Trapped) & GTRP & & & BGTRP & WAG hysteresis
only. \\
Property & Gas Viscosity & GVIS & FGVIS & RGVIS & BGVIS & \\
Property & Gas-Oil Ratio (Saturated) & RSSAT & & & BRSSAT & \\
Property & Initial Contact Corrected Potential & PPC & FPPC & RPPC &
BPPC & \\
Property & Oil Reservoir Density & ODEN & FODEN & RODEN & BODEN & \\
Property & Oil Reservoir Density & DENO & & & BDENO & Same as BODEN. \\
Property & Oil Saturation & OSAT & FOSAT & ROSAT & BOSAT & Also
BSOIL. \\
Property & Oil Viscosity & OVIS & FOVIS & ROVIS & BOVIS & Also BVOIL. \\
Property & Oil-Gas Ratio (Saturated) & RVSAT & & & BRVSAT & \\
Property & Relative Permeability (Gas in the R- Direction) & GKRR- & & &
BGKRR- & \\
Property & Relative Permeability (Gas in the R+ Direction) & GKRR & & &
BGKRR & \\
Property & Relative Permeability (Gas in the Theta- Direction) & GKRT- &
& & BGKRT- & \\
Property & Relative Permeability (Gas in the Theta+ Direction) & GKRT &
& & BGKRT & \\
Property & Relative Permeability (Gas in the X- Direction) & GKRX- & & &
BGKRX- & Also BGKRI-. \\
Property & Relative Permeability (Gas in the X+ Direction) & GKRX & & &
BGKRX & Also BGKRI. \\
Property & Relative Permeability (Gas in the Y- Direction) & GKRY- & & &
BGKRY- & Also BGKRJ-. \\
Property & Relative Permeability (Gas in the Y+ Direction) & GKRY & & &
BGKRY & Also BGKRJ. \\
Property & Relative Permeability (Gas in the Z- Direction) & GKRZ- & & &
BGKRZ- & Also BGKRK-. \\
Property & Relative Permeability (Gas in the Z+ Direction) & GKRZ & & &
BGKRZ & Also BGKRK. \\
Property & Relative Permeability (Gas) & KRG & & & BKRG & Also BGWR. \\
Property & Relative Permeability (Krw Reduction Factor) & RK & & & BRK &
Due to Polymer. \\
Property & Relative Permeability (Oil Two Phase to Gas) & KROG & & &
BKROG & \\
Property & Relative Permeability (Oil Two Phase to Water) & KROW & & &
BKROW & \\
Property & Relative Permeability (Oil in the R- Direction) & OKRR- & & &
BOKRR- & \\
Property & Relative Permeability (Oil in the R+ Direction) & OKRR & & &
BOKRR & \\
Property & Relative Permeability (Oil in the Theta- Direction) & OKRT- &
& & BOKRT- & \\
Property & Relative Permeability (Oil in the Theta+ Direction) & OKRT &
& & BOKRT & \\
Property & Relative Permeability (Oil in the X- Direction) & OKRX- & & &
BOKRX- & Also BOKRI-. \\
Property & Relative Permeability (Oil in the X+ Direction) & OKRX & & &
BOKRX & Also BOKRI. \\
Property & Relative Permeability (Oil in the Y- Direction) & OKRY- & & &
BOKRY- & Also BOKRJ-. \\
Property & Relative Permeability (Oil in the Y+ Direction) & OKRY & & &
BOKRY & Also BOKRJ. \\
Property & Relative Permeability (Oil in the Z- Direction) & OKRZ- & & &
BOKRZ- & Also BOKRK-. \\
Property & Relative Permeability (Oil in the Z+ Direction) & OKRZ & & &
BOKRZ & Also BOKRK. \\
Property & Relative Permeability (Oil) & KRO & & & BKRO & Also BOKR. \\
Property & Relative Permeability (Water in the R- Direction) & WKRR- & &
& BWKRR- & \\
Property & Relative Permeability (Water in the R+ Direction) & WKRR & &
& BWKRR & \\
Property & Relative Permeability (Water in the Theta- Direction) & WKRT-
& & & BWKRT- & \\
Property & Relative Permeability (Water in the Theta+ Direction) & WKRT
& & & BWKRT & \\
Property & Relative Permeability (Water in the X- Direction) & WKRX- & &
& BWKRX- & Also BWKRI-. \\
Property & Relative Permeability (Water in the X+ Direction) & WKRX & &
& BWKRX & Also BWKRI. \\
Property & Relative Permeability (Water in the Y- Direction) & WKRY- & &
& BWKRY- & Also BWKRJ-. \\
Property & Relative Permeability (Water in the Y+ Direction) & WKRY & &
& BWKRY & Also BWKRJ. \\
Property & Relative Permeability (Water in the Z- Direction) & WKRZ- & &
& BWKRZ- & Also BWKRK-. \\
Property & Relative Permeability (Water in the Z+ Direction) & WKRZ & &
& BWKRZ & Also BWKRK. \\
Property & Relative Permeability (Water) & KRW & & & BKRW & Also
BWKR. \\
Property & Water Reservoir Density & WDEN & FWDEN & RWDEN & BWDEN & \\
Property & Water Reservoir Density & DENW & & & BDENW & Same as
BWDEN. \\
Property & Water Saturation & WSAT & FWSAT & RWSAT & BWSAT & Also
BSWAT. \\
Property & Water Saturation (Drainage to Imbibition) & WSHY & & & BWSHY
& Leaving water saturation for gas capillary pressure hysteresis. \\
Property & Water Saturation (Maximum Wetting) & WSMA & & & BWSMA &
Capillary pressure hysteresis maximum water saturation. \\
Property & Water Viscosity & WVIS & FWVIS & RWVIS & BWVIS & \\
Ratio & Gas-Oil Ratio & RS & FRS & RRS & BRS & \\
Ratio & Oil-Gas Ratio & RV & FRV & RRV & BRV & \\
Volumes & Gas In-Place & GIP & FGIP & RGIP & BGIP & Liquid and gas
phases. \\
Volumes & Gas In-Place (Gas Phase) & GIPG & FGIPG & RGIPG & BGIPG & \\
Volumes & Gas In-Place (Liquid Phase) & GIPL & FGIPL & RGIPL & BGIPL &
Only for cases that have dissolved gas in the water phase. \\
Volumes & Gas In-Place (Reservoir Conditions) & GIPR & FGIPR & & &
Compositional vector implemented in OPM Flow. \\
Volumes & Oil In-Place & OIP & FOIP & ROIP & BOIP & Liquid and gas
phases. \\
Volumes & Oil In-Place (Gas Phase) & OIPG & FOIPG & ROIPG & BOIPG & \\
Volumes & Oil In-Place (Liquid Phase) & OIPL & FOIPL & ROIPL & BOIPL
& \\
Volumes & Oil In-Place (Reservoir Conditions) & OIPR & FOIPR & & &
Compositional vector implemented in OPM Flow. \\
Volumes & Pore Volume (Gas) & GPV & FGPV & RGPV & BGPV & \\
Volumes & Pore Volume (HCPV) & HPV & FHPV & RHPV & BHPV & \\
Volumes & Pore Volume (Oil) & OPV & FOPV & ROPV & BOPV & \\
Volumes & Pore Volume (Reservoir Conditions) & RPV & FRPV & RRPV & BRPV
& Also BPORV. \\
Volumes & Pore Volume (Water) & WPV & FWPV & RWPV & BWPV & \\
Volumes & Rock Compaction (Trans. Multiplier X Direction) & PERMMDX & &
& BPERMMDX & Also BPERMOD. \\
Volumes & Rock Compaction (Trans. Multiplier Y Direction) & PERMMDY & &
& BPERMMDY & \\
Volumes & Rock Compaction (Trans. Multiplier Z Direction) & PERMMDZ & &
& BPERMMDZ & \\
Volumes & Rock Compaction (Trans. Multiplier) & RTM & FRTM & RRTM & BRTM
& \\
Volumes & Rock Compaction (Dual Porosity SIGMA Multiplier) & SIGMMOD & &
& BSIGMMOD & \\
Volumes & Rock Compaction (PV Multiplier) & PORVMOD & & & BPORVMOD & \\
Volumes & Water In-Place & WIP & FWIP & RWIP & BWIP & \\
Volumes & Water In-Place (Reservoir Conditions) & WIPR & FWIPR & & &
Compositional vector implemented in OPM Flow. \\
Notes: & & & & & & \\
\end{longtable}

\emph{Table 11.2.9: Block, Field, and Region Summary Variables}

\subsection{Field and Region Recovery
Variables}\label{field-and-region-recovery-variables}

Oil recovery factors (or efficiency) may be reported at both the field
and the region levels using a variety of parameters as summarized in
Table 11.2.10 and equations (11.2.1) to (11.2.6). In addition to the
actual recovery factors the simulator can also report the oil volumes
produced by various drive mechanisms as described in the table.

See also section 11.2.4 Field, Group, Well, Well Connection, and
Completion Variables that outlines the field production data summary
variables, and section 11.2.6 Field, Region, and Block Variables that
outlines the field, region and block production and injection data
summary variables.

\begin{longtable}[]{@{}
  >{\raggedright\arraybackslash}p{(\linewidth - 10\tabcolsep) * \real{0.1667}}
  >{\raggedright\arraybackslash}p{(\linewidth - 10\tabcolsep) * \real{0.1667}}
  >{\raggedright\arraybackslash}p{(\linewidth - 10\tabcolsep) * \real{0.1667}}
  >{\raggedright\arraybackslash}p{(\linewidth - 10\tabcolsep) * \real{0.1667}}
  >{\raggedright\arraybackslash}p{(\linewidth - 10\tabcolsep) * \real{0.1667}}
  >{\raggedright\arraybackslash}p{(\linewidth - 10\tabcolsep) * \real{0.1667}}@{}}
\toprule\noalign{}
\begin{minipage}[b]{\linewidth}\raggedright
Field and Region Summary Recovery Variables
\end{minipage} & \begin{minipage}[b]{\linewidth}\raggedright
\end{minipage} & \begin{minipage}[b]{\linewidth}\raggedright
\end{minipage} & \begin{minipage}[b]{\linewidth}\raggedright
\end{minipage} & \begin{minipage}[b]{\linewidth}\raggedright
\end{minipage} & \begin{minipage}[b]{\linewidth}\raggedright
\end{minipage} \\
\midrule\noalign{}
\endhead
\bottomrule\noalign{}
\endlastfoot
Type & Variable & Root & Field & Region & Comment \\
Recovery & Oil Recovery Efficiency (STOIIP) & OE & FOE & ROE & Based on
STOIIP. \\
Recovery & Oil Recovery Efficiency (Wells) & OEW & FOEW & ROEW & Based
on well production and STOIIP. \\
Recovery & Oil Recovery Efficiency (STOIIP Mobile with Respect to Water)
& OEIW & FOEIW & ROEIW & Based on STOIIP and initial mobile oil saturate
with respect to water. \\
Recovery & Oil Recovery Efficiency Wells Mobile with Respect to Water) &
OEWW & FOEWW & ROEWW & Based on well production and initial mobile oil
saturate with respect to water. \\
Recovery & Oil Recovery Efficiency (STOIIP Mobile with Respect to Gas) &
OEIG & FOEIG & ROEIG & Based on STOIIP and initial mobile oil saturate
with respect to gas. \\
Recovery & Oil Recovery Efficiency (Wells Mobile with Respect to Gas) &
OEWG & FOEWG & ROEWG & Based on well production and initial mobile oil
saturate with respect to gas. \\
Recovery & Oil Production Total (Free Gas Gas) & OTMF & FOTMF & ROTMF &
Volumes reported at stock tank conditions. \\
Recovery & Oil Production Total (Gas Influx) & ORMG & FORMG & RORMG & \\
Recovery & Oil Production Total (Oil Expansion) & ORME & FORME & RORME
& \\
Recovery & Oil Production Total (Rock Compaction) & ORMR & FORMR & RORMR
& \\
Recovery & Oil Production Total (Solution Gas) & ORMS & FORMS & RORMS
& \\
Recovery & Oil Production Total (Water Influx) & ORMW & FORMW & RORMW
& \\
Notes: & & & & & \\
\end{longtable}

\emph{Table 11.2.10: Field and Region Summary Recovery Variables}

These summary variables are computed as:

\begin{equation}\phantomsection\label{eq-11-2-1}{
\text{OE}=\frac{\left({\mathit{STOIIP}}_{t=0} - {\mathit{STOIIP}}_{t=i}\right)}{{\mathit{STOIIP}}_{t=0}}
}\end{equation}

\begin{equation}\phantomsection\label{eq-11-2-2}{
\text{OEW}=\frac{\left({\text{Well Oil Production}}_{t=i}\right)}{{\mathit{STOIIP}}_{t=0}}
}\end{equation}

\begin{equation}\phantomsection\label{eq-11-2-3}{
\text{OEIW}=\frac{\left({\mathit{STOIIP}}_{t=0} - {\mathit{STOIIP}}_{t=i}\right)}{\left({\text{Initial Mobile Oil Saturation with respect to Water}}_{t=0}\right)}
}\end{equation}

\begin{equation}\phantomsection\label{eq-11-2-4}{
\text{OEWW}=\frac{\left({\text{Well Oil Production}}_{t=i}\right)}{\left({\text{Initial Mobile Oil Saturation with respect to Water}}_{t=0}\right)}
}\end{equation}

\begin{equation}\phantomsection\label{eq-11-2-5}{
\text{OEIG}=\frac{\left({\mathit{STOIIP}}_{t=0} - {\mathit{STOIIP}}_{t=i}\right)}{\left({\text{Initial Mobile Oil Saturation with respect to Gas}}_{t=0}\right)}
}\end{equation}

\begin{equation}\phantomsection\label{eq-11-2-6}{
\text{OEWG}=\frac{\left({\text{Well Oil Production}}_{t=i}\right)}{\left({\text{Initial Mobile Oil Saturation with respect to Gas}}_{t=0}\right)}
}\end{equation}

In addition to all the aforementioned summary variables, there are
variables associated with specific features that may have been enabled
in the input deck, for instance, the Brine Model, or the Polymer Model.
The next set of sections describe the variables associated with these
specific options.

\subsection{Option Specific Variables -- Alkaline
Model}\label{option-specific-variables-alkaline-model}

This feature has not been implemented in OPM Flow.

\subsection{Option Specific Variables -- API and Tracer
Tracking}\label{option-specific-variables-api-and-tracer-tracking}

The summary variables in this section are for use for when either the
API tracking model has been activated by the API keyword in the RUNSPEC
section so that the various ``oil types'' are tracked in the model, or
for when the Tracer Model has been requested via the TRACERS keyword in
the RUNSPEC section, that defines the number of tracers in the model and
the various passive tracer tracking options. Currently, only the Tracer
model is implemented in OPM Flow.

\begin{longtable}[]{@{}
  >{\raggedright\arraybackslash}p{(\linewidth - 16\tabcolsep) * \real{0.1111}}
  >{\raggedright\arraybackslash}p{(\linewidth - 16\tabcolsep) * \real{0.1111}}
  >{\raggedright\arraybackslash}p{(\linewidth - 16\tabcolsep) * \real{0.1111}}
  >{\raggedright\arraybackslash}p{(\linewidth - 16\tabcolsep) * \real{0.1111}}
  >{\raggedright\arraybackslash}p{(\linewidth - 16\tabcolsep) * \real{0.1111}}
  >{\raggedright\arraybackslash}p{(\linewidth - 16\tabcolsep) * \real{0.1111}}
  >{\raggedright\arraybackslash}p{(\linewidth - 16\tabcolsep) * \real{0.1111}}
  >{\raggedright\arraybackslash}p{(\linewidth - 16\tabcolsep) * \real{0.1111}}
  >{\raggedright\arraybackslash}p{(\linewidth - 16\tabcolsep) * \real{0.1111}}@{}}
\toprule\noalign{}
\begin{minipage}[b]{\linewidth}\raggedright
API and Tracer Tracking Summary Variables
\end{minipage} & \begin{minipage}[b]{\linewidth}\raggedright
\end{minipage} & \begin{minipage}[b]{\linewidth}\raggedright
\end{minipage} & \begin{minipage}[b]{\linewidth}\raggedright
\end{minipage} & \begin{minipage}[b]{\linewidth}\raggedright
\end{minipage} & \begin{minipage}[b]{\linewidth}\raggedright
\end{minipage} & \begin{minipage}[b]{\linewidth}\raggedright
\end{minipage} & \begin{minipage}[b]{\linewidth}\raggedright
\end{minipage} & \begin{minipage}[b]{\linewidth}\raggedright
\end{minipage} \\
\midrule\noalign{}
\endhead
\bottomrule\noalign{}
\endlastfoot
Type & Variable & Root & Field & Group & Well & Well Connection & Region
& Block \\
Flow & Oil API & API & FAPI & GAPI & WAPI & CAPI & RAPI & BAPI \\
Flow & Tracer Flow Rate & TFR & & & & CTFR & & \\
Flow & Tracer Flow Total (Inter Region) & TFTT & & & & & RTFTT & \\
Flow & Tracer Flow Total (Free Inter Region) & TFTF & & & & & RTFTF & \\
Flow & Tracer Flow Total (Solution Inter Region) & TFTS & & & & & RTFTS
& \\
Flow & Tracer Injection Concentration & TIC & FTIC & GTIC & WTIC & CTIC
& & \\
Flow & Tracer Injection Concentration (Free) & TICF & FTICF & GTICF &
WTICF & CTICF & & \\
Flow & Tracer Injection Concentration (Solution) & TICS & FTICS & GTICS
& WTICS & CTICS & & \\
Flow & Tracer Injection Rate & TIR & FTIR & GTIR & WTIR & CTIR & & \\
Flow & Tracer Injection Rate (Free) & TIRF & FTIRF & GTIRF & WTIRF &
CTIRF & & \\
Flow & Tracer Injection Rate (Solution) & TIRS & FTIRS & GTIRS & WTIRS &
CTIRS & & \\
Flow & Tracer Injection Total & TIT & FTIT & GTIT & WTIT & CTIT & & \\
Flow & Tracer Injection Total (Free) & TITF & FTITF & GTITF & WTITF &
CTITF & & \\
Flow & Tracer Injection Total (Solution) & TITS & FTITS & GTITS & WTITS
& CTITS & & \\
Flow & Tracer Production Concentration & TPC & FTPC & GTPC & WTPC & CTPC
& & \\
Flow & Tracer Production Concentration (Free) & TPCF & FTPCF & GTPCF &
WTPCF & CTPCF & & \\
Flow & Tracer Production Concentration (Solution) & TPCS & FTPCS & GTPCS
& WTPCS & CTPCS & & \\
Flow & Tracer Production Rate & TPR & FTPR & GTPR & WTPR & CTPR & & \\
Flow & Tracer Production Rate (Free) & TPRF & FTPRF & GTPRF & WTPRF &
CTPRF & & \\
Flow & Tracer Production Rate (Solution) & TPRS & FTPRS & GTPRS & WTPRS
& CTPRS & & \\
Flow & Tracer Production Total & TPT & FTPT & GTPT & WTPT & CTPT & & \\
Flow & Tracer Production Total (Free) & TPTF & FTPTF & GTPTF & WTPTF &
CTPTF & & \\
Flow & Tracer Production Total (Solution) & TPTS & FTPTS & GTPTS & WTPTS
& CTPTS & & \\
Volume & Tracer Concentration & TCN & & & & & & BTCN \\
Volume & Tracer Concentration (Free) & TCNF & & & & & & BTCNF \\
Volume & Tracer Concentration (Solution) & TCNS & & & & & & BTCNS \\
Volume & Tracer In-Place & TIPT & FTIPT & & & & RTIPT & BTIPT \\
Volume & Tracer In-Place (Free) & TIPF & FTIPF & & & & RTIPF & BTIPF \\
Volume & Tracer In-Place (Solution) & TIPS & FTIPS & & & & RTIPS &
BTIPS \\
Notes: & & & & & & & & \\
\end{longtable}

\emph{Table 11.2.11: API and Tracer Tracking Summary Variables}

Note that for the tracer summary variables the mnemonic should be
concatenated with the tracer name as defined by the TRACER keyword in
the PROPS section. This keyword defines a series of passive tracers that
are associated with a phase (oil, water, or gas) in the model, see the
example for reference.

\subsubsection{Example}\label{example-490}

In the PROPS section the TRACER keyword defines four passive tracers one
for a gas injection well, one for tracking the dissolved gas, and two to
track the injected water from two water injection wells.

\begin{verbatim}
-- ==============================================================================
--
-- PROPS SECTION
--
-- ==============================================================================
PROPS
--
--       DEFINE TRACER NAMES
--
--       TRACER   TRACER
--       NAME     PHASE
--       ------   ------
TRACER
        'IGS'     'GAS'                                    / GAS INJECTOR
        'DGS'     'GAS'                                    / DISOLVED GAS
        'IW1'     'WAT'                                    / WAT INJECTOR 1
        'Iw2'     'WAT'                                    / WAT INJECTOR 2
/
\end{verbatim}

If we wish to report the field total amount of the tracers in-place we
would use the following in the SUMMARY section:

\begin{verbatim}
-- ==============================================================================
--
-- SUMMARY SECTION
--
-- ==============================================================================
SUMMARY

FTIPTIGS
FTIPTDGS
FTIPTIW1
FTIPTIW2
\end{verbatim}

Note that there is no terminating ``/'' for field summary variables.

\subsection{Option Specific Variables -- Asphaltene
Model}\label{option-specific-variables-asphaltene-model}

This feature has not been implemented in OPM Flow.

\subsection{Option Specific Variables -- Brine
Model}\label{option-specific-variables-brine-model}

The available Brine model summary keywords are dependent on which form
of the Brine model has been invoked in the simulation input deck. In the
RUNSPEC section the BRINE keyword activates the standard Brine Tracking
model and optionally defines the water phase to have various salinities
if the ECLMC keyword in the RUNSPEC section has been used to activate
the Multi-Component Brine model. The Multi-Component Brine model allows
for the water phase to have multiple water salinities.

Note that the Multi-Component Brine model is currently not supported by
OPM Flow.

If OPM Flow's Salt Precipitation Model has been activated in the input
deck via the PRECSALT keyword in the RUNSPEC section, then the Brine
model summary keywords can also be employed with this option.

\begin{longtable}[]{@{}
  >{\raggedright\arraybackslash}p{(\linewidth - 18\tabcolsep) * \real{0.1000}}
  >{\raggedright\arraybackslash}p{(\linewidth - 18\tabcolsep) * \real{0.1000}}
  >{\raggedright\arraybackslash}p{(\linewidth - 18\tabcolsep) * \real{0.1000}}
  >{\raggedright\arraybackslash}p{(\linewidth - 18\tabcolsep) * \real{0.1000}}
  >{\raggedright\arraybackslash}p{(\linewidth - 18\tabcolsep) * \real{0.1000}}
  >{\raggedright\arraybackslash}p{(\linewidth - 18\tabcolsep) * \real{0.1000}}
  >{\raggedright\arraybackslash}p{(\linewidth - 18\tabcolsep) * \real{0.1000}}
  >{\raggedright\arraybackslash}p{(\linewidth - 18\tabcolsep) * \real{0.1000}}
  >{\raggedright\arraybackslash}p{(\linewidth - 18\tabcolsep) * \real{0.1000}}
  >{\raggedright\arraybackslash}p{(\linewidth - 18\tabcolsep) * \real{0.1000}}@{}}
\toprule\noalign{}
\begin{minipage}[b]{\linewidth}\raggedright
Brine and Multi-Component Brine Model Summary Variables
\end{minipage} & \begin{minipage}[b]{\linewidth}\raggedright
\end{minipage} & \begin{minipage}[b]{\linewidth}\raggedright
\end{minipage} & \begin{minipage}[b]{\linewidth}\raggedright
\end{minipage} & \begin{minipage}[b]{\linewidth}\raggedright
\end{minipage} & \begin{minipage}[b]{\linewidth}\raggedright
\end{minipage} & \begin{minipage}[b]{\linewidth}\raggedright
\end{minipage} & \begin{minipage}[b]{\linewidth}\raggedright
\end{minipage} & \begin{minipage}[b]{\linewidth}\raggedright
\end{minipage} & \begin{minipage}[b]{\linewidth}\raggedright
\end{minipage} \\
\midrule\noalign{}
\endhead
\bottomrule\noalign{}
\endlastfoot
Type & Variable & Root & Field & Group & Well & Well Connection & Region
& Block & Comment \\
Flow & Salt Flow Rate & SFR & & & & CSFR & & & Brine and Salt
Precipitation Models. \\
Flow & Salt Injection Rate & SIR & FSIR & GSIR & WSIR & CSIR & & & \\
Flow & Salt Injection Total & SIT & FSIT & GSIT & WSIT & CSIT & & & \\
Flow & Salt Production Rate & SPR & FSPR & GSPR & WSPR & CSPR & & & \\
Flow & Salt Production Total & SPT & FSPT & GSPT & WSPT & CSPT & & & \\
Property & Salt Corrected Water Viscosity & EMV\_SAL & & & & & &
BEMV\_SAL & \\
Volume & Salt Concentration (Block) & SCN & & & & & & BSCN & \\
Volume & Salt Concentration (Injection) & SIC & FSIC & GSIC & WSIC &
CSIC & & & \\
Volume & Salt Concentration (Production) & SPC & FSPC & GSPC & WSPC &
CSPC & & & \\
Volume & Salt In-Place & SIP & FSIP & & & & RSIP & BSIP & \\
Flow & Anion Concentration (Block) & TCNFANI & & & & & & BTCNFANI &
Multi-Component Brine Model. \\
Flow & Anion Flow Rate & TFRANI & & & & CTFRANI & & & \\
Flow & Anion Injection Rate & TIRANI & FTIRANI & GTIRANI & WTIRANI & & &
& \\
Flow & Anion Injection Total & TITANI & FTITANI & GTITANI & WTITANI &
CTITANI & & & \\
Flow & Anion Production Rate & TPRANI & FTPRANI & GTPRANI & WTPRANI & &
& & \\
Flow & Anion Production Total & TPTANI & FTPTANI & GTPTANI & WTPTANI &
CTPTANI & & & \\
Flow & Cation Flow Rate & TFRCAT & & & & CTFRCAT & & & \\
Flow & Cation Injection Rate & TIRCAT & FTIRCAT & GTIRCAT & WTIRCAT & &
& & \\
Flow & Cation Injection Total & TITCAT & FTITCAT & GTITCAT & WTITCAT &
CTITCAT & & & \\
Flow & Cation Production Rate & TPRCAT & FTPRCAT & GTPRCAT & WTPRCAT & &
& & \\
Flow & Cation Production Total & TPTCAT & FTPTCAT & GTPTCAT & WTPTCAT &
CTPTCAT & & & \\
Property & Effective Salinity (Polymer) & ESALPLY & & & & & & BESALPLY
& \\
Property & Effective Salinity (Surfactant) & ESALSUR & & & & & &
BESALSUR & \\
Volume & Cation Concentration (Block) & TCNFCAT & & & & & & BTCNFCAT
& \\
Volume & Cation Concentration (Rock Associated) & TRADCAT & & & & & &
BTRADCAT & \\
Volume & Cation Concentration (Surfactant Associated) & TSADCAT & & & &
& & BTSADCAT & \\
Notes: & & & & & & & & & \\
\end{longtable}

\emph{Table 11.2.12: Brine and Multi-Component Brine Model Summary
Variables}

Note that not all these variables are available in OPM Flow; however,
the simulator will issue warning messages if this is indeed the case. It
is anticipated that the number of recognized summary variables will
increase in future releases of OPM Flow.

\subsection{Option Specific Variables - CO2STORE/H2STORE
Model}\label{option-specific-variables---co2storeh2store-model}

The variables in this section are for OPM Flow's black-oil
CO2STORE/\hyperref[REF_HEADING_KEYWORD_H2STORE]{H2STORE} Model, and are
based on the commercial simulator's compositional vectors (GCDI, GCDM,
and WCD series), as the model is only available in the commercial
simulator's compositional simulator.

\begin{longtable}[]{@{}
  >{\raggedright\arraybackslash}p{(\linewidth - 14\tabcolsep) * \real{0.1250}}
  >{\raggedright\arraybackslash}p{(\linewidth - 14\tabcolsep) * \real{0.1250}}
  >{\raggedright\arraybackslash}p{(\linewidth - 14\tabcolsep) * \real{0.1250}}
  >{\raggedright\arraybackslash}p{(\linewidth - 14\tabcolsep) * \real{0.1250}}
  >{\raggedright\arraybackslash}p{(\linewidth - 14\tabcolsep) * \real{0.1250}}
  >{\raggedright\arraybackslash}p{(\linewidth - 14\tabcolsep) * \real{0.1250}}
  >{\raggedright\arraybackslash}p{(\linewidth - 14\tabcolsep) * \real{0.1250}}
  >{\raggedright\arraybackslash}p{(\linewidth - 14\tabcolsep) * \real{0.1250}}@{}}
\toprule\noalign{}
\begin{minipage}[b]{\linewidth}\raggedright
CO2STORE/H2STORE Field, Group and Well Summary Variables
\end{minipage} & \begin{minipage}[b]{\linewidth}\raggedright
\end{minipage} & \begin{minipage}[b]{\linewidth}\raggedright
\end{minipage} & \begin{minipage}[b]{\linewidth}\raggedright
\end{minipage} & \begin{minipage}[b]{\linewidth}\raggedright
\end{minipage} & \begin{minipage}[b]{\linewidth}\raggedright
\end{minipage} & \begin{minipage}[b]{\linewidth}\raggedright
\end{minipage} & \begin{minipage}[b]{\linewidth}\raggedright
\end{minipage} \\
\midrule\noalign{}
\endhead
\bottomrule\noalign{}
\endlastfoot
Type & Variable & Root & Field & Group & Well & Well Connection & Well
Completion \\
Flow & CO2/H2 Mass Injection Rate & GMIR & FGMIR & GGMIR & WGMIR & CGMIR
& CGMIRL \\
Flow & CO2/H2 Mass Injection Total & GMIT & FGMIT & GGMIT & WGMIT &
CGMIT & CGMITL \\
Flow & CO2/H2 Mass Production Rate & GMPR & FGMPR & GGMPR & WGMPR &
CGMPR & CGMPRL \\
Flow & CO2/H2 Mass Production Total & GMPT & FGMPT & GGMPT & WGMPT &
CGMPT & CGMPTL \\
Notes: & & & & & & & \\
\end{longtable}

\emph{Table 11.2.13: CO2STORE/H2STORE Field, Group and Well Summary
Variables}

\begin{longtable}[]{@{}
  >{\raggedright\arraybackslash}p{(\linewidth - 10\tabcolsep) * \real{0.1667}}
  >{\raggedright\arraybackslash}p{(\linewidth - 10\tabcolsep) * \real{0.1667}}
  >{\raggedright\arraybackslash}p{(\linewidth - 10\tabcolsep) * \real{0.1667}}
  >{\raggedright\arraybackslash}p{(\linewidth - 10\tabcolsep) * \real{0.1667}}
  >{\raggedright\arraybackslash}p{(\linewidth - 10\tabcolsep) * \real{0.1667}}
  >{\raggedright\arraybackslash}p{(\linewidth - 10\tabcolsep) * \real{0.1667}}@{}}
\toprule\noalign{}
\begin{minipage}[b]{\linewidth}\raggedright
CO2STORE/H2STORE Field, Region, and Block Summary Variables
\end{minipage} & \begin{minipage}[b]{\linewidth}\raggedright
\end{minipage} & \begin{minipage}[b]{\linewidth}\raggedright
\end{minipage} & \begin{minipage}[b]{\linewidth}\raggedright
\end{minipage} & \begin{minipage}[b]{\linewidth}\raggedright
\end{minipage} & \begin{minipage}[b]{\linewidth}\raggedright
\end{minipage} \\
\midrule\noalign{}
\endhead
\bottomrule\noalign{}
\endlastfoot
Type & Variable & Root & Field & Region & Block \\
Mass & CO2/H2 In-place & GMIP & FGMIP & RGMIP & BGMIP \\
Mass & CO2/H2 In-place (Gas Phase) & GMGP & FGMGP & RGMGP & BGMGP \\
Mass & CO2/H2 In-place (Liquid Phase) & GMDS & FGMDS & RGMDS & BGMDS \\
Mass & CO2/H2 Effectively Residual Trapped (`Stranded') in the gas phase
& GMST & FGMST & RGMST & BGMST \\
Mass & CO2/H2 Effectively Residual Untrapped (`Unstranded') in the gas
phase & GMUS & FGMUS & RGMUS & BGMUS \\
Mass & CO2/H2 Maximum Potential Residual Trapped in the gas phase & GMTR
& FGMTR & RGMTR & BGMTR \\
Mass & CO2/H2 Maximum Potential Residual Untrapped in the gas phase &
GMMO & FGMMO & RGMMO & BGMMO \\
Mass & CO2/H2 Immobile in the gas phase (Non-wetting relative
permeability equals zero) & GKTR & FGKTR & RGKTR & BGKTR \\
Mass & CO2/H2 Mobile in the gas phase (Non-wetting relative permeability
greater than zero) & GKMO & FGKMO & RGKMO & BGKMO \\
Moles & CO2/H2 Immobile in the gas phase & GCDI & FGCDI & RGCDI &
BGCDI \\
Moles & CO2/H2 Mobile in the gas phase & GCDM & FGCDM & RGCDM & BGCDM \\
Moles & CO2/H2 Immobile in the gas phase (Non-wetting relative
permeability equals zero) & GKDI & FGKDI & RGKDI & BGKDI \\
Moles & CO2/H2 Mobile in the gas phase (Non-wetting relative
permeability greater than zero) & GKDM & FGKDM & RGKDM & BGKDM \\
Moles & CO2/H2 Dissolved in the water phase & WCD & FWCD & RWCD &
BWCD \\
Volume & Water in the gas phase & WIPG & FWIPG & RWIPG & BWIPG \\
Volume & Water in the water phase & WIPL & FWIPL & RWIPL & BWIPL \\
Notes: & & & & & \\
\end{longtable}

\emph{Table 11.2.13: CO2STORE/H2STORE Field, Region, and Block Summary
Variables}

See the OpenGoSim documentation
(https://docs.opengosim.com/cirrus/theory/residual\_and\_stranded.html)
for more details of how the trapped gas saturation and the stranded gas
saturation contribute to the maximum potential residual trapping (GMTR)
and the effective residual trapping (GMST) of CO2 in saline aquifers
when relative permeability hysteresis is modelled.

The 11th Society of Petroleum Engineers Comparative Solution Project
(http://spe.org/csp) defines the ``immobile free-phase CO2'' as ``CO2 at
saturations for which the non-wetting phase relative permeability equals
zero'' (GKDI, GKTR).

Note that the GK* and GM* series of variables are OPM Flow specific
variables.

Note that the WIPG and WIPL series of variables are OPM Flow specific
variables.

\subsection{Option Specific Variables -- Coal Bed Methane
Model}\label{option-specific-variables-coal-bed-methane-model}

The Coal Bed Methane model is not supported by OPM Flow.

\subsection{Option Specific Variables -- Conductive Faults
Model}\label{option-specific-variables-conductive-faults-model}

Conductive faults have not been implemented in OPM Flow.

\subsection{Option Specific Variables -- Environmental
Tracers}\label{option-specific-variables-environmental-tracers}

Environmental tracers are not supported by OPM Flow.

\subsection{Option Specific Variables -- Formation Damage
Model}\label{option-specific-variables-formation-damage-model}

The variables in this section are for the OPM Flow's Formation Damage
Model. The \hyperref[REF_HEADING_KEYWORD_WINJDAM]{WINJDAM} keyword sets
up the filter cake properties for the specific water injection well in
the SCHEDULE section.

\begin{longtable}[]{@{}
  >{\raggedright\arraybackslash}p{(\linewidth - 16\tabcolsep) * \real{0.1111}}
  >{\raggedright\arraybackslash}p{(\linewidth - 16\tabcolsep) * \real{0.1111}}
  >{\raggedright\arraybackslash}p{(\linewidth - 16\tabcolsep) * \real{0.1111}}
  >{\raggedright\arraybackslash}p{(\linewidth - 16\tabcolsep) * \real{0.1111}}
  >{\raggedright\arraybackslash}p{(\linewidth - 16\tabcolsep) * \real{0.1111}}
  >{\raggedright\arraybackslash}p{(\linewidth - 16\tabcolsep) * \real{0.1111}}
  >{\raggedright\arraybackslash}p{(\linewidth - 16\tabcolsep) * \real{0.1111}}
  >{\raggedright\arraybackslash}p{(\linewidth - 16\tabcolsep) * \real{0.1111}}
  >{\raggedright\arraybackslash}p{(\linewidth - 16\tabcolsep) * \real{0.1111}}@{}}
\toprule\noalign{}
\begin{minipage}[b]{\linewidth}\raggedright
Formation Damage Model Summary Variables
\end{minipage} & \begin{minipage}[b]{\linewidth}\raggedright
\end{minipage} & \begin{minipage}[b]{\linewidth}\raggedright
\end{minipage} & \begin{minipage}[b]{\linewidth}\raggedright
\end{minipage} & \begin{minipage}[b]{\linewidth}\raggedright
\end{minipage} & \begin{minipage}[b]{\linewidth}\raggedright
\end{minipage} & \begin{minipage}[b]{\linewidth}\raggedright
\end{minipage} & \begin{minipage}[b]{\linewidth}\raggedright
\end{minipage} & \begin{minipage}[b]{\linewidth}\raggedright
\end{minipage} \\
\midrule\noalign{}
\endhead
\bottomrule\noalign{}
\endlastfoot
Type & Variable & Root & Field & Group & Well & Well Connection & Region
& Block \\
Flow & Filtrate Injection Concentration & INJFC & & & WINFC & & & \\
Flow & Filtrate Volume Injection Rate & INJFVR & & & WINJFVR & CINJFVR &
& \\
Flow & Filtrate Volume Injection Total & INJFVT & & & WINJFVT & CINJFVT
& & \\
Property & Filter Cake Area Of Flow & FCAOF & & & & CFCAOF & & \\
Property & Filter Cake Permeability & FCPERM & & & & CFCPERM & & \\
Property & Filter Cake Porosity & FCPORO & & & & CFCPORO & & \\
Property & Filter Cake Radius & FCRAD & & & & CFCRAD & & \\
Property & Filter Cake Skin & FCSKIN & & & & CFCSKIN & & \\
Property & Filter Cake Width & FCWIDTH & & & & CFCWIDTH & & \\
Notes: & & & & & & & & \\
\end{longtable}

\emph{Table 11.2.14: Formation Damage Model Summary Variables}

Note that these are OPM Flow specific variables.

\subsection{Option Specific Variables -- Foam Model
Tracers}\label{option-specific-variables-foam-model-tracers}

The Foam phase and model are activated via the FOAM keyword in the
RUNSPEC section. Note in the commercial simulator the FOAM phase and
model can be used in conjunction with the POLYMER and SURFACT phases;
this is not the case for OPM Flow. OPM Flow's FOAM phase and model is a
standalone implementation and cannot be used in conjunction with the
either the POLYMER or SURFACT phases.

\begin{longtable}[]{@{}
  >{\raggedright\arraybackslash}p{(\linewidth - 18\tabcolsep) * \real{0.1000}}
  >{\raggedright\arraybackslash}p{(\linewidth - 18\tabcolsep) * \real{0.1000}}
  >{\raggedright\arraybackslash}p{(\linewidth - 18\tabcolsep) * \real{0.1000}}
  >{\raggedright\arraybackslash}p{(\linewidth - 18\tabcolsep) * \real{0.1000}}
  >{\raggedright\arraybackslash}p{(\linewidth - 18\tabcolsep) * \real{0.1000}}
  >{\raggedright\arraybackslash}p{(\linewidth - 18\tabcolsep) * \real{0.1000}}
  >{\raggedright\arraybackslash}p{(\linewidth - 18\tabcolsep) * \real{0.1000}}
  >{\raggedright\arraybackslash}p{(\linewidth - 18\tabcolsep) * \real{0.1000}}
  >{\raggedright\arraybackslash}p{(\linewidth - 18\tabcolsep) * \real{0.1000}}
  >{\raggedright\arraybackslash}p{(\linewidth - 18\tabcolsep) * \real{0.1000}}@{}}
\toprule\noalign{}
\begin{minipage}[b]{\linewidth}\raggedright
Foam Model Tracers Summary Variables
\end{minipage} & \begin{minipage}[b]{\linewidth}\raggedright
\end{minipage} & \begin{minipage}[b]{\linewidth}\raggedright
\end{minipage} & \begin{minipage}[b]{\linewidth}\raggedright
\end{minipage} & \begin{minipage}[b]{\linewidth}\raggedright
\end{minipage} & \begin{minipage}[b]{\linewidth}\raggedright
\end{minipage} & \begin{minipage}[b]{\linewidth}\raggedright
\end{minipage} & \begin{minipage}[b]{\linewidth}\raggedright
\end{minipage} & \begin{minipage}[b]{\linewidth}\raggedright
\end{minipage} & \begin{minipage}[b]{\linewidth}\raggedright
\end{minipage} \\
\midrule\noalign{}
\endhead
\bottomrule\noalign{}
\endlastfoot
Type & Variable & Root & Field & Group & Well & Well Connection & Region
& Block & Comment \\
Flow & Foam Flow Rate & TFRFOA & & & & CTFRFOA & & & \\
Flow & Foam Injection Rate & TIRFOA & FTIRFOA & GTIRFOA & WTIRFOA & & &
& \\
Flow & Foam Injection Total & TITFOA & FTITFOA & GTITFOA & WTITFOA &
CTITFOA & & & \\
Flow & Foam Production Rate & TPRFOA & FTPRFOA & GTPRFOA & WTPRFOA & & &
& \\
Flow & Foam Production Total & TPTFOA & FTPTFOA & GTPTFOA & WTPTFOA &
CTPTFOA & & & \\
Property & Foam Adsorption Total & TADSFOA & FTADSFOA & & & CTADSFOA &
RTADSFOA & BTADSFOA & Block is current value. \\
Property & Foam Capillary Number & TCNMFOA & & & & & & BTCNMFOA & \\
Property & Foam Decayed Half Life & THLFFOA & & & & & & BTHLFFOA & \\
Property & Foam Decayed Tracer & TDCYFOA & FTDCYFOA & & & & RTDCYFOA &
BTDCYFOA & \\
Property & Foam Gas Mobility Factor & TMOBFOA & FTMOBFOA & & & & & &
Excluding shear. \\
Property & Foam in Solution & TIPTFOA & FTIPTFOA & & & CTIPTFOA &
RTIPTFOA & BTIPTFOA & \\
Volume & Foam Concentration & TCNFFOA & & & & & & BTCNFFOA & \\
Notes: & & & & & & & & & \\
\end{longtable}

\emph{Table 11.2.15: Foam Model Tracers Summary Variables}

Note that not all these variables are available in OPM Flow; however,
the simulator will issue warning messages if this is indeed the case. It
is anticipated that the number of recognized summary variables will
increase in future releases of OPM Flow.

\subsection{Option Specific Variables -- Gas Field Operations
Model}\label{option-specific-variables-gas-field-operations-model}

The Gas Field Operations model has not been implemented in OPM Flow.

\subsection{Option Specific Variables -- Gas Lift Optimization
Model}\label{option-specific-variables-gas-lift-optimization-model}

For the Gas Lift Optimization model only a few SUMMARY vectors are
available, both in OPM Flow and the commercial simulator, as depicted in
Table 11.2.16. Gas lift optimization is activated by the LIFTOPT keyword
in the SCHEDULE section as gas lift activation can vary through time
depending on the flow characteristics of the wells.

\begin{longtable}[]{@{}
  >{\raggedright\arraybackslash}p{(\linewidth - 16\tabcolsep) * \real{0.1111}}
  >{\raggedright\arraybackslash}p{(\linewidth - 16\tabcolsep) * \real{0.1111}}
  >{\raggedright\arraybackslash}p{(\linewidth - 16\tabcolsep) * \real{0.1111}}
  >{\raggedright\arraybackslash}p{(\linewidth - 16\tabcolsep) * \real{0.1111}}
  >{\raggedright\arraybackslash}p{(\linewidth - 16\tabcolsep) * \real{0.1111}}
  >{\raggedright\arraybackslash}p{(\linewidth - 16\tabcolsep) * \real{0.1111}}
  >{\raggedright\arraybackslash}p{(\linewidth - 16\tabcolsep) * \real{0.1111}}
  >{\raggedright\arraybackslash}p{(\linewidth - 16\tabcolsep) * \real{0.1111}}
  >{\raggedright\arraybackslash}p{(\linewidth - 16\tabcolsep) * \real{0.1111}}@{}}
\toprule\noalign{}
\begin{minipage}[b]{\linewidth}\raggedright
Gas Lift Optimization Model Summary Variables
\end{minipage} & \begin{minipage}[b]{\linewidth}\raggedright
\end{minipage} & \begin{minipage}[b]{\linewidth}\raggedright
\end{minipage} & \begin{minipage}[b]{\linewidth}\raggedright
\end{minipage} & \begin{minipage}[b]{\linewidth}\raggedright
\end{minipage} & \begin{minipage}[b]{\linewidth}\raggedright
\end{minipage} & \begin{minipage}[b]{\linewidth}\raggedright
\end{minipage} & \begin{minipage}[b]{\linewidth}\raggedright
\end{minipage} & \begin{minipage}[b]{\linewidth}\raggedright
\end{minipage} \\
\midrule\noalign{}
\endhead
\bottomrule\noalign{}
\endlastfoot
Type & Variable & Root & Field & Group & Well & Well Connection & Region
& Block \\
Flow & Gas lift gas rate, based on ALQ & GLIR & FGLIR & GGLIR & WGLIR &
& & \\
Flow & Incremental oil rate per unit of incremental gas lift gas
quantity, that is the well's gradient as defined by equation (11.2.7)
below & OGLR & & & WOGLR & & & \\
Notes: & & & & & & & & \\
\end{longtable}

\emph{Table 11.2.16: Gas Lift Optimization Model Summary Variables}

The incremental gas lift gas supply is allocated to a well based on the
well's current incremental gradient using the following formulae:

\begin{equation}\phantomsection\label{eq-11-2-7}{
\mathit{Gradient} = \left(\frac{\left({β}_{w}\times \mathrm{Δ}{Q}_{\mathit{Oil}}\right)}{\left(\mathrm{Δ}{Q}_{\mathit{GasLift}} + {β}_{g}\times \mathrm{Δ}{Q}_{\mathit{Gas}}\right)}\right)
}\end{equation}

where

βw = a weighting factor for the preferential allocation of lift gas,

βg = the gas production rate weighting factor,

∆QOil = the increment/decrement in oil production rate,

∆QGas = the increment/decrement in gas production rate, and

∆QGasLift = the increment/decrement in gas lift gas rate.

Normally βg is defaulted to zero and in which case equation (11.2.7)
reduces to:

\begin{equation}\phantomsection\label{eq-11-2-8}{
\mathit{Gradient} = \left(\frac{\left({β}_{w}\times \mathrm{Δ}{Q}_{\mathit{Oil}}\right)}{\mathrm{Δ}{Q}_{\mathit{GasLift}}}\right)
}\end{equation}

Note that if gas lift optimization has been activated in the model then
the gas production SUMMARY variables (FGPR, GGPR, WGPR, FGPT, etc.) only
contain the produced gas volumes, that is the reported values exclude
the gas associated with gas lift gas.

Secondly, the well gas lift rates are adjusted by the well efficiencies
in the calculating the group and field gas lift rates.

\subsection{Option Specific Variables -- Gas Calorific Value
Reporting}\label{option-specific-variables-gas-calorific-value-reporting}

This feature has not been implemented in OPM Flow.

\subsection{Option Specific Variables -- Gi Psuedo-Compositional
Model}\label{option-specific-variables-gi-psuedo-compositional-model}

This feature has not been implemented in OPM Flow.

\subsection{Option Specific Variables -- LGR and
Coarsening}\label{option-specific-variables-lgr-and-coarsening}

Local Grid Refinement (``LGR'') and cell coarsening have not been
implemented in OPM Flow.

\subsection{Option Specific Variables -- Biofilm And MICP
Model}\label{option-specific-variables-biofilm-and-micp-model}

The MICP keyword in the RUNSPEC section activates the Microbially
Induced Calcite Precipitation (``MICP'') model for the run. The summary
variables for this option are outlined in Table 11.2.33. Summary
variables related to microbial and biofilm quantities can be also used
with the \hyperref[REF_HEADING_KEYWORD_BIOFILM]{BIOFILM} model.

\begin{longtable}[]{@{}
  >{\raggedright\arraybackslash}p{(\linewidth - 18\tabcolsep) * \real{0.1000}}
  >{\raggedright\arraybackslash}p{(\linewidth - 18\tabcolsep) * \real{0.1000}}
  >{\raggedright\arraybackslash}p{(\linewidth - 18\tabcolsep) * \real{0.1000}}
  >{\raggedright\arraybackslash}p{(\linewidth - 18\tabcolsep) * \real{0.1000}}
  >{\raggedright\arraybackslash}p{(\linewidth - 18\tabcolsep) * \real{0.1000}}
  >{\raggedright\arraybackslash}p{(\linewidth - 18\tabcolsep) * \real{0.1000}}
  >{\raggedright\arraybackslash}p{(\linewidth - 18\tabcolsep) * \real{0.1000}}
  >{\raggedright\arraybackslash}p{(\linewidth - 18\tabcolsep) * \real{0.1000}}
  >{\raggedright\arraybackslash}p{(\linewidth - 18\tabcolsep) * \real{0.1000}}
  >{\raggedright\arraybackslash}p{(\linewidth - 18\tabcolsep) * \real{0.1000}}@{}}
\toprule\noalign{}
\begin{minipage}[b]{\linewidth}\raggedright
Biofilm and MICP Model Summary Variables
\end{minipage} & \begin{minipage}[b]{\linewidth}\raggedright
\end{minipage} & \begin{minipage}[b]{\linewidth}\raggedright
\end{minipage} & \begin{minipage}[b]{\linewidth}\raggedright
\end{minipage} & \begin{minipage}[b]{\linewidth}\raggedright
\end{minipage} & \begin{minipage}[b]{\linewidth}\raggedright
\end{minipage} & \begin{minipage}[b]{\linewidth}\raggedright
\end{minipage} & \begin{minipage}[b]{\linewidth}\raggedright
\end{minipage} & \begin{minipage}[b]{\linewidth}\raggedright
\end{minipage} & \begin{minipage}[b]{\linewidth}\raggedright
\end{minipage} \\
\midrule\noalign{}
\endhead
\bottomrule\noalign{}
\endlastfoot
Type & Variable & Root & Field & Group & Well & Well Connection & Well
Completion & Region & Block \\
Flow & Microbial Mass Injection Rate & MMIR & FMMIR & GMMIR & WMMIR &
CMMIR & CMMIRL & & \\
Flow & Oxygen Mass Injection Rate & MOIR & FMOIR & GMOIR & WMOIR & CMOIR
& CMOIRL & & \\
Flow & Urea Mass Injection Rate & MUIR & FMUIR & GMUIR & WMUIR & CMUIR &
CMUIRL & & \\
Flow & Microbial Mass Injection Total & MMIT & FMMIT & GMMIT & WMMIT &
CMMIT & CMMITL & & \\
Flow & Oxygen Mass Injection Total & MOIT & FMOIT & GMOIT & WMOIT &
CMOIT & CMOITL & & \\
Flow & Urea Mass Injection Total & MUIT & FMUIT & GMUIT & WMUIT & CMUIT
& CMUITL & & \\
Flow & Microbial Mass Production Rate & MMPR & FMMPR & GMMPR & WMMPR &
CMMPR & CMMPRL & & \\
Flow & Oxygen Mass Production Rate & MOPR & FMOPR & GMOPR & WMOPR &
CMOPR & CMOPRL & & \\
Flow & Urea Mass Production Rate & MUPR & FMUPR & GMUPR & WMUPR & CMUPR
& CMUPRL & & \\
Flow & Microbial Mass Production Total & MMPT & FMMPT & GMMPT & WMMPT &
CMMPT & CMMPTL & & \\
Flow & Oxygen Mass Production Total & MOPT & FMOPT & GMOPT & WMOPT &
CMOPT & CMOPTL & & \\
Flow & Urea Mass Production Total & MUPT & FMUPT & GMUPT & WMUPT & CMUPT
& CMUPTL & & \\
Mass & Microbes In-place & MMIP & FMMIP & & & & & RMMIP & BMMIP \\
Mass & Oxygen In-place & MOIP & FMOIP & & & & & RMOIP & BMOIP \\
Mass & Urea In-place & MUIP & FMUIP & & & & & RMUIP & BMUIP \\
Mass & Biofilm In-place & MBIP & FMBIP & & & & & RMBIP & BMBIP \\
Mass & Calcite In-place & MCIP & FMCIP & & & & & RMCIP & BMCIP \\
Notes: Cells colored in gray with no mnemonic indicate that the summary
variable is not available for this combination. & & & & & & & & & \\
\end{longtable}

\emph{Table 11.2.21: Biofilm and MICP Model Summary Variables}

\subsection{Option Specific Variables -- Multi-Segment
Wells}\label{option-specific-variables-multi-segment-wells}

The summary variables in this section are for use for when multi-segment
wells are active in the model, and the variables only apply to the well
segments requested for a multi-segment well, not to the well, well
connection, or well completion.

Note that not all these variables are available in OPM Flow; however,
the simulator will issue warning messages if this is indeed the case. It
is anticipated that the number of recognized summary variables will
increase in future releases of OPM Flow.

\begin{longtable}[]{@{}
  >{\raggedright\arraybackslash}p{(\linewidth - 8\tabcolsep) * \real{0.2000}}
  >{\raggedright\arraybackslash}p{(\linewidth - 8\tabcolsep) * \real{0.2000}}
  >{\raggedright\arraybackslash}p{(\linewidth - 8\tabcolsep) * \real{0.2000}}
  >{\raggedright\arraybackslash}p{(\linewidth - 8\tabcolsep) * \real{0.2000}}
  >{\raggedright\arraybackslash}p{(\linewidth - 8\tabcolsep) * \real{0.2000}}@{}}
\toprule\noalign{}
\begin{minipage}[b]{\linewidth}\raggedright
Multi-Segment Wells Summary Variables
\end{minipage} & \begin{minipage}[b]{\linewidth}\raggedright
\end{minipage} & \begin{minipage}[b]{\linewidth}\raggedright
\end{minipage} & \begin{minipage}[b]{\linewidth}\raggedright
\end{minipage} & \begin{minipage}[b]{\linewidth}\raggedright
\end{minipage} \\
\midrule\noalign{}
\endhead
\bottomrule\noalign{}
\endlastfoot
Type & Variable & Root & Well Segment & Comment \\
Flow & Artificial Lift Quantity & ALQ & SALQ & See keyword WSEGTABL. \\
Flow & Brine Flow Rate & FR & SFR & Surfactant and Polymer model. \\
Flow & Gas Flow Rate & GFR & SGFR & \\
Flow & Gas Flow Rate (Free) & GFRF & SGFRF & \\
Flow & Gas Flow Rate (Solution) & GFRS & SGFRS & \\
Flow & Gas Flow Total (Absolute) & GFTA & SGFTA & \\
Flow & Gas Flow Total & GFT & SGFT & \\
Flow & Gas Flow Velocity & GFV & SGFV & \\
Flow & Gas Holdup Fraction & GHF & SGHF & \\
Flow & Gas Import Rate & GIMR & SGIMR & See keyword WSEGEXSS. \\
Flow & Gas Import Total & GIMT & SGIMT & See keyword WSEGEXSS. \\
Flow & Oil Absolute Flow Total & OFTA & SOFTA & \\
Flow & Oil Flow Rate & OFR & SOFR & \\
Flow & Oil Flow Rate (Free) & OFRF & SOFRF & \\
Flow & Oil Flow Rate (Solution) & OFRS & SOFRS & \\
Flow & Oil Flow Total & OFT & SOFT & \\
Flow & Oil Flow Velocity & OFV & SOFV & \\
Flow & Oil Holdup Fraction & OHF & SOHF & \\
Flow & Polymer Flow Rate & CFR & SCFR & Polymer model. \\
Flow & Tracer Flow Rate & TFR & STFR & Tracer model. \\
Flow & Water Flow Rate & WFR & SWFR & \\
Flow & Water Flow Total & WFT & SWFT & \\
Flow & Water Flow Total (Absolute) & WFTA & SWFTA & \\
Flow & Water Flow Velocity & WFV & SWFV & \\
Flow & Water Holdup Fraction & WHF & SWHF & \\
Flow & Water Import Rate & WIMR & SWIMR & See keyword WSEGEXSS. \\
Flow & Water Import Total & WIMT & SWIMT & See keyword WSEGEXSS. \\
Pressure & Pressure & PR & SPR & \\
Pressure & Pressure Drop & PRD & SPRD & \\
Pressure & Pressure Drop Component due to Friction & PRDF & SPRDF & \\
Pressure & Pressure Drop Component due to Hydrostatic head & PRDH &
SPRDH & \\
Pressure & Pressure Drop due to Acceleration head & PRDA & SPRDA & \\
Pressure & Pressure Drop Frictional Multiplier & PRDM & SPRDM & See
keyword WSEGMULT. \\
Property & API & API & SAPI & API model. \\
Property & Fluid Mixture Density (Excludes Exponents) & DENM & SDENM &
Excludes exponents of flowing fractions. \\
Property & Fluid Mixture Density (Includes Exponents) & MDEN & SMDEN &
Includes exponents of flowing fractions. \\
Property & Fluid Mixture Viscosity (Includes Exponents) & MVIS & SMVIS &
Includes exponents of flowing fractions. \\
Property & Fluid Viscosity (Effective Mixture) & EMVIS & SEMVIS &
Water/polymer fluid. \\
Property & Gas Density & GDEN & SGDEN & \\
Property & Gas Viscosity & GVIS & SGVIS & \\
Property & Gas-Liquid Drift Velocity (Vd) & GLVD & SGLVD & Drift flux
slip model. \\
Property & Gas-Liquid Profile Parameter (C0) & GLPP & SGLPP & Drift flux
slip model. \\
Property & Oil Density & ODEN & SODEN & \\
Property & Oil Viscosity & OVIS & SOVIS & \\
Property & Oil-Water Drift Velocity (Vd) & OWVD & SOWVD & Drift flux
slip model. \\
Property & Oil-Water Profile Parameter (C0) & OWPP & SOWPP & Drift flux
slip model. \\
Property & Pump Working Power & PPOW & SPPOW & See keyword WSEGPULL. \\
Property & Scale Deposition Diameter (Karst Conduit Calcite Dissolution)
& FD & SFD & Scale Deposition model. \\
Property & Setting of Segment & FOPN & SFOPN & See keywords WSEGVALV,
WSEGAICD, WSEGSICD, WSEGTABL, WSEGLABY, and WSEGFLIM. \\
Property & Strength of ICD on Segment & STR & SSTR & See keywords
WSEGVALV, WSEGAICD, and WSEGSICD. \\
Property & Water Density (Pure Water) & WDEN & SWDEN & \\
Property & Water Viscosity (Pure Water) & WVIS & SWVIS & \\
Ratio & Gas-Oil Ratio & GOR & SGOR & \\
Ratio & Oil-Gas Ratio & OGR & SOGR & \\
Ratio & Water Cut & WCT & SWCT & \\
Ratio & Water Gas Ratio & WGR & SWGR & \\
Volume & Brine Concentration & SCN & SSCN & Surfactant and Polymer
model \\
Volume & Polymer Concentration & CCN & SCCN & Polymer model \\
Volume & Tracer Concentration & TFC & STFC & Tracer model \\
Notes: & & & & \\
\end{longtable}

\emph{Table 11.2.17: Multi-Segment Wells Summary Variables}

\subsubsection{Examples}\label{examples-62}

For example, to get the segment oil flow rates for well OP01 segments
two, three, and four and for segment two for OP02 one would use:

\begin{verbatim}
-- ==============================================================================
--
-- SUMMARY SECTION
--
-- ==============================================================================
SUMMARY

SOFR
‘OP01’ 2 /
‘OP01’ 3 /
‘OP01’ 4 /
‘OP02’ 2 /
/
\end{verbatim}

or to get all the oil flow rates for all the segments for well OP01:

\begin{verbatim}
-- ==============================================================================
--
-- SUMMARY SECTION
--
-- ==============================================================================
SUMMARY

SOFR
‘OP01’ /
/
\end{verbatim}

Care should be exercised when defaulting the list of wells and segments
as there is the potential to generate large volumes of data.

\subsection{Option Specific Variables -- Network
Model}\label{option-specific-variables-network-model}

For the Network models only a few SUMMARY vectors for groups and nodes
are available, both in OPM Flow and the commercial simulator, as
depicted in Table 11.2.18. There are two types of network option in the
simulator: a Standard Network and an Extended Network option, the latter
is activated by the NETWORK keyword in the RUNSPEC section. The summary
vectors apply to both the Standard and Extended Network options.

\begin{longtable}[]{@{}
  >{\raggedright\arraybackslash}p{(\linewidth - 16\tabcolsep) * \real{0.1111}}
  >{\raggedright\arraybackslash}p{(\linewidth - 16\tabcolsep) * \real{0.1111}}
  >{\raggedright\arraybackslash}p{(\linewidth - 16\tabcolsep) * \real{0.1111}}
  >{\raggedright\arraybackslash}p{(\linewidth - 16\tabcolsep) * \real{0.1111}}
  >{\raggedright\arraybackslash}p{(\linewidth - 16\tabcolsep) * \real{0.1111}}
  >{\raggedright\arraybackslash}p{(\linewidth - 16\tabcolsep) * \real{0.1111}}
  >{\raggedright\arraybackslash}p{(\linewidth - 16\tabcolsep) * \real{0.1111}}
  >{\raggedright\arraybackslash}p{(\linewidth - 16\tabcolsep) * \real{0.1111}}
  >{\raggedright\arraybackslash}p{(\linewidth - 16\tabcolsep) * \real{0.1111}}@{}}
\toprule\noalign{}
\begin{minipage}[b]{\linewidth}\raggedright
Network Model Summary Variables
\end{minipage} & \begin{minipage}[b]{\linewidth}\raggedright
\end{minipage} & \begin{minipage}[b]{\linewidth}\raggedright
\end{minipage} & \begin{minipage}[b]{\linewidth}\raggedright
\end{minipage} & \begin{minipage}[b]{\linewidth}\raggedright
\end{minipage} & \begin{minipage}[b]{\linewidth}\raggedright
\end{minipage} & \begin{minipage}[b]{\linewidth}\raggedright
\end{minipage} & \begin{minipage}[b]{\linewidth}\raggedright
\end{minipage} & \begin{minipage}[b]{\linewidth}\raggedright
\end{minipage} \\
\midrule\noalign{}
\endhead
\bottomrule\noalign{}
\endlastfoot
Type & Variable & Root & Field & Group & Well & Well Connection & Region
& Block \\
Pressure & Group/Node pressure in a production network based on rates at
end of timestep (OPM Flow specific). An alias for NPR. & & & GNETPR & &
& & \\
Pressure & Group/Node pressure in a production network. & & & GPR & & &
& \\
Pressure & Group /Node pressure in a gas injection network. & & & GPRG &
& & & \\
Pressure & Group/Node pressure in a water injection network. & & & GPRW
& & & & \\
Pressure & Pressure drop along the group's or node's outlet branch in a
production network. & & & GPRB & & & & \\
Pressure & Pressure drop along the group's or node's inlet branch in a
gas injection network. & & & GPRBG & & & & \\
Pressure & Pressure drop along the group's or node's inlet branch in a
water injection network. & & & GPRBW & & & & \\
Pressure & Group/Node pressure in a production network based on rates at
end of timestep (OPM Flow specific). & & & NPR & & & & \\
Flow & Artificial Lift Quantity (``ALQ'') in the group's or node's
outlet branch in a production network. Note no units are stated as ALQ
units are variable, for example Hz for ESP or rate for gas lift gas. & &
& GALQ & & & & \\
Flow & Oil flow rate along the group's or node's outlet branch in a
production network. & & & GOPRNB & & & & \\
Flow & Water flow rate along the group's or node's outlet branch in a
production network. & & & GWPRNB & & & & \\
Flow & Gas flow rate along the group's or node's outlet branch in a
production network. & & & GGPRNB & & & & \\
Flow & Liquid flow rate along the group's or node's outlet branch in a
production network. & & & GLPRNB & & & & \\
Flow & Water flow rate along the group's or node's inlet branch in water
injection network. & & & GWIRNB & & & & \\
Flow & Gas flow rate along the group's or node's inlet branch in gas
injection network. & & & GGIRNB & & & & \\
Notes: & & & & & & & & \\
\end{longtable}

\emph{Table 11.2.18: Network Model Summary Variables}

Note that not all these variables are available in OPM Flow; however,
the simulator will issue warning messages if this is indeed the case. It
is anticipated that the number of recognized summary variables will
increase in future releases of OPM Flow.

\subsection{Option Specific Variables -- OPM Flow Simulation
Performance}\label{option-specific-variables-opm-flow-simulation-performance}

The following table (Table 11.2.19) lists the OPM Flow simulation
performance summary variables that can be written to the SUMMARY file.
Note that not all these variables are available in OPM Flow; however,
the simulator will issue warning messages if this is indeed the case. It
is anticipated that the number of recognized summary variables will
increase in future releases of OPM Flow.

\begin{longtable}[]{@{}
  >{\raggedright\arraybackslash}p{(\linewidth - 4\tabcolsep) * \real{0.3333}}
  >{\raggedright\arraybackslash}p{(\linewidth - 4\tabcolsep) * \real{0.3333}}
  >{\raggedright\arraybackslash}p{(\linewidth - 4\tabcolsep) * \real{0.3333}}@{}}
\toprule\noalign{}
\begin{minipage}[b]{\linewidth}\raggedright
OPM Flow Simulation Performance
\end{minipage} & \begin{minipage}[b]{\linewidth}\raggedright
\end{minipage} & \begin{minipage}[b]{\linewidth}\raggedright
\end{minipage} \\
\midrule\noalign{}
\endhead
\bottomrule\noalign{}
\endlastfoot
Variable Description & Variable & Comment \\
Elapsed - Cumulative waiting time per process in seconds. & DTMWAIT & \\
Elapsed - Elapsed time in seconds. & ELAPSED & No data written to
file. \\
Iterations - Number of linear iterations for each gradient calculation
(Gradient Option). & HLINEARS & \\
Iterations - Total number of linear iterations for each gradient
pressure (Gradient Option). & HSUMLINS & \\
Maximum Courant-Friedrichs-Lewy value for each time step. & MAXCFL &
Compositional keyword. \\
Maximum variation in energy for each time step. & MAXDE & Compositional
keyword. \\
Maximum variation in pressure for each time step. & MAXDPR & \\
Maximum variation in gas saturation for each time step. & MAXDSG & \\
Maximum variation in oil saturation for each time step. & MAXDSO & \\
Maximum variation in water saturation for each time step. & MAXDSW & \\
Maximum variation in temperature for each time step. & MAXDT &
Compositional keyword. \\
Memory - Operating system reported maximum current memory usage across
all processors. & MEMORYTS & \\
CPR Solver - Number of pressure iterations for each time step. &
MLINEARP & \\
Iterations - Number of linear iterations for each time step. & MLINEARS
& \\
Messages - Cumulative number of BUG messages. & MSUMBUG & \\
Messages - Cumulative number of COMMENT messages. & MSUMCOMM & \\
Messages - Cumulative number of ERROR messages. & MSUMERR & \\
CPR Solver - Cumulative number of pressure iterations. & MSUMLINP & \\
Iterations - Cumulative number of linear iterations. & MSUMLINS & \\
Messages - Cumulative number of MESSAGES messages. & MSUMMESS & \\
Iterations - Cumulative number of Newton iterations. & MSUMNEWT & \\
Messages - Cumulative number of PROBLEM messages. & MSUMPROB & \\
Messages - Cumulative number of WARNING messages. & MSUMWARN & \\
Number of flash calculations for each time step & NBAKFL & Compositional
keyword. \\
Memory - Peak usage of dynamically allocated memory reported by OPM
Flow. For parallel cases this is the maximum across all processors. &
NBYTOT & \\
CPR Solver - Average number of pressure iterations per linear iteration
for each time step. & NCPRLINS & \\
Iterations -- Cumulative average number of Newton iterations per cell in
flash calculations.. & NEWTFL & Compositional keyword. \\
Iterations - Number of Newton iterations used for each time step. &
NEWTON & \\
CPR Solver - Average number of pressure iterations per Newton iteration
per time step. & NLINEARP & \\
Iterations - Average number of linear iterations per Newton iteration
for each time step. & NLINEARS & For runs with LGRs, LLINEARS will
automatically be exported for each LGR. \\
Iterations - Maximum number of linear iterations in the Newton
iterations per time step. & NLINSMAX & \\
Iterations - Minimum number of linear iterations in the Newton
iterations per time step. & NLINSMIN & \\
Cumulative number of flash calculations. & NNUMFL & Compositional
keyword. \\
Cumulative number of stability tests. & NNUMST & Compositional
keyword. \\
Non-linear residual total. & NLRESSUM & Compositional keyword. \\
Non-linear residual maximum value. & NLRESMAX & Compositional
keyword. \\
Number of time steps -- taken. & NTS & Compositional keyword. \\
Number of time steps -- pressure converged. & NTSPCL & Compositional
keyword. \\
Number of time steps -- molar density converged. & NTSMCL &
Compositional keyword. \\
Number of time steps -- energy density converged. & NTSECL &
Compositional keyword. \\
Time Step -- Criteria used to select the length of the time step, set to
one of the following: & STEPTYPE & No data written to file. Note if the
RUNSUM keyword has been activated in the SUMMARY section, then the
mnemonics are written to the RSM file instead of integer values. \\
Value & Criteria & Definition \\
1 & INIT & Initial time step. \\
2 & TRNC & Controlled via time truncation error. \\
3 & MINF & Minimum ratio between one time step and the next. \\
4 & MAXF & Maximum ratio between one time step and the next. \\
5 & MINS & Minimum permitted time step. \\
6 & MAXS & Maximum permitted time step. \\
7 & REPT & Report time step. \\
8 & HALF & Half of the time set to reach the next report time step. \\
9 & CHOP & Time step chopped due non-convergence of non-linear
equations. \\
10 & SATM & Maximum saturation change. \\
11 & PCHP & Time step chopped due to maximum pressure change in IMPES
formulation. \\
12 & DIFF & Difficult time step after a time step CHOP. \\
13 & LGRC & Determined by LGR fluid in-place error targets. \\
14 & SURF & Set by maximum expected grid block surfactant concentration
change. \\
15 & NETW & Time step determined by network balancing controls. \\
16 & THRP & Time step controlled by maximum through put ratio. \\
17 & EMTH & Time step selected to match the end of a month. \\
18 & MAXP & Set by maximum expected grid block pressure change. \\
19 & WCYC & Determined by well cycling on/off time. \\
20 & MAST & Chosen by the master simulation in a reservoir coupling
run. \\
21 & SLVR & Selected by a slave simulation to match a slave reporting
time step, in a reservoir coupling run. \\
22 & SLVC & Time step selected by a slave simulation due to a slave's
expected flow rate change in a reservoir coupling run. \\
23 & MAXW & Maximum time step size after a well control event. \\
24 & EFF+ & Selected by ZIPPY optimum time selection algorithm. \\
25 & EFF- & \\
26 & NLTR & \\
27 & EFFT & \\
28 & DLYA & \\
29 & ACTN & \\
30 & RAIN & \\
CPU - Current CPU usage in seconds. & TCPU & It does not consider the
time taken by inter-process communications in parallel runs, whereas
ELAPSED does. Thus, for parallel jobs, ELAPSED is the most relevant time
measurement. \\
CPU - CPU time per day (or hour in lab units depending on run units
system). & TCPUDAY & \\
CPU -- Cumulative CPU time for each gradient calculation (Gradient
Option). & TCPUH & \\
CPU - Cumulative CPU time for all gradient calculations (Gradient
Option). & TCPUHT & \\
CPU - Cumulative CPU time used in SCHEDULE section. & TCPUSCH & \\
CPU - CPU time per time step in seconds. & TCPUTS & \\
CPU - CPU time per time step for each gradient calculation (Gradient
Option). & TCPUTSH & \\
CPU - CPU time per time step for all gradient calculations (Gradient
Option). & TCPUTSHT & \\
CPU -- Cumulative Tracer CPU time. & TCPUTR & Compositional keyword. \\
CPU -- Cumulative Tracer implicit CPU time. & TCPUTRSV & Compositional
keyword. \\
CPU -- Cumulative Tracer Lagrangian tracer stramline CPU time. &
TCPULGTR & Compositional keyword. \\
CPU -- Cumulative Tracer Lagrangian tracer solver CPU time. & TCPULGSV &
Compositional keyword. \\
Elapsed - Elapsed time per day (or hour in lab units). & TELAPDAY & \\
Elapsed - Elapsed time per linear iteration in seconds. & TELAPLIN & No
data written to file. \\
Elapsed - Elapsed time per time step in seconds. & TELAPTS & \\
Time Step - Length of time step. & TIMESTEP & \\
Iterations - Number of well Newton iterations taken in the last global
Newton iteration. A negative value indicates that the well failed to
converge. & WNEWTON & \\
Time Step - Predicted efficiency of the time step (developer use). &
ZIPEFF & \\
Time Step - Predicted efficiency of the time step divided by the actual
efficiency (developer use). & ZIPEFFC & \\
Notes: & & \\
\end{longtable}

\emph{Table 11.2.19: OPM Flow Simulator Performance Summary Variables}

See also the PERFORMA and NMESSAGE keywords in section Error: Reference
source not found Error: Reference source not found that write out a
selection of the variables in Table 11.2.19.

\subsection{Option Specific Variables -- Polymer
Model}\label{option-specific-variables-polymer-model}

The summary variables in this section are associated with Polymer phase
and the polymer flood model. The feature is activated by the POLYMER
keyword in the RUNSPEC section.

Note that not all these variables are available in OPM Flow; however,
the simulator will issue warning messages if this is indeed the case. It
is anticipated that the number of recognized summary variables will
increase in future releases of OPM Flow.

\begin{longtable}[]{@{}
  >{\raggedright\arraybackslash}p{(\linewidth - 18\tabcolsep) * \real{0.1000}}
  >{\raggedright\arraybackslash}p{(\linewidth - 18\tabcolsep) * \real{0.1000}}
  >{\raggedright\arraybackslash}p{(\linewidth - 18\tabcolsep) * \real{0.1000}}
  >{\raggedright\arraybackslash}p{(\linewidth - 18\tabcolsep) * \real{0.1000}}
  >{\raggedright\arraybackslash}p{(\linewidth - 18\tabcolsep) * \real{0.1000}}
  >{\raggedright\arraybackslash}p{(\linewidth - 18\tabcolsep) * \real{0.1000}}
  >{\raggedright\arraybackslash}p{(\linewidth - 18\tabcolsep) * \real{0.1000}}
  >{\raggedright\arraybackslash}p{(\linewidth - 18\tabcolsep) * \real{0.1000}}
  >{\raggedright\arraybackslash}p{(\linewidth - 18\tabcolsep) * \real{0.1000}}
  >{\raggedright\arraybackslash}p{(\linewidth - 18\tabcolsep) * \real{0.1000}}@{}}
\toprule\noalign{}
\begin{minipage}[b]{\linewidth}\raggedright
Polymer Model Summary Variables
\end{minipage} & \begin{minipage}[b]{\linewidth}\raggedright
\end{minipage} & \begin{minipage}[b]{\linewidth}\raggedright
\end{minipage} & \begin{minipage}[b]{\linewidth}\raggedright
\end{minipage} & \begin{minipage}[b]{\linewidth}\raggedright
\end{minipage} & \begin{minipage}[b]{\linewidth}\raggedright
\end{minipage} & \begin{minipage}[b]{\linewidth}\raggedright
\end{minipage} & \begin{minipage}[b]{\linewidth}\raggedright
\end{minipage} & \begin{minipage}[b]{\linewidth}\raggedright
\end{minipage} & \begin{minipage}[b]{\linewidth}\raggedright
\end{minipage} \\
\midrule\noalign{}
\endhead
\bottomrule\noalign{}
\endlastfoot
Type & Variable & Root & Field & Group & Well & Well Connection & Region
& Block & Comment \\
Flow & Polymer Flow Rate & CFR & & & & CCFR & & & \\
Flow & Polymer Flow Total (Inter-Region) & CFT & & & & & RCFT & & \\
Flow & Polymer Injection Rate & CIR & FCIR & GCIR & WCIR & CCIR & & & \\
Flow & Polymer Injection Total & CIT & FCIT & GCIT & WCIT & CCIT & &
& \\
Flow & Polymer Production Rate & CPR & FCPR & GCPR & WCPR & CCPR & &
& \\
Flow & Polymer Production Total & CPT & FCPT & GCPT & WCPT & CCPT & &
& \\
Flow & Salt Flow Rate & SFR & & & & CSPR & & & \\
Flow & Salt Injection Rate & SIR & FSIR & GSIR & WSIR & CSIR & & & \\
Flow & Salt Injection Total & SIT & FSIT & GSIT & WSIT & CSIT & & & \\
Flow & Salt Production Rate & SPR & FSPR & GSPR & WSPR & CSPR & & & \\
Flow & Salt Production Total & SPT & FSPT & GSPT & WSPT & CSPT & & & \\
Property & Polymer Adsorbed (PLYTRRFA) & CABnnn & & & & & & BCABnnn &
The highest temperature band at which residual resistance factor was
calculated, see keyword PLYTRRFA. The band number nnn can range from 001
to 999, but must be less than or equal to the argument on the PLYTRRFA
keyword. \\
Property & Polymer Adsorption Total & CAD & FCAD & & & & RCAD & BCAD &
Block is concentration. \\
Property & Polymer Solution (Effective Viscosity) & EPVIS & & & & & &
BEPVIS & Also BVPOLY. \\
Property & Polymer Solution Shear Viscosity (Inter-Block I+ Direction) &
SHWVISI & & & & & & BSHWVISI & \\
Property & Polymer Solution Shear Viscosity (Inter-Block J+ Direction) &
SHWVISJ & & & & & & BSHWVISJ & \\
Property & Polymer Solution Shear Viscosity (Inter-Block K+ Direction) &
SHWVISK & & & & & & BSHWVISK & \\
Property & Polymer Thermal Degradation & CDCS & & & & & & BCDCS & Total
mass degraded in previous time step. \\
Property & Polymer Thermal Degradation Rate & CDCR & & & & & & BCDCR
& \\
Property & Polymer Thermal Degradation Rate (Adsorbed) & CDCA & & & & &
& BCDCA & \\
Property & Polymer Thermal Degradation Rate (Solution) & CDCP & & & & &
& BCDCP & \\
Property & Polymer Water (Effective Viscosity) & EMVIS & & & & & &
BEMVIS & Based on block center properties. \\
Property & Water Flow Rate Times Shear Multiplier (Inter-Block I+
Direction) & FLOW0I & & & & & & BFLOW0I & Multiplied by the
corresponding shear multiplier (PLYSHLOG and PLYSHEAR options only). \\
Property & Water Flow Rate Times Shear Multiplier (Inter-Block J+
Direction) & FLOW0J & & & & & & BFLOW0J & \\
Property & Water Flow Rate Times Shear Multiplier (Inter-Block K+
Direction) & FLOW0K & & & & & & BFLOW0K & \\
Property & Water Shear Rate After Shear Effects (Inter-Block I+
Direction) & SRTWI & & & & & & BSRTWI & \\
Property & Water Shear Rate After Shear Effects (Inter-Block J+
Direction) & SRTWJ & & & & & & BSRTWJ & \\
Property & Water Shear Rate After Shear Effects (Inter-Block K+
Direction) & SRTWK & & & & & & BSRTWK & \\
Property & Water Shear Rate Prior to Shear Effects (Inter-Block I+
Direction) & SRTW0I & & & & & & BSRTW0I & \\
Property & Water Shear Rate Prior to Shear Effects (Inter-Block J+
Direction) & SRTW0J & & & & & & BSRTW0J & \\
Property & Water Shear Rate Prior to Shear Effects (Inter-Block K+
Direction) & SRTW0K & & & & & & BSRTW0K & \\
Property & Water Velocity (Inter-Block I+ Direction) & VELW0I & & & & &
& BVELW0I & \\
Property & Water Velocity (Inter-Block J+ Direction) & VELW0J & & & & &
& BVELW0J & \\
Property & Water Velocity (Inter-Block K+ Direction) & VELW0K & & & & &
& BVELW0K & \\
Property & Water Viscosity (Effective) & EWV\_POL & & & & & & BEWV\_POL
& \\
Property & Water Viscosity Sheared Factor (Inter-Block I+ Direction) &
PSHLZI & & & & & & BPSHLZI & \\
Property & Water Viscosity Sheared Factor (Inter-Block J+ Direction) &
PSHLZJ & & & & & & BPSHLZJ & \\
Property & Water Viscosity Sheared Factor (Inter-Block K+ Direction) &
PSHLZK & & & & & & BPSHLZK & \\
Volume & Polymer Concentration (Block) & CCN & & & & & & BCCN & \\
Volume & Polymer Concentration (Injection) & CIC & FCIC & GCIC & WCIC &
CCIC & & & \\
Volume & Polymer Concentration (Production) & CPC & FCPC & GCPC & WCPC &
CCPC & & & \\
Volume & Polymer Solution (In Solution) & CIP & FCIP & & & & RCIP & BCIP
& \\
Volume & Salt Cell Concentration & SCN & & & & & & BSCN & \\
Volume & Salt In Place & SIP & FSIP & & & & RSIP & BSIP & \\
Volume & Salt inter-region Flow Total & SFT & & & & & RSFT & & \\
Notes: & & & & & & & & & \\
\end{longtable}

\emph{Table 11.2.20: Polymer Model Summary Variables}

\subsection{Option Specific Variables -- Pseudo Steady
State}\label{option-specific-variables-pseudo-steady-state}

This feature has not been implemented in OPM Flow.

\subsection{Option Specific Variables -- Reservoir
Coupling}\label{option-specific-variables-reservoir-coupling}

This feature has not been implemented in OPM Flow.

\subsection{Option Specific Variables -- Solvent
Model}\label{option-specific-variables-solvent-model}

The SOLVENT keyword in the RUNSPEC section activates the solvent phase
in the model and activates the four component solvent model for the run.
In addition to this keyword, the oil, water, and gas phases should also
be declared for the run using the OIL, WATER, and GAS keywords. The
summary variables for this option are outlined in Table 11.2.21.

Note that not all these variables are available in OPM Flow; however,
the simulator will issue warning messages if this is indeed the case. It
is anticipated that the number of recognized summary variables will
increase in future releases of OPM Flow.

\begin{longtable}[]{@{}
  >{\raggedright\arraybackslash}p{(\linewidth - 18\tabcolsep) * \real{0.1000}}
  >{\raggedright\arraybackslash}p{(\linewidth - 18\tabcolsep) * \real{0.1000}}
  >{\raggedright\arraybackslash}p{(\linewidth - 18\tabcolsep) * \real{0.1000}}
  >{\raggedright\arraybackslash}p{(\linewidth - 18\tabcolsep) * \real{0.1000}}
  >{\raggedright\arraybackslash}p{(\linewidth - 18\tabcolsep) * \real{0.1000}}
  >{\raggedright\arraybackslash}p{(\linewidth - 18\tabcolsep) * \real{0.1000}}
  >{\raggedright\arraybackslash}p{(\linewidth - 18\tabcolsep) * \real{0.1000}}
  >{\raggedright\arraybackslash}p{(\linewidth - 18\tabcolsep) * \real{0.1000}}
  >{\raggedright\arraybackslash}p{(\linewidth - 18\tabcolsep) * \real{0.1000}}
  >{\raggedright\arraybackslash}p{(\linewidth - 18\tabcolsep) * \real{0.1000}}@{}}
\toprule\noalign{}
\begin{minipage}[b]{\linewidth}\raggedright
Solvent Model Summary Variables
\end{minipage} & \begin{minipage}[b]{\linewidth}\raggedright
\end{minipage} & \begin{minipage}[b]{\linewidth}\raggedright
\end{minipage} & \begin{minipage}[b]{\linewidth}\raggedright
\end{minipage} & \begin{minipage}[b]{\linewidth}\raggedright
\end{minipage} & \begin{minipage}[b]{\linewidth}\raggedright
\end{minipage} & \begin{minipage}[b]{\linewidth}\raggedright
\end{minipage} & \begin{minipage}[b]{\linewidth}\raggedright
\end{minipage} & \begin{minipage}[b]{\linewidth}\raggedright
\end{minipage} & \begin{minipage}[b]{\linewidth}\raggedright
\end{minipage} \\
\midrule\noalign{}
\endhead
\bottomrule\noalign{}
\endlastfoot
Type & Variable & Root & Field & Group & Well & Well Connection & Region
& Block & Comment \\
Flow & Solvent Flow Rate & NFR & & & & CNFR & & & Field, group, and well
gas rates and cumulative volumes include the solvent gas. \\
Flow & Solvent Flow Total (Inter-Region) & NFT & & & & & RNFT & & \\
Flow & Solvent Injection Rate & NIR & FNIR & GNIR & WNIR & & & & \\
Flow & Solvent Injection Total & NIT & FNIT & GNIT & WNIT & CNIT & &
& \\
Flow & Solvent Production Rate & NPR & FNPR & GNPR & WNPR & & & & \\
Flow & Solvent Production Total & NPT & FNPT & GNPT & WNPT & CNPT & &
& \\
Property & Solvent Relative Permeability & NKR & & & & & & BNKR & \\
Volume & Solvent In-Place & NIP & FNIP & & & & RNIP & BNIP & \\
Volume & Solvent Saturation & NSAT & & & & & & BNSAT & \\
Notes: & & & & & & & & & \\
\end{longtable}

\emph{Table 11.2.21: Solvent Model Summary Variables}

\subsection{Option Specific Variables -- Surfactant
Model}\label{option-specific-variables-surfactant-model}

This feature has not been implemented in OPM Flow.

\subsection{Option Specific Variables -- Thermal
Model}\label{option-specific-variables-thermal-model}

The temperature option (TEMP keyword in the RUNSPEC section) and the
thermal option (THERMAL keyword in the RUNSPEC section) are two separate
modeling facilities in the commercial simulator, although some keywords
can be used by both options. OPM Flow's thermal implementation is based
on solving the energy equation fully coupled with the black-oil
equations so the results are not directly equivalent to the commercial
simulator's black-oil TEMP or compositional THERMAL formulations. The
THERMAL keyword is used to invoke OPM Flow's thermal option. The summary
variables for this option are listed in Table 11.2.22.

Note that not all these variables are available in OPM Flow; however,
the simulator will issue warning messages if this is indeed the case. It
is anticipated that the number of recognized summary variables will
increase in future releases of OPM Flow.

\begin{longtable}[]{@{}
  >{\raggedright\arraybackslash}p{(\linewidth - 18\tabcolsep) * \real{0.1000}}
  >{\raggedright\arraybackslash}p{(\linewidth - 18\tabcolsep) * \real{0.1000}}
  >{\raggedright\arraybackslash}p{(\linewidth - 18\tabcolsep) * \real{0.1000}}
  >{\raggedright\arraybackslash}p{(\linewidth - 18\tabcolsep) * \real{0.1000}}
  >{\raggedright\arraybackslash}p{(\linewidth - 18\tabcolsep) * \real{0.1000}}
  >{\raggedright\arraybackslash}p{(\linewidth - 18\tabcolsep) * \real{0.1000}}
  >{\raggedright\arraybackslash}p{(\linewidth - 18\tabcolsep) * \real{0.1000}}
  >{\raggedright\arraybackslash}p{(\linewidth - 18\tabcolsep) * \real{0.1000}}
  >{\raggedright\arraybackslash}p{(\linewidth - 18\tabcolsep) * \real{0.1000}}
  >{\raggedright\arraybackslash}p{(\linewidth - 18\tabcolsep) * \real{0.1000}}@{}}
\toprule\noalign{}
\begin{minipage}[b]{\linewidth}\raggedright
Thermal Model Summary Variables
\end{minipage} & \begin{minipage}[b]{\linewidth}\raggedright
\end{minipage} & \begin{minipage}[b]{\linewidth}\raggedright
\end{minipage} & \begin{minipage}[b]{\linewidth}\raggedright
\end{minipage} & \begin{minipage}[b]{\linewidth}\raggedright
\end{minipage} & \begin{minipage}[b]{\linewidth}\raggedright
\end{minipage} & \begin{minipage}[b]{\linewidth}\raggedright
\end{minipage} & \begin{minipage}[b]{\linewidth}\raggedright
\end{minipage} & \begin{minipage}[b]{\linewidth}\raggedright
\end{minipage} & \begin{minipage}[b]{\linewidth}\raggedright
\end{minipage} \\
\midrule\noalign{}
\endhead
\bottomrule\noalign{}
\endlastfoot
Type & Variable & Root & Field & Group & Well & Well Connection & Region
& Block & Comment \\
Flow & Energy Injection Rate & TIRHEA & FTIRHEA & GTIRHEA & WTIRHEA & &
& & \\
Flow & Energy Injection Total & TITHEA & FTITHEA & GTITHEA & WTITHEA & &
& & \\
Flow & Energy Production Rate & TPRHEA & FTPRHEA & GTPRHEA & WTPRHEA & &
& & \\
Flow & Energy Production Total & TPTHEA & FTPTHEA & GTPTHEA & WTPTHEA &
& & & \\
Property & Temperature (Block) & TCNFHEA & & & & & & BTCNFHEA & \\
Property & Temperature (Injection) & TICHEA & FTICHEA & GTICHEA &
WTICHEA & & & & \\
Property & Temperature (Production) & TPCHEA & FTPCHEA & GTPCHEA &
WTPCHEA & & & BTPCHEA & \\
Volume & Energy In-Place Difference & TIPTHEA & FTIPTHEA & & & &
RTIPTHEA4 & BTIPTHEA & Difference in energy in- place from start of run
to current time. \\
Notes: & & & & & & & & & \\
\end{longtable}

\emph{Table 11.2.22: Thermal Model Summary Variables}

See the THERMAL keyword in the RUNSPEC section for additional
information on this feature.

\subsection{Option Specific Variables -- User Define
Quantities}\label{option-specific-variables-user-define-quantities}

The UDQ keyword in the SCHEDULE section defines variables and operations
used to access the User Defined Quantities features in OPM Flow. UDQ
variables can be constants, summary variables, as defined in the SUMMARY
section, or a formula using various mathematical functions together with
constants and summary variables.

The variable names defined by the UDQ keyword in the SCHEDULE section
consist of a character string of up to eight characters in length, which
stipulates the name of the user defined variable based on the type of
variable. The variable type is defined by the first two characters of
the variable name and must be set to one of the following:

\begin{verbatim}
- CU: For variables that are associated with connections, for example summary variable COFR (Connection Oil Flow Rate).
- FU: For variables that are associated with field data, for example summary variable FOPR (Field Oil Production Rate).
- GU: For variables that are associated with groups, for example summary variable GLPR (Group Liquid Production Rate).
- RU: For variables that are associated with regions, for example summary variable RPR (Region Pressure).
- SU: For variables that are associated with multi-segment wells, for example summary variable SOFR (Segment Oil Flow Rate).
- WU: For variables that are associated with wells, for example summary variable WWCT (Well Water Cut).
- AU: For variables that are associated with aquifers, for example summary variable AAQP (Analytical Aquifer Pressure).
- BU: For variables that are associated with blocks, for example summary variable BPR (Block oil phase Pressure).
\end{verbatim}

The summary variable format for UDQ defined variables is presented in
Table 11.2.23.

\begin{longtable}[]{@{}
  >{\raggedright\arraybackslash}p{(\linewidth - 14\tabcolsep) * \real{0.1250}}
  >{\raggedright\arraybackslash}p{(\linewidth - 14\tabcolsep) * \real{0.1250}}
  >{\raggedright\arraybackslash}p{(\linewidth - 14\tabcolsep) * \real{0.1250}}
  >{\raggedright\arraybackslash}p{(\linewidth - 14\tabcolsep) * \real{0.1250}}
  >{\raggedright\arraybackslash}p{(\linewidth - 14\tabcolsep) * \real{0.1250}}
  >{\raggedright\arraybackslash}p{(\linewidth - 14\tabcolsep) * \real{0.1250}}
  >{\raggedright\arraybackslash}p{(\linewidth - 14\tabcolsep) * \real{0.1250}}
  >{\raggedright\arraybackslash}p{(\linewidth - 14\tabcolsep) * \real{0.1250}}@{}}
\toprule\noalign{}
\begin{minipage}[b]{\linewidth}\raggedright
User Define Quantities Summary Variables
\end{minipage} & \begin{minipage}[b]{\linewidth}\raggedright
\end{minipage} & \begin{minipage}[b]{\linewidth}\raggedright
\end{minipage} & \begin{minipage}[b]{\linewidth}\raggedright
\end{minipage} & \begin{minipage}[b]{\linewidth}\raggedright
\end{minipage} & \begin{minipage}[b]{\linewidth}\raggedright
\end{minipage} & \begin{minipage}[b]{\linewidth}\raggedright
\end{minipage} & \begin{minipage}[b]{\linewidth}\raggedright
\end{minipage} \\
\midrule\noalign{}
\endhead
\bottomrule\noalign{}
\endlastfoot
Variable & Root & Field & Group & Well & Well Connection & Region & Well
Segment \\
UDQ Variable & XXXXXX & FUXXXXXX & GUXXXXXX & WUXXXXXX & CUXXXXXX &
RUXXXXXX & SUXXXXXX \\
Notes: & & & & & & & \\
\end{longtable}

\emph{Table 11.2.23: User Define Quantities Summary Variables}

The following example illustrates the use of this type of summary
variable.

\subsubsection{Examples}\label{examples-63}

The example shows how to define some constant field variables used for
calculating facilities corrected condensate yields in a wet gas model in
the SCHEDULE section. The corrected condensate rate is stored as
variable named FU\_FNGLR.

\begin{verbatim}
-- ==============================================================================
--
-- SCHEDULE SECTION
--
-- ==============================================================================
SCHEDULE
--
-- DEFINE START OF USER DEFINED QUANTITY SECTION
--
UDQ
--
-- OPERATOR VARIABLE  EXPRESSION
--
ASSIGN      FUNGLYLD  1.100000                 / Condensate Yield Correction
--
-- OPERATOR VARIABLE  EXPRESSION
--
DEFINE      FU_FNGLR  FGPR *(FOGR * FUNGLYLD)  / Calculate Condensate Rate Field
UPDATE      FU_FNGLR  ON                       /
UNITS       FU_FNGLR  STBD                     /

/  DEFINE END OF USER DEFINED QUANTITY SECTION
\end{verbatim}

In order to report the field corrected condensate values, one would
declare FU\_FNGLR in the SUMMARY section, as shown below.

\begin{verbatim}
-- ==============================================================================
--
-- SUMMARY SECTION
--
-- ==============================================================================
SUMMARY

FU_FNGLR
\end{verbatim}

Note that the normal summary variables syntax rules still apply to
groups, regions, wells, etc. So for example, in addition, if the same
calculations were performed on two wet gas producers, GP01 and GP02,
then one would use:

\begin{verbatim}
-- ==============================================================================
--
-- SUMMARY SECTION
--
-- ==============================================================================
SUMMARY

FU_FNGLR

WU_FNGLR
GPO1
GPO2
/
\end{verbatim}

Assuming of course that WU\_FNGLR was defined in the SCHEDULE section.

\subsection{Option Specific Variables -- Wellbore Friction
model}\label{option-specific-variables-wellbore-friction-model}

This feature has not been implemented in OPM Flow.

\section{Keyword Definitions}\label{keyword-definitions-16}

\subsection{ALL -- Export Standard Summary Variable Vectors to
File}\label{all-export-standard-summary-variable-vectors-to-file}

\begin{longtable}[]{@{}
  >{\raggedright\arraybackslash}p{(\linewidth - 14\tabcolsep) * \real{0.1250}}
  >{\raggedright\arraybackslash}p{(\linewidth - 14\tabcolsep) * \real{0.1250}}
  >{\raggedright\arraybackslash}p{(\linewidth - 14\tabcolsep) * \real{0.1250}}
  >{\raggedright\arraybackslash}p{(\linewidth - 14\tabcolsep) * \real{0.1250}}
  >{\raggedright\arraybackslash}p{(\linewidth - 14\tabcolsep) * \real{0.1250}}
  >{\raggedright\arraybackslash}p{(\linewidth - 14\tabcolsep) * \real{0.1250}}
  >{\raggedright\arraybackslash}p{(\linewidth - 14\tabcolsep) * \real{0.1250}}
  >{\raggedright\arraybackslash}p{(\linewidth - 14\tabcolsep) * \real{0.1250}}@{}}
\toprule\noalign{}
\begin{minipage}[b]{\linewidth}\raggedright
RUNSPEC
\end{minipage} & \begin{minipage}[b]{\linewidth}\raggedright
GRID
\end{minipage} & \begin{minipage}[b]{\linewidth}\raggedright
EDIT
\end{minipage} & \begin{minipage}[b]{\linewidth}\raggedright
PROPS
\end{minipage} & \begin{minipage}[b]{\linewidth}\raggedright
REGIONS
\end{minipage} & \begin{minipage}[b]{\linewidth}\raggedright
SOLUTION
\end{minipage} & \begin{minipage}[b]{\linewidth}\raggedright
SUMMARY
\end{minipage} & \begin{minipage}[b]{\linewidth}\raggedright
SCHEDULE
\end{minipage} \\
\midrule\noalign{}
\endhead
\bottomrule\noalign{}
\endlastfoot
\end{longtable}

\subsubsection{Description}\label{description-847}

This keyword activates the writing out of a standard set of summary
production and injection data vectors for the field, group and well
objects to the SUMMARY (\emph{.SMSPEC and }.UNSMRY) and RSM (*.RSM)
files. Table 11.26 lists the production, injection, pressure and volume
summary variables written out by the ALL keyword, and Table 11.27 list
the aquifer variables.

\begin{longtable}[]{@{}
  >{\raggedright\arraybackslash}p{(\linewidth - 10\tabcolsep) * \real{0.1667}}
  >{\raggedright\arraybackslash}p{(\linewidth - 10\tabcolsep) * \real{0.1667}}
  >{\raggedright\arraybackslash}p{(\linewidth - 10\tabcolsep) * \real{0.1667}}
  >{\raggedright\arraybackslash}p{(\linewidth - 10\tabcolsep) * \real{0.1667}}
  >{\raggedright\arraybackslash}p{(\linewidth - 10\tabcolsep) * \real{0.1667}}
  >{\raggedright\arraybackslash}p{(\linewidth - 10\tabcolsep) * \real{0.1667}}@{}}
\toprule\noalign{}
\begin{minipage}[b]{\linewidth}\raggedright
Standard Production, Injection, and Pressures Summary Variables
\end{minipage} & \begin{minipage}[b]{\linewidth}\raggedright
\end{minipage} & \begin{minipage}[b]{\linewidth}\raggedright
\end{minipage} & \begin{minipage}[b]{\linewidth}\raggedright
\end{minipage} & \begin{minipage}[b]{\linewidth}\raggedright
\end{minipage} & \begin{minipage}[b]{\linewidth}\raggedright
\end{minipage} \\
\midrule\noalign{}
\endhead
\bottomrule\noalign{}
\endlastfoot
Variable & Root & Field & Group & Well & Comment \\
Gas Injection Rate & GIR & FGIR & GGIR & WGIR & \\
Gas Injection Total & GIT & FGIT & GGIT & WGIT & \\
Gas Production Rate & GPR & FGPR & GGPR & WGPR & Produced reservoir gas
only, gas lift gas is excluded. \\
Gas Production Total & GPT & FGPT & GGPT & WGPT & \\
Oil Injection Rate & OIR & FOIR & GOIR & WOIR & \\
Oil Injection Total & OIT & FOIT & GOIT & WOIT & \\
Oil Production Rate & OPR & FOPR & GOPR & WOPR & \\
Oil Production Total & OPT & FOPT & GOPT & WOPT & \\
Res. Vol. Injection Rate & VIR & FVIR & GVIR & WVIR & \\
Res. Vol. Injection Total & VIT & FVIT & GVIT & WVIT & \\
Res. Vol. Production Rate & VPR & FVPR & GVPR & WVPR & \\
Res. Vol. Production Total & VPT & FVPT & GVPT & WVPT & \\
Water Injection Rate & WIR & FWIR & GWIR & WWIR & \\
Water Injection Total & WIT & FWIT & GWIT & WWIT & \\
Water Production Rate & WPR & FWPR & GWPR & WWPR & \\
Water Production Total & WPT & FWPT & GWPT & WWPT & \\
Bottom-Hole Pressure & BHP & & & WBHP & \\
Productivity & PI & & & WPI & Preferred Phase \\
Tubing Head Pressure & THP & & & WTHP & \\
Gas-Oil Ratio & GOR & FGOR & GGOR & WGOR & \\
Water -Gas Ratio & WGR & FWGR & GWGR & WWGR & \\
Water Cut & WCT & FWCT & GWCT & WWCT & \\
Pressures and Volumes & & & & & \\
Variable & Root & Field & Group & Well & Comment \\
Oil Phase Pressure & PR & FPR & & & Field and Region HCPV weighted \\
Gas In-Place & GIP & FGIP & & & Liquid and gas phases \\
Gas In-Place (Gas Phase) & GIPG & FGIPG & & & \\
Gas In-Place (Liquid Phase) & GIPL & FGIPL & & & Only for cases that
have dissolved gas in the water phase \\
Oil In-Place & OIP & FOIP & & & Liquid and gas phases \\
Oil In-Place (Gas Phase) & OIPG & FOIPG & & & \\
Oil In-Place (Liquid Phase) & OIPL & FOIPL & & & \\
Water In-Place & WIP & FWIP & & & \\
Notes: & & & & & \\
\end{longtable}

\emph{Table 11.26: Standard Production, Injection, and Pressure Summary
Variables}

\begin{longtable}[]{@{}
  >{\raggedright\arraybackslash}p{(\linewidth - 12\tabcolsep) * \real{0.1429}}
  >{\raggedright\arraybackslash}p{(\linewidth - 12\tabcolsep) * \real{0.1429}}
  >{\raggedright\arraybackslash}p{(\linewidth - 12\tabcolsep) * \real{0.1429}}
  >{\raggedright\arraybackslash}p{(\linewidth - 12\tabcolsep) * \real{0.1429}}
  >{\raggedright\arraybackslash}p{(\linewidth - 12\tabcolsep) * \real{0.1429}}
  >{\raggedright\arraybackslash}p{(\linewidth - 12\tabcolsep) * \real{0.1429}}
  >{\raggedright\arraybackslash}p{(\linewidth - 12\tabcolsep) * \real{0.1429}}@{}}
\toprule\noalign{}
\begin{minipage}[b]{\linewidth}\raggedright
Standard Aquifer Summary Variables
\end{minipage} & \begin{minipage}[b]{\linewidth}\raggedright
\end{minipage} & \begin{minipage}[b]{\linewidth}\raggedright
\end{minipage} & \begin{minipage}[b]{\linewidth}\raggedright
\end{minipage} & \begin{minipage}[b]{\linewidth}\raggedright
\end{minipage} & \begin{minipage}[b]{\linewidth}\raggedright
\end{minipage} & \begin{minipage}[b]{\linewidth}\raggedright
\end{minipage} \\
\midrule\noalign{}
\endhead
\bottomrule\noalign{}
\endlastfoot
Variable & Root & Field & Analytical Aquifer & Analytical Aquifer List &
Numerical Aquifer & Comment \\
Aquifer Influx Rate (Water Aquifers) & QR & FAQR & AAQR & & & \\
Aquifer Influx Total (Water Aquifers) & QT & FAQT & AAQT & & & \\
Aquifer Influx Rate (Gas Aquifers) & QRG & FAQRG & AAQRG & & & \\
Aquifer Influx Total (Gas Aquifers) & QTG & FAQTG & AAQTG & & & \\
Notes: & & & & & & \\
\end{longtable}

\emph{Table 11.27: Standard Aquifer Summary Variables}

There is no data required for this keyword and there is no terminating
``/'' for this keyword.

\subsubsection{Examples}\label{examples-64}

\begin{verbatim}
-- ==============================================================================
--
-- SUMMARY SECTION
--
-- ==============================================================================
SUMMARY
--
--       EXPORT STANDARD SUMMARY VARIABLE VECTORS TO FILE
--
ALL
--
--       ACTIVATE COLUMNAR SUMMARY DATA REPORTING OPTION
--
RUNSUM
--
--       ACTIVATE SUMMARY DATA RSM FILE OUTPUT OPTION
--
SEPARATE
\end{verbatim}

Note the SEPARATE keyword is not required for OPM Flow as this is the
default behavior; however, it is probably good practice to include it if
the same input decks are being run with commercial simulator.

\subsection{COLUMNS -- Define Input File Column
Margins}\label{columns-define-input-file-column-margins-7}

The COLUMNS keyword defines the input file column margins; characters
outside the margins are ignored by the input parser.

See COLUMNS -- Define Input File Column Margins in the GLOBAL section
for a full description.

\subsection{DATE -- Activate the DATE Option for the SUMMARY
File}\label{date-activate-the-date-option-for-the-summary-file}

\begin{longtable}[]{@{}
  >{\raggedright\arraybackslash}p{(\linewidth - 14\tabcolsep) * \real{0.1250}}
  >{\raggedright\arraybackslash}p{(\linewidth - 14\tabcolsep) * \real{0.1250}}
  >{\raggedright\arraybackslash}p{(\linewidth - 14\tabcolsep) * \real{0.1250}}
  >{\raggedright\arraybackslash}p{(\linewidth - 14\tabcolsep) * \real{0.1250}}
  >{\raggedright\arraybackslash}p{(\linewidth - 14\tabcolsep) * \real{0.1250}}
  >{\raggedright\arraybackslash}p{(\linewidth - 14\tabcolsep) * \real{0.1250}}
  >{\raggedright\arraybackslash}p{(\linewidth - 14\tabcolsep) * \real{0.1250}}
  >{\raggedright\arraybackslash}p{(\linewidth - 14\tabcolsep) * \real{0.1250}}@{}}
\toprule\noalign{}
\begin{minipage}[b]{\linewidth}\raggedright
RUNSPEC
\end{minipage} & \begin{minipage}[b]{\linewidth}\raggedright
GRID
\end{minipage} & \begin{minipage}[b]{\linewidth}\raggedright
EDIT
\end{minipage} & \begin{minipage}[b]{\linewidth}\raggedright
PROPS
\end{minipage} & \begin{minipage}[b]{\linewidth}\raggedright
REGIONS
\end{minipage} & \begin{minipage}[b]{\linewidth}\raggedright
SOLUTION
\end{minipage} & \begin{minipage}[b]{\linewidth}\raggedright
SUMMARY
\end{minipage} & \begin{minipage}[b]{\linewidth}\raggedright
SCHEDULE
\end{minipage} \\
\midrule\noalign{}
\endhead
\bottomrule\noalign{}
\endlastfoot
\end{longtable}

\subsubsection{Description}\label{description-848}

This keyword activates the writing out of the date of each time step to
the SUMMARY file. Normally only the time in days and decimal years are
written out to the SUMMARY file, activating the DATE option also results
in the DATE being written out to the SUMMARY file as well. This option
is normally used when the RUNSUM keyword in the SUMMARY section has been
activated to produce a RSM file.

There is no data required for this keyword and there is no terminating
``/'' for this keyword.

\subsubsection{Example}\label{example-491}

The following example shows an example RSM file output when the DATE
option has not been activated.

\begin{verbatim}
 -------------------------------------------------------------------------------
 SUMMARY OF RUN NO-DATE-KEYWORD
 -------------------------------------------------------------------------------
 TIME         YEARS        FPR          FOEW         FOPR         FOPT
 DAYS         YEARS        PSIA                      STB/DAY      STB


 -------------------------------------------------------------------------------
        0            0     4467.125            0            0            0
 1.000000     0.002738     4466.943     0.000239     3235.662     3235.662
 31.00000     0.084873     4464.476     0.007407     3230.117     100256.4
 60.00000     0.164271     4462.717     0.014291     3193.902     193421.5
 91.00000     0.249144     4460.813     0.021523     3127.557     291306.3
 121.0000     0.331280     4458.909     0.028362     3055.878     383879.7
 152.0000     0.416153     4456.914     0.035262     2982.212     477271.4
……………………………………………….
\end{verbatim}

And activating the SUMMARY file DATE option with:

\begin{verbatim}
--
--       ACTIVATE DATE SUMMARY FILE OPTION
--
DATE
\end{verbatim}

Results in the following example RSM file output.

\begin{verbatim}
 -------------------------------------------------------------------------------
 SUMMARY OF RUN WITH-DATE-KEYWORD
 -------------------------------------------------------------------------------
 DATE         YEARS    DAY  MONTH  YEAR  FPR          FOEW         FOPR
              YEARS                      PSIA                      STB/DAY


 -------------------------------------------------------------------------------
  1-JAN-98           0  19   10    1992  4467.125            0            0
  2-JAN-98    0.002738  20   10    1992  4466.943     0.000239     3235.662
 31-JAN-98    0.084873  21   10    1992  4464.476     0.007407     3230.117
 28-FEB-98    0.164271  24   10    1992  4462.717     0.014291     3193.902
 31-MAR-98    0.249144  28   10    1992  4460.813     0.021523     3127.557
 30-APR-98    0.331280   3   11    1992  4458.909     0.028362     3055.878
 31-MAY-98    0.416153  14   11    1992  4456.914     0.035262     2982.212
……………………………………………….
\end{verbatim}

\subsection{DEBUG -- Define the Debug Data to be Printed to
File}\label{debug-define-the-debug-data-to-be-printed-to-file-7}

This keyword defines the debug data to be written to the debug file
(*.DBG). This keyword is not supported by OPM Flow but has no effect on
the results so it will be ignored.

See DEBUG -- Define the Debug Data to be Printed to File in the GLOBAL
section for a full description.

\subsection{ECHO -- Activate Echoing of User Input Files to the Print
File}\label{echo-activate-echoing-of-user-input-files-to-the-print-file-7}

Turns on echoing of all the input files to the print file; note that
this keyword is activated by default and can subsequently be switched
off by the NOECHO activation keyword.

See ECHO -- Activate Echoing of User Input Files to the Print File in
the GLOBAL section for a full description.

\subsection{END -- Define the End of the Input
File}\label{end-define-the-end-of-the-input-file-7}

This keyword marks the end of the input file and can occur in any
section. Any keywords and data after the END keyword are ignored.

See END -- Define the End of the Input File in the GLOBAL section for a
full description.

\subsection{ENDINC -- Define the End of an Include
File}\label{endinc-define-the-end-of-an-include-file-7}

This keyword marks the end of an include file specified on the INCLUDE
keyword. When the ENDINC keyword is encountered in the INCLUDE file,
input data is read from the next keyword in the current file. Any
keywords and data after the ENDINC keyword in the INCLUDE file are
ignored.

See ENDINC -- Define the End of an Include File in the GLOBAL section
for a full description.

\subsection{ENDSKIP -- Deactivate Skipping of Keywords and Input
Data}\label{endskip-deactivate-skipping-of-keywords-and-input-data-7}

The ENDSKIP keyword deactivates the skipping of keywords that was
activated by the SKIP, SKIP100, or SKIP300 keywords. Each SKIP keyword
should be paired with an ENDSKIP keyword.

See ENDSKIP -- Deactivate Skipping of Keywords and Input Data in the
GLOBAL section for a full description.

\subsection{EXCEL -- Activate the EXCEL Option for the SUMMARY
File}\label{excel-activate-the-excel-option-for-the-summary-file}

\begin{longtable}[]{@{}
  >{\raggedright\arraybackslash}p{(\linewidth - 14\tabcolsep) * \real{0.1250}}
  >{\raggedright\arraybackslash}p{(\linewidth - 14\tabcolsep) * \real{0.1250}}
  >{\raggedright\arraybackslash}p{(\linewidth - 14\tabcolsep) * \real{0.1250}}
  >{\raggedright\arraybackslash}p{(\linewidth - 14\tabcolsep) * \real{0.1250}}
  >{\raggedright\arraybackslash}p{(\linewidth - 14\tabcolsep) * \real{0.1250}}
  >{\raggedright\arraybackslash}p{(\linewidth - 14\tabcolsep) * \real{0.1250}}
  >{\raggedright\arraybackslash}p{(\linewidth - 14\tabcolsep) * \real{0.1250}}
  >{\raggedright\arraybackslash}p{(\linewidth - 14\tabcolsep) * \real{0.1250}}@{}}
\toprule\noalign{}
\begin{minipage}[b]{\linewidth}\raggedright
RUNSPEC
\end{minipage} & \begin{minipage}[b]{\linewidth}\raggedright
GRID
\end{minipage} & \begin{minipage}[b]{\linewidth}\raggedright
EDIT
\end{minipage} & \begin{minipage}[b]{\linewidth}\raggedright
PROPS
\end{minipage} & \begin{minipage}[b]{\linewidth}\raggedright
REGIONS
\end{minipage} & \begin{minipage}[b]{\linewidth}\raggedright
SOLUTION
\end{minipage} & \begin{minipage}[b]{\linewidth}\raggedright
SUMMARY
\end{minipage} & \begin{minipage}[b]{\linewidth}\raggedright
SCHEDULE
\end{minipage} \\
\midrule\noalign{}
\endhead
\bottomrule\noalign{}
\endlastfoot
\end{longtable}

\subsubsection{Description}\label{description-849}

This keyword activates the writing out of the RSM file data in a format
that can easily be loaded into Microsoft's EXCEL spreadsheet program or
LibreOffice's CALC spreadsheet program. The RSM file output is activated
by the RUNSUM keyword in the SUMMARY section.

There is no data required for this keyword and there is no terminating
``/'' for this keyword.

This keyword is not supported by OPM Flow but has no effect on the
results so it will be ignored.

\subsubsection{Examples}\label{examples-65}

\begin{verbatim}
-- ==============================================================================
--
-- SUMMARY SECTION
--
-- ==============================================================================
SUMMARY
--
--       ACTIVATE COLUMNAR SUMMARY DATA REPORTING OPTION
--
RUNSUM
--
--       ACTIVATE EXCEL SUMMARY FILE OPTION
--
EXCEL
\end{verbatim}

The above example activates the SUMMARY file EXCEL option for directly
loading the RSM file into either Microsoft's EXCEL or LibreOffice's CALC
spreadsheet programs.

\subsection{EXTRAPMS -- Activate Extrapolation Warning
Messages}\label{extrapms-activate-extrapolation-warning-messages-7}

The EXTRAPMS keyword activates extrapolation warning messages for when
OPM Flow extrapolates the PVT or VFP tables. Frequent extrapolation
warning messages should be investigated and resolved as this would
indicate possible incorrect data and may result in the simulator
extrapolating to unrealistic values.

See EXTRAPMS -- Activate Extrapolation Warning Messages in the GLOBAL
section for a full description.

\begin{verbatim}
\end{verbatim}

\subsection{FMWSET -- Export Well Status Vectors for the Field to
File}\label{fmwset-export-well-status-vectors-for-the-field-to-file}

\begin{longtable}[]{@{}
  >{\raggedright\arraybackslash}p{(\linewidth - 14\tabcolsep) * \real{0.1250}}
  >{\raggedright\arraybackslash}p{(\linewidth - 14\tabcolsep) * \real{0.1250}}
  >{\raggedright\arraybackslash}p{(\linewidth - 14\tabcolsep) * \real{0.1250}}
  >{\raggedright\arraybackslash}p{(\linewidth - 14\tabcolsep) * \real{0.1250}}
  >{\raggedright\arraybackslash}p{(\linewidth - 14\tabcolsep) * \real{0.1250}}
  >{\raggedright\arraybackslash}p{(\linewidth - 14\tabcolsep) * \real{0.1250}}
  >{\raggedright\arraybackslash}p{(\linewidth - 14\tabcolsep) * \real{0.1250}}
  >{\raggedright\arraybackslash}p{(\linewidth - 14\tabcolsep) * \real{0.1250}}@{}}
\toprule\noalign{}
\begin{minipage}[b]{\linewidth}\raggedright
RUNSPEC
\end{minipage} & \begin{minipage}[b]{\linewidth}\raggedright
GRID
\end{minipage} & \begin{minipage}[b]{\linewidth}\raggedright
EDIT
\end{minipage} & \begin{minipage}[b]{\linewidth}\raggedright
PROPS
\end{minipage} & \begin{minipage}[b]{\linewidth}\raggedright
REGIONS
\end{minipage} & \begin{minipage}[b]{\linewidth}\raggedright
SOLUTION
\end{minipage} & \begin{minipage}[b]{\linewidth}\raggedright
SUMMARY
\end{minipage} & \begin{minipage}[b]{\linewidth}\raggedright
SCHEDULE
\end{minipage} \\
\midrule\noalign{}
\endhead
\bottomrule\noalign{}
\endlastfoot
\end{longtable}

\subsubsection{Description}\label{description-850}

This keyword is similar to the ALL keyword in the SUMMARY section, in
that it results in a group of summary variables to be written out to the
SUMMARY file. In this case the keyword activates the writing out of a
set of data vectors that give the production and injections status of
all the wells in the model, as well as the number of wells in the
drilling queue and the number of workover events occurring within a time
step. Both instantaneous and cumulative well counts and events are
written out as listed in Table 11.28.

There is no data required for this keyword.

This keyword is not supported by OPM Flow but has no effect on the
results so it will be ignored.

\begin{longtable}[]{@{}
  >{\raggedright\arraybackslash}p{(\linewidth - 8\tabcolsep) * \real{0.2000}}
  >{\raggedright\arraybackslash}p{(\linewidth - 8\tabcolsep) * \real{0.2000}}
  >{\raggedright\arraybackslash}p{(\linewidth - 8\tabcolsep) * \real{0.2000}}
  >{\raggedright\arraybackslash}p{(\linewidth - 8\tabcolsep) * \real{0.2000}}
  >{\raggedright\arraybackslash}p{(\linewidth - 8\tabcolsep) * \real{0.2000}}@{}}
\toprule\noalign{}
\begin{minipage}[b]{\linewidth}\raggedright
Field and Group Well Status Summary Variables
\end{minipage} & \begin{minipage}[b]{\linewidth}\raggedright
\end{minipage} & \begin{minipage}[b]{\linewidth}\raggedright
\end{minipage} & \begin{minipage}[b]{\linewidth}\raggedright
\end{minipage} & \begin{minipage}[b]{\linewidth}\raggedright
\end{minipage} \\
\midrule\noalign{}
\endhead
\bottomrule\noalign{}
\endlastfoot
Variable & Root & Field & Group & Comment \\
Number of abandoned injection wells & MWIA & FMWIA & GMWIA & \\
Number of abandoned production wells & MWPA & FMWPA & GMWPA & \\
Number of drilling events in total & MWDT & FMWDT & GMWDT & \\
Number of drilling events this time step & MWDR & FMWDR & GMWDR & \\
Number of injection wells currently flowing & MWIN & FMWIN & GMWIN & \\
Number of injectors on group control & MWIG & FMWIG & GMWIG & \\
Number of injectors on own reservoir volume rate limit control & MWIV &
FMWIV & GMWIV & \\
Number of injectors on own surface rate limit control & MWIS & FMWIS &
GMWIS & \\
Number of injectors on pressure control & MWIP & FMWIP & GMWIP & \\
Number of producers controlled by own oil rate limit & MWPO & FMWPO &
GMWPO & \\
Number of producers on group control & MWPG & FMWPG & GMWPG & \\
Number of producers on own reservoir volume rate limit control & MWPV &
FMWPV & GMWPV & \\
Number of producers on own surface rate limit control & MWPS & FMWPS &
GMWPS & \\
Number of producers on pressure control & MWPP & FMWPP & GMWPP & \\
Number of producers using artificial lift (with ALQ \textgreater{} 0.0)
& MWPL & FMWPL & GMWPL & \\
Number of production wells currently flowing & MWPR & FMWPR & GMWPR & \\
Number of unused injection wells & MWIU & FMWIU & GMWIU & \\
Number of unused production wells & MWPU & FMWPU & GMWPU & \\
Number of workover events in total & MWWT & FMWWT & GMWWT & \\
Number of workover events this time step. & MWWO & FMWWO & GMWWO & \\
Total number of injection wells & MWIT & FMWIT & GMWIT & \\
Total number of production wells & MWPT & FMWPT & GMWPT & \\
Notes: & & & & \\
\end{longtable}

\emph{Table 11.28: FMWSET - Standard Field and Group Well Status Summary
Variables}

See also the GMWSET keyword in the SUMMARY schedule that has similar
functionality but at a group level.

\subsubsection{Example}\label{example-492}

\begin{verbatim}
-- ==============================================================================
--
-- SUMMARY SECTION
--
-- ==============================================================================
SUMMARY
--
--       EXPORT WELL STATUS VECTORS FOR THE FIELD TO FILE
--
FMWSET
--
--       ACTIVATE COLUMNAR SUMMARY DATA REPORTING OPTION
--
RUNSUM
--
--       ACTIVATE SUMMARY DATA RSM FILE OUTPUT OPTION
--
SEPARATE
\end{verbatim}

The above example exports the field standard well status variables to
the SUMMARY file.

\subsection{FORMFEED -- Defined the Print File Form-Feed
Character}\label{formfeed-defined-the-print-file-form-feed-character-7}

The FORMFEED keyword defines the form-feed character, or carriage
control character, for the output print (\emph{.PRT) run summary (}.RSM)
files. The keyword should be place at the very top of the input file.

See FORMFEED -- Defined the Print File Form-Feed Character in the GLOBAL
section for a dull description.

\subsection{GMWSET -- Export Well Status Vectors by Group to
File}\label{gmwset-export-well-status-vectors-by-group-to-file}

\begin{longtable}[]{@{}
  >{\raggedright\arraybackslash}p{(\linewidth - 14\tabcolsep) * \real{0.1250}}
  >{\raggedright\arraybackslash}p{(\linewidth - 14\tabcolsep) * \real{0.1250}}
  >{\raggedright\arraybackslash}p{(\linewidth - 14\tabcolsep) * \real{0.1250}}
  >{\raggedright\arraybackslash}p{(\linewidth - 14\tabcolsep) * \real{0.1250}}
  >{\raggedright\arraybackslash}p{(\linewidth - 14\tabcolsep) * \real{0.1250}}
  >{\raggedright\arraybackslash}p{(\linewidth - 14\tabcolsep) * \real{0.1250}}
  >{\raggedright\arraybackslash}p{(\linewidth - 14\tabcolsep) * \real{0.1250}}
  >{\raggedright\arraybackslash}p{(\linewidth - 14\tabcolsep) * \real{0.1250}}@{}}
\toprule\noalign{}
\begin{minipage}[b]{\linewidth}\raggedright
RUNSPEC
\end{minipage} & \begin{minipage}[b]{\linewidth}\raggedright
GRID
\end{minipage} & \begin{minipage}[b]{\linewidth}\raggedright
EDIT
\end{minipage} & \begin{minipage}[b]{\linewidth}\raggedright
PROPS
\end{minipage} & \begin{minipage}[b]{\linewidth}\raggedright
REGIONS
\end{minipage} & \begin{minipage}[b]{\linewidth}\raggedright
SOLUTION
\end{minipage} & \begin{minipage}[b]{\linewidth}\raggedright
SUMMARY
\end{minipage} & \begin{minipage}[b]{\linewidth}\raggedright
SCHEDULE
\end{minipage} \\
\midrule\noalign{}
\endhead
\bottomrule\noalign{}
\endlastfoot
\end{longtable}

\subsubsection{Description}\label{description-851}

This keyword is similar to the ALL keyword in the SUMMARY section, in
that it results in a group of summary variables to be written out to the
SUMMARY file. In this case the keyword activates the writing out of a
set of data vectors that give the production and injections status of
all the wells in the named groups, as well as the number of wells in the
drilling queue and the number of workover events occurring within a time
step for the requested groups. Both instantaneous and cumulative well
counts and events for the groups are written out as tabulated in Table
11.29.

Note that GMWSET should be followed by a list of group names enclosed in
quotes and therefore a terminating ``/'' is required to end the list of
groups. A blank list requests output for all groups.

This keyword is not supported by OPM Flow but has no effect on the
results so it will be ignored.

\begin{longtable}[]{@{}
  >{\raggedright\arraybackslash}p{(\linewidth - 8\tabcolsep) * \real{0.2000}}
  >{\raggedright\arraybackslash}p{(\linewidth - 8\tabcolsep) * \real{0.2000}}
  >{\raggedright\arraybackslash}p{(\linewidth - 8\tabcolsep) * \real{0.2000}}
  >{\raggedright\arraybackslash}p{(\linewidth - 8\tabcolsep) * \real{0.2000}}
  >{\raggedright\arraybackslash}p{(\linewidth - 8\tabcolsep) * \real{0.2000}}@{}}
\toprule\noalign{}
\begin{minipage}[b]{\linewidth}\raggedright
Field and Group Well Status Summary Variables
\end{minipage} & \begin{minipage}[b]{\linewidth}\raggedright
\end{minipage} & \begin{minipage}[b]{\linewidth}\raggedright
\end{minipage} & \begin{minipage}[b]{\linewidth}\raggedright
\end{minipage} & \begin{minipage}[b]{\linewidth}\raggedright
\end{minipage} \\
\midrule\noalign{}
\endhead
\bottomrule\noalign{}
\endlastfoot
Variable & Root & Field & Group & Comment \\
Number of abandoned injection wells & MWIA & FMWIA & GMWIA & \\
Number of abandoned production wells & MWPA & FMWPA & GMWPA & \\
Number of drilling events in total & MWDT & FMWDT & GMWDT & \\
Number of drilling events this timestep & MWDR & FMWDR & GMWDR & \\
Number of injection wells currently flowing & MWIN & FMWIN & GMWIN & \\
Number of injectors on group control & MWIG & FMWIG & GMWIG & \\
Number of injectors on own reservoir volume rate limit control & MWIV &
FMWIV & GMWIV & \\
Number of injectors on own surface rate limit control & MWIS & FMWIS &
GMWIS & \\
Number of injectors on pressure control & MWIP & FMWIP & GMWIP & \\
Number of producers controlled by own oil rate limit & MWPO & FMWPO &
GMWPO & \\
Number of producers on group control & MWPG & FMWPG & GMWPG & \\
Number of producers on own reservoir volume rate limit control & MWPV &
FMWPV & GMWPV & \\
Number of producers on own surface rate limit control & MWPS & FMWPS &
GMWPS & \\
Number of producers on pressure control & MWPP & FMWPP & GMWPP & \\
Number of producers using artificial lift (with ALQ \textgreater{} 0.0)
& MWPL & FMWPL & GMWPL & \\
Number of production wells currently flowing & MWPR & FMWPR & GMWPR & \\
Number of unused injection wells & MWIU & FMWIU & GMWIU & \\
Number of unused production wells & MWPU & FMWPU & GMWPU & \\
Number of workover events in total & MWWT & FMWWT & GMWWT & \\
Number of workover events this time step. & MWWO & FMWWO & GMWWO & \\
Total number of injection wells & MWIT & FMWIT & GMWIT & \\
Total number of production wells & MWPT & FMWPT & GMWPT & \\
Notes: & & & & \\
\end{longtable}

\emph{Table 11.29: GMWSET - Standard Field and Group Well Status Summary
Variables}

See also the FMWSET keyword in the SUMMARY schedule that has similar
functionality but at the field level.

\subsubsection{Examples}\label{examples-66}

The first example below exports all the group standard well status
variables to the SUMMARY file.

\begin{verbatim}
-- ==============================================================================
--
-- SUMMARY SECTION
--
-- ==============================================================================
SUMMARY
--
--       EXPORT WELL STATUS VECTORS FOR NAMED GROUPS TO FILE
--
GMWSET
/
--
--       ACTIVATE COLUMNAR SUMMARY DATA REPORTING OPTION
--
RUNSUM
--
--       ACTIVATE SUMMARY DATA RSM FILE OUTPUT OPTION
--
SEPARATE
\end{verbatim}

The second example exports all the group standard well status variables
for just the PLAT1 and PLT2 groups only to the SUMMARY file.

\begin{verbatim}
--
--       EXPORT WELL STATUS VECTORS FOR NAMED GROUPS TO FILE
--
GMWSET
‘PLAT1’ ‘PLAT2’
/
--
--       ACTIVATE COLUMNAR SUMMARY DATA REPORTING OPTION
--
RUNSUM
--
--       ACTIVATE SUMMARY DATA RSM FILE OUTPUT OPTION
--
SEPARATE
\end{verbatim}

\subsection{INCLUDE -- Load Another Data File at the Current
Position}\label{include-load-another-data-file-at-the-current-position-7}

The INCLUDE keyword informs OPM Flow to continue reading input data from
the specified INCLUDE file. When the end of the INCLUDE file is reached,
or the ENDINC keyword is encountered in the included file, input data is
read from the next keyword in the current file.

See INCLUDE -- Load Another Data File at the Current Position in the
GLOBAL section for a full description.

\subsection{MESSAGE -- Output User
Message}\label{message-output-user-message-7}

The MESSAGE keyword outputs a user message to the terminal, as well as
to the print (\emph{.PRT) and debug (}.DBG) files. Note this is
different to the MESSAGES keyword, that defines OPM Flows message print
limits and stop limits generated by the simulator.

See MESSAGE -- Output User Message in the GLOBAL section for a full
description.

\subsection{MESSAGES -- Define Message Print Limits and Stop
Limits}\label{messages-define-message-print-limits-and-stop-limits-7}

The MESSAGES keyword defines the print and stops levels for various
messages. The ``print limits'' set the maximum number of messages that
will be printed, after which no more messages will be printed and the
``stop limits'' terminate the run when these limits are exceeded. There
are six levels of message that increase in severity from informative all
the way to programming errors, as outlined in Table 4.5.

See MESSAGES -- Define Message Print Limits and Stop Limits in the
GLOBAL section for a full description.

\subsection{MONITOR -- Activate Output of the Monitoring Data and
File}\label{monitor-activate-output-of-the-monitoring-data-and-file-1}

The MONITOR keyword activates the writing out of the run time monitoring
information used by post-processing graphics software to display run
time information, for example the simulated production and injection
rates and cumulative values. OPM Flow does not have this functionality.

See MONITOR -- Activate Output of the Monitoring Data and File in the
RUNSPEC section for a full description.

\subsection{NARROW -- Activate Run Summary Narrow Column Output
Option}\label{narrow-activate-run-summary-narrow-column-output-option}

\begin{longtable}[]{@{}
  >{\raggedright\arraybackslash}p{(\linewidth - 14\tabcolsep) * \real{0.1250}}
  >{\raggedright\arraybackslash}p{(\linewidth - 14\tabcolsep) * \real{0.1250}}
  >{\raggedright\arraybackslash}p{(\linewidth - 14\tabcolsep) * \real{0.1250}}
  >{\raggedright\arraybackslash}p{(\linewidth - 14\tabcolsep) * \real{0.1250}}
  >{\raggedright\arraybackslash}p{(\linewidth - 14\tabcolsep) * \real{0.1250}}
  >{\raggedright\arraybackslash}p{(\linewidth - 14\tabcolsep) * \real{0.1250}}
  >{\raggedright\arraybackslash}p{(\linewidth - 14\tabcolsep) * \real{0.1250}}
  >{\raggedright\arraybackslash}p{(\linewidth - 14\tabcolsep) * \real{0.1250}}@{}}
\toprule\noalign{}
\begin{minipage}[b]{\linewidth}\raggedright
RUNSPEC
\end{minipage} & \begin{minipage}[b]{\linewidth}\raggedright
GRID
\end{minipage} & \begin{minipage}[b]{\linewidth}\raggedright
EDIT
\end{minipage} & \begin{minipage}[b]{\linewidth}\raggedright
PROPS
\end{minipage} & \begin{minipage}[b]{\linewidth}\raggedright
REGIONS
\end{minipage} & \begin{minipage}[b]{\linewidth}\raggedright
SOLUTION
\end{minipage} & \begin{minipage}[b]{\linewidth}\raggedright
SUMMARY
\end{minipage} & \begin{minipage}[b]{\linewidth}\raggedright
SCHEDULE
\end{minipage} \\
\midrule\noalign{}
\endhead
\bottomrule\noalign{}
\endlastfoot
\end{longtable}

\subsubsection{Description}\label{description-852}

The NARROW keyword activates the Run Summary Narrow Column Output
option, for when printed SUMMARY data has been requested by the RUNSUM
keyword in the SUMMARY section. The option increases the number of
columns ``printed on the page''.

There is no data required for this keyword and there is no terminating
``/'' for this keyword.

This keyword is not supported by OPM Flow but would change the results
if supported so the simulation will be stopped.

\subsection{NEWTON -- Activate Newton Iteration SUMMARY
Output}\label{newton-activate-newton-iteration-summary-output}

\begin{longtable}[]{@{}
  >{\raggedright\arraybackslash}p{(\linewidth - 14\tabcolsep) * \real{0.1250}}
  >{\raggedright\arraybackslash}p{(\linewidth - 14\tabcolsep) * \real{0.1250}}
  >{\raggedright\arraybackslash}p{(\linewidth - 14\tabcolsep) * \real{0.1250}}
  >{\raggedright\arraybackslash}p{(\linewidth - 14\tabcolsep) * \real{0.1250}}
  >{\raggedright\arraybackslash}p{(\linewidth - 14\tabcolsep) * \real{0.1250}}
  >{\raggedright\arraybackslash}p{(\linewidth - 14\tabcolsep) * \real{0.1250}}
  >{\raggedright\arraybackslash}p{(\linewidth - 14\tabcolsep) * \real{0.1250}}
  >{\raggedright\arraybackslash}p{(\linewidth - 14\tabcolsep) * \real{0.1250}}@{}}
\toprule\noalign{}
\begin{minipage}[b]{\linewidth}\raggedright
RUNSPEC
\end{minipage} & \begin{minipage}[b]{\linewidth}\raggedright
GRID
\end{minipage} & \begin{minipage}[b]{\linewidth}\raggedright
EDIT
\end{minipage} & \begin{minipage}[b]{\linewidth}\raggedright
PROPS
\end{minipage} & \begin{minipage}[b]{\linewidth}\raggedright
REGIONS
\end{minipage} & \begin{minipage}[b]{\linewidth}\raggedright
SOLUTION
\end{minipage} & \begin{minipage}[b]{\linewidth}\raggedright
SUMMARY
\end{minipage} & \begin{minipage}[b]{\linewidth}\raggedright
SCHEDULE
\end{minipage} \\
\midrule\noalign{}
\endhead
\bottomrule\noalign{}
\endlastfoot
\end{longtable}

\subsubsection{Description}\label{description-853}

This keyword activates the writing out of the Newton iteration vector
(the number of non-linear iterations per time step) to the SUMMARY file,
and the RSM file if the RSM file option has been requested by the RUNSUM
keyword in the SUMMARY section,

There is no data required for this keyword and there is no terminating
``/'' for this keyword.

Although the keyword is recognized by OPM Flow only zeros are written to
the SUMMARY file.

\subsubsection{Example}\label{example-493}

\begin{verbatim}
-- ==============================================================================
--
-- SUMMARY SECTION
--
-- ==============================================================================
SUMMARY
--
--       ACTIVATE COLUMNAR SUMMARY DATA REPORTING OPTION
--
RUNSUM
--
--       ACTIVATE SUMMARY DATA RSM FILE OUTPUT OPTION
--
SEPARATE       --
--       ACTIVATE NEWTON ITERATION SUMMARY OUTPUT
--
NEWTON
\end{verbatim}

The above example actives the writing out of the Newton iteration vector
to the SUMMARY file.

\subsection{NMESSAGE -- Export Cumulative Message Summary Variables to
File}\label{nmessage-export-cumulative-message-summary-variables-to-file}

\begin{longtable}[]{@{}
  >{\raggedright\arraybackslash}p{(\linewidth - 14\tabcolsep) * \real{0.1250}}
  >{\raggedright\arraybackslash}p{(\linewidth - 14\tabcolsep) * \real{0.1250}}
  >{\raggedright\arraybackslash}p{(\linewidth - 14\tabcolsep) * \real{0.1250}}
  >{\raggedright\arraybackslash}p{(\linewidth - 14\tabcolsep) * \real{0.1250}}
  >{\raggedright\arraybackslash}p{(\linewidth - 14\tabcolsep) * \real{0.1250}}
  >{\raggedright\arraybackslash}p{(\linewidth - 14\tabcolsep) * \real{0.1250}}
  >{\raggedright\arraybackslash}p{(\linewidth - 14\tabcolsep) * \real{0.1250}}
  >{\raggedright\arraybackslash}p{(\linewidth - 14\tabcolsep) * \real{0.1250}}@{}}
\toprule\noalign{}
\begin{minipage}[b]{\linewidth}\raggedright
RUNSPEC
\end{minipage} & \begin{minipage}[b]{\linewidth}\raggedright
GRID
\end{minipage} & \begin{minipage}[b]{\linewidth}\raggedright
EDIT
\end{minipage} & \begin{minipage}[b]{\linewidth}\raggedright
PROPS
\end{minipage} & \begin{minipage}[b]{\linewidth}\raggedright
REGIONS
\end{minipage} & \begin{minipage}[b]{\linewidth}\raggedright
SOLUTION
\end{minipage} & \begin{minipage}[b]{\linewidth}\raggedright
SUMMARY
\end{minipage} & \begin{minipage}[b]{\linewidth}\raggedright
SCHEDULE
\end{minipage} \\
\midrule\noalign{}
\endhead
\bottomrule\noalign{}
\endlastfoot
\end{longtable}

\subsubsection{Description}\label{description-854}

This keyword activates the writing out of a standard set of summary OPM
Flow simulation performance summary variables to the SUMMARY
(\emph{.SMSPEC and }.UNSMRY) and RSM (*.RSM) files, namely the number of
messages written per message class. Table 11.30 lists the summary
variables written out by the NMESSAGE keyword.

The keyword is recognized by OPM Flow but none of the variables are
currently supported.

\begin{longtable}[]{@{}
  >{\raggedright\arraybackslash}p{(\linewidth - 4\tabcolsep) * \real{0.3333}}
  >{\raggedright\arraybackslash}p{(\linewidth - 4\tabcolsep) * \real{0.3333}}
  >{\raggedright\arraybackslash}p{(\linewidth - 4\tabcolsep) * \real{0.3333}}@{}}
\toprule\noalign{}
\begin{minipage}[b]{\linewidth}\raggedright
OPM Flow Simulator Performance Summary Variables Cumulative Message
Variables
\end{minipage} & \begin{minipage}[b]{\linewidth}\raggedright
\end{minipage} & \begin{minipage}[b]{\linewidth}\raggedright
\end{minipage} \\
\midrule\noalign{}
\endhead
\bottomrule\noalign{}
\endlastfoot
Variable Description & Variable & Comment \\
Messages - Cumulative number of BUG messages. & MSUMBUG &
Unsupported. \\
Messages - Cumulative number of COMMENT messages. & MSUMCOMM &
Unsupported. \\
Messages - Cumulative number of ERROR messages. & MSUMERR &
Unsupported. \\
Messages - Cumulative number of MESSAGES messages. & MSUMMESS &
Unsupported. \\
Messages - Cumulative number of PROBLEM messages. & MSUMPROB &
Unsupported. \\
Messages - Cumulative number of WARNING messages. & MSUMWARN &
Unsupported. \\
Notes: & & \\
\end{longtable}

\emph{Table 11.30: Simulator Performance Summary Variables (Cumulative
Messages)}

\subsubsection{Example}\label{example-494}

\begin{verbatim}
-- ==============================================================================
--
-- SUMMARY SECTION
--
-- ==============================================================================
SUMMARY
--
--       EXPORT PERFORMANCE CUMULATIVE MESSAGE VARIABLE VECTORS TO FILE
--
NMESSAGE
--
--       ACTIVATE COLUMNAR SUMMARY DATA REPORTING OPTION
--
RUNSUM
--
--       ACTIVATE SUMMARY DATA RSM FILE OUTPUT OPTION
--
SEPARATE
\end{verbatim}

Note the SEPARATE keyword is not required for OPM Flow as this is the
default behavior; however, it is probably good practice to include it if
the same input decks are being run with the commercial simulator.

\subsection{NOECHO -- Deactivate Echoing of User Input Files to the
Print
File}\label{noecho-deactivate-echoing-of-user-input-files-to-the-print-file-7}

Turns off echoing of all the input files to the print file. Note by
default echoing of the inputs files is active. but can subsequently be
switched off by the NOECHO activation keyword.

See NOECHO -- Deactivate Echoing of User Input Files to the Print File
in the GLOBAL section for a full description.

\subsection{NOMONITO -- Deactivate Output of the Monitoring Data and
File}\label{nomonito-deactivate-output-of-the-monitoring-data-and-file-1}

The NOMONITO keyword deactivates the writing out of the run time
monitoring information used by post-processing graphics software to
display run time information, for example the simulated production and
injection rates and cumulative values. OPM Flow does not have this
functionality.

See NOMONITO -- Deactivate Output of the Monitoring Data and File in the
RUNSPEC section for a full description.

\subsection{NOWARN -- Deactivate Warning
Messages}\label{nowarn-deactivate-warning-messages-7}

Turns off warning messages to be printed to the print file; note that
this keyword is deactivated by default and can subsequently be switched
off by the WARN activation keyword. The warning messages may be turned
on and off using keywords WARN and NOWARN.

See NOWARN -- Deactivate Warning Messages in the GLOBAL section for a
full description.

\subsection{OFM -- Activate OFM File Output of the SUMMARY
Data}\label{ofm-activate-ofm-file-output-of-the-summary-data}

\begin{longtable}[]{@{}
  >{\raggedright\arraybackslash}p{(\linewidth - 14\tabcolsep) * \real{0.1250}}
  >{\raggedright\arraybackslash}p{(\linewidth - 14\tabcolsep) * \real{0.1250}}
  >{\raggedright\arraybackslash}p{(\linewidth - 14\tabcolsep) * \real{0.1250}}
  >{\raggedright\arraybackslash}p{(\linewidth - 14\tabcolsep) * \real{0.1250}}
  >{\raggedright\arraybackslash}p{(\linewidth - 14\tabcolsep) * \real{0.1250}}
  >{\raggedright\arraybackslash}p{(\linewidth - 14\tabcolsep) * \real{0.1250}}
  >{\raggedright\arraybackslash}p{(\linewidth - 14\tabcolsep) * \real{0.1250}}
  >{\raggedright\arraybackslash}p{(\linewidth - 14\tabcolsep) * \real{0.1250}}@{}}
\toprule\noalign{}
\begin{minipage}[b]{\linewidth}\raggedright
RUNSPEC
\end{minipage} & \begin{minipage}[b]{\linewidth}\raggedright
GRID
\end{minipage} & \begin{minipage}[b]{\linewidth}\raggedright
EDIT
\end{minipage} & \begin{minipage}[b]{\linewidth}\raggedright
PROPS
\end{minipage} & \begin{minipage}[b]{\linewidth}\raggedright
REGIONS
\end{minipage} & \begin{minipage}[b]{\linewidth}\raggedright
SOLUTION
\end{minipage} & \begin{minipage}[b]{\linewidth}\raggedright
SUMMARY
\end{minipage} & \begin{minipage}[b]{\linewidth}\raggedright
SCHEDULE
\end{minipage} \\
\midrule\noalign{}
\endhead
\bottomrule\noalign{}
\endlastfoot
\end{longtable}

\subsubsection{Description}\label{description-855}

This keyword activates the writing out of the SUMMARY file data in the
Oil Field Manager (``OFM'') file format to enable the simulated data to
be directly loaded into Oil Field Manager.

There is no data required for this keyword and there is no terminating
``/'' for this keyword.

See also the EXCEL, RPTONLY, RUNSUM and SEPARATE keywords in the SUMMARY
section.

This keyword is not supported by OPM Flow but would change the results
if supported so the simulation will be stopped.

\subsection{PERFORMA -- Export Standard Simulator Performance Summary
Variables to
File}\label{performa-export-standard-simulator-performance-summary-variables-to-file}

\begin{longtable}[]{@{}
  >{\raggedright\arraybackslash}p{(\linewidth - 14\tabcolsep) * \real{0.1250}}
  >{\raggedright\arraybackslash}p{(\linewidth - 14\tabcolsep) * \real{0.1250}}
  >{\raggedright\arraybackslash}p{(\linewidth - 14\tabcolsep) * \real{0.1250}}
  >{\raggedright\arraybackslash}p{(\linewidth - 14\tabcolsep) * \real{0.1250}}
  >{\raggedright\arraybackslash}p{(\linewidth - 14\tabcolsep) * \real{0.1250}}
  >{\raggedright\arraybackslash}p{(\linewidth - 14\tabcolsep) * \real{0.1250}}
  >{\raggedright\arraybackslash}p{(\linewidth - 14\tabcolsep) * \real{0.1250}}
  >{\raggedright\arraybackslash}p{(\linewidth - 14\tabcolsep) * \real{0.1250}}@{}}
\toprule\noalign{}
\begin{minipage}[b]{\linewidth}\raggedright
RUNSPEC
\end{minipage} & \begin{minipage}[b]{\linewidth}\raggedright
GRID
\end{minipage} & \begin{minipage}[b]{\linewidth}\raggedright
EDIT
\end{minipage} & \begin{minipage}[b]{\linewidth}\raggedright
PROPS
\end{minipage} & \begin{minipage}[b]{\linewidth}\raggedright
REGIONS
\end{minipage} & \begin{minipage}[b]{\linewidth}\raggedright
SOLUTION
\end{minipage} & \begin{minipage}[b]{\linewidth}\raggedright
SUMMARY
\end{minipage} & \begin{minipage}[b]{\linewidth}\raggedright
SCHEDULE
\end{minipage} \\
\midrule\noalign{}
\endhead
\bottomrule\noalign{}
\endlastfoot
\end{longtable}

\subsubsection{Description}\label{description-856}

The PERFORMA keyword activates the writing out of a standard set of OPM
Flow simulation numerical performance summary variables to the SUMMARY
(\emph{.SMSPEC and .}UNSMRY) and RSM (*.RSM) files. Table 11.31 lists
the summary variables written out by the keyword.

Note that not all these variables are available in OPM Flow; however,
the simulator will issue a warning message if this is indeed the case.
It is anticipated that the number of recognized summary variables will
increase in future releases of OPM Flow.

\begin{longtable}[]{@{}
  >{\raggedright\arraybackslash}p{(\linewidth - 4\tabcolsep) * \real{0.3333}}
  >{\raggedright\arraybackslash}p{(\linewidth - 4\tabcolsep) * \real{0.3333}}
  >{\raggedright\arraybackslash}p{(\linewidth - 4\tabcolsep) * \real{0.3333}}@{}}
\toprule\noalign{}
\begin{minipage}[b]{\linewidth}\raggedright
OPM Flow Simulator Performance Summary Variables Numerical Performance
Variables
\end{minipage} & \begin{minipage}[b]{\linewidth}\raggedright
\end{minipage} & \begin{minipage}[b]{\linewidth}\raggedright
\end{minipage} \\
\midrule\noalign{}
\endhead
\bottomrule\noalign{}
\endlastfoot
Variable Description & Variable & Comment \\
Elapsed - Elapsed time in seconds. & ELAPSED & No data written to
file. \\
Iterations - Number linear iterations for each time step. & MLINEARS
& \\
Iterations - Cumulative number of linear iterations. & MSUMLINS & \\
Iterations - Cumulative number of Newton iterations. & MSUMNEWT & \\
Iterations - Number of Newton iterations used for each time step. &
NEWTON & \\
Iterations - Average number of linear iterations per Newton iteration
for each time step. & NLINEARS & For runs with LGRs, LLINEARS will
automatically be exported for each LGR. \\
Iterations - Maximum number of linear iterations in the Newton
iterations per time step. & NLINSMAX & \\
Iterations - Minimum number of linear iterations in the Newton
iterations per time step. & NLINSMIN & \\
Time Step -- Criteria used to select the length of the time step. See
section 11.2.24 Option Specific Variables -- OPM Flow Simulation
Performance for the definition of of the STEPTYPE mnemonics. & STEPTYPE
& No data written to file. \\
CPU - Current CPU usage in seconds. & TCPU & Does not consider the time
taken by inter-process communications in parallel runs, whereas ELAPSED
does. Thus, for parallel jobs, ELAPSED is the most relevant time
measurement. \\
CPU - CPU time per day (or hour in lab units depending on run units
system). & TCPUDAY & No data written to file. \\
CPU - CPU time per time step in seconds. & TCPUTS & No data written to
file. \\
Elapsed - Elapsed time per linear iteration in seconds. & TELAPLIN & No
data written to file. \\
Time Step - Length of time step. & TIMESTEP & \\
Notes: & & \\
\end{longtable}

\emph{Table 11.31: Simulator Performance Summary Variables (Numerical
Performance)}

\begin{tcolorbox}[enhanced jigsaw, title=\textcolor{quarto-callout-note-color}{\faInfo}\hspace{0.5em}{Note}, coltitle=black, breakable, toprule=.15mm, colbacktitle=quarto-callout-note-color!10!white, leftrule=.75mm, opacitybacktitle=0.6, colback=white, bottomtitle=1mm, toptitle=1mm, titlerule=0mm, left=2mm, colframe=quarto-callout-note-color-frame, arc=.35mm, rightrule=.15mm, bottomrule=.15mm, opacityback=0]

If the summary vector data is unavailable then zeros are written at each
time step to the SUMMARY and RSM file.

\end{tcolorbox}

\subsubsection{Example}\label{example-495}

\begin{verbatim}
-- ==============================================================================
--
-- SUMMARY SECTION
--
-- ==============================================================================
SUMMARY
--
--       EXPORT NUMERICAL PERFORMANCE SUMMARY VARIABLES TO FILE
--
PERFORMA
--
--       ACTIVATE COLUMNAR SUMMARY DATA REPORTING OPTION
--
RUNSUM
--
--       ACTIVATE SUMMARY DATA RSM FILE OUTPUT OPTION
--
SEPARATE
\end{verbatim}

Note the SEPARATE keyword is not required for OPM Flow as this is the
default behavior; however, it is probably good practice to include it if
the same input decks are being run with commercial simulator.

\subsection{PYEND -- End the Definition of a PYINPUT
Section}\label{pyend-end-the-definition-of-a-pyinput-section-5}

The PYINPUT and PYEND keywords are a part of OPM Flow's Python scripting
facility that processes standard Python commands that can be used to
manipulate and define the simulators input parameters during processing
of the input deck. The main purpose of the facility is to script the
construction of the various keywords. A PYINPUT Definition Section is
terminated by a PYEND keyword on a separate single line.

See PYEND -- End the Definition of a PYINPUT Section in the GRID section
for a full description.

\subsection{PYINPUT -- Define the Start of a PYINPUT
Section}\label{pyinput-define-the-start-of-a-pyinput-section-5}

The PYINPUT and PYEND keywords are a part of OPM Flow's Python scripting
facility that processes standard Python commands that can be used to
manipulate and define the simulators input parameters during processing
of the input deck. The main purpose of the facility is to script the
construction of the various keywords used by the simulator. PYINPUT
declares the start of a PYINPUT Definition Section on a single separate
line, which is then followed by various standard Python commands, with
one command per line.

See PYINPUT -- Define the Start of a PYINPUT Section in the GRID section
for a full description.

\subsection{RPTONLY -- Activate the Report Time Steps Only Option for
the SUMMARY
File}\label{rptonly-activate-the-report-time-steps-only-option-for-the-summary-file}

\begin{longtable}[]{@{}
  >{\raggedright\arraybackslash}p{(\linewidth - 14\tabcolsep) * \real{0.1250}}
  >{\raggedright\arraybackslash}p{(\linewidth - 14\tabcolsep) * \real{0.1250}}
  >{\raggedright\arraybackslash}p{(\linewidth - 14\tabcolsep) * \real{0.1250}}
  >{\raggedright\arraybackslash}p{(\linewidth - 14\tabcolsep) * \real{0.1250}}
  >{\raggedright\arraybackslash}p{(\linewidth - 14\tabcolsep) * \real{0.1250}}
  >{\raggedright\arraybackslash}p{(\linewidth - 14\tabcolsep) * \real{0.1250}}
  >{\raggedright\arraybackslash}p{(\linewidth - 14\tabcolsep) * \real{0.1250}}
  >{\raggedright\arraybackslash}p{(\linewidth - 14\tabcolsep) * \real{0.1250}}@{}}
\toprule\noalign{}
\begin{minipage}[b]{\linewidth}\raggedright
RUNSPEC
\end{minipage} & \begin{minipage}[b]{\linewidth}\raggedright
GRID
\end{minipage} & \begin{minipage}[b]{\linewidth}\raggedright
EDIT
\end{minipage} & \begin{minipage}[b]{\linewidth}\raggedright
PROPS
\end{minipage} & \begin{minipage}[b]{\linewidth}\raggedright
REGIONS
\end{minipage} & \begin{minipage}[b]{\linewidth}\raggedright
SOLUTION
\end{minipage} & \begin{minipage}[b]{\linewidth}\raggedright
SUMMARY
\end{minipage} & \begin{minipage}[b]{\linewidth}\raggedright
SCHEDULE
\end{minipage} \\
\midrule\noalign{}
\endhead
\bottomrule\noalign{}
\endlastfoot
\end{longtable}

\subsubsection{Description}\label{description-857}

This keyword activates the writing out of the SUMMARY file and the RSM
file data, if the RSM file option has been requested by the RUNSUM
keyword in the SUMMARY section, at report time steps only. The default
is for the data to be written out for all time steps to the SUMMARY
files. This keyword reduces the file size at the expense of lower
resolution in the time domain.

There is no data required for this keyword and there is no terminating
``/'' for this keyword.

The option can be deactivated by the RPTONLYO keyword in the SUMMARY
section.

\subsubsection{Example}\label{example-496}

\begin{verbatim}
-- ==============================================================================
--
-- SUMMARY SECTION
--
-- ==============================================================================
SUMMARY
--
--       EXPORT STANDARD SUMMARY VARIABLE VECTORS TO FILE
--
ALL
--
--       ACTIVATE COLUMNAR SUMMARY DATA REPORTING OPTION
--
RUNSUM
--
--       ACTIVATE SUMMARY DATA RSM FILE OUTPUT OPTION
--
SEPARATE
--
--       ACTIVATE REPORT TIME STEPS ONLY SUMMARY FILE OPTION
--
RPTONLY
\end{verbatim}

The above example activates the writing out of the SUMMARY file at
report time steps only.

\subsection{RPTONLYO -- Deactivate the Report Time Steps Only Option for
the SUMMARY
File}\label{rptonlyo-deactivate-the-report-time-steps-only-option-for-the-summary-file}

\begin{longtable}[]{@{}
  >{\raggedright\arraybackslash}p{(\linewidth - 14\tabcolsep) * \real{0.1250}}
  >{\raggedright\arraybackslash}p{(\linewidth - 14\tabcolsep) * \real{0.1250}}
  >{\raggedright\arraybackslash}p{(\linewidth - 14\tabcolsep) * \real{0.1250}}
  >{\raggedright\arraybackslash}p{(\linewidth - 14\tabcolsep) * \real{0.1250}}
  >{\raggedright\arraybackslash}p{(\linewidth - 14\tabcolsep) * \real{0.1250}}
  >{\raggedright\arraybackslash}p{(\linewidth - 14\tabcolsep) * \real{0.1250}}
  >{\raggedright\arraybackslash}p{(\linewidth - 14\tabcolsep) * \real{0.1250}}
  >{\raggedright\arraybackslash}p{(\linewidth - 14\tabcolsep) * \real{0.1250}}@{}}
\toprule\noalign{}
\begin{minipage}[b]{\linewidth}\raggedright
RUNSPEC
\end{minipage} & \begin{minipage}[b]{\linewidth}\raggedright
GRID
\end{minipage} & \begin{minipage}[b]{\linewidth}\raggedright
EDIT
\end{minipage} & \begin{minipage}[b]{\linewidth}\raggedright
PROPS
\end{minipage} & \begin{minipage}[b]{\linewidth}\raggedright
REGIONS
\end{minipage} & \begin{minipage}[b]{\linewidth}\raggedright
SOLUTION
\end{minipage} & \begin{minipage}[b]{\linewidth}\raggedright
SUMMARY
\end{minipage} & \begin{minipage}[b]{\linewidth}\raggedright
SCHEDULE
\end{minipage} \\
\midrule\noalign{}
\endhead
\bottomrule\noalign{}
\endlastfoot
\end{longtable}

\subsubsection{Description}\label{description-858}

This keyword deactivates the writing out of the SUMMARY file and the RSM
file data, if the RSM file option has been requested by the RUNSUM
keyword in the SUMMARY section, at report time steps only, and switches
on writing out all the time steps to the files. This option is the
default behavior for when RPTONLY has not been activated.

There is no data required for this keyword and there is no terminating
``/'' for this keyword.

The option can be deactivated by the RPTONLY keyword in the SUMMARY
section that will switch on writing the data at every report time step
instead of every time step.

\subsubsection{Example}\label{example-497}

\begin{verbatim}
-- ==============================================================================
--
-- SUMMARY SECTION
--
-- ==============================================================================
SUMMARY
--
--       EXPORT STANDARD SUMMARY VARIABLE VECTORS TO FILE
--
ALL
--
--       ACTIVATE COLUMNAR SUMMARY DATA REPORTING OPTION
--
RUNSUM
--
--       ACTIVATE SUMMARY DATA RSM FILE OUTPUT OPTION
--
SEPARATE
--
--       DEACTIVATE REPORT TIME STEPS ONLY SUMMARY FILE OPTION
--
RPTONLYO
\end{verbatim}

The above example deactivates the writing out of the SUMMARY file at
report time steps only, and switches on writing out all the time steps
to the file.

\subsection{RPTSMRY -- Activate or Deactivate Summary List
Report}\label{rptsmry-activate-or-deactivate-summary-list-report}

\begin{longtable}[]{@{}
  >{\raggedright\arraybackslash}p{(\linewidth - 14\tabcolsep) * \real{0.1250}}
  >{\raggedright\arraybackslash}p{(\linewidth - 14\tabcolsep) * \real{0.1250}}
  >{\raggedright\arraybackslash}p{(\linewidth - 14\tabcolsep) * \real{0.1250}}
  >{\raggedright\arraybackslash}p{(\linewidth - 14\tabcolsep) * \real{0.1250}}
  >{\raggedright\arraybackslash}p{(\linewidth - 14\tabcolsep) * \real{0.1250}}
  >{\raggedright\arraybackslash}p{(\linewidth - 14\tabcolsep) * \real{0.1250}}
  >{\raggedright\arraybackslash}p{(\linewidth - 14\tabcolsep) * \real{0.1250}}
  >{\raggedright\arraybackslash}p{(\linewidth - 14\tabcolsep) * \real{0.1250}}@{}}
\toprule\noalign{}
\begin{minipage}[b]{\linewidth}\raggedright
RUNSPEC
\end{minipage} & \begin{minipage}[b]{\linewidth}\raggedright
GRID
\end{minipage} & \begin{minipage}[b]{\linewidth}\raggedright
EDIT
\end{minipage} & \begin{minipage}[b]{\linewidth}\raggedright
PROPS
\end{minipage} & \begin{minipage}[b]{\linewidth}\raggedright
REGIONS
\end{minipage} & \begin{minipage}[b]{\linewidth}\raggedright
SOLUTION
\end{minipage} & \begin{minipage}[b]{\linewidth}\raggedright
SUMMARY
\end{minipage} & \begin{minipage}[b]{\linewidth}\raggedright
SCHEDULE
\end{minipage} \\
\midrule\noalign{}
\endhead
\bottomrule\noalign{}
\endlastfoot
\end{longtable}

\subsubsection{Description}\label{description-859}

This keyword activates or deactivates a listing of all the summary
variables that are going to be written to the SUMMARY file and the RSM
file, if the RSM file option has been requested by the RUNSUM keyword in
the SUMMARY section.

This keyword is not supported by OPM Flow but has no effect on the
results so it will be ignored.

\begin{longtable}[]{@{}
  >{\raggedright\arraybackslash}p{(\linewidth - 6\tabcolsep) * \real{0.2500}}
  >{\raggedright\arraybackslash}p{(\linewidth - 6\tabcolsep) * \real{0.2500}}
  >{\raggedright\arraybackslash}p{(\linewidth - 6\tabcolsep) * \real{0.2500}}
  >{\raggedright\arraybackslash}p{(\linewidth - 6\tabcolsep) * \real{0.2500}}@{}}
\toprule\noalign{}
\begin{minipage}[b]{\linewidth}\raggedright
No.
\end{minipage} & \begin{minipage}[b]{\linewidth}\raggedright
Name
\end{minipage} & \begin{minipage}[b]{\linewidth}\raggedright
Description
\end{minipage} & \begin{minipage}[b]{\linewidth}\raggedright
Default
\end{minipage} \\
\midrule\noalign{}
\endhead
\bottomrule\noalign{}
\endlastfoot
1 & RPTSMRY & An integer value set to zero for no report, or one to
produce the report. & 0 \\
Notes: & & & \\
\end{longtable}

\emph{Table 11.32: RPTSMRY Keyword Description}

\subsubsection{Example}\label{example-498}

\begin{verbatim}
-- ==============================================================================
--
-- SUMMARY SECTION
--
-- ==============================================================================
SUMMARY
--
--       EXPORT STANDARD SUMMARY VARIABLE VECTORS TO FILE
--
ALL
--
--       ACTIVATE COLUMNAR SUMMARY DATA REPORTING OPTION
--
RUNSUM
--
--       ACTIVATE SUMMARY DATA RSM FILE OUTPUT OPTION
--
SEPARATE
--
--       ACTIVATE OR DEACTIVATE SUMMARY LIST REPORT
--
RPTSMRY
         1                                                 /
\end{verbatim}

The example switches on the summary list report.

\subsection{RUNSUM -- Activate RSM File Output of the SUMMARY
Data}\label{runsum-activate-rsm-file-output-of-the-summary-data}

\begin{longtable}[]{@{}
  >{\raggedright\arraybackslash}p{(\linewidth - 14\tabcolsep) * \real{0.1250}}
  >{\raggedright\arraybackslash}p{(\linewidth - 14\tabcolsep) * \real{0.1250}}
  >{\raggedright\arraybackslash}p{(\linewidth - 14\tabcolsep) * \real{0.1250}}
  >{\raggedright\arraybackslash}p{(\linewidth - 14\tabcolsep) * \real{0.1250}}
  >{\raggedright\arraybackslash}p{(\linewidth - 14\tabcolsep) * \real{0.1250}}
  >{\raggedright\arraybackslash}p{(\linewidth - 14\tabcolsep) * \real{0.1250}}
  >{\raggedright\arraybackslash}p{(\linewidth - 14\tabcolsep) * \real{0.1250}}
  >{\raggedright\arraybackslash}p{(\linewidth - 14\tabcolsep) * \real{0.1250}}@{}}
\toprule\noalign{}
\begin{minipage}[b]{\linewidth}\raggedright
RUNSPEC
\end{minipage} & \begin{minipage}[b]{\linewidth}\raggedright
GRID
\end{minipage} & \begin{minipage}[b]{\linewidth}\raggedright
EDIT
\end{minipage} & \begin{minipage}[b]{\linewidth}\raggedright
PROPS
\end{minipage} & \begin{minipage}[b]{\linewidth}\raggedright
REGIONS
\end{minipage} & \begin{minipage}[b]{\linewidth}\raggedright
SOLUTION
\end{minipage} & \begin{minipage}[b]{\linewidth}\raggedright
SUMMARY
\end{minipage} & \begin{minipage}[b]{\linewidth}\raggedright
SCHEDULE
\end{minipage} \\
\midrule\noalign{}
\endhead
\bottomrule\noalign{}
\endlastfoot
\end{longtable}

\subsubsection{Description}\label{description-860}

This keyword activates the writing out of the SUMMARY file data in a
columnar format to the PRT file. Normally the SEPARATE keyword in the
SUMMARY section is invoked in the same run to direct the data stream to
a separate RSM file for easy loading into other programs, for example,
Microsoft's EXCEL or LibreOffice's CALC spreadsheet programs.

There is no data required for this keyword and there is no terminating
``/'' for this keyword.

See also the EXCEL, RPTONLY and SEPARATE keywords in the SUMMARY
section.

\subsubsection{Example}\label{example-499}

\begin{verbatim}
-- ==============================================================================
--
-- SUMMARY SECTION
--
-- ==============================================================================
SUMMARY
--
--       ACTIVATE COLUMNAR SUMMARY DATA REPORTING OPTION
--
RUNSUM
--
--       ACTIVATE SUMMARY DATA RSM FILE OUTPUT OPTION
--
SEPARATE
\end{verbatim}

Note unlike the commercial simulator, OPM Flow always writes out the
data to a separate file.

\subsection{SEPARATE -- Activate the Separate RSM File Output
Option}\label{separate-activate-the-separate-rsm-file-output-option}

\begin{longtable}[]{@{}
  >{\raggedright\arraybackslash}p{(\linewidth - 14\tabcolsep) * \real{0.1250}}
  >{\raggedright\arraybackslash}p{(\linewidth - 14\tabcolsep) * \real{0.1250}}
  >{\raggedright\arraybackslash}p{(\linewidth - 14\tabcolsep) * \real{0.1250}}
  >{\raggedright\arraybackslash}p{(\linewidth - 14\tabcolsep) * \real{0.1250}}
  >{\raggedright\arraybackslash}p{(\linewidth - 14\tabcolsep) * \real{0.1250}}
  >{\raggedright\arraybackslash}p{(\linewidth - 14\tabcolsep) * \real{0.1250}}
  >{\raggedright\arraybackslash}p{(\linewidth - 14\tabcolsep) * \real{0.1250}}
  >{\raggedright\arraybackslash}p{(\linewidth - 14\tabcolsep) * \real{0.1250}}@{}}
\toprule\noalign{}
\begin{minipage}[b]{\linewidth}\raggedright
RUNSPEC
\end{minipage} & \begin{minipage}[b]{\linewidth}\raggedright
GRID
\end{minipage} & \begin{minipage}[b]{\linewidth}\raggedright
EDIT
\end{minipage} & \begin{minipage}[b]{\linewidth}\raggedright
PROPS
\end{minipage} & \begin{minipage}[b]{\linewidth}\raggedright
REGIONS
\end{minipage} & \begin{minipage}[b]{\linewidth}\raggedright
SOLUTION
\end{minipage} & \begin{minipage}[b]{\linewidth}\raggedright
SUMMARY
\end{minipage} & \begin{minipage}[b]{\linewidth}\raggedright
SCHEDULE
\end{minipage} \\
\midrule\noalign{}
\endhead
\bottomrule\noalign{}
\endlastfoot
\end{longtable}

\subsubsection{Description}\label{description-861}

This keyword activates the writing out of the SUMMARY file date in a
columnar format to the RSM file, if the RUNSUM keyword has also been
activated in the SUMMARY section. Both the SEPARATE and the RUNSUM
keywords need to be invoked. If the SEPARATE option is not activated
then the RSM output is directed to the end of the PRT file. Normally the
both the SEPARATE and RUNSUM keywords are invoked in the same run to
enable easy loading of the data into Microsoft's EXCEL or LibreOffice's
CALC spreadsheet programs.

There is no data required for this keyword and there is no terminating
``/'' for this keyword.

See also the EXCEL, RPTONLY and RUNSUM keywords in the SUMMARY section.

\subsubsection{Example}\label{example-500}

\begin{verbatim}
-- ==============================================================================
--
-- SUMMARY SECTION
--
-- ==============================================================================
SUMMARY
--
--       ACTIVATE COLUMNAR SUMMARY DATA REPORTING OPTION
--
RUNSUM
--
--       ACTIVATE SUMMARY DATA RSM FILE OUTPUT OPTION
--
SEPARATE
\end{verbatim}

Note unlike the commercial simulator, OPM Flow always writes out the
data to a separate file; however, it is probably good practice to
include it if the same input decks are being run with commercial
simulator.

\subsection{SKIP -- Activate Skipping of All Keywords and Input
Data}\label{skip-activate-skipping-of-all-keywords-and-input-data-7}

The SKIP keyword activates skipping of all keywords and input data until
the ENDSKIP keyword is encountered. All keywords between the SKIP and
ENDSKIP keywords are ignored.

See Error: Reference source not found in the GLOBAL section for a full
description.

\subsection{SKIP100 -- Activate Skipping of Black-Oil Keywords and Input
Data}\label{skip100-activate-skipping-of-black-oil-keywords-and-input-data-6}

This keyword turns on skipping of black-oil keywords until the ENDSKIP
activation keyword is encountered. All black-oil keywords between the
SKIP100 and ENDSKIP keywords are ignored by OPM Flow.

See SKIP100 -- Activate Skipping of Black-Oil Keywords and Input Data in
the GLOBAL section for a full description.

\subsection{SKIP300 -- Activate Skipping of Compositional Keywords and
Input
Data}\label{skip300-activate-skipping-of-compositional-keywords-and-input-data}

Turns on skipping of ``compositional'' keywords until the ENDSKIP
activation keyword is encountered. All ``compositional'' keywords
between the SKIP300 and ENDSKIP keywords are ignored by OPM Flow.

See Error: Reference source not found in the GLOBAL section for a full
description.

\subsection{SUMMARY -- Define the Start of the SUMMARY Section of
Keywords}\label{summary-define-the-start-of-the-summary-section-of-keywords}

\begin{longtable}[]{@{}
  >{\raggedright\arraybackslash}p{(\linewidth - 14\tabcolsep) * \real{0.1250}}
  >{\raggedright\arraybackslash}p{(\linewidth - 14\tabcolsep) * \real{0.1250}}
  >{\raggedright\arraybackslash}p{(\linewidth - 14\tabcolsep) * \real{0.1250}}
  >{\raggedright\arraybackslash}p{(\linewidth - 14\tabcolsep) * \real{0.1250}}
  >{\raggedright\arraybackslash}p{(\linewidth - 14\tabcolsep) * \real{0.1250}}
  >{\raggedright\arraybackslash}p{(\linewidth - 14\tabcolsep) * \real{0.1250}}
  >{\raggedright\arraybackslash}p{(\linewidth - 14\tabcolsep) * \real{0.1250}}
  >{\raggedright\arraybackslash}p{(\linewidth - 14\tabcolsep) * \real{0.1250}}@{}}
\toprule\noalign{}
\begin{minipage}[b]{\linewidth}\raggedright
RUNSPEC
\end{minipage} & \begin{minipage}[b]{\linewidth}\raggedright
GRID
\end{minipage} & \begin{minipage}[b]{\linewidth}\raggedright
EDIT
\end{minipage} & \begin{minipage}[b]{\linewidth}\raggedright
PROPS
\end{minipage} & \begin{minipage}[b]{\linewidth}\raggedright
REGIONS
\end{minipage} & \begin{minipage}[b]{\linewidth}\raggedright
SOLUTION
\end{minipage} & \begin{minipage}[b]{\linewidth}\raggedright
SUMMARY
\end{minipage} & \begin{minipage}[b]{\linewidth}\raggedright
SCHEDULE
\end{minipage} \\
\midrule\noalign{}
\endhead
\bottomrule\noalign{}
\endlastfoot
\end{longtable}

\subsubsection{Description}\label{description-862}

The SUMMARY activation keyword marks the end of the SOLUTION section and
the start of the SUMMARY section that defines the variables to be
written out to the SUMMARY file for reporting and plotting of grid block
data, production data, etc.

There is no data required for this keyword and there is no terminating
``/'' for this keyword.

\subsubsection{Example}\label{example-501}

\begin{verbatim}
-- ==============================================================================
--
-- SUMMARY SECTION
--
-- ==============================================================================
SUMMARY
\end{verbatim}

The above example marks the end of the SOLUTION section and the start of
the SUMMARY section in the OPM Flow data input file.

\subsection{SUMTHIN -- Define SUMMARY Data Reporting Time
Steps}\label{sumthin-define-summary-data-reporting-time-steps}

\begin{longtable}[]{@{}
  >{\raggedright\arraybackslash}p{(\linewidth - 14\tabcolsep) * \real{0.1250}}
  >{\raggedright\arraybackslash}p{(\linewidth - 14\tabcolsep) * \real{0.1250}}
  >{\raggedright\arraybackslash}p{(\linewidth - 14\tabcolsep) * \real{0.1250}}
  >{\raggedright\arraybackslash}p{(\linewidth - 14\tabcolsep) * \real{0.1250}}
  >{\raggedright\arraybackslash}p{(\linewidth - 14\tabcolsep) * \real{0.1250}}
  >{\raggedright\arraybackslash}p{(\linewidth - 14\tabcolsep) * \real{0.1250}}
  >{\raggedright\arraybackslash}p{(\linewidth - 14\tabcolsep) * \real{0.1250}}
  >{\raggedright\arraybackslash}p{(\linewidth - 14\tabcolsep) * \real{0.1250}}@{}}
\toprule\noalign{}
\begin{minipage}[b]{\linewidth}\raggedright
RUNSPEC
\end{minipage} & \begin{minipage}[b]{\linewidth}\raggedright
GRID
\end{minipage} & \begin{minipage}[b]{\linewidth}\raggedright
EDIT
\end{minipage} & \begin{minipage}[b]{\linewidth}\raggedright
PROPS
\end{minipage} & \begin{minipage}[b]{\linewidth}\raggedright
REGIONS
\end{minipage} & \begin{minipage}[b]{\linewidth}\raggedright
SOLUTION
\end{minipage} & \begin{minipage}[b]{\linewidth}\raggedright
SUMMARY
\end{minipage} & \begin{minipage}[b]{\linewidth}\raggedright
SCHEDULE
\end{minipage} \\
\midrule\noalign{}
\endhead
\bottomrule\noalign{}
\endlastfoot
\end{longtable}

\subsubsection{Description}\label{description-863}

This keyword defines a time interval for writing out the SUMMARY data to
the SUMMARY file and the RSM file, if the RUNSUM keyword has also been
activated in the SUMMARY section. Only the data for the first time step
in the time interval is written out and the other time steps are skipped
until the next time interval, except for report time steps. Note that
report time steps data are always written out regardless of the setting
on this keyword. This enables the size of the SUMMARY files to be
reduced depending on the size of the time interval. However, the keyword
will produce irregular time steps reports of the SUMMARY data.

\begin{longtable}[]{@{}
  >{\raggedright\arraybackslash}p{(\linewidth - 6\tabcolsep) * \real{0.2500}}
  >{\raggedright\arraybackslash}p{(\linewidth - 6\tabcolsep) * \real{0.2500}}
  >{\raggedright\arraybackslash}p{(\linewidth - 6\tabcolsep) * \real{0.2500}}
  >{\raggedright\arraybackslash}p{(\linewidth - 6\tabcolsep) * \real{0.2500}}@{}}
\toprule\noalign{}
\begin{minipage}[b]{\linewidth}\raggedright
No.
\end{minipage} & \begin{minipage}[b]{\linewidth}\raggedright
Name
\end{minipage} & \begin{minipage}[b]{\linewidth}\raggedright
Description
\end{minipage} & \begin{minipage}[b]{\linewidth}\raggedright
Default
\end{minipage} \\
\midrule\noalign{}
\endhead
\bottomrule\noalign{}
\endlastfoot
Field & Metric & Laboratory & \\
1 & SUMSTEP & SUMSTEP is a real positive number that defines the time
interval for which the first time step of data will be written to the
SUMMARY file (and the RSM file if RSM output has been activated). For
example, if SUMSTEP is set to 30 days, and if the simulator takes time
steps of 0, 5, 10, 16, 24, 30, 40, 45, 60, 90 days. Then the SUMMARY
data will be written out at time steps 0, 30, 40 and 60 days. & None \\
days & days & hours & \\
Notes: & & & \\
\end{longtable}

\emph{Table 11.33: SUMTHIN Keyword Description}

See also the RPTONLY keyword in the SUMMARY section that forces the
SUMMARY data to be only written out at report time steps, as oppose to
all time steps or SUMSTEP time intervals.

\subsubsection{Example}\label{example-502}

\begin{verbatim}
--
--       DEFINE SUMMARY DATA REPORTING TIME STEP INTERVAL
--
--       SUMSTEP
SUMTHIN
         30.0                                              /                                                                  /
\end{verbatim}

The above example defines the SUMMARY file time step interval to be 30
days for both field and metric units.

\subsection{WARN -- Activate Warning
Messages}\label{warn-activate-warning-messages-7}

Turns on warning messages to be printed to the print file (*.PRT); note
that this keyword is activated by default and can subsequently be
switched off by the NOWARN activation keyword. The warning messages may
be turned on and off using keywords WARN and NOWARN. OPM Flow always
prints error messages.

See WARN -- Activate Warning Messages in the GLOBAL section for a full
description.

\bookmarksetup{startatroot}

\chapter{SCHEDULE SECTION}\label{schedule-section-1}

\section{Introduction}\label{introduction-10}

The SCHEDULE section is used to specify the production and injection
targets and constraints for the entire model, advance the model through
time, and stipulate any other data that depends on time. For example,
drilling wells at a particular time, incorporating production and/or
injection rate changes as the model progresses through time, and opening
and shutting wells and well connections/completions at particular times
during the history of the reservoir. Field development and operating
strategy are also incorporated in the SCHEDULE section, for example the
timing and effect of upgrading of the field facilities liquid and water
handling capacities. All these actions must be specified in the SCHEDULE
section at the time they occur, that is all the events in the SCHEDULE
section are time dependent. This is different to the other sections of
the input deck were the order of keywords is largely unimportant; in the
SCHEDULE section it is essential because the order in which events take
place in the field have to be preserved. For instance, the keyword
defining wellhead locations via the WELSPECS keyword must precede the
keyword defining the locations and properties of the well connections
using the COMPDAT keyword. The rules on keyword order generally follow
the same sequence as the order of events in a production plan, that is
we drill wells first (WELSPECS), then complete them (COMPDAT), and put
them on production (WCONHIST, WCONPROD, WCONINJH, or WCONINJE).

In most cases a significant part of the contents of the SCHEDULE section
is often imported from other applications. For instance, the Vertical
Lift Performance (``VLP'') tables used to convert from Bottom-Hole
Pressure (``BHP'') to Tubing Head Pressure (``THP'') and entered via the
VFPPROD keyword, are always generated by an external Nodal Analysis
application. These tables are then included in the input deck via the
INCLUDE keyword. Note as the VFP tables are assigned to a well via the
WCONHIST, WCONPROD, WCONINJH, or WCONINJE keywords, then the tables need
to have been previously entered prior to assigning them to a well, that
is the data is order dependent. Well historical production data is also
normally ``included'' from an external source, typically the production
data base, and may have to be edited to ensure that the correct
connections/completions are opened and shut at the appropriate time.

\begin{tcolorbox}[enhanced jigsaw, title=\textcolor{quarto-callout-note-color}{\faInfo}\hspace{0.5em}{Note}, coltitle=black, breakable, toprule=.15mm, colbacktitle=quarto-callout-note-color!10!white, leftrule=.75mm, opacitybacktitle=0.6, colback=white, bottomtitle=1mm, toptitle=1mm, titlerule=0mm, left=2mm, colframe=quarto-callout-note-color-frame, arc=.35mm, rightrule=.15mm, bottomrule=.15mm, opacityback=0]

The term well connection is used to describe individual connections from
the wellbore to a reservoir grid cell, as opposed to a well completion.
A well completion is used to describe a set of connections that are
``grouped'' into well completions using the COMPLUMP keyword in the
SCHEDULE section. For instance, a well may consist of several
completions with each completion consisting of multiple connections.

\end{tcolorbox}

The reservoir management flexibility offered in the SCHEDULE section is
indicated by the more than 300 keywords available in the section and
thus makes quality assurance of the data challenging. This means that
one should not assume that the data included from third party
applications are error free and therefore suitable error checking should
be performed.

The majority of the keywords are related to controlling how wells and
groups are operated, by applying production/injection targets and
constraints, economic criteria and associated actions if a criteria is
violated, drilling of additional wells to meet production targets or
constraints, as well as the relationship between the different wells and
groups in the model. In addition, the DATES, TIME, and TSTEP keywords
are used to progress the simulation through time and the TUNING series
of keywords control the numerical parameters in defining numerical
convergence of each time step. Note that OPM Flow can be instructed to
read the TUNING keyword if the appropriate command line parameter has
been activated (see section 2.2 Running OPM Flow 2023-04 From The
Command Line).

\begin{tcolorbox}[enhanced jigsaw, title=\textcolor{quarto-callout-note-color}{\faInfo}\hspace{0.5em}{Note}, coltitle=black, breakable, toprule=.15mm, colbacktitle=quarto-callout-note-color!10!white, leftrule=.75mm, opacitybacktitle=0.6, colback=white, bottomtitle=1mm, toptitle=1mm, titlerule=0mm, left=2mm, colframe=quarto-callout-note-color-frame, arc=.35mm, rightrule=.15mm, bottomrule=.15mm, opacityback=0]

When the DATES, TIME, and TSTEP keywords are used to advance the
simulation through time, this ``new'' time step is known as a ``report
time step'', as either an action (open a well for example) or a report
request will be made. This is different to a simulator's automatic time
step used to advance the simulator to the next ``reporting time step''.
For instance, if the simulation is currently at January 1, 2020 and the
next reporting time step is January 1, 2021, then the simulator may not
be able to advance directly in one time step to January 1, 2021.
Instead, a series of time steps of various duration will be
automatically selected by the simulator based on the current set of
numerical convergence controls until the required ``report time step''
is reached.

\end{tcolorbox}

The next section summarizes the data requirements for the SCHEDULE
section and offers some guidance on structuring the format of the
keywords in order to increase readability of the input deck and to
minimize potential input errors.

\section{Keyword Definitions}\label{keyword-definitions-17}

\section{Data Requirements}\label{data-requirements-8}

Apart from the keywords that advance the simulator through time (the
DATES, TIME, and TSTEP keywords), the minimum required data is
associated with defining a well (WELSPECS), defining a well's connection
to the reservoir (COMPDAT), and the operating and production targets and
constraints for the well (WCONHIST, WCONPROD, WCONINJH, or WCONINJE).
Well connections can be ``grouped'' into completions using the COMPLUMP
keyword to represent actual physical well completions in the model.
Wells can either operate independently or under group control. That is
when a well is allocated to a group, then the group can dictate how the
wells in the group under group control are operated. For example, a
group may have production targets and constraints and all wells under
group control within the group will be operated in such a manner as to
satisfy the group's targets and constraints. Note that wells can belong
to a group but do not necessary have to be under group control. The top
level group, level one, has the name FIELD and under this level can be
wells, groups, and sub groups to the higher level groups. By default
three group levels are defined that sets the wells as level three,
reporting directly to defined groups at level two, and the level two
groups reporting to the FIELD group at level one. If a different
configuration is required, then the GRUPTREE keyword should be used to
define the group hierarchy by defining a lower level group that reports
directly to a higher level group.

The major well specification keywords are summarized in Table 12.1 for
ease of reference.

\begin{longtable}[]{@{}
  >{\raggedright\arraybackslash}p{(\linewidth - 4\tabcolsep) * \real{0.3333}}
  >{\raggedright\arraybackslash}p{(\linewidth - 4\tabcolsep) * \real{0.3333}}
  >{\raggedright\arraybackslash}p{(\linewidth - 4\tabcolsep) * \real{0.3333}}@{}}
\toprule\noalign{}
\begin{minipage}[b]{\linewidth}\raggedright
Major Well Specification Keywords
\end{minipage} & \begin{minipage}[b]{\linewidth}\raggedright
\end{minipage} & \begin{minipage}[b]{\linewidth}\raggedright
\end{minipage} \\
\midrule\noalign{}
\endhead
\bottomrule\noalign{}
\endlastfoot
Purpose & Keywords & Description \\
Vertical Flow Tubing Performance Table Specification & VFPINJ & VFPINJ
-- Define Injection Vertical Flow Performance Tables. VFPINJ declares
similar data to VFPPROD but for injection wells. \\
VFPPROD & VFPPROD -- Define Production Vertical Flow Performance Tables.
The VFPPROD keyword defines production Vertical Flow Performance
(``VFP'') tables that are used to determine the outflow or downstream
pressure based on the inlet or upstream pressure and the phases flowing
through the system. For a well this means the table relates the flowing
bottom-hole pressure (``FBHP'') to the well's tubing head pressure
(``THP'') based on the oil, gas, and water rates (and any artificial
lift quantities like gas lift gas), or phases ratios, flowing up the
wellbore. & \\
Well Specification ‍ & WELSPECS & WELSPECS -- Define Well Specifications.
The WELSPECS keyword defines the general well specification data for all
well types, and must be used for all wells before any other well
specification keywords are used in the input file. The keyword declares
the name of the well, the initial group the well belongs to, the
wellhead location and other key parameters. \\
Well Completion Data ‍ ‍ & COMPDAT & COMPDAT -- Define Well Connections to
the Grid. COMPDAT defines how a well is connected to the reservoir by
defining or modifying existing well connections. \\
COMPLUMP & COMPLUMP -- Assign Well Connections to Completions. The
COMPLUMP keyword assigns connections, as defined by the COMPDAT keyword,
to completion intervals. This ``lumping'' or ``grouping'' of the
connections to various completion intervals allows automatic workovers
and economic criteria to be applied to the completions (that is a set of
connections) as opposed to the individual connections. & \\
COMPTRAJ & COMPTRAJ -- Define Well Trajectory Connections to the Grid.
COMPTRAJ keyword defines how a well that has been declared as a
trajectory well, using the WELTRAJ keyword, is connected to the
reservoir model by defining or modifying existing well perforation
depths. The keyword can only be used for wells defined by the WELTRAJ
keyword, This is an OPM Flow specific keyword and will therefore cause
an error in the commercial simulator. & \\
WELTRAJ & WELTRAJ -- Define Well Trajectory Data. WELTRAJ defines a
trajectory well together with the well trajectory data, and is used in
conjunction with the COMPTRAJ keyword to define the well connections to
the simulation grid blocks. The keyword can only be used for trajectory
wells that employ the COMPTRAJ keyword to define the connections to the
grid, This is an OPM Flow specific keyword and will therefore cause an
error in the commercial simulator. & \\
Historical Production and Injection Data & WCONHIST & WCONHIST -- Define
Well Historical Production Rates and Pressures. WCONHIST defines
production rates and pressures for wells that have been declared history
matching wells by the use of this keyword. History matching wells are
handled differently than ordinary wells that use the WCONPROD keyword
for controlling their production targets and constraints. \\
WCONINJH & WCONINJH -- Well Historical Observed Injection Rates and
Pressures. WCONINJH declares similar data as the WCONHIST keyword but
for injection wells. & \\
Predictive Production and Injection Data & WCONPROD & WCONPROD -- Define
Well Production Targets and Constraints. The WCONPROD keyword defines
production targets and constraints for wells that have previously been
defined by the WELSPECS keyword. \\
WCONINJE & WCONINJE -- Well Injection Targets and Constraints. WCONINJE
declares similar data as the WCONPROD keyword but for injection wells.
& \\
Well Control & WELCNTL & WELCNTL -- Modify Well Control and Targets. The
WELCNTL keyword modifies a well's target control and value, both rates
and pressures, for previously defined wells without having to define all
the variables on the well control keywords: WCONPROD, WCONHIST,
WCONINJE, or WCONINJH keywords. See also the WELTARG keyword below. \\
WECON & WECON -- Well Economic Criteria for Production Wells. WECON
defines the economic criteria for production wells that have previously
been defined by the WELSPECS and WCONPROD keywords. & \\
WEFAC & WEFAC -- Define Well Efficiency. WEFAC defines a well's
efficiency or up-time factor as opposed to setting the efficiency
factors at the group level. & \\
WHISTCTL & WHISTCTL - Define Well Historical Target Phase. The WHISTCTL
keyword changes the target control phase for wells declared as history
match wells via the WCONHIST keyword. The target phase is set on the
WCONHIST keyword and WHISTCTL overrides this value for all subsequent
entries on the WCONHIST keyword. & \\
WELOPEN & WELOPEN -- Define Well and Well Connections Flowing Status.
WELOPEN defines the status of wells and the well connections, and is
used to open and shut previously defined wells and well connections
without having to re-specify all the data on the well control keywords:
WCONPROD, WCONHIST, WCONINJE, or WCONINJH keywords. & \\
WELTARG & WELTARG -- Modify Well Target and Constraint Values. The
WELTARG keyword modifies the target and constraints values of both rates
and pressures for previously defined wells without having to define all
the variables on the well control keywords: WCONPROD, WCONHIST,
WCONINJE, or WCONINJH keywords. See also the WELCNTL keyword above. & \\
Well Productivity Adjustments (Normally used in History Matching) &
WELPI & WELPI -- Define Well Productivity and Injectivity Indices. WELPI
is used to define a well's productivity or injectivity index and the
values entered on this keyword for a given well will override any
previously calculated values, or values previously entered using this
keyword. \\
WPIMULT & WPIMULT -- Define Well Connection Multipliers. The WPIMULT
keyword defines a well connection multiplier factor that scales the
existing well connection values. The resulting effect is to scale the
well's productivity at the reporting time step the keyword is entered.
& \\
Notes: & & \\
\end{longtable}

\emph{Table 12.1: Major Well Specification Keywords}

Wells are initially allocated to groups via the WELSPECS keyword and
groups have a similar set of keywords as for wells, as outlined in Table
12.2. However only a limited set of keywords have been implemented in
OPM Flow compared with the commercial simulator's set of keywords.

\begin{longtable}[]{@{}
  >{\raggedright\arraybackslash}p{(\linewidth - 4\tabcolsep) * \real{0.3333}}
  >{\raggedright\arraybackslash}p{(\linewidth - 4\tabcolsep) * \real{0.3333}}
  >{\raggedright\arraybackslash}p{(\linewidth - 4\tabcolsep) * \real{0.3333}}@{}}
\toprule\noalign{}
\begin{minipage}[b]{\linewidth}\raggedright
Major Group Specification Keywords
\end{minipage} & \begin{minipage}[b]{\linewidth}\raggedright
\end{minipage} & \begin{minipage}[b]{\linewidth}\raggedright
\end{minipage} \\
\midrule\noalign{}
\endhead
\bottomrule\noalign{}
\endlastfoot
Purpose & Keywords & Description \\
Group Specification & GRUPTREE & GRUPTREE -- Define Group Tree
Hierarchy. GRUPTREE defines the hierarchy of groups that have been
created by having wells assigned to them via the WELSPECS keyword. By
default three group levels are defined that sets the wells as level
three, reporting directly to defined groups at level two, and the level
two groups reporting to the FIELD group at level one. \\
Predictive Production and Injection Data & GCONPROD & GCONPROD -- Group
Production Targets and Constraints. The GCONPROD keyword defines
production targets and constraints for groups, including the top most
group in the group hierarchy known as the FIELD group. \\
GCONINJE & GCONINJE -- Group Injection Targets and Constraints. GCONINJE
declares similar data as for GCONPROD but for group injection targets
and constraints. & \\
Group Control & GECON & GECON -- Group Economic Criteria for Production
Groups. The GECON keyword defines economic criteria for production
groups, including the field level group FIELD, that have previously been
defined by the WELSPECS and GCONPROD keyword. \\
GEFAC & GEFAC -- Define Group Efficiency. GEFAC defines a group's
efficiency or up-time factor as opposed to setting the efficiency
factors for individual wells. & \\
GRUPTARG & GRUPTARG -- Modify Group Targets and Constraints Values.
GRUPTARG keyword modifies the production targets and constraints of both
rates and pressures for previously defined groups without having to
define all the variables on the group production control keywords:
GCONPROD or GCONPRI keywords. & \\
Notes: & & \\
\end{longtable}

\emph{Table 12.2: Major Group Specification Keywords}

The final set of keywords define various controls for the model, the
keywords available to advance the simulator through time, and the
available reporting keywords, as illustrated in Table 12.3.

\begin{longtable}[]{@{}
  >{\raggedright\arraybackslash}p{(\linewidth - 4\tabcolsep) * \real{0.3333}}
  >{\raggedright\arraybackslash}p{(\linewidth - 4\tabcolsep) * \real{0.3333}}
  >{\raggedright\arraybackslash}p{(\linewidth - 4\tabcolsep) * \real{0.3333}}@{}}
\toprule\noalign{}
\begin{minipage}[b]{\linewidth}\raggedright
Schedule Advancement, Control, and Reporting Keywords
\end{minipage} & \begin{minipage}[b]{\linewidth}\raggedright
\end{minipage} & \begin{minipage}[b]{\linewidth}\raggedright
\end{minipage} \\
\midrule\noalign{}
\endhead
\bottomrule\noalign{}
\endlastfoot
Purpose & Keywords & Description \\
Control RS and RV Behavior & DRSDT & DRSDT -- Solution Gas (Rs) Maximum
Rate of Increase Parameters. DRSDT defines the maximum rate at which the
solution gas-oil ratio (Rs) can be increased in a grid cell. The keyword
is similar in functionality to the DRSDTR keyword, that defines the
maximum rate at which Rs can be increased in a grid cell by region. \\
DRSDTR & DRSDTR -- Solution Gas (Rs) Maximum Rate of Increase Parameters
by Region. DRSDTR defines the maximum rate at which the solution gas-oil
ratio (Rs) can be increased in a grid cell for various regions in the
model. The keyword is similar in functionality to the DRSDT keyword,
that defines the maximum rate at which Rs can be increased in a grid
cell for all cells in the model. & \\
DRVDT & DRVDT -- Solution Oil (Rv) Maximum Rate of Increase Parameters.
DRVDT defines the maximum rate at which the solution oil-gas ratio or
condensate-gas ratio (Rv) can be increased in a grid cell. The keyword
is similar in functionality to the DRVDTR keyword, that defines the
maximum rate at which Rv can be increased in a grid cell by region. & \\
DRVDTR & DRVDTR -- Solution Oil (Rv) Maximum Rate of Increase Parameters
by Region. DRVDTR defines the maximum rate at which the solution oil-gas
ratio or condensate-gas ratio (Rv) can be increased in a grid cell for
various regions in the model. The keyword is similar in functionality to
the DRVDT keyword, that defines the maximum rate at which Rv can be
increased in a grid cell for all cells in the model. & \\
Schedule Advancement & DATES & DATES -- Advance Simulation by Reporting
Date. DATES advances the simulation to a given report date after which
additional keywords may be entered to instruct OPM Flow to perform
additional functions via the SCHEDULE section keywords, or further DATES
keywords may be entered to advance the simulator to the next report
date. \\
TIME & TIME -- Advance Simulation by Cumulative Reporting Time. TIME
advances the simulation to a given cumulative report time after which
additional keywords may be entered to instruct OPM Flow to perform
additional functions via the SCHEDULE section keywords, or further TIME
keywords may be entered to advance the simulator to the next report
time. & \\
TSTEP & TSTEP -- Advance Simulation by Reporting Time. TSTEP keyword
advances the simulation to a given report time after which additional
keywords may be entered to instruct OPM Flow to perform additional
functions via the SCHEDULE section keywords, or further TSTEP keywords
may be entered to advance the simulator to the next report time & \\
Schedule Advancement Control & TUNING & TUNING - Numerical Tuning
Control. TUNING defines the parameters used for controlling the
commercial simulator's numerical convergence parameters for the global
grid. The keyword is mostly ignored by OPM Flow; however, the simulator
can be instructed to read the TUNING keyword if the appropriate command
line parameter has been activated (see section 2.2 Running OPM Flow
2023-04 From The Command Line). \\
NEXTSTEP & NEXTSTEP -- Maximum Next Time Step Size . NEXTSTEP defines
the maximum time step size the simulator should take for the next time
step. This keyword can be used to reset the time step for when known
large changes to the model are taking place that may result in time step
chops. & \\
Reporting & RPTSCHED & RPTSCHED -- Define SCHEDULE Section Reporting.
RPTSCHED keyword defines the data in the SCHEDULE section that is to be
printed to the output print file in human readable format. \\
RPTRST & RPTRST -- Define Data to be Written to the RESTART File. RPTRST
keyword defines the data and frequency of the data to be written to the
RESTART file at each requested restart point. In addition to the
solution data arrays required to restart a run, the user may request
additional data to be written to the restart file for visualization in
OPM ResInsight. & \\
Notes: & & \\
\end{longtable}

\emph{Table 12.3: Schedule Advancement, Control, and Reporting Keywords}

In terms of structuring the format of the keywords in the SCHEDULE
section it is advisable to declared all the VLP tables, wells and groups
at the start of SCHEDULE section as oppose to declaring the items as
they are needed, or when they come on stream at various times during the
simulation. This produces a cleaner input deck and tends to limit
unforeseen errors, as all items are declared upfront and only the
operational changes (opening wells, changing group and well targets,
etc.) are needed. The example SCHEDULE section given on the following
pages illustrates a typical SCHEDULE based on this philosophy.

The first segment of the example shows the start of the SCHEDULE section
and the group definitions and controls. Here there are group controls
only at the FIELD level and there are both production and injection rate
targets, as well as water and liquid handling constraints applied on the
FIELD level. These keywords are activated from the start of the
simulation as set by the START keyword in the RUNSPEC section, in this
case January 1, 2020.

\begin{verbatim}
-- ==============================================================================
--
-- SCHEDULE SECTION
--
-- ==============================================================================
SCHEDULE
-- ------------------------------------------------------------------------------
-- GROUP PRODUCTION AND INJECTION CONTROLS
-- ------------------------------------------------------------------------------
--
--       DEFINE GROUP TREE HIERARCHY
--
--       LOWER     HIGHER
--       GROUP     GROUP
GRUPTREE
         FLTBLK1  FIELD                                                        /
         FLTBLK2  FIELD                                                        /
         FLTBLK3  FLTBK2                                                       /
/
--
--       GROUP PRODUCTION CONTROLS
--
-- GRUP  CNTL  OIL    WAT    GAS    LIQ    CNTL  GRUP  GUIDE  GUIDE  CNTL
-- NAME  MODE  RATE   RATE   RATE   RATE   OPT   CNTL  RATE   DEF    WAT
GCONPROD
FIELD    GRAT  1*     8E3   125E3   10E3   1*     1*    1*     1*     1*       /
/

--
--       GROUP INJECTION TARGETS AND CONSTRAINTS
--
-- GRUP  FLUID CNTL   SURF   RESV   REINJ  VOID  GRUP  GUIDE  GUIDE GRUP  GRUP
-- NAME  TYPE  MODE   RATE   RATE   FRAC   FRAC  CNTL  RATE   DEF   REINJ RESV
GCONINJE
FIELD    GAS   RATE  100E3   1*     1*     1.0   YES   1*     1*    1*    1*   /
/
\end{verbatim}

The next segment covers the well specification and includes the loading
of the VFP tables via include files, declaring the wells using the
WELSPECS keyword, and connecting the wells to the reservoir using the
COMPDAT keyword. In addition, the COMPLUMP keyword is invoked to assign
the well connections to well completions.

In multi-stacked reservoirs it is a good idea to associate the same
completion number with a given reservoir for all the wells. So for
example completion number one is always associated with the Lower Talang
Akar Formation Unit A, completion number two with the Lower Talang Akar
Formation Unit B, and completion number three with the Upper Talang Akar
Formation Unit C, etc. In this way one can easily identify which zone a
well is producing from or completed in.

Finally, the WCONPROD keyword is used to define the operating conditions
for the gas producers (GP01 and GP02) and to assign the VFP tables to
the producing wells. Whereas the WCONINJE keyword performs a similar
function for the single gas injector, GI01.

Notice also the well naming convention that easily identifies the type
of well: the letter G for a gas well, O for an oil well, and W for a
water well, and the function of the well, P for a producer and I for an
injector.

\begin{verbatim}
-- ------------------------------------------------------------------------------
-- WELL SPECIFICATIONS AND COMPLETIONS
-- ------------------------------------------------------------------------------
--
--       LOAD INCLUDE FILES WITH VLP TABLES
--
INCLUDE
         'WEL-P50-VLP01.inc'     / VFP TABLE #1 5 1/2 inch tubing

INCLUDE
         'WEL-P50-VLP02.inc'     / VFP TABLE #2 4 1/2 inch tubing
--
--       WELL SPECIFICATION DATA
--
-- WELL  GROUP     LOCATION  BHP    PHASE  DRAIN  INFLOW  OPEN  CROSS  PVT
-- NAME  NAME        I    J  DEPTH  FLUID  AREA   EQUANS  SHUT  FLOW   TABLE
WELSPECS
GI01     FIELD      14   13   1*     GAS    1*     GPP    SHUT   NO     1*     /
GP01     FLTBLK1    64   80   1*     GAS    1*     GPP    SHUT   NO     1*     /
GP02     FLTBLK1    24  110   1*     GAS    1*     GPP    SHUT   NO     1*     /
/
--
--       WELL CONNECTION DATA
--
-- WELL  --- LOCATION ---  OPEN   SAT   CONN   WELL   KH    SKIN   D     DIR
-- NAME   II  JJ  K1  K2   SHUT   TAB   FACT   DIA    FACT  FACT   FACT  PEN
COMPDAT
GI01      1*  1*   1  60   SHUT   1*    1*    0.708   1*    0.0    1*    'Z'   /
GP01      1*  1*   1  60   SHUT   1*    1*    0.708   1*    0.0    1*    'Z'   /
GP01      1*  1*   1  60   SHUT   1*    1*    0.708   1*    0.0    1*    'Z'   /
/
--
--       ASSIGN WELL CONNECTIONS TO COMPLETIONS
--
-- WELL  --- LOCATION ---  COMPL
-- NAME   II  JJ  K1  K2   NO.
COMPLUMP
GI01      14  13   1  20    1                              / COMPLETION NO. 01
GI01      14  13  35  60    2                              / COMPLETION NO. 02
GP01      64  80   1  20    1                              / COMPLETION NO. 01
GP01      64  80  22  60    2                              / COMPLETION NO. 02
GP02      24 110   1  20    1                              / COMPLETION NO. 01
GP02      24 110  25  60    2                              / COMPLETION NO. 02
/
--
--       WELL PRODUCTION WELL CONTROLS
--
-- WELL  OPEN/  CNTL   OIL    WAT    GAS   LIQ    RES    BHP   THP   VFP    VFP
-- NAME  SHUT   MODE   RATE   RATE   RATE  RATE   RATE   PRES  PRES  TABLE  ALFQ
WCONPROD
GP01     SHUT   GRUP   1*     1*     70E3  1*     1*     500.0 120.0 1         /
GP02     SHUT   GRUP   1*     1*     70E3  1*     1*     500.0 120.0 2         /
/
--
--       WELL INJECTION CONTROLS
--
-- WELL  FLUID  OPEN/  CNTL  SURF   RESV   BHP   THP   VFP
-- NAME  TYPE   SHUT   MODE  RATE   RATE   PRSES PRES  TABLE
WCONINJE
GI01     GAS    SHUT   RATE 125E3   1*     10E3   1*    1*                     /
/
\end{verbatim}

The final part of the well specification segment is shown below and sets
the well efficiency for the wells via the WEFAC keyword, the producing
wells economic limits via the WECON keyword, and then advances the
simulation to January 25, 2020. Thus, there is no production up to
January 25, 2020 as all the wells are shut-in.

\begin{verbatim}
--
--       WELL EFFICIENCY FACTORS
--
-- WELL  EFF    NETWK
-- NAME  FACT   OPTN
WEFAC
'GI*'    0.950                                                                 /
'GP*'    0.950                                                                 /
/
--
--       WELL ECONOMIC CRITERIA FOR PRODUCTION WELLS
-- WELL  MIN    MIN    MAX    MAX    MAX    CNTL    END   WELL
-- NAME  ORAT   GRAT   WCUT   GOR    WGR    MODE    RUN   NAME
WECON
'GP*'    1*     5.0E3  1*     1*     1*    'WELL'  'NO'                        /
/
DATES
 25 JAN   2020            /
/
\end{verbatim}

On January 25, 2020 yearly restarts are requested (RPTRST) together with
various printed reports (RPTSCHED), as well as the maximum time step
size of 31 days being set via the TUNING keyword. More importantly, well
GP01 is opened together with completion number one to put the well on
production. Note that one needs to open both the well and the
completion; hence, the two lines in the WELOPEN keyword.

\begin{verbatim}
-- ------------------------------------------------------------------------------
-- SCHEDULE SECTION FOR PHASE 1 DEVELOPMENT
-- ------------------------------------------------------------------------------
--
--       RESTART CONTROL BASIC = 4 (ALL=2, YEARLY=4, MONTHLY=5, TSTEP=6)
--
RPTRST
         BASIC=4                                                               /
--
--       DEFINE SCHEDULE SECTION REPORT OPTION
--
RPTSCHED
         'WELLS=2'    'WELSPECS'    'CPU=2'     'FIP=2'                        /
--
-- DEFAULT TUNING PARAMETERS (SET MAX TSTEP = 31 days)
--
TUNING
1.0      31.0     0.1     0.15     3     0.3     0.1     1.25     0.75         /
                                                                               / /
--
--       DEFINE WELL AND WELL CONNECTIONS FLOWING STATUS
--
--  WELL WELL   --LOCATION--  COMPLETION
--  NAME STAT     I   J    K  FIRST LAST
WELOPEN
GP01     OPEN                                                                  /
GP01     OPEN     0   0    0     1     1                                       /
/

DATES
 1  FEB   2020  /
/
\end{verbatim}

On February 1, 2020 further SCHEDULE keywords are processed. The first
keyword is the RPTSCHED keyword that in this case switches off all
reports printed to the *.PRT file, as only annual reports are required
for this particular model.

Next the gas producer GP02 is put on production using two statements in
the WELOPEN keyword, and the simulator is requested to advance to report
time step of March 1, 2020.

\begin{verbatim}
--
--       DEFINE SCHEDULE SECTION REPORT OPTION
--
RPTSCHED
'NOTHING'                                                                      /
--
--       DEFINE WELL AND WELL CONNECTIONS FLOWING STATUS
--
--  WELL WELL   --LOCATION--  COMPLETION
--  NAME STAT     I   J    K  FIRST LAST
WELOPEN
GP02     OPEN                                                                  /
GP02     OPEN     0   0    0     1     1                                       /
/
DATES
 1  MAR   2020  /
/
\end{verbatim}

Then finally on March 1, 2020 the gas injector is placed on stream and
the simulation advances on a monthly report basis to the end of the
year.

\begin{verbatim}
--
--       DEFINE WELL AND WELL CONNECTIONS FLOWING STATUS
--
--  WELL WELL   --LOCATION--  COMPLETION
--  NAME STAT     I   J    K  FIRST LAST
WELOPEN
GI01     OPEN                                                                  /
GI01     OPEN     0   0    0     1     1                                       /
/
DATES
 1  APR   2020  /
 1  MAY   2020  /
 1  JUN   2020  /
 1  JLY   2020  /
 1  AUG   2020  /
 1  SEP   2020  /
 1  OCT   2020  /
 1  NOV   2020  /
 1  DEC   2020  /
/
\end{verbatim}

Next a detailed report is produced as of January 1, 2021, and reporting
is then switch off and the simulation continues to the end of the run.

\begin{verbatim}
-- ------------------------------------------------------------------------------
-- SCHEDULE SECTION - 2021-01-01
-- ------------------------------------------------------------------------------
RPTSCHED
'WELLS=2'    'WELSPECS'    'CPU=2'     'FIP=2'                                 /

DATES
 1  JAN   2021  /
/

RPTSCHED
'NOTHING'                                                                      /

DATES
 1  FEB   2021  /
 1  MAR   2021  /
 1  APR   2021  /
 ...
/
\end{verbatim}

\subsection{History Matching Runs}\label{history-matching-runs}

History matching is the process of matching actual historical production
data versus the model's production performance. During history matching
the engineer specifies the actual field measured production and
injection rates for a reservoir which has been producing for a period of
time. Although other reservoir characteristics such as permeability,
layering structure, aquifer strength, and individual well performance
are specified as normal in the remainder of the data file, they are only
one realization out of possibly thousands of realizations. There is
therefore significant uncertainty associated with the data set. Thus in
most cases, the simulated past performance of the reservoir, that is the
well rates calculated by OPM Flow up to the present day, will not in
general correspond to the measured rates specified by the user. History
matching then becomes an exercise in identifying those reservoir
properties subject to the greatest uncertainty and adjusting them to
bring the simulated and measured rates to an acceptable degree of
agreement.

Prior to beginning the matching process the production data should be
reviewed and quality checked. This step is needed mainly to:

\begin{itemize}
\tightlist
\item
  Allocate production data areally and vertically to the pertinent areas
  and zones, as well as ensuring that the correct well completions have
  been implemented into the model's schedule section. Particular care
  should be taken with multiple completed wells to ensure that the wells
  produce from the correct zones at the correct times. Wells that are
  completed across fault planes also need to be checked so that the
  completions are on the right side of the fault plane.
\item
  Correct field production data to separator conditions consistent with
  model conditions, or vice versa.
\item
  Correct reported pressures for simulation use (e.g., Peaceman
  correction), or use the weighted areal average well bottom-hole
  pressures from the simulator to match against actual data.
\end{itemize}

Note also that not all production volumes have the same level of quality
assurance. Product streams that are sold are going to have a higher
level of quality assurance because they have value, compared with
production streams that are disposed off or re-injected, as they have no
value. Thus, in an oil field selling the produced oil but flaring or
re-injecting the produced gas, the oil volumes will be back allocated
based on ``ticketed sales'', and thus the field volumes are reasonable
certain. However, even here there are uncertainties associated with
allocating production volumes back to the individual wells and well
completions. As the gas volumes are not sold they have no value and thus
have greater uncertainty than the oil volumes, especially when back
allocating back to the wells.

History matching is an iterative process, in which steps are repeated a
number of times with variations in reservoir characterization
(permeability, pore volume, relative permeability modifications, etc.),
until a reasonable match is obtained. There are no precise rules for
conducting a history match but the methodology is well established.
Generally one first matches pressures globally, then overall saturations
(and well rates), and repeat until an overall match is achieved, as the
process of matching pressures also effects the saturations and vice
versa. Which is why this is an iterative process. This methodology is
then repeated on a regional level until the process is complete and
reasonable matches have been obtained at the field, region, and well
levels.

Generally, group control, targets, and constraints are not utilized in
the history matching part of a simulation study. Instead, the WCONHIST
keyword is used to enter the well production (oil, gas, and water
rates), pressure data (BHP and THP when available), and VFP table
assignment, on a well by well basis and on discrete time intervals,
normally monthly. WCONHIST allows one to set the ``target'' phase and
rate the well will attempt to produce at, if possible, and to produce
the other phases that come with the target rate. That is there are no
constraints applied to the other phases, unlike when using the WCONPROD
keyword. If well productivity has been matched then the target phase
rates will be satisfied and the quality of the history match is
determined by the accuracy of the other phases, together with the
pressure match. Thus, in an oil field the target phase will normally be
the oil phase and the oil rate, and the quality of the history would be
determined by the produced gas and water volumes plus the pressure match
at various levels (field, region, reservoir, and well). If water
breakthrough has not occurred and the field is producing above the
saturation pressure then there is very little to match, as the gas
production in this case is a direct function of the oil rate.

The WCONHIST keyword also has a RESV target phase where the target is
set to the in situ reservoir volume rate which is calculated by the
simulator using the oil, water, gas, and liquid rates declared on the
keyword. This is useful in the initial part of the history matching
study where one is trying to match the overall pressure behavior, as the
option forces the simulator to produce the correct number of reservoir
barrels (or m3), although the actual oil, gas, and water rates will most
likely not match at this stage of the study. Note that it is not
necessary to edit the WCONHIST keywords for a run to change the target
phase, as the WHISTCNTL keyword can be used do this from the time this
keyword is invoked, thus avoiding changing the control mode on all
subsequent WCONHIST keywords.

As mentioned above, the rates and volumes of production streams that are
sold are more reliable than those that are being re-injected, flared, or
otherwise disposed of, and therefore in an oil field the preferred
target phase is oil and not liquid, which some engineers prefer to use.

The WCONINJH keyword has similar functionality as the WCONHIST keyword,
but is used for history matching injection wells instead.

More recent developments in history matching incorporate Assisted
History Matching (``AHM'') techniques {[}History Matching and
Uncertainty Quantification: Multi-objective Particle Swarm Optimisation
Approach (SPE 143067), L. Mohamed, M. Christie, V. Demyanov, Vienna,
Austria, 23--26 May 2011.{]}, {[}Field-Scale Assisted History Matching
Using a Systematic, Massively Parallel Ensemble Kalman Smoother
Procedure (SPE00182617), Binghuai Lin, Paul I Crumpton, and Ali H.
Dogru, Society of Petroleum Engineers (February 2017).{]}, {[}Enhancing
the Geological Models Consistency in Ensemble-Based History Matching -
An Integrated Approach (SPE00186049), A. Perrone, F Pennadoro, A Tiani,
E Della Rossa, and J. Saetrom, Society of Petroleum Engineers (May,
2017).{]}, {[}Correlation-Based Adaptive Localization for Ensemble-Based
History Matching - Applied to the Norne Field Case Study (SPE00191305),
Xiaodong Luo, Rolf Lorentzen, Randi Valestrand, and Geir Evensen, ,
Society of Petroleum Engineers (April, 2018).{]}, {[}Rapid Forecast
Calibration Using Nonlinear Simulation Regression with Localization
(SPE-193845-MS), Jincong He, Wenyue Sun, and Xian-Huan Wen, SPE
Reservoir Simulation Conference held in Galveston, Texas, USA, 10--11
April 2019.{]}, {[}Multilevel Strategies and Geological
Parameterizations for History Matching Complex Reservoir Models
(SPE-193895-MS), Yimin Liu and Louis J. Durlofsky,SPE Reservoir
Simulation Conference held in Galveston, Texas, USA, 10--11 April
2019.{]}and {[}Improved Estimation and Forecast Through Model Error
Estimation -- Norne Field Example (IPTC-19142-MS), Minjie Lu and Yan
Chen, International Petroleum Technology Conference held in Beijing,
China, 26 -- 28 March 2019.{]} in which various statistical techniques
are applied in specialized software to match the field data by the user
supplying the parameters to vary, the range of the parameters values,
and the objective function to minimize. This type of software can also
generate a range of ``reasonable'' history match models that can then be
used to quantify the uncertainty range in the prediction phase.

\subsection{Prediction Runs}\label{prediction-runs}

Prediction runs are used to optimize future production from the field
and can take two forms:

\begin{itemize}
\tightlist
\item
  Prediction runs may follow directly from a history match case using
  the reservoir description estimated during the history matching phase.
  The history matching and prediction runs are frequently in separate
  data files, with the prediction cases ``restarting'' from the final
  history match run, as the history match restart file contains
  sufficient data to continue the simulation into the future. Therefore,
  the prediction data set contains the same reservoir description as the
  history match but is initialized using the restart file output from
  the history match case. In general, the base case scenario is the No
  Further Activity (``NFA'') scenario. This is just a continuation of
  the history match into the future using the well controls in force at
  the end of the history matching period. This is the simplest way to
  run the field since the surface facilities remain unchanged and no
  workovers take place. However, the switch from history matching mode
  to prediction mode should be treated with care in order to ensure a
  smooth transition between the two phases. Subsequent cases will then
  be based on the NFA scenario, and will attempt to optimize recovery by
  examining the impact of drilling infill wells, modifying facilities
  constraints, varying the reservoir management strategy, etc.
\item
  Prediction cases may also be based on predicting production from
  structures that have not been put on production at the time of the
  study, that is there is no historical production available. In this
  case, the reservoir description is based on the best available data
  set and interpretation, although there may be various variations of
  the models to cover the range of uncertainties in the static earth
  modeling. Typically, Low, Best, and High models are generated to
  capture the range of uncertainty in the static modeling, although full
  uncertainty analysis is becoming more common as computational computer
  power becomes more available. Additional dynamic sensitivities are
  then run to determine the Low, Best, and High models that address both
  the static and dynamic uncertainties. These models then form the basis
  of the supplementary predictions using a variety of different
  production strategies that are used to improve the petroleum recovery
  and reduce production costs. For each prediction case, an aspect of
  the surface facilities or well management policy is changed and the
  case is compared to the base prediction. Most simulators, including
  OPM Flow, have a wide range of well, group and surface facility
  modeling features and it is likely that several will be invoked during
  recovery optimization to model events such as drilling infill wells,
  modifying facilities constraints, investigating various depletion
  strategies, etc.
\end{itemize}

In both scenarios more recent developments within the industry apply a
full uncertainty analysis to the models, in which multiple cases are run
under various parameter changes to derive a distribution of recoverable
volumes, which in turn are used to select the Low, Best, and High
scenarios. This can result in hundreds of cases being run and is
computationally expensive.

As mentioned above, the switch from history matching mode to prediction
mode should be treated with care in order to ensure a smooth transition
between the two phases. History matching of a well's BHP has always been
somewhat problematic, as the field measured data consists of shutting-in
the well for a specific period for the BHP to build backup to a well's
average drainage pressure (usually assumed to be the reservoir pressure
in the surrounding area). History matching time steps are normally based
on monthly time steps to coincide with the monthly back allocated field
reported production volumes. Thus, to correctly match the measured
Shut-in Bottom-Hole Pressure (``SBHP'') data in the model, one would
have to shut-in the well in the model as per the field, and then open up
the well back up at the end of the BHP survey. One would also have to
re-adjust the well rates to match the cumulative volumes produced for
the time the well was shut-in. These steps are impracticable and
engineers originally handled the problem by assigning the wells a high
productivity index, so that the difference between the Flowing
Bottom-Hole Pressure (``FBHP'') and the SBHP was small.

This would lead to a mismatch in rates when switching from history
matching mode to prediction mode as shown in Figure 12.1. Thus, the
final step in history matching was to calibrate the well rates based on
adjusting the Productivity Indices (``PI'') of the wells, in order to
get a smooth transition between the two modes of operating the wells.

\pandocbounded{\includegraphics[keepaspectratio]{chapters/images/Frame41_f7479abd5940.png}}
\pandocbounded{\includegraphics[keepaspectratio]{chapters/images/graphics15_f7479abd5940.png}}

Calibration is achieved by running the model at the back pressures held
against the wells at one or several time-steps. The well PIs and/or
wellbore flow parameters are adjusted to duplicate field-observed rates.
Although the calibration step was usually done after the
history-matching step, there are times when it must be done beforehand
(e.g., when modeling multi-layer flow into a well). This is because in
stratified reservoirs the correct allocation of rates among layers is
dependent on the absolute value of the well's PI (and not solely on the
ratios of PIs of the various layers).

Although well PI matching is still performed as part of the general
history matching methodology, it is no longer necessary to set a well's
PI to a high value in order to mimic the field measured SBHP. This is
because most simulators are now capable of reporting an ``average''
bottom-hole pressure around the well, including one, four, five, and
nine grid cell averaging, as well as alternative averaging schemes.
Thus, the engineer now matches the field measured SBHP with one of the
simulator's average grid cell BHPs, normally either the five or nine
grid cell average BHP values.

This leaves the question on what controls should be used when moving
from history matching mode to prediction mode and ensuring a smooth
transition between the two, and at the same time ensuring the resulting
predictive forecast is reasonable. There are basically three approaches
to this problem:

\begin{verbatim}
- set the well’s prediction constraint to the last observed rate,
- set the well’s prediction constraint to the average FBHP and adjusting a well’s PI, and thirdly,
- set the well’s prediction constraint to the last FBHP or THP.
\end{verbatim}

The next three sections discusses these options in some detail.

\begin{enumerate}
\def\labelenumi{(\arabic{enumi})}
\tightlist
\item
  Prediction Constraint Based on a Well's Last Observed Rate
\end{enumerate}

In the history matching mode the wells are controlled by a rate target,
so using the last observed well rate would seem a sensible approach.
Thus, this approach uses the last observed well rates as well
constraints in the prediction run. The typical approach is to constrain
production based on a preferred phase (oil rate, gas rate, or liquid
rate), and then set a reasonable flowing BHP value as a limit for when
the well can no longer meet the rate target set. The main requirement
for using this approach is therefore that the history-matched model
reproduces the rates and flowing bottom-hole pressures fairly exactly.

\pandocbounded{\includegraphics[keepaspectratio]{chapters/images/Frame48_a66cfeb4d167.png}}
\pandocbounded{\includegraphics[keepaspectratio]{chapters/images/Image38_a66cfeb4d167.png}}

However, the are some drawbacks to this approach as one can observe in
Figure 12.2, which shows the oil rate and cumulative oil for a well, the
symbols on the plot are the actual historical oil production rates. If
the last period of production history were not matched successfully, a
smooth transition from the history simulation to prediction phase will
not be achieved, as shown by the orange lines in Figure 12.2.

A second drawback to this approach, is that the well and field
production rate profiles will tend to have a flat trend for a few months
(or even years) when the well can still deliver the rate target. In real
field management such rate profiles can be optimistic if the specified
boundary conditions are not representative of actual operating
conditions in the field.

\begin{enumerate}
\def\labelenumi{(\arabic{enumi})}
\setcounter{enumi}{1}
\tightlist
\item
  Prediction Constraint Based on a Well's Average FBHP and Adjusting a
  Well's PI
\end{enumerate}

This approach has been adopted by some reservoir engineers as a result
of the issues mentioned for the first approach.

Here, an average FBHP of each well, based on most recent production
history, is used as a target to control the wells during the prediction
phase. Usually, this will lead to abrupt changes in rate and pressure
profiles during the transition. The next step is therefore to dampen
these abrupt changes in the well rate profiles by applying well PI
multipliers (at the well and/or completion level) until a smooth
transition is obtained in the rate profiles, similar to that shown by
the green lines after the history matching period in Figure 12.1.

The problem with this approach is that it is synonymous to
re-calibrating a supposedly history-matched simulation model (with a
given permeability distribution based on a certain geological
characterization), and thus this approach is rather questionable and is
not recommended.

\begin{enumerate}
\def\labelenumi{(\arabic{enumi})}
\setcounter{enumi}{2}
\tightlist
\item
  Prediction Constraint Based on a Well's Last Flowing BHP or THP
\end{enumerate}

This methodology is the recommended approach for switching between
historical and prediction controls in a reservoir model. It involves
automatically applying the last flowing bottom-hole pressure at the
instant the transition takes place as the production target, and
therefore setting the latest flow rates to be constraints. This ensures
a smooth transition between history and prediction phases without having
to resort to unreasonable changes to the model. In OPM Flow this is
accomplished by using the WELTARG keyword to specify that all the wells
should use the bottom-hole pressure as the target at the beginning of
the restart run, and defaulting the actual value for the BHP. By
defaulting the BHP value the simulator will use the current value of
FBHP for a well as the constraint. Note that the WELTARG keyword only
defines the variable to be changed, it does not change how a well is
controlled. However, by setting the BHP to equal to the current FBHP
this has a similar effect as changing the operating target.

Using this approach, all wells will continue flowing at the last flowing
bottom-hole pressure from the previous time step, using this value of
FBHP as the effective target while the last historical rates are kept as
limits (if BHP is requested as shown in the example below). This
approach achieves two important objectives, firstly it results in a
smooth transition from the history-match phase to prediction phase in
well and field production profiles And secondly, it gives a smooth
decline in well production rate, as is expected of real wells in the
field.

Although this discussion has been based on using the BHP as the
effective control mechanism, one could also use THP, and to a lesser
extent, liquid rate, instead.

The following example below shows how to apply the recommended method to
obtain a smooth transition to the predictive phase using the BHP as the
controlling mechanism.

\subsubsection{Example}\label{example-503}

In the example, December 2019 is the last month of historical data and
therefore the predication phase starts from January 1, 2020.

\begin{verbatim}
DATES
01 DEC 2019 /
/
--
--       WELL HISTORICAL PRODUCTION CONTROLS
--
-- WELL  OPEN/  CNTL   OIL    WAT    GAS    VFP    VFP   THP   BHP
-- NAME  SHUT   MODE   RATE   RATE   RATE   TABLE  ALFQ  PRES  PRES
WCONHIST                                                                                                                          OP01     OPEN   ORAT  15.5E3  100.0  1550   10      1*   900.0 1*     /
OP02     OPEN   ORAT  10.3E3  500.0  1550   10      1*   700.0 1*     /
OP03     OPEN   ORAT   5.4E3 2100.0  1550   10      1*   500.0 1*     /
/
-- ------------------------------------------------------------------------------
-- 01 JAN 2020 START OF PREDICTION SCHEDULE SECTION
-- ------------------------------------------------------------------------------
DATES
01 JAN 2020 /
/
--
--       WELL PRODUCTION AND INJECTION TARGETS
--
--  WELL WELL   TARGET
--  NAME TARG   VALUE
WELTARG
 OP*     BHP                                                /
/
\end{verbatim}

From December 1, 2019 to January 1, 2020 the three oil producing wells
will be under history matching control with the oil rate being the
target rate for the wells. On January 1, 2020 the wells will be
converted to normal producers as the WELTARG keyword changes a well
status from a history match well to a regular well. At the same time the
WELTARG keyword sets all the wells to have their current FBHP as a
constraint, by setting the target BHP values to the default values,
which effectively makes the FBHP a target as it is now the ``active''
constraint.

The alternative keyword for implementing the same approach is the
WELCNTL keyword, but unfortunately this is currently not implemented in
OPM Flow.

\subsection{Well Productivity}\label{well-productivity}

Throughout this section the equations are presented in terms of oil
production or liquid flow (i.e., oil and water). Note that similar
equations exists for gas flow in terms of p2 and pseudo pressure
formulations (ψ(p)).

Under the assumption of (1) one-dimensional radial flow into the
wellbore, (2) steady-state pressure behavior, and (3) single-phase flow,
it can be shown that the flow rate of a well is given by:

\begin{equation}\phantomsection\label{eq-12-1}{
q = (0.001127) (2π) \frac{\mathit{kh}}{μ}\left[\frac{{p}_{e} - {p}_{w}}{\left(ln \left(\frac{{r}_{e}}{{r}_{w}}\right) + S\right)}\right]
}\end{equation}

where

q = production rate (rb/day),

pe = pressure at the drainage radius of the well (psia),

pw = bottom-hole well pressure (psia),

re = drainage radius (feet),

rw = wellbore radius (feet),

S = skin factor (dimensionless), and

kh = effective permeability thickness (mD.ft).

This is extended to multi-phase flow by introduction of relative
permeabilities and adjusted to surface conditions by use of the
formation volume factor resulting in the following expression for oil
production rate:

\begin{equation}\phantomsection\label{eq-12-2}{
{q}_{o} = (0.001127) (2π) \frac{{\mathit{kk}}_{\mathit{ro}} h}{{B}_{o} {μ}_{o}}\left[\frac{{p}_{e} - {p}_{w}}{\left(ln \left(\frac{{r}_{e}}{{r}_{w}}\right) + S\right)}\right]
}\end{equation}

If the modified radial flow equation is written for each perforated
interval and summed over all intervals, the resulting equation is the
one used to calculate oil production rates in the simulator. One way to
specify the flow capacities of wells is to enter re, rw, S, and pw
values for the well, with the kh values for each perforated interval.
The simulator then determines the other parameter values from the
calculated pressure and saturation distributions.

Because the pressure at the drainage radius is not readily determined in
the simulator, it substitutes pb (a representative pressure of each
perforation grid block) for p.~That is readily accommodated by replacing
the drainage radius by a probe radius as defined by Peaceman
{[}Peaceman, D.W. ``Interpretation of Well-Block Pressures in Numerical
Reservoir Simulation with Nonsquare Grid Blocks and Anisotropic
Permeability.'' SPE 10528 presented at the Sixth SPE Symposium on
Reservoir Simulation, New Orleans, (February 1982), pp.~553-569.{]} for
a vertical well:

\begin{equation}\phantomsection\label{eq-12-3}{
{r}_{b} = 0.28 \frac{{\left[\sqrt{\frac{{K}_{y}}{{K}_{x}}} \mathrm{Δ}{x}^{2} + \sqrt{\frac{{K}_{x}}{{K}_{y}}} \mathrm{Δ}{y}^{2}\right] }^{0.5}}{{\left(\frac{{K}_{y}}{{K}_{x}}\right)}^{0.25} + {\left(\frac{{K}_{x}}{{K}_{y}}\right)}^{0.25}}
}\end{equation}

where

Kx and Ky = x- and y- direction absolute permeabilities, and

Δx and Δy = grid block sizes (feet).

The simulator then uses the modified radial flow equation, written for
each connection:

\begin{equation}\phantomsection\label{eq-12-4}{
{q}_{o} = \frac{(0.001127) (2π)}{\left(ln \left(\frac{{r}_{b}}{{r}_{w}}\right) + S\right)}\left(\frac{k{k}_{\mathit{ro}}h}{{B}_{o} {μ}_{o}} \left( {p}_{b} - {p}_{w} \right)\right)
}\end{equation}

The connection term on the COMPDAT keyword in the SCHEDULE section,
CONFACT, is simply the transmissibility portion of equation (12.4) for a
given connection that is:

\begin{equation}\phantomsection\label{eq-12-5}{
{T}_{\mathit{xy}} = \frac{(0.001127) (2π)(\mathit{kh})}{\left(ln \left(\frac{{r}_{b}}{{r}_{w}}\right) + S\right)}
}\end{equation}

And the oil phase mobility term is defined as:

\begin{equation}\phantomsection\label{eq-12-6}{
{M}_{o} = \left(\frac{{k}_{\mathit{ro}}}{{B}_{o} {μ}_{o}}\right)
}\end{equation}

Substituting (12.5) and (12.6) into (12.4) and summing all the
connections for the well to obtain the oil rate for the well gives:

\begin{equation}\phantomsection\label{eq-12-7}{
{q}_{o} = {\sum }_{1}^{N}\left({T}_{\mathit{xy}} {M}_{o} \left({p}_{b} - {p}_{w}\right)\right)
}\end{equation}

Note in the case of partially completed wells, the effective
permeability thickness of a perforated interval may be substantially
greater than the permeability thickness of the formation along the
perforation.

Engineers commonly use the Productivity Index (``PI'') as a parameter to
defined an oil well's deliverability, as it is derived relatively easy
from field measurements, that is:

\begin{equation}\phantomsection\label{eq-12-8}{
{q}_{o} = PI\left( {p}_{e} - {p}_{w} \right)
}\end{equation}

or in terms of the productivity index:

\begin{equation}\phantomsection\label{eq-12-9}{
PI = \frac{{q}_{o}}{\left( {p}_{e} - {p}_{w} \right)}
}\end{equation}

It is significant that this equation for the Productivity Index is
written in terms of the pressure at the well drainage radius, pe, rather
than grid block pressure, pb. It must be written this way since PI is a
property that is determined by field production tests. One can also
relate the above equation to the multilayered extension of the
steady-state radial flow equation (12.7), to give the grid block
Productivity Index:

\begin{equation}\phantomsection\label{eq-12-10}{
PI = \frac{{q}_{o}}{\left({p}_{b} - {p}_{w}\right)} = {\sum }_{1}^{N}\left({T}_{\mathit{xy}} {M}_{o} \right)
}\end{equation}

One very important point to note is that the productivity index is, in
general, not constant for a given well since both relative permeability
and drainage radius can be time dependent. This makes PI an
inappropriate quantity for computer prediction of production rates, but
it is often the only type of data available on well performance. To
overcome this limitation, most simulators accept well drainage area PI
data as input, and then immediately convert this assuming steady state
radial flow with a uniform mobility throughout the drainage area,
resulting in the following equation.

\begin{equation}\phantomsection\label{eq-12-11}{
PI = {{\sum }_{n=1}^{N}\left({T}_{\mathit{xy}} {M}_{o} \left(\frac{ln \left(\frac{{r}_{b}}{{r}_{w}}\right) + S}{ln \left(\frac{{r}_{e}}{{r}_{w}}\right) + S}\right)\right)}_{n}
}\end{equation}

Since the drainage radius, re, is generally different from the grid
block radius, rb, the ratio of logarithms may be significant; hence, the
use of the PI option normally requires specification of both re and rb.

Equation (12.11) is used to calculate the PI for the well and is used to
print the PI on the WELLS production report requested via the RPTSCHED
keyword in the SCHEDULE section. Note that for gas wells with non-zero
D-factors, the non-Darcy skin factor is added to to the skin (S) term in
equation (12.11) in both the numerator and the denominator. Secondly, re
is often undetermined or unknown and it is common to default the DRADIUS
parameter on the WELSPECS keyword, in the SCHEDULE section. If the well
drainage radius (DRADIUS) is defaulted then equation (12.11) simplifies
to:

\begin{equation}\phantomsection\label{eq-12-12}{
PI = {{\sum }_{n=1}^{N}\left({T}_{\mathit{xy}} {M}_{o} \right)}_{n}
}\end{equation}

In this case the PI should be considered a grid block productivity index
and not a well drainage area productivity index. Note also that if
DRADIUS is set to a negative number, then this results in a well's
potential being written out to the WELLS production report instead of
the productivity index. In this instance this is the rate based on only
applying the BHP and THP constraints only, all other constraints are
ignored.

\section{Keyword Definitions}\label{keyword-definitions-18}

\subsection{ACTION -- Define Action Conditions and Command Processing
(Field)}\label{action-define-action-conditions-and-command-processing-field}

\begin{longtable}[]{@{}
  >{\raggedright\arraybackslash}p{(\linewidth - 14\tabcolsep) * \real{0.1250}}
  >{\raggedright\arraybackslash}p{(\linewidth - 14\tabcolsep) * \real{0.1250}}
  >{\raggedright\arraybackslash}p{(\linewidth - 14\tabcolsep) * \real{0.1250}}
  >{\raggedright\arraybackslash}p{(\linewidth - 14\tabcolsep) * \real{0.1250}}
  >{\raggedright\arraybackslash}p{(\linewidth - 14\tabcolsep) * \real{0.1250}}
  >{\raggedright\arraybackslash}p{(\linewidth - 14\tabcolsep) * \real{0.1250}}
  >{\raggedright\arraybackslash}p{(\linewidth - 14\tabcolsep) * \real{0.1250}}
  >{\raggedright\arraybackslash}p{(\linewidth - 14\tabcolsep) * \real{0.1250}}@{}}
\toprule\noalign{}
\begin{minipage}[b]{\linewidth}\raggedright
RUNSPEC
\end{minipage} & \begin{minipage}[b]{\linewidth}\raggedright
GRID
\end{minipage} & \begin{minipage}[b]{\linewidth}\raggedright
EDIT
\end{minipage} & \begin{minipage}[b]{\linewidth}\raggedright
PROPS
\end{minipage} & \begin{minipage}[b]{\linewidth}\raggedright
REGIONS
\end{minipage} & \begin{minipage}[b]{\linewidth}\raggedright
SOLUTION
\end{minipage} & \begin{minipage}[b]{\linewidth}\raggedright
SUMMARY
\end{minipage} & \begin{minipage}[b]{\linewidth}\raggedright
SCHEDULE
\end{minipage} \\
\midrule\noalign{}
\endhead
\bottomrule\noalign{}
\endlastfoot
\end{longtable}

\subsubsection{Description}\label{description-864}

The ACTION keyword defines a series of conditions that invoke run time
processing of ACTION functions and is similar to executing a run time
script for conditions and variables at the field level.

This keyword is not supported by OPM Flow but would change the results
if supported so the simulation will be stopped; use the ACTIONX keyword
instead.

\subsection{ACTIONG -- Define Action Conditions and Command Processing
(Groups)}\label{actiong-define-action-conditions-and-command-processing-groups}

\begin{longtable}[]{@{}
  >{\raggedright\arraybackslash}p{(\linewidth - 14\tabcolsep) * \real{0.1250}}
  >{\raggedright\arraybackslash}p{(\linewidth - 14\tabcolsep) * \real{0.1250}}
  >{\raggedright\arraybackslash}p{(\linewidth - 14\tabcolsep) * \real{0.1250}}
  >{\raggedright\arraybackslash}p{(\linewidth - 14\tabcolsep) * \real{0.1250}}
  >{\raggedright\arraybackslash}p{(\linewidth - 14\tabcolsep) * \real{0.1250}}
  >{\raggedright\arraybackslash}p{(\linewidth - 14\tabcolsep) * \real{0.1250}}
  >{\raggedright\arraybackslash}p{(\linewidth - 14\tabcolsep) * \real{0.1250}}
  >{\raggedright\arraybackslash}p{(\linewidth - 14\tabcolsep) * \real{0.1250}}@{}}
\toprule\noalign{}
\begin{minipage}[b]{\linewidth}\raggedright
RUNSPEC
\end{minipage} & \begin{minipage}[b]{\linewidth}\raggedright
GRID
\end{minipage} & \begin{minipage}[b]{\linewidth}\raggedright
EDIT
\end{minipage} & \begin{minipage}[b]{\linewidth}\raggedright
PROPS
\end{minipage} & \begin{minipage}[b]{\linewidth}\raggedright
REGIONS
\end{minipage} & \begin{minipage}[b]{\linewidth}\raggedright
SOLUTION
\end{minipage} & \begin{minipage}[b]{\linewidth}\raggedright
SUMMARY
\end{minipage} & \begin{minipage}[b]{\linewidth}\raggedright
SCHEDULE
\end{minipage} \\
\midrule\noalign{}
\endhead
\bottomrule\noalign{}
\endlastfoot
\end{longtable}

\subsubsection{Description}\label{description-865}

The ACTIONG keyword defines a series of conditions that invoke run time
processing of ACTION functions and is similar to executing a run time
script for conditions and variables at the group level

This keyword is not supported by OPM Flow but would change the results
if supported so the simulation will be stopped; use the ACTIONX keyword
instead.

\subsection{ACTIONR -- Define Action Conditions and Command Processing
(Regions)}\label{actionr-define-action-conditions-and-command-processing-regions}

\begin{longtable}[]{@{}
  >{\raggedright\arraybackslash}p{(\linewidth - 14\tabcolsep) * \real{0.1250}}
  >{\raggedright\arraybackslash}p{(\linewidth - 14\tabcolsep) * \real{0.1250}}
  >{\raggedright\arraybackslash}p{(\linewidth - 14\tabcolsep) * \real{0.1250}}
  >{\raggedright\arraybackslash}p{(\linewidth - 14\tabcolsep) * \real{0.1250}}
  >{\raggedright\arraybackslash}p{(\linewidth - 14\tabcolsep) * \real{0.1250}}
  >{\raggedright\arraybackslash}p{(\linewidth - 14\tabcolsep) * \real{0.1250}}
  >{\raggedright\arraybackslash}p{(\linewidth - 14\tabcolsep) * \real{0.1250}}
  >{\raggedright\arraybackslash}p{(\linewidth - 14\tabcolsep) * \real{0.1250}}@{}}
\toprule\noalign{}
\begin{minipage}[b]{\linewidth}\raggedright
RUNSPEC
\end{minipage} & \begin{minipage}[b]{\linewidth}\raggedright
GRID
\end{minipage} & \begin{minipage}[b]{\linewidth}\raggedright
EDIT
\end{minipage} & \begin{minipage}[b]{\linewidth}\raggedright
PROPS
\end{minipage} & \begin{minipage}[b]{\linewidth}\raggedright
REGIONS
\end{minipage} & \begin{minipage}[b]{\linewidth}\raggedright
SOLUTION
\end{minipage} & \begin{minipage}[b]{\linewidth}\raggedright
SUMMARY
\end{minipage} & \begin{minipage}[b]{\linewidth}\raggedright
SCHEDULE
\end{minipage} \\
\midrule\noalign{}
\endhead
\bottomrule\noalign{}
\endlastfoot
\end{longtable}

\subsubsection{Description}\label{description-866}

The ACTIONR keyword defines a series of conditions that invoke run time
processing of ACTION functions and is similar to executing a run time
script for conditions and variables at the region level.

This keyword is not supported by OPM Flow but would change the results
if supported so the simulation will be stopped; use the ACTIONX keyword
instead.

\subsection{ACTIONS -- Define Action Conditions and Command Processing
(Well
Segments)}\label{actions-define-action-conditions-and-command-processing-well-segments}

\begin{longtable}[]{@{}
  >{\raggedright\arraybackslash}p{(\linewidth - 14\tabcolsep) * \real{0.1250}}
  >{\raggedright\arraybackslash}p{(\linewidth - 14\tabcolsep) * \real{0.1250}}
  >{\raggedright\arraybackslash}p{(\linewidth - 14\tabcolsep) * \real{0.1250}}
  >{\raggedright\arraybackslash}p{(\linewidth - 14\tabcolsep) * \real{0.1250}}
  >{\raggedright\arraybackslash}p{(\linewidth - 14\tabcolsep) * \real{0.1250}}
  >{\raggedright\arraybackslash}p{(\linewidth - 14\tabcolsep) * \real{0.1250}}
  >{\raggedright\arraybackslash}p{(\linewidth - 14\tabcolsep) * \real{0.1250}}
  >{\raggedright\arraybackslash}p{(\linewidth - 14\tabcolsep) * \real{0.1250}}@{}}
\toprule\noalign{}
\begin{minipage}[b]{\linewidth}\raggedright
RUNSPEC
\end{minipage} & \begin{minipage}[b]{\linewidth}\raggedright
GRID
\end{minipage} & \begin{minipage}[b]{\linewidth}\raggedright
EDIT
\end{minipage} & \begin{minipage}[b]{\linewidth}\raggedright
PROPS
\end{minipage} & \begin{minipage}[b]{\linewidth}\raggedright
REGIONS
\end{minipage} & \begin{minipage}[b]{\linewidth}\raggedright
SOLUTION
\end{minipage} & \begin{minipage}[b]{\linewidth}\raggedright
SUMMARY
\end{minipage} & \begin{minipage}[b]{\linewidth}\raggedright
SCHEDULE
\end{minipage} \\
\midrule\noalign{}
\endhead
\bottomrule\noalign{}
\endlastfoot
\end{longtable}

\subsubsection{Description}\label{description-867}

The ACTIONS keyword defines a series of conditions that invoke run time
processing of ACTION functions and is similar to executing a run time
script for conditions and variables associated with well segments.

This keyword is not supported by OPM Flow but would change the results
if supported so the simulation will be stopped; use the ACTIONX keyword
instead.

\subsection{ACTIONW -- Define Action Conditions and Command Processing
(Wells)}\label{actionw-define-action-conditions-and-command-processing-wells}

\begin{longtable}[]{@{}
  >{\raggedright\arraybackslash}p{(\linewidth - 14\tabcolsep) * \real{0.1250}}
  >{\raggedright\arraybackslash}p{(\linewidth - 14\tabcolsep) * \real{0.1250}}
  >{\raggedright\arraybackslash}p{(\linewidth - 14\tabcolsep) * \real{0.1250}}
  >{\raggedright\arraybackslash}p{(\linewidth - 14\tabcolsep) * \real{0.1250}}
  >{\raggedright\arraybackslash}p{(\linewidth - 14\tabcolsep) * \real{0.1250}}
  >{\raggedright\arraybackslash}p{(\linewidth - 14\tabcolsep) * \real{0.1250}}
  >{\raggedright\arraybackslash}p{(\linewidth - 14\tabcolsep) * \real{0.1250}}
  >{\raggedright\arraybackslash}p{(\linewidth - 14\tabcolsep) * \real{0.1250}}@{}}
\toprule\noalign{}
\begin{minipage}[b]{\linewidth}\raggedright
RUNSPEC
\end{minipage} & \begin{minipage}[b]{\linewidth}\raggedright
GRID
\end{minipage} & \begin{minipage}[b]{\linewidth}\raggedright
EDIT
\end{minipage} & \begin{minipage}[b]{\linewidth}\raggedright
PROPS
\end{minipage} & \begin{minipage}[b]{\linewidth}\raggedright
REGIONS
\end{minipage} & \begin{minipage}[b]{\linewidth}\raggedright
SOLUTION
\end{minipage} & \begin{minipage}[b]{\linewidth}\raggedright
SUMMARY
\end{minipage} & \begin{minipage}[b]{\linewidth}\raggedright
SCHEDULE
\end{minipage} \\
\midrule\noalign{}
\endhead
\bottomrule\noalign{}
\endlastfoot
\end{longtable}

\subsubsection{Description}\label{description-868}

This keyword starts the definition of an ACTIONW section that stipulates
the Boolean conditions to test the nominated well parameters, and the
resulting SCHEDULE keywords to be executed, if the Boolean condition
evaluates to true. An ACTIONW Definition Section is terminated by an
ENDACTIO keyword on a separate line. Here, the keyword defines a series
of conditions applied to wells only, that invoke run time processing of
ACTION functions, and is similar to executing a run time script for
conditions and variables at the well level. The ACTION series of
keywords (ACTION, ACTIONG, ACTIONR, ACTIONS and ACTIONW) can apply
Boolean conditional tests to variables at the field, group, region, well
segment and well levels.

Note that one should use caution when using this facility as it may
result in the simulator aborting, because the ACTIONW keyword enables
the user to implement complex functionality and therefore it is
advisable to start with simple expressions before adding the desired
complexity.

See also the PYACTION keyword in the SCHEDULE that implements OPM Flow's
Python scripting facility using the Python scripting language.

This keyword is not supported by OPM Flow but would change the results
if supported so the simulation will be stopped; use the ACTIONX keyword
instead.

\begin{longtable}[]{@{}
  >{\raggedright\arraybackslash}p{(\linewidth - 6\tabcolsep) * \real{0.2500}}
  >{\raggedright\arraybackslash}p{(\linewidth - 6\tabcolsep) * \real{0.2500}}
  >{\raggedright\arraybackslash}p{(\linewidth - 6\tabcolsep) * \real{0.2500}}
  >{\raggedright\arraybackslash}p{(\linewidth - 6\tabcolsep) * \real{0.2500}}@{}}
\toprule\noalign{}
\begin{minipage}[b]{\linewidth}\raggedright
No.
\end{minipage} & \begin{minipage}[b]{\linewidth}\raggedright
Name
\end{minipage} & \begin{minipage}[b]{\linewidth}\raggedright
Description
\end{minipage} & \begin{minipage}[b]{\linewidth}\raggedright
Default
\end{minipage} \\
\midrule\noalign{}
\endhead
\bottomrule\noalign{}
\endlastfoot
ACTIONW & Defines the start of an ACTIONW Definition Section. This is
followed on a new line by the ACTIONW record that stipulates the
conditions for which the defined action will be executed, and then this
is followed by the various operations to be performed if the condition
is satisfied. & & \\
1-1 & ACTNAME & A character sting of up to eight characters in length
that defines the name of this action definition. If ACTNAME has
previously been used by any ACTION series keyword, then the previous
ACTION definitions will be replaced by the definition declared by this
ACTIONW Definition Section. & None \\
1-2 & ACTWELL & A character string of up to eight characters in length
that defines the well name for which the ACTIONW Definition Section is
being defined. Note that the well name (ACTWELL) must have been declared
previously using the WELSPECS keyword in the SCHEDULE section, otherwise
an error may occur. & None \\
1-3 & ACTLHS & A character string of up to eight characters in length
that stipulates a well SUMMARY variable on the left hand side of a
Boolean conditional test. For example, WOPR that stands for the Well Oil
Production Rate. In addition, ACTLHS can also be a well User Defined
Quantity (``UDQ'') defined by the UDQ - Declare User Define Quantities
(``UDQ'') keyword in the SCHEDULE section. In this case, the first two
characters of ACTLHS must be set to WU and the UDQ variable must have
previously been declared with the UDQ keyword. A complete list of well
SUMMARY vectors that can be used with ACTIONW keyword by OPM Flow and
the commercial black-oil simulator is summarized in Table 12.5. & Not
Applicable \\
1-4 & ACTTEST & A defined character string that the defines the Boolean
operator and must be set to one of the following Boolean conditionals:
Note that the OPM Flow implementation of ACTTEST has been enhanced to
support all of the above conditions. Whereas, the commercial simulator
only supports (1) and (2). Thus, one should be mindful of this fact if
the input deck is to be run by both OPM Flow and the commercial
simulator. & Not Applicable \\
1-5 & ACTRHS & A numeric value that defines the constant on the right
hand side of the Boolean conditional test. The set of wells for which
the Boolean condition is true, is captured and can be used as a general
``well list'' with the symbol `?' in subsequent well keywords. For
example, to shut wells with a gas production rate less than 5 MMscf/d,
one would use: ACTIONW ACT01 'GP*' WGPR \textless{} 5000.0 10000 /
WELOPEN `?' SHUT / / \ldots{} ENDACTIO & Not Applicable \\
1-6 & ACTNSTEP & A positive integer that defines the maximum number
times that the ACTNAME action is to be executed. The ACTIONW action is
excuted at the end of each time step if at least one well satisfies the
Boolean condition until the action has been executed the specified
number of times. The default value of one means that the definition will
be executed only once. One can use a large value, for example 10,000 for
the definition to be executed at every time step. Note that the counter
only affects successful evaluations; i.e.~if ACTNSTEP is set equal to
one (the default), then the simulator will test the action at the end of
every time step until it evaluates to true. & 1 \\
1-3 & ACTINCR & ACTINCR is a real negative or positive value that
specifies a value to increment ACTRHS every time the Boolean condition
evaluates to true. For example, if ACTINCR is set to 500.0: ACTIONW
ACT02 `GP\emph{' WGPR \textless{} 5000.0 10000 500.0 / WELOPEN '?' SHUT
/ / \ldots{} ENDACTIO Then after the third time the Boolean conditional
has evaluated to true, the condition would effectively be: ACTIONW ACT02
'GP}' WGPR \textless{} 6500.0 10000 500.0 / \ldots{} ENDACTIO & 0.0 \\
1-4 & / & Record terminated by a ``/'' & Not Applicable \\
& & The next section contains any number of standard SCHEDULE keywords
that will be executed if the Boolean expression evaluates to true. For
example, to reduce the tubing head pressure constraint when any of the
oil producers' oil rate drop below 100 stb/d then one could use: ACTIONW
ACT03 `OP\emph{' WOPR \textless{} 100.0 10000 / -- -- FLOW WELLS THROUGH
LOW PRESSURE SEPARATOR -- -- WELL WELL TARGET -- NAME TARG VALUE WELTARG
'OP}' THP 150 / / ENDACTIO In theory, most SCHEDULE keywords can be used
in an ACTIONW Definition Section here, except for the time stepping
keywords, i.e, TSTEP and DATES. See the ACTIONX keyword for a list of
the SCHEDULE keywords that are known to work with the ACTION series of
keywords. & Not Applicable \\
ENDACTIO & Define the end of ACTIONW Definition Section. & Not
Applicable & \\
Notes: & & & \\
\end{longtable}

\emph{Table 12.4: ACTIONW Keyword Description}

The well SUMMARY vectors that can be used with the ACTIONW keyword by
OPM Flow and the commercial black-oil simulator are outlined in Table
12.5.

\begin{longtable}[]{@{}
  >{\raggedright\arraybackslash}p{(\linewidth - 4\tabcolsep) * \real{0.3333}}
  >{\raggedright\arraybackslash}p{(\linewidth - 4\tabcolsep) * \real{0.3333}}
  >{\raggedright\arraybackslash}p{(\linewidth - 4\tabcolsep) * \real{0.3333}}@{}}
\toprule\noalign{}
\begin{minipage}[b]{\linewidth}\raggedright
No.
\end{minipage} & \begin{minipage}[b]{\linewidth}\raggedright
Well SUMMARY Vector Description
\end{minipage} & \begin{minipage}[b]{\linewidth}\raggedright
SUMMARY Mnemonic
\end{minipage} \\
\midrule\noalign{}
\endhead
\bottomrule\noalign{}
\endlastfoot
1 & Bottom-Hole Pressure & WBHP \\
2 & Gas Injection Rate & WGIR \\
3 & Gas Injection Total & WGIT \\
4 & Gas-Liquid Ratio & WGLR \\
5 & Gas-Liquid Ratio (Bottom Hole) & WBGLR \\
6 & Gas-Oil Ratio & WGOR \\
7 & Gas Production Rate & WGPR \\
8 & Gas Production Total & WGPT \\
9 & Liquid Production Rate & WLPR \\
10 & Liquid Production Rate & WLPT \\
11 & Oil Injection Rate & WOIR \\
12 & Oil Injection Total & WOIT \\
13 & Oil Production Rate & WOPR \\
14 & Oil Production Total & WOPT \\
15 & Polymer Production Concentration & WCPC \\
16 & Polymer Production Rate & WCPR \\
17 & Salt Production Concentration & WSPC \\
18 & Salt Production Rate & WSPR \\
19 & Tracer Production Concentration (Tracer name should be added to
WTPC) & WTPC \\
20 & Tracer Production Rate (Tracer name should be added to WTPC) &
WTPR \\
21 & Tubing Head Pressure & WTHP \\
22 & Voidage Injection Rate & WVIR \\
23 & Voidage Injection Total & WVIT \\
24 & Voidage Production Rate & WVPR \\
25 & Voidage Production Total & WVPT \\
26 & Water Cut & WWCT \\
27 & Water Injection Rate & WWIR \\
28 & Water Injection Total & WWIT \\
29 & Water Production Rate & WWPR \\
30 & Water Production Total & WWPT \\
31 & Water-Gas Ratio & WWGR \\
32 & User Defined Quantity & WUXXXXXX \\
Notes: & & \\
\end{longtable}

\emph{Table 12.5: ACTIONW Supported Well SUMMARY Variables}

See also the ACTDIMS and UDADIMS keywords in the RUNSPEC section to
define the dimensions for the ACTION series of keywords and associated
variables. In addition, the EXIT keyword in the SCHEDULE section that
allows for terminating the simulation for when a condition within an
ACTIONW definition is satisfied

Although most SCHEDULE section keywords should work with the ACTIONW
keyword, Table 12.8 under the description of ACTIONX -- Define Action
Conditions and Command Processing shows the status of keywords that have
been tested and known to work, together with keywords that are currently
planned to be implemented.

As mentioned previously, the UDQ keyword stipulates the variables and
operations used to access the User Defined Quantities features in OPM
Flow. UDQ variables can be constants, SUMMARY variables, as defined in
the SUMMARY section, or a formula using various mathematical functions
together with constants and SUMMARY variables.

\begin{tcolorbox}[enhanced jigsaw, title=\textcolor{quarto-callout-note-color}{\faInfo}\hspace{0.5em}{Note}, coltitle=black, breakable, toprule=.15mm, colbacktitle=quarto-callout-note-color!10!white, leftrule=.75mm, opacitybacktitle=0.6, colback=white, bottomtitle=1mm, toptitle=1mm, titlerule=0mm, left=2mm, colframe=quarto-callout-note-color-frame, arc=.35mm, rightrule=.15mm, bottomrule=.15mm, opacityback=0]

Within an ACTIONW Definition Section any UDQ variables utilizing well
variables, must have their associated wells previously fully defined in
the commercial simulator, otherwise an error will occur. For example, if
a well's GOR is being used as part of a UDQ definition, then the well
must be fully characterized prior to declaring the UDQ definition. This
restriction does not apply to OPM Flow; however, it should be considered
if the same deck is to be run with both simulators.

\end{tcolorbox}

User Defined Quantities can also be used as User Defined Arguments
(``UDA'') in the SCHEDULE section with various group, well, and
connection keywords. In this case, the UDA variables are used to replace
numerical values on these keywords by UDA variables that have been
defined by the UDQ keyword. For example, if we wish to make the oil rate
for certain wells to be a function of their water cut, then one can
define the function using the UDQ keyword that results in a UDQ
variable, WU\_WOPR say, and then use WU\_WOPR as a UDA variable on the
WCONPROD keyword for the ORAT parameter. See Table 12.76 for a list of
keywords that can be used with UDA variables in the UDQ - Declare User
Define Quantities (``UDQ'') keyword section.

\subsubsection{Examples}\label{examples-67}

The first example uses the ACTIONW keyword to re-complete a gas
injection well, GI01, when the well's bottom-hole pressure exceeds the
fracture pressure of the formation, and to open up a structurally higher
zone in the well. The keyword NEXTSTEP is used to set the next step to
0.1 days to avoid convergence issues due to a well event.

\begin{verbatim}
--
--       ACTIONW WELL COMMANDS
--
ACTIONW
--       ACTION     WELL     ACTION   ACTION  ACTION  MAX      INCREMENT
--       NAME       NAME     LHS      TEST    RHS     ACTIONS  RHS
         ACTW-1     GI01     WBHP      >      6800.0                            /
--
--       ACTION COMMANDS TO BE EXECUTED
--
--       NEXT   ALL
--       STEP   TIME
NEXTSTEP
         0.1    'NO'                                                           /
--
-- WELL PRODUCTION STATUS
--
--  WELL WELL   --LOCATION--  COMPLETION
--  NAME STAT     I   J    K  FIRST LAST
WELOPEN
GI01     SHUT                                                                  /
GI01     SHUT     0   0    0     2     2                                       /
GI01     OPEN                                                                  /
GI01     OPEN     0   0    0     1     1                                       /
/
ENDACTIO
\end{verbatim}

The second example shows how to reduce a producing well's deliverability
due to water production. Here, well OP01's productivity index is reduced
by 2\% when the well's liquid production total is greater than 100,000
stb. The reduction is repeated 300 times and each time the action is
executed the ACTRHS constant is increased by 10,000 stb.

\begin{verbatim}
--
--       ACTIONW WELL COMMANDS
--
ACTIONW
--       ACTION     WELL     ACTION   ACTION  ACTION  MAX      ICREMENT
--       NAME       NAME     LHS      TEST    RHS     ACTIONS  RHS
         ACTW-01    OP01     WLPT       >     100E3   300      10E3     /
--
--       DEFINE WELL CONNECTION MULTIPLIERS
--
-- WELL  PI     --LOCATION--  COMPLETION
-- NAME  MULT     I   J    K  FIRST LAST
WPIMULT
OP01     0.980    1*  1*  1*    1*    1*                                /
/

ENDACTIO
\end{verbatim}

The final example is similar to the previous example, and shows how to
nest ACTIONW blocks and how to apply ACTIONW to all the oil producers
when the Boolean condition are satisfied.

\begin{verbatim}
--
--       ACTIONW WELL COMMANDS
--
ACTIONW
--       ACTION     WELL     ACTION   ACTION  ACTION  MAX      ICREMENT
--       NAME       NAME     LHS      TEST    RHS     ACTIONS  RHS
         ACTW-01A   'OP*'    WLPT       >     100E3   300      10E3     /

--
--       ACTIONW WELL COMMANDS
--
ACTIONW
--       ACTION     WELL     ACTION   ACTION  ACTION  MAX      ICREMENT
--       NAME       NAME     LHS      TEST    RHS     ACTIONS  RHS
         ACTW-01B   'OP*'    WPI       >      5.0     1        0.0      /

--
--       DEFINE WELL CONNECTION MULTIPLIERS
--
-- WELL  PI     --LOCATION--  COMPLETION
-- NAME  MULT     I   J    K  FIRST LAST
WPIMULT
'?'      0.980    1*  1*  1*    1*    1*                                /
/

ENDACTIO
ENDACTIO
\end{verbatim}

In this case the outer ACTIONW checks if the liquid production for each
OP* well is great than the calculated ACTRHS constant, if it is, then
inner ACTIONW reduces all the selected wells' productivity indices by
2\%, provided the well's productivity index is greater than five.

\begin{verbatim}
\end{verbatim}

\subsection{ACTIONX -- Define Action Conditions and Command
Processing}\label{actionx-define-action-conditions-and-command-processing}

\begin{longtable}[]{@{}
  >{\raggedright\arraybackslash}p{(\linewidth - 14\tabcolsep) * \real{0.1250}}
  >{\raggedright\arraybackslash}p{(\linewidth - 14\tabcolsep) * \real{0.1250}}
  >{\raggedright\arraybackslash}p{(\linewidth - 14\tabcolsep) * \real{0.1250}}
  >{\raggedright\arraybackslash}p{(\linewidth - 14\tabcolsep) * \real{0.1250}}
  >{\raggedright\arraybackslash}p{(\linewidth - 14\tabcolsep) * \real{0.1250}}
  >{\raggedright\arraybackslash}p{(\linewidth - 14\tabcolsep) * \real{0.1250}}
  >{\raggedright\arraybackslash}p{(\linewidth - 14\tabcolsep) * \real{0.1250}}
  >{\raggedright\arraybackslash}p{(\linewidth - 14\tabcolsep) * \real{0.1250}}@{}}
\toprule\noalign{}
\begin{minipage}[b]{\linewidth}\raggedright
RUNSPEC
\end{minipage} & \begin{minipage}[b]{\linewidth}\raggedright
GRID
\end{minipage} & \begin{minipage}[b]{\linewidth}\raggedright
EDIT
\end{minipage} & \begin{minipage}[b]{\linewidth}\raggedright
PROPS
\end{minipage} & \begin{minipage}[b]{\linewidth}\raggedright
REGIONS
\end{minipage} & \begin{minipage}[b]{\linewidth}\raggedright
SOLUTION
\end{minipage} & \begin{minipage}[b]{\linewidth}\raggedright
SUMMARY
\end{minipage} & \begin{minipage}[b]{\linewidth}\raggedright
SCHEDULE
\end{minipage} \\
\midrule\noalign{}
\endhead
\bottomrule\noalign{}
\endlastfoot
\end{longtable}

\subsubsection{Description}\label{description-869}

The ACTIONX keyword defines a series of conditions that invoke run time
processing of ACTION functions and is similar to executing a run time
script. This is the general purpose version of the ACTION series of
keywords that can apply Boolean conditional tests to variables at the
field, group, region, well segment and well levels. The ACTION, ACTIONG,
ACTIONR, ACTIONS and ACTIONW keywords are not implemented in OPM Flow
and are unlikely to be so, as the ACTIONX keyword implements their
functionality with greater flexibility.

This keyword starts the definition of an ACTIONX section that stipulates
the Boolean conditions to test and the resulting SCHEDULE keywords to be
executed if the Boolean condition evaluates to true. An ACTIONX
Definition Section is terminated by an ENDACTIO keyword on a separate
single line.

Although this keyword is read by OPM Flow and the ACTION and UDQ
computational logic and calculations have been implemented, one should
use caution when using this facility as it may result in OPM Flow
aborting. This is because the ACTIONX keyword enables the user to
implement complex functionality and therefore it is advisable to start
with simple expressions before adding the desired complexity.

See also the PYACTION keyword in the SCHEDULE section that implements
OPM Flow's Python scripting facility using the standard Python scripting
language.

\begin{longtable}[]{@{}
  >{\raggedright\arraybackslash}p{(\linewidth - 6\tabcolsep) * \real{0.2500}}
  >{\raggedright\arraybackslash}p{(\linewidth - 6\tabcolsep) * \real{0.2500}}
  >{\raggedright\arraybackslash}p{(\linewidth - 6\tabcolsep) * \real{0.2500}}
  >{\raggedright\arraybackslash}p{(\linewidth - 6\tabcolsep) * \real{0.2500}}@{}}
\toprule\noalign{}
\begin{minipage}[b]{\linewidth}\raggedright
No.
\end{minipage} & \begin{minipage}[b]{\linewidth}\raggedright
Name
\end{minipage} & \begin{minipage}[b]{\linewidth}\raggedright
Description
\end{minipage} & \begin{minipage}[b]{\linewidth}\raggedright
Default
\end{minipage} \\
\midrule\noalign{}
\endhead
\bottomrule\noalign{}
\endlastfoot
ACTIONX & Define the start of ACTIONX Definition Section. This is then
followed on a new line by any number of ACTIONX records that define the
conditions for which the defined action will be executed and the various
operations to be performed if the conditions are satisfied. & & \\
1-1 & ACTNAME & ACTNAME is a character sting of up to eight characters
in length, that defines the name of this action definition. If ACTNAME
has previously been used by any ACTION series keyword, then the previous
ACTION series definition will be replaced by the definition declared by
this ACTIONX Definition Section. & \\
1-2 & ACTNSTEP & ACTNSTEP is a positive integer that defines the number
times that the ACTNAME definition is executed. ACTIONX definitions are
activated at the end of a time step and this parameter is used to set
how many time steps the ACTNAME definition will be invoked. The default
value of one means that the definition will be executed only once. Use a
large value, for example 10,000 for the definition to be executed at
every time step. Note that the counter only affects successful
evaluations; i.e.~if ACTNSTEP is set equal to one (the default), then
the simulator will test the action at the end of every time step until
it evaluates to true. & 1 \\
1-3 & ACTDELTA & ACTDELTA is a real positive value that defines the
minimum duration of time after the conditions defined on the second
record have been satisfied before the ACTIONX actions are executed. For
example, if ACTDELTA is defaulted the actions will be executed at the
end of the time step for which the conditions are met. If set to say 30,
then a minimum of 30 days will pass before the actions are executed
(assuming field or metric units). & 0.0 \\
days & days & hours & \\
1-4 & / & Record terminated by a ``/'' & Not Applicable \\
2-1 & ACTLHS & ACTLHS is a series of character strings, each up to eight
characters in length, that defines a constant, UDQ defined value, or a
SUMMARY variable on the left hand side of a Boolean conditional test.
The format for ACTLHS is dependent on the SUMMARY variable type:
Aquifer, Block, Field, Group, Region, Time, Well, Well Connection, Well
Local Grid Refinement Connection, or a Well Segment. In addition to
SUMMARY variables, an UDQ defined value or a Constant variable can be
used. The format for the various data types is given in Table 12.7. &
Not Applicable \\
2-2 & ACTTEST & ACTTEST is a defined character string that states the
Boolean operator and must be set to one of the following Boolean
conditionals: For example to test if the field's gas production rate is
less than 600 MMscf/d then one would use: ACTIONX PHASE2 1 / GGPR
`FIELD' \textless{} 600E3 / / \ldots{} ENDACTIO & Not Applicable \\
2-3 & ACTRHS & ACTRHS is a numeric value or a series of character
strings, each up to eight characters in length, that defines a constant,
an UDQ defined value, or a SUMMARY variable on the right hand side of a
Boolean conditional test, as outlined in Table 12.7 (see also ACTLHS).
In the case of well quantities the set of matching wells is captured and
can be used as a general ``well list'' with the symbol `?' in subsequent
well keywords. For example, to shut-in all oil producing wells
(`OP\emph{') with a water cut greater than 90\% for every time the field
water production rate exceeds 60,000 stb/d one would use: ACTIONX
MXWATER 10000 / GWPR 'FIELD' \textgreater{} 60E3 AND / WWCT 'OP}'
\textgreater{} 0.90 / / -- WELL PRODUCTION STATUS -- -- WELL WELL
--LOCATION-- COMPLETION -- NAME STAT I J K FIRST LAST WELOPEN `?' SHUT /
/ ENDACTIO & Not Applicable \\
2-4 & ANDOR & An optional defined character string that specifies a
Boolean operator that must be set to either AND or OR if included on
this record, that links this record with additional records of this
type. For example, to test if the field's gas production rate is less
than 600 MMscf/d after 2020 then one would use: ACTIONX PHASE2 1 / GGPR
`FIELD' \textless{} 600E3 AND / YEAR \textgreater{} 2020 / / \ldots{}
ENDACTIO This item should be left blank if not required. & Not
Applicable \\
2.5 & / & Termination of an ACTIONX Boolean condition record. Note that
multiple numbers of records of this type can be entered with each record
terminated by a ``/'', as illustrated above. & Not Applicable \\
3-1 & / & The Boolean condition section of the ACTIONX keyword is
terminated by an empty line with a single ``/''. & Not Applicable \\
& & The next section contains any number of standard SCHEDULE keywords
that will be executed if the Boolean expression evaluates to true. For
example, to test if the field's gas production rate is less than 600
MMscf/d after 2020 and to open up additional wells if this occurs, then
one would use: ACTIONX PHASE2 1 / GGPR `FIELD' \textless{} 600E3 AND /
YEAR \textgreater{} 2020 / / -- WELL PRODUCTION STATUS -- -- WELL WELL
--LOCATION-- COMPLETION -- NAME STAT I J K FIRST LAST WELOPEN GP10 OPEN
/ GP11 OPEN / / ENDACTIO In theory, most SCHEDULE keywords can be used
in an ACTIONX Definition Section here, except for the time stepping
keywords, i.e, TSTEP and DATES. See Table 12.8 for a list of the
SCHEDULE keywords that are known to work with the ACTIONX keyword. & Not
Applicable \\
ENDACTIO & Define the end of ACTIONX Definition Section. & Not
Applicable & \\
Notes: & & & \\
\end{longtable}

\emph{Table 12.6: ACTIONX Keyword Description}

The variable types and the associated definitions that are available for
use with Boolean conditionals are outlined in Table 12.7.

\begin{longtable}[]{@{}
  >{\raggedright\arraybackslash}p{(\linewidth - 2\tabcolsep) * \real{0.5000}}
  >{\raggedright\arraybackslash}p{(\linewidth - 2\tabcolsep) * \real{0.5000}}@{}}
\toprule\noalign{}
\begin{minipage}[b]{\linewidth}\raggedright
Variable Type
\end{minipage} & \begin{minipage}[b]{\linewidth}\raggedright
Description
\end{minipage} \\
\midrule\noalign{}
\endhead
\bottomrule\noalign{}
\endlastfoot
AQUIFER & AQUIFER variable consists of two parameters: \\
BLOCK & BLOCK variable consists of four parameters: The NX, NY, and NZ
parameters are defined on the DIMENS keyword in the RUNSPEC section. \\
CONSTANTS & CONSTANTS can consist of one or optionally two
parameters: \\
FIELD & The FIELD variable consists of any field SUMMARY variable; for
example the Field average Pressure, as shown below: ACTIONX WIPHASE 1 /
FPR \textless{} 2500 / / \ldots{} ENDACTIO The above would action a set
of SCHEDULE keywords if the field average pressure fell below 2,500 psia
for a run using FIELD units. \\
GROUP & GROUP variable definition consists of: To enable an action for
when the field's oil production rate drops below 20,000 stb/d then one
could use. ACTIONX OILMIN 1 / GOPR `FIELD' \textless{} 20.0E3 / /
\ldots{} ENDACTIO \\
REGION & REGION variable definition consists of: For example, ACTIONX
WIPHASE 1 / RPR 0 \textless{} 2500 / / \ldots{} ENDACTIO Would action a
set of SCHEDULE keywords if the field average pressure fell below 2,500
psia for a run using FIELD units. \\
TIME & TIME variable definition consists of one parameter that can have
three values: DAY for the current simulation day of the month, MNTH for
the current simulation month, and YEAR for the current simulation year.
Thus, to set an action for April 1, 2025 one would use: ACTIONX ACT01 1
/ DAY = 1 AND / MNTH =`APR' AND / YEAR = 2025 / / \ldots{} ENDACTIO Note
that the value for the MNTH variable, `APR', in the example, can also be
entered without the quotes, that is: MNTH = APR AND / In addition,
numerical values for MNTH, similar to DAY and YEAR, are also permitted
and are converted to the nearest integer for comparison, so for example:
ACTIONX ACT01 1 / DAY = 0.95 AND / MNTH = 4.40 AND / YEAR = 2024.9 / /
\ldots{} ENDACTIO Would again result in the action taking place on April
1, 2025. \\
WELL & WELL variable definition consists of: Note that the use of well
lists is an OPM Flow specific feature. To reduce the tubing head
pressure constraint for when any of the oil producers' oil rate drop
below 100 stb/d then one could use. ACTIONX WOILMIN 1 / WOPR `OP\emph{'
\textless{} 100.0 / / -- -- FLOW WELLS THROUGH LOW PRESSURE SEPARATOR --
-- WELL WELL TARGET -- NAME TARG VALUE WELTARG 'OP}' THP 150 / /
ENDACTIO \\
WELL CONNECTION & WELL CONNECTION variable definition is comprised of:
The NX, NY, and NZ parameters are defined on the DIMENS keyword in the
RUNSPEC section. \\
WELL COMPLETION & WELL COMPLETION variable definition is comprised
of: \\
WELL LOCAL GRID REFINEMENT CONNECTION & WELL LOCAL GRID REFINEMENT
CONNECTION variable definition is comprised of: The NX, NY, and NZ
parameters are defined on either the CARFIN or RADFIN keywords in the
GRID section depending upon whether a Cartesian or radial local grid
refinement is being utilized. Note Local Grid Refinements are currently
not implemented in OPM Flow. \\
WELL SEGMENT & WELL SEGMENT variable definition consists of: Note that
the total number of wells should be defined via the WELLSDIMS keyword
and the number of multi-segment wells should be declared on the WSEGDIMS
keyword, both keywords are in the RUNSPEC section. \\
\end{longtable}

\emph{Table 12.7: ACTIONX Variable Definitions}

See also the ACTDIMS and UDADIMS keyword in the RUNSPEC section to
define the dimensions for the ACTIONX keyword and associated variables.
In addition, the EXIT keyword in the SCHEDULE section that allows for
terminating the simulation for when a condition within an ACTIONX
definition is satisfied.

Although most SCHEDULE keywords should work with the ACTIONX keyword,
Table 12.8 shows the status of keywords that have been tested and known
to work, together with keywords that are currently planned to be
implemented.

\begin{longtable}[]{@{}
  >{\raggedright\arraybackslash}p{(\linewidth - 10\tabcolsep) * \real{0.1667}}
  >{\raggedright\arraybackslash}p{(\linewidth - 10\tabcolsep) * \real{0.1667}}
  >{\raggedright\arraybackslash}p{(\linewidth - 10\tabcolsep) * \real{0.1667}}
  >{\raggedright\arraybackslash}p{(\linewidth - 10\tabcolsep) * \real{0.1667}}
  >{\raggedright\arraybackslash}p{(\linewidth - 10\tabcolsep) * \real{0.1667}}
  >{\raggedright\arraybackslash}p{(\linewidth - 10\tabcolsep) * \real{0.1667}}@{}}
\toprule\noalign{}
\begin{minipage}[b]{\linewidth}\raggedright
ACTIONX Schedule Section Keywords Status
\end{minipage} & \begin{minipage}[b]{\linewidth}\raggedright
\end{minipage} & \begin{minipage}[b]{\linewidth}\raggedright
\end{minipage} & \begin{minipage}[b]{\linewidth}\raggedright
\end{minipage} & \begin{minipage}[b]{\linewidth}\raggedright
\end{minipage} & \begin{minipage}[b]{\linewidth}\raggedright
\end{minipage} \\
\midrule\noalign{}
\endhead
\bottomrule\noalign{}
\endlastfoot
Action Keywords & Group Keywords & Well Keywords & Connection Keywords &
Miscellaneous Keywords & \\
ACTIONX & GCONINJE & WCONHIST & WPIMULT & COMPDAT & BOX \\
UDQ & GCONPROD & WCONINJE & WSEGVALV & COMPLUMP & BRANPROP \\
& GCONSUMP & WCONINJH & WTEST & COMPSEGS & ECHO \\
& GEFAC & WCONPROD & WTMULT & & ENDBOX \\
& GLIFTOPT & WDFAC & & & EXIT5 \\
& GRUPNET & WECON & & & INCLUDE7 \\
& GRUPTARG & WEFAC & & & MULTX \\
& GRUPTREE & WELOPEN & & & MULTX- \\
& GSATINJE & WELPI & & & MULTY \\
& GSATPROD & WELSEGS & & & MULTY- \\
& & WELSPECS6 & & & MULTZ \\
& & WELTARG & & & MULTZ- \\
& & WGRUPCON & & & NEXT \\
& & WINJMULT & & & NEXTSTEP \\
& & WLIST & & & NOECHO \\
Notes: & & & & & \\
\end{longtable}

\emph{Table 12.8: ACTIONX Schedule Section Keywords Status}

As mentioned previously, the UDQ keyword stipulates the variables and
operations used to access the User Defined Quantities features in OPM
Flow. UDQ variables can be constants, SUMMARY variables, as defined in
the SUMMARY section, or a formula using various mathematical functions
together with constants and SUMMARY variables.

\begin{tcolorbox}[enhanced jigsaw, title=\textcolor{quarto-callout-note-color}{\faInfo}\hspace{0.5em}{Note}, coltitle=black, breakable, toprule=.15mm, colbacktitle=quarto-callout-note-color!10!white, leftrule=.75mm, opacitybacktitle=0.6, colback=white, bottomtitle=1mm, toptitle=1mm, titlerule=0mm, left=2mm, colframe=quarto-callout-note-color-frame, arc=.35mm, rightrule=.15mm, bottomrule=.15mm, opacityback=0]

Within an ACTIONX Definition Section any UDQ variables utilizing group
and well variables, must have their associated groups and wells
previously fully defined in the commercial simulator, otherwise an error
will occur. For example, if a well's GOR is being used as part of a UDQ
definition, then the well must be fully characterized prior to declaring
the UDQ definition. This restriction does not apply to OPM Flow;
however, it should be considered if the same deck is to be run with both
simulators.

\end{tcolorbox}

User Defined Quantities can also be used as User Defined Arguments
(``UDA'') in the SCHEDULE section with various group, well, and
connection keywords. In this case, the UDA variables are used to replace
numerical values on these keywords by UDA variables that have been
defined by the UDQ keyword. For example, if we wish to make the oil rate
for certain wells to be a function of their water cut, then one can
define the function using the UDQ keyword that results in a UDQ
variable, WU\_OPR say, and then use WU\_OPR as a UDA variable on the
WCONPROD keyword for the ORAT parameter. See Table 12.76 for a list of
keywords that can be used with UDA variables in the UDQ - Declare User
Define Quantities (``UDQ'') keyword section.

\subsubsection{Examples}\label{examples-68}

The first example uses the UDQ keyword to sort the oil wells from high
water cut to low, via the WU\_WLIST variable, and then use the ACTIONX
keyword to shut-in the worst offending well when the field's water
production is greater than 30,000 stb/d.

\begin{verbatim}
--
--      DEFINE START OF USER DEFINED QUANTITY SECTION
--
UDQ
--
--      OPERATOR VARIABLE  EXPRESSION
--
DEFINE  WU_WCUT   WWCT 'OP*'                        / WELL WWCT LIST
DEFINE  WU_LIST   SORTD(WU_WCUT)                    / WELL WWCT LIST SORTED
/                                                     END OF UDQ SECTION
--
--      DEFINE START OF ACTIONX SECTION
--
ACTIONX
        WSHUTIN        10                                                        /
        GWPR  'FIELD' > 30E3  AND                                              /
        WU_LIST 'OP*' = 1                                                      /
/

--
--      DEFINE WELL AND WELL CONNECTIONS FLOWING STATUS
--
--  WELL WELL   --LOCATION--  COMPLETION
--  NAME STAT     I   J    K  FIRST LAST
WELOPEN
'?'      SHUT                                                                  /
'?'      SHUT     0   0    0     0     0                                       /
/

ENDACTIO
\end{verbatim}

Apart from checking that the field's water production rate is greater
than 30,000 stb/d the Boolean conditional also checks that there is at
least one well in the sorted well list. Notice also the use of `?'
symbol as a substitution of the well name and that the ACTIONX WSHUTIN
series of commands will be executed a total of ten times.

The second example checks to see if the field's gas rate is below 600
MMscf/d and if the simulation time is greater that July 1, 2030. If it
is, then compression is installed by re-setting all the gas producing
well's THP and BHP pressures to 450 psia and 300 psia respectively. In
addition all gas wells currently shut-in are tested to see if they can
be opened up under the new THP and BHP constraints.

\begin{verbatim}
--
--      START ACTIONX FIELD PHASE-3 AUTOMATIC COMPRESSION
--
ACTIONX
        PHASE-3      1                                                         /
        GGPR  'FIELD' < 600E3 AND                                              /
        DAY           >= 1     AND                                             /
        MNTH          >= JUL   AND                                             /
        YEAR          >= 2030                                                  /
/

--
--      INSTALL COMPRESSION AND RESET WELL THP AND BHPS
--
--  WELL    WELL   TARGET
--  NAME    TARG   VALUE
WELTARG
'GP*'       THP     450                                                        /
'GP*'       BHP     300                                                        /
/

--
--      TEST AND OPEN ALL WELLS UNDER COMPRESSION CONSTRAINTS
--
--  WELL    TEST   CLOSE   NO.     START
--  NAME    INTV   CHECK   CHECK   TIME
WTEST
'GP*'       1.0    PE      1       3                                           /
/

--
--      END OF ACTIONX FIELD PHASE-3 AUTOMATIC COMPRESSION DEFINITION
--
ENDACTIO

\end{verbatim}

\subsection{AITS -- Activate Intelligent Time
Stepping}\label{aits-activate-intelligent-time-stepping-1}

Turns on the commercial simulator's intelligent time stepping.

See AITS -- Activate Intelligent Time Stepping in the RUNSPEC section.

\subsection{AITSOFF -- Deactivate Intelligent Time
Stepping}\label{aitsoff-deactivate-intelligent-time-stepping-1}

Turns off the commercial simulator's intelligent time stepping.

See AITSOFF -- Deactivate Intelligent Time Stepping in the RUNSPEC
Section

\subsection{APILIM -- Define API Tracking Grid Block
Limits}\label{apilim-define-api-tracking-grid-block-limits}

\begin{longtable}[]{@{}
  >{\raggedright\arraybackslash}p{(\linewidth - 14\tabcolsep) * \real{0.1250}}
  >{\raggedright\arraybackslash}p{(\linewidth - 14\tabcolsep) * \real{0.1250}}
  >{\raggedright\arraybackslash}p{(\linewidth - 14\tabcolsep) * \real{0.1250}}
  >{\raggedright\arraybackslash}p{(\linewidth - 14\tabcolsep) * \real{0.1250}}
  >{\raggedright\arraybackslash}p{(\linewidth - 14\tabcolsep) * \real{0.1250}}
  >{\raggedright\arraybackslash}p{(\linewidth - 14\tabcolsep) * \real{0.1250}}
  >{\raggedright\arraybackslash}p{(\linewidth - 14\tabcolsep) * \real{0.1250}}
  >{\raggedright\arraybackslash}p{(\linewidth - 14\tabcolsep) * \real{0.1250}}@{}}
\toprule\noalign{}
\begin{minipage}[b]{\linewidth}\raggedright
RUNSPEC
\end{minipage} & \begin{minipage}[b]{\linewidth}\raggedright
GRID
\end{minipage} & \begin{minipage}[b]{\linewidth}\raggedright
EDIT
\end{minipage} & \begin{minipage}[b]{\linewidth}\raggedright
PROPS
\end{minipage} & \begin{minipage}[b]{\linewidth}\raggedright
REGIONS
\end{minipage} & \begin{minipage}[b]{\linewidth}\raggedright
SOLUTION
\end{minipage} & \begin{minipage}[b]{\linewidth}\raggedright
SUMMARY
\end{minipage} & \begin{minipage}[b]{\linewidth}\raggedright
SCHEDULE
\end{minipage} \\
\midrule\noalign{}
\endhead
\bottomrule\noalign{}
\endlastfoot
\end{longtable}

\subsubsection{Description}\label{description-870}

The APILIM keyword defines API Tracking grid block limits for when API
Tracking has been activated via the API keyword in the RUNSPEC section.
The keyword enables the simulator to monitor the grid blocks outside the
limits defined on the keyword, as well as to optionally contrain the
values within a given range.

This keyword is not supported by OPM Flow but would change the results
if supported so the simulation will be stopped.

\subsection{AQUCHGAS -- Define Constant Pressure Gas Analytical Aquifer
Properties}\label{aquchgas-define-constant-pressure-gas-analytical-aquifer-properties-1}

The AQUCHGAS keyword defines the properties of constant pressure gas
analytical aquifers.

See AQUCHGAS -- Define Constant Pressure Gas Analytical Aquifer
Properties in the SOLUTION section for a full description.

\subsection{AQUCHWAT -- Define Constant Pressure Water Analytical
Aquifer
Properties}\label{aquchwat-define-constant-pressure-water-analytical-aquifer-properties-1}

The AQUCHWAT keyword defines the properties of constant pressure water
analytical aquifers.

See AQUCHWAT -- Define Constant Pressure Water Analytical Aquifer
Properties in the SOLUTION section for a full description.

\subsection{AQUCT -- Define Carter-Tracy Analytical
Aquifers}\label{aquct-define-carter-tracy-analytical-aquifers-3}

The AQUCT keyword defines Carter-Tracy {[}Carter, R. D. and Tracy, G. W.
``An Improved Method for Calculating Water Influx,'' Transactions of
AIME (1960) 219, 215-417.{]} analytical aquifers, the properties of the
aquifer, including the Carter-Tracy aquifer influence function
associated with the aquifer, defined by the AQUTAB keyword in the PROPS
section.

See AQUCT -- Define Carter-Tracy Analytical Aquifers in the GRID
section.

\subsection{AQUCWFAC -- Modify Constant Pressure Water Analytical
Aquifer
Properties}\label{aqucwfac-modify-constant-pressure-water-analytical-aquifer-properties}

\begin{longtable}[]{@{}
  >{\raggedright\arraybackslash}p{(\linewidth - 14\tabcolsep) * \real{0.1250}}
  >{\raggedright\arraybackslash}p{(\linewidth - 14\tabcolsep) * \real{0.1250}}
  >{\raggedright\arraybackslash}p{(\linewidth - 14\tabcolsep) * \real{0.1250}}
  >{\raggedright\arraybackslash}p{(\linewidth - 14\tabcolsep) * \real{0.1250}}
  >{\raggedright\arraybackslash}p{(\linewidth - 14\tabcolsep) * \real{0.1250}}
  >{\raggedright\arraybackslash}p{(\linewidth - 14\tabcolsep) * \real{0.1250}}
  >{\raggedright\arraybackslash}p{(\linewidth - 14\tabcolsep) * \real{0.1250}}
  >{\raggedright\arraybackslash}p{(\linewidth - 14\tabcolsep) * \real{0.1250}}@{}}
\toprule\noalign{}
\begin{minipage}[b]{\linewidth}\raggedright
RUNSPEC
\end{minipage} & \begin{minipage}[b]{\linewidth}\raggedright
GRID
\end{minipage} & \begin{minipage}[b]{\linewidth}\raggedright
EDIT
\end{minipage} & \begin{minipage}[b]{\linewidth}\raggedright
PROPS
\end{minipage} & \begin{minipage}[b]{\linewidth}\raggedright
REGIONS
\end{minipage} & \begin{minipage}[b]{\linewidth}\raggedright
SOLUTION
\end{minipage} & \begin{minipage}[b]{\linewidth}\raggedright
SUMMARY
\end{minipage} & \begin{minipage}[b]{\linewidth}\raggedright
SCHEDULE
\end{minipage} \\
\midrule\noalign{}
\endhead
\bottomrule\noalign{}
\endlastfoot
\end{longtable}

\subsubsection{Description}\label{description-871}

The AQUCWFAC keyword modifies the datum depth and pressure for all
aquifers specified by the AQUCHWAT keyword in the SOLUTION or SCHEDULE
sections.

This keyword is not supported by OPM Flow but would change the results
if supported so the simulation will be stopped.

\subsection{AQUFETP -- Define Fetkovich Analytical
Aquifers}\label{aqufetp-define-fetkovich-analytical-aquifers-1}

The AQUFETP keyword defines Fetkovich Analytical aquifers and the
aquifer properties. Each row entry in the AQUFETP keyword defines one
Fetkovich Analytical aquifer. In order to fully define this type of
aquifer, the aquifer must be connected to the reservoir using the
AQUANCON keyword in the GRID or SOLUTION sections.

See AQUFETP -- Define Fetkovich Analytical Aquifers in the SOLUTION
section for a full description.

\subsection{AQUFLUX -- Define Constant Flux Analytical
Aquifer}\label{aquflux-define-constant-flux-analytical-aquifer-1}

The AQUFLUX keyword defines the properties of Constant Flux Analytical
Aquifers, that allows for a constant water influx to the connected grid
blocks

See AQUFLUX - Define Constant Flux Analytical Aquifer in the SOLUTION
section for a full description.

\subsection{BCPROP -- Define Boundary Conditions
Properties}\label{bcprop-define-boundary-conditions-properties}

\begin{longtable}[]{@{}
  >{\raggedright\arraybackslash}p{(\linewidth - 14\tabcolsep) * \real{0.1250}}
  >{\raggedright\arraybackslash}p{(\linewidth - 14\tabcolsep) * \real{0.1250}}
  >{\raggedright\arraybackslash}p{(\linewidth - 14\tabcolsep) * \real{0.1250}}
  >{\raggedright\arraybackslash}p{(\linewidth - 14\tabcolsep) * \real{0.1250}}
  >{\raggedright\arraybackslash}p{(\linewidth - 14\tabcolsep) * \real{0.1250}}
  >{\raggedright\arraybackslash}p{(\linewidth - 14\tabcolsep) * \real{0.1250}}
  >{\raggedright\arraybackslash}p{(\linewidth - 14\tabcolsep) * \real{0.1250}}
  >{\raggedright\arraybackslash}p{(\linewidth - 14\tabcolsep) * \real{0.1250}}@{}}
\toprule\noalign{}
\begin{minipage}[b]{\linewidth}\raggedright
RUNSPEC
\end{minipage} & \begin{minipage}[b]{\linewidth}\raggedright
GRID
\end{minipage} & \begin{minipage}[b]{\linewidth}\raggedright
EDIT
\end{minipage} & \begin{minipage}[b]{\linewidth}\raggedright
PROPS
\end{minipage} & \begin{minipage}[b]{\linewidth}\raggedright
REGIONS
\end{minipage} & \begin{minipage}[b]{\linewidth}\raggedright
SOLUTION
\end{minipage} & \begin{minipage}[b]{\linewidth}\raggedright
SUMMARY
\end{minipage} & \begin{minipage}[b]{\linewidth}\raggedright
SCHEDULE
\end{minipage} \\
\midrule\noalign{}
\endhead
\bottomrule\noalign{}
\endlastfoot
\end{longtable}

\subsubsection{Description}\label{description-872}

The \hyperref[REF_HEADING_KEYWORD_BCPROP]{BCPROP} keyword defines the
type and properties of the boundary conditions.

Together the \hyperref[REF_HEADING_KEYWORD_BCCON]{BCCON} and
\hyperref[REF_HEADING_KEYWORD_BCPROP]{BCPROP} keywords define the
boundary conditions for the model, and can be used to set boundary
conditions for when external influx or efflux volumes are influencing
the reservoir pressure and production history. For example, when the
average reservoir pressure remains constant throughout the production
period due to water influx, or gas migration from an external source.

\begin{tcolorbox}[enhanced jigsaw, title=\textcolor{quarto-callout-note-color}{\faInfo}\hspace{0.5em}{Note}, coltitle=black, breakable, toprule=.15mm, colbacktitle=quarto-callout-note-color!10!white, leftrule=.75mm, opacitybacktitle=0.6, colback=white, bottomtitle=1mm, toptitle=1mm, titlerule=0mm, left=2mm, colframe=quarto-callout-note-color-frame, arc=.35mm, rightrule=.15mm, bottomrule=.15mm, opacityback=0]

This is an OPM Flow specific keyword and will therefore cause an error
in the commercial simulator.

\end{tcolorbox}

\begin{longtable}[]{@{}
  >{\raggedright\arraybackslash}p{(\linewidth - 6\tabcolsep) * \real{0.2500}}
  >{\raggedright\arraybackslash}p{(\linewidth - 6\tabcolsep) * \real{0.2500}}
  >{\raggedright\arraybackslash}p{(\linewidth - 6\tabcolsep) * \real{0.2500}}
  >{\raggedright\arraybackslash}p{(\linewidth - 6\tabcolsep) * \real{0.2500}}@{}}
\toprule\noalign{}
\begin{minipage}[b]{\linewidth}\raggedright
No.
\end{minipage} & \begin{minipage}[b]{\linewidth}\raggedright
Name
\end{minipage} & \begin{minipage}[b]{\linewidth}\raggedright
Description
\end{minipage} & \begin{minipage}[b]{\linewidth}\raggedright
Default
\end{minipage} \\
\midrule\noalign{}
\endhead
\bottomrule\noalign{}
\endlastfoot
Field & Metric & Laboratory & \\
1 & INDEX & A positive integer that identifies the boundary condition. &
None \\
2 & TYPE & A defined character string that defines the type of boundary
condition to be applied, and should be set to one of the following
character strings: & None \\
3 & COMPONENT & A defined character string that sets fluid type used in
the boundary calculations, and should be set to one of the following
character strings: COMPONENT must be declared and not equal to NONE if
TYPE has been set to either DIRICHLET or RATE. MICR requires the
\hyperref[REF_HEADING_KEYWORD_BIOFILM]{BIOFILM} or MICP keyword, and
OXYG or UREA requires the MICP keyword. & NONE \\
4 & RATE & A real value that defines the constant mass rate per unit
area of the specified COMPONENT to be injected or withdrawn at the
boundary, when TYPE has been set to RATE. Note a negative value implies
an influx rate, whereas, a positive value indicates an efflux. & 0.0 \\
lb/day/ft2 & kg/day/m2 & gm/hour/cm2 & \\
5 & PRESS & PRESS is a real positive value that defines the constant
pressure boundary condition. PRESS should only be entered if TYPE has
been set to DIRICHLET. If the pressure at the boundary is less than
PRESS, then the fluid type declared via COMPONENT will flow across the
boundary. The default value of 1* will use the simulator's calculated
value based on data entered via the EQUIL keyword in the SOLUTION
section. & 1* \\
psia & barsa & atma & \\
6 & TEMP & TEMP is a real positive number that defines the constant
temperature boundary condition. TEMP should only be entered if TYPE has
been set to DIRICHLET or THERMAL. The default value of 1* will use the
simulator's calculated value based on data entered via one of the
following reservoir temperature keywords: RTEMP, RTEMPA, RTEMPVD, TEMPI,
or TEMPVD, in the SOLUTION section. Note that all of the aforementioned
reservoir temperature keywords, except for TEMPI, may also be used in
the PROPS section as well. & 1* \\
oF & oC & oC & \\
7 & MECHTYPE & A defined character string that defines the type of
geo-mechanical boundary condition to be applied, and should be set to
one of the following character strings: & NONE \\
8 & FIXEDX & A positive integer that identifies the whether the boundary
is free or fixed in the x-direction. A value of 0 implies the boundary
is free to move in the x-direction. A value \textgreater{} 0 implies the
boundary is fixed in the x-direction with displacement defined by DISPX,
& 1 \\
9 & FIXEDY & A positive integer that identifies the whether the boundary
is free or fixed in the y-direction. A value of 0 implies the boundary
is free to move in the y-direction. A value \textgreater{} 0 implies the
boundary is fixed in the y-direction with displacement defined by DISPY,
& 1 \\
10 & FIXEDZ & A positive integer that identifies the whether the
boundary is free or fixed in the z-direction. A value of 0 implies the
boundary is free to move in the z-direction. A value \textgreater{} 0
implies the boundary is fixed in the z-direction with displacement
defined by DISPZ, & 1 \\
& STRESSXX & A real value that defines the diagonal component
\({σ}_{\mathit{xx}}\) of the stress tensor. & 0 \\
psia & barsa & atma & \\
& STRESSYY & A real value that defines the diagonal component
\({σ}_{\mathit{yy}}\) of the stress tensor. & 0 \\
psia & barsa & atma & \\
& STRESSZZ & A real value that defines the diagonal component
\({σ}_{\mathit{zz}}\) of the stress tensor. & 0 \\
psia & barsa & atma & \\
& DISPX & A real value that defines the x-direction displacement
boundary condition. & 0 \\
feet & m & cm & \\
& DISPY & A real value that defines the y-direction displacement
boundary condition. & 0 \\
feet & m & cm & \\
& DISPZ & A real value that defines the z-direction displacement
boundary condition. & 0 \\
feet & m & cm & \\
Notes: & & & \\
\end{longtable}

\emph{Table 12.3.16.1: BCPROP Keyword Description}

See also the AQUFLUX keyword that is supported by OPM Flow in both the
SOLUTION and SCHEDULE sections, to define a constant flux analytical
aquifer.

If the \hyperref[REF_HEADING_KEYWORD_BCCON]{BCCON} and
\hyperref[REF_HEADING_KEYWORD_BCPROP]{BCPROP} keywords are not present
in the input deck, then the boundary conditions for the model are set to
be no flow, which is the normal behavior in both OPM Flow and the
commercial simulator.

The BC keyword has been replaced by the
\hyperref[REF_HEADING_KEYWORD_BCCON]{BCCON} and
\hyperref[REF_HEADING_KEYWORD_BCPROP]{BCPROP} keywords.

\subsubsection{Examples}\label{examples-69}

The first example shows how to set a constant pressure boundary using
TYPE equal to FREE:

\begin{verbatim}
--
--       DEFINE BOUNDARY CONDITIONS CONNECTION (OPM FLOW GRID KEYWORD)
--
--INDEX  ---------- BOX ---------   BC
--       I1  I2   J1  J2   K1  K2   DIRC
BCCON
1        1   1    1   1*   1   1*   X-    /
2        1   1*   1   1    1   1*   Y     /
/

--
--       DEFINE BOUNDARY CONDITIONS PROPERTIES (OPM FLOW SCHEDULE KEYWORD)
--
--INDEX  BC         BC     BC    BC     BC
--       TYPE       COMPT  RATE  PRESS  TEMP
BCPROP
1        FREE       1*     1*    1*     1*    /
2        FREE                                 /
/
\end{verbatim}

With this option it is only necessary to define the boundary cells and
all the other parameters (COMPONENT, RATE, PRESS, and TEMP) can be
defaulted, as they are ignored when TYPE equals FREE.

The next example is based on NX, NY and NZ equal to 20, 1, 10
respectively, on the DIMENS keyword in the RUNSPEC section, and shows
how different boundary types can be assigned to different parts of the
model.

\begin{verbatim}
--
--       DEFINE BOUNDARY CONDITIONS CONNECTION (OPM FLOW GRID KEYWORD)
--
--INDEX  ---------- BOX ---------   BC
--       I1  I2   J1  J2   K1  K2   DIRC
BCCON
1        1   1    1   1    1   10   X-    /
2        20  20   1   1    1   10   X     /
/

--
--       DEFINE BOUNDARY CONDITIONS PROPERTIES (OPM FLOW SCHEDULE KEYWORD)
--
--INDEX  BC         BC     BC    BC     BC
--       TYPE       COMPT  RATE  PRESS  TEMP
BCPROP
1        DIRICHLET  GAS    1*    256.0  100.0 /
2        FREE       4*                        /
/
\end{verbatim}

The last example shows how the DIRICHLET boundary condition option may
be used:

\begin{verbatim}
--
--       DEFINE BOUNDARY CONDITIONS CONNECTION (OPM FLOW GRID KEYWORD)
--
--INDEX  ---------- BOX ---------   BC
--       I1  I2   J1  J2   K1  K2   DIRC
BCCON
1        1   1    1   1*   1   1*   X-    /
2        1   1*   1   1    1   1*   Y     /
/

--
--       DEFINE BOUNDARY CONDITIONS PROPERTIES (OPM FLOW SCHEDULE KEYWORD)
--
--INDEX  BC         BC     BC    BC     BC
--       TYPE       COMPT  RATE  PRESS  TEMP
BCPROP
1        DIRICHLET  WAT    1*    256.0  100.0 /
2        DIRICHLET  WAT    1*    1*     100.0 /
/
\end{verbatim}

Here, the first line sets both the pressure and temperature at the
boundary, and the second line defaults the pressure entry, so that the
simulator calculated initial boundary pressure will be used.

\subsection{BOUNDARY -- Define a Boundary Box for
Printing}\label{boundary-define-a-boundary-box-for-printing-4}

The BOUNDARY keyword defines a rectangular grid for printing various
arrays to the output print file (*.PRT); thus, avoiding printing all the
elements in the selected array.

See BOUNDARY -- Define a Boundary Box for Printing in the GRID section
for a full description.

\subsection{BOX -- Define a Range of Grid Blocks to Enter Property
Data}\label{box-define-a-range-of-grid-blocks-to-enter-property-data-5}

BOX defines a range of grid blocks for which subsequent data is assigned
for all the cells in the defined BOX. Note that the BOX grid is reset by
the keyword ENDBOX by resetting the current defined BOX to be the whole
grid. The keyword can be used for any array and for all grid types.

See BOX - Define a Range of Grid Blocks to Enter Property Data in the
GRID section for a full description.

\subsection{BRANPROP -- Define Network Branch Properties for Extended
Network
Option}\label{branprop-define-network-branch-properties-for-extended-network-option}

\begin{longtable}[]{@{}
  >{\raggedright\arraybackslash}p{(\linewidth - 14\tabcolsep) * \real{0.1250}}
  >{\raggedright\arraybackslash}p{(\linewidth - 14\tabcolsep) * \real{0.1250}}
  >{\raggedright\arraybackslash}p{(\linewidth - 14\tabcolsep) * \real{0.1250}}
  >{\raggedright\arraybackslash}p{(\linewidth - 14\tabcolsep) * \real{0.1250}}
  >{\raggedright\arraybackslash}p{(\linewidth - 14\tabcolsep) * \real{0.1250}}
  >{\raggedright\arraybackslash}p{(\linewidth - 14\tabcolsep) * \real{0.1250}}
  >{\raggedright\arraybackslash}p{(\linewidth - 14\tabcolsep) * \real{0.1250}}
  >{\raggedright\arraybackslash}p{(\linewidth - 14\tabcolsep) * \real{0.1250}}@{}}
\toprule\noalign{}
\begin{minipage}[b]{\linewidth}\raggedright
RUNSPEC
\end{minipage} & \begin{minipage}[b]{\linewidth}\raggedright
GRID
\end{minipage} & \begin{minipage}[b]{\linewidth}\raggedright
EDIT
\end{minipage} & \begin{minipage}[b]{\linewidth}\raggedright
PROPS
\end{minipage} & \begin{minipage}[b]{\linewidth}\raggedright
REGIONS
\end{minipage} & \begin{minipage}[b]{\linewidth}\raggedright
SOLUTION
\end{minipage} & \begin{minipage}[b]{\linewidth}\raggedright
SUMMARY
\end{minipage} & \begin{minipage}[b]{\linewidth}\raggedright
SCHEDULE
\end{minipage} \\
\midrule\noalign{}
\endhead
\bottomrule\noalign{}
\endlastfoot
\end{longtable}

\subsubsection{Description}\label{description-873}

BRANPROP defines network branch properties for the extended network
option for when the Extended Network Model has been activated by the
NETWORK keyword in the RUNSPEC section. There are two types of network
facilities in the simulator, the Standard Network model, which is
defined with the GRUPNET keyword in the SCHEDULE section and the
Extended Network Model defined by the BRANPROP and NODEPROP keywords,
again in the SCHEDULE section.

For the Extended Network Model the group hierarchy can be different to
that defined by the GRUPTREE keyword; however, the bottom most nodes in
the network tree associated with wells, must be the same as that defined
by the GRUPTREE keyword.

\begin{longtable}[]{@{}
  >{\raggedright\arraybackslash}p{(\linewidth - 6\tabcolsep) * \real{0.2500}}
  >{\raggedright\arraybackslash}p{(\linewidth - 6\tabcolsep) * \real{0.2500}}
  >{\raggedright\arraybackslash}p{(\linewidth - 6\tabcolsep) * \real{0.2500}}
  >{\raggedright\arraybackslash}p{(\linewidth - 6\tabcolsep) * \real{0.2500}}@{}}
\toprule\noalign{}
\begin{minipage}[b]{\linewidth}\raggedright
No.
\end{minipage} & \begin{minipage}[b]{\linewidth}\raggedright
Name
\end{minipage} & \begin{minipage}[b]{\linewidth}\raggedright
Description
\end{minipage} & \begin{minipage}[b]{\linewidth}\raggedright
Default
\end{minipage} \\
\midrule\noalign{}
\endhead
\bottomrule\noalign{}
\endlastfoot
Field & Metric & Laboratory & \\
1 & DOWNNODE & A character string of up to eight characters in length
that defines the down stream node name for this branch, that is the node
closest to the wells. Thus for a production network, this will be an
inlet node as the wells are importing fluid into the branch node.
Whereas for an injection node, this is an outlet node as injection fluid
is being exported to the wells. & None \\
2 & UPNODE & A character string of up to eight characters in length that
defines the up stream node name for this branch, that is the node
furthermost from the wells. Thus for a production network, this will be
an outlet node as the wells are exporting fluid from the branch node.
Whereas for an injection node, this is an inlet node as the wells are
importing the injection fluid. & None \\
3 & VFPTAB & A positive integer greater than or equal to zero that
defines the vertical lift performance table to be used for calculating
the pressure behavior between the inlet (for production) or outlet (for
injection) node (DOWNNODE) and the outlet (for production) or inlet (for
injection) node (UPNODE). For a production network this must reference a
table associated with the VFPPROD keyword, and for an injection network
a table declared via the VFPINJ keyword. Both keywords are in the
SCHEDULE section. If the pressure behavior between the two nodes is zero
(no pressure loss in the network branch), then a value of 9999 should be
entered for this variable. A value of zero for VFPTAB removes the branch
from the extended network and the resulting flows are ignored in the
network flow stream. & None \\
4 & ALQ-NODE & A real positive value that defines the artificial lift
quantity to be used in conjunction with the VFPPROD assigned to the
branch via the VPFTAB variable. VFPTAB vertical lift performance table
and the artificial lift quantity ALQ-NODE are used with the branch fluid
rates to calculate the pressure behavior through the branch. For a
network this can be considered to be either a pump to pump fluid through
the network or a compressor to compress gas to a higher export pressure.
Basically, ALQ-NODE is used to reduce the pressure loss through the
branch. Note that the units for ALQ-NODE are dependent on the associated
variable on the VFPPROD keyword. Should be set to zero if ALQ-DEN is set
to either DENO or DENG, or if the branch is associated with an automatic
compressor. & 0.0 \\
5 & ALQ-DEN & A defined character string that defines that ALQ-NODE
variable represents either as a surface density or as an artificial lift
quantity for a pump or a compressor, and should be set to one of the
following: The VFPPROD keyword should be consistent with this variable,
that is, if ALQ-DEN is set to to DENO then the surface oil should be
used as the ALQ variable on the VFPPROD keyword. & NONE \\
Notes: & & & \\
\end{longtable}

\emph{Table 12.9: BRANPROP Keyword Description}

See also the NETWORK keyword in the RUNSPEC section and the NODEPROP
keyword in the SCHEDULE section.

\subsubsection{Example}\label{example-504}

Given the following Extended Network model in Figure 12.3.

\pandocbounded{\includegraphics[keepaspectratio]{chapters/images/Frame66_c40b193fe195.png}}
\pandocbounded{\includegraphics[keepaspectratio]{chapters/images/Image60_c40b193fe195.png}}

First the Extended Network model should be used invoked in the RUNSPEC
section, and then the BRANPROP keyword should be used to define the
branch network, and finally the NODEPROP keyword is used to describe the
node properties.

\begin{verbatim}
-- ==============================================================================
--
-- RUNSPEC SECTION
--
-- ==============================================================================
RUNSPEC
--
--       ACTIVATE THE EXTENDED NETWORK OPTION AND DEFINE PARAMETERS
--
--       MAX.    MAX     NOT
--       NODE    LINK    USED
NETWORK
         3        2      1*                                                    /
…..…………..
..……..…..
-- ==============================================================================
--
-- SCHEDULE SECTION
--
-- ==============================================================================
SCHEDULE
--
--       EXTENDED NETWORK BRANCH PROPERTIES
--
-- DOWN  UP       VFP    VFP
-- NODE  NODE     TABLE  ALFQ
BRANPROP
B1       PLAT-A   5      1*                                                   /
C1       PLAT-A   4      1*                                                   /
/

--
--       EXTENDED NETWORK NODE PROPERTIES
--
-- NODE  NODE   CHOKE  GAS   CHOKE  SOURCE  NETWORK
-- NAME  PRESS  OPTN   LIFT  GROUP  SINK    TYPE
NODEPROP
PLAT-A   21.0   NO     NO                                                     /
B1       1*     NO     NO                                                     /
C1       1*     NO     NO                                                     /
/
\end{verbatim}

Here the main platform for the field, PLAT-A, has a fixed 21 barsa
pressure applied as an operating constraint.

\subsection{CALTRAC -- Define a Gas Calorfic Value
Tracer}\label{caltrac-define-a-gas-calorfic-value-tracer}

\begin{longtable}[]{@{}
  >{\raggedright\arraybackslash}p{(\linewidth - 14\tabcolsep) * \real{0.1250}}
  >{\raggedright\arraybackslash}p{(\linewidth - 14\tabcolsep) * \real{0.1250}}
  >{\raggedright\arraybackslash}p{(\linewidth - 14\tabcolsep) * \real{0.1250}}
  >{\raggedright\arraybackslash}p{(\linewidth - 14\tabcolsep) * \real{0.1250}}
  >{\raggedright\arraybackslash}p{(\linewidth - 14\tabcolsep) * \real{0.1250}}
  >{\raggedright\arraybackslash}p{(\linewidth - 14\tabcolsep) * \real{0.1250}}
  >{\raggedright\arraybackslash}p{(\linewidth - 14\tabcolsep) * \real{0.1250}}
  >{\raggedright\arraybackslash}p{(\linewidth - 14\tabcolsep) * \real{0.1250}}@{}}
\toprule\noalign{}
\begin{minipage}[b]{\linewidth}\raggedright
RUNSPEC
\end{minipage} & \begin{minipage}[b]{\linewidth}\raggedright
GRID
\end{minipage} & \begin{minipage}[b]{\linewidth}\raggedright
EDIT
\end{minipage} & \begin{minipage}[b]{\linewidth}\raggedright
PROPS
\end{minipage} & \begin{minipage}[b]{\linewidth}\raggedright
REGIONS
\end{minipage} & \begin{minipage}[b]{\linewidth}\raggedright
SOLUTION
\end{minipage} & \begin{minipage}[b]{\linewidth}\raggedright
SUMMARY
\end{minipage} & \begin{minipage}[b]{\linewidth}\raggedright
SCHEDULE
\end{minipage} \\
\midrule\noalign{}
\endhead
\bottomrule\noalign{}
\endlastfoot
\end{longtable}

\subsubsection{Description}\label{description-874}

The CALTRAC keyword is used to assign a gas calorfic value to a tracer,
for when the Tracer option has been invoked by the TRACER keyword in the
RUNSPEC section.

This keyword is not supported by OPM Flow but would change the results
if supported so the simulation will be stopped.

\subsection{CECON -- Define Well Connections Economic Limit
Criteria}\label{cecon-define-well-connections-economic-limit-criteria}

\begin{longtable}[]{@{}
  >{\raggedright\arraybackslash}p{(\linewidth - 14\tabcolsep) * \real{0.1250}}
  >{\raggedright\arraybackslash}p{(\linewidth - 14\tabcolsep) * \real{0.1250}}
  >{\raggedright\arraybackslash}p{(\linewidth - 14\tabcolsep) * \real{0.1250}}
  >{\raggedright\arraybackslash}p{(\linewidth - 14\tabcolsep) * \real{0.1250}}
  >{\raggedright\arraybackslash}p{(\linewidth - 14\tabcolsep) * \real{0.1250}}
  >{\raggedright\arraybackslash}p{(\linewidth - 14\tabcolsep) * \real{0.1250}}
  >{\raggedright\arraybackslash}p{(\linewidth - 14\tabcolsep) * \real{0.1250}}
  >{\raggedright\arraybackslash}p{(\linewidth - 14\tabcolsep) * \real{0.1250}}@{}}
\toprule\noalign{}
\begin{minipage}[b]{\linewidth}\raggedright
RUNSPEC
\end{minipage} & \begin{minipage}[b]{\linewidth}\raggedright
GRID
\end{minipage} & \begin{minipage}[b]{\linewidth}\raggedright
EDIT
\end{minipage} & \begin{minipage}[b]{\linewidth}\raggedright
PROPS
\end{minipage} & \begin{minipage}[b]{\linewidth}\raggedright
REGIONS
\end{minipage} & \begin{minipage}[b]{\linewidth}\raggedright
SOLUTION
\end{minipage} & \begin{minipage}[b]{\linewidth}\raggedright
SUMMARY
\end{minipage} & \begin{minipage}[b]{\linewidth}\raggedright
SCHEDULE
\end{minipage} \\
\midrule\noalign{}
\endhead
\bottomrule\noalign{}
\endlastfoot
\end{longtable}

\subsubsection{Description}\label{description-875}

CECON sets the economic cut-off criteria for a well's connection to the
simulation grid.

This keyword is not supported by OPM Flow but would change the results
if supported so the simulation will be stopped.

\subsection{CECONT -- Define Well Connections Tracer Economic Limit
Criteria}\label{cecont-define-well-connections-tracer-economic-limit-criteria}

\begin{longtable}[]{@{}
  >{\raggedright\arraybackslash}p{(\linewidth - 14\tabcolsep) * \real{0.1250}}
  >{\raggedright\arraybackslash}p{(\linewidth - 14\tabcolsep) * \real{0.1250}}
  >{\raggedright\arraybackslash}p{(\linewidth - 14\tabcolsep) * \real{0.1250}}
  >{\raggedright\arraybackslash}p{(\linewidth - 14\tabcolsep) * \real{0.1250}}
  >{\raggedright\arraybackslash}p{(\linewidth - 14\tabcolsep) * \real{0.1250}}
  >{\raggedright\arraybackslash}p{(\linewidth - 14\tabcolsep) * \real{0.1250}}
  >{\raggedright\arraybackslash}p{(\linewidth - 14\tabcolsep) * \real{0.1250}}
  >{\raggedright\arraybackslash}p{(\linewidth - 14\tabcolsep) * \real{0.1250}}@{}}
\toprule\noalign{}
\begin{minipage}[b]{\linewidth}\raggedright
RUNSPEC
\end{minipage} & \begin{minipage}[b]{\linewidth}\raggedright
GRID
\end{minipage} & \begin{minipage}[b]{\linewidth}\raggedright
EDIT
\end{minipage} & \begin{minipage}[b]{\linewidth}\raggedright
PROPS
\end{minipage} & \begin{minipage}[b]{\linewidth}\raggedright
REGIONS
\end{minipage} & \begin{minipage}[b]{\linewidth}\raggedright
SOLUTION
\end{minipage} & \begin{minipage}[b]{\linewidth}\raggedright
SUMMARY
\end{minipage} & \begin{minipage}[b]{\linewidth}\raggedright
SCHEDULE
\end{minipage} \\
\midrule\noalign{}
\endhead
\bottomrule\noalign{}
\endlastfoot
\end{longtable}

\subsubsection{Description}\label{description-876}

This keyword, CECONT, sets the tracer economic cut-off criteria for a
well's connection to the simulation grid.

This keyword is not supported by OPM Flow but would change the results
if supported so the simulation will be stopped.

\subsection{COLUMNS -- Define Input File Column
Margins}\label{columns-define-input-file-column-margins-8}

The COLUMNS keyword defines the input file column margins; characters
outside the margins are ignored by the input parser.

See COLUMNS -- Define Input File Column Margins in the GLOBAL section
for a full description.

\subsection{COMPDAT -- Define Well Connections to the
Grid}\label{compdat-define-well-connections-to-the-grid}

\begin{longtable}[]{@{}
  >{\raggedright\arraybackslash}p{(\linewidth - 14\tabcolsep) * \real{0.1250}}
  >{\raggedright\arraybackslash}p{(\linewidth - 14\tabcolsep) * \real{0.1250}}
  >{\raggedright\arraybackslash}p{(\linewidth - 14\tabcolsep) * \real{0.1250}}
  >{\raggedright\arraybackslash}p{(\linewidth - 14\tabcolsep) * \real{0.1250}}
  >{\raggedright\arraybackslash}p{(\linewidth - 14\tabcolsep) * \real{0.1250}}
  >{\raggedright\arraybackslash}p{(\linewidth - 14\tabcolsep) * \real{0.1250}}
  >{\raggedright\arraybackslash}p{(\linewidth - 14\tabcolsep) * \real{0.1250}}
  >{\raggedright\arraybackslash}p{(\linewidth - 14\tabcolsep) * \real{0.1250}}@{}}
\toprule\noalign{}
\begin{minipage}[b]{\linewidth}\raggedright
RUNSPEC
\end{minipage} & \begin{minipage}[b]{\linewidth}\raggedright
GRID
\end{minipage} & \begin{minipage}[b]{\linewidth}\raggedright
EDIT
\end{minipage} & \begin{minipage}[b]{\linewidth}\raggedright
PROPS
\end{minipage} & \begin{minipage}[b]{\linewidth}\raggedright
REGIONS
\end{minipage} & \begin{minipage}[b]{\linewidth}\raggedright
SOLUTION
\end{minipage} & \begin{minipage}[b]{\linewidth}\raggedright
SUMMARY
\end{minipage} & \begin{minipage}[b]{\linewidth}\raggedright
SCHEDULE
\end{minipage} \\
\midrule\noalign{}
\endhead
\bottomrule\noalign{}
\endlastfoot
\end{longtable}

\subsubsection{Description}\label{description-877}

The COMPDAT keyword defines how a well is connected to the reservoir by
defining or modifying existing well connections. Ideally the connections
should be declared in the correct sequence, starting with the connection
nearest the well head and then working along the wellbore towards the
bottom or toe of the well, however this may not be possible or
convenient, for example when connections are added or removed from a
well during the simulation (see COMPORD in the SCHEDULE section for
options regarding connection ordering).

\begin{longtable}[]{@{}
  >{\raggedright\arraybackslash}p{(\linewidth - 6\tabcolsep) * \real{0.2500}}
  >{\raggedright\arraybackslash}p{(\linewidth - 6\tabcolsep) * \real{0.2500}}
  >{\raggedright\arraybackslash}p{(\linewidth - 6\tabcolsep) * \real{0.2500}}
  >{\raggedright\arraybackslash}p{(\linewidth - 6\tabcolsep) * \real{0.2500}}@{}}
\toprule\noalign{}
\begin{minipage}[b]{\linewidth}\raggedright
No.
\end{minipage} & \begin{minipage}[b]{\linewidth}\raggedright
Name
\end{minipage} & \begin{minipage}[b]{\linewidth}\raggedright
Description
\end{minipage} & \begin{minipage}[b]{\linewidth}\raggedright
Default
\end{minipage} \\
\midrule\noalign{}
\endhead
\bottomrule\noalign{}
\endlastfoot
Field & Metric & Laboratory & \\
1 & WELNAME & A character string of up to eight characters in length
that defines the well name for which the well connection data are being
defined. Note that the well name (WELNAME) must have been declared
previously using the WELSPECS keyword in the SCHEDULE section, otherwise
an error may occur. & None \\
2 & I & A positive integer greater than or equal to zero and less than
or equal to NX that defines the connection location in the I-direction.
If set to zero or defaulted with 1* the location is taken from the
wellhead location I-direction value on the WELSPECS keyword in the
SCHEDULE section. & 0 \\
3 & J & A positive integer greater than or equal to zero and less than
or equal to NY that defines the connection location in the J-direction.
If set to zero or defaulted with 1* the location is taken from the
wellhead location J-direction value on the WELSPECS keyword in the
SCHEDULE section. & 0 \\
4 & K1 & A positive integer greater than or equal to one and less than
or equal to NZ that defines the UPPER connection location in the
K-direction. & None \\
5 & K2 & A positive integer greater than or equal to K1 and less than or
equal to NZ that defines the LOWER connection location in the
K-direction. & None \\
6 & STATUS & A defined character string of length four that defines the
connections' operational status, STATUS should be set to one of the
following character strings: & OPEN \\
7 & SATNUM & An integer greater than or equal to zero and less than
NTSFUN as declared on the TABDIMS keyword in the RUNSPEC section, that
defines the saturation table number to be used for flow between the
reservoir grid block and the well connections. If SATNUM is set to zero
or defaulted with 1* then: & 0 \\
8 & CONFACT & A real value greater than or equal to zero that defines
the transmissibility connection factor between the well bore and the
reservoir grid block. If set to zero or defaulted with 1* then items (9)
through (13) are used to calculate CONFACT. & Defined \\
cP.rb/day/psia 0 & cP.rm3/day/bars 0 & cP.rcc/hr/atm 0 & \\
9 & DW & A real positive value that defines the well bore diameter of
the connections for the well. DW is used in calculating a well's
productivity or injectivity index; however the value will be ignored in
calculating the connections CONFACT value if CONFACT has been directly
entered. & None \\
feet & m & cm & \\
10 & KH & A real value that defines the effective KH (permeability x
length) for the connections. If less than or equal to zero, or defaulted
by 1*, then KH is calculated from the connected grid blocks. KH is
ignored if CONFAC has been directly entered. & Calculated from connected
grid blocks \\
mD.ft & mD.m & mD.cm & \\
11 & SKIN & A real value that defines the connections dimensionless skin
factor. SKIN is used in calculating a well's productivity or injectivity
index; however, the value will be ignored in calculating the connections
CONFACT value if CONFAC has been directly entered. & 0.0 \\
dimensionless & dimensionless & dimensionless & \\
12 & DFACT & A real value that defines the non-Darcy D factor
coefficient for gas wells. It is recommended that the value should be
defaulted with 1* and the non-Darcy D factor coefficient for gas wells
defined via the WDFAC keyword in the SCHEDULE section. If the value is
greater than or equal to zero then the value is used to define the well
D factor and the simulator evaluates the connection D factors
internally. If the value is less than zero then its sign is reversed and
then used directly to define the connection D factors. In this case any
D factor values previously defined using the WDFAC keyword for other
connections will retained, whereas, any connection D factors previously
defined for the entire well using the WDFACCOR keyword will be
discarded. & 0.0 \\
day/Mscf & day/m3 & hour/sc & \\
13 & DIRECT & A defined character string of length one that defines the
orientation of the connections and should be set to either X, Y, or Z.
The direction of connections also determines the length of the
connection used to calculate the connection factor if CONFACT has not
been entered directly. The default value is for a vertical connection,
that is DIRECT is defaulted to Z. & Z \\
Notes: & & & \\
\end{longtable}

\emph{Table 12.3.23.1: COMPDAT Keyword Description}

Multiple grid block connections can be defined on one record for
vertical wells by assigning different values to K1 and K2, for deviated
and horizontal wells this may not be possible and therefore each grid
block connection must be separately defined by using one record per
connection, setting K1 equal to K2 in each record.

See also the WELSPECS keyword to define wells, the COMPIMB to reset the
imbibition relative permeability table allocation, and the COMPORD to
re-order the completions along the well trajectory. In addition, the
COMPLUMP keyword groups well connections together to form well
completions for a well. All the aforementioned keywords are described in
the SCHEDULE section.

\begin{tcolorbox}[enhanced jigsaw, title=\textcolor{quarto-callout-note-color}{\faInfo}\hspace{0.5em}{Note}, coltitle=black, breakable, toprule=.15mm, colbacktitle=quarto-callout-note-color!10!white, leftrule=.75mm, opacitybacktitle=0.6, colback=white, bottomtitle=1mm, toptitle=1mm, titlerule=0mm, left=2mm, colframe=quarto-callout-note-color-frame, arc=.35mm, rightrule=.15mm, bottomrule=.15mm, opacityback=0]

The term well connection is used to describe individual connections from
the wellbore to the reservoir grid, as opposed to well completions. A
well completion is used to describe a set of connections, for example, a
well may consist of several completions with each completion consisting
of multiple connections.

\end{tcolorbox}

\subsubsection{Example}\label{example-505}

The following example defines two vertical oil wells using the WELSPECS
keyword and their associated connection data.

\begin{verbatim}
--
--       WELL SPECIFICATION DATA
--
-- WELL  GROUP     LOCATION  BHP    PHASE  DRAIN  INFLOW  OPEN  CROSS  PRESS
-- NAME  NAME        I    J  DEPTH  FLUID  AREA   EQUANS  SHUT  FLOW   TABLE       WELSPECS
OP01     PLATFORM   14   13   1*     OIL    1*     STD    SHUT   NO     1*  /
OP02     PLATFORM   35   96   1*     OIL    1*     STD    SHUT   NO     1*  /
/
--
--       WELL CONNECTION DATA
--
-- WELL  --- LOCATION ---  OPEN   SAT   CONN   WELL   KH    SKIN   D     DIR
-- NAME   II  JJ  K1  K2   SHUT   TAB   FACT   DIA    FACT  FACT   FACT  PEN
COMPDAT
OP01      1*  1*  20  56   OPEN   1*    1*    0.708   1*    0.0    1*    'Z' /
OP01      1*  1*  75 100   SHUT   1*    1*    0.708   1*    0.0    1*    'Z' /
OP02      35  96  75 100   OPEN   1*    1*    0.708   1*    0.0    1*    'Z' /
\end{verbatim}

Well OP01 has two sets of connections; the first one connects grid cells
(14, 13, 20) to (14, 13, 56) to the well and is open to flow and the
second connecting grid cells (14, 13, 75) to (14, 13, 100) is shut. Well
OP02 has only one open set of connections from cells (35, 96, 75) to
cells (35, 96, 100).

\subsection{COMPDATL -- Define Well Connections to a LGR
Grid}\label{compdatl-define-well-connections-to-a-lgr-grid}

\begin{longtable}[]{@{}
  >{\raggedright\arraybackslash}p{(\linewidth - 14\tabcolsep) * \real{0.1250}}
  >{\raggedright\arraybackslash}p{(\linewidth - 14\tabcolsep) * \real{0.1250}}
  >{\raggedright\arraybackslash}p{(\linewidth - 14\tabcolsep) * \real{0.1250}}
  >{\raggedright\arraybackslash}p{(\linewidth - 14\tabcolsep) * \real{0.1250}}
  >{\raggedright\arraybackslash}p{(\linewidth - 14\tabcolsep) * \real{0.1250}}
  >{\raggedright\arraybackslash}p{(\linewidth - 14\tabcolsep) * \real{0.1250}}
  >{\raggedright\arraybackslash}p{(\linewidth - 14\tabcolsep) * \real{0.1250}}
  >{\raggedright\arraybackslash}p{(\linewidth - 14\tabcolsep) * \real{0.1250}}@{}}
\toprule\noalign{}
\begin{minipage}[b]{\linewidth}\raggedright
RUNSPEC
\end{minipage} & \begin{minipage}[b]{\linewidth}\raggedright
GRID
\end{minipage} & \begin{minipage}[b]{\linewidth}\raggedright
EDIT
\end{minipage} & \begin{minipage}[b]{\linewidth}\raggedright
PROPS
\end{minipage} & \begin{minipage}[b]{\linewidth}\raggedright
REGIONS
\end{minipage} & \begin{minipage}[b]{\linewidth}\raggedright
SOLUTION
\end{minipage} & \begin{minipage}[b]{\linewidth}\raggedright
SUMMARY
\end{minipage} & \begin{minipage}[b]{\linewidth}\raggedright
SCHEDULE
\end{minipage} \\
\midrule\noalign{}
\endhead
\bottomrule\noalign{}
\endlastfoot
\end{longtable}

\subsubsection{Description}\label{description-878}

The COMPDATL keyword defines how a well in a Local Grid Refinement
(``LGR'') is connected to the reservoir by declaring the LGR and
defining or modifying existing well connections. Ideally the connections
should be declared in the correct sequence, starting with the connection
nearest the well head and then working along the wellbore towards the
bottom or toe of the well, however this may not be possible or
convenient, for example when connections are added or removed from a
well during the simulation (see the COMPORD keyword in the SCHEDULE
section for options regarding connection ordering).

This keyword is not supported by OPM Flow but would change the results
if supported so the simulation will be stopped.

\begin{longtable}[]{@{}
  >{\raggedright\arraybackslash}p{(\linewidth - 6\tabcolsep) * \real{0.2500}}
  >{\raggedright\arraybackslash}p{(\linewidth - 6\tabcolsep) * \real{0.2500}}
  >{\raggedright\arraybackslash}p{(\linewidth - 6\tabcolsep) * \real{0.2500}}
  >{\raggedright\arraybackslash}p{(\linewidth - 6\tabcolsep) * \real{0.2500}}@{}}
\toprule\noalign{}
\begin{minipage}[b]{\linewidth}\raggedright
No.
\end{minipage} & \begin{minipage}[b]{\linewidth}\raggedright
Name
\end{minipage} & \begin{minipage}[b]{\linewidth}\raggedright
Description
\end{minipage} & \begin{minipage}[b]{\linewidth}\raggedright
Default
\end{minipage} \\
\midrule\noalign{}
\endhead
\bottomrule\noalign{}
\endlastfoot
Field & Metric & Laboratory & \\
1 & WELNAME & A character string of up to eight characters in length
that defines the well name for which the well connection data are being
defined. Note that the well name (WELNAME) must have been declared
previously using the WELSPECS keyword in the SCHEDULE section, otherwise
an error may occur. & None \\
2 & LGRNAME & A character string of up to eight characters in length
that defines the LGR name for which the well LGR connection data are
being defined. Note that the well name (LGRNAME) must have been declared
previously using the WELSPECL keyword in the SCHEDULE section, otherwise
an error may occur. If defaulted with 1* the LGR on the WELSPECL keyword
will be utilized. & Defined \\
3 & I & A positive integer greater than or equal to zero and less than
or equal to NX that defines the connection location in the I-direction.
If set to zero or defaulted with 1* the location is taken from the
wellhead location I-direction value on the WELSPECS keyword in the
SCHEDULE section. & 0 \\
4 & J & A positive integer greater than or equal to zero and less than
or equal to NY that defines the connection location in the J-direction.
If set to zero or defaulted with 1* the location is taken from the
wellhead location J-direction value on the WELSPECS keyword in the
SCHEDULE section. & 0 \\
5 & K1 & A positive integer greater than or equal to one and less than
or equal to NZ that defines the UPPER connection location in the
K-direction. & None \\
6 & K2 & A positive integer greater than or equal to K1 and less than or
equal to NZ that defines the LOWER connection location in the
K-direction. & None \\
7 & STATUS & A character string of length four that defines the
connections' operational status, STATUS should be set to one of the
following character strings: & OPEN \\
8 & SATNUM & An integer greater than or equal to zero and less than
NTSFUN as declared on the TABDIMS keyword in the RUNSPEC, that defines
the saturation table number to be used for flow between the reservoir
grid block and the well connections. If SATNUM is set to zero or
defaulted with 1* then: & 0 \\
9 & CONFACT & A real value greater than or equal to zero that defines
the transmissibility connection factor between the well bore and the
reservoir grid block. If set to zero or defaulted with 1* then items (9)
through (13) are used to calculate CONFACT. & Defined \\
cP.rb/day/psia 0 & cP.rm3/day/bars 0 & cP.rcc/hr/atm 0 & \\
10 & DW & A real positive value that defines the well bore diameter of
the connections for the well. DW is used in calculating a well's
productivity or injectivity index; however the value will be ignored in
calculating the connections CONFACT value if CONFAC has been directly
entered. & None \\
feet & m & cm & \\
11 & KH & A real value that defines the effective KH (permeability x
length) for the connections. If less than or equal to zero or defaulted
by 1* then KH is calculated from the connected grid blocks. KH is
ignored if CONFAC has been directly entered. & Calculated from connected
grid blocks \\
mD.ft & mD.m & mD.cm & \\
12 & SKIN & A real value that defines the connections dimensionless skin
factor. SKIN is used in calculating a well's productivity or injectivity
index; however, the value will be ignored in calculating the connections
CONFACT value if CONFAC has been directly entered. & 0.0 \\
dimensionless & dimensionless & dimensionless & \\
13 & DFACT & A real value that defines the non-Darcy D factor
coefficient for gas wells. This value should be defaulted with 1* and
the non-Darcy D factor coefficient for gas wells defined via the WDFAC
keyword in the SCHEDULE section. Currently this option is not supported
by OPM Flow. & 1* \\
day/Mscf & day/m3 & hour/sc & \\
14 & DIRECT & A one letter character string that defines the orientation
of the connections and should be set to either X, Y, or Z. The direction
of connections also determines the length of the connection used to
calculate the connection factor if CONFAC has not been entered directly.
The default value is for a vertical connection, that is DIRECT is
defaulted to Z. & Z \\
Notes: & & & \\
\end{longtable}

\emph{Table 12.11: COMPDATL Keyword Description}

Multiple grid block connections can be defined on one record for
vertical wells by assigning different values to K1 and K2, for deviated
and horizontal wells this may not be possible and therefore each grid
block connection must be separately defined by using one record per
connection, setting K1 equal to K2 in each record.

See also the WELSPECS keyword to define wells, the COMPIMB to reset the
imbibition relative permeability table allocation, and the COMPORD to
re-order the completions along the well trajectory. In addition, the
COMPLUMP keyword groups well connections together to form well
completions for a well. All the aforementioned keywords are described in
the SCHEDULE section.

\begin{tcolorbox}[enhanced jigsaw, title=\textcolor{quarto-callout-note-color}{\faInfo}\hspace{0.5em}{Note}, coltitle=black, breakable, toprule=.15mm, colbacktitle=quarto-callout-note-color!10!white, leftrule=.75mm, opacitybacktitle=0.6, colback=white, bottomtitle=1mm, toptitle=1mm, titlerule=0mm, left=2mm, colframe=quarto-callout-note-color-frame, arc=.35mm, rightrule=.15mm, bottomrule=.15mm, opacityback=0]

The term well connection is used to describe individual connections from
the wellbore to the reservoir grid, as opposed to well completions. A
well completion is used to describe a set of connections, for example, a
well may consist of several completions with each completion consisting
of multiple connections.

\end{tcolorbox}

\subsubsection{Example}\label{example-506}

The following example defines two vertical oil wells using the WELSPECS
keyword and their associated connection data.

\begin{verbatim}
--
--       WELL LGR SPECIFICATION DATA
--
--       WELL GROUP  LGR    -LOCATION- BHP    PHASE DRAIN INFLOW SHUT CROSS PVT
--       NAME NAME   NAME     I    J   DEPTH  FLUID AREA  EQUA.  IN   FLOW  TABLE
WELSPECL
         OP01  PLAT OP01LGR  14   13   1*     OIL   1*    STD    SHUT  NO   1*  /
         OP02  PLAT OP02LGR  28   96   1*     OIL   1*    STD    SHUT  NO   1*  /
/
--
--       WELL LGR CONNECTION DATA
--
--       WELL  LGR   ---LOCATION---   OPEN   SAT CONN  WELL KH   SKIN  D    DIR
--       NAME  NAME   II  JJ  K1  K2  SHUT   TAB FACT  DIA  FACT FACT  FACT PEN
COMPDATL
         OP01 OP01LGR 1*  1*  20  56  OPEN   1*  1*  0.708   1*   1*   1*   Z /
         OP01 OP01LGR 1*  1*  75 100  SHUT   1*  1*  0.708   1*   1*   1*   Z /                  OP02 OP02LGR 35  96  75 100  OPEN   1*  1*  0.708   1*   1*   1*   Z /
/
\end{verbatim}

Well OP01 has two sets of connections; the first one connects grid cells
(14, 13, 20) to (14, 13, 56) to the well and is open to flow and the
second connecting grid cells (14, 13, 75) to (14, 13, 100) is shut. Well
OP02 has only one open connection from cells (35, 96, 75) to cells (35,
96, 100).

\subsection{COMPDATM -- Define Well Connections to an Amalgamated LGR
Grid}\label{compdatm-define-well-connections-to-an-amalgamated-lgr-grid}

\begin{longtable}[]{@{}
  >{\raggedright\arraybackslash}p{(\linewidth - 14\tabcolsep) * \real{0.1250}}
  >{\raggedright\arraybackslash}p{(\linewidth - 14\tabcolsep) * \real{0.1250}}
  >{\raggedright\arraybackslash}p{(\linewidth - 14\tabcolsep) * \real{0.1250}}
  >{\raggedright\arraybackslash}p{(\linewidth - 14\tabcolsep) * \real{0.1250}}
  >{\raggedright\arraybackslash}p{(\linewidth - 14\tabcolsep) * \real{0.1250}}
  >{\raggedright\arraybackslash}p{(\linewidth - 14\tabcolsep) * \real{0.1250}}
  >{\raggedright\arraybackslash}p{(\linewidth - 14\tabcolsep) * \real{0.1250}}
  >{\raggedright\arraybackslash}p{(\linewidth - 14\tabcolsep) * \real{0.1250}}@{}}
\toprule\noalign{}
\begin{minipage}[b]{\linewidth}\raggedright
RUNSPEC
\end{minipage} & \begin{minipage}[b]{\linewidth}\raggedright
GRID
\end{minipage} & \begin{minipage}[b]{\linewidth}\raggedright
EDIT
\end{minipage} & \begin{minipage}[b]{\linewidth}\raggedright
PROPS
\end{minipage} & \begin{minipage}[b]{\linewidth}\raggedright
REGIONS
\end{minipage} & \begin{minipage}[b]{\linewidth}\raggedright
SOLUTION
\end{minipage} & \begin{minipage}[b]{\linewidth}\raggedright
SUMMARY
\end{minipage} & \begin{minipage}[b]{\linewidth}\raggedright
SCHEDULE
\end{minipage} \\
\midrule\noalign{}
\endhead
\bottomrule\noalign{}
\endlastfoot
\end{longtable}

\subsubsection{Description}\label{description-879}

The COMPDATM keyword is an alias for the COMPDATL keyword. COMPDATM
defines how a well in an amalgamated Local Grid Refinement (``LGR'') is
connected to the reservoir by declaring the LGR and defining or
modifying existing well connections. Ideally the connections should be
declared in the correct sequence, starting with the connection nearest
the well head and then working along the wellbore towards the bottom or
toe of the well, however this may not be possible or convenient, for
example when connections are added or removed from a well during the
simulation (see the COMPORD keyword in the SCHEDULE section for options
regarding connection ordering).

This keyword is not supported by OPM Flow but would change the results
if supported so the simulation will be stopped.

\subsection{COMPFLSH -- Assign COMPDAT Differential-Flash Liberation
Ratios}\label{compflsh-assign-compdat-differential-flash-liberation-ratios}

\begin{longtable}[]{@{}
  >{\raggedright\arraybackslash}p{(\linewidth - 14\tabcolsep) * \real{0.1250}}
  >{\raggedright\arraybackslash}p{(\linewidth - 14\tabcolsep) * \real{0.1250}}
  >{\raggedright\arraybackslash}p{(\linewidth - 14\tabcolsep) * \real{0.1250}}
  >{\raggedright\arraybackslash}p{(\linewidth - 14\tabcolsep) * \real{0.1250}}
  >{\raggedright\arraybackslash}p{(\linewidth - 14\tabcolsep) * \real{0.1250}}
  >{\raggedright\arraybackslash}p{(\linewidth - 14\tabcolsep) * \real{0.1250}}
  >{\raggedright\arraybackslash}p{(\linewidth - 14\tabcolsep) * \real{0.1250}}
  >{\raggedright\arraybackslash}p{(\linewidth - 14\tabcolsep) * \real{0.1250}}@{}}
\toprule\noalign{}
\begin{minipage}[b]{\linewidth}\raggedright
RUNSPEC
\end{minipage} & \begin{minipage}[b]{\linewidth}\raggedright
GRID
\end{minipage} & \begin{minipage}[b]{\linewidth}\raggedright
EDIT
\end{minipage} & \begin{minipage}[b]{\linewidth}\raggedright
PROPS
\end{minipage} & \begin{minipage}[b]{\linewidth}\raggedright
REGIONS
\end{minipage} & \begin{minipage}[b]{\linewidth}\raggedright
SOLUTION
\end{minipage} & \begin{minipage}[b]{\linewidth}\raggedright
SUMMARY
\end{minipage} & \begin{minipage}[b]{\linewidth}\raggedright
SCHEDULE
\end{minipage} \\
\midrule\noalign{}
\endhead
\bottomrule\noalign{}
\endlastfoot
\end{longtable}

\subsubsection{Description}\label{description-880}

COMPFLSH is used to assign saturated PVT differential-flash liberation
ratios to individual well completions for both the oil formation volume
factor and the gas-oil ratio. Only the saturated oil properties are
changed by this keyword, that is only data entered via the PVCO or PVTO
keywords in the PROPS section will be used in the calculation. The other
fluid PVT property keywords: PVTW, PVTG, DENSITY, etc., are not used.

The differential liberation expansion (``DLE'') experiment is designed
to approximate the depletion process of an oil reservoir, and thereby
provide suitable PVT data for calculating reservoir performance. This is
then coupled with surface faculties via the multi-stage separator flash
data. The multi-stage separator test is performed on an oil sample
primarily to provide a basis for converting differential liberation data
from a residual oil to a stock-tank oil basis. Occasionally, several
separator tests are run to help choose separator conditions that
maximize stock-tank oil production. Usually two or three stages of
separation are used, with the last stage being at atmospheric pressure
and near-ambient temperature (60 to 80 oF).

Note, normally the data on the PVCO and PVTO keywords should have
already taken the two datasets into account and therefore this keyword
should not normally be used in practice, unless there is a specific
requirement to adjust the oil PVT data on a well connection basis.

This keyword is not supported by OPM Flow but would change the results
if supported so the simulation will be stopped.

\subsection{COMPIMB -- Assign Imbibition Saturation Tables to Well
Connections}\label{compimb-assign-imbibition-saturation-tables-to-well-connections}

\begin{longtable}[]{@{}
  >{\raggedright\arraybackslash}p{(\linewidth - 14\tabcolsep) * \real{0.1250}}
  >{\raggedright\arraybackslash}p{(\linewidth - 14\tabcolsep) * \real{0.1250}}
  >{\raggedright\arraybackslash}p{(\linewidth - 14\tabcolsep) * \real{0.1250}}
  >{\raggedright\arraybackslash}p{(\linewidth - 14\tabcolsep) * \real{0.1250}}
  >{\raggedright\arraybackslash}p{(\linewidth - 14\tabcolsep) * \real{0.1250}}
  >{\raggedright\arraybackslash}p{(\linewidth - 14\tabcolsep) * \real{0.1250}}
  >{\raggedright\arraybackslash}p{(\linewidth - 14\tabcolsep) * \real{0.1250}}
  >{\raggedright\arraybackslash}p{(\linewidth - 14\tabcolsep) * \real{0.1250}}@{}}
\toprule\noalign{}
\begin{minipage}[b]{\linewidth}\raggedright
RUNSPEC
\end{minipage} & \begin{minipage}[b]{\linewidth}\raggedright
GRID
\end{minipage} & \begin{minipage}[b]{\linewidth}\raggedright
EDIT
\end{minipage} & \begin{minipage}[b]{\linewidth}\raggedright
PROPS
\end{minipage} & \begin{minipage}[b]{\linewidth}\raggedright
REGIONS
\end{minipage} & \begin{minipage}[b]{\linewidth}\raggedright
SOLUTION
\end{minipage} & \begin{minipage}[b]{\linewidth}\raggedright
SUMMARY
\end{minipage} & \begin{minipage}[b]{\linewidth}\raggedright
SCHEDULE
\end{minipage} \\
\midrule\noalign{}
\endhead
\bottomrule\noalign{}
\endlastfoot
\end{longtable}

\subsubsection{Description}\label{description-881}

The COMPIMB keyword assigns imbibition saturation tables to well
connections. The COMPDAT keyword in the SCHEDULE section also assigns
imbibition saturation tables to connections, but in this case the table
number is the same as for the drainage curve. If this is not the
required assignment then the COMPIMB keyword can be used to reset the
imbibition saturation table number. For this to be effective the COMPIMB
keyword must precede the COMPDAT keyword, otherwise it will have no
effect.

The COMPIMB keyword should only be used if the hysteresis option has
been activated via the HYSTER variable on the SATOPTS keyword in the
RUNSPEC section.

This keyword is not supported by OPM Flow but would change the results
if supported so the simulation will be stopped.

\begin{longtable}[]{@{}
  >{\raggedright\arraybackslash}p{(\linewidth - 6\tabcolsep) * \real{0.2500}}
  >{\raggedright\arraybackslash}p{(\linewidth - 6\tabcolsep) * \real{0.2500}}
  >{\raggedright\arraybackslash}p{(\linewidth - 6\tabcolsep) * \real{0.2500}}
  >{\raggedright\arraybackslash}p{(\linewidth - 6\tabcolsep) * \real{0.2500}}@{}}
\toprule\noalign{}
\begin{minipage}[b]{\linewidth}\raggedright
No.
\end{minipage} & \begin{minipage}[b]{\linewidth}\raggedright
Name
\end{minipage} & \begin{minipage}[b]{\linewidth}\raggedright
Description
\end{minipage} & \begin{minipage}[b]{\linewidth}\raggedright
Default
\end{minipage} \\
\midrule\noalign{}
\endhead
\bottomrule\noalign{}
\endlastfoot
Field & Metric & Laboratory & \\
1 & WELNAME & A character string of up to eight characters in length
that defines the well name for which the well connection data is being
defined. Note that the well name (WELNAME) must have been declared
previously using the WELSPECS keyword in the SCHEDULE section, otherwise
an error may occur. & None \\
2 & I & A positive integer greater than or equal to zero and less than
or equal to NX that defines the connection location in the I-direction.
If set to zero or defaulted with 1* then all connections in the
I-direction that also satisfy J, K1 and K2 criteria are assigned the
IMBNUM imbibition table number. & 0 \\
3 & J & A positive integer greater than or equal to zero and less than
or equal to NY that defines the connection location in the J-direction.
If set to zero or defaulted with 1* then all connections in the
J-direction that also satisfy I, K1 and K2 criteria are assigned the
IMBNUM imbibition table number. & 0 \\
4 & K1 & A positive integer greater than or equal to one and less than
or equal to NZ that defines the UPPER connection location in the
K-direction. If set to zero or defaulted with 1* then the upper most
connection in the well is used. & 0 \\
5 & K2 & A positive integer greater than or equal to K1 and less than or
equal to NZ that defines the LOWER connection location in the
K-direction. If set to zero or defaulted with 1* then the lowest most
connection in the well is used. & 0 \\
6 & IMBNUM & An integer greater than or equal to zero and less than
NTSFUN as declared on the TABDIMS keyword in the RUNSPEC, that defines
the imbibition saturation table number to be used for flow between the
reservoir grid block and the well connections. If IMBNUM is set to zero
or defaulted with 1* then the inhibition saturation table allocated to
the grid block that the connections are located within is used. If I, J,
K1, K2 are all set to zero or defaulted to 1*, then IMBNUM is allocated
to all connections in the well. & 0 \\
Notes: & & & \\
\end{longtable}

\emph{Table 12.12: COMPIMB Keyword Description}

Multiple grid block connections can be defined on one record for
vertical wells by assigning different values to K1 and K2, for deviated
and horizontal wells this may not be possible and therefore each grid
block connection must be separately defines by setting K1 equal to K2.

See also the COMPDAT keyword in the SCHEDULE section.

\subsubsection{Example}\label{example-507}

The following example defines the connections for two vertical oil wells
using the COMPDAT keyword and then re-sets the imbibition saturation
functions using the COMPIMP keyword.

\begin{verbatim}
--
--       WELL CONNECTION DATA
--
-- WELL  --- LOCATION ---  OPEN   SAT   CONN   WELL   KH    SKIN   D     DIR
-- NAME   II  JJ  K1  K2   SHUT   TAB   FACT   DIA    FACT  FACT   FACT  PEN
COMPDAT
OP01      1*  1*  20  56   OPEN   1     1*    0.708   1*    0.0    1*    'Z' /
OP01      1*  1*  75 100   SHUT   2     1*    0.708   1*    0.0    1*    'Z' /
OP02      35  96  75 100   OPEN   1     1*    0.708   1*    0.0    1*    'Z' /
--
-- ASSIGN IMBIBITION SATURATION TABLES TO CONNECTIONS
--
-- WELL     ---LOCATION---   SAT
-- NAME     II  JJ  K1  K2   TAB
COMPIMB
OP01        1*  1*  20  56   11                      /
OP01        1*  1*  75 100   12                      /
OP02        1*  1*  1*  1*   11                      /
/
\end{verbatim}

Well OP01 has two sets of COMPIMB records to overwrite the imbibition
saturation tables, one for connections (14, 13, 20) to (14, 13, 56)
resetting the imbibition saturation table number from one to 11 and one
for connections (14, 13, 75) to (14, 13, 100) that resets the imbibition
table number from 2 to 12. Well OP02 has only one connection from cells
(35, 96, 75) to cells (35, 96, 100), so all the default values for I, J,
K1, and K2 can be used to set the imbibition table numbers from 2 to 11.
Note in all cases the drainage saturation table retains the value as
specified by the COMPDAT keyword, that is one, two and one.

\subsection{COMPINJK -- Assign Injection well Relative Permeability
Values}\label{compinjk-assign-injection-well-relative-permeability-values}

\begin{longtable}[]{@{}
  >{\raggedright\arraybackslash}p{(\linewidth - 14\tabcolsep) * \real{0.1250}}
  >{\raggedright\arraybackslash}p{(\linewidth - 14\tabcolsep) * \real{0.1250}}
  >{\raggedright\arraybackslash}p{(\linewidth - 14\tabcolsep) * \real{0.1250}}
  >{\raggedright\arraybackslash}p{(\linewidth - 14\tabcolsep) * \real{0.1250}}
  >{\raggedright\arraybackslash}p{(\linewidth - 14\tabcolsep) * \real{0.1250}}
  >{\raggedright\arraybackslash}p{(\linewidth - 14\tabcolsep) * \real{0.1250}}
  >{\raggedright\arraybackslash}p{(\linewidth - 14\tabcolsep) * \real{0.1250}}
  >{\raggedright\arraybackslash}p{(\linewidth - 14\tabcolsep) * \real{0.1250}}@{}}
\toprule\noalign{}
\begin{minipage}[b]{\linewidth}\raggedright
RUNSPEC
\end{minipage} & \begin{minipage}[b]{\linewidth}\raggedright
GRID
\end{minipage} & \begin{minipage}[b]{\linewidth}\raggedright
EDIT
\end{minipage} & \begin{minipage}[b]{\linewidth}\raggedright
PROPS
\end{minipage} & \begin{minipage}[b]{\linewidth}\raggedright
REGIONS
\end{minipage} & \begin{minipage}[b]{\linewidth}\raggedright
SOLUTION
\end{minipage} & \begin{minipage}[b]{\linewidth}\raggedright
SUMMARY
\end{minipage} & \begin{minipage}[b]{\linewidth}\raggedright
SCHEDULE
\end{minipage} \\
\midrule\noalign{}
\endhead
\bottomrule\noalign{}
\endlastfoot
\end{longtable}

\subsubsection{Description}\label{description-882}

The COMPINJK keyword assigns injection well relative permeability values
to well connections.

This keyword is not supported by OPM Flow but would change the results
if supported so the simulation will be stopped.

\subsection{COMPLMPL -- Assign Well LGR Connections to
Completions}\label{complmpl-assign-well-lgr-connections-to-completions}

\begin{longtable}[]{@{}
  >{\raggedright\arraybackslash}p{(\linewidth - 14\tabcolsep) * \real{0.1250}}
  >{\raggedright\arraybackslash}p{(\linewidth - 14\tabcolsep) * \real{0.1250}}
  >{\raggedright\arraybackslash}p{(\linewidth - 14\tabcolsep) * \real{0.1250}}
  >{\raggedright\arraybackslash}p{(\linewidth - 14\tabcolsep) * \real{0.1250}}
  >{\raggedright\arraybackslash}p{(\linewidth - 14\tabcolsep) * \real{0.1250}}
  >{\raggedright\arraybackslash}p{(\linewidth - 14\tabcolsep) * \real{0.1250}}
  >{\raggedright\arraybackslash}p{(\linewidth - 14\tabcolsep) * \real{0.1250}}
  >{\raggedright\arraybackslash}p{(\linewidth - 14\tabcolsep) * \real{0.1250}}@{}}
\toprule\noalign{}
\begin{minipage}[b]{\linewidth}\raggedright
RUNSPEC
\end{minipage} & \begin{minipage}[b]{\linewidth}\raggedright
GRID
\end{minipage} & \begin{minipage}[b]{\linewidth}\raggedright
EDIT
\end{minipage} & \begin{minipage}[b]{\linewidth}\raggedright
PROPS
\end{minipage} & \begin{minipage}[b]{\linewidth}\raggedright
REGIONS
\end{minipage} & \begin{minipage}[b]{\linewidth}\raggedright
SOLUTION
\end{minipage} & \begin{minipage}[b]{\linewidth}\raggedright
SUMMARY
\end{minipage} & \begin{minipage}[b]{\linewidth}\raggedright
SCHEDULE
\end{minipage} \\
\midrule\noalign{}
\endhead
\bottomrule\noalign{}
\endlastfoot
\end{longtable}

\subsubsection{Description}\label{description-883}

The COMPLMPL keyword assigns well connections in a LGR, as defined by
the COMPDATL keyword in the SCHEDULE section, to completion intervals.
This ``lumping'' of the connections to various completion intervals
allows automatic workovers and economic criteria to be applied to the
completions (that is a set of connections) as opposed to the
connections. This allows for a more realistic approach for workovers
operations.

For example, if the water cut criteria for working over a well was set
to 95\%, and the average grid block connection thickness was one meter,
then once a well's water cut reached 95\% the worst offending one meter
connection would be shut-in. If the well's actual perforation interval
was 10 meters and the 10 connections were lumped as one completion, then
when the water cut limit of 95\% is reach, the completion would be
shut-in, that is all of the 10 connections within the completion would
be shut-in.

As the keyword is used to lump connections into a completions, the
simulator adds together the contribution from all connections in the
completion and uses the total values to test the economic limits. Note
that a connection can only belong to one completion. In addition,
completions can be used instead of connections in the WELOPEN and
WPIMULT keywords if the completions have been defined by COMPLUMP for a
well.

This keyword is not supported by OPM Flow but would change the results
if supported so the simulation will be stopped.

\begin{longtable}[]{@{}
  >{\raggedright\arraybackslash}p{(\linewidth - 6\tabcolsep) * \real{0.2500}}
  >{\raggedright\arraybackslash}p{(\linewidth - 6\tabcolsep) * \real{0.2500}}
  >{\raggedright\arraybackslash}p{(\linewidth - 6\tabcolsep) * \real{0.2500}}
  >{\raggedright\arraybackslash}p{(\linewidth - 6\tabcolsep) * \real{0.2500}}@{}}
\toprule\noalign{}
\begin{minipage}[b]{\linewidth}\raggedright
No.
\end{minipage} & \begin{minipage}[b]{\linewidth}\raggedright
Name
\end{minipage} & \begin{minipage}[b]{\linewidth}\raggedright
Description
\end{minipage} & \begin{minipage}[b]{\linewidth}\raggedright
Default
\end{minipage} \\
\midrule\noalign{}
\endhead
\bottomrule\noalign{}
\endlastfoot
Field & Metric & Laboratory & \\
1 & WELNAME & A character string of up to eight characters in length
that defines the well name for which the well connection data are being
defined. Note that the well name (WELNAME) must have been declared
previously using the WELSPECS keyword in the SCHEDULE section, otherwise
an error may occur. & None \\
2 & LGRNAME & A character string of up to eight characters in length
that defines the LGR name for which the well LGR connection data are
being defined. Note that the well name (LGRNAME) must have been declared
previously using the WELSPECL keyword in the SCHEDULE section, otherwise
an error may occur. If defaulted with 1* the LGR on the WELSPECL keyword
will be utilized. & Defined \\
3 & I & A positive integer greater than or equal to zero and less than
or equal to NX that defines the connection location in the I-direction.
If set to zero or defaulted with 1* then all connections in the
I-direction that also satisfy J, K1 and K2 criteria are assigned the
ICOMP completion number. & 0 \\
4 & J & A positive integer greater than or equal to zero and less than
or equal to NY that defines the connection location in the J-direction.
If set to zero or defaulted with 1* then all connections in the
J-direction that also satisfy I, K1 and K2 criteria are assigned the
ICOMP completion number. & 0 \\
5 & K1 & A positive integer greater than or equal to one and less than
or equal to NZ that defines the UPPER connection location in the
K-direction. If set to zero or defaulted with 1* then the upper most
connection in the well is used. & 0 \\
6 & K2 & A positive integer greater than or equal to K1 and less than or
equal to NZ that defines the LOWER connection location in the
K-direction. If set to zero or defaulted with 1* then the low most
connection in the well is used. & 0 \\
7 & ICOMP & An integer greater than or equal to one and less than or
equal to MXCONS as defined on the WELLDIMS keyword in the RUNSPEC
section, that defines the completion number of the currently defined set
of connections. If I, J, K1, K2 are all set to zero or defaulted to 1*,
then all connections in the well have the same completion number of
ICOMP. & None \\
Notes: & & & \\
\end{longtable}

\emph{Table 12.13: COMPLUPL Keyword Description}

Multiple grid block connections can be defined on one record for
vertical wells by assigning different values to K1 and K2, for deviated
and horizontal wells this may not be possible and therefore each grid
block connection must be separately defined by setting K1 equal to K2.

See also the COMPDATL keyword in the SCHEDULE section.

\subsubsection{Example}\label{example-508}

The following example defines the connections for two vertical oil wells
using the COMPDATL keyword and the re-allocation of the connections to
completions intervals using the COMPLMPL keyword.

\begin{verbatim}
--
--       WELL CONNECTION DATA FOR LGR WELLS
--
--       WELL    LGR      --- LOCATION ---  OPEN   SAT   CONN   WELL   D     DIR
--       NAME    NAME      II  JJ  K1  K2   SHUT   TAB   FACT   DIA    FACT  PEN
COMPDATL
         OP01    OP01LGR    14  13  20  56  OPEN   1*    1*    0.708   3*     Z /
         OP01    OP01LGR    14  13  75 100  SHUT   1*    1*    0.708   3*     Z /
         OP02    OP02LGR    35  96  75 100  OPEN   1*    1*    0.708   3*     Z /
/
--
--       ASSIGN WELL CONNECTIONS TO COMPLETIONS
--
--       WELL    LGR      --- LOCATION ---  COMPL
--       NAME    NAME      II  JJ  K1  K2   NO.                                                                                                                                                                                                                                                                              COMPLMPL                                                                                                                                                                                                                                OP01    OP01LGR   1*  1*  20  56    1          / COMPLETION NO. 01
             OP01    OP01LGR   1*  1*  75 100    2          / COMPLETION NO. 02
             OP02    OP02LGR   1*  1*  75  85    1          / COMPLETION NO. 01
             OP02    OP21LGR   1*  1*  86 100    2          / COMPLETION NO. 02
/
\end{verbatim}

Here the well OP01 connections (14, 13, 20) to (14, 13, 56) are assigned
to completion number one and connections (14, 13, 75) to (14, 13, 100)
are assigned to completion number two. Well OP02 has only one set of
connection data from cells (35, 96, 75) to cells (35, 96, 100), but they
have split into two separate completion intervals, with connections (35.
96, 75) to (35. 96, 85) assigned to completion interval number one and
(35, 96, 86) to (35, 96, 100) to completion number two.

\subsection{COMPLUMP -- Assign Well Connections to
Completions}\label{complump-assign-well-connections-to-completions}

\begin{longtable}[]{@{}
  >{\raggedright\arraybackslash}p{(\linewidth - 14\tabcolsep) * \real{0.1250}}
  >{\raggedright\arraybackslash}p{(\linewidth - 14\tabcolsep) * \real{0.1250}}
  >{\raggedright\arraybackslash}p{(\linewidth - 14\tabcolsep) * \real{0.1250}}
  >{\raggedright\arraybackslash}p{(\linewidth - 14\tabcolsep) * \real{0.1250}}
  >{\raggedright\arraybackslash}p{(\linewidth - 14\tabcolsep) * \real{0.1250}}
  >{\raggedright\arraybackslash}p{(\linewidth - 14\tabcolsep) * \real{0.1250}}
  >{\raggedright\arraybackslash}p{(\linewidth - 14\tabcolsep) * \real{0.1250}}
  >{\raggedright\arraybackslash}p{(\linewidth - 14\tabcolsep) * \real{0.1250}}@{}}
\toprule\noalign{}
\begin{minipage}[b]{\linewidth}\raggedright
RUNSPEC
\end{minipage} & \begin{minipage}[b]{\linewidth}\raggedright
GRID
\end{minipage} & \begin{minipage}[b]{\linewidth}\raggedright
EDIT
\end{minipage} & \begin{minipage}[b]{\linewidth}\raggedright
PROPS
\end{minipage} & \begin{minipage}[b]{\linewidth}\raggedright
REGIONS
\end{minipage} & \begin{minipage}[b]{\linewidth}\raggedright
SOLUTION
\end{minipage} & \begin{minipage}[b]{\linewidth}\raggedright
SUMMARY
\end{minipage} & \begin{minipage}[b]{\linewidth}\raggedright
SCHEDULE
\end{minipage} \\
\midrule\noalign{}
\endhead
\bottomrule\noalign{}
\endlastfoot
\end{longtable}

\subsubsection{Description}\label{description-884}

The COMPLUMP keyword assigns connections, as defined by the COMPDAT
keyword in the SCHEDULE section, to completion intervals. This
``lumping'' of the connections to various completion intervals allows
automatic workovers and economic criteria to be applied to the
completions (that is a set of connections) as opposed to the
connections. This allows for a more realistic approach for workovers
operations.

For example, if the water cut criteria for working over a well was set
to 95\%, and the average grid block connection thickness was one meter,
then once a well's water cut reached 95\% the worst offending one meter
connection would be shut-in. If the well's actual perforation interval
was 10 meters and the 10 connections were lumped as one completion, then
when the water cut limit of 95\% is reach, the completion would be
shut-in, that is all of the 10 connections within the completion would
be shut-in.

As the keyword is used to lump connections into a completions, the
simulator adds together the contribution from all connections in the
completion and uses the total values to test the economic limits. Note
that a connection can only belong to one completion. In addition,
completions can be used instead of connections in the WELOPEN and
WPIMULT keywords if the completions have been defined by COMPLUMP for a
well.

\begin{longtable}[]{@{}
  >{\raggedright\arraybackslash}p{(\linewidth - 6\tabcolsep) * \real{0.2500}}
  >{\raggedright\arraybackslash}p{(\linewidth - 6\tabcolsep) * \real{0.2500}}
  >{\raggedright\arraybackslash}p{(\linewidth - 6\tabcolsep) * \real{0.2500}}
  >{\raggedright\arraybackslash}p{(\linewidth - 6\tabcolsep) * \real{0.2500}}@{}}
\toprule\noalign{}
\begin{minipage}[b]{\linewidth}\raggedright
No.
\end{minipage} & \begin{minipage}[b]{\linewidth}\raggedright
Name
\end{minipage} & \begin{minipage}[b]{\linewidth}\raggedright
Description
\end{minipage} & \begin{minipage}[b]{\linewidth}\raggedright
Default
\end{minipage} \\
\midrule\noalign{}
\endhead
\bottomrule\noalign{}
\endlastfoot
Field & Metric & Laboratory & \\
1 & WELNAME & A character string of up to eight characters in length
that defines the well name for which the well connection data are being
defined. Note that the well name (WELNAME) must have been declared
previously using the WELSPECS keyword in the SCHEDULE section, otherwise
an error may occur. & None \\
2 & I & A positive integer greater than or equal to zero and less than
or equal to NX that defines the connection location in the I-direction.
If set to zero or defaulted with 1* then all connections in the
I-direction that also satisfy J, K1 and K2 criteria are assigned the
ICOMP completion number. & 0 \\
3 & J & A positive integer greater than or equal to zero and less than
or equal to NY that defines the connection location in the J-direction.
If set to zero or defaulted with 1* then all connections in the
J-direction that also satisfy I, K1 and K2 criteria are assigned the
ICOMP completion number. & 0 \\
4 & K1 & A positive integer greater than or equal to one and less than
or equal to NZ that defines the UPPER connection location in the
K-direction. If set to zero or defaulted with 1* then the uppermost
connection in the well is used. & 0 \\
5 & K2 & A positive integer greater than or equal to K1 and less than or
equal to NZ that defines the LOWER connection location in the
K-direction. If set to zero or defaulted with 1* then the lowermost
connection in the well is used. & 0 \\
6 & ICOMP & An integer greater than or equal to one and less than or
equal to MXCONS as defined on the WELLDIMS keyword in the RUNSPEC
section, that defines the completion number of the currently defined set
of connections. If I, J, K1, K2 are all set to zero or defaulted to 1*,
then all connections in the well have the same completion number of
ICOMP. & None \\
Notes: & & & \\
\end{longtable}

\emph{Table 12.14: COMPLUMP Keyword Description}

Multiple grid block connections can be defined on one record for
vertical wells by assigning different values to K1 and K2, for deviated
and horizontal wells this may not be possible and therefore each grid
block connection must be separately defined by setting K1 equal to K2.

See also the COMPDAT keyword in the SCHEDULE section.

\subsubsection{Example}\label{example-509}

The following example defines the connections for two vertical oil wells
using the COMPDAT keyword and the re-allocation of the connections to
completions intervals using the COMPLUMP keyword.

\begin{verbatim}
--
--       WELL CONNECTION DATA
--
-- WELL  --- LOCATION ---  OPEN   SAT   CONN   WELL   KH    SKIN   D     DIR
-- NAME   II  JJ  K1  K2   SHUT   TAB   FACT   DIA    FACT  FACT   FACT  PEN
COMPDAT
OP01      14  13  20  56   OPEN   1*    1*    0.708   1*    0.0    1*    'Z' /
OP01      14  13  75 100   SHUT   1*    1*    0.708   1*    0.0    1*    'Z' /
OP02      35  96  75 100   OPEN   1*    1*    0.708   1*    0.0    1*    'Z' /
/
--
--       ASSIGN WELL CONNECTIONS TO COMPLETIONS
--
-- WELL  --- LOCATION ---  COMPL
-- NAME   II  JJ  K1  K2   NO.                                                                                                                                                                                                                                                                              COMPLUMP                                                                                                                                                                                                       OP01      1*  1*  20  56    1                              / COMPLETION NO. 01
OP01      1*  1*  75 100    2                              / COMPLETION NO. 02
OP02      1*  1*  75  85    1                              / COMPLETION NO. 01
OP02      1*  1*  86 100    2                              / COMPLETION NO. 02
/
\end{verbatim}

Here the well OP01 connections (14, 13, 20) to (14, 13, 56) are assigned
to completion number one and connections (14, 13, 75) to (14, 13, 100)
are assigned to completion number two. Well OP02 has only one set of
connection data from cells (35, 96, 75) to cells (35, 96, 100), but they
have been split into two separate completion intervals, with connections
(35, 96, 75) to (35. 96, 85) assigned to completion interval number one
and (35, 96, 86) to (35, 96, 100) assigned to completion number two.

\subsection{COMPOFF -- Deactivate Network Automatic
Compressors}\label{compoff-deactivate-network-automatic-compressors}

\begin{longtable}[]{@{}
  >{\raggedright\arraybackslash}p{(\linewidth - 14\tabcolsep) * \real{0.1250}}
  >{\raggedright\arraybackslash}p{(\linewidth - 14\tabcolsep) * \real{0.1250}}
  >{\raggedright\arraybackslash}p{(\linewidth - 14\tabcolsep) * \real{0.1250}}
  >{\raggedright\arraybackslash}p{(\linewidth - 14\tabcolsep) * \real{0.1250}}
  >{\raggedright\arraybackslash}p{(\linewidth - 14\tabcolsep) * \real{0.1250}}
  >{\raggedright\arraybackslash}p{(\linewidth - 14\tabcolsep) * \real{0.1250}}
  >{\raggedright\arraybackslash}p{(\linewidth - 14\tabcolsep) * \real{0.1250}}
  >{\raggedright\arraybackslash}p{(\linewidth - 14\tabcolsep) * \real{0.1250}}@{}}
\toprule\noalign{}
\begin{minipage}[b]{\linewidth}\raggedright
RUNSPEC
\end{minipage} & \begin{minipage}[b]{\linewidth}\raggedright
GRID
\end{minipage} & \begin{minipage}[b]{\linewidth}\raggedright
EDIT
\end{minipage} & \begin{minipage}[b]{\linewidth}\raggedright
PROPS
\end{minipage} & \begin{minipage}[b]{\linewidth}\raggedright
REGIONS
\end{minipage} & \begin{minipage}[b]{\linewidth}\raggedright
SOLUTION
\end{minipage} & \begin{minipage}[b]{\linewidth}\raggedright
SUMMARY
\end{minipage} & \begin{minipage}[b]{\linewidth}\raggedright
SCHEDULE
\end{minipage} \\
\midrule\noalign{}
\endhead
\bottomrule\noalign{}
\endlastfoot
\end{longtable}

\subsubsection{Description}\label{description-885}

The COMPOFF keyword deactivates network automatic compressors defined
via the GASFCOMP keyword in the SCHEDULE section.

This keyword is not supported by OPM Flow but would change the results
if supported so the simulation will be stopped.

\subsection{COMPORD -- Define Well Connection
Ordering}\label{compord-define-well-connection-ordering}

\begin{longtable}[]{@{}
  >{\raggedright\arraybackslash}p{(\linewidth - 14\tabcolsep) * \real{0.1250}}
  >{\raggedright\arraybackslash}p{(\linewidth - 14\tabcolsep) * \real{0.1250}}
  >{\raggedright\arraybackslash}p{(\linewidth - 14\tabcolsep) * \real{0.1250}}
  >{\raggedright\arraybackslash}p{(\linewidth - 14\tabcolsep) * \real{0.1250}}
  >{\raggedright\arraybackslash}p{(\linewidth - 14\tabcolsep) * \real{0.1250}}
  >{\raggedright\arraybackslash}p{(\linewidth - 14\tabcolsep) * \real{0.1250}}
  >{\raggedright\arraybackslash}p{(\linewidth - 14\tabcolsep) * \real{0.1250}}
  >{\raggedright\arraybackslash}p{(\linewidth - 14\tabcolsep) * \real{0.1250}}@{}}
\toprule\noalign{}
\begin{minipage}[b]{\linewidth}\raggedright
RUNSPEC
\end{minipage} & \begin{minipage}[b]{\linewidth}\raggedright
GRID
\end{minipage} & \begin{minipage}[b]{\linewidth}\raggedright
EDIT
\end{minipage} & \begin{minipage}[b]{\linewidth}\raggedright
PROPS
\end{minipage} & \begin{minipage}[b]{\linewidth}\raggedright
REGIONS
\end{minipage} & \begin{minipage}[b]{\linewidth}\raggedright
SOLUTION
\end{minipage} & \begin{minipage}[b]{\linewidth}\raggedright
SUMMARY
\end{minipage} & \begin{minipage}[b]{\linewidth}\raggedright
SCHEDULE
\end{minipage} \\
\midrule\noalign{}
\endhead
\bottomrule\noalign{}
\endlastfoot
\end{longtable}

\subsubsection{Description}\label{description-886}

The COMPORD keyword defines how the well connection data entered on the
COMPDAT keyword in the SCHEDULE section are to be ordered for a well.

\begin{longtable}[]{@{}
  >{\raggedright\arraybackslash}p{(\linewidth - 6\tabcolsep) * \real{0.2500}}
  >{\raggedright\arraybackslash}p{(\linewidth - 6\tabcolsep) * \real{0.2500}}
  >{\raggedright\arraybackslash}p{(\linewidth - 6\tabcolsep) * \real{0.2500}}
  >{\raggedright\arraybackslash}p{(\linewidth - 6\tabcolsep) * \real{0.2500}}@{}}
\toprule\noalign{}
\begin{minipage}[b]{\linewidth}\raggedright
No.
\end{minipage} & \begin{minipage}[b]{\linewidth}\raggedright
Name
\end{minipage} & \begin{minipage}[b]{\linewidth}\raggedright
Description
\end{minipage} & \begin{minipage}[b]{\linewidth}\raggedright
Default
\end{minipage} \\
\midrule\noalign{}
\endhead
\bottomrule\noalign{}
\endlastfoot
Field & Metric & Laboratory & \\
1 & WELNAME & A character string of up to eight characters in length
that defines the well name for which the well connection data are being
defined. Note that the well name (WELNAME) must have been declared
previously using the WELSPECS keyword in the SCHEDULE section, otherwise
an error may occur. & None \\
2 & COMPORD & A character string that defines the method for ordering
the well connections given on the COMPDAT keyword, and should be set to
DEPTH, INPUT, or TRACK. All options are now supported by OPM Flow. &
TRACK \\
Notes: & & & \\
\end{longtable}

\emph{Table 12.15: COMPORD Keyword Description}

See also the COMPDAT keyword in the SCHEDULE section.

\begin{tcolorbox}[enhanced jigsaw, title=\textcolor{quarto-callout-note-color}{\faInfo}\hspace{0.5em}{Note}, coltitle=black, breakable, toprule=.15mm, colbacktitle=quarto-callout-note-color!10!white, leftrule=.75mm, opacitybacktitle=0.6, colback=white, bottomtitle=1mm, toptitle=1mm, titlerule=0mm, left=2mm, colframe=quarto-callout-note-color-frame, arc=.35mm, rightrule=.15mm, bottomrule=.15mm, opacityback=0]

If visual inspection of the well trajectories in the model indicate
problematic or unrealistic well connections, the options on this keyword
may be useful in correcting the issue.

\end{tcolorbox}

\subsubsection{Example}\label{example-510}

The following example defines the connections for two vertical oil wells
using the COMPDAT keyword and the COMPORD to defined the connection
ordering for the wells.

\begin{verbatim}
--
--       WELL CONNECTION DATA
--
-- WELL  --- LOCATION ---  OPEN   SAT   CONN   WELL   KH    SKIN   D     DIR
-- NAME   II  JJ  K1  K2   SHUT   TAB   FACT   DIA    FACT  FACT   FACT  PEN
COMPDAT
OP01      1*  1*  20  56   OPEN   1*    1*    0.708   1*    0.0    1*    'Z' /
OP01      1*  1*  75 100   SHUT   1*    1*    0.708   1*    0.0    1*    'Z' /
OP02      35  96  75 100   OPEN   1*    1*    0.708   1*    0.0    1*    'Z' /
--
--       DEFINE WELL CONNECTION ORDERING
--
-- WELL  COMPL
-- NAME  ORDER
COMPORD
OP01     DEPTH                                             /
OP02     DEPTH                                             /
/
\end{verbatim}

The DEPTH option has been chosen because both wells are vertical. Also
one could use the following format instead for the COMPORD:

\begin{verbatim}
--
--       DEFINE WELL CONNECTION ORDERING
--
-- WELL  COMPL
-- NAME  ORDER
COMPORD
*        DEPTH                                             /
/
\end{verbatim}

as both wells should utilize the DEPTH option. This version would set
all wells in the model to DEPTH connection ordering.

\subsection{COMPRIV -- Define Grid Cell Connections to a
River}\label{compriv-define-grid-cell-connections-to-a-river}

\begin{longtable}[]{@{}
  >{\raggedright\arraybackslash}p{(\linewidth - 14\tabcolsep) * \real{0.1250}}
  >{\raggedright\arraybackslash}p{(\linewidth - 14\tabcolsep) * \real{0.1250}}
  >{\raggedright\arraybackslash}p{(\linewidth - 14\tabcolsep) * \real{0.1250}}
  >{\raggedright\arraybackslash}p{(\linewidth - 14\tabcolsep) * \real{0.1250}}
  >{\raggedright\arraybackslash}p{(\linewidth - 14\tabcolsep) * \real{0.1250}}
  >{\raggedright\arraybackslash}p{(\linewidth - 14\tabcolsep) * \real{0.1250}}
  >{\raggedright\arraybackslash}p{(\linewidth - 14\tabcolsep) * \real{0.1250}}
  >{\raggedright\arraybackslash}p{(\linewidth - 14\tabcolsep) * \real{0.1250}}@{}}
\toprule\noalign{}
\begin{minipage}[b]{\linewidth}\raggedright
RUNSPEC
\end{minipage} & \begin{minipage}[b]{\linewidth}\raggedright
GRID
\end{minipage} & \begin{minipage}[b]{\linewidth}\raggedright
EDIT
\end{minipage} & \begin{minipage}[b]{\linewidth}\raggedright
PROPS
\end{minipage} & \begin{minipage}[b]{\linewidth}\raggedright
REGIONS
\end{minipage} & \begin{minipage}[b]{\linewidth}\raggedright
SOLUTION
\end{minipage} & \begin{minipage}[b]{\linewidth}\raggedright
SUMMARY
\end{minipage} & \begin{minipage}[b]{\linewidth}\raggedright
SCHEDULE
\end{minipage} \\
\midrule\noalign{}
\endhead
\bottomrule\noalign{}
\endlastfoot
\end{longtable}

\subsubsection{Description}\label{description-887}

The COMPRIV keyword defines grid cell connections to a river, for when
the River option has been activated via the RIVRDIMS keyword in the
RUNSPEC section.

This keyword is not supported by OPM Flow but would change the results
if supported so the simulation will be stopped.

\subsection{COMPRP -- Re-Scale Fluid Saturations of Well
Connections}\label{comprp-re-scale-fluid-saturations-of-well-connections}

\begin{longtable}[]{@{}
  >{\raggedright\arraybackslash}p{(\linewidth - 14\tabcolsep) * \real{0.1250}}
  >{\raggedright\arraybackslash}p{(\linewidth - 14\tabcolsep) * \real{0.1250}}
  >{\raggedright\arraybackslash}p{(\linewidth - 14\tabcolsep) * \real{0.1250}}
  >{\raggedright\arraybackslash}p{(\linewidth - 14\tabcolsep) * \real{0.1250}}
  >{\raggedright\arraybackslash}p{(\linewidth - 14\tabcolsep) * \real{0.1250}}
  >{\raggedright\arraybackslash}p{(\linewidth - 14\tabcolsep) * \real{0.1250}}
  >{\raggedright\arraybackslash}p{(\linewidth - 14\tabcolsep) * \real{0.1250}}
  >{\raggedright\arraybackslash}p{(\linewidth - 14\tabcolsep) * \real{0.1250}}@{}}
\toprule\noalign{}
\begin{minipage}[b]{\linewidth}\raggedright
RUNSPEC
\end{minipage} & \begin{minipage}[b]{\linewidth}\raggedright
GRID
\end{minipage} & \begin{minipage}[b]{\linewidth}\raggedright
EDIT
\end{minipage} & \begin{minipage}[b]{\linewidth}\raggedright
PROPS
\end{minipage} & \begin{minipage}[b]{\linewidth}\raggedright
REGIONS
\end{minipage} & \begin{minipage}[b]{\linewidth}\raggedright
SOLUTION
\end{minipage} & \begin{minipage}[b]{\linewidth}\raggedright
SUMMARY
\end{minipage} & \begin{minipage}[b]{\linewidth}\raggedright
SCHEDULE
\end{minipage} \\
\midrule\noalign{}
\endhead
\bottomrule\noalign{}
\endlastfoot
\end{longtable}

\subsubsection{Description}\label{description-888}

The COPMPRP keyword re-scales the fluid saturations of a well's
connection to the grid block.

This keyword is not supported by OPM Flow but would change the results
if supported so the simulation will be stopped.

\subsection{COMPRPL -- Re-Scale Fluid Saturations of Well LGR
Connections}\label{comprpl-re-scale-fluid-saturations-of-well-lgr-connections}

\begin{longtable}[]{@{}
  >{\raggedright\arraybackslash}p{(\linewidth - 14\tabcolsep) * \real{0.1250}}
  >{\raggedright\arraybackslash}p{(\linewidth - 14\tabcolsep) * \real{0.1250}}
  >{\raggedright\arraybackslash}p{(\linewidth - 14\tabcolsep) * \real{0.1250}}
  >{\raggedright\arraybackslash}p{(\linewidth - 14\tabcolsep) * \real{0.1250}}
  >{\raggedright\arraybackslash}p{(\linewidth - 14\tabcolsep) * \real{0.1250}}
  >{\raggedright\arraybackslash}p{(\linewidth - 14\tabcolsep) * \real{0.1250}}
  >{\raggedright\arraybackslash}p{(\linewidth - 14\tabcolsep) * \real{0.1250}}
  >{\raggedright\arraybackslash}p{(\linewidth - 14\tabcolsep) * \real{0.1250}}@{}}
\toprule\noalign{}
\begin{minipage}[b]{\linewidth}\raggedright
RUNSPEC
\end{minipage} & \begin{minipage}[b]{\linewidth}\raggedright
GRID
\end{minipage} & \begin{minipage}[b]{\linewidth}\raggedright
EDIT
\end{minipage} & \begin{minipage}[b]{\linewidth}\raggedright
PROPS
\end{minipage} & \begin{minipage}[b]{\linewidth}\raggedright
REGIONS
\end{minipage} & \begin{minipage}[b]{\linewidth}\raggedright
SOLUTION
\end{minipage} & \begin{minipage}[b]{\linewidth}\raggedright
SUMMARY
\end{minipage} & \begin{minipage}[b]{\linewidth}\raggedright
SCHEDULE
\end{minipage} \\
\midrule\noalign{}
\endhead
\bottomrule\noalign{}
\endlastfoot
\end{longtable}

\subsubsection{Description}\label{description-889}

The COPMPRPL keyword re-scales the fluid saturations of a well's
connection to an LGR grid block, for when the Local Grid Refinement
(``LGR'') option has been activated by the LGR keyword in the RUNSPEC
section.

This keyword is not supported by OPM Flow but would change the results
if supported so the simulation will be stopped.

\subsection{COMPSEGL -- Define Well Connections for Multi-Segment Wells
in a
LGR}\label{compsegl-define-well-connections-for-multi-segment-wells-in-a-lgr}

\begin{longtable}[]{@{}
  >{\raggedright\arraybackslash}p{(\linewidth - 14\tabcolsep) * \real{0.1250}}
  >{\raggedright\arraybackslash}p{(\linewidth - 14\tabcolsep) * \real{0.1250}}
  >{\raggedright\arraybackslash}p{(\linewidth - 14\tabcolsep) * \real{0.1250}}
  >{\raggedright\arraybackslash}p{(\linewidth - 14\tabcolsep) * \real{0.1250}}
  >{\raggedright\arraybackslash}p{(\linewidth - 14\tabcolsep) * \real{0.1250}}
  >{\raggedright\arraybackslash}p{(\linewidth - 14\tabcolsep) * \real{0.1250}}
  >{\raggedright\arraybackslash}p{(\linewidth - 14\tabcolsep) * \real{0.1250}}
  >{\raggedright\arraybackslash}p{(\linewidth - 14\tabcolsep) * \real{0.1250}}@{}}
\toprule\noalign{}
\begin{minipage}[b]{\linewidth}\raggedright
RUNSPEC
\end{minipage} & \begin{minipage}[b]{\linewidth}\raggedright
GRID
\end{minipage} & \begin{minipage}[b]{\linewidth}\raggedright
EDIT
\end{minipage} & \begin{minipage}[b]{\linewidth}\raggedright
PROPS
\end{minipage} & \begin{minipage}[b]{\linewidth}\raggedright
REGIONS
\end{minipage} & \begin{minipage}[b]{\linewidth}\raggedright
SOLUTION
\end{minipage} & \begin{minipage}[b]{\linewidth}\raggedright
SUMMARY
\end{minipage} & \begin{minipage}[b]{\linewidth}\raggedright
SCHEDULE
\end{minipage} \\
\midrule\noalign{}
\endhead
\bottomrule\noalign{}
\endlastfoot
\end{longtable}

\subsubsection{Description}\label{description-890}

The COMSEGSL keyword defines how a multi-segment well is connected to
the reservoir by defining or modifying existing well connections in an
LGR. Note that well must have been previously define by the WELSPECL
keyword in the SCHEDULE section and the well connections must have been
previously defined via the COMPDATL keyword in the SCHEDULE section. The
COMPSEGL keyword should be repeated for each multi-segment well in the
model.

This keyword is not supported by OPM Flow but would change the results
if supported so the simulation will be stopped.

\begin{longtable}[]{@{}
  >{\raggedright\arraybackslash}p{(\linewidth - 6\tabcolsep) * \real{0.2500}}
  >{\raggedright\arraybackslash}p{(\linewidth - 6\tabcolsep) * \real{0.2500}}
  >{\raggedright\arraybackslash}p{(\linewidth - 6\tabcolsep) * \real{0.2500}}
  >{\raggedright\arraybackslash}p{(\linewidth - 6\tabcolsep) * \real{0.2500}}@{}}
\toprule\noalign{}
\begin{minipage}[b]{\linewidth}\raggedright
No.
\end{minipage} & \begin{minipage}[b]{\linewidth}\raggedright
Name
\end{minipage} & \begin{minipage}[b]{\linewidth}\raggedright
Description
\end{minipage} & \begin{minipage}[b]{\linewidth}\raggedright
Default
\end{minipage} \\
\midrule\noalign{}
\endhead
\bottomrule\noalign{}
\endlastfoot
Field & Metric & Laboratory & \\
1-1 & WELNAME & A character string of up to eight characters in length
that defines the well name for which a multi-segment well is being
defined. Note that the well name (WELNAME) must have been declared
previously using the WELSPECL keyword in the SCHEDULE section, otherwise
an error may occur. & None \\
1-2 & / & Record terminated by a ``/'' & Not Applicable \\
2-1 & LGRNAME & A character string of up to eight characters in length
that defines the name of the local grid refinement for which the well is
assigned to. & None \\
2-2 & I & A positive integer greater than or equal to one and less than
or equal to NX that defines the connection location in the I-direction.
& None \\
2-3 & J & A positive integer greater than or equal to zero and less than
or equal to NY that defines the connection location in the J-direction.
& None \\
2-4 & K & A positive integer greater than or equal to zero and less than
or equal to NZ that defines the connection location in the K-direction.
& None \\
2-5 & IBRANCH & A positive integer greater than or equal to one and less
than or equal to MXBRAN on WSEGDIMS keyword in the RUNSPEC section that
defines the branch number of the defined I, J and K connection. &
None \\
2.6 & DEPTH1 & DEPTH1 is a real positive value that defines the length
of the tubing from the tubing head or wellhead at the surface to the
start of the connection in the I, J, K cell. & None \\
feet & m & cm & \\
2-7 & DEPTH2 & DEPTH2 is a real positive value that defines the length
of the tubing from the tubing head or wellhead at the surface to the end
of the connection in the I, J, K cell. & \\
feet & m & cm & None \\
2-8 & DIRECT & A one letter character string that defines the
orientation of the connections and should be set to either X, Y, or Z.
The direction of connections also determines the length of the
connection The default value is for a vertical connection, that is
DIRECT is defaulted to Z. & Z \\
2-9 & IEND & IEND is positive or negative integer, that is not equal to
zero that is set to one of the following: that defines the end of the
range of the connections depending on the value of DIRECT. For example,
if DIRECT is equal to Y or J then the IEND will be associated with the
J-direction. The value may be positive or negative but must be
calculated to remain within the grid. For example for NY is set 100 on
the DIMENS keyword in the RUNSPEC section and J on this record set to
50, then IEND most range between -49 to +50. Currently this option is
not supported by OPM Flow. & None \\
2-10 & DEPTH3 & DEPTH3 is a real positive value that defines the datum
depth for this set of connection, normally taken as the mid-point of the
perforations associated with this set of connections. Currently this
option is not supported by OPM Flow. & None \\
feet & m & cm & \\
2-11 & LENGTH & LENGTH is a real positive value that defines the length
of the well for this set of completions that is used in thermal
calculations.. Currently this option is not supported by OPM Flow. &
None \\
feet & m & cm & \\
2-12 & ISEG & A real positive values equal to or greater than zero that
defines the coordinate in the x-direction of the nodal point of this
segment that is used for display purposes only. Currently this option is
not supported by OPM Flow. & None \\
feet & m & cm & \\
2-13 & / & Record terminated by a ``/'' & Not Applicable \\
Notes: & & & \\
\end{longtable}

\emph{Table 12.16: COMPSEGL Keyword Description}

The total number of wells and completions should be defined via the
WELLSDIMS keyword and the number of multi-segment wells and completions
should be declared on the WSEGDIMS keyword, both keywords are in the
RUNSPEC section.

See also the WELSPECL keyword to define wells in an LGR, the COMPDATL
keyword to define the well connections for both ordinary wells and
multi-segment wells with an LGR, and the COMPSEGS keyword to define a
multi-segment connections in the global grid. All the aforementioned
keywords are described in the SCHEDULE section.

\subsubsection{Example}\label{example-511}

The following example defines the completions for two oil producing
segment oil wells (OP01 and OP02) using the COMPSEGS keywords.

\begin{verbatim}
--       COMPLETION SEGMENT SPECIFICATION DATA
--
-- WELL
-- NAME
COMPSEGS
OP01                                                                           /
--      LGR  --LOCATION--  BRAN  TUBING  NODAL   DIR  LOC    MID    COMP    ISEG
--      NAME  II  JJ  K1   NO    LENGTH  DEPTH   PEN  I,J,K  PERFS  LENGTH
        LGR01 10  10   1    1    2512.5  2525.0                                /
        LGR01 10  10   2    1    2525.0  2550.0                                /
        LGR01 10  10   3    1    2550.0  2575.0                                /
        LGR01 10  10   4    1    2575.0  2600.0                                /
        LGR01 10  10   5    1    2600.0  2625.0                                /
        LGR01 10  10   6    1    2625.0  2650.0                                /

        LGR01  9  10   2    2    2637.5  2837.5                                /
        LGR01  8  10   2    2    2837.5  3037.5                                /
        LGR01  7  10   2    2    3037.5  3237.5                                /
        LGR01  6  10   2    2    3237.5  3437.5                                /
        LGR01  5  10   2    2    3437.5  3637.5                                /
/
\end{verbatim}

Note that the COMPDATL keyword in the SCHEDULE section must also be
defines for this well.

\subsection{COMPSEGS -- Define Well Connections for Multi-Segment
Wells}\label{compsegs-define-well-connections-for-multi-segment-wells}

\begin{longtable}[]{@{}
  >{\raggedright\arraybackslash}p{(\linewidth - 14\tabcolsep) * \real{0.1250}}
  >{\raggedright\arraybackslash}p{(\linewidth - 14\tabcolsep) * \real{0.1250}}
  >{\raggedright\arraybackslash}p{(\linewidth - 14\tabcolsep) * \real{0.1250}}
  >{\raggedright\arraybackslash}p{(\linewidth - 14\tabcolsep) * \real{0.1250}}
  >{\raggedright\arraybackslash}p{(\linewidth - 14\tabcolsep) * \real{0.1250}}
  >{\raggedright\arraybackslash}p{(\linewidth - 14\tabcolsep) * \real{0.1250}}
  >{\raggedright\arraybackslash}p{(\linewidth - 14\tabcolsep) * \real{0.1250}}
  >{\raggedright\arraybackslash}p{(\linewidth - 14\tabcolsep) * \real{0.1250}}@{}}
\toprule\noalign{}
\begin{minipage}[b]{\linewidth}\raggedright
RUNSPEC
\end{minipage} & \begin{minipage}[b]{\linewidth}\raggedright
GRID
\end{minipage} & \begin{minipage}[b]{\linewidth}\raggedright
EDIT
\end{minipage} & \begin{minipage}[b]{\linewidth}\raggedright
PROPS
\end{minipage} & \begin{minipage}[b]{\linewidth}\raggedright
REGIONS
\end{minipage} & \begin{minipage}[b]{\linewidth}\raggedright
SOLUTION
\end{minipage} & \begin{minipage}[b]{\linewidth}\raggedright
SUMMARY
\end{minipage} & \begin{minipage}[b]{\linewidth}\raggedright
SCHEDULE
\end{minipage} \\
\midrule\noalign{}
\endhead
\bottomrule\noalign{}
\endlastfoot
\end{longtable}

\subsubsection{Description}\label{description-891}

This keyword defines how a multi-segment well is connected to the
reservoir by allocating each well connection to a well segement. Note
that the well must have previously been defined as a multi-segment well
using the WELSEGS keyword in the SCHEDULE section and the well
connections must have previously been defined via the COMPDAT keyword in
the SCHEDULE section.

The COMPSEGS keyword should be repeated for each multi-segment well in
the model.

\begin{longtable}[]{@{}
  >{\raggedright\arraybackslash}p{(\linewidth - 6\tabcolsep) * \real{0.2500}}
  >{\raggedright\arraybackslash}p{(\linewidth - 6\tabcolsep) * \real{0.2500}}
  >{\raggedright\arraybackslash}p{(\linewidth - 6\tabcolsep) * \real{0.2500}}
  >{\raggedright\arraybackslash}p{(\linewidth - 6\tabcolsep) * \real{0.2500}}@{}}
\toprule\noalign{}
\begin{minipage}[b]{\linewidth}\raggedright
No.
\end{minipage} & \begin{minipage}[b]{\linewidth}\raggedright
Name
\end{minipage} & \begin{minipage}[b]{\linewidth}\raggedright
Description
\end{minipage} & \begin{minipage}[b]{\linewidth}\raggedright
Default
\end{minipage} \\
\midrule\noalign{}
\endhead
\bottomrule\noalign{}
\endlastfoot
Field & Metric & Laboratory & \\
1-1 & WELNAME & A character string of up to eight characters in length
that defines the well name for which a multi-segment well is being
defined. Note that the well name (WELNAME) must have been declared
previously using the WELSPECS keyword in the SCHEDULE section, otherwise
an error may occur. & None \\
1-2 & / & Record terminated by a ``/'' & Not Applicable \\
2-1 & I & A positive integer greater than or equal to one and less than
or equal to NX that defines the connection location in the I-direction.
& None \\
2-2 & J & A positive integer greater than or equal to one and less than
or equal to NY that defines the connection location in the J-direction.
& None \\
2-3 & K & A positive integer greater than or equal to one and less than
or equal to NZ that defines the connection location in the K-direction.
& None \\
2-4 & IBRANCH & A positive integer greater than or equal to one and less
than or equal to MXBRAN on WSEGDIMS keyword in the RUNSPEC section that
defines the branch number of the defined I, J and K connection. &
None \\
2-5 & LENGTH1 & A real positive value that defines the length of the
tubing from the tubing reference point (for example Measure Depth
Relative to Kelly Bushing, MDRKB) to the start of the connection in the
I, J, K cell. & None \\
feet & m & cm & \\
2-6 & LENGTH2 & A real positive value that defines the length of the
tubing from the tubing reference point to the end of the connection in
the I, J, K cell. & \\
feet & m & cm & None \\
2-7 & DIRECT & A one letter character string that defines the
orientation of the connections and should be set to either X, Y or Z, or
either I, J or K. The direction of connections may also determine the
length of the connection. Currently this option is not supported by OPM
Flow. & None \\
2-8 & IJKEND & A positive integer that is set to one of the following
depending on the orientation of the connections (DIRECT): For example,
if DIRECT is equal to Y or J then the IJKEND will be associated with the
J-direction. The value may be less than or greater than Y or J but must
be between 1 and NY. Currently this option is not supported by OPM Flow.
& None \\
2-9 & DEPTH & A real positive value that defines the datum depth for
this set of connections, normally taken as the mid-point of the
perforations associated with this set of connections. Currently this
option is not supported by OPM Flow. & None \\
feet & m & cm & \\
2-10 & THLEN & A real positive value that defines the length of the well
in the connection cells for this set of completions that is used in
thermal calculations. If this value is defaulted then the thickness of
the cell in the direction of the well orientation (DIRECT) is used.
Currently this option is not supported by OPM Flow. & None \\
feet & m & cm & \\
2-11 & ISEG & A real positive value that defines the segment number to
assign to this range of connections. If this value is defaulted then
each connection is assigned to the segment with the nodal point nearest
to the centre of the connection. Currently this option is not supported
by OPM Flow. & None \\
feet & m & cm & \\
2-12 & / & Record terminated by a ``/'' & Not Applicable \\
Notes: & & & \\
\end{longtable}

\emph{Table 12.17: COMPSEGS Keyword Description}

The maximum number of wells and completions should be defined via the
WELLDIMS keyword and the maximum number of multi-segment wells, segments
and branches should be declared on the WSEGDIMS keyword, both keywords
are in the RUNSPEC section.

See also the WELSPECS keyword to define wells, the COMPDAT keyword to
define the well connections for both standard wells and multi-segment
wells, and the COMPSEGL keyword to define multi-segment connections in a
LGR. All the aforementioned keywords are described in the SCHEDULE
section.

\subsubsection{Example}\label{example-512}

The following example defines the completions for two multi-segment oil
production wells (OP01 and OP02) using the COMPSEGS keyword.

\begin{verbatim}
--       COMPLETION SEGMENT SPECIFICATION DATA
--
-- WELL
-- NAME
COMPSEGS
OP01                                                                           /
--      --LOCATION--  BRAN  START   END     DIR  LOC    MID    COMP    ISEG
--       II  JJ  K1   NO    LENGTH  LENGTH  PEN  I,J,K  PERFS  LENGTH
         10  10   1    1    2512.5  2525.0                                     /
         10  10   2    1    2525.0  2550.0                                     /
         10  10   3    1    2550.0  2575.0                                     /
         10  10   4    1    2575.0  2600.0                                     /
         10  10   5    1    2600.0  2625.0                                     /
         10  10   6    1    2625.0  2650.0                                     /

          9  10   2    2    2637.5  2837.5                                     /
          8  10   2    2    2837.5  3037.5                                     /
          7  10   2    2    3037.5  3237.5                                     /
          6  10   2    2    3237.5  3437.5                                     /
          5  10   2    2    3437.5  3637.5                                     /
/
--       COMPLETION SEGMENT SPECIFICATION DATA
--
-- WELL
-- NAME
COMPSEGS
OP02                                                                           /
--      --LOCATION--  BRAN  START   END     DIR  LOC    MID    COMP    ISEG
--       II  JJ  K1   NO    LENGTH  LENGTH  PEN  I,J,K  PERFS  LENGTH
          1   9   3    1    2662.5  2862.5                                     /
          1   8   3    1    2862.5  3062.5                                     /
          1   7   3    1    3062.5  3262.5                                     /
          1   6   3    1    3262.5  3462.5                                     /
          1   5   3    1    3462.5  3662.5                                     /

          2  10   5    2    2712.5  2912.5                                     /
          2  10   5    2    2912.5  3112.5                                     /
          4  10   5    2    3112.5  3312.5                                     /
          5  10   5    2    3312.5  3512.5                                     /
          6  10   5    2    3512.5  3712.5                                     /

          1   9   6    3    2737.5  2937.5                                     /
          1   8   6    3    2937.5  3137.5                                     /
          1   7   6    3    3137.5  3337.5                                     /
          1   6   6    3    3337.5  3537.5                                     /
          1   5   6    3    3537.5  3737.5                                     /
/
\end{verbatim}

Note that the COMPDAT keyword in the SCHEDULE section must also be
defined for these two wells.

\subsection{COMPTRAJ -- Define Well Trajectory Connections to the
Grid}\label{comptraj-define-well-trajectory-connections-to-the-grid}

\begin{longtable}[]{@{}
  >{\raggedright\arraybackslash}p{(\linewidth - 14\tabcolsep) * \real{0.1250}}
  >{\raggedright\arraybackslash}p{(\linewidth - 14\tabcolsep) * \real{0.1250}}
  >{\raggedright\arraybackslash}p{(\linewidth - 14\tabcolsep) * \real{0.1250}}
  >{\raggedright\arraybackslash}p{(\linewidth - 14\tabcolsep) * \real{0.1250}}
  >{\raggedright\arraybackslash}p{(\linewidth - 14\tabcolsep) * \real{0.1250}}
  >{\raggedright\arraybackslash}p{(\linewidth - 14\tabcolsep) * \real{0.1250}}
  >{\raggedright\arraybackslash}p{(\linewidth - 14\tabcolsep) * \real{0.1250}}
  >{\raggedright\arraybackslash}p{(\linewidth - 14\tabcolsep) * \real{0.1250}}@{}}
\toprule\noalign{}
\begin{minipage}[b]{\linewidth}\raggedright
RUNSPEC
\end{minipage} & \begin{minipage}[b]{\linewidth}\raggedright
GRID
\end{minipage} & \begin{minipage}[b]{\linewidth}\raggedright
EDIT
\end{minipage} & \begin{minipage}[b]{\linewidth}\raggedright
PROPS
\end{minipage} & \begin{minipage}[b]{\linewidth}\raggedright
REGIONS
\end{minipage} & \begin{minipage}[b]{\linewidth}\raggedright
SOLUTION
\end{minipage} & \begin{minipage}[b]{\linewidth}\raggedright
SUMMARY
\end{minipage} & \begin{minipage}[b]{\linewidth}\raggedright
SCHEDULE
\end{minipage} \\
\midrule\noalign{}
\endhead
\bottomrule\noalign{}
\endlastfoot
\end{longtable}

\subsubsection{Description}\label{description-892}

The COMPTRAJ keyword defines how a well that has been declared as a
trajectory well, using the WELTRAJ keyword in the SCHEDULE section, is
connected to the reservoir model by defining or modifying existing well
perforation depths. The keyword can only be used for wells defined by
the WELTRAJ keyword, and WELTRAJ defined wells must use the COMPTRAJ
keyword to define the connections to the grid, that is one cannot use
COMPDAT for these type of wells.

\begin{tcolorbox}[enhanced jigsaw, title=\textcolor{quarto-callout-note-color}{\faInfo}\hspace{0.5em}{Note}, coltitle=black, breakable, toprule=.15mm, colbacktitle=quarto-callout-note-color!10!white, leftrule=.75mm, opacitybacktitle=0.6, colback=white, bottomtitle=1mm, toptitle=1mm, titlerule=0mm, left=2mm, colframe=quarto-callout-note-color-frame, arc=.35mm, rightrule=.15mm, bottomrule=.15mm, opacityback=0]

This is an OPM Flow specific keyword and will therefore cause an error
in the commercial simulator.

\end{tcolorbox}

\begin{longtable}[]{@{}
  >{\raggedright\arraybackslash}p{(\linewidth - 6\tabcolsep) * \real{0.2500}}
  >{\raggedright\arraybackslash}p{(\linewidth - 6\tabcolsep) * \real{0.2500}}
  >{\raggedright\arraybackslash}p{(\linewidth - 6\tabcolsep) * \real{0.2500}}
  >{\raggedright\arraybackslash}p{(\linewidth - 6\tabcolsep) * \real{0.2500}}@{}}
\toprule\noalign{}
\begin{minipage}[b]{\linewidth}\raggedright
No.
\end{minipage} & \begin{minipage}[b]{\linewidth}\raggedright
Name
\end{minipage} & \begin{minipage}[b]{\linewidth}\raggedright
Description
\end{minipage} & \begin{minipage}[b]{\linewidth}\raggedright
Default
\end{minipage} \\
\midrule\noalign{}
\endhead
\bottomrule\noalign{}
\endlastfoot
Field & Metric & Laboratory & \\
1 & WELNAME & A character string of up to eight characters in length
that defines the well name for which the well connection data are being
defined. Note that the well name (WELNAME) must have been declared
previously using the WELSPECS keyword in the SCHEDULE section, otherwise
an error may occur. & None \\
2 & IBRANCH & A positive integer greater than or equal to one and less
than or equal to MXBRAN on WSEGDIMS keyword in the RUNSPEC section that
defines the branch number of a segment. All segments on the main stem
must have IBRANCH set to one and lateral branches should have values
between two and MXSEGS on the WSEGDIMS keyword in the RUNSPEC section.
Only the default value of one is currently supported, that is only the
main branch of a multi-segment well is supported, or a single trajectory
for a conventional well. & 1 \\
3 ‍ & TOP & A real positive value that defines the depth of the top of
the perforation interval, and should be less than the value entered for
the BOT parameter (the bottom perforation depth). & None \\
feet & m & cm & \\
4 ‍ & BOT & A real positive value that defines the depth of the base of
the perforation interval, and should be greater than the value entered
for the TOP parameter (the top perforation interval). & None \\
feet & m & cm & \\
5 & REF & REF is a defined character string that defines the reference
depth type for TOP and BOT, and should be set to either: Only measured
depth is currently supported, that is MD. & MD \\
6 & ICOMP & An integer greater than or equal to one and less than or
equal to MXCONS as defined on the WELLDIMS keyword in the RUNSPEC
section, that defines the completion number of the currently defined
perforation interval (connection interval). If defaulted with 1*, then
ICOMP is set equal to one. & 1 \\
7 & STATUS & A defined character string of length four that defines the
connections' operational status within the perforation interval, STATUS
should be set to one of the following character strings: & OPEN \\
8 & SATNUM & An integer greater than or equal to zero and less than
NTSFUN as declared on the TABDIMS keyword in the RUNSPEC, that defines
the saturation table number to be used for flow between the reservoir
grid block and the well connections. If SATNUM is set to zero or
defaulted with 1* then: & 0 \\
9 & CONFACT & A real value greater than or equal to zero that defines
the transmissibility connection factor between the well bore and the
reservoir grid block. If set to zero or defaulted with 1* then items
(10) through (13) are used to calculate CONFACT. & Defined \\
cP.rb/day/psia 0 & cP.rm3/day/bars 0 & cP.rcc/hr/atm 0 & \\
10 & DW & A real positive value that defines the well bore diameter of
the connections for the well. DW is used in calculating a well's
productivity or injectivity index; however the value will be ignored in
calculating the connections CONFACT value if CONFAC has been directly
entered. & None \\
feet & m & cm & \\
11 & KH & A real value that defines the effective KH (permeability x
length) for the connections within the perforation interval. If less
than or equal to zero, or defaulted by 1*, then KH is calculated from
the connected grid blocks. KH is ignored if CONFACT has been directly
entered. & Calculated from connected grid blocks \\
mD.ft & mD.m & mD.cm & \\
12 & SKIN & A real value that defines the connections dimensionless skin
factor. SKIN is used in calculating a well's productivity or injectivity
index; however, the value will be ignored in calculating the connections
CONFACT value if CONFACT has been directly entered. & 0.0 \\
dimensionless & dimensionless & dimensionless & \\
13 & DFACT & A real value that defines the non-Darcy D factor
coefficient for gas wells. This value should be defaulted with 1* and
the non-Darcy D factor coefficient for gas wells defined via the WDFAC
keyword in the SCHEDULE section. Currently this option is not supported
by OPM Flow. & 1* \\
day/Mscf & day/m3 & hour/sc & \\
Notes: & & & \\
\end{longtable}

\emph{Table 12.18: COMPTRAJ Keyword Description}

Using the WELTRAJ and COMPTRAJ keywords to define wells and how they are
connected to grid, offers several advantages compared to the
conventional approach based on the (I, J, K) co-ordinates of the grid.
The approach allows for the wells to be independent of the grid, which
is particularly useful when running ensemble cases, as the well
connections are no longer required to be re-calculated for each ensemble
case. In addition, quality control of the model is improved by using
consistent perforation data in both the static and dynamic models.

\begin{tcolorbox}[enhanced jigsaw, title=\textcolor{quarto-callout-note-color}{\faInfo}\hspace{0.5em}{Note}, coltitle=black, breakable, toprule=.15mm, colbacktitle=quarto-callout-note-color!10!white, leftrule=.75mm, opacitybacktitle=0.6, colback=white, bottomtitle=1mm, toptitle=1mm, titlerule=0mm, left=2mm, colframe=quarto-callout-note-color-frame, arc=.35mm, rightrule=.15mm, bottomrule=.15mm, opacityback=0]

The term well connection is used to describe individual connections from
the wellbore to the reservoir grid, as opposed to well completions. A
well completion is used to describe a set of connections, for example, a
well may consist of several completions with each completion consisting
of multiple connections. For wells defined using the WELTRAJ and
COMPTRAJ keywords, the WELTRAJ keyword defines the trajectory of the
well within the model, and the COMPTRAJ defines the perforation
intervals in the well. A perforation interval will automatically
generate various well connections to the grid, and in addition multiple
perforation intervals may be grouped into a completion.

\end{tcolorbox}

\subsubsection{Example}\label{example-513}

The following example defines two trajectory wells oil wells, OP01 and
OP02, using the WELTRAJ keyword, together with their perforations using
the COMPTRAJ keyword.

\begin{verbatim}
--
--       WELL SPECIFICATION DATA
--
-- WELL  GROUP     LOCATION  BHP    PHASE  DRAIN  INFLOW  OPEN  CROSS  PVT
-- NAME  NAME        I    J  DEPTH  FLUID  AREA   EQUANS  SHUT  FLOW   TABLE
WELSPECS
OP01     PLATFORM   1*   1*   1*     OIL    1*     STD    SHUT   NO     1*     /
OP02     PLATFORM   1*   1*   1*     OIL    1*     STD    SHUT   NO     1*     /
/
--
--       WELL TRAJECTORY DATA
--
-- WELL   BRAN  XCORD         YCORD         TVDSS         MD
-- NAME   NO                                DEPTH         DEPTH
-- -----  ----  ------------  ------------  ------------  ------------
WELTRAJ
OP01      1*    2.805445e+06  3.602948e+06  -100.000000   0.0         /
OP01      1*    2.805445e+06  3.602948e+06  877.0000000   977.0       /
OP01      1*    2.805445e+06  3.602948e+06  957.9950240   1058.0      /
OP01      1*    2.805444e+06  3.602946e+06  1051.976081   1152.0      /
…………………...                                                            /
OP02      1*    2.810828e+06  3.604507e+06  9371.792711   11418.0     /
OP02      1*    2.810885e+06  3.604525e+06  9443.657000   11511.0     /
OP02      1*    2.810952e+06  3.604546e+06  9531.966162   11624.0     /
OP02      1*    2.810973e+06  3.604553e+06  9560.411742   11660.0     /
/
--
--
--       WELL TRAJECTORY CONNECTION DATA
--
-- WELL  BRAN  -- PERFORATION --  COMPL OPEN  SAT CONN  WELL   KH    SKIN  D
-- NAME  NO.    TOP    BOT   REF  NO.   SHUT  TAB FACT  DIA    FACT  FACT  FACT
COMPTRAJ
OP01     1*     8230   8244  MD    1    SHUT  1*  1*    0.708  1*    0.0   1*  /
OP01     1*     8352   8380  MD    1    SHUT  1*  1*    0.708  1*    0.0   1*  /
OP01     1*     9070   9100  MD    1    SHUT  1*  1*    0.708  1*    0.0   1*  /
OP01     1*     9220   9250  MD    2    SHUT  1*  1*    0.708  1*    0.0   1*  /
OP01     1*     9266   9280  MD    2    SHUT  1*  1*    0.708  1*    0.0   1*  /
OP01     1*     9693   9703  MD    3    SHUT  1*  1*    0.708  1*    0.0   1*  /
OP01     1*     9940   9974  MD    3    SHUT  1*  1*    0.708  1*    0.0   1*  /
OP02     1*     9979   9985  TVD   1*   SHUT  1*  1*    0.708  1*    0.0   1*  /
OP02     1*    10173  10183  TVD   1*   SHUT  1*  1*    0.708  1*    0.0   1*  /
OP02     1*    10190  10204  TVD   1*   SHUT  1*  1*    0.708  1*    0.0   1*  /
OP02     1*    10327  10333  TVD   1*   SHUT  1*  1*    0.708  1*    0.0   1*  /
OP02     1*    10339  10345  TVD   1*   SHUT  1*  1*    0.708  1*    0.0   1*  /
OP02     1*    11528  11538  TVD   1*   SHUT  1*  1*    0.708  1*    0.0   1*  /
/
\end{verbatim}

Well OP01 has eight perforation intervals, with the intervals one to
three grouped into one completion, perforation intervals four to five
grouped into completion number two, and finally the bottom three
perforations are grouped into completion number three. In contrast, OP02
has six perforated intervals with their completion interval defaulted to
one.

\subsection{COMPVE -- Re-Define Well Connection
Depths}\label{compve-re-define-well-connection-depths}

\begin{longtable}[]{@{}
  >{\raggedright\arraybackslash}p{(\linewidth - 14\tabcolsep) * \real{0.1250}}
  >{\raggedright\arraybackslash}p{(\linewidth - 14\tabcolsep) * \real{0.1250}}
  >{\raggedright\arraybackslash}p{(\linewidth - 14\tabcolsep) * \real{0.1250}}
  >{\raggedright\arraybackslash}p{(\linewidth - 14\tabcolsep) * \real{0.1250}}
  >{\raggedright\arraybackslash}p{(\linewidth - 14\tabcolsep) * \real{0.1250}}
  >{\raggedright\arraybackslash}p{(\linewidth - 14\tabcolsep) * \real{0.1250}}
  >{\raggedright\arraybackslash}p{(\linewidth - 14\tabcolsep) * \real{0.1250}}
  >{\raggedright\arraybackslash}p{(\linewidth - 14\tabcolsep) * \real{0.1250}}@{}}
\toprule\noalign{}
\begin{minipage}[b]{\linewidth}\raggedright
RUNSPEC
\end{minipage} & \begin{minipage}[b]{\linewidth}\raggedright
GRID
\end{minipage} & \begin{minipage}[b]{\linewidth}\raggedright
EDIT
\end{minipage} & \begin{minipage}[b]{\linewidth}\raggedright
PROPS
\end{minipage} & \begin{minipage}[b]{\linewidth}\raggedright
REGIONS
\end{minipage} & \begin{minipage}[b]{\linewidth}\raggedright
SOLUTION
\end{minipage} & \begin{minipage}[b]{\linewidth}\raggedright
SUMMARY
\end{minipage} & \begin{minipage}[b]{\linewidth}\raggedright
SCHEDULE
\end{minipage} \\
\midrule\noalign{}
\endhead
\bottomrule\noalign{}
\endlastfoot
\end{longtable}

\subsubsection{Description}\label{description-893}

The COMPVE keyword is used to re-define the well connection depths in
the global grid.

This keyword is not supported by OPM Flow but would change the results
if supported so the simulation will be stopped.

\subsection{COMPVEL -- Re-Define Well LGR Connection
Depths}\label{compvel-re-define-well-lgr-connection-depths}

\begin{longtable}[]{@{}
  >{\raggedright\arraybackslash}p{(\linewidth - 14\tabcolsep) * \real{0.1250}}
  >{\raggedright\arraybackslash}p{(\linewidth - 14\tabcolsep) * \real{0.1250}}
  >{\raggedright\arraybackslash}p{(\linewidth - 14\tabcolsep) * \real{0.1250}}
  >{\raggedright\arraybackslash}p{(\linewidth - 14\tabcolsep) * \real{0.1250}}
  >{\raggedright\arraybackslash}p{(\linewidth - 14\tabcolsep) * \real{0.1250}}
  >{\raggedright\arraybackslash}p{(\linewidth - 14\tabcolsep) * \real{0.1250}}
  >{\raggedright\arraybackslash}p{(\linewidth - 14\tabcolsep) * \real{0.1250}}
  >{\raggedright\arraybackslash}p{(\linewidth - 14\tabcolsep) * \real{0.1250}}@{}}
\toprule\noalign{}
\begin{minipage}[b]{\linewidth}\raggedright
RUNSPEC
\end{minipage} & \begin{minipage}[b]{\linewidth}\raggedright
GRID
\end{minipage} & \begin{minipage}[b]{\linewidth}\raggedright
EDIT
\end{minipage} & \begin{minipage}[b]{\linewidth}\raggedright
PROPS
\end{minipage} & \begin{minipage}[b]{\linewidth}\raggedright
REGIONS
\end{minipage} & \begin{minipage}[b]{\linewidth}\raggedright
SOLUTION
\end{minipage} & \begin{minipage}[b]{\linewidth}\raggedright
SUMMARY
\end{minipage} & \begin{minipage}[b]{\linewidth}\raggedright
SCHEDULE
\end{minipage} \\
\midrule\noalign{}
\endhead
\bottomrule\noalign{}
\endlastfoot
\end{longtable}

\subsubsection{Description}\label{description-894}

The COMPVEL keyword is used to re-define the well connection depths in a
Local Grid Refinement (``LGR'') grid.

This keyword is not supported by OPM Flow but would change the results
if supported so the simulation will be stopped.

\subsection{CPIFACT -- Define Well Connection Transmissibility
Multipliers}\label{cpifact-define-well-connection-transmissibility-multipliers}

\begin{longtable}[]{@{}
  >{\raggedright\arraybackslash}p{(\linewidth - 14\tabcolsep) * \real{0.1250}}
  >{\raggedright\arraybackslash}p{(\linewidth - 14\tabcolsep) * \real{0.1250}}
  >{\raggedright\arraybackslash}p{(\linewidth - 14\tabcolsep) * \real{0.1250}}
  >{\raggedright\arraybackslash}p{(\linewidth - 14\tabcolsep) * \real{0.1250}}
  >{\raggedright\arraybackslash}p{(\linewidth - 14\tabcolsep) * \real{0.1250}}
  >{\raggedright\arraybackslash}p{(\linewidth - 14\tabcolsep) * \real{0.1250}}
  >{\raggedright\arraybackslash}p{(\linewidth - 14\tabcolsep) * \real{0.1250}}
  >{\raggedright\arraybackslash}p{(\linewidth - 14\tabcolsep) * \real{0.1250}}@{}}
\toprule\noalign{}
\begin{minipage}[b]{\linewidth}\raggedright
RUNSPEC
\end{minipage} & \begin{minipage}[b]{\linewidth}\raggedright
GRID
\end{minipage} & \begin{minipage}[b]{\linewidth}\raggedright
EDIT
\end{minipage} & \begin{minipage}[b]{\linewidth}\raggedright
PROPS
\end{minipage} & \begin{minipage}[b]{\linewidth}\raggedright
REGIONS
\end{minipage} & \begin{minipage}[b]{\linewidth}\raggedright
SOLUTION
\end{minipage} & \begin{minipage}[b]{\linewidth}\raggedright
SUMMARY
\end{minipage} & \begin{minipage}[b]{\linewidth}\raggedright
SCHEDULE
\end{minipage} \\
\midrule\noalign{}
\endhead
\bottomrule\noalign{}
\endlastfoot
\end{longtable}

\subsubsection{Description}\label{description-895}

The CPIFACT keyword is used to define well connection transmissibility
multipliers for well connections to the global grid.

This keyword is not supported by OPM Flow but would change the results
if supported so the simulation will be stopped.

\subsection{CPIFACTL -- Define Well Connection Transmissibility
Multipliers in a
LGR}\label{cpifactl-define-well-connection-transmissibility-multipliers-in-a-lgr}

\begin{longtable}[]{@{}
  >{\raggedright\arraybackslash}p{(\linewidth - 14\tabcolsep) * \real{0.1250}}
  >{\raggedright\arraybackslash}p{(\linewidth - 14\tabcolsep) * \real{0.1250}}
  >{\raggedright\arraybackslash}p{(\linewidth - 14\tabcolsep) * \real{0.1250}}
  >{\raggedright\arraybackslash}p{(\linewidth - 14\tabcolsep) * \real{0.1250}}
  >{\raggedright\arraybackslash}p{(\linewidth - 14\tabcolsep) * \real{0.1250}}
  >{\raggedright\arraybackslash}p{(\linewidth - 14\tabcolsep) * \real{0.1250}}
  >{\raggedright\arraybackslash}p{(\linewidth - 14\tabcolsep) * \real{0.1250}}
  >{\raggedright\arraybackslash}p{(\linewidth - 14\tabcolsep) * \real{0.1250}}@{}}
\toprule\noalign{}
\begin{minipage}[b]{\linewidth}\raggedright
RUNSPEC
\end{minipage} & \begin{minipage}[b]{\linewidth}\raggedright
GRID
\end{minipage} & \begin{minipage}[b]{\linewidth}\raggedright
EDIT
\end{minipage} & \begin{minipage}[b]{\linewidth}\raggedright
PROPS
\end{minipage} & \begin{minipage}[b]{\linewidth}\raggedright
REGIONS
\end{minipage} & \begin{minipage}[b]{\linewidth}\raggedright
SOLUTION
\end{minipage} & \begin{minipage}[b]{\linewidth}\raggedright
SUMMARY
\end{minipage} & \begin{minipage}[b]{\linewidth}\raggedright
SCHEDULE
\end{minipage} \\
\midrule\noalign{}
\endhead
\bottomrule\noalign{}
\endlastfoot
\end{longtable}

\subsubsection{Description}\label{description-896}

The CPIFACT keyword is used to define well Local Grid Refinement
(``LGR'') connection transmissibility multipliers for well connections
to a LGR, for when the LGR option has been invoked by the LGR keyword in
the RUNSPEC section.

This keyword is not supported by OPM Flow but would change the results
if supported so the simulation will be stopped.

\subsection{CSKIN -- Re-Define Well Connection Skin
Factors}\label{cskin-re-define-well-connection-skin-factors}

\begin{longtable}[]{@{}
  >{\raggedright\arraybackslash}p{(\linewidth - 14\tabcolsep) * \real{0.1250}}
  >{\raggedright\arraybackslash}p{(\linewidth - 14\tabcolsep) * \real{0.1250}}
  >{\raggedright\arraybackslash}p{(\linewidth - 14\tabcolsep) * \real{0.1250}}
  >{\raggedright\arraybackslash}p{(\linewidth - 14\tabcolsep) * \real{0.1250}}
  >{\raggedright\arraybackslash}p{(\linewidth - 14\tabcolsep) * \real{0.1250}}
  >{\raggedright\arraybackslash}p{(\linewidth - 14\tabcolsep) * \real{0.1250}}
  >{\raggedright\arraybackslash}p{(\linewidth - 14\tabcolsep) * \real{0.1250}}
  >{\raggedright\arraybackslash}p{(\linewidth - 14\tabcolsep) * \real{0.1250}}@{}}
\toprule\noalign{}
\begin{minipage}[b]{\linewidth}\raggedright
RUNSPEC
\end{minipage} & \begin{minipage}[b]{\linewidth}\raggedright
GRID
\end{minipage} & \begin{minipage}[b]{\linewidth}\raggedright
EDIT
\end{minipage} & \begin{minipage}[b]{\linewidth}\raggedright
PROPS
\end{minipage} & \begin{minipage}[b]{\linewidth}\raggedright
REGIONS
\end{minipage} & \begin{minipage}[b]{\linewidth}\raggedright
SOLUTION
\end{minipage} & \begin{minipage}[b]{\linewidth}\raggedright
SUMMARY
\end{minipage} & \begin{minipage}[b]{\linewidth}\raggedright
SCHEDULE
\end{minipage} \\
\midrule\noalign{}
\endhead
\bottomrule\noalign{}
\endlastfoot
\end{longtable}

\subsubsection{Description}\label{description-897}

This keyword, CSKIN, is used to re-define a well's connection skin
factors and as such will result in the well's connection
transmissibility factors being updated accordingly. The well connections
must already be defined using the COMPDAT keyword in the SCHEDULE
section before this keyword can be used.

\begin{longtable}[]{@{}
  >{\raggedright\arraybackslash}p{(\linewidth - 6\tabcolsep) * \real{0.2500}}
  >{\raggedright\arraybackslash}p{(\linewidth - 6\tabcolsep) * \real{0.2500}}
  >{\raggedright\arraybackslash}p{(\linewidth - 6\tabcolsep) * \real{0.2500}}
  >{\raggedright\arraybackslash}p{(\linewidth - 6\tabcolsep) * \real{0.2500}}@{}}
\toprule\noalign{}
\begin{minipage}[b]{\linewidth}\raggedright
No.
\end{minipage} & \begin{minipage}[b]{\linewidth}\raggedright
Name
\end{minipage} & \begin{minipage}[b]{\linewidth}\raggedright
Description
\end{minipage} & \begin{minipage}[b]{\linewidth}\raggedright
Default
\end{minipage} \\
\midrule\noalign{}
\endhead
\bottomrule\noalign{}
\endlastfoot
Field & Metric & Laboratory & \\
1 & WELNAME & A character string of up to eight characters in length
that defines the well name for which the well connection skin is being
re-defined. Note that the well name (WELNAME) must have been declared
previously using the WELSPECS keyword in the SCHEDULE section, otherwise
an error may occur. & None \\
2 & I & A positive integer greater than or equal to one and less than or
equal to NX that defines the connection location in the I-direction. If
set to zero or defaulted with 1* then connections in any I-direction
location will be modified. & 0 \\
3 & J & A positive integer greater than or equal to one and less than or
equal to NY that defines the connection location in the J-direction. If
set to zero or defaulted with 1* then connections in any J-direction
location will be modified. & 0 \\
4 & K1 & A positive integer greater than or equal to one and less than
or equal to K2 and NZ that defines the UPPER connection location in the
K-direction. If set to zero or defaulted with 1* then this will be taken
from the top connection of the well. & 0 \\
5 & K2 & A positive integer greater than or equal to K1 and less than or
equal to NZ that defines the LOWER connection location in the
K-direction. If set to zero or defaulted with 1* then this will be taken
from the bottom connection of the well. & 0 \\
6 & SKIN & A real value that defines the connections' dimensionless skin
factors. & 0.0 \\
dimensionless & dimensionless & dimensionless & \\
Notes: & & & \\
\end{longtable}

\emph{Table 12.3.43.1: CSKIN Keyword Description}

\subsubsection{Example}\label{example-514}

The following example re-defines the skin factor for all well OP01
connections in layers 4 to 6.

\begin{verbatim}
--
--       CONNECTION SKIN DATA
--
-- WELL  --- LOCATION ---  SKIN
-- NAME   II  JJ  K1  K2   FACT
CSKIN
OP01      1*  1*   4   6   2.0    /
/
\end{verbatim}

\subsection{DATES -- Advance Simulation by Reporting
Date}\label{dates-advance-simulation-by-reporting-date}

\begin{longtable}[]{@{}
  >{\raggedright\arraybackslash}p{(\linewidth - 14\tabcolsep) * \real{0.1250}}
  >{\raggedright\arraybackslash}p{(\linewidth - 14\tabcolsep) * \real{0.1250}}
  >{\raggedright\arraybackslash}p{(\linewidth - 14\tabcolsep) * \real{0.1250}}
  >{\raggedright\arraybackslash}p{(\linewidth - 14\tabcolsep) * \real{0.1250}}
  >{\raggedright\arraybackslash}p{(\linewidth - 14\tabcolsep) * \real{0.1250}}
  >{\raggedright\arraybackslash}p{(\linewidth - 14\tabcolsep) * \real{0.1250}}
  >{\raggedright\arraybackslash}p{(\linewidth - 14\tabcolsep) * \real{0.1250}}
  >{\raggedright\arraybackslash}p{(\linewidth - 14\tabcolsep) * \real{0.1250}}@{}}
\toprule\noalign{}
\begin{minipage}[b]{\linewidth}\raggedright
RUNSPEC
\end{minipage} & \begin{minipage}[b]{\linewidth}\raggedright
GRID
\end{minipage} & \begin{minipage}[b]{\linewidth}\raggedright
EDIT
\end{minipage} & \begin{minipage}[b]{\linewidth}\raggedright
PROPS
\end{minipage} & \begin{minipage}[b]{\linewidth}\raggedright
REGIONS
\end{minipage} & \begin{minipage}[b]{\linewidth}\raggedright
SOLUTION
\end{minipage} & \begin{minipage}[b]{\linewidth}\raggedright
SUMMARY
\end{minipage} & \begin{minipage}[b]{\linewidth}\raggedright
SCHEDULE
\end{minipage} \\
\midrule\noalign{}
\endhead
\bottomrule\noalign{}
\endlastfoot
\end{longtable}

\subsubsection{Description}\label{description-898}

This keyword advances the simulation to a given report date after which
additional keywords may be entered to instruct OPM Flow to perform
additional functions via the SCHEDULE section keywords, or further DATES
data sets or keywords may be entered to advance the simulator to the
next report date.

If the DATES keyword is to be used during the simulation, then the START
keyword in the RUNSPEC section must be declared to set the start date
for the run.

\begin{longtable}[]{@{}
  >{\raggedright\arraybackslash}p{(\linewidth - 6\tabcolsep) * \real{0.2500}}
  >{\raggedright\arraybackslash}p{(\linewidth - 6\tabcolsep) * \real{0.2500}}
  >{\raggedright\arraybackslash}p{(\linewidth - 6\tabcolsep) * \real{0.2500}}
  >{\raggedright\arraybackslash}p{(\linewidth - 6\tabcolsep) * \real{0.2500}}@{}}
\toprule\noalign{}
\begin{minipage}[b]{\linewidth}\raggedright
No.
\end{minipage} & \begin{minipage}[b]{\linewidth}\raggedright
Name
\end{minipage} & \begin{minipage}[b]{\linewidth}\raggedright
Description
\end{minipage} & \begin{minipage}[b]{\linewidth}\raggedright
Default
\end{minipage} \\
\midrule\noalign{}
\endhead
\bottomrule\noalign{}
\endlastfoot
1 & DAY & A positive integer that defines the day of the month for the
data set, the value should be greater than or equal to one and less than
or equal to 31. & None \\
2 & MONTH & Character string for the month for the data set and should
be one of the following `JAN', `FEB', `MAR', `APR', `MAY', `JUN', `JUL'
(or `JLY'), `AUG', `SEP', `OCT', `NOV', or `DEC' & None \\
3 & YEAR & A positive four digit integer value representing the year for
the data set, which must be specified fully by four digits, that is
1986. & None \\
4 & TIME & A numeric character string that defines the time for the data
set in the form of: HH:MM:SS.SSSS The default value means in most cases
this parameter can be defaulted. TIME is normally used when detailed DST
matching is performed to enable the pressures and rates to be stated at
specific dates and times. & 00:00:00 \\
Notes: & & & \\
\end{longtable}

\emph{Table 12.19: DATES Keyword Description}

Note that OPM Flow uses the standard Gregorian calendar and therefore
leap years are accounted for in the DATES keyword. Thus, it is more
accurate to use the DATES keyword to progress the simulator through time
if one is matching actual production data.

See also the TIME and TSTEP keywords in the SCHEDULE section.

Whenever possible it is a good idea to always set the start date to be
at the beginning of the year, as like most simulators, OPM Flow reports
are always stated at the number of days from the start date (and
sometimes at a given date). If the start date is at the beginning of the
year, then calculating the actual date is relatively straight forward
and simple.

\subsubsection{Examples}\label{examples-70}

Given a start date of January 1, 2020 set via the START keyword in the
RUNSPEC section, the following example advances the simulator from the
start date of January 1, 2020 to January 1, 2021, using quarterly
reporting time steps.

\begin{verbatim}
-- ==============================================================================
--
-- SCHEDULE SECTION
--
-- ==============================================================================
SCHEDULE

-- ------------------------------------------------------------------------------
-- SCHEDULE SECTION - 2020-01-01
-- ------------------------------------------------------------------------------
RPTSCHED
          'WELLS=2'    'WELSPECS'    'CPU=2'     FIP=2'                         /

DATES
          2  JAN   2020  /
/

RPTSCHED
          'NOTHING'                                                             /

DATES
          1  APR   2020  /
          1  JLY   2020  /
          1  OCT   2020  /
/
-- ------------------------------------------------------------------------------
-- SCHEDULE SECTION - 2021-01-01
-- ------------------------------------------------------------------------------
RPTSCHED
          'WELLS=2'    'WELSPECS'    'CPU=2'     FIP=2'                         /

DATES
          1  JAN   2021  /
/
RPTSCHED
          'NOTHING'                                                             /

DATES
          1  APR   2021  /
          1  JLY   2021  /
          1  OCT   2021  /
/
\end{verbatim}

The above example writes out a series of report at the start of the run
and then advances the simulation one day to January 2, 2020 and switches
off the reporting. The simulation then advances to April 1, July 1 and
October 1, 2020 with no further changes to the run. After October 1,
2020 reporting is switched on again to enable a report on January 1,
2021, which is then subsequently switched off after the January 1, 2021
report time step has been completed.

Note if one wishes to terminate the run at the end of year (as opposes
to the beginning of the year and get a final report for the year, then
the next example demonstrates the keyword sequence to enable this.

\begin{verbatim}
-- ------------------------------------------------------------------------------
-- SCHEDULE SECTION - 2021-01-01
-- ------------------------------------------------------------------------------
RPTSCHED
          'WELLS=2'    'WELSPECS'    'CPU=2'     FIP=2'                         /

DATES
          2  JAN   2021  /
/

RPTSCHED
         'NOTHING'                                                             /

DATES
          1  FEB   2021  /
          1  MAR   2021  /
          1  APR   2021  /
          1  MAY   2021  /
          1  JUN   2021  /
          1  JLY   2021  /
          1  AUG   2021  /
          1  SEP   2021  /
          1  OCT   2021  /
          1  NOV   2021  /
          1  DEC   2021  /
/
--
-- FINAL REPORT AND RESTART AT YEAR END
--
RPTSCHED
          'WELLS=2'    'WELSPECS'    'CPU=2'     FIP=2'                         /

RPTRST
          'BASIC=2'                                                             /

DATES
           31 DEC   2021  /
/
\end{verbatim}

In the above example monthly reporting time steps have been used instead
of quarterly and report is requested after the December 1, 2021 time
step and is therefore written out on December 31, 2021.

\subsection{DCQDEFN -- Define Gas DCQ Units as Rate or
Energy}\label{dcqdefn-define-gas-dcq-units-as-rate-or-energy}

\begin{longtable}[]{@{}
  >{\raggedright\arraybackslash}p{(\linewidth - 14\tabcolsep) * \real{0.1250}}
  >{\raggedright\arraybackslash}p{(\linewidth - 14\tabcolsep) * \real{0.1250}}
  >{\raggedright\arraybackslash}p{(\linewidth - 14\tabcolsep) * \real{0.1250}}
  >{\raggedright\arraybackslash}p{(\linewidth - 14\tabcolsep) * \real{0.1250}}
  >{\raggedright\arraybackslash}p{(\linewidth - 14\tabcolsep) * \real{0.1250}}
  >{\raggedright\arraybackslash}p{(\linewidth - 14\tabcolsep) * \real{0.1250}}
  >{\raggedright\arraybackslash}p{(\linewidth - 14\tabcolsep) * \real{0.1250}}
  >{\raggedright\arraybackslash}p{(\linewidth - 14\tabcolsep) * \real{0.1250}}@{}}
\toprule\noalign{}
\begin{minipage}[b]{\linewidth}\raggedright
RUNSPEC
\end{minipage} & \begin{minipage}[b]{\linewidth}\raggedright
GRID
\end{minipage} & \begin{minipage}[b]{\linewidth}\raggedright
EDIT
\end{minipage} & \begin{minipage}[b]{\linewidth}\raggedright
PROPS
\end{minipage} & \begin{minipage}[b]{\linewidth}\raggedright
REGIONS
\end{minipage} & \begin{minipage}[b]{\linewidth}\raggedright
SOLUTION
\end{minipage} & \begin{minipage}[b]{\linewidth}\raggedright
SUMMARY
\end{minipage} & \begin{minipage}[b]{\linewidth}\raggedright
SCHEDULE
\end{minipage} \\
\midrule\noalign{}
\endhead
\bottomrule\noalign{}
\endlastfoot
\end{longtable}

\subsubsection{Description}\label{description-899}

The DCQDEFN keyword defines the DCQ units to be rate or energy
(calorific value) when using the Gas Field Operation model and the Gas
Calorific Value control option. The gas DCQ rates are controlled by the
GASYEAR, GASPERIO, GDCQ, GASFTARG or GASFDECR keywords in the SCHEDULE
section.

This keyword is not supported by OPM Flow but would change the results
if supported so the simulation will be stopped.

\subsection{DEBUG -- Define the Debug Data to be Printed to
File}\label{debug-define-the-debug-data-to-be-printed-to-file-8}

This keyword defines the debug data to be written to the debug file
(*.DBG). This keyword is not supported by OPM Flow but has no effect on
the results so it will be ignored.

See DEBUG -- Define the Debug Data to be Printed to File in the GLOBAL
section for a full description.

\subsection{DELAYACT -- Define Delayed Action
Keywords}\label{delayact-define-delayed-action-keywords}

\begin{longtable}[]{@{}
  >{\raggedright\arraybackslash}p{(\linewidth - 14\tabcolsep) * \real{0.1250}}
  >{\raggedright\arraybackslash}p{(\linewidth - 14\tabcolsep) * \real{0.1250}}
  >{\raggedright\arraybackslash}p{(\linewidth - 14\tabcolsep) * \real{0.1250}}
  >{\raggedright\arraybackslash}p{(\linewidth - 14\tabcolsep) * \real{0.1250}}
  >{\raggedright\arraybackslash}p{(\linewidth - 14\tabcolsep) * \real{0.1250}}
  >{\raggedright\arraybackslash}p{(\linewidth - 14\tabcolsep) * \real{0.1250}}
  >{\raggedright\arraybackslash}p{(\linewidth - 14\tabcolsep) * \real{0.1250}}
  >{\raggedright\arraybackslash}p{(\linewidth - 14\tabcolsep) * \real{0.1250}}@{}}
\toprule\noalign{}
\begin{minipage}[b]{\linewidth}\raggedright
RUNSPEC
\end{minipage} & \begin{minipage}[b]{\linewidth}\raggedright
GRID
\end{minipage} & \begin{minipage}[b]{\linewidth}\raggedright
EDIT
\end{minipage} & \begin{minipage}[b]{\linewidth}\raggedright
PROPS
\end{minipage} & \begin{minipage}[b]{\linewidth}\raggedright
REGIONS
\end{minipage} & \begin{minipage}[b]{\linewidth}\raggedright
SOLUTION
\end{minipage} & \begin{minipage}[b]{\linewidth}\raggedright
SUMMARY
\end{minipage} & \begin{minipage}[b]{\linewidth}\raggedright
SCHEDULE
\end{minipage} \\
\midrule\noalign{}
\endhead
\bottomrule\noalign{}
\endlastfoot
\end{longtable}

\subsubsection{Description}\label{description-900}

The DELAYACT keyword defines a series of keywords that should be
executed after an ACTION keyword has been triggered by the ACTION,
ACTIONG, ACTIONR, ACTIONW, ACTIONS, or ACTIONX keywords.

This keyword is not supported by OPM Flow but would change the results
if supported so the simulation will be stopped.

\subsection{DIFFMMF -- Define Diffusivity Multipliers for
Matrix-Fractures}\label{diffmmf-define-diffusivity-multipliers-for-matrix-fractures-1}

This keyword, DIFFMMF, defines the diffusivity multipliers for
matrix-fractures for when the Dual Porosity option has be activated by
either the DUALPORO or DUALPERM keywords, or the Coal Bed Methane option
is selected by the COAL keyword, and the Diffusivity option has been
activated by the DIFFUSE keywords; all four keywords are in the RUNSPEC
section.

See DIFFMMF -- Define Diffusivity Multipliers for Matrix-Fractures in
the GRID section for a full description.

\subsection{DIMPES -- Define IMPES Dynamic Solution
Parameters}\label{dimpes-define-impes-dynamic-solution-parameters}

\begin{longtable}[]{@{}
  >{\raggedright\arraybackslash}p{(\linewidth - 14\tabcolsep) * \real{0.1250}}
  >{\raggedright\arraybackslash}p{(\linewidth - 14\tabcolsep) * \real{0.1250}}
  >{\raggedright\arraybackslash}p{(\linewidth - 14\tabcolsep) * \real{0.1250}}
  >{\raggedright\arraybackslash}p{(\linewidth - 14\tabcolsep) * \real{0.1250}}
  >{\raggedright\arraybackslash}p{(\linewidth - 14\tabcolsep) * \real{0.1250}}
  >{\raggedright\arraybackslash}p{(\linewidth - 14\tabcolsep) * \real{0.1250}}
  >{\raggedright\arraybackslash}p{(\linewidth - 14\tabcolsep) * \real{0.1250}}
  >{\raggedright\arraybackslash}p{(\linewidth - 14\tabcolsep) * \real{0.1250}}@{}}
\toprule\noalign{}
\begin{minipage}[b]{\linewidth}\raggedright
RUNSPEC
\end{minipage} & \begin{minipage}[b]{\linewidth}\raggedright
GRID
\end{minipage} & \begin{minipage}[b]{\linewidth}\raggedright
EDIT
\end{minipage} & \begin{minipage}[b]{\linewidth}\raggedright
PROPS
\end{minipage} & \begin{minipage}[b]{\linewidth}\raggedright
REGIONS
\end{minipage} & \begin{minipage}[b]{\linewidth}\raggedright
SOLUTION
\end{minipage} & \begin{minipage}[b]{\linewidth}\raggedright
SUMMARY
\end{minipage} & \begin{minipage}[b]{\linewidth}\raggedright
SCHEDULE
\end{minipage} \\
\midrule\noalign{}
\endhead
\bottomrule\noalign{}
\endlastfoot
\end{longtable}

\subsubsection{Description}\label{description-901}

This keyword, DIMPES, defines the Implicit in Pressure Explicit
Saturation (``IMPES'') dynamic solution parameters and results in the
simulator switching from the current solution formulation to the IMPES
formulation. OPM Flow uses a different numerical scheme which makes this
keyword redundant; hence, OPM Flow ignores this keyword. It is
documented here for completeness.

There is no data required for this keyword.

See section 2.2 Running OPM Flow 2023-04 From The Command Line on how to
invoke various numerical schemes via the OPM Flow command line
interface.

\subsubsection{Example}\label{example-515}

\begin{verbatim}
--
--       ACTIVATE THE IMPES SOLUTION OPTION
--
DIMPES
\end{verbatim}

The above example switches on the IMPES solution option; however, this
has no effect in OPM Flow input decks.

\subsection{DIMPLICT -- Activate Fully Implicit Dynamic Solution
Formulation}\label{dimplict-activate-fully-implicit-dynamic-solution-formulation}

\begin{longtable}[]{@{}
  >{\raggedright\arraybackslash}p{(\linewidth - 14\tabcolsep) * \real{0.1250}}
  >{\raggedright\arraybackslash}p{(\linewidth - 14\tabcolsep) * \real{0.1250}}
  >{\raggedright\arraybackslash}p{(\linewidth - 14\tabcolsep) * \real{0.1250}}
  >{\raggedright\arraybackslash}p{(\linewidth - 14\tabcolsep) * \real{0.1250}}
  >{\raggedright\arraybackslash}p{(\linewidth - 14\tabcolsep) * \real{0.1250}}
  >{\raggedright\arraybackslash}p{(\linewidth - 14\tabcolsep) * \real{0.1250}}
  >{\raggedright\arraybackslash}p{(\linewidth - 14\tabcolsep) * \real{0.1250}}
  >{\raggedright\arraybackslash}p{(\linewidth - 14\tabcolsep) * \real{0.1250}}@{}}
\toprule\noalign{}
\begin{minipage}[b]{\linewidth}\raggedright
RUNSPEC
\end{minipage} & \begin{minipage}[b]{\linewidth}\raggedright
GRID
\end{minipage} & \begin{minipage}[b]{\linewidth}\raggedright
EDIT
\end{minipage} & \begin{minipage}[b]{\linewidth}\raggedright
PROPS
\end{minipage} & \begin{minipage}[b]{\linewidth}\raggedright
REGIONS
\end{minipage} & \begin{minipage}[b]{\linewidth}\raggedright
SOLUTION
\end{minipage} & \begin{minipage}[b]{\linewidth}\raggedright
SUMMARY
\end{minipage} & \begin{minipage}[b]{\linewidth}\raggedright
SCHEDULE
\end{minipage} \\
\midrule\noalign{}
\endhead
\bottomrule\noalign{}
\endlastfoot
\end{longtable}

\subsubsection{Description}\label{description-902}

This keyword, DIMPLICT, activates the Fully Implicit Formulation and
results in the simulator switching from the current solution formulation
to the fully implicit formulation. OPM Flow uses a different numerical
scheme which makes this keyword redundant; hence, OPM Flow ignores this
keyword. It is documented here for completeness.

There is no data required for this keyword.

See section 2.2 Running OPM Flow 2023-04 From The Command Line on how to
invoke various numerical schemes via the OPM Flow command line
interface.

\subsubsection{Example}\label{example-516}

\begin{verbatim}
--
--       ACTIVATES THE FULLY IMPLICIT SOLUTION OPTION
--
DIMPLICT
\end{verbatim}

The above example switches on the fully implicit solution option;
however, this has no effect in OPM Flow input decks.

\subsection{DRILPRI -- Define Prioritized Drilling Queue Priority
Parameters}\label{drilpri-define-prioritized-drilling-queue-priority-parameters}

\begin{longtable}[]{@{}
  >{\raggedright\arraybackslash}p{(\linewidth - 14\tabcolsep) * \real{0.1250}}
  >{\raggedright\arraybackslash}p{(\linewidth - 14\tabcolsep) * \real{0.1250}}
  >{\raggedright\arraybackslash}p{(\linewidth - 14\tabcolsep) * \real{0.1250}}
  >{\raggedright\arraybackslash}p{(\linewidth - 14\tabcolsep) * \real{0.1250}}
  >{\raggedright\arraybackslash}p{(\linewidth - 14\tabcolsep) * \real{0.1250}}
  >{\raggedright\arraybackslash}p{(\linewidth - 14\tabcolsep) * \real{0.1250}}
  >{\raggedright\arraybackslash}p{(\linewidth - 14\tabcolsep) * \real{0.1250}}
  >{\raggedright\arraybackslash}p{(\linewidth - 14\tabcolsep) * \real{0.1250}}@{}}
\toprule\noalign{}
\begin{minipage}[b]{\linewidth}\raggedright
RUNSPEC
\end{minipage} & \begin{minipage}[b]{\linewidth}\raggedright
GRID
\end{minipage} & \begin{minipage}[b]{\linewidth}\raggedright
EDIT
\end{minipage} & \begin{minipage}[b]{\linewidth}\raggedright
PROPS
\end{minipage} & \begin{minipage}[b]{\linewidth}\raggedright
REGIONS
\end{minipage} & \begin{minipage}[b]{\linewidth}\raggedright
SOLUTION
\end{minipage} & \begin{minipage}[b]{\linewidth}\raggedright
SUMMARY
\end{minipage} & \begin{minipage}[b]{\linewidth}\raggedright
SCHEDULE
\end{minipage} \\
\midrule\noalign{}
\endhead
\bottomrule\noalign{}
\endlastfoot
\end{longtable}

\subsubsection{Description}\label{description-903}

This keyword, DRILPRI, defines the prioritized drilling queue priority
parameters used in the priority formulae for this type of drilling
queue.

This keyword is not supported by OPM Flow but would change the results
if supported so the simulation will be stopped.

\subsection{DRSDT -- Solution Gas (Rs) Maximum Rate of Increase
Parameters}\label{drsdt-solution-gas-rs-maximum-rate-of-increase-parameters}

\begin{longtable}[]{@{}
  >{\raggedright\arraybackslash}p{(\linewidth - 14\tabcolsep) * \real{0.1250}}
  >{\raggedright\arraybackslash}p{(\linewidth - 14\tabcolsep) * \real{0.1250}}
  >{\raggedright\arraybackslash}p{(\linewidth - 14\tabcolsep) * \real{0.1250}}
  >{\raggedright\arraybackslash}p{(\linewidth - 14\tabcolsep) * \real{0.1250}}
  >{\raggedright\arraybackslash}p{(\linewidth - 14\tabcolsep) * \real{0.1250}}
  >{\raggedright\arraybackslash}p{(\linewidth - 14\tabcolsep) * \real{0.1250}}
  >{\raggedright\arraybackslash}p{(\linewidth - 14\tabcolsep) * \real{0.1250}}
  >{\raggedright\arraybackslash}p{(\linewidth - 14\tabcolsep) * \real{0.1250}}@{}}
\toprule\noalign{}
\begin{minipage}[b]{\linewidth}\raggedright
RUNSPEC
\end{minipage} & \begin{minipage}[b]{\linewidth}\raggedright
GRID
\end{minipage} & \begin{minipage}[b]{\linewidth}\raggedright
EDIT
\end{minipage} & \begin{minipage}[b]{\linewidth}\raggedright
PROPS
\end{minipage} & \begin{minipage}[b]{\linewidth}\raggedright
REGIONS
\end{minipage} & \begin{minipage}[b]{\linewidth}\raggedright
SOLUTION
\end{minipage} & \begin{minipage}[b]{\linewidth}\raggedright
SUMMARY
\end{minipage} & \begin{minipage}[b]{\linewidth}\raggedright
SCHEDULE
\end{minipage} \\
\midrule\noalign{}
\endhead
\bottomrule\noalign{}
\endlastfoot
\end{longtable}

\subsubsection{Description}\label{description-904}

DRSDT defines the maximum rate at which the solution gas-oil ratio (Rs)
can be increased in a grid cell. The keyword is similar in functionality
to the DRSDTR keyword, that defines the maximum rate at which Rs can be
increased in a grid cell by region. Both keywords should only be used if
the OIL, GAS, and DISGAS keywords in the RUNSPEC section have been
invoked to allow oil, gas and dissolved gas to be present in the model.
The keyword only affects the behavior of an increasing Rs, for example
when gas is being injected into an oil reservoir, and is subject to the
availability of free gas and the ability of the undersaturated oil to
adsorb this gas.

\begin{longtable}[]{@{}
  >{\raggedright\arraybackslash}p{(\linewidth - 6\tabcolsep) * \real{0.2500}}
  >{\raggedright\arraybackslash}p{(\linewidth - 6\tabcolsep) * \real{0.2500}}
  >{\raggedright\arraybackslash}p{(\linewidth - 6\tabcolsep) * \real{0.2500}}
  >{\raggedright\arraybackslash}p{(\linewidth - 6\tabcolsep) * \real{0.2500}}@{}}
\toprule\noalign{}
\begin{minipage}[b]{\linewidth}\raggedright
No.
\end{minipage} & \begin{minipage}[b]{\linewidth}\raggedright
Name
\end{minipage} & \begin{minipage}[b]{\linewidth}\raggedright
Description
\end{minipage} & \begin{minipage}[b]{\linewidth}\raggedright
Default
\end{minipage} \\
\midrule\noalign{}
\endhead
\bottomrule\noalign{}
\endlastfoot
Field & Metric & Laboratory & \\
1 & DRSDT1 & DRSDT1 is a real positive number that defines the maximum
rate at which the solution gas-oil ratio is allowed to increase in a
grid cell, that is the maximum rate the gas can dissolve into the
available undersaturated oil. A value of zero means that Rs cannot
increase and free gas cannot dissolve into the unsaturated oil in a grid
cell. Alternatively a very large value of DRSDT1 allows Rs to increase
rapidly until there is no free gas or the oil within the grid block is
fully saturated. Note if the keyword is not present in the input deck
then DRSDT1 is assumed to be a very large number resulting in complete
re-solution of the gas into the available undersaturated oil. & None \\
Mscf/stb/d & sm3/sm3/day & scc/scc/day & \\
2 & DRSDT2 & DRSDT2 is a defined character string that defines whether
the DRSDT1 is applied to either all grid blocks or just those grid
blocks containing free gas: Note if the keyword is not present in the
input deck then DRSDT2 is set to the default value of ALL. & ALL \\
Notes: & & & \\
\end{longtable}

\emph{Table 12.20: DRSDT Keyword Description}

Note this keyword can be used in history matching field performance to
control the availability of the movable gas phase.

See also the VAPPARS keyword in the SOLUTION section and the DRSDTR,
DRVDT and DRVDTR keywords in the SCHEDULE section that controls how
vaporized oil is treated and the rate at which the dissolved phase ratio
increases within a grid block.

\subsubsection{Examples}\label{examples-71}

The first example prevents the solution gas-oil ratio from increasing
and applies this to all grid cells.

\begin{verbatim}
--
--       SOLUTION GAS (RS) MAXIMUM RATE OF INCREASE FOR MODEL                                            --
DRSDT
--       MAX RS    ALL/FREE
--       DRSDT1    DRSDT2
--       -------   --------
         0.000     ALL                                     /
\end{verbatim}

And the second example below applies 0.0005 Mscf/stb/d as the maximum
rate at which the solution gas-oil ratio is allowed to increase in a
grid cell, and applies this to only cells containing free gas.

\begin{verbatim}
--
--       SOLUTION GAS (RS) MAXIMUM RATE OF INCREASE FOR MODEL                                            --
DRSDT
--       MAX RS    ALL/FREE
--       DRSDT1    DRSDT2
--       -------   --------
         0.0005    FREE                                    /
\end{verbatim}

Again, the keyword parameters when applied are subject to the
availability of free gas and the ability of the undersaturated oil to
adsorb this gas.

\subsection{DRSDTCON -- CO2 Convective Dissolution
Parameters}\label{drsdtcon-co2-convective-dissolution-parameters}

\begin{longtable}[]{@{}
  >{\raggedright\arraybackslash}p{(\linewidth - 14\tabcolsep) * \real{0.1250}}
  >{\raggedright\arraybackslash}p{(\linewidth - 14\tabcolsep) * \real{0.1250}}
  >{\raggedright\arraybackslash}p{(\linewidth - 14\tabcolsep) * \real{0.1250}}
  >{\raggedright\arraybackslash}p{(\linewidth - 14\tabcolsep) * \real{0.1250}}
  >{\raggedright\arraybackslash}p{(\linewidth - 14\tabcolsep) * \real{0.1250}}
  >{\raggedright\arraybackslash}p{(\linewidth - 14\tabcolsep) * \real{0.1250}}
  >{\raggedright\arraybackslash}p{(\linewidth - 14\tabcolsep) * \real{0.1250}}
  >{\raggedright\arraybackslash}p{(\linewidth - 14\tabcolsep) * \real{0.1250}}@{}}
\toprule\noalign{}
\begin{minipage}[b]{\linewidth}\raggedright
RUNSPEC
\end{minipage} & \begin{minipage}[b]{\linewidth}\raggedright
GRID
\end{minipage} & \begin{minipage}[b]{\linewidth}\raggedright
EDIT
\end{minipage} & \begin{minipage}[b]{\linewidth}\raggedright
PROPS
\end{minipage} & \begin{minipage}[b]{\linewidth}\raggedright
REGIONS
\end{minipage} & \begin{minipage}[b]{\linewidth}\raggedright
SOLUTION
\end{minipage} & \begin{minipage}[b]{\linewidth}\raggedright
SUMMARY
\end{minipage} & \begin{minipage}[b]{\linewidth}\raggedright
SCHEDULE
\end{minipage} \\
\midrule\noalign{}
\endhead
\bottomrule\noalign{}
\endlastfoot
\end{longtable}

\subsubsection{Description}\label{description-905}

The DRSDTCON keyword defines the parameters that control the convective
dissolution of carbon dioxide (CO2) into the in-situ brine within a grid
cell, as described by Mykkeltvedt et al.~{[}Mykkeltvedt, T.S., Sandve,
T.H. \& Gasda, S.E. New Sub-grid Model for Convective Mixing in
Field-Scale CO2 Storage Simulation. Transp Porous Med 152, 9 (2025).
https://doi.org/10.1007/s11242-024-02141-5{]}, based on an assumption of
vertical equilibrium. The keyword internally causes the simulator to
calculate the solution gas-oil ratio (Rs), normally controlled by the
DRSDT1 parameter on the DRSDT keyword in the SCHEDULE section, making
the DRSDT keyword redundant.

The keyword should only be used if the CO2STORE keyword in the RUNSPEC
section has been activated to model CO2 storage via OPM Flow's CO2-Brine
PVT model. Note that the CO2STORE keyword must be used together with
either: (1) the GAS, WATER and DISGASW keywords (or alternatively the
GASWAT and DISGASW keywords), or (2) the GAS, OIL and DISGAS keywords in
the RUNSPEC section. See the CO2STORE keyword for details.

The DRSDTCON keyword is an OPM Flow specific keyword.

\begin{longtable}[]{@{}
  >{\raggedright\arraybackslash}p{(\linewidth - 6\tabcolsep) * \real{0.2500}}
  >{\raggedright\arraybackslash}p{(\linewidth - 6\tabcolsep) * \real{0.2500}}
  >{\raggedright\arraybackslash}p{(\linewidth - 6\tabcolsep) * \real{0.2500}}
  >{\raggedright\arraybackslash}p{(\linewidth - 6\tabcolsep) * \real{0.2500}}@{}}
\toprule\noalign{}
\begin{minipage}[b]{\linewidth}\raggedright
No.
\end{minipage} & \begin{minipage}[b]{\linewidth}\raggedright
Name
\end{minipage} & \begin{minipage}[b]{\linewidth}\raggedright
Description
\end{minipage} & \begin{minipage}[b]{\linewidth}\raggedright
Default
\end{minipage} \\
\midrule\noalign{}
\endhead
\bottomrule\noalign{}
\endlastfoot
Field & Metric & Laboratory & \\
1 & CHI & A real positive value () that defines the proportionality
constant related to the maximum rate of increase of CO2 solution gas-oil
ratio (Rs) in the Linear regime. A value of zero means that convective
dissolution of CO2 into in-situ brine does not occur and free CO2 cannot
dissolve into the brine. Alternatively, a non-zero value of CHI allows
convective dissolution of CO2. Note if the CO2STORE keyword is present
but the DRSDTCON keyword is absent from the input deck, then this
results in instantaneous dissolution of CO2 into the available
undersaturated in-situ brine. & 0.04 \\
dimensionless & dimensionless & dimensionless & \\
1 & PSI & A real positive value () that defines the normalised CO2
solution gas-oil ratio (\(\overset{~}{X}\)) at the transition between
the Linear and the Steady-State regimes. & 0.34 \\
dimensionless & dimensionless & dimensionless & \\
1 & OMEGA & A real positive value () that defines the maximum rate of
increase in CO2 solution gas-oil ratio (Rs) in the Steady-State regime.
& 3.0e-9 \\
1/s & 1/s & 1/s & \\
1 & OPTION & A defined character string that specifies in which cells
the convective dissolution rate limit is applied, and should be set to
one of the following character strings: & ALL \\
Notes: & & & \\
\end{longtable}

\emph{Table 12.3.53.1: DRSDTCON Keyword Description}

Mykkeltvedt et al.~describe four regimes that characterise the evolution
of the CO2 solution gas-oil ratio in a control volume:

\begin{itemize}
\tightlist
\item
  Initial Phase: Near instantaneous equilibrium partitioning and
  diffusion.
\item
  Linear Phase: Fingers of CO2-rich brine propagate downwards towards
  the base of the control volume.
\item
  Steady-State Phase: Fingers have passed through the base of the
  control volume.
\item
  Decline Phase: Gas phase CO2 has been completely dissolved in the
  brine.
\end{itemize}

In the Initial phase a near instantaneous jump in dissolved CO2 is
assumed if the control volume spans the CO2-brine interface (otherwise
the solution gas-oil ratio is zero). The initial jump comes from the
capillary transition zone, the height of which depends on the capillary
pressure function. Capillary equilibrium happens almost immediately and
therefore the CO2 solution gas-oil ratio is allowed to increase at
unlimited rate until capillary equilibrium is reached. Capillary
equilibrium is assumed when:

\begin{equation}\phantomsection\label{eq-12-3-53-1}{
{R}_{s} > {R}_{s,\mathit{sat}}{S}_{o}
}\end{equation}

Where:

Rs = solution gas-oil ratio,

Rs,sat = saturated solution gas-oil ratio,

So = oil saturation.

This is followed by the Linear phase where the finger speed is assumed
to be relatively constant and therefore it is reasonable to assume that
the CO2 gas-oil ratio will increase at a constant rate (Linear regime).
The maximum dissolution rate, as per the DRSDT keyword in the SCHEDULE
section during the Linear regime (Flin) is given by:

\begin{equation}\phantomsection\label{eq-12-3-53-2}{
{F}_{\mathit{lin}} = χ\left(\frac{{R}_{s,\mathit{sat}}{K}_{z}\mathrm{Δ}{ρ}_{c}g}{{μ}_{o}{S}_{o} {D}_{z}ϕ}\right)
}\end{equation}

Where:

χ = the non-dimensional parameter CHI controlling the dissolution rate,

Rs,sat = the maximum gas-oil ratio at the CO2 solubility limit in the
brine,

Kz = the vertical permeability,

∆ρc = the difference between the brine density saturated with CO2 at
Rs,sat and

the brine density without dissolved CO2,

g = the gravity constant,

μo = the oil (brine) viscosity, and

So = the oil (brine) saturation.

The third phase is characterised by a change in slope of the rate of
dissolution as the system transitions from linear build-up to
quasi-steady state (Steady-State regime). The parameter PSI defines the
normalised gas-oil ratio at the transition from the Linear to the
Steady-State regime, where the normalised gas-oil ratio is given by:

\begin{equation}\phantomsection\label{eq-12-3-53-3}{
\overset{~}{X} = \frac{{R}_{s}-{R}_{s,\mathit{sat}}{S}_{g}}{{R}_{s,\mathit{sat}}(1-{S}_{g})}
}\end{equation}

The rate of dissolution during the Steady-State regime is defined by the
parameter OMEGA.

The final Decline phase occurs when the CO2 has been completely
dissolved in the brine at which point the rate of dissolution is zero.

The values of CHI, PSI and OMEGA can either be estimated from fine-scale
simulation, historical data, or from laboratory tests {[}Taheri, A.,
Torsæter, O., Lindeberg, E., Hadia, N. J., \& Wessel-Berg, D. (2018).
Qualitative and quantitative experimental study of convective mixing
process during storage of CO2 in heterogeneous saline aquifers.
International Journal of Greenhouse Gas Control, 71, 212-226.{]}.
Analysis of fine-scale simulation results by Mykkeltvedt et al indicates
that 0.04 +/-0.01 is a reasonable value for CHI. Elenius et
al.~{[}Elenius, M. T., Nordbotten, J. M., \& Kalisch, H. (2014).
Convective mixing influenced by the capillary transition zone.
Computational Geosciences, 18(3-4), 417-431{]} also state that 0.04 is a
reasonable value for CHI for the Utsira formation.

\begin{tcolorbox}[enhanced jigsaw, title=\textcolor{quarto-callout-note-color}{\faInfo}\hspace{0.5em}{Note}, coltitle=black, breakable, toprule=.15mm, colbacktitle=quarto-callout-note-color!10!white, leftrule=.75mm, opacitybacktitle=0.6, colback=white, bottomtitle=1mm, toptitle=1mm, titlerule=0mm, left=2mm, colframe=quarto-callout-note-color-frame, arc=.35mm, rightrule=.15mm, bottomrule=.15mm, opacityback=0]

In the commercial simulator a constant or regional value for DRSDT can
be given as an input parameter. This can be used to include the effect
of convective mixing as shown by Thibeau and Dutin {[}Thibeau, S., \&
Dutin, A. (2011). Large scale CO2 storage in unstructured aquifers:
Modeling study of the ultimate CO2 migration distance. Energy Procedia,
4, 4230-4237.{]}. OPM Flow's approach differs from this in that the
DRSDT value is computed internally and is dependent on both the static
and dynamic cell properties.

\end{tcolorbox}

\subsubsection{Example}\label{example-517}

The example below is similar to that shown under the CO2STORE keyword in
the RUNSPEC section. In the RUNSPEC section one declares that the carbon
dioxide storage model is active for the run to account for both carbon
dioxide and water phase solubility using OPM Flow's CO2-Brine PVT model.

\begin{verbatim}
-- ==============================================================================
--
-- RUNSPEC SECTION
--
-- ==============================================================================
RUNSPEC
-- ------------------------------------------------------------------------------
--  FLUID TYPES AND TRACER OPTIONS
-- ------------------------------------------------------------------------------
--
--       ACTIVATE CO2 STORAGE IN THE MODEL (OPM FLOW CO2 STORAGE KEYWORD)
--
CO2STORE
--
--       ACTIVATE GAS-WATER THE MODEL (OPM FLOW KEYWORD)
--
GASWAT
--
--       DISSOLVED GAS IN WATER IS PRESENT IN THE RUN (OPM FLOW KEYWORD)
--
DISGASW
--
--       VAPORIZED WATER IN DRY/WET GAS IS PRESENT IN THE RUN (OPM FLOW KEYWORD)
--
VAPWAT
\end{verbatim}

The second part of the example sets the maximum dissolution rate for
convective CO2 mixing via the DRSDTCON keyword in the SCHEDULE section
using the base case parameters calculated by Mykkeltvedt et al from
fine-scale simulations.

\begin{verbatim}
-- ==============================================================================
--
-- SCHEDULE SECTION
--
-- ==============================================================================
SCHEDULE
--
--       CO2 CONVECTIVE DISSOLUTION PARAMETERS
--
DRSDTCON
--       CHI      PSI      OMEGA    OPTION
         0.04     0.34     3.0E-09  ALL                             /
\end{verbatim}

See also the CO2STORE keyword in the RUNSPEC section for further
information on OPM Flow's CO2 storage facility.

\subsection{DRSDTR -- Solution Gas (Rs) Maximum Rate of Increase
Parameters by
Region}\label{drsdtr-solution-gas-rs-maximum-rate-of-increase-parameters-by-region}

\begin{longtable}[]{@{}
  >{\raggedright\arraybackslash}p{(\linewidth - 14\tabcolsep) * \real{0.1250}}
  >{\raggedright\arraybackslash}p{(\linewidth - 14\tabcolsep) * \real{0.1250}}
  >{\raggedright\arraybackslash}p{(\linewidth - 14\tabcolsep) * \real{0.1250}}
  >{\raggedright\arraybackslash}p{(\linewidth - 14\tabcolsep) * \real{0.1250}}
  >{\raggedright\arraybackslash}p{(\linewidth - 14\tabcolsep) * \real{0.1250}}
  >{\raggedright\arraybackslash}p{(\linewidth - 14\tabcolsep) * \real{0.1250}}
  >{\raggedright\arraybackslash}p{(\linewidth - 14\tabcolsep) * \real{0.1250}}
  >{\raggedright\arraybackslash}p{(\linewidth - 14\tabcolsep) * \real{0.1250}}@{}}
\toprule\noalign{}
\begin{minipage}[b]{\linewidth}\raggedright
RUNSPEC
\end{minipage} & \begin{minipage}[b]{\linewidth}\raggedright
GRID
\end{minipage} & \begin{minipage}[b]{\linewidth}\raggedright
EDIT
\end{minipage} & \begin{minipage}[b]{\linewidth}\raggedright
PROPS
\end{minipage} & \begin{minipage}[b]{\linewidth}\raggedright
REGIONS
\end{minipage} & \begin{minipage}[b]{\linewidth}\raggedright
SOLUTION
\end{minipage} & \begin{minipage}[b]{\linewidth}\raggedright
SUMMARY
\end{minipage} & \begin{minipage}[b]{\linewidth}\raggedright
SCHEDULE
\end{minipage} \\
\midrule\noalign{}
\endhead
\bottomrule\noalign{}
\endlastfoot
\end{longtable}

\subsubsection{Description}\label{description-906}

DRSDTR defines the maximum rate at which the solution gas-oil ratio (Rs)
can be increased in a grid cell for various regions in the model. The
keyword is similar in functionality to the DRSDT keyword, that defines
the maximum rate at which Rs can be increased in a grid cell for all
cells in the model. The number of DRSDTR vector data sets is defined by
the NTPVT parameter on the TABDIMS keyword in the RUNSPEC section and
the allocation of the DRSDTR records to different grid blocks in the
model is done via the PVTNUM keyword in the REGION section. One data set
consists of one record or line which is terminated by a ``/''.

DRSDTR should only be used if the OIL, GAS, and DISGAS keywords in the
RUNSPEC section have been invoked to allow oil, gas and dissolved gas to
be present in the model. The keyword only affects the behavior of an
increasing Rs, for example when gas is being injected into an oil
reservoir, and is subject to the availability of free gas and the
ability of the undersaturated oil to adsorb this gas.

\begin{longtable}[]{@{}
  >{\raggedright\arraybackslash}p{(\linewidth - 6\tabcolsep) * \real{0.2500}}
  >{\raggedright\arraybackslash}p{(\linewidth - 6\tabcolsep) * \real{0.2500}}
  >{\raggedright\arraybackslash}p{(\linewidth - 6\tabcolsep) * \real{0.2500}}
  >{\raggedright\arraybackslash}p{(\linewidth - 6\tabcolsep) * \real{0.2500}}@{}}
\toprule\noalign{}
\begin{minipage}[b]{\linewidth}\raggedright
No.
\end{minipage} & \begin{minipage}[b]{\linewidth}\raggedright
Name
\end{minipage} & \begin{minipage}[b]{\linewidth}\raggedright
Description
\end{minipage} & \begin{minipage}[b]{\linewidth}\raggedright
Default
\end{minipage} \\
\midrule\noalign{}
\endhead
\bottomrule\noalign{}
\endlastfoot
Field & Metric & Laboratory & \\
1 & DRSDT1 & DRSDT1 is a real positive number that defines the maximum
rate at which the solution gas-oil ratio is allowed to increase in a
grid cell, that is the maximum rate the gas can dissolve into the
available undersaturated oil. A value of zero means that Rs cannot
increase and free gas cannot dissolve into the unsaturated oil in a grid
cell. Alternatively a very large value of DRSDT1 allows Rs to increase
rapidly until there is no free gas or the oil within the grid block is
fully saturated. Note if the keyword is not present in the input deck
then DRSDT1 is assumed to be a very large number resulting in complete
re-solution of the gas into the available undersaturated oil. & None \\
Mscf/stb/d & sm3/sm3/day & scc/scc/day & \\
2 & DRSDT2 & DRSDT2 is a defined character string that defines whether
the DRSDT1 is applied to either all grid blocks or just those grid
blocks containing free gas: Note if the keyword is not present in the
input deck then DRSDT2 is set to the default value of ALL. & ALL \\
Notes: & & & \\
\end{longtable}

\emph{Table 12.22: DRSDTR Keyword Description}

Note this keyword can be used in history matching field performance to
control the availability of the movable gas phase.

See also the VAPPARS keyword in the SOLUTION section and the DRSDT,
DRVDT and DRVDTR keywords in the SCHEDULE section that controls how
vaporized oil is treated and the rate at which the dissolved phase ratio
increases within a grid block.

This keyword is not supported by OPM Flow but has no effect on the
results so it will be ignored but is documented here for completeness,
as it is expected to be available in the next release of OPM Flow.

\subsubsection{Examples}\label{examples-72}

The first example prevents the solution gas-oil ratio from increasing
and applies this to all regions for when NTPVT on the TABDIMS keyword in
the RUNSPEC section is set to three.

\begin{verbatim}
--
--       SOLUTION GAS (RS) MAXIMUM RATE OF INCREASE BY REGION                                            --
DRSDTR
--       MAX RS    ALL/FREE
--       DRSDT1    DRSDT2
--       -------   --------
         0.0000    ALL                                     /
         0.0000    ALL                                     /
         0.0000    ALL                                     /
\end{verbatim}

The second example below prevents the solution gas-oil ratio from
increasing and applies this to all grid cells in PVTNUM region one. For
PVTNUM regions one and two the keyword applies 0.0005 Mscf/stb/d as the
maximum rate at which the solution gas-oil ratio is allowed to increase
in a grid cell, and applies this to only cells containing free gas.

\begin{verbatim}
--
--       SOLUTION GAS (RS) MAXIMUM RATE OF INCREASE BY REGION                                            --
DRSDTR
--       MAX RS    ALL/FREE
--       DRSDT1    DRSDT2
--       -------   --------
         0.0000    ALL                                     /
         0.0005    FREE                                    /
         0.0005    FREE                                    /
\end{verbatim}

Again, the keyword parameters when applied are subject to the
availability of free gas and the ability of the undersaturated oil to
adsorb this gas.

\subsection{DRVDT -- Vaporized Oil (Rv) Maximum Rate of Increase
Parameters}\label{drvdt-vaporized-oil-rv-maximum-rate-of-increase-parameters}

\begin{longtable}[]{@{}
  >{\raggedright\arraybackslash}p{(\linewidth - 14\tabcolsep) * \real{0.1250}}
  >{\raggedright\arraybackslash}p{(\linewidth - 14\tabcolsep) * \real{0.1250}}
  >{\raggedright\arraybackslash}p{(\linewidth - 14\tabcolsep) * \real{0.1250}}
  >{\raggedright\arraybackslash}p{(\linewidth - 14\tabcolsep) * \real{0.1250}}
  >{\raggedright\arraybackslash}p{(\linewidth - 14\tabcolsep) * \real{0.1250}}
  >{\raggedright\arraybackslash}p{(\linewidth - 14\tabcolsep) * \real{0.1250}}
  >{\raggedright\arraybackslash}p{(\linewidth - 14\tabcolsep) * \real{0.1250}}
  >{\raggedright\arraybackslash}p{(\linewidth - 14\tabcolsep) * \real{0.1250}}@{}}
\toprule\noalign{}
\begin{minipage}[b]{\linewidth}\raggedright
RUNSPEC
\end{minipage} & \begin{minipage}[b]{\linewidth}\raggedright
GRID
\end{minipage} & \begin{minipage}[b]{\linewidth}\raggedright
EDIT
\end{minipage} & \begin{minipage}[b]{\linewidth}\raggedright
PROPS
\end{minipage} & \begin{minipage}[b]{\linewidth}\raggedright
REGIONS
\end{minipage} & \begin{minipage}[b]{\linewidth}\raggedright
SOLUTION
\end{minipage} & \begin{minipage}[b]{\linewidth}\raggedright
SUMMARY
\end{minipage} & \begin{minipage}[b]{\linewidth}\raggedright
SCHEDULE
\end{minipage} \\
\midrule\noalign{}
\endhead
\bottomrule\noalign{}
\endlastfoot
\end{longtable}

\subsubsection{Description}\label{description-907}

DRVDT defines the maximum rate at which the vaporized oil-gas ratio or
condensate-gas ratio (Rv) can be increased in a grid cell. The keyword
is similar in functionality to the DRVDTR keyword, that defines the
maximum rate at which Rv can be increased in a grid cell by region. Both
keywords should only be used if the OIL, GAS, and VAPOIL (condensate)
keywords in the RUNSPEC section have been invoked to allow oil, gas and
condensate to be present in the model. The keyword only affects the
behavior of an increasing Rv, for example when gas is being injected
into a gas condensate reservoir as part of as gas re-cycling scheme, and
is subject to the availability of free oil (condensate) and the ability
of the undersaturated gas to adsorb this condensate.

\begin{longtable}[]{@{}
  >{\raggedright\arraybackslash}p{(\linewidth - 6\tabcolsep) * \real{0.2500}}
  >{\raggedright\arraybackslash}p{(\linewidth - 6\tabcolsep) * \real{0.2500}}
  >{\raggedright\arraybackslash}p{(\linewidth - 6\tabcolsep) * \real{0.2500}}
  >{\raggedright\arraybackslash}p{(\linewidth - 6\tabcolsep) * \real{0.2500}}@{}}
\toprule\noalign{}
\begin{minipage}[b]{\linewidth}\raggedright
No.
\end{minipage} & \begin{minipage}[b]{\linewidth}\raggedright
Name
\end{minipage} & \begin{minipage}[b]{\linewidth}\raggedright
Description
\end{minipage} & \begin{minipage}[b]{\linewidth}\raggedright
Default
\end{minipage} \\
\midrule\noalign{}
\endhead
\bottomrule\noalign{}
\endlastfoot
Field & Metric & Laboratory & \\
1 & DRVDT1 & DRVDT1 is a real positive number that defines the maximum
rate at which the vaporized oil-gas ratio is allowed to increase in a
grid cell, that is the maximum rate at which the oil can dissolve into
the available undersaturated gas. A value of zero means that Rv cannot
increase and free oil cannot vaporize into the unsaturated gas in a grid
cell. Alternatively a very large value of DRVDT1 allows Rv to increase
rapidly until there is no free oil or the gas within the grid block is
fully saturated. Note if the keyword is not present in the input deck
then DRVDT1 is assumed to be a very large number resulting in complete
re-vaporization of the oil into the available undersaturated gas. &
None \\
stb/Mscf/d & sm3/sm3/day & scc/scc/day & \\
Notes: & & & \\
\end{longtable}

\emph{Table 12.23: DRVDT Keyword Description}

Note this keyword can be used in history matching field performance to
control the availability of the movable gas phase. See also the VAPPARS
keyword in the SOLUTION section and the DRVDTR, DRSDT and DRSDTR
keywords in the SCHEDULE section that controls how solution gas is
treated and the rate at which the vaporized phase ratio increases within
a grid block.

\subsubsection{Example}\label{example-518}

The example prevents the solution oil-gas ratio from increasing.

\begin{verbatim}
--
--       VAPORIZED OIL (RV) MAXIMUM RATE OF INCREASE FOR MODEL                                           --
DRVDT
--       MAX RV
--       DRVDT1
--       -------
         0.000                                             /
\end{verbatim}

Again, the keyword parameters when applied are subject to the
availability of free oil and the ability of the undersaturated gas to
adsorb this oil.

\subsection{DRVDTR -- Vaporized Oil (Rv) Maximum Rate of Increase
Parameters by
Region}\label{drvdtr-vaporized-oil-rv-maximum-rate-of-increase-parameters-by-region}

\begin{longtable}[]{@{}
  >{\raggedright\arraybackslash}p{(\linewidth - 14\tabcolsep) * \real{0.1250}}
  >{\raggedright\arraybackslash}p{(\linewidth - 14\tabcolsep) * \real{0.1250}}
  >{\raggedright\arraybackslash}p{(\linewidth - 14\tabcolsep) * \real{0.1250}}
  >{\raggedright\arraybackslash}p{(\linewidth - 14\tabcolsep) * \real{0.1250}}
  >{\raggedright\arraybackslash}p{(\linewidth - 14\tabcolsep) * \real{0.1250}}
  >{\raggedright\arraybackslash}p{(\linewidth - 14\tabcolsep) * \real{0.1250}}
  >{\raggedright\arraybackslash}p{(\linewidth - 14\tabcolsep) * \real{0.1250}}
  >{\raggedright\arraybackslash}p{(\linewidth - 14\tabcolsep) * \real{0.1250}}@{}}
\toprule\noalign{}
\begin{minipage}[b]{\linewidth}\raggedright
RUNSPEC
\end{minipage} & \begin{minipage}[b]{\linewidth}\raggedright
GRID
\end{minipage} & \begin{minipage}[b]{\linewidth}\raggedright
EDIT
\end{minipage} & \begin{minipage}[b]{\linewidth}\raggedright
PROPS
\end{minipage} & \begin{minipage}[b]{\linewidth}\raggedright
REGIONS
\end{minipage} & \begin{minipage}[b]{\linewidth}\raggedright
SOLUTION
\end{minipage} & \begin{minipage}[b]{\linewidth}\raggedright
SUMMARY
\end{minipage} & \begin{minipage}[b]{\linewidth}\raggedright
SCHEDULE
\end{minipage} \\
\midrule\noalign{}
\endhead
\bottomrule\noalign{}
\endlastfoot
\end{longtable}

\subsubsection{Description}\label{description-908}

DRVDTR defines the maximum rate at which the vaporized oil-gas ratio or
condensate-gas ratio (Rv) can be increased in a grid cell for various
regions in the model. The keyword is similar in functionality to the
DRVDT keyword, that defines the maximum rate at which Rv can be
increased in a grid cell for all cells in the model. The number of
DRVDTR vector data sets is defined by the NTPVT parameter on the TABDIMS
keyword in the RUNSPEC section and the allocation of the DRVDTR records
to different grid blocks in the model is done via the PVTNUM keyword in
the REGION section. One data set consists of one record or line which is
terminated by a ``/''.

This keyword should only be used if the OIL, GAS, and VAPOIL
(condensate) keywords in the RUNSPEC section have been invoked to allow
oil, gas and condensate to be present in the model. The keyword only
affects the behavior of an increasing Rv, for example when gas is being
injected into a gas condensate reservoir as part of as gas re-cycling
scheme, and is subject to the availability of free oil (condensate) and
the ability of the undersaturated gas to adsorb this condensate.

\begin{longtable}[]{@{}
  >{\raggedright\arraybackslash}p{(\linewidth - 6\tabcolsep) * \real{0.2500}}
  >{\raggedright\arraybackslash}p{(\linewidth - 6\tabcolsep) * \real{0.2500}}
  >{\raggedright\arraybackslash}p{(\linewidth - 6\tabcolsep) * \real{0.2500}}
  >{\raggedright\arraybackslash}p{(\linewidth - 6\tabcolsep) * \real{0.2500}}@{}}
\toprule\noalign{}
\begin{minipage}[b]{\linewidth}\raggedright
No.
\end{minipage} & \begin{minipage}[b]{\linewidth}\raggedright
Name
\end{minipage} & \begin{minipage}[b]{\linewidth}\raggedright
Description
\end{minipage} & \begin{minipage}[b]{\linewidth}\raggedright
Default
\end{minipage} \\
\midrule\noalign{}
\endhead
\bottomrule\noalign{}
\endlastfoot
Field & Metric & Laboratory & \\
1 & DRVDT1 & DRVDT1 is a real positive number that defines the maximum
rate at which the vaporized oil-gas ratio is allowed to increase in a
grid cell, that is the maximum rate at which the oil can vaporize into
the available undersaturated gas. A value of zero means that Rv cannot
increase and free oil cannot vaporize into the unsaturated gas in a grid
cell. Alternatively a very large value of DRVDT1 allows Rv to increase
rapidly until there is no free oil or the gas within the grid block is
fully saturated. Note if the keyword is not present in the input deck
then DRVDT1 is assumed to be a very large number resulting in complete
re-vaporization of the oil into the available undersaturated gas. &
None \\
stb/Mscf/d & sm3/sm3/day & scc/scc/day & \\
Notes: & & & \\
\end{longtable}

\emph{Table 12.24: DRVDTR Keyword Description}

Note this keyword can be used in history matching field performance to
control the availability of the movable gas phase.

See also the VAPPARS keyword in the SOLUTION section and DRVDT, DRSDT,
and DRSDTR keywords in the SCHEDULE section that controls how dissolved
gas is treated and the rate at which the vaporized phase ratio increases
within a grid block.

\subsubsection{Examples}\label{examples-73}

The first example prevents the vaporized oil-gas ratio from increasing
and applies this to all regions for when NTPVT on the TABDIMS keyword in
the RUNSPEC section is set to three.

\begin{verbatim}
--
--       VAPORIZED OIL (RV) MAXIMUM RATE OF INCREASE PARAMETERS BY REGION
--
DRVDTR
--       MAX RV
--       DRVDT1
--       -------          -
         0.000                                             /
         0.000                                             /
         0.000                                             /
\end{verbatim}

The second example below prevents the vaporized oil-gas ratio from
increasing and applies this to all grid cells in PVTNUM region one. For
PVTNUM regions one and two the keyword applies 0.005 stb//Mscf/d as the
maximum rate at which the vaporized oil-gas ratio is allowed to increase
in a grid cell,

\begin{verbatim}
--
--       VAPORIZED OIL (RV) MAXIMUM RATE OF INCREASE PARAMETERS BY REGION
--
DRVDTR
--       MAX RV
--       DRVDT1
--       -------
         0.0000                                            /
         0.0005                                            /
         0.0005                                            /
\end{verbatim}

Again, the keyword parameters when applied are subject to the
availability of free oil and the ability of the undersaturated gas to
adsorb this oil.

\subsection{DUMPCUPL -- Activate Output to the Reservoir Coupling
File}\label{dumpcupl-activate-output-to-the-reservoir-coupling-file}

\begin{longtable}[]{@{}
  >{\raggedright\arraybackslash}p{(\linewidth - 14\tabcolsep) * \real{0.1250}}
  >{\raggedright\arraybackslash}p{(\linewidth - 14\tabcolsep) * \real{0.1250}}
  >{\raggedright\arraybackslash}p{(\linewidth - 14\tabcolsep) * \real{0.1250}}
  >{\raggedright\arraybackslash}p{(\linewidth - 14\tabcolsep) * \real{0.1250}}
  >{\raggedright\arraybackslash}p{(\linewidth - 14\tabcolsep) * \real{0.1250}}
  >{\raggedright\arraybackslash}p{(\linewidth - 14\tabcolsep) * \real{0.1250}}
  >{\raggedright\arraybackslash}p{(\linewidth - 14\tabcolsep) * \real{0.1250}}
  >{\raggedright\arraybackslash}p{(\linewidth - 14\tabcolsep) * \real{0.1250}}@{}}
\toprule\noalign{}
\begin{minipage}[b]{\linewidth}\raggedright
RUNSPEC
\end{minipage} & \begin{minipage}[b]{\linewidth}\raggedright
GRID
\end{minipage} & \begin{minipage}[b]{\linewidth}\raggedright
EDIT
\end{minipage} & \begin{minipage}[b]{\linewidth}\raggedright
PROPS
\end{minipage} & \begin{minipage}[b]{\linewidth}\raggedright
REGIONS
\end{minipage} & \begin{minipage}[b]{\linewidth}\raggedright
SOLUTION
\end{minipage} & \begin{minipage}[b]{\linewidth}\raggedright
SUMMARY
\end{minipage} & \begin{minipage}[b]{\linewidth}\raggedright
SCHEDULE
\end{minipage} \\
\midrule\noalign{}
\endhead
\bottomrule\noalign{}
\endlastfoot
\end{longtable}

\subsubsection{Description}\label{description-909}

This keyword, DUMPCUPL, activates output to the Reservoir Coupling file
from the reservoir coupling file in the master run for when reservoir
coupling is invoked by the GRUPMAST and SLAVES keywords in the SCHEDULE
section.

This keyword is not supported by OPM Flow but would change the results
if supported so the simulation will be stopped.

\subsection{DYNAMICR -- Define Dynamic Region
Parameters}\label{dynamicr-define-dynamic-region-parameters}

The DYNAMICR keyword marks the start of a Dynamic Region section and
defines the parameters used for Dynamic Regions that allows for property
and reporting regions to vary as the run progresses, based on the
parameters and logic defined by this keyword and section. A Dynamic
Region section is terminated by the ENDDYN keyword in the SOLUTION or
SCHEDULE sections.

See DYNAMICR -- Start of Dynamic Region Parameter Definition in the
SOLUTION section for a full description.

\subsection{ECHO -- Activate Echoing of User Input Files to the Print
File}\label{echo-activate-echoing-of-user-input-files-to-the-print-file-8}

Turns on echoing of all the input files to the print file; note that
this keyword is activated by default and can subsequently be switched
off by the NOECHO activation keyword.

See ECHO -- Activate Echoing of User Input Files to the Print File in
the GLOBAL section for a full description.

\subsection{END -- Define the End of the Input
File}\label{end-define-the-end-of-the-input-file-8}

This keyword marks the end of the input file and can occur in any
section. Any keywords and data after the END keyword are ignored.

See END -- Define the End of the Input File in the GLOBAL section for a
full description.

\subsection{ENDACTIO -- End the Definition of ACTION
Commands}\label{endactio-end-the-definition-of-action-commands}

\begin{longtable}[]{@{}
  >{\raggedright\arraybackslash}p{(\linewidth - 14\tabcolsep) * \real{0.1250}}
  >{\raggedright\arraybackslash}p{(\linewidth - 14\tabcolsep) * \real{0.1250}}
  >{\raggedright\arraybackslash}p{(\linewidth - 14\tabcolsep) * \real{0.1250}}
  >{\raggedright\arraybackslash}p{(\linewidth - 14\tabcolsep) * \real{0.1250}}
  >{\raggedright\arraybackslash}p{(\linewidth - 14\tabcolsep) * \real{0.1250}}
  >{\raggedright\arraybackslash}p{(\linewidth - 14\tabcolsep) * \real{0.1250}}
  >{\raggedright\arraybackslash}p{(\linewidth - 14\tabcolsep) * \real{0.1250}}
  >{\raggedright\arraybackslash}p{(\linewidth - 14\tabcolsep) * \real{0.1250}}@{}}
\toprule\noalign{}
\begin{minipage}[b]{\linewidth}\raggedright
RUNSPEC
\end{minipage} & \begin{minipage}[b]{\linewidth}\raggedright
GRID
\end{minipage} & \begin{minipage}[b]{\linewidth}\raggedright
EDIT
\end{minipage} & \begin{minipage}[b]{\linewidth}\raggedright
PROPS
\end{minipage} & \begin{minipage}[b]{\linewidth}\raggedright
REGIONS
\end{minipage} & \begin{minipage}[b]{\linewidth}\raggedright
SOLUTION
\end{minipage} & \begin{minipage}[b]{\linewidth}\raggedright
SUMMARY
\end{minipage} & \begin{minipage}[b]{\linewidth}\raggedright
SCHEDULE
\end{minipage} \\
\midrule\noalign{}
\endhead
\bottomrule\noalign{}
\endlastfoot
\end{longtable}

\subsubsection{Description}\label{description-910}

The ENDACTIO keyword defines the end of a series of conditions that
invoke run time processing of the ACTION series of keywords, namely:
ACTION, ACTIONG, ACTIONR, ACTIONS, ACTIONW and ACTIONX. Only the ACTIONX
keyword is implemented in OPM Flow as this keyword implements the
ACTION, ACTIONG, ACTIONR, ACTIONS, ACTIONW functionality with greater
flexibility. See the ACTIONX keyword in the SCHEDULE section for a full
description of the ACTION facility.

There is no data required for this keyword and there is no terminating
``/'' for this keyword.

\subsubsection{Example}\label{example-519}

The example shows the use of the ACTIONX and ENDACTIO keywords to test
if the field's gas production rate is less than 600 MMscf/d after 2020
and to open up additional wells if this occurs.

\begin{verbatim}
--
-- START OF ACTIONX FIELD PHASE-2 DEVELOPMENT DEFINITION
--
ACTIONX
     PHASE2       1               /
     GGPR  'FIELD' < 600E3 AND    /
     YEAR > 2020                  /
/
-- WELL PRODUCTION STATUS
--
--  WELL    WELL   --LOCATION--  COMPLETION
--  NAME    STAT     I   J    K  FIRST LAST
WELOPEN
GP10        OPEN                           /
GP11        OPEN                           /
/
--
-- END OF ACTIONX FIELD PHASE-2 DEVELOPMENT DEFINITION
--
ENDACTIO
\end{verbatim}

\subsection{ENDBOX -- Define the End of the BOX Defined
Grid}\label{endbox-define-the-end-of-the-box-defined-grid-5}

This keyword marks the end of a previously defined BOX sub-grid as
defined by a previously entered BOX keyword. The keyword resets the
input grid to be the full grid as defined by the NX, NY, and NZ
variables on the DIMENS keyword in the RUNSPEC section.

See ENDBOX -- Define the End of the BOX Defined Grid in the GRID section
for a full description.

\subsection{ENDDYN -- End of Dynamic Region Parameter
Definition}\label{enddyn-end-of-dynamic-region-parameter-definition-1}

The ENDDYN keyword marks the end of a Dynamic Region section that was
started with the DYNAMICR keyword in the SOLUTION or SCHEDULE sections.
Dynamic Regions allow for property and reporting regions to vary as the
run progresses, based on the parameters and logic defined within the
section.

See ENDDYN-- End of Dynamic Region Parameter Definition in the SOLUTION
for a full description.

\subsection{ENDFIN -- End the Definition of a Local Grid
Refinement}\label{endfin-end-the-definition-of-a-local-grid-refinement-5}

ENDFIN defines the end of a Cartesian or radial local grid refinement
(``LGR'') definition and a LGR property definition data set.

See ENDFIN -- End the Definition of a Local Grid Refinement in the GRID
section for a full description.

\subsection{ENDINC -- Define the End of an Include
File}\label{endinc-define-the-end-of-an-include-file-8}

This keyword marks the end of an include file specified on the INCLUDE
keyword. When the ENDINC keyword is encountered in the INCLUDE file,
input data is read from the next keyword in the current file. Any
keywords and data after the ENDINC keyword in the INCLUDE file are
ignored.

See ENDINC -- Define the End of an Include File in the GLOBAL section
for a full description.

\subsection{ENDSKIP -- Deactivate Skipping of Keywords and Input
Data}\label{endskip-deactivate-skipping-of-keywords-and-input-data-8}

The ENDSKIP keyword deactivates the skipping of keywords that was
activated by the SKIP, SKIP100, or SKIP300 keywords. Each SKIP keyword
should be paired with an ENDSKIP keyword.

See ENDSKIP -- Deactivate Skipping of Keywords and Input Data in the
GLOBAL section for a full description.

\subsection{EPSDBGS -- Write End-Point Debug Data to the DEBUG File
(Multiple)}\label{epsdbgs-write-end-point-debug-data-to-the-debug-file-multiple-1}

This keyword, EPSDBGS, defines the end-point debug data for multiple
grid blocks that should be written to the debug file (*.DBG) for when
the End-Point Scaling option has been activated by the ENDSCALE keyword
in the RONSPEC section.

See EPSDBGS - Write End-Point Debug Data to the DEBUG File (Multiple) in
the PROPS section for afull description.

\subsection{EPSDEBUG -- Write End-Point Debug Data to the DEBUG File
(Individual)}\label{epsdebug-write-end-point-debug-data-to-the-debug-file-individual-1}

This keyword, EPSDEBUG, defines the end-point debug data for individual
grid blocks that should be written to the Debug file (*.DBG) for when
the End-Point Scaling option has been activated by the ENDSCALE keyword
in the RUNSPEC section.

See EPSDEBUG - Write End-Point Debug Data to the DEBUG File
(Individual)n the PROPS section for afull description.

\subsection{EXCAVATE -- Set the Status of a Grid Block To Active or
Excavate}\label{excavate-set-the-status-of-a-grid-block-to-active-or-excavate}

\begin{longtable}[]{@{}
  >{\raggedright\arraybackslash}p{(\linewidth - 14\tabcolsep) * \real{0.1250}}
  >{\raggedright\arraybackslash}p{(\linewidth - 14\tabcolsep) * \real{0.1250}}
  >{\raggedright\arraybackslash}p{(\linewidth - 14\tabcolsep) * \real{0.1250}}
  >{\raggedright\arraybackslash}p{(\linewidth - 14\tabcolsep) * \real{0.1250}}
  >{\raggedright\arraybackslash}p{(\linewidth - 14\tabcolsep) * \real{0.1250}}
  >{\raggedright\arraybackslash}p{(\linewidth - 14\tabcolsep) * \real{0.1250}}
  >{\raggedright\arraybackslash}p{(\linewidth - 14\tabcolsep) * \real{0.1250}}
  >{\raggedright\arraybackslash}p{(\linewidth - 14\tabcolsep) * \real{0.1250}}@{}}
\toprule\noalign{}
\begin{minipage}[b]{\linewidth}\raggedright
RUNSPEC
\end{minipage} & \begin{minipage}[b]{\linewidth}\raggedright
GRID
\end{minipage} & \begin{minipage}[b]{\linewidth}\raggedright
EDIT
\end{minipage} & \begin{minipage}[b]{\linewidth}\raggedright
PROPS
\end{minipage} & \begin{minipage}[b]{\linewidth}\raggedright
REGIONS
\end{minipage} & \begin{minipage}[b]{\linewidth}\raggedright
SOLUTION
\end{minipage} & \begin{minipage}[b]{\linewidth}\raggedright
SUMMARY
\end{minipage} & \begin{minipage}[b]{\linewidth}\raggedright
SCHEDULE
\end{minipage} \\
\midrule\noalign{}
\endhead
\bottomrule\noalign{}
\endlastfoot
\end{longtable}

\subsubsection{Description}\label{description-911}

This keyword, EXCAVATE, sets the status of global and LGR grid blocks to
active or excavate. Excavated grid blocks have all the
transmissibilities set to zero thus disabling flow between the
surrounding grid blocks.

This keyword is not supported by OPM Flow but would change the results
if supported so the simulation will be stopped.

\subsection{EXIT -- Exit Simulation from within an Action
Section}\label{exit-exit-simulation-from-within-an-action-section}

\begin{longtable}[]{@{}
  >{\raggedright\arraybackslash}p{(\linewidth - 14\tabcolsep) * \real{0.1250}}
  >{\raggedright\arraybackslash}p{(\linewidth - 14\tabcolsep) * \real{0.1250}}
  >{\raggedright\arraybackslash}p{(\linewidth - 14\tabcolsep) * \real{0.1250}}
  >{\raggedright\arraybackslash}p{(\linewidth - 14\tabcolsep) * \real{0.1250}}
  >{\raggedright\arraybackslash}p{(\linewidth - 14\tabcolsep) * \real{0.1250}}
  >{\raggedright\arraybackslash}p{(\linewidth - 14\tabcolsep) * \real{0.1250}}
  >{\raggedright\arraybackslash}p{(\linewidth - 14\tabcolsep) * \real{0.1250}}
  >{\raggedright\arraybackslash}p{(\linewidth - 14\tabcolsep) * \real{0.1250}}@{}}
\toprule\noalign{}
\begin{minipage}[b]{\linewidth}\raggedright
RUNSPEC
\end{minipage} & \begin{minipage}[b]{\linewidth}\raggedright
GRID
\end{minipage} & \begin{minipage}[b]{\linewidth}\raggedright
EDIT
\end{minipage} & \begin{minipage}[b]{\linewidth}\raggedright
PROPS
\end{minipage} & \begin{minipage}[b]{\linewidth}\raggedright
REGIONS
\end{minipage} & \begin{minipage}[b]{\linewidth}\raggedright
SOLUTION
\end{minipage} & \begin{minipage}[b]{\linewidth}\raggedright
SUMMARY
\end{minipage} & \begin{minipage}[b]{\linewidth}\raggedright
SCHEDULE
\end{minipage} \\
\midrule\noalign{}
\endhead
\bottomrule\noalign{}
\endlastfoot
\end{longtable}

\subsubsection{Description}\label{description-912}

The EXIT keyword is part of OPM Flow's ACTION facility that allows for
terminating the simulation for when a condition within an ACTIONX
definition is satisfied. Invoking the keyword within an ACTIONX
definition will result in the simulation terminating with an exit status
code. The ACTION facility allows the user to enter computational logic
to the simulation run based on the how the simulation run is proceeding
-- see the ACTIONX keyword in the SCHEDULE section.

\begin{tcolorbox}[enhanced jigsaw, title=\textcolor{quarto-callout-note-color}{\faInfo}\hspace{0.5em}{Note}, coltitle=black, breakable, toprule=.15mm, colbacktitle=quarto-callout-note-color!10!white, leftrule=.75mm, opacitybacktitle=0.6, colback=white, bottomtitle=1mm, toptitle=1mm, titlerule=0mm, left=2mm, colframe=quarto-callout-note-color-frame, arc=.35mm, rightrule=.15mm, bottomrule=.15mm, opacityback=0]

This is an OPM Flow specific keyword for the simulator's ACTION facility
and will therefore cause an error if used in the commercial simulator.

\end{tcolorbox}

\begin{longtable}[]{@{}
  >{\raggedright\arraybackslash}p{(\linewidth - 6\tabcolsep) * \real{0.2500}}
  >{\raggedright\arraybackslash}p{(\linewidth - 6\tabcolsep) * \real{0.2500}}
  >{\raggedright\arraybackslash}p{(\linewidth - 6\tabcolsep) * \real{0.2500}}
  >{\raggedright\arraybackslash}p{(\linewidth - 6\tabcolsep) * \real{0.2500}}@{}}
\toprule\noalign{}
\begin{minipage}[b]{\linewidth}\raggedright
No.
\end{minipage} & \begin{minipage}[b]{\linewidth}\raggedright
Name
\end{minipage} & \begin{minipage}[b]{\linewidth}\raggedright
Description
\end{minipage} & \begin{minipage}[b]{\linewidth}\raggedright
Default
\end{minipage} \\
\midrule\noalign{}
\endhead
\bottomrule\noalign{}
\endlastfoot
1 & EXITCODE & An optional integer that sets the exit code printed to
the *.PRT file, if not not defined the default value of zero will be
used. & 0 \\
Notes: & & & \\
\end{longtable}

\emph{Table 12.25: EXIT Keyword Description}

The EXIT keyword should only be used as part of an ACTIONX block, if
found elsewhere in the input deck it will be ignored.

\subsubsection{Examples}\label{examples-74}

The first example uses the ACTIONX keyword to define a condition for
when the Field Oil Production Rate (``FOPR) falls below 1,000 stb/d (or
1,000 m3) using the default value for the EXITCODE.

\begin{verbatim}
--
--       DEFINE START OF ACTIONX SECTION
--
ACTIONX
         'CHECK_FOPR' 100000            /
         FOPR < 1000                    /
/
--
--       TERMINATE AND EXIT SIMULATION
--
EXIT
                                        /

ENDACTIO
\end{verbatim}

The next example terminates the simulation with EXITCODE one when the
Field Pressure (``FPR'') falls below 200 psia (or 200 barsa).

\begin{verbatim}
--
--       DEFINE START OF ACTIONX SECTION
--
ACTIONX
         'CHECK_FPR'  100000             /
          FPR < 200                      /
/
--
--       TERMINATE AND EXIT SIMULATION
--
EXIT
         1                               /

ENDACTIO
\end{verbatim}

Note is is probably good practice to always set the EXITCODE to be able
to identify the reason for the simulation stopping.

\begin{verbatim}
\end{verbatim}

\subsection{EXTRAPMS -- Activate Extrapolation Warning
Messages}\label{extrapms-activate-extrapolation-warning-messages-8}

The EXTRAPMS keyword activates extrapolation warning messages for when
OPM Flow extrapolates the PVT or VFP tables. Frequent extrapolation
warning messages should be investigated and resolved as this would
indicate possible incorrect data and may result in the simulator
extrapolating to unrealistic values.

See EXTRAPMS -- Activate Extrapolation Warning Messages in the GLOBAL
section for a full description.

\subsection{FBHPDEF -- Define Well Default BHP Target and
Constraints}\label{fbhpdef-define-well-default-bhp-target-and-constraints}

\begin{longtable}[]{@{}
  >{\raggedright\arraybackslash}p{(\linewidth - 14\tabcolsep) * \real{0.1250}}
  >{\raggedright\arraybackslash}p{(\linewidth - 14\tabcolsep) * \real{0.1250}}
  >{\raggedright\arraybackslash}p{(\linewidth - 14\tabcolsep) * \real{0.1250}}
  >{\raggedright\arraybackslash}p{(\linewidth - 14\tabcolsep) * \real{0.1250}}
  >{\raggedright\arraybackslash}p{(\linewidth - 14\tabcolsep) * \real{0.1250}}
  >{\raggedright\arraybackslash}p{(\linewidth - 14\tabcolsep) * \real{0.1250}}
  >{\raggedright\arraybackslash}p{(\linewidth - 14\tabcolsep) * \real{0.1250}}
  >{\raggedright\arraybackslash}p{(\linewidth - 14\tabcolsep) * \real{0.1250}}@{}}
\toprule\noalign{}
\begin{minipage}[b]{\linewidth}\raggedright
RUNSPEC
\end{minipage} & \begin{minipage}[b]{\linewidth}\raggedright
GRID
\end{minipage} & \begin{minipage}[b]{\linewidth}\raggedright
EDIT
\end{minipage} & \begin{minipage}[b]{\linewidth}\raggedright
PROPS
\end{minipage} & \begin{minipage}[b]{\linewidth}\raggedright
REGIONS
\end{minipage} & \begin{minipage}[b]{\linewidth}\raggedright
SOLUTION
\end{minipage} & \begin{minipage}[b]{\linewidth}\raggedright
SUMMARY
\end{minipage} & \begin{minipage}[b]{\linewidth}\raggedright
SCHEDULE
\end{minipage} \\
\midrule\noalign{}
\endhead
\bottomrule\noalign{}
\endlastfoot
\end{longtable}

\subsubsection{Description}\label{description-913}

This keyword, FBHPDEF, defines the default well BHP target for
production wells and the default BHP constraint for injection wells.

\begin{longtable}[]{@{}
  >{\raggedright\arraybackslash}p{(\linewidth - 6\tabcolsep) * \real{0.2500}}
  >{\raggedright\arraybackslash}p{(\linewidth - 6\tabcolsep) * \real{0.2500}}
  >{\raggedright\arraybackslash}p{(\linewidth - 6\tabcolsep) * \real{0.2500}}
  >{\raggedright\arraybackslash}p{(\linewidth - 6\tabcolsep) * \real{0.2500}}@{}}
\toprule\noalign{}
\begin{minipage}[b]{\linewidth}\raggedright
No.
\end{minipage} & \begin{minipage}[b]{\linewidth}\raggedright
Name
\end{minipage} & \begin{minipage}[b]{\linewidth}\raggedright
Description
\end{minipage} & \begin{minipage}[b]{\linewidth}\raggedright
Default
\end{minipage} \\
\midrule\noalign{}
\endhead
\bottomrule\noalign{}
\endlastfoot
Field & Metric & Laboratory & \\
1 & TARGET\_BHP & A real positive value that defines the default well
BHP target for production wells. & 1.01325 barsa \\
psia & barsa & atma & \\
2 & LIMIT\_BHP & A real positive value that defines the default well BHP
limit for injection wells. & 6895 barsa \\
psia & barsa & atma & \\
Notes: & & & \\
\end{longtable}

\emph{Table 12.3.72.1: FBHPDEF Keyword Description}

See also the WELTARG keyword in the SCHEDULE section that modifies a
well's target and constraint values.

\subsubsection{Example}\label{example-520}

The following examples sets the default BHP target for production wells
to 1.01325, and the default BHP limit for injection wells to 6895:

\begin{verbatim}
--
--       DEFAULT BHP TARGET AND CONSTRAINTS
--
FBHPDEF
         1.01325  6895.  /
\end{verbatim}

\subsection{FILEUNIT -- Activate Unit Consistency
Verification}\label{fileunit-activate-unit-consistency-verification-4}

The FILEUNIT keyword defines the units of the data set, and is used to
verify that the units in the input deck and any associated include files
are consistent. The keyword does not provide for the conversion between
different sets of units.

See FILEUNIT -- Activate Unit Consistency Checking in the GRID section
for a full description.

\subsection{FORMFEED -- Defined the Print File Form-Feed
Character}\label{formfeed-defined-the-print-file-form-feed-character-8}

The FORMFEED keyword defines the form-feed character, or carriage
control character, for the output print (\emph{.PRT) run summary (}.RSM)
files. The keyword should be place at the very top of the input file.

See FORMFEED -- Defined the Print File Form-Feed Character in the GLOBAL
section for a dull description.

\subsection{GASBEGIN -- Define Start of Annual Scheduling
Section}\label{gasbegin-define-start-of-annual-scheduling-section}

\begin{longtable}[]{@{}
  >{\raggedright\arraybackslash}p{(\linewidth - 14\tabcolsep) * \real{0.1250}}
  >{\raggedright\arraybackslash}p{(\linewidth - 14\tabcolsep) * \real{0.1250}}
  >{\raggedright\arraybackslash}p{(\linewidth - 14\tabcolsep) * \real{0.1250}}
  >{\raggedright\arraybackslash}p{(\linewidth - 14\tabcolsep) * \real{0.1250}}
  >{\raggedright\arraybackslash}p{(\linewidth - 14\tabcolsep) * \real{0.1250}}
  >{\raggedright\arraybackslash}p{(\linewidth - 14\tabcolsep) * \real{0.1250}}
  >{\raggedright\arraybackslash}p{(\linewidth - 14\tabcolsep) * \real{0.1250}}
  >{\raggedright\arraybackslash}p{(\linewidth - 14\tabcolsep) * \real{0.1250}}@{}}
\toprule\noalign{}
\begin{minipage}[b]{\linewidth}\raggedright
RUNSPEC
\end{minipage} & \begin{minipage}[b]{\linewidth}\raggedright
GRID
\end{minipage} & \begin{minipage}[b]{\linewidth}\raggedright
EDIT
\end{minipage} & \begin{minipage}[b]{\linewidth}\raggedright
PROPS
\end{minipage} & \begin{minipage}[b]{\linewidth}\raggedright
REGIONS
\end{minipage} & \begin{minipage}[b]{\linewidth}\raggedright
SOLUTION
\end{minipage} & \begin{minipage}[b]{\linewidth}\raggedright
SUMMARY
\end{minipage} & \begin{minipage}[b]{\linewidth}\raggedright
SCHEDULE
\end{minipage} \\
\midrule\noalign{}
\endhead
\bottomrule\noalign{}
\endlastfoot
\end{longtable}

\subsubsection{Description}\label{description-914}

This keyword, GASBEGIN, defines the start of an Annual Scheduling
section set of keywords used when the Gas Field Operations option has
been activated by the GASFIELD keyword in the RUNSPEC section. An Annual
Scheduling section starts with the GASBEGIN keyword and is terminated by
the GASEND keyword, with keywords in between used to control and write
reports at selected times between the start and end of a contract
period. Only one Annual Scheduling section is activate at a time, that
is, a subsequent Annual Scheduling section overwrites the previous set
of entries. To clear the current Annual Schedule section enter the
GASBEGIN keyword followed by the GASEND keyword word with no other
keywords in between.

There is no data required for this keyword and there is no terminating
``/'' for this keyword.

This keyword is not supported by OPM Flow but would change the results
if supported so the simulation will be stopped.

\subsection{GASEND -- Define End of Annual Scheduling
Section}\label{gasend-define-end-of-annual-scheduling-section}

\begin{longtable}[]{@{}
  >{\raggedright\arraybackslash}p{(\linewidth - 14\tabcolsep) * \real{0.1250}}
  >{\raggedright\arraybackslash}p{(\linewidth - 14\tabcolsep) * \real{0.1250}}
  >{\raggedright\arraybackslash}p{(\linewidth - 14\tabcolsep) * \real{0.1250}}
  >{\raggedright\arraybackslash}p{(\linewidth - 14\tabcolsep) * \real{0.1250}}
  >{\raggedright\arraybackslash}p{(\linewidth - 14\tabcolsep) * \real{0.1250}}
  >{\raggedright\arraybackslash}p{(\linewidth - 14\tabcolsep) * \real{0.1250}}
  >{\raggedright\arraybackslash}p{(\linewidth - 14\tabcolsep) * \real{0.1250}}
  >{\raggedright\arraybackslash}p{(\linewidth - 14\tabcolsep) * \real{0.1250}}@{}}
\toprule\noalign{}
\begin{minipage}[b]{\linewidth}\raggedright
RUNSPEC
\end{minipage} & \begin{minipage}[b]{\linewidth}\raggedright
GRID
\end{minipage} & \begin{minipage}[b]{\linewidth}\raggedright
EDIT
\end{minipage} & \begin{minipage}[b]{\linewidth}\raggedright
PROPS
\end{minipage} & \begin{minipage}[b]{\linewidth}\raggedright
REGIONS
\end{minipage} & \begin{minipage}[b]{\linewidth}\raggedright
SOLUTION
\end{minipage} & \begin{minipage}[b]{\linewidth}\raggedright
SUMMARY
\end{minipage} & \begin{minipage}[b]{\linewidth}\raggedright
SCHEDULE
\end{minipage} \\
\midrule\noalign{}
\endhead
\bottomrule\noalign{}
\endlastfoot
\end{longtable}

\subsubsection{Description}\label{description-915}

This keyword, GASEND, defines the end of an Annual Scheduling section
set of keywords used when the Gas Field Operations option has been
activated by the GASFIELD keyword in the RUNSPEC section. An Annual
Scheduling section starts with the GASBEGIN keyword and is terminated by
the GASEND keyword, with keywords in between used to control and write
reports at selected times between the start and end of a contract
period. Only one Annual Scheduling section is activate at a time, that
is, a subsequent Annual Scheduling section overwrites the previous set
of entries. To clear the current Annual Schedule section enter the
GASBEGIN keyword followed by the GASEND keyword word with no other
keywords in between.

There is no data required for this keyword and there is no terminating
``/'' for this keyword.

This keyword is not supported by OPM Flow but would change the results
if supported so the simulation will be stopped.

\subsection{GASFCOMP -- Define Automatic Gas
Compressors}\label{gasfcomp-define-automatic-gas-compressors}

\begin{longtable}[]{@{}
  >{\raggedright\arraybackslash}p{(\linewidth - 14\tabcolsep) * \real{0.1250}}
  >{\raggedright\arraybackslash}p{(\linewidth - 14\tabcolsep) * \real{0.1250}}
  >{\raggedright\arraybackslash}p{(\linewidth - 14\tabcolsep) * \real{0.1250}}
  >{\raggedright\arraybackslash}p{(\linewidth - 14\tabcolsep) * \real{0.1250}}
  >{\raggedright\arraybackslash}p{(\linewidth - 14\tabcolsep) * \real{0.1250}}
  >{\raggedright\arraybackslash}p{(\linewidth - 14\tabcolsep) * \real{0.1250}}
  >{\raggedright\arraybackslash}p{(\linewidth - 14\tabcolsep) * \real{0.1250}}
  >{\raggedright\arraybackslash}p{(\linewidth - 14\tabcolsep) * \real{0.1250}}@{}}
\toprule\noalign{}
\begin{minipage}[b]{\linewidth}\raggedright
RUNSPEC
\end{minipage} & \begin{minipage}[b]{\linewidth}\raggedright
GRID
\end{minipage} & \begin{minipage}[b]{\linewidth}\raggedright
EDIT
\end{minipage} & \begin{minipage}[b]{\linewidth}\raggedright
PROPS
\end{minipage} & \begin{minipage}[b]{\linewidth}\raggedright
REGIONS
\end{minipage} & \begin{minipage}[b]{\linewidth}\raggedright
SOLUTION
\end{minipage} & \begin{minipage}[b]{\linewidth}\raggedright
SUMMARY
\end{minipage} & \begin{minipage}[b]{\linewidth}\raggedright
SCHEDULE
\end{minipage} \\
\midrule\noalign{}
\endhead
\bottomrule\noalign{}
\endlastfoot
\end{longtable}

\subsubsection{Description}\label{description-916}

This keyword, GASFCOMP, defines automatic gas compressors for when the
Gas Field Operations option has been activated by the GASFIELD keyword
in the RUNSPEC section and the Standard Network option has been
specified by the GRUPTREE, GRUPNET and GNETINJE series of keywords in
the SCHEDULE section. Automatic gas compressors are automatically switch
on for a group if a group's gas production target cannot be satisfied.
In addition, if a group's gas target is reduced then the automatic
compressors are initially switch off to test that the reduced gas rate
target can be met without compression, if not, compression is switched
back on. Note that all automatic compressors are ``switch on'' when
calculating a field's gas deliverability.

There is no data required for this keyword and there is no terminating
``/'' for this keyword.

This keyword is not supported by OPM Flow but would change the results
if supported so the simulation will be stopped.

\subsection{GASFDECR -- Define Field Gas Sales Contract Monthly
Reduction}\label{gasfdecr-define-field-gas-sales-contract-monthly-reduction}

\begin{longtable}[]{@{}
  >{\raggedright\arraybackslash}p{(\linewidth - 14\tabcolsep) * \real{0.1250}}
  >{\raggedright\arraybackslash}p{(\linewidth - 14\tabcolsep) * \real{0.1250}}
  >{\raggedright\arraybackslash}p{(\linewidth - 14\tabcolsep) * \real{0.1250}}
  >{\raggedright\arraybackslash}p{(\linewidth - 14\tabcolsep) * \real{0.1250}}
  >{\raggedright\arraybackslash}p{(\linewidth - 14\tabcolsep) * \real{0.1250}}
  >{\raggedright\arraybackslash}p{(\linewidth - 14\tabcolsep) * \real{0.1250}}
  >{\raggedright\arraybackslash}p{(\linewidth - 14\tabcolsep) * \real{0.1250}}
  >{\raggedright\arraybackslash}p{(\linewidth - 14\tabcolsep) * \real{0.1250}}@{}}
\toprule\noalign{}
\begin{minipage}[b]{\linewidth}\raggedright
RUNSPEC
\end{minipage} & \begin{minipage}[b]{\linewidth}\raggedright
GRID
\end{minipage} & \begin{minipage}[b]{\linewidth}\raggedright
EDIT
\end{minipage} & \begin{minipage}[b]{\linewidth}\raggedright
PROPS
\end{minipage} & \begin{minipage}[b]{\linewidth}\raggedright
REGIONS
\end{minipage} & \begin{minipage}[b]{\linewidth}\raggedright
SOLUTION
\end{minipage} & \begin{minipage}[b]{\linewidth}\raggedright
SUMMARY
\end{minipage} & \begin{minipage}[b]{\linewidth}\raggedright
SCHEDULE
\end{minipage} \\
\midrule\noalign{}
\endhead
\bottomrule\noalign{}
\endlastfoot
\end{longtable}

\subsubsection{Description}\label{description-917}

This keyword, GASFDECR, defines the field's monthly reduction in the gas
sales contract quantity for when the Gas Field Operations option has
been activated by the GASFIELD keyword in the RUNSPEC section. Gas
contracts are commonly based on a Daily Contract Quantity (``DCQ'') that
determines the gas rate that the field should be produced at, which is
normally expressed as a multiple of the DCQ, for example 1.33, and is
often referred to as the ``swing factor''. Some gas contracts also
define a maximum DCQ (``Max DCQ'') and/or a minimum take or pay DCQ
(``Min DCQ''), as well as seasonal demand characteristics. For example,
gas rates may be set higher in the winter months in order to meet
heating demand compared with summer months in colder climates, and the
opposite in warmer climates where air conditioning demand is high.

Thus, the DCQ must be calculated first over a contract period, where the
contract period is commonly contract years, or in some instances
contract quarters. This is performed by the simulator using the current
DCQ and checking to see if the (DCQ x Swing Factor) can be satisfied
throughout the current contract period, if not the DCQ is re-calculated
so that the (DCQ x Swing Factor) condition is satisfied. Once this
condition is met, the second and final pass uses the calculated DCQ in
conjunction with the monthly scaling profile data to set the monthly gas
rate for the field:

\begin{equation}\phantomsection\label{eq-12-19}{
{Q}_{\mathit{month}} = \mathit{DCQ}\times {\mathit{SWINGFAC}}_{\mathit{month}}
}\end{equation}

Where:

Qmonth = the monthly gas production target

DCQ = Daily Contract Quantity

SWINGFACmonth = monthly rate scaling factor that takes into account
seasonal demand, etc.

Here the GASFDECR keyword allows for a reduction in the calculated final
pass monthly gas rates and thus equation (12.19) become:

\begin{equation}\phantomsection\label{eq-12-20}{
{Q}_{\mathit{month}} = \left(\mathit{DCQ}\times {\mathit{SWINGFAC}}_{\mathit{month}}\right) - {\mathit{GASFDECR}}_{\mathit{month}}
}\end{equation}

Where:

Qmonth = the monthly gas production target

DCQ = Daily Contract Quantity

SWINGFACmonth = monthly rate scaling factor that takes into account
seasonal demand,

etc.

GASFDECRmonth = monthly gas rate reduction.

Since the simulator must make two passes to calculate the final rates
this will naturally decrease computational efficiency.

This keyword is not supported by OPM Flow but would change the results
if supported so the simulation will be stopped.

\subsection{GASFDELC -- Define Gas Deliverability
Calculation}\label{gasfdelc-define-gas-deliverability-calculation}

\begin{longtable}[]{@{}
  >{\raggedright\arraybackslash}p{(\linewidth - 14\tabcolsep) * \real{0.1250}}
  >{\raggedright\arraybackslash}p{(\linewidth - 14\tabcolsep) * \real{0.1250}}
  >{\raggedright\arraybackslash}p{(\linewidth - 14\tabcolsep) * \real{0.1250}}
  >{\raggedright\arraybackslash}p{(\linewidth - 14\tabcolsep) * \real{0.1250}}
  >{\raggedright\arraybackslash}p{(\linewidth - 14\tabcolsep) * \real{0.1250}}
  >{\raggedright\arraybackslash}p{(\linewidth - 14\tabcolsep) * \real{0.1250}}
  >{\raggedright\arraybackslash}p{(\linewidth - 14\tabcolsep) * \real{0.1250}}
  >{\raggedright\arraybackslash}p{(\linewidth - 14\tabcolsep) * \real{0.1250}}@{}}
\toprule\noalign{}
\begin{minipage}[b]{\linewidth}\raggedright
RUNSPEC
\end{minipage} & \begin{minipage}[b]{\linewidth}\raggedright
GRID
\end{minipage} & \begin{minipage}[b]{\linewidth}\raggedright
EDIT
\end{minipage} & \begin{minipage}[b]{\linewidth}\raggedright
PROPS
\end{minipage} & \begin{minipage}[b]{\linewidth}\raggedright
REGIONS
\end{minipage} & \begin{minipage}[b]{\linewidth}\raggedright
SOLUTION
\end{minipage} & \begin{minipage}[b]{\linewidth}\raggedright
SUMMARY
\end{minipage} & \begin{minipage}[b]{\linewidth}\raggedright
SCHEDULE
\end{minipage} \\
\midrule\noalign{}
\endhead
\bottomrule\noalign{}
\endlastfoot
\end{longtable}

\subsubsection{Description}\label{description-918}

The GASFDELC keyword defines how the field's gas deliverability
calculation should be performed for when the Gas Field Operations option
has been activated by the GASFIELD keyword in the RUNSPEC section.

This keyword is not supported by OPM Flow but would change the results
if supported so the simulation will be stopped.

\subsection{GASFTARG -- Define Field Gas Sales Contract Monthly
Target}\label{gasftarg-define-field-gas-sales-contract-monthly-target}

\begin{longtable}[]{@{}
  >{\raggedright\arraybackslash}p{(\linewidth - 14\tabcolsep) * \real{0.1250}}
  >{\raggedright\arraybackslash}p{(\linewidth - 14\tabcolsep) * \real{0.1250}}
  >{\raggedright\arraybackslash}p{(\linewidth - 14\tabcolsep) * \real{0.1250}}
  >{\raggedright\arraybackslash}p{(\linewidth - 14\tabcolsep) * \real{0.1250}}
  >{\raggedright\arraybackslash}p{(\linewidth - 14\tabcolsep) * \real{0.1250}}
  >{\raggedright\arraybackslash}p{(\linewidth - 14\tabcolsep) * \real{0.1250}}
  >{\raggedright\arraybackslash}p{(\linewidth - 14\tabcolsep) * \real{0.1250}}
  >{\raggedright\arraybackslash}p{(\linewidth - 14\tabcolsep) * \real{0.1250}}@{}}
\toprule\noalign{}
\begin{minipage}[b]{\linewidth}\raggedright
RUNSPEC
\end{minipage} & \begin{minipage}[b]{\linewidth}\raggedright
GRID
\end{minipage} & \begin{minipage}[b]{\linewidth}\raggedright
EDIT
\end{minipage} & \begin{minipage}[b]{\linewidth}\raggedright
PROPS
\end{minipage} & \begin{minipage}[b]{\linewidth}\raggedright
REGIONS
\end{minipage} & \begin{minipage}[b]{\linewidth}\raggedright
SOLUTION
\end{minipage} & \begin{minipage}[b]{\linewidth}\raggedright
SUMMARY
\end{minipage} & \begin{minipage}[b]{\linewidth}\raggedright
SCHEDULE
\end{minipage} \\
\midrule\noalign{}
\endhead
\bottomrule\noalign{}
\endlastfoot
\end{longtable}

\subsubsection{Description}\label{description-919}

This keyword, GASFTARG, defines the field's monthly gas sales contract
quantity for when the Gas Field Operations option has been activated by
the GASFIELD keyword in the RUNSPEC section. Gas contracts are commonly
based on a Daily Contract Quantity (``DCQ'') that determines the gas
rate that the field should be produced at, which is normally expressed
as a multiple of the DCQ, for example 1.33, and is often referred to as
the ``swing factor''. Some gas contracts also define a maximum DCQ
(``Max DCQ'') and/or a minimum take or pay DCQ (``Min DCQ''), as well as
seasonal demand characteristics. For example, gas rates may be set
higher in the winter months in order to meet heating demand compared
with summer months in colder climates, and the opposite in warmer
climates where air conditioning demand is high.

Thus, the DCQ must be calculated first over a contract period, where the
contract period is commonly contract years, or in some instances
contract quarters. This is performed by the simulator using the current
DCQ and checking to see if the (DCQ x Swing Factor) can be satisfied
throughout the current contract period, if not the DCQ is re-calculated
so that the (DCQ x Swing Factor) condition is satisfied. Once this
condition is met, the second and final pass uses the calculated DCQ in
conjunction with the monthly scaling profile data to set the monthly gas
rate for the field:

\begin{equation}\phantomsection\label{eq-12-21}{
{Q}_{\mathit{month}} = \mathit{DCQ}\times {\mathit{SWINGFAC}}_{\mathit{month}}
}\end{equation}

Where:

Qmonth = the monthly gas production target

DCQ = Daily Contract Quantity

SWINGFACmonth = monthly rate scaling factor that takes into account
seasonal demand, etc.

Here the GASFTARG keyword sets a minimum target rate in the calculated
final pass monthly gas rates and thus equation (12.21) become:

\begin{equation}\phantomsection\label{eq-12-22}{
{Q}_{\mathit{month}} = \mathit{Minimum} \left(\left(\mathit{DCQ}\times {\mathit{SWINGFAC}}_{\mathit{month}}\right), {\mathit{GASFTARG}}_{\mathit{month}}\right)
}\end{equation}

Where:

Qmonth = the monthly gas production target

DCQ = Daily Contract Quantity

SWINGFACmonth = monthly rate scaling factor that takes into account
seasonal demand, etc.

GASFTARGmonth = minimum monthly gas rate target.

Since the simulator must make two passes to calculate the final rates
this will naturally decrease computational efficiency.

This keyword is not supported by OPM Flow but would change the results
if supported so the simulation will be stopped.

\subsection{GASMONTH -- Define Start of Annual Scheduling
Event}\label{gasmonth-define-start-of-annual-scheduling-event}

\begin{longtable}[]{@{}
  >{\raggedright\arraybackslash}p{(\linewidth - 14\tabcolsep) * \real{0.1250}}
  >{\raggedright\arraybackslash}p{(\linewidth - 14\tabcolsep) * \real{0.1250}}
  >{\raggedright\arraybackslash}p{(\linewidth - 14\tabcolsep) * \real{0.1250}}
  >{\raggedright\arraybackslash}p{(\linewidth - 14\tabcolsep) * \real{0.1250}}
  >{\raggedright\arraybackslash}p{(\linewidth - 14\tabcolsep) * \real{0.1250}}
  >{\raggedright\arraybackslash}p{(\linewidth - 14\tabcolsep) * \real{0.1250}}
  >{\raggedright\arraybackslash}p{(\linewidth - 14\tabcolsep) * \real{0.1250}}
  >{\raggedright\arraybackslash}p{(\linewidth - 14\tabcolsep) * \real{0.1250}}@{}}
\toprule\noalign{}
\begin{minipage}[b]{\linewidth}\raggedright
RUNSPEC
\end{minipage} & \begin{minipage}[b]{\linewidth}\raggedright
GRID
\end{minipage} & \begin{minipage}[b]{\linewidth}\raggedright
EDIT
\end{minipage} & \begin{minipage}[b]{\linewidth}\raggedright
PROPS
\end{minipage} & \begin{minipage}[b]{\linewidth}\raggedright
REGIONS
\end{minipage} & \begin{minipage}[b]{\linewidth}\raggedright
SOLUTION
\end{minipage} & \begin{minipage}[b]{\linewidth}\raggedright
SUMMARY
\end{minipage} & \begin{minipage}[b]{\linewidth}\raggedright
SCHEDULE
\end{minipage} \\
\midrule\noalign{}
\endhead
\bottomrule\noalign{}
\endlastfoot
\end{longtable}

\subsubsection{Description}\label{description-920}

This keyword, GASMONTH, states the month for which subsequent scheduling
events take place within an Annual Schedule section for when the Gas
Field Operations option has been activated by the GASFIELD keyword in
the RUNSPEC section. The keyword must lie in between the GASBEGIN, that
defines the start of an Annual Scheduling section and the GASEND keyword
that ends the section. Optionally, the keyword can be used to write a
report to the print file (*.PRT) at the requested month.

See also the GASBEGIN and GASEND keywords in the RUNSPEC section.

This keyword is not supported by OPM Flow but would change the results
if supported so the simulation will be stopped.

\subsection{GASPERIO -- Advance Simulation by Gas Contract
Period}\label{gasperio-advance-simulation-by-gas-contract-period}

\begin{longtable}[]{@{}
  >{\raggedright\arraybackslash}p{(\linewidth - 14\tabcolsep) * \real{0.1250}}
  >{\raggedright\arraybackslash}p{(\linewidth - 14\tabcolsep) * \real{0.1250}}
  >{\raggedright\arraybackslash}p{(\linewidth - 14\tabcolsep) * \real{0.1250}}
  >{\raggedright\arraybackslash}p{(\linewidth - 14\tabcolsep) * \real{0.1250}}
  >{\raggedright\arraybackslash}p{(\linewidth - 14\tabcolsep) * \real{0.1250}}
  >{\raggedright\arraybackslash}p{(\linewidth - 14\tabcolsep) * \real{0.1250}}
  >{\raggedright\arraybackslash}p{(\linewidth - 14\tabcolsep) * \real{0.1250}}
  >{\raggedright\arraybackslash}p{(\linewidth - 14\tabcolsep) * \real{0.1250}}@{}}
\toprule\noalign{}
\begin{minipage}[b]{\linewidth}\raggedright
RUNSPEC
\end{minipage} & \begin{minipage}[b]{\linewidth}\raggedright
GRID
\end{minipage} & \begin{minipage}[b]{\linewidth}\raggedright
EDIT
\end{minipage} & \begin{minipage}[b]{\linewidth}\raggedright
PROPS
\end{minipage} & \begin{minipage}[b]{\linewidth}\raggedright
REGIONS
\end{minipage} & \begin{minipage}[b]{\linewidth}\raggedright
SOLUTION
\end{minipage} & \begin{minipage}[b]{\linewidth}\raggedright
SUMMARY
\end{minipage} & \begin{minipage}[b]{\linewidth}\raggedright
SCHEDULE
\end{minipage} \\
\midrule\noalign{}
\endhead
\bottomrule\noalign{}
\endlastfoot
\end{longtable}

\subsubsection{Description}\label{description-921}

This keyword advances the simulation over one or more gas contract
periods for when the Gas Field Operations option has been activated by
the GASFIELD keyword in the RUNSPEC section. A contract period in this
case is the period over which the Daily Contract Quantity is fixed, this
can a be year or one or more months. If the contract period is a year
then the GASYEAR keyword in the SCHEDULE section can be used instead of
GASPERIOD.

GASPERIO is an alternative to the DATES, TIME and TSTEP keywords in the
SCHEDULE section that advances the simulation to a given report time
after which additional keywords may be entered to instruct OPM Flow to
perform additional functions via the SCHEDULE section keywords.

This keyword is not supported by OPM Flow but would change the results
if supported so the simulation will be stopped.

\subsection{GASYEAR -- Advance Simulation by Gas Contract
Year}\label{gasyear-advance-simulation-by-gas-contract-year}

\begin{longtable}[]{@{}
  >{\raggedright\arraybackslash}p{(\linewidth - 14\tabcolsep) * \real{0.1250}}
  >{\raggedright\arraybackslash}p{(\linewidth - 14\tabcolsep) * \real{0.1250}}
  >{\raggedright\arraybackslash}p{(\linewidth - 14\tabcolsep) * \real{0.1250}}
  >{\raggedright\arraybackslash}p{(\linewidth - 14\tabcolsep) * \real{0.1250}}
  >{\raggedright\arraybackslash}p{(\linewidth - 14\tabcolsep) * \real{0.1250}}
  >{\raggedright\arraybackslash}p{(\linewidth - 14\tabcolsep) * \real{0.1250}}
  >{\raggedright\arraybackslash}p{(\linewidth - 14\tabcolsep) * \real{0.1250}}
  >{\raggedright\arraybackslash}p{(\linewidth - 14\tabcolsep) * \real{0.1250}}@{}}
\toprule\noalign{}
\begin{minipage}[b]{\linewidth}\raggedright
RUNSPEC
\end{minipage} & \begin{minipage}[b]{\linewidth}\raggedright
GRID
\end{minipage} & \begin{minipage}[b]{\linewidth}\raggedright
EDIT
\end{minipage} & \begin{minipage}[b]{\linewidth}\raggedright
PROPS
\end{minipage} & \begin{minipage}[b]{\linewidth}\raggedright
REGIONS
\end{minipage} & \begin{minipage}[b]{\linewidth}\raggedright
SOLUTION
\end{minipage} & \begin{minipage}[b]{\linewidth}\raggedright
SUMMARY
\end{minipage} & \begin{minipage}[b]{\linewidth}\raggedright
SCHEDULE
\end{minipage} \\
\midrule\noalign{}
\endhead
\bottomrule\noalign{}
\endlastfoot
\end{longtable}

\subsubsection{Description}\label{description-922}

This keyword advances the simulation over one or more gas contract years
for when the Gas Field Operations option has been activated by the
GASFIELD keyword in the RUNSPEC section. A contract year in this case is
the period over which the Daily Contract Quantity is fixed, this can a
be year, this keyword or the GASPERIO keyword in the SCHEDULE section,
or one or more months. If the contract period is over one or more months
then the GASPERIO keyword in the SCHEDULE section can be used instead of
GASYEAR.

GASYEAR is an alternative to the DATES, TIME and TSTEP keywords in the
SCHEDULE section that advances the simulation to a given report time
after which additional keywords may be entered to instruct OPM Flow to
perform additional functions via the SCHEDULE section keywords.

This keyword is not supported by OPM Flow but would change the results
if supported so the simulation will be stopped.

\subsection{GCALECON -- Group Economic Criteria for Production Calorific
Groups}\label{gcalecon-group-economic-criteria-for-production-calorific-groups}

\begin{longtable}[]{@{}
  >{\raggedright\arraybackslash}p{(\linewidth - 14\tabcolsep) * \real{0.1250}}
  >{\raggedright\arraybackslash}p{(\linewidth - 14\tabcolsep) * \real{0.1250}}
  >{\raggedright\arraybackslash}p{(\linewidth - 14\tabcolsep) * \real{0.1250}}
  >{\raggedright\arraybackslash}p{(\linewidth - 14\tabcolsep) * \real{0.1250}}
  >{\raggedright\arraybackslash}p{(\linewidth - 14\tabcolsep) * \real{0.1250}}
  >{\raggedright\arraybackslash}p{(\linewidth - 14\tabcolsep) * \real{0.1250}}
  >{\raggedright\arraybackslash}p{(\linewidth - 14\tabcolsep) * \real{0.1250}}
  >{\raggedright\arraybackslash}p{(\linewidth - 14\tabcolsep) * \real{0.1250}}@{}}
\toprule\noalign{}
\begin{minipage}[b]{\linewidth}\raggedright
RUNSPEC
\end{minipage} & \begin{minipage}[b]{\linewidth}\raggedright
GRID
\end{minipage} & \begin{minipage}[b]{\linewidth}\raggedright
EDIT
\end{minipage} & \begin{minipage}[b]{\linewidth}\raggedright
PROPS
\end{minipage} & \begin{minipage}[b]{\linewidth}\raggedright
REGIONS
\end{minipage} & \begin{minipage}[b]{\linewidth}\raggedright
SOLUTION
\end{minipage} & \begin{minipage}[b]{\linewidth}\raggedright
SUMMARY
\end{minipage} & \begin{minipage}[b]{\linewidth}\raggedright
SCHEDULE
\end{minipage} \\
\midrule\noalign{}
\endhead
\bottomrule\noalign{}
\endlastfoot
\end{longtable}

\subsubsection{Description}\label{description-923}

The GCALECON keyword defines economic criteria for production groups,
including the field level group FIELD, that have previously been defined
by the GCONPROD keyword in the SCHEDULE section and have had their rate
targets and constraints set by calorific value via the GCONVAL keyword
in the SCHEDULE section.

This keyword is not supported by OPM Flow but would change the results
if supported so the simulation will be stopped.

Note that wells are allocated to a group when they are specified by the
WELSPECS keyword and wells can also have economic controls. Wells under
group control are therefore subject to the economic criteria set via the
GCONPROD and CECON keywords in the SCHEDULE section and the controls
specified by the WECON keyword.

\begin{longtable}[]{@{}
  >{\raggedright\arraybackslash}p{(\linewidth - 6\tabcolsep) * \real{0.2500}}
  >{\raggedright\arraybackslash}p{(\linewidth - 6\tabcolsep) * \real{0.2500}}
  >{\raggedright\arraybackslash}p{(\linewidth - 6\tabcolsep) * \real{0.2500}}
  >{\raggedright\arraybackslash}p{(\linewidth - 6\tabcolsep) * \real{0.2500}}@{}}
\toprule\noalign{}
\begin{minipage}[b]{\linewidth}\raggedright
No.
\end{minipage} & \begin{minipage}[b]{\linewidth}\raggedright
Name
\end{minipage} & \begin{minipage}[b]{\linewidth}\raggedright
Description
\end{minipage} & \begin{minipage}[b]{\linewidth}\raggedright
Default
\end{minipage} \\
\midrule\noalign{}
\endhead
\bottomrule\noalign{}
\endlastfoot
Field & Metric & Laboratory & \\
1 & GRPNAME & A character string of up to eight characters in length
that defines the group name for which the group target and constraints
are being defined. The group named FIELD is the top most group and
should be used to set targets and constraints for the field. Note that
the group hierarchy should be defined by the GRUPTREE keyword when there
is more than one level of groups, otherwise all the groups will sit
directly under the FIELD group in the group tree hierarchy. & None \\
2 & ENEVAL & A real positive value that defines the minimum economic
surface energy production rate, below which an economic action of
shutting in or stopping all the wells in the group, as requested by item
(9) of the WELSPECS keyword. A value less than or equal to zero switches
of this criteria. & 0.0 \\
BTU/day & kJ/day & J/hour & \\
3 & CALVAL & A real positive value that defines the minimum economic
surface calorific value, below which an economic action of shutting in
or stopping all the wells in the group, as requested by item (9) of the
WELSPECS keyword. A value less than or equal to zero switches of this
criteria, & 0.0 \\
Btu/Mscf & kJ/sm3 & J/scc & \\
8 & END & A defined character string that defines if the simulation
should terminate if all the producing wells in the group, including the
FIELD group, are shut or stopped. END should be set to one of the
following character strings: & NO \\
Notes: & & & \\
\end{longtable}

\emph{Table 12.26: GCALECON Keyword Description}

See also the WELSPECS keyword to define a wells shut-in or stop options,
GCONPROD, GCONCAL, GCONENG for group controls, and WECON for setting a
well's economic criteria. All the aforementioned keywords are described
in the SCHEDULE section.

\subsubsection{Example}\label{example-521}

The following example defines the economic criteria for the field with a
minimum economic surface energy production rate of 5 x 109 BTU/day and a
minimum economic surface calorific value of900 Btu/Mscf

\begin{verbatim}
--
--       GROUP ECONOMIC CRITERIA FOR PRODUCTION GROUPS UNDER CALORIFIC CONTROL
--
-- GRUP  ENERGY  CALORIFIC  END
-- NAME  RATE    VALUE      RUN
GCALECON
FIELD    5E9     900.0     'YES'                                               /
/
\end{verbatim}

If the economic limits are violated then the run will stop at the next
report time step.

\subsection{GCONCAL -- Group Production Calorific
Targets}\label{gconcal-group-production-calorific-targets}

\begin{longtable}[]{@{}
  >{\raggedright\arraybackslash}p{(\linewidth - 14\tabcolsep) * \real{0.1250}}
  >{\raggedright\arraybackslash}p{(\linewidth - 14\tabcolsep) * \real{0.1250}}
  >{\raggedright\arraybackslash}p{(\linewidth - 14\tabcolsep) * \real{0.1250}}
  >{\raggedright\arraybackslash}p{(\linewidth - 14\tabcolsep) * \real{0.1250}}
  >{\raggedright\arraybackslash}p{(\linewidth - 14\tabcolsep) * \real{0.1250}}
  >{\raggedright\arraybackslash}p{(\linewidth - 14\tabcolsep) * \real{0.1250}}
  >{\raggedright\arraybackslash}p{(\linewidth - 14\tabcolsep) * \real{0.1250}}
  >{\raggedright\arraybackslash}p{(\linewidth - 14\tabcolsep) * \real{0.1250}}@{}}
\toprule\noalign{}
\begin{minipage}[b]{\linewidth}\raggedright
RUNSPEC
\end{minipage} & \begin{minipage}[b]{\linewidth}\raggedright
GRID
\end{minipage} & \begin{minipage}[b]{\linewidth}\raggedright
EDIT
\end{minipage} & \begin{minipage}[b]{\linewidth}\raggedright
PROPS
\end{minipage} & \begin{minipage}[b]{\linewidth}\raggedright
REGIONS
\end{minipage} & \begin{minipage}[b]{\linewidth}\raggedright
SOLUTION
\end{minipage} & \begin{minipage}[b]{\linewidth}\raggedright
SUMMARY
\end{minipage} & \begin{minipage}[b]{\linewidth}\raggedright
SCHEDULE
\end{minipage} \\
\midrule\noalign{}
\endhead
\bottomrule\noalign{}
\endlastfoot
\end{longtable}

\subsubsection{Description}\label{description-924}

The GCONCAL keyword defines calorific production targets and constraints
for groups, including the top most group in the group hierarchy known as
the FIELD group. Wells are allocated to groups when the wells are
specified by the WELSPECS keyword in the SCHEDULE section. Wells defined
to be under group control will have their production rates controlled by
the group to which they belong, in addition to any well constraints
defined for the wells.

This keyword is not supported by OPM Flow but would change the results
if supported so the simulation will be stopped.

\begin{longtable}[]{@{}
  >{\raggedright\arraybackslash}p{(\linewidth - 6\tabcolsep) * \real{0.2500}}
  >{\raggedright\arraybackslash}p{(\linewidth - 6\tabcolsep) * \real{0.2500}}
  >{\raggedright\arraybackslash}p{(\linewidth - 6\tabcolsep) * \real{0.2500}}
  >{\raggedright\arraybackslash}p{(\linewidth - 6\tabcolsep) * \real{0.2500}}@{}}
\toprule\noalign{}
\begin{minipage}[b]{\linewidth}\raggedright
No.
\end{minipage} & \begin{minipage}[b]{\linewidth}\raggedright
Name
\end{minipage} & \begin{minipage}[b]{\linewidth}\raggedright
Description
\end{minipage} & \begin{minipage}[b]{\linewidth}\raggedright
Default
\end{minipage} \\
\midrule\noalign{}
\endhead
\bottomrule\noalign{}
\endlastfoot
Field & Metric & Laboratory & \\
1 & GRPNAME & A character string of up to eight characters in length
that defines the group name for which the group's target calorific value
is being defined. The group named FIELD is the top most group and should
be used to set targets and constraints for the field. Note that the
group hierarchy should be defined by the GRUPTREE keyword when there is
more than one level of groups, otherwise all the groups will sit
directly under the FIELD group in the group tree hierarchy. & None \\
2 & CALVAL & A real positive value that defines the target surface
calorific value of the produced gas. The default value of 1 x 1020
switches off calorific control for the group. & 1.0 x 1020 \\
Btu/Mscf & kJ/sm3 & J/hour & \\
3 & ACTION & A defined character string that defines the action to be
taken if the CALVAL is violated. ACTION should be set to one of the
following character strings: The corrective action takes places at the
end of the time step in which the constraint is violated. & None \\
4 & FACTOR & A real positive value that is less than or equal to one
that defines the amount wells can be scaled back in order to satisfy
CALVAL. Note this assumes that there are wells within the group that are
producing with higher and lower calorific values, and the simulator is
thus able to fine a combination of wells that satisfy the group's CALVAL
target. & None \\
dimensionless & dimensionless & dimensionless & \\
Notes: & & & \\
\end{longtable}

\emph{Table 12.27: GCONCAL Keyword Description}

See also the GRUPTREE keyword to define the hierarchy of the groups
below the FIELD level, the GCONPROD and GCONINJE keywords to define a
group's production and injection rate targets and constraints, the
WCONPROD keyword to define a production well's targets and constraints,
and the WCONINJE keyword to define an injection well's targets and
constraints. All the aforementioned keywords are described in the
SCHEDULE section.

\subsubsection{Example}\label{example-522}

The following example defines the calorific production target for the
field.

\begin{verbatim}
---
--       GROUP CALORIFIC PRODUCTION CONTROLS
--
-- GRUP  CALORIFIC  ACTION  CUT
-- NAME  VALUE              BACK
GCONCAL
FIELD    1010E3     RATE    0.95                                       /
/
\end{verbatim}

Here the calorific production target has been set to 1,010 x 103
Btu/Mscf for the field and if the target cannot be met then the well
rates are reduced by 0.95 at each iteration until the target is
satisfied.

\subsection{GCONENG -- Group Production Energy
Targets}\label{gconeng-group-production-energy-targets}

\begin{longtable}[]{@{}
  >{\raggedright\arraybackslash}p{(\linewidth - 14\tabcolsep) * \real{0.1250}}
  >{\raggedright\arraybackslash}p{(\linewidth - 14\tabcolsep) * \real{0.1250}}
  >{\raggedright\arraybackslash}p{(\linewidth - 14\tabcolsep) * \real{0.1250}}
  >{\raggedright\arraybackslash}p{(\linewidth - 14\tabcolsep) * \real{0.1250}}
  >{\raggedright\arraybackslash}p{(\linewidth - 14\tabcolsep) * \real{0.1250}}
  >{\raggedright\arraybackslash}p{(\linewidth - 14\tabcolsep) * \real{0.1250}}
  >{\raggedright\arraybackslash}p{(\linewidth - 14\tabcolsep) * \real{0.1250}}
  >{\raggedright\arraybackslash}p{(\linewidth - 14\tabcolsep) * \real{0.1250}}@{}}
\toprule\noalign{}
\begin{minipage}[b]{\linewidth}\raggedright
RUNSPEC
\end{minipage} & \begin{minipage}[b]{\linewidth}\raggedright
GRID
\end{minipage} & \begin{minipage}[b]{\linewidth}\raggedright
EDIT
\end{minipage} & \begin{minipage}[b]{\linewidth}\raggedright
PROPS
\end{minipage} & \begin{minipage}[b]{\linewidth}\raggedright
REGIONS
\end{minipage} & \begin{minipage}[b]{\linewidth}\raggedright
SOLUTION
\end{minipage} & \begin{minipage}[b]{\linewidth}\raggedright
SUMMARY
\end{minipage} & \begin{minipage}[b]{\linewidth}\raggedright
SCHEDULE
\end{minipage} \\
\midrule\noalign{}
\endhead
\bottomrule\noalign{}
\endlastfoot
\end{longtable}

\subsubsection{Description}\label{description-925}

The GCONCAL keyword defines energy production targets and constraints
for groups, including the top most group in the group hierarchy known as
the FIELD group. Wells are allocated to groups when the wells are
specified by the WELSPECS keyword in the SCHEDULE section. Wells defined
to be under group control will have their production rates controlled by
the group to which they belong, in addition to any well constraints
defined for the wells.

This keyword is not supported by OPM Flow but would change the results
if supported so the simulation will be stopped.

\begin{longtable}[]{@{}
  >{\raggedright\arraybackslash}p{(\linewidth - 6\tabcolsep) * \real{0.2500}}
  >{\raggedright\arraybackslash}p{(\linewidth - 6\tabcolsep) * \real{0.2500}}
  >{\raggedright\arraybackslash}p{(\linewidth - 6\tabcolsep) * \real{0.2500}}
  >{\raggedright\arraybackslash}p{(\linewidth - 6\tabcolsep) * \real{0.2500}}@{}}
\toprule\noalign{}
\begin{minipage}[b]{\linewidth}\raggedright
No.
\end{minipage} & \begin{minipage}[b]{\linewidth}\raggedright
Name
\end{minipage} & \begin{minipage}[b]{\linewidth}\raggedright
Description
\end{minipage} & \begin{minipage}[b]{\linewidth}\raggedright
Default
\end{minipage} \\
\midrule\noalign{}
\endhead
\bottomrule\noalign{}
\endlastfoot
Field & Metric & Laboratory & \\
1 & GRPNAME & A character string of up to eight characters in length
that defines the group name for which the group's target calorific value
is being defined. The group named FIELD is the top most group and should
be used to set targets and constraints for the field. Note that the
group hierarchy should be defined by the GRUPTREE keyword when there is
more than one level of groups, otherwise all the groups will sit
directly under the FIELD group in the group tree hierarchy. & None \\
2 & ENGVAL & A real positive value that defines the surface energy
target for the group. The default value of 1 x 1020 switches off the
energy target for the group. & 1.0 x 1020 \\
BTU/day & kJ/day & J/hour & \\
Notes: & & & \\
\end{longtable}

\emph{Table 12.28: GCONENG Keyword Description}

See also the GRUPTREE keyword to define the hierarchy of the groups
below the FIELD level, the GCONPROD and GCONINJE keywords to define a
group's production and injection rate targets and constraints, the
WCONPROD keyword to define a production well's targets and constraints,
and the WCONINJE keyword to define an injection well's targets and
constraints. All the aforementioned keywords are described in the
SCHEDULE section.

\subsubsection{Example}\label{example-523}

The following example defines the energy production target for the
field.

\begin{verbatim}
--
--       GROUP ENERGY PRODUCTION CONTROLS
--
-- GRUP  ENERGY
-- NAME  VALUE
GCONENG
FIELD    1010E9                                                        /
/
\end{verbatim}

Here the energy production target has been set to 1,010 x 109 Btu/day
for the field.

\subsection{GCONINJE -- Group Injection Targets and
Constraints}\label{gconinje-group-injection-targets-and-constraints}

\begin{longtable}[]{@{}
  >{\raggedright\arraybackslash}p{(\linewidth - 14\tabcolsep) * \real{0.1250}}
  >{\raggedright\arraybackslash}p{(\linewidth - 14\tabcolsep) * \real{0.1250}}
  >{\raggedright\arraybackslash}p{(\linewidth - 14\tabcolsep) * \real{0.1250}}
  >{\raggedright\arraybackslash}p{(\linewidth - 14\tabcolsep) * \real{0.1250}}
  >{\raggedright\arraybackslash}p{(\linewidth - 14\tabcolsep) * \real{0.1250}}
  >{\raggedright\arraybackslash}p{(\linewidth - 14\tabcolsep) * \real{0.1250}}
  >{\raggedright\arraybackslash}p{(\linewidth - 14\tabcolsep) * \real{0.1250}}
  >{\raggedright\arraybackslash}p{(\linewidth - 14\tabcolsep) * \real{0.1250}}@{}}
\toprule\noalign{}
\begin{minipage}[b]{\linewidth}\raggedright
RUNSPEC
\end{minipage} & \begin{minipage}[b]{\linewidth}\raggedright
GRID
\end{minipage} & \begin{minipage}[b]{\linewidth}\raggedright
EDIT
\end{minipage} & \begin{minipage}[b]{\linewidth}\raggedright
PROPS
\end{minipage} & \begin{minipage}[b]{\linewidth}\raggedright
REGIONS
\end{minipage} & \begin{minipage}[b]{\linewidth}\raggedright
SOLUTION
\end{minipage} & \begin{minipage}[b]{\linewidth}\raggedright
SUMMARY
\end{minipage} & \begin{minipage}[b]{\linewidth}\raggedright
SCHEDULE
\end{minipage} \\
\midrule\noalign{}
\endhead
\bottomrule\noalign{}
\endlastfoot
\end{longtable}

\subsubsection{Description}\label{description-926}

The GCONINJE keyword defines injection targets and constraints for
groups, including the top most group in the group hierarchy known as the
FIELD group. Wells are allocated to groups when the wells are specified
by the WELSPECS keyword in the SCHEDULE section. Wells defined to be
under group control will have their injection rates controlled by the
group to which they belong, in addition to any well constraints defined
for the wells.

\begin{longtable}[]{@{}
  >{\raggedright\arraybackslash}p{(\linewidth - 6\tabcolsep) * \real{0.2500}}
  >{\raggedright\arraybackslash}p{(\linewidth - 6\tabcolsep) * \real{0.2500}}
  >{\raggedright\arraybackslash}p{(\linewidth - 6\tabcolsep) * \real{0.2500}}
  >{\raggedright\arraybackslash}p{(\linewidth - 6\tabcolsep) * \real{0.2500}}@{}}
\toprule\noalign{}
\begin{minipage}[b]{\linewidth}\raggedright
No.
\end{minipage} & \begin{minipage}[b]{\linewidth}\raggedright
Name
\end{minipage} & \begin{minipage}[b]{\linewidth}\raggedright
Description
\end{minipage} & \begin{minipage}[b]{\linewidth}\raggedright
Default
\end{minipage} \\
\midrule\noalign{}
\endhead
\bottomrule\noalign{}
\endlastfoot
Field & Metric & Laboratory & \\
1 & GRPNAME & A character string of up to eight characters in length
that defines the group name for which the group target and constraints
are being defined. The group named FIELD is the top most group and
should be used to set targets and constraints for the whole field. Note
that the group hierarchy should be defined by the GRUPTREE keyword in
the SCHEDULE, when there is more than one level of groups, otherwise all
the groups will sit directly under the FIELD group in the group tree
hierarchy. & None \\
2 & TYPE & A defined character string that defines the type of injection
fluid. TYPE should be set to one of the following character strings: &
None \\
3 & TARGET & A defined character string that sets the target injection
control for the group, all the other phases will therefore act as
constraints. The simulator will attempt to meet the TARGET based on the
phase rate stated in items (4) to (7) on this keyword. TARGET should be
set to one of the following character strings: & None \\
4 & RATE & A real positive value that defines the maximum surface
injection rate target or constraint for the phase declared by the TYPE
variable. This value may be specified using a User Defined Argument
(UDA). & None \\
Liquid stb/d Gas Mscf/d & Liquid sm3/day Gas sm3/day & Liquid scc/hour
Gas scc/hour & \\
5 & RESV & A real positive value that defines the maximum reservoir
volume injection rate target or constraint. This value may be specified
using a User Defined Argument (UDA). Note setting a value here other
than the default means that TYPE, item (2) will be the supplement or
``make up'' phase. & None \\
rtb/d & rm3/day & rcc/hour & \\
6 & REIN & A real positive value that defines the target or constraint
re-injection fraction for the produced phase defined by the TYPE
variable. This value may be specified using a User Defined Argument
(UDA). For example, if TYPE is equal to GAS and REIN is equal to 0.85,
then 85\% of the produced gas will be re-injected. & None \\
dimensionless & dimensionless & dimensionless & \\
7 & VREP & A real positive value that defines the target or constraint
of the voidage replacement ratio based on all the produced fluids. This
value may be specified using a User Defined Argument (UDA). For example,
if TYPE is equal to WAT and VREP is equal to 1.00, then 100\% of the
produced reservoir volume will be re-inject as an equivalent water
volume. Note setting a value here other than the default means that
TYPE, item (2) will be the supplement phase. & None \\
dimensionless & dimensionless & dimensionless & \\
8 & GRPCNTL & A defined character string that determines if this group
is subject to higher level group control. This variable is ignored if
GRPNAME is equal to FIELD. & YES \\
9 & GRPGUIDE & A real positive value that defines a group's injection
guide rate expressed as a dimensionless number. A group requires a value
for GRPGUIDE only if it is required to produce a specified proportion of
a higher level group's rate. Defaulting GRPGUIDE results in the
subordinate groups and wells under guide control having their rates
dictated by any higher level groups under guide rate control. In other
words the GRPNAME is masked out. Setting GRPGUIDE to a real positive
value and GUIPHASE to either RATE or RESV will result in a constant
injection guide rate. & None \\
dimensionless & dimensionless & dimensionless & \\
10 & GUIPHASE & A defined character string that sets the guide phase to
which the guide rate in item (9) applies. GUIPHASE should be set to one
of the following character strings: OPM Flow now supports all guide rate
options. The default value of 1* means that the group has no injection
guide rate for this phase. & None \\
11 & GRPREIN & A character string of up to eight characters in length
that defines the group name whose production rate should be used for
applying the REIN quantity to be injected into GRPNAME. This variable is
used to re-inject the REIN production faction from another group
(GRPREIN) via this group (GRPNAME). If GRPREIN is defaulted then the
re-injection quantity for GRPNAME will be based on the production from
GRPNAME itself. & GRPNAME \\
12 & GRPVREP & A character string of up to eight characters in length
that defines the group name whose production rate should be used for
applying the VREP quantity to be injected into GRPNAME. This variable is
used to re-inject the VREP production faction from another group
(GRPVREP) via this group (GRPNAME). If GRPVREP is defaulted then the
voidage quantity for GRPNAME will be based on the production from
GRPNAME itself. & GRPNAME \\
13 & WGASRATE & Wet gas injection rate used in the commercial
compositional simulator. Not used and should be defaulted with 1*. &
1* \\
Notes: & & & \\
\end{longtable}

\emph{Table 12.29: GCONINJE Keyword Description}

See also the GRUPTREE keyword to define the hierarchy of the groups
below the FIELD level, the GCONPROD keyword to define a group's
production targets and constraints, the WCONPROD keyword to define a
production well's targets and constraints, and the WCONINJE keyword to
define an injection well's targets and constraints. All the
aforementioned keywords are described in the SCHEDULE section.

\subsubsection{Example}\label{example-524}

The following example defines the injection targets and constraints for
the field and two groups that are one level below the field group, since
the GRUPTREE keyword has not been entered to define the group hierarchy.

\begin{verbatim}
--
--       GROUP INJECTION TARGETS AND CONSTRAINTS
--
-- GRUP  FLUID CNTL   SURF   RESV   REINJ  VOID  GRUP  GUIDE  GUIDE GRUP  GRUP
-- NAME  TYPE  MODE   RATE   RATE   FRAC   FRAC  CNTL  RATE   DEF   REINJ RESV
GCONINJE
FIELD    WAT   VREP   35E3   1*     1*     1*     NO   1*     1*    1*    1*   /
GRP01    WAT   VREP   1*     1*     1*     1.0    YES  1*     1*    1*    1*   /
GRPO2    WAT   VREP   1*     1*     1*     1.0    YES  1*     1*    1*    1*   /
/
\end{verbatim}

In this example, group GRP01 and GRP02 are injecting water via voidage
replacement with a voidage replacement of one and are under the control
on the field group, that imposes a 35,000 m3/day total water injection
limit.

\subsection{GCONPRI -- Group Production Priority Targets and
Constraints}\label{gconpri-group-production-priority-targets-and-constraints}

\begin{longtable}[]{@{}
  >{\raggedright\arraybackslash}p{(\linewidth - 14\tabcolsep) * \real{0.1250}}
  >{\raggedright\arraybackslash}p{(\linewidth - 14\tabcolsep) * \real{0.1250}}
  >{\raggedright\arraybackslash}p{(\linewidth - 14\tabcolsep) * \real{0.1250}}
  >{\raggedright\arraybackslash}p{(\linewidth - 14\tabcolsep) * \real{0.1250}}
  >{\raggedright\arraybackslash}p{(\linewidth - 14\tabcolsep) * \real{0.1250}}
  >{\raggedright\arraybackslash}p{(\linewidth - 14\tabcolsep) * \real{0.1250}}
  >{\raggedright\arraybackslash}p{(\linewidth - 14\tabcolsep) * \real{0.1250}}
  >{\raggedright\arraybackslash}p{(\linewidth - 14\tabcolsep) * \real{0.1250}}@{}}
\toprule\noalign{}
\begin{minipage}[b]{\linewidth}\raggedright
RUNSPEC
\end{minipage} & \begin{minipage}[b]{\linewidth}\raggedright
GRID
\end{minipage} & \begin{minipage}[b]{\linewidth}\raggedright
EDIT
\end{minipage} & \begin{minipage}[b]{\linewidth}\raggedright
PROPS
\end{minipage} & \begin{minipage}[b]{\linewidth}\raggedright
REGIONS
\end{minipage} & \begin{minipage}[b]{\linewidth}\raggedright
SOLUTION
\end{minipage} & \begin{minipage}[b]{\linewidth}\raggedright
SUMMARY
\end{minipage} & \begin{minipage}[b]{\linewidth}\raggedright
SCHEDULE
\end{minipage} \\
\midrule\noalign{}
\endhead
\bottomrule\noalign{}
\endlastfoot
\end{longtable}

\subsubsection{Description}\label{description-927}

The GCONPRI keyword defines production targets and constraints for
groups, including the top most group in the group hierarchy known as the
FIELD group, for when groups and their associated wells are operating
under priority control as oppose to guide rate control. Priority control
is activated by the PRIORITY keyword in the SCHEDULE section. Wells are
allocated to groups when the wells are specified by the WELSPECS keyword
in the SCHEDULE section. Wells defined to be under group control will
have their production rates controlled by the group to which they
belong, in addition to any well constraints defined for the wells.

This keyword is not supported by OPM Flow but would change the results
if supported so the simulation will be stopped.

\begin{longtable}[]{@{}
  >{\raggedright\arraybackslash}p{(\linewidth - 6\tabcolsep) * \real{0.2500}}
  >{\raggedright\arraybackslash}p{(\linewidth - 6\tabcolsep) * \real{0.2500}}
  >{\raggedright\arraybackslash}p{(\linewidth - 6\tabcolsep) * \real{0.2500}}
  >{\raggedright\arraybackslash}p{(\linewidth - 6\tabcolsep) * \real{0.2500}}@{}}
\toprule\noalign{}
\begin{minipage}[b]{\linewidth}\raggedright
No.
\end{minipage} & \begin{minipage}[b]{\linewidth}\raggedright
Name
\end{minipage} & \begin{minipage}[b]{\linewidth}\raggedright
Description
\end{minipage} & \begin{minipage}[b]{\linewidth}\raggedright
Default
\end{minipage} \\
\midrule\noalign{}
\endhead
\bottomrule\noalign{}
\endlastfoot
Field & Metric & Laboratory & \\
1 & GRPNAME & A character string of up to eight characters in length
that defines the group name for which the group target and constraints
are being defined. The group named FIELD is the top most group and
should be used to set targets and constraints for the field. Note that
the group hierarchy should be defined by the GRUPTREE keyword when there
is more than one level of groups, otherwise all the groups will sit
directly under the FIELD group in the group tree hierarchy. & None \\
2 & ORAT & A real positive value that defines the maximum surface oil
production rate constraint. This value may be specified using a User
Defined Argument (UDA). & None \\
stb/d & sm3/day & scc/hour & \\
3 & OILACT & A defined character string that defines the action to be
taken if ORAT is exceeded. OILACT should be set to one of the following
character strings: & ``NONE'' \\
4 & WRAT & A real positive value that defines the maximum surface water
production rate constraint. This value may be specified using a User
Defined Argument (UDA). & None \\
stb/d & sm3/day & scc/hour & \\
5 & WATACT & A defined character string that defines the action to be
taken if WRAT is exceeded. WATACT should be set to a character string
described by the OILACT parameter on this record. & None \\
6 & GRAT & A real positive value that defines the maximum surface gas
production rate constraint This value may be specified using a User
Defined Argument (UDA). & None \\
Mscf/d & sm3/day & scc/hour & \\
7 & GASACT & A defined character string that defines the action to be
taken if GRAT is exceeded. GASACT should be set to a character string
described by the OILACT parameter on this record. & None \\
8 & LRAT & A real positive value that defines the maximum surface liquid
(oil plus water) production rate constraint. This value may be specified
using a User Defined Argument (UDA). & None \\
stb/d & sm3/day & scc/hour & \\
9 & LIQACT & A defined character string that defines the action to be
taken if LRAT is exceeded. LIQACT should be set to a character string
described by the OILACT parameter on this record. & None \\
10 & RESV & A real positive value that defines the maximum reservoir
volume production rate constraint. This value may be specified using a
User Defined Argument (UDA). & None \\
rb/d & rm3/day & rcc/hour & \\
11 & RESVFRAC & A real positive value that defines a group's maximum
production balancing fraction constraint. This value may be specified
using a User Defined Argument (UDA). & None \\
dimensionless & dimensionless & dimensionless & \\
12 & & Not used should be defaulted with 1*. & 1* \\
13 & & Not used should be defaulted with 1*. & 1* \\
14 & & Not used should be defaulted with 1*. & 1* \\
15 & & Not used should be defaulted with 1*. & 1* \\
16 & LINCOMB & A real positive value that defines the linearly combined
maximum surface target rate or constraint, as per the LINCOM keyword in
the SCHEDULE section. & 1* \\
17 & LINACT & A defined character string that defines the action to be
taken if LINCOMB is exceeded. LINACT should be set to a character string
described by the OILACT parameter on this record. & None \\
Notes: & & & \\
\end{longtable}

\emph{Table 12.30: GCONPRI Keyword Description}

See also the GRUPTREE keyword to define the hierarchy of the groups
below the FIELD level, the GCONPROD keyword to define a group's
production targets and constraints, the GCONINJE keyword to define a
group's injection targets and constraints, the WCONPROD keyword to
define a production well's targets and constraints, and the WCONINJE
keyword to define an injection well's targets and constraints. All the
aforementioned keywords are described in the SCHEDULE section.

\subsubsection{Example}\label{example-525}

The following example defines the production targets and constraints for
the field and two groups that are one level below the field group, since
the GRUPTREE keyword has not been entered to define the group hierarchy.

\begin{verbatim}
--
--       GROUP PRIORITY PRODUCTION CONTROLS
--
-- GRUP  ORAT   ORAT   WRAT   WRAT   GRAT   GRAT   LRAT   LRAT   RVOL   RVOL
-- NAME  LIMT   ACTN   LIMT   ACTN   LIMT   ACTN   LIMT   ACTN   LIMT   ACTN
GCONPRI
FIELD    40E3   PRI    30E3   +CON   125E3  PRI    1*     1*     1*     1*     /
GRP01    25E3   PRI    1*     1*     1*     1*     1*     1*     1*     1*     /
GRP02    25E3   PRI    1*     1*     1*     1*     1*     1*     1*     1*     /
/
\end{verbatim}

All groups are controlled by oil constraints, but only the field level
has water and gas constraints to reflect the actual production facility
constraints. The oil production constraint of 40,000, 25,000 and 25,000
stb/d are defined for the field, GRP01 and GRP02 groups, respectively.
If the oil rate constraint is exceeded then the wells will be controlled
using the priority formulae one, as defined on the PRIORITY keyword in
the SCHEDULE section. Similarly for the field, for when the gas
constraint is exceeded. Finally, if the field water constraint is
surpassed then the worst offending connection and below in the worst
offending well are shut.

\subsection{GCONPROD -- Group Production Targets and
Constraints}\label{gconprod-group-production-targets-and-constraints}

\begin{longtable}[]{@{}
  >{\raggedright\arraybackslash}p{(\linewidth - 14\tabcolsep) * \real{0.1250}}
  >{\raggedright\arraybackslash}p{(\linewidth - 14\tabcolsep) * \real{0.1250}}
  >{\raggedright\arraybackslash}p{(\linewidth - 14\tabcolsep) * \real{0.1250}}
  >{\raggedright\arraybackslash}p{(\linewidth - 14\tabcolsep) * \real{0.1250}}
  >{\raggedright\arraybackslash}p{(\linewidth - 14\tabcolsep) * \real{0.1250}}
  >{\raggedright\arraybackslash}p{(\linewidth - 14\tabcolsep) * \real{0.1250}}
  >{\raggedright\arraybackslash}p{(\linewidth - 14\tabcolsep) * \real{0.1250}}
  >{\raggedright\arraybackslash}p{(\linewidth - 14\tabcolsep) * \real{0.1250}}@{}}
\toprule\noalign{}
\begin{minipage}[b]{\linewidth}\raggedright
RUNSPEC
\end{minipage} & \begin{minipage}[b]{\linewidth}\raggedright
GRID
\end{minipage} & \begin{minipage}[b]{\linewidth}\raggedright
EDIT
\end{minipage} & \begin{minipage}[b]{\linewidth}\raggedright
PROPS
\end{minipage} & \begin{minipage}[b]{\linewidth}\raggedright
REGIONS
\end{minipage} & \begin{minipage}[b]{\linewidth}\raggedright
SOLUTION
\end{minipage} & \begin{minipage}[b]{\linewidth}\raggedright
SUMMARY
\end{minipage} & \begin{minipage}[b]{\linewidth}\raggedright
SCHEDULE
\end{minipage} \\
\midrule\noalign{}
\endhead
\bottomrule\noalign{}
\endlastfoot
\end{longtable}

\subsubsection{Description}\label{description-928}

The GCONPROD keyword defines production targets and constraints for
groups, including the top most group in the group hierarchy known as the
FIELD group. Wells are allocated to groups when the wells are specified
by the WELSPECS keyword in the SCHEDULE section. Wells defined to be
under group control will have their production rates controlled by the
group to which they belong, in addition to any well constraints defined
for the wells.

\begin{longtable}[]{@{}
  >{\raggedright\arraybackslash}p{(\linewidth - 6\tabcolsep) * \real{0.2500}}
  >{\raggedright\arraybackslash}p{(\linewidth - 6\tabcolsep) * \real{0.2500}}
  >{\raggedright\arraybackslash}p{(\linewidth - 6\tabcolsep) * \real{0.2500}}
  >{\raggedright\arraybackslash}p{(\linewidth - 6\tabcolsep) * \real{0.2500}}@{}}
\toprule\noalign{}
\begin{minipage}[b]{\linewidth}\raggedright
No.
\end{minipage} & \begin{minipage}[b]{\linewidth}\raggedright
Name
\end{minipage} & \begin{minipage}[b]{\linewidth}\raggedright
Description
\end{minipage} & \begin{minipage}[b]{\linewidth}\raggedright
Default
\end{minipage} \\
\midrule\noalign{}
\endhead
\bottomrule\noalign{}
\endlastfoot
Field & Metric & Laboratory & \\
1 & GRPNAME & A character string of up to eight characters in length
that defines the group name for which the group target and constraints
are being defined. The group named FIELD is the top most group and
should be used to set targets and constraints for the field. Note that
the group hierarchy should be defined by the GRUPTREE keyword in the
SCHEDULE section, when there is more than one level of groups, otherwise
all the groups will sit directly under the FIELD group in the group tree
hierarchy. & None \\
2 & TARGET & A defined character string that specifies the production
rate control mode for the group. The simulator will attempt to meet the
TARGET rate as defined by remaining items on this keyword. TARGET should
be set to one of the following character strings: Note that the
commercial simulator includes additional options (not listed above) that
are not currently supported by OPM Flow. & None \\
3 & ORAT & A real positive value that defines the maximum surface oil
production rate target or constraint. This value may be specified using
a User Defined Argument (UDA). & None \\
stb/d & sm3/day & scc/hour & \\
4 & WRAT & A real positive value that defines the maximum surface water
production rate target or constraint. This value may be specified using
a User Defined Argument (UDA). & None \\
stb/d & sm3/day & scc/hour & \\
5 & GRAT & A real positive value that defines the maximum surface gas
production rate target or constraint. This value may be specified using
a User Defined Argument (UDA). & None \\
Mscf/d & sm3/day & scc/hour & \\
6 & LRAT & A real positive value that defines the maximum surface liquid
(oil plus water) production rate target or constraint. This value may be
specified using a User Defined Argument (UDA). & None \\
stb/d & sm3/day & scc/hour & \\
7 & ACTION & A defined character string that specifies the action to be
taken if the constraints in (3) to (6) are violated. ACTION should be
set to one of the following character strings: The corrective action
takes places at the end of the time step in which the constraint is
violated. Note that only the NONE, WELL and RATE options are currently
supported by OPM Flow. & None \\
8 & GRPCNTL & A defined character string that determines if this group
is subject to higher level group control. GRPCNTL will be ignored for
the FIELD group. & None \\
9 & GRPGUIDE & A real positive value that defines a group's production
guide rate expressed as a dimensionless number. A group requires a value
for GRPGUIDE only if it is required to produce a specified proportion of
a higher level group's rate & None \\
dimensionless & dimensionless & dimensionless & \\
10 & GUIPHASE & A defined character string that sets the guide phase to
which the guide rate in item (9) applies. GUIPHASE should be set to one
of the following character strings: Note that only the OIL, WAT, GAS,
LIQ and default options are currently supported by OPM Flow. & 1* \\
11 & ACTWAT & A defined character string that defines the action to be
taken if the WRAT constraint, item (4), is violated. ACTWAT should be
set to one of the following character strings: If defaulted then
procedure defined by ACTION, item (7), is applied. The corrective action
takes places at the end of the time step in which the constraint is
violated. Note that only the NONE and RATE options are currently
supported by OPM Flow. & 1* \\
12 & ACTGAS & A defined character string that defines the action to be
taken if the GRAT constraint, item (5), is violated. ACTGAS should be
set to one of the following character strings: If defaulted then
procedure defined by ACTION, item (7), is applied. The corrective action
takes places at the end of the time step in which the constraint is
violated Note that only the NONE and RATE options are currently
supported by OPM Flow. & 1* \\
13 & ACTLIQ & A defined character string that defines the action to be
taken if the LRAT constraint, item (6), is violated. ACLIQT should be
set to one of the following character strings: If defaulted then
procedure defined by ACTION, item (7), is applied. The corrective action
takes places at the end of the time step in which the constraint is
violated. Note that only the NONE and RATE options are currently
supported by OPM Flow. & 1* \\
14 & RESV & A real positive value that defines the maximum reservoir
volume production rate target or constraint. This value may be specified
using a User Defined Argument (UDA). & None \\
rb/d & rm3/day & rcc/hour & \\
15 & RESVFRAC & A real positive value that defines the maximum reservoir
volume production balancing fraction. Not used and should be defaulted
with 1*. & 1* \\
16 & WGASRATE & Wet gas production rate used in the commercial
compositional simulator. Not used and should be defaulted with 1*. &
1* \\
17 & CALRATE & Calorific production rate used in the commercial
compositional simulator. Not used and should be defaulted with 1*. &
1* \\
18 & GASFRAC & Surface gas production fraction used in the commercial
compositional simulator Not used and should be defaulted with 1*. &
1* \\
19 & WATFRAC & Surface water production fraction used in the commercial
compositional simulator. Not used and should be defaulted with 1*. &
1* \\
20 & COMBRATE & Linearly combined production rate used in the commercial
compositional simulator. Not used and should be defaulted with 1*. &
1* \\
21 & COMBPROC & Linearly combined procure for when exceeding COMBRATE,
used in the commercial black-oil simulator. Not used and should be
defaulted with 1*. & 1* \\
Notes: & & & \\
\end{longtable}

\emph{Table 12.3.89.1: GCONPROD Keyword Description}

See also the GRUPTREE keyword to define the hierarchy of the groups
below the FIELD level, the GCONINJE keyword to define a group's
injection targets and constraints, the WCONPROD keyword to define a
production well's targets and constraints, and the WCONINJE keyword to
define an injection well's targets and constraints. All the
aforementioned keywords are described in the SCHEDULE section.

\subsubsection{Example}\label{example-526}

The following example defines the production targets and constraints for
the field and two groups that are one level below the field group, since
the GRUPTREE keyword has not been entered to define the group hierarchy.

\begin{verbatim}
--
--       GROUP PRODUCTION CONTROLS
--
-- GRUP  CNTL  OIL    WAT    GAS    LIQ    CNTL  GRUP  GUIDE  GUIDE  CNTL
-- NAME  MODE  RATE   RATE   RATE   RATE   OPT   CNTL  RATE   DEF    WAT
GCONPROD
FIELD    ORAT  40E3   60E3   300E3  60E3   1*     1*    1*     1*     1*      /
GRP01    FLD   25E3   1*     1*     1*     1*     1*    1*     1*     1*      /
GRP02    FLD   25E3   1*     1*     1*     1*     1*    1*     1*     1*      /
/
\end{verbatim}

All groups are controlled by oil rate targets or constraints, but only
the field level has water, gas and liquid constraints to reflect the
actual production facility constraints. The wells under group control
will be produced based on oil potential of each of the wells under group
control, such that the field oil production target of 40,000 stb/d is
honored and subject to the other phase fluid constraints. In addition,
GRP01 and GRP02 oil rate values of 25,000 stb/d are constraints as these
two groups are subject to the FIELD level targets and constraints.

\subsection{GCONSALE -- Define Group Sales Gas Production Targets and
Constraints}\label{gconsale-define-group-sales-gas-production-targets-and-constraints}

\begin{longtable}[]{@{}
  >{\raggedright\arraybackslash}p{(\linewidth - 14\tabcolsep) * \real{0.1250}}
  >{\raggedright\arraybackslash}p{(\linewidth - 14\tabcolsep) * \real{0.1250}}
  >{\raggedright\arraybackslash}p{(\linewidth - 14\tabcolsep) * \real{0.1250}}
  >{\raggedright\arraybackslash}p{(\linewidth - 14\tabcolsep) * \real{0.1250}}
  >{\raggedright\arraybackslash}p{(\linewidth - 14\tabcolsep) * \real{0.1250}}
  >{\raggedright\arraybackslash}p{(\linewidth - 14\tabcolsep) * \real{0.1250}}
  >{\raggedright\arraybackslash}p{(\linewidth - 14\tabcolsep) * \real{0.1250}}
  >{\raggedright\arraybackslash}p{(\linewidth - 14\tabcolsep) * \real{0.1250}}@{}}
\toprule\noalign{}
\begin{minipage}[b]{\linewidth}\raggedright
RUNSPEC
\end{minipage} & \begin{minipage}[b]{\linewidth}\raggedright
GRID
\end{minipage} & \begin{minipage}[b]{\linewidth}\raggedright
EDIT
\end{minipage} & \begin{minipage}[b]{\linewidth}\raggedright
PROPS
\end{minipage} & \begin{minipage}[b]{\linewidth}\raggedright
REGIONS
\end{minipage} & \begin{minipage}[b]{\linewidth}\raggedright
SOLUTION
\end{minipage} & \begin{minipage}[b]{\linewidth}\raggedright
SUMMARY
\end{minipage} & \begin{minipage}[b]{\linewidth}\raggedright
SCHEDULE
\end{minipage} \\
\midrule\noalign{}
\endhead
\bottomrule\noalign{}
\endlastfoot
\end{longtable}

\subsubsection{Description}\label{description-929}

GCONSALE defines group sales gas production targets and constraints for
when the gas production from an oil field group is exported under a Gas
Sales Agreement (``GSA'') and the oil field group also has oil
production targets and constraints.

Note that the keyword should not be used to control sales gas for a gas
field group, as the gas injection rate is used to control the sales gas
production with this keyword, that is:

\begin{equation}\phantomsection\label{eq-12-23}{
\begin{matrix}\text{Gas Sales Rate} = \text{Total Group Gas Production Rate} \\  - \text{Group Gas Injection Rate} \\  + \text{Total Group Gas Import Rate} \\  - \text{Total Group Gas Consumption}\end{matrix}
}\end{equation}

Thus, surplus gas that cannot be sold is re-injected, which requires
that there are active gas injectors in the model that are subordinate to
groups with gas sales targets. Note that the surplus gas re-injection
rates are automatically calculated by OPM Flow at each time step.

\begin{longtable}[]{@{}
  >{\raggedright\arraybackslash}p{(\linewidth - 6\tabcolsep) * \real{0.2500}}
  >{\raggedright\arraybackslash}p{(\linewidth - 6\tabcolsep) * \real{0.2500}}
  >{\raggedright\arraybackslash}p{(\linewidth - 6\tabcolsep) * \real{0.2500}}
  >{\raggedright\arraybackslash}p{(\linewidth - 6\tabcolsep) * \real{0.2500}}@{}}
\toprule\noalign{}
\begin{minipage}[b]{\linewidth}\raggedright
No.
\end{minipage} & \begin{minipage}[b]{\linewidth}\raggedright
Name
\end{minipage} & \begin{minipage}[b]{\linewidth}\raggedright
Description
\end{minipage} & \begin{minipage}[b]{\linewidth}\raggedright
Default
\end{minipage} \\
\midrule\noalign{}
\endhead
\bottomrule\noalign{}
\endlastfoot
Field & Metric & Laboratory & \\
1 & GRPNAME & A character string of up to eight characters in length
that defines the group name for which the group gas sales target and
constraints are being defined. The group named FIELD is the top most
group and should be used to set the gas sales targets and constraints
for the field. Note that the group hierarchy should be defined by the
GRUPTREE keyword when there is more than one level of groups, otherwise
all the groups will sit directly under the FIELD group in the group tree
hierarchy. & None \\
2 & GSALE & GSALE should either be set to: This value may be specified
using a User Defined Argument (UDA). Note that if GSALE has been set to
switch off both gas sales and gas re-injection, then the GCONINJE
keyword in the SCHEDULE section may be used to re-enable gas
re-injection again. & None \\
Mscf/d & sm3/day & scc/hour & \\
3 & GSALEMAX & A real positive value that must be greater than GSALE
that defines the maximum allowed gas sales rate. If GSALE exceeds
GSALEMAX then the action defined by the ACTION variable on this keyword
is implemented at the end of the current time step. This value may be
specified using a User Defined Argument (UDA). & 1 x 1020 \\
Mscf/d & sm3/day & scc/hour & \\
4 & GSALEMIN & A real positive value that must be less than GSALE that
defines the minimum allowed gas sales rate. If GSALE is less than
GSALEMIN then one of the following actions will be implemented at the
end of the current time step: If none of the above actions can be
implemented then the minimum gas sales rate will not be satisfied. This
value may be specified using a User Defined Argument (UDA). & -1 x
1020 \\
5 & ACTION & A defined character string that defines the action to be
taken if the maximum gas sales rate, GSALEMAX, is violated. ACTION
should be set to one of the following character strings: The corrective
action takes places at the end of the time step in which the constraint
is violated. & None \\
Notes: & & & \\
\end{longtable}

\emph{Table 12.32: GCONSALE Keyword Description}

GCONSALE can also be used with the GCONINJE keyword in the SCHEDULE
section in order to apply additional limits, for example by applying a
maximum group injection rate. In this scenario, the TARGET variable on
the GCONINJE keyword must be set to ``REIN'', and if desired, a
re-injection fraction (REIN on the GCONINJE keyword), or any other
constraint.

See also the GCONSUMP in the SCHEDULE section that defines the fuel gas
requirements for groups.

\subsubsection{Example}\label{example-527}

The following examples sets the field gas sales target rate:

\begin{verbatim}
--
--       GROUP GAS SALES FOR OIL FIELDS
--
-- GRUP  GAS    MAX    MIN    CNTL
-- NAME  SALES  RATE   RATE   ACTN
GCONSALE
FIELD    40E3   50E3   20E3   RATE                                            /
/
\end{verbatim}

Here the field has a gas sales target of 40 MMscf/d, with a maximum rate
of 50 MMscf/d and a minimum of 20 MMscf/d.~If the maximum gas sales rate
is exceeded then the group's gas production rate target is reduced to
equal GSALEMAX, after accounting for fuel gas and the current rate of
re-injection. This will also place the group on gas production control.

\begin{verbatim}
\end{verbatim}

\subsection{GCONSUMP -- Define Group Gas Consumption and Gas Import
Targets}\label{gconsump-define-group-gas-consumption-and-gas-import-targets}

\begin{longtable}[]{@{}
  >{\raggedright\arraybackslash}p{(\linewidth - 14\tabcolsep) * \real{0.1250}}
  >{\raggedright\arraybackslash}p{(\linewidth - 14\tabcolsep) * \real{0.1250}}
  >{\raggedright\arraybackslash}p{(\linewidth - 14\tabcolsep) * \real{0.1250}}
  >{\raggedright\arraybackslash}p{(\linewidth - 14\tabcolsep) * \real{0.1250}}
  >{\raggedright\arraybackslash}p{(\linewidth - 14\tabcolsep) * \real{0.1250}}
  >{\raggedright\arraybackslash}p{(\linewidth - 14\tabcolsep) * \real{0.1250}}
  >{\raggedright\arraybackslash}p{(\linewidth - 14\tabcolsep) * \real{0.1250}}
  >{\raggedright\arraybackslash}p{(\linewidth - 14\tabcolsep) * \real{0.1250}}@{}}
\toprule\noalign{}
\begin{minipage}[b]{\linewidth}\raggedright
RUNSPEC
\end{minipage} & \begin{minipage}[b]{\linewidth}\raggedright
GRID
\end{minipage} & \begin{minipage}[b]{\linewidth}\raggedright
EDIT
\end{minipage} & \begin{minipage}[b]{\linewidth}\raggedright
PROPS
\end{minipage} & \begin{minipage}[b]{\linewidth}\raggedright
REGIONS
\end{minipage} & \begin{minipage}[b]{\linewidth}\raggedright
SOLUTION
\end{minipage} & \begin{minipage}[b]{\linewidth}\raggedright
SUMMARY
\end{minipage} & \begin{minipage}[b]{\linewidth}\raggedright
SCHEDULE
\end{minipage} \\
\midrule\noalign{}
\endhead
\bottomrule\noalign{}
\endlastfoot
\end{longtable}

\subsubsection{Description}\label{description-930}

GCONSUMP defines the group gas consumption rate either as an actual rate
or as a percentage of the group's production. In both oil and gas fields
produced gas is commonly used as fuel to support the processing and
utility facilities needed to run the plant.

In addition to defining gas consumption, the keyword can also be used to
define the group's gas import rate, if required. This is used to import
gas into the model from other sources (fields, reservoirs etc.) that are
not included in the current run. For example, if a several fields are
supplying gas to a power plant (field A, B and C), but only one is being
modeled in the current import deck (A), then production from the other
two fields (B and C) can be incorporated into model in order to meet the
plant demand. Note in this case the import gas rates from fields B and C
are fixed, and therefore field A acts like a ``swing'' producer to match
the gas demand target.

\begin{longtable}[]{@{}
  >{\raggedright\arraybackslash}p{(\linewidth - 6\tabcolsep) * \real{0.2500}}
  >{\raggedright\arraybackslash}p{(\linewidth - 6\tabcolsep) * \real{0.2500}}
  >{\raggedright\arraybackslash}p{(\linewidth - 6\tabcolsep) * \real{0.2500}}
  >{\raggedright\arraybackslash}p{(\linewidth - 6\tabcolsep) * \real{0.2500}}@{}}
\toprule\noalign{}
\begin{minipage}[b]{\linewidth}\raggedright
No.
\end{minipage} & \begin{minipage}[b]{\linewidth}\raggedright
Name
\end{minipage} & \begin{minipage}[b]{\linewidth}\raggedright
Description
\end{minipage} & \begin{minipage}[b]{\linewidth}\raggedright
Default
\end{minipage} \\
\midrule\noalign{}
\endhead
\bottomrule\noalign{}
\endlastfoot
Field & Metric & Laboratory & \\
1 & GRPNAME & A character string of up to eight characters in length
that defines the group name for which the group gas consumption is being
defined. The group named FIELD is the top most group and should be used
to set the fuel consumption for the field. Note that the group hierarchy
should be defined by the GRUPTREE keyword when there is more than one
level of groups, otherwise all the groups will sit directly under the
FIELD group in the group tree hierarchy. & None \\
2 & GASFUEL & A real value that defines the gas consumption, that is the
fuel gas consumed by the group, either defined as a volumetric rate or
as a fraction of the group's gas production. The two options are
implemented by: This value may be specified using a User Defined
Argument (UDA). & 0.0 \\
Mscf/d & sm3/day & scc/hour & \\
3 & GASIMP & A real positive value greater than zero that defines the
amount of gas to be imported into the group. This value may be specified
using a User Defined Argument (UDA). This option is currently not
supported by OPM Flow & 0.0 \\
stb/d & sm3/day & scc/hour & \\
4 & GASNODE & A character string of up to eight characters in length
that defines the network node in the Extended Network Model, for which
the fuel gas should be removed (GASFUEL) or the imported gas (GASIMP)
assigned. This option is currently not supported by OPM Flow & None \\
Notes: & & & \\
\end{longtable}

\emph{Table 12.33: GCONSUMP Keyword Description}

If the group is acting under Group Gas Sales control via the GCONSALE
keyword in the SCHEDULE section, then the sales gas is calculated by:

\begin{equation}\phantomsection\label{eq-12-24}{
\begin{matrix}\text{Gas Sales Rate} = \text{Total Group Gas Production Rate} \\  - \text{Group Gas Injection Rate} \\  + \text{Total Group Gas Import Rate} \\  - \text{Total Group Gas Consumption}\end{matrix}
}\end{equation}

If the group is acting under Group Gas Re-Injection control via the
GCONINJE keyword in the SCHEDULE section, then the group gas injection
rate calculated by:

\begin{equation}\phantomsection\label{eq-12-25}{
\begin{matrix}\text{Group Gas Injection Rate} = \text{Group Gas Injection Rate} \times  \text{Group Re-Injection Fraction} \\  + \text{Total Group Gas Import Rate} \\  - \text{Total Group Gas Consumption}\end{matrix}
}\end{equation}

\begin{tcolorbox}[enhanced jigsaw, title=\textcolor{quarto-callout-note-color}{\faInfo}\hspace{0.5em}{Note}, coltitle=black, breakable, toprule=.15mm, colbacktitle=quarto-callout-note-color!10!white, leftrule=.75mm, opacitybacktitle=0.6, colback=white, bottomtitle=1mm, toptitle=1mm, titlerule=0mm, left=2mm, colframe=quarto-callout-note-color-frame, arc=.35mm, rightrule=.15mm, bottomrule=.15mm, opacityback=0]

In oil fields with no gas compression typical values of fuel gas range
from three to five percent.

\end{tcolorbox}

\subsubsection{Example}\label{example-528}

The first example sets the fuel gas consumption to 3.0 MMscf/d for the
field.

\begin{verbatim}
--
--       GROUP GAS CONSUMPTION (FUEL) AND IMPORT
--
-- GRUP   GAS      GAS
-- NAME   FUEL    IMPORT
--       ------   -------
GCONSUMP
FIELD    3.0E3                                                                 /
/
\end{verbatim}

The second example sets group PLAT-WST's fuel consumption to be 5\% of
the platform's produced gas and group PLAT-EST's to a constant 1.0
MMscf/d.

\begin{verbatim}
--
--       GROUP GAS CONSUMPTION (FUEL) AND IMPORT
--
-- GRUP   GAS      GAS
-- NAME   FUEL    IMPORT
--       ------   -------
GCONSUMP
PLAT-WST -0.050                                                                /
PLAT-EST  1.0E3                                                                /
/
\end{verbatim}

\subsection{GCONTOL -- Define Group Constraint
Tolerance}\label{gcontol-define-group-constraint-tolerance}

\begin{longtable}[]{@{}
  >{\raggedright\arraybackslash}p{(\linewidth - 14\tabcolsep) * \real{0.1250}}
  >{\raggedright\arraybackslash}p{(\linewidth - 14\tabcolsep) * \real{0.1250}}
  >{\raggedright\arraybackslash}p{(\linewidth - 14\tabcolsep) * \real{0.1250}}
  >{\raggedright\arraybackslash}p{(\linewidth - 14\tabcolsep) * \real{0.1250}}
  >{\raggedright\arraybackslash}p{(\linewidth - 14\tabcolsep) * \real{0.1250}}
  >{\raggedright\arraybackslash}p{(\linewidth - 14\tabcolsep) * \real{0.1250}}
  >{\raggedright\arraybackslash}p{(\linewidth - 14\tabcolsep) * \real{0.1250}}
  >{\raggedright\arraybackslash}p{(\linewidth - 14\tabcolsep) * \real{0.1250}}@{}}
\toprule\noalign{}
\begin{minipage}[b]{\linewidth}\raggedright
RUNSPEC
\end{minipage} & \begin{minipage}[b]{\linewidth}\raggedright
GRID
\end{minipage} & \begin{minipage}[b]{\linewidth}\raggedright
EDIT
\end{minipage} & \begin{minipage}[b]{\linewidth}\raggedright
PROPS
\end{minipage} & \begin{minipage}[b]{\linewidth}\raggedright
REGIONS
\end{minipage} & \begin{minipage}[b]{\linewidth}\raggedright
SOLUTION
\end{minipage} & \begin{minipage}[b]{\linewidth}\raggedright
SUMMARY
\end{minipage} & \begin{minipage}[b]{\linewidth}\raggedright
SCHEDULE
\end{minipage} \\
\midrule\noalign{}
\endhead
\bottomrule\noalign{}
\endlastfoot
\end{longtable}

\subsubsection{Description}\label{description-931}

The GCONTOL keyword defines the tolerance and parameters used to control
the accuracy of group targets and constraints, including the field's
targets and constraints. The keyword sets the tolerance and number of
Newton iterations for each time step so that the wells under group
control can match the desired group targets and constraints.

This keyword is not supported by OPM Flow but would change the results
if supported so the simulation will be stopped.

\begin{longtable}[]{@{}
  >{\raggedright\arraybackslash}p{(\linewidth - 6\tabcolsep) * \real{0.2500}}
  >{\raggedright\arraybackslash}p{(\linewidth - 6\tabcolsep) * \real{0.2500}}
  >{\raggedright\arraybackslash}p{(\linewidth - 6\tabcolsep) * \real{0.2500}}
  >{\raggedright\arraybackslash}p{(\linewidth - 6\tabcolsep) * \real{0.2500}}@{}}
\toprule\noalign{}
\begin{minipage}[b]{\linewidth}\raggedright
No.
\end{minipage} & \begin{minipage}[b]{\linewidth}\raggedright
Name
\end{minipage} & \begin{minipage}[b]{\linewidth}\raggedright
Description
\end{minipage} & \begin{minipage}[b]{\linewidth}\raggedright
Default
\end{minipage} \\
\midrule\noalign{}
\endhead
\bottomrule\noalign{}
\endlastfoot
Field & Metric & Laboratory & \\
1 & GRPCNV & GRPCNV is a real positive value less than one that sets the
group tolerance criteria used to define if convergence has been
satisfied for group's under rate control. Here, GRPCNV is an acceptable
fraction of the group's rate target. Thus, the a numerical quantity 0.01
means that values must be with 0.01 (or 1.0\%) of the groups' production
target. For groups under priority control, as per GCONPRI keyword in the
SCHEDULE section, GRPCNV is ignored. Note that this criteria may not be
satisfied if the number of Newton iterations used in updating the well
targets, as set by the NUPCOL keyword in the RUNSPEC section, or the
NUPCOL parameter on this keyword, is exceeded. In this case, and only if
the well potentials allow, the well rates are re-calculated ignoring the
NUPCOL limit. Group tolerance criteria can also be set by the GRPCNV
parameter on the NETBALAN keyword in the SCHEDULE section. If both
values of GRPCNV have been entered, then the minimum of the two is used.
The default value means there is no convergence criteria. & 1.0 x
1020 \\
dimensionless & dimensionless & dimensionless & \\
2 & NUPCOL & A positive integer that defines the maximum number of
Newton iterations used to update well targets within a time step. Wells
under group control may suffer from some dependency with other wells in
the same group that are under group control. This may cause some
oscillation in the production and injection well rates within the group.
In order to avoid this, after the number Newton iterations within a time
step surpasses NUPCOL, the group well rates are frozen until the time
step has converged. Reducing the potential of well rate oscillations
within the time step may result in the group targets and limits not
being exactly being met in this case. Increasing the value of NUPCOL
will improve the accuracy of the group targets and limits at the expense
of computational efficiency. NUPCOL can also be set by the NUPCOL
keyword in the RUNSPEC section. A value entered on this keyword
overwrites any previously entered values of NUPCOL, including the value
set by the NUPCOL keyword. The default value of 1* invokes the last
previously entered value of either by the NUPCOL keyword or this
keyword. & 1* \\
dimensionless & dimensionless & dimensionless & \\
3 & GASCNV & In the commercial compositional simulator GASCNV is a real
positive value less than one that sets the tolerance criteria used to
define if convergence has been satisfied for well gas injection rates.
GASCNV is also used in conjunction with well availability gas fraction
quantities in the commercial compositional simulator The value is
ignored by OPM Flow. & 1.0 x 10-3 \\
dimensionless & dimensionless & dimensionless & \\
4 & GASMXITE & In the commercial compositional simulator GASMXITE is a
positive integer value that defines the maximum number of iterations for
the gas injection computation. The value is ignored by OPM Flow. & 5 \\
dimensionless & dimensionless & dimensionless & \\
Notes: & & & \\
\end{longtable}

\emph{Table 12.34: GCONTOL Keyword Description}

See also the NUPCOL keyword in the RUNSPEC section, and the NETBALAN and
WLIMTOL keywords in the SCHEDULE section that control the network
balancing convergence criteria and the tolerance parameters for wells,
respectively.

\subsubsection{Examples}\label{examples-75}

The example sets the group target convergence criteria to 1\% (0.01)
with a maximum of four Newton iterations (NUPCOL).

\begin{verbatim}
--
--       GROUP CONSTRAINT TOLERANCE
--
--       GROUP  NEWTON  INJEC  NEWTON
--       TOL    ITERS   TOL    ITERS
--       -----  -----   -----  ------
GCONTOL
         0.01    4       1*     1*                         /
\end{verbatim}

\subsection{GCUTBACK -- Define Group Cutback Limits and
Parameters}\label{gcutback-define-group-cutback-limits-and-parameters}

\begin{longtable}[]{@{}
  >{\raggedright\arraybackslash}p{(\linewidth - 14\tabcolsep) * \real{0.1250}}
  >{\raggedright\arraybackslash}p{(\linewidth - 14\tabcolsep) * \real{0.1250}}
  >{\raggedright\arraybackslash}p{(\linewidth - 14\tabcolsep) * \real{0.1250}}
  >{\raggedright\arraybackslash}p{(\linewidth - 14\tabcolsep) * \real{0.1250}}
  >{\raggedright\arraybackslash}p{(\linewidth - 14\tabcolsep) * \real{0.1250}}
  >{\raggedright\arraybackslash}p{(\linewidth - 14\tabcolsep) * \real{0.1250}}
  >{\raggedright\arraybackslash}p{(\linewidth - 14\tabcolsep) * \real{0.1250}}
  >{\raggedright\arraybackslash}p{(\linewidth - 14\tabcolsep) * \real{0.1250}}@{}}
\toprule\noalign{}
\begin{minipage}[b]{\linewidth}\raggedright
RUNSPEC
\end{minipage} & \begin{minipage}[b]{\linewidth}\raggedright
GRID
\end{minipage} & \begin{minipage}[b]{\linewidth}\raggedright
EDIT
\end{minipage} & \begin{minipage}[b]{\linewidth}\raggedright
PROPS
\end{minipage} & \begin{minipage}[b]{\linewidth}\raggedright
REGIONS
\end{minipage} & \begin{minipage}[b]{\linewidth}\raggedright
SOLUTION
\end{minipage} & \begin{minipage}[b]{\linewidth}\raggedright
SUMMARY
\end{minipage} & \begin{minipage}[b]{\linewidth}\raggedright
SCHEDULE
\end{minipage} \\
\midrule\noalign{}
\endhead
\bottomrule\noalign{}
\endlastfoot
\end{longtable}

\subsubsection{Description}\label{description-932}

This keyword, GCUTBACK, defines a production group's cutback limits and
parameters. See also the WCUTBACK keyword in the SCHEDULE section that
provides similar functionality for wells.

This keyword is not supported by OPM Flow but would change the results
if supported so the simulation will be stopped.

\subsection{GCUTBACT -- Define Group Tracer Cutback Limits and
Parameters}\label{gcutbact-define-group-tracer-cutback-limits-and-parameters}

\begin{longtable}[]{@{}
  >{\raggedright\arraybackslash}p{(\linewidth - 14\tabcolsep) * \real{0.1250}}
  >{\raggedright\arraybackslash}p{(\linewidth - 14\tabcolsep) * \real{0.1250}}
  >{\raggedright\arraybackslash}p{(\linewidth - 14\tabcolsep) * \real{0.1250}}
  >{\raggedright\arraybackslash}p{(\linewidth - 14\tabcolsep) * \real{0.1250}}
  >{\raggedright\arraybackslash}p{(\linewidth - 14\tabcolsep) * \real{0.1250}}
  >{\raggedright\arraybackslash}p{(\linewidth - 14\tabcolsep) * \real{0.1250}}
  >{\raggedright\arraybackslash}p{(\linewidth - 14\tabcolsep) * \real{0.1250}}
  >{\raggedright\arraybackslash}p{(\linewidth - 14\tabcolsep) * \real{0.1250}}@{}}
\toprule\noalign{}
\begin{minipage}[b]{\linewidth}\raggedright
RUNSPEC
\end{minipage} & \begin{minipage}[b]{\linewidth}\raggedright
GRID
\end{minipage} & \begin{minipage}[b]{\linewidth}\raggedright
EDIT
\end{minipage} & \begin{minipage}[b]{\linewidth}\raggedright
PROPS
\end{minipage} & \begin{minipage}[b]{\linewidth}\raggedright
REGIONS
\end{minipage} & \begin{minipage}[b]{\linewidth}\raggedright
SOLUTION
\end{minipage} & \begin{minipage}[b]{\linewidth}\raggedright
SUMMARY
\end{minipage} & \begin{minipage}[b]{\linewidth}\raggedright
SCHEDULE
\end{minipage} \\
\midrule\noalign{}
\endhead
\bottomrule\noalign{}
\endlastfoot
\end{longtable}

\subsubsection{Description}\label{description-933}

This keyword, GCUTBACT, defines a production group's cutback limits and
parameters based on the named produced tracer from the group. See also
the WCUTBACT keyword in the SCHEDULE section that provides similar
functionality for groups.

This keyword is not supported by OPM Flow but would change the results
if supported so the simulation will be stopped.

\subsection{GDCQ -- Define Group Multiple Daily Contract
Quantities}\label{gdcq-define-group-multiple-daily-contract-quantities}

\begin{longtable}[]{@{}
  >{\raggedright\arraybackslash}p{(\linewidth - 14\tabcolsep) * \real{0.1250}}
  >{\raggedright\arraybackslash}p{(\linewidth - 14\tabcolsep) * \real{0.1250}}
  >{\raggedright\arraybackslash}p{(\linewidth - 14\tabcolsep) * \real{0.1250}}
  >{\raggedright\arraybackslash}p{(\linewidth - 14\tabcolsep) * \real{0.1250}}
  >{\raggedright\arraybackslash}p{(\linewidth - 14\tabcolsep) * \real{0.1250}}
  >{\raggedright\arraybackslash}p{(\linewidth - 14\tabcolsep) * \real{0.1250}}
  >{\raggedright\arraybackslash}p{(\linewidth - 14\tabcolsep) * \real{0.1250}}
  >{\raggedright\arraybackslash}p{(\linewidth - 14\tabcolsep) * \real{0.1250}}@{}}
\toprule\noalign{}
\begin{minipage}[b]{\linewidth}\raggedright
RUNSPEC
\end{minipage} & \begin{minipage}[b]{\linewidth}\raggedright
GRID
\end{minipage} & \begin{minipage}[b]{\linewidth}\raggedright
EDIT
\end{minipage} & \begin{minipage}[b]{\linewidth}\raggedright
PROPS
\end{minipage} & \begin{minipage}[b]{\linewidth}\raggedright
REGIONS
\end{minipage} & \begin{minipage}[b]{\linewidth}\raggedright
SOLUTION
\end{minipage} & \begin{minipage}[b]{\linewidth}\raggedright
SUMMARY
\end{minipage} & \begin{minipage}[b]{\linewidth}\raggedright
SCHEDULE
\end{minipage} \\
\midrule\noalign{}
\endhead
\bottomrule\noalign{}
\endlastfoot
\end{longtable}

\subsubsection{Description}\label{description-934}

The GDCQ keyword defines the Daily Contract Quantities (``DCQ'') for
when multiple group contracts are required when the Gas Field Operations
model has been activated by the GASFIELD keyword in the RUNSPEC section,
or the GWSINGF has been invoked to define multiple group contracts in
the SCHEDULE section. The group contracts must first be defined by the
GSWINGF keyword, followed by the GCDQ keyword, and then the GASYEAR or
GASPERIO keywords. GCDQ may be repeated in the SCHEDULE section to reset
group DCQs.

See also the SWINGFAC keyword that set a single group DCQ at the field
level, as opposed to having multiple DCQ group contracts using the GDCQ
keyword.

This keyword is not supported by OPM Flow but would change the results
if supported so the simulation will be stopped.

\subsection{GDCQECON -- Group Economic Criteria for DCQ Production
Groups}\label{gdcqecon-group-economic-criteria-for-dcq-production-groups}

\begin{longtable}[]{@{}
  >{\raggedright\arraybackslash}p{(\linewidth - 14\tabcolsep) * \real{0.1250}}
  >{\raggedright\arraybackslash}p{(\linewidth - 14\tabcolsep) * \real{0.1250}}
  >{\raggedright\arraybackslash}p{(\linewidth - 14\tabcolsep) * \real{0.1250}}
  >{\raggedright\arraybackslash}p{(\linewidth - 14\tabcolsep) * \real{0.1250}}
  >{\raggedright\arraybackslash}p{(\linewidth - 14\tabcolsep) * \real{0.1250}}
  >{\raggedright\arraybackslash}p{(\linewidth - 14\tabcolsep) * \real{0.1250}}
  >{\raggedright\arraybackslash}p{(\linewidth - 14\tabcolsep) * \real{0.1250}}
  >{\raggedright\arraybackslash}p{(\linewidth - 14\tabcolsep) * \real{0.1250}}@{}}
\toprule\noalign{}
\begin{minipage}[b]{\linewidth}\raggedright
RUNSPEC
\end{minipage} & \begin{minipage}[b]{\linewidth}\raggedright
GRID
\end{minipage} & \begin{minipage}[b]{\linewidth}\raggedright
EDIT
\end{minipage} & \begin{minipage}[b]{\linewidth}\raggedright
PROPS
\end{minipage} & \begin{minipage}[b]{\linewidth}\raggedright
REGIONS
\end{minipage} & \begin{minipage}[b]{\linewidth}\raggedright
SOLUTION
\end{minipage} & \begin{minipage}[b]{\linewidth}\raggedright
SUMMARY
\end{minipage} & \begin{minipage}[b]{\linewidth}\raggedright
SCHEDULE
\end{minipage} \\
\midrule\noalign{}
\endhead
\bottomrule\noalign{}
\endlastfoot
\end{longtable}

\subsubsection{Description}\label{description-935}

The GDCQECON keyword defines economic criteria for DCQ production
groups, including the field level group FIELD, that have previously been
defined by the GCONPROD keywords in the SCHEDULE section. Note that
wells are allocated to a group when they are specified by the WELSPECS
keyword and wells can also have economic controls. Wells under group
control are therefore subject to the economic criteria set via the
GCONPROD and CECON keywords in the SCHEDULE section and the controls
specified by the WECON keyword.

This keyword is not supported by OPM Flow but would change the results
if supported so the simulation will be stopped.

\begin{longtable}[]{@{}
  >{\raggedright\arraybackslash}p{(\linewidth - 6\tabcolsep) * \real{0.2500}}
  >{\raggedright\arraybackslash}p{(\linewidth - 6\tabcolsep) * \real{0.2500}}
  >{\raggedright\arraybackslash}p{(\linewidth - 6\tabcolsep) * \real{0.2500}}
  >{\raggedright\arraybackslash}p{(\linewidth - 6\tabcolsep) * \real{0.2500}}@{}}
\toprule\noalign{}
\begin{minipage}[b]{\linewidth}\raggedright
No.
\end{minipage} & \begin{minipage}[b]{\linewidth}\raggedright
Name
\end{minipage} & \begin{minipage}[b]{\linewidth}\raggedright
Description
\end{minipage} & \begin{minipage}[b]{\linewidth}\raggedright
Default
\end{minipage} \\
\midrule\noalign{}
\endhead
\bottomrule\noalign{}
\endlastfoot
Field & Metric & Laboratory & \\
1 & GRPNAME & A character string of up to eight characters in length
that defines the group name for which the group target and constraints
are being defined. The group named FIELD is the top most group and
should be used to set targets and constraints for the field. Note that
the group hierarchy should be defined by the GRUPTREE keyword when there
is more than one level of groups, otherwise all the groups will sit
directly under the FIELD group in the group tree hierarchy. & None \\
2 & DCQ & A real positive value that defines the minimum economic DCQ
gas production rate, below which an economic action of shutting in or
stopping all the wells in the group, as requested by item (9) of the
WELSPECS keyword. Note if GRPNAME is equal to FIELD then the run will be
terminated. A value less than or equal to zero switches of this
criteria. & 0.0 \\
Mscf/d & sm3/day & scc/hour & \\
Notes: & & & \\
\end{longtable}

\emph{Table 12.35: GDCQECON Keyword Description}

See also the WELSPECS keyword to define a wells shut-in or stop options,
GCONPROD for group controls, and WECON for setting a well's economic
criteria. All the aforementioned keywords are described in the SCHEDULE
section.

\subsubsection{Example}\label{example-529}

The following example defines the minimum DCQ for the field to be 10
MMscf/d.

\begin{verbatim}
--
--       GROUP ECONOMIC CRITERIA FOR DCQ PRODUCTION GROUPS
--
-- GRUP  GAS
-- NAME  DCQ
GDCQECON
FIELD    10E3                                                                 /
/
\end{verbatim}

\subsection{GDRILPOT -- Define Group Potential Rates for Automatic
Drilling}\label{gdrilpot-define-group-potential-rates-for-automatic-drilling}

\begin{longtable}[]{@{}
  >{\raggedright\arraybackslash}p{(\linewidth - 14\tabcolsep) * \real{0.1250}}
  >{\raggedright\arraybackslash}p{(\linewidth - 14\tabcolsep) * \real{0.1250}}
  >{\raggedright\arraybackslash}p{(\linewidth - 14\tabcolsep) * \real{0.1250}}
  >{\raggedright\arraybackslash}p{(\linewidth - 14\tabcolsep) * \real{0.1250}}
  >{\raggedright\arraybackslash}p{(\linewidth - 14\tabcolsep) * \real{0.1250}}
  >{\raggedright\arraybackslash}p{(\linewidth - 14\tabcolsep) * \real{0.1250}}
  >{\raggedright\arraybackslash}p{(\linewidth - 14\tabcolsep) * \real{0.1250}}
  >{\raggedright\arraybackslash}p{(\linewidth - 14\tabcolsep) * \real{0.1250}}@{}}
\toprule\noalign{}
\begin{minipage}[b]{\linewidth}\raggedright
RUNSPEC
\end{minipage} & \begin{minipage}[b]{\linewidth}\raggedright
GRID
\end{minipage} & \begin{minipage}[b]{\linewidth}\raggedright
EDIT
\end{minipage} & \begin{minipage}[b]{\linewidth}\raggedright
PROPS
\end{minipage} & \begin{minipage}[b]{\linewidth}\raggedright
REGIONS
\end{minipage} & \begin{minipage}[b]{\linewidth}\raggedright
SOLUTION
\end{minipage} & \begin{minipage}[b]{\linewidth}\raggedright
SUMMARY
\end{minipage} & \begin{minipage}[b]{\linewidth}\raggedright
SCHEDULE
\end{minipage} \\
\midrule\noalign{}
\endhead
\bottomrule\noalign{}
\endlastfoot
\end{longtable}

\subsubsection{Description}\label{description-936}

This keyword, GDRILPOT, defines the minimum group potential rate that
will result in a well from the one of the automatic drilling queues, as
defined by either the QDRILL or WDRILPRI keywords in the SCHEDULE
section, to be drilled and placed on production. The advantage of using
a group's potential, as oppose to a minimum rate limit, is that setting
the potential greater than the group's minimum flow rate, will result in
well being drilled in time to support the desired production rate.

This keyword is not supported by OPM Flow but would change the results
if supported so the simulation will be stopped.

\subsection{GECON -- Group Economic Criteria for Production
Groups}\label{gecon-group-economic-criteria-for-production-groups}

\begin{longtable}[]{@{}
  >{\raggedright\arraybackslash}p{(\linewidth - 14\tabcolsep) * \real{0.1250}}
  >{\raggedright\arraybackslash}p{(\linewidth - 14\tabcolsep) * \real{0.1250}}
  >{\raggedright\arraybackslash}p{(\linewidth - 14\tabcolsep) * \real{0.1250}}
  >{\raggedright\arraybackslash}p{(\linewidth - 14\tabcolsep) * \real{0.1250}}
  >{\raggedright\arraybackslash}p{(\linewidth - 14\tabcolsep) * \real{0.1250}}
  >{\raggedright\arraybackslash}p{(\linewidth - 14\tabcolsep) * \real{0.1250}}
  >{\raggedright\arraybackslash}p{(\linewidth - 14\tabcolsep) * \real{0.1250}}
  >{\raggedright\arraybackslash}p{(\linewidth - 14\tabcolsep) * \real{0.1250}}@{}}
\toprule\noalign{}
\begin{minipage}[b]{\linewidth}\raggedright
RUNSPEC
\end{minipage} & \begin{minipage}[b]{\linewidth}\raggedright
GRID
\end{minipage} & \begin{minipage}[b]{\linewidth}\raggedright
EDIT
\end{minipage} & \begin{minipage}[b]{\linewidth}\raggedright
PROPS
\end{minipage} & \begin{minipage}[b]{\linewidth}\raggedright
REGIONS
\end{minipage} & \begin{minipage}[b]{\linewidth}\raggedright
SOLUTION
\end{minipage} & \begin{minipage}[b]{\linewidth}\raggedright
SUMMARY
\end{minipage} & \begin{minipage}[b]{\linewidth}\raggedright
SCHEDULE
\end{minipage} \\
\midrule\noalign{}
\endhead
\bottomrule\noalign{}
\endlastfoot
\end{longtable}

\subsubsection{Description}\label{description-937}

The GECON keyword defines economic criteria for production groups,
including the top level FIELD group, that have been created by having
wells assigned to them via the WELSPECS keyword in the SCHEDULE section.

Note that wells are allocated to a group when they are specified by the
WELSPECS keyword and wells can also have well level economic controls.
Wells under group control are therefore subject to the economic criteria
set via the GECON and WECON keywords and the controls specified by the
GCONPROD keyword in the SCHEDULE section.

\begin{longtable}[]{@{}
  >{\raggedright\arraybackslash}p{(\linewidth - 6\tabcolsep) * \real{0.2500}}
  >{\raggedright\arraybackslash}p{(\linewidth - 6\tabcolsep) * \real{0.2500}}
  >{\raggedright\arraybackslash}p{(\linewidth - 6\tabcolsep) * \real{0.2500}}
  >{\raggedright\arraybackslash}p{(\linewidth - 6\tabcolsep) * \real{0.2500}}@{}}
\toprule\noalign{}
\begin{minipage}[b]{\linewidth}\raggedright
No.
\end{minipage} & \begin{minipage}[b]{\linewidth}\raggedright
Name
\end{minipage} & \begin{minipage}[b]{\linewidth}\raggedright
Description
\end{minipage} & \begin{minipage}[b]{\linewidth}\raggedright
Default
\end{minipage} \\
\midrule\noalign{}
\endhead
\bottomrule\noalign{}
\endlastfoot
Field & Metric & Laboratory & \\
1 & GRPNAME & A character string of up to eight characters in length
that defines the group name for which the group target and constraints
are being defined. The group named FIELD is the top most group and
should be used to set targets and constraints for the field. Note that
the group hierarchy should be defined by the GRUPTREE keyword when there
is more than one level of groups, otherwise all the groups will sit
directly under the FIELD group in the group tree hierarchy. & None \\
2 & ORAT & A real positive value that defines the minimum economic
surface oil production rate, below which an economic action of shutting
in or stopping all the wells in the group, as requested by the AUTO
variable item (9) of the WELSPECS keyword. This value may be specified
using a User Defined Argument (UDA). A value less than or equal to zero
switches off this criteria. & 0.0 \\
stb/d & sm3/day & scc/hour & \\
3 & GRAT & A real positive value that defines the minimum economic
surface gas production rate, below which an economic action of shutting
in or stopping all the wells in the group, as requested by the AUTO
variable item (9) of the WELSPECS keyword. This value may be specified
using a User Defined Argument (UDA). A value less than or equal to zero
switches off this criteria, & 0.0 \\
Mscf/d & sm3/day & scc/hour & \\
4 & WCUT & A real positive value that defines the maximum economic
surface water cut, above which an economic action will take place. This
value may be specified using a User Defined Argument (UDA). Water cut is
defined as:\({f}_{w} = \frac{{q}_{w}}{{q}_{w} + {q}_{o}}\), and the
various actions that are available if the water cut limit is exceeded
are described in item (7). A value less than or equal to zero switches
off this criteria. & 0.0 \\
dimensionless & dimensionless & dimensionless & \\
5 & GOR & A real positive value that defines the maximum economic
surface gas-oil ratio, above which an economic action will take place,
as defined by item (7). This value may be specified using a User Defined
Argument (UDA). A value less than or equal to zero switches off this
criteria. & 0.0 \\
Mscf/stb & sm3/sm3 & scc/scc & \\
6 & WGR & A real positive value that defines the maximum economic
surface water-gas ratio, above which an economic action will take place,
as defined by item (7). This value may be specified using a User Defined
Argument (UDA). A value less than or equal to zero switches off this
criteria. & 0.0 \\
stb/Mscf & sm3/sm3 & scc/scc & \\
7 & ACTION & A defined character string that defines the action to be
taken if the economic WCUT, GOR, or WGR limits are violated. ACTION
should be set to one of the following character strings: The corrective
action takes places at the end of the time step in which the constraint
is violated. Only the NONE option is currently supported by the
simulator. & None \\
8 & END & A defined character string that defines if the simulation
should terminate if all the producing wells in the group, including the
FIELD group, are shut or stopped. END should be set to one of the
following character strings: Only the NO option is currently supported
by the simulator. & NO \\
9 & MXWELS & A positive integer defining the maximum number of producing
and injecting wells for this this group and any subordinate groups. The
default value of zero implies that there is no limit to the number of
wells. This is not currently supported by the simulator and must be
defaulted. & 0 \\
Notes: & & & \\
\end{longtable}

\emph{Table 12.3.98.1: GECON Keyword Description}

See also the WELSPECS keyword to define a well's shut-in or stop
options, GCONPROD for group controls, and WECON for setting a well's
economic criteria. All the aforementioned keywords are described in the
SCHEDULE section.

\subsubsection{Example}\label{example-530}

The following example defines the economic criteria for the field with a
minimum oil rate of 2,000 m3/day and a maximum water cut of 95\%.

\begin{verbatim}
--
--       GROUP ECONOMIC CRITERIA FOR PRODUCTION GROUPS
--
-- GRUP  OIL   GAS    WCT    GOR    WGR    WORK   END   MAX
-- NAME  MIN   MIN    MAX    MAX    MAX    OVER   RUN   WELLS
GECON
FIELD    2E3   1*     0.95     1*    1*     CON   'YES'  1*                     /
/
\end{verbatim}

If the economic limits are violated then the run will stop at the next
report time step.

\subsection{GECONT -- Group Tracer Economic Criteria for Production
Groups}\label{gecont-group-tracer-economic-criteria-for-production-groups}

\begin{longtable}[]{@{}
  >{\raggedright\arraybackslash}p{(\linewidth - 14\tabcolsep) * \real{0.1250}}
  >{\raggedright\arraybackslash}p{(\linewidth - 14\tabcolsep) * \real{0.1250}}
  >{\raggedright\arraybackslash}p{(\linewidth - 14\tabcolsep) * \real{0.1250}}
  >{\raggedright\arraybackslash}p{(\linewidth - 14\tabcolsep) * \real{0.1250}}
  >{\raggedright\arraybackslash}p{(\linewidth - 14\tabcolsep) * \real{0.1250}}
  >{\raggedright\arraybackslash}p{(\linewidth - 14\tabcolsep) * \real{0.1250}}
  >{\raggedright\arraybackslash}p{(\linewidth - 14\tabcolsep) * \real{0.1250}}
  >{\raggedright\arraybackslash}p{(\linewidth - 14\tabcolsep) * \real{0.1250}}@{}}
\toprule\noalign{}
\begin{minipage}[b]{\linewidth}\raggedright
RUNSPEC
\end{minipage} & \begin{minipage}[b]{\linewidth}\raggedright
GRID
\end{minipage} & \begin{minipage}[b]{\linewidth}\raggedright
EDIT
\end{minipage} & \begin{minipage}[b]{\linewidth}\raggedright
PROPS
\end{minipage} & \begin{minipage}[b]{\linewidth}\raggedright
REGIONS
\end{minipage} & \begin{minipage}[b]{\linewidth}\raggedright
SOLUTION
\end{minipage} & \begin{minipage}[b]{\linewidth}\raggedright
SUMMARY
\end{minipage} & \begin{minipage}[b]{\linewidth}\raggedright
SCHEDULE
\end{minipage} \\
\midrule\noalign{}
\endhead
\bottomrule\noalign{}
\endlastfoot
\end{longtable}

\subsubsection{Description}\label{description-938}

The GECONT keyword defines tracer economic criteria for production
groups, including the field level group FIELD, that have previously been
defined by the GCONPROD keywords in the SCHEDULE section, for tracers
define by the TRACER keyword in the PROPS section.

Note that wells are allocated to a group when they are specified by the
WELSPECS keyword and wells can also have economic controls. Wells under
group control are therefore subject to the economic criteria set via the
GCONPROD and CECON keywords in the SCHEDULE section and the controls
specified by the WECON keyword.

This keyword is not supported by OPM Flow but would change the results
if supported so the simulation will be stopped.

\begin{longtable}[]{@{}
  >{\raggedright\arraybackslash}p{(\linewidth - 6\tabcolsep) * \real{0.2500}}
  >{\raggedright\arraybackslash}p{(\linewidth - 6\tabcolsep) * \real{0.2500}}
  >{\raggedright\arraybackslash}p{(\linewidth - 6\tabcolsep) * \real{0.2500}}
  >{\raggedright\arraybackslash}p{(\linewidth - 6\tabcolsep) * \real{0.2500}}@{}}
\toprule\noalign{}
\begin{minipage}[b]{\linewidth}\raggedright
No.
\end{minipage} & \begin{minipage}[b]{\linewidth}\raggedright
Name
\end{minipage} & \begin{minipage}[b]{\linewidth}\raggedright
Description
\end{minipage} & \begin{minipage}[b]{\linewidth}\raggedright
Default
\end{minipage} \\
\midrule\noalign{}
\endhead
\bottomrule\noalign{}
\endlastfoot
Field & Metric & Laboratory & \\
1-1 & GRPNAME & A character string of up to eight characters in length
that defines the group name for which the group target and constraints
are being defined. The group named FIELD is the top most group and
should be used to set targets and constraints for the field. Note that
the group hierarchy should be defined by the GRUPTREE keyword when there
is more than one level of groups, otherwise all the groups will sit
directly under the FIELD group in the group tree hierarchy. & None \\
1-2 & ACTION & A defined character string that defines the action to be
taken if the economic WCUT, GOR, or WGR limits are violated. ACTION
should be set to one of the following character strings: The corrective
action takes places at the end of the time step in which the constraint
is violated. & None \\
1-3 & END & A defined character string that defines if the simulation
should terminate if all the producing wells in the group, including the
FIELD group, are shut or stopped. END should be set to one of the
following character strings: & NO \\
1-4 & MXWELS & A positive integer defining the maximum number of
producing and injecting wells for this group and any subordinate groups.
The default value of zero implies that there is no limit to the number
of wells. & 0 \\
1-5 & / & Record one terminated by a ``/'' & Not Applicable \\
2-1 & NAME & A three letter character string defining the tracer's name.
Note it is best to void names beginning with the letters F, S, and T as
these names may create naming issues in post-processing software. &
None \\
2-2 & MXTOTAL & A real positive value that defines the maximum total
(free plus solution) tracer rate. Tracer units are per those defined by
the carrying fluid, oil, gas, water, etc. & None \\
2-3 & MXFREET & A real positive value that defines the maximum total
(free plus solution) tracer concentration. Tracer units are per those
defined by the carrying fluid, oil, gas, water, etc. & None \\
2-4 & MXFREEQ & A real positive value that defines the maximum free
tracer rate. Tracer units are per those defined by the carrying fluid,
oil, gas, water, etc. & None \\
2-5 & MXCONC & A real positive value that defines the maximum free
tracer concentration. Tracer units are per those defined by the carrying
fluid, oil, gas, water, etc. & None \\
2-6 & MXSOLNQ & A real positive value that defines the maximum solution
rate. Tracer units are per those defined by the carrying fluid, oil,
gas, water, etc. & None \\
2-7 & MXSOLNC & A real positive value that defines the maximum solution
concentration. Tracer units are per those defined by the carrying fluid,
oil, gas, water, etc. & None \\
2-8 & / & Record two terminated by a ``/'' & Not Applicable \\
3-1 & / & Group terminated by a ``/'' & Not Applicable \\
Notes: & & & \\
\end{longtable}

\emph{Table 12.37: GECONT Keyword Description}

See also the WELSPECS keyword to define a wells shut-in or stop options,
GCONPROD for group controls, and WECON for setting a well's economic
criteria. All the aforementioned keywords are described in the SCHEDULE
section.

\subsubsection{Example}\label{example-531}

The following example defines the tracer economic criteria for the field
and two groups, FLTBLK1 and FLTBLK,.

\begin{verbatim}
--
--       GROUP TRACER ECONOMIC CRITERIA FOR PRODUCTION GROUPS
--
-- GRUP  WORK    END    MAX
-- NAME  OVER    RUN    WELLS
GECONT
FIELD    +CON   'YES'   1*                                     / START OF GROUP
--
--       TRACER  TRACER  TRACER  TRACER  TRACER  TRACER  TRACER
--       NAME    TOTAL   TOTAL   FREE    FREE    SOLN    SOLN
--               RATE    CONCEN  RATE    CONCEN  RATE    CONCEN
         PLY     1000.0                                        /
         BRI     1000.0                                        /
         TR1     1*      0.7500                                /
                                                               /
FLTBLK1  +CON   'YES'   1*                                     / START OF GROUP
--
--       TRACER  TRACER  TRACER  TRACER  TRACER  TRACER  TRACER
--       NAME    TOTAL   TOTAL   FREE    FREE    SOLN    SOLN
--               RATE    CONCEN  RATE    CONCEN  RATE    CONCEN
         PLY     800.0                                         /
         BRI     800.0                                         /
                                                                      /
FLTBLK2  +CON   'YES'   1*                                     / START OF GROUP
--
--       TRACER  TRACER  TRACER  TRACER  TRACER  TRACER  TRACER
--       NAME    TOTAL   TOTAL   FREE    FREE    SOLN    SOLN
--               RATE    CONCEN  RATE    CONCEN  RATE    CONCEN
         PLY     800.0                                         /
         BRI     800.0                                         / END OF GROUP
                                                               / END OF KEYWORD
\end{verbatim}

If the economic limits are violated then the worst offending connection
and all below it in the worst offending well will be closed, If
connections have been grouped as completions then the worst offending
completion and all below it in the worst offending well will be closed.

\subsection{GEFAC -- Define Group
Efficiency}\label{gefac-define-group-efficiency}

\begin{longtable}[]{@{}
  >{\raggedright\arraybackslash}p{(\linewidth - 14\tabcolsep) * \real{0.1250}}
  >{\raggedright\arraybackslash}p{(\linewidth - 14\tabcolsep) * \real{0.1250}}
  >{\raggedright\arraybackslash}p{(\linewidth - 14\tabcolsep) * \real{0.1250}}
  >{\raggedright\arraybackslash}p{(\linewidth - 14\tabcolsep) * \real{0.1250}}
  >{\raggedright\arraybackslash}p{(\linewidth - 14\tabcolsep) * \real{0.1250}}
  >{\raggedright\arraybackslash}p{(\linewidth - 14\tabcolsep) * \real{0.1250}}
  >{\raggedright\arraybackslash}p{(\linewidth - 14\tabcolsep) * \real{0.1250}}
  >{\raggedright\arraybackslash}p{(\linewidth - 14\tabcolsep) * \real{0.1250}}@{}}
\toprule\noalign{}
\begin{minipage}[b]{\linewidth}\raggedright
RUNSPEC
\end{minipage} & \begin{minipage}[b]{\linewidth}\raggedright
GRID
\end{minipage} & \begin{minipage}[b]{\linewidth}\raggedright
EDIT
\end{minipage} & \begin{minipage}[b]{\linewidth}\raggedright
PROPS
\end{minipage} & \begin{minipage}[b]{\linewidth}\raggedright
REGIONS
\end{minipage} & \begin{minipage}[b]{\linewidth}\raggedright
SOLUTION
\end{minipage} & \begin{minipage}[b]{\linewidth}\raggedright
SUMMARY
\end{minipage} & \begin{minipage}[b]{\linewidth}\raggedright
SCHEDULE
\end{minipage} \\
\midrule\noalign{}
\endhead
\bottomrule\noalign{}
\endlastfoot
\end{longtable}

\subsubsection{Description}\label{description-939}

Defines a group's efficiency or up-time factor as opposed to setting the
efficiency factors for individual wells.

Note that wells are allocated to a group when they are specified by the
WELSPECS keyword and wells can also have efficiency factors.

\begin{longtable}[]{@{}
  >{\raggedright\arraybackslash}p{(\linewidth - 6\tabcolsep) * \real{0.2500}}
  >{\raggedright\arraybackslash}p{(\linewidth - 6\tabcolsep) * \real{0.2500}}
  >{\raggedright\arraybackslash}p{(\linewidth - 6\tabcolsep) * \real{0.2500}}
  >{\raggedright\arraybackslash}p{(\linewidth - 6\tabcolsep) * \real{0.2500}}@{}}
\toprule\noalign{}
\begin{minipage}[b]{\linewidth}\raggedright
No.
\end{minipage} & \begin{minipage}[b]{\linewidth}\raggedright
Name
\end{minipage} & \begin{minipage}[b]{\linewidth}\raggedright
Description
\end{minipage} & \begin{minipage}[b]{\linewidth}\raggedright
Default
\end{minipage} \\
\midrule\noalign{}
\endhead
\bottomrule\noalign{}
\endlastfoot
Field & Metric & Laboratory & \\
1 & GRPNAME & A character string of up to eight characters in length
that defines the group name for which the group efficiency factor is
being defined. The group named FIELD is the top most group and cannot
have an efficiency factor set. Note that the group hierarchy should be
defined by the GRUPTREE keyword when there is more than one level of
groups, otherwise all the groups will sit directly under the FIELD group
in the group tree hierarchy. & None \\
2 & FACTOR & A real positive value less than or equal to one that
defines the efficiency factor for the group. If a group's down time is
5\% then FACTOR should be set to 0.95 (1.0 -- 0.05). & 1.0 \\
dimensionless & dimensionless & dimensionless & \\
3 & GRPNETWK & A defined character string that determines if the GRPNAME
efficiency factor should be transferred to the equivalent Extended
Network Model node, and should be set to either: This option is only
applicable for the Extended Network Model. In the Standard Network Model
groups' reported flow rates are always used in the calculation of
pressure drops. & YES \\
Notes: & & & \\
\end{longtable}

\emph{Table 12.38: GEFAC Keyword Description}

See also the WEFAC and NEFAC keywords in the SCHEDULE section to define
the efficiency factor for a well or a node.

\subsubsection{Example}\label{example-532}

\begin{verbatim}
--
--       GROUP EFFICIENCY FACTORS
--
-- GRUP  EFF      NETWK
-- NAME  FACT     OPTN
--      ------   ------
GEFAC
PLATFORM 0.950                                                                 /
SUBSEA1  0.860                                                                 /
/
\end{verbatim}

In the above example the group PLATFORM has its up-time factor set to
95\% and the subsea group SUBSEA1 has an up-time factor of 86\%.

\subsection{GLIFTLIM -- Group Artificial Lift
Constraints}\label{gliftlim-group-artificial-lift-constraints}

\begin{longtable}[]{@{}
  >{\raggedright\arraybackslash}p{(\linewidth - 14\tabcolsep) * \real{0.1250}}
  >{\raggedright\arraybackslash}p{(\linewidth - 14\tabcolsep) * \real{0.1250}}
  >{\raggedright\arraybackslash}p{(\linewidth - 14\tabcolsep) * \real{0.1250}}
  >{\raggedright\arraybackslash}p{(\linewidth - 14\tabcolsep) * \real{0.1250}}
  >{\raggedright\arraybackslash}p{(\linewidth - 14\tabcolsep) * \real{0.1250}}
  >{\raggedright\arraybackslash}p{(\linewidth - 14\tabcolsep) * \real{0.1250}}
  >{\raggedright\arraybackslash}p{(\linewidth - 14\tabcolsep) * \real{0.1250}}
  >{\raggedright\arraybackslash}p{(\linewidth - 14\tabcolsep) * \real{0.1250}}@{}}
\toprule\noalign{}
\begin{minipage}[b]{\linewidth}\raggedright
RUNSPEC
\end{minipage} & \begin{minipage}[b]{\linewidth}\raggedright
GRID
\end{minipage} & \begin{minipage}[b]{\linewidth}\raggedright
EDIT
\end{minipage} & \begin{minipage}[b]{\linewidth}\raggedright
PROPS
\end{minipage} & \begin{minipage}[b]{\linewidth}\raggedright
REGIONS
\end{minipage} & \begin{minipage}[b]{\linewidth}\raggedright
SOLUTION
\end{minipage} & \begin{minipage}[b]{\linewidth}\raggedright
SUMMARY
\end{minipage} & \begin{minipage}[b]{\linewidth}\raggedright
SCHEDULE
\end{minipage} \\
\midrule\noalign{}
\endhead
\bottomrule\noalign{}
\endlastfoot
\end{longtable}

\subsubsection{Description}\label{description-940}

The GLIFTLIM keyword defines the maximum number of wells on artificial
lift and the maximum amount of the artificial lift that is available for
a group, including the top most group in the group hierarchy known as
the FIELD group. Wells are allocated to groups when the wells are
specified by the WELSPECS keyword in the SCHEDULE section. Wells defined
to be under group control will have their production rates controlled by
the group to which they belong, in addition to any well constraints
defined for the wells.

This keyword is not supported by OPM Flow but would change the results
if supported so the simulation will be stopped.

\begin{longtable}[]{@{}
  >{\raggedright\arraybackslash}p{(\linewidth - 6\tabcolsep) * \real{0.2500}}
  >{\raggedright\arraybackslash}p{(\linewidth - 6\tabcolsep) * \real{0.2500}}
  >{\raggedright\arraybackslash}p{(\linewidth - 6\tabcolsep) * \real{0.2500}}
  >{\raggedright\arraybackslash}p{(\linewidth - 6\tabcolsep) * \real{0.2500}}@{}}
\toprule\noalign{}
\begin{minipage}[b]{\linewidth}\raggedright
No.
\end{minipage} & \begin{minipage}[b]{\linewidth}\raggedright
Name
\end{minipage} & \begin{minipage}[b]{\linewidth}\raggedright
Description
\end{minipage} & \begin{minipage}[b]{\linewidth}\raggedright
Default
\end{minipage} \\
\midrule\noalign{}
\endhead
\bottomrule\noalign{}
\endlastfoot
Field & Metric & Laboratory & \\
1 & GRPNAME & A character string of up to eight characters in length
that defines the group name for which the group's artificial lift
constraints are being defined. The group named FIELD is the top most
group and should be used to set targets and constraints for the field.
Note that the group hierarchy should be defined by the GRUPTREE keyword
when there is more than one level of groups, otherwise all the groups
will sit directly under the FIELD group in the group tree hierarchy. &
None \\
2 & MXLIFT & A real positive value that defines the total amount of
artificial lift available for this group and any subordinate groups. The
units for MXLIFT are the same as that defined by the ALQ parameter on
the VFPPROD keyword in the SCHEDULE section. For example, if ALQ has
been set to GRAT on the VFPPROD keyword, then MXLIFT would be the
maximum amount of gas lift gas available for this group and any
subordinate groups, and the units would Mscf, assuming FIELD units had
been activated in the RUNSPEC section. The default value of zero implies
that there is no limit applied to the group and its subordinate groups.
& 0 \\
See VFPPROD (ALQ) & See VFPPROD (ALQ) & See VFPPROD (ALQ) & \\
3 & MXWELS & A positive integer defining the maximum number of producing
wells on artificial lift for this group and any subordinate groups. The
default value of zero implies that there is no limit to the number of
wells. & 0 \\
dimensionless & dimensionless & dimensionless & \\
Notes: & & & \\
\end{longtable}

\emph{Table 12.39: GLIFTLIM Keyword Description}

See also the GRUPTREE keyword to define the hierarchy of the groups
below the FIELD level, the GCONPROD and GCONINJE keywords to define a
group's production and injection rate targets and constraints, the
WCONPROD keyword to define a production well's targets and constraints,
and the WCONINJE keyword to define an injection well's targets and
constraints. All the aforementioned keywords are described in the
SCHEDULE section.

\subsubsection{Example}\label{example-533}

The following example defines the artificial lift constraints for the
field, assuming all the wells are on gas lift.

\begin{verbatim}
---
--       GROUP ARTIFICIAL LIFT CONSTRAINTS
--
-- GRUP  MAX   MAX
-- NAME  ALQ   WELLS
GLIFTLIM
FIELD    20E3  20                                                             /
/
\end{verbatim}

Here the maximum amount of gas lift gas for the field is set to 20.0
MMscf/f and a maximum of 20 wells can utilize gas lift at a time.

\subsection{GLIFTOPT -- Define Group Gas Optimization
Limits}\label{gliftopt-define-group-gas-optimization-limits}

\begin{longtable}[]{@{}
  >{\raggedright\arraybackslash}p{(\linewidth - 14\tabcolsep) * \real{0.1250}}
  >{\raggedright\arraybackslash}p{(\linewidth - 14\tabcolsep) * \real{0.1250}}
  >{\raggedright\arraybackslash}p{(\linewidth - 14\tabcolsep) * \real{0.1250}}
  >{\raggedright\arraybackslash}p{(\linewidth - 14\tabcolsep) * \real{0.1250}}
  >{\raggedright\arraybackslash}p{(\linewidth - 14\tabcolsep) * \real{0.1250}}
  >{\raggedright\arraybackslash}p{(\linewidth - 14\tabcolsep) * \real{0.1250}}
  >{\raggedright\arraybackslash}p{(\linewidth - 14\tabcolsep) * \real{0.1250}}
  >{\raggedright\arraybackslash}p{(\linewidth - 14\tabcolsep) * \real{0.1250}}@{}}
\toprule\noalign{}
\begin{minipage}[b]{\linewidth}\raggedright
RUNSPEC
\end{minipage} & \begin{minipage}[b]{\linewidth}\raggedright
GRID
\end{minipage} & \begin{minipage}[b]{\linewidth}\raggedright
EDIT
\end{minipage} & \begin{minipage}[b]{\linewidth}\raggedright
PROPS
\end{minipage} & \begin{minipage}[b]{\linewidth}\raggedright
REGIONS
\end{minipage} & \begin{minipage}[b]{\linewidth}\raggedright
SOLUTION
\end{minipage} & \begin{minipage}[b]{\linewidth}\raggedright
SUMMARY
\end{minipage} & \begin{minipage}[b]{\linewidth}\raggedright
SCHEDULE
\end{minipage} \\
\midrule\noalign{}
\endhead
\bottomrule\noalign{}
\endlastfoot
\end{longtable}

\subsubsection{Description}\label{description-941}

The GLIFTOPT keyword defines the maximum amount of gas lift gas
available and the maximum amount of gas the group can produce, including
the top most group in the group hierarchy known as the FIELD group, for
when gas lift optimization has been activated via the LIFTOPT keyword in
the SCHEDULE section. Note that the LIFTOPT keyword should precede the
GLIFTOPT keyword in the SCHEDULE section in order to activate the gas
lift optimization facility.

Wells are allocated to groups when the wells are specified by the
WELSPECS keyword in the SCHEDULE section. Wells defined to be under
group control will have their production rates controlled by the group
to which they belong, in addition to any well constraints defined for
the wells, including any well gas lift optimization parameters on the
WLIFTOPT keyword.

\begin{longtable}[]{@{}
  >{\raggedright\arraybackslash}p{(\linewidth - 6\tabcolsep) * \real{0.2500}}
  >{\raggedright\arraybackslash}p{(\linewidth - 6\tabcolsep) * \real{0.2500}}
  >{\raggedright\arraybackslash}p{(\linewidth - 6\tabcolsep) * \real{0.2500}}
  >{\raggedright\arraybackslash}p{(\linewidth - 6\tabcolsep) * \real{0.2500}}@{}}
\toprule\noalign{}
\begin{minipage}[b]{\linewidth}\raggedright
No.
\end{minipage} & \begin{minipage}[b]{\linewidth}\raggedright
Name
\end{minipage} & \begin{minipage}[b]{\linewidth}\raggedright
Description
\end{minipage} & \begin{minipage}[b]{\linewidth}\raggedright
Default
\end{minipage} \\
\midrule\noalign{}
\endhead
\bottomrule\noalign{}
\endlastfoot
Field & Metric & Laboratory & \\
1 & GRPNAME & A character string of up to eight characters in length
that defines the group name for which the group's gas lift optimization
parameters are being defined. The group named FIELD is the top most
group and should be used to set targets and constraints for the field.
Note that the group hierarchy should be defined by the GRUPTREE keyword
when there is more than one level of groups, otherwise all the groups
will sit directly under the FIELD group in the group tree hierarchy. &
None \\
2 & MXLIFT & A real value that defines the total amount of gas lift gas
available for this group and any subordinate groups, multiplied by their
respective efficiency factors. The units for MXLIFT are the same as that
defined by the ALQ parameter on the VFPPROD keyword in the SCHEDULE
section. In this case ALQ should be GRAT on the VFPPROD keyword, as
MXLIFT applies to the maximum amount of gas lift gas available. The
default value of zero, or a negative value implies that there is no
limit applied to the group and its subordinate groups. & 0 \\
Mscf/d & sm3/day & scc/hour & \\
3 & MXGAS & A real value that defines the total amount of gas the group
can process. This is the sum of the gas lift gas plus the produced gas
for this group and any subordinate groups, multiplied by their
respective efficiency factors. The default value of zero, or a negative
value implies that there is no limit applied to the group and its
subordinate groups & 0 \\
Mscf/d & sm3/day & scc/hour & \\
Notes: & & & \\
\end{longtable}

\emph{Table 12.40: GLIFTOPT Keyword Description}

See also the LIFTOPT keyword to activate gas lift optimization, the
WLIFTOPT keyword to define the wells under gas lift optimization
control, the GRUPTREE keyword to define the hierarchy of the groups
below the FIELD level, the GCONPROD and GCONINJE keywords to define a
group's production and injection rate targets and constraints. All the
aforementioned keywords are described in the SCHEDULE section.

\subsubsection{Example}\label{example-534}

The following example first switches on gas lift optimization via the
LIFTOPT keyword and then defines the artificial lift constraints for the
field, assuming all the well are on gas lift, using the GLIFTOPT
keyword.

\begin{verbatim}
--       ACTIVATE GAS LIFT OPTIMIZATION AND PARAMETERS
--
-- INCR   INCR   TSTEP    NEWTON
-- GAS    OIL    INTVAL   OPTN
LIFTOPT
12.5E3    5E-3   0.0      YES                                                  /
/
--
--       GROUP GAS LIFT OPTIMIZATION CONSTRAINTS
--
-- GRUP  MAX       MAX
-- NAME  GAS ALQ   TOTAL GAS
GLIFTOPT
FIELD    200E3     1*                                                          /
/
\end{verbatim}

Here the LIFTOPT keyword defines the maximum incremental gas lift gas
quantity to be 12.5 x 103 m3, the minimum incremental oil gain per m3 of
gas lift gas is set to 5.0 x 10-3 m3, the time step interval is set to
zero to perform the gas optimization every time step, and finally the
gas lift optimization will be performed NUPCOL Newton iterations for the
time step.

The GLIFTOPT sets the maximum amount of gas lift gas for the field to
200,000 m3 and there is no maximum limit for the total maximum amount of
gas that the group can process.

\subsection{GNETDP -- Group Network Pressure and Rate
Controls}\label{gnetdp-group-network-pressure-and-rate-controls}

\begin{longtable}[]{@{}
  >{\raggedright\arraybackslash}p{(\linewidth - 14\tabcolsep) * \real{0.1250}}
  >{\raggedright\arraybackslash}p{(\linewidth - 14\tabcolsep) * \real{0.1250}}
  >{\raggedright\arraybackslash}p{(\linewidth - 14\tabcolsep) * \real{0.1250}}
  >{\raggedright\arraybackslash}p{(\linewidth - 14\tabcolsep) * \real{0.1250}}
  >{\raggedright\arraybackslash}p{(\linewidth - 14\tabcolsep) * \real{0.1250}}
  >{\raggedright\arraybackslash}p{(\linewidth - 14\tabcolsep) * \real{0.1250}}
  >{\raggedright\arraybackslash}p{(\linewidth - 14\tabcolsep) * \real{0.1250}}
  >{\raggedright\arraybackslash}p{(\linewidth - 14\tabcolsep) * \real{0.1250}}@{}}
\toprule\noalign{}
\begin{minipage}[b]{\linewidth}\raggedright
RUNSPEC
\end{minipage} & \begin{minipage}[b]{\linewidth}\raggedright
GRID
\end{minipage} & \begin{minipage}[b]{\linewidth}\raggedright
EDIT
\end{minipage} & \begin{minipage}[b]{\linewidth}\raggedright
PROPS
\end{minipage} & \begin{minipage}[b]{\linewidth}\raggedright
REGIONS
\end{minipage} & \begin{minipage}[b]{\linewidth}\raggedright
SOLUTION
\end{minipage} & \begin{minipage}[b]{\linewidth}\raggedright
SUMMARY
\end{minipage} & \begin{minipage}[b]{\linewidth}\raggedright
SCHEDULE
\end{minipage} \\
\midrule\noalign{}
\endhead
\bottomrule\noalign{}
\endlastfoot
\end{longtable}

\subsubsection{Description}\label{description-942}

The GNETDP keyword sets a group's minimum and maximum network pressure
and rate controls for when the either the Standard Network or the
Extended Network options have been activated, and the group is part of a
network. The keywords allows for the pressure of the group to vary in
order to satisfy the rate conditions declared by this keyword. The
Standard Network option is invoked if the GRUPTREE, GRUPNET, GNETINJE,
GNETPUMP, etc. series of keywords have been used in the SCHEDULE
section. Whereas, the Extended Network option is activated by the
NETWORK keyword in the RUNSPEC section. Several keywords, including,
GNETDP, can be used by both network options.

This keyword is not supported by OPM Flow but would change the results
if supported so the simulation will be stopped.

\subsection{GNETINJE -- Define Group Injection Network
Configuration}\label{gnetinje-define-group-injection-network-configuration}

\begin{longtable}[]{@{}
  >{\raggedright\arraybackslash}p{(\linewidth - 14\tabcolsep) * \real{0.1250}}
  >{\raggedright\arraybackslash}p{(\linewidth - 14\tabcolsep) * \real{0.1250}}
  >{\raggedright\arraybackslash}p{(\linewidth - 14\tabcolsep) * \real{0.1250}}
  >{\raggedright\arraybackslash}p{(\linewidth - 14\tabcolsep) * \real{0.1250}}
  >{\raggedright\arraybackslash}p{(\linewidth - 14\tabcolsep) * \real{0.1250}}
  >{\raggedright\arraybackslash}p{(\linewidth - 14\tabcolsep) * \real{0.1250}}
  >{\raggedright\arraybackslash}p{(\linewidth - 14\tabcolsep) * \real{0.1250}}
  >{\raggedright\arraybackslash}p{(\linewidth - 14\tabcolsep) * \real{0.1250}}@{}}
\toprule\noalign{}
\begin{minipage}[b]{\linewidth}\raggedright
RUNSPEC
\end{minipage} & \begin{minipage}[b]{\linewidth}\raggedright
GRID
\end{minipage} & \begin{minipage}[b]{\linewidth}\raggedright
EDIT
\end{minipage} & \begin{minipage}[b]{\linewidth}\raggedright
PROPS
\end{minipage} & \begin{minipage}[b]{\linewidth}\raggedright
REGIONS
\end{minipage} & \begin{minipage}[b]{\linewidth}\raggedright
SOLUTION
\end{minipage} & \begin{minipage}[b]{\linewidth}\raggedright
SUMMARY
\end{minipage} & \begin{minipage}[b]{\linewidth}\raggedright
SCHEDULE
\end{minipage} \\
\midrule\noalign{}
\endhead
\bottomrule\noalign{}
\endlastfoot
\end{longtable}

\subsubsection{Description}\label{description-943}

The GNETINJE keyword defines the configuration of a group injection
network for when the either the Standard Network or the Extended Network
options have been activated. The Standard Network option is invoked if
the GRUPTREE, GRUPNET, GNETINJE, GNETPUMP, etc. series of keywords have
been used in the SCHEDULE section. Whereas, the Extended Network option
is activated by the NETWORK keyword in the RUNSPEC section. Several
keywords, including, GNETINJE, can be used by both network options.

This keyword is not supported by OPM Flow but would change the results
if supported so the simulation will be stopped.

\subsection{GNETPUMP -- Standard Network Automatic Compressor and
Pumps}\label{gnetpump-standard-network-automatic-compressor-and-pumps}

\begin{longtable}[]{@{}
  >{\raggedright\arraybackslash}p{(\linewidth - 14\tabcolsep) * \real{0.1250}}
  >{\raggedright\arraybackslash}p{(\linewidth - 14\tabcolsep) * \real{0.1250}}
  >{\raggedright\arraybackslash}p{(\linewidth - 14\tabcolsep) * \real{0.1250}}
  >{\raggedright\arraybackslash}p{(\linewidth - 14\tabcolsep) * \real{0.1250}}
  >{\raggedright\arraybackslash}p{(\linewidth - 14\tabcolsep) * \real{0.1250}}
  >{\raggedright\arraybackslash}p{(\linewidth - 14\tabcolsep) * \real{0.1250}}
  >{\raggedright\arraybackslash}p{(\linewidth - 14\tabcolsep) * \real{0.1250}}
  >{\raggedright\arraybackslash}p{(\linewidth - 14\tabcolsep) * \real{0.1250}}@{}}
\toprule\noalign{}
\begin{minipage}[b]{\linewidth}\raggedright
RUNSPEC
\end{minipage} & \begin{minipage}[b]{\linewidth}\raggedright
GRID
\end{minipage} & \begin{minipage}[b]{\linewidth}\raggedright
EDIT
\end{minipage} & \begin{minipage}[b]{\linewidth}\raggedright
PROPS
\end{minipage} & \begin{minipage}[b]{\linewidth}\raggedright
REGIONS
\end{minipage} & \begin{minipage}[b]{\linewidth}\raggedright
SOLUTION
\end{minipage} & \begin{minipage}[b]{\linewidth}\raggedright
SUMMARY
\end{minipage} & \begin{minipage}[b]{\linewidth}\raggedright
SCHEDULE
\end{minipage} \\
\midrule\noalign{}
\endhead
\bottomrule\noalign{}
\endlastfoot
\end{longtable}

\subsubsection{Description}\label{description-944}

The GNETPUMP keyword defines the configuration of automatic compressors
and pumps in a production Standard Network, for when the Standard
Network option is invoked by the GRUPTREE, GRUPNET, GNETINJE, GNETPUMP,
etc., series of keywords in the SCHEDULE section. Although several
keywords can be used by both the Standard and Extended Network options,
GNETPUMP can only be used with the Standard Network option. The
equivalent keyword for the Extended Network option is the NETCOMPA
keyword in the RUNSPEC section.

This keyword is not supported by OPM Flow but would change the results
if supported so the simulation will be stopped.

\subsection{GPMAINT -- Define Group Pressure Maintenance Targets and
Controls}\label{gpmaint-define-group-pressure-maintenance-targets-and-controls}

\begin{longtable}[]{@{}
  >{\raggedright\arraybackslash}p{(\linewidth - 14\tabcolsep) * \real{0.1250}}
  >{\raggedright\arraybackslash}p{(\linewidth - 14\tabcolsep) * \real{0.1250}}
  >{\raggedright\arraybackslash}p{(\linewidth - 14\tabcolsep) * \real{0.1250}}
  >{\raggedright\arraybackslash}p{(\linewidth - 14\tabcolsep) * \real{0.1250}}
  >{\raggedright\arraybackslash}p{(\linewidth - 14\tabcolsep) * \real{0.1250}}
  >{\raggedright\arraybackslash}p{(\linewidth - 14\tabcolsep) * \real{0.1250}}
  >{\raggedright\arraybackslash}p{(\linewidth - 14\tabcolsep) * \real{0.1250}}
  >{\raggedright\arraybackslash}p{(\linewidth - 14\tabcolsep) * \real{0.1250}}@{}}
\toprule\noalign{}
\begin{minipage}[b]{\linewidth}\raggedright
RUNSPEC
\end{minipage} & \begin{minipage}[b]{\linewidth}\raggedright
GRID
\end{minipage} & \begin{minipage}[b]{\linewidth}\raggedright
EDIT
\end{minipage} & \begin{minipage}[b]{\linewidth}\raggedright
PROPS
\end{minipage} & \begin{minipage}[b]{\linewidth}\raggedright
REGIONS
\end{minipage} & \begin{minipage}[b]{\linewidth}\raggedright
SOLUTION
\end{minipage} & \begin{minipage}[b]{\linewidth}\raggedright
SUMMARY
\end{minipage} & \begin{minipage}[b]{\linewidth}\raggedright
SCHEDULE
\end{minipage} \\
\midrule\noalign{}
\endhead
\bottomrule\noalign{}
\endlastfoot
\end{longtable}

\subsubsection{Description}\label{description-945}

The GPMAINT keyword defines the groups under pressure maintenance
control, the associated flow rate and pressure targets, and fluid
in-place regions associated with pressure maintenance, as well as
various pressure maintenance controls. GPMAINT allows for various
regions, as defined by the FIPNUM or FIP keywords in the REGIONS
section, to have their average reservoir pressure maintained at a
specified value.

\begin{longtable}[]{@{}
  >{\raggedright\arraybackslash}p{(\linewidth - 6\tabcolsep) * \real{0.2500}}
  >{\raggedright\arraybackslash}p{(\linewidth - 6\tabcolsep) * \real{0.2500}}
  >{\raggedright\arraybackslash}p{(\linewidth - 6\tabcolsep) * \real{0.2500}}
  >{\raggedright\arraybackslash}p{(\linewidth - 6\tabcolsep) * \real{0.2500}}@{}}
\toprule\noalign{}
\begin{minipage}[b]{\linewidth}\raggedright
No.
\end{minipage} & \begin{minipage}[b]{\linewidth}\raggedright
Name
\end{minipage} & \begin{minipage}[b]{\linewidth}\raggedright
Description
\end{minipage} & \begin{minipage}[b]{\linewidth}\raggedright
Default
\end{minipage} \\
\midrule\noalign{}
\endhead
\bottomrule\noalign{}
\endlastfoot
Field & Metric & Laboratory & \\
1 & GRPNAME & A character string of up to eight characters in length
that defines the group name for which the group's associated FIPNUM
region will have a targeted average reservoir pressure maintained. The
group named FIELD is the top most group and can also be used to set
pressure target for the whole field. Note that the group hierarchy
should be defined by the GRUPTREE keyword when there is more than one
level of groups, otherwise all the groups will sit directly under the
FIELD group in the group tree hierarchy. & None \\
2 & GRPCNTL & A defined character string of length four, that sets the
production or injection control for the group, used to maintain the
average pressure for the FIPNUM region. GRPCNTL should be set to one of
the following character strings: & None \\
3 & FIPNUM & A positive integer value that defines the region for which
pressure maintenance is to be applied by GRPNAME, where FIPNUM may be
define as per the: The maximum number of FIPNUM and FIP regions is set
by the REGDIMS(NMFIPR) or the TABDIMS(NTFIP) keywords(variables) in the
RUNSPEC section. If both REGDIMS(NMFIPR) and TABDIMS(NTFIP) have been
defined then the maximum of the two is used. & None \\
0 & & & \\
4 & FIPNAME & A character string of up to five characters in length,
defining the fluid in-place's name, as defined by the FIP keyword in the
REGIONS section. For example, if the FIP keyword has been used to define
a FIP group or family, named FIPBLK-A, then FIPNAME would be set to
BLK-A. The default value of 1* means that FIPNUM on this keyword applies
to the standard FIPNUM array and not to the FIP defined group regions. &
1* \\
5 & TARGET & A real positive value that defines the average hydrocarbon
pore volume weighted reservoir pressure target for the region, for which
GRPNAME should attempt to satisfy. Note that the average pressure report
in the RPTSCHED series of reports is a pore volume weight average
reservoir pressure, not the average hydrocarbon pore volume weighted
reservoir pressure specified by TARGET on this keyword. Thus, there will
be small differences between the two numbers. & None \\
0 & & & \\
psia & barsa & atma & \\
6 & ALPHA & A real positive value that defines the proportionality
constant used to control the flow rates in order to maintain/reach the
TARGET average hydrocarbon pore volume pressure for the region. See
equation (12.26). Larger values of ALPHA accelerate the time for the
region's pressure to satisfy the required target pressure (TARGET).
However, ALPHA values that are too large may cause the region's pressure
to oscillate. Thus, the value of ALPHA should be a value that gives the
expected steady state flow rate (TARGET) of the group divided by a
reasonable transient pressure error. & None \\
Liquid: stb/d/psi Gas: Mscf/d/psi RESV: rb/d/psi & Liquid: sm3/day/bars
Gas: sm3/day/bars RESV: rm3/day/bars & Liquid: scc/hour/atm Gas:
scc/hour/atm RESV: rcc/hour/atm & \\
7 & BETA & A real positive value that defines the time interval used to
control the flow rates in order to maintain the TARGET average
hydrocarbon pore volume pressure for the region. See equation (12.26).
Larger values of BETA reduce the tendency for the regional pressure to
oscillate, but consequently decrease the rate at which the pressure
reaches its target value. This is because the pressure error is measured
at the end of the previous time step, resulting in a ``delay'' in the
response. Also, the larger the time steps the greater the propensity for
the pressure to oscillate for given ALPHA and BETA. As a guide BETA
should be set to a value that is at least as large as the maximum time
step size (see the TUNING keyword in the SCHEDULE section and section
2.2 Running OPM Flow 2023-04 From The Command Line for setting the
maximum time step size). & 1* \\
day & day & day & \\
Notes: & & & \\
\end{longtable}

\emph{Table 12.41: GPMAINT Keyword Description}

The GPMAINT keyword utilizes control theory {[}John Doyle, Bruce
Francis, Allen Tannenbaum, Feedback Control Theory, Macmillan Publishing
Co., 1990.{]}, {[}Karl J. Åström; Richard M. Murray (2008). Feedback
Systems: An Introduction for Scientists and Engineers. Princeton
University Press. ISBN 978-0-691-13576-2.{]}, {[}Andrei, Neculai
(2005)., Modern Control Theory -- A Historical Perspective, Research
Institute for Informatics, Center for Advanced Modeling and
Optimization,8-10, Averescu Avenue, Bucharest 1, Romania.{]} and
{[}Sontag, Eduardo (1998). Mathematical Control Theory: Deterministic
Finite Dimensional Systems. Second Edition, Springer. ISBN
978-0-387-98489-6.{]} to calculate the required injection and production
rates to ensure some stability in the behavior of the average
hydrocarbon pore volume reservoir pressure for a given region and the
injection rates over time. For example, if the group's flow control
(GRPCNTL) is set to one of the in situ reservoir volume injection rates
(GINJ, OINJ or WINJ), say Q, then the group's flow injection rates for
each time step is calculated by:

\begin{equation}\phantomsection\label{eq-12-26}{
Q = {Q}_{i} + \mathrm{α} \left(\left({P}_{\mathit{TARGET}} - {P}_{i-1}\right) + \frac{{\sum }_{i=1}^{i-1}\left({P}_{\mathit{TARGET}} - {P}_{i-1}\right)\times \mathrm{Δ}{t}_{i}}{\mathrm{β}}\right)
}\end{equation}

Where:

P TARGET = the regions average hydrocarbon pore volume weighted
reservoir

pressure target, the TARGET parameter in Table 12.41.

Pi-1 = the resulting average hydrocarbon pore volume weighted reservoir

pressure for the region from the preceding completed time step.

Q = the new in situ reservoir volume rate.

Qi = the initial in situ reservoir volume rate at the time the keyword
was

activated in the input. Subsequent GPMAINT keywords will reset Qi.

α = the proportionality constant used to control the flow rates, the
ALPHA

parameter in Table 12.41.

β = the time interval parameter used to control the flow rates, the BETA

parameter in Table 12.41.

The
\({\sum }_{i=1}^{i-1}\left({P}_{\mathit{TARGET}} - {P}_{i-1}\right)\times \mathrm{Δ}{t}_{i}\)term
in equation (12.26) is the cumulative sum of the average hydrocarbon
pore volume weighted reservoir pressure error for the region, times the
time step length, from when the last GPMAINT keyword was entered, up to
the preceding completed time step.

The objective of equation (12.26) is to stabilize the
injection/production rates for the group for when there are changes to
the average hydrocarbon pore volume weighted reservoir pressure target
set by the initial and subsequent GPMAINT keywords, or by any group
constraints applied afterwards. In addition, if the pressure error is
small \(\left({P}_{\mathit{TARGET}} - {P}_{i-1}\right)\)then the
equation should result in approximately the same rates as the previous
time step, if the GPMAINT has been previously entered.

However, if the first time the keyword is entered coincides with the
start of injection, then Qi in equation (12.26) by definition will be
zero, and the injection rate will start from zero and increase as
calculated by the equation. The ALPHA and BETA parameters in Table 12.41
can be used to control how the rate is built up to meet the desired
average hydrocarbon pore volume weighted reservoir pressure target.

\subsubsection{Examples}\label{examples-76}

The first example uses the gas surface rate to maintain the fields'
average hydrocarbon pore volume pressure at 225 barsa, with the ALPHA
and BETA parameters set to 40.0 and 70.0 respectively.

\begin{verbatim}
--
--       GROUP PRESSURE MAINTENANCE TARGETS AND CONTROLS
--
-- GRUP  CNTL   FIPNUM  FIP     PRESS   ALPHA  BETA
-- NAME  MODE   REGION  FIPNAME TARGET  CONST  CONST
GPMAINT
FIELD    GINS   0       1*      225     40.0   70.0        /
\end{verbatim}

/

The second example uses group's BLK-A gas surface rate to maintain the
average hydrocarbon pore volume pressure at 200 barsa for FIPNUM regions
one to four, with the ALPHA and BETA parameters set to 40.0 and 70.0
respectively.

\begin{verbatim}
--
--       GROUP PRESSURE MAINTENANCE TARGETS AND CONTROLS
--
-- GRUP  CNTL   FIPNUM  FIP     PRESS   ALPHA  BETA
-- NAME  MODE   REGION  FIPNAME TARGET  CONST  CONST
GPMAINT
BLK-A    GINS   1       1*      200     40.0   70.0        /
BLK-A    GINS   2       1*      200     40.0   70.0        /
BLK-A    GINS   3       1*      200     40.0   70.0        /
BLK-A    GINS   4       1*      200     40.0   70.0        /

BLK-B    WINJ   1       FLT-B   215     30.0   65.0        /
BLK-B    WINJ   2       FLT-B   215     30.0   65.0        /
BLK-B    WINJ   3       FLT-B   215     30.0   65.0        /
BLK-B    WINJ   4       FLT-B   215     30.0   65.0        /
/
\end{verbatim}

For group BLK-B, water in situ reservoir volume injection rate (RESV) is
used to maintain FIP group/family FLT-B's average reservoir pressure at
215 barsa for FIPFLT-B regions one to four. Here the ALPHA and BETA
parameters set to 30.0 and 65.0 respectively.

\subsection{GRADGRUP -- Define Group History Match Gradient File
Output}\label{gradgrup-define-group-history-match-gradient-file-output}

\begin{longtable}[]{@{}
  >{\raggedright\arraybackslash}p{(\linewidth - 14\tabcolsep) * \real{0.1250}}
  >{\raggedright\arraybackslash}p{(\linewidth - 14\tabcolsep) * \real{0.1250}}
  >{\raggedright\arraybackslash}p{(\linewidth - 14\tabcolsep) * \real{0.1250}}
  >{\raggedright\arraybackslash}p{(\linewidth - 14\tabcolsep) * \real{0.1250}}
  >{\raggedright\arraybackslash}p{(\linewidth - 14\tabcolsep) * \real{0.1250}}
  >{\raggedright\arraybackslash}p{(\linewidth - 14\tabcolsep) * \real{0.1250}}
  >{\raggedright\arraybackslash}p{(\linewidth - 14\tabcolsep) * \real{0.1250}}
  >{\raggedright\arraybackslash}p{(\linewidth - 14\tabcolsep) * \real{0.1250}}@{}}
\toprule\noalign{}
\begin{minipage}[b]{\linewidth}\raggedright
RUNSPEC
\end{minipage} & \begin{minipage}[b]{\linewidth}\raggedright
GRID
\end{minipage} & \begin{minipage}[b]{\linewidth}\raggedright
EDIT
\end{minipage} & \begin{minipage}[b]{\linewidth}\raggedright
PROPS
\end{minipage} & \begin{minipage}[b]{\linewidth}\raggedright
REGIONS
\end{minipage} & \begin{minipage}[b]{\linewidth}\raggedright
SOLUTION
\end{minipage} & \begin{minipage}[b]{\linewidth}\raggedright
SUMMARY
\end{minipage} & \begin{minipage}[b]{\linewidth}\raggedright
SCHEDULE
\end{minipage} \\
\midrule\noalign{}
\endhead
\bottomrule\noalign{}
\endlastfoot
\end{longtable}

\subsubsection{Description}\label{description-946}

The GRADGRUP keyword defines the SUMMARY field and group vectors that
should be written to the History Match Gradient output file, for when
the History Match Gradient option has been activated by the HMDIMS
keyword in the RUNSPEC section.

This keyword is not supported by OPM Flow but would change the results
if supported so the simulation will be stopped.

\subsection{GRADRESV -- Define Solution Derivative History Match
Gradient
Output}\label{gradresv-define-solution-derivative-history-match-gradient-output}

\begin{longtable}[]{@{}
  >{\raggedright\arraybackslash}p{(\linewidth - 14\tabcolsep) * \real{0.1250}}
  >{\raggedright\arraybackslash}p{(\linewidth - 14\tabcolsep) * \real{0.1250}}
  >{\raggedright\arraybackslash}p{(\linewidth - 14\tabcolsep) * \real{0.1250}}
  >{\raggedright\arraybackslash}p{(\linewidth - 14\tabcolsep) * \real{0.1250}}
  >{\raggedright\arraybackslash}p{(\linewidth - 14\tabcolsep) * \real{0.1250}}
  >{\raggedright\arraybackslash}p{(\linewidth - 14\tabcolsep) * \real{0.1250}}
  >{\raggedright\arraybackslash}p{(\linewidth - 14\tabcolsep) * \real{0.1250}}
  >{\raggedright\arraybackslash}p{(\linewidth - 14\tabcolsep) * \real{0.1250}}@{}}
\toprule\noalign{}
\begin{minipage}[b]{\linewidth}\raggedright
RUNSPEC
\end{minipage} & \begin{minipage}[b]{\linewidth}\raggedright
GRID
\end{minipage} & \begin{minipage}[b]{\linewidth}\raggedright
EDIT
\end{minipage} & \begin{minipage}[b]{\linewidth}\raggedright
PROPS
\end{minipage} & \begin{minipage}[b]{\linewidth}\raggedright
REGIONS
\end{minipage} & \begin{minipage}[b]{\linewidth}\raggedright
SOLUTION
\end{minipage} & \begin{minipage}[b]{\linewidth}\raggedright
SUMMARY
\end{minipage} & \begin{minipage}[b]{\linewidth}\raggedright
SCHEDULE
\end{minipage} \\
\midrule\noalign{}
\endhead
\bottomrule\noalign{}
\endlastfoot
\end{longtable}

\subsubsection{Description}\label{description-947}

The GRADRESV keyword defines the SOLUTION derivative arrays that should
be written to the History Match Gradient output file, for when the
History Match Gradient option has been activated by the HMDIMS keyword
in the RUNSPEC section.

This keyword is not supported by OPM Flow but would change the results
if supported so the simulation will be stopped.

\subsection{GRADRFT -- Define RFT Derivative History Match Gradient
Output}\label{gradrft-define-rft-derivative-history-match-gradient-output}

\begin{longtable}[]{@{}
  >{\raggedright\arraybackslash}p{(\linewidth - 14\tabcolsep) * \real{0.1250}}
  >{\raggedright\arraybackslash}p{(\linewidth - 14\tabcolsep) * \real{0.1250}}
  >{\raggedright\arraybackslash}p{(\linewidth - 14\tabcolsep) * \real{0.1250}}
  >{\raggedright\arraybackslash}p{(\linewidth - 14\tabcolsep) * \real{0.1250}}
  >{\raggedright\arraybackslash}p{(\linewidth - 14\tabcolsep) * \real{0.1250}}
  >{\raggedright\arraybackslash}p{(\linewidth - 14\tabcolsep) * \real{0.1250}}
  >{\raggedright\arraybackslash}p{(\linewidth - 14\tabcolsep) * \real{0.1250}}
  >{\raggedright\arraybackslash}p{(\linewidth - 14\tabcolsep) * \real{0.1250}}@{}}
\toprule\noalign{}
\begin{minipage}[b]{\linewidth}\raggedright
RUNSPEC
\end{minipage} & \begin{minipage}[b]{\linewidth}\raggedright
GRID
\end{minipage} & \begin{minipage}[b]{\linewidth}\raggedright
EDIT
\end{minipage} & \begin{minipage}[b]{\linewidth}\raggedright
PROPS
\end{minipage} & \begin{minipage}[b]{\linewidth}\raggedright
REGIONS
\end{minipage} & \begin{minipage}[b]{\linewidth}\raggedright
SOLUTION
\end{minipage} & \begin{minipage}[b]{\linewidth}\raggedright
SUMMARY
\end{minipage} & \begin{minipage}[b]{\linewidth}\raggedright
SCHEDULE
\end{minipage} \\
\midrule\noalign{}
\endhead
\bottomrule\noalign{}
\endlastfoot
\end{longtable}

\subsubsection{Description}\label{description-948}

The GRADRFT keyword defines the derivative well RFT data, the SOLUTION
pressure and saturations at a well's connected grid block, that should
be written to the History Match Gradient output file, for when the
History Match Gradient option has been activated by the HMDIMS keyword
in the RUNSPEC section.

This keyword is not supported by OPM Flow but would change the results
if supported so the simulation will be stopped.

\subsection{GRADWELL -- Define Well History Match Gradient File
Output}\label{gradwell-define-well-history-match-gradient-file-output}

\begin{longtable}[]{@{}
  >{\raggedright\arraybackslash}p{(\linewidth - 14\tabcolsep) * \real{0.1250}}
  >{\raggedright\arraybackslash}p{(\linewidth - 14\tabcolsep) * \real{0.1250}}
  >{\raggedright\arraybackslash}p{(\linewidth - 14\tabcolsep) * \real{0.1250}}
  >{\raggedright\arraybackslash}p{(\linewidth - 14\tabcolsep) * \real{0.1250}}
  >{\raggedright\arraybackslash}p{(\linewidth - 14\tabcolsep) * \real{0.1250}}
  >{\raggedright\arraybackslash}p{(\linewidth - 14\tabcolsep) * \real{0.1250}}
  >{\raggedright\arraybackslash}p{(\linewidth - 14\tabcolsep) * \real{0.1250}}
  >{\raggedright\arraybackslash}p{(\linewidth - 14\tabcolsep) * \real{0.1250}}@{}}
\toprule\noalign{}
\begin{minipage}[b]{\linewidth}\raggedright
RUNSPEC
\end{minipage} & \begin{minipage}[b]{\linewidth}\raggedright
GRID
\end{minipage} & \begin{minipage}[b]{\linewidth}\raggedright
EDIT
\end{minipage} & \begin{minipage}[b]{\linewidth}\raggedright
PROPS
\end{minipage} & \begin{minipage}[b]{\linewidth}\raggedright
REGIONS
\end{minipage} & \begin{minipage}[b]{\linewidth}\raggedright
SOLUTION
\end{minipage} & \begin{minipage}[b]{\linewidth}\raggedright
SUMMARY
\end{minipage} & \begin{minipage}[b]{\linewidth}\raggedright
SCHEDULE
\end{minipage} \\
\midrule\noalign{}
\endhead
\bottomrule\noalign{}
\endlastfoot
\end{longtable}

\subsubsection{Description}\label{description-949}

The GRADWELL keyword defines the SUMMARY well vectors that should be
written to the History Match Gradient output file, for when the History
Match Gradient option has been activated by the HMDIMS keyword in the
RUNSPEC section.

This keyword is not supported by OPM Flow but would change the results
if supported so the simulation will be stopped.

\subsection{GRDREACH -- Define River and Grid Block
Connections}\label{grdreach-define-river-and-grid-block-connections}

\begin{longtable}[]{@{}
  >{\raggedright\arraybackslash}p{(\linewidth - 14\tabcolsep) * \real{0.1250}}
  >{\raggedright\arraybackslash}p{(\linewidth - 14\tabcolsep) * \real{0.1250}}
  >{\raggedright\arraybackslash}p{(\linewidth - 14\tabcolsep) * \real{0.1250}}
  >{\raggedright\arraybackslash}p{(\linewidth - 14\tabcolsep) * \real{0.1250}}
  >{\raggedright\arraybackslash}p{(\linewidth - 14\tabcolsep) * \real{0.1250}}
  >{\raggedright\arraybackslash}p{(\linewidth - 14\tabcolsep) * \real{0.1250}}
  >{\raggedright\arraybackslash}p{(\linewidth - 14\tabcolsep) * \real{0.1250}}
  >{\raggedright\arraybackslash}p{(\linewidth - 14\tabcolsep) * \real{0.1250}}@{}}
\toprule\noalign{}
\begin{minipage}[b]{\linewidth}\raggedright
RUNSPEC
\end{minipage} & \begin{minipage}[b]{\linewidth}\raggedright
GRID
\end{minipage} & \begin{minipage}[b]{\linewidth}\raggedright
EDIT
\end{minipage} & \begin{minipage}[b]{\linewidth}\raggedright
PROPS
\end{minipage} & \begin{minipage}[b]{\linewidth}\raggedright
REGIONS
\end{minipage} & \begin{minipage}[b]{\linewidth}\raggedright
SOLUTION
\end{minipage} & \begin{minipage}[b]{\linewidth}\raggedright
SUMMARY
\end{minipage} & \begin{minipage}[b]{\linewidth}\raggedright
SCHEDULE
\end{minipage} \\
\midrule\noalign{}
\endhead
\bottomrule\noalign{}
\endlastfoot
\end{longtable}

\subsubsection{Description}\label{description-950}

The GRDREACH keyword defines the location of grid blocks connecting to a
previously defined river, for when the River option has been activated
via the RIVRDIMS keyword in the RUNSPEC section.

This keyword is not supported by OPM Flow but would change the results
if supported so the simulation will be stopped.

\subsection{GRUPMAST -- Define Master and Slave
Groups}\label{grupmast-define-master-and-slave-groups}

\begin{longtable}[]{@{}
  >{\raggedright\arraybackslash}p{(\linewidth - 14\tabcolsep) * \real{0.1250}}
  >{\raggedright\arraybackslash}p{(\linewidth - 14\tabcolsep) * \real{0.1250}}
  >{\raggedright\arraybackslash}p{(\linewidth - 14\tabcolsep) * \real{0.1250}}
  >{\raggedright\arraybackslash}p{(\linewidth - 14\tabcolsep) * \real{0.1250}}
  >{\raggedright\arraybackslash}p{(\linewidth - 14\tabcolsep) * \real{0.1250}}
  >{\raggedright\arraybackslash}p{(\linewidth - 14\tabcolsep) * \real{0.1250}}
  >{\raggedright\arraybackslash}p{(\linewidth - 14\tabcolsep) * \real{0.1250}}
  >{\raggedright\arraybackslash}p{(\linewidth - 14\tabcolsep) * \real{0.1250}}@{}}
\toprule\noalign{}
\begin{minipage}[b]{\linewidth}\raggedright
RUNSPEC
\end{minipage} & \begin{minipage}[b]{\linewidth}\raggedright
GRID
\end{minipage} & \begin{minipage}[b]{\linewidth}\raggedright
EDIT
\end{minipage} & \begin{minipage}[b]{\linewidth}\raggedright
PROPS
\end{minipage} & \begin{minipage}[b]{\linewidth}\raggedright
REGIONS
\end{minipage} & \begin{minipage}[b]{\linewidth}\raggedright
SOLUTION
\end{minipage} & \begin{minipage}[b]{\linewidth}\raggedright
SUMMARY
\end{minipage} & \begin{minipage}[b]{\linewidth}\raggedright
SCHEDULE
\end{minipage} \\
\midrule\noalign{}
\endhead
\bottomrule\noalign{}
\endlastfoot
\end{longtable}

\subsubsection{Description}\label{description-951}

The GRUPMAST keyword defines master groups and their associated slave
groups for when the Reservoir Coupling option as been activated by the
GRUPMAST and SLAVES keywords in the SCHEDULE section. Reservoir coupling
allows for independent reservoir simulation decks (SLAVES) to be
controlled by a separate master run file. For example, if there are five
separate reservoir models each representing one field, one of the four
would be used as the master and the other four would be the subordinate
SLAVES.

This keyword is not supported by OPM Flow but would change the results
if supported so the simulation will be stopped.

\subsection{Atgeirr Rasmussen}\label{atgeirr-rasmussen}

\begin{verbatim}
  2017-09-20T13:28:08
  AFR
  I don’t think this is supported.
 GRUPNET
  David Baxendale
  2017-10-03T16:14:46.395000000
  DBx
  Reply to Atgeirr Rasmussen (09/20/2017, 13:28): "..."
  But it is in the Norne model to model the surface network, will switch to orange for now.
  – Define Group Standard Network Parameters
\end{verbatim}

\begin{longtable}[]{@{}
  >{\raggedright\arraybackslash}p{(\linewidth - 14\tabcolsep) * \real{0.1250}}
  >{\raggedright\arraybackslash}p{(\linewidth - 14\tabcolsep) * \real{0.1250}}
  >{\raggedright\arraybackslash}p{(\linewidth - 14\tabcolsep) * \real{0.1250}}
  >{\raggedright\arraybackslash}p{(\linewidth - 14\tabcolsep) * \real{0.1250}}
  >{\raggedright\arraybackslash}p{(\linewidth - 14\tabcolsep) * \real{0.1250}}
  >{\raggedright\arraybackslash}p{(\linewidth - 14\tabcolsep) * \real{0.1250}}
  >{\raggedright\arraybackslash}p{(\linewidth - 14\tabcolsep) * \real{0.1250}}
  >{\raggedright\arraybackslash}p{(\linewidth - 14\tabcolsep) * \real{0.1250}}@{}}
\toprule\noalign{}
\begin{minipage}[b]{\linewidth}\raggedright
RUNSPEC
\end{minipage} & \begin{minipage}[b]{\linewidth}\raggedright
GRID
\end{minipage} & \begin{minipage}[b]{\linewidth}\raggedright
EDIT
\end{minipage} & \begin{minipage}[b]{\linewidth}\raggedright
PROPS
\end{minipage} & \begin{minipage}[b]{\linewidth}\raggedright
REGIONS
\end{minipage} & \begin{minipage}[b]{\linewidth}\raggedright
SOLUTION
\end{minipage} & \begin{minipage}[b]{\linewidth}\raggedright
SUMMARY
\end{minipage} & \begin{minipage}[b]{\linewidth}\raggedright
SCHEDULE
\end{minipage} \\
\midrule\noalign{}
\endhead
\bottomrule\noalign{}
\endlastfoot
\end{longtable}

\subsubsection{Description}\label{description-952}

The GRUPNET keyword defines the standard group network parameters used
to model the flow and pressures behavior within the network. The group
hierarchy is defined by the GRUPTREE keyword and wells are assigned to
groups using the WELSPECS keyword, both keywords are in the SCHEDULE
section.

Group pressure values are optionally entered for each group in the
network together with a vertical flow (lift) performance (``VFP'') table
that the determines the pipeline pressure behavior from the LOWER to the
HIGHER group given the current flowing conditions; the group
relationship is defined by the GRUPTREE keyword. The VFP table is
entered via the VFPPROD keyword for production pipelines and VFPINJ for
injection pipelines. These are the same keywords as used for well
modeling, however, the manner in which they are generated by an external
program is completely different when used for pipeline modeling.

\begin{longtable}[]{@{}
  >{\raggedright\arraybackslash}p{(\linewidth - 6\tabcolsep) * \real{0.2500}}
  >{\raggedright\arraybackslash}p{(\linewidth - 6\tabcolsep) * \real{0.2500}}
  >{\raggedright\arraybackslash}p{(\linewidth - 6\tabcolsep) * \real{0.2500}}
  >{\raggedright\arraybackslash}p{(\linewidth - 6\tabcolsep) * \real{0.2500}}@{}}
\toprule\noalign{}
\begin{minipage}[b]{\linewidth}\raggedright
No.
\end{minipage} & \begin{minipage}[b]{\linewidth}\raggedright
Name
\end{minipage} & \begin{minipage}[b]{\linewidth}\raggedright
Description
\end{minipage} & \begin{minipage}[b]{\linewidth}\raggedright
Default
\end{minipage} \\
\midrule\noalign{}
\endhead
\bottomrule\noalign{}
\endlastfoot
Field & Metric & Laboratory & \\
1 & GRPNAME & A character string of up to eight characters in length
that defines the group name for which the network parameters are being
defined. The group named FIELD is the top most group and may be used as
a GRPNAME. & None \\
2 & PRES & A real value that defines the fixed pressure for this group
when the group is a terminating group. If the group is not a terminating
group then PRES should be defaulted with 1* or set to a negative number.
& 1* \\
psia & barsa & atma & \\
3 & VFPTAB & A positive integer greater than or equal to zero that
defines the VFPPROD or VFPINJ vertical lift performance table to be used
for calculating the pipeline pressures connecting the LOWER and HIGHER
groups in the network. Note that: If a non-zero value is entered then
the vertical lift performance tables must be entered via the VFPPROD or
VFPINJ keyword in the SCHEDULE section. & 0 \\
4 & ALQ-PIPE & A real positive value that defines the artificial lift
quantity to be used in conjunction with the vertical lift performance
table (VFPPROD) assigned to the group via VPFTAB variable. The vertical
lift performance table and the artificial lift quantity ALQ-WELL are
used with the pipeline fluid rates to calculate the pipeline change
pressure between the LOWER and HIGHER groups. Note that the units for
ALQ-PIPE are dependent on the associated variable on the VFPPROD keyword
and may represent a pump or a compressor depending how the VFPPROD table
was generated by an external program. & 0.0 \\
5 & OPTION1 & A defined character string that defines if a group's
production target should be achieved by adjusting the tubing pressure of
the wells within the group or by the adjusting the well rates by their
guide rate. OPTION1 should be set to one of the following character
strings: If a group is using this option and has a higher group with
production targets or constraints, then this group should have it's
guide rate set via the GCONPROD keyword in the SCHEDULE section, to
ensure that the well's within this group operate at the same tubing head
pressure. Only groups containing wells can use OPTION1 equal to YES or
NO, a group without wells should set OPTION1 to NO. Numerical
convergence controls and iteration limits for wells using OPTION1 set
equal to YES are defined via the NETBALAN keyword in the SCHEDULE
section. Only the NO option is currently supported by the simulator. &
NO \\
6 & OPTION2 & A defined character string that defines if well gas lift
gas flows through the group's pipeline. OPTION2 should be set to one of
the following character strings: If either FLO or ALQ have been selected
then artificial lift quantity for the pipeline (ALQ-PIPE) and the wells
(ALQ-WELL) must be defined as gas lift gas on the VFPPROD tables. A
well's specific gas lift gas quantity is set via the ALQ-WELL variable
on the WCONPROD keyword in the SCHEDULE section. Only the NO and FLO
options are currently supported by the simulator. & NO \\
7 & OPTION3 & A defined character string that defines if the ALQ-PIPE
variable should be reset to an equivalent surface oil or gas density
flowing along the pipeline. OPTION3 should be set to one of the
following character strings: If either DENO or DENG have been selected
then artificial lift quantity on the VFPPROD tables must be based on the
same density parameter. These options are normally used when a mixture
of oil or gas with different surface densities flows into the network.
Only the NONE option is currently supported by the simulator. & NONE \\
Notes: & & & \\
\end{longtable}

\emph{Table 12.3.113.1: GRUPNET Keyword Description}

See also the WELSPECS keyword to define wells, the VFPPROD and VFPINJ
keywords that the define vertical lift performance table to be used for
calculating the pipeline pressures connecting the LOWER and HIGHER group
in the network. All the aforementioned keywords are described in the
SCHEDULE section.

\subsubsection{Examples}\label{examples-77}

The following example defines a network based on two groups

\begin{verbatim}
--
--       DEFINE GROUP STANDARD NETWORK PARAMETERS
--
-- GRUP  CNTL   VFP    PUMP   MANIFOLD  INCLUDE   ALQ
-- NAME  PRES   TABLE  POWER  GROUP     LIFT GAS  DENS
GRUPNET
PROD-A   1200.  1*                                                             /
PROD-B   1*     1      1*    'YES'      1*        1*                           /
/
\end{verbatim}

The next example is more complex and is taken form the Norne model.

\begin{verbatim}
--
--       DEFINE GROUP STANDARD NETWORK PARAMETERS
--
-- GRUP  CNTL   VFP    PUMP   MANIFOLD  INCLUDE   ALQ
-- NAME  PRES   TABLE  POWER  GROUP     LIFT GAS  DENS
GRUPNET
FIELD    20.0   5*                                                             /
PROD     20.0   5*                                                             /
MANI-B2  1*     8      1*      NO       2*                                     /
MANI-B1  1*     8      1*      NO       2*                                     /
MANI-K1  1*    9999    4*                                                      /
B1-DUMMY 1*    9999    4*                                                      /
MANI-D1  1*     8      1*      NO       2*                                     /
MANI-D2  1*     8      1*      NO       2*                                     /
MANI-K2  1*    9999    4*                                                      /
D2-DUMMY 1*    9999    4*                                                      /
MANI-E1  1*     9      1*      NO       2*                                     /
MANI-E2  1*     9      4*                                                      /
/
\end{verbatim}

Here the FIELD controlling pressure is set at 20 barsa and the same
limit is used for group PROD which sits directly under the FIELD group
(see Figure 12.4).

\subsection{GRUPRIG -- Group Drilling and Workover Rig
Specifications}\label{gruprig-group-drilling-and-workover-rig-specifications}

\begin{longtable}[]{@{}
  >{\raggedright\arraybackslash}p{(\linewidth - 14\tabcolsep) * \real{0.1250}}
  >{\raggedright\arraybackslash}p{(\linewidth - 14\tabcolsep) * \real{0.1250}}
  >{\raggedright\arraybackslash}p{(\linewidth - 14\tabcolsep) * \real{0.1250}}
  >{\raggedright\arraybackslash}p{(\linewidth - 14\tabcolsep) * \real{0.1250}}
  >{\raggedright\arraybackslash}p{(\linewidth - 14\tabcolsep) * \real{0.1250}}
  >{\raggedright\arraybackslash}p{(\linewidth - 14\tabcolsep) * \real{0.1250}}
  >{\raggedright\arraybackslash}p{(\linewidth - 14\tabcolsep) * \real{0.1250}}
  >{\raggedright\arraybackslash}p{(\linewidth - 14\tabcolsep) * \real{0.1250}}@{}}
\toprule\noalign{}
\begin{minipage}[b]{\linewidth}\raggedright
RUNSPEC
\end{minipage} & \begin{minipage}[b]{\linewidth}\raggedright
GRID
\end{minipage} & \begin{minipage}[b]{\linewidth}\raggedright
EDIT
\end{minipage} & \begin{minipage}[b]{\linewidth}\raggedright
PROPS
\end{minipage} & \begin{minipage}[b]{\linewidth}\raggedright
REGIONS
\end{minipage} & \begin{minipage}[b]{\linewidth}\raggedright
SOLUTION
\end{minipage} & \begin{minipage}[b]{\linewidth}\raggedright
SUMMARY
\end{minipage} & \begin{minipage}[b]{\linewidth}\raggedright
SCHEDULE
\end{minipage} \\
\midrule\noalign{}
\endhead
\bottomrule\noalign{}
\endlastfoot
\end{longtable}

\subsubsection{Description}\label{description-953}

Defines a groups drilling and workover specifications.

This keyword is not supported by OPM Flow but would change the results
if supported so the simulation will be stopped.

\subsection{GRUPSLAV -- Define Slave Groups in Slave
Reservoirs}\label{grupslav-define-slave-groups-in-slave-reservoirs}

\begin{longtable}[]{@{}
  >{\raggedright\arraybackslash}p{(\linewidth - 14\tabcolsep) * \real{0.1250}}
  >{\raggedright\arraybackslash}p{(\linewidth - 14\tabcolsep) * \real{0.1250}}
  >{\raggedright\arraybackslash}p{(\linewidth - 14\tabcolsep) * \real{0.1250}}
  >{\raggedright\arraybackslash}p{(\linewidth - 14\tabcolsep) * \real{0.1250}}
  >{\raggedright\arraybackslash}p{(\linewidth - 14\tabcolsep) * \real{0.1250}}
  >{\raggedright\arraybackslash}p{(\linewidth - 14\tabcolsep) * \real{0.1250}}
  >{\raggedright\arraybackslash}p{(\linewidth - 14\tabcolsep) * \real{0.1250}}
  >{\raggedright\arraybackslash}p{(\linewidth - 14\tabcolsep) * \real{0.1250}}@{}}
\toprule\noalign{}
\begin{minipage}[b]{\linewidth}\raggedright
RUNSPEC
\end{minipage} & \begin{minipage}[b]{\linewidth}\raggedright
GRID
\end{minipage} & \begin{minipage}[b]{\linewidth}\raggedright
EDIT
\end{minipage} & \begin{minipage}[b]{\linewidth}\raggedright
PROPS
\end{minipage} & \begin{minipage}[b]{\linewidth}\raggedright
REGIONS
\end{minipage} & \begin{minipage}[b]{\linewidth}\raggedright
SOLUTION
\end{minipage} & \begin{minipage}[b]{\linewidth}\raggedright
SUMMARY
\end{minipage} & \begin{minipage}[b]{\linewidth}\raggedright
SCHEDULE
\end{minipage} \\
\midrule\noalign{}
\endhead
\bottomrule\noalign{}
\endlastfoot
\end{longtable}

\subsubsection{Description}\label{description-954}

The GRUPSLAV keyword defines slave groups in a slave input deck and
their associated master groups in the master run, for when the Reservoir
Coupling option as been activated by the GRUPMAST and SLAVES keywords in
the SCHEDULE section. This keyword is required for every slave input
deck. Reservoir coupling allows for independent reservoir simulation
decks (SLAVES) to be controlled by a separate master run file. For
example, if there are five separate reservoir models each representing
one field, one of the four would be used as the master and the other
four would be the subordinate SLAVES.

This keyword is not supported by OPM Flow but would change the results
if supported so the simulation will be stopped.

\subsection{GRUPTARG -- Modify Group Targets and Constraints
Values}\label{gruptarg-modify-group-targets-and-constraints-values}

\begin{longtable}[]{@{}
  >{\raggedright\arraybackslash}p{(\linewidth - 14\tabcolsep) * \real{0.1250}}
  >{\raggedright\arraybackslash}p{(\linewidth - 14\tabcolsep) * \real{0.1250}}
  >{\raggedright\arraybackslash}p{(\linewidth - 14\tabcolsep) * \real{0.1250}}
  >{\raggedright\arraybackslash}p{(\linewidth - 14\tabcolsep) * \real{0.1250}}
  >{\raggedright\arraybackslash}p{(\linewidth - 14\tabcolsep) * \real{0.1250}}
  >{\raggedright\arraybackslash}p{(\linewidth - 14\tabcolsep) * \real{0.1250}}
  >{\raggedright\arraybackslash}p{(\linewidth - 14\tabcolsep) * \real{0.1250}}
  >{\raggedright\arraybackslash}p{(\linewidth - 14\tabcolsep) * \real{0.1250}}@{}}
\toprule\noalign{}
\begin{minipage}[b]{\linewidth}\raggedright
RUNSPEC
\end{minipage} & \begin{minipage}[b]{\linewidth}\raggedright
GRID
\end{minipage} & \begin{minipage}[b]{\linewidth}\raggedright
EDIT
\end{minipage} & \begin{minipage}[b]{\linewidth}\raggedright
PROPS
\end{minipage} & \begin{minipage}[b]{\linewidth}\raggedright
REGIONS
\end{minipage} & \begin{minipage}[b]{\linewidth}\raggedright
SOLUTION
\end{minipage} & \begin{minipage}[b]{\linewidth}\raggedright
SUMMARY
\end{minipage} & \begin{minipage}[b]{\linewidth}\raggedright
SCHEDULE
\end{minipage} \\
\midrule\noalign{}
\endhead
\bottomrule\noalign{}
\endlastfoot
\end{longtable}

\subsubsection{Description}\label{description-955}

The GRUPTARG keyword modifies the target and constraints values of both
rates and pressures for previously defined groups without having to
define all the variables on the group control keywords: GCONPROD or
GCONPRI keywords. Variables not changed by the GRUPTARG keyword remain
the same as those previously entered via the group control keywords or
previously entered GRUPTARG keywords. Note that the group must still be
initially be fully defined using the GCONPROD or GCONPRI keywords. All
the aforementioned keywords are described in the SCHEDULE section.

This keyword is not supported by OPM Flow but would change the results
if supported so the simulation will be stopped.

Note that wells are allocated to groups when the wells are specified by
the WELSPECS keyword in the SCHEDULE section. Wells defined to be under
group control will have their production rates controlled by the group
to which they belong, in addition to any well constraints defined for
the wells.

\begin{longtable}[]{@{}
  >{\raggedright\arraybackslash}p{(\linewidth - 6\tabcolsep) * \real{0.2500}}
  >{\raggedright\arraybackslash}p{(\linewidth - 6\tabcolsep) * \real{0.2500}}
  >{\raggedright\arraybackslash}p{(\linewidth - 6\tabcolsep) * \real{0.2500}}
  >{\raggedright\arraybackslash}p{(\linewidth - 6\tabcolsep) * \real{0.2500}}@{}}
\toprule\noalign{}
\begin{minipage}[b]{\linewidth}\raggedright
No.
\end{minipage} & \begin{minipage}[b]{\linewidth}\raggedright
Name
\end{minipage} & \begin{minipage}[b]{\linewidth}\raggedright
Description
\end{minipage} & \begin{minipage}[b]{\linewidth}\raggedright
Default
\end{minipage} \\
\midrule\noalign{}
\endhead
\bottomrule\noalign{}
\endlastfoot
Field & Metric & Laboratory & \\
1 & GRPNAME & A character string of up to eight characters in length
that defines the group name for which the group target and constraints
are being defined. The group named FIELD is the top most group and
should be used to set targets and constraints for the field. Note that
the group hierarchy should be defined by the GRUPTREE keyword when there
is more than one level of groups, otherwise all the groups will sit
directly under the FIELD group in the group tree hierarchy. & None \\
2 & TARGET & A defined character string that sets the item to be changed
for the group the value of the item is set by item (3). & None \\
3 & VALUE Liquid Gas Res Vol Pressure & A real positive value that
defines the value of the variable declared by TARGET & None \\
stb/d Mscf/d rb/d psia & sm3/day sm3/day rm3/day barsa & scc/hour
scc/hour rcc/hour atma & \\
Notes: & & & \\
\end{longtable}

\emph{Table 12.43: GRUPTARG Keyword Description}

See also the WELTARG and WELCNTL keyword, in the SCHEDULE section that
can be used to reset a well's control mode, as well as a well's target
and constraints of both rates and pressures.

\subsubsection{Example}\label{example-535}

The following example below shows the oil rates for the field at the
start of the schedule section (January 1, 2000).

\begin{verbatim}
-- ------------------------------------------------------------------------------
-- 01 JAN 2000 START OF SCHEDULE SECTION
-- ------------------------------------------------------------------------------
--
--       GROUP PRODUCTION CONTROLS
--
-- GRUP  CNTL  OIL    WAT    GAS    LIQ    CNTL  GRUP  GUIDE  GUIDE  CNTL
-- NAME  MODE  RATE   RATE   RATE   RATE   OPT   CNTL  RATE   DEF    WAT
GCONPROD
FIELD    ORAT  40E3   60E3   30E3   65E3   1*     1*    1*     1*     1*      /
/                                                                               DATES
01 FEB 2000 /
/
--
--       GROUP PRODUCTION AND INJECTION TARGETS
--
--       GROUP    GROUP  TARGET
--       NAME     TARG   VALUE
GRUPTARG
         FIELD    ORAT   45E3                                                  /
         FIELD    LIQ    75E3                                                  /
/
\end{verbatim}

From January 1, 2000 to February 1, 2000 the field is on oil rate
control and has a target oil rate of 40,000 stb/d, a maximum water
handling capacity of 60,000 stb/d, a maximum liquid capacity of 65,000
stb/d, and a maximum gas constraint of 30 MMscf/d.~After February 1,
2000 the field's target oil rate is increased to 45,000 stb/d and the
maximum liquid constraint is increased to 75,000 stb/s; all the other
parameters remain unchanged.

\subsection{GRUPTREE -- Define Group Tree
Hierarchy}\label{gruptree-define-group-tree-hierarchy}

\begin{longtable}[]{@{}
  >{\raggedright\arraybackslash}p{(\linewidth - 14\tabcolsep) * \real{0.1250}}
  >{\raggedright\arraybackslash}p{(\linewidth - 14\tabcolsep) * \real{0.1250}}
  >{\raggedright\arraybackslash}p{(\linewidth - 14\tabcolsep) * \real{0.1250}}
  >{\raggedright\arraybackslash}p{(\linewidth - 14\tabcolsep) * \real{0.1250}}
  >{\raggedright\arraybackslash}p{(\linewidth - 14\tabcolsep) * \real{0.1250}}
  >{\raggedright\arraybackslash}p{(\linewidth - 14\tabcolsep) * \real{0.1250}}
  >{\raggedright\arraybackslash}p{(\linewidth - 14\tabcolsep) * \real{0.1250}}
  >{\raggedright\arraybackslash}p{(\linewidth - 14\tabcolsep) * \real{0.1250}}@{}}
\toprule\noalign{}
\begin{minipage}[b]{\linewidth}\raggedright
RUNSPEC
\end{minipage} & \begin{minipage}[b]{\linewidth}\raggedright
GRID
\end{minipage} & \begin{minipage}[b]{\linewidth}\raggedright
EDIT
\end{minipage} & \begin{minipage}[b]{\linewidth}\raggedright
PROPS
\end{minipage} & \begin{minipage}[b]{\linewidth}\raggedright
REGIONS
\end{minipage} & \begin{minipage}[b]{\linewidth}\raggedright
SOLUTION
\end{minipage} & \begin{minipage}[b]{\linewidth}\raggedright
SUMMARY
\end{minipage} & \begin{minipage}[b]{\linewidth}\raggedright
SCHEDULE
\end{minipage} \\
\midrule\noalign{}
\endhead
\bottomrule\noalign{}
\endlastfoot
\end{longtable}

\subsubsection{Description}\label{description-956}

GRUPTREE defines the group hierarchy of groups that have been created by
having wells assigned to them via the WELSPECS keyword in the SCHEDULE
section, By default three group levels are defined that sets the wells
as level three, reporting directly to defined groups at level two, and
the level two groups reporting to the FIELD group at level one. If a
different configuration is required then the GRUPTREE keyword should be
used to define the group hierarchy by defining a lower level group that
reports directly to a higher level group.

\begin{longtable}[]{@{}
  >{\raggedright\arraybackslash}p{(\linewidth - 6\tabcolsep) * \real{0.2500}}
  >{\raggedright\arraybackslash}p{(\linewidth - 6\tabcolsep) * \real{0.2500}}
  >{\raggedright\arraybackslash}p{(\linewidth - 6\tabcolsep) * \real{0.2500}}
  >{\raggedright\arraybackslash}p{(\linewidth - 6\tabcolsep) * \real{0.2500}}@{}}
\toprule\noalign{}
\begin{minipage}[b]{\linewidth}\raggedright
No.
\end{minipage} & \begin{minipage}[b]{\linewidth}\raggedright
Name
\end{minipage} & \begin{minipage}[b]{\linewidth}\raggedright
Description
\end{minipage} & \begin{minipage}[b]{\linewidth}\raggedright
Default
\end{minipage} \\
\midrule\noalign{}
\endhead
\bottomrule\noalign{}
\endlastfoot
1 & LOWER & A character string of up to eight characters in length that
defines the group name which belongs to the HIGHER group. The group
named FIELD is the top most group and should NOT be used as as a group
name for the LOWER group name. Undefined group relationships are
automatically assigned to the FIELD group. & None \\
2 & HIGHER & A character string of up to eight characters in length that
defines the HIGHER group name that the LOWER group belongs to. The group
named FIELD is the top most group and can be used as as the HIGHER group
name. Undefined group relationships are automatically assigned to the
FIELD group. & None \\
Notes: & & & \\
\end{longtable}

\emph{Table 12.44: GRUPTREE Keyword Description}

A group hierarchy can have any number of levels but groups that have
other groups as LOWER groups cannot also have wells for the HIGHER
group. Thus, a group either contains wells or has LOWER groups

See also the GCONPROD and GCONINJE for defining group production and
injection volumes, and the WELSPECS keywords to allocate wells to
groups. All the aforementioned keywords are described in the SCHEDULE
section.

\subsubsection{Examples}\label{examples-78}

The first example defines PLAT01 and PLAT03 reporting to the FIELD level
(default if these records are omitted) and PLAT02 reporting to PLAT01.

\begin{verbatim}
--
--       DEFINE GROUP TREE HIERARCHY
--
--       LOWER     HIGHER
--       GROUP     GROUP
GRUPTREE
         PLAT01    FIELD                                                       /
         PLAT02    PLAT01                                                      /
         PLAT03    FIELD                                                       /
/
\end{verbatim}

The next example is more complex and is taken form the Norne model.

\begin{verbatim}
--
--       DEFINE GROUP TREE HIERARCHY
--
--       LOWER      HIGHER
--       GROUP      GROUP
GRUPTREE
        'INJE'     'FIELD'                                                     /
        'PROD'     'FIELD'                                                     /
        'MANI-B2'  'PROD'                                                      /
        'MANI-B1'  'PROD'                                                      /
        'MANI-D1'  'PROD'                                                      /
        'MANI-D2'  'PROD'                                                      /
        'MANI-E1'  'PROD'                                                      /
        'MANI-E2'  'PROD'                                                      /
        'MANI-K1'  'MANI-B1'                                                   /
        'MANI-K2'  'MANI-D2'                                                   /
        'MANI-C'   'INJE'                                                      /
        'MANI-F'   'INJE'                                                      /
        'WI-GSEG'  'INJE'                                                      /
        'B1-DUMMY' 'MANI-B1'                                                   /
        'D2-DUMMY' 'MANI-D2'                                                   /
/
\end{verbatim}

The group hierarchy for this example is shown below.

\pandocbounded{\includegraphics[keepaspectratio]{chapters/images/Frame156_e13710c958a3.png}}
\pandocbounded{\includegraphics[keepaspectratio]{chapters/images/Image149_e13710c958a3.png}}

Here groups PROD, INJ, MAN-B1, and MAN-D2 report to higher level groups
and the other remaining groups all have individual wells allocated to
them instead.

\subsection{GSATINJE -- Define Group Satellite Injection
Rates}\label{gsatinje-define-group-satellite-injection-rates}

\begin{longtable}[]{@{}
  >{\raggedright\arraybackslash}p{(\linewidth - 14\tabcolsep) * \real{0.1250}}
  >{\raggedright\arraybackslash}p{(\linewidth - 14\tabcolsep) * \real{0.1250}}
  >{\raggedright\arraybackslash}p{(\linewidth - 14\tabcolsep) * \real{0.1250}}
  >{\raggedright\arraybackslash}p{(\linewidth - 14\tabcolsep) * \real{0.1250}}
  >{\raggedright\arraybackslash}p{(\linewidth - 14\tabcolsep) * \real{0.1250}}
  >{\raggedright\arraybackslash}p{(\linewidth - 14\tabcolsep) * \real{0.1250}}
  >{\raggedright\arraybackslash}p{(\linewidth - 14\tabcolsep) * \real{0.1250}}
  >{\raggedright\arraybackslash}p{(\linewidth - 14\tabcolsep) * \real{0.1250}}@{}}
\toprule\noalign{}
\begin{minipage}[b]{\linewidth}\raggedright
RUNSPEC
\end{minipage} & \begin{minipage}[b]{\linewidth}\raggedright
GRID
\end{minipage} & \begin{minipage}[b]{\linewidth}\raggedright
EDIT
\end{minipage} & \begin{minipage}[b]{\linewidth}\raggedright
PROPS
\end{minipage} & \begin{minipage}[b]{\linewidth}\raggedright
REGIONS
\end{minipage} & \begin{minipage}[b]{\linewidth}\raggedright
SOLUTION
\end{minipage} & \begin{minipage}[b]{\linewidth}\raggedright
SUMMARY
\end{minipage} & \begin{minipage}[b]{\linewidth}\raggedright
SCHEDULE
\end{minipage} \\
\midrule\noalign{}
\endhead
\bottomrule\noalign{}
\endlastfoot
\end{longtable}

\subsubsection{Description}\label{description-957}

The GSATINJE keyword defines a satellite group's oil, gas and water
injection rates in the model. Satellite groups are not connected to the
reservoir model and therefore have no wells or subordinate groups
associated with them, they are nevertheless connected to other higher
level groups and higher level groups within a network model (if
activated). Thus, they provide a means to ``add-in'' outside injection
and production to the model without modeling the ``add-in'' reservoirs.

The keyword is used to define injection rates into the model from other
sources (fields, reservoirs etc.) that are not defined in the current
run. For example, if Fields A, B and C are all producing through a
common plant, but only Field A is being modeled in the current input
deck, then injection and production from Fields B and C can be
incorporated into the model in order to meet the plant injection rates
for Field A. Note in this case the import rates from Fields B and C are
fixed, and therefore Field A would act like a ``swing'' injector to
match the overall injection requirement.

\begin{longtable}[]{@{}
  >{\raggedright\arraybackslash}p{(\linewidth - 6\tabcolsep) * \real{0.2500}}
  >{\raggedright\arraybackslash}p{(\linewidth - 6\tabcolsep) * \real{0.2500}}
  >{\raggedright\arraybackslash}p{(\linewidth - 6\tabcolsep) * \real{0.2500}}
  >{\raggedright\arraybackslash}p{(\linewidth - 6\tabcolsep) * \real{0.2500}}@{}}
\toprule\noalign{}
\begin{minipage}[b]{\linewidth}\raggedright
No.
\end{minipage} & \begin{minipage}[b]{\linewidth}\raggedright
Name
\end{minipage} & \begin{minipage}[b]{\linewidth}\raggedright
Description
\end{minipage} & \begin{minipage}[b]{\linewidth}\raggedright
Default
\end{minipage} \\
\midrule\noalign{}
\endhead
\bottomrule\noalign{}
\endlastfoot
Field & Metric & Laboratory & \\
1 & GRPNAME & A character string of up to eight characters in length
that defines the satellite group name for which the group injection
rates are being defined. The group named FIELD is the top most group and
should not be used to define a satellite group. Note that the group
hierarchy should be defined by the GRUPTREE keyword in the SCHEDULE
section, when there is more than one level of groups, otherwise all the
groups will sit directly under the FIELD group in the group tree
hierarchy. Note that a satellite group cannot have subordinate groups or
wells. & None \\
2 & TYPE & A defined character string that specifies the type of
injection fluid. TYPE should be set to one of the following character
strings: If the satellite group injects more than one phase then the
rates should be specified in separate records. & None \\
3 & RATE & A positive real value that defines the surface injection rate
for the phase declared by the TYPE variable. & 0.0 \\
Liquid: stb/d Gas: Mscf/d & Liquid: sm3/day Gas: sm3/day & Liquid:
scc/hour Gas: scc/hour & \\
4 & RESV & A positive real value that defines the reservoir volume
injection rate for the phase declared by the TYPE variable for the
satellite group. Generally, RESV should be set to zero, unless one
wishes to have the satellite reservoir volume rate added to the group's
superordinate groups. If RESV is set to zero for all satellite groups
and if the superordinate group has a reservoir volume injection target,
a voidage replacement target, or a production balancing fraction target,
then only the non-satellite group volumes will be used. This item is not
supported by OPM Flow and should be defaulted (1*) or set to zero. &
0.0 \\
rb/d & rm3/day & rcc/hour & \\
5 & CALRATE & A positive real value that defines the calorific injection
rate used in the commercial compositional simulator. This item is not
supported by OPM Flow and should be defaulted (1*) or set to zero. &
0.0 \\
Notes: & & & \\
\end{longtable}

\emph{Table 12.45: GSATINJE Keyword Description}

See also the GSATPROD and GRUPTREE keywords to define satellite
production rates and the group hierarchy, respectively. For
non-satellite groups see the GCONINJE and GCONPROD keywords. All the
aforementioned keywords are in the SCHEDULE section.

\begin{tcolorbox}[enhanced jigsaw, title=\textcolor{quarto-callout-note-color}{\faInfo}\hspace{0.5em}{Note}, coltitle=black, breakable, toprule=.15mm, colbacktitle=quarto-callout-note-color!10!white, leftrule=.75mm, opacitybacktitle=0.6, colback=white, bottomtitle=1mm, toptitle=1mm, titlerule=0mm, left=2mm, colframe=quarto-callout-note-color-frame, arc=.35mm, rightrule=.15mm, bottomrule=.15mm, opacityback=0]

Once a group has been defined to be a satellite group, via the GSATINJE
and GSATPROD keywords, then the equivalent modeled group keywords,
GCONINJE and GCONPROD in the SCHEDULE section, cannot be used to set the
operating conditions for satellite groups, only the GSATINJE and
GSATPROD keywords may be used.

\end{tcolorbox}

\subsubsection{Example}\label{example-536}

The example below is based on an oil field, with the main field FLD-A
being modeled in the input deck with two groups (FLD-A1 and FLD-A2),
combined with one satellite field supplementing oil production and water
injection (FLD-B).

\begin{verbatim}
-- ==============================================================================
--
-- SCHEDULE SECTION
--
-- ==============================================================================
SCHEDULE
--
--       DEFINE GROUP TREE HIERARCHY
--       LOWER     HIGHER
--       GROUP     GROUP
GRUPTREE
         FLD-A     FIELD                                                       /
         FLD-A1    FLD-A                                                       /
         FLD-A2    FLD-A                                                       /
         FLD-B     FIELD                                                       /
/
--
--       GROUP PRODUCTION CONTROLS
--
-- GRUP  CNTL  OIL    WAT    GAS    LIQ    CNTL  GRUP  GUIDE  GUIDE  CNTL
-- NAME  MODE  RATE   RATE   RATE   RATE   OPT   CNTL  RATE   DEF    WAT
GCONPROD
FIELD    ORAT  20E3   30E3    50E3  1*     1*    1*                            /
FLD-A    FLD   20E3   30E3    50E3  1*     1*    1*                            /
/                                                                                --
--       LOAD INCLUDE FILE - LOAD ALL WELLS AND VFP TABLES
--
INCLUDE
         'FLD-A-P50-WELLS.INC'  /
--
--       GROUP INJECTION TARGETS AND CONSTRAINTS
--
-- GRUP  FLUID CNTL   SURF   RESV   REINJ  VOID  GRUP  GUIDE  GUIDE GRUP  GRUP
-- NAME  TYPE  MODE   RATE   RATE   FRAC   FRAC  CNTL  RATE   DEF   REINJ RESV
GCONINJE
FIELD    WAT   REIN   35E3   1*     1*     1.0    NO   1*     1*    1*    1*   /
/                                                                               --
--       SATELLITE GROUP PRODUCTION CONTROLS
--
-- GRUP  OIL    WAT    GAS    RESV   GAS   CALORIFIC
-- NAME  RATE   RATE   RATE   RATE   LIFT  RATE
GSATPROD
FLD-B    5E3    0.00   2E3    1*     1*                    /
/
--
--       SATELLITE GROUP INJECTION CONTROLS
--
-- GRUP  FLUID  SURF   RESV   CALORIFIC
-- NAME  TYPE   RATE   RATE   RATEV
GSATINJE
FLD-B    WAT    10E3   1*     1*                           /
/
--
--       ADVANCE SIMULATION BY REPORTING DATE
--
DATES
         1  FEB   2021  /
/
--
--       SATELLITE GROUP PRODUCTION CONTROLS
--
-- GRUP  OIL    WAT    GAS    RESV   GAS   CALORIFIC
-- NAME  RATE   RATE   RATE   RATE   LIFT  RATE
GSATPROD
FLD-B    5E3    0.00   2E3    0.0    1*                    /
/
--
--       SATELLITE GROUP INJECTION CONTROLS
--
-- GRUP  FLUID  SURF   RESV   CALORIFIC
-- NAME  TYPE   RATE   RATE   RATEV
GSATINJE
FLD-B    WAT    10E3   1*     1*                           /
/
\end{verbatim}

Here FLD-B will supplement FLD-A's water re-injection rate by 10,000
stb/d, subject to a maximum field water injection rate of 35,000 stb/d
on the GCONINJE keyword. This means that FLD-A's maximum water injection
rate cannot exceed 25,000 stb/d.~In terms of oil rates, GCONPROD defines
a maximum oil rate of 20,000 stb/d of which 5,000 stb/d is from FLD-B.
If any of the constraints on the GCONINJE and GCONPROD keywords are
violated, then the appropriate action will occur only for FLD-A, in
order to ensure the group constraints are honored.

\subsection{GSATPROD -- Define Group Satellite Production
Rates}\label{gsatprod-define-group-satellite-production-rates}

\begin{longtable}[]{@{}
  >{\raggedright\arraybackslash}p{(\linewidth - 14\tabcolsep) * \real{0.1250}}
  >{\raggedright\arraybackslash}p{(\linewidth - 14\tabcolsep) * \real{0.1250}}
  >{\raggedright\arraybackslash}p{(\linewidth - 14\tabcolsep) * \real{0.1250}}
  >{\raggedright\arraybackslash}p{(\linewidth - 14\tabcolsep) * \real{0.1250}}
  >{\raggedright\arraybackslash}p{(\linewidth - 14\tabcolsep) * \real{0.1250}}
  >{\raggedright\arraybackslash}p{(\linewidth - 14\tabcolsep) * \real{0.1250}}
  >{\raggedright\arraybackslash}p{(\linewidth - 14\tabcolsep) * \real{0.1250}}
  >{\raggedright\arraybackslash}p{(\linewidth - 14\tabcolsep) * \real{0.1250}}@{}}
\toprule\noalign{}
\begin{minipage}[b]{\linewidth}\raggedright
RUNSPEC
\end{minipage} & \begin{minipage}[b]{\linewidth}\raggedright
GRID
\end{minipage} & \begin{minipage}[b]{\linewidth}\raggedright
EDIT
\end{minipage} & \begin{minipage}[b]{\linewidth}\raggedright
PROPS
\end{minipage} & \begin{minipage}[b]{\linewidth}\raggedright
REGIONS
\end{minipage} & \begin{minipage}[b]{\linewidth}\raggedright
SOLUTION
\end{minipage} & \begin{minipage}[b]{\linewidth}\raggedright
SUMMARY
\end{minipage} & \begin{minipage}[b]{\linewidth}\raggedright
SCHEDULE
\end{minipage} \\
\midrule\noalign{}
\endhead
\bottomrule\noalign{}
\endlastfoot
\end{longtable}

\subsubsection{Description}\label{description-958}

The GSATPROD keyword defines a satellite group's oil, water and gas
production rates in the model. Satellite groups are not connected to the
reservoir model and therefore have no wells or subordinate groups
associated with them, they are nevertheless connected to other higher
level groups and higher level groups within a network model (if
activated). They provide a means to ``add-in'' outside injection and
production to the model without modeling the ``add-in'' reservoirs.

This keyword is used to import additional fluid streams into the model
from other sources (fields, reservoirs, etc.) that are not defined in
the current run. For example, if Fields A, B and C are supplying gas to
a power plant, but only Field A is being modeled in the current input
deck, then production from Fields B and C can be incorporated into model
in order to meet the plant demand. Note in this case the import gas
rates from Fields B and C are fixed, and therefore Field A acts like a
``swing'' producer to match the gas demand target.

\begin{longtable}[]{@{}
  >{\raggedright\arraybackslash}p{(\linewidth - 6\tabcolsep) * \real{0.2500}}
  >{\raggedright\arraybackslash}p{(\linewidth - 6\tabcolsep) * \real{0.2500}}
  >{\raggedright\arraybackslash}p{(\linewidth - 6\tabcolsep) * \real{0.2500}}
  >{\raggedright\arraybackslash}p{(\linewidth - 6\tabcolsep) * \real{0.2500}}@{}}
\toprule\noalign{}
\begin{minipage}[b]{\linewidth}\raggedright
No.
\end{minipage} & \begin{minipage}[b]{\linewidth}\raggedright
Name
\end{minipage} & \begin{minipage}[b]{\linewidth}\raggedright
Description
\end{minipage} & \begin{minipage}[b]{\linewidth}\raggedright
Default
\end{minipage} \\
\midrule\noalign{}
\endhead
\bottomrule\noalign{}
\endlastfoot
Field & Metric & Laboratory & \\
1 & GRPNAME & A character string of up to eight characters in length
that defines the satellite group name for which the group production
import rates are being defined. The group named FIELD is the top most
group and should not be used with this keyword. Note that the group
hierarchy should be defined by the GRUPTREE keyword in the SCHEDULE
section when there is more than one level of groups, otherwise all the
groups will sit directly under the FIELD group in the group tree
hierarchy. Note that a satellite group cannot have subordinate groups or
wells. & None \\
2 & ORAT & A real value, greater than or equal to zero, that defines the
satellite's surface oil production rate to be imported into the model.
This value may be specified using a User Defined Argument (UDA). &
0.0 \\
stb/d & sm3/day & scc/hour & \\
3 & WRAT & A real value, greater than or equal to zero, that defines the
satellite's surface water production rate to be imported into the model.
This value may be specified using a User Defined Argument (UDA). &
0.0 \\
stb/d & sm3/day & scc/hour & \\
4 & GRAT & A real value, greater than or equal to zero, that defines the
satellite's surface gas production rate to be imported into the model.
This value may be specified using a User Defined Argument (UDA). &
0.0 \\
Mscf/d & sm3/day & scc/hour & \\
5 & RESV & A real value, greater than or equal to zero, that defines the
satellite's reservoir volume production rate to be imported into the
model. This item is not supported by OPM Flow and should be defaulted
(1*) or set to zero. & 0.0 \\
rb/d & rm3/day & rcc/hour & \\
6 & GASLIFT & A real value, greater than or equal to zero, that defines
the satellite's gas lift gas to be imported for this group. A value for
GASLIFT is only required if externally supplied gas lift is being
imported into the model, and: This item is not supported by OPM Flow and
should be defaulted (1*) or set to zero. & 0.0 \\
Mscf/d & sm3/day & scc/hour & \\
7 & CALRATE & A real value greater than or equal to zero that defines
the satellite's calorific production rate. This item is not supported by
OPM Flow and should be defaulted (1*) or set to zero. & 0.0 \\
Notes: & & & \\
\end{longtable}

\emph{Table 12.46: GSATPROD Keyword Description}

See also the GSATINJE and GRUPTREE keywords to define satellite
injection rates and the group hierarchy, respectively. For non-satellite
groups see the GCONINJE and GCONPROD keywords. All the aforementioned
keywords are in the SCHEDULE section.

\begin{tcolorbox}[enhanced jigsaw, title=\textcolor{quarto-callout-note-color}{\faInfo}\hspace{0.5em}{Note}, coltitle=black, breakable, toprule=.15mm, colbacktitle=quarto-callout-note-color!10!white, leftrule=.75mm, opacitybacktitle=0.6, colback=white, bottomtitle=1mm, toptitle=1mm, titlerule=0mm, left=2mm, colframe=quarto-callout-note-color-frame, arc=.35mm, rightrule=.15mm, bottomrule=.15mm, opacityback=0]

Once a group has been defined to be a satellite group, via the GSATINJE
and GSATPROD keywords, then the equivalent modeled group keywords,
GCONINJE and GCONPROD in the SCHEDULE section, cannot be used to set the
operating conditions for satellite groups, only the GSATINJE and
GSATPROD keywords may be used.

\end{tcolorbox}

\subsubsection{Example}\label{example-537}

The example below is based on a dry gas field, with the main field FLD-A
being modeled in the input deck with two groups (FLD-A1 and FLD-A2),
combined with two satellite gas fields, FLD-B and FLD-C.

\begin{verbatim}
-- ==============================================================================
--
-- SCHEDULE SECTION
--
-- ==============================================================================
SCHEDULE
--
--       DEFINE GROUP TREE HIERARCHY
--       LOWER     HIGHER
--       GROUP     GROUP
GRUPTREE
         FLD-A     FIELD                                                       /
         FLD-A1    FLD-A                                                       /
         FLD-A2    FLD-A                                                       /
         FLD-B     FIELD                                                       /
         FLD-C     FIELD                                                       /
/
--
--       GROUP PRODUCTION CONTROLS
--
-- GRUP  CNTL  OIL    WAT    GAS    LIQ    CNTL  GRUP  GUIDE  GUIDE  CNTL
-- NAME  MODE  RATE   RATE   RATE   RATE   OPT   CNTL  RATE   DEF    WAT
GCONPROD
FIELD    GRAT  1*     1*     450E3  1*     1*    1*                            /
FLD-A    FLD   1*     1*     450E3  1*     1*    1*                            /
/
--
--       LOAD INCLUDE FILE - LOAD ALL WELLS AND VFP TABLES
--
INCLUDE
         'FLD-A-P50-WELLS.INC'  /
--
--       SATELLITE GROUP PRODUCTION CONTROLS
--
-- GRUP  OIL    WAT    GAS    RESV   GAS   CALORIFIC
-- NAME  RATE   RATE   RATE   RATE   LIFT  RATE
GSATPROD
FLD-B    10.0   0.00   30E3   1*     1*                    /
FLD-C    20.0   5.50   20E3   0.0    1*                    /
/
--
--       ADVANCE SIMULATION BY REPORTING DATE
--
DATES
         1  FEB   2021  /
/
--
--       SATELLITE GROUP PRODUCTION CONTROLS
--
-- GRUP  OIL    WAT    GAS    RESV   GAS   CALORIFIC
-- NAME  RATE   RATE   RATE   RATE   LIFT  RATE
GSATPROD
FLD-B    10.0   0.00   30E3   0.0    1*                    /
FLD-C    20.0   9.00   20E3   0.0    1*                    /
/
\end{verbatim}

Since the field gas rate is set to 450 MMscf/d and the satellite
production is 50 MMscf/d, then FLD-A will produce only 400 MMscf/d and
not the stipulated 450 MMscf/d on the GCONPROD keyword.

\subsection{GSEPCOND -- Assign Group
Separators}\label{gsepcond-assign-group-separators}

\begin{longtable}[]{@{}
  >{\raggedright\arraybackslash}p{(\linewidth - 14\tabcolsep) * \real{0.1250}}
  >{\raggedright\arraybackslash}p{(\linewidth - 14\tabcolsep) * \real{0.1250}}
  >{\raggedright\arraybackslash}p{(\linewidth - 14\tabcolsep) * \real{0.1250}}
  >{\raggedright\arraybackslash}p{(\linewidth - 14\tabcolsep) * \real{0.1250}}
  >{\raggedright\arraybackslash}p{(\linewidth - 14\tabcolsep) * \real{0.1250}}
  >{\raggedright\arraybackslash}p{(\linewidth - 14\tabcolsep) * \real{0.1250}}
  >{\raggedright\arraybackslash}p{(\linewidth - 14\tabcolsep) * \real{0.1250}}
  >{\raggedright\arraybackslash}p{(\linewidth - 14\tabcolsep) * \real{0.1250}}@{}}
\toprule\noalign{}
\begin{minipage}[b]{\linewidth}\raggedright
RUNSPEC
\end{minipage} & \begin{minipage}[b]{\linewidth}\raggedright
GRID
\end{minipage} & \begin{minipage}[b]{\linewidth}\raggedright
EDIT
\end{minipage} & \begin{minipage}[b]{\linewidth}\raggedright
PROPS
\end{minipage} & \begin{minipage}[b]{\linewidth}\raggedright
REGIONS
\end{minipage} & \begin{minipage}[b]{\linewidth}\raggedright
SOLUTION
\end{minipage} & \begin{minipage}[b]{\linewidth}\raggedright
SUMMARY
\end{minipage} & \begin{minipage}[b]{\linewidth}\raggedright
SCHEDULE
\end{minipage} \\
\midrule\noalign{}
\endhead
\bottomrule\noalign{}
\endlastfoot
\end{longtable}

\subsubsection{Description}\label{description-959}

The GSEPCOND keyword assigns previously defined separators to a group.
Group separators are specified by the SEPVALS keyword in the SCHEDULE
section. The facility is used in black-oil modeling to re-scale the PVT
data entered via the PROPS section, based on the saturation point oil
formation volume factor (Bob) and the initial saturated gas-oil ratio
(Rsi) entered on the SEVPALS keyword.

This keyword is not supported by OPM Flow but would change the results
if supported so the simulation will be stopped.

\subsection{GSSCPTST -- Perform Sustainable Capacity
Test}\label{gsscptst-perform-sustainable-capacity-test}

\begin{longtable}[]{@{}
  >{\raggedright\arraybackslash}p{(\linewidth - 14\tabcolsep) * \real{0.1250}}
  >{\raggedright\arraybackslash}p{(\linewidth - 14\tabcolsep) * \real{0.1250}}
  >{\raggedright\arraybackslash}p{(\linewidth - 14\tabcolsep) * \real{0.1250}}
  >{\raggedright\arraybackslash}p{(\linewidth - 14\tabcolsep) * \real{0.1250}}
  >{\raggedright\arraybackslash}p{(\linewidth - 14\tabcolsep) * \real{0.1250}}
  >{\raggedright\arraybackslash}p{(\linewidth - 14\tabcolsep) * \real{0.1250}}
  >{\raggedright\arraybackslash}p{(\linewidth - 14\tabcolsep) * \real{0.1250}}
  >{\raggedright\arraybackslash}p{(\linewidth - 14\tabcolsep) * \real{0.1250}}@{}}
\toprule\noalign{}
\begin{minipage}[b]{\linewidth}\raggedright
RUNSPEC
\end{minipage} & \begin{minipage}[b]{\linewidth}\raggedright
GRID
\end{minipage} & \begin{minipage}[b]{\linewidth}\raggedright
EDIT
\end{minipage} & \begin{minipage}[b]{\linewidth}\raggedright
PROPS
\end{minipage} & \begin{minipage}[b]{\linewidth}\raggedright
REGIONS
\end{minipage} & \begin{minipage}[b]{\linewidth}\raggedright
SOLUTION
\end{minipage} & \begin{minipage}[b]{\linewidth}\raggedright
SUMMARY
\end{minipage} & \begin{minipage}[b]{\linewidth}\raggedright
SCHEDULE
\end{minipage} \\
\midrule\noalign{}
\endhead
\bottomrule\noalign{}
\endlastfoot
\end{longtable}

\subsubsection{Description}\label{description-960}

The GSSCPTST keyword instructs the simulator to perform a sustainable
capacity test. This causes the model to be saved in it's current state
via the RESTART file, and the test performed by running the simulation
under the current conditions combine with the parameters on this
keyword. After the test is perform, the simulator will restart from the
point prior to the test by loading in the RESTART file. This type of
testing is normally applied to gas fields for which the gas sales
contracts stipulate that the gas sales rates are based on a sustainable
capacity rate over a fixed period of time.

This keyword is not supported by OPM Flow but would change the results
if supported so the simulation will be stopped.

\subsection{GSWINGF -- Define Group Multiple Gas Contract
Parameters}\label{gswingf-define-group-multiple-gas-contract-parameters}

\begin{longtable}[]{@{}
  >{\raggedright\arraybackslash}p{(\linewidth - 14\tabcolsep) * \real{0.1250}}
  >{\raggedright\arraybackslash}p{(\linewidth - 14\tabcolsep) * \real{0.1250}}
  >{\raggedright\arraybackslash}p{(\linewidth - 14\tabcolsep) * \real{0.1250}}
  >{\raggedright\arraybackslash}p{(\linewidth - 14\tabcolsep) * \real{0.1250}}
  >{\raggedright\arraybackslash}p{(\linewidth - 14\tabcolsep) * \real{0.1250}}
  >{\raggedright\arraybackslash}p{(\linewidth - 14\tabcolsep) * \real{0.1250}}
  >{\raggedright\arraybackslash}p{(\linewidth - 14\tabcolsep) * \real{0.1250}}
  >{\raggedright\arraybackslash}p{(\linewidth - 14\tabcolsep) * \real{0.1250}}@{}}
\toprule\noalign{}
\begin{minipage}[b]{\linewidth}\raggedright
RUNSPEC
\end{minipage} & \begin{minipage}[b]{\linewidth}\raggedright
GRID
\end{minipage} & \begin{minipage}[b]{\linewidth}\raggedright
EDIT
\end{minipage} & \begin{minipage}[b]{\linewidth}\raggedright
PROPS
\end{minipage} & \begin{minipage}[b]{\linewidth}\raggedright
REGIONS
\end{minipage} & \begin{minipage}[b]{\linewidth}\raggedright
SOLUTION
\end{minipage} & \begin{minipage}[b]{\linewidth}\raggedright
SUMMARY
\end{minipage} & \begin{minipage}[b]{\linewidth}\raggedright
SCHEDULE
\end{minipage} \\
\midrule\noalign{}
\endhead
\bottomrule\noalign{}
\endlastfoot
\end{longtable}

\subsubsection{Description}\label{description-961}

This keyword, GSWINGF, defines the gas contract parameters, swing factor
and the monthly seasonal profile factor, for when there are multiple gas
contracts being used in the model. The keyword is used with the Gas
Field Operations option which is activated by the GASFIELD keyword in
the RUNSPEC section. Gas contracts are commonly based on a Daily
Contract Quantity (``DCQ'') that determines the gas rate that the field
should be produced at, which is normally expressed as a multiple of the
DCQ, for example 1.33, and is often referred to as the ``swing factor''.
Some gas contracts also define a maximum DCQ (``Max DCQ'') and/or a
minimum take or pay DCQ (``Min DCQ''), as well as seasonal demand
characteristics. For example, gas rates may be set higher in the winter
months in order to meet heating demand compared with summer months in
colder climates, and the opposite in warmer climates where air
conditioning demand is high.

Thus, the DCQ must be calculated first over a contract period, where the
contract period is commonly contract years, or in some instances
contract quarters. This is performed by the simulator using the current
DCQ and checking to see if the (DCQ x Swing Factor) can be satisfied
throughout the current contract period, if not the DCQ is re-calculated
so that the (DCQ x Swing Factor) condition is satisfied. Once this
condition is met, the second and final pass uses the calculated DCQ in
conjunction with the monthly scaling profile data to set the monthly gas
rate for the field:

\begin{equation}\phantomsection\label{eq-12-27}{
{Q}_{\mathit{month}} = \mathit{DCQ}\times {\mathit{SWINGFAC}}_{\mathit{month}}
}\end{equation}

Where:

Qmonth = the monthly gas production target

DCQ = Daily Contract Quantity

SWINGFACmonth = monthly rate scaling factor that takes into account
seasonal demand, etc.

Here the GSWINGF keyword allows for different gas contract parameters to
be assign to different groups and is mutually exclusive to the SWINGFAC
keyword in the SCHEDULE section, that sets the gas contract parameters
for a single contract at the FIELD group level.

This keyword is not supported by OPM Flow but would change the results
if supported so the simulation will be stopped.

\subsection{GTADD -- Add a Constant to a Group Target or
Constraint}\label{gtadd-add-a-constant-to-a-group-target-or-constraint}

\begin{longtable}[]{@{}
  >{\raggedright\arraybackslash}p{(\linewidth - 14\tabcolsep) * \real{0.1250}}
  >{\raggedright\arraybackslash}p{(\linewidth - 14\tabcolsep) * \real{0.1250}}
  >{\raggedright\arraybackslash}p{(\linewidth - 14\tabcolsep) * \real{0.1250}}
  >{\raggedright\arraybackslash}p{(\linewidth - 14\tabcolsep) * \real{0.1250}}
  >{\raggedright\arraybackslash}p{(\linewidth - 14\tabcolsep) * \real{0.1250}}
  >{\raggedright\arraybackslash}p{(\linewidth - 14\tabcolsep) * \real{0.1250}}
  >{\raggedright\arraybackslash}p{(\linewidth - 14\tabcolsep) * \real{0.1250}}
  >{\raggedright\arraybackslash}p{(\linewidth - 14\tabcolsep) * \real{0.1250}}@{}}
\toprule\noalign{}
\begin{minipage}[b]{\linewidth}\raggedright
RUNSPEC
\end{minipage} & \begin{minipage}[b]{\linewidth}\raggedright
GRID
\end{minipage} & \begin{minipage}[b]{\linewidth}\raggedright
EDIT
\end{minipage} & \begin{minipage}[b]{\linewidth}\raggedright
PROPS
\end{minipage} & \begin{minipage}[b]{\linewidth}\raggedright
REGIONS
\end{minipage} & \begin{minipage}[b]{\linewidth}\raggedright
SOLUTION
\end{minipage} & \begin{minipage}[b]{\linewidth}\raggedright
SUMMARY
\end{minipage} & \begin{minipage}[b]{\linewidth}\raggedright
SCHEDULE
\end{minipage} \\
\midrule\noalign{}
\endhead
\bottomrule\noalign{}
\endlastfoot
\end{longtable}

\subsubsection{Description}\label{description-962}

This keyword, GTADD, adds a numerical constant to a group's target or
constraint value. The group must have been initially fully defined using
the GCONPROD or GCONPRI keywords for producers or GCONINJE for
injectors. Variables not changed by the GTADD keyword remain the same as
those previously entered via the group control keywords or previously
entered GTADD keywords. See also the GRUPTARG keyword that sets the
values for a group's target and constraints and the GTMULT keyword that
multiplies a group target or constraint by a constant. All the
aforementioned keywords are described in the SCHEDULE section.

This keyword is not supported by OPM Flow but would change the results
if supported so the simulation will be stopped.

\subsection{GTMULT -- Multiply Group Target or Constraint by a
Constant}\label{gtmult-multiply-group-target-or-constraint-by-a-constant}

\begin{longtable}[]{@{}
  >{\raggedright\arraybackslash}p{(\linewidth - 14\tabcolsep) * \real{0.1250}}
  >{\raggedright\arraybackslash}p{(\linewidth - 14\tabcolsep) * \real{0.1250}}
  >{\raggedright\arraybackslash}p{(\linewidth - 14\tabcolsep) * \real{0.1250}}
  >{\raggedright\arraybackslash}p{(\linewidth - 14\tabcolsep) * \real{0.1250}}
  >{\raggedright\arraybackslash}p{(\linewidth - 14\tabcolsep) * \real{0.1250}}
  >{\raggedright\arraybackslash}p{(\linewidth - 14\tabcolsep) * \real{0.1250}}
  >{\raggedright\arraybackslash}p{(\linewidth - 14\tabcolsep) * \real{0.1250}}
  >{\raggedright\arraybackslash}p{(\linewidth - 14\tabcolsep) * \real{0.1250}}@{}}
\toprule\noalign{}
\begin{minipage}[b]{\linewidth}\raggedright
RUNSPEC
\end{minipage} & \begin{minipage}[b]{\linewidth}\raggedright
GRID
\end{minipage} & \begin{minipage}[b]{\linewidth}\raggedright
EDIT
\end{minipage} & \begin{minipage}[b]{\linewidth}\raggedright
PROPS
\end{minipage} & \begin{minipage}[b]{\linewidth}\raggedright
REGIONS
\end{minipage} & \begin{minipage}[b]{\linewidth}\raggedright
SOLUTION
\end{minipage} & \begin{minipage}[b]{\linewidth}\raggedright
SUMMARY
\end{minipage} & \begin{minipage}[b]{\linewidth}\raggedright
SCHEDULE
\end{minipage} \\
\midrule\noalign{}
\endhead
\bottomrule\noalign{}
\endlastfoot
\end{longtable}

\subsubsection{Description}\label{description-963}

This keyword, GTMULT, multiplies a group's target or constraint value by
a numerical constant. The group must have been initially fully defined
using the GCONPROD or GCONPRI keywords for producers or GCONINJE for
injectors. Variables not changed by the GTMULT keyword remain the same
as those previously entered via the group control keywords or previously
entered GTMULT keywords. See also the GRUPTARG keyword that sets the
values for a group's target and constraints, and the GTADD keyword that
adds a constant to a group's target or constraint. All the
aforementioned keywords are described in the SCHEDULE section.

This keyword is not supported by OPM Flow but would change the results
if supported so the simulation will be stopped.

\subsection{GUIDECAL -- Scale Guide Rates Based on Gas Calorific
Value}\label{guidecal-scale-guide-rates-based-on-gas-calorific-value}

\begin{longtable}[]{@{}
  >{\raggedright\arraybackslash}p{(\linewidth - 14\tabcolsep) * \real{0.1250}}
  >{\raggedright\arraybackslash}p{(\linewidth - 14\tabcolsep) * \real{0.1250}}
  >{\raggedright\arraybackslash}p{(\linewidth - 14\tabcolsep) * \real{0.1250}}
  >{\raggedright\arraybackslash}p{(\linewidth - 14\tabcolsep) * \real{0.1250}}
  >{\raggedright\arraybackslash}p{(\linewidth - 14\tabcolsep) * \real{0.1250}}
  >{\raggedright\arraybackslash}p{(\linewidth - 14\tabcolsep) * \real{0.1250}}
  >{\raggedright\arraybackslash}p{(\linewidth - 14\tabcolsep) * \real{0.1250}}
  >{\raggedright\arraybackslash}p{(\linewidth - 14\tabcolsep) * \real{0.1250}}@{}}
\toprule\noalign{}
\begin{minipage}[b]{\linewidth}\raggedright
RUNSPEC
\end{minipage} & \begin{minipage}[b]{\linewidth}\raggedright
GRID
\end{minipage} & \begin{minipage}[b]{\linewidth}\raggedright
EDIT
\end{minipage} & \begin{minipage}[b]{\linewidth}\raggedright
PROPS
\end{minipage} & \begin{minipage}[b]{\linewidth}\raggedright
REGIONS
\end{minipage} & \begin{minipage}[b]{\linewidth}\raggedright
SOLUTION
\end{minipage} & \begin{minipage}[b]{\linewidth}\raggedright
SUMMARY
\end{minipage} & \begin{minipage}[b]{\linewidth}\raggedright
SCHEDULE
\end{minipage} \\
\midrule\noalign{}
\endhead
\bottomrule\noalign{}
\endlastfoot
\end{longtable}

\subsubsection{Description}\label{description-964}

The GUIDECAL keyword defines a well or group's guide rate as a function
of their calorific values, for when the individual wells and groups are
under guide rate control. Group and well guide rates that have not been
directly defined are set equal to their production potentials at the
start of each time step. In this case the GUIDECAL keyword can be used
to specify the coefficients of a function that takes into account the
calorific value of the produced gas, effectively scaling the guide rates
based on the calorific value of the gas being produced.

This keyword is not supported by OPM Flow but would change the results
if supported so the simulation will be stopped.

\subsection{GUIDERAT -- Define Group Guide Rate
Formula}\label{guiderat-define-group-guide-rate-formula}

\begin{longtable}[]{@{}
  >{\raggedright\arraybackslash}p{(\linewidth - 14\tabcolsep) * \real{0.1250}}
  >{\raggedright\arraybackslash}p{(\linewidth - 14\tabcolsep) * \real{0.1250}}
  >{\raggedright\arraybackslash}p{(\linewidth - 14\tabcolsep) * \real{0.1250}}
  >{\raggedright\arraybackslash}p{(\linewidth - 14\tabcolsep) * \real{0.1250}}
  >{\raggedright\arraybackslash}p{(\linewidth - 14\tabcolsep) * \real{0.1250}}
  >{\raggedright\arraybackslash}p{(\linewidth - 14\tabcolsep) * \real{0.1250}}
  >{\raggedright\arraybackslash}p{(\linewidth - 14\tabcolsep) * \real{0.1250}}
  >{\raggedright\arraybackslash}p{(\linewidth - 14\tabcolsep) * \real{0.1250}}@{}}
\toprule\noalign{}
\begin{minipage}[b]{\linewidth}\raggedright
RUNSPEC
\end{minipage} & \begin{minipage}[b]{\linewidth}\raggedright
GRID
\end{minipage} & \begin{minipage}[b]{\linewidth}\raggedright
EDIT
\end{minipage} & \begin{minipage}[b]{\linewidth}\raggedright
PROPS
\end{minipage} & \begin{minipage}[b]{\linewidth}\raggedright
REGIONS
\end{minipage} & \begin{minipage}[b]{\linewidth}\raggedright
SOLUTION
\end{minipage} & \begin{minipage}[b]{\linewidth}\raggedright
SUMMARY
\end{minipage} & \begin{minipage}[b]{\linewidth}\raggedright
SCHEDULE
\end{minipage} \\
\midrule\noalign{}
\endhead
\bottomrule\noalign{}
\endlastfoot
\end{longtable}

\subsubsection{Description}\label{description-965}

This keyword defines a general formulae used to define a group's and
well's guide rate as a function of the their potential. The default
behavior, that is when this keyword is not invoked, is to set the target
control mode and rate via the GCONPROD keyword in the SCHEDULE section.
In this case the target rate is distributed between the group's wells
that are under group control using a well's guide rate. If a well's
guide rate has not been defined, for example by this keyword, then the
well potential of the group controlling phase at the beginning of the
time step is used. For example, if the group target rate and phase is
oil, then the well's under group control will have their oil rates
determined by their oil rate potential {[}Production and injection
potentials are based on rates that are unrestricted. For wells this
implies that well potential is calculated based on either the BHP or THP
limit, which ever is the more constraining.{]}. The GUIDERAT keyword
substitutes the potential calculation with a more general formulae in
the aforementioned distribution and allocation of the rates:

\begin{equation}\phantomsection\label{eq-12-28}{
\mathit{Phase}\mathit{Guide}\mathit{Rate} =\frac{{({\mathit{Potential}}_{\mathit{Phase}})}^{A}}{B + C{(\mathit{Potential}{\mathit{Ratio}}_{1})}^{D} + E{(\mathit{Potential}{\mathit{Ratio}}_{2})}^{F}}
}\end{equation}

Where:

PotentialPhase = the potential of the phase,

A to F = constants defined on this keyword,

Potential Ratio1 = the potential phase ratio as defined by this keyword,

Potential Ratio2 = the potential phase ratio as defined by this keyword.

The formulae can be used to control high water cut or high GOR wells in
an oil field, such as the offending wells are given progressively
smaller guide rates as they water out or gas out.

Note that groups can only have potential guide rates if they are
subordinate in another group and required to produce a proportion of the
superior group's target rate. In this case the GUIDERAT keyword can
optionally be applied by setting GUIPHASE variable on the GCONPROD
keyword in the SCHEDULE section. Group potentials are the sum of the
potentials of their subordinate open wells.

\begin{longtable}[]{@{}
  >{\raggedright\arraybackslash}p{(\linewidth - 6\tabcolsep) * \real{0.2500}}
  >{\raggedright\arraybackslash}p{(\linewidth - 6\tabcolsep) * \real{0.2500}}
  >{\raggedright\arraybackslash}p{(\linewidth - 6\tabcolsep) * \real{0.2500}}
  >{\raggedright\arraybackslash}p{(\linewidth - 6\tabcolsep) * \real{0.2500}}@{}}
\toprule\noalign{}
\begin{minipage}[b]{\linewidth}\raggedright
No.
\end{minipage} & \begin{minipage}[b]{\linewidth}\raggedright
Name
\end{minipage} & \begin{minipage}[b]{\linewidth}\raggedright
Description
\end{minipage} & \begin{minipage}[b]{\linewidth}\raggedright
Default
\end{minipage} \\
\midrule\noalign{}
\endhead
\bottomrule\noalign{}
\endlastfoot
Field & Metric & Laboratory & \\
1 & TSTEP & A real positive value that defines the minimum time interval
to re-calculate the guide rates. The guide rates are calculated at the
start of a time step and the default value of zero means that the guide
rates are calculated for each time step. A non-zero value for TSTEP
resets the minimum interval, for example setting TSTEP equal to 30 would
mean the guide rates are calculate every 30 days, or to the nearest
associated time step. Calculating guide rates every time step may cause
issues due to the rate dependent behavior, for example gas cusping or
water coning causing the well rates to oscillate. In this case using a
non-zero value of TSTEP may eliminate this oscillating behavior. &
0.0 \\
days & days & hours & \\
2 & PHASE & A defined character string that sets the potential phase
guide rate for the group and well, the resulting Phase Guide Rate in
equation (12.28). PHASE should be set to one of the following character
strings: For reference, the units for the various options is given below
. & None \\
WOR: dimensionless WCT: dimensionless WGR: stb/Mscf & dimensionless
dimensionless dimensionless & dimensionless dimensionless dimensionless
& \\
GOR: Mscf/stb GLR: Mscf/stb OGR: stb/Mscf & dimensionless dimensionless
dimensionless & dimensionless dimensionless dimensionless & \\
3 & A & A real value greater than or equal to -3 and less than or equal
to 3, that defines coefficient A in equation (12.28). & 0.0 \\
4 & B & B is a real positive value that defines coefficient B in
equation (12.28). & 0.0 \\
5 & C & C is a real value that defines coefficient C in equation
(12.28). & 0.0 \\
6 & D & D is a real value greater than or equal to -3 and less than or
equal to 3, that defines coefficient D in equation (12.28). & 0.0 \\
7 & E & E is a real value that defines coefficient E in equation
(12.28). & 0.0 \\
8 & F & F is a real value greater than or equal to -3 and less than or
equal to 3, that defines coefficient F in equation (12.28). & 0.0 \\
9 & GROPT01 & A defined character string that determines if calculated
phase guide rates should be allowed to increase (YES) or not (NO), and
should be set to one of the following: Note only the default value is
currently supported by OPM Flow. & YES \\
10 & GROPT02 & A real positive value greater than or equal to zero and
less than or equal to one that ``dampens'' the calculated phase guide
rate based on the following formula:
\(\begin{matrix}{(\mathit{Phase} \mathit{Guide} \mathit{Rate})}_{t}^{\mathit{new}} = \mathit{GROPT}02\times {(\mathit{Phase} \mathit{Guide} \mathit{Rate})}_{t} +  \\ (1 - \mathit{GROPT}02)\times {(\mathit{Phase} \mathit{Guide} \mathit{Rate})}_{(t-1)}\end{matrix}\)
The option is intended to have a similar effect as the GROPT01 NO option
to reduce oscillations as a result of either the water cut or GOR being
rate dependent. Values approaching one allows the calculated phase guide
rates to change instantaneously with the phase potentials, whereas
values approaching zero dampen the potential guide rates towards the
previously calculated values, thereby reducing the potential for
oscillating behavior. & 1.0 \\
11 & GROPT03 & A defined character string that determines if ``free''
gas potential rates for the Potential Ratio2 variable in equation
(12.28) should be used (YES), or if ``free and associated'' gas should
be used (NO), and should be set to one of the following: & NO \\
12 & GROPT04 & A real positive value that sets the minimum potential
guide rate. If the calculated potential guide is below this value it
will be reset to GROPT04. The option is meant to avoid groups and wells
being ignored due to the calculated potential guide rates being
minuscule. & 1.0 x10-6 \\
Notes: & & & \\
\end{longtable}

\emph{Table 12.47: GUIDERAT Keyword Description}

Note that the GUIDERAT keyword only applies to production groups and
wells. Injection groups and wells are still controlled by their
potential guide rates.

Finally, as mentioned previously, if the GUIDERAT or WGRUPCON keywords
are not present in the input deck then the group and well potential
guide rates will be calculated using the well's potential rates. The
WGRUPCON keyword in the SCHEDULE section can be used to set a constant
potential guide rate for a well.

\begin{tcolorbox}[enhanced jigsaw, title=\textcolor{quarto-callout-note-color}{\faInfo}\hspace{0.5em}{Note}, coltitle=black, breakable, toprule=.15mm, colbacktitle=quarto-callout-note-color!10!white, leftrule=.75mm, opacitybacktitle=0.6, colback=white, bottomtitle=1mm, toptitle=1mm, titlerule=0mm, left=2mm, colframe=quarto-callout-note-color-frame, arc=.35mm, rightrule=.15mm, bottomrule=.15mm, opacityback=0]

GUIDERAT can be used to penalize wells producing excessive water by
utilizing the C and D coefficients in equation (12.28), and to
discriminate against wells that are gassing out by setting the E and F
coefficients. Note that the value range through which Potential Ratio1
and Potential Ratio2 vary is variable. For example, if Potential Ratio1
is water cut, then the value should be between zero and one, whereas for
the water-oil ratio the value can vary between zero and infinity. The
same applies to the units of Potential Ratio2 which are dependent on if
the GOR, GLR or OGR ratio is used in the calculation. One can use the C
and E coefficients to scale these terms to the required relative
magnitudes in the denominator and the D and F powers to influence how
quickly the penalty increases with increasing water and gas fractions.
High positive value for D and F coefficients will make production fall
off rapidly as the water or gas fraction increases, while a negative
values will favor producing these type of wells. Note that the B
coefficient should always be positive to prevent the denominator's going
to zero. Finally, if one wishes each well to produce in proportion to
its potential when the water fraction and gas fraction are equal (the
usual case), then the A coefficient should be set to one.

\end{tcolorbox}

\subsubsection{Examples}\label{examples-79}

The first example sets the guide phase to oil and the resulting Phase
Guide Rate based on oil potential based on setting the A and B
coefficients to to one, that is:

\begin{equation}\phantomsection\label{eq-12-29}{
\begin{matrix}\mathit{Phase}\mathit{Guide}\mathit{Rate} =\frac{{({\mathit{Potental}}_{\mathit{Phase}})}^{A}}{B + C{(\mathit{Potential}{\mathit{Ratio}}_{1})}^{D} + E{(\mathit{Potential}{\mathit{Ratio}}_{2})}^{F}} \\  \\  = \frac{\mathit{Oil}{\mathit{Potential}}^{1.0}}{1.0}\end{matrix}
}\end{equation}

with all the other parameters defaulted, except for the minimum time
interval to re-calculate the guide rates which is set to 30 days.

\begin{verbatim}
--
--       SETS GUIDE RATES FOR GROUPS AND WELLS UNDER GUIDE RATE CONTROL
--
--       TIME  GUIDE  A      B     C     D     E    F     INCR   DAMP   FREE
--       STEP  PHASE  POW    CON   CON   POW   CON  POW   OPTN   OPTN   GAS
GUIDERAT
         30    'OIL'  1.0    1.0    1*   1*    1*   1*    1*     1*     1*      /
\end{verbatim}

The next example sets the Phase Guide Rate to the reservoir fluid volume
rate, with preference given to low GOR wells and with high GOR wells
penalized, based on setting A and B to one, C and D to zero, E equal to
10 and F equal to two, that is:

\begin{equation}\phantomsection\label{eq-12-30}{
\begin{matrix}\mathit{Phase}\mathit{Guide}\mathit{Rate} =\frac{{({\mathit{Potental}}_{\mathit{Phase}})}^{A}}{B + C{(\mathit{Potential}{\mathit{Ratio}}_{1})}^{D} + E{(\mathit{Potential}{\mathit{Ratio}}_{2})}^{F}}  \\  \\  = \frac{\mathit{Reservoir} \mathit{Fluid} \mathit{Volume} {\mathit{Potential}}^{1.0}}{1.0 + 10\times {(\mathit{GOR})}^{2}}\end{matrix}
}\end{equation}

with all the other parameters defaulted.

\begin{verbatim}
--
--       SETS GUIDE RATES FOR GROUPS AND WELLS UNDER GUIDE RATE CONTROL
--
--       TIME  GUIDE  A      B     C     D     E    F     INCR   DAMP   FREE
--       STEP  PHASE  POW    CON   CON   POW   CON  POW   OPTN   OPTN   GAS
GUIDERAT
         1*    'RES'  1.0    1.0    1*   1*    10   2     1*     1*     1*      /
\end{verbatim}

The GUIDERAT keyword is very flexible but can also lead to unexpected
results, thus it is probably useful to perform some manual calculations
outside of the simulator before implementing the selected scheme in the
input deck.

\subsection{GUPFREQ -- Instantaneous Gradient Option Update
Frequency}\label{gupfreq-instantaneous-gradient-option-update-frequency}

\begin{longtable}[]{@{}
  >{\raggedright\arraybackslash}p{(\linewidth - 14\tabcolsep) * \real{0.1250}}
  >{\raggedright\arraybackslash}p{(\linewidth - 14\tabcolsep) * \real{0.1250}}
  >{\raggedright\arraybackslash}p{(\linewidth - 14\tabcolsep) * \real{0.1250}}
  >{\raggedright\arraybackslash}p{(\linewidth - 14\tabcolsep) * \real{0.1250}}
  >{\raggedright\arraybackslash}p{(\linewidth - 14\tabcolsep) * \real{0.1250}}
  >{\raggedright\arraybackslash}p{(\linewidth - 14\tabcolsep) * \real{0.1250}}
  >{\raggedright\arraybackslash}p{(\linewidth - 14\tabcolsep) * \real{0.1250}}
  >{\raggedright\arraybackslash}p{(\linewidth - 14\tabcolsep) * \real{0.1250}}@{}}
\toprule\noalign{}
\begin{minipage}[b]{\linewidth}\raggedright
RUNSPEC
\end{minipage} & \begin{minipage}[b]{\linewidth}\raggedright
GRID
\end{minipage} & \begin{minipage}[b]{\linewidth}\raggedright
EDIT
\end{minipage} & \begin{minipage}[b]{\linewidth}\raggedright
PROPS
\end{minipage} & \begin{minipage}[b]{\linewidth}\raggedright
REGIONS
\end{minipage} & \begin{minipage}[b]{\linewidth}\raggedright
SOLUTION
\end{minipage} & \begin{minipage}[b]{\linewidth}\raggedright
SUMMARY
\end{minipage} & \begin{minipage}[b]{\linewidth}\raggedright
SCHEDULE
\end{minipage} \\
\midrule\noalign{}
\endhead
\bottomrule\noalign{}
\endlastfoot
\end{longtable}

\subsubsection{Description}\label{description-966}

The GUPFREQ keyword sets the update frequency of the Instantaneous
Gradient option for when this option has been activated by the GDIMS
keyword in the RUNSPEC section. The Instantaneous Gradient option
calculates derivatives of solution quantities at the current time step
with respect to variations in the variables at the current time step.
This is different to Gradient option that calculates the derivatives of
solution quantities at the current time step with respect to variations
in the variables at the initial time step, that is a time equal to zero.
Consequently, the Instantaneous Gradient option can be switched on and
off by the GUPFREQ keyword in the SCHEDULE section, whereas the Gradient
option cannot.

This keyword is not supported by OPM Flow but would change the results
if supported so the simulation will be stopped.

\subsection{GWRTWCV -- Instantaneous Gradient Option Well
Variables}\label{gwrtwcv-instantaneous-gradient-option-well-variables}

\begin{longtable}[]{@{}
  >{\raggedright\arraybackslash}p{(\linewidth - 14\tabcolsep) * \real{0.1250}}
  >{\raggedright\arraybackslash}p{(\linewidth - 14\tabcolsep) * \real{0.1250}}
  >{\raggedright\arraybackslash}p{(\linewidth - 14\tabcolsep) * \real{0.1250}}
  >{\raggedright\arraybackslash}p{(\linewidth - 14\tabcolsep) * \real{0.1250}}
  >{\raggedright\arraybackslash}p{(\linewidth - 14\tabcolsep) * \real{0.1250}}
  >{\raggedright\arraybackslash}p{(\linewidth - 14\tabcolsep) * \real{0.1250}}
  >{\raggedright\arraybackslash}p{(\linewidth - 14\tabcolsep) * \real{0.1250}}
  >{\raggedright\arraybackslash}p{(\linewidth - 14\tabcolsep) * \real{0.1250}}@{}}
\toprule\noalign{}
\begin{minipage}[b]{\linewidth}\raggedright
RUNSPEC
\end{minipage} & \begin{minipage}[b]{\linewidth}\raggedright
GRID
\end{minipage} & \begin{minipage}[b]{\linewidth}\raggedright
EDIT
\end{minipage} & \begin{minipage}[b]{\linewidth}\raggedright
PROPS
\end{minipage} & \begin{minipage}[b]{\linewidth}\raggedright
REGIONS
\end{minipage} & \begin{minipage}[b]{\linewidth}\raggedright
SOLUTION
\end{minipage} & \begin{minipage}[b]{\linewidth}\raggedright
SUMMARY
\end{minipage} & \begin{minipage}[b]{\linewidth}\raggedright
SCHEDULE
\end{minipage} \\
\midrule\noalign{}
\endhead
\bottomrule\noalign{}
\endlastfoot
\end{longtable}

\subsubsection{Description}\label{description-967}

The GWRTWCV keyword defines the wells and instantaneous gradient
parameters to be calculated and exported as SUMMARY vectors to the
summary file, for when the Instantaneous Gradient option has been
activated by the GDIMS keyword in the RUNSPEC section. The Instantaneous
Gradient option calculates derivatives of solution quantities at the
current time step with respect to variations in the variables at the
current time step. This is different to Gradient option that calculates
the derivatives of solution quantities at the current time step with
respect to variations in the variables at the initial time step, that is
a time equal to zero. Consequently, the Instantaneous Gradient option
can be switched on and off by the GUPFREQ keyword in the SCHEDULE
section, whereas the Gradient option cannot.

This keyword is not supported by OPM Flow but would change the results
if supported so the simulation will be stopped.

\subsection{HMWPIMLT -- History Match Well Productivity Index
Parameters}\label{hmwpimlt-history-match-well-productivity-index-parameters}

\begin{longtable}[]{@{}
  >{\raggedright\arraybackslash}p{(\linewidth - 14\tabcolsep) * \real{0.1250}}
  >{\raggedright\arraybackslash}p{(\linewidth - 14\tabcolsep) * \real{0.1250}}
  >{\raggedright\arraybackslash}p{(\linewidth - 14\tabcolsep) * \real{0.1250}}
  >{\raggedright\arraybackslash}p{(\linewidth - 14\tabcolsep) * \real{0.1250}}
  >{\raggedright\arraybackslash}p{(\linewidth - 14\tabcolsep) * \real{0.1250}}
  >{\raggedright\arraybackslash}p{(\linewidth - 14\tabcolsep) * \real{0.1250}}
  >{\raggedright\arraybackslash}p{(\linewidth - 14\tabcolsep) * \real{0.1250}}
  >{\raggedright\arraybackslash}p{(\linewidth - 14\tabcolsep) * \real{0.1250}}@{}}
\toprule\noalign{}
\begin{minipage}[b]{\linewidth}\raggedright
RUNSPEC
\end{minipage} & \begin{minipage}[b]{\linewidth}\raggedright
GRID
\end{minipage} & \begin{minipage}[b]{\linewidth}\raggedright
EDIT
\end{minipage} & \begin{minipage}[b]{\linewidth}\raggedright
PROPS
\end{minipage} & \begin{minipage}[b]{\linewidth}\raggedright
REGIONS
\end{minipage} & \begin{minipage}[b]{\linewidth}\raggedright
SOLUTION
\end{minipage} & \begin{minipage}[b]{\linewidth}\raggedright
SUMMARY
\end{minipage} & \begin{minipage}[b]{\linewidth}\raggedright
SCHEDULE
\end{minipage} \\
\midrule\noalign{}
\endhead
\bottomrule\noalign{}
\endlastfoot
\end{longtable}

\subsubsection{Description}\label{description-968}

This keyword, HMWPIMLT, defines the history match gradient parameters
for well productiviity indices, for when the History Match Gradient
option has been activated by the HMDIMS keyword in the RUNSPEC section.
Wells must be specified using the WELPSECS keyword in the SCHEDULE
section and their connections defined by the COMPDAT and/or COMPDATL
keywords, also in the SCHEDULE section

See also the HMDIMS keyword in the RUNSPEC section that specifies the
dimensions for the gradient option, including the maximum number of
gradient wellss that can be used with the History Match Gradient option.

This keyword is not supported by OPM Flow but would change the results
if supported so the simulation will be stopped.

\subsection{INCLUDE -- Load Another Data File at the Current
Position}\label{include-load-another-data-file-at-the-current-position-8}

The INCLUDE keyword informs OPM Flow to continue reading input data from
the specified INCLUDE file. When the end of the INCLUDE file is reached,
or the ENDINC keyword is encountered in the included file, input data is
read from the next keyword in the current file.

See INCLUDE -- Load Another Data File at the Current Position in the
GLOBAL section for a full description.

\subsection{LGRFREE -- Activate Local Grid Refinement Independent Time
Steps}\label{lgrfree-activate-local-grid-refinement-independent-time-steps}

\begin{longtable}[]{@{}
  >{\raggedright\arraybackslash}p{(\linewidth - 14\tabcolsep) * \real{0.1250}}
  >{\raggedright\arraybackslash}p{(\linewidth - 14\tabcolsep) * \real{0.1250}}
  >{\raggedright\arraybackslash}p{(\linewidth - 14\tabcolsep) * \real{0.1250}}
  >{\raggedright\arraybackslash}p{(\linewidth - 14\tabcolsep) * \real{0.1250}}
  >{\raggedright\arraybackslash}p{(\linewidth - 14\tabcolsep) * \real{0.1250}}
  >{\raggedright\arraybackslash}p{(\linewidth - 14\tabcolsep) * \real{0.1250}}
  >{\raggedright\arraybackslash}p{(\linewidth - 14\tabcolsep) * \real{0.1250}}
  >{\raggedright\arraybackslash}p{(\linewidth - 14\tabcolsep) * \real{0.1250}}@{}}
\toprule\noalign{}
\begin{minipage}[b]{\linewidth}\raggedright
RUNSPEC
\end{minipage} & \begin{minipage}[b]{\linewidth}\raggedright
GRID
\end{minipage} & \begin{minipage}[b]{\linewidth}\raggedright
EDIT
\end{minipage} & \begin{minipage}[b]{\linewidth}\raggedright
PROPS
\end{minipage} & \begin{minipage}[b]{\linewidth}\raggedright
REGIONS
\end{minipage} & \begin{minipage}[b]{\linewidth}\raggedright
SOLUTION
\end{minipage} & \begin{minipage}[b]{\linewidth}\raggedright
SUMMARY
\end{minipage} & \begin{minipage}[b]{\linewidth}\raggedright
SCHEDULE
\end{minipage} \\
\midrule\noalign{}
\endhead
\bottomrule\noalign{}
\endlastfoot
\end{longtable}

\subsubsection{Description}\label{description-969}

The LGRFREE keyword activates the Local Grid Refinement (``LGR'')
Independent Time Step option that allows the LGR to have solution time
steps independent of the host grid for the stated LGR, and for when LGRs
have been declared by the LGR keyword in the RUNSPEC section, and
defined by the CARFIN (Cartesian LGR grid) or RADIN/RADIN4 (radial LGR
grid) keywords in the GRID section. LGR independent solution time
stepping can be deactivated by the LGRLOCK keyword in the SCHEDULE
section.

Currently, OPM Flow does not support the local grid refinement feature
and therefore this keyword is ignored by the simulator.

\begin{longtable}[]{@{}
  >{\raggedright\arraybackslash}p{(\linewidth - 6\tabcolsep) * \real{0.2500}}
  >{\raggedright\arraybackslash}p{(\linewidth - 6\tabcolsep) * \real{0.2500}}
  >{\raggedright\arraybackslash}p{(\linewidth - 6\tabcolsep) * \real{0.2500}}
  >{\raggedright\arraybackslash}p{(\linewidth - 6\tabcolsep) * \real{0.2500}}@{}}
\toprule\noalign{}
\begin{minipage}[b]{\linewidth}\raggedright
No.
\end{minipage} & \begin{minipage}[b]{\linewidth}\raggedright
Name
\end{minipage} & \begin{minipage}[b]{\linewidth}\raggedright
Description
\end{minipage} & \begin{minipage}[b]{\linewidth}\raggedright
Default
\end{minipage} \\
\midrule\noalign{}
\endhead
\bottomrule\noalign{}
\endlastfoot
1 & LGRNAME & A character string of up to eight characters in length
that defines the LGR name for which independent solution time stepping
is to be activated. The LGR must have been previously defined by the
CARFIN (Cartesian LGR grid) or RADIN/RADIN4 (radial LGR grid) keywords
in the GRID section. & None \\
Notes: & & & \\
\end{longtable}

\emph{Table 12.48: LGRFREE Keyword Description}

\subsubsection{Example}\label{example-538}

The example below defines three oil LGRs(LGR-OP01,-OP02, and -OP03) and
all the gas well LGRs (LGR-GP*) that should use independent solution
time steps.

\begin{verbatim}
--
--       ACTIVATE LOCAL GRID REFINEMENT INDEPENDENT TIME STEPS
--
--       LGRNAME
LGRFREE
         LGR-OP01                                                              /
         LGR-OP02                                                              /
         LGR-OP03                                                              /
         LGR-GP*                                                               /
/
\end{verbatim}

\subsection{LGRLOCK -- Deactivate Local Grid Refinement Independent Time
Steps}\label{lgrlock-deactivate-local-grid-refinement-independent-time-steps}

\begin{longtable}[]{@{}
  >{\raggedright\arraybackslash}p{(\linewidth - 14\tabcolsep) * \real{0.1250}}
  >{\raggedright\arraybackslash}p{(\linewidth - 14\tabcolsep) * \real{0.1250}}
  >{\raggedright\arraybackslash}p{(\linewidth - 14\tabcolsep) * \real{0.1250}}
  >{\raggedright\arraybackslash}p{(\linewidth - 14\tabcolsep) * \real{0.1250}}
  >{\raggedright\arraybackslash}p{(\linewidth - 14\tabcolsep) * \real{0.1250}}
  >{\raggedright\arraybackslash}p{(\linewidth - 14\tabcolsep) * \real{0.1250}}
  >{\raggedright\arraybackslash}p{(\linewidth - 14\tabcolsep) * \real{0.1250}}
  >{\raggedright\arraybackslash}p{(\linewidth - 14\tabcolsep) * \real{0.1250}}@{}}
\toprule\noalign{}
\begin{minipage}[b]{\linewidth}\raggedright
RUNSPEC
\end{minipage} & \begin{minipage}[b]{\linewidth}\raggedright
GRID
\end{minipage} & \begin{minipage}[b]{\linewidth}\raggedright
EDIT
\end{minipage} & \begin{minipage}[b]{\linewidth}\raggedright
PROPS
\end{minipage} & \begin{minipage}[b]{\linewidth}\raggedright
REGIONS
\end{minipage} & \begin{minipage}[b]{\linewidth}\raggedright
SOLUTION
\end{minipage} & \begin{minipage}[b]{\linewidth}\raggedright
SUMMARY
\end{minipage} & \begin{minipage}[b]{\linewidth}\raggedright
SCHEDULE
\end{minipage} \\
\midrule\noalign{}
\endhead
\bottomrule\noalign{}
\endlastfoot
\end{longtable}

\subsubsection{Description}\label{description-970}

The LGRLOCK keyword deactivates the Local Grid Refinement (``LGR'')
Independent Time Step option that allows the LGR to have solution time
steps independent of the host grid for the stated LGR, that is the LGR
will now follow the global grid solution time steps. LGRs must have
declared by the LGR keyword in the RUNSPEC section, and defined by the
CARFIN (Cartesian LGR grid) or RADIN/RADIN4 (radial LGR grid) keywords
in the GRID section. LGR independent solution time stepping can be
activated by the LGRFREE keyword in the SCHEDULE section.

Currently, OPM Flow does not support the local grid refinement feature
and therefore this keyword is ignored by the simulator.

\begin{longtable}[]{@{}
  >{\raggedright\arraybackslash}p{(\linewidth - 6\tabcolsep) * \real{0.2500}}
  >{\raggedright\arraybackslash}p{(\linewidth - 6\tabcolsep) * \real{0.2500}}
  >{\raggedright\arraybackslash}p{(\linewidth - 6\tabcolsep) * \real{0.2500}}
  >{\raggedright\arraybackslash}p{(\linewidth - 6\tabcolsep) * \real{0.2500}}@{}}
\toprule\noalign{}
\begin{minipage}[b]{\linewidth}\raggedright
No.
\end{minipage} & \begin{minipage}[b]{\linewidth}\raggedright
Name
\end{minipage} & \begin{minipage}[b]{\linewidth}\raggedright
Description
\end{minipage} & \begin{minipage}[b]{\linewidth}\raggedright
Default
\end{minipage} \\
\midrule\noalign{}
\endhead
\bottomrule\noalign{}
\endlastfoot
1 & LGRNAME & A character string of up to eight characters in length
that defines the LGR name for which independent solution time stepping
is to be deactivated. The LGR must have been previously defined by the
CARFIN (Cartesian LGR grid) or RADIN/RADIN4 (radial LGR grid) keywords
in the GRID section. & None \\
Notes: & & & \\
\end{longtable}

\emph{Table 12.49: LGRLOCK Keyword Description}

\subsubsection{Example}\label{example-539}

The example below defines three oil LGRs(LGR-OP01,-OP02, and -OP03) and
all the gas well LGRs (LGR-GP*) that should have their independent
solution time steps deactivated.

\begin{verbatim}
--
--       DEACTIVATE LOCAL GRID REFINEMENT INDEPENDENT TIME STEPS
--
--       LGRNAME
LGRLOCK
         LGR-OP01                                                              /
         LGR-OP02                                                              /
         LGR-OP03                                                              /
         LGR-GP*                                                               /
/
\end{verbatim}

\subsection{LGROFF -- Deactivate a Local Grid
Refinement}\label{lgroff-deactivate-a-local-grid-refinement}

\begin{longtable}[]{@{}
  >{\raggedright\arraybackslash}p{(\linewidth - 14\tabcolsep) * \real{0.1250}}
  >{\raggedright\arraybackslash}p{(\linewidth - 14\tabcolsep) * \real{0.1250}}
  >{\raggedright\arraybackslash}p{(\linewidth - 14\tabcolsep) * \real{0.1250}}
  >{\raggedright\arraybackslash}p{(\linewidth - 14\tabcolsep) * \real{0.1250}}
  >{\raggedright\arraybackslash}p{(\linewidth - 14\tabcolsep) * \real{0.1250}}
  >{\raggedright\arraybackslash}p{(\linewidth - 14\tabcolsep) * \real{0.1250}}
  >{\raggedright\arraybackslash}p{(\linewidth - 14\tabcolsep) * \real{0.1250}}
  >{\raggedright\arraybackslash}p{(\linewidth - 14\tabcolsep) * \real{0.1250}}@{}}
\toprule\noalign{}
\begin{minipage}[b]{\linewidth}\raggedright
RUNSPEC
\end{minipage} & \begin{minipage}[b]{\linewidth}\raggedright
GRID
\end{minipage} & \begin{minipage}[b]{\linewidth}\raggedright
EDIT
\end{minipage} & \begin{minipage}[b]{\linewidth}\raggedright
PROPS
\end{minipage} & \begin{minipage}[b]{\linewidth}\raggedright
REGIONS
\end{minipage} & \begin{minipage}[b]{\linewidth}\raggedright
SOLUTION
\end{minipage} & \begin{minipage}[b]{\linewidth}\raggedright
SUMMARY
\end{minipage} & \begin{minipage}[b]{\linewidth}\raggedright
SCHEDULE
\end{minipage} \\
\midrule\noalign{}
\endhead
\bottomrule\noalign{}
\endlastfoot
\end{longtable}

\subsubsection{Description}\label{description-971}

The LGROFF keyword deactivates a stated Local Grid Refinement (``LGR'')
and optionally sets the minimum number of wells below which the LGR will
be automatically deactivated. LGRs must have declared by the LGR keyword
in the RUNSPEC section, and defined by the CARFIN (Cartesian LGR grid)
or RADIN/RADIN4 (radial LGR grid) keywords in the GRID section. LGRs can
subsequently be activated by the LGRON keyword in the SCHEDULE section.

Currently, OPM Flow does not support the local grid refinement feature
and therefore this keyword is ignored by the simulator.

\begin{longtable}[]{@{}
  >{\raggedright\arraybackslash}p{(\linewidth - 6\tabcolsep) * \real{0.2500}}
  >{\raggedright\arraybackslash}p{(\linewidth - 6\tabcolsep) * \real{0.2500}}
  >{\raggedright\arraybackslash}p{(\linewidth - 6\tabcolsep) * \real{0.2500}}
  >{\raggedright\arraybackslash}p{(\linewidth - 6\tabcolsep) * \real{0.2500}}@{}}
\toprule\noalign{}
\begin{minipage}[b]{\linewidth}\raggedright
No.
\end{minipage} & \begin{minipage}[b]{\linewidth}\raggedright
Name
\end{minipage} & \begin{minipage}[b]{\linewidth}\raggedright
Description
\end{minipage} & \begin{minipage}[b]{\linewidth}\raggedright
Default
\end{minipage} \\
\midrule\noalign{}
\endhead
\bottomrule\noalign{}
\endlastfoot
1 & LGRNAME & A character string of up to eight characters in length
that defines the LGR name for which the LGR is being deactivated. The
LGR must have been previously defined by the CARFIN (Cartesian LGR grid)
or RADIN/RADIN4 (radial LGR grid) keywords in the GRID section. &
None \\
2 & MNWELLS & A positive integer greater than or equal to zero that
defines the minimum number of active wells, below which the LGR will be
automatically deactivated. The default value of zero implies that there
is no limit to the number of wells and results in the LGR being
unconditionally being deactivated. & 0 \\
Notes: & & & \\
\end{longtable}

\emph{Table 12.50: LGROFF Keyword Description}

\subsubsection{Example}\label{example-540}

The example below unconditionally deactivates LGR-OP01, and sets the
minimum number of active wells for deactivating LGR-OP02 and LGR-OP03 to
one. For all the gas well LGRs (LGR-GP*) the minimum number of wells for
deactivation is set to two.

\begin{verbatim}
--
--       DEACTIVATE LOCAL GRID REFINEMENTS
--
--       LGRNAME   MNWELLS
LGROFF
         LGR-OP01                                                              /
         LGR-OP02  1                                                           /
         LGR-OP03  1                                                           /
         LGR-GP*   2                                                           /
/
\end{verbatim}

\subsection{LGRON -- Activate a Local Grid
Refinement}\label{lgron-activate-a-local-grid-refinement}

\begin{longtable}[]{@{}
  >{\raggedright\arraybackslash}p{(\linewidth - 14\tabcolsep) * \real{0.1250}}
  >{\raggedright\arraybackslash}p{(\linewidth - 14\tabcolsep) * \real{0.1250}}
  >{\raggedright\arraybackslash}p{(\linewidth - 14\tabcolsep) * \real{0.1250}}
  >{\raggedright\arraybackslash}p{(\linewidth - 14\tabcolsep) * \real{0.1250}}
  >{\raggedright\arraybackslash}p{(\linewidth - 14\tabcolsep) * \real{0.1250}}
  >{\raggedright\arraybackslash}p{(\linewidth - 14\tabcolsep) * \real{0.1250}}
  >{\raggedright\arraybackslash}p{(\linewidth - 14\tabcolsep) * \real{0.1250}}
  >{\raggedright\arraybackslash}p{(\linewidth - 14\tabcolsep) * \real{0.1250}}@{}}
\toprule\noalign{}
\begin{minipage}[b]{\linewidth}\raggedright
RUNSPEC
\end{minipage} & \begin{minipage}[b]{\linewidth}\raggedright
GRID
\end{minipage} & \begin{minipage}[b]{\linewidth}\raggedright
EDIT
\end{minipage} & \begin{minipage}[b]{\linewidth}\raggedright
PROPS
\end{minipage} & \begin{minipage}[b]{\linewidth}\raggedright
REGIONS
\end{minipage} & \begin{minipage}[b]{\linewidth}\raggedright
SOLUTION
\end{minipage} & \begin{minipage}[b]{\linewidth}\raggedright
SUMMARY
\end{minipage} & \begin{minipage}[b]{\linewidth}\raggedright
SCHEDULE
\end{minipage} \\
\midrule\noalign{}
\endhead
\bottomrule\noalign{}
\endlastfoot
\end{longtable}

\subsubsection{Description}\label{description-972}

The LGRON keyword activates a stated Local Grid Refinement (``LGR'') and
optionally sets the minimum number of wells above which the LGR will
remain active. LGRs must have declared by the LGR keyword in the RUNSPEC
section, and defined by the CARFIN (Cartesian LGR grid) or RADIN/RADIN4
(radial LGR grid) keywords in the GRID section. LGRs can subsequently be
deactivated by the LGROFF keyword in the SCHEDULE section.

Currently, OPM Flow does not support the local grid refinement feature
and therefore this keyword is ignored by the simulator.

\begin{longtable}[]{@{}
  >{\raggedright\arraybackslash}p{(\linewidth - 6\tabcolsep) * \real{0.2500}}
  >{\raggedright\arraybackslash}p{(\linewidth - 6\tabcolsep) * \real{0.2500}}
  >{\raggedright\arraybackslash}p{(\linewidth - 6\tabcolsep) * \real{0.2500}}
  >{\raggedright\arraybackslash}p{(\linewidth - 6\tabcolsep) * \real{0.2500}}@{}}
\toprule\noalign{}
\begin{minipage}[b]{\linewidth}\raggedright
No.
\end{minipage} & \begin{minipage}[b]{\linewidth}\raggedright
Name
\end{minipage} & \begin{minipage}[b]{\linewidth}\raggedright
Description
\end{minipage} & \begin{minipage}[b]{\linewidth}\raggedright
Default
\end{minipage} \\
\midrule\noalign{}
\endhead
\bottomrule\noalign{}
\endlastfoot
1 & LGRNAME & A character string of up to eight characters in length
that defines the LGR name for which the LGR is being activated. The LGR
must have been previously defined by the CARFIN (Cartesian LGR grid) or
RADIN/RADIN4 (radial LGR grid) keywords in the GRID section. & None \\
2 & MNWELLS & A positive integer greater than or equal to zero that
defines the minimum number of active wells, below which the LGR will be
automatically deactivated. The default value of zero implies that there
is no limit to the number of wells and results in the LGR being
unconditionally being activated. & 0 \\
Notes: & & & \\
\end{longtable}

\emph{Table 12.51: LGRON Keyword Description}

\subsubsection{Example}\label{example-541}

The example below unconditionally activates LGR-OP01, and sets the
minimum number of active wells for activating LGR-OP02 and LGR-OP03 to
one. For all the gas well LGRs (LGR-GP*) the minimum number of wells for
activating these LGRs is set to two.

\begin{verbatim}
--
--       ACTIVATE LOCAL GRID REFINEMENTS
--
--       LGRNAME   MNWELLS
LGRON
         LGR-OP01                                                              /
         LGR-OP02  1                                                           /
         LGR-OP03  1                                                           /
         LGR-GP*   2                                                           /
/
\end{verbatim}

\subsection{LIFTOPT -- Activate Gas Lift
Optimization}\label{liftopt-activate-gas-lift-optimization}

\begin{longtable}[]{@{}
  >{\raggedright\arraybackslash}p{(\linewidth - 14\tabcolsep) * \real{0.1250}}
  >{\raggedright\arraybackslash}p{(\linewidth - 14\tabcolsep) * \real{0.1250}}
  >{\raggedright\arraybackslash}p{(\linewidth - 14\tabcolsep) * \real{0.1250}}
  >{\raggedright\arraybackslash}p{(\linewidth - 14\tabcolsep) * \real{0.1250}}
  >{\raggedright\arraybackslash}p{(\linewidth - 14\tabcolsep) * \real{0.1250}}
  >{\raggedright\arraybackslash}p{(\linewidth - 14\tabcolsep) * \real{0.1250}}
  >{\raggedright\arraybackslash}p{(\linewidth - 14\tabcolsep) * \real{0.1250}}
  >{\raggedright\arraybackslash}p{(\linewidth - 14\tabcolsep) * \real{0.1250}}@{}}
\toprule\noalign{}
\begin{minipage}[b]{\linewidth}\raggedright
RUNSPEC
\end{minipage} & \begin{minipage}[b]{\linewidth}\raggedright
GRID
\end{minipage} & \begin{minipage}[b]{\linewidth}\raggedright
EDIT
\end{minipage} & \begin{minipage}[b]{\linewidth}\raggedright
PROPS
\end{minipage} & \begin{minipage}[b]{\linewidth}\raggedright
REGIONS
\end{minipage} & \begin{minipage}[b]{\linewidth}\raggedright
SOLUTION
\end{minipage} & \begin{minipage}[b]{\linewidth}\raggedright
SUMMARY
\end{minipage} & \begin{minipage}[b]{\linewidth}\raggedright
SCHEDULE
\end{minipage} \\
\midrule\noalign{}
\endhead
\bottomrule\noalign{}
\endlastfoot
\end{longtable}

\subsubsection{Description}\label{description-973}

The LIFTOPT keyword actives the gas lift optimization option and defines
the gas lift gas increment size, the minimum incremental oil
improvement, as well as the timing of the calculations. Note that the
LIFTOPT keyword should precede any GLIFTOPT and WLIFTOPT keywords in the
SCHEDULE section in order to activate the gas lift optimization
facility.

\begin{longtable}[]{@{}
  >{\raggedright\arraybackslash}p{(\linewidth - 6\tabcolsep) * \real{0.2500}}
  >{\raggedright\arraybackslash}p{(\linewidth - 6\tabcolsep) * \real{0.2500}}
  >{\raggedright\arraybackslash}p{(\linewidth - 6\tabcolsep) * \real{0.2500}}
  >{\raggedright\arraybackslash}p{(\linewidth - 6\tabcolsep) * \real{0.2500}}@{}}
\toprule\noalign{}
\begin{minipage}[b]{\linewidth}\raggedright
No.
\end{minipage} & \begin{minipage}[b]{\linewidth}\raggedright
Name
\end{minipage} & \begin{minipage}[b]{\linewidth}\raggedright
Description
\end{minipage} & \begin{minipage}[b]{\linewidth}\raggedright
Default
\end{minipage} \\
\midrule\noalign{}
\endhead
\bottomrule\noalign{}
\endlastfoot
Field & Metric & Laboratory & \\
1 & GASLIFT & A real positive number that defines the gas lift gas size
increment that is used to increase the gas lift size quantity in steps.
For example, if GASLIFT is set to 0.5 MMscf/d then gas lift gas will be
allocated in step of 0.5 MMscf/d to each well during the optimization
process. A zero or negative value switches off gas lift optimization. &
None \\
Mscf/d & sm3/day & scc/hour & \\
2 & MINOIL & MINOIL is a real positive value that defines the minimum
increase in oil rate for a given quantity of gas lift gas, for when gas
lift gas should be applied to a well. Additional GASLIFT will only be
assigned to a well if:
\(\frac{\mathrm{Δ}{Q}_{\mathit{Oil}} \times  \mathit{OPTWGT}}{\mathit{GASLIFT}} > \mathit{MINOIL}\)
Where ΔQOil is the incremental oil and OPTWGT is the well's weighting
factor defined by the OPTWGT variable on the WLIFTOPT keyword in the
SCHEDULE section. & None \\
stb/Mscf & sm3/sm3 & scc/scc & \\
3 & TSTEP & TSTEP is a real positive value that defines the frequency of
the gas lift optimization calculations, for example setting TSTEP equal
to 30 days would result in the gas lift optimization calculation being
performed approximately every 30 days. The default value of zero will
result in the calculations being performed every time step. Note if the
group or well is part of a production network then gas lift optimization
is performed at the same time as the network is being balance, that is
this parameter is ignored in this scenario. See the NETBALAN keyword in
the SCHEDULE section to set the network balancing frequency in this
case. & 0.0 \\
days & days & hours & \\
4 & OPTLIFT & A defined character string that determines if the gas lift
optimization iterations should be performed for the same number of
Newton iterations within a time step as used to update well targets, or
to just use the first Newton iteration only. The NUPCOL keyword in the
RUNSPEC section determines the number of Newton iterations used to
update well targets during a time step. OPTLIFT should be set to one of
the following: & \\
Notes: & & & \\
\end{longtable}

\emph{Table 12.52: LIFTOPT Keyword Description}

See also the GLIFTOPT keyword to define the group gas lift optimization
controls and the WLIFTOPT keyword to define the wells under gas lift
optimization control, both keywords are described in the SCHEDULE
section. The NUPCOL keyword in the RUNSPEC section that determines the
number of Newton iterations used to update well targets and gas lift
optimization calculations during a time step may also be of interest.

\subsubsection{Example}\label{example-542}

The following example activates gas lift optimization for the field and
defines the optimization parameters.

\begin{verbatim}
--
--       ACTIVATE GAS LIFT OPTIMIZATION AND PARAMETERS
--
--       INCR   INCR   TSTEP    NEWTON
--       GAS    OIL    INTVAL   OPTN
LIFTOPT
        12.5E3  5E-3   0.0      YES    /
\end{verbatim}

Here the maximum incremental gas lift gas quantity is set to 12.5 x 103
m3, the minimum incremental oil gain per m3 of gas lift gas is set to
5.0 x 10-3 m3, the time step interval is set to zero to perform the gas
optimization every time step, and finally the gas lift optimization will
be performed NUPCOL Newton iterations for the time step.

\subsection{LINCOM -- Define Linear Combination
Coefficients}\label{lincom-define-linear-combination-coefficients}

\begin{longtable}[]{@{}
  >{\raggedright\arraybackslash}p{(\linewidth - 14\tabcolsep) * \real{0.1250}}
  >{\raggedright\arraybackslash}p{(\linewidth - 14\tabcolsep) * \real{0.1250}}
  >{\raggedright\arraybackslash}p{(\linewidth - 14\tabcolsep) * \real{0.1250}}
  >{\raggedright\arraybackslash}p{(\linewidth - 14\tabcolsep) * \real{0.1250}}
  >{\raggedright\arraybackslash}p{(\linewidth - 14\tabcolsep) * \real{0.1250}}
  >{\raggedright\arraybackslash}p{(\linewidth - 14\tabcolsep) * \real{0.1250}}
  >{\raggedright\arraybackslash}p{(\linewidth - 14\tabcolsep) * \real{0.1250}}
  >{\raggedright\arraybackslash}p{(\linewidth - 14\tabcolsep) * \real{0.1250}}@{}}
\toprule\noalign{}
\begin{minipage}[b]{\linewidth}\raggedright
RUNSPEC
\end{minipage} & \begin{minipage}[b]{\linewidth}\raggedright
GRID
\end{minipage} & \begin{minipage}[b]{\linewidth}\raggedright
EDIT
\end{minipage} & \begin{minipage}[b]{\linewidth}\raggedright
PROPS
\end{minipage} & \begin{minipage}[b]{\linewidth}\raggedright
REGIONS
\end{minipage} & \begin{minipage}[b]{\linewidth}\raggedright
SOLUTION
\end{minipage} & \begin{minipage}[b]{\linewidth}\raggedright
SUMMARY
\end{minipage} & \begin{minipage}[b]{\linewidth}\raggedright
SCHEDULE
\end{minipage} \\
\midrule\noalign{}
\endhead
\bottomrule\noalign{}
\endlastfoot
\end{longtable}

\subsubsection{Description}\label{description-974}

The LINCOM keyword defines the oil, gas and water coefficients for the
Linear Combination facility which allows for a linear combination of the
aforementioned phase rates and volumes to be used as targets and
constraints in controlling group and well production and injection data.
See also the LCUNIT in the PROPS section that defines the units for
linear combination equation.

This keyword is not supported by OPM Flow but would change the results
if supported so the simulation will be stopped.

\subsection{MATCORR -- Activate the Material Balance Correction
Option}\label{matcorr-activate-the-material-balance-correction-option}

\begin{longtable}[]{@{}
  >{\raggedright\arraybackslash}p{(\linewidth - 14\tabcolsep) * \real{0.1250}}
  >{\raggedright\arraybackslash}p{(\linewidth - 14\tabcolsep) * \real{0.1250}}
  >{\raggedright\arraybackslash}p{(\linewidth - 14\tabcolsep) * \real{0.1250}}
  >{\raggedright\arraybackslash}p{(\linewidth - 14\tabcolsep) * \real{0.1250}}
  >{\raggedright\arraybackslash}p{(\linewidth - 14\tabcolsep) * \real{0.1250}}
  >{\raggedright\arraybackslash}p{(\linewidth - 14\tabcolsep) * \real{0.1250}}
  >{\raggedright\arraybackslash}p{(\linewidth - 14\tabcolsep) * \real{0.1250}}
  >{\raggedright\arraybackslash}p{(\linewidth - 14\tabcolsep) * \real{0.1250}}@{}}
\toprule\noalign{}
\begin{minipage}[b]{\linewidth}\raggedright
RUNSPEC
\end{minipage} & \begin{minipage}[b]{\linewidth}\raggedright
GRID
\end{minipage} & \begin{minipage}[b]{\linewidth}\raggedright
EDIT
\end{minipage} & \begin{minipage}[b]{\linewidth}\raggedright
PROPS
\end{minipage} & \begin{minipage}[b]{\linewidth}\raggedright
REGIONS
\end{minipage} & \begin{minipage}[b]{\linewidth}\raggedright
SOLUTION
\end{minipage} & \begin{minipage}[b]{\linewidth}\raggedright
SUMMARY
\end{minipage} & \begin{minipage}[b]{\linewidth}\raggedright
SCHEDULE
\end{minipage} \\
\midrule\noalign{}
\endhead
\bottomrule\noalign{}
\endlastfoot
\end{longtable}

\subsubsection{Description}\label{description-975}

The MATCORR keyword activates the Material Balance Correction option
used to adjust the accumulated material balance error in the simulation.

This keyword is not supported by OPM Flow but would change the results
if supported so the simulation will be stopped.

\subsection{MESSAGE -- Output User
Message}\label{message-output-user-message-8}

The MESSAGE keyword outputs a user message to the terminal, as well as
to the print (\emph{.PRT) and debug (}.DBG) files. Note this is
different to the MESSAGES keyword, that defines OPM Flows message print
limits and stop limits generated by the simulator.

See MESSAGE -- Output User Message in the GLOBAL section for a full
description.

\subsection{MESSAGES -- Define Message Print Limits and Stop
Limits}\label{messages-define-message-print-limits-and-stop-limits-8}

The MESSAGES keyword defines the print and stops levels for various
messages. The ``print limits'' set the maximum number of messages that
will be printed, after which no more messages will be printed and the
``stop limits'' terminate the run when these limits are exceeded. There
are six levels of message that increase in severity from informative all
the way to programming errors, as outlined in Table 4.5.

See MESSAGES -- Define Message Print Limits and Stop Limits in the
GLOBAL section for a full description.

\subsection{MESSOPTS -- Reset Severity Level for Forced Time
Steps}\label{messopts-reset-severity-level-for-forced-time-steps}

\begin{longtable}[]{@{}
  >{\raggedright\arraybackslash}p{(\linewidth - 14\tabcolsep) * \real{0.1250}}
  >{\raggedright\arraybackslash}p{(\linewidth - 14\tabcolsep) * \real{0.1250}}
  >{\raggedright\arraybackslash}p{(\linewidth - 14\tabcolsep) * \real{0.1250}}
  >{\raggedright\arraybackslash}p{(\linewidth - 14\tabcolsep) * \real{0.1250}}
  >{\raggedright\arraybackslash}p{(\linewidth - 14\tabcolsep) * \real{0.1250}}
  >{\raggedright\arraybackslash}p{(\linewidth - 14\tabcolsep) * \real{0.1250}}
  >{\raggedright\arraybackslash}p{(\linewidth - 14\tabcolsep) * \real{0.1250}}
  >{\raggedright\arraybackslash}p{(\linewidth - 14\tabcolsep) * \real{0.1250}}@{}}
\toprule\noalign{}
\begin{minipage}[b]{\linewidth}\raggedright
RUNSPEC
\end{minipage} & \begin{minipage}[b]{\linewidth}\raggedright
GRID
\end{minipage} & \begin{minipage}[b]{\linewidth}\raggedright
EDIT
\end{minipage} & \begin{minipage}[b]{\linewidth}\raggedright
PROPS
\end{minipage} & \begin{minipage}[b]{\linewidth}\raggedright
REGIONS
\end{minipage} & \begin{minipage}[b]{\linewidth}\raggedright
SOLUTION
\end{minipage} & \begin{minipage}[b]{\linewidth}\raggedright
SUMMARY
\end{minipage} & \begin{minipage}[b]{\linewidth}\raggedright
SCHEDULE
\end{minipage} \\
\midrule\noalign{}
\endhead
\bottomrule\noalign{}
\endlastfoot
\end{longtable}

\subsubsection{Description}\label{description-976}

This keyword, MESSOPTS, resets the severity level for time steps that
are forced to be accepted by the simulator. The normal severity level
for this type of simulator generated message is PROBLEM and this can
result in the run stopping depending on the parameters entered on the
MESSAGES keyword. MESSOPTS can be used to reset the severity level to
MESSAGE, COMMENT, WARNING, or PROBLEM; for example, to avoid the run
terminating due to too many PROBLEM messages.

Note that the MESSAGES keyword is a global keyword can therefore be used
in any section; however, only the last instance of the keywords is
active. The MESSOPTS keyword can only be used in the SCHEDULE section
but can be used multiple times to change the severity level for forced
time steps. Again, only the last occurrence of the keyword is active.

This keyword is not supported by OPM Flow but has no effect on the
results so it will be ignored.

\subsection{MULSGGD -- Multiply Matrix-Fracture Coupling for Oil-Gas
Gravity Drainage for All
Cells}\label{mulsggd-multiply-matrix-fracture-coupling-for-oil-gas-gravity-drainage-for-all-cells}

\begin{longtable}[]{@{}
  >{\raggedright\arraybackslash}p{(\linewidth - 14\tabcolsep) * \real{0.1250}}
  >{\raggedright\arraybackslash}p{(\linewidth - 14\tabcolsep) * \real{0.1250}}
  >{\raggedright\arraybackslash}p{(\linewidth - 14\tabcolsep) * \real{0.1250}}
  >{\raggedright\arraybackslash}p{(\linewidth - 14\tabcolsep) * \real{0.1250}}
  >{\raggedright\arraybackslash}p{(\linewidth - 14\tabcolsep) * \real{0.1250}}
  >{\raggedright\arraybackslash}p{(\linewidth - 14\tabcolsep) * \real{0.1250}}
  >{\raggedright\arraybackslash}p{(\linewidth - 14\tabcolsep) * \real{0.1250}}
  >{\raggedright\arraybackslash}p{(\linewidth - 14\tabcolsep) * \real{0.1250}}@{}}
\toprule\noalign{}
\begin{minipage}[b]{\linewidth}\raggedright
RUNSPEC
\end{minipage} & \begin{minipage}[b]{\linewidth}\raggedright
GRID
\end{minipage} & \begin{minipage}[b]{\linewidth}\raggedright
EDIT
\end{minipage} & \begin{minipage}[b]{\linewidth}\raggedright
PROPS
\end{minipage} & \begin{minipage}[b]{\linewidth}\raggedright
REGIONS
\end{minipage} & \begin{minipage}[b]{\linewidth}\raggedright
SOLUTION
\end{minipage} & \begin{minipage}[b]{\linewidth}\raggedright
SUMMARY
\end{minipage} & \begin{minipage}[b]{\linewidth}\raggedright
SCHEDULE
\end{minipage} \\
\midrule\noalign{}
\endhead
\bottomrule\noalign{}
\endlastfoot
\end{longtable}

\subsubsection{Description}\label{description-977}

This keyword, MULSGGD, defines a constant multiplier to modify the
matrix-fracture coupling transmissibility for dual porosity and dual
permeability models, for when the alternative matrix-fracture coupling
transmissibilities for oil-gas gravity drainage has been selected. The
alternative matrix-fracture coupling transmissibilities for oil-gas
gravity drainage option is activated via the SIGMAGD or SIGMAGDV
keywords in the GRID section, and the dual porosity or dual permeability
models are activated by the DUALPORO or DUALPERM keywords in the RUNSPEC
section, respectively.

This keyword applies the multiplier for all cells in model; whereas, the
MULSGGDV keyword in the SCHEDULE section applies the multiplier to
individual grid blocks.

This keyword is not supported by OPM Flow but would change the results
if supported so the simulation will be stopped.

\subsection{MULSGGDV -- Multiply Matrix-Fracture Coupling for Oil-Gas
Gravity Drainage for Individual
Cells}\label{mulsggdv-multiply-matrix-fracture-coupling-for-oil-gas-gravity-drainage-for-individual-cells}

\begin{longtable}[]{@{}
  >{\raggedright\arraybackslash}p{(\linewidth - 14\tabcolsep) * \real{0.1250}}
  >{\raggedright\arraybackslash}p{(\linewidth - 14\tabcolsep) * \real{0.1250}}
  >{\raggedright\arraybackslash}p{(\linewidth - 14\tabcolsep) * \real{0.1250}}
  >{\raggedright\arraybackslash}p{(\linewidth - 14\tabcolsep) * \real{0.1250}}
  >{\raggedright\arraybackslash}p{(\linewidth - 14\tabcolsep) * \real{0.1250}}
  >{\raggedright\arraybackslash}p{(\linewidth - 14\tabcolsep) * \real{0.1250}}
  >{\raggedright\arraybackslash}p{(\linewidth - 14\tabcolsep) * \real{0.1250}}
  >{\raggedright\arraybackslash}p{(\linewidth - 14\tabcolsep) * \real{0.1250}}@{}}
\toprule\noalign{}
\begin{minipage}[b]{\linewidth}\raggedright
RUNSPEC
\end{minipage} & \begin{minipage}[b]{\linewidth}\raggedright
GRID
\end{minipage} & \begin{minipage}[b]{\linewidth}\raggedright
EDIT
\end{minipage} & \begin{minipage}[b]{\linewidth}\raggedright
PROPS
\end{minipage} & \begin{minipage}[b]{\linewidth}\raggedright
REGIONS
\end{minipage} & \begin{minipage}[b]{\linewidth}\raggedright
SOLUTION
\end{minipage} & \begin{minipage}[b]{\linewidth}\raggedright
SUMMARY
\end{minipage} & \begin{minipage}[b]{\linewidth}\raggedright
SCHEDULE
\end{minipage} \\
\midrule\noalign{}
\endhead
\bottomrule\noalign{}
\endlastfoot
\end{longtable}

\subsubsection{Description}\label{description-978}

This keyword, MULSGGDV, defines a constant multiplier to modify the
matrix-fracture coupling transmissibility for dual porosity and dual
permeability models, for when the alternative matrix-fracture coupling
transmissibilities for oil-gas gravity drainage has been selected. The
alternative matrix-fracture coupling transmissibilities for oil-gas
gravity drainage option is activated via the SIGMAGD or SIGMAGDV
keywords in the GRID section, and the dual porosity or dual permeability
models are activated by the DUALPORO or DUALPERM keywords in the RUNSPEC
section, respectively.

This keyword applies the multiplier for individual cells in model;
whereas, the MULSGGD keyword in the SCHEDULE section applies the
multiplier to all grid blocks.

This keyword is not supported by OPM Flow but would change the results
if supported so the simulation will be stopped.

\subsection{MULTFLT -- Multiply the Transmissibility of a Defined Fault
by a
Constant}\label{multflt-multiply-the-transmissibility-of-a-defined-fault-by-a-constant-2}

MULTFLT enables the transmissibilities across defined faults, as
declared by the FAULTS keyword, to be modified. They keyword allows for
the re-scaling of the existing fault transmissibilities calculated by
OPM Flow, for example setting a fault to be completely sealing by
setting the multiplier to zero..

See MULTFLT -- Multiply the Transmissibility of a Defined Fault by a
Constant in the GRID section for a full description.

\subsection{MULTPV -- Multiply Cell Pore Volumes by a
Constant}\label{multpv-multiply-cell-pore-volumes-by-a-constant-2}

MULTPV multiples the pore volumes of a cell by a real positive constant
for all the cells in the model via an array. An alternative to defining
the complete array is to use the BOX keyword to define an area of the
grid and then use the MULTPV keyword to set the multipliers just for the
area defined by the BOX keyword (see the example).

See MULTPV -- Multiply Cell Pore Volumes by a Constantin the GRID
section for a full description.

\subsection{MULTR- -- Multiply Cell Transmissibility in the -R
Direction}\label{multr--multiply-cell-transmissibility-in-the--r-direction-2}

MULTX- multiples the transmissibility between two cell faces in the -X
direction for all the cells in the model via an array, that is the
keyword sets the transmissibility multiplier of block (I, J, K) between
the cells (I-1, J, K) and (I, J, K). The keyword should only be used
with radial grids, as declared by the RADIAL keyword in the RUNSPEC
section.

See MULTR- - Multiply Cell Transmissibility in the -R Direction in the
GRID section for a full description.

\subsection{MULTR -- Multiply Cell Transmissibility in the +R
Direction}\label{multr-multiply-cell-transmissibility-in-the-r-direction-2}

MULTR multiples the transmissibility between two cell faces in the +R
direction for all the cells in the model via an array, that is the
keyword sets the transmissibility multiplier of block (I, J, K) between
the cells (I, J, K) and (I+I, J, K). The keyword should only be used
with radial grids, as declared by the RADIAL keyword in the RUNSPEC
section.

See MULTR - Multiply Cell Transmissibility in the +R Direction in the
GRID section for a full description.

\subsection{MULTREGT -- Multiply Transmissibilities Between
Regions}\label{multregt-multiply-transmissibilities-between-regions-2}

The MULTREGT keyword multiplies the transmissibility between two regions
by a constant. The region number array can be FLUXNUM, MULTNUM or
OPERNUM and these arrays must be defined and be available before the
MULTREGT keyword is read by the simulator. The constant should be a real
number.

See MULTREGT -- Multiply Transmissibilities Between Regions in the GRID
section for a full description.

\subsection{MULTSIG -- Multiply Matrix-Fracture Coupling for All
Cells}\label{multsig-multiply-matrix-fracture-coupling-for-all-cells}

\begin{longtable}[]{@{}
  >{\raggedright\arraybackslash}p{(\linewidth - 14\tabcolsep) * \real{0.1250}}
  >{\raggedright\arraybackslash}p{(\linewidth - 14\tabcolsep) * \real{0.1250}}
  >{\raggedright\arraybackslash}p{(\linewidth - 14\tabcolsep) * \real{0.1250}}
  >{\raggedright\arraybackslash}p{(\linewidth - 14\tabcolsep) * \real{0.1250}}
  >{\raggedright\arraybackslash}p{(\linewidth - 14\tabcolsep) * \real{0.1250}}
  >{\raggedright\arraybackslash}p{(\linewidth - 14\tabcolsep) * \real{0.1250}}
  >{\raggedright\arraybackslash}p{(\linewidth - 14\tabcolsep) * \real{0.1250}}
  >{\raggedright\arraybackslash}p{(\linewidth - 14\tabcolsep) * \real{0.1250}}@{}}
\toprule\noalign{}
\begin{minipage}[b]{\linewidth}\raggedright
RUNSPEC
\end{minipage} & \begin{minipage}[b]{\linewidth}\raggedright
GRID
\end{minipage} & \begin{minipage}[b]{\linewidth}\raggedright
EDIT
\end{minipage} & \begin{minipage}[b]{\linewidth}\raggedright
PROPS
\end{minipage} & \begin{minipage}[b]{\linewidth}\raggedright
REGIONS
\end{minipage} & \begin{minipage}[b]{\linewidth}\raggedright
SOLUTION
\end{minipage} & \begin{minipage}[b]{\linewidth}\raggedright
SUMMARY
\end{minipage} & \begin{minipage}[b]{\linewidth}\raggedright
SCHEDULE
\end{minipage} \\
\midrule\noalign{}
\endhead
\bottomrule\noalign{}
\endlastfoot
\end{longtable}

\subsubsection{Description}\label{description-979}

This keyword, MULTSIG, defines a constant multiplier to modify the
matrix-fracture coupling transmissibility for dual porosity and dual
permeability models, for when the matrix-fracture coupling
transmissibilities have been specified via the SIGMAGD or SIGMAGDV
keywords in the GRID section, and the dual porosity or dual permeability
models have been activated by the DUALPORO or DUALPERM keywords in the
RUNSPEC section, respectively.

This keyword applies the multiplier for all cells in model; whereas, the
MULTSIGV keyword in the SCHEDULE section applies the multiplier to
individual grid blocks.

This keyword is not supported by OPM Flow but would change the results
if supported so the simulation will be stopped.

\subsection{MULTSIGV -- Multiply Matrix-Fracture Coupling for Individual
Cells}\label{multsigv-multiply-matrix-fracture-coupling-for-individual-cells}

\begin{longtable}[]{@{}
  >{\raggedright\arraybackslash}p{(\linewidth - 14\tabcolsep) * \real{0.1250}}
  >{\raggedright\arraybackslash}p{(\linewidth - 14\tabcolsep) * \real{0.1250}}
  >{\raggedright\arraybackslash}p{(\linewidth - 14\tabcolsep) * \real{0.1250}}
  >{\raggedright\arraybackslash}p{(\linewidth - 14\tabcolsep) * \real{0.1250}}
  >{\raggedright\arraybackslash}p{(\linewidth - 14\tabcolsep) * \real{0.1250}}
  >{\raggedright\arraybackslash}p{(\linewidth - 14\tabcolsep) * \real{0.1250}}
  >{\raggedright\arraybackslash}p{(\linewidth - 14\tabcolsep) * \real{0.1250}}
  >{\raggedright\arraybackslash}p{(\linewidth - 14\tabcolsep) * \real{0.1250}}@{}}
\toprule\noalign{}
\begin{minipage}[b]{\linewidth}\raggedright
RUNSPEC
\end{minipage} & \begin{minipage}[b]{\linewidth}\raggedright
GRID
\end{minipage} & \begin{minipage}[b]{\linewidth}\raggedright
EDIT
\end{minipage} & \begin{minipage}[b]{\linewidth}\raggedright
PROPS
\end{minipage} & \begin{minipage}[b]{\linewidth}\raggedright
REGIONS
\end{minipage} & \begin{minipage}[b]{\linewidth}\raggedright
SOLUTION
\end{minipage} & \begin{minipage}[b]{\linewidth}\raggedright
SUMMARY
\end{minipage} & \begin{minipage}[b]{\linewidth}\raggedright
SCHEDULE
\end{minipage} \\
\midrule\noalign{}
\endhead
\bottomrule\noalign{}
\endlastfoot
\end{longtable}

\subsubsection{Description}\label{description-980}

This keyword, MULTSIGV, defines a constant multiplier to modify the
matrix-fracture coupling transmissibility for dual porosity and dual
permeability models, for when the matrix-fracture coupling
transmissibilities have been specified via the SIGMAGD or SIGMAGDV
keywords in the GRID section, and the dual porosity or dual permeability
models have activated by the DUALPORO or DUALPERM keywords in the
RUNSPEC section, respectively.

This keyword applies the multiplier for individual cells in model;
whereas, the MULTSIG keyword in the SCHEDULE section applies the
multiplier to all grid blocks.

This keyword is not supported by OPM Flow but would change the results
if supported so the simulation will be stopped.

\subsection{MULTTHT- -- Multiply Cell Transmissibility in the -Theta
Direction}\label{multtht--multiply-cell-transmissibility-in-the--theta-direction-2}

MULTTHT- multiples the transmissibility between two cell faces in the
-Theta direction for all the cells in the model via an array, that is
the keyword sets the transmissibility multiplier of block (I, J, K)
between the cells (I, J-1, K) and (I, J, K). The keyword should only be
used with radial grids, as declared by the RADIAL keyword in the RUNSPEC
section.

See MULTTHT- - Multiply Cell Transmissibility in the -Theta Direction in
the GRID section for a full description.

\subsection{MULTTHT -- Multiply Cell Transmissibility in the +Theta
Direction}\label{multtht-multiply-cell-transmissibility-in-the-theta-direction-2}

MULTTHT multiples the transmissibility between two cell faces in the
+Theta direction for all the cells in the model via an array, that is
the keyword sets the transmissibility multiplier of block (I, J, K)
between the cells (I, J, K) and (I, J+1, K). The keyword should only be
used with radial grids, as declared by the RADIAL keyword in the RUNSPEC
section.

See MULTTHT - Multiply Cell Transmissibility in the +Theta Direction in
the GRID section for a full description.

\subsection{MULTX- -- Multiply Cell Transmissibility in the -X
Direction}\label{multx--multiply-cell-transmissibility-in-the--x-direction-2}

MULTX- multiples the transmissibility between two cell faces in the -X
direction for all the cells in the model via an array, that is the
keyword sets the transmissibility multiplier of block (I, J, K) between
the cells (I-1, J, K) and (I, J, K).

See MULTX- - Multiply Cell Transmissibility in the -X Direction in the
GRID section for a full description.

\subsection{MULTX -- Multiply Cell Transmissibility in the +X
Direction}\label{multx-multiply-cell-transmissibility-in-the-x-direction-2}

MULTX multiples the transmissibility between two cell faces in the +X
direction for all the cells in the model via an array, that is the
keyword sets the transmissibility multiplier of block (I, J, K) between
the cells (I, J, K) and (I+I, J, K).

See MULTX - Multiply Cell Transmissibility in the +X Direction in the
GRID section for a full description.

\subsection{MULTY- -- Multiply Cell Transmissibility in the -Y
Direction}\label{multy--multiply-cell-transmissibility-in-the--y-direction-2}

MULTY- multiples the transmissibility between two cell faces in the -Y
direction for all the cells in the model via an array, that is the
keyword sets the transmissibility multiplier of block (I, J, K) between
the cells (I, J-1, K) and (I, J, K).

See MULTY- - Multiply Cell Transmissibility in the -Y Direction in the
GRID section for a full description.

\subsection{MULTY -- Multiply Cell Transmissibility in the +Y
Direction}\label{multy-multiply-cell-transmissibility-in-the-y-direction-2}

MULTY multiples the transmissibility between two cell faces in the +Y
direction for all the cells in the model via an array, that is the
keyword sets the transmissibility multiplier of block (I, J, K) between
the cells (I, J, K) and (I, J+1, K).

See MULTY - Multiply Cell Transmissibility in the +Y Direction in the
GRID section for a full description.

\subsection{MULTZ- -- Multiply Cell Transmissibility in the -Z
Direction}\label{multz--multiply-cell-transmissibility-in-the--z-direction-2}

MULTZ- multiples the transmissibility between two cell faces in the -Z
direction for all the cells in the model via an array, that is the
keyword sets the transmissibility multiplier of block (I, J, K) between
the cells (1, J, K-1) and (I, J, K).

See MULTZ- - Multiply Cell Transmissibility in the -Z Direction in the
GRID section for a full description.

\subsection{MULTZ -- Multiply Cell Transmissibility in the +Z
Direction}\label{multz-multiply-cell-transmissibility-in-the-z-direction-2}

MULTZ multiples the transmissibility between two cell faces in the +Z
direction for all the cells in the model via an array, that is the
keyword sets the transmissibility multiplier of block (I, J, K) between
the cells (I, J, K) and (I, J, K+1).

See MULTZ - Multiply Cell Transmissibility in the +Z Direction in the
GRID section for a full description.

\subsection{NCONSUMP -- Node Gas Consumption (Extended
Network)}\label{nconsump-node-gas-consumption-extended-network}

\begin{longtable}[]{@{}
  >{\raggedright\arraybackslash}p{(\linewidth - 14\tabcolsep) * \real{0.1250}}
  >{\raggedright\arraybackslash}p{(\linewidth - 14\tabcolsep) * \real{0.1250}}
  >{\raggedright\arraybackslash}p{(\linewidth - 14\tabcolsep) * \real{0.1250}}
  >{\raggedright\arraybackslash}p{(\linewidth - 14\tabcolsep) * \real{0.1250}}
  >{\raggedright\arraybackslash}p{(\linewidth - 14\tabcolsep) * \real{0.1250}}
  >{\raggedright\arraybackslash}p{(\linewidth - 14\tabcolsep) * \real{0.1250}}
  >{\raggedright\arraybackslash}p{(\linewidth - 14\tabcolsep) * \real{0.1250}}
  >{\raggedright\arraybackslash}p{(\linewidth - 14\tabcolsep) * \real{0.1250}}@{}}
\toprule\noalign{}
\begin{minipage}[b]{\linewidth}\raggedright
RUNSPEC
\end{minipage} & \begin{minipage}[b]{\linewidth}\raggedright
GRID
\end{minipage} & \begin{minipage}[b]{\linewidth}\raggedright
EDIT
\end{minipage} & \begin{minipage}[b]{\linewidth}\raggedright
PROPS
\end{minipage} & \begin{minipage}[b]{\linewidth}\raggedright
REGIONS
\end{minipage} & \begin{minipage}[b]{\linewidth}\raggedright
SOLUTION
\end{minipage} & \begin{minipage}[b]{\linewidth}\raggedright
SUMMARY
\end{minipage} & \begin{minipage}[b]{\linewidth}\raggedright
SCHEDULE
\end{minipage} \\
\midrule\noalign{}
\endhead
\bottomrule\noalign{}
\endlastfoot
\end{longtable}

\subsubsection{Description}\label{description-981}

The NCONSUMP keyword defines an extended network node's gas consumption
rate, for when the Extended Network option has been activated by the
NETWORK keyword in the RUNSPEC section. The keyword can also be used to
attribute the gas consumption to a previously defined group. See also
the GCONSUMP keyword in the SCHEDULE section that overs more flexibility
and can also be used with the Extended Network option.

This keyword is not supported by OPM Flow but would change the results
if supported so the simulation will be stopped.

\subsection{NEFAC -- Node Efficiency Factors (Extended
Network)}\label{nefac-node-efficiency-factors-extended-network}

\begin{longtable}[]{@{}
  >{\raggedright\arraybackslash}p{(\linewidth - 14\tabcolsep) * \real{0.1250}}
  >{\raggedright\arraybackslash}p{(\linewidth - 14\tabcolsep) * \real{0.1250}}
  >{\raggedright\arraybackslash}p{(\linewidth - 14\tabcolsep) * \real{0.1250}}
  >{\raggedright\arraybackslash}p{(\linewidth - 14\tabcolsep) * \real{0.1250}}
  >{\raggedright\arraybackslash}p{(\linewidth - 14\tabcolsep) * \real{0.1250}}
  >{\raggedright\arraybackslash}p{(\linewidth - 14\tabcolsep) * \real{0.1250}}
  >{\raggedright\arraybackslash}p{(\linewidth - 14\tabcolsep) * \real{0.1250}}
  >{\raggedright\arraybackslash}p{(\linewidth - 14\tabcolsep) * \real{0.1250}}@{}}
\toprule\noalign{}
\begin{minipage}[b]{\linewidth}\raggedright
RUNSPEC
\end{minipage} & \begin{minipage}[b]{\linewidth}\raggedright
GRID
\end{minipage} & \begin{minipage}[b]{\linewidth}\raggedright
EDIT
\end{minipage} & \begin{minipage}[b]{\linewidth}\raggedright
PROPS
\end{minipage} & \begin{minipage}[b]{\linewidth}\raggedright
REGIONS
\end{minipage} & \begin{minipage}[b]{\linewidth}\raggedright
SOLUTION
\end{minipage} & \begin{minipage}[b]{\linewidth}\raggedright
SUMMARY
\end{minipage} & \begin{minipage}[b]{\linewidth}\raggedright
SCHEDULE
\end{minipage} \\
\midrule\noalign{}
\endhead
\bottomrule\noalign{}
\endlastfoot
\end{longtable}

\subsubsection{Description}\label{description-982}

The NEFAC keyword defines an extended network node's efficiency factor,
for when the Extended Network option has been activated by the NETWORK
keyword in the RUNSPEC section. See also the GEFAC keyword in the
SCHEDULE section that can also be used with the Extended Network option.

\begin{longtable}[]{@{}
  >{\raggedright\arraybackslash}p{(\linewidth - 6\tabcolsep) * \real{0.2500}}
  >{\raggedright\arraybackslash}p{(\linewidth - 6\tabcolsep) * \real{0.2500}}
  >{\raggedright\arraybackslash}p{(\linewidth - 6\tabcolsep) * \real{0.2500}}
  >{\raggedright\arraybackslash}p{(\linewidth - 6\tabcolsep) * \real{0.2500}}@{}}
\toprule\noalign{}
\begin{minipage}[b]{\linewidth}\raggedright
No.
\end{minipage} & \begin{minipage}[b]{\linewidth}\raggedright
Name
\end{minipage} & \begin{minipage}[b]{\linewidth}\raggedright
Description
\end{minipage} & \begin{minipage}[b]{\linewidth}\raggedright
Default
\end{minipage} \\
\midrule\noalign{}
\endhead
\bottomrule\noalign{}
\endlastfoot
Field & Metric & Laboratory & \\
1 & NODE & A character string of up to eight characters in length that
defines the node name for which the node efficiency factor is being
defined. & None \\
2 & FACTOR & A real positive value that is less than or equal to one
that defines the efficiency factor for the node. If a node's down time
is 5\% then FACTOR should be set to 0.95 (1.0 -- 0.05). & 1.0 \\
dimensionless & dimensionless & dimensionless & \\
Notes: & & & \\
\end{longtable}

\emph{Table 12.3.159.1: NEFAC Keyword Description}

See also the WEFAC keyword in the SCHEDULE section to define a well's'
efficiency factor.

\subsection{NETBALAN -- Network Balancing
Parameters}\label{netbalan-network-balancing-parameters}

\begin{longtable}[]{@{}
  >{\raggedright\arraybackslash}p{(\linewidth - 14\tabcolsep) * \real{0.1250}}
  >{\raggedright\arraybackslash}p{(\linewidth - 14\tabcolsep) * \real{0.1250}}
  >{\raggedright\arraybackslash}p{(\linewidth - 14\tabcolsep) * \real{0.1250}}
  >{\raggedright\arraybackslash}p{(\linewidth - 14\tabcolsep) * \real{0.1250}}
  >{\raggedright\arraybackslash}p{(\linewidth - 14\tabcolsep) * \real{0.1250}}
  >{\raggedright\arraybackslash}p{(\linewidth - 14\tabcolsep) * \real{0.1250}}
  >{\raggedright\arraybackslash}p{(\linewidth - 14\tabcolsep) * \real{0.1250}}
  >{\raggedright\arraybackslash}p{(\linewidth - 14\tabcolsep) * \real{0.1250}}@{}}
\toprule\noalign{}
\begin{minipage}[b]{\linewidth}\raggedright
RUNSPEC
\end{minipage} & \begin{minipage}[b]{\linewidth}\raggedright
GRID
\end{minipage} & \begin{minipage}[b]{\linewidth}\raggedright
EDIT
\end{minipage} & \begin{minipage}[b]{\linewidth}\raggedright
PROPS
\end{minipage} & \begin{minipage}[b]{\linewidth}\raggedright
REGIONS
\end{minipage} & \begin{minipage}[b]{\linewidth}\raggedright
SOLUTION
\end{minipage} & \begin{minipage}[b]{\linewidth}\raggedright
SUMMARY
\end{minipage} & \begin{minipage}[b]{\linewidth}\raggedright
SCHEDULE
\end{minipage} \\
\midrule\noalign{}
\endhead
\bottomrule\noalign{}
\endlastfoot
\end{longtable}

\subsubsection{Description}\label{description-983}

NETBALAN keyword causes the simulator to perform a network balancing
operation at the next time step. In addition, the keyword defines the
network balancing parameters used to control how network balancing is
performed on a network, as well as the frequency of the network
balancing calculations. If any of the parameters on the keyword are
defaulted, then either the previously entered values are used, or if
there are no previously entered values, then the keyword default values
are employed.

\begin{longtable}[]{@{}
  >{\raggedright\arraybackslash}p{(\linewidth - 6\tabcolsep) * \real{0.2500}}
  >{\raggedright\arraybackslash}p{(\linewidth - 6\tabcolsep) * \real{0.2500}}
  >{\raggedright\arraybackslash}p{(\linewidth - 6\tabcolsep) * \real{0.2500}}
  >{\raggedright\arraybackslash}p{(\linewidth - 6\tabcolsep) * \real{0.2500}}@{}}
\toprule\noalign{}
\begin{minipage}[b]{\linewidth}\raggedright
No.
\end{minipage} & \begin{minipage}[b]{\linewidth}\raggedright
Name
\end{minipage} & \begin{minipage}[b]{\linewidth}\raggedright
Description
\end{minipage} & \begin{minipage}[b]{\linewidth}\raggedright
Default
\end{minipage} \\
\midrule\noalign{}
\endhead
\bottomrule\noalign{}
\endlastfoot
Field & Metric & Laboratory & \\
1 & NTSTEP & NSTEP is a real value that defines the criteria for the
network balancing interval, used to define the frequency of the network
balancing algorithm. NSTEP may be a negative number, a value of zero, or
a positive number, as described below: Note only negative or zero values
are currently supported. & 0.0 \\
days or dimensionless & days or dimensionless & hours or dimensionless
& \\
2 & NCNV & NCNV is a real positive value that defines the nodal pressure
convergence variance, used to define if convergence has been satisfied.
& Defined \\
psia 1.45 & barsa 0.10 & atma 0.09869 & \\
3 & NMXITER & A positive integer value that defines the maximum number
of network balancing operations that may be performed during a network
balancing computation. & 10 \\
dimensionless & dimensionless & dimensionless & \\
4 & GRPCNV & GRPCNV is a real positive value less than one that sets the
group tolerance criteria used to define if convergence has been
satisfied for: In both cases GRPCNV is an acceptable fraction of the
group's rate target. Thus, the default value 0.01 means that network
balance calculated rate must be within 0.01 (or 1.0\%) of the groups'
production target. Note that this criteria may not be satisfied if the
number of Newton iterations used in updating the well targets, as set by
the NUPCOL keyword in the RUNSPEC section, or the NUPCOL parameter on
the GCONTOL keyword in the SCHEDULE section, is exceeded. Group
convergence criteria can also be set by the GRPCNV parameter on the
GCONTOL keyword in the SCHEDULE section. If both values of GRPCNV have
been entered, then the minimum of the two is used. & 0.01 \\
dimensionless & dimensionless & dimensionless & \\
5 & THPMXITE & A positive integer value that defines the maximum number
of well THP iterations for wells in subsea completion.anifold groups
under rate control. There is no equivalent iteration maximum for
automatic chokes in groups under rate control, as the pressure drop
across the chokes is automatically calculated in the network balancing
calculation. Thus, NMXITER is used for these calculations. & 10 \\
dimensionless & dimensionless & dimensionless & \\
6 & NTRGERR & A real positive value that stipulates the maximum target
branch error in a network balance calculation, at the end of the time
step. NTRGERR is compared with the branch error (Error), where Error is
the difference between the pressure drop along the branch from the
previous network balance and the current calculation, again, at the end
of the time step. Subject to existing targets and constraints, the
simulator will attempt to adjust the time step size in order to roughly
match NTRGERR. Note also that NTRGERR should be a value that is
sufficient to allow convergence as defined by NCNV and GRPNCN
parameters. The default value of 1.0 x 1020 means that this parameter
has no effect in the selection of the time step size. & 1.0 x 1020 \\
psia & barsa & atma & \\
7 & NMAXERR & A real positive value that stipulates the maximum branch
error in a network balance calculation, at the end of the time step.
Again, NMAXERR is compared against the branch error (Error), where Error
is the difference between the pressure drop along the branch from the
previous network balance and the current calculation, at the end of the
time step. If Error is greater than NMAXERR then the simulator will
enforce a time step chop. Time step chops are computationally expensive
and should therefore be minimized. Hence, care should be used in setting
this value, which should be significantly greater than NTRGERR. The
default value of 1.0 x 1020 means that this parameter has no effect in
the selection of the time step size. & 1.0 x 1020 \\
psia & barsa & atma & \\
8 & NTSMIN & NTSMIN is a real positive value that sets the minimum time
step size for when NTRGERR and NMAXERR have values. Large network
balancing errors (Error) can take place for various reasons, a common
occurrence is for wells producing near their operating limit which may
shut in during a network balance operation. If NTRGERR and NMAXERR have
values, then this will result in the simulator enforcing time step chops
to perhaps very small time steps. NTSMIN can therefore be used to set
the minimum time step size under, and only, these circumstances. The
default value, TSMINZ, is taken from TSMINZ parameter on the TUNING
keyword in the SCHEDULE section. & TSMINZ \\
days & days & hour & \\
Notes: & & & \\
\end{longtable}

\emph{Table 12.53: NETBALAN Keyword Description}

\subsubsection{Examples}\label{examples-80}

The first example sets network balancing to occur at the beginning of
each time step using a 1.0 psia convergence criteria, and a maximum of
10 iterations. All the other parameters are defaulted.

\begin{verbatim}
--
--       NETWORK BALANCE CONTROL OPTIONS
--
--       BAL   CONV  MAX   WTHP  WTHP  ERR   ERR  TSTEP
--       INTV  TOL   ITRS  TOL   ITS   TARG  MAX  MIN
NETBALAN
         0.0   1.0   10    1*    1*    1*    1*   1*       /
\end{verbatim}

The next example sets network balancing to occur for the first NUPCOL
Newton iterations for each time step. All the other parameters are
defaulted.

\begin{verbatim}
--
--       NETWORK BALANCE CONTROL OPTIONS
--
--       BAL   CONV  MAX   WTHP  WTHP  ERR   ERR  TSTEP
--       INTV  TOL   ITRS  TOL   ITS   TARG  MAX  MIN
NETBALAN
         -1    1*    1*    1*    1*    1*    1*   1*       /
\end{verbatim}

\subsection{NETCOMPA -- Define Automatic Compressors (Extended
Network)}\label{netcompa-define-automatic-compressors-extended-network}

\begin{longtable}[]{@{}
  >{\raggedright\arraybackslash}p{(\linewidth - 14\tabcolsep) * \real{0.1250}}
  >{\raggedright\arraybackslash}p{(\linewidth - 14\tabcolsep) * \real{0.1250}}
  >{\raggedright\arraybackslash}p{(\linewidth - 14\tabcolsep) * \real{0.1250}}
  >{\raggedright\arraybackslash}p{(\linewidth - 14\tabcolsep) * \real{0.1250}}
  >{\raggedright\arraybackslash}p{(\linewidth - 14\tabcolsep) * \real{0.1250}}
  >{\raggedright\arraybackslash}p{(\linewidth - 14\tabcolsep) * \real{0.1250}}
  >{\raggedright\arraybackslash}p{(\linewidth - 14\tabcolsep) * \real{0.1250}}
  >{\raggedright\arraybackslash}p{(\linewidth - 14\tabcolsep) * \real{0.1250}}@{}}
\toprule\noalign{}
\begin{minipage}[b]{\linewidth}\raggedright
RUNSPEC
\end{minipage} & \begin{minipage}[b]{\linewidth}\raggedright
GRID
\end{minipage} & \begin{minipage}[b]{\linewidth}\raggedright
EDIT
\end{minipage} & \begin{minipage}[b]{\linewidth}\raggedright
PROPS
\end{minipage} & \begin{minipage}[b]{\linewidth}\raggedright
REGIONS
\end{minipage} & \begin{minipage}[b]{\linewidth}\raggedright
SOLUTION
\end{minipage} & \begin{minipage}[b]{\linewidth}\raggedright
SUMMARY
\end{minipage} & \begin{minipage}[b]{\linewidth}\raggedright
SCHEDULE
\end{minipage} \\
\midrule\noalign{}
\endhead
\bottomrule\noalign{}
\endlastfoot
\end{longtable}

\subsubsection{Description}\label{description-984}

The NETCOMPA keyword defines automatic compressors in an extended
network, for when the Extended Network option has been activated by the
NETWORK keyword in the RUNSPEC section.

This keyword is not supported by OPM Flow but would change the results
if supported so the simulation will be stopped.

\subsection{NEXT -- Maximum Next Time Step Size (Alias for
NEXTSTEP)}\label{next-maximum-next-time-step-size-alias-for-nextstep}

\begin{longtable}[]{@{}
  >{\raggedright\arraybackslash}p{(\linewidth - 14\tabcolsep) * \real{0.1250}}
  >{\raggedright\arraybackslash}p{(\linewidth - 14\tabcolsep) * \real{0.1250}}
  >{\raggedright\arraybackslash}p{(\linewidth - 14\tabcolsep) * \real{0.1250}}
  >{\raggedright\arraybackslash}p{(\linewidth - 14\tabcolsep) * \real{0.1250}}
  >{\raggedright\arraybackslash}p{(\linewidth - 14\tabcolsep) * \real{0.1250}}
  >{\raggedright\arraybackslash}p{(\linewidth - 14\tabcolsep) * \real{0.1250}}
  >{\raggedright\arraybackslash}p{(\linewidth - 14\tabcolsep) * \real{0.1250}}
  >{\raggedright\arraybackslash}p{(\linewidth - 14\tabcolsep) * \real{0.1250}}@{}}
\toprule\noalign{}
\begin{minipage}[b]{\linewidth}\raggedright
RUNSPEC
\end{minipage} & \begin{minipage}[b]{\linewidth}\raggedright
GRID
\end{minipage} & \begin{minipage}[b]{\linewidth}\raggedright
EDIT
\end{minipage} & \begin{minipage}[b]{\linewidth}\raggedright
PROPS
\end{minipage} & \begin{minipage}[b]{\linewidth}\raggedright
REGIONS
\end{minipage} & \begin{minipage}[b]{\linewidth}\raggedright
SOLUTION
\end{minipage} & \begin{minipage}[b]{\linewidth}\raggedright
SUMMARY
\end{minipage} & \begin{minipage}[b]{\linewidth}\raggedright
SCHEDULE
\end{minipage} \\
\midrule\noalign{}
\endhead
\bottomrule\noalign{}
\endlastfoot
\end{longtable}

\subsubsection{Description}\label{description-985}

This keyword defines the maximum time step size the simulator should
take for the next time step. This keyword can be used to reset the time
step for when known large changes to the model are taking place that may
result in time step chops. For example, if the reporting time size is
using monthly reporting steps via the DATES keyword in the SCHEDULE
section, then if for example, a group of wells start production at a
given date, then the NEXT keyword can be used to shorten the next step
in order to avoid a time step chop.

Time steps chops are computationally expensive as the simulator cannot
solve the current time step at the given tolerance, and therefore has to
reduce the time step size. For example, if the previous completed time
step was at day 365 and the current time step ending at 396 days cannot
be solved, then the simulator will reduce the current time step to
perhaps end at day 370, if this still cannot solved then the time step
will be chopped back again to perhaps to less then one day. Using the
NEXT or NEXTSTEP keyword, the simulator is instructed to take a small
time step in the anticipation that this will avoid time step chops and
thus improve computational performance.

\begin{longtable}[]{@{}
  >{\raggedright\arraybackslash}p{(\linewidth - 6\tabcolsep) * \real{0.2500}}
  >{\raggedright\arraybackslash}p{(\linewidth - 6\tabcolsep) * \real{0.2500}}
  >{\raggedright\arraybackslash}p{(\linewidth - 6\tabcolsep) * \real{0.2500}}
  >{\raggedright\arraybackslash}p{(\linewidth - 6\tabcolsep) * \real{0.2500}}@{}}
\toprule\noalign{}
\begin{minipage}[b]{\linewidth}\raggedright
No.
\end{minipage} & \begin{minipage}[b]{\linewidth}\raggedright
Name
\end{minipage} & \begin{minipage}[b]{\linewidth}\raggedright
Description
\end{minipage} & \begin{minipage}[b]{\linewidth}\raggedright
Default
\end{minipage} \\
\midrule\noalign{}
\endhead
\bottomrule\noalign{}
\endlastfoot
Field & Metric & Laboratory & \\
1 & NSTEP1 & NSTEP1 is a real positive value that defines the maximum
length of the next time step. & None \\
days & days & hours & \\
2 & NSTEP2 & NSTEP2 is a character string that should be set to either
NO or YES to state if the NSTEP1 should be applied to future reporting
time steps. The default value of NO means that NSTEP1 will only be
applied once. & NO \\
Notes: & & & \\
\end{longtable}

\emph{Table 12.54: NEXT Keyword Description}

See also the DATES and TSTEP keywords in the RUNSPEC section that are
used to advance the simulation through time.

This keyword is not supported by OPM Flow but has no effect on the
results so it will be ignored. See section 2.2 Running OPM Flow 2023-04
From The Command Line on how to control time stepping for OPM Flow.

\subsubsection{Example}\label{example-543}

\begin{verbatim}
--       NEXT   ALL
--       STEP   TIME
--       ----   ----
NEXT
         1     'NO'                                        /
\end{verbatim}

Here the next step size is set to one day and should only be used once.

\subsection{NEXTSTEP -- Maximum Next Time Step
Size}\label{nextstep-maximum-next-time-step-size}

\begin{longtable}[]{@{}
  >{\raggedright\arraybackslash}p{(\linewidth - 14\tabcolsep) * \real{0.1250}}
  >{\raggedright\arraybackslash}p{(\linewidth - 14\tabcolsep) * \real{0.1250}}
  >{\raggedright\arraybackslash}p{(\linewidth - 14\tabcolsep) * \real{0.1250}}
  >{\raggedright\arraybackslash}p{(\linewidth - 14\tabcolsep) * \real{0.1250}}
  >{\raggedright\arraybackslash}p{(\linewidth - 14\tabcolsep) * \real{0.1250}}
  >{\raggedright\arraybackslash}p{(\linewidth - 14\tabcolsep) * \real{0.1250}}
  >{\raggedright\arraybackslash}p{(\linewidth - 14\tabcolsep) * \real{0.1250}}
  >{\raggedright\arraybackslash}p{(\linewidth - 14\tabcolsep) * \real{0.1250}}@{}}
\toprule\noalign{}
\begin{minipage}[b]{\linewidth}\raggedright
RUNSPEC
\end{minipage} & \begin{minipage}[b]{\linewidth}\raggedright
GRID
\end{minipage} & \begin{minipage}[b]{\linewidth}\raggedright
EDIT
\end{minipage} & \begin{minipage}[b]{\linewidth}\raggedright
PROPS
\end{minipage} & \begin{minipage}[b]{\linewidth}\raggedright
REGIONS
\end{minipage} & \begin{minipage}[b]{\linewidth}\raggedright
SOLUTION
\end{minipage} & \begin{minipage}[b]{\linewidth}\raggedright
SUMMARY
\end{minipage} & \begin{minipage}[b]{\linewidth}\raggedright
SCHEDULE
\end{minipage} \\
\midrule\noalign{}
\endhead
\bottomrule\noalign{}
\endlastfoot
\end{longtable}

\subsubsection{Description}\label{description-986}

This keyword defines the maximum time step size the simulator should
take for the next time step. This keyword can be used to limit the time
step size when known large changes to the model are taking place that
may otherwise result in time step chops. For example, if the simulator
is using monthly reporting steps defined by the DATES keyword in the
SCHEDULE section, then if a group of wells start production at a given
date, then the NEXTSTEP keyword can be used to limit the next time step
in order to avoid a time step chop.

Time steps chops are computationally expensive as the simulator cannot
solve the current time step at the given tolerance, and therefore has to
reduce the time step size. For example, if the previous completed time
step was at day 365 and the current time step ending at 396 days cannot
be solved, then the simulator will reduce the current time step to
perhaps end at day 370, if this still cannot solved then the time step
will be chopped back again and so on, possibly reducing the time step to
less then one day. Using the NEXT or NEXTSTEP keyword, the simulator is
instructed to take a small time step in the anticipation that this will
avoid time step chops and thus improve computational performance.

\begin{longtable}[]{@{}
  >{\raggedright\arraybackslash}p{(\linewidth - 6\tabcolsep) * \real{0.2500}}
  >{\raggedright\arraybackslash}p{(\linewidth - 6\tabcolsep) * \real{0.2500}}
  >{\raggedright\arraybackslash}p{(\linewidth - 6\tabcolsep) * \real{0.2500}}
  >{\raggedright\arraybackslash}p{(\linewidth - 6\tabcolsep) * \real{0.2500}}@{}}
\toprule\noalign{}
\begin{minipage}[b]{\linewidth}\raggedright
No.
\end{minipage} & \begin{minipage}[b]{\linewidth}\raggedright
Name
\end{minipage} & \begin{minipage}[b]{\linewidth}\raggedright
Description
\end{minipage} & \begin{minipage}[b]{\linewidth}\raggedright
Default
\end{minipage} \\
\midrule\noalign{}
\endhead
\bottomrule\noalign{}
\endlastfoot
Field & Metric & Laboratory & \\
1 & NSTEP1 & A real positive value that defines the maximum length of
the next time step. & None \\
days & days & hours & \\
2 & NSTEP2 & A defined character string that should be set to either NO
or YES to state if the NSTEP1 should be applied to future reporting time
steps. The default value of NO means that NSTEP1 will only be applied
once. & NO \\
Notes: & & & \\
\end{longtable}

\emph{Table 12.55: NEXTSTEP Keyword Description}

See also the DATES and TSTEP keywords in the RUNSPEC section that are
used to advance the simulation through time.

See section 2.2 Running OPM Flow 2023-04 From The Command Line on how to
control time stepping for OPM Flow that allows for greater flexibility.

\subsubsection{Examples}\label{examples-81}

The first example shows the direct use of the NEXTSTEP keyword. Here the
next time step size is limited to one day for the first time step
following the next report time only.

\begin{verbatim}
--       NEXT   ALL
--       STEP   TIME
--       ----   ----
NEXTSTEP
         1     'NO'                                        /
\end{verbatim}

The next example shows a more complete use of the keyword for when the
field oil production has increased dramatically from 10,000 stb/d to
50,000 stb/d as indicated by the two GCONPROD keywords.

\begin{verbatim}
-- ------------------------------------------------------------------------------
-- SCHEDULE SECTION - 2021-01-01
-- ------------------------------------------------------------------------------
--
-- GROUP PRODUCTION CONTROLS
--
-- GRUP    CNTL  OIL    WAT    GAS    LIQ    CNTL  GRUP  GUIDE  GUIDE  CNTL
-- NAME    MODE  RATE   RATE   RATE   RATE   OPT   CNTL  RATE   DEF    WAT
GCONPROD
'FIELD'   'ORAT' 10E3   60E3   300E3  60E3   1*     1*    1*     1*     1*     /
/
RPTSCHED
'WELLS=2'    'WELSPECS'    'CPU=2'     'FIP=2'                                 /

DATES
 2  JAN   2021  /
/
RPTSCHED
'NOTHING'                                                                      /

DATES
 1  FEB   2021  /
 1  MAR   2021  /
/
-- GROUP PRODUCTION CONTROLS
--
-- GRUP    CNTL  OIL    WAT    GAS    LIQ    CNTL  GRUP  GUIDE  GUIDE  CNTL
-- NAME    MODE  RATE   RATE   RATE   RATE   OPT   CNTL  RATE   DEF    WAT
GCONPROD
'FIELD'   'ORAT' 50E3   90E3   300E3  90E3   1*     1*    1*     1*     1*     /
/
--
--       NEXT   ALL
--       STEP   TIME
--       ----   ----
NEXTSTEP
         1     'NO'   /
DATES
 1  APR   2021  /
 1  MAY   2021  /
 1  JUN   2021  /
 1  JLY   2021  /
 1  AUG   2021  /
 1  SEP   2021  /
 1  OCT   2021  /
 1  NOV   2021  /
 1  DEC   2021  /
/
\end{verbatim}

Given a start date of January 1, 2020 set via the START keyword in the
RUNSPEC section, the above example shows the initial oil production of
10,000 stb/d starting in January 1, 2020, and continuing up to March 1,
2021. At the March 1, 2021 report time the field oil production rate is
increased to 50,000 stb/d and the maximum next time step is set to one
day. After the one day time step is completed (March 2, 2012), the
simulator will progressively in increase the time step size until a
maximum of 31 days is reached. The 31 day maximum is a result of
requesting monthly time steps via the DATES keyword. The intent of using
the NEXTSTEP keyword in this case is to prevent time step chops
occurring due to the ``shock'' to the system caused by the large
increase in oil production.

\subsection{NEXTSTPL -- Maximum Next Time Step Size
(LGR)}\label{nextstpl-maximum-next-time-step-size-lgr}

\begin{longtable}[]{@{}
  >{\raggedright\arraybackslash}p{(\linewidth - 14\tabcolsep) * \real{0.1250}}
  >{\raggedright\arraybackslash}p{(\linewidth - 14\tabcolsep) * \real{0.1250}}
  >{\raggedright\arraybackslash}p{(\linewidth - 14\tabcolsep) * \real{0.1250}}
  >{\raggedright\arraybackslash}p{(\linewidth - 14\tabcolsep) * \real{0.1250}}
  >{\raggedright\arraybackslash}p{(\linewidth - 14\tabcolsep) * \real{0.1250}}
  >{\raggedright\arraybackslash}p{(\linewidth - 14\tabcolsep) * \real{0.1250}}
  >{\raggedright\arraybackslash}p{(\linewidth - 14\tabcolsep) * \real{0.1250}}
  >{\raggedright\arraybackslash}p{(\linewidth - 14\tabcolsep) * \real{0.1250}}@{}}
\toprule\noalign{}
\begin{minipage}[b]{\linewidth}\raggedright
RUNSPEC
\end{minipage} & \begin{minipage}[b]{\linewidth}\raggedright
GRID
\end{minipage} & \begin{minipage}[b]{\linewidth}\raggedright
EDIT
\end{minipage} & \begin{minipage}[b]{\linewidth}\raggedright
PROPS
\end{minipage} & \begin{minipage}[b]{\linewidth}\raggedright
REGIONS
\end{minipage} & \begin{minipage}[b]{\linewidth}\raggedright
SOLUTION
\end{minipage} & \begin{minipage}[b]{\linewidth}\raggedright
SUMMARY
\end{minipage} & \begin{minipage}[b]{\linewidth}\raggedright
SCHEDULE
\end{minipage} \\
\midrule\noalign{}
\endhead
\bottomrule\noalign{}
\endlastfoot
\end{longtable}

\subsubsection{Description}\label{description-987}

This keyword defines the maximum time step size the simulator should
take for the next time step for all Local Grid Refinements (``LGR'').
This keyword can be used to reset the time step for when known large
changes to the model are taking place that may result in time step
chops. For example, if the reporting time size is using monthly
reporting steps via the DATES keyword in the SCHEDULE section, then if
for example, a group of wells start production at a given date, then the
NEXTSTPL keyword can be used to shorten the next step in all the LGRs in
order to avoid a time step chop.

\begin{longtable}[]{@{}
  >{\raggedright\arraybackslash}p{(\linewidth - 6\tabcolsep) * \real{0.2500}}
  >{\raggedright\arraybackslash}p{(\linewidth - 6\tabcolsep) * \real{0.2500}}
  >{\raggedright\arraybackslash}p{(\linewidth - 6\tabcolsep) * \real{0.2500}}
  >{\raggedright\arraybackslash}p{(\linewidth - 6\tabcolsep) * \real{0.2500}}@{}}
\toprule\noalign{}
\begin{minipage}[b]{\linewidth}\raggedright
No.
\end{minipage} & \begin{minipage}[b]{\linewidth}\raggedright
Name
\end{minipage} & \begin{minipage}[b]{\linewidth}\raggedright
Description
\end{minipage} & \begin{minipage}[b]{\linewidth}\raggedright
Default
\end{minipage} \\
\midrule\noalign{}
\endhead
\bottomrule\noalign{}
\endlastfoot
Field & Metric & Laboratory & \\
1 & NSTEP1 & NSTEP1 is a real positive value that defines the maximum
length of the next time step. & None \\
days & days & hours & \\
2 & NSTEP2 & NSTEP2 is a character string that should be set to either
NO or YES to state if the NSTEP1 should be applied to future reporting
time steps. The default value of NO means that NSTEP1 will only be
applied once. & NO \\
Notes: & & & \\
\end{longtable}

\emph{Table 12.56: NEXTSTPL Keyword Description}

See also the NEXT and NEXTSTEP keywords in the SCHEDULE section that are
used to control the global grid's next time step.

This keyword is not supported by OPM Flow but would change the results
if supported so the simulation will be stopped. See section 2.2 Running
OPM Flow 2023-04 From The Command Line on how to control time stepping
for OPM Flow.

\subsubsection{Example}\label{example-544}

\begin{verbatim}
--
--       NEXT   ALL
--       STEP   TIME
--       ----   ----
NEXTSTEP
         1     'NO'                                        /
\end{verbatim}

Here the next step size for all LGRs is set to one day and should only
be used once.

\subsection{NODEPROP -- Define Network Node Properties for Extended
Network}\label{nodeprop-define-network-node-properties-for-extended-network}

\begin{longtable}[]{@{}
  >{\raggedright\arraybackslash}p{(\linewidth - 14\tabcolsep) * \real{0.1250}}
  >{\raggedright\arraybackslash}p{(\linewidth - 14\tabcolsep) * \real{0.1250}}
  >{\raggedright\arraybackslash}p{(\linewidth - 14\tabcolsep) * \real{0.1250}}
  >{\raggedright\arraybackslash}p{(\linewidth - 14\tabcolsep) * \real{0.1250}}
  >{\raggedright\arraybackslash}p{(\linewidth - 14\tabcolsep) * \real{0.1250}}
  >{\raggedright\arraybackslash}p{(\linewidth - 14\tabcolsep) * \real{0.1250}}
  >{\raggedright\arraybackslash}p{(\linewidth - 14\tabcolsep) * \real{0.1250}}
  >{\raggedright\arraybackslash}p{(\linewidth - 14\tabcolsep) * \real{0.1250}}@{}}
\toprule\noalign{}
\begin{minipage}[b]{\linewidth}\raggedright
RUNSPEC
\end{minipage} & \begin{minipage}[b]{\linewidth}\raggedright
GRID
\end{minipage} & \begin{minipage}[b]{\linewidth}\raggedright
EDIT
\end{minipage} & \begin{minipage}[b]{\linewidth}\raggedright
PROPS
\end{minipage} & \begin{minipage}[b]{\linewidth}\raggedright
REGIONS
\end{minipage} & \begin{minipage}[b]{\linewidth}\raggedright
SOLUTION
\end{minipage} & \begin{minipage}[b]{\linewidth}\raggedright
SUMMARY
\end{minipage} & \begin{minipage}[b]{\linewidth}\raggedright
SCHEDULE
\end{minipage} \\
\midrule\noalign{}
\endhead
\bottomrule\noalign{}
\endlastfoot
\end{longtable}

\subsubsection{Description}\label{description-988}

This keyword defines the network node properties for the extended
network option for when the Extended Network Model has been invoked by
the NETWORK keyword in the RUNSPEC section. There are two types of
network facilities in the simulator, the Standard Network model, which
is defined with the GRUPNET keyword in the SCHEDULE section and the
Extended Network Model defined by the BRANPROP and NODEPROP keywords,
again in the SCHEDULE section.

For the Extended Network Model the group hierarchy can be different to
that defined by the GRUPTREE keyword; however, the bottom most nodes in
the network tree associated with wells, must be the same as that defined
by the GRUPTREE keyword.

\begin{longtable}[]{@{}
  >{\raggedright\arraybackslash}p{(\linewidth - 6\tabcolsep) * \real{0.2500}}
  >{\raggedright\arraybackslash}p{(\linewidth - 6\tabcolsep) * \real{0.2500}}
  >{\raggedright\arraybackslash}p{(\linewidth - 6\tabcolsep) * \real{0.2500}}
  >{\raggedright\arraybackslash}p{(\linewidth - 6\tabcolsep) * \real{0.2500}}@{}}
\toprule\noalign{}
\begin{minipage}[b]{\linewidth}\raggedright
No.
\end{minipage} & \begin{minipage}[b]{\linewidth}\raggedright
Name
\end{minipage} & \begin{minipage}[b]{\linewidth}\raggedright
Description
\end{minipage} & \begin{minipage}[b]{\linewidth}\raggedright
Default
\end{minipage} \\
\midrule\noalign{}
\endhead
\bottomrule\noalign{}
\endlastfoot
Field & Metric & Laboratory & \\
1 & NODE & A character string of up to eight characters in length that
defines the node name for the data on this keyword record. & None \\
2 & PRESS & A real value that sets the terminal fixed pressure for the
node, this should be set to: & 1* \\
psia & bars & atm & \\
3 & CHOKE & CHOKE is a defined character string that sets if the
downstream branch (or uptree branch, being the one closer towards the
terminal node) from this node should have the capability to choke back
the flow rate in order impose a flow constraint. Here the downstream
branch is the node furthermost from the wells. Thus for a production
network, this will be an outlet branch as the wells are exporting fluid
from the branch node. Whereas for an injection node, this is an inlet
branch as the wells are importing the injection fluid. & NO \\
4 & GASLIFT & A defined character string that sets if the associated
subordinate well's produced gas lift gas should be included in the
node's flow stream (YES), or not (NO). GASLIFT should be set to either:
If NODE does not have any subordinate wells or satellite groups (see the
GSATPROD keyword in the SCHEDULE section) directly attached to the node,
then GASLIFT should be defaulted (1*) or set to NO. The option is only
valid for producing networks & NO \\
5 & GROUP & A character string of up to eight characters in length that
defines the group for which the automatic choke will be applied in order
to match the group's target rate (the TARGET variable on the GCONPROD
keyword in the SCHEDULE section). The target rate is matched by
adjusting the pressure drops across the automatic choke. The default
value of 1* uses NODE (item one) as the group if it exists in the run.
In addition, if NODE is just connected to subordinate wells then GROUP
should also be defaulted. & 1* \\
6 & GRPNAME & GRPNAME defines the name of the source or sink group in
the commercial compositional simulator; however, the variable is not
used as sources and sinks nodes must have the same name as their
comparable groups in the Extended Network Model in both OPM Flow and the
commercial black-oil simulator. & 1* \\
7 & NETYPE & NETYPE defines the network type in the commercial
compositional simulator: PROD, WINJ or GINJ. However, the variable is
not used as only production networks are supported in the Extended
Network Model in both OPM Flow and the commercial black-oil simulator. &
1* \\
Notes: & & & \\
\end{longtable}

\emph{Table 12.57: NODEPROP Keyword Description}

See also the NETWORK keyword in the RUNSPEC section and the BRANPROP in
the SCHEDULE section.

\subsubsection{Example}\label{example-545}

Given the following Extended Network model in Figure 12.5.

\pandocbounded{\includegraphics[keepaspectratio]{chapters/images/Frame67_c40b193fe195.png}}
\pandocbounded{\includegraphics[keepaspectratio]{chapters/images/Image61_c40b193fe195.png}}

First the Extended Network model should be used invoked in the RUNSPEC
section, and then the BRANPROP keyword should be used to define the
branch network, and finally the NODEPROP keyword is used to describe the
node properties with the network.

\begin{verbatim}
-- ==============================================================================
--
-- RUNSPEC SECTION
--
-- ==============================================================================
RUNSPEC
--
--       ACTIVATE THE EXTENDED NETWORK OPTION AND DEFINE PARAMETERS
--
--       MAX.    MAX     NOT
--       NODE    LINK    USED
NETWORK
         3        2      1*                                                    /
…..…………..
..……..…..
-- ==============================================================================
--
-- SCHEDULE SECTION
--
-- ==============================================================================
SCHEDULE
--
--       EXTENDED NETWORK BRANCH PROPERTIES
--
-- DOWN  UP       VFP    VFP
-- NODE  NODE     TABLE  ALFQ
BRANPROP
B1       PLAT-A   5      1*                                                   /
C1       PLAT-A   4      1*                                                   /
/
--
--       EXTENDED NETWORK NODE PROPERTIES
--
-- NODE  NODE   CHOKE  GAS   CHOKE  SOURCE  NETWORK
-- NAME  PRESS  OPTN   LIFT  GROUP  SINK    TYPE
NODEPROP
PLAT-A   21.0   NO     NO                                                     /
B1       1*     NO     NO                                                     /
C1       1*     NO     NO                                                     /
/
\end{verbatim}

Here the main platform for the field, PLAT-A, has a fixed 21 barsa
pressure applied as an operating constraint.

\subsection{NOECHO -- Deactivate Echoing of User Input Files to the
Print
File}\label{noecho-deactivate-echoing-of-user-input-files-to-the-print-file-8}

Turns off echoing of all the input files to the print file. Note by
default echoing of the inputs files is active. but can subsequently be
switched off by the NOECHO activation keyword.

See NOECHO -- Deactivate Echoing of User Input Files to the Print File
in the GLOBAL section for a full description.

\subsection{NOHMD -- Deactivate History Match Gradient Derivative
Calculations}\label{nohmd-deactivate-history-match-gradient-derivative-calculations-1}

The NOHMD deactivates various history match gradient derivative
calculations for when the History Match Gradient option has been
activated by the HMDIMS keyword in the RUNSPEC section. The keyword
consists of a series of character strings that define which derivative
should be switch off based on the keyword that requested the derivatives
to be calculated, for example HMFAULTS keyword in the GRID section. If
an empty list is entered then all the gradient derivative calculations
previously requested are switch off. The keyword is useful for changing
from history matching runs to predication cases, as the prediction cases
will be more computationally efficient without the burden of the
gradient derivative calculations.

See NOHMD -- Deactivate History Match Gradient Derivative Calculations
in the SOLUTION section for a full description.

\subsection{NOHMO -- Deactivate History Match Gradient Derivative
Calculations
(Alias)}\label{nohmo-deactivate-history-match-gradient-derivative-calculations-alias-1}

The NOHMO deactivates various history match gradient derivative
calculations for when the History Match Gradient option has been
activated by the HMDIMS keyword in the RUNSPEC section. The keyword
consists of a series of character strings that define which derivative
should be switch off based on the keyword that requested the derivatives
to be calculated, for example HMFAULTS keyword in the GRID section. If
an empty list is entered then all the gradient derivative calculations
previously requested are switch off. The keyword is useful for changing
from history matching runs to predication cases, as the prediction cases
will be more computationally efficient without the burden of the
gradient derivative calculations.

See NOHMO -- Deactivate History Match Gradient Derivative Calculations
(Alias) in the SOLUTION section for a full description.

\subsection{NOSIM -- Activate the No Simulation Mode for Data File
Checking}\label{nosim-activate-the-no-simulation-mode-for-data-file-checking-1}

NOSIM switches the mode of OPM Flow to data input checking mode. In this
mode the input file is read and all messages and print instructions are
sent to the respective output files. The SCHEDULE section is read but
the simulation is not performed.

See NOSIM -- Activate the No Simulation Mode for Data File Checking in
the RUNSPEC section for a full description

\subsection{NOWARN -- Deactivate Warning
Messages}\label{nowarn-deactivate-warning-messages-8}

Turns off warning messages to be printed to the print file; note that
this keyword is deactivated by default and can subsequently be switched
off by the WARN activation keyword. The warning messages may be turned
on and off using keywords WARN and NOWARN.

See NOWARN -- Deactivate Warning Messages in the GLOBAL section for a
full description.

\subsection{NUPCOL -- Define the Number of Newton Iterations Used to
Update Well
Targets}\label{nupcol-define-the-number-of-newton-iterations-used-to-update-well-targets}

The NUPCOL keyword defines the maximum number of Newton iterations
within a time step that may be used to update the well production and
injection targets, after which the well targets will be frozen until the
time step calculations have converged and the time step is complete.

See NUPCOL -- Define the Maximum Number of Newton Iterations Used to
Update Well Targets in the RUNSPEC section for a full description and
also section 2.2 Running OPM Flow 2023-04 From The Command Line on how
to set various other numerical control parameters for OPM Flow.

\subsection{NWATREM -- Node Water Removal (Extended
Network)}\label{nwatrem-node-water-removal-extended-network}

\begin{longtable}[]{@{}
  >{\raggedright\arraybackslash}p{(\linewidth - 14\tabcolsep) * \real{0.1250}}
  >{\raggedright\arraybackslash}p{(\linewidth - 14\tabcolsep) * \real{0.1250}}
  >{\raggedright\arraybackslash}p{(\linewidth - 14\tabcolsep) * \real{0.1250}}
  >{\raggedright\arraybackslash}p{(\linewidth - 14\tabcolsep) * \real{0.1250}}
  >{\raggedright\arraybackslash}p{(\linewidth - 14\tabcolsep) * \real{0.1250}}
  >{\raggedright\arraybackslash}p{(\linewidth - 14\tabcolsep) * \real{0.1250}}
  >{\raggedright\arraybackslash}p{(\linewidth - 14\tabcolsep) * \real{0.1250}}
  >{\raggedright\arraybackslash}p{(\linewidth - 14\tabcolsep) * \real{0.1250}}@{}}
\toprule\noalign{}
\begin{minipage}[b]{\linewidth}\raggedright
RUNSPEC
\end{minipage} & \begin{minipage}[b]{\linewidth}\raggedright
GRID
\end{minipage} & \begin{minipage}[b]{\linewidth}\raggedright
EDIT
\end{minipage} & \begin{minipage}[b]{\linewidth}\raggedright
PROPS
\end{minipage} & \begin{minipage}[b]{\linewidth}\raggedright
REGIONS
\end{minipage} & \begin{minipage}[b]{\linewidth}\raggedright
SOLUTION
\end{minipage} & \begin{minipage}[b]{\linewidth}\raggedright
SUMMARY
\end{minipage} & \begin{minipage}[b]{\linewidth}\raggedright
SCHEDULE
\end{minipage} \\
\midrule\noalign{}
\endhead
\bottomrule\noalign{}
\endlastfoot
\end{longtable}

\subsubsection{Description}\label{description-989}

The NWATREM keyword defines an extended network node as a point where
water is removed from the network, for when the Extended Network option
has been activated by the NETWORK keyword in the RUNSPEC section. The
water to be removed can be specified as a rate or as a fraction of the
total volume passing through the node.

This keyword is not supported by OPM Flow but would change the results
if supported so the simulation will be stopped.

\subsection{OUTSOL -- Define Data to be Written to the RESTART File
(Retired)}\label{outsol-define-data-to-be-written-to-the-restart-file-retired-1}

This keyword defines the data and frequency of the data to be written to
the RESTART file at each requested restart point. The keyword has been
replaced by the RPTRST keyword in the SOLUTION and SCHEDULE sections and
is therefore considered retired.

See OUTSOL -- Define Data to be Written to the RESTART File (Retired) in
the SOLUTION section for a full description.

\subsection{PICOND -- Define the Generalized Pseudo Pressure
Parameters}\label{picond-define-the-generalized-pseudo-pressure-parameters}

\begin{longtable}[]{@{}
  >{\raggedright\arraybackslash}p{(\linewidth - 14\tabcolsep) * \real{0.1250}}
  >{\raggedright\arraybackslash}p{(\linewidth - 14\tabcolsep) * \real{0.1250}}
  >{\raggedright\arraybackslash}p{(\linewidth - 14\tabcolsep) * \real{0.1250}}
  >{\raggedright\arraybackslash}p{(\linewidth - 14\tabcolsep) * \real{0.1250}}
  >{\raggedright\arraybackslash}p{(\linewidth - 14\tabcolsep) * \real{0.1250}}
  >{\raggedright\arraybackslash}p{(\linewidth - 14\tabcolsep) * \real{0.1250}}
  >{\raggedright\arraybackslash}p{(\linewidth - 14\tabcolsep) * \real{0.1250}}
  >{\raggedright\arraybackslash}p{(\linewidth - 14\tabcolsep) * \real{0.1250}}@{}}
\toprule\noalign{}
\begin{minipage}[b]{\linewidth}\raggedright
RUNSPEC
\end{minipage} & \begin{minipage}[b]{\linewidth}\raggedright
GRID
\end{minipage} & \begin{minipage}[b]{\linewidth}\raggedright
EDIT
\end{minipage} & \begin{minipage}[b]{\linewidth}\raggedright
PROPS
\end{minipage} & \begin{minipage}[b]{\linewidth}\raggedright
REGIONS
\end{minipage} & \begin{minipage}[b]{\linewidth}\raggedright
SOLUTION
\end{minipage} & \begin{minipage}[b]{\linewidth}\raggedright
SUMMARY
\end{minipage} & \begin{minipage}[b]{\linewidth}\raggedright
SCHEDULE
\end{minipage} \\
\midrule\noalign{}
\endhead
\bottomrule\noalign{}
\endlastfoot
\end{longtable}

\subsubsection{Description}\label{description-990}

The PICOND keyword defines the Generalized Pseudo Pressure (``GPP'')
{[}Whitson, C. H. and Fevang, Ø. ``Generalised Pseudopressure Well
Treatment in Reservoir Simulation,'' Presented at the IBC Technical
Services Conference on Optimisation of Gas Condensate Fields, Aberdeen,
UK (June 26-27, 1997){]} and {[}Whitson, C. H. and Fevang, Ø. ``Modeling
Gas Condensate Well Deliverability,'' paper SPE 30714, SPE Reservoir
Engineering (1996) 11, No.~4, 221-230; also presented at the SPE Annual
Technical Conference and Exhibition, Dallas, Texas, USA (October 22-25,
1995).{]} parameters used in a gas condensate well connection inflow
equations. GPP accounts for both the impact of condensate drop out and
compressibility in the mobility inflow term . If the keyword is absent
from the input deck then the default values are applied.

This keyword is not supported by OPM Flow but would change the results
if supported so the simulation will be stopped.

\subsection{Atgeirr Rasmussen}\label{atgeirr-rasmussen-1}

\begin{verbatim}
  2017-09-20T14:53:08
  AFR
  Not supported. (But perhaps it should be!)
 PIMULTAB
  David Baxendale
  2017-10-02T18:37:05.154000000
  DBx
  Reply to Atgeirr Rasmussen (09/20/2017, 14:53): "..."
  Are you sure, it is in the Norne deck for history matching.

  Changed.
  – Define Well Productivity Index versus Water Cut Tables
\end{verbatim}

\begin{longtable}[]{@{}
  >{\raggedright\arraybackslash}p{(\linewidth - 14\tabcolsep) * \real{0.1250}}
  >{\raggedright\arraybackslash}p{(\linewidth - 14\tabcolsep) * \real{0.1250}}
  >{\raggedright\arraybackslash}p{(\linewidth - 14\tabcolsep) * \real{0.1250}}
  >{\raggedright\arraybackslash}p{(\linewidth - 14\tabcolsep) * \real{0.1250}}
  >{\raggedright\arraybackslash}p{(\linewidth - 14\tabcolsep) * \real{0.1250}}
  >{\raggedright\arraybackslash}p{(\linewidth - 14\tabcolsep) * \real{0.1250}}
  >{\raggedright\arraybackslash}p{(\linewidth - 14\tabcolsep) * \real{0.1250}}
  >{\raggedright\arraybackslash}p{(\linewidth - 14\tabcolsep) * \real{0.1250}}@{}}
\toprule\noalign{}
\begin{minipage}[b]{\linewidth}\raggedright
RUNSPEC
\end{minipage} & \begin{minipage}[b]{\linewidth}\raggedright
GRID
\end{minipage} & \begin{minipage}[b]{\linewidth}\raggedright
EDIT
\end{minipage} & \begin{minipage}[b]{\linewidth}\raggedright
PROPS
\end{minipage} & \begin{minipage}[b]{\linewidth}\raggedright
REGIONS
\end{minipage} & \begin{minipage}[b]{\linewidth}\raggedright
SOLUTION
\end{minipage} & \begin{minipage}[b]{\linewidth}\raggedright
SUMMARY
\end{minipage} & \begin{minipage}[b]{\linewidth}\raggedright
SCHEDULE
\end{minipage} \\
\midrule\noalign{}
\endhead
\bottomrule\noalign{}
\endlastfoot
\end{longtable}

\subsubsection{Description}\label{description-991}

PIMULTAB defines productivity index multiplier versus water cut tables
that are used to scaled a well's connection factors based on
connection's current producing water cut. The tables are used for
modeling the productivity decline due to increasing water cut.
Allocation of the tables to a well is via the WPITAB keyword in the
SCHEDULE section.

\begin{longtable}[]{@{}
  >{\raggedright\arraybackslash}p{(\linewidth - 6\tabcolsep) * \real{0.2500}}
  >{\raggedright\arraybackslash}p{(\linewidth - 6\tabcolsep) * \real{0.2500}}
  >{\raggedright\arraybackslash}p{(\linewidth - 6\tabcolsep) * \real{0.2500}}
  >{\raggedright\arraybackslash}p{(\linewidth - 6\tabcolsep) * \real{0.2500}}@{}}
\toprule\noalign{}
\begin{minipage}[b]{\linewidth}\raggedright
No.
\end{minipage} & \begin{minipage}[b]{\linewidth}\raggedright
Name
\end{minipage} & \begin{minipage}[b]{\linewidth}\raggedright
Description
\end{minipage} & \begin{minipage}[b]{\linewidth}\raggedright
Default
\end{minipage} \\
\midrule\noalign{}
\endhead
\bottomrule\noalign{}
\endlastfoot
Field & Metric & Laboratory & \\
1 & WCUT & A real monotonically increasing positive columnar vector that
defines the maximum surface water cut for the corresponding PIMULT
vector. Water cut is defined
as\({f}_{w} = \frac{{q}_{w}}{{q}_{w} + {q}_{o}}\). & None \\
dimensionless & dimensionless & dimensionless & \\
2 & PIMULT & A real positive decreasing columnar vector that defines the
productivity index multiplier used to scale a well's connection factors,
for the corresponding WCUT vector. & None \\
dimensionless & dimensionless & dimensionless & \\
Notes: & & & \\
\end{longtable}

\emph{Table 12.58: PIMULTAB Keyword Description}

See also the WPITAB keyword that allocates the tables to the wells, and
also the WPIMULT keyword that scales a well's productivity index by a
constant value, both of which are in the SCHEDULE section.

This keyword is not supported by OPM Flow but would change the results
if supported so the simulation will be stopped.

\subsubsection{Examples}\label{examples-82}

Given NTPIMT equals two and NRPIMT equals four on PIMTDIMS keyword in
the RUNSPEC section, then:

\begin{verbatim}
--       DEFINE WELL PRODUCTIVITY INDEX VERSUS WATER CUT TABLES
--
--       MAX      PI
--       WCUT     MULT
--       ------   ------
PIMULTAB
         0.0000   1.0000
         0.2500   0.9500
         0.5000   0.8500
         0.7500   0.7500                                   /
--
--
         0.0000   1.0000
         0.2500   0.9500
         0.5000   0.8500
         0.7500   0.7500                                   /
\end{verbatim}

The next example is summarized from the Norne model with NTPIMT equals
one and NRPIMT equals to 51 on the PIMTDIMS keyword in the RUNSPEC
section.

\begin{verbatim}
--
--       DEFINE WELL PRODUCTIVITY INDEX VERSUS WATER CUT TABLES
--       The following is the reviewed model in Aug-2006, low-high case
--       a=0.25, b=0.1; PIMULT=(1-a)/exp(fw/b)+a
--
--       MAX      PI
--       WCUT     MULT
--       ------   ------
PIMULTAB
         0.000    1.0000
         0.025    0.8341
         0.050    0.7049
         0.075    0.6043
         0.100    0.5259
         0.125    0.4649
         0.150    0.4173
         0.175    0.3803
         0.200    0.3515
         0.225    0.3290
         0.250    0.3116
         0.275    0.2979
         0.300    0.2873
         0.325    0.2791
         0.350    0.2726
         0.375    0.2676
         0.400    0.2637
         0.425    0.2607
         0.450    0.2583
         0.475    0.2565
         0.500    0.2551
         0.525    0.2539
         0.550    0.2531
         0.575    0.2524
         0.600    0.2519
         0.625    0.2514
         0.650    0.2511
         0.675    0.2509
         0.700    0.2507
         0.725    0.2505
         0.750    0.2504
         0.775    0.2503
         0.800    0.2503
         0.825    0.2502
         0.850    0.2502
         0.875    0.2501
         0.900    0.2501
         0.925    0.2501
         0.950    0.2501
         0.975    0.2500
         1.000    0.2500 /

\end{verbatim}

\subsection{PLYADS -- Define Polymer Rock Adsorption
Tables}\label{plyads-define-polymer-rock-adsorption-tables-1}

The PLYADS keyword defines the rock polymer adsorption tables for when
the polymer option has been activated by the POLYMER keyword in the
RUNSPEC section. Alternatively, the functions can be entered via the
PLYADSS keyword in the PROPS section for when salt sensitivity is to be
considered.

See PLYADS - Define Polymer Rock Adsorption Tables in the PROPS section
for a full description.

\subsection{PLYDHFLF -- Define Polymer Thermal Degradation Half-Life
Tables}\label{plydhflf-define-polymer-thermal-degradation-half-life-tables-1}

The PLYDHFLF keyword defines the polymer thermal degradation half-life
with respect to temperature functions for when the polymer option has
been activated by the POLYMER keyword in the RUNSPEC section.

See PLYDHFLF - Define Polymer Thermal Degradation Half-Life Tables in
the PROPS section for a full description.

\subsection{PLYMAX -- Define Polymer-Salt Viscosity Mixing
Concentrations}\label{plymax-define-polymer-salt-viscosity-mixing-concentrations-1}

The PLYMAX keyword defines maximum polymer and salt concentrations that
are to be used in the mixing parameter calculation of the fluid
component viscosities, for when the polymer option has been activated by
the POLYMER keyword in the RUNSPEC section.

See PLYMAX - Define Polymer-Salt Viscosity Mixing Concentrations in the
PROPS section for a full description.

\subsection{PLYROCKM -- Modify Polymer-Rock
Properties}\label{plyrockm-modify-polymer-rock-properties}

\begin{longtable}[]{@{}
  >{\raggedright\arraybackslash}p{(\linewidth - 14\tabcolsep) * \real{0.1250}}
  >{\raggedright\arraybackslash}p{(\linewidth - 14\tabcolsep) * \real{0.1250}}
  >{\raggedright\arraybackslash}p{(\linewidth - 14\tabcolsep) * \real{0.1250}}
  >{\raggedright\arraybackslash}p{(\linewidth - 14\tabcolsep) * \real{0.1250}}
  >{\raggedright\arraybackslash}p{(\linewidth - 14\tabcolsep) * \real{0.1250}}
  >{\raggedright\arraybackslash}p{(\linewidth - 14\tabcolsep) * \real{0.1250}}
  >{\raggedright\arraybackslash}p{(\linewidth - 14\tabcolsep) * \real{0.1250}}
  >{\raggedright\arraybackslash}p{(\linewidth - 14\tabcolsep) * \real{0.1250}}@{}}
\toprule\noalign{}
\begin{minipage}[b]{\linewidth}\raggedright
RUNSPEC
\end{minipage} & \begin{minipage}[b]{\linewidth}\raggedright
GRID
\end{minipage} & \begin{minipage}[b]{\linewidth}\raggedright
EDIT
\end{minipage} & \begin{minipage}[b]{\linewidth}\raggedright
PROPS
\end{minipage} & \begin{minipage}[b]{\linewidth}\raggedright
REGIONS
\end{minipage} & \begin{minipage}[b]{\linewidth}\raggedright
SOLUTION
\end{minipage} & \begin{minipage}[b]{\linewidth}\raggedright
SUMMARY
\end{minipage} & \begin{minipage}[b]{\linewidth}\raggedright
SCHEDULE
\end{minipage} \\
\midrule\noalign{}
\endhead
\bottomrule\noalign{}
\endlastfoot
\end{longtable}

\subsubsection{Description}\label{description-992}

The PLYROCKM keyword modifies rock properties entered via the PLYCAMAX,
PLYKRRF, PLYRMDEN, and PLYROCK keywords in the PROPS section, for when
the Polymer option has been activated by the POLYMER keyword in the
RUNSPEC section.

This keyword is not supported by OPM Flow but would change the results
if supported so the simulation will be stopped.

\subsection{PLYSHEAR -- Activate and Define Polymer Shearing
Parameters}\label{plyshear-activate-and-define-polymer-shearing-parameters-1}

The PLYSHEAR keyword activates and the defines the polymer shear
thinning-thickening option for when the polymer option has been
activated by the POLYMER keyword in the RUNSPEC section.

See PLYSHEAR -- Activate and Define Polymer Shearing Parameters in the
PROPS section for a full description.

\subsection{PLYSHLOG -- Activate and Define the Polymer Shearing
Logarithmic
Parameters}\label{plyshlog-activate-and-define-the-polymer-shearing-logarithmic-parameters-1}

This keyword activates and defines the parameters for the
logarithm-based polymer shear thinning/thickening option.

See PLYSHLOG - Activate and Define the Polymer Shearing Logarithmic
Parameters in the PROPS section for a full description.

\subsection{PLYVISC -- Define Polymer Viscosity Scaling
Fact}\label{plyvisc-define-polymer-viscosity-scaling-fact}

PLYSVISC defines the polymer viscosity scaling factors used to determine
the relationship of pure water viscosity with respect to increasing
polymer saturation within a grid block. The polymer option must be
activated by the POLYMER keyword in the RUNSPEC section in order to use
this keyword.

See PLYVISC -- Define Polymer Viscosity Scaling Factors in the PROPS
section for a full description.

\subsection{PLYVISCS -- Define Polymer-Salt Viscosity Scaling
Factors}\label{plyviscs-define-polymer-salt-viscosity-scaling-factors-1}

PLYSVISCS defines the polymer-salt viscosity scaling factor tables
applied to pure water that are used to determine the viscosity of a
polymer-salt mixture with respect to increasing polymer saturation
within a grid block. The polymer option must be activated by the POLYMER
keyword, as well as the brine phase declared by the BRINE keyword in the
RUNSPEC section in order to use this keyword. However the ECLM keyword
in the RUNSPEC must not be used with this keyword.

See PLYVISCS -- Define Polymer-Salt Viscosity Scaling Factors in the
PROPS section for a full description.

\subsection{PLYVISCT -- Define Polymer-Temperature Viscosity Scaling
Factors}\label{plyvisct-define-polymer-temperature-viscosity-scaling-factors-1}

PLYSVISCT defines the polymer-temperature viscosity scaling factor
tables applied to pure water that are used to determine the viscosity of
the polymer at a given temperature with respect to increasing polymer
saturation within a grid block. Both the polymer option must be
activated by the POLYMER keyword and the temperature option invoked by
the TEMP keyword in the RUNSPEC section in order to use this keyword.
However the BRINE keyword in the RUNSPEC must not be used with this
keyword, that is the salt sensitivity options should be deactivated.

See PLYVISCT -- Define Polymer-Temperature Viscosity Scaling Factors in
the PROPS section for a full description.

\subsection{PLYVSCST -- Define Polymer-Salt-Temperature Viscosity
Scaling
Factors}\label{plyvscst-define-polymer-salt-temperature-viscosity-scaling-factors-1}

PLYVSCST defines the polymer-salt-temperature viscosity scaling factor
tables applied to pure water that are used to determine the viscosity of
the polymer at a given salt concentration and for a given temperature,
with respect to increasing polymer saturation within a grid block. Both
the polymer option must be activated by the POLYMER keyword and the
temperature option invoked by the TEMP keyword in the RUNSPEC section in
order to use this keyword. In addition, the BRINE keyword in the RUNSPEC
must also be invoked. The keyword is used in conjunction with the
SALTNODE keyword to define the various salt concentrations and the
TEMPNODE keyword to define the various reservoir temperatures. Both
keywords are in the PROPS section.

See PLYVSCST -- Define Polymer-Salt-Temperature Viscosity Scaling
Factors in the PROPS section for a full description.

\subsection{PRIORITY -- Activate and Define Well Prioritization
Coefficients}\label{priority-activate-and-define-well-prioritization-coefficients}

\begin{longtable}[]{@{}
  >{\raggedright\arraybackslash}p{(\linewidth - 14\tabcolsep) * \real{0.1250}}
  >{\raggedright\arraybackslash}p{(\linewidth - 14\tabcolsep) * \real{0.1250}}
  >{\raggedright\arraybackslash}p{(\linewidth - 14\tabcolsep) * \real{0.1250}}
  >{\raggedright\arraybackslash}p{(\linewidth - 14\tabcolsep) * \real{0.1250}}
  >{\raggedright\arraybackslash}p{(\linewidth - 14\tabcolsep) * \real{0.1250}}
  >{\raggedright\arraybackslash}p{(\linewidth - 14\tabcolsep) * \real{0.1250}}
  >{\raggedright\arraybackslash}p{(\linewidth - 14\tabcolsep) * \real{0.1250}}
  >{\raggedright\arraybackslash}p{(\linewidth - 14\tabcolsep) * \real{0.1250}}@{}}
\toprule\noalign{}
\begin{minipage}[b]{\linewidth}\raggedright
RUNSPEC
\end{minipage} & \begin{minipage}[b]{\linewidth}\raggedright
GRID
\end{minipage} & \begin{minipage}[b]{\linewidth}\raggedright
EDIT
\end{minipage} & \begin{minipage}[b]{\linewidth}\raggedright
PROPS
\end{minipage} & \begin{minipage}[b]{\linewidth}\raggedright
REGIONS
\end{minipage} & \begin{minipage}[b]{\linewidth}\raggedright
SOLUTION
\end{minipage} & \begin{minipage}[b]{\linewidth}\raggedright
SUMMARY
\end{minipage} & \begin{minipage}[b]{\linewidth}\raggedright
SCHEDULE
\end{minipage} \\
\midrule\noalign{}
\endhead
\bottomrule\noalign{}
\endlastfoot
\end{longtable}

\subsubsection{Description}\label{description-993}

The PRIORITY keyword activates the Well Priority option and defines the
coefficients in the well priority equation. Wells under group control
are ranked based on their well potential in order to satisfy group
controls. For example if a group's oil target is exceeded, then the
group may shut-in the lease productive oil wells based on their well
potential. The Priority option is an alternative form of ranking the
wells based on the following equation:

\begin{equation}\phantomsection\label{eq-12-31}{
\text{Priority} = \frac{{a}_{1} + {a}_{2}{Q}_{\mathit{oil}} +{a}_{3}{Q}_{\mathit{water}} +{a}_{4}{Q}_{\mathit{gas}}}{{b}_{1} + {b}_{2}{Q}_{\mathit{oil}} +{b}_{3}{Q}_{\mathit{water}} +{b}_{4}{Q}_{\mathit{gas}}}
}\end{equation}

Where:

Qoil = well oil potential

Qwater = well water potential

Qgas = well gas potential

a1-4 = priority coefficients supplied by this keyword

b1-4 = priority coefficients supplied by this keyword

This keyword is not supported by OPM Flow but would change the results
if supported so the simulation will be stopped.

\begin{longtable}[]{@{}
  >{\raggedright\arraybackslash}p{(\linewidth - 6\tabcolsep) * \real{0.2500}}
  >{\raggedright\arraybackslash}p{(\linewidth - 6\tabcolsep) * \real{0.2500}}
  >{\raggedright\arraybackslash}p{(\linewidth - 6\tabcolsep) * \real{0.2500}}
  >{\raggedright\arraybackslash}p{(\linewidth - 6\tabcolsep) * \real{0.2500}}@{}}
\toprule\noalign{}
\begin{minipage}[b]{\linewidth}\raggedright
No.
\end{minipage} & \begin{minipage}[b]{\linewidth}\raggedright
Name
\end{minipage} & \begin{minipage}[b]{\linewidth}\raggedright
Description
\end{minipage} & \begin{minipage}[b]{\linewidth}\raggedright
Default
\end{minipage} \\
\midrule\noalign{}
\endhead
\bottomrule\noalign{}
\endlastfoot
Field & Metric & Laboratory & \\
1 & TIME & A real positive integer that defines the minimum time
interval between executing the well priority calculation. The
calculation is performed at the beginning of the time step that exceeds
the previous calculation (t0) by a minimum of TIME, that is for when tn
≥ ( t0 + TIME). Note that the default value of zero means that the
calculation is performed at each time step. As a consequence, this may
result in some oscillation as well wells are switched on/off at
subsequent time steps. & 0 \\
days & days & hours & \\
2 & A1 & A real positive integer greater than or equal to zero that
defines a1 priority coefficient in equation (12.31). & 0 \\
3 & A2 & A real positive integer greater than or equal to zero that
defines a2 priority coefficient in equation (12.31). & 0 \\
4 & A3 & A real positive integer greater than or equal to zero that
defines a3 priority coefficient in equation (12.31). & 0 \\
5 & A4 & A real positive integer greater than or equal to zero that
defines a4 priority coefficient in equation (12.31). & 0 \\
6 & B1 & A real positive integer greater than or equal to zero that
defines b1 priority coefficient in equation (12.31). & 0 \\
7 & B2 & A real positive integer greater than or equal to zero that
defines b2 priority coefficient in equation (12.31). & 0 \\
8 & B3 & A real positive integer greater than or equal to zero that
defines b3 priority coefficient in equation (12.31). & 0 \\
9 & B4 & A real positive integer greater than or equal to zero that
defines b4 priority coefficient in equation (12.31). & 0 \\
Notes: & & & \\
\end{longtable}

\emph{Table 12.59: PRIORITY Keyword Description}

\subsubsection{Example}\label{example-546}

\begin{verbatim}
--
--       SETS COEFFICIENTS FOR WELL PRIORITIZATION OPTION
--
--       TIME   A     B     C     D     E     F     G     H
--       STEP         Qo    Qw    Qg          Qo    Qw    Qg
PRIORITY
         0.0    0.0   1.0   0.0   0.0   1.0   0.0   0.0   0.0 / High Oil Pot
--       0.0    0.0   1.0   1.0   0.0   0.0   0.0   1.0   0.0 / Low Water Cut Pot
--       0.0    0.0   1.0   0.0   0.0   0.0   0.0   0.0   1.0 / Low GOR Pot
\end{verbatim}

The above example defines the well priority calculation to be based on a
well's oil potential, with calculation to be performed at each time
step. Note that the low water cut and low GOR options are given for
reference but are commented out and therefore ignored by the simulator.

\subsection{PRORDER -- Define a Group Production Rules
Sequence}\label{prorder-define-a-group-production-rules-sequence}

\begin{longtable}[]{@{}
  >{\raggedright\arraybackslash}p{(\linewidth - 14\tabcolsep) * \real{0.1250}}
  >{\raggedright\arraybackslash}p{(\linewidth - 14\tabcolsep) * \real{0.1250}}
  >{\raggedright\arraybackslash}p{(\linewidth - 14\tabcolsep) * \real{0.1250}}
  >{\raggedright\arraybackslash}p{(\linewidth - 14\tabcolsep) * \real{0.1250}}
  >{\raggedright\arraybackslash}p{(\linewidth - 14\tabcolsep) * \real{0.1250}}
  >{\raggedright\arraybackslash}p{(\linewidth - 14\tabcolsep) * \real{0.1250}}
  >{\raggedright\arraybackslash}p{(\linewidth - 14\tabcolsep) * \real{0.1250}}
  >{\raggedright\arraybackslash}p{(\linewidth - 14\tabcolsep) * \real{0.1250}}@{}}
\toprule\noalign{}
\begin{minipage}[b]{\linewidth}\raggedright
RUNSPEC
\end{minipage} & \begin{minipage}[b]{\linewidth}\raggedright
GRID
\end{minipage} & \begin{minipage}[b]{\linewidth}\raggedright
EDIT
\end{minipage} & \begin{minipage}[b]{\linewidth}\raggedright
PROPS
\end{minipage} & \begin{minipage}[b]{\linewidth}\raggedright
REGIONS
\end{minipage} & \begin{minipage}[b]{\linewidth}\raggedright
SOLUTION
\end{minipage} & \begin{minipage}[b]{\linewidth}\raggedright
SUMMARY
\end{minipage} & \begin{minipage}[b]{\linewidth}\raggedright
SCHEDULE
\end{minipage} \\
\midrule\noalign{}
\endhead
\bottomrule\noalign{}
\endlastfoot
\end{longtable}

\subsubsection{Description}\label{description-994}

PRORDER defines the order of group production rules to be implemented
fore when a group's target is not satisfied.

This keyword is not supported by OPM Flow but would change the results
if supported so the simulation will be stopped.

\subsection{PYACTION -- Define Python Based Action Conditions and
Command
Processing}\label{pyaction-define-python-based-action-conditions-and-command-processing}

\begin{longtable}[]{@{}
  >{\raggedright\arraybackslash}p{(\linewidth - 14\tabcolsep) * \real{0.1250}}
  >{\raggedright\arraybackslash}p{(\linewidth - 14\tabcolsep) * \real{0.1250}}
  >{\raggedright\arraybackslash}p{(\linewidth - 14\tabcolsep) * \real{0.1250}}
  >{\raggedright\arraybackslash}p{(\linewidth - 14\tabcolsep) * \real{0.1250}}
  >{\raggedright\arraybackslash}p{(\linewidth - 14\tabcolsep) * \real{0.1250}}
  >{\raggedright\arraybackslash}p{(\linewidth - 14\tabcolsep) * \real{0.1250}}
  >{\raggedright\arraybackslash}p{(\linewidth - 14\tabcolsep) * \real{0.1250}}
  >{\raggedright\arraybackslash}p{(\linewidth - 14\tabcolsep) * \real{0.1250}}@{}}
\toprule\noalign{}
\begin{minipage}[b]{\linewidth}\raggedright
RUNSPEC
\end{minipage} & \begin{minipage}[b]{\linewidth}\raggedright
GRID
\end{minipage} & \begin{minipage}[b]{\linewidth}\raggedright
EDIT
\end{minipage} & \begin{minipage}[b]{\linewidth}\raggedright
PROPS
\end{minipage} & \begin{minipage}[b]{\linewidth}\raggedright
REGIONS
\end{minipage} & \begin{minipage}[b]{\linewidth}\raggedright
SOLUTION
\end{minipage} & \begin{minipage}[b]{\linewidth}\raggedright
SUMMARY
\end{minipage} & \begin{minipage}[b]{\linewidth}\raggedright
SCHEDULE
\end{minipage} \\
\midrule\noalign{}
\endhead
\bottomrule\noalign{}
\endlastfoot
\end{longtable}

\subsubsection{Description}\label{description-995}

The PYACTION keyword is part of OPM Flow's Python scripting facility
that loads a standard Python script file that can be used to define a
series of conditions and actions as the simulation proceeds through
time. The ``included'' Python script file is executed by the standard
Python interpreter. Thus, OPM Flow's Python scripting facility offers
greater flexibility compared to the commercial simulator's ACTION series
of keywords (ACTION, ACTIONG, ACTIONR, ACTIONS, ACTIONW and ACTIONX)
that can apply Boolean conditional tests to variables at the field,
group, region, well segment and well levels. Note that OPM Flow has also
implemented the commercial simulator's ACTIONX keyword, but not the
ACTION, ACTIONG, ACTIONR, ACTIONS and ACTIONW keywords, as the ACTIONX
keyword implements their functionality with greater flexibility. For a
python class documentation and for further documentation and examples,
also refer to the
\href{https://opm.github.io/opm-python-documentation/index.html}{OPM
Online Python Documentation}.

This keyword starts the definition of a PYACTION section that specifies
the name of the action and a string indicating the number of times the
action should be run; this is then followed by the name of the file
containing the Python script. Note that unlike an ACTIONX block that is
terminated by the ENDACTIO keyword the PYACTION section has no
terminating keyword.

Although this keyword is read by OPM Flow and the script processing has
been implemented, one should use caution when using this facility as it
may result in OPM Flow aborting. This is because the PYACTION facility
allows the user to implement complex functionality and the
implementation was new for the 2020-04 release. Users should therefore
use caution when using this facility.

\begin{tcolorbox}[enhanced jigsaw, title=\textcolor{quarto-callout-note-color}{\faInfo}\hspace{0.5em}{Note}, coltitle=black, breakable, toprule=.15mm, colbacktitle=quarto-callout-note-color!10!white, leftrule=.75mm, opacitybacktitle=0.6, colback=white, bottomtitle=1mm, toptitle=1mm, titlerule=0mm, left=2mm, colframe=quarto-callout-note-color-frame, arc=.35mm, rightrule=.15mm, bottomrule=.15mm, opacityback=0]

This is an OPM Flow specific keyword for the simulator's scripting
facility using the standard Python interpreter, as such it gives more
flexibility than the commercial simulator's ACTIONX keyword, although
OPM Flow also supports this as well. The keyword should be considered
experimental as details of the OPM Flow - Python interface might change
for future releases. In particular, the current implementation is quite
minimal; however, future releases are expected to add more entry points
in the Schedule class which can be used to manipulate the reservoir
model as the simulation progresses. Users are encouraged to make
suggestions for new features in this regard. The PYACTION keyword is a
very powerful keyword and allows any piece of Python code to be included
and run, including potentially malicious code. The important point is to
scrutinize any PYACTION keyword in a deck you receive from other
parties.

\end{tcolorbox}

\begin{longtable}[]{@{}
  >{\raggedright\arraybackslash}p{(\linewidth - 6\tabcolsep) * \real{0.2500}}
  >{\raggedright\arraybackslash}p{(\linewidth - 6\tabcolsep) * \real{0.2500}}
  >{\raggedright\arraybackslash}p{(\linewidth - 6\tabcolsep) * \real{0.2500}}
  >{\raggedright\arraybackslash}p{(\linewidth - 6\tabcolsep) * \real{0.2500}}@{}}
\toprule\noalign{}
\begin{minipage}[b]{\linewidth}\raggedright
No.
\end{minipage} & \begin{minipage}[b]{\linewidth}\raggedright
Name
\end{minipage} & \begin{minipage}[b]{\linewidth}\raggedright
Description
\end{minipage} & \begin{minipage}[b]{\linewidth}\raggedright
Default
\end{minipage} \\
\midrule\noalign{}
\endhead
\bottomrule\noalign{}
\endlastfoot
PYACTION & PYACTION declares the start of a PYACTION Definition Section.
This is then followed by one record that defines the name of the action
and a string indicating the number of times the action should be run;
this is then followed by a second record indicating the file containing
the Python script. & Not Applicable & \\
1-1 & ACTNAME & ACTNAME is a character sting of any length enclose in
quotes that defines the name of this action definition. & None \\
1-2 & ACTNSTEP & ACTNSTEP is a defined character string that indicates
the number of times the action should be performed, and should be set to
one of the following: Note that the FIRST\_TRUE option is only supported
for back compatibility when using the deprecated run() function style
Python scripts. & SINGLE \\
1-4 & / & Record terminated by a ``/'' & Not Applicable \\
2-1 & FILENAME & A character string enclosed in quotes that defines the
Python module/script file to read in and to be processed by OPM Flow. &
None \\
2-2 & / & Record terminated by a ``/'' & Not Applicable \\
Notes: & & & \\
\end{longtable}

\emph{Table 12.60: PYACTION Keyword Description}

The PYACTION keyword is a result of combining two programming languages,
the interactive Python interpreter and OPM Flow's source code language
C++. When combing two languages one extends and embeds one into the
other. When extending Python with C++ the functionality implemented in
C++ is made available to Python applications, when embedding Python in
C++ one can call Python functions from within C++. The PYACTION keyword
is based on embedding a Python interpreter in the C++ OPM Flow
simulator, such that Python code is run from within the simulator.

Note that to use the PYACTION keyword OPM Flow must be built with
embedded python enabled in opm-common/CMakeLists.txt.

\begin{verbatim}
option(OPM_ENABLE_EMBEDDED_PYTHON "Enable embedded python?" ON)
\end{verbatim}

In order to fully benefit from the power of the PYACTION keyword one
should familiarize oneself with the Python wrapper classes for the
various OPM Flow C++ classes. These classes are essential to share the
simulator state with the PYACTION script. That is, the function provides
the interface between the two programming languages.

To learn more about these classes you can use the pydoc utility, for
example, to learn more about the SummaryState class, in a Linux terminal
one would issue the following bash command:

\begin{verbatim}
bash% pydoc opm.io.sim.SummaryState
\end{verbatim}

The Python script file (FILENAME in the PYACTION keyword) should be a
standard Python module {[}A Python module is a file containing Python
definitions and statements. The file name is the module name with the
suffix .py appended. Within a module, the module's name (as a string) is
available as the value of the global variable \textbf{name}.{]} that
defines the PYACTION script and should consist of 100\% pure Python. The
PYACTION Python module (FILENAME) is imported during processing of the
input deck and as such this implies:

\begin{verbatim}
- Basic Python syntax checking is performed when the Python module (FILENAME) is read in.
- It is verified that the module has the correct format.
\end{verbatim}

The syntax of the Python module (FILENAME) is given in Table 12.61
together with a description of the Python module opm\_embedded.

\begin{longtable}[]{@{}
  >{\raggedright\arraybackslash}p{(\linewidth - 4\tabcolsep) * \real{0.3333}}
  >{\raggedright\arraybackslash}p{(\linewidth - 4\tabcolsep) * \real{0.3333}}
  >{\raggedright\arraybackslash}p{(\linewidth - 4\tabcolsep) * \real{0.3333}}@{}}
\toprule\noalign{}
\begin{minipage}[b]{\linewidth}\raggedright
Python Script Definition
\end{minipage} & \begin{minipage}[b]{\linewidth}\raggedright
\end{minipage} & \begin{minipage}[b]{\linewidth}\raggedright
\end{minipage} \\
\midrule\noalign{}
\endhead
\bottomrule\noalign{}
\endlastfoot
\# \# OPM Flow PYACTION Module Script \# import opm\_embedded ecl\_state
= opm\_embedded.current\_ecl\_state schedule =
opm\_embedded.current\_schedule report\_step =
opm\_embedded.current\_report\_step summary\_state =
opm\_embedded.current\_summary\_state \ldots{} & & \\
No. & Name & Description \\
1 & ecl\_state & The ecl\_state variable is an instance of the
opm.io.ecl\_state.EclipseState class which is initialized with all the
static information in the simulation, for example the 3D grid properties
and the PVT and saturation tables. This object can be read by the Python
script but the available data should not be modified as any updates will
not be passed back to the simulator. \\
2 & schedule & The schedule variable is an instance of the
opm.io.sched.Schedule class which has all the SCHEDULE information
internalized (Schedule object). For a complete documentation of the
available call points of opm.io.schedule.Schedule, we refer to
\url{https://opm.github.io/opm-python-documentation/master/common.html\#opm.io.schedule.Schedule}.
The most notable ones are: schedule.open\_well(well\_name, report\_step)
schedule.shut\_well(well\_name, report\_step)
schedule.stop\_well(well\_name, report\_step) Which can be used to open,
shut or stop a well at a particular report step (or at the current
report step if the report\_step argument is not specified), and:
schedule.insert\_keywords(kw) Which can be used to insert various
keywords where kw is a string defining the keyword and associated
records. For the 2024.10 release this has been tested with the following
keywords: BOX, COMPSEGS, ENDBOX, EXIT, FIELD, INCLUDE, GCONINJE,
GCONPROD, GCONSUMP, GEFAC, GRUPTREE, METRIC, MULTX, MULTX-, MULTY,
MULTY-, MULTZ, MULTZ-, NEXT, NEXTSTEP, WCONINJE, WCONPROD, WECON, WEFAC,
WELOPEN, WELPI, WELTARG, WGRUPCON, WELSEGS, WELSPECS, WSEGVALV, WLIST,
WPIMULT, WTEST, WTMULT This list is subject to change for future
releases. The parsing strictness for the ACTIONX and PYACTION keywords
can be set using the --action-parsing-strictness command line option. If
the default option --action-parsing-strictness=``normal'' has been set
then OPM Flow will only apply keywords that have been tested, whereas if
the option --action-parsing-strictness=``low'' has been set then OPM
Flow will attempt to apply all keywords (however the simulation results
may be incorrect). Adding additional call points to the
opm.io.schedule.Schedule instance in order to update the Schedule object
in more ways is an obvious candidate for improving the PYACTION
functionality. \\
3 & report\_step & The current report step. \\
4 & summary\_state & The summary\_state variable is an instance of the
opm.io.sim.SummaryState class, the purpose of this class is to serve as
a container for SUMMARY variables that contain the simulation results;
for example, the Field Oil Rate (FOPR) or a Well's Water CuT (WWCT). See
the SUMMARY SECTION for a detailed description of the variables
available. The summary\_state variable will typically be the variable
one will use to access the state of the simulation. For example, to
check if the water cut in well OP01 exceeds 0.5 one would use the
following statement: if summary\_state.well\_var(``OP01'', ``WWCT'')
\textgreater{} 0.50: \ldots{} The summary\_state variable can also be
use to update variables including UDQ variables, i.e:
summary\_state.update\_well\_var(``OP01'', ``WUXX'', 0.25) The above
assigns a value of 0.25 to the well UDQ variable WUXX for well OP01. \\
\end{longtable}

\emph{Table 12.61: PYACTION Module Script Definition}

See also the PYINPUT and PYEND keywords in the GRID {[}Note the PYINPUT
and PYEND keywords can be used in the GRID, EDIT, PROPS, SOLUTION,
SUMMARY and SCHEDULE sections, but are described in the GRID section.{]}
section which are also part of OPM Flow's Python scripting facility,
that process standard Python commands that can be used to manipulate and
define the simulators input parameters during processing of the input
deck. The main purpose of the facility is to script the construction of
the various keywords.

\subsubsection{Examples}\label{examples-83}

The first example checks if well OP01 has a water cut greater than 0.8
and if so then the well is shut in. In the input deck we would have:

\begin{verbatim}
--
--       START OF PYACTION SECTION
--
--       ACTNAME       ACTNSTEP
PYACTION
         'MAXWCUT'     'UNLIMITED'                          /
         'pthon/script/MAXWCUT.py'                          /
--
--       END OF PYACTION SECTION
--
\end{verbatim}

And then in the Python module file `pthon/script/MAXWCUT.py' one would
have:

\begin{verbatim}
#
# OPM Flow PYACTION Module Script
#
import opm_embedded
schedule = opm_embedded.current_schedule
summary_state = opm_embedded.current_summary_state

if (not 'setup_done' in locals()):
    executed = False
    setup_done = True

if (not executed and summary_state.well_var("OP1", "WWCT") > 0.80):
    message = "Well OP01 has been shut-in due to WWCT > 0.80\n"
    opm_embedded.OpmLog.info(message)
    schedule.shut_well("OP1")
    executed = True
\end{verbatim}

Note how the `executed' variable is used to ensure that the action is
only executed once the first time the water cut is greater than 0.8.

The next example is based on the first example from the ACTIONX keyword
(ACTIONX -- Define Action Conditions and Command Processing). The Python
script first checks if the field's water production is greater than
30,000 stb/d, and if not returns control back to the simulator. If the
field water production is greater than 30,000 stb/d then the script uses
a Python variable count to keep track of the number of times the script
has been executed, and then sorts the wells from high water cut to low,
via the wct\_list variable, and then shuts in the worst offending well.
If a well is shut-in the count variable is increased by one and control
is passed back to the simulator. The script is executed as maximum of
ten times.

\begin{verbatim}
--
--       START OF PYACTION SECTION
--
--       ACTNAME       ACTNSTEP
PYACTION
         'WSHUTIN'     'UNLIMITED'                          /
         'pthon/script/WSHUTIN.py'                          /
--
--       END OF PYACTION SECTION
--
\end{verbatim}

And then in the Python module file `pthon/script/WSHUTIN.py' one would
have:

\begin{verbatim}
#
# OPM Flow PYACTION Module Script
#
import opm_embedded

schedule = opm_embedded.current_schedule
summary_state = opm_embedded.current_summary_state

if (not 'setup_done' in locals()):
    execution_counter = 0
    setup_done = True

if (execution_counter < 10 and summary_state["FWPR"] >= 30000):
    #
    # Get Sorted Well List
    #
    wct_list = sorted( [ (well, summary_state.well_var(well, "WWCT"))
                       for well in summary_state.wells], reverse=True )
    #
    # Shut-in Well with Highest Water Cut
    #
    well, wwct = wct_list[0]
    if wwct > 0:
        schedule.shut_well(well)
        execution_counter += 1
\end{verbatim}

Note that by using PYACTION it makes sense to combine the UDQ variable
into the PYACTION statement, like shown in the example above, although
this is not necessary, as one could in principle use a normal UDQ
statement and then access the variables in the PYACTION script using the
summary\_state variable.

The final example checks to see if the field's gas rate is below 600
MMscf/d and if the simulation time is greater that January 1, 2030. If
it is, then compression is installed by re-setting all the gas producing
well's THP and BHP pressures to 450 psia and 300 psia respectively. In
addition all gas wells currently shut-in are tested to see if they can
be opened up under the new THP and BHP constraints.

\begin{verbatim}
--
--       START OF PYACTION SECTION
--
--       ACTNAME       ACTNSTEP
PYACTION
         'WSHUTIN'     'UNLIMITED'                          /
         'pthon/script/PHASE3.py'                           /
--
--       END OF PYACTION SECTION
--
\end{verbatim}

And then in the Python module file `pthon/script/PHASE3.py' one would
have:

\begin{verbatim}
#
# OPM Flow PYACTION Module Script
#
import datetime
import opm_embedded

schedule = opm_embedded.current_schedule
summary_state = opm_embedded.current_summary_state

if (not 'setup_done' in locals()):
    executed = False
    setup_done = True

if not executed:
    sim_time = schedule.start +
               datetime.timedelta( seconds = summary_state.elapsed() )

    if (summary_state["FGPR"] < 600000 and
        sim_time > datetime.datetime(2030, 1, 1)):
        #
        # Do WELTARG and WTEST action
        #
        ...
        executed = True
\end{verbatim}

Note how the current simulation time (sim\_time) is evaluated from the
simulation start time (schedule.start) and elapsed time
(summary\_state.elapsed); and how it is compared with 00:00 on 1 January
2030 using the Python datetime module.

\subsection{PYEND -- End the Definition of a PYINPUT
Section}\label{pyend-end-the-definition-of-a-pyinput-section-6}

The PYINPUT and PYEND keywords are a part of OPM Flow's Python scripting
facility that processes standard Python commands that can be used to
manipulate and define the simulators input parameters during processing
of the input deck. The main purpose of the facility is to script the
construction of the various keywords. A PYINPUT Definition Section is
terminated by a PYEND keyword on a separate single line.

See PYEND -- End the Definition of a PYINPUT Section in the GRID section
for a full description.

\subsection{PYINPUT -- Define the Start of a PYINPUT
Section}\label{pyinput-define-the-start-of-a-pyinput-section-6}

The PYINPUT and PYEND keywords are a part of OPM Flow's Python scripting
facility that processes standard Python commands that can be used to
manipulate and define the simulators input parameters during processing
of the input deck. The main purpose of the facility is to script the
construction of the various keywords used by the simulator. PYINPUT
declares the start of a PYINPUT Definition Section on a single separate
line, which is then followed by various standard Python commands, with
one command per line.

See PYINPUT -- Define the Start of a PYINPUT Section in the GRID section
for a full description.

\subsection{QDRILL -- Define Sequential Drilling Queue
Wells}\label{qdrill-define-sequential-drilling-queue-wells}

\begin{longtable}[]{@{}
  >{\raggedright\arraybackslash}p{(\linewidth - 14\tabcolsep) * \real{0.1250}}
  >{\raggedright\arraybackslash}p{(\linewidth - 14\tabcolsep) * \real{0.1250}}
  >{\raggedright\arraybackslash}p{(\linewidth - 14\tabcolsep) * \real{0.1250}}
  >{\raggedright\arraybackslash}p{(\linewidth - 14\tabcolsep) * \real{0.1250}}
  >{\raggedright\arraybackslash}p{(\linewidth - 14\tabcolsep) * \real{0.1250}}
  >{\raggedright\arraybackslash}p{(\linewidth - 14\tabcolsep) * \real{0.1250}}
  >{\raggedright\arraybackslash}p{(\linewidth - 14\tabcolsep) * \real{0.1250}}
  >{\raggedright\arraybackslash}p{(\linewidth - 14\tabcolsep) * \real{0.1250}}@{}}
\toprule\noalign{}
\begin{minipage}[b]{\linewidth}\raggedright
RUNSPEC
\end{minipage} & \begin{minipage}[b]{\linewidth}\raggedright
GRID
\end{minipage} & \begin{minipage}[b]{\linewidth}\raggedright
EDIT
\end{minipage} & \begin{minipage}[b]{\linewidth}\raggedright
PROPS
\end{minipage} & \begin{minipage}[b]{\linewidth}\raggedright
REGIONS
\end{minipage} & \begin{minipage}[b]{\linewidth}\raggedright
SOLUTION
\end{minipage} & \begin{minipage}[b]{\linewidth}\raggedright
SUMMARY
\end{minipage} & \begin{minipage}[b]{\linewidth}\raggedright
SCHEDULE
\end{minipage} \\
\midrule\noalign{}
\endhead
\bottomrule\noalign{}
\endlastfoot
\end{longtable}

\subsubsection{Description}\label{description-996}

The QDRILL keyword places previously defined wells in the Sequential
Drilling Queue. Wells in this type of queue will be automatically
drilled and completed in the sequence entered in order to satisfy group
targets, as defined by the GCONPROD, GCONINJE and GCONSALE keywords. or
a group's production potential as per the GDRILLPOT keyword. All the
previously mentioned keywords are in the SCHEDULE section.

This keyword is not supported by OPM Flow but would change the results
if supported so the simulation will be stopped.

\subsection{RAINFALL -- Constant Flux Aquifer Rainfall Flux by
Month}\label{rainfall-constant-flux-aquifer-rainfall-flux-by-month-1}

This keyword, RAINFALL, defines the month by month rainfall flux for
constant flux aquifers.

See RAINFALL -- Constant Flux Aquifer Rainfall Flux by Month in the GRID
section for a full description.

\subsection{RCMASTS -- Reservoir Coupling Group Minimum Time Step for
Flow
Restriction}\label{rcmasts-reservoir-coupling-group-minimum-time-step-for-flow-restriction}

\begin{longtable}[]{@{}
  >{\raggedright\arraybackslash}p{(\linewidth - 14\tabcolsep) * \real{0.1250}}
  >{\raggedright\arraybackslash}p{(\linewidth - 14\tabcolsep) * \real{0.1250}}
  >{\raggedright\arraybackslash}p{(\linewidth - 14\tabcolsep) * \real{0.1250}}
  >{\raggedright\arraybackslash}p{(\linewidth - 14\tabcolsep) * \real{0.1250}}
  >{\raggedright\arraybackslash}p{(\linewidth - 14\tabcolsep) * \real{0.1250}}
  >{\raggedright\arraybackslash}p{(\linewidth - 14\tabcolsep) * \real{0.1250}}
  >{\raggedright\arraybackslash}p{(\linewidth - 14\tabcolsep) * \real{0.1250}}
  >{\raggedright\arraybackslash}p{(\linewidth - 14\tabcolsep) * \real{0.1250}}@{}}
\toprule\noalign{}
\begin{minipage}[b]{\linewidth}\raggedright
RUNSPEC
\end{minipage} & \begin{minipage}[b]{\linewidth}\raggedright
GRID
\end{minipage} & \begin{minipage}[b]{\linewidth}\raggedright
EDIT
\end{minipage} & \begin{minipage}[b]{\linewidth}\raggedright
PROPS
\end{minipage} & \begin{minipage}[b]{\linewidth}\raggedright
REGIONS
\end{minipage} & \begin{minipage}[b]{\linewidth}\raggedright
SOLUTION
\end{minipage} & \begin{minipage}[b]{\linewidth}\raggedright
SUMMARY
\end{minipage} & \begin{minipage}[b]{\linewidth}\raggedright
SCHEDULE
\end{minipage} \\
\midrule\noalign{}
\endhead
\bottomrule\noalign{}
\endlastfoot
\end{longtable}

\subsubsection{Description}\label{description-997}

RCMASTS is used when reservoir coupling is invoked by the GRUPMAST and
SLAVES keywords in the SCHEDULE section. The keyword should be placed
within the master file and it sets the minimum time step size for groups
for when a group is being restricted by a group's limiting flow rate
fractional change (see the GRUPMAST keyword in the SCHEDULE section).

This keyword is not supported by OPM Flow but would change the results
if supported so the simulation will be stopped.

\subsection{REACHES -- Define River Reaches
Structure}\label{reaches-define-river-reaches-structure}

\begin{longtable}[]{@{}
  >{\raggedright\arraybackslash}p{(\linewidth - 14\tabcolsep) * \real{0.1250}}
  >{\raggedright\arraybackslash}p{(\linewidth - 14\tabcolsep) * \real{0.1250}}
  >{\raggedright\arraybackslash}p{(\linewidth - 14\tabcolsep) * \real{0.1250}}
  >{\raggedright\arraybackslash}p{(\linewidth - 14\tabcolsep) * \real{0.1250}}
  >{\raggedright\arraybackslash}p{(\linewidth - 14\tabcolsep) * \real{0.1250}}
  >{\raggedright\arraybackslash}p{(\linewidth - 14\tabcolsep) * \real{0.1250}}
  >{\raggedright\arraybackslash}p{(\linewidth - 14\tabcolsep) * \real{0.1250}}
  >{\raggedright\arraybackslash}p{(\linewidth - 14\tabcolsep) * \real{0.1250}}@{}}
\toprule\noalign{}
\begin{minipage}[b]{\linewidth}\raggedright
RUNSPEC
\end{minipage} & \begin{minipage}[b]{\linewidth}\raggedright
GRID
\end{minipage} & \begin{minipage}[b]{\linewidth}\raggedright
EDIT
\end{minipage} & \begin{minipage}[b]{\linewidth}\raggedright
PROPS
\end{minipage} & \begin{minipage}[b]{\linewidth}\raggedright
REGIONS
\end{minipage} & \begin{minipage}[b]{\linewidth}\raggedright
SOLUTION
\end{minipage} & \begin{minipage}[b]{\linewidth}\raggedright
SUMMARY
\end{minipage} & \begin{minipage}[b]{\linewidth}\raggedright
SCHEDULE
\end{minipage} \\
\midrule\noalign{}
\endhead
\bottomrule\noalign{}
\endlastfoot
\end{longtable}

\subsubsection{Description}\label{description-998}

The REACHES keyword defines the reach structure of a previously
characterized river system using the RIVERSYS keyword in the SCHEDULE
section. The River option must be activated via the RIVRDIMS keyword in
the RUNSPEC section in order to use these keywords.

This keyword is not supported by OPM Flow but would change the results
if supported so the simulation will be stopped.

\subsection{READDATA -- Read Schedule Data Based on Current Time
Step}\label{readdata-read-schedule-data-based-on-current-time-step}

\begin{longtable}[]{@{}
  >{\raggedright\arraybackslash}p{(\linewidth - 14\tabcolsep) * \real{0.1250}}
  >{\raggedright\arraybackslash}p{(\linewidth - 14\tabcolsep) * \real{0.1250}}
  >{\raggedright\arraybackslash}p{(\linewidth - 14\tabcolsep) * \real{0.1250}}
  >{\raggedright\arraybackslash}p{(\linewidth - 14\tabcolsep) * \real{0.1250}}
  >{\raggedright\arraybackslash}p{(\linewidth - 14\tabcolsep) * \real{0.1250}}
  >{\raggedright\arraybackslash}p{(\linewidth - 14\tabcolsep) * \real{0.1250}}
  >{\raggedright\arraybackslash}p{(\linewidth - 14\tabcolsep) * \real{0.1250}}
  >{\raggedright\arraybackslash}p{(\linewidth - 14\tabcolsep) * \real{0.1250}}@{}}
\toprule\noalign{}
\begin{minipage}[b]{\linewidth}\raggedright
RUNSPEC
\end{minipage} & \begin{minipage}[b]{\linewidth}\raggedright
GRID
\end{minipage} & \begin{minipage}[b]{\linewidth}\raggedright
EDIT
\end{minipage} & \begin{minipage}[b]{\linewidth}\raggedright
PROPS
\end{minipage} & \begin{minipage}[b]{\linewidth}\raggedright
REGIONS
\end{minipage} & \begin{minipage}[b]{\linewidth}\raggedright
SOLUTION
\end{minipage} & \begin{minipage}[b]{\linewidth}\raggedright
SUMMARY
\end{minipage} & \begin{minipage}[b]{\linewidth}\raggedright
SCHEDULE
\end{minipage} \\
\midrule\noalign{}
\endhead
\bottomrule\noalign{}
\endlastfoot
\end{longtable}

\subsubsection{Description}\label{description-999}

The READDATA keyword enables the simulator to read SCHEDULE data files
generated by external programs on the fly, that is as the run is
progressing. The external program can ``hold'' the simulation by using a
file lock, in order for the external program to evaluate the current
simulation results, then write out a SCHEDULE data file for the next
time step, and finally releasing the ``hold'' by deleting the file lock
and continuing with newly written SCHEDULE data file. The mechanism can
be repeated so that the external program is dictating how the simulation
progresses through time,

This keyword is not supported by OPM Flow but would change the results
if supported so the simulation will be stopped.

\subsection{REFINE -- Start the Definition of a Local Grid
Refinement}\label{refine-start-the-definition-of-a-local-grid-refinement-5}

The REFINE keyword defines the start of a Cartesian or radial local grid
refinement (``LGR'') definition that sets the properties of the selected
LGR. The keyword is then followed by the property keywords associated
with the section where the keyword is being invoked. For example, if the
REFINE keyword is used in the GRID section then most of the keywords in
that section can be used to set the grid properties for the LGR.

See REFINE -- Start the Definition of a Local Grid Refinement in the
GRID section for a full description.

\subsection{RIVDEBUG -- Define the Debug Data to be Printed to File
(Rivers)}\label{rivdebug-define-the-debug-data-to-be-printed-to-file-rivers}

\begin{longtable}[]{@{}
  >{\raggedright\arraybackslash}p{(\linewidth - 14\tabcolsep) * \real{0.1250}}
  >{\raggedright\arraybackslash}p{(\linewidth - 14\tabcolsep) * \real{0.1250}}
  >{\raggedright\arraybackslash}p{(\linewidth - 14\tabcolsep) * \real{0.1250}}
  >{\raggedright\arraybackslash}p{(\linewidth - 14\tabcolsep) * \real{0.1250}}
  >{\raggedright\arraybackslash}p{(\linewidth - 14\tabcolsep) * \real{0.1250}}
  >{\raggedright\arraybackslash}p{(\linewidth - 14\tabcolsep) * \real{0.1250}}
  >{\raggedright\arraybackslash}p{(\linewidth - 14\tabcolsep) * \real{0.1250}}
  >{\raggedright\arraybackslash}p{(\linewidth - 14\tabcolsep) * \real{0.1250}}@{}}
\toprule\noalign{}
\begin{minipage}[b]{\linewidth}\raggedright
RUNSPEC
\end{minipage} & \begin{minipage}[b]{\linewidth}\raggedright
GRID
\end{minipage} & \begin{minipage}[b]{\linewidth}\raggedright
EDIT
\end{minipage} & \begin{minipage}[b]{\linewidth}\raggedright
PROPS
\end{minipage} & \begin{minipage}[b]{\linewidth}\raggedright
REGIONS
\end{minipage} & \begin{minipage}[b]{\linewidth}\raggedright
SOLUTION
\end{minipage} & \begin{minipage}[b]{\linewidth}\raggedright
SUMMARY
\end{minipage} & \begin{minipage}[b]{\linewidth}\raggedright
SCHEDULE
\end{minipage} \\
\midrule\noalign{}
\endhead
\bottomrule\noalign{}
\endlastfoot
\end{longtable}

\subsubsection{Description}\label{description-1000}

This keyword defines the debug data associated with rivers to be written
to the debug file (*.DBG), for when the River option has been activated
via the RIVRDIMS keyword in the RUNSPEC section.

This keyword is not supported by OPM Flow but would change the results
if supported so the simulation will be stopped.

\subsection{RIVERSYS -- Define River System (Branch Structure and
Boundary
Conditions)}\label{riversys-define-river-system-branch-structure-and-boundary-conditions-1}

RIVERSYS defines a river system by specifying the branch structure of
the river together with the branch's associated boundary conditions, for
when the River option has been activated via the RIVRDIMS keyword in the
RUNSPEC section.

See RIVERSYS - Define River System (Branch Structure and Boundary
Conditions) in the SOLUTION section for a full description.

\subsection{RIVRPROP -- Modify River Reaches
Properties}\label{rivrprop-modify-river-reaches-properties}

\begin{longtable}[]{@{}
  >{\raggedright\arraybackslash}p{(\linewidth - 14\tabcolsep) * \real{0.1250}}
  >{\raggedright\arraybackslash}p{(\linewidth - 14\tabcolsep) * \real{0.1250}}
  >{\raggedright\arraybackslash}p{(\linewidth - 14\tabcolsep) * \real{0.1250}}
  >{\raggedright\arraybackslash}p{(\linewidth - 14\tabcolsep) * \real{0.1250}}
  >{\raggedright\arraybackslash}p{(\linewidth - 14\tabcolsep) * \real{0.1250}}
  >{\raggedright\arraybackslash}p{(\linewidth - 14\tabcolsep) * \real{0.1250}}
  >{\raggedright\arraybackslash}p{(\linewidth - 14\tabcolsep) * \real{0.1250}}
  >{\raggedright\arraybackslash}p{(\linewidth - 14\tabcolsep) * \real{0.1250}}@{}}
\toprule\noalign{}
\begin{minipage}[b]{\linewidth}\raggedright
RUNSPEC
\end{minipage} & \begin{minipage}[b]{\linewidth}\raggedright
GRID
\end{minipage} & \begin{minipage}[b]{\linewidth}\raggedright
EDIT
\end{minipage} & \begin{minipage}[b]{\linewidth}\raggedright
PROPS
\end{minipage} & \begin{minipage}[b]{\linewidth}\raggedright
REGIONS
\end{minipage} & \begin{minipage}[b]{\linewidth}\raggedright
SOLUTION
\end{minipage} & \begin{minipage}[b]{\linewidth}\raggedright
SUMMARY
\end{minipage} & \begin{minipage}[b]{\linewidth}\raggedright
SCHEDULE
\end{minipage} \\
\midrule\noalign{}
\endhead
\bottomrule\noalign{}
\endlastfoot
\end{longtable}

\subsubsection{Description}\label{description-1001}

The RIVRPROP keyword modifies the individual reaches in a river
structure of a previously characterized river system using the RIVERSYS
and the REACHES keywords in the SCHEDULE section. The River option must
be activated via the RIVRDIMS keyword in the RUNSPEC section in order to
use this keyword. RIVRPROP is an alternative and a more concise way to
changing the individual reaches in a river structure than the REACHES
keyword.

This keyword is not supported by OPM Flow but would change the results
if supported so the simulation will be stopped.

\subsection{RIVSALT -- Define River Upstream Flow Salt
Concentrations}\label{rivsalt-define-river-upstream-flow-salt-concentrations}

\begin{longtable}[]{@{}
  >{\raggedright\arraybackslash}p{(\linewidth - 14\tabcolsep) * \real{0.1250}}
  >{\raggedright\arraybackslash}p{(\linewidth - 14\tabcolsep) * \real{0.1250}}
  >{\raggedright\arraybackslash}p{(\linewidth - 14\tabcolsep) * \real{0.1250}}
  >{\raggedright\arraybackslash}p{(\linewidth - 14\tabcolsep) * \real{0.1250}}
  >{\raggedright\arraybackslash}p{(\linewidth - 14\tabcolsep) * \real{0.1250}}
  >{\raggedright\arraybackslash}p{(\linewidth - 14\tabcolsep) * \real{0.1250}}
  >{\raggedright\arraybackslash}p{(\linewidth - 14\tabcolsep) * \real{0.1250}}
  >{\raggedright\arraybackslash}p{(\linewidth - 14\tabcolsep) * \real{0.1250}}@{}}
\toprule\noalign{}
\begin{minipage}[b]{\linewidth}\raggedright
RUNSPEC
\end{minipage} & \begin{minipage}[b]{\linewidth}\raggedright
GRID
\end{minipage} & \begin{minipage}[b]{\linewidth}\raggedright
EDIT
\end{minipage} & \begin{minipage}[b]{\linewidth}\raggedright
PROPS
\end{minipage} & \begin{minipage}[b]{\linewidth}\raggedright
REGIONS
\end{minipage} & \begin{minipage}[b]{\linewidth}\raggedright
SOLUTION
\end{minipage} & \begin{minipage}[b]{\linewidth}\raggedright
SUMMARY
\end{minipage} & \begin{minipage}[b]{\linewidth}\raggedright
SCHEDULE
\end{minipage} \\
\midrule\noalign{}
\endhead
\bottomrule\noalign{}
\endlastfoot
\end{longtable}

\subsubsection{Description}\label{description-1002}

The RIVSALT keyword defines the injected salt concentration in
individual river branches in a previously characterized river system
using the RIVERSYS and the REACHES keywords in the SCHEDULE section. The
River option must be activated via the RIVRDIMS keyword in the RUNSPEC
section in order to use this keyword. In addition, the Brine option must
also be enabled via the BRINE keyword in the RUNSPEC section.

This keyword is not supported by OPM Flow but would change the results
if supported so the simulation will be stopped.

\subsection{RIVTRACE -- Define River Upstream Flow Tracer
Concentrations}\label{rivtrace-define-river-upstream-flow-tracer-concentrations}

\begin{longtable}[]{@{}
  >{\raggedright\arraybackslash}p{(\linewidth - 14\tabcolsep) * \real{0.1250}}
  >{\raggedright\arraybackslash}p{(\linewidth - 14\tabcolsep) * \real{0.1250}}
  >{\raggedright\arraybackslash}p{(\linewidth - 14\tabcolsep) * \real{0.1250}}
  >{\raggedright\arraybackslash}p{(\linewidth - 14\tabcolsep) * \real{0.1250}}
  >{\raggedright\arraybackslash}p{(\linewidth - 14\tabcolsep) * \real{0.1250}}
  >{\raggedright\arraybackslash}p{(\linewidth - 14\tabcolsep) * \real{0.1250}}
  >{\raggedright\arraybackslash}p{(\linewidth - 14\tabcolsep) * \real{0.1250}}
  >{\raggedright\arraybackslash}p{(\linewidth - 14\tabcolsep) * \real{0.1250}}@{}}
\toprule\noalign{}
\begin{minipage}[b]{\linewidth}\raggedright
RUNSPEC
\end{minipage} & \begin{minipage}[b]{\linewidth}\raggedright
GRID
\end{minipage} & \begin{minipage}[b]{\linewidth}\raggedright
EDIT
\end{minipage} & \begin{minipage}[b]{\linewidth}\raggedright
PROPS
\end{minipage} & \begin{minipage}[b]{\linewidth}\raggedright
REGIONS
\end{minipage} & \begin{minipage}[b]{\linewidth}\raggedright
SOLUTION
\end{minipage} & \begin{minipage}[b]{\linewidth}\raggedright
SUMMARY
\end{minipage} & \begin{minipage}[b]{\linewidth}\raggedright
SCHEDULE
\end{minipage} \\
\midrule\noalign{}
\endhead
\bottomrule\noalign{}
\endlastfoot
\end{longtable}

\subsubsection{Description}\label{description-1003}

The RIVTRACE keyword defines the injected tracer concentration in
individual river branches in a previously characterized river system
using the RIVERSYS and the REACHES keywords in the SCHEDULE section. The
River option must be activated via the RIVRDIMS keyword in the RUNSPEC
section in order to use this keyword. In addition, the Tracer option
must also be enabled by the TRACER keyword in the RUNSPEC section.

This keyword is not supported by OPM Flow but would change the results
if supported so the simulation will be stopped.

\subsection{RPTHMG -- Define Well History Match Gradient Reporting
(Groups)}\label{rpthmg-define-well-history-match-gradient-reporting-groups}

\begin{longtable}[]{@{}
  >{\raggedright\arraybackslash}p{(\linewidth - 14\tabcolsep) * \real{0.1250}}
  >{\raggedright\arraybackslash}p{(\linewidth - 14\tabcolsep) * \real{0.1250}}
  >{\raggedright\arraybackslash}p{(\linewidth - 14\tabcolsep) * \real{0.1250}}
  >{\raggedright\arraybackslash}p{(\linewidth - 14\tabcolsep) * \real{0.1250}}
  >{\raggedright\arraybackslash}p{(\linewidth - 14\tabcolsep) * \real{0.1250}}
  >{\raggedright\arraybackslash}p{(\linewidth - 14\tabcolsep) * \real{0.1250}}
  >{\raggedright\arraybackslash}p{(\linewidth - 14\tabcolsep) * \real{0.1250}}
  >{\raggedright\arraybackslash}p{(\linewidth - 14\tabcolsep) * \real{0.1250}}@{}}
\toprule\noalign{}
\begin{minipage}[b]{\linewidth}\raggedright
RUNSPEC
\end{minipage} & \begin{minipage}[b]{\linewidth}\raggedright
GRID
\end{minipage} & \begin{minipage}[b]{\linewidth}\raggedright
EDIT
\end{minipage} & \begin{minipage}[b]{\linewidth}\raggedright
PROPS
\end{minipage} & \begin{minipage}[b]{\linewidth}\raggedright
REGIONS
\end{minipage} & \begin{minipage}[b]{\linewidth}\raggedright
SOLUTION
\end{minipage} & \begin{minipage}[b]{\linewidth}\raggedright
SUMMARY
\end{minipage} & \begin{minipage}[b]{\linewidth}\raggedright
SCHEDULE
\end{minipage} \\
\midrule\noalign{}
\endhead
\bottomrule\noalign{}
\endlastfoot
\end{longtable}

\subsubsection{Description}\label{description-1004}

This keyword, RPTHMG, either enables or disables history match output
reporting to the history match file (*.HMD) for the named group, for
when the History Match Gradient option has been activated by the HMDIMS
keyword in the RUNSPEC section.

This keyword is not supported by OPM Flow but has no effect on the
results so it will be ignored.

\subsection{RPTHMW -- Define Well History Match Gradient Reporting
(Wells)}\label{rpthmw-define-well-history-match-gradient-reporting-wells}

\begin{longtable}[]{@{}
  >{\raggedright\arraybackslash}p{(\linewidth - 14\tabcolsep) * \real{0.1250}}
  >{\raggedright\arraybackslash}p{(\linewidth - 14\tabcolsep) * \real{0.1250}}
  >{\raggedright\arraybackslash}p{(\linewidth - 14\tabcolsep) * \real{0.1250}}
  >{\raggedright\arraybackslash}p{(\linewidth - 14\tabcolsep) * \real{0.1250}}
  >{\raggedright\arraybackslash}p{(\linewidth - 14\tabcolsep) * \real{0.1250}}
  >{\raggedright\arraybackslash}p{(\linewidth - 14\tabcolsep) * \real{0.1250}}
  >{\raggedright\arraybackslash}p{(\linewidth - 14\tabcolsep) * \real{0.1250}}
  >{\raggedright\arraybackslash}p{(\linewidth - 14\tabcolsep) * \real{0.1250}}@{}}
\toprule\noalign{}
\begin{minipage}[b]{\linewidth}\raggedright
RUNSPEC
\end{minipage} & \begin{minipage}[b]{\linewidth}\raggedright
GRID
\end{minipage} & \begin{minipage}[b]{\linewidth}\raggedright
EDIT
\end{minipage} & \begin{minipage}[b]{\linewidth}\raggedright
PROPS
\end{minipage} & \begin{minipage}[b]{\linewidth}\raggedright
REGIONS
\end{minipage} & \begin{minipage}[b]{\linewidth}\raggedright
SOLUTION
\end{minipage} & \begin{minipage}[b]{\linewidth}\raggedright
SUMMARY
\end{minipage} & \begin{minipage}[b]{\linewidth}\raggedright
SCHEDULE
\end{minipage} \\
\midrule\noalign{}
\endhead
\bottomrule\noalign{}
\endlastfoot
\end{longtable}

\subsubsection{Description}\label{description-1005}

This keyword, RPTHMG, either enables or disables history match output
reporting to the history match file (*.HMD) for the named well, for when
the History Match Gradient option has been activated by the HMDIMS
keyword in the RUNSPEC section.

This keyword is not supported by OPM Flow but has no effect on the
results so it will be ignored.

\subsection{RPTONLY -- Activate the Report Time Steps Only Option for
the SUMMARY
File}\label{rptonly-activate-the-report-time-steps-only-option-for-the-summary-file-1}

This keyword activates the writing out of the SUMMARY file and RSM file
data, if the RSM file option has been requested by the RUNSUM keyword in
the SUMMARY section, at report time steps only. The default is for all
time steps to be written out to the files. This keyword reduces the file
size at the expense of lower resolution in the time domain. There is no
data required for this keyword.

See RPTONLY - Activate the Report Time Steps Only Option for the SUMMARY
File in the SUMMARY section for a full description.

\subsection{RPTONLYO -- Deactivate the Report Time Steps Only Option for
the SUMMARY
File}\label{rptonlyo-deactivate-the-report-time-steps-only-option-for-the-summary-file-1}

This keyword deactivates the writing out of the SUMMARY file and RSM
file data, if the RSM file option has been requested by the RUNSUM
keyword in the SUMMARY section, at report time steps only, and switches
on writing out all the time steps to the files. This option is the
default behavior for when RPTONLY has not been activated. There is no
data required for this keyword.

See RPTONLYO - Deactivate the Report Time Steps Only Option for the
SUMMARY File in the SUMMARY section for a full description.

\subsection{RPTRST -- Define Data to be Written to the RESTART
File}\label{rptrst-define-data-to-be-written-to-the-restart-file-1}

This keyword defines the data to be written to the RESTART file at each
requested restart point.

See RPTRST -- Define Data to be Written to the RESTART File in the
SOLUTION section for a full description.

\subsection{RPTSCHED -- Define SCHEDULE Section
Reporting}\label{rptsched-define-schedule-section-reporting}

\begin{longtable}[]{@{}
  >{\raggedright\arraybackslash}p{(\linewidth - 14\tabcolsep) * \real{0.1250}}
  >{\raggedright\arraybackslash}p{(\linewidth - 14\tabcolsep) * \real{0.1250}}
  >{\raggedright\arraybackslash}p{(\linewidth - 14\tabcolsep) * \real{0.1250}}
  >{\raggedright\arraybackslash}p{(\linewidth - 14\tabcolsep) * \real{0.1250}}
  >{\raggedright\arraybackslash}p{(\linewidth - 14\tabcolsep) * \real{0.1250}}
  >{\raggedright\arraybackslash}p{(\linewidth - 14\tabcolsep) * \real{0.1250}}
  >{\raggedright\arraybackslash}p{(\linewidth - 14\tabcolsep) * \real{0.1250}}
  >{\raggedright\arraybackslash}p{(\linewidth - 14\tabcolsep) * \real{0.1250}}@{}}
\toprule\noalign{}
\begin{minipage}[b]{\linewidth}\raggedright
RUNSPEC
\end{minipage} & \begin{minipage}[b]{\linewidth}\raggedright
GRID
\end{minipage} & \begin{minipage}[b]{\linewidth}\raggedright
EDIT
\end{minipage} & \begin{minipage}[b]{\linewidth}\raggedright
PROPS
\end{minipage} & \begin{minipage}[b]{\linewidth}\raggedright
REGIONS
\end{minipage} & \begin{minipage}[b]{\linewidth}\raggedright
SOLUTION
\end{minipage} & \begin{minipage}[b]{\linewidth}\raggedright
SUMMARY
\end{minipage} & \begin{minipage}[b]{\linewidth}\raggedright
SCHEDULE
\end{minipage} \\
\midrule\noalign{}
\endhead
\bottomrule\noalign{}
\endlastfoot
\end{longtable}

\subsubsection{Description}\label{description-1006}

This keyword defines the data in the SCHEDULE section that is to be
printed to the output print file in human readable format. The keyword
has two distinct forms, the first of which consists of the keyword
followed by a series of integers on the next line indicating the data to
be printed (see the first example). This is the original format in the
commercial simulator and was subsequently superseded by the second
format. The second format consists of the keyword followed by a series
of character strings that indicate the data to be printed. In most cases
the character string is the keyword used to define the data in the OPM
Flow input deck, for example WELSPECS to defined the basic well
definitions.

\begin{longtable}[]{@{}
  >{\raggedright\arraybackslash}p{(\linewidth - 6\tabcolsep) * \real{0.2500}}
  >{\raggedright\arraybackslash}p{(\linewidth - 6\tabcolsep) * \real{0.2500}}
  >{\raggedright\arraybackslash}p{(\linewidth - 6\tabcolsep) * \real{0.2500}}
  >{\raggedright\arraybackslash}p{(\linewidth - 6\tabcolsep) * \real{0.2500}}@{}}
\toprule\noalign{}
\begin{minipage}[b]{\linewidth}\raggedright
No.
\end{minipage} & \begin{minipage}[b]{\linewidth}\raggedright
Name
\end{minipage} & \begin{minipage}[b]{\linewidth}\raggedright
Description
\end{minipage} & \begin{minipage}[b]{\linewidth}\raggedright
Default
\end{minipage} \\
\midrule\noalign{}
\endhead
\bottomrule\noalign{}
\endlastfoot
1 & FIP & Print the fluid in-place report. The parameter is assigned a
value, OPTION, using the form FIP = OPTION, where OPTION is an integer
variable set to: & FIP=2 \\
2 & FIPRESV & Print the reservoir volumes in-place report. & None \\
3 & NOTHING & Switches off all printed reports in the SCHEDULE section.
& N/A \\
4 & SALT & Print grid block salt concentration values. Note this is an
OPM Flow specific keyword. & N/A \\
5 & RESTART & RESTART defines the frequency at which the restart data
for restarting a run is written to the RESTART file. The parameter is
assigned a value, OPTION, using the form RESTART = OPTION, where OPTION
is an integer variable set to: See the RPTRST keyword in the SOLUTION
section for a more flexible way to write out restart files. & \\
6 & WELLS & The WELLS option turns on production and injection rate and
cumulative volume reporting for produced and injected fluids. The
parameter has several levels of reporting details set by the assigned
OPTION value, using the form WELLS = OPTION, where OPTION is an integer
variable set to: Only OPTION equal to one is supported by OPM Flow. &
WELLS=1 \\
7 & WELSPECS & WELSPECS switches on reporting of the well connections,
wells and groups at each report time step. There are numerous reports
associated with this option. Unlike the other reporting parameters that
produce a report for each reporting time step, the WELSPECS report
option only produces a report if an associated keyword has been
activated at the current reporting time step. For example, if the
reporting time steps are January, February, and March 2020, and the
RPTSCHED WELSPECS option is activated in January, with wells OP01 and
OP02 being declared via the WELSPECS and COMPDAT keywords, then a report
will be printed for January for these two wells. If there are no further
well activations until March, with well OP03 being declared, then there
will be no report for February, and only well OP03 will reported at the
March reporting time step. & \\
Notes: & & & \\
\end{longtable}

\emph{Table 12.62: RPTSCHED Keyword Description}

Development is current progressing on developing reports in a similar
format to the commercial simulator and this section will be updated as
additional reports are added to OPM Flow's functionality.

\begin{tcolorbox}[enhanced jigsaw, title=\textcolor{quarto-callout-note-color}{\faInfo}\hspace{0.5em}{Note}, coltitle=black, breakable, toprule=.15mm, colbacktitle=quarto-callout-note-color!10!white, leftrule=.75mm, opacitybacktitle=0.6, colback=white, bottomtitle=1mm, toptitle=1mm, titlerule=0mm, left=2mm, colframe=quarto-callout-note-color-frame, arc=.35mm, rightrule=.15mm, bottomrule=.15mm, opacityback=0]

Unlike the other reporting keywords in the RUNSPEC, GRID, EDIT, PROPS
and SOLUTION keywords, the requested reports on the this keyword remain
in effect until they are switched off by this keyword, that is, the
reports are written out every report time step until requested to stop.
Use the `NOTHING' parameter to switch off all reporting.

\end{tcolorbox}

An example FIP report is shown in Figure 12.6 from the Norne field, note
only the field and the first two region reports are shown.

\pandocbounded{\includegraphics[keepaspectratio]{chapters/images/Frame52_fc8503037684.png}}
\pandocbounded{\includegraphics[keepaspectratio]{chapters/images/Image48_fc8503037684.png}}

Figure 12.7 illustrates the reservoir volumes in-place report for the
first two regions from the Norne field.

\pandocbounded{\includegraphics[keepaspectratio]{chapters/images/Frame72_6cf6993cfc73.png}}
\pandocbounded{\includegraphics[keepaspectratio]{chapters/images/Image65_6cf6993cfc73.png}}

\begin{tcolorbox}[enhanced jigsaw, title=\textcolor{quarto-callout-note-color}{\faInfo}\hspace{0.5em}{Note}, coltitle=black, breakable, toprule=.15mm, colbacktitle=quarto-callout-note-color!10!white, leftrule=.75mm, opacitybacktitle=0.6, colback=white, bottomtitle=1mm, toptitle=1mm, titlerule=0mm, left=2mm, colframe=quarto-callout-note-color-frame, arc=.35mm, rightrule=.15mm, bottomrule=.15mm, opacityback=0]

Note that the ``PORV'' quantity in the FIP (Balance) Report, as shown in
Figure 12.6, is reported at reference conditions, meaning there is no
pressure dependence involved. However, the ``TOTAL PORE VOLUME'' values
in the Reservoir Volumes Report (Figure 12.7) are pressure dependent
pore volumes. Thus, for region one the ``PORV'' value is 44,729,956 rm3
(Figure 12.6) and the ``TOTAL PORE VOLUME'' (Figure 12.7) is 44,719,142
rm3. This is the same as the commercial simulator.

\end{tcolorbox}

The WELLS report consists of several sub-reports depending on the
selected option for this report type. Figure 12.8 and Figure 12.9 show
example Injection and Production sub-reports

\pandocbounded{\includegraphics[keepaspectratio]{chapters/images/Frame45_d1266921f1e1.png}}
\pandocbounded{\includegraphics[keepaspectratio]{chapters/images/Image42_d1266921f1e1.png}}

\pandocbounded{\includegraphics[keepaspectratio]{chapters/images/Frame44_a374ccaa0dbb.png}}
\pandocbounded{\includegraphics[keepaspectratio]{chapters/images/Image41_a374ccaa0dbb.png}}

The third and final report is the Cumulative Production and Injection
sub-report, shown in Figure 12.10.

\pandocbounded{\includegraphics[keepaspectratio]{chapters/images/Frame53_07b5c206910b.png}}
\pandocbounded{\includegraphics[keepaspectratio]{chapters/images/Image49_07b5c206910b.png}}

Similarly as for the WELLS report, the WELSPECS report consists of
several sub-reports, including the Well Production Control report shown
in Figure 12.11 for the Volve field.

\pandocbounded{\includegraphics[keepaspectratio]{chapters/images/Frame40_8ae9f6d6efda.png}}
\pandocbounded{\includegraphics[keepaspectratio]{chapters/images/Image37_8ae9f6d6efda.png}}

The COMPDAT keyword data is listed on the Well Connection Data
sub-report as depicted in Figure 12.12 for the I-F-5 well. Note that the
data is repeated for all connections and for all wells declared at the
reporting time step.

\pandocbounded{\includegraphics[keepaspectratio]{chapters/images/Frame46_ab9afd0d1f16.png}}
\pandocbounded{\includegraphics[keepaspectratio]{chapters/images/Image43_ab9afd0d1f16.png}}

For multi-segment wells both the Production Well Control and Well
Connection sub-reports are printed as per Figure 12.11 and Figure 12.12,
and in addition the equivalent multi-segment well data is printed as
well, as shown in Figure 12.13 and Figure 12.14, as shown on the
following page.

Figure 12.13 shows the Multi-Segment Well Segment Structure sub-report
for a single multi-segment well, OP01. See the WELSEGS keyword Example
in the SCHEDULE section to see how the OP01 well is defined.

\pandocbounded{\includegraphics[keepaspectratio]{chapters/images/Frame42_d3dd8ff34480.png}}
\pandocbounded{\includegraphics[keepaspectratio]{chapters/images/Image39_d3dd8ff34480.png}}

And Figure 12.14 depicts Multi-Segment Well Connection Data sub-report
for the same well.

\pandocbounded{\includegraphics[keepaspectratio]{chapters/images/Frame43_40aa1237d71d.png}}
\pandocbounded{\includegraphics[keepaspectratio]{chapters/images/Image40_40aa1237d71d.png}}

\subsubsection{Example}\label{example-547}

The first example shows the original format of this keyword.

\begin{verbatim}
--
--       DEFINE SCHEDULE SECTION REPORT OPTION (ORIGINAL FORMAT)
--
RPTSCHED
         1        2*0      1        3*1                                        /
\end{verbatim}

The next example shows the second format of the keyword.

\begin{verbatim}
-- ==============================================================================
--
-- SCHEDULE SECTION
--
-- ==============================================================================
SCHEDULE

-- ------------------------------------------------------------------------------
-- SCHEDULE SECTION - 2000-01-01
-- ------------------------------------------------------------------------------
RPTSCHED
         'WELLS=2'    'WELSPECS'    'CPU=2'     'FIP=2'                        /

DATES
         1  JAN   2000  /
/

RPTSCHED
         'NOTHING'                                                             /

DATES
         1  FEB   2000  /
         1  MAR   2000  /
         1  APR   2000  /
         1  MAY   2000  /
         1  JUN   2000  /
         1  JLY   2000  /
         1  AUG   2000  /
         1  SEP   2000  /
         1  OCT   2000  /
         1  NOV   2000  /
         1  DEC   2000  /
/
\end{verbatim}

In the above example monthly reporting time steps have been used with a
SCHEDULE section report on the January 1, 2000; after which all reports
are switched off for the subsequent reporting time steps.

\subsection{SAVE -- Activate Output of a SAVE File for Fast
Restarts}\label{save-activate-output-of-a-save-file-for-fast-restarts-1}

This keyword activates output of a SAVE file for fast restarts. There is
no data required for this keyword.

See SAVE -- Activate Output of a SAVE File for Fast Restarts in the
RUNSPEC section for a full description.

This keyword is not supported by OPM Flow but has no effect on the
results so it will be ignored.

\subsection{SCDATAB -- Well Connection PI Multipliers versus Scale
Deposit}\label{scdatab-well-connection-pi-multipliers-versus-scale-deposit}

\begin{longtable}[]{@{}
  >{\raggedright\arraybackslash}p{(\linewidth - 14\tabcolsep) * \real{0.1250}}
  >{\raggedright\arraybackslash}p{(\linewidth - 14\tabcolsep) * \real{0.1250}}
  >{\raggedright\arraybackslash}p{(\linewidth - 14\tabcolsep) * \real{0.1250}}
  >{\raggedright\arraybackslash}p{(\linewidth - 14\tabcolsep) * \real{0.1250}}
  >{\raggedright\arraybackslash}p{(\linewidth - 14\tabcolsep) * \real{0.1250}}
  >{\raggedright\arraybackslash}p{(\linewidth - 14\tabcolsep) * \real{0.1250}}
  >{\raggedright\arraybackslash}p{(\linewidth - 14\tabcolsep) * \real{0.1250}}
  >{\raggedright\arraybackslash}p{(\linewidth - 14\tabcolsep) * \real{0.1250}}@{}}
\toprule\noalign{}
\begin{minipage}[b]{\linewidth}\raggedright
RUNSPEC
\end{minipage} & \begin{minipage}[b]{\linewidth}\raggedright
GRID
\end{minipage} & \begin{minipage}[b]{\linewidth}\raggedright
EDIT
\end{minipage} & \begin{minipage}[b]{\linewidth}\raggedright
PROPS
\end{minipage} & \begin{minipage}[b]{\linewidth}\raggedright
REGIONS
\end{minipage} & \begin{minipage}[b]{\linewidth}\raggedright
SOLUTION
\end{minipage} & \begin{minipage}[b]{\linewidth}\raggedright
SUMMARY
\end{minipage} & \begin{minipage}[b]{\linewidth}\raggedright
SCHEDULE
\end{minipage} \\
\midrule\noalign{}
\endhead
\bottomrule\noalign{}
\endlastfoot
\end{longtable}

\subsubsection{Description}\label{description-1007}

SCDATAB defines well connection Productivity Index (``PI'') reduction
multipliers versus scale deposited per unit length of the perforated
interval tables, for when the Scale Deposition option has been activated
by declaring the dimensions of the scaling deposition tables using the
SCDPDIMS keyword in the RUNSPEC section. The SCDATAB tables are
allocated to individual wells using the WSCTAB keyword and the rate of
scale accumulation around the well connections is given by the SCDPTAB
keyword; both keywords are in the SCHEDULE section.

This keyword is not supported by OPM Flow but would change the results
if supported so the simulation will be stopped.

\subsection{SCDETAB -- Well Connection Karst Aquifer Properties for
Scale
Deposit}\label{scdetab-well-connection-karst-aquifer-properties-for-scale-deposit}

\begin{longtable}[]{@{}
  >{\raggedright\arraybackslash}p{(\linewidth - 14\tabcolsep) * \real{0.1250}}
  >{\raggedright\arraybackslash}p{(\linewidth - 14\tabcolsep) * \real{0.1250}}
  >{\raggedright\arraybackslash}p{(\linewidth - 14\tabcolsep) * \real{0.1250}}
  >{\raggedright\arraybackslash}p{(\linewidth - 14\tabcolsep) * \real{0.1250}}
  >{\raggedright\arraybackslash}p{(\linewidth - 14\tabcolsep) * \real{0.1250}}
  >{\raggedright\arraybackslash}p{(\linewidth - 14\tabcolsep) * \real{0.1250}}
  >{\raggedright\arraybackslash}p{(\linewidth - 14\tabcolsep) * \real{0.1250}}
  >{\raggedright\arraybackslash}p{(\linewidth - 14\tabcolsep) * \real{0.1250}}@{}}
\toprule\noalign{}
\begin{minipage}[b]{\linewidth}\raggedright
RUNSPEC
\end{minipage} & \begin{minipage}[b]{\linewidth}\raggedright
GRID
\end{minipage} & \begin{minipage}[b]{\linewidth}\raggedright
EDIT
\end{minipage} & \begin{minipage}[b]{\linewidth}\raggedright
PROPS
\end{minipage} & \begin{minipage}[b]{\linewidth}\raggedright
REGIONS
\end{minipage} & \begin{minipage}[b]{\linewidth}\raggedright
SOLUTION
\end{minipage} & \begin{minipage}[b]{\linewidth}\raggedright
SUMMARY
\end{minipage} & \begin{minipage}[b]{\linewidth}\raggedright
SCHEDULE
\end{minipage} \\
\midrule\noalign{}
\endhead
\bottomrule\noalign{}
\endlastfoot
\end{longtable}

\subsubsection{Description}\label{description-1008}

SCDETAB defines well connection karst {[}Karst is a topography formed
from the dissolution of soluble rocks such as limestone, dolomite, and
gypsum. Karst aquifers are characterized by a network of conduits and
caves, with the conduits and caves draining the pore space between the
limestone grains (intergranular or primary porosity) and the fractures
(secondary porosity) formed by joints, bedding planes, and faults.{]}
aquifer properties for modeling scale deposited by dissolution of
calcite from the aquifer water, for when the Scale Deposition option has
been activated by declaring the dimensions of the scaling deposition
tables using the SCDPDIMS keyword in the RUNSPEC section. The SCDETAB
tables are allocated to individual wells using the WSCTAB keyword in the
SCHEDULE section.

This keyword is not supported by OPM Flow but would change the results
if supported so the simulation will be stopped.

\subsection{SCDPTAB -- Well Connection Scale Deposition
Tables}\label{scdptab-well-connection-scale-deposition-tables}

\begin{longtable}[]{@{}
  >{\raggedright\arraybackslash}p{(\linewidth - 14\tabcolsep) * \real{0.1250}}
  >{\raggedright\arraybackslash}p{(\linewidth - 14\tabcolsep) * \real{0.1250}}
  >{\raggedright\arraybackslash}p{(\linewidth - 14\tabcolsep) * \real{0.1250}}
  >{\raggedright\arraybackslash}p{(\linewidth - 14\tabcolsep) * \real{0.1250}}
  >{\raggedright\arraybackslash}p{(\linewidth - 14\tabcolsep) * \real{0.1250}}
  >{\raggedright\arraybackslash}p{(\linewidth - 14\tabcolsep) * \real{0.1250}}
  >{\raggedright\arraybackslash}p{(\linewidth - 14\tabcolsep) * \real{0.1250}}
  >{\raggedright\arraybackslash}p{(\linewidth - 14\tabcolsep) * \real{0.1250}}@{}}
\toprule\noalign{}
\begin{minipage}[b]{\linewidth}\raggedright
RUNSPEC
\end{minipage} & \begin{minipage}[b]{\linewidth}\raggedright
GRID
\end{minipage} & \begin{minipage}[b]{\linewidth}\raggedright
EDIT
\end{minipage} & \begin{minipage}[b]{\linewidth}\raggedright
PROPS
\end{minipage} & \begin{minipage}[b]{\linewidth}\raggedright
REGIONS
\end{minipage} & \begin{minipage}[b]{\linewidth}\raggedright
SOLUTION
\end{minipage} & \begin{minipage}[b]{\linewidth}\raggedright
SUMMARY
\end{minipage} & \begin{minipage}[b]{\linewidth}\raggedright
SCHEDULE
\end{minipage} \\
\midrule\noalign{}
\endhead
\bottomrule\noalign{}
\endlastfoot
\end{longtable}

\subsubsection{Description}\label{description-1009}

SCDATAB defines the well connection scale deposition rate as a function
of sea water flow rate, for when the Scale Deposition option has been
activated by declaring the dimensions of the scaling deposition tables
using the SCDPDIMS keyword in the RUNSPEC section. The SCDATAB tables
are allocated to individual wells using the WSCTAB keyword and the sea
water fraction is based on a water tracer entered via the SCDPTRAC
keyword; both keywords are in the SCHEDULE section.

This keyword is not supported by OPM Flow but would change the results
if supported so the simulation will be stopped.

\subsection{SCDPTRAC -- Allocate Sea Water Tracer for Scale
Deposition}\label{scdptrac-allocate-sea-water-tracer-for-scale-deposition}

\begin{longtable}[]{@{}
  >{\raggedright\arraybackslash}p{(\linewidth - 14\tabcolsep) * \real{0.1250}}
  >{\raggedright\arraybackslash}p{(\linewidth - 14\tabcolsep) * \real{0.1250}}
  >{\raggedright\arraybackslash}p{(\linewidth - 14\tabcolsep) * \real{0.1250}}
  >{\raggedright\arraybackslash}p{(\linewidth - 14\tabcolsep) * \real{0.1250}}
  >{\raggedright\arraybackslash}p{(\linewidth - 14\tabcolsep) * \real{0.1250}}
  >{\raggedright\arraybackslash}p{(\linewidth - 14\tabcolsep) * \real{0.1250}}
  >{\raggedright\arraybackslash}p{(\linewidth - 14\tabcolsep) * \real{0.1250}}
  >{\raggedright\arraybackslash}p{(\linewidth - 14\tabcolsep) * \real{0.1250}}@{}}
\toprule\noalign{}
\begin{minipage}[b]{\linewidth}\raggedright
RUNSPEC
\end{minipage} & \begin{minipage}[b]{\linewidth}\raggedright
GRID
\end{minipage} & \begin{minipage}[b]{\linewidth}\raggedright
EDIT
\end{minipage} & \begin{minipage}[b]{\linewidth}\raggedright
PROPS
\end{minipage} & \begin{minipage}[b]{\linewidth}\raggedright
REGIONS
\end{minipage} & \begin{minipage}[b]{\linewidth}\raggedright
SOLUTION
\end{minipage} & \begin{minipage}[b]{\linewidth}\raggedright
SUMMARY
\end{minipage} & \begin{minipage}[b]{\linewidth}\raggedright
SCHEDULE
\end{minipage} \\
\midrule\noalign{}
\endhead
\bottomrule\noalign{}
\endlastfoot
\end{longtable}

\subsubsection{Description}\label{description-1010}

The SCDPTRAC keyword is used to allocate an existing passive water
tracer defined by the TRACER keyword in the PROPS section, to represent
the sea water flowing into a well connection as a fraction of the total
water influx. The keyword is used together with the SCDPTAB keyword in
the SCHEDULE section to calculated the volume of scale deposited around
the well connections.

This keyword is not supported by OPM Flow but would change the results
if supported so the simulation will be stopped.

\begin{longtable}[]{@{}
  >{\raggedright\arraybackslash}p{(\linewidth - 6\tabcolsep) * \real{0.2500}}
  >{\raggedright\arraybackslash}p{(\linewidth - 6\tabcolsep) * \real{0.2500}}
  >{\raggedright\arraybackslash}p{(\linewidth - 6\tabcolsep) * \real{0.2500}}
  >{\raggedright\arraybackslash}p{(\linewidth - 6\tabcolsep) * \real{0.2500}}@{}}
\toprule\noalign{}
\begin{minipage}[b]{\linewidth}\raggedright
No.
\end{minipage} & \begin{minipage}[b]{\linewidth}\raggedright
Name
\end{minipage} & \begin{minipage}[b]{\linewidth}\raggedright
Description
\end{minipage} & \begin{minipage}[b]{\linewidth}\raggedright
Default
\end{minipage} \\
\midrule\noalign{}
\endhead
\bottomrule\noalign{}
\endlastfoot
Field & Metric & Laboratory & \\
1 & NAME & A three letter character string defining the tracer's name
that has previously been defined by the TRACER keyword in the PROPS
section & None \\
Notes: & & & \\
\end{longtable}

\emph{Table 12.63: SCDPTRAC Keyword Description}

\subsubsection{Example}\label{example-548}

In the PROPS section define a tracer in the water phase, for example:

\begin{verbatim}
--
--       DEFINE TRACER NAMES
--
--       TRACER   TRACER
--       NAME     PHASE
--       ------   ------
TRACER
        'SEA'     'WAT'                                    / SEA WATER TRACER

/
\end{verbatim}

Then in the SCHEDULE section allocate the previously defined water
tracer as a sea water tracer to be used with the scale deposition
facility, that is:

\begin{verbatim}
--
--       ALLOCATE SEA WATER TRACER FOR SCALE DEPOSITION
--
--       TRACER
--       NAME
--       ------
SCDPTRAC
        'SEA'                                              / SEA WATER TRACER

/
\end{verbatim}

\subsection{SCHEDULE -- Define the Start of the SCHEDULE Section of
Keywords}\label{schedule-define-the-start-of-the-schedule-section-of-keywords}

\begin{longtable}[]{@{}
  >{\raggedright\arraybackslash}p{(\linewidth - 14\tabcolsep) * \real{0.1250}}
  >{\raggedright\arraybackslash}p{(\linewidth - 14\tabcolsep) * \real{0.1250}}
  >{\raggedright\arraybackslash}p{(\linewidth - 14\tabcolsep) * \real{0.1250}}
  >{\raggedright\arraybackslash}p{(\linewidth - 14\tabcolsep) * \real{0.1250}}
  >{\raggedright\arraybackslash}p{(\linewidth - 14\tabcolsep) * \real{0.1250}}
  >{\raggedright\arraybackslash}p{(\linewidth - 14\tabcolsep) * \real{0.1250}}
  >{\raggedright\arraybackslash}p{(\linewidth - 14\tabcolsep) * \real{0.1250}}
  >{\raggedright\arraybackslash}p{(\linewidth - 14\tabcolsep) * \real{0.1250}}@{}}
\toprule\noalign{}
\begin{minipage}[b]{\linewidth}\raggedright
RUNSPEC
\end{minipage} & \begin{minipage}[b]{\linewidth}\raggedright
GRID
\end{minipage} & \begin{minipage}[b]{\linewidth}\raggedright
EDIT
\end{minipage} & \begin{minipage}[b]{\linewidth}\raggedright
PROPS
\end{minipage} & \begin{minipage}[b]{\linewidth}\raggedright
REGIONS
\end{minipage} & \begin{minipage}[b]{\linewidth}\raggedright
SOLUTION
\end{minipage} & \begin{minipage}[b]{\linewidth}\raggedright
SUMMARY
\end{minipage} & \begin{minipage}[b]{\linewidth}\raggedright
SCHEDULE
\end{minipage} \\
\midrule\noalign{}
\endhead
\bottomrule\noalign{}
\endlastfoot
\end{longtable}

\subsubsection{Description}\label{description-1011}

The SCHEDULE activation keyword marks the end of the SUMMARY section and
the start of the SCHEDULE section that defines the group and well
definitions, operating and economic constraints, as well as how OPM Flow
should advance through time. Numerical controls are also defined in this
section and all parameters can be varied through time.

There is no data required for this keyword.

\subsubsection{Example}\label{example-549}

\begin{verbatim}
-- ==============================================================================
--
-- SCHEDULE SECTION
--
-- ==============================================================================
SCHEDULE
\end{verbatim}

The above example marks the end of the SUMMARY section and the start of
the SCHEDULE section in the OPM Flow data input file.

\subsection{SEPVALS -- Define Separator Oil Formation Volume Factor and
GOR}\label{sepvals-define-separator-oil-formation-volume-factor-and-gor}

\begin{longtable}[]{@{}
  >{\raggedright\arraybackslash}p{(\linewidth - 14\tabcolsep) * \real{0.1250}}
  >{\raggedright\arraybackslash}p{(\linewidth - 14\tabcolsep) * \real{0.1250}}
  >{\raggedright\arraybackslash}p{(\linewidth - 14\tabcolsep) * \real{0.1250}}
  >{\raggedright\arraybackslash}p{(\linewidth - 14\tabcolsep) * \real{0.1250}}
  >{\raggedright\arraybackslash}p{(\linewidth - 14\tabcolsep) * \real{0.1250}}
  >{\raggedright\arraybackslash}p{(\linewidth - 14\tabcolsep) * \real{0.1250}}
  >{\raggedright\arraybackslash}p{(\linewidth - 14\tabcolsep) * \real{0.1250}}
  >{\raggedright\arraybackslash}p{(\linewidth - 14\tabcolsep) * \real{0.1250}}@{}}
\toprule\noalign{}
\begin{minipage}[b]{\linewidth}\raggedright
RUNSPEC
\end{minipage} & \begin{minipage}[b]{\linewidth}\raggedright
GRID
\end{minipage} & \begin{minipage}[b]{\linewidth}\raggedright
EDIT
\end{minipage} & \begin{minipage}[b]{\linewidth}\raggedright
PROPS
\end{minipage} & \begin{minipage}[b]{\linewidth}\raggedright
REGIONS
\end{minipage} & \begin{minipage}[b]{\linewidth}\raggedright
SOLUTION
\end{minipage} & \begin{minipage}[b]{\linewidth}\raggedright
SUMMARY
\end{minipage} & \begin{minipage}[b]{\linewidth}\raggedright
SCHEDULE
\end{minipage} \\
\midrule\noalign{}
\endhead
\bottomrule\noalign{}
\endlastfoot
\end{longtable}

\subsubsection{Description}\label{description-1012}

The SEPVALS keyword defines the initial and subsequent separator oil
formation volume factor (Bo) and Gas Oil Ratio (``GOR'' or Rs). The
facility is used in black-oil modeling to re-scale the PVT data entered
via the PROPS section, based on the saturation point oil formation
volume factor (Bob) and the initial saturated gas-oil ratio (Rsi)
entered on the SEPVALS keyword. The first occurrence of this keyword
sets the initial conditions and must be followed by the GSEPCOND keyword
that assigns previously defined separators to a group.

Note that the keyword can only be used in runs with oil and dissolve gas
only, with no vaporized oil (condensate) in the gas phase.

This keyword is not supported by OPM Flow but would change the results
if supported so the simulation will be stopped.

\subsection{SHRATE -- Activate and Define the Polymer Shearing
Logarithmic
Parameters}\label{shrate-activate-and-define-the-polymer-shearing-logarithmic-parameters}

This keyword activates the logarithm-based polymer shear
thinning/thickening option and defines the shear rate constant. This
keyword can only be used in conjunction with the PLYSHLOG in the PROPS
section

See SHRATE - Activate and Define the Polymer Shearing Logarithmic
Parameters in the PROPS section for a full description.

\subsection{SIMULATE -- Activate the Simulation
Mode}\label{simulate-activate-the-simulation-mode}

\begin{longtable}[]{@{}
  >{\raggedright\arraybackslash}p{(\linewidth - 14\tabcolsep) * \real{0.1250}}
  >{\raggedright\arraybackslash}p{(\linewidth - 14\tabcolsep) * \real{0.1250}}
  >{\raggedright\arraybackslash}p{(\linewidth - 14\tabcolsep) * \real{0.1250}}
  >{\raggedright\arraybackslash}p{(\linewidth - 14\tabcolsep) * \real{0.1250}}
  >{\raggedright\arraybackslash}p{(\linewidth - 14\tabcolsep) * \real{0.1250}}
  >{\raggedright\arraybackslash}p{(\linewidth - 14\tabcolsep) * \real{0.1250}}
  >{\raggedright\arraybackslash}p{(\linewidth - 14\tabcolsep) * \real{0.1250}}
  >{\raggedright\arraybackslash}p{(\linewidth - 14\tabcolsep) * \real{0.1250}}@{}}
\toprule\noalign{}
\begin{minipage}[b]{\linewidth}\raggedright
RUNSPEC
\end{minipage} & \begin{minipage}[b]{\linewidth}\raggedright
GRID
\end{minipage} & \begin{minipage}[b]{\linewidth}\raggedright
EDIT
\end{minipage} & \begin{minipage}[b]{\linewidth}\raggedright
PROPS
\end{minipage} & \begin{minipage}[b]{\linewidth}\raggedright
REGIONS
\end{minipage} & \begin{minipage}[b]{\linewidth}\raggedright
SOLUTION
\end{minipage} & \begin{minipage}[b]{\linewidth}\raggedright
SUMMARY
\end{minipage} & \begin{minipage}[b]{\linewidth}\raggedright
SCHEDULE
\end{minipage} \\
\midrule\noalign{}
\endhead
\bottomrule\noalign{}
\endlastfoot
\end{longtable}

\subsubsection{Description}\label{description-1013}

SIMULATE switches the mode of the simulation to run simulation mode from
the data input checking mode activated by the NOSIM keyword in the
SCHEDULE section. Note that if NOSIM has been used in the RUNSPEC
section then SIMULATE will have no effect.

There is no data required for this keyword and there is no terminating
``/'' for this keyword.

This keyword is not supported by OPM Flow but would change the results
if supported so the simulation will be stopped.

\subsubsection{Example}\label{example-550}

The example below switches OPM Flow to no simulation mode for data
checking of the input deck.

\begin{verbatim}
--
--       ACTIVATE SIMULATION MODE TO RUN THE MODEL
--
SIMULATE
\end{verbatim}

\subsection{SKIP -- Activate Skipping of All Keywords and Input
Data}\label{skip-activate-skipping-of-all-keywords-and-input-data-8}

The SKIP keyword activates skipping of all keywords and input data until
the ENDSKIP keyword is encountered. All keywords between the SKIP and
ENDSKIP keywords are ignored.

See SKIP -- Activate Skipping of All Keywords and Input Data in the
GLOBAL section for a full description.

\subsection{SKIP100 -- Activate Skipping of Black-Oil Keywords and Input
Data}\label{skip100-activate-skipping-of-black-oil-keywords-and-input-data-7}

The SKIP100 keyword activates skipping of all keywords and input data by
the commercial black-oil simulator until the ENDSKIP keyword is
encountered. All keywords between the SKIP100 and ENDSKIP keywords are
ignored by the commercial black-oil simulator. The SKIP100 keyword is
ignored by the commercial compositional simulator.

See SKIP100 -- Activate Skipping of Keywords by Black-Oil Simulator in
the GLOBAL section for a full description.

\subsection{SKIP300 -- Activate Skipping of Keywords by Compositional
Simulator}\label{skip300-activate-skipping-of-keywords-by-compositional-simulator-7}

The SKIP300 keyword activates skipping of all keywords and input data by
the commercial compositional simulator until the ENDSKIP keyword is
encountered. All keywords between the SKIP300 and ENDSKIP keywords are
ignored by the commercial compositional simulator. The SKIP300 keyword
is ignored by the commercial black-oil simulator.

See Error: Reference source not found in the GLOBAL section for a full
description.

\subsection{SKIPREST -- Activate Skipping of Restart Schedule
Data}\label{skiprest-activate-skipping-of-restart-schedule-data}

\begin{longtable}[]{@{}
  >{\raggedright\arraybackslash}p{(\linewidth - 14\tabcolsep) * \real{0.1250}}
  >{\raggedright\arraybackslash}p{(\linewidth - 14\tabcolsep) * \real{0.1250}}
  >{\raggedright\arraybackslash}p{(\linewidth - 14\tabcolsep) * \real{0.1250}}
  >{\raggedright\arraybackslash}p{(\linewidth - 14\tabcolsep) * \real{0.1250}}
  >{\raggedright\arraybackslash}p{(\linewidth - 14\tabcolsep) * \real{0.1250}}
  >{\raggedright\arraybackslash}p{(\linewidth - 14\tabcolsep) * \real{0.1250}}
  >{\raggedright\arraybackslash}p{(\linewidth - 14\tabcolsep) * \real{0.1250}}
  >{\raggedright\arraybackslash}p{(\linewidth - 14\tabcolsep) * \real{0.1250}}@{}}
\toprule\noalign{}
\begin{minipage}[b]{\linewidth}\raggedright
RUNSPEC
\end{minipage} & \begin{minipage}[b]{\linewidth}\raggedright
GRID
\end{minipage} & \begin{minipage}[b]{\linewidth}\raggedright
EDIT
\end{minipage} & \begin{minipage}[b]{\linewidth}\raggedright
PROPS
\end{minipage} & \begin{minipage}[b]{\linewidth}\raggedright
REGIONS
\end{minipage} & \begin{minipage}[b]{\linewidth}\raggedright
SOLUTION
\end{minipage} & \begin{minipage}[b]{\linewidth}\raggedright
SUMMARY
\end{minipage} & \begin{minipage}[b]{\linewidth}\raggedright
SCHEDULE
\end{minipage} \\
\midrule\noalign{}
\endhead
\bottomrule\noalign{}
\endlastfoot
\end{longtable}

\subsubsection{Description}\label{description-1014}

This keyword turns on skipping of keywords up to the start of the
restart point, as defined on the RESTART keyword in the RUNSPEC section.
The RESTART keyword defines the parameters to restart the simulation
from a previous run that has written a RESTART file out to disk.
Activating the SKIPREST keyword causes the simulator to only read in
data it requires for restarting the run up to the RESTART point (RSNUM
on the RESTART keyword in the SOLUTION section). Note that certain
keywords always need to be present in a restart run in the SCHEDULE
section as the data is not stored on the RESTART file, for example the
VFP tables (VFPPROD and VFPINJ keywords). The SKIPREST keyword
automatically processes the input deck and reads the required data.

There is no data required for this keyword.

\subsubsection{Example}\label{example-551}

The example below defines a restart from the previously run NOR-OPM-A01
case at time step number 40.

\begin{verbatim}
-- ==============================================================================
--
-- SOLUTION SECTION
--
-- ==============================================================================
SOLUTION
--
--       FLEXIBLE RESTART FROM PREVIOUS SIMULATION RUN
--
--       FILE                   RESTART   RESTART   FILE
--       NAME                   NUMBER    TYPE      FORMAT
RESTART
         'NOR-OPM-A01'          40        1*        1*     /
\end{verbatim}

Then in the SCHEDULE section the SKIPREST keyword is used to correctly
read in the schedule data up to the RESTART point.

\begin{verbatim}
-- ==============================================================================
--
-- SCHEDULE SECTION
--
-- ==============================================================================
SCHEDULE
--
--       ACTIVATE SKIPREST OPTION TO AVOID MODIFYING SCHEDULE SECTION
--
SKIPREST
\end{verbatim}

\subsection{SLAVES -- Define Slave Reservoir Simulation
Parameters}\label{slaves-define-slave-reservoir-simulation-parameters}

\begin{longtable}[]{@{}
  >{\raggedright\arraybackslash}p{(\linewidth - 14\tabcolsep) * \real{0.1250}}
  >{\raggedright\arraybackslash}p{(\linewidth - 14\tabcolsep) * \real{0.1250}}
  >{\raggedright\arraybackslash}p{(\linewidth - 14\tabcolsep) * \real{0.1250}}
  >{\raggedright\arraybackslash}p{(\linewidth - 14\tabcolsep) * \real{0.1250}}
  >{\raggedright\arraybackslash}p{(\linewidth - 14\tabcolsep) * \real{0.1250}}
  >{\raggedright\arraybackslash}p{(\linewidth - 14\tabcolsep) * \real{0.1250}}
  >{\raggedright\arraybackslash}p{(\linewidth - 14\tabcolsep) * \real{0.1250}}
  >{\raggedright\arraybackslash}p{(\linewidth - 14\tabcolsep) * \real{0.1250}}@{}}
\toprule\noalign{}
\begin{minipage}[b]{\linewidth}\raggedright
RUNSPEC
\end{minipage} & \begin{minipage}[b]{\linewidth}\raggedright
GRID
\end{minipage} & \begin{minipage}[b]{\linewidth}\raggedright
EDIT
\end{minipage} & \begin{minipage}[b]{\linewidth}\raggedright
PROPS
\end{minipage} & \begin{minipage}[b]{\linewidth}\raggedright
REGIONS
\end{minipage} & \begin{minipage}[b]{\linewidth}\raggedright
SOLUTION
\end{minipage} & \begin{minipage}[b]{\linewidth}\raggedright
SUMMARY
\end{minipage} & \begin{minipage}[b]{\linewidth}\raggedright
SCHEDULE
\end{minipage} \\
\midrule\noalign{}
\endhead
\bottomrule\noalign{}
\endlastfoot
\end{longtable}

\subsubsection{Description}\label{description-1015}

This keyword, SLAVES, defines the name of the slave reservoirs and their
associated simulation input files, for when the Reservoir Coupling
option has been declared active by the GRUPMAST and SLAVES keywords in
the SCHEDULE section.

This keyword is not supported by OPM Flow but would change the results
if supported so the simulation will be stopped.

\subsection{SOURCE -- Define Source
Term}\label{source-define-source-term}

\begin{longtable}[]{@{}
  >{\raggedright\arraybackslash}p{(\linewidth - 14\tabcolsep) * \real{0.1250}}
  >{\raggedright\arraybackslash}p{(\linewidth - 14\tabcolsep) * \real{0.1250}}
  >{\raggedright\arraybackslash}p{(\linewidth - 14\tabcolsep) * \real{0.1250}}
  >{\raggedright\arraybackslash}p{(\linewidth - 14\tabcolsep) * \real{0.1250}}
  >{\raggedright\arraybackslash}p{(\linewidth - 14\tabcolsep) * \real{0.1250}}
  >{\raggedright\arraybackslash}p{(\linewidth - 14\tabcolsep) * \real{0.1250}}
  >{\raggedright\arraybackslash}p{(\linewidth - 14\tabcolsep) * \real{0.1250}}
  >{\raggedright\arraybackslash}p{(\linewidth - 14\tabcolsep) * \real{0.1250}}@{}}
\toprule\noalign{}
\begin{minipage}[b]{\linewidth}\raggedright
RUNSPEC
\end{minipage} & \begin{minipage}[b]{\linewidth}\raggedright
GRID
\end{minipage} & \begin{minipage}[b]{\linewidth}\raggedright
EDIT
\end{minipage} & \begin{minipage}[b]{\linewidth}\raggedright
PROPS
\end{minipage} & \begin{minipage}[b]{\linewidth}\raggedright
REGIONS
\end{minipage} & \begin{minipage}[b]{\linewidth}\raggedright
SOLUTION
\end{minipage} & \begin{minipage}[b]{\linewidth}\raggedright
SUMMARY
\end{minipage} & \begin{minipage}[b]{\linewidth}\raggedright
SCHEDULE
\end{minipage} \\
\midrule\noalign{}
\endhead
\bottomrule\noalign{}
\endlastfoot
\end{longtable}

\subsubsection{Description}\label{description-1016}

The \hyperref[REF_HEADING_KEYWORD_SOURCE_12_3]{SOURCE} keyword defines a
mass and heat source term for a given pseudo component in the specified
grid block.

One source term can be added to each grid block for each pseudo
component active in the model. If the source term for a given pseudo
component is later redefined in the same block this will replace the
previous source term.

This is an OPM Flow specific keyword that is not supported by the
commercial simulator.

\begin{longtable}[]{@{}
  >{\raggedright\arraybackslash}p{(\linewidth - 6\tabcolsep) * \real{0.2500}}
  >{\raggedright\arraybackslash}p{(\linewidth - 6\tabcolsep) * \real{0.2500}}
  >{\raggedright\arraybackslash}p{(\linewidth - 6\tabcolsep) * \real{0.2500}}
  >{\raggedright\arraybackslash}p{(\linewidth - 6\tabcolsep) * \real{0.2500}}@{}}
\toprule\noalign{}
\begin{minipage}[b]{\linewidth}\raggedright
No.
\end{minipage} & \begin{minipage}[b]{\linewidth}\raggedright
Name
\end{minipage} & \begin{minipage}[b]{\linewidth}\raggedright
Description
\end{minipage} & \begin{minipage}[b]{\linewidth}\raggedright
Default
\end{minipage} \\
\midrule\noalign{}
\endhead
\bottomrule\noalign{}
\endlastfoot
Field & Metric & Laboratory & \\
1 & I & A positive integer greater than zero and less than or equal to
NX that defines the location of the source term in the I-direction. &
None \\
2 & J & A positive integer greater than zero and less than or equal to
NY that defines the location of the source term in the J-direction. &
None \\
3 & K & A positive integer greater than zero and less than or equal to
NZ that defines the location of the source term in the K-direction. &
None \\
4 & COMP & A defined character string that defines the pseudo component
to be injected by the source term, COMP should be set to one of the
following character strings: The SOLVENT or POLYMER option should only
be specified if the solvent or polymer model has been activated by the
SOLVENT or POLYMER keyword respectively in the RUNSPEC section. MICR,
OXYG, or UREA requires the MICP keyword. A pseudo component must be
specified if the mass injection rate is non-zero. & NONE \\
5 & RATE & A real value that defines the source term's mass injection
rate of the specified component COMP. & 0.0 \\
lb/day & kg/day & gm/hour & \\
6 & HRATE & A real value that defines the source term's heat injection
rate1. The heat injection rate should only be specified if the thermal
model has been activated by the THERMAL keyword in the RUNSPEC section.
& 0.0 \\
Btu/day & kJ/day & J/hour & \\
7 & TEMP & A real value that defines the temperature of the injected
component1. The temperature of the source term should only be specified
if the thermal model has been activated by the THERMAL keyword in the
RUNSPEC section. & None \\
°F & °C & °C & \\
Notes: & & & \\
\end{longtable}

\emph{Table 12.3.223.1: SOURCE Keyword Description}

See also the \hyperref[REF_HEADING_KEYWORD_BCCON]{BCCON} keyword in the
GRID section and the \hyperref[REF_HEADING_KEYWORD_BCPROP]{BCPROP}
keyword in the SCHEDULE section that can be used to specify boundary
conditions.

\subsubsection{Example}\label{example-552}

The first example adds two source terms to block (I=2, J=4, K=6) that
inject gas at a mass flow rate of 25,000 lb/day and inject water at a
mass flow rate of 1,000 lb/day respectively.

\begin{verbatim}
--
--       DEFINE SOURCE TERM
--
-- I    J    K    COMP     MASS     HEAT    TEMP
-- --------------          RATE     RATE
SOURCE
   2    4    6    GAS      25000    1*      1*       /
   2    4    6    WATER     1000    1*      1*       /
/
\end{verbatim}

The second example adds a source term to block (I=2, J=4, K=6) that
injects gas at a mass flow rate of 25,000 lb/day and a heat injection
rate of 1.25 × 106 Btu/day in a thermal black-oil simulation.

\begin{verbatim}
--
--       DEFINE SOURCE TERM
--
-- I    J    K    COMP     MASS     HEAT    TEMP
-- --------------          RATE     RATE
SOURCE
   2    4    6    GAS      25000    1.25E6      1*       /
/
\end{verbatim}

The third example adds a source term to block (I=2, J=4, K=6) that
injects gas at a mass flow rate of 25,000 lb/day and a temperature of
200°F in a thermal black-oil simulation.

\begin{verbatim}
--
--       DEFINE SOURCE TERM
--
-- I    J    K    COMP     MASS     HEAT    TEMP
-- --------------          RATE     RATE
SOURCE
   2    4    6    GAS      25000    1*      200      /
/
\end{verbatim}

\subsection{SUMTHIN -- Define SUMMARY DATA Reporting Time
Steps}\label{sumthin-define-summary-data-reporting-time-steps-1}

This keyword defines a time interval for writing out the SUMMARY data to
the SUMMARY file and the Atgeirr Rasmussen 2017-09-22T12:19:56.742970000
AFR We don't write any RSM files (what are those?) RSM David Baxendale
2017-10-02T18:55:45.272000000 DBx Reply to Atgeirr Rasmussen
(09/22/2017, 12:19): ``\ldots{}'' We should write RSM files, as they
contain the production and injection data from the SUMMARY files in
ASCII format. Engineers normally load this file into EXCEL to generate
plots and economic production forecasts. file, if the RUNSUM keyword has
been has also been activated in the SUMMARY section. Only the data for
the first time step in the time interval is written out and the other
time steps are skipped until the next time interval. This enable the
size of the SUMMARY files to be reduced depending on the size of the
time interval. However, the keyword will produce irregular time steps
reports of the SUMMARY data.

See SUMTHIN -- Define SUMMARY Data Reporting Time Steps in the SUMMARY
section for a full description.

\subsection{SURFVISC -- Surfactant Solution Viscosity versus
Concentration}\label{surfvisc-surfactant-solution-viscosity-versus-concentration-1}

SURFSVISC defines the surfactant viscosity relationship of solution
water viscosity with respect to increasing surfactant concentration
within a grid block. The surfactant option must be activated by the
SURFACT keyword in the RUNSPEC section in order to use this keyword.

See SURFVISC -- Surfactant Solution Viscosity versus Concentration in
the PROPS section for a full description.

\subsection{SWINGFAC -- Define Field Gas Contract
Parameters}\label{swingfac-define-field-gas-contract-parameters}

\begin{longtable}[]{@{}
  >{\raggedright\arraybackslash}p{(\linewidth - 14\tabcolsep) * \real{0.1250}}
  >{\raggedright\arraybackslash}p{(\linewidth - 14\tabcolsep) * \real{0.1250}}
  >{\raggedright\arraybackslash}p{(\linewidth - 14\tabcolsep) * \real{0.1250}}
  >{\raggedright\arraybackslash}p{(\linewidth - 14\tabcolsep) * \real{0.1250}}
  >{\raggedright\arraybackslash}p{(\linewidth - 14\tabcolsep) * \real{0.1250}}
  >{\raggedright\arraybackslash}p{(\linewidth - 14\tabcolsep) * \real{0.1250}}
  >{\raggedright\arraybackslash}p{(\linewidth - 14\tabcolsep) * \real{0.1250}}
  >{\raggedright\arraybackslash}p{(\linewidth - 14\tabcolsep) * \real{0.1250}}@{}}
\toprule\noalign{}
\begin{minipage}[b]{\linewidth}\raggedright
RUNSPEC
\end{minipage} & \begin{minipage}[b]{\linewidth}\raggedright
GRID
\end{minipage} & \begin{minipage}[b]{\linewidth}\raggedright
EDIT
\end{minipage} & \begin{minipage}[b]{\linewidth}\raggedright
PROPS
\end{minipage} & \begin{minipage}[b]{\linewidth}\raggedright
REGIONS
\end{minipage} & \begin{minipage}[b]{\linewidth}\raggedright
SOLUTION
\end{minipage} & \begin{minipage}[b]{\linewidth}\raggedright
SUMMARY
\end{minipage} & \begin{minipage}[b]{\linewidth}\raggedright
SCHEDULE
\end{minipage} \\
\midrule\noalign{}
\endhead
\bottomrule\noalign{}
\endlastfoot
\end{longtable}

\subsubsection{Description}\label{description-1017}

This keyword, SWINGFAC, defines the gas contract parameters, swing
factor and the monthly seasonal profile factor, for when there is a
single gas contract being used in the model. The keyword is used with
the Gas Field Operations option which is activated by the GASFIELD
keyword in the RUNSPEC section. Gas contracts are commonly based on a
Daily Contract Quantity (``DCQ'') that determines the gas rate that the
field should be produced at, which is normally expressed as a multiple
of the DCQ, for example 1.33, and is often referred to as the ``swing
factor''. Some gas contracts also define a maximum DCQ (``Max DCQ'')
and/or a minimum take or pay DCQ (``Min DCQ''), as well as seasonal
demand characteristics. For example, gas rates may be set higher in the
winter months in order to meet heating demand compared with summer
months in colder climates, and the opposite in warmer climates where air
conditioning demand is high.

Thus, the DCQ must be calculated first over a contract period, where the
contract period is commonly contract years, or in some instances
contract quarters. This is performed by the simulator using the current
DCQ and checking to see if the (DCQ x Swing Factor) can be satisfied
throughout the current contract period, if not the DCQ is re-calculated
so that the (DCQ x Swing Factor) condition is satisfied. Once this
condition is met, the second and final pass uses the calculated DCQ in
conjunction with the monthly scaling profile data to set the monthly gas
rate for the field:

\begin{equation}\phantomsection\label{eq-12-32}{
{Q}_{\mathit{month}} = \mathit{DCQ}\times {\mathit{SWINGFAC}}_{\mathit{month}}
}\end{equation}

Where:

Qmonth = the monthly gas production target

DCQ = Daily Contract Quantity

SWINGFACmonth = monthly rate scaling factor that takes into account
seasonal demand, etc.

Here the SWINGFAC keyword sets the gas contract parameters for a single
contract at the FIELD group level and is mutually exclusive to GSWINGF
keyword in the SCHEDULE section that allows for different gas contract
parameters to be assign to different groups.

This keyword is not supported by OPM Flow but would change the results
if supported so the simulation will be stopped.

\subsection{TIGHTEN -- Tighten and Relax Numerical
Controls}\label{tighten-tighten-and-relax-numerical-controls}

\begin{longtable}[]{@{}
  >{\raggedright\arraybackslash}p{(\linewidth - 14\tabcolsep) * \real{0.1250}}
  >{\raggedright\arraybackslash}p{(\linewidth - 14\tabcolsep) * \real{0.1250}}
  >{\raggedright\arraybackslash}p{(\linewidth - 14\tabcolsep) * \real{0.1250}}
  >{\raggedright\arraybackslash}p{(\linewidth - 14\tabcolsep) * \real{0.1250}}
  >{\raggedright\arraybackslash}p{(\linewidth - 14\tabcolsep) * \real{0.1250}}
  >{\raggedright\arraybackslash}p{(\linewidth - 14\tabcolsep) * \real{0.1250}}
  >{\raggedright\arraybackslash}p{(\linewidth - 14\tabcolsep) * \real{0.1250}}
  >{\raggedright\arraybackslash}p{(\linewidth - 14\tabcolsep) * \real{0.1250}}@{}}
\toprule\noalign{}
\begin{minipage}[b]{\linewidth}\raggedright
RUNSPEC
\end{minipage} & \begin{minipage}[b]{\linewidth}\raggedright
GRID
\end{minipage} & \begin{minipage}[b]{\linewidth}\raggedright
EDIT
\end{minipage} & \begin{minipage}[b]{\linewidth}\raggedright
PROPS
\end{minipage} & \begin{minipage}[b]{\linewidth}\raggedright
REGIONS
\end{minipage} & \begin{minipage}[b]{\linewidth}\raggedright
SOLUTION
\end{minipage} & \begin{minipage}[b]{\linewidth}\raggedright
SUMMARY
\end{minipage} & \begin{minipage}[b]{\linewidth}\raggedright
SCHEDULE
\end{minipage} \\
\midrule\noalign{}
\endhead
\bottomrule\noalign{}
\endlastfoot
\end{longtable}

\subsubsection{Description}\label{description-1018}

The TIGHTEN keyword tightens up or slackens the numerical controls for
the linear, non-linear and material balance convergence targets and also
tightens or relaxes the maximum values for the aforementioned
parameters. The keyword should be used with caution as it may result in
significantly increasing the run times.

Note that any subsequent use of the TUNING keyword in the SCHEDULE
section will result in resetting the numerical controls. See also the
TIGHTENP in the SCHEDULE section that allows for greater flexibility in
modifying the numerical controls.

This keyword is not supported by OPM Flow but would change the results
if supported so the simulation will be stopped.

\subsection{TIGHTENP -- Tighten and Relax Numerical Controls
Individually}\label{tightenp-tighten-and-relax-numerical-controls-individually}

\begin{longtable}[]{@{}
  >{\raggedright\arraybackslash}p{(\linewidth - 14\tabcolsep) * \real{0.1250}}
  >{\raggedright\arraybackslash}p{(\linewidth - 14\tabcolsep) * \real{0.1250}}
  >{\raggedright\arraybackslash}p{(\linewidth - 14\tabcolsep) * \real{0.1250}}
  >{\raggedright\arraybackslash}p{(\linewidth - 14\tabcolsep) * \real{0.1250}}
  >{\raggedright\arraybackslash}p{(\linewidth - 14\tabcolsep) * \real{0.1250}}
  >{\raggedright\arraybackslash}p{(\linewidth - 14\tabcolsep) * \real{0.1250}}
  >{\raggedright\arraybackslash}p{(\linewidth - 14\tabcolsep) * \real{0.1250}}
  >{\raggedright\arraybackslash}p{(\linewidth - 14\tabcolsep) * \real{0.1250}}@{}}
\toprule\noalign{}
\begin{minipage}[b]{\linewidth}\raggedright
RUNSPEC
\end{minipage} & \begin{minipage}[b]{\linewidth}\raggedright
GRID
\end{minipage} & \begin{minipage}[b]{\linewidth}\raggedright
EDIT
\end{minipage} & \begin{minipage}[b]{\linewidth}\raggedright
PROPS
\end{minipage} & \begin{minipage}[b]{\linewidth}\raggedright
REGIONS
\end{minipage} & \begin{minipage}[b]{\linewidth}\raggedright
SOLUTION
\end{minipage} & \begin{minipage}[b]{\linewidth}\raggedright
SUMMARY
\end{minipage} & \begin{minipage}[b]{\linewidth}\raggedright
SCHEDULE
\end{minipage} \\
\midrule\noalign{}
\endhead
\bottomrule\noalign{}
\endlastfoot
\end{longtable}

\subsubsection{Description}\label{description-1019}

The TIGHTENP keyword is similar to the TIGHTEN keyword in the SCHEDULE
section, in that it tightens up or slackens the numerical controls for
the linear, non-linear and material balance convergence targets and also
tightens or relaxes the maximum values for the aforementioned
parameters. However, TIGHTENP allows for greater flexibility as there
are four parameters on this keyword, as opposed to just one on the
TIGHTEN keyword, that can be used to modify the numerical controls. The
keyword should be used with caution as it may result in significantly
increasing the run times.

Note that any subsequent use of the TUNING keyword in the SCHEDULE
section will result in resetting the numerical controls. See also the
TIGHTEN keyword in the SCHEDULE section that has more limited
flexibility in modifying the numerical controls.

This keyword is not supported by OPM Flow but would change the results
if supported so the simulation will be stopped.

\subsection{TIME -- Advance Simulation by Cumulative Reporting
Time}\label{time-advance-simulation-by-cumulative-reporting-time}

\begin{longtable}[]{@{}
  >{\raggedright\arraybackslash}p{(\linewidth - 14\tabcolsep) * \real{0.1250}}
  >{\raggedright\arraybackslash}p{(\linewidth - 14\tabcolsep) * \real{0.1250}}
  >{\raggedright\arraybackslash}p{(\linewidth - 14\tabcolsep) * \real{0.1250}}
  >{\raggedright\arraybackslash}p{(\linewidth - 14\tabcolsep) * \real{0.1250}}
  >{\raggedright\arraybackslash}p{(\linewidth - 14\tabcolsep) * \real{0.1250}}
  >{\raggedright\arraybackslash}p{(\linewidth - 14\tabcolsep) * \real{0.1250}}
  >{\raggedright\arraybackslash}p{(\linewidth - 14\tabcolsep) * \real{0.1250}}
  >{\raggedright\arraybackslash}p{(\linewidth - 14\tabcolsep) * \real{0.1250}}@{}}
\toprule\noalign{}
\begin{minipage}[b]{\linewidth}\raggedright
RUNSPEC
\end{minipage} & \begin{minipage}[b]{\linewidth}\raggedright
GRID
\end{minipage} & \begin{minipage}[b]{\linewidth}\raggedright
EDIT
\end{minipage} & \begin{minipage}[b]{\linewidth}\raggedright
PROPS
\end{minipage} & \begin{minipage}[b]{\linewidth}\raggedright
REGIONS
\end{minipage} & \begin{minipage}[b]{\linewidth}\raggedright
SOLUTION
\end{minipage} & \begin{minipage}[b]{\linewidth}\raggedright
SUMMARY
\end{minipage} & \begin{minipage}[b]{\linewidth}\raggedright
SCHEDULE
\end{minipage} \\
\midrule\noalign{}
\endhead
\bottomrule\noalign{}
\endlastfoot
\end{longtable}

\subsubsection{Description}\label{description-1020}

This keyword advances the simulation to a given cumulative report time
after which additional keywords may be entered to instruct OPM Flow to
perform additional functions via the SCHEDULE section keywords, or
further TIME keywords may be entered to advance the simulator to the
next report time.

This keyword is not supported by OPM Flow but would change the results
if supported so the simulation will be stopped.

\begin{longtable}[]{@{}
  >{\raggedright\arraybackslash}p{(\linewidth - 6\tabcolsep) * \real{0.2500}}
  >{\raggedright\arraybackslash}p{(\linewidth - 6\tabcolsep) * \real{0.2500}}
  >{\raggedright\arraybackslash}p{(\linewidth - 6\tabcolsep) * \real{0.2500}}
  >{\raggedright\arraybackslash}p{(\linewidth - 6\tabcolsep) * \real{0.2500}}@{}}
\toprule\noalign{}
\begin{minipage}[b]{\linewidth}\raggedright
No.
\end{minipage} & \begin{minipage}[b]{\linewidth}\raggedright
Name
\end{minipage} & \begin{minipage}[b]{\linewidth}\raggedright
Description
\end{minipage} & \begin{minipage}[b]{\linewidth}\raggedright
Default
\end{minipage} \\
\midrule\noalign{}
\endhead
\bottomrule\noalign{}
\endlastfoot
Field & Metric & Laboratory & \\
1 & TIME & A vector of real positive numbers that define the cumulative
length of the of report times. & None \\
days & days & hours & \\
Notes: & & & \\
\end{longtable}

\emph{Table 12.64: TIME Keyword Description}

See also the DATES and TSTEP keyword in the SCHEDULE section. Note that
since OPM Flow uses the standard Gregorian calendar and therefore leap
years are accounted for in the DATES keyword. Thus, it is more accurate
to use the DATES keyword to progress the simulator through time if one
is matching actual production data.

Whenever possible it is a good idea to always set the start date to be
at the beginning of the year, as like most simulators, OPM Flow reports
are always stated at the number of days from the start date (and
sometimes at a given date). If the start date is at the beginning of the
year, then calculating the actual date is relatively straight forward
and simple.

\subsubsection{Examples}\label{examples-84}

The fist example shows how to advance the simulation three years using
the TIME keyword, from the given start date of January 1, 2022 set via
the START keyword in the RUNSPEC section.

\begin{verbatim}
-- ==============================================================================
--
-- SCHEDULE SECTION
--
-- ==============================================================================
SCHEDULE

-- ------------------------------------------------------------------------------
-- SCHEDULE SECTION - 2022-01-01
-- ------------------------------------------------------------------------------
--
--       ADVANCE SIMULATION BY REPORTING TIME
--
TIME
         365.25  730.50  1095.75
/
\end{verbatim}

The second example shows the same advance but using the TSTEP keyword
instead. Atgeirr Rasmussen 2017-09-22T12:22:06.652621000 AFR Again, 2020
has 366 days.

\begin{verbatim}
  David Baxendale
  2017-10-02T19:01:11.673000000
  DBx
  Reply to Atgeirr Rasmussen (09/22/2017, 12:22): "..."
  I changed the year to simplify the example.
\end{verbatim}

\begin{verbatim}
-- ==============================================================================
--
-- SCHEDULE SECTION
--
-- ==============================================================================
SCHEDULE

-- ------------------------------------------------------------------------------
-- SCHEDULE SECTION - 2022-01-01
-- ------------------------------------------------------------------------------
RPTSCHED
          'WELLS=2'    'WELSPECS'    'CPU=2'     FIP=2'                         /
--
--       ADVANCE SIMULATION BY REPORTING TIME
--
TSTEP
         3*365.25
/
\end{verbatim}

Again, if the simulated production targets are actual production data or
the results are going to be used in economic evaluations then the DATES
keyword may be more useful in advancing the simulation via the reporting
time steps, as the exact dates will be honored.

\subsection{TSTEP -- Advance Simulation by Reporting
Time}\label{tstep-advance-simulation-by-reporting-time}

\begin{longtable}[]{@{}
  >{\raggedright\arraybackslash}p{(\linewidth - 14\tabcolsep) * \real{0.1250}}
  >{\raggedright\arraybackslash}p{(\linewidth - 14\tabcolsep) * \real{0.1250}}
  >{\raggedright\arraybackslash}p{(\linewidth - 14\tabcolsep) * \real{0.1250}}
  >{\raggedright\arraybackslash}p{(\linewidth - 14\tabcolsep) * \real{0.1250}}
  >{\raggedright\arraybackslash}p{(\linewidth - 14\tabcolsep) * \real{0.1250}}
  >{\raggedright\arraybackslash}p{(\linewidth - 14\tabcolsep) * \real{0.1250}}
  >{\raggedright\arraybackslash}p{(\linewidth - 14\tabcolsep) * \real{0.1250}}
  >{\raggedright\arraybackslash}p{(\linewidth - 14\tabcolsep) * \real{0.1250}}@{}}
\toprule\noalign{}
\begin{minipage}[b]{\linewidth}\raggedright
RUNSPEC
\end{minipage} & \begin{minipage}[b]{\linewidth}\raggedright
GRID
\end{minipage} & \begin{minipage}[b]{\linewidth}\raggedright
EDIT
\end{minipage} & \begin{minipage}[b]{\linewidth}\raggedright
PROPS
\end{minipage} & \begin{minipage}[b]{\linewidth}\raggedright
REGIONS
\end{minipage} & \begin{minipage}[b]{\linewidth}\raggedright
SOLUTION
\end{minipage} & \begin{minipage}[b]{\linewidth}\raggedright
SUMMARY
\end{minipage} & \begin{minipage}[b]{\linewidth}\raggedright
SCHEDULE
\end{minipage} \\
\midrule\noalign{}
\endhead
\bottomrule\noalign{}
\endlastfoot
\end{longtable}

\subsubsection{Description}\label{description-1021}

This keyword advances the simulation to a given report time after which
additional keywords may be entered to instruct OPM Flow to perform
additional functions via the SCHEDULE section keywords, or further TSTEP
keywords may be entered to advance the simulator to the next report
time.

\begin{longtable}[]{@{}
  >{\raggedright\arraybackslash}p{(\linewidth - 6\tabcolsep) * \real{0.2500}}
  >{\raggedright\arraybackslash}p{(\linewidth - 6\tabcolsep) * \real{0.2500}}
  >{\raggedright\arraybackslash}p{(\linewidth - 6\tabcolsep) * \real{0.2500}}
  >{\raggedright\arraybackslash}p{(\linewidth - 6\tabcolsep) * \real{0.2500}}@{}}
\toprule\noalign{}
\begin{minipage}[b]{\linewidth}\raggedright
No.
\end{minipage} & \begin{minipage}[b]{\linewidth}\raggedright
Name
\end{minipage} & \begin{minipage}[b]{\linewidth}\raggedright
Description
\end{minipage} & \begin{minipage}[b]{\linewidth}\raggedright
Default
\end{minipage} \\
\midrule\noalign{}
\endhead
\bottomrule\noalign{}
\endlastfoot
Field & Metric & Laboratory & \\
1 & TSTEP & A vector of real positive numbers that define the length of
the time intervals to subsequent report steps. Repeat counts may be
used, for example 10*365.25. & None \\
days & days & hours & \\
Notes: & & & \\
\end{longtable}

\emph{Table 12.65: TSTEP Keyword Description}

See also the DATES and TIME keywords in the SCHEDULE section. Note that
since OPM Flow uses the standard Gregorian calendar and therefore leap
years are accounted for in the DATES keyword. Thus, it is more accurate
to use the DATES keyword to progress the simulator through time if one
is matching actual production data.

Whenever possible it is a good idea to always set the start date to be
at the beginning of the year, as like most simulators, OPM Flow reports
are always stated at the number of days from the start date (and
sometimes at a given date). If the start date is at the beginning of the
year, then calculating the actual date is relatively straight forward
and simple.

\subsubsection{Examples}\label{examples-85}

The first example shows how to advance the simulation via the reporting
time steps from the given start date of January 1, 2022 set via the
START keyword in the RUNSPEC section, to the next year, without any
actions or reporting taking place.

\begin{verbatim}
-- ==============================================================================
--
-- SCHEDULE SECTION
--
-- ==============================================================================
SCHEDULE

-- ------------------------------------------------------------------------------
-- SCHEDULE SECTION - 2022-01-01
-- ------------------------------------------------------------------------------
--
--       ADVANCE SIMULATION BY REPORTING TIME
--
--       JAN  FEB  MAR  APR  MAY  JUN  JLY  AUG  SEP  OCT  NOV  DEC
TSTEP
         31   28   31   30   31   30   31   31   30   31   30   31
/
\end{verbatim}

The second example is similar to the previous example but with quarterly
reporting time steps used instead based on
\(\frac{365.25}{4}=91.3125\)days per quarter

\begin{verbatim}
-- ==============================================================================
--
-- SCHEDULE SECTION
--
-- ==============================================================================
SCHEDULE

-- ------------------------------------------------------------------------------
-- SCHEDULE SECTION - 2022-01-01
-- ------------------------------------------------------------------------------
RPTSCHED
          'WELLS=2'    'WELSPECS'    'CPU=2'     FIP=2'                         /

--
--       ADVANCE SIMULATION BY REPORTING TIME
--
--       QUARTERLY
TSTEP
         4*91.3125
/
\end{verbatim}

Again, if the simulated production targets are actual production data or
the results are going to be used in economic evaluations then the DATES
keyword may be more useful in advancing the simulation via the reporting
time steps, as the exact dates will be honored.

\subsection{TUNING -- Numerical Tuning
Control}\label{tuning-numerical-tuning-control}

\begin{longtable}[]{@{}
  >{\raggedright\arraybackslash}p{(\linewidth - 14\tabcolsep) * \real{0.1250}}
  >{\raggedright\arraybackslash}p{(\linewidth - 14\tabcolsep) * \real{0.1250}}
  >{\raggedright\arraybackslash}p{(\linewidth - 14\tabcolsep) * \real{0.1250}}
  >{\raggedright\arraybackslash}p{(\linewidth - 14\tabcolsep) * \real{0.1250}}
  >{\raggedright\arraybackslash}p{(\linewidth - 14\tabcolsep) * \real{0.1250}}
  >{\raggedright\arraybackslash}p{(\linewidth - 14\tabcolsep) * \real{0.1250}}
  >{\raggedright\arraybackslash}p{(\linewidth - 14\tabcolsep) * \real{0.1250}}
  >{\raggedright\arraybackslash}p{(\linewidth - 14\tabcolsep) * \real{0.1250}}@{}}
\toprule\noalign{}
\begin{minipage}[b]{\linewidth}\raggedright
RUNSPEC
\end{minipage} & \begin{minipage}[b]{\linewidth}\raggedright
GRID
\end{minipage} & \begin{minipage}[b]{\linewidth}\raggedright
EDIT
\end{minipage} & \begin{minipage}[b]{\linewidth}\raggedright
PROPS
\end{minipage} & \begin{minipage}[b]{\linewidth}\raggedright
REGIONS
\end{minipage} & \begin{minipage}[b]{\linewidth}\raggedright
SOLUTION
\end{minipage} & \begin{minipage}[b]{\linewidth}\raggedright
SUMMARY
\end{minipage} & \begin{minipage}[b]{\linewidth}\raggedright
SCHEDULE
\end{minipage} \\
\midrule\noalign{}
\endhead
\bottomrule\noalign{}
\endlastfoot
\end{longtable}

\subsubsection{Description}\label{description-1022}

The TUNING keyword defines the parameters used to control the commercial
simulator's time stepping and numerical convergence for the global grid.
The keyword is similar to the TUNINGDP keyword in the SCHEDULE section
that is optimized for high throughput runs. The keyword is mostly
ignored by OPM Flow; however, the simulator can be instructed to read
some parameters from the TUNING keyword if the appropriate command line
parameter has been activated (see section 2.2 Running OPM Flow 2023-04
From The Command Line).

\begin{longtable}[]{@{}
  >{\raggedright\arraybackslash}p{(\linewidth - 6\tabcolsep) * \real{0.2500}}
  >{\raggedright\arraybackslash}p{(\linewidth - 6\tabcolsep) * \real{0.2500}}
  >{\raggedright\arraybackslash}p{(\linewidth - 6\tabcolsep) * \real{0.2500}}
  >{\raggedright\arraybackslash}p{(\linewidth - 6\tabcolsep) * \real{0.2500}}@{}}
\toprule\noalign{}
\begin{minipage}[b]{\linewidth}\raggedright
No.
\end{minipage} & \begin{minipage}[b]{\linewidth}\raggedright
Name
\end{minipage} & \begin{minipage}[b]{\linewidth}\raggedright
Description
\end{minipage} & \begin{minipage}[b]{\linewidth}\raggedright
Default
\end{minipage} \\
\midrule\noalign{}
\endhead
\bottomrule\noalign{}
\endlastfoot
Field & Metric & Laboratory & \\
1-1 & TSINIT & TSINIT is a positive real value that defines the maximum
length of the next time step. Note that whenever the keyword is used
TSINIT is always set back to the default value of one, unless explicitly
over written. & 1.0 \\
days & days & hours & \\
1-2 & TSMAXZ & TSMAXZ is a positive real value that defines the maximum
length of time steps following the next time step (TSINIT). & 365.0 \\
days & days & hours & \\
1-3 & TSMINZ & TSMINZ is a positive real value that defines the minimum
length of all time steps. & 0.1 \\
days & days & hours & \\
1-4 & TSMCHP & TSMCHP is a positive real value that sets the minimum
time step length that can be chopped. & 0.15 \\
days & days & hours & \\
1-5 & TSFMAX & TSFMAX is a positive real value that specifies the
maximum growth rate a time step can be increased by, subject to the
maximum allowable time step size set by TSMAXZ. For example, if the
current time step has converged at 10 days and TSFMAX is set to the
default value, then the next time step will be 3.0 x 10 days, that is 30
days provided it is less than TSMAXZ. & 3.0 \\
dimensionless & dimensionless & dimensionless & \\
1-6 & TSFMIN & TSFMIN is a positive real value that specifies the
minimum decay rate a time step can be decreased by, subject to the
minimum allowable time step size set by TSMINZ. For example, if the
current time step has not converged at 10 days and TSFMIN is set to the
default value, then the next time step will be 0.3 x 10 days, that is
the maximum of 3 days and TSMINZ. & 0.3 \\
dimensionless & dimensionless & dimensionless & \\
1-7 & TSFCNV & TSFCNV is a positive real value that specifies the decay
rate a time step can be decreased by after a convergence failure. &
0.1 \\
dimensionless & dimensionless & dimensionless & \\
1-8 & TFDIFF & TFDIFF is a positive real value that sets the time step
growth factor of the time step after a convergence failure. For example,
if the chopped current convergent time step is 10 days and TFDIFF is set
to the default value, then the time step will be increased to 1.25 x 10
days, that is the minimum of 12.5 days and TSMAXZ. & 1.25 \\
dimensionless & dimensionless & dimensionless & \\
1-9 & THRUPT & THRURT is a positive real value that specifies the
maximum throughput ratio over a time step. & 1.0 x 1020 \\
dimensionless & dimensionless & dimensionless & \\
1-10 & TMAXWC & TMAXWC is a positive real value that defines maximum
allowed time step after a well event; for example, when a well is opened
or closed, etc. & None \\
days & days & hours & \\
1-11 & / & Record terminated by a ``/'' & Not Applicable \\
2-1 & TRGTTE & TRGTTE is a positive real value that sets the target time
truncation error. & 0.1 \\
dimensionless & dimensionless & dimensionless & \\
2-2 & TRGCNV & TRGCNV is a positive real value that defines the target
non-linear convergence error. & 0.001 \\
dimensionless & dimensionless & dimensionless & \\
2-3 & TRGMBE & TRGMBE is a positive real value that specifies the target
material balance error. & 1.0 x 10-7 \\
dimensionless & dimensionless & dimensionless & \\
2-4 & TRGLCV & TRGLCV is a positive real value that specifies the target
linear convergence error. & 0.0001 \\
dimensionless & dimensionless & dimensionless & \\
2-5 & XXXTTE & XXXTTE is a positive real value that sets the maximum
time truncation error. & 10.0 \\
dimensionless & dimensionless & dimensionless & \\
2-6 & XXXCNV & XXXCNV is a positive real value that defines the maximum
non-linear convergence error. & 0.01 \\
dimensionless & dimensionless & dimensionless & \\
2-7 & XXXMBE & XXXMBE is a positive real value that specifies the
maximum mass balance error, that is the tolerated mass balance error
relative to total mass present. & 1.0 x 10-6 \\
dimensionless & dimensionless & dimensionless & \\
2-8 & XXXLCV & XXXLCV is a positive real values that sets the maximum
linear convergence error. & 0.0001 \\
dimensionless & dimensionless & dimensionless & \\
2-9 & XXXWFL & XXXWFL is a positive real values that fixes the maximum
well flow convergence error. & 0.001 \\
dimensionless & dimensionless & dimensionless & \\
2-10 & TRGFIP & TRGFIP is a positive real value that stipulates the
target fluid in-place error in Local Grid Refinements. & 0.025 \\
dimensionless & dimensionless & dimensionless & \\
2-11 & TRGSFT & TRGSFT is a positive real values that defines the target
surfactant change when the Surfactant Model is active in the run. &
None \\
dimensionless & dimensionless & dimensionless & \\
2-12 & THIONX & THIONX is a positive real value used to set the
threshold for damping in the ion echange calculation for when the Brine
Model is active in the run. & 0.01 \\
dimensionless & dimensionless & dimensionless & \\
2-13 & TRWGHT & TRWGHT is an integer that stipulates the implicitness
for active tracer updates within the Newton iterations, and should be
set to: & 1 \\
dimensionless & dimensionless & dimensionless & \\
2-14 & / & Record terminated by a ``/'' & Not Applicable \\
3-1 & NEWTMX & NEWTMX is a positive integer greater or equal to NEWTMN
that stipulates the maximum number of Newton iterations for a time step.
& 12 \\
dimensionless & dimensionless & dimensionless & \\
3-2 & NEWTMN & NEWTMN is a positive integer that is less or equal to
NEWTMX that defines the minimum number of Newton iterations for a time
step. & 1 \\
dimensionless & dimensionless & dimensionless & \\
3-3 & LITMAX & LITMAX is a positive integer greater or equal to LITMIN
that sets the maximum number of linear iterations within a Newton
iteration. & 25 \\
dimensionless & dimensionless & dimensionless & \\
3-4 & LITMIN & LITMIN is a positive integer less or equal to LITMAX that
sets the minimum number of linear iterations within a Newton iteration.
& 1 \\
dimensionless & dimensionless & dimensionless & \\
3-5 & MXWSIT & MXWSIT is a positive integer that defines the maximum
number of iterations within a well flow calculation. & 8 \\
dimensionless & dimensionless & dimensionless & \\
3-6 & MXWPIT & MXWPIT is a positive integer that stipulates the maximum
number of iterations for solving the bottom-hole pressure for wells
under tubing head pressure control within a well flow calculation. &
8 \\
dimensionless & dimensionless & dimensionless & \\
3-7 & DDPLIM & DDPLIM is a positive real value that stipulates the
maximum pressure change at the last Newton iteration. & 1.0 x 106 \\
psia & barsa & atma & \\
3-8 & DDSLIM & DDSLIM is a positive real value that sets the maximum
saturation change at the last Newton iteration. & 1.0 x 106 \\
dimensionless & dimensionless & dimensionless & \\
3-9 & TRGDPR & TRGDPR is a positive real value that defines the target
maximum pressure change within a time step. & 1.0 x 106 \\
psia & barsa & atma & \\
3-10 & XXXDPR & XXXDPR is a positive real value that stipulates the
maximum tolerable pressure change within a time step. & None \\
psia & barsa & atma & \\
3-11 & MNWRFP & MNWRFP is a positive integer greater than one and less
than NEWTMX that defines the minimum number of Newton iterations before
invoking the bisection algorithm for when the polymer phase is active in
the model via the POLYMER keyword in the RUNSPEC section. & 4 \\
dimensionless & dimensionless & dimensionless & \\
3-12 & / & Record terminated by a ``/'' & Not Applicable \\
Notes: & & & \\
\end{longtable}

\emph{Table 12.66: TUNING Keyword Description}

Note that for record number two (items 2-1 to 2-13) the maximum values
should always be greater than the associated target value; for example,
XXXCNV should be greater than TRGCNV. Also note that the TUNING keyword
is stored on the restart files (see RPTRST -- Define Data to be Written
to the RESTART File) enabling the parameters to be utilized in a restart
run without re-specifying the keyword.

\subsubsection{Example}\label{example-553}

\begin{verbatim}
--
--       DEFAULT TUNING PARAMETERS
--
TUNING
         1.0    365.0     0.1     0.15     3     0.3     0.1   1.25  1E20  1*   /
                                                                                /
                                                                                /
\end{verbatim}

The above example explicitly sets the default parameters for OPM Flow
for when the appropriate command line parameter has been activated (see
section 2.2 Running OPM Flow 2023-04 From The Command Line) to instruct
the simulator to read the first record of the TUNING keyword.
Alternatively one could just use the following to accomplish the same
thing.

\begin{verbatim}
TUNING
/
/
/
\end{verbatim}

\subsection{TUNINGDP -- Numerical Tuning Control for High Throughput
Cases}\label{tuningdp-numerical-tuning-control-for-high-throughput-cases}

\begin{longtable}[]{@{}
  >{\raggedright\arraybackslash}p{(\linewidth - 14\tabcolsep) * \real{0.1250}}
  >{\raggedright\arraybackslash}p{(\linewidth - 14\tabcolsep) * \real{0.1250}}
  >{\raggedright\arraybackslash}p{(\linewidth - 14\tabcolsep) * \real{0.1250}}
  >{\raggedright\arraybackslash}p{(\linewidth - 14\tabcolsep) * \real{0.1250}}
  >{\raggedright\arraybackslash}p{(\linewidth - 14\tabcolsep) * \real{0.1250}}
  >{\raggedright\arraybackslash}p{(\linewidth - 14\tabcolsep) * \real{0.1250}}
  >{\raggedright\arraybackslash}p{(\linewidth - 14\tabcolsep) * \real{0.1250}}
  >{\raggedright\arraybackslash}p{(\linewidth - 14\tabcolsep) * \real{0.1250}}@{}}
\toprule\noalign{}
\begin{minipage}[b]{\linewidth}\raggedright
RUNSPEC
\end{minipage} & \begin{minipage}[b]{\linewidth}\raggedright
GRID
\end{minipage} & \begin{minipage}[b]{\linewidth}\raggedright
EDIT
\end{minipage} & \begin{minipage}[b]{\linewidth}\raggedright
PROPS
\end{minipage} & \begin{minipage}[b]{\linewidth}\raggedright
REGIONS
\end{minipage} & \begin{minipage}[b]{\linewidth}\raggedright
SOLUTION
\end{minipage} & \begin{minipage}[b]{\linewidth}\raggedright
SUMMARY
\end{minipage} & \begin{minipage}[b]{\linewidth}\raggedright
SCHEDULE
\end{minipage} \\
\midrule\noalign{}
\endhead
\bottomrule\noalign{}
\endlastfoot
\end{longtable}

\subsubsection{Description}\label{description-1023}

The TUNINGDP keyword defines the parameters used for controlling the
commercial simulator's numerical convergence parameters. This keyword is
similar to the TUNING keyword in the SCHEDULE section, but the defaults
on this keyword are optimized for high throughput runs.

This keyword is generally ignored by OPM Flow; however, the simulator
can be instructed to read some parameters from the TUNINGDP keyword if
the command line option --enable-tuning=true has been used (see section
2.2 Running OPM Flow 2023-04 From The Command Line).

\begin{longtable}[]{@{}
  >{\raggedright\arraybackslash}p{(\linewidth - 6\tabcolsep) * \real{0.2500}}
  >{\raggedright\arraybackslash}p{(\linewidth - 6\tabcolsep) * \real{0.2500}}
  >{\raggedright\arraybackslash}p{(\linewidth - 6\tabcolsep) * \real{0.2500}}
  >{\raggedright\arraybackslash}p{(\linewidth - 6\tabcolsep) * \real{0.2500}}@{}}
\toprule\noalign{}
\begin{minipage}[b]{\linewidth}\raggedright
No.
\end{minipage} & \begin{minipage}[b]{\linewidth}\raggedright
Name
\end{minipage} & \begin{minipage}[b]{\linewidth}\raggedright
Description
\end{minipage} & \begin{minipage}[b]{\linewidth}\raggedright
Default
\end{minipage} \\
\midrule\noalign{}
\endhead
\bottomrule\noalign{}
\endlastfoot
Field & Metric & Laboratory & \\
1 & TRGLCV & TRGLCV is a positive real value that specifies the linear
convergence error target. The default value is ten times lower than the
default value on the TUNING keyword. & 0.00001 \\
dimensionless & dimensionless & dimensionless & \\
2 & XXXLCV & XXXLCV is a positive real values that sets the maximum
linear convergence error. The default value is ten times lower than the
default value on the TUNING keyword. & 0.0001 \\
dimensionless & dimensionless & dimensionless & \\
3 & TRGDDP & TRGDDP a positive real value that stipulates the maximum
pressure change during a Newton iteration that enables the solution to
be accepted when the residual pressure is still outside its convergence
criteria. & 1.0 \\
psia & barsa & atma & \\
4 & TRGDDS & TRGDDS a positive real value that sets the maximum
saturation change during a Newton iteration that enables the solution to
be accepted when the residual saturation is still outside its
convergence criteria. & 0.01 \\
dimensionless & dimensionless & dimensionless & \\
5 & TRGDDRS & TRGDDRS a positive real value that sets the maximum gas
dissolution factor change during a Newton iteration that enables the
solution to be accepted when the residual gas dissolution factor is
still outside its convergence criteria. This is an OPM Flow specific
item that is not supported by the commercial simulator. & 0.0 \\
Mscf/stb & sm3/sm3 & scc/scc & \\
6 & TRGDDRV & TRGDDRV a positive real value that sets the maximum oil
dissolution factor change during a Newton iteration that enables the
solution to be accepted when the residual oil dissolution factor change
is still outside its convergence criteria. This is an OPM Flow specific
item that is not supported by the commercial simulator. & 0.0 \\
stb/Mscf & sm3/sm3 & scc/scc & \\
Notes: & & & \\
\end{longtable}

\emph{Table 12.67: TUNINGDP Keyword Description}

Note that the TUNING keyword is stored on the restart files (see RPTRST
-- Define Data to be Written to the RESTART File) enabling the
parameters to be utilized in a restart run without re-specifying the
keyword.

\subsubsection{Example}\label{example-554}

\begin{verbatim}
--
--       DEFAULT TUNINGDP PARAMETERS
--
TUNINGDP
       /
\end{verbatim}

The above example explicitly sets the default parameters.

\subsection{TUNINGH -- Numerical Tuning Control for History Match
Gradient
Calculations}\label{tuningh-numerical-tuning-control-for-history-match-gradient-calculations}

\begin{longtable}[]{@{}
  >{\raggedright\arraybackslash}p{(\linewidth - 14\tabcolsep) * \real{0.1250}}
  >{\raggedright\arraybackslash}p{(\linewidth - 14\tabcolsep) * \real{0.1250}}
  >{\raggedright\arraybackslash}p{(\linewidth - 14\tabcolsep) * \real{0.1250}}
  >{\raggedright\arraybackslash}p{(\linewidth - 14\tabcolsep) * \real{0.1250}}
  >{\raggedright\arraybackslash}p{(\linewidth - 14\tabcolsep) * \real{0.1250}}
  >{\raggedright\arraybackslash}p{(\linewidth - 14\tabcolsep) * \real{0.1250}}
  >{\raggedright\arraybackslash}p{(\linewidth - 14\tabcolsep) * \real{0.1250}}
  >{\raggedright\arraybackslash}p{(\linewidth - 14\tabcolsep) * \real{0.1250}}@{}}
\toprule\noalign{}
\begin{minipage}[b]{\linewidth}\raggedright
RUNSPEC
\end{minipage} & \begin{minipage}[b]{\linewidth}\raggedright
GRID
\end{minipage} & \begin{minipage}[b]{\linewidth}\raggedright
EDIT
\end{minipage} & \begin{minipage}[b]{\linewidth}\raggedright
PROPS
\end{minipage} & \begin{minipage}[b]{\linewidth}\raggedright
REGIONS
\end{minipage} & \begin{minipage}[b]{\linewidth}\raggedright
SOLUTION
\end{minipage} & \begin{minipage}[b]{\linewidth}\raggedright
SUMMARY
\end{minipage} & \begin{minipage}[b]{\linewidth}\raggedright
SCHEDULE
\end{minipage} \\
\midrule\noalign{}
\endhead
\bottomrule\noalign{}
\endlastfoot
\end{longtable}

\subsubsection{Description}\label{description-1024}

Defines the parameters used for controlling the commercial simulator's
numerical convergence parameters. The keyword is similar to the TUNING
keyword in the SCHEDULE section, but the defaults on this keyword are
optimized for high throughput runs. See section 2.2 Running OPM Flow
2023-04 From The Command Line on how to invoke various numerical schemes
via the OPM Flow command line interface.

This keyword is not supported by OPM Flow but has no effect on the
results so it will be ignored.

\begin{longtable}[]{@{}
  >{\raggedright\arraybackslash}p{(\linewidth - 6\tabcolsep) * \real{0.2500}}
  >{\raggedright\arraybackslash}p{(\linewidth - 6\tabcolsep) * \real{0.2500}}
  >{\raggedright\arraybackslash}p{(\linewidth - 6\tabcolsep) * \real{0.2500}}
  >{\raggedright\arraybackslash}p{(\linewidth - 6\tabcolsep) * \real{0.2500}}@{}}
\toprule\noalign{}
\begin{minipage}[b]{\linewidth}\raggedright
No.
\end{minipage} & \begin{minipage}[b]{\linewidth}\raggedright
Name
\end{minipage} & \begin{minipage}[b]{\linewidth}\raggedright
Description
\end{minipage} & \begin{minipage}[b]{\linewidth}\raggedright
Default
\end{minipage} \\
\midrule\noalign{}
\endhead
\bottomrule\noalign{}
\endlastfoot
Field & Metric & Laboratory & \\
1 & GRGLCV & GRGLCV is a real positive value that specifies the linear
convergence error target. & 0.0001 \\
dimensionless & dimensionless & dimensionless & \\
2 & GXXLCV & GXXLCV is a real positive values that sets the maximum
linear convergence error. & 0.001 \\
dimensionless & dimensionless & dimensionless & \\
3 & GMSLCV & GMSLCV is a real positive value that specifies the linear
convergence residual reduction. & 1.0 x 10-20 \\
dimensionless & dimensionless & dimensionless & \\
4 & LGTMIN & LGTMIN is a positive integer less or equal to LGTMAX that
sets the minimum number of linear iterations within a Newton iteration.
& 1 \\
dimensionless & dimensionless & dimensionless & \\
5 & LGTMAX & LGTMAX is a positive integer greater or equal to LGTMIN
that sets the maximum number of linear iterations within a Newton
iteration. & 25 \\
dimensionless & dimensionless & dimensionless & \\
Notes: & & & \\
\end{longtable}

\emph{Table 12.68: TUNINGH Keyword Description}

\subsubsection{Example}\label{example-555}

\begin{verbatim}
--
--       DEFAULT TUNINGH PARAMETERS
--
TUNINGH
       /
\end{verbatim}

The above example explicitly sets the default parameters.

\subsection{TUNINGL -- Numerical Tuning Control for All
LGRs}\label{tuningl-numerical-tuning-control-for-all-lgrs}

\begin{longtable}[]{@{}
  >{\raggedright\arraybackslash}p{(\linewidth - 14\tabcolsep) * \real{0.1250}}
  >{\raggedright\arraybackslash}p{(\linewidth - 14\tabcolsep) * \real{0.1250}}
  >{\raggedright\arraybackslash}p{(\linewidth - 14\tabcolsep) * \real{0.1250}}
  >{\raggedright\arraybackslash}p{(\linewidth - 14\tabcolsep) * \real{0.1250}}
  >{\raggedright\arraybackslash}p{(\linewidth - 14\tabcolsep) * \real{0.1250}}
  >{\raggedright\arraybackslash}p{(\linewidth - 14\tabcolsep) * \real{0.1250}}
  >{\raggedright\arraybackslash}p{(\linewidth - 14\tabcolsep) * \real{0.1250}}
  >{\raggedright\arraybackslash}p{(\linewidth - 14\tabcolsep) * \real{0.1250}}@{}}
\toprule\noalign{}
\begin{minipage}[b]{\linewidth}\raggedright
RUNSPEC
\end{minipage} & \begin{minipage}[b]{\linewidth}\raggedright
GRID
\end{minipage} & \begin{minipage}[b]{\linewidth}\raggedright
EDIT
\end{minipage} & \begin{minipage}[b]{\linewidth}\raggedright
PROPS
\end{minipage} & \begin{minipage}[b]{\linewidth}\raggedright
REGIONS
\end{minipage} & \begin{minipage}[b]{\linewidth}\raggedright
SOLUTION
\end{minipage} & \begin{minipage}[b]{\linewidth}\raggedright
SUMMARY
\end{minipage} & \begin{minipage}[b]{\linewidth}\raggedright
SCHEDULE
\end{minipage} \\
\midrule\noalign{}
\endhead
\bottomrule\noalign{}
\endlastfoot
\end{longtable}

\subsubsection{Description}\label{description-1025}

TUNINGL defines the parameters used for controlling the commercial
simulator's numerical convergence parameters for all Local Grid
Refinements (``LGR''). The keyword is the same as the TUNING keyword in
the SCHEDULE section that applies the tuning parameters to the global
grid. See section 2.2 Running OPM Flow 2023-04 From The Command Line on
how to invoke various numerical schemes via the OPM Flow command line
interface.

This keyword is not supported by OPM Flow but has no effect on the
results so it will be ignored.

\begin{longtable}[]{@{}
  >{\raggedright\arraybackslash}p{(\linewidth - 6\tabcolsep) * \real{0.2500}}
  >{\raggedright\arraybackslash}p{(\linewidth - 6\tabcolsep) * \real{0.2500}}
  >{\raggedright\arraybackslash}p{(\linewidth - 6\tabcolsep) * \real{0.2500}}
  >{\raggedright\arraybackslash}p{(\linewidth - 6\tabcolsep) * \real{0.2500}}@{}}
\toprule\noalign{}
\begin{minipage}[b]{\linewidth}\raggedright
No.
\end{minipage} & \begin{minipage}[b]{\linewidth}\raggedright
Name
\end{minipage} & \begin{minipage}[b]{\linewidth}\raggedright
Description
\end{minipage} & \begin{minipage}[b]{\linewidth}\raggedright
Default
\end{minipage} \\
\midrule\noalign{}
\endhead
\bottomrule\noalign{}
\endlastfoot
Field & Metric & Laboratory & \\
1-1 & TSINIT & TSINT is a real positive value that defines the maximum
length of the next time step. Note that whenever the keyword is used
TSINIT is always set back to the default value of one, unless explicitly
over written. & 1.0 \\
days & days & hours & \\
1-2 & TSMAXZ & TSMAXZ is a real positive value that defines the maximum
length of the next time step following TSINIT. & 365.0 \\
days & days & hours & \\
1-3 & TSMINZ & TSMINZ is a real positive value that defines the minimum
length of all time steps. & 0.1 \\
days & days & hours & \\
1-4 & TSMCHP & TSMCHP is a real positive values that sets the minimum
length of all chopped time steps. & 0.15 \\
days & days & hours & \\
1-5 & TSFMAX & TSFMAX is a real positive value that specifies the
maximum growth rate a time step can be increased by, subject to the
maximum allowable time step size set by TSMAXZ. For example, if the
current time step has converged at 10 days and TSFMAX is set to the
default value, then the next time step will be 3.0 x 10 days, that is 30
days provided it is less than TSMAXZ. & 3.0 \\
dimensionless & dimensionless & dimensionless & \\
1-6 & TSFMIN & TSFMIN is a real positive value that specifies the
minimum decay rate a time step can be decreased by, subject to the
minimum allowable time step size set by TSMINZ. For example, if the
current time step has not converged at 10 days and TSFMAX is set to the
default value, then the next time step will be 0.3 x 10 days, that is
the maximum of 0.3 days and TSMINZ. & 0.3 \\
dimensionless & dimensionless & dimensionless & \\
1-7 & TSFCNV & TSFCNV real positive value that specifies the decay rate
a time step can be decreased by after the number of target iterations
has been exceeded. & 0.1 \\
dimensionless & dimensionless & dimensionless & \\
1-8 & TFDIFF & TFDIFFA is a real positive value that sets the time step
growth factor of the time step after a convergence failure. For example,
if the chopped current convergent time step is 10 days and TFDIFF is set
to the default value, then the time step will be increased to 1.25 x 10
days, that is the minimum of 11.25 days and TSMAXZ. & 1.25 \\
dimensionless & dimensionless & dimensionless & \\
1-9 & THRURPT & THRURPT is a real positive value that specifies the
maximum throughput ratio over a time step. & 1.0 x 1020 \\
dimensionless & dimensionless & dimensionless & \\
1-10 & TMAXWC & TMAXWC is a real double precision value that defines
maximum allowed time step after a well event; for example, when a well
is opened or closed, etc. & None \\
days & days & hours & \\
1-11 & / & Record terminated by a ``/'' & Not Applicable \\
2-1 & TRGTTE & TRGTTE is a real positive value that sets the time
truncation error target. & 0.1 \\
dimensionless & dimensionless & dimensionless & \\
2-2 & TRGCNV & TRGCNV a real positive value that defines the non-linear
convergence error. & 0.001 \\
dimensionless & dimensionless & dimensionless & \\
2-3 & TRGMBE & TRGMBE is a real positive value that specifies then
target material balance error. & 1.0 x 10-7 \\
dimensionless & dimensionless & dimensionless & \\
2-4 & TRGLCV & TRGLCV is a real positive value that specifies the linear
convergence error target. & 0.00001 \\
dimensionless & dimensionless & dimensionless & \\
2-5 & XXXTTE & XXXTTE is a real positive value that sets the maximum
time truncation error. & 10.0 \\
dimensionless & dimensionless & dimensionless & \\
2-6 & XXXCNV & XXXCNV is a real positive value that defines the maximum
non-linear convergence error. & 0.01 \\
dimensionless & dimensionless & dimensionless & \\
2-7 & XXXMBE & XXXMBE is a real positive value that specifies the
maximum mass balance error, that is the tolerated mass balance error
relative to total mass present. & 1.0 x 10-6 \\
dimensionless & dimensionless & dimensionless & \\
2-8 & XXXLCV & XXXLCV is a real positive values that sets the maximum
linear convergence error. & 0.0001 \\
dimensionless & dimensionless & dimensionless & \\
2-9 & XXXWFL & XXXWFL is a real positive values that fixes the maximum
well flow convergence error. & 0.001 \\
dimensionless & dimensionless & dimensionless & \\
2-10 & TRGFIP & TRGFIP is a real positive value that stipulates the
target fluid in-place error in Local Grid Refinements. & 0.025 \\
dimensionless & dimensionless & dimensionless & \\
2-11 & TRGSFT & TRGSFT is a real positive values that defines the target
surfactant change when the Surfactant Model is active in the run. &
None \\
dimensionless & dimensionless & dimensionless & \\
2-12 & THIONX & THIONX is a positive real value used to set the
threshold for damping in the ion echange calculation for when the Brine
Model is active in the run. & 0.01 \\
dimensionless & dimensionless & dimensionless & \\
2-13 & TRWGHT & TRWGHT is a positive integer that stipulates the
implicitness for active tracer updates within the Newton iterations, and
should be set to: & 1 \\
dimensionless & dimensionless & dimensionless & \\
2-14 & / & Record terminated by a ``/'' & Not Applicable \\
3-1 & NEWTMX & NEWTMX is a positive integer greater or equal to NEWTMN
that stipulates the maximum number of Newton iterations for a time step.
& 12 \\
dimensionless & dimensionless & dimensionless & \\
3-2 & NEWTMN & NEWTMN is a positive integer that is less or equal to
NEWTMX that defines the minimum number of Newton iterations for a time
step. & 1 \\
dimensionless & dimensionless & dimensionless & \\
3-3 & LITMAX & LITMAX is a positive integer greater or equal to LITMIN
that sets the maximum number of linear iterations within a Newton
iteration. & 25 \\
dimensionless & dimensionless & dimensionless & \\
3-4 & LITMIN & LITMIN is a positive integer less or equal to LITMAX that
sets the minimum number of linear iterations within a Newton iteration.
& 1 \\
dimensionless & dimensionless & dimensionless & \\
3-5 & MXWSIT & MXWSIT is a positive integer that defines the maximum
number of iterations within a well flow calculation. & 8 \\
dimensionless & dimensionless & dimensionless & \\
3-6 & MXWPIT & MXWPIT is a positive integer that stipulates the maximum
number of iterations for solving the bottom-hole pressure for wells
under tubing head pressure control within a well flow calculation. &
8 \\
dimensionless & dimensionless & dimensionless & \\
3-7 & DDPLIM & DDPLIM a real positive value that stipulates the maximum
pressure change at the last Newton iteration. & 1.0 x 10-6 \\
psia & barsa & atma & \\
3-8 & DDSLIM & DDSLIM a real positive value that sets the maximum
saturation change at the last Newton iteration. & 1.0 x 10-6 \\
dimensionless & dimensionless & dimensionless & \\
3-9 & TRGDPR & TRGDP is a real positive value that defines the target
pressure change within a time step. & 1.0 x 10-6 \\
psia & barsa & atma & \\
3-10 & XXXDPR & XXXDPR is a real positive value that stipulates the
maximum tolerable pressure change within a time step. & 1.0 x 10-6 \\
psia & barsa & atma & \\
3-11 & MNWRFP & MNWRFP is a positive integer greater than one and less
than NEWTMX that defines the minimum number of Newton iterations before
invoking the bisection algorithm for when the polymer phase is active in
the model via the POLYMER keyword in the RUNSPEC section. & 4 \\
dimensionless & dimensionless & dimensionless & \\
3-12 & / & Record terminated by a ``/'' & Not Applicable \\
Notes: & & & \\
\end{longtable}

\emph{Table 12.69: TUNINGL Keyword Description}

Note that for record number two (items 2-1 to 2-13) the maximum values
should always be greater than the associated target value; for example,
XXXCNV should be greater than TRGCNV. Also note that the TUNINGS keyword
is stored on the restart files (see RPTRST -- Define Data to be Written
to the RESTART File) enabling the parameters to be utilized in a restart
run without re-specifying the keyword.

See also the TUNINGS keyword in the SCHEDULE section that sets the
tuning parameters for individual LGRs.

\subsubsection{Example}\label{example-556}

\begin{itemize}
\tightlist
\item
\end{itemize}

\begin{verbatim}
--       DEFAULT TUNINGL PARAMETERS
--
TUNINGL
       /
/
/
\end{verbatim}

The above example explicitly sets the default parameters for all LGRs.

\subsection{TUNINGS -- Numerical Tuning Control for Individual
LGRs}\label{tunings-numerical-tuning-control-for-individual-lgrs}

\begin{longtable}[]{@{}
  >{\raggedright\arraybackslash}p{(\linewidth - 14\tabcolsep) * \real{0.1250}}
  >{\raggedright\arraybackslash}p{(\linewidth - 14\tabcolsep) * \real{0.1250}}
  >{\raggedright\arraybackslash}p{(\linewidth - 14\tabcolsep) * \real{0.1250}}
  >{\raggedright\arraybackslash}p{(\linewidth - 14\tabcolsep) * \real{0.1250}}
  >{\raggedright\arraybackslash}p{(\linewidth - 14\tabcolsep) * \real{0.1250}}
  >{\raggedright\arraybackslash}p{(\linewidth - 14\tabcolsep) * \real{0.1250}}
  >{\raggedright\arraybackslash}p{(\linewidth - 14\tabcolsep) * \real{0.1250}}
  >{\raggedright\arraybackslash}p{(\linewidth - 14\tabcolsep) * \real{0.1250}}@{}}
\toprule\noalign{}
\begin{minipage}[b]{\linewidth}\raggedright
RUNSPEC
\end{minipage} & \begin{minipage}[b]{\linewidth}\raggedright
GRID
\end{minipage} & \begin{minipage}[b]{\linewidth}\raggedright
EDIT
\end{minipage} & \begin{minipage}[b]{\linewidth}\raggedright
PROPS
\end{minipage} & \begin{minipage}[b]{\linewidth}\raggedright
REGIONS
\end{minipage} & \begin{minipage}[b]{\linewidth}\raggedright
SOLUTION
\end{minipage} & \begin{minipage}[b]{\linewidth}\raggedright
SUMMARY
\end{minipage} & \begin{minipage}[b]{\linewidth}\raggedright
SCHEDULE
\end{minipage} \\
\midrule\noalign{}
\endhead
\bottomrule\noalign{}
\endlastfoot
\end{longtable}

\subsubsection{Description}\label{description-1026}

TUNINGS defines the parameters used for controlling the commercial
simulator's numerical convergence parameters for individual Local Grid
Refinements (``LGR''). The keyword is similar to the TUNINGL keyword in
the SCHEDULE section that applies the tuning parameters to the all LGRs,
except for an additional first record that includes the LGR name.

This keyword is not supported by OPM Flow but has no effect on the
results so it will be ignored.

See section 2.2 Running OPM Flow 2023-04 From The Command Line on how to
invoke various numerical schemes via the OPM Flow command line
interface.

\begin{longtable}[]{@{}
  >{\raggedright\arraybackslash}p{(\linewidth - 6\tabcolsep) * \real{0.2500}}
  >{\raggedright\arraybackslash}p{(\linewidth - 6\tabcolsep) * \real{0.2500}}
  >{\raggedright\arraybackslash}p{(\linewidth - 6\tabcolsep) * \real{0.2500}}
  >{\raggedright\arraybackslash}p{(\linewidth - 6\tabcolsep) * \real{0.2500}}@{}}
\toprule\noalign{}
\begin{minipage}[b]{\linewidth}\raggedright
No.
\end{minipage} & \begin{minipage}[b]{\linewidth}\raggedright
Name
\end{minipage} & \begin{minipage}[b]{\linewidth}\raggedright
Description
\end{minipage} & \begin{minipage}[b]{\linewidth}\raggedright
Default
\end{minipage} \\
\midrule\noalign{}
\endhead
\bottomrule\noalign{}
\endlastfoot
Field & Metric & Laboratory & \\
0-1 & LGRNAME & A character string of up to eight characters in length
that defines the LGR name for which the tuning data is being being
defined. & None \\
0-2 & / & Record terminated by a ``/'' & Not Applicable \\
1-1 & TSINIT & TSINT is a real positive value that defines the maximum
length of the next time step. Note that whenever the keyword is used
TSINIT is always set back to the default value of one, unless explicitly
over written. & 1.0 \\
days & days & hours & \\
1-2 & TSMAXZ & TSMAXZ is a real positive value that defines the maximum
length of the next time step following TSINIT. & 365.0 \\
days & days & hours & \\
1-3 & TSMINZ & TSMINZ is a real positive value that defines the minimum
length of all time steps. & 0.1 \\
days & days & hours & \\
1-4 & TSMCHP & TSMCHP is a real positive values that sets the minimum
length of all chopped time steps. & 0.15 \\
days & days & hours & \\
1-5 & TSFMAX & TSFMAX is a real positive value that specifies the
maximum growth rate a time step can be increased by, subject to the
maximum allowable time step size set by TSMAXZ. For example, if the
current time step has converged at 10 days and TSFMAX is set to the
default value, then the next time step will be 3.0 x 10 days, that is 30
days provided it is less than TSMAXZ. & 3.0 \\
dimensionless & dimensionless & dimensionless & \\
1-6 & TSFMIN & TSFMIN is a real positive value that specifies the
minimum decay rate a time step can be decreased by, subject to the
minimum allowable time step size set by TSMINZ. For example, if the
current time step has not converged at 10 days and TSFMAX is set to the
default value, then the next time step will be 0.3 x 10 days, that is
the maximum of 0.3 days and TSMINZ. & 0.3 \\
dimensionless & dimensionless & dimensionless & \\
1-7 & TSFCNV & TSFCNV real positive value that specifies the decay rate
a time step can be decreased by after the number of target iterations
has been exceeded. & 0.1 \\
dimensionless & dimensionless & dimensionless & \\
1-8 & TFDIFF & TFDIFFA is a real positive value that sets the time step
growth factor of the time step after a convergence failure. For example,
if the chopped current convergent time step is 10 days and TFDIFF is set
to the default value, then the time step will be increased to 1.25 x 10
days, that is the minimum of 11.25 days and TSMAXZ. & 1.25 \\
dimensionless & dimensionless & dimensionless & \\
1-9 & THRURPT & THRURPT is a real positive value that specifies the
maximum throughput ratio over a time step. & 1.0 x 1020 \\
dimensionless & dimensionless & dimensionless & \\
1-10 & TMAXWC & TMAXWC is a real double precision value that defines
maximum allowed time step after a well event; for example, when a well
is opened or closed, etc. & None \\
days & days & hours & \\
1-11 & / & Record terminated by a ``/'' & Not Applicable \\
2-1 & TRGTTE & TRGTTE is a real positive value that sets the time
truncation error target. & 0.1 \\
dimensionless & dimensionless & dimensionless & \\
2-2 & TRGCNV & TRGCNV a real positive value that defines the non-linear
convergence error. & 0.001 \\
dimensionless & dimensionless & dimensionless & \\
2-3 & TRGMBE & TRGMBE is a real positive value that specifies then
target material balance error. & 1.0 x 10-7 \\
dimensionless & dimensionless & dimensionless & \\
2-4 & TRGLCV & TRGLCV is a real positive value that specifies the linear
convergence error target. & 0.00001 \\
dimensionless & dimensionless & dimensionless & \\
2-5 & XXXTTE & XXXTTE is a real positive value that sets the maximum
time truncation error. & 10.0 \\
dimensionless & dimensionless & dimensionless & \\
2-6 & XXXCNV & XXXCNV is a real positive value that defines the maximum
non-linear convergence error. & 0.01 \\
dimensionless & dimensionless & dimensionless & \\
2-7 & XXXMBE & XXXMBE is a real positive value that specifies the
maximum mass balance error, that is the tolerated mass balance error
relative to total mass present. & 1.0 x 10-6 \\
dimensionless & dimensionless & dimensionless & \\
2-8 & XXXLCV & XXXLCV is a real positive values that sets the maximum
linear convergence error. & 0.0001 \\
dimensionless & dimensionless & dimensionless & \\
2-9 & XXXWFL & XXXWFL is a real positive values that fixes the maximum
well flow convergence error. & 0.001 \\
dimensionless & dimensionless & dimensionless & \\
2-10 & TRGFIP & TRGFIP is a real positive value that stipulates the
target fluid in-place error in Local Grid Refinements. & 0.025 \\
dimensionless & dimensionless & dimensionless & \\
2-11 & TRGSFT & TRGSFT is a real positive values that defines the target
surfactant change when the Surfactant Model is active in the run. &
None \\
dimensionless & dimensionless & dimensionless & \\
2-12 & THIONX & THIONX is a positive real value used to set the
threshold for damping in the ion echange calculation for when the Brine
Model is active in the run. & 0.01 \\
dimensionless & dimensionless & dimensionless & \\
2-13 & TRWGHT & TRWGHT is a positive integer that stipulates the
implicitness for active tracer updates within the Newton iterations, and
should be set to: & 1 \\
dimensionless & dimensionless & dimensionless & \\
2-14 & / & Record terminated by a ``/'' & Not Applicable \\
3-1 & NEWTMX & NEWTMX is a positive integer greater or equal to NEWTMN
that stipulates the maximum number of Newton iterations for a time step.
& 12 \\
dimensionless & dimensionless & dimensionless & \\
3-2 & NEWTMN & NEWTMN is a positive integer that is less or equal to
NEWTMX that defines the minimum number of Newton iterations for a time
step. & 1 \\
dimensionless & dimensionless & dimensionless & \\
3-3 & LITMAX & LITMAX is a positive integer greater or equal to LITMIN
that sets the maximum number of linear iterations within a Newton
iteration. & 25 \\
dimensionless & dimensionless & dimensionless & \\
3-4 & LITMIN & LITMIN is a positive integer less or equal to LITMAX that
sets the minimum number of linear iterations within a Newton iteration.
& 1 \\
dimensionless & dimensionless & dimensionless & \\
3-5 & MXWSIT & MXWSIT is a positive integer that defines the maximum
number of iterations within a well flow calculation. & 8 \\
dimensionless & dimensionless & dimensionless & \\
3-6 & MXWPIT & MXWPIT is a positive integer that stipulates the maximum
number of iterations for solving the bottom-hole pressure for wells
under tubing head pressure control within a well flow calculation. &
8 \\
dimensionless & dimensionless & dimensionless & \\
3-7 & DDPLIM & DDPLIM a real positive value that stipulates the maximum
pressure change at the last Newton iteration. & 1.0 x 10-6 \\
psia & barsa & atma & \\
3-8 & DDSLIM & DDSLIM a real positive value that sets the maximum
saturation change at the last Newton iteration. & 1.0 x 10-6 \\
dimensionless & dimensionless & dimensionless & \\
3-9 & TRGDPR & TRGDP is a real positive value that defines the target
pressure change within a time step. & 1.0 x 10-6 \\
psia & barsa & atma & \\
3-10 & XXXDPR & XXXDPR is a real positive value that stipulates the
maximum tolerable pressure change within a time step. & 1.0 x 10-6 \\
psia & barsa & atma & \\
3-11 & MNWRFP & MNWRFP is a positive integer greater than one and less
than NEWTMX that defines the minimum number of Newton iterations before
invoking the bisection algorithm for when the polymer phase is active in
the model via the POLYMER keyword in the RUNSPEC section. & 4 \\
dimensionless & dimensionless & dimensionless & \\
3-12 & / & Record terminated by a ``/'' & Not Applicable \\
Notes: & & & \\
\end{longtable}

\emph{Table 12.70: TUNINGS Keyword Description}

Note that for record number two (items 2-1 to 2-13) the maximum values
should always be greater than the associated target value; for example,
XXXCNV should be greater than TRGCNV. Also note that the TUNING keyword
is stored on the restart files (see RPTRST -- Define Data to be Written
to the RESTART File) enabling the parameters to be utilized in a restart
run without re-specifying the keyword.

See also the TUNINGL keyword in the SCHEDULE section that sets the
tuning parameters for all LGRs.

\subsubsection{Example}\label{example-557}

\begin{verbatim}
--
--       DEFAULT TUNINGS PARAMETERS
--
TUNINGS
OP01-LGR
/
/
/
\end{verbatim}

The above example explicitly sets the default parameters for the LGR
named OP01-LGR

\begin{verbatim}
\end{verbatim}

\subsection{UDQ -- Declare User Define Quantities
(``UDQ'')}\label{udq-declare-user-define-quantities-udq}

\begin{longtable}[]{@{}
  >{\raggedright\arraybackslash}p{(\linewidth - 14\tabcolsep) * \real{0.1250}}
  >{\raggedright\arraybackslash}p{(\linewidth - 14\tabcolsep) * \real{0.1250}}
  >{\raggedright\arraybackslash}p{(\linewidth - 14\tabcolsep) * \real{0.1250}}
  >{\raggedright\arraybackslash}p{(\linewidth - 14\tabcolsep) * \real{0.1250}}
  >{\raggedright\arraybackslash}p{(\linewidth - 14\tabcolsep) * \real{0.1250}}
  >{\raggedright\arraybackslash}p{(\linewidth - 14\tabcolsep) * \real{0.1250}}
  >{\raggedright\arraybackslash}p{(\linewidth - 14\tabcolsep) * \real{0.1250}}
  >{\raggedright\arraybackslash}p{(\linewidth - 14\tabcolsep) * \real{0.1250}}@{}}
\toprule\noalign{}
\begin{minipage}[b]{\linewidth}\raggedright
RUNSPEC
\end{minipage} & \begin{minipage}[b]{\linewidth}\raggedright
GRID
\end{minipage} & \begin{minipage}[b]{\linewidth}\raggedright
EDIT
\end{minipage} & \begin{minipage}[b]{\linewidth}\raggedright
PROPS
\end{minipage} & \begin{minipage}[b]{\linewidth}\raggedright
REGIONS
\end{minipage} & \begin{minipage}[b]{\linewidth}\raggedright
SOLUTION
\end{minipage} & \begin{minipage}[b]{\linewidth}\raggedright
SUMMARY
\end{minipage} & \begin{minipage}[b]{\linewidth}\raggedright
SCHEDULE
\end{minipage} \\
\midrule\noalign{}
\endhead
\bottomrule\noalign{}
\endlastfoot
\end{longtable}

\subsubsection{Description}\label{description-1027}

This keyword starts the definition of a UDQ section that stipulates the
variables and operations used to access the User Defined Quantities
features in OPM Flow. UDQ variables can be constants, SUMMARY variables,
as defined in the SUMMARY section, or a formula using various
mathematical functions together with constants and SUMMARY variables.
Available operations include the ASSIGN, DEFINE, UNITS and UPDATE
commands that are sub-keywords to the UDQ section keyword.

User Defined Quantities can be output to the summmary file and used as
User Defined Arguments (``UDA'') in the SCHEDULE section with various
group, well, and connection keywords.

Although this keyword is read by OPM Flow and the ACTION and UDQ
computational logic and calculations have been implemented, one should
use caution when using this facility as it may result in OPM Flow
aborting.

\begin{longtable}[]{@{}
  >{\raggedright\arraybackslash}p{(\linewidth - 6\tabcolsep) * \real{0.2500}}
  >{\raggedright\arraybackslash}p{(\linewidth - 6\tabcolsep) * \real{0.2500}}
  >{\raggedright\arraybackslash}p{(\linewidth - 6\tabcolsep) * \real{0.2500}}
  >{\raggedright\arraybackslash}p{(\linewidth - 6\tabcolsep) * \real{0.2500}}@{}}
\toprule\noalign{}
\begin{minipage}[b]{\linewidth}\raggedright
No.
\end{minipage} & \begin{minipage}[b]{\linewidth}\raggedright
Name
\end{minipage} & \begin{minipage}[b]{\linewidth}\raggedright
Description
\end{minipage} & \begin{minipage}[b]{\linewidth}\raggedright
Default
\end{minipage} \\
\midrule\noalign{}
\endhead
\bottomrule\noalign{}
\endlastfoot
UDQ & Define the start of UDQ Definition Section. This is then followed
on a new line by any number of UDQ records that define the various
operations to be performed using the ASSIGN, DEFINE UNITS and UPDATE
sub-keywords for the OPERATOR. & & \\
1 & OPERATOR & OPERATOR is a defined character string that specifies the
type of operation to perform, and should be one of the following: & \\
2 & VARIABLE & VARIABLE is a character string of length eight that
stipulates the name of the user defined variable that will processed by
the OPERATOR command. The first two characters of VARIABLE must be set
based on the type of variable being defined, that is: OPM Flow currently
only supports field, group, segment and well variables (FU\emph{, GU},
SU* and WU*). & \\
3 & EXPRESSION & The data type for EXPRESSION is based on the OPERATOR
option above, namely if OPERATOR is set to: & \\
& / & Termination of a UDQ record. Note that multiple numbers of records
can be entered within a UDQ section with each record terminated by a
``/''. & \\
/ & Define the end of UDQ Definition Section & & \\
Notes: & & & \\
\end{longtable}

\emph{Table 12.71: UDQ Keyword Description}

All the functions available for use in the DEFINE EXPRESSION in the
commercial simulator are listed in Table 12.71 below.

\begin{longtable}[]{@{}
  >{\raggedright\arraybackslash}p{(\linewidth - 6\tabcolsep) * \real{0.2500}}
  >{\raggedright\arraybackslash}p{(\linewidth - 6\tabcolsep) * \real{0.2500}}
  >{\raggedright\arraybackslash}p{(\linewidth - 6\tabcolsep) * \real{0.2500}}
  >{\raggedright\arraybackslash}p{(\linewidth - 6\tabcolsep) * \real{0.2500}}@{}}
\toprule\noalign{}
\begin{minipage}[b]{\linewidth}\raggedright
UDQ - Description of Functions
\end{minipage} & \begin{minipage}[b]{\linewidth}\raggedright
\end{minipage} & \begin{minipage}[b]{\linewidth}\raggedright
\end{minipage} & \begin{minipage}[b]{\linewidth}\raggedright
\end{minipage} \\
\midrule\noalign{}
\endhead
\bottomrule\noalign{}
\endlastfoot
Function & Prece-dence & Type & Description \\
( & & & Open bracket. \\
) & & & Close bracket. \\
{[} & & & Start of user defined table argument list. \\
{]} & & & End of user defined table argument list. \\
, & & & Separator for multi-dimensional user defined table argument
list. \\
- & 6 & unary elemental & Negation of defined elements. \\
ABS() & 6 & unary elemental & Absolute value of defined elements. \\
AVEA() & 6 & unary scalar & Arithmetic average of defined elements. \\
AVEG() & 6 & unary scalar & Geometric average of defined elements. \\
AVEH() & 6 & unary scalar & Harmonic average of defined elements. \\
DEF() & 6 & unary elemental & Returns 1 if the element is defined,
otherwise returns undefined. \\
EXP() & 6 & unary elemental & Exponential of defined elements. \\
IDV() & 6 & unary elemental & Returns 1 if the element is defined,
otherwise returns 0. \\
LN() & 6 & unary elemental & Natural logarithm of defined elements. \\
LOG() & 6 & unary elemental & Logarithm base 10 of defined elements. \\
MAX() & 6 & unary scalar & Maximum of defined elements. \\
MIN() & 6 & unary scalar & Minimum of defined elements. \\
NORM1() & 6 & unary scalar & 1 norm of defined elements. \\
NORM2() & 6 & unary scalar & 2 norm of defined elements. \\
NORMI() & 6 & unary scalar & Infinity norm of defined elements. \\
NINT() & 6 & unary elemental & Nearest integer to defined elements. \\
PROD() & 6 & unary scalar & Product of defined elements. \\
RANDN() & 6 & unary elemental & Random sample from Normal distribution
N(0,1), where the seed is specified by the UDQPARAM keyword Item 1. \\
RANDU() & 6 & unary elemental & Random sample from Uniform distribution
U(-1,1), where the seed is specified by the UDQPARAM keyword Item 1. \\
RRNDN() & 6 & unary elemental & Random sample from Normal distribution
N(0,1), where the seed is determined from the current time. By default
restart simulations will use the seed from the base simulation (see
UDQDIMS keyword Item 11). \\
RRNDU() & 6 & unary elemental & Random sample from Uniform distribution
U(-1,1), where the seed is determined from the current time. By default
restart simulations will use the seed from the base simulation (see
UDQDIMS keyword Item 11). \\
SORTA() & 6 & unary elemental & Position of the element in an ascending
sort of defined elements. \\
SORTD() & 6 & unary elemental & Position of the element in an descending
sort of defined elements. \\
SUM() & 6 & unary scalar & Sum of defined elements. \\
UNDEF() & 6 & unary elemental & Returns 1 if element is undefined,
otherwise returns undefined. \\
TU*{[}{]} & 6 & unary elemental & Lookup value of the user defined table
based on the arguments in the square brackets. All arguments must have
the same type or be scalar. A scalar is returned if all the arguments
are scalar, or the same type of UDQ set is returned as the UDQ set
arguments. \\
\textless= & 5 & binary intersection & Returns 1 if the LHS is less than
or equal to the RHS, otherwise return 0. The tolerance for equality is
specified by the UDQPARAM keyword Item 4. \\
\textgreater= & 5 & binary intersection & Returns 1 if the LHS is
greater than or equal to the RHS, otherwise return 0. The tolerance for
equality is specified by the UDQPARAM keyword Item 4. \\
\textless{} & 5 & binary intersection & Returns 1 if the LHS is less
than the RHS, otherwise return 0. \\
\textgreater{} & 5 & binary intersection & Returns 1 if the LHS is
greater than the RHS, otherwise return 0. \\
== & 5 & binary intersection & Returns 1 if the LHS is equal to the RHS.
The tolerance for equality is specified by the UDQPARAM keyword Item
4. \\
!= & 5 & binary intersection & Returns 1 if the LHS is not equal to the
RHS. The tolerance for equality is specified by the UDQPARAM keyword
Item 4. \\
\^{} & 4 & binary intersection & Exponentiation. \\
* & 3 & binary intersection & Multiplication. \\
/ & 3 & binary intersection & Division. \\
+ & 2 & binary intersection & Addition. \\
- & 2 & binary intersection & Subtraction. \\
UADD & 1 & binary union & Returns sum of intersecting elements from the
union of two UDQ sets. \\
UMAX & 1 & binary union & Returns maximum of intersecting elements from
the union of two UDQ sets. \\
UMIN & 1 & binary union & Returns minimum of intersecting elements from
the union of two UDQ sets. \\
UMUL & 1 & binary union & Returns product of intersecting elements from
the union of two UDQ sets. \\
Notes: & & & \\
\end{longtable}

\emph{Table 12.71: UDQ - Description of Functions}

See also the UDADIMS, UDQDIMS and UDQPARAM keywords in the RUNSPEC
section to define the dimensions for the UDQ keyword and associated
variables.

\begin{tcolorbox}[enhanced jigsaw, title=\textcolor{quarto-callout-note-color}{\faInfo}\hspace{0.5em}{Note}, coltitle=black, breakable, toprule=.15mm, colbacktitle=quarto-callout-note-color!10!white, leftrule=.75mm, opacitybacktitle=0.6, colback=white, bottomtitle=1mm, toptitle=1mm, titlerule=0mm, left=2mm, colframe=quarto-callout-note-color-frame, arc=.35mm, rightrule=.15mm, bottomrule=.15mm, opacityback=0]

Wells and/or groups needed in a UDQ DEFINE statement must be present at
the point of definition. In particular, wells must have been introduced
using the WELSPECS keyword prior to the UDQ keyword if any specific well
name is used in the defining expression. Similarly, well lists must have
been introduced and populated through the WLIST keyword before using the
well list name in the UDQ definition. Finally, groups must be introduced
through the WELSPECS or GRUPTREE keywords before using any specific
group names in the defining expression.

\end{tcolorbox}

User Defined Quantities can also be used as User Defined Arguments
(``UDA'') in the SCHEDULE section with various group, well, and
connection keywords. In this case, the UDA variables are used to replace
numerical values on these keywords by UDA variables that have been
defined by the UDQ keyword. For example, if we wish to make the oil rate
for certain wells be a function of their water cut, then one can define
the function using the UDQ keyword that results in a UDQ variable,
WU\_WCUT say, and then use WU\_WCUT as a UDA variable on the WCONPROD
keyword for the ORAT parameter. Table 12.72 lists the keywords that can
be used with UDA variables.

\begin{longtable}[]{@{}
  >{\raggedright\arraybackslash}p{(\linewidth - 10\tabcolsep) * \real{0.1667}}
  >{\raggedright\arraybackslash}p{(\linewidth - 10\tabcolsep) * \real{0.1667}}
  >{\raggedright\arraybackslash}p{(\linewidth - 10\tabcolsep) * \real{0.1667}}
  >{\raggedright\arraybackslash}p{(\linewidth - 10\tabcolsep) * \real{0.1667}}
  >{\raggedright\arraybackslash}p{(\linewidth - 10\tabcolsep) * \real{0.1667}}
  >{\raggedright\arraybackslash}p{(\linewidth - 10\tabcolsep) * \real{0.1667}}@{}}
\toprule\noalign{}
\begin{minipage}[b]{\linewidth}\raggedright
UDQ - User Defined Argument Supported Keywords Schedule Section Keywords
Status
\end{minipage} & \begin{minipage}[b]{\linewidth}\raggedright
\end{minipage} & \begin{minipage}[b]{\linewidth}\raggedright
\end{minipage} & \begin{minipage}[b]{\linewidth}\raggedright
\end{minipage} & \begin{minipage}[b]{\linewidth}\raggedright
\end{minipage} & \begin{minipage}[b]{\linewidth}\raggedright
\end{minipage} \\
\midrule\noalign{}
\endhead
\bottomrule\noalign{}
\endlastfoot
Number & Group Keywords & Well Keywords & Connection Keywords &
Miscellaneous Keywords & \\
1 & GCONINJE & WALKALIN & WINJEDET & CECON & LINCOM \\
2 & GCONPRI & WAPI & \hyperref[REF_HEADING_KEYWORD_WINJFCNC]{WINJFCNC} &
CPIFACT & \\
3 & GCONPROD & WCONHIST & WINJTEMP & CPIFACTL & \\
4 & GCONSALE & WCONINJE & WPOLYMER & & \\
5 & GCONSUMP & WCONPROD & WSALT & & \\
6 & GECON & WECON & WSOLVENT & & \\
7 & GRUPFUEL & WECONCMF & WSURFACT & & \\
8 & GRUPSALE & WELDRAW & WTADD & & \\
9 & GSATPROD & WELLSTRE & WTMULT & & \\
10 & GTADD & WELTARG & WTRACER & & \\
11 & GTMULT & WFOAM & & & \\
& & Multi-Segment Well Keywords & & & \\
1 & & WSEGTABL & WSEGVALV & & \\
Notes: & & & & & \\
\end{longtable}

\emph{Table 12.72: UDQ - User Defined Argument Supported Keywords}

\begin{tcolorbox}[enhanced jigsaw, title=\textcolor{quarto-callout-note-color}{\faInfo}\hspace{0.5em}{Note}, coltitle=black, breakable, toprule=.15mm, colbacktitle=quarto-callout-note-color!10!white, leftrule=.75mm, opacitybacktitle=0.6, colback=white, bottomtitle=1mm, toptitle=1mm, titlerule=0mm, left=2mm, colframe=quarto-callout-note-color-frame, arc=.35mm, rightrule=.15mm, bottomrule=.15mm, opacityback=0]

Note that after the terminating ``/'' for the ASSIGN operator normally
any comments can be entered; however, if there is ``/'' within the
comment field, as per: ASSIGN FUNGLYLD 1.196 / Condensate Yield (63.5 /
56.7) / (1.0 -- 0.065) then the simulator will abort. The work around is
to manually place the comment characters ``--'' after the ASSIGN
terminating ``/'', like so: ASSIGN FUNGLYLD 1.196 / -- Condensate Yield
(63.5 / 56.7) / (1.0 -- 0.06)

\end{tcolorbox}

\subsubsection{Examples}\label{examples-86}

The first example shows how to define some constant field variables used
for calculating facilities corrected condensate and Liquefied Petroleum
Gas {[}Liquefied Petroleum Gas or LPG consists mainly of propane,
propylene, butane, and butylene in various mixtures. It is produced as a
by-product of natural gas processing and petroleum refining. The
components of LPG are gases at standard conditions.{]} (``LPG ``) yields
in a wet gas model:

\begin{verbatim}
--
-- DEFINE START OF USER DEFINED QUANTITY SECTION
--
UDQ
--
-- OPERATOR VARIABLE  EXPRESSION
--
ASSIGN      FUNGLYLD  1.100000         / -- Condensate  Yield (stb/Mscf)
ASSIGN      FUNGLSHK  0.000000         /    Condensate  Shrinkage Factor to Zero
ASSIGN      FULPGYLD  0.065775         / -- LPG Sep Gas Yield (stb/Mscf)
ASSIGN      FULPGSHK  0.080410         /    LPG         Shrinkage Factor
ASSIGN      FUFACSHK  0.000935         /    Facilities  Shrinkage Factor
ASSIGN      FUFULSHK  0.052924         /    Fuel        Utilization
ASSIGN      FUDELTA   1E-10            /    Value to avoid dividing by zero errors

/  DEFINE END OF USER DEFINED QUANTITY SECTION
\end{verbatim}

The next example is a continuation of this example by showing how one
can calculate the adjusted field condensate and LPG rates. Note both
examples could be merged into a single UDQ definition but have been
stated separately for ease of reference.

\begin{verbatim}
--
-- DEFINE START OF USER DEFINED QUANTITY SECTION
--
UDQ
--
-- OPERATOR VARIABLE  EXPRESSION
--
DEFINE      FU_FNGLR  FGPR * (FOGR * FUNGLYLD)  / Calculate Condensate Rate Field
UPDATE      FU_FNGLR  ON                       /
UNITS       FU_FNGLR  STBD                     /

DEFINE      FU_FLPGR  FU_FWGPR * FULPGYLD      / Calculate LPG Rate Field
UPDATE      FU_FLPGR  ON                       /
UNITS       FU_FLPGR  STBD                     /

/  DEFINE END OF USER DEFINED QUANTITY SECTION
\end{verbatim}

In the above the DEFINE operator is use to define the equations to
calculate the corrected condensate (FU\_FNGLR) and LPG rates (FU\_FLPGR)
with the UPDATE operator set to ON so that the rates are calculate at
every time step, and finally, the UNITS operator is used to set the
units of the calculated rates.

The final example show the use of the UDADIMS and UDQDIMS keywords in
the RUNSPEC section, followed by the keywords in the SCHEDULE section
that define a UDQ definition that uses the DEFINE operator to calculate
adjusted well rates based on an expression. The final set of keywords
show how the UDQ defined variables are employed on the WCONPROD keyword
to control the production constraints for several wells.

\begin{verbatim}
RUNSPEC SECTION KEYWORDS
------------------------
--
--       USER DEFINED ARGUMENT DIMENSIONS
--       NO.     NOT     TOTAL
--       ARGS    USED    UDQ
UDADIMS
         10       1*     10                                                    /
--
--       USER DEFINED ARGUMENT DIMENSIONS FACILITY
--       MAX     MAX     MAX    MAX    MAX    MAX   MAX   MAX   MAX  MAX    RAND
--       FUNCS   ITEMS   CONNS  FIELD  GROUP  REGS  SEGTM WELL  AQUF BLCKS  OPT
UDQDIMS
         50      25      0      50     50     0     0     0     0    0      N  /
\end{verbatim}

And the SCHEDULE section part of the example is shown below.

\begin{verbatim}

SCHEDULE SECTION KEYWORDS
--------------------------
--
-- DEFINE START OF USER DEFINED QUANTITY SECTION
--
UDQ
--
-- OPERATOR VARIABLE EXPRESSION
--
DEFINE      WUOPRL (WOPR OPL01 - 150) * 0.90 / OIL & LIQ CAPACITIES
DEFINE      WULPRL (WLPR OPL01 - 200) * 0.90 /  at GEFAC = 0.8995
DEFINE      WUOPRU (WOPR OPU01 - 250) * 0.80 /
DEFINE      WULPRU (WLPR OPU01 - 300) * 0.80 /
--
UNITS       WUOPRL SM3/DAY                   / DEFINE REPORTING UNITS
UNITS       WULPRL SM3/DAY                   / FOR UDQ VARIABLES
UNITS       WUOPRU SM3/DAY                   /
UNITS       WULPRU SM3/DAY                   /
/  DEFINE END OF USER DEFINED QUANTITY SECTION
--
--       WELL PRODUCTION WELL CONTROLS
--
-- WELL  OPEN/  CNTL   OIL    WAT    GAS   LIQ    RES    BHP   THP   VFP    VFP
-- NAME  SHUT   MODE   RATE   RATE   RATE  RATE   RATE   PRES  PRES  TABLE  ALFQ
WCONPROD
OP01     SHUT   GRUP   1*     1*     1*    1*     1*     200.0                 /
OP02     SHUT   GRUP   1*     1*     1*    1*     1*     200.0                 /
/
DATES
         1  FEB   2020  /
--
--       WELL PRODUCTION WELL CONTROLS
--
-- WELL  OPEN/  CNTL   OIL    WAT    GAS   LIQ    RES    BHP   THP   VFP    VFP
-- NAME  SHUT   MODE   RATE   RATE   RATE  RATE   RATE   PRES  PRES  TABLE  ALFQ
WCONPROD
OP01     OPEN   GRUP   WUOPRL 1*     1*    WULPRL 1*     60.0                  /
OP02     OPEN   GRUP   WUOPRL 1*     1*    WULPRL 1*     00.0                  /
/
DATES
         1  MAR   2020  /
         1  APR   2020  /
         1  MAY   2020  /
         1  JUN   2020  /
         1  JLY   2020  /
         1  AUG   2020  /
         1  SEP   2020  /
/
\end{verbatim}

\subsection{UDT -- Declare User Define Tables
(``UDT'')}\label{udt-declare-user-define-tables-udt}

\begin{longtable}[]{@{}
  >{\raggedright\arraybackslash}p{(\linewidth - 14\tabcolsep) * \real{0.1250}}
  >{\raggedright\arraybackslash}p{(\linewidth - 14\tabcolsep) * \real{0.1250}}
  >{\raggedright\arraybackslash}p{(\linewidth - 14\tabcolsep) * \real{0.1250}}
  >{\raggedright\arraybackslash}p{(\linewidth - 14\tabcolsep) * \real{0.1250}}
  >{\raggedright\arraybackslash}p{(\linewidth - 14\tabcolsep) * \real{0.1250}}
  >{\raggedright\arraybackslash}p{(\linewidth - 14\tabcolsep) * \real{0.1250}}
  >{\raggedright\arraybackslash}p{(\linewidth - 14\tabcolsep) * \real{0.1250}}
  >{\raggedright\arraybackslash}p{(\linewidth - 14\tabcolsep) * \real{0.1250}}@{}}
\toprule\noalign{}
\begin{minipage}[b]{\linewidth}\raggedright
RUNSPEC
\end{minipage} & \begin{minipage}[b]{\linewidth}\raggedright
GRID
\end{minipage} & \begin{minipage}[b]{\linewidth}\raggedright
EDIT
\end{minipage} & \begin{minipage}[b]{\linewidth}\raggedright
PROPS
\end{minipage} & \begin{minipage}[b]{\linewidth}\raggedright
REGIONS
\end{minipage} & \begin{minipage}[b]{\linewidth}\raggedright
SOLUTION
\end{minipage} & \begin{minipage}[b]{\linewidth}\raggedright
SUMMARY
\end{minipage} & \begin{minipage}[b]{\linewidth}\raggedright
SCHEDULE
\end{minipage} \\
\midrule\noalign{}
\endhead
\bottomrule\noalign{}
\endlastfoot
\end{longtable}

\subsubsection{Description}\label{description-1028}

The UDT keyword defines a single multi-dimensional User Defined Table
(``UDT'').

User Defined Quantities (``UDQ'') defined using the UDQ keyword in the
SCHEDULE section can be assigned values by looking up a UDT.

This keyword is partially supported by OPM Flow; only one-dimensional
lookup tables are supported.

Each User Defined Table must be defined with a separate UDT keyword that
consists of a number of records defining the table name and dimensions;
the interpolation type and interpolation points for each dimension; and
the table values for each set of interpolation points. The total number
of records depends on the table dimensions and the number of
interpolation points.

\begin{longtable}[]{@{}
  >{\raggedright\arraybackslash}p{(\linewidth - 6\tabcolsep) * \real{0.2500}}
  >{\raggedright\arraybackslash}p{(\linewidth - 6\tabcolsep) * \real{0.2500}}
  >{\raggedright\arraybackslash}p{(\linewidth - 6\tabcolsep) * \real{0.2500}}
  >{\raggedright\arraybackslash}p{(\linewidth - 6\tabcolsep) * \real{0.2500}}@{}}
\toprule\noalign{}
\begin{minipage}[b]{\linewidth}\raggedright
No.
\end{minipage} & \begin{minipage}[b]{\linewidth}\raggedright
Name
\end{minipage} & \begin{minipage}[b]{\linewidth}\raggedright
Description
\end{minipage} & \begin{minipage}[b]{\linewidth}\raggedright
Default
\end{minipage} \\
\midrule\noalign{}
\endhead
\bottomrule\noalign{}
\endlastfoot
1-1 & NAME & A character string of up to eight characters in length
beginning with `TU' that defines the table name. & None \\
1-2 & NDIMS & An integer value between one and MXDIMS that defines the
number of dimensions for the table. MXDIMS is the maximum number of UDT
table dimensions defined using the UDTDIMS keyword in the RUNSPEC
section. OPM Flow only supports one-dimensional tables (NDIMS = 1). &
None \\
1-3 & / & Record terminated by a ``/'' & Not Applicable \\
2-1 & TYPE & A defined character string that defines the type of
interpolation for this dimension and must one of the following: &
None \\
2-2 & POINTS & A monotonically increasing vector of real values that
defines the numerical value of the interpolation points for this
dimension if the specified type of interpolation TYPE is either `NV',
`LC' or `LL'. Otherwise, a vector of character strings of up to eight
characters in length defining the `ID'. & None \\
2-2 & / & One record for each of the NDIMS dimensions each terminated by
a ``/'' & Not Applicable \\
3-1 & VALUES & A vector of real values corresponding to each of the
interpolations points POINTS in the first dimension. & None \\
3-2 & / & Record (3-1) is included once if the number of dimensions
NDIMS equals one. Otherwise, if the number of dimensions is greater than
one then Record (3-1) is repeated for each of the interpolation POINTS
in the second dimension. Each record is terminated by a ``/''. & Not
Applicable \\
3-3 & / & Item (3-2) is followed by an additional ``/''. & Not
Applicable \\
3-4 & / & If the number of dimensions NDIMS is greater than or equal to
three then records (3-1) to (3-3) are repeated for each of the
interpolation POINTS in the third dimension, and then followed by an
additional ``/''. & Not Applicable \\
3-5 & / & If the number of dimensions NDIMS is equal to four then
records (3-1) to (3-4) are repeated for each of the interpolation POINTS
in the forth dimension, and then followed by an additional ``/''. & Not
Applicable \\
3-6 & / & The keyword is terminated by a ``/'' & Not Applicable \\
Notes: & & & \\
\end{longtable}

\emph{Table 12.3.235.1: UDT Keyword Description}

The format for a one-dimensional table with m interpolation points is
outlined below:

\begin{verbatim}
UDT
(NAME) 1 / ----------------------- NAME, NDIMS
--
-- Interpolation points for each dimension:
--
(TYPE1) a1 a2 ... am / ----------- TYPE, POINTS
--
-- VALUES: fi = f(ai)
--
f1 f2 ... fm / ------------------- end of 1D f(a1) ... f(am)
/ -------------------------------- end of 2D f(*)
/ -------------------------------- end of keyword
\end{verbatim}

The format for a two-dimensional table with m and n interpolation points
in the first and second dimensions respectively is outlined below:

\begin{verbatim}
UDT
(NAME) 2 / ----------------------- NAME, NDIMS
--
-- Interpolation points for each dimension:
--
(TYPE1) a1 a2 ... am / ----------- TYPE, POINTS
(TYPE2) b1 b2 ... bn /
--
-- VALUES: fij = f(ai,bj)
--
f11 f21 ... fm1 / ---------------- end of 1D f(a1,b1) ... f(am,b1)
f12 f22 ... fm2 /
...
f1n f2n ... fmn / ---------------- end of 1D f(a1,bn) ... f(am,bn)
/ -------------------------------- end of 2D f(*, b1) ... f(*, bn)
/ -------------------------------- end of keyword
\end{verbatim}

The format for a three-dimensional table with m, n and p interpolation
points in the first, second and third dimensions respectively is
outlined below:

\begin{verbatim}
UDT
(NAME) 3 / ----------------------- NAME, NDIMS
--
-- Interpolation points for each dimension:
--
(TYPE1) a1 a2 ... am / ----------- TYPE, POINTS
(TYPE2) b1 b2 ... bn /
(TYPE3) c1 c2 ... cp /
--
-- VALUES: fijk = f(ai,bj,ck)
--
f111 f211 ... fm11 / ------------- end of 1D f(a1,b1,c1) ... f(am,b1,c1)
f121 f221 ... fm21 /
...
f1n1 f2n1 ... fmn1 / ------------- end of 1D f(a1,bn,c1) ... f(am,bn,c1)
/ -------------------------------- end of 2D f(*, b1,c1) ... f(*, bn,c1)
f112 f212 ... fm12 /
f122 f222 ... fm22 /
...
f1n2 f2n2 ... fmn2 /
/

...

f11p f21p ... fm1p /
f12p f22p ... fm2p /
...
f1np f2np ... fmnp / ------------- end of 1D f(a1,bn,cp) ... f(am,bn,cp)
/ -------------------------------- end of 2D f(*, b1,cp) ... f(*, bn,cp)
/ -------------------------------- end of 3D f(*, *, c1) ... f(*, *, cp)
/ -------------------------------- end of keyword
\end{verbatim}

The format for a four-dimensional table with m, n, p and q interpolation
points in the first, second, third and forth dimensions respectively is
outlined below:

\begin{verbatim}
UDT
(NAME) 4 / ----------------------- NAME, NDIMS
--
-- Interpolation points for each dimension:
--
(TYPE1) a1 a2 ... am / ----------- TYPE, POINTS
(TYPE2) b1 b2 ... bn /
(TYPE3) c1 c2 ... cp /
(TYPE4) d1 d2 ... dq /
--
-- VALUES: fijkl = f(ai,bj,ck,dl)
--
--  f(*, *, *,d1):
--
f1111 f2111 ... fm111 / ---------- end of 1D f(a1,b1,c1,d1) ... f(am,b1,c1,d1)
f1211 f2211 ... fm211 /
...
f1n11 f2n11 ... fmn11 / ---------- end of 1D f(a1,bn,c1,d1) ... f(am,bn,c1,d1)
/ -------------------------------- end of 2D f(*, b1,c1,d1) ... f(*, bn,c1,d1)
f1121 f2121 ... fm121 /
f1221 f2221 ... fm221 /
...
f1n21 f2n21 ... fmn21 /
/

...

f11p1 f21p1 ... fm1p1 /
f12p1 f22p1 ... fm2p1 /
...
f1np1 f2np1 ... fmnp1 / ---------- end of 1D f(a1,bn,cp,d1) ... f(am,bn,cp,d1)
/ -------------------------------- end of 2D f(*, b1,cp,d1) ... f(*, bn,cp,d1)
/ -------------------------------- end of 3D f(*, *, c1,d1) ... f(*, *, cp d1)
--
--  f(*, *, *,d2):
--
f1112 f2112 ... fm112 / ---------- end of 1D f(a1,b1,c1,d2) ... f(am,b1,c1,d2)
f1212 f2212 ... fm212 /
...
f1n12 f2n12 ... fmn12 / ---------- end of 1D f(a1,bn,c1,d2) ... f(am,bn,c1,d2)
/ -------------------------------- end of 2D f(*, b1,c1,d2) ... f(*, bn,c1,d2)
f1122 f2122 ... fm122 /
f1222 f2222 ... fm222 /
...
f1n22 f2n22 ... fmn22 /
/

...

f11p2 f21p2 ... fm1p2 /
f12p2 f22p2 ... fm2p2 /
...
f1np2 f2np2 ... fmnp2 / ---------- end of 1D f(a1,bn,cp,d2) ... f(am,bn,cp,d2)
/ -------------------------------- end of 2D f(*, b1,cp,d2) ... f(*, bn,cp,d2)
/ -------------------------------- end of 3D f(*, *, c1,d2) ... f(*, *, cp,d2)
--
--  f(*, *, *,d3):
--

...

--
--  f(*, *, *,dq):
--
f111q f211q ... fm11q / ---------- end of 1D f(a1,b1,c1,dq) ... f(am,b1,c1,dq)
f121q f221q ... fm21q /
...
f1n1q f2n1q ... fmn1q / ---------- end of 1D f(a1,bn,c1,dq) ... f(am,bn,c1,dq)
/ -------------------------------- end of 2D f(*, b1,c1,dq) ... f(*, bn,c1,dq)
f112q f212q ... fm12q /
f122q f222q ... fm22q /
...
f1n2q f2n2q ... fmn2q /
/
...
f11pq f21pq ... fm1pq /
f12pq f22pq ... fm2pq /
...
f1npq f2npq ... fmnpq / ---------- end of 1D f(a1,bn,cp,dq) ... f(am,bn,cp,dq)
/ -------------------------------- end of 2D f(*, b1,cp,dq) ... f(*, bn,cp,dq)
/ -------------------------------- end of 3D f(*, *, c1,dq) ... f(*, *, cp,dq)
/ -------------------------------- end of 4D f(*, *, *, d1) ... f(*, *, *, dq)
/ -------------------------------- end of keyword
\end{verbatim}

\subsubsection{Examples}\label{examples-87}

The following example shows a one-dimension User Defined Table that will
linearly interpolate between values of 100 and 180 for values of FOPR
between 100 and 500 and will clamp to the end points for values of FOPR
outside this range.

\begin{verbatim}
--
-- DECLARE USER DEFINED TABLE
--
UDT
-- NAME  NDIMS
 'TU_FBHP'  1 /
-- TYPE  POINTS
 'LC'       100.0  500.0 / -- FOPR values
-- VALUES
            100.0  180.0 / -- FBHP values
/
/
--
-- DECLARE USER DEFINED QUANTITIES
--
UDQ
ASSIGN FU_WBHP 0 /
DEFINE FU_WBHP0 FU_WBHP /
DEFINE FU_WBHP (TU_FBHP[FOPR] UMIN WBHP 'PROD') UMIN FU_WBHP0 /
/
--
-- WELL PRODUCTION TARGETS AND CONSTRAINTS
--
WCONPROD
  'PROD'  'OPEN'  ORAT   500.0   4*   FU_WBHP /
/
\end{verbatim}

In the example above, the UDT specifies how the bottom hole pressure
limit should reduce as the field oil production rate decreases (note the
interpolation point values must be monotone increasing). The User
Defined Quantity FU\_WBHP uses a lookup value from the UDT in the form
TU\_BHP{[}FOPR{]}.

\subsection{USECUPL -- Load a Reservoir Coupling
File}\label{usecupl-load-a-reservoir-coupling-file}

\begin{longtable}[]{@{}
  >{\raggedright\arraybackslash}p{(\linewidth - 14\tabcolsep) * \real{0.1250}}
  >{\raggedright\arraybackslash}p{(\linewidth - 14\tabcolsep) * \real{0.1250}}
  >{\raggedright\arraybackslash}p{(\linewidth - 14\tabcolsep) * \real{0.1250}}
  >{\raggedright\arraybackslash}p{(\linewidth - 14\tabcolsep) * \real{0.1250}}
  >{\raggedright\arraybackslash}p{(\linewidth - 14\tabcolsep) * \real{0.1250}}
  >{\raggedright\arraybackslash}p{(\linewidth - 14\tabcolsep) * \real{0.1250}}
  >{\raggedright\arraybackslash}p{(\linewidth - 14\tabcolsep) * \real{0.1250}}
  >{\raggedright\arraybackslash}p{(\linewidth - 14\tabcolsep) * \real{0.1250}}@{}}
\toprule\noalign{}
\begin{minipage}[b]{\linewidth}\raggedright
RUNSPEC
\end{minipage} & \begin{minipage}[b]{\linewidth}\raggedright
GRID
\end{minipage} & \begin{minipage}[b]{\linewidth}\raggedright
EDIT
\end{minipage} & \begin{minipage}[b]{\linewidth}\raggedright
PROPS
\end{minipage} & \begin{minipage}[b]{\linewidth}\raggedright
REGIONS
\end{minipage} & \begin{minipage}[b]{\linewidth}\raggedright
SOLUTION
\end{minipage} & \begin{minipage}[b]{\linewidth}\raggedright
SUMMARY
\end{minipage} & \begin{minipage}[b]{\linewidth}\raggedright
SCHEDULE
\end{minipage} \\
\midrule\noalign{}
\endhead
\bottomrule\noalign{}
\endlastfoot
\end{longtable}

\subsubsection{Description}\label{description-1029}

The USECUPL keyword causes the simulator to read a Reservoir Coupling
file that has been previously created in a master run using the DUMPCUPL
keyword in the SCHEDULE section, for when reservoir coupling is invoked
by the GRUPMAST and SLAVES keywords in the SCHEDULE section.

This keyword is not supported by OPM Flow but would change the results
if supported so the simulation will be stopped.

\subsection{VAPPARS -- Oil Vaporization
Parameters}\label{vappars-oil-vaporization-parameters-1}

VAPPARS defines the rate of oil vaporization in the presence of
undersaturated gas and the rate at which the remaining oil gets
``heavier'' via the reduction in the solution gas-oil ratio (``Rs'').
This keyword should only be used if the OIL, GAS, DISGAS and VAPOIL
keywords in the RUNSPEC section have been invoked to allow oil, gas,
dissolved gas and vaporized oil to be present in the model.

In OPM Flow, VAPPARS can only be set once, and not changed multiple
times in the SCHEDULE section.

See VAPPARS -- Oil Vaporization Parameters in the SOLUTION section for a
full description.

\subsection{VFPCHK -- Define Production Vertical Flow Performance BHP
Check}\label{vfpchk-define-production-vertical-flow-performance-bhp-check}

\begin{longtable}[]{@{}
  >{\raggedright\arraybackslash}p{(\linewidth - 14\tabcolsep) * \real{0.1250}}
  >{\raggedright\arraybackslash}p{(\linewidth - 14\tabcolsep) * \real{0.1250}}
  >{\raggedright\arraybackslash}p{(\linewidth - 14\tabcolsep) * \real{0.1250}}
  >{\raggedright\arraybackslash}p{(\linewidth - 14\tabcolsep) * \real{0.1250}}
  >{\raggedright\arraybackslash}p{(\linewidth - 14\tabcolsep) * \real{0.1250}}
  >{\raggedright\arraybackslash}p{(\linewidth - 14\tabcolsep) * \real{0.1250}}
  >{\raggedright\arraybackslash}p{(\linewidth - 14\tabcolsep) * \real{0.1250}}
  >{\raggedright\arraybackslash}p{(\linewidth - 14\tabcolsep) * \real{0.1250}}@{}}
\toprule\noalign{}
\begin{minipage}[b]{\linewidth}\raggedright
RUNSPEC
\end{minipage} & \begin{minipage}[b]{\linewidth}\raggedright
GRID
\end{minipage} & \begin{minipage}[b]{\linewidth}\raggedright
EDIT
\end{minipage} & \begin{minipage}[b]{\linewidth}\raggedright
PROPS
\end{minipage} & \begin{minipage}[b]{\linewidth}\raggedright
REGIONS
\end{minipage} & \begin{minipage}[b]{\linewidth}\raggedright
SOLUTION
\end{minipage} & \begin{minipage}[b]{\linewidth}\raggedright
SUMMARY
\end{minipage} & \begin{minipage}[b]{\linewidth}\raggedright
SCHEDULE
\end{minipage} \\
\midrule\noalign{}
\endhead
\bottomrule\noalign{}
\endlastfoot
\end{longtable}

\subsubsection{Description}\label{description-1030}

The VFPPROD keyword defines production Vertical Flow Performance
(``VFP'') tables that are used to determine the outflow or downstream
pressure based on the inlet or upstream pressure and the phases flowing
through the system. For a well this means the table relates the flowing
bottom-hole pressure (``BHP'') to the well's tubing head pressure
(``THP'') based on the oil, gas and water rates (and any artificial lift
quantities like gas lift gas), or phases ratios, flowing up the
wellbore.

The VFP tables are generated by external programs and may result in some
inconsistencies. A common inconsistency is that the curves of BHP versus
flow rate at adjacent THP values cross, that is the BHP is increasing
with decreasing THP. This will result in numerical convergence problems
and should therefore be avoided; which is why the simulator checks for
this particular inconsistency. However, there are cases when the
external software generating the VFP table outputs ``valid'' high BHP
values, for example, when the flow exceeds the erosion velocity limits,
or the flow is supersonic, indicating a no flow condition. The VPFCHK
keyword sets the BHP check pressure (VFPCHK) for subsequent VFPPROD
tables, so that crossing BHP values above VFPCHK will be ignored.

This keyword is not supported by OPM Flow but would change the results
if supported so the simulation will be stopped.

\begin{longtable}[]{@{}
  >{\raggedright\arraybackslash}p{(\linewidth - 6\tabcolsep) * \real{0.2500}}
  >{\raggedright\arraybackslash}p{(\linewidth - 6\tabcolsep) * \real{0.2500}}
  >{\raggedright\arraybackslash}p{(\linewidth - 6\tabcolsep) * \real{0.2500}}
  >{\raggedright\arraybackslash}p{(\linewidth - 6\tabcolsep) * \real{0.2500}}@{}}
\toprule\noalign{}
\begin{minipage}[b]{\linewidth}\raggedright
No.
\end{minipage} & \begin{minipage}[b]{\linewidth}\raggedright
Name
\end{minipage} & \begin{minipage}[b]{\linewidth}\raggedright
Description
\end{minipage} & \begin{minipage}[b]{\linewidth}\raggedright
Default
\end{minipage} \\
\midrule\noalign{}
\endhead
\bottomrule\noalign{}
\endlastfoot
Field & Metric & Laboratory & \\
1 & VFPCHK & VFPCHK is a real positive value that defines the BHP
pressure above which crossing VFP curves will be ignored. Setting VFPCHK
to a large number like the default value number will cause all crossing
curves to be checked. Also if the keyword is omitted from the input deck
then the check is performed using the default value. & 1.010 \\
psia & barsa & atma & \\
Notes: & & & \\
\end{longtable}

\emph{Table 12.73: VFPCHK Keyword Description}

\begin{tcolorbox}[enhanced jigsaw, title=\textcolor{quarto-callout-note-color}{\faInfo}\hspace{0.5em}{Note}, coltitle=black, breakable, toprule=.15mm, colbacktitle=quarto-callout-note-color!10!white, leftrule=.75mm, opacitybacktitle=0.6, colback=white, bottomtitle=1mm, toptitle=1mm, titlerule=0mm, left=2mm, colframe=quarto-callout-note-color-frame, arc=.35mm, rightrule=.15mm, bottomrule=.15mm, opacityback=0]

One reason for external programs generating crossing VFP curves is that
the curves have been generated with too much resolution. For example, if
the GOR entries has been generated with values of 100, 150, 200, 250,
300, 350, 400, 450 and 500, then use a geometric spacing instead to
generated the VFP table, that is: 100, 300, 900. This will enable the
simulator to interpolate the curves consistently and avoid crossing VFP
curves.

\end{tcolorbox}

\subsubsection{Example}\label{example-558}

Here the example sets the maximum BHP to be 1.0 x106 above which
crossing VFP curves will be ignored.

\begin{verbatim}
--
--       DEFINE PRODUCTION VFP CHECK MAX BHP
--
--       MXBHP
VFPCHK
         1.0E6
/
\end{verbatim}

\subsection{VFPINJ -- Define Injection Vertical Flow Performance
Tables}\label{vfpinj-define-injection-vertical-flow-performance-tables}

\begin{longtable}[]{@{}
  >{\raggedright\arraybackslash}p{(\linewidth - 14\tabcolsep) * \real{0.1250}}
  >{\raggedright\arraybackslash}p{(\linewidth - 14\tabcolsep) * \real{0.1250}}
  >{\raggedright\arraybackslash}p{(\linewidth - 14\tabcolsep) * \real{0.1250}}
  >{\raggedright\arraybackslash}p{(\linewidth - 14\tabcolsep) * \real{0.1250}}
  >{\raggedright\arraybackslash}p{(\linewidth - 14\tabcolsep) * \real{0.1250}}
  >{\raggedright\arraybackslash}p{(\linewidth - 14\tabcolsep) * \real{0.1250}}
  >{\raggedright\arraybackslash}p{(\linewidth - 14\tabcolsep) * \real{0.1250}}
  >{\raggedright\arraybackslash}p{(\linewidth - 14\tabcolsep) * \real{0.1250}}@{}}
\toprule\noalign{}
\begin{minipage}[b]{\linewidth}\raggedright
RUNSPEC
\end{minipage} & \begin{minipage}[b]{\linewidth}\raggedright
GRID
\end{minipage} & \begin{minipage}[b]{\linewidth}\raggedright
EDIT
\end{minipage} & \begin{minipage}[b]{\linewidth}\raggedright
PROPS
\end{minipage} & \begin{minipage}[b]{\linewidth}\raggedright
REGIONS
\end{minipage} & \begin{minipage}[b]{\linewidth}\raggedright
SOLUTION
\end{minipage} & \begin{minipage}[b]{\linewidth}\raggedright
SUMMARY
\end{minipage} & \begin{minipage}[b]{\linewidth}\raggedright
SCHEDULE
\end{minipage} \\
\midrule\noalign{}
\endhead
\bottomrule\noalign{}
\endlastfoot
\end{longtable}

\subsubsection{Description}\label{description-1031}

The VFPINJ keyword defines injection Vertical Flow Performance (``VFP'')
tables that are used to determine the outflow or downstream pressure
based on the inlet or upstream pressure and the phases being injected
into the system. For an injection well this means the table relates the
flowing bottom-hole pressure (``BHP'') to the well's tubing head
pressure (``THP'') based on the oil, gas or water injection rates. The
table is also used to describe the pressure relationship when the
network option is being used, although the Network option is not
currently implemented in OPM Flow. In this case the table describes the
pipeline pressure behavior from the HIGHER group (inlet node) to the
LOWER group (outlet node) given the current flowing conditions (the
group relationship is defined by the GRUPTREE keyword in SCHEDULE
section).

The data for this keyword is generated by an external program and is
normally included into the input deck using the INCLUDE keyword as
described in section CHAPTER 4: GLOBAL SECTION KEYWORDS, as the data can
be quite voluminous.

Each VFPINJ table must be entered with a separate VFPINJ keyword that
consists of four records, with 1-1 to 1-6 representing record one items
(1) to (6), 2-1 representing record number two item (1), and so on in
the ``No.'' column in Table 12.74. Each record is terminated by a ``/''.
The fourth record must be repeated to give BHP data as a function of FLO
for all THP values.

\begin{longtable}[]{@{}
  >{\raggedright\arraybackslash}p{(\linewidth - 6\tabcolsep) * \real{0.2500}}
  >{\raggedright\arraybackslash}p{(\linewidth - 6\tabcolsep) * \real{0.2500}}
  >{\raggedright\arraybackslash}p{(\linewidth - 6\tabcolsep) * \real{0.2500}}
  >{\raggedright\arraybackslash}p{(\linewidth - 6\tabcolsep) * \real{0.2500}}@{}}
\toprule\noalign{}
\begin{minipage}[b]{\linewidth}\raggedright
No.
\end{minipage} & \begin{minipage}[b]{\linewidth}\raggedright
Name
\end{minipage} & \begin{minipage}[b]{\linewidth}\raggedright
Description
\end{minipage} & \begin{minipage}[b]{\linewidth}\raggedright
Default
\end{minipage} \\
\midrule\noalign{}
\endhead
\bottomrule\noalign{}
\endlastfoot
Field & Metric & Laboratory & \\
1-1 & VFPTAB & A positive integer greater than zero and less than or
equal to the MXVFPTAB variable as defined on the VFPIDIMS keyword in the
RUNSPEC section, that defines the vertical flow performance table
number. & None \\
1-2 & VFPREF & A real positive value that defines the reference depth
used to generate this VFPINJ table data set. OPM Flow automatically
corrects any difference between VFPREF and the BHPREF on the WELSPECS
and WPAVDEP keywords in the SCHEDULE section, using the current
hydrostatic head. & None \\
1-3 & FLO & A defined character string that defines the injection
phases, and should be set to one of the following character strings: &
None \\
1-4 & VFPTYPE & A defined character string that should be defaulted or
set equal to THP. & THP \\
1-5 & VFPUNITS & A defined character string that specifies the units
system for the VFP table. An error message is output if this is not the
same as the model units. & Model units \\
FIELD & METRIC & LAB & \\
1-6 & VFPVALUE & A defined character string that defines the type of
data in the VFP-DATA vector. This should be defaulted or set equal to
BHP. & BHP \\
1-7 & / & Record terminated by a ``/'' & Not Applicable \\
2-1 & FLO-DATA & A real positive monotonically increasing vector that
defines the numerical values of the injection phase declared by the FLO
variable. The number of entries must greater than two and less than or
equal to MXMFLO as defined on the VFPIDIMS keyword in the RUNSPEC
section. & None \\
Liquid: stb Gas: Mscf & Liquid: sm3 Gas: sm3 & Liquid: scc Gas: scc & \\
2-2 & / & Record terminated by a ``/'' & Not Applicable \\
3-1 & THP-DATA & A real positive monotonically increasing vector that
defines the numerical values of the tubing head pressure values. The
number of entries must be greater than two and less than or equal to
MXMTHP as defined on the VFPIDIMS keyword in the RUNSPEC section. &
None \\
psia & barsa & atma & \\
3-2 & / & Record terminated by a ``/'' & Not Applicable \\
4-1 & NTHP & This data record consists of an integer value that defines
the index of THP values entered via the THP-DATA records on this
keyword. For example, if THP-DATA is equal to 1000, 2000, 3000 and 3500
and NTHP is equal to three then NTHP refers to third entry, that is THP
equal to 3000. & \\
BHP-DATA & NTHP is then followed by a real vector of BHP values for each
FLO injection rate for the corresponding index value (NTHP) and is then
terminated with a''/'' The (4-1) record, which consists of both NTHP and
BHP-DATA data, is then repeated, until all combinations of (NTHP and
FLO) and the associate BHP data has been entered. & None & \\
psia & barsa & atma & \\
4-2 & / & Each Index (NTHP, BHP-DATA) data set is terminated by a ``/''
& Not Applicable \\
Notes: & & & \\
\end{longtable}

\emph{Table 12.74: VFPINJ Keyword Description}

See also the WELSPECS keyword to define wells and the WCONINJE keyword
that is used to allocate the VFPINJ tables to specific wells.

Note that one VFPINJ table can be allocated to one or more wells,
provided the wells in question have a similar trajectory and similar
flow characteristics, for example vertical water injection wells
injecting into the same reservoir.

The VFPPROD keyword is used to enter VFP tables for production wells or
to describe the pressure relationship when the network option is being
used. In this case the table describes the pipeline pressure behavior
from the LOWER group (inlet node) to the HIGHER group (outlet node)
given the current production conditions.

All the aforementioned keywords are described in the SCHEDULE section.

\subsubsection{Examples}\label{examples-88}

The following example shows the VFPINJ table for a water injection well
and is taken from the Norne OPM Flow model.

\begin{verbatim}
VFPINJ

-- Table   Datum Depth   Rate Type
-- -----   -----------   ---------
      12       2718.07       'WAT' /


-- 'WAT' units - SM3/DAY
  500.0  1263.2  2026.3  2789.5  3552.6
 4315.8  5078.9  5842.1  6605.3  7368.4
 8131.6  8894.7  9657.9 10421.1 11184.2
11947.4 12710.5 13473.7 14236.8 15000.0 /

-- 'THP' units - BARSA
  21.01   63.24  105.46  147.68  189.90
 232.12  274.35  316.57  358.79  401.01 /

1  254.51  253.95  252.27  249.83  246.69
   242.88  238.42  233.32  227.59  221.22
   214.23  206.62  198.38  189.53  180.06
   169.97  159.26  147.95  136.00  123.46
/
2  297.02  296.49  294.82  292.39  289.26
   285.47  281.01  275.92  270.20  263.84
   256.87  249.28  241.05  232.22  222.76
   212.70  202.01  190.71  178.79  166.27
/
……………………………………….………………………….……………………………………………………………………
……………………………………….………………………….……………………………………………………………………

9  594.67  594.29  592.70  590.34  587.29
   583.57  579.16  574.17  568.55  562.25
   555.40  547.92  539.79  531.09  521.74
   511.82  501.25  490.13  478.34  466.01
/
10 637.19  636.83  635.26  632.91  629.86
   626.16  621.76  616.78  611.17  604.89
   598.05  590.59  582.47  573.79  564.45
   554.56  544.01  532.91  521.14  508.83
/
\end{verbatim}

The example shows the first two and the last two records of the fourth
kind, as the data is too voluminous to be included.

\begin{tcolorbox}[enhanced jigsaw, title=\textcolor{quarto-callout-note-color}{\faInfo}\hspace{0.5em}{Note}, coltitle=black, breakable, toprule=.15mm, colbacktitle=quarto-callout-note-color!10!white, leftrule=.75mm, opacitybacktitle=0.6, colback=white, bottomtitle=1mm, toptitle=1mm, titlerule=0mm, left=2mm, colframe=quarto-callout-note-color-frame, arc=.35mm, rightrule=.15mm, bottomrule=.15mm, opacityback=0]

The VFPTAB variable defines the table number of the VFPINJ data set; if
more then one VFPINJ keyword is entered with the same VFPTAB number then
the VFPINJ data set will be overwritten by the last VFPINJ keyword with
the same VFPTAB number. The same comment is also applicable to the
VFPPROD keyword.

\end{tcolorbox}

\subsection{VFPPROD -- Define Production Vertical Flow Performance
Tables}\label{vfpprod-define-production-vertical-flow-performance-tables}

\begin{longtable}[]{@{}
  >{\raggedright\arraybackslash}p{(\linewidth - 14\tabcolsep) * \real{0.1250}}
  >{\raggedright\arraybackslash}p{(\linewidth - 14\tabcolsep) * \real{0.1250}}
  >{\raggedright\arraybackslash}p{(\linewidth - 14\tabcolsep) * \real{0.1250}}
  >{\raggedright\arraybackslash}p{(\linewidth - 14\tabcolsep) * \real{0.1250}}
  >{\raggedright\arraybackslash}p{(\linewidth - 14\tabcolsep) * \real{0.1250}}
  >{\raggedright\arraybackslash}p{(\linewidth - 14\tabcolsep) * \real{0.1250}}
  >{\raggedright\arraybackslash}p{(\linewidth - 14\tabcolsep) * \real{0.1250}}
  >{\raggedright\arraybackslash}p{(\linewidth - 14\tabcolsep) * \real{0.1250}}@{}}
\toprule\noalign{}
\begin{minipage}[b]{\linewidth}\raggedright
RUNSPEC
\end{minipage} & \begin{minipage}[b]{\linewidth}\raggedright
GRID
\end{minipage} & \begin{minipage}[b]{\linewidth}\raggedright
EDIT
\end{minipage} & \begin{minipage}[b]{\linewidth}\raggedright
PROPS
\end{minipage} & \begin{minipage}[b]{\linewidth}\raggedright
REGIONS
\end{minipage} & \begin{minipage}[b]{\linewidth}\raggedright
SOLUTION
\end{minipage} & \begin{minipage}[b]{\linewidth}\raggedright
SUMMARY
\end{minipage} & \begin{minipage}[b]{\linewidth}\raggedright
SCHEDULE
\end{minipage} \\
\midrule\noalign{}
\endhead
\bottomrule\noalign{}
\endlastfoot
\end{longtable}

\subsubsection{Description}\label{description-1032}

The VFPPROD keyword defines production Vertical Flow Performance
(``VFP'') tables that are used to determine the outflow or downstream
pressure based on the inlet or upstream pressure and the phases flowing
through the system. For a production well this means the table relates
the flowing bottom-hole pressure (``BHP'') to the well's tubing head
pressure (``THP'') based on the oil, gas and water rates (and any
artificial lift quantities (``ALQ'') like gas lift gas), or phases
ratios, flowing up the wellbore. The table is also used to describe the
pressure relationship when the network option is being used. In this
case the table describes the pipeline pressure behavior from the LOWER
group (inlet node) to the HIGHER group (outlet node) given the current
flowing conditions (the group relationship is defined by the GRUPTREE
keyword in SCHEDULE section).

Each VFPPROD table must be entered with a separate VFPPROD keyword that
consists of seven records, with 1-1 to 1-9 representing record one items
(1) to (9), 2-1 representing record two item (1), and so on in the
``No.'' column in Table 12.75. Each record is terminated by a ``/''.

The seventh record should be repeated to give BHP data as a function of
FLO for each combination of THP, WFR, GFR, and ALQ values. If record
seven is only entered once then the BHP values will depend only on flow
rate (FLO), and will be replicated across all combinations of THP, WFR,
GFR, and ALQ values.

\begin{longtable}[]{@{}
  >{\raggedright\arraybackslash}p{(\linewidth - 6\tabcolsep) * \real{0.2500}}
  >{\raggedright\arraybackslash}p{(\linewidth - 6\tabcolsep) * \real{0.2500}}
  >{\raggedright\arraybackslash}p{(\linewidth - 6\tabcolsep) * \real{0.2500}}
  >{\raggedright\arraybackslash}p{(\linewidth - 6\tabcolsep) * \real{0.2500}}@{}}
\toprule\noalign{}
\begin{minipage}[b]{\linewidth}\raggedright
No.
\end{minipage} & \begin{minipage}[b]{\linewidth}\raggedright
Name
\end{minipage} & \begin{minipage}[b]{\linewidth}\raggedright
Description
\end{minipage} & \begin{minipage}[b]{\linewidth}\raggedright
Default
\end{minipage} \\
\midrule\noalign{}
\endhead
\bottomrule\noalign{}
\endlastfoot
Field & Metric & Laboratory & \\
1-1 & VFPTAB & A positive integer greater than zero and less than or
equal to the MXVFPTAB variable as defined on the VFPPDIMS keyword in the
RUNSPEC section, that defines the vertical lift performance table
number. & None \\
1-2 & VFPREF & A real positive value that defines the reference depth
used to generate this VFPPROD table data set. OPM Flow automatically
corrects any difference between VFPREF and the BHPREF on the WELSPECS
and WPAVEDEP keywords in the SCHEDULE section, using the current
hydrostatic head. & None \\
1-3 & FLO & A defined character string that defines the flowing phases,
and should be set to one of the following character strings: & None \\
1-4 & WFR & A defined character string that defines the flowing water
fraction and should be set to one of the following character strings: &
None \\
1-5 & GFR & A defined character string that defines the flowing gas
fraction and should be set to one of the following character strings: &
None \\
1-6 & VFPTYPE & A defined character string that should be defaulted or
set equal to THP. & THP \\
1-7 & ALQ & A defined character string that defines the artificial lift
quantity and should be set to one of the following character strings:
The DENO and DENG options are not supported by OPM Flow. The ALQ
parameter is just another variable used to interpolate the outflow
pressure based on the inlet pressure, together with the phases flowing
through the system. As such, ALQ can represent any parameter; however,
the units should be consistent with that used to set the value. For
example, if pump speed (Hz) is used to set a well's artificial lift
quantity via the WCONPROD(ALQ-WELL) parameter in the SCHEDULE section,
then the ALQ-DATA should represent pump speed in Hz. The default value
is ' ' or undefined, that covers the case when the ALQ variable is not
entered, except for when gas lift is employed in the model. When gas
lift is active then the default value for ALQ is set to GRAT provided:
In addition, if any of the above is true than ALQ must be set to GRAT. &
' ' \\
1-8 & VFPUNITS & A defined character string that specifies the units
system for the VFP table. An error message is output if this is not the
same as the model units. & Model units \\
FIELD & METRIC & LAB & \\
1-9 & VFPVALUE & A defined character string that defines the type of
data in the VFP-DATA vector. This should be set equal to BHP if the
vector contains bottom-hole pressure data, or TEMP if the vector
contains Tubing Head Temperature (THT) data. OPM Flow only supports the
(default) BHP option. & BHP \\
1-10 & / & Record terminated by a ``/'' & Not Applicable \\
2-1 & FLO-DATA & A real positive monotonically increasing vector that
defines the numerical values of the flowing phase declared by the FLO
variable. The number of entries must greater than two and less than or
equal to MXMFLO as defined on the VFPPDIMS keyword in the RUNSPEC
section. & None \\
Liquid: stb Gas: Mscf & Liquid: sm3 Gas: sm3 & Liquid: scc Gas: scc & \\
2-2 & / & Record terminated by a ``/'' & Not Applicable \\
3-1 & THP-DATA & A real positive monotonically increasing vector that
defines the numerical values of the tubing head pressure values. The
number of entries must greater than two and less than or equal to MXMTHP
as defined on the VFPPDIMS keyword in the RUNSPEC section. & None \\
psia & barsa & atma & \\
3-2 & / & Record terminated by a ``/'' & Not Applicable \\
4-1 & WFR-DATA & A real positive monotonically increasing vector that
defines the numerical values of the flowing water fraction declared by
the WFR variable. The number of entries must greater than two and less
than or equal to MXMWFR as defined on the VFPPDIMS keyword in the
RUNSPEC section. & None \\
WOR: stb/stb WCT: stb/stb WGR: stb/Mscf & sm3/sm3 sm3/sm3 sm3/sm3 &
scc/scc scc/scc scc/scc & \\
4-2 & / & Record terminated by a ``/'' & Not Applicable \\
5-1 & GFR-DATA & A real positive monotonically increasing vector that
defines the numerical values of the flowing gas fraction declared by the
GFR variable. The number of entries must greater than two and less than
or equal to MXMGFR as defined on the VFPPDIMS keyword in the RUNSPEC
section. & None \\
GOR: Mscf/stb GLR: Mscf/stb OGR: stb/Mscf & sm3/sm3 sm3/sm3 sm3/sm3 &
scc/scc scc/scc scc/scc & \\
5-2 & / & Record terminated by a ``/'' & Not Applicable \\
6-1 & ALQ-DATA & A real positive monotonically increasing vector that
defines the numerical values of the artificial lift quantity declared by
the ALQ variable. The number of entries must greater than two and less
than or equal to MXMALQ as defined on the VFPPDIMS keyword in the
RUNSPEC section. & None \\
GRAT: Mscf/d IGLR: Mscf/stb TGLR: Mscf/stb DENO: lb/ft3 DENG: lb/ft3
BEAN: 1/64 inch & sm3/day sm3/sm3 sm3/sm3 kg/m3 kg/m3 mm & scc/hour
scc/scc scc/scc gm/cc gm/cc mm & \\
6-2 & / & Record terminated by a ``/'' & Not Applicable \\
7-1 & NTHP & This data record consists of a series of integer values
that defines the index of THP, WFR, GFR, ALQ entered via the those
records on this keyword. The first index, NTHP, is an integer value that
defines the index of THP values entered via the THP-DATA records on this
keyword. For example, if THP-DATA is equal to 100, 200, 300 and 350 and
NTHP is equal to three then NTHP refers to third entry, that is THP
equal to 300. & None \\
NWFR & The second index, NWFR, is an integer value that defines the
index of the water fraction values entered via the WFR-DATA records on
this keyword. For example, if WFR-DATA is equal to 0.00, 0.25, 0.50 and
0.75 and NWFR is equal to two then NWFR refers to second entry, that is
WFR equal to 0.25. & None & \\
NGFR & The third index, NGFR, is an integer value that defines the index
of the gas fraction values entered via the GFR-DATA records on this
keyword. For example, if GFR-DATA is equal to 100.0, 200.0, 500.0 and
750.0 and NGFR is equal to three then NGFR refers to third entry, that
is GFR equal to 500.0. & None & \\
NALQ & The fourth and final index entry, NALQ, is an integer value that
defines the index of artificial lift values via the ALQ-DATA records on
this keyword. For example, if ALQ-DATA is equal to 50, 100, 200 and 300
and NALQ is equal to one then NALQ refers to first entry, that is ALQ
equal to 50. The fourth index is then followed by the VFP-DATA,
containing either the BHP or THT vector values. & None & \\
VFP-DATA & VFP-DATA is a real vector of either BHP or THT values for
each FLO production rate for the corresponding index value (NTHP, NWFR,
NGFR, NALQ) and is then terminated with a''/''. The (7-1) record, which
consists of the four indices and BHP (or THT) data, is then repeated
until all combinations of (NTHP, NWFR, NGFR, NALQ) and the associate BHP
(or THT) data has been entered. Note that the VFPVALUE parameter,
determines the type of VFP-DATA, that is BHP for bottom-hole pressure,
or THT for tubing head temperature. & None & \\
psia & barsa & atma & \\
7-2 & / & Each Index (NTHP, NWFR, NGFR, NALQ. VFP-DATA) data set is
terminated by a ``/'' & Not Applicable \\
Notes: & & & \\
\end{longtable}

\emph{Table 12.75: VFPPROD Keyword Description}

The data for this keyword is generated by an external program and is
normally included into the input deck using the INCLUDE keyword as
described in section CHAPTER 4: GLOBAL SECTION KEYWORDS, as the data can
be quite voluminous.

Note that for equivalent two phase runs:

\begin{verbatim}
- For example oil-water runs with only the OIL and WATER keywords in the RUNSPEC section, or runs that model dead oil [“Dead” oil is oil that it contains no dissolved gas or a relatively thick oil or residue that has lost its volatile components.] with a constant solution gas-oil ratio value defined by the RSCONST keyword in the PROPS section, then the FLO parameter in Table 12.75 must be set to either OIL or LIQ, WFR to either WCT or WOR, and GFR to GOR.
- Gas-water models with only the GAS and WATER keywords in the RUNSPEC section, or models that only have dry gas [Natural gas that occurs in the absence of condensate or liquid hydrocarbons, or gas that had condensable hydrocarbons removed, is called dry gas. It is primarily methane with some intermediates. The hydrocarbon mixture is solely gas in the reservoir and there is no liquid (condensate surface liquid) formed either in the reservoir or at surface. The term dry indicates that the gas does not contain heavier hydrocarbons to form liquids at the surface conditions. Dry gas typically has GOR's greater than 100,000 scf/stb or 18,000 Sm3/m3.] with a constant condensate-gas ratio defined by the RVCONST keyword in the PROPS section, then the FLO parameter in Table 12.75 must be set to GAS, WFR to WGR, and GFR to OGR.
\end{verbatim}

\begin{tcolorbox}[enhanced jigsaw, title=\textcolor{quarto-callout-note-color}{\faInfo}\hspace{0.5em}{Note}, coltitle=black, breakable, toprule=.15mm, colbacktitle=quarto-callout-note-color!10!white, leftrule=.75mm, opacitybacktitle=0.6, colback=white, bottomtitle=1mm, toptitle=1mm, titlerule=0mm, left=2mm, colframe=quarto-callout-note-color-frame, arc=.35mm, rightrule=.15mm, bottomrule=.15mm, opacityback=0]

It is possible to have only the OIL and WATER keywords in the RUNSPEC
section and to use gas lift for the wells, without declaring the GAS
phase in the RUNSPEC section. In this case, the FLO parameter in Table
12.75 must be set to either OIL or LIQ, WFR to either WCT or WOR, and
GFR to GOR. In this case the ALQ parameter is optional, but if present
must set to GRAT. If the ALQ and ALQ-DATA parameters are absent then the
GFR-DATA will be used based on the flowing GOR plus the stipulated gas
lift gas.

\end{tcolorbox}

See also the WELSPECS keyword to define wells and the WCONPROD keyword
that is used to allocate the VFPPROD tables to specific wells. Note that
one VFPPROD table can be allocated to one or more wells, provided the
wells in question have a similar trajectory and similar flow
characteristics, for example vertical oil wells producing from the same
reservoir, or different reservoirs with similar PVT properties.

The VFPINJ keyword is used to enter VFP tables for injection wells or to
describe the pressure relationship when the network option is being
used. In this case the table describes the pipeline pressure behavior
from the HIGHER group (inlet node) to the LOWER group (outlet node)
given the current injection conditions.

All the aforementioned keywords are described in the SCHEDULE section.

\subsubsection{Examples}\label{examples-89}

The following example shows the VFPPROD table for a production gas well
and is taken from the Norne OPM Flow model. Here WFR has been set to
water-gas ratio and GFR has been set to the oil-gas ratio, and the ALQ
value is defaulted.

\begin{verbatim}
VFPPROD

-- Table   Datum Depth   Rate Type   WFR Type   GFR Type
-- -----   -----------   ---------   --------   --------
       5       2623.39       'GAS'      'WGR'      'OGR' /


-- 'GAS' units - SM3/DAY
  50000.0  100000.0  200000.0  400000.0  800000.0
1200000.0 1600000.0 1999999.9 3000000.0 3999999.8
5000000.5 /

-- 'THP' units - BARSA
  10.00   20.00   40.00   80.00  120.00
 150.00  200.00  250.00 /

-- 'WGR' units - SM3/SM3
      0    1e-9    1e-6    1e-5  0.0001
  0.001    0.01     0.1 /

-- 'OGR' units - SM3/SM3
   1e-7    1e-6    1e-5  0.0001   0.001
   0.01 /

-- 'ALQ' units -
      0 /

 1  1  1  1   11.93   12.22   13.35   17.24   27.93
              39.83   52.06   64.38   95.20  125.89
             156.52
/
 1  1  2  1   11.93   12.22   13.35   17.24   27.94
              39.84   52.07   64.39   95.21  125.91
             156.55
/

……………………………………….………………………….……………………………………………………………………
……………………………………….………………………….……………………………………………………………………

 8  8  5  1  483.75  511.15  614.09 1044.78 2757.56
            5592.55 9528.36 14567.24 32005.79 56375.24
              87684
/
 8  8  6  1  487.68  516.24  624.74 1075.40 2860.16
            5803.92 9880.58 15093.76 33119.59 58297.57
              90639
/
\end{verbatim}

The example shows the first two and the last two records of type seven,
as the data is too voluminous to be included.

The next example below shows an example oil producing well VFPPROD,
again taken from Norne OPM Flow model. Here WFR has been set to water
cut and GFR has been set to the gas-oil ratio, and the ALQ value is
defaulted.

\begin{verbatim}
VFPPROD

-- Table   Datum Depth   Rate Type   WFR Type   GFR Type   TAB Type
-- -----   -----------   ---------   --------   --------   --------
--    37       2641.02       'LIQ'      'WCT'      'GOR'            /

-- Prosper files are corrected from RKB to MSL depth. lmarr

-- Table   Datum Depth   Rate Type   WFR Type   GFR Type   TAB Type
-- -----   -----------   ---------   --------   --------   --------
      37       2617.02       'LIQ'      'WCT'      'GOR'            /

-- 'LIQ' units - SM3/DAY
  200.0   500.0  1000.0  1500.0  2000.0
 2500.0  3000.0  3500.0  4000.0  4500.0
 5000.0  5500.0  6000.0  6500.0  7000.0
 7500.0  8000.0 10000.0 14000.0 /

-- 'THP' units - BARSA
  21.01   51.01   61.01   81.01  101.01
 121.01  141.01  161.01  181.01  201.01 /

-- 'WCT' units - FRACTION
      0     0.1     0.2     0.3     0.4
    0.5     0.6     0.7     0.8       1 /

-- 'GOR' units - SM3/SM3
     90     100     150     200     500
   1000    2000 /

-- 'ALQ' units -
      0 /

 1  1  1  1  160.82  136.70  119.79  115.86  117.38
             121.16  126.08  131.56  137.48  143.74
             150.29  157.07  164.02  171.07  178.13
             185.11  192.09  220.38  280.86
/
 1  1  2  1  155.63  129.40  112.32  108.64  110.44
             114.74  120.15  126.09  132.47  139.05
             146.02  153.41  160.67  167.91  175.13
             182.34  189.55  218.81  281.02
/
……………………………………….………………………….……………………………………………………………………
……………………………………….………………………….……………………………………………………………………

10 10  6  1  439.30  437.95  437.53  437.79  438.39
             439.26  440.36  441.67  443.19  444.92
             446.85  448.99  451.32  453.85  456.58
             459.51  462.64  477.11  515.47
/
10 10  7  1  439.30  437.95  437.53  437.79  438.39
             439.26  440.36  441.67  443.19  444.92
             446.85  448.99  451.32  453.85  456.58
             459.51  462.64  477.11  515.47
/
\end{verbatim}

The example shows the first two and the last two records of type seven,
as the data is too voluminous to be included.

\subsection{VFPTABL -- Define Production Vertical Flow Performance ALQ
Interpolation}\label{vfptabl-define-production-vertical-flow-performance-alq-interpolation}

\begin{longtable}[]{@{}
  >{\raggedright\arraybackslash}p{(\linewidth - 14\tabcolsep) * \real{0.1250}}
  >{\raggedright\arraybackslash}p{(\linewidth - 14\tabcolsep) * \real{0.1250}}
  >{\raggedright\arraybackslash}p{(\linewidth - 14\tabcolsep) * \real{0.1250}}
  >{\raggedright\arraybackslash}p{(\linewidth - 14\tabcolsep) * \real{0.1250}}
  >{\raggedright\arraybackslash}p{(\linewidth - 14\tabcolsep) * \real{0.1250}}
  >{\raggedright\arraybackslash}p{(\linewidth - 14\tabcolsep) * \real{0.1250}}
  >{\raggedright\arraybackslash}p{(\linewidth - 14\tabcolsep) * \real{0.1250}}
  >{\raggedright\arraybackslash}p{(\linewidth - 14\tabcolsep) * \real{0.1250}}@{}}
\toprule\noalign{}
\begin{minipage}[b]{\linewidth}\raggedright
RUNSPEC
\end{minipage} & \begin{minipage}[b]{\linewidth}\raggedright
GRID
\end{minipage} & \begin{minipage}[b]{\linewidth}\raggedright
EDIT
\end{minipage} & \begin{minipage}[b]{\linewidth}\raggedright
PROPS
\end{minipage} & \begin{minipage}[b]{\linewidth}\raggedright
REGIONS
\end{minipage} & \begin{minipage}[b]{\linewidth}\raggedright
SOLUTION
\end{minipage} & \begin{minipage}[b]{\linewidth}\raggedright
SUMMARY
\end{minipage} & \begin{minipage}[b]{\linewidth}\raggedright
SCHEDULE
\end{minipage} \\
\midrule\noalign{}
\endhead
\bottomrule\noalign{}
\endlastfoot
\end{longtable}

\subsubsection{Description}\label{description-1033}

The VFPTABL keyword defines the interpolation method for production
Vertical Flow Performance (``VFP'') tables for the Artificial Lift
Quantity (``ALQ''). Production VFP data is entered via the VFPPROD
keyword in the SCHEDULE section. By default the simulator interpolates
all the variables in the VFP tables using linear interpolation,
including the ALQ quantity. However, if the ALQ values represent gas
lift, then linear interpolation may not be insufficient, as the gradient
change between the tabulated ALQ values may result in sudden changes.
This is particularly important in gas lift optimization studies where
the available gas lift gas is being allocated to a group of wells in
order to maximize oil production rates. To overcome this issue the
VFPTABL keyword allows the ALQ values to be interpolated using cubic
spline interpolation, and results in a smother transition between the
various ALQ entries.

This keyword is not supported by OPM Flow but would change the results
if supported so the simulation will be stopped.

\begin{longtable}[]{@{}
  >{\raggedright\arraybackslash}p{(\linewidth - 6\tabcolsep) * \real{0.2500}}
  >{\raggedright\arraybackslash}p{(\linewidth - 6\tabcolsep) * \real{0.2500}}
  >{\raggedright\arraybackslash}p{(\linewidth - 6\tabcolsep) * \real{0.2500}}
  >{\raggedright\arraybackslash}p{(\linewidth - 6\tabcolsep) * \real{0.2500}}@{}}
\toprule\noalign{}
\begin{minipage}[b]{\linewidth}\raggedright
No.
\end{minipage} & \begin{minipage}[b]{\linewidth}\raggedright
Name
\end{minipage} & \begin{minipage}[b]{\linewidth}\raggedright
Description
\end{minipage} & \begin{minipage}[b]{\linewidth}\raggedright
Default
\end{minipage} \\
\midrule\noalign{}
\endhead
\bottomrule\noalign{}
\endlastfoot
Field & Metric & Laboratory & \\
1 & VFPTABL & VFPTABL is a defined positive integer that specifies the
interpolation method to be used with the ALQ quantity in the VFP
production tables, and should be set to one of the following: & 1 \\
dimensionless & dimensionless & dimensionless & \\
Notes: & & & \\
\end{longtable}

\emph{Table 12.76: VFPTABL Keyword Description}

\subsubsection{Example}\label{example-559}

The example sets cubic spline interpolation for the ALQ quantity in the
VFPPROD tables, with linear interpolation used for all the variables.

\begin{verbatim}
--
--       ALQ INTERPOLATION OPTION
--
--       OPTION
VFPTABL
         2
/
\end{verbatim}

\subsection{WAITBAL -- Wait On Network Balance Before Allowing Further
Actions}\label{waitbal-wait-on-network-balance-before-allowing-further-actions}

\begin{longtable}[]{@{}
  >{\raggedright\arraybackslash}p{(\linewidth - 14\tabcolsep) * \real{0.1250}}
  >{\raggedright\arraybackslash}p{(\linewidth - 14\tabcolsep) * \real{0.1250}}
  >{\raggedright\arraybackslash}p{(\linewidth - 14\tabcolsep) * \real{0.1250}}
  >{\raggedright\arraybackslash}p{(\linewidth - 14\tabcolsep) * \real{0.1250}}
  >{\raggedright\arraybackslash}p{(\linewidth - 14\tabcolsep) * \real{0.1250}}
  >{\raggedright\arraybackslash}p{(\linewidth - 14\tabcolsep) * \real{0.1250}}
  >{\raggedright\arraybackslash}p{(\linewidth - 14\tabcolsep) * \real{0.1250}}
  >{\raggedright\arraybackslash}p{(\linewidth - 14\tabcolsep) * \real{0.1250}}@{}}
\toprule\noalign{}
\begin{minipage}[b]{\linewidth}\raggedright
RUNSPEC
\end{minipage} & \begin{minipage}[b]{\linewidth}\raggedright
GRID
\end{minipage} & \begin{minipage}[b]{\linewidth}\raggedright
EDIT
\end{minipage} & \begin{minipage}[b]{\linewidth}\raggedright
PROPS
\end{minipage} & \begin{minipage}[b]{\linewidth}\raggedright
REGIONS
\end{minipage} & \begin{minipage}[b]{\linewidth}\raggedright
SOLUTION
\end{minipage} & \begin{minipage}[b]{\linewidth}\raggedright
SUMMARY
\end{minipage} & \begin{minipage}[b]{\linewidth}\raggedright
SCHEDULE
\end{minipage} \\
\midrule\noalign{}
\endhead
\bottomrule\noalign{}
\endlastfoot
\end{longtable}

\subsubsection{Description}\label{description-1034}

This keyword sets the network balance option for all networks when
networks are active in the model. Basically, the keyword either
activates the PRORDER and GDRILPOT stipulated actions before or after
the network has been balanced

The network option is normally used to ensure that the tubing head
pressure (``THP'') of a group of wells flowing into a common network
node is consistent with a group's flow rates, that is each well's THP is
flowing at the same THP and at the same time satisfying well and group
targets and constraints. This is accomplished by calculating the well
THP limits dynamically by balancing the flow rates and pressure losses
in the network.

This keyword is not supported by OPM Flow but would change the results
if supported so the simulation will be stopped.

\subsection{WALKALIN -- Define Water Injection Alkaline
Concentration}\label{walkalin-define-water-injection-alkaline-concentration}

\begin{longtable}[]{@{}
  >{\raggedright\arraybackslash}p{(\linewidth - 14\tabcolsep) * \real{0.1250}}
  >{\raggedright\arraybackslash}p{(\linewidth - 14\tabcolsep) * \real{0.1250}}
  >{\raggedright\arraybackslash}p{(\linewidth - 14\tabcolsep) * \real{0.1250}}
  >{\raggedright\arraybackslash}p{(\linewidth - 14\tabcolsep) * \real{0.1250}}
  >{\raggedright\arraybackslash}p{(\linewidth - 14\tabcolsep) * \real{0.1250}}
  >{\raggedright\arraybackslash}p{(\linewidth - 14\tabcolsep) * \real{0.1250}}
  >{\raggedright\arraybackslash}p{(\linewidth - 14\tabcolsep) * \real{0.1250}}
  >{\raggedright\arraybackslash}p{(\linewidth - 14\tabcolsep) * \real{0.1250}}@{}}
\toprule\noalign{}
\begin{minipage}[b]{\linewidth}\raggedright
RUNSPEC
\end{minipage} & \begin{minipage}[b]{\linewidth}\raggedright
GRID
\end{minipage} & \begin{minipage}[b]{\linewidth}\raggedright
EDIT
\end{minipage} & \begin{minipage}[b]{\linewidth}\raggedright
PROPS
\end{minipage} & \begin{minipage}[b]{\linewidth}\raggedright
REGIONS
\end{minipage} & \begin{minipage}[b]{\linewidth}\raggedright
SOLUTION
\end{minipage} & \begin{minipage}[b]{\linewidth}\raggedright
SUMMARY
\end{minipage} & \begin{minipage}[b]{\linewidth}\raggedright
SCHEDULE
\end{minipage} \\
\midrule\noalign{}
\endhead
\bottomrule\noalign{}
\endlastfoot
\end{longtable}

\subsubsection{Description}\label{description-1035}

WAKALIN keyword defines the water injection alkaline concentration for
water injection wells for when the surfactant and/or polymer models have
been activated by the SURFACT, SURFACTW, or the POLYMER keywords in the
RUNSPEC section, combined with the ALKALINE keyword which is also in the
RUNSPEC section.

This keyword is not supported by OPM Flow but would change the results
if supported so the simulation will be stopped.

\subsection{WALQCALC -- Define Well VFP Surface ALQ Phase
Density}\label{walqcalc-define-well-vfp-surface-alq-phase-density}

\begin{longtable}[]{@{}
  >{\raggedright\arraybackslash}p{(\linewidth - 14\tabcolsep) * \real{0.1250}}
  >{\raggedright\arraybackslash}p{(\linewidth - 14\tabcolsep) * \real{0.1250}}
  >{\raggedright\arraybackslash}p{(\linewidth - 14\tabcolsep) * \real{0.1250}}
  >{\raggedright\arraybackslash}p{(\linewidth - 14\tabcolsep) * \real{0.1250}}
  >{\raggedright\arraybackslash}p{(\linewidth - 14\tabcolsep) * \real{0.1250}}
  >{\raggedright\arraybackslash}p{(\linewidth - 14\tabcolsep) * \real{0.1250}}
  >{\raggedright\arraybackslash}p{(\linewidth - 14\tabcolsep) * \real{0.1250}}
  >{\raggedright\arraybackslash}p{(\linewidth - 14\tabcolsep) * \real{0.1250}}@{}}
\toprule\noalign{}
\begin{minipage}[b]{\linewidth}\raggedright
RUNSPEC
\end{minipage} & \begin{minipage}[b]{\linewidth}\raggedright
GRID
\end{minipage} & \begin{minipage}[b]{\linewidth}\raggedright
EDIT
\end{minipage} & \begin{minipage}[b]{\linewidth}\raggedright
PROPS
\end{minipage} & \begin{minipage}[b]{\linewidth}\raggedright
REGIONS
\end{minipage} & \begin{minipage}[b]{\linewidth}\raggedright
SOLUTION
\end{minipage} & \begin{minipage}[b]{\linewidth}\raggedright
SUMMARY
\end{minipage} & \begin{minipage}[b]{\linewidth}\raggedright
SCHEDULE
\end{minipage} \\
\midrule\noalign{}
\endhead
\bottomrule\noalign{}
\endlastfoot
\end{longtable}

\subsubsection{Description}\label{description-1036}

The WALQCALC keyword defines the well VFP surface ALQ phase density use
in the VFP table lookup and interpolation to be gas surface density, oil
surface density, or neither. Note that the user should ensure that
generated VFP tables have been generated consistent with the setting on
this keyword.

This keyword is not supported by OPM Flow but would change the results
if supported so the simulation will be stopped.

\subsection{WAPI -- Define Oil Well Injection API
Gravity}\label{wapi-define-oil-well-injection-api-gravity}

\begin{longtable}[]{@{}
  >{\raggedright\arraybackslash}p{(\linewidth - 14\tabcolsep) * \real{0.1250}}
  >{\raggedright\arraybackslash}p{(\linewidth - 14\tabcolsep) * \real{0.1250}}
  >{\raggedright\arraybackslash}p{(\linewidth - 14\tabcolsep) * \real{0.1250}}
  >{\raggedright\arraybackslash}p{(\linewidth - 14\tabcolsep) * \real{0.1250}}
  >{\raggedright\arraybackslash}p{(\linewidth - 14\tabcolsep) * \real{0.1250}}
  >{\raggedright\arraybackslash}p{(\linewidth - 14\tabcolsep) * \real{0.1250}}
  >{\raggedright\arraybackslash}p{(\linewidth - 14\tabcolsep) * \real{0.1250}}
  >{\raggedright\arraybackslash}p{(\linewidth - 14\tabcolsep) * \real{0.1250}}@{}}
\toprule\noalign{}
\begin{minipage}[b]{\linewidth}\raggedright
RUNSPEC
\end{minipage} & \begin{minipage}[b]{\linewidth}\raggedright
GRID
\end{minipage} & \begin{minipage}[b]{\linewidth}\raggedright
EDIT
\end{minipage} & \begin{minipage}[b]{\linewidth}\raggedright
PROPS
\end{minipage} & \begin{minipage}[b]{\linewidth}\raggedright
REGIONS
\end{minipage} & \begin{minipage}[b]{\linewidth}\raggedright
SOLUTION
\end{minipage} & \begin{minipage}[b]{\linewidth}\raggedright
SUMMARY
\end{minipage} & \begin{minipage}[b]{\linewidth}\raggedright
SCHEDULE
\end{minipage} \\
\midrule\noalign{}
\endhead
\bottomrule\noalign{}
\endlastfoot
\end{longtable}

\subsubsection{Description}\label{description-1037}

This keyword defines an oil injection well's API gravity for when API
tracking has been made active via the API keyword in the RUNSPEC
section. The American Petroleum Institute (API) classifies oils based on
an API gravity (γAPI), or degrees API (oAPI), the relationship between
relative density (γo) of oil and API gravity (γAPI) is given by:

\begin{equation}\phantomsection\label{eq-12-33}{
{\mathrm{γ}}_{\mathit{API}} = \frac{141.5}{{\mathrm{γ}}_{o}} - 131.5
}\end{equation}

This keyword is not supported by OPM Flow but would change the results
if supported so the simulation will be stopped.

\subsection{WARN -- Activate Warning
Messages}\label{warn-activate-warning-messages-8}

Turns on warning messages to be printed to the print file (*.PRT); note
that this keyword is activated by default and can subsequently be
switched off by the NOWARN activation keyword. The warning messages may
be turned on and off using keywords WARN and NOWARN. OPM Flow always
prints error messages.

See WARN -- Activate Warning Messages in the GLOBAL section for a full
description.

\subsection{WBHGLR -- Define Well Bottom-Hole GLR
Constraint}\label{wbhglr-define-well-bottom-hole-glr-constraint}

\begin{longtable}[]{@{}
  >{\raggedright\arraybackslash}p{(\linewidth - 14\tabcolsep) * \real{0.1250}}
  >{\raggedright\arraybackslash}p{(\linewidth - 14\tabcolsep) * \real{0.1250}}
  >{\raggedright\arraybackslash}p{(\linewidth - 14\tabcolsep) * \real{0.1250}}
  >{\raggedright\arraybackslash}p{(\linewidth - 14\tabcolsep) * \real{0.1250}}
  >{\raggedright\arraybackslash}p{(\linewidth - 14\tabcolsep) * \real{0.1250}}
  >{\raggedright\arraybackslash}p{(\linewidth - 14\tabcolsep) * \real{0.1250}}
  >{\raggedright\arraybackslash}p{(\linewidth - 14\tabcolsep) * \real{0.1250}}
  >{\raggedright\arraybackslash}p{(\linewidth - 14\tabcolsep) * \real{0.1250}}@{}}
\toprule\noalign{}
\begin{minipage}[b]{\linewidth}\raggedright
RUNSPEC
\end{minipage} & \begin{minipage}[b]{\linewidth}\raggedright
GRID
\end{minipage} & \begin{minipage}[b]{\linewidth}\raggedright
EDIT
\end{minipage} & \begin{minipage}[b]{\linewidth}\raggedright
PROPS
\end{minipage} & \begin{minipage}[b]{\linewidth}\raggedright
REGIONS
\end{minipage} & \begin{minipage}[b]{\linewidth}\raggedright
SOLUTION
\end{minipage} & \begin{minipage}[b]{\linewidth}\raggedright
SUMMARY
\end{minipage} & \begin{minipage}[b]{\linewidth}\raggedright
SCHEDULE
\end{minipage} \\
\midrule\noalign{}
\endhead
\bottomrule\noalign{}
\endlastfoot
\end{longtable}

\subsubsection{Description}\label{description-1038}

This keyword, WBHGLR, defines a well's bottom-hole Gas Liquid Ratio
(``GLR'') constraint, where the GLR is the is the ratio of the ``free''
gas rate and liquid rate at bottom-hole conditions. The reference depth
for bottom-hole conditions is given by the BHPREF variable on the
WELSPECS keyword in the SCHEDULE section.

Normally this type of well control is applied to pumping wells to avoid
the well ``pumping off'', that is when the liquid column above the pump
is low, resulting in an increase in gas intake and an associated loss in
pump efficiency.

This keyword is not supported by OPM Flow but would change the results
if supported so the simulation will be stopped.

\subsection{WBOREVOL -- Define Effective Wellbore Storage
Volume}\label{wborevol-define-effective-wellbore-storage-volume}

\begin{longtable}[]{@{}
  >{\raggedright\arraybackslash}p{(\linewidth - 14\tabcolsep) * \real{0.1250}}
  >{\raggedright\arraybackslash}p{(\linewidth - 14\tabcolsep) * \real{0.1250}}
  >{\raggedright\arraybackslash}p{(\linewidth - 14\tabcolsep) * \real{0.1250}}
  >{\raggedright\arraybackslash}p{(\linewidth - 14\tabcolsep) * \real{0.1250}}
  >{\raggedright\arraybackslash}p{(\linewidth - 14\tabcolsep) * \real{0.1250}}
  >{\raggedright\arraybackslash}p{(\linewidth - 14\tabcolsep) * \real{0.1250}}
  >{\raggedright\arraybackslash}p{(\linewidth - 14\tabcolsep) * \real{0.1250}}
  >{\raggedright\arraybackslash}p{(\linewidth - 14\tabcolsep) * \real{0.1250}}@{}}
\toprule\noalign{}
\begin{minipage}[b]{\linewidth}\raggedright
RUNSPEC
\end{minipage} & \begin{minipage}[b]{\linewidth}\raggedright
GRID
\end{minipage} & \begin{minipage}[b]{\linewidth}\raggedright
EDIT
\end{minipage} & \begin{minipage}[b]{\linewidth}\raggedright
PROPS
\end{minipage} & \begin{minipage}[b]{\linewidth}\raggedright
REGIONS
\end{minipage} & \begin{minipage}[b]{\linewidth}\raggedright
SOLUTION
\end{minipage} & \begin{minipage}[b]{\linewidth}\raggedright
SUMMARY
\end{minipage} & \begin{minipage}[b]{\linewidth}\raggedright
SCHEDULE
\end{minipage} \\
\midrule\noalign{}
\endhead
\bottomrule\noalign{}
\endlastfoot
\end{longtable}

\subsubsection{Description}\label{description-1039}

The WBOREVOL defines a well's effective wellbore storage volume. The
primary purpose of the keyword is to enable matching of the wellbore
storage effects in well tests and the corresponding pressure response
observed in the test. Normally, as part of well test interpretation, the
pressure, permeability, effective wellbore storage, etc., are derived
from the analytical interpretation of the test. This keyword therefore
allows the engineer to enter the analytical derived effective wellbore
storage.

Wellbore storage, in terms of well testing, is an important variable
when the well is shut-in at the surface, as the well continues to flow
down-hole until the fluids obtain equilibrium. Most well tests are now
conducted using specialized tools that shut-in the well down-hole, thus
eliminating, or mostly eliminating, wellbore storage effects.

This keyword is not supported by OPM Flow but would change the results
if supported so the simulation will be stopped.

\subsection{WCALCVAL -- Define Gas Well Calorific
Value}\label{wcalcval-define-gas-well-calorific-value}

\begin{longtable}[]{@{}
  >{\raggedright\arraybackslash}p{(\linewidth - 14\tabcolsep) * \real{0.1250}}
  >{\raggedright\arraybackslash}p{(\linewidth - 14\tabcolsep) * \real{0.1250}}
  >{\raggedright\arraybackslash}p{(\linewidth - 14\tabcolsep) * \real{0.1250}}
  >{\raggedright\arraybackslash}p{(\linewidth - 14\tabcolsep) * \real{0.1250}}
  >{\raggedright\arraybackslash}p{(\linewidth - 14\tabcolsep) * \real{0.1250}}
  >{\raggedright\arraybackslash}p{(\linewidth - 14\tabcolsep) * \real{0.1250}}
  >{\raggedright\arraybackslash}p{(\linewidth - 14\tabcolsep) * \real{0.1250}}
  >{\raggedright\arraybackslash}p{(\linewidth - 14\tabcolsep) * \real{0.1250}}@{}}
\toprule\noalign{}
\begin{minipage}[b]{\linewidth}\raggedright
RUNSPEC
\end{minipage} & \begin{minipage}[b]{\linewidth}\raggedright
GRID
\end{minipage} & \begin{minipage}[b]{\linewidth}\raggedright
EDIT
\end{minipage} & \begin{minipage}[b]{\linewidth}\raggedright
PROPS
\end{minipage} & \begin{minipage}[b]{\linewidth}\raggedright
REGIONS
\end{minipage} & \begin{minipage}[b]{\linewidth}\raggedright
SOLUTION
\end{minipage} & \begin{minipage}[b]{\linewidth}\raggedright
SUMMARY
\end{minipage} & \begin{minipage}[b]{\linewidth}\raggedright
SCHEDULE
\end{minipage} \\
\midrule\noalign{}
\endhead
\bottomrule\noalign{}
\endlastfoot
\end{longtable}

\subsubsection{Description}\label{description-1040}

This keyword defines a gas well's calorific value for when the Gas
Calorific Value option has been activated by specifying a target
calorific value for a group via the GCONCAL keyword in the SCHEDULE
section. If this option is invoked then the gas calorific value must be
set either by this keyword for a well by well allocation of the
calorific value, or by using the Tracer Tracking option (activated by
the TRACER keyword in the RUNSPEC section) combined with CALTRAC keyword
in the SCHEDULE section that defines the tracer for the calorific value.

This keyword is not supported by OPM Flow but would change the results
if supported so the simulation will be stopped.

\subsection{WCONHIST -- Define Well Historical Production Rates and
Pressures}\label{wconhist-define-well-historical-production-rates-and-pressures}

\begin{longtable}[]{@{}
  >{\raggedright\arraybackslash}p{(\linewidth - 14\tabcolsep) * \real{0.1250}}
  >{\raggedright\arraybackslash}p{(\linewidth - 14\tabcolsep) * \real{0.1250}}
  >{\raggedright\arraybackslash}p{(\linewidth - 14\tabcolsep) * \real{0.1250}}
  >{\raggedright\arraybackslash}p{(\linewidth - 14\tabcolsep) * \real{0.1250}}
  >{\raggedright\arraybackslash}p{(\linewidth - 14\tabcolsep) * \real{0.1250}}
  >{\raggedright\arraybackslash}p{(\linewidth - 14\tabcolsep) * \real{0.1250}}
  >{\raggedright\arraybackslash}p{(\linewidth - 14\tabcolsep) * \real{0.1250}}
  >{\raggedright\arraybackslash}p{(\linewidth - 14\tabcolsep) * \real{0.1250}}@{}}
\toprule\noalign{}
\begin{minipage}[b]{\linewidth}\raggedright
RUNSPEC
\end{minipage} & \begin{minipage}[b]{\linewidth}\raggedright
GRID
\end{minipage} & \begin{minipage}[b]{\linewidth}\raggedright
EDIT
\end{minipage} & \begin{minipage}[b]{\linewidth}\raggedright
PROPS
\end{minipage} & \begin{minipage}[b]{\linewidth}\raggedright
REGIONS
\end{minipage} & \begin{minipage}[b]{\linewidth}\raggedright
SOLUTION
\end{minipage} & \begin{minipage}[b]{\linewidth}\raggedright
SUMMARY
\end{minipage} & \begin{minipage}[b]{\linewidth}\raggedright
SCHEDULE
\end{minipage} \\
\midrule\noalign{}
\endhead
\bottomrule\noalign{}
\endlastfoot
\end{longtable}

\subsubsection{Description}\label{description-1041}

The WCONHIST keyword defines production rates and pressures for wells
that have been declared history matching wells by the use of this
keyword. History matching wells are handled differently than ordinary
wells that use the WCONPROD keyword for controlling their production
targets and constraints. However, the wells still need to be defined
like ordinary production wells using the WELSPECS keyword in the
SCHEDULE section.

Note that although wells can be allocated to a group when they are
specified by the WELSPECS keyword, history matching wells cannot operate
under group control. Field and group reporting is still consistent for
all wells allocated to a group.

\begin{longtable}[]{@{}
  >{\raggedright\arraybackslash}p{(\linewidth - 6\tabcolsep) * \real{0.2500}}
  >{\raggedright\arraybackslash}p{(\linewidth - 6\tabcolsep) * \real{0.2500}}
  >{\raggedright\arraybackslash}p{(\linewidth - 6\tabcolsep) * \real{0.2500}}
  >{\raggedright\arraybackslash}p{(\linewidth - 6\tabcolsep) * \real{0.2500}}@{}}
\toprule\noalign{}
\begin{minipage}[b]{\linewidth}\raggedright
No.
\end{minipage} & \begin{minipage}[b]{\linewidth}\raggedright
Name
\end{minipage} & \begin{minipage}[b]{\linewidth}\raggedright
Description
\end{minipage} & \begin{minipage}[b]{\linewidth}\raggedright
Default
\end{minipage} \\
\midrule\noalign{}
\endhead
\bottomrule\noalign{}
\endlastfoot
Field & Metric & Laboratory & \\
1 & WELNAME & A character string of up to eight characters in length
that defines the well name for which the wells observed production rates
and pressures are being defined. Note that the well name (WELNAME) must
have been declared previously using the WELSPECS keyword in the SCHEDULE
section, otherwise an error may occur. & None \\
2 & STATUS & A defined character string that declares the status of the
well. STATUS should be set to one of the following character strings:
Note a well's STATUS should always be set either STOP or SHUT if the
well's production is to be set to zero. Just setting a well's production
rate to zero means that the well is open to flow with a zero rate. &
OPEN \\
3 & TARGET & A defined character string that sets the observed target
production phase for the well, all the other phases are calculated
unconstrained and used for reporting only. The simulator will attempt to
meet the TARGET based on the phase rate stated in items (4) to (6) and
(10) on this keyword. TARGET should be set to one of the following
character strings: & None \\
4 & ORAT & A real positive value that defines the observed surface oil
production rate target or constraint. This value may be specified using
a User Defined Argument (UDA). & Defined \\
stb/d 0.0 & sm3/day 0.0 & scc/hour 0.0 & \\
5 & WRAT & A real positive value that defines the observed surface water
production rate target or constraint. This value may be specified using
a User Defined Argument (UDA). & Defined \\
stb/d 0.0 & sm3/day 0.0 & scc/hour 0.0 & \\
6 & GRAT & A real positive value that defines the observed surface gas
production rate target or constraint This value may be specified using a
User Defined Argument (UDA). & Defined \\
Mscf/d 0.0 & sm3/day 0.0 & scc/hour 0.0 & \\
7 & VFPTAB & A positive integer greater than or equal to zero that
defines the vertical lift performance tables to be used for calculating
the tubing head pressure for the well. If a non-zero value is entered
then the vertical lift performance tables must be entered via the
VFPPROD keyword in the SCHEDULE section and allocated to the well via
this item. The default value of zero implies no vertical lift
performance table initially. If this value is then reset to be greater
than zero then the table will be used to calculate the well's tubing
head pressure. Subsequently, the default is to use the previously
declared table number. & 0 \\
8 & ALQ-WELL & A real positive value that defines the artificial lift
quantity to be used in conjunction with the VFPPROD assigned to the well
via that keyword's VFPTAB variable. This value may be specified using a
User Defined Argument (UDA). VFPTAB vertical lift performance table and
the artificial lift quantity ALQ-WELL are used with the well fluid rates
to calculate the well's tubing head pressures values from the
bottom-hole pressure. Note that the units for ALQ-WELL is dependent on
the associated variable on the VFPPROD keyword. & None \\
9 & THP & A real positive value that defines the observed tubing head
pressure. This value may be specified using a User Defined Argument
(UDA). This parameter is only used for comparing the actual tubing head
pressure given here with those calculated by the simulator, that is
history matching wells can only be controlled by either the surface
injection rate or their bottom-hole pressure. & Defined \\
psia 0.0 & barsa 0.0 & atma 0.0 & \\
10 & BHP & A real positive value that defines the observed bottom-hole
pressure. This value may be specified using a User Defined Argument
(UDA). & Defined \\
psia 0.0 & barsa 0.0 & atma 0.0 & \\
11 & WGRA & A real positive value that defines the observed wet gas rate
in the commercial compositional simulator. The parameter is ignored by
OPM Flow and should be defaulted or set to the default value of zero. &
0.0 \\
12 & NGL & A real positive value that defines the observed Natural Gas
Liquid (``NGL'') rate in the commercial compositional simulator. The
parameter is ignored by OPM Flow and should be defaulted or set to the
default value of zero. & 0.0 \\
Notes: & & & \\
\end{longtable}

\emph{Table 12.77: WCONHIST Keyword Description}

See also the WHISTCTL keyword that can be used to reset the TARGET
phase, the GCONPROD and GCONINJE keywords to define a group's production
and injection targets and constraints, and the WCONPROD keyword to
define a production well's targets and constraints. All the
aforementioned keywords are described in the SCHEDULE section.

\begin{tcolorbox}[enhanced jigsaw, title=\textcolor{quarto-callout-note-color}{\faInfo}\hspace{0.5em}{Note}, coltitle=black, breakable, toprule=.15mm, colbacktitle=quarto-callout-note-color!10!white, leftrule=.75mm, opacitybacktitle=0.6, colback=white, bottomtitle=1mm, toptitle=1mm, titlerule=0mm, left=2mm, colframe=quarto-callout-note-color-frame, arc=.35mm, rightrule=.15mm, bottomrule=.15mm, opacityback=0]

One can use TARGET set to RESV in the initial history matching runs to
get a ``reasonable'' pressure match, this ensures that the total
reservoir withdrawals are correct, although the individual phase
withdrawals will not match. Once a reasonable pressure match is achieved
for the reservoir then one can reset TARGET to the sales phase, OIL or
GAS, and continue with the matching of all the phases. In oil reservoirs
some engineers prefer to use LIQ rather than OIL as the TARGET phase,
although one should consider that as the water phase has no commercial
value, the measurement accuracy is significantly less than the oil sales
phase.

\end{tcolorbox}

History matching wells are converted to ordinary wells by restating a
well's control mode using either the WCONPROD or WELTARG keywords in the
SCHEDULE section.

\subsubsection{Examples}\label{examples-90}

The following example below shows the observed production rates for the
OP01 oil producer for the first quarter of 2000.

\begin{verbatim}
-- ------------------------------------------------------------------------------
-- 01 JAN 2000 START OF SCHEDULE SECTION
-- ------------------------------------------------------------------------------
--
--       WELL HISTORICAL PRODUCTION CONTROLS
--
-- WELL  OPEN/  CNTL   OIL    WAT    GAS    VFP    VFP   THP   BHP
-- NAME  SHUT   MODE   RATE   RATE   RATE   TABLE  ALFQ  PRES  PRES
WCONHIST                                                                                                                          OP01     OPEN   ORAT  15.5E3  100.0  1550   10      1*   900.0 1*     /
/                                                                               DATES
01 FEB 2000 /
/
--
--       WELL HISTORICAL PRODUCTION CONTROLS
--
-- WELL  OPEN/  CNTL   OIL    WAT    GAS    VFP    VFP   THP   BHP
-- NAME  SHUT   MODE   RATE   RATE   RATE   TABLE  ALFQ  PRES  PRES
WCONHIST                                                                                                                          OP01     OPEN   ORAT  15.2E3  150.0  1520   1*      1*   875.0 3250.0 /
/                                                                               DATES
01 MAR 2000 /
/
--
--       WELL HISTORICAL PRODUCTION CONTROLS
--
-- WELL  OPEN/  CNTL   OIL    WAT    GAS    VFP    VFP   THP   BHP
-- NAME  SHUT   MODE   RATE   RATE   RATE   TABLE  ALFQ  PRES  PRES
WCONHIST                                                                                                                          OP01     OPEN   ORAT  15.0E3  200.0  1500   1*      1*   850.0 1*     /
/
\end{verbatim}

From January 1, 2000 well OP01 is open and is on oil rate control, and
produces 15,500 stb/d oil, with the observed rates of 100 stb/d of water
and 15.5 MMscf/d of gas. The well uses VFPPROD vertical lift table
number 10 so that OPM Flow can calculate the tubing head pressures based
on the fluids produced and the calculated pressures in the simulator.

The next example illustrates how to convert OP01 from a history match
well to a normal production well at the start for the forecast run at
August 1, 2017 using the WELTARG keyword.

\begin{verbatim}
DATES
01 AUG 2017 /
/
--
--       WELL PRODUCTION AND INJECTION TARGETS
--
--  WELL WELL   TARGET
--  NAME TARG   VALUE
WELTARG
OP01     THP    1*                                                  /
/
\end{verbatim}

Here by defaulting the bottom-hole pressure via 1* OPM Flow
automatically applies the last bottom-hole pressure from the previous
time step as the ``constraining phase'' together with the last
historical rates as constraints. This ensures a smooth transition
between history and prediction without having to resort to unreasonable
changes to the model. This option is currently not implemented in OPM
Flow but is expected to be incorporated in a future release.

\subsection{WCONINJ -- Well Injection Targets and
Constraints}\label{wconinj-well-injection-targets-and-constraints}

\begin{longtable}[]{@{}
  >{\raggedright\arraybackslash}p{(\linewidth - 14\tabcolsep) * \real{0.1250}}
  >{\raggedright\arraybackslash}p{(\linewidth - 14\tabcolsep) * \real{0.1250}}
  >{\raggedright\arraybackslash}p{(\linewidth - 14\tabcolsep) * \real{0.1250}}
  >{\raggedright\arraybackslash}p{(\linewidth - 14\tabcolsep) * \real{0.1250}}
  >{\raggedright\arraybackslash}p{(\linewidth - 14\tabcolsep) * \real{0.1250}}
  >{\raggedright\arraybackslash}p{(\linewidth - 14\tabcolsep) * \real{0.1250}}
  >{\raggedright\arraybackslash}p{(\linewidth - 14\tabcolsep) * \real{0.1250}}
  >{\raggedright\arraybackslash}p{(\linewidth - 14\tabcolsep) * \real{0.1250}}@{}}
\toprule\noalign{}
\begin{minipage}[b]{\linewidth}\raggedright
RUNSPEC
\end{minipage} & \begin{minipage}[b]{\linewidth}\raggedright
GRID
\end{minipage} & \begin{minipage}[b]{\linewidth}\raggedright
EDIT
\end{minipage} & \begin{minipage}[b]{\linewidth}\raggedright
PROPS
\end{minipage} & \begin{minipage}[b]{\linewidth}\raggedright
REGIONS
\end{minipage} & \begin{minipage}[b]{\linewidth}\raggedright
SOLUTION
\end{minipage} & \begin{minipage}[b]{\linewidth}\raggedright
SUMMARY
\end{minipage} & \begin{minipage}[b]{\linewidth}\raggedright
SCHEDULE
\end{minipage} \\
\midrule\noalign{}
\endhead
\bottomrule\noalign{}
\endlastfoot
\end{longtable}

\subsubsection{Description}\label{description-1042}

The WCONINJ is a legacy keyword that is no longer used in the commercial
simulator and is not supported by OPM Flow. Instead well injection
targets and constraints should be defined using the WCONINJE keyword in
the SCHEDULE section.

This keyword is not supported by OPM Flow but would change the results
if supported so the simulation will be stopped.

\subsection{WCONINJE -- Well Injection Targets and
Constraints}\label{wconinje-well-injection-targets-and-constraints}

\begin{longtable}[]{@{}
  >{\raggedright\arraybackslash}p{(\linewidth - 14\tabcolsep) * \real{0.1250}}
  >{\raggedright\arraybackslash}p{(\linewidth - 14\tabcolsep) * \real{0.1250}}
  >{\raggedright\arraybackslash}p{(\linewidth - 14\tabcolsep) * \real{0.1250}}
  >{\raggedright\arraybackslash}p{(\linewidth - 14\tabcolsep) * \real{0.1250}}
  >{\raggedright\arraybackslash}p{(\linewidth - 14\tabcolsep) * \real{0.1250}}
  >{\raggedright\arraybackslash}p{(\linewidth - 14\tabcolsep) * \real{0.1250}}
  >{\raggedright\arraybackslash}p{(\linewidth - 14\tabcolsep) * \real{0.1250}}
  >{\raggedright\arraybackslash}p{(\linewidth - 14\tabcolsep) * \real{0.1250}}@{}}
\toprule\noalign{}
\begin{minipage}[b]{\linewidth}\raggedright
RUNSPEC
\end{minipage} & \begin{minipage}[b]{\linewidth}\raggedright
GRID
\end{minipage} & \begin{minipage}[b]{\linewidth}\raggedright
EDIT
\end{minipage} & \begin{minipage}[b]{\linewidth}\raggedright
PROPS
\end{minipage} & \begin{minipage}[b]{\linewidth}\raggedright
REGIONS
\end{minipage} & \begin{minipage}[b]{\linewidth}\raggedright
SOLUTION
\end{minipage} & \begin{minipage}[b]{\linewidth}\raggedright
SUMMARY
\end{minipage} & \begin{minipage}[b]{\linewidth}\raggedright
SCHEDULE
\end{minipage} \\
\midrule\noalign{}
\endhead
\bottomrule\noalign{}
\endlastfoot
\end{longtable}

\subsubsection{Description}\label{description-1043}

The WCONINJE keyword defines injection targets and constraints for wells
that have previously been defined by the WELSPECS keyword in the
SCHEDULE section. Note that wells can be allocated to a group when they
are specified by the WELSPECS keyword. Wells defined to be under group
control will have their injection rates controlled by the group to which
they belong, in addition to any well constraints defined for the wells
using this keyword.

\begin{longtable}[]{@{}
  >{\raggedright\arraybackslash}p{(\linewidth - 6\tabcolsep) * \real{0.2500}}
  >{\raggedright\arraybackslash}p{(\linewidth - 6\tabcolsep) * \real{0.2500}}
  >{\raggedright\arraybackslash}p{(\linewidth - 6\tabcolsep) * \real{0.2500}}
  >{\raggedright\arraybackslash}p{(\linewidth - 6\tabcolsep) * \real{0.2500}}@{}}
\toprule\noalign{}
\begin{minipage}[b]{\linewidth}\raggedright
No.
\end{minipage} & \begin{minipage}[b]{\linewidth}\raggedright
Name
\end{minipage} & \begin{minipage}[b]{\linewidth}\raggedright
Description
\end{minipage} & \begin{minipage}[b]{\linewidth}\raggedright
Default
\end{minipage} \\
\midrule\noalign{}
\endhead
\bottomrule\noalign{}
\endlastfoot
Field & Metric & Laboratory & \\
1 & WELNAME & A character string of up to eight characters in length
that defines the well name for which the well injection targets and
constraints data are being defined. Note that the well name (WELNAME)
must have been declared previously using the WELSPECS keyword in the
SCHEDULE section, otherwise an error may occur. & None \\
2 & TYPE & A defined character string that defines the type of injection
well. TYPE should be set to one of the following character strings: &
None \\
3 & STATUS & A defined character string that declares the status of the
well. STATUS should be set to one of the following character strings:
Note a well's STATUS should always be set either STOP or SHUT if the
well's injection is to be set to zero. Just setting a well's injection
rate to zero means that the well is open for injection with a zero rate,
this will cause numerical issues especially for wells under THP control.
& OPEN \\
4 & TARGET & A defined character string that sets the target injection
control mode for the well. TARGET should be set to one of the following
character strings: & None \\
5 & RATE & A real positive value that defines the maximum surface
injection rate target or constraint. This value may be specified using a
User Defined Argument (UDA). & None \\
Liquid stb/d Gas Mscf/d & Liquid sm3/day Gas sm3/day & Liquid scc/hour
Gas scc/hour & \\
6 & RESV & A real positive value that defines the maximum reservoir
volume injection rate target or constraint. This value may be specified
using a User Defined Argument (UDA). & None \\
rb/d & rm3/day & rcc/hour & \\
7 & BHP & A real positive value that defines the maximum bottom-hole
pressure target or constraint. This value may be specified using a User
Defined Argument (UDA). Note the default value basically means unlimited
injection or no constraint and should therefore be avoided as the BHP
will result in unrealistic well potentials as well as optimistic
injection forecasts for the well. & Defined \\
psia 10,0000 & barsa 6,895 & atma 6,803 & \\
8 & THP & A real positive value that defines the maximum tubing head
pressure target or constraint. This value may be specified using a User
Defined Argument (UDA). & None \\
psia & barsa & atma & \\
9 & VFPTAB & A positive integer greater than or equal to zero that
defines the vertical lift performance tables to be used for calculating
the tubing head pressure for the well. If a non-zero value is entered
then the vertical lift performance tables must be entered via the VFPINJ
keyword in the SCHEDULE section and allocated to the well via this item.
The default value of zero implies no vertical lift performance tables
and in this case TARGET cannot be set to THP and in addition item (10)
should be defaulted or set to zero. & 0 \\
10 & RSRVINJ & The dissolved gas-oil ratio in the injected oil, or the
vaporized oil-gas ratio in the injected gas. & 0.0 \\
Gas Injection: stb/Mscf Oil Injection: Mscf/stb & Gas Injection: sm3/sm3
Oil Injection: sm3/sm3 & Gas Injection: scc/scc Oil Injection: scc/scc
& \\
11 & RSSTEAM & Thermal/Temperature gas-steam ratio for steam-gas
injectors. The parameter is ignored by OPM Flow and should be defaulted
or set to the default value of zero. & 0.0 \\
12 & OILFRAC & Surface oil fraction in a multi-phase injector. The
parameter is ignored by OPM Flow and should be defaulted or set to the
default value of zero. & 0.0 \\
13 & WATFRAC & Surface water fraction in a multi-phase injector. The
parameter is ignored by OPM Flow and should be defaulted or set to the
default value of zero. & 0.0 \\
14 & GASFRAC & Surface gas fraction in a multi-phase injector. The
parameter is ignored by OPM Flow and should be defaulted or set to the
default value of zero. & 0.0 \\
15 & OILSTEAM & Surface oil volume to steam volume ratio in a steam-oil
injector. The parameter is ignored by OPM Flow and should be defaulted
or set to the default value of zero. & 0.0 \\
Notes: & & & \\
\end{longtable}

\emph{Table 12.78: WCONINJE Keyword Description}

See also the GCONPROD and the GCONINJE keywords to define a group's
production and injection targets and constraints, and the WCONPROD
keyword to define a production well's targets and constraints. All the
aforementioned keywords are described in the SCHEDULE section.

\subsubsection{Example}\label{example-560}

The following example defines the injection targets and constraints for
one gas injection well and one water injection well as follows:

\begin{verbatim}
--
--       WELL INJECTION CONTROLS
--
-- WELL  FLUID  OPEN/  CNTL  SURF   RESV   BHP   THP   VFP
-- NAME  TYPE   SHUT   MODE  RATE   RATE   PRSES PRES  TABLE
WCONINJE
GI01     GAS    OPEN   GRUP  50E3   1*     1*     1*    1*  /
WI01     WAT    OPEN   RATE  25E3   1*     5000.  1*    1*  /
/
\end{verbatim}

Well GI01 is a gas injection well directly under group control
constrained by a maximum surface gas injection rate of 50 MMscf/d and
well WI01 is an open water injection well with a surface water injection
rate target of 25,000 stb/d, subject to a maximum bottom-hole pressure
constraint 5,000 psia.

\subsection{WCONINJH -- Well Historical Observed Injection Rates and
Pressures}\label{wconinjh-well-historical-observed-injection-rates-and-pressures}

\begin{longtable}[]{@{}
  >{\raggedright\arraybackslash}p{(\linewidth - 14\tabcolsep) * \real{0.1250}}
  >{\raggedright\arraybackslash}p{(\linewidth - 14\tabcolsep) * \real{0.1250}}
  >{\raggedright\arraybackslash}p{(\linewidth - 14\tabcolsep) * \real{0.1250}}
  >{\raggedright\arraybackslash}p{(\linewidth - 14\tabcolsep) * \real{0.1250}}
  >{\raggedright\arraybackslash}p{(\linewidth - 14\tabcolsep) * \real{0.1250}}
  >{\raggedright\arraybackslash}p{(\linewidth - 14\tabcolsep) * \real{0.1250}}
  >{\raggedright\arraybackslash}p{(\linewidth - 14\tabcolsep) * \real{0.1250}}
  >{\raggedright\arraybackslash}p{(\linewidth - 14\tabcolsep) * \real{0.1250}}@{}}
\toprule\noalign{}
\begin{minipage}[b]{\linewidth}\raggedright
RUNSPEC
\end{minipage} & \begin{minipage}[b]{\linewidth}\raggedright
GRID
\end{minipage} & \begin{minipage}[b]{\linewidth}\raggedright
EDIT
\end{minipage} & \begin{minipage}[b]{\linewidth}\raggedright
PROPS
\end{minipage} & \begin{minipage}[b]{\linewidth}\raggedright
REGIONS
\end{minipage} & \begin{minipage}[b]{\linewidth}\raggedright
SOLUTION
\end{minipage} & \begin{minipage}[b]{\linewidth}\raggedright
SUMMARY
\end{minipage} & \begin{minipage}[b]{\linewidth}\raggedright
SCHEDULE
\end{minipage} \\
\midrule\noalign{}
\endhead
\bottomrule\noalign{}
\endlastfoot
\end{longtable}

\subsubsection{Description}\label{description-1044}

The WCONINJH keyword defines injection rates and pressures for wells
that have been declared history matching wells by the use of this
keyword. History matching wells are handled differently then ordinary
wells that use the WCONINJE keyword for controlling their injection
targets and constraints. However, the wells still need to be defined
like ordinary injection wells using the WELSPECS keyword in the SCHEDULE
section.

Note that although wells can be allocated to a group when they are
specified by the WELSPECS keyword, history matching wells cannot operate
under group control. Field and group reporting is still consistent for
all wells allocated to a group, but history matching wells cannot be
under group control.

\begin{longtable}[]{@{}
  >{\raggedright\arraybackslash}p{(\linewidth - 6\tabcolsep) * \real{0.2500}}
  >{\raggedright\arraybackslash}p{(\linewidth - 6\tabcolsep) * \real{0.2500}}
  >{\raggedright\arraybackslash}p{(\linewidth - 6\tabcolsep) * \real{0.2500}}
  >{\raggedright\arraybackslash}p{(\linewidth - 6\tabcolsep) * \real{0.2500}}@{}}
\toprule\noalign{}
\begin{minipage}[b]{\linewidth}\raggedright
No.
\end{minipage} & \begin{minipage}[b]{\linewidth}\raggedright
Name
\end{minipage} & \begin{minipage}[b]{\linewidth}\raggedright
Description
\end{minipage} & \begin{minipage}[b]{\linewidth}\raggedright
Default
\end{minipage} \\
\midrule\noalign{}
\endhead
\bottomrule\noalign{}
\endlastfoot
Field & Metric & Laboratory & \\
1 & WELNAME & A character string of up to eight characters in length
that defines the well name for which the wells observed injection rates
and pressures are being defined. Note that the well name (WELNAME) must
have been declared previously using the WELSPECS keyword in the SCHEDULE
section, otherwise an error may occur. & None \\
2 & TYPE & A defined character string that defines the type of injection
well. TYPE should be set to one of the following character strings: &
None \\
3 & STATUS & A defined character string that declares the status of the
well. STATUS should be set to one of the following character strings:
Note a well's STATUS should always be set either STOP or SHUT if the
well's injection is to be set to zero. Just setting a well's injection
rate to zero means that the well is open to flow with a zero injection
rate, this may cause numerical issues. & OPEN \\
4 & RATE & A real positive value that defines the observed surface
injection rate. & 0.0 \\
Liquid stb/d Gas Mscf/d & Liquid sm3/day Gas sm3/day & Liquid scc/hour
Gas scc/hour & \\
5 & BHP & A real positive value that defines the observed bottom-hole
pressure. & 0.0 \\
psia & barsa & atma & \\
6 & THP & A real positive value that defines the observed tubing head
pressure. This parameter is only used for comparing the actual tubing
head pressure given here with those calculated by the simulator, that is
history matching wells can only be controlled by either the surface
injection rate or their bottom-hole pressure. & None \\
psia & barsa & atma & \\
7 & VFPTAB & A positive integer greater than or equal to zero that
defines the vertical lift performance tables to be used for calculating
the tubing head pressure for the well. If a non-zero value is entered
then the vertical lift performance tables must be entered via the VFPINJ
keyword in the SCHEDULE section and allocated to the well via this item.
The default value of zero implies no vertical lift performance table
initially. If this value is then reset to be greater than zero then the
table will be used to calculate the well's tubing head pressure.
Subsequently, the default is to use the previously declared table
number. & 0 \\
8 & RSRVINJ & Dissolved gas fraction in injected oil or vaporized oil
fraction in injected gas. The parameter is ignored by OPM Flow and
should be defaulted or set to the default value of zero. & 0.0 \\
9 & OILFRAC & Surface oil fraction in a multi-phase injector. The
parameter is ignored by OPM Flow and should be defaulted or set to the
default value of zero. & 0.0 \\
10 & WATFRAC & Surface water fraction in a multi-phase injector. The
parameter is ignored by OPM Flow and should be defaulted or set to the
default value of zero. & 0.0 \\
11 & GASFRAC & Surface gas fraction in a multi-phase injector. The
parameter is ignored by OPM Flow and should be defaulted or set to the
default value of zero. & 0.0 \\
12 & TARGET & A defined character string that sets the target injection
control mode for the well. TARGET should be set to one of the following
character strings: & RATE \\
Notes: & & & \\
\end{longtable}

\emph{Table 12.79: WCONINJH Keyword Description}

This keyword should be repeated at various time steps to fully describe
the historical injection performance of the wells. For example, as most
production and injection data is reconciled on a monthly basis, then
monthly time steps covering the injection history of the wells should be
used with WCONINJH keyword entered on a monthly basis.

History matching well are converted to ordinary wells by restating a
well's control mode using either the WCONINJE or WELTARG keywords in the
SCHEDULE section.

\subsubsection{Example}\label{example-561}

The following example below shows the observed gas rates for the GI01
gas injector for the first quarter of 2000.

\begin{verbatim}
-- ------------------------------------------------------------------------------
-- 01 JAN 2000 START OF SCHEDULE SECTION
-- ------------------------------------------------------------------------------
--
--       WELL HISTORICAL INJECTION CONTROLS
--
-- WELL  FLUID  OPEN/  SURF   BHP   THP   VFP    NOT   CNTL
-- NAME  TYPE   SHUT   RATE   PRES  PRES  TABLE  USED  MODE
WCONINJH
GI01     GAS    OPEN  15.5E3  1*    5462   12    4*    1*  /
/                                                                               DATES
01 FEB 2000 /
/
--
--       WELL HISTORICAL INJECTION CONTROLS
--
-- WELL  FLUID  OPEN/  SURF   BHP   THP   VFP    NOT   CNTL
-- NAME  TYPE   SHUT   RATE   PRES  PRES  TABLE  USED  MODE
WCONINJH
GI01     GAS    OPEN  15.9E3  1*    5468   1*    4*    1*  /
/                                                                               DATES
01 MAR 2000 /
/

--
--       WELL HISTORICAL INJECTION CONTROLS
--
-- WELL  FLUID  OPEN/  SURF   BHP   THP   VFP    NOT   CNTL
-- NAME  TYPE   SHUT   RATE   PRES  PRES  TABLE  USED  MODE
WCONINJH
GI01     GAS    OPEN  17.2E3  1*    5489   1*    4*    1*  /
/
\end{verbatim}

Well GI01is declared as a gas injection well under gas rate control as
TARGET variable is defaulted to rate control by using 1* (the last entry
on the record). In addition, the well uses vertical lift table VFPINJ
number 12 (as shown at January 1, 2000) to calculate the tubing head
pressures for the well. Note that it is not necessary to declare the
VFPINJ table number if it remains the same for subsequent time steps and
thus the default 1* is used to indicate the last entry should be used.

\subsection{WCONINJP -- Define Well Injection Targets and Constraints
for Pattern Flood
Wells}\label{wconinjp-define-well-injection-targets-and-constraints-for-pattern-flood-wells}

\begin{longtable}[]{@{}
  >{\raggedright\arraybackslash}p{(\linewidth - 14\tabcolsep) * \real{0.1250}}
  >{\raggedright\arraybackslash}p{(\linewidth - 14\tabcolsep) * \real{0.1250}}
  >{\raggedright\arraybackslash}p{(\linewidth - 14\tabcolsep) * \real{0.1250}}
  >{\raggedright\arraybackslash}p{(\linewidth - 14\tabcolsep) * \real{0.1250}}
  >{\raggedright\arraybackslash}p{(\linewidth - 14\tabcolsep) * \real{0.1250}}
  >{\raggedright\arraybackslash}p{(\linewidth - 14\tabcolsep) * \real{0.1250}}
  >{\raggedright\arraybackslash}p{(\linewidth - 14\tabcolsep) * \real{0.1250}}
  >{\raggedright\arraybackslash}p{(\linewidth - 14\tabcolsep) * \real{0.1250}}@{}}
\toprule\noalign{}
\begin{minipage}[b]{\linewidth}\raggedright
RUNSPEC
\end{minipage} & \begin{minipage}[b]{\linewidth}\raggedright
GRID
\end{minipage} & \begin{minipage}[b]{\linewidth}\raggedright
EDIT
\end{minipage} & \begin{minipage}[b]{\linewidth}\raggedright
PROPS
\end{minipage} & \begin{minipage}[b]{\linewidth}\raggedright
REGIONS
\end{minipage} & \begin{minipage}[b]{\linewidth}\raggedright
SOLUTION
\end{minipage} & \begin{minipage}[b]{\linewidth}\raggedright
SUMMARY
\end{minipage} & \begin{minipage}[b]{\linewidth}\raggedright
SCHEDULE
\end{minipage} \\
\midrule\noalign{}
\endhead
\bottomrule\noalign{}
\endlastfoot
\end{longtable}

\subsubsection{Description}\label{description-1045}

This keyword, WCONINJP, defines well injection targets and constraints
for pattern flood wells. The keyword is similar to the WCONINJE keyword
in the SCHEDULE section except that the injection control is applied to
a group of wells defined by the first record of this keyword, combined
with a second record that defines the wells in the pattern and their
contribution to the pattern.

This keyword is not supported by OPM Flow but would change the results
if supported so the simulation will be stopped.

\subsection{WCONPROD -- Define Well Production Targets and
Constraints}\label{wconprod-define-well-production-targets-and-constraints}

\begin{longtable}[]{@{}
  >{\raggedright\arraybackslash}p{(\linewidth - 14\tabcolsep) * \real{0.1250}}
  >{\raggedright\arraybackslash}p{(\linewidth - 14\tabcolsep) * \real{0.1250}}
  >{\raggedright\arraybackslash}p{(\linewidth - 14\tabcolsep) * \real{0.1250}}
  >{\raggedright\arraybackslash}p{(\linewidth - 14\tabcolsep) * \real{0.1250}}
  >{\raggedright\arraybackslash}p{(\linewidth - 14\tabcolsep) * \real{0.1250}}
  >{\raggedright\arraybackslash}p{(\linewidth - 14\tabcolsep) * \real{0.1250}}
  >{\raggedright\arraybackslash}p{(\linewidth - 14\tabcolsep) * \real{0.1250}}
  >{\raggedright\arraybackslash}p{(\linewidth - 14\tabcolsep) * \real{0.1250}}@{}}
\toprule\noalign{}
\begin{minipage}[b]{\linewidth}\raggedright
RUNSPEC
\end{minipage} & \begin{minipage}[b]{\linewidth}\raggedright
GRID
\end{minipage} & \begin{minipage}[b]{\linewidth}\raggedright
EDIT
\end{minipage} & \begin{minipage}[b]{\linewidth}\raggedright
PROPS
\end{minipage} & \begin{minipage}[b]{\linewidth}\raggedright
REGIONS
\end{minipage} & \begin{minipage}[b]{\linewidth}\raggedright
SOLUTION
\end{minipage} & \begin{minipage}[b]{\linewidth}\raggedright
SUMMARY
\end{minipage} & \begin{minipage}[b]{\linewidth}\raggedright
SCHEDULE
\end{minipage} \\
\midrule\noalign{}
\endhead
\bottomrule\noalign{}
\endlastfoot
\end{longtable}

\subsubsection{Description}\label{description-1046}

The WCONPROD keyword defines production targets and constraints for
wells that have previously been defined by the WELSPECS keyword in the
SCHEDULE section. Note that wells can be allocated to a group when they
are specified by the WELSPECS keyword. Wells defined to be under group
control will have their production rates controlled by the group to
which they belong, in addition to any well constraints defined for the
wells using this keyword.

\begin{longtable}[]{@{}
  >{\raggedright\arraybackslash}p{(\linewidth - 6\tabcolsep) * \real{0.2500}}
  >{\raggedright\arraybackslash}p{(\linewidth - 6\tabcolsep) * \real{0.2500}}
  >{\raggedright\arraybackslash}p{(\linewidth - 6\tabcolsep) * \real{0.2500}}
  >{\raggedright\arraybackslash}p{(\linewidth - 6\tabcolsep) * \real{0.2500}}@{}}
\toprule\noalign{}
\begin{minipage}[b]{\linewidth}\raggedright
No.
\end{minipage} & \begin{minipage}[b]{\linewidth}\raggedright
Name
\end{minipage} & \begin{minipage}[b]{\linewidth}\raggedright
Description
\end{minipage} & \begin{minipage}[b]{\linewidth}\raggedright
Default
\end{minipage} \\
\midrule\noalign{}
\endhead
\bottomrule\noalign{}
\endlastfoot
Field & Metric & Laboratory & \\
1 & WELNAME & A character string of up to eight characters in length
that defines the well name for which the well production targets and
constraints data are being defined. Note that the well name (WELNAME)
must have been declared previously using the WELSPECS keyword in the
SCHEDULE section, otherwise an error may occur. & None \\
2 & STATUS & A defined character string that declares the status of the
well. STATUS should be set to one of the following character strings:
Note a well's STATUS should always be set either STOP or SHUT if the
well's production is to be set to zero. Just setting a well's production
rate to zero means that the well is open to flow with a zero rate, this
will cause numerical issues especially for wells under THP control. &
OPEN \\
3 & TARGET & A defined character string that sets the target production
phase for the well, all the other phases will therefore act as
constraints. The simulator will attempt to meet the TARGET based on the
phase rate stated in items (4) to (10) on this keyword. TARGET should be
set to one of the following character strings: Here the default tokens
of 1* or '\,' may be used, and both default tokens behave in the same
manner; however, the actual default value used when TARGET is defaulted
is dependent on the data entered on this keyword as described below:
Note the default value of one atmosphere should be avoided as the BHP
will result in unrealistic well potentials as well as optimistic
production forecasts for the well. & 1* \\
4 & ORAT & A real positive value that defines the maximum surface oil
production rate target or constraint. This value may be specified using
a User Defined Argument (UDA). & None \\
stb/day & sm3/day & scc/hour & \\
5 & WRAT & A real positive value that defines the maximum surface water
production rate target or constraint. This value may be specified using
a User Defined Argument (UDA). & None \\
stb/day & sm3/day & scc/hour & \\
6 & GRAT & A real positive value that defines the maximum surface gas
production rate target or constraint This value may be specified using a
User Defined Argument (UDA). & None \\
Mscf/day & sm3/day & scc/hour & \\
7 & LRAT & A real positive value that defines the maximum surface liquid
(oil plus water) production rate target or constraint. This value may be
specified using a User Defined Argument (UDA). & None \\
stb/day & sm3/day & scc/hour & \\
8 & RESV & A real positive value that defines the maximum reservoir
volume production rate target or constraint. This value may be specified
using a User Defined Argument (UDA). & None \\
rtb/day & rm3/day & rcc/hour & \\
9 & BHP & A real positive value that defines the minimum bottom-hole
pressure target or constraint. This value may be specified using a User
Defined Argument (UDA). Note the default value of one atmosphere should
be avoided as the BHP will result in unrealistic well potentials as well
as optimistic production forecasts for the well. & Defined \\
psia 14.70 & barsa 1.01325. & atma 1.0 & \\
10 & THP & A real positive value that defines the minimum tubing head
pressure target or constraint. This value may be specified using a User
Defined Argument (UDA). Note the default value of zero should be avoided
if the well's control TARGET has been set to THP, as this will result in
optimistic production forecasts for a well, since a well must flow
against a back pressure imposed by the surface facilities. & Defined \\
psia 0.0 & barsa 0.0 & atma 0.0 & \\
11 & VFPTAB & A positive integer greater than or equal to zero that
defines the vertical lift performance tables to be used for calculating
the tubing head pressure for the well. If a non-zero value is entered
then the vertical lift performance tables must be entered via the
VFPPROD keyword in the SCHEDULE section and allocated to the well via
this item. The default value of zero implies no vertical lift
performance tables and in this case TARGET cannot be set to THP and in
addition item (10) should be defaulted or set to zero. & 0 \\
12 & ALQ-WELL & A real positive value that defines the artificial lift
quantity to be used in conjunction with the VFPPROD assigned to the well
via the VFPTAB variable. This value may be specified using a User
Defined Argument (UDA). VFPTAB vertical lift performance table and the
artificial lift quantity ALQ-WELL are used with the well fluid rates to
calculate the well's tubing head pressures values from the bottom-hole
pressure. Note that the units for ALQ-WELL are dependent on the
associated variable on the VFPPROD keyword. & 0.0 \\
13 & WGASRATE & Wet gas production rate used in the commercial
compositional simulator. Not used and should be defaulted with 1*. &
1* \\
14 & MOLARATE & Total molar rate used in the commercial compositional
simulator. Not used and should be defaulted with 1*. & 1* \\
15 & STEAMRAT & Thermal/Temperature steam rate (Cold Water Equivalent)
for steam producers used in the commercial compositional simulator. Not
used and should be defaulted with 1*. & 1* \\
16 & DELTAP & Thermal/Temperature delta pressure offset for steam
producers used in the commercial compositional simulator. Not used and
should be defaulted with 1*. & 1* \\
17 & DELTAT & Thermal/Temperature delta temperature offset for steam
producer used in the commercial compositional simulator. Not used and
should be defaulted with 1*. & 1* \\
18 & CALRATE & Calorific production rate used in the commercial
compositional simulator. Not used and should be defaulted with 1*. &
1* \\
19 & COMBPROC & Linearly combined procedure for when exceeding COMBRATE,
used in the commercial compositional simulator. Not used and should be
defaulted with 1*. & 1* \\
20 & NGL & A real positive value that defines the observed Natural Gas
Liquid (``NGL'') rate in the commercial compositional simulator. Not
used and should be defaulted with 1*. & 1* \\
Notes: & & & \\
\end{longtable}

\emph{Table 12.80: WCONPROD Keyword Description}

See also the GCONPROD and GCONINJE keywords to define a group's
production and injection targets and constraints, and the WCONINJE
keyword to define an injection well's targets and constraints. All the
aforementioned keywords are described in the SCHEDULE section.

\subsubsection{Example}\label{example-562}

The following example defines the production targets and constraints for
five wells as follows:

\begin{verbatim}
--
--       WELL PRODUCTION WELL CONTROLS
--
-- WELL  OPEN/  CNTL  OIL   WAT   GAS   LIQ    RES    BHP    THP   VFP    VFP
-- NAME  SHUT   MODE  RATE  RATE  RATE  RATE   RATE   PRES   PRES  TABLE  ALFQ
WCONPROD
OP01     OPEN   GRUP  5E3   1*     1*    1*    1*     500.0                    /
OP02     OPEN   GRUP 10E3   1*     1*    1*    1*     200.0  500.0  2     0.0  /
OP03     OPEN   GRUP 15E3   1*     1*    1*    1*     200.0  500.0  3    10.0  /
OP04     OPEN   ORAT 20E3   1*     1*    1*    1*     500.0                    /
OP05     SHUT   GRUP 20E3   1*     1*    1*    1*     500.0                    /
/
\end{verbatim}

Well OP01 is open and is on group control, subject to a maximum oil rate
constraint of 5,000 stb/d and a minimum bottom-hole pressure of 500
psia. OP02 is also open and on group control but it's maximum oil rate
constraint has been set 10,000 stb/d, and is subject to a minimum
bottom-hole pressure limit of 200 psia and a minimum tubing head
pressure limit of 500 psia using VFPPROD vertical lift table number two.
Well OP03 is very similar to OP02, but with a 15,000 stb/d maximum oil
constraint and using VFPPROD vertical lift table number three with an
artificial lift parameter of 10. The next well is not on group control.
Well OP04 is open and has an oil rate target of 20,000 stb/d, subject to
a minimum bottom-hole pressure of 500 psia. Finally, well OP05 is shut
and will not be brought back on production despite being put under group
control, as the well has been declared shut.

The next example defines the production targets and constraints for five
wells, of which well OP01 is under group control as the well's group
target and constraints have been set with the GCONPROD keyword, and
wells OP02 to OP05 belong to groups that have not had their group
constraints set by GCONPROD.

\begin{verbatim}
--
--       GROUP PRODUCTION CONTROLS
--
-- GRUP  CNTL  OIL    WAT    GAS    LIQ    CNTL  GRUP  GUIDE  GUIDE  CNTL
-- NAME  MODE  RATE   RATE   RATE   RATE   OPT   CNTL  RATE   DEF    WAT
GCONPROD
GRP01    SAT   25E3   1*     1*     1*     1*     1*    1*     1*     1*       /
/
--
--       WELL SPECIFICATION DATA
--
-- WELL  GROUP     LOCATION  BHP    PHASE  DRAIN  INFLOW  OPEN  CROSS  PVT
-- NAME  NAME        I    J  DEPTH  FLUID  AREA   EQUANS  SHUT  FLOW   TABLE
WELSPECS
OP01     SAT        14   13   1*     OIL    1*     P-P    SHUT   NO     1*     /
0P02     PLATFORM   64   80   1*     OIL    1*     GPP    SHUT   NO     1*     /
OP03     PLATFORM   24  110   1*     OIL    1*     STD    SHUT   NO     1*     /
OP04     PLATFORM   24  110   1*     OIL    1*     STD    SHUT   NO     1*     /
OP05     PLATFORM   24  110   1*     OIL    1*     STD    SHUT   NO     1*     /
/
--
--       WELL PRODUCTION WELL CONTROLS
--
-- WELL  OPEN/  CNTL  OIL   WAT   GAS   LIQ    RES    BHP    THP   VFP    VFP
-- NAME  SHUT   MODE  RATE  RATE  RATE  RATE   RATE   PRES   PRES  TABLE  ALFQ
WCONPROD
OP01     OPEN   1*   15E3    1*    1*    1*    1*     500.0                    /
OP02     OPEN   1*    1*    5E3*  15E6   1*    1*     200.0  500.0  2      0.0 /
OP03     OPEN   1*    1*     1*    1*    1*    1*     1*     500.0  3     10.0 /
OP04     OPEN   ''    1*     1*    1*    1*    1*     1*     1*     1*     1*  /
OP05     SHUT   1*    1*     1*    1*    1*    1*     1*     1*     1*     1*  /
/
\end{verbatim}

Here, well OP01's control mode, TARGET, will be set to group control
(GRUP) and the well's production parameters will all act as constraints.
For well OP02, the well's control mode is set to GRAT, and the other
production parameters all act as constraints. Well OP03 control mode is
set to THP and if the well's BHP value has been previously defined then
that value will be used as a constraint; otherwise the default value of
one atmosphere will be used instead. Finally, wells OP04 and OP05 will
have their control mode set to BHP control. Note in this case, OP04's
status has been set to OPEN; thus the well will produce at a target BHP
of one atmosphere, which is unrealistic.

\subsection{WCUTBACK -- Define Well Cutback Limits and
Parameters}\label{wcutback-define-well-cutback-limits-and-parameters}

\begin{longtable}[]{@{}
  >{\raggedright\arraybackslash}p{(\linewidth - 14\tabcolsep) * \real{0.1250}}
  >{\raggedright\arraybackslash}p{(\linewidth - 14\tabcolsep) * \real{0.1250}}
  >{\raggedright\arraybackslash}p{(\linewidth - 14\tabcolsep) * \real{0.1250}}
  >{\raggedright\arraybackslash}p{(\linewidth - 14\tabcolsep) * \real{0.1250}}
  >{\raggedright\arraybackslash}p{(\linewidth - 14\tabcolsep) * \real{0.1250}}
  >{\raggedright\arraybackslash}p{(\linewidth - 14\tabcolsep) * \real{0.1250}}
  >{\raggedright\arraybackslash}p{(\linewidth - 14\tabcolsep) * \real{0.1250}}
  >{\raggedright\arraybackslash}p{(\linewidth - 14\tabcolsep) * \real{0.1250}}@{}}
\toprule\noalign{}
\begin{minipage}[b]{\linewidth}\raggedright
RUNSPEC
\end{minipage} & \begin{minipage}[b]{\linewidth}\raggedright
GRID
\end{minipage} & \begin{minipage}[b]{\linewidth}\raggedright
EDIT
\end{minipage} & \begin{minipage}[b]{\linewidth}\raggedright
PROPS
\end{minipage} & \begin{minipage}[b]{\linewidth}\raggedright
REGIONS
\end{minipage} & \begin{minipage}[b]{\linewidth}\raggedright
SOLUTION
\end{minipage} & \begin{minipage}[b]{\linewidth}\raggedright
SUMMARY
\end{minipage} & \begin{minipage}[b]{\linewidth}\raggedright
SCHEDULE
\end{minipage} \\
\midrule\noalign{}
\endhead
\bottomrule\noalign{}
\endlastfoot
\end{longtable}

\subsubsection{Description}\label{description-1047}

This keyword, WCUTBACK, defines a well's cutback limits and parameters
for both production and injection wells. See also the GCUTBACK keyword
in the SCHEDULE section that provides similar functionality for groups.

This keyword is not supported by OPM Flow but would change the results
if supported so the simulation will be stopped.

\subsection{WCUTBACT -- Define Well Tracer Cutback Limits and
Parameters}\label{wcutbact-define-well-tracer-cutback-limits-and-parameters}

\begin{longtable}[]{@{}
  >{\raggedright\arraybackslash}p{(\linewidth - 14\tabcolsep) * \real{0.1250}}
  >{\raggedright\arraybackslash}p{(\linewidth - 14\tabcolsep) * \real{0.1250}}
  >{\raggedright\arraybackslash}p{(\linewidth - 14\tabcolsep) * \real{0.1250}}
  >{\raggedright\arraybackslash}p{(\linewidth - 14\tabcolsep) * \real{0.1250}}
  >{\raggedright\arraybackslash}p{(\linewidth - 14\tabcolsep) * \real{0.1250}}
  >{\raggedright\arraybackslash}p{(\linewidth - 14\tabcolsep) * \real{0.1250}}
  >{\raggedright\arraybackslash}p{(\linewidth - 14\tabcolsep) * \real{0.1250}}
  >{\raggedright\arraybackslash}p{(\linewidth - 14\tabcolsep) * \real{0.1250}}@{}}
\toprule\noalign{}
\begin{minipage}[b]{\linewidth}\raggedright
RUNSPEC
\end{minipage} & \begin{minipage}[b]{\linewidth}\raggedright
GRID
\end{minipage} & \begin{minipage}[b]{\linewidth}\raggedright
EDIT
\end{minipage} & \begin{minipage}[b]{\linewidth}\raggedright
PROPS
\end{minipage} & \begin{minipage}[b]{\linewidth}\raggedright
REGIONS
\end{minipage} & \begin{minipage}[b]{\linewidth}\raggedright
SOLUTION
\end{minipage} & \begin{minipage}[b]{\linewidth}\raggedright
SUMMARY
\end{minipage} & \begin{minipage}[b]{\linewidth}\raggedright
SCHEDULE
\end{minipage} \\
\midrule\noalign{}
\endhead
\bottomrule\noalign{}
\endlastfoot
\end{longtable}

\subsubsection{Description}\label{description-1048}

This keyword, WCUTBACT, defines a production well's cutback limits and
parameters based on the named produced tracer from the well. See also
the GCUTBACT keyword in the SCHEDULE section that provides similar
functionality for groups.

This keyword is not supported by OPM Flow but would change the results
if supported so the simulation will be stopped.

\subsection{WCYCLE -- Define Automatic Well Opening and Closing Cycling
Parameters}\label{wcycle-define-automatic-well-opening-and-closing-cycling-parameters}

\begin{longtable}[]{@{}
  >{\raggedright\arraybackslash}p{(\linewidth - 14\tabcolsep) * \real{0.1250}}
  >{\raggedright\arraybackslash}p{(\linewidth - 14\tabcolsep) * \real{0.1250}}
  >{\raggedright\arraybackslash}p{(\linewidth - 14\tabcolsep) * \real{0.1250}}
  >{\raggedright\arraybackslash}p{(\linewidth - 14\tabcolsep) * \real{0.1250}}
  >{\raggedright\arraybackslash}p{(\linewidth - 14\tabcolsep) * \real{0.1250}}
  >{\raggedright\arraybackslash}p{(\linewidth - 14\tabcolsep) * \real{0.1250}}
  >{\raggedright\arraybackslash}p{(\linewidth - 14\tabcolsep) * \real{0.1250}}
  >{\raggedright\arraybackslash}p{(\linewidth - 14\tabcolsep) * \real{0.1250}}@{}}
\toprule\noalign{}
\begin{minipage}[b]{\linewidth}\raggedright
RUNSPEC
\end{minipage} & \begin{minipage}[b]{\linewidth}\raggedright
GRID
\end{minipage} & \begin{minipage}[b]{\linewidth}\raggedright
EDIT
\end{minipage} & \begin{minipage}[b]{\linewidth}\raggedright
PROPS
\end{minipage} & \begin{minipage}[b]{\linewidth}\raggedright
REGIONS
\end{minipage} & \begin{minipage}[b]{\linewidth}\raggedright
SOLUTION
\end{minipage} & \begin{minipage}[b]{\linewidth}\raggedright
SUMMARY
\end{minipage} & \begin{minipage}[b]{\linewidth}\raggedright
SCHEDULE
\end{minipage} \\
\midrule\noalign{}
\endhead
\bottomrule\noalign{}
\endlastfoot
\end{longtable}

\subsubsection{Description}\label{description-1049}

The WCYCLE keyword defines automatic well opening and closing cycling
parameters. These are used to model for example ``Huff and Puff'' cyclic
steam injection in heavy oil reservoirs or Water-Alternating-Gas
(``WAG'') processes in enhanced oil recovery modeling. The keyword
defines specific time periods for automatically cycling wells on and
off. For example in a WAG scheme the water injection wells would have
one set of cycling parameters and the gas injection wells another, such
that only one type of well is active at a time.

\begin{longtable}[]{@{}
  >{\raggedright\arraybackslash}p{(\linewidth - 6\tabcolsep) * \real{0.2500}}
  >{\raggedright\arraybackslash}p{(\linewidth - 6\tabcolsep) * \real{0.2500}}
  >{\raggedright\arraybackslash}p{(\linewidth - 6\tabcolsep) * \real{0.2500}}
  >{\raggedright\arraybackslash}p{(\linewidth - 6\tabcolsep) * \real{0.2500}}@{}}
\toprule\noalign{}
\begin{minipage}[b]{\linewidth}\raggedright
No.
\end{minipage} & \begin{minipage}[b]{\linewidth}\raggedright
Name
\end{minipage} & \begin{minipage}[b]{\linewidth}\raggedright
Description
\end{minipage} & \begin{minipage}[b]{\linewidth}\raggedright
Default
\end{minipage} \\
\midrule\noalign{}
\endhead
\bottomrule\noalign{}
\endlastfoot
Field & Metric & Laboratory & \\
1 & WELNAME & A character string of up to eight characters in length
that defines the well name for which the well automatic cycling
parameters are being defined. Note that the well name (WELNAME) must
have been declared previously using the WELSPECS keyword in the SCHEDULE
section, otherwise an error may occur. & None \\
4 & ONTIME & A real positive value that defines the length of time for
the on period. The well will be turned off at the beginning of the first
time step after the on period has elapsed. If ONTIME is zero or negative
then the well will not be turned on by the automatic well cycling. &
0.0 \\
day & day & hour & \\
5 & OFFTIME & A real positive value that defines the length of time for
the off period. If the well was turned off by automatic cycling then it
will be turned on at the beginning of the first time step after the off
period has elasped. If the well was turned off other than by automatic
cycling then the well will remain turned off. If OFFTIME is zero or
negative then the well will not be turned on by the automatic well
cycling. & 0.0 \\
day & day & hour & \\
6 & STARTTIME & A real positive value that defines the length of time
for the start-up period during which the well efficiency factor will be
ramped up. The start-up period begins at the start of the on period. The
well's efficiency factor is ramped up linearly from zero at the start of
the on period to the value specified previously by the WEFAC keyword
after the start-up period has elasped. The well efficiency factor to
apply during each time step is linearly interpolated based on the time
at the end of the time step. If the time step size is small compared
with STARTTIME then the rate will be ramped up gradually. Whereas, if
the time step size is greater than STARTTIME then the rate will ramp up
instantaneously potentially causing convergence problems. & 0.0 \\
day & day & hour & \\
7 & MAXTS & A real positive value that defines the maximum time step
length when the well is turned on by automatic cycling. This overrides
the standard maximum time step length. This can help to reduce potential
convergence problems caused when a high rate well is opened. If MAXTS is
zero then the standard maximum time step length applies. & 0.0 \\
day & day & hour & \\
13 & CTRLTS & A defined character string that specifies whether the time
steps should be controlled to coincide exactly with the automatic
cycling times for this well. CTRLTS should be set to one of the
following: & NO \\
Notes: & & & \\
\end{longtable}

\emph{Table 12.3.258.1: WCYCLE Keyword Description}

\subsubsection{Example}\label{example-563}

The following example defines the cycle lengths for a
Water-Alternating-Gas injection scheme with a water-gas slug length
ratio of 1:2, and initiates automatic cycling in the water injection
well. A maximum time step of 5 days has been specified following the
start of each injection period. The simulation then proceeds to the end
of the first water injection cycle when automatic cycling is initiated
in the gas injection well.

\begin{verbatim}
--
-- AUTOMATIC CYCLING OF WELLS ON AND OFF
--
WCYCLE
--WELL     ON       OFF      STARTUP  STARTUP  CTRL
--NAME     PERIOD   PERIOD   TIME     MAXTS    TIMESTEP
  I01W     30.0     60.0     1*       5.0      YES  /
  I01G     60.0     30.0     1*       5.0      YES  /
--
-- OPEN WELLS
--
WELOPEN
--WELL     OPEN
--NAME     SHUT
  I01W     OPEN /
  I01G     SHUT /
/
--
-- REPORT STEPS
--
TSTEP
  30.0     /
--
-- OPEN WELLS
--
WELOPEN
--WELL     OPEN
--NAME     SHUT
  I01G     OPEN /
/
\end{verbatim}

Note that the wells initially need to be opened manually in order to
start the automatic cycling. In addition, this needs to be done at the
correct times to ensure that the water and gas injection wells are
syncronised.

\subsection{WDFAC -- Define Gas Flow Dependent Skin
Factor}\label{wdfac-define-gas-flow-dependent-skin-factor}

\begin{longtable}[]{@{}
  >{\raggedright\arraybackslash}p{(\linewidth - 14\tabcolsep) * \real{0.1250}}
  >{\raggedright\arraybackslash}p{(\linewidth - 14\tabcolsep) * \real{0.1250}}
  >{\raggedright\arraybackslash}p{(\linewidth - 14\tabcolsep) * \real{0.1250}}
  >{\raggedright\arraybackslash}p{(\linewidth - 14\tabcolsep) * \real{0.1250}}
  >{\raggedright\arraybackslash}p{(\linewidth - 14\tabcolsep) * \real{0.1250}}
  >{\raggedright\arraybackslash}p{(\linewidth - 14\tabcolsep) * \real{0.1250}}
  >{\raggedright\arraybackslash}p{(\linewidth - 14\tabcolsep) * \real{0.1250}}
  >{\raggedright\arraybackslash}p{(\linewidth - 14\tabcolsep) * \real{0.1250}}@{}}
\toprule\noalign{}
\begin{minipage}[b]{\linewidth}\raggedright
RUNSPEC
\end{minipage} & \begin{minipage}[b]{\linewidth}\raggedright
GRID
\end{minipage} & \begin{minipage}[b]{\linewidth}\raggedright
EDIT
\end{minipage} & \begin{minipage}[b]{\linewidth}\raggedright
PROPS
\end{minipage} & \begin{minipage}[b]{\linewidth}\raggedright
REGIONS
\end{minipage} & \begin{minipage}[b]{\linewidth}\raggedright
SOLUTION
\end{minipage} & \begin{minipage}[b]{\linewidth}\raggedright
SUMMARY
\end{minipage} & \begin{minipage}[b]{\linewidth}\raggedright
SCHEDULE
\end{minipage} \\
\midrule\noalign{}
\endhead
\bottomrule\noalign{}
\endlastfoot
\end{longtable}

\subsubsection{Description}\label{description-1050}

The WDFAC keyword defines a gas well's D-factor (flow dependent skin
factor), which is normally derived from well tests or calculated
analytically based on the coefficient of inertial resistance, usually
known as β, in Forchheimer's flow equation, {[}Geertsma, J., 1974.
Estimating the Coefficient of Inertial Resistance in Fluid Flow Through
Porous Media. Soc.Pet.Eng.J., October: 445-450.{]}, {[}Gewers, C.W.W.
and Nichol, L.R., 1969. Gas Turbulence Factor in a Microvugular
Carbonate. J.Can.Pet.Tech., April.{]} and {[}Wong, S.W., 1970. Effects
of Liquid Saturation on Turbulence Factors for Gas Liquid Systems.
J.Can.Pet.Tech., October{]}.

Internally the simulator evaluates a D-factor for each well connection
by assuming the D-factors are inversely proportional to the connection
transmissibility factors. This is done at the time the well D-factor is
specified. If a connection is later closed then the well D-factor is
likely to change but the connection D-factors for each of the remaining
open connections should not be affected.

\begin{longtable}[]{@{}
  >{\raggedright\arraybackslash}p{(\linewidth - 6\tabcolsep) * \real{0.2500}}
  >{\raggedright\arraybackslash}p{(\linewidth - 6\tabcolsep) * \real{0.2500}}
  >{\raggedright\arraybackslash}p{(\linewidth - 6\tabcolsep) * \real{0.2500}}
  >{\raggedright\arraybackslash}p{(\linewidth - 6\tabcolsep) * \real{0.2500}}@{}}
\toprule\noalign{}
\begin{minipage}[b]{\linewidth}\raggedright
No.
\end{minipage} & \begin{minipage}[b]{\linewidth}\raggedright
Name
\end{minipage} & \begin{minipage}[b]{\linewidth}\raggedright
Description
\end{minipage} & \begin{minipage}[b]{\linewidth}\raggedright
Default
\end{minipage} \\
\midrule\noalign{}
\endhead
\bottomrule\noalign{}
\endlastfoot
Field & Metric & Laboratory & \\
1 & WELNAME & A character string of up to eight characters in length
that defines the well name for which the well D-factor is being defined.
Note that the well name (WELNAME) must have been declared previously
using the WELSPECS keyword in the SCHEDULE section. & None \\
2 & DFACTOR & A real positive value greater than or equal to zero that
defines the D-factor for the well. & 0.0 \\
day/Mscf & day/m3 & hour/scc & \\
Notes: & & & \\
\end{longtable}

\emph{Table 12.3.259.1: WDFAC Keyword Description}

See also the WDFACCOR keyword in the SCHEDULE section that uses Dake's
{[}Dake, L.P. Fundamentals of Reservoir Engineering, Amsterdam, The
Netherlands, Elsevier Science BV (1978) Chapter 8.6, pages 252-257.{]}
correlation to calculate the well D-factor; and the COMPDAT keyword in
the SCHEDULE section that can be used to specify the connection
D-factors directly.

\subsubsection{Example}\label{example-564}

\begin{verbatim}
--
--       WELL D-FACTORS
--
-- WELL  D-
-- NAME  FACTOR
WEFAC
  GP01   1.0E-03 /
/
\end{verbatim}

In the above example a well D-factor of 10-3 is specified for gas well
GP01.

\subsection{WDFACCOR -- Gas Flow Dependent Skin Factor
(Correlation)}\label{wdfaccor-gas-flow-dependent-skin-factor-correlation}

\begin{longtable}[]{@{}
  >{\raggedright\arraybackslash}p{(\linewidth - 14\tabcolsep) * \real{0.1250}}
  >{\raggedright\arraybackslash}p{(\linewidth - 14\tabcolsep) * \real{0.1250}}
  >{\raggedright\arraybackslash}p{(\linewidth - 14\tabcolsep) * \real{0.1250}}
  >{\raggedright\arraybackslash}p{(\linewidth - 14\tabcolsep) * \real{0.1250}}
  >{\raggedright\arraybackslash}p{(\linewidth - 14\tabcolsep) * \real{0.1250}}
  >{\raggedright\arraybackslash}p{(\linewidth - 14\tabcolsep) * \real{0.1250}}
  >{\raggedright\arraybackslash}p{(\linewidth - 14\tabcolsep) * \real{0.1250}}
  >{\raggedright\arraybackslash}p{(\linewidth - 14\tabcolsep) * \real{0.1250}}@{}}
\toprule\noalign{}
\begin{minipage}[b]{\linewidth}\raggedright
RUNSPEC
\end{minipage} & \begin{minipage}[b]{\linewidth}\raggedright
GRID
\end{minipage} & \begin{minipage}[b]{\linewidth}\raggedright
EDIT
\end{minipage} & \begin{minipage}[b]{\linewidth}\raggedright
PROPS
\end{minipage} & \begin{minipage}[b]{\linewidth}\raggedright
REGIONS
\end{minipage} & \begin{minipage}[b]{\linewidth}\raggedright
SOLUTION
\end{minipage} & \begin{minipage}[b]{\linewidth}\raggedright
SUMMARY
\end{minipage} & \begin{minipage}[b]{\linewidth}\raggedright
SCHEDULE
\end{minipage} \\
\midrule\noalign{}
\endhead
\bottomrule\noalign{}
\endlastfoot
\end{longtable}

\subsubsection{Description}\label{description-1051}

WDFACCOR keyword defines the parameters to calculate a gas well's
connection D-factors (flow dependent skin factor) for each connection
based on a correlation for the coefficient of inertial resistance,
usually known as β, in Forchheimer's flow equation {[}Dake, L.P.
Fundamentals of Reservoir Engineering, Amsterdam, The Netherlands,
Elsevier Science BV (1978) Chapter 8.6, pages 252-257.{]}, {[}Geertsma,
J., 1974. Estimating the Coefficient of Inertial Resistance in Fluid
Flow Through Porous Media. Soc.Pet.Eng.J., October: 445-450.{]},
{[}Gewers, C.W.W. and Nichol, L.R., 1969. Gas Turbulence Factor in a
Microvugular Carbonate. J.Can.Pet.Tech., April.{]} and {[}Wong, S.W.,
1970. Effects of Liquid Saturation on Turbulence Factors for Gas Liquid
Systems. J.Can.Pet.Tech., October{]}. This keyword uses Dake's
correlation to calculate the D-factor.

\begin{longtable}[]{@{}
  >{\raggedright\arraybackslash}p{(\linewidth - 6\tabcolsep) * \real{0.2500}}
  >{\raggedright\arraybackslash}p{(\linewidth - 6\tabcolsep) * \real{0.2500}}
  >{\raggedright\arraybackslash}p{(\linewidth - 6\tabcolsep) * \real{0.2500}}
  >{\raggedright\arraybackslash}p{(\linewidth - 6\tabcolsep) * \real{0.2500}}@{}}
\toprule\noalign{}
\begin{minipage}[b]{\linewidth}\raggedright
No.
\end{minipage} & \begin{minipage}[b]{\linewidth}\raggedright
Name
\end{minipage} & \begin{minipage}[b]{\linewidth}\raggedright
Description
\end{minipage} & \begin{minipage}[b]{\linewidth}\raggedright
Default
\end{minipage} \\
\midrule\noalign{}
\endhead
\bottomrule\noalign{}
\endlastfoot
Field & Metric & Laboratory & \\
1 & WELNAME & A character string of up to eight characters in length
that defines the well name for which the well D-factor correlation is
being defined. Note that the well name (WELNAME) must have been declared
previously using the WELSPECS keyword in the SCHEDULE section. & None \\
2 & A & A real value greater than or equal to zero that defines the
coefficient A in the D-factor correlation. & 0.0 \\
cP.day.ft2/mD/Mscf & cP.day.m2/mD/sm3 & cP.hour.cm2/mD/scc & \\
3 & B & A real value that defines the exponent B of the grid block
permeability in the D-factor correlation. & 0.0 \\
dimensionless & dimensionless & dimensionless & \\
4 & C & A real value that defines the exponent C of the grid block
porosity in the D-factor correlation. & 0.0 \\
dimensionless & dimensionless & dimensionless & \\
Notes: & & & \\
\end{longtable}

\emph{Table 12.3.260.1: WDFACCOR Keyword Description}

The simulator evaluates the connection D-factors using the following
expression based on Dake's correlation:

\begin{equation}\phantomsection\label{eq-12-3-260-1}{
D = A{K}_{e}^{B}{ϕ}^{C}\frac{{K}_{e}}{h}\frac{1}{{r}_{w}}\frac{{γ}_{g}}{{μ}_{g}}
}\end{equation}

Where:

\({K}_{e}\) = effective permeability of the grid block,

\(ϕ\) = porosity of the grid block,

\(h\) = connection length,

\({r}_{w}\) = wellbore radius,

\({γ}_{g}\) = gas gravity at surface conditions,

\({μ}_{g}\) = gas viscosity at bottom hole conditions.

Note that since \({μ}_{g}\) is dependent on pressure the D-factor will
also vary with pressure.

See also the WDFAC keyword in the SCHEDULE section that can be used to
specify the well D-factor; and the COMPDAT keyword in the SCHEDULE
section that can be used to specify the connection D-factors directly.

\subsubsection{Examples}\label{examples-91}

In the first example the connection D-factors for gas well GP01 are
evaluated for correlation coefficients specified in field units based on
the laboratory determined relationship between β and the absolute
permeability presented by Dake in Fig. 8.8 (variation of β with porosity
has been neglected).

\begin{verbatim}
--
--       WELL D-FACTOR CORRELATIONS
--
-- WELL
-- NAME    A         B         C
WEFAC
  GP01     6.04E-05  -1.1045   0.0 /
/
\end{verbatim}

In the second example the same correlation coefficients are specified in
metric units.

\begin{verbatim}
--
--       WELL D-FACTOR CORRELATIONS
--
-- WELL
-- NAME    A         B         C
WEFAC
  GP01     1.20E-07  -1.1045   0.0 /
/
\end{verbatim}

\subsection{WDRILPRI -- Add Wells to the Drilling Priority Drilling
Queue}\label{wdrilpri-add-wells-to-the-drilling-priority-drilling-queue}

\begin{longtable}[]{@{}
  >{\raggedright\arraybackslash}p{(\linewidth - 14\tabcolsep) * \real{0.1250}}
  >{\raggedright\arraybackslash}p{(\linewidth - 14\tabcolsep) * \real{0.1250}}
  >{\raggedright\arraybackslash}p{(\linewidth - 14\tabcolsep) * \real{0.1250}}
  >{\raggedright\arraybackslash}p{(\linewidth - 14\tabcolsep) * \real{0.1250}}
  >{\raggedright\arraybackslash}p{(\linewidth - 14\tabcolsep) * \real{0.1250}}
  >{\raggedright\arraybackslash}p{(\linewidth - 14\tabcolsep) * \real{0.1250}}
  >{\raggedright\arraybackslash}p{(\linewidth - 14\tabcolsep) * \real{0.1250}}
  >{\raggedright\arraybackslash}p{(\linewidth - 14\tabcolsep) * \real{0.1250}}@{}}
\toprule\noalign{}
\begin{minipage}[b]{\linewidth}\raggedright
RUNSPEC
\end{minipage} & \begin{minipage}[b]{\linewidth}\raggedright
GRID
\end{minipage} & \begin{minipage}[b]{\linewidth}\raggedright
EDIT
\end{minipage} & \begin{minipage}[b]{\linewidth}\raggedright
PROPS
\end{minipage} & \begin{minipage}[b]{\linewidth}\raggedright
REGIONS
\end{minipage} & \begin{minipage}[b]{\linewidth}\raggedright
SOLUTION
\end{minipage} & \begin{minipage}[b]{\linewidth}\raggedright
SUMMARY
\end{minipage} & \begin{minipage}[b]{\linewidth}\raggedright
SCHEDULE
\end{minipage} \\
\midrule\noalign{}
\endhead
\bottomrule\noalign{}
\endlastfoot
\end{longtable}

\subsubsection{Description}\label{description-1052}

This keyword, WDRILPRI, adds wells to the Drilling Priority Drilling
Queue and defines the well priority and drilling unit number or batch
queue sequence for the well. The batch queue sequence number enables all
wells with the same sequence number to drilled at the same time.

This keyword is not supported by OPM Flow but would change the results
if supported so the simulation will be stopped.

\subsection{WDRILRES -- Activate Prevention of Multi-Completions in the
Same Cell for Queued
Wells}\label{wdrilres-activate-prevention-of-multi-completions-in-the-same-cell-for-queued-wells}

\begin{longtable}[]{@{}
  >{\raggedright\arraybackslash}p{(\linewidth - 14\tabcolsep) * \real{0.1250}}
  >{\raggedright\arraybackslash}p{(\linewidth - 14\tabcolsep) * \real{0.1250}}
  >{\raggedright\arraybackslash}p{(\linewidth - 14\tabcolsep) * \real{0.1250}}
  >{\raggedright\arraybackslash}p{(\linewidth - 14\tabcolsep) * \real{0.1250}}
  >{\raggedright\arraybackslash}p{(\linewidth - 14\tabcolsep) * \real{0.1250}}
  >{\raggedright\arraybackslash}p{(\linewidth - 14\tabcolsep) * \real{0.1250}}
  >{\raggedright\arraybackslash}p{(\linewidth - 14\tabcolsep) * \real{0.1250}}
  >{\raggedright\arraybackslash}p{(\linewidth - 14\tabcolsep) * \real{0.1250}}@{}}
\toprule\noalign{}
\begin{minipage}[b]{\linewidth}\raggedright
RUNSPEC
\end{minipage} & \begin{minipage}[b]{\linewidth}\raggedright
GRID
\end{minipage} & \begin{minipage}[b]{\linewidth}\raggedright
EDIT
\end{minipage} & \begin{minipage}[b]{\linewidth}\raggedright
PROPS
\end{minipage} & \begin{minipage}[b]{\linewidth}\raggedright
REGIONS
\end{minipage} & \begin{minipage}[b]{\linewidth}\raggedright
SOLUTION
\end{minipage} & \begin{minipage}[b]{\linewidth}\raggedright
SUMMARY
\end{minipage} & \begin{minipage}[b]{\linewidth}\raggedright
SCHEDULE
\end{minipage} \\
\midrule\noalign{}
\endhead
\bottomrule\noalign{}
\endlastfoot
\end{longtable}

\subsubsection{Description}\label{description-1053}

The WDRILRES keyword activates the prevention of multi-completions being
completed in the same cell for wells in a drilling queue. Setting this
option stops any well defined as a queued well via the QDRILL and
WDRILLPRI keywords in the SCHEDULE section, or any wells set to
automatic opening by setting the STATUS variable to AUTO on the WCONPROD
keyword in the RUNSPEC section, from opening if there is an already
existing active well connection to a cell.

This keyword is not supported by OPM Flow but would change the results
if supported so the simulation will be stopped.

\subsection{WDRILTIM -- Define Drilling Parameters for Automatic
Drilling of New
Wells}\label{wdriltim-define-drilling-parameters-for-automatic-drilling-of-new-wells}

\begin{longtable}[]{@{}
  >{\raggedright\arraybackslash}p{(\linewidth - 14\tabcolsep) * \real{0.1250}}
  >{\raggedright\arraybackslash}p{(\linewidth - 14\tabcolsep) * \real{0.1250}}
  >{\raggedright\arraybackslash}p{(\linewidth - 14\tabcolsep) * \real{0.1250}}
  >{\raggedright\arraybackslash}p{(\linewidth - 14\tabcolsep) * \real{0.1250}}
  >{\raggedright\arraybackslash}p{(\linewidth - 14\tabcolsep) * \real{0.1250}}
  >{\raggedright\arraybackslash}p{(\linewidth - 14\tabcolsep) * \real{0.1250}}
  >{\raggedright\arraybackslash}p{(\linewidth - 14\tabcolsep) * \real{0.1250}}
  >{\raggedright\arraybackslash}p{(\linewidth - 14\tabcolsep) * \real{0.1250}}@{}}
\toprule\noalign{}
\begin{minipage}[b]{\linewidth}\raggedright
RUNSPEC
\end{minipage} & \begin{minipage}[b]{\linewidth}\raggedright
GRID
\end{minipage} & \begin{minipage}[b]{\linewidth}\raggedright
EDIT
\end{minipage} & \begin{minipage}[b]{\linewidth}\raggedright
PROPS
\end{minipage} & \begin{minipage}[b]{\linewidth}\raggedright
REGIONS
\end{minipage} & \begin{minipage}[b]{\linewidth}\raggedright
SOLUTION
\end{minipage} & \begin{minipage}[b]{\linewidth}\raggedright
SUMMARY
\end{minipage} & \begin{minipage}[b]{\linewidth}\raggedright
SCHEDULE
\end{minipage} \\
\midrule\noalign{}
\endhead
\bottomrule\noalign{}
\endlastfoot
\end{longtable}

\subsubsection{Description}\label{description-1054}

WDRILTIM defines the automatic drilling parameters used to describe the
numbers of days taken to drill a well, the drilling status of the well,
and status of other wells when drilling an automatically drilled well.

This keyword is not supported by OPM Flow but would change the results
if supported so the simulation will be stopped.

\subsection{WECON -- Well Economic Criteria for Production
Wells}\label{wecon-well-economic-criteria-for-production-wells}

\begin{longtable}[]{@{}
  >{\raggedright\arraybackslash}p{(\linewidth - 14\tabcolsep) * \real{0.1250}}
  >{\raggedright\arraybackslash}p{(\linewidth - 14\tabcolsep) * \real{0.1250}}
  >{\raggedright\arraybackslash}p{(\linewidth - 14\tabcolsep) * \real{0.1250}}
  >{\raggedright\arraybackslash}p{(\linewidth - 14\tabcolsep) * \real{0.1250}}
  >{\raggedright\arraybackslash}p{(\linewidth - 14\tabcolsep) * \real{0.1250}}
  >{\raggedright\arraybackslash}p{(\linewidth - 14\tabcolsep) * \real{0.1250}}
  >{\raggedright\arraybackslash}p{(\linewidth - 14\tabcolsep) * \real{0.1250}}
  >{\raggedright\arraybackslash}p{(\linewidth - 14\tabcolsep) * \real{0.1250}}@{}}
\toprule\noalign{}
\begin{minipage}[b]{\linewidth}\raggedright
RUNSPEC
\end{minipage} & \begin{minipage}[b]{\linewidth}\raggedright
GRID
\end{minipage} & \begin{minipage}[b]{\linewidth}\raggedright
EDIT
\end{minipage} & \begin{minipage}[b]{\linewidth}\raggedright
PROPS
\end{minipage} & \begin{minipage}[b]{\linewidth}\raggedright
REGIONS
\end{minipage} & \begin{minipage}[b]{\linewidth}\raggedright
SOLUTION
\end{minipage} & \begin{minipage}[b]{\linewidth}\raggedright
SUMMARY
\end{minipage} & \begin{minipage}[b]{\linewidth}\raggedright
SCHEDULE
\end{minipage} \\
\midrule\noalign{}
\endhead
\bottomrule\noalign{}
\endlastfoot
\end{longtable}

\subsubsection{Description}\label{description-1055}

The WECON keyword defines the economic criteria for production wells
that have previously been defined by the WELSPECS and WCONPROD keywords
in the SCHEDULE section.

Note that wells can be allocated to a group when they are specified by
the WELSPECS keyword and groups can also have economic controls. Wells
under group control are therefore subject to the economic criteria set
via the GCONPROD and GECON keywords in the SCHEDULE section and the
economic criteria specified by the WECON keyword.

\begin{longtable}[]{@{}
  >{\raggedright\arraybackslash}p{(\linewidth - 6\tabcolsep) * \real{0.2500}}
  >{\raggedright\arraybackslash}p{(\linewidth - 6\tabcolsep) * \real{0.2500}}
  >{\raggedright\arraybackslash}p{(\linewidth - 6\tabcolsep) * \real{0.2500}}
  >{\raggedright\arraybackslash}p{(\linewidth - 6\tabcolsep) * \real{0.2500}}@{}}
\toprule\noalign{}
\begin{minipage}[b]{\linewidth}\raggedright
No.
\end{minipage} & \begin{minipage}[b]{\linewidth}\raggedright
Name
\end{minipage} & \begin{minipage}[b]{\linewidth}\raggedright
Description
\end{minipage} & \begin{minipage}[b]{\linewidth}\raggedright
Default
\end{minipage} \\
\midrule\noalign{}
\endhead
\bottomrule\noalign{}
\endlastfoot
Field & Metric & Laboratory & \\
1 & WELNAME & A character string of up to eight characters in length
that defines the well name for which the well economic criteria data is
being defined. Note that the well name (WELNAME) must have been declared
previously using the WELSPECS keyword in the SCHEDULE section, otherwise
an error may occur. & None \\
2 & ORAT & A real positive value that defines the minimum economic
surface oil production rate, below which an economic action will take
place, as outlined below: Only option (2) is supported by OPM Flow as
STATUS equals AUTO on the COMPDAT keyword is currently not supported by
the simulator. Hence, the well be either shut or stopped. A value less
than or equal to zero switches off this criterion. This value may be
specified using a User Defined Argument (UDA). & 0.0 \\
stb/d & sm3/day & scc/hour & \\
3 & GRAT & A real positive value that defines the minimum economic
surface gas production rate, below which an economic action will take
place, as outlined below: Only option (2) is supported by OPM Flow as
STATUS equals AUTO on the COMPDAT keyword is currently not supported by
the simulator. Hence, the well be either shut or stopped. A value less
than or equal to zero switches off this criterion. This value may be
specified using a User Defined Argument (UDA). & 0.0 \\
Mscf/d & sm3/day & scc/hour & \\
4 & WCUT & A real positive value that defines the maximum economic
surface water cut, above which an economic action will take place. This
value may be specified using a User Defined Argument (UDA). Water cut is
defined as:\({f}_{w} = \frac{{q}_{w}}{{q}_{w} + {q}_{o}}\), and the
various actions that are available if the water cut limit is exceeded
are described in item (7). A value less than or equal to zero switches
off this criterion. & 0.0 \\
dimensionless & dimensionless & dimensionless & \\
5 & GOR & A real positive value that defines the maximum economic
surface gas-oil ratio, above which an economic action will take place,
as defined by item (7). A value less than or equal to zero switches off
this criterion. This value may be specified using a User Defined
Argument (UDA). & 0.0 \\
Mscf/stb & sm3/sm3 & scc/scc & \\
6 & WGR & A real positive value that defines the maximum economic
surface water-gas ratio, above which an economic action will take place,
as defined by item (7). A value less than or equal to zero switches off
this criterion. This value may be specified using a User Defined
Argument (UDA). & 0.0 \\
stb/Mscf & sm3/sm3 & scc/scc & \\
7 & ACTION & A defined character string that defines the action to be
taken if the economic WCUT, GOR, or WGR limits are violated. ACTION
should be set to one of the following character strings: The corrective
action takes places at the end of the time step in which the constraint
is violated. The +CON option is not currently supported by OPM Flow. &
NONE \\
8 & END & A defined character string that defines if the simulation
should terminate if the well is shut or stopped. END should be set to
one of the following character strings: The YES option is not currently
supported by OPM Flow. & NO \\
9 & WELOPEN & A character string of up to eight characters in length
that defines the name of the well, which will be opened when the current
well (WELNAME) has been automatically closed/shut by the simulator.
Wells closed manually do not invoke this opton. The well name (WELOPEN)
must have been declared previously using the WELSPECS keyword in the
SCHEDULE section, otherwise an error may occur. & None \\
10 & ELTOPT & A defined character string that defines the type of
quantity used to test the economic limit, and should be set to one of
the following character strings: Wells under group control via the
GCONPROD keyword in the SCHEDULE section will be subject to the ELT,
except for when the group's production target rate is set to zero. Wells
under group control via the GCONPRI keyword in the SCHEDULE section will
be subject to the ELT, except for when the group has curtailed a well's
production. & RATE \\
11 & WCUT2 & A real positive value that defines the secondary maximum
economic surface water cut, above which an economic action will take
place. A value less than or equal to zero switches off this criterion.
This item is not supported by OPM Flow and should be defaulted (1*) or
set to zero. & 0.0 \\
12 & ACTION2 & A defined character string that defines the action to be
taken if the secondary economic WCUT2 limits is violated. This item is
not supported by OPM Flow and should be defaulted (1*). & 1* \\
13 & GLR & A real positive value that defines the maximum economic
surface gas-liquid ratio, above which an economic action will take
place, as defined by item (7). A value less than or equal to zero
switches off this criterion. This item is not supported by OPM Flow and
should be defaulted (1*). & None \\
14 & LRAT & A real positive value that defines the minimum liquid rate,
below which the well (WELNAME) is shut-in. This item is not supported by
OPM Flow and should be defaulted (1*). & 0.0 \\
15 & TEMP & A real positive value that defines the maximum economic
temperature in the commercial compositional Thermal/Temperature model
for a well. This item is not supported by OPM Flow and should be
defaulted (1*). & None \\
16 & RESV & A real positive value that defines the minimum reservoir
volume rate, below which the well (WELNAME) is shut-in. This item is not
supported by OPM Flow and should be defaulted (1*). & 0.0 \\
Notes: & & & \\
\end{longtable}

\emph{Table 12.81: WECON Keyword Description}

See also the WELSPECS keyword to define a wells shut-in or stop options,
GCONPROD for group controls, and GECON for setting a group's economic
criteria. All the aforementioned keywords are described in the SCHEDULE
section.

\subsubsection{Example}\label{example-565}

The following example defines one oil well and one gas well using the
WELSPECS keyword, together with their economic criteria.

\begin{verbatim}
--
-- WELL SPECIFICATION DATA
--
-- WELL     GROUP      LOCATION  BHP    PHASE  DRAIN  INFLOW  SHUT  CROSS  PRESS
-- NAME     NAME        I    J   DEPTH  FLUID  AREA   EQUA.   IN    FLOW   TABLE
WELSPECS
GP01      PLATFORM     14   13   1*      GAS   1*     GPP     SHUT   NO    1*  /
OP01      PLATFORM     28   96   1*      OIL   1*     STD     SHUT   NO    1*  /
/
--
--       WELL ECONOMIC CRITERIA FOR PRODUCTION WELLS
-- WELL  MIN    MIN    MAX    MAX    MAX    CNTL    END
-- NAME  ORAT   GRAT   WCUT   GOR    WGR    MODE    RUN
WECON
GP01     1*     5.0E3  1*     1*     1*    'WELL'  'NO'                        /
OP01     500    1*     0.95   15     1*    'WELL'  'YES'                       /
/
\end{verbatim}

Well GP01 has a minimum economic gas rate of 5 MMscf/d and will shut-in
if the gas rate falls below this rate, but the simulation will continue
even if this occurs. Well OP01 has a minimum economic oil rate of 500
stb/d, a maximum water cut limit of 95\%, and a maximum GOR of 15
Mscf/stb, if any any of these limits are violated the well will be
shut-in and the run should be terminated at the next reporting time
step, however the option to end the run is not currently supported by
OPM Flow.

\subsection{WECONINJ -- Well Economic Criteria for Injection
Wells}\label{weconinj-well-economic-criteria-for-injection-wells}

\begin{longtable}[]{@{}
  >{\raggedright\arraybackslash}p{(\linewidth - 14\tabcolsep) * \real{0.1250}}
  >{\raggedright\arraybackslash}p{(\linewidth - 14\tabcolsep) * \real{0.1250}}
  >{\raggedright\arraybackslash}p{(\linewidth - 14\tabcolsep) * \real{0.1250}}
  >{\raggedright\arraybackslash}p{(\linewidth - 14\tabcolsep) * \real{0.1250}}
  >{\raggedright\arraybackslash}p{(\linewidth - 14\tabcolsep) * \real{0.1250}}
  >{\raggedright\arraybackslash}p{(\linewidth - 14\tabcolsep) * \real{0.1250}}
  >{\raggedright\arraybackslash}p{(\linewidth - 14\tabcolsep) * \real{0.1250}}
  >{\raggedright\arraybackslash}p{(\linewidth - 14\tabcolsep) * \real{0.1250}}@{}}
\toprule\noalign{}
\begin{minipage}[b]{\linewidth}\raggedright
RUNSPEC
\end{minipage} & \begin{minipage}[b]{\linewidth}\raggedright
GRID
\end{minipage} & \begin{minipage}[b]{\linewidth}\raggedright
EDIT
\end{minipage} & \begin{minipage}[b]{\linewidth}\raggedright
PROPS
\end{minipage} & \begin{minipage}[b]{\linewidth}\raggedright
REGIONS
\end{minipage} & \begin{minipage}[b]{\linewidth}\raggedright
SOLUTION
\end{minipage} & \begin{minipage}[b]{\linewidth}\raggedright
SUMMARY
\end{minipage} & \begin{minipage}[b]{\linewidth}\raggedright
SCHEDULE
\end{minipage} \\
\midrule\noalign{}
\endhead
\bottomrule\noalign{}
\endlastfoot
\end{longtable}

\subsubsection{Description}\label{description-1056}

The WECONINJ keyword defines economic criteria for injection wells that
have previously been defined by the WELSPECS and WCONINJE keywords in
the SCHEDULE section.

Note that wells can be allocated to a group when they are specified by
the WELSPECS keyword and groups can also have economic controls. Wells
under group control are therefore subject to the economic criteria set
via the GCONINJE keyword in the SCHEDULE section and the controls
specified by this keyword.

This keyword is not supported by OPM Flow but would change the results
if supported so the simulation will be stopped.

\begin{longtable}[]{@{}
  >{\raggedright\arraybackslash}p{(\linewidth - 6\tabcolsep) * \real{0.2500}}
  >{\raggedright\arraybackslash}p{(\linewidth - 6\tabcolsep) * \real{0.2500}}
  >{\raggedright\arraybackslash}p{(\linewidth - 6\tabcolsep) * \real{0.2500}}
  >{\raggedright\arraybackslash}p{(\linewidth - 6\tabcolsep) * \real{0.2500}}@{}}
\toprule\noalign{}
\begin{minipage}[b]{\linewidth}\raggedright
No.
\end{minipage} & \begin{minipage}[b]{\linewidth}\raggedright
Name
\end{minipage} & \begin{minipage}[b]{\linewidth}\raggedright
Description
\end{minipage} & \begin{minipage}[b]{\linewidth}\raggedright
Default
\end{minipage} \\
\midrule\noalign{}
\endhead
\bottomrule\noalign{}
\endlastfoot
Field & Metric & Laboratory & \\
1 & WELNAME & A character string of up to eight characters in length
that defines the well name for which the well economic injection
criteria data is being defined. Note that the well name (WELNAME) must
have been declared previously using the WELSPECS keyword in the SCHEDULE
section, otherwise an error may occur. & None \\
2 & MINVALUE & A real positive value that defines the minimum economic
injection value, below which an economic action will take place, as
defined by the AUTO parameter on the WELSPECS keyword (SHUT or STOP).
Note that TYPE determines if the minimum value is applied to the well's
actual injection rate or the well's potential. A value less than or
equal to zero switches off this criterion. & 0.0 \\
Liquid stb/d Gas Mscf/d & Liquid sm3/day Gas sm3/day & Liquid scc/hour
Gas scc/hour & \\
3 & TYPE & A defined character string that determines if MINVALUE is
applied to a well's actual rate or potential, and should be set to one
of the following: The default value is RATE. & RATE \\
Notes: & & & \\
\end{longtable}

\emph{Table 12.82: WECONINJ Keyword Description}

See also the WELSPECS keyword to define a wells shut-in or stop options
and GCONINJE for group controls in the SCHEDULE section.

\subsubsection{Example}\label{example-566}

The following example defines the economic injection parameters for all
gas and water injection wells.

\begin{verbatim}
--
--       WELL ECONOMIC LIMIT DATA FOR INJECTION WELLS
--
-- WELL  MIN    RATE
-- NAME  VALUE  POTN
WECONINJ
GI*      2.0E3  RATE                                                           /
WI*      5.0E3  POTN                                                           /
/
\end{verbatim}

Here all the gas injection wells have a minimum economic gas injection
rate of 2 MMscf/d and the water injection wells have a minimum water
potential rate of 5,000 stb/d.~The AUTO parameter on the WELSPECS
keyword will determine if the wells will be shut-in or stopped.

\subsection{WECONT -- Well Economic Tracer Criteria for Production
Wells}\label{wecont-well-economic-tracer-criteria-for-production-wells}

\begin{longtable}[]{@{}
  >{\raggedright\arraybackslash}p{(\linewidth - 14\tabcolsep) * \real{0.1250}}
  >{\raggedright\arraybackslash}p{(\linewidth - 14\tabcolsep) * \real{0.1250}}
  >{\raggedright\arraybackslash}p{(\linewidth - 14\tabcolsep) * \real{0.1250}}
  >{\raggedright\arraybackslash}p{(\linewidth - 14\tabcolsep) * \real{0.1250}}
  >{\raggedright\arraybackslash}p{(\linewidth - 14\tabcolsep) * \real{0.1250}}
  >{\raggedright\arraybackslash}p{(\linewidth - 14\tabcolsep) * \real{0.1250}}
  >{\raggedright\arraybackslash}p{(\linewidth - 14\tabcolsep) * \real{0.1250}}
  >{\raggedright\arraybackslash}p{(\linewidth - 14\tabcolsep) * \real{0.1250}}@{}}
\toprule\noalign{}
\begin{minipage}[b]{\linewidth}\raggedright
RUNSPEC
\end{minipage} & \begin{minipage}[b]{\linewidth}\raggedright
GRID
\end{minipage} & \begin{minipage}[b]{\linewidth}\raggedright
EDIT
\end{minipage} & \begin{minipage}[b]{\linewidth}\raggedright
PROPS
\end{minipage} & \begin{minipage}[b]{\linewidth}\raggedright
REGIONS
\end{minipage} & \begin{minipage}[b]{\linewidth}\raggedright
SOLUTION
\end{minipage} & \begin{minipage}[b]{\linewidth}\raggedright
SUMMARY
\end{minipage} & \begin{minipage}[b]{\linewidth}\raggedright
SCHEDULE
\end{minipage} \\
\midrule\noalign{}
\endhead
\bottomrule\noalign{}
\endlastfoot
\end{longtable}

\subsubsection{Description}\label{description-1057}

The WECONT keyword defines the tracer economic criteria for production
wells that have previously been defined by the WELSPECS and WCONPROD
keywords in the SCHEDULE section, for tracers define by the TRACER
keyword in the PROPS section.

Note that wells can be allocated to a group when they are specified by
the WELSPECS keyword and groups can also have economic controls. Wells
under group control are therefore subject to the economic criteria set
via the GECONT keyword in the SCHEDULE section and the controls
specified by this keyword. Note that GECONT is not supported by OPM Flow
in the current release

This keyword is not supported by OPM Flow but would change the results
if supported so the simulation will be stopped.

\begin{longtable}[]{@{}
  >{\raggedright\arraybackslash}p{(\linewidth - 6\tabcolsep) * \real{0.2500}}
  >{\raggedright\arraybackslash}p{(\linewidth - 6\tabcolsep) * \real{0.2500}}
  >{\raggedright\arraybackslash}p{(\linewidth - 6\tabcolsep) * \real{0.2500}}
  >{\raggedright\arraybackslash}p{(\linewidth - 6\tabcolsep) * \real{0.2500}}@{}}
\toprule\noalign{}
\begin{minipage}[b]{\linewidth}\raggedright
No.
\end{minipage} & \begin{minipage}[b]{\linewidth}\raggedright
Name
\end{minipage} & \begin{minipage}[b]{\linewidth}\raggedright
Description
\end{minipage} & \begin{minipage}[b]{\linewidth}\raggedright
Default
\end{minipage} \\
\midrule\noalign{}
\endhead
\bottomrule\noalign{}
\endlastfoot
Field & Metric & Laboratory & \\
1-1 & WELNAME & A character string of up to eight characters in length
that defines the well name for which the well target and constraints are
being defined. & None \\
1-2 & ACTION & A defined character string that defines the action to be
taken if the economic WCUT, GOR, or WGR limits are violated. ACTION
should be set to one of the following character strings: The corrective
action takes places at the end of the time step in which the constraint
is violated. & None \\
1-3 & END & A defined character string that defines if the simulation
should terminate if the well is shut or stopped. END should be set to
one of the following character strings: & NO \\
1-4 & WELL & A character string of up to eight characters in length that
defines the well name of a fully defined well that will be ``opened''
when the well WELNAME is shut-in or stopped. & None \\
1-5 & / & Record one terminated by a ``/'' & Not Applicable \\
2-1 & NAME & A three letter character string defining the tracer's name.
Note it is best to void names beginning with the letters F, S, and T as
these names may create naming issues in post-processing software. &
None \\
2-2 & MXTOTAL & A real positive value that defines the maximum total
(free plus solution) tracer rate. Tracer units are per those defined by
the carrying fluid, oil, gas, water, etc. & None \\
2-3 & MXFREET & A real positive value that defines the maximum total
(free plus solution) tracer concentration. Tracer units are per those
defined by the carrying fluid, oil, gas, water, etc. & None \\
2-4 & MXFREEQ & A real positive value that defines the maximum free
tracer rate. Tracer units are per those defined by the carrying fluid,
oil, gas, water, etc. & None \\
2-5 & MXCONC & A real positive value that defines the maximum free
tracer concentration. Tracer units are per those defined by the carrying
fluid, oil, gas, water, etc. & None \\
2-6 & MXSOLNQ & A real positive value that defines the maximum solution
rate. Tracer units are per those defined by the carrying fluid, oil,
gas, water, etc. & None \\
2-7 & MXSOLNC & A real positive value that defines the maximum solution
concentration. Tracer units are per those defined by the carrying fluid,
oil, gas, water, etc. & \\
2-8 & / & Record two terminated by a ``/'' & Not Applicable \\
3-1 & / & Well terminated by a ``/'' & Not Applicable \\
Notes: & & & \\
\end{longtable}

\emph{Table 12.83: WECONT Keyword Description}

See also the WELSPECS keyword to define a wells shut-in or stop options,
and WECON for setting a well's economic criteria. Both the
aforementioned keywords are described in the SCHEDULE section.

\subsubsection{Example}\label{example-567}

The following example defines the tracer economic criteria for the field
and two wells, OP01 and OP02.

\begin{verbatim}
--
--       WELL TRACER ECONOMIC CRITERIA FOR PRODUCTION WELLS
--
-- WELL  WORK    END    MAX
-- NAME  OVER    RUN    WELLS
WECONT
OP01     +CON   'YES'   1*                                     / START OF WELL
--
--       TRACER  TRACER  TRACER  TRACER  TRACER  TRACER  TRACER
--       NAME    TOTAL   TOTAL   FREE    FREE    SOLN    SOLN
--               RATE    CONCEN  RATE    CONCEN  RATE    CONCEN
         PLY     800.0                                         /
         BRI     800.0                                         /
                                                                      /
OP02     +CON   'YES'   1*                                     / START OF WELL
--
--       TRACER  TRACER  TRACER  TRACER  TRACER  TRACER  TRACER
--       NAME    TOTAL   TOTAL   FREE    FREE    SOLN    SOLN
--               RATE    CONCEN  RATE    CONCEN  RATE    CONCEN
         PLY     800.0                                         /
         BRI     800.0                                         / END OF WELL
                                                               / END OF KEYWORD
\end{verbatim}

If the economic limits are violated then the worst offending connection
and all below it in the worst offending well will be closed, If
connections have been grouped as completions then the worst offending
completion and all below it in the worst offending well will be closed

\subsection{WEFAC -- Define Well
Efficiency}\label{wefac-define-well-efficiency}

\begin{longtable}[]{@{}
  >{\raggedright\arraybackslash}p{(\linewidth - 14\tabcolsep) * \real{0.1250}}
  >{\raggedright\arraybackslash}p{(\linewidth - 14\tabcolsep) * \real{0.1250}}
  >{\raggedright\arraybackslash}p{(\linewidth - 14\tabcolsep) * \real{0.1250}}
  >{\raggedright\arraybackslash}p{(\linewidth - 14\tabcolsep) * \real{0.1250}}
  >{\raggedright\arraybackslash}p{(\linewidth - 14\tabcolsep) * \real{0.1250}}
  >{\raggedright\arraybackslash}p{(\linewidth - 14\tabcolsep) * \real{0.1250}}
  >{\raggedright\arraybackslash}p{(\linewidth - 14\tabcolsep) * \real{0.1250}}
  >{\raggedright\arraybackslash}p{(\linewidth - 14\tabcolsep) * \real{0.1250}}@{}}
\toprule\noalign{}
\begin{minipage}[b]{\linewidth}\raggedright
RUNSPEC
\end{minipage} & \begin{minipage}[b]{\linewidth}\raggedright
GRID
\end{minipage} & \begin{minipage}[b]{\linewidth}\raggedright
EDIT
\end{minipage} & \begin{minipage}[b]{\linewidth}\raggedright
PROPS
\end{minipage} & \begin{minipage}[b]{\linewidth}\raggedright
REGIONS
\end{minipage} & \begin{minipage}[b]{\linewidth}\raggedright
SOLUTION
\end{minipage} & \begin{minipage}[b]{\linewidth}\raggedright
SUMMARY
\end{minipage} & \begin{minipage}[b]{\linewidth}\raggedright
SCHEDULE
\end{minipage} \\
\midrule\noalign{}
\endhead
\bottomrule\noalign{}
\endlastfoot
\end{longtable}

\subsubsection{Description}\label{description-1058}

Defines a well's efficiency or up-time as opposed to setting the
efficiency at the group level.

Note that wells are allocated to a group when they are specified by the
WELSPECS keyword and groups can also have efficiency factors.

\begin{longtable}[]{@{}
  >{\raggedright\arraybackslash}p{(\linewidth - 6\tabcolsep) * \real{0.2500}}
  >{\raggedright\arraybackslash}p{(\linewidth - 6\tabcolsep) * \real{0.2500}}
  >{\raggedright\arraybackslash}p{(\linewidth - 6\tabcolsep) * \real{0.2500}}
  >{\raggedright\arraybackslash}p{(\linewidth - 6\tabcolsep) * \real{0.2500}}@{}}
\toprule\noalign{}
\begin{minipage}[b]{\linewidth}\raggedright
No.
\end{minipage} & \begin{minipage}[b]{\linewidth}\raggedright
Name
\end{minipage} & \begin{minipage}[b]{\linewidth}\raggedright
Description
\end{minipage} & \begin{minipage}[b]{\linewidth}\raggedright
Default
\end{minipage} \\
\midrule\noalign{}
\endhead
\bottomrule\noalign{}
\endlastfoot
Field & Metric & Laboratory & \\
1 & WELNAME & A character string of up to eight characters in length
that defines the well name for which the well efficiency factor is being
defined. Note that the well name (WELNAME) must have been declared
previously using the WELSPECS keyword in the SCHEDULE section, otherwise
an error may occur. & None \\
2 & FACTOR & A real positive value greater than zero and less than or
equal to one that defines the efficiency factor for the well. If a
well's down time is 5\% then FACTOR should be set to 0.95 (1.0 -- 0.05).
Note that well pressures and rates are calculated at their full flowing
conditions but subject to any operating constraints, that is without the
well efficiency being applied (FACTOR), in order to represent the actual
flowing conditions in the field. The effective rates and volumes are
calculated by applying FACTOR when summing individual well rates to
their group level and higher, including summing to the top most group
FIELD. In terms of a well's cumulative production, FACTOR is applied to
the well rate times the time interval for the time step. This ensures
that correct effective volume is withdrawn from (or injected to) the
reservoir. This approach means that wells are effectively arbitrarily
offline for a period during a time step, as opposed to all wells going
offline concurrently. And thus the group and field rates and volumes are
the effective rates and volumes for the field. & 1.0 \\
dimensionless & dimensionless & dimensionless & \\
3 & WELNETWK & A defined character string that determines if the WELNAME
efficiency factor should be applied when calculating the flows and
pressure losses in the Extended Network Model, and should be set to
either: This option is only applicable for the Extended Network Model,
as in the Standard Network Model groups flow rates are always used in
the calculation of pressure drops (this is equivalent to the default
option, YES). & YES \\
Notes: & & & \\
\end{longtable}

\emph{Table 12.84: WEFAC Keyword Description}

See also the GEFAC keyword in the SCHEDULE section to set the efficiency
at the group level, as opposed to applying the efficiency to individual
wells.

\begin{tcolorbox}[enhanced jigsaw, title=\textcolor{quarto-callout-note-color}{\faInfo}\hspace{0.5em}{Note}, coltitle=black, breakable, toprule=.15mm, colbacktitle=quarto-callout-note-color!10!white, leftrule=.75mm, opacitybacktitle=0.6, colback=white, bottomtitle=1mm, toptitle=1mm, titlerule=0mm, left=2mm, colframe=quarto-callout-note-color-frame, arc=.35mm, rightrule=.15mm, bottomrule=.15mm, opacityback=0]

One can also apply plant efficiencies through the GEFAC keyword in the
SCHEDULE section. If all the wells in a group are flowing through a
facility that has an overall efficiency factor, then it is more
appropriate to apply the efficiency factor at the group level. This of
course does not preclude applying additional well efficiencies to
individual wells. For example, subsea wells (wet trees) may have
additional down time compared to platform wells (dry trees) even though
both sets of well are flowing through the same platform. Another example
would be gas lift wells and wells using electrical submersible pumps, as
their artificial lift mechanisms.

\end{tcolorbox}

\subsubsection{Example}\label{example-568}

\begin{verbatim}
--
--       WELL EFFICIENCY FACTORS
--
-- WELL  EFF    NETWK
-- NAME  FACT   OPTN
WEFAC
'GP*  '  0.950                                                                 /
'OP*  '  0.862                                                                 /
/
\end{verbatim}

In the above example the all the gas wells are are defined as having a
well efficiency factor (up time) of 0.950 and all the oil wells have a
lower efficiency factor of 0.862.

\subsection{WELCNTL -- Modify Well Control and
Targets}\label{welcntl-modify-well-control-and-targets}

\begin{longtable}[]{@{}
  >{\raggedright\arraybackslash}p{(\linewidth - 14\tabcolsep) * \real{0.1250}}
  >{\raggedright\arraybackslash}p{(\linewidth - 14\tabcolsep) * \real{0.1250}}
  >{\raggedright\arraybackslash}p{(\linewidth - 14\tabcolsep) * \real{0.1250}}
  >{\raggedright\arraybackslash}p{(\linewidth - 14\tabcolsep) * \real{0.1250}}
  >{\raggedright\arraybackslash}p{(\linewidth - 14\tabcolsep) * \real{0.1250}}
  >{\raggedright\arraybackslash}p{(\linewidth - 14\tabcolsep) * \real{0.1250}}
  >{\raggedright\arraybackslash}p{(\linewidth - 14\tabcolsep) * \real{0.1250}}
  >{\raggedright\arraybackslash}p{(\linewidth - 14\tabcolsep) * \real{0.1250}}@{}}
\toprule\noalign{}
\begin{minipage}[b]{\linewidth}\raggedright
RUNSPEC
\end{minipage} & \begin{minipage}[b]{\linewidth}\raggedright
GRID
\end{minipage} & \begin{minipage}[b]{\linewidth}\raggedright
EDIT
\end{minipage} & \begin{minipage}[b]{\linewidth}\raggedright
PROPS
\end{minipage} & \begin{minipage}[b]{\linewidth}\raggedright
REGIONS
\end{minipage} & \begin{minipage}[b]{\linewidth}\raggedright
SOLUTION
\end{minipage} & \begin{minipage}[b]{\linewidth}\raggedright
SUMMARY
\end{minipage} & \begin{minipage}[b]{\linewidth}\raggedright
SCHEDULE
\end{minipage} \\
\midrule\noalign{}
\endhead
\bottomrule\noalign{}
\endlastfoot
\end{longtable}

\subsubsection{Description}\label{description-1059}

The WELCNTL keyword modifies a well's target control and value, both
rates and pressures, for previously defined wells without having to
define all the variables on the well control keywords: WCONPROD,
WCONHIST, WCONINJE, or WCONINJH keywords. Variables not changed by the
WELCNTL keyword remain the same as those previously entered via the well
control keywords or previously entered WELCNTL keywords. Note that the
well must still be initially be fully defined using the WCONPROD or
WCONINJE keywords. All the aforementioned keywords are described in the
SCHEDULE section.

This keyword is not supported by OPM Flow but would change the results
if supported so the simulation will be stopped.

\begin{longtable}[]{@{}
  >{\raggedright\arraybackslash}p{(\linewidth - 6\tabcolsep) * \real{0.2500}}
  >{\raggedright\arraybackslash}p{(\linewidth - 6\tabcolsep) * \real{0.2500}}
  >{\raggedright\arraybackslash}p{(\linewidth - 6\tabcolsep) * \real{0.2500}}
  >{\raggedright\arraybackslash}p{(\linewidth - 6\tabcolsep) * \real{0.2500}}@{}}
\toprule\noalign{}
\begin{minipage}[b]{\linewidth}\raggedright
No.
\end{minipage} & \begin{minipage}[b]{\linewidth}\raggedright
Name
\end{minipage} & \begin{minipage}[b]{\linewidth}\raggedright
Description
\end{minipage} & \begin{minipage}[b]{\linewidth}\raggedright
Default
\end{minipage} \\
\midrule\noalign{}
\endhead
\bottomrule\noalign{}
\endlastfoot
Field & Metric & Laboratory & \\
1 & WELNAME & A character string of up to eight characters in length
that defines the well name for which the well production rates and
pressures data are being redefined. Note that the well name (WELNAME)
must have been declared previously using the WELSPECS and WCONPROD (or
WCONINJE) keywords in the SCHEDULE section, otherwise an error may
occur. & None \\
2 & TARGET & A defined character string that sets the item to be changed
for the well the value of the item is set by item (3). & None \\
3 & VALUE Liquid Gas Res Vol Pressure VFP LIFT & A real positive value
that defines the value of the variable declared by TARGET & None \\
stb/d Mscf/d rb/d psia dimensionless same as VFPPROD or VFPINJ & sm3/day
sm3/day rm3/day barsa dimensionless same as VFPPROD or VFPINJ & scc/hour
scc/hour rcc/hour atma dimensionless same as VFPPROD or VFPINJ & \\
Notes: & & & \\
\end{longtable}

\emph{Table 12.85: WELCNTL Keyword Description}

If a well is currently a history matching well, then WELCNTL can be used
to change the well to a standard well. See also the WELTARG keyword, in
the SCHEDULE section that can be used to reset a well's target and
constraints of both rates and pressures.

\subsubsection{Example}\label{example-569}

The following example below shows the oil rates for the OP01 oil
producer at the start of the schedule section (January 1, 2000).

\begin{verbatim}
-- ------------------------------------------------------------------------------
-- 01 JAN 2000 START OF SCHEDULE SECTION
-- ------------------------------------------------------------------------------
--
--       WELL PRODUCTION WELL CONTROLS
--
-- WELL  OPEN/  CNTL   OIL    WAT    GAS   LIQ    RES    BHP    THP   VFP    VFP
-- NAME  SHUT   MODE   RATE   RATE   RATE  RATE   RATE   PRES   PRES  TABLE  ALFQ
WCONPROD
OP01     OPEN   ORAT   3000   1*     1*    1*     1*     750.0  500.  9      1* /
/                                                                               DATES
01 FEB 2000 /
/
--
--       WELL CONTROL MODE AND OPERATING TARGET
--
--  WELL WELL   TARGET
--  NAME CNTL   VALUE                                                        WELCNTL
OP01     LRAT   5000                                        /
/
\end{verbatim}

From January 1, 2000 to February 1, 2000 well OP01 is open and is on oil
rate control and has a target oil rate of 3,000 stb/d and uses VFPPROD
vertical lift table number 9 with a minimum tubing head pressure
constraint of 500 psia. After February 1, 2000 the well is changed to
liquid control with a target rate of 5,000 stb/d of liquid and all the
other parameters remain unchanged.

\subsection{WELDEBUG -- Define the Well Debug Data to be Printed to
File}\label{weldebug-define-the-well-debug-data-to-be-printed-to-file}

\begin{longtable}[]{@{}
  >{\raggedright\arraybackslash}p{(\linewidth - 14\tabcolsep) * \real{0.1250}}
  >{\raggedright\arraybackslash}p{(\linewidth - 14\tabcolsep) * \real{0.1250}}
  >{\raggedright\arraybackslash}p{(\linewidth - 14\tabcolsep) * \real{0.1250}}
  >{\raggedright\arraybackslash}p{(\linewidth - 14\tabcolsep) * \real{0.1250}}
  >{\raggedright\arraybackslash}p{(\linewidth - 14\tabcolsep) * \real{0.1250}}
  >{\raggedright\arraybackslash}p{(\linewidth - 14\tabcolsep) * \real{0.1250}}
  >{\raggedright\arraybackslash}p{(\linewidth - 14\tabcolsep) * \real{0.1250}}
  >{\raggedright\arraybackslash}p{(\linewidth - 14\tabcolsep) * \real{0.1250}}@{}}
\toprule\noalign{}
\begin{minipage}[b]{\linewidth}\raggedright
RUNSPEC
\end{minipage} & \begin{minipage}[b]{\linewidth}\raggedright
GRID
\end{minipage} & \begin{minipage}[b]{\linewidth}\raggedright
EDIT
\end{minipage} & \begin{minipage}[b]{\linewidth}\raggedright
PROPS
\end{minipage} & \begin{minipage}[b]{\linewidth}\raggedright
REGIONS
\end{minipage} & \begin{minipage}[b]{\linewidth}\raggedright
SOLUTION
\end{minipage} & \begin{minipage}[b]{\linewidth}\raggedright
SUMMARY
\end{minipage} & \begin{minipage}[b]{\linewidth}\raggedright
SCHEDULE
\end{minipage} \\
\midrule\noalign{}
\endhead
\bottomrule\noalign{}
\endlastfoot
\end{longtable}

\subsubsection{Description}\label{description-1060}

This keyword defines the well debug data to be written to the debug file
(*.DBG). This keyword is not supported by OPM Flow but would change the
results if supported so the simulation will be stopped.

\subsection{WELDRAW -- Define Maximum Draw Down for Production
Wells}\label{weldraw-define-maximum-draw-down-for-production-wells}

\begin{longtable}[]{@{}
  >{\raggedright\arraybackslash}p{(\linewidth - 14\tabcolsep) * \real{0.1250}}
  >{\raggedright\arraybackslash}p{(\linewidth - 14\tabcolsep) * \real{0.1250}}
  >{\raggedright\arraybackslash}p{(\linewidth - 14\tabcolsep) * \real{0.1250}}
  >{\raggedright\arraybackslash}p{(\linewidth - 14\tabcolsep) * \real{0.1250}}
  >{\raggedright\arraybackslash}p{(\linewidth - 14\tabcolsep) * \real{0.1250}}
  >{\raggedright\arraybackslash}p{(\linewidth - 14\tabcolsep) * \real{0.1250}}
  >{\raggedright\arraybackslash}p{(\linewidth - 14\tabcolsep) * \real{0.1250}}
  >{\raggedright\arraybackslash}p{(\linewidth - 14\tabcolsep) * \real{0.1250}}@{}}
\toprule\noalign{}
\begin{minipage}[b]{\linewidth}\raggedright
RUNSPEC
\end{minipage} & \begin{minipage}[b]{\linewidth}\raggedright
GRID
\end{minipage} & \begin{minipage}[b]{\linewidth}\raggedright
EDIT
\end{minipage} & \begin{minipage}[b]{\linewidth}\raggedright
PROPS
\end{minipage} & \begin{minipage}[b]{\linewidth}\raggedright
REGIONS
\end{minipage} & \begin{minipage}[b]{\linewidth}\raggedright
SOLUTION
\end{minipage} & \begin{minipage}[b]{\linewidth}\raggedright
SUMMARY
\end{minipage} & \begin{minipage}[b]{\linewidth}\raggedright
SCHEDULE
\end{minipage} \\
\midrule\noalign{}
\endhead
\bottomrule\noalign{}
\endlastfoot
\end{longtable}

\subsubsection{Description}\label{description-1061}

This keyword, WELDRAW, defines the maximum draw down for production
wells The keyword may be useful in wells that are subject to fines or
sand production to limit the draw down between the sand face and the
well in order to limit or avoid sand production.

This keyword is not supported by OPM Flow but would change the results
if supported so the simulation will be stopped.

\subsection{WELEVNT -- Define Well WPWEM Summary Mnemonic Output
Value}\label{welevnt-define-well-wpwem-summary-mnemonic-output-value}

\begin{longtable}[]{@{}
  >{\raggedright\arraybackslash}p{(\linewidth - 14\tabcolsep) * \real{0.1250}}
  >{\raggedright\arraybackslash}p{(\linewidth - 14\tabcolsep) * \real{0.1250}}
  >{\raggedright\arraybackslash}p{(\linewidth - 14\tabcolsep) * \real{0.1250}}
  >{\raggedright\arraybackslash}p{(\linewidth - 14\tabcolsep) * \real{0.1250}}
  >{\raggedright\arraybackslash}p{(\linewidth - 14\tabcolsep) * \real{0.1250}}
  >{\raggedright\arraybackslash}p{(\linewidth - 14\tabcolsep) * \real{0.1250}}
  >{\raggedright\arraybackslash}p{(\linewidth - 14\tabcolsep) * \real{0.1250}}
  >{\raggedright\arraybackslash}p{(\linewidth - 14\tabcolsep) * \real{0.1250}}@{}}
\toprule\noalign{}
\begin{minipage}[b]{\linewidth}\raggedright
RUNSPEC
\end{minipage} & \begin{minipage}[b]{\linewidth}\raggedright
GRID
\end{minipage} & \begin{minipage}[b]{\linewidth}\raggedright
EDIT
\end{minipage} & \begin{minipage}[b]{\linewidth}\raggedright
PROPS
\end{minipage} & \begin{minipage}[b]{\linewidth}\raggedright
REGIONS
\end{minipage} & \begin{minipage}[b]{\linewidth}\raggedright
SOLUTION
\end{minipage} & \begin{minipage}[b]{\linewidth}\raggedright
SUMMARY
\end{minipage} & \begin{minipage}[b]{\linewidth}\raggedright
SCHEDULE
\end{minipage} \\
\midrule\noalign{}
\endhead
\bottomrule\noalign{}
\endlastfoot
\end{longtable}

\subsubsection{Description}\label{description-1062}

This keyword, WELEVNT, defines an integer value to be assigned to an
individual well's WPWEM summary variable that is written to the SUMMARY
file. The value is set to zero after the current time step.

This keyword is not supported by OPM Flow but would change the results
if supported so the simulation will be stopped.

\subsection{WELLSTRE -- Define Injection Stream
Composition}\label{wellstre-define-injection-stream-composition}

\begin{longtable}[]{@{}
  >{\raggedright\arraybackslash}p{(\linewidth - 14\tabcolsep) * \real{0.1250}}
  >{\raggedright\arraybackslash}p{(\linewidth - 14\tabcolsep) * \real{0.1250}}
  >{\raggedright\arraybackslash}p{(\linewidth - 14\tabcolsep) * \real{0.1250}}
  >{\raggedright\arraybackslash}p{(\linewidth - 14\tabcolsep) * \real{0.1250}}
  >{\raggedright\arraybackslash}p{(\linewidth - 14\tabcolsep) * \real{0.1250}}
  >{\raggedright\arraybackslash}p{(\linewidth - 14\tabcolsep) * \real{0.1250}}
  >{\raggedright\arraybackslash}p{(\linewidth - 14\tabcolsep) * \real{0.1250}}
  >{\raggedright\arraybackslash}p{(\linewidth - 14\tabcolsep) * \real{0.1250}}@{}}
\toprule\noalign{}
\begin{minipage}[b]{\linewidth}\raggedright
RUNSPEC
\end{minipage} & \begin{minipage}[b]{\linewidth}\raggedright
GRID
\end{minipage} & \begin{minipage}[b]{\linewidth}\raggedright
EDIT
\end{minipage} & \begin{minipage}[b]{\linewidth}\raggedright
PROPS
\end{minipage} & \begin{minipage}[b]{\linewidth}\raggedright
REGIONS
\end{minipage} & \begin{minipage}[b]{\linewidth}\raggedright
SOLUTION
\end{minipage} & \begin{minipage}[b]{\linewidth}\raggedright
SUMMARY
\end{minipage} & \begin{minipage}[b]{\linewidth}\raggedright
SCHEDULE
\end{minipage} \\
\midrule\noalign{}
\endhead
\bottomrule\noalign{}
\endlastfoot
\end{longtable}

\subsubsection{Description}\label{description-1063}

The WELLSTRE keyword defines a well stream together with the
compositional component mole factions, associated with the well stream,
as such it should have the same number of entries as that declared via
the COMPS keyword in the RUNSPEC section, and the NCOMPS keyword in the
PROPS section. Once a gas well stream has been defined, it can be used
with either the WINJGAS or GINJGAS keywords in the SCHEDULE section, to
set the injected gas composition. Similarly, if an oil well stream has
been defined by WELLSTRE, then the well stream can be used with the
WINJOIL keyword in the SCHEDULE section, to specify the injected oil
composition.

The keyword should only be used if the CO2STORE and GASWAT keywords in
the RUNSPEC section have also be activated for the gas-water two
component model.

\begin{tcolorbox}[enhanced jigsaw, title=\textcolor{quarto-callout-note-color}{\faInfo}\hspace{0.5em}{Note}, coltitle=black, breakable, toprule=.15mm, colbacktitle=quarto-callout-note-color!10!white, leftrule=.75mm, opacitybacktitle=0.6, colback=white, bottomtitle=1mm, toptitle=1mm, titlerule=0mm, left=2mm, colframe=quarto-callout-note-color-frame, arc=.35mm, rightrule=.15mm, bottomrule=.15mm, opacityback=0]

This is an OPM Flow keyword used with OPM Flow's CO2STORE and GASWAT
keywords in the RUNSPEC section, and should not be confused with the
more general version of the WELLSTRE keyword used in the commercial
compositional simulator. Secondly, although OPM Flow parses the keyword,
the simulator currently ignores the data for this keyword.

\end{tcolorbox}

\begin{longtable}[]{@{}
  >{\raggedright\arraybackslash}p{(\linewidth - 6\tabcolsep) * \real{0.2500}}
  >{\raggedright\arraybackslash}p{(\linewidth - 6\tabcolsep) * \real{0.2500}}
  >{\raggedright\arraybackslash}p{(\linewidth - 6\tabcolsep) * \real{0.2500}}
  >{\raggedright\arraybackslash}p{(\linewidth - 6\tabcolsep) * \real{0.2500}}@{}}
\toprule\noalign{}
\begin{minipage}[b]{\linewidth}\raggedright
No.
\end{minipage} & \begin{minipage}[b]{\linewidth}\raggedright
Name
\end{minipage} & \begin{minipage}[b]{\linewidth}\raggedright
Description
\end{minipage} & \begin{minipage}[b]{\linewidth}\raggedright
Default
\end{minipage} \\
\midrule\noalign{}
\endhead
\bottomrule\noalign{}
\endlastfoot
Field & Metric & Laboratory & \\
1 & STREAM & STREAM is a character string of up to eight characters in
length, representing the well stream name. The WELLDIMS(MXSTRMS)
parameter in the RUNSPEC section determines the maximum number of well
streams allowed in the model. & None \\
2 & ZCOMP & A row vector, with each item representing a compositional
component mole fraction for a given component. In addition, the sum of
the compositional component mole fractions must sum to one, otherwise an
error will occur. Note that the number ZCOMP values, should be the same
as that entered via the NCOMPS keyword in the PROPS section, and the
COMPS keyword in the RUNSPEC section. However, for a given ZCOMP
component, the mole fraction may be defaulted with 1*, in which case the
mole fraction is set to zero. Secondly, ZCOMP may be terminated early,
in this case the undefined mole fractions will be set to zero. Finally,
only the default value of two components are currently supported by OPM
Flow. & None \\
mole fraction & mole fraction & mole fraction & \\
Notes: & & & \\
\end{longtable}

\emph{Table 12.86: WELLSTRE Keyword Description}

\subsubsection{Example}\label{example-570}

The following example defines how to specify a two component
formulation, together with defining the names of the composition
components, to be used with the CO2STORE and GASWAT options.

\begin{verbatim}
-- ==============================================================================
--
-- PROPS SECTION
--
-- ==============================================================================
PROPS             --
--       CONFIRM NUMBER OF COMPOSITIONAL COMPONENTS (OPM FLOW KEYWORD)
--
NCOMPS
                2                                                                      /
--
--       DEFINE COMPOSITIONAL COMPONENTS NAMES (OPM FLOW KEYWORD)
--
CNAMES
         'CO2'
         'H2O'                                                                 /
\end{verbatim}

The second part of the example, defines the well stream for the above
two component CO2 water system.

\begin{verbatim}
-- ==============================================================================
--
-- SCHEDULE SECTION
--
-- ==============================================================================
SCHEDULE
--
--       WELL STREAM INJECTION COMPOSITION (OPM FLOW Keyword)
--
-- WELL     -- WELL STREAM COMPOSITIONAL COMPONENT       --
-- STREAM   --          MOLE FRACTIONS                   --
WELLSTRE
'C02STREAM'    1.000   0.000                               /
/
\end{verbatim}

Here the well stream consists of 100\% CO2 and zero water.

\subsection{WELMOVEL -- Move Global Well Into an
LGR}\label{welmovel-move-global-well-into-an-lgr}

\begin{longtable}[]{@{}
  >{\raggedright\arraybackslash}p{(\linewidth - 14\tabcolsep) * \real{0.1250}}
  >{\raggedright\arraybackslash}p{(\linewidth - 14\tabcolsep) * \real{0.1250}}
  >{\raggedright\arraybackslash}p{(\linewidth - 14\tabcolsep) * \real{0.1250}}
  >{\raggedright\arraybackslash}p{(\linewidth - 14\tabcolsep) * \real{0.1250}}
  >{\raggedright\arraybackslash}p{(\linewidth - 14\tabcolsep) * \real{0.1250}}
  >{\raggedright\arraybackslash}p{(\linewidth - 14\tabcolsep) * \real{0.1250}}
  >{\raggedright\arraybackslash}p{(\linewidth - 14\tabcolsep) * \real{0.1250}}
  >{\raggedright\arraybackslash}p{(\linewidth - 14\tabcolsep) * \real{0.1250}}@{}}
\toprule\noalign{}
\begin{minipage}[b]{\linewidth}\raggedright
RUNSPEC
\end{minipage} & \begin{minipage}[b]{\linewidth}\raggedright
GRID
\end{minipage} & \begin{minipage}[b]{\linewidth}\raggedright
EDIT
\end{minipage} & \begin{minipage}[b]{\linewidth}\raggedright
PROPS
\end{minipage} & \begin{minipage}[b]{\linewidth}\raggedright
REGIONS
\end{minipage} & \begin{minipage}[b]{\linewidth}\raggedright
SOLUTION
\end{minipage} & \begin{minipage}[b]{\linewidth}\raggedright
SUMMARY
\end{minipage} & \begin{minipage}[b]{\linewidth}\raggedright
SCHEDULE
\end{minipage} \\
\midrule\noalign{}
\endhead
\bottomrule\noalign{}
\endlastfoot
\end{longtable}

\subsubsection{Description}\label{description-1064}

This keyword, WELMOVEL, moves a previously defined gloabal well into a
previously declared Local Grid Refinement (``LGR''), in a RESTART run.
The keyword should only be used in RESTART runs.

This keyword is not supported by OPM Flow but would change the results
if supported so the simulation will be stopped.

\subsection{WELOPEN -- Define Well and Well Connection Flowing
Status}\label{welopen-define-well-and-well-connection-flowing-status}

\begin{longtable}[]{@{}
  >{\raggedright\arraybackslash}p{(\linewidth - 14\tabcolsep) * \real{0.1250}}
  >{\raggedright\arraybackslash}p{(\linewidth - 14\tabcolsep) * \real{0.1250}}
  >{\raggedright\arraybackslash}p{(\linewidth - 14\tabcolsep) * \real{0.1250}}
  >{\raggedright\arraybackslash}p{(\linewidth - 14\tabcolsep) * \real{0.1250}}
  >{\raggedright\arraybackslash}p{(\linewidth - 14\tabcolsep) * \real{0.1250}}
  >{\raggedright\arraybackslash}p{(\linewidth - 14\tabcolsep) * \real{0.1250}}
  >{\raggedright\arraybackslash}p{(\linewidth - 14\tabcolsep) * \real{0.1250}}
  >{\raggedright\arraybackslash}p{(\linewidth - 14\tabcolsep) * \real{0.1250}}@{}}
\toprule\noalign{}
\begin{minipage}[b]{\linewidth}\raggedright
RUNSPEC
\end{minipage} & \begin{minipage}[b]{\linewidth}\raggedright
GRID
\end{minipage} & \begin{minipage}[b]{\linewidth}\raggedright
EDIT
\end{minipage} & \begin{minipage}[b]{\linewidth}\raggedright
PROPS
\end{minipage} & \begin{minipage}[b]{\linewidth}\raggedright
REGIONS
\end{minipage} & \begin{minipage}[b]{\linewidth}\raggedright
SOLUTION
\end{minipage} & \begin{minipage}[b]{\linewidth}\raggedright
SUMMARY
\end{minipage} & \begin{minipage}[b]{\linewidth}\raggedright
SCHEDULE
\end{minipage} \\
\midrule\noalign{}
\endhead
\bottomrule\noalign{}
\endlastfoot
\end{longtable}

\subsubsection{Description}\label{description-1065}

The WELOPEN keyword defines the status of wells and well connections,
and is used to open and shut previously defined wells and well
connections without having to re-specify all the data on the well
control keywords: WCONPROD, WCONHIST, WCONINJE, or WCONINJH keywords.
Note that the well must still be initially fully defined using the
WCONPROD or WCONINJE keywords. All the aforementioned keywords are
described in the SCHEDULE section.

\begin{longtable}[]{@{}
  >{\raggedright\arraybackslash}p{(\linewidth - 6\tabcolsep) * \real{0.2500}}
  >{\raggedright\arraybackslash}p{(\linewidth - 6\tabcolsep) * \real{0.2500}}
  >{\raggedright\arraybackslash}p{(\linewidth - 6\tabcolsep) * \real{0.2500}}
  >{\raggedright\arraybackslash}p{(\linewidth - 6\tabcolsep) * \real{0.2500}}@{}}
\toprule\noalign{}
\begin{minipage}[b]{\linewidth}\raggedright
No.
\end{minipage} & \begin{minipage}[b]{\linewidth}\raggedright
Name
\end{minipage} & \begin{minipage}[b]{\linewidth}\raggedright
Description
\end{minipage} & \begin{minipage}[b]{\linewidth}\raggedright
Default
\end{minipage} \\
\midrule\noalign{}
\endhead
\bottomrule\noalign{}
\endlastfoot
Field & Metric & Laboratory & \\
1 & WELNAME & A character string of up to eight characters in length
that defines the well name for which the well or well connection status
data are being defined. Note that the well name (WELNAME) must have been
declared previously using the WELSPECS keyword in the SCHEDULE section,
otherwise an error may occur. & None \\
2 & STATUS & A defined character string of length four that defines the
well or well connection operational status, STATUS should be set to one
of the following character strings: & OPEN \\
3 & I & An integer less than or equal to NX that defines the connection
location in the I-direction. & Negative \\
4 & J & An integer less than or equal to NY that defines the connection
location in the J-direction. & Negative \\
5 & K & An integer less than or equal to NZ that defines the connection
location in the K-direction. & Negative \\
6 & C1 & An integer less than or equal to the total number of
completions that defines the first completion in the range. &
Negative \\
7 & C2 & An integer less than or equal to the total number of
completions that defines the last completion in the range. & Negative \\
Notes: & & & \\
\end{longtable}

\emph{Table 12.87: WELOPEN Keyword Description}

If variables I, J, K, C1 and C2 are all set to a negative number or
defaulted with 1* then STATUS is applied to the well and the well
connections remain unchanged.

If any of the variables I, J, K, C1 and C2 are set to zero or a positive
value then STATUS is applied to the defined well connections and the
well status remains unchanged. The defined well connections are those
with the specified I, J, K locations and that are in the completion
number range specified by C1 and C2. If a zero or negative value is
specified for I, J or K then this matches any value for that index. If a
zero or negative value is specified for C1 or C2 respectively then no
lower or upper bound is applied to the completion range.

See also the COMPDAT keyword to define a well's connections, the
COMPLUMP keyword to group well connections into well completions, the
WCONPROD and WCONINJE keywords to define a well's production and
injections targets and constraints. All the aforementioned keywords are
described in the SCHEDULE section.

\subsubsection{Examples}\label{examples-92}

The following example defines two vertical oil wells using the WELSPECS
keyword and their associated connection data.

\begin{verbatim}
-- WELL SPECIFICATION DATA
--
-- WELL     GROUP      LOCATION  BHP    PHASE  DRAIN  INFLOW  SHUT  CROSS  PRESS
-- NAME     NAME        I    J   DEPTH  FLUID  AREA   EQUA.   IN    FLOW   TABLE
WELSPECS
OP01      PLATFORM     14   13   1*      OIL   1*     STD     OPEN   NO    1*  /
OP02      PLATFORM     28   96   1*      OIL   1*     STD     OPEN   NO    1*  /
/
--
--       WELL PRODUCTION WELL CONTROLS
--
-- WELL  OPEN/  CNTL   OIL    WAT    GAS   LIQ    RES    BHP    THP   VFP    VFP
-- NAME  SHUT   MODE   RATE   RATE   RATE  RATE   RATE   PRES   PRES  TABLE  ALFQ
WCONPROD
‘*’      SHUT   GRUP   1*     1*     1*    1*     1*     200.0                  /
/
--       WELL CONNECTION DATA
--
-- WELL  --- LOCATION ---  OPEN   SAT   CONN   WELL   KH    SKIN   D     DIR
-- NAME   II  JJ  K1  K2   SHUT   TAB   FACT   DIA    FACT  FACT   FACT  PEN
COMPDAT
OP01      1*  1*   1  10   SHUT   1*    1*    0.708  1*    0.0    1*    'Z' /
OP01      1*  1*  15  30   SHUT   1*    1*    0.708  1*    0.0    1*    'Z' /
OP01      1*  1*  35  90   SHUT   1*    1*    0.708  1*    0.0    1*    'Z' /
OP02      1*  1*   1  10   SHUT   1*    1*    0.708  1*    0.0    1*    'Z' /
/
--
--       DEFINE WELL AND WELL CONNECTIONS FLOWING STATUS
--
--  WELL WELL   --LOCATION--  COMPLETION
--  NAME STAT     I   J    K  FIRST LAST
WELOPEN
OP01     OPEN                                              /
OP01     OPEN     0   0    1                               /
OP01     OPEN     0   0    2                               /
OP01     OPEN     0   0    3                               /
OP01     OPEN     0   0    4                               /
OP02     OPEN                                              /
OP02     OPEN     0   0    0                               /
/
\end{verbatim}

In this example the first record for each well in the WELOPEN keyword
changes the well status from shut (as per the WCONPROD keyword) to open.
Then for well OP01 well connections in layers one to four are opened for
flow and all the connections for well OP02 are opened, that is the
connections in layers one to ten.

The next example shows the use of the COMPLUMP keyword to group the well
connections into well completions for wells OP01, OP02 and OP03, and
then use the WELOPEN keyword to open the well and the well completions.

\begin{verbatim}
--
--       ASSIGN WELL CONNECTIONS TO COMPLETIONS
--
-- WELL  --- LOCATION ---  COMPL
-- NAME   II  JJ  K1  K2   NO.
COMPLUMP
OP01       0   0   1  10    1                              / COMPLETION NO. 01
OP01       0   0  15  30    2                              / COMPLETION NO. 02
OP01       0   0  35  90    3                              / COMPLETION NO. 03
OP02       0   0  15  30    2                              / COMPLETION NO. 02
OP02       0   0  35  90    3                              / COMPLETION NO. 03
OP03       0   0   1  10    1                              / COMPLETION NO. 01
OP03       0   0  35  90    3                              / COMPLETION NO. 03
/
--
--       DEFINE WELL AND WELL CONNECTIONS FLOWING STATUS
--
--  WELL WELL   --LOCATION--  COMPLETION
--  NAME STAT     I   J    K  FIRST LAST
WELOPEN
OP01     OPEN                                              /
OP01     OPEN     0   0    0     3     3                   /
OP02     OPEN                                              /
OP02     OPEN     0   0    0     2     2                   /
OP03     OPEN                                              /
OP03     OPEN     0   0    0     0     0                   /
/
\end{verbatim}

Again, the first record of each well WELOPEN keyword changes the well
status from shut (as per the WCONPROD keyword) to open. Then for well
OP01 well completion number three is opened (connections 35 to 90),
completion two (connections 15 to 30) for well OP02 and completion
numbers one to three (all the connections) for well OP03. Note that
there is no completion number two for OP03 so connections 15 to 30 will
not be opened.

Note the last completion number for well OP03 was named completion
number three, but it could have been named number two as well. The
reason why it was named number three instead of two was because it was
assumed (for the example) that layers 35 to 90 represent a particular
reservoir, and therefore allowing for the tracking of completions for
individual reservoirs, as shown in the example.

This example shows how one can open all the wells and well completions
for a given reservoir.

\begin{verbatim}
--
--       DEFINE WELL AND WELL CONNECTIONS FLOWING STATUS
--
--  WELL WELL   --LOCATION--  COMPLETION
--  NAME STAT     I   J    K  FIRST LAST
WELOPEN
‘*’      OPEN                                              /
‘*’      OPEN     0   0    0     3     3                   /
OP02     SHUT     0   0    0     0     0                   /
OP02     OPEN     0   0    0     2     2                   /
/
\end{verbatim}

In this case wells OP01, OP02 and OP03 are opened via completion number
three. All the connections for OP02 are then shut, and the connections
associated with completion two are opened instead for this well.

\subsection{WELOPENL -- Define Well and Well Connections Flowing Status
(LGR)}\label{welopenl-define-well-and-well-connections-flowing-status-lgr}

\begin{longtable}[]{@{}
  >{\raggedright\arraybackslash}p{(\linewidth - 14\tabcolsep) * \real{0.1250}}
  >{\raggedright\arraybackslash}p{(\linewidth - 14\tabcolsep) * \real{0.1250}}
  >{\raggedright\arraybackslash}p{(\linewidth - 14\tabcolsep) * \real{0.1250}}
  >{\raggedright\arraybackslash}p{(\linewidth - 14\tabcolsep) * \real{0.1250}}
  >{\raggedright\arraybackslash}p{(\linewidth - 14\tabcolsep) * \real{0.1250}}
  >{\raggedright\arraybackslash}p{(\linewidth - 14\tabcolsep) * \real{0.1250}}
  >{\raggedright\arraybackslash}p{(\linewidth - 14\tabcolsep) * \real{0.1250}}
  >{\raggedright\arraybackslash}p{(\linewidth - 14\tabcolsep) * \real{0.1250}}@{}}
\toprule\noalign{}
\begin{minipage}[b]{\linewidth}\raggedright
RUNSPEC
\end{minipage} & \begin{minipage}[b]{\linewidth}\raggedright
GRID
\end{minipage} & \begin{minipage}[b]{\linewidth}\raggedright
EDIT
\end{minipage} & \begin{minipage}[b]{\linewidth}\raggedright
PROPS
\end{minipage} & \begin{minipage}[b]{\linewidth}\raggedright
REGIONS
\end{minipage} & \begin{minipage}[b]{\linewidth}\raggedright
SOLUTION
\end{minipage} & \begin{minipage}[b]{\linewidth}\raggedright
SUMMARY
\end{minipage} & \begin{minipage}[b]{\linewidth}\raggedright
SCHEDULE
\end{minipage} \\
\midrule\noalign{}
\endhead
\bottomrule\noalign{}
\endlastfoot
\end{longtable}

\subsubsection{Description}\label{description-1066}

The WELOPENL keyword defines the status of wells and well connection in
Local Grid Refinement Grids (``LGR'') and is used to open and shut
previously defined well and well connections without having to
re-specify all the data on the well control keywords: WCONPROD,
WCONHIST, WCONINJE, or WCONINJH keywords. Note that the well must still
be initially be fully defined using the WCONPROD or WCONINJE keywords.
All the aforementioned keywords are described in the SCHEDULE section

This keyword is not supported by OPM Flow but would change the results
if supported so the simulation will be stopped.

\begin{longtable}[]{@{}
  >{\raggedright\arraybackslash}p{(\linewidth - 6\tabcolsep) * \real{0.2500}}
  >{\raggedright\arraybackslash}p{(\linewidth - 6\tabcolsep) * \real{0.2500}}
  >{\raggedright\arraybackslash}p{(\linewidth - 6\tabcolsep) * \real{0.2500}}
  >{\raggedright\arraybackslash}p{(\linewidth - 6\tabcolsep) * \real{0.2500}}@{}}
\toprule\noalign{}
\begin{minipage}[b]{\linewidth}\raggedright
No.
\end{minipage} & \begin{minipage}[b]{\linewidth}\raggedright
Name
\end{minipage} & \begin{minipage}[b]{\linewidth}\raggedright
Description
\end{minipage} & \begin{minipage}[b]{\linewidth}\raggedright
Default
\end{minipage} \\
\midrule\noalign{}
\endhead
\bottomrule\noalign{}
\endlastfoot
Field & Metric & Laboratory & \\
1 & WELNAME & A character string of up to eight characters in length
that defines the well name for which the well and well connection status
data is being defined. Note that the well name (WELNAME) must have been
declared previously using the WELSPECS keyword in the SCHEDULE section,
otherwise an error may occur. & None \\
2 & LGRNAME & A character string of up to eight characters in length
that defines the LGR name for which the well LGR connection data are
being defined. Note that the well name (LGRNAME) must have been declared
previously using the WELSPECL keyword in the SCHEDULE section, otherwise
an error may occur. If defaulted with 1* the LGR on the WELSPECL keyword
will be utilized. & Defined \\
3 & STATUS & A character string of length four that defines the well and
a well's connections' operational status, STATUS should be set to one of
the following character strings: & OPEN \\
4 & I & An integer less than or equal to NX that defines the connection
location in the I-direction. & 1* \\
5 & J & An integer less than or equal to NY that defines the connection
location in the J-direction. & 1* \\
6 & K & An integer less than or equal to NZ that defines the connection
location in the K-direction. & 1* \\
7 & K1 & An integer less than or equal to NZ that defines the UPPER
completion location in the K-direction. Connections are lumped into
completions via the COMPLUMP keyword, and K1 refers to the first
completion number, as defined by the COMPLUMP keyword, and all the
connections contained within the K1 completion. & 1* \\
8 & K2 & An integer less than or equal to NZ that defines the LOWER
completion location in the K-direction. Connections are lumped into
completions via the COMPLUMP keyword, and K2 refers to the last
completion number, as defined by the COMPLUMP keyword, and all the
connections contained within the K2 completion. & 1* \\
Notes: & & & \\
\end{longtable}

\emph{Table 12.88: WELOPENL Keyword Description}

If variables I, J K, K1 and K2 are all set to a negative number or
defaulted with 1* then STATUS is applied to the well and the well
connection remain unchanged.

If variables I, J K, K1 and K2 are all set to zero or a positive value
then STATUS is applied to the defined connections and the well status
remains unchanged. The defined connections are those values with the I,
J, K specified locations and connections within the completion number
range specified by Atgeirr Rasmussen 2017-09-22T13:24:11.685209000 AFR
Changed K to C here and above because I think it refers to a connection
number and not a grid layer. (Unless as usual I misunderstood
something.) K1 and K2. David Baxendale 2017-10-03T12:14:26.106000000 DBx
Reply to Atgeirr Rasmussen (09/22/2017, 13:24): ``\ldots{}'' I reverted
back to K1 and K2 as I have consistently used this on other keywords.

See also the COMPDAT keyword to define a well's connections, the
COMPLUMP keyword to group well connections into well completions, the
WCONPROD and WCONINJE keywords to define a well's production and
injections targets and constraints. All the aforementioned keywords are
described in the SCHEDULE section.

\subsubsection{Example}\label{example-571}

The following example shows the use of the COMPLMPL keyword to group the
well connections into well completions for well OP01 and then use the
WELOPEN keyword to open the well and the well connections.

\begin{verbatim}
--
--       ASSIGN WELL LGR CONNECTIONS TO COMPLETIONS
--
--       WELL   LGR       ---LOCATION---   COMPL
--       NAME   NAME      II  JJ  K1  K2   NO.
COMPLMPL
         OP01   LGR1      26  58   1   3    1              /
         OP01   LGR1      26  58   4  10    2              /
         OP01   LGR1      26  58  11  12    3              /
/
--
--       WELL PRODUCTION STATUS FOR LGR WELLS
--
--  WELL   LGR      WELL   --LOCATION--  COMPLETION
--  NAME   NAME     STAT     I   J    K  FIRST LAST
WELOPENL
OP01       LGR1     OPEN                                                       /
OP01       LGR1     OPEN     0   0    0     1     2                            /
OP01       LGR2     SHUT     0   0    0     3     3                            /
/
\end{verbatim}

The first record of the WELOPENL keyword changes the well status from
shut (as per the WCONPROD keyword) to open, in case it has been shut-in.
Then well completion number one and two are opened (connections 1 to
10), and completion number three shut-in (connections 11 to 12).

\subsection{WELPI -- Define Well Productivity and Injectivity
Indices}\label{welpi-define-well-productivity-and-injectivity-indices}

\begin{longtable}[]{@{}
  >{\raggedright\arraybackslash}p{(\linewidth - 14\tabcolsep) * \real{0.1250}}
  >{\raggedright\arraybackslash}p{(\linewidth - 14\tabcolsep) * \real{0.1250}}
  >{\raggedright\arraybackslash}p{(\linewidth - 14\tabcolsep) * \real{0.1250}}
  >{\raggedright\arraybackslash}p{(\linewidth - 14\tabcolsep) * \real{0.1250}}
  >{\raggedright\arraybackslash}p{(\linewidth - 14\tabcolsep) * \real{0.1250}}
  >{\raggedright\arraybackslash}p{(\linewidth - 14\tabcolsep) * \real{0.1250}}
  >{\raggedright\arraybackslash}p{(\linewidth - 14\tabcolsep) * \real{0.1250}}
  >{\raggedright\arraybackslash}p{(\linewidth - 14\tabcolsep) * \real{0.1250}}@{}}
\toprule\noalign{}
\begin{minipage}[b]{\linewidth}\raggedright
RUNSPEC
\end{minipage} & \begin{minipage}[b]{\linewidth}\raggedright
GRID
\end{minipage} & \begin{minipage}[b]{\linewidth}\raggedright
EDIT
\end{minipage} & \begin{minipage}[b]{\linewidth}\raggedright
PROPS
\end{minipage} & \begin{minipage}[b]{\linewidth}\raggedright
REGIONS
\end{minipage} & \begin{minipage}[b]{\linewidth}\raggedright
SOLUTION
\end{minipage} & \begin{minipage}[b]{\linewidth}\raggedright
SUMMARY
\end{minipage} & \begin{minipage}[b]{\linewidth}\raggedright
SCHEDULE
\end{minipage} \\
\midrule\noalign{}
\endhead
\bottomrule\noalign{}
\endlastfoot
\end{longtable}

\subsubsection{Description}\label{description-1067}

The WELPI keyword is used to define a well's productivity or injectivity
index at the time the keyword is activated. Productivity and injectivity
indices are a function of both bottom-hole pressure and mobility and
thus will vary in time as the bottom-hole pressure and fluids produce by
the well are changing. Thus, the values enter on this keyword for a
given well will override any previously calculated values, or values
previously entered, using this keyword. This keyword should only be
invoked after a well's connection factors have been declared via the
COMPDAT keyword in the SCHEDULE section.

\begin{longtable}[]{@{}
  >{\raggedright\arraybackslash}p{(\linewidth - 6\tabcolsep) * \real{0.2500}}
  >{\raggedright\arraybackslash}p{(\linewidth - 6\tabcolsep) * \real{0.2500}}
  >{\raggedright\arraybackslash}p{(\linewidth - 6\tabcolsep) * \real{0.2500}}
  >{\raggedright\arraybackslash}p{(\linewidth - 6\tabcolsep) * \real{0.2500}}@{}}
\toprule\noalign{}
\begin{minipage}[b]{\linewidth}\raggedright
No.
\end{minipage} & \begin{minipage}[b]{\linewidth}\raggedright
Name
\end{minipage} & \begin{minipage}[b]{\linewidth}\raggedright
Description
\end{minipage} & \begin{minipage}[b]{\linewidth}\raggedright
Default
\end{minipage} \\
\midrule\noalign{}
\endhead
\bottomrule\noalign{}
\endlastfoot
Field & Metric & Laboratory & \\
1 & WELNAME & A character string of up to eight characters in length
that defines the well name for which the well productivity index is
being defined. Note that the well name (WELNAME) must have been declared
previously using the WELSPECS keyword in the SCHEDULE section, otherwise
an error may occur. & None \\
2 & WELPI & A real positive value that defines the productivity or
injectivity of a well at the time the keyword is activated in the input
deck. Subsequent use of the keyword for a given well will override all
previously entered data, including the original productivity index
calculated from the well's COMPDAT connection data. Note that:
Subsequent usage of the keyword for the same well may lead to unintended
values; use the WPI series of variables in the SUMMARY section to output
and verify that the values are consistent with the desired results. &
1.0 \\
Liquid stb/d/psia Gas Mscf/d/psia & Liquid sm3/day/bars Gas sm3/day/bars
& Liquid scc/hour/atm Gas scc/hour/atm & \\
Notes: & & & \\
\end{longtable}

\emph{Table 12.89: WELPI Keyword Description}

Note that only a well's currently opened connection factors, as entered
via the COMPDAT keyword, are re-scaled to match the required
productivity and injectivity index. See section Well Productivity on
well productivity for further information on well productivity.

\begin{tcolorbox}[enhanced jigsaw, title=\textcolor{quarto-callout-note-color}{\faInfo}\hspace{0.5em}{Note}, coltitle=black, breakable, toprule=.15mm, colbacktitle=quarto-callout-note-color!10!white, leftrule=.75mm, opacitybacktitle=0.6, colback=white, bottomtitle=1mm, toptitle=1mm, titlerule=0mm, left=2mm, colframe=quarto-callout-note-color-frame, arc=.35mm, rightrule=.15mm, bottomrule=.15mm, opacityback=0]

One should not modify the productivity and injectivity index of gas
wells that use:

\end{tcolorbox}

See also the WPIMULT keyword in the SCHEDULE section to set the
productivity and injectivity index for a well by scaling the current
index by a scaling factor.

\subsubsection{Example}\label{example-572}

\begin{verbatim}
--
--       DEFINE WELL PRODUCTIVITY/INJECTIVITY INDEX
--
-- WELL  PI
-- NAME  MULT
WELPI
OP01     1.250                                             /
OP02     2.750                                             /
OP03     1.100                                             /
/
\end{verbatim}

In the above example the oil wells are are defined as having a well
productivity indices of 1.250, 2.750 and 1.100 for the OP01, OP02 and
OP03 wells, respectively.

\subsection{WELPRI -- Assign Well
Priority}\label{welpri-assign-well-priority}

\begin{longtable}[]{@{}
  >{\raggedright\arraybackslash}p{(\linewidth - 14\tabcolsep) * \real{0.1250}}
  >{\raggedright\arraybackslash}p{(\linewidth - 14\tabcolsep) * \real{0.1250}}
  >{\raggedright\arraybackslash}p{(\linewidth - 14\tabcolsep) * \real{0.1250}}
  >{\raggedright\arraybackslash}p{(\linewidth - 14\tabcolsep) * \real{0.1250}}
  >{\raggedright\arraybackslash}p{(\linewidth - 14\tabcolsep) * \real{0.1250}}
  >{\raggedright\arraybackslash}p{(\linewidth - 14\tabcolsep) * \real{0.1250}}
  >{\raggedright\arraybackslash}p{(\linewidth - 14\tabcolsep) * \real{0.1250}}
  >{\raggedright\arraybackslash}p{(\linewidth - 14\tabcolsep) * \real{0.1250}}@{}}
\toprule\noalign{}
\begin{minipage}[b]{\linewidth}\raggedright
RUNSPEC
\end{minipage} & \begin{minipage}[b]{\linewidth}\raggedright
GRID
\end{minipage} & \begin{minipage}[b]{\linewidth}\raggedright
EDIT
\end{minipage} & \begin{minipage}[b]{\linewidth}\raggedright
PROPS
\end{minipage} & \begin{minipage}[b]{\linewidth}\raggedright
REGIONS
\end{minipage} & \begin{minipage}[b]{\linewidth}\raggedright
SOLUTION
\end{minipage} & \begin{minipage}[b]{\linewidth}\raggedright
SUMMARY
\end{minipage} & \begin{minipage}[b]{\linewidth}\raggedright
SCHEDULE
\end{minipage} \\
\midrule\noalign{}
\endhead
\bottomrule\noalign{}
\endlastfoot
\end{longtable}

\subsubsection{Description}\label{description-1068}

The WELPRI keyword is used to re-assign a priority number to a well for
when the PRIORITY keyword has been used in the SCHEDULE section. The
PRIORITY keyword activates the Well Priority option and defines the
coefficients in the well priority equation; WELPRI keyword can be used
to over write these calculated priority numbers

This keyword is not supported by OPM Flow but would change the results
if supported so the simulation will be stopped.

\subsection{WELSEGS -- Define Multi-Segment Wells and Their Segment
Structure}\label{welsegs-define-multi-segment-wells-and-their-segment-structure}

\begin{longtable}[]{@{}
  >{\raggedright\arraybackslash}p{(\linewidth - 14\tabcolsep) * \real{0.1250}}
  >{\raggedright\arraybackslash}p{(\linewidth - 14\tabcolsep) * \real{0.1250}}
  >{\raggedright\arraybackslash}p{(\linewidth - 14\tabcolsep) * \real{0.1250}}
  >{\raggedright\arraybackslash}p{(\linewidth - 14\tabcolsep) * \real{0.1250}}
  >{\raggedright\arraybackslash}p{(\linewidth - 14\tabcolsep) * \real{0.1250}}
  >{\raggedright\arraybackslash}p{(\linewidth - 14\tabcolsep) * \real{0.1250}}
  >{\raggedright\arraybackslash}p{(\linewidth - 14\tabcolsep) * \real{0.1250}}
  >{\raggedright\arraybackslash}p{(\linewidth - 14\tabcolsep) * \real{0.1250}}@{}}
\toprule\noalign{}
\begin{minipage}[b]{\linewidth}\raggedright
RUNSPEC
\end{minipage} & \begin{minipage}[b]{\linewidth}\raggedright
GRID
\end{minipage} & \begin{minipage}[b]{\linewidth}\raggedright
EDIT
\end{minipage} & \begin{minipage}[b]{\linewidth}\raggedright
PROPS
\end{minipage} & \begin{minipage}[b]{\linewidth}\raggedright
REGIONS
\end{minipage} & \begin{minipage}[b]{\linewidth}\raggedright
SOLUTION
\end{minipage} & \begin{minipage}[b]{\linewidth}\raggedright
SUMMARY
\end{minipage} & \begin{minipage}[b]{\linewidth}\raggedright
SCHEDULE
\end{minipage} \\
\midrule\noalign{}
\endhead
\bottomrule\noalign{}
\endlastfoot
\end{longtable}

\subsubsection{Description}\label{description-1069}

The WELSEGS keyword defines a well to be a multi-segment well and
defines the well's segment structure. Note that the well must have
previously been defined by the WELSPECS keyword in the SCHEDULE section
and that the WELSEGS keyword should be repeated for each multi-segment
well in the model.

\begin{longtable}[]{@{}
  >{\raggedright\arraybackslash}p{(\linewidth - 6\tabcolsep) * \real{0.2500}}
  >{\raggedright\arraybackslash}p{(\linewidth - 6\tabcolsep) * \real{0.2500}}
  >{\raggedright\arraybackslash}p{(\linewidth - 6\tabcolsep) * \real{0.2500}}
  >{\raggedright\arraybackslash}p{(\linewidth - 6\tabcolsep) * \real{0.2500}}@{}}
\toprule\noalign{}
\begin{minipage}[b]{\linewidth}\raggedright
No.
\end{minipage} & \begin{minipage}[b]{\linewidth}\raggedright
Name
\end{minipage} & \begin{minipage}[b]{\linewidth}\raggedright
Description
\end{minipage} & \begin{minipage}[b]{\linewidth}\raggedright
Default
\end{minipage} \\
\midrule\noalign{}
\endhead
\bottomrule\noalign{}
\endlastfoot
Field & Metric & Laboratory & \\
1-1 & WELNAME & A character string of up to eight characters in length
that defines the well name for which a multi-segment well is being
defined. Note that the well name (WELNAME) must have been declared
previously using the WELSPECS keyword in the SCHEDULE section, otherwise
an error may occur. & None \\
1-2 & TOPDEP & A real value that defines the depth of the nodal point of
the top segment. This is used as the reference depth for reporting the
bottom-hole pressure for the multi-segment well. If the keyword is
entered multiple times for the same well, due to for example the well
configuration changing through time, then it is only necessary to enter
this data the first time the keyword is used for a well. & None \\
feet & m & cm & \\
1-3 & TOPLEN & A real positive value that defines the length of the
tubing from the tubing reference point (for example Measured Depth
Relative to Kelly Bushing, MDRKB) to the nodal point of the top segment.
Tubing pressures between the nodal point of the top segment and the well
head are not calculated by the multi-segment well option as these are
taken into account by the VFP tables allocated to the well and entered
via the VFPPROD and VFPINJ keywords in the SCHEDULE section. If TOPLEN
is set to zero or defaulted then the tubing length is measured from the
nodal point of the top segment. & 0 \\
feet & m & cm & \\
1-4 & WBORVOL & A real positive value that defines the effective
wellbore volume for the top segment, that is from the tubing head or
wellhead at the surface to the nodal point of the top segment. The
default value of 1.0 x 10-5 results in minimal wellbore storage. &
1.0E-5 \\
ft3 & m3 & cm3 & \\
1-5 & TUBOPT & A defined character string that specifies the type of
length and depth data entered for LENGTH and DEPTH on the second record
and should be set to one of the following: As there is no default value
for TUBOPT one of the above options must be explicitly defined. &
None \\
1-6 & PRESOPT & A defined character string that specifies the pressure
drop calculation used for each well segment and should be set to one of
the following: The default value for PRESOPT of HFA sets the pressure
calculation to include the hydrostatic, friction and acceleration terms.
& HFA \\
1-7 & FLOWOPT & A defined character string that specifies the type of
multi-phase calculation used for each well segment and should be set to
one of the following: OPM Flow only supports the default value of HO. &
HO \\
1-8 & TOPX & A real positive value equal to or greater than zero that
defines the coordinate in the x-direction of the nodal point of the top
segment that is used for display purposes only. Currently this option is
not supported by OPM Flow. & None \\
feet & m & cm & \\
1-9 & TOPY & A real positive value equal to or greater than zero that
defines the coordinate in the y-direction of the nodal point of the top
segment that is used for display purposes only. Currently this option is
not supported by OPM Flow. & None \\
feet & m & cm & \\
1-10 & XAREATH & A real positive value equal to or greater than zero
that defines the cross-sectional area of the pipe wall used in thermal
conductivity calculations for when the temperature calculation is
activated by the TEMP keyword in the RUNSPEC section. Currently this
option is not supported by OPM Flow. & None \\
ft2 & m2 & cm2 & \\
1-11 & VHEATCAP & A real positive value equal to or greater than zero
that defines the volumetric heat capacity of the pipe wall used in
thermal conductivity calculations for when the temperature calculation
is activated by the TEMP keyword in the RUNSPEC section. Currently this
option is not supported by OPM Flow. & None \\
Btu/ft/day/°R & kJ/m/day/K & J/cm/hr/K & \\
1-12 & THCON & A real positive value equal to or greater than zero that
defines the thermal conductivity of the pipe wall used in thermal
conductivity calculations for when the temperature calculation is
activated by the TEMP keyword in the RUNSPEC section. Currently this
option is not supported by OPM Flow. & None \\
Btu/ft/day/°R & kJ/m/day/K & J/cm/hr/K & \\
1-13 & / & Record terminated by a ``/'' & Not Applicable \\
2-1 & ISEG1 & A positive integer greater than or equal to two and less
than or equal to MXSEGS on WSEGDIMS keyword in the RUNSPEC section that
defines the start of the segment range. & None \\
2-2 & ISEG2 & A positive integer greater than or equal to ISEG1 on this
record and less than or equal to MXSEGS on the WSEGDIMS keyword in the
RUNSPEC section that defines the end of the segment range. & None \\
2-3 & IBRANCH & A positive integer greater than or equal to one and less
than or equal to MXBRAN on WSEGDIMS keyword in the RUNSPEC section that
defines the branch number of a segment. All segments on the main stem
must have IBRANCH set to one, and those on lateral branches should have
IBRANCH set to between two and MXBRAN. & None \\
2-4 & ISEG3 & A positive integer greater than or equal to one and less
than or equal to MXSEGS on the WSEGDIMS keyword in the RUNSPEC section
that defines the outlet segment for the segment at the start of the
segment range (ISEG1). & None \\
2-5 & LENGTH & A real positive value that: & None \\
feet & m & cm & \\
2-6 & DEPTH & A real positive value that: & \\
feet & m & cm & None \\
2-7 & ID & A real positive value that defines the tubing internal
diameter of each segment in the segment range. & None \\
feet & m & cm & \\
2-8 & EPSILON & A real positive value that defines the tubing absolute
roughness of each segment in the segment range. & None \\
feet & m & cm & \\
2-9 & XAREA & A real positive value equal to or greater than zero that
defines the cross-sectional area for fluid flow. Currently this option
is not supported by OPM Flow. & None \\
ft2 & m2 & cm2 & \\
2-10 & VOLSEG & VOLSEG is a real positive value that defines the
effective segment volume for the this segment. Currently this option is
not supported by OPM Flow. & None \\
ft3 & m3 & cm3 & \\
2-11 & XCORDS & A real positive value that: Currently this option is not
supported by OPM Flow. & None \\
feet & m & cm & \\
2-12 & YCORDS & A real positive value that: Currently this option is not
supported by OPM Flow. & None \\
feet & m & cm & \\
2-13 & XAREAS & A real positive value equal to or greater than zero that
defines the cross-sectional area of the pipe wall for this segment, that
is used in thermal conductivity calculations for when the temperature
calculation is activated by the TEMP keyword in the RUNSPEC section.
Currently this option is not supported by OPM Flow. & None \\
ft2 & m2 & cm2 & \\
2-14 & VHEATSEG & A real positive value equal to or greater than zero
that defines the volumetric heat capacity of the pipe wall for this
segment, that is used in thermal conductivity calculations for when the
temperature calculation is activated by the TEMP keyword in the RUNSPEC
section. Currently this option is not supported by OPM Flow. & None \\
Btu/ft/day/°R & kJ/m/day/K & J/cm/hr/K & \\
2-15 & THCSEG & A real positive value equal to or greater than zero that
defines the thermal conductivity of the pipe wall for this segment, that
is used in thermal conductivity calculations for when the temperature
calculation is activated by the TEMP keyword in the RUNSPEC section.
Currently this option is not supported by OPM Flow. & None \\
Btu/ft/day/°R & kJ/m/day/K & J/cm/hr/K & \\
2-16 & / & Record terminated by a ``/'' & Not Applicable \\
Notes: & & & \\
\end{longtable}

\emph{Table 12.90: WELSEGS Keyword Description}

The maximum number of wells should be defined via the WELLDIMS keyword
and the maximum number of multi-segment wells should be declared on the
WSEGDIMS keyword, both keywords are in the RUNSPEC section.

See also the WELSPECS keyword to define wells, the COMPDAT keyword to
define the well completions for both ordinary wells and multi-segment
wells, and the COMPSEGS keyword to define multi-segment well segment
completions. All the aforementioned keywords are described in the
SCHEDULE section.

\subsubsection{Example}\label{example-573}

The following example defines one multi-segment oil production well
(OP01) using the WELSPECS, WELSEGS, COMPDAT and COMPSEGS keywords, and
one standard water injection well (WI01) using the WELSPECS and COMPDAT
keywords.

\begin{verbatim}
--
--       WELL SPECIFICATION DATA
--
-- WELL  GROUP     LOCATION  BHP    PHASE  DRAIN  INFLOW  OPEN  CROSS  PVT
-- NAME  NAME        I    J  DEPTH  FLUID  AREA   EQUANS  SHUT  FLOW   TABLE
WELSPECS
OP01     PLATFORM   10   10   1*     OIL                                       /
WI01     PLATFORM    1    1   1*     WATER                                     /
/
--
--       WELL CONNECTION DATA
--
-- WELL  --- LOCATION ---  OPEN   SAT   CONN   WELL   KH    SKIN   D     DIR
-- NAME   II  JJ  K1  K2   SHUT   TAB   FACT   DIA    FACT  FACT   FACT  PEN
COMPDAT
OP01      10  10   1   1   OPEN   1*    200.   0.5                             /
OP01      10  10   2   2   OPEN   1*    200.   0.5                             /
OP01      10  10   3   3   OPEN   1*    200.   0.4                             /
OP01      10  10   4   4   OPEN   1*    200.   0.4                             /
OP01      10  10   5   5   OPEN   1*    200.   0.4                             /
OP01      10  10   6   6   OPEN   1*    200.   0.4                             /

OP01       9  10   2   2   OPEN   1*    200.   0.4                             /
OP01       8  10   2   2   OPEN   1*    200.   0.4                             /
OP01       7  10   2   2   OPEN   1*    200.   0.4                             /
OP01       6  10   2   2   OPEN   1*    200.   0.4                             /
OP01       5  10   2   2   OPEN   1*    200.   0.4                             /

OP01      10   9   3   3   OPEN   1*    200.   0.4                             /
OP01      10   8   3   3   OPEN   1*    200.   0.4                             /
OP01      10   7   3   3   OPEN   1*    200.   0.4                             /
OP01      10   6   3   3   OPEN   1*    200.   0.4                             /
OP01      10   5   3   3   OPEN   1*    200.   0.4                             /

OP01       9  10   5   5   OPEN   1*    200.   0.4                             /
OP01       8  10   5   5   OPEN   1*    200.   0.4                             /
OP01       7  10   5   5   OPEN   1*    200.   0.4                             /
OP01       6  10   5   5   OPEN   1*    200.   0.4                             /
OP01       5  10   5   5   OPEN   1*    200.   0.4                             /

OP01      10   9   6   6   OPEN   1*    200.   0.4                             /
OP01      10   8   6   6   OPEN   1*    200.   0.4                             /
OP01      10   7   6   6   OPEN   1*    200.   0.4                             /
OP01      10   6   6   6   OPEN   1*    200.   0.4                             /
OP01      10   5   6   6   OPEN   1*    200.   0.4                             /

WI01       1   1   7   9   OPEN   1*    200.   0.5                             /
/
--
--       WELL SEGMENT SPECIFICATION DATA
--
-- WELL  NODAL       LEN     WELL    DEPH    PRESS   FLOW
-- NAME  DEPTH       TUBING  VOLM    OPTN    CALC    MODEL
WELSEGS
OP01     2512.5      2512.5  1.0E-5  ABS     HFA     HO                        /
--
--       SEG   SEG   BRAN    SEG     TUBING  NODAL   TUBE   TUBE     XSEC  VOL
--       ISTR  IEND  NO      NO      LENGTH  DEPTH   ID     ROUGH    AREA  SEG
           2     2    1       1      2537.5  2534.5  0.3    0.00010            /
           3     3    1       2      2562.5  2560.5  0.3    0.00010            /
           4     4    1       3      2587.5  2593.5  0.3    0.00010            /
           5     5    1       4      2612.5  2614.5  0.3    0.00010            /
           6     6    1       5      2637.5  2635.5  0.3    0.00010            /

           7     7    2       2      2737.5  2538.5  0.2    0.00010            /
           8     8    2       7      2937.5  2537.5  0.2    0.00010            /
           9     9    2       8      3137.5  2539.5  0.2    0.00010            /
          10    10    2       9      3337.5  2535.5  0.2    0.00010            /
          11    11    2      10      3537.5  2536.5  0.2    0.00010            /

          12    12    3       3      2762.5  2563.5  0.2    0.00010            /
          13    13    3      12      2962.5  2562.5  0.1    0.00010            /
          14    14    3      13      3162.5  2562.5  0.1    0.00010            /
          15    15    3      14      3362.5  2564.5  0.1    0.00010            /
          16    16    3      15      3562.5  2562.5  0.1    0.00010            /

          17    17    4       5      2812.5  2613.5  0.2    0.00010            /
          18    18    4      17      3012.5  2612.5  0.1    0.00010            /
          19    19    4      18      3212.5  2612.5  0.1    0.00010            /
          20    20    4      19      3412.5  2612.5  0.1    0.00010            /
          21    21    4      20      3612.5  2613.5  0.1    0.00010            /

          22    22    5       6      2837.5  2634.5  0.2    0.00010            /
          23    23    5      22      3037.5  2637.5  0.2    0.00010            /
          24    24    5      23      3237.5  2638.5  0.2    0.00010            /
          25    25    5      24      3437.5  2639.5  0.1    0.00010            /
          26    26    5      25      3637.5  2639.5  0.1    0.00010            /
/
--
--       COMPLETION SEGMENT SPECIFICATION DATA
--
-- WELL
-- NAME
COMPSEGS
OP01                                                                           /
--
--      --LOCATION--  BRAN  START   END     DIR  LOC    MID    COMP    ISEG
--       II  JJ  K1   NO    LENGTH  LENGTH  PEN  I,J,K  PERFS  LENGTH  NO.
         10  10   1    1    2512.5  2525.0                                     /
         10  10   2    1    2525.0  2550.0                                     /
         10  10   3    1    2550.0  2575.0                                     /
         10  10   4    1    2575.0  2600.0                                     /
         10  10   5    1    2600.0  2625.0                                     /
         10  10   6    1    2625.0  2650.0                                     /

          9  10   2    2    2637.5  2837.5                                     /
          8  10   2    2    2837.5  3037.5                                     /
          7  10   2    2    3037.5  3237.5                                     /
          6  10   2    2    3237.5  3437.5                                     /
          5  10   2    2    3437.5  3637.5                                     /

         10   9   3    3    2662.5  2862.5                                     /
         10   8   3    3    2862.5  3062.5                                     /
         10   7   3    3    3062.5  3262.5                                     /
         10   6   3    3    3262.5  3462.5                                     /
         10   5   3    3    3462.5  3662.5                                     /

          9  10   5    4    2712.5  2912.5                                     /
          8  10   5    4    2912.5  3112.5                                     /
          7  10   5    4    3112.5  3312.5                                     /
          6  10   5    4    3312.5  3512.5                                     /
          5  10   5    4    3512.5  3712.5                                     /

         10   9   6    5    2737.5  2937.5                                     /
         10   8   6    5    2937.5  3137.5                                     /
         10   7   6    5    3137.5  3337.5                                     /
         10   6   6    5    3337.5  3537.5                                     /
         10   5   6    5    3537.5  3737.5                                     /
/
\end{verbatim}

Note the use of both the COMPDAT and COMPSEGS keywords to fully define a
multi-segment well's completion.

Finally, Figure 12.15 depicts the resulting well configuration for both
wells, with the conventional water injection well shown in blue and the
multi-segment oil producer shown in green.

\pandocbounded{\includegraphics[keepaspectratio]{chapters/images/Frame8_47683e16b52f.png}}
\pandocbounded{\includegraphics[keepaspectratio]{chapters/images/Image8_47683e16b52f.png}}

\subsection{WELSOMIN -- Define Well Connection Minimum Oil Saturation
for
Opening}\label{welsomin-define-well-connection-minimum-oil-saturation-for-opening}

\begin{longtable}[]{@{}
  >{\raggedright\arraybackslash}p{(\linewidth - 14\tabcolsep) * \real{0.1250}}
  >{\raggedright\arraybackslash}p{(\linewidth - 14\tabcolsep) * \real{0.1250}}
  >{\raggedright\arraybackslash}p{(\linewidth - 14\tabcolsep) * \real{0.1250}}
  >{\raggedright\arraybackslash}p{(\linewidth - 14\tabcolsep) * \real{0.1250}}
  >{\raggedright\arraybackslash}p{(\linewidth - 14\tabcolsep) * \real{0.1250}}
  >{\raggedright\arraybackslash}p{(\linewidth - 14\tabcolsep) * \real{0.1250}}
  >{\raggedright\arraybackslash}p{(\linewidth - 14\tabcolsep) * \real{0.1250}}
  >{\raggedright\arraybackslash}p{(\linewidth - 14\tabcolsep) * \real{0.1250}}@{}}
\toprule\noalign{}
\begin{minipage}[b]{\linewidth}\raggedright
RUNSPEC
\end{minipage} & \begin{minipage}[b]{\linewidth}\raggedright
GRID
\end{minipage} & \begin{minipage}[b]{\linewidth}\raggedright
EDIT
\end{minipage} & \begin{minipage}[b]{\linewidth}\raggedright
PROPS
\end{minipage} & \begin{minipage}[b]{\linewidth}\raggedright
REGIONS
\end{minipage} & \begin{minipage}[b]{\linewidth}\raggedright
SOLUTION
\end{minipage} & \begin{minipage}[b]{\linewidth}\raggedright
SUMMARY
\end{minipage} & \begin{minipage}[b]{\linewidth}\raggedright
SCHEDULE
\end{minipage} \\
\midrule\noalign{}
\endhead
\bottomrule\noalign{}
\endlastfoot
\end{longtable}

\subsubsection{Description}\label{description-1070}

WELSOMIN defines a minimum oil saturation for a well connection above
which the connection will be opened automatically. If the grid block
connection is below WELSOMIN then connection will not be automatically
opened. Automatic opening of connection is controlled by the STATUS
parameter on the COMPDAT keyword in the SCHEDULE section. Note that if
the COMPLUMP keyword in the SCHEDULE section has been used to lump
connections into completions then WELSOMIN is compared to the average
oil saturation of the completion.

This keyword is not supported by OPM Flow but would change the results
if supported so the simulation will be stopped.

\subsection{WELSPECL -- Define Well Specifications for Local Grid
Refinements}\label{welspecl-define-well-specifications-for-local-grid-refinements}

\begin{longtable}[]{@{}
  >{\raggedright\arraybackslash}p{(\linewidth - 14\tabcolsep) * \real{0.1250}}
  >{\raggedright\arraybackslash}p{(\linewidth - 14\tabcolsep) * \real{0.1250}}
  >{\raggedright\arraybackslash}p{(\linewidth - 14\tabcolsep) * \real{0.1250}}
  >{\raggedright\arraybackslash}p{(\linewidth - 14\tabcolsep) * \real{0.1250}}
  >{\raggedright\arraybackslash}p{(\linewidth - 14\tabcolsep) * \real{0.1250}}
  >{\raggedright\arraybackslash}p{(\linewidth - 14\tabcolsep) * \real{0.1250}}
  >{\raggedright\arraybackslash}p{(\linewidth - 14\tabcolsep) * \real{0.1250}}
  >{\raggedright\arraybackslash}p{(\linewidth - 14\tabcolsep) * \real{0.1250}}@{}}
\toprule\noalign{}
\begin{minipage}[b]{\linewidth}\raggedright
RUNSPEC
\end{minipage} & \begin{minipage}[b]{\linewidth}\raggedright
GRID
\end{minipage} & \begin{minipage}[b]{\linewidth}\raggedright
EDIT
\end{minipage} & \begin{minipage}[b]{\linewidth}\raggedright
PROPS
\end{minipage} & \begin{minipage}[b]{\linewidth}\raggedright
REGIONS
\end{minipage} & \begin{minipage}[b]{\linewidth}\raggedright
SOLUTION
\end{minipage} & \begin{minipage}[b]{\linewidth}\raggedright
SUMMARY
\end{minipage} & \begin{minipage}[b]{\linewidth}\raggedright
SCHEDULE
\end{minipage} \\
\midrule\noalign{}
\endhead
\bottomrule\noalign{}
\endlastfoot
\end{longtable}

\subsubsection{Description}\label{description-1071}

The WELSPECL keyword defines the general well specification data for all
well types and must be used for all wells contained within a Local Grid
Refinement (``LGR'') instead of the WELSPECS keyword. WELSPECL must
declare wells first before any other LGR well specification keywords are
used in the input file. The keyword declares the name of well, the group
the well belongs to, the LGR the well is incorporated into, the wellhead
location and other key parameters.

This keyword is not supported by OPM Flow but would change the results
if supported so the simulation will be stopped.

\begin{longtable}[]{@{}
  >{\raggedright\arraybackslash}p{(\linewidth - 6\tabcolsep) * \real{0.2500}}
  >{\raggedright\arraybackslash}p{(\linewidth - 6\tabcolsep) * \real{0.2500}}
  >{\raggedright\arraybackslash}p{(\linewidth - 6\tabcolsep) * \real{0.2500}}
  >{\raggedright\arraybackslash}p{(\linewidth - 6\tabcolsep) * \real{0.2500}}@{}}
\toprule\noalign{}
\begin{minipage}[b]{\linewidth}\raggedright
No.
\end{minipage} & \begin{minipage}[b]{\linewidth}\raggedright
Name
\end{minipage} & \begin{minipage}[b]{\linewidth}\raggedright
Description
\end{minipage} & \begin{minipage}[b]{\linewidth}\raggedright
Default
\end{minipage} \\
\midrule\noalign{}
\endhead
\bottomrule\noalign{}
\endlastfoot
Field & Metric & Laboratory & \\
1 & WELNAME & A character string of up to eight characters in length
that defines the well name for which the well specification data is
being defined. & None \\
2 & GRPNAME & A character string of up to eight characters in length
that defines the group name for which the well is assigned to. The group
named FIELD is the top most group and thus GRPNAME cannot be set to
FIELD, although this is allowed in the commercial compositional
simulator but not the commercial black-oil simulator. Note that the
group hierarchy should be defined by the GRUPTREE keyword when there is
more than one level of groups, otherwise all the groups will sit
directly under the FIELD group in the group tree hierarchy. Secondly,
groups defined by the GRUPTREE keyword cannot contain other groups and
wells; that is, groups must either contain other groups or wells but not
both. If necessary, wells can be re-allocated to a different group by
re-entering a well's WELSPECS data together with a new value for
GRPNAME. & None \\
3 & LGRNAME & A character string of up to eight characters in length
that defines the name of the local grid refinement for which the well is
assigned to. & None \\
4 & I & A positive integer greater than or equal to zero and less than
or equal to NX on the CARFIN keyword for Cartesian grids, that defines
the wellhead location for a vertical or deviated well, or the heel for a
horizontal well in the I-direction within the LGR. For radial LGRs this
parameter should be set to one. & None \\
5 & J & A positive integer greater than or equal to zero and less than
or equal to NY on the CARFIN keyword for Cartesian grids, that defines
the wellhead location for a vertical or deviated well, or the heel for a
horizontal well in the J-direction within the LGR. For radial LGRs this
parameter should be set to one. & None \\
6 & BHPREF & A real value that defines the reference depth for reporting
the bottom-hole pressure for the well. Ideally this value should be set
to the midpoint of the perforations as defined by the COMPDATL keyword
in the SCHEDULE section. If defaulted by 1* or set to a value less than
or equal to zero, then the mid-point of shallowest connection defined by
the COMPDATL keyword will be used. & Mid-point of shallowest connection
defined by the COMPDATL keyword \\
feet & m & cm & \\
7 & TYPE & A defined character that defines the ``main'' phase for the
well, and should be set to one of the following character strings: This
parameter defines the phase used to calculate a well's productivity or
injectivity index and the type of well, or a well's connection, to close
when a group's production constraints, as defined on the GCONPROD
keyword in the SCHEDULE section, have been violated. For example, if the
well is declared as an oil well, then excessive gas and water
connections will be subject to closure. & None \\
8 & DRADIUS & A real value that defines the well drainage radius for the
well used to calculate a well's productivity or injectivity index. A
default of zero results in the pressure equivalent radius of the grid
blocks containing the well connections are used. & 0.0 \\
feet & m & cm & \\
9 & INFLOW & A defined character string that defines the inflow equation
to be used for the well in calculating the well's flow rates. INFLOW
should be set to one of the following character strings: For oil and
water wells the INFLOW should be set to STD, why for dry gas wells
INFLOW can be set to either R-G or P-P; however, the P-P option is
preferred for dry gas wells due to the more rigorous treatment of gas
flow. For wet gas wells, that is gas condensate wells, INFLOW should be
set to GPP. Only INFLOW equal to STD and NO are currently implemented in
OPM Flow. & STD \\
10 & AUTO & A defined character string that defines the automatic action
to be taken if the economic WCUT, GOR, or WGR limits are violated and
the well is to cease production. AUTO should be set to one of the
following character strings: The corrective action takes places at the
end of the time step in which the constraint is violated. & SHUT \\
11 & XFLOW & A defined character string that defines the if cross flow
should occur within the wellbore, and should be set to either: In some
cases numerical issues can occur if this variable is set to YES, and
resetting it to NO may resolve the issue; however the results may not
represent the physical process in this case. & YES \\
12 & PVTNUM & A positive integer greater than or equal to zero that
defines the PVT table used to calculate the wellbore fluid properties
that define the relationship between reservoir and surface volume rates.
The default value of zero sets PVTNUM to be the PVT table of the deepest
connection in the well. & 0 \\
13 & DENOPT & A defined character string that sets the type of density
calculation used in calculating the wellbore hydrostatic head, and
should be set to one of the following character strings: The default
option of 1* invokes the SEG option and is the only option implemented
in OPM Flow. & SEG \\
14 & FIPNUM & An integer value defines the FIPNUM region used to
determine the reservoir conditions in calculating the well's reservoir
volumes. If set to a negative integer value then the FIPNUM region of
the deepest connection in the well will be used. If set to zero, the
default value, then the average properties for the field will be used.
If set to an integer value greater than zero, then the FIPNUM indicated
by this value will be used. & 0 \\
15 & & Not used. & \\
16 & & Not used. & \\
17 & & Not used. & \\
18 & & Not used. & \\
Notes: & & & \\
\end{longtable}

\emph{Table 12.91: WELSPECL Keyword Description}

See also the COMPDATL keyword to define a well's connections in a LGR,
the WCONPROD and WCONINJE keywords to define a well's production and
injections targets and constraints. All the aforementioned keywords are
described in the SCHEDULE section.

\subsubsection{Example}\label{example-574}

The following example defines three wells using the WELSPECL keyword

\begin{verbatim}
--
-- WELL SPECIFICATION DATA FOR LGR WELLS
--
-- WELL  GROUP    LGR    LOCATION  BHP    PHASE  DRAIN  INFLOW  SHUT  CROSS  PVT
-- NAME  NAME     NAME    I    J   DEPTH  FLUID  AREA   EQUA.   IN    FLOW  TABLE
WELSPECL
GI01     PLATFORM LGR01  14   13   1*     GAS    1*     P-P    SHUT   NO     1* /
GP01     PLATFORM LGR01  64   80   1*     GAS    1*     GPP    SHUT   NO     1* /
OP01     PLATFORM LGR02  24   10   1*     OIL    1*     STD    SHUT   NO     1* /
/
\end{verbatim}

Here, well GI01 and GP01 are in the same LGR named LGR01 and OP01 is in
a separate LGR named LGR02. GI01 is a dry gas injection well that uses
the dry gas pseudo inflow equation, GP01 is a gas condensate well that
uses the generalized gas pseudo pressure inflow equation, and finally,
OP01 is an oil well that uses the standard inflow equation. All wells:
will be shut if they are required to cease production, all wells
disallow cross flow, and the hydrostatic head calculation is defaulted
to the segment option for all wells.

\subsection{WELSPECS -- Define Well
Specifications}\label{welspecs-define-well-specifications}

\begin{longtable}[]{@{}
  >{\raggedright\arraybackslash}p{(\linewidth - 14\tabcolsep) * \real{0.1250}}
  >{\raggedright\arraybackslash}p{(\linewidth - 14\tabcolsep) * \real{0.1250}}
  >{\raggedright\arraybackslash}p{(\linewidth - 14\tabcolsep) * \real{0.1250}}
  >{\raggedright\arraybackslash}p{(\linewidth - 14\tabcolsep) * \real{0.1250}}
  >{\raggedright\arraybackslash}p{(\linewidth - 14\tabcolsep) * \real{0.1250}}
  >{\raggedright\arraybackslash}p{(\linewidth - 14\tabcolsep) * \real{0.1250}}
  >{\raggedright\arraybackslash}p{(\linewidth - 14\tabcolsep) * \real{0.1250}}
  >{\raggedright\arraybackslash}p{(\linewidth - 14\tabcolsep) * \real{0.1250}}@{}}
\toprule\noalign{}
\begin{minipage}[b]{\linewidth}\raggedright
RUNSPEC
\end{minipage} & \begin{minipage}[b]{\linewidth}\raggedright
GRID
\end{minipage} & \begin{minipage}[b]{\linewidth}\raggedright
EDIT
\end{minipage} & \begin{minipage}[b]{\linewidth}\raggedright
PROPS
\end{minipage} & \begin{minipage}[b]{\linewidth}\raggedright
REGIONS
\end{minipage} & \begin{minipage}[b]{\linewidth}\raggedright
SOLUTION
\end{minipage} & \begin{minipage}[b]{\linewidth}\raggedright
SUMMARY
\end{minipage} & \begin{minipage}[b]{\linewidth}\raggedright
SCHEDULE
\end{minipage} \\
\midrule\noalign{}
\endhead
\bottomrule\noalign{}
\endlastfoot
\end{longtable}

\subsubsection{Description}\label{description-1072}

The WELSPECS keyword defines the general well specification data for all
well types, and must be used for all wells before any other well
specification keywords are used in the input file. The keyword declares
the well, the name of the well, the group the well initial belongs to,
the wellhead location and other key parameters.

OPM Flow also supports the use of the WELSPECS keyword to change the
properties of one or more wells previously specified using the WELSPECS
keyword. In particular, this supports changing the controlling group
(GRPNAME) without affecting any other well properties such as the
location of the well head (I, J) or the well reference depth (BHPREF),
since defaulted properties are not changed in this case. However, the
well reference depth can be re-defaulted by specifying a negative value
for BHPREF. The wells to be changed can be specified using a well name,
well name pattern or well list (WELNAME).

\begin{longtable}[]{@{}
  >{\raggedright\arraybackslash}p{(\linewidth - 6\tabcolsep) * \real{0.2500}}
  >{\raggedright\arraybackslash}p{(\linewidth - 6\tabcolsep) * \real{0.2500}}
  >{\raggedright\arraybackslash}p{(\linewidth - 6\tabcolsep) * \real{0.2500}}
  >{\raggedright\arraybackslash}p{(\linewidth - 6\tabcolsep) * \real{0.2500}}@{}}
\toprule\noalign{}
\begin{minipage}[b]{\linewidth}\raggedright
No.
\end{minipage} & \begin{minipage}[b]{\linewidth}\raggedright
Name
\end{minipage} & \begin{minipage}[b]{\linewidth}\raggedright
Description
\end{minipage} & \begin{minipage}[b]{\linewidth}\raggedright
Default
\end{minipage} \\
\midrule\noalign{}
\endhead
\bottomrule\noalign{}
\endlastfoot
Field & Metric & Laboratory & \\
1 & WELNAME & A character string of up to eight characters in length
that defines the well name for which the well specification data are
being defined. & None \\
2 & GRPNAME & A character string of up to eight characters in length
that defines the group name for which the well is assigned to. The group
named FIELD is the top most group. GRPNAME can be set to FIELD although
this is discouraged and a warning message will be issued. This is
allowed in the commercial compositional simulator but not the commercial
black-oil simulator. The FIELD group cannot contain both groups and
wells. Note that the group hierarchy should be defined by the GRUPTREE
keyword when there is more than one level of groups, otherwise all the
groups will sit directly under the FIELD group in the group tree
hierarchy. Secondly, groups defined by the GRUPTREE keyword cannot
contain other groups and wells; that is, groups must either contain
other groups or wells but not both. If necessary, wells can be
re-allocated to a different group by re-entering a well's WELSPECS data
together with a new value for GRPNAME. & None \\
3 & I & A positive integer greater than or equal to zero and less than
or equal to NX that defines the wellhead location for a vertical or
deviated well, or the heel for a horizontal well in the I-direction. For
wells being specified with the COMPTRAJ and WELTRAJ keywords in SCHEDULE
section, that allow for an alternative manner to define the well
connections to the simulation grid blocks, this parameter should be
defaulted with 1*. Since the simulator will calculate the wellhead
location from the trajectory data on the WELTRAJ keyword. Note that the
COMPTRAJ and WELTRAJ keywords are OPM Flow specific keywords, and will
cause an error in the commercial simulator. & None \\
4 & J & A positive integer greater than or equal to zero and less than
or equal to NY that defines the wellhead location for a vertical or
deviated well, or the heel for a horizontal well in the J-direction. For
wells being specified with the COMPTRAJ and WELTRAJ keywords in SCHEDULE
section, that allows for an alternative manner to define the well
connections to the simulation grid blocks, this parameter should be
defaulted with 1*. Since the simulator will calculate the wellhead
location from the trajectory data on the WELTRAJ keyword. Note that the
COMPTRAJ and WELTRAJ keywords are OPM Flow specific keywords, and will
cause an error in the commercial simulator. & None \\
5 & BHPREF & A real value that defines the reference depth for reporting
the bottom-hole pressure for the well. Ideally this value should be set
to the midpoint of the perforations as defined by the COMPDAT keyword in
the SCHEDULE section. If defaulted by 1* or set to a value less than or
equal to zero, then the mid-point of shallowest connection defined by
the COMPDAT keyword will be used. & Mid-point of shallowest connection
defined by the COMPDAT keyword \\
feet & m & cm & \\
6 & TYPE & A defined character string that defines the ``main'' phase
for the well, and should be set to one of the following character
strings: This parameter defines the phase used to calculate a well's
productivity or injectivity index and the type of well, or a well's
connection, to close when a group's production constraints, as defined
on the GCONPROD keyword in the SCHEDULE section, have been violated. For
example, if the well is declared as an oil well, then excessive gas and
water connections will be subject to closure. Note OPM Flow only
currently supports options one to three, that is option four (LIQ) is
not supported. For producing wells this mostly matters if one plots the
WPI summary vector (productivity index for well's preferred phase). In
the current treatment WPI will not have contributions from the water
phase if the declared preferred phase is LIQ. For injecting wells
WELSPECS's preferred phase does not matter, since the preferred phase is
(typically) reset to the injected phase via the WCONINJE and WCONINJH
keywords. & None \\
7 & DRADIUS & A real value that defines the well drainage radius for the
well used to calculate a well's productivity or injectivity index. A
default of zero results in the pressure equivalent radius of the grid
blocks containing the well connections are used. & 0.0 \\
feet & m & cm & \\
8 & INFLOW & A defined character string that defines the inflow equation
to be used for the well in calculating the well's flow rates. INFLOW
should be set to one of the following character strings: For oil and
water wells the INFLOW should be set to STD, whereas for dry gas wells
INFLOW can be set to either R-G or P-P; however, the P-P option is
preferred for dry gas wells due to the more rigorous treatment of gas
flow. For wet gas wells, that is gas condensate wells, INFLOW should be
set to GPP. Only INFLOW equal to STD and NO are currently implemented in
OPM Flow. & STD \\
9 & AUTO & A defined character string that specifies the action to be
taken if the well is automatically closed/shut by the simulator, and
should be set to one of the following: The corrective action takes
places at the end of the time step in which the constraint is violated.
& SHUT \\
10 & XFLOW & A defined character string that defines the if cross flow
should occur within the wellbore, and should be set to either: In some
cases numerical issues can occur if this variable is set to YES, and
resetting it to NO may resolve the issue; however, the results may not
represent the physical down hole process in this case. & YES \\
11 & PVTNUM & A positive integer greater than or equal to zero that
defines the PVT table used to calculate the wellbore fluid properties
that define the relationship between reservoir and surface volume rates.
The default value of zero sets PVTNUM to be the PVT table of the deepest
connection in the well. & 0 \\
12 & DENOPT & A defined character string that sets the type of density
calculation used in calculating the wellbore hydrostatic head, and
should be set to one of the following character strings: The default
option of 1* invokes the SEG option and is the only option implemented
in OPM Flow. & SEG \\
13 & FIPNUM & An integer value defines the FIPNUM region used to
determine the reservoir conditions in calculating the well's reservoir
volumes and is determined by: & 0 \\
14 & STRMLIN1 & Not used and should be defaulted with 1*. & 1* \\
15 & STRMLIN2 & Not used and should be defaulted with 1*. & 1* \\
16 & TYPECOMP & Commercial compositional simulator well type model
option that is not used and should be defaulted with either STD or 1*. &
STD \\
17 & POLYTAB & A positive integer greater than or equal to zero that
defines the polymer mixing table, as defined by the PLMIXPAR and PLYMAX
keywords, to be used in calculating the well's well bore properties. The
default value of zero means the table allocated via the PLMIXNUM array
for the deepest connection in the well bore is utilized. Only the
default value of zero is supported by OPM Flow. & 0 \\
Notes: & & & \\
\end{longtable}

\emph{Table 12.3.281.1: WELSPECS Keyword Description}

See also the COMPDAT keyword to define a well's connections, the
WCONPROD and WCONINJE keywords to define a well's production and
injection targets and constraints. All the aforementioned keywords are
described in the SCHEDULE section.

\subsubsection{Example}\label{example-575}

The following example defines three wells using the WELSPECS keyword

\begin{verbatim}
--
--       WELL SPECIFICATION DATA
--
-- WELL  GROUP     LOCATION  BHP    PHASE  DRAIN  INFLOW  OPEN  CROSS  PVT
-- NAME  NAME        I    J  DEPTH  FLUID  AREA   EQUANS  SHUT  FLOW   TABLE
WELSPECS
GI01     PLATFORM   14   13   1*     GAS    1*     P-P    SHUT   NO     1*     /
GP01     PLATFORM   64   80   1*     GAS    1*     GPP    SHUT   NO     1*     /
OP01     PLATFORM   24  110   1*     OIL    1*     STD    SHUT   NO     1*     /
/
\end{verbatim}

Here, well GI01 is a dry gas injection well that uses the dry gas pseudo
inflow equation, GP01 is a gas condensate well that uses the generalized
gas pseudo pressure inflow equation, and finally, OP01 is an oil well
that uses the standard inflow equation. All wells will be shut if they
are required to cease production, all wells disallow cross flow, and the
hydrostatic head calculation is defaulted to the segment option for all
wells.

If the same three wells are using the COMPTRAJ and WELTRAJ keywords to
specify the connections to the simulation grid, then the WELSPECS
keyword should be;

\begin{verbatim}
--
--       WELL SPECIFICATION DATA
--
-- WELL  GROUP     LOCATION  BHP    PHASE  DRAIN  INFLOW  OPEN  CROSS  PVT
-- NAME  NAME        I    J  DEPTH  FLUID  AREA   EQUANS  SHUT  FLOW   TABLE
WELSPECS
GI01     PLATFORM   1*   1*   1*     GAS    1*     P-P    SHUT   NO     1*     /
GP01     PLATFORM   1*   1*   1*     GAS    1*     GPP    SHUT   NO     1*     /
OP01     PLATFORM   1*   1*   1*     OIL    1*     STD    SHUT   NO     1*     /
/
\end{verbatim}

Notice how the well location parameters have been defaulted with 1* in
this case.

The following example illustrates how the WELSPECS keyword can be used
to change the controlling group for oil production wells previously
specified by the WELSPECS keyword that have an oil production rate less
than 100.

\begin{verbatim}
--
--       ACTIONX BLOCK
--
ACTIONX
  ACT01    1             /
  WOPR     'OP*' < 100.0 /
/
WELSPECS
  '?'      'LOWPRESS'   /
/
ENDACTIO
\end{verbatim}

\subsection{WELTARG -- Modify Well Target and Constraint
Values}\label{weltarg-modify-well-target-and-constraint-values}

\begin{longtable}[]{@{}
  >{\raggedright\arraybackslash}p{(\linewidth - 14\tabcolsep) * \real{0.1250}}
  >{\raggedright\arraybackslash}p{(\linewidth - 14\tabcolsep) * \real{0.1250}}
  >{\raggedright\arraybackslash}p{(\linewidth - 14\tabcolsep) * \real{0.1250}}
  >{\raggedright\arraybackslash}p{(\linewidth - 14\tabcolsep) * \real{0.1250}}
  >{\raggedright\arraybackslash}p{(\linewidth - 14\tabcolsep) * \real{0.1250}}
  >{\raggedright\arraybackslash}p{(\linewidth - 14\tabcolsep) * \real{0.1250}}
  >{\raggedright\arraybackslash}p{(\linewidth - 14\tabcolsep) * \real{0.1250}}
  >{\raggedright\arraybackslash}p{(\linewidth - 14\tabcolsep) * \real{0.1250}}@{}}
\toprule\noalign{}
\begin{minipage}[b]{\linewidth}\raggedright
RUNSPEC
\end{minipage} & \begin{minipage}[b]{\linewidth}\raggedright
GRID
\end{minipage} & \begin{minipage}[b]{\linewidth}\raggedright
EDIT
\end{minipage} & \begin{minipage}[b]{\linewidth}\raggedright
PROPS
\end{minipage} & \begin{minipage}[b]{\linewidth}\raggedright
REGIONS
\end{minipage} & \begin{minipage}[b]{\linewidth}\raggedright
SOLUTION
\end{minipage} & \begin{minipage}[b]{\linewidth}\raggedright
SUMMARY
\end{minipage} & \begin{minipage}[b]{\linewidth}\raggedright
SCHEDULE
\end{minipage} \\
\midrule\noalign{}
\endhead
\bottomrule\noalign{}
\endlastfoot
\end{longtable}

\subsubsection{Description}\label{description-1073}

The WELTARG keyword modifies the target and constraints values of both
rates and pressures for previously defined wells without having to
define all the variables on the well control keywords: WCONPROD,
WCONHIST, WCONINJE, or WCONINJH. Variables not changed by the WELTARG
keyword remain the same as those previously entered via the well control
keywords or previously entered WELTARG keywords. Note that the well must
still be initially fully defined using the WCONPROD or WCONINJE
keywords. All the aforementioned keywords are described in the SCHEDULE
section.

\begin{longtable}[]{@{}
  >{\raggedright\arraybackslash}p{(\linewidth - 6\tabcolsep) * \real{0.2500}}
  >{\raggedright\arraybackslash}p{(\linewidth - 6\tabcolsep) * \real{0.2500}}
  >{\raggedright\arraybackslash}p{(\linewidth - 6\tabcolsep) * \real{0.2500}}
  >{\raggedright\arraybackslash}p{(\linewidth - 6\tabcolsep) * \real{0.2500}}@{}}
\toprule\noalign{}
\begin{minipage}[b]{\linewidth}\raggedright
No.
\end{minipage} & \begin{minipage}[b]{\linewidth}\raggedright
Name
\end{minipage} & \begin{minipage}[b]{\linewidth}\raggedright
Description
\end{minipage} & \begin{minipage}[b]{\linewidth}\raggedright
Default
\end{minipage} \\
\midrule\noalign{}
\endhead
\bottomrule\noalign{}
\endlastfoot
Field & Metric & Laboratory & \\
1 & WELNAME & A character string of up to eight characters in length
that defines the well name for which the well production rates and
pressures data are being redefined. Note that the well name (WELNAME)
must have been declared previously using the WELSPECS and WCONPROD (or
WCONINJE) keywords in the SCHEDULE section, otherwise an error may
occur. & None \\
2 & TARGET & A defined character string that sets the item to be changed
for the well the value of the item is set by item (3). The commercial
compositional simulator options: WGRA, NGL, CVAL, REIN, STRA, SATP and
SATT are not applicable. Note that TARGET only defines the variable to
be changed, it does not change how a well is controlled. For example, if
a well is operating on ORAT control, as defined by the previously
entered WCONPROD keyword, entering TARGET equal to LRAT with a value,
changes the liquid constraint but the well still remains on ORAT
control. Use the WELCNTL keyword in the SCHEDULE section to change the
control mode of a well. & None \\
3 & VALUE Liquid Gas Res Vol Pressure VFP LIFT & A real positive value
that defines the value of the variable declared by TARGET This value may
be specified using a User Defined Argument (UDA). & None \\
stb/d Mscf/d rb/d psia dimensionless same as VFPPROD or VFPINJ & sm3/day
sm3/day rm3/day barsa dimensionless same as VFPPROD or VFPINJ & scc/hour
scc/hour rcc/hour atma dimensionless same as VFPPROD or VFPINJ & \\
Notes: & & & \\
\end{longtable}

\emph{Table 12.93: WELTARG Keyword Description}

If a well is currently a history matching well, then WELTARG should only
be used to change a well's bottom-hole pressure limit, vertical flow
performance table number or the artificial lift quantity.

See also the WELCNTL keyword, in the SCHEDULE section that can be used
to reset the control mode, as well as a well's target and constraints of
both rates and pressures.

\subsubsection{Example}\label{example-576}

The following example below shows the oil rates for the OP01 oil
producer at the start of the schedule section (January 1, 2000).

\begin{verbatim}
-- ------------------------------------------------------------------------------
-- 01 JAN 2000 START OF SCHEDULE SECTION
-- ------------------------------------------------------------------------------
--
--       WELL PRODUCTION WELL CONTROLS
--
-- WELL  OPEN/  CNTL   OIL    WAT    GAS   LIQ    RES    BHP    THP   VFP    VFP
-- NAME  SHUT   MODE   RATE   RATE   RATE  RATE   RATE   PRES   PRES  TABLE  ALFQ
WCONPROD
OP01     OPEN   ORAT   3000   1*     1*    1*     1*     750.0  500.  9      1* /
/                                                                               DATES
01 FEB 2000 /
/
--
--       WELL PRODUCTION AND INJECTION TARGETS
--
--  WELL WELL   TARGET
--  NAME TARG   VALUE
WELTARG
OP01     ORAT   2000                                        /
/
\end{verbatim}

From January 1, 2000 to February 1, 2000 well OP01 is open and is on oil
rate control and has a target oil rate of 3,000 stb/d, and uses VFPPROD
vertical lift table number 9 with a minimum tubing head pressure
constraint of 500 psia. After February 1, 2000 the well's oil rate is
reduced to 2,000 stb/d and all the other parameters remain unchanged.

\subsection{WELTRAJ -- Define Well Trajectory
Data}\label{weltraj-define-well-trajectory-data}

\begin{longtable}[]{@{}
  >{\raggedright\arraybackslash}p{(\linewidth - 14\tabcolsep) * \real{0.1250}}
  >{\raggedright\arraybackslash}p{(\linewidth - 14\tabcolsep) * \real{0.1250}}
  >{\raggedright\arraybackslash}p{(\linewidth - 14\tabcolsep) * \real{0.1250}}
  >{\raggedright\arraybackslash}p{(\linewidth - 14\tabcolsep) * \real{0.1250}}
  >{\raggedright\arraybackslash}p{(\linewidth - 14\tabcolsep) * \real{0.1250}}
  >{\raggedright\arraybackslash}p{(\linewidth - 14\tabcolsep) * \real{0.1250}}
  >{\raggedright\arraybackslash}p{(\linewidth - 14\tabcolsep) * \real{0.1250}}
  >{\raggedright\arraybackslash}p{(\linewidth - 14\tabcolsep) * \real{0.1250}}@{}}
\toprule\noalign{}
\begin{minipage}[b]{\linewidth}\raggedright
RUNSPEC
\end{minipage} & \begin{minipage}[b]{\linewidth}\raggedright
GRID
\end{minipage} & \begin{minipage}[b]{\linewidth}\raggedright
EDIT
\end{minipage} & \begin{minipage}[b]{\linewidth}\raggedright
PROPS
\end{minipage} & \begin{minipage}[b]{\linewidth}\raggedright
REGIONS
\end{minipage} & \begin{minipage}[b]{\linewidth}\raggedright
SOLUTION
\end{minipage} & \begin{minipage}[b]{\linewidth}\raggedright
SUMMARY
\end{minipage} & \begin{minipage}[b]{\linewidth}\raggedright
SCHEDULE
\end{minipage} \\
\midrule\noalign{}
\endhead
\bottomrule\noalign{}
\endlastfoot
\end{longtable}

\subsubsection{Description}\label{description-1074}

The WELTRAJ keyword defines a trajectory well together with the well
trajectory data (simplified directional survey data), and is used in
conjunction with the COMPTRAJ keyword in the SCHEDULE section to define
the well connections to the simulation grid blocks. The keyword can only
be used for trajectory wells that employ the COMPTRAJ keyword to define
the connections to the grid, that is, one cannot use COMPDAT keyword in
the SCHEDULE section for declaring the connections to the grid for these
type of wells.

Although WELTRAJ and COMPTRAJ keywords are sufficient to define the
wellbore path and connections to the grid, it is still necessary to
defined the general well specification parameters using the WELSPECS
keyword in the SCHEDULE section. In this case, the wellhead location
parameters, WELSPECS(I, J), should be defaulted with 1*.

\begin{tcolorbox}[enhanced jigsaw, title=\textcolor{quarto-callout-note-color}{\faInfo}\hspace{0.5em}{Note}, coltitle=black, breakable, toprule=.15mm, colbacktitle=quarto-callout-note-color!10!white, leftrule=.75mm, opacitybacktitle=0.6, colback=white, bottomtitle=1mm, toptitle=1mm, titlerule=0mm, left=2mm, colframe=quarto-callout-note-color-frame, arc=.35mm, rightrule=.15mm, bottomrule=.15mm, opacityback=0]

This is an OPM Flow specific keyword and will therefore cause an error
in the commercial simulator.

\end{tcolorbox}

\begin{longtable}[]{@{}
  >{\raggedright\arraybackslash}p{(\linewidth - 6\tabcolsep) * \real{0.2500}}
  >{\raggedright\arraybackslash}p{(\linewidth - 6\tabcolsep) * \real{0.2500}}
  >{\raggedright\arraybackslash}p{(\linewidth - 6\tabcolsep) * \real{0.2500}}
  >{\raggedright\arraybackslash}p{(\linewidth - 6\tabcolsep) * \real{0.2500}}@{}}
\toprule\noalign{}
\begin{minipage}[b]{\linewidth}\raggedright
No.
\end{minipage} & \begin{minipage}[b]{\linewidth}\raggedright
Name
\end{minipage} & \begin{minipage}[b]{\linewidth}\raggedright
Description
\end{minipage} & \begin{minipage}[b]{\linewidth}\raggedright
Default
\end{minipage} \\
\midrule\noalign{}
\endhead
\bottomrule\noalign{}
\endlastfoot
Field & Metric & Laboratory & \\
1 & WELNAME & A character string of up to eight characters in length
that defines the well name for which the well trajectory data are being
defined. Note that the well name (WELNAME) must have been declared
previously using the WELSPECS keyword in the SCHEDULE section, otherwise
an error may occur. Secondly, the wellhead location parameters on the
WELSPECS keyword, WELSPECS(I, J), should be defaulted with 1* for
trajectory wells, as the well location will be calculated by the
simulator. & None \\
2 & IBRANCH & A positive integer greater than or equal to one and less
than or equal to MXBRAN on WSEGDIMS keyword in the RUNSPEC section that
defines the branch number of a segment. All segments on the main stem
must have IBRANCH set to one and lateral branches should have values
between two and MXSEGS on the WSEGDIMS keyword in the RUNSPEC section.
Only the default value of one is currently supported, that is only the
main branch of a multi-segment well is supported, or a single trajectory
for a conventional well. & 1 \\
3 ‍ & XCORD & A real positive value representing the X coordinate in
three dimensional space of the well trajectory path. The Geodesy
Reference System should be same as the static model used to generate the
dynamic model. & None \\
feet & m & cm & \\
4 ‍ & YCORD & A real positive value representing the Y coordinate in
three dimensional space of the well trajectory path. Again, the Geodesy
Reference System should be same as the static model used to generate the
dynamic model. & None \\
feet & m & cm & \\
5 & TVD & A real negative or positive value that defines the True
Vertical Depth along the wellbore path at XCORD and YCORD. Normally, TVD
is referenced to the subsea elevation, that is TVDSS. & None \\
feet & m & cm & \\
6 ‍ & MD & A real positive value that defines the equivalent measured
depth along the wellbore path at TVD. Normally, MD is referenced to a
reference point on the drilling rig that drilled the well, for example,
Relative to Kelly Bushing, MDRKB. & None \\
feet & m & cm & \\
Notes: & & & \\
\end{longtable}

\emph{Table 12.94: WELTRAJ Keyword Description}

Using the WELTRAJ and COMPTRAJ keywords to define wells and how they are
connected to grid, offers several advantages compared to the
conventional approach based on the (I, J, K) co-ordinates of the grid.
The approach allows for the wells to be independent of the grid, which
is particularly useful when running ensemble cases, as the well
connections are no longer required to be re-calculated for each ensemble
case. In addition, quality control of the model is improved by using
consistent perforation data in both the static and dynamic models.

See the example on the following page.

\subsubsection{Example}\label{example-577}

The following example defines two trajectory wells oil wells, OP01 and
OP02, using the WELSPECS and WELTRAJ keywords, together with their
perforations using the COMPTRAJ keyword.

\begin{verbatim}
--
--       WELL SPECIFICATION DATA
--
-- WELL  GROUP     LOCATION  BHP    PHASE  DRAIN  INFLOW  OPEN  CROSS  PVT
-- NAME  NAME        I    J  DEPTH  FLUID  AREA   EQUANS  SHUT  FLOW   TABLE
WELSPECS
OP01     PLATFORM   1*   1*   1*     OIL    1*     STD    SHUT   NO     1*     /
OP02     PLATFORM   1*   1*   1*     OIL    1*     STD    SHUT   NO     1*     /
/
--
--       WELL TRAJECTORY DATA
--
-- WELL   BRAN  XCORD         YCORD         TVDSS         MD
-- NAME   NO                                DEPTH         DEPTH
-- -----  ----  ------------  ------------  ------------  ------------
WELTRAJ
OP01      1*    2.805445e+06  3.602948e+06  -100.000000   0.0         /
OP01      1*    2.805445e+06  3.602948e+06  877.0000000   977.0       /
OP01      1*    2.805445e+06  3.602948e+06  957.9950240   1058.0      /
OP01      1*    2.805444e+06  3.602946e+06  1051.976081   1152.0      /
…………………...                                                            /
OP02      1*    2.810828e+06  3.604507e+06  9371.792711   11418.0     /
OP02      1*    2.810885e+06  3.604525e+06  9443.657000   11511.0     /
OP02      1*    2.810952e+06  3.604546e+06  9531.966162   11624.0     /
OP02      1*    2.810973e+06  3.604553e+06  9560.411742   11660.0     /
/
--
--
--       WELL TRAJECTORY CONNECTION DATA
--
-- WELL  BRAN  -- PERFORATION --  COMPL OPEN  SAT CONN  WELL   KH    SKIN  D
-- NAME  NO.    TOP    BOT   REF  NO.   SHUT  TAB FACT  DIA    FACT  FACT  FACT
COMPTRAJ
OP01     1*     8230   8244  MD    1    SHUT  1*  1*    0.708  1*    0.0   1*  /
OP01     1*     8352   8380  MD    1    SHUT  1*  1*    0.708  1*    0.0   1*  /
OP01     1*     9070   9100  MD    1    SHUT  1*  1*    0.708  1*    0.0   1*  /
OP01     1*     9220   9250  MD    2    SHUT  1*  1*    0.708  1*    0.0   1*  /
OP01     1*     9266   9280  MD    2    SHUT  1*  1*    0.708  1*    0.0   1*  /
OP01     1*     9693   9703  MD    3    SHUT  1*  1*    0.708  1*    0.0   1*  /
OP01     1*     9940   9974  MD    3    SHUT  1*  1*    0.708  1*    0.0   1*  /
OP02     1*     9979   9985  TVD   1*   SHUT  1*  1*    0.708  1*    0.0   1*  /
OP02     1*    10173  10183  TVD   1*   SHUT  1*  1*    0.708  1*    0.0   1*  /
OP02     1*    10190  10204  TVD   1*   SHUT  1*  1*    0.708  1*    0.0   1*  /
OP02     1*    10327  10333  TVD   1*   SHUT  1*  1*    0.708  1*    0.0   1*  /
OP02     1*    10339  10345  TVD   1*   SHUT  1*  1*    0.708  1*    0.0   1*  /
OP02     1*    11528  11538  TVD   1*   SHUT  1*  1*    0.708  1*    0.0   1*  /
\end{verbatim}

Here, well OP01 has eight perforation intervals, with the intervals one
to three grouped into one completion, perforation intervals four to five
grouped into completion number two, and finally the bottom three
perforations are grouped into completion number three. In contrast, OP02
has six perforated intervals with their completion interval defaulted to
one.

\subsection{WFOAM -- Define Well Foam Injection
Concentrations}\label{wfoam-define-well-foam-injection-concentrations}

\begin{longtable}[]{@{}
  >{\raggedright\arraybackslash}p{(\linewidth - 14\tabcolsep) * \real{0.1250}}
  >{\raggedright\arraybackslash}p{(\linewidth - 14\tabcolsep) * \real{0.1250}}
  >{\raggedright\arraybackslash}p{(\linewidth - 14\tabcolsep) * \real{0.1250}}
  >{\raggedright\arraybackslash}p{(\linewidth - 14\tabcolsep) * \real{0.1250}}
  >{\raggedright\arraybackslash}p{(\linewidth - 14\tabcolsep) * \real{0.1250}}
  >{\raggedright\arraybackslash}p{(\linewidth - 14\tabcolsep) * \real{0.1250}}
  >{\raggedright\arraybackslash}p{(\linewidth - 14\tabcolsep) * \real{0.1250}}
  >{\raggedright\arraybackslash}p{(\linewidth - 14\tabcolsep) * \real{0.1250}}@{}}
\toprule\noalign{}
\begin{minipage}[b]{\linewidth}\raggedright
RUNSPEC
\end{minipage} & \begin{minipage}[b]{\linewidth}\raggedright
GRID
\end{minipage} & \begin{minipage}[b]{\linewidth}\raggedright
EDIT
\end{minipage} & \begin{minipage}[b]{\linewidth}\raggedright
PROPS
\end{minipage} & \begin{minipage}[b]{\linewidth}\raggedright
REGIONS
\end{minipage} & \begin{minipage}[b]{\linewidth}\raggedright
SOLUTION
\end{minipage} & \begin{minipage}[b]{\linewidth}\raggedright
SUMMARY
\end{minipage} & \begin{minipage}[b]{\linewidth}\raggedright
SCHEDULE
\end{minipage} \\
\midrule\noalign{}
\endhead
\bottomrule\noalign{}
\endlastfoot
\end{longtable}

\subsubsection{Description}\label{description-1075}

The WFOAM keyword defines an injection wells foam concentration. The
foam option must be activated by the FOAM keyword in the RUNSPEC section
in order to use this keyword. Note if a well's foam concentration is not
set with this keyword then default value of zero is assigned to a well.

\begin{longtable}[]{@{}
  >{\raggedright\arraybackslash}p{(\linewidth - 6\tabcolsep) * \real{0.2500}}
  >{\raggedright\arraybackslash}p{(\linewidth - 6\tabcolsep) * \real{0.2500}}
  >{\raggedright\arraybackslash}p{(\linewidth - 6\tabcolsep) * \real{0.2500}}
  >{\raggedright\arraybackslash}p{(\linewidth - 6\tabcolsep) * \real{0.2500}}@{}}
\toprule\noalign{}
\begin{minipage}[b]{\linewidth}\raggedright
No.
\end{minipage} & \begin{minipage}[b]{\linewidth}\raggedright
Name
\end{minipage} & \begin{minipage}[b]{\linewidth}\raggedright
Description
\end{minipage} & \begin{minipage}[b]{\linewidth}\raggedright
Default
\end{minipage} \\
\midrule\noalign{}
\endhead
\bottomrule\noalign{}
\endlastfoot
Field & Metric & Laboratory & \\
1 & WELNAME & A character string of up to eight characters in length
that defines the well name for which the well injection foam
concentration is being defined. & None \\
2 & FOAMCON & A real positive value that defines the well's injection
foam concentration. This value may be specified using a User Defined
Argument (UDA). Units are dependent on the transport phase specified via
the FOAMOPT1 variable on the FOAMOPTS keyword in the PROPS section.
FOAMOPT1 should be set to either GAS or WATER. Currently OPM Flow only
supports injecting foam via the GAS phase. & None \\
Gas: lb/Mscf Water: lb/stb & Gas: kg/sm3 Water: kg/sm3 & Gas: gm/scc
Water: gm/scc & \\
Notes: & & & \\
\end{longtable}

\emph{Table 12.95: WFOAM Keyword Description}

See also the FOAM keyword in the RUNSPEC section, the FOAMADS, FOAMOPTS
and FOAMROCK keywords in the PROPS section.

\subsubsection{Example}\label{example-578}

\begin{verbatim}
--
--       WELL INJECTION FOAM CONCENTRATION
--
-- WELL  FOAM
-- NAME  FOAMCON
WFOAM
GI01     0.020                                                                 /
GI02     0.020                                                                 /
GI03     0.020                                                                 /
/
\end{verbatim}

Here three gas wells are given an injection foam concentration of 0.020
lb/Mscf, assuming field units.

\subsection{WFRICSEG -- Convert Friction Well to Multi-Segment
Well}\label{wfricseg-convert-friction-well-to-multi-segment-well}

\begin{longtable}[]{@{}
  >{\raggedright\arraybackslash}p{(\linewidth - 14\tabcolsep) * \real{0.1250}}
  >{\raggedright\arraybackslash}p{(\linewidth - 14\tabcolsep) * \real{0.1250}}
  >{\raggedright\arraybackslash}p{(\linewidth - 14\tabcolsep) * \real{0.1250}}
  >{\raggedright\arraybackslash}p{(\linewidth - 14\tabcolsep) * \real{0.1250}}
  >{\raggedright\arraybackslash}p{(\linewidth - 14\tabcolsep) * \real{0.1250}}
  >{\raggedright\arraybackslash}p{(\linewidth - 14\tabcolsep) * \real{0.1250}}
  >{\raggedright\arraybackslash}p{(\linewidth - 14\tabcolsep) * \real{0.1250}}
  >{\raggedright\arraybackslash}p{(\linewidth - 14\tabcolsep) * \real{0.1250}}@{}}
\toprule\noalign{}
\begin{minipage}[b]{\linewidth}\raggedright
RUNSPEC
\end{minipage} & \begin{minipage}[b]{\linewidth}\raggedright
GRID
\end{minipage} & \begin{minipage}[b]{\linewidth}\raggedright
EDIT
\end{minipage} & \begin{minipage}[b]{\linewidth}\raggedright
PROPS
\end{minipage} & \begin{minipage}[b]{\linewidth}\raggedright
REGIONS
\end{minipage} & \begin{minipage}[b]{\linewidth}\raggedright
SOLUTION
\end{minipage} & \begin{minipage}[b]{\linewidth}\raggedright
SUMMARY
\end{minipage} & \begin{minipage}[b]{\linewidth}\raggedright
SCHEDULE
\end{minipage} \\
\midrule\noalign{}
\endhead
\bottomrule\noalign{}
\endlastfoot
\end{longtable}

\subsubsection{Description}\label{description-1076}

WFRICSEG converts a previously defined friction well, as per the WFRICTN
keyword in the SCHEDULE section, to a multi-segment well. The keyword
thus acts as a replacement for the WELSEGS and COMPSEGS keywords for
multi-segment wells. See also the WFRICSGL keyword in the SCHEDULE
section that performs similar functionality for wells in Local Grid
Refinements.

This keyword is not supported by OPM Flow but would change the results
if supported so the simulation will be stopped.

\subsection{WFRICSGL -- Convert Friction Well to Multi-Segment Well
(LGR)}\label{wfricsgl-convert-friction-well-to-multi-segment-well-lgr}

\begin{longtable}[]{@{}
  >{\raggedright\arraybackslash}p{(\linewidth - 14\tabcolsep) * \real{0.1250}}
  >{\raggedright\arraybackslash}p{(\linewidth - 14\tabcolsep) * \real{0.1250}}
  >{\raggedright\arraybackslash}p{(\linewidth - 14\tabcolsep) * \real{0.1250}}
  >{\raggedright\arraybackslash}p{(\linewidth - 14\tabcolsep) * \real{0.1250}}
  >{\raggedright\arraybackslash}p{(\linewidth - 14\tabcolsep) * \real{0.1250}}
  >{\raggedright\arraybackslash}p{(\linewidth - 14\tabcolsep) * \real{0.1250}}
  >{\raggedright\arraybackslash}p{(\linewidth - 14\tabcolsep) * \real{0.1250}}
  >{\raggedright\arraybackslash}p{(\linewidth - 14\tabcolsep) * \real{0.1250}}@{}}
\toprule\noalign{}
\begin{minipage}[b]{\linewidth}\raggedright
RUNSPEC
\end{minipage} & \begin{minipage}[b]{\linewidth}\raggedright
GRID
\end{minipage} & \begin{minipage}[b]{\linewidth}\raggedright
EDIT
\end{minipage} & \begin{minipage}[b]{\linewidth}\raggedright
PROPS
\end{minipage} & \begin{minipage}[b]{\linewidth}\raggedright
REGIONS
\end{minipage} & \begin{minipage}[b]{\linewidth}\raggedright
SOLUTION
\end{minipage} & \begin{minipage}[b]{\linewidth}\raggedright
SUMMARY
\end{minipage} & \begin{minipage}[b]{\linewidth}\raggedright
SCHEDULE
\end{minipage} \\
\midrule\noalign{}
\endhead
\bottomrule\noalign{}
\endlastfoot
\end{longtable}

\subsubsection{Description}\label{description-1077}

WFRICSGL converts a previously defined Local Grid Refinement (``LGR'')
friction well, as per the WFRICTNL keyword in the SCHEDULE section, to a
multi-segment LGR well. The keyword thus acts as a replacement for the
WELSEGS and COMPSEGL keywords for LGR multi-segment wells. See also the
WFRICSEG keyword in the SCHEDULE section that performs similar
functionality for wells in the global grid.

This keyword is not supported by OPM Flow but would change the results
if supported so the simulation will be stopped.

\subsection{WFRICTN -- Define Well as a Friction
Well}\label{wfrictn-define-well-as-a-friction-well}

\begin{longtable}[]{@{}
  >{\raggedright\arraybackslash}p{(\linewidth - 14\tabcolsep) * \real{0.1250}}
  >{\raggedright\arraybackslash}p{(\linewidth - 14\tabcolsep) * \real{0.1250}}
  >{\raggedright\arraybackslash}p{(\linewidth - 14\tabcolsep) * \real{0.1250}}
  >{\raggedright\arraybackslash}p{(\linewidth - 14\tabcolsep) * \real{0.1250}}
  >{\raggedright\arraybackslash}p{(\linewidth - 14\tabcolsep) * \real{0.1250}}
  >{\raggedright\arraybackslash}p{(\linewidth - 14\tabcolsep) * \real{0.1250}}
  >{\raggedright\arraybackslash}p{(\linewidth - 14\tabcolsep) * \real{0.1250}}
  >{\raggedright\arraybackslash}p{(\linewidth - 14\tabcolsep) * \real{0.1250}}@{}}
\toprule\noalign{}
\begin{minipage}[b]{\linewidth}\raggedright
RUNSPEC
\end{minipage} & \begin{minipage}[b]{\linewidth}\raggedright
GRID
\end{minipage} & \begin{minipage}[b]{\linewidth}\raggedright
EDIT
\end{minipage} & \begin{minipage}[b]{\linewidth}\raggedright
PROPS
\end{minipage} & \begin{minipage}[b]{\linewidth}\raggedright
REGIONS
\end{minipage} & \begin{minipage}[b]{\linewidth}\raggedright
SOLUTION
\end{minipage} & \begin{minipage}[b]{\linewidth}\raggedright
SUMMARY
\end{minipage} & \begin{minipage}[b]{\linewidth}\raggedright
SCHEDULE
\end{minipage} \\
\midrule\noalign{}
\endhead
\bottomrule\noalign{}
\endlastfoot
\end{longtable}

\subsubsection{Description}\label{description-1078}

The WFRICTN keyword is used to declare a previously defined well as a
friction well and to set the characteristics for this type of well
including: tubing size, pipe roughness, and the connections to the grid.
Wellbore friction is important in horizontal and multi-lateral wells
where the pressure loss along the pipe can effect a well's performance.
Note that unlike other SCHEDULE section well keywords, multiple wells
cannot be entered with one WFRICTN keyword, that is, the keyword must be
repeated for each well.

See also the WFRICTNL keyword in the SCHEDULE section that performs
similar functionality for wells in Local Grid Refinements.

This keyword is not supported by OPM Flow but would change the results
if supported so the simulation will be stopped.

\subsection{WFRICTNL -- Define Well as a Friction Well
(LGR)}\label{wfrictnl-define-well-as-a-friction-well-lgr}

\begin{longtable}[]{@{}
  >{\raggedright\arraybackslash}p{(\linewidth - 14\tabcolsep) * \real{0.1250}}
  >{\raggedright\arraybackslash}p{(\linewidth - 14\tabcolsep) * \real{0.1250}}
  >{\raggedright\arraybackslash}p{(\linewidth - 14\tabcolsep) * \real{0.1250}}
  >{\raggedright\arraybackslash}p{(\linewidth - 14\tabcolsep) * \real{0.1250}}
  >{\raggedright\arraybackslash}p{(\linewidth - 14\tabcolsep) * \real{0.1250}}
  >{\raggedright\arraybackslash}p{(\linewidth - 14\tabcolsep) * \real{0.1250}}
  >{\raggedright\arraybackslash}p{(\linewidth - 14\tabcolsep) * \real{0.1250}}
  >{\raggedright\arraybackslash}p{(\linewidth - 14\tabcolsep) * \real{0.1250}}@{}}
\toprule\noalign{}
\begin{minipage}[b]{\linewidth}\raggedright
RUNSPEC
\end{minipage} & \begin{minipage}[b]{\linewidth}\raggedright
GRID
\end{minipage} & \begin{minipage}[b]{\linewidth}\raggedright
EDIT
\end{minipage} & \begin{minipage}[b]{\linewidth}\raggedright
PROPS
\end{minipage} & \begin{minipage}[b]{\linewidth}\raggedright
REGIONS
\end{minipage} & \begin{minipage}[b]{\linewidth}\raggedright
SOLUTION
\end{minipage} & \begin{minipage}[b]{\linewidth}\raggedright
SUMMARY
\end{minipage} & \begin{minipage}[b]{\linewidth}\raggedright
SCHEDULE
\end{minipage} \\
\midrule\noalign{}
\endhead
\bottomrule\noalign{}
\endlastfoot
\end{longtable}

\subsubsection{Description}\label{description-1079}

The WFRICTNL keyword is used to declare a previously defined Local Grid
Refinement (``LGR'') well as a LGR friction well and to set the
characteristics for this type of well including: tubing size, pipe
roughness, and the connections to the grid. Wellbore friction is
important in horizontal and multi-lateral wells where the pressure loss
along the pipe can effect a well's performance. Note that unlike other
SCHEDULE section well keywords, multiple wells cannot be entered with
one WFRICTNL keyword, that is, the keyword must be repeated for each
well.

See also the WFRICTN keyword in the SCHEDULE section that performs
similar functionality for wells in the global grid.

This keyword is not supported by OPM Flow but would change the results
if supported so the simulation will be stopped.

\subsection{WGASPROD -- Define Sale Gas Well Production
Targets}\label{wgasprod-define-sale-gas-well-production-targets}

\begin{longtable}[]{@{}
  >{\raggedright\arraybackslash}p{(\linewidth - 14\tabcolsep) * \real{0.1250}}
  >{\raggedright\arraybackslash}p{(\linewidth - 14\tabcolsep) * \real{0.1250}}
  >{\raggedright\arraybackslash}p{(\linewidth - 14\tabcolsep) * \real{0.1250}}
  >{\raggedright\arraybackslash}p{(\linewidth - 14\tabcolsep) * \real{0.1250}}
  >{\raggedright\arraybackslash}p{(\linewidth - 14\tabcolsep) * \real{0.1250}}
  >{\raggedright\arraybackslash}p{(\linewidth - 14\tabcolsep) * \real{0.1250}}
  >{\raggedright\arraybackslash}p{(\linewidth - 14\tabcolsep) * \real{0.1250}}
  >{\raggedright\arraybackslash}p{(\linewidth - 14\tabcolsep) * \real{0.1250}}@{}}
\toprule\noalign{}
\begin{minipage}[b]{\linewidth}\raggedright
RUNSPEC
\end{minipage} & \begin{minipage}[b]{\linewidth}\raggedright
GRID
\end{minipage} & \begin{minipage}[b]{\linewidth}\raggedright
EDIT
\end{minipage} & \begin{minipage}[b]{\linewidth}\raggedright
PROPS
\end{minipage} & \begin{minipage}[b]{\linewidth}\raggedright
REGIONS
\end{minipage} & \begin{minipage}[b]{\linewidth}\raggedright
SOLUTION
\end{minipage} & \begin{minipage}[b]{\linewidth}\raggedright
SUMMARY
\end{minipage} & \begin{minipage}[b]{\linewidth}\raggedright
SCHEDULE
\end{minipage} \\
\midrule\noalign{}
\endhead
\bottomrule\noalign{}
\endlastfoot
\end{longtable}

\subsubsection{Description}\label{description-1080}

The WGASPROD keyword declares wells to be Sales Gas producers and sets
the incremental gas rate for a well and the maximum number of increments
that this rate can be increased. Wells must have been previously been
defined via the WELSPECS and WCONPROD keywords in the SCHEDULE section
and are subject to any targets or constraints on WCONPROD keyword.

This keyword is not supported by OPM Flow but would change the results
if supported so the simulation will be stopped.

\subsection{WGORPEN -- Define Well GOR Penalty
Parameters}\label{wgorpen-define-well-gor-penalty-parameters}

\begin{longtable}[]{@{}
  >{\raggedright\arraybackslash}p{(\linewidth - 14\tabcolsep) * \real{0.1250}}
  >{\raggedright\arraybackslash}p{(\linewidth - 14\tabcolsep) * \real{0.1250}}
  >{\raggedright\arraybackslash}p{(\linewidth - 14\tabcolsep) * \real{0.1250}}
  >{\raggedright\arraybackslash}p{(\linewidth - 14\tabcolsep) * \real{0.1250}}
  >{\raggedright\arraybackslash}p{(\linewidth - 14\tabcolsep) * \real{0.1250}}
  >{\raggedright\arraybackslash}p{(\linewidth - 14\tabcolsep) * \real{0.1250}}
  >{\raggedright\arraybackslash}p{(\linewidth - 14\tabcolsep) * \real{0.1250}}
  >{\raggedright\arraybackslash}p{(\linewidth - 14\tabcolsep) * \real{0.1250}}@{}}
\toprule\noalign{}
\begin{minipage}[b]{\linewidth}\raggedright
RUNSPEC
\end{minipage} & \begin{minipage}[b]{\linewidth}\raggedright
GRID
\end{minipage} & \begin{minipage}[b]{\linewidth}\raggedright
EDIT
\end{minipage} & \begin{minipage}[b]{\linewidth}\raggedright
PROPS
\end{minipage} & \begin{minipage}[b]{\linewidth}\raggedright
REGIONS
\end{minipage} & \begin{minipage}[b]{\linewidth}\raggedright
SOLUTION
\end{minipage} & \begin{minipage}[b]{\linewidth}\raggedright
SUMMARY
\end{minipage} & \begin{minipage}[b]{\linewidth}\raggedright
SCHEDULE
\end{minipage} \\
\midrule\noalign{}
\endhead
\bottomrule\noalign{}
\endlastfoot
\end{longtable}

\subsubsection{Description}\label{description-1081}

The WGORPEN keyword defines a well's Gas-Oil Ratio (``GOR'') penalty
parameters used to calculate a well's oil production target for the
current month, as a function of the well's previous month's average GOR.
The WGORPEN calculated oil rate overwrites any oil targets set by the
WCONPROD and WELTARG keywords in the SCHEDULE section. In North
American, it is common practice for the regulator to enforce GOR
penalties, in order to control gas production in depletion drive oil
reservoirs, with the stated intention to maximize oil recovery by
limiting the energy loss from the reservoir by excessive gas production.

This keyword is not supported by OPM Flow but would change the results
if supported so the simulation will be stopped.

\subsection{WGRUPCON -- Define Well Guide Rates for Group
Control}\label{wgrupcon-define-well-guide-rates-for-group-control}

\begin{longtable}[]{@{}
  >{\raggedright\arraybackslash}p{(\linewidth - 14\tabcolsep) * \real{0.1250}}
  >{\raggedright\arraybackslash}p{(\linewidth - 14\tabcolsep) * \real{0.1250}}
  >{\raggedright\arraybackslash}p{(\linewidth - 14\tabcolsep) * \real{0.1250}}
  >{\raggedright\arraybackslash}p{(\linewidth - 14\tabcolsep) * \real{0.1250}}
  >{\raggedright\arraybackslash}p{(\linewidth - 14\tabcolsep) * \real{0.1250}}
  >{\raggedright\arraybackslash}p{(\linewidth - 14\tabcolsep) * \real{0.1250}}
  >{\raggedright\arraybackslash}p{(\linewidth - 14\tabcolsep) * \real{0.1250}}
  >{\raggedright\arraybackslash}p{(\linewidth - 14\tabcolsep) * \real{0.1250}}@{}}
\toprule\noalign{}
\begin{minipage}[b]{\linewidth}\raggedright
RUNSPEC
\end{minipage} & \begin{minipage}[b]{\linewidth}\raggedright
GRID
\end{minipage} & \begin{minipage}[b]{\linewidth}\raggedright
EDIT
\end{minipage} & \begin{minipage}[b]{\linewidth}\raggedright
PROPS
\end{minipage} & \begin{minipage}[b]{\linewidth}\raggedright
REGIONS
\end{minipage} & \begin{minipage}[b]{\linewidth}\raggedright
SOLUTION
\end{minipage} & \begin{minipage}[b]{\linewidth}\raggedright
SUMMARY
\end{minipage} & \begin{minipage}[b]{\linewidth}\raggedright
SCHEDULE
\end{minipage} \\
\midrule\noalign{}
\endhead
\bottomrule\noalign{}
\endlastfoot
\end{longtable}

\subsubsection{Description}\label{description-1082}

The WGRUPCON keyword defines a well's production or injection guide rate
for when a well is under group control. The guide rate is used to
determine a well's production target under group control in order to
satisfy a group's targets and constraints, including any higher level
related groups as well as the FIELD group.

Wells must have been previously defined and allocated to a group by the
WELSPECS keyword in the SCHEDULE section. Wells defined to be under
group control will have their production rates controlled by the group
to which they belong, in addition to any well constraints defined for
the wells via the WCONPROD and WCONINJE keywords in the SCHEDULE
section.

\begin{longtable}[]{@{}
  >{\raggedright\arraybackslash}p{(\linewidth - 6\tabcolsep) * \real{0.2500}}
  >{\raggedright\arraybackslash}p{(\linewidth - 6\tabcolsep) * \real{0.2500}}
  >{\raggedright\arraybackslash}p{(\linewidth - 6\tabcolsep) * \real{0.2500}}
  >{\raggedright\arraybackslash}p{(\linewidth - 6\tabcolsep) * \real{0.2500}}@{}}
\toprule\noalign{}
\begin{minipage}[b]{\linewidth}\raggedright
No.
\end{minipage} & \begin{minipage}[b]{\linewidth}\raggedright
Name
\end{minipage} & \begin{minipage}[b]{\linewidth}\raggedright
Description
\end{minipage} & \begin{minipage}[b]{\linewidth}\raggedright
Default
\end{minipage} \\
\midrule\noalign{}
\endhead
\bottomrule\noalign{}
\endlastfoot
Field & Metric & Laboratory & \\
1 & WELNAME & A character string of up to eight characters in length
that defines the well name for which the well production targets and
constraints data are being defined. Note that the well name (WELNAME)
must have been declared previously using the WELSPECS keyword in the
SCHEDULE section, otherwise an error may occur. & None \\
2 & STATUS & A defined character string that declares the status of the
well to be under group control or not under group control. STATUS should
be set to one of the following character strings: Note the default value
of YES puts all wells under group control unless specified otherwise by
the STATUS variable, or the TARGET variable on the WCONPROD and WCONINJE
keywords in the SCHEDULE section. & YES \\
3 & GUIDERAT & A dimensionless real number that determines the well's
share of its group production (or injection) target rate. If GUIDERAT is
a positive number then the guide rate for the well is fixed until
modified by this keyword at a subsequent time. If TARGET variable on
this keyword is not equal to the group's controlling phase, then the
GUIDERAT is converted into the groups' controlling phase and is updated
every time step. If GUIDERAT is less than or equal to zero then the
well's guide rate is based on the well's potential (unrestricted flow)
and the potential is calculated every time step. & -1.0 \\
dimensionless & dimensionless & dimensionless & \\
4 & TARGET & A defined character string that sets the well's guide rate
phase that the GUIDERAT value should be applied to. TARGET should be set
to one of the following character strings: TARGET may be defaulted if
GUIDERAT has been defaulted, either by 1* or a value less than or equal
to zero. & None \\
5 & SCALE & A real value that is used to multiple the GUIDERAT or the
calculated well potentials to determine the final GUIDERAT for the well.
& 1.0 \\
dimensionless & dimensionless & dimensionless & \\
Notes: & & & \\
\end{longtable}

\emph{Table 12.96: WGRUPCON Keyword Description}

See also the GCONPROD the GCONINJE keywords to define a group's
production and injection targets and constraints, and the WCONPROD and
WCONINJE keyword to define a well's production and injection
characteristics. All the aforementioned keywords are described in the
SCHEDULE section.

\subsubsection{Example}\label{example-579}

The following example defines the guides rates for all oil and gas
producers and the gas injectors as follows:

\begin{verbatim}
--
--       DEFINE WELL GUIDES FOR GROUP CONTROL
--
-- WELL  GRUP  GUIDE  GUIDE  SCALE
-- NAME  CNTL  RATE   PHASE  FACT
WGRUPCON
'GI*'    YES   0      RAT    1.0                           /
'GP*'    YES   0      GAS    1.0                           /
'OP*'    NO    2      OIL    1.0                           /
/
\end{verbatim}

Both the gas producers ('GP\emph{') and injectors (`GI'}) are under
group control with their guide rates based on their potentials. The gas
injectors are controlled based on their potential surface gas injection
rates and the gas producers on their potential surface gas production
rates. In comparison, the oil wells (OP*) are controlled by their own
targets and constraints.

\subsection{WHEDREFD -- Define Well Hydraulic Head Reference
Depth}\label{whedrefd-define-well-hydraulic-head-reference-depth}

\begin{longtable}[]{@{}
  >{\raggedright\arraybackslash}p{(\linewidth - 14\tabcolsep) * \real{0.1250}}
  >{\raggedright\arraybackslash}p{(\linewidth - 14\tabcolsep) * \real{0.1250}}
  >{\raggedright\arraybackslash}p{(\linewidth - 14\tabcolsep) * \real{0.1250}}
  >{\raggedright\arraybackslash}p{(\linewidth - 14\tabcolsep) * \real{0.1250}}
  >{\raggedright\arraybackslash}p{(\linewidth - 14\tabcolsep) * \real{0.1250}}
  >{\raggedright\arraybackslash}p{(\linewidth - 14\tabcolsep) * \real{0.1250}}
  >{\raggedright\arraybackslash}p{(\linewidth - 14\tabcolsep) * \real{0.1250}}
  >{\raggedright\arraybackslash}p{(\linewidth - 14\tabcolsep) * \real{0.1250}}@{}}
\toprule\noalign{}
\begin{minipage}[b]{\linewidth}\raggedright
RUNSPEC
\end{minipage} & \begin{minipage}[b]{\linewidth}\raggedright
GRID
\end{minipage} & \begin{minipage}[b]{\linewidth}\raggedright
EDIT
\end{minipage} & \begin{minipage}[b]{\linewidth}\raggedright
PROPS
\end{minipage} & \begin{minipage}[b]{\linewidth}\raggedright
REGIONS
\end{minipage} & \begin{minipage}[b]{\linewidth}\raggedright
SOLUTION
\end{minipage} & \begin{minipage}[b]{\linewidth}\raggedright
SUMMARY
\end{minipage} & \begin{minipage}[b]{\linewidth}\raggedright
SCHEDULE
\end{minipage} \\
\midrule\noalign{}
\endhead
\bottomrule\noalign{}
\endlastfoot
\end{longtable}

\subsubsection{Description}\label{description-1083}

The WHEDREFD keyword sets the hydraulic head reference depth for
reporting the hydraulic head pressure for the well, for wells that have
previously been defined by the WELSPECS keyword in the SCHEDULE section.

This keyword is not supported by OPM Flow but would change the results
if supported so the simulation will be stopped.

\begin{longtable}[]{@{}
  >{\raggedright\arraybackslash}p{(\linewidth - 6\tabcolsep) * \real{0.2500}}
  >{\raggedright\arraybackslash}p{(\linewidth - 6\tabcolsep) * \real{0.2500}}
  >{\raggedright\arraybackslash}p{(\linewidth - 6\tabcolsep) * \real{0.2500}}
  >{\raggedright\arraybackslash}p{(\linewidth - 6\tabcolsep) * \real{0.2500}}@{}}
\toprule\noalign{}
\begin{minipage}[b]{\linewidth}\raggedright
No.
\end{minipage} & \begin{minipage}[b]{\linewidth}\raggedright
Name
\end{minipage} & \begin{minipage}[b]{\linewidth}\raggedright
Description
\end{minipage} & \begin{minipage}[b]{\linewidth}\raggedright
Default
\end{minipage} \\
\midrule\noalign{}
\endhead
\bottomrule\noalign{}
\endlastfoot
Field & Metric & Laboratory & \\
1 & WELNAME & A character string of up to eight characters in length
that defines the well name for which the well hydraulic head reference
depth data is being defined. & None \\
2 & HYDREF & A real value that defines the hydraulic head reference
depth for reporting the hydraulic head pressure for the well. HYDREF
cannot be defaulted on the keyword; however if a well has not been set
by this keyword HYDREF is set equal to the value on the HYDRHEAD
keyword. & None \\
feet & m & cm & \\
Notes: & & & \\
\end{longtable}

\emph{Table 12.97: WHEDREFD Keyword Description}

See also the HYDRHEAD keyword in the PROPS section.

\subsubsection{Example}\label{example-580}

The following example defines three wells hydraulic head reference
depths for reporting, using the WHEDREFD keyword

\begin{verbatim}
--
--       WELL HYDRAULIC HEAD REFERENCE DEPTH
--
-- WELL  HYDREF
-- NAME  DEPTH
WHEDREFD
OP01     150.0                                                                 /
OP02     175.0                                                                 /
OP03     150.0                                                                 /
/
\end{verbatim}

Here, well OP01 and OP03 have their hydraulic head reference depths set
to 150.0 ft and well OP02's hydraulic head reference depth is set to
175.0 ft.

\subsection{WHISTCTL -- Define Well Historical Target
Phase}\label{whistctl-define-well-historical-target-phase}

\begin{longtable}[]{@{}
  >{\raggedright\arraybackslash}p{(\linewidth - 14\tabcolsep) * \real{0.1250}}
  >{\raggedright\arraybackslash}p{(\linewidth - 14\tabcolsep) * \real{0.1250}}
  >{\raggedright\arraybackslash}p{(\linewidth - 14\tabcolsep) * \real{0.1250}}
  >{\raggedright\arraybackslash}p{(\linewidth - 14\tabcolsep) * \real{0.1250}}
  >{\raggedright\arraybackslash}p{(\linewidth - 14\tabcolsep) * \real{0.1250}}
  >{\raggedright\arraybackslash}p{(\linewidth - 14\tabcolsep) * \real{0.1250}}
  >{\raggedright\arraybackslash}p{(\linewidth - 14\tabcolsep) * \real{0.1250}}
  >{\raggedright\arraybackslash}p{(\linewidth - 14\tabcolsep) * \real{0.1250}}@{}}
\toprule\noalign{}
\begin{minipage}[b]{\linewidth}\raggedright
RUNSPEC
\end{minipage} & \begin{minipage}[b]{\linewidth}\raggedright
GRID
\end{minipage} & \begin{minipage}[b]{\linewidth}\raggedright
EDIT
\end{minipage} & \begin{minipage}[b]{\linewidth}\raggedright
PROPS
\end{minipage} & \begin{minipage}[b]{\linewidth}\raggedright
REGIONS
\end{minipage} & \begin{minipage}[b]{\linewidth}\raggedright
SOLUTION
\end{minipage} & \begin{minipage}[b]{\linewidth}\raggedright
SUMMARY
\end{minipage} & \begin{minipage}[b]{\linewidth}\raggedright
SCHEDULE
\end{minipage} \\
\midrule\noalign{}
\endhead
\bottomrule\noalign{}
\endlastfoot
\end{longtable}

\subsubsection{Description}\label{description-1084}

The WHISTCTL keyword changes the target control for wells declared as
history match wells via the WCONHIST keyword in the SCHEDULE section.
The target phase is set on the WCONHIST keyword and WHISTCTL overrides
this value for all subsequent entries on the WCONHIST keyword.

\begin{longtable}[]{@{}
  >{\raggedright\arraybackslash}p{(\linewidth - 6\tabcolsep) * \real{0.2500}}
  >{\raggedright\arraybackslash}p{(\linewidth - 6\tabcolsep) * \real{0.2500}}
  >{\raggedright\arraybackslash}p{(\linewidth - 6\tabcolsep) * \real{0.2500}}
  >{\raggedright\arraybackslash}p{(\linewidth - 6\tabcolsep) * \real{0.2500}}@{}}
\toprule\noalign{}
\begin{minipage}[b]{\linewidth}\raggedright
No.
\end{minipage} & \begin{minipage}[b]{\linewidth}\raggedright
Name
\end{minipage} & \begin{minipage}[b]{\linewidth}\raggedright
Description
\end{minipage} & \begin{minipage}[b]{\linewidth}\raggedright
Default
\end{minipage} \\
\midrule\noalign{}
\endhead
\bottomrule\noalign{}
\endlastfoot
Field & Metric & Laboratory & \\
1 & TARGET & A defined character string that sets the observed target
production phase for the well, all the other phases are calculated
unconstrained and used for reporting only. The simulator will attempt to
meet the TARGET based on the phase rate stated in items (4) to (6) and
(10) on the WCONHIST keyword. TARGET should be set to one of the
following character strings: & None \\
2 & END & A defined character string that defines if the simulation
should terminate if the well has switch to BHP control by the simulator,
and should be set to one of the following character strings: Wells set
to BHP control via the WCONHIST or WHISCTL keywords are ignored. Only
END equal to NO is currently supported in OPM Flow. & NO \\
Notes: & & & \\
\end{longtable}

\emph{Table 12.98: WHISTCTL Keyword Description}

History matching wells are handled differently then ordinary wells that
use the WCONPROD keyword for controlling their production targets and
constraints. However, the wells still need to be defined like ordinary
production wells using the WELSPECS keyword in the SCHEDULE section.
History matching wells are converted to ordinary wells by restating a
well's control mode using either the WCONPROD or WELTARG keywords in the
SCHEDULE section.

See also the WCONHIST and WCONINJH keywords that are used to define the
historical production and injection data, All the aforementioned
keywords are described in the SCHEDULE section.

\subsubsection{Example}\label{example-581}

The example below shows the observed gas rates for the OP01 oil producer
for the first quarter of 2000.

\begin{verbatim}
-- ------------------------------------------------------------------------------
-- 01 JAN 2000 START OF SCHEDULE SECTION
-- ------------------------------------------------------------------------------
--
--       DEFINE WELL HISTORICAL TARGET PHASE
--
--       CNTL    BHP
--       MODE    STOP
WHISTCTL
         RESV    NO                                                   /
--
--       WELL HISTORICAL PRODUCTION CONTROLS
--
-- WELL  OPEN/  CNTL   OIL    WAT    GAS    VFP    VFP   THP   BHP
-- NAME  SHUT   MODE   RATE   RATE   RATE   TABLE  ALFQ  PRES  PRES
WCONHIST                                                                                                                          OP01     OPEN   ORAT  15.5E3  100.0  1550   10      1*   900.0 1*     /
/                                                                               DATES
01 FEB 2000 /
/
--
--       WELL HISTORICAL PRODUCTION CONTROLS
--
-- WELL  OPEN/  CNTL   OIL    WAT    GAS    VFP    VFP   THP   BHP
-- NAME  SHUT   MODE   RATE   RATE   RATE   TABLE  ALFQ  PRES  PRES
WCONHIST                                                                                                                          OP01     OPEN   ORAT  15.2E3  150.0  1520   1*      1*   875.0 3250.0 /
/                                                                               DATES
01 MAR 2000 /
/
--
--       WELL HISTORICAL PRODUCTION CONTROLS
--
-- WELL  OPEN/  CNTL   OIL    WAT    GAS    VFP    VFP   THP   BHP
-- NAME  SHUT   MODE   RATE   RATE   RATE   TABLE  ALFQ  PRES  PRES
WCONHIST                                                                                                                          OP01     OPEN   ORAT  15.0E3  200.0  1500   1*      1*   850.0 1*     /
/
\end{verbatim}

From January 1, 2000 the WCONHIST keyword defines well OP01, which is
open and is on oil rate control, to produce 15,500 stb/d oil, with the
observed rates of 100 stb/d of water and 15.5 MMscf/d of gas. However
the WHISCTL keyword resets the target control to reservoir voidage from
January 1, 2000 and onward. This is useful in initial history matching
runs to get a ``reasonable'' pressure match, by ensuring that the total
reservoir withdrawals are correct, although the individual phase
withdrawals will not match. Once a reasonable pressure match is achieved
for the reservoir then one can reset TARGET to the sales phase, OIL or
GAS, and continue with the matching of all the phases.

\subsection{WHTEMP -- Define Well Tubing Head Temperature
Parameters}\label{whtemp-define-well-tubing-head-temperature-parameters}

\begin{longtable}[]{@{}
  >{\raggedright\arraybackslash}p{(\linewidth - 14\tabcolsep) * \real{0.1250}}
  >{\raggedright\arraybackslash}p{(\linewidth - 14\tabcolsep) * \real{0.1250}}
  >{\raggedright\arraybackslash}p{(\linewidth - 14\tabcolsep) * \real{0.1250}}
  >{\raggedright\arraybackslash}p{(\linewidth - 14\tabcolsep) * \real{0.1250}}
  >{\raggedright\arraybackslash}p{(\linewidth - 14\tabcolsep) * \real{0.1250}}
  >{\raggedright\arraybackslash}p{(\linewidth - 14\tabcolsep) * \real{0.1250}}
  >{\raggedright\arraybackslash}p{(\linewidth - 14\tabcolsep) * \real{0.1250}}
  >{\raggedright\arraybackslash}p{(\linewidth - 14\tabcolsep) * \real{0.1250}}@{}}
\toprule\noalign{}
\begin{minipage}[b]{\linewidth}\raggedright
RUNSPEC
\end{minipage} & \begin{minipage}[b]{\linewidth}\raggedright
GRID
\end{minipage} & \begin{minipage}[b]{\linewidth}\raggedright
EDIT
\end{minipage} & \begin{minipage}[b]{\linewidth}\raggedright
PROPS
\end{minipage} & \begin{minipage}[b]{\linewidth}\raggedright
REGIONS
\end{minipage} & \begin{minipage}[b]{\linewidth}\raggedright
SOLUTION
\end{minipage} & \begin{minipage}[b]{\linewidth}\raggedright
SUMMARY
\end{minipage} & \begin{minipage}[b]{\linewidth}\raggedright
SCHEDULE
\end{minipage} \\
\midrule\noalign{}
\endhead
\bottomrule\noalign{}
\endlastfoot
\end{longtable}

\subsubsection{Description}\label{description-1085}

This keyword, WHTEMP, sets the parameters for the Tubing Head
Temperature calculation, which can either be a constant value, or from a
table lookup using a VFPPROD table, via the VFPPROD keyword in the
SCHEDULE section, containing tubing head temperature data.

This keyword can only be used if the thermal option has been activated
by the THERMAL keyword in the RUNSPEC section. Note this is different to
the commercial simulator that uses the TEMP keyword in the RUNSPEC
section to activate the black-oil thermal model.

This keyword is not supported by OPM Flow but would change the results
if supported so the simulation will be stopped.

\begin{longtable}[]{@{}
  >{\raggedright\arraybackslash}p{(\linewidth - 6\tabcolsep) * \real{0.2500}}
  >{\raggedright\arraybackslash}p{(\linewidth - 6\tabcolsep) * \real{0.2500}}
  >{\raggedright\arraybackslash}p{(\linewidth - 6\tabcolsep) * \real{0.2500}}
  >{\raggedright\arraybackslash}p{(\linewidth - 6\tabcolsep) * \real{0.2500}}@{}}
\toprule\noalign{}
\begin{minipage}[b]{\linewidth}\raggedright
No.
\end{minipage} & \begin{minipage}[b]{\linewidth}\raggedright
Name
\end{minipage} & \begin{minipage}[b]{\linewidth}\raggedright
Description
\end{minipage} & \begin{minipage}[b]{\linewidth}\raggedright
Default
\end{minipage} \\
\midrule\noalign{}
\endhead
\bottomrule\noalign{}
\endlastfoot
Field & Metric & Laboratory & \\
1 & WELNAME & A character string of up to eight characters in length
that defines the well name for which the well production targets and
constraints data are being defined. Note that the well name (WELNAME)
must have been declared previously using the WELSPECS keyword in the
SCHEDULE section, otherwise an error may occur. & None \\
2 & VFPTAB & A positive integer greater than or equal to zero that
references the production vertical lift performance table (VFPPROD),
containing the tubing head temperature data for the well. Note, a well
must have both a VFPPROD pressure and a VFPPROD temperature table, if
the tubing head temperatures are to be calculated. Alternatively, if a
constant tubing head temperature for a production well is to be defined
via the TEMP parameter, then VFPTAB should be defaulted with 1* instead.
& None \\
dimensionless & dimensionless & dimensionless & \\
3 & TEMP & A real positive value greater than zero that defines a
constant tubing head temperature for a production well. In this case the
VFPTAB parameter should be defaulted with 1* if a constant tubing head
temperature for a production well. & None \\
oF & oC & oC & \\
Notes: & & & \\
\end{longtable}

\emph{Table 12.99: WHTEMP Keyword Description}

See also the VFPPROD keyword in the SCHEDULE section.

\subsubsection{Example}\label{example-582}

The following example defines three wells tubing head temperature
parameters using the WHTEMP keyword

\begin{verbatim}
--
--       DEFINE WELL TUBING HEAD TEMPERATURE PARAMETERS
--
-- WELL  VFP    TUB
-- NAME  TABLE  TEMP
WHTEMP
OP01     5                                                                     /
OP02     1*     150                                                            /
OP03     5                                                                     /
/
\end{verbatim}

Here, well OP01 and OP03 used VFPPROD table number five to calculate the
tubing head temperature, and well OP02's uses a constant 150o tubing
head temperature.

\subsection{WINJCLN -- Clean a Fraction of a Deposited Filter
Cake}\label{winjcln-clean-a-fraction-of-a-deposited-filter-cake}

\begin{longtable}[]{@{}
  >{\raggedright\arraybackslash}p{(\linewidth - 14\tabcolsep) * \real{0.1250}}
  >{\raggedright\arraybackslash}p{(\linewidth - 14\tabcolsep) * \real{0.1250}}
  >{\raggedright\arraybackslash}p{(\linewidth - 14\tabcolsep) * \real{0.1250}}
  >{\raggedright\arraybackslash}p{(\linewidth - 14\tabcolsep) * \real{0.1250}}
  >{\raggedright\arraybackslash}p{(\linewidth - 14\tabcolsep) * \real{0.1250}}
  >{\raggedright\arraybackslash}p{(\linewidth - 14\tabcolsep) * \real{0.1250}}
  >{\raggedright\arraybackslash}p{(\linewidth - 14\tabcolsep) * \real{0.1250}}
  >{\raggedright\arraybackslash}p{(\linewidth - 14\tabcolsep) * \real{0.1250}}@{}}
\toprule\noalign{}
\begin{minipage}[b]{\linewidth}\raggedright
RUNSPEC
\end{minipage} & \begin{minipage}[b]{\linewidth}\raggedright
GRID
\end{minipage} & \begin{minipage}[b]{\linewidth}\raggedright
EDIT
\end{minipage} & \begin{minipage}[b]{\linewidth}\raggedright
PROPS
\end{minipage} & \begin{minipage}[b]{\linewidth}\raggedright
REGIONS
\end{minipage} & \begin{minipage}[b]{\linewidth}\raggedright
SOLUTION
\end{minipage} & \begin{minipage}[b]{\linewidth}\raggedright
SUMMARY
\end{minipage} & \begin{minipage}[b]{\linewidth}\raggedright
SCHEDULE
\end{minipage} \\
\midrule\noalign{}
\endhead
\bottomrule\noalign{}
\endlastfoot
\end{longtable}

\subsubsection{Description}\label{description-1086}

The \hyperref[REF_HEADING_KEYWORD_WINJCLN]{WINJCLN} keyword signals that
a filter cake should be completely or partially cleaned -- this
effectively multiplies the accumulated filter cake skin.

\begin{longtable}[]{@{}
  >{\raggedright\arraybackslash}p{(\linewidth - 6\tabcolsep) * \real{0.2500}}
  >{\raggedright\arraybackslash}p{(\linewidth - 6\tabcolsep) * \real{0.2500}}
  >{\raggedright\arraybackslash}p{(\linewidth - 6\tabcolsep) * \real{0.2500}}
  >{\raggedright\arraybackslash}p{(\linewidth - 6\tabcolsep) * \real{0.2500}}@{}}
\toprule\noalign{}
\begin{minipage}[b]{\linewidth}\raggedright
No.
\end{minipage} & \begin{minipage}[b]{\linewidth}\raggedright
Name
\end{minipage} & \begin{minipage}[b]{\linewidth}\raggedright
Description
\end{minipage} & \begin{minipage}[b]{\linewidth}\raggedright
Default
\end{minipage} \\
\midrule\noalign{}
\endhead
\bottomrule\noalign{}
\endlastfoot
Field & Metric & Laboratory & \\
1 & WELNAME & A character string of up to eight characters in length
that defines the well name for which the filter cake properties are
being defined. Note that the well name (WELNAME) must have been declared
previously using the WELSPECS keyword in the SCHEDULE section, otherwise
an error may occur. & None \\
2 & FCLNFRAC & A real positive value between 0 and 1 that defines the
fraction of filter cake permeability (skin factor) to be removed. The
accumulated filter cake skin factor for matching connections will be
multiplied by (1 -- FCLNFRAC), so the default value of 1 will completely
clean the filter cake. & 1 \\
dimensionless & dimensionless & dimensionless & \\
3 & I & An integer that defines the matching connection location in the
I-direction. If set to \textless{} 1 then all connections in the
I-direction that also satisfy J and K criteria are selected. & -1 \\
4 & J & An integer that defines the matching connection location in the
J-direction. If set to \textless{} 1 then all connections in the
J-direction that also satisfy I and K criteria are selected. & -1 \\
5 & K & An integer that defines the matching connection location in the
K-direction. If set to \textless{} 1 then all connections in the
K-direction that also satisfy I and J criteria are selected. & -1 \\
Notes: & & & \\
\end{longtable}

\emph{Table 12.3.295.1: WINJCLN Keyword Description}

See also the WINJDAN keyword to define the filter cake properties and
the \hyperref[REF_HEADING_KEYWORD_WINJFCNC]{WINJFCNC} keyword to define
a well's injected filtrate concentration. All the aforementioned
keywords are described in the SCHEDULE section.

\subsubsection{Example}\label{example-583}

The following example signals the filter cake clean up for water
injection wells using \hyperref[REF_HEADING_KEYWORD_WINJCLN]{WINJCLN}:

\begin{verbatim}
--
--       WELL FILTER CAKE CLEAN UP
--
-- WELL  CLEAN  --LOCATION--
-- NAME  FRAC    II  JJ  KK
WINJDAM
INJ-A1   0.4      0   0   3  /
INJ-A1   0.4      0   0   4  /
INJ-B*   0.9    /
/
\end{verbatim}

In well INJ-A1 forty percent of the filter cake is cleaned up in well
connections in layers 3 and 4 (filter cake skin is multiplied by 0.6).
In wells matching INJ-B* ninety percent of the filter cake is cleaned up
in all well connections (filter cake skin is multiplied by 0.1).

\subsection{WINJDAM -- Define Well Filter Cake
Properties}\label{winjdam-define-well-filter-cake-properties}

\begin{longtable}[]{@{}
  >{\raggedright\arraybackslash}p{(\linewidth - 14\tabcolsep) * \real{0.1250}}
  >{\raggedright\arraybackslash}p{(\linewidth - 14\tabcolsep) * \real{0.1250}}
  >{\raggedright\arraybackslash}p{(\linewidth - 14\tabcolsep) * \real{0.1250}}
  >{\raggedright\arraybackslash}p{(\linewidth - 14\tabcolsep) * \real{0.1250}}
  >{\raggedright\arraybackslash}p{(\linewidth - 14\tabcolsep) * \real{0.1250}}
  >{\raggedright\arraybackslash}p{(\linewidth - 14\tabcolsep) * \real{0.1250}}
  >{\raggedright\arraybackslash}p{(\linewidth - 14\tabcolsep) * \real{0.1250}}
  >{\raggedright\arraybackslash}p{(\linewidth - 14\tabcolsep) * \real{0.1250}}@{}}
\toprule\noalign{}
\begin{minipage}[b]{\linewidth}\raggedright
RUNSPEC
\end{minipage} & \begin{minipage}[b]{\linewidth}\raggedright
GRID
\end{minipage} & \begin{minipage}[b]{\linewidth}\raggedright
EDIT
\end{minipage} & \begin{minipage}[b]{\linewidth}\raggedright
PROPS
\end{minipage} & \begin{minipage}[b]{\linewidth}\raggedright
REGIONS
\end{minipage} & \begin{minipage}[b]{\linewidth}\raggedright
SOLUTION
\end{minipage} & \begin{minipage}[b]{\linewidth}\raggedright
SUMMARY
\end{minipage} & \begin{minipage}[b]{\linewidth}\raggedright
SCHEDULE
\end{minipage} \\
\midrule\noalign{}
\endhead
\bottomrule\noalign{}
\endlastfoot
\end{longtable}

\subsubsection{Description}\label{description-1087}

The \hyperref[REF_HEADING_KEYWORD_WINJDAM]{WINJDAM} keyword defines
filter cake properties for wells that have previously been defined by
the WELSPECS keyword in the SCHEDULE section. This keyword must be
accompanied by the \hyperref[REF_HEADING_KEYWORD_WINJFCNC]{WINJFCNC}
keyword to define the injected filtrate concentration.

\begin{longtable}[]{@{}
  >{\raggedright\arraybackslash}p{(\linewidth - 6\tabcolsep) * \real{0.2500}}
  >{\raggedright\arraybackslash}p{(\linewidth - 6\tabcolsep) * \real{0.2500}}
  >{\raggedright\arraybackslash}p{(\linewidth - 6\tabcolsep) * \real{0.2500}}
  >{\raggedright\arraybackslash}p{(\linewidth - 6\tabcolsep) * \real{0.2500}}@{}}
\toprule\noalign{}
\begin{minipage}[b]{\linewidth}\raggedright
No.
\end{minipage} & \begin{minipage}[b]{\linewidth}\raggedright
Name
\end{minipage} & \begin{minipage}[b]{\linewidth}\raggedright
Description
\end{minipage} & \begin{minipage}[b]{\linewidth}\raggedright
Default
\end{minipage} \\
\midrule\noalign{}
\endhead
\bottomrule\noalign{}
\endlastfoot
Field & Metric & Laboratory & \\
1 & WELNAME & A character string of up to eight characters in length
that defines the well name for which the filter cake properties are
being defined. Note that the well name (WELNAME) must have been declared
previously using the WELSPECS keyword in the SCHEDULE section, otherwise
an error may occur. & None \\
2 & GEOMETRY & A defined character string that the filter cake geometry
of the well. GEOMETRY must be set to one of the following character
strings: & None \\
3 & FCPERM & A real positive value that defines the filter cake
permeability. & None \\
mD & mD & mD & \\
4 & FCPORO & A real positive value that defines the filter cake
porosity. & 0.3 \\
dimensionless & dimensionless & dimensionless & \\
5 & FCRADIUS & Well radius to use in skin factor calculations. If
FCRADIUS is defaulted then half the diameter from the COMPDAT keyword
item 9 will be used if it is defined otherwise 0.5 feet will be used. &
0.5 feet \\
feet & m & cm & \\
6 & FCAOF & A real positive value that defines the flow area for each
connection. If FCAOF is defaulted then the value will be evaluated as
2πrwh, where rw is equal to FCRADIUS, and h is evaluated using the
permeability thickness kh value from the COMPDAT keyword item 10 and the
permeability FCPERM item 3 above. & 2πrwh \\
ft2 & m2 & cm2 & \\
7 & I & An integer that defines the matching connection location in the
I-direction. If set to \textless{} 1 then all connections in the
I-direction that also satisfy J and K criteria are selected. & -1 \\
8 & J & An integer that defines the matching connection location in the
J-direction. If set to \textless{} 1 then all connections in the
J-direction that also satisfy I and K criteria are selected. & -1 \\
9 & K & An integer that defines the matching connection location in the
K-direction. If set to \textless{} 1 then all connections in the
K-direction that also satisfy I and J criteria are selected. & -1 \\
Notes: & & & \\
\end{longtable}

\emph{Table 12.3.295.1: WINJDAM Keyword Description}

See also the \hyperref[REF_HEADING_KEYWORD_WINJFCNC]{WINJFCNC} keyword
to define a well's injected filtrate concentration and the
\hyperref[REF_HEADING_KEYWORD_WINJCLN]{WINJCLN} keyword to signal that a
filter cake should be completely or partially cleaned. All the
aforementioned keywords are described in the SCHEDULE section.

\subsubsection{Example}\label{example-584}

The following example defines the filter cake properties for a water
injection well using \hyperref[REF_HEADING_KEYWORD_WINJDAM]{WINJDAM}:

\begin{verbatim}
--
--       WELL FILTER CAKE PROPERTIES
--
-- WELL  LINEAR/  PERM   PORO   WELL   FLOW   --LOCATION--
-- NAME  RADIAL                 RAD    AREA    II  JJ  KK
WINJDAM
INJ      LINEAR   10     0.32   /
INJ      RADIAL   1      0.25   1*     1*       0   0   3 /
/
\end{verbatim}

Well INJ initially has a linear filter cake model for all completions
with permeability of 10 mD, porosity of 0.32, and default well radius
and flow area based on data from the COMPDAT keyword. This is then
overwritten for completions in layer 3 with a radial filter cake model
with permeability of 1 mD and porosity of 0.25 (again with default well
radius and flow area).

\subsection{WINJFCNC -- Define Injection Well Filtrate
Concentration}\label{winjfcnc-define-injection-well-filtrate-concentration}

\begin{longtable}[]{@{}
  >{\raggedright\arraybackslash}p{(\linewidth - 14\tabcolsep) * \real{0.1250}}
  >{\raggedright\arraybackslash}p{(\linewidth - 14\tabcolsep) * \real{0.1250}}
  >{\raggedright\arraybackslash}p{(\linewidth - 14\tabcolsep) * \real{0.1250}}
  >{\raggedright\arraybackslash}p{(\linewidth - 14\tabcolsep) * \real{0.1250}}
  >{\raggedright\arraybackslash}p{(\linewidth - 14\tabcolsep) * \real{0.1250}}
  >{\raggedright\arraybackslash}p{(\linewidth - 14\tabcolsep) * \real{0.1250}}
  >{\raggedright\arraybackslash}p{(\linewidth - 14\tabcolsep) * \real{0.1250}}
  >{\raggedright\arraybackslash}p{(\linewidth - 14\tabcolsep) * \real{0.1250}}@{}}
\toprule\noalign{}
\begin{minipage}[b]{\linewidth}\raggedright
RUNSPEC
\end{minipage} & \begin{minipage}[b]{\linewidth}\raggedright
GRID
\end{minipage} & \begin{minipage}[b]{\linewidth}\raggedright
EDIT
\end{minipage} & \begin{minipage}[b]{\linewidth}\raggedright
PROPS
\end{minipage} & \begin{minipage}[b]{\linewidth}\raggedright
REGIONS
\end{minipage} & \begin{minipage}[b]{\linewidth}\raggedright
SOLUTION
\end{minipage} & \begin{minipage}[b]{\linewidth}\raggedright
SUMMARY
\end{minipage} & \begin{minipage}[b]{\linewidth}\raggedright
SCHEDULE
\end{minipage} \\
\midrule\noalign{}
\endhead
\bottomrule\noalign{}
\endlastfoot
\end{longtable}

\subsubsection{Description}\label{description-1088}

The \hyperref[REF_HEADING_KEYWORD_WINJFCNC]{WINJFCNC} keyword defines
the injected filtrate concentration for wells that have previously been
defined by the WELSPECS keyword in the SCHEDULE section.

\begin{longtable}[]{@{}
  >{\raggedright\arraybackslash}p{(\linewidth - 6\tabcolsep) * \real{0.2500}}
  >{\raggedright\arraybackslash}p{(\linewidth - 6\tabcolsep) * \real{0.2500}}
  >{\raggedright\arraybackslash}p{(\linewidth - 6\tabcolsep) * \real{0.2500}}
  >{\raggedright\arraybackslash}p{(\linewidth - 6\tabcolsep) * \real{0.2500}}@{}}
\toprule\noalign{}
\begin{minipage}[b]{\linewidth}\raggedright
No.
\end{minipage} & \begin{minipage}[b]{\linewidth}\raggedright
Name
\end{minipage} & \begin{minipage}[b]{\linewidth}\raggedright
Description
\end{minipage} & \begin{minipage}[b]{\linewidth}\raggedright
Default
\end{minipage} \\
\midrule\noalign{}
\endhead
\bottomrule\noalign{}
\endlastfoot
Field & Metric & Laboratory & \\
1 & WELNAME & A character string of up to eight characters in length
that defines the well name for which the filter cake properties are
being defined. Note that the well name (WELNAME) must have been declared
previously using the WELSPECS keyword in the SCHEDULE section, otherwise
an error may occur. & None \\
2 & FCONCPPM & A real positive value that defines the volumetric
concentration of filtrate in the injected water. This value may be
specified using a User Defined Argument (UDA). & 0 \\
ppm & ppm & ppm & \\
Notes: & & & \\
\end{longtable}

\emph{Table 12.3.296.1: WINJFCNC Keyword Description}

See also the \hyperref[REF_HEADING_KEYWORD_WINJDAM]{WINJDAM} keyword to
define the filter cake properties and the
\hyperref[REF_HEADING_KEYWORD_WINJCLN]{WINJCLN} keyword to signal that a
filter cake should be completely or partially cleaned. All the
aforementioned keywords are described in the SCHEDULE section.

\subsubsection{Example}\label{example-585}

The following example defines the filtrate injection concentration for a
water injection well using
\hyperref[REF_HEADING_KEYWORD_WINJFCNC]{WINJFCNC}:

\begin{verbatim}
--
--       WELL FILTRATE INJECTION CONCENTRATION
--
-- WELL  FILTRATE
-- NAME  CONC
WINJFCNC
INJ-A1   10   /
INJ-B*   30   /
/
\end{verbatim}

Well INJ-A1 has a filtrate concentration of 10 ppm in the injection
water, and wells matching INJ-B* have a concentration of 30 ppm.

\subsection{WINJGAS -- Define Gas Injection Properties for a
Well}\label{winjgas-define-gas-injection-properties-for-a-well}

\begin{longtable}[]{@{}
  >{\raggedright\arraybackslash}p{(\linewidth - 14\tabcolsep) * \real{0.1250}}
  >{\raggedright\arraybackslash}p{(\linewidth - 14\tabcolsep) * \real{0.1250}}
  >{\raggedright\arraybackslash}p{(\linewidth - 14\tabcolsep) * \real{0.1250}}
  >{\raggedright\arraybackslash}p{(\linewidth - 14\tabcolsep) * \real{0.1250}}
  >{\raggedright\arraybackslash}p{(\linewidth - 14\tabcolsep) * \real{0.1250}}
  >{\raggedright\arraybackslash}p{(\linewidth - 14\tabcolsep) * \real{0.1250}}
  >{\raggedright\arraybackslash}p{(\linewidth - 14\tabcolsep) * \real{0.1250}}
  >{\raggedright\arraybackslash}p{(\linewidth - 14\tabcolsep) * \real{0.1250}}@{}}
\toprule\noalign{}
\begin{minipage}[b]{\linewidth}\raggedright
RUNSPEC
\end{minipage} & \begin{minipage}[b]{\linewidth}\raggedright
GRID
\end{minipage} & \begin{minipage}[b]{\linewidth}\raggedright
EDIT
\end{minipage} & \begin{minipage}[b]{\linewidth}\raggedright
PROPS
\end{minipage} & \begin{minipage}[b]{\linewidth}\raggedright
REGIONS
\end{minipage} & \begin{minipage}[b]{\linewidth}\raggedright
SOLUTION
\end{minipage} & \begin{minipage}[b]{\linewidth}\raggedright
SUMMARY
\end{minipage} & \begin{minipage}[b]{\linewidth}\raggedright
SCHEDULE
\end{minipage} \\
\midrule\noalign{}
\endhead
\bottomrule\noalign{}
\endlastfoot
\end{longtable}

\subsubsection{Description}\label{description-1089}

The WINJGAS keyword defines the properties of the injection gas stream,
for a given well. Once a gas well stream has been defined via the
WELLSTRE keyword in the RUNSPEC section, it can be used with either the
WINJGAS or GINJGAS keywords, to set the injected gas composition.
Similarly, if an oil well stream has been defined by WELLSTRE, then the
well stream can be used with the WINJOIL keyword in the SCHEDULE
section, to specify the injected oil composition. Note that, it is
unnecessary to use WINJGAS for wells subordinate to a group having gas
injection control, with the gas properties set by GINJGAS keyword in the
RUNSPEC section. In this case the injection stream is defined by the
GINJGAS keyword. However, if a gas injection well under group control
users the WINJGAS keyword, then this fluid, and not the group's fluid
will be injected instead, at a rate controlled by the group.

The keyword should only be used if the CO2STORE and GASWAT keywords in
the RUNSPEC section have also be activated for the gas-water two
component model.

\begin{tcolorbox}[enhanced jigsaw, title=\textcolor{quarto-callout-note-color}{\faInfo}\hspace{0.5em}{Note}, coltitle=black, breakable, toprule=.15mm, colbacktitle=quarto-callout-note-color!10!white, leftrule=.75mm, opacitybacktitle=0.6, colback=white, bottomtitle=1mm, toptitle=1mm, titlerule=0mm, left=2mm, colframe=quarto-callout-note-color-frame, arc=.35mm, rightrule=.15mm, bottomrule=.15mm, opacityback=0]

This is an OPM Flow keyword used with OPM Flow's CO2STORE and GASWAT
keywords in the RUNSPEC section, and should not be confused with the
more general version of the WINJGAS keyword used in the commercial
compositional simulator. Secondly, although OPM Flow parses the keyword,
the simulator currently ignores the data for this keyword.

\end{tcolorbox}

\begin{longtable}[]{@{}
  >{\raggedright\arraybackslash}p{(\linewidth - 6\tabcolsep) * \real{0.2500}}
  >{\raggedright\arraybackslash}p{(\linewidth - 6\tabcolsep) * \real{0.2500}}
  >{\raggedright\arraybackslash}p{(\linewidth - 6\tabcolsep) * \real{0.2500}}
  >{\raggedright\arraybackslash}p{(\linewidth - 6\tabcolsep) * \real{0.2500}}@{}}
\toprule\noalign{}
\begin{minipage}[b]{\linewidth}\raggedright
No.
\end{minipage} & \begin{minipage}[b]{\linewidth}\raggedright
Name
\end{minipage} & \begin{minipage}[b]{\linewidth}\raggedright
Description
\end{minipage} & \begin{minipage}[b]{\linewidth}\raggedright
Default
\end{minipage} \\
\midrule\noalign{}
\endhead
\bottomrule\noalign{}
\endlastfoot
1 & WELNAME & A character string of up to eight characters in length
that defines the injection well name, for which the gas injection
properties are being specified. Note that the well name (WELNAME) must
have been declared previously using the WELSPECS keyword in the SCHEDULE
section, otherwise an error may occur. & None \\
2 & STREAM & A defined character string that determines the properties
of the injection gas, and as should be set to one of the following
values: Only the STREAM option is supported by OPM Flow. & GRUP \\
3 & \hyperref[REF_HEADING_KEYWORD_SOURCE_12_3]{SOURCE} & A character
string of up to eight characters in length, that defines the source of
the gas injection stream, based on the value of STREAM. If STREAM equals
GV, then source should be set to a group name. For STREAM equal to MIX
or STREAM, then \hyperref[REF_HEADING_KEYWORD_SOURCE_12_3]{SOURCE}
should be set to the name of the gas stream, as defined by the WINJMIX,
WINJORD, or WELLSTRE keywords. & None \\
4 & MAKEUP & The name of the well stream used for the make-up gas, if
make-up gas is required for WELNAME to match the injection target for
the well. This option is not supported by OPM Flow. & None \\
5 & STAGE & STAGE defines the separator stage from which the injection
gas should be taken from. In this case, the vapor phase from any stage
may be used, and the default value of zero users the total vapor phase
from the separator. This option is not supported by OPM Flow. & 0 \\
Notes: & & & \\
\end{longtable}

\emph{Table 12.100: WINJGAS Keyword Description}

\subsubsection{Example}\label{example-586}

The following example defines how to specify a two component
formulation, together with defining the names of the composition
components, to be used with the CO2STORE and GASWAT options.

\begin{verbatim}
-- ==============================================================================
--
-- PROPS SECTION
--
-- ==============================================================================
PROPS             --
--       CONFIRM NUMBER OF COMPOSITIONAL COMPONENTS (OPM FLOW KEYWORD)
--
NCOMPS
                2                                                                      /
--
--       DEFINE COMPOSITIONAL COMPONENTS NAMES (OPM FLOW KEYWORD)
--
CNAMES
         'CO2'
         'H2O'                                                                 /
\end{verbatim}

The second part of the example, defines the well stream for the above
two component CO2 water system.

\begin{verbatim}
-- ==============================================================================
--
-- SCHEDULE SECTION
--
-- ==============================================================================
SCHEDULE
--
--       WELL STREAM INJECTION COMPOSITION (OPM FLOW Keyword)
--
-- WELL     -- WELL STREAM COMPOSITIONAL COMPONENT       --
-- STREAM   --          MOLE FRACTIONS                   --
WELLSTRE
'C02STREAM'    1.000   0.000                               /
/
--
--       WELL GAS INJECTION PROPERTIES
--
-- WELL  STREAM   SOURCE      MAKEUP   SEP
-- NAME  OPTION   DEPTH       GAS      STAGE
WINJGAS
GI01     STREAM   C02STREAM   1*       1* /
/
\end{verbatim}

Here the well stream consists of 100\% CO2 and zero water, with well
GI01 using the gas injection properties as defined by the WELLSTRE
keyword and allocated via the WINJGAS keyword.

Finally, the gas injection rate is set via the WCONINJE keyword as shown
below.

\begin{verbatim}
--
--       WELL INJECTION CONTROLS
--
-- WELL  FLUID  OPEN/  CNTL  SURF   RESV   BHP   THP   VFP
-- NAME  TYPE   SHUT   MODE  RATE   RATE   PRES  PRES  TABLE
WCONINJE
GI01     GAS    OPEN   RATE  10E4   1*     300    1*    1*                     /
/
\end{verbatim}

Thus, gas injector GI01 will inject 10 x 104 m3 of CO2 per day, assuming
metric units.

\subsection{WINJMULT -- Define Well Pressure Dependent Injectivity
Multipliers}\label{winjmult-define-well-pressure-dependent-injectivity-multipliers}

\begin{longtable}[]{@{}
  >{\raggedright\arraybackslash}p{(\linewidth - 14\tabcolsep) * \real{0.1250}}
  >{\raggedright\arraybackslash}p{(\linewidth - 14\tabcolsep) * \real{0.1250}}
  >{\raggedright\arraybackslash}p{(\linewidth - 14\tabcolsep) * \real{0.1250}}
  >{\raggedright\arraybackslash}p{(\linewidth - 14\tabcolsep) * \real{0.1250}}
  >{\raggedright\arraybackslash}p{(\linewidth - 14\tabcolsep) * \real{0.1250}}
  >{\raggedright\arraybackslash}p{(\linewidth - 14\tabcolsep) * \real{0.1250}}
  >{\raggedright\arraybackslash}p{(\linewidth - 14\tabcolsep) * \real{0.1250}}
  >{\raggedright\arraybackslash}p{(\linewidth - 14\tabcolsep) * \real{0.1250}}@{}}
\toprule\noalign{}
\begin{minipage}[b]{\linewidth}\raggedright
RUNSPEC
\end{minipage} & \begin{minipage}[b]{\linewidth}\raggedright
GRID
\end{minipage} & \begin{minipage}[b]{\linewidth}\raggedright
EDIT
\end{minipage} & \begin{minipage}[b]{\linewidth}\raggedright
PROPS
\end{minipage} & \begin{minipage}[b]{\linewidth}\raggedright
REGIONS
\end{minipage} & \begin{minipage}[b]{\linewidth}\raggedright
SOLUTION
\end{minipage} & \begin{minipage}[b]{\linewidth}\raggedright
SUMMARY
\end{minipage} & \begin{minipage}[b]{\linewidth}\raggedright
SCHEDULE
\end{minipage} \\
\midrule\noalign{}
\endhead
\bottomrule\noalign{}
\endlastfoot
\end{longtable}

\subsubsection{Description}\label{description-1090}

The WINJMULT keyword defines pressure dependent injectivity multipliers
for injection wells and can be used to approximate the increase or
decrease in a well's injectivity due to hydraulic fracturing in water
injection wells. Only injection wells are processed by this keyword,
even if production wells have been entered by the keyword.

\begin{longtable}[]{@{}
  >{\raggedright\arraybackslash}p{(\linewidth - 6\tabcolsep) * \real{0.2500}}
  >{\raggedright\arraybackslash}p{(\linewidth - 6\tabcolsep) * \real{0.2500}}
  >{\raggedright\arraybackslash}p{(\linewidth - 6\tabcolsep) * \real{0.2500}}
  >{\raggedright\arraybackslash}p{(\linewidth - 6\tabcolsep) * \real{0.2500}}@{}}
\toprule\noalign{}
\begin{minipage}[b]{\linewidth}\raggedright
No.
\end{minipage} & \begin{minipage}[b]{\linewidth}\raggedright
Name
\end{minipage} & \begin{minipage}[b]{\linewidth}\raggedright
Description
\end{minipage} & \begin{minipage}[b]{\linewidth}\raggedright
Default
\end{minipage} \\
\midrule\noalign{}
\endhead
\bottomrule\noalign{}
\endlastfoot
Field & Metric & Laboratory & \\
1 & WELNAME & A character string of up to eight characters in length
that defines the well name for which the well connection data are being
defined. Note that the well name (WELNAME) must have been declared
previously using the WELSPECS keyword in the SCHEDULE section, otherwise
an error may occur. & None \\
2 & FRACPRES & FRACPRES is the fracture opening pressure (Pfractue) used
in equation 12.3.296.1. & None \\
psia & barsa & atma & \\
3 & ALPHA & ALPHA is the multiplier gradient, α, in equation 12.3.296.1.
& Defined \\
1/psia 0.0 & 1/barsa 0.0 & 1/atma 0.0 & \\
4 & OPTION & A defined character string that determines how the data on
this keyword is applied, and should be set to one of the following
character strings: This option means that even if the pressures in the
wellbore, and therefore the sand face pressures, later declines,
injectivity remains unchanged for all the connections. The reversibility
of this option means that if the pressures in the wellbore, and
therefore the sand face pressures, later declines, injectivity will also
decline for all the connections. The reversibility of this option means
that if the pressures in the wellbore, later declines, injectivity will
also decline for all the connections. & WREV \\
5 & I & An integer value less than or equal to NX that defines the
connection location in the I-direction. If set to zero, a negative
value, or defaulted with 1* then all connections in the I-direction will
be multiplied by ALPHA, depending on the selected OPTION value. & 1* \\
6 & J & An integer value less than or equal to NY that defines the
connection location in the J-direction. If set to zero, a negative
value, or defaulted with 1* then all connections in the J-direction will
be multiplied by ALPHA, depending on the selected OPTION value. & 1* \\
7 & K & An integer value less than or equal to NZ that defines the
connection location in the K-direction. If set to zero, a negative
value, or defaulted with 1* then all connections in the K-direction will
be multiplied by ALPHA, depending on the selected OPTION value. & 1* \\
Notes: & & & \\
\end{longtable}

\emph{Table 12.3.296.1: WINJMULT Keyword Description}

The methodology for applying the well pressure dependent injectivity
multipliers is outlined in equation 12.3.296.1.

\begin{equation}\phantomsection\label{eq-12-3-296-1}{
\begin{matrix}\mathit{Multiplier}= & 1.0 + α \left({P}_{\mathit{WBHP}} - {P}_{\mathit{fracture}}\right) & \text{ for }{P}_{\mathit{WBHP}} > {P}_{\mathit{fracture}} \\ \mathit{Multiplier}= & 1.0 & \text{ for }{P}_{\mathit{WBHP}} < {P}_{\mathit{fracture}}\end{matrix}
}\end{equation}

where:

Multiplier = the resulting multiplier to be applied to the selected
connections.

α = the ALPHA multiplier gradient in Table 12.3.296.1.

PWBHP = either the well's current flowing bottom-hole pressure, or the
selected

individual connection's sand face pressure.

Pfracture = the effective fracture opening pressure, FRACPRES in Table
12.3.296.1.

The equation is applied every time there is a calculation to determine a
well's flow rate, this results in the calculated mobility rates being
scaled up by the Multiplier value in equation 12.3.296.1. Note also that
since the scaling is performed on the connection fluid mobility values,
then the reported connection transmissibilities in the print file etc.
remain unchanged.

\begin{tcolorbox}[enhanced jigsaw, title=\textcolor{quarto-callout-note-color}{\faInfo}\hspace{0.5em}{Note}, coltitle=black, breakable, toprule=.15mm, colbacktitle=quarto-callout-note-color!10!white, leftrule=.75mm, opacitybacktitle=0.6, colback=white, bottomtitle=1mm, toptitle=1mm, titlerule=0mm, left=2mm, colframe=quarto-callout-note-color-frame, arc=.35mm, rightrule=.15mm, bottomrule=.15mm, opacityback=0]

If all the connection parameters (I, J, K) are defaulted, or OPTION is
set equal to WREV, then the multiplier is applied to all connections in
the well. If any of the connection parameters (I, J, K) have positive
values and OPTION is set equal to CIRR or CREV, then the multiplier is
applied to the selected connections as determined by the (I, J, K)
parameters.

\end{tcolorbox}

\subsubsection{Example}\label{example-587}

The example below show the WINJMULT keyword for three water injection
wells.

\begin{verbatim}
--
--       DEFINE WELL CONNECTION MULTIPLIERS
--
-- WELL  FRAC  MULT    FRAC  --LOCATION--
-- NAME  PRES  VALUE   OPTN    I   J   K
WINJMULT
WI01     4200  0.0250  1*      1*  1*  1*        /
WI02     4250  0.0025  CIRR    1*  1* 145        /
WI02     4250  0.0025  CIRR    1*  1* 146        /
WI02     4250  0.0025  CIRR    1*  1* 147        /
WI03     4400  0.0055  CREV    1*  1* 160        /
WI03     4400  0.0055  CREV    1*  1* 165        /
/
\end{verbatim}

The first well, WI01, uses the default value for OPTION, that is WREV,
which means that the injectivity multiplier will be applied to all the
connections in the well and the well's bottom-hole pressure is used in
the calculation. In this case the process is reversible. The second
well, WI02, applies the injectivity multiplier to all connections in
layers 145 to 147 using the sand face pressures in the calculation, and
the process is irreversible. Finally for well WI03, the multiplier is
applied to all connections in layers 160 and 165 using the sand face
pressures in the calculation, and in this case the process is
reversible.

\subsection{WINJTEMP -- Define Injection Fluid Thermal
Properties}\label{winjtemp-define-injection-fluid-thermal-properties}

\begin{longtable}[]{@{}
  >{\raggedright\arraybackslash}p{(\linewidth - 14\tabcolsep) * \real{0.1250}}
  >{\raggedright\arraybackslash}p{(\linewidth - 14\tabcolsep) * \real{0.1250}}
  >{\raggedright\arraybackslash}p{(\linewidth - 14\tabcolsep) * \real{0.1250}}
  >{\raggedright\arraybackslash}p{(\linewidth - 14\tabcolsep) * \real{0.1250}}
  >{\raggedright\arraybackslash}p{(\linewidth - 14\tabcolsep) * \real{0.1250}}
  >{\raggedright\arraybackslash}p{(\linewidth - 14\tabcolsep) * \real{0.1250}}
  >{\raggedright\arraybackslash}p{(\linewidth - 14\tabcolsep) * \real{0.1250}}
  >{\raggedright\arraybackslash}p{(\linewidth - 14\tabcolsep) * \real{0.1250}}@{}}
\toprule\noalign{}
\begin{minipage}[b]{\linewidth}\raggedright
RUNSPEC
\end{minipage} & \begin{minipage}[b]{\linewidth}\raggedright
GRID
\end{minipage} & \begin{minipage}[b]{\linewidth}\raggedright
EDIT
\end{minipage} & \begin{minipage}[b]{\linewidth}\raggedright
PROPS
\end{minipage} & \begin{minipage}[b]{\linewidth}\raggedright
REGIONS
\end{minipage} & \begin{minipage}[b]{\linewidth}\raggedright
SOLUTION
\end{minipage} & \begin{minipage}[b]{\linewidth}\raggedright
SUMMARY
\end{minipage} & \begin{minipage}[b]{\linewidth}\raggedright
SCHEDULE
\end{minipage} \\
\midrule\noalign{}
\endhead
\bottomrule\noalign{}
\endlastfoot
\end{longtable}

\subsubsection{Description}\label{description-1091}

WINJTEMP defines the injection fluid thermal properties for when the
thermal option has been activated by the THERMAL keyword in the RUNSPEC
section. Only water and gas injection is supported.

This keyword can only be used if OPM Flow's thermal option has been
activated by the THERMAL keyword in the RUNSPEC section. Note this is
different to the commercial simulator that uses the TEMP keyword in the
RUNSPEC section to activate the black-oil thermal model.

\begin{longtable}[]{@{}
  >{\raggedright\arraybackslash}p{(\linewidth - 6\tabcolsep) * \real{0.2500}}
  >{\raggedright\arraybackslash}p{(\linewidth - 6\tabcolsep) * \real{0.2500}}
  >{\raggedright\arraybackslash}p{(\linewidth - 6\tabcolsep) * \real{0.2500}}
  >{\raggedright\arraybackslash}p{(\linewidth - 6\tabcolsep) * \real{0.2500}}@{}}
\toprule\noalign{}
\begin{minipage}[b]{\linewidth}\raggedright
No.
\end{minipage} & \begin{minipage}[b]{\linewidth}\raggedright
Name
\end{minipage} & \begin{minipage}[b]{\linewidth}\raggedright
Description
\end{minipage} & \begin{minipage}[b]{\linewidth}\raggedright
Default
\end{minipage} \\
\midrule\noalign{}
\endhead
\bottomrule\noalign{}
\endlastfoot
Field & Metric & Laboratory & \\
1 & WELNAME & A character string of up to eight characters in length
that defines the well name for which the well injection fluid thermal
properties are being defined. Note that the well name (WELNAME) must
have been declared previously using the WELSPECS keyword in the SCHEDULE
section, otherwise an error may occur. & None \\
2 & STEAMQAL & STEAMQAL is a real positive value greater than or equal
to zero and less than or equal to one that defines the steam quality of
the injected fluid for the defined well. This parameter should be
defaulted using 1* as STEAMQAL is not used by OPM FLOW, as only water
and gas injection is supported. This data is used by the commercial
simulator's THERMAL option and is not supported by OPM Flow's THERMAL
option. & 1* \\
dimensionless & dimensionless & dimensionless & \\
3 & TEMP & TEMP is a real positive value that defines the temperature of
the injected fluid for the defined well. & None \\
oF & oC & oC & \\
4 & PRES & PRES is a real positive value that defines the pressure of
the injected fluid for the defined well. & None \\
psia & barsa & atma & \\
5 & ENTHALPY & ENTHALPY is a real positive value that defines the
specific enthalpy of the injected fluid for the defined well. This is
data is used by the commercial simulator's THERMAL option and is not
supported by OPM Flow's THERMAL option. & None \\
Btu/lbs-M & kJ/kg-M & J/gm-M & \\
Notes: & & & \\
\end{longtable}

\emph{Table 12.102: WINJTEMP Keyword Description}

\subsubsection{Example}\label{example-588}

The following example shows the WINJTEMP keyword for when OPM Flow's
temperature option has been activated by the THERMAL keyword in the
RUNSPEC section.

\begin{verbatim}
--
--       INJECTION FLUID THERMAL PROPERTIES
--
-- WELL  STEAM  INJ    INJ    SPEC
-- NAME  QUAL   TEMP   PRES   ENTH
WINJTEMP
WI01     1*     68.0   220.0    1*                         /
WI02     1*     70.0   230.0    1*                         /
/
\end{verbatim}

Here the water injection fluid's temperature and pressure, in field
units, for two water injections well are defined. Notice that both the
steam quality and the specific enthalpy of the injected fluid for the
defined wells are defaulted (or skipped), as OPM Flow's THERMAL option
does not support this data.

\subsection{WLIFT -- Define Well Re-Tubing, THP and Lift Switching
Workover
Operations}\label{wlift-define-well-re-tubing-thp-and-lift-switching-workover-operations}

\begin{longtable}[]{@{}
  >{\raggedright\arraybackslash}p{(\linewidth - 14\tabcolsep) * \real{0.1250}}
  >{\raggedright\arraybackslash}p{(\linewidth - 14\tabcolsep) * \real{0.1250}}
  >{\raggedright\arraybackslash}p{(\linewidth - 14\tabcolsep) * \real{0.1250}}
  >{\raggedright\arraybackslash}p{(\linewidth - 14\tabcolsep) * \real{0.1250}}
  >{\raggedright\arraybackslash}p{(\linewidth - 14\tabcolsep) * \real{0.1250}}
  >{\raggedright\arraybackslash}p{(\linewidth - 14\tabcolsep) * \real{0.1250}}
  >{\raggedright\arraybackslash}p{(\linewidth - 14\tabcolsep) * \real{0.1250}}
  >{\raggedright\arraybackslash}p{(\linewidth - 14\tabcolsep) * \real{0.1250}}@{}}
\toprule\noalign{}
\begin{minipage}[b]{\linewidth}\raggedright
RUNSPEC
\end{minipage} & \begin{minipage}[b]{\linewidth}\raggedright
GRID
\end{minipage} & \begin{minipage}[b]{\linewidth}\raggedright
EDIT
\end{minipage} & \begin{minipage}[b]{\linewidth}\raggedright
PROPS
\end{minipage} & \begin{minipage}[b]{\linewidth}\raggedright
REGIONS
\end{minipage} & \begin{minipage}[b]{\linewidth}\raggedright
SOLUTION
\end{minipage} & \begin{minipage}[b]{\linewidth}\raggedright
SUMMARY
\end{minipage} & \begin{minipage}[b]{\linewidth}\raggedright
SCHEDULE
\end{minipage} \\
\midrule\noalign{}
\endhead
\bottomrule\noalign{}
\endlastfoot
\end{longtable}

\subsubsection{Description}\label{description-1092}

The WLIFT defines the automatic workovers parameters for changing out
wellbore tubing, changing the THP limit (for example switching from the
high stage pressure separator to the low stage pressure separator), or
changing the artificial lift parameters, for wells.

OPM Flow does not have this feature and hence this keyword is ignored by
OPM Flow and the WLIFT keyword has no effect on the simulation.

\subsection{WLIFTOPT -- Define Well Gas Lift Optimization
Parameters}\label{wliftopt-define-well-gas-lift-optimization-parameters}

\begin{longtable}[]{@{}
  >{\raggedright\arraybackslash}p{(\linewidth - 14\tabcolsep) * \real{0.1250}}
  >{\raggedright\arraybackslash}p{(\linewidth - 14\tabcolsep) * \real{0.1250}}
  >{\raggedright\arraybackslash}p{(\linewidth - 14\tabcolsep) * \real{0.1250}}
  >{\raggedright\arraybackslash}p{(\linewidth - 14\tabcolsep) * \real{0.1250}}
  >{\raggedright\arraybackslash}p{(\linewidth - 14\tabcolsep) * \real{0.1250}}
  >{\raggedright\arraybackslash}p{(\linewidth - 14\tabcolsep) * \real{0.1250}}
  >{\raggedright\arraybackslash}p{(\linewidth - 14\tabcolsep) * \real{0.1250}}
  >{\raggedright\arraybackslash}p{(\linewidth - 14\tabcolsep) * \real{0.1250}}@{}}
\toprule\noalign{}
\begin{minipage}[b]{\linewidth}\raggedright
RUNSPEC
\end{minipage} & \begin{minipage}[b]{\linewidth}\raggedright
GRID
\end{minipage} & \begin{minipage}[b]{\linewidth}\raggedright
EDIT
\end{minipage} & \begin{minipage}[b]{\linewidth}\raggedright
PROPS
\end{minipage} & \begin{minipage}[b]{\linewidth}\raggedright
REGIONS
\end{minipage} & \begin{minipage}[b]{\linewidth}\raggedright
SOLUTION
\end{minipage} & \begin{minipage}[b]{\linewidth}\raggedright
SUMMARY
\end{minipage} & \begin{minipage}[b]{\linewidth}\raggedright
SCHEDULE
\end{minipage} \\
\midrule\noalign{}
\endhead
\bottomrule\noalign{}
\endlastfoot
\end{longtable}

\subsubsection{Description}\label{description-1093}

The WLIFTOPT defines which wells should use the Gas Lift Optimization
facility in order to maximize oil production, as well as defining the
associated gas lift optimization parameters for a given well. The
keyword can also be used to switch off gas lift optimization for a well.
Gas lift optimization is invoked via the LIFTOPT keyword in the SCHEDULE
section. Note that the LIFTOPT keyword should precede the WLIFTOPT
keyword in the SCHEDULE section in order to activate the gas lift
optimization facility.

Wells are allocated to groups when the wells are specified by the
WELSPECS keyword in the SCHEDULE section. Wells defined to be under
group control will have their production rates and gas lift gas
constraints (GLIFTOPT keyword in the SCHEDULE section) controlled by the
group to which they belong, in addition to any well constraints defined
for the wells, including the gas optimization parameters on the WLIFTOPT
keyword.

\begin{longtable}[]{@{}
  >{\raggedright\arraybackslash}p{(\linewidth - 6\tabcolsep) * \real{0.2500}}
  >{\raggedright\arraybackslash}p{(\linewidth - 6\tabcolsep) * \real{0.2500}}
  >{\raggedright\arraybackslash}p{(\linewidth - 6\tabcolsep) * \real{0.2500}}
  >{\raggedright\arraybackslash}p{(\linewidth - 6\tabcolsep) * \real{0.2500}}@{}}
\toprule\noalign{}
\begin{minipage}[b]{\linewidth}\raggedright
No.
\end{minipage} & \begin{minipage}[b]{\linewidth}\raggedright
Name
\end{minipage} & \begin{minipage}[b]{\linewidth}\raggedright
Description
\end{minipage} & \begin{minipage}[b]{\linewidth}\raggedright
Default
\end{minipage} \\
\midrule\noalign{}
\endhead
\bottomrule\noalign{}
\endlastfoot
Field & Metric & Laboratory & \\
1 & WELNAME & A character string of up to eight characters in length
that defines the well name for which the well gas lift optimization
parameters are being defined. Note that the well name (WELNAME) must
have been declared previously using the WELSPECS keyword in the SCHEDULE
section, otherwise an error may occur. & None \\
2 & OPTLIFT & A defined character string that sets if a well's gas lift
gas rate should be calculated by the gas lift optimization facility or
not, and should be set to: & None \\
3 & MXLIFT & A real value that defines the total amount of gas lift gas
available for this well, multiplied by the well's efficiency factor.
Note that the value entered here should be in the range entered in the
VFPPROD table allocated to the well, otherwise errors may occur when
optimizing the gas lift gas injection rate for the well. & 1* \\
Mscf/d & sm3/day & scc/hour & \\
4 & OPTWGT & OPTWGT (βw) is real positive value that defines a weighting
factor for allocating the available gas lift gas to a well. An increment
of gas lift gas supply is allocated to a well based on the well's
current incremental gradient multiplied by OPTWGT using the following
formulae:
\(\mathit{Gradient} = \left(\frac{\left({β}_{w}\times \mathrm{Δ}{Q}_{\mathit{Oil}}\right)}{\left(\mathrm{Δ}{Q}_{\mathit{GasLift}} + {β}_{g}\times \mathrm{Δ}{Q}_{\mathit{Gas}}\right)}\right)\)
Where: βw = is OPTWGT, the weighting factor for the preferential
allocation of lift gas, βg = is the gas production rate weighting
factor, ∆QOil = is the increment/decrement in oil production rate, ∆Qgas
= is the increment/decrement in gas production rate, and ∆QGasLift = is
the increment/decrement in gas lift gas rate. Note by default βg, the
gas production rate weighting factor, is set to zero, and therefore the
gradient equation simplifies to:
\(\mathit{Gradient} = \left(\frac{{β}_{w}\times \mathrm{Δ}{Q}_{\mathit{Oil}}}{\mathrm{Δ}{Q}_{\mathit{GasLift}}}\right)\)
OPTWGT is ignored if OPTLIFT is equal to NO. & 1 \\
dimensionless & dimensionless & dimensionless & \\
5 & MINGAS & A real value that defines the minimum amount of gas lift
gas available to the well, multiplied by the well's efficiency factor.
The allocation of the gas lift gas is determined by: This parameter is
ignored if OPTLIFT is defined as NO. & 0.0 \\
Mscf/d & sm3/day & scc/hour & \\
6 & OPTGAS & OPTGAS (βg) is real positive value that defines the
incremental gas weighting factor for allocating the available gas lift
gas to a well. An increment of gas lift gas supply is allocated to a
well based on the well's current incremental gradient as described by
the definition of the OPTWGT (βg) variable above, that is by the
following formulae:
\(\mathit{Gradient} = \left(\frac{\left({β}_{w}\times \mathrm{Δ}{Q}_{\mathit{Oil}}\right)}{\left(\mathrm{Δ}{Q}_{\mathit{GasLift}} + {β}_{g}\times \mathrm{Δ}{Q}_{\mathit{Gas}}\right)}\right)\)
See OPTWGT for a definition of the variables in the equation. This
parameter is ignored if OPTLIFT is defined as NO. & 0.0 \\
dimensionless1 & dimensionless & dimensionless & \\
7 & OPTLIMIT & A defined character string that defines if additional gas
lift gas should be applied to the well, if the well's group gas target
has been satisfied but the group's oil rate limit has not been achieved.
This parameter is ignored if OPTLIFT is defined as NO. & NO \\
Notes: & & & \\
\end{longtable}

\emph{Table 12.103: WLIFTOPT Keyword Description}

See also the LIFTOPT keyword to activate gas lift optimization, the
GLIFTOPT keyword to define the group gas lift optimization controls,
GRUPTREE keyword to define the hierarchy of the groups below the FIELD
level, the WCONPROD and WCONINJE keywords to define a well's production
and injection rate targets and constraints. All the aforementioned
keywords are described in the SCHEDULE section.

\subsubsection{Example}\label{example-589}

The following example first switches on gas lift optimization via the
LIFTOPT keyword and then defines the artificial lift constraints for
PLAT-A, using the GLIFTOPT keyword and sets the well gas lift parameters
using the WLIFTOPT keyword.

\begin{verbatim}
--       ACTIVATE GAS LIFT OPTIMIZATION AND PARAMETERS
--
-- INCR   INCR   TSTEP    NEWTON
-- GAS    OIL    INTVAL   OPTN
LIFTOPT
12.5E3    5E-3   0.0      YES                                                  /
/
--
--       GROUP GAS LIFT OPTIMIZATION CONSTRAINTS
--
-- GRUP  MAX       MAX
-- NAME  GAS ALQ   TOTAL GAS
GLIFTOPT
PLAT-A    200E3    1*                                                          /
/
--
--       WELL GAS LIFT OPTIMIZATION PARAMETERS
--
-- WELL  OPTN   MAX    WEIGHT  MIN    GAS      OPTN
-- NAME  LIFT   LIFT   FACTOR  LIFT   FACTOR   LIMIT
WLIFTOPT
OP01     YES    150E3  1.01    -1.0                                            /
OP02     YES    150E3  1.01    -1.0                                            /
OPO3     YES    150E3  1.01    -1.0                                            /
OP04     YES    150E3  1.01    -1.0                                            /
OP05     YES    150E3  1.01    -1.0                                            /
/
\end{verbatim}

Here the LIFTOPT keyword defines the maximum incremental gas lift gas
quantity to be 12.5 x 103 m3, the minimum incremental oil gain per m3 of
gas lift gas is set to 5.0 x 10-3 m3, the time step interval is set to
zero to perform the gas optimization every time step, and finally the
gas lift optimization will be performed NUPCOL Newton iterations for the
time step.

The GLIFTOPT keyword sets the maximum amount of gas lift gas for PLAT-A
to 200,000 m3 and there is no maximum limit for the total maximum amount
of gas that the group can process. In addition, WLIFTOPT sets all five
wells to have gas lift gas optimization implemented with a maximum gas
lift gas value of 150,000 m3 per well, with equal weighting factors and
all wells are supplied with sufficient gas lift gas to allow the wells
to flow, (minimum lift set to a negative value) subject to the maximum
allowed gas lift quantity for the well (150,000 m3).

\subsection{WLIMTOL -- Define Well Constraint
Tolerance}\label{wlimtol-define-well-constraint-tolerance}

\begin{longtable}[]{@{}
  >{\raggedright\arraybackslash}p{(\linewidth - 14\tabcolsep) * \real{0.1250}}
  >{\raggedright\arraybackslash}p{(\linewidth - 14\tabcolsep) * \real{0.1250}}
  >{\raggedright\arraybackslash}p{(\linewidth - 14\tabcolsep) * \real{0.1250}}
  >{\raggedright\arraybackslash}p{(\linewidth - 14\tabcolsep) * \real{0.1250}}
  >{\raggedright\arraybackslash}p{(\linewidth - 14\tabcolsep) * \real{0.1250}}
  >{\raggedright\arraybackslash}p{(\linewidth - 14\tabcolsep) * \real{0.1250}}
  >{\raggedright\arraybackslash}p{(\linewidth - 14\tabcolsep) * \real{0.1250}}
  >{\raggedright\arraybackslash}p{(\linewidth - 14\tabcolsep) * \real{0.1250}}@{}}
\toprule\noalign{}
\begin{minipage}[b]{\linewidth}\raggedright
RUNSPEC
\end{minipage} & \begin{minipage}[b]{\linewidth}\raggedright
GRID
\end{minipage} & \begin{minipage}[b]{\linewidth}\raggedright
EDIT
\end{minipage} & \begin{minipage}[b]{\linewidth}\raggedright
PROPS
\end{minipage} & \begin{minipage}[b]{\linewidth}\raggedright
REGIONS
\end{minipage} & \begin{minipage}[b]{\linewidth}\raggedright
SOLUTION
\end{minipage} & \begin{minipage}[b]{\linewidth}\raggedright
SUMMARY
\end{minipage} & \begin{minipage}[b]{\linewidth}\raggedright
SCHEDULE
\end{minipage} \\
\midrule\noalign{}
\endhead
\bottomrule\noalign{}
\endlastfoot
\end{longtable}

\subsubsection{Description}\label{description-1094}

WLIMTOL keyword defines the tolerance to be used for various constraints
applied to connections, completions (if connections have been lumped via
the COMPLUMP keyword in the SCHEDULE section), wells, and groups,
including the field group. See also the GCONTOL keyword in the SCHEDULE
section that sets the tolerance parameters for groups.

This keyword is not supported by OPM Flow but would change the results
if supported so the simulation will be stopped.

\subsection{WLIST -- Define Well Lists
(Static)}\label{wlist-define-well-lists-static}

\begin{longtable}[]{@{}
  >{\raggedright\arraybackslash}p{(\linewidth - 14\tabcolsep) * \real{0.1250}}
  >{\raggedright\arraybackslash}p{(\linewidth - 14\tabcolsep) * \real{0.1250}}
  >{\raggedright\arraybackslash}p{(\linewidth - 14\tabcolsep) * \real{0.1250}}
  >{\raggedright\arraybackslash}p{(\linewidth - 14\tabcolsep) * \real{0.1250}}
  >{\raggedright\arraybackslash}p{(\linewidth - 14\tabcolsep) * \real{0.1250}}
  >{\raggedright\arraybackslash}p{(\linewidth - 14\tabcolsep) * \real{0.1250}}
  >{\raggedright\arraybackslash}p{(\linewidth - 14\tabcolsep) * \real{0.1250}}
  >{\raggedright\arraybackslash}p{(\linewidth - 14\tabcolsep) * \real{0.1250}}@{}}
\toprule\noalign{}
\begin{minipage}[b]{\linewidth}\raggedright
RUNSPEC
\end{minipage} & \begin{minipage}[b]{\linewidth}\raggedright
GRID
\end{minipage} & \begin{minipage}[b]{\linewidth}\raggedright
EDIT
\end{minipage} & \begin{minipage}[b]{\linewidth}\raggedright
PROPS
\end{minipage} & \begin{minipage}[b]{\linewidth}\raggedright
REGIONS
\end{minipage} & \begin{minipage}[b]{\linewidth}\raggedright
SOLUTION
\end{minipage} & \begin{minipage}[b]{\linewidth}\raggedright
SUMMARY
\end{minipage} & \begin{minipage}[b]{\linewidth}\raggedright
SCHEDULE
\end{minipage} \\
\midrule\noalign{}
\endhead
\bottomrule\noalign{}
\endlastfoot
\end{longtable}

\subsubsection{Description}\label{description-1095}

WLIST declares a group of wells to belong to a named static well list.
Wells in a named well list are treated as a group of wells for which the
standard well keywords can be applied. For example, instead of repeating
a well keyword for each well, the keyword only needs to have the named
well list instead, for the action to be applied to all wells in the
named well list. In general any well keyword that allows well name roots
as a well name, for example, PROD*, can use a named well list.

\begin{longtable}[]{@{}
  >{\raggedright\arraybackslash}p{(\linewidth - 6\tabcolsep) * \real{0.2500}}
  >{\raggedright\arraybackslash}p{(\linewidth - 6\tabcolsep) * \real{0.2500}}
  >{\raggedright\arraybackslash}p{(\linewidth - 6\tabcolsep) * \real{0.2500}}
  >{\raggedright\arraybackslash}p{(\linewidth - 6\tabcolsep) * \real{0.2500}}@{}}
\toprule\noalign{}
\begin{minipage}[b]{\linewidth}\raggedright
No.
\end{minipage} & \begin{minipage}[b]{\linewidth}\raggedright
Name
\end{minipage} & \begin{minipage}[b]{\linewidth}\raggedright
Description
\end{minipage} & \begin{minipage}[b]{\linewidth}\raggedright
Default
\end{minipage} \\
\midrule\noalign{}
\endhead
\bottomrule\noalign{}
\endlastfoot
1 & WLIST & A character string of up to eight characters in length,
enclosed in quotes, that defines the well list name for the WELNAMES
declared by this record. Note the first character must be asterisk
(``\emph{'') and the second character must be a letter, for example,
}PROD. & None \\
2 & ACTION & A defined character string that determines how the WELNAMES
should be handled with respect to the named well list (WLIST). ACTION
should be set to one of the following:: & \\
3-52 & WELNAMES & A character string of up to eight characters in length
that defines the well name that belongs to the named well list (WLIST).
A total of 50 well names can be added to WLIST at a time. If additional
wells are needed to added then use the ADD option of ACTION to add
additional wells. Well names roots may all be used in WELNAMES as long
as they are enclosed in quotes and end with an asterisk (``\emph{''). In
this case all wells that match the specification will be added to the
list. For example, wells named OP01, OP02 and OP03, can be added as
group by using ``OP}'' as the well name. Note that the well names must
have been declared previously using the WELSPECS keyword in the SCHEDULE
section, otherwise an error may occur & \\
Notes: & & & \\
\end{longtable}

\emph{Table 12.104: WLIST Keyword Description}

Note that named well list production data in the SUMMARY file is well
dependent, that is, if the wells belonging to a named well list is
changed through time, the SUMMARY data will be based on the wells in the
named well list group at the end of the run. Thus, if there are three
wells in a named well list called \emph{PROD1 at the beginning of a run;
OP01, OP02 and OP03, and during the run OP03 is moved to a well named
list called }PROD2, then the SUMMARY data for \emph{PROD1will only
contain the production data for wells OP01 and OP02 and }PROD2 will only
contain the SUMMARY data for OP03 from the start to the end of the run.

\subsubsection{Example}\label{example-590}

The following example defines two named well lists using the WLIST
keyword.

\begin{verbatim}
--
--          WELL LIST SPECIFICATION
--
-- LIST     OPER     WELL NAME LIST
-- NAME
WLIST
'*BLK-1'    NEW      WEL-01M WEL-02M WEL-03M WEL-04M WEL-05M WEL-06M WEL-07M   /
'*BLK-1'    ADD      WEL-08M WEL-09M WEL-10M WEL-11M WEL-12M WEL-13M WEL-14M   /
'*BLK-1'    ADD      WEL-15M WEL-16M WEL-17M WEL-18M WEL-19M WEL-20M WEL-23M   /
'*BLK-1'    ADD      WEL-24M WEL-25M WEL-26M WEL-28M                           /

'*BLK-2'    NEW      WEL-03U WEL-05U WEL-06U WEL-10U WEL-11U WEL-13U WEL-14U   /
'*BLK-2'    ADD      WEL-15U WEL-16U WEL-17U WEL-18U WEL-19U WEL-25U WEL-27U   /
/
DATES
          1  JAN   2020  /
/
--
--       DEFINE WELL AND WELL CONNECTIONS FLOWING STATUS
--
--  WELL WELL   --LOCATION--  COMPLETION
--  NAME STAT     I   J    K  FIRST LAST
WELOPEN
'*BLK-1' OPEN                                              /
'*BLK-1' OPEN     0   0    0     0     0                   /
/
DATES
          1  JAN   2021  /
          1  JLY   2021  /
          1  OCT   2021  /
/
--
--       DEFINE WELL AND WELL CONNECTIONS FLOWING STATUS
--
--  WELL WELL   --LOCATION--  COMPLETION
--  NAME STAT     I   J    K  FIRST LAST
WELOPEN
'*BLK-2' OPEN                                              /
'*BLK-2' OPEN     0   0    0     0     0                   /
/
\end{verbatim}

In this example the wells in named well list '\,'\emph{BLK-1'' are
opened on January 1, 2020 and wells in named well list '}BLK-2'' are
opened October 1, 2021.

\subsection{WLISTARG -- Modify Well List Target and Constraint Values
(Static)}\label{wlistarg-modify-well-list-target-and-constraint-values-static}

\begin{longtable}[]{@{}
  >{\raggedright\arraybackslash}p{(\linewidth - 14\tabcolsep) * \real{0.1250}}
  >{\raggedright\arraybackslash}p{(\linewidth - 14\tabcolsep) * \real{0.1250}}
  >{\raggedright\arraybackslash}p{(\linewidth - 14\tabcolsep) * \real{0.1250}}
  >{\raggedright\arraybackslash}p{(\linewidth - 14\tabcolsep) * \real{0.1250}}
  >{\raggedright\arraybackslash}p{(\linewidth - 14\tabcolsep) * \real{0.1250}}
  >{\raggedright\arraybackslash}p{(\linewidth - 14\tabcolsep) * \real{0.1250}}
  >{\raggedright\arraybackslash}p{(\linewidth - 14\tabcolsep) * \real{0.1250}}
  >{\raggedright\arraybackslash}p{(\linewidth - 14\tabcolsep) * \real{0.1250}}@{}}
\toprule\noalign{}
\begin{minipage}[b]{\linewidth}\raggedright
RUNSPEC
\end{minipage} & \begin{minipage}[b]{\linewidth}\raggedright
GRID
\end{minipage} & \begin{minipage}[b]{\linewidth}\raggedright
EDIT
\end{minipage} & \begin{minipage}[b]{\linewidth}\raggedright
PROPS
\end{minipage} & \begin{minipage}[b]{\linewidth}\raggedright
REGIONS
\end{minipage} & \begin{minipage}[b]{\linewidth}\raggedright
SOLUTION
\end{minipage} & \begin{minipage}[b]{\linewidth}\raggedright
SUMMARY
\end{minipage} & \begin{minipage}[b]{\linewidth}\raggedright
SCHEDULE
\end{minipage} \\
\midrule\noalign{}
\endhead
\bottomrule\noalign{}
\endlastfoot
\end{longtable}

\subsubsection{Description}\label{description-1096}

The WLISTARG keyword modifies the target and constraint values of both
rates and pressures for wells previously defined in a well list by the
WLIST or WLISTNAM keywords. WLISTARG is similar to the WELTARG keyword
in it that allows for modifying targets and constraints without having
to define all the variables on the well control keywords: WCONPROD,
WCONHIST, WCONINJE, or WCONINJH keywords. Variables not changed by the
WLISTARG keyword remain the same as those previously entered via the
well control keywords or previously entered WLISTARG keywords. Note that
the well must still be initially be fully defined using the WCONPROD or
WCONINJE keywords. All the aforementioned keywords are described in the
SCHEDULE section.

This keyword is not supported by OPM Flow but would change the results
if supported so the simulation will be stopped.

\begin{longtable}[]{@{}
  >{\raggedright\arraybackslash}p{(\linewidth - 6\tabcolsep) * \real{0.2500}}
  >{\raggedright\arraybackslash}p{(\linewidth - 6\tabcolsep) * \real{0.2500}}
  >{\raggedright\arraybackslash}p{(\linewidth - 6\tabcolsep) * \real{0.2500}}
  >{\raggedright\arraybackslash}p{(\linewidth - 6\tabcolsep) * \real{0.2500}}@{}}
\toprule\noalign{}
\begin{minipage}[b]{\linewidth}\raggedright
No.
\end{minipage} & \begin{minipage}[b]{\linewidth}\raggedright
Name
\end{minipage} & \begin{minipage}[b]{\linewidth}\raggedright
Description
\end{minipage} & \begin{minipage}[b]{\linewidth}\raggedright
Default
\end{minipage} \\
\midrule\noalign{}
\endhead
\bottomrule\noalign{}
\endlastfoot
Field & Metric & Laboratory & \\
1 & WLIST & A character string of up to eight characters in length,
enclosed in quotes, that defines the well list name declared by the
WLIST keyword. Note the first character must be asterisk (``\emph{'')
and the second character must be a letter, for example, }PROD. & None \\
2 & TARGET & A defined character string that sets the item to be changed
for the well the value of the item is set by item (3). & None \\
3 & VALUE Liquid Gas Res Vol Pressure VFP LIFT & A real positive vector
of values that defines the value of the variable declared by TARGET for
all the wells contained in WLIST. For example if there are four wells in
WLIST then there must four real numbers for VALUE. The vector should be
terminated by a ``/'' as indicated in the notes below. & None \\
stb/d Mscf/d rb/d psia dimensionless same as VFPPROD or VFPINJ & sm3/day
sm3/day rm3/day barsa dimensionless same as VFPPROD or VFPINJ & scc/hour
scc/hour rcc/hour atma dimensionless same as VFPPROD or VFPINJ & \\
Notes: & & & \\
\end{longtable}

\emph{Table 12.105: WLISTARG Keyword Description}

If a well is currently a history matching well, then WLISELTARG should
only be used to change a wells bottom-hole pressure limit, vertical flow
performance table number or the artificial lift quantity.

See also the WELCNTL keyword, in the SCHEDULE section that can be used
to reset the control mode, as well as a well's target and constraints of
both rates and pressures.

\subsubsection{Example}\label{example-591}

The following example defines two named well lists using the WLIST
keyword.

\begin{verbatim}
--
--       WELL LIST SPECIFICATION
--
-- LIST  OPER     WELL NAME LIST
-- NAME
WLIST
'*BLK-1' NEW      WEL-01M WEL-02M WEL-03M                                   /
'*BLK-2' NEW      WEL-03U WEL-05U WEL-06U WEL-10U                           /
/
--
--       WELL PRODUCTION AND INJECTION TARGETS
--
--  WELL WELL   TARGET
--  NAME TARG   VALUE
WLISTARG
‘*BLK-1’ ORAT   2000.0  2000.00  2000.0                     /
‘*BLK-2’ ORAT   3000.0  3500.00  4000.0  2000.0             /
/
\end{verbatim}

The wells in the `\emph{BLK-1' well list are all given an oil rate of
2,000 stb/d and wells in the '}BLK-2' well list are given rates of
3,000, 3,500, 4,000 and 2,000 stb/d.

\subsection{WLISTNAM -- Define Well Lists
(WLISTARG)}\label{wlistnam-define-well-lists-wlistarg}

\begin{longtable}[]{@{}
  >{\raggedright\arraybackslash}p{(\linewidth - 14\tabcolsep) * \real{0.1250}}
  >{\raggedright\arraybackslash}p{(\linewidth - 14\tabcolsep) * \real{0.1250}}
  >{\raggedright\arraybackslash}p{(\linewidth - 14\tabcolsep) * \real{0.1250}}
  >{\raggedright\arraybackslash}p{(\linewidth - 14\tabcolsep) * \real{0.1250}}
  >{\raggedright\arraybackslash}p{(\linewidth - 14\tabcolsep) * \real{0.1250}}
  >{\raggedright\arraybackslash}p{(\linewidth - 14\tabcolsep) * \real{0.1250}}
  >{\raggedright\arraybackslash}p{(\linewidth - 14\tabcolsep) * \real{0.1250}}
  >{\raggedright\arraybackslash}p{(\linewidth - 14\tabcolsep) * \real{0.1250}}@{}}
\toprule\noalign{}
\begin{minipage}[b]{\linewidth}\raggedright
RUNSPEC
\end{minipage} & \begin{minipage}[b]{\linewidth}\raggedright
GRID
\end{minipage} & \begin{minipage}[b]{\linewidth}\raggedright
EDIT
\end{minipage} & \begin{minipage}[b]{\linewidth}\raggedright
PROPS
\end{minipage} & \begin{minipage}[b]{\linewidth}\raggedright
REGIONS
\end{minipage} & \begin{minipage}[b]{\linewidth}\raggedright
SOLUTION
\end{minipage} & \begin{minipage}[b]{\linewidth}\raggedright
SUMMARY
\end{minipage} & \begin{minipage}[b]{\linewidth}\raggedright
SCHEDULE
\end{minipage} \\
\midrule\noalign{}
\endhead
\bottomrule\noalign{}
\endlastfoot
\end{longtable}

\subsubsection{Description}\label{description-1097}

WLISTNAM declares a group of wells to belong to a named WLISTARG well
list for use with the WLISTARG keyword. Only the WLISTARG keyword can be
used with this type of well list, and therefore it is better to use the
WLIST keyword instead, that defines a static well list but offers more
flexibility than a WLISTNAM well list.

This keyword is not supported by OPM Flow but would change the results
if supported so the simulation will be stopped.

\begin{longtable}[]{@{}
  >{\raggedright\arraybackslash}p{(\linewidth - 6\tabcolsep) * \real{0.2500}}
  >{\raggedright\arraybackslash}p{(\linewidth - 6\tabcolsep) * \real{0.2500}}
  >{\raggedright\arraybackslash}p{(\linewidth - 6\tabcolsep) * \real{0.2500}}
  >{\raggedright\arraybackslash}p{(\linewidth - 6\tabcolsep) * \real{0.2500}}@{}}
\toprule\noalign{}
\begin{minipage}[b]{\linewidth}\raggedright
No.
\end{minipage} & \begin{minipage}[b]{\linewidth}\raggedright
Name
\end{minipage} & \begin{minipage}[b]{\linewidth}\raggedright
Description
\end{minipage} & \begin{minipage}[b]{\linewidth}\raggedright
Default
\end{minipage} \\
\midrule\noalign{}
\endhead
\bottomrule\noalign{}
\endlastfoot
1 & WLIST & A character string of up to eight characters in length,
enclosed in quotes, that defines the well list name for the WELLNAMES
declared by this record. Note the first character must be asterisk
(``\emph{'') and the second character must be a letter, for example,
}PROD. & None \\
2-51 & WELNAMES & A character string of up to eight characters in length
that defines the well name that belongs to the named well list (WLIST).
A total of 50 well names can be added to WLISTNAM at a time. If the
first well name in the list is the default value (``\emph{1''), then the
list is first cleared of all wells, before adding the subsequent wells
in WELLNAMES. Well names roots may all be used in WELLNAMES as long as
they are enclosed in quotes and end with an asterisk (``}''). In this
case all wells that match the specification will be added to the list.
For example, wells named OP01, OP02 and OP03, can be added as group by
using ``OP*'' as the well name. Note that the well names must have been
declared previously using the WELSPECS keyword in the SCHEDULE section,
otherwise an error may occur. & 1* \\
Notes: & & & \\
\end{longtable}

\emph{Table 12.106: WLISTNAM Keyword Description}

Note that named well list production data in the SUMMARY file is well
dependent, that is, if the wells belonging to a named well list is
changed through time, the SUMMARY data will be based on the wells in the
named well list group at the end of the run. Thus, if there are three
wells in a named well list called \emph{PROD1 at the beginning of a run;
OP01, OP02 and OP03, and during the run OP03 is moved to a well named
list called }PROD2, then the SUMMARY data for \emph{PROD1will only
contain the production data for wells OP01 and OP02 and }PROD2 will only
contain the SUMMARY data for OP03 from the start to the end of the run.

\subsubsection{Example}\label{example-592}

The following example defines two named well lists using the WLISTNAM
keyword.

\begin{verbatim}
--
--          WELL LIST SPECIFICATION
--
-- LIST     WELL NAME LIST
-- NAME
WLISTNAM
'*BLK-1'    WEL-01M WEL-02M WEL-03M WEL-04M WEL-05M WEL-06M WEL-07M            /
'*BLK-1'    WEL-08M WEL-09M WEL-10M WEL-11M WEL-12M WEL-13M WEL-14M            /
'*BLK-1'    WEL-15M WEL-16M WEL-17M WEL-18M WEL-19M WEL-20M WEL-23M            /
'*BLK-1'    WEL-24M WEL-25M WEL-26M WEL-28M                                    /

'*BLK-2'    1*      WEL-03U WEL-05U WEL-06U WEL-10U WEL-11U WEL-13U WEL-14U    /
'*BLK-2'            WEL-15U WEL-16U WEL-17U WEL-18U WEL-19U WEL-25U WEL-27U    /
/
\end{verbatim}

Here well list '\,'\emph{BLK-1' contains 28 wells, that is wells WEL-01M
to WEL-28M. For the '}BLK-2' well list all wells are first deleted due
to the ``1*'' default value and then wells WEL-03U to WEL-27U are added
to the list.

\subsection{WMICP -- Define Water Injection Well's Microbial, Oxygen,
And Urea
Concentrations}\label{wmicp-define-water-injection-wells-microbial-oxygen-and-urea-concentrations}

\begin{longtable}[]{@{}
  >{\raggedright\arraybackslash}p{(\linewidth - 14\tabcolsep) * \real{0.1250}}
  >{\raggedright\arraybackslash}p{(\linewidth - 14\tabcolsep) * \real{0.1250}}
  >{\raggedright\arraybackslash}p{(\linewidth - 14\tabcolsep) * \real{0.1250}}
  >{\raggedright\arraybackslash}p{(\linewidth - 14\tabcolsep) * \real{0.1250}}
  >{\raggedright\arraybackslash}p{(\linewidth - 14\tabcolsep) * \real{0.1250}}
  >{\raggedright\arraybackslash}p{(\linewidth - 14\tabcolsep) * \real{0.1250}}
  >{\raggedright\arraybackslash}p{(\linewidth - 14\tabcolsep) * \real{0.1250}}
  >{\raggedright\arraybackslash}p{(\linewidth - 14\tabcolsep) * \real{0.1250}}@{}}
\toprule\noalign{}
\begin{minipage}[b]{\linewidth}\raggedright
RUNSPEC
\end{minipage} & \begin{minipage}[b]{\linewidth}\raggedright
GRID
\end{minipage} & \begin{minipage}[b]{\linewidth}\raggedright
EDIT
\end{minipage} & \begin{minipage}[b]{\linewidth}\raggedright
PROPS
\end{minipage} & \begin{minipage}[b]{\linewidth}\raggedright
REGIONS
\end{minipage} & \begin{minipage}[b]{\linewidth}\raggedright
SOLUTION
\end{minipage} & \begin{minipage}[b]{\linewidth}\raggedright
SUMMARY
\end{minipage} & \begin{minipage}[b]{\linewidth}\raggedright
SCHEDULE
\end{minipage} \\
\midrule\noalign{}
\endhead
\bottomrule\noalign{}
\endlastfoot
\end{longtable}

\subsubsection{Description}\label{description-1098}

The WMICP keyword defines a water injection well's microbial, growth,
and cementation injection stream solutions, where the rate-limiting
components are suspended microbes, oxygen, and urea concentrations
respectively. These concentrations are used when the MICP keyword in the
RUNSPEC section has been used to activate OPM Flow's Microbially Induced
Calcite Precipitation model. See Landa-Marbán et al {[}Landa-Marbán, D.,
Tveit, S., Kumar, K., Gasda, S.E., 2021. Practical approaches to study
microbially induced calcite precipitation at the eld scale. Int. J.
Greenh. Gas Control 106, 103256.
https://doi.org/10.1016/j.ijggc.2021.103256.{]} and {[}Landa-Marbán, D.,
Kumar, K., Tveit, S., Gasda, S.E., 2021. Numerical studies of CO2
leakage remediation by micp-based plugging technology. In: Røkke, N.A.
and Knuutila, H.K. (Eds) Short Papers from the 11th International
Trondheim CCS conference, ISBN: 978-82-536-1714-5, 284-290.{]} for a
description of the model.

Note the keyword should only be used for wells declared as water
injection wells via the WCONINJE keyword in the SCHEDULE section.

\begin{tcolorbox}[enhanced jigsaw, title=\textcolor{quarto-callout-note-color}{\faInfo}\hspace{0.5em}{Note}, coltitle=black, breakable, toprule=.15mm, colbacktitle=quarto-callout-note-color!10!white, leftrule=.75mm, opacitybacktitle=0.6, colback=white, bottomtitle=1mm, toptitle=1mm, titlerule=0mm, left=2mm, colframe=quarto-callout-note-color-frame, arc=.35mm, rightrule=.15mm, bottomrule=.15mm, opacityback=0]

This is an OPM Flow specific keyword.

\end{tcolorbox}

\begin{longtable}[]{@{}
  >{\raggedright\arraybackslash}p{(\linewidth - 6\tabcolsep) * \real{0.2500}}
  >{\raggedright\arraybackslash}p{(\linewidth - 6\tabcolsep) * \real{0.2500}}
  >{\raggedright\arraybackslash}p{(\linewidth - 6\tabcolsep) * \real{0.2500}}
  >{\raggedright\arraybackslash}p{(\linewidth - 6\tabcolsep) * \real{0.2500}}@{}}
\toprule\noalign{}
\begin{minipage}[b]{\linewidth}\raggedright
No.
\end{minipage} & \begin{minipage}[b]{\linewidth}\raggedright
Name
\end{minipage} & \begin{minipage}[b]{\linewidth}\raggedright
Description
\end{minipage} & \begin{minipage}[b]{\linewidth}\raggedright
Default
\end{minipage} \\
\midrule\noalign{}
\endhead
\bottomrule\noalign{}
\endlastfoot
Field & Metric & Laboratory & \\
1 & WELNAME & A character string of up to eight characters in length
that defines the well name for which the well connection data is being
defined. Note that the well name (WELNAME) must have been declared
previously using the WELSPECS keyword in the SCHEDULE section, otherwise
an error may occur. & None \\
2 & MICRCON & MICRCON is a real positive value that defines the
microbial concentration of the well's injection stream. & 0.0 \\
lb/stb & kg/sm3 & gm/scc & \\
3 & OXYGCON & A real positive value that defines the oxygen
concentration of the well's injection stream. & 0.0 \\
lb/stb & kg/sm3 & gm/scc & \\
4 & UREACON & UREACON is a real positive value that defines the urea
concentration of the well's injection stream. & 0.0 \\
lb/stb & kg/sm3 & gm/scc & \\
Notes: & & & \\
\end{longtable}

\emph{Table 12.107: WMICP Keyword Description}

Water injection wells that are not declared via this keyword have their
concentrations defaulted to zero.

\subsubsection{Example}\label{example-593}

\begin{verbatim}
--
--       DEFINE WATER INJECTION WELL MICROBIAL, OXYGEN, AND UREA CONCENTRATIONS
--
-- WELL  MICROBIAL  OXYGEN     UREA
-- NAME  MICRCON    OXYGCON    UREACON
--       --------   --------   --------
WMICP
WI01     0.01                                    /
WI02     1*         0.04                         /
WI03     1*         1*         60.0              /
WI04     1*         0.04       60.0              /
/
\end{verbatim}

Here the microbial concentration for well WI01 is set to 0.01, the
oxygen concentration for well WI02 is set to 0.04, the urea
concentration for well WI03 is set to 60, and the oxygen and urea
concentrations for well WI04 are set to 0.04 and 60 respectively.

\subsection{WNETCTRL -- Define Well Control for Network Control
Option}\label{wnetctrl-define-well-control-for-network-control-option}

\begin{longtable}[]{@{}
  >{\raggedright\arraybackslash}p{(\linewidth - 14\tabcolsep) * \real{0.1250}}
  >{\raggedright\arraybackslash}p{(\linewidth - 14\tabcolsep) * \real{0.1250}}
  >{\raggedright\arraybackslash}p{(\linewidth - 14\tabcolsep) * \real{0.1250}}
  >{\raggedright\arraybackslash}p{(\linewidth - 14\tabcolsep) * \real{0.1250}}
  >{\raggedright\arraybackslash}p{(\linewidth - 14\tabcolsep) * \real{0.1250}}
  >{\raggedright\arraybackslash}p{(\linewidth - 14\tabcolsep) * \real{0.1250}}
  >{\raggedright\arraybackslash}p{(\linewidth - 14\tabcolsep) * \real{0.1250}}
  >{\raggedright\arraybackslash}p{(\linewidth - 14\tabcolsep) * \real{0.1250}}@{}}
\toprule\noalign{}
\begin{minipage}[b]{\linewidth}\raggedright
RUNSPEC
\end{minipage} & \begin{minipage}[b]{\linewidth}\raggedright
GRID
\end{minipage} & \begin{minipage}[b]{\linewidth}\raggedright
EDIT
\end{minipage} & \begin{minipage}[b]{\linewidth}\raggedright
PROPS
\end{minipage} & \begin{minipage}[b]{\linewidth}\raggedright
REGIONS
\end{minipage} & \begin{minipage}[b]{\linewidth}\raggedright
SOLUTION
\end{minipage} & \begin{minipage}[b]{\linewidth}\raggedright
SUMMARY
\end{minipage} & \begin{minipage}[b]{\linewidth}\raggedright
SCHEDULE
\end{minipage} \\
\midrule\noalign{}
\endhead
\bottomrule\noalign{}
\endlastfoot
\end{longtable}

\subsubsection{Description}\label{description-1099}

The WNETCNTL keyword sets a well's control mode that should remain fixed
after each network balancing calculation, for when the either the
Standard Network or the Extended Network options have been activated,
and the well is part of a network. The keyword allows for a well's
Tubing Head Pressure (``THP''), oil, gas, liquid, or water rate to be
selected as fixed after each network balance calculation. Normally this
should be the THP, and if the keyword is absent from the input deck then
THP will be used as the default value. The Standard Network option is
invoked if the GRUPTREE, GRUPNET, GNETINJE, GNETPUMP, etc. series of
keywords have been used in the SCHEDULE section. Whereas, the Extended
Network option is activated by the NETWORK keyword in the RUNSPEC
section.

This keyword is not supported by OPM Flow but would change the results
if supported so the simulation will be stopped.

\subsection{WNETDP -- Define Well THP to Network Pressure
Drop}\label{wnetdp-define-well-thp-to-network-pressure-drop}

\begin{longtable}[]{@{}
  >{\raggedright\arraybackslash}p{(\linewidth - 14\tabcolsep) * \real{0.1250}}
  >{\raggedright\arraybackslash}p{(\linewidth - 14\tabcolsep) * \real{0.1250}}
  >{\raggedright\arraybackslash}p{(\linewidth - 14\tabcolsep) * \real{0.1250}}
  >{\raggedright\arraybackslash}p{(\linewidth - 14\tabcolsep) * \real{0.1250}}
  >{\raggedright\arraybackslash}p{(\linewidth - 14\tabcolsep) * \real{0.1250}}
  >{\raggedright\arraybackslash}p{(\linewidth - 14\tabcolsep) * \real{0.1250}}
  >{\raggedright\arraybackslash}p{(\linewidth - 14\tabcolsep) * \real{0.1250}}
  >{\raggedright\arraybackslash}p{(\linewidth - 14\tabcolsep) * \real{0.1250}}@{}}
\toprule\noalign{}
\begin{minipage}[b]{\linewidth}\raggedright
RUNSPEC
\end{minipage} & \begin{minipage}[b]{\linewidth}\raggedright
GRID
\end{minipage} & \begin{minipage}[b]{\linewidth}\raggedright
EDIT
\end{minipage} & \begin{minipage}[b]{\linewidth}\raggedright
PROPS
\end{minipage} & \begin{minipage}[b]{\linewidth}\raggedright
REGIONS
\end{minipage} & \begin{minipage}[b]{\linewidth}\raggedright
SOLUTION
\end{minipage} & \begin{minipage}[b]{\linewidth}\raggedright
SUMMARY
\end{minipage} & \begin{minipage}[b]{\linewidth}\raggedright
SCHEDULE
\end{minipage} \\
\midrule\noalign{}
\endhead
\bottomrule\noalign{}
\endlastfoot
\end{longtable}

\subsubsection{Description}\label{description-1100}

The WNETDP keyword allows for a constant pressure drop between a well's
Tubing Head Pressure (``THP'') and the well's connecting network node,
for when the either the Standard Network or the Extended Network options
have been activated, and the well is part of a network. For production
wells in a production network, WNETDP is added to the well's connecting
network node pressure to arrive at the well's THP value. Whereas for
injection wells in an injection network, WNETDP is subtracted from the
well's connecting network node pressure to arrive at the well's THP
value. The Standard Network option is invoked if the GRUPTREE, GRUPNET,
GNETINJE, GNETPUMP, etc. series of keywords have been used in the
SCHEDULE section. The Extended Network option is activated by the
NETWORK keyword in the RUNSPEC section.

This keyword is not supported by OPM Flow but would change the results
if supported so the simulation will be stopped.

\subsection{WORKLIM -- Define Well Workover
Time}\label{worklim-define-well-workover-time}

\begin{longtable}[]{@{}
  >{\raggedright\arraybackslash}p{(\linewidth - 14\tabcolsep) * \real{0.1250}}
  >{\raggedright\arraybackslash}p{(\linewidth - 14\tabcolsep) * \real{0.1250}}
  >{\raggedright\arraybackslash}p{(\linewidth - 14\tabcolsep) * \real{0.1250}}
  >{\raggedright\arraybackslash}p{(\linewidth - 14\tabcolsep) * \real{0.1250}}
  >{\raggedright\arraybackslash}p{(\linewidth - 14\tabcolsep) * \real{0.1250}}
  >{\raggedright\arraybackslash}p{(\linewidth - 14\tabcolsep) * \real{0.1250}}
  >{\raggedright\arraybackslash}p{(\linewidth - 14\tabcolsep) * \real{0.1250}}
  >{\raggedright\arraybackslash}p{(\linewidth - 14\tabcolsep) * \real{0.1250}}@{}}
\toprule\noalign{}
\begin{minipage}[b]{\linewidth}\raggedright
RUNSPEC
\end{minipage} & \begin{minipage}[b]{\linewidth}\raggedright
GRID
\end{minipage} & \begin{minipage}[b]{\linewidth}\raggedright
EDIT
\end{minipage} & \begin{minipage}[b]{\linewidth}\raggedright
PROPS
\end{minipage} & \begin{minipage}[b]{\linewidth}\raggedright
REGIONS
\end{minipage} & \begin{minipage}[b]{\linewidth}\raggedright
SOLUTION
\end{minipage} & \begin{minipage}[b]{\linewidth}\raggedright
SUMMARY
\end{minipage} & \begin{minipage}[b]{\linewidth}\raggedright
SCHEDULE
\end{minipage} \\
\midrule\noalign{}
\endhead
\bottomrule\noalign{}
\endlastfoot
\end{longtable}

\subsubsection{Description}\label{description-1101}

WORKLIM sets the numbers of days taken to complete a workover.

This keyword is not supported by OPM Flow but would change the results
if supported so the simulation will be stopped.

\subsection{WORKTHP -- Define Well Workover Options for THP Killed
Wells}\label{workthp-define-well-workover-options-for-thp-killed-wells}

\begin{longtable}[]{@{}
  >{\raggedright\arraybackslash}p{(\linewidth - 14\tabcolsep) * \real{0.1250}}
  >{\raggedright\arraybackslash}p{(\linewidth - 14\tabcolsep) * \real{0.1250}}
  >{\raggedright\arraybackslash}p{(\linewidth - 14\tabcolsep) * \real{0.1250}}
  >{\raggedright\arraybackslash}p{(\linewidth - 14\tabcolsep) * \real{0.1250}}
  >{\raggedright\arraybackslash}p{(\linewidth - 14\tabcolsep) * \real{0.1250}}
  >{\raggedright\arraybackslash}p{(\linewidth - 14\tabcolsep) * \real{0.1250}}
  >{\raggedright\arraybackslash}p{(\linewidth - 14\tabcolsep) * \real{0.1250}}
  >{\raggedright\arraybackslash}p{(\linewidth - 14\tabcolsep) * \real{0.1250}}@{}}
\toprule\noalign{}
\begin{minipage}[b]{\linewidth}\raggedright
RUNSPEC
\end{minipage} & \begin{minipage}[b]{\linewidth}\raggedright
GRID
\end{minipage} & \begin{minipage}[b]{\linewidth}\raggedright
EDIT
\end{minipage} & \begin{minipage}[b]{\linewidth}\raggedright
PROPS
\end{minipage} & \begin{minipage}[b]{\linewidth}\raggedright
REGIONS
\end{minipage} & \begin{minipage}[b]{\linewidth}\raggedright
SOLUTION
\end{minipage} & \begin{minipage}[b]{\linewidth}\raggedright
SUMMARY
\end{minipage} & \begin{minipage}[b]{\linewidth}\raggedright
SCHEDULE
\end{minipage} \\
\midrule\noalign{}
\endhead
\bottomrule\noalign{}
\endlastfoot
\end{longtable}

\subsubsection{Description}\label{description-1102}

The WORKTHP keyword defines workover options for when a well dies, that
is unable to produce at the current operating conditions, when under
Tubing Head Pressure (``THP'') control. For example, if a well is
producing to the high pressure separator and therefore has a high THP
constraint, then the WORKTHP keyword can be used to switch the well to
the lower pressure separator via re-setting the THP constraint.

This keyword is not supported by OPM Flow but would change the results
if supported so the simulation will be stopped.

\subsection{Atgeirr Rasmussen}\label{atgeirr-rasmussen-2}

\begin{verbatim}
  2017-09-22T14:01:41.387549000
  AFR
  Not supported. It is used on Norne, but investigations showed no significant effect so we did not implement it.
 WPAVE
  David Baxendale
  2017-10-03T12:19:22.164000000
  DBx
  Reply to Atgeirr Rasmussen (09/22/2017, 14:01): "..."
  Okay, but if we done the work then perhaps we should implement it as it is one more feature we can say is implemented versus the commercial simulator.

  Will change to not implemented.
  – Well Block Average Pressure Calculation Parameters for All Wells
\end{verbatim}

\begin{longtable}[]{@{}
  >{\raggedright\arraybackslash}p{(\linewidth - 14\tabcolsep) * \real{0.1250}}
  >{\raggedright\arraybackslash}p{(\linewidth - 14\tabcolsep) * \real{0.1250}}
  >{\raggedright\arraybackslash}p{(\linewidth - 14\tabcolsep) * \real{0.1250}}
  >{\raggedright\arraybackslash}p{(\linewidth - 14\tabcolsep) * \real{0.1250}}
  >{\raggedright\arraybackslash}p{(\linewidth - 14\tabcolsep) * \real{0.1250}}
  >{\raggedright\arraybackslash}p{(\linewidth - 14\tabcolsep) * \real{0.1250}}
  >{\raggedright\arraybackslash}p{(\linewidth - 14\tabcolsep) * \real{0.1250}}
  >{\raggedright\arraybackslash}p{(\linewidth - 14\tabcolsep) * \real{0.1250}}@{}}
\toprule\noalign{}
\begin{minipage}[b]{\linewidth}\raggedright
RUNSPEC
\end{minipage} & \begin{minipage}[b]{\linewidth}\raggedright
GRID
\end{minipage} & \begin{minipage}[b]{\linewidth}\raggedright
EDIT
\end{minipage} & \begin{minipage}[b]{\linewidth}\raggedright
PROPS
\end{minipage} & \begin{minipage}[b]{\linewidth}\raggedright
REGIONS
\end{minipage} & \begin{minipage}[b]{\linewidth}\raggedright
SOLUTION
\end{minipage} & \begin{minipage}[b]{\linewidth}\raggedright
SUMMARY
\end{minipage} & \begin{minipage}[b]{\linewidth}\raggedright
SCHEDULE
\end{minipage} \\
\midrule\noalign{}
\endhead
\bottomrule\noalign{}
\endlastfoot
\end{longtable}

\subsubsection{Description}\label{description-1103}

The WPAVE keyword defines the method and parameters for calculating a
well's block average pressures for all wells in the model. The resulting
average pressure can be written out to the SUMMARY and RSM files in
order to compare with field observed data via the WBP, WBP4, WBP5 and
WBP9 vectors in the SUMMARY section.

\begin{longtable}[]{@{}
  >{\raggedright\arraybackslash}p{(\linewidth - 6\tabcolsep) * \real{0.2500}}
  >{\raggedright\arraybackslash}p{(\linewidth - 6\tabcolsep) * \real{0.2500}}
  >{\raggedright\arraybackslash}p{(\linewidth - 6\tabcolsep) * \real{0.2500}}
  >{\raggedright\arraybackslash}p{(\linewidth - 6\tabcolsep) * \real{0.2500}}@{}}
\toprule\noalign{}
\begin{minipage}[b]{\linewidth}\raggedright
No.
\end{minipage} & \begin{minipage}[b]{\linewidth}\raggedright
Name
\end{minipage} & \begin{minipage}[b]{\linewidth}\raggedright
Description
\end{minipage} & \begin{minipage}[b]{\linewidth}\raggedright
Default
\end{minipage} \\
\midrule\noalign{}
\endhead
\bottomrule\noalign{}
\endlastfoot
Field & Metric & Laboratory & \\
1 & WPAVE1 & A real dimensionless value that defines the weighting
factor between the inner block and the surrounding blocks used in the
calculation of the connection factor weighted average pressure. If
WPAVE1 is greater than or equal to zero and less than or equal to one,
then the average pressure for each well connection is calculated based
on this weighting factor. A value of zero indicates only the surrounding
blocks should be used in the calculation; and a value of one indicates
only the inner blocks should be used. If WPAVE1 is less than zero, then
the average pressure for each well connection is weighted based on the
pore volumes of the inner and surrounding blocks. & 0.5 \\
2 & WPAVE2 & A real dimensionless value greater than or equal to zero
and less than or equal to one, that defines the weighting factor between
the connection factor weighted average pressures and the pore volume
weighted average pressures. If WPAVE2 is equal to one, then the average
pressures are calculated based only on the connection factor weighted
average pressures. If WPAVE2 is equal to zero, then average pressures
are calculated based only on the pore volumes weighted average
pressures. & 1.0 \\
3 & WPAVE3 & A defined character string that determines how the
hydrostatic head calculation is performed in correcting the pressures to
the BHP reference depth on the WELSPECS or WPAVEDEP keywords in the
SCHEDULE section. WPAVE3 should be set to one of the following character
strings: & WELL \\
4 & WPAVE4 & A defined character string that determines which
connections should be used in the calculations, WPAVE4 should be set to
one of the following character strings: Only the OPEN option is
currently supported by the simulator. & OPEN \\
Notes: & & & \\
\end{longtable}

\emph{Table 12.3.309.1: WPAVE Keyword Description}

The keyword is not applicable and should not be used with radial and
spider grid geometries.

See also the WELSPECS keyword that defines a well and a well's
bottom-hole pressure reference depth, the WPAVEDEP keyword that also
defines a well's bottom-hole pressure reference depth, and the COMPDAT
keyword to define a well's connections. All the aforementioned keywords
are described in the SCHEDULE section.

\subsubsection{Examples}\label{examples-93}

The following example defines the default well block average pressure
calculation parameters

\begin{verbatim}
--
--       DEFINE WELL BLOCK AVERAGE PRESSURE CALCULATION PARAMETERS
--
--       INNER   PORV   WELL  OPEN
--       OUTER   CONN   RES   ALL
WPAVE
         0.5     1.0    WELL  ALL                                              /
\end{verbatim}

And the next example shows the parameters used in the Norne model.

\begin{verbatim}
--
--       DEFINE WELL BLOCK AVERAGE PRESSURE CALCULATION PARAMETERS
--
--       INNER   PORV   WELL  OPEN
--       OUTER   CONN   RES   ALL
WPAVE
         1*      0.0    WELL  ALL                                              /
\end{verbatim}

Here only pore volume weighting is used instead of connection weighting.

\subsection{WPAVEDEP -- Define Well Reference Depth for Pressure
Calculations}\label{wpavedep-define-well-reference-depth-for-pressure-calculations}

\begin{longtable}[]{@{}
  >{\raggedright\arraybackslash}p{(\linewidth - 14\tabcolsep) * \real{0.1250}}
  >{\raggedright\arraybackslash}p{(\linewidth - 14\tabcolsep) * \real{0.1250}}
  >{\raggedright\arraybackslash}p{(\linewidth - 14\tabcolsep) * \real{0.1250}}
  >{\raggedright\arraybackslash}p{(\linewidth - 14\tabcolsep) * \real{0.1250}}
  >{\raggedright\arraybackslash}p{(\linewidth - 14\tabcolsep) * \real{0.1250}}
  >{\raggedright\arraybackslash}p{(\linewidth - 14\tabcolsep) * \real{0.1250}}
  >{\raggedright\arraybackslash}p{(\linewidth - 14\tabcolsep) * \real{0.1250}}
  >{\raggedright\arraybackslash}p{(\linewidth - 14\tabcolsep) * \real{0.1250}}@{}}
\toprule\noalign{}
\begin{minipage}[b]{\linewidth}\raggedright
RUNSPEC
\end{minipage} & \begin{minipage}[b]{\linewidth}\raggedright
GRID
\end{minipage} & \begin{minipage}[b]{\linewidth}\raggedright
EDIT
\end{minipage} & \begin{minipage}[b]{\linewidth}\raggedright
PROPS
\end{minipage} & \begin{minipage}[b]{\linewidth}\raggedright
REGIONS
\end{minipage} & \begin{minipage}[b]{\linewidth}\raggedright
SOLUTION
\end{minipage} & \begin{minipage}[b]{\linewidth}\raggedright
SUMMARY
\end{minipage} & \begin{minipage}[b]{\linewidth}\raggedright
SCHEDULE
\end{minipage} \\
\midrule\noalign{}
\endhead
\bottomrule\noalign{}
\endlastfoot
\end{longtable}

\subsubsection{Description}\label{description-1104}

The WPAVEDEP keyword defines the reference depth to be used to calculate
and report grid block average bottom-hole pressures for a well. This
keyword can be used to override the values entered or defaulted on the
WELSPECS keyword in the SCHEDULE section. The simulator corrects the
grid block calculated pressures to a well's reference depth using the
hydrostatic well of the producing fluids.

\begin{longtable}[]{@{}
  >{\raggedright\arraybackslash}p{(\linewidth - 6\tabcolsep) * \real{0.2500}}
  >{\raggedright\arraybackslash}p{(\linewidth - 6\tabcolsep) * \real{0.2500}}
  >{\raggedright\arraybackslash}p{(\linewidth - 6\tabcolsep) * \real{0.2500}}
  >{\raggedright\arraybackslash}p{(\linewidth - 6\tabcolsep) * \real{0.2500}}@{}}
\toprule\noalign{}
\begin{minipage}[b]{\linewidth}\raggedright
No.
\end{minipage} & \begin{minipage}[b]{\linewidth}\raggedright
Name
\end{minipage} & \begin{minipage}[b]{\linewidth}\raggedright
Description
\end{minipage} & \begin{minipage}[b]{\linewidth}\raggedright
Default
\end{minipage} \\
\midrule\noalign{}
\endhead
\bottomrule\noalign{}
\endlastfoot
Field & Metric & Laboratory & \\
1 & WELNAME & A character string of up to eight characters in length
that defines the well name for which the well and well connection status
data is being defined. Note that the well name (WELNAME) must have been
declared previously using the WELSPECS keyword in the SCHEDULE section,
otherwise an error may occur. & None \\
2 & BHPREF & A real value that defines the reference depth for reporting
the bottom-hole pressure for the well. Ideally this value should be set
to the midpoint of the perforations as defined by the COMPDAT keyword in
the SCHEDULE section. If defaulted by 1* or set to a value less than or
equal to zero, then the mid-point of shallowest connection defined by
the COMPDAT keyword will be used. & Mid-point of shallowest connection
defined by the COMPDAT keyword \\
feet & m & cm & \\
Notes: & & & \\
\end{longtable}

\emph{Table 12.109: WPAVDEP Keyword Description}

See also the WELSPECS keyword that defines a well, the COMPDAT keyword
to define a well's connections, and the WPAVE for defining how the
average bottom-hole pressure should be calculated. All the
aforementioned keywords are described in the SCHEDULE section.

\begin{tcolorbox}[enhanced jigsaw, title=\textcolor{quarto-callout-note-color}{\faInfo}\hspace{0.5em}{Note}, coltitle=black, breakable, toprule=.15mm, colbacktitle=quarto-callout-note-color!10!white, leftrule=.75mm, opacitybacktitle=0.6, colback=white, bottomtitle=1mm, toptitle=1mm, titlerule=0mm, left=2mm, colframe=quarto-callout-note-color-frame, arc=.35mm, rightrule=.15mm, bottomrule=.15mm, opacityback=0]

The keyword is normally used to reset a well's bottom-hole pressure
depth to match the pressure gauge depth for when observed pressure is
available, for example when conducting a history match for a well test,
or when attempting to match static bottom-hole surveys conducted on a
well.

\end{tcolorbox}

\subsubsection{Example}\label{example-594}

The following example illustrates how to set the bottom-hole reference
depth for wells completed in different reservoirs that have different
datum depths. Here it is assumed that all wells in a reservoir A have
RES-A as part of their well name, and similarly for reservoirs B and C.

\begin{verbatim}
--
-- WELL SPECIFICATION DATA
--
-- WELL     GROUP      LOCATION  BHP    PHASE  DRAIN  INFLOW  SHUT  CROSS  PRESS
-- NAME     NAME        I    J   DEPTH  FLUID  AREA   EQUA.   IN    FLOW   TABLE
WELSPECS
RES-AOP1  PLATFORM     14   13   1*      OIL   1*     STD     OPEN   NO    1*  /
RES-AOP2  PLATFORM     17   16   1*      OIL   1*     STD     OPEN   NO    1*  /
RES-AOP3  PLATFORM     21   19   1*      OIL   1*     STD     OPEN   NO    1*  /
RES-BOP4  PLATFORM     28   96   1*      OIL   1*     STD     OPEN   NO    1*  /
RES-BOP5  PLATFORM     34   89   1*      OIL   1*     STD     OPEN   NO    1*  /
RES-COP6  PLATFORM    128   52   1*      OIL   1*     STD     OPEN   NO    1*  /
RES-COP7  PLATFORM    134   56   1*      OIL   1*     STD     OPEN   NO    1*  /
RES-COP8  PLATFORM    138   50   1*      OIL   1*     STD     OPEN   NO    1*  /
RES-COP9  PLATFORM    120   52   1*      OIL   1*     STD     OPEN   NO    1*  /
/
--
--       DEFINE WELL REFERENCE DEPTH FOR PRESSURE CALCULATIONS
--
-- WELL  REF
-- NAME  DEPTH
-- ----  ------
WPAVEDEP
'RES-A*’ 3100.0                                                                 /
'RES-B*’ 3300.0                                                                 /
'RES-C*’ 5909.0                                                                 /
/
\end{verbatim}

In the example the all wells dedicated to RES-A will have their
bottom-hole reference depth set to 3,000 ft. TVDSS, RES-B wells to 3,300
ft. TVDSS and well RES-C wells to 5909 ft. TVDSS.

\subsection{WPIMULT -- Define Well Connection
Multipliers}\label{wpimult-define-well-connection-multipliers}

\begin{longtable}[]{@{}
  >{\raggedright\arraybackslash}p{(\linewidth - 14\tabcolsep) * \real{0.1250}}
  >{\raggedright\arraybackslash}p{(\linewidth - 14\tabcolsep) * \real{0.1250}}
  >{\raggedright\arraybackslash}p{(\linewidth - 14\tabcolsep) * \real{0.1250}}
  >{\raggedright\arraybackslash}p{(\linewidth - 14\tabcolsep) * \real{0.1250}}
  >{\raggedright\arraybackslash}p{(\linewidth - 14\tabcolsep) * \real{0.1250}}
  >{\raggedright\arraybackslash}p{(\linewidth - 14\tabcolsep) * \real{0.1250}}
  >{\raggedright\arraybackslash}p{(\linewidth - 14\tabcolsep) * \real{0.1250}}
  >{\raggedright\arraybackslash}p{(\linewidth - 14\tabcolsep) * \real{0.1250}}@{}}
\toprule\noalign{}
\begin{minipage}[b]{\linewidth}\raggedright
RUNSPEC
\end{minipage} & \begin{minipage}[b]{\linewidth}\raggedright
GRID
\end{minipage} & \begin{minipage}[b]{\linewidth}\raggedright
EDIT
\end{minipage} & \begin{minipage}[b]{\linewidth}\raggedright
PROPS
\end{minipage} & \begin{minipage}[b]{\linewidth}\raggedright
REGIONS
\end{minipage} & \begin{minipage}[b]{\linewidth}\raggedright
SOLUTION
\end{minipage} & \begin{minipage}[b]{\linewidth}\raggedright
SUMMARY
\end{minipage} & \begin{minipage}[b]{\linewidth}\raggedright
SCHEDULE
\end{minipage} \\
\midrule\noalign{}
\endhead
\bottomrule\noalign{}
\endlastfoot
\end{longtable}

\subsubsection{Description}\label{description-1105}

The WPIMULT keyword defines a well connection factor multiplier that
scales the existing well connection factor values. The resulting effect
is to scale the well's productivity at the reporting time step the
keyword is entered.

\begin{longtable}[]{@{}
  >{\raggedright\arraybackslash}p{(\linewidth - 6\tabcolsep) * \real{0.2500}}
  >{\raggedright\arraybackslash}p{(\linewidth - 6\tabcolsep) * \real{0.2500}}
  >{\raggedright\arraybackslash}p{(\linewidth - 6\tabcolsep) * \real{0.2500}}
  >{\raggedright\arraybackslash}p{(\linewidth - 6\tabcolsep) * \real{0.2500}}@{}}
\toprule\noalign{}
\begin{minipage}[b]{\linewidth}\raggedright
No.
\end{minipage} & \begin{minipage}[b]{\linewidth}\raggedright
Name
\end{minipage} & \begin{minipage}[b]{\linewidth}\raggedright
Description
\end{minipage} & \begin{minipage}[b]{\linewidth}\raggedright
Default
\end{minipage} \\
\midrule\noalign{}
\endhead
\bottomrule\noalign{}
\endlastfoot
Field & Metric & Laboratory & \\
1 & WELNAME & A character string of up to eight characters in length
that defines the well name for which the well and well connection status
data is being defined. Note that the well name (WELNAME) must have been
declared previously using the WELSPECS keyword in the SCHEDULE section,
otherwise an error may occur. & None \\
2 & PIMULT & A real positive value that will be used to scale the well
connection factors defined by I, J, K, C1 and C2 below. & 1.0 \\
3 & I & An integer less than or equal to NX that defines the connection
locations in the I-direction. & 1* \\
4 & J & An integer less than or equal to NY that defines the connection
locations in the J-direction. & 1* \\
5 & K & An integer less than or equal to NZ that defines the connection
locations in the K-direction. & 1* \\
6 & C1 & An integer value that defines the first completion number in
the range. Connections are lumped into completions via the COMPLUMP
keyword, and C1 refers to the first completion number, as defined by the
COMPLUMP keyword, and all the connections contained within the C1
completion. & 1* \\
7 & C2 & An integer that defines the last completion in the range.
Connections are lumped into completions via the COMPLUMP keyword, and C2
refers to the last completion number, as defined by the COMPLUMP
keyword, and all the connections contained within the C2 completion. &
1* \\
Notes: & & & \\
\end{longtable}

\emph{Table 12.110: WPIMULT Keyword Description}

If variables I, J, K, C1 and C2 are all negative values or defaulted
with 1*, then PIMULT is applied to all the well connections in the well.

If any of the variables I, J, K, C1 and C2 are set to zero (meaning any
or all values), or a positive value then PIMULT is applied to the
defined connections. The defined connections are those with the I, J, K
variables in the specified location and a completion number in the range
specified by C1 and C2.

Note, that PIMULT is applied at the time the WPIMULT keyword is entered
and is cumulative if applied to the same well connections, provided
there are intervening report time steps between consecutive WPIMULT
keywords. Consequently, if there are no intervening report time steps
between consecutive WPIMULT keywords utilizing the same well
connections, then only the last set is applied.

See also the WELPI keyword to set a well's productivity or injectivity
index at the time the keyword is activated, and also the PIMULTAB
keyword that defines productivity index multiplier versus water cut
tables that are used to scaled a well's connection factors based on a
wells connection current producing water cut. Both keywords are
documented in the SCHEDULE section.

\subsubsection{Example}\label{example-595}

The following example defines three vertical oil wells using the
WELSPECS keyword and their associated connection data.

\begin{verbatim}
--
-- WELL SPECIFICATION DATA
--
-- WELL     GROUP      LOCATION  BHP    PHASE  DRAIN  INFLOW  SHUT  CROSS  PRESS
-- NAME     NAME        I    J   DEPTH  FLUID  AREA   EQUA.   IN    FLOW   TABLE
WELSPECS
OP01      PLATFORM     14   13   1*      OIL   1*     STD     OPEN   NO    1*  /
OP02      PLATFORM     28   96   1*      OIL   1*     STD     OPEN   NO    1*  /
OP03      PLATFORM    128   56   1*      OIL   1*     STD     OPEN   NO    1*  /
/
--
--       WELL PRODUCTION WELL CONTROLS
--
-- WELL  OPEN/  CNTL   OIL    WAT    GAS   LIQ    RES    BHP    THP   VFP    VFP
-- NAME  SHUT   MODE   RATE   RATE   RATE  RATE   RATE   PRES   PRES  TABLE  ALFQ
WCONPROD
‘*’      SHUT   GRUP   1*     1*     1*    1*     1*     200.0                  /
/
--
--       WELL CONNECTION DATA
--
-- WELL  --- LOCATION ---  OPEN   SAT   CONN   WELL   KH    SKIN   D     DIR
-- NAME   II  JJ  K1  K2   SHUT   TAB   FACT   DIA    FACT  FACT   FACT  PEN
COMPDAT
OP01      1*  1*   1  10   OPEN   1*    1*    0.708   1*    0.0    1*    'Z' /
OP01      1*  1*  15  30   OPEN   1*    1*    0.708   1*    0.0    1*    'Z' /
OP01      1*  1*  35  90   OPEN   1*    1*    0.708   1*    0.0    1*    'Z' /
OP02      1*  1*   1  10   OPEN   1*    1*    0.708   1*    0.0    1*    'Z' /
OP03      1*  1*  35  90   OPEN   1*    1*    0.708   1*    0.0    1*    'Z' /
/
--       ASSIGN WELL CONNECTIONS TO COMPLETIONS
--
-- WELL  --- LOCATION ---  COMPL
-- NAME   II  JJ  K1  K2   NO.                                                                                                                                                                                                                                                                              COMPLUMP                                                                                                                                                                                                       OP03      1*  1*  35  45    1                              / COMPLETION NO. 01
OP03      1*  1*  50  90    2                              / COMPLETION NO. 02
/
--
--       DEFINE WELL CONNECTION MULTIPLIERS
--
-- WELL  PI     --LOCATION--  COMPLETION
-- NAME  MULT     I   J    K  FIRST LAST
WPIMULT
OP01     1.250    1*  1*   1*   1*    1*                  /
OP02     0.750    1*  1*   10   1*    1*                  /
OP03     1.100    1*  1*   1*   1     2                   /
/
\end{verbatim}

In this example the WPIMULT scales the well productivity of well OP01 by
1.25, and scales all the well connection factors in layer 10 only by
0.75 for well OP02. For well OP03, WPIMULT scales all the connections in
completions one and two by 1.1.

\subsection{WPIMULTL -- Define Well Connection Multipliers
(LGR)}\label{wpimultl-define-well-connection-multipliers-lgr}

\begin{longtable}[]{@{}
  >{\raggedright\arraybackslash}p{(\linewidth - 14\tabcolsep) * \real{0.1250}}
  >{\raggedright\arraybackslash}p{(\linewidth - 14\tabcolsep) * \real{0.1250}}
  >{\raggedright\arraybackslash}p{(\linewidth - 14\tabcolsep) * \real{0.1250}}
  >{\raggedright\arraybackslash}p{(\linewidth - 14\tabcolsep) * \real{0.1250}}
  >{\raggedright\arraybackslash}p{(\linewidth - 14\tabcolsep) * \real{0.1250}}
  >{\raggedright\arraybackslash}p{(\linewidth - 14\tabcolsep) * \real{0.1250}}
  >{\raggedright\arraybackslash}p{(\linewidth - 14\tabcolsep) * \real{0.1250}}
  >{\raggedright\arraybackslash}p{(\linewidth - 14\tabcolsep) * \real{0.1250}}@{}}
\toprule\noalign{}
\begin{minipage}[b]{\linewidth}\raggedright
RUNSPEC
\end{minipage} & \begin{minipage}[b]{\linewidth}\raggedright
GRID
\end{minipage} & \begin{minipage}[b]{\linewidth}\raggedright
EDIT
\end{minipage} & \begin{minipage}[b]{\linewidth}\raggedright
PROPS
\end{minipage} & \begin{minipage}[b]{\linewidth}\raggedright
REGIONS
\end{minipage} & \begin{minipage}[b]{\linewidth}\raggedright
SOLUTION
\end{minipage} & \begin{minipage}[b]{\linewidth}\raggedright
SUMMARY
\end{minipage} & \begin{minipage}[b]{\linewidth}\raggedright
SCHEDULE
\end{minipage} \\
\midrule\noalign{}
\endhead
\bottomrule\noalign{}
\endlastfoot
\end{longtable}

\subsubsection{Description}\label{description-1106}

The WPIMULTL keyword defines a well connection factor multiplier that
scales the existing well connection factor values, for a well in a Local
Grid Refinement (``LGR''). The resulting effect is scale the well's
productivity at the reporting time step the keyword is entered.

This keyword is not supported by OPM Flow but would change the results
if supported so the simulation will be stopped.

\begin{longtable}[]{@{}
  >{\raggedright\arraybackslash}p{(\linewidth - 6\tabcolsep) * \real{0.2500}}
  >{\raggedright\arraybackslash}p{(\linewidth - 6\tabcolsep) * \real{0.2500}}
  >{\raggedright\arraybackslash}p{(\linewidth - 6\tabcolsep) * \real{0.2500}}
  >{\raggedright\arraybackslash}p{(\linewidth - 6\tabcolsep) * \real{0.2500}}@{}}
\toprule\noalign{}
\begin{minipage}[b]{\linewidth}\raggedright
No.
\end{minipage} & \begin{minipage}[b]{\linewidth}\raggedright
Name
\end{minipage} & \begin{minipage}[b]{\linewidth}\raggedright
Description
\end{minipage} & \begin{minipage}[b]{\linewidth}\raggedright
Default
\end{minipage} \\
\midrule\noalign{}
\endhead
\bottomrule\noalign{}
\endlastfoot
Field & Metric & Laboratory & \\
1 & WELNAME & A character string of up to eight characters in length
that defines the well name for which the well and well connection status
data is being defined. Note that the well name (WELNAME) must have been
declared previously using the WELSPECL keyword in the SCHEDULE section,
otherwise an error may occur. & None \\
2 & PIMULT & A real positive value that will be used to scale the well
connection factors defined by I, J, K, C1 and C2 below. & 1.0 \\
3 & LGRNAME & A character string of up to eight characters in length
that defines the LGR name for which the well LGR connection multiplier
factor (PIMULT) is being defined. Note that LGRNAME must have been
declared previously using the WELSPECL keyword in the SCHEDULE section,
otherwise an error may occur. If defaulted with 1* the LGR on the
WELSPECL keyword will be utilized. & Defined \\
3 & I & An integer less than or equal to NX that defines the connection
location in the I-direction. & 1* \\
4 & J & An integer less than or equal to NY that defines the connection
location in the J-direction. & 1* \\
5 & K & An integer less than or equal to NZ that defines the connection
location in the K-direction. & 1* \\
6 & C1 & An integer value that defines the first completion in the
range. Connections are lumped into completions via the COMPLUMP keyword,
and C1 refers to the first completion number, as defined by the COMPLUMP
keyword, and all the connections contained within the C1 completion &
1* \\
7 & C2 & An integer value that defines the last completion in the range.
Connections are lumped into completions via the COMPLUMP keyword, and C2
refers to the last completion number, as defined by the COMPLUMP
keyword, and all the connections contained within the C2 completion &
1* \\
Notes: & & & \\
\end{longtable}

\emph{Table 12.111: WPIMULT Keyword Description}

If variables I, J K, C1 and C2 are all defaulted with a negative value
or 1*, then PIMULT is applied to all the well connections in the well.

If variables I, J K, C1 and C2 are set to zero (meaning any or all
values), or a positive value then PIMULT is applied to the defined
connections. The defined connections are those with the I, J, K
variables in the specified location and a completion number in the range
specified by C1 and C2.

Note, that PIMULT is applied at the time the WPIMULTL keyword is entered
and is cumulative if applied to the same well connections, provided
there are intervening report time steps between consecutive WPIMULTL
keywords. Consequently, if there are no intervening report time steps
between consecutive WPIMULTL keywords utilizing the same well
connections, then only the last set is applied.

See also the PIMULTAB keyword that defines productivity index multiplier
versus water cut tables that are used to scaled a well's connection
factors based on a wells connection current producing water cut. The
keyword is documented in the SCHEDULE section.

\subsubsection{Example}\label{example-596}

The following example defines two vertical oil wells using the WELSPECL
keyword and their associated connection data.

\begin{verbatim}
--
--       WELL LGR SPECIFICATION DATA
--
--       WELL GROUP  LGR    -LOCATION- BHP    PHASE DRAIN INFLOW SHUT CROSS PVT
--       NAME NAME   NAME     I    J   DEPTH  FLUID AREA  EQUA.  IN   FLOW  TABLE
WELSPECL
         OP01  PLAT OP01LGR  14   13   1*     OIL   1*    STD    SHUT  NO   1*  /
         OP02  PLAT OP02LGR  28   96   1*     OIL   1*    STD    SHUT  NO   1*  /
/
--
--       WELL LGR CONNECTION DATA
--
--       WELL  LGR   ---LOCATION---   OPEN   SAT CONN  WELL KH   SKIN  D    DIR
--       NAME  NAME   II  JJ  K1  K2  SHUT   TAB FACT  DIA  FACT FACT  FACT PEN
COMPDATL
         OP01 OP01LGR 1*  1*  20  56  OPEN   1*  1*  0.708   1*   1*   1*   Z /
         OP01 OP01LGR 1*  1*  75 100  SHUT   1*  1*  0.708   1*   1*   1*   Z /                  OP02 OP02LGR 1*  1*  75 100  OPEN   1*  1*  0.708   1*   1*   1*   Z /
         OP03 OP02LGR 1*  1*  75 100  OPEN   1*  1*  0.708   1*   1*   1*   Z /
/
--
--       ASSIGN WELL CONNECTIONS TO COMPLETIONS
--
--       WELL    LGR      --- LOCATION ---  COMPL
--       NAME    NAME      II  JJ  K1  K2   NO.                                                                                                                                                                                                                                                                              COMPLMPL                                                                                                                                                                                                                                OP03    OP02LGR   1*  1*  75  85    1          / COMPLETION NO. 01
             OP03    OP21LGR   1*  1*  86 100    2          / COMPLETION NO. 02
/
--
--       DEFINE WELL CONNECTION MULTIPLIERS
--
-- WELL  PI     LGR     --LOCATION--  COMPLETION
-- NAME  MULT   NAME      I   J    K  FIRST LAST
WPIMULTL
OP01     1.250  OP01LGR   1*  1*   1*    1*    1*          /
OP02     0.750  OP01LGR   1*  1*   10    1*    1*          /
OP03     1.100  OP02LGR   1*  1*   1*    1     2           /
/
\end{verbatim}

In this example the WPIMULTL keyword scales the well productivity of
well OP01 by 1.25, scales all the well connection factors in layer 10
only by 0.75 for well OP02, and for OP03, scales all the connections in
completions one and two by 1.100.

\subsection{WPITAB -- Assign Well Productivity Index versus Water Cut
Tables}\label{wpitab-assign-well-productivity-index-versus-water-cut-tables}

\begin{longtable}[]{@{}
  >{\raggedright\arraybackslash}p{(\linewidth - 14\tabcolsep) * \real{0.1250}}
  >{\raggedright\arraybackslash}p{(\linewidth - 14\tabcolsep) * \real{0.1250}}
  >{\raggedright\arraybackslash}p{(\linewidth - 14\tabcolsep) * \real{0.1250}}
  >{\raggedright\arraybackslash}p{(\linewidth - 14\tabcolsep) * \real{0.1250}}
  >{\raggedright\arraybackslash}p{(\linewidth - 14\tabcolsep) * \real{0.1250}}
  >{\raggedright\arraybackslash}p{(\linewidth - 14\tabcolsep) * \real{0.1250}}
  >{\raggedright\arraybackslash}p{(\linewidth - 14\tabcolsep) * \real{0.1250}}
  >{\raggedright\arraybackslash}p{(\linewidth - 14\tabcolsep) * \real{0.1250}}@{}}
\toprule\noalign{}
\begin{minipage}[b]{\linewidth}\raggedright
RUNSPEC
\end{minipage} & \begin{minipage}[b]{\linewidth}\raggedright
GRID
\end{minipage} & \begin{minipage}[b]{\linewidth}\raggedright
EDIT
\end{minipage} & \begin{minipage}[b]{\linewidth}\raggedright
PROPS
\end{minipage} & \begin{minipage}[b]{\linewidth}\raggedright
REGIONS
\end{minipage} & \begin{minipage}[b]{\linewidth}\raggedright
SOLUTION
\end{minipage} & \begin{minipage}[b]{\linewidth}\raggedright
SUMMARY
\end{minipage} & \begin{minipage}[b]{\linewidth}\raggedright
SCHEDULE
\end{minipage} \\
\midrule\noalign{}
\endhead
\bottomrule\noalign{}
\endlastfoot
\end{longtable}

\subsubsection{Description}\label{description-1107}

The WPITAB keyword assigns the well productivity index multiplier versus
water cut tables, that are used to scaled a well's connection factors
based on the connection's current producing water cut, to a well. The
tables are defined via the PIMULTAB keyword in the SCHEDULE section.

This keyword is not supported by OPM Flow but would change the results
if supported so the simulation will be stopped.

\begin{longtable}[]{@{}
  >{\raggedright\arraybackslash}p{(\linewidth - 6\tabcolsep) * \real{0.2500}}
  >{\raggedright\arraybackslash}p{(\linewidth - 6\tabcolsep) * \real{0.2500}}
  >{\raggedright\arraybackslash}p{(\linewidth - 6\tabcolsep) * \real{0.2500}}
  >{\raggedright\arraybackslash}p{(\linewidth - 6\tabcolsep) * \real{0.2500}}@{}}
\toprule\noalign{}
\begin{minipage}[b]{\linewidth}\raggedright
No.
\end{minipage} & \begin{minipage}[b]{\linewidth}\raggedright
Name
\end{minipage} & \begin{minipage}[b]{\linewidth}\raggedright
Description
\end{minipage} & \begin{minipage}[b]{\linewidth}\raggedright
Default
\end{minipage} \\
\midrule\noalign{}
\endhead
\bottomrule\noalign{}
\endlastfoot
Field & Metric & Laboratory & \\
1 & WELNAME & A character string of up to eight characters in length
that defines the well name for which the well productivity index
multiplier versus water cut table, PIMULTAB, is being assigned. Note
that the well name (WELNAME) must have been declared previously using
the WELSPECS keyword in the SCHEDULE section, otherwise an error may
occur. & None \\
2 & PIMULTAB & A positive integer value that defines the corresponding
PIMULTAB table to be allocated to the well. A value less than or equal
to zero means that no PIMULTAB table is allocated to the well & 0 \\
Notes: & & & \\
\end{longtable}

\emph{Table 12.112: WPITAB Keyword Description}

See also the PIMULTAB keyword that defines productivity index multiplier
versus water cut tables and also the WPIMULT keyword that scales a
well's productivity index by a constant value, both of which are in the
SCHEDULE section.

\subsubsection{Example}\label{example-597}

Given NTPIMT equals two on the PIMTDIMS keyword in the RUNSPEC section,
then:

\begin{verbatim}
--
--       ASSIGN WELL PRODUCTIVITY INDEX VS WATER CUT TABLE
--
-- WELL  PI
-- NAME  TABLE
WPITAB
OP01     1                                                        /
OP02     1                                                        /
OP03     2                                                        /
/
\end{verbatim}

Assigns PIMULTAB table one to wells OP01 and OP02 and table two to OP03.

\subsection{WPLUG -- Define Well Plug Back
Length}\label{wplug-define-well-plug-back-length}

\begin{longtable}[]{@{}
  >{\raggedright\arraybackslash}p{(\linewidth - 14\tabcolsep) * \real{0.1250}}
  >{\raggedright\arraybackslash}p{(\linewidth - 14\tabcolsep) * \real{0.1250}}
  >{\raggedright\arraybackslash}p{(\linewidth - 14\tabcolsep) * \real{0.1250}}
  >{\raggedright\arraybackslash}p{(\linewidth - 14\tabcolsep) * \real{0.1250}}
  >{\raggedright\arraybackslash}p{(\linewidth - 14\tabcolsep) * \real{0.1250}}
  >{\raggedright\arraybackslash}p{(\linewidth - 14\tabcolsep) * \real{0.1250}}
  >{\raggedright\arraybackslash}p{(\linewidth - 14\tabcolsep) * \real{0.1250}}
  >{\raggedright\arraybackslash}p{(\linewidth - 14\tabcolsep) * \real{0.1250}}@{}}
\toprule\noalign{}
\begin{minipage}[b]{\linewidth}\raggedright
RUNSPEC
\end{minipage} & \begin{minipage}[b]{\linewidth}\raggedright
GRID
\end{minipage} & \begin{minipage}[b]{\linewidth}\raggedright
EDIT
\end{minipage} & \begin{minipage}[b]{\linewidth}\raggedright
PROPS
\end{minipage} & \begin{minipage}[b]{\linewidth}\raggedright
REGIONS
\end{minipage} & \begin{minipage}[b]{\linewidth}\raggedright
SOLUTION
\end{minipage} & \begin{minipage}[b]{\linewidth}\raggedright
SUMMARY
\end{minipage} & \begin{minipage}[b]{\linewidth}\raggedright
SCHEDULE
\end{minipage} \\
\midrule\noalign{}
\endhead
\bottomrule\noalign{}
\endlastfoot
\end{longtable}

\subsubsection{Description}\label{description-1108}

Various keywords in the SCHEDULE section (WECON, GECON etc.) allow for a
well to be automatically plugged back if the well violates a constraint,
that is to close existing perforations (well connections). For example
if the water cut exceeds 90\%, then plug back the well. The WPLUG
keyword defines for automatic plug backs the length of the perforations
(length of connections) to be closed each time an automatic plug back is
performed, together with various options on how the workover should be
performed, top down, bottom up, etc.

This keyword is not supported by OPM Flow but would change the results
if supported so the simulation will be stopped.

\subsection{WPMITAB -- Assign Well Polymer Molecular Model Injection
Tables}\label{wpmitab-assign-well-polymer-molecular-model-injection-tables}

\begin{longtable}[]{@{}
  >{\raggedright\arraybackslash}p{(\linewidth - 14\tabcolsep) * \real{0.1250}}
  >{\raggedright\arraybackslash}p{(\linewidth - 14\tabcolsep) * \real{0.1250}}
  >{\raggedright\arraybackslash}p{(\linewidth - 14\tabcolsep) * \real{0.1250}}
  >{\raggedright\arraybackslash}p{(\linewidth - 14\tabcolsep) * \real{0.1250}}
  >{\raggedright\arraybackslash}p{(\linewidth - 14\tabcolsep) * \real{0.1250}}
  >{\raggedright\arraybackslash}p{(\linewidth - 14\tabcolsep) * \real{0.1250}}
  >{\raggedright\arraybackslash}p{(\linewidth - 14\tabcolsep) * \real{0.1250}}
  >{\raggedright\arraybackslash}p{(\linewidth - 14\tabcolsep) * \real{0.1250}}@{}}
\toprule\noalign{}
\begin{minipage}[b]{\linewidth}\raggedright
RUNSPEC
\end{minipage} & \begin{minipage}[b]{\linewidth}\raggedright
GRID
\end{minipage} & \begin{minipage}[b]{\linewidth}\raggedright
EDIT
\end{minipage} & \begin{minipage}[b]{\linewidth}\raggedright
PROPS
\end{minipage} & \begin{minipage}[b]{\linewidth}\raggedright
REGIONS
\end{minipage} & \begin{minipage}[b]{\linewidth}\raggedright
SOLUTION
\end{minipage} & \begin{minipage}[b]{\linewidth}\raggedright
SUMMARY
\end{minipage} & \begin{minipage}[b]{\linewidth}\raggedright
SCHEDULE
\end{minipage} \\
\midrule\noalign{}
\endhead
\bottomrule\noalign{}
\endlastfoot
\end{longtable}

\subsubsection{Description}\label{description-1109}

The WPMITAB keyword assigns the well polymer molecular injection tables
to water injection wells in OPM Flow's Polymer Molecular Weight
Transport option, that uses the polymer molecular weight in calculating
the polymer viscosity, as well as accounting for formation damage due to
the water and polymer injection, by adjusting the wellbore skin
pressure. This keyword should only be used if the POLYMER and POLYMW
keywords in the RUNSPEC section are also activated. The keyword assigns
the PLYMWINJ tables that are defined via the PLYMWINJ keyword in the
PROPS section.

\begin{tcolorbox}[enhanced jigsaw, title=\textcolor{quarto-callout-note-color}{\faInfo}\hspace{0.5em}{Note}, coltitle=black, breakable, toprule=.15mm, colbacktitle=quarto-callout-note-color!10!white, leftrule=.75mm, opacitybacktitle=0.6, colback=white, bottomtitle=1mm, toptitle=1mm, titlerule=0mm, left=2mm, colframe=quarto-callout-note-color-frame, arc=.35mm, rightrule=.15mm, bottomrule=.15mm, opacityback=0]

This is an OPM Flow specific keyword that employs an alternative polymer
flood model based on a Polymer Molecular Weight Transport equation, that
is not available in the commercial simulator.

\end{tcolorbox}

\begin{longtable}[]{@{}
  >{\raggedright\arraybackslash}p{(\linewidth - 6\tabcolsep) * \real{0.2500}}
  >{\raggedright\arraybackslash}p{(\linewidth - 6\tabcolsep) * \real{0.2500}}
  >{\raggedright\arraybackslash}p{(\linewidth - 6\tabcolsep) * \real{0.2500}}
  >{\raggedright\arraybackslash}p{(\linewidth - 6\tabcolsep) * \real{0.2500}}@{}}
\toprule\noalign{}
\begin{minipage}[b]{\linewidth}\raggedright
No.
\end{minipage} & \begin{minipage}[b]{\linewidth}\raggedright
Name
\end{minipage} & \begin{minipage}[b]{\linewidth}\raggedright
Description
\end{minipage} & \begin{minipage}[b]{\linewidth}\raggedright
Default
\end{minipage} \\
\midrule\noalign{}
\endhead
\bottomrule\noalign{}
\endlastfoot
Field & Metric & Laboratory & \\
1 & WELNAME & A character string of up to eight characters in length,
that defines the water injection well name, for which the well polymer
molecular injection table, PLYMWINJ, is to be assigned. Note that the
well name (WELNAME) must have been declared previously using the
WELSPECS keyword in the SCHEDULE section, otherwise an error may occur.
& None \\
2 & PLYMWINJ & A positive integer value that defines the corresponding
PLYMWINJ table to be allocated to the water injection well. A value less
than or equal to zero means that no PLYMWIN table is allocated to the
well & 0 \\
Notes: & & & \\
\end{longtable}

\emph{Table 12.113: WPMITAB Keyword Description}

See also the PLYMWINJ keyword in the PROPS section, that describes the
relationship of the injected polymer molecular weight as a function of
polymer throughput and polymer velocity, for the simulator's Polymer
Molecular Weight Transport option, as well as the SKPRWAT, SKPRPOLY, and
PLYVMH keywords, also in the PROPS section, that are the additional
keywords required for the Polymer Molecular Weight Transport option.

The WSKPTAB keyword in the SCHEDULE section may be used to assign the
SKPRWAT and SKPRPOL tables to water injections wells, that enable the
calculation of the wellbore skin pressure based on the fluids being
injected.

Note that the standard polymer property data keywords: PLYROCK, PLYADS,
PLYMAX, etc., are still required to fully describe the polymer fluid.

\subsubsection{Example}\label{example-598}

Given NTPMWINJ equals two on the PINTDIMS keyword in the RUNSPEC
section, then:

\begin{verbatim}
--
--       ASSIGN WELL POLYMER MOLECULAR MODEL INJECTION TABLES
--
-- WELL  PLYMWINJ
-- NAME  TABLE
WPMITAB
WI01     1                                                        /
WI02     1                                                        /
WI03     2                                                        /
/
\end{verbatim}

Assigns PLYMWINJ table one to wells WI01 and WI02 and table two to WI03.

\subsection{WPOLYMER -- Define Water Injection Well Polymer and Salt
Concentrations}\label{wpolymer-define-water-injection-well-polymer-and-salt-concentrations}

\begin{longtable}[]{@{}
  >{\raggedright\arraybackslash}p{(\linewidth - 14\tabcolsep) * \real{0.1250}}
  >{\raggedright\arraybackslash}p{(\linewidth - 14\tabcolsep) * \real{0.1250}}
  >{\raggedright\arraybackslash}p{(\linewidth - 14\tabcolsep) * \real{0.1250}}
  >{\raggedright\arraybackslash}p{(\linewidth - 14\tabcolsep) * \real{0.1250}}
  >{\raggedright\arraybackslash}p{(\linewidth - 14\tabcolsep) * \real{0.1250}}
  >{\raggedright\arraybackslash}p{(\linewidth - 14\tabcolsep) * \real{0.1250}}
  >{\raggedright\arraybackslash}p{(\linewidth - 14\tabcolsep) * \real{0.1250}}
  >{\raggedright\arraybackslash}p{(\linewidth - 14\tabcolsep) * \real{0.1250}}@{}}
\toprule\noalign{}
\begin{minipage}[b]{\linewidth}\raggedright
RUNSPEC
\end{minipage} & \begin{minipage}[b]{\linewidth}\raggedright
GRID
\end{minipage} & \begin{minipage}[b]{\linewidth}\raggedright
EDIT
\end{minipage} & \begin{minipage}[b]{\linewidth}\raggedright
PROPS
\end{minipage} & \begin{minipage}[b]{\linewidth}\raggedright
REGIONS
\end{minipage} & \begin{minipage}[b]{\linewidth}\raggedright
SOLUTION
\end{minipage} & \begin{minipage}[b]{\linewidth}\raggedright
SUMMARY
\end{minipage} & \begin{minipage}[b]{\linewidth}\raggedright
SCHEDULE
\end{minipage} \\
\midrule\noalign{}
\endhead
\bottomrule\noalign{}
\endlastfoot
\end{longtable}

\subsubsection{Description}\label{description-1110}

The WPOLYMER keyword defines a water injection well's polymer and salt
injection stream concentrations that are to be used for when the polymer
and salt options have been activated by the POLYMER and BRINE keywords
in the RUNSPEC section.

Note that if the Brine option has not be activated by the BRINE keyword
in the RUNSPEC section, then the salt concentrations in the third column
are ignored. Secondly, if the brine phase is declared but the polymer
phase has not been made active, then the WSALT keyword in the SCHEDULE
section can be used to set the salt concentration.

Currently the Brine option is not implemented in OPM Flow and therefore
both the SALTCON and GRPSALT variables on this keyword are ignored.

\begin{longtable}[]{@{}
  >{\raggedright\arraybackslash}p{(\linewidth - 6\tabcolsep) * \real{0.2500}}
  >{\raggedright\arraybackslash}p{(\linewidth - 6\tabcolsep) * \real{0.2500}}
  >{\raggedright\arraybackslash}p{(\linewidth - 6\tabcolsep) * \real{0.2500}}
  >{\raggedright\arraybackslash}p{(\linewidth - 6\tabcolsep) * \real{0.2500}}@{}}
\toprule\noalign{}
\begin{minipage}[b]{\linewidth}\raggedright
No.
\end{minipage} & \begin{minipage}[b]{\linewidth}\raggedright
Name
\end{minipage} & \begin{minipage}[b]{\linewidth}\raggedright
Description
\end{minipage} & \begin{minipage}[b]{\linewidth}\raggedright
Default
\end{minipage} \\
\midrule\noalign{}
\endhead
\bottomrule\noalign{}
\endlastfoot
Field & Metric & Laboratory & \\
1 & WELNAME & A character string of up to eight characters in length
that defines the well name for which the well connection data is being
defined. Note that the well name (WELNAME) must have been declared
previously using the WELSPECS keyword in the SCHEDULE section, otherwise
an error may occur. & None \\
2 & POLCON & A real positive value that defines the polymer
concentration of the well's injection stream. This value may be
specified using a User Defined Argument (UDA). & None \\
lb/stb & kg/sm3 & gm/scc & \\
3 & SALTCON & A real positive value that defines the salt concentration
of the well's injection stream. This value may be specified using a User
Defined Argument (UDA). This keyword is not supported by OPM Flow but
has no effect on the results so it will be ignored. & None \\
lb/stb & kg/sm3 & gm/scc & \\
4 & GRPPOL & A character string of up to eight characters in length that
defines the group name for which the group's produced polymer
concentration should be used instead of the well's POLCON value stated
on this keyword. & None \\
5 & GRPSALT & A character string of up to eight characters in length
that defines the group name for which the group's produced salt
concentration should be used instead of the well's SALTCON value stated
on this keyword. This keyword is not supported by OPM Flow but has no
effect on the results so it will be ignored. & None \\
Notes: & & & \\
\end{longtable}

\emph{Table 12.114: WPOLYMER Keyword Description}

Water injection wells that are not declared via this keyword have their
concentrations defaulted to zero.

See also the GCONPROD and GCONINJE keywords to define a group's
production and injection targets and constraints, and the WCONINJE
keyword to define an injection well's targets and constraints. All the
aforementioned keywords are described in the SCHEDULE section.

\subsubsection{Example}\label{example-599}

The following example defines the polymer and salt injection stream
concentrations for three water injection wells for when the polymer
option has been activated by the POLYMER keyword in the RUNSPEC section.

\begin{verbatim}
--
--       DEFINE WATER INJECTION WELL POLYMER AND SALT CONCENTRATIONS
--
-- WELL  POLYMER    SALT       POLYMER    SALT
-- NAME  POLCON     SALTCON    GROUP      GROUP
--       -------    --------   --------   --------
WPOLYMER
WI01     0.2500                                            /
WI02     1*         1*         GRPINJ1                     /
WI03     0.2500     1*         GRPINJ1                     /
/
\end{verbatim}

The polymer concentration for well WI01 is set to 0.25 and the stated
polymer concentration for well WI02 will be ignored, as both WI02 and
WI03 will re-inject the produced polymer from the GRPINJ1 group.

\subsection{WPOLYRED -- Define Well Polymer-Water Viscosity Reduction
Factor}\label{wpolyred-define-well-polymer-water-viscosity-reduction-factor}

\begin{longtable}[]{@{}
  >{\raggedright\arraybackslash}p{(\linewidth - 14\tabcolsep) * \real{0.1250}}
  >{\raggedright\arraybackslash}p{(\linewidth - 14\tabcolsep) * \real{0.1250}}
  >{\raggedright\arraybackslash}p{(\linewidth - 14\tabcolsep) * \real{0.1250}}
  >{\raggedright\arraybackslash}p{(\linewidth - 14\tabcolsep) * \real{0.1250}}
  >{\raggedright\arraybackslash}p{(\linewidth - 14\tabcolsep) * \real{0.1250}}
  >{\raggedright\arraybackslash}p{(\linewidth - 14\tabcolsep) * \real{0.1250}}
  >{\raggedright\arraybackslash}p{(\linewidth - 14\tabcolsep) * \real{0.1250}}
  >{\raggedright\arraybackslash}p{(\linewidth - 14\tabcolsep) * \real{0.1250}}@{}}
\toprule\noalign{}
\begin{minipage}[b]{\linewidth}\raggedright
RUNSPEC
\end{minipage} & \begin{minipage}[b]{\linewidth}\raggedright
GRID
\end{minipage} & \begin{minipage}[b]{\linewidth}\raggedright
EDIT
\end{minipage} & \begin{minipage}[b]{\linewidth}\raggedright
PROPS
\end{minipage} & \begin{minipage}[b]{\linewidth}\raggedright
REGIONS
\end{minipage} & \begin{minipage}[b]{\linewidth}\raggedright
SOLUTION
\end{minipage} & \begin{minipage}[b]{\linewidth}\raggedright
SUMMARY
\end{minipage} & \begin{minipage}[b]{\linewidth}\raggedright
SCHEDULE
\end{minipage} \\
\midrule\noalign{}
\endhead
\bottomrule\noalign{}
\endlastfoot
\end{longtable}

\subsubsection{Description}\label{description-1111}

The WPOLYRED keyword defines the polymer-water reduction factor for
injection wells, for when the polymer phase has been activated by the
POLYMER keyword in the RUNSPEC section. WPOLYRED should be set to a
value greater than or equal to zero and less than or equal to one that
determines the injection mixture's viscosity. A value of zero indicates
for pure water injection and a value of one will use the simulator's
valuated mixture viscosity. A value between zero and one will use an
interpolated mixture viscosity.

This keyword is not supported by OPM Flow but would change the results
if supported so the simulation will be stopped.

\subsection{WREGROUP -- Automatic Re-Assignment of Wells to
Groups}\label{wregroup-automatic-re-assignment-of-wells-to-groups}

\begin{longtable}[]{@{}
  >{\raggedright\arraybackslash}p{(\linewidth - 14\tabcolsep) * \real{0.1250}}
  >{\raggedright\arraybackslash}p{(\linewidth - 14\tabcolsep) * \real{0.1250}}
  >{\raggedright\arraybackslash}p{(\linewidth - 14\tabcolsep) * \real{0.1250}}
  >{\raggedright\arraybackslash}p{(\linewidth - 14\tabcolsep) * \real{0.1250}}
  >{\raggedright\arraybackslash}p{(\linewidth - 14\tabcolsep) * \real{0.1250}}
  >{\raggedright\arraybackslash}p{(\linewidth - 14\tabcolsep) * \real{0.1250}}
  >{\raggedright\arraybackslash}p{(\linewidth - 14\tabcolsep) * \real{0.1250}}
  >{\raggedright\arraybackslash}p{(\linewidth - 14\tabcolsep) * \real{0.1250}}@{}}
\toprule\noalign{}
\begin{minipage}[b]{\linewidth}\raggedright
RUNSPEC
\end{minipage} & \begin{minipage}[b]{\linewidth}\raggedright
GRID
\end{minipage} & \begin{minipage}[b]{\linewidth}\raggedright
EDIT
\end{minipage} & \begin{minipage}[b]{\linewidth}\raggedright
PROPS
\end{minipage} & \begin{minipage}[b]{\linewidth}\raggedright
REGIONS
\end{minipage} & \begin{minipage}[b]{\linewidth}\raggedright
SOLUTION
\end{minipage} & \begin{minipage}[b]{\linewidth}\raggedright
SUMMARY
\end{minipage} & \begin{minipage}[b]{\linewidth}\raggedright
SCHEDULE
\end{minipage} \\
\midrule\noalign{}
\endhead
\bottomrule\noalign{}
\endlastfoot
\end{longtable}

\subsubsection{Description}\label{description-1112}

WREGROUP defines the criteria to automatically re-assign wells to
various other groups. This can be used, for example, to move wells on
THP control flowing through a high pressure separator group to a low
pressure separator group in order for the wells to be under different
group controls for low pressure wells.

This keyword is not supported by OPM Flow but would change the results
if supported so the simulation will be stopped.

\subsection{WRFT -- Activate Well RFT Reporting to the RFT
File}\label{wrft-activate-well-rft-reporting-to-the-rft-file}

\begin{longtable}[]{@{}
  >{\raggedright\arraybackslash}p{(\linewidth - 14\tabcolsep) * \real{0.1250}}
  >{\raggedright\arraybackslash}p{(\linewidth - 14\tabcolsep) * \real{0.1250}}
  >{\raggedright\arraybackslash}p{(\linewidth - 14\tabcolsep) * \real{0.1250}}
  >{\raggedright\arraybackslash}p{(\linewidth - 14\tabcolsep) * \real{0.1250}}
  >{\raggedright\arraybackslash}p{(\linewidth - 14\tabcolsep) * \real{0.1250}}
  >{\raggedright\arraybackslash}p{(\linewidth - 14\tabcolsep) * \real{0.1250}}
  >{\raggedright\arraybackslash}p{(\linewidth - 14\tabcolsep) * \real{0.1250}}
  >{\raggedright\arraybackslash}p{(\linewidth - 14\tabcolsep) * \real{0.1250}}@{}}
\toprule\noalign{}
\begin{minipage}[b]{\linewidth}\raggedright
RUNSPEC
\end{minipage} & \begin{minipage}[b]{\linewidth}\raggedright
GRID
\end{minipage} & \begin{minipage}[b]{\linewidth}\raggedright
EDIT
\end{minipage} & \begin{minipage}[b]{\linewidth}\raggedright
PROPS
\end{minipage} & \begin{minipage}[b]{\linewidth}\raggedright
REGIONS
\end{minipage} & \begin{minipage}[b]{\linewidth}\raggedright
SOLUTION
\end{minipage} & \begin{minipage}[b]{\linewidth}\raggedright
SUMMARY
\end{minipage} & \begin{minipage}[b]{\linewidth}\raggedright
SCHEDULE
\end{minipage} \\
\midrule\noalign{}
\endhead
\bottomrule\noalign{}
\endlastfoot
\end{longtable}

\subsubsection{Description}\label{description-1113}

This keyword activates reporting of a well's pressure and saturation
profile versus depth for the connected grid blocks, to the RFT file for
the requested wells at the time the keyword is activated. Data written
out by OPM Flow is used to match the field measured data collected from
a Repeat Formation Tester (``RFT'') tool.

The RFT tool is an open hole device which is an updated version of the
Formation Interval Tester (``FIT''), both of which are run on wire line.
Both tools take multiple pressure readings (at various depths) thus
enabling a pressure depth profile to be obtained from the formation,
and, in addition, they can also take fluid samples from the formation.
The latest tool available from Schlumberger is the Modular Formation
Dynamics Tester (``MDT''), which, as its name suggests, is a modular
tool that can be assembled in different configurations depending on what
are the objectives for running the tool. Note other vendors have similar
wire line logging tools with alternative names for the tools. Throughout
this section the term RFT applies to all tools that measure a pressure
profile versus depth (RFT/FIT/MDT etc.).

\begin{longtable}[]{@{}
  >{\raggedright\arraybackslash}p{(\linewidth - 6\tabcolsep) * \real{0.2500}}
  >{\raggedright\arraybackslash}p{(\linewidth - 6\tabcolsep) * \real{0.2500}}
  >{\raggedright\arraybackslash}p{(\linewidth - 6\tabcolsep) * \real{0.2500}}
  >{\raggedright\arraybackslash}p{(\linewidth - 6\tabcolsep) * \real{0.2500}}@{}}
\toprule\noalign{}
\begin{minipage}[b]{\linewidth}\raggedright
No.
\end{minipage} & \begin{minipage}[b]{\linewidth}\raggedright
Name
\end{minipage} & \begin{minipage}[b]{\linewidth}\raggedright
Description
\end{minipage} & \begin{minipage}[b]{\linewidth}\raggedright
Default
\end{minipage} \\
\midrule\noalign{}
\endhead
\bottomrule\noalign{}
\endlastfoot
Field & Metric & Laboratory & \\
1 & WELNAME & A columnar vector of character strings of up to eight
characters in length for each item, that defines the well name for which
the RFT data should be written to the RFT file. Note that the WELNAME
must have been declared previously using the WELSPECS keyword in the
SCHEDULE section, otherwise an error may occur. If the WELNAME is left
blank then the data is written out for all wells at the time they are
first opened to flow. If the WELNAME is given, then the RFT data for the
well at the time step the keyword is invoked is written out. & None \\
Notes: & & & \\
\end{longtable}

\emph{Table 12.115: WRFT Keyword Description}

See also the WRFTPLT keyword in the SCHEDULE section that has more
flexible reporting options.

\subsubsection{Examples}\label{examples-94}

The first example activates RFT reporting for all wells at the time a
well is first opened to flow:

\begin{verbatim}
--
--       ACTIVATE WELL RFT REPORTING TO THE RFT FILE
--
-- WELL
-- NAME
WRFT
/
\end{verbatim}

Ideally, this version of the keyword should be place at the beginning of
the SCHEDULE section to obtain the data for the wells in the run before
they are opened up through time.

The next example shows how to use the keyword to request the output for
several wells at different reporting time steps.

\begin{verbatim}
-- ------------------------------------------------------------------------------
-- 01 JAN 2000 START OF SCHEDULE SECTION
-- ------------------------------------------------------------------------------
DATES
15 JAN 2000 /
/
--
--       WELL HISTORICAL PRODUCTION CONTROLS
--
-- WELL  OPEN/  CNTL   OIL    WAT    GAS    VFP    VFP   THP   BHP
-- NAME  SHUT   MODE   RATE   RATE   RATE   TABLE  ALFQ  PRES  PRES
WCONHIST                                                                                                                          OP01     OPEN   ORAT  15.5E3  0.0    1550   10      1*   900.0 1*     /
OP02     SHUT                                                         /
/
--
--       ACTIVATE WELL RFT REPORTING TO THE RFT FILE
--
-- WELL
-- NAME
WRFT
OP01                                                                  /
OP02                                                                  /
/
DATES
01 FEB 2000 /
/
--
--       WELL HISTORICAL PRODUCTION CONTROLS
--
-- WELL  OPEN/  CNTL   OIL    WAT    GAS    VFP    VFP   THP   BHP
-- NAME  SHUT   MODE   RATE   RATE   RATE   TABLE  ALFQ  PRES  PRES
WCONHIST                                                                                                                          OP01     OPEN   ORAT  15.5E3  0.0    1550   10      1*   900.0 1*     /
OP02     SHUT                                                         /
/                                                                               --
--       ACTIVATE WELL RFT REPORTING TO THE RFT FILE
--
-- WELL
-- NAME
WRFT
OP01                                                                  /
OP02                                                                  /
/
DATES
01 MAR 2000 /
/
--
--       WELL HISTORICAL PRODUCTION CONTROLS
--
-- WELL  OPEN/  CNTL   OIL    WAT    GAS    VFP    VFP   THP   BHP
-- NAME  SHUT   MODE   RATE   RATE   RATE   TABLE  ALFQ  PRES  PRES
WCONHIST                                                                                                                          OP01     OPEN   ORAT  15.5E3  0.0    1550   10      1*   900.0 1*     /
OP02     OPEN   ORAT  10.5E3  0.0    1000   10      1*   900.0 1*     /
/
\end{verbatim}

In this example, both well's have their RFT written out on February 1
and March 1 2000.

\subsection{WRFTPLT -- Activate Well RFT and PLT Reporting to the RFT
File}\label{wrftplt-activate-well-rft-and-plt-reporting-to-the-rft-file}

\begin{longtable}[]{@{}
  >{\raggedright\arraybackslash}p{(\linewidth - 14\tabcolsep) * \real{0.1250}}
  >{\raggedright\arraybackslash}p{(\linewidth - 14\tabcolsep) * \real{0.1250}}
  >{\raggedright\arraybackslash}p{(\linewidth - 14\tabcolsep) * \real{0.1250}}
  >{\raggedright\arraybackslash}p{(\linewidth - 14\tabcolsep) * \real{0.1250}}
  >{\raggedright\arraybackslash}p{(\linewidth - 14\tabcolsep) * \real{0.1250}}
  >{\raggedright\arraybackslash}p{(\linewidth - 14\tabcolsep) * \real{0.1250}}
  >{\raggedright\arraybackslash}p{(\linewidth - 14\tabcolsep) * \real{0.1250}}
  >{\raggedright\arraybackslash}p{(\linewidth - 14\tabcolsep) * \real{0.1250}}@{}}
\toprule\noalign{}
\begin{minipage}[b]{\linewidth}\raggedright
RUNSPEC
\end{minipage} & \begin{minipage}[b]{\linewidth}\raggedright
GRID
\end{minipage} & \begin{minipage}[b]{\linewidth}\raggedright
EDIT
\end{minipage} & \begin{minipage}[b]{\linewidth}\raggedright
PROPS
\end{minipage} & \begin{minipage}[b]{\linewidth}\raggedright
REGIONS
\end{minipage} & \begin{minipage}[b]{\linewidth}\raggedright
SOLUTION
\end{minipage} & \begin{minipage}[b]{\linewidth}\raggedright
SUMMARY
\end{minipage} & \begin{minipage}[b]{\linewidth}\raggedright
SCHEDULE
\end{minipage} \\
\midrule\noalign{}
\endhead
\bottomrule\noalign{}
\endlastfoot
\end{longtable}

\subsubsection{Description}\label{description-1114}

This keyword activates reporting of a well's depth pressure and fluid
rates profile to the RFT file for the requested wells at the time the
keyword is activated. Data written out by the simulator is used to match
the field measured data collected from both the Repeat Formation Tester
(``RFT'') tool and various Production Logging Tools (``PLT'').

See the WRFT keyword in the SCHEDULE section for a brief description of
the RFT data set. This keyword also actives the writing out of each well
connection's fluid rates, connection factors and KH data, etc., as the
PLT data. The PLT data is used to compare with measured data from wire
line production logging tools.

\begin{longtable}[]{@{}
  >{\raggedright\arraybackslash}p{(\linewidth - 6\tabcolsep) * \real{0.2500}}
  >{\raggedright\arraybackslash}p{(\linewidth - 6\tabcolsep) * \real{0.2500}}
  >{\raggedright\arraybackslash}p{(\linewidth - 6\tabcolsep) * \real{0.2500}}
  >{\raggedright\arraybackslash}p{(\linewidth - 6\tabcolsep) * \real{0.2500}}@{}}
\toprule\noalign{}
\begin{minipage}[b]{\linewidth}\raggedright
No.
\end{minipage} & \begin{minipage}[b]{\linewidth}\raggedright
Name
\end{minipage} & \begin{minipage}[b]{\linewidth}\raggedright
Description
\end{minipage} & \begin{minipage}[b]{\linewidth}\raggedright
Default
\end{minipage} \\
\midrule\noalign{}
\endhead
\bottomrule\noalign{}
\endlastfoot
Field & Metric & Laboratory & \\
1 & WELNAME & A columnar vector of character strings of up to eight
characters in length for each item, that defines the well name for which
the RFT data should be written to the RFT file. Note that the WELNAME
must have been declared previously using the WELSPECS keyword in the
SCHEDULE section, otherwise an error may occur. If the WELNAME is left
blank then the data is written out for all wells at the time they are
first opened to flow. If the WELNAME is given, then the RFT data for the
well at the time step the keyword is invoked is written out. & None \\
2 & RFT & A defined character string that sets the RFT data set output
options and should be set to one of the following character strings. &
NO \\
3 & PLT & A defined character string that sets the PLT data set output
options and should be set to one of the following character strings. &
NO \\
4 & MULTISEG & A defined character string that sets the output options
for multi-segment wells, that is the flow rates and pressures through
each well segment, and should be set to one of the following character
strings. Note the commercial simulator also uses MULTISEG to control the
output of ``rivers'' for when the RIVERS Model has been enabled via the
RIVRDIMS keyword in the RUNSPEC section. OPM Flow does not support the
RIVERS Model. & NO \\
Notes: & & & \\
\end{longtable}

\emph{Table 12.116: WRFTPLT Keyword Description}

See also the WRFT keyword in the SCHEDULE section that has less flexible
reporting options.

\subsubsection{Examples}\label{examples-95}

The first example activates RFT output at the current reporting time
step for all the wells that are opened to flow, otherwise the RFT data
is written out the first time a well is opened.

\begin{verbatim}
--
--       WELL RFT, PLT AND SEGMENT DATA
--
-- WELL  RFT   PLT   SEGMENT
-- NAME  DATA  DATA  DATA
WRFTPLT
'*'      FOPN                                                                  /
/
\end{verbatim}

The next example writes out the RFT and PLT data for two wells at the
current reporting time step.

\begin{verbatim}
--
--       WELL RFT, PLT AND SEGMENT DATA
--
-- WELL  RFT   PLT   SEGMENT
-- NAME  DATA  DATA  DATA
WRFTPLT
OP01     YES   YES                                                             /
OP02     YES   YES                                                             /
/
\end{verbatim}

The final example is shown below:

\begin{verbatim}
--
--       WELL RFT, PLT AND SEGMENT DATA
--
-- WELL  RFT   PLT   SEGMENT
-- NAME  DATA  DATA  DATA
WRFTPLT
OP01     REPT  NO                                                              /
OP02     NO    YES                                                             /
/
\end{verbatim}

In this case the RFT data for well OP01 is written out at the current
reporting time step and all subsequent reporting time steps. For well
OP02, no RFT is written out but the PLT data is written out for the
current report time step only.

\subsection{WSALT -- Define Water Injection Well Salt
Concentrations}\label{wsalt-define-water-injection-well-salt-concentrations}

\begin{longtable}[]{@{}
  >{\raggedright\arraybackslash}p{(\linewidth - 14\tabcolsep) * \real{0.1250}}
  >{\raggedright\arraybackslash}p{(\linewidth - 14\tabcolsep) * \real{0.1250}}
  >{\raggedright\arraybackslash}p{(\linewidth - 14\tabcolsep) * \real{0.1250}}
  >{\raggedright\arraybackslash}p{(\linewidth - 14\tabcolsep) * \real{0.1250}}
  >{\raggedright\arraybackslash}p{(\linewidth - 14\tabcolsep) * \real{0.1250}}
  >{\raggedright\arraybackslash}p{(\linewidth - 14\tabcolsep) * \real{0.1250}}
  >{\raggedright\arraybackslash}p{(\linewidth - 14\tabcolsep) * \real{0.1250}}
  >{\raggedright\arraybackslash}p{(\linewidth - 14\tabcolsep) * \real{0.1250}}@{}}
\toprule\noalign{}
\begin{minipage}[b]{\linewidth}\raggedright
RUNSPEC
\end{minipage} & \begin{minipage}[b]{\linewidth}\raggedright
GRID
\end{minipage} & \begin{minipage}[b]{\linewidth}\raggedright
EDIT
\end{minipage} & \begin{minipage}[b]{\linewidth}\raggedright
PROPS
\end{minipage} & \begin{minipage}[b]{\linewidth}\raggedright
REGIONS
\end{minipage} & \begin{minipage}[b]{\linewidth}\raggedright
SOLUTION
\end{minipage} & \begin{minipage}[b]{\linewidth}\raggedright
SUMMARY
\end{minipage} & \begin{minipage}[b]{\linewidth}\raggedright
SCHEDULE
\end{minipage} \\
\midrule\noalign{}
\endhead
\bottomrule\noalign{}
\endlastfoot
\end{longtable}

\subsubsection{Description}\label{description-1115}

The WSALT keyword defines a water injection well's salt injection stream
concentration that is to be used for when the salt option has been
activated by the BRINE keywords in the RUNSPEC section. Note that if the
Polymer option has also been activated by the POLYMER keyword in the
RUNSPEC section, then the WPOLYMER keyword in the SCHEDULE section
should be used to enter both the polymer and salt concentrations.

\begin{longtable}[]{@{}
  >{\raggedright\arraybackslash}p{(\linewidth - 6\tabcolsep) * \real{0.2500}}
  >{\raggedright\arraybackslash}p{(\linewidth - 6\tabcolsep) * \real{0.2500}}
  >{\raggedright\arraybackslash}p{(\linewidth - 6\tabcolsep) * \real{0.2500}}
  >{\raggedright\arraybackslash}p{(\linewidth - 6\tabcolsep) * \real{0.2500}}@{}}
\toprule\noalign{}
\begin{minipage}[b]{\linewidth}\raggedright
No.
\end{minipage} & \begin{minipage}[b]{\linewidth}\raggedright
Name
\end{minipage} & \begin{minipage}[b]{\linewidth}\raggedright
Description
\end{minipage} & \begin{minipage}[b]{\linewidth}\raggedright
Default
\end{minipage} \\
\midrule\noalign{}
\endhead
\bottomrule\noalign{}
\endlastfoot
Field & Metric & Laboratory & \\
1 & WELNAME & A character string of up to eight characters in length
that defines the well name for which the injection salt concentrations
are being defined. Note that the well name (WELNAME) must have been
declared previously using the WELSPECS keyword in the SCHEDULE section,
otherwise an error may occur. & None \\
2 & SALTCON & A real positive vector of values that defines the salt
concentration of the well's injection stream and consists of: Only
options (1) and (2) are currently supported. This value may be specified
using a User Defined Argument (UDA). Note if SALTCON is defaulted (1*)
then the well's salt concentration will be equal to the well's group
salt concentration. & 1* \\
lb/stb & kg/sm3 & gm/scc & \\
Notes: & & & \\
\end{longtable}

\emph{Table 12.117: WSALT Keyword Description}

Water injection wells that are not declared via this keyword have their
concentrations defaulted to zero.

See also the GCONPROD and GCONINJE keywords to define a group's
production and injection targets and constraints, and the WCONINJE
keyword to define an injection well's targets and constraints. All the
aforementioned keywords are described in the SCHEDULE section.

\subsubsection{Examples}\label{examples-96}

The following example defines the salt injection stream concentration
for three water injection wells for when the brine phase has been
activated by the BRINE keyword in the RUNSPEC section.

\begin{verbatim}
--
--       DEFINE WATER INJECTION WELL SALT CONCENTRATIONS (STANDARD)
--
-- WELL  SALT-1     SALT-2     SALT-3     SALT-4
-- NAME  SALTCON    SALTCON    SALTCON    SALTCON
--       -------    -------    -------    -------
WSALT
WI01     0.2500                                            /
WI02     1*                                                /
WI03     0.2500                                            /
/
\end{verbatim}

The salt concentration for both well WI01 and WI03 is set to 0.25, and
for well WI02 the salt concentration will be taken from the well's group
salt concentration.

The next example is based on using the Multi-Component Brine option,
that is the BRINE and ECLMC keywords have been used in the RUNSPEC
section, and assuming three salts.

\begin{verbatim}
--
--       DEFINE WATER INJECTION WELL SALT CONCENTRATIONS (MULTIPLE)
--
-- WELL  SALT-1     SALT-2     SALT-3     SALT-4
-- NAME  SALTCON    SALTCON    SALTCON    SALTCON
--       -------    -------    -------    -------
WSALT
WI01     0.1500     0.0500     0.0500                      /
WI02     0.1500     0.0500     0.0500                      /
WI03     0.2000     0.0500     0.0600                      /
/
\end{verbatim}

Here the salt concentrations for both well WI01 and WI02 are set to
0.1500, 0.0500, 0.0500 for the three salts and for well WI03 the salt
concentrations are 0.2000, 0.0500 and 0.0600.

Note that OPM Flow does not currently support the Multi-Component brine
model.

\subsection{WSCCLEAN -- Well Deposited Scale
Adjustment}\label{wscclean-well-deposited-scale-adjustment}

\begin{longtable}[]{@{}
  >{\raggedright\arraybackslash}p{(\linewidth - 14\tabcolsep) * \real{0.1250}}
  >{\raggedright\arraybackslash}p{(\linewidth - 14\tabcolsep) * \real{0.1250}}
  >{\raggedright\arraybackslash}p{(\linewidth - 14\tabcolsep) * \real{0.1250}}
  >{\raggedright\arraybackslash}p{(\linewidth - 14\tabcolsep) * \real{0.1250}}
  >{\raggedright\arraybackslash}p{(\linewidth - 14\tabcolsep) * \real{0.1250}}
  >{\raggedright\arraybackslash}p{(\linewidth - 14\tabcolsep) * \real{0.1250}}
  >{\raggedright\arraybackslash}p{(\linewidth - 14\tabcolsep) * \real{0.1250}}
  >{\raggedright\arraybackslash}p{(\linewidth - 14\tabcolsep) * \real{0.1250}}@{}}
\toprule\noalign{}
\begin{minipage}[b]{\linewidth}\raggedright
RUNSPEC
\end{minipage} & \begin{minipage}[b]{\linewidth}\raggedright
GRID
\end{minipage} & \begin{minipage}[b]{\linewidth}\raggedright
EDIT
\end{minipage} & \begin{minipage}[b]{\linewidth}\raggedright
PROPS
\end{minipage} & \begin{minipage}[b]{\linewidth}\raggedright
REGIONS
\end{minipage} & \begin{minipage}[b]{\linewidth}\raggedright
SOLUTION
\end{minipage} & \begin{minipage}[b]{\linewidth}\raggedright
SUMMARY
\end{minipage} & \begin{minipage}[b]{\linewidth}\raggedright
SCHEDULE
\end{minipage} \\
\midrule\noalign{}
\endhead
\bottomrule\noalign{}
\endlastfoot
\end{longtable}

\subsubsection{Description}\label{description-1116}

The WSCCLEAN keyword adjusts the amount of scale currently accumulated
around a well's well connections for wells located in the global grid.
For example, if a workover has been performed on a well to remove (or
reduce) the deposited scale over the perforations, then this keyword can
be used to implement the effects of the workover. Scale deposits reduce
the productivity of well and this relationship is defined in the SCDPTAB
and SCDATAB keywords in SCHEDULE section. The tables are allocated to a
well via the WSCTAB keyword, which is also in the SCHEDULE section. Note
that the Scale Deposition option must have been activated by declaring
the dimensions of the scaling deposition tables using the SCDPDIMS
keyword in the RUNSPEC section.

See also the WSSCLENL keyword in the SCHEDULE section that performs
similar functionality for wells located in a Local Grid Refinement.

This keyword is not supported by OPM Flow but would change the results
if supported so the simulation will be stopped.

\subsection{WSCCLENL -- Well Deposited Scale Adjustment
(LGR)}\label{wscclenl-well-deposited-scale-adjustment-lgr}

\begin{longtable}[]{@{}
  >{\raggedright\arraybackslash}p{(\linewidth - 14\tabcolsep) * \real{0.1250}}
  >{\raggedright\arraybackslash}p{(\linewidth - 14\tabcolsep) * \real{0.1250}}
  >{\raggedright\arraybackslash}p{(\linewidth - 14\tabcolsep) * \real{0.1250}}
  >{\raggedright\arraybackslash}p{(\linewidth - 14\tabcolsep) * \real{0.1250}}
  >{\raggedright\arraybackslash}p{(\linewidth - 14\tabcolsep) * \real{0.1250}}
  >{\raggedright\arraybackslash}p{(\linewidth - 14\tabcolsep) * \real{0.1250}}
  >{\raggedright\arraybackslash}p{(\linewidth - 14\tabcolsep) * \real{0.1250}}
  >{\raggedright\arraybackslash}p{(\linewidth - 14\tabcolsep) * \real{0.1250}}@{}}
\toprule\noalign{}
\begin{minipage}[b]{\linewidth}\raggedright
RUNSPEC
\end{minipage} & \begin{minipage}[b]{\linewidth}\raggedright
GRID
\end{minipage} & \begin{minipage}[b]{\linewidth}\raggedright
EDIT
\end{minipage} & \begin{minipage}[b]{\linewidth}\raggedright
PROPS
\end{minipage} & \begin{minipage}[b]{\linewidth}\raggedright
REGIONS
\end{minipage} & \begin{minipage}[b]{\linewidth}\raggedright
SOLUTION
\end{minipage} & \begin{minipage}[b]{\linewidth}\raggedright
SUMMARY
\end{minipage} & \begin{minipage}[b]{\linewidth}\raggedright
SCHEDULE
\end{minipage} \\
\midrule\noalign{}
\endhead
\bottomrule\noalign{}
\endlastfoot
\end{longtable}

\subsubsection{Description}\label{description-1117}

The WSCCLENL keyword adjusts the amount of scale currently accumulated
around a well's well connections for wells located in a Local Grid
Refinement (``LGR''). For example, if a workover has been performed on a
well to remove (or reduce) the deposited scale over the perforations,
then this keyword can be used to implement the effects of the workover.
Scale deposits reduce the productivity of well and this relationship is
defined in the SCDPTAB and SCDATAB keywords in SCHEDULE section. The
tables are allocated to a well via the WSCTAB keyword, which is also in
the SCHEDULE section. Note that the Scale Deposition option must have
been activated by declaring the dimensions of the scaling deposition
tables using the SCDPDIMS keyword in the RUNSPEC section.

See also the WSCCLEAN keyword in the SCHEDULE section that performs
similar functionality for wells located in the global grid.

This keyword is not supported by OPM Flow but would change the results
if supported so the simulation will be stopped.

\subsection{WSCTAB -- Assign Well Scale Deposition and Scale Damage
Tables}\label{wsctab-assign-well-scale-deposition-and-scale-damage-tables}

\begin{longtable}[]{@{}
  >{\raggedright\arraybackslash}p{(\linewidth - 14\tabcolsep) * \real{0.1250}}
  >{\raggedright\arraybackslash}p{(\linewidth - 14\tabcolsep) * \real{0.1250}}
  >{\raggedright\arraybackslash}p{(\linewidth - 14\tabcolsep) * \real{0.1250}}
  >{\raggedright\arraybackslash}p{(\linewidth - 14\tabcolsep) * \real{0.1250}}
  >{\raggedright\arraybackslash}p{(\linewidth - 14\tabcolsep) * \real{0.1250}}
  >{\raggedright\arraybackslash}p{(\linewidth - 14\tabcolsep) * \real{0.1250}}
  >{\raggedright\arraybackslash}p{(\linewidth - 14\tabcolsep) * \real{0.1250}}
  >{\raggedright\arraybackslash}p{(\linewidth - 14\tabcolsep) * \real{0.1250}}@{}}
\toprule\noalign{}
\begin{minipage}[b]{\linewidth}\raggedright
RUNSPEC
\end{minipage} & \begin{minipage}[b]{\linewidth}\raggedright
GRID
\end{minipage} & \begin{minipage}[b]{\linewidth}\raggedright
EDIT
\end{minipage} & \begin{minipage}[b]{\linewidth}\raggedright
PROPS
\end{minipage} & \begin{minipage}[b]{\linewidth}\raggedright
REGIONS
\end{minipage} & \begin{minipage}[b]{\linewidth}\raggedright
SOLUTION
\end{minipage} & \begin{minipage}[b]{\linewidth}\raggedright
SUMMARY
\end{minipage} & \begin{minipage}[b]{\linewidth}\raggedright
SCHEDULE
\end{minipage} \\
\midrule\noalign{}
\endhead
\bottomrule\noalign{}
\endlastfoot
\end{longtable}

\subsubsection{Description}\label{description-1118}

WSCTAB assigns scale deposition and scale damage tables to a well, for
when the Scale Deposition option has been activated by declaring the
dimensions of the scaling deposition tables using the SCDPDIMS keyword
in the RUNSPEC section. Scale deposits reduce the productivity of well
and this relationship is defined in the SCDPTAB and SCDATAB keywords in
the SCHEDULE section, and are allocated to a well by the WSCTAB keyword.

See also the WPIMULT keyword in the SCHEDULE section that adjusts a
well's productivity index by a constant value.

This keyword is not supported by OPM Flow but would change the results
if supported so the simulation will be stopped.

\subsection{WSEGAICD -- Define Multi-Segment Well Autonomous ICD
Connections}\label{wsegaicd-define-multi-segment-well-autonomous-icd-connections}

\begin{longtable}[]{@{}
  >{\raggedright\arraybackslash}p{(\linewidth - 14\tabcolsep) * \real{0.1250}}
  >{\raggedright\arraybackslash}p{(\linewidth - 14\tabcolsep) * \real{0.1250}}
  >{\raggedright\arraybackslash}p{(\linewidth - 14\tabcolsep) * \real{0.1250}}
  >{\raggedright\arraybackslash}p{(\linewidth - 14\tabcolsep) * \real{0.1250}}
  >{\raggedright\arraybackslash}p{(\linewidth - 14\tabcolsep) * \real{0.1250}}
  >{\raggedright\arraybackslash}p{(\linewidth - 14\tabcolsep) * \real{0.1250}}
  >{\raggedright\arraybackslash}p{(\linewidth - 14\tabcolsep) * \real{0.1250}}
  >{\raggedright\arraybackslash}p{(\linewidth - 14\tabcolsep) * \real{0.1250}}@{}}
\toprule\noalign{}
\begin{minipage}[b]{\linewidth}\raggedright
RUNSPEC
\end{minipage} & \begin{minipage}[b]{\linewidth}\raggedright
GRID
\end{minipage} & \begin{minipage}[b]{\linewidth}\raggedright
EDIT
\end{minipage} & \begin{minipage}[b]{\linewidth}\raggedright
PROPS
\end{minipage} & \begin{minipage}[b]{\linewidth}\raggedright
REGIONS
\end{minipage} & \begin{minipage}[b]{\linewidth}\raggedright
SOLUTION
\end{minipage} & \begin{minipage}[b]{\linewidth}\raggedright
SUMMARY
\end{minipage} & \begin{minipage}[b]{\linewidth}\raggedright
SCHEDULE
\end{minipage} \\
\midrule\noalign{}
\endhead
\bottomrule\noalign{}
\endlastfoot
\end{longtable}

\subsubsection{Description}\label{description-1119}

The WSEGAICD keyword defines a multi-segment well segment to be an
autonomous Inflow Control Device (``ICD'') as part of a completion for a
multi-segment well. Note that the well must have been previously defined
by the WELSPECS and WELSEGS keywords in the SCHEDULE section and that
the data for the keyword should be repeated for each multi-segment
completion that contains an autonomous ICD.

An ICD is a well completion component usually installed along the
producing section of a well to minimize the unwanted water and gas
breakthrough in an oil well, or early water production in a gas well,
due to an uneven inflow profile over the completed interval.
Permeability variations over the producing interval cause the high
permeability zones to produce higher quantities of fluids than the lower
permeability zones and this uneven producing fluid profile may result in
bypassed hydrocarbons. Secondly, for horizontal wells, the pressure loss
from the ``toe'' to the ``heel'' of the well again results in an uneven
fluid profile over the producing interval. In order to rectify this ICDs
can be installed so that the well fluids have to flow through an ICD
before entering the tubing; thus, creating an additional ``designed''
pressure loss.

An autonomous ICD inhibits the production of high-mobility fluids such
as water and gas since the pressure drop in each unit is dependent on
fluid properties and mobility, the device automatically increases the
pressure differential across zones with high water or gas saturations,
thus choking back production from these zones.

\begin{longtable}[]{@{}
  >{\raggedright\arraybackslash}p{(\linewidth - 6\tabcolsep) * \real{0.2500}}
  >{\raggedright\arraybackslash}p{(\linewidth - 6\tabcolsep) * \real{0.2500}}
  >{\raggedright\arraybackslash}p{(\linewidth - 6\tabcolsep) * \real{0.2500}}
  >{\raggedright\arraybackslash}p{(\linewidth - 6\tabcolsep) * \real{0.2500}}@{}}
\toprule\noalign{}
\begin{minipage}[b]{\linewidth}\raggedright
No.
\end{minipage} & \begin{minipage}[b]{\linewidth}\raggedright
Name
\end{minipage} & \begin{minipage}[b]{\linewidth}\raggedright
Description
\end{minipage} & \begin{minipage}[b]{\linewidth}\raggedright
Default
\end{minipage} \\
\midrule\noalign{}
\endhead
\bottomrule\noalign{}
\endlastfoot
Field & Metric & Laboratory & \\
1 & WELNAME & A character string of up to eight characters in length
that defines the well name for which a multi-segment well is being
defined. Note that the well name (WELNAME) must have been declared
previously using both the WELSPECS and WELSEGS keywords in the SCHEDULE
section, otherwise an error may occur. & None \\
2 & ISEG1 & An integer greater than or equal to two and less than or
equal to MXSEGS on the WSEGDIMS keyword in the RUNSPEC section that
defines the start of the segment range. & None \\
3 & ISEG2 & An integer greater than or equal to two and not less then
ISEG1 on this record and less than or equal to MXSEGS on the WSEGDIMS
keyword in the RUNSPEC section, that defines the end of the segment
range. & None \\
4 & ICDSTREN & A positive real value that defines an empirical constant
for the strength of the given ICD as determined from measurements using
the calibrated fluid. & None \\
psi/((lb/ft3)(rft3/day)2) & bars/((kg/m3)(rm3/day)2) &
atm/((gm/cc)(rcc/hr)2) & \\
5 & ICDLEN & A real value that defines the length of the ICD used in
conjunction with NSCALFAC to calculate a scaling factor to be applied to
the reservoir flow to adjust the flow through each ICD, that is:
NSCALFAC explicitly sets which of the above three options is used. If
NSCALFAC is defaulted, then option 1) is used whenever ICDLEN is
positive and option 2) when ICDLEN is negative. & Defined \\
feet 39.37 & m 12.00 & cm 1,2000 & \\
6 & CALDEN & A positive real value that defines the density of the
calibrating fluid at surface conditions. & Defined \\
lb/ft3 62.416 & kg/m3 1000.25 & gm/cc 1.00025 & \\
7 & CALVISC & A positive real value that defines the viscosity of the
calibrating fluid at surface conditions. & 1.45 \\
cP & cP & cP & \\
8 & EMLCRT & A positive real value that defines the ``local water'' in
liquid fraction used to determine whether the ``water-in-oil'' or
``oil-in-water'' emulsion viscosity equation should be applied. & 0.5 \\
dimensionless & dimensionless & dimensionless & \\
9 & EMLTRANS & A positive real value that defines the width of the
transition zone around EMLCRT and is used to ensure that the calculated
viscosity forms a continuous function of water in liquid fraction.
Within this region, the emulsion viscosity is a linear interpolation
between the ``water-in-oil'' and ``oil-in-water'' viscosity values
either side of the region. & 0.05 \\
dimensionless & dimensionless & dimensionless & \\
10 & EMLMAX & A positive real value that defines the maximum emulsion
viscosity to continuous phase viscosity (oil or water) ratio. & 5.0 \\
dimensionless & dimensionless & dimensionless & \\
11 & NSCALFAC & An integer value greater than or equal to zero, that
specifies the method to be used when applying the scaling factor and
should be set to one of the following: NSCALFAC explicitly sets which of
the above three options is used. If NSCALFAC is negative, then option 1)
is used whenever ICDLEN is positive and option 2) when ICDLEN is
negative. & -1 \\
dimensionless & dimensionless & dimensionless & \\
12 & CALRATE & A positive real value that defines the maximum surface
flow rate for which the ICD was calibrated. Values calculated greater
than CALRATE will use linear extrapolation. & None \\
scf/d & sm3/day & scc/hour & \\
13 & RATEXP & A real value greater than or equal to zero that defines
the flow rate exponent in equation (12.35). & None \\
dimensionless & dimensionless & dimensionless & \\
14 & VISCEXP & A real value that defines the viscosity exponent in
equation (12.35). & None \\
dimensionless & dimensionless & dimensionless & \\
15 & STATUS & A defined character string that specifies the ICD's
operational status, STATUS should be set to one of the following
character strings: & OPEN \\
16 & A1 & A real value that defines the density mixture OIL flowing
fraction exponent in equation (12.36) & 1.0 \\
17 & A2 & A real value that defines the density mixture WATER flowing
fraction exponent in equation (12.36) & 1.0 \\
18 & A3 & A real value that defines the density mixture GAS flowing
fraction exponent in equation (12.36) & 1.0 \\
19 & B1 & A real value that defines the viscosity mixture OIL flowing
fraction exponent in equation (12.37) & 1.0 \\
20 & B2 & A real value that defines the viscosity mixture WATER flowing
fraction exponent in equation (12.37) & 1.0 \\
21 & B3 & A real value that defines the viscosity mixture GAS flowing
fraction exponent in equation (12.37) & 1.0 \\
22 & DENEXP & A real value that defines the density exponent in equation
(12.35). & 1.0 \\
dimensionless & dimensionless & dimensionless & \\
Notes: & & & \\
\end{longtable}

\emph{Table 12.118: WSEGAICD Keyword Description}

The total number of wells should be defined via the WELLDIMS keyword and
the number of multi-segment wells should be declared on the WSEGDIMS
keyword, both keywords are in the RUNSPEC section. In addition, the
WELSPECS keyword should be used to define wells, the COMPDAT keyword to
define the well completions for both standard wells and multi-segment
wells, and the COMPSEGS keyword to define a multi-segment well segment
completions. Finally, the WSEGSICD keyword can be used to define spiral
ICD connections for the well. All the aforementioned keywords are
described in the SCHEDULE section.

The equations used to calculate the pressure drop across the ICD are
given below and illustrate how the pressure reduction is dependent on
the density and viscosity of the fluid flowing through the device.

\begin{equation}\phantomsection\label{eq-12-35}{
\mathrm{Δ}P = \left(\frac{{\mathrm{ρ}}_{\mathit{mixture}}}{{\mathrm{ρ}}_{\mathit{calibrated}}}\right){\mathrm{ρ}}_{exp}⋅{\left(\frac{{\mathrm{μ}}_{\mathit{calibrated}}}{{\mathrm{μ}}_{\mathit{mixture}}}\right)}^{{\mathrm{μ}}_{exp}}⋅ {\mathrm{ρ}}_{\mathit{mixture}} ⋅ \mathrm{β} ⋅ {\left(\frac{q}{{q}_{\mathit{calibrated}}}\right)}^{({Q}_{exp}- 2)} ⋅ {q}^{2}
}\end{equation}

Where:

ΔP = the pressure drop across the device.

ρmixture = the density of the mixture, as per:

\begin{equation}\phantomsection\label{eq-12-36}{
{\mathrm{ρ}}_{\mathit{mixture}} = \left({\mathrm{α}}_{\mathit{oil}}^{{a}_{1}}⋅{\mathrm{ρ}}_{\mathit{oil}}\right) + \left({\mathrm{α}}_{\mathit{wat}}^{{a}_{2}}⋅{\mathrm{ρ}}_{\mathit{wat}}\right) + \left({\mathrm{α}}_{\mathit{gas}}^{{a}_{3}}⋅{\mathrm{ρ}}_{\mathit{gas}}\right)
}\end{equation}

ρcalibrated = CALDEN, the density of the calibrating fluid at surface
conditions.

μcalibrated = CALVISC, the viscosity of the calibrating fluid at surface
conditions.

μmixture = the viscosity of the mixture, as per:

\begin{equation}\phantomsection\label{eq-12-37}{
{\mathrm{μ}}_{\mathit{mixture}} = \left({\mathrm{α}}_{\mathit{oil}}^{{b}_{1}}⋅{\mathrm{μ}}_{\mathit{oil}}\right) + \left({\mathrm{α}}_{\mathit{wat}}^{{b}_{2}}⋅{\mathrm{μ}}_{\mathit{wat}}\right) + \left({\mathrm{α}}_{\mathit{gas}}^{{b}_{3}}⋅{\mathrm{μ}}_{\mathit{gas}}\right)
}\end{equation}

β = ICDSTREN, the strength of the ICD as measured using the calibrated
fluid.

q = the local flow rate through the ICD at local conditions adjusted for
scaling

based on the ICDLEN and NSCALFAC parameters in Table 12.118.

qcalibrated = CALRATE, the maximum surface flow rate for which the ICD
was

calibrated.

ρexp = DENEXP, the density exponent in Table 12.118.

μexp = VISCEXP, the viscosity exponent in Table 12.118.

Qexp = RATEXP, the flow rate exponent in Table 12.118.

In equation (12.8) αi represent the volume fractions of oil, water and
gas at local conditions and ai the density exponents of the three phases
(A1, A2 and A3 in Table 12.118). Similarity for equation (12.37), αi
represent the volume fractions of oil, water and gas at local conditions
and bi the viscosity exponents of the three phases (B1, B2 and B3 in
Table 12.118).

See also the WSEGSICD keyword in the SCHEDULE section for spiral ICDs,
that operate in a similar fashion to autonomous ICDs.

\subsubsection{Example}\label{example-600}

The following example defines a multi-segment oil production well (OP01)
using the WELSPECS, WELSEGS COMPDAT and COMPSEGS keywords, followed by
the WSEGAICD keyword to define the autonomous inflow control devices for
the well.

\begin{verbatim}
--
--       WELL SPECIFICATION DATA
--
-- WELL  GROUP   LOCATION  BHP   PHASE  DRAIN  INFLOW  OPEN  CROSS PVT   DEN  FIP
-- NAME  NAME      I    J  DEPTH FLUID  AREA   EQUANS  SHUT  FLOW  TABLE CAL  NUM
WELSPECS
OP01     PLATFORM 14    8  1*    OIL    0.0    STD     SHUT  YES   0     SEG  0 /
/
--
--       WELL CONNECTION DATA
--
-- WELL  --- LOCATION ---  OPEN   SAT  CONN         WELL  KH    SKIN   D    DIR
-- NAME   II  JJ  K1  K2   SHUT   TAB  FACT         DIA   FACT  FACT   FACT PEN
COMPDAT
OP01      14   8   1    1  OPEN   1*   1.972960E+2 0.216  1*    0.00   1*   'Z' /
OP01      14   7   1    1  OPEN   1*   1.619450E+2 0.216  1*    0.00   1*   'Z' /
OP01      14   7   2    2  OPEN   1*   4.126449E+2 0.216  1*    0.00   1*   'Z' /
OP01      14   7   3    3  OPEN   1*   2.033290E+1 0.216  1*    0.00   1*   'Z' /
OP01      15   7   3    3  OPEN   1*   9.095613E+1 0.216  1*    0.00   1*   'Z' /
OP01      15   6   3    3  OPEN   1*   2.090607E+2 0.216  1*    0.00   1*   'Y' /
OP01      15   6   4    4  OPEN   1*   3.010669E+1 0.216  1*    0.00   1*   'Y' /
OP01      16   6   4    4  OPEN   1*   7.123814E+1 0.216  1*    0.00   1*   'Y' /
OP01      16   5   4    4  OPEN   1*   4.414386E+2 0.216  1*    0.00   1*   'Y' /
OP01      16   4   4    4  OPEN   1*   4.345126E+2 0.216  1*    0.00   1*   'Y' /
OP01      16   3   4    4  OPEN   1*   2.894573E+2 0.216  1*    0.00   1*   'Y' /
OP01      17   3   4    4  OPEN   1*   1.329523E+2 0.216  1*    0.00   1*   'Y' /
OP01      17   2   4    4  OPEN   1*   6.981252E+1 0.216  1*    0.00   1*   'Y' /
OP01      17   2   5    5  OPEN   1*   1.392382E+2 0.216  1*    0.00   1*   'Y' /
/
--
--       WELL SEGMENT SPECIFICATION DATA
--
-- WELL  NODAL       LEN     WELL    DEPH    PRESS   FLOW
-- NAME  DEPTH       TUBING  VOLM    OPTN    CALC    MODEL
WELSEGS
   OP01  2041.56259  0.0000  1*      INC     'HF-'          /
--
--       SEG   SEG   BRAN    SEG     TUBING  NODAL   TUBE   TUBE     XSEC  VOL
--       ISTR  IEND  NO      NO      LENGTH  DEPTH   ID     ROUGH    AREA  SEG
          2     2    1        1     9.56005  0.61168 0.152  0.0000100         /
          3     3    1        2    17.40714  1.11376 0.152  0.0000100         /
          4     4    1        3    41.24996  2.38255 0.152  0.0000100         /
          5     5    1        4    38.35922  2.06899 0.152  0.0000100         /
          6     6    1        5    27.13029  1.02248 0.152  0.0000100         /
          7     7    1        6    73.45099  2.40389 0.152  0.0000100         /
          8     8    1        7    54.95304  1.64832 0.152  0.0000100         /
          9     9    1        8    12.37381  0.23781 0.152  0.0000100         /
         10    10    1        9    62.61459  0.73403 0.152  0.0000100         /
         11    11    1       10    106.9805  1.37749 0.152  0.0000100         /
         12    12    1       11    88.42739  0.90931 0.152  0.0000100         /
         13    13    1       12    51.59899  0.27327 0.152  0.0000100         /
         14    14    1       13    24.75582  0.30814 0.152  0.0000100         /
         15    15    1       14    29.77334  0.49371 0.152  0.0000100         /

--
-- PERFORATION VALVE SEGMENTS
--
--       SEG   SEG   BRAN    SEG     TUBING  NODAL   TUBE   TUBE     XSEC  VOL
--       ISTR  IEND  NO      NO      LENGTH  DEPTH   ID     ROUGH    AREA  SEG
         16    16    2        2    0.10000   0       0.152  0.0000100         /
         17    17    3        3    0.10000   0       0.152  0.0000100         /
         18    18    4        4    0.10000   0       0.152  0.0000100         /
         19    19    5        5    0.10000   0       0.152  0.0000100         /
         20    20    6        6    0.10000   0       0.152  0.0000100         /
         21    21    7        7    0.10000   0       0.152  0.0000100         /
         22    22    8        8    0.10000   0       0.152  0.0000100         /
         23    23    9        9    0.10000   0       0.152  0.0000100         /
         24    24   10       10    0.10000   0       0.152  0.0000100         /
         25    25   11       11    0.10000   0       0.152  0.0000100         /
         26    26   12       12    0.10000   0       0.152  0.0000100         /
         27    27   13       13    0.10000   0       0.152  0.0000100         /
         28    28   14       14    0.10000   0       0.152  0.0000100         /
         29    29   15       15    0.10000   0       0.152  0.0000100         /
/
--
--       MULTISEGMENT WELL COMPLETION SEGMENT SPECIFICATION DATA
--
-- WELL
-- NAME
COMPSEGS
OP01                                                                           /
--
--      --LOCATION--  BRAN  TUBING     NODAL    DIR  LOC    MID    COMP    ISEG
--       II  JJ  K1   NO    LENGTH     DEPTH    PEN  I,J,K  PERFS  LENGTH  NO.
         14   8   1   2     0.000000   0.10000                                 /
         14   7   1   3     0.000000   0.10000                                 /
         14   7   2   4     0.000000   0.10000                                 /
         14   7   3   5     0.000000   0.10000                                 /
         15   7   3   6     0.000000   0.10000                                 /
         15   6   3   7     0.000000   0.10000                                 /
         15   6   4   8     0.000000   0.10000                                 /
         16   6   4   9     0.000000   0.10000                                 /
         16   5   4   10    46.76168  46.86168                                 /
         16   4   4   11    0.000000   0.10000                                 /
         16   3   4   12    0.000000   0.10000                                 /
         17   3   4   13    0.000000   0.10000                                 /
         17   2   4   14    0.000000   0.10000                                 /
         17   2   5   15    42.50573  42.60573                                 /
/
--
--       MULTI-SEGMENT WELL ICD SEGMENT SPECIFICATION DATA
--
-- WELL SEG  SEG  ICD    ICD      CAL  CAL  EML  EML  EML SCAL CAL RATE VISC OPEN
-- NAME ISTR IEND STREN  LEN      DEN  VISC CRIT TRAN MAX FAC  RAT EXP  EXP CLOSE
WSEGAICD
OP01    16   16  7.5e-5 -0.028975 1020 0.48 0.7  1*   1*  1*   1*  2.2  0.5  7* /
OP01    17   17  7.5e-5 -0.023783 1020 0.48 0.7  1*   1*  1*   1*  2.2  0.5  7* /
OP01    18   18  7.5e-5 -0.101240 1020 0.48 0.7  1*   1*  1*   1*  2.2  0.5  7* /
OP01    19   19  7.5e-5 -0.015022 1020 0.48 0.7  1*   1*  1*   1*  2.2  0.5  7* /
OP01    20   20  7.5e-5 -0.067205 1020 0.48 0.7  1*   1*  1*   1*  2.2  0.5  7* /
OP01    21   21  7.5e-5 -0.155410 1020 0.48 0.7  1*   1*  1*   1*  2.2  0.5  7* /
OP01    22   22  7.5e-5 -0.011141 1020 0.48 0.7  1*   1*  1*   1*  2.2  0.5  7* /
OP01    23   23  7.5e-5 -0.026362 1020 0.48 0.7  1*   1*  1*   1*  2.2  0.5  7* /
OP01    24   24  7.5e-5 -0.163410 1020 0.48 0.7  1*   1*  1*   1*  2.2  0.5  7* /
OP01    25   25  7.5e-5 -0.160830 1020 0.48 0.7  1*   1*  1*   1*  2.2  0.5  7* /
OP01    26   26  7.5e-5 -0.107180 1020 0.48 0.7  1*   1*  1*   1*  2.2  0.5  7* /
--
-- WELL SEG  SEG  ICD    ICD      CAL  CAL  EML  EML  EML SCAL CAL RATE VISC OPEN
-- NAME ISTR IEND STREN  LEN      DEN  VISC CRIT TRAN MAX FAC  RAT EXP  EXP CLOSE
OP01    27   27  7.5e-5 -0.049206 1020 0.48 0.7  1*   1*  1*   1*  2.2  0.5  7* /
OP01    28   28  7.5e-5 -0.025824 1020 0.48 0.7  1*   1*  1*   1*  2.2  0.5  7* /
OP01    29   29  7.5e-5 -0.064414 1020 0.48 0.7  1*   1*  1*   1*  2.2  0.5  7* /
/
\end{verbatim}

The WSEGAICD keyword defines the autonomous ICD over segments 16 to 29
using the default value for NSCALFAC of -1, which sets the scale factor
equal to the absolute value of ICDLEN, as all the ICDLEN values are
negative. Note that the WSEGAICD should be repeated for each
multi-segment well that has this type of ICD. So for example for two
wells one would have:

\begin{verbatim}
--
--       MULTI-SEGMENT WELL ICD SEGMENT SPECIFICATION DATA
--
-- WELL SEG  SEG  ICD    ICD      CAL  CAL  EML  EML  EML SCAL CAL RATE VISC OPEN
-- NAME ISTR IEND STREN  LEN      DEN  VISC CRIT TRAN MAX FAC  RAT EXP  EXP CLOSE
WSEGAICD
OP01    16   16  7.5e-5 -0.028975 1020 0.48 0.7  1*   1*  1*   1*  2.2  0.5  7* /
...
OP01    28   28  7.5e-5 -0.025824 1020 0.48 0.7  1*   1*  1*   1*  2.2  0.5  7* /
/

-- WELL SEG  SEG  ICD    ICD      CAL  CAL  EML  EML  EML SCAL CAL RATE VISC OPEN
-- NAME ISTR IEND STREN  LEN      DEN  VISC CRIT TRAN MAX FAC  RAT EXP  EXP CLOSE
WSEGAICD
OP02    17   17  1.75E-4 .0049945 1020 0.48 1*   1*   1*  1*   1*  2.1  1.2  7* /
...
OP02    31   31  1.75E-4 .0038583 1020 0.48 1*   1*   1*  1*   1*  2.1  1.2  7* /
/
\end{verbatim}

Figure 12.16 shows the layout of the two example wells.

\pandocbounded{\includegraphics[keepaspectratio]{chapters/images/Frame73_6b237212f230.png}}
\pandocbounded{\includegraphics[keepaspectratio]{chapters/images/Image66_6b237212f230.png}}

\subsection{WSEGDFIN -- Define Multi-Segment Well Drift Flux Slip Model
Input
Data}\label{wsegdfin-define-multi-segment-well-drift-flux-slip-model-input-data}

\begin{longtable}[]{@{}
  >{\raggedright\arraybackslash}p{(\linewidth - 14\tabcolsep) * \real{0.1250}}
  >{\raggedright\arraybackslash}p{(\linewidth - 14\tabcolsep) * \real{0.1250}}
  >{\raggedright\arraybackslash}p{(\linewidth - 14\tabcolsep) * \real{0.1250}}
  >{\raggedright\arraybackslash}p{(\linewidth - 14\tabcolsep) * \real{0.1250}}
  >{\raggedright\arraybackslash}p{(\linewidth - 14\tabcolsep) * \real{0.1250}}
  >{\raggedright\arraybackslash}p{(\linewidth - 14\tabcolsep) * \real{0.1250}}
  >{\raggedright\arraybackslash}p{(\linewidth - 14\tabcolsep) * \real{0.1250}}
  >{\raggedright\arraybackslash}p{(\linewidth - 14\tabcolsep) * \real{0.1250}}@{}}
\toprule\noalign{}
\begin{minipage}[b]{\linewidth}\raggedright
RUNSPEC
\end{minipage} & \begin{minipage}[b]{\linewidth}\raggedright
GRID
\end{minipage} & \begin{minipage}[b]{\linewidth}\raggedright
EDIT
\end{minipage} & \begin{minipage}[b]{\linewidth}\raggedright
PROPS
\end{minipage} & \begin{minipage}[b]{\linewidth}\raggedright
REGIONS
\end{minipage} & \begin{minipage}[b]{\linewidth}\raggedright
SOLUTION
\end{minipage} & \begin{minipage}[b]{\linewidth}\raggedright
SUMMARY
\end{minipage} & \begin{minipage}[b]{\linewidth}\raggedright
SCHEDULE
\end{minipage} \\
\midrule\noalign{}
\endhead
\bottomrule\noalign{}
\endlastfoot
\end{longtable}

\subsubsection{Description}\label{description-1120}

The WSEGDFIN keyword defines a multi-segment well's drift flux slip
model parameters. A slip model {[}Shi, H., Holmes, J.A., Durlofsky, L.
J., Aziz, K., Diaz, L. R., Alkaya, B., and Oddie, G. ``Drift-Flux
Modeling of Two-Phase Flow in Wellbores,'' paper SPE 84228, Society of
Petroleum Engineers Journal (2005) 10, No.~1, 24-33; also presented as
``Drift-Flux Modeling of Multiphase Flow in Wellbores,'' at the SPE
Annual Technical Conference and Exhibition, Denver, Colorado,
USA(October 5-8, 2003).{]} and {[}Shi, H., Holmes, J.A., Diaz, L. R.,
Durlofsky, L. J., and Aziz, K. ``Drift-Flux Parameters for Three-Phase
Steady-State Flow in Wellbores,'' paper SPE 89836, Society of Petroleum
Engineers Journal(2005) 10, No.~2, 130-137; also presented at the SPE
Annual Technical Conference and Exhibition, Houston, Texas, USA
(September 26-29, 2004).{]} enables the different phases in the wellbore
to flow at different velocities, for example gas will flow up the tubing
at a higher velocity than oil and water. The option is activated by the
FLOWOPT parameter on the WELSEGS keyword in the SCHEDULE section;
however, the slip model flow calculation is not available in OPM Flow.

See also the WSEGFMD keyword in the SCHEDULE section that sets which
type of slip model should be used.

This keyword is not supported by OPM Flow but would change the results
if supported so the simulation will be stopped.

\subsection{WSEGDFMD -- Define Multi-Segment Well Drift Flux Slip
Model}\label{wsegdfmd-define-multi-segment-well-drift-flux-slip-model}

\begin{longtable}[]{@{}
  >{\raggedright\arraybackslash}p{(\linewidth - 14\tabcolsep) * \real{0.1250}}
  >{\raggedright\arraybackslash}p{(\linewidth - 14\tabcolsep) * \real{0.1250}}
  >{\raggedright\arraybackslash}p{(\linewidth - 14\tabcolsep) * \real{0.1250}}
  >{\raggedright\arraybackslash}p{(\linewidth - 14\tabcolsep) * \real{0.1250}}
  >{\raggedright\arraybackslash}p{(\linewidth - 14\tabcolsep) * \real{0.1250}}
  >{\raggedright\arraybackslash}p{(\linewidth - 14\tabcolsep) * \real{0.1250}}
  >{\raggedright\arraybackslash}p{(\linewidth - 14\tabcolsep) * \real{0.1250}}
  >{\raggedright\arraybackslash}p{(\linewidth - 14\tabcolsep) * \real{0.1250}}@{}}
\toprule\noalign{}
\begin{minipage}[b]{\linewidth}\raggedright
RUNSPEC
\end{minipage} & \begin{minipage}[b]{\linewidth}\raggedright
GRID
\end{minipage} & \begin{minipage}[b]{\linewidth}\raggedright
EDIT
\end{minipage} & \begin{minipage}[b]{\linewidth}\raggedright
PROPS
\end{minipage} & \begin{minipage}[b]{\linewidth}\raggedright
REGIONS
\end{minipage} & \begin{minipage}[b]{\linewidth}\raggedright
SOLUTION
\end{minipage} & \begin{minipage}[b]{\linewidth}\raggedright
SUMMARY
\end{minipage} & \begin{minipage}[b]{\linewidth}\raggedright
SCHEDULE
\end{minipage} \\
\midrule\noalign{}
\endhead
\bottomrule\noalign{}
\endlastfoot
\end{longtable}

\subsubsection{Description}\label{description-1121}

The WSEGDFMD keyword defines a multi-segment well's drift flux slip
model definition that sets the type of slip model to be used. A slip
model {[}Shi, H., Holmes, J.A., Durlofsky, L. J., Aziz, K., Diaz, L. R.,
Alkaya, B., and Oddie, G. ``Drift-Flux Modeling of Two-Phase Flow in
Wellbores,'' paper SPE 84228, Society of Petroleum Engineers Journal
(2005) 10, No.~1, 24-33; also presented as ``Drift-Flux Modeling of
Multiphase Flow in Wellbores,'' at the SPE Annual Technical Conference
and Exhibition, Denver, Colorado, USA(October 5-8, 2003).{]} and {[}Shi,
H., Holmes, J.A., Diaz, L. R., Durlofsky, L. J., and Aziz, K.
``Drift-Flux Parameters for Three-Phase Steady-State Flow in
Wellbores,'' paper SPE 89836, Society of Petroleum Engineers
Journal(2005) 10, No.~2, 130-137; also presented at the SPE Annual
Technical Conference and Exhibition, Houston, Texas, USA (September
26-29, 2004).{]} enables the different phases in the wellbore to flow at
different velocities, for example gas will flow up the tubing at a
higher velocity than oil and water. The option is activated by the
FLOWOPT parameter on the WELSEGS keyword in the SCHEDULE section;
however, the slip model flow calculation is not available in OPM Flow.

See also the WSEGDFIN keyword that sets the slip model's input
parameters and the WSEGDFPA keyword that allows the model default values
employed by WSEGDFMD to be modified..

This keyword is not supported by OPM Flow but would change the results
if supported so the simulation will be stopped.

\subsection{WSEGDFPA -- Define Multi-Segment Well Drift Flux Slip Model
Parameters}\label{wsegdfpa-define-multi-segment-well-drift-flux-slip-model-parameters}

\begin{longtable}[]{@{}
  >{\raggedright\arraybackslash}p{(\linewidth - 14\tabcolsep) * \real{0.1250}}
  >{\raggedright\arraybackslash}p{(\linewidth - 14\tabcolsep) * \real{0.1250}}
  >{\raggedright\arraybackslash}p{(\linewidth - 14\tabcolsep) * \real{0.1250}}
  >{\raggedright\arraybackslash}p{(\linewidth - 14\tabcolsep) * \real{0.1250}}
  >{\raggedright\arraybackslash}p{(\linewidth - 14\tabcolsep) * \real{0.1250}}
  >{\raggedright\arraybackslash}p{(\linewidth - 14\tabcolsep) * \real{0.1250}}
  >{\raggedright\arraybackslash}p{(\linewidth - 14\tabcolsep) * \real{0.1250}}
  >{\raggedright\arraybackslash}p{(\linewidth - 14\tabcolsep) * \real{0.1250}}@{}}
\toprule\noalign{}
\begin{minipage}[b]{\linewidth}\raggedright
RUNSPEC
\end{minipage} & \begin{minipage}[b]{\linewidth}\raggedright
GRID
\end{minipage} & \begin{minipage}[b]{\linewidth}\raggedright
EDIT
\end{minipage} & \begin{minipage}[b]{\linewidth}\raggedright
PROPS
\end{minipage} & \begin{minipage}[b]{\linewidth}\raggedright
REGIONS
\end{minipage} & \begin{minipage}[b]{\linewidth}\raggedright
SOLUTION
\end{minipage} & \begin{minipage}[b]{\linewidth}\raggedright
SUMMARY
\end{minipage} & \begin{minipage}[b]{\linewidth}\raggedright
SCHEDULE
\end{minipage} \\
\midrule\noalign{}
\endhead
\bottomrule\noalign{}
\endlastfoot
\end{longtable}

\subsubsection{Description}\label{description-1122}

This keyword, WSEGDFPA, enables modification of a multi-segment well's
drift flux slip model default parameters used by the WSEGDFMA keyword in
the SCHEDULE section to define the model. A slip model {[}Shi, H.,
Holmes, J.A., Durlofsky, L. J., Aziz, K., Diaz, L. R., Alkaya, B., and
Oddie, G. ``Drift-Flux Modeling of Two-Phase Flow in Wellbores,'' paper
SPE 84228, Society of Petroleum Engineers Journal (2005) 10, No.~1,
24-33; also presented as ``Drift-Flux Modeling of Multiphase Flow in
Wellbores,'' at the SPE Annual Technical Conference and Exhibition,
Denver, Colorado, USA(October 5-8, 2003).{]} and {[}Shi, H., Holmes,
J.A., Diaz, L. R., Durlofsky, L. J., and Aziz, K. ``Drift-Flux
Parameters for Three-Phase Steady-State Flow in Wellbores,'' paper SPE
89836, Society of Petroleum Engineers Journal(2005) 10, No.~2, 130-137;
also presented at the SPE Annual Technical Conference and Exhibition,
Houston, Texas, USA (September 26-29, 2004).{]} enables the different
phases in the wellbore to flow at different velocities, for example gas
will flow up the tubing at a higher velocity than oil and water. The
option is activated by the FLOWOPT parameter on the WELSEGS keyword in
the SCHEDULE section; however, the slip model flow calculation is not
available in OPM Flow.

See also the WSEGDFIN keyword that sets the slip model's input
parameters and the WSEGDFMD keyword that sets which drift flux slip
model should be used. Note if the WSEGDFPA keyword is used than it must
be placed after the WSEGDFMD keyword, but before the WELSEGS keyword
that defines a multi-segment well. All the aforementioned keywords are
in the SCHEDULE section.

This keyword is not supported by OPM Flow but would change the results
if supported so the simulation will be stopped.

\subsection{WSEGEXSS -- Define Multi-Segment Well Import-Export Segment
Volumes}\label{wsegexss-define-multi-segment-well-import-export-segment-volumes}

\begin{longtable}[]{@{}
  >{\raggedright\arraybackslash}p{(\linewidth - 14\tabcolsep) * \real{0.1250}}
  >{\raggedright\arraybackslash}p{(\linewidth - 14\tabcolsep) * \real{0.1250}}
  >{\raggedright\arraybackslash}p{(\linewidth - 14\tabcolsep) * \real{0.1250}}
  >{\raggedright\arraybackslash}p{(\linewidth - 14\tabcolsep) * \real{0.1250}}
  >{\raggedright\arraybackslash}p{(\linewidth - 14\tabcolsep) * \real{0.1250}}
  >{\raggedright\arraybackslash}p{(\linewidth - 14\tabcolsep) * \real{0.1250}}
  >{\raggedright\arraybackslash}p{(\linewidth - 14\tabcolsep) * \real{0.1250}}
  >{\raggedright\arraybackslash}p{(\linewidth - 14\tabcolsep) * \real{0.1250}}@{}}
\toprule\noalign{}
\begin{minipage}[b]{\linewidth}\raggedright
RUNSPEC
\end{minipage} & \begin{minipage}[b]{\linewidth}\raggedright
GRID
\end{minipage} & \begin{minipage}[b]{\linewidth}\raggedright
EDIT
\end{minipage} & \begin{minipage}[b]{\linewidth}\raggedright
PROPS
\end{minipage} & \begin{minipage}[b]{\linewidth}\raggedright
REGIONS
\end{minipage} & \begin{minipage}[b]{\linewidth}\raggedright
SOLUTION
\end{minipage} & \begin{minipage}[b]{\linewidth}\raggedright
SUMMARY
\end{minipage} & \begin{minipage}[b]{\linewidth}\raggedright
SCHEDULE
\end{minipage} \\
\midrule\noalign{}
\endhead
\bottomrule\noalign{}
\endlastfoot
\end{longtable}

\subsubsection{Description}\label{description-1123}

This keyword, WSEGEXSS, enables the import or export of fluids from a
segment in a multi-segment well. This can be used to, for example, model
gas lift injection for oil wells under artificial lift, or to
approximate the behavior of a down-hole separator. The import-export
fluid volumes can either be expressed as rates or defined as a function
of a segment's pressure value.

This keyword is not supported by OPM Flow but would change the results
if supported so the simulation will be stopped.

\subsection{WSEGFLIM -- Define Multi-Segment Well Artificial Choke
Connections}\label{wsegflim-define-multi-segment-well-artificial-choke-connections}

\begin{longtable}[]{@{}
  >{\raggedright\arraybackslash}p{(\linewidth - 14\tabcolsep) * \real{0.1250}}
  >{\raggedright\arraybackslash}p{(\linewidth - 14\tabcolsep) * \real{0.1250}}
  >{\raggedright\arraybackslash}p{(\linewidth - 14\tabcolsep) * \real{0.1250}}
  >{\raggedright\arraybackslash}p{(\linewidth - 14\tabcolsep) * \real{0.1250}}
  >{\raggedright\arraybackslash}p{(\linewidth - 14\tabcolsep) * \real{0.1250}}
  >{\raggedright\arraybackslash}p{(\linewidth - 14\tabcolsep) * \real{0.1250}}
  >{\raggedright\arraybackslash}p{(\linewidth - 14\tabcolsep) * \real{0.1250}}
  >{\raggedright\arraybackslash}p{(\linewidth - 14\tabcolsep) * \real{0.1250}}@{}}
\toprule\noalign{}
\begin{minipage}[b]{\linewidth}\raggedright
RUNSPEC
\end{minipage} & \begin{minipage}[b]{\linewidth}\raggedright
GRID
\end{minipage} & \begin{minipage}[b]{\linewidth}\raggedright
EDIT
\end{minipage} & \begin{minipage}[b]{\linewidth}\raggedright
PROPS
\end{minipage} & \begin{minipage}[b]{\linewidth}\raggedright
REGIONS
\end{minipage} & \begin{minipage}[b]{\linewidth}\raggedright
SOLUTION
\end{minipage} & \begin{minipage}[b]{\linewidth}\raggedright
SUMMARY
\end{minipage} & \begin{minipage}[b]{\linewidth}\raggedright
SCHEDULE
\end{minipage} \\
\midrule\noalign{}
\endhead
\bottomrule\noalign{}
\endlastfoot
\end{longtable}

\subsubsection{Description}\label{description-1124}

WSEGFLIM enables an artificial choke that chokes a given phase flow rate
for a segment in a multi-segment well. This can be used, for example, to
constraint unwanted production phase through a section of tubing, or to
model a down-hole choke. The keyword provides coefficients that are
applied to the frictional pressure drop across a multi-segment well's
segment in order to inhibit production from that particular zone or
segment. As such, the keyword does not actually model a down-hole choke;
hence, the term artificial.

This keyword is not supported by OPM Flow but would change the results
if supported so the simulation will be stopped.

\subsection{WSEGFMOD -- Define Multi-Segment Well
Model}\label{wsegfmod-define-multi-segment-well-model}

\begin{longtable}[]{@{}
  >{\raggedright\arraybackslash}p{(\linewidth - 14\tabcolsep) * \real{0.1250}}
  >{\raggedright\arraybackslash}p{(\linewidth - 14\tabcolsep) * \real{0.1250}}
  >{\raggedright\arraybackslash}p{(\linewidth - 14\tabcolsep) * \real{0.1250}}
  >{\raggedright\arraybackslash}p{(\linewidth - 14\tabcolsep) * \real{0.1250}}
  >{\raggedright\arraybackslash}p{(\linewidth - 14\tabcolsep) * \real{0.1250}}
  >{\raggedright\arraybackslash}p{(\linewidth - 14\tabcolsep) * \real{0.1250}}
  >{\raggedright\arraybackslash}p{(\linewidth - 14\tabcolsep) * \real{0.1250}}
  >{\raggedright\arraybackslash}p{(\linewidth - 14\tabcolsep) * \real{0.1250}}@{}}
\toprule\noalign{}
\begin{minipage}[b]{\linewidth}\raggedright
RUNSPEC
\end{minipage} & \begin{minipage}[b]{\linewidth}\raggedright
GRID
\end{minipage} & \begin{minipage}[b]{\linewidth}\raggedright
EDIT
\end{minipage} & \begin{minipage}[b]{\linewidth}\raggedright
PROPS
\end{minipage} & \begin{minipage}[b]{\linewidth}\raggedright
REGIONS
\end{minipage} & \begin{minipage}[b]{\linewidth}\raggedright
SOLUTION
\end{minipage} & \begin{minipage}[b]{\linewidth}\raggedright
SUMMARY
\end{minipage} & \begin{minipage}[b]{\linewidth}\raggedright
SCHEDULE
\end{minipage} \\
\midrule\noalign{}
\endhead
\bottomrule\noalign{}
\endlastfoot
\end{longtable}

\subsubsection{Description}\label{description-1125}

The WSEGFMOD declares the multi-phase flow model to be used to calculate
the pressure drop within an individual segment for multi-segment wells.
The FLOWOPT parameter on the WELSEGS keyword in the SCHEDULE section
sets the default multi-segment well model. FLOWOPT is a character string
that can be set to HO that activates the homogeneous model, that is all
phases flow at the same velocity, or DF that invokes the Drift Flux Slip
model (note OPM Flow only supports the default value of HO for the
homogeneous model). Here WSEGFMOD can be used to set the flow model for
a segment to either the homogeneous model or the Drift Flux Slip model,
and addition a: VLP table allocated via the WSEGTABL keyword, or a
specific model as defined by the WSEGVALV, WSEGFLIM and WSEGLABY
keywords. All the aforementioned keywords are in the SCHEDULE section.

This keyword is not supported by OPM Flow but would change the results
if supported so the simulation will be stopped.

\subsection{WSEGINIT -- Define Multi-Segment Well Initial
Conditions}\label{wseginit-define-multi-segment-well-initial-conditions}

\begin{longtable}[]{@{}
  >{\raggedright\arraybackslash}p{(\linewidth - 14\tabcolsep) * \real{0.1250}}
  >{\raggedright\arraybackslash}p{(\linewidth - 14\tabcolsep) * \real{0.1250}}
  >{\raggedright\arraybackslash}p{(\linewidth - 14\tabcolsep) * \real{0.1250}}
  >{\raggedright\arraybackslash}p{(\linewidth - 14\tabcolsep) * \real{0.1250}}
  >{\raggedright\arraybackslash}p{(\linewidth - 14\tabcolsep) * \real{0.1250}}
  >{\raggedright\arraybackslash}p{(\linewidth - 14\tabcolsep) * \real{0.1250}}
  >{\raggedright\arraybackslash}p{(\linewidth - 14\tabcolsep) * \real{0.1250}}
  >{\raggedright\arraybackslash}p{(\linewidth - 14\tabcolsep) * \real{0.1250}}@{}}
\toprule\noalign{}
\begin{minipage}[b]{\linewidth}\raggedright
RUNSPEC
\end{minipage} & \begin{minipage}[b]{\linewidth}\raggedright
GRID
\end{minipage} & \begin{minipage}[b]{\linewidth}\raggedright
EDIT
\end{minipage} & \begin{minipage}[b]{\linewidth}\raggedright
PROPS
\end{minipage} & \begin{minipage}[b]{\linewidth}\raggedright
REGIONS
\end{minipage} & \begin{minipage}[b]{\linewidth}\raggedright
SOLUTION
\end{minipage} & \begin{minipage}[b]{\linewidth}\raggedright
SUMMARY
\end{minipage} & \begin{minipage}[b]{\linewidth}\raggedright
SCHEDULE
\end{minipage} \\
\midrule\noalign{}
\endhead
\bottomrule\noalign{}
\endlastfoot
\end{longtable}

\subsubsection{Description}\label{description-1126}

Normally the simulator calculates the initial conditions for
multi-segment wells, that is the pressure and fluid distributions in
each segment. However, there are occasions when manually setting the
pressures and phase distributions for each segment to investigate
certain flow conditions may be useful. In this case the WSEGINIT keyword
may be used to specify the initial conditions manually. Note that
segments not initialized by this keyword will be automatically
initialized by the simulator

This keyword is not supported by OPM Flow but would change the results
if supported so the simulation will be stopped.

\subsection{WSEGITER -- Define Multi-Segment Wells Iteration
Parameters}\label{wsegiter-define-multi-segment-wells-iteration-parameters}

\begin{longtable}[]{@{}
  >{\raggedright\arraybackslash}p{(\linewidth - 14\tabcolsep) * \real{0.1250}}
  >{\raggedright\arraybackslash}p{(\linewidth - 14\tabcolsep) * \real{0.1250}}
  >{\raggedright\arraybackslash}p{(\linewidth - 14\tabcolsep) * \real{0.1250}}
  >{\raggedright\arraybackslash}p{(\linewidth - 14\tabcolsep) * \real{0.1250}}
  >{\raggedright\arraybackslash}p{(\linewidth - 14\tabcolsep) * \real{0.1250}}
  >{\raggedright\arraybackslash}p{(\linewidth - 14\tabcolsep) * \real{0.1250}}
  >{\raggedright\arraybackslash}p{(\linewidth - 14\tabcolsep) * \real{0.1250}}
  >{\raggedright\arraybackslash}p{(\linewidth - 14\tabcolsep) * \real{0.1250}}@{}}
\toprule\noalign{}
\begin{minipage}[b]{\linewidth}\raggedright
RUNSPEC
\end{minipage} & \begin{minipage}[b]{\linewidth}\raggedright
GRID
\end{minipage} & \begin{minipage}[b]{\linewidth}\raggedright
EDIT
\end{minipage} & \begin{minipage}[b]{\linewidth}\raggedright
PROPS
\end{minipage} & \begin{minipage}[b]{\linewidth}\raggedright
REGIONS
\end{minipage} & \begin{minipage}[b]{\linewidth}\raggedright
SOLUTION
\end{minipage} & \begin{minipage}[b]{\linewidth}\raggedright
SUMMARY
\end{minipage} & \begin{minipage}[b]{\linewidth}\raggedright
SCHEDULE
\end{minipage} \\
\midrule\noalign{}
\endhead
\bottomrule\noalign{}
\endlastfoot
\end{longtable}

\subsubsection{Description}\label{description-1127}

The WSEGITER keyword defines the multi-segment well solution iteration
sequence and solution controls.

OPM Flow uses a different numerical scheme which makes this keyword
redundant; hence, OPM Flow ignores this keyword. It is documented here
for completeness. See section 2.2 Running OPM Flow 2023-04 From The
Command Line on how to set the numerical control parameters for OPM
Flow.

\subsection{WSEGLABY -- Define Multi-Segment Well Labyrinth ICD
Connections}\label{wseglaby-define-multi-segment-well-labyrinth-icd-connections}

\begin{longtable}[]{@{}
  >{\raggedright\arraybackslash}p{(\linewidth - 14\tabcolsep) * \real{0.1250}}
  >{\raggedright\arraybackslash}p{(\linewidth - 14\tabcolsep) * \real{0.1250}}
  >{\raggedright\arraybackslash}p{(\linewidth - 14\tabcolsep) * \real{0.1250}}
  >{\raggedright\arraybackslash}p{(\linewidth - 14\tabcolsep) * \real{0.1250}}
  >{\raggedright\arraybackslash}p{(\linewidth - 14\tabcolsep) * \real{0.1250}}
  >{\raggedright\arraybackslash}p{(\linewidth - 14\tabcolsep) * \real{0.1250}}
  >{\raggedright\arraybackslash}p{(\linewidth - 14\tabcolsep) * \real{0.1250}}
  >{\raggedright\arraybackslash}p{(\linewidth - 14\tabcolsep) * \real{0.1250}}@{}}
\toprule\noalign{}
\begin{minipage}[b]{\linewidth}\raggedright
RUNSPEC
\end{minipage} & \begin{minipage}[b]{\linewidth}\raggedright
GRID
\end{minipage} & \begin{minipage}[b]{\linewidth}\raggedright
EDIT
\end{minipage} & \begin{minipage}[b]{\linewidth}\raggedright
PROPS
\end{minipage} & \begin{minipage}[b]{\linewidth}\raggedright
REGIONS
\end{minipage} & \begin{minipage}[b]{\linewidth}\raggedright
SOLUTION
\end{minipage} & \begin{minipage}[b]{\linewidth}\raggedright
SUMMARY
\end{minipage} & \begin{minipage}[b]{\linewidth}\raggedright
SCHEDULE
\end{minipage} \\
\midrule\noalign{}
\endhead
\bottomrule\noalign{}
\endlastfoot
\end{longtable}

\subsubsection{Description}\label{description-1128}

The WSEGSICD keyword defines a multi-segment well segment to be a
labyrinth Inflow Control Device. (``ICD'') as part of a completion for a
multi-segment well. Note that the well must have been previously define
by the WELSPECS and WELSEGS keywords in the SCHEDULE section and that
the data for the keyword should be repeated for each multi-segment
completion that contains a labyrinth ICD.

An ICD is a well completion component usually installed along the
producing section of a well to minimize the unwanted water and gas
breakthrough in an oil well, or early water production in a gas well,
due to an uneven flow profile over the completed interval. Permeability
variations over the producing interval cause the high permeability zones
to produce higher quantities of fluids than the lower permeability zones
and this uneven producing fluid profile may result in bypassed
hydrocarbons. Secondly, for horizontal wells, the pressure loss from the
``toe'' to the ``heel'' of the well again results in an uneven fluid
profile over the producing interval. In order to rectify this ICDs can
be installed so that the well fluids have to flow through an ICD before
entering the tubing; thus, creating an additional ``designed'' pressure
loss.

A labyrinth ICD is a type of frictional ICD that adds an additional
pressure loss by directing the fluid along a series of channels before
entering the tubing. The channel flow path is designed in such a manner
as to create the desired pressure loss for a given ICD. By placing
various ICD's over the production interval one can design a completion
that results in a more uniform producing fluid profile throughout the
length of the producing interval. Although this type of ICD is not
implemented in OPM Flow, it works in a similar fashion to how a spiral
ICD works. Spiral ICDs are implemented in OPM Flow and the data is
entered via the WSEGSICD keyword in the SCHEDULE section.

This keyword is not supported by OPM Flow but would change the results
if supported so the simulation will be stopped.

\subsection{WSEGLINK -- Define Multi-Segment Well Looped Flow
Paths}\label{wseglink-define-multi-segment-well-looped-flow-paths}

\begin{longtable}[]{@{}
  >{\raggedright\arraybackslash}p{(\linewidth - 14\tabcolsep) * \real{0.1250}}
  >{\raggedright\arraybackslash}p{(\linewidth - 14\tabcolsep) * \real{0.1250}}
  >{\raggedright\arraybackslash}p{(\linewidth - 14\tabcolsep) * \real{0.1250}}
  >{\raggedright\arraybackslash}p{(\linewidth - 14\tabcolsep) * \real{0.1250}}
  >{\raggedright\arraybackslash}p{(\linewidth - 14\tabcolsep) * \real{0.1250}}
  >{\raggedright\arraybackslash}p{(\linewidth - 14\tabcolsep) * \real{0.1250}}
  >{\raggedright\arraybackslash}p{(\linewidth - 14\tabcolsep) * \real{0.1250}}
  >{\raggedright\arraybackslash}p{(\linewidth - 14\tabcolsep) * \real{0.1250}}@{}}
\toprule\noalign{}
\begin{minipage}[b]{\linewidth}\raggedright
RUNSPEC
\end{minipage} & \begin{minipage}[b]{\linewidth}\raggedright
GRID
\end{minipage} & \begin{minipage}[b]{\linewidth}\raggedright
EDIT
\end{minipage} & \begin{minipage}[b]{\linewidth}\raggedright
PROPS
\end{minipage} & \begin{minipage}[b]{\linewidth}\raggedright
REGIONS
\end{minipage} & \begin{minipage}[b]{\linewidth}\raggedright
SOLUTION
\end{minipage} & \begin{minipage}[b]{\linewidth}\raggedright
SUMMARY
\end{minipage} & \begin{minipage}[b]{\linewidth}\raggedright
SCHEDULE
\end{minipage} \\
\midrule\noalign{}
\endhead
\bottomrule\noalign{}
\endlastfoot
\end{longtable}

\subsubsection{Description}\label{description-1129}

This keyword, WSEGLINK, specifies multi-segment well looped flow paths
as part of a completion for a multi-segment well. A looped segment
results in the nodes of the two individual segments that are looped (or
connected) having the same solution pressures and oil, water and gas
flowing rates.

This keyword is not supported by OPM Flow but would change the results
if supported so the simulation will be stopped.

\subsection{WSEGMULT -- Define Multi-Segment Well Frictional Pressure
Loss
Multipliers}\label{wsegmult-define-multi-segment-well-frictional-pressure-loss-multipliers}

\begin{longtable}[]{@{}
  >{\raggedright\arraybackslash}p{(\linewidth - 14\tabcolsep) * \real{0.1250}}
  >{\raggedright\arraybackslash}p{(\linewidth - 14\tabcolsep) * \real{0.1250}}
  >{\raggedright\arraybackslash}p{(\linewidth - 14\tabcolsep) * \real{0.1250}}
  >{\raggedright\arraybackslash}p{(\linewidth - 14\tabcolsep) * \real{0.1250}}
  >{\raggedright\arraybackslash}p{(\linewidth - 14\tabcolsep) * \real{0.1250}}
  >{\raggedright\arraybackslash}p{(\linewidth - 14\tabcolsep) * \real{0.1250}}
  >{\raggedright\arraybackslash}p{(\linewidth - 14\tabcolsep) * \real{0.1250}}
  >{\raggedright\arraybackslash}p{(\linewidth - 14\tabcolsep) * \real{0.1250}}@{}}
\toprule\noalign{}
\begin{minipage}[b]{\linewidth}\raggedright
RUNSPEC
\end{minipage} & \begin{minipage}[b]{\linewidth}\raggedright
GRID
\end{minipage} & \begin{minipage}[b]{\linewidth}\raggedright
EDIT
\end{minipage} & \begin{minipage}[b]{\linewidth}\raggedright
PROPS
\end{minipage} & \begin{minipage}[b]{\linewidth}\raggedright
REGIONS
\end{minipage} & \begin{minipage}[b]{\linewidth}\raggedright
SOLUTION
\end{minipage} & \begin{minipage}[b]{\linewidth}\raggedright
SUMMARY
\end{minipage} & \begin{minipage}[b]{\linewidth}\raggedright
SCHEDULE
\end{minipage} \\
\midrule\noalign{}
\endhead
\bottomrule\noalign{}
\endlastfoot
\end{longtable}

\subsubsection{Description}\label{description-1130}

WSEGMULT supplies a set of constants used to modify (or scale) a
multi-segment well's segment frictional pressure drop between connecting
segments The constants enable either a constant pressure to be applied,
or for the pressure drop to vary as a function of the Gas-Oil Ratio
(``GOR'') or the Water-Oil Ratio (``WOR''). The simulator calculated
pressure drop is multiplied by the following resulting value:

\begin{equation}\phantomsection\label{eq-12-38}{
\mathit{Frictional} \mathit{Loss} \mathit{Multipler} = \mathit{min}\left({x}_{1} + {x}_{2}{(\mathit{WOR})}^{{x}_{3}} + {x}_{4}{\left(\frac{\mathit{GOR}}{{\mathit{GOR}}_{\mathit{min}}}\right)}^{{x}^{5}}, 1.0\right)
}\end{equation}

Where the constants x1 to x5 are defined by the values on this keyword.

This keyword is not supported by OPM Flow but would change the results
if supported so the simulation will be stopped.

\subsection{WSEGPROP -- Modify Multi-Segment Wells and Their Segment
Structure}\label{wsegprop-modify-multi-segment-wells-and-their-segment-structure}

\begin{longtable}[]{@{}
  >{\raggedright\arraybackslash}p{(\linewidth - 14\tabcolsep) * \real{0.1250}}
  >{\raggedright\arraybackslash}p{(\linewidth - 14\tabcolsep) * \real{0.1250}}
  >{\raggedright\arraybackslash}p{(\linewidth - 14\tabcolsep) * \real{0.1250}}
  >{\raggedright\arraybackslash}p{(\linewidth - 14\tabcolsep) * \real{0.1250}}
  >{\raggedright\arraybackslash}p{(\linewidth - 14\tabcolsep) * \real{0.1250}}
  >{\raggedright\arraybackslash}p{(\linewidth - 14\tabcolsep) * \real{0.1250}}
  >{\raggedright\arraybackslash}p{(\linewidth - 14\tabcolsep) * \real{0.1250}}
  >{\raggedright\arraybackslash}p{(\linewidth - 14\tabcolsep) * \real{0.1250}}@{}}
\toprule\noalign{}
\begin{minipage}[b]{\linewidth}\raggedright
RUNSPEC
\end{minipage} & \begin{minipage}[b]{\linewidth}\raggedright
GRID
\end{minipage} & \begin{minipage}[b]{\linewidth}\raggedright
EDIT
\end{minipage} & \begin{minipage}[b]{\linewidth}\raggedright
PROPS
\end{minipage} & \begin{minipage}[b]{\linewidth}\raggedright
REGIONS
\end{minipage} & \begin{minipage}[b]{\linewidth}\raggedright
SOLUTION
\end{minipage} & \begin{minipage}[b]{\linewidth}\raggedright
SUMMARY
\end{minipage} & \begin{minipage}[b]{\linewidth}\raggedright
SCHEDULE
\end{minipage} \\
\midrule\noalign{}
\endhead
\bottomrule\noalign{}
\endlastfoot
\end{longtable}

\subsubsection{Description}\label{description-1131}

The WSEGPROP keyword allows for the editing of exiting multi-segment
wells created by WELSEGS keyword in the SCHEDULE section without having
to re-define all the information that is on the WELSEGS keyword. Note
that the well must have been previously define by both the WELSPECS and
WELSEGS keywords in the SCHEDULE section to use the WSEGPROP keyword.

This keyword is not supported by OPM Flow but would change the results
if supported so the simulation will be stopped.

\begin{longtable}[]{@{}
  >{\raggedright\arraybackslash}p{(\linewidth - 6\tabcolsep) * \real{0.2500}}
  >{\raggedright\arraybackslash}p{(\linewidth - 6\tabcolsep) * \real{0.2500}}
  >{\raggedright\arraybackslash}p{(\linewidth - 6\tabcolsep) * \real{0.2500}}
  >{\raggedright\arraybackslash}p{(\linewidth - 6\tabcolsep) * \real{0.2500}}@{}}
\toprule\noalign{}
\begin{minipage}[b]{\linewidth}\raggedright
No.
\end{minipage} & \begin{minipage}[b]{\linewidth}\raggedright
Name
\end{minipage} & \begin{minipage}[b]{\linewidth}\raggedright
Description
\end{minipage} & \begin{minipage}[b]{\linewidth}\raggedright
Default
\end{minipage} \\
\midrule\noalign{}
\endhead
\bottomrule\noalign{}
\endlastfoot
Field & Metric & Laboratory & \\
1 & WELNAME & A character string of up to eight characters in length
that defines the well name for which a multi-segment well is being
defined. Note that the well name (WELNAME) must have been declared
previously using the WELSPECS keyword in the SCHEDULE section, otherwise
an error may occur. & None \\
2 & ISEG1 & A positive integer greater than or equal to two and less
than or equal to MXSEGS on WSEGDIMS keyword in the RUNSPEC section that
defines the start of a segment & None \\
3 & ISEG2 & A positive integer greater than or equal to two and less
than or equal to ISEG1 on this record and MXSEGS on the WSEGDIMS keyword
in the RUNSPEC section that defines the end of a segment. & None \\
4 & ID & A real positive value that defines the tubing internal diameter
of the segment for the well. & Previous Entered Value \\
feet & m & cm & \\
5 & EPSILON & A real positive value that defines the tubing absolute
roughness of the segment for the well. & Previous Entered Value \\
feet & m & cm & \\
6 & XAREA & XAREA is real positive value equal to or greater than zero
that defines the cross sectional area for fluid flow. Currently this
option is not supported by OPM Flow. & Previous Entered Value \\
ft2 & m2 & cm2 & \\
7 & VOLSEG & VOLSEG is a real positive value that defines the effective
segment volume for the this segment. Currently this option is not
supported by OPM Flow. & Previous Entered Value \\
ft3 & m3 & cm3 & \\
8 & XAREAS & XAREAS is real positive value equal to or greater than zero
that defines the cross sectional area of the pipe wall for this segment,
that is used in thermal conductivity calculations for when the
temperature calculation is activated by the TEMP keyword in the RUNSPEC
section. Currently this option is not supported by OPM Flow. & Previous
Entered Value \\
& ft2 & m2 & cm2 \\
9 & VHEATCAP & VHEATCAP is real positive value equal to or greater than
zero that defines the volumetric heat capacity of the pipe wall used in
thermal conductivity calculations for when the temperature calculation
is activated by the TEMP keyword in the RUNSPEC section. Currently this
option is not supported by OPM Flow. & Previous Entered Value \\
Btu/ft/day/°R & kJ/m/day/K & J/cm/hr/K & \\
10 & THCON & THCON is real positive value equal to or greater than zero
that defines the thermal conductivity of the pipe wall used in thermal
conductivity calculations for when the temperature calculation is
activated by the TEMP keyword in the RUNSPEC section. Currently this
option is not supported by OPM Flow. & Previous Entered Value \\
Btu/ft/day/°R & kJ/m/day/K & J/cm/hr/K & \\
Notes: & & & \\
\end{longtable}

\emph{Table 12.119: WSEGPROP Keyword Description}

The total number of wells should be defined via the WELLDIMS keyword and
the number of multi-segment wells should be declared on the WSEGDIMS
keyword, both keywords are in the RUNSPEC section.

See also the WELSPECS keyword to define wells, the COMPDAT keyword to
define the well completions for both ordinary wells and multi-segment
wells, and the COMPSEGS keyword to define a multi-segment well segment
completions. All the aforementioned keywords are described in the
SCHEDULE section.

\subsubsection{Example}\label{example-601}

The following example modifies two segments in well OP01 and one segment
in well OP02.

\begin{verbatim}
--
--       WELL SEGMENT SPECIFICATION DATA
--
-- WELL  SEG   SEG   TUBE   TUBE     XSEC  VOL
-- NAME  ISTR  IEND  ID     ROUGH    AREA  SEG
WSEGPROP
OPO1      12    14   0.3    0.00010                                            /
OP01      13    15   0.275  0.00010                                            /
OP02      14    14   0.275  0.00010                                            /
/
\end{verbatim}

Note that the two multi-segment wells and their respective segments must
have been previously defined by the WELSEGS keyword.

\subsection{WSEGPULL -- Define a Multi-Segment Well Down-Hole Separator
Pump}\label{wsegpull-define-a-multi-segment-well-down-hole-separator-pump}

\begin{longtable}[]{@{}
  >{\raggedright\arraybackslash}p{(\linewidth - 14\tabcolsep) * \real{0.1250}}
  >{\raggedright\arraybackslash}p{(\linewidth - 14\tabcolsep) * \real{0.1250}}
  >{\raggedright\arraybackslash}p{(\linewidth - 14\tabcolsep) * \real{0.1250}}
  >{\raggedright\arraybackslash}p{(\linewidth - 14\tabcolsep) * \real{0.1250}}
  >{\raggedright\arraybackslash}p{(\linewidth - 14\tabcolsep) * \real{0.1250}}
  >{\raggedright\arraybackslash}p{(\linewidth - 14\tabcolsep) * \real{0.1250}}
  >{\raggedright\arraybackslash}p{(\linewidth - 14\tabcolsep) * \real{0.1250}}
  >{\raggedright\arraybackslash}p{(\linewidth - 14\tabcolsep) * \real{0.1250}}@{}}
\toprule\noalign{}
\begin{minipage}[b]{\linewidth}\raggedright
RUNSPEC
\end{minipage} & \begin{minipage}[b]{\linewidth}\raggedright
GRID
\end{minipage} & \begin{minipage}[b]{\linewidth}\raggedright
EDIT
\end{minipage} & \begin{minipage}[b]{\linewidth}\raggedright
PROPS
\end{minipage} & \begin{minipage}[b]{\linewidth}\raggedright
REGIONS
\end{minipage} & \begin{minipage}[b]{\linewidth}\raggedright
SOLUTION
\end{minipage} & \begin{minipage}[b]{\linewidth}\raggedright
SUMMARY
\end{minipage} & \begin{minipage}[b]{\linewidth}\raggedright
SCHEDULE
\end{minipage} \\
\midrule\noalign{}
\endhead
\bottomrule\noalign{}
\endlastfoot
\end{longtable}

\subsubsection{Description}\label{description-1132}

This keyword, WSEGPULL, specifies a multi-segment well segment to be a
pull-through pump for a down-hole water separator, defined by the
WSEGSEP keyword in the SCHEDULE section, and defines the various
parameters for this type of pump. Down-hole separators are used to
separate water or free gas from the in situ fluid entering the wellbore
in order to increase hydrocarbon recovery.

This keyword is not supported by OPM Flow but would change the results
if supported so the simulation will be stopped.

\subsection{WSEGSEP -- Define a Multi-Segment Well Down-Hole
Separator}\label{wsegsep-define-a-multi-segment-well-down-hole-separator}

\begin{longtable}[]{@{}
  >{\raggedright\arraybackslash}p{(\linewidth - 14\tabcolsep) * \real{0.1250}}
  >{\raggedright\arraybackslash}p{(\linewidth - 14\tabcolsep) * \real{0.1250}}
  >{\raggedright\arraybackslash}p{(\linewidth - 14\tabcolsep) * \real{0.1250}}
  >{\raggedright\arraybackslash}p{(\linewidth - 14\tabcolsep) * \real{0.1250}}
  >{\raggedright\arraybackslash}p{(\linewidth - 14\tabcolsep) * \real{0.1250}}
  >{\raggedright\arraybackslash}p{(\linewidth - 14\tabcolsep) * \real{0.1250}}
  >{\raggedright\arraybackslash}p{(\linewidth - 14\tabcolsep) * \real{0.1250}}
  >{\raggedright\arraybackslash}p{(\linewidth - 14\tabcolsep) * \real{0.1250}}@{}}
\toprule\noalign{}
\begin{minipage}[b]{\linewidth}\raggedright
RUNSPEC
\end{minipage} & \begin{minipage}[b]{\linewidth}\raggedright
GRID
\end{minipage} & \begin{minipage}[b]{\linewidth}\raggedright
EDIT
\end{minipage} & \begin{minipage}[b]{\linewidth}\raggedright
PROPS
\end{minipage} & \begin{minipage}[b]{\linewidth}\raggedright
REGIONS
\end{minipage} & \begin{minipage}[b]{\linewidth}\raggedright
SOLUTION
\end{minipage} & \begin{minipage}[b]{\linewidth}\raggedright
SUMMARY
\end{minipage} & \begin{minipage}[b]{\linewidth}\raggedright
SCHEDULE
\end{minipage} \\
\midrule\noalign{}
\endhead
\bottomrule\noalign{}
\endlastfoot
\end{longtable}

\subsubsection{Description}\label{description-1133}

WSEGSEP specifies a multi-segment well segment to be a down-hole
separator, that enables the separation of fluids down-hole. Down-hole
separators are used to separate water or free gas from the in situ fluid
entering the wellbore in order to increase hydrocarbon recovery. See
also the WSEGPULL keyword in the SCHEDULE section that specifies a
pull-through pump for a down-hole water separator.

This keyword is not supported by OPM Flow but would change the results
if supported so the simulation will be stopped.

\subsection{WSEGSICD -- Define Multi-Segment Well Spiral ICD
Connections}\label{wsegsicd-define-multi-segment-well-spiral-icd-connections}

\begin{longtable}[]{@{}
  >{\raggedright\arraybackslash}p{(\linewidth - 14\tabcolsep) * \real{0.1250}}
  >{\raggedright\arraybackslash}p{(\linewidth - 14\tabcolsep) * \real{0.1250}}
  >{\raggedright\arraybackslash}p{(\linewidth - 14\tabcolsep) * \real{0.1250}}
  >{\raggedright\arraybackslash}p{(\linewidth - 14\tabcolsep) * \real{0.1250}}
  >{\raggedright\arraybackslash}p{(\linewidth - 14\tabcolsep) * \real{0.1250}}
  >{\raggedright\arraybackslash}p{(\linewidth - 14\tabcolsep) * \real{0.1250}}
  >{\raggedright\arraybackslash}p{(\linewidth - 14\tabcolsep) * \real{0.1250}}
  >{\raggedright\arraybackslash}p{(\linewidth - 14\tabcolsep) * \real{0.1250}}@{}}
\toprule\noalign{}
\begin{minipage}[b]{\linewidth}\raggedright
RUNSPEC
\end{minipage} & \begin{minipage}[b]{\linewidth}\raggedright
GRID
\end{minipage} & \begin{minipage}[b]{\linewidth}\raggedright
EDIT
\end{minipage} & \begin{minipage}[b]{\linewidth}\raggedright
PROPS
\end{minipage} & \begin{minipage}[b]{\linewidth}\raggedright
REGIONS
\end{minipage} & \begin{minipage}[b]{\linewidth}\raggedright
SOLUTION
\end{minipage} & \begin{minipage}[b]{\linewidth}\raggedright
SUMMARY
\end{minipage} & \begin{minipage}[b]{\linewidth}\raggedright
SCHEDULE
\end{minipage} \\
\midrule\noalign{}
\endhead
\bottomrule\noalign{}
\endlastfoot
\end{longtable}

\subsubsection{Description}\label{description-1134}

The WSEGSICD keyword defines a multi-segment well segment to be a spiral
Inflow Control Device (``ICD'') as part of a completion for a
multi-segment well. Note that the well must have been previously defined
by the WELSPECS and WELSEGS keywords in the SCHEDULE section and that
the data for the keyword should be repeated for each multi-segment
completion that contains a spiral ICD.

An ICD is a well completion component usually installed along the
producing section of a well to minimize the unwanted water and gas
breakthrough in an oil well, or early water production in a gas well,
due to an uneven flow profile over the completed interval. Permeability
variations over the producing interval cause the high permeability zones
to produce higher quantities of fluids than the lower permeability zones
and this uneven producing fluid profile may result in bypassed
hydrocarbons. Secondly, for horizontal wells, the pressure loss from the
``toe'' to the ``heel'' of the well again results in an uneven fluid
profile over the producing interval. In order to rectify this ICDs can
be installed so that the well fluids have to flow through an ICD before
entering the tubing; thus, creating an additional ``designed'' pressure
loss.

A spiral ICD is a type of frictional ICD that adds an additional
pressure loss by directing the fluid along a ``channel'' (spiral/helix)
before entering the tubing. The channel flow path is designed in such a
manner as to create the desired pressure loss for a given ICD. By
placing various ICD's over the production interval one can design a
completion that results in a more uniform producing fluid profile
throughout the length of the producing interval.

\begin{longtable}[]{@{}
  >{\raggedright\arraybackslash}p{(\linewidth - 6\tabcolsep) * \real{0.2500}}
  >{\raggedright\arraybackslash}p{(\linewidth - 6\tabcolsep) * \real{0.2500}}
  >{\raggedright\arraybackslash}p{(\linewidth - 6\tabcolsep) * \real{0.2500}}
  >{\raggedright\arraybackslash}p{(\linewidth - 6\tabcolsep) * \real{0.2500}}@{}}
\toprule\noalign{}
\begin{minipage}[b]{\linewidth}\raggedright
No.
\end{minipage} & \begin{minipage}[b]{\linewidth}\raggedright
Name
\end{minipage} & \begin{minipage}[b]{\linewidth}\raggedright
Description
\end{minipage} & \begin{minipage}[b]{\linewidth}\raggedright
Default
\end{minipage} \\
\midrule\noalign{}
\endhead
\bottomrule\noalign{}
\endlastfoot
Field & Metric & Laboratory & \\
1 & WELNAME & A character string of up to eight characters in length
that defines the well name for which a multi-segment well is being
defined. Note that the well name (WELNAME) must have been declared
previously using both the WELSPECS and WELSEGS keywords in the SCHEDULE
section, otherwise an error may occur. & None \\
2 & ISEG1 & A positive integer greater than or equal to two and less
than or equal to MXSEGS on the WSEGDIMS keyword in the RUNSPEC section
that defines the start of a segment & None \\
3 & ISEG2 & A positive integer greater than or equal to two and not less
then ISEG1 on this record and less than or equal to MXSEGS on the
WSEGDIMS keyword in the RUNSPEC section, that defines the end of a
segment & None \\
4 & ICDSTREN & A real positive value greater than zero that defines an
empirical constant for the strength of the given ICD as determined from
measurements using the calibrated fluid. & None \\
psia(rft3/day)2 & barsa/(rm3/day)2 & atma/(rcc/hr)2 & \\
5 & ICDLEN & A real value that defines the length of the ICD used in
conjunction with NSCALFAC to calculate a scaling factor to be applied to
the reservoir flow to adjust the flow through each ICD, that is:
NSCALFAC explicitly sets which of the above three options is used. If
NSCALFAC is defaulted, then option 1) is used whenever ICDLEN is
positive and option 2) when ICDLEN is negative. & Defined \\
feet 39.37 & m 12.00 & cm 1,2000 & \\
6 & CALDEN & CALDEN is a real positive value greater than zero that
defines the density of the calibrating fluid at surface conditions. &
Defined \\
lb/ft3 62.416 & kg/m3 1000.25 & gm/cc 1.00025 & \\
7 & CALVISC & CALVISC is a real positive value greater than zero that
defines the viscosity of the calibrating fluid at surface conditions. &
1.45 \\
cP & cP & cP & \\
8 & EMLCRT & EMLCRT is a real positive value greater than zero that
defines the ``local water'' in liquid fraction used to determine whether
the ``water-in-oil'' or ``oil-in-water'' viscosity emulation equation
should be applied. & 0.5 \\
dimensionless & dimensionless & dimensionless & \\
9 & EMLTRANS & EMLTRANS is a real positive value greater than zero that
defines the width of the transition zone around EMLCRT and is used to
ensure that the calculated viscosity forms a continuous function of
water in liquid fraction. Within this region, the emulsion viscosity is
a linear interpolation between the ``water-in-oil'' and ``oil-in-water''
viscosity values either side of the region. & 0.05 \\
dimensionless & dimensionless & dimensionless & \\
10 & EMLMAX & EMLMAX is a real positive value greater than zero that
defines the maximum emulsion viscosity to continuous phase viscosity
(oil or water) ratio. & 5.0 \\
dimensionless & dimensionless & dimensionless & \\
11 & NSCALFAC & NSCALFAC is a positive integer value that is greater
than or equal to zero, that sets the method to be used when applying the
scaling factor and should be set to one of the following: NSCALFAC
explicitly sets which of the above three options is used. If NSCALFAC is
defaulted, then option 1) is used whenever ICDLEN is positive and option
2) when ICDLEN is negative. & -1 \\
dimensionless & dimensionless & dimensionless & \\
12 & CALRATE & A real positive value that defines the maximum surface
flow rate for which the ICD was calibrated. & None \\
scf/d & sm3/day & scc/hour & \\
13 & STATUS & A defined character string of length four that defines the
ICD's operational status, STATUS should be set to one of the following
character strings: & OPEN \\
Notes: & & & \\
\end{longtable}

\emph{Table 12.120: WSEGSICD Keyword Description}

The total number of wells should be defined via the WELLDIMS keyword and
the number of multi-segment wells should be declared on the WSEGDIMS
keyword, both keywords are in the RUNSPEC section. In addition, the
WELSPECS keyword should be used to define wells, the COMPDAT keyword to
define the well completions for both ordinary wells and multi-segment
wells, and the COMPSEGS keyword to define a multi-segment segment
completions. Finally, the WSEGSICD keyword can then be use to define ICD
connections for the well. All the aforementioned keywords are described
in the SCHEDULE section.

The equations used to calculate the pressure drop across the ICD are
given below and illustrate how the pressure reduction is dependent on
the density and viscosity of the fluid flowing through the device.

\begin{equation}\phantomsection\label{eq-12-39}{
\mathrm{Δ}P = {\left(\frac{{\mathrm{ρ}}_{\mathit{calibrated}}}{{\mathrm{ρ}}_{\mathit{mixture}}}⋅\frac{{\mathrm{μ}}_{\mathit{mixture}}}{{\mathrm{μ}}_{\mathit{calibrated}}}\right)}^{\frac{1}{4}}⋅ \frac{{\mathrm{ρ}}_{\mathit{mixture}}}{{\mathrm{ρ}}_{\mathit{calibrated}}} ⋅ \mathrm{β} ⋅ {q}^{2}
}\end{equation}

Where:

ΔP = the pressure drop across the device.

ρmixture = the density of the mixture, as per:

\begin{equation}\phantomsection\label{eq-12-40}{
{\mathrm{ρ}}_{\mathit{mixture}} = \left({\mathrm{α}}_{\mathit{oil}} ⋅ {\mathrm{ρ}}_{\mathit{oil}}\right) + \left({\mathrm{α}}_{\mathit{wat}} ⋅ {\mathrm{ρ}}_{\mathit{wat}}\right) + \left({\mathrm{α}}_{\mathit{gas}} ⋅ {\mathrm{ρ}}_{\mathit{gas}}\right)
}\end{equation}

ρcalibrated = CALDEN, the density of the calibrating fluid at surface
conditions.

μcalibrated = CALVISC, the viscosity of the calibrating fluid at surface
conditions.

μmixture = the viscosity of the mixture, as per:

\begin{equation}\phantomsection\label{eq-12-41}{
{\mathrm{μ}}_{\mathit{mixture}} = \left({\mathrm{α}}_{\mathit{oil}} + {\mathrm{α}}_{\mathit{wat}}\right) ⋅ {\mathrm{μ}}_{\mathit{emulsion}} + \left({\mathrm{α}}_{\mathit{gas}} ⋅ {\mathrm{μ}}_{\mathit{gas}}\right)
}\end{equation}

μemulsion = the viscosity of the oil-water emulsion at local conditions.

β = ICDSTREN, the strength of the ICD as measured using the calibrated
fluid.

q = the local flow rate through the ICD at local conditions adjusted for
scaling

based on ICDLEN and NSCALFAC parameters in Table 12.120.

In equation (12.40), αi represents the volume fraction of oil, water and
gas at local conditions and ρi the density of the three phases.
Similarity for equation (12.41), αi represents the volume fraction of
oil, water and gas at local conditions and μi the viscosity of the three
phases.

See also the WSEGAICD keyword in the SCHEDULE section for autonomous
ICDs, that work in a similar fashion to how sprial ICDs work.

\subsubsection{Example}\label{example-602}

The following example defines one producing well segment oil well (OP01)
using the WELSPECS, WELSEGS COMPDAT and COMPSEGS keywords, followed by
the WSEGSICD keyword to define the spiral inflow control devices for the
well.

\begin{verbatim}
--
--       WELL SPECIFICATION DATA
--
-- WELL  GROUP     LOCATION  BHP    PHASE  DRAIN  INFLOW  OPEN  CROSS  PVT
-- NAME  NAME        I    J  DEPTH  FLUID  AREA   EQUANS  SHUT  FLOW   TABLE
WELSPECS
OP01     PLATFORM   10   10   1*     OIL                                       /
/
--
--       WELL CONNECTION DATA
--
-- WELL  --- LOCATION ---  OPEN   SAT   CONN   WELL   KH    SKIN   D     DIR
-- NAME   II  JJ  K1  K2   SHUT   TAB   FACT   DIA    FACT  FACT   FACT  PEN
COMPDAT
OP01      10  10   1   1   OPEN   1*    200.   0.5                             /
OP01      10  10   2   2   OPEN   1*    200.   0.5                             /
OP01      10  10   3   3   OPEN   1*    200.   0.4                             /
OP01      10  10   4   4   OPEN   1*    200.   0.4                             /
OP01      10  10   5   5   OPEN   1*    200.   0.4                             /
OP01      10  10   6   6   OPEN   1*    200.   0.4                             /

OP01       9  10   2   2   OPEN   1*    200.   0.4                             /
OP01       8  10   2   2   OPEN   1*    200.   0.4                             /
OP01       7  10   2   2   OPEN   1*    200.   0.4                             /
OP01       6  10   2   2   OPEN   1*    200.   0.4                             /
OP01       5  10   2   2   OPEN   1*    200.   0.4                             /

OP01      10   9   3   3   OPEN   1*    200.   0.4                             /
OP01      10   8   3   3   OPEN   1*    200.   0.4                             /
OP01      10   7   3   3   OPEN   1*    200.   0.4                             /
OP01      10   6   3   3   OPEN   1*    200.   0.4                             /
OP01      10   5   3   3   OPEN   1*    200.   0.4                             /

OP01       9  10   5   5   OPEN   1*    200.   0.4                             /
OP01       8  10   5   5   OPEN   1*    200.   0.4                             /
OP01       7  10   5   5   OPEN   1*    200.   0.4                             /
OP01       6  10   5   5   OPEN   1*    200.   0.4                             /
OP01       5  10   5   5   OPEN   1*    200.   0.4                             /

OP01      10   9   6   6   OPEN   1*    200.   0.4                             /
OP01      10   8   6   6   OPEN   1*    200.   0.4                             /
OP01      10   7   6   6   OPEN   1*    200.   0.4                             /
OP01      10   6   6   6   OPEN   1*    200.   0.4                             /
OP01      10   5   6   6   OPEN   1*    200.   0.4                             /
/
--
--       WELL SEGMENT SPECIFICATION DATA
--
-- WELL  NODAL       LEN     WELL    DEPH    PRESS   FLOW
-- NAME  DEPTH       TUBING  VOLM    OPTN    CALC    MODEL
WELSEGS
OP01     2512.5      2512.5  1.0E-5  ABS     HFA     HO                        /
--
--       SEG   SEG   BRAN    SEG     TUBING  NODAL   TUBE   TUBE     XSEC  VOL
--       ISTR  IEND  NO      NO      LENGTH  DEPTH   ID     ROUGH    AREA  SEG
           2     2    1       1      2537.5  2534.5  0.3    0.00010            /
           3     3    1       2      2562.5  2560.5  0.3    0.00010            /
           4     4    1       3      2587.5  2593.5  0.3    0.00010            /
           5     5    1       4      2612.5  2614.5  0.3    0.00010            /
           6     6    1       5      2637.5  2635.5  0.3    0.00010            /

           7     7    2       2      2737.5  2538.5  0.2    0.00010            /
           8     8    2       7      2937.5  2537.5  0.2    0.00010            /
           9     9    2       8      3137.5  2539.5  0.2    0.00010            /
          10    10    2       9      3337.5  2535.5  0.2    0.00010            /
          11    11    2      10      3537.5  2536.5  0.2    0.00010            /

          12    12    3       3      2762.5  2563.5  0.2    0.00010            /
          13    13    3      12      2962.5  2562.5  0.1    0.00010            /
          14    14    3      13      3162.5  2562.5  0.1    0.00010            /
          15    15    3      14      3362.5  2564.5  0.1    0.00010            /
          16    16    3      15      3562.5  2562.5  0.1    0.00010            /

          17    17    4       5      2812.5  2613.5  0.2    0.00010            /
          18    18    4      17      3012.5  2612.5  0.1    0.00010            /
          19    19    4      18      3212.5  2612.5  0.1    0.00010            /
          20    20    4      19      3412.5  2612.5  0.1    0.00010            /
          21    21    4      20      3612.5  2613.5  0.1    0.00010            /

          22    22    5       6      2837.5  2634.5  0.2    0.00010            /
          23    23    5      22      3037.5  2637.5  0.2    0.00010            /
          24    24    5      23      3237.5  2638.5  0.2    0.00010            /
          25    25    5      24      3437.5  2639.5  0.1    0.00010            /
          26    26    5      25      3637.5  2639.5  0.1    0.00010            /
/
--
--       COMPLETION SEGMENT SPECIFICATION DATA
--
-- WELL
-- NAME
COMPSEGS
OP01                                                                           /
--
--      --LOCATION--  BRAN  TUBING  NODAL   DIR  LOC    MID    COMP    ISEG
--       II  JJ  K1   NO    LENGTH  DEPTH   PEN  I,J,K  PERFS  LENGTH  NO.
         10  10   1    1    2512.5  2525.0                                     /
         10  10   2    1    2525.0  2550.0                                     /
         10  10   3    1    2550.0  2575.0                                     /
         10  10   4    1    2575.0  2600.0                                     /
         10  10   5    1    2600.0  2625.0                                     /
         10  10   6    1    2625.0  2650.0                                     /

          9  10   2    2    2637.5  2837.5                                     /
          8  10   2    2    2837.5  3037.5                                     /
          7  10   2    2    3037.5  3237.5                                     /
          6  10   2    2    3237.5  3437.5                                     /
          5  10   2    2    3437.5  3637.5                                     /

         10   9   3    3    2662.5  2862.5                                     /
         10   8   3    3    2862.5  3062.5                                     /
         10   7   3    3    3062.5  3262.5                                     /
         10   6   3    3    3262.5  3462.5                                     /
         10   5   3    3    3462.5  3662.5                                     /

          9  10   5    4    2712.5  2912.5                                     /
          8  10   5    4    2912.5  3112.5                                     /
          7  10   5    4    3112.5  3312.5                                     /
          6  10   5    4    3312.5  3512.5                                     /
          5  10   5    4    3512.5  3712.5                                     /

         10   9   6    5    2737.5  2937.5                                     /
         10   8   6    5    2937.5  3137.5                                     /
         10   7   6    5    3137.5  3337.5                                     /
         10   6   6    5    3337.5  3537.5                                     /
         10   5   6    5    3537.5  3737.5                                     /
--
--       MULTI-SEGMENT WELL ICD SEGMENT SPECIFICATION DATA
--
-- WELL  SEG   SEG   ICD     ICD  CAL  CAL   EML   EML    EML  SCAL  CAL   OPEN
-- NAME  ISTR  IEND  STRNEN  LEN  DEN  VISC  CRIT  TRANS  MAX  FAC   RATE  CLOSE
WSEGSICD
OP01       7    10  0.00025   1*  1.0  0.45  0.50  0.05   5.0   2    1*    OPEN /

OP01      12    15  0.00025   1*  1.0  0.45  0.50  0.05   5.0   2    1*    OPEN /

OP01      17    20  0.00025   1*  1.0  0.45  0.50  0.05   5.0   2    1*    OPEN /

OP01      22    22  0.00025   1*  1.0  0.45  0.50  0.05   5.0   2    1*    OPEN /
OP01      23    23  0.00025   1*  1.0  0.45  0.50  0.05   5.0   2    1*    OPEN /
OP01      24    24  0.00025   1*  1.0  0.45  0.50  0.05   5.0   2    1*    OPEN /
OP01      25    25  0.00050 10.0  1.0  0.45  0.50  0.05   5.0   2    1*    OPEN /
/
\end{verbatim}

Branch number two uses one ICD for segments seven to ten, branch number
three again uses only one ICD for segments 12 to 15 and similarly branch
number three uses one ICD. The fifth and final branch has a total of
four ICDs with the last interval having a 10 foot length. Since NSCALFAC
equals two for the ICDs, then the scale factor is equal to the length of
ICDLEN, divided by the total length of the completions which supply the
ICD. Where ICDLEN is defaulted, the default value of 39.37 ft will be
used for field units and 12 m for SI units.

\subsection{WSEGSOLV -- Define Multi-Segment Well Iterative Linear
Solver
Parameters}\label{wsegsolv-define-multi-segment-well-iterative-linear-solver-parameters}

\begin{longtable}[]{@{}
  >{\raggedright\arraybackslash}p{(\linewidth - 14\tabcolsep) * \real{0.1250}}
  >{\raggedright\arraybackslash}p{(\linewidth - 14\tabcolsep) * \real{0.1250}}
  >{\raggedright\arraybackslash}p{(\linewidth - 14\tabcolsep) * \real{0.1250}}
  >{\raggedright\arraybackslash}p{(\linewidth - 14\tabcolsep) * \real{0.1250}}
  >{\raggedright\arraybackslash}p{(\linewidth - 14\tabcolsep) * \real{0.1250}}
  >{\raggedright\arraybackslash}p{(\linewidth - 14\tabcolsep) * \real{0.1250}}
  >{\raggedright\arraybackslash}p{(\linewidth - 14\tabcolsep) * \real{0.1250}}
  >{\raggedright\arraybackslash}p{(\linewidth - 14\tabcolsep) * \real{0.1250}}@{}}
\toprule\noalign{}
\begin{minipage}[b]{\linewidth}\raggedright
RUNSPEC
\end{minipage} & \begin{minipage}[b]{\linewidth}\raggedright
GRID
\end{minipage} & \begin{minipage}[b]{\linewidth}\raggedright
EDIT
\end{minipage} & \begin{minipage}[b]{\linewidth}\raggedright
PROPS
\end{minipage} & \begin{minipage}[b]{\linewidth}\raggedright
REGIONS
\end{minipage} & \begin{minipage}[b]{\linewidth}\raggedright
SOLUTION
\end{minipage} & \begin{minipage}[b]{\linewidth}\raggedright
SUMMARY
\end{minipage} & \begin{minipage}[b]{\linewidth}\raggedright
SCHEDULE
\end{minipage} \\
\midrule\noalign{}
\endhead
\bottomrule\noalign{}
\endlastfoot
\end{longtable}

\subsubsection{Description}\label{description-1135}

The WSEGSOLV keyword defines the numerical control parameters for the
iterative linear solver for multi-segment well looped flow paths, as
defined by the WSEGLINK keyword in the SCHEDULE section. A looped
segment results in the nodes of the two individual segments that are
looped (or connected) having the same solution pressures and oil, water
and gas flowing rates.

This keyword is not supported by OPM Flow but would change the results
if supported so the simulation will be stopped.

\subsection{WSEGTABL -- Assign Multi-Segment Well VLP Tables to
Segments}\label{wsegtabl-assign-multi-segment-well-vlp-tables-to-segments}

\begin{longtable}[]{@{}
  >{\raggedright\arraybackslash}p{(\linewidth - 14\tabcolsep) * \real{0.1250}}
  >{\raggedright\arraybackslash}p{(\linewidth - 14\tabcolsep) * \real{0.1250}}
  >{\raggedright\arraybackslash}p{(\linewidth - 14\tabcolsep) * \real{0.1250}}
  >{\raggedright\arraybackslash}p{(\linewidth - 14\tabcolsep) * \real{0.1250}}
  >{\raggedright\arraybackslash}p{(\linewidth - 14\tabcolsep) * \real{0.1250}}
  >{\raggedright\arraybackslash}p{(\linewidth - 14\tabcolsep) * \real{0.1250}}
  >{\raggedright\arraybackslash}p{(\linewidth - 14\tabcolsep) * \real{0.1250}}
  >{\raggedright\arraybackslash}p{(\linewidth - 14\tabcolsep) * \real{0.1250}}@{}}
\toprule\noalign{}
\begin{minipage}[b]{\linewidth}\raggedright
RUNSPEC
\end{minipage} & \begin{minipage}[b]{\linewidth}\raggedright
GRID
\end{minipage} & \begin{minipage}[b]{\linewidth}\raggedright
EDIT
\end{minipage} & \begin{minipage}[b]{\linewidth}\raggedright
PROPS
\end{minipage} & \begin{minipage}[b]{\linewidth}\raggedright
REGIONS
\end{minipage} & \begin{minipage}[b]{\linewidth}\raggedright
SOLUTION
\end{minipage} & \begin{minipage}[b]{\linewidth}\raggedright
SUMMARY
\end{minipage} & \begin{minipage}[b]{\linewidth}\raggedright
SCHEDULE
\end{minipage} \\
\midrule\noalign{}
\endhead
\bottomrule\noalign{}
\endlastfoot
\end{longtable}

\subsubsection{Description}\label{description-1136}

WSEGTABL assigns previously defined Vertical Lift Performance (``VLP'')
tables as specified by the VFPPROD keyword in the SCHEDULE section, to
multi-segment well segments, as well as stipulating how the tables are
to be applied.

The FLOWOPT parameter on the WELSEGS keyword in the SCHEDULE section
sets the default multi-segment well model. FLOWOPT either activates the
homogeneous model, that is all phases flow at the same velocity, or the
Drift Flux Slip model. However, the WSEGFMOD keyword in the SCHEDULE
section, can be used to set the flow model for a segment to either the
homogeneous model or the Drift Flux Slip model, and addition a: VLP
table allocated via the WSEGTABL keyword, or a specific model as defined
by the WSEGVALV, WSEGFLIM and WSEGLABY keywords.

This keyword is not supported by OPM Flow but would change the results
if supported so the simulation will be stopped.

\subsection{WSEGVALV -- Define Multi-Segment Well Sub-Critical
Valve}\label{wsegvalv-define-multi-segment-well-sub-critical-valve}

\begin{longtable}[]{@{}
  >{\raggedright\arraybackslash}p{(\linewidth - 14\tabcolsep) * \real{0.1250}}
  >{\raggedright\arraybackslash}p{(\linewidth - 14\tabcolsep) * \real{0.1250}}
  >{\raggedright\arraybackslash}p{(\linewidth - 14\tabcolsep) * \real{0.1250}}
  >{\raggedright\arraybackslash}p{(\linewidth - 14\tabcolsep) * \real{0.1250}}
  >{\raggedright\arraybackslash}p{(\linewidth - 14\tabcolsep) * \real{0.1250}}
  >{\raggedright\arraybackslash}p{(\linewidth - 14\tabcolsep) * \real{0.1250}}
  >{\raggedright\arraybackslash}p{(\linewidth - 14\tabcolsep) * \real{0.1250}}
  >{\raggedright\arraybackslash}p{(\linewidth - 14\tabcolsep) * \real{0.1250}}@{}}
\toprule\noalign{}
\begin{minipage}[b]{\linewidth}\raggedright
RUNSPEC
\end{minipage} & \begin{minipage}[b]{\linewidth}\raggedright
GRID
\end{minipage} & \begin{minipage}[b]{\linewidth}\raggedright
EDIT
\end{minipage} & \begin{minipage}[b]{\linewidth}\raggedright
PROPS
\end{minipage} & \begin{minipage}[b]{\linewidth}\raggedright
REGIONS
\end{minipage} & \begin{minipage}[b]{\linewidth}\raggedright
SOLUTION
\end{minipage} & \begin{minipage}[b]{\linewidth}\raggedright
SUMMARY
\end{minipage} & \begin{minipage}[b]{\linewidth}\raggedright
SCHEDULE
\end{minipage} \\
\midrule\noalign{}
\endhead
\bottomrule\noalign{}
\endlastfoot
\end{longtable}

\subsubsection{Description}\label{description-1137}

The WSEGVALV keyword defines a multi-segment well segment to be a
sub-critical valve Inflow Control Device (``ICD'') as part of a
completion for a multi-segment well. Note that the well must have been
previously defined by the WELSPECS and WELSEGS keywords in the SCHEDULE
section and that the data for the keyword should be repeated for each
multi-segment completion that contains a sub-critical valve ICD.

An ICD is a well completion component usually installed along the
producing section of a well to minimize the unwanted water and gas
breakthrough in an oil well, or early water production in a gas well,
due to an uneven flow profile over the completed interval. Permeability
variations over the producing interval cause the high permeability zones
to produce higher quantities of fluids than the lower permeability zones
and this uneven producing fluid profile may result in bypassed
hydrocarbons. Secondly, for horizontal wells, the pressure loss from the
``toe'' to the ``heel'' of the well again results in an uneven fluid
profile over the producing interval. In order to rectify this ICDs can
be installed so that the well fluids have to flow through an ICD before
entering the tubing; thus, creating an additional ``designed'' pressure
loss.

A sub-critical valve ICD is a type of frictional ICD that adds an
additional pressure loss by directing the fluid through a constriction
before entering the tubing. The pressure drop is a function of the fluid
density, the volumetric flow rate, and the diameter of the constriction.
By placing various ICD's over the production interval one can design a
completion that results in a more uniform producing fluid profile
throughout the length of the producing interval.

\begin{longtable}[]{@{}
  >{\raggedright\arraybackslash}p{(\linewidth - 6\tabcolsep) * \real{0.2500}}
  >{\raggedright\arraybackslash}p{(\linewidth - 6\tabcolsep) * \real{0.2500}}
  >{\raggedright\arraybackslash}p{(\linewidth - 6\tabcolsep) * \real{0.2500}}
  >{\raggedright\arraybackslash}p{(\linewidth - 6\tabcolsep) * \real{0.2500}}@{}}
\toprule\noalign{}
\begin{minipage}[b]{\linewidth}\raggedright
No.
\end{minipage} & \begin{minipage}[b]{\linewidth}\raggedright
Name
\end{minipage} & \begin{minipage}[b]{\linewidth}\raggedright
Description
\end{minipage} & \begin{minipage}[b]{\linewidth}\raggedright
Default
\end{minipage} \\
\midrule\noalign{}
\endhead
\bottomrule\noalign{}
\endlastfoot
Field & Metric & Laboratory & \\
1 & WELNAME & A character string of up to eight characters in length
that defines the well name for which a multi-segment well is being
defined. Note that the well name (WELNAME) must have been declared
previously using both the WELSPECS and WELSEGS keywords in the SCHEDULE
section, otherwise an error may occur. & None \\
2 & ISEG1 & A positive integer greater than or equal to two and less
than or equal to MXSEGS on the WSEGDIMS keyword in the RUNSPEC section
that defines the segment containing the sub-critical valve. & None \\
3 & ICDCV & A real positive value greater than zero that defines the
dimensionless flow coefficient for the valve (Cv). This a vendor
specific value for a given vendor's ICD. & None \\
dimensionless & dimensionless & dimensionless & \\
4 & AREAREST & A real positive value that defines the cross-sectional
area of flow in the restricted section of the valve (Ar) and should have
a minimum value of 1.0 x 10-10. This value may be specified using a User
Defined Argument (UDA). AREAREST is used to convert the segment
volumetric flow rate into the flow velocity at the constriction (υr). &
None \\
ft2 & m2 & cm2 & \\
5 & SEGLEN & A real positive value greater than or equal to zero that
defines the additional pipe length for the frictional pressure drop (L).
If set to zero then there is no additional pressure loss due to
friction, whereas, if set to the default (1*), then the segment pipe
length is calculated from the corresponding WELSEGS keyword. &
Defined \\
ft & m & cm & \\
6 & ID & A real positive value that defines the pipe internal diameter
of the segment used to calculate the pressure drop due to friction (D).
The value is used to replace the segment pipe internal diameter defined
on the WELSEGS keyword in record 2-7, also named ID. If ID is defaulted
on this keyword then the equivalent value on the WELSEGS keyword will be
used instead. Note for non-circular pipe segments use the equivalent
diameter instead, that is:
\(\mathit{Equivalent}\mathit{ID} = \frac{4.0\times (\mathit{Cross}-\mathit{Sectional}\mathit{Area})}{\mathit{Perimeter}}\)
& Defined \\
feet & m & cm & \\
7 & EPSILON & A real positive value that defines the pipe absolute
roughness for this segment. The value is used to replace the segment
pipe absolute roughness defined on the WELSEGS keyword in record 2-8,
also named EPSILON. If EPSILON is defaulted on this keyword then the
equivalent value on the WELSEGS keyword will be used instead. &
Defined \\
feet & m & cm & \\
8 & AREAPIPE & A real positive value that defines the cross-sectional
area of flow in the pipe (Ap), as opposed to the restricted section of
the valve (Ar). AREAPIPE is used to convert the segment volumetric flow
rate into the flow velocity through the pipe (υp). The value is used to
replace the segment pipe cross-sectional area of flow defined on the
WELSEGS keyword in record 2-9, named XAREA. If AREAPIPE is defaulted on
this keyword then the equivalent value on the WELSEGS keyword will be
used instead. & Defined \\
ft2 & m2 & cm2 & \\
9 & STATUS & A character string of length four that defines the ICD's
operational status, STATUS should be set to one of the following
character strings: & OPEN \\
10 & AREAMAX & A real positive value that defines the maximum
cross-sectional area of flow in the restricted section of the valve
(Amax). AREAMAX is used to convert the segment volumetric flow rate into
the maximum flow velocity at the constriction (υr). If defaulted then
AREAPIPE will be used if defined, otherwise XAREA (item 2-9) on the
WELSEGS keyword will be used. & Defined \\
ft2 & m2 & cm2 & \\
Notes: & & & \\
\end{longtable}

\emph{Table 12.121: WSEGVALV Keyword Description}

The total number of wells should be defined via the WELLDIMS keyword and
the number of multi-segment wells should be declared on the WSEGDIMS
keyword, both keywords are in the RUNSPEC section. In addition, the
WELSPECS keyword should be used to define wells, the COMPDAT keyword to
define the well completions for both ordinary wells and multi-segment
wells, and the COMPSEGS keyword to define a multi-segment segment
completions. Finally, the WSEGVALV keyword can then be use to define ICD
connections for the well. All the aforementioned keywords are described
in the SCHEDULE section.

OPM flow calculates the pressure drop across the ICD using a homogeneous
sub-critical flow through a constriction model. The model consists of
two variables, the pressure drop due to the
constriction,\(δ{P}_{\mathit{restriction}}\), and the pressure drop due
to friction,\(δ{P}_{\mathit{friction}}\), as shown in equation (12.42).

\begin{equation}\phantomsection\label{eq-12-42}{
\begin{matrix}\mathrm{δ}P = \mathrm{δ}{P}_{\mathit{restriction}} + \mathrm{δ}{P}_{\mathit{friction}} \\  \\ \mathit{where} \\ \mathrm{δ}{P}_{\mathit{restriction}} = {C}_{1}\frac{\mathrm{ρ}{{\mathrm{υ}}^{2}}_{r}}{{2{C}_{v}}^{2}} \\  \\ \mathrm{δ}{P}_{\mathit{friction}} = 2{C}_{2}f\frac{L}{D}{\mathrm{υ}}_{p}^{2}\end{matrix}
}\end{equation}

Where:

C1 and C2 = Unit conversion constants as listed in Table 12.122.

Cv = Vendor supplied dimensionless flow coefficient for the valve.

D = Diameter of the pipe.

f = Fanning friction factor.

L = Segment pipe length.

ρ = Fluid mixture density.

υr = The flow velocity of the mixture through the constriction.

υp = The flow velocity of the mixture through the segment pipe.

\begin{longtable}[]{@{}llllll@{}}
\toprule\noalign{}
Conversion Factor Constants and Variable Units & & & & & \\
\midrule\noalign{}
\endhead
\bottomrule\noalign{}
\endlastfoot
No. & Name & Field & Metric & Laboratory & Property \\
1 & C1 & 2.159 x 10-4 & 1.0 x 10-5 & 9.869 x 10-7 & Constant \\
2 & C2 & 2.892 x 10-14 & 1.340 x 10-15 & 7.615 x 10-14 & \\
3 & Density & lb/ft3 & kg/m3 & gm/cc & Units \\
4 & Pressure & psia & bars & atm & \\
5 & Velocity & ft/s & m/s & cm/s & \\
6 & Volumetric Flow & ft3/day & m3/day & cc/hour & \\
\end{longtable}

\emph{Table 12.122: WELVALV Conversion Factor Constants and Units}

In addition as both υr and υp are dependent on their respective
cross-section areas then the volumetric flow (q) through the device
requires that:

\begin{equation}\phantomsection\label{eq-12-43}{
q = {\mathrm{υ}}_{r}{A}_{r} = {\mathrm{υ}}_{p}{A}_{p}
}\end{equation}

Substituting equation (12.43) for\(δ{P}_{\mathit{restriction}}\)in
equation (12.42) one obtains:

\begin{equation}\phantomsection\label{eq-12-44}{
\mathrm{δ}{P}_{\mathit{restriction}} = {C}_{2}\frac{\mathrm{ρ}{q}^{2}}{{2{C}_{v}}^{2}{A}_{r}^{2}}
}\end{equation}

Where:

Ar = Cross-sectional area of the constriction.

C2 = Unit conversion constants as listed in Table 12.122.

Cv = Vendor supplied dimensionless flow coefficient for the valve.

ρ = Fluid mixture density.

q = Volumetric flow rate.

The base strength of the device, K is defined using equation (12.44) as
follows:

\begin{equation}\phantomsection\label{eq-12-45}{
K = \frac{ {C}_{2}}{{2{C}_{v}}^{2}{A}_{r}^{2}}
}\end{equation}

Note if K is greater than 0.1 then the device will be shut.

The setting of the device, that is how open the device is, is related to
the restricted area and the maximum restricted area of the device as
follows:

\begin{equation}\phantomsection\label{eq-12-46}{
\mathit{Seting}\mathit{of}\mathit{Device} = \frac{{A}_{r}}{{A}_{\mathit{max}}}
}\end{equation}

Where:

Ar = Cross-sectional area of the constriction.

Amax = Maximum cross-sectional area of the constriction.

\subsubsection{Example}\label{example-603}

The following example is based on one producing well segment oil well
(OP01) using the WELSPECS, WELSEGS COMPDAT and COMPSEGS keywords, as per
the WSEGSICD keyword example (Example), and is therefore not repeated
here.

\begin{verbatim}
--
--       MULTI-SEGMENT WELL ICD SEGMENT SPECIFICATION DATA
--
-- WELL  SEG  DEVICE  AREA   PIPE   PIPE  PIPE  PIPE  OPEN  MAX
-- NAME  NO   CV      REST   LENG   ID    EPSI  AREA  SHUT  AREA
WSEGVALV                                                                                                                                    OP01       7  0.960   0.012  1*     1*    1*    1*    OPEN  1*  / DEFAULT VALUES
OP01      12  0.960   0.012  1*     1*    1*    1*    OPEN  1*  / TAKEN FROM
OP01      17  0.960   0.012  1*     1*    1*    1*    OPEN  1*  / WELSEGS KEYWORD

OP01      22  0.850   0.100  1*     1*    1*    1*    OPEN  1*  /
OP01      23  0.850   0.100  1*     1*    1*    1*    OPEN  1*  /
OP01      24  0.850   0.100  1*     1*    1*    1*    OPEN  1*  /
OP01      25  0.850   0.100  1*     1*    1*    1*    OPEN  1*  /
/
\end{verbatim}

Here segments 7, 12 and 17 have the same type of sub-critical valves
with their pipe properties taken from the WELSEGS keyword used to define
well OP01 as a multi-segment well. Similarly, segments 22 to 25 have the
same ICD properties, and again the pipe properties are taken from the
WELSEGS keyword.

\subsection{WSKPTAB -- Assign Well Polymer Molecular Model Skin
Tables}\label{wskptab-assign-well-polymer-molecular-model-skin-tables}

\begin{longtable}[]{@{}
  >{\raggedright\arraybackslash}p{(\linewidth - 14\tabcolsep) * \real{0.1250}}
  >{\raggedright\arraybackslash}p{(\linewidth - 14\tabcolsep) * \real{0.1250}}
  >{\raggedright\arraybackslash}p{(\linewidth - 14\tabcolsep) * \real{0.1250}}
  >{\raggedright\arraybackslash}p{(\linewidth - 14\tabcolsep) * \real{0.1250}}
  >{\raggedright\arraybackslash}p{(\linewidth - 14\tabcolsep) * \real{0.1250}}
  >{\raggedright\arraybackslash}p{(\linewidth - 14\tabcolsep) * \real{0.1250}}
  >{\raggedright\arraybackslash}p{(\linewidth - 14\tabcolsep) * \real{0.1250}}
  >{\raggedright\arraybackslash}p{(\linewidth - 14\tabcolsep) * \real{0.1250}}@{}}
\toprule\noalign{}
\begin{minipage}[b]{\linewidth}\raggedright
RUNSPEC
\end{minipage} & \begin{minipage}[b]{\linewidth}\raggedright
GRID
\end{minipage} & \begin{minipage}[b]{\linewidth}\raggedright
EDIT
\end{minipage} & \begin{minipage}[b]{\linewidth}\raggedright
PROPS
\end{minipage} & \begin{minipage}[b]{\linewidth}\raggedright
REGIONS
\end{minipage} & \begin{minipage}[b]{\linewidth}\raggedright
SOLUTION
\end{minipage} & \begin{minipage}[b]{\linewidth}\raggedright
SUMMARY
\end{minipage} & \begin{minipage}[b]{\linewidth}\raggedright
SCHEDULE
\end{minipage} \\
\midrule\noalign{}
\endhead
\bottomrule\noalign{}
\endlastfoot
\end{longtable}

\subsubsection{Description}\label{description-1138}

The WSKPTAB keyword assigns the well polymer molecular water and polymer
skin tables to water injection wells in OPM Flow's Polymer Molecular
Weight Transport option, that uses the polymer molecular weight in
calculating the polymer viscosity, as well as accounting for formation
damage due to the water and polymer injection, by adjusting the wellbore
skin pressure. This keyword should only be used if the POLYMER and
POLYMW keywords in the RUNSPEC section are also activated. The keyword
assigns the water SKPRWAT tables, that are defined via the SKPRWAT
keyword in the PROPS section, that are used to calculable the wellbore
skin pressure during water injection. As well as the polymer SKPRPOLY
tables, that are defined via the SKPRPOLY keyword in the PROPS section,
that are used to calculable the wellbore skin pressure during polymer
injection.

\begin{tcolorbox}[enhanced jigsaw, title=\textcolor{quarto-callout-note-color}{\faInfo}\hspace{0.5em}{Note}, coltitle=black, breakable, toprule=.15mm, colbacktitle=quarto-callout-note-color!10!white, leftrule=.75mm, opacitybacktitle=0.6, colback=white, bottomtitle=1mm, toptitle=1mm, titlerule=0mm, left=2mm, colframe=quarto-callout-note-color-frame, arc=.35mm, rightrule=.15mm, bottomrule=.15mm, opacityback=0]

This is an OPM Flow specific keyword that employs an alternative polymer
flood model based on a Polymer Molecular Weight Transport equation, that
is not available in the commercial simulator.

\end{tcolorbox}

\begin{longtable}[]{@{}
  >{\raggedright\arraybackslash}p{(\linewidth - 6\tabcolsep) * \real{0.2500}}
  >{\raggedright\arraybackslash}p{(\linewidth - 6\tabcolsep) * \real{0.2500}}
  >{\raggedright\arraybackslash}p{(\linewidth - 6\tabcolsep) * \real{0.2500}}
  >{\raggedright\arraybackslash}p{(\linewidth - 6\tabcolsep) * \real{0.2500}}@{}}
\toprule\noalign{}
\begin{minipage}[b]{\linewidth}\raggedright
No.
\end{minipage} & \begin{minipage}[b]{\linewidth}\raggedright
Name
\end{minipage} & \begin{minipage}[b]{\linewidth}\raggedright
Description
\end{minipage} & \begin{minipage}[b]{\linewidth}\raggedright
Default
\end{minipage} \\
\midrule\noalign{}
\endhead
\bottomrule\noalign{}
\endlastfoot
Field & Metric & Laboratory & \\
1 & WELNAME & A character string of up to eight characters in length,
that defines the water injection well name, for which the well water
injection skin table, SKKPRWAT, and the polymer skin injection table,
SKPRPOLY, are to be assigned. Note that the well name (WELNAME) must
have been declared previously using the WELSPECS keyword in the SCHEDULE
section, otherwise an error may occur. & None \\
2 & SKPRWAT & A positive integer value that defines the corresponding
SKPRWAT table to be allocated to the water injection well. A value less
than or equal to zero means that no SKPRWAT table is allocated to the
well & 0 \\
3 & SKPRPOLY & A positive integer value that defines the corresponding
SKPRPOLY table to be allocated to the water injection well. A value less
than or equal to zero means that no SKPRPOLY table is allocated to the
well & 0 \\
Notes: & & & \\
\end{longtable}

\emph{Table 12.123: WSKPTAB Keyword Description}

See also the SKPRWAT and SKPRPOLY keywords in the PROPS section, that
describe the relationship of the wellbore skin pressure with respect to
the injected water and polymer throughput and fluid velocity, for the
simulator's Polymer Molecular Weight Transport option. As well as the
PLYMWINJ and PLYVMH keywords, also in the PROPS section, that are the
additional keywords required for the Polymer Molecular Weight Transport
option.

The WPMITAB keyword in the SCHEDULE section may be used to assign the
PLYMWINJ keyword in the PROPS section. Where PLYMWINJ, describes the
relationship of the injected polymer molecular weight as a function of
polymer throughput and polymer velocity, for the simulator's Polymer
Molecular Weight Transport option

Note that the standard polymer property data keywords: PLYROCK, PLYADS,
PLYMAX, etc., are still required to fully describe the polymer fluid.

\subsubsection{Example}\label{example-604}

Given NTSKWAT equals two and NTSKPOLY equals three on the PINTDIMS
keyword in the RUNSPEC section, then:

\begin{verbatim}
--
--       ASSIGN WELL POLYMER MOLECULAR MODEL SKIN TABLES
--
-- WELL  SKPRWAT   SKPRPOLY
-- NAME  TABLE     TABLE
WSKPTAB
WI01     1         1                                              /
WI02     1         3                                              /
WI03     2         2                                              /
/
\end{verbatim}

Assigns SKPRWAT table one to wells WI01 and WI02 and table two to WI03,
and SKPRPOLY tables one, two and three to wells WI01, WI03, and WI02,
respectively.

\subsection{WSOLVENT -- Define Gas Injection Well Solvent
Fraction}\label{wsolvent-define-gas-injection-well-solvent-fraction}

\begin{longtable}[]{@{}
  >{\raggedright\arraybackslash}p{(\linewidth - 14\tabcolsep) * \real{0.1250}}
  >{\raggedright\arraybackslash}p{(\linewidth - 14\tabcolsep) * \real{0.1250}}
  >{\raggedright\arraybackslash}p{(\linewidth - 14\tabcolsep) * \real{0.1250}}
  >{\raggedright\arraybackslash}p{(\linewidth - 14\tabcolsep) * \real{0.1250}}
  >{\raggedright\arraybackslash}p{(\linewidth - 14\tabcolsep) * \real{0.1250}}
  >{\raggedright\arraybackslash}p{(\linewidth - 14\tabcolsep) * \real{0.1250}}
  >{\raggedright\arraybackslash}p{(\linewidth - 14\tabcolsep) * \real{0.1250}}
  >{\raggedright\arraybackslash}p{(\linewidth - 14\tabcolsep) * \real{0.1250}}@{}}
\toprule\noalign{}
\begin{minipage}[b]{\linewidth}\raggedright
RUNSPEC
\end{minipage} & \begin{minipage}[b]{\linewidth}\raggedright
GRID
\end{minipage} & \begin{minipage}[b]{\linewidth}\raggedright
EDIT
\end{minipage} & \begin{minipage}[b]{\linewidth}\raggedright
PROPS
\end{minipage} & \begin{minipage}[b]{\linewidth}\raggedright
REGIONS
\end{minipage} & \begin{minipage}[b]{\linewidth}\raggedright
SOLUTION
\end{minipage} & \begin{minipage}[b]{\linewidth}\raggedright
SUMMARY
\end{minipage} & \begin{minipage}[b]{\linewidth}\raggedright
SCHEDULE
\end{minipage} \\
\midrule\noalign{}
\endhead
\bottomrule\noalign{}
\endlastfoot
\end{longtable}

\subsubsection{Description}\label{description-1139}

WSOLVENT defines a gas injection well's solvent fraction in the
injection stream that is to be used when the Solvent option has been
activated by the SOLVENT keyword in the RUNSPEC section.

\begin{longtable}[]{@{}
  >{\raggedright\arraybackslash}p{(\linewidth - 6\tabcolsep) * \real{0.2500}}
  >{\raggedright\arraybackslash}p{(\linewidth - 6\tabcolsep) * \real{0.2500}}
  >{\raggedright\arraybackslash}p{(\linewidth - 6\tabcolsep) * \real{0.2500}}
  >{\raggedright\arraybackslash}p{(\linewidth - 6\tabcolsep) * \real{0.2500}}@{}}
\toprule\noalign{}
\begin{minipage}[b]{\linewidth}\raggedright
No.
\end{minipage} & \begin{minipage}[b]{\linewidth}\raggedright
Name
\end{minipage} & \begin{minipage}[b]{\linewidth}\raggedright
Description
\end{minipage} & \begin{minipage}[b]{\linewidth}\raggedright
Default
\end{minipage} \\
\midrule\noalign{}
\endhead
\bottomrule\noalign{}
\endlastfoot
Field & Metric & Laboratory & \\
1 & WELNAME & A character string of up to eight characters in length
that defines the well name of a gas injection well for which the solvent
fraction data is being defined. Note that the well name (WELNAME) must
have been declared previously using the WELSPECS keyword in the SCHEDULE
section, otherwise an error may occur. & None \\
4 & SOLFRA & A real positive value greater than or equal to zero and
less than or equal to one that defines the fraction of solvent in the
gas well's injection stream. This value may be specified using a User
Defined Argument (UDA). & None \\
fraction & fraction & fraction & \\
Notes: & & & \\
\end{longtable}

\emph{Table 12.124: WSOLVENT Keyword Description}

Gas injection wells that are not declared via this keyword have their
solvent fractions set to zero.

See also the GCONINJE keyword to define a group's injection targets and
constraints, and the WCONINJE keyword to define an injection well's
targets and constraints. All the aforementioned keywords are described
in the SCHEDULE section.

\subsubsection{Example}\label{example-605}

The following example defines the solvent fractions for three gas
injection wells for when the solvent option has been activated by the
SOLVENT keyword in the RUNSPEC section.

\begin{verbatim}
--
--       DEFINE GAS INJECTION WELL SOLVENT FRACTION
--
-- WELL  SOLVENT
-- NAME  FRACTION
--       --------
WSOLVENT
GI01     0.0000                                            /
GI02     0.5000                                            /
GI03     0.5000                                            /
/
\end{verbatim}

The solvent fraction for the GI01 gas injector is set to zero and both
GI02 and GI03 gas injectors have solvent fraction values of 0.5 for
their injection streams.

\subsection{WSURFACT -- Define Water Injection Well Surfactant
Concentration}\label{wsurfact-define-water-injection-well-surfactant-concentration}

\begin{longtable}[]{@{}
  >{\raggedright\arraybackslash}p{(\linewidth - 14\tabcolsep) * \real{0.1250}}
  >{\raggedright\arraybackslash}p{(\linewidth - 14\tabcolsep) * \real{0.1250}}
  >{\raggedright\arraybackslash}p{(\linewidth - 14\tabcolsep) * \real{0.1250}}
  >{\raggedright\arraybackslash}p{(\linewidth - 14\tabcolsep) * \real{0.1250}}
  >{\raggedright\arraybackslash}p{(\linewidth - 14\tabcolsep) * \real{0.1250}}
  >{\raggedright\arraybackslash}p{(\linewidth - 14\tabcolsep) * \real{0.1250}}
  >{\raggedright\arraybackslash}p{(\linewidth - 14\tabcolsep) * \real{0.1250}}
  >{\raggedright\arraybackslash}p{(\linewidth - 14\tabcolsep) * \real{0.1250}}@{}}
\toprule\noalign{}
\begin{minipage}[b]{\linewidth}\raggedright
RUNSPEC
\end{minipage} & \begin{minipage}[b]{\linewidth}\raggedright
GRID
\end{minipage} & \begin{minipage}[b]{\linewidth}\raggedright
EDIT
\end{minipage} & \begin{minipage}[b]{\linewidth}\raggedright
PROPS
\end{minipage} & \begin{minipage}[b]{\linewidth}\raggedright
REGIONS
\end{minipage} & \begin{minipage}[b]{\linewidth}\raggedright
SOLUTION
\end{minipage} & \begin{minipage}[b]{\linewidth}\raggedright
SUMMARY
\end{minipage} & \begin{minipage}[b]{\linewidth}\raggedright
SCHEDULE
\end{minipage} \\
\midrule\noalign{}
\endhead
\bottomrule\noalign{}
\endlastfoot
\end{longtable}

\subsubsection{Description}\label{description-1140}

WSURFACT defines a water injection well's surfactant concentration in
the injection stream that is to be used when the Surfactant phase has
been activated by either the SURFACT or SURFACTW keywords in the RUNSPEC
section.

This keyword is not supported by OPM Flow but would change the results
if supported so the simulation will be stopped.

\begin{longtable}[]{@{}
  >{\raggedright\arraybackslash}p{(\linewidth - 6\tabcolsep) * \real{0.2500}}
  >{\raggedright\arraybackslash}p{(\linewidth - 6\tabcolsep) * \real{0.2500}}
  >{\raggedright\arraybackslash}p{(\linewidth - 6\tabcolsep) * \real{0.2500}}
  >{\raggedright\arraybackslash}p{(\linewidth - 6\tabcolsep) * \real{0.2500}}@{}}
\toprule\noalign{}
\begin{minipage}[b]{\linewidth}\raggedright
No.
\end{minipage} & \begin{minipage}[b]{\linewidth}\raggedright
Name
\end{minipage} & \begin{minipage}[b]{\linewidth}\raggedright
Description
\end{minipage} & \begin{minipage}[b]{\linewidth}\raggedright
Default
\end{minipage} \\
\midrule\noalign{}
\endhead
\bottomrule\noalign{}
\endlastfoot
Field & Metric & Laboratory & \\
1 & WELNAME & A character string of up to eight characters in length
that defines the well name of a gas injection well for which the solvent
fraction data is being defined. Note that the well name (WELNAME) must
have been declared previously using the WELSPECS keyword in the SCHEDULE
section, otherwise an error may occur. & None \\
2 & SURCON & A real positive value that defines the surfactant
concentration of the well's injection stream. & None \\
lb/stb & kg/sm3 & gm/scc & \\
Notes: & & & \\
\end{longtable}

\emph{Table 12.125: WSURFACT Keyword Description}

Water injection wells that are not declared via this keyword have their
surfactant concentrations set to zero.

See also the GCONINJE keyword to define a group's injection targets and
constraints, and the WCONINJE keyword to define an injection well's
targets and constraints. All the aforementioned keywords are described
in the SCHEDULE section.

\subsubsection{Example}\label{example-606}

The following example defines the surfactant concentrations for three
water injection wells for when the surfactant phase option has been
activated by either the SURFACT or SURFACTW keywords in the RUNSPEC
section. Here the surfactant concentration has been set to 0.200 for all
three wells.

\begin{verbatim}
--
--       DEFINE WATER INJECTION WELL SURFACTANT CONCENTRATION
--
-- WELL  SURFACT
-- NAME  SURCON
--       --------
WSURFACT
WI01     0.2000                                            /
WI02     0.2000                                            /
WI03     0.2000                                            /
/
\end{verbatim}

\subsection{WTADD -- Add a Constant to a Well Target or
Constraint}\label{wtadd-add-a-constant-to-a-well-target-or-constraint}

\begin{longtable}[]{@{}
  >{\raggedright\arraybackslash}p{(\linewidth - 14\tabcolsep) * \real{0.1250}}
  >{\raggedright\arraybackslash}p{(\linewidth - 14\tabcolsep) * \real{0.1250}}
  >{\raggedright\arraybackslash}p{(\linewidth - 14\tabcolsep) * \real{0.1250}}
  >{\raggedright\arraybackslash}p{(\linewidth - 14\tabcolsep) * \real{0.1250}}
  >{\raggedright\arraybackslash}p{(\linewidth - 14\tabcolsep) * \real{0.1250}}
  >{\raggedright\arraybackslash}p{(\linewidth - 14\tabcolsep) * \real{0.1250}}
  >{\raggedright\arraybackslash}p{(\linewidth - 14\tabcolsep) * \real{0.1250}}
  >{\raggedright\arraybackslash}p{(\linewidth - 14\tabcolsep) * \real{0.1250}}@{}}
\toprule\noalign{}
\begin{minipage}[b]{\linewidth}\raggedright
RUNSPEC
\end{minipage} & \begin{minipage}[b]{\linewidth}\raggedright
GRID
\end{minipage} & \begin{minipage}[b]{\linewidth}\raggedright
EDIT
\end{minipage} & \begin{minipage}[b]{\linewidth}\raggedright
PROPS
\end{minipage} & \begin{minipage}[b]{\linewidth}\raggedright
REGIONS
\end{minipage} & \begin{minipage}[b]{\linewidth}\raggedright
SOLUTION
\end{minipage} & \begin{minipage}[b]{\linewidth}\raggedright
SUMMARY
\end{minipage} & \begin{minipage}[b]{\linewidth}\raggedright
SCHEDULE
\end{minipage} \\
\midrule\noalign{}
\endhead
\bottomrule\noalign{}
\endlastfoot
\end{longtable}

\subsubsection{Description}\label{description-1141}

This keyword, WTADD, adds a constant to a previously define well's
target or constraint, as stated on the WCONPROD, WCONINJE, or WELTARG
keywords, but not for the history matching wells using the WCONHIST or
WCONINJH keywords. All the aforementioned keywords are in the SCHEDULE
section. The constant can be positive or negative.

This keyword is not supported by OPM Flow but would change the results
if supported so the simulation will be stopped.

\subsection{WTEMP -- Define An Injection Well's Fluid
Temperature}\label{wtemp-define-an-injection-wells-fluid-temperature}

\begin{longtable}[]{@{}
  >{\raggedright\arraybackslash}p{(\linewidth - 14\tabcolsep) * \real{0.1250}}
  >{\raggedright\arraybackslash}p{(\linewidth - 14\tabcolsep) * \real{0.1250}}
  >{\raggedright\arraybackslash}p{(\linewidth - 14\tabcolsep) * \real{0.1250}}
  >{\raggedright\arraybackslash}p{(\linewidth - 14\tabcolsep) * \real{0.1250}}
  >{\raggedright\arraybackslash}p{(\linewidth - 14\tabcolsep) * \real{0.1250}}
  >{\raggedright\arraybackslash}p{(\linewidth - 14\tabcolsep) * \real{0.1250}}
  >{\raggedright\arraybackslash}p{(\linewidth - 14\tabcolsep) * \real{0.1250}}
  >{\raggedright\arraybackslash}p{(\linewidth - 14\tabcolsep) * \real{0.1250}}@{}}
\toprule\noalign{}
\begin{minipage}[b]{\linewidth}\raggedright
RUNSPEC
\end{minipage} & \begin{minipage}[b]{\linewidth}\raggedright
GRID
\end{minipage} & \begin{minipage}[b]{\linewidth}\raggedright
EDIT
\end{minipage} & \begin{minipage}[b]{\linewidth}\raggedright
PROPS
\end{minipage} & \begin{minipage}[b]{\linewidth}\raggedright
REGIONS
\end{minipage} & \begin{minipage}[b]{\linewidth}\raggedright
SOLUTION
\end{minipage} & \begin{minipage}[b]{\linewidth}\raggedright
SUMMARY
\end{minipage} & \begin{minipage}[b]{\linewidth}\raggedright
SCHEDULE
\end{minipage} \\
\midrule\noalign{}
\endhead
\bottomrule\noalign{}
\endlastfoot
\end{longtable}

\subsubsection{Description}\label{description-1142}

The WTEMP keyword defines the temperature of the injection fluid being
injected by an injection well.

This keyword can only be used if OPM Flow's thermal option has been
activated by the THERMAL keyword in the RUNSPEC section. Note this is
different to the commercial simulator that uses the TEMP keyword in the
RUNSPEC section to activate the black-oil thermal model.

\begin{longtable}[]{@{}
  >{\raggedright\arraybackslash}p{(\linewidth - 6\tabcolsep) * \real{0.2500}}
  >{\raggedright\arraybackslash}p{(\linewidth - 6\tabcolsep) * \real{0.2500}}
  >{\raggedright\arraybackslash}p{(\linewidth - 6\tabcolsep) * \real{0.2500}}
  >{\raggedright\arraybackslash}p{(\linewidth - 6\tabcolsep) * \real{0.2500}}@{}}
\toprule\noalign{}
\begin{minipage}[b]{\linewidth}\raggedright
No.
\end{minipage} & \begin{minipage}[b]{\linewidth}\raggedright
Name
\end{minipage} & \begin{minipage}[b]{\linewidth}\raggedright
Description
\end{minipage} & \begin{minipage}[b]{\linewidth}\raggedright
Default
\end{minipage} \\
\midrule\noalign{}
\endhead
\bottomrule\noalign{}
\endlastfoot
Field & Metric & Laboratory & \\
1 & WELNAME & A character string of up to eight characters in length
that defines the well name for an injection well for which the injection
well fluid's temperature data is being defined. Note that the well name
(WELNAME) must have been declared previously using the WELSPECS keyword
in the SCHEDULE section, otherwise an error may occur. & None \\
2 & TEMP & A real positive value greater than zero that defines the
temperature of the injected fluid. & None \\
oF & oC & oC & \\
Notes: & & & \\
\end{longtable}

\emph{Table 12.126: WTEMP Keyword Description}

See also the GCONINJE keyword to define a group's injection targets and
constraints, and the WCONINJE keyword to define an injection well's
targets and constraints. All the aforementioned keywords are described
in the SCHEDULE section.

\subsubsection{Example}\label{example-607}

The following example defines the injected fluid temperatures for three
water injection wells for when the thermal option has been activated by
the THERMAL keyword in the RUNSPEC section.

\begin{verbatim}
--
--       DEFINE INJECTION WELL FLUID TEMPERATURE
--
-- WELL  FLUID
-- NAME  TEMP.
--       --------
WTEMP
WI01     39.00                                             /
WI02     37.00                                             /
WI03     39.00                                             /
/
\end{verbatim}

Here wells WI01 and WI03 inject water with a water temperature of 39 oF
and well WI02's injection water temperature is 37 oF.

\subsection{WTEMPQ -- Output Well Names and Well Lists to the Print
File}\label{wtempq-output-well-names-and-well-lists-to-the-print-file}

\begin{longtable}[]{@{}
  >{\raggedright\arraybackslash}p{(\linewidth - 14\tabcolsep) * \real{0.1250}}
  >{\raggedright\arraybackslash}p{(\linewidth - 14\tabcolsep) * \real{0.1250}}
  >{\raggedright\arraybackslash}p{(\linewidth - 14\tabcolsep) * \real{0.1250}}
  >{\raggedright\arraybackslash}p{(\linewidth - 14\tabcolsep) * \real{0.1250}}
  >{\raggedright\arraybackslash}p{(\linewidth - 14\tabcolsep) * \real{0.1250}}
  >{\raggedright\arraybackslash}p{(\linewidth - 14\tabcolsep) * \real{0.1250}}
  >{\raggedright\arraybackslash}p{(\linewidth - 14\tabcolsep) * \real{0.1250}}
  >{\raggedright\arraybackslash}p{(\linewidth - 14\tabcolsep) * \real{0.1250}}@{}}
\toprule\noalign{}
\begin{minipage}[b]{\linewidth}\raggedright
RUNSPEC
\end{minipage} & \begin{minipage}[b]{\linewidth}\raggedright
GRID
\end{minipage} & \begin{minipage}[b]{\linewidth}\raggedright
EDIT
\end{minipage} & \begin{minipage}[b]{\linewidth}\raggedright
PROPS
\end{minipage} & \begin{minipage}[b]{\linewidth}\raggedright
REGIONS
\end{minipage} & \begin{minipage}[b]{\linewidth}\raggedright
SOLUTION
\end{minipage} & \begin{minipage}[b]{\linewidth}\raggedright
SUMMARY
\end{minipage} & \begin{minipage}[b]{\linewidth}\raggedright
SCHEDULE
\end{minipage} \\
\midrule\noalign{}
\endhead
\bottomrule\noalign{}
\endlastfoot
\end{longtable}

\subsubsection{Description}\label{description-1143}

The WTEMPQ prints out a user defined selected list of currently defined
wells and well lists to the print file (\emph{.PRT). The keyword allows
for sub-setting the well names etc., using the normal well and well list
naming conventions. For example to list all wells beginning with the
characters ``OP'' then one would use ``OP}'' as the well name on this
keyword.

This keyword is not supported by OPM Flow but would change the results
if supported so the simulation will be stopped.

\subsection{WTEST -- Well Testing Criteria for Re-Opening Closed
Wells}\label{wtest-well-testing-criteria-for-re-opening-closed-wells}

\begin{longtable}[]{@{}
  >{\raggedright\arraybackslash}p{(\linewidth - 14\tabcolsep) * \real{0.1250}}
  >{\raggedright\arraybackslash}p{(\linewidth - 14\tabcolsep) * \real{0.1250}}
  >{\raggedright\arraybackslash}p{(\linewidth - 14\tabcolsep) * \real{0.1250}}
  >{\raggedright\arraybackslash}p{(\linewidth - 14\tabcolsep) * \real{0.1250}}
  >{\raggedright\arraybackslash}p{(\linewidth - 14\tabcolsep) * \real{0.1250}}
  >{\raggedright\arraybackslash}p{(\linewidth - 14\tabcolsep) * \real{0.1250}}
  >{\raggedright\arraybackslash}p{(\linewidth - 14\tabcolsep) * \real{0.1250}}
  >{\raggedright\arraybackslash}p{(\linewidth - 14\tabcolsep) * \real{0.1250}}@{}}
\toprule\noalign{}
\begin{minipage}[b]{\linewidth}\raggedright
RUNSPEC
\end{minipage} & \begin{minipage}[b]{\linewidth}\raggedright
GRID
\end{minipage} & \begin{minipage}[b]{\linewidth}\raggedright
EDIT
\end{minipage} & \begin{minipage}[b]{\linewidth}\raggedright
PROPS
\end{minipage} & \begin{minipage}[b]{\linewidth}\raggedright
REGIONS
\end{minipage} & \begin{minipage}[b]{\linewidth}\raggedright
SOLUTION
\end{minipage} & \begin{minipage}[b]{\linewidth}\raggedright
SUMMARY
\end{minipage} & \begin{minipage}[b]{\linewidth}\raggedright
SCHEDULE
\end{minipage} \\
\midrule\noalign{}
\endhead
\bottomrule\noalign{}
\endlastfoot
\end{longtable}

\subsubsection{Description}\label{description-1144}

The WTEST keyword outlines the testing procedures to be applied to wells
that are closed for various reasons to see if the wells are capable
flowing under the current operating conditions. The keyword can be
applied to single wells or groups of wells.

\begin{longtable}[]{@{}
  >{\raggedright\arraybackslash}p{(\linewidth - 6\tabcolsep) * \real{0.2500}}
  >{\raggedright\arraybackslash}p{(\linewidth - 6\tabcolsep) * \real{0.2500}}
  >{\raggedright\arraybackslash}p{(\linewidth - 6\tabcolsep) * \real{0.2500}}
  >{\raggedright\arraybackslash}p{(\linewidth - 6\tabcolsep) * \real{0.2500}}@{}}
\toprule\noalign{}
\begin{minipage}[b]{\linewidth}\raggedright
No.
\end{minipage} & \begin{minipage}[b]{\linewidth}\raggedright
Name
\end{minipage} & \begin{minipage}[b]{\linewidth}\raggedright
Description
\end{minipage} & \begin{minipage}[b]{\linewidth}\raggedright
Default
\end{minipage} \\
\midrule\noalign{}
\endhead
\bottomrule\noalign{}
\endlastfoot
Field & Metric & Laboratory & \\
1 & WELNAME & A character string of up to eight characters in length
that defines the well name of the well to be tested. Note that the well
name (WELNAME) must have been declared previously using the WELSPECS
keyword in the SCHEDULE section, otherwise an error may occur. & None \\
2 & TIME & A real value greater than zero that defines the minimum
period of time that must elapse between well tests. The next well test
is performed at the beginning of the next time step after the specified
period of time has elapsed since the previous well test, for example if
TIME is set equal to 365.25 (days), the test is performed approximately
every year. & None \\
days & days & hours & \\
3 & TEST & A character string of up to five characters in length that
defines a list of reasons why a well could have been closed. If a well
was closed for one of the specified reasons then the well is tested to
see if it can be put back on production. The characters that can be used
to define TEST are as follows: Here, if the well has at least one
connection open then the well is opened. If a well has no open
connections, then those connections that have been closed due to an
economic limit test (``ELT''), are all individually tested and reopened
if they meet the ELT criteria as defined by the WECON, CECON and
WECONINJ keywords in the SCHEDULE section. And the well is opened if it
has at least one open connection. Where the ELT applies to ORAT, GRAT,
WCUT, GOR and WGR etc. The default value is an empty string `` '' that
switches off testing. Note that only the P, E and G options are
currently supported in OPM Flow. & `` '' \\
4 & NTIME & A positive integer greater than or equal to zero that define
the number of times a well can be tested. The default value of zero
means an infinite number of times. & 0 \\
5 & START & A real positive value that defines the start up time used to
prorate the rate at which the well is put back on production. If START
is large compared to the time step size, then the well is brought on
gradually, if it is less than the time step size, then the well is
opened faster. The rate is prorated based on the following:
\({Q}_{i} = {Q}_{\mathit{final}}\times \left(\frac{{T}_{\mathit{End}\mathit{of}\mathit{Time}\mathit{Step}} - {T}_{\mathit{opened}}}{\mathit{START}}\right)\)
The default value of 0.0 means the well is opened immediately. & 0.0 \\
days & days & hours & \\
Notes: & & & \\
\end{longtable}

\emph{Table 12.127: WTEST Keyword Description}

See also the WELSPECS keyword to define a wells shut-in or stop options,
WECON for setting a well's economic criteria, GCONPROD and GCONINJE for
group controls, and GECON for setting a group's economic criteria. All
the aforementioned keywords are described in the SCHEDULE section.

\subsubsection{Example}\label{example-608}

The following example defines test criteria for all gas wells ('GP*')
and three oil wells (OP01, OP02, and OP03).

\begin{verbatim}
--
--
--       WELL TESTING CRITERIA FOR RE-OPENING CLOSED WELLS
--
-- WELL  TST    TST    NO.    STRT
-- NAME  INTV   TYPE   TSTS   TIME
-- ----  ----   ----   ----   ----
WTEST
‘GP*’    365.25 P      5      0.0
OP01     30.0   PEG    0      0.0                                              /
OP02     30.0   PEG    0      0.0                                              /
OPO3     30.0   PEG    0      0.0                                              /
/
\end{verbatim}

All the gas wells are test annually if they have been shut-in due to a
bottom-hole or tubing head pressure limit, are tested five times after
they have been closed, and are opened up immediately. The oil wells are
tested every 30 days if they have been closed due to a bottom-hole or
tubing head pressure limit, a well economic limit or a group economic
limit. All the oil wells are tested an infinite amount of times and are
opened up immediately. Note that only the P, E and G options are
currently supported in OPM Flow.

\subsection{WTHPMAX -- Define a Well's Maximum Flowing THP for
Shut-In}\label{wthpmax-define-a-wells-maximum-flowing-thp-for-shut-in}

\begin{longtable}[]{@{}
  >{\raggedright\arraybackslash}p{(\linewidth - 14\tabcolsep) * \real{0.1250}}
  >{\raggedright\arraybackslash}p{(\linewidth - 14\tabcolsep) * \real{0.1250}}
  >{\raggedright\arraybackslash}p{(\linewidth - 14\tabcolsep) * \real{0.1250}}
  >{\raggedright\arraybackslash}p{(\linewidth - 14\tabcolsep) * \real{0.1250}}
  >{\raggedright\arraybackslash}p{(\linewidth - 14\tabcolsep) * \real{0.1250}}
  >{\raggedright\arraybackslash}p{(\linewidth - 14\tabcolsep) * \real{0.1250}}
  >{\raggedright\arraybackslash}p{(\linewidth - 14\tabcolsep) * \real{0.1250}}
  >{\raggedright\arraybackslash}p{(\linewidth - 14\tabcolsep) * \real{0.1250}}@{}}
\toprule\noalign{}
\begin{minipage}[b]{\linewidth}\raggedright
RUNSPEC
\end{minipage} & \begin{minipage}[b]{\linewidth}\raggedright
GRID
\end{minipage} & \begin{minipage}[b]{\linewidth}\raggedright
EDIT
\end{minipage} & \begin{minipage}[b]{\linewidth}\raggedright
PROPS
\end{minipage} & \begin{minipage}[b]{\linewidth}\raggedright
REGIONS
\end{minipage} & \begin{minipage}[b]{\linewidth}\raggedright
SOLUTION
\end{minipage} & \begin{minipage}[b]{\linewidth}\raggedright
SUMMARY
\end{minipage} & \begin{minipage}[b]{\linewidth}\raggedright
SCHEDULE
\end{minipage} \\
\midrule\noalign{}
\endhead
\bottomrule\noalign{}
\endlastfoot
\end{longtable}

\subsubsection{Description}\label{description-1145}

WTHPMAX stipulates a well's maximum flowing Tubing Head Pressure
(``THP''), above which the well will be shut-in. The facility is useful
if the THP exceeds the wellhead maximum design pressure, which can occur
if excessive gas invades the wellbore. In addition to setting the
maximum THP, the keyword defines the criteria for re-testing the well to
see if the THP has fallen below the maximum value.

This keyword is not supported by OPM Flow but would change the results
if supported so the simulation will be stopped.

\subsection{WTMULT -- Multiply a Well Target or Constraint by a
Constant}\label{wtmult-multiply-a-well-target-or-constraint-by-a-constant}

\begin{longtable}[]{@{}
  >{\raggedright\arraybackslash}p{(\linewidth - 14\tabcolsep) * \real{0.1250}}
  >{\raggedright\arraybackslash}p{(\linewidth - 14\tabcolsep) * \real{0.1250}}
  >{\raggedright\arraybackslash}p{(\linewidth - 14\tabcolsep) * \real{0.1250}}
  >{\raggedright\arraybackslash}p{(\linewidth - 14\tabcolsep) * \real{0.1250}}
  >{\raggedright\arraybackslash}p{(\linewidth - 14\tabcolsep) * \real{0.1250}}
  >{\raggedright\arraybackslash}p{(\linewidth - 14\tabcolsep) * \real{0.1250}}
  >{\raggedright\arraybackslash}p{(\linewidth - 14\tabcolsep) * \real{0.1250}}
  >{\raggedright\arraybackslash}p{(\linewidth - 14\tabcolsep) * \real{0.1250}}@{}}
\toprule\noalign{}
\begin{minipage}[b]{\linewidth}\raggedright
RUNSPEC
\end{minipage} & \begin{minipage}[b]{\linewidth}\raggedright
GRID
\end{minipage} & \begin{minipage}[b]{\linewidth}\raggedright
EDIT
\end{minipage} & \begin{minipage}[b]{\linewidth}\raggedright
PROPS
\end{minipage} & \begin{minipage}[b]{\linewidth}\raggedright
REGIONS
\end{minipage} & \begin{minipage}[b]{\linewidth}\raggedright
SOLUTION
\end{minipage} & \begin{minipage}[b]{\linewidth}\raggedright
SUMMARY
\end{minipage} & \begin{minipage}[b]{\linewidth}\raggedright
SCHEDULE
\end{minipage} \\
\midrule\noalign{}
\endhead
\bottomrule\noalign{}
\endlastfoot
\end{longtable}

\subsubsection{Description}\label{description-1146}

This keyword, WTMULT, multiplies a defined well's target or constraint
by a constant, for the target and constraints previously stipulated on
the WCONPROD, WCONINJE, or WELTARG keywords, but not for the history
matching wells using the WCONHIST or WCONINJH keywords. All the
aforementioned keywords are in the SCHEDULE section. The constant should
be positive value.

\begin{longtable}[]{@{}
  >{\raggedright\arraybackslash}p{(\linewidth - 6\tabcolsep) * \real{0.2500}}
  >{\raggedright\arraybackslash}p{(\linewidth - 6\tabcolsep) * \real{0.2500}}
  >{\raggedright\arraybackslash}p{(\linewidth - 6\tabcolsep) * \real{0.2500}}
  >{\raggedright\arraybackslash}p{(\linewidth - 6\tabcolsep) * \real{0.2500}}@{}}
\toprule\noalign{}
\begin{minipage}[b]{\linewidth}\raggedright
No.
\end{minipage} & \begin{minipage}[b]{\linewidth}\raggedright
Name
\end{minipage} & \begin{minipage}[b]{\linewidth}\raggedright
Description
\end{minipage} & \begin{minipage}[b]{\linewidth}\raggedright
Default
\end{minipage} \\
\midrule\noalign{}
\endhead
\bottomrule\noalign{}
\endlastfoot
Field & Metric & Laboratory & \\
1 & WELNAME & A character string of up to eight characters in length
that defines the well name for which the well target or constraint
(CONTROL) is being adjusted by the multipler (FACTOR). Note that the
well name (WELNAME) must have been declared previously using the
WELSPECS keyword in the SCHEDULE section, otherwise an error may occur.
& None \\
2 & CONTROL & A defined character string that declares the well target
or constraint that will be adjusted by multiplying its current value by
the FACTOR defined in (3). CONTROL should be set to one of the following
character strings: Note that the CRAT, CVAL and the NGL options are not
available in OPM Flow and should not be used. & None \\
3 & FACTOR & A postive real value that defines the FACTOR that is used
to multiply the CONTROL specified in (2). This value may be specified
using a User Defined Argument (UDA). & None \\
dimensionless & dimensionless & dimensionless & \\
4 & NTIME & A positive integer greater than or equal to one that defines
the number of report time steps for which the well target or constraint
(CONTROL) is multiplied by the FACTOR. This is only applied when the
FACTOR is specified by a User Defined Argument (UDA). The default value
of one means that the multiplication is applied only at the current time
step for the UDA variable. This option is not supported by OPM Flow. &
1 \\
dimensionless & dimensionless & dimensionless & \\
Notes: & & & \\
\end{longtable}

\emph{Table 12.128: WTMULT Keyword Description}

See also the WELTARG and WELCNTL keywords in the SCHEDULE section, that
can be used to reset the well's target and constraints of both rates and
pressures, as well as the well's control mode.

\subsubsection{Example}\label{example-609}

The example shows three oil wells having the flow streams adjusted.

\begin{verbatim}
--
-- WELL TARGET/LIMIT MULTIPLIER
--
--  WELL    WELL   MULT    REPORT
--  NAME    CNTL   FACTOR  TIMES
WTMULT
OP01        ORAT   0.90                /
OP02        BHP    0.95                /
OP03        LIFT   1.25                /
/
\end{verbatim}

Well OP01 has its current oil rate target multiplied by 0.90, well OP02
has its bottom-hole pressure constraint multiplied by 0.95, and well
OP03 has its artificial lift quantity increased by 1.25 times.

\subsection{WTRACER -- Define An Injection Well's Tracer
Concentration}\label{wtracer-define-an-injection-wells-tracer-concentration}

\begin{longtable}[]{@{}
  >{\raggedright\arraybackslash}p{(\linewidth - 14\tabcolsep) * \real{0.1250}}
  >{\raggedright\arraybackslash}p{(\linewidth - 14\tabcolsep) * \real{0.1250}}
  >{\raggedright\arraybackslash}p{(\linewidth - 14\tabcolsep) * \real{0.1250}}
  >{\raggedright\arraybackslash}p{(\linewidth - 14\tabcolsep) * \real{0.1250}}
  >{\raggedright\arraybackslash}p{(\linewidth - 14\tabcolsep) * \real{0.1250}}
  >{\raggedright\arraybackslash}p{(\linewidth - 14\tabcolsep) * \real{0.1250}}
  >{\raggedright\arraybackslash}p{(\linewidth - 14\tabcolsep) * \real{0.1250}}
  >{\raggedright\arraybackslash}p{(\linewidth - 14\tabcolsep) * \real{0.1250}}@{}}
\toprule\noalign{}
\begin{minipage}[b]{\linewidth}\raggedright
RUNSPEC
\end{minipage} & \begin{minipage}[b]{\linewidth}\raggedright
GRID
\end{minipage} & \begin{minipage}[b]{\linewidth}\raggedright
EDIT
\end{minipage} & \begin{minipage}[b]{\linewidth}\raggedright
PROPS
\end{minipage} & \begin{minipage}[b]{\linewidth}\raggedright
REGIONS
\end{minipage} & \begin{minipage}[b]{\linewidth}\raggedright
SOLUTION
\end{minipage} & \begin{minipage}[b]{\linewidth}\raggedright
SUMMARY
\end{minipage} & \begin{minipage}[b]{\linewidth}\raggedright
SCHEDULE
\end{minipage} \\
\midrule\noalign{}
\endhead
\bottomrule\noalign{}
\endlastfoot
\end{longtable}

\subsubsection{Description}\label{description-1147}

The WTRACER keyword defines the tracer concentration of the injection
fluid being injected by an injection well. This keyword should only be
used if the Tracer option has been invoked by the TRACERS keyword in the
RUNSPEC section and the tracers have be declared via the TRACER keyword
in the PROPS section.

\begin{longtable}[]{@{}
  >{\raggedright\arraybackslash}p{(\linewidth - 6\tabcolsep) * \real{0.2500}}
  >{\raggedright\arraybackslash}p{(\linewidth - 6\tabcolsep) * \real{0.2500}}
  >{\raggedright\arraybackslash}p{(\linewidth - 6\tabcolsep) * \real{0.2500}}
  >{\raggedright\arraybackslash}p{(\linewidth - 6\tabcolsep) * \real{0.2500}}@{}}
\toprule\noalign{}
\begin{minipage}[b]{\linewidth}\raggedright
No.
\end{minipage} & \begin{minipage}[b]{\linewidth}\raggedright
Name
\end{minipage} & \begin{minipage}[b]{\linewidth}\raggedright
Description
\end{minipage} & \begin{minipage}[b]{\linewidth}\raggedright
Default
\end{minipage} \\
\midrule\noalign{}
\endhead
\bottomrule\noalign{}
\endlastfoot
Field & Metric & Laboratory & \\
1 & WELNAME & A character string of up to eight characters in length
that defines the well name of an injection well for which the tracer
fraction data is being defined. Note that the well name (WELNAME) must
have been declared previously using the WELSPECS keyword in the SCHEDULE
section, otherwise an error may occur. & None \\
2 & NAME & A three letter character string defining the tracer's name
which has previously been defined via the TRACER keyword in the PROPS
section. Note it is best to avoid names beginning with the letters F, S,
and T as these names may create naming issues in post-processing
software. & None \\
3 & TRCON & A real positive value that defines the tracer concentration
of the well's injection stream. This value may be specified using a User
Defined Argument (UDA). & None \\
Liquid: 1/stb Gas: 1/Mscf & Liquid: 1/sm3 Gas: 1/sm3 & Liquid: 1/scc
Gas: 1/scc & \\
4 & TRCUM & A real positive value that defines the cumulative tracer
concentration factor of the well's injection stream. This value may be
specified using a User Defined Argument (UDA). This feature is currently
not supported by OPM Flow. & None \\
Liquid: 1/stb Gas: 1/Mscf & Liquid: 1/sm3 Gas: 1/sm3 & Liquid: 1/scc
Gas: 1/scc & \\
5 & GRPNAME & A character string of up to eight characters in length
that defines the group from which the produced tracer concentration
should be used for the well's injection stream. GRPNAME must have been
previously defined via the GCONPROD keyword in the SCHEDULE section,
unless the FIELD group has been specified here. Note if GRPNAME is not
defined then TRCON will be used for the tracer concentration of the
well's injection stream. This feature is currently not supported by OPM
Flow. & None \\
Notes: & & & \\
\end{longtable}

\emph{Table 12.129: WTRACER Keyword Description}

Injection wells that are not declared via this keyword have their tracer
concentrations set to zero.

See also the GCONINJE keyword to define a group's injection targets and
constraints, and the WCONINJE keyword to define an injection well's
targets and constraints. All the aforementioned keywords are described
in the SCHEDULE section.

\subsubsection{Example}\label{example-610}

The following example defines the tracer concentrations for two gas
injectors and three water injection wells, with water injection well
WI02 having no `WAT' tracer injected in the water phase.

\begin{verbatim}
--
--       DEFINE CONCENTRATION OF TRACERS IN THE INJECTION STREAMS,
--       INJECTION TRACER CONCENTRATIONS NOT DEFINED USING THE WTRACER
--       KEYWORD ARE ASSUMED TO BE ZERO.
--
-- WELL  NAME    TRACER  TRACER  TRACER
-- NAME  TRACER  CONC    CUM     GROUP
WTRACER
GI01     'GAS'   1.0                                       /
GI02     'GAS'   1.0                                       /

WI01     'WAT'   1.0                                       /
WI02     'WAT'   0.0                                       /
WI03     'WAT'   1.0                                       /
/
\end{verbatim}

Note the terminating ``/'' for the keyword.

\subsection{WVFPDP -- Modify Well BHP Obtained from VFP
Tables}\label{wvfpdp-modify-well-bhp-obtained-from-vfp-tables}

\begin{longtable}[]{@{}
  >{\raggedright\arraybackslash}p{(\linewidth - 14\tabcolsep) * \real{0.1250}}
  >{\raggedright\arraybackslash}p{(\linewidth - 14\tabcolsep) * \real{0.1250}}
  >{\raggedright\arraybackslash}p{(\linewidth - 14\tabcolsep) * \real{0.1250}}
  >{\raggedright\arraybackslash}p{(\linewidth - 14\tabcolsep) * \real{0.1250}}
  >{\raggedright\arraybackslash}p{(\linewidth - 14\tabcolsep) * \real{0.1250}}
  >{\raggedright\arraybackslash}p{(\linewidth - 14\tabcolsep) * \real{0.1250}}
  >{\raggedright\arraybackslash}p{(\linewidth - 14\tabcolsep) * \real{0.1250}}
  >{\raggedright\arraybackslash}p{(\linewidth - 14\tabcolsep) * \real{0.1250}}@{}}
\toprule\noalign{}
\begin{minipage}[b]{\linewidth}\raggedright
RUNSPEC
\end{minipage} & \begin{minipage}[b]{\linewidth}\raggedright
GRID
\end{minipage} & \begin{minipage}[b]{\linewidth}\raggedright
EDIT
\end{minipage} & \begin{minipage}[b]{\linewidth}\raggedright
PROPS
\end{minipage} & \begin{minipage}[b]{\linewidth}\raggedright
REGIONS
\end{minipage} & \begin{minipage}[b]{\linewidth}\raggedright
SOLUTION
\end{minipage} & \begin{minipage}[b]{\linewidth}\raggedright
SUMMARY
\end{minipage} & \begin{minipage}[b]{\linewidth}\raggedright
SCHEDULE
\end{minipage} \\
\midrule\noalign{}
\endhead
\bottomrule\noalign{}
\endlastfoot
\end{longtable}

\subsubsection{Description}\label{description-1148}

The WVFPDP keyword modifies a well's Bottom-Hole Pressure (``BHP'')
estimated by the simulator by interpolation of the Vertical Flow
Performance (``VFP'') tables. The production VFP tables are entered via
the VFPPROD keyword and the injection tables by the VFPINJ keyword; both
keywords are in the SCHEDULE section.

Note that simulator automatically adjusts the interpolated BHP to
account for hydrostatic head using the density of the wellbore fluid and
the difference between a well's BHP reference depth, as per the BHPREF
parameter on the WELSPECS or WELSPECL keywords in the SCHEDULE section,
and the VFPREF parameter reference depth on the VFPPROD and VFPINJ
keywords. Thus, WVFPDP applies an additional adjustment in order to
match a well's flow rate to a given tubing head pressure, by adjusting
the BHP.

\begin{longtable}[]{@{}
  >{\raggedright\arraybackslash}p{(\linewidth - 6\tabcolsep) * \real{0.2500}}
  >{\raggedright\arraybackslash}p{(\linewidth - 6\tabcolsep) * \real{0.2500}}
  >{\raggedright\arraybackslash}p{(\linewidth - 6\tabcolsep) * \real{0.2500}}
  >{\raggedright\arraybackslash}p{(\linewidth - 6\tabcolsep) * \real{0.2500}}@{}}
\toprule\noalign{}
\begin{minipage}[b]{\linewidth}\raggedright
No.
\end{minipage} & \begin{minipage}[b]{\linewidth}\raggedright
Name
\end{minipage} & \begin{minipage}[b]{\linewidth}\raggedright
Description
\end{minipage} & \begin{minipage}[b]{\linewidth}\raggedright
Default
\end{minipage} \\
\midrule\noalign{}
\endhead
\bottomrule\noalign{}
\endlastfoot
Field & Metric & Laboratory & \\
1 & WELNAME & A character string of up to eight characters in length
that defines the well name for which VFP interpolated BHP adjustment is
to be applied. Note that the well name (WELNAME) must have been declared
previously using the WELSPECS and WCONPROD (or WCONINJE) keywords in the
SCHEDULE section, otherwise an error may occur. & None \\
2 & DELTAP & A real positive or negative value that is added to the VFP
interpolated BHP value (BHPVFP). A positive value of DELTAP increases
the BHP and therefore makes a production well less productive; whereas,
a negative value is subtracted from the BHP and therefore increases the
productivity of a production well. Consequently, the opposite effect
occurs for injection wells, that is, a positive value of BHPVFP
increases the BHP and therefore increases an injection well's
injectivity; whereas, a negative value is subtracted from the BHP and
therefore decreases the injectivity of an injection well. & 0.0 \\
psia & barsa & atma & \\
3 & MULTP & MULTP is a real positive or negative value that scales the
tubing pressure loss by the following equation;
\({\mathit{BHP}}_{\mathit{Adjusted}} = \mathit{THP} + \mathit{MULTP}\left({\mathit{BHP}}_{\mathit{VFP}} - \mathit{THP}\right)\)
Thus, a MULTP value greater then 1.0 increases the BHP and therefore
makes a production well less productive; whereas, a value less than 1.0
increases the productivity of a production well. Consequently, the
opposite effect occurs for injection wells & 1.0 \\
dimensionless & dimensionless & dimensionless & \\
Notes: & & & \\
\end{longtable}

\emph{Table 12.130: WVFPDP Keyword Description}

\subsubsection{Example}\label{example-611}

The following example below shows three oils operating under THP
control.

\begin{verbatim}
-- ------------------------------------------------------------------------------
-- 01 JAN 2000 START OF SCHEDULE SECTION
-- ------------------------------------------------------------------------------
--
--       WELL PRODUCTION WELL CONTROLS
--
-- WELL  OPEN/  CNTL   OIL    WAT    GAS   LIQ    RES    BHP    THP   VFP    VFP
-- NAME  SHUT   MODE   RATE   RATE   RATE  RATE   RATE   PRES   PRES  TABLE  ALFQ
WCONPROD
OP01     OPEN   THP    1*     1*     1*    5000   1*     750.0  500.  9      1* /
OP02     OPEN   THP    1*     1*     1*    5000   1*     750.0  500.  9      1* /
OP03     OPEN   THP    1*     1*     1*    5000   1*     750.0  500.  9      1* /
/
--
--       WELL VFP BHP-THP CORRECTION DATA
--
--  WELL BHP     BHP
--  NAME DELTAP  MULTP
WVFPDP
OP01     20.O    1*                                         /
OP01     -5.O    1*                                         /
OP01      0.O    1.10                                       /
/
\end{verbatim}

Well OP01 has a delta pressure correction of 20 psia applied to it's BHP
resulting in a reduction in the well's productivity for the given 500.0
psia THP operating target. For well OP02, the well's productivity is
increased by subtracting 5.0 psia from the BHP. And finally for well
OP03, the MULTP value of 1.10 decreases the well's productivity by
increasing the pressure loss between the THP and BHP by 10\%.

\subsection{WVFPEXP -- Define Well VFP Interpolation
Options}\label{wvfpexp-define-well-vfp-interpolation-options}

\begin{longtable}[]{@{}
  >{\raggedright\arraybackslash}p{(\linewidth - 14\tabcolsep) * \real{0.1250}}
  >{\raggedright\arraybackslash}p{(\linewidth - 14\tabcolsep) * \real{0.1250}}
  >{\raggedright\arraybackslash}p{(\linewidth - 14\tabcolsep) * \real{0.1250}}
  >{\raggedright\arraybackslash}p{(\linewidth - 14\tabcolsep) * \real{0.1250}}
  >{\raggedright\arraybackslash}p{(\linewidth - 14\tabcolsep) * \real{0.1250}}
  >{\raggedright\arraybackslash}p{(\linewidth - 14\tabcolsep) * \real{0.1250}}
  >{\raggedright\arraybackslash}p{(\linewidth - 14\tabcolsep) * \real{0.1250}}
  >{\raggedright\arraybackslash}p{(\linewidth - 14\tabcolsep) * \real{0.1250}}@{}}
\toprule\noalign{}
\begin{minipage}[b]{\linewidth}\raggedright
RUNSPEC
\end{minipage} & \begin{minipage}[b]{\linewidth}\raggedright
GRID
\end{minipage} & \begin{minipage}[b]{\linewidth}\raggedright
EDIT
\end{minipage} & \begin{minipage}[b]{\linewidth}\raggedright
PROPS
\end{minipage} & \begin{minipage}[b]{\linewidth}\raggedright
REGIONS
\end{minipage} & \begin{minipage}[b]{\linewidth}\raggedright
SOLUTION
\end{minipage} & \begin{minipage}[b]{\linewidth}\raggedright
SUMMARY
\end{minipage} & \begin{minipage}[b]{\linewidth}\raggedright
SCHEDULE
\end{minipage} \\
\midrule\noalign{}
\endhead
\bottomrule\noalign{}
\endlastfoot
\end{longtable}

\subsubsection{Description}\label{description-1149}

This keyword, WVFPEXP, defines how Vertical Flow Performance (``VFP'')
tables are interpolated and can be used to resolve certain issues with
wells operating under tubing head pressure control. For example, setting
the VFP table to interpolate explicitly, that is using the previous time
step results of the gas and water ratios for an oil well, may improve
convergence. The default is to use implicit interpolation that uses the
current time step values and may result in solution convergence
oscillations in solving the linear equations. The WCONPROD keyword is
used to allocate the VFPPROD tables to specific production wells. Note
that one VFP table can be allocated to one or more wells; however,
WVFPEXP is applied to a well's allocated VFP table, not to all wells
that use the same table, unless specifically requested. All the
aforementioned keywords are in the SCHEDULE section.

\begin{longtable}[]{@{}
  >{\raggedright\arraybackslash}p{(\linewidth - 6\tabcolsep) * \real{0.2500}}
  >{\raggedright\arraybackslash}p{(\linewidth - 6\tabcolsep) * \real{0.2500}}
  >{\raggedright\arraybackslash}p{(\linewidth - 6\tabcolsep) * \real{0.2500}}
  >{\raggedright\arraybackslash}p{(\linewidth - 6\tabcolsep) * \real{0.2500}}@{}}
\toprule\noalign{}
\begin{minipage}[b]{\linewidth}\raggedright
No.
\end{minipage} & \begin{minipage}[b]{\linewidth}\raggedright
Name
\end{minipage} & \begin{minipage}[b]{\linewidth}\raggedright
Description
\end{minipage} & \begin{minipage}[b]{\linewidth}\raggedright
Default
\end{minipage} \\
\midrule\noalign{}
\endhead
\bottomrule\noalign{}
\endlastfoot
Field & Metric & Laboratory & \\
1 & WELNAME & A character string of up to eight characters in length
that defines the well name for which the well's VFP interpolation
options are being defined. Note that the well name (WELNAME) must have
been declared previously using the WELSPECS keyword in the SCHEDULE
section, otherwise an error may occur. & None \\
2 & IMPEXP & A defined character string of length three that defines the
how the well's VFPPROD tables are to be interpolated and should be set
to one of the following character strings: If a well's WCUT and GOR is
varying as a function of BHP inside a Newton/Linear iteration (as for
example when gas cusping or water coning is occurring), the implicit
lookup of the VFP tables might indicate that the well has died, due to
the interpolation/extrapolation. If a well dies using implicit WCUT and
GOR VFP lookup then the simulator automatically switches to explicit
(previous time step) VFP lookup to prevent the well from premature
closure, and writes a message to the screen and print file stating the
fact. Note the explicit treatment is just for the VFP lookup only. Using
the IMPEXP option allows one to select wells that should use either the
implicit or explicit lookup option, which may be useful in improving run
time performance. Note that if the default value of IMP is used, the
simulator will still switch to explicit VFP lookup to prevent a well
from premature closure if the condition occurs. & IMP \\
3 & STATUS & A defined character string that defines if the well's
operating condition should be checked to see if the VFP table lookup
indicates if the well is operating on the flat or horizontal part of
stabilized portion of the VLP table (curve), and should be set to one of
the following character strings: This option is specific to the
commercial simulator's VFPi program that can generate VFPF tables where
the unstable part of the VFP curves to the left of the minimum is
replaced with a horizontal line at the minimum BHP value. This option is
not supported by OPM Flow. & NO \\
4 & CONTROL & A defined character string that defines the behavior of
rate control wells operating on the unstabilized portion of the VFP
table, and should be set to one of the following character strings: In
the commercial simulator this option will also print a message the first
time a well is prevented from changing its control mode, whereas OPM
Flow always prints a message every time a well is prevented from
changing its control mode. The default value of NO allows wells to die
when they are constrained by a lower rate that forces them to operate on
the unstable section of the VFP curve. This can occur if a well's rate
is being set to match its group target rate or constraint limit. Using
either the YES1 or YES2 option will help prevent the ``hunting between
control modes'' issue or the well prematurely being shut-in. & NO \\
5 & EXTRAP & The EXTRAP parameter is a defined character string in the
commercial composition simulator that declares how the extrapolation of
a well's gas fraction, water fraction, and ALQ values is to be
conducted. This option is not supported by OPM Flow. & WG \\
Notes: & & & \\
\end{longtable}

\emph{Table 12.131: WVFPEXP Keyword Description}

\subsubsection{Example}\label{example-612}

The following example defines two oil wells using the WELSPECS and
WCONPROD keywords, together with the WVFPEXP keyword that declares how
the VFPPROD table lookup should be performed.

\begin{verbatim}
--
-- WELL SPECIFICATION DATA
--
-- WELL     GROUP      LOCATION  BHP    PHASE  DRAIN  INFLOW  SHUT  CROSS  PRESS
-- NAME     NAME        I    J   DEPTH  FLUID  AREA   EQUA.   IN    FLOW   TABLE
WELSPECS
OP01      PLATFORM     14   13   1*      OIL   1*     STD     SHUT   NO    1*  /
OP02      PLATFORM     28   96   1*      OIL   1*     STD     SHUT   NO    1*  /
/
--
--       WELL PRODUCTION WELL CONTROLS
--
-- WELL  OPEN/  CNTL   OIL    WAT    GAS   LIQ    RES    BHP   THP   VFP    VFP
-- NAME  SHUT   MODE   RATE   RATE   RATE  RATE   RATE   PRES  PRES  TABLE  ALFQ
WCONPROD
OP01     SHUT   GRUP   1*     1*     1*    1*     1*     500.0 100.0   1       /
OP02     SHUT   GRUP   1*     1*     1*    1*     1*     500.0 100.0   1       /
/
--
-- WELL OPTIONS FOR PROBLEMATIC THP CONTROLLED WELLS
--
-- WELL   IMP   CLSE   RATE   VFP
-- NAME   EXP   WELL   CNTL   EXT
WVFPEXP
'OP*   '  1*    1*     YES1   1*                                               /
/
\end{verbatim}

Here both wells are declared as initially shut on the WELSPECS keyword
and use VFPPROD table number one as declared on the WCONPROD keyword.
The WVFPEXP keyword declares all wells with a name beginning with OP to
use implicit lookup and the wells are prevented from changing from rate
control to THP control provided they can produce at a higher rate under
THP control.

\subsection{WWPAVE -- Well Block Average Pressure Calculation Parameters
for Individual
Wells}\label{wwpave-well-block-average-pressure-calculation-parameters-for-individual-wells}

\begin{longtable}[]{@{}
  >{\raggedright\arraybackslash}p{(\linewidth - 14\tabcolsep) * \real{0.1250}}
  >{\raggedright\arraybackslash}p{(\linewidth - 14\tabcolsep) * \real{0.1250}}
  >{\raggedright\arraybackslash}p{(\linewidth - 14\tabcolsep) * \real{0.1250}}
  >{\raggedright\arraybackslash}p{(\linewidth - 14\tabcolsep) * \real{0.1250}}
  >{\raggedright\arraybackslash}p{(\linewidth - 14\tabcolsep) * \real{0.1250}}
  >{\raggedright\arraybackslash}p{(\linewidth - 14\tabcolsep) * \real{0.1250}}
  >{\raggedright\arraybackslash}p{(\linewidth - 14\tabcolsep) * \real{0.1250}}
  >{\raggedright\arraybackslash}p{(\linewidth - 14\tabcolsep) * \real{0.1250}}@{}}
\toprule\noalign{}
\begin{minipage}[b]{\linewidth}\raggedright
RUNSPEC
\end{minipage} & \begin{minipage}[b]{\linewidth}\raggedright
GRID
\end{minipage} & \begin{minipage}[b]{\linewidth}\raggedright
EDIT
\end{minipage} & \begin{minipage}[b]{\linewidth}\raggedright
PROPS
\end{minipage} & \begin{minipage}[b]{\linewidth}\raggedright
REGIONS
\end{minipage} & \begin{minipage}[b]{\linewidth}\raggedright
SOLUTION
\end{minipage} & \begin{minipage}[b]{\linewidth}\raggedright
SUMMARY
\end{minipage} & \begin{minipage}[b]{\linewidth}\raggedright
SCHEDULE
\end{minipage} \\
\midrule\noalign{}
\endhead
\bottomrule\noalign{}
\endlastfoot
\end{longtable}

\subsubsection{Description}\label{description-1150}

The WWPAVE keyword defines the method and parameters for calculating a
well's block average pressures for individual wells that can be written
to the SUMMARY and RSM files via the WBP, WBP4, WBP5 and WBP9 vectors in
the SUMMARY section. The resulting average pressure can be written out
to the summary file in order to compared with field observed data. The
keyword is similar to the WPAVE keyword in the SCHEDULE section that has
similar functionality, but is applied to all wells in the model.

Note that WWPAVE will overwrite any parameters on the WPAVE keyword for
a given well, and that WWPAVE can also be overwritten by any subsequent
WPAVE keyword.

\begin{longtable}[]{@{}
  >{\raggedright\arraybackslash}p{(\linewidth - 6\tabcolsep) * \real{0.2500}}
  >{\raggedright\arraybackslash}p{(\linewidth - 6\tabcolsep) * \real{0.2500}}
  >{\raggedright\arraybackslash}p{(\linewidth - 6\tabcolsep) * \real{0.2500}}
  >{\raggedright\arraybackslash}p{(\linewidth - 6\tabcolsep) * \real{0.2500}}@{}}
\toprule\noalign{}
\begin{minipage}[b]{\linewidth}\raggedright
No.
\end{minipage} & \begin{minipage}[b]{\linewidth}\raggedright
Name
\end{minipage} & \begin{minipage}[b]{\linewidth}\raggedright
Description
\end{minipage} & \begin{minipage}[b]{\linewidth}\raggedright
Default
\end{minipage} \\
\midrule\noalign{}
\endhead
\bottomrule\noalign{}
\endlastfoot
Field & Metric & Laboratory & \\
1 & WELNAME & A character string of up to eight characters in length
that defines the well name for which the well average pressure
calculation parameters are being defined. Note that the well name
(WELNAME) must have been declared previously using the WELSPECS keyword
in the SCHEDULE section, otherwise an error may occur. & None \\
2 & WPAVE1 & A real dimensionless value that defines the weighting
factor between the inner block and the surrounding blocks used in the
calculation of the connection factor weighted average pressure. If
WPAVE1 is greater than or equal to zero and less than or equal to one,
then the average pressure for each well connection is calculated based
on this weighting factor. A value of zero indicates only the surrounding
blocks should be used in the calculation; and a value of one indicates
only the inner blocks should be used. If WPAVE1 is less than zero, then
the average pressure for each well connection is weighted based on the
pore volumes of the inner and surrounding blocks. & 0.5 \\
3 & WPAVE2 & A real dimensionless value greater than or equal to zero
and less than or equal to one, that defines the weighting factor between
the connection weighted average pressures and the pore volume weighted
average pressures. If WPAVE2 is equal to one, then the average pressures
are calculate based only using the connection factor calculated
pressures. If WPAVE2 is equal to zero, then average pressures are
calculate based on only using the pore volumes calculated pressures. &
1.0 \\
4 & WPAVE3 & A defined character string that determines how the
hydrostatic head calculation is performed in correcting the pressures to
the BHP reference depth on the WELSPECS or WPAVEDEP keywords in the
SCHEDULE section. WPAVE3 should be set to one of the following character
strings: & WELL \\
5 & WPAVE4 & A defined character string that determines which
connections should be used in the calculations, WPAVE4 should be set to
one of the following character strings: Only the OPEN option is
currently supported by the simulator. & OPEN \\
Notes: & & & \\
\end{longtable}

\emph{Table 12.3.356.1: WWPAVE Keyword Description}

The keyword is not applicable and should not be used with radial and
spider grid geometries.

See also the WELSPECS keyword that defines a well and a well's
bottom-hole pressure reference depth, the WPAVEDEP keyword that also
defines a well's bottom-hole pressure reference depth, and the COMPDAT
keyword to define a well's connections. All the aforementioned keywords
are described in the SCHEDULE section.

\subsubsection{Examples}\label{examples-97}

The following example defines the default well block average pressure
calculation parameters for three oil wells: OP01, OP02 and OP03.

\begin{verbatim}
--
--       DEFINE WELL BLOCK AVERAGE PRESSURE CALCULATION PARAMETERS
--
-- WELL  INNER   PORV   WELL  OPEN
-- NAME  OUTER   CONN   RES   ALL
WWPAVE
OP01     0.5     1.0    WELL  ALL                                               /
OP02     0.5     1.0    WELL  ALL                                               /
OP03     0.5     1.0    1*    1*                                                /
\end{verbatim}

And the next example shows the parameters used in the Norne model.

\begin{verbatim}
--
--       DEFINE WELL BLOCK AVERAGE PRESSURE CALCULATION PARAMETERS
--
-- WELL  INNER   PORV   WELL  OPEN
-- NAME  OUTER   CONN   RES   ALL
WPAVE
OP01     1*      0.0    WELL  ALL                                               /
OP02     1*      0.0    WELL  ALL                                               /
OP03     1*      0.0    WELL  ALL                                               /
\end{verbatim}

Here only pore volume weighting is used instead of connection weighting.

\subsection{ZIPP2OFF -- Deactivate Automatic Time Step
Control}\label{zipp2off-deactivate-automatic-time-step-control}

\begin{longtable}[]{@{}
  >{\raggedright\arraybackslash}p{(\linewidth - 14\tabcolsep) * \real{0.1250}}
  >{\raggedright\arraybackslash}p{(\linewidth - 14\tabcolsep) * \real{0.1250}}
  >{\raggedright\arraybackslash}p{(\linewidth - 14\tabcolsep) * \real{0.1250}}
  >{\raggedright\arraybackslash}p{(\linewidth - 14\tabcolsep) * \real{0.1250}}
  >{\raggedright\arraybackslash}p{(\linewidth - 14\tabcolsep) * \real{0.1250}}
  >{\raggedright\arraybackslash}p{(\linewidth - 14\tabcolsep) * \real{0.1250}}
  >{\raggedright\arraybackslash}p{(\linewidth - 14\tabcolsep) * \real{0.1250}}
  >{\raggedright\arraybackslash}p{(\linewidth - 14\tabcolsep) * \real{0.1250}}@{}}
\toprule\noalign{}
\begin{minipage}[b]{\linewidth}\raggedright
RUNSPEC
\end{minipage} & \begin{minipage}[b]{\linewidth}\raggedright
GRID
\end{minipage} & \begin{minipage}[b]{\linewidth}\raggedright
EDIT
\end{minipage} & \begin{minipage}[b]{\linewidth}\raggedright
PROPS
\end{minipage} & \begin{minipage}[b]{\linewidth}\raggedright
REGIONS
\end{minipage} & \begin{minipage}[b]{\linewidth}\raggedright
SOLUTION
\end{minipage} & \begin{minipage}[b]{\linewidth}\raggedright
SUMMARY
\end{minipage} & \begin{minipage}[b]{\linewidth}\raggedright
SCHEDULE
\end{minipage} \\
\midrule\noalign{}
\endhead
\bottomrule\noalign{}
\endlastfoot
\end{longtable}

\subsubsection{Description}\label{description-1151}

The ZIPP2OFF keyword deactivates the commercial simulator's alternative
automatic time step selection algorithm that assumes no prior knowledge
of the problem, as opposed to the standard time step algorithm that is
controlled via the TUNING keyword in the SCHEDULE section, combined with
posterior knowledge gained from previous time steps.

This keyword is not supported by OPM Flow but has no effect on the
results so it will be ignored.

See section 2.2 Running OPM Flow 2023-04 From The Command Line on how to
control time stepping for OPM Flow.

\subsection{ZIPPY2 -- Activate Automatic Time Step
Control}\label{zippy2-activate-automatic-time-step-control}

\begin{longtable}[]{@{}
  >{\raggedright\arraybackslash}p{(\linewidth - 14\tabcolsep) * \real{0.1250}}
  >{\raggedright\arraybackslash}p{(\linewidth - 14\tabcolsep) * \real{0.1250}}
  >{\raggedright\arraybackslash}p{(\linewidth - 14\tabcolsep) * \real{0.1250}}
  >{\raggedright\arraybackslash}p{(\linewidth - 14\tabcolsep) * \real{0.1250}}
  >{\raggedright\arraybackslash}p{(\linewidth - 14\tabcolsep) * \real{0.1250}}
  >{\raggedright\arraybackslash}p{(\linewidth - 14\tabcolsep) * \real{0.1250}}
  >{\raggedright\arraybackslash}p{(\linewidth - 14\tabcolsep) * \real{0.1250}}
  >{\raggedright\arraybackslash}p{(\linewidth - 14\tabcolsep) * \real{0.1250}}@{}}
\toprule\noalign{}
\begin{minipage}[b]{\linewidth}\raggedright
RUNSPEC
\end{minipage} & \begin{minipage}[b]{\linewidth}\raggedright
GRID
\end{minipage} & \begin{minipage}[b]{\linewidth}\raggedright
EDIT
\end{minipage} & \begin{minipage}[b]{\linewidth}\raggedright
PROPS
\end{minipage} & \begin{minipage}[b]{\linewidth}\raggedright
REGIONS
\end{minipage} & \begin{minipage}[b]{\linewidth}\raggedright
SOLUTION
\end{minipage} & \begin{minipage}[b]{\linewidth}\raggedright
SUMMARY
\end{minipage} & \begin{minipage}[b]{\linewidth}\raggedright
SCHEDULE
\end{minipage} \\
\midrule\noalign{}
\endhead
\bottomrule\noalign{}
\endlastfoot
\end{longtable}

\subsubsection{Description}\label{description-1152}

The ZIPPY2 keyword activates the commercial simulator's alternative
automatic time step selection algorithm that assumes no prior knowledge
of the problem, as opposed to the standard time step algorithm that is
controlled via the TUNNING keyword in the SCHEDULE section, combined
with posterior knowledge gained from previous time steps.

This keyword is not supported by OPM Flow but has no effect on the
results so it will be ignored.

See section 2.2 Running OPM Flow 2023-04 From The Command Line on how to
control time stepping for OPM Flow.

\cleardoublepage
\phantomsection
\addcontentsline{toc}{part}{Appendices}
\appendix

\chapter{Keyword Summary}\label{keyword-summary}

\begin{longtable}[]{@{}
  >{\raggedright\arraybackslash}p{(\linewidth - 14\tabcolsep) * \real{0.1250}}
  >{\raggedright\arraybackslash}p{(\linewidth - 14\tabcolsep) * \real{0.1250}}
  >{\raggedright\arraybackslash}p{(\linewidth - 14\tabcolsep) * \real{0.1250}}
  >{\raggedright\arraybackslash}p{(\linewidth - 14\tabcolsep) * \real{0.1250}}
  >{\raggedright\arraybackslash}p{(\linewidth - 14\tabcolsep) * \real{0.1250}}
  >{\raggedright\arraybackslash}p{(\linewidth - 14\tabcolsep) * \real{0.1250}}
  >{\raggedright\arraybackslash}p{(\linewidth - 14\tabcolsep) * \real{0.1250}}
  >{\raggedright\arraybackslash}p{(\linewidth - 14\tabcolsep) * \real{0.1250}}@{}}
\toprule\noalign{}
\begin{minipage}[b]{\linewidth}\raggedright
RUNSPEC
\end{minipage} & \begin{minipage}[b]{\linewidth}\raggedright
GRID
\end{minipage} & \begin{minipage}[b]{\linewidth}\raggedright
EDIT
\end{minipage} & \begin{minipage}[b]{\linewidth}\raggedright
PROPS
\end{minipage} & \begin{minipage}[b]{\linewidth}\raggedright
REGIONS
\end{minipage} & \begin{minipage}[b]{\linewidth}\raggedright
SOLUTION
\end{minipage} & \begin{minipage}[b]{\linewidth}\raggedright
SUMMARY
\end{minipage} & \begin{minipage}[b]{\linewidth}\raggedright
SCHEDULE
\end{minipage} \\
\midrule\noalign{}
\endhead
\bottomrule\noalign{}
\endlastfoot
Alphabetic Listing of Keywords Starting with the Letter A & Status & & &
& & & \\
\hyperref[0.0.0.ACF]{ACF -- Define Acentric Factors} & & & & & & & \\
\hyperref[0.0.0.ACTCO2S]{ACTCO2S -- ACTCO2S -- Activity Model for CO2} &
& & & & & & \\
ACTDIMS -- ACTION Keyword Dimensions & A & & & & & & \\
ACTION -- Define Action Conditions and Command Processing (Field) & & &
& & & & \\
ACTIONG -- Define Action Conditions and Command Processing (Groups) & &
& & & & & \\
ACTIONR -- Define Action Conditions and Command Processing (Regions) & &
& & & & & \\
ACTIONS -- Define Action Conditions and Command Processing (Well
Segments) & & & & & & & \\
ACTIONW -- Define Action Conditions and Command Processing (Wells) & & &
& & & & \\
ACTIONX -- Define Action Conditions and Command Processing & & & & & &
& \\
ACTNUM -- Set the Status of a Grid Block To Active or Inactive & & & & &
& & \\
ACTPARAM -- Define Action Facility Target and Tolerance Parameters & & &
& & & & \\
ADD -- Add a Constant to a Specified Array & & & & & & & \\
ADDREG -- Add a Constant to an Array based on a Region Number & & & & &
& & \\
ADDZCORN -- Add a Constant to the ZCORN Depth Array & & & & & & & \\
ADSALNOD -- Salt Concentration Based on SATNUM Array & & & & & & & \\
ADSORP -- Define Generalized Langmuir Adsorption Function & & & & & &
& \\
AITS -- Activate Intelligent Time Stepping & & & & & & & \\
AITSOFF -- Deactivate Intelligent Time Stepping & & & & & & & \\
ALKADS -- Define Alkaline Adsorption Functions & & & & & & & \\
ALKALINE -- Activate the Alkaline Phase and Model & & & & & & & \\
ALKROCK -- Define Rock Alkaline Properties & & & & & & & \\
ALL -- Export Standard Summary Variable Vectors to File & & & & & & & \\
ALPOLADS -- Polymer Adsorption versus Alkaline Concentration Multipliers
& & & & & & & \\
ALSURFAD -- Surfactant Adsorption versus Alkaline Concentration
Multipliers & & & & & & & \\
ALSURFST -- Water-Oil Surface Tension versus Alkaline Concentration
Multipliers & & & & & & & \\
AMALGAM -- Define LGR Amalgamations & & & & & & & \\
API -- Activate API Tracking & & & & & & & \\
APIGROUP -- Define API Tracking Number of Grouped Oil PVT Tables & & & &
& & & \\
APILIM -- Define API Tracking Grid Block Limits & & & & & & & \\
APIVD - Equilibration Oil API Gravity versus Depth Tables & & & & & &
& \\
‍AQANCONL -- Define Analytic Connections to a LGR Grid & & & & & & & \\
‍AQANNC -- Define Analytic Aquifer Non-Neighbor Connections & & & & & &
& \\
AQANTRC - Define Analytic Aquifer Initial Tracer Concentrations & & & &
& & & \\
AQUALIST -- Define An Analytic Aquifer Name to Aquifer Numbers & & & & &
& & \\
AQUANCON -- Define Analytical Connections to the Grid‍ & & & & & & & \\
AQUCHGAS -- Define Constant Pressure Gas Analytical Aquifer Properties &
& & & & & & \\
AQUCHWAT -- Define Constant Pressure Water Analytical Aquifer Properties
& & & & & & & \\
AQUCON -- Define Numerical Aquifer Connections to the Grid & & & & & &
& \\
AQUCT -- Define Carter-Tracy Analytical Aquifers & & & & & & & \\
AQUCWFAC -- Modify Constant Pressure Water Analytical Aquifer Properties
& & & & & & & \\
AQUDIMS -- Define Aquifer Dimensions & & & & & & & \\
AQUFET -- Define Fetkovich Analytical Aquifer and Connections & & & & &
& & \\
AQUFETP -- Define Fetkovich Analytical Aquifers & & & & & & & \\
AQUFLUX - Define Constant Flux Analytical Aquifer & & & & & & & \\
AQUNNC -- Define Numerical Aquifer Non-Neighbor Connections & & & & & &
& \\
AQUNUM -- Define Numerical Aquifer Properties & & & & & & & \\
AQUTAB -- Define Carter-Tracy Aquifer Influence Functions & & & & & &
& \\
AUTOCOAR - Define Auto Refinement Grid Coarsen Area & & & & & & & \\
AUTOREF - Define Auto Refinement Options & & & & & & & \\
\end{longtable}

\begin{longtable}[]{@{}
  >{\raggedright\arraybackslash}p{(\linewidth - 14\tabcolsep) * \real{0.1250}}
  >{\raggedright\arraybackslash}p{(\linewidth - 14\tabcolsep) * \real{0.1250}}
  >{\raggedright\arraybackslash}p{(\linewidth - 14\tabcolsep) * \real{0.1250}}
  >{\raggedright\arraybackslash}p{(\linewidth - 14\tabcolsep) * \real{0.1250}}
  >{\raggedright\arraybackslash}p{(\linewidth - 14\tabcolsep) * \real{0.1250}}
  >{\raggedright\arraybackslash}p{(\linewidth - 14\tabcolsep) * \real{0.1250}}
  >{\raggedright\arraybackslash}p{(\linewidth - 14\tabcolsep) * \real{0.1250}}
  >{\raggedright\arraybackslash}p{(\linewidth - 14\tabcolsep) * \real{0.1250}}@{}}
\toprule\noalign{}
\begin{minipage}[b]{\linewidth}\raggedright
RUNSPEC
\end{minipage} & \begin{minipage}[b]{\linewidth}\raggedright
GRID
\end{minipage} & \begin{minipage}[b]{\linewidth}\raggedright
EDIT
\end{minipage} & \begin{minipage}[b]{\linewidth}\raggedright
PROPS
\end{minipage} & \begin{minipage}[b]{\linewidth}\raggedright
REGIONS
\end{minipage} & \begin{minipage}[b]{\linewidth}\raggedright
SOLUTION
\end{minipage} & \begin{minipage}[b]{\linewidth}\raggedright
SUMMARY
\end{minipage} & \begin{minipage}[b]{\linewidth}\raggedright
SCHEDULE
\end{minipage} \\
\midrule\noalign{}
\endhead
\bottomrule\noalign{}
\endlastfoot
Alphabetic Listing of Keywords Starting with the Letter B & Status & & &
& & & \\
BC -- Define Boundary Conditions & OPM Flow & & & & & & \\
BCCON -- Define Boundary Conditions Connections & OPM Flow & & & & &
& \\
BCPROP -- Define Boundary Conditions Properties & OPM Flow & & & & &
& \\
BDENSITY -- Define the Surface Brine Density for the Fluid & B & & & & &
& \\
BGGI - Define Gi Gas Formation Volume Factor Pressure Tables & & & & & &
& \\
\hyperref[0.0.0.BIC]{BIC -- Binary Interaction Coefficients} & & & & & &
& \\
BIGMODEL -- Activate Big Model Option (Retired) & & & & & & & \\
\hyperref[0.0.0.BIOFILM]{BIOFILM -- Activate the Biofilm Model} & OPM
Flow & & & & & & \\
\hyperref[0.0.0.BIOFPARA]{BIOFPARA -- Define Properties for Models with
Biofilms} & OPM Flow & & & & & & \\
BLACKOIL -- Activate Black-Oil Phases & & & & & & & \\
BOGI - Define Gi Oil Formation Volume Factor Pressure Tables & & & & & &
& \\
BOUNDARY -- Define a Boundary Box for Printing & & & & & & & \\
BOX - Define a Range of Grid Blocks to Enter Property Data & AFR & & & &
& & \\
BPARA -- Activate Block Parallel License & & & & & & & \\
BPIDIMS -- Define the Dimensions of the Interpolated Block Quantities &
& & & & & & \\
BRANPROP -- Define Network Branch Properties for Extended Network Option
& & & & & & & \\
BRINE -- Activate Brine Tracking Option & & & & & & & \\
BTOBALFA -- Dual Porosity Matrix to Fracture Multiplier (All Cells) & &
& & & & & \\
BTOBALFV -- Dual Porosity Matrix to Fracture Multiplier (Individual
Cells) & & & & & & & \\
\end{longtable}

\begin{longtable}[]{@{}
  >{\raggedright\arraybackslash}p{(\linewidth - 14\tabcolsep) * \real{0.1250}}
  >{\raggedright\arraybackslash}p{(\linewidth - 14\tabcolsep) * \real{0.1250}}
  >{\raggedright\arraybackslash}p{(\linewidth - 14\tabcolsep) * \real{0.1250}}
  >{\raggedright\arraybackslash}p{(\linewidth - 14\tabcolsep) * \real{0.1250}}
  >{\raggedright\arraybackslash}p{(\linewidth - 14\tabcolsep) * \real{0.1250}}
  >{\raggedright\arraybackslash}p{(\linewidth - 14\tabcolsep) * \real{0.1250}}
  >{\raggedright\arraybackslash}p{(\linewidth - 14\tabcolsep) * \real{0.1250}}
  >{\raggedright\arraybackslash}p{(\linewidth - 14\tabcolsep) * \real{0.1250}}@{}}
\toprule\noalign{}
\begin{minipage}[b]{\linewidth}\raggedright
RUNSPEC
\end{minipage} & \begin{minipage}[b]{\linewidth}\raggedright
GRID
\end{minipage} & \begin{minipage}[b]{\linewidth}\raggedright
EDIT
\end{minipage} & \begin{minipage}[b]{\linewidth}\raggedright
PROPS
\end{minipage} & \begin{minipage}[b]{\linewidth}\raggedright
REGIONS
\end{minipage} & \begin{minipage}[b]{\linewidth}\raggedright
SOLUTION
\end{minipage} & \begin{minipage}[b]{\linewidth}\raggedright
SUMMARY
\end{minipage} & \begin{minipage}[b]{\linewidth}\raggedright
SCHEDULE
\end{minipage} \\
\midrule\noalign{}
\endhead
\bottomrule\noalign{}
\endlastfoot
Alphabetic Listing of Keywords Starting with the Letter C & Status & & &
& & & \\
CALTRAC -- Define a Gas Calorfic Value Tracer & C & & & & & & \\
CARFIN -- Define a Cartesian Local Grid Refinement & & & & & & & \\
CART -- Activate Cartesian Geometry & & & & & & & \\
CBMOPTS -- Define Coal Bed Methane Options & & & & & & & \\
CECON -- Define Well Connections Economic Limit Criteria & & & & & &
& \\
CECONT -- Define Well Connections Tracer Economic Limit Criteria & & & &
& & & \\
CNAMES -- Define Compositional Component Names & OPM Flow & & & & & & \\
\hyperref[0.0.0.CO2SOL]{CO2SOL -- Activate Dissolved CO2 in the Water
Phase} & OPM Flow & & & & & & \\
CO2STORE -- Activate the CO2 Storage Model & OPM Flow & & & & & & \\
COAL -- Activate the Coal Phase (CBM Model) & & & & & & & \\
COALADS -- Define Gas and Solvent Relative Adsorption Tables & & & & & &
& \\
COALNUM -- Define the Coal Region Numbers & & & & & & & \\
COALPP -- Define Gas and Solvent Partial Pressure Adsorption Tables & &
& & & & & \\
COARSEN -- Define Grid Coarsening Cells & & & & & & & \\
COLLAPSE -- Define Compressed Vertical Equilibrium Cells & & & & & &
& \\
COLUMNS -- Define Input File Column Margins & & & & & & & \\
COMPDAT -- Define Well Connections to the Grid & & & & & & & \\
COMPDATL -- Define Well Connections to a LGR Grid & & & & & & & \\
COMPDATM -- Define Well Connections to an Amalgamated LGR Grid & & & & &
& & \\
COMPFLSH -- Assign COMPDAT Differential-Flash Liberation Ratios & & & &
& & & \\
COMPIMB -- Assign Imbibition Saturation Tables to Well Connections & & &
& & & & \\
COMPINJK -- Assign Injection well Relative Permeability Values & & & & &
& & \\
COMPLMPL -- Assign Well LGR Connections to Completions & & & & & & & \\
COMPLUMP -- Assign Well Connections to Completions & & & & & & & \\
COMPOFF -- Deactivate Network Automatic Compressors & & & & & & & \\
COMPORD - Define Well Connection Ordering & & & & & & & \\
COMPRIV -- Define Grid Cell Connections to a River & & & & & & & \\
COMPRP -- Re-Scale Fluid Saturations of Well Connections & & & & & &
& \\
COMPRPL -- Re-Scale Fluid Saturations of Well LGR Connections & & & & &
& & \\
COMPS -- Activate Compositional Modeling Formulation & OPM Flow & & & &
& & \\
COMPSEGL -- Define Well Connections for Multi-Segment Wells in a LGR & &
& & & & & \\
COMPSEGS -- Define Well Connections for Multi-Segment Wells & & & & & &
& \\
‍COMPTRAJ -- Define Well Trajectory Connections to the Grid & OPM Flow &
& & & & & \\
COMPVE -- Re-Define Well Connection Depths & & & & & & & \\
COMPVEL -- Re-Define Well LGR Connection Depths & & & & & & & \\
COORD -- Define a Set of Coordinates Lines for a Reservoir Grid & & & &
& & & \\
COORDSYS -- Define Coordinate Grid Options & & & & & & & \\
COPY -- Copy Array Data to Another Array & & & & & & & \\
COPYBOX -- Copy Array Data Defined by a Box & & & & & & & \\
COPYREG -- Copy an Array to Another Array based on a Region Number & & &
& & & & \\
CPIFACT -- Define Well Connection Transmissibility Multipliers & & & & &
& & \\
CPIFACTL -- Define Well Connection Transmissibility Multipliers in a LGR
& & & & & & & \\
CPR -- Activate Constrained Pressure Residual (``CPR'') Linear Solver &
& & & & & & \\
CRITPERM -- Define Minimum Permeability for Vertical Equilibrium Grid
Cell Compression & & & & & & & \\
CSKIN -- Re-Define Well Connection Skin Factors & & & & & & & \\
\end{longtable}

\begin{longtable}[]{@{}
  >{\raggedright\arraybackslash}p{(\linewidth - 14\tabcolsep) * \real{0.1250}}
  >{\raggedright\arraybackslash}p{(\linewidth - 14\tabcolsep) * \real{0.1250}}
  >{\raggedright\arraybackslash}p{(\linewidth - 14\tabcolsep) * \real{0.1250}}
  >{\raggedright\arraybackslash}p{(\linewidth - 14\tabcolsep) * \real{0.1250}}
  >{\raggedright\arraybackslash}p{(\linewidth - 14\tabcolsep) * \real{0.1250}}
  >{\raggedright\arraybackslash}p{(\linewidth - 14\tabcolsep) * \real{0.1250}}
  >{\raggedright\arraybackslash}p{(\linewidth - 14\tabcolsep) * \real{0.1250}}
  >{\raggedright\arraybackslash}p{(\linewidth - 14\tabcolsep) * \real{0.1250}}@{}}
\toprule\noalign{}
\begin{minipage}[b]{\linewidth}\raggedright
RUNSPEC
\end{minipage} & \begin{minipage}[b]{\linewidth}\raggedright
GRID
\end{minipage} & \begin{minipage}[b]{\linewidth}\raggedright
EDIT
\end{minipage} & \begin{minipage}[b]{\linewidth}\raggedright
PROPS
\end{minipage} & \begin{minipage}[b]{\linewidth}\raggedright
REGIONS
\end{minipage} & \begin{minipage}[b]{\linewidth}\raggedright
SOLUTION
\end{minipage} & \begin{minipage}[b]{\linewidth}\raggedright
SUMMARY
\end{minipage} & \begin{minipage}[b]{\linewidth}\raggedright
SCHEDULE
\end{minipage} \\
\midrule\noalign{}
\endhead
\bottomrule\noalign{}
\endlastfoot
Alphabetic Listing of Keywords Starting with the Letter D & Status & & &
& & & \\
DATE - Activate the DATE Option for the SUMMARY File & D & & & & & & \\
DATES -- Advance Simulation by Reporting Date & & & & & & & \\
DATUM -- Define the Datum Depth for the Model & & & & & & & \\
DATUMR -- Define Datum Depths for the FIPNUM Regions & & & & & & & \\
DATUMRX -- Define Datum Depths for the FIP Allocated Regions & & & & & &
& \\
DCQDEFN -- Define Gas DCQ Units as Rate or Energy & & & & & & & \\
DEBUG -- Define the Debug Data to be Printed to File & & & & & & & \\
DELAYACT - Define Delayed Action Keywords & & & & & & & \\
\hyperref[0.0.0.DENAQA]{DENAQA -- Specify Ezrokhi Coefficients for
Aqueous Density} & & & & & & & \\
DENSITY -- Define the Surface Oil, Water Gas Densities for the Fluids &
& & & & & & \\
DEPTH - Edits the Depth at the Center of Each Cell & & & & & & & \\
DEPTHTAB -- River Time and Depth Tables & & & & & & & \\
DIAGDISP -- Activate Alternate Form of Tracer Dispersion & & & & & &
& \\
\hyperref[0.0.0.DIFFAGAS]{DIFFAGAS -- Define PVT Region Gas Component
Diffusion Coefficients (Mass Fraction Formulation)} & OPM Flow & & & & &
& \\
\hyperref[0.0.0.DIFFAWAT]{DIFFAWAT -- Define PVT Region Water Component
Diffusion Coefficients (Mass Fraction Formulation)} & OPM Flow & & & & &
& \\
DIFFC -- Define PVT Region Molecular Diffusion Tables & & & & & & & \\
DIFFCGAS -- Define PVT Region Gas Component Diffusion Coefficients & OPM
Flow & & & & & & \\
DIFFCOAL -- Define Coal Bed Methane Gas Diffusion Data & & & & & & & \\
DIFFCWAT -- Define PVT Region Water Component Diffusion Coefficients &
OPM Flow & & & & & & \\
DIFFDP -- Activate Dual Porosity Molecular Diffusion for Matrix-Fracture
Flow Only & & & & & & & \\
\hyperref[0.0.0.DIFFMICP]{DIFFMICP -- Define PVT Region MICP Diffusion
Coefficients (Mass Concentration Formulation)} & OPM Flow & & & & & & \\
DIFFMMF -- Define Diffusivity Multipliers for Matrix-Fractures & & & & &
& & \\
DIFFMR -- Define Grid Block Radial Direction Diffusivity Multipliers & &
& & & & & \\
DIFFMR- -- Define Grid Block Negative Radial Direction Diffusivity
Multipliers & & & & & & & \\
DIFFMTHT -- Define Grid Block Theta Direction Diffusivity Multipliers &
& & & & & & \\
DIFFMTH- -- Define Grid Block Negative Theta Direction Diffusivity
Multipliers & & & & & & & \\
DIFFMX -- Define Grid Block X-Direction Diffusivity Multipliers & & & &
& & & \\
DIFFMX- -- Define Grid Block Negative X-Direction Diffusivity
Multipliers & & & & & & & \\
DIFFMY -- Define Grid Block Y-Direction Diffusivity Multipliers & & & &
& & & \\
DIFFMY- -- Define Grid Block Negative Y-Direction Diffusivity
Multipliers & & & & & & & \\
DIFFMZ -- Define Grid Block Z-Direction Diffusivity Multipliers & & & &
& & & \\
DIFFMZ- -- Define Grid Block Negative Z-Direction Diffusivity
Multipliers & & & & & & & \\
DIFFR -- Define Grid Block Radial Direction Diffusivity Values & & & & &
& & \\
DIFFTHT -- Define Grid Block Theta Direction Diffusivity Values & & & &
& & & \\
DIFFUSE -- Activate Molecular Diffusion Option & & & & & & & \\
DIFFX -- Define Grid Block X-Direction Diffusivity Values & & & & & &
& \\
DIFFY -- Define Grid Block Y-Direction Diffusivity Values & & & & & &
& \\
DIFFZ -- Define Grid Block Z-Direction Diffusivity Values & & & & & &
& \\
DIMENS -- Define the Dimension of the Model & & & & & & & \\
DIMPES -- Define IMPES Dynamic Solution Parameters & & & & & & & \\
DIMPLICT -- Activate Fully Implicit Dynamic Solution Formulation & & & &
& & & \\
DISGAS -- Activate the Dissolved Gas Phase in the Model & & & & & & & \\
‍DISGASW -- Activate Dissolved Gas in the Water Phase in the Model & OPM
Flow & & & & & & \\
DISPDIMS -- Define the Maximum Number of Dispersion Tables & & & & & &
& \\
\hyperref[0.0.0.DISPERC]{DISPERC -- Define the Mechanical Dispersivity
for All Cells} & OPM Flow & & & & & & \\
DISPERSE -- Define Dispersion Tables & & & & & & & \\
DOMAINS -- Define the Parallel Domain Properties & & & & & & & \\
DPGRID -- Activate The Matrix Cell to Fracture Cell Option & & & & & &
& \\
DPKRMOD -- Modify Matrix Oil Relative Permeability Data & & & & & & & \\
DPNUM -- Define Dual and Single Porosity Grid Block Array & & & & & &
& \\
DR - Define the Size of Grid Blocks in the R Direction for All Cells & &
& & & & & \\
DRILPRI -- Define Prioritized Drilling Queue Priority Parameters & & & &
& & & \\
DRSDT -- Solution Gas (Rs) Maximum Rate of Increase Parameters & & & & &
& & \\
DRSDTCON -- CO2 Convective Dissolution Parameter (χ) & OPM Flow & & & &
& & \\
DRSDTR -- Solution Gas (Rs) Maximum Rate of Increase Parameters by
Region & & & & & & & \\
DRV - Define the Size of Grid Blocks in the R Direction via a Vector & &
& & & & & \\
DRVDT -- Solution Oil (Rv) Maximum Rate of Increase Parameters & & & & &
& & \\
DRVDTR -- Solution Oil (Rv) Maximum Rate of Increase Parameters by
Region & & & & & & & \\
DSPDEINT -- Activate Brine Tracer Dispersion Interpolation by Water
Density & & & & & & & \\
DTHETA - Define the Size of Grid Blocks in the THETA Direction for All
Cells & & & & & & & \\
DTHETAV - Sets the Size of Grid Blocks in THETA Direction via a Vector &
& & & & & & \\
DUALPERM -- Activate Dual Permeability Model & & & & & & & \\
DUALPORO -- Activate Dual Porosity Model & & & & & & & \\
DUMPCUPL -- Activate Output to the Reservoir Coupling File & & & & & &
& \\
DUMPFLUX -- Activate Writing Out of a Flux File & & & & & & & \\
DX - Define the Size of Grid Blocks in the X Direction for All Cells & &
& & & & & \\
DXV - Define the Size of Grid Blocks in the X Direction via a Vector & &
& & & & & \\
DY - Define the Size of Grid Blocks in the Y Direction for All Cells & &
& & & & & \\
DYNAMICR -- Start of Dynamic Region Parameter Definition & & & & & &
& \\
DYNRDIMS -- Define Dynamic Region Dimensions & & & & & & & \\
DYV - Define the Size of Grid Blocks in the Y Direction via a Vector & &
& & & & & \\
DZ - Define the Size of Grid Blocks in the Z Direction for All Cells & &
& & & & & \\
DZMATRIX - Matrix Block Height for Gravity Drainage Model For All Cells
& & & & & & & \\
DZMTRX - Matrix Block Height for Gravity Drainage Model for the Grid & &
& & & & & \\
DZMTRXV - Matrix Block Height for Gravity Drainage Model For All Cells &
& & & & & & \\
DZNET -- Define Grid Block Net Thickness for All Cells & & & & & & & \\
DZV - Define the Size of Grid Blocks in the Z Direction via a Vector & &
& & & & & \\
\end{longtable}

\begin{longtable}[]{@{}
  >{\raggedright\arraybackslash}p{(\linewidth - 14\tabcolsep) * \real{0.1250}}
  >{\raggedright\arraybackslash}p{(\linewidth - 14\tabcolsep) * \real{0.1250}}
  >{\raggedright\arraybackslash}p{(\linewidth - 14\tabcolsep) * \real{0.1250}}
  >{\raggedright\arraybackslash}p{(\linewidth - 14\tabcolsep) * \real{0.1250}}
  >{\raggedright\arraybackslash}p{(\linewidth - 14\tabcolsep) * \real{0.1250}}
  >{\raggedright\arraybackslash}p{(\linewidth - 14\tabcolsep) * \real{0.1250}}
  >{\raggedright\arraybackslash}p{(\linewidth - 14\tabcolsep) * \real{0.1250}}
  >{\raggedright\arraybackslash}p{(\linewidth - 14\tabcolsep) * \real{0.1250}}@{}}
\toprule\noalign{}
\begin{minipage}[b]{\linewidth}\raggedright
RUNSPEC
\end{minipage} & \begin{minipage}[b]{\linewidth}\raggedright
GRID
\end{minipage} & \begin{minipage}[b]{\linewidth}\raggedright
EDIT
\end{minipage} & \begin{minipage}[b]{\linewidth}\raggedright
PROPS
\end{minipage} & \begin{minipage}[b]{\linewidth}\raggedright
REGIONS
\end{minipage} & \begin{minipage}[b]{\linewidth}\raggedright
SOLUTION
\end{minipage} & \begin{minipage}[b]{\linewidth}\raggedright
SUMMARY
\end{minipage} & \begin{minipage}[b]{\linewidth}\raggedright
SCHEDULE
\end{minipage} \\
\midrule\noalign{}
\endhead
\bottomrule\noalign{}
\endlastfoot
Alphabetic Listing of Keywords Starting with the Letter E & Status & & &
& & & \\
ECHO -- Activate Echoing of User Input Files to the Print File & E & & &
& & & \\
ECLMC -- Activate Multi-Component Brine Model & & & & & & & \\
EDIT - Define the Start of the EDIT Section of Keywords & & & & & & & \\
EDITNNC -- Scale Non-Neighbor Connections Between Cells Manually & & & &
& & & \\
EDITNNCR -- Reset Non-Neighbor Connections Between Cells Manually & & &
& & & & \\
EHYSTR -- Define Hysteresis Model and Parameters & & & & & & & \\
EHYSTRR -- Define Hysteresis Model and Parameters via SATNUM & & & & & &
& \\
END -- Define the End of the Input File & & & & & & & \\
ENDACTIO -- End the Definition of ACTION Commands & & & & & & & \\
ENDBOX -- Define the End of the BOX Defined Grid & & & & & & & \\
ENDDYN-- End of Dynamic Region Parameter Definition & & & & & & & \\
ENDFIN -- End the Definition of a Local Grid Refinement & & & & & & & \\
ENDINC -- Define the End of an Include File & & & & & & & \\
ENDNUM -- Define the End-Point Scaling Depth Region Numbers & & & & & &
& \\
ENDSCALE -- Activate Relative Permeability End-Point Scaling Option & &
& & & & & \\
ENDSKIP -- Deactivate Skipping of Keywords and Input Data & & & & & &
& \\
ENKRVD -- Define Relative Permeability End-Points versus Depth Functions
& & & & & & & \\
ENKRVDX & & & & & & & \\
ENKRVDX- & & & & & & & \\
ENKRVDY & & & & & & & \\
ENKRVDY- & & & & & & & \\
ENKRVDZ & & & & & & & \\
ENKRVDZ- & & & & & & & \\
ENPCVD -- Define Maximum Capillary Pressure versus Depth Functions & & &
& & & & \\
ENPTVD -- Define Relative Permeability Saturation End-Points versus
Depth & & & & & & & \\
ENPTVDX & & & & & & & \\
ENPTVDX- & & & & & & & \\
ENPTVDY & & & & & & & \\
ENPTVDY- & & & & & & & \\
ENPTVDZ & & & & & & & \\
ENPTVDZ- & & & & & & & \\
ENSPCVD -- Define Capillary Pressure End-Points versus Depth & & & & & &
& \\
\hyperref[0.0.0.EOS]{EOS -- Specify Equation of State} & & & & & & & \\
\hyperref[0.0.0.EOSNUM]{EOSNUM -- Define the Equation of State Region
Numbers} & & & & & & & \\
EPSDBGS - Write End-Point Debug Data to the DEBUG File (Multiple) & & &
& & & & \\
EPSDEBUG - Write End-Point Debug Data to the DEBUG File (Individual) & &
& & & & & \\
EQLDIMS -- Define the Equilibration Data Dimensions & & & & & & & \\
EQLNUM -- Define the Equilibration Region Numbers & & & & & & & \\
EQLOPTS -- Activate the Equilibration Options & & & & & & & \\
EQLZCORN - Modify the Depth of the Corner-Point Depth Array & & & & & &
& \\
EQUALREG -- Sets an Array to a Constant by Region Number & & & & & &
& \\
EQUALS -- Sets a Specified Array to a Constant & & & & & & & \\
EQUIL -- Define the Equilibration Initialization Data & & & & & & & \\
ESSNODE -- Define Salt Concentration Data for Water-Oil Surface Tension
& & & & & & & \\
EXCAVATE - Set the Status of a Grid Block To Active or Excavate & & & &
& & & \\
EXCEL - Activate the EXCEL Option for the SUMMARY File & & & & & & & \\
EXIT -- Exit Simulation from within an Action Section & OPM Flow & & & &
& & \\
EXTFIN - Define an External Unstructured Local Grid Refinement & & & & &
& & \\
EXTHOST - Define Host Cells for External LGRs & & & & & & & \\
EXTRAPMS -- Activate Extrapolation Warning Messages & & & & & & & \\
EXTREPGL - Define Host Cells for External Unstructured LGRs & & & & & &
& \\
\end{longtable}

\begin{longtable}[]{@{}
  >{\raggedright\arraybackslash}p{(\linewidth - 14\tabcolsep) * \real{0.1250}}
  >{\raggedright\arraybackslash}p{(\linewidth - 14\tabcolsep) * \real{0.1250}}
  >{\raggedright\arraybackslash}p{(\linewidth - 14\tabcolsep) * \real{0.1250}}
  >{\raggedright\arraybackslash}p{(\linewidth - 14\tabcolsep) * \real{0.1250}}
  >{\raggedright\arraybackslash}p{(\linewidth - 14\tabcolsep) * \real{0.1250}}
  >{\raggedright\arraybackslash}p{(\linewidth - 14\tabcolsep) * \real{0.1250}}
  >{\raggedright\arraybackslash}p{(\linewidth - 14\tabcolsep) * \real{0.1250}}
  >{\raggedright\arraybackslash}p{(\linewidth - 14\tabcolsep) * \real{0.1250}}@{}}
\toprule\noalign{}
\begin{minipage}[b]{\linewidth}\raggedright
RUNSPEC
\end{minipage} & \begin{minipage}[b]{\linewidth}\raggedright
GRID
\end{minipage} & \begin{minipage}[b]{\linewidth}\raggedright
EDIT
\end{minipage} & \begin{minipage}[b]{\linewidth}\raggedright
PROPS
\end{minipage} & \begin{minipage}[b]{\linewidth}\raggedright
REGIONS
\end{minipage} & \begin{minipage}[b]{\linewidth}\raggedright
SOLUTION
\end{minipage} & \begin{minipage}[b]{\linewidth}\raggedright
SUMMARY
\end{minipage} & \begin{minipage}[b]{\linewidth}\raggedright
SCHEDULE
\end{minipage} \\
\midrule\noalign{}
\endhead
\bottomrule\noalign{}
\endlastfoot
Alphabetic Listing of Keywords Starting with the Letter F & Status & & &
& & & \\
FAULTDIM -- Define the Number of Fault Segments & F & & & & & & \\
FAULTS -- Define Faults in the Grid Geometry & & & & & & & \\
FBHPDEF -- Define Well Default BHP Target and Constraints & & & & & &
& \\
FHERCHBL -- Define Herschel-Bulkley Data versus Polymer Concentration &
& & & & & & \\
FIELD -- Activate the Oil Field System of Units for the Model & & & & &
& & \\
FILEUNIT -- Activate Unit Consistency Checking & & & & & & & \\
FILLEPS -- Activate Saturation End-Point Export to the INIT File & & & &
& & & \\
FIP -- Define the Fluid In-Place Names and Region and Numbers & & & & &
& & \\
FIPNUM -- Define the Fluid In-Place Region Numbers & & & & & & & \\
FIPOWG -- Activate Oil, Gas, and Water FIP Zone Reporting & & & & & &
& \\
FLUXNUM -- Define the Flux Regions & & & & & & & \\
FLUXREG -- Define Active Flux Regions & & & & & & & \\
FLUXTYPE -- Defines the Flux Boundary Type & & & & & & & \\
FMTHMD -- Activate The Format History Match Gradient File Option & & & &
& & & \\
FMTIN -- Activate The Format Input File Option & & & & & & & \\
FMTOUT -- Activate The Format Output File Option & & & & & & & \\
FMWSET - Export Well Status Vectors for the Field to File & & & & & &
& \\
FOAM -- Activate the Foam Phase and Model & & & & & & & \\
FOAMADS - Define Foam Rock Adsorption Tables & & & & & & & \\
FOAMDCYO -- Define Foam Decay versus Oil Saturation Tables & & & & & &
& \\
FOAMDCYW -- Define Foam Decay versus Water Saturation Tables & & & & & &
& \\
FOAMFCN -- Define Foam Gas Mobility Reduction versus Capillary Number &
& & & & & & \\
FOAMFRM -- Define Foam Gas Mobility Reduction versus Reference Mobility
& & & & & & & \\
FOAMFSC -- Define Foam Gas Mobility versus Surfactant Concentration
Functions & & & & & & & \\
FOAMFSO -- Define Foam Gas Mobility Reduction versus Oil Saturation & &
& & & & & \\
FOAMFST -- Define Foam Gas-Water Surface Tension versus Surfactant
Concentration & & & & & & & \\
FOAMFSW -- Define Foam Gas Mobility Reduction versus Water Saturation &
& & & & & & \\
FOAMMOB - Define Foam Gas Mobility versus Foam Concentration Tables & &
& & & & & \\
FOAMMOBP -- Define Foam Mobility Reduction versus Oil Pressure & & & & &
& & \\
FOAMMOBS -- Define Foam Mobility Reduction versus Shear & & & & & & & \\
FOAMOPTS - Define Foam Model Options & & & & & & & \\
FOAMROCK - Define Foam Rock Properties & & & & & & & \\
FORMFEED -- Defined the Print File Form-Feed Character & & & & & & & \\
FRICTION -- Activate Wellbore Friction Option & & & & & & & \\
FULLIMP -- Activate Fully Implicit Solution Option & & & & & & & \\
\end{longtable}

\begin{longtable}[]{@{}
  >{\raggedright\arraybackslash}p{(\linewidth - 14\tabcolsep) * \real{0.1250}}
  >{\raggedright\arraybackslash}p{(\linewidth - 14\tabcolsep) * \real{0.1250}}
  >{\raggedright\arraybackslash}p{(\linewidth - 14\tabcolsep) * \real{0.1250}}
  >{\raggedright\arraybackslash}p{(\linewidth - 14\tabcolsep) * \real{0.1250}}
  >{\raggedright\arraybackslash}p{(\linewidth - 14\tabcolsep) * \real{0.1250}}
  >{\raggedright\arraybackslash}p{(\linewidth - 14\tabcolsep) * \real{0.1250}}
  >{\raggedright\arraybackslash}p{(\linewidth - 14\tabcolsep) * \real{0.1250}}
  >{\raggedright\arraybackslash}p{(\linewidth - 14\tabcolsep) * \real{0.1250}}@{}}
\toprule\noalign{}
\begin{minipage}[b]{\linewidth}\raggedright
RUNSPEC
\end{minipage} & \begin{minipage}[b]{\linewidth}\raggedright
GRID
\end{minipage} & \begin{minipage}[b]{\linewidth}\raggedright
EDIT
\end{minipage} & \begin{minipage}[b]{\linewidth}\raggedright
PROPS
\end{minipage} & \begin{minipage}[b]{\linewidth}\raggedright
REGIONS
\end{minipage} & \begin{minipage}[b]{\linewidth}\raggedright
SOLUTION
\end{minipage} & \begin{minipage}[b]{\linewidth}\raggedright
SUMMARY
\end{minipage} & \begin{minipage}[b]{\linewidth}\raggedright
SCHEDULE
\end{minipage} \\
\midrule\noalign{}
\endhead
\bottomrule\noalign{}
\endlastfoot
Alphabetic Listing of Keywords Starting with the Letter G & Status & & &
& & & \\
GAS -- Activate the Gas Phase in the Model & G & & & & & & \\
GASBEGIN -- Define Start of Annual Scheduling Section & & & & & & & \\
GASCONC -- Define the Initial Equilibration Coal Gas Concentration for
All Grid Blocks & & & & & & & \\
GASDENT -- Define Gas Density Temperature Coefficients & OPM Flow & & &
& & & \\
GASEND -- Define End of Annual Scheduling Section & & & & & & & \\
GASFCOMP -- Define Automatic Gas Compressors & & & & & & & \\
GASFDECR -- Define Field Gas Sales Contract Monthly Reduction & & & & &
& & \\
GASFDELC -- Define Gas Deliverability Calculation & & & & & & & \\
GASFIELD -- Define Gas Field Operations Options & & & & & & & \\
GASFTARG -- Define Field Gas Sales Contract Monthly Target & & & & & &
& \\
GASJT -- Define Gas Joule-Thomson Coefficient & OPM Flow & & & & & & \\
GASMONTH -- Define Start of Annual Scheduling Event & & & & & & & \\
GASPERIO -- Advance Simulation by Gas Contract Period & & & & & & & \\
GASSATC -- Define the Initial Equilibration Saturated Coal Gas
Concentration for All Grid Blocks & & & & & & & \\
GASVISCT - Define Gas Viscosity versus Temperature Functions & OPM Flow
& & & & & & \\
GASWAT -- Activate the Gas-Water Model Formulation & OPM Flow & & & & &
& \\
GASYEAR -- Advance Simulation by Gas Contract Year & & & & & & & \\
GCALECON -- Group Economic Criteria for Production Calorific Groups & &
& & & & & \\
GCONCAL -- Group Production Calorific Targets & & & & & & & \\
GCONENG -- Group Production Energy Targets & & & & & & & \\
GCONINJE -- Group Injection Targets and Constraints & & & & & & & \\
GCONPRI -- Group Production Priority Targets and Constraints & & & & & &
& \\
GCONPROD -- Group Production Targets and Constraints & & & & & & & \\
GCONSALE -- Define Group Sales Gas Production Targets and Constraints &
& & & & & & \\
GCONSUMP -- Define Group Gas Consumption and Gas Import Targets & & & &
& & & \\
GCONTOL -- Define Group Constraint Tolerance & & & & & & & \\
GCUTBACK -- Define Group Cutback Limits and Parameters & & & & & & & \\
GCUTBACT -- Define Group Tracer Cutback Limits and Parameters & & & & &
& & \\
GCVD -- Define Equilibration Coal Gas Concentration versus Depth Tables
& & & & & & & \\
GDCQ -- Define Group Multiple Daily Contract Quantities & & & & & & & \\
GDCQECON -- Group Economic Criteria for DCQ Production Groups & & & & &
& & \\
GDFILE -- Load a Grid File & & & & & & & \\
GDIMS -- Activate Instantaneous Gradient Option and Define Dimensions &
& & & & & & \\
GDORIENT - Define Grid Orientation Parameters & & & & & & & \\
GDRILPOT -- Define Group Potential Rates for Automatic Drilling & & & &
& & & \\
GECON -- Group Economic Criteria for Production Groups & & & & & & & \\
GECONT -- Group Tracer Economic Criteria for Production Groups & & & & &
& & \\
GEFAC -- Define Group Efficiency & & & & & & & \\
GETDATA -- Load and Assign Data Array from INIT or RESTART Files & & & &
& & & \\
GETGLOB -- Activate Loading of Global Grid Restart Data Option & & & & &
& & \\
GI - Define the Initial Equilibration Gi Values for All Grid Blocks & &
& & & & & \\
GIALL -- Define Gi Values and PVT Properties versus Pressure & & & & & &
& \\
GIMODEL -- Activate Gi Pseudo Compositional Option & & & & & & & \\
GINODE -- Define Gi Node Values & & & & & & & \\
GLIFTLIM -- Group Artificial Lift Constraints & & & & & & & \\
GLIFTOPT -- Define Group Gas Optimization Limits & & & & & & & \\
GMWSET - Export Well Status Vectors by Group to File & & & & & & & \\
GNETDP -- Group Network Pressure and Rate Controls & & & & & & & \\
GNETINJE -- Define Group Injection Network Configuration & & & & & &
& \\
GNETPUMP -- Standard Network Automatic Compressor and Pumps & & & & & &
& \\
GPMAINT -- Define Group Pressure Maintenance Targets and Controls & & &
& & & & \\
GRADGRUP -- Define Group History Match Gradient File Output & & & & & &
& \\
GRADRESV -- Define Solution Derivative History Match Gradient Output & &
& & & & & \\
GRADRFT -- Define RFT Derivative History Match Gradient Output & & & & &
& & \\
GRADWELL -- Define Well History Match Gradient File Output & & & & & &
& \\
GRAVCONS -- Re-Define Gravity Constant & & & & & & & \\
GRAVDR -- Activate Gravity Drainage and Imbibition for Dual Porosity
Model & & & & & & & \\
GRAVDRB - Activate Vertical Discretized Gravity Drainage and Imbibition
for Dual Porosity Model & & & & & & & \\
GRAVDRM - Activate Alternative Gravity Drainage and Imbibition for Dual
Porosity Model & & & & & & & \\
GRAVITY-- Define the Surface Oil, Water Gas Gravities for the Fluids & &
& & & & & \\
GRDREACH -- Define River and Grid Block Connections & & & & & & & \\
GRID - Define the Start of the GRID Section of Keywords & & & & & & & \\
GRIDFILE -- Set the Grid File Output Options & & & & & & & \\
GRIDOPTS - Grid Processing Options & & & & & & & \\
GRIDUNIT -- Define the Grid Units & & & & & & & \\
GRUPMAST -- Define Master and Slave Groups & & & & & & & \\
GRUPNET -- Define Group Standard Network Parameters & & & & & & & \\
GRUPRIG -- Group Drilling and Workover Rig Specifications & & & & & &
& \\
GRUPSLAV -- Define Slave Groups in Slave Reservoirs & & & & & & & \\
GRUPTARG -- Modify Group Targets and Constraints Values & & & & & & & \\
GRUPTREE -- Define Group Tree Hierarchy & & & & & & & \\
GSATINJE -- Define Group Satellite Injection Rates & & & & & & & \\
GSATPROD -- Define Group Satellite Production Rates & & & & & & & \\
GSEPCOND -- Assign Group Separators & & & & & & & \\
GSF - Gas Saturation Tables versus Gas Saturation (Gas-Water and
CO2STORE Systems) & OPM Flow & & & & & & \\
GSSCPTST -- Perform Sustainable Capacity Test & & & & & & & \\
GSWINGF -- Define Group Multiple Gas Contract Parameters & & & & & &
& \\
GTADD -- Add a Constant to a Group Target or Constraint & & & & & & & \\
GTMULT -- Multiply Group Target or Constraint by a Constant & & & & & &
& \\
GUIDECAL -- Scale Guide Rates Based on Gas Calorific Value & & & & & &
& \\
GUIDERAT -- Define Group Guide Rate Formula & & & & & & & \\
GUPFREQ -- Instantaneous Gradient Option Update Frequency & & & & & &
& \\
GWRTWCV -- Instantaneous Gradient Option Well Variables & & & & & & & \\
\end{longtable}

\begin{longtable}[]{@{}
  >{\raggedright\arraybackslash}p{(\linewidth - 14\tabcolsep) * \real{0.1250}}
  >{\raggedright\arraybackslash}p{(\linewidth - 14\tabcolsep) * \real{0.1250}}
  >{\raggedright\arraybackslash}p{(\linewidth - 14\tabcolsep) * \real{0.1250}}
  >{\raggedright\arraybackslash}p{(\linewidth - 14\tabcolsep) * \real{0.1250}}
  >{\raggedright\arraybackslash}p{(\linewidth - 14\tabcolsep) * \real{0.1250}}
  >{\raggedright\arraybackslash}p{(\linewidth - 14\tabcolsep) * \real{0.1250}}
  >{\raggedright\arraybackslash}p{(\linewidth - 14\tabcolsep) * \real{0.1250}}
  >{\raggedright\arraybackslash}p{(\linewidth - 14\tabcolsep) * \real{0.1250}}@{}}
\toprule\noalign{}
\begin{minipage}[b]{\linewidth}\raggedright
RUNSPEC
\end{minipage} & \begin{minipage}[b]{\linewidth}\raggedright
GRID
\end{minipage} & \begin{minipage}[b]{\linewidth}\raggedright
EDIT
\end{minipage} & \begin{minipage}[b]{\linewidth}\raggedright
PROPS
\end{minipage} & \begin{minipage}[b]{\linewidth}\raggedright
REGIONS
\end{minipage} & \begin{minipage}[b]{\linewidth}\raggedright
SOLUTION
\end{minipage} & \begin{minipage}[b]{\linewidth}\raggedright
SUMMARY
\end{minipage} & \begin{minipage}[b]{\linewidth}\raggedright
SCHEDULE
\end{minipage} \\
\midrule\noalign{}
\endhead
\bottomrule\noalign{}
\endlastfoot
Alphabetic Listing of Keywords Starting with the Letter H & Status & & &
& & & \\
\hyperref[0.0.0.H2SOL]{H2SOL -- Activate Dissolved H2 in the Water
Phase} & OPM Flow & & & & & & \\
\hyperref[0.0.0.H2STORE]{H2STORE -- Activate the H2 Storage Model} & OPM
Flow & & & & & & \\
HALFTRAN -- Activate Half Block Transmissibility Calculations & H & & &
& & & \\
HA -- History Match End-Point Gradient Additive Modifier & & & & & &
& \\
HAKRG & & & & & & & \\
HAKRGR & & & & & & & \\
HAKRO & & & & & & & \\
HAKRORG & & & & & & & \\
HAKRORW & & & & & & & \\
HAKRW & & & & & & & \\
HAKRWR & & & & & & & \\
HAPCG & & & & & & & \\
HAPCG & & & & & & & \\
HASGLPC & & & & & & & \\
HASOGCR & & & & & & & \\
HASOWCR & & & & & & & \\
HASWL & & & & & & & \\
HASWLPC & & & & & & & \\
HBNUM -- Define Herschel-Bulkley Region Numbers & & & & & & & \\
HDISP -- Define Tracer Mechanical Dispersivity Parameters & & & & & &
& \\
HEATCR -- Define Reservoir Rock Heat Capacity for All Cells & & & & & &
& \\
HEATCRT -- Define Reservoir Rock Heat Capacity Temperature Dependence
for All Cells & & & & & & & \\
HM -- History Match End-Point Gradient Multiplicative Modifier & & & & &
& & \\
HMKRG & & & & & & & \\
HMKRGR & & & & & & & \\
HMKRO & & & & & & & \\
HMKRORG & & & & & & & \\
HMKRORW & & & & & & & \\
HMKRW & & & & & & & \\
HMKRWR & & & & & & & \\
HMPCG & & & & & & & \\
HMPCG & & & & & & & \\
HMSGLPC & & & & & & & \\
HMSOGCR & & & & & & & \\
HMSOWCR & & & & & & & \\
HMSWL & & & & & & & \\
HMSWLPC & & & & & & & \\
HM -- History Match Region Gradient Parameters & & & & & & & \\
HMKRG & & & & & & & \\
HMKRGR & & & & & & & \\
HMKRO & & & & & & & \\
HMKRORG & & & & & & & \\
HMKRORW & & & & & & & \\
HMKRW & & & & & & & \\
HMKRWR & & & & & & & \\
HMPERMX & & & & & & & \\
HMPERMY & & & & & & & \\
HMPERMZ & & & & & & & \\
HMPORVM & & & & & & & \\
HMPRMXY & & & & & & & \\
HMSGCR & & & & & & & \\
HMSGL & & & & & & & \\
HMSGLPC & & & & & & & \\
HMSIGMA & & & & & & & \\
HMSOGCR & & & & & & & \\
HMSOWCR & & & & & & & \\
HMSWCR & & & & & & & \\
HMSWL & & & & & & & \\
HMSWLPC & & & & & & & \\
HMTRANX & & & & & & & \\
HMTRANY & & & & & & & \\
HMTRANZ & & & & & & & \\
HMTRNXY & & & & & & & \\
HMAQUCT -- History Match Carter-Tracy Aquifer Gradient Parameters & & &
& & & & \\
HMAQUFET -- History Match Fetkovich Aquifer Gradient Parameters & & & &
& & & \\
HMAQUNUM - History Match Numerical Aquifer Gradient Parameters & & & & &
& & \\
HMDIMS -- Define History Match Gradient Parameter Dimensions & & & & & &
& \\
HMFAULTS -- History Match Fault Gradient Parameters & & & & & & & \\
HMMLAQUN -- History Match Numerical Aquifer Gradient Multipliers & & & &
& & & \\
HMMLCTAQ -- History Match Carter-Tracy Aquifer Gradient Multipliers & &
& & & & & \\
HMMLFTAQ -- History Match Fetkovich Aquifer Gradient Multipliers & & & &
& & & \\
HMMLT -- History Match Grid Permeability Gradient Cumulative Multipliers
& & & & & & & \\
HMMLTPR & & & & & & & \\
HMMLTPTH & & & & & & & \\
HMMLTPX & & & & & & & \\
HMMLTPXY & & & & & & & \\
HMMLTPY & & & & & & & \\
HMMLTPZ & & & & & & & \\
HMMLTWCN -- History Match Well Connection and Skin Multipliers & & & & &
& & \\
HMMMREGT - History Match Region Transmissibility Gradient Cumulative
Multipliers & & & & & & & \\
HMMROCK -- History Match Rock Compressibility Gradient Cumulative
Multipliers & & & & & & & \\
HMMROCKT -- History Match Rock Compaction Gradient Cumulative
Multipliers & & & & & & & \\
HMMULRGT -- History Match Region Transmissibility Parameters & & & & & &
& \\
HMMULT -- History Match Grid Transmissibility \& Pore Volume Gradient
Cumulative Multipliers & & & & & & & \\
HMULTPV & & & & & & & \\
HMMULTR & & & & & & & \\
HMMULTTH & & & & & & & \\
HMMULTX & & & & & & & \\
HMMULTXY & & & & & & & \\
HMMULTY & & & & & & & \\
HMMULTZ & & & & & & & \\
HMMULTFT -- History Match Fault Transmissibility Gradient Cumulative
Multipliers & & & & & & & \\
HMMULTSG -- History Match Dual porosity Sigma Gradient Cumulative
Multipliers & & & & & & & \\
HMPROPS -- History Match End-Point Section Start & & & & & & & \\
HMROCK -- History Match Rock Compressibility Gradient Parameters & & & &
& & & \\
HMROCKT -- History Match Rock Compaction Gradient Parameters & & & & & &
& \\
HMRREF -- History Match Rock Table Reference Pressure Values & & & & & &
& \\
HMWELCON -- History Match Well Connection and Skin Parameters & & & & &
& & \\
HMWPIMLT -- History Match Well Productivity Index Parameters & & & & & &
& \\
HRFIN - Define the Ratio of LGR Grid Blocks in the R-Direction & & & & &
& & \\
HWKRO -- End-Point Scaling of Grid Cell Kro(Swl) (High Salinity and
Water Wet) & & & & & & & \\
HWKRORG -- End-Point Scaling of Grid Cell Kro(Sgcr) (High Salinity and
Water Wet) & & & & & & & \\
HWKRORW -- End-Point Scaling of Grid Cell Kro(Swcr) (High Salinity and
Water Wet) & & & & & & & \\
HWKRW -- End-Point Scaling of Grid Cell Krw(Sw =1.0) (High Salinity and
Water Wet) & & & & & & & \\
HWKRWR -- End-Point Scaling of Grid Cell KRWR(Sw=1.0) (High Salinity and
Water Wet) & & & & & & & \\
HWPCW -- End-Point Scaling of Grid Cell Water Capillary Pressure (High
Salinity and Water Wet) & & & & & & & \\
HWSNUM -- Define the Saturation Table Region Numbers (High Salinity and
Water Wet) & & & & & & & \\
HWSOGCR -- End-Point Scaling Grid Cell SOGCR (High Salinity and Water
Wet) & & & & & & & \\
HWSOWCR -- End-Point Scaling Grid Cell SOWCR (High Salinity and Water
Wet) & & & & & & & \\
HWSWCR -- End-Point Scaling Grid Cell SWCR (High Salinity and Water Wet)
& & & & & & & \\
HWSWL -- End-Point Scaling Grid Cell SWL (High Salinity and Water Wet) &
& & & & & & \\
HWSWLPC -- End-Point Scaling Grid Cell SWLPC (High Salinity and Water
Wet) & & & & & & & \\
HWSWU -- End-Point Scaling Grid Cell SWU (High Salinity and Water Wet) &
& & & & & & \\
HXFIN - Define the Ratio of LGR Grid Blocks in the X-Direction & & & & &
& & \\
HYDRHEAD -- Define Hydraulic Head Output Reference Data & & & & & & & \\
HYFIN - Define the Ratio of LGR Grid Blocks in the Y-Direction & & & & &
& & \\
HYMOBGDR -- Activate Carlson and Killough Alternative Drainage
Hysteresis Option & & & & & & & \\
HYST -- Activate the Hysteresis Option (Retired) & & & & & & & \\
HYSTCHCK - Activate Hysteresis Inhibition and Drainage End-Point
Validation & & & & & & & \\
HZFIN - Define the Ratio of LGR Grid Blocks in the Z-Direction & & & & &
& & \\
\end{longtable}

\begin{longtable}[]{@{}
  >{\raggedright\arraybackslash}p{(\linewidth - 14\tabcolsep) * \real{0.1250}}
  >{\raggedright\arraybackslash}p{(\linewidth - 14\tabcolsep) * \real{0.1250}}
  >{\raggedright\arraybackslash}p{(\linewidth - 14\tabcolsep) * \real{0.1250}}
  >{\raggedright\arraybackslash}p{(\linewidth - 14\tabcolsep) * \real{0.1250}}
  >{\raggedright\arraybackslash}p{(\linewidth - 14\tabcolsep) * \real{0.1250}}
  >{\raggedright\arraybackslash}p{(\linewidth - 14\tabcolsep) * \real{0.1250}}
  >{\raggedright\arraybackslash}p{(\linewidth - 14\tabcolsep) * \real{0.1250}}
  >{\raggedright\arraybackslash}p{(\linewidth - 14\tabcolsep) * \real{0.1250}}@{}}
\toprule\noalign{}
\begin{minipage}[b]{\linewidth}\raggedright
RUNSPEC
\end{minipage} & \begin{minipage}[b]{\linewidth}\raggedright
GRID
\end{minipage} & \begin{minipage}[b]{\linewidth}\raggedright
EDIT
\end{minipage} & \begin{minipage}[b]{\linewidth}\raggedright
PROPS
\end{minipage} & \begin{minipage}[b]{\linewidth}\raggedright
REGIONS
\end{minipage} & \begin{minipage}[b]{\linewidth}\raggedright
SOLUTION
\end{minipage} & \begin{minipage}[b]{\linewidth}\raggedright
SUMMARY
\end{minipage} & \begin{minipage}[b]{\linewidth}\raggedright
SCHEDULE
\end{minipage} \\
\midrule\noalign{}
\endhead
\bottomrule\noalign{}
\endlastfoot
Alphabetic Listing of Keywords Starting with the Letter I & Status & & &
& & & \\
IHOST -- Assign LGRs to Parallel Process Number & I & & & & & & \\
IKRG -- End-Point Scaling of Grid Cell Krg(Sgu) (Imbibition) & & & & & &
& \\
IKRGX & & & & & & & \\
IKRGX- & & & & & & & \\
IKRGY & & & & & & & \\
IKRGY- & & & & & & & \\
IKRGZ & & & & & & & \\
IKRGZ- & & & & & & & \\
IKRGR -- End-Point Scaling of Grid Cell Krgr(1-Sogcr) (Imbibition) & & &
& & & & \\
IKRGRX & & & & & & & \\
IKRGRX- & & & & & & & \\
IKRGRY & & & & & & & \\
IKRGRY- & & & & & & & \\
IKRGRZ & & & & & & & \\
IKRGRZ- & & & & & & & \\
IKRO -- End-Point Scaling of Grid Cell Kro(Swl) (Imbibition) & & & & & &
& \\
IKROX & & & & & & & \\
IKROX- & & & & & & & \\
IKROY & & & & & & & \\
IKROY- & & & & & & & \\
IKROZ & & & & & & & \\
KROZ- & & & & & & & \\
IKRORG -- End-Point Scaling of Grid Cell Kro(Sgcr) (Imbibition) & & & &
& & & \\
IKRORGX & & & & & & & \\
IKRORGX- & & & & & & & \\
IKRORGY & & & & & & & \\
IKRORGY- & & & & & & & \\
IKRORGZ & & & & & & & \\
IKRORGZ- & & & & & & & \\
IKRORW -- End-Point Scaling of Grid Cell Kro(Swcr) (Imbibition) & & & &
& & & \\
IKRORWX & & & & & & & \\
IKRORWX- & & & & & & & \\
IKRORWY & & & & & & & \\
IKRORWY- & & & & & & & \\
IKRORWZ & & & & & & & \\
IKRORWZ- & & & & & & & \\
IKRW -- End-Point Scaling of Grid Cell Krw(Sw =1.0) (Imbibition) & & & &
& & & \\
IKRWX & & & & & & & \\
IKRWX- & & & & & & & \\
IKRWY & & & & & & & \\
IKRWY- & & & & & & & \\
IKRWZ & & & & & & & \\
IKRWZ- & & & & & & & \\
IKRWR -- End-Point Scaling of Grid Cell KRWR(Sowcr) (Imbibition) & & & &
& & & \\
IKRWRX & & & & & & & \\
IKRWRX- & & & & & & & \\
IKRWRY & & & & & & & \\
IKRWRY- & & & & & & & \\
IKRWRZ & & & & & & & \\
IKRWRZ- & & & & & & & \\
IMBNUM -- Define the Imbibition Saturation Table Region Numbers & & & &
& & & \\
IMBNUMX & & & & & & & \\
‍IMBNUMX- & & & & & & & \\
IMBNUMY & & & & & & & \\
IMBNUMY- & & & & & & & \\
‍IMBNUMZ & & & & & & & \\
IMBNUMZ- & & & & & & & \\
IMBNUMMF -- Define the Imbibition Saturation Table Region Numbers
(Matrix-Fracture) & & & & & & & \\
IMKRVD -- Imbibition Relative Permeability End-Points versus Depth
Functions & & & & & & & \\
IMKRVDX & & & & & & & \\
IMKRVDX- & & & & & & & \\
IMKRVDY & & & & & & & \\
IMKRVDY- & & & & & & & \\
IMKRVDX & & & & & & & \\
IMKRVDZ- & & & & & & & \\
IMPCVD -- Imbibition Maximum Capillary Pressure versus Depth Functions &
& & & & & & \\
IMPES -- Activate Implicit Pressure Explicit Saturation Solution Option
& & & & & & & \\
IMPLICIT -- Activate Fully Implicit Solution Option & & & & & & & \\
IMPORT -- Import Grid File Data at the Current Position & & & & & & & \\
IMPTVD -- Imbibition Relative Permeability Saturation End-Points versus
Depth & & & & & & & \\
IMPTVDX & & & & & & & \\
IMPTVDX- & & & & & & & \\
IMPTVDY & & & & & & & \\
IMPTVDY- & & & & & & & \\
IMPTVDZ & & & & & & & \\
IMPTVDZ- & & & & & & & \\
IMSPCVD -- Imbibition Capillary Pressure Connate Saturations versus
Depth & & & & & & & \\
INCLUDE -- Load Another Data File at the Current Position & & & & & &
& \\
INIT -- Activate the INIT File Option & & & & & & & \\
INRAD -- Define the Inner Radius of a Radial Grid & & & & & & & \\
INSPEC -- Activate the INSPEC File Option & & & & & & & \\
INTPC -- Activate Dual Porosity Integrated Capillary Pressure Option & &
& & & & & \\
IONROCK -- Define the Ion Exchange Capacity for All the Cells & & & & &
& & \\
IONXROCK - Define Ion Exchange Constant by Saturation Table Regions & &
& & & & & \\
IONXSURF - Define Surfactant Ion Exchange Constant by Saturation Table
Regions & & & & & & & \\
IPCG -- End-Point Scaling of Grid Cell Gas Capillary Pressure
(Imbibition) & & & & & & & \\
IPCW -- End-Point Scaling of Grid Cell Water Capillary Pressure
(Imbibition) & & & & & & & \\
ISGCR -- End-Point Scaling of Grid Cell Critical Gas Saturation
(Imbibition) & & & & & & & \\
ISGCRX & & & & & & & \\
ISGCRX- & & & & & & & \\
ISGCRY & & & & & & & \\
ISGCRY- & & & & & & & \\
ISGCRZ & & & & & & & \\
ISGCRZ- & & & & & & & \\
ISGL -- End-Point Scaling of Grid Cell Connate Gas Saturation
(Imbibition) & & & & & & & \\
ISGLX & & & & & & & \\
ISGLX- & & & & & & & \\
ISGLY & & & & & & & \\
ISGLY- & & & & & & & \\
ISGLZ & & & & & & & \\
ISGLZ- & & & & & & & \\
ISGLPC -- End-Point Scaling of Grid Cell Capillary Pressure Connate Gas
Saturation (Imbibition) & & & & & & & \\
ISGU -- End-Point Scaling of Grid Cell Maximum Gas Saturation
(Imbibition) & & & & & & & \\
ISGUX & & & & & & & \\
ISGUX- & & & & & & & \\
ISGUY & & & & & & & \\
ISGUY- & & & & & & & \\
ISGUZ & & & & & & & \\
ISGUZ- & & & & & & & \\
ISOGCR -- End-Point Scaling of Grid Cell Critical Oil Saturation to Gas
(Imbibition) & & & & & & & \\
ISOGCRX & & & & & & & \\
ISOGCRX- & & & & & & & \\
ISOGCRY & & & & & & & \\
ISOGCRY- & & & & & & & \\
ISOGCRZ & & & & & & & \\
ISOGCRZ- & & & & & & & \\
ISOLNUM -- Define the Independent Reservoir Regions & & & & & & & \\
ISOWCR -- End-Point Scaling of Grid Cell Critical Oil Saturation to
Water (Imbib.) & & & & & & & \\
ISOWCRX & & & & & & & \\
ISOWCRX- & & & & & & & \\
ISOWCRY & & & & & & & \\
ISOWCRY- & & & & & & & \\
ISOWCRZ & & & & & & & \\
ISOWCRZ- & & & & & & & \\
ISWCR -- End-Point Scaling of Grid Cell Critical Water Saturation
(Imbibition) & & & & & & & \\
ISWCRX & & & & & & & \\
ISWCRX- & & & & & & & \\
ISWCRY & & & & & & & \\
ISWCRY- & & & & & & & \\
ISWCRZ & & & & & & & \\
ISWCRZ- & & & & & & & \\
ISWL -- End-Point Scaling of Grid Cell Connate Water Saturation
(Imbibition) & & & & & & & \\
ISWLX & & & & & & & \\
ISWLX- & & & & & & & \\
ISWLY & & & & & & & \\
ISWLY- & & & & & & & \\
ISWLZ & & & & & & & \\
ISWLZ- & & & & & & & \\
ISWLPC -- End-Point Scaling Grid Cell Capillary Pressure Connate Water
Saturations (Imbibition) & & & & & & & \\
ISWU -- End-Point Scaling of Grid Cell Maximum Water Saturation
(Imbibition) & & & & & & & \\
ISWUX & & & & & & & \\
ISWUX- & & & & & & & \\
ISWUY & & & & & & & \\
ISWUY- & & & & & & & \\
ISWUZ & & & & & & & \\
ISWUZ- & & & & & & & \\
\end{longtable}

\begin{longtable}[]{@{}
  >{\raggedright\arraybackslash}p{(\linewidth - 14\tabcolsep) * \real{0.1250}}
  >{\raggedright\arraybackslash}p{(\linewidth - 14\tabcolsep) * \real{0.1250}}
  >{\raggedright\arraybackslash}p{(\linewidth - 14\tabcolsep) * \real{0.1250}}
  >{\raggedright\arraybackslash}p{(\linewidth - 14\tabcolsep) * \real{0.1250}}
  >{\raggedright\arraybackslash}p{(\linewidth - 14\tabcolsep) * \real{0.1250}}
  >{\raggedright\arraybackslash}p{(\linewidth - 14\tabcolsep) * \real{0.1250}}
  >{\raggedright\arraybackslash}p{(\linewidth - 14\tabcolsep) * \real{0.1250}}
  >{\raggedright\arraybackslash}p{(\linewidth - 14\tabcolsep) * \real{0.1250}}@{}}
\toprule\noalign{}
\begin{minipage}[b]{\linewidth}\raggedright
RUNSPEC
\end{minipage} & \begin{minipage}[b]{\linewidth}\raggedright
GRID
\end{minipage} & \begin{minipage}[b]{\linewidth}\raggedright
EDIT
\end{minipage} & \begin{minipage}[b]{\linewidth}\raggedright
PROPS
\end{minipage} & \begin{minipage}[b]{\linewidth}\raggedright
REGIONS
\end{minipage} & \begin{minipage}[b]{\linewidth}\raggedright
SOLUTION
\end{minipage} & \begin{minipage}[b]{\linewidth}\raggedright
SUMMARY
\end{minipage} & \begin{minipage}[b]{\linewidth}\raggedright
SCHEDULE
\end{minipage} \\
\midrule\noalign{}
\endhead
\bottomrule\noalign{}
\endlastfoot
Alphabetic Listing of Keywords Starting with the Letter J & Status & & &
& & & \\
JFUNC - Activate the Leverett J-function Option & J & & & & & & \\
JFUNCR - Activate the Leverett J-function Saturation Table Option & & &
& & & & \\
\end{longtable}

\begin{longtable}[]{@{}
  >{\raggedright\arraybackslash}p{(\linewidth - 14\tabcolsep) * \real{0.1250}}
  >{\raggedright\arraybackslash}p{(\linewidth - 14\tabcolsep) * \real{0.1250}}
  >{\raggedright\arraybackslash}p{(\linewidth - 14\tabcolsep) * \real{0.1250}}
  >{\raggedright\arraybackslash}p{(\linewidth - 14\tabcolsep) * \real{0.1250}}
  >{\raggedright\arraybackslash}p{(\linewidth - 14\tabcolsep) * \real{0.1250}}
  >{\raggedright\arraybackslash}p{(\linewidth - 14\tabcolsep) * \real{0.1250}}
  >{\raggedright\arraybackslash}p{(\linewidth - 14\tabcolsep) * \real{0.1250}}
  >{\raggedright\arraybackslash}p{(\linewidth - 14\tabcolsep) * \real{0.1250}}@{}}
\toprule\noalign{}
\begin{minipage}[b]{\linewidth}\raggedright
RUNSPEC
\end{minipage} & \begin{minipage}[b]{\linewidth}\raggedright
GRID
\end{minipage} & \begin{minipage}[b]{\linewidth}\raggedright
EDIT
\end{minipage} & \begin{minipage}[b]{\linewidth}\raggedright
PROPS
\end{minipage} & \begin{minipage}[b]{\linewidth}\raggedright
REGIONS
\end{minipage} & \begin{minipage}[b]{\linewidth}\raggedright
SOLUTION
\end{minipage} & \begin{minipage}[b]{\linewidth}\raggedright
SUMMARY
\end{minipage} & \begin{minipage}[b]{\linewidth}\raggedright
SCHEDULE
\end{minipage} \\
\midrule\noalign{}
\endhead
\bottomrule\noalign{}
\endlastfoot
Alphabetic Listing of Keywords Starting with the Letter K & Status & & &
& & & \\
KRG -- End-Point Scaling of Grid Cell Krg(Sgu) (Drainage) & K & & & & &
& \\
KRGX & & & & & & & \\
KRGX- & & & & & & & \\
KRGY & & & & & & & \\
KRGY- & & & & & & & \\
KRGZ & & & & & & & \\
KRGZ- & & & & & & & \\
KRGR -- End-Point Scaling of Grid Cell Krgr(1-Sogcr) (Drainage) & & & &
& & & \\
KRGRX & & & & & & & \\
KRGRX- & & & & & & & \\
KRGRY & & & & & & & \\
KRGRY- & & & & & & & \\
KRGRZ & & & & & & & \\
KRGRZ- & & & & & & & \\
KRNUM -- Define the Directional Saturation Table Region Numbers & & & &
& & & \\
KRNUMR & & & & & & & \\
KRNUMR- & & & & & & & \\
KRNUMT & & & & & & & \\
KRNUMT- & & & & & & & \\
KRNUMX & & & & & & & \\
KRNUMX- & & & & & & & \\
KRNUMY & & & & & & & \\
KRNUMY- & & & & & & & \\
KRNUMZ & & & & & & & \\
KRNUMZ- & & & & & & & \\
KRNUMMF -- Define the Saturation Table Region Numbers (Matrix-Fracture)
& & & & & & & \\
KRO -- End-Point Scaling of Grid Cell Kro(Swl) (Drainage) & & & & & &
& \\
KROX & & & & & & & \\
KROX- & & & & & & & \\
KROY & & & & & & & \\
KROY- & & & & & & & \\
KROZ & & & & & & & \\
KROZ- & & & & & & & \\
KRORG -- End-Point Scaling of Grid Cell Kro(Sgcr) (Drainage) & & & & & &
& \\
KRORGX & & & & & & & \\
KRORGX- & & & & & & & \\
KRORGY & & & & & & & \\
KRORGY- & & & & & & & \\
KRORGZ & & & & & & & \\
KRORGZ- & & & & & & & \\
KRORW -- End-Point Scaling of Grid Cell Kro(Swcr) (Drainage) & & & & & &
& \\
KRORWX & & & & & & & \\
KRORWX- & & & & & & & \\
KRORWY & & & & & & & \\
KRORWY- & & & & & & & \\
KRORWZ & & & & & & & \\
KRORWZ- & & & & & & & \\
KRW -- End-Point Scaling of Grid Cell Krw(Sw =1.0) (Drainage) & & & & &
& & \\
KRWX & & & & & & & \\
KRWX- & & & & & & & \\
KRWY & & & & & & & \\
KRWY- & & & & & & & \\
KRWZ & & & & & & & \\
KRWZ- & & & & & & & \\
KRWR -- End-Point Scaling of Grid Cell KRWR(Sowcr) (Drainage) & & & & &
& & \\
KRWRX & & & & & & & \\
KRWRX- & & & & & & & \\
KRWRY & & & & & & & \\
KRWRY- & & & & & & & \\
KRWRZ & & & & & & & \\
KRWRZ- & & & & & & & \\
\end{longtable}

\begin{longtable}[]{@{}
  >{\raggedright\arraybackslash}p{(\linewidth - 14\tabcolsep) * \real{0.1250}}
  >{\raggedright\arraybackslash}p{(\linewidth - 14\tabcolsep) * \real{0.1250}}
  >{\raggedright\arraybackslash}p{(\linewidth - 14\tabcolsep) * \real{0.1250}}
  >{\raggedright\arraybackslash}p{(\linewidth - 14\tabcolsep) * \real{0.1250}}
  >{\raggedright\arraybackslash}p{(\linewidth - 14\tabcolsep) * \real{0.1250}}
  >{\raggedright\arraybackslash}p{(\linewidth - 14\tabcolsep) * \real{0.1250}}
  >{\raggedright\arraybackslash}p{(\linewidth - 14\tabcolsep) * \real{0.1250}}
  >{\raggedright\arraybackslash}p{(\linewidth - 14\tabcolsep) * \real{0.1250}}@{}}
\toprule\noalign{}
\begin{minipage}[b]{\linewidth}\raggedright
RUNSPEC
\end{minipage} & \begin{minipage}[b]{\linewidth}\raggedright
GRID
\end{minipage} & \begin{minipage}[b]{\linewidth}\raggedright
EDIT
\end{minipage} & \begin{minipage}[b]{\linewidth}\raggedright
PROPS
\end{minipage} & \begin{minipage}[b]{\linewidth}\raggedright
REGIONS
\end{minipage} & \begin{minipage}[b]{\linewidth}\raggedright
SOLUTION
\end{minipage} & \begin{minipage}[b]{\linewidth}\raggedright
SUMMARY
\end{minipage} & \begin{minipage}[b]{\linewidth}\raggedright
SCHEDULE
\end{minipage} \\
\midrule\noalign{}
\endhead
\bottomrule\noalign{}
\endlastfoot
Alphabetic Listing of Keywords Starting with the Letter L & Status & & &
& & & \\
LAB - Activate the Laboratory System of Units for the Model & L & & & &
& & \\
LANGMPL -- Define Langmuir Pressure Grid Cell Multiplier & & & & & &
& \\
LANGMUIR -- Langmuir Adsorption Isotherm Tables & & & & & & & \\
LANGSOLV -- Langmuir Adsorption Isotherm Solvent Tables & & & & & & & \\
LCUNIT -- Define Linear Combination Rate and Volume Units & & & & & &
& \\
LGR -- Define Local Grid Refinement Dimensions and Parameters & & & & &
& & \\
LGRCOPY -- Activate Local Grid Refinement Inheritance & & & & & & & \\
LGRFREE -- Activate Local Grid Refinement Independent Time Steps & & & &
& & & \\
LGRLOCK -- Deactivate Local Grid Refinement Independent Time Steps & & &
& & & & \\
LGROFF -- Deactivate a Local Grid Refinement & & & & & & & \\
LGRON -- Activate a Local Grid Refinement & & & & & & & \\
LICENSES -- Define Required Licenses for Run & & & & & & & \\
LIFTOPT -- Activate Gas Lift Optimization & & & & & & & \\
LINCOM -- Define Linear Combination Coefficients & & & & & & & \\
LINKPERM - Assign Cell Permeabilities to Cell Faces & & & & & & & \\
LIVEOIL -- Activate the Live Oil Phase (Oil with Free and Dissolved Gas)
& & & & & & & \\
LKRO -- End-Point Scaling of Grid Cell Kro(Swl) (Low Salinity and Oil
Wet) & & & & & & & \\
LKRORG -- End-Point Scaling of Grid Cell Kro(Sgcr) (Low Salinity and Oil
Wet) & & & & & & & \\
LKRORW -- End-Point Scaling of Grid Cell Kro(Swcr) (Low Salinity and Oil
Wet) & & & & & & & \\
LKRW -- End-Point Scaling of Grid Cell Krw(Sw =1.0) (Low Salinity and
Oil Wet) & & & & & & & \\
LKRWR -- End-Point Scaling of Grid Cell KRWR(Sowcr) (Low Salinity and
Oil Wet) & & & & & & & \\
LOAD -- Load a SAVE File for a Fast Restart & & & & & & & \\
LOWSALT -- Activate the Low Salt Brine Phase in the Brine Model & & & &
& & & \\
LPCW -- End-Point Scaling of Grid Cell Water Capillary Pressure (Low
Salinity and Oil Wet) & & & & & & & \\
LSALTFNC - Define Low Salt Weighting Factors versus Salt Concentration
Functions & & & & & & & \\
LSLTWNUM -- Define the Low Salt Water Wet Saturation Table Region
Numbers & & & & & & & \\
LSNUM -- Define the Low Salt Oil Wet Saturation Table Region Numbers & &
& & & & & \\
LSOGCR -- End-Point Scaling Grid Cell SOGCR (Low Salinity and Oil Wet) &
& & & & & & \\
LSOWCR -- End-Point Scaling Grid Cell SOWCR (Low Salinity and Oil Wet) &
& & & & & & \\
LSWCR -- End-Point Scaling Grid Cell SWCR (Low Salinity and Oil Wet) & &
& & & & & \\
LSWL -- End-Point Scaling Grid Cell SWL (Low Salinity and Oil Wet) & & &
& & & & \\
LSWLPC -- End-Point Scaling Grid Cell SWLPC (Low Salinity and Oil Wet) &
& & & & & & \\
LSWU -- End-Point Scaling Grid Cell SWU (Low Salinity and Oil Wet) & & &
& & & & \\
LTOSIGMA - Dual Porosity Viscous Displacement Sigma Parameters & & & & &
& & \\
LWKRO -- End-Point Scaling of Grid Cell Kro(Swl) (Low Salinity and Water
Wet) & & & & & & & \\
LWKRORG -- End-Point Scaling of Grid Cell Kro(Sgcr) (Low Salinity and
Water Wet) & & & & & & & \\
LWKRORW -- End-Point Scaling of Grid Cell Kro(Swcr) (Low Salinity and
Water Wet) & & & & & & & \\
LWKRW -- End-Point Scaling of Grid Cell Krw(Sw =1.0) (Low Salinity and
Water Wet) & & & & & & & \\
LWKRWR -- End-Point Scaling of Grid Cell KRWR(Sw=1.0) (Low Salinity and
Water Wet) & & & & & & & \\
LWPCW -- End-Point Scaling of Grid Cell Water Capillary Pressure (Low
Salinity and Water Wet) & & & & & & & \\
LWSLTNUM -- Define the Low Salt Oil Wet Saturation Table Region Numbers
& & & & & & & \\
LWSNUM -- Define the Low Salt Water Wet Saturation Table Region Numbers
& & & & & & & \\
LWSOGCR -- End-Point Scaling Grid Cell SOGCR (Low Salinity and Water
Wet) & & & & & & & \\
LWSOWCR -- End-Point Scaling Grid Cell SOWCR (Low Salinity and Water
Wet) & & & & & & & \\
LWSWCR -- End-Point Scaling Grid Cell SWCR (Low Salinity and Water Wet)
& & & & & & & \\
LWSWL -- End-Point Scaling Grid Cell SWL (Low Salinity and Water Wet) &
& & & & & & \\
LWSWLPC -- End-Point Scaling Grid Cell SWLPC (Low Salinity and Water
Wet) & & & & & & & \\
LWSWU -- End-Point Scaling Grid Cell SWU (Low Salinity and Water Wet) &
& & & & & & \\
LX - Dual Porosity Viscous Displacement X Direction Matrix Size for All
Cells & & & & & & & \\
LXFIN -- Define Logarithmic LGR Grid Block Spacing in the X-Direction &
& & & & & & \\
LY - Dual Porosity Viscous Displacement Y Direction Matrix Size for All
Cells & & & & & & & \\
LYFIN -- Define Logarithmic LGR Grid Block Spacing in the Y-Direction &
& & & & & & \\
LZ - Dual Porosity Viscous Displacement Z Direction Matrix Size for All
Cells & & & & & & & \\
LZFIN -- Define Logarithmic LGR Grid Block Spacing in the Z-Direction &
& & & & & & \\
\end{longtable}

\begin{longtable}[]{@{}
  >{\raggedright\arraybackslash}p{(\linewidth - 14\tabcolsep) * \real{0.1250}}
  >{\raggedright\arraybackslash}p{(\linewidth - 14\tabcolsep) * \real{0.1250}}
  >{\raggedright\arraybackslash}p{(\linewidth - 14\tabcolsep) * \real{0.1250}}
  >{\raggedright\arraybackslash}p{(\linewidth - 14\tabcolsep) * \real{0.1250}}
  >{\raggedright\arraybackslash}p{(\linewidth - 14\tabcolsep) * \real{0.1250}}
  >{\raggedright\arraybackslash}p{(\linewidth - 14\tabcolsep) * \real{0.1250}}
  >{\raggedright\arraybackslash}p{(\linewidth - 14\tabcolsep) * \real{0.1250}}
  >{\raggedright\arraybackslash}p{(\linewidth - 14\tabcolsep) * \real{0.1250}}@{}}
\toprule\noalign{}
\begin{minipage}[b]{\linewidth}\raggedright
RUNSPEC
\end{minipage} & \begin{minipage}[b]{\linewidth}\raggedright
GRID
\end{minipage} & \begin{minipage}[b]{\linewidth}\raggedright
EDIT
\end{minipage} & \begin{minipage}[b]{\linewidth}\raggedright
PROPS
\end{minipage} & \begin{minipage}[b]{\linewidth}\raggedright
REGIONS
\end{minipage} & \begin{minipage}[b]{\linewidth}\raggedright
SOLUTION
\end{minipage} & \begin{minipage}[b]{\linewidth}\raggedright
SUMMARY
\end{minipage} & \begin{minipage}[b]{\linewidth}\raggedright
SCHEDULE
\end{minipage} \\
\midrule\noalign{}
\endhead
\bottomrule\noalign{}
\endlastfoot
Alphabetic Listing of Keywords Starting with the Letter M & Status & & &
& & & \\
MAPAXES - Define the Map Origin Input Data & M & & & & & & \\
MAPUNITS -- Define the Map Axes Units & & & & & & & \\
MASSFLOW -- Define River Mass Flow versus Time Tables & & & & & & & \\
MATCORR -- Activate the Material Balance Correction Option & & & & & &
& \\
MAXVALUE -- Sets a Maximum Value for an Array Element & & & & & & & \\
MEMORY -- Define Allocated Memory (Retired) & & & & & & & \\
MESSAGE -- Output User Message & & & & & & & \\
MESSAGES -- Define Message Print Limits and Stop Limits & & & & & & & \\
MESSOPTS -- Reset Severity Level for Forced Time Steps & & & & & & & \\
MESSSRVC - Activate or Deactivate Database Message File Output & & & & &
& & \\
METRIC -- Activate the Metric System of Units for the Model & & & & & &
& \\
MICP -- Activate the Microbially Induced Calcite Precipitation Model &
OPM Flow & & & & & & \\
MICPPARA -- Define Microbially Induced Calcite Precipitation Parameters
(Retired) & OPM Flow & & & & & & \\
MINNNCT -- Set a Minimum Non-Neighbor Connection Transmissibility & & &
& & & & \\
MINNPCOL - Define the Minimum Number of Newton Iterations Used to Update
Well Targets & OPM Flow & & & & & & \\
MINPORV -- Set a Minimum Grid Block Pore Volume Threshold for All Cells
& & & & & & & \\
MINPV -- Set a Minimum Grid Block Pore Volume Threshold for All Cells &
& & & & & & \\
MINPVV -- Set a Minimum Grid Block Pore Volume Threshold for Individual
Cells & & & & & & & \\
MINVALUE -- Set a Minimum Value for an Array Element & & & & & & & \\
MISC -- Define Solvent Miscibility-Immiscibility Transform Functions & &
& & & & & \\
MISCIBLE -- Define Miscibility Todd-Longstaff Parameters & & & & & &
& \\
MISCNUM -- Define the Miscibility Region Numbers & & & & & & & \\
MLANG -- Define Langmuir Maximum Gas Concentration for All Grid Cells &
& & & & & & \\
MLANGSLV -- Define Langmuir Maximum Solvent Concentration for All Grid
Cells & & & & & & & \\
MONITOR -- Activate Output of the Monitoring Data and File & & & & & &
& \\
MPFANUM -- Define Multi-Point Flux Discretization Regions & & & & & &
& \\
MPFNNC -- Define Multi-Point Flux Non-Neighbor Connections & & & & & &
& \\
MSFN -- Miscible Normalized Relative Permeability Tables & & & & & &
& \\
MSGFILE -- Activate or Deactivate Message File Output & & & & & & & \\
MULSGGD -- Multiply Matrix-Fracture Coupling for Oil-Gas Gravity
Drainage for All Cells & & & & & & & \\
MULSGGDV -- Multiply Matrix-Fracture Coupling for Oil-Gas Gravity
Drainage for Individual Cells & & & & & & & \\
MULTFLT -- Multiply the Transmissibility of a Defined Fault by a
Constant & & & & & & & \\
MULTIN -- Activate the Non-Unified Multiple Input File Option & & & & &
& & \\
MULTIPLY -- Multiply a Specified Array by a Constant & & & & & & & \\
MULTIREG -- Multiply an Array by a Constant based on a Region Number & &
& & & & & \\
MULTNUM -- Define the Multiple Transmissibility Regions & & & & & & & \\
MULTOUT -- Activate the Non-Unified Multiple Output File Option & & & &
& & & \\
MULTOUTS -- Activate Non-Unified Multiple Summary Output File Option & &
& & & & & \\
MULTPV -- Multiply Cell Pore Volumes by a Constant & & & & & & & \\
MULTR - Multiply Cell Transmissibility in the +R Direction & & & & & &
& \\
MULTR- - Multiply Cell Transmissibility in the -R Direction & & & & & &
& \\
MULTREAL -- Activate Commercial Simulator's Multi-Realization License &
& & & & & & \\
MULTREGD -- Multiply Diffusivities Between Regions & & & & & & & \\
MULTREGH -- Multiply Thermal Conductivities Between Regions & & & & & &
& \\
MULTREGP -- Multiply Pore Volumes Based On Region Number & & & & & &
& \\
MULTREGT -- Multiply Transmissibilities Between Regions & & & & & & & \\
MULTSIG -- Multiply Matrix-Fracture Coupling for All Cells & & & & & &
& \\
MULTSIGV -- Multiply Matrix-Fracture Coupling for Individual Cells & & &
& & & & \\
MULTTHT - Multiply Cell Transmissibility in the +Theta Direction & & & &
& & & \\
MULTTHT- - Multiply Cell Transmissibility in the -Theta Direction & & &
& & & & \\
MULTX - Multiply Cell Transmissibility in the +X Direction & & & & & &
& \\
MULTX- - Multiply Cell Transmissibility in the -X Direction & & & & & &
& \\
MULTY - Multiply Cell Transmissibility in the +Y Direction & & & & & &
& \\
MULTY- - Multiply Cell Transmissibility in the -Y Direction & & & & & &
& \\
MULTZ - Multiply Cell Transmissibility in the +Z Direction & & & & & &
& \\
MULTZ- - Multiply Cell Transmissibility in the -Z Direction & & & & & &
& \\
\hyperref[0.0.0.MW]{MW -- Molecular Weights} & & & & & & & \\
\end{longtable}

\begin{longtable}[]{@{}
  >{\raggedright\arraybackslash}p{(\linewidth - 14\tabcolsep) * \real{0.1250}}
  >{\raggedright\arraybackslash}p{(\linewidth - 14\tabcolsep) * \real{0.1250}}
  >{\raggedright\arraybackslash}p{(\linewidth - 14\tabcolsep) * \real{0.1250}}
  >{\raggedright\arraybackslash}p{(\linewidth - 14\tabcolsep) * \real{0.1250}}
  >{\raggedright\arraybackslash}p{(\linewidth - 14\tabcolsep) * \real{0.1250}}
  >{\raggedright\arraybackslash}p{(\linewidth - 14\tabcolsep) * \real{0.1250}}
  >{\raggedright\arraybackslash}p{(\linewidth - 14\tabcolsep) * \real{0.1250}}
  >{\raggedright\arraybackslash}p{(\linewidth - 14\tabcolsep) * \real{0.1250}}@{}}
\toprule\noalign{}
\begin{minipage}[b]{\linewidth}\raggedright
RUNSPEC
\end{minipage} & \begin{minipage}[b]{\linewidth}\raggedright
GRID
\end{minipage} & \begin{minipage}[b]{\linewidth}\raggedright
EDIT
\end{minipage} & \begin{minipage}[b]{\linewidth}\raggedright
PROPS
\end{minipage} & \begin{minipage}[b]{\linewidth}\raggedright
REGIONS
\end{minipage} & \begin{minipage}[b]{\linewidth}\raggedright
SOLUTION
\end{minipage} & \begin{minipage}[b]{\linewidth}\raggedright
SUMMARY
\end{minipage} & \begin{minipage}[b]{\linewidth}\raggedright
SCHEDULE
\end{minipage} \\
\midrule\noalign{}
\endhead
\bottomrule\noalign{}
\endlastfoot
Alphabetic Listing of Keywords Starting with the Letter N & Status & & &
& & & \\
NARROW -- Activate Run Summary Narrow Column Output Option & N & & & & &
& \\
NCOMPS -- Confirm Number of Compositional Components & OPM Flow & & & &
& & \\
NCONSUMP - Node Gas Consumption (Extended Network) & & & & & & & \\
NEFAC - Node Efficiency Factors (Extended Network) & & & & & & & \\
NETBALAN -- Network Balancing Parameters & & & & & & & \\
NETCOMPA -- Define Automatic Compressors (Extended Network) & & & & & &
& \\
NETWORK -- Activate the Extended Network Option and Define Parameters &
& & & & & & \\
NEWTON -- Activate Newton Iteration SUMMARY Output & & & & & & & \\
NEWTRAN -- Activate Irregular Corner-Point Grid Transmissibilities & & &
& & & & \\
NEXT -- Maximum Next Time Step Size (Alias for NEXTSTEP) & & & & & &
& \\
NEXTSTEP -- Maximum Next Time Step Size & & & & & & & \\
NEXTSTPL -- Maximum Next Time Step Size (LGR) & & & & & & & \\
NINENUM -- Define the Nine-Point Discretization Region & & & & & & & \\
NINEPOIN -- Activate the Nine-Point Discretization Option & & & & & &
& \\
NMATOPTS -- Define the Discretized Matrix Dual Porosity Parameters & & &
& & & & \\
NMATRIX -- Activate the Discretized Matrix Dual Porosity Option & & & &
& & & \\
NMESSAGE -- Export Cumulative Message Summary Variables to File & & & &
& & & \\
NNC -- Define Non-Neighbor Connections Between Cells Manually & & & & &
& & \\
NNEWTF -- Activate the Non-Newtonian Fluid Model & & & & & & & \\
NOCASC -- Activate Linear Solver Tracer Algorithm & & & & & & & \\
NODEPROP -- Define Network Node Properties for Extended Network & & & &
& & & \\
NODPPM -- Deactivate Fracture Porosity-Permeability Calculation & & & &
& & & \\
NOECHO -- Deactivate Echoing of User Input Files to the Print File & & &
& & & & \\
NOGGF -- Deactivate Output of Grid Geometry File & & & & & & & \\
NOHMD -- Deactivate History Match Gradient Derivative Calculations & & &
& & & & \\
NOHMO -- Deactivate History Match Gradient Derivative Calculations
(Alias) & & & & & & & \\
NOHYST - Deactivate the Hysteresis Option & & & & & & & \\
NOINSPEC -- Deactivate Output of the INIT Index File & & & & & & & \\
NOMONITO -- Deactivate Output of the Monitoring Data and File & & & & &
& & \\
NONNC -- Deactivate Non-Neighbor Connections & & & & & & & \\
NORSSPEC -- Deactivate Output of the RESTART Index File & & & & & & & \\
NOSIM -- Activate the No Simulation Mode for Data File Checking & & & &
& & & \\
NOWARN -- Deactivate Warning Messages & & & & & & & \\
NOWARNEP -- Deactivate End-Point Scaling Warning Messages & & & & & &
& \\
NRSOUT -- Defined Maximum Number of RESTART Elements & & & & & & & \\
NSTACK -- Define the Stack Length for the Iterative Linear Solver & & &
& & & & \\
NTG -- Define the Net-to-Gross Ratio for All the Cells & & & & & & & \\
NUMRES -- Define the Number of Reservoir Grids & & & & & & & \\
NUPCOL -- Define the Maximum Number of Newton Iterations Used to Update
Well Targets & & & & & & & \\
NWATREM -- Node Water Removal (Extended Network) & & & & & & & \\
NXFIN -- Define the Number of LGR Grid Blocks in the X-Direction & & & &
& & & \\
NYFIN -- Define the Number of LGR Grid Blocks in the Y-Direction & & & &
& & & \\
NZFIN -- Define the Number of LGR Grid Blocks in the Z-Direction & & & &
& & & \\
\end{longtable}

\begin{longtable}[]{@{}
  >{\raggedright\arraybackslash}p{(\linewidth - 14\tabcolsep) * \real{0.1250}}
  >{\raggedright\arraybackslash}p{(\linewidth - 14\tabcolsep) * \real{0.1250}}
  >{\raggedright\arraybackslash}p{(\linewidth - 14\tabcolsep) * \real{0.1250}}
  >{\raggedright\arraybackslash}p{(\linewidth - 14\tabcolsep) * \real{0.1250}}
  >{\raggedright\arraybackslash}p{(\linewidth - 14\tabcolsep) * \real{0.1250}}
  >{\raggedright\arraybackslash}p{(\linewidth - 14\tabcolsep) * \real{0.1250}}
  >{\raggedright\arraybackslash}p{(\linewidth - 14\tabcolsep) * \real{0.1250}}
  >{\raggedright\arraybackslash}p{(\linewidth - 14\tabcolsep) * \real{0.1250}}@{}}
\toprule\noalign{}
\begin{minipage}[b]{\linewidth}\raggedright
RUNSPEC
\end{minipage} & \begin{minipage}[b]{\linewidth}\raggedright
GRID
\end{minipage} & \begin{minipage}[b]{\linewidth}\raggedright
EDIT
\end{minipage} & \begin{minipage}[b]{\linewidth}\raggedright
PROPS
\end{minipage} & \begin{minipage}[b]{\linewidth}\raggedright
REGIONS
\end{minipage} & \begin{minipage}[b]{\linewidth}\raggedright
SOLUTION
\end{minipage} & \begin{minipage}[b]{\linewidth}\raggedright
SUMMARY
\end{minipage} & \begin{minipage}[b]{\linewidth}\raggedright
SCHEDULE
\end{minipage} \\
\midrule\noalign{}
\endhead
\bottomrule\noalign{}
\endlastfoot
Alphabetic Listing of Keywords Starting with the Letter O & Status & & &
& & & \\
OFM -- Activate OFM File Output of the SUMMARY Data & O & & & & & & \\
OIL -- Activate the Oil Phase in the Model & & & & & & & \\
OILAPI -- Define the Initial Equilibration Oil API for All Grid Blocks &
& & & & & & \\
OILDENT -- Define Oil Density Temperature Coefficients & OPM Flow & & &
& & & \\
OILJT -- Define Oil Joule-Thomson Coefficient & OPM Flow & & & & & & \\
OILVISCT -- Define Oil Viscosity versus Temperature Functions & OPM Flow
& & & & & & \\
OLDTRAN -- Activate Cartesian Regular Grid Transmissibilities & & & & &
& & \\
OLDTRANR -- Activate Radial Regular Grid Transmissibilities & & & & & &
& \\
OPERATE -- Define Mathematical Operations on Arrays & & & & & & & \\
OPERATER -- Define Mathematical Operations on Arrays by Region & & & & &
& & \\
OPERNUM -- Define Regions for Mathematical Operations on Arrays & & & &
& & & \\
OPTIONS -- Activate Various Program Options & & & & & & & \\
OUTRAD - Define the Outer Radius of a Radial Grid & & & & & & & \\
OUTSOL -- Define Data to be Written to the RESTART File (Retired) & & &
& & & & \\
OVERBURD -- Define Rock Overburden Pressure versus Depth Tables & & & &
& & & \\
\end{longtable}

\begin{longtable}[]{@{}
  >{\raggedright\arraybackslash}p{(\linewidth - 14\tabcolsep) * \real{0.1250}}
  >{\raggedright\arraybackslash}p{(\linewidth - 14\tabcolsep) * \real{0.1250}}
  >{\raggedright\arraybackslash}p{(\linewidth - 14\tabcolsep) * \real{0.1250}}
  >{\raggedright\arraybackslash}p{(\linewidth - 14\tabcolsep) * \real{0.1250}}
  >{\raggedright\arraybackslash}p{(\linewidth - 14\tabcolsep) * \real{0.1250}}
  >{\raggedright\arraybackslash}p{(\linewidth - 14\tabcolsep) * \real{0.1250}}
  >{\raggedright\arraybackslash}p{(\linewidth - 14\tabcolsep) * \real{0.1250}}
  >{\raggedright\arraybackslash}p{(\linewidth - 14\tabcolsep) * \real{0.1250}}@{}}
\toprule\noalign{}
\begin{minipage}[b]{\linewidth}\raggedright
RUNSPEC
\end{minipage} & \begin{minipage}[b]{\linewidth}\raggedright
GRID
\end{minipage} & \begin{minipage}[b]{\linewidth}\raggedright
EDIT
\end{minipage} & \begin{minipage}[b]{\linewidth}\raggedright
PROPS
\end{minipage} & \begin{minipage}[b]{\linewidth}\raggedright
REGIONS
\end{minipage} & \begin{minipage}[b]{\linewidth}\raggedright
SOLUTION
\end{minipage} & \begin{minipage}[b]{\linewidth}\raggedright
SUMMARY
\end{minipage} & \begin{minipage}[b]{\linewidth}\raggedright
SCHEDULE
\end{minipage} \\
\midrule\noalign{}
\endhead
\bottomrule\noalign{}
\endlastfoot
Alphabetic Listing of Keywords Starting with the Letter P & Status & & &
& & & \\
PARALLEL -- Define Parallel Run Configuration & P & & & & & & \\
PARAOPTS -- Define Parallel Run Options & & & & & & & \\
PARTTRAC -- Activate and Define Partitioned Tracer Option & & & & & &
& \\
PATHS -- Define Filename Directory Path Aliases & & & & & & & \\
PBUB -- Define the Initial Equilibration Bubble-Point Pressure for All
Grid Blocks & & & & & & & \\
PBVD -- Equilibration Bubble-Point versus Depth Tables & & & & & & & \\
\hyperref[0.0.0.PCFACT]{PCFACT -- Capillary Pressure Multiplication
Factor as a Function of Porosity Change} & OPM Flow & & & & & & \\
\hyperref[0.0.0.PCRIT]{PCRIT -- Critical Pressures} & & & & & & & \\
PCG -- End-Point Scaling of Grid Cell Maximum Gas Capillary Pressure
(Drainage) & & & & & & & \\
PCG32D -- Gas-Oil Capillary Pressure versus Oil and Water Saturation
Tables & & & & & & & \\
PCW -- End-Point Scaling of Grid Cell Water Capillary Pressure
(Drainage) & & & & & & & \\
PCW32D -- Water-Oil Capillary Pressure versus Oil and Gas Saturation
Tables & & & & & & & \\
PDEW -- Define the Initial Equilibration Dew-Point Pressure for All Grid
Blocks & & & & & & & \\
PDVD -- Define Equilibration Dew-Point versus Depth Tables & & & & & &
& \\
PEBI -- Activate and Defined PEBI Grid Options & & & & & & & \\
PECOEFS -- Define Petro-Elastic Model Coefficients & & & & & & & \\
PEDIMS -- Define Petro-Elastic Model Regions and Table Dimensions & & &
& & & & \\
PEGTAB -- Petro-Elastic Pressure Shear Modulus Table & & & & & & & \\
PEGTAB0 & & & & & & & \\
PEGTAB1 & & & & & & & \\
PEGTAB2 & & & & & & & \\
PEGTAB3 & & & & & & & \\
PEGTAB4 & & & & & & & \\
PEGTAB5 & & & & & & & \\
PEGTAB6 & & & & & & & \\
PEGTAB7 & & & & & & & \\
PEKTAB -- Petro-Elastic Pressure Bulk Modulus Table & & & & & & & \\
PEKTAB0 & & & & & & & \\
PEKTAB1 & & & & & & & \\
PEKTAB2 & & & & & & & \\
PEKTAB3 & & & & & & & \\
PEKTAB4 & & & & & & & \\
PEKTAB5 & & & & & & & \\
PEKTAB6 & & & & & & & \\
PEKTAB7 & & & & & & & \\
PENUM -- Define the Petro-Elastic Region Numbers & & & & & & & \\
PERFORMA -- Export Standard Simulator Performance Summary Variables to
File & & & & & & & \\
PERMAVE -- Define Average Transmissibility Coefficients & & & & & & & \\
PERMFACT -- Permeability Multiplication Factor as a Function of Porosity
Change & OPM Flow & & & & & & \\
PERMJFUN -- Define Leverett J-Function Permeability for All Cells & & &
& & & & \\
PERMR -- Define the Permeability for Each Cell in the R Direction & & &
& & & & \\
PERMTHT -- Define the Permeability for Each Cell in the THETA Direction
& & & & & & & \\
PERMX - Define the Permeability in the X Direction for All the Cells & &
& & & & & \\
PERMY - Define the Permeability in the Y Direction for All the Cells & &
& & & & & \\
PERMZ - Define the Permeability in the Z Direction for All the Cells & &
& & & & & \\
PETGRID -- Load a Generic Simulation Grid File & & & & & & & \\
PETOPTS -- Define Petrel and Generic Simulation File Options & & & & & &
& \\
PICOND -- Define the Generalized Pseudo Pressure Parameters & & & & & &
& \\
PIMTDIMS -- Define Well Productivity Scaling Table Dimensions & & & & &
& & \\
PIMULTAB -- Define Well Productivity Index versus Water Cut Tables & & &
& & & & \\
PINCH -- Define Pinch-Out Layer Options & & & & & & & \\
PINCHNUM -- Define Pinch-Out Regions for the PINCHREG Keyword & & & & &
& & \\
PINCHOUT - Define Pinch-Out Layers Option (Fixed) & & & & & & & \\
PINCHREG - Define Pinch-Out Region Options & & & & & & & \\
PINCHXY -- Define Pinch-Out Areal Options & & & & & & & \\
PINTDIMS -- Define Polymer Molecular Weight Model Table Dimensions & OPM
Flow & & & & & & \\
PLMIXNUM -- Define the Polymer Region Numbers & & & & & & & \\
PLMIXPAR -- Define the Polymer Todd-Longstaff Mixing Parameters & & & &
& & & \\
PLYADS - Define Polymer Rock Adsorption Tables & & & & & & & \\
PLYADSS - Define Polymer Rock Adsorption with Salt Dependence Tables & &
& & & & & \\
PLYATEMP -- Define Polymer Adsorption Table Temperature & & & & & & & \\
PLYCAMAX - Define Polymer-Rock Maximum Adsorption by Cell & & & & & &
& \\
PLYDHFLF - Define Polymer Thermal Degradation Half-Life Tables & & & & &
& & \\
PLYESAL -- Define Polymer Effective Salinity Coefficient & & & & & &
& \\
PLYKRRF -- Define Polymer Rock Permeability Reduction by Cell & & & & &
& & \\
PLYMAX - Define Polymer-Salt Viscosity Mixing Concentrations & & & & & &
& \\
PLYMWINJ -- Polymer Molecular Weight Model Throughput and Velocity Table
& OPM Flow & & & & & & \\
PLYRMDEN - Define Polymer Model In Situ Rock Density & & & & & & & \\
PLYROCK - Define Polymer-Rock Properties & & & & & & & \\
PLYROCKM - Modify Polymer-Rock Properties & & & & & & & \\
PLYSHEAR -- Activate and Define Polymer Shearing Parameters & & & & & &
& \\
PLYSHLOG - Activate and Define the Polymer Shearing Logarithmic
Parameters & & & & & & & \\
PLYTRRF -- Define Polymer Rock Permeability Reduction versus Temperature
& & & & & & & \\
PLYTRRFA -- Define Polymer Rock Permeability Reduction versus
Temperature Option & & & & & & & \\
PLYVISC -- Define Polymer Viscosity Scaling Factors & & & & & & & \\
PLYVISCS -- Define Polymer-Salt Viscosity Scaling Factors & & & & & &
& \\
PLYVISCT -- Define Polymer-Temperature Viscosity Scaling Factors & & & &
& & & \\
PLYVMH - Polymer Molecular Weight Model Polymer Viscosity Constants &
OPM Flow & & & & & & \\
PLYVSCST -- Define Polymer-Salt-Temperature Viscosity Scaling Factors &
& & & & & & \\
PMAX -- Maximum and Minimum Pressure for Total Compressibility Check & &
& & & & & \\
PMISC -- Define Miscibility versus Pressure Tables & & & & & & & \\
POLYMER -- Activate the Polymer Phase in the Model & & & & & & & \\
POLYMW -- Activate the Polymer Molecular Weight Transport Option & OPM
Flow & & & & & & \\
PORO - Define the Porosity Values for All the Cells & & & & & & & \\
PORV - Define the Pore Volumes for All the Cells & & & & & & & \\
PPCWMAX -- Define SWATINIT Calculated Capillary Pressure Constraints & &
& & & & & \\
PRECSALT -- Activate the OPM Flow Salt Precipitation Model & OPM Flow &
& & & & & \\
PRESSURE -- Define the Initial Equilibration Pressures for All Grid
Blocks & & & & & & & \\
PRIORITY -- Activate and Define Well Prioritization Coefficients & & & &
& & & \\
PROPS - Define the Start of the PROPS Section of Keywords & & & & & &
& \\
PRORDER -- Define a Group Production Rules Sequence & & & & & & & \\
PRVD -- Define the Initial Equilibration Pressures versus Depth & & & &
& & & \\
PSTEADY -- Activate Pseudo Steady State Flow Calculation Option & & & &
& & & \\
PVCDO - Oil PVT Properties for Dead Oil (Constant Compressibility) & & &
& & & & \\
PVCO - Oil PVT Properties for Live Oil & & & & & & & \\
PVDG - Gas PVT Properties for Dry Gas & & & & & & & \\
PVDO -- Oil PVT Properties for Dead Oil & & & & & & & \\
PVDS - Solvent PVT Properties for the Solvent Model & & & & & & & \\
PVTG - Gas PVT Properties for Wet Gas with Vaporized Oil & & & & & &
& \\
PVTGW - Gas PVT Properties for Dry Gas with Vaporized Water & OPM Flow &
& & & & & \\
PVTGWO - Gas PVT Properties for Wet Gas with Vaporized Water and Oil &
OPM Flow & & & & & & \\
PVTNUM -- Define the PVT Regions & & & & & & & \\
PVTO - Oil PVT Properties for Live Oil & & & & & & & \\
PVTSOL - Oil PVT Properties for Live Oil versus CO2 Mass Fraction & OPM
Flow & & & & & & \\
PVTW - Define Water Fluid Properties for Various Regions & & & & & &
& \\
PVTWSALT - Define Brine Water Fluid Properties for Various Regions & & &
& & & & \\
PVZG - Gas PVT Properties for Dry Gas (Z-Factor) & & & & & & & \\
PYACTION -- Define Python Based Action Conditions and Command Processing
& OPM Flow & & & & & & \\
PYEND -- End the Definition of a PYINPUT Section & OPM Flow & & & & &
& \\
PYINPUT -- Define the Start of a PYINPUT Section & OPM Flow & & & & &
& \\
\end{longtable}

\begin{longtable}[]{@{}
  >{\raggedright\arraybackslash}p{(\linewidth - 14\tabcolsep) * \real{0.1250}}
  >{\raggedright\arraybackslash}p{(\linewidth - 14\tabcolsep) * \real{0.1250}}
  >{\raggedright\arraybackslash}p{(\linewidth - 14\tabcolsep) * \real{0.1250}}
  >{\raggedright\arraybackslash}p{(\linewidth - 14\tabcolsep) * \real{0.1250}}
  >{\raggedright\arraybackslash}p{(\linewidth - 14\tabcolsep) * \real{0.1250}}
  >{\raggedright\arraybackslash}p{(\linewidth - 14\tabcolsep) * \real{0.1250}}
  >{\raggedright\arraybackslash}p{(\linewidth - 14\tabcolsep) * \real{0.1250}}
  >{\raggedright\arraybackslash}p{(\linewidth - 14\tabcolsep) * \real{0.1250}}@{}}
\toprule\noalign{}
\begin{minipage}[b]{\linewidth}\raggedright
RUNSPEC
\end{minipage} & \begin{minipage}[b]{\linewidth}\raggedright
GRID
\end{minipage} & \begin{minipage}[b]{\linewidth}\raggedright
EDIT
\end{minipage} & \begin{minipage}[b]{\linewidth}\raggedright
PROPS
\end{minipage} & \begin{minipage}[b]{\linewidth}\raggedright
REGIONS
\end{minipage} & \begin{minipage}[b]{\linewidth}\raggedright
SOLUTION
\end{minipage} & \begin{minipage}[b]{\linewidth}\raggedright
SUMMARY
\end{minipage} & \begin{minipage}[b]{\linewidth}\raggedright
SCHEDULE
\end{minipage} \\
\midrule\noalign{}
\endhead
\bottomrule\noalign{}
\endlastfoot
Alphabetic Listing of Keywords Starting with the Letter Q & Status & & &
& & & \\
QDRILL -- Define Sequential Drilling Queue Wells & Q & & & & & & \\
QHRATING -- Define River Mass Flow versus Depth Tables & & & & & & & \\
QMOBIL Activate or Deactivate LGR End-Point Mobility Correction & & & &
& & & \\
\end{longtable}

\begin{longtable}[]{@{}
  >{\raggedright\arraybackslash}p{(\linewidth - 14\tabcolsep) * \real{0.1250}}
  >{\raggedright\arraybackslash}p{(\linewidth - 14\tabcolsep) * \real{0.1250}}
  >{\raggedright\arraybackslash}p{(\linewidth - 14\tabcolsep) * \real{0.1250}}
  >{\raggedright\arraybackslash}p{(\linewidth - 14\tabcolsep) * \real{0.1250}}
  >{\raggedright\arraybackslash}p{(\linewidth - 14\tabcolsep) * \real{0.1250}}
  >{\raggedright\arraybackslash}p{(\linewidth - 14\tabcolsep) * \real{0.1250}}
  >{\raggedright\arraybackslash}p{(\linewidth - 14\tabcolsep) * \real{0.1250}}
  >{\raggedright\arraybackslash}p{(\linewidth - 14\tabcolsep) * \real{0.1250}}@{}}
\toprule\noalign{}
\begin{minipage}[b]{\linewidth}\raggedright
RUNSPEC
\end{minipage} & \begin{minipage}[b]{\linewidth}\raggedright
GRID
\end{minipage} & \begin{minipage}[b]{\linewidth}\raggedright
EDIT
\end{minipage} & \begin{minipage}[b]{\linewidth}\raggedright
PROPS
\end{minipage} & \begin{minipage}[b]{\linewidth}\raggedright
REGIONS
\end{minipage} & \begin{minipage}[b]{\linewidth}\raggedright
SOLUTION
\end{minipage} & \begin{minipage}[b]{\linewidth}\raggedright
SUMMARY
\end{minipage} & \begin{minipage}[b]{\linewidth}\raggedright
SCHEDULE
\end{minipage} \\
\midrule\noalign{}
\endhead
\bottomrule\noalign{}
\endlastfoot
Alphabetic Listing of Keywords Starting with the Letter R & Status & & &
& & & \\
RADFIN -- Define a Radial Local Grid Refinement with One Column & R & &
& & & & \\
RADFIN4 -- Define a Radial Local Grid Refinement with Four Columns & & &
& & & & \\
RADIAL -- Radial Grid Activation Option & & & & & & & \\
RAINFALL -- Constant Flux Aquifer Rainfall Flux by Month & & & & & &
& \\
RBEDCONT -- Define River Grid Block Contact Area versus Depth & & & & &
& & \\
RCMASTS -- Reservoir Coupling Group Minimum Time Step for Flow
Restriction & & & & & & & \\
REACHES -- Define River Reaches Structure & & & & & & & \\
READDATA -- Read Schedule Data Based on Current Time Step & & & & & &
& \\
REFINE -- Start the Definition of a Local Grid Refinement & & & & & &
& \\
REGDIMS -- Define the Maximum Number of Regions for a Region Array & & &
& & & & \\
REGIONS - Define the Start of the REGIONS Section of Keywords & & & & &
& & \\
RESIDNUM -- Define Vertical Equilibrium Residual Flow Region Numbers & &
& & & & & \\
RESTART -- Restart Run From an Existing Restart File & & & & & & & \\
RESVNUM -- Define Reservoir Coordinate Data Set & & & & & & & \\
RIVDEBUG -- Define the Debug Data to be Printed to File (Rivers) & & & &
& & & \\
RIVERSYS - Define River System (Branch Structure and Boundary
Conditions) & & & & & & & \\
RIVRDIMS -- Define the River Dimensions and Associated Data & & & & & &
& \\
RIVRPROP -- Modify River Reaches Properties & & & & & & & \\
RIVRXSEC -- Define River Cross-Section versus Depth Parameters & & & & &
& & \\
RIVSALT -- Define River Upstream Flow Salt Concentrations & & & & & &
& \\
RIVTRACE -- Define River Upstream Flow Tracer Concentrations & & & & & &
& \\
RKTRMDIR - Activate ROCKTAB Keyword Directional Transmissibility
Multipliers & & & & & & & \\
ROCK - Define the Rock Compressibility for Various Regions & & & & & &
& \\
ROCK2D -- Pore Volume Compaction versus Pressure and Sw Tables & & & & &
& & \\
ROCK2DTR -- Transmissibility Compaction versus Pressure and Sw Tables &
& & & & & & \\
ROCKCOMP -- Activate Rock Compaction & & & & & & & \\
ROCKFRAC - Define the Rock Volume to Bulk Volume Fraction for All the
Cells & & & & & & & \\
ROCKNUM -- Define Rock Compaction Table Region Numbers & & & & & & & \\
ROCKOPTS -- Define Rock Compaction and Compressibility Options & & & & &
& & \\
ROCKPAMA -- Define Coal Palmer-Mansorri Rock Model Parameters & & & & &
& & \\
ROCKTAB -- Rock Compaction Tables & & & & & & & \\
ROCKTABH -- Rock Compaction Hysteresis Tables & & & & & & & \\
ROCKTABW -- Rock Compaction Tables (Water Induced) & & & & & & & \\
ROCKTHSG -- Rock Compaction Hysteresis Tables (Dual Porosity) & & & & &
& & \\
ROCKTSIG -- Rock Compaction Tables (Dual Porosity) & & & & & & & \\
ROCKWNOD -- Water Saturation Values for Compaction Pressure-Sw Tables &
& & & & & & \\
RPTCPL -- Activate Couple Simulation Reporting & & & & & & & \\
RPTGRID -- Define GRID Section Reporting & & & & & & & \\
RPTGRIDL -- Define GRID Section Reporting for LGRs & & & & & & & \\
RPTHMD - Define Well History Match Gradient Reporting Options & & & & &
& & \\
RPTHMG - Define Well History Match Gradient Reporting (Groups) & & & & &
& & \\
RPTHMW - Define Well History Match Gradient Reporting (Wells) & & & & &
& & \\
RPTINIT -- Define Output to the INIT File & & & & & & & \\
RPTISOL -- Activate Isolated Reservoir Number Reporting & & & & & & & \\
RPTONLY - Activate the Report Time Steps Only Option for the SUMMARY
File & & & & & & & \\
RPTONLYO - Deactivate the Report Time Steps Only Option for the SUMMARY
File & & & & & & & \\
RPTPROPS -- Define PROPS Section Reporting & & & & & & & \\
RPTREGS -- Define REGIONS Section Reporting & & & & & & & \\
RPTRST -- Define Data to be Written to the RESTART File & & & & & & & \\
RPTRUNSP -- Activate RUNSPEC Reporting & & & & & & & \\
RPTSCHED -- Define SCHEDULE Section Reporting & & & & & & & \\
RPTSMRY - Activate or Deactivate Summary List Report & & & & & & & \\
RPTSOL -- Define SOLUTION Section Reporting & & & & & & & \\
RS -- Define the Initial Equilibration GOR (Rs) for All Grid Blocks & &
& & & & & \\
RSCONST -- Define Constant GOR (Rs) for All Dead Oil PVT Fluids & & & &
& & & \\
RSCONSTT -- Define Constant GOR (Rs) for Each Dead Oil PVT Fluid & & & &
& & & \\
RSGI -- Define Gas-Oil Ratio versus Pressure and Gi Tables & & & & & &
& \\
RSSPEC -- Activate Output of the RESTART Index File & & & & & & & \\
RSVD -- Equilibration Dissolved Gas-Oil Ratio (Rs) versus Depth Tables &
& & & & & & \\
\hyperref[0.0.0.RSW]{RSW -- Define the Initial Equilibrium Solution Gas
in Water Ratio for All Grid Blocks} & OPM Flow & & & & & & \\
RTEMP - Define the Initial Reservoir Temperature for the Model & & & & &
& & \\
RTEMPA - Define the Initial Reservoir Temperature for the Model & & & &
& & & \\
RTEMPVD - Define the Initial Reservoir Temperature versus Depth Tables &
& & & & & & \\
RUNSPEC - Define the Start of the RUNSPEC Section of Keywords & & & & &
& & \\
RUNSUM -- Activate RSM File Output of the SUMMARY Data & & & & & & & \\
RV -- Define the Initial Equilibration CGR (Rv) for All Grid Blocks & &
& & & & & \\
RVCONST -- Define Constant CGR (Rv) for All Dry Gas PVT Fluids & & & & &
& & \\
RVCONSTT -- Define Constant CGR (Rv) for Each Dry Gas PVT Fluid & & & &
& & & \\
RVGI -- Define Condensate-Gas Ratio versus Pressure and Gi Tables & & &
& & & & \\
RVVD -- Equilibration Vaporized Oil-Gas Ratio (Rv) versus Depth Tables &
& & & & & & \\
RVW -- Define the Initial Equilibration Vaporized Water in Gas Ratio for
All Grid Blocks & OPM Flow & & & & & & \\
RVWVD -- Equilibration Vaporized Water-Gas Ratio (Rvw) versus Depth
Tables & OPM Flow & & & & & & \\
RWGSALT -- Water Vaporization versus Pressure and Salt Concentration &
OPM Flow & & & & & & \\
\end{longtable}

\begin{longtable}[]{@{}
  >{\raggedright\arraybackslash}p{(\linewidth - 14\tabcolsep) * \real{0.1250}}
  >{\raggedright\arraybackslash}p{(\linewidth - 14\tabcolsep) * \real{0.1250}}
  >{\raggedright\arraybackslash}p{(\linewidth - 14\tabcolsep) * \real{0.1250}}
  >{\raggedright\arraybackslash}p{(\linewidth - 14\tabcolsep) * \real{0.1250}}
  >{\raggedright\arraybackslash}p{(\linewidth - 14\tabcolsep) * \real{0.1250}}
  >{\raggedright\arraybackslash}p{(\linewidth - 14\tabcolsep) * \real{0.1250}}
  >{\raggedright\arraybackslash}p{(\linewidth - 14\tabcolsep) * \real{0.1250}}
  >{\raggedright\arraybackslash}p{(\linewidth - 14\tabcolsep) * \real{0.1250}}@{}}
\toprule\noalign{}
\begin{minipage}[b]{\linewidth}\raggedright
RUNSPEC
\end{minipage} & \begin{minipage}[b]{\linewidth}\raggedright
GRID
\end{minipage} & \begin{minipage}[b]{\linewidth}\raggedright
EDIT
\end{minipage} & \begin{minipage}[b]{\linewidth}\raggedright
PROPS
\end{minipage} & \begin{minipage}[b]{\linewidth}\raggedright
REGIONS
\end{minipage} & \begin{minipage}[b]{\linewidth}\raggedright
SOLUTION
\end{minipage} & \begin{minipage}[b]{\linewidth}\raggedright
SUMMARY
\end{minipage} & \begin{minipage}[b]{\linewidth}\raggedright
SCHEDULE
\end{minipage} \\
\midrule\noalign{}
\endhead
\bottomrule\noalign{}
\endlastfoot
Alphabetic Listing of Keywords Starting with the Letter S & Status & & &
& & & \\
SALINITY -- Define the Reservoir Salinity for All Cells & OPM Flow & & &
& & & \\
SALT -- Define the Initial Equilibration Salt Concentration for All Grid
Blocks & S & & & & & & \\
SALTNODE -- Salt Concentration Based PVTNUM Array & & & & & & & \\
\hyperref[0.0.0.SALTMF]{SALTMF -- Define the Salt Liquid-Phase Mole
Fraction for All Cells} & OPM Flow & & & & & & \\
SALTP -- Define the Initial Precipitated Salt Volume Fraction for All
Grid Blocks & OPM Flow & & & & & & \\
SALTPVD -- Initial Precipitated Salt Volume Fraction versus Depth Tables
& OPM Flow & & & & & & \\
SALTREST -- Define the Restart Salt Concentration for All Grid Blocks &
& & & & & & \\
SALTSOL -- Define the Salt Solubility Limit by Region & OPM Flow & & & &
& & \\
SALTVD -- Equilibration Salt Concentration versus Depth Tables & & & & &
& & \\
SAMG -- Activate Algebraic Multi-Grid Linear Solver & & & & & & & \\
SATNUM -- Define the Saturation Table Region Numbers & & & & & & & \\
SATOPTS -- Activate Relative Permeability Assignment Options & & & & & &
& \\
SAVE -- Activate Output of a SAVE File for Fast Restarts & & & & & &
& \\
SBIOF - Define The Initial Equilibration Biofilm Volume Fraction For All
Grid Blocks & OPM Flow & & & & & & \\
SCALC -- Define The Initial Equilibration Calcite Volume Fraction For
All Grid Blocks & OPM Flow & & & & & & \\
SCALECRS -- Define End-Point Scaling Option & & & & & & & \\
SCALELIM -- End-Point Scaling versus Depth Maximum Water Saturation & &
& & & & & \\
SCDATAB -- Well Connection PI Multipliers versus Scale Deposit & & & & &
& & \\
SCDETAB -- Well Connection Karst Aquifer Properties for Scale Deposit &
& & & & & & \\
SCDPDIMS -- Define Scale Deposition and Damage Table Dimensions & & & &
& & & \\
SCDPTAB -- Well Connection Scale Deposition Tables & & & & & & & \\
SCDPTRAC -- Allocate Sea Water Tracer for Scale Deposition & & & & & &
& \\
SCHEDULE - Define the Start of the SCHEDULE Section of Keywords & & & &
& & & \\
SCVD -- Define Equilibration Coal Solvent Concentration versus Depth
Tables & & & & & & & \\
SDENSITY -- Define the Miscible or Solvent Surface Gas Density & & & & &
& & \\
SEPARATE -- Activate the Separate RSM File Output Option & & & & & &
& \\
SEPVALS -- Define Separator Oil Formation Volume Factor and GOR & & & &
& & & \\
SFOAM -- Define the Initial Equilibration Foam Concentration for All
Grid Blocks & & & & & & & \\
SGAS -- Define the Initial Equilibration Gas Saturation for All Grid
Blocks & & & & & & & \\
SGCR -- End-Point Scaling Grid Cell Critical Gas Saturations & & & & & &
& \\
SGCRX & & & & & & & \\
SGCRX- & & & & & & & \\
SGCRY & & & & & & & \\
SGCRY- & & & & & & & \\
SGCRZ & & & & & & & \\
SGCRZ- & & & & & & & \\
SGCWMIS -- Miscible Critical Gas versus Water Saturation Functions & & &
& & & & \\
SGF32D -- Gas Saturation Tables versus Oil and Water Saturations & & & &
& & & \\
SGFN -- Gas Saturation Tables (Format Type 2) & & & & & & & \\
SGL -- End-Point Scaling Grid Cell Connate Gas Saturations & & & & & &
& \\
SGLX & & & & & & & \\
SGLX- & & & & & & & \\
SGLY & & & & & & & \\
SGLY- & & & & & & & \\
SGLZ & & & & & & & \\
SGLZ- & & & & & & & \\
SGLPC -- End-Point Scaling Grid Cell Capillary Pressure Connate Gas
Saturations & & & & & & & \\
SGOF -- Gas-Oil Saturation Tables versus Gas (Format Type 1) & & & & & &
& \\
SGOFLET-- Gas-Oil LET Relative Permeability Functions & OPM Flow & & & &
& & \\
\hyperref[8.3.280.SGUux5cux2520End-Pointux5cux2520Scalingux5cux2520Gridux5cux2520Cellux5cux2520Gasux5cux2520Saturation]{SGU
-- End-Point Scaling Grid Cell Gas Saturation} & & & & & & & \\
SGUX & & & & & & & \\
SGUX- & & & & & & & \\
SGUY & & & & & & & \\
SGUY- & & & & & & & \\
SGUZ & & & & & & & \\
SGUZ- & & & & & & & \\
SGWFLET -- Gas-Water LET Relative Permeability Functions & OPM Flow & &
& & & & \\
SGWFN -- Gas-Water Saturation Tables (Format Type 2) & & & & & & & \\
SHRATE - Activate and Define the Polymer Shearing Logarithmic Parameters
& & & & & & & \\
SIGMA -- Dual Porosity Matrix to Fracture Sigma (All Cells) & & & & & &
& \\
SIGMAGD -- Dual Porosity Matrix to Fracture Sigma for Gravity Drainage
(All Cells) & & & & & & & \\
SIGMAGDV -- Dual Porosity Matrix to Fracture Sigma Gravity Drainage
(Individual Cells) & & & & & & & \\
SIGMAV -- Dual Porosity Matrix to Fracture Sigma (Individual Cells) & &
& & & & & \\
SIMULATE - Activate the Simulation Mode & & & & & & & \\
SKIP -- Activate Skipping of All Keywords and Input Data & & & & & &
& \\
SKIP100 -- Activate Skipping of Keywords by Black-Oil Simulator & & & &
& & & \\
SKIP300 -- Activate Skipping of Keywords by Compositional Simulator & &
& & & & & \\
SKIPREST -- Activate Skipping of Restart Schedule Data & & & & & & & \\
SKPRPOLY -- Polymer Molecular Weight Model Polymer Injection Skin Table
& OPM Flow & & & & & & \\
SKPRWAT -- Polymer Molecular Weight Model Water Injection Skin Table &
OPM Flow & & & & & & \\
SKRO -- End-Point Scaling of Grid Cell Kro(Swl) (Surfactant) & & & & & &
& \\
SKRORG -- End-Point Scaling of Grid Cell Kro(Sgcr) (Surfactant) & & & &
& & & \\
SKRORW -- End-Point Scaling of Grid Cell Kro(Swcr) (Surfactant) & & & &
& & & \\
SKRW -- End-Point Scaling of Grid Cell Krw(Sw =1.0) (Surfactant) & & & &
& & & \\
SKRWR -- End-Point Scaling of Grid Cell KRWR(Sowcr) (Surfactant) & & & &
& & & \\
SLAVES -- Define Slave Reservoir Simulation Parameters & & & & & & & \\
SLGOF -- Gas-Oil Saturation Tables versus Gas (Format Type 1) & & & & &
& & \\
SMICR -- Define The Initial Equilibration Microbial Concentration For
All Grid Blocks & OPM Flow & & & & & & \\
SMRYDIMS -- Define Maximum Number of Summary Vectors to be Written & & &
& & & & \\
SMULTX - Multiply Cell Transmissibility in the +X Direction
(Auto-Refinement) & & & & & & & \\
SMULTY - Multiply Cell Transmissibility in the +Y Direction
(Auto-Refinement) & & & & & & & \\
SMULTZ - Multiply Cell Transmissibility in the +Z Direction
(Auto-Refinement) & & & & & & & \\
SOCRS -- End-Point Scaling Grid Cell Miscible Critical Oil Saturation
with Respect to Water & & & & & & & \\
SOF2 -- Oil Saturation Tables with Respect to Gas or Water (Format Type
2) & & & & & & & \\
SOF3 -- Oil Saturation Tables with Respect to Gas and Water (Format Type
2) & & & & & & & \\
SOF32D -- Oil Saturation Tables with Respect to Water and Gas (Three
Phase) & & & & & & & \\
SOGCR -- End-Point Scaling Grid Cell Critical Oil Saturation with
Respect to Gas & & & & & & & \\
SOGCRX & & & & & & & \\
SOGCRX- & & & & & & & \\
SOGCRY & & & & & & & \\
SOGCRY- & & & & & & & \\
SOGCRZ & & & & & & & \\
SOGCRZ- & & & & & & & \\
SOIL -- Define the Initial Equilibration Oil Saturation for All Grid
Blocks & & & & & & & \\
SOLUTION - Define the Start of the SOLUTION Section of Keywords & & & &
& & & \\
SOLVCONC -- Define the Initial Coal Solvent Concentration for All Grid
Blocks & & & & & & & \\
SOLVDIMS -- Define PEBI Grid Nested Factorization Solver Dimensions & &
& & & & & \\
SOLVDIRS -- Define Linear Solver Principal Directions & & & & & & & \\
SOLVENT -- Activate the SOLVENT Phase in the Model & & & & & & & \\
SOLVFRAC -- Define the Initial Gas Solvent Fraction for All Grid Blocks
& & & & & & & \\
SOLVNUM -- Define PEBI Grid Correspondence to Solver Order & & & & & &
& \\
SOMGAS -- STONE1 Model Minimum Oil Saturation versus Gas Saturation & &
& & & & & \\
SOMWAT -- STONE1 Model Minimum Oil Saturation versus Water Saturation &
& & & & & & \\
SORWMIS -- Miscible Residual Oil versus Water Saturation Functions & & &
& & & & \\
\hyperref[0.0.0.SOURCE]{SOURCE -- Define Source Term} & OPM Flow & & & &
& & \\
SOWCR -- End-Point Scaling Grid Cell Critical Oil Saturation with
Respect to Water & & & & & & & \\
SOWCRX & & & & & & & \\
SOWCRX- & & & & & & & \\
SOWCRY & & & & & & & \\
SOWCRY- & & & & & & & \\
SOWCRZ & & & & & & & \\
SOWCRZ- & & & & & & & \\
SOXYG - Define The Initial Equilibration Oxygen Concentration For All
Grid Blocks & OPM Flow & & & & & & \\
SPECGRID- Define the Dimensions of a Corner-Point Grid & & & & & & & \\
SPECHEAT -- Define the Specific Heat of Oil, Water and Gas & & & & & &
& \\
SPECROCK -- Define the Specific Heat of the Reservoir Rock & & & & & &
& \\
SPIDER -- Spider Grid Activation Option & OPM Flow & & & & & & \\
SPOLY -- Define the Initial Equilibration Polymer Concentration for All
Grid Blocks & & & & & & & \\
SPOLYMW -- Define The Initial Equilibration Polymer Molecular Weights
For All Grid Blocks & OPM Flow & & & & & & \\
SSFN -- Solvent and Gas Relative Permeability Tables & & & & & & & \\
SSGCR -- End-Point Scaling Grid Cell Surfactant Critical Gas Saturations
& & & & & & & \\
SSGL -- End-Point Scaling Grid Cell Surfactant Connate Gas Saturations &
& & & & & & \\
SSOGCR -- End-Point Scaling Grid Cell Surfactant Critical Oil Saturation
with Respect to Gas & & & & & & & \\
SSOL -- Define the Initial Equilibration Solvent Saturation for All Grid
Blocks & & & & & & & \\
SSOWCR -- End-Point Scaling Grid Cell Surfactant Critical Oil Saturation
with Respect to Water & & & & & & & \\
SSWCR -- End-Point Scaling Grid Cell Critical Water Saturation & & & & &
& & \\
SSWL -- End-Point Scaling Grid Cell Surfactant Connate Water Saturation
& & & & & & & \\
SSWU -- End-Point Scaling Grid Cell Surfactant Maximum Water Saturation
& & & & & & & \\
START -- Simulation Start Date & & & & & & & \\
STOG Define Capillary Pressure Oil-Gas Surface Tension versus Pressure &
& & & & & & \\
STONE -- Activate Stone's Second Three Phase Oil Relative Permeability
Model (Alias for STONE2) & & & & & & & \\
STONE1 -- Activate Stone's First Three Phase Oil Relative Permeability
Model & & & & & & & \\
STONE1EX -- Define Stone's First Three Phase Oil Relative Permeability
Parameter & & & & & & & \\
STONE2 -- Activate Stone's Second Three Phase Oil Relative Permeability
Model & & & & & & & \\
STOW Define Capillary Pressure Oil-Water Surface Tension versus Pressure
& & & & & & & \\
STWG Define Capillary Pressure Water-Gas Surface Tension versus Pressure
& & & & & & & \\
SUMMARY - Define the Start of the SUMMARY Section of Keywords & & & & &
& & \\
SUMTHIN -- Define SUMMARY Data Reporting Time Steps & & & & & & & \\
SUREA - Define The Initial Equilibration Urea Concentration For All Grid
Blocks & OPM Flow & & & & & & \\
SURF -- Define the Initial Equilibration Polymer Concentration for All
Grid Blocks & & & & & & & \\
SURFACT -- Activate the Surfactant Phase in the Model & & & & & & & \\
SURFACTW -- Activate the Surfactant Phase with Wettability Changes in
the Model & & & & & & & \\
SURFADDW -- Defined Surfactant Adsorbed Concentration versus Wettability
Fraction & & & & & & & \\
SURFADS - Define Surfactant Rock Adsorption Tables & & & & & & & \\
SURFCAPD -- Capillary Number versus Miscibility Tables & & & & & & & \\
SURFESAL -- Define Surfactant Effective Salinity Coefficient & & & & & &
& \\
SURFNUM -- Define the Surfactant Miscible Saturation Table Region
Numbers & & & & & & & \\
SURFROCK - Define Surfactant-Rock Properties & & & & & & & \\
SURFST - Surfactant Water-Oil Surface Tension versus Surfactant
Concentration & & & & & & & \\
SURFSTES - Surfactant Water-Oil Surface Tension versus Surfactant and
Salt Concentrations & & & & & & & \\
SURFVISC -- Surfactant Solution Viscosity versus Concentration & & & & &
& & \\
SURFWNUM -- Define the Saturation Table Region Numbers (High Salinity
and Water Wet) & & & & & & & \\
SWAT -- Define the Initial Equilibration Water Saturation for All Grid
Blocks & & & & & & & \\
SWATINIT -- Define the Initial Water Saturation Array for Capillary
Pressure Scaling & & & & & & & \\
SWCR -- End-Point Scaling Grid Cell Critical Water Saturation & & & & &
& & \\
SWCRX & & & & & & & \\
SWCRX- & & & & & & & \\
SWCRY & & & & & & & \\
SWCRY- & & & & & & & \\
SWCRZ & & & & & & & \\
SWCRZ- & & & & & & & \\
SWF32D -- Water Saturation Tables with Respect to Oil and Gas (Three
Phase) & & & & & & & \\
SWFN -- Water Saturation Tables (Format Type 2) & & & & & & & \\
SWINGFAC -- Define Field Gas Contract Parameters & & & & & & & \\
SWL -- End-Point Scaling Grid Cell Connate Water Saturation & & & & & &
& \\
SWLX & & & & & & & \\
SWLX- & & & & & & & \\
SWLY & & & & & & & \\
SWLY- & & & & & & & \\
SWLZ & & & & & & & \\
SWLZ- & & & & & & & \\
SWLPC -- End-Point Scaling Grid Cell Capillary Pressure Connate Water
Saturations & & & & & & & \\
SWOF -- Water-Oil Saturation Tables (Format Type 1) & & & & & & & \\
SWU -- End-Point Scaling Grid Cell Maximum Water Saturation & & & & & &
& \\
SWUX & & & & & & & \\
SWUX- & & & & & & & \\
SWUY & & & & & & & \\
SWUY- & & & & & & & \\
SWUZ & & & & & & & \\
SWUZ- & & & & & & & \\
\end{longtable}

\begin{longtable}[]{@{}
  >{\raggedright\arraybackslash}p{(\linewidth - 14\tabcolsep) * \real{0.1250}}
  >{\raggedright\arraybackslash}p{(\linewidth - 14\tabcolsep) * \real{0.1250}}
  >{\raggedright\arraybackslash}p{(\linewidth - 14\tabcolsep) * \real{0.1250}}
  >{\raggedright\arraybackslash}p{(\linewidth - 14\tabcolsep) * \real{0.1250}}
  >{\raggedright\arraybackslash}p{(\linewidth - 14\tabcolsep) * \real{0.1250}}
  >{\raggedright\arraybackslash}p{(\linewidth - 14\tabcolsep) * \real{0.1250}}
  >{\raggedright\arraybackslash}p{(\linewidth - 14\tabcolsep) * \real{0.1250}}
  >{\raggedright\arraybackslash}p{(\linewidth - 14\tabcolsep) * \real{0.1250}}@{}}
\toprule\noalign{}
\begin{minipage}[b]{\linewidth}\raggedright
RUNSPEC
\end{minipage} & \begin{minipage}[b]{\linewidth}\raggedright
GRID
\end{minipage} & \begin{minipage}[b]{\linewidth}\raggedright
EDIT
\end{minipage} & \begin{minipage}[b]{\linewidth}\raggedright
PROPS
\end{minipage} & \begin{minipage}[b]{\linewidth}\raggedright
REGIONS
\end{minipage} & \begin{minipage}[b]{\linewidth}\raggedright
SOLUTION
\end{minipage} & \begin{minipage}[b]{\linewidth}\raggedright
SUMMARY
\end{minipage} & \begin{minipage}[b]{\linewidth}\raggedright
SCHEDULE
\end{minipage} \\
\midrule\noalign{}
\endhead
\bottomrule\noalign{}
\endlastfoot
Alphabetic Listing of Keywords Starting with the Letter T & Status & & &
& & & \\
TABDIMS -- Define the Number of Tables and the Table Dimensions & T & &
& & & & \\
TBLK -- Define Tracer Initial Grid Block Concentrations & & & & & & & \\
\hyperref[0.0.0.TCRIT]{TCRIT -- Critical Temperatures} & & & & & & & \\
TEMP -- Activate the Temperature Modeling Option & & & & & & & \\
TEMPI -- Define the Initial Temperature Values for All Cells & & & & & &
& \\
TEMPNODE - Temperature Table for Polymer Solution Viscosity & & & & & &
& \\
TEMPTVD -- Activate Temperature Flux Limited Transport Option & & & & &
& & \\
TEMPVD - Define the Initial Reservoir Temperature versus Depth Tables &
& & & & & & \\
\hyperref[0.0.0.THCO2MIX]{THCO2MIX -- Specify Thermal Mixing Models} & &
& & & & & \\
THCGAS -- Define Gas Phase Thermal Conductivity for All Cells & & & & &
& & \\
THCOIL -- Define Oil Phase Thermal Conductivity for All Cells & & & & &
& & \\
THCONR -- Define Rock and Fluid Thermal Conductivity for All Cells & & &
& & & & \\
THCONSF -- Define Gas Saturation Dependent Thermal Conductivity Scaling
Factor for All Cells & & & & & & & \\
THCROCK -- Define Reservoir Rock Thermal Conductivity for All Cells & &
& & & & & \\
THCWATER -- Define Water Phase Thermal Conductivity for All Cells & & &
& & & & \\
THERMAL-- Activate the Thermal Modeling Option & & & & & & & \\
THPRES - Define Equilibration Region Threshold Pressures & & & & & &
& \\
THPRESFT - Define Fault Threshold Pressures & & & & & & & \\
TIGHTEN -- Tighten and Relax Numerical Controls & & & & & & & \\
TIGHTENP -- Tighten and Relax Numerical Controls Individually & & & & &
& & \\
TIME -- Advance Simulation by Cumulative Reporting Time & & & & & & & \\
TITLE -- Define the Title for the Input Deck & & & & & & & \\
TLMIXPAR -- Define the Miscible Todd-Longstaff Mixing Parameters & & & &
& & & \\
TNUM -- Define Passive Tracer Concentration Regions & & & & & & & \\
TOLCRIT -- Define The Critical Saturation Tolerance & & & & & & & \\
TOPS - Define the Depth at the Center of the Top Face for Each Cell & &
& & & & & \\
TPAMEPS -- Volumetric Strain versus Coal Gas Concentration Tables & & &
& & & & \\
TPAMEPSS - Volumetric Strain versus Coal Solvent Concentration Tables &
& & & & & & \\
TRACER -- Define Passive Tracer Variables & & & & & & & \\
TRACERKM -- Multi-Partitioned Tracer Option K(P) Tables & & & & & & & \\
TRACERKP -- Standard Partitioned Tracer Option K(P) Tables & & & & & &
& \\
TRACERS -- Activate Tracer Options and Set Tracer Array Dimensions & & &
& & & & \\
TRACITVD -- Activate and Define Tracer Implicit Flux Limited Transport
Option & & & & & & & \\
TRACTVD -- Activate Tracer Explicit Flux Limited Transport Option & & &
& & & & \\
TRADS -- Environmental Tracer Adsorption Tables & & & & & & & \\
TRANGL -- Define Non-Neighbor Connections Between Global and LGR Cells
Manually & & & & & & & \\
TRANR - Define the Transmissibility in the +R Direction for All the
Cells & & & & & & & \\
TRANTHT - Define the Transmissibility in the +Theta Direction for All
the Cells & & & & & & & \\
TRANX - Define the Transmissibility in the X Direction for All the Cells
& & & & & & & \\
TRANY - Define the Transmissibility in the Y Direction for All the Cells
& & & & & & & \\
TRANZ - Define the Transmissibility in the Z Direction for All the Cells
& & & & & & & \\
TRDCY -- Environmental Tracer Decay Tables & & & & & & & \\
TRDIF -- Tracer Diffusion Tables & & & & & & & \\
TRDIS -- Tracer Dispersion Table Number Allocation & & & & & & & \\
TRKPF -- Define Partitioned Tracer Regions & & & & & & & \\
TRNHD -- Activate Dispersion Non-Homogeneous Diffusion Option & & & & &
& & \\
TRPLPORO -- Activate the Triple Porosity Model Option & & & & & & & \\
TRROCK -- Environmental Tracer-Rock Property Data & & & & & & & \\
TSTEP -- Advance Simulation by Reporting Time & & & & & & & \\
TUNING - Numerical Tuning Control & & & & & & & \\
TUNINGDP -- Numerical Tuning Control for High Throughput Cases & & & & &
& & \\
TUNINGH -- Numerical Tuning Control for History Match Gradient
Calculations & & & & & & & \\
TUNINGL - Numerical Tuning Control for All LGRs & & & & & & & \\
TUNINGS - Numerical Tuning Control for Individual LGRs & & & & & & & \\
TVDP -- Define the Initial Equilibration Tracer Saturation versus Depth
Functions & & & & & & & \\
TZONE -- End-Point Scaling Transition Zone Options & & & & & & & \\
\end{longtable}

\begin{longtable}[]{@{}
  >{\raggedright\arraybackslash}p{(\linewidth - 14\tabcolsep) * \real{0.1250}}
  >{\raggedright\arraybackslash}p{(\linewidth - 14\tabcolsep) * \real{0.1250}}
  >{\raggedright\arraybackslash}p{(\linewidth - 14\tabcolsep) * \real{0.1250}}
  >{\raggedright\arraybackslash}p{(\linewidth - 14\tabcolsep) * \real{0.1250}}
  >{\raggedright\arraybackslash}p{(\linewidth - 14\tabcolsep) * \real{0.1250}}
  >{\raggedright\arraybackslash}p{(\linewidth - 14\tabcolsep) * \real{0.1250}}
  >{\raggedright\arraybackslash}p{(\linewidth - 14\tabcolsep) * \real{0.1250}}
  >{\raggedright\arraybackslash}p{(\linewidth - 14\tabcolsep) * \real{0.1250}}@{}}
\toprule\noalign{}
\begin{minipage}[b]{\linewidth}\raggedright
RUNSPEC
\end{minipage} & \begin{minipage}[b]{\linewidth}\raggedright
GRID
\end{minipage} & \begin{minipage}[b]{\linewidth}\raggedright
EDIT
\end{minipage} & \begin{minipage}[b]{\linewidth}\raggedright
PROPS
\end{minipage} & \begin{minipage}[b]{\linewidth}\raggedright
REGIONS
\end{minipage} & \begin{minipage}[b]{\linewidth}\raggedright
SOLUTION
\end{minipage} & \begin{minipage}[b]{\linewidth}\raggedright
SUMMARY
\end{minipage} & \begin{minipage}[b]{\linewidth}\raggedright
SCHEDULE
\end{minipage} \\
\midrule\noalign{}
\endhead
\bottomrule\noalign{}
\endlastfoot
Alphabetic Listing of Keywords Starting with the Letter U & Status & & &
& & & \\
UDADIMS -- Define the Dimensions of the User Defined Arguments & U & & &
& & & \\
UDQ - Declare User Define Quantities (``UDQ'') & & & & & & & \\
UDQDIMS -- Define the Dimensions of the User Defined UDQ Feature & & & &
& & & \\
UDQPARAM -- Define Parameters for the User Defined Quantity Feature & &
& & & & & \\
UDT - Declare User Define Tables (``UDT'') & & & & & & & \\
UDTDIMS -- Define the Dimensions of the User Defined Tables & & & & & &
& \\
UNCODHMD -- Activate History Match Gradient Unencoded Output & & & & & &
& \\
UNIFIN -- Activate The Unified Input File Option & & & & & & & \\
UNIFOUT -- Activate The Unified Output File Option & & & & & & & \\
UNIFOUTS -- Activate The Unified Output Summary File Option & & & & & &
& \\
UNIFSAVE -- Activate The Unified Output Save File Option & & & & & &
& \\
USECUPL -- Load a Reservoir Coupling File & & & & & & & \\
USEFLUX -- Activate Flux Boundary Model and Define Flux File & & & & & &
& \\
USENOFLO -- Activate Flux Boundary Model Without a Flux File & & & & & &
& \\
\end{longtable}

\begin{longtable}[]{@{}
  >{\raggedright\arraybackslash}p{(\linewidth - 14\tabcolsep) * \real{0.1250}}
  >{\raggedright\arraybackslash}p{(\linewidth - 14\tabcolsep) * \real{0.1250}}
  >{\raggedright\arraybackslash}p{(\linewidth - 14\tabcolsep) * \real{0.1250}}
  >{\raggedright\arraybackslash}p{(\linewidth - 14\tabcolsep) * \real{0.1250}}
  >{\raggedright\arraybackslash}p{(\linewidth - 14\tabcolsep) * \real{0.1250}}
  >{\raggedright\arraybackslash}p{(\linewidth - 14\tabcolsep) * \real{0.1250}}
  >{\raggedright\arraybackslash}p{(\linewidth - 14\tabcolsep) * \real{0.1250}}
  >{\raggedright\arraybackslash}p{(\linewidth - 14\tabcolsep) * \real{0.1250}}@{}}
\toprule\noalign{}
\begin{minipage}[b]{\linewidth}\raggedright
RUNSPEC
\end{minipage} & \begin{minipage}[b]{\linewidth}\raggedright
GRID
\end{minipage} & \begin{minipage}[b]{\linewidth}\raggedright
EDIT
\end{minipage} & \begin{minipage}[b]{\linewidth}\raggedright
PROPS
\end{minipage} & \begin{minipage}[b]{\linewidth}\raggedright
REGIONS
\end{minipage} & \begin{minipage}[b]{\linewidth}\raggedright
SOLUTION
\end{minipage} & \begin{minipage}[b]{\linewidth}\raggedright
SUMMARY
\end{minipage} & \begin{minipage}[b]{\linewidth}\raggedright
SCHEDULE
\end{minipage} \\
\midrule\noalign{}
\endhead
\bottomrule\noalign{}
\endlastfoot
Alphabetic Listing of Keywords Starting with the Letter V & Status & & &
& & & \\
VAPOIL -- Activate the Vaporize Oil in Wet Gas Phase in the Model & V &
& & & & & \\
VAPPARS -- Oil Vaporization Parameters & & & & & & & \\
VAPWAT -- Activate Vaporize Water in the Dry and Wet Gas Phases & OPM
Flow & & & & & & \\
\hyperref[0.0.0.VCRIT]{VCRIT -- Critical Volumes} & & & & & & & \\
VDFLOW -- Velocity Dependent Flow Coefficient for Grid Block Flow (Grid)
& & & & & & & \\
VDFLOWR -- Velocity Dependent Flow Coefficient for Grid Block Flow
(Region) & & & & & & & \\
VE -- Activate Vertical Equilibrium Model (Global) & & & & & & & \\
VEDEBUG -- Vertical Equilibrium Debug Data Output & & & & & & & \\
VEFIN -- Activate Vertical Equilibrium Model (LGR) & & & & & & & \\
VEFRAC -- Vertical Equilibrium Relative Permeability Fraction (Grid) & &
& & & & & \\
VEFRACP -- Vertical Equilibrium Capillary Pressure Fraction (Grid) & & &
& & & & \\
VEFRACPV -- Vertical Equilibrium Capillary Pressure Fraction (Cell) & &
& & & & & \\
VEFRACV -- Vertical Equilibrium Relative Permeability Fraction (Cell) &
& & & & & & \\
VFPCHK -- Define Production Vertical Flow Performance BHP Check & & & &
& & & \\
VFPIDIMS -- Injection Vertical Flow Performance Table Dimensions & & & &
& & & \\
VFPINJ -- Define Injection Vertical Flow Performance Tables & & & & & &
& \\
VFPPDIMS -- Production Vertical Flow Performance Table Dimensions & & &
& & & & \\
VFPPROD -- Define Production Vertical Flow Performance Tables & & & & &
& & \\
VFPTABL -- Define Production Vertical Flow Performance ALQ Interpolation
& & & & & & & \\
VISAGE - Activate External Reservoir Geo-Mechanics VISAGE Option & & & &
& & & \\
\hyperref[0.0.0.VISCAQA]{VISCAQA -- Specify Ezrokhi Coefficients for
Aqueous Viscosity} & & & & & & & \\
VISCD -- Activate Dual Porosity Viscous Displacement Option & & & & & &
& \\
VISCREF - Define Viscosity-Temperature Reference Conditions & & & & & &
& \\
VISDATES -- Define External Reservoir Geo-Mechanics VISAGE Stress Dates
& & & & & & & \\
VISOPTS -- Define External Reservoir Geo-Mechanics VISAGE Options & & &
& & & & \\
\end{longtable}

\begin{longtable}[]{@{}
  >{\raggedright\arraybackslash}p{(\linewidth - 14\tabcolsep) * \real{0.1250}}
  >{\raggedright\arraybackslash}p{(\linewidth - 14\tabcolsep) * \real{0.1250}}
  >{\raggedright\arraybackslash}p{(\linewidth - 14\tabcolsep) * \real{0.1250}}
  >{\raggedright\arraybackslash}p{(\linewidth - 14\tabcolsep) * \real{0.1250}}
  >{\raggedright\arraybackslash}p{(\linewidth - 14\tabcolsep) * \real{0.1250}}
  >{\raggedright\arraybackslash}p{(\linewidth - 14\tabcolsep) * \real{0.1250}}
  >{\raggedright\arraybackslash}p{(\linewidth - 14\tabcolsep) * \real{0.1250}}
  >{\raggedright\arraybackslash}p{(\linewidth - 14\tabcolsep) * \real{0.1250}}@{}}
\toprule\noalign{}
\begin{minipage}[b]{\linewidth}\raggedright
RUNSPEC
\end{minipage} & \begin{minipage}[b]{\linewidth}\raggedright
GRID
\end{minipage} & \begin{minipage}[b]{\linewidth}\raggedright
EDIT
\end{minipage} & \begin{minipage}[b]{\linewidth}\raggedright
PROPS
\end{minipage} & \begin{minipage}[b]{\linewidth}\raggedright
REGIONS
\end{minipage} & \begin{minipage}[b]{\linewidth}\raggedright
SOLUTION
\end{minipage} & \begin{minipage}[b]{\linewidth}\raggedright
SUMMARY
\end{minipage} & \begin{minipage}[b]{\linewidth}\raggedright
SCHEDULE
\end{minipage} \\
\midrule\noalign{}
\endhead
\bottomrule\noalign{}
\endlastfoot
Alphabetic Listing of Keywords Starting with the Letter W & Status & & &
& & & \\
WAGHYSTR -- Define Water-Alternating-Gas Hysteresis Parameters & W & & &
& & & \\
WAITBAL -- Wait On Network Balance Before Allowing Further Actions & & &
& & & & \\
WALKALIN -- Define Water Injection Alkaline Concentration & & & & & &
& \\
WALQCALC -- Define Well VFP Surface ALQ Phase Density & & & & & & & \\
WAPI -- Define Oil Well Injection API Gravity & & & & & & & \\
WARN -- Activate Warning Messages & & & & & & & \\
WATDENT -- Define Water Density Temperature Coefficients & & & & & &
& \\
WATER -- Activate the Water Phase in the Model & & & & & & & \\
WATJT -- Define Water Joule-Thomson Coefficient & OPM Flow & & & & &
& \\
WATVISCT -- Define Water Viscosity versus Temperature Functions & & & &
& & & \\
WBHGLR -- Define Well Bottom-Hole GLR Constraint & & & & & & & \\
WBOREVOL -- Define Effective Wellbore Storage Volume & & & & & & & \\
WCALCVAL -- Define Gas Well Calorific Value & & & & & & & \\
WCONHIST -- Define Well Historical Production Rates and Pressures & & &
& & & & \\
WCONINJ -- Well Injection Targets and Constraints & & & & & & & \\
WCONINJE -- Well Injection Targets and Constraints & & & & & & & \\
WCONINJH -- Well Historical Observed Injection Rates and Pressures & & &
& & & & \\
WCONINJP -- Define Well Injection Targets and Constraints for Pattern
Flood Wells & & & & & & & \\
WCONPROD -- Define Well Production Targets and Constraints & & & & & &
& \\
WCUTBACK -- Define Well Cutback Limits and Parameters & & & & & & & \\
WCUTBACT -- Define Well Tracer Cutback Limits and Parameters & & & & & &
& \\
WCYCLE -- Define Automatic Well Opening and Closing Cycling Parameters &
& & & & & & \\
WDFAC -- Define Gas Flow Dependent Skin Factor & & & & & & & \\
WDFACCOR -- Gas Flow Dependent Skin Factor (Correlation) & & & & & &
& \\
WDRILPRI -- Add Wells to the Drilling Priority Drilling Queue & & & & &
& & \\
WDRILRES -- Activate Prevention of Multi-Completions in the Same Cell
for Queued Wells & & & & & & & \\
WDRILTIM -- Define Drilling Parameters for Automatic Drilling of New
Wells & & & & & & & \\
WECON -- Well Economic Criteria for Production Wells & & & & & & & \\
WECONINJ -- Well Economic Criteria for Injection Wells & & & & & & & \\
WECONT -- Well Economic Tracer Criteria for Production Wells & & & & & &
& \\
WEFAC -- Define Well Efficiency & & & & & & & \\
WELCNTL -- Modify Well Control and Targets & & & & & & & \\
WELDEBUG -- Define the Well Debug Data to be Printed to File & & & & & &
& \\
WELDRAW -- Define Maximum Draw Down for Production Wells & & & & & &
& \\
WELEVNT -- Define Well WPWEM Summary Mnemonic Output Value & & & & & &
& \\
WELLDIMS -- Define the Wells and Group Dimensions & & & & & & & \\
WELLSTRE -- Define Injection Stream Composition & OPM Flow & & & & &
& \\
WELMOVEL -- Move Global Well Into an LGR & & & & & & & \\
WELOPEN -- Define Well and Well Connections Flowing Status & & & & & &
& \\
WELOPENL -- Define Well and Well Connections Flowing Status (LGR) & & &
& & & & \\
WELPI -- Define Well Productivity and Injectivity Indices & & & & & &
& \\
WELPRI -- Assign Well Priority & & & & & & & \\
WELSEGS -- Define Multi-Segment Wells and Their Segment Structure & & &
& & & & \\
WELSOMIN -- Define Well Connection Minimum Oil Saturation for Opening &
& & & & & & \\
WELSPECL -- Define Well Specifications for Local Grid Refinements & & &
& & & & \\
WELSPECS -- Define Well Specifications & & & & & & & \\
WELTARG -- Modify Well Target and Constraint Values & & & & & & & \\
‍WELTRAJ -- Define Well Trajectory Data & OPM Flow & & & & & & \\
WFOAM - Define Well Foam Injection Concentrations & & & & & & & \\
WFRICSEG -- Convert Friction Well to Multi-Segment Well & & & & & & & \\
WFRICSGL -- Convert Friction Well to Multi-Segment Well (LGR) & & & & &
& & \\
WFRICTN -- Define Well as a Friction Well & & & & & & & \\
WFRICTNL -- Define Well as a Friction Well (LGR) & & & & & & & \\
WGASPROD -- Define Sale Gas Well Production Targets & & & & & & & \\
WGORPEN -- Define Well GOR Penalty Parameters & & & & & & & \\
WGRUPCON -- Define Well Guide Rates for Group Control & & & & & & & \\
WH2NUM -- Define WAG Hysteresis Saturation Table Region Numbers (Two
Phase) & & & & & & & \\
WH3NUM -- Define WAG Hysteresis Saturation Table Region Numbers (Three
Phase) & & & & & & & \\
WHEDREFD -- Define Well Hydraulic Head Reference Depth & & & & & & & \\
WHISTCTL - Define Well Historical Target Phase & & & & & & & \\
WHTEMP -- Define Well Tubing Head Temperature Parameters & & & & & &
& \\
\hyperref[12.3.0.WINJCLN]{WINJCLN -- Clean a Fraction of a Deposited
Filter Cake} & & & & & & & \\
\hyperref[12.3.0.WINJDAM]{WINJDAM -- Define Well Filter Cake Properties}
& & & & & & & \\
\hyperref[12.3.0.WINJFCNC]{WINJFCNC -- Define Injection Well Filtrate
Concentration} & & & & & & & \\
WINJGAS -- Define Gas Injection Properties for a Well & OPM Flow & & & &
& & \\
WINJMULT -- Define Well Pressure Dependent Injectivity Multipliers & & &
& & & & \\
WINJTEMP -- Define Injection Fluid Thermal Properties & & & & & & & \\
WLIFT -- Define Well Re-Tubing, THP and Lift Switching Workover
Operations & & & & & & & \\
WLIFTOPT -- Define Well Gas Lift Optimization Parameters & & & & & &
& \\
WLIMTOL -- Define Well Constraint Tolerance & & & & & & & \\
WLIST -- Define Well Lists (Static) & & & & & & & \\
WLISTARG -- Modify Well List Target and Constraint Values (Static) & & &
& & & & \\
WLISTNAM -- Define Well Lists (WLISTARG) & & & & & & & \\
WMICP -- Define Water Injection Well's Microbial, Oxygen, And Urea
Concentrations & OPM Flow & & & & & & \\
WNETCTRL -- Define Well Control for Network Control Option & & & & & &
& \\
WNETDP -- Define Well THP to Network Pressure Drop & & & & & & & \\
WORKLIM -- Define Well Workover Time & & & & & & & \\
WORKTHP -- Define Well Workover Options for THP Killed Wells & & & & & &
& \\
WPAVE -- Well Block Average Pressure Calculation Parameters for All
Wells & & & & & & & \\
WPAVEDEP -- Define Well Reference Depth for Pressure Calculations & & &
& & & & \\
WPIMULT -- Define Well Connection Multipliers & & & & & & & \\
WPIMULTL -- Define Well Connection Multipliers (LGR) & & & & & & & \\
WPITAB - Assign Well Productivity Index versus Water Cut Tables & & & &
& & & \\
WPLUG -- Define Well Plug Back Length & & & & & & & \\
WPMITAB - Assign Well Polymer Molecular Model Injection Tables & OPM
Flow & & & & & & \\
WPOLYMER - Define Water Injection Well Polymer and Salt Concentrations &
& & & & & & \\
WPOLYRED -- Define Well Polymer-Water Viscosity Reduction Factor & & & &
& & & \\
WPOTCALC -- Well Potential Calculation Options & & & & & & & \\
WREGROUP -- Automatic Re-Assignment of Wells to Groups & & & & & & & \\
WRFT -- Activate Well RFT Reporting to the RFT File & & & & & & & \\
WRFTPLT -- Activate Well RFT and PLT Reporting to the RFT File & & & & &
& & \\
WSALT - Define Water Injection Well Salt Concentrations & & & & & & & \\
WSCCLEAN -- Well Deposited Scale Adjustment & & & & & & & \\
WSCCLENL -- Well Deposited Scale Adjustment (LGR) & & & & & & & \\
WSCTAB -- Assign Well Scale Deposition and Scale Damage Tables & & & & &
& & \\
WSEGAICD -- Define Multi-Segment Well Autonomous ICD Connections & & & &
& & & \\
WSEGDFIN -- Define Multi-Segment Well Drift Flux Slip Model Input Data &
& & & & & & \\
WSEGDFMD -- Define Multi-Segment Well Drift Flux Slip Model & & & & & &
& \\
WSEGDFPA -- Define Multi-Segment Well Drift Flux Slip Model Parameters &
& & & & & & \\
WSEGDIMS -- Define Multi-Segment Well Dimensions & & & & & & & \\
WSEGEXSS -- Define Multi-Segment Well Import-Export Segment Volumes & &
& & & & & \\
WSEGFLIM - Define Multi-Segment Well Artificial Choke Connections & & &
& & & & \\
WSEGFMOD -- Define Multi-Segment Well Model & & & & & & & \\
WSEGINIT - Define Multi-Segment Well Initial Conditions & & & & & & & \\
WSEGITER -- Define Multi-Segment Wells Iteration Parameters & & & & & &
& \\
WSEGLABY - Define Multi-Segment Well Labyrinth ICD Connections & & & & &
& & \\
WSEGLINK - Define Multi-Segment Well Looped Flow Paths & & & & & & & \\
WSEGMULT -- Define Multi-Segment Well Frictional Pressure Loss
Multipliers & & & & & & & \\
WSEGPROP -- Modify Multi-Segment Wells and Their Segment Structure & & &
& & & & \\
WSEGPULL -- Define a Multi-Segment Well Down-Hole Separator Pump & & & &
& & & \\
WSEGSEP -- Define a Multi-Segment Well Down-Hole Separator & & & & & &
& \\
WSEGSICD -- Define Multi-Segment Well Spiral ICD Connections & & & & & &
& \\
WSEGSOLV Define Multi-Segment Well Iterative Linear Solver Parameters &
& & & & & & \\
WSEGTABL -- Assign Multi-Segment Well VLP Tables to Segments & & & & & &
& \\
WSEGVALV -- Define Multi-Segment Well Sub-Critical Valve & & & & & &
& \\
WSF - Water Saturation Tables versus Water Saturation (Gas-Water and
CO2STORE Systems) & OPM Flow & & & & & & \\
‍WSKPTAB - Assign Well Polymer Molecular Model Skin Tables & OPM Flow & &
& & & & \\
WSOLVENT - Define Gas Injection Well Solvent Fraction & & & & & & & \\
WSURFACT - Define Water Injection Well Surfactant Concentration & & & &
& & & \\
WTADD -- Add a Constant to a Well Target or Constraint & & & & & & & \\
WTEMP -- Define An Injection Well's Fluid Temperature & & & & & & & \\
WTEMPQ -- Output Well Names and Well Lists to the Print File & & & & & &
& \\
WTEST -- Well Testing Criteria for Re-Opening Closed Wells & & & & & &
& \\
WTHPMAX -- Define a Well's Maximum Flowing THP for Shut-In & & & & & &
& \\
WTMULT -- Multiple a Well Target or Constraint by a Constant & & & & & &
& \\
WTRACER -- Define An Injection Well's Tracer Concentration & & & & & &
& \\
WVFPDP -- Modify Well BHP Obtained from VFP Tables & & & & & & & \\
WVFPEXP -- Define Well VFP Interpolation Options & & & & & & & \\
WWPAVE -- Well Block Average Pressure Calculation Parameters for
Individual Wells & & & & & & & \\
\end{longtable}

\begin{longtable}[]{@{}
  >{\raggedright\arraybackslash}p{(\linewidth - 14\tabcolsep) * \real{0.1250}}
  >{\raggedright\arraybackslash}p{(\linewidth - 14\tabcolsep) * \real{0.1250}}
  >{\raggedright\arraybackslash}p{(\linewidth - 14\tabcolsep) * \real{0.1250}}
  >{\raggedright\arraybackslash}p{(\linewidth - 14\tabcolsep) * \real{0.1250}}
  >{\raggedright\arraybackslash}p{(\linewidth - 14\tabcolsep) * \real{0.1250}}
  >{\raggedright\arraybackslash}p{(\linewidth - 14\tabcolsep) * \real{0.1250}}
  >{\raggedright\arraybackslash}p{(\linewidth - 14\tabcolsep) * \real{0.1250}}
  >{\raggedright\arraybackslash}p{(\linewidth - 14\tabcolsep) * \real{0.1250}}@{}}
\toprule\noalign{}
\begin{minipage}[b]{\linewidth}\raggedright
RUNSPEC
\end{minipage} & \begin{minipage}[b]{\linewidth}\raggedright
GRID
\end{minipage} & \begin{minipage}[b]{\linewidth}\raggedright
EDIT
\end{minipage} & \begin{minipage}[b]{\linewidth}\raggedright
PROPS
\end{minipage} & \begin{minipage}[b]{\linewidth}\raggedright
REGIONS
\end{minipage} & \begin{minipage}[b]{\linewidth}\raggedright
SOLUTION
\end{minipage} & \begin{minipage}[b]{\linewidth}\raggedright
SUMMARY
\end{minipage} & \begin{minipage}[b]{\linewidth}\raggedright
SCHEDULE
\end{minipage} \\
\midrule\noalign{}
\endhead
\bottomrule\noalign{}
\endlastfoot
Alphabetic Listing of Keywords Starting with the Letter X & Status & & &
& & & \\
There are No Keywords Beginning with the Letter X & X & & & & & & \\
\end{longtable}

\begin{longtable}[]{@{}
  >{\raggedright\arraybackslash}p{(\linewidth - 14\tabcolsep) * \real{0.1250}}
  >{\raggedright\arraybackslash}p{(\linewidth - 14\tabcolsep) * \real{0.1250}}
  >{\raggedright\arraybackslash}p{(\linewidth - 14\tabcolsep) * \real{0.1250}}
  >{\raggedright\arraybackslash}p{(\linewidth - 14\tabcolsep) * \real{0.1250}}
  >{\raggedright\arraybackslash}p{(\linewidth - 14\tabcolsep) * \real{0.1250}}
  >{\raggedright\arraybackslash}p{(\linewidth - 14\tabcolsep) * \real{0.1250}}
  >{\raggedright\arraybackslash}p{(\linewidth - 14\tabcolsep) * \real{0.1250}}
  >{\raggedright\arraybackslash}p{(\linewidth - 14\tabcolsep) * \real{0.1250}}@{}}
\toprule\noalign{}
\begin{minipage}[b]{\linewidth}\raggedright
RUNSPEC
\end{minipage} & \begin{minipage}[b]{\linewidth}\raggedright
GRID
\end{minipage} & \begin{minipage}[b]{\linewidth}\raggedright
EDIT
\end{minipage} & \begin{minipage}[b]{\linewidth}\raggedright
PROPS
\end{minipage} & \begin{minipage}[b]{\linewidth}\raggedright
REGIONS
\end{minipage} & \begin{minipage}[b]{\linewidth}\raggedright
SOLUTION
\end{minipage} & \begin{minipage}[b]{\linewidth}\raggedright
SUMMARY
\end{minipage} & \begin{minipage}[b]{\linewidth}\raggedright
SCHEDULE
\end{minipage} \\
\midrule\noalign{}
\endhead
\bottomrule\noalign{}
\endlastfoot
Alphabetic Listing of Keywords Starting with the Letter Y & Status & & &
& & & \\
There are No Keywords Beginning with the Letter Y & Y & & & & & & \\
\end{longtable}

\begin{longtable}[]{@{}
  >{\raggedright\arraybackslash}p{(\linewidth - 14\tabcolsep) * \real{0.1250}}
  >{\raggedright\arraybackslash}p{(\linewidth - 14\tabcolsep) * \real{0.1250}}
  >{\raggedright\arraybackslash}p{(\linewidth - 14\tabcolsep) * \real{0.1250}}
  >{\raggedright\arraybackslash}p{(\linewidth - 14\tabcolsep) * \real{0.1250}}
  >{\raggedright\arraybackslash}p{(\linewidth - 14\tabcolsep) * \real{0.1250}}
  >{\raggedright\arraybackslash}p{(\linewidth - 14\tabcolsep) * \real{0.1250}}
  >{\raggedright\arraybackslash}p{(\linewidth - 14\tabcolsep) * \real{0.1250}}
  >{\raggedright\arraybackslash}p{(\linewidth - 14\tabcolsep) * \real{0.1250}}@{}}
\toprule\noalign{}
\begin{minipage}[b]{\linewidth}\raggedright
RUNSPEC
\end{minipage} & \begin{minipage}[b]{\linewidth}\raggedright
GRID
\end{minipage} & \begin{minipage}[b]{\linewidth}\raggedright
EDIT
\end{minipage} & \begin{minipage}[b]{\linewidth}\raggedright
PROPS
\end{minipage} & \begin{minipage}[b]{\linewidth}\raggedright
REGIONS
\end{minipage} & \begin{minipage}[b]{\linewidth}\raggedright
SOLUTION
\end{minipage} & \begin{minipage}[b]{\linewidth}\raggedright
SUMMARY
\end{minipage} & \begin{minipage}[b]{\linewidth}\raggedright
SCHEDULE
\end{minipage} \\
\midrule\noalign{}
\endhead
\bottomrule\noalign{}
\endlastfoot
Alphabetic Listing of Keywords Starting with the Letter Z & Status & & &
& & & \\
ZCORN -- Define the Depth of Each Corner-Point of a Grid Block & Z & & &
& & & \\
ZIPP2OFF -- Deactivate Automatic Time Step Control & & & & & & & \\
ZIPPY2 -- Activate Automatic Time Step Control & & & & & & & \\
ZMFVD -- Define Compositional Components versus Depth & OPM Flow & & & &
& & \\
\end{longtable}

\chapter{Release Notes}\label{release-notes}

The 2025-10 release consists of some new features and various
improvements and bug fixes. These include the following highlights:

\begin{verbatim}
- Added support for the biofilm model in two-phase gas-water systems.
- The peak memory allocation has been reduced when recreating the linear solver.
- Improved GPU support for blackoil fluid systems.
- Memory consumption has been reduced when the hysteresis model is not being used.
- Support has been added for link time optimization.
- Optimized update calls for CPR-type preconditioners.
- Memory usage is now reduced when energy, diffusion, or dispersion storage is not used.
- The intensive quantities cache now stores values for the current time step only and not the previous time step. This reduces the peak memory usage of the entire simulator by 10 %, without any performance degradation.
- Added a OPM Flow specific option to the Killough hysteresis model (EHYSTR keyword item 14) to correct the construction of the capillary pressure scanning curve.
- The new saturation function consistency checks have been enabled by default.
- Added support for running OPM Flow single phase (water) simulations from Python code using the Python bindings.
- Added support for running OPM Flow gas-water simulations, for example [H2STORE](#REF_HEADING_KEYWORD_H2STORE) or CO2STORE cases, from Python code using the Python bindings.
- Added support for additional TUNING keyword items: NEWTMX, NEWTMN, TRGMBE, TRGCNV, and XXXCNV.
- Improved support for satellite injection and production groups.
- A user defined argument can now be used for group reservoir volume production rate (RESV) controls when using the GCONPROD keyword.
- Restart file support has been added for the number of Newton iterations required by each cell to satisfy the solution change convergence criteria at the last time step (CONV_NEW).
- Support has been added for running simulations beyond 31-Dec-9999 (up to 292 million years).
- Added support for WECON keyword item 9, the name of the follow-on well to be open if the current well is shut.
- Added support for WSEGAICD keyword item 22, the density function exponent.
- Added UDA support for items 3 and 4 of the WTRACER keyword.
- Support has been added for the use of UDAs for ORAT, WRAT, and GRAT in the GSATPROD keyword.
- Full support has been added for running the Constrained Pressure Residual with Algebraic Multigrid (CPR-AMG) preconditioner entirely on the GPU for serial or single-GPU runs, using the GpuIstl framework.
- Support has been added for the Hypre preconditioner with GPU input types.
- The new GPU aware ISTL solver now has MPI support.
- An experimental Hybrid Newton method has been added that enhances nonlinear preconditioning using machine learning.
- The NLDD (Nonlinear Domain Decomposition) solver has reduced the time spent on local solves by skipping subdomains that haven't changed significantly between iterations.
- Partitioner zoltanwell (default) no longer ignores the command line argument --allow-distributed-wells. This makes its behaviour more like that of the other partitioners.
- Added support for water mass in place and water mass injection/production rate/total summary vectors.
- Added support for gas sales summary vectors.
- The well's preferred phase is now used to calculate the well productivity index.
- Improved robustness of well solves and well group behaviour.
\end{verbatim}

The major command line changes made for this release are summarized in
Table A.1

\begin{longtable}[]{@{}
  >{\raggedright\arraybackslash}p{(\linewidth - 6\tabcolsep) * \real{0.2500}}
  >{\raggedright\arraybackslash}p{(\linewidth - 6\tabcolsep) * \real{0.2500}}
  >{\raggedright\arraybackslash}p{(\linewidth - 6\tabcolsep) * \real{0.2500}}
  >{\raggedright\arraybackslash}p{(\linewidth - 6\tabcolsep) * \real{0.2500}}@{}}
\toprule\noalign{}
\begin{minipage}[b]{\linewidth}\raggedright
OPM Flow 2025-10 New and Deprecated Command Line Options
\end{minipage} & \begin{minipage}[b]{\linewidth}\raggedright
\end{minipage} & \begin{minipage}[b]{\linewidth}\raggedright
\end{minipage} & \begin{minipage}[b]{\linewidth}\raggedright
\end{minipage} \\
\midrule\noalign{}
\endhead
\bottomrule\noalign{}
\endlastfoot
No. & Variable Name & Description & Default \\
1 & --check-satfunc-consistency & Saturation function consistency
checking is now activated by default. A boolean value. Whether or not to
check saturation function consistency requirements. & true \\
2 & --cpr-weights-thread-parallel & A boolean value. Enable OpenMP
thread parallelization of CPR weight calculation. This can improve
performance for large models but is disabled by default. & false \\
3 & --debug-verbosity-level & An integer value. Set debug verbosity
level globally. Default is 1, increasing values give additional output,
and 0 disables most messages to the .DBG file. & 1 \\
4 & --edge-conformal & A boolean value. Edge conformal cornerpoint
processing. & false \\
5 & --edge-weights-method & The edge-weighting method is now chosen by
specifying a string rather than an integer. A character string. Choose
edge-weighing strategy: `uniform', `transmissibility', or `logtrans'
(logarithm of transmissibility). & ``transmissibility'' \\
6 & --gpu-aware-mpi & A boolean value. MPI communication use GPU aware
MPI in the sense that it will use GPU direct communication. Setting this
to true will require that the MPI implementation supports GPU direct
communication. If you are unsure, set this to false. & false \\
7 & --hy-ne-config-file & A character string. Use config files for
Hybrid Newton. & ``hybridNewtonConfig.json'' \\
8 & --linear-solver-accelerator & A character string. Choose the backend
for the linear solver, either cpu or gpu. & ``cpu'' \\
9 & --max-well-status-switch-for-wells & An integer value. Maximum
number of status switching (stop\textless-\textgreater open) for a well
during a time-step. & 99 \\
10 & --max-well-status-switch-in-inner-iter-wells & An integer value.
Maximum number of status switching (stop\textless-\textgreater open) for
a well during inner iterations. & 99 \\
11 & --milu-variant & The names of the possible variants have been
converted to lower case. A character string. Specify which variant of
the modified-ILU preconditioner ought to be used. Possible variants are:
ilu (default, plain ILU), milu\_1 (lump diagonal with dropped row
entries), milu\_2 (lump diagonal with the sum of the absolute values of
the dropped row entries), milu\_3 (if diagonal is positive add sum of
dropped row entries, otherwise subtract them), milu\_4 (if diagonal is
positive add sum of dropped row entries, otherwise do nothing). &
``ilu'' \\
12 & --nldd-relative-mobility-change-tol & A positive real value.
Threshold for single cell relative mobility change in the NLDD solver. &
0.1 \\
13 & --partition-method & The names of the partitioning methods have
been converted to lower case. A character string. Choose partitioning
method: `simple', `zoltan', `metis', or `zoltanwell' (Zoltan with all
cells perforated by a well represented by a single vertex). &
``zoltanwell'' \\
14 & --time-step-control-max-reduction-time-step & A positive real value
(only applicable for the general 3rd order controller, using
`control-error-filtering' as time-step-control-tolerance-test-version).
If the (proposed) relative change in time step size is larger than this
parameter, the time step will be rejected. & 0.1 \\
15 & --time-step-control-parameters & A character string (only
applicable for the general 3rd order controller). Parameters for the
general 3rd order controller. Should be given as
`beta\_1;beta\_2;beta\_3;alpha\_2;alpha\_3'. &
``0.125;0.25;0.125;0.75;0.25'' \\
16 & --time-step-control-reject-completed-step & A boolean value (only
applicable for the general 3rd order controller). Include rejection of
completed time steps if the relative change is larger than the time step
control tolerance. & false \\
17 & --time-step-control-tolerance-test-version & A character string
(only applicable for the general 3rd order controller). Ways to decide
if the time step should be rejected. Options: `standard' and
`control-error-filtering'. The standard version compares relative change
to tolerance directly to decide if the time step should be rejected,
while the control-error-filtering version compares the relative change
in time step size to the max-reduction-time-step parameter. &
``standard'' \\
18 & --use-hy-ne & A boolean value. Wheter or not to use Hybrid Newton.
& false \\
19 & --verify-gpu-aware-mpi & A boolean value. Verify that the MPI
implementation supports GPU aware MPI. If this is set to true \emph{and}
--gpu-aware-mpi=true, the simulation will fail if the MPI implementation
does not support GPU aware MPI. Note that the verification is not
exhaustive, and some configurations might not be verified, but will work
in practice. & false \\
20 & --well-group-constraints-max-iterations & An integer value. Maximum
number of iterations in the well/group switching algorithm. & 1 \\
Notes: & & & \\
\end{longtable}

\emph{Table A.1: OPM Flow 2025-10 New and Deprecated Command Line
Options}

In addition to the changes to the command line options the following new
features have been added to the simulator:

\begin{verbatim}
- Added support for the biofilm model in two-phase gas-water systems, which can be activated using the [BIOFILM](#REF_HEADING_KEYWORD_BIOFILM) keyword in the RUNSPEC section in CO2STORE or [H2STORE](#REF_HEADING_KEYWORD_H2STORE) runs ([#6462](https://github.com/OPM/opm-simulators/pull/6462)). The model includes biofilm as a solid phase and as microbes suspended in the water. This can be used in applications such as underground storage where dissolved hydrogen consumed by microbes leads to injectivity and hydrogen loss.
- Added an optimized version of the preconditioner update calls, at the cost of apply times, to tailor the implementation more towards CPR-type preconditioners with smaller apply/update ratios ([#6310](https://github.com/OPM/opm-simulators/pull/6310)).
- Added support for edge conformal grid processing, activated using the --edge-conformal command line option ([#5806](https://github.com/OPM/opm-simulators/pull/5806), [#814](https://github.com/OPM/opm-grid/pull/814)). Edge conformal grids are very useful in the context of geo-mechanical fracturing processes.
- The new saturation function consistency checks have been enabled by default ([#6252](https://github.com/OPM/opm-simulators/pull/6252)).
- Added support for running OPM Flow single phase (water) simulations from Python code using the Python bindings ([#6482](https://github.com/OPM/opm-simulators/pull/6482)). This has many potential applications including coupling to external geochemistry and geomechanics solvers.
- Added additional restart file items to maintain compatibility with user defined arguments ([#4439](https://github.com/OPM/opm-common/pull/4439)).
- Added support for specifying the maximum and minimum number of Newton iterations for a time step (NEWTMX and NEWTMN) on the TUNING keyword in the SCHEDULE section ([#6378](https://github.com/OPM/opm-simulators/pull/6378)).
- Added parser support for the GSATINJE keyword to specify satellite injection at the group level ([#4700](https://github.com/OPM/opm-common/pull/4700)).
- A check has been added to ensure that all wells and/or groups needed in a UDQ DEFINE statement are present at the point of definition ([#4619](https://github.com/OPM/opm-common/pull/4619)). An error message will be generated for any missing well objects, for example:
\end{verbatim}

\begin{longtable}[]{@{}
  >{\raggedright\arraybackslash}p{(\linewidth - 0\tabcolsep) * \real{1.0000}}@{}}
\toprule\noalign{}
\begin{minipage}[b]{\linewidth}\raggedright
Error: Unrecoverable errors while loading input: Problem with UDQ In
CASE.DATA line 251 DEFINE cannot be evaluated: Missing object in
defining expression for FUBHPGI -\textgreater{} No existing well matches
the name XD-1H
\end{minipage} \\
\midrule\noalign{}
\endhead
\bottomrule\noalign{}
\endlastfoot
\end{longtable}

\begin{verbatim}
- Restart file support has been added for the number of Newton iterations required by each cell to satisfy the solution change convergence criteria at the last time step (CONV_NEW) ([#6374](https://github.com/OPM/opm-simulators/pull/6374)). This output can be requested by including the CONV mnemonic on the RPTRST keyword. Note that simulations using NLDD are not currently supported.
- Added support for WSEGAICD keyword item 22, the density function exponent ([#6473](https://github.com/OPM/opm-simulators/pull/6473), [#4736](https://github.com/OPM/opm-common/pull/4736)).
- Experimental support has been added for the WELSPECL and COMPDATL keywords ([#6297](https://github.com/OPM/opm-simulators/pull/6297)).
- Added a command line option (--linear-solver-accelerator=cpu or gpu) to optionally carry out the linear solver operations on the  GPU ([#6270](https://github.com/OPM/opm-simulators/pull/6270)). By default all operations are carried out on the CPU. When using the gpu option, the linear solvers can be specified with the normal --linear-solver argument, which can point to a JSON file which can have the same format as the one for the CPU. Note that some options and preconditioners are currently limited to the CPU.
- Full support has been added for running the Constrained Pressure Residual with Algebraic Multigrid (CPR-AMG) preconditioner entirely on the GPU for serial or single-GPU runs, using the GpuIstl framework ([#6308](https://github.com/OPM/opm-simulators/pull/6308), [Lye et al. (2025)](https://doi.org/10.21105/joss.07740)). It supports both AMGX and Hypre BoomerAMG solvers, allows consistent configuration across CPU and GPU, and introduces fast GPU-native implementations for CPR operations and quasiimpes weights. The setup minimizes memory usage and transfer overhead by sharing a single GPU matrix across solver components. Initial benchmarks show comparable or improved convergence versus CPU runs.
- It is now possible to run the identical CPR-AMG preconditioner setup on both CPU and GPU (with different hardware specific AMG solvers) by using the command line option linear-solver-accelerator=cpu and gpu respectively.
- An experimental Hybrid Newton method has been added that enhances nonlinear preconditioning using machine learning. Machine learning-based corrections have been integrated into the standard Newton solver of the black-oil simulator ([#6424](https://github.com/OPM/opm-simulators/pull/6424), [Lechevallier, et al. (2025)](https://doi.org/10.2118/223843-MS)). One or more trained Hybrid Newton models can now be used to adjust the initial guess of the nonlinear solver at specific timesteps. Two new command line options have been added: setting --use-hy-ne=true activates the Hybrid Newton method, and --hy-ne-config-file=<path> specifies the path of the JSON configuration file that allows OPM Flow to correctly interpret how to precondition Newton's method.
- The NLDD (Nonlinear Domain Decomposition) solver has reduced the time spent on local solves by skipping subdomains that haven't changed significantly between iterations ([#6346](https://github.com/OPM/opm-simulators/pull/6346)). Mobility changes indicate significant alterations in fluid flow characteristics, making it an effective indicator for determining when subdomain re-solving is necessary. The threshold for relative mobility change detection (default: 0.1) is specified using the --nldd-relative-mobility-change-tol command line option. By setting this parameter to 0, the mobility thresholding will be disabled.
- Partitioner zoltanwell (default) no longer ignores the command line argument --allow-distributed-wells ([#874](https://github.com/OPM/opm-grid/pull/874)). This makes its behaviour more like that of the other partitioners.
- Added option to limit the relative change in the solution variables ([#6009](https://github.com/OPM/opm-simulators/pull/6009)). If the relative change in the solution variables from one time step to the next becomes too large, it might be best to discard the computations even though the time step completed without numerical errors. The time step size would then be adjusted down before retrying the same time step. This option is only available when using the general3rdorder option for time-step-control. (See the new --time-step-control-reject-completed-step command line argument).
- Added the option to use control error filtering as the tolerance test, instead of comparing relative change directly, to the general third order time step controller ([#6280](https://github.com/OPM/opm-simulators/pull/6280)).
- Added support for water mass in place summary vectors at the field, region and block level [FRB]AMIP, and water mass injection/production rates/totals at the field, group, well, connection, and completion levels [FGWC]AM[IP][RT][L] ([#6181](https://github.com/OPM/opm-simulators/pull/6181), [#4569](https://github.com/OPM/opm-common/pull/4569)). Note that AM has been used to indicate water (aqueous) mass since WM, for example GWMIR, is already used to indicate water molar mass.
- Satellite injection rates have been added to group level summary vectors ([#4728](https://github.com/OPM/opm-common/pull/4728)).
- A new application networkgraph has been created ([#4535](https://github.com/OPM/opm-common/pull/4535)). This can be used to visualize a production network (extended network model) including wells.
- Added the option to perform additional iterations in the well/group switching algorithm ([#5933](https://github.com/OPM/opm-simulators/pull/5933)). See the command line argument --well-group-constraints-max-iterations, which has a default value of 1.
\end{verbatim}

Improvements to the simulator in this release include the following:

\begin{verbatim}
- Both the linear solver command line options (e.g. linear-solver=ilu) and the parameters in JSON configuration files, that select the solvers and preconditioners, have been made case-insensitive by internally transforming all input into lowercase ([#6388](https://github.com/OPM/opm-simulators/pull/6388)).
- The following command line options have changed from taking integer arguments to taking string arguments ([#6264](https://github.com/OPM/opm-simulators/pull/6264)). The defaults and behaviour have not been changed.
- --edge-weights-method allowed values are: uniform, transmissibility, or logtrans. Replacing 0, 1, and 2.
- --partition-method allowed values are: simple, zoltan, metis, or zoltanwell. Replacing 0-3.
- The verbosity of the output from the time step controllers can be specified using the --time-step-verbosity=BOOLEAN command line argument ([#6320](https://github.com/OPM/opm-simulators/pull/6320)). Setting this to true (default) increases the verbosity, or to false decreases the verbosity.
- The code has been simplified eliminating ambiguous calls to broadcast() ([#6180](https://github.com/OPM/opm-simulators/pull/6180)).
- The compatibility of the interface with Dune-Grid has been improved ([#892](https://github.com/OPM/opm-grid/pull/892)).
- Reduced peak memory consumption when the command line option --cpr-reuse-setup is set to recreate the linear solver (which it does by default), see that option and --cpr-reuse-interval for details ([#6345](https://github.com/OPM/opm-simulators/pull/6345)).
- Improved GPU support for blackoil fluid systems ([#6365](https://github.com/OPM/opm-simulators/pull/6365)). The BlackoilIntensiveQuantities::update() function can now be called without an ElementContext. The motivation is to make it a better fit for the GPU code, but it should also allow some more optimizations on the CPU, such as static dispatch for the "multiplexed" (runtime switched) functions of the fluid and material subsystems.
- Memory consumption has been reduced when the hysteresis model is not being used ([#6510](https://github.com/OPM/opm-simulators/pull/6510)). A new variant of the simulator (flow_blackoil_nohyst) has been added. This variant is automatically chosen whenever there is no hysteresis in the case and the black-oil variant (flow_blackoil) would otherwise be chosen.
- Support has been added for link time optimization ([#4666](https://github.com/OPM/opm-common/pull/4666)).
- The code now avoids using an ElementContext just to evaluate IntensiveQuantities and demonstrates the use of compile-time dispatch to select a particular three-phase model avoiding run-time switching per cell ([#6367](https://github.com/OPM/opm-simulators/pull/6367), [#4658](https://github.com/OPM/opm-common/pull/4658)). This gives a 5-10% speedup for the property calculations.
- Improved the performance of the Extended Summary file (ESMRY) output routine and fixed the RSTEP vector output ([#4677](https://github.com/OPM/opm-common/pull/4677)).
- Detailed instrumentation using Tracy has been enabled for individual subsystems of the simulator ([#4751](https://github.com/OPM/opm-common/pull/4751)). This is used by developers to profile the application and improve performance, and not an end user feature.
- Improved the efficiency of the formation volume factor and viscosity evaluation ([#4674](https://github.com/OPM/opm-common/pull/4674), [#6386](https://github.com/OPM/opm-simulators/pull/6386)).  The "Props/update time" for Norne was reduced from 3.03s to 2.16s (subject to some variation).
- Memory usage is now reduced when energy, diffusion, or dispersion storage is not used ([#6351](https://github.com/OPM/opm-simulators/pull/6351)).
- The efficiency of the CPR weight calculation has been improved ([#6457](https://github.com/OPM/opm-simulators/pull/6457)).
- The intensive quantities cache now stores values for the current time step only and not the previous time step ([#6357](https://github.com/OPM/opm-simulators/pull/6357)). This reduces the peak memory usage of the entire simulator by 10 %, without any performance degradation.
- The memory consumption during simulation has been reduced by 96 MB per million cells by removing unused capillary pressure storage ([#4650](https://github.com/OPM/opm-common/pull/4650)).
- There is a speed improvement in some cases where a single PYACTION block generates multiple keywords ([#4624](https://github.com/OPM/opm-common/pull/4624)).
- GPU matrix updates have been optimized by using direct memory transfer and pinned memory for Gpu-ISTL, providing significant performance improvements for GPU linear solver operations ([#6406](https://github.com/OPM/opm-simulators/pull/6406)).
- Improved error messages when invalid grid block ranges are defined using the BOX keyword ([#4600](https://github.com/OPM/opm-common/pull/4600)).
- Messages reporting that a well has converged in zero iterations (for example, 'Well P-1 converged in 0 inner iterations') will now only be output to the Debug (DBG) file if an elevated verbosity level has been specified by the user ([#6350](https://github.com/OPM/opm-simulators/pull/6350)).
- Added support for running OPM Flow gas-water simulations, for example [H2STORE](#REF_HEADING_KEYWORD_H2STORE) or CO2STORE cases, from Python code using the Python bindings ([#6075](https://github.com/OPM/opm-simulators/pull/6075)).
- The Python documentation for OPM Flow has been updated ([#4615](https://github.com/OPM/opm-common/pull/4615), [OPM Python Documentation](https://opm.github.io/opm-python-documentation/master/index.html)).
- Added effect of satellite surface injection rates, and updated satellite surface production rates ([#6454](https://github.com/OPM/opm-simulators/pull/6454)).
- Added support for additional TUNING keyword items: target non-linear convergence error (TRGCNV), maximum non-linear convergence error (XXXCNV), and target material balance error (TRGMBE) ([#6413](https://github.com/OPM/opm-simulators/pull/6413)).
- A user defined argument can now be used for group reservoir volume production rate (RESV) controls when using the GCONPROD keyword ([#4503](https://github.com/OPM/opm-common/pull/4503)). In addition, restart support has been added for this control.
- Support has been added for running simulations beyond 31-Dec-9999 (up to 292 million years) ([#6336](https://github.com/OPM/opm-simulators/pull/6336)).
- Support has been added for underscore characters in user defined region set names ([#4611](https://github.com/OPM/opm-common/pull/4611)). These were previously rejected at the input level.
- The implementation of the exponentiation operator "^" has been extended to support combinations of UDQ sets and UDQ scalars in the definition of UDQs (see the UDQ keyword) ([#4757](https://github.com/OPM/opm-common/pull/4757)). Previously, this operator only supported UDQ sets of the same size.
- Added support for WECON keyword item 9, the name of the follow-on well to be open if the current well is shut ([#6050](https://github.com/OPM/opm-simulators/pull/6050), [#4507](https://github.com/OPM/opm-common/pull/4507)).
- Added UDA support for items 3 and 4 of the WTRACER keyword ([#6549](https://github.com/OPM/opm-simulators/pull/6549), [#4800](https://github.com/OPM/opm-common/pull/4800)).
- If the GSATPROD keyword specifies a group that doesn't exist at this point in the schedule then it will be created ([#4689](https://github.com/OPM/opm-common/pull/4689)). Groups created in this way will be parented directly to the FIELD group.
- Support has been added for the use of UDAs for ORAT, WRAT, and GRAT in the GSATPROD keyword ([#6429](https://github.com/OPM/opm-simulators/pull/6429), [#4701](https://github.com/OPM/opm-common/pull/4701)).
- Added the option to output FIP and RESV data to a separate CSV file ([#6329](https://github.com/OPM/opm-simulators/pull/6329)). This is requested by including the CSVFIP mnemonic in the RPTSOL keyword.
- A more user friendly message is now output if an attempt is made to restart a simulation from report step zero, which is not supported by OPM Flow ([#4638](https://github.com/OPM/opm-common/pull/4638)).
- Added support for the parallel Hypre preconditioner for the solution of scalar systems ([#6384](https://github.com/OPM/opm-simulators/pull/6384)). This enables the use of Hypre for CPR, CPRW and possibly other systems written as scalar compressed row storage (CRS) matrices, for example, geomechanics or separate temperature equations.
- Support has been added for the Hypre preconditioner with GPU input types by adding a unified, backend-agnostic HypreInterface that enables seamless integration with both CPU and GPU data types ([#6444](https://github.com/OPM/opm-simulators/pull/6444)). Using the functionality from [#6384](https://github.com/OPM/opm-simulators/pull/6384), this also makes it possible to use Hypre MPI parallel on the GPU. This is an important functionality for the broader CPR-AMG on GPU implementation in [#6308](https://github.com/OPM/opm-simulators/pull/6308).
- A new AmgxInterface class has been added that provides a unified, backend-agnostic interface to the AMGX functionality, enabling seamless usage with both CPU and GPU data types ([#6423](https://github.com/OPM/opm-simulators/pull/6423)). This provides important functionality for the broader CPR-AMG on GPU implementation in [#6308](https://github.com/OPM/opm-simulators/pull/6308) that is essential for achieving good GPU performance, as it ensures minimal data transfers.
- A new sparse matrix implementation has been added that uses cuSPARSE's modern Generic API instead of the legacy API ([#6401](https://github.com/OPM/opm-simulators/pull/6401)). This addition is essential for supporting block size 1 operations in the GPU CPR-AMG implementation in [#6308](https://github.com/OPM/opm-simulators/pull/6308).
- The CPR weight calculation has been parallelized using OpenMP threading to improve performance, particularly important for GPU configurations where the CPU is idle and the CPU-based weight calculation can become a significant proportion of linear solve time ([#6419](https://github.com/OPM/opm-simulators/pull/6419)).
- The new GPU aware ISTL solver now has MPI support ([#6379](https://github.com/OPM/opm-simulators/pull/6379)). This setup only makes sense if you have multiple GPUs either on one node or spread across multiple nodes in a cluster. To use this simply run flow through MPI:
- mpirun -np 2 flow --linear-solver-accelerator=gpu --matrix-add-well-contributions=true CASENAME
- We additionally support using GPUDirect/GPU aware MPI/CUDA aware MPI should the cluster and MPI libraries used for compilation support this. To enable, use the --gpu-aware-mpi=BOOLEAN option. Using GPU aware MPI is expected to give a speedup.
- Furthermore, currently limited to OpenMPI systems, we can verify if GPU aware MPI support is available. To enable or disable this check, use the --verify-gpu-aware-mpi=BOOLEAN option. Note that this check is not always conclusive, and for most systems it is recommended to turn it off and rather test correctness through other means.
- The Gpu-ISTL code structure has been reorganised to reduce memory consumption and facilitate fewer host-to-device copies in the future ([#6241](https://github.com/OPM/opm-simulators/pull/6241)).
- Improved robustness by checking that the specified external partition file exists before trying to read it ([#6326](https://github.com/OPM/opm-simulators/pull/6326)). If the external partition file does not exist, an error message will be reported.
- Some types of run-time errors have been converted to compile-time errors, making it easier to avoid certain types of programming mistakes ([#4651](https://github.com/OPM/opm-common/pull/4651)).
- Added support for gas sales summary vectors [FG]GS[RT] and [FG]SG[RT] ([#4731](https://github.com/OPM/opm-common/pull/4731)).
- The well's preferred phase is now used to calculate the well productivity index (WPI, WPI4, WPI5, and WPI9) summary vectors ([#4770](https://github.com/OPM/opm-common/pull/4770)). Previously, the oil phase was used.
- The performance of the TracerModel has been improved by adding OpenMP parallelization over grid elements ([#6206](https://github.com/OPM/opm-simulators/pull/6206)).
- A new command line parameter (--max-well-status-switch-iter-wells) has been added to limit the maximum number of well switches between stopped and open during a time step ([#6456](https://github.com/OPM/opm-simulators/pull/6456)). This parameter should be used with care as it could lead to a well being shut too early.
- Added a command line option (--max-well-status-switch-in-inner-iter-wells=99) to specify the maximum number of status switches (between shut and open) in inner iterations for wells ([#6338](https://github.com/OPM/opm-simulators/pull/6338)).
- If the IPR-line lies above all (FLO, BHP) points in the lift table (VFP), the intersection is now determined by extrapolation to higher rates based on the last two points ([#6187](https://github.com/OPM/opm-simulators/pull/6187)).
- Only surface rates are allowed to be assigned to satellite production or injection groups ([#6480](https://github.com/OPM/opm-simulators/pull/6480)). If higher level groups containing satellite groups have a reservoir volume rate (RESV) or voidage replacement (VREP) control, the treatment of the satellite rates is not well defined. A check has been added and an error message is now output if this is the case.
- Previously well control changes weren't logged during local switching ([#6145](https://github.com/OPM/opm-simulators/pull/6145)). Well control changes are now logged when the NUPCOL iteration threshold is reached.
- The initialization procedure for production wells is now more robust, in particular it ensures that valid gas and water fractions are obtained for the VFP table lookup ([#6114](https://github.com/OPM/opm-simulators/pull/6114)). Also, the gas and water fractions of a well are stored across iterations to better deal with situations such as stopped wells, where it is not possible to reconstruct gas and water fractions from the rates.
- When detecting oscillations only the number of group switches after NUPCOL non-linear iterations are now considered ([#6358](https://github.com/OPM/opm-simulators/pull/6358)). This makes the behaviour of groups similar to that of wells.
\end{verbatim}

The following bug fixes have been incorporated into this release:

\begin{verbatim}
- Corrected the application of region to region transmissibility multipliers (MULTREGT keyword) to NNCs created by pinched-out cells (PINCH keyword) in some specific cases ([#6498](https://github.com/OPM/opm-simulators/pull/6498)).
- Fixed a small bug for grid topology that did not affect OPM Flow, but did affect the geomechanics codes ([#913](https://github.com/OPM/opm-grid/pull/913)).
- Added a OPM Flow specific option to the Killough hysteresis model (EHYSTR keyword item 14) to correct the construction of the capillary pressure scanning curve ([#4676](https://github.com/OPM/opm-common/pull/4676)). This option is currently inactive by default but this is likely to change in the future as the correction is believed to be a bug fix.
- The old style PYACTION scripts are now assumed to have run successfully if they do not return a value ([#4725](https://github.com/OPM/opm-common/pull/4725)). This fixes backwards compatibility.
- The TUNING keyword now correctly sets the minimum length of all time steps Item 1-3 TSMINZ instead of Item 1-6 TSFMIN ([#6332](https://github.com/OPM/opm-simulators/pull/6332)).
- A bug has been fixed in the unit handling of the OPERATE keyword ([#4765](https://github.com/OPM/opm-common/pull/4765)). This resolves an reported issue with the units when performing operations on permeability arrays ([#4597](https://github.com/OPM/opm-simulators/pull/4597)).
- Allow the WELTARG keyword to accept an empty well list ([#4654](https://github.com/OPM/opm-common/pull/4654)). In this case no well targets will be changed and a warning message will be output.
- A bug has been fixed where previously DRSDT was activated in restart runs even if the original simulation did not contain the DRSDT keyword ([#4655](https://github.com/OPM/opm-common/pull/4655)).
- Fixed the handling of the group production control (GCONPROD) keyword item 2 equals FLD ([#4635](https://github.com/OPM/opm-common/pull/4635)). This is now correctly interpreted as inheriting the control from the group's parent(s).
- If WLIST keyword add (ADD) or move (MOV) operations are performed on a non-existent well list then a new list will now be created ([#4768](https://github.com/OPM/opm-common/pull/4768)). Also, if an attempt is made to delete (DEL) a non-existent well list then an error message is output.
- OPM Flow now respects GCONPROD keyword items 11 to 13 properly ([#6331](https://github.com/OPM/opm-simulators/pull/6331)). These are the actions to be taken if water, gas, and liquid rate constraints are violated. Previously the action would default to that specified by item 7 whichever of the constraints was violated.
- The code now protects against possible division by zero when applying the UDQ comparison operators, which have a relative tolerance defined by the UDQPARAMS keyword item 4 ([#4758](https://github.com/OPM/opm-common/pull/4758)).
- UDQ values are now only updated once at the end of the next time step if the UDQ keyword UPDATE flag is set to NEXT ([#4767](https://github.com/OPM/opm-common/pull/4767)). Previously, the UDQ values were also updated at the beginning of the next report step.
- An error is now reported if the specified initial gas-oil contact is deeper than the initial oil-water contact specified on the EQUIL keyword in the SOLUTION section ([#4771](https://github.com/OPM/opm-common/pull/4771)).
- Fixed the summary (RSM) file generation failure in cases with formatted output ([#4653](https://github.com/OPM/opm-common/pull/4653)). Also, corrected the generation of the SMSPEC file name used as the source when creating RSM file output in certain special cases that were previously not handled.
- Fixed the time step counter in cases where the SUMTHIN or RPTONLY keywords are used ([#4644](https://github.com/OPM/opm-common/pull/4644)).
- Invalid region index pairs specified when requesting inter-region flow summary vectors (RxFT) are now ignored and a warning message is output ([#4704](https://github.com/OPM/opm-common/pull/4704)).
- Since the TIME vector in the summary file is stored as a single precision floating-point number, small time steps can give TIME values that are not strictly increasing especially at long elapsed times ([#4647](https://github.com/OPM/opm-common/pull/4647)). This is known to cause problems for some post-processing tools. Summary output is now filtered to prevent this occurring.
- The summary vector calculation of group and field level cumulative production volumes ([FG]xPT) have been amended to incorporate group level efficiency factors (GEFAC) for satellite production ([#4683](https://github.com/OPM/opm-common/pull/4683)).
- An incorrect reference solTracerRate has been fixed in TracerModel ([#6207](https://github.com/OPM/opm-simulators/pull/6207)).
- A well's preferred phase is now correctly restored following a restart ([#4740](https://github.com/OPM/opm-common/pull/4740)). Previously, all producing wells would be restarted as oil producers, the common case, but this would be wrong for gas producers.
- The last valid network pressure is now stored and used for initialization if the current iteration fails ([#6200](https://github.com/OPM/opm-simulators/pull/6200)).
- The new operability check introduced for wells under THP control has now also been introduced for wells not under THP control ([#6317](https://github.com/OPM/opm-simulators/pull/6317)). This fixes convergence issues seen in a couple of test cases.
- Well group targets are now stored: to make group targets consistent for local switching during well solves, to avoid using global vectors with name mapping from Wellstate, to only use singleWellState in the well switching algorithm, and to avoid recomputing well group targets ([#6311](https://github.com/OPM/opm-simulators/pull/6311)).
- Well status checks based on the well objects in the schedule have been replaced with the ones based on the well state ([#6051](https://github.com/OPM/opm-simulators/pull/6051)).
- A bug has been fixed in the LeftExtreme interpolation of 2D tables (where the y sample values are entered in ascending order) ([#4668](https://github.com/OPM/opm-common/pull/4668)). Only small effects have been observed but could be significant in corner cases.
- Allow restart runs for cases with binary (IMPORT'ed) input files ([#4708](https://github.com/OPM/opm-common/pull/4708)).
- Avoid very small denominator values when the product Rs.Rv or Rsw.Rvw is too large ([#6227](https://github.com/OPM/opm-simulators/pull/6227)). This fixes the spikes seen in the gas injection potentials in a few regression tests.
- Compilation issues have been resolved on architectures where long double maps to __float128 with QuadMath ([#4756](https://github.com/OPM/opm-common/pull/4756), [#6497](https://github.com/OPM/opm-simulators/pull/6497)). This resolves issue [#4702](https://github.com/OPM/opm-common/pull/4702).
- Diagnostic messages for DELAYED_EXIT1 will now be printed at the end of the parsing process ([#4613](https://github.com/OPM/opm-common/pull/4613)).
- Several bugs in the rocSPARSE backend for the GPU have been fixed including a memory leakage that caused a segmentation fault for a big model, and the rocSPARSE based CPR implementation has been updated to be inline with the OPM default CPR implementation ([#6339](https://github.com/OPM/opm-simulators/pull/6339)).
- The baseName has been separated from the file removal regular expression since it's a user-controlled string that might contain regular expression metacharacters like '+' or '-' ([#6293](https://github.com/OPM/opm-simulators/pull/6293)). In rare corner cases, those characters might lead to matching too many or not enough files. The simulator now uses regular string matching for the baseName. Symbolic links and directories are omitted from consideration. Grid files are also omitted from file removal since it's possible for CASE.DATA to use CASE.EGRID as an input file (see the GDFILE keyword).
\end{verbatim}

Known issues in this release of the simulator include:

\begin{verbatim}
- There is an error with the RSM header for summary vectors whose NUMS entry in the SMSPEC file is derived from more than a single number source (e.g., single region or segment ID). This applies to all block vectors (BGPV, BOPV, BWPV, etc.), connection level quantities (COPT, etc.), and inter-region flows such as ROFT etc ([#3078](https://github.com/OPM/opm-common/issues/3078)). The work around is to plot the data in OPM ResInsight and right-click on the plot to view and copy the data.
- As per previous releases of the radial model, the COORDSYS keyword item three must be set to COMP to complete the circle, this has not been implemented in this release. Also there appears to be a bug for full radial models when a well goes on BHP control that causes the well not to respect the BHP constraint, this eventually causes the well to die prematurely. See [#2640](https://github.com/OPM/opm-common/pull/2640) for a discussion on the topic.
- As in previous releases there are some issues with the OPERATE and OPERATER keywords associated with the input parsing; for various reasons a few of the fields require special case treatment in the grid processing, including (at least) MULTZ, PORV, and ACTNUM, and for those keywords the OPERATE/OPERATER keyword does not work. The work around is to use the MULTIPLY keyword instead.
- For the UDQ ASSIGN operator after the terminating “/” normally any comments can be entered; however, if there is “/” within the comment field, as per:
- ASSIGN FUNGLYLD 1.196   /    Condensate Yield (63.5/56.7)/(1.0 – 0.065)
- then the simulator will abort. The work around is to manually place the comment characters “--” after the ASSIGN terminating “/”, like so:
- ASSIGN FUNGLYLD 1.196   / -- Condensate Yield (63.5/56.7)/(1.0 – 0.06)
- At the moment, one cannot initialize tracers using the EQUALS keyword. Instead use the array format, that is the keyword followed by the required number of values, or the TVDP keyword in the SOLUTION section to set the initial tracer concentrations as a function of depth.
- Currently, gas tracers cannot be used if the dissolved gas phase, as per the DISGAS keyword in the RUNSPEC section, is active in the model.
- The summary vector RTIPTHEA, that defines the energy in-place between the initial and the current time for regions, is not supported unlike the FTIPTHEA and BTIPTHEA vectors. Secondly, the error message:
\end{verbatim}

\begin{longtable}[]{@{}
  >{\raggedright\arraybackslash}p{(\linewidth - 0\tabcolsep) * \real{1.0000}}@{}}
\toprule\noalign{}
\begin{minipage}[b]{\linewidth}\raggedright
Warning: Problem with summary keyword RTIPTHEA In RSM-THERMAL.data line
492 FIP region FIPHEA not defined in REGIONS section - RTIPTHEA ignored
\end{minipage} \\
\midrule\noalign{}
\endhead
\bottomrule\noalign{}
\endlastfoot
\end{longtable}

is incorrect, as the message indicates that it is being treated like a
named region, as per the FIP keyword, when it is actually a SUMMARY
vector
(\href{https://github.com/OPM/opm-simulators/issues/3870}{\#3870}).

\begin{verbatim}
- If there are cells that are very distorted, which can occur near fault planes, then the simulator may abort because it cannot calculate the pore volume of such cells. The work around is to re-generate the grid in the static model, taking care that the cells around the fault planes are more or less orthogonal ([#2992](https://github.com/OPM/opm-common/issues/2992) and [#3770](https://github.com/OPM/opm-simulators/issues/3770)).
- Currently the OPERATER keyword in the EDIT section does not work with the DEPTH, TRANX, TRANY, and TRANZ property arrays ([#2994](https://github.com/OPM/opm-common/issues/2994) and [#748](https://github.com/OPM/opm-tests/pull/748)).
- If a standard well is fully declared in an ACTIONX block which is then activated at a later date, and later the well is modified to be a multi-segment well using the WELSEGS and COMPSEGS keywords, then this will cause the simulator to abort with an assert failure. The solution to this issue is to not use this type of work flow in declaring wells ([#2891](https://github.com/OPM/opm-common/issues/2891) and [#2895](https://github.com/OPM/opm-common/pull/2895)).
- Although the ACTIONX EXIT command works as expected, it does not write out the requested RSM file at the end of the run. However, the other SUMMARY and RESTART files are written out ([#2877](https://github.com/OPM/opm-common/issues/2877)).
- There are small differences in the behavior of the NEXTSTEP keyword in the RUNSPEC section between OPM Flow and the commercial simulator that remain unresolved ([#3745](https://github.com/OPM/opm-simulators/issues/3745)).
- OPM cannot be built with dune-fem version 2.9 or later ([#4934](https://github.com/OPM/opm-simulators/issues/4934)). Please use a previous version.
- The simulator uses an irregular corner-point grid geometry with adjusted pore volumes to represent radial grids so it is not possible to create a full ring (360 degree disk) with only one cell in the theta direction (NY=1). The work around is to model a slice (say DTHETA=60 degrees). Note that as the angle increases larger pore volume adjustments are required ([#4755](https://github.com/OPM/opm-simulators/issues/4755)).
- In some cases with the network option the simulator can wrongly report that a well has no THP constraints, for example
\end{verbatim}

\begin{longtable}[]{@{}l@{}}
\toprule\noalign{}
GLIFT WTEST: Well S-3H does not have THP constraints \\
\midrule\noalign{}
\endhead
\bottomrule\noalign{}
\endlastfoot
\end{longtable}

when THP constraints have been defined
(\href{https://github.com/OPM/opm-simulators/issues/4887}{\#4887}).

\begin{verbatim}
- Dispersion in the gas phase leads to convergence issues for the 11th SPE CSP ([https://spe.org/csp](https://spe.org/csp)) models Version 11B and Version 11C (and is not fully tested). Therefore dispersion in the gas phase has been removed and a warning added ([#5101](https://github.com/OPM/opm-simulators/pull/5101), [#859](https://github.com/OPM/opm-models/pull/859)).
\end{verbatim}

The 2025-04 release consists of some new features and various
improvements and bug fixes. These include the following highlights:

\begin{verbatim}
- A command line option --check-satfnc-consistency has been added that enables improved saturation function consistency checking.
- A new hybrid linear-radial filter cake modelling option LINRAD has been added to the [WINJDAM](#REF_HEADING_KEYWORD_WINJDAM) keyword.
- The simulator now prints a message to the log files and standard output for all errors that occur while parsing  the input deck, previously no message was reported for some errors before aborting.
- Gas consumption (GCONSUMP keyword) is now supported at multiple levels in the group hierarchy.
- The linear setup time has been reduced when using two-stage (CPR-like) preconditioners.
- Network solver sub-iterations have been added to avoid updating the gas lift, group controls and so on at every network update. This gives a significant speed up in cases including gas lift and network.
- The ability to pass standard command line parameters to the simulator from Python has been added.
- The handling of the NUPCOL keyword has been updated.
- To avoid group control oscillations the code now limits the maximum number of times a group can switch to the same control.
- Support has been added for using the Hypre library's BoomerAMG solver in serial runs, providing a high-performance algebraic multigrid solver option particularly beneficial for large-scale simulations and the possibility to run the algebraic multigrid (AMG) solver on the GPU.
- Support has been added for the WCYCLE keyword in the SCHEDULE section that defines automatic well opening and closing cycle parameters.
- Support has been added for using the AmgX library in serial runs. AmgX is a high-performance algebraic multigrid solver for NVIDIA GPUs.
- Added support for ROCKOPTS keyword item 2 equal to STORE. If ROCKOPTS item 2 is equal to STORE, the initial pressure is used in ROCK and ROCKTAB keywords.
- The simulator now uses the TPFA linearizer for Black Oil Thermal runs.
- The default load balancer has been changed to Zoltan with the grid represented by a graph. All cells potentially perforated by a well are represented by a single vertex, therefore the partitioning is unable to split a well across more than one process. This has resulted in a general improvement and has resolved convergence issues in some hard cases.
- The WELPI and WPIMULT keywords can now be used with the "insert keywords" function in PYACTION blocks.
- Added partial support for the GSATPROD keyword in the SCHEDULE section to define group satellite production rates.
- Added command line parameters for network solver sub-iterations and pressure update damping.
- Added the following command line arguments to specify load balancing parameters: add corners to partition (--add-corners=false), and numbers of layers overlap in parallel partition (--num-overlap=1).
- Support has been added for instantaneous flow rates in extended network models (WEFAC and GEFAC item 3). The WEFAC and GEFAC keywords are now fully supported.
- The permeability multiplication factor as a function of porosity change (PERMFACT keyword) can now be defined for each saturation table region.
- Support has been added for the WCONHIST and WCONINJH keywords in ACTIONX blocks.
- Added support for multi-segment well output to the PRT file for RPTSCHED mnemonic WELLS=x.
- Support has been added for the [RSW](#REF_HEADING_KEYWORD_RSW_10_3) keyword in the SOLUTION section to explicitly initialise the dissolved gas-water ratio in each cell when dissolved gas in the water phase has been activated.
- Defaulted tables are now allowed in the SWOF, SGOF, SWFN, SGFN and SLGOF keywords provided the first table has been defined. A defaulted table is set equal to the previous table.
- Added support for outputting saturation function tables defined using the SLGOF keyword to the initial data (INIT) file.
- The WELTRAJ and COMPTRAJ keywords now handle any MAPAXES specifications.
- Support has been added for using the special name selector '?' in well level UDQ assignments within an ACTIONX block.
- Support has been added for gas consumption and gas import summary vectors.
- Threshold pressures (THPRES keyword) are now restored from the RESTART file.
- Added support for the WTMULT and WLIST keywords in PYACTION blocks.
- Support has been added for the GEFAC keyword in ACTIONX blocks and in PYACTION code.
- The use of named regions is now supported in ACTIONX conditions.
- Support has been added for summary output of group and field production and injection targets.
- Support has been added for the CO2STORE/[H2STORE](#REF_HEADING_KEYWORD_H2STORE) specific block level CO2/H2 mass, moles and volumes in place summary variables; mass production rate (GMPR) and mass production total (GMPT) variables; and mass injection/production variables at the connection and completion level (for example, CGMIR and CGMIRL).
- The microbially induced calcite precipitation model (see the MICP keyword) has been updated to use the new lineariser, support dispersion, diffusion, sources, boundary conditions, and new summary variables. In addition, the new [BIOFPARA](#REF_HEADING_KEYWORD_BIOFPARA) keyword (which can be set per SATNUM region) is now used to set the parameters instead of the MICPPARA keyword.
\end{verbatim}

The major command line changes made for this release are summarized in
Table A.1

\begin{longtable}[]{@{}
  >{\raggedright\arraybackslash}p{(\linewidth - 6\tabcolsep) * \real{0.2500}}
  >{\raggedright\arraybackslash}p{(\linewidth - 6\tabcolsep) * \real{0.2500}}
  >{\raggedright\arraybackslash}p{(\linewidth - 6\tabcolsep) * \real{0.2500}}
  >{\raggedright\arraybackslash}p{(\linewidth - 6\tabcolsep) * \real{0.2500}}@{}}
\toprule\noalign{}
\begin{minipage}[b]{\linewidth}\raggedright
OPM Flow 2025-04 New and Deprecated Command Line Options
\end{minipage} & \begin{minipage}[b]{\linewidth}\raggedright
\end{minipage} & \begin{minipage}[b]{\linewidth}\raggedright
\end{minipage} & \begin{minipage}[b]{\linewidth}\raggedright
\end{minipage} \\
\midrule\noalign{}
\endhead
\bottomrule\noalign{}
\endlastfoot
No. & Variable Name & Description & Default \\
1 & --add-corners & A boolean value. Add corners to partition. &
false \\
2 & --allow-splitting-inactive-wells & A boolean value. Allow inactive
(never non-shut) wells to be split across multiple domains. & true \\
3 & --bda-device-id & Renamed as --gpu-device-id. & 0 \\
4 & --check-satfunc-consistency & A boolean value. Whether or not to
check saturation function consistency requirements. & false \\
5 & --gpu-device-id & An integet value. Choose device ID for
cusparseSolver or openclSolver, use `nvidia-smi' or `clinfo' to
determine valid IDs. & 0 \\
6 & --local-domains-partition-well-neighbor-levels & An integer value.
Number of neighbor levels around wells to include in the same domain
during NLDD partitioning. & 1 \\
7 & --maximum-number-of-group-switches & An integer value. Maximum
number of times a group can switch to the same control. & 3 \\
8 & --network-max-iterations & Replaced by
--network-max-outer-iterations. & 200 \\
9 & --network-max-outer-iterations & An integer value. Maximum outer
number of iterations in the network solver before giving up. & 10 \\
10 & --network-max-pressure-update-in-bars & A real value. Maximum
pressure update in the inner network pressure update iterations. & 5 \\
11 & --network-max-strict-iterations & Replaced by
--network-max-strict-outer-iterations. & 100 \\
12 & --network-max-strict-outer-iterations & An integer value. Maximum
outer iterations in network solver before relaxing tolerance. & 10 \\
13 & --network-max-sub-iterations & An integer value. Maximum number of
sub-iterations to update network pressures (within a single well/group
control update). & 20 \\
14 & --network-pressure-update-damping-factor & A real value. Damping
factor in the inner network pressure update iterations. & 0.1 \\
15 & --num-overlap & An integer value. Numbers of layers overlap in
parallel partition. & 1 \\
16 & --num-satfunc-consistency-sample-points & Maximum number of
reported failures for each individual saturation function consistency
check. & 5 \\
17 & --nupcol-group-rate-tolerance & A real value. Tolerance for
acceptable changes in VREP/RAIN group rates. & 0.001 \\
18 & --partition-method & An integer value. Choose partitioning
strategy: 0=simple, 1=Zoltan, 2=METIS, 3=Zoltan with all cells of well
represented by one vertex. & 3 \\
19 & --pre-solve-network & A boolean value. Pre solve and iterate the
network model at start-up. & true \\
20 & --slave & A boolean value. Specify if the simulation is a slave
simulation in a master-slave simulation. & false \\
21 & --threads-per-process & An integer value. The maximum number of
threads to be instantiated per process (`-1' means `automatic'). & 2 \\
22 & --time-step-control-safety-factor & A real value to be multiplied
with the time step control tolerance to ensure that the target relative
change is lower than the tolerance. & 0.8 \\
23 & --zoltan-phg-edge-size-threshold & A real value. Low-level
threshold fraction in the range {[}0,1{]} controlling which hypergraph
edge to omit. Used if --zoltan-params=``graph'' or if
--zoltan-params=``hypergraph''. & 0.35 \\
Notes: & & & \\
\end{longtable}

\emph{Table A.1: OPM Flow 2025-04 New and Deprecated Command Line
Options}

In addition to the changes to the command line options the following new
features have been added to the simulator:

\begin{verbatim}
- A command line option --check-satfnc-consistency has been added that enables improved saturation function consistency checking. In addition, the new command line option --num-satfnc-consistency-sample-points allows the user to select the maximum number of reported failures for each individual consistency check. By default the simulator will report at most five failures for each check. If end-point scaling has been activated then the scaled curves are checked otherwise the unscaled curves are checked ([#5596](https://github.com/OPM/opm-simulators/pull/5596)).
- A new hybrid linear-radial filter cake modelling option LINRAD has been added to the [WINJDAM](#REF_HEADING_KEYWORD_WINJDAM) keyword in the SCHEDULE section ([#5717](https://github.com/OPM/opm-simulators/pull/5717)). The cake thickness is computed assuming linear geometry. This thickness is then used to compute a damage skin in radial geometry.
- Running flow with the parameter --partition-method=3 will use Zoltan partitioner with the grid represented by a graph ([#5749](https://github.com/OPM/opm-simulators/pull/5749)). All cells potentially perforated by a well are represented by a single vertex, therefore the partitioning is unable to split a well into several processes.
\end{verbatim}

This should lead to better partitioning results on problems with a
complex well structure, where other partitioners might split wells over
several processes after which some well cells would have to be remapped
(when distributed wells are not allowed). - An experimental feature has
been added that enables multi-segment wells to be distributed across
several processes when using the command line argument
--allow-distributed-wells=true
(\href{https://github.com/OPM/opm-simulators/pull/5746}{\#5746}). -
Support has been added for using the Hypre library's BoomerAMG solver,
providing a high-performance algebraic multigrid solver option
particularly beneficial for large-scale simulations and the possibility
to run the algebraic multigrid (AMG) solver on the GPU
(\href{https://github.com/OPM/opm-simulators/pull/5762}{\#5762}). -
Support has been added for the WCYCLE keyword in the SCHEDULE section
that defines automatic well opening and closing cycle parameters
(\href{https://github.com/OPM/opm-simulators/pull/5792}{\#5792}). -
Support has been added for using the AmgX library in serial runs
(\href{https://github.com/OPM/opm-simulators/pull/5808}{\#5808}). AmgX
is a high-performance algebraic multigrid solver for NVIDIA GPUs. -
Added support for ROCKOPTS keyword item 2 equal to STORE
(\href{https://github.com/OPM/opm-simulators/pull/5785}{\#5785}). If
ROCKOPTS item 2 is equal to STORE, the initial pressure is used in ROCK
and ROCKTAB keywords. - A mixed precision option has been added to GPU
DILU similar to that for GPU ILU0
(\href{https://github.com/OPM/opm-simulators/pull/5674}{\#5674}). - A
mixed precision option where off-diagonals are stored as float and
diagonals are stored as double has been added to both GPU DILU and GPU
ILU0 (\href{https://github.com/OPM/opm-simulators/pull/5688}{\#5688}). -
The WELPI and WPIMULT keywords can now be used with the ``insert
keywords'' function in PYACTION blocks
(\href{https://github.com/OPM/opm-simulators/pull/5892}{\#5892}). - The
following cubic equations of state can now be selected using the
\hyperref[REF_HEADING_KEYWORD_EOS_5_3]{EOS} keyword: Peng-Robinson (PR),
modified Peng-Robinson (PRCORR), Redlich-Kwong (RK) and
Soave-Redlich-Kwong (SRK) when using the experimental compositional
simulator
(\href{https://github.com/OPM/opm-simulators/pull/5886}{\#5886}). -
Added partial support for the GSATPROD keyword in the SCHEDULE section
to define group satellite production rates
(\href{https://github.com/OPM/opm-simulators/pull/5894}{\#5894}). -
Added command line parameters for network solver sub-iterations and
pressure update damping
(\href{https://github.com/OPM/opm-simulators/pull/5918}{\#5918}). See
the following command line arguments: maximum number of sub-iterations
to update network pressures, within a single well/group control update
(--network-max-sub-iterations), damping factor in the inner network
pressure update iterations (--network-pressure-update-damping-factor),
and maximum pressure update in the inner network pressure update
iterations (--network-max-pressure-update-in-bars). - Added the
following command line arguments to specify load balancing parameters:
add corners to partition (--add-corners=false), and numbers of layers
overlap in parallel partition (--num-overlap=1)
(\href{https://github.com/OPM/opm-simulators/pull/5942}{\#5942}). - The
experimental compositional simulator has been expanded to three phases
by the addition of a (``dummy'') immiscible water phase
(\href{https://github.com/OPM/opm-simulators/pull/5851}{\#5851}). -
Support has been added for instantaneous flow rates in extended network
models (WEFAC and GEFAC item 3). The WEFAC and GEFAC keywords are now
fully supported
(\href{https://github.com/OPM/opm-simulators/pull/5410}{\#5410}). - A
third-order time step controller has been added that can be used as an
alternative to the PID controller
(\href{https://github.com/OPM/opm-simulators/pull/5974}{\#5974}). The
new controller can be specified by using the
--time-step-control=``general3rdorder'' command line argument. - Added a
command line parameter (--pre-solve-network=true) to make the pre-step
network rebalance optional
(\href{https://github.com/OPM/opm-simulators/pull/6048}{\#6048}). -
Added command line parameter (--zoltan-phg-edge-size-threshold=0.35)
that controls which hypergraph edges to omit/discard by the Zoltan
partitioner
(\href{https://github.com/OPM/opm-simulators/pull/6036}{\#6036}). This
is to enable runtime experimentation with this aspect of the
partitioning algorithm, the default value should not normally be
altered. - The permeability multiplication factor as a function of
porosity change (PERMFACT keyword) can now be defined for each
saturation table region
(\href{https://github.com/OPM/opm-simulators/pull/6074}{\#6074}). -
Added support for RPTSCHED mnemonic WELLS=2. In particular, this appends
connection level results to the injection/production output sent to the
print (PRT) file
(\href{https://github.com/OPM/opm-simulators/pull/6099}{\#6099}). -
Support has been added for the WCONHIST and WCONINJH keywords in ACTIONX
blocks (\href{https://github.com/OPM/opm-simulators/pull/6092}{\#6092}).
- Added support for multi-segment well output to the PRT file for
RPTSCHED mnemonic WELLS=x
(\href{https://github.com/OPM/opm-simulators/pull/6110}{\#6110}). -
Support has been added for the
\hyperref[REF_HEADING_KEYWORD_RSW_10_3]{RSW} keyword in the SOLUTION
section to explicitly initialise the dissolved gas-water ratio in each
cell when dissolved gas in the water phase has been activated by
specifying the DISGASW keyword in the RUNSPEC section
(\href{https://github.com/OPM/opm-simulators/pull/6111}{\#6111}). -
Added support for block summary vectors related to mass and volume
in-place for CO2STORE and
\hyperref[REF_HEADING_KEYWORD_H2STORE]{H2STORE} runs including: BGMIP,
BGMGP, BGMDS, BGMST, BGMUS, BGMTR, BGMMO, BGKTR, BGKMO, BGCDI, BGCDM,
BGKDI, BGKDM, BWCD, BWIPG, BWIPL
(\href{https://github.com/OPM/opm-simulators/pull/6127}{\#6127}). -
Added the option to merge each well with a layer(s) of cells surrounding
it, in order to keep the well further from the subdomain boundary, when
partitioning the grid
(\href{https://github.com/OPM/opm-grid/pull/811}{\#811}). To use this
option, flow must be called with parameter --zoltan-params=*.json taking
a json file. The json file must contain the keyword EnvelopeWellLayers
with the number of layers (default 0). The feature is supported only by
the partitioner using the GraphOfGrid (--partition-method=3). - Added an
option to do serial partitioning with Zoltan using the graph
representation of grid (GraphOfGrid)
(\href{https://github.com/OPM/opm-grid/pull/821}{\#821}). The parallel
partitioning with GraphOfGrid is also available and is currently the
default partitioner. - The zoltanGoG (Zoltan Graph of Grid) partitioner
has replaced zoltan as the default partitioner (if there are no wells
zoltanGoG is equivalent to zoltan)
(\href{https://github.com/OPM/opm-grid/pull/840}{\#840}). - Defaulted
tables are now allowed in the SWOF, SGOF, SWFN, SGFN and SLGOF keywords
provided the first table has been defined
(\href{https://github.com/OPM/opm-common/pull/4308}{\#4308}). A
defaulted table is set equal to the previous table. - Added support for
`LINRAD' as a valid option for item 2 (GEOMETRY) of the keyword
\hyperref[REF_HEADING_KEYWORD_WINJDAM]{WINJDAM}
(\href{https://github.com/OPM/opm-common/pull/4320}{\#4320}). This
option specifies a hybrid filter cake model where the cake thickness is
computed assuming linear geometry (as with item 2 = LINEAR). This
thickness is then used to compute a damage skin in radial geometry (as
with item 2 = RADIAL). - Support has been added for gas consumption and
gas import summary vectors GGCR, GGIMR, GGCT, GGIMT, FGCR, FGIMR, FGCT
and FGIMT and for the corresponding restart input/output
(\href{https://github.com/OPM/opm-common/pull/4342}{\#4342}). - Support
has been added for group name pattern matching in UDQ assignments
(\href{https://github.com/OPM/opm-common/pull/4433}{\#4433}). - Support
has been added for summary output of group (and field) production and
injection targets
(\href{https://github.com/OPM/opm-common/pull/4500}{\#4500}). The
following summary vectors are now supported: G{[}OWGLV{]}PRT,
G{[}WGV{]}IRT, F{[}OWGLV{]}PRT, and F{[}WGV{]}IRT. - Support has been
added for the CO2STORE/\hyperref[REF_HEADING_KEYWORD_H2STORE]{H2STORE}
specific block level CO2/H2 mass, moles and volumes in place summary
variables; mass production rate (GMPR) and mass production total (GMPT)
variables; and mass injection/production variables at the connection and
completion level (for example, CGMIR and CGMIRL)
(\href{https://github.com/OPM/opm-common/pull/4550}{\#4550}). - The
microbially induced calcite precipitation model (see the MICP keyword)
has been updated to use the new lineariser, support dispersion,
diffusion, sources, boundary conditions, and new summary variables
(\href{https://github.com/OPM/opm-common/pull/4538}{\#4538}). In
addition, the new \hyperref[REF_HEADING_KEYWORD_BIOFPARA]{BIOFPARA}
keyword (which can be set per SATNUM region) is now used to set the
parameters instead of the MICPPARA keyword.

Improvements to simulator in this release include the following:

\begin{verbatim}
- When partitioning, the simulator now only accounts for wells that might be active at some point in time during the run ([#5609](https://github.com/OPM/opm-simulators/pull/5609)).
- Logging has been added for PYACTION blocks similar to the current logging for ACTIONX blocks ([#5699](https://github.com/OPM/opm-simulators/pull/5699)).
- The experimental compositional simulator now outputs INIT, EGRID, SMSPEC, UNRST and UNSMRY files ([#5687](https://github.com/OPM/opm-simulators/pull/5687)).
- The device name and compute capability are now output for every rank ([#5611](https://github.com/OPM/opm-simulators/pull/5611)).
- Since the code can now switch to/from GRUP control in the inner well iterations this is problematic because the group reduction rates (the combined rates of all wells that are not on group rate) are not updated, which in turn may result in excessive switching ([#5695](https://github.com/OPM/opm-simulators/pull/5695)). As a temporary solution a new command line argument has been added (--check-group-constrains-well-iterations=false). Setting --check-group-constrains-well-iterations=true means that individual well constraints will not be checked during the inner well iterations if the well is on GRUP control.
- The code no longer repeatedly recomputes the well BHP from the THP when evaluating well potentials for gaslift optimisation ([#5696](https://github.com/OPM/opm-simulators/pull/5696)). This gives a 10 times speed up of the gaslift optimisation for many cases since the BHP is not repeatedly computed for ALQ values that don't converge.
- The simulator now prints a message to the log files and standard output for all errors that occur while parsing  the input deck, previously no message was reported for some errors before aborting ([#5707](https://github.com/OPM/opm-simulators/pull/5707)).
- If one well operating on a well constraint (typically THP or BHP) in a group struggles to converge it may lead to problems for the other wells in the group. In this case, the simulator now only shuts the unconverged well rather than all the wells in the group ([#5735](https://github.com/OPM/opm-simulators/pull/5735)).
- Previously, using the WTEST keyword on a gas lifted well typically resulted in the maximum artificial lift quantity (ALQ) being used after the well was reopened. The simulator now optimises the ALQ to give the maximum incremental gradient. This has reduced the large spikes in ALQ seen after reopening wells in the gas lift test cases ([#5757](https://github.com/OPM/opm-simulators/pull/5757)).
- Debug output from THP calculations previously output to the PRT file will now be output to the DBG file ([#5761](https://github.com/OPM/opm-simulators/pull/5761)). The volume of output can be large especially for cases with gas lift.
- Gas consumption (GCONSUMP keyword) is now supported at multiple levels in the group hierarchy ([#5739](https://github.com/OPM/opm-simulators/pull/5739)).
- The simulator now avoids duplicating the output of warnings and errors to the standard error stream (stderr) ([#5765](https://github.com/OPM/opm-simulators/pull/5765)).
- The linear setup time has been reduced when using two-stage (CPR-like) preconditioners. Previously, the second stage ILU0 preconditioner was unnecessarily recreated every Newton iteration. The simulator now only updates the preconditioner ([#5758](https://github.com/OPM/opm-simulators/pull/5758)). This will be especially important when using GPU preconditioners with CPR, as these GPU preconditioners typically involve expensive GPU allocations, matrix analysis and autotuning.
- Network solver sub-iterations have been added to avoid updating the gas lift, group controls and so on at every network update ([#5767](https://github.com/OPM/opm-simulators/pull/5767)). This gives a significant speed up in cases including gas lift and network.
- The ability to pass standard command line parameters to the simulator from Python has been added ([#5760](https://github.com/OPM/opm-simulators/pull/5760)).
- The handling of the NUPCOL keyword has been updated ([#5724](https://github.com/OPM/opm-simulators/pull/5724)). Wells are no longer allowed to change to group control when the number of iterations is greater than NUPCOL. Voidage replacement and reinjection target rates are updated even if the number of iterations is greater than NUPCOL provided the change is greater than a specified relative tolerance (see --nupcol-group-rate-tolerance).
- The fine-level smoothing parameters (pre_smooth and post_smooth) for the CPR AMG implementation in OPM have been exposed by adding them to the linear solver JSON configuration tree output to the DBG file ([#5786](https://github.com/OPM/opm-simulators/pull/5786)).
- To avoid group control oscillations the code now limits the maximum number of times a group can switch to the same control ([#5753](https://github.com/OPM/opm-simulators/pull/5753)). The number of switches is only checked when the number of iterations is greater than NUPCOL (that is multiple switches are allowed whilst the network is iterating). See the command line argument --maximum-number-of-group-switches.
- The code now ensures that tracer flows for cross-flowing injector connections are included ([#5789](https://github.com/OPM/opm-simulators/pull/5789)).
- The 16 day restriction on timesteps in prediction mode with THP constraints has been removed from the adaptive time stepping algorithm ([#2970](https://github.com/OPM/opm-simulators/pull/2970)). The time-stepping, convergence and well models have all improved since this hard-coded limit was originally put in place.
- Group guide rates are now allowed in cases where a group has an automatic choke in an extended network model ([#5754](https://github.com/OPM/opm-simulators/pull/5754)).
- The memory requirement has been reduced by 15 bytes per cell by reducing the size of enums ([#5841](https://github.com/OPM/opm-simulators/pull/5841)).
- The simulator now uses the Two Point Flux Approximation (TPFA) linearizer for Black Oil Thermal runs ([#5843](https://github.com/OPM/opm-simulators/pull/5843)).
- The default load balancer has been changed to Zoltan with the grid represented by a graph ([#5800](https://github.com/OPM/opm-simulators/pull/5800)). All cells potentially perforated by a well are represented by a single vertex, therefore the partitioning is unable to split a well across more than one process. Also the weights of the faces will be added up due to the merging of cells. A cell that is not perforated by any well is still represented by one vertex with weight one. This has resulted in a general improvement and has resolved convergence issues in some hard cases.
- If the partition has been loaded from a file containing an externally generated partition (--external-partition="filename") and the output of  binary files compatible with the commercial simulator has been deactivated (--enable-ecl-output=false), then the expensive routine allocTrans() is now skipped ([#5817](https://github.com/OPM/opm-simulators/pull/5817)).
- Repeated calculation of the cell centroids has been avoided by caching the values ([#5888](https://github.com/OPM/opm-simulators/pull/5888)). This can significantly reduce the setup time for large models.
- The default maximum number of threads per process has been changed from -1 ('automatic') to 2 (see command line argument --threads-per-process=2). If OpenMP is available the maximum number of threads will be defaulted to 2 unless the environment variable OMP_NUM_THREADS is set or the command line argument is used. Note that the OMP_NUM_THREADS takes precedence over the --threads-per-process command line argument ([#5856](https://github.com/OPM/opm-simulators/pull/5856)).
- Empty command line arguments are now ignored, this avoids error messages of the form "Parameter 'EclDeckFileName' specified multiple times as a command line parameter" ([#5679](https://github.com/OPM/opm-simulators/pull/5679)).
- In an extended network model the bottommost (leaf) nodes associated with wells must be the same as those defined by GRUPTREE keyword. However, the extended network model will now allow empty leaf nodes (with no wells) that do not exist in the GRUPTREE ([#5919](https://github.com/OPM/opm-simulators/pull/5919), [#6010](https://github.com/OPM/opm-simulators/pull/6010)).
- Improved the logic for shutting failed wells ([#5900](https://github.com/OPM/opm-simulators/pull/5900)). In particular, allow for failed wells to shut even though they haven't failed consistently in the last 3 attempts if the time step has already been reduced enough.
- A check has been added to guard against high roughness values in the WELSEGS keyword (that can result in a singularity) whilst reading the input deck ([#5945](https://github.com/OPM/opm-simulators/pull/5945)).
- Information about wells being stopped or revived before the timestep has converged will now only be sent to the debug (DBG) file. However, changes to the well status when the timestep has converged will still be output to the terminal and the print (PRT) file ([#6012](https://github.com/OPM/opm-simulators/pull/6012)).
- The stagnation and oscillation code for multi-segment wells has been modified. If the well stagnates five times with full damping then the well is marked as unconverged and the simulator tries again at next Newton iteration ([#6020](https://github.com/OPM/opm-simulators/pull/6020)).
- The output of some NaN's has been prevented for ratios of rates, for example produced gas-oil ratio ([#6077](https://github.com/OPM/opm-simulators/pull/6077)). If the numerator is zero then zero is now output instead of dividing it by the denominator that might also be zero.
- The gaslift optimiser information messages are now output to the debug (DBG) file rather than the print (PRT) file to avoid exceeding the message limit ([#6082](https://github.com/OPM/opm-simulators/pull/6082)).
- Updates generated by PYACTION or ACTIONX blocks are now collected and applied together ([#6070](https://github.com/OPM/opm-simulators/pull/6070)). Previously, the updates would have been applied one at a time, potentially causing unnecessary recalculations.
- When testing a well (see WTEST keyword) belonging to a group that corresponds to a node in an extended network model (see NETWORK keyword) the well's dynamic THP constraint will now be set based on the network's nodal pressure ([#6116](https://github.com/OPM/opm-simulators/pull/6116)).
- Previously, the original fluid in place values in a restart run would be zero ([#6119](https://github.com/OPM/opm-simulators/pull/6119)). If a restart file is available for report step zero then this will now be used to generate the original fluids in place.
- The well rates used to avoid oscillations are now taken from the network iterations (see NUPCOL keyword) rather than using the previous well rates ([#6128](https://github.com/OPM/opm-simulators/pull/6128)). This is expected to give more accurate results in particular when wells startup.
- The maximum number of Newton iterations used to updated well targets (see NUPCOL keyword) is now applied correctly, rather than the specified value plus one ([#6100](https://github.com/OPM/opm-simulators/pull/6100)). Also, gas lift is no longer optimized in the last iteration to allow the network to converge (unless the maximum number of iterations is less than two).
- The value of the saturated dissolution factor is now stored giving a significant speed up in cases with convective mixing (see DRSDTCON keyword) ([#5774](https://github.com/OPM/opm-simulators/pull/5774)).
- Some bugs that were present in the Non-Linear Domain Decomposition (NLDD) reporting have been fixed, and the way statistics from the NLDD solver are collected, displayed, and analysed has been improved ([#6124](https://github.com/OPM/opm-simulators/pull/6124)).
- The well equation is no longer re-solved when computing the initial guess for BHP at the THP constraint if the well state is updated ([#6081](https://github.com/OPM/opm-simulators/pull/6081)).
- Non-neighbour connections generated due to pinched out cells (see PINCH keyword) have been added to the log output ([#798](https://github.com/OPM/opm-grid/pull/798)).
- The code now uses unsigned integers when processing the grid to allow larger grid sizes ([#818](https://github.com/OPM/opm-grid/pull/818)).
- Support for lowercase file extensions has been added to the file conversion utility (convertECL) ([#4274](https://github.com/OPM/opm-common/pull/4274)).
- Added support for outputting saturation function tables defined using the SLGOF keyword to the initial data (INIT) file ([#4307](https://github.com/OPM/opm-common/pull/4307)).
- Added GPU support for evaluation of CO2 PVT functions ([#4222](https://github.com/OPM/opm-common/pull/4222)).
- The simulator now allows a non-zero artificial lift quantity (ALQ) to be used in the defining expression of a user defined quantity (UDQ) in the case where no vertical flow performance (VFP) tables have been provided ([#4324](https://github.com/OPM/opm-common/pull/4324)).
- The simulator now gives a more meaningful error message if the well connections defined in a COMPDAT keyword within an ACTIONX block are not part of the grid ([#4336](https://github.com/OPM/opm-common/pull/4336)).
- Added support for restart runs with the standard network model, where the network dimensions have not been defined by the NETWORK keyword in the RUNSPEC section ([#4337](https://github.com/OPM/opm-common/pull/4337)).
- Support has been added for using the special name selector '?' in well level UDQ assignments within an ACTIONX block ([#4343](https://github.com/OPM/opm-common/pull/4343)). In particular, OPM Flow now supports ACTIONX blocks like
- ACTIONX
- TESTW 25 1 /
- DAY > 10 AND /
- WMCTL '*PROD?' = 0 /
- /
- UDQ
- ASSIGN WU_TEST '?' 1 /
- /
- ENDACTIO
- in which the selector '?' represents all wells from all well lists matching the well list template '*PROD?' for which the WMCTL is zero, meaning the well is shut.
- Support has been added for using a leading backslash (\) to disambiguate a leading wildcard characters (* or ?) in a name pattern as being truly a wildcard, for example wells matching '\*PROD' rather than the well list '*PRODS' ([#4345](https://github.com/OPM/opm-common/pull/4345)).
- Threshold pressures (THPRES keyword) are now restored from the RESTART file ([#5743](https://github.com/OPM/opm-simulators/pull/5743), [#4347](https://github.com/OPM/opm-common/pull/4347)).
- The code will now issue a warning message if the THPRES keyword is present in a restart run (and it will be ignored) ([#4347](https://github.com/OPM/opm-common/pull/4347)). When performing a restart the solution section should be updated, and in particular the THPRES keyword should be removed. The THPRES values should be read from the restart file instead.
- Added restart file support for field level user defined arguments (UDA) ([#4357](https://github.com/OPM/opm-common/pull/4357)).
- The default well injection temperature (see the WTEMP keyword in the SCHEDULE section) has been changed from standard conditions to zero degrees Celsius in runs using the temperature modelling option (see the TEMP keyword in the RUNSPEC section) ([#4042](https://github.com/OPM/opm-common/pull/4042)).
- The depth correction algorithm for the WBPn well block pressure summary vectors has been revised ([#4376](https://github.com/OPM/opm-common/pull/4376)). The revised approach uses the common mixture density, but corrects each cell pressure individually based on that cell's centre depth before aggregating the per-connection contributions.
- Added support for the WTMULT and WLIST keywords in PYACTION blocks ([#4397](https://github.com/OPM/opm-common/pull/4397), [#4402](https://github.com/OPM/opm-common/pull/4402)).
- Variable and function names in the Python Schedule class have been changed from "timestep" to "reportstep" where these refer to report steps rather than time steps ([#4412](https://github.com/OPM/opm-common/pull/4412)).
- Added documentation to Python classes: SummayState, Schedule, Connection, Deck, DeckKeyword, DeckItem, DeckRecord, UDAValue, EclipseGrid, EclipseState, ESmry and EGrid ([#4405](https://github.com/OPM/opm-common/pull/4405), [#4413](https://github.com/OPM/opm-common/pull/4413), [#4414](https://github.com/OPM/opm-common/pull/4414), [#4428](https://github.com/OPM/opm-common/pull/4428), [#4497](https://github.com/OPM/opm-common/pull/4497)).
- Support has been added for the GEFAC keyword in ACTIONX blocks and in PYACTION code ([#4408](https://github.com/OPM/opm-common/pull/4408)).
- The use of named regions is now supported in ACTIONX conditions, for example 'RPR 1 RE3 < 215 /' ([#4449](https://github.com/OPM/opm-common/pull/4449)).
- Added support for writing the combined thermal conductivity of rock plus fluid array (THCONR) to the initial data (INIT) file ([#4474](https://github.com/OPM/opm-common/pull/4474)).
- If necessary the gas PVT tables entered using either the PVDG or PVTG keywords are now extended down to a pressure of 1 barsa ([#3779](https://github.com/OPM/opm-common/pull/3779)). This change guards against potential negative formation volume factors (densities) and/or viscosities when extrapolating the original tables to pressures in the order of one atmosphere.
- The surface densities of CO2 and brine are now output to the terminal in the same units as the input deck for CO2STORE/[CO2SOL](#REF_HEADING_KEYWORD_CO2SOL) runs ([#4510](https://github.com/OPM/opm-common/pull/4510)).
- The simulator now avoids calculating the vapour pressure when not required for the CO2-Brine PVT model ([#4515](https://github.com/OPM/opm-common/pull/4515)). Note that in Spycher and Pruess 2010 they use 31 °C which is inconsistent with the critical temperature (30.95 °C) used in OPM Flow.
- The simulator now properly restores gas-lift optimisation parameters following a restart ([#4517](https://github.com/OPM/opm-common/pull/4517)).
- NODEPROP keyword Item 4 (GASLIFT) for leaf nodes in an extended network model is now saved to and restored from the restart file ([#4499](https://github.com/OPM/opm-common/pull/4499)).
- The COMPSEGS keyword now allows the distance to the start of the connection (Item 5) to be equal to the distance to the end of the connection (Item 6) to define very short completion intervals ([#4265](https://github.com/OPM/opm-common/pull/4265)).
- A warning is now issued if the depth change along a well segment is larger than the segment length ([#3976](https://github.com/OPM/opm-common/pull/3976)).
- Invalid keyword combinations encountered whilst parsing ACTIONX blocks are now ignored ([#4511](https://github.com/OPM/opm-common/pull/4511)).
- If a well is completed in a numerical aquifer cell then the connection properties are now taken from the numerical aquifer rather than the grid cell ([#4547](https://github.com/OPM/opm-common/pull/4547)). Previously, this could give not-a-number (NaN) values for the permeability-thickness (kh) product and the connection length since the grid's permeability would be explicitly reset to zero in numerical aquifer cells. This impacts the WELSPECS report generated by the RPTSCHED keyword and the SCON restart file array.
- Report step zero will now always be written to the restart file (unless this has been explicitly turned off) ([#4545](https://github.com/OPM/opm-common/pull/4545)). This is used to generate initial fluid in place (FIP) reports in restart runs.
- Support has been restored for the old-style "integer control" format for PRT file report generation using the RPTSCHED keyword ([#4526](https://github.com/OPM/opm-common/pull/4526)).
- The lookup of source terms defined by the [SOURCE](#REF_HEADING_KEYWORD_SOURCE_12_3) keyword is now much faster ([#4553](https://github.com/OPM/opm-common/pull/4553)). Previously, the lookup could be very slow when using the [SOURCE](#REF_HEADING_KEYWORD_SOURCE_12_3) keyword to couple OPM Flow with another simulator where there could be a source term in every cell.
\end{verbatim}

The following bug fixes and improvements have been incorporated into
this release:

\begin{verbatim}
- A bug that resulted in an out-of-bounds error when restarting parallel runs with tracers has been fixed ([#5665](https://github.com/OPM/opm-simulators/pull/5665)). The tracer values are now the same in a parallel restart run as in the equivalent sequential restart run.
- A memory bug that affected runs where the NLDD nonlinear solver was used and where there was an unconverged local NLDD domain, which had a multi-segment well on that domain, has been fixed ([#5675](https://github.com/OPM/opm-simulators/pull/5675)).
- The bottom hole temperature for stopped injection wells is now assumed to be equal to the reservoir temperature ([#5605](https://github.com/OPM/opm-simulators/pull/5605)). Previously, the injection temperature was used for stopped injection wells.
- A number of fixes have been applied to the gaslift code for cases with complex group controls with liquid or water constraints ([#5694](https://github.com/OPM/opm-simulators/pull/5694)).
- Fixed a bug where a local index was used instead of a global index when extracting well contributions to add to the matrix for NLDD with the GPU bridge activated ([#5720](https://github.com/OPM/opm-simulators/pull/5720)).
- The simulator no longer tries to compute the initial solution for inactive multi-segment wells split across processes ([#5751](https://github.com/OPM/opm-simulators/pull/5751)).
- Fixed a bug where the total group rate was used to check a higher level group constraint rather than just the portion subject to the higher level control ([#5729](https://github.com/OPM/opm-simulators/pull/5729)).
- Avoid updating the state of wells that are not under group control when updating and communicating group data ([#5748](https://github.com/OPM/opm-simulators/pull/5748)).
- A potential deadlock that could occur if the Unsymmetric MultiFront (UMFPACK) solver encountered a singular matrix has been fixed ([#5773](https://github.com/OPM/opm-simulators/pull/5773)).
- The simulator now aways calculates the initial in place volumes since these may be used by various keywords in the SCHEDULE section even if FIP is not requested by the RPTSOL keyword ([#5784](https://github.com/OPM/opm-simulators/pull/5784)).
- The code now allows efficiency scaling factors to be queried in off-process wells where required ([#5816](https://github.com/OPM/opm-simulators/pull/5816)). Previously, this could cause MPI to hang due to an exception thrown by well_state[name].
- When calculating the voidage replacement (VREP) injection rate for a given phase the simulator now accounts for possible injection rate controls on the other phases ([#5783](https://github.com/OPM/opm-simulators/pull/5783)).
- The RFT output has been fixed for shut wells (and in the case of partition method 3 for inactive wells) split across processors ([#5824](https://github.com/OPM/opm-simulators/pull/5824)). Previously, zeroes would be output for some connections.
- Several issues related to zero group net voidage rates (NETV) have been fixed ([#5822](https://github.com/OPM/opm-simulators/pull/5822)).
- The current well state is now used as an initial guess for the implicit well potential calculations ([#5929](https://github.com/OPM/opm-simulators/pull/5929)). This avoids some occasional convergence issues.
- Fixed the output of residuals for gas-water systems, previously the residuals were only output for the gas phase ([#6002](https://github.com/OPM/opm-simulators/pull/6002)).
- Summary output has been corrected for cumulative quantities (MSUMLINS and MSUMNEWT) at time steps between report steps ([#6039](https://github.com/OPM/opm-simulators/pull/6039)). Previously, this was only output correctly at report steps.
- Corrected application of THP constraints for wells subject to group constraints ([#6018](https://github.com/OPM/opm-simulators/pull/6018)).
- Fixed VTK output for dry runs (requested using --enable-vtk-output=true and --enable-dry-run=true) in cases without the oil phase present ([#6057](https://github.com/OPM/opm-simulators/pull/6057)) or with the DRSDT keyword in the SCHEDULE section ([#6060](https://github.com/OPM/opm-simulators/pull/6060)).
- Fixed problem with group pressure maintenance targets and controls applied at the Field level using the GPMAINT keyword ([#5916](https://github.com/OPM/opm-simulators/pull/5916)).
- A potential division by zero has been prevented that could occur when using the CPRW linear solver and a well has no active perforations on some rank ([#6053](https://github.com/OPM/opm-simulators/pull/6053)).
- Previously, a small non-zero oil saturation could be generated when the model was initialised by enumerated gas and water saturations due to the numerical precision. This caused an issue when setting up the primary variables. This issue has been resolved by setting initial oil saturations less than a small tolerance value to zero ([#6079](https://github.com/OPM/opm-simulators/pull/6079)).
- The list of well events is now cleared out once they have been processed at the start of a report step preventing them from potentially being executed again ([#6078](https://github.com/OPM/opm-simulators/pull/6078)).
- The simulator now re-iterates the network solver if the well group controls are changed, the gaslift quantity is changed, or the inner iterations did not converge ([#6062](https://github.com/OPM/opm-simulators/pull/6062)).
- The MICP model now works with the default linear solver settings ([#6109](https://github.com/OPM/opm-simulators/pull/6109)).
- Groups subject to higher level controls are now allowed to switch to the FIELD group controls ([#6135](https://github.com/OPM/opm-simulators/pull/6135)).
- Fixed an issue where the wrong phase index could be used for restart output if output was only requested for one phase ([#6134](https://github.com/OPM/opm-simulators/pull/6134)).
- The correct ALQ dimension is now restored when restarting a run ([#4271](https://github.com/OPM/opm-common/pull/4271)).
- Specifying UPDATE NEXT
\end{verbatim}

in an ACTIONX block now causes the UDQ to be updated on the next time
step rather than the next report step
(\href{https://github.com/OPM/opm-common/pull/4325}{\#4325}). - A few
issues with the WELTRAJ and COMPTRAJ keywords have been fixed
(\href{https://github.com/OPM/opm-common/pull/4279}{\#4279}). Any
MAPAXES specifications were not handled, and cases with the origin not
at zero failed. To solve this, the well coordinates are now transformed
when they are read. There was a bug in COMPTRAJ where defaulting the
saturation table column was not handled, leading to a segmentation
fault. The code for converting the trajectory to cell indices only took
the top and bottom coordinates into account, effectively assuming linear
trajectories. Now intermediate points are added, allowing arbitrary
trajectories. - All group constraints are now written to the restart
file even if the group is controlled by a higher-level group
(\href{https://github.com/OPM/opm-common/pull/4346}{\#4346}).
Previously, production groups controlled by a higher-level group could
be given the wrong group control following a restart. - The minimum grid
block pore volume thresholds for individual cells specified using the
MINPVV keyword will now take precedence over the threshold set of all
cells using the MINPV or MINPORV keyword
(\href{https://github.com/OPM/opm-common/pull/4360}{\#4360}).
Previously, MINPVV would be ignored if either MINPV or MINPORV was used.
- The simulator can now handle deck files with sub-second report steps
with a resolution of 1ms
(\href{https://github.com/OPM/opm-common/pull/4374}{\#4374}). - A
potential divide by zero error at report step zero has been avoided
(\href{https://github.com/OPM/opm-common/pull/4392}{\#4392}). - The
GRUPTARG, GSATINJE, GSATPROD, GRUPNET and TEST keywords have been
removed from the list of supported keywords in ACTIONX blocks since they
are not supported by OPM Flow
(\href{https://github.com/OPM/opm-common/pull/4400}{\#4400},
\href{https://github.com/OPM/opm-common/pull/4436}{\#4436}). Note that
TEST is not a keyword. - Previously, there was a risk that the defining
expression for a UDQ could overflow the 128 character limit when stored
in the restart file
(\href{https://github.com/OPM/opm-common/pull/4422}{\#4422}). The
simulator now stores a ``fused'' version of the defining expression that
honours the character limit. - The code now ensures that any events that
might happen during the first time step after the restart time are
preserved (\href{https://github.com/OPM/opm-common/pull/4417}{\#4417}).
- Fixed the handling of the group (GRUP) control mode for wells in
auto-choke groups in extended network models
(\href{https://github.com/OPM/opm-common/pull/4435}{\#4435}). - Fixed a
bug where the simulator tried to restore UDQ unit strings as UDQ
variables after a restart
(\href{https://github.com/OPM/opm-common/pull/4457}{\#4457}). - The code
now uses the branch number to correctly determine if a segment of a
multi-segment well is active at restart
(\href{https://github.com/OPM/opm-common/pull/4483}{\#4483}). - When a
keyword is expected the simulator will now ignore any line where the
first non-whitespace character is a slash
(\href{https://github.com/OPM/opm-common/pull/4462}{\#4462}).
Previously, if the slash was followed by any non-whitespace characters
on the same line then an exception would have been raised. - The
simulator now updates the well efficiency factor immediately if the
WEFAC keyword is run from an ACTIONX block during a report step rather
than at the end of the report step
(\href{https://github.com/OPM/opm-common/pull/4495}{\#4495}). - The
handling of daylight-saving time when importing the start date from a
summary file (or enhanced summary file) using Python bindings has been
fixed (\href{https://github.com/OPM/opm-common/pull/4509}{\#4509}).

Known issues in this release of the simulator include:

\begin{verbatim}
- There is an error with the RSM header for summary vectors whose NUMS entry in the SMSPEC file is derived from more than a single number source (e.g., single region or segment ID). This applies to all block vectors (BGPV, BOPV, BWPV, etc.), connection level quantities (COPT, etc.), and inter-region flows such as ROFT etc ([#3078](https://github.com/OPM/opm-common/issues/3078)). The work around is to plot the data in OPM ResInsight and right-click on the plot to view and copy the data.
- As per previous releases of the radial model, the COORDSYS keyword item three must be set to COMP to complete the circle, this has not been implemented in this release. Also there appears to be a bug for full radial models when a well goes on BHP control that causes the well not to respect the BHP constraint, this eventually causes the well to die prematurely. See [#2640](https://github.com/OPM/opm-common/pull/2640) for a discussion on the topic.
- As in previous releases there are some issues with the OPERATE and OPERATER keywords associated with the input parsing; for various reasons a few of the fields require special case treatment in the grid processing, including (at least) MULTZ, PORV and ACTNUM, and for those keywords the OPERATE/OPERATER keyword does not work. The work around is to use the MULTIPLY keyword instead.
- For the UDQ ASSIGN operator after the terminating “/” normally any comments can be entered; however, if there is “/” within the comment field, as per:
\end{verbatim}

\begin{verbatim}
ASSIGN FUNGLYLD 1.196   /    Condensate Yield (63.5/56.7)/(1.0 – 0.065)
\end{verbatim}

then the simulator will abort. The work around is to manually place the
comment characters ``--'' after the ASSIGN terminating ``/'', like so:

\begin{verbatim}
ASSIGN FUNGLYLD 1.196   / -- Condensate Yield (63.5/56.7)/(1.0 – 0.06)
\end{verbatim}

\begin{verbatim}
- At the moment, one cannot initialize tracers using the EQUALS keyword. Instead use the array format, that is the keyword followed by the required number of values, or the TVDP keyword in the SOLUTION section to set the initial tracer concentrations as a function of depth.
- Currently, gas tracers cannot be used if the dissolved gas phase, as per the DISGAS keyword in the RUNSPEC section, is active in the model.
- The summary vector RTIPTHEA, that defines the energy in-place between the initial and the current time for regions, is not supported unlike the FTIPTHEA and BTIPTHEA vectors. Secondly, the error message:
\end{verbatim}

\begin{longtable}[]{@{}
  >{\raggedright\arraybackslash}p{(\linewidth - 0\tabcolsep) * \real{1.0000}}@{}}
\toprule\noalign{}
\begin{minipage}[b]{\linewidth}\raggedright
Warning: Problem with summary keyword RTIPTHEA In RSM-THERMAL.data line
492 FIP region FIPHEA not defined in REGIONS section - RTIPTHEA ignored
\end{minipage} \\
\midrule\noalign{}
\endhead
\bottomrule\noalign{}
\endlastfoot
\end{longtable}

is incorrect, as the message indicates that it is being treated like a
named region, as per the FIP keyword, when it is actually a SUMMARY
vector
(\href{https://github.com/OPM/opm-simulators/issues/3870}{\#3870}).

\begin{verbatim}
- If there are cells that are very distorted, which can occur near fault planes, then the simulator may abort because it cannot calculate the pore volume of such cells. The work around is to re-generate the grid in the static model, taking care that the cells around the fault planes are more or less orthogonal ([#2992](https://github.com/OPM/opm-common/issues/2992) and [#3770](https://github.com/OPM/opm-simulators/issues/3770)).
- Currently the OPERATER keyword in the EDIT section does not work with the DEPTH, TRANX, TRANY and TRANZ property arrays ([#2994](https://github.com/OPM/opm-common/issues/2994) and [#748](https://github.com/OPM/opm-tests/pull/748)).
- If a standard well is fully declared in an ACTIONX block which is then activated at a later date, and later the well is modified to be a multi-segment well using the WELSEGS and COMPSEGS keywords, then this will cause the simulator to abort with an assert failure. The solution to this issue is to not use this type of work flow in declaring wells ([#2891](https://github.com/OPM/opm-common/issues/2891) and [#2895](https://github.com/OPM/opm-common/pull/2895)).
- Although the ACTIONX EXIT command works as expected, it does not write out the requested RSM file at the end of the run. However, the other SUMMARY and RESTART files are written out ([#2877](https://github.com/OPM/opm-common/issues/2877)).
- Although the GCONSUMP keyword in the SCHEDULE section is fully implemented as documented, it is not possible to verify the output as the associated SUMMARY vectors are not written out, that is the SUMMARY sales gas vectors FGSR, FGST, GGSR and GGST, and fuel vectors FGCR, FGCT, GGCR, and GGCT have not been implemented ([#2679](https://github.com/OPM/opm-common/issues/2679)).
- There are small differences in the behavior of the NEXTSTEP keyword in the RUNSPEC section between OPM Flow and the commercial simulator that remain unresolved ([#3745](https://github.com/OPM/opm-simulators/issues/3745)).
- There is a unit handling issue associated with OPERATE keyword. If the OPERATE(X) parameter has units, as for example PERMX, then the conversion is always done in SI units, despite the input deck declaring the deck to be fields units, as per FIELD keyword in the RUNSPEC section. Note that OPM Flow performs all of its calculations internally in SI and performs unit handling only when inputting the *.DATA file and when outputting result files. Thanks to [lrijkels](https://github.com/lrijkels) for reporting the issue. See [#4597](https://github.com/OPM/opm-simulators/issues/4597) for details.
- There is an issue associated with restarting from a restart file with the solution gas (Rs) maximum rate of increase, as defined by the DRSDT keyword in the SCHEDULE, that has been set to zero. This is because, the simulator does not save/restore this setting in simulator's restart files, which means that simulator misses the essential value zero upon restarting the case. As a work-around one can use the option --sched-restart=true, when running the restart case. This will initialize the restarted simulation based on information from the complete SCHEDULE section, instead of just the parts that we're going to simulate and the rest from the restart file. Thanks to [goncalvesmachadoc](https://github.com/goncalvesmachadoc) for reporting the issue. See [#4272](https://github.com/OPM/opm-simulators/issues/4272) for details.
- OPM cannot be built with dune-fem version 2.9 or later ([#4934](https://github.com/OPM/opm-simulators/issues/4934)). Please use a previous version.
- The simulator uses an irregular corner-point grid geometry with adjusted pore volumes to represent radial grids so it is not possible to create a full ring (360 degree disk) with only one cell in the theta direction (NY=1). The work around is to model a slice (say DTHETA=60 degrees). Note that as the angle increases larger pore volume adjustments are required ([#4755](https://github.com/OPM/opm-simulators/issues/4755)).
- In some cases with the network option the simulator can wrongly report that a well has no THP constraints, for example
\end{verbatim}

\begin{longtable}[]{@{}l@{}}
\toprule\noalign{}
GLIFT WTEST: Well S-3H does not have THP constraints \\
\midrule\noalign{}
\endhead
\bottomrule\noalign{}
\endlastfoot
\end{longtable}

when THP constraints have been defined
(\href{https://github.com/OPM/opm-simulators/issues/4887}{\#4887}).

\begin{verbatim}
- Dispersion in the gas phase leads to convergence issues for the 11th SPE CSP ([https://spe.org/csp](https://spe.org/csp)) models Version 11B and Version 11C (and is not fully tested). Therefore dispersion in the gas phase has been temporarily removed and a warning added ([#5101](https://github.com/OPM/opm-simulators/pull/5101), [#859](https://github.com/OPM/opm-models/pull/859)).
\end{verbatim}

The 2024-10 release consists of some new features and various
improvements and bug fixes. These include the following highlights:

\begin{verbatim}
- Linear solver default parameters have been updated to improve accuracy and performance including.
- Add summary output vectors for stranded CO2 in the gas phase (GMST) and maximum potential residual trapping (GMTR).
- Added PSI and OMEGA parameters to DRSDTCON to control the regime and rate of convective mixing. Now supports GAS/WATER systems in addition to OIL/GAS systems.
- Support has been added for MPI communication being performed directly between GPUs instead of being sent via the CPU when using a multi-gpu linear solver.
- Added support for WTMULT to multiply a well target or constraint by a constant.
- Simulator is now able to make better usage of the Damaris middleware if enabled (https://project.inria.fr/damaris/).
- Extended the tracer model to include partitioning of tracers into free and solution states whenever dissolved gas/vaporized oil exist.
- Added support for automatic chokes (NODEPROP item 3 CHOKE) in extended network models.
- Added support for Killough's hysteresis model for both wetting and non-wetting phase hysteresis in a water wet system (EHYSTR item 2 HYSTMOD = 4).
- Added the option to specify the linear solver for subdomains when using the NLDD nonlinear solver.
- Added foundations for future experiments using the four-byte float type instead of the eight-byte double type as the primary array element in order to save memory.
- Support added for RFIP and SFIP mnemonics for the RPTRST keyword.
- Support added for mnemonic WELLS=N for the RPTSCHED keyword.
- Previously the THP was set to zero if a well was stopped or set to a zero rate (well or group). The THP will now be calculated provided a VFP table is active.
- Corrected transmissibilities calculated for pinched out cells when using PINCH option 4 equals ALL.
- Support added for the EQUALREG keyword.
- Added support for using METIS in addition to Zoltan to partition corner point grids.
- Proper INIT file table output for the saturation functions in LET format.
- Add support for outputting MULTPV to the Restart file.
- Add support for additional keywords in PYACTION blocks.
- Added support for [DENAQA](#REF_HEADING_KEYWORD_DENAQA_8_3) and [VISCAQA](#REF_HEADING_KEYWORD_VISCAQA_8_3) keywords when using CO2STORE. Also added support for [SALTMF](#REF_HEADING_KEYWORD_SALTMF_8_3) keyword.
- Support has been expanded for saving/restoring scaling factors for segment level devices (SICDs, AICDs and Valves) to/from the Restart file.
- Added support for output of well level target or constraint summary vectors.
- Added support for field, group and well gas mass injection rate and injection total summary vectors for use with CO2STORE.
- Proper handling of the WHISTCTL in restarted simulations.
- Expanded consistency checks for array operations such as ADD, COPY, and MULTIPLY, as well as the per-region counterparts.
- Performance improvements in well name matching on platforms which support the Posix fnmatch() function.
- Initial and experimental support for compositional simulations–currently supporting 2 to 7 components.
- Fixed two-phase gas/water and gas/oil model initialisation (equilibration).
\end{verbatim}

The major command line changes made for this release are summarized in
Table A.1

\begin{longtable}[]{@{}
  >{\raggedright\arraybackslash}p{(\linewidth - 6\tabcolsep) * \real{0.2500}}
  >{\raggedright\arraybackslash}p{(\linewidth - 6\tabcolsep) * \real{0.2500}}
  >{\raggedright\arraybackslash}p{(\linewidth - 6\tabcolsep) * \real{0.2500}}
  >{\raggedright\arraybackslash}p{(\linewidth - 6\tabcolsep) * \real{0.2500}}@{}}
\toprule\noalign{}
\begin{minipage}[b]{\linewidth}\raggedright
OPM Flow 2024-10 New and Deprecated Command Line Options
\end{minipage} & \begin{minipage}[b]{\linewidth}\raggedright
\end{minipage} & \begin{minipage}[b]{\linewidth}\raggedright
\end{minipage} & \begin{minipage}[b]{\linewidth}\raggedright
\end{minipage} \\
\midrule\noalign{}
\endhead
\bottomrule\noalign{}
\endlastfoot
No. & Variable Name & Description & Default \\
1 & --action-parsing-strictness & Added a new command line parameter
that specifies the parsing strictness for ACTIONX and PYACTION blocks
(\href{https://github.com/OPM/opm-simulators/pull/5547}{\#5547},
\href{https://github.com/OPM/opm-common/pull/4167}{\#4167}). A defined
character string. Set strictness of parsing process for ActionX and
PyAction. Available options are normal (do not apply keywords that have
not been tested for ActionX or PyAction) and low (try to apply all
keywords, beware: the simulation outcome might be incorrect). &
``normal'' \\
2 & --check-group-constraints-inner-well-iterations & A boolean value.
Allow checking of group constraints during inner well iterations. &
true \\
3 & --convergence-monitoring & A new convergence monitoring feature
(\href{https://github.com/OPM/opm-simulators/pull/5590}{\#5590}) has
been introduced to improve the robustness and efficiency of the
simulator based on Moyner et al., 2024
(https://doi.org/10.3997/2214-4609.202437057). A boolean value. Enable
convergence monitoring. & false \\
4 & --convergence-monitoring-cut-off & An integer value. Cut off limit
for convergence monitoring. & 6 \\
5 & --convergence-monitoring-decay-factor & A real value. Decay factor
for convergence monitoring. & 0.75 \\
6 & --enable-drift-compensation & Linear solver default parameters have
been updated to improve accuracy and performance including: material
balance tolerance reduced to 1.0e-7, local well solve control switching
is turned on, use implicit IPR is turned on, and enable drift
compensation is turned off
(\href{https://github.com/OPM/opm-simulators/pull/5157}{\#5157}). A
boolean value. Enable partial compensation of systematic mass losses via
the source term of the next time step. & false \\
7 & --imbalance-tol & A real value. Tolerable imbalance of the load
balancing. & 1.1 \\
8 & --inj-mult-damp-mult & Added damping to the injectivity multipliers
defined by the WINJMULT keyword
(\href{https://github.com/OPM/opm-simulators/pull/5631}{\#5631}). A real
value. Injection multiplier dampening factor (dampening multiplied by
this each time oscillation is detected). & 0.9 \\
9 & --inj-mult-min-damp-factor & A real value. Minimum injection
multiplier dampening factor (maximum dampening level). & 0.05 \\
10 & --inj-mult-osc-threshold & A real value. Injection multiplier
oscillation threshold (used for multiplier dampening). & 0.1 \\
11 & --input-skip-mode & Added a command line argument to control the
way OPM Flow treats SKIP100 and SKIP300 keywords
(\href{https://github.com/OPM/opm-simulators/pull/5329}{\#5329}). By
default OPM Flow behaves like the commercial black-oil simulator. A
defined character string. Set compatibility mode for the SKIP100/SKIP300
keywords. Options are: 100 (skip SKIP100..ENDSKIP, keep
SKIP300..ENDSKIP), 300 (skip SKIP300..ENDSKIP, keep SKIP100..ENDSKIP) or
all (skip both SKIP100..ENDSKIP and SKIP300..ENDSKIP). & ``100'' \\
12 & --linear-solver & A defined character string. Configuration of
solver. Valid options are: cprw (default), ilu0, dilu, cpr (an alias for
cprw), cpr\_quasiimpes, cpr\_trueimpes, cpr\_trueimpesanalytic, amg or
hybrid (experimental). Alternatively, you can request a configuration to
be read from a JSON file by giving the filename here, ending with
`.json.'. & ``cprw'' \\
13 & --local-well-solve-control-switching & A boolean value. Allow
control switching during local well solutions. & true \\
14 & --metis-params & A character string. Configuration of Metis
partitioner. You can request a configuration to be read from a JSON file
by giving the filename here, ending with `.json.' See
http://glaros.dtc.umn.edu/gkhome/fetch/sw/metis/manual.pdf for available
METIS options. & ``default'' \\
15 & --min-strict-cnv-iter & An integer value. Minimum number of Newton
iterations before relaxed tolerances can be used for the CNV convergence
criterion. & -1 \\
16 & --nldd-local-linear-solver & Added the option to specify the linear
solver for subdomains when using the NLDD nonlinear solver
(\href{https://github.com/OPM/opm-simulators/pull/5553}{\#5553}).
Previously, the simulator was hardcoded to use the ILU0 solver with the
default maximum number of iterations and default reduction tolerance. A
defined character string. Configuration of NLDD local linear solver.
Valid options are: ilu0 (default), dilu, cpr\_quasiimpes and amg.
Alternatively, you can request a configuration to be read from a JSON
file by giving the filename here, ending with `.json.'. & ``ilu0'' \\
17 & --nldd-local-linear-solver-max-iter & An integer value. The maximum
number of iterations of the NLDD local linear solver. & 200 \\
18 & --nldd-local-linear-solver-reduction & A real value. The minimum
reduction of the residual which the NLDD local linear solver must
achieve. & 0.01 \\
19 & --partition-method & Added support for using METIS in addition to
Zoltan to partition corner point grids
(\href{https://github.com/OPM/opm-grid/pull/725}{\#725}). An integer
value. Choose partitioning strategy: 0=simple, 1=Zoltan, 2=METIS. & 1 \\
20 & --tolerance-cnv-energy & Added separate energy convergence
tolerance and energy balance tolerance
(\href{https://github.com/OPM/opm-simulators/pull/5321}{\#5321}).
Previously the tolerances for mass were used. A real value. Local energy
convergence tolerance (maximum of local energy errors). & 0.01 \\
21 & --tolerance-cnv-energy-relaxed & A real value. Relaxed local energy
convergence tolerance that applies for iterations after the iterations
with the strict tolerance. & 1.0 \\
22 & --tolerance-energy-balance & A real value. Tolerated energy balance
error relative to (scaled) total energy present. & 1e-07 \\
23 & --tolerance-energy-balance-relaxed & A real value. Relaxed
tolerated energy balance error that applies for iterations after the
iterations with the strict tolerance. & 1e-06 \\
24 & --tolerance-mb & A real value. Tolerated mass balance error
relative to total mass present. & 1e-07 \\
25 & --use-implicit-ipr & A boolean value. Compute implict IPR for
stability checks and stable solution search. & true \\
26 & --print-properties & Support for this command line option has been
removed
(\href{https://github.com/OPM/opm-simulators/pull/5514}{\#5514}). An
integer value. Print the values of the compile time properties at the
start of the simulation. & 2 \\
27 & --zoltan-imbalance-tol & A real value. Tolerable imbalance of the
loadbalancing provided by Zoltan. This command line option has been
deprecated; --imbalance-tol should be used instead. & 1.1 \\
Notes: & & & \\
\end{longtable}

\emph{Table A.1: OPM Flow 2024-10 New and Deprecated Command Line
Options}

In addition to the changes to the command line options the following new
features have been added to the simulator:

\begin{verbatim}
- Add summary output vectors for stranded CO2 in the gas phase (GMST) and maximum potential residual trapping (GMTR) ([#5344](https://github.com/OPM/opm-simulators/pull/5344), https://docs.opengosim.com/theory/residual_and_stranded/).
- Added support for WTMULT to multiply a well target or constraint by a constant ([#5411](https://github.com/OPM/opm-simulators/pull/5411)).
- Simulator is now able to make better usage of the Damaris middleware if enabled (https://project.inria.fr/damaris/). Added the ability to pass multiple variables to Damaris using the DamarisWriter class ([#5352](https://github.com/OPM/opm-simulators/pull/5352)). A command line option has also been added to limit the variables to be passed using --damaris-limit-variables=<CSV list>.
- Restart files can now be output every time step if the command line option --enable-write-all-solutions=true ([#5426](https://github.com/OPM/opm-simulators/pull/5426)). The numbering of the restart and summary files will follow the time step index and not the report index. This will not be compatible with restart, but it is sometimes useful to investigate the results at every time step. Note that this doesn't add support for RPTRST BASIC=6.
- Added support for automatic chokes (NODEPROP item 3 CHOKE) in extended network models ([#4935](https://github.com/OPM/opm-simulators/pull/4935)).
- Added support for using CPR on AMD GPU using the rocSPARSE library ([#5408](https://github.com/OPM/opm-simulators/pull/5408)).
- Added support for outputting the RSWSAT and RVWSAT arrays to the Restart file ([#5483](https://github.com/OPM/opm-simulators/pull/5483)).
- Added PSI and OMEGA parameters to DRSDTCON to control the regime and rate of convective mixing ([#5491](https://github.com/OPM/opm-simulators/pull/5491)). Now supports GAS/WATER systems in addition to OIL/GAS systems.
- Added an OPM implementation of the ILU0 preconditioner in CUDA similar to that implemented in CuDILU ([#5441](https://github.com/OPM/opm-simulators/pull/5441)).
- Added support for Killough's hysteresis model for both wetting and non-wetting phase hysteresis in a water wet system (EHYSTR item 2 HYSTMOD = 4) ([#5273](https://github.com/OPM/opm-simulators/pull/5273)).
- Extended implementation of DRSDTCON to apply to multiple PVT regions ([#5527](https://github.com/OPM/opm-simulators/pull/5527)). Also added support for deactivation of individual DRSDT, DRVDT and DRSDTCON regions.
- Added foundations for future experiments using the four-byte float type instead of the eight-byte double type as the primary array element in order to save memory ([#5560](https://github.com/OPM/opm-simulators/pull/5560)).
- Support added for RFIP and SFIP mnemonics for the RPTRST keyword to request output of fluid in place volumes at surface and reservoir conditions to the Restart file ([#5621](https://github.com/OPM/opm-simulators/pull/5621)). Note that FIP is an alias for SFIP.
- Support added for mnemonic WELLS=N for the RPTSCHED keyword to request output of well reports to the Print file ([#5626](https://github.com/OPM/opm-simulators/pull/5626)).
- Support added for the EQUALREG keyword to specify property arrays by region ([#5648](https://github.com/OPM/opm-simulators/pull/5648)).
- Added a failure flag to the TaskletRunner allowing errors in a tasklet to be caught and logged ([#909](https://github.com/OPM/opm-models/pull/909)).
- Add support for outputting MULTPV to the Restart file if it is specified in the GRID or EDIT section ([#4021](https://github.com/OPM/opm-common/pull/4021)).
- Add support for WELSPECS, COMPSEGS, WELSEGS, WSEGVALV, WECON, WTEST, WGRUPCON, GCONSUMP, GRUPTREE and EXIT keywords in PYACTION blocks ([#4059](https://github.com/OPM/opm-common/pull/4059), [#4061](https://github.com/OPM/opm-common/pull/4061), [#4069](https://github.com/OPM/opm-common/pull/4069), [#4070](https://github.com/OPM/opm-common/pull/4070), [#4071](https://github.com/OPM/opm-common/pull/4071)).
- Added support for CNAMES to specify the component names ([#5405](https://github.com/OPM/opm-simulators/pull/5405), [#4091](https://github.com/OPM/opm-common/pull/4091)).
- Added support for [DENAQA](#REF_HEADING_KEYWORD_DENAQA_8_3) and [VISCAQA](#REF_HEADING_KEYWORD_VISCAQA_8_3) keywords to specify parameters for calculating brine density and viscosity using Ezrokhi's method when using CO2STORE ([#4091](https://github.com/OPM/opm-common/pull/4091)). Also added support for [SALTMF](#REF_HEADING_KEYWORD_SALTMF_8_3) to allow the initial salt concentration to be input as a mole fraction.
- Extended the tracer model to include partitioning of tracers into free and solution states whenever dissolved gas/vaporized oil exist ([#5268](https://github.com/OPM/opm-simulators/pull/5268), [#4001](https://github.com/OPM/opm-common/pull/4001)). Essentially the tracer linear system has been extended with equations for free and solution tracers with simple mass transfer coupling between them. Both free and solution tracers are now output to the Restart file.
- Added support for output of well level target or constraint summary vectors: WBHPT, W[OGLVW]PRT, and W[OGVW]IRT where O, G, L, V and W represent oil, gas, liquid, reservoir volume and water respectively ([#4093](https://github.com/OPM/opm-common/pull/4093)).
- Initial support for enumerated compositional initialisation ([#4110](https://github.com/OPM/opm-common/pull/4110)).
- Added support for [FGW]GMI[RT] field, group and well gas mass injection rate and injection total summary vectors for use with CO2STORE ([#4162](https://github.com/OPM/opm-common/pull/4162)).
- A mixing model has been added to the thermal black-oil simulator ([#4014](https://github.com/OPM/opm-common/pull/4014)). The internal energy of each phase now also includes the energy for the dissolved component. The meaning of the TEMP keyword has been changed. In the Thermal model (THERMAL) enthalpy equals internal energy plus work. Whereas, the Temperature model (TEMP) now assumes enthalpy equals internal energy.
- Added support for the COPY, OPERATE, ADDREG, COPYREG, EQUALREG, and MULTIREG keywords for data with global storage ([#4236](https://github.com/OPM/opm-common/pull/4236), [#4237](https://github.com/OPM/opm-common/pull/4237), [#4238](https://github.com/OPM/opm-common/pull/4238)).
- Extended API
\end{verbatim}

to allow a source term to be added to a cell
(\href{https://github.com/OPM/opm-common/pull/4241}{\#4241}). In
particular this allows a source term to be generated by Python code
without needing to pass in a DeckRecord. - Added parsing support for ZMF
keyword (\href{https://github.com/OPM/opm-common/pull/4200}{\#4200}). -
The convertECL utility has been updated to support conversion of LGR
output between the unformatted (.LGR) and formatted (.FLGR) file types
(\href{https://github.com/OPM/opm-common/pull/4249}{\#4249}).

Improvements to simulator in this release include the following:

\begin{verbatim}
- Support has been added for MPI communication being performed directly between GPUs instead of being sent via the CPU when using a multi-gnu linear solver ([#5145](https://github.com/OPM/opm-simulators/pull/5145)).
- Python bindings can now control whether MPI_Init() and MPI_Finalize() will be called when creating an Opm::Main object ([#5325](https://github.com/OPM/opm-simulators/pull/5325)). This allows the Python unit test files to run more than a single unit test.
- The fraction of pore-volume whose CNV targets are violated as a new per-iteration quantity is now included in the INFOITER file (--output-extra-convergence-info=iteration), with the column header CnvErrPvFrac ([#5302](https://github.com/OPM/opm-simulators/pull/5302)).
- Support has been added for using hipify on the CUDA code to allow for GPU support on both AMD and Nvidia architectures ([#5253](https://github.com/OPM/opm-simulators/pull/5253)).
- The simulator now checks if relaxed tolerances are tighter than the strict tolerances and if so adjusts the relaxed tolerances ([#5337](https://github.com/OPM/opm-simulators/pull/5337)).
- The code now resets the group pressure maintenance state to avoid oscillations when the target pressure is reached and the pressure maintenance rate is set to zero ([#5350](https://github.com/OPM/opm-simulators/pull/5350)).
- Optimized the memory operations time by overlapping the creation (which includes a host memory copy) of the Jacobian matrix used in the block Jacobi ILU with copying the matrix data to the GPU ([#5256](https://github.com/OPM/opm-simulators/pull/5256)).
- SummaryState Objects are now Aware of Undefined UDQ Values ([#5334](https://github.com/OPM/opm-simulators/pull/5334)).
- Previously the code would only start handling oscillations when two phases had oscillating residuals. The code will now start handling oscillations when just one phase is oscillating ([#5398](https://github.com/OPM/opm-simulators/pull/5398)).
- Added Newton update damping for the situation where the solution for a domain oscillates ([#5342](https://github.com/OPM/opm-simulators/pull/5342)).
- Reduced execution time for parallel AMG/CPR simulations. Allowed use of AMG with command line option --matrix-add-well-contributions=false ([#5182](https://github.com/OPM/opm-simulators/pull/5182)).
- The code has been updated to only issue warnings the first time the transmissibilities are calculated thus removing spurious and duplicate warnings that can occur when using load balancing ([#5351](https://github.com/OPM/opm-simulators/pull/5351)).
- Nearly all exceptions thrown when computing wells will not abort the simulator but result in a timestep chop. Hence those should not be counted as errors and be reported in the PRT file. The simulator now only reports problems for exceptions thrown when updating the well potentials ([#5423](https://github.com/OPM/opm-simulators/pull/5423)).
- The well economic (WECON) minimum rates are no longer applied if the well is subject to a zero group rate target or constraint ([#5424](https://github.com/OPM/opm-simulators/pull/5424)).
- Directly output the SGMAX, SHMAX, SOMAX, SOMIN and SWHY1 arrays rather than the KRNSW_GO, KRNSW_OW, PCSWM_GO, PCSWM_OW and SOMAX arrays to the Restart file for cases with relative permeability and/or capillary pressure hysteresis ([#5402](https://github.com/OPM/opm-simulators/pull/5402)).
- Sphinx documentation ([https://opm.github.io/opm-python-documentation/master/index.html](https://opm.github.io/opm-python-documentation/master/index.html)) has been added for the opm.simulators.BlackOilSimulator Python module ([#5243](https://github.com/OPM/opm-simulators/pull/5243)).
- The simulator now more dynamically picks the thread block size when using CUDA giving up to a 50% speed up on linear iterations ([#5433](https://github.com/OPM/opm-simulators/pull/5433)).
- If OpenMP is available then the simulator will now default to two threads per process unless the OMP_NUM_THREADS environment variable is set or the command line option --threads-per-process is set. A warning will be issued if both the environment variable and the command line  option are set, but the environment variable will take precedence ([#5454](https://github.com/OPM/opm-simulators/pull/5454)).
- Visualization Toolkit (VTK) files can now be output for dry runs by using the --enable-vtk-output=1 --enable-dry-run=1 command line options ([#5499](https://github.com/OPM/opm-simulators/pull/5499)).
- Linear solver default parameters have been updated to improve accuracy and performance including: material balance tolerance reduced to 1.0e-7, local well solve control switching is turned on, use implicit IPR is turned on, and enable drift compensation is turned off ([#5157](https://github.com/OPM/opm-simulators/pull/5157)).
- A more helpful message is now output if the PETGRID keyword is used to try to load a generic simulation grid (*.GSG) file ([#5541](https://github.com/OPM/opm-simulators/pull/5541)): "OPM Flow cannot read GSG format grid input, re-export case from Petrel using GRDECL instead."
- The check for active wells as a condition for drift compensation has been removed ([#5110](https://github.com/OPM/opm-simulators/pull/5110)). For CO2 injection simulation the drift compensation is still beneficial even when wells are no longer active.
- The command line option --cpr-reuse-setup now only applies to the CPR preconditioner ([#5445](https://github.com/OPM/opm-simulators/pull/5445)).
- The simulator now uses a mobility weighted average of cell level densities for the connection level mixture densities in no-flow producing wells ([#5479](https://github.com/OPM/opm-simulators/pull/5479)). This mobility weighted average approach gives a more monotone pressure build up.
- Output cell counts and pore volume fractions to INFOITER file  for cells that: (1) satisfy the strict convergence tolerance, (2) satisfy the relaxed but not the strict convergence tolerance, (3) do not satisfy the relaxed convergence tolerance ([#5338](https://github.com/OPM/opm-simulators/pull/5338)).
- Consistency checks have been added for oil, gas and water phase saturation functions ([#5570](https://github.com/OPM/opm-simulators/pull/5570), [#5571](https://github.com/OPM/opm-simulators/pull/5571), [#5572](https://github.com/OPM/opm-simulators/pull/5572)).
- The simulator now avoids computing the default number of domains based on a specific target number of cells per domain if the user choses a specific number of NLDD domains ([#5579](https://github.com/OPM/opm-simulators/pull/5579)).
- The simulator now avoids doing unnecessary updates in the numerical aquifer calculations ([#5632](https://github.com/OPM/opm-simulators/pull/5632)).
- Previously the THP was set to zero if a well was stopped or set to a zero rate (well or group). The THP will now be calculated provided a VFP table is active ([#5634](https://github.com/OPM/opm-simulators/pull/5634)).
- Consistency checks have been added for three-point saturation scaling SCALECRS=YES ([#5573](https://github.com/OPM/opm-simulators/pull/5573), [#5640](https://github.com/OPM/opm-simulators/pull/5640)). These checks collectively guarantee a mobile displacing oil saturation in the two phase gas-oil and oil-water systems.
- The dependencies on the OPM PyPi packages have been updated to add support for newer versions of Python, versions 3.6 to 3.12 are now supported ([#5644](https://github.com/OPM/opm-simulators/pull/5644)).
- A new strategy for preventing wells from oscillating between being stopped and revived has been introduced ([#5659](https://github.com/OPM/opm-simulators/pull/5659)). The cycle of oscillations are broken by requiring that for a well to remain stopped after it has been stopped in the local solve, that it should have had zero rate when it entered the local solve, meaning that both the global and local solve agree that the well should be stopped.
- It is no longer necessary to use add well contributions to matrix for NLDD. Also, the CPRW preconditioner is now supported for NLDD ([#5341](https://github.com/OPM/opm-simulators/pull/5341), [#5676](https://github.com/OPM/opm-simulators/pull/5676)).
- The number of pinch-out connections generated will now be reported to the print file ([#767](https://github.com/OPM/opm-grid/pull/767)).
- An error message will now be generated if the COMPSEGS keyword is used before the multi-segment well is defined using the WELSEGS keyword ([#4031](https://github.com/OPM/opm-common/pull/4031)).
- The reservoir volume rate (RESV) constraint for a history matching producer is evaluated by converting the observed oil, water and gas surface rates to the current reservoir conditions hence its value is not fixed. However, to test if the well is subject to a zero reservoir volume rate constraint it is sufficient to test if all of the observed rates are zero ([#4039](https://github.com/OPM/opm-common/pull/4039)).
- The path of the Restart file (either absolute or relative to the .DATA file) as specified by the RESTART keyword is now saved for writing to the SMSPEC file ([#4056](https://github.com/OPM/opm-common/pull/4056)).
- Restart file paths of between 73 and 132 characters in length are now also written to the SMSPEC file ([#4057](https://github.com/OPM/opm-common/pull/4057)).
- Support has been added for formatted FSMSPEC files in the summary utility ([#4080](https://github.com/OPM/opm-common/pull/4080)).
- Improved the reading of possibly corrupt summary files ([#4079](https://github.com/OPM/opm-common/pull/4079)).
- The ACTNUM array is now updated immediately after scanning the GRID and EDIT sections ([#4082](https://github.com/OPM/opm-common/pull/4082)). This avoids problems where cells are deactivated by setting PORO equal to zero while at the same time setting regions numbers for these cells to invalid values (which is not currently supported).
- The compatibility with restart files has been improved specifically reading the well closure reason for history matching wells, and not using a guide rate model for wells with a target type of NONE ([#4090](https://github.com/OPM/opm-common/pull/4090)).
- The INIT file output code's assumption that tabulated saturation functions have at most NSSFUN (TABDIMS keyword item 3) rows per saturation region has been relaxed ([#4066](https://github.com/OPM/opm-common/pull/4066)). This allows larger saturation function tables to be output to the INIT file without an error being generated.
- Added support for restarting cases with group injection controlled voidage replacement ([#4108](https://github.com/OPM/opm-common/pull/4108)).
- Support has been expanded for saving/restoring scaling factors for segment level devices (SICDs, AICDs and Valves) to/from the Restart file ([#3978](https://github.com/OPM/opm-common/pull/3978)).
- Avoid unnecessarily checking if the fluid is saturated in CO2STORE and [H2STORE](#REF_HEADING_KEYWORD_H2STORE) cases ([#4125](https://github.com/OPM/opm-common/pull/4125)).
- If the injection fluid temperature has not been set for an injection well then it now defaults to the cell temperature in the upper most completion ([#5490](https://github.com/OPM/opm-simulators/pull/5490), [#4155](https://github.com/OPM/opm-common/pull/4155), [#4152](https://github.com/OPM/opm-common/pull/4152)).
- A warning message will now be generated and the well connection will be skipped if the well connection is in a cell that will be deactivated due a pore volume cut off specified by the MINPV or MINPVV keyword ([#4174](https://github.com/OPM/opm-common/pull/4174)). This also avoids a crash that would have occurred if any such well was re-completed at a later stage in the simulation.
- Restart support has been added for DRSDT ([#4166](https://github.com/OPM/opm-common/pull/4166)).
- The code now uses the failure flag of the tasklet runner to detect a failure while writing output and throw an error if this happens ([#5478](https://github.com/OPM/opm-simulators/pull/5478), [#4178](https://github.com/OPM/opm-common/pull/4178)).
- Added ESMRY methods to get a list of dates for summary results and obtain the units for a summary vector have been exported to Python ([#4184](https://github.com/OPM/opm-common/pull/4184)).
- Add support for converting between formatted (.FGRID) and unformatted (.GRID) grid files to the convertECL utility ([#4208](https://github.com/OPM/opm-common/pull/4208)).
- An error message will now be generated if an attempt is made to copy an undefined array (for example using the COPY keyword). Similarly, operations on an array (for example the ADD keyword) other than assignment (for example using the EQUALS keyword) will now generate an error message if the array being operated on has not been previously been defined in the input deck ([#4109](https://github.com/OPM/opm-common/pull/4109), [#4235](https://github.com/OPM/opm-common/pull/4235)). However, operations are allowed on multiplier arrays (MULTX, MULTY, etc.) that have not been explicitly defined in the input deck but have a default value of one ([#4235](https://github.com/OPM/opm-common/pull/4235)).
- Performance improvements in well name matching on platforms which support the Posix fnmatch() function ([#4243](https://github.com/OPM/opm-common/pull/4243)).
\end{verbatim}

The following bug fixes and improvements have been incorporated into
this release:

\begin{verbatim}
- Avoid sending and receiving from the same Buffer to fix errors with MPICH on Red Hat Linux ([#5297](https://github.com/OPM/opm-simulators/pull/5297)).
- Redistribute gaslift for groups with an ALQ limit ([#5296](https://github.com/OPM/opm-simulators/pull/5296)). Previousy, gaslift was only redistributed for groups limited by rate.
- Fixed output of CO2 mass in place ([#5324](https://github.com/OPM/opm-simulators/pull/5324)).
- Initialise the reservoir temperature in CO2STORE restart runs using the initial reservoir temperature if no temperatures are stored in the restart file, as is now the case for isothermal runs ([#5331](https://github.com/OPM/opm-simulators/pull/5331)).
- Mark the option 'BASIC=0' as allowed for the RPTRST keyword ([#5335](https://github.com/OPM/opm-simulators/pull/5335)).
- Corrected definitions of GMDS and GMGP summary vectors ([#5292](https://github.com/OPM/opm-simulators/pull/5292)).
- Fixed a typo that would cause problems when compiling using CUDA-aware MPI ([#5368](https://github.com/OPM/opm-simulators/pull/5368)).
- When checking for zero target rates, also check wells under group control where required ([#5232](https://github.com/OPM/opm-simulators/pull/5232)).
- Avoid possible use of out-of-scope temporary ([#5393](https://github.com/OPM/opm-simulators/pull/5393)).
- Fixed an issue where a UDQ assignment would not be performed  if it was triggered on the last time step of report step ([#5379](https://github.com/OPM/opm-simulators/pull/5379)).
- Three small fixes: Make sure that the same explicit/implicit fraction is used when computing THP(BHP) and when computing BHP(THP). Reduced risk of switching control in case of extreme VFP-extrapolation. Disallow local well control switching when the maximum number of well switches has been reached ([#5395](https://github.com/OPM/opm-simulators/pull/5395)).
- Fix the modification of transmissibility arrays (for example, using the MULTIPLY keyword in the EDIT section to modify TRANX) when running in parallel ([#5414](https://github.com/OPM/opm-simulators/pull/5414)).
- Avoid recursively deleting results in all subfolders if the command line argument --output-dir is empty ([#5427](https://github.com/OPM/opm-simulators/pull/5427)).
- Avoid mass creation when using DRSDT or DRSDTCON if hysteresis was deactivated ([#5407](https://github.com/OPM/opm-simulators/pull/5407)).
- If convergence fails due to NaN then report this as a problem rather than an error, then chop the time step and continue ([#5434](https://github.com/OPM/opm-simulators/pull/5434)).
- Support was previously added for GCONPROD items 11 to 13 but generally did not take effect. This has been now been fixed ([#5432](https://github.com/OPM/opm-simulators/pull/5432)).
- Check that a well exists before closing it due to a group action ([#5462](https://github.com/OPM/opm-simulators/pull/5462)).
- Switch to using size_t instead of int for buffer position to deal with buffers exceeding size intmax in MPI calls ([#5473](https://github.com/OPM/opm-simulators/pull/5473)).
- Avoid generating an exception from std::stoi() if the environment variable OMP_NUM_THREADS is not set ([#5504](https://github.com/OPM/opm-simulators/pull/5504), [#5516](https://github.com/OPM/opm-simulators/pull/5516)).
- Ensure default ALQ is set for all production wells on all processes. Fixes an MPI bug for cases where a well switches from producer to injector at the restart step ([#5534](https://github.com/OPM/opm-simulators/pull/5534)).
- Avoids an exception that would be thrown if one or more wells are shut in a CO2STORE run ([#5546](https://github.com/OPM/opm-simulators/pull/5546)).
- Avoid deleting derivatives in the case of zero pressure or zero flux that would have resulted in an incorrect matrix ([#5236](https://github.com/OPM/opm-simulators/pull/5236)).
- Fixed conditional expression in DRVDT initialisation code ([#5565](https://github.com/OPM/opm-simulators/pull/5565)).
- Set the correct number of target iterations in the simple iteration timestep controller ([#5476](https://github.com/OPM/opm-simulators/pull/5476)).
- Ensure well and accumulated group efficiency factors are applied to the connection energy rates ([#5323](https://github.com/OPM/opm-simulators/pull/5323)).
- Fixed calculation of injection well BHP for CO2STORE and [H2STORE](#REF_HEADING_KEYWORD_H2STORE) runs ([#5602](https://github.com/OPM/opm-simulators/pull/5602)).
- Fixed output of well temperature ([#5305](https://github.com/OPM/opm-simulators/pull/5305)).
- Avoid adding gas lift gas of non-open wells to the network ([#5601](https://github.com/OPM/opm-simulators/pull/5601)). Also allow non-zero artificial lift quantities (ALQ) for extended network models.
- Fixed an indexing error in the PID controller formula as used in the time stepping algorithm ([#4855](https://github.com/OPM/opm-simulators/pull/4855)).
- SFC (Space-Filling Curve) reordering for cells in AluGrid is now the default choice to optimise data access patterns for adaptive mesh refinements/coarsening ([#5630](https://github.com/OPM/opm-simulators/pull/5630)).
- Actually use fixed ordering when constructing matrices in the AMG hierarchy to get reproducible results in parallel runs ([#5646](https://github.com/OPM/opm-simulators/pull/5646)).
- Corrected transmissibilities calculated for pinched out cells when using PINCH option 4 equals ALL ([#5577](https://github.com/OPM/opm-simulators/pull/5577)).
- Call method with the correct precision in the rocSPARSE solver ([#5655](https://github.com/OPM/opm-simulators/pull/5655)).
- The lookup index for tracer concentrations when restarting parallel runs has been fixed ([#5670](https://github.com/OPM/opm-simulators/pull/5670)).
- Fixed a memory bug that only occurred when using the NLDD non-linear solver when there was an unconverged local NLDD domain, which had a multi-segment well on that domain ([#5676](https://github.com/OPM/opm-simulators/pull/5676)).
- Fixed a bug which prevented restarts for gas-water cases ([#5681](https://github.com/OPM/opm-simulators/pull/5681)).
- Solution data is now marked as invalid at the end of a time step in the experimental black-oil simulator (flowexp_blackoil). Previously, many quantities were not evaluated or updated and output to the Restart file ([#5685](https://github.com/OPM/opm-simulators/pull/5685)).
- Initialise the water saturation to one if the primary variable is switched from water saturation to gas-water ratio ([#908](https://github.com/OPM/opm-models/pull/908)).
- Add scoped and named variable to avoid using out-of-scope temporary ([#724](https://github.com/OPM/opm-grid/pull/724)).
- Make methods that handle well connections aware of possible future connections that might be opened in ACTIONX ([#742](https://github.com/OPM/opm-grid/pull/742)).
- Updated the molar mass of n-Decane ("C10") defined in the code from 0.0142 to 0.142 kg/gmol ([#4037](https://github.com/OPM/opm-common/pull/4037)).
- Change the dimension for critical temperature ([TCRIT](#REF_HEADING_KEYWORD_TCRIT)) from temperature to absolute temperature ([#4041](https://github.com/OPM/opm-common/pull/4041)).
- Previously three modules were exported to Python for use with PYACTION, now only opmcommon_python is exported ([#4029](https://github.com/OPM/opm-common/pull/4029)).
- Avoid computing a well THP if the VFP table is defaulted or explicitly set to zero ([#4058](https://github.com/OPM/opm-common/pull/4058)).
- Proper INIT file table output for the saturation functions in LET format ([#4051](https://github.com/OPM/opm-common/pull/4051)).
- Fix deadlock when outputting MULTPV to INIT file in parallel ([#4068](https://github.com/OPM/opm-common/pull/4068)).
- The value specified by the UDQPARAM keyword for undefined UDQs is now used as a fallback default for undefined or missing UDQs, in particular well, group or segment UDQs ([#4054](https://github.com/OPM/opm-common/pull/4054)).
- Fixed comparison and serialization from Runspec object where multiple members were ignored ([#4086](https://github.com/OPM/opm-common/pull/4086)).
- Fix output of LET curves by taking into account the connate saturations ([#4085](https://github.com/OPM/opm-common/pull/4085)).
- Trigger an update of the well targets when running WTMULT in an ACTIONX  block ([#4095](https://github.com/OPM/opm-common/pull/4095)).
- Fix salinity input in CO2STORE module and allow CO2STORE to be used with multiple PVT regions ([#4099](https://github.com/OPM/opm-common/pull/4099)).
- Wells modified by user-defined arguments (UDA) of a WSEGVALV keyword will be updated during time steps not only report steps ([#4121](https://github.com/OPM/opm-common/pull/4121)).
- Use gas-water ratio ([RSW](#REF_HEADING_KEYWORD_RSW_10_3)) rather than saturated gas-water ratio (RSWSAT) when calculating water density ([#4124](https://github.com/OPM/opm-common/pull/4124)).
- Hysteresis parameter get and set methods now work directly on min/max saturations rather than relative permeability and capillary pressure values to avoid conversion in the output layer ([#4089](https://github.com/OPM/opm-common/pull/4089)). Also, imbibition and drainage curves were mixed up, and connate water saturation was incorrectly added to the drainage wetting phase critical saturation.
- Correctly iterate through file deck container while erasing. This would fail if for example the last keyword of a block is erased. Also the iteration endpoint was not updated correctly ([#4122](https://github.com/OPM/opm-common/pull/4122)).
- Avoid referencing object until its index is known to be good ([#4127](https://github.com/OPM/opm-common/pull/4127)). This is the minimal change needed to allow the use of field level user defined quantities to define group input/production targets. However, such models cannot be reliably restarted.
- Fix warning message when resetting control modes other than 'RATE','BHP' to 'RATE' for historical injection wells ([#4138](https://github.com/OPM/opm-common/pull/4138)).
- Avoid duplicating objects on pointer serialization/de-serialization  ([#4141](https://github.com/OPM/opm-common/pull/4141)).
- Avoid division by zero when evaluating corner point coordinates if coordinate lines (COORD) have equal top and bottom depths ([#4145](https://github.com/OPM/opm-common/pull/4145)).
- Ensure that the unit conversion offset is not discarded when parsing Temperature units ([#4155](https://github.com/OPM/opm-common/pull/4155)). Also the units for the Joule-Thomson coefficient have been changed from Temperature/Pressure to AbsoluteTemperature/Pressure on the GASJT, OILJT amd WATJT keywords, and the units for the thermal expansion ratio for a cell in mechanics models have been changed from 1/Temperature to 1/AbsoluteTemperature on the THERMEXR keyword.
- The summary tool now prints the help message and returns failure if you run it with no parameters (or no parameters after the options). Previously, the  tool would generate a segmentation fault ([#4156](https://github.com/OPM/opm-common/pull/4156)).
- Wetting phase hystresis model has been fixed. The formulation as in Killough 1976 is followed more faithfully, in particular, the non-wetting phase saturation is now used to determine the turning points ([#4160](https://github.com/OPM/opm-common/pull/4160)).
- Use correct bhp limit for production wells under historic control when restarting a run ([#4112](https://github.com/OPM/opm-common/pull/4112)).
- Fix assignment of selected well variables using the UDQ keyword, and improve handling of quoted well names ([#4177](https://github.com/OPM/opm-common/pull/4177)).
- Proper handling of the WHISTCTL in restarted simulations ([#4119](https://github.com/OPM/opm-common/pull/4119)).
- The serializer nolonger passes null pointers to the packers, and a workaround for Dune 2.6 has been removed ([#4192](https://github.com/OPM/opm-common/pull/4192)).
- Flag indicating whether the parameters for the logarithm-based polymer shear thinning/thickening option include a reference temperature is now assigned to the correct member of the PLYSHLOG table object ([#4195](https://github.com/OPM/opm-common/pull/4195)).
- Fixed serialization and output of Ezrokhi density coefficients for parallel runs ([#4209](https://github.com/OPM/opm-common/pull/4209)).
- Fix limit on maximum number of iterations in flash calculation ([#4214](https://github.com/OPM/opm-common/pull/4214)).
- Added missing throw statements for logic errors ([#4203](https://github.com/OPM/opm-common/pull/4203)).
- If the VFP table number is defaulted for history wells using WCONHIST or WCONINJH then the previously defined value now will be used ([#4228](https://github.com/OPM/opm-common/pull/4228)).
- Previously, a warning 'well is not connected to the grid' is output if any connection is not connected to the grid. Now a warning message is only output if the well is completely unconnected from the grid after the entire COMPDAT record has been processed ([#4223](https://github.com/OPM/opm-common/pull/4223)).
- Added missing UDA Dimension for group production (GCONPROD) gas rate target or constraint (GRAT). Without this a restarted simulation with UDAs for GCONPROD's GRAT target or constraint  with fail to load ([#4248](https://github.com/OPM/opm-common/pull/4248)).
- Fixed a bug where well and group rates that can be specified by UDAs were converted to SI units twice ([#4269](https://github.com/OPM/opm-common/pull/4269)).
- The code now correctly identifies when a new well object is encountered ([#4278](https://github.com/OPM/opm-common/pull/4278)). Previously, the simulator would produce inconsistent results for polymer cases run in parallel.
\end{verbatim}

Known issues in this release of the simulator include:

\begin{verbatim}
- There is an error with the RSM header for summary vectors whose NUMS entry in the SMSPEC file is derived from more than a single number source (e.g., single region or segment ID). This applies to all block vectors (BGPV, BOPV, BWPV, etc.), connection level quantities (COPT, etc.), and inter-region flows such as ROFT etc ([#3078](https://github.com/OPM/opm-common/issues/3078)). The work around is to plot the data in OPM ResInsight and right-click on the plot to view and copy the data.
- As per previous releases of the radial model, the COORDSYS keyword item three must be set to COMP to complete the circle, this has not been implemented in this release. Also there appears to be a bug for full radial models when a well goes on BHP control that causes the well not to respect the BHP constraint, this eventually causes the well to die prematurely. See [#2640](https://github.com/OPM/opm-common/pull/2640) for a discussion on the topic.
- As in previous releases there are some issues with the OPERATE and OPERATER keywords associated with the input parsing; for various reasons a few of the fields require special case treatment in the grid processing, including (at least) MULTZ, PORV and ACTNUM, and for those keywords the OPERATE/OPERATER keyword doe not work. The work around is to use the MULTIPLY keyword instead.
- For the UDQ ASSIGN operator after the terminating “/” normally any comments can be entered; however, if there is “/” within the comment field, as per:
\end{verbatim}

\begin{verbatim}
ASSIGN FUNGLYLD 1.196   /    Condensate Yield (63.5/56.7)/(1.0 – 0.065)
\end{verbatim}

then the simulator will abort. The work around is to manually place the
comment characters ``--'' after the ASSIGN terminating ``/'', like so:

\begin{verbatim}
ASSIGN FUNGLYLD 1.196   / -- Condensate Yield (63.5/56.7)/(1.0 – 0.06)
\end{verbatim}

\begin{verbatim}
- At the moment, one cannot initialize tracers using the EQUALS keyword. Instead use the array format, that is the keyword followed by the required number of values, or the TVDP keyword in the SOLUTION section to set the initial tracer concentrations as a function of depth.
- Currently, gas tracers cannot be used if the dissolved gas phase, as per the DISGAS keyword in the RUNSPEC section, is active in the model.
- The summary vector RTIPTHEA, that defines the energy in-place between the initial and the current time for regions, is not supported unlike the FTIPTHEA and BTIPTHEA vectors. Secondly, the error message:
\end{verbatim}

\begin{longtable}[]{@{}
  >{\raggedright\arraybackslash}p{(\linewidth - 0\tabcolsep) * \real{1.0000}}@{}}
\toprule\noalign{}
\begin{minipage}[b]{\linewidth}\raggedright
Warning: Problem with summary keyword RTIPTHEA In RSM-THERMAL.data line
492 FIP region FIPHEA not defined in REGIONS section - RTIPTHEA ignored
\end{minipage} \\
\midrule\noalign{}
\endhead
\bottomrule\noalign{}
\endlastfoot
\end{longtable}

is incorrect, as the message indicates that it is being treated like a
named region, as per the FIP keyword, when it is actually a SUMMARY
vector
(\href{https://github.com/OPM/opm-simulators/issues/3870}{\#3870}).

\begin{verbatim}
- If there are cells that are very distorted, which can occur near fault planes, then the simulator may abort because it cannot calculate the pore volume of such cells. The work around is to re-generate the grid in the static model, taking care that the cells around the fault planes are more or less orthogonal ([#2992](https://github.com/OPM/opm-common/issues/2992) and [#3770](https://github.com/OPM/opm-simulators/issues/3770)).
- Currently the OPERATER keyword in the EDIT section does not work with the DEPTH, TRANX, TRANY and TRANZ property arrays ([#2994](https://github.com/OPM/opm-common/issues/2994) and [#748](https://github.com/OPM/opm-tests/pull/748)).
- If a standard well is fully declared in an ACTIONX block which is then activated at a later date, and later the well is modified to be a multi-segment well using the WELSEGS and COMPSEGS keywords, then this will cause the simulator to abort with an assert failure. The solution to this issue is to not use this type of work flow in declaring wells ([#2891](https://github.com/OPM/opm-common/issues/2891) and [#2895](https://github.com/OPM/opm-common/pull/2895)).
- Although the ACTIONX EXIT command works as expected, it does not write out the requested RSM file at the end of the run. However, the other SUMMARY and RESTART files are written out ([#2877](https://github.com/OPM/opm-common/issues/2877)).
- Although the GCONSUMP keyword in the SCHEDULE section is fully implemented as documented, it is not possible to verify the output as the associated SUMMARY vectors are not written out, that is the SUMMARY sales gas vectors FGSR, FGST, GGSR and GGST, and fuel vectors FGCR, FGCT, GGCR, and GGCT have not been implemented ([#2679](https://github.com/OPM/opm-common/issues/2679)).
- There are small differences in the behavior of the NEXTSTEP keyword in the RUNSPEC section between OPM Flow and the commercial simulator that remain unresolved ([#3745](https://github.com/OPM/opm-simulators/issues/3745)).
- There is a unit handling issue associated with OPERATE keyword. If the OPERATE(X) parameter has units, as for example PERMX, then the conversion is always done in SI units, despite the input deck declaring the deck to be fields units, as per FIELD keyword in the RUNSPEC section. Note that OPM Flow performs all of its calculations internally in SI and performs unit handling only when inputting the *.DATA file and when outputting result files. Thanks to [lrijkels](https://github.com/lrijkels) for reporting the issue. See [#4597](https://github.com/OPM/opm-simulators/issues/4597) for details.
- There is an issue associated with restarting from a restart file with the solution gas (Rs) maximum rate of increase, as defined by the DRSDT keyword in the SCHEDULE, that has been set to zero. This is because, the simulator does not save/restore this setting in simulator's restart files, which means that simulator misses the essential value zero upon restarting the case. As a work-around one can use the option --sched-restart=true, when running the restart case. This will initialize the restarted simulation based on information from the complete SCHEDULE section, instead of just the parts that we're going to simulate and the rest from the restart file. Thanks to [goncalvesmachadoc](https://github.com/goncalvesmachadoc) for reporting the issue. See [#4272](https://github.com/OPM/opm-simulators/issues/4272) for details.
- OPM cannot be built with dune-fem version 2.9 or later ([#4934](https://github.com/OPM/opm-simulators/issues/4934)). Please use a previous version.
- The simulator uses an irregular corner-point grid geometry with adjusted pore volumes to represent radial grids so it is not possible to create a full ring (360 degree disk) with only one cell in the theta direction (NY=1). The work around is to model a slice (say DTHETA=60 degrees). Note that as the angle increases larger pore volume adjustments are required ([#4755](https://github.com/OPM/opm-simulators/issues/4755)).
- In some cases with the network option the simulator can wrongly report that a well has no THP constraints, for example
\end{verbatim}

\begin{longtable}[]{@{}l@{}}
\toprule\noalign{}
GLIFT WTEST: Well S-3H does not have THP constraints \\
\midrule\noalign{}
\endhead
\bottomrule\noalign{}
\endlastfoot
\end{longtable}

when THP constraints have been defined
(\href{https://github.com/OPM/opm-simulators/issues/4887}{\#4887}).

\begin{verbatim}
- Dispersion in the gas phase leads to convergence issues for the 11th SPE CSP ([https://spe.org/csp](https://spe.org/csp)) models Version 11B and Version 11C (and is not fully tested). Therefore dispersion in the gas phase has been temporarily removed and a warning added ([#5101](https://github.com/OPM/opm-simulators/pull/5101), [#859](https://github.com/OPM/opm-models/pull/859)).
\end{verbatim}

The 2024-04 release consists of some new features and various
improvements and bug fixes. The new features and improvements include
the following highlights:

\begin{verbatim}
- Added support for the [CO2SOL](#REF_HEADING_KEYWORD_CO2SOL) and [H2SOL](#REF_HEADING_KEYWORD_H2SOL) keywords for modeling CO2 or H2 injection in hydrocarbon reservoirs. These are similar to the existing CO2STORE and [H2STORE](#REF_HEADING_KEYWORD_H2STORE) keywords for modeling CO2 or H2 injection in saline aquifers.
- The option has been added to output the liquid phase H2 mole fraction and vapor phase water mole fraction to the RESTART file when using the [H2STORE](#REF_HEADING_KEYWORD_H2STORE) option (RPTRST keyword mnemonics XMFH2 and YMFWAT).
- Added CO2 mass in place summary vectors for the field and region level when using CO2STORE, including trapped masses based on the definition of immobile gas in the 11th SPE Comparative Solution Project.
- Added support for the [SOURCE](#REF_HEADING_KEYWORD_SOURCE_12_3) keyword for modeling a simple source term.
- Support has been added for WDFAC and WDFACCOR keywords in the SCHEDULE section.  The option of setting connection specific D factors (Item 12 in COMPDAT) is now also supported.
- Added support for setting GCONPROD keyword item 7 (ACTION) equal to WELL. Added support for reason equals G (group) to the WTEST keyword item 3 (TEST).
- Added support for the [DIFFAGAS](#REF_HEADING_KEYWORD_DIFFAGAS_8_3) and [DIFFAWAT](#REF_HEADING_KEYWORD_DIFFAWAT_8_3) keywords to define respectively the gas phase and water phase diffusion coefficients based on a mass fraction formulation.
- Support has been added for the FBHPDEF keyword in the SCHEDULE section to define the default well BHP target for production wells and the default BHP constraint for injection wells.
- Support has been added for the SKIP, SKIP100, SKIP300 and ENDSKIP keywords.
- Support has been added for the NONNC keyword to deactivate non-neighbor connections.
- Added partial support for DATUMR and DATUMRX keywords in the SOLUTION section for use with block level summary keywords (BPP*).
- Limited support has been added for User Defined Tables (UDT).
- Added support for using region-level summary vectors in the defining expressions of field-level user defined quantities (UDQs).
- Support has been added for output of addition summary vectors, initialization arrays and restart arrays.
- A new brine-CO2 mutual solubility model has been implemented from Spycher & Pruess, Transp. Porous Med., 2010. This model improves the accuracy of the solubility calculations in CO2STORE at temperatures above 100 degrees Celsius. In addition, a new keyword, [ACTCO2S](#REF_HEADING_KEYWORD_ACTCO2S), has been added to choose between activity models.
- Added support for the [PCFACT](#REF_HEADING_KEYWORD_PCFACT_8_3) keyword in the PROPS sections to define the capillary pressure multiplication factor as a function of porosity change; used in conjunction with OPM Flow’s Salt Precipitation model.
- Added support for the [THCO2MIX](#REF_HEADING_KEYWORD_THCO2MIX_8_3) keyword in the PROPS section to specify the thermal mixing models for salt in water, CO2 dissolved in brine, and water vaporized in CO2.
- Added support for using well lists (WLIST) in ACTIONX blocks with the following keywords: COMPDAT, COMPLUMP, WCONINJE, WCONPROD, WECON, WEFAC, WELOPEN, WELPI, WELSPECS, WELTARG, WGRUPCON, WPIMULT, WSEGVALV, WTEST and WTMULT.
- Added support for an expanded set of summary vectors to be used in defining expressions for UDQs and in the condition blocks of the ACTIONX keyword.
- Removed the need for a run-function in PYACTION blocks. EclipseState, Schedule, ReportStep and SummaryState are now available as attributes of the module opm_embedded. The code maintains backwards compatibility with the previous usage.
- Tooltips can now be enabled when writing code in a Python IDE by typing "import opm_embedded".
- An error will now be reported and the simulation stopped when the PYACTION keyword is used if flow was built without embedded Python.
- Repaired the PYACTION functions open_well(), shut_well() and stop_well(), which open/shut/stop a well at the specified report step, and added the option to use these functions at the current report step (if no report step is specified).
- Added PYACTION function insert_keyword() to insert a keyword at the specified later report step or at the current report step (if no report step is specified). This functionality is available for the keywords: FIELD, ENDBOX, GCONINJE, GCONPROD, METRIC, MULTX, MULTX-, MULTY, MULTY-, MULTZ, MULTZ-, NEXT, NEXTSTEP, WCONINJE, WCONPROD, WEFAC, WELOPEN and WELTARG.
\end{verbatim}

The major command line changes made for this release are summarized in
Table A.2

\begin{longtable}[]{@{}
  >{\raggedright\arraybackslash}p{(\linewidth - 6\tabcolsep) * \real{0.2500}}
  >{\raggedright\arraybackslash}p{(\linewidth - 6\tabcolsep) * \real{0.2500}}
  >{\raggedright\arraybackslash}p{(\linewidth - 6\tabcolsep) * \real{0.2500}}
  >{\raggedright\arraybackslash}p{(\linewidth - 6\tabcolsep) * \real{0.2500}}@{}}
\toprule\noalign{}
\begin{minipage}[b]{\linewidth}\raggedright
OPM Flow 2024-04 New and Deprecated Command Line Options
\end{minipage} & \begin{minipage}[b]{\linewidth}\raggedright
\end{minipage} & \begin{minipage}[b]{\linewidth}\raggedright
\end{minipage} & \begin{minipage}[b]{\linewidth}\raggedright
\end{minipage} \\
\midrule\noalign{}
\endhead
\bottomrule\noalign{}
\endlastfoot
No. & Variable Name & Description & Default \\
1 & --accelerator-mode & Added rocsparse option. A defined character
string that selects the linear solver, usage:
`--accelerator-mode={[}none\textbar cusparse\textbar opencl\textbar amgcl\textbar rocalution\textbar rocsparse{]}'.
& ``none'' \\
2 & --debug-emit-cell-partition & Added new command line option. A
boolean value that determines whether (true) or not (false) to emit cell
partitions as a debugging aid
(\href{https://github.com/OPM/opm-simulators/pull/4938}{\#4938}). &
false \\
3 & --explicit-rock-compaction & Added new command line option
(\href{https://github.com/OPM/opm-simulators/pull/5089}{\#5089}). A
boolean value. Use pressure from end of the last time step when
evaluating rock compaction (true) or not (false). & false \\
4 & --external-partition & Added new experimental command line option
(\href{https://github.com/OPM/opm-simulators/pull/4882}{\#4882}). A
quoted character string. Name of file from which to load an externally
generated partitioning of the model's active cells for MPI distribution
purposes. If empty, the built-in partitioning method will be employed. &
``\,'' \\
5 & --linear-solver & Added the dilu and cpr\_trueimpesanalytic options
(\href{https://github.com/OPM/opm-simulators/pull/5119}{\#5119},
\href{https://github.com/OPM/opm-simulators/pull/5002}{\#5002},
\href{https://github.com/OPM/opm-simulators/pull/4899}{\#4899}). A
defined character string that selects the configuration of solver. Valid
options are: ilu0, dilu, cprw, cpr (an alias for cprw), cpr\_quasiimpes,
cpr\_trueimpes, cpr\_trueimpesanalytic, amg or hybrid (experimental).
Alternatively, you can request a configuration to be read from a JSON
file by giving the filename here, ending with `.json.'. & ``ilu0'' \\
6 & --local-domains-ordering-measure & Replaced the ``pressure'' option
with the ``maxpressure'' and ``averagepressure'' options. A defined
character string that selects the subdomain ordering measure. Allowed
values are `maxpressure', `averagepressure' and `residual'. &
``maxpressure'' \\
7 & --local-solve-approach & Changed the default value from ``jacobi''
to ``gauss-seidel''. A defined character string to choose the local
solve approach. Valid choices are jacobi and gauss-seidel. &
``gauss-seidel'' \\
8 & --local-tolerance-scaling-cnv & Changed the default value from 0.01
to 0.1. A positive real number. Set to a lower than 1.0 to use stricter
convergence tolerance for local solves. & 0.1 \\
9 & --min-strict-mb-iter & Added new command line option
(\href{https://github.com/OPM/opm-simulators/pull/5173}{\#5173}). An
integer value defining the number of Newton iterations before relaxed
tolerances can be used for the MB convergence criterion. Default -1
means that the relaxed tolerance is used when maximum number of Newton
iterations are reached. & -1 \\
10 & --newton-max-error & Changed the default value from 1e+10 to
1e+100. A positive real value. The maximum error tolerated by the Newton
method to which does not cause an abort. & 1e+100 \\
11 & --newton-min-iterations & Changed the default value from 1 to 2. A
positive integer value. The minimum number of Newton iterations per time
step. & 2 \\
12 & --pressure-max & Added new command line option. A positive real
value. Maximum absolute pressure. & 1e+99 \\
13 & --pressure-min & Added new command line option. A real value.
Minimum absolute pressure. & -1e+99 \\
14 & --pressure-scale & Added new command line option. A positive real
value. Scaling of pressure primary variable. & 1 \\
15 & --tolerance-mb-relaxed & Added new command line option
(\href{https://github.com/OPM/opm-simulators/pull/5173}{\#5173}). A
positive real value. Relaxed tolerated mass balance error that applies
for iterations after the iterations with the strict tolerance. &
1e-06 \\
16 & --use-implicit-ipr & Added new command line option. A boolean
value. Compute implicit IPR for stability checks and stable solution
search. & false \\
17 & --ecl-enable-drift-compensation & Renamed to
--enable-drift-compensation & true \\
18 & --newton-target-iterations & Changed the default value from 6 to
10. A positive integer value. The `optimum' number of Newton iterations
per time step. This is an experimental option that is not used by the
standard version of OPM Flow; the
--time-step-control-target-newton-iterations option should be used
instead. & 10 \\
19 & --vtk-write-ecl-tracer-concentration & Renamed to
--vtk-write-tracer-concentration & false \\
20 & --ecl-enable-tuning & Deprecated
(\href{https://github.com/OPM/opm-simulators/pull/5213}{\#5213}) & \\
21 & --ecl-newton-relaxed-tolerance & Deprecated & \\
22 & --ecl-newton-relaxed-volume-fraction & Deprecated & \\
23 & --ecl-newton-strict-iterations & Deprecated & \\
24 & --ecl-newton-sum-tolerance & Deprecated & \\
25 & --ecl-newton-sum-tolerance-exponent & Deprecated & \\
Notes: & & & \\
\end{longtable}

\emph{Table A.2: OPM Flow 2024-04 New and Deprecated Command Line
Options}

In addition to the above the following new features have been added to
the simulator:

\begin{verbatim}
- The Forchheimer term has been added as an additional skin term that depends (explicitly) on the rate ([#4832](https://github.com/OPM/opm-simulators/pull/4832)). Support has been added for WDFAC and WDFACCOR keywords in the SCHEDULE section. These keywords can be used to model a Forchheimer term in the well model either by specifying the D factor directly or by using the Dake model ([#3661](https://github.com/OPM/opm-common/pull/3661)). The option of setting connection specific D factors (Item 12 in COMPDAT) is now also supported ([#4971](https://github.com/OPM/opm-simulators/pull/4971)).
- Support has been added for the FBHPDEF keyword in the SCHEDULE section to define the default well BHP target for production wells and the default BHP constraint for injection wells ([#4969](https://github.com/OPM/opm-simulators/pull/4969), [#3747](https://github.com/OPM/opm-common/pull/3747)).
- Added partial support for fluid in place region (FIP) output to the RPTSOL and RPTSCHED keywords ([#4982](https://github.com/OPM/opm-simulators/pull/4982), [#4978](https://github.com/OPM/opm-simulators/pull/4978), [#4980](https://github.com/OPM/opm-simulators/pull/4980), [#3756](https://github.com/OPM/opm-common/pull/3756)). The simulator supports mnemonic FIP = 1 or 2, and partially supports FIP =3.
- Added support for mechanical dispersion as a compile time option ([#4960](https://github.com/OPM/opm-simulators/pull/4960), [#3737](https://github.com/OPM/opm-common/pull/3737), [#847](https://github.com/OPM/opm-models/pull/847)). This adds support for the linear dispersion model described in the 11th SPE CSP benchmark [document](https://github.com/Simulation-Benchmarks/11thSPE-CSP/blob/main/description/spe_csp11_description.pdf) (Eq. 2.3).
- Support has been added for specifying WCONPROD keyword item 12 (ALQ) with a user defined argument (UDA) ([#5033](https://github.com/OPM/opm-simulators/pull/5033)).
- Support has been added for specifying WSEGVALV keyword item 4 (AREAREST) with a user defined argument (UDA) ([#4873](https://github.com/OPM/opm-simulators/pull/4873)).
- Support has been added for the SKIP, SKIP100, SKIP300 and ENDSKIP keywords ([#5085](https://github.com/OPM/opm-simulators/pull/5085), [#3817](https://github.com/OPM/opm-common/pull/3817)). Note that OPM Flow will skip data between the SKIP300 and ENDSKIP keywords whereas the commercial black-oil simulator will ignore the SKIP300 keyword.
- Added support for defining a fluid mass (and enthalpy) source term within a grid cell using the new OPM Flow specific [SOURCE](#REF_HEADING_KEYWORD_SOURCE_12_3) keyword in the SCHEDULE section ([#5050](https://github.com/OPM/opm-simulators/pull/5050), [#3835](https://github.com/OPM/opm-common/pull/3835), [#1106](https://github.com/OPM/opm-tests/pull/1106), [#5107](https://github.com/OPM/opm-simulators/pull/5107), [#3864](https://github.com/OPM/opm-common/pull/3864)).
- Support has been added for the NONNC keyword in the RUNSPEC section to deactivate non-neighbor connections ([#5134](https://github.com/OPM/opm-simulators/pull/5134)).
- Added support for the [DIFFAGAS](#REF_HEADING_KEYWORD_DIFFAGAS_8_3) and [DIFFAWAT](#REF_HEADING_KEYWORD_DIFFAWAT_8_3) keywords in the PROPS section to define respectively the gas phase and water phase diffusion coefficients for each pseudo component and for each PVT region. The [DIFFAGAS](#REF_HEADING_KEYWORD_DIFFAGAS_8_3) and [DIFFAWAT](#REF_HEADING_KEYWORD_DIFFAWAT_8_3) keywords are equivalent to DIFFCGAS and DIFFCWAT keywords, but assume a mass fraction formulation of diffusion rather than the default mole fraction formulation ([#5143](https://github.com/OPM/opm-simulators/pull/5143), [#3878](https://github.com/OPM/opm-common/pull/3878)).
- Added partial support for DATUMR and DATUMRX keywords in the SOLUTION section for use with block level summary keywords (BPP*) ([#5241](https://github.com/OPM/opm-simulators/pull/5241), [#5240](https://github.com/OPM/opm-simulators/pull/5240), [#3958](https://github.com/OPM/opm-common/pull/3958)).
- Added support for using region-level summary vectors in the defining expressions of field-level UDQs ([#5250](https://github.com/OPM/opm-simulators/pull/5250), [#5247](https://github.com/OPM/opm-simulators/pull/5247), [#3977](https://github.com/OPM/opm-common/pull/3977), [#3980](https://github.com/OPM/opm-common/pull/3980)). For example:
\end{verbatim}

\begin{longtable}[]{@{}l@{}}
\toprule\noalign{}
UDQ DEFINE FURE2 (ROIP\_RE2 2)+(ROIP\_RE2 4)+(ROIP\_RE2 5) / / \\
\midrule\noalign{}
\endhead
\bottomrule\noalign{}
\endlastfoot
\end{longtable}

\begin{verbatim}
- Added support for GCONPROD item 7 (ACTION) equal to WELL. Added support for reason WTEST item 3 (TEST) equals G (group) ([#5172](https://github.com/OPM/opm-simulators/pull/5172), [#3894](https://github.com/OPM/opm-common/pull/3894)).
- Limited support has been added for User Defined Tables (see the UDTDIMS keyword in the RUNSPEC section and UDT keyword in the SCHEDULE section), one-dimensional tables are now supported ([#4940](https://github.com/OPM/opm-simulators/pull/4940), [#3728](https://github.com/OPM/opm-common/pull/3728)). The maximum number of dimensions in any given User Defined Table is specified by item 4 (MXDIMS) of the UDTDIMS keyword.
- Added support for setting a no flow and constant temperature boundary condition for the 11th SPE CSP benchmark cases B/C ([#865](https://github.com/OPM/opm-models/pull/865)).
- Added support for updating individual well properties for one or more wells (based on a well name, well name pattern or well list) using the WELSPECS keyword in the SCHEDULE section. In particular, this revised logic enables the controlling group to be changed without affecting any other well properties such as the location of the well head or the well reference depth ([#3703](https://github.com/OPM/opm-common/pull/3703)). As an example, this new WELSPECS behavior enables the following usage:
\end{verbatim}

\begin{longtable}[]{@{}
  >{\raggedright\arraybackslash}p{(\linewidth - 0\tabcolsep) * \real{1.0000}}@{}}
\toprule\noalign{}
\begin{minipage}[b]{\linewidth}\raggedright
ACTIONX A 1 / WOPR 'P*' \textless{} 123.4 / / WELSPECS `?' `LOWPRESS' /
/ ENDACTIO
\end{minipage} \\
\midrule\noalign{}
\endhead
\bottomrule\noalign{}
\endlastfoot
\end{longtable}

\begin{verbatim}
- A new brine-CO2 mutual solubility model has been implemented from Spycher & Pruess [Spycher, N., Pruess, K. A Phase-Partitioning Model for CO2–Brine Mixtures at Elevated Temperatures and Pressures: Application to CO2-Enhanced Geothermal Systems. Transp Porous Med 82, 173–196 (2010).]. This model improves the accuracy of the solubility calculations in CO2STORE at temperatures above 100 degrees Celsius. In addition, a new keyword, [ACTCO2S](#REF_HEADING_KEYWORD_ACTCO2S), has been added to choose between activity models. The original CO2STORE solubility and salt activity models are currently the default behavior (i.e., [ACTCO2S](#REF_HEADING_KEYWORD_ACTCO2S) keyword model number 3) ([#3694](https://github.com/OPM/opm-common/pull/3694)).
- Added the [CO2SOL](#REF_HEADING_KEYWORD_CO2SOL) and [H2SOL](#REF_HEADING_KEYWORD_H2SOL) keywords in the RUNSPEC section which activate either dissolved carbon dioxide (CO2) or hydrogen (H2) in the water phase, where CO2 or H2 is represented by the SOLVENT pseudo component, using the simulator’s CO2-Brine or H2-Brine PVT model. The [CO2SOL](#REF_HEADING_KEYWORD_CO2SOL) and [H2SOL](#REF_HEADING_KEYWORD_H2SOL) keywords can be used when modeling CO2 or H2 injection in depleted hydrocarbon reservoirs ([#3764](https://github.com/OPM/opm-common/pull/3764), [#851](https://github.com/OPM/opm-models/pull/851), [#4991](https://github.com/OPM/opm-simulators/pull/4991)).
- Added support for the [PCFACT](#REF_HEADING_KEYWORD_PCFACT_8_3) keyword in the PROPS sections to define the capillary pressure multiplication factor as a function of porosity change ([#3857](https://github.com/OPM/opm-common/pull/3857)). Currently the keyword is used in conjunction with OPM Flow’s Salt Precipitation model.
- Added methods to the Python module opm.simulators.BlackOilSimulator to access primary variables and fluid state variables ([#4950](https://github.com/OPM/opm-simulators/pull/4950)).
- Added docstrings for the Python bindings. Currently, the docstrings can be used to get context help in editors or in IPython. However, the plan is to also use these docstrings to generate sphinx documentation for the Python bindings ([#5242](https://github.com/OPM/opm-simulators/pull/5242)).
- Added PYACTION functions insert_keyword(const std::string& deck_string) and insert_keyword(const std::string& deck_string, std::size_t report_step) to insert a keyword. A keyword can be inserted at a later report step or at the current report step; inserting a keyword at a past report step or at a report step exceeding the total number of report steps throws an error ([#3994](https://github.com/OPM/opm-common/pull/3994)).
- Added support for the [THCO2MIX](#REF_HEADING_KEYWORD_THCO2MIX_8_3) keyword in the PROPS section to specify the thermal mixing models for salt in water, CO2 dissolved in brine, and water vaporized in CO2 ([#4012](https://github.com/OPM/opm-common/pull/4012)).
\end{verbatim}

Improvements include:

\begin{verbatim}
- Previously, the simulator would perform unnecessary network iterations if all wells belonging to the network were in history mode.  The simulator has been updated to avoid this and only compute the network pressures for reporting purposes ([#4881](https://github.com/OPM/opm-simulators/pull/4881)).
- Previously, if no well THP constraint had been specified then the THP would be set to zero. The simulator has been updated so that the THP would only be set to zero if no valid VFP table has been provided. If a VFP table has been provided but no THP constraint has been specified then the THP value is calculated for reporting purposes. In addition, if the well is part of a network then the THP value is required for a constraint check ([#4930](https://github.com/OPM/opm-simulators/pull/4930)).
- Updated code to output more information and avoid harsh termination due to assert in StandardWellPrimaryVariables ([#4946](https://github.com/OPM/opm-simulators/pull/4946)).
- The simulator has been updated to reduce well connection transmissibility factor due to salt precipitation ([#4919](https://github.com/OPM/opm-simulators/pull/4919)).
- The simulator has been updated to compute the reservoir volume rates based on the average pressure and temperature in the well for injectors in THERMAL cases ([#4826](https://github.com/OPM/opm-simulators/pull/4826)).
- The code has been updated to ensure there is a non-zero initial guess for well rates, and that non-zero water/gas fractions are available from the previous time step for the (explicit) VFP table interpolation. Previously, if this was not the case then the well was likely to converge to the zero rate and/or zero fractions case and would be deemed to be inoperable ([#4952](https://github.com/OPM/opm-simulators/pull/4952)).
- The simulator has been updated to allow the creating or opening of a well as a result of an ACTIONX keyword block in the middle rather than at the end of a report step ([#4749](https://github.com/OPM/opm-simulators/pull/4749)).
- Previously, numerical problems were reported as errors in the print file (*.PRT), before continuing the simulation. This resulted in successful runs being incorrectly reported as having errors. The code has been updated to report them as problems rather than errors ([#4977](https://github.com/OPM/opm-simulators/pull/4977), [#4957](https://github.com/OPM/opm-simulators/pull/4957)).
- Rerouting is now allowed in the standard and extended network models ([#3695](https://github.com/OPM/opm-common/pull/3695), [#4945](https://github.com/OPM/opm-simulators/pull/4945)).
- Added support for networks with multiple fixed-pressure nodes ([#3702](https://github.com/OPM/opm-common/pull/3702), [#4915](https://github.com/OPM/opm-simulators/pull/4915)).
- Added extra information about well convergence failures to the INFOITER file (if this file has been requested, the default is for this file not to be output) ([#4975](https://github.com/OPM/opm-simulators/pull/4975)).
- The simulator will now log a problem if asked to continue with the run after a convergence failure of the non-linear solver has occurred, by specifying the command line option --solver-continue-on-convergence-failure=true ([#4979](https://github.com/OPM/opm-simulators/pull/4979)).
- If parsing strictness is set to low with the command line option --parsing-strictness=low and no matching wells are found by various keywords in the SCHEDULE section then a warning will be issued and the run will continue ([#4905](https://github.com/OPM/opm-simulators/pull/4905), [#3698](https://github.com/OPM/opm-common/pull/3698)).
- Previously, numerical problems during well testing would be caught at a higher level, causing time step chops in serial mode and possibly MPI communication errors in parallel. With this change wells that experience such problems during testing will simply not open, and the MPI errors are no longer present ([#5032](https://github.com/OPM/opm-simulators/pull/5032)).
- A guard has been added against not a number (NaN) from certain expressions (e.g. evaluating a power of a negative number) in the valve models ([#5041](https://github.com/OPM/opm-simulators/pull/5041)).
- If Hierarchical Data Format (HDF5) input/output is requested but no support is available then an exception is thrown early instead of logging later an error when save/load is to be performed ([#5042](https://github.com/OPM/opm-simulators/pull/5042)).
- Functionality has been added for computing an implicit Inflow Performance Relationship (IPR). Moreover several functions are now utilizing this for stability checking, estimating the operability of stopped/unconverged wells and searching for solutions (VFP/IPR intersections) in problematic cases. In addition, direct computation of potentials is included. For the new code to take effect, currently simulations must be run with both options --local-well-solve-control-switching=true and --use-implicit-ipr=true ([#4986](https://github.com/OPM/opm-simulators/pull/4986)).
- The code has been updated to apply MULTX, MULTY and MULTZ transmissibility multipliers to boundary transmissibilities ([#5047](https://github.com/OPM/opm-simulators/pull/5047)).
- The name of the destination file is now included in warning messages about unhandled report keywords ([#5048](https://github.com/OPM/opm-simulators/pull/5048)).
- Reformatted production and injection reports in the Print file (*.PRT) by removing dashed lines between records ([#5049](https://github.com/OPM/opm-simulators/pull/5049)).
- Previously information about parallel runs was only output to the standard out. To help with debugging and replicating cases (e.g. in case of crashes) without saved standard output, the information about MPI processes and OMP threads is now also sent to the PRT file ([#5031](https://github.com/OPM/opm-simulators/pull/5031)).
- The test of whether the gas phase at the well connection is saturated or not has been made more robust by simply comparing the oil-gas ratio with saturated oil-gas ratio at the connection pressure ([#5075](https://github.com/OPM/opm-simulators/pull/5075)).
- Output of reservoir volumes to the Print file (*.PRT) has been added for custom fluid in place (FIP) regions. In addition, the compatibility of the output with the commercial simulator has been improved ([#5068](https://github.com/OPM/opm-simulators/pull/5068)).
- Previously the flow reports in the PRT file were restricted to wells active on the current MPI rank. The simulator now outputs the flow reports for all wells active at the current report step ([#5090](https://github.com/OPM/opm-simulators/pull/5090)).
- The output of Newton progress has been improved for the NLDD option ([#5104](https://github.com/OPM/opm-simulators/pull/5104)).
- Added parallel well domain and error logging for NLDD ([#5116](https://github.com/OPM/opm-simulators/pull/5116)).
- Currently the NLDD option requires the command line option --matrix-add-well-contributions=true. An error message will be now be generated if NLDD is used without this option set ([#5120](https://github.com/OPM/opm-simulators/pull/5120)).
- Improved handling of wells with zero rate where cross-flow is turned on and off ([#5139](https://github.com/OPM/opm-simulators/pull/5139), [#3880](https://github.com/OPM/opm-common/pull/3880)).
- The simulator has been updated to ignore the COORDSYS keyword in the GRID section ([#5161](https://github.com/OPM/opm-simulators/pull/5161)). Note that the simulator will stop if the number of reservoir grids to process, set using the NUMRES keyword in the GRID section, is greater than one.
- Improved error message to be more informative when an exception occurs in computeWellPotentials() ([#5170](https://github.com/OPM/opm-simulators/pull/5170)).
- The simulator now removes possible old simulation output files from the specified output directory before logging starts ([#5168](https://github.com/OPM/opm-simulators/pull/5168)).
- A number of minor improvements have been made to the time reporting output including the addition of setup time ([#5174](https://github.com/OPM/opm-simulators/pull/5174)).
\end{verbatim}

\begin{longtable}[]{@{}
  >{\raggedright\arraybackslash}p{(\linewidth - 0\tabcolsep) * \real{1.0000}}@{}}
\toprule\noalign{}
\begin{minipage}[b]{\linewidth}\raggedright
================ End of simulation =============== Number of MPI
processes: 1 Threads per MPI process: 2 Setup time: 0.29 s Deck input:
0.21 s Number of timesteps: 123 Simulation time: 0.98 s Assembly time:
0.31 s (Wasted: 0.0 s; 0.0\%) Well assembly: 0.04 s (Wasted: 0.0 s;
0.0\%) Linear solve time: 0.23 s (Wasted: 0.0 s; 0.0\%) Linear setup:
0.06 s (Wasted: 0.0 s; 0.0\%) Props/update time: 0.19 s (Wasted: 0.0 s;
0.0\%) Pre/post step: 0.14 s (Wasted: 0.0 s; 0.0\%) Output write time:
0.07 s Overall Linearizations: 483 (Wasted: 0; 0.0\%) Overall Newton
Iterations: 360 (Wasted: 0; 0.0\%) Overall Linear Iterations: 1311
(Wasted: 0; 0.0\%)
\end{minipage} \\
\midrule\noalign{}
\endhead
\bottomrule\noalign{}
\endlastfoot
\end{longtable}

\begin{verbatim}
- The code now avoids unnecessarily looking up the artificial lift quantity (ALQ) for injection wells ([#5221](https://github.com/OPM/opm-simulators/pull/5221)).
- Improved formatting of small numbers in adaptive time stepping output ([#5237](https://github.com/OPM/opm-simulators/pull/5237)).
- With multi-segment wells, the segment pressure can be lower than the BHP, so it is not desirable to use the BHP lower limit to limit the segment pressures. The simulator has been updated so that the BHP lower limit is only used to limit the BHP; a lower limit of zero is imposed on the segment pressure ([#5234](https://github.com/OPM/opm-simulators/pull/5234)).
- The simulator has been updated to allow for the possibility that not all well connections are in active cells for the purpose of calculating the WBPn summary vectors. This situation might arise for example when "re-parsing" portions of the SCHEDULE section following the successful triggering of an ACTIONX block ([#5258](https://github.com/OPM/opm-simulators/pull/5258), [#3988](https://github.com/OPM/opm-common/pull/3988)).
- Avoid repeatedly getting a parameter value in a performance-critical section of the code, which was causing a significant performance loss for the Norne model, by getting the parameter once in the constructor instead ([#5318](https://github.com/OPM/opm-simulators/pull/5318)).
- Avoid expensive call to FaceDir::FromIntersectionIndex() in a performance-critical section of the code, this fixes a significant performance regression on the order of 5% or more in the flux calculations that was introduced earlier ([#898](https://github.com/OPM/opm-models/pull/898)).
- Added support for using well lists (WLIST) in ACTIONX blocks with the following keywords: COMPDAT, COMPLUMP, WCONINJE, WCONPROD, WECON, WEFAC, WELOPEN, WELPI, WELSPECS, WELTARG, WGRUPCON, WPIMULT, WSEGVALV, WTEST and WTMULT ([#3741](https://github.com/OPM/opm-common/pull/3741)).
- Output computed surface densities for CO2 and Brine to the PRT file if the CO2STORE option is enabled ([#3730](https://github.com/OPM/opm-common/pull/3730)).
- The simulator will now output the time step in units of hours when using laboratory (LAB) units. Previously units of days would have been used for all unit systems ([#3829](https://github.com/OPM/opm-common/pull/3829)).
- Improved the handling of errors related to the grid keywords. Missing or ambiguous keywords are detected and the error messages list the valid options ([#3849](https://github.com/OPM/opm-common/pull/3849), [#3551](https://github.com/OPM/opm-common/pull/3551)).
- Added support for pattern matching in the MULTFLT keyword ([#3851](https://github.com/OPM/opm-common/pull/3851)).
- Improved the error reporting for SALTSOL tables when a column is missing ([#3850](https://github.com/OPM/opm-common/pull/3850)).
- Added support for setting the same region for the source and target in the MULTREGT keyword. In this case the multipliers will be applied to all connections within the specified region as well as to connections between the specified region and all other regions ([#3845](https://github.com/OPM/opm-common/pull/3845)).
- The code has been updated so that for mixed wettability systems it no longer scales the negative part of the oil-water capillary pressure curve when using the SWATINIT keyword ([#3720](https://github.com/OPM/opm-common/pull/3720)). This behavior is controlled by an option in the commercial simulator.
- The CO2STORE option could lead to slow runs if a lot of extrapolation was done resulting in numerous warning messages. These messages have been turned off unless the debug option has been set. Note that extrapolation of tables can be avoided by limiting the temperatures and pressures using the --temperature-min, --temperature-max, --pressure-min and --pressure-max command line options ([#3858](https://github.com/OPM/opm-common/pull/3858)).
- The simulator now outputs a meaningful error message in the case where it fails to open a potentially corrupt file ([#3870](https://github.com/OPM/opm-common/pull/3870)).
- Improved parsing of sections in the input deck and support for optional sections EDIT, REGIONS and SUMMARY ([#3863](https://github.com/OPM/opm-common/pull/3863)).
- Enforce consistency checks on keywords specifying the standard and extended network models respectively ([#3885](https://github.com/OPM/opm-common/pull/3885)).
- The simulator has been updated to check whether the standard network model is active when initializing a network balancing operation.  Previously, the simulator would only check for the extended network model. This was an issue because the pressure convergence tolerance is set to zero unless a network model is active ([#3874](https://github.com/OPM/opm-common/pull/3874)).
- Add support for operations (e.g. EQUALS) on the MINPVV array ([#3897](https://github.com/OPM/opm-common/pull/3897)).
- Added support for numbers in user defined region set names, for example FIPAB1 and FIPAB2 ([#3979](https://github.com/OPM/opm-common/pull/3979)).
- Allow the use of the FIELD group for selected group level keywords when handling UDQs and especially ACTIONX condition blocks, for example GGOR ‘FIELD’ > 123.4 ([#3993](https://github.com/OPM/opm-common/pull/3993)).
- Allow an expanded set of summary vectors to be used in defining expressions for UDQs and in the condition blocks of the ACTIONX keyword. In principle, this change enables using all known summary vector categories as part of the defining expressions in a UDQ, but additional testing is needed before claiming to fully support such usage ([#3968](https://github.com/OPM/opm-common/pull/3968)).
- The simulator will now report an error if the PYACTION keyword is used when flow has been built without embedded Python ([#3998](https://github.com/OPM/opm-common/pull/3998)). Previously, only a warning message would have been reported.
- Removed the need for a run-function in the Python code in PYACTION. The EclipseState, Schedule, ReportStep and SummaryState have been made available as attributes of the module opm_embedded instead. This enables tooltips for writing code in a Python IDE, after importing the library as "import opm_embedded", typing "opm_embedded." will show the tooltips. The code will still be backwards compatible ([#3986](https://github.com/OPM/opm-common/pull/3986)).
- Added check that when a non-default THP constraint has been specified for a well that a valid VFP table has been defined ([#4005](https://github.com/OPM/opm-common/pull/4005)).
- Previously, when temperature results were not available 0 K would be converted to the units system specified in the deck and output to the SUMMARY file (for example -273.15 or -459.67 depending on the unit system). The simulator will now output a zero temperature in the specified units system when temperature results are not available ([#4013](https://github.com/OPM/opm-common/pull/4013)).
- Added PYACTION functions to open, close or shut a well without specifying the report step ([#4019](https://github.com/OPM/opm-common/pull/4019)).
- Added an error message if an attempt is made to use the COMPSEGS keyword to define the well segment connections before the structure of the multi-segment well has been defined using the WELSEGS keyword ([#4031](https://github.com/OPM/opm-common/pull/4031)).
- The code now determines if a history matching producer with a reservoir volume rate (RESV) target is subject to a zero rate constraint by checking if the specified historic rates for oil, gas and water are all zero ([#4039](https://github.com/OPM/opm-common/pull/4039)).
\end{verbatim}

In addition, the following new SUMMARY keywords are now recognized as
described by the comments in Table A.3.

\begin{longtable}[]{@{}
  >{\raggedright\arraybackslash}p{(\linewidth - 4\tabcolsep) * \real{0.3333}}
  >{\raggedright\arraybackslash}p{(\linewidth - 4\tabcolsep) * \real{0.3333}}
  >{\raggedright\arraybackslash}p{(\linewidth - 4\tabcolsep) * \real{0.3333}}@{}}
\toprule\noalign{}
\begin{minipage}[b]{\linewidth}\raggedright
No.
\end{minipage} & \begin{minipage}[b]{\linewidth}\raggedright
Summary Keyword
\end{minipage} & \begin{minipage}[b]{\linewidth}\raggedright
Comment
\end{minipage} \\
\midrule\noalign{}
\endhead
\bottomrule\noalign{}
\endlastfoot
1 & NPR, GNETPR & Added support for output of network pressures based on
rates at end of the time step
(\href{https://github.com/OPM/opm-common/pull/3701}{\#3701},
\href{https://github.com/OPM/opm-simulators/pull/4914}{\#4914}). Note
that GNETPR is an alias for NPR included since NPR can cause problems
for certain post-processors. \\
2 & CGFRF, CGFRS, COFRF and COFRS & Added support for output of well
connection free gas, solution gas, liquid oil and vaporized oil flow
rate vectors
(\href{https://github.com/OPM/opm-simulators/pull/5010}{\#5010},
\href{https://github.com/OPM/opm-common/pull/3782}{\#3782}). \\
3 & CDFAC & Added support for output of well connection D factor
(\href{https://github.com/OPM/opm-simulators/pull/4832}{\#4832}). \\
4 & FGMIP, FGMGP, FGMDS, FGMTR and FGMMO & Added CO2 mass in place
summary vectors for the field when using the CO2STORE option: total mass
in place, mass in gas phase, mass dissolved in water phase, mass in
trapped gas phase and mass in mobile gas phase
(\href{https://github.com/OPM/opm-simulators/pull/5114}{\#5114},
\href{https://github.com/OPM/opm-common/pull/3868}{\#3868}). \\
5 & RGMIP, RGMGP, RGMDS, RGMTR and RGMMO & Added CO2 mass in place
summary vectors for FIP regions when using the CO2STORE option: total
mass in place, mass in gas phase, mass dissolved in water phase, mass in
trapped gas phase and mass in mobile gas phase
(\href{https://github.com/OPM/opm-simulators/pull/5114}{\#5114},
\href{https://github.com/OPM/opm-common/pull/3868}{\#3868}). \\
6 & BGIP, BGIPL & Added support for block total gas in place (BGIP) and
block gas in place in the water phase (BGIPL) for gas-water systems
(\href{https://github.com/OPM/opm-simulators/pull/5188}{\#5188}). \\
7 & BPPO, BPPG and BPPW & Added support for output of block level datum
depth corrected phase pressures
(\href{https://github.com/OPM/opm-simulators/pull/5240}{\#5240},
\href{https://github.com/OPM/opm-common/pull/3959}{\#3959},
\href{https://github.com/OPM/opm-common/pull/3958}{\#3958}). \\
8 & FGKDI, FGKDM, FGKTR and FGKMO RGKDI, RGKDM, RGKTR and RGKMO & The
11th Society of Petroleum Engineers Comparative Solution Project
(http://spe.org/csp) defines the ``immobile free-phase CO2'' as ``CO2 at
saturations for which the non-wetting phase relative permeability equals
zero''. Field and Region level summary vectors have been added when
using the CO2STORE option based on this definition as follows
(\href{https://github.com/OPM/opm-simulators/pull/5281}{\#5281},
\href{https://github.com/OPM/opm-common/pull/4010}{\#4010}): Gas Moles
in the Immobile Gas Phase (Non-wetting relative permeability equals
zero), Gas Moles in the Mobile Gas Phase (Non-wetting relative
permeability greater than zero), Gas Mass in the Immobile Gas Phase, and
Gas Mass in the Mobile Gas Phase. \\
\end{longtable}

\emph{Table A.3: New SUMMARY Keywords for the 2024-04 Release}

Support has been added for the following initialization output as
described by the comments in Table A.4.

\begin{longtable}[]{@{}
  >{\raggedright\arraybackslash}p{(\linewidth - 4\tabcolsep) * \real{0.3333}}
  >{\raggedright\arraybackslash}p{(\linewidth - 4\tabcolsep) * \real{0.3333}}
  >{\raggedright\arraybackslash}p{(\linewidth - 4\tabcolsep) * \real{0.3333}}@{}}
\toprule\noalign{}
\begin{minipage}[b]{\linewidth}\raggedright
No.
\end{minipage} & \begin{minipage}[b]{\linewidth}\raggedright
Mnemonic
\end{minipage} & \begin{minipage}[b]{\linewidth}\raggedright
Comment
\end{minipage} \\
\midrule\noalign{}
\endhead
\bottomrule\noalign{}
\endlastfoot
1 & MULTX, MULTY and MULTZ MULTX-, MULTY- and MULTZ- & Added support for
outputting transmissibility multiplier arrays to the INIT file. The
MULTX, MULTY and MULTZ arrays are always written. If the first item of
the GRIDOPTS keyword is NO then only the MULTX-, MULTY- and MULTZ-
arrays that have be specified in the deck will be written. If the first
item of the GRIDOPTS keyword is YES then all transmissibility multiplier
arrays are written
(\href{https://github.com/OPM/opm-common/pull/3997}{\#3997}). \\
2 & SWATINIT & The simulator will now output the SWATINIT array to INIT
file if it exists in the input deck
(\href{https://github.com/OPM/opm-common/pull/3832}{\#3832}). \\
\end{longtable}

\emph{Table A.4: New Initialization Output for the 2024-04 Release}

Supported has been added for the following restart output as described
by the comments in Table A.5.

\begin{longtable}[]{@{}
  >{\raggedright\arraybackslash}p{(\linewidth - 4\tabcolsep) * \real{0.3333}}
  >{\raggedright\arraybackslash}p{(\linewidth - 4\tabcolsep) * \real{0.3333}}
  >{\raggedright\arraybackslash}p{(\linewidth - 4\tabcolsep) * \real{0.3333}}@{}}
\toprule\noalign{}
\begin{minipage}[b]{\linewidth}\raggedright
No.
\end{minipage} & \begin{minipage}[b]{\linewidth}\raggedright
Mnemonic
\end{minipage} & \begin{minipage}[b]{\linewidth}\raggedright
Comment
\end{minipage} \\
\midrule\noalign{}
\endhead
\bottomrule\noalign{}
\endlastfoot
1 & RESIDUAL & Added option to output the residuals at the end of the
time step. To activate this output the command line option
--enable-opm-rst-file=true should be specified and the mnemonic RESIDUAL
should be added to the RPTRST keyword
(\href{https://github.com/OPM/opm-simulators/pull/4937}{\#4937}). \\
2 & XMFH2, YMFWAT & The option has been added to output the liquid phase
H2 mole fraction and vapor phase water mole fraction (XMFH2 and YMFWAT)
when using \hyperref[REF_HEADING_KEYWORD_H2STORE]{H2STORE}. To activate
this output the command line option --enable-opm-rst-file=true should be
specified
(\href{https://github.com/OPM/opm-simulators/pull/4964}{\#4964}).
Previously the mnemonics XMFCO2 and YMFWAT were only available for
CO2STORE. \\
3 & RSWSOL & Support has been added for output of dissolved solvent to
water ratio RSWSOL
(\href{https://github.com/OPM/opm-simulators/pull/4991}{\#4991,}
\href{https://github.com/OPM/opm-common/pull/3764}{\#3764},
\href{https://github.com/OPM/opm-models/pull/851}{\#851}). \\
4 & FLOWS-, FLORES- & Support has been added for output of negative
direction inter-block flows at surface conditions (FLOWS-) and at
reservoir conditions (FLORES-)
(\href{https://github.com/OPM/opm-simulators/pull/5038}{\#5038},
\href{https://github.com/OPM/opm-common/pull/3823}{\#3823},
\href{https://github.com/OPM/opm-models/pull/854}{\#854}). \\
5 & RPORV & Support has been added for output of pore volumes at
reservoir conditions using the RPORV mnemonic of the RPTRST keyword
(\href{https://github.com/OPM/opm-simulators/pull/5092}{\#5092}). \\
6 & CONV & Added partial support for output of cells causing convergence
problems for the residuals of the oil, water, gas, polymer, brine and
solvent equations to the restart file using the CONV mnemonic of the
RPTRST keyword
(\href{https://github.com/OPM/opm-simulators/pull/5054}{\#5054},
\href{https://github.com/OPM/opm-simulators/pull/5112}{\#5112}). \\
7 & PCGW & Added support for outputting gas-water capillary pressure
(Pcgw) to the restart file using the PCGW mnemonic of the RPTRST keyword
(\href{https://github.com/OPM/opm-simulators/pull/5126}{\#5126}). \\
8 & FIPOIL, FIPWAT, FIPGAS, RFIPOIL, RFIPWAT, RFIPGAS & Added support
for output of fluid in place arrays for oil, gas and water to the
RESTART file. The SFIPOIL, SFIPWAT and SFIPGAS are the surface
conditions volumes for oil, water and gas respectively (these are
aliases for FIPOIL, FIPWAT and FIPGAS). The RFIPOIL, RFIPWAT and RFIPGAS
are the reservoir conditions volumes for oil, water and gas respectively
(\href{https://github.com/OPM/opm-simulators/pull/5224}{\#5224},
\href{https://github.com/OPM/opm-common/pull/3952}{\#3952}). \\
\end{longtable}

\emph{Table A.5: New Restart Output for the 2024-04 Release}

The following bug fixes and improvements have been incorporated into
this release.

\begin{verbatim}
- The simulator has been updated to use the applyRestartSwatInit() member function instead of directly mutating a data member of the EpsInfo structure. This way we defer the PCOW re-scaling for SWATINIT to a context with a more complete view of the changes needed to convey the information to all components involved. With this change the simulator is mostly able to reproduce the base run in a restart run from the end of the historic period for a real field case that has strong SWATINIT scaling of the oil-water capillary pressure function ([#4944](https://github.com/OPM/opm-simulators/pull/4944), [#3731](https://github.com/OPM/opm-common/pull/3731)).
- The code was updated to include a new event to address an issue with a reported case where WELOPEN failed to open a well previously shut due to an economic limit ([#4924](https://github.com/OPM/opm-simulators/pull/4924), [#3707](https://github.com/OPM/opm-common/pull/3707)).
- The logger is now initialized before parsing ([#5007](https://github.com/OPM/opm-simulators/pull/5007)). Previously, uninitialized values would have been used when the parser logged some errors.
- Fixed an issue where [BCPROP](#REF_HEADING_KEYWORD_BCPROP) would be ignored if it was not set in the first time step. This update makes it possible to set [BCPROP](#REF_HEADING_KEYWORD_BCPROP) later in the schedule ([#4912](https://github.com/OPM/opm-simulators/pull/4912), [#833](https://github.com/OPM/opm-models/pull/833)).
- Previously, a well with shut-in instruction "shut" that changes to non-operable during an iteration will give the message well xxx gets SHUT during iteration, but no further actions are triggered. This is due to the flag changed_to_stopped_this_step_ not being set. The simulator has been updated so that a non-operable well will give the message well xxx gets STOPPED during iteration (irrespective of it's shut-in instructions), and at the end of the time-step get closed according to it's shut/stop instruction (with accompanied message) ([#5014](https://github.com/OPM/opm-simulators/pull/5014), [#4902](https://github.com/OPM/opm-simulators/pull/4902)).
- Fixed an issue where there was an inconsistency between well BHP and perforation pressures, and between well rate and perforation rate ([#5015](https://github.com/OPM/opm-simulators/pull/5015)).
- Fixed a bug in the Gauss-Seidel NLDD approach that would actually do a Jacobi initialization for the very first iteration of a local domain solve. We add the maxpressure option for Gauss-Seidel NLDD domain ordering, and make it the new default ordering, and make Gauss-Seidel the default approach. Finally, we make the default number of domains equal to one domain per 1000 cells, instead of just a single domain, and tweak the default convergence tolerance scaling for local solves. These changes are intended to lower the bar to testing the --nonlinear-solver=nldd option. You still need to use --matrix-add-well-contributions=true to run with NLDD, otherwise the defaults should be sufficient for reasonable behavior ([#5021](https://github.com/OPM/opm-simulators/pull/5021)).
- Previously, combining --load-step with different --output-dir triggered a warning. The code has been updated to not issue a warning since these options can safely be used together ([#5024](https://github.com/OPM/opm-simulators/pull/5024)).
- Previously, when singular matrices occurred when solving multi-segment wells an exception would be thrown only on one process resulting in MPI errors about truncated messages. The code has now been updated so that all MPI processes will throw, terminate the non-linear solver and cut the time step to continue the simulation ([#5025](https://github.com/OPM/opm-simulators/pull/5025), [#5036](https://github.com/OPM/opm-simulators/pull/5036)).
- Fixed the sign of the production rate for group pressure maintenance ([#5026](https://github.com/OPM/opm-simulators/pull/5026)).
- Get a complete list of fluid in place (FIP) regions from FieldPropsManager rather than from SummaryConfig, which only holds regions with summary keywords ([#5034](https://github.com/OPM/opm-simulators/pull/5034), [#3813](https://github.com/OPM/opm-common/pull/3813)).
- Previously an exception thrown by WetGasPVT::saturationPressure due to failure to converge within 20 iterations was only seen on one MPI rank and the others continued. The code now communicates the problem and throws on all MPI processes. This results in the time step being cut as a result and simulation continuing ([#5046](https://github.com/OPM/opm-simulators/pull/5046)).
- Ensure the input temperature is used in the density calculations in the initialization ([#5058](https://github.com/OPM/opm-simulators/pull/5058)).
- Avoid modifying outputDir, leave empty if defaulted ([#5073](https://github.com/OPM/opm-simulators/pull/5073)).
- The code has been updated to only output FLOWS and FLORES for interior cells. This fixes an issue with parallel output of FLOWS and FLORES ([#5043](https://github.com/OPM/opm-simulators/pull/5043)).
- The setting of solvent values has been made conditional on the SOLVENT being active ([#5086](https://github.com/OPM/opm-simulators/pull/5086)).
- Previously when assembling the acceleration pressure drop terms for multi-segment wells certain derivatives for reverse flow were skipped. The code has been rewritten to fix this issue and improves convergence in some cases ([#5040](https://github.com/OPM/opm-simulators/pull/5040)).
- Code has been updated to use the gas pressure rather than the oil pressure in gas-water simulations when modifying rock properties due to compaction ([#5081](https://github.com/OPM/opm-simulators/pull/5081)).
- The code has been updated to honor the default minimum number of Newtons iterations. Previously the simulator would perform a minimum of two Newton iterations even though the default was one ([#5094](https://github.com/OPM/opm-simulators/pull/5094)). Note that also the default has been changed from one to two.
- The code has been updated to use the arithmetic average of compaction transmissibility multipliers rather than the upstream value ([#5088](https://github.com/OPM/opm-simulators/pull/5088), [#858](https://github.com/OPM/opm-models/pull/858)).
- A fix has been provided for a crash when cleaning up the Damaris mesh offset variable due to the shared memory buffer size being off by one ([#5109](https://github.com/OPM/opm-simulators/pull/5109)).
- The simulator now supports at least NTFIP (REGDIMS keyword item 1) distinct regions for inter-region flow to avoid a possible segmentation failure ([#5117](https://github.com/OPM/opm-simulators/pull/5117)).
- A couple of changes have been made to fix network convergence issues. The lower bound on the network pressure update has been removed as this could lead to convergence issues. In addition, the convergence tolerance has been tightened by a factor of 10 for wells that are controlled under a dynamic THP-limit since the default (standard) well tolerance was not sufficiently strict to get network convergence in some cases ([#5128](https://github.com/OPM/opm-simulators/pull/5128)).
- Previously the solvent production rate was not included in the total gas production rate when applying the gas production rate limit (GRAT) in the well control keyword WCONPROD. The gas production rate limit is now applied to the total gas production rate including both the gas and solvent pseudo components ([#5140](https://github.com/OPM/opm-simulators/pull/5140)).
- A bug revealed by (#5124) in the local well control switching option where previously stopped wells would re-open during the local solve has been fixed ([#5129](https://github.com/OPM/opm-simulators/pull/5129)).
- Previously, NEXTSTEP was ignored if the user did not run OPM Flow explicitly setting the command line parameter --enable-tuning=true. Now NEXTSTEP will be applied regardless of which parameter value was set ([#5215](https://github.com/OPM/opm-simulators/pull/5215)).
- The handling of the NEXTSTEP keyword in an ACTIONX block has now been fixed ([#5222](https://github.com/OPM/opm-simulators/pull/5222)).
- The code has been updated to ensure that the simulator does not advance beyond the current report step and that the next time step is not too small if the NEXTSTEP keyword is used in an ACTIONX block ([#5259](https://github.com/OPM/opm-simulators/pull/5259)).
- Fixed an issue where previously the low parsing strictness command line option was not being respected when reprocessing the schedule section after the application of an action ([#5220](https://github.com/OPM/opm-simulators/pull/5220)).
- Fixed a crash in the Python module due to duplicate symbols that depended on the order of the object libraries in the linker command during the build ([#5235](https://github.com/OPM/opm-simulators/pull/5235)).
- Added missing convergence metrics to the domain convergence report. Previously, for NLDD local solves the metrics were always zero and non-converged local solves were accepted ([#5249](https://github.com/OPM/opm-simulators/pull/5249)).
- The code has been updated to call loadRestartData on all processes even if there are no local wells since group data is still required ([#5252](https://github.com/OPM/opm-simulators/pull/5252)).
- Corrected formatting of debug message in MultisegmentWellSegments.cpp to avoid throwing exception in libfmt ([#5261](https://github.com/OPM/opm-simulators/pull/5261)).
- Fixed a potential out-of-bounds reference that could occur if there are no wells ([#5264](https://github.com/OPM/opm-simulators/pull/5264)).
- Fixed a segmentation fault due to taking a reference to a temporary variable in the WBP calculation service ([#5267](https://github.com/OPM/opm-simulators/pull/5267)).
- Fixed a bug introduced when the ROCm dependency on OpenCL was removed (#4883), which disabled the option to run the rocSPARSE and cuSPARSE backends with the optimized version of the block Jacobi incomplete lower-upper (ILU) preconditioner ([#5255](https://github.com/OPM/opm-simulators/pull/5255)).
- Fixed dangling references in Logger for test_LogOutputHelper that can result in a segmentation fault that has been seen on the ppc64el architecture ([#5263](https://github.com/OPM/opm-simulators/pull/5263)).
- The code now makes sure that the fast AMG preconditioner is reconstructed when performing an update ([#5284](https://github.com/OPM/opm-simulators/pull/5284)).
- Fixed an issue where the send and receive buffers were the same which caused an MPICH error on MPICH based Red Hat Enterprise Linux 7/8 builds ([#5297](https://github.com/OPM/opm-simulators/pull/5297)).
- Fixed a bug which caused an error whenever item 4 (PINCHCAL) of the PINCH keyword was set to ALL. The case where one of the pinched out cells has zero z-direction permeability (or z-direction transmissibility multiplier) is now supported. However, the ALL option is still not fully supported ([#707](https://github.com/OPM/opm-grid/pull/707), [#701](https://github.com/OPM/opm-grid/pull/701)).
- Fixed an integer overflow in the initial processing of grids with more than 134 million cells that would cause a segmentation fault ([#710](https://github.com/OPM/opm-grid/pull/710)).
- The simulator now correctly allows for solvent flux when the dissolved gas in water option has been activated by the DISGASW keyword in the RUNSPEC section ([#866](https://github.com/OPM/opm-models/pull/866)).
- The ‘meaning’ variables are now initialized to avoid potential out-of-bounds dereferences ([#880](https://github.com/OPM/opm-models/pull/880)).
- Fixed a bug where setting the saturation table number (SATNUM) in the COMPDAT keyword (item 7) to a value other than the cell SATNUM resulted in a segmentation fault in two-phase cases. The simulator now checks which two-phase approach is active and sets the correct relative permeability parameters accordingly ([#3739](https://github.com/OPM/opm-common/pull/3739)).
- Fixed an issue where the simulator interpreted the month ‘FEB’ in an ACTIONX condition as a field-level summary vector and wrote a value of zero rather two (for February) to the restart file. This resulted in an unhelpful error message when later trying to perform a restart ([#3826](https://github.com/OPM/opm-common/pull/3826)).
- The simulator now assumes that the reservoir density and internal energy of CO2 and H2 are not affected by their water content ([#3806](https://github.com/OPM/opm-common/pull/3806)). This avoids the use of calculations based on the ideal gas assumptions implemented in simpleHuDuanH20.hpp. The formulation is now as in the 11th SPE CSP.
- The simulator previously assumed that all UDQs loaded from the restart file had an UPDATE status of ON ([#3808](https://github.com/OPM/opm-common/pull/3808)). The simulator now inspects the actual update status and uses that to configure the UDQ when restarting the simulation ([#3833](https://github.com/OPM/opm-common/pull/3833)).
- Previously the simulator conflated history control with the WHISTCTL keyword setting, thus failing to respect any control mode setting in the WCONHIST keyword in the case of a simulation restart in the historic period. The WHISTCTL keyword setting has been decoupled from that of the control mode setting specified in the WCONHIST keyword thereby supporting both modes when a simulation is restarted in the historic period. This, in turn, fixes a simulation restart on a real field ([#3854](https://github.com/OPM/opm-common/pull/3854)).
- Fixed output of the group hierarchy chart to the PRT and DBG files as requested by the WELSPECS mnemonic of the RPTSCHED keyword ([#3856](https://github.com/OPM/opm-common/pull/3856)).
- Corrected the units for region level CO2 summary vectors RGCDI, RGCDM and RWCD to be moles instead of surface volumes ([#3866](https://github.com/OPM/opm-common/pull/3866)).
- Fixed support for the not equals operators ‘.ne.’ and ‘!=’ in ACTIONX keywords ([#3893](https://github.com/OPM/opm-common/pull/3893)).
- Account for well and/or group efficiency factors when evaluating rates for non-default region sets, for example ROPR_ABC and ROPT_ABC ([#3955](https://github.com/OPM/opm-common/pull/3955)).
- The code has been updated so that TSINIT (item 1) of the TUNING keyword only applies to the next time step ([#3953](https://github.com/OPM/opm-common/pull/3953)).
- Fixed the handling of the MULTX, MULTX-, MULTY, MULTY-, MULTZ, MULTZ- keywords and EQUALS MULT* keywords in the EDIT section. Previously multiple occurrences of these keywords in the EDIT section would have been applied cumulatively. Now the EDIT section has its own set of multipliers which are applied to the ones in the GRID section when the end of the EDIT section is reached ([#4002](https://github.com/OPM/opm-common/pull/4002)).
- Added a check that a keyword is valid for use with the PYACTION keyword before it is inserted into the deck ([#4008](https://github.com/OPM/opm-common/pull/4008)).
- Fixed the initialization of gas-water simulations with a transition zone ([#5293](https://github.com/OPM/opm-simulators/pull/5293)).


- There is an error with the RSM header for summary vectors whose NUMS entry in the SMSPEC file is derived from more than a single number source (e.g., single region or segment ID). This applies to all block vectors (BGPV, BOPV, BWPV, etc.), connection level quantities (COPT, etc.), and inter-region flows such as ROFT etc ([#3078](https://github.com/OPM/opm-common/issues/3078)). The work around is to plot the data in OPM ResInsight and right-click on the plot to view and copy the data.
- The GDFILE keyword in the GRID section loads a grid file in various formats, with the FMTOPT parameter setting the format type of the file. If the variable FMTOPT is omitted then the default is for binary file input for the commercial simulator; whereas, OPM Flow derives FMTOPT from the file extension (*.EGRID or *.FEGRID), making FMTOPT superfluous. However, if the extension is lower case then OPM Flow may incorrectly determine the file type. The work around is change the extension to upper case.
- As per previous releases of the radial model, the COORDSYS keyword item three must be set to COMP to complete the circle, this has not been implemented in this release. Also there appears to be a bug for full radial models when a well goes on BHP control that causes the well not to respect the BHP constraint, this eventually causes the well to die prematurely. See [#2640](https://github.com/OPM/opm-common/pull/2640) for a discussion on the topic.
- As in previous releases there are some issues with the OPERATE and OPERATER keywords associated with the input parsing; for various reasons a few of the fields require special case treatment in the grid processing, including (at least) MULTZ, PORV and ACTNUM, and for those keywords the OPERATE/OPERATER keyword doe not work. The work around is to use the MULTIPLY keyword instead.
- For the UDQ ASSIGN operator after the terminating “/” normally any comments can be entered; however, if there is “/” within the comment field, as per:
\end{verbatim}

\begin{verbatim}
ASSIGN FUNGLYLD 1.196   /    Condensate Yield (63.5/56.7)/(1.0 – 0.065)
\end{verbatim}

then the simulator will abort. The work around is to manually place the
comment characters ``--'' after the ASSIGN terminating ``/'', like so:

\begin{verbatim}
ASSIGN FUNGLYLD 1.196   / -- Condensate Yield (63.5/56.7)/(1.0 – 0.06)
\end{verbatim}

\begin{verbatim}
- At the moment, one cannot initialize tracers using the EQUALS keyword. Instead use the array format, that is the keyword followed by the required number of values, or the TVDP keyword in the SOLUTION section to set the initial tracer concentrations as a function of depth.
- Currently, gas tracers cannot be used if the dissolved gas phase, as per the DISGAS keyword in the RUNSPEC section, is active in the model.
- The summary vector RTIPTHEA, that defines the energy in-place between the initial and the current time for regions, is not supported unlike the FTIPTHEA and BTIPTHEA vectors. Secondly, the error message:
\end{verbatim}

\begin{longtable}[]{@{}
  >{\raggedright\arraybackslash}p{(\linewidth - 0\tabcolsep) * \real{1.0000}}@{}}
\toprule\noalign{}
\begin{minipage}[b]{\linewidth}\raggedright
Warning: Problem with summary keyword RTIPTHEA In RSM-THERMAL.data line
492 FIP region FIPHEA not defined in REGIONS section - RTIPTHEA ignored
\end{minipage} \\
\midrule\noalign{}
\endhead
\bottomrule\noalign{}
\endlastfoot
\end{longtable}

is incorrect, as the message indicates that it is being treated like a
named region, as per the FIP keyword, when it is actually a SUMMARY
vector
(\href{https://github.com/OPM/opm-simulators/issues/3870}{\#3870}).

\begin{verbatim}
- If there are cells that are very distorted, which can occur near fault planes, then the simulator may abort because it cannot calculate the pore volume of such cells. The work around is to re-generate the grid in the static model, taking care that the cells around the fault planes are more or less orthogonal ([#2992](https://github.com/OPM/opm-common/issues/2992) and [#3770](https://github.com/OPM/opm-simulators/issues/3770)).
- Currently the OPERATER keyword in the EDIT section does not work with the DEPTH, TRANX, TRANY and TRANZ property arrays ([#2994](https://github.com/OPM/opm-common/issues/2994) and [#748](https://github.com/OPM/opm-tests/pull/748)).
- If a standard well is fully declared in an ACTIONX block which is then activated at a later date, and later the well is modified to be a multi-segment well using the WELSEGS and COMPSEGS keywords, then this will cause the simulator to abort with an assert failure. The solution to this issue is to not use this type of work flow in declaring wells ([#2891](https://github.com/OPM/opm-common/issues/2891) and [#2895](https://github.com/OPM/opm-common/pull/2895)).
- Although the ACTIONX EXIT command works as expected, it does not write out the requested RSM file at the end of the run. However, the other SUMMARY and RESTART files are written out ([#2877](https://github.com/OPM/opm-common/issues/2877)).
- Although the GCONSUMP keyword in the SCHEDULE section is fully implemented as documented, it is not possible to verify the output as the associated SUMMARY vectors are not written out, that is the SUMMARY sales gas vectors FGSR, FGST, GGSR and GGST, and fuel vectors FGCR, FGCT, GGCR, and GGCT have not been implemented ([#2679](https://github.com/OPM/opm-common/issues/2679)).
- There are small differences in the behavior of the NEXTSTEP keyword in the RUNSPEC section between OPM Flow and the commercial simulator that remain unresolved ([#3745](https://github.com/OPM/opm-simulators/issues/3745)).
- There is a unit handling issue associated with OPERATE keyword. If the OPERATE(X) parameter has units, as for example PERMX, then the conversion is always done in SI units, despite the input deck declaring the deck to be fields units, as per FIELD keyword in the RUNSPEC section. Note that OPM Flow performs all of its calculations internally in SI and performs unit handling only when inputting the *.DATA file and when outputting result files. Thanks to [lrijkels](https://github.com/lrijkels) for reporting the issue. See [#4597](https://github.com/OPM/opm-simulators/issues/4597) for details.
- There is an issue associated with restarting from a restart file with the solution gas (Rs) maximum rate of increase, as defined by the DRSDT keyword in the SCHEDULE, that has been set to zero. This is because, the simulator does not save/restore this setting in simulator's restart files, which means that simulator misses the essential value zero upon restarting the case. As a work-around one can use the option --sched-restart=true, when running the restart case. This will initialize the restarted simulation based on information from the complete SCHEDULE section, instead of just the parts that we're going to simulate and the rest from the restart file. Thanks to [goncalvesmachadoc](https://github.com/goncalvesmachadoc) for reporting the issue. See [#4272](https://github.com/OPM/opm-simulators/issues/4272) for details.
- OPM cannot be built with dune-fem version 2.9 or later ([#4934](https://github.com/OPM/opm-simulators/issues/4934)). Please use a previous version.
- The simulator uses an irregular corner-point grid geometry with adjusted pore volumes to represent radial grids so it is not possible to create a full ring (360 degree disk) with only one cell in the theta direction (NY=1). The work around is to model a slice (say DTHETA=60 degrees). Note that as the angle increases larger pore volume adjustments are required ([#4755](https://github.com/OPM/opm-simulators/issues/4755)).
- In some cases with the network option the simulator can wrongly report that a well has no THP constraints, for example
\end{verbatim}

\begin{longtable}[]{@{}l@{}}
\toprule\noalign{}
GLIFT WTEST: Well S-3H does not have THP constraints \\
\midrule\noalign{}
\endhead
\bottomrule\noalign{}
\endlastfoot
\end{longtable}

when THP constraints have been defined
(\href{https://github.com/OPM/opm-simulators/issues/4887}{\#4887}).

\begin{verbatim}
- Dispersion in the gas phase leads to convergence issues for the 11th SPE CSP ([https://spe.org/csp](https://spe.org/csp)) models Version 11B and Version 11C (and is not fully tested). Therefore dispersion in the gas phase has been temporarily removed and a warning added ([#5101](https://github.com/OPM/opm-simulators/pull/5101), [#859](https://github.com/OPM/opm-models/pull/859)).
\end{verbatim}

We dedicate this OPM 2023-10 release to our colleague and friend David
Baxendale. David passed away in late June 2023 after a short severe
illness. Our thoughts are with his wife and son. The OPM community is
very thankful for all his contributions to OPM and the fruitful
discussions with him about issues with and advancements of the
simulator.

David started contributing to OPM in 2016 as OPMUSER on github and
continued his good work until his very last days. We owe the OPM
Reference manual to him. He started this heroic effort in 2017 and it
now has thousands of pages. We, his colleagues and friends, are and will
be surely missing him with his reservoir engineering expertise and
know-how, his enthusiasm, and humor.

The 2023-10 release consists of some new features and various
improvements and bug fixes. Our main target was to support more keywords
used for relevant fields and reduce differences between OPM flow and the
commercial simulator. These improvements include:

\begin{itemize}
\tightlist
\item
  Added support for temperature (THERMAL) plus salt precipitation
  (PRECSALT) modeling in gas-water-brine (GAS-WATER-BRINE) systems.
\item
  Added support for modeling dissolved gas in water (DISGASW) and
  vaporized water in the gas phase (VAPWAT) in the thermal-gas-water
  simulator.
\item
  Support for modeling FOAM combined with SOLVENT.
\item
  Partial support for WAGHYSTR keyword (Water-Alternating-Gas
  hysteresis).
\item
  Improvements to many user-facing error messages.
\item
  More graceful exits for problems in parallel runs.
\item
  Temperature is output to the RESTART file if requested via RPTRST.
\item
  Added support for WBP, WBP4, WBP5 and WBP9 in the SUMMARY section to
  output well block averaged pressures for open completions.
\item
  Added support for initializing constant flux aquifers from a restart.
\item
  Faster two-point flux-approximation introduced in the last release is
  now also used for linearizing gas-oil cases with energy (with
  diffusion) and gas-oil diffusion.
\end{itemize}

The major command line changes made for this release are summarized in
Table A.6

\begin{longtable}[]{@{}
  >{\raggedright\arraybackslash}p{(\linewidth - 6\tabcolsep) * \real{0.2500}}
  >{\raggedright\arraybackslash}p{(\linewidth - 6\tabcolsep) * \real{0.2500}}
  >{\raggedright\arraybackslash}p{(\linewidth - 6\tabcolsep) * \real{0.2500}}
  >{\raggedright\arraybackslash}p{(\linewidth - 6\tabcolsep) * \real{0.2500}}@{}}
\toprule\noalign{}
\begin{minipage}[b]{\linewidth}\raggedright
OPM Flow 2023-10 New and Deprecated Command Line Options
\end{minipage} & \begin{minipage}[b]{\linewidth}\raggedright
\end{minipage} & \begin{minipage}[b]{\linewidth}\raggedright
\end{minipage} & \begin{minipage}[b]{\linewidth}\raggedright
\end{minipage} \\
\midrule\noalign{}
\endhead
\bottomrule\noalign{}
\endlastfoot
No. & Variable Name & Description & Default \\
1 & --linear-solver & Added option to use the experimental hybrid solver
configuration. Valid options are: ilu0 (default), cprw, cpr (an alias
for cprw), cpr\_quasiimpes, cpr\_trueimpes, amg or hybrid
(experimental). Alternatively, you can request a configuration to be
read from a JSON file by giving the filename here, ending with `.json.'.
& ilu0 \\
2 & --linear-solver-print-json-definition & Added option to write the
JSON definition of the linear solver setup to the DBG file. & true \\
3 & --linear-solver-reduction & The minimum reduction of the residual
which the linear solver must achieve. & 0.01 \\
4 & --local-domain-ordering-measure & Added parameter to specify the
domain ordering measure as either residual or pressure when using the
NLDD solver
(\href{https://github.com/OPM/opm-simulators/pull/4738}{\#4738}). &
pressure \\
5 & --local-domains-partitioning-imbalance & Subdomain partitioning
imbalance tolerance for the NLDD solver. 1.03 is 3 percent imbalance. &
1.03 \\
6 & --local-domains-partitioning-method & Subdomain partitioning method
for the NLDD solver. Allowed values are `zoltan', `simple', and the name
of a partition file ending with `.partition'. & zoltan \\
7 & --local-solve-approach & Choose local solve approach for the NLDD
solver. Valid choices are jacobi and gauss-seidel. & jacobi \\
8 & --local-tolerance-scaling-cnv & Set lower than 1.0 to use stricter
convergence tolerance for local solves when using the NLDD solver. &
0.01 \\
9 & --local-tolerance-scaling-mb & Set lower than 1.0 to use stricter
convergence tolerance for local solves when using the NLDD solver. &
1 \\
10 & --local-well-solve-control-switching & Added trigger to enable
(true) or disable (false) well control/status switching during local
well equation solves when using the NLDD solver
(\href{https://github.com/OPM/opm-simulators/pull/4895}{\#4895}). &
false \\
11 & --max-local-solve-iterations & Max iterations for local solves with
NLDD nonlinear solver. & 20 \\
12 & --maximum-water-saturation & Maximum water saturation. & 1 \\
13 & --network-max-iterations & Maximum number of iterations in the
network solver before giving up. & 200 \\
14 & --network-max-strict-iterations & Maximum iterations in network
solver before relaxing tolerance. & 100 \\
15 & --nldd-num-initial-newton-iter & Added option to specify number of
global non-linear (Newton) iterations performed by the NLDD solver
before starting the local non-linear iterations. The default value of 1
preserves the current behavior
(\href{https://github.com/OPM/opm-simulators/pull/4922}{\#4922}). & 1 \\
16 & --nonlinear-solver & Choose nonlinear solver. Valid choices are
newton or nldd. & newton \\
17 & --num-local-domains & Number of local domains for NLDD nonlinear
solver. Note this is an experimental feature in the current release
2023.10 that is expected to be more complete and tested by the 2024.10
release. & 0 \\
18 & --save-step & Added new options to keep only the last step. Save
serialized state to .OPMRST file. Either a specific report step, ``all''
to save all report steps or ``:x'' to save every x'th step. Use negative
values of ``x'' to keep only the last written step, or ``last'' to save
every step, keeping only the last
(\href{https://github.com/OPM/opm-simulators/pull/4807}{\#4807}). &
``\,'' \\
19 & --water-only-threshold & Cells with water saturation above or equal
to this value are considered one-phase water only. & 1 \\
20 & --load-file & FileName for .OPMRST file used to load serialized
state. If empty (``\,''), CASENAME.OPMRST is used. & ``\,'' \\
21 & --num-pressure-points-equil & Number of pressure points (in each
direction) in tables used for equilibration (\#4718). & 2000 \\
22 & --save-file & FileName for .OPMRST file used for saving serialized
state. If empty, CASENAME.OPMRST is used. & ``\,'' \\
23 & --use-average-density-ms-wells & Approximate segment densitities by
averaging over segment and its outlet. & false \\
24 & --ecl-max-time-step-size-after-well-event & Maximum time step size
after a well event (seconds). Default is equivalent to 1 year.
Depreciated. & 3.15576e+07 \\
25 & --ecl-restart-shrink-factor & Factor by which the time step is
reduced after convergence failure. Depreciated & 3 \\
Notes: & & & \\
\end{longtable}

\emph{Table A.6: OPM Flow 2023-10 New and Deprecated Command Line
Options}

In addition to the above the following new features have been added to
the simulator:

\begin{verbatim}
- When the general well specification data is defined by the WELSPECS keyword in the SCHEDULE section the well can now be assigned directly to the FIELD group in item 2 GRPNAME ([#3485](https://github.com/OPM/opm-common/pull/3485) and [#4608](https://github.com/OPM/opm-simulators/pull/4608)). Previously there was a restriction preventing wells from being parented directly to FIELD. Although wells are now allowed to be parented directly to FIELD, this is discouraged and a warning message will be issued. Mixing wells and groups as children of a single group is still forbidden.
- Added support for dissolved gas in water (DISGASW) in the gas-water simulator with salt precipitation (PRECSALT) and vaporized water in the gas phase (VAPWAT) ([#4623](https://github.com/OPM/opm-simulators/pull/4623)). However, input of gas solubility in water is not currently supported so this option is only useful in combination with CO2STORE.
- Added support for modeling dissolved gas in water (DISGASW) and vaporized water in the gas phase (VAPWAT) in the thermal-gas-water simulator ([#4661](https://github.com/OPM/opm-simulators/pull/4661)).
- Added support for temperature (THERMAL) plus salt precipitation (PRECSALT) modeling in gas-water-brine (GAS-WATER-BRINE) systems as specified in the RUNSPEC section ([#4650](https://github.com/OPM/opm-simulators/pull/4650)). This has to date only been tested in combination with CO2STORE. Note that generally, if the PRECSALT keyword has been activated in the input deck then the VAPWAT keyword should also be activated.
- Added support for modeling foam (FOAM) plus solvent (SOLVENT) in the simulator ([#4654](https://github.com/OPM/opm-simulators/pull/4654), [#3523](https://github.com/OPM/opm-common/pull/3523) and [#805](https://github.com/OPM/opm-models/pull/805)). In addition, gas or water is now allowed as the transport phase for the foam as specified by FOAMOPTS item 1 (previously only gas was supported as the foam transport phase).
- Added support for WVFPDP in the SCHEDULE section ([#4620](https://github.com/OPM/opm-models/pull/4620) and [#3504](https://github.com/OPM/opm-common/pull/3504)). The WVFPDP keyword modifies a well’s Bottom-Hole Pressure (“BHP”) estimated by the simulator by interpolation of the Vertical Flow Performance (“VFP”) tables.
- Added support for GCONPROD item 7 ACTION equal to NONE in the SCHEDULE section, which specifies that no action is to be taken if the oil, water, gas or liquid rate constraints are violated. Previously only the RATE option was supported by the simulator ([#4658](https://github.com/OPM/opm-simulators/pull/4658)).
- Added partial support for WAGHYSTR keyword in the PROPS section ([#4710](https://github.com/OPM/opm-simulators/pull/4710) and [#3542](https://github.com/OPM/opm-common/pull/3542)). This keyword defines the parameters for the Water-Alternating-Gas (“WAG”) hysteresis option, when the hysteresis option has been activated by the HYSTER variable on the SATOPTS keyword in the RUNSPEC section. Only gas phase hysteresis is currently supported by the WAGHYSTR keyword. The residual oil modification fraction, which would only be active when the STONE1 three-phase oil relative permeability model is used, is not currently supported.
- Support added for PPCWMAX to limit the maximum capillary pressure scaling when initializing the model using SWATINIT, and the option to modify the connate water saturation to match the input water saturation if the capillary pressure is exceeded ([#4707](https://github.com/OPM/opm-simulators/pull/4707) and [#3570](https://github.com/OPM/opm-common/pull/3570)).
- Added support for WINJMULT in the SCHEDULE section to define pressure dependent injectivity multipliers for injection wells and can be used to approximate the change in injectivity due to hydraulic fracturing ([#4686](https://github.com/OPM/opm-simulators/pull/4686)).
- Added partial support for WPAVE and WWPAVE defining the method and parameters for calculating a well’s block average pressures for either all wells or specific wells ([#4695](https://github.com/OPM/opm-simulators/pull/4695), [#4694](https://github.com/OPM/opm-simulators/pull/4694) and [#4693](https://github.com/OPM/opm-simulators/pull/4693)). Calculation of block averages pressures is currently only supported for OPEN completions (not for ALL completions) as specified in WPAVE item 4 and WWPAVE item 5.
- Aquifer keywords have been enabled when [H2STORE](#REF_HEADING_KEYWORD_H2STORE) is specified in the RUNSPEC section with brine modeled by the OIL phase ([#4791](https://github.com/OPM/opm-simulators/pull/4791)). This is similar to the behavior when CO2STORE is specified.
- Added a model for formation damage due to suspended solids in the injection water ([#4346](https://github.com/OPM/opm-simulators/pull/4346) and [#3313](https://github.com/OPM/opm-common/pull/3313)).  New OPM specific well keywords [WINJDAM](#REF_HEADING_KEYWORD_WINJDAM), [WINJFCNC](#REF_HEADING_KEYWORD_WINJFCNC) and [WINJCLN](#REF_HEADING_KEYWORD_WINJCLN) have been added to the SCHEDULE section. [WINJDAM](#REF_HEADING_KEYWORD_WINJDAM) sets up the filter cake properties for the specified water injection well, [WINJFCNC](#REF_HEADING_KEYWORD_WINJFCNC) defines the filtrate concentration in the injection water, and [WINJCLN](#REF_HEADING_KEYWORD_WINJCLN) specifies that a fraction of the filter cake has been cleaned up.
- User defined arguments (UDA) have been enabled for the [WINJFCNC](#REF_HEADING_KEYWORD_WINJFCNC) keyword item 2 (FCONCPPM) the injection concentration parameter ([#4763](https://github.com/OPM/opm-simulators/pull/4763) and [#3614](https://github.com/OPM/opm-common/pull/3614)).
- Added partial support for GRUPNET, which defines the standard group network parameters used to model the flow and pressure behavior within the network ([#4760](https://github.com/OPM/opm-simulators/pull/4760), [#3609](https://github.com/OPM/opm-common/pull/3609), [#4815](https://github.com/OPM/opm-simulators/pull/4815) and [#3651](https://github.com/OPM/opm-common/pull/3651)). GRUPNET item 5 OPTION1 equals NO only is supported (defines how the groups production target should be achieved). Item 6 OPTION2 equals NO or FLO is supported but ALQ is not supported (defines how gas lift gas flows through the group’s pipeline). Item 7 OPTION3 equals NONE only is supported (defines if the ALQ-PIPE variable should be reset).
- Added partial support for GECON ([#4819](https://github.com/OPM/opm-simulators/pull/4819) and [#3657](https://github.com/OPM/opm-common/pull/3657)). Item 7 WORKOVER only supports NONE, workover procedures are not currently implemented. Item 8 ENDRUN only supports NO, end run is not currently implemented. Item 9 MXWELLS is not supported and must be defaulted.
- Support has been added for the CSKIN keyword in the SCHEDULE section to modify the connection skin factor ([#4871](https://github.com/OPM/opm-simulators/pull/4871) and [#3681](https://github.com/OPM/opm-common/pull/3681)).
- Support for WSF and GSF keywords in the PROPS section has been added for [H2STORE](#REF_HEADING_KEYWORD_H2STORE) runs (as for CO2STORE) ([#4528](https://github.com/OPM/opm-simulators/pull/4528)). The WSF and GSF keywords define the water relative permeability data versus water saturation tables and the gas relative permeability data versus gas saturation tables for when only gas and water are present in the input deck. These keywords can only be used with either the CO2STORE or [H2STORE](#REF_HEADING_KEYWORD_H2STORE) models.
- Support has been added for setting the WGRUPCON keyword Item 2 STATUS in the SCHEDULE section equal to NO ([#4575](https://github.com/OPM/opm-simulators/pull/4575)).
- Support has been added for setting the PINCH item 2 PINCHOPT parameter equal to NOGAP in the GRID section. This parameter is now fully supported by the simulator. Previously only the default GAP option was supported ([#4603](https://github.com/OPM/opm-simulators/pull/4603)).
- Added support for initializing constant flux aquifers from a restart ([#4520](https://github.com/OPM/opm-simulators/pull/4520) and [#4519](https://github.com/OPM/opm-simulators/pull/4519)). In particular, form constant flux aquifer objects from the restart step's collection if available and properly initialize their total produced volume.
- Added support for gas-water-solvent systems ([#4568](https://github.com/OPM/opm-simulators/pull/4568) and [#4548](https://github.com/OPM/opm-simulators/pull/4548)).
- Added partial support for setting items 11, 12 and 13 of GCONPROD in the SCHEDULE section: ACTWAT, ACTGAS and ACTLIQ respectively define the action to be taken if the water rate (WRAT), gas rate (GRAT) and liquid rate (LRAT) constraints defined by GCONPROD are violated. Supported options are now NONE or RATE. Options CON, +CON, WELL, and PLUG are not currently supported ([#4748](https://github.com/OPM/opm-simulators/pull/4748)).
- Added support for Non-linear Domain Decomposition (NLDD) as the non-linear solver.
- The option has been added to allow well control or well status to be updated during the iteration process during local solve for the well equations ([#4895](https://github.com/OPM/opm-simulators/pull/4895)). In this way the converged well is given the correct control/status for the current reservoir state. A command line argument –local-well-solve-control-switching=true has been added to trigger the use of the function (this is false by default).
\end{verbatim}

Improvements include:

\begin{verbatim}
- The correct version of flow is now selected for two-phase gas-water systems with either water vaporized in the gas phase (VAPWAT) and/or gas dissolved in the water phase (DISGASW) as defined in the RUNSPEC section ([#4592](https://github.com/OPM/opm-simulators/pull/4592)). Previously the wrong version of flow was selected for gas-water systems with water vaporized in the gas phase (but no gas dissolved in the water phase).
- The boundary conditions keyword BC has been deprecated and been split into two parts: [BCCON](#REF_HEADING_KEYWORD_BCCON) in the GRID section where the block indices and direction are set, and [BCPROP](#REF_HEADING_KEYWORD_BCPROP) in the SCHEDULE section where the type and value of the boundary condition are set ([#3482](https://github.com/OPM/opm-common/pull/3482), [#802](https://github.com/OPM/opm-models/pull/802) and [#949](https://github.com/OPM/opm-tests/pull/949)).
- The maximum number of allowable EQLNUM regions has been increased from 255 to 65525 ([#4726](https://github.com/OPM/opm-simulators/pull/4726)).
- A stricter convergence tolerance has been imposed on standard (not multi-segment) wells with a zero rate target. This change improves the results of a reported case with zero reservoir volume rate (RESV) control and potentially improves the material balance of the whole system ([#4572](https://github.com/OPM/opm-simulators/pull/4572)).
- When testing a gas lift well under THP control, if the well does not converge with the maximum artificial lift quantity (ALQ) then the simulator will now try to reduce ALQ in increments to check if the well equations converge with a smaller ALQ ([#4579](https://github.com/OPM/opm-simulators/pull/4579)).
- The bisection algorithm has been updated to keep track of the size of its search interval. The algorithm will now stop iterating if the minimum interval size is reached ([#4617](https://github.com/OPM/opm-simulators/pull/4617)).
- The simulator has been modified to ensure that THP constraints are only applied to prediction wells and not history match wells ([#4615](https://github.com/OPM/opm-simulators/pull/4615)).
- The simulator has been updated to avoid requesting the ALQ value for injectors, the ALQ value is only requested for producers ([#4648](https://github.com/OPM/opm-simulators/pull/4648)).
- The simulator has been modified to output temperature information to the restart file if TEMP is specified in the RPTRST keyword or if THERMAL is set in the RUNSPEC section even if the command line option --enable-opm-rst-file=false is set ([#4646](https://github.com/OPM/opm-simulators/pull/4646)). The command line option --enable-opm-rst-file is a Boolean value that controls the output of OPM specific data sets to the commercial simulator’s restart file to enable restart of OPM Flow runs by OPM Flow (true), or not to output the data (false).
- Update the calculation of immobile gas saturation to take into the account trapped saturation (calculated when the hysteresis model is used) in the calculation of mobile and immobile fluid volumes in place ([#4642](https://github.com/OPM/opm-simulators/pull/4642) and [#3517](https://github.com/OPM/opm-common/pull/3517)). The hysteresis model is activated by the specifying the HYSTER parameter of the SATOPTS keyword in the RUNSPEC section.
- An error message will now be generated if invalid region numbers SATNUM, PVTNUM, IMBNUM or EQLNUM are input ([#4705](https://github.com/OPM/opm-simulators/pull/4705)). Valid region numbers are positive integers less than or equal to the maximum number of regions specified in the RUNSPEC section.
- Added setVapPars() at the start of the report step to properly handle oil vaporization parameters (VAPPARS) updates in the SCHEDULE section ([#4677](https://github.com/OPM/opm-simulators/pull/4677)).
- If vaporised water (VAPWAT) is present in the model but the initial equilibrium vaporized water in gas ratio (RVW) for an equilibration region is not explicitly defined by a value versus depth table (RVWVD), and the datum depth is not at either the gas-oil contact (if oil is present) or the gas-water contact (if oil is not present) then the values of RVW will be initialised to zero ([#4647](https://github.com/OPM/opm-simulators/pull/4647) and [#4688](https://github.com/OPM/opm-simulators/pull/4688)). This makes it easier to include VAPWAT in CO2STORE cases.
- The mass balance limit (XXXMBE) from the keyword TUNING is now used if it is explicitly specified (not defaulted) and the command line argument --enable-tuning=true is used ([#4621](https://github.com/OPM/opm-simulators/pull/4621) and [#3522](https://github.com/OPM/opm-common/pull/3522)).
- The simulator now gives a more informative error message if the input grid has no active cells at all. Often this points to an error in the input data and this change might help in finding the problem ([#4735](https://github.com/OPM/opm-simulators/pull/4735)).
- A more user-friendly error message is now reported when the time step is cut too often or too much. ([#4746](https://github.com/OPM/opm-simulators/pull/4746)).
- Perform a more graceful exit instead of MPI_Abort for expected exceptions in parallel runs. Instead of unconditionally issuing MPI_Abort if a fatal exception is encountered, the simulator will try to test whether all processes have experienced this exception and if this is the case just terminate normally with an exit code that signals an error. MPI_Abort is still used if not all processes get an exception as this is the only way to ensure that the program aborts ([#4750](https://github.com/OPM/opm-simulators/pull/4750)). This approach also works around issues in some MPI implementations that might not correctly return the error.
- Hydrostatic and acceleration pressure losses have been included for well segments representing an inflow control device (ICD) ([#4824](https://github.com/OPM/opm-simulators/pull/4824)). Valves may be placed in long segments with significant depth differences (for example, at branch inlets), where the hydrostatic contribution in particular may be significant.
- An error message is now reported if a negative oil or water saturation is passed to the RV and RVW initialization routine ([#4675](https://github.com/OPM/opm-simulators/pull/4675)).
- When approximately zero well rates are encountered during iterations (pre-convergence), gas/water fractions become highly inaccurate which in turn may lead to the solver getting stuck or the well getting shut prematurely. The code has been updated to switch to explicit VFP table lookup whenever the rate drops below the lowest value in the table. The logic around explicit lookup is also updated so that this approach also works for just opened wells ([#4669](https://github.com/OPM/opm-simulators/pull/4669)).
- Information about linear system sizes is now output to the DBG file ([#4734](https://github.com/OPM/opm-simulators/pull/4734) and [#4754](https://github.com/OPM/opm-simulators/pull/4754)).
- For stopped or zero-rate-target wells the alternative_well_rate_init procedure previously returned unscaled well-rates resulting from a zero-bhp condition.  This could lead to convergence failures for network balancing since the initialized rates could be off by orders of magnitude. The code now skips this procedure for stopped or zero-rate-target wells reducing the risk of convergence failures ([#4817](https://github.com/OPM/opm-simulators/pull/4817)).
- The two-point flux-approximation (TPFA) has been added as a linearizer for gas-oil cases with energy (with diffusion) and gas-oil diffusion ([#4825](https://github.com/OPM/opm-simulators/pull/4825) and [#4816](https://github.com/OPM/opm-simulators/pull/4816)).  All energy cases now have diffusion enabled.
- Added RESTART file output for the geomechanical module ([#4803](https://github.com/OPM/opm-simulators/pull/4803) and [#4588](https://github.com/OPM/opm-simulators/pull/4588)).
- Relaxation factors slightly outside the interval [0, 1] in the standard well model are reset to the interval limits, while keeping the assertion for factors further outside the interval to possibly pick up failure cases ([#4862](https://github.com/OPM/opm-simulators/pull/4862)).
- Connections between reservoir cells and numerical aquifer cells, or between numerical aquifer cells when multiple such cells define a single numerical aquifer, are now always treated as NNCs for output purposes ([#4821](https://github.com/OPM/opm-simulators/pull/4821)).
- Regional transmissibility multipliers such as those entered in the MULTREGT keyword are now applied to explicit input non-neighbor connections (NNC). This now implements all known connection behaviors for inter-region connections. In the commercial simulator, if the region numbers specified in MULTREGT items 1 and 2 are equal and positive then the transmissibilities within the region as well as any transmissibilities connecting any other regions are multiplied by the given value. The option to specify transmissibility multipliers in this way using MULTREGT is not yet supported ([#4822](https://github.com/OPM/opm-simulators/pull/4822) and [#4821](https://github.com/OPM/opm-simulators/pull/4821)).
- The name of the missing OPM restart file is reported in the error message when it cannot be located ([#4870](https://github.com/OPM/opm-simulators/pull/4870)).
\end{verbatim}

\begin{longtable}[]{@{}l@{}}
\toprule\noalign{}
Error locating serialized restart file PATH/FILENAME.OPMRST \\
\midrule\noalign{}
\endhead
\bottomrule\noalign{}
\endlastfoot
\end{longtable}

\begin{verbatim}
- Added PRT file output when group economic criteria for production groups GECON is activated ([#4866](https://github.com/OPM/opm-simulators/pull/4866)).
- Reduced repetitive output of network information when running in parallel ([#4879](https://github.com/OPM/opm-simulators/pull/4879)).
- The lower limit for bottom hole pressure (BHP) in Newton updates has been reduced slightly from 1 bar. This allows for cases that might have a defaulted BHP constraint of 1 bar ([#4877](https://github.com/OPM/opm-simulators/pull/4877)).
- The region set name matching algorithm has been changed to using unique prefixes. This enables the simulator to recognize that the region set name FIPUNI should match the user defined region name FIPUNIT ([#4868](https://github.com/OPM/opm-simulators/pull/4868)).
- Debug output sent to the DBG file is no longer also echoed to the console as well ([#4955](https://github.com/OPM/opm-simulators/pull/4955) and [#4941](https://github.com/OPM/opm-simulators/pull/4941)).
- Updated the derivatives and also the Jacobian matrix for the multisegment well pressure equations ([#4640](https://github.com/OPM/opm-simulators/pull/4640)).
- The intensive quantities for timeIdx==1 corresponding to the start of the current timestep are now updated if we do not use the storage cache, or if we cannot recycle the first iteration storage ([#4662](https://github.com/OPM/opm-simulators/pull/4662)).
- Improved message issued when shutting a well because it fails to converge ([#4741](https://github.com/OPM/opm-simulators/pull/4741)).
- Damaris was updated in version 1.8.0 to support the HDF5 H5Sselect_elements() capability to rewrite data in memory to another order on disk. This allows (MPI decomposed) local simulation data to be written back to its original global position on disk. To support this a new element was added to the Damaris XML <variable...> type, named "select-file", along with some other options (not required by OPM Flow). This keyword is added to the basic in-built Damaris XML file of OPM Flow ([#4830](https://github.com/OPM/opm-simulators/pull/4830)).
\end{verbatim}

In addition, the following new SUMMARY keywords are now recognized as
described by the comments in Table A.7.

\begin{longtable}[]{@{}
  >{\raggedright\arraybackslash}p{(\linewidth - 4\tabcolsep) * \real{0.3333}}
  >{\raggedright\arraybackslash}p{(\linewidth - 4\tabcolsep) * \real{0.3333}}
  >{\raggedright\arraybackslash}p{(\linewidth - 4\tabcolsep) * \real{0.3333}}@{}}
\toprule\noalign{}
\begin{minipage}[b]{\linewidth}\raggedright
No.
\end{minipage} & \begin{minipage}[b]{\linewidth}\raggedright
Summary Keyword
\end{minipage} & \begin{minipage}[b]{\linewidth}\raggedright
Comment
\end{minipage} \\
\midrule\noalign{}
\endhead
\bottomrule\noalign{}
\endlastfoot
1 & WBP, WBP4, WBP5 and WBP9 & Added support for SUMMARY vectors to
output well block averaged pressures for open completions
(\href{https://github.com/OPM/opm-simulators/pull/4695}{\#4695},
\href{https://github.com/OPM/opm-simulators/pull/4694}{\#4694} and
\href{https://github.com/OPM/opm-simulators/pull/4693}{\#4693}). \\
2 & WINJFVR, WINJFVT and WINJFC & Added support for SUMMARY keywords
relating to filtration injection modelling
(\href{https://github.com/OPM/opm-simulators/pull/4790}{\#4790},
\href{https://github.com/OPM/opm-common/pull/3628}{\#3628} and
\href{https://github.com/OPM/opm-common/pull/3626}{\#3626}). Added
injection well SUMMARY vectors for filtrate volume injection rate
(volume rate), filtrate volume injection total (volume), and filtrate
injection concentration (ppm). \\
3 & CINJFVR, CINJFVT, CFCSKIN, CFCWIDTH, CFCPERM, CFCPORO, CFCRAD and
CFCAOF & Added support for SUMMARY keywords relating to filtration
injection modelling
(\href{https://github.com/OPM/opm-simulators/pull/4790}{\#4790},
\href{https://github.com/OPM/opm-common/pull/3628}{\#3628} and
\href{https://github.com/OPM/opm-common/pull/3626}{\#3626}). Added
injection well connection SUMMARY vectors for filtrate volume injection
rate (volume rate), filtrate volume injection total (volume), skin
factor due to filtration (dimensionless), thickness of filter cake
(Length), permeability of filter cake (MD), porosity of filter cake
(dimensionless), well bore radius used in the filtration modeling
(Length), and well bore area of flow used in the filtration modeling
(Area). \\
4 & SxDEN, SDENM and SMDEN & Added support for SUMMARY vectors to output
fluid densities at the block and well segment levels
(\href{https://github.com/OPM/opm-simulators/pull/4744}{\#4744},
\href{https://github.com/OPM/opm-common/pull/3594}{\#3594} and
\href{https://github.com/OPM/opm-common/pull/3593}{\#3593}). Added well
segment SUMMARY vectors for phase density of phase x (segment
conditions), fluid mixture density weighted by phase flowing fractions,
and fluid mixture density weighted by phase flowing fractions with
exponents. Phase x is one of O (oil), G (gas), or W (water). \\
5 & BxDEN and BDENx & Added support for SUMMARY vectors to output fluid
densities at the block and well segment levels
(\href{https://github.com/OPM/opm-simulators/pull/4744}{\#4744},
\href{https://github.com/OPM/opm-common/pull/3594}{\#3594} and
\href{https://github.com/OPM/opm-common/pull/3593}{\#3593}). Added block
SUMMARY vectors for phase density of phase x. Phase x is one of O (oil),
G (gas), or W (water). \\
6 & BFLOWI, BFLOWJ, and BFLOWK & Added support for SUMMARY vectors
BFLOWI, BFLOWJ and BFLOWK to request block oil/gas/water flow rates at
surface conditions
(\href{https://github.com/OPM/opm-simulators/pull/4867}{\#4867},
\href{https://github.com/OPM/opm-common/pull/3675}{\#3675} and
\href{https://github.com/OPM/opm-models/pull/827}{\#827}). This does not
require FLOWS to be specified in RPTRST. \\
\end{longtable}

\emph{Table A.7: New SUMMARY Keywords for the 2023-10 Release}

The following bug fixes and improvements have been incorporated into
this release.

\begin{verbatim}
- Fixed the issue of incorrectly using the gas-oil contact depth for initialization of vaporized water (RVW) in the special case of a two-phase gas-water system. The gas-water contact depth is now used in this case ([#4688](https://github.com/OPM/opm-simulators/pull/4688) and [#4647](https://github.com/OPM/opm-simulators/pull/4647)).
- Fixed a bug in the group pressure maintenance (GPMAINT) code ([#4664](https://github.com/OPM/opm-simulators/pull/4664)). Previously, if negative group pressure maintenance rates were calculated, they would first be used to incorrectly update the GPMAINT error integral, before the rates were later set to zero. Now the simulator only calculates rates for injectors if the pressure is below the regional pressure target (or for producers if the pressure is above the regional pressure target).
- Bug fixes related to the handling of gas dissolved in water ratio ([RSW](#REF_HEADING_KEYWORD_RSW_10_3)) and water vaporized in gas ratio (RVW) in the standard well model ([#4591](https://github.com/OPM/opm-simulators/pull/4591)).
- The max size of the next time step is now only used when specified in TUNING or NEXTSTEP ([#4660](https://github.com/OPM/opm-simulators/pull/4660)).
- If the data file contains the CPR keyword then the "CPR" preconditioner should be used unless it is overridden by a command line argument. Previously this was incorrectly overridden by the command line argument --linear-solver-max-iter, the correct command line argument --linear-solver is now used ([#4700](https://github.com/OPM/opm-simulators/pull/4700)).
- The user is now able to specify any number of threads, this prevents only one thread being used on some hardware where the number of processors would always be reported as one when using MPI and OpenMP (irrespective of the actual number). The behaviour is now: (a) if nothing is specified then 2 threads are used, (b) if OMP_NUM_THREADS is specified then this number is used regardless of the hardware and the command line argument --threads-per-process is ignored, and (c) if --threads-per-process is used and OMP_NUM_THREADS is not set then the number specified on the command line is used ([#4709](https://github.com/OPM/opm-simulators/pull/4709)). There is a check to make sure that the number of threads used by the linearizer is consistent ([#811](https://github.com/OPM/opm-models/pull/811)).
- Fixed a bug related to indexing in the temperature boundary condition ([#4761](https://github.com/OPM/opm-simulators/pull/4761)).
- Gas dissolved in oil (RS) and oil vaporized in gas (RV) ratio initialized using value versus depth tables (RSVD, RVVD) should be limited by their respective saturated values. Fixed a bug where this limit was not applied at depths outside the range of depths in the RSVD or RVVD table ([#4723](https://github.com/OPM/opm-simulators/pull/4723)).
- For producers where all perforations have zero rates the perforation mixture fraction is approximated using the (inverse formation volume factor times mobility) ratio, and weight the perforation rates using the well transmissibility. The perforation mixture fraction was previously approximated using only the mobility ratio ([#4681](https://github.com/OPM/opm-simulators/pull/4681)).
- A case under investigation has some cells containing only oil and water (zero gas saturation), which have zero relative permeability (and mobility) for all the three phases. This caused problems in part of the code which assumed the total mobility was non-zero. For perforations having zero mobility for all the phases, the simulator now uses a small value to generate small perforation rates for those perforations, at the same time, the simulator can use these rates to recover the mixing ratios for those perforations ([#4682](https://github.com/OPM/opm-simulators/pull/4682) and [#4681](https://github.com/OPM/opm-simulators/pull/4681)).
- The correct formation volume factors are now used in the case of zero phase rates resulting in zero RS, RV, [RSW](#REF_HEADING_KEYWORD_RSW_10_3) or RVW. Previously the saturated formation volume factor was used in these cases ([#4590](https://github.com/OPM/opm-simulators/pull/4590)).
- The reservoir volume rate constraint (RESV) in GCONPROD was not honoured. This has been fixed ([#4687](https://github.com/OPM/opm-simulators/pull/4687)).
- Previously, in the case of zero threshold pressure and zero pressure difference, the code would set the pressure difference explicitly to zero. This would also set any derivatives to zero, which could disconnect the corresponding matrix rows. The code no longer sets the pressure difference to zero when not necessary ([#4701](https://github.com/OPM/opm-simulators/pull/4701)).
- Code has been updated to avoid dividing by zero when scaling the well rates in updateWellStateRates() ([#4715](https://github.com/OPM/opm-simulators/pull/4715) and [#4649](https://github.com/OPM/opm-simulators/pull/4649)).
- The code has been updated to avoid round off errors in phase saturations leading to the initialization of dissolved or vaporized fluid ratios (for example saturated RS) with saturated values instead of taking values from the input ratio versus depth table (for example RSVD) ([#4720](https://github.com/OPM/opm-simulators/pull/4720)).
- Code has been updated to avoid writing beyond array limits ([#4753](https://github.com/OPM/opm-simulators/pull/4753)).
- The well closure reason is now set to GROUP rather than ECONOMIC if a group economic constraint (GECON) is reached. This prevents the closed well being reopened with WTEST item 3 TEST equal to ‘E’ ([#4854](https://github.com/OPM/opm-simulators/pull/4854)).
- The simulator now sets the well THP to be zero in the WellState only if there is no VFP table associated with the well ([#4932](https://github.com/OPM/opm-simulators/pull/4932)). Previously, this would be done if the well had no target THP or THP limit.
- The code now checks whether LIFTOPT is active first to avoid unnecessarily running routines relating to gas lift optimization ([#4956](https://github.com/OPM/opm-simulators/pull/4956)).
- MPI is now initialized before creating the communicator and a couple of fields are now set to ROOT_ONLY in opmrst_inspect() ([#4601](https://github.com/OPM/opm-simulators/pull/4601)).
- The updateWaterMobilityWithPolymer() method updates values, it does not rewrite them. Hence we have to feed it the scalar values not zeros. Scalar values are now used ([#4655](https://github.com/OPM/opm-simulators/pull/4655)).
- The restriction on the length of the next timestep following an event will now be applied following production or injection updates (WCONPROD, WCONINJE keywords) ([#4781](https://github.com/OPM/opm-simulators/pull/4781)). The maximum length of the next timestep following an event can be set using either the TUNING keyword item 10 (TMAXWC) with the command line argument --enable-tuning=true, or by using the command line argument --time-step-after-event-in-days=x (where x is the number of days).
- Added perforation data (perf_data) comparison in the equality operator for SingleWellState ([#4783](https://github.com/OPM/opm-simulators/pull/4783)).
- Added check that rock compaction table indices (ROCKNUM) are within the bounds of the tables defined by ROCKTAB ([#4788](https://github.com/OPM/opm-simulators/pull/4788)).
- Code modified to avoid a possible segmentation fault in the cleanup routine if the simulator has not been set up ([#4794](https://github.com/OPM/opm-simulators/pull/4794)).
- Fixed a bug where the code could attempt to access entries in a zero element array due to not dereferencing a pointer to inspect array elements when deciding whether to apply TRANX, TRANY, etc. keywords ([#4801](https://github.com/OPM/opm-simulators/pull/4801)).
- Code has been modified to prevent slightly negative oil fractions in well segments occurring due to round-off. This could have lead to failure if the segment represented for example an autonomous Inflow Control Device (ICD) valve in a multi-segment well ([#4834](https://github.com/OPM/opm-simulators/pull/4834)).
- Fixed a bug with the DIRICHLET option in the boundary condition property definition ([BCPROP](#REF_HEADING_KEYWORD_BCPROP)) not working in gas-water runs. Additionally, fixed a bug where [BCPROP](#REF_HEADING_KEYWORD_BCPROP) had to be defined at every report step ([#4835](https://github.com/OPM/opm-simulators/pull/4835)).
- If a well was SHUT due to economic or physical reasons the well-state quantities were mostly set to zero (including the BHP). The code has been modified to set the BHP to the BHP limit when its value has not been initialized when under well target rate control ([#4841](https://github.com/OPM/opm-simulators/pull/4841)). This helps to prevent the Newton update from stagnating under certain circumstances (for example, a reported issue introduced by [#4772](https://github.com/OPM/opm-simulators/pull/4772)).
- Fixed a bug which occurred when using non-neighbor connections (NNC) in thermal simulations ([#4900](https://github.com/OPM/opm-simulators/pull/4900)).
- Fixed computation of temperature for distributed wells ([#4888](https://github.com/OPM/opm-simulators/pull/4888)). This was not an issue for non-distributed wells.
- Code has been modified to always allocate buffers for storing flows if BFLOWI, BFLOWJ or BFLOWK is requested. This fixes a segmentation fault that occurred if these keywords were requested without RPTRST in the SOLUTION section ([#4904](https://github.com/OPM/opm-simulators/pull/4904)).
- Code has been modified so that there will be no vertical connection between cells if the layers in between are inactive or collapsed unless the PINCH keyword has been specified ([#4901](https://github.com/OPM/opm-simulators/pull/4901)).
- Fixed a bug which occurred in parallel runs where all processes other than rank zero would keep iterating beyond the maximum number of allowed iterations ([#4909](https://github.com/OPM/opm-simulators/pull/4909)).
- Fixed a bug where regional multipliers (MULTREGT) were applied twice to non-neighbor connections entered in EDITNNC ([#4921](https://github.com/OPM/opm-simulators/pull/4921)).
- Fixed a bug where if --enable-vtk-output=true and --enable-write-all-solutions=false (it is false by default), then Visualization Toolkit (VTK) files were written for all solutions. The code has been modified so that VTK files are only written for report steps (unless –enable-write-all-solutions=true) ([#4916](https://github.com/OPM/opm-simulators/pull/4916)).
- After DUNE version 2.7, the method getCollectiveCommunication() was changed to getCommunication(). The getCollectiveCommunication() method issues the deprecation warning and assumes all process take part in the communication. However, this might not be the case, namely, only a subset might be used. This fix consists of getting the communication object used by the grid as this uses only processes that are actually computing and not only outputting ([#4962](https://github.com/OPM/opm-simulators/pull/4962)).
- Fixed computation of dissolved gas in water ratio (Rsw) used in calculation of wellbore density for standard wells ([#4976](https://github.com/OPM/opm-simulators/pull/4976)).


- There is an error with the RSM header for summary vectors whose NUMS entry in the SMSPEC file is derived from more than a single number source (e.g., single region or segment ID). This applies to all block vectors (BGPV, BOPV, BWPV, etc.), connection level quantities (COPT, etc.), and inter-region flows such as ROFT etc ([#3078](https://github.com/OPM/opm-common/issues/3078)). The work around is to plot the data in OPM ResInsight and right-click on the plot to view and copy the data.
- OPM Flow does not support using LIQ as a well's preferred phase with the WELSPECS. keyword, that is WELSPECS(TYPE) equals LIQ. This is a long-standing bug/omission in the simulator stemming from a somewhat naive internal notion of phases so we don't have an entry for a liquid phase, only for the distinct oil and water phases. For producing wells this mostly matters if you plot the WPI summary vector (productivity index for well's preferred phase). In the current treatment WPI will not have contributions from the water phase if the declared preferred phase is LIQ. For injecting wells WELSPECS's preferred phase doesn't really matter at all since the preferred phase is (typically) reset to the injected phase in WCONINJE/WCONINJH anyway. See issue [#3075](https://github.com/OPM/opm-common/issues/3075).
- If the simulator finds well connection being declared as connections via the COMPDAT keyword in the SCHEDULE section, then it writes out a warning message:
\end{verbatim}

\begin{longtable}[]{@{}
  >{\raggedright\arraybackslash}p{(\linewidth - 0\tabcolsep) * \real{1.0000}}@{}}
\toprule\noalign{}
\begin{minipage}[b]{\linewidth}\raggedright
Warning: Problem with COMPDAT keyword In SIM\_PEE\_Basic\_Simulator.DATA
line 535 The cell (0,6,4) in well INJ1 is not active and the connection
will be ignored Warning: Problem with COMPDAT keyword In
SIM\_PEE\_Basic\_Simulator.DATA line 535 The cell (16,0,4) in well INJ2
is not active and the connection will be ignored
\end{minipage} \\
\midrule\noalign{}
\endhead
\bottomrule\noalign{}
\endlastfoot
\end{longtable}

However, the reported cell references are offset by minus one, meaning
the correct warning messages should be:

\begin{longtable}[]{@{}
  >{\raggedright\arraybackslash}p{(\linewidth - 0\tabcolsep) * \real{1.0000}}@{}}
\toprule\noalign{}
\begin{minipage}[b]{\linewidth}\raggedright
Warning: Problem with COMPDAT keyword In SIM\_PEE\_Basic\_Simulator.DATA
line 535 The cell (1,7,5) in well INJ1 is not active and the connection
will be ignored Warning: Problem with COMPDAT keyword In
SIM\_PEE\_Basic\_Simulator.DATA line 535 The cell (17,1,5) in well INJ2
is not active and the connection will be ignored
\end{minipage} \\
\midrule\noalign{}
\endhead
\bottomrule\noalign{}
\endlastfoot
\end{longtable}

See issue \href{https://github.com/OPM/opm-common/issues/3167}{\#3167}
for details.

\begin{verbatim}
- The GDFILE keyword in the GRID section loads a grid file in various formats, with the FMTOPT parameter setting the format type of the file. If the variable FMTOPT is omitted then the default is for binary file input for the commercial simulator; whereas, OPM Flow derives FMTOPT from the file extension (*.EGRID or *.FEGRID), making FMTOPT superfluous. However, if the extension is lower case then OPM Flow may incorrectly determine the file type. The work around is change the extension to upper case.
- As per previous releases of the radial model, the COORDSYS keyword item three must be set to COMP to complete the circle, this has not been implemented in this release. Also there appears to be a bug for full radial models when a well goes on BHP control that causes the well not to respect the BHP constraint, this eventually causes the well to die prematurely. See [#2640](https://github.com/OPM/opm-common/pull/2640) for a discussion on the topic.
- As in previous releases there are some issues with the OPERATE and OPERATER keywords associated with the input parsing; for various reasons a few of the fields require special case treatment in the grid processing, including (at least) MULTZ, PORV and ACTNUM, and for those keywords the OPERATE/OPERATER keyword doe not work. The work around is to use the MULTIPLY keyword instead.
- For the UDQ ASSIGN operator after the terminating “/” normally any comments can be entered; however, if there is “/” within the comment field, as per:
\end{verbatim}

\begin{verbatim}
ASSIGN FUNGLYLD 1.196   /    Condensate Yield (63.5/56.7)/(1.0 – 0.065)
\end{verbatim}

then the simulator will abort. The work around is to manually place the
comment characters ``--'' after the ASSIGN terminating ``/'', like so:

\begin{verbatim}
ASSIGN FUNGLYLD 1.196   / -- Condensate Yield (63.5/56.7)/(1.0 – 0.06)
\end{verbatim}

\begin{verbatim}
- At the moment, one cannot initialize tracers using the EQUALS keyword. Instead use the array format, that is the keyword followed by the required number of values, or the TVDP keyword in the SOLUTION section to set the initial tracer concentrations as a function of depth.
- Currently, gas tracers cannot be used if the dissolved gas phase, as per the DISGAS keyword in the RUNSPEC section, is active in the model.
- The summary vector RTIPTHEA, that defines the energy in-place between the initial and the current time for regions, is not supported unlike the FTIPTHEA and BTIPTHEA vectors. Secondly, the error message:
\end{verbatim}

\begin{longtable}[]{@{}
  >{\raggedright\arraybackslash}p{(\linewidth - 0\tabcolsep) * \real{1.0000}}@{}}
\toprule\noalign{}
\begin{minipage}[b]{\linewidth}\raggedright
Warning: Problem with summary keyword RTIPTHEA In RSM-THERMAL.data line
492 FIP region FIPHEA not defined in REGIONS section - RTIPTHEA ignored
\end{minipage} \\
\midrule\noalign{}
\endhead
\bottomrule\noalign{}
\endlastfoot
\end{longtable}

is incorrect, as the message indicates that it is being treated like a
named region, as per the FIP keyword, when it is actually a SUMMARY
vector
(\href{https://github.com/OPM/opm-simulators/issues/3870}{\#3870}).

\begin{verbatim}
- If there are cells that are very distorted, which can occur near fault planes, then the simulator may abort because it cannot calculate the pore volume of such cells. The work around is to re-generate the grid in the static model, taking care that the cells around the fault planes are more or less orthogonal ([#2992](https://github.com/OPM/opm-common/issues/2992) and [#3770](https://github.com/OPM/opm-simulators/issues/3770)).
- Currently the OPERATER keyword in the EDIT section does not work with the DEPTH, TRANX, TRANY and TRANZ property arrays ([#2994](https://github.com/OPM/opm-common/issues/2994) and [#748](https://github.com/OPM/opm-tests/pull/748)).
- If a standard well is fully declared in an ACTIONX block which is then activated at a later date, and later the well is modified to be a multi-segment well using the WELSEGS and COMPSEGS keywords, then this will cause the simulator to abort with an assert failure. The solution to this issue is to not use this type of work flow in declaring wells ([#2891](https://github.com/OPM/opm-common/issues/2891) and [#2895](https://github.com/OPM/opm-common/pull/2895)).
- Although the ACTIONX EXIT command works as expected, it does not write out the requested RSM file at the end of the run. However, the other SUMMARY and RESTART files are written out ([#2877](https://github.com/OPM/opm-common/issues/2877)).
- Although the GCONSUMP keyword in the SCHEDULE section is fully implemented as documented, it is not possible to verify the output as the associated SUMMARY vectors are not written out, that is the SUMMARY sales gas vectors FGSR, FGST, GGSR and GGST, and fuel vectors FGCR, FGCT, GGCR, and GGCT have not been implemented ([#2679](https://github.com/OPM/opm-common/issues/2679)).
- There are small differences in the behavior of the NEXTSTEP keyword in the RUNSPEC section between OPM Flow and the commercial simulator that remain unresolved ([#3745](https://github.com/OPM/opm-simulators/issues/3745)).
- There is a unit handling issue associated with OPERATE keyword. If the OPERATE(X) parameter has units, as for example PERMX, then the conversion is always done in SI units, despite the input deck declaring the deck to be fields units, as per FIELD keyword in the RUNSPEC section. Note that OPM Flow performs all of its calculations internally in SI and performs unit handling only when inputting the *.DATA file and when outputting result files. Thanks to [lrijkels](https://github.com/lrijkels) for reporting the issue. See [#4597](https://github.com/OPM/opm-simulators/issues/4597) for details.
- There is an issue associated with restarting from a restart file with the solution gas (Rs) maximum rate of increase, as defined by the DRSDT keyword in the SCHEDULE, that has been set to zero. This is because, the simulator does not save/restore this setting in simulator's restart files, which means that simulator misses the essential value zero upon restarting the case. As a work-around one can use the option --sched-restart=true, when running the restart case. This will initialize the restarted simulation based on information from the complete SCHEDULE section, instead of just the parts that we're going to simulate and the rest from the restart file. Thanks to [goncalvesmachadoc](https://github.com/goncalvesmachadoc) for reporting the issue. See [#4272](https://github.com/OPM/opm-simulators/issues/4272) for details.
- OPM cannot be built with dune-fem version 2.9 or later. Please use a previous version.
- The simulator uses an irregular corner-point grid geometry with adjusted pore volumes to represent radial grids so it is not possible to create a full ring (360 degree disk) with only one cell in the theta direction (NY=1). The work around is to model a slice (say DTHETA=60 degrees). Note that as the angle increases larger pore volume adjustments are required ([#4755](https://github.com/OPM/opm-simulators/issues/4755)).
- In principle the PYACTION code can control anything, however the simulator generally deals poorly with wells not explicitly controlled by the input file (for example opening or closing wells). It is recommended to utilize the normal ACTIONX machinery for well controls by constructing an ACTIONX block in the Python code and then passing that back to the Schedule object ([#4810](https://github.com/OPM/opm-simulators/issues/4810) and [#4813](https://github.com/OPM/opm-simulators/issues/4813)).
- If the --save-step command line option is used to request the serialized state is saved to an .OPMRST file and HDF5 support is missing then an error is generated at the report step where the save was requested. This error should be generated during startup to save time and resources ([#4812](https://github.com/OPM/opm-simulators/issues/4812)).
- In some cases with the network option the simulator can wrongly report that a well has no THP constraints, for example
\end{verbatim}

\begin{longtable}[]{@{}l@{}}
\toprule\noalign{}
GLIFT WTEST: Well S-3H does not have THP constraints \\
\midrule\noalign{}
\endhead
\bottomrule\noalign{}
\endlastfoot
\end{longtable}

when THP constraints have been defined
(\href{https://github.com/OPM/opm-simulators/issues/4887}{\#4887}).

\chapter{OPMRUN}\label{opmrun}

OPMRUN is graphical user interface to Flow that has similar
functionality to the commercial simulator's ECLRUN program. Target
audience are Reservoir Engineers in a production environment. Developers
and experienced Linux users will already have compatible work flows. The
application performs the following:

\begin{itemize}
\tightlist
\item
  Allows editing and management of OPM Flow's run time parameters.
  Default parameters are automatically loaded from OPM Flow, and the
  user can reset the default parameter set either from a parameter or
  PRT file. Editing of a job's parameter file is also available.
\item
  Runs under Linux and Windows 10. For Windows 10 OPM Flow is run via
  the Windows Subsystem for Linux (``WSL'').
\item
  Allows simulation jobs to be queued and run in either foreground
  (under OPMRUN), or in background in a xterm terminal session in Lunux
  or WSL under Windows 10. Jobs in the queue can be set to run in NOSIM
  mode or RUN mode.
\item
  Foreground jobs can be killed from OPMRUN, with the option of killing
  all the jobs in the queue.
\item
  Queues can be edited, saved and loaded.
\end{itemize}

Various additional simulation input generation and conversion utilities
are available including:

\begin{itemize}
\tightlist
\item
  Compressing a job to save space (DATA, and all OPM Flow output files)
  and uncompressing previously compressed jobs,
\item
  Keywords, a keyword generator based on the Apache Velocity Template
  Language (``VTL''). The templates can therefore also be used with any
  editor that supports VTL, jEdit for example. There is one template per
  keyword, with the formatting the same as the OPM Flow manual. Over 450
  templates are currently implemented. One can also customize the
  existing templates as well as creating User defined templates. The
  keywords are examples, one still has to edit the resulting deck with
  the actual required data, but the format with comments should make
  this a straight forward process.
\item
  A Production Schedule application that takes a comma delimited CSV
  file containing historical production and injection data and converts
  the data to an OPM Flow SCHEDULE file using the WCONHIST series of
  keywords. Currently only production data is supported.
\item
  Sensitivities application that generates sensitivity cases based on a
  ``Base'' case file. The Base file contains ``Factors'' (variable
  names), \$X01, \$X02, etc., that are substituted with user defined
  values using the data entered and the type of Sensitivity Scenario
  selected.
\item
  A Well Specification application that uses the standard well export
  files from OPM ResInsight to reformat the data in a more user-friendly
  manner for the WELSPECS and COMPDAT keywords. Optionally, the
  application can generate the COMPLUMP keyword based on the OPM
  ResInsight layers file, with one completion per defined reservoir
  layer.
\item
  Calling OPM ResInsight and loading the currently selected job into OPM
  ResInsight for viewing.
\item
  A Well Trajectory Conversion application that converts a Schlumberger
  Petrel exported well trajectory file into an OPM ResInsight file,
  containing all the wells.
\end{itemize}

The software can be downloaded from the following link:

\url{https://github.com/OPM/opm-utilities/tree/master/opmrun}

OPMRUN is written in Python 3 and tested under various Ubuntu
distributions. Note that only Python 3 is supported and tested, Python2
support has been deprecated.

The program requires the following standard module Python libraries:

\begin{itemize}
\tightlist
\item
  datetime, getpass, importlib, os, numpy, pkg\_resources, pandas,
  pathlib, platform, psutil, sys, re, subprocess, and tkinter as tk.
\end{itemize}

In addition, the following non-standard Python modules are required:

\begin{itemize}
\tightlist
\item
  airspeed, notify-py, pyDOE2, and PySimpleGUI.
\end{itemize}

For some Linux systems the relevant package manager may have to be used
to install tkinter as tk; whereas for Windows 10 users the tkinter
package is pre-installed with Python.

OPMRUN enables the editing and management of OPM Flow's run time
parameters, setting up job queues to run a series of simulation jobs
sequentially, as well as the management of the job queues. Figure
C.1shows the initial display.

\pandocbounded{\includegraphics[keepaspectratio]{appendices/images/Frame120_f5dfac27d930.png}}
\pandocbounded{\includegraphics[keepaspectratio]{appendices/images/Image120_f5dfac27d930.png}}

Upon launch the program runs OPM Flow to get a list of command line
parameters from the current version of OPM Flow. These default
parameters can be edited for each case, or alternative default parameter
sets can be loaded from an existing parameter file from another job, or
a *.PRT file from a completed simulation.

As can be seen in Figure C.2 the program has upper and lower display
elements. The upper element shows a list of simulation jobs that are in
the job queue and the lower element consists of two elements, one for
the OPM Flow Output (the terminal output from OPM Flow) and a second
element (OPM Run Log) that is a session log of the jobs run by OPMRUN.
Clicking the OPM Flow Output and OPM Run Log tabs switches the display
on the lowered element between two display types.

To add jobs to the queue use the Add Job button or load an existing job
queue using the Load Queue button. Jobs can be edited or deleted from
the queue using the Edit Job and Delete Job buttons, and a series of
jobs can be saved as a job queue by using the Save Queue button. The
Clear Queue button deletes all jobs from the queue.

Pressing the Add Job button will display the following dialogue box:

\pandocbounded{\includegraphics[keepaspectratio]{appendices/images/Frame121_557019afb7dc.png}}
\pandocbounded{\includegraphics[keepaspectratio]{appendices/images/Image121_557019afb7dc.png}}

Use the Browse button to select the input file to add to the queue, then
select the Run Parameters for this input file, then press the Submit
button to add the input file to the job queue (Figure C.4).

The Overwrite Parameter File Options allow for different default
treatments of existing *.PARAM files, which is particularly useful when
adding multiple jobs at the same time.

\pandocbounded{\includegraphics[keepaspectratio]{appendices/images/Frame75_c76ed472a410.png}}
\pandocbounded{\includegraphics[keepaspectratio]{appendices/images/Image68_c76ed472a410.png}}

The reason for this is because different versions of OPM Flow have
different parameter sets and if a newer version of OPM Flow runs with a
previous version's *.PARAM file then the simulator will stop with an
error message for the various invalid parameters for the current version
of the simulator.

Jobs in the queue can be edited by selecting the Edit Job button that
will display two options (Figure C.5): one to edit the input file using
the defined editor and the second to edit the OPM Flow Parameter File.

\pandocbounded{\includegraphics[keepaspectratio]{appendices/images/Frame76_2483217ee1aa.png}}
\pandocbounded{\includegraphics[keepaspectratio]{appendices/images/Image69_2483217ee1aa.png}}

If the second option is selected OPMRUN will display a dialog box that
shows a list of the OPM Flow command line parameters together with the
parameter help information Figure C.7). Selecting a parameter from the
list and selecting the Edit button will display the setting for the
selected parameter (alternatively one can double click the required
entry). One can then edit the parameter as required. Use the Save button
to save the change and use the Exit button to save all the changes to
the parameter file. The Cancel button will cancel all changes to the
parameter file.

\pandocbounded{\includegraphics[keepaspectratio]{appendices/images/Frame122_30b2cadad0e0.png}}
\pandocbounded{\includegraphics[keepaspectratio]{appendices/images/Image122_30b2cadad0e0.png}}

Alternatively one can use the:

\begin{verbatim}
- Edit OPM Flow Parameter menu option to edit the parameter file for a job.
- List OPM Flow Parameters menu option to list the commands in the parameter file for a job.
- Set OPM Flow Default Parameters to set the default parameters for all subsequent jobs added to the queue. This option allows the user to load a default set of parameters from (1) OPM Flow, (2) an OPM Flow Parameter File, or (3) an OPM Flow print file (*.PRT).
\end{verbatim}

One can also right-click on a job and select one of the available
options.

To load a previously saved job queue, press the Load Queue button this
will display a dialog box allowing the user to select a queue file
(*.que), after pressing the OK button the jobs will be displayed in the
Job List Element as illustrated in Figure C.7.

\pandocbounded{\includegraphics[keepaspectratio]{appendices/images/Frame77_0d01d65c6763.png}}
\pandocbounded{\includegraphics[keepaspectratio]{appendices/images/Image70_0d01d65c6763.png}}

Queue files allow for various jobs to be load efficiently, especially
for ensemble and sensitivity cases and may contain a large number of
cases.

\begin{tcolorbox}[enhanced jigsaw, title=\textcolor{quarto-callout-note-color}{\faInfo}\hspace{0.5em}{Note}, coltitle=black, breakable, toprule=.15mm, colbacktitle=quarto-callout-note-color!10!white, leftrule=.75mm, opacitybacktitle=0.6, colback=white, bottomtitle=1mm, toptitle=1mm, titlerule=0mm, left=2mm, colframe=quarto-callout-note-color-frame, arc=.35mm, rightrule=.15mm, bottomrule=.15mm, opacityback=0]

When running under Windows 10 the Job names will follow the Windows 10
file naming convention and OPMRUN automatically handles the file names
for running the jobs under the Window Subsystem for Linux. However, care
is needed for any ``included'' file in the input deck. In this case the
PATHS -- Define Filename Directory Path Aliases keyword in the RUNSPEC
section may be of use.

\end{tcolorbox}

Reset Job Queue Parameters allows jobs run under Windows 10 WSL to be
renamed for running under Linux, and changing jobs from serial to
parallel and vice versa.

\pandocbounded{\includegraphics[keepaspectratio]{appendices/images/Frame78_eaf810db0798.png}}
\pandocbounded{\includegraphics[keepaspectratio]{appendices/images/Image71_eaf810db0798.png}}

The application will attempt to list or Linux mount points and Windows
drives depending on the host operating system, once the two systems
mount and drive points have been selected then the files in the queue
will be renamed from the previous host system to the current host
system. This is only performed for the *.DATA files, included files in
the input deck are currently not converted.

\pandocbounded{\includegraphics[keepaspectratio]{appendices/images/Frame123_29f24db28c27.png}}
\pandocbounded{\includegraphics[keepaspectratio]{appendices/images/Image123_29f24db28c27.png}}

As previously mentioned, one can also change the queue run time
parameters from serial to parallel or vice versa.

Selecting the Run Jobs button displays the Select Run Option dialog box
shown in Figure C.10 and Figure C.11.

\pandocbounded{\includegraphics[keepaspectratio]{appendices/images/Frame113_601e0d9ae485.png}}
\pandocbounded{\includegraphics[keepaspectratio]{appendices/images/Image72_601e0d9ae485.png}}

On the Select Run Dialog, the Run in No Simulation Mode option is
equivalent to setting the NOSIM option in the input deck for all jobs in
the queue (see section 5.3.96 NOSIM -- Activate the No Simulation Mode
for Data File Checking and the --enable-dry-run command line parameter
in Table 2.1 in section 2.2 Running OPM Flow From The Command Line).
This allows for checking all the jobs at once.

\pandocbounded{\includegraphics[keepaspectratio]{appendices/images/Frame24_e58f1a43b3ac.png}}
\pandocbounded{\includegraphics[keepaspectratio]{appendices/images/Image22_e58f1a43b3ac.png}}

Selecting Run in Standard Simulation Mode will run all the jobs in the
queue sequentially, with the OPM Flow terminal output directed to OPM
Flow Output Element, as shown in Figure C.12. The terminal output is
also directed to a *.LOG file as well, similar to what the commercial
simulator does.

\pandocbounded{\includegraphics[keepaspectratio]{appendices/images/Frame114_7d770e7418a7.png}}
\pandocbounded{\includegraphics[keepaspectratio]{appendices/images/Image73_7d770e7418a7.png}}

Clicking the OPM Run Log tab displays the OPMRUN's session log file that
records the time and date of the major events that have occurred,
including the start and end times of each run. Notice also how OPMRUN
deletes all the existing output files for a given job, if they exists,
before running OPM Flow. as well as creating a Schedule Log for tracking
progress (Figure C.13).

\pandocbounded{\includegraphics[keepaspectratio]{appendices/images/Frame115_6dee5c16282c.png}}
\pandocbounded{\includegraphics[keepaspectratio]{appendices/images/Image74_6dee5c16282c.png}}

The Kill button will ask the user if the current running job should be
killed, and if the job is to killed, the application will prompt as
whether or not all the jobs in the queue should be killed.

The Clear clears the OPM Flow Output Element from the currently
displayed tab (Output or Log) and the Copy button copies the data to the
clipboard.

\subsubsection{File Menu Options}\label{file-menu-options}

Enables open and saving the job queue, switching projects and listing
OPMRUN's user properties.

\subsubsection{Edit Menu Options}\label{edit-menu-options}

Lets one add jobs, add jobs recursively (all jobs in the selected
directory and below), edit the data file (Figure C.14) and the parameter
file for the selected job,(see section C.2.2 Edit Job Data and Parameter
File) edit, list and set the default parameters for running jobs that
will be added to the queue, set OPMRUN options, and set the project's
project directories.

\pandocbounded{\includegraphics[keepaspectratio]{appendices/images/Frame116_d38791aed54c.png}}
\pandocbounded{\includegraphics[keepaspectratio]{appendices/images/Image77_d38791aed54c.png}}
\pandocbounded{\includegraphics[keepaspectratio]{appendices/images/Image76_5bb837d10094.png}}

The options are also available by right-clicking a job in the Job List
Element

The Edit Parameters, List Parameters and Set Parameters relate to the
default parameter set, not the parameter set for a particular job. When
a new job is added to the queue the application checks if a *.PARAM file
exists for the job, if not then the default parameter set is used for
the job. Editing the default parameter set is the same as editing a job
parameter set (see section C.2.2 Edit Job Data and Parameter File for
more information).

One can also define the default parameter set by using the OPM Flow
default values, which is the set created when OPMRUN is first
initialized. In addition, one can load an existing parameter set from an
existing \emph{.PARAM file or from an OPM Flow }.PRT file (Figure C.15).

\pandocbounded{\includegraphics[keepaspectratio]{appendices/images/Frame118_c77302b21bac.png}}
\pandocbounded{\includegraphics[keepaspectratio]{appendices/images/Image79_c77302b21bac.png}}

OPMRUN has several configuration options that can be set via the
Edit/Options menu option as illustrated in Figure C.16. These include
setting:

\begin{itemize}
\tightlist
\item
  The location of the OPM Flow manual.
\item
  The Keyword Generator Template Directory, one of tools supplied with
  OPMRUN, see section C.3.2 OPMRUN Tools: Simulator Input/Keywords for
  further information on this application.
\item
  The location of the ResInsight program, for loading the results of a
  simulation run for viewing.
\item
  The editor command to used to edit the input deck and view the
  resulting simulator output files.
\item
  The terminal console to be be used for background jobs. WSL (``Windows
  Subsystem for Linux) should be selected if OPMRUN is running under
  Windows 10 to enable jobs to be submitted to the installed Linux
  distribution.
\item
  The User Information series of fields are used by various supplied
  tools and in some templates used by the Keyword Generator application
  (see C.3.2 OPMRUN Tools: Simulator Input/Keywords for more information
  on this application). Note if a User Information'' field is not
  defined then the template variable will be output instead -- this can
  easily be deleted in the application.
\item
  One can also define the main OPMRUN windows configuration parameters
  define: input (Job List Element) and output panel's (OPM Flow Output
  Element) size, font and font size.
\end{itemize}

\pandocbounded{\includegraphics[keepaspectratio]{appendices/images/Frame119_19e55a063cf3.png}}
\pandocbounded{\includegraphics[keepaspectratio]{appendices/images/Image80_19e55a063cf3.png}}

In addition to the aforementioned options, the Edit/Projects menu item
enables the setting of project directories that allows the user to set a
default directory for loading and saving files within OPMRUN and the
auxiliary applications (Figure C.17).

\pandocbounded{\includegraphics[keepaspectratio]{appendices/images/Frame117_adea6795ab92.png}}
\pandocbounded{\includegraphics[keepaspectratio]{appendices/images/Image81_adea6795ab92.png}}

\subsubsection{View Menu Options.}\label{view-menu-options.}

Allows the user to view the results of an OPM Flow simulation run using
the default editor. The options are also available by right-clicking a
job in the Job List Element.

\subsubsection{Tools Menu Options}\label{tools-menu-options}

Contains various tool that may be useful in building a simulation model.

\pandocbounded{\includegraphics[keepaspectratio]{appendices/images/Image84_fc144c3ce032.png}}
\pandocbounded{\includegraphics[keepaspectratio]{appendices/images/Image83_155309479934.png}}

See the section C.3 OPMRUN Tools for further details on the available
tools.

Help Menu Options

\begin{figure}[H]

{\centering \pandocbounded{\includegraphics[keepaspectratio]{appendices/images/Image85_9d7215479111.png}}

}

\caption{Image85}

\end{figure}%

Use the Edit / Options menu option to select the location of the OPM
Flow Manual.

Simulation input and output files can be extremely large, especially for
large full field models. Running multiple cases in these circumstances
can easily use up all available disk space, especially if multiple users
are running multiple cases. The Tools/Compression Jobs option allows the
user to compress a series of jobs into individual zip files (one zip
file per job), as well as uncompressing previously zip job files.

\pandocbounded{\includegraphics[keepaspectratio]{appendices/images/Frame92_155309479934.png}}
\pandocbounded{\includegraphics[keepaspectratio]{appendices/images/Image86_155309479934.png}}

\pandocbounded{\includegraphics[keepaspectratio]{appendices/images/Frame93_2ddce23345cf.png}}
\pandocbounded{\includegraphics[keepaspectratio]{appendices/images/Image87_2ddce23345cf.png}}

Note there is a similar application for uncompressing zip files and that
the tool uses the Linux zip and unzip programs both on Linux host
systems and Windows 10 systems using WSL.

The OPMRUN Keyword generator is an application that generates OPM Flow
keywords that can be cut and pasted into any editor or saved to a
separate file to form the basis of a new input deck (*.DATA files). The
application utilizes templates based on the Apache Velocity Template
Language (``VTL''), a commonly used template language used by software
engineers. The templates can therefore also be used with any editor that
supports VTL, jEdit for example, a popular open source Java based editor
and PyCharm, a Python integrated development environment used in
computer programming, specifically for the Python language.

There is one template per keyword, with formatting of the keywords being
the same as the OPM Flow manual. Currently there are over 450 templates
implemented and the intention is for additional keywords to be added as
the their usage is implemented within the simulator and documented
within the manual. The application allows one to customize the existing
templates as well as creating user defined templates by including the
templates in the template directory and following the VTL language
syntax. One still has to edit the resulting keywords to match the data
require, but the structure and comments are provided by the application.

The application is accessed via the Tool/Simulator Input/Keywords menu
item and the tool can generate specific keywords, as well as complete
sections (Figure C.20)

\pandocbounded{\includegraphics[keepaspectratio]{appendices/images/Frame94_c3ec4b8e078d.png}}
\pandocbounded{\includegraphics[keepaspectratio]{appendices/images/Image88_c3ec4b8e078d.png}}

The application consists of several elements, a conventional menu system
at the top, a Deck Element Area that will contain the resulting
generated keywords, a Keyword Element Area for the user to select the
keyword, data, models or user templates, and finally a series of
buttons, HEADER, GLOBAL, etc., that are used to select the keywords in a
OPM Flow section, specific data sets, models or user defined templates.
The selection will appear in the Keyword Element Area.

Clicking on an item in the Keyword Element Area will generate the data
for the item in the Deck Element Area, as shown below for the OPM Flow
copyright header in Figure C.21.

\pandocbounded{\includegraphics[keepaspectratio]{appendices/images/Frame79_dd0a15097acd.png}}
\pandocbounded{\includegraphics[keepaspectratio]{appendices/images/Image89_dd0a15097acd.png}}

The Deck Element Area is editable by simply clicking anywhere in the
element and making changes. Use the Clear button to clear the Deck
Element Area display, the Copy button to copy the Deck Element Area data
to the clipboard, and the Save button to save the data to a file. The
Load button allows one to load an existing file into the Deck Element
Area for additional editing.

Note that the HEADER section is not an OPM Flow section, but a list of
various comment block headers used to make the deck more readable.

\subsubsection{Keywords: Menu Options}\label{keywords-menu-options}

The various menu options include the File Menu

\pandocbounded{\includegraphics[keepaspectratio]{appendices/images/Frame95_23f062fbd2e6.png}}
\pandocbounded{\includegraphics[keepaspectratio]{appendices/images/Image90_23f062fbd2e6.png}}

Where the Open and Save options load and save a file, and the Properties
displays OPMRUN's properties.

\pandocbounded{\includegraphics[keepaspectratio]{appendices/images/Frame96_5a8b0dbf9da5.png}}
\pandocbounded{\includegraphics[keepaspectratio]{appendices/images/Image91_5a8b0dbf9da5.png}}

The Edit Menu provides some basic standard editing facilities

\pandocbounded{\includegraphics[keepaspectratio]{appendices/images/Frame97_b575b69a5512.png}}
\pandocbounded{\includegraphics[keepaspectratio]{appendices/images/Image92_b575b69a5512.png}}

Next, the Generate Menu options allows one to generate a complete
section of keywords, as described below. These options are equivalent to
selecting the equivalent section keyword (RUNSPEC, GRID, etc.) in the
Keyword Element Area.

\pandocbounded{\includegraphics[keepaspectratio]{appendices/images/Frame98_e1811c16b086.png}}
\pandocbounded{\includegraphics[keepaspectratio]{appendices/images/Image93_e1811c16b086.png}}

Finally, the Help Menu option display the Keyword Help information
(Figure C.26):

\pandocbounded{\includegraphics[keepaspectratio]{appendices/images/Frame99_b2f2dd5a5681.png}}
\pandocbounded{\includegraphics[keepaspectratio]{appendices/images/Image94_b2f2dd5a5681.png}}

And the Velocity Template Help (Figure C.27).

\pandocbounded{\includegraphics[keepaspectratio]{appendices/images/Frame100_517ee8307d51.png}}
\pandocbounded{\includegraphics[keepaspectratio]{appendices/images/Image95_517ee8307d51.png}}

As mentioned previously, the tool uses the Apache Velocity Template
Language (``VTL'') for the templates, and therefore the templates can
also be used directly with an editor, provided the editor supports VTL.

\subsubsection{Keywords: File Imports}\label{keywords-file-imports}

If a keyword requires a file, for example, the INCLUDE and LOAD
keywords, then a dialog box is presented to enable the file to be
selected. The application will also allow one to select the file name
format, after the file has been selected.

\pandocbounded{\includegraphics[keepaspectratio]{appendices/images/Frame101_ce729ab5152b.png}}
\pandocbounded{\includegraphics[keepaspectratio]{appendices/images/Image96_ce729ab5152b.png}}

Note that COMMENT template is not an actual keyword, but a comment block
to make the deck more readable for the user.

\subsubsection{Keywords: Section Standard Set of
Keywords}\label{keywords-section-standard-set-of-keywords}

Selecting a Generate Menu option or a Section keyword (RUNSPEC, GRID,
EDIT, PROPS, SOLUTION, SUMMARY, and SCHEDULE) in the Keyword Element
Area will give an option to generate a representative set of keywords
for that section, as per the RUNSPEC example in Figure C.29.

\pandocbounded{\includegraphics[keepaspectratio]{appendices/images/Frame102_ab2ed7a58bcd.png}}
\pandocbounded{\includegraphics[keepaspectratio]{appendices/images/Image97_ab2ed7a58bcd.png}}

One can therefore generate a complete input deck in a matter of minutes;
however, you still have to edit the generated input deck with your
actual data.

\subsubsection{Keywords: SUMMARY Section
Variables}\label{keywords-summary-section-variables}

For the SUMMARY section keyword, one can also generate various sets of
summary variables based on the options being used in the model. Note
that not all the variables are currently available in OPM Flow, but
additional variables are added at each release.

\pandocbounded{\includegraphics[keepaspectratio]{appendices/images/Frame103_c72bf3b63b70.png}}
\pandocbounded{\includegraphics[keepaspectratio]{appendices/images/Image98_c72bf3b63b70.png}}

For SUMMARY variables not recognized by OPM Flow, the simulator will
issue a warning message and ignore those variables not implemented.

\subsubsection{Keywords: SCHEDULE Section Keywords and Date
Schedule}\label{keywords-schedule-section-keywords-and-date-schedule}

For the SCHEDULE Section keyword, one can also generate a representative
set of SCHEDULE keywords, plus a date schedule from a start year to an
end year, using Annual, Quarterly, or Monthly time steps.

\pandocbounded{\includegraphics[keepaspectratio]{appendices/images/Frame104_6da10d143395.png}}
\pandocbounded{\includegraphics[keepaspectratio]{appendices/images/Image99_6da10d143395.png}}

This option also writes a standard report using the RPTSCHED keyword at
the beginning of each year which is subsequently switch off for the
intermediate Quarterly and Monthly time steps. A final report is written
at the end of the run.

\subsubsection{Keywords: DATA (Sets)
Option}\label{keywords-data-sets-option}

There is also a DATA option which is not an OPM Flow section, but a
series of data sets, as shown in Figure C.32.

\pandocbounded{\includegraphics[keepaspectratio]{appendices/images/Frame105_15944593cc79.png}}
\pandocbounded{\includegraphics[keepaspectratio]{appendices/images/Image100_15944593cc79.png}}

The data sets are complete examples for a given type of data used in OPM
Flow, for example a PVT data set for a wet gas reservoir, or three phase
relative permeability data set. The data sets are intended to be used as
a guide for generating ones own keyword input, or for building models
for testing.

The intention is to expand the collection of data sets over time as more
data becomes available.

\subsubsection{Keywords: MODEL Option}\label{keywords-model-option}

Like the DATA option, the MODEL option is not an OPM Flow section, but
is instead a collection of working models, as illustrated in Figure
C.33.

The purpose of the models is to illustrate the functionality of various
features implemented in OPM Flow and to act as guide for users in
building their own models.

\pandocbounded{\includegraphics[keepaspectratio]{appendices/images/Frame106_df5522b4bbe3.png}}
\pandocbounded{\includegraphics[keepaspectratio]{appendices/images/Image101_df5522b4bbe3.png}}

Additional models will be added when available.

\subsubsection{Keywords: USER Templates
Option}\label{keywords-user-templates-option}

Finally, the USER option is where users can store their own templates.
USER templates with the ``vm'' extension will automatically be listed by
the USER button. To use this feature, after selecting a keyword, right
clicking on the keyword allows one to load the actual template for the
keyword. One can then edit the template and save the changes back to the
same template or another template using the Save button.

\pandocbounded{\includegraphics[keepaspectratio]{appendices/images/Frame107_92c3f14cad63.png}}
\pandocbounded{\includegraphics[keepaspectratio]{appendices/images/Image102_92c3f14cad63.png}}

The Template Help option displays a brief introduction to VTL for
further reference.

The Tools/Simulator Input/Production Schedule application takes a comma
delimited CSV file containing historical production and injection data
and converts the data to an OPM Flow SCHEDULE file using the WCONHIST
series of keywords. An example input file is shown below:

\pandocbounded{\includegraphics[keepaspectratio]{appendices/images/Frame108_8e4aef9b7c61.png}}
\pandocbounded{\includegraphics[keepaspectratio]{appendices/images/Image103_8e4aef9b7c61.png}}

The first row in the input file is a header row that declares the data
type for a column, the example shows typical Oil Field Manager (``OFM'')
header variable names, but various variable names can be used to define
the data type.

The tool can convert daily production data to a: daily production
schedule, monthly average, or monthly on-stream average production
schedule, as shown below:

\pandocbounded{\includegraphics[keepaspectratio]{appendices/images/Frame109_5ccfa8fa22e8.png}}
\pandocbounded{\includegraphics[keepaspectratio]{appendices/images/Image104_5ccfa8fa22e8.png}}

Notice that the application checks various variable names for the column
headers. For example for the BHP data, the column names can be: bhp,
bottom-hole pressure, BHP, or BOTTOM-HOLE PRESSURE.

A sample of the generated output file is shown in Figure C.37.

\pandocbounded{\includegraphics[keepaspectratio]{appendices/images/Frame110_e1ad4ec6e910.png}}
\pandocbounded{\includegraphics[keepaspectratio]{appendices/images/Image105_e1ad4ec6e910.png}}

\begin{tcolorbox}[enhanced jigsaw, title=\textcolor{quarto-callout-note-color}{\faInfo}\hspace{0.5em}{Note}, coltitle=black, breakable, toprule=.15mm, colbacktitle=quarto-callout-note-color!10!white, leftrule=.75mm, opacitybacktitle=0.6, colback=white, bottomtitle=1mm, toptitle=1mm, titlerule=0mm, left=2mm, colframe=quarto-callout-note-color-frame, arc=.35mm, rightrule=.15mm, bottomrule=.15mm, opacityback=0]

Note the current release only support production data via the WCONHIST
keyword, injection data via WCONINJH keyword is not supported.

\end{tcolorbox}

The Tools/Simulator Input/Sensitivities option generates sensitivity
cases based on a ``Base'' case file. The Base file contains ``Factors''
(variable names), \$X01, \$X02, etc., that are substituted with user
defined values using the data entered and the type of Sensitivity
Scenario selected. Thus, the first step is to configure the Base file in
a text editor by replacing actual values by the variable names,
previously mentioned.

\pandocbounded{\includegraphics[keepaspectratio]{appendices/images/Frame111_543ef82c00ed.png}}
\pandocbounded{\includegraphics[keepaspectratio]{appendices/images/Image106_543ef82c00ed.png}}

After editing the Base file, the next step is to load the Base file into
the application using the Base button, which will prompt the user for
the file to load and then display the file in the Base tab, as shown in
Figure C.39

\pandocbounded{\includegraphics[keepaspectratio]{appendices/images/Frame80_dcc9bd1c4244.png}}
\pandocbounded{\includegraphics[keepaspectratio]{appendices/images/Image107_dcc9bd1c4244.png}}

Limited editing of the Base file is supported on the above screen.

The next step is to define the ``Factors'' and the factor values. A
total of 20 factors are available and each factor consist of a Low, Best
and High estimates. Note it is not necessary to enter all three
estimates, if one wishes just to generate a limited sensitivity case.
For example, if on wishes to only run a Low Scenario sensitivity then it
is only necessary to enter data for the Low factor values.

Previously saved factor data can be loaded via the Load button, as shown
below:

\pandocbounded{\includegraphics[keepaspectratio]{appendices/images/Frame81_1bd085a6a12f.png}}
\pandocbounded{\includegraphics[keepaspectratio]{appendices/images/Image108_1bd085a6a12f.png}}

Selecting a Factor Description row allows one to define a description
for the factor variable, so for \$X01 in the above figure the
description is GRID - PERMX. When selecting a Factor Description, a
popup dialog will be displayed to enter the data, and if one
right-clicks on the popup's Factor Description field one can select a
description for one of the pre-defined descriptions as illustrated in
the next figure.

\pandocbounded{\includegraphics[keepaspectratio]{appendices/images/Frame82_8ef499b5c7ea.png}}
\pandocbounded{\includegraphics[keepaspectratio]{appendices/images/Image109_8ef499b5c7ea.png}}

After the Sensitivity Factors have been entered one can then select the
Sensitivity Scenario that one wishes to use generate the sensitivity
cases. In the example in figure Figure C.42 the Factorial: Low, Best and
High Box-Behnken DOE (Design of Experiments) has been selected.
Selecting the Generate button, runs a series of checks, and if there are
no errors the program will inquire if you wish to generate the set of
cases (Figure C.42).

\pandocbounded{\includegraphics[keepaspectratio]{appendices/images/Frame83_b3269ce01e60.png}}
\pandocbounded{\includegraphics[keepaspectratio]{appendices/images/Image110_b3269ce01e60.png}}

If the Yes option is selected then the cases will be generated and the
application will ask for the name of OPMRUN Queue file to write the jobs
to, as depicted in

\pandocbounded{\includegraphics[keepaspectratio]{appendices/images/Frame112_15ae4dea63b4.png}}
\pandocbounded{\includegraphics[keepaspectratio]{appendices/images/Image111_15ae4dea63b4.png}}

This allows the user to load the queue file into OPMRUN and to run all
the jobs.

This tool, Tools/Simulator Input/Well Specification, uses the standard
well export files from OPM ResInsight to reformat the data in a more
user-friendly manner for the WELSPECS and COMPDAT keywords. Optionally,
the application can generate the COMPLUMP keyword based on the OPM
ResInsight layers file, with one completion per defined reservoir layer.

An example OPM ResInsight Exported Well Completion File Format(*.exp) is
shown in Figure C.44

\pandocbounded{\includegraphics[keepaspectratio]{appendices/images/Frame84_c3e08ebf4dd4.png}}
\pandocbounded{\includegraphics[keepaspectratio]{appendices/images/Image112_c3e08ebf4dd4.png}}

And an OPM ResInsight Imported Formation Layer File (.Lyr) example is
illustrated in Figure C.45

\pandocbounded{\includegraphics[keepaspectratio]{appendices/images/Frame85_5b63ed26d33a.png}}
\pandocbounded{\includegraphics[keepaspectratio]{appendices/images/Image113_5b63ed26d33a.png}}

The application also can generate a well a OPM ResInsight perforation
file with the formation names for cross-checking the perforations.

The application user interface is shown in Figure C.46. Note that in
Figure C.46 the Output Header options are used for comments only, no
unit conversion is performed.

\pandocbounded{\includegraphics[keepaspectratio]{appendices/images/Frame86_0088f772bfbe.png}}
\pandocbounded{\includegraphics[keepaspectratio]{appendices/images/Image114_0088f772bfbe.png}}

In terms of output, the next figure shows the resulting well completion
file to be used with OPM Flow, showing the WELSPECS and COMPDAT keywords
(the COMPLUMP keyword is not shown in this example)

\pandocbounded{\includegraphics[keepaspectratio]{appendices/images/Frame87_fa267410371c.png}}
\pandocbounded{\includegraphics[keepaspectratio]{appendices/images/Image115_fa267410371c.png}}

The final figure for this tool shows the resulting generated OPM
ResInsight perforation file.

\pandocbounded{\includegraphics[keepaspectratio]{appendices/images/Frame88_9b039e8d5921.png}}
\pandocbounded{\includegraphics[keepaspectratio]{appendices/images/Image116_9b039e8d5921.png}}

This option, Tools/ResInsight, loads the currently selected job into OPM
ResInsight for viewing, this done via a Python sub-process call in
OPMRUN, rather than using OPM ResInsight's Python API.

OPM ResInsight can read well trajectories in a given format into the
program, the Tools/Well Trajectory Conversion option converts a
Schlumberger Petrel exported well trajectory file, as shown Figure C.49,
into a OPM ResInsight well trajectory file containing all the wells.

\pandocbounded{\includegraphics[keepaspectratio]{appendices/images/Frame89_21aee4be8c4f.png}}
\pandocbounded{\includegraphics[keepaspectratio]{appendices/images/Image117_21aee4be8c4f.png}}

The utility allows for the multiple wells to be converted at once and
for conversion of units.

Note in some areas of the world it is not uncommon for the areal units
to be in UTM and the depth to be in feet. This configuration is also
handled by the application.

\pandocbounded{\includegraphics[keepaspectratio]{appendices/images/Frame90_46df760852cd.png}}
\pandocbounded{\includegraphics[keepaspectratio]{appendices/images/Image118_46df760852cd.png}}

An example output file is shown in Figure C.51

\pandocbounded{\includegraphics[keepaspectratio]{appendices/images/Frame91_9025a30f6d3d.png}}
\pandocbounded{\includegraphics[keepaspectratio]{appendices/images/Image119_9025a30f6d3d.png}}

\chapter{Python Interface}\label{python-interface}

Some of the C++ code in OPM Flow has been ``wrapped'' in Python which
means that one can invoke the simulator's C++ code from one's own Python
programs. At the moment the wrappers for the input layer and the code
for working with the result files is quite complete and usable. For a
python class documentation and for further documentation and examples,
please refer to the
\href{https://opm.github.io/opm-python-documentation/index.html}{OPM
Online Python Documentation}. There is also the option to interact with
running simulations from within Python using the PYACTION keyword.

The goal is for the Python code to be structured like a native Python
API, but some design decisions are certainly influenced by the
underlying C++ implementation, and to get a deeper understanding of what
is possible and how to achieve specialized tasks one might need to
consult the simulator's C++ code.

\begin{tcolorbox}[enhanced jigsaw, title=\textcolor{quarto-callout-note-color}{\faInfo}\hspace{0.5em}{Note}, coltitle=black, breakable, toprule=.15mm, colbacktitle=quarto-callout-note-color!10!white, leftrule=.75mm, opacitybacktitle=0.6, colback=white, bottomtitle=1mm, toptitle=1mm, titlerule=0mm, left=2mm, colframe=quarto-callout-note-color-frame, arc=.35mm, rightrule=.15mm, bottomrule=.15mm, opacityback=0]

Observe that the Python bindings described here are not very mature. The
API might change in future releases. It is in general quite simple to
expose new functionality to Python.

\end{tcolorbox}

The reading of the input deck in OPM Flow is a two step process, first a
data structure called a ``Deck'' is created - the ``Deck'' is
essentially a collection of keywords where all elements have been
converted to the correct type and default values have been injected into
the keyword. The ``Deck'' is a quite low level data structure, and the
actual simulation is based on higher level data structures where the
origin in ``keywords'' from a *.DATA file is no longer so apparent. The
most notable high level objects are the ``EclipseState'' and
``Schedule'' classes. Large parts of this functionality is available
from Python.

The initial building block in the parsing process is the ``Parser''
object which knows how to interpret the input keywords and create data
structures. Virtually all scripts working with the input files will
start with creating a Parser object:

\begin{verbatim}
#!/usr/bin/env python3
from opm.io.parser import Parser
# Create a parser object which can be used to parse a string or input files
parser = Parser()
\end{verbatim}

In most cases one will just create a default parser object and be done
with it, but it is possible to both add your own keyword definitions to
the parser and alternatively create a custom parser which only accepts a
subset of keywords. These features will be demonstrated briefly in
section Special Parsing.

The ``Deck'' data structure is essentially a list of keywords which have
been loaded as ``DeckKeyword'' instances. The Deck can either be loaded
from an input data file or from a string. The following example shows
how to load two SCHEDULE section keywords: WELSPECS and COMPDAT:

\begin{verbatim}
#!/usr/bin/env python3
from opm.io.parser import Parser
deck_string = """
WELSPECS
   'W1'  'G1'   19    4      1*   WATER   1*   1*   SHUT  1*  1*    1*      /
/
COMPDAT
  'W1'   11   3    1    5   OPEN   1*   1*   0.216   1*   0    1*   Z    1* /
/
"""
parser = Parser()
deck = parser.parseString( deck_string )
\end{verbatim}

After running this small script the deck variable will contain the two
keywords WELSPECS and COMPDAT. For these keywords the '*' have been
replaced with the correct default value and the items have been
converted to the correct type. In addition, all numerical values have
been converted to SI units, more about units can be found in section
Units. Alternatively you can use the Parser.parseFile() method to parse
an entire input file. In this case INCLUDE and IMPORT keywords will be
resolved and everything will be coalesced into one large ``DeckKeyword''
data structure:

\begin{verbatim}
#!/usr/bin/env python3
 import os.path
 import sys
 from opm.io.parser import Parser
 data_file = sys.argv[1]
 print(f"Loading deck from {data_file}")
 parser = Parser()
 deck = parser.parseFile( data_file )

\end{verbatim}

The code above loads a *.DATA file given by calling the Python script.
For example, if the name of the input deck, that may contain a complete
input deck or just a selection of keywords, is called NORNE.DATA, and
the name of the above script is LoadDeck.py, then:

\begin{verbatim}
python LoadDeck NORNE.DATA
\end{verbatim}

would load the NORNE.DATA into the simulator.

Internally in OPM Flow all quantities are managed as SI units, whereas
input decks use the unit systems METRIC, FIELD, or LAB (laboratory). The
METRIC unit system is comparable to SI units, the most notable
difference is that time is expressed in days instead of seconds, the
FIELD unit system is based on historical units like barrels of oil (stb)
and feet for length.

The important point is that when you access the elements in the Deck the
values you will get are in SI units, whereas the input you provided has
been in one of other set of units. Consider the example:

\begin{verbatim}
#!/usr/bin/env python3
from opm.io.parser import Parser
deck_string = """
WCONPROD
 'W1'  'OPEN'   'ORAT'  86400 /
/
"""
parser = Parser()
deck = parser.parseString( deck_string )
kw = deck[0]
record0 = kw[0]
orat = record0["ORAT"].value
well = record0["WELL"].value
print("==> The Oil rate in well:{} is {} m3/s".format(orat, well))
\end{verbatim}

When this script is run it will print:

\begin{verbatim}
==> The Oil rate in well:W1 is 1.0 m3/s
\end{verbatim}

because the input value of 86,400 m3/day is internally converted to 1.0
m3/second.

The default input unit system is METRIC, that means that if the input is
to be interpreted in either FIELD or LAB units, then one needs to add a
keyword to declare the units before entering the data. Thus, to
interpret the WCONPROD keyword's oil rate as stb/d in the previous
example, the FIELD keyword needs to be added to the example, that is:

\begin{verbatim}
#!/usr/bin/env python3
   from opm.io.parser import Parser
   deck_string = """

   FIELD

   WCONPROD
    'W1'  'OPEN'   'ORAT'  86400 /
   /
   """
\end{verbatim}

After you have loaded a Deck you can query it with normal Python
functions:

\begin{verbatim}
# Check if deck has keyword:
if "GRID" in deck:
    print("Deck contains 'GRID' keyword")
else
    print("Deck does not have 'GRID' keyword")

# Loop through all the keywords:
for kw in deck:
    print("kw: {}".format(kw.name)
\end{verbatim}

To get a keyword from the deck you can access it with index, name, or a
combination of name and occurrence index:

\begin{verbatim}
# Get keyword 10
kw10 = deck[10]

# Get the last DATES keyword
last_dates = deck["DATES"]

# Get the third WELSPECS keyword
welspecs3 = deck[("WELSPECS", 3)]
\end{verbatim}

A deck keyword is composed of ``records'', which are again composed of
``items''. The example below shows how one can iterate through the wells
in a WELSPECS keyword:

\begin{verbatim}
kw = deck["WELSPECS"]
print(f"Keyword has {len(kw)} records")
for record in kw:
    well = record["WELL"].get_str()
    group = record[1].get_str()

    print(f"Well {num}: {well} is part of group{group})
\end{verbatim}

The keywords that OPM Flow recognizes are configured via JSON files
which are embedded in the source distribution. The JSON keywords for the
simulator are in the following public GitHub repository:

\url{https://github.com/OPM/opm-common/tree/master/opm/input/eclipse/share/keywords}

By default the Parser class will recognize all the keywords which are
known to OPM Flow, but one can create your own custom parser with only a
subset of keywords:

\begin{verbatim}
from opm.io.parser import Builtin
# Create a special parser which only recognizes the corner point grid               # keywords
parser = Parser(add_default = False)
builtin = Builtin()
parser.add_keyword( builtin.COORD )
parser.add_keyword( builtin.ZCORN )
parser.add_keyword( builtin.ACTNUM )
\end{verbatim}

The opposite is also possible, here is how one can add a special keyword
GCLOSE which could be used to close all the wells in the group:

\begin{verbatim}
# Create a fictitious home mode keyword GCLOSE:
GCLOSE = {"name" : "GCLOSE",
          "sections" : ["SCHEDULE"],
          "items" : [
             {"name" : "GROUP", "value_type" : "STRING"}
           ]}

parser = Parser()
# Add the keyword description via a JSON string
parser.add_keyword(json.dumps(GCLOSE))
\end{verbatim}

By default the parser used in OPM Flow is quite strict - a tad stricter
than the one used by the commercial simulator. However, one can
configure how the parser should react to certain error conditions. By
default the parser will fail with an exception if an INCLUDE file is not
found, but one can for instance change this behavior to ignore that
error condition, for example:

\begin{verbatim}
parse_context = ParseContext([('PARSE_MISSING_INCLUDE',
                                opm.io.action.ignore)])
parser = Parser()
deck = parser.parseFile(self.norne_fname, parse_context)
\end{verbatim}

The ``Deck'' is a quite raw structure and not very user friendly. Before
actually used in the simulator the properties from the ``Deck'' are
coalesced to a form more easily usable format. Consider, for example the
PORO array configured below where the configuration is given in three
steps:

\begin{verbatim}
-- GRID Dimensions 10 x 10 x 3
DIMENS
  10 10 3 /
--
-- Set A Global Value 0.10 For All Of The Grid
--
PORO
   300*0.10 /
--
-- Set A Value Of 0.15 For All The Cells In The Middle Layer
--
BOX
   1 10 1 10 2 2 /
PORO
   100*0.15 /
ENDBOX
--
-- Scale The Bottom Layer By A Factor OF Two
--
MULTIPLY
   PORO 2 1 10 1 10 3 3 /
/
\end{verbatim}

In the deck representation this will be a collection of five different
keywords: \{``PORO'', ``BOX'', ``PORO'', ``ENDBOX'', ``MULTIPLY''\},
whereas in the ``EclipseState'' this will be one property ``PORO'' where
all the modifications have been completed. In addition to the grid
properties like ``PERMX'' and ``PORO'' the ``EclipseState'' object
contains numerous other objects like the fault properties, PVT tables,
and more.

\begin{verbatim}
#!/usr/bin/env python3
import syshe
form opm.io.parser import Parser
from opm.io.ecl_state import EclipseState
parser = Parser()
data_file = sys.argv[1]
deck = parser.parseFile( data_file )
es = EclipseState( deck )
#
# Get the FieldPropsManager from the EclipseState. The FieldPropsManager
# is then used to query and look up grid properties like PERMX and PORO.
# Observe that the properties you get from the FieldPropsManager only have
# the active cells only.
#
fp = es.field_props()
poro = fp["PORO"]
satnum = fp["SATNUM"]
#
# Get the input grid from the EclipseState:
#
grid = es.grid()
# Get a manager for all the tables
tables = es.tables()
\end{verbatim}

The Schedule object contains all the dynamic information in the model,
in particular that includes all information about wells and groups. The
Schedule object is constructed from the ``Deck'' and the
``EclipseState'' objects.:

\begin{verbatim}
#!/usr/bin/env python3
import sys
from opm.io.parser import Parser
from opm.io.ecl_state import EclipseState
parser = Parser()
data_file = sys.argv[1]
deck = parser.parseFile( data_file )
es = EclipseState( deck )
schedule = Schedule(deck, es)
print(f"Schedule file has {len(schedule)} report steps")'
print("List of wells")
for well in schedule.well_names():
    print(well)
\end{verbatim}

The example above prints the number of report time steps and all the
well names in the SCHEDULE section.

In addition to the functionality to load and inspect the input deck, OPM
Flow also has a Python interface to load and inspect the various result
files, SUMMARY and RESTART files. There is also some functionality to
create files with the correct formatting.

Using the class ``EclipseGrid'' one can load a grid representation from
a *.EGRID file on disk:

\begin{verbatim}
#!/usr/bin/env python3
import sys
from opm.io.ecl import EGrid
grid_file = sys.argv[1]
grid = EGrid(grid_file)
\end{verbatim}

The return value from ``EclipseState.ecl\_grid()'' is also of this type.

Using the class ``ESMry'' one can load the Extended SUMMARY file for a
simulations that contains all the SUMMARY vectors.

\begin{verbatim}
#!/usr/bin/env python3
import sys
from opm.io.ecl import ESMry
case = sys.argv[1]
summary = Esmry(case)
\end{verbatim}

Finally, one load the solution arrays stored in the *.RESTART file using
the Erst class, as shown below.

\begin{verbatim}
#!/usr/bin/env python
import sys
from opm.io.ecl import ERst

rst_file = sys.argv[1]
rst = Erst(rst_file)
\end{verbatim}

Note that a restart file can be large, so care must be taken on how
these type of files are loaded into Python.

\chapter{Command Line Options}\label{command-line-options}

The command line syntax for running OPM Flow is:

\begin{verbatim}

flow [OPTIONS] CASENAME
\end{verbatim}

and typing the following command from your terminal:

\begin{verbatim}

flow CASENAME
\end{verbatim}

will start the simulator and run the case specified by CASENAME.DATA.

OPM Flow accepts command line options to control various aspects of the
simulator's run time behavior, as described in the table below. To give
a command line option use ``key=value'' syntax, with no spaces around
the equals sign. It is also possible to put multiple options together in
a parameter file. To do so, put one option on each line of the file with
``.param'' as the extension and pass that filename as a command line
parameter to OPM Flow using:

\begin{verbatim}

flow --parameter-file=CASENAME.param
\end{verbatim}

to instruct OPM Flow to read the parameter file.

\begin{longtable}[]{@{}
  >{\raggedright\arraybackslash}p{(\linewidth - 6\tabcolsep) * \real{0.2500}}
  >{\raggedright\arraybackslash}p{(\linewidth - 6\tabcolsep) * \real{0.2500}}
  >{\raggedright\arraybackslash}p{(\linewidth - 6\tabcolsep) * \real{0.2500}}
  >{\raggedright\arraybackslash}p{(\linewidth - 6\tabcolsep) * \real{0.2500}}@{}}
\toprule\noalign{}
\begin{minipage}[b]{\linewidth}\raggedright
OPM Flow 2025-04 Command Line Options
\end{minipage} & \begin{minipage}[b]{\linewidth}\raggedright
\end{minipage} & \begin{minipage}[b]{\linewidth}\raggedright
\end{minipage} & \begin{minipage}[b]{\linewidth}\raggedright
\end{minipage} \\
\midrule\noalign{}
\endhead
\bottomrule\noalign{}
\endlastfoot
No. & Variable Name & Description & Default \\
1 & -h or --help & A character string that causes OPM Flow to print a
help message that gives a brief description of the available command
line parameters. This now only prints the active user facing command
line options. & N/A \\
2 & --help-all & Prints all the command line options included in the
release, including obsolete, hidden and deprecated options. & N/A \\
3 & --accelerator-mode & A defined character string that selects the
linear solver, usage:
`--accelerator-mode={[}none\textbar cusparse\textbar opencl\textbar amgcl\textbar rocalution\textbar rocsparse{]}'.
& ``none'' \\
4 & --action-parsing-strictness & A defined character string. Set
strictness of parsing process for ActionX and PyAction. Available
options are normal (do not apply keywords that have not been tested for
ActionX or PyAction) and low (try to apply all keywords, beware: the
simulation outcome might be incorrect). & ``normal'' \\
5 & --add-corners & A boolean value. Add corners to partition. &
false \\
6 & --allow-distributed-wells & Allow the perforations of a well to be
distributed to interior of multiple processes. & false \\
7 & --allow-splitting-inactive-wells & A boolean value. Allow inactive
(never non-shut) wells to be split across multiple domains. & true \\
8 & --alternative-well-rate-init & Use alternative well rate
initialization procedure. & true \\
9 & --bda-device-id & Renamed as --gpu-device-id. & 0 \\
10 & --check-group-constraints-inner-well-iterations & A boolean value.
Allow checking of group constraints during inner well iterations. &
true \\
11 & --check-satfunc-consistency & A boolean value. Whether or not to
check saturation function consistency requirements. & false \\
12 & --continue-on-convergence-error & Continue with a non-converged
solution instead of giving up if we encounter a time step size smaller
than the minimum time step size. & false \\
13 & --convergence-monitoring & A boolean value. Enable convergence
monitoring. & false \\
14 & --convergence-monitoring-cut-off & An integer value. Cut off limit
for convergence monitoring. & 6 \\
15 & --convergence-monitoring-decay-factor & A real value. Decay factor
for convergence monitoring. & 0.75 \\
16 & --cpr-reuse-interval & A positive integer that sets the reuse
preconditioner interval. Used when --cpr-reuse-setup is set to 4, then
the preconditioner will be fully recreated instead of reused every N
linear solve, where N is this parameter. & 30 \\
17 & --cpr-reuse-setup & A positive integer that defines if the CPR
solver should re-use the AMG setup. Valid options are: Changed the
default value from three to four. & 4 \\
18 & --dbph-max-rel & A real positive value that sets the maximum
relative change of the bottom-hole pressure in a single iteration. &
1.0 \\
19 & --debug-emit-cell-partition & A boolean value that determines
whether (true) or not (false) to emit cell partitions as a debugging
aid. & false \\
20 & --dp-max-rel & A real positive double precision value that sets the
maximum allowed relative change of pressure per iteration. & 0.3 \\
21 & --ds-max & A real positive double precision value that sets the
maximum allowed change in saturation per iteration. & 0.2 \\
22 & --dwell-fraction-max & A real positive value that sets the maximum
allowed change in well's volume fraction per iteration. & 0.2 \\
23 & --ecl-deck-file-name & A character string that defines the name of
the OPM Flow input file which contains the simulator's ECLIPSE formatted
input deck to be simulated. Note that there must be no spaces in the
filename. If there are spaces, which can typically occur on Windows
based systems, then the filename should be enclosed in quotes; thus
instead of FILE NAME use ``FILE NAME''. & ``\,'' \\
24 & --ecl-output-double-precision & A boolean value that switches on
(true) or off (false) double precision in restart files. Useful for
`perfect' restarts. & false \\
25 & --ecl-output-interval & An integer value defining the number of
report steps that ought to be skipped between two writes of restart
files results. & -1 \\
26 & --edge-weights-method & A defined positive integer value that
defines the edge-weighing strategy: & 1 \\
27 & --enable-adaptive-time-stepping & A boolean value that turns on
(true) or off (false) adaptive time stepping. If set to false the report
time steps provided by the input deck are used. & true \\
28 & --enable-async-ecl-output & A boolean value that sets the output
mode to be asynchronous (true), letting the simulator continue computing
the next time step while writing results to restart and summary files in
the commercial simulator's format on a separate thread. & true \\
29 & --enable-async-vtk-output & A boolean value that sets the output
mode to be asynchronous (true), letting the simulator continue computing
the next time step while writing results to the VTK output files on a
separate thread. & true \\
30 & --enable-drift-compensation & A boolean value that enables (true)
or disables (false) partial compensation of systematic mass losses via
the source term of the next time step. & false \\
31 & --enable-dry-run & A boolean value that specifies if the simulation
should actually run (true), or just check the input deck (false). This
option is equivalent to activating the NOSIM keyword in the RUNSPEC
section of the input deck (see section Error: Reference source not found
Error: Reference source not found). The default value of ``auto' will
use whatever is stipulated in the input deck via the NOSIM keyword. &
``auto'' \\
32 & --enable-ecl-output & A boolean value that specifies if the binary
output files (restart, summary files, etc.) should be written in the
commercial simulator's format (true), or OPM Flow's format (false). &
true \\
33 & --enable-esmry & A boolean value that switches on (true) or off
(false) the output of SUMMARY vectors to the ESMRY file for fast loading
of summary data. & false \\
34 & --enable-gravity & A boolean value that switches on (true) or off
(false) the use of the gravity correction for the pressure gradients. &
true \\
35 & --enable-grid-adaptation & A boolean value that enables (true) or
disables (false) adaptive grid refinement/coarsening. Note that Eclipse
grid output will be disabled if adaptive grid refinement/coarsening is
enabled. & false \\
36 & --enable-intensive-quantity-cache & A boolean value that switches
on (true) or off (false) the caching of intensive quantities. & true \\
37 & --enable-logging-fallout-warning & A boolean value that sets the
developer option to see whether logging was on non-root processors
(true). If set to true output will be appended to the \emph{.DBG or
}.PRT files. & false \\
38 & --enable-opm-rst-file & A boolean value set to true or false to
write OPM specific data sets to the commercial simulators restart file
to enable restart of OPM Flow runs by OPM Flow (true), or not to write
the data (false). & false \\
39 & --enable-storage-cache & A boolean value that turns on (true) or
off (false) storing previous storage terms and avoid re-calculating
them. & true \\
40 & --enable-terminal-output & A boolean value that turns on (true) or
off (false) high-level information about the simulation's progress to
the terminal & true \\
41 & --enable-thermodynamic-hints & A boolean value that enables (true)
or disables (false) thermodynamic hints. & false \\
42 & --enable-tuning & A boolean value that instructs OPM Flow to read
the time stepping and convergence parameters from the TUNING keyword in
the input deck, if set to true. Note that only some of the TUNING
keyword parameters are processed. & false \\
43 & --enable-vtk-output & A boolean value that turns on (true) or off
(false) a global switch for writing VTK files. & false \\
44 & --enable-well-operability-check & A boolean value that enables
(true) checking of a well's operating status, or disables (false) the
checking. & true \\
45 & --enable-well-operability-check-iter & A boolean value that enables
(true) checking of a well's operating status during iterations, or
disables (false) the checking during iterations. & false \\
46 & --enable-write-all-solutions & A boolean value that turns on (true)
or off (false) the writing of all solutions to disk instead of only the
ones for the report steps. & false \\
47 & --end-time & A real value. The simulation time at which the
simulation is finished {[}s{]}. & 1e+100 \\
48 & --explicit-rock-compaction & A boolean value. Use pressure from end
of the last time step when evaluating rock compaction (true) or not
(false). & false \\
49 & --external-partition & A quoted character string. Name of file from
which to load an externally generated partitioning of the model's active
cells for MPI distribution purposes. If empty, the built-in partitioning
method will be employed. & ``\,'' \\
50 & --force-disable-fluid-in-place-output & A boolean value that
instructs OPM Flow not to print the Fluid In-Place report after each
report time step (true) or not (false). Note this parameter will
override the print request in the input deck. & false \\
51 & --force-disable-resv-fluid-in-place-output & A boolean value that
instructs OPM Flow to not to print the Reservoir Volume Fluid In-Place
report after each report time step (true) or not (false). Note this
parameter will override the print request in the input deck. & false \\
52 & --full-time-step-initially & A boolean value that instructs OPM
Flow to always attempt to finish a report step using a single time step
(true) or not (false). & false \\
53 & --gpu-device-id & An integet value. Choose device ID for
cusparseSolver or openclSolver, use `nvidia-smi' or `clinfo' to
determine valid IDs. & 0 \\
54 & --ignore-keywords & A character string that defines a list of
keywords which should be ignored by OPM Flow. The keyword string should
be enclosed in quotes and each keyword separated by a colon, that is
`:'. & ``\,'' \\
55 & --ilu-fillin-level & A positive integer value that sets the fill in
level for the ILU preconditioner. & 0 \\
56 & --ilu-redblack & A boolean value that instructs OPM Flow to use
red-black partitioning for the ILU preconditioner. & false \\
57 & --ilu-relaxation & A real positive double precision value that sets
the relaxation factor of the linear solver's ILU preconditioner. &
0.9 \\
58 & --ilu-reorder-spheres & A boolean value that specifies OPM Flow to
reorder the entries of the matrix in the red-black ILU preconditioner in
spheres starting at an edge (true) or not (false). If false the original
ordering is preserved in each color. Otherwise try to ensure D4 ordering
(in a 2D structured grid, the diagonal elements are consecutive). &
false \\
59 & --imbalance-tol & A real value. Tolerable imbalance of the load
balancing. & 1.1 \\
60 & --initial-time-step-in-days & A real double precision value that
sets the size of initial time step in days. & 1.0 \\
61 & --initial-time-step-size & A real double precision value that sets
the size of initial time step in seconds. & 86400 \\
62 & --inj-mult-damp-mult & A real value. Injection multiplier dampening
factor (dampening multiplied by this each time oscillation is detected).
& 0.9 \\
63 & --inj-mult-min-damp-factor & A real value. Minimum injection
multiplier dampening factor (maximum dampening level). & 0.05 \\
64 & --inj-mult-osc-threshold & A real value. Injection multiplier
oscillation threshold (used for multiplier dampening). & 0.1 \\
65 & --input-skip-mode & A defined character string. Set compatibility
mode for the
\hyperref[REF_HEADING_KEYWORD_SKIP100_11_3]{SKIP100}/\hyperref[REF_HEADING_KEYWORD_SKIP300_11_3]{SKIP300}
keywords. Options are: 100 (skip
\hyperref[REF_HEADING_KEYWORD_SKIP100_11_3]{SKIP100}..ENDSKIP, keep
\hyperref[REF_HEADING_KEYWORD_SKIP300_11_3]{SKIP300}..ENDSKIP), 300
(skip \hyperref[REF_HEADING_KEYWORD_SKIP300_11_3]{SKIP300}..ENDSKIP,
keep \hyperref[REF_HEADING_KEYWORD_SKIP100_11_3]{SKIP100}..ENDSKIP) or
all (skip both
\hyperref[REF_HEADING_KEYWORD_SKIP100_11_3]{SKIP100}..ENDSKIP and
\hyperref[REF_HEADING_KEYWORD_SKIP300_11_3]{SKIP300}..ENDSKIP). &
``100'' \\
66 & --linear-solver & A defined quoted character string that sets the
configuration of the linear solver, valid values are: Option (1) extends
the existing Constrained Pressure Residual (``CPR'') preconditioner to
include wells. This option can also be invoked via the CPR keyword in
the RUNSPEC section; however, the command line parameter takes
precedence. For option (10) one enters a character string enclosed in
quotes that defines the filename of a JSON configuration file for the
flexible linear solver system. In this case the file extension should be
`.json'. Note that the *.PRT file contains the ``Property tree for the
linear solver'' listing, which is the JSON specification of the current
case, and can be used to configure a user specific linear solver JSON
file. & ``cprw'' \\
67 & --linear-solver-ignore-convergence-failure & A boolean value that
if set to true convergences failures in the linear solver are ignored.
This option should be used with care, as the results might be
unreliable. & false \\
68 & --linear-solver-max-iter & A positive integer value that defines
the maximum number of linear iterations. & 200 \\
69 & --linear-solver-print-json-definition & Write the JSON definition
of the linear solver setup to the DBG file. & true \\
70 & --linear-solver-reduction & A real positive double precision value
that sets the minimum reduction of the residual for the linear solver.
The linear solver convergences when the residual is reduced
sufficiently. The simulator now overrides the default reduction to be
0.005 instead of 0.01, if the linear solver has been set to one of the
cpr options, in a similar manner as how the default maximum number of
linear iterations for the cpr and cprw options are changed to 20 instead
of 200, unless specified by the command line option by the user. &
0.01 \\
71 & --linear-solver-restart & A positive integer value that sets the
number of iterations after which GMRES is restarted. & 40 \\
72 & --linear-solver-verbosity & A positive integer value that defines
the output from linear solver: & 0 \\
73 & --load-file & A character string enclosed in quotes that defines
the FileName for the .OPMRST file used to load the serialized state. If
empty, CASENAME.OPMRST is used. & ``\,'' \\
74 & --load-step & An integer value that determines if the simulator
should load the serialized state from OPM Flow's version of the restart
file, and should be set to either a specific report step, or 0 to load
the last stored report step. The default value of -1 does not load the
data from the OPM Flow specific restart file. OPM Flow's version of the
restart file, is written using the --save-step=N option. & -1 \\
75 & --local-domains-ordering-measure & A character string enclosed in
quotes that defines the subdomain ordering measure. Allowed values are
`maxpressure', `averagepressure' and `residual'. & ``maxpressure'' \\
76 & --local-domains-partition-well-neighbor-levels & An integer value.
Number of neighbor levels around wells to include in the same domain
during NLDD partitioning. & 1 \\
77 & --local-domains-partitioning-imbalance & A real positive double
precision value that sets the subdomain partitioning imbalance
tolerance. 1.03 is 3 percent imbalance. & 1.03 \\
78 & --local-domains-partitioning-method & A character string enclosed
in quotes that defines the subdomain partitioning method. Allowed values
are `zoltan', `simple', and the name of a partition file ending with
`.partition'. & ``zoltan'' \\
79 & --local-solve-approach & A character string enclosed in quotes that
defines the local solve approach. Valid choices are `jacobi' and
`gauss-seidel'. & ``gauss-seidel'' \\
80 & --local-tolerance-scaling-cnv & A real positive double precision
value that sets the convergence tolerance for local solves. Set to lower
than 1.0 to use stricter convergence tolerance for local solves. &
0.1 \\
81 & --local-tolerance-scaling-mb & A real positive double precision
value that sets the convergence tolerance for local solves. Set to lower
than 1.0 to use stricter convergence tolerance for local solves. & 1 \\
82 & --local-well-solve-control-switching & A boolean value that allows
(true) or disallows (false) control switching during local well
solutions. & true \\
83 & --matrix-add-well-contributions & A boolean value that if set to
true explicitly specifies the influences of wells between cells in the
Jacobian and preconditioner matrices. & false \\
84 & --max-inner-iter-ms-wells & A positive integer value that defines
the maximum number of inner iterations for multi-segment wells. & 100 \\
85 & --max-inner-iter-wells & A positive integer value that defines the
maximum number of inner iterations for standard wells. & 50 \\
86 & --max-local-solve-iterations & A positive integer value that
defines the maximum number of iterations for local solves with NLDD
nonlinear solver. & 20 \\
87 & --max-newton-iterations-with-inner-well-iterations & A positive
integer that specifies the maximum Newton iterations with inner well
iterations. & 8 \\
89 & --max-pressure-change-ms-wells & A real positive value that defines
the maximum relative pressure change for a single iteration of the
multi-segment well model. & 1e+06 \\
90 & --max-residual-allowed & A real positive value that sets the
absolute maximum tolerance for residuals without cutting the time step
size. & 1e+07 \\
91 & --max-single-precision-days & A real positive value that sets the
maximum time step size where single precision floating point arithmetic
can be used for solving the linear systems of equations. & 20.0 \\
92 & --max-temperature-change & A real positive value that stipulates
the maximum absolute change of temperature in a single iteration. & 5 \\
93 & --max-time-step-divisions & A positive integer that specifies the
maximum number of divisions by two of the timestep size before the
simulation bails out. & 10 \\
94 & --max-time-step-size & A real positive value that defines the
maximum size to which all time steps are limited to in seconds. & inf \\
95 & --max-welleq-iter & A positive integer that defines the maximum
number of iterations to determine the solution to the well equations. &
30 \\
96 & --maximum-number-of-group-switches & An integer value. Maximum
number of times a group can switch to the same control. & 3 \\
97 & --maximum-number-of-well-switches & A positive integer values that
stipulates the maximum number of times a well can switch to the same
control. & 3 \\
98 & --maximum-water-saturation & A real positive value that defines the
maximum water saturation. & 1.0 \\
99 & --metis-params & A character string. Configuration of Metis
partitioner. You can request a configuration to be read from a JSON file
by giving the filename here, ending with `.json.' See
http://glaros.dtc.umn.edu/gkhome/fetch/sw/metis/manual.pdf for available
METIS options. & ``default'' \\
100 & --milu-variant & A defined character string that specifies which
variant of the modified ILU preconditioner ought to be used. Possible
variants are: & ``ILU'' \\
101 & --min-strict-cnv-iter & A positive integer that sets the minimum
number of Newton iterations before relaxed tolerances are used for the
CNV convergence criterion. & -1 \\
102 & --min-strict-mb-iter & An integer value defining the number of
Newton iterations before relaxed tolerances can be used for the MB
convergence criterion. Default -1 means that the relaxed tolerance is
used when maximum number of Newton iterations are reached. & -1 \\
103 & --min-time-step-based-on-newton-iterations & A real positive value
that sets the minimum time step size (in days for field and metric unit
and hours for lab unit) can be reduced to based on Newton iteration
counts. & 0.0 \\
104 & --min-time-step-before-shutting-problematic-wells-in-days & A real
positive value that sets the minimum time step size in days for which
problematic wells are not shut. Time steps below this value will result
in problematic wells being shut-in by the simulator. & 0.01 \\
105 & --min-time-step-size & A real positive value that sets the minimum
size to which all time steps are limited to in seconds. & 0.0 \\
106 & --network-max-iterations & Replaced by
--network-max-outer-iterations. & 200 \\
107 & --network-max-outer-iterations & An integer value. Maximum outer
number of iterations in the network solver before giving up. & 10 \\
108 & --network-max-pressure-update-in-bars & A real value. Maximum
pressure update in the inner network pressure update iterations. & 5 \\
109 & --network-max-strict-iterations & Replaced by
--network-max-strict-outer-iterations. & 100 \\
110 & --network-max-strict-outer-iterations & An integer value. Maximum
outer iterations in network solver before relaxing tolerance. & 10 \\
111 & --network-max-sub-iterations & An integer value. Maximum number of
sub-iterations to update network pressures (within a single well/group
control update). & 20 \\
112 & --newton-max-error & A real positive value that sets the maximum
error tolerated by the Newton method to which does not cause an abort. &
1e+100 \\
113 & --newton-max-iterations & A positive integer that defines the
maximum number of Newton iterations per time step used by the simulator.
& 20 \\
114 & --newton-max-relax & A real positive value that sets the maximum
relaxation factor of a Newton iteration used by the simulator. & 0.5 \\
115 & --newton-min-iterations & A positive integer that sets the minimum
number of Newton iterations per time step used by the simulator. & 2 \\
116 & --network-pressure-update-damping-factor & A real value. Damping
factor in the inner network pressure update iterations. & 0.1 \\
117 & --newton-relaxation-type & A character string. The type of
relaxation used by the Newton method. Valid options are: dampen or sor.
& ``dampen'' \\
118 & --newton-target-iterations & A positive integer that sets the
`optimum' number of Newton iterations per time step. & 10 \\
119 & --newton-tolerance & A real positive value that sets the maximum
raw error tolerated by the Newtonmethod for considering a solution to be
converged. & 0.01 \\
120 & --newton-verbose & A boolean value that specifies whether the
Newton method should inform the user about its progress (true) or not
(false). & true \\
121 & --newton-write-convergence & A boolean value that specifies
whether to write the convergence behaviour of the Newton method to a VTK
file (true) or not (false). & false \\
122 & --nldd-local-linear-solver & A defined character string.
Configuration of NLDD local linear solver. Valid options are: ilu0
(default), dilu, cpr\_quasiimpes and amg. Alternatively, you can request
a configuration to be read from a JSON file by giving the filename here,
ending with `.json.'. & ``ilu0'' \\
123 & --nldd-local-linear-solver-max-iter & An integer value. The
maximum number of iterations of the NLDD local linear solver. & 200 \\
124 & --nldd-local-linear-solver-reduction & A real value. The minimum
reduction of the residual which the NLDD local linear solver must
achieve. & 0.01 \\
125 & --nldd-num-initial-newton-iter & A positive integer that sets the
number of initial global Newton iterations when running the NLDD
nonlinear solver. & 1 \\
126 & --nonlinear-solver & A character string that specifies the
nonlinear solver. Valid choices are newton or nldd. & ``newton'' \\
127 & --num-local-domains & A positive integer that sets the number of
local domains for NLDD nonlinear solver. Note this is an experimental
feature in the current release 2023.10 that is expected to be more
complete and tested by the 2024.10 release. & 0 \\
128 & --num-overlap & An integer value. Numbers of layers overlap in
parallel partition. & 1 \\
129 & --num-pressure-points-equil & A positive integer that sets the
number of pressure points (in each direction) in tables used for
equilibration. & 2000 \\
130 & --num-satfunc-consistency-sample-points & An integer value.
Maximum number of reported failures for each individual saturation
function consistency check. & 5 \\
131 & --nupcol-group-rate-tolerance & A real value. Tolerance for
acceptable changes in VREP/RAIN group rates. & 0.001 \\
132 & --opencl-ilu-parallel & A boolean value that if set to true then
parallelize the ILU decomposition and application on GPU, or not
(false). & true \\
133 & --opencl-platform-id & A positive integer that specifies the
platform identification (``ID'') for the openSolver. Use the ``clinfo''
command to determine valid IDs. & 0 \\
134 & --output-dir & A character string that defines the directory to
which OPM Flow is to write the commercial simulator compatible output
files (restart, summary files etc.). The default value results in the
files be written to the same directory as the input file. & ``\,'' \\
135 & --output-extra-convergence-info & A defined character string that
specifies whether to provide additional convergence output to separate
files for diagnostic purposes. The available options are: Options can be
combined with commas, e.g.''steps,iterations'' for multiple outputs. The
default value of ``none'' prevents the two files from being written out,
for better compatibility with the commercial simulator. & ``none'' \\
136 & --output-interval & A positive integer that specifies the number
of report steps between two consecutive writes of restart data. & 1 \\
137 & --output-mode & A defined character string that defines the output
to the \emph{.PRT and }.DBG files: For example to just output logging
information use: --output-mode=``log''. or --output-mode=false &
``all'' \\
138 & --owner-cells-first & A boolean value that determines if the cells
owned by a rank should be ordered before ghost/overlapping the cells
(true), or not (false). & true \\
139 & --parameter-file & A character string that defines the name of a
parameter file which contains the simulator's set of run-time
parameters, as listed in this table. & ``\,'' \\
140 & --parsing-strictness & Set strictness of parsing process.
Available options are: Default: ``normal'' & ``normal'' \\
141 & --partition-method & An integer value. Choose partitioning
strategy: 0=simple, 1=Zoltan, 2=METIS, 3=Zoltan with all cells of well
represented by one vertex. & 3 \\
142 & --pre-solve-network & A boolean value. Pre solve and iterate the
network model at start-up. & true \\
143 & --predetermined-time-steps-file & A file with a list of
predetermined time step sizes (one time step per line). & ``\,'' \\
144 & --pressure-max & A positive real value. Maximum absolute pressure.
& 1e+99 \\
145 & --pressure-min & A real value. Minimum absolute pressure. &
-1e+99 \\
146 & --pressure-scale & A positive real value. Scaling of pressure
primary variable. & 1 \\
147 & --pri-var-oscilation-threshold & A real positive value that
defines the threshold value for the primary variable switching
conditions after its meaning has switched to hinder oscillations. &
1e-05 \\
148 & --print-parameters & A positive integer value that request that
the run time parameters be printed at the start of the run: & 2 \\
149 & --project-saturations & A boolean value that determines if all the
fluid saturations should be scaled to ensure the values are in the
interval (0, 1), including runs that use solvents (true), or not
(false). & false \\
150 & --regularization-factor-wells & A real positive value that defines
the ``regularization factor'' for wells. & 100 \\
151 & --relaxed-linear-solver-reduction & A real positive value that
defines the minimum reduction of the residual which the linear solver
need to achieve for the solution to be accepted. & 0.01 \\
152 & --relaxed-max-pv-fraction & A real positive value that defines The
fraction of the pore volume of the reservoir where the volumetric error
(CNV) might be violated during strict Newton iterations. & 0.03 \\
153 & --relaxed-pressure-tol-msw & A real positive value that sets the
relaxation tolerance for the multi-segment well pressure solution in
Pascals. & 10000 \\
154 & --relaxed-well-flow-tol & A real positive value that sets the
relaxation tolerance for the well flow residual in reservoir cubic
metres (rm3). & 0.001 \\
155 & --restart-time & A real positive value that sets the simulation
time at which a restart should be attempted {[}s{]}. & -1e+35 \\
156 & --restart-writing-interval & A positive integer value that sets
the frequency at which time steps are serialized to disk. & 16777215 \\
157 & --save-file & A character string that specifies the FileName for
.OPMRST file used for saving serialized state. If empty, CASENAME.OPMRST
is used. & ``\,'' \\
158 & --save-step & A character string that determines if the simulator
should save the serialized state of the OPM Flow simulator at one or
more report steps to a special *.OPMRST file. This is in addition to the
normal restart files written, and consumes significantly more space than
the normal restart files, but restarting OPM Flow from this file using
the --load-step option deviates less from the original run, compared to
restarting using the normal restart features. The files produced are not
compatible with other simulators, and also will not be compatible
between different releases of OPM Flow The parameter should be set to
one of the following: The default value of ``\,'' does not write
anything to the OPM Flow specific restart file. & ``\,'' \\
159 & --scale-linear-system & A boolean value that switches on (true) or
off (false) the scaling of linear system of equations according to the
equation scale and primary variable types. & false \\
160 & --sched-restart & A boolean value that determines for a restart
run if the case should initialize wells and groups from the historical
SCHEDULE section (true), or from the well and group data on the restart
file (false). Note that the commercial simulator always uses data from
the restart file (false). & false \\
161 & --serial-partitioning & A boolean value that determines if
partitioning for parallel runs on a single process (true), or not
(false). & false \\
162 & --shut-unsolvable-wells & A boolean value that determines if the
simulator should shut unsolvable wells (true), or not (false). & true \\
163 & --slave & A boolean value. Specify if the simulation is a slave
simulation in a master-slave simulation. Note that reservoir coupling is
not currently supported. This is a developer option that should not be
used. & false \\
164 & --solve-welleq-initially & A boolean value that determines if the
simulator should fully solve the well equations before each iteration of
the reservoir model (true), or not (false). Note that the well equations
are always added to the full system and solved until convergence. &
true \\
165 & --solver-continue-on-convergence-failure & A boolean value that
stipulates if the simulator should continue (true) instead of stopping
(false) when the minimum solver time step is reached. & false \\
166 & --solver-growth-factor & A real positive value that specifies the
growth factor a time step can be increased by when recovering from one
or more time step chops, subject to the maximum allowable time step size
set by the --solver-max-time-step-in-days parameter. For example, if the
current time step has converged at 5 days after a time step chop, and
--solver-growth-factor is set to the default value of 2.0, then the next
time step will be 2.0 * 5 days, that is at 10 days. & 2.0 \\
167 & --solver-max-growth & A real positive value that specifies the
maximum growth factor a time step can be increased by after a report
time step, subject to the maximum allowable time step size set by the
--solver-max-time-step-in-days parameter. Thus, if the current time step
has converged at 5 days after at a report step, and --solver-max-growth
is set to the default value of 3.0, then the next time step will be 3.0
* 5 days, that is at 15 days. & 3.0 \\
168 & --solver-max-restarts & A positive integer that defines the number
of allowed consecutive restarts (or time step chops) before the
simulation is terminated. & 10 \\
169 & --solver-max-time-step-in-days & A real positive double precision
value that specifies the maximum allowed time step size in days. &
365 \\
170 & --solver-min-time-step & A real positive double precision value
that specifies the minimum size of a time step in days for field and
metric units, and hours for laboratory units If a time step cannot
converge without getting cut below this time step size the simulator
will stop. & 1e-12 \\
171 & --solver-restart-factor & A real positive double precision value
that sets the time step chop factor of the time step after a convergence
failure. For example, if the current non-convergent time step is 30 days
and --solver-restart-factor is set to the default value, then the time
step will be repeated using 0.33 * 30 days as the time step, that is 9.9
days. & 0.33 \\
172 & --solver-verbosity & A positive integer that specifies the
``chattiness'' of the non-linear solver. & 1 \\
173 & --strict-inner-iter-wells & A positive integer that specifies the
number of inner well iterations with strict tolerance. & 40 \\
174 & --strict-outer-iter-wells & A positive integer that specifies the
number of Newton iterations for which wells are checked with strict
tolerance. & 6 \\
175 & --temperature-max & A real positive value that sets the maximum
absolute temperature. & 1e+09 \\
176 & --temperature-min & A real positive value that sets the minimum
absolute temperature. & 0 \\
177 & --threads-per-process & A positive integer value that stipulates
the maximum number of threads to be instantiated per process (`-1' means
`automatic'). & 2 \\
178 & --time-step-after-event-in-days & A real double precision value
that sets the maximum allowed time step after an event; for example,
when a well is open or closed etc. The default value of -1 means that
events to do effect the time stepping. & -1 \\
179 & --time-step-control & A defined character string that defines the
time stepping control algorithm and is set to one of the following: &
``pid+newtoniteration'' \\
180 & --time-step-control-decay-damping-factor & A real positive value
that specifies the decay rate a time step can be decreased by after the
number target iterations has been exceeded & 1.0 \\
181 & --time-step-control-decay-rate & A real positive value that
specifies the decay rate a time step size can be decreased by after the
number of target iterations has been exceeded & 0.75 \\
182 & --time-step-control-file-name & A character string enclosed in
quotes that specifies a filename where time steps are specified. The
default is the character string timesteps. For instance time steps can
be generated by the ecl\_summary application in libecl as per the
following UNIX command line: path\_to\_libecl\_applications/ecl\_summary
DECK TIME \textgreater{} filename Where: DECK is the name of the data
deck you want to get the time steps from, TIME tells the application to
return the timing for the run, and ``filename'' is the name of the file
the times are piped to. & ``timesteps'' \\
183 & --time-step-control-growth-damping-factor & A real positive value
that specifies the growth rate of the time step increase when the number
of target iterations are undercut. & 3.2 \\
184 & --time-step-control-growth-rate & A real positive value that
specifies the growth rate of the time step size when the number of
target iterations are undercut. & 1.25 \\
185 & --time-step-control-safety-factor & A real value to be multiplied
with the time step control tolerance to ensure that the target relative
change is lower than the tolerance. & 0.8 \\
186 & --time-step-control-target-iterations & A positive integer that
defines the number of linear iterations which the time step control
scheme should aim for (if applicable). & 30 \\
187 & --time-step-control-target-newton-iterations & A positive integer
that specifies the number of Newton iterations which the time step
control scheme should aim for (if applicable). & 8 \\
188 & --time-step-control-tolerance & A real double precision value that
sets the tolerance for PID (only used with the pid and pid+ options
defined by the --time-step-control option). & 0.1 \\
189 & --time-step-verbosity & A positive integer that specifies the
``chattiness'' during the time integration. & 1 \\
190 & --tolerance-cnv & A real positive double precision value that
specifies the maximum non-linear tolerance error. This is the local
convergence tolerance (maximum of local saturation errors). & 0.01 \\
191 & --tolerance-cnv-energy & A real value. Local energy convergence
tolerance (maximum of local energy errors). & 0.01 \\
192 & --tolerance-cnv-energy-relaxed & A real value. Relaxed local
energy convergence tolerance that applies for iterations after the
iterations with the strict tolerance. & 1.0 \\
193 & --tolerance-cnv-relaxed & A real positive value that defines the
relaxed local convergence tolerance that applies for iterations after
the iterations with the strict tolerance. & 1.0 \\
194 & --tolerance-energy-balance & A real value. Tolerated energy
balance error relative to (scaled) total energy present. & 1e-07 \\
195 & --tolerance-energy-balance-relaxed & A real value. Relaxed
tolerated energy balance error that applies for iterations after the
iterations with the strict tolerance. & 1e-06 \\
196 & --tolerance-mb & A real positive double precision value that sets
the maximum mass balance error, that is the tolerated mass balance error
relative to total mass present. & 1e-07 \\
197 & --tolerance-mb-relaxed & A positive real value. Relaxed tolerated
mass balance error that applies for iterations after the iterations with
the strict tolerance. & 1e-06 \\
198 & --tolerance-pressure-ms-wells & A real positive double precision
value that specifies the tolerance for the pressure equations for
multi-segment wells. & 1000 \\
199 & --tolerance-well-control & A real positive double precision value
that sets the maximum tolerance for the well control equations. &
1e-07 \\
200 & --tolerance-wells & A real positive double precision value that
defines the maximum non-linear error for the well equations. & 1e-04 \\
201 & --update-equations-scaling & A boolean value that switches on
(true) or off (false) the updating of the scaling factors for mass
balance equations during the simulation. & false \\
202 & --use-average-density-ms-wells & A boolean value that specifies
whether to approximate segment densitities by averaging over the segment
and its outlet (true) or not (false). & false \\
203 & --use-gmres & A boolean value that when set to true OPM Flow will
use Generalized Minimal Residual (``GMRES'') {[}Y. Saad, A flexible
inner-outer preconditioned GMRES algorithm, SIAM J. Sci. Statist.
Comput.,14, (1993).{]} and {[}Y. Saad and M.H. Schultz, GMRES: A
generalized minimum residual algorithm for solving nonsymmetric linear
systems, SIAM J. Sci. Statist. Comput., 7 (1986), pp.~856\{869.{]}
solver instead of Biconjugate Gradient Stabilized (``BiCGSTAB'') {[}Van
der Vorst, H. A. (1992). ``Bi-CGSTAB: A Fast and Smoothly Converging
Variant of Bi-CG for the Solution of Nonsymmetric Linear Systems''. SIAM
J. Sci. Stat. Comput. 13 (2): 631--644. doi:10.1137/0913035.
hdl:10338.dmlcz/104566{]} and {[}Sleijpen, G. L. G.; Fokkema, D. R.
(November 1993). ``BiCGstab(l) for linear equations involving
unsymmetric matrices with complex spectrum'' (PDF). Electronic
Transactions on Numerical Analysis. Kent, OH: Kent State University. 1:
11--32. ISSN 1068-9613.{]} as the linear solver within the Newton
iterations. & false \\
204 & --use-implicit-ipr & A boolean value. Compute implicit IPR for
stability checks and stable solution search. & true \\
205 & --use-multisegment-well & A boolean value that when set to true
the simulator will use the well model for multi-segment wells instead of
the one for single-segment wells. & true \\
206 & --use-update-stabilization & A boolean value that switches on
(true) or off (false) the stabilized Newton option, that attempts to
detect and correct oscillations or stagnation during the Newton
iterations. This option may improve convergence for some cases. &
true \\
207 & --water-only-threshold & A real positive value that defines the
saturation threshold, for which cells with water saturations above or
equal to this threshold are considered one-phase water only. & 1.0 \\
208 & --zoltan-imbalance-tol & Replaced by --imbalance-tol. & 1.1 \\
208 & --zoltan-params & A character string that specifies the
configuration of the Zoltan partitioner. Valid options are: graph,
hypergraph or scotch. Alternatively, you can request a configuration to
be read from a JSON file by giving the filename here, ending with
`.json.' See https://sandialabs.github.io/Zoltan/ug\_html/ug.html for
available Zoltan options. & ``graph'' \\
209 & --zoltan-phg-edge-size-threshold & A real value. Low-level
threshold fraction in the range {[}0,1{]} controlling which hypergraph
edge to omit. Used if --zoltan-params=``graph'' or if
--zoltan-params=``hypergraph''. & 0.35 \\
Notes: & & & \\
\end{longtable}

\emph{Table A.1: OPM Flow 2025-04 Command Line Options}

The command line syntax for running OPM Flow is:

\begin{verbatim}

flow [OPTIONS] CASENAME
\end{verbatim}

and typing the following command from your terminal:

\begin{verbatim}

flow CASENAME
\end{verbatim}

will start the simulator and run the case specified by CASENAME.DATA.

OPM Flow accepts command line options to control various aspects of the
simulator's run time behavior, as described in the table below. To give
a command line option use ``key=value'' syntax, with no spaces around
the equals sign. It is also possible to put multiple options together in
a parameter file. To do so, put one option on each line of the file with
``.param'' as the extension and pass that filename as a command line
parameter to OPM Flow using:

\begin{verbatim}

flow --parameter-file=CASENAME.param
\end{verbatim}

to instruct OPM Flow to read the parameter file.

\begin{longtable}[]{@{}
  >{\raggedright\arraybackslash}p{(\linewidth - 6\tabcolsep) * \real{0.2500}}
  >{\raggedright\arraybackslash}p{(\linewidth - 6\tabcolsep) * \real{0.2500}}
  >{\raggedright\arraybackslash}p{(\linewidth - 6\tabcolsep) * \real{0.2500}}
  >{\raggedright\arraybackslash}p{(\linewidth - 6\tabcolsep) * \real{0.2500}}@{}}
\toprule\noalign{}
\begin{minipage}[b]{\linewidth}\raggedright
OPM Flow 2024-10 Command Line Options
\end{minipage} & \begin{minipage}[b]{\linewidth}\raggedright
\end{minipage} & \begin{minipage}[b]{\linewidth}\raggedright
\end{minipage} & \begin{minipage}[b]{\linewidth}\raggedright
\end{minipage} \\
\midrule\noalign{}
\endhead
\bottomrule\noalign{}
\endlastfoot
No. & Variable Name & Description & Default \\
1 & -h or --help & A character string that causes OPM Flow to print a
help message that gives a brief description of the available command
line parameters. This now only prints the active user facing command
line options. & N/A \\
2 & --help-all & Prints all the command line options included in the
release, including obsolete, hidden and deprecated options. & N/A \\
3 & --accelerator-mode & A defined character string that selects the
linear solver, usage:
`--accelerator-mode={[}none\textbar cusparse\textbar opencl\textbar amgcl\textbar rocalution\textbar rocsparse{]}'.
& ``none'' \\
4 & --action-parsing-strictness & A defined character string. Set
strictness of parsing process for ActionX and PyAction. Available
options are normal (do not apply keywords that have not been tested for
ActionX or PyAction) and low (try to apply all keywords, beware: the
simulation outcome might be incorrect). & ``normal'' \\
5 & --allow-distributed-wells & Allow the perforations of a well to be
distributed to interior of multiple processes. & false \\
6 & --alternative-well-rate-init & Use alternative well rate
initialization procedure. & true \\
7 & --bda-device-id & Choose device ID for cusparseSolver or
openclSolver, use `nvidia-smi' or `clinfo' commands to determine valid
IDs. & 0 \\
8 & --check-group-constraints-inner-well-iterations & A boolean value.
Allow checking of group constraints during inner well iterations. &
true \\
9 & --continue-on-convergence-error & Continue with a non-converged
solution instead of giving up if we encounter a time step size smaller
than the minimum time step size. & false \\
10 & --convergence-monitoring & A boolean value. Enable convergence
monitoring. & false \\
11 & --convergence-monitoring-cut-off & An integer value. Cut off limit
for convergence monitoring. & 6 \\
12 & --convergence-monitoring-decay-factor & A real value. Decay factor
for convergence monitoring. & 0.75 \\
13 & --cpr-reuse-interval & A positive integer that sets the reuse
preconditioner interval. Used when --cpr-reuse-setup is set to 4, then
the preconditioner will be fully recreated instead of reused every N
linear solve, where N is this parameter. & 30 \\
14 & --cpr-reuse-setup & A positive integer that defines if the CPR
solver should re-use the AMG setup. Valid options are: Changed the
default value from three to four. & 4 \\
15 & --dbph-max-rel & A real positive value that sets the maximum
relative change of the bottom-hole pressure in a single iteration. &
1.0 \\
16 & --debug-emit-cell-partition & A boolean value that determines
whether (true) or not (false) to emit cell partitions as a debugging
aid. & false \\
17 & --dp-max-rel & A real positive double precision value that sets the
maximum allowed relative change of pressure per iteration. & 0.3 \\
18 & --ds-max & A real positive double precision value that sets the
maximum allowed change in saturation per iteration. & 0.2 \\
19 & --dwell-fraction-max & A real positive value that sets the maximum
allowed change in well's volume fraction per iteration. & 0.2 \\
20 & --ecl-deck-file-name & A character string that defines the name of
the OPM Flow input file which contains the simulator's ECLIPSE formatted
input deck to be simulated. Note that there must be no spaces in the
filename. If there are spaces, which can typically occur on Windows
based systems, then the filename should be enclosed in quotes; thus
instead of FILE NAME use ``FILE NAME''. & ``\,'' \\
21 & --ecl-output-double-precision & A boolean value that switches on
(true) or off (false) double precision in restart files. Useful for
`perfect' restarts. & false \\
22 & --ecl-output-interval & An integer value defining the number of
report steps that ought to be skipped between two writes of restart
files results. & -1 \\
23 & --edge-weights-method & A defined positive integer value that
defines the edge-weighing strategy: & 1 \\
24 & --enable-adaptive-time-stepping & A boolean value that turns on
(true) or off (false) adaptive time stepping. If set to false the report
time steps provided by the input deck are used. & true \\
25 & --enable-async-ecl-output & A boolean value that sets the output
mode to be asynchronous (true), letting the simulator continue computing
the next time step while writing results to restart and summary files in
the commercial simulator's format on a separate thread. & true \\
26 & --enable-async-vtk-output & A boolean value that sets the output
mode to be asynchronous (true), letting the simulator continue computing
the next time step while writing results to the VTK output files on a
separate thread. & true \\
27 & --enable-drift-compensation & A boolean value that enables (true)
or disables (false) partial compensation of systematic mass losses via
the source term of the next time step. & false \\
28 & --enable-dry-run & A boolean value that specifies if the simulation
should actually run (true), or just check the input deck (false). This
option is equivalent to activating the NOSIM keyword in the RUNSPEC
section of the input deck (see section Error: Reference source not found
Error: Reference source not found). The default value of ``auto' will
use whatever is stipulated in the input deck via the NOSIM keyword. &
``auto'' \\
29 & --enable-ecl-output & A boolean value that specifies if the binary
output files (restart, summary files, etc.) should be written in the
commercial simulator's format (true), or OPM Flow's format (false). &
true \\
30 & --enable-esmry & A boolean value that switches on (true) or off
(false) the output of SUMMARY vectors to the ESMRY file for fast loading
of summary data. & false \\
31 & --enable-gravity & A boolean value that switches on (true) or off
(false) the use of the gravity correction for the pressure gradients. &
true \\
32 & --enable-grid-adaptation & A boolean value that enables (true) or
disables (false) adaptive grid refinement/coarsening. Note that Eclipse
grid output will be disabled if adaptive grid refinement/coarsening is
enabled. & false \\
33 & --enable-intensive-quantity-cache & A boolean value that switches
on (true) or off (false) the caching of intensive quantities. & true \\
34 & --enable-logging-fallout-warning & A boolean value that sets the
developer option to see whether logging was on non-root processors
(true). If set to true output will be appended to the \emph{.DBG or
}.PRT files. & false \\
35 & --enable-opm-rst-file & A boolean value set to true or false to
write OPM specific data sets to the commercial simulators restart file
to enable restart of OPM Flow runs by OPM Flow (true), or not to write
the data (false). & false \\
36 & --enable-storage-cache & A boolean value that turns on (true) or
off (false) storing previous storage terms and avoid re-calculating
them. & true \\
37 & --enable-terminal-output & A boolean value that turns on (true) or
off (false) high-level information about the simulation's progress to
the terminal & true \\
38 & --enable-thermodynamic-hints & A boolean value that enables (true)
or disables (false) thermodynamic hints. & false \\
39 & --enable-tuning & A boolean value that instructs OPM Flow to read
the time stepping parameters from the Error: Reference source not found
keyword in the input deck, if set to true. Note that only the first
record of the TUNING keyword is processed. & false \\
40 & --enable-vtk-output & A boolean value that turns on (true) or off
(false) a global switch for writing VTK files. & false \\
41 & --enable-well-operability-check & A boolean value that enables
(true) checking of a well's operating status, or disables (false) the
checking. & true \\
42 & --enable-well-operability-check-iter & A boolean value that enables
(true) checking of a well's operating status during iterations, or
disables (false) the checking during iterations. & false \\
43 & --enable-write-all-solutions & A boolean value that turns on (true)
or off (false) the writing of all solutions to disk instead of only the
ones for the report steps. & false \\
44 & --end-time & A real value. The simulation time at which the
simulation is finished {[}s{]}. & 1e+100 \\
45 & --explicit-rock-compaction & A boolean value. Use pressure from end
of the last time step when evaluating rock compaction (true) or not
(false). & false \\
46 & --external-partition & A quoted character string. Name of file from
which to load an externally generated partitioning of the model's active
cells for MPI distribution purposes. If empty, the built-in partitioning
method will be employed. & ``\,'' \\
47 & --force-disable-fluid-in-place-output & A boolean value that
instructs OPM Flow not to print the Fluid In-Place report after each
report time step (true) or not (false). Note this parameter will
override the print request in the input deck. & false \\
48 & --force-disable-resv-fluid-in-place-output & A boolean value that
instructs OPM Flow to not to print the Reservoir Volume Fluid In-Place
report after each report time step (true) or not (false). Note this
parameter will override the print request in the input deck. & false \\
49 & --full-time-step-initially & A boolean value that instructs OPM
Flow to always attempt to finish a report step using a single time step
(true) or not (false). & false \\
50 & --ignore-keywords & A character string that defines a list of
keywords which should be ignored by OPM Flow. The keyword string should
be enclosed in quotes and each keyword separated by a colon, that is
`:'. & ``\,'' \\
51 & --ilu-fillin-level & A positive integer value that sets the fill in
level for the ILU preconditioner. & 0 \\
52 & --ilu-redblack & A boolean value that instructs OPM Flow to use
red-black partitioning for the ILU preconditioner. & false \\
53 & --ilu-relaxation & A real positive double precision value that sets
the relaxation factor of the linear solver's ILU preconditioner. &
0.9 \\
54 & --ilu-reorder-spheres & A boolean value that specifies OPM Flow to
reorder the entries of the matrix in the red-black ILU preconditioner in
spheres starting at an edge (true) or not (false). If false the original
ordering is preserved in each color. Otherwise try to ensure D4 ordering
(in a 2D structured grid, the diagonal elements are consecutive). &
false \\
55 & --imbalance-tol & A real value. Tolerable imbalance of the load
balancing. & 1.1 \\
56 & --initial-time-step-in-days & A real double precision value that
sets the size of initial time step in days. & 1.0 \\
57 & --initial-time-step-size & A real double precision value that sets
the size of initial time step in seconds. & 86400 \\
58 & --inj-mult-damp-mult & A real value. Injection multiplier dampening
factor (dampening multiplied by this each time oscillation is detected).
& 0.9 \\
59 & --inj-mult-min-damp-factor & A real value. Minimum injection
multiplier dampening factor (maximum dampening level). & 0.05 \\
60 & --inj-mult-osc-threshold & A real value. Injection multiplier
oscillation threshold (used for multiplier dampening). & 0.1 \\
61 & --input-skip-mode & A defined character string. Set compatibility
mode for the
\hyperref[REF_HEADING_KEYWORD_SKIP100_11_3]{SKIP100}/\hyperref[REF_HEADING_KEYWORD_SKIP300_11_3]{SKIP300}
keywords. Options are: 100 (skip
\hyperref[REF_HEADING_KEYWORD_SKIP100_11_3]{SKIP100}..ENDSKIP, keep
\hyperref[REF_HEADING_KEYWORD_SKIP300_11_3]{SKIP300}..ENDSKIP), 300
(skip \hyperref[REF_HEADING_KEYWORD_SKIP300_11_3]{SKIP300}..ENDSKIP,
keep \hyperref[REF_HEADING_KEYWORD_SKIP100_11_3]{SKIP100}..ENDSKIP) or
all (skip both
\hyperref[REF_HEADING_KEYWORD_SKIP100_11_3]{SKIP100}..ENDSKIP and
\hyperref[REF_HEADING_KEYWORD_SKIP300_11_3]{SKIP300}..ENDSKIP). &
``100'' \\
62 & --linear-solver & A defined quoted character string that sets the
configuration of the linear solver, valid values are: Option (1) extends
the existing Constrained Pressure Residual (``CPR'') preconditioner to
include wells. This option can also be invoked via the CPR keyword in
the RUNSPEC section; however, the command line parameter takes
precedence. For option (10) one enters a character string enclosed in
quotes that defines the filename of a JSON configuration file for the
flexible linear solver system. In this case the file extension should be
`.json'. Note that the *.PRT file contains the ``Property tree for the
linear solver'' listing, which is the JSON specification of the current
case, and can be used to configure a user specific linear solver JSON
file. & ``cprw'' \\
63 & --linear-solver-ignore-convergence-failure & A boolean value that
if set to true convergences failures in the linear solver are ignored.
This option should be used with care, as the results might be
unreliable. & false \\
64 & --linear-solver-max-iter & A positive integer value that defines
the maximum number of linear iterations. & 200 \\
65 & --linear-solver-print-json-definition & Write the JSON definition
of the linear solver setup to the DBG file. & true \\
66 & --linear-solver-reduction & A real positive double precision value
that sets the minimum reduction of the residual for the linear solver.
The linear solver convergences when the residual is reduced
sufficiently. The simulator now overrides the default reduction to be
0.005 instead of 0.01, if the linear solver has been set to one of the
cpr options, in a similar manner as how the default maximum number of
linear iterations for the cpr and cprw options are changed to 20 instead
of 200, unless specified by the command line option by the user. &
0.01 \\
67 & --linear-solver-restart & A positive integer value that sets the
number of iterations after which GMRES is restarted. & 40 \\
68 & --linear-solver-verbosity & A positive integer value that defines
the output from linear solver: & 0 \\
69 & --load-file & A character string enclosed in quotes that defines
the FileName for the .OPMRST file used to load the serialized state. If
empty, CASENAME.OPMRST is used. & ``\,'' \\
70 & --load-step & An integer value that determines if the simulator
should load the serialized state from OPM Flow's version of the restart
file, and should be set to either a specific report step, or 0 to load
the last stored report step. The default value of -1 does not load the
data from the OPM Flow specific restart file. OPM Flow's version of the
restart file, is written using the --save-step=N option. & -1 \\
71 & --local-domains-ordering-measure & A character string enclosed in
quotes that defines the subdomain ordering measure. Allowed values are
`maxpressure', `averagepressure' and `residual'. & ``maxpressure'' \\
72 & --local-domains-partitioning-imbalance & A real positive double
precision value that sets the subdomain partitioning imbalance
tolerance. 1.03 is 3 percent imbalance. & 1.03 \\
73 & --local-domains-partitioning-method & A character string enclosed
in quotes that defines the subdomain partitioning method. Allowed values
are `zoltan', `simple', and the name of a partition file ending with
`.partition'. & ``zoltan'' \\
74 & --local-solve-approach & A character string enclosed in quotes that
defines the local solve approach. Valid choices are `jacobi' and
`gauss-seidel'. & ``gauss-seidel'' \\
75 & --local-tolerance-scaling-cnv & A real positive double precision
value that sets the convergence tolerance for local solves. Set to lower
than 1.0 to use stricter convergence tolerance for local solves. &
0.1 \\
76 & --local-tolerance-scaling-mb & A real positive double precision
value that sets the convergence tolerance for local solves. Set to lower
than 1.0 to use stricter convergence tolerance for local solves. & 1 \\
77 & --local-well-solve-control-switching & A boolean value that allows
(true) or disallows (false) control switching during local well
solutions. & true \\
78 & --matrix-add-well-contributions & A boolean value that if set to
true explicitly specifies the influences of wells between cells in the
Jacobian and preconditioner matrices. & false \\
79 & --max-inner-iter-ms-wells & A positive integer value that defines
the maximum number of inner iterations for multi-segment wells. & 100 \\
80 & --max-inner-iter-wells & A positive integer value that defines the
maximum number of inner iterations for standard wells. & 50 \\
81 & --max-local-solve-iterations & A positive integer value that
defines the maximum number of iterations for local solves with NLDD
nonlinear solver. & 20 \\
82 & --max-newton-iterations-with-inner-well-iterations & A positive
integer that specifies the maximum Newton iterations with inner well
iterations. & 8 \\
83 & --max-pressure-change-ms-wells & A real positive value that defines
the maximum relative pressure change for a single iteration of the
multi-segment well model. & 1e+06 \\
84 & --max-residual-allowed & A real positive value that sets the
absolute maximum tolerance for residuals without cutting the time step
size. & 1e+07 \\
85 & --max-single-precision-days & A real positive value that sets the
maximum time step size where single precision floating point arithmetic
can be used for solving the linear systems of equations. & 20.0 \\
86 & --max-temperature-change & A real positive value that stipulates
the maximum absolute change of temperature in a single iteration. & 5 \\
87 & --max-time-step-divisions & A positive integer that specifies the
maximum number of divisions by two of the timestep size before the
simulation bails out. & 10 \\
88 & --max-time-step-size & A real positive value that defines the
maximum size to which all time steps are limited to in seconds. & inf \\
89 & --max-welleq-iter & A positive integer that defines the maximum
number of iterations to determine the solution to the well equations. &
30 \\
90 & --maximum-number-of-well-switches & A positive integer values that
stipulates the maximum number of times a well can switch to the same
control. & 3 \\
91 & --maximum-water-saturation & A real positive value that defines the
maximum water saturation. & 1.0 \\
92 & --metis-params & A character string. Configuration of Metis
partitioner. You can request a configuration to be read from a JSON file
by giving the filename here, ending with `.json.' See
http://glaros.dtc.umn.edu/gkhome/fetch/sw/metis/manual.pdf for available
METIS options. & ``default'' \\
93 & --milu-variant & A defined character string that specifies which
variant of the modified ILU preconditioner ought to be used. Possible
variants are: & ``ILU'' \\
94 & --min-strict-cnv-iter & A positive integer that sets the minimum
number of Newton iterations before relaxed tolerances are used for the
CNV convergence criterion. & -1 \\
95 & --min-strict-mb-iter & An integer value defining the number of
Newton iterations before relaxed tolerances can be used for the MB
convergence criterion. Default -1 means that the relaxed tolerance is
used when maximum number of Newton iterations are reached. & -1 \\
96 & --min-time-step-based-on-newton-iterations & A real positive value
that sets the minimum time step size (in days for field and metric unit
and hours for lab unit) can be reduced to based on Newton iteration
counts. & 0.0 \\
97 & --min-time-step-before-shutting-problematic-wells-in-days & A real
positive value that sets the minimum time step size in days for which
problematic wells are not shut. Time steps below this value will result
in problematic wells being shut-in by the simulator. & 0.01 \\
98 & --min-time-step-size & A real positive value that sets the minimum
size to which all time steps are limited to in seconds. & 0.0 \\
99 & --network-max-iterations & A positive integer that sets the maximum
number of iterations in the network solver before giving up. & 200 \\
100 & --network-max-strict-iterations & A positive integer that sets the
maximum iterations in network solver before relaxing tolerance. & 100 \\
101 & --newton-max-error & A real positive value that sets the maximum
error tolerated by the Newton method to which does not cause an abort. &
1e+100 \\
102 & --newton-max-iterations & A positive integer that defines the
maximum number of Newton iterations per time step used by the simulator.
& 20 \\
103 & --newton-max-relax & A real positive value that sets the maximum
relaxation factor of a Newton iteration used by the simulator. & 0.5 \\
104 & --newton-min-iterations & A positive integer that sets the minimum
number of Newton iterations per time step used by the simulator. & 2 \\
105 & --newton-relaxation-type & A character string. The type of
relaxation used by the Newton method. Valid options are: dampen or sor.
& ``dampen'' \\
106 & --newton-target-iterations & A positive integer that sets the
`optimum' number of Newton iterations per time step. & 10 \\
107 & --newton-tolerance & A real positive value that sets the maximum
raw error tolerated by the Newtonmethod for considering a solution to be
converged. & 0.01 \\
108 & --newton-verbose & A boolean value that specifies whether the
Newton method should inform the user about its progress (true) or not
(false). & true \\
109 & --newton-write-convergence & A boolean value that specifies
whether to write the convergence behaviour of the Newton method to a VTK
file (true) or not (false). & false \\
110 & --nldd-local-linear-solver & A defined character string.
Configuration of NLDD local linear solver. Valid options are: ilu0
(default), dilu, cpr\_quasiimpes and amg. Alternatively, you can request
a configuration to be read from a JSON file by giving the filename here,
ending with `.json.'. & ``ilu0'' \\
111 & --nldd-local-linear-solver-max-iter & An integer value. The
maximum number of iterations of the NLDD local linear solver. & 200 \\
112 & --nldd-local-linear-solver-reduction & A real value. The minimum
reduction of the residual which the NLDD local linear solver must
achieve. & 0.01 \\
113 & --nldd-num-initial-newton-iter & A positive integer that sets the
number of initial global Newton iterations when running the NLDD
nonlinear solver. & 1 \\
114 & --nonlinear-solver & A character string that specifies the
nonlinear solver. Valid choices are newton or nldd. & ``newton'' \\
115 & --num-local-domains & A positive integer that sets the number of
local domains for NLDD nonlinear solver. Note this is an experimental
feature in the current release 2023.10 that is expected to be more
complete and tested by the 2024.10 release. & 0 \\
116 & --num-pressure-points-equil & A positive integer that sets the
number of pressure points (in each direction) in tables used for
equilibration. & 2000 \\
117 & --opencl-ilu-parallel & A boolean value that if set to true then
parallelize the ILU decomposition and application on GPU, or not
(false). & true \\
118 & --opencl-platform-id & A positive integer that specifies the
platform identification (``ID'') for the openSolver. Use the ``clinfo''
command to determine valid IDs. & 0 \\
119 & --output-dir & A character string that defines the directory to
which OPM Flow is to write the commercial simulator compatible output
files (restart, summary files etc.). The default value results in the
files be written to the same directory as the input file. & ``\,'' \\
120 & --output-extra-convergence-info & A defined character string that
specifies whether to provide additional convergence output to separate
files for diagnostic purposes. The available options are: Options can be
combined with commas, e.g.''steps,iterations'' for multiple outputs. The
default value of ``none'' prevents the two files from being written out,
for better compatibility with the commercial simulator. & ``none'' \\
121 & --output-interval & A positive integer that specifies the number
of report steps between two consecutive writes of restart data. & 1 \\
122 & --output-mode & A defined character string that defines the output
to the \emph{.PRT and }.DBG files: For example to just output logging
information use: --output-mode=``log''. or --output-mode=false &
``all'' \\
123 & --owner-cells-first & A boolean value that determines if the cells
owned by a rank should be ordered before ghost/overlapping the cells
(true), or not (false). & true \\
124 & --parameter-file & A character string that defines the name of a
parameter file which contains the simulator's set of run-time
parameters, as listed in this table. & ``\,'' \\
125 & --parsing-strictness & Set strictness of parsing process.
Available options are: Default: ``normal'' & ``normal'' \\
126 & --partition-method & An integer value. Choose partitioning
strategy: 0=simple, 1=Zoltan, 2=METIS. & 1 \\
127 & --predetermined-time-steps-file & A file with a list of
predetermined time step sizes (one time step per line). & ``\,'' \\
128 & --pressure-max & A positive real value. Maximum absolute pressure.
& 1e+99 \\
129 & --pressure-min & A real value. Minimum absolute pressure. &
-1e+99 \\
130 & --pressure-scale & A positive real value. Scaling of pressure
primary variable. & 1 \\
131 & --pri-var-oscilation-threshold & A real positive value that
defines the threshold value for the primary variable switching
conditions after its meaning has switched to hinder oscillations. &
1e-05 \\
132 & --print-parameters & A positive integer value that request that
the run time parameters be printed at the start of the run: & 2 \\
133 & --print-properties & A positive integer value that request that
the compile time parameters be printed at the start of the run: & 2 \\
134 & --project-saturations & A boolean value that determines if all the
fluid saturations should be scaled to ensure the values are in the
interval (0, 1), including runs that use solvents (true), or not
(false). & false \\
135 & --regularization-factor-wells & A real positive value that defines
the ``regularization factor'' for wells. & 100 \\
136 & --relaxed-linear-solver-reduction & A real positive value that
defines the minimum reduction of the residual which the linear solver
need to achieve for the solution to be accepted. & 0.01 \\
137 & --relaxed-max-pv-fraction & A real positive value that defines The
fraction of the pore volume of the reservoir where the volumetric error
(CNV) might be violated during strict Newton iterations. & 0.03 \\
138 & --relaxed-pressure-tol-msw & A real positive value that sets the
relaxation tolerance for the multi-segment well pressure solution in
Pascals. & 10000 \\
139 & --relaxed-well-flow-tol & A real positive value that sets the
relaxation tolerance for the well flow residual in reservoir cubic
metres (rm3). & 0.001 \\
140 & --restart-time & A real positive value that sets the simulation
time at which a restart should be attempted {[}s{]}. & -1e+35 \\
141 & --restart-writing-interval & A positive integer value that sets
the frequency at which time steps are serialized to disk. & 16777215 \\
142 & --save-file & A character string that specifies the FileName for
.OPMRST file used for saving serialized state. If empty, CASENAME.OPMRST
is used. & ``\,'' \\
143 & --save-step & A character string that determines if the simulator
should save the serialized state of the OPM Flow simulator at one or
more report steps to a special *.OPMRST file. This is in addition to the
normal restart files written, and consumes significantly more space than
the normal restart files, but restarting OPM Flow from this file using
the --load-step option deviates less from the original run, compared to
restarting using the normal restart features. The files produced are not
compatible with other simulators, and also will not be compatible
between different releases of OPM Flow The parameter should be set to
one of the following: The default value of ``\,'' does not write
anything to the OPM Flow specific restart file. & ``\,'' \\
144 & --scale-linear-system & A boolean value that switches on (true) or
off (false) the scaling of linear system of equations according to the
equation scale and primary variable types. & false \\
145 & --sched-restart & A boolean value that determines for a restart
run if the case should initialize wells and groups from the historical
SCHEDULE section (true), or from the well and group data on the restart
file (false). Note that the commercial simulator always uses data from
the restart file (false). & false \\
146 & --serial-partitioning & A boolean value that determines if
partitioning for parallel runs on a single process (true), or not
(false). & false \\
147 & --shut-unsolvable-wells & A boolean value that determines if the
simulator should shut unsolvable wells (true), or not (false). & true \\
148 & --solve-welleq-initially & A boolean value that determines if the
simulator should fully solve the well equations before each iteration of
the reservoir model (true), or not (false). Note that the well equations
are always added to the full system and solved until convergence. &
true \\
149 & --solver-continue-on-convergence-failure & A boolean value that
stipulates if the simulator should continue (true) instead of stopping
(false) when the minimum solver time step is reached. & false \\
150 & --solver-growth-factor & A real positive value that specifies the
growth factor a time step can be increased by when recovering from one
or more time step chops, subject to the maximum allowable time step size
set by the --solver-max-time-step-in-days parameter. For example, if the
current time step has converged at 5 days after a time step chop, and
--solver-growth-factor is set to the default value of 2.0, then the next
time step will be 2.0 * 5 days, that is at 10 days. & 2.0 \\
151 & --solver-max-growth & A real positive value that specifies the
maximum growth factor a time step can be increased by after a report
time step, subject to the maximum allowable time step size set by the
--solver-max-time-step-in-days parameter. Thus, if the current time step
has converged at 5 days after at a report step, and --solver-max-growth
is set to the default value of 3.0, then the next time step will be 3.0
* 5 days, that is at 15 days. & 3.0 \\
152 & --solver-max-restarts & A positive integer that defines the number
of allowed consecutive restarts (or time step chops) before the
simulation is terminated. & 10 \\
153 & --solver-max-time-step-in-days & A real positive double precision
value that specifies the maximum allowed time step size in days. &
365 \\
154 & --solver-min-time-step & A real positive double precision value
that specifies the minimum size of a time step in days for field and
metric units, and hours for laboratory units If a time step cannot
converge without getting cut below this time step size the simulator
will stop. & 1e-12 \\
155 & --solver-restart-factor & A real positive double precision value
that sets the time step chop factor of the time step after a convergence
failure. For example, if the current non-convergent time step is 30 days
and --solver-restart-factor is set to the default value, then the time
step will be repeated using 0.33 * 30 days as the time step, that is 9.9
days. & 0.33 \\
156 & --solver-verbosity & A positive integer that specifies the
``chattiness'' of the non-linear solver. & 1 \\
157 & --strict-inner-iter-wells & A positive integer that specifies the
number of inner well iterations with strict tolerance. & 40 \\
158 & --strict-outer-iter-wells & A positive integer that specifies the
number of Newton iterations for which wells are checked with strict
tolerance. & 6 \\
159 & --temperature-max & A real positive value that sets the maximum
absolute temperature. & 1e+09 \\
160 & --temperature-min & A real positive value that sets the minimum
absolute temperature. & 0 \\
161 & --threads-per-process & A positive integer value that stipulates
the maximum number of threads to be instantiated per process (`-1' means
`automatic'). & -1 \\
162 & --time-step-after-event-in-days & A real double precision value
that sets the maximum allowed time step after an event; for example,
when a well is open or closed etc. The default value of -1 means that
events to do effect the time stepping. & -1 \\
163 & --time-step-control & A defined character string that defines the
time stepping control algorithm and is set to one of the following: &
``pid+newtoniteration'' \\
164 & --time-step-control-decay-damping-factor & A real positive value
that specifies the decay rate a time step can be decreased by after the
number target iterations has been exceeded & 1.0 \\
165 & --time-step-control-decay-rate & A real positive value that
specifies the decay rate a time step size can be decreased by after the
number of target iterations has been exceeded & 0.75 \\
166 & --time-step-control-file-name & A character string enclosed in
quotes that specifies a filename where time steps are specified. The
default is the character string timesteps. For instance time steps can
be generated by the ecl\_summary application in libecl as per the
following UNIX command line: path\_to\_libecl\_applications/ecl\_summary
DECK TIME \textgreater{} filename Where: DECK is the name of the data
deck you want to get the time steps from, TIME tells the application to
return the timing for the run, and ``filename'' is the name of the file
the times are piped to. & ``timesteps'' \\
167 & --time-step-control-growth-damping-factor & A real positive value
that specifies the growth rate of the time step increase when the number
of target iterations are undercut. & 3.2 \\
168 & --time-step-control-growth-rate & A real positive value that
specifies the growth rate of the time step size when the number of
target iterations are undercut. & 1.25 \\
169 & --time-step-control-target-iterations & A positive integer that
defines the number of linear iterations which the time step control
scheme should aim for (if applicable). & 30 \\
170 & --time-step-control-target-newton-iterations & A positive integer
that specifies the number of Newton iterations which the time step
control scheme should aim for (if applicable). & 8 \\
171 & --time-step-control-tolerance & A real double precision value that
sets the tolerance for PID (only used with the pid and pid+ options
defined by the --time-step-control option). & 0.1 \\
172 & --time-step-verbosity & A positive integer that specifies the
``chattiness'' during the time integration. & 1 \\
173 & --tolerance-cnv & A real positive double precision value that
specifies the maximum non-linear tolerance error. This is the local
convergence tolerance (maximum of local saturation errors). & 0.01 \\
174 & --tolerance-cnv-energy & A real value. Local energy convergence
tolerance (maximum of local energy errors). & 0.01 \\
175 & --tolerance-cnv-energy-relaxed & A real value. Relaxed local
energy convergence tolerance that applies for iterations after the
iterations with the strict tolerance. & 1.0 \\
176 & --tolerance-cnv-relaxed & A real positive value that defines the
relaxed local convergence tolerance that applies for iterations after
the iterations with the strict tolerance. & 1.0 \\
177 & --tolerance-energy-balance & A real value. Tolerated energy
balance error relative to (scaled) total energy present. & 1e-07 \\
178 & --tolerance-energy-balance-relaxed & A real value. Relaxed
tolerated energy balance error that applies for iterations after the
iterations with the strict tolerance. & 1e-06 \\
179 & --tolerance-mb & A real positive double precision value that sets
the maximum mass balance error, that is the tolerated mass balance error
relative to total mass present. & 1e-07 \\
180 & --tolerance-mb-relaxed & A positive real value. Relaxed tolerated
mass balance error that applies for iterations after the iterations with
the strict tolerance. & 1e-06 \\
180 & --tolerance-pressure-ms-wells & A real positive double precision
value that specifies the tolerance for the pressure equations for
multi-segment wells. & 1000 \\
181 & --tolerance-well-control & A real positive double precision value
that sets the maximum tolerance for the well control equations. &
1e-07 \\
182 & --tolerance-wells & A real positive double precision value that
defines the maximum non-linear error for the well equations. & 1e-04 \\
183 & -update-equations-scaling & A boolean value that switches on
(true) or off (false) the updating of the scaling factors for mass
balance equations during the simulation. & false \\
184 & --use-average-density-ms-wells & A boolean value that specifies
whether to approximate segment densitities by averaging over the segment
and its outlet (true) or not (false). & false \\
185 & --use-gmres & A boolean value that when set to true OPM Flow will
use Generalized Minimal Residual (``GMRES'') {[}Y. Saad, A flexible
inner-outer preconditioned GMRES algorithm, SIAM J. Sci. Statist.
Comput.,14, (1993).{]} and {[}Y. Saad and M.H. Schultz, GMRES: A
generalized minimum residual algorithm for solving nonsymmetric linear
systems, SIAM J. Sci. Statist. Comput., 7 (1986), pp.~856\{869.{]}
solver instead of Biconjugate Gradient Stabilized (``BiCGSTAB'') {[}Van
der Vorst, H. A. (1992). ``Bi-CGSTAB: A Fast and Smoothly Converging
Variant of Bi-CG for the Solution of Nonsymmetric Linear Systems''. SIAM
J. Sci. Stat. Comput. 13 (2): 631--644. doi:10.1137/0913035.
hdl:10338.dmlcz/104566{]} and {[}Sleijpen, G. L. G.; Fokkema, D. R.
(November 1993). ``BiCGstab(l) for linear equations involving
unsymmetric matrices with complex spectrum'' (PDF). Electronic
Transactions on Numerical Analysis. Kent, OH: Kent State University. 1:
11--32. ISSN 1068-9613.{]} as the linear solver within the Newton
iterations. & false \\
186 & --use-implicit-ipr & A boolean value. Compute implicit IPR for
stability checks and stable solution search. & true \\
187 & --use-multisegment-well & A boolean value that when set to true
the simulator will use the well model for multi-segment wells instead of
the one for single-segment wells. & true \\
188 & --use-update-stabilization & A boolean value that switches on
(true) or off (false) the stabilized Newton option, that attempts to
detect and correct oscillations or stagnation during the Newton
iterations. This option may improve convergence for some cases. &
true \\
189 & --water-only-threshold & A real positive value that defines the
saturation threshold, for which cells with water saturations above or
equal to this threshold are considered one-phase water only. & 1.0 \\
190 & --zoltan-imbalance-tol & A real value. Tolerable imbalance of the
loadbalancing provided by Zoltan. This command line option has been
deprecated; --imbalance-tol should be used instead. & 1.1 \\
191 & --zoltan-params & A character string that specifies the
configuration of the Zoltan partitioner. Valid options are: graph,
hypergraph or scotch. Alternatively, you can request a configuration to
be read from a JSON file by giving the filename here, ending with
`.json.' See https://sandialabs.github.io/Zoltan/ug\_html/ug.html for
available Zoltan options. & ``graph'' \\
Notes: flow --help instead use: flow --help-all to get a list of
supported command line Parameters. & & & \\
\end{longtable}

\emph{Table A.1: OPM Flow 2024-10 Command Line Options}

\begin{itemize}
\tightlist
\item
  The command line syntax for running OPM Flow is:
\end{itemize}

\begin{verbatim}

flow [OPTIONS] CASENAME
\end{verbatim}

and typing the following command from your terminal:

\begin{verbatim}

flow CASENAME
\end{verbatim}

will start the simulator and run the case specified by CASENAME.DATA.

OPM Flow accepts command line options to control various aspects of the
simulator's run time behavior, as described in the table below. To give
a command line option use ``key=value'' syntax, with no spaces around
the equals sign. It is also possible to put multiple options together in
a parameter file. To do so, put one option on each line of the file with
``.param'' as the extension and pass that filename as a command line
parameter to OPM Flow using:

\begin{verbatim}

flow --parameter-file=CASENAME.param
\end{verbatim}

to instruct OPM Flow to read the parameter file.

\begin{longtable}[]{@{}
  >{\raggedright\arraybackslash}p{(\linewidth - 6\tabcolsep) * \real{0.2500}}
  >{\raggedright\arraybackslash}p{(\linewidth - 6\tabcolsep) * \real{0.2500}}
  >{\raggedright\arraybackslash}p{(\linewidth - 6\tabcolsep) * \real{0.2500}}
  >{\raggedright\arraybackslash}p{(\linewidth - 6\tabcolsep) * \real{0.2500}}@{}}
\toprule\noalign{}
\begin{minipage}[b]{\linewidth}\raggedright
OPM Flow 2024-04 Command Line Options
\end{minipage} & \begin{minipage}[b]{\linewidth}\raggedright
\end{minipage} & \begin{minipage}[b]{\linewidth}\raggedright
\end{minipage} & \begin{minipage}[b]{\linewidth}\raggedright
\end{minipage} \\
\midrule\noalign{}
\endhead
\bottomrule\noalign{}
\endlastfoot
No. & Variable Name & Description & Default \\
1 & -h or --help & A character string that causes OPM Flow to print a
help message that gives a brief description of the available command
line parameters. This now only prints the active user facing command
line options. & N/A \\
2 & --help-all & Prints all the command line options included in the
release, including obsolete, hidden and deprecated options. & N/A \\
3 & --accelerator-mode & A defined character string that selects the
linear solver, usage:
`--accelerator-mode={[}none\textbar cusparse\textbar opencl\textbar amgcl\textbar rocalution\textbar rocsparse{]}'.
& ``none'' \\
4 & --allow-distributed-wells & Allow the perforations of a well to be
distributed to interior of multiple processes. & false \\
5 & --alternative-well-rate-init & Use alternative well rate
initialization procedure. & true \\
6 & --bda-device-id & Choose device ID for cusparseSolver or
openclSolver, use `nvidia-smi' or `clinfo' commands to determine valid
IDs. & 0 \\
7 & --continue-on-convergence-error & Continue with a non-converged
solution instead of giving up if we encounter a time step size smaller
than the minimum time step size. & false \\
8 & --cpr-reuse-interval & A positive integer that sets the reuse
preconditioner interval. Used when --cpr-reuse-setup is set to 4, then
the preconditioner will be fully recreated instead of reused every N
linear solve, where N is this parameter. & 30 \\
9 & --cpr-reuse-setup & A positive integer that defines if the CPR
solver should re-use the AMG setup. Valid options are: Changed the
default value from three to four. & 4 \\
10 & --dbph-max-rel & A real positive value that sets the maximum
relative change of the bottom-hole pressure in a single iteration. &
1.0 \\
11 & --dp-max-rel & A real positive double precision value that sets the
maximum allowed relative change of pressure per iteration. & 0.3 \\
12 & --ds-max & A real positive double precision value that sets the
maximum allowed change in saturation per iteration. & 0.2 \\
13 & --dwell-fraction-max & A real positive value that sets the maximum
allowed change in well's volume fraction per iteration. & 0.2 \\
14 & --debug-emit-cell-partition & A boolean value that determines
whether (true) or not (false) to emit cell partitions as a debugging
aid. & false \\
15 & --ecl-deck-file-name & A character string that defines the name of
the OPM Flow input file which contains the simulator's ECLIPSE formatted
input deck to be simulated. Note that there must be no spaces in the
filename. If there are spaces, which can typically occur on Windows
based systems, then the filename should be enclosed in quotes; thus
instead of FILE NAME use ``FILE NAME''. & ``\,'' \\
16 & --enable-drift-compensation & A boolean value that enables (true)
or disables (false) partial compensation of systematic mass losses via
the source term of the next time step. & true \\
17 & --ecl-enable-tuning & Honor some aspects of the TUNING keyword from
the ECL deck. & false \\
18 & --ecl-newton-relaxed-tolerance & The maximum error which the
volumetric residual may exhibit if it is in a `relaxed' region during a
strict iteration. & 1e+09 \\
19 & --ecl-newton-relaxed-volume-fraction & The fraction of the pore
volume of the reservoir where the volumetric error may be voilated
during strict Newton iterations. & 0.03 \\
20 & --ecl-newton-strict-iterations & The number of Newton iterations
where the volumetric error is considered. & 8 \\
21 & --ecl-newton-sum-tolerance & The maximum error tolerated by the
Newton method for considering a solution to be converged. & 0.0001 \\
22 & --ecl-newton-sum-tolerance-exponent & The exponent used to scale
the sum tolerance by the total pore volume of the reservoir. &
0.333333 \\
23 & --ecl-output-double-precision & A boolean value that switches on
(true) or off (false) double precision in restart files. Useful for
`perfect' restarts. & false \\
24 & --ecl-output-interval & An integer value defining the number of
report steps that ought to be skipped between two writes of restart
files results. & -1 \\
25 & --edge-weights-method & A defined positive integer value that
defines the edge-weighing strategy: & 1 \\
26 & --enable-adaptive-time-stepping & A boolean value that turns on
(true) or off (false) adaptive time stepping. If set to false the report
time steps provided by the input deck are used. & true \\
27 & --enable-async-ecl-output & A boolean value that sets the output
mode to be asynchronous (true), letting the simulator continue computing
the next time step while writing results to restart and summary files in
the commercial simulator's format on a separate thread. & true \\
28 & --enable-async-vtk-output & A boolean value that sets the output
mode to be asynchronous (true), letting the simulator continue computing
the next time step while writing results to the VTK output files on a
separate thread. & true \\
29 & --enable-dry-run & A boolean value that specifies if the simulation
should actually run (true), or just check the input deck (false). This
option is equivalent to activating the NOSIM keyword in the RUNSPEC
section of the input deck (see section Error: Reference source not found
Error: Reference source not found). The default value of ``auto' will
use whatever is stipulated in the input deck via the NOSIM keyword. &
``auto'' \\
30 & --enable-ecl-output & A boolean value that specifies if the binary
output files (restart, summary files, etc.) should be written in the
commercial simulator's format (true), or OPM Flow's format (false). &
true \\
31 & --enable-esmry & A boolean value that switches on (true) or off
(false) the output of SUMMARY vectors to the ESMRY file for fast loading
of summary data. & false \\
32 & --enable-gravity & A boolean value that switches on (true) or off
(false) the use of the gravity correction for the pressure gradients. &
true \\
33 & --enable-grid-adaptation & A boolean value that enables (true) or
disables (false) adaptive grid refinement/coarsening. Note that Eclipse
grid output will be disabled if adaptive grid refinement/coarsening is
enabled. & false \\
34 & --enable-intensive-quantity-cache & A boolean value that switches
on (true) or off (false) the caching of intensive quantities. & true \\
35 & --enable-logging-fallout-warning & A boolean value that sets the
developer option to see whether logging was on non-root processors
(true). If set to true output will be appended to the \emph{.DBG or
}.PRT files. & false \\
36 & --enable-opm-rst-file & A boolean value set to true or false to
write OPM specific data sets to the commercial simulators restart file
to enable restart of OPM Flow runs by OPM Flow (true), or not to write
the data (false). & false \\
37 & --enable-storage-cache & A boolean value that turns on (true) or
off (false) storing previous storage terms and avoid re-calculating
them. & true \\
38 & --enable-terminal-output & A boolean value that turns on (true) or
off (false) high-level information about the simulation's progress to
the terminal & true \\
39 & --enable-thermodynamic-hints & A boolean value that enables (true)
or disables (false) thermodynamic hints. & false \\
40 & --enable-tuning & A boolean value that instructs OPM Flow to read
the time stepping parameters from the Error: Reference source not found
keyword in the input deck, if set to true. Note that only the first
record of the TUNING keyword is processed. & false \\
41 & --enable-vtk-output & A boolean value that turns on (true) or off
(false) a global switch for writing VTK files. & false \\
42 & --enable-well-operability-check & A boolean value that enables
(true) checking of a well's operating status, or disables (false) the
checking. & true \\
43 & --enable-well-operability-check-iter & A boolean value that enables
(true) checking of a well's operating status during iterations, or
disables (false) the checking during iterations. & false \\
44 & --enable-write-all-solutions & A boolean value that turns on (true)
or off (false) the writing of all solutions to disk instead of only the
ones for the report steps. & false \\
45 & --end-time & The simulation time at which the simulation is
finished {[}s{]}. & 1e+100 \\
46 & --explicit-rock-compaction & A boolean value. Use pressure from end
of the last time step when evaluating rock compaction (true) or not
(false). & false \\
47 & --external-partition & A quoted character string. Name of file from
which to load an externally generated partitioning of the model's active
cells for MPI distribution purposes. If empty, the built-in partitioning
method will be employed. & ``\,'' \\
48 & --force-disable-fluid-in-place-output & A boolean value that
instructs OPM Flow not to print the Fluid In-Place report after each
report time step (true) or not (false). Note this parameter will
override the print request in the input deck. & false \\
49 & --force-disable-resv-fluid-in-place-output & A boolean value that
instructs OPM Flow to not to print the Reservoir Volume Fluid In-Place
report after each report time step (true) or not (false). Note this
parameter will override the print request in the input deck. & false \\
50 & --full-time-step-initially & A boolean value that instructs OPM
Flow to always attempt to finish a report step using a single time step
(true) or not (false). & false \\
51 & --ignore-keywords & A character string that defines a list of
keywords which should be ignored by OPM Flow. The keyword string should
be enclosed in quotes and each keyword separated by a colon, that is
`:'. & ``\,'' \\
52 & --ilu-fillin-level & A positive integer value that sets the fill in
level for the ILU preconditioner. & 0 \\
53 & --ilu-redblack & A boolean value that instructs OPM Flow to use
red-black partitioning for the ILU preconditioner. & false \\
54 & --ilu-relaxation & A real positive double precision value that sets
the relaxation factor of the linear solver's ILU preconditioner. &
0.9 \\
55 & --ilu-reorder-spheres & A boolean value that specifies OPM Flow to
reorder the entries of the matrix in the red-black ILU preconditioner in
spheres starting at an edge (true) or not (false). If false the original
ordering is preserved in each color. Otherwise try to ensure D4 ordering
(in a 2D structured grid, the diagonal elements are consecutive). &
false \\
56 & --initial-time-step-in-days & A real double precision value that
sets the size of initial time step in days. & 1.0 \\
57 & --initial-time-step-size & A real double precision value that sets
the size of initial time step in seconds. & 86400 \\
58 & --linear-solver & A defined quoted character string that sets the
configuration of the linear solver, valid values are: Option (3) extends
the existing Constrained Pressure Residual (``CPR'') preconditioner to
include wells. This option can also be invoked via the CPR keyword in
the RUNSPEC section; however, the command line parameter takes
precedence. For option (9) one enters a character string enclosed in
quotes that defines the filename of a JSON configuration file for the
flexible linear solver system. In this case the file extension should be
`.json'. Note that the *.PRT file contains the ``Property tree for the
linear solver'' listing, which is the JSON specification of the current
case, and can be used to configure a user specific linear solver JSON
file. The option ``cpr'' now is an alias for ``cprw'' instead of
``cpr\_trueimpes''. & ``ilu0'' \\
59 & --linear-solver-ignore-convergence-failure & A boolean value that
if set to true convergences failures in the linear solver are ignored.
This option should be used with care, as the results might be
unreliable. & false \\
60 & --linear-solver-max-iter & A positive integer value that defines
the maximum number of linear iterations. & 200 \\
61 & --linear-solver-print-json-definition & Write the JSON definition
of the linear solver setup to the DBG file. & true \\
62 & --linear-solver-reduction & A real positive double precision value
that sets the minimum reduction of the residual for the linear solver.
The linear solver convergences when the residual is reduced
sufficiently. The simulator now overrides the default reduction to be
0.005 instead of 0.01, if the linear solver has been set to one of the
cpr options, in a similar manner as how the default maximum number of
linear iterations for the cpr and cprw options are changed to 20 instead
of 200, unless specified by the command line option by the user. &
0.01 \\
63 & --linear-solver-restart & A positive integer value that sets the
number of iterations after which GMRES is restarted. & 40 \\
64 & --linear-solver-verbosity & A positive integer value that defines
the output from linear solver: & 0 \\
65 & --load-file & A character string enclosed in quotes that defines
the FileName for the .OPMRST file used to load the serialized state. If
empty, CASENAME.OPMRST is used. & ``\,'' \\
66 & --load-step & An integer value that determines if the simulator
should load the serialized state from OPM Flow's version of the restart
file, and should be set to either a specific report step, or 0 to load
the last stored report step. The default value of -1 does not load the
data from the OPM Flow specific restart file. OPM Flow's version of the
restart file, is written using the --save-step=N option. & -1 \\
67 & --local-domains-ordering-measure & A character string enclosed in
quotes that defines the subdomain ordering measure. Allowed values are
`maxpressure', `averagepressure' and `residual'. & ``maxpressure'' \\
68 & --local-domains-partitioning-imbalance & A real positive double
precision value that sets the subdomain partitioning imbalance
tolerance. 1.03 is 3 percent imbalance. & 1.03 \\
69 & --local-domains-partitioning-method & A character string enclosed
in quotes that defines the subdomain partitioning method. Allowed values
are `zoltan', `simple', and the name of a partition file ending with
`.partition'. & ``zoltan'' \\
70 & --local-solve-approach & A character string enclosed in quotes that
defines the local solve approach. Valid choices are `jacobi' and
`gauss-seidel'. & ``gauss-seidel'' \\
71 & --local-tolerance-scaling-cnv & A real positive double precision
value that sets the convergence tolerance for local solves. Set to lower
than 1.0 to use stricter convergence tolerance for local solves. &
0.1 \\
72 & --local-tolerance-scaling-mb & A real positive double precision
value that sets the convergence tolerance for local solves. Set to lower
than 1.0 to use stricter convergence tolerance for local solves. & 1 \\
73 & --local-well-solve-control-switching & A boolean value that allows
(true) or disallows (false) control switching during local well
solutions. & false \\
74 & --matrix-add-well-contributions & A boolean value that if set to
true explicitly specifies the influences of wells between cells in the
Jacobian and preconditioner matrices. & false \\
75 & --max-inner-iter-ms-wells & A positive integer value that defines
the maximum number of inner iterations for multi-segment wells. & 100 \\
76 & --max-inner-iter-wells & A positive integer value that defines the
maximum number of inner iterations for standard wells. & 50 \\
77 & --max-local-solve-iterations & A positive integer value that
defines the maximum number of iterations for local solves with NLDD
nonlinear solver. & 20 \\
78 & --max-newton-iterations-with-inner-well-iterations & A positive
integer that specifies the maximum Newton iterations with inner well
iterations. & 8 \\
79 & --max-pressure-change-ms-wells & A real positive value that defines
the maximum relative pressure change for a single iteration of the
multi-segment well model. & 1.0 x 106 \\
80 & --max-residual-allowed & A real positive value that sets the
absolute maximum tolerance for residuals without cutting the time step
size. & 1.0 x 107 \\
81 & --max-single-precision-days & A real positive value that sets the
maximum time step size where single precision floating point arithmetic
can be used for solving the linear systems of equations. & 20.0 \\
82 & --max-temperature-change & A real positive value that stipulates
the maximum absolute change of temperature in a single iteration. & 5 \\
83 & --max-time-step-divisions & A positive integer that specifies the
maximum number of divisions by two of the timestep size before the
simulation bails out. & 10 \\
84 & --max-time-step-size & A real positive value that defines the
maximum size to which all time steps are limited to in seconds. & inf \\
85 & --max-welleq-iter & A positive integer that defines the maximum
number of iterations to determine the solution to the well equations. &
30 \\
86 & --maximum-number-of-well-switches & A positive integer values that
stipulates the maximum number of times a well can switch to the same
control. & 3 \\
87 & --maximum-water-saturation & A real positive value that defines the
maximum water saturation. & 1.0 \\
88 & --milu-variant & A defined character string that specifies which
variant of the modified ILU preconditioner ought to be used. Possible
variants are: & ``ILU'' \\
89 & --min-strict-cnv-ite & A positive integer that sets the minimum
number of Newton iterations before relaxed tolerances are used for the
CNV convergence criterion. & 0 \\
90 & --min-strict-mb-iter & An integer value defining the number of
Newton iterations before relaxed tolerances can be used for the MB
convergence criterion. Default -1 means that the relaxed tolerance is
used when maximum number of Newton iterations are reached. & -1 \\
91 & --min-time-step-based-on-newton-iterations & A real positive value
that sets the minimum time step size (in days for field and metric unit
and hours for lab unit) can be reduced to based on Newton iteration
counts. & 0.0 \\
92 & --min-time-step-before-shutting-problematic-wells-in-days & A real
positive value that sets the minimum time step size in days for which
problematic wells are not shut. Time steps below this value will result
in problematic wells being shut-in by the simulator. & 0.01 \\
93 & --min-time-step-size & A real positive value that sets the minimum
size to which all time steps are limited to in seconds. & 0.0 \\
94 & --network-max-iterations & A positive integer that sets the maximum
number of iterations in the network solver before giving up. & 200 \\
95 & --network-max-strict-iterations & A positive integer that sets the
maximum iterations in network solver before relaxing tolerance. & 100 \\
96 & --newton-max-error & A real positive value that sets the maximum
error tolerated by the Newton method to which does not cause an abort. &
1e+100 \\
97 & --newton-max-iterations & A positive integer that defines the
maximum number of Newton iterations per time step used by the simulator.
& 20 \\
98 & --newton-max-relax & A real positive value that sets the maximum
relaxation factor of a Newton iteration used by the simulator. & 0.5 \\
99 & --newton-min-iterations & A positive integer that sets the minimum
number of Newton iterations per time step used by the simulator. & 2 \\
100 & --newton-relaxation-type & A character string. The type of
relaxation used by the Newton method. Valid options are: dampen or sor.
& ``dampen'' \\
101 & --newton-target-iterations & A positive integer that sets the
`optimum' number of Newton iterations per time step. & 10 \\
102 & --newton-tolerance & A real positive value that sets the maximum
raw error tolerated by the Newtonmethod for considering a solution to be
converged. & 0.01 \\
103 & --newton-verbose & A boolean value that specifies whether the
Newton method should inform the user about its progress (true) or not
(false). & true \\
104 & --newton-write-convergence & A boolean value that specifies
whether to write the convergence behaviour of the Newton method to a VTK
file (true) or not (false). & false \\
105 & --nldd-num-initial-newton-iter & A positive integer that sets the
number of initial global Newton iterations when running the NLDD
nonlinear solver. & 1 \\
106 & --nonlinear-solver & A character string that specifies the
nonlinear solver. Valid choices are newton or nldd. & ``newton'' \\
107 & --num-local-domains & A positive integer that sets the number of
local domains for NLDD nonlinear solver. Note this is an experimental
feature in the current release 2023.10 that is expected to be more
complete and tested by the 2024.10 release. & 0 \\
108 & --num-pressure-points-equil & A positive integer that sets the
number of pressure points (in each direction) in tables used for
equilibration. & 2000 \\
109 & --opencl-ilu-parallel & A boolean value that if set to true then
parallelize the ILU decomposition and application on GPU, or not
(false). & true \\
110 & --opencl-platform-id & A positive integer that specifies the
platform identification (``ID'') for the openSolver. Use the ``clinfo''
command to determine valid IDs. & 0 \\
111 & --output-dir & A character string that defines the directory to
which OPM Flow is to write the commercial simulator compatible output
files (restart, summary files etc.). The default value results in the
files be written to the same directory as the input file. & ``\,'' \\
112 & --output-extra-convergence-info & A defined character string that
specifies whether to provide additional convergence output to separate
files for diagnostic purposes. The available options are: Options can be
combined with commas, e.g.''steps,iterations'' for multiple outputs. The
default value of ``none'' prevents the two files from being written out,
for better compatibility with the commercial simulator. & ``none'' \\
113 & --output-interval & A positive integer that specifies the number
of report steps between two consecutive writes of restart data. & 1 \\
114 & --output-mode & A defined character string that defines the output
to the \emph{.PRT and }.DBG files: For example to just output logging
information use: --output-mode=``log''. or --output-mode=false &
``all'' \\
115 & --owner-cells-first & A boolean value that determines if the cells
owned by a rank should be ordered before ghost/overlapping the cells
(true), or not (false). & true \\
116 & --parameter-file & A character string that defines the name of a
parameter file which contains the simulator's set of run-time
parameters, as listed in this table. & ``\,'' \\
117 & --parsing-strictness & Set strictness of parsing process.
Available options are: Default: ``normal'' & ``normal'' \\
118 & --predetermined-time-steps-file & A file with a list of
predetermined time step sizes (one time step per line). & ``\,'' \\
119 & --pressure-max & A positive real value. Maximum absolute pressure.
& 1e+99 \\
120 & --pressure-min & A real value. Minimum absolute pressure. &
-1e+99 \\
121 & --pressure-scale & A positive real value. Scaling of pressure
primary variable. & 1 \\
122 & --pri-var-oscilation-threshold & A real positive value that
defines the threshold value for the primary variable switching
conditions after its meaning has switched to hinder oscillations. & 1.0
x 10-5 \\
123 & --print-parameters & A positive integer value that request that
the run time parameters be printed at the start of the run: & 2 \\
124 & --print-properties & A positive integer value that request that
the compile time parameters be printed at the start of the run: & 2 \\
125 & --project-saturations & A boolean value that determines if all the
fluid saturations should be scaled to ensure the values are in the
interval (0, 1), including runs that use solvents (true), or not
(false). & false \\
126 & --regularization-factor-wells & A real positive value that defines
the ``regularization factor'' for wells. & 100 \\
127 & --relaxed-linear-solver-reduction & A real positive value that
defines the minimum reduction of the residual which the linear solver
need to achieve for the solution to be accepted. & 0.01 \\
128 & --relaxed-max-pv-fraction & A real positive value that defines The
fraction of the pore volume of the reservoir where the volumetric error
(CNV) might be violated during strict Newton iterations. & 0.03 \\
129 & --relaxed-pressure-tol-msw & A real positive value that sets the
relaxation tolerance for the multi-segment well pressure solution in
Pascals. & 10000 \\
130 & --relaxed-well-flow-tol & A real positive value that sets the
relaxation tolerance for the well flow residual in reservoir cubic
metres (rm3). & 0.001 \\
131 & --restart-time & A real positive value that sets the simulation
time at which a restart should be attempted {[}s{]}. & -1e+35 \\
132 & --restart-writing-interval & A positive integer value that sets
the frequency at which time steps are serialized to disk. & 16777215 \\
133 & --save-file & A character string that specifies the FileName for
.OPMRST file used for saving serialized state. If empty, CASENAME.OPMRST
is used. & ``\,'' \\
134 & --save-step & A character string that determines if the simulator
should save the serialized state of the OPM Flow simulator at one or
more report steps to a special *.OPMRST file. This is in addition to the
normal restart files written, and consumes significantly more space than
the normal restart files, but restarting OPM Flow from this file using
the --load-step option deviates less from the original run, compared to
restarting using the normal restart features. The files produced are not
compatible with other simulators, and also will not be compatible
between different releases of OPM Flow The parameter should be set to
one of the following: The default value of ``\,'' does not write
anything to the OPM Flow specific restart file. & ``\,'' \\
135 & --scale-linear-system & A boolean value that switches on (true) or
off (false) the scaling of linear system of equations according to the
equation scale and primary variable types. & false \\
136 & --sched-restart & A boolean value that determines for a restart
run if the case should initialize wells and groups from the historical
SCHEDULE section (true), or from the well and group data on the restart
file (false). Note that the commercial simulator always uses data from
the restart file (false). & false \\
137 & --serial-partitioning & A boolean value that determines if
partitioning for parallel runs on a single process (true), or not
(false). & false \\
138 & --shut-unsolvable-wells & A boolean value that determines if the
simulator should shut unsolvable wells (true), or not (false). & true \\
139 & --solve-welleq-initially & A boolean value that determines if the
simulator should fully solve the well equations before each iteration of
the reservoir model (true), or not (false). Note that the well equations
are always added to the full system and solved until convergence. &
true \\
140 & --solver-continue-on-convergence-failure & A boolean value that
stipulates if the simulator should continue (true) instead of stopping
(false) when the minimum solver time step is reached. & false \\
141 & --solver-growth-factor & A real positive value that specifies the
growth factor a time step can be increased by when recovering from one
or more time step chops, subject to the maximum allowable time step size
set by the --solver-max-time-step-in-days parameter. For example, if the
current time step has converged at 5 days after a time step chop, and
--solver-growth-factor is set to the default value of 2.0, then the next
time step will be 2.0 * 5 days, that is at 10 days. & 2.0 \\
142 & --solver-max-growth & A real positive value that specifies the
maximum growth factor a time step can be increased by after a report
time step, subject to the maximum allowable time step size set by the
--solver-max-time-step-in-days parameter. Thus, if the current time step
has converged at 5 days after at a report step, and --solver-max-growth
is set to the default value of 3.0, then the next time step will be 3.0
* 5 days, that is at 15 days. & 3.0 \\
143 & --solver-max-restarts & A positive integer that defines the number
of allowed consecutive restarts (or time step chops) before the
simulation is terminated. & 10 \\
144 & --solver-max-time-step-in-days & A real positive double precision
value that specifies the maximum allowed time step size in days. &
365 \\
145 & --solver-min-time-step & A real positive double precision value
that specifies the minimum size of a time step in days for field and
metric units, and hours for laboratory units If a time step cannot
converge without getting cut below this time step size the simulator
will stop. & 1.0 x 10-12 \\
146 & --solver-restart-factor & A real positive double precision value
that sets the time step chop factor of the time step after a convergence
failure. For example, if the current non-convergent time step is 30 days
and --solver-restart-factor is set to the default value, then the time
step will be repeated using 0.33 * 30 days as the time step, that is 9.9
days. & 0.33 \\
147 & --solver-verbosity & A positive integer that specifies the
``chattiness'' of the non-linear solver. & 1 \\
148 & --strict-inner-iter-wells & A positive integer that specifies the
number of inner well iterations with strict tolerance. & 40 \\
149 & --strict-outer-iter-wells & A positive integer that specifies the
number of Newton iterations for which wells are checked with strict
tolerance. & 6 \\
150 & --temperature-max & A real positive value that sets the maximum
absolute temperature. & 1.0 x 109 \\
151 & --temperature-min & A real positive value that sets the minimum
absolute temperature. & 0 \\
152 & --threads-per-process & A positive integer value that stipulates
the maximum number of threads to be instantiated per process (`-1' means
`automatic'). & -1 \\
153 & --tolerance-mb-relaxed & A positive real value. Relaxed tolerated
mass balance error that applies for iterations after the iterations with
the strict tolerance. & 1e-06 \\
154 & --time-step-after-event-in-days & A real double precision value
that sets the maximum allowed time step after an event; for example,
when a well is open or closed etc. The default value of -1 means that
events to do effect the time stepping. & -1 \\
155 & --time-step-control & A defined character string that defines the
time stepping control algorithm and is set to one of the following: &
``pid+newtoniteration'' \\
156 & --time-step-control-decay-damping-factor & A real positive value
that specifies the decay rate a time step can be decreased by after the
number target iterations has been exceeded & 1.0 \\
157 & --time-step-control-decay-rate & A real positive value that
specifies the decay rate a time step size can be decreased by after the
number of target iterations has been exceeded & 0.75 \\
158 & --time-step-control-file-name & A character string enclosed in
quotes that specifies a filename where time steps are specified. The
default is the character string timesteps. For instance time steps can
be generated by the ecl\_summary application in libecl as per the
following UNIX command line: path\_to\_libecl\_applications/ecl\_summary
DECK TIME \textgreater{} filename Where: DECK is the name of the data
deck you want to get the time steps from, TIME tells the application to
return the timing for the run, and ``filename'' is the name of the file
the times are piped to. & ``timesteps'' \\
159 & --time-step-control-growth-damping-factor & A real positive value
that specifies the growth rate of the time step increase when the number
of target iterations are undercut. & 3.2 \\
160 & --time-step-control-growth-rate & A real positive value that
specifies the growth rate of the time step size when the number of
target iterations are undercut. & 1.25 \\
161 & --time-step-control-target-iterations & A positive integer that
defines the number of linear iterations which the time step control
scheme should aim for (if applicable). & 30 \\
162 & --time-step-control-target-newton-iterations & A positive integer
that specifies the number of Newton iterations which the time step
control scheme should aim for (if applicable). & 8 \\
163 & --time-step-control-tolerance & A real double precision value that
sets the tolerance for PID (only used with the pid and pid+ options
defined by the --time-step-control option). & 0.1 \\
164 & --time-step-verbosity & A positive integer that specifies the
``chattiness'' during the time integration. & 1 \\
165 & --tolerance-cnv & A real positive double precision value that
specifies the maximum non-linear tolerance error. This is the local
convergence tolerance (maximum of local saturation errors). & 0.01 \\
166 & --tolerance-cnv-relaxed & A real positive value that defines the
relaxed local convergence tolerance that applies for iterations after
the iterations with the strict tolerance. & 1.0 \\
167 & --tolerance-mb & A real positive double precision value that sets
the maximum mass balance error, that is the tolerated mass balance error
relative to total mass present. & 1.0 x 10-6 \\
168 & --tolerance-pressure-ms-wells & A real positive double precision
value that specifies the tolerance for the pressure equations for
multi-segment wells. & 1000 \\
169 & --tolerance-well-control & A real positive double precision value
that sets the maximum tolerance for the well control equations. & 1.0 x
10-7 \\
170 & --tolerance-wells & A real positive double precision value that
defines the maximum non-linear error for the well equations. & 1.0 x
10-4 \\
171 & -update-equations-scaling & A boolean value that switches on
(true) or off (false) the updating of the scaling factors for mass
balance equations during the simulation. & false \\
172 & --use-average-density-ms-wells & A boolean value that specifies
whether to approximate segment densitities by averaging over the segment
and its outlet (true) or not (false). & false \\
173 & --use-gmres & A boolean value that when set to true OPM Flow will
use Generalized Minimal Residual (``GMRES'') {[}Y. Saad, A flexible
inner-outer preconditioned GMRES algorithm, SIAM J. Sci. Statist.
Comput.,14, (1993).{]} and {[}Y. Saad and M.H. Schultz, GMRES: A
generalized minimum residual algorithm for solving nonsymmetric linear
systems, SIAM J. Sci. Statist. Comput., 7 (1986), pp.~856\{869.{]}
solver instead of Biconjugate Gradient Stabilized (``BiCGSTAB'') {[}Van
der Vorst, H. A. (1992). ``Bi-CGSTAB: A Fast and Smoothly Converging
Variant of Bi-CG for the Solution of Nonsymmetric Linear Systems''. SIAM
J. Sci. Stat. Comput. 13 (2): 631--644. doi:10.1137/0913035.
hdl:10338.dmlcz/104566{]} and {[}Sleijpen, G. L. G.; Fokkema, D. R.
(November 1993). ``BiCGstab(l) for linear equations involving
unsymmetric matrices with complex spectrum'' (PDF). Electronic
Transactions on Numerical Analysis. Kent, OH: Kent State University. 1:
11--32. ISSN 1068-9613.{]} as the linear solver within the Newton
iterations. & false \\
174 & --use-implicit-ipr & A boolean value. Compute implicit IPR for
stability checks and stable solution search. & false \\
175 & --use-multisegment-well & A boolean value that when set to true
the simulator will use the well model for multi-segment wells instead of
the one for single-segment wells. & true \\
176 & --use-update-stabilization & A boolean value that switches on
(true) or off (false) the stabilized Newton option, that attempts to
detect and correct oscillations or stagnation during the Newton
iterations. This option may improve convergence for some cases. &
true \\
177 & --water-only-threshold & A real positive value that defines the
saturation threshold, for which cells with water saturations above or
equal to this threshold are considered one-phase water only. & 1.0 \\
178 & --zoltan-imbalance-tol & A real positive value that defines the
tolerable imbalance of the load balancing provided by Zoltan package. &
1.1 \\
179 & --zoltan-params & A character string that specifies the
configuration of the Zoltan partitioner. Valid options are: graph,
hypergraph or scotch. Alternatively, you can request a configuration to
be read from a JSON file by giving the filename here, ending with
`.json.' See https://sandialabs.github.io/Zoltan/ug\_html/ug.html for
available Zoltan options. & ``graph'' \\
Notes: flow --help instead use: flow --help-all to get a list of
supported command line Parameters. & & & \\
\end{longtable}

\emph{Table A.3: OPM Flow 2024-04 Command Line Options}

\begin{itemize}
\tightlist
\item
  The command line syntax for running OPM Flow is:
\end{itemize}

\begin{verbatim}

flow [OPTIONS] CASENAME
\end{verbatim}

and typing the following command from your terminal:

\begin{verbatim}

flow CASENAME
\end{verbatim}

will start the simulator and run the case specified by CASENAME.DATA.

OPM Flow accepts command line options to control various aspects of the
simulator's run time behavior, as described in the table below. To give
a command line option use ``key=value'' syntax, with no spaces around
the equals sign. It is also possible to put multiple options together in
a parameter file. To do so, put one option on each line of the file with
``.param'' as the extension and pass that filename as a command line
parameter to OPM Flow using:

\begin{verbatim}

flow --parameter-file=CASENAME.param
\end{verbatim}

to instruct OPM Flow to read the parameter file.

\begin{longtable}[]{@{}
  >{\raggedright\arraybackslash}p{(\linewidth - 6\tabcolsep) * \real{0.2500}}
  >{\raggedright\arraybackslash}p{(\linewidth - 6\tabcolsep) * \real{0.2500}}
  >{\raggedright\arraybackslash}p{(\linewidth - 6\tabcolsep) * \real{0.2500}}
  >{\raggedright\arraybackslash}p{(\linewidth - 6\tabcolsep) * \real{0.2500}}@{}}
\toprule\noalign{}
\begin{minipage}[b]{\linewidth}\raggedright
OPM Flow 2023-10 Command Line Options
\end{minipage} & \begin{minipage}[b]{\linewidth}\raggedright
\end{minipage} & \begin{minipage}[b]{\linewidth}\raggedright
\end{minipage} & \begin{minipage}[b]{\linewidth}\raggedright
\end{minipage} \\
\midrule\noalign{}
\endhead
\bottomrule\noalign{}
\endlastfoot
No. & Variable Name & Description & Default \\
1 & -h or --help & A character string that causes OPM Flow to print a
help message that gives a brief description of the available command
line parameters. This now only prints the active user facing command
line options. & N/A \\
2 & --help-all & Prints all the command line options included in the
release, including obsolete, hidden and deprecated options. & N/A \\
3 & --accelerator-mode & Choose a linear solver, usage:
`--accelerator-mode={[} none \textbar{} cusparse \textbar{} opencl
\textbar{} amgcl \textbar{} rocalution {]}'. & ``none'' \\
4 & --allow-distributed-wells & Allow the perforations of a well to be
distributed to interior of multiple processes. & false \\
5 & --alternative-well-rate-init & Use alternative well rate
initialization procedure. & true \\
6 & --bda-device-id & Choose device ID for cusparseSolver or
openclSolver, use `nvidia-smi' or `clinfo' commands to determine valid
IDs. & 0 \\
7 & --continue-on-convergence-error & Continue with a non-converged
solution instead of giving up if we encounter a time step size smaller
than the minimum time step size. & false \\
8 & --cpr-reuse-interval & A positive integer that sets the reuse
preconditioner interval. Used when --cpr-reuse-setup is set to 4, then
the preconditioner will be fully recreated instead of reused every N
linear solve, where N is this parameter. & 30 \\
9 & --cpr-reuse-setup & A positive integer that defines if the CPR
solver should re-use the AMG setup. Valid options are: Changed the
default value from three to four. & 4 \\
10 & --dbph-max-rel & A real positive value that sets the maximum
relative change of the bottom-hole pressure in a single iteration. &
1.0 \\
11 & --dp-max-rel & A real positive double precision value that sets the
maximum allowed relative change of pressure per iteration. & 0.3 \\
12 & --ds-max & A real positive double precision value that sets the
maximum allowed change in saturation per iteration. & 0.2 \\
13 & --dwell-fraction-max & A real positive value that sets the maximum
allowed change in well's volume fraction per iteration. & 0.2 \\
14 & --ecl-deck-file-name & A character string that defines the name of
the OPM Flow input file which contains the simulator's ECLIPSE formatted
input deck to be simulated. Note that there must be no spaces in the
filename. If there are spaces, which can typically occur on Windows
based systems, then the filename should be enclosed in quotes; thus
instead of FILE NAME use ``FILE NAME''. & ``\,'' \\
15 & --ecl-enable-drift-compensation & A boolean value that enables
(true) or disables (false) partial compensation of systematic mass
losses via the source term of the next time step. & true \\
16 & --ecl-enable-tuning & Honor some aspects of the TUNING keyword from
the ECL deck. & false \\
17 & --ecl-newton-relaxed-tolerance & The maximum error which the
volumetric residual may exhibit if it is in a `relaxed' region during a
strict iteration. & 1e+09 \\
18 & --ecl-newton-relaxed-volume-fraction & The fraction of the pore
volume of the reservoir where the volumetric error may be voilated
during strict Newton iterations. & 0.03 \\
19 & --ecl-newton-strict-iterations & The number of Newton iterations
where the volumetric error is considered. & 8 \\
20 & --ecl-newton-sum-tolerance & The maximum error tolerated by the
Newton method for considering a solution to be converged. & 0.0001 \\
21 & --ecl-newton-sum-tolerance-exponent & The exponent used to scale
the sum tolerance by the total pore volume of the reservoir. &
0.333333 \\
22 & --ecl-output-double-precision & A boolean value that switches on
(true) or off (false) double precision in restart files. Useful for
`perfect' restarts. & false \\
23 & --ecl-output-interval & An integer value defining the number of
report steps that ought to be skipped between two writes of restart
files results. & -1 \\
24 & --edge-weights-method & A defined positive integer value that
defines the edge-weighing strategy: & 1 \\
25 & --enable-adaptive-time-stepping & A boolean value that turns on
(true) or off (false) adaptive time stepping. If set to false the report
time steps provided by the input deck are used. & true \\
26 & --enable-async-ecl-output & A boolean value that sets the output
mode to be asynchronously (true), letting the simulator continue
computing the next time step while writing results to restart and
summary files in the commercial simulators format. & true \\
27 & --enable-async-vtk-output & A boolean value that sets the output
mode to be asynchronously (true), letting the simulator continue
computing the next time step while writing results to the VTK output
files. & true \\
28 & --enable-dry-run & A boolean value that specifies if the simulation
should actually run, (true) or just check the input deck (false). This
option is equivalent to activating the NOSIM keyword in the RUNSPEC
section of the input deck (see section Error: Reference source not found
Error: Reference source not found). The default value of ``auto' will
use whatever is stipulated in the input deck via the NOSIM keyword. &
``auto'' \\
29 & --enable-ecl-output & A boolean value that specifies if the binary
output files (restart, summary files, etc.) should be written in the
commercial simulator's format (true), or OPM Flow's format (false). &
true \\
30 & --enable-esmry & A boolean value that switches on (true) or off
(false) the output SUMMARY vectors to the ESMRY file for fast loading of
summary data. & false \\
31 & --enable-gravity & A boolean value that switches on (true) or off
(false) the use of the gravity correction for the pressure gradients. &
true \\
32 & --enable-grid-adaptation & A boolean value that enables (true) or
disables (false) adaptive grid refinement/coarsening. Note that Eclipse
grid output will be disabled if adaptive grid refinement/coarsening is
enabled. & false \\
33 & --enable-intensive-quantity-cache & A boolean value that switches
on (true) or off (false) the caching of intensive quantities. & true \\
34 & --enable-logging-fallout-warning & A boolean value that sets the
developer option to see whether logging was on non-root processors
(true). If set to true output will be appended to the \emph{.DBG or
}.PRT files. & false \\
35 & --enable-opm-rst-file & A boolean value set to true or false to
write OPM specific data sets to the commercial simulators restart file
to enable restart of an OPM Flow runs (true) by OPM Flow, or not to
write the data (false). & false \\
36 & --enable-storage-cache & A boolean value that turns on (true) or
off (false) storing previous storage terms and avoid re-calculating
them. & true \\
37 & --enable-terminal-output & A boolean value that turns on (true) or
off (false) high-level information about the simulation's progress to
the terminal & true \\
38 & --enable-thermodynamic-hints & A boolean value that enables (true)
or disables (false) thermodynamic hints. & false \\
39 & --enable-tuning & A boolean value that instructs OPM Flow to read
the time stepping parameters from the Error: Reference source not found
keyword in the input deck, if set to true. Note that only the first
record of the TUNING keyword is processed. & false \\
40 & --enable-vtk-output & A boolean value that turns on (true) or off
(false) a global switch for writing VTK files. & false \\
41 & --enable-well-operability-check & A boolean value that enables
(true) checking of a well's operating status, or disables (false) the
checking. & true \\
42 & --enable-well-operability-check-iter & A boolean value that enables
(true) checking of a well's operating status during iterations, or
disables (false) the checking during iterations. & false \\
43 & --enable-write-all-solutions & A boolean value that turns on (true)
or off (false) the writing of all solutions to disk instead of only the
ones for the report steps. & false \\
44 & --end-time & The simulation time at which the simulation is
finished {[}s{]}. & 1e+100 \\
45 & --force-disable-fluid-in-place-output & A boolean value that
instructs OPM Flow not to print the Fluid In-Place report after each
report time step (true) or not (false). Note this parameter will
override the print request in the input deck. & false \\
46 & --force-disable-resv-fluid-in-place-output & A boolean value that
instructs OPM Flow to not to print the Reservoir Volume Fluid In-Place
report after each report time step (true) or not (false). Note this
parameter will override the print request in the input deck. & false \\
47 & --full-time-step-initially & A boolean value that instructs OPM
Flow to always attempt to finish a report step using a single time step
(true) or not (false). & false \\
48 & --ignore-keywords & A character string that defines a list of
keywords which should be ignored by OPM Flow. The keyword string should
be enclosed in quotes and each keyword separated by a colon, that is
`:'. & ``\,'' \\
49 & --ilu-fillin-level & A positive integer value that sets the fill in
level for the ILU preconditioner. & 0 \\
50 & --ilu-redblack & A boolean value that instructs OPM Flow to use
red-black partitioning for the ILU preconditioner. & false \\
51 & --ilu-relaxation & A real positive double precision value that sets
the relaxation factor of the linear solver's ILU preconditioner. &
0.9 \\
52 & --ilu-reorder-spheres & A boolean value that specifies OPM Flow to
reorder the entries of the matrix in the red-black ILU preconditioner in
spheres starting at an edge (true) or not (false). If false the original
ordering is preserved in each color. Otherwise try to ensure D4 ordering
(in a 2D structured grid, the diagonal elements are consecutive). &
false \\
53 & --initial-time-step-in-days & A real double precision value that
sets the size of initial time step in days. & 1.0 \\
54 & --initial-time-step-size & A real double precision value that sets
the size of initial time step in seconds. & 86400 \\
55 & --linear-solver & A defined quoted character string that sets the
configuration of the linear solver, valid values are: Option (2) extends
the existing Constrained Pressure Residual (``CPR'') preconditioner to
include wells. This option can also be invoked via the CPR keyword in
the RUNSPEC section; however, the command line parameter takes
precedence. For option (7) one enters a character string enclosed in
quotes that defines the filename of a JSON configuration file for the
flexible linear solver system. In this case the file extension should be
`.json'. Note that the *.PRT file contains the ``Property tree for the
linear solver'' listing, which is the JSON specification of the current
case, and can be used to configure a user specific linear solver JSON
file. The option ``cpr'' now is an alias for ``cprw'' instead of
``cpr\_trueimpes''. & ``ilu0'' \\
56 & --linear-solver-ignore-convergence-failure & A boolean value that
if set to true convergences failures in the linear solver are ignored.
This option should be used with care, as the results might be
unreliable. & false \\
57 & --linear-solver-max-iter & A positive integer value that defines
the maximum number of linear iterations. & 200 \\
58 & --linear-solver-print-json-definition & Write the JSON definition
of the linear solver setup to the DBG file. & true \\
59 & --linear-solver-reduction & A real positive double precision value
that sets the minimum reduction of the residual for the linear solver.
The linear solver convergences when the residual is reduced
sufficiently. The simulator now overrides the default reduction to be
0.005 instead of 0.01, if the linear solver has been set to one of the
cpr options, in a similar manner as how the default maximum number of
linear iterations for the cpr and cprw options are changed to 20 instead
of 200, unless specified by the command line option by the user. &
0.01 \\
60 & --linear-solver-restart & A positive integer value that sets the
number of iterations after which GMRES is restarted. & 40 \\
61 & --linear-solver-verbosity & A positive integer value that defines
the output from linear solver: & 0 \\
62 & --load-file & A character string enclosed in quotes that defines
the FileName for the .OPMRST file used to load the serialized state. If
empty, CASENAME.OPMRST is used. & ``\,'' \\
63 & --load-step & An integer value that determines if the simulator
should load the serialized state from OPM Flow's version of the restart
file, and should be set to either a specific report step, or 0 to load
the last stored report step. The default value of -1 does not load the
data from the OPM Flow specific restart file. OPM Flow's version of the
restart file, is written using the --save-step=N option. & -1 \\
64 & --local-domains-ordering-measure & A character string enclosed in
quotes that defines the subdomain ordering measure. Allowed values are
`pressure' and `residual'. & ``pressure'' \\
65 & --local-domains-partitioning-imbalance & A real positive double
precision value that sets the subdomain partitioning imbalance
tolerance. 1.03 is 3 percent imbalance. & 1.03 \\
66 & --local-domains-partitioning-method & A character string enclosed
in quotes that defines the subdomain partitioning method. Allowed values
are `zoltan', `simple', and the name of a partition file ending with
`.partition'. & ``zoltan'' \\
67 & --local-solve-approach & A character string enclosed in quotes that
defines the local solve approach. Valid choices are `jacobi' and
`gauss-seidel'. & ``jacobi'' \\
68 & --local-tolerance-scaling-cnv & A real positive double precision
value that sets the convergence tolerance for local solves. Set to lower
than 1.0 to use stricter convergence tolerance for local solves. &
0.01 \\
69 & --local-tolerance-scaling-mb & A real positive double precision
value that sets the convergence tolerance for local solves. Set to lower
than 1.0 to use stricter convergence tolerance for local solves. & 1 \\
70 & --local-well-solve-control-switching & A boolean value that allows
(true) or disallows (false) control switching during local well
solutions. & false \\
71 & --matrix-add-well-contributions & A boolean value that if set to
true explicitly specifies the influences of wells between cells in the
Jacobian and preconditioner matrices. & false \\
72 & --max-inner-iter-ms-wells & A positive integer value that defines
the maximum number of inner iterations for multi-segment wells. & 100 \\
73 & --max-inner-iter-wells & A positive integer value that defines the
maximum number of inner iterations for standard wells. & 50 \\
74 & --max-local-solve-iterations & A positive integer value that
defines the maximum number of iterations for local solves with NLDD
nonlinear solver. & 20 \\
75 & --max-newton-iterations-with-inner-well-iterations & A positive
integer that specifies the maximum Newton iterations with inner well
iterations. & 8 \\
76 & --max-pressure-change-ms-wells & A real positive value that defines
the maximum relative pressure change for a single iteration of the
multi-segment well model. & 1.0 x 106 \\
77 & --max-residual-allowed & A real positive value that sets the
absolute maximum tolerance for residuals without cutting the time step
size. & 1.0 x 107 \\
78 & --max-single-precision-days & A real positive value that sets the
maximum time step size where single precision floating point arithmetic
can be used for solving the linear systems of equations. & 20.0 \\
79 & --max-temperature-change & A real positive value that stipulates
the maximum absolute change of temperature in a single iteration. & 5 \\
80 & --max-time-step-divisions & A positive integer that specifies the
maximum number of divisions by two of the timestep size before the
simulation bails out. & 10 \\
81 & --max-time-step-size & A real positive value that defines the
maximum size to which all time steps are limited to in seconds. & inf \\
82 & --max-welleq-iter & A positive integer that defines the maximum
number of iterations to determine the solution to the well equations. &
30 \\
83 & --maximum-number-of-well-switches & A positive integer values that
stipulates the maximum number of times a well can switch to the same
control. & 3 \\
84 & --maximum-water-saturation & A real positive value that defines the
maximum water saturation. & 1.0 \\
85 & --milu-variant & A defined character string that specifies which
variant of the modified ILU preconditioner ought to be used. Possible
variants are: The default is ``ILU'' & ``ILU'' \\
86 & --min-strict-cnv-ite & A positive integer that sets the minimum
number of Newton iterations before relaxed tolerances are used for the
CNV convergence criterion. & 0 \\
87 & --min-time-step-based-on-newton-iterations & A real positive value
that sets the minimum time step size (in days for field and metric unit
and hours for lab unit) can be reduced to based on Newton iteration
counts. & 0.0 \\
88 & --min-time-step-before-shutting-problematic-wells-in-days & A real
positive value that sets the minimum time step size in days for which
problematic wells are not shut. Time steps below this value will result
in problematic wells being shut-in by the simulator. & 0.01 \\
89 & --min-time-step-size & A real positive value that sets the minimum
size to which all time steps are limited to in seconds. & 0.0 \\
90 & --network-max-iterations & A positive integer that sets the maximum
number of iterations in the network solver before giving up. & 200 \\
91 & --network-max-strict-iterations & A positive integer that sets the
maximum iterations in network solver before relaxing tolerance. & 100 \\
92 & --newton-max-error & A real positive value that sets the maximum
error tolerated by the Newton method to which does not cause an abort. &
1e+10 \\
93 & --newton-max-iterations & A positive integer that defines the
maximum number of Newton iterations per time step used by the simulator.
& 20 \\
94 & --newton-max-relax & A real positive value that sets the maximum
relaxation factor of a Newton iteration used by the simulator. & 0.5 \\
95 & --newton-min-iterations & A positive integer that sets the minimum
number of Newton iterations per time step used by the simulator. The
default value of one ensures that at least one Newton iteration is
performed after the previous time step. & 1 \\
96 & --newton-relaxation-type & A character string. The type of
relaxation used by the Newton method. Valid options are: dampen or sor.
& ``dampen'' \\
97 & --newton-target-iterations & A positive integer that sets the
`optimum' number of Newton iterations per time step. & 6 \\
98 & --newton-tolerance & A real positive value that sets the maximum
raw error tolerated by the Newtonmethod for considering a solution to be
converged. & 0.01 \\
99 & --newton-verbose & A boolean value that specifies whether the
Newton method should inform the user about its progress (true) or not
(false). & true \\
100 & --newton-write-convergence & A boolean value that specifies
whether to write the convergence behaviour of the Newton method to a VTK
file (true) or not (false). & false \\
101 & --nldd-num-initial-newton-iter & A positive integer that sets the
number of initial global Newton iterations when running the NLDD
nonlinear solver. & 1 \\
102 & --nonlinear-solver & A character string that specifies the
nonlinear solver. Valid choices are newton or nldd. & ``newton'' \\
103 & --num-local-domains & A positive integer that sets the number of
local domains for NLDD nonlinear solver. Note this is an experimental
feature in the current release 2023.10 that is expected to be more
complete and tested by the 2024.10 release. & 0 \\
104 & --num-pressure-points-equil & A positive integer that sets the
number of pressure points (in each direction) in tables used for
equilibration. & 2000 \\
105 & --opencl-ilu-parallel & A boolean value that if set to true then
parallelize the ILU decomposition and application on GPU, or not
(false). & true \\
106 & --opencl-platform-id & A positive integer that specifies the
platform identification (``ID'') for the openSolver. Use the ``clinfo''
command to determine valid IDs. & 0 \\
107 & --output-dir & A character string that defines the directory to
which OPM Flow is to write the commercial simulator compatible output
files (restart, summary files etc.). The default value results in the
files be written to the same directory as the input file. & ``\,'' \\
108 & --output-extra-convergence-info & A defined character string that
specifies whether to provide additional convergence output to separate
files for diagnostic purposes. The available options are: Options can be
combined with commas, e.g.''steps,iterations'' for multiple outputs. The
default value of ``none'' prevents the two files from being written out,
for better compatibility with the commercial simulator. & ``none'' \\
109 & --output-interval & A positive integer that specifies the number
of report steps between two consecutive writes of restart data. & 1 \\
110 & --output-mode & A defined character string that defines the output
to the \emph{.PRT and }.DBG files: For example to just output logging
information use: --output-mode=``log''. or --output-mode=false &
``all'' \\
111 & --owner-cells-first & A boolean value that determines if the cells
owned by a rank should be ordered before ghost/overlapping the cells
(true), or not (false). & true \\
112 & --parameter-file & A character string that defines the name of a
parameter file which contains the simulator's set of run-time
parameters, as listed in this table. & ``\,'' \\
113 & --parsing-strictness & Set strictness of parsing process.
Available options are: Default: ``normal'' & ``normal'' \\
114 & --predetermined-time-steps-file & A file with a list of
predetermined time step sizes (one time step per line). & ``\,'' \\
115 & --pri-var-oscilation-threshold & A real positive value that
defines the threshold value for the primary variable switching
conditions after its meaning has switched to hinder oscillations. & 1.0
x 10-5 \\
116 & --print-parameters & A positive integer value that request that
the run time parameters be printed at the start of the run: & 2 \\
117 & --print-properties & A positive integer value that request that
the compile time parameters be printed at the start of the run: & 2 \\
118 & --project-saturations & A boolean value that determines if all the
fluid saturations should be scaled to ensure the values are in the
interval (0, 1), including runs that use solvents (true), or not
(false). & false \\
119 & --regularization-factor-wells & A real positive value that defines
the ``regularization factor'' for wells. & 100 \\
120 & --relaxed-linear-solver-reduction & A real positive value that
defines the minimum reduction of the residual which the linear solver
need to achieve for the solution to be accepted. & 0.01 \\
121 & --relaxed-max-pv-fraction & A real positive value that defines The
fraction of the pore volume of the reservoir where the volumetric error
(CNV) might be violated during strict Newton iterations. & 0.03 \\
122 & --relaxed-pressure-tol-msw & A real positive value that sets the
relaxation tolerance for the multi-segment well pressure solution in
Pascals. & 10000 \\
123 & --relaxed-well-flow-tol & A real positive value that sets the
relaxation tolerance for the well flow residual in reservoir cubic
metres (rm3). & 0.001 \\
124 & --restart-time & A real positive value that sets the simulation
time at which a restart should be attempted {[}s{]}. & -1e+35 \\
125 & --restart-writing-interval & A positive integer value that sets
the frequency at which time steps are serialized to disk. & 16777215 \\
126 & --save-file & A character string that specifies the FileName for
.OPMRST file used for saving serialized state. If empty, CASENAME.OPMRST
is used. & ``\,'' \\
127 & --save-step & A character string that determines if the simulator
should save the serialized state of the OPM Flow simulator at one or
more report steps to a special *.OPMRST file. This is in addition to the
normal restart files written, and consumes significantly more space than
the normal restart files, but restarting OPM Flow from this file using
the --load-step option deviates less from the original run, compared to
restarting using the normal restart features. The files produced are not
compatible with other simulators, and also will not be compatible
between different releases of OPM Flow The parameter should be set to
one of the following: The default value of ``\,'' does not write
anything to the OPM Flow specific restart file. & ``\,'' \\
128 & --scale-linear-system & A boolean value that switches on (true) or
off (false) the scaling of linear system of equations according to the
equation scale and primary variable types. & false \\
129 & --sched-restart & A boolean value that determines for a restart
run if the case should initialize wells and groups from the historical
SCHEDULE section (true), or from the well and group data on the restart
file (false). Note that the commercial simulator always uses data from
the restart file (false). & FALSe \\
130 & --serial-partitioning & A boolean value that determines if
partitioning for parallel runs on a single process (true), or not
(false). & false \\
131 & --shut-unsolvable-wells & A boolean value that determines if the
simulator should shut unsolvable wells (true), or not (false). & true \\
132 & --solve-welleq-initially & A boolean value that determines if the
simulator should fully solve the well equations before each iteration of
the reservoir model (true), or not (false). Note that the well equations
are always added to the full system and solved until convergence. &
true \\
133 & --solver-continue-on-convergence-failure & A boolean value that
stipulates if the simulator should continue (true) instead of stopping
(false) when the minimum solver time step is reached. & false \\
134 & --solver-growth-factor & A real positive value that specifies the
growth factor a time step can be increased by when recovering from one
or more time step chops, subject to the maximum allowable time step size
set by the --solver-max-time-step-in-days parameter. For example, if the
current time step has converged at 5 days after a time step chop, and
--solver-growth-factor is set to the default value of 2.0, then the next
time step will be 2.0 * 5 days, that is at 10 days. & 2.0 \\
135 & --solver-max-growth & A real positive value that specifies the
maximum growth factor a time step can be increased by after a report
time step, subject to the maximum allowable time step size set by the
--solver-max-time-step-in-days parameter. Thus, if the current time step
has converged at 5 days after at a report step, and --solver-max-growth
is set to the default value of 3.0, then the next time step will be 3.0
* 5 days, that is at 15 days. & 3.0 \\
136 & --solver-max-restarts & A positive integer that defines the number
of allowed consecutive restarts (or time step chops) before the
simulation is terminated. & 10 \\
137 & --solver-max-time-step-in-days & A real positive double precision
value that specifies the maximum allowed time step size in days. &
365 \\
138 & --solver-min-time-step & A real positive double precision value
that specifies the minimum size of a time step in days for field and
metric units, and hours for laboratory units If a time step cannot
converge without getting cut below this time step size the simulator
will stop. & 1.0 x 10-12 \\
139 & --solver-restart-factor & A real positive double precision value
that sets the time step chop factor of the time step after a convergence
failure. For example, if the current non-convergent time step is 30 days
and --solver-restart-factor is set to the default value, then the time
step will be repeated using 0.33 * 30 days as the time step, that is 9.9
days. & 0.33 \\
140 & --solver-verbosity & A positive integer that specifies the
``chattiness'' of the non-linear solver. & 1 \\
141 & --strict-inner-iter-wells & A positive integer that specifies the
number of inner well iterations with strict tolerance. & 40 \\
142 & --strict-outer-iter-wells & A positive integer that specifies the
number of Newton iterations for which wells are checked with strict
tolerance. & 6 \\
143 & --temperature-max & A real positive value that sets the maximum
absolute temperature. & 1.0 x 109 \\
144 & --temperature-min & A real positive value that sets the minimum
absolute temperature. & 0 \\
145 & --threads-per-process & A positive integer value that stipulates
the maximum number of threads to be instantiated per process (`-1' means
`automatic'). & -1 \\
146 & --time-step-after-event-in-days & A real double precision value
that sets the maximum allowed time step after an event; for example,
when a well is open or closed etc. The default value of -1 means that
events to do effect the time stepping. & -1 \\
147 & --time-step-control & A defined character string that defines the
time stepping control algorithm and is set to one of the following: &
``pid+newtoniteration'' \\
148 & --time-step-control-decay-damping-factor & A real positive value
that specifies the decay rate a time step can be decreased by after the
number target iterations has been exceeded & 1.0 \\
149 & --time-step-control-decay-rate & A real positive value that
specifies the decay rate a time step size can be decreased by after the
number of target iterations has been exceeded & 0.75 \\
150 & --time-step-control-file-name & A character string enclosed in
quotes that specifies a filename where time steps are specified. The
default is the character string timesteps. For instance time steps can
be generated by the ecl\_summary application in libecl as per the
following UNIX command line: path\_to\_libecl\_applications/ecl\_summary
DECK TIME \textgreater{} filename Where: DECK is the name of the data
deck you want to get the time steps from, TIME tells the application to
return the timing for the run, and ``filename'' is the name of the file
the times are piped to. & ``timesteps'' \\
151 & --time-step-control-growth-damping-factor & A real positive value
that specifies the growth rate of the time step increase when the number
of target iterations are undercut. & 3.2 \\
152 & --time-step-control-growth-rate & A real positive value that
specifies the growth rate of the time step size when the number of
target iterations are undercut. & 1.25 \\
153 & --time-step-control-target-iterations & A positive integer that
defines the number of linear iterations which the time step control
scheme should aim for (if applicable). & 30 \\
154 & --time-step-control-target-newton-iterations & A positive integer
that specifies the number of Newton iterations which the time step
control scheme should aim for (if applicable). & 8 \\
155 & --time-step-control-tolerance & A real double precision value that
sets the tolerance for PID (only used with the pid and pid+ options
defined by the --time-step-control option). & 0.1 \\
156 & --time-step-verbosity & A positive integer that specifies the
``chattiness'' during the time integration. & 1 \\
157 & --tolerance-cnv & A real positive double precision value that
specifies the maximum non-linear tolerance error. This is the local
convergence tolerance (maximum of local saturation errors). & 0.01 \\
158 & --tolerance-cnv-relaxed & A real positive value that defines the
relaxed local convergence tolerance that applies for iterations after
the iterations with the strict tolerance. & 1.0 \\
159 & --tolerance-mb & A real positive double precision value that sets
the maximum mass balance error, that is the tolerated mass balance error
relative to total mass present. & 1.0 x 10-6 \\
160 & --tolerance-pressure-ms-wells & A real positive double precision
value that specifies the tolerance for the pressure equations for
multi-segment wells. & 1000 \\
161 & --tolerance-well-control & A real positive double precision value
that sets the maximum tolerance for the well control equations. & 1.0 x
10-7 \\
162 & --tolerance-wells & A real positive double precision value that
defines the maximum non-linear error for the well equations. & 1.0 x
10-4 \\
163 & -update-equations-scaling & A boolean value that switches on
(true) or off (false) the updating of the scaling factors for mass
balance equations during the simulation. & false \\
164 & --use-average-density-ms-wells & A boolean value that specifies
whether to approximate segment densitities by averaging over the segment
and its outlet (true) or not (false). & false \\
165 & --use-gmres & A boolean value that when set to true OPM Flow will
use Generalized Minimal Residual (``GMRES'') {[}Y. Saad, A flexible
inner-outer preconditioned GMRES algorithm, SIAM J. Sci. Statist.
Comput.,14, (1993).{]} and {[}Y. Saad and M.H. Schultz, GMRES: A
generalized minimum residual algorithm for solving nonsymmetric linear
systems, SIAM J. Sci. Statist. Comput., 7 (1986), pp.~856\{869.{]}
solver instead of Biconjugate Gradient Stabilized (``BiCGSTAB'') {[}Van
der Vorst, H. A. (1992). ``Bi-CGSTAB: A Fast and Smoothly Converging
Variant of Bi-CG for the Solution of Nonsymmetric Linear Systems''. SIAM
J. Sci. Stat. Comput. 13 (2): 631--644. doi:10.1137/0913035.
hdl:10338.dmlcz/104566{]} and {[}Sleijpen, G. L. G.; Fokkema, D. R.
(November 1993). ``BiCGstab(l) for linear equations involving
unsymmetric matrices with complex spectrum'' (PDF). Electronic
Transactions on Numerical Analysis. Kent, OH: Kent State University. 1:
11--32. ISSN 1068-9613.{]} as the linear solver within the Newton
iterations. & false \\
166 & --use-multisegment-well & A boolean value that when set to true
the simulator will use the well model for multi-segment wells instead of
the one for single-segment wells. & true \\
167 & --use-update-stabilization & A boolean value that switches on
(true) or off (false) the stabilized Newton option, that attempts to
detect and correct oscillations or stagnation during the Newton
iterations. This option may improve convergence for some cases. &
true \\
168 & --water-only-threshold & A real positive value that defines the
saturation threshold, for which cells with water saturations above or
equal to this threshold are considered one-phase water only. & 1.0 \\
169 & --zoltan-imbalance-tol & A real positive value that defines the
tolerable imbalance of the load balancing provided by Zoltan package. &
1.1 \\
170 & --zoltan-params & A character string that specifies the
configuration of the Zoltan partitioner. Valid options are: graph,
hypergraph or scotch. Alternatively, you can request a configuration to
be read from a JSON file by giving the filename here, ending with
`.json.' See https://sandialabs.github.io/Zoltan/ug\_html/ug.html for
available Zoltan options. & ``graph'' \\
Notes: flow --help instead use: flow --help-all to get a list of
supported command line Parameters. & & & \\
\end{longtable}

\emph{Table A.4: OPM Flow 2023-10 Command Line Options}

\chapter{Output File Formats}\label{output-file-formats}

This section of the manual describes the various output files generated
by OPM Flow and attempts to outline the format of the various files, in
order for third party software to read and write these type of files.
From the 2019-04 release a substantial effort has been undertaken to
make OPM Flow's output files compatible with the commercial simulator's
output files. This compatibility enables OMP Flow to ``restart'' from
the commercial simulator's generated RESTART files as well as the
commercial simulator to ``restart'' from the OPM Flow's RESTART files.
Although, this might not be the case for all models, as restarting a
simulation model from another simulators restart file is complicated and
may not be possible in all cases.

Where applicable, files written by OPM Flow can be loaded into OPM
ResInsight post processing software for further analysis and for
displaying the results. Please see section Error: Reference source not
found Error: Reference source not found for the various command line
options for setting the output format type.

OPM Flow, similar to the commercial simulator, writes out various files,
some of which are used by post processing software (OPM ResInsight) and
some that are directly used by the user, for example the *.PRT file that
contains various reports. Table A.1 summaries the various file formats
and the status of the file formats currently supported by OPM Flow.

\begin{longtable}[]{@{}
  >{\raggedright\arraybackslash}p{(\linewidth - 6\tabcolsep) * \real{0.2500}}
  >{\raggedright\arraybackslash}p{(\linewidth - 6\tabcolsep) * \real{0.2500}}
  >{\raggedright\arraybackslash}p{(\linewidth - 6\tabcolsep) * \real{0.2500}}
  >{\raggedright\arraybackslash}p{(\linewidth - 6\tabcolsep) * \real{0.2500}}@{}}
\toprule\noalign{}
\begin{minipage}[b]{\linewidth}\raggedright
File Type
\end{minipage} & \begin{minipage}[b]{\linewidth}\raggedright
Data Type
\end{minipage} & \begin{minipage}[b]{\linewidth}\raggedright
Description
\end{minipage} & \begin{minipage}[b]{\linewidth}\raggedright
OPM Flow Status
\end{minipage} \\
\midrule\noalign{}
\endhead
\bottomrule\noalign{}
\endlastfoot
DATA & Input Data & DATA files contain the input data in ASCII format
used to run OPM Flow. & Fully Supported, as outlined in the manual \\
DBG & Debug Data & This file contains ASCII developer debug output
specific to OPM Flow, that is there no compatibility with the commercial
simulator's DBG file. & OPM Flow Specific \\
EGRID & Structure Data & EGRID files contain the structural information
for the model via the COORD, ZCORN, etc. keywords, and employs the
extensible GRID format, which is more compact than the GRID formatted
files and is the only format supported by OPM Flow. The output for this
file type must be requested via the Error: Reference source not found in
the GRID section. The data is used in post-processing software to view
the grid and is not in itself used by OPM Flow. & Fully Supported \\
GRID & Structure Data & This file type contains the structural
information for the model via the COORD, ZCORN, etc. keywords, and
consists of two data formats: Neither of these two formats are currently
supported by OPM Flow. & Not Supported \\
INIT Index & Static Property Index & The INIT index file type specifies
and defines the format and data type written to the INIT Data file. This
file is not required by OPM Flow or OPM ResInsight to be able to read
and write the INIT data file. & Not Supported and Not Required \\
INIT Data & Static Property Data & This file type contains static model
properties, that is, In order to request this type of output see the
Error: Reference source not found in the GRID section. Activating this
keyword results in both the INIT static property index and data files
being written out by the commercial simulator. OPM Flow only writes out
this file as the index file is not required by OPM ResInsight. The data
is used in post-processing software to view the grid property data and
is not in itself used by OPM Flow. & Fully Supported \\
LOG & Output File & The LOG file for the commercial simulator contains
an ASCII copy of the output from the terminal. The file itself is not
generated by the commercial simulator but by an auxiliary program,
ECLRUN, that is used to call and execute all the commercial simulators.
OPM Flow does not generate a LOG file; however, the file can be
generated by using the Linux tee command, for example: flow CASENAME
\textbar{} tee CASENAME.LOG Which will copy the terminal output to the
CASENAME.LOG file. Alternatively, one can use the OPMRUN, the Graphical
User Interface (``GUI'') program for OPM Flow, that provides similar
functionality to the commercial simulator's auxiliary program. & OPM
Flow Specific \\
PRT & Output File & This file contains ASCII user output reports
specific to OPM Flow, although there is a strong desire to make the OPM
Flow reports similar to those that are produced by the commercial
simulator. & OPM Flow Specific \\
RESTART Index & Dynamic Solution Index & The RESTART index file type
specifies and defines the format and data type written to the RESTART
Data file. This file is not required by OPM Flow or OPM ResInsight to be
able to read and write the RESTART data file. & Not Supported and Not
Required \\
RESTART Data & Dynamic Solution Data & This file type contains dynamic
solution data for each active cell in the model at the user requested
time steps, that is, The RPTRST and the RPTSCHED keywords in the
SCHEDULE defines the data and frequency of the data to be written to the
RESTART file at each requested restart point. Activating this keyword
results in both the RESTART index and data files being written out by
the commercial simulator; OPM Flow only writes this file. In addition to
the solution data, this file type also includes group and well
configuration data (number of connections, connections open and closed,
etc.) This information is required to ensure that the correct group and
well configuration is available at a given time in order to be able to
restart the simulation at a restart point. The data is used to visualize
the simulation results of the model in two and three dimensional space
using post-processing software. The file is also used by OPM Flow to
``restart'' from a previous simulation case. & Fully Supported \\
RFT & Dynamic Wellbore Vector Data & Data written to the RFT file
consists of wellbore vector data, for example, pressure and saturation
versus wellbore depth at various time steps. The data written out is not
restricted to just Repeat Formation Tester (``RFT'') data, but can
contain any Production Logging Tool (``PLT'') data made available in the
simulator. Note that only the RFT data set is currently supported by OPM
Flow. The keywords WRFT and WRFTPLT in the SCHEDULE section are used to
define the wells, data, and time step at which the data should be
written to the RFT file. The data is used in post-processing software to
compare the actual wire line logging data with the simulation derived
results. The file is not in itself used by OPM Flow. & Fully Supported
(RFT) \\
RSM & Output File & The RSM file contains the variables requested to be
written to the SUMMARY file via the keywords described in the Error:
Reference source not found, in a tabulated ASCII output format. The
format of the file enables the data to be easily loaded into the
LibreOffice Calc program for further processing, as each vector
represents one column. The report is written at the end of the
simulation run by parsing the SUMMARY Index and SUMMARY Data files. &
Fully Supported \\
SAVE & Staic and Dynamic Data & SAVE files are basically a combination
of both the INIT and RESTART files, except for the static initialization
data that also includes the processed data, e.g., PORV, TRANX, TRANY,
TRANZ, etc. This enables ``fast restarts'' as the pore volumes and
transmissibility arrays do not have to be re-calculated in a ``restart''
run. This file type is currently not supported by OPM Flow. & Not
Supported \\
SUMMARY ESMRY & Dynamic Vector Data & This is an OPM FLOW specific
Enhanced SUMMARY (\emph{.ESMRY) output file format that is optimized for
fast loading of selected vectors by post-processing applications. The
}.ESMRY file is re-written for every time step and does not have a
performance impact on the simulation. Load times are in the order of 30
times faster compared with fetching selected vectors using the *.UNSMRY
file format in post processing software. & Fully Supported \\
SUMMARY Index & Dynamic Vector Index & The SUMMARY index file type
specifies and defines the format and data type written to the SUMMARY
Data file. & Fully Supported \\
SUMMARY Data & Dynamic Vector Data & The SUMMARY data file contains the
variables requested to be written to the file via the keywords described
in the Error: Reference source not found. The data consists of vectors
that are used to generate line graphs of properties such as oil flow
rate versus time, grid plot pressure versus time, etc. The properties to
be stored on the SUMMARY file are written to the summary file at the end
of each successful time step. The data can be used to compare actual
production data with the simulation derived results in post processing
software. Note that although the file format is fully supported, not all
of the SUMMARY variables are are available. & Fully Supported \\
Notes: & & & \\
\end{longtable}

\emph{Table A.1: OPM Flow Output File Types Summary}

As mentioned in Table A.1 all files can be written out in either ASCII
or binary formats. In addition, the SUMMARY and RESTART files can be
loaded or written out in either unified or non-unified formats. The file
type (ASCII or binary) and file structure format (unified or non-unified
formats) are set via various keywords in the RUNSPEC section, as
outlined in Table A.2.

\begin{longtable}[]{@{}
  >{\raggedright\arraybackslash}p{(\linewidth - 6\tabcolsep) * \real{0.2500}}
  >{\raggedright\arraybackslash}p{(\linewidth - 6\tabcolsep) * \real{0.2500}}
  >{\raggedright\arraybackslash}p{(\linewidth - 6\tabcolsep) * \real{0.2500}}
  >{\raggedright\arraybackslash}p{(\linewidth - 6\tabcolsep) * \real{0.2500}}@{}}
\toprule\noalign{}
\begin{minipage}[b]{\linewidth}\raggedright
Process
\end{minipage} & \begin{minipage}[b]{\linewidth}\raggedright
RUNSPEC Keyword
\end{minipage} & \begin{minipage}[b]{\linewidth}\raggedright
Description
\end{minipage} & \begin{minipage}[b]{\linewidth}\raggedright
Files
\end{minipage} \\
\midrule\noalign{}
\endhead
\bottomrule\noalign{}
\endlastfoot
Input & FMTIN & The keyword defines the input files to be formatted as
ASCII, i.e., text files, as oppose to binary files. The input deck file
is always of this type. The option relates to the OPM Flow derived files
that are used as input, for example when restarting from another case.
If the keyword is omitted then the default is for binary file input. &
\emph{.FEGRID }.FINSPEC \emph{.FINIT }.FRSSPEC \emph{.FUNRST }.FSMSPEC
*.FUNSMRY \\
MULTIN & A character string that defines the input files to be
non-unified multiple files, as opposed to unified files. In this case,
one file is read in per reporting time step, as opposed to all time
steps reports being read from one file. If the keyword is omitted then
the default is for one file per reporting time step. & \emph{.RSSPEC
}.X0001 \emph{.SMSPEC }.S0001 & \\
UNIFIN & UNIFIN defines the input files to be unified files, as opposed
to non-unified multiple files. A unified file is a single file
containing output for each reporting time step. For this option a single
summary file and a single restart file will be read. If the keyword is
omitted then the default is for one file per report time step. &
\emph{.RSSPEC }.UNRST \emph{.SMSPEC }.UNSMRY & \\
Output & FMTOUT & The keyword sets all output files to be formatted as
ASCII i.e.~text files, as opposed to binary files. The \emph{.DBG,
}.LOG, \emph{.PRT, and }.RSM files are always of this type. The option
relates to the OPM Flow output files only. In this case the files will
be portable across operating systems, but will also be very large in
terms of hard disk space. For this reason it is recommend that the
default option is used so that only binary files are written out. If the
keyword is omitted then the default is for binary file input. &
\emph{.FEGRID }.FINSPEC \emph{.FINIT }.FRSSPEC \emph{.FUNRST }.FSMSPEC
*.FUNSMRY \\
MULTOUT & A character string that defines the output files to be
non-unified multiple files, as opposed to unified files. In this case,
one file is written for each reporting time step, as opposed to all time
steps reports being written in one file. If the keyword is omitted then
the default is for one file per reporting time step. & \emph{.RSSPEC
}.X0001 \emph{.SMSPEC }.S0001 & \\
UNIFOUT & UNIFOUT defines the output files to be unified files, as
opposed to non-unified multiple files. A unified file is a single file
containing output for each reporting time step. Here a single summary
file and a single restart file will be generated, as opposed to one file
per report time step. If the keyword is omitted then the default is for
one file per report time step . & \emph{.RSSPEC }.UNRST \emph{.SMSPEC
}.UNSMRY & \\
Notes: & & & \\
\end{longtable}

\emph{Table A.2: RUNSPEC Input and Output File Format Keywords}

OPM Flow automatically generates the file names based on the input file
name and the output options selected via the keywords in the RUNSPEC
summarized in Table A.2. For example, the command line syntax for
running OPM Flow is:

\begin{verbatim}

flow [OPTIONS] CASENAME
\end{verbatim}

and typing the following command from the terminal:

\begin{verbatim}
flow CASENAME
\end{verbatim}

will start the simulator and run the case specified by CASENAME.DATA.
Here CASENAME is the ``root'' of the filename and DATA is the extension
of the filename. OPM Flow will generate the output files based on the
CASENAME with the extension based on the type (ASCII or binary) and file
structure format (unified or non-unified formats), as outlined in Table
A.3. Note that if CASENAME includes a relative or absolute path to the
data file, the run directory will be set so that all output files are
generated in the directory holding the data file.

\begin{longtable}[]{@{}
  >{\raggedright\arraybackslash}p{(\linewidth - 6\tabcolsep) * \real{0.2500}}
  >{\raggedright\arraybackslash}p{(\linewidth - 6\tabcolsep) * \real{0.2500}}
  >{\raggedright\arraybackslash}p{(\linewidth - 6\tabcolsep) * \real{0.2500}}
  >{\raggedright\arraybackslash}p{(\linewidth - 6\tabcolsep) * \real{0.2500}}@{}}
\toprule\noalign{}
\begin{minipage}[b]{\linewidth}\raggedright
File Type
\end{minipage} & \begin{minipage}[b]{\linewidth}\raggedright
Unformatted Extension
\end{minipage} & \begin{minipage}[b]{\linewidth}\raggedright
Description
\end{minipage} & \begin{minipage}[b]{\linewidth}\raggedright
Formatted Extension
\end{minipage} \\
\midrule\noalign{}
\endhead
\bottomrule\noalign{}
\endlastfoot
DATA & & Input data file. & *.DATA \\
DBG & & Output debug file. & *.DBG \\
EGRID & *.EGRID & EGRID files containing the structural grid data and is
the only format supported by OPM Flow. This output for this file type
must be requested via the Error: Reference source not found in the GRID
section. & *.FEGRID \\
GRID & *.GRID & Grid file in either the standard GRID file format or the
extended GRID file format Neither of these two formats are currently
supported by OPM Flow. & *.FGRID \\
INIT Index & *.INSPEC & The INIT index file for both unified and
non-unified formats & *.FINSPEC \\
INIT & *.INIT & This file type contains static model properties. &
*.FINIT \\
LOG & & Output log file that contains a copy of the output from the
terminal. The file itself is not generated by the commercial simulator
but by an auxiliary program, ECLRUN, that is used to call and execute
all the commercial simulators. & *.LOG \\
PRT & & Output print file. & *.PRT \\
RESTART Index & *.RSSPEC & The RESTART index file for both unified and
non-unified formats. This file is not required by OPM Flow or OPM
ResInsight to be able to read and write the RESTART data file. &
*.FRSSPEC \\
RESTART Data Non-Unified & *.Xnnnn & The RESTART data files containing
the solution arrays requested to be written to the RESTART files. For
example: CASENAME.X0001 CASENAME.X0002 CASENAME.X0003 etc.
CASENAME.RSSPEC & *.Fnnnn \\
RESTART Data Unified & *.UNRST & The RESTART data file containing the
solution arrays requested to be written to the RESTART file. For
example: CASENAME.UNRST CASENAME.RSSPEC & *.FUNRST \\
RFT & *.RFT & The RFT data file containing wellbore vector data
requested to be written to the RFT file. For example: CASENAME.RFT &
*.FRFT \\
RSM & & Output RSM file. & *.RSM \\
SAVE & *.SAVE & The SAVE file type is currently not supported by OPM
Flow. & *.FSAVE \\
SUMMARY ESMRY & *.ESMRY & OPM Flow Enhanced SUMMARY (*.ESMRY) output
file format that is optimized for fast loading of selected vectors by
post-processing applications. Note that the FMTOUT, MULTOUT, and UNIFOUT
keywords have no impact on this file type. & \\
SUMMARY Index & *.SMSPEC & The SUMMARY index file for both unified and
non-unified formats. & *.FSMSPEC \\
SUMMARY Data Non-Unified & *.Snnnn & The SUMMARY data files containing
the variables requested to be written to the SUMMARY files. For example:
CASENAME.S0001 CASENAME.S0002 CASENAME.S0003 etc. CASENAME.SMSPEC &
*.Annnn \\
SUMMARY Data Unified & *.UNSMRY & The SUMMARY data file containing the
variables requested to be written to the SUMMARY file. For example:
CASENAME.UNSMRY CASENAME.SMSPEC & *.FUNSMRY \\
Notes: & & & \\
\end{longtable}

\emph{Table A.3: OPM Flow File Naming Conventions}

Originally the commercial simulator was written in FORTRAN 77 and the
current version is still written in a version of FORTRAN. FORTRAN
unformatted file output is dependent on the FORTRAN compiler used to
generated the executable program. Typically an unformatted record
consists of a four byte prefix outlining the length of the record in
bytes, then the actual data record, followed by suffix containing the
length of the record in bytes. Most but not all compilers use four
bytes. This aids in reading records, for example, the four byte record
size at the end of the record assists with a backspace operation. If the
record size is greater than two Gigabytes (2\^{}30 bytes), the record is
divide into sub-records. In this case the sign bit of the prefix inform
that the record is continued by a sub-record or not and the sign bit of
the suffix indicate whether or not there is a preceding sub record.

OPM Flow unformatted files are written using the big-endian mode, that
is the most significant value in the sequence is stored at the lowest
storage address, that is first. This is opposite to the little-endian
mode, where the least significant value in the sequence is stored first.
For example, consider the number 1025 (2 to the tenth power plus one)
stored in a four byte integer:

\begin{verbatim}
00000000 00000000 00000100 00000001
\end{verbatim}

\begin{longtable}[]{@{}
  >{\raggedright\arraybackslash}p{(\linewidth - 4\tabcolsep) * \real{0.3333}}
  >{\raggedright\arraybackslash}p{(\linewidth - 4\tabcolsep) * \real{0.3333}}
  >{\raggedright\arraybackslash}p{(\linewidth - 4\tabcolsep) * \real{0.3333}}@{}}
\toprule\noalign{}
\begin{minipage}[b]{\linewidth}\raggedright
Address
\end{minipage} & \begin{minipage}[b]{\linewidth}\raggedright
Big-Endian representation of 1025
\end{minipage} & \begin{minipage}[b]{\linewidth}\raggedright
Little-Endian representation of 1025
\end{minipage} \\
\midrule\noalign{}
\endhead
\bottomrule\noalign{}
\endlastfoot
00 & 00000000 & 00000001 \\
01 & 00000000 & 00000100 \\
02 & 00000100 & 00000000 \\
03 & 00000001 & 00000000 \\
\end{longtable}

\emph{Table A.4: Big-Endian and Little-Endian Representation}

OPM Flow is written in C++ using the open source GNU C++ compiler and
therefore developers using C++ need to ensure that unformatted files are
read or written using the correct big-endian mode.

All the files consist of a series of ``header'' keywords that defined
various parameters based on their format type: integer, real, Boolean,
double precision, and character. For example for the INIT and RESTART
files the header keywords are:

\begin{verbatim}
- SEQNUM: Report time step keyword (RESTART file only).
- INTEHEAD: Integer (four byte) keyword that contains constant values for example, the dimensions of the model, as well as the location and length of data contained in the subsequent property keywords.
- LOGIHEAD: Boolean or logical (four byte) keyword that consists of a set of either T (true) or F (false) character strings, for example, if dissolved gas is present in the model then set to T for true or F for false.
- DOUBHEAD: Real double precision keyword (eight byte) that contains real values, for example, the current time step and TUNING keyword parameters. Note that the integer parameters on the TUNING keyword are stored in the INTEHEAD keyword data array.
\end{verbatim}

These are then followed by a series of ``property'' keywords that
outline the various additional data. For the RESTART files the header
keywords are repeated for each time step a RESTART record is written.
There are also additional header keywords for Local Grid Refinements
(``LGR'') that are repeated for each LGR and for time step.

All the keywords follow the same format, including the property
keywords, as outlined in Table A.5, which shows a portion of the
INTEHEAD header keyword definition for the RESTART file.

\begin{longtable}[]{@{}
  >{\raggedright\arraybackslash}p{(\linewidth - 10\tabcolsep) * \real{0.1667}}
  >{\raggedright\arraybackslash}p{(\linewidth - 10\tabcolsep) * \real{0.1667}}
  >{\raggedright\arraybackslash}p{(\linewidth - 10\tabcolsep) * \real{0.1667}}
  >{\raggedright\arraybackslash}p{(\linewidth - 10\tabcolsep) * \real{0.1667}}
  >{\raggedright\arraybackslash}p{(\linewidth - 10\tabcolsep) * \real{0.1667}}
  >{\raggedright\arraybackslash}p{(\linewidth - 10\tabcolsep) * \real{0.1667}}@{}}
\toprule\noalign{}
\begin{minipage}[b]{\linewidth}\raggedright
No.
\end{minipage} & \begin{minipage}[b]{\linewidth}\raggedright
Description
\end{minipage} & \begin{minipage}[b]{\linewidth}\raggedright
Output File Keyword Format
\end{minipage} & \begin{minipage}[b]{\linewidth}\raggedright
Status or Value
\end{minipage} & \begin{minipage}[b]{\linewidth}\raggedright
\end{minipage} & \begin{minipage}[b]{\linewidth}\raggedright
\end{minipage} \\
\midrule\noalign{}
\endhead
\bottomrule\noalign{}
\endlastfoot
Keyword & No.~of Entries & Data Type & & & \\
& INTEHEAD & Global grid header keyword for the RESTART Data file that
defines the integer variables for this time step. NIHEAD is the number
of entries associated with this keyword, and is dependent on the
features being used in the model. Undefined values should be set to
zero. & Required & & \\
1-1 & Format & INTEHEAD & NIHEAD & INTE & \\
2-1 & Data & An encoded integer corresponding to the time the file was
created. For files not originating from the commercial simulator, this
value may be set to zero. & ISNUM & & \\
2-2 & Data & Simulator version, should be set to zero. & VERSION & & \\
2-3 & Data & Unit system used in the simulation, set to 1 for metric, 2
for field, and 3 for laboratory & & & \\
2-4 & Data & Undefined. & 0 & & \\
Notes: & & & & & \\
\end{longtable}

\emph{Table A.5: Output File Keyword Format}

The row colored gray in the table specifies the start of a keyword
definition and shows the format for the keyword (INTEHEAD) together with
a description of the keyword. The next row, item 1-1 in the No.~column,
outlines the format of the keyword, that is the keyword name (INITHEAD),
the number of data entries for the keyword (NIHEAD), and the type of
entries (INTE for integers). The next set of rows specify the individual
data entries, for example, item 2-3 in the No.~column is the third
integer value that declares the unit system for the file. Note that
Cells colored red in the No.~column indicate that this item is either
``Not Used'' or ``Undefined''.

\begin{tcolorbox}[enhanced jigsaw, title=\textcolor{quarto-callout-note-color}{\faInfo}\hspace{0.5em}{Note}, coltitle=black, breakable, toprule=.15mm, colbacktitle=quarto-callout-note-color!10!white, leftrule=.75mm, opacitybacktitle=0.6, colback=white, bottomtitle=1mm, toptitle=1mm, titlerule=0mm, left=2mm, colframe=quarto-callout-note-color-frame, arc=.35mm, rightrule=.15mm, bottomrule=.15mm, opacityback=0]

Under the No.~column in Table A.5 the count base is one (offset one), as
per FORTRAN; however, in C++ the base is zero (offset zero) and
therefore C++ programmers must subtract one from this column to obtain
the correct reference. This reference adjustment should be applied to
all references in this Appendix. OPM Flow is written in C++, then
developers should use offset zero in reviewing or modifying the code.

\end{tcolorbox}

The Data Type variable in row 1-1 in Table A.5 should have one of the
following values:

\begin{verbatim}
- CHAR for eight character words enclosed in single quotes (ASCII character decimal 39) for formatted output.
- DOUB for double precision eight byte real values.
- INTE for standard four byte integers.
- LOGI for standard four byte Boolean (logical) values.
- MESS keyword that has no associated data.
- REAL for single precision four byte reals values.
\end{verbatim}

All keywords follow the same format including the property array like
PORO and PRESSURE. So for example, the INTEHEAD keyword would be like
this in a typical formatted RESTART file:

\begin{verbatim}
'INTEHEAD'         249 'INTE'
 -955283513      200400           2       -2345       -2345       -2345
      -2345       -2345          10          10           1         100
         10       -2345           7       -2345           0           1
          0           2           2           0           0           0
        110         108         109           3          97          93
      -2345       -2345          19          38          53       -2345
         97          93         146           5           0           1
         15          24           8           5           2           4
 ………………………….……………..
\end{verbatim}

And the pressure and water saturation arrays would be of the same form
in a typical formatted RESTART file:

\begin{verbatim}

 'PRESSURE'       44431 'REAL'
   0.26889725E+03   0.26838983E+03   0.26826810E+03   0.26820352E+03
   0.26814493E+03   0.26809378E+03   0.26807767E+03   0.26806363E+03
   0.26802625E+03   0.26798474E+03   0.26795001E+03   0.26791434E+03
   0.26787915E+03   0.26783920E+03   0.26777118E+03   0.26769208E+03         ………………………….……………..
 'SWAT    '       44431 'REAL'
   0.10500000E+00   0.10500000E+00   0.10500000E+00   0.10500000E+00
   0.10500000E+00   0.10500000E+00   0.10500000E+00   0.14000000E+00
   0.14000000E+00   0.14500000E+00   0.14500000E+00   0.14500000E+00
   0.14500000E+00   0.14500000E+00   0.16000000E+00   0.16000000E+00
………………………….……………..
\end{verbatim}

Secondly, the No.~Entries is the number of values following the keyword
and as such cannot be used to determine the size of an array that is not
associated with the grid, say for example the number of active
connections in a well. Instead, the array size is given in another
property keyword. In this case, the number of elements for each well in
the IWEL property array is given in the INTEHEAD keyword as shown in
Table A.6 (see Table A.21 for a full description of the RESTART INTEHEAD
keyword).

\begin{longtable}[]{@{}
  >{\raggedright\arraybackslash}p{(\linewidth - 10\tabcolsep) * \real{0.1667}}
  >{\raggedright\arraybackslash}p{(\linewidth - 10\tabcolsep) * \real{0.1667}}
  >{\raggedright\arraybackslash}p{(\linewidth - 10\tabcolsep) * \real{0.1667}}
  >{\raggedright\arraybackslash}p{(\linewidth - 10\tabcolsep) * \real{0.1667}}
  >{\raggedright\arraybackslash}p{(\linewidth - 10\tabcolsep) * \real{0.1667}}
  >{\raggedright\arraybackslash}p{(\linewidth - 10\tabcolsep) * \real{0.1667}}@{}}
\toprule\noalign{}
\begin{minipage}[b]{\linewidth}\raggedright
No.
\end{minipage} & \begin{minipage}[b]{\linewidth}\raggedright
keyword Name
\end{minipage} & \begin{minipage}[b]{\linewidth}\raggedright
RESTART Data - Header keywords
\end{minipage} & \begin{minipage}[b]{\linewidth}\raggedright
Status or Value
\end{minipage} & \begin{minipage}[b]{\linewidth}\raggedright
\end{minipage} & \begin{minipage}[b]{\linewidth}\raggedright
\end{minipage} \\
\midrule\noalign{}
\endhead
\bottomrule\noalign{}
\endlastfoot
Keyword & No.~of Entries & Data Type & & & \\
& INTEHEAD & Global grid header keyword for the RESTART Data file. &
Required & & \\
1-1 & Format & INTEHEAD & NIHEAD & INTE & \\
2-17 & Data & Number of wells in the model. & NWELLS & & \\
2-25 & Data & Number of values per well in the IWEL array (default value
is 97). & NIWELZ & & \\
\end{longtable}

\emph{Table A.6: INTEHEAD Keyword - IWEL Example}

Now if INTEHEAD(17) was equal to three and INTEHEAD(5) was equal to five
then the IWEL property array for the three wells would have the data
arranged as shown in Figure A.1.

Here the Items are the values affiliated with a specific well, and
Pos(ition) is the actual position in the IWEL property array, shown as
gray boxes in Figure A.1.

Similarly for well connections and completions. In this case the maximum
number of connections for all the wells is given in INTEHEAD(18) element
and a well's current number of connections is given via the IWEL(5)
element in the IWEL property array (see Table A.27 for a detailed
description of the IWEL keyword).

Now if the number of wells equal two (INTEHEAD(17)) and INTEHEAD(18) is
equal to four, that is the maximum number of connections is set to four,
together with IWEL(5) equal to three for the first well and four for the
second well, the ICON property array for the two wells would have the
data arranged as shown in Figure A.2.

Here the connection co-ordinates are given by (I, J, K) and the actual
position in the ICON array is shown in the gray colored boxes. Note that
Well-1 has the required maximum number of connections but the fourth
connection has the default values of zero, as the IWEL(5) value for this
well only has three active connections.

Note also that the number of data elements per connection in the ICON
vector is given by NICONZ which is given by INTEHEAD(33) and that the
size of the ICON vector is NICONZ x NCWMAX x NWELLS, where NCWMAX is set
by INTEHEAD (18) and NWELLS by INTEHEAD(17). In the example, NICONZ was
assumed to be three for illustrative purposes, while the actual value is
a constant 25 in the most recent versions of the commercial simulator.
These arbitrary constants, like NICONZ, use to determine vector lengths,
appear to be constant for various versions of the commercial simulator.

The following sections outline the format of the various individual
files supported by OPM Flow, except for the DBG, LOG, PRT, and RSM ASCII
format files.

EGRID files contain the structural information for the model via the
COORD, ZCORN, etc. keywords, and employs the extensible GRID format,
which is more compact than the GRID formatted files and is the only
format supported by OPM Flow.

The output for this file type must be requested via the Error: Reference
source not found in the GRID section, as shown below:

\begin{verbatim}
--
--       GRID FILE OUTPUT OPTIONS
--       GRID    EGRID
--       OPTN    OPTN
GRIDFILE
         0       1                                                             /
\end{verbatim}

The above example defines that no GRID file will be written out and that
the extensible GRID (that is the EGRID geometry format) file will be
produced. This is the only configuration that OPM Flow supports.

The FMTIN and FMTOUT keywords in the RUNSPEC section defines the input
and output files to be formatted as ASCII i.e.~text files,
(\emph{.FEGRID), as oppose to unformatted binary files (}.EGRID). If
these keywords are omitted then the default is for unformatted binary
file input.

The EGRID file structure consists of a series header keywords that
define the header keyword data type, the number of data entries, and
then the associated header keyword data. There are a total of four major
keyword header types for this file type: (1) File Header, (2) Global
Grid Header, (3) Local Grid Refinement Header {[}Several keywords are
used for both global and LGR grids, for example the GRIDHEAD, COORD,
ZCORN keywords.{]} (repeated for each local grid refinement in the
model), and (4) Non-Neighbor Connection Header.

In addition to the structural data, this file type also contains
information describing the geometry relationship between the global grid
and any Local Grid Refinements (``LGR'') that are present in the model,
together with the Non-Neighbor Connections (``NNC'') that may be present
due to structural discontinuities. The overall structure of this file is
is outlined in Table A.7

\begin{longtable}[]{@{}
  >{\raggedright\arraybackslash}p{(\linewidth - 12\tabcolsep) * \real{0.1429}}
  >{\raggedright\arraybackslash}p{(\linewidth - 12\tabcolsep) * \real{0.1429}}
  >{\raggedright\arraybackslash}p{(\linewidth - 12\tabcolsep) * \real{0.1429}}
  >{\raggedright\arraybackslash}p{(\linewidth - 12\tabcolsep) * \real{0.1429}}
  >{\raggedright\arraybackslash}p{(\linewidth - 12\tabcolsep) * \real{0.1429}}
  >{\raggedright\arraybackslash}p{(\linewidth - 12\tabcolsep) * \real{0.1429}}
  >{\raggedright\arraybackslash}p{(\linewidth - 12\tabcolsep) * \real{0.1429}}@{}}
\toprule\noalign{}
\begin{minipage}[b]{\linewidth}\raggedright
Reference
\end{minipage} & \begin{minipage}[b]{\linewidth}\raggedright
Section
\end{minipage} & \begin{minipage}[b]{\linewidth}\raggedright
Global Section keywords
\end{minipage} & \begin{minipage}[b]{\linewidth}\raggedright
Global Data keywords
\end{minipage} & \begin{minipage}[b]{\linewidth}\raggedright
LGR Section keywords
\end{minipage} & \begin{minipage}[b]{\linewidth}\raggedright
LGR Data keywords
\end{minipage} & \begin{minipage}[b]{\linewidth}\raggedright
Notes
\end{minipage} \\
\midrule\noalign{}
\endhead
\bottomrule\noalign{}
\endlastfoot
A.2.1 & FILEHEAD MAPUNITS MAPAXES GRIDUNIT GDORIENT & & & & EGRID header
keyword that defines various Global Grid properties. & \\
A.2.2 & GRIDHEAD & COORD & & & EGRID global irregular corner point grid
keyword that defines the grid geometry. & \\
ZCORN & & & & & & \\
etc. & & & & & & \\
A.2.2 & ENDGRID & & & & Marks the end of the EGRID global irregular
corner point grid keyword. & \\
A.2.3 & & & LGR LGRPARNT LGRSGRID & & EGRID LGR irregular corner point
grid keyword that defines the LGR name and optionally the nested or
parent LGRs associated with this LGR. & \\
A.2.3 & & & GRIDHEAD & COORD & EGRID LGR irregular corner point grid
keyword that defines the grid geometry. & \\
ZCORN & & & & & & \\
etc. & & & & & & \\
A.2.3 & & & ENDGRID & & Marks the end of the EGRID LGR irregular corner
point grid keyword. & \\
A.2.3 & & & ENDLGR & & Sets the end of current LGR section. & \\
A.2.3 & & & LGR LGRPARNT LGRSGRID & & EGRID LGR irregular corner point
grid keyword that defines the LGR name and optionally the nested or
parent LGRs associated with this LGR & \\
A.2.3 & & & GRIDHEAD & COORD & EGRID LGR irregular corner point grid
keyword that defines the grid geometry. & \\
ZCORN & & & & & & \\
etc. & & & & & & \\
A.2.3 & & & ENDGRID & & End of the EGRID LGR grid section. & \\
A.2.3 & & & ENDLGR & & Sets the end of current LGR section. & \\
A.2.4 & NNCHEAD & NNC1 & & & EGRID Global Non-Neighbor Connections for
Irregular Corner Point Grid keywords. & \\
NNC2 & & & & & & \\
A.2.4 & & & NNCHEAD & NNC1 & First LGR EGRID global and local
non-neighbor connections for irregular corner point grid keywords. & \\
NNC2 & & & & & & \\
NNCL & & & & & & \\
NNCG & & & & & & \\
A.2.4 & & & NNCHEAD & NNC1 & Second LGR EGRID global and local
non-neighbor connections for irregular corner point grid keywords. & \\
NNC2 & & & & & & \\
NNCL & & & & & & \\
NNCG & & & & & & \\
A.2.4 & & & NNCHEADA & NNA1 & EGRID amalgamated and local non-neighbor
connections for irregular corner point grid keywords. & \\
NNA2 & & & & & & \\
Notes: & & & & & & \\
\end{longtable}

\emph{Table A.7: EGRID - Model Structural Data for Irregular
Corner-Point Grids}

Each major header keyword is subdivided into a series of sub-header
keywords that define a particular data set. The general format for a
header keyword is the keyword (enclosed in single quotes), followed by
the number of data entries (an integer value), followed by the data type
which is set to CHAR, INTE, or REAL (enclosed in quotes). This will then
be followed by the data for the keyword.

The following sections describe the format of the various keywords and
the associated data sets.

The FILEHEAD keyword defines the key attributes of the file via a series
of integer constants, this is then followed by the MAPUNITS, MAPAXES,
GRIDUNIT, and GDORIENT keywords and their data sets. Note that some
keywords are optional and are not used by OPM Flow per se, but are used
by pre- and post processing software. Secondly, some features are not
supported by OPM Flow, for example local grid refinements, and thus OPM
will not write out the data associated with these features.

Table A.8 outlines the structure of the EGRID Header Keywords and their
affiliated data.

\begin{longtable}[]{@{}
  >{\raggedright\arraybackslash}p{(\linewidth - 10\tabcolsep) * \real{0.1667}}
  >{\raggedright\arraybackslash}p{(\linewidth - 10\tabcolsep) * \real{0.1667}}
  >{\raggedright\arraybackslash}p{(\linewidth - 10\tabcolsep) * \real{0.1667}}
  >{\raggedright\arraybackslash}p{(\linewidth - 10\tabcolsep) * \real{0.1667}}
  >{\raggedright\arraybackslash}p{(\linewidth - 10\tabcolsep) * \real{0.1667}}
  >{\raggedright\arraybackslash}p{(\linewidth - 10\tabcolsep) * \real{0.1667}}@{}}
\toprule\noalign{}
\begin{minipage}[b]{\linewidth}\raggedright
No.
\end{minipage} & \begin{minipage}[b]{\linewidth}\raggedright
Keyword Name
\end{minipage} & \begin{minipage}[b]{\linewidth}\raggedright
EGRID Header Keyword Table A.8
\end{minipage} & \begin{minipage}[b]{\linewidth}\raggedright
Status or Value
\end{minipage} & \begin{minipage}[b]{\linewidth}\raggedright
\end{minipage} & \begin{minipage}[b]{\linewidth}\raggedright
\end{minipage} \\
\midrule\noalign{}
\endhead
\bottomrule\noalign{}
\endlastfoot
Keyword & No.~of Entries & Data Type & & & \\
& FILEHEAD & This keyword marks the start of the EGRID file and defines
various parameters associated with this file type. & Required & & \\
1-1 & Format & FILEHEAD & 100 & INTE & \\
2-1 & Data & Version number for this file type, for example 3. & 3 &
& \\
2-2 & Data & The year this version of the file format was released, for
example 2004. & 2004 & & \\
2-3 & Data & Not used. & & & \\
2-4 & Data & Version number of earliest this file format is supported,
normally set to 0. & 0 & & \\
2-5 & Data & Type of grid, set to 0 for Irregular Corner-Point Grids, 1
for Unstructured Grids, or 2 for a mixture of Irregular Corner-Point and
Unstructured Grids (Mixed Grids). & 0 & & \\
2-6 & Data & Type of rock model, set to 0 for a single
porosity/permeability rock model, 1 for dual porosity rock model, or 3
for a dual permeability rock model. & 0 & & \\
2-7 & Data & Original grid format, set to 0 for unknown, 1 for Irregular
Corner-Point Grids, or 2 for Cartesian Regular Grids. & & & \\
2-8 to 2-100 & Data & Not used & 0 & & \\
& MAPUNITS & This keyword defines the MAPUNITS. & Optional & & \\
1-1 & Format & MAPUNITS & 1 & CHAR & \\
2-1 & Data & The units of the map data, normally set to FEET or METRES.
& & & \\
& MAPAXES & The MAPAXES keyword defines the grid position relative to
the map coordinates. The six values in the data fields define the
relative map (x, y) coordinates for three locations to enable conversion
from grid to map coordinate in pre- and post processing software. &
Optional & & \\
1-1 & Format & MAPAXES & 6 & REAL & \\
2-1 & Data & X-coordinate on the y-axis located at the end of the
y-axis. & & & \\
2-2 & Data & Y-coordinate on the y-axis located at the end of the
y-axis. & & & \\
2-3 & Data & X-coordinate at the origin. & & & \\
2-4 & Data & Y-coordinate at the origin. & & & \\
2-5 & Data & X-coordinate on the x-axis located at the end of the
x-axis. & & & \\
2-6 & Data & Y-coordinate on the x-axis located at the end of the
x-axis. & & & \\
& GRIDUNIT & The GRIDUNIT keyword defines units of the COORD and ZCORN
arrays and the grid coordinate space & Required & & \\
1-2 & Format & GRIDUNIT & 2 & CHAR & \\
2-1 & Data & The first entry defines units of the COORD and ZCORN arrays
and should be normally set to FEET or METRES (enclosed in single
quotes). & & & \\
2-2 & Data & The second entry defines the grid coordinate space (MAP or
MAPFT if COORD and ZCORN are defined in map coordinates, or blank if in
grid coordinates (` '). & & & \\
& GDORIENT & Grid orientation keyword defines the property ordering in
the I, J, and K dimensions, the direction of the Z-direction, and the
``handedness of the grid. & Optional & & \\
1-1 & Format & GDORIENT & 5 & CHAR & \\
2-1 & Data & Define the property ordering in the I dimension (set to
either INC for increasing or DEC for decreasing). & & & \\
2-2 & Data & Define the property ordering in the J dimension (set to
either INC for increasing or DEC for decreasing). & & & \\
2-3 & Data & Define the property ordering in the K dimension (set to
either INC for increasing or DEC for decreasing). & & & \\
2-4 & Data & Declare the direction of the Z-direction, should be set to
either UP or DOWN. & & & \\
2-5 & Data & Set the ``handedness of the grid., should be set to either
LEFT or RIGHT. & & & \\
Notes: & & & & & \\
\end{longtable}

\emph{Table A.8: EGRID Header Keywords}

\subsubsection{Example: EGRID Header
Keywords}\label{example-egrid-header-keywords}

The following example shows a typical formatted file EGRID header
keyword data set taken from the Norne model.

\begin{verbatim}
'FILEHEAD'         100 'INTE'
          3        2004           0           0           0           0
          0           0           0           0           0           0
          0           0           0           0           0           0
          0           0           0           0           0           0
          0           0           0           0           0           0
          0           0           0           0           0           0
          0           0           0           0           0           0
          0           0           0           0           0           0
          0           0           0           0           0           0
          0           0           0           0           0           0
          0           0           0           0           0           0
          0           0           0           0           0           0
          0           0           0           0           0           0
          0           0           0           0           0           0
          0           0           0           0           0           0
          0           0           0           0           0           0
          0           0           0           0
'MAPAXES '           6 'REAL'
  0.00000000E+00   0.10000000E+03   0.00000000E+00   0.00000000E+00
  0.10000000E+03   0.00000000E+00
'GRIDUNIT'           2 'CHAR'
'METRES  ' '        '
\end{verbatim}

Note that the GDORIENT keyword is not written out.

Following the EGRID Header Keyword data set, the EGRID Global Irregular
Corner Point Grid Keywords should be read or written to the EGRID file.
This data set includes the actual structural data for the model, that is
the COORD and ZCORN keyword data, as well as the ACTNUM keyword data
that defines if a given cell is active (set to 1), or inactive (set to
zero). The structure for this type data is defined in Table A.9.

\begin{longtable}[]{@{}
  >{\raggedright\arraybackslash}p{(\linewidth - 10\tabcolsep) * \real{0.1667}}
  >{\raggedright\arraybackslash}p{(\linewidth - 10\tabcolsep) * \real{0.1667}}
  >{\raggedright\arraybackslash}p{(\linewidth - 10\tabcolsep) * \real{0.1667}}
  >{\raggedright\arraybackslash}p{(\linewidth - 10\tabcolsep) * \real{0.1667}}
  >{\raggedright\arraybackslash}p{(\linewidth - 10\tabcolsep) * \real{0.1667}}
  >{\raggedright\arraybackslash}p{(\linewidth - 10\tabcolsep) * \real{0.1667}}@{}}
\toprule\noalign{}
\begin{minipage}[b]{\linewidth}\raggedright
No.
\end{minipage} & \begin{minipage}[b]{\linewidth}\raggedright
keyword Name
\end{minipage} & \begin{minipage}[b]{\linewidth}\raggedright
EGRID Global Irregular Corner Point Grid Keywords Table A.9
\end{minipage} & \begin{minipage}[b]{\linewidth}\raggedright
Status or Value
\end{minipage} & \begin{minipage}[b]{\linewidth}\raggedright
\end{minipage} & \begin{minipage}[b]{\linewidth}\raggedright
\end{minipage} \\
\midrule\noalign{}
\endhead
\bottomrule\noalign{}
\endlastfoot
Keyword & No.~of Entries & Data Type & & & \\
& GRIDHEAD & Global grid keyword that defines the start of this keyword
type. & Required & & \\
1-1 & Format & GRIDHEAD & 100 & INTE & \\
2-1 & Data & Type of grid, set to 0 for a mixture of Irregular
Corner-Point and Unstructured Grids, 1 for Irregular Corner-Point Grids,
and 2 for Unstructured Grids. Only the default value of one is supported
by OPM Flow. & 1 & & \\
2-2 & Data & The number of grid blocks in the x-direction (NX). & & & \\
2-3 & Data & The number of grid blocks in the y-direction (NY). & & & \\
2-4 & Data & The number of grid blocks in the z-direction (NZ). & & & \\
2-5 & Data & The grid reference number set to zero for the global grid
or a value greater than zero to represent a LGR. For example if there
are five LGR's, then this complete data set is repeated six times, once
for the global grid, and five times for the five LGRs., with the first
LGR having a reference number of 1 and the last LGR having a reference
number of 5. & & & \\
2-6 to 2-24 & Data & Not used. & 0 & & \\
2-25 & Data & NUMRES the number of reservoirs, that is the number of
COORD data sets to be processed by OPM Flow. This should be set to one
as the only value supported by OPM Flow & 1 & & \\
2-26 & Data & NSEG the number of coordinate line segments, should be set
to 1. & 1 & & \\
2-27 & Data & Cartesian/Radial grid indicator set to 0 for Cartesian
grids and greater than zero for radial grids. Currently OPM Flow does
not support radial geometries so this value should always be set to
zero. & 0 & & \\
2-28 & Data & The LGR location of the lower I-index LGR in the global
grid. This value should be set to zero if there is no LGR grids. & 0 &
& \\
2-29 & Data & The LGR location of the lower J-index LGR in the global
grid. This value should be set to zero if there is no LGR grids. & 0 &
& \\
2-30 & Data & The LGR location of the lower K-index LGR in the global
grid. This value should be set to zero if there is no LGR grids. & 0 &
& \\
2-31 & Data & The LGR location of the upperI -index LGR in the global
grid. This value should be set to zero if there is no LGR grids. & 0 &
& \\
2-32 & Data & The LGR location of the upper J-index LGR in the global
grid. This value should be set to zero if there is no LGR grids. & 0 &
& \\
2-33 & Data & The LGR location of the uper K-index LGR in the global
grid. This value should be set to zero if there is no LGR grids. & 0 &
& \\
2-34 to 2-100 & Data & Not used by OPM Flow. & 0 & & \\
& BOXORIG & Not supported. & Optional & & \\
1-1 & Format & BOXORIG & 3 & INTE & \\
2-1 to 2-3 & Data & Ignored by OPM Flow. & & & \\
& COORD & Property keyword that defines the start of the COORD data set.
The COORD data defines a set of coordinate lines or pillars for a
reservoir grid via an array. The number of data values for this keyword
should be 6 x (NX+1) x (NY+1) x NUMRES. Where NUMRES represents the
number of separate coordinate lines (or reservoirs). In OPM Flow NUMRES
can only be set to one. & Required & & \\
1-1 & Format & COORD & 6 x (NX+1) x (NY+1) & REAL & \\
2-1 & Data & COORD data in the same format as the COORD keyword
described in the GRID section of the manual. & & & \\
& COORDSYS & Coordinate definitions for each NUMRES reservoir. The
number of data values for this keyword should be 6 x NUMRES. Where
NUMRES represents the number of separate coordinate lines (or
reservoirs). In OPM Flow NUMRES can only be set to one. & Optional &
& \\
1-1 & Format & COORDSYS & 6 & INTE & \\
2-1 to 2-6 & Data & Ignored by OPM Flow. & & & \\
& ZCORN & Property keyword that defines the start of the ZCORN data set.
ZCORN defines the depth of each corner point of a grid block on the
pillars defining the reservoir grid. The number of data values for this
keyword should be 8 x NX x NY x NZ. & Required & & \\
1-2 & Format & ZCORN & 8 x NX x NY x NZ & REAL & \\
1-3 & Type & Type of data in the Data fields. & & & \\
2-1 & Data & ZCORN data in the same format as the ZCORN keyword
described in the GRID section of the manual. & & & \\
& ACTNUM & Property keyword that defines the start of the ACTNUM data
set. ACTNUM specifies which grid blocks are active or inactive. A value
of 1 indicates the block is active and a value of 0 indicates the block
is inactive. Although this data set status is set to Optional, it is
normally always written out by pre-processing software. The number of
data values for this keyword should be NX x NY x NZ integer values. &
Optional & & \\
1-1 & Format & ACTNUM & NX x NY x NZ & INTE & \\
2-1 & Data & ACTNUM data in the same format as the ACTNUM keyword
described in the GRID section of the manual. & & & \\
& CORSNUM & Property keyword that defines the start of the CORSNUM data
set for when the grid has been coarsened. The number of data values for
this keyword should be NX x NY x NZ integer values. Grid coarsening is
currently not supported in OPM Flow. & Optional & & \\
1-1 & Format & CORSNUM & NX x NY x NZ & INTE & \\
2-1 & Data & CORSNUM data set. & & & \\
1-1 & ENDGRID & The ENDGRID keyword marks the end of the these set of
keywords. Always set to 0 to indicate that there is no data section
following this keyword. & Required & & \\
1-2 & Format & ENDGRID & 0 & INTE & 0 \\
Notes: & & & & & \\
\end{longtable}

\emph{Table A.9: EGRID Global Irregular Corner Point Grid Keywords}

\subsubsection{Example: EGRID Global Irregular Corner Point Grid
Keywords}\label{example-egrid-global-irregular-corner-point-grid-keywords}

The following example shows a typical formatted EGRID Global Irregular
Corner Point Grid Keywords data set, taken from the Norne model.

\begin{verbatim}
'GRIDHEAD'         100 'INTE'
           1          46         112          22           0           0
           0           0           0           0           0           0
           0           0           0           0           0           0
           0           0           0           0           0           0
           1           1           0           0           0           0
           0           0           0           0           0           0
           0           0           0           0           0           0
           0           0           0           0           0           0
           0           0           0           0           0           0
           0           0           0           0           0           0
           0           0           0           0           0           0
           0           0           0           0           0           0
           0           0           0           0           0           0
           0           0           0           0           0           0
           0           0           0           0           0           0
           0           0           0           0           0           0
           0           0           0           0
 'COORD   '       31866 'REAL'
   0.45311400E+06   0.73199210E+07   0.30374729E+04   0.45311400E+06
   0.73199210E+07   0.31328311E+04   0.45315503E+06   0.73198400E+07
   0.29839331E+04   0.45314275E+06   0.73198640E+07   0.31735720E+04
   0.45319609E+06   0.73197590E+07   0.30059690E+04   0.45317150E+06
   0.73198075E+07   0.32158359E+04   0.45323716E+06   0.73196780E+07
   0.30002649E+04   0.45320025E+06   0.73197510E+07   0.32172500E+04
   0.45327819E+06   0.73195970E+07   0.29893479E+04   0.45322903E+06
   0.73196940E+07   0.32139509E+04   0.45331925E+06   0.73195165E+07
………………………….……………..
 'ZCORN   '      906752 'REAL'
   0.30374729E+04   0.29839331E+04   0.29839331E+04   0.30059690E+04
   0.30059690E+04   0.30002649E+04   0.30002649E+04   0.29893479E+04
   0.29893479E+04   0.29956799E+04   0.29956799E+04   0.30008550E+04
   0.30008550E+04   0.30052520E+04   0.30052520E+04   0.30308621E+04
   0.30308621E+04   0.30368701E+04   0.30368701E+04   0.30380171E+04
   0.30380171E+04   0.30450271E+04   0.30450271E+04   0.30554099E+04
   0.30554099E+04   0.30665410E+04   0.30665410E+04   0.30766240E+04
   0.30766240E+04   0.30869380E+04   0.30869380E+04   0.30961531E+04
………………………….……………..
 'ACTNUM  '      113344 'INTE'
           0           0           0           0           0           0
           0           0           0           0           0           0
           0           0           0           0           0           0
           0           0           0           0           0           0
           0           0           0           0           0           0
           0           0           0           0           0           0
           0           0           0           0           0           0
           0           0           0           0           0           0
           0           0           0           0           0           0
           0           0           0           0           0           0
………………………….……………..
 'ENDGRID '           0 'INTE'
\end{verbatim}

End of Example

The data structure for Local Grid Refinement (``LGR'') Irregular Corner
Point Grids is similar to the global grid format described in the
previous section. Additional data that defines the LGR properties (LGR
Name for example) are included in this definition and the LGR keywords
are repeated for each LGR in the model. The keyword description is
outlined in Table A.10.

Note that currently OPM Flow does not support LGR grids and therefore
this series of keywords cannot be used by OPM Flow.

\begin{longtable}[]{@{}
  >{\raggedright\arraybackslash}p{(\linewidth - 10\tabcolsep) * \real{0.1667}}
  >{\raggedright\arraybackslash}p{(\linewidth - 10\tabcolsep) * \real{0.1667}}
  >{\raggedright\arraybackslash}p{(\linewidth - 10\tabcolsep) * \real{0.1667}}
  >{\raggedright\arraybackslash}p{(\linewidth - 10\tabcolsep) * \real{0.1667}}
  >{\raggedright\arraybackslash}p{(\linewidth - 10\tabcolsep) * \real{0.1667}}
  >{\raggedright\arraybackslash}p{(\linewidth - 10\tabcolsep) * \real{0.1667}}@{}}
\toprule\noalign{}
\begin{minipage}[b]{\linewidth}\raggedright
No.
\end{minipage} & \begin{minipage}[b]{\linewidth}\raggedright
Keyword Name
\end{minipage} & \begin{minipage}[b]{\linewidth}\raggedright
EGRID LGR Irregular Corner Point Grid Keywords Table A.10 (Repeated for
Each LGR in the Model)
\end{minipage} & \begin{minipage}[b]{\linewidth}\raggedright
Status or Value
\end{minipage} & \begin{minipage}[b]{\linewidth}\raggedright
\end{minipage} & \begin{minipage}[b]{\linewidth}\raggedright
\end{minipage} \\
\midrule\noalign{}
\endhead
\bottomrule\noalign{}
\endlastfoot
Keyword & No.~of Entries & Data Type & & & \\
& LGR & LGR keyword defines the start of the LGR section of keywords and
is required when LGRs are in the model. & Required & & \\
1-1 & Format & LGR & 1 & CHAR & \\
2-1 & Data & Name of the LGR & & & \\
& LGRPARNT & Start of LGRPARNT keyword that defines the parent LGR for
when nested LGR have been used. & Optional & & \\
1-1 & Format & LGRPARNT & 1 & CHAR & \\
2-1 & Data & Name of parent LGR & & & \\
& LGRSGRID & Start of the LGRSGRID keyword for LGRs parented from a
sub-grid. & Optional & & \\
1-1 & Format & LGRSGRID & 1 & CHAR & \\
2-1 & Data & Name of parent LGR & & & \\
& GRIDHEAD & LGR grid keyword that defines the start of this keyword
type for the LGR. The format is the same as the Global grid keyword -
see Table A.9. & Required & & \\
& BOXORIG & Not supported. The format is the same as the Global grid
keyword - see Table A.9. & Optional & & \\
& COORD & Property keyword that defines the start of the LGR COORD data
set. The COORD data defines a set of coordinate lines or pillars for LGR
grid via an array. The format is the same as the Global grid keyword -
see Table A.9. & Required & & \\
& COORDSYS & Coordinate definitions for the LGR. The format is the same
as the Global grid keyword - see Table A.9. & Optional & & \\
& ZCORN & Property keyword that defines the start of the LGR ZCORN data
set. The format is the same as the Global grid keyword - see Table A.9.
& Required & & \\
& ACTNUM & Property keyword that defines the start of the LGR ACTNUM
data set. The format is the same as the Global grid keyword - see Table
A.9. & Optional & & \\
& HOSTNUM & Property keyword that defines the start of the HOSTNUM data
set. which is only applicable to LGR grids. Number of data values for
this keyword should be NX x NY x NZ integer values. & Required & & \\
1-1 & Format & HOSTNUM & NX x NY x NZ & INTE & \\
2-1 & Data & HOSTNUM data set. & & & \\
& ENDGRID & The ENDGRID keyword marks the end of the LGR Grid Header
section. Always set to 0 to indicate that there is no Data section for
this keyword. & Required & & \\
1-1 & Format & ENDGRID & 0 & INTE & 0 \\
& ENDLGR & ENDLGR keyword marks the end of the LGR Header section.
Always set to 0 to indicate that there is no Data section for this
keyword. & Required & & \\
1-1 & Format & ENDLGR & 0 & INTE & 0 \\
Notes: & & & & & \\
\end{longtable}

\emph{Table A.10: EGRID LGR Irregular Corner Point Grid Keywords}

Note that this data set is repeated for each LGR in the model as
demonstrated in the following example.

\subsubsection{Example: EGRID LGR Irregular Corner Point Grid
Keywords}\label{example-egrid-lgr-irregular-corner-point-grid-keywords}

The following example shows a typical formatted data set for this series
of keywords with two Cartesian LGR grids, named LGR-1 and LGR-2 define
via the following CARFIN keyword statements:

\begin{verbatim}
--
--       CARFIN LGR GRID COMMANDS
--
--       LGR        ----- FINE GRID ------   -- CARFIN GRID --  MAX     HOST
--       NAME       I1  I2  J1  J2  K1  K2     NX    NY    NZ   WELLS   NAME
CARFIN
        'LGR-1'      2   2   2   2   1   1      2     2      2     1    GLOBAL / CARFIN LGR GRID PARAMETERS

--
--       CARFIN LGR GRID COMMANDS
--
--       LGR        ----- FINE GRID ------   -- CARFIN GRID --  MAX     HOST
--       NAME       I1  I2  J1  J2  K1  K2     NX    NY    NZ   WELLS   NAME
CARFIN
        'LGR-2'      9   9   9   9   1   1      2     2      2     1    GLOBAL / CARFIN LGR GRID PARAMETERS

ENDFIN
\end{verbatim}

The resulting keywords are as follows:

\begin{verbatim}
'LGR     '           1 'CHAR'
'LGR-1   '
'LGRPARNT'           1 'CHAR'
'        '
'GRIDHEAD'         100 'INTE'
          1           2           2           2           1           0
          0           0           0           0           0           0
          0           0           0           0           0           0
          0           0           0           0           0           0
          1           1           0           2           2           1
 ………………………….……………..
           0           0           0           0
'COORD   '          54 'REAL'
  0.10000000E+04   0.10000000E+04   0.00000000E+00   0.10000000E+04
  0.10000000E+04   0.10000000E+04   0.15000000E+04   0.10000000E+04
  0.00000000E+00   0.15000000E+04   0.10000000E+04   0.10000000E+04
  0.20000000E+04   0.10000000E+04   0.00000000E+00   0.20000000E+04
  0.10000000E+04   0.10000000E+04   0.10000000E+04   0.15000000E+04
 ………………………….……………..
   0.20000000E+04   0.10000000E+04
'ZCORN   '          64 'REAL'
  0.83250000E+04   0.83250000E+04   0.83250000E+04   0.83250000E+04
  0.83250000E+04   0.83250000E+04   0.83250000E+04   0.83250000E+04
  0.83250000E+04   0.83250000E+04   0.83250000E+04   0.83250000E+04
  0.83250000E+04   0.83250000E+04   0.83250000E+04   0.83250000E+04
  0.83350000E+04   0.83350000E+04   0.83350000E+04   0.83350000E+04
 ………………………….……………..
   0.83450000E+04   0.83450000E+04   0.83450000E+04   0.83450000E+04
'ACTNUM  '           8 'INTE'
          1           1           1           1           1           1
          1           1
'HOSTNUM '           8 'INTE'
         12          12          12          12          12          12
         12          12
'ENDGRID '           0 'INTE'
'ENDLGR  '           0 'INTE'
'LGR     '           1 'CHAR'
'LGR-2   '
'LGRPARNT'           1 'CHAR'
'        '
'GRIDHEAD'         100 'INTE'
          1           2           2           2           2           0
  ………………………….……………..
          0           0           0           0
'COORD   '          54 'REAL'
  0.80000000E+04   0.80000000E+04   0.00000000E+00   0.80000000E+04
  0.80000000E+04   0.10000000E+04   0.85000000E+04   0.80000000E+04
  ………………………….……………..
  0.90000000E+04   0.10000000E+04
'ZCORN   '          64 'REAL'
  0.83250000E+04   0.83250000E+04   0.83250000E+04   0.83250000E+04
  ………………………….……………..
  0.83450000E+04   0.83450000E+04   0.83450000E+04   0.83450000E+04
'ACTNUM  '           8 'INTE'
          1           1           1           1           1           1
          1           1
'HOSTNUM '           8 'INTE'
         89          89          89          89          89          89
         89          89
'ENDGRID '           0 'INTE'
'ENDLGR  '           0 'INTE'
\end{verbatim}

End of Example

The keywords in this section are used to define Non-Neighbor Connections
(``NNC'') for both the global grid and any LGR's included in the model.
This data set is only written out if there are NNCs in the model and can
be skipped if this is the case; however, apart from very simple models,
nearly all full field models will have NNCs as a result of modeling the
structural faults in the model. Likewise, the keywords affiliate with
LGRs are also skipped if the there are no LGRs in the model.

Note that actual NNC transmissibility data for the NNCs (TRANX, TRANY,
and TRANZ) are stored in the INIT Data File and not in the EGRID File.

The keyword definitions are tabulated in Table A.11.

\begin{longtable}[]{@{}
  >{\raggedright\arraybackslash}p{(\linewidth - 10\tabcolsep) * \real{0.1667}}
  >{\raggedright\arraybackslash}p{(\linewidth - 10\tabcolsep) * \real{0.1667}}
  >{\raggedright\arraybackslash}p{(\linewidth - 10\tabcolsep) * \real{0.1667}}
  >{\raggedright\arraybackslash}p{(\linewidth - 10\tabcolsep) * \real{0.1667}}
  >{\raggedright\arraybackslash}p{(\linewidth - 10\tabcolsep) * \real{0.1667}}
  >{\raggedright\arraybackslash}p{(\linewidth - 10\tabcolsep) * \real{0.1667}}@{}}
\toprule\noalign{}
\begin{minipage}[b]{\linewidth}\raggedright
No.
\end{minipage} & \begin{minipage}[b]{\linewidth}\raggedright
Keyword Name
\end{minipage} & \begin{minipage}[b]{\linewidth}\raggedright
EGRID Non-Neighbor Connections Keywords for Irregular Corner Point Grids
Table A.11
\end{minipage} & \begin{minipage}[b]{\linewidth}\raggedright
Status or Value
\end{minipage} & \begin{minipage}[b]{\linewidth}\raggedright
\end{minipage} & \begin{minipage}[b]{\linewidth}\raggedright
\end{minipage} \\
\midrule\noalign{}
\endhead
\bottomrule\noalign{}
\endlastfoot
Keyword & No.~of Entries & Data Type & & & \\
& NNCHEAD & This keyword marks the start of the non-neighbor keywords
and defines the number of non-neighbor connections. & Required & & \\
1-1 & Format & NNCHEAD & 10 & INTE & \\
2-1 & Data & The number of non-neighbor connections (NUMNNC). & & & \\
2-2 & Data & Grid identifier set to zero for the global grid, or a value
greater than zero to represent a LGR. & & & \\
2-3 to 2-10 & Data & Not used. & & & \\
& NNC1 & NNC1 defines the cell numbers for the UPSTREAM non-neighbor
connections, where NUMNNC is the number of non-neighbor connections. &
Required & & \\
1-1 & Format & NNC1 & NUMNNC & INTE & \\
2-1 & Data & NUMNNC integer values defining the cell numbers for the
UPSTREAM non-neighbor connections. & & & \\
& NNC2 & NNC2 defines the cell numbers for the DOWNSTREAM non-neighbor
connections, where NUMNNC is the number of non-neighbor connections. &
Required & & \\
1-1 & Format & NNC2 & NUMNNC & INTE & \\
2-1 & Data & NUMNNC integer values defining the cell numbers for the
DOWNSTREAM non-neighbor connections. & & & \\
& NNCL & This keyword defines the local cells connected to the global
grid, where NCONGL is the number of entries. & Optional & & \\
1-1 & Format & NNCL & NCONCL & INTE & \\
2-1 & Data & NCONCL integer values defining the local cells connected to
the global grid. & & & \\
& NNCG & This keyword defines the global cells connected to the current
local grid being defined by this section of keywords, where NCONGL is
the number of entries. & Optional & & \\
1-1 & Format & NNCG & NCONGL & INTE & \\
2-1 & Data & NCONGL integer values defining the global cells connected
to the current local grid. & & & \\
& NNCHEADA & This keyword defines the connection between two LGRs that
have been amalgamated. & Optional & & \\
1-1 & Format & NNCHEADA & 10 & INTE & \\
2-1 & Data & LGR index of the first LGR in the amalgamation (ILOC1). & &
& \\
2-2 & Data & LGR index of the second LGR in the amalgamation (ILOC2). &
& & \\
2-3 to 2-10 & & Not used. & & & \\
& NNA1 & This keyword defines the ILOC1 cells connected in the
amalgamation, where NUMNCA is the number of entries. & Optional & & \\
1-1 & Format & NNA1 & NUMNCA & INTE & \\
2-1 & Data & NUMNCA integer values defining the ILOC1 cells connected in
the amalgamation. & & & \\
& NNA2 & This keyword defines the ILOC2 cells connected in the
amalgamation, where NUMNCA is the number of entries. & & & \\
1-1 & Format & NNA2 & NUMNCA & INTE & \\
2-1 & Data & NUMNCA integer values defining the ILOC2 cells connected in
the amalgamation. & & & \\
Notes: & & & & & \\
\end{longtable}

\emph{Table A.11: EGRID NNC Keywords for Irregular Corner Point Grids}

Note that if LGR's are present in the model then the keywords relating
to LGRs are repeated for each LGR in the model.

\subsubsection{Example: EGRID Grid Keywords for Non-Neighbor Connections
for Irregular Corner Point
Grids}\label{example-egrid-grid-keywords-for-non-neighbor-connections-for-irregular-corner-point-grids}

The example shows the NNCs for the global grid and two LGRs.

\begin{verbatim}
 'NNCHEAD '          10 'INTE'
       11363           0           0           0           0           0
           0           0           0           0
 'NNC1    '       11363 'INTE'
        5717        5717       21173       26325       31477       36629
       41781       46933       52085       57237       62389       67541
       72693       77845       82997       88149       93301      103605
………………………….……………..
 'NNC2    '       11363 'INTE'
         566       10870       26326       31478       36630       41782
       46934       52086       57238       62390       67542       72694
       77846       82998       88150       93302       98454       98454
      103606         612        5764       26372       21220       31524
………………………….……………..

 'NNCHEAD '          10 'INTE'
           0           1           0           0           0           0
           0           0           0           0
 'NNC1    '           0 'INTE'
 'NNC2    '           0 'INTE'
 'NNCL    '          16 'INTE'
           1           5           3           7           1           5
           2           6           2           6           4           8
           3           7           4           8
 'NNCG    '          16 'INTE'
          11          11          11          11           2           2
           2           2          13          13          13          13
          22          22          22          22
 'NNCHEAD '          10 'INTE'
           0           2           0           0           0           0
           0           0           0           0
 'NNC1    '           0 'INTE'
 'NNC2    '           0 'INTE'
 'NNCL    '          16 'INTE'
           1           5           3           7           1           5
           2           6           2           6           4           8
           3           7           4           8
 'NNCG    '          16 'INTE'
          88          88          88          88          79          79
          79          79          90          90          90          90
          99          99          99          99
\end{verbatim}

End of Example

This file format is currently not supported by OPM Flow.

This file format is currently not supported by OPM Flow.

This file format is currently not supported by OPM Flow.

This set of keywords define an LGR's properties, LGR Name for example,
and always proceeds an LGR data set for a given LGR; that is the
keywords are repeated for each LGR data set in the model and each time
step. Note that a given LGR data set is terminated by the ENDLGR keyword
and the SEQNUM keyword terminates the input for a time step, which
switches the input back to the global grid properties. The keyword
description is outlined in Table A.12.

Note that currently OPM Flow does not support LGR grids and therefore
this set of keywords cannot be used by OPM Flow.

\begin{longtable}[]{@{}
  >{\raggedright\arraybackslash}p{(\linewidth - 10\tabcolsep) * \real{0.1667}}
  >{\raggedright\arraybackslash}p{(\linewidth - 10\tabcolsep) * \real{0.1667}}
  >{\raggedright\arraybackslash}p{(\linewidth - 10\tabcolsep) * \real{0.1667}}
  >{\raggedright\arraybackslash}p{(\linewidth - 10\tabcolsep) * \real{0.1667}}
  >{\raggedright\arraybackslash}p{(\linewidth - 10\tabcolsep) * \real{0.1667}}
  >{\raggedright\arraybackslash}p{(\linewidth - 10\tabcolsep) * \real{0.1667}}@{}}
\toprule\noalign{}
\begin{minipage}[b]{\linewidth}\raggedright
No.
\end{minipage} & \begin{minipage}[b]{\linewidth}\raggedright
keyword Name
\end{minipage} & \begin{minipage}[b]{\linewidth}\raggedright
RESTART Data - LGR Grid Header Keywords Table A.12 (Repeated for Each
LGR in the Model)
\end{minipage} & \begin{minipage}[b]{\linewidth}\raggedright
Status or Value
\end{minipage} & \begin{minipage}[b]{\linewidth}\raggedright
\end{minipage} & \begin{minipage}[b]{\linewidth}\raggedright
\end{minipage} \\
\midrule\noalign{}
\endhead
\bottomrule\noalign{}
\endlastfoot
Keyword & No.~of Entries & Data Type & & & \\
& LGR & LGR grid header keyword that defines the name of the LGR. &
Always Required & & \\
1-1 & Keyword & LGR & 1 & CHAR & \\
2-1 & Data & Name of the LGR. & & & \\
& LGRHEADI & LGR that defines the integer variables for this time step.
& Always Required & & \\
1-1 & Format & LGRHEADI & 45 & INTE & \\
2-1 to 2-45 & Data & Undefined. & & & \\
& LGRHEADQ & LGR grid keyword that defines the logical variables (T for
true and F for false) for this time step. & Always Required & & \\
1-1 & Format & LGRHEADQ & 5 & LOGI & \\
2-1 to 2-5 & Data & Undefined. & & & \\
& LGRHEADD & LGR grid keyword that defines the double precision REAL
variables for this time step. & Always Required & & \\
1-1 & Format & LGRHEADD & 5 & DOUB & 5 \\
2-1 to 2-5 & Data & Undefined. & & & \\
& ENDLGR & The ENDLGR keyword marks the end of an LGR section for a
given LGR. & Always Required & & \\
1-1 & Format & ENDLGR & 1 & MESG & \\
& Data & LGR number. & & & \\
Notes: & & & & & \\
\end{longtable}

\emph{Table A.12: RESTART Data - LGR Grid Header Keywords}

This keyword type is repeated for each LGR in the model with the last
LGR data set terminated by the LGRSGONE keyword to mark the end of the
LGR input.

\begin{tcolorbox}[enhanced jigsaw, title=\textcolor{quarto-callout-note-color}{\faInfo}\hspace{0.5em}{Note}, coltitle=black, breakable, toprule=.15mm, colbacktitle=quarto-callout-note-color!10!white, leftrule=.75mm, opacitybacktitle=0.6, colback=white, bottomtitle=1mm, toptitle=1mm, titlerule=0mm, left=2mm, colframe=quarto-callout-note-color-frame, arc=.35mm, rightrule=.15mm, bottomrule=.15mm, opacityback=0]

The LGR, LGRHEADI, LGRHEADQ, and LGRHEADD always proceed LGR data sets
and the ENDLGR keyword terminates the end of a data section of LGR data.

\end{tcolorbox}

\subsubsection{Example: RESTART Data - LGR Grid Header
Keywords}\label{example-restart-data---lgr-grid-header-keywords}

The following example shows a typical formatted LGR grid header data set
for a single LGR grid named LGR-1.

\begin{verbatim}
'LGR     '           1 'CHAR'
'LGR-1   '
'LGRHEADI'          45 'INTE'
          1         100       -2345       -2345       -2345       -2345
      -2345       -2345       -2345       -2345       -2345       -2345
      -2345       -2345       -2345       -2345       -2345       -2345
      -2345       -2345       -2345       -2345       -2345       -2345
      -2345       -2345       -2345       -2345       -2345       -2345
      -2345       -2345       -2345       -2345       -2345       -2345
      -2345       -2345       -2345       -2345       -2345       -2345
      -2345       -2345       -2345
'LGRHEADQ'           5 'LOGI'
 F  F  F  F  F
'LGRHEADD'           5 'DOUB'
  0.00000000000000D+00  -0.10000000200409D+21  -0.10000000200409D+21
 -0.10000000200409D+21  -0.10000000200409D+21
\end{verbatim}

End of Example

The INIT files contain the static data specified in the GRID, PROPS, and
REGIONS sections. For example, the PORO, PERMX, and NTG arrays from the
GRID section and the resulting PORV, TRANX, TRANY, and TRANZ arrays. The
data is used in post-processing software, for example OPM ResInsight, to
visualize the static grid properties and the tabular PVT and relative
permeability data, and is not used per se by the simulator.

The data is written to file if the INIT keyword in the GRID section has
been activated (Error: Reference source not found). Similar to the EGRID
file, the INIT file can either be written out in formatted form as ASCII
i.e.~text files, if the FMTOUT keyword has been activated, or binary
format if the FMTOUT keyword has not been activated. Normally, this
option is always activated by the user and when activated the binary
form of the file is used.

There are two files types associated with the INIT data:

\begin{verbatim}
- an index file (INIT Index) that contains a list of variables, the array type (character, integer, or real) and the size of the arrays that are written out to the initialization data file and,
- the initialization data file (INIT Data) that contains the actual data written out at the beginning of the run once the model has been initialized.
\end{verbatim}

The commercial simulator writes out both files, whereas OPM Flow only
writes out the INIT Data file, as the data can be read by most
post-processing software, including OPM ResInsight, using only the INIT
Data file as input.

The following sections outline the format for these two file types.

This file type is not supported or required by OPM Flow or OPM
ResInsight to read or write the static data written to the INIT Data
file.

This file type contains the global and any LGR grid property data
present in the model {[}Several keywords are used for both global and
LGR grids, for example the INTEHEAD, LOGIHEAD, DOUBHEAD, PROPERTY, and
REGION keywords and arrays.{]}, for example porosity (PORO) and
permeability data arrays (PERMX, PERMY, and PERMZ) for properties
allocated to each grid cell, as well as the fluid and rock property
functions tabular data. The overall structure of this file is similar to
the RESTART file with the individual keywords being structured as shown
in Table A.13.

\begin{longtable}[]{@{}
  >{\raggedright\arraybackslash}p{(\linewidth - 12\tabcolsep) * \real{0.1429}}
  >{\raggedright\arraybackslash}p{(\linewidth - 12\tabcolsep) * \real{0.1429}}
  >{\raggedright\arraybackslash}p{(\linewidth - 12\tabcolsep) * \real{0.1429}}
  >{\raggedright\arraybackslash}p{(\linewidth - 12\tabcolsep) * \real{0.1429}}
  >{\raggedright\arraybackslash}p{(\linewidth - 12\tabcolsep) * \real{0.1429}}
  >{\raggedright\arraybackslash}p{(\linewidth - 12\tabcolsep) * \real{0.1429}}
  >{\raggedright\arraybackslash}p{(\linewidth - 12\tabcolsep) * \real{0.1429}}@{}}
\toprule\noalign{}
\begin{minipage}[b]{\linewidth}\raggedright
Reference
\end{minipage} & \begin{minipage}[b]{\linewidth}\raggedright
Section
\end{minipage} & \begin{minipage}[b]{\linewidth}\raggedright
Global Section keywords
\end{minipage} & \begin{minipage}[b]{\linewidth}\raggedright
Global Data keywords
\end{minipage} & \begin{minipage}[b]{\linewidth}\raggedright
LGR Section keywords
\end{minipage} & \begin{minipage}[b]{\linewidth}\raggedright
LGR Data keywords
\end{minipage} & \begin{minipage}[b]{\linewidth}\raggedright
Notes
\end{minipage} \\
\midrule\noalign{}
\endhead
\bottomrule\noalign{}
\endlastfoot
A.6.3 & INTEHEAD LOGIHEAD DOUBHEAD & & & & Global grid header keywords
that defines the integer, logical, and double precision variables for
this header keyword. & \\
A.6.4 & Property Keywords & PORV & & & Global grid property keywords
that defines the global grid property data defined in the GRID and Edit
sections. For example, PORV, NTG, PERMX, etc.) & \\
PORO & & & & & & \\
etc. & & & & & & \\
A.6.5 & & & LGR LGRHEADI LGRHEADQ LGRHEADD & & Defines the start of an
LGR section of keywords and the LGR header keywords. & \\
A.6.5 & & & INTEHEAD LOGIHEAD DOUBHEAD & & LGR grid header keywords
defining the same data as the global keywords, except for the given LGR.
This set of keywords are only written once for each LGR. & \\
A.6.4 & & & Property Keywords & PORV & LGR property keywords; same
format as the global grid keywords. & \\
PORO & & & & & & \\
etc. & & & & & & \\
A.6.5 & & & LGR LGRHEADI LGRHEADQ LGRHEADD & & The LGR header keywords.
& \\
A.6.5 & & & INTEHEAD LOGIHEAD DOUBHEA & & The LGR header keywords which
are only written once for each LGR. & \\
A.6.4 & & & Property Keywords & PORV & LGR property keywords. & \\
PORO & & & & & & \\
etc. & & & & & & \\
& & & etc. & & Repeated for each LGR in the model. & \\
A.6.5 & & & LGRSGONE & & Marks the end of LGR section of keywords. & \\
A.6.7 & TABDIMS & & & & TABDIMS defines the tables dimensions, for the
subsequent TAB keywords. & \\
A.6.8 & TAB & & & & Rock and fluid tabular property data keywords. & \\
A.6.8 & CON & & & & CON keyword defines the constant PVT dead oil data
(the data associated with RSCONST, RSCONSTT, RVCONST, and RVCONSTT
keywords). & \\
A.6.4 & Region Keywords & PVTNUM & & & Region property data for the
global grid as defined in the REGION section (PVTNUM, SATNUM, etc.). The
keyword is repeated to account for all region property data. & \\
SATNUM & & & & & & \\
etc. & & & & & & \\
A.6.5 & & & LGR LGRHEADI LGRHEADQ LGRHEADD & & The LGR header keywords.
& \\
A.6.4 & & & Region Keywords & PVTNUM & LGR region property data for the
LGR grid. & \\
SATNUM & & & & & & \\
etc. & & & & & & \\
A.6.5 & & & LGR LGRHEADI LGRHEADQ LGRHEADD & & The LGR header keywords.
& \\
A.6.4 & & & Region Keywords & PVTNUM & LGR region property data for the
LGR grid. & \\
SATNUM & & & & & & \\
etc. & & & & & & \\
& & & etc. & & Repeated for each LGR in the model. & \\
A.6.5 & & & LGRSGONE & & Marks the end of LGR section of keywords. & \\
A.6.4 & End-Point Keywords & SWATINIT & & & Global grid saturation and
end-point data keywords. The keyword uses the same keywords outlined in
the GRID and PROPS section, for example, the SWL, SWATINIT, KRG, PCW,
etc. arrays & \\
SWL & & & & & & \\
etc. & & & & & & \\
A.6.5 & & & LGR LGRHEADI LGRHEADQ LGRHEADD & & The LGR header keywords.
& \\
A.6.4 & & & End-Point Keywords & SWATINIT & LGR grid saturation and
end-point data keywords. & \\
SWL & & & & & & \\
etc. & & & & & & \\
A.6.5 & & & LGR LGRHEADI LGRHEADQ LGRHEADD & & The LGR header keywords.
& \\
A.6.4 & & & End-Point Keywords & SWATINIT & LGR grid saturation and
end-point data keywords. & \\
SWL & & & & & & \\
etc. & & & & & & \\
& & & etc. & & & \\
A.6.5 & & & LGRSGONE & & Marks the end of LGR section of keywords. & \\
A.6.4 & NNC Keywords & TRANNC & & & Non-neighbor connection property
data (``NNC'') for the global grid. The keyword is repeated to account
for various NNC arrays. & \\
HEATNNC & & & & & & \\
etc. & & & & & & \\
A.6.5 & & & LGR LGRHEADI LGRHEADQ LGRHEADD & & The LGR header keywords.
& \\
A.6.4 & & & NNC Keywords & TRANNC & LGR NNC property data. & \\
TRANGL & & & & & & \\
etc. & & & & & & \\
A.6.5 & & & LGR LGRHEADI LGRHEADQ LGRHEADD & & The LGR header keywords.
& \\
A.6.4 & & & NNC Keywords & TRANNC & LGR NNC property data. & \\
TRANGL & & & & & & \\
etc. & & & & & & \\
& & & etc. & & Repeated for each LGR in the model. & \\
A.6.5 & & & LGRSGONE & & Marks the end of LGR section of keywords. & \\
Notes: & & & & & & \\
\end{longtable}

\emph{Table A.13: INIT Data File Structure}

The INIT Data -- Static Data Header Keyword section is the first data
set that should be read or written to the INIT Data file. This data set
includes INTEHEAD, LOGIHEAD, and DOUBHEAD keywords that define versus
parameters used in subsequent keywords in the INIT Data file. This is
then followed by a series of global grid PROPERTY data keywords that
define the various global grid property arrays for each grid cell in the
model, as well as the static tabular data (relative permeability tables
and PVT tables), etc.

The structure for this set of keywords is defined in Table A.14.

\begin{longtable}[]{@{}
  >{\raggedright\arraybackslash}p{(\linewidth - 10\tabcolsep) * \real{0.1667}}
  >{\raggedright\arraybackslash}p{(\linewidth - 10\tabcolsep) * \real{0.1667}}
  >{\raggedright\arraybackslash}p{(\linewidth - 10\tabcolsep) * \real{0.1667}}
  >{\raggedright\arraybackslash}p{(\linewidth - 10\tabcolsep) * \real{0.1667}}
  >{\raggedright\arraybackslash}p{(\linewidth - 10\tabcolsep) * \real{0.1667}}
  >{\raggedright\arraybackslash}p{(\linewidth - 10\tabcolsep) * \real{0.1667}}@{}}
\toprule\noalign{}
\begin{minipage}[b]{\linewidth}\raggedright
No.
\end{minipage} & \begin{minipage}[b]{\linewidth}\raggedright
keyword Name
\end{minipage} & \begin{minipage}[b]{\linewidth}\raggedright
INIT Data -- Static Data Header Keywords Table A.14
\end{minipage} & \begin{minipage}[b]{\linewidth}\raggedright
Status or Value
\end{minipage} & \begin{minipage}[b]{\linewidth}\raggedright
\end{minipage} & \begin{minipage}[b]{\linewidth}\raggedright
\end{minipage} \\
\midrule\noalign{}
\endhead
\bottomrule\noalign{}
\endlastfoot
Keyword & No.~of Entries & Data Type & & & \\
& INTEHEAD & Global static header keyword for the INIT Data file that
defines the start of INIT static data set, and contains the integer
variables for this header keyword. NIHEAD is the number of entries
associated with this keyword, and is dependent on the features being
used in the model. & Required & & \\
1-1 & Format & INTEHEAD & NIHEAD & INTE & \\
2-1 & & Undefined. & 0 & & \\
2-2 & & Undefined. & 0 & & \\
2-3 & Data & Unit system used in the simulation, set to 1 for metric, 2
for field, and 3 for laboratory. & & & \\
2-4 to 2-8 & & Undefined. & 0 & & \\
2-9 & Data & DIMENS key word in the RUNSPEC section - the number of grid
blocks in the x-direction. & NX & & \\
2-10 & Data & DIMENS key word in the RUNSPEC section - the number of
grid blocks in the y-direction. & NY & & \\
2-11 & Data & DIMENS key word in the RUNSPEC section - the number of
grid blocks in the z-direction. & NZ & & \\
2-12 & Data & Number of global active cells in the model. & NACTIV &
& \\
2-13 & Data & Not used. & 0 & & \\
2-14 & Data & Type of grid, set to 0 for Irregular Corner-Point Grids, 1
for Unstructured Grids, or 2 for a mixture of Irregular Corner-Point,
Unstructured Grids (Mixed Grids) and 3 for Cartesian grids. & 1 & & \\
2-15 & Data & Type of phases in the model, set to 1 for oil, 2 for
water, 3 for oil-water, 4 for gas, 5 for oil-gas, 6 for gas-water and 7
for oil-water-gas. & PHASE & & \\
2-16 to 2-64 & & Undefined. & 0 & & \\
2-65 & Data & First part of start date of the run, DAY, the value should
be greater than or equal to one and less than or equal to 31 for the
calendar days. & DAY & & \\
2-66 & Data & Second part of start date of the run, MONTH, the value
should be greater than or equal to one and less than or equal to 12 for
the calendar month. & MONTH & & \\
2-77 & Data & Third part of start date of run, YEAR, a positive four
digit integer value of the start year, which must be specified fully by
four digits, that is 1986. & YEAR & & \\
2-68 to 2-94 & & Undefined. & 0 & & \\
2-95 & Data & Simulator code used to generate the file, set to 100 for
Schlumberger's ECLIPSE 100, 300 for Schlumberger's ECLIPSE 300, 500 for
Schlumberger's ECLIPSE 300 Thermal, 700 for Schlumberger's INTERSECT
simulator, 800 for Schlumberger's FrontSim simulator, or a negative
value for other simulators. OPM Flow uses a value of 100. & IPROG & & \\
2-96 to 206 & & Undefined. & 0 & & \\
2-207 & Data & The first part of the current simulation time in the form
HH:MM:SS, that is HH for hours. The value should be greater than or
equal to zero and less than or equal to 23. & IHOURZ & & \\
2-208 & Data & The second part of the current simulation time in the
form HH:MM:SS, that is MM for minutes. The value should be greater than
or equal to zero and less than or equal to 59. & IMINTS & & \\
2-209 to 410 & & Undefined. & 0 & & \\
411 & Data & The third part of the current simulation time in the form
HH:MM:SS, that is SS for seconds. The value should be greater than or
equal to zero and less than or equal to 59. & ISECND & & \\
& LOGIHEAD & Global static data keyword for the INIT Data file that
defines the logical variables (T for true and F for false) for this
keyword. NLHEAD is the number of entries associated with this keyword,
and is dependent on the features being used in the model. Undefined
items should be set to F (False). & Required & & \\
1-1 & Format & LOGIHEAD & NLHEAD & LOGI & \\
2-1 & Data & Dissolved gas is present in the model indicator, set to T
if present in the model else set to F. & DISGAS & & \\
2-2 & Data & Vaporized oil is present in the model indicator, set to T
if present in the model else set to F. & VAPOIL & & \\
2-3 & Data & SATOPTS in the RUNSPEC section - directional relative
permeability curves are active in the model. & DIRECT & & \\
2-4 & Data & SATOPTS keyword in the RUNSPEC section - reversible
relative permeability (black-oil) are active in the model. & IRREVERS &
& \\
Radial grid geometry has been used for a compositional model indicator,
set to T or F. Note that currently OPM Flow does not support radial
grids and therefore this item should be set to F. & RADIAL & & & & \\
2.5 & Data & Radial grid geometry has been used for a black-oil model
indicator, set to T or F. Note that currently OPM Flow does not support
radial grids and therefore this item should be set to F. & RADIAL & & \\
Reversible relative permeability (compositional) are active in the
model. & IRREVERS & & & & \\
2.6 & & Undefined. & F & & \\
2-7 & Data & The hysteresis option has been activated in the model
indicator, set to T or F. & HYSTER & & \\
2-8 to 2-14 & Data & Undefined. & F & & \\
2-15 & Data & The dual porosity option has been activated in the model
indicator, set to T or F. Note that currently OPM Flow does not support
dual porosity grids and therefore this item should be set to F. &
DUALPORO & & \\
2-16 & & Undefined. & F & & \\
2-17 & Data & ENDSCALE keyword in the RUNSPEC section - enable end-point
scaling indicator. & ENDSCALE & & \\
2-18 & Data & ENDSCALE keyword in the RUNSPEC section - directional
end-point scaling indicator. & DIRECT & & \\
2.19 & Data & ENDSCALE keyword in the RUNSPEC section - reversible
end-point scaling indicator. & IRREVERS & & \\
2-20 & Data & SCALECRS keyword in the PROPS section - alternative three
point end-point scaling indicator. & SCALECRS & & \\
2-21 to 2-35 & & Undefined. & F & & \\
2-36 & Data & Miscible displacement has been activated in the model
indicator, set to T or F. See the Error: Reference source not found
keyword in the RUNSPEC section for further information. & & & \\
2-37 to 2-55 & & Undefined. & F & & \\
2-56 & Data & The scale water capillary pressure at maximum water
saturation (that is minimum pressure) has been activated in the model,
set to T or F. & & & \\
2-57 & Data & The scale water capillary pressure at maximum water
saturation (that is minimum pressure) has been activated in the model,
set to T or F. & & & \\
2-58 to 2-127 & & Undefined. & 0 & & \\
2-128 & Data & Coal Bed Methane option has been activated in the
compositional model indicator, set to T or F. Note that currently OPM
Flow does not support the Coal Bed Methane option and therefore this
item should be set to F. & F & & \\
& DOUBHEAD & Global static data keyword for the INIT Data file that
defines the double precision REAL variables for this keyword. NDHEAD is
the number of entries associated with this keyword, and is dependent on
the features being used in the model. & Required & & \\
1-1 & Format & DOUBHEAD & NDHEAD & DOUB & \\
2-1 to 2-185 & & Undefined. & 0.0 & & \\
Notes: & & & & & \\
\end{longtable}

\emph{Table A.14: INIT Data -- Static Data Header Keywords}

\subsubsection{Example: INIT Data -- Static Header
Keywords}\label{example-init-data-static-header-keywords}

The following example shows a typical formatted INIT Data header set of
keywords.

\begin{verbatim}
'INTEHEAD'         249 'INTE'
 -955283513      200400           2       -2345       -2345       -2345
      -2345       -2345          10          10           1         100
         10       -2345           7       -2345           0           1
          0           2           2           0           0           0
        110         108         109           3          97          93
      -2345       -2345          19          38          53       -2345
         97          93         146           5           0           1
         15          24           8           5           2           4
          0           0           0           0           0           0
          0           0           0           0           0           0
          0           0           0           0          19          10
       1982           0           0           0           1           0
          0           0           0           0           1          10
          0           0          12           1          25           1
 ………………………….……………..
'LOGIHEAD'          79 'LOGI'
 T  F  F  T  F  F  F  F  T  F  F  F  F  F  F  F  T  F  T  F  F  F  F  F  F
 F  F  F  F  F  F  F  F  F  F  F  F  F  F  F  F  F  F  F  F  F  F  F  F  F
 F  F  F  F  F  F  F  F  F  F  F  F  F  F  F  F  F  F  F  F  F  F  F  F  F
 F  F  F  F
'DOUBHEAD'         185 'DOUB'
  0.00000000000000D+00   0.10000000000000D+01   0.36500000000000D+03
  0.10000000149012D+00   0.15000000596046D+00   0.30000000000000D+01
  0.30000001192093D+00   0.10000000149012D+00   0.10000000149012D+00
 -0.10000000200409D+21   0.10000000000000D+01   0.10000000000000D+01
  0.10000000000000D+01   0.00000000000000D+00   0.00000000000000D+00
  0.10000000000000D+01   0.10000000000000D-03
 ………………………….……………..
\end{verbatim}

End of Example

This series of keywords define various static global and LGR grid
property arrays for each grid cell in the model. The grid array keywords
use the same grid property keywords described in the GRID section of the
manual. For example, the DEPTH, PORO, PERMX, PERMY, PERMZ, NTG, TOPS
arrays etc. In addition, the keywords also include the calculated pore
volumes (PORV) and transmissibility arrays (TRANX, TRANY, and TRANZ)
after processing the GRID and EDIT sections. Note the first property
should always be the PORV array for both global and LGR grid types.

This format is also used to define the global grid REGION, NNC, and
ENPOINTS section arrays, with the keyword name corresponding to the
array name. For example, for the REGION keywords the same keywords
outlined in the REGION section of the manual; that is EQLNUM, FIPNUM,
PVTNUM, ROCKNUM, and SATNUM are utilized as the keywords in this
section.

The structure for this type of data and the keywords are defined in
Table A.15.

\begin{longtable}[]{@{}
  >{\raggedright\arraybackslash}p{(\linewidth - 10\tabcolsep) * \real{0.1667}}
  >{\raggedright\arraybackslash}p{(\linewidth - 10\tabcolsep) * \real{0.1667}}
  >{\raggedright\arraybackslash}p{(\linewidth - 10\tabcolsep) * \real{0.1667}}
  >{\raggedright\arraybackslash}p{(\linewidth - 10\tabcolsep) * \real{0.1667}}
  >{\raggedright\arraybackslash}p{(\linewidth - 10\tabcolsep) * \real{0.1667}}
  >{\raggedright\arraybackslash}p{(\linewidth - 10\tabcolsep) * \real{0.1667}}@{}}
\toprule\noalign{}
\begin{minipage}[b]{\linewidth}\raggedright
No.
\end{minipage} & \begin{minipage}[b]{\linewidth}\raggedright
keyword Name
\end{minipage} & \begin{minipage}[b]{\linewidth}\raggedright
INIT Data -- Static Grid Array Data keywords Table A.15
\end{minipage} & \begin{minipage}[b]{\linewidth}\raggedright
Status or Value
\end{minipage} & \begin{minipage}[b]{\linewidth}\raggedright
\end{minipage} & \begin{minipage}[b]{\linewidth}\raggedright
\end{minipage} \\
\midrule\noalign{}
\endhead
\bottomrule\noalign{}
\endlastfoot
Keyword & No.~of Entries & Data Type & & & \\
& PROPERTY & The PROPERTY keyword defines the start of a grid property
array section of the INIT file. PROPERTY is the keyword name of the
array being written out, as outlined in the GRID section. For example,
DEPTH, PORO, PERMX, PERMY, PERMZ, NTG, TOPS keywords and the calculated
PORV, TRANX, TRANY, and TRANZ arrays. The PROPERTY keyword (items 1-1
and 2-1) is repeated for each grid property. The PORV property must be
the first data set for this keyword type and is written out in full for
each cell in the model (that is NXYZ = NX x NY x NZ values). Note the
PORV array values for inactive cells are set to a constant value of
zero. Subsequent grid property arrays can be in any order and only the
active cells are written out (NACTIV). & Required & & \\
1-1 & Format & PROPERTY & NXYZ or NACTIV & REAL & \\
2-1 & Data & PROPERTY data set. & & & \\
LGR Property Data for Each LGR Keyword and Data & & & & & \\
Table and Constant Keywords (TABDIMS, TAB, and CON) & & & & & \\
& REGION & The REGION keyword defines the start of the grid region array
section. REGION is the keyword name of the region array being written
out, as outlined in the REGIONS section. For example, EQLNUM, FIPNUM,
PVTNUM, ROCKNUM, SATNUM, etc. For analytical aquifers in the model, the
AQUIFERA variable name is used to indicate which cells are connected to
a given aquifer. If a cell is connected to an aquifer number N, then the
cell value will be set to 2(N-1) in the array. The REGION keyword (items
1-1 and 2-1) is repeated for each global region array and the arrays can
be written out in any order. & Required & & \\
1-1 & Format & REGION & NACTIV & INTE & \\
2-1 & Data & REGION data set. & & & \\
LGR Region Keywords and Data for Each Grid & & & & & \\
& ENDPOINT & The ENDPOINT keywords define the start of a grid's
saturation and end-point data arrays. ENDPOINT is the keyword name of
the saturation (SWL, SWATINIT, etc.) array or end-point array (PCW,
KRORW, etc.) being written out, as described in the GRID and PROPS
sections. The type of data for the keyword should be set to REAL for
arrays with real values (SWL, SWATINIT, etc) or INTE for integer arrays
(ENDNUM, etc). ENDPOINT arrays can be in any order and only the active
cells are written out (NACTIV). This keyword (items 1-1 and 2-1) is
repeated for each saturation and end-point array being written out. The
data written out is dependent on the various options declared in the
RUNSPEC section as well as the data declared in the GRID and PROPS
sections. & Required & & \\
1-1 & Keyword & ENDPOINTS & NACTIV & INTE or REAL & \\
2-1 & Data & ENDPOINT data set. & & & \\
LGR Saturation End-Point Keywords and Data for Each LGR & & & & & \\
& NNC & The NNC keyword defines the start of a grid's non-neighbor
connection array data. NNC is the keyword name of the non-neighbor
connection array being written out, and should be set to one of the
following keywords: The following two additional NNC keywords are
required for when the data for an LGR is being defined: The number of
data values for this keyword data set should set to the number of
non-neighbor connection to be read or written out (NUMNNC) for the
TRANNC, DIFFNNC, HEATNCC, TRANCC, and TRANGL arrays. The NNC keyword
(items 1-1 and 2-1) are repeated for each non-neighbor connection array
being written out. & Required & & \\
1-1 & Format & NNC & NUMNCC & REAL & \\
2-1 & Data & NNC data set. & & & \\
LGR Non-Neighbor Connection Keywords for Each LGR & & & & & \\
Notes: & & & & & \\
\end{longtable}

\emph{Table A.15: INIT Data -- Static Grid Array Data Keywords}

\subsubsection{Example: INIT Data -- Static Grid Array Data
Keywords}\label{example-init-data-static-grid-array-data-keywords}

The following example shows a typical formatted global grid property
data keyword.

\begin{verbatim}
 'PORV    '         100 'REAL'
   0.10686456E+07   0.10686456E+07   0.10686456E+07   0.10686456E+07
   0.10686456E+07   0.10686456E+07   0.10686456E+07   0.10686456E+07
  ………………………….……………..
   0.10686456E+07   0.10686456E+07   0.10686456E+07   0.10686456E+07
   0.10686456E+07   0.10686456E+07   0.10686456E+07   0.10686456E+07
 'DX      '         100 'REAL'
   0.10000000E+04   0.10000000E+04   0.10000000E+04   0.10000000E+04
   0.10000000E+04   0.10000000E+04   0.10000000E+04   0.10000000E+04
  ………………………….……………..
   0.10000000E+04   0.10000000E+04   0.10000000E+04   0.10000000E+04
   0.10000000E+04   0.10000000E+04   0.10000000E+04   0.10000000E+04
 'DY      '         100 'REAL'
   0.10000000E+04   0.10000000E+04   0.10000000E+04   0.10000000E+04
   0.10000000E+04   0.10000000E+04   0.10000000E+04   0.10000000E+04
  ………………………….……………..
   0.10000000E+04   0.10000000E+04   0.10000000E+04   0.10000000E+04
   0.10000000E+04   0.10000000E+04   0.10000000E+04   0.10000000E+04
 'DZ      '         100 'REAL'
   0.20000000E+02   0.20000000E+02   0.20000000E+02   0.20000000E+02
   0.20000000E+02   0.20000000E+02   0.20000000E+02   0.20000000E+02
   ………………………….……………..
   0.20000000E+02   0.20000000E+02   0.20000000E+02   0.20000000E+02
   0.20000000E+02   0.20000000E+02   0.20000000E+02   0.20000000E+02
 'PERMX   '         100 'REAL'
   0.50000000E+03   0.50000000E+03   0.50000000E+03   0.50000000E+03
   ………………………….……………..
   0.50000000E+03   0.50000000E+03   0.50000000E+03   0.50000000E+03
   0.50000000E+03   0.50000000E+03   0.50000000E+03   0.50000000E+03
………………………….……………..
TABULAR DATA
………………………….……………..
 'PVTNUM  '         100 'INTE'
           1           1           1           1           1           1
           1           1           1           1           1           1
………………………….……………..
           1           1           1           1           1           1
 'SATNUM  '         100 'INTE'
           1           1           1           1           1           2
           1           1           1           2           2           2
………………………….……………..
 'EQLNUM  '         100 'INTE'
           1           1           1           1           1           1
           1           1           1           1           1           1
 ………………………….……………..
\end{verbatim}

End of Example

The file structure for this series of keywords consists of two sets of
keywords, the first set is repeated each time an LGR data set is written
out and the second set, which is the same as global static header
keywords in section INIT Data -- Static Data Header Keywords, is only
written once. The first set of keywords defines an LGR's properties, LGR
Name for example, and always proceeds an LGR data set for a given LGR;
that is this series of keywords are repeated for each LGR data set in
the model. This is then followed by the second set of static data header
keywords, or the actual array or table keywords. Termination of the LGR
data set occurs after the last LGR data set has been written out using
the LGRSGONE keyword, which switches input back to the global grid
properties. Table A.16 outlines the structure of the INIT Data - LGR
Static Header Keywords and their affiliated data.

Note that currently OPM Flow does not support LGR grids and therefore
this set of keywords cannot be used by OPM Flow.

\begin{longtable}[]{@{}
  >{\raggedright\arraybackslash}p{(\linewidth - 10\tabcolsep) * \real{0.1667}}
  >{\raggedright\arraybackslash}p{(\linewidth - 10\tabcolsep) * \real{0.1667}}
  >{\raggedright\arraybackslash}p{(\linewidth - 10\tabcolsep) * \real{0.1667}}
  >{\raggedright\arraybackslash}p{(\linewidth - 10\tabcolsep) * \real{0.1667}}
  >{\raggedright\arraybackslash}p{(\linewidth - 10\tabcolsep) * \real{0.1667}}
  >{\raggedright\arraybackslash}p{(\linewidth - 10\tabcolsep) * \real{0.1667}}@{}}
\toprule\noalign{}
\begin{minipage}[b]{\linewidth}\raggedright
No.
\end{minipage} & \begin{minipage}[b]{\linewidth}\raggedright
Keyword Name
\end{minipage} & \begin{minipage}[b]{\linewidth}\raggedright
INIT Data - LGR Static Header Keywords Table A.16 (Repeated Each Time an
LGR Each Data Set is Written Out)
\end{minipage} & \begin{minipage}[b]{\linewidth}\raggedright
Status or Value
\end{minipage} & \begin{minipage}[b]{\linewidth}\raggedright
\end{minipage} & \begin{minipage}[b]{\linewidth}\raggedright
\end{minipage} \\
\midrule\noalign{}
\endhead
\bottomrule\noalign{}
\endlastfoot
Keyword & No.~of Entries & Data Type & & & \\
& LGR & LGR static data header keyword for the INIT Data file that
defines the name of the LGR. & Always Required & & \\
1-1 & Keyword & LGR & 1 & CHAR & \\
2-1 & Data & Name of the LGR. & & & \\
& LGRHEADI & LGR that defines the integer variables for this keyword. &
Always Required & & \\
1-1 & Format & LGRHEADI & 45 & INTE & \\
2-1 to 2-45 & Data & Undefined. & & & \\
& LGRHEADQ & LGR grid keyword for the INIT Data file that defines the
logical variables (T for true and F for false) for this keyword. &
Always Required & & \\
1-1 & Format & LGRHEADQ & 5 & LOGI & \\
2-1 to 2-5 & Data & Undefined. & & & \\
& LGRHEADD & LGR grid keyword for the INIT Data file that defines the
double precision REAL variables for this keyword. & Always Required &
& \\
1-1 & Format & LGRHEADD & 5 & DOUB & 5 \\
2-1 to 2-5 & Data & Undefined. & & & \\
& LGRSGONE & The LGRSGONE keyword marks the end of the LGR section.
Number of entries is always set to zero. There is no data set associated
with this keyword. & Always Required & & \\
1-1 & Format & LGRSGONE & 0 & MESG & \\
& INTEHEAD & The keywords are the same as the INIT Data Static Data
Header keywords but the data pertains to the current LGR. See Table A.14
for a detailed description of this keyword type. & Required Only Once &
& \\
& LOGIHEAD & See Table A.14 for a detailed description of this keyword
type. & Required Only Once & & \\
& DOUBHEAD & See Table A.14 for a detailed description of this keyword
type. & Required Only Once & & \\
Notes: & & & & & \\
\end{longtable}

\emph{Table A.16: INIT Data - LGR Static Data Header Keywords}

This series of keywords are repeated for each LGR in the model followed
by the pertinent data and in the order outlined in Table A.13 and with
the last LGR data set terminated by the LGRSGONE keyword to mark the end
of the LGR input section.

\begin{tcolorbox}[enhanced jigsaw, title=\textcolor{quarto-callout-note-color}{\faInfo}\hspace{0.5em}{Note}, coltitle=black, breakable, toprule=.15mm, colbacktitle=quarto-callout-note-color!10!white, leftrule=.75mm, opacitybacktitle=0.6, colback=white, bottomtitle=1mm, toptitle=1mm, titlerule=0mm, left=2mm, colframe=quarto-callout-note-color-frame, arc=.35mm, rightrule=.15mm, bottomrule=.15mm, opacityback=0]

The LGR, LGRHEADI, LGRHEADQ, and LGRHEADD keywords always proceeds and
LGR data sets and the LGRSGONE keyword terminates the end of the LGR
data section, which causes the data to revert back to the global grid
data.

\end{tcolorbox}

\subsubsection{Example: INIT Data - LGR Static Header
Keywords}\label{example-init-data---lgr-static-header-keywords}

The following example shows a typical formatted INIT Data - LGR Grid
Header Keyword data set for a single LGR grid named LGR-1. As this is
the first time LGR data is being written out both sets of LGR header
keywords are written out followed by the property array data.

\begin{verbatim}
 'LGR     '           1 'CHAR'
 'LGR-1   '
 'LGRHEADI'          45 'INTE'
           1         100       -2345       -2345       -2345       -2345
       -2345       -2345       -2345       -2345       -2345       -2345
       -2345       -2345       -2345       -2345       -2345       -2345
       -2345       -2345       -2345       -2345       -2345       -2345
       -2345       -2345       -2345       -2345       -2345       -2345
       -2345       -2345       -2345       -2345       -2345       -2345
       -2345       -2345       -2345       -2345       -2345       -2345
       -2345       -2345       -2345
 'LGRHEADQ'           5 'LOGI'
  F  F  F  F  F
 'LGRHEADD'           5 'DOUB'
   0.00000000000000D+00  -0.10000000200409D+21  -0.10000000200409D+21
  -0.10000000200409D+21  -0.10000000200409D+21
………………………….……………..

 'INTEHEAD'         249 'INTE'
  -955283513      200400           2       -2345       -2345       -2345
       -2345       -2345           2           2           2           8
………………………….……………..
           1           6           1          10           1          10
           1           1           1           0          30           3
          18          10           9
 'LOGIHEAD'          79 'LOGI'
  T  F  F  T  F  F  F  F  T  F  F  F  F  F  F  F  T  F  T  F  F  F  F  F  F
  F  F  F  F  F  F  F  F  F  F  F  F  F  F  F  F  F  F  F  F  F  F  F  F  F
  F  F  F  F  F  F  F  F  F  F  F  F  F  F  F  F  F  F  F  F  F  F  F  F  F
  F  F  F  F
 'DOUBHEAD'         185 'DOUB'
   0.00000000000000D+00   0.10000000000000D+01   0.36500000000000D+03
   0.10000000149012D+00   0.15000000596046D+00   0.30000000000000D+01
………………………….……………..
   0.10000000000000D+01   0.00000000000000D+00   0.00000000000000D+00
   0.10000000000000D+01   0.10000000000000D-03
………………………….……………..
 'PORV    '           8 'REAL'
   0.13358070E+06   0.13358070E+06   0.13358070E+06   0.13358070E+06
   0.13358070E+06   0.13358070E+06   0.13358070E+06   0.13358070E+06
 'DX      '           8 'REAL'
   0.50000000E+03   0.50000000E+03   0.50000000E+03   0.50000000E+03
   0.50000000E+03   0.50000000E+03   0.50000000E+03   0.50000000E+03
 'DY      '           8 'REAL'
   0.50000000E+03   0.50000000E+03   0.50000000E+03   0.50000000E+03
   0.50000000E+03   0.50000000E+03   0.50000000E+03   0.50000000E+03
 'DZ      '           8 'REAL'
   0.10000000E+02   0.10000000E+02   0.10000000E+02   0.10000000E+02
   0.10000000E+02   0.10000000E+02   0.10000000E+02   0.10000000E+02
 'PERMX   '           8 'REAL'
   0.50000000E+03   0.50000000E+03   0.50000000E+03   0.50000000E+03
   0.50000000E+03   0.50000000E+03   0.50000000E+03   0.50000000E+03
………………………….……………..
\end{verbatim}

End of Example

The file structure for this set of keywords is identical to the Static
Grid Array Data Keywords described in Table A.15, except that LGR header
keywords prefix the actual data keywords. Both the LGR header and data
keywords are repeated for each LGR in the model, and the last LGR data
set is terminated by the LGRSGONE keyword to mark the end of the LGR
input section. The keyword description for this type of data set is
outlined in Table A.17.

Note that currently OPM Flow does not support LGR grids and therefore
this series of keywords cannot be used by OPM Flow.

\begin{longtable}[]{@{}
  >{\raggedright\arraybackslash}p{(\linewidth - 6\tabcolsep) * \real{0.2500}}
  >{\raggedright\arraybackslash}p{(\linewidth - 6\tabcolsep) * \real{0.2500}}
  >{\raggedright\arraybackslash}p{(\linewidth - 6\tabcolsep) * \real{0.2500}}
  >{\raggedright\arraybackslash}p{(\linewidth - 6\tabcolsep) * \real{0.2500}}@{}}
\toprule\noalign{}
\begin{minipage}[b]{\linewidth}\raggedright
No.
\end{minipage} & \begin{minipage}[b]{\linewidth}\raggedright
Keyword Name
\end{minipage} & \begin{minipage}[b]{\linewidth}\raggedright
INIT Data - LGR Grid Array Data Keywords Table A.17 (Repeated for Each
LGR in the Model)
\end{minipage} & \begin{minipage}[b]{\linewidth}\raggedright
Status or Value
\end{minipage} \\
\midrule\noalign{}
\endhead
\bottomrule\noalign{}
\endlastfoot
1-1 & LGR & LGR Header Keywords (see Table A.16). & Required \\
LGRHEADI & & & \\
LGRHEADQ & & & \\
LGRHEADD & & & \\
1-2 & INTEHEAD & LGR Header Keywords (see Table A.16). & Required Only
Once \\
LOGIHEAD & & & \\
DOUBHEAD & & & \\
1-3 & PROPERTY REGION NNC or ENPOINIT & LGR array data set for the
current LGR -- same format as the INIT Data - Grid Array Data Keywords
(see Table A.15). & Required \\
1-4 & LGRSGONE & LGR Header Section termination keyword (see Table
A.16). & Required \\
Notes: & & & \\
\end{longtable}

\emph{Table A.17: INIT Data - LGR Grid Array Data Keywords}

This keyword data set is repeated for each LGR in the model with the
last LGR data set terminated by the LGRSGONE keyword to mark the end of
the LGR data section.

\subsubsection{Example: INIT Data - LGR Static Grid Array Data
Keywords}\label{example-init-data---lgr-static-grid-array-data-keywords}

The following example shows a typical formatted INIT Data - LGR Static
Grid Array Data Keywords for the region data set with two LGR grids
named LGR-1 and LGR-2.

\begin{verbatim}
'LGR     '           1 'CHAR'
'LGR-1   '
'LGRHEADI'          45 'INTE'
          1         100       -2345       -2345       -2345       -2345
      -2345       -2345       -2345       -2345       -2345       -2345
      -2345       -2345       -2345       -2345       -2345       -2345
      -2345       -2345       -2345       -2345       -2345       -2345
      -2345       -2345       -2345       -2345       -2345       -2345
      -2345       -2345       -2345       -2345       -2345       -2345
      -2345       -2345       -2345       -2345       -2345       -2345
      -2345       -2345       -2345
'LGRHEADQ'           5 'LOGI'
 F  F  F  F  F
'LGRHEADD'           5 'DOUB'
  0.00000000000000D+00  -0.10000000200409D+21  -0.10000000200409D+21
 -0.10000000200409D+21  -0.10000000200409D+21
'PVTNUM  '           8 'INTE'
          1           1           1           1           1           1
          1           1
'SATNUM  '           8 'INTE'
          1           1           1           1           2           2
          2           2
'EQLNUM  '           8 'INTE'
          1           1           1           1           1           1
          1           1
'FIPNUM  '           8 'INTE'
          1           1           1           1           1           1
          1           1
'LGR     '           1 'CHAR'
'LGR-2   '
'LGRHEADI'          45 'INTE'
          2         100       -2345       -2345       -2345       -2345
      -2345       -2345       -2345       -2345       -2345       -2345
      -2345       -2345       -2345       -2345       -2345       -2345
      -2345       -2345       -2345       -2345       -2345       -2345
      -2345       -2345       -2345       -2345       -2345       -2345
      -2345       -2345       -2345       -2345       -2345       -2345
      -2345       -2345       -2345       -2345       -2345       -2345
      -2345       -2345       -2345
'LGRHEADQ'           5 'LOGI'
 F  F  F  F  F
'LGRHEADD'           5 'DOUB'
  0.00000000000000D+00  -0.10000000200409D+21  -0.10000000200409D+21
 -0.10000000200409D+21  -0.10000000200409D+21
'PVTNUM  '           8 'INTE'
          1           1           1           1           1           1
          1           1
'SATNUM  '           8 'INTE'
          1           2           1           2           2           1
          2           1
'EQLNUM  '           8 'INTE'
          1           1           1           1           1           1
          1           1
'FIPNUM  '           8 'INTE'
          1           1           1           1           1           1
          1           1
'LGRSGONE'           0 'MESS'
\end{verbatim}

End of Example

The INIT Data -- Static Tabular Data Header Keyword precedes the TAB
keyword and consists of just one keyword followed by a series of integer
values that define the dimensions of the data in the TAB keyword.
TABDIMS is then followed by the TAB keyword that contains the actual
tabular data in one continuous array.

The structure for this keyword is defined in Table A.18.

\begin{longtable}[]{@{}
  >{\raggedright\arraybackslash}p{(\linewidth - 10\tabcolsep) * \real{0.1667}}
  >{\raggedright\arraybackslash}p{(\linewidth - 10\tabcolsep) * \real{0.1667}}
  >{\raggedright\arraybackslash}p{(\linewidth - 10\tabcolsep) * \real{0.1667}}
  >{\raggedright\arraybackslash}p{(\linewidth - 10\tabcolsep) * \real{0.1667}}
  >{\raggedright\arraybackslash}p{(\linewidth - 10\tabcolsep) * \real{0.1667}}
  >{\raggedright\arraybackslash}p{(\linewidth - 10\tabcolsep) * \real{0.1667}}@{}}
\toprule\noalign{}
\begin{minipage}[b]{\linewidth}\raggedright
No.
\end{minipage} & \begin{minipage}[b]{\linewidth}\raggedright
keyword Name
\end{minipage} & \begin{minipage}[b]{\linewidth}\raggedright
INIT Data -- Static Tabular Data Header Keyword Table A.18
\end{minipage} & \begin{minipage}[b]{\linewidth}\raggedright
Status or Value
\end{minipage} & \begin{minipage}[b]{\linewidth}\raggedright
\end{minipage} & \begin{minipage}[b]{\linewidth}\raggedright
\end{minipage} \\
\midrule\noalign{}
\endhead
\bottomrule\noalign{}
\endlastfoot
Keyword & No.~of Entries & Data Type & & & \\
& TABDIMS & Tabular header keyword for the INIT Data file that defines
the start of the tabular data and contains the base addresses and
dimension of subsequent tabular array keyword TAB. Missing values should
be defaulted to one. & Required & & \\
1-1 & Format & TABDIMS & 100 & INTE & \\
2-1 & Data & Total size of TAB data array, value should be the same as
that on the TAB keyword. & NTABDA & & \\
2-2 & Data & Location of the first entry in the TAB array for the rock
property table. & IBROCK & & \\
2-3 & Data & Maximum number of rock property tables. & NTROCK & & \\
2-4 & Data & Location of the first entry in the TAB array of rock
compaction data. & IBROCC & & \\
2-5 & Data & Maximum number of pressure nodes in rock compaction data
table. & NPROCC & & \\
2-6 & Data & Maximum number of rock compaction data tables. & NTROCC &
& \\
2-7 & Data & Location of the first entry in the TAB array for the oil
PVT property tables. & IBPVTO & & \\
2-8 & Data & Location of the first entry in the TAB array for the oil
pressure nodes. & JBPVTO & & \\
2-9 & Data & Maximum number of composition nodes in oil PVT tables. &
NRPVTO & & \\
2-10 & Data & Maximum number of pressure nodes in oil PVT tables. &
NPPVTO & & \\
2-11 & Data & Maximum number of oil PVT property tables. & NTPVTO & & \\
2-12 & Data & Location of the first entry in the TAB array for water PVT
property tables. & IBPVTW & & \\
2-13 & Data & Maximum number of water property tables. & NTPVTW & & \\
2-14 & Data & Location of the first entry in the TAB array for gas PVT
property tables. & IBPVTG & & \\
2-15 & Data & Location of the first entry in the TAB array for the gas
pressure nodes. & JBPVTG & & \\
2-16 & Data & Maximum number of composition nodes in the gas PVT tables.
& NRPVTG & & \\
2-17 & Data & Maximum number of pressure nodes in the gas PVT tables. &
NPPVTG & & \\
2-18 & Data & Maximum number of gas PVT property tables. & NTPVTG & & \\
2-19 & Data & Location of the first entry in the TAB array for surface
density tables. & IBDENS & & \\
2-20 & Data & Maximum number of surface densities tables. & NTDENS &
& \\
2-21 & Data & Location of the first entry in the TAB array for water
saturation tables. & IBSWFN & & \\
2-22 & Data & Maximum number of saturation nodes in water saturation
tables. & NSSWFN & & \\
2-23 & Data & Maximum number of water saturation tables. & NTSWFN & & \\
2-24 & Data & Location of the first entry in the TAB array for gas
saturation tables. & IBSGFN & & \\
2-25 & Data & Maximum number of saturation nodes in gas saturation
tables. & NSSGFN & & \\
2-26 & Data & Maximum number of gas saturation tables. & NTSGFN & & \\
2-27 & Data & Location of the first entry in the TAB array for oil
saturation tables. & IBSOFN & & \\
2-28 & Data & Location of the first entry in the TAB array for array of
connate water saturations. & IBSWCO & & \\
2-29 & Data & Maximum number of saturation nodes in the oil saturation
tables. & NSSOFN & & \\
2-30 & Data & Maximum number of oil saturation tables. & NTSOFN & & \\
2-31 to 40 & & Undefined. & 1 & & \\
2-41 & Data & Location of the first entry in the TAB array to the VE
table data. & IBVETB & & \\
2-42 & Data & Maximum number of items in a VE table. & NSVETB & & \\
2-43 & Data & Maximum number of VE tables. & NTVETB & & \\
2-44 & Data & Location of the first entry in the TAB array to the
threshold pressure array. & IBTHPR & & \\
2.45 & Data & Location of the first entry in the TAB array to SCALELIM
data & IBSLIM & & \\
2.46 & Data & Maximum number of end-point versus depth nodes. & NSENDP &
& \\
2.47 & Data & Maximum number of end-point versus depth tables. & NTENDP
& & \\
2.48 & Data & Location of the first entry in the TAB array to the
reference temperature RTEMP used in PVZG keyword in the PROPD section. &
IBRTEM & & \\
2.49 & Data & Location of the first entry in the TAB array to the value
specified by the TOLCRIT keyword. Note that this is a single value
applied to all the saturation tables. & IBCTOL & & \\
2.50 & & Undefined. & 1 & & \\
2-51 & Data & Location of the first entry in the TAB array to LANGMUIR
table. & IBLANG & & \\
2-52 & Data & Maximum number of columns in LANGMUIR tables. & NCLANG &
& \\
2-53 & Data & Maximum number of rows in the LANGMUIR tables. & NSLANG &
& \\
2-54 & Data & Maximum number of LANGMUIR tables & NTLANG & & \\
2-55 & Data & Location of the first entry in the TAB array to a LANGSOLV
table. & IBLNG2 & & \\
2-56 & Data & Location of the first entry in the TAB array to a COALPP
table. & IBCADP & & \\
2-57 & Data & Location of the first entry in the TAB array to a COALADS
table. & IBCADS & & \\
2-58 & Data & Location of the first entry in the TAB array to a ROCKPAMA
table. & IBROCP & & \\
2-59 & Data & Maximum number of ROCKPAMA tables. & NTRPMA & & \\
2-60 to 2-100 & & Undefined. & 1 & & \\
Notes: & & & & & \\
\end{longtable}

\emph{Table A.18: INIT Data -- Static Tabular Data Header Keyword}

\subsubsection{Example: INIT Data -- Static Tabular Data Header
Keyword}\label{example-init-data-static-tabular-data-header-keyword}

The following example shows a typical formatted INIT Data -- Static
Tabular Data Header Keyword and the affiliated data.

\begin{verbatim}
'TABDIMS '         100 'INTE'
      77006           1           2           1          60           1
         17       36017          60          60           2       36799
          2       37139       73139          60          60           2
          5           2       36809          33           2       73259
         33           2       36137       36797          66           2
      73589           5       73644          20           5       74044
         20           5           0           0       74444          25
          2       74494       75835          20           1       77003
          0           0           0           0           0           0
          0           0           0           0           0           0
          0           0           0           0           0           0
          0           0           0           0           0           0
          0           0           0           0           0           0
          0           0           0           0           0           0
          0           0           0           0           0           0
          0           0           0           0           0           0
          0           0           0           0
\end{verbatim}

End of Example

The INIT Data -- Static Tabular and Constant Data Array Keyword section
contains the actual tabular data in one continuous array, followed by
the constant data array. The tabular array data is parse according to
the phases present in the model and the location and dimension data
given by the TABDIMS keyword.

The structure for this set of keywords is defined in Table A.19.

\begin{longtable}[]{@{}
  >{\raggedright\arraybackslash}p{(\linewidth - 10\tabcolsep) * \real{0.1667}}
  >{\raggedright\arraybackslash}p{(\linewidth - 10\tabcolsep) * \real{0.1667}}
  >{\raggedright\arraybackslash}p{(\linewidth - 10\tabcolsep) * \real{0.1667}}
  >{\raggedright\arraybackslash}p{(\linewidth - 10\tabcolsep) * \real{0.1667}}
  >{\raggedright\arraybackslash}p{(\linewidth - 10\tabcolsep) * \real{0.1667}}
  >{\raggedright\arraybackslash}p{(\linewidth - 10\tabcolsep) * \real{0.1667}}@{}}
\toprule\noalign{}
\begin{minipage}[b]{\linewidth}\raggedright
No.
\end{minipage} & \begin{minipage}[b]{\linewidth}\raggedright
keyword or Variable Name
\end{minipage} & \begin{minipage}[b]{\linewidth}\raggedright
INIT Data -- Static Tabular and Constant Data Array Keyword Table A.19
\end{minipage} & \begin{minipage}[b]{\linewidth}\raggedright
Status or Value
\end{minipage} & \begin{minipage}[b]{\linewidth}\raggedright
\end{minipage} & \begin{minipage}[b]{\linewidth}\raggedright
\end{minipage} \\
\midrule\noalign{}
\endhead
\bottomrule\noalign{}
\endlastfoot
Keyword & No.~of Entries & Data Type & & & \\
& TAB & Tabular data array keyword for the INIT Data file that defines
the start of the tabular data array. The total number of entries is
given by the NTABDA and the value should be the same as declared on the
TABDIMS keyword. & Required & & \\
1-1 & Format & TABDIMS & NTABDA & DOUB & \\
2-1 & Data & Tabular data array. The data should be written out in
double precision and missing data should be set to less than or equal to
-10.0 x 1020 or greater than or equal to 10.0 x 1020. & & & \\
& CON & Constant data array that contains the constant dead oil, dead
vaporized oil, and saturation pressure PVT entered via the
RSCONST/RSCONSTT and RVCONST/RVCONSTT keywords in the PROPS section. &
Required & & \\
1-1 & Format & CON & NPVT x 3 & DOUB & \\
2-1 & Data & The data consists of three values for each PVT data set
(NPVT), RS, RV, and the saturation pressure of the active phase in the
model (bubble-point or dew point). & & & \\
Notes: & & & & & \\
\end{longtable}

\emph{Table A.19: INIT Data -- Static Tabular and Constant Data Array
Keywords}

The phases present in the model are given on the INTEHEAD keyword (Table
A.14) in location 2-15 where the type of phases in the model are set via
an integer: 1 for oil, 2 for water, 3 for oil-water, 4 for gas, 5 for
oil-gas, 6 for gas-water, and 7 for oil-water-gas. This data is used to
access the phase dependent data (PVT, saturation functions tables, etc.)
combine with the location (or base address) and the dimension data given
on the TABDIMS keyword in Table A.18.

Note that data given in Table A.18 refers to the maximum values and
therefore missing data should be set to the default value.

\subsubsection{Example: INIT Data -- Static Tabular and Constant Data
Array
Keywords}\label{example-init-data-static-tabular-and-constant-data-array-keywords}

The following example shows a typical formatted data set for this type
set of keywords and the affiliated data. Here the CON array has two PVT
tables but the values are all set to zero indicating that dead oil and
vaporized dead oil are not active phases in the model.

\begin{verbatim}
'TAB     '       73576 'DOUB'
   0.49700000000000D-05   0.24800000000000D-05   0.00000000000000D+00
  -0.20000000000000D+21  -0.20000000000000D+21  -0.20000000000000D+21
  -0.20000000000000D+21  -0.20000000000000D+21  -0.20000000000000D+21
 ………………………….……………..
  -0.20000000000000D+21  -0.20000000000000D+21  -0.20000000000000D+21
   0.52100000000000D-05   0.26100000000000D-05   0.00000000000000D+00
  -0.20000000000000D+21  -0.20000000000000D+21  -0.20000000000000D+21
'CON     '           6 'DOUB'
   0.00000000000000D+00   0.00000000000000D+00   0.00000000000000D+00
   0.00000000000000D+00   0.00000000000000D+00   0.00000000000000D+00
 ………………………….……………..
\end{verbatim}

End of Example

The LOG file for the commercial simulator contains an ASCII copy of the
output from the terminal. The file itself is not generated by the
commercial simulator but by an auxiliary program, ECLRUN, that is used
to call and execute all the commercial simulators.

OPM Flow does not generate a LOG file; however, the file can be
generated by using the Linux tee command, for example:

\begin{verbatim}
flow CASENAME | tee CASENAME.LOG
\end{verbatim}

Which will copy the terminal output to the CASENAME.LOG file.
Alternatively, one can use OPMRUN, the Graphical User Interface
(``GUI'') program for OPM Flow, that provides similar functionality to
the commercial simulator's auxiliary program.

The OPM Flow terminal output differs from that generated by the
commercial simulator. This section outlines the format of the OPM Flow
specific Terminal Output (LOG File) illustrated with examples.

Note that the OPM Flow terminal output is generally also sent in the
Print (PRT) file. In addition to the terminal output, the Print file
also includes:

\begin{itemize}
\tightlist
\item
  the name of the machine the simulation was run on,
\item
  the operating system,
\item
  the start time for the simulation,
\item
  a list of the command line arguments specified when running OPM Flow,
\item
  the property tree for the linear solver,
\item
  any output reports requested by the user (for example, using the
  RPTPROPS keyword in the PROPS section, the RPTSOL keyword in the
  SOLUTION section, the RPTSCHED keyword in the SCHEDULE section, and so
  forth), and finally
\item
  an error summary giving the total number of Warnings, Problems,
  Errors, and so forth reported during the simulation.
\end{itemize}

Also note that the contents of the LOG file will vary depending on the
model being simulated, the command line parameters, and the numerical
performance of the simulator. The main simulation loop for example may
contain additional messages other than those illustrated including: when
the simulator changes the control mode for a well, when a time step
fails to converge, and so forth.

At the start of the LOG file OPM Flow writes a header giving the
software version followed by information about the number of MPI
processes if applicable. The following example is taken from the LOG
file generated for the Norne model
(\href{https://github.com/OPM/opm-tests/blob/master/norne/NORNE_ATW2013.DATA}{NORNE\_ATW2013}):

\begin{verbatim}
**********************************************************************
*                                                                    *
*                        This is flow 2024.04                        *
*                                                                    *
* Flow is a simulator for fully implicit three-phase black-oil flow, *
*             including solvent and polymer capabilities.            *
*          For more information, see https://opm-project.org         *
*                                                                    *
**********************************************************************

Using 8 MPI processes with 2 OMP threads on each
\end{verbatim}

In this example OPM Flow version 2024.04 has been run using eight MPI
processes and each MPI process may itself use 2 OpenMP threads.

Next OPM Flow parses the input deck including any files referenced with
the INCLUDE keyword. As each keyword is parsed and it's data records are
read, the simulator reports the keyword name, the file it is located in,
and the corresponding line number. Any warning messages about
unsupported keywords or unsupported keyword items found in the deck are
then reported. See the following shortened example:

\begin{verbatim}
Reading deck file 'NORNE_ATW2013.data'
    0 Reading RUNSPEC  in NORNE_ATW2013.data line 9
    1 Reading DIMENS   in NORNE_ATW2013.data line 11
...
 1342 Reading DATES    in /.../norne/INCLUDE/BC0407_HIST01122006.SCH line 9565
#[34;1m
Warning: Unsupported keywords or keyword items:

  NSTACK: keyword not supported
  In file: NORNE_ATW2013.data, line 102
...
  RPTRST: invalid value '0' for item 1
  In file: /.../norne/INCLUDE/BC0407_HIST01122006.SCH, line 9245
  RPTRST(RPTRST): invalid option or unsupported integer control format#[0m
\end{verbatim}

In this example OPM Flow reports that the NSTACK keyword in the main
data file is not currently supported, and that an invalid option has
been specified for item 1 on the RPTRST keyword in an include file.

Note that OPM Flow outputs text formatting to the terminal to provide
color highlighting of warning messages etc. (for example,
``\#{[}34;1m''). This formatting may not be displayed in a text editor.

The input that was read in is now processed. First the grid structure
and the fault model is created, and in this example the tracers are
initialized:

\begin{verbatim}
Creating corner-point grid from keywords COORD, ZCORN and others

Initializing tracers from TRACER in NORNE_ATW2013.data line 268
...

Loading faults from FAULTS in /.../norne/INCLUDE/FAULT/FAULT_JUN_05.INC line 9

Applying MULTFLT in /.../norne/INCLUDE/FAULT/FAULTMULT_AUG-2006.INC line 8
...
\end{verbatim}

In the above example, the fault structure is read in from an include
file (FAULT\_JUN\_05.INC), and then the corresponding fault multipliers
are read from different include files (FAULTMULT\_AUG-2006.INC, etc.).

Next the dynamic information from the SCHEDULE section is processed to
prepare the report steps:

\begin{verbatim}
Processing dynamic information from
NORNE_ATW2013.data line 449
Initializing report step 0/247 at 1997-11-06 0 DAYS line 449
Processing keyword DRSDT at line 452
...
Reading from: /.../norne/INCLUDE/VFP/NEW_D2_GAS_0.00003.VFP line 62
Processing keyword VFPPROD at line 62
#[34;1m
Warning: VFPPROD table 5 has 7 non-monotonic points of BHP(THP)
In /.../norne/INCLUDE/VFP/NEW_D2_GAS_0.00003.VFP line 62
This may cause convergence issues due to switching between BHP and THP control.
#[0m
...
Complete report step 1 (8 DAYS) at 1997-11-14 (0 DAYS)

Initializing report step 2/247 at 1997-11-14 (0 DAYS) line 204
Processing keyword WCONHIST at line 208
Processing keyword WCONINJE at line 212
Complete report step 2 (17 DAYS) at 1997-12-01 (8 DAYS)
...
\end{verbatim}

In the above example where key values have been highlighted in yellow,
the simulation will start the first report step (out of a total of 247)
at the simulation start time of 6 November 1997; and will complete it
after advancing the simulation by 8 days to 14 November 1997.

Whilst, initializing the first report step, lift table (VFPPROD table 5)
is processed and OPM Flow reports that it has found 7 non-monotonic
points of BHP versus THP, and warns that this may cause convergence
issues. At this point the user may decide to adjust the table or
continue with the unmodified table.

The simulation will then continue with report step 2 of 247 starting on
14 November 1997; and will complete report step 2 after advancing the
simulation by a further 17 days to 1 December 1997 (elapsed time 25
days).

During report step 2 the WCONHIST and WCONINJE keywords have been
processed. These modify the well production and injection controls
effective from the start of the report step.

Note that OPM Flow only reports the initialization of the first five
report steps.

Note at this point the dynamic data is only being processed to set up
the report step; and the simulation has not yet started.

The grid is then processed after which the simulator reports the number
of active cells in the model along with their total pore volume. The
number of cells removed from the model due to the cut-off values applied
and the resulting reduction in total pore volume are also reported. For
example:

\begin{verbatim}
Processing grid
Total number of active cells: 44431 / total pore volume: 673248728 RM3
Removed 496 cells with a pore volume of 100450 RM3 (0.015 %) due to MINPV/MINPVV
\end{verbatim}

For parallel runs information about the load balancing between the
processes is then reported:

\begin{verbatim}
Load balancing distributes 44431 active cells on 8 processes as follows:
  rank   owned cells   overlap cells   total cells
--------------------------------------------------
     0          5821            3382          9203
     1          6104            1771          7875
...
--------------------------------------------------
   sum         44431           13568         57999


Load balancing distributed the wells as follows:
  well name            ranks with perforated cells
---------------------------------------------------
           B-1AH       4
            B-2H       1
...
\end{verbatim}

Finally, warning messages about unsupported keywords are reported, for
example:

\begin{verbatim}
#[34;1m
Warning: Unhandled summary keyword DATE/DAY
In /.../norne/INCLUDE/SUMMARY/summary.data line 9#[0m
...
\end{verbatim}

Having processed the input data the simulator reports diagnostic
information about the saturation functions including: the fluid system,
the saturation function family, and the number of saturation function
regions. For example:

\begin{verbatim}
===============Saturation Functions Diagnostics===============

System:  Black-oil system.
Relative permeability input format: Saturation Family I.
Number of saturation regions: 2
\end{verbatim}

In the above example a black-oil system has been specified, saturation
function family I has been used, and there are two saturation function
regions.

The simulator now enters the main simulation loop and reports the
progress of the simulation as it advances from the start date to each of
the report times specified in the input deck, for example:

\begin{verbatim}
================ Starting main simulation loop ===============


Report step  0/247 at day 0/3312, date = 06-Nov-1997
Using Newton nonlinear solver.

Starting time step 0, stepsize 1 days, at day 0/8, date = 06-Nov-1997
Restart file written for report step   0/247, date = 06-Nov-1997 00:00:00
 Newton its= 3, linearizations= 4 (0.1sec), linear its= 15 (0.1sec)

Starting time step 1, stepsize 3 days, at day 1/8, date = 07-Nov-1997
 Newton its= 2, linearizations= 3 (0.0sec), linear its= 21 (0.1sec)

Starting time step 2, stepsize 4 days, at day 4/8, date = 10-Nov-1997
 Newton its= 2, linearizations= 3 (0.0sec), linear its= 33 (0.1sec)

Report step  1/247 at day 8/3312, date = 14-Nov-1997
...
Report step 246/247 at day 3298/3312, date = 17-Nov-2006

Starting time step 0, stepsize 14 days, at day 3298/3312, date = 17-Nov-2006
 Newton its= 3, linearizations= 4 (0.0sec), linear its= 70 (0.2sec)

\end{verbatim}

In the example above, the simulation starts the first report step number
0 (out of a total of 247 report steps) on 6 November 1997 an elapsed
time of 0 days (of a total of 3312 days).

The simulator reports that it is using the Newton nonlinear solver.

A restart file is written at the start of the first report step number 0
(out of a total of 247 report steps) on 6 November 1997 at 00:00:00.

The simulator then attempts to advance the simulation by a first time
step of 1 day starting from 6 November 1997 an elapsed time of 0 days
(towards the next report time at an elapsed time of 8 days). The time
step is successfully completed after three Newton iterations, four
linearizations (taking 0.1s) and 15 linear iterations (taking 0.1s).

The simulator then attempts to advance the simulation by a second time
step of 3 days starting on 7 November 1997 an elapsed time of 1 day
(towards the next report time at an elapsed time of 8 days). The time
step is successfully completed after two Newton iterations, three
linearizations (taking 0.0s), and 21 linear iterations (taking 0.1s).
Note that since the previous time step successfully converged the time
step size has been increased by OPM Flow's time stepping algorithm.

The simulator then attempts to advance the simulation by a third time
step of 4 days starting on 10 November 1997 an elapsed time of 4 days
(towards the next report time at an elapsed time of 8 days). The time
step is successfully completed after two Newton iterations, three
linearizations (taking 0.0s), and 33 linear iterations (taking 0.1s).
Note the time step size has been constrained to advance the simulation
to the next report time.

The simulation has now reached the start of report step 1 (of 247) on 14
November 1997 an elapsed time of 8 days (out of a total of 3312 days).

The above example now skips the intermediate report steps and jumps to
the end of the simulation.

The simulation starts the final report step number 246 (of 247) on 17
November 2006 an elapsed time of 3298 days (out of a total of 3312
days).

The simulator then attempts to advance the simulation by a single time
step of 14 days starting on 17 November 2006 an elapsed time of 3298
days (towards the final report time at an elapsed time of 3312 days).
The time step is successfully completed after three Newton iterations,
four linearizations (taking 0.0s), and 70 linear iterations (taking
0.2s). The simulation ends on 1 December 2006.

The final section of the LOG file reports the numerical performance of
the simulation including the number of processors, the simulation time,
and the total numbers of linearizations, Newton iterations, and linear
iterations, for example:

\begin{verbatim}
================    End of simulation     ===============

Number of MPI processes:         8
Threads per MPI process:         2
Setup time:                      8.40 s
  Deck input:                    5.18 s
Number of timesteps:           369
Simulation time:               240.47 s
  Assembly time:                22.44 s (Wasted: 0.0 s; 0.0%)
    Well assembly:               2.74 s (Wasted: 0.0 s; 0.0%)
  Linear solve time:            83.22 s (Wasted: 0.0 s; 0.0%)
    Linear setup:                8.16 s (Wasted: 0.0 s; 0.0%)
  Props/update time:            21.69 s (Wasted: 0.0 s; 0.0%)
  Pre/post step:                35.75 s (Wasted: 0.0 s; 0.0%)
  Output write time:            69.24 s
Overall Linearizations:       1903      (Wasted:     0; 0.0%)
Overall Newton Iterations:    1534      (Wasted:     0; 0.0%)
Overall Linear Iterations:   22713      (Wasted:     0; 0.0%)

\end{verbatim}

In the above example, the simulation took 8.4s to setup followed by
240.47s of simulation time during which it performed at total of 369
time steps, which required a total of 1534 Newton iterations. The
summary includes a breakdown of the simulation time; of which the linear
solve time was the largest component in this example.

The total linear solve time was 83.22s of which 0.0s (or 0.0\% of the
total) was wasted due to convergence failures (since there were no
convergence failures in this example). Similarly, there were a total of
1,534 Newton iterations of which zero were wasted due to convergence
failures.

RESTART files contain the solution data arrays for each active cell in
the model, and represents a ``snapshot'' of the reservoir state at a
reporting time step, for example the pressure (PRESSURE, PPCW, and
PPCG), fluid saturations (SGAS, SOIL, and SWAT), and fluid state (RS and
RV) for all active cells. The PPCW and PPCG arrays are the water and
capillary solution arrays and the RS and RV arrays are the gas-oil and
vaporized-oil ratio (condensate-gas ratio) solution arrays. The data
written to the RESTART file is dependent on the fluid phases and options
invoked in the model, as well as any additional user requested solution
data. The default arrays written out are always sufficient to enable the
simulator to restart from a previous run {[}For example in a three phase
(oil, gas, and water) runs only two phase saturations are needed for a
``restart'' run, as the third phase can be calculated in by summing two
phases and subtracting from one.{]}. This was the original purpose of
the RESTART file, that is saving the reservoir state in such a manner,
is to enable a ``restart'' of the simulation.

RESTART files also contain the necessary information to do restart runs
for Groups, Wells, Connections, Aquifer data, etc. This information
enables restart runs based on Schedule information read from the Restart
file only, using the SKIPREST keyword in the SCHEDULE section.

As of the OPM Flow 2019-04 release the simulator can ``restart'' from
the commercial simulator's generated RESTART files as well as the
commercial simulator being to ``restart'' from the OPM Flow's generated
RESTART files. See section Error: Reference source not found Error:
Reference source not found for additional information.

As RESTART files contain a complete description of the reservoir state,
they are also used in post processing software to visualize the
reservoir solutions through time, for example by OPM ResInsight. And as
the functionality of the both OPM Flow and the commercial simulator have
expanded over time additional solution arrays have been added to the
RESTART file to enable ``restarts'' (POLYMER -- polymer saturations), as
well as to write out additional user specific solution arrays (FIPOIL -
fluid in-place oil array), that are not necessary for a ``restart'', but
are consider useful in understanding the reservoir performance.

The frequency and type of data written to the RESTART file is controlled
via the RPTSCHED and RPTRST keywords in the SCHEDULE, with the latter
having greater functionality and flexibility. For example, to request
that the standard restart data be written out every month using the
RPTRST keyword:

\begin{verbatim}
--
--       RESTART CONTROL BASIC = 4 (YEARLY) 5 (MONTHLY)
--
RPTRST
         BASIC=5                                                               /
\end{verbatim}

The next example requests that the standard restart data be written at
every report time step until this switch is reset and all the restarts
are kept. In addition to the standard data, the gas, oil, and water
relative permeability data will also be written out at each report time
step.

\begin{verbatim}
--
--       RESTART CONTROL BASIC = 4 (YEARLY) 5 (MONTHLY)
--
RPTRST
         BASIC=2  KRG   KRO   KRW                                             /
\end{verbatim}

RESTART solution data is written out via two separate files; the RESTART
Index file that specifies and defines the format and data type written
to the RESTART Data file, and the RESTART Data file that contains the
solution data for each active cell in the model for the requested report
time steps. The commercial simulator writes out both files, where as OPM
Flow only writes out the RESTART Data file, as the data can be read by
most post-processing software, including OPM ResInsight, using only the
RESTART Data file as input. Secondly, both the commercial simulator and
OPM Flow can be ``restarted'' using just the RESTART Data file.

The following sections outline the format for these two file types.

This file type is not supported or required by OPM Flow or OPM
ResInsight to read or write restart files.

This file type contains the global and LGR grid solution data property
data {[}Several keywords are used for both global and LGR grids, for
example the INTEHEAD, LOGIHEAD, and DOUBHEAD keywords and SOLUTION
arrays.{]} (PRESSURE, SGAS, SOIL, and SWAT) and the group, well, and
connection data {[}The group, well, and connection production data is
not stored on the RESTART file, but is instead stored on the SUMMARY
file. The RESTART file contains the group, well, and connection
configurations as the they change versus time. For example, the status
of a well (OPEN or SHUT), or the connections open to flow in a well.{]}.
The overall structure of this file is similar to the INIT file, and the
individual data keywords are structured as presented in Table A.20.

\begin{longtable}[]{@{}
  >{\raggedright\arraybackslash}p{(\linewidth - 12\tabcolsep) * \real{0.1429}}
  >{\raggedright\arraybackslash}p{(\linewidth - 12\tabcolsep) * \real{0.1429}}
  >{\raggedright\arraybackslash}p{(\linewidth - 12\tabcolsep) * \real{0.1429}}
  >{\raggedright\arraybackslash}p{(\linewidth - 12\tabcolsep) * \real{0.1429}}
  >{\raggedright\arraybackslash}p{(\linewidth - 12\tabcolsep) * \real{0.1429}}
  >{\raggedright\arraybackslash}p{(\linewidth - 12\tabcolsep) * \real{0.1429}}
  >{\raggedright\arraybackslash}p{(\linewidth - 12\tabcolsep) * \real{0.1429}}@{}}
\toprule\noalign{}
\begin{minipage}[b]{\linewidth}\raggedright
Reference
\end{minipage} & \begin{minipage}[b]{\linewidth}\raggedright
Section
\end{minipage} & \begin{minipage}[b]{\linewidth}\raggedright
Global Section keywords
\end{minipage} & \begin{minipage}[b]{\linewidth}\raggedright
Global Data keywords
\end{minipage} & \begin{minipage}[b]{\linewidth}\raggedright
LGR Data keywords
\end{minipage} & \begin{minipage}[b]{\linewidth}\raggedright
Global Section keywords
\end{minipage} & \begin{minipage}[b]{\linewidth}\raggedright
Notes
\end{minipage} \\
\midrule\noalign{}
\endhead
\bottomrule\noalign{}
\endlastfoot
A.8.3 & SEQNUM & & & & Global grid header keyword that defines the start
of a RESTART global grid time step data set. & \\
A.8.3 & INTEHEAD LOGIHEAD DOUBHEAD & & & & Global grid header keyword
that defines the integer, logical, and double precision variables for
this header keyword. & \\
A.8.4 & IGRP SGRP etc. & & & & Group, well, and connection data status
keywords for this reporting time step. & \\
A.8.6 & IAAQ SAAQ ect. & & & & Aquifer definition arrays. & \\
A.8.7 & HIDDEN & & & & Defines HIDDEN solution data, that is a list of
solution arrays NOT to be read by post processing software. & \\
A.8.8 & STARTSOL & & & & This keyword marks the start of the solution
variable section for the global grid. & \\
A.8.8 & & PRESSURE & & & Solution array section (PRESSURE, SGAS, SOIL,
SWAT, etc.) for the Global grid. & \\
SWAT & & & & & & \\
etc. & & & & & & \\
A.8.8 & LGRNAMES & & & & A list of LGRs in this report step. & \\
A.8.8 & ENDSOL & & & & Marks the end of the solution variable section
for the global grid. & \\
A.5.1 & & & LGR LGRHEADI LGRHEADQ LGRHEADD & & LGR grid header keyword
for that defines the start of a RESTART LGR data set. & \\
A.8.3 & & & INTEHEAD LOGIHEAD DOUBHEAD & & LGR header keyword and
defines the integer, logical, and double precision variables for this
header keyword. & \\
A.8.4 & & & IGRP SGRP etc. & & Group, well, and connection data status
keywords for this reporting time step. & \\
A.8.7 & & & HIDDEN & & Defines HIDDEN solution data section. & \\
A.8.8 & & & STARTSOL & & Start of LGR solution array section. & \\
A.8.8 & & & Solution Keywords & PRESSURE & LGR solution array section --
same as global grid. & \\
SWAT & & & & & & \\
etc. & & & & & & \\
A.8.8 & & & ENDSOL & & End of the solution variable section for the
current LGR grid. & \\
A.5.1 & & & ENDLGR & & End of current LGR section. & \\
A.5.1 & & & LGR LGRHEADI LGRHEADQ LGRHEADD & & LGR grid header keywords.
& \\
A.8.4 & & & IGRP SGRP etc. & & Group, well, and connection data status
keywords for this reporting time step. & \\
A.8.7 & & & HIDDEN & & Defines HIDDEN solution data section. & \\
A.8.8 & & & STARTSOL & & Start of solution array section. & \\
A.8.8 & & & Solution Keywords & PRESSURE & LGR solution array section.
& \\
& SWAT & & & & & \\
etc. & & & & & & \\
A.8.8 & & & ENDSOL & & End of the solution variable section. & \\
A.5.1 & & & ENDLGR & & End of current LGR section. & \\
A.8.3 & SEQNUM & & & & Global grid header keyword that defines the start
of a RESTART global grid time step data set. & \\
Notes: & & & & & & \\
\end{longtable}

\emph{Table A.20: RESTART Data File Structure}

The individual keywords are described in the following sections.

This set of keywords {[}The RESTART Data -- Header keywords are used
both for global and LGR grids with the global grid data for the global
grid and the individual LGR data for the LGR grids. The LGR data is
preceded by a series of LGR header keywords and terminated by an LGR
termination keyword.{]} are the first set of keywords that should be
read or written to the RESTART file. This data set includes INTEHEAD,
LOGIHEAD, and DOUBHEAD keywords that define versus parameters used in
subsequent keywords in the RESTART file for the global and LGR grids.

The structure for this type of keyword is defined in Table A.21.

\begin{longtable}[]{@{}
  >{\raggedright\arraybackslash}p{(\linewidth - 10\tabcolsep) * \real{0.1667}}
  >{\raggedright\arraybackslash}p{(\linewidth - 10\tabcolsep) * \real{0.1667}}
  >{\raggedright\arraybackslash}p{(\linewidth - 10\tabcolsep) * \real{0.1667}}
  >{\raggedright\arraybackslash}p{(\linewidth - 10\tabcolsep) * \real{0.1667}}
  >{\raggedright\arraybackslash}p{(\linewidth - 10\tabcolsep) * \real{0.1667}}
  >{\raggedright\arraybackslash}p{(\linewidth - 10\tabcolsep) * \real{0.1667}}@{}}
\toprule\noalign{}
\begin{minipage}[b]{\linewidth}\raggedright
No.
\end{minipage} & \begin{minipage}[b]{\linewidth}\raggedright
keyword Name
\end{minipage} & \begin{minipage}[b]{\linewidth}\raggedright
RESTART Data - Header keywords Table A.21
\end{minipage} & \begin{minipage}[b]{\linewidth}\raggedright
Status or Value
\end{minipage} & \begin{minipage}[b]{\linewidth}\raggedright
\end{minipage} & \begin{minipage}[b]{\linewidth}\raggedright
\end{minipage} \\
\midrule\noalign{}
\endhead
\bottomrule\noalign{}
\endlastfoot
Keyword & No.~of Entries & Data Type & & & \\
& SEQNUM & Global grid header keyword for the RESTART Data file that
defines the start of a RESTART global grid time step data set for
unified restart files only. & Optional & & \\
1-1 & Format & SEQNUM & 1 & INTE & \\
2-1 & Data & Report time step number. & 0 & & \\
& INTEHEAD & Global grid header keyword for the RESTART Data file that
defines the integer variables for this time step. NIHEAD is the number
of entries associated with this keyword, and is dependent on the
features being used in the model. Undefined values should be set to
zero. & Required & & \\
1-1 & Format & INTEHEAD & NIHEAD & INTE & \\
2-1 & Data & An encoded integer corresponding to the time the file was
created. For files not originating from the commercial simulator, this
value may be set to zero. & ISNUM & & \\
2-2 & Data & Simulator version, should be set to zero. & VERSION & & \\
2-3 & Data & Unit system used in the simulation, set to 1 for metric, 2
for field, and 3 for laboratory & & & \\
2-4 to 2-8 & Data & Undefined. & 0 & & \\
2-9 & Data & DIMENS key word in the RUNSPEC section - the number of grid
blocks in the x-direction. & NX & & \\
2-10 & Data & DIMENS key word in the RUNSPEC section - the number of
grid blocks in the y-direction. & NY & & \\
2-11 & Data & DIMENS key word in the RUNSPEC section - the number of
grid blocks in the z-direction. & NZ & & \\
2-12 & Data & Number of global active cells in the model . & NACTIV &
& \\
2-13 & Data & Not used. & 0 & & \\
2-14 & Data & Not used. & 0 & & \\
2-15 & Data & Type of phases in the model, set to 1 for oil, 2 for
water, 3 for oil-water, 4 for gas, 5 for oil-gas, 6 for gas-water, and 7
for oil-water-gas. & PHASE & & \\
2-16 & Data & Undefined. & 0 & & \\
2-17 & Data & Number of wells in the model. & NWELLS & & \\
2-18 & Data & Maximum number of grid block connections per well. &
NCWMAX & & \\
2-19 & Data & Actual number of groups. & NGRP & & \\
2-20 & Data & Maximum number of wells belonging to a group in the model.
& NWGMAX & & \\
2-21 & Data & Maximum number of groups for this model. & NGMAXZ & & \\
2-22 to 2-24 & Data & Not used. & & & \\
2-25 & Data & Number of values per well in the IWEL array (default value
is 97). & NIWELZ & & \\
2-26 & Data & Number of values per well in the SWEL array. & NSWELZ &
& \\
2.27 & Data & Number of values per well in the XWEL array. & NXWELZ &
& \\
2-28 & Data & Number of eight character words per well in the ZWEL
array. & NZWELZ & & \\
2-29 to 2-32 & Data & Not used. & & & \\
& NICONZ & & & & \\
2-33 & Data & Number of values per completion in ICON array (default
19). & NICONZ & & \\
2-34 & Data & Number of values per completion in SCON array. & NSCONZ &
& \\
2-35 & Data & Number of values per completion in XCON array. & NXCONZ &
& \\
2-36 & & Not used. & 0 & & \\
2-37 & Data & Number of values per group in IGRP array. & NIGRPZ & & \\
2-38 & Data & Number of values per group in SGRP array. & NISGRPZ & & \\
2-39 & Data & Number of values per group in XGRP array. & NIXGRPZ & & \\
2-40 & Data & Number of values per group in ZGRP array. & NIZGRPZ & & \\
2-41 & & Not used. & & & \\
2-42 & Data & Maximum number of analytical aquifer connections. & NCAMAX
& & \\
2-43 & Data & Number of values per aquifer in the IAAQ array. & NIAAQZ &
& \\
2-44 & Data & Number of values per aquifer in the SAAQ array. & NSAAQZ &
& \\
2-45 & Data & Number of values per aquifer in the XAAQ array. & NXAAQZ &
& \\
2-46 & Data & Number of values per aquifer connection in the ICAQ array.
& NICAQZ & & \\
2-47 & Data & Number of values per aquifer connection in the SCAQ array.
& NSCAQZ & & \\
2-48 & Data & Number of values per aquifer connection in the ACAQ array.
& NXCAQZ & & \\
2-49 to 2-51 & & Not used. & & & \\
2-52 & Data & Index indicating if group control is used or not (1 -- for
GCONPROD group control, 2 for GCONINJE control, or 0 if no group
control) & NGCTRL & & \\
2-53 to 2.58 & & Undefined. & & & \\
2-59 & Data & Parameter to determine the nominated phase for the guide
rate. & NGRNPH & & \\
2-60 to 2-64 & & Undefined. & & & \\
2-65 & Data & Calendar day of report step, the value should be greater
than or equal to one and less than or equal to 31 for the calendar days.
& DAY & & \\
2-66 & Data & Calendar month of report step, the value should be greater
than or equal to one and less than or equal to 12 for the calendar
month. & MONTH & & \\
2-67 & Data & Calendar year of report step, a positive four digit
integer value of the start year, which must be specified fully by four
digits, that is 2019. & YEAR & & \\
2-68 & Data & Numerical solver step, that is the number of solver steps
the simulator has performed so far. & NUM\_SOLVER\_STEP & & \\
2-69 & Data & Report step. & REPORT\_STEP & & \\
2-70 to 2-71 & & Undefined. & 0 & & \\
2-72 & Data & Index for WHISTCTL keyword. & WHISTC & & \\
2-73 to 2-80 & & Undefined. & & & \\
2-81 & Data & TUNING keyword in the SCHEDULE section - the maximum
number of Newton iterations for a time step. & NEWTMX & & \\
2-82 & Data & TUNING keyword in the SCHEDULE section - the minimum
number of Newton iterations for a time step & NEWTMN & & \\
2-83 & Data & TUNING keyword in the SCHEDULE section - the maximum
number of linear iterations within a Newton iteration. & LITMAX & & \\
2-84 & Data & TUNING keyword in the SCHEDULE section - the minimum
number of linear iterations within a Newton iteration. & LITMIN & & \\
2-85 & & Undefined. & & & \\
2-86 & & Undefined. & & & \\
2-87 & Data & TUNING keyword in the SCHEDULE section - the maximum
number of iterations within a well flow calculation. & MXWSIT & & \\
2-88 & Data & TUNING keyword in the SCHEDULE section - the maximum
number of iterations for solving the bottom-hole pressure for wells
under tubing head pressure control within a well flow calculation. &
MXWPIT & & \\
2-89 & & Undefined. & & & \\
2-90 & Data & REGDIMS and TABDIMS keywords in the RUNSPEC section. The
maximum number of regions in the FIPNUM region array. Note that this
parameter may also be set on the TABDIMS keyword as well. If NTFIP is
set in both places then the maximum value should be used. & NTFIP & & \\
2-91 to 2-94 & & Undefined. & & & \\
2-95 & Data & Simulator code used to generate the file, set to 100 for
Schlumberger's ECLIPSE 100, 300 for Schlumberger's ECLIPSE 300, 500 for
Schlumberger's ECLIPSE 300 Thermal, 700 for Schlumberger's INTERSECT
simulator, 800 for Schlumberger's FrontSim simulator, or a negative
value for other simulators. OPM Flow uses a value of 100. & IPROG & & \\
2-96 to 2-99 & & Undefined. & 0 & & \\
2-100 & Data & REGDIMS keyword in the RUNSPEC section. The total maximum
number of regions. The FIPNUM regions are defined by (2-90). If
additional sets of fluid in-place regions are used, as per the FIPxxx
series of fluid in-place region keywords, then NMFIPR is the sum of all
FIP regions. & NMFIPR & & \\
2-101 to 2- 131 & & & & & \\
2-132 & Data & NETWORK keyword in the RUNSPEC section - maximum number
of nodes in an extended network model or zero if the extended network
option has not been activated. & NODMAX & & \\
2-133 & Data & NETWORK keyword in the RUNSPEC section - maximum number
of branches in an extended network model or zero if the extended network
option has not been activated. & NBRMAX & & \\
2-134 & Data & Number of items per branch in the IBRAN array when the
extended network model is active, or zero if the extended network option
has not been activated. & NIBRAN & & \\
2-135 & Data & Number of items per branch in the RBRAN array when the
extended network model is active, or zero if the extended network option
has not been activated. & NRBRAN & & \\
2-136 & Data & Number of items per node in the INODE array when the
extended network model is active, or zero if the extended network option
has not been activated. & NINODE & & \\
2-137 & Data & Number of items per node in the RNODE array when the
extended network model is active, or zero if the extended network option
has not been activated. & NRNODE & & \\
2-138 & Data & Number of items per node in the ZNODE array when the
extended network model is active, or zero if the extended network option
has not been activated. & NZNODE & & \\
2-139 & Data & Number of items in the INOBR array. & NINOBR & & \\
2-140 to 2-156 & & Undefined. & 0 & & \\
2-157 & Data & The number of actions in the data set. & NO\_ACT & & \\
2-158 & Data & Maximum number of lines of schedule data for ACTION
keyword - including ENDACTIO. & MAX\_LINES & & \\
2-159 & Data & Maximum number of eight character strings per input line
of Action data (rounded up from input deck). & MAXSPRLINE & & \\
2-160 to 2-162 & & Undefined. & & & \\
2-163 & Data & Maximum number of aquifer connections actually used in
the model. & NGCAUS & & \\
2-164 & Data & Maximum number of wells in the model. & NWMAXZ & & \\
2-165 to 2-174 & & Undefined. & 0 & & \\
2-175 & Data & Number of multi-segmented wells defined with the WELSEGS
keyword in the SCHEDULE section, for when multi-segment wells have been
activated, or zero otherwise. & NSEGWL & & \\
2-176 & Data & WSEGDIMS keyword in the RUNSPEC section - maximum number
of multi-segment wells, for when multi-segment wells have been
activated, or zero otherwise. & NSWLMX & & \\
2-177 & Data & WSEGDIMS keyword in the RUNSPEC section - maximum number
of segments per multi-segment well, for when multi-segment wells have
been activated, or zero otherwise. & NSEGMX & & \\
2-178 & Data & WSEGDIMS keyword in the RUNSPEC section - maximum number
of branches per multi-segment well, including the main branch, for when
multi-segment wells have been activated, or zero otherwise. & NLBRMX &
& \\
2-179 & Data & Number of entries per segment in the multi-segment well
ISEG array, for when multi-segment wells have been activated, or zero
otherwise. & NISEGZ & & \\
2-180 & Data & Number of entries per segment in the multi-segment well
RSEG array, for when multi-segment wells have been activated, or zero
otherwise. The number of data items per well segment in the RSEG restart
vector is dependent on the number of active phases (oil, gas, water) in
the simulation run. In particular, the following relations hold: &
NRSEGZ & & \\
Scenario & Phases & NRSEGZ & & & \\
Single Phase & Oil, Gas, or Water & 126 & & & \\
Two Phase & Oil/Gas, Oil/Water, Gas/Water & 134 & & & \\
Three Phase & Oil/Gas/Water & 146 & & & \\
2-181 & Data & Number of entries per segment in the multi-segment well
ILBR array, for when multi-segment wells have been activated, or zero
otherwise. & NILBRZ & & \\
2-181 to 2-206 & & Undefined. & 0 & & \\
2-207 & Data & The first part of the current simulation time in the form
HH:MM:SS, that is HH for hours. The value should be greater than or
equal to zero and less than or equal to 23. & IHOURZ & & \\
2-208 & Data & The second part of the current simulation time in the
form HH:MM:SS, that is MM for minutes. The value should be greater than
or equal to zero and less than or equal to 59. & IMINTS & & \\
2-209 to 2-223 & & Undefined. & 0 & & \\
2-224 & Data & Number of INTEGER lines in the AQUNUM data set. & NILAQN
& & \\
2-225 & Data & Number of REAL lines in the AQUNUM data set. & NIRAQN &
& \\
2-226 & & Undefined. & 0 & & \\
2-227 & Data & Number of entries in the AQUNUM data set. & NUMAQN & & \\
2-228 to 2-234 & & Undefined. & 0 & & \\
2-235 & Data & Number of items in the ICOT array. & NICOTZ & & \\
2-236 & Data & Number of items in the XCOT array. & NXCOTZ & & \\
2-237 & Data & Number of items in the IWET array. & NIWETZ & & \\
2-238 & Data & Number of items in the XWET array. & NXWETZ & & \\
2-239 & Data & Number of items in the IGRT array. & NIGRTZ & & \\
2-240 & Data & Number of items in the XGRT array. & NXGRTZ & & \\
2-241 & Data & Number of tracers in the model plus two. & NSTRA2 & & \\
2-242 to 2-245 & & Undefined. & 0 & & \\
2-246 & Data & Maximum number of conditions per action. & MAXACTC & & \\
2-253 & Data & Maximum number of analytical aquifers in the model. &
MAAQID & & \\
2.254 to 2-262 & & Undefined. & 0 & & \\
2-263 & Data & Number of Field UDQ data (parameters). & NOFUDQS & & \\
2-264 & Data & Number of Group UDQ data (parameters). & NOGUDQS & & \\
2-265 & & Undefined. & & & \\
2-266 & & Undefined. & & & \\
2-267 & Data & Number of Well UDQ data (parameters). & NOWUDQS & & \\
2-268 & Data & Integer seed value for the RAND. & UDQPAR\_1 & & \\
2-269 & & Undefined. & & & \\
2-270 & & Undefined. & & & \\
2.271 & Data & Maximum number of Chord segment links per multi-segment
well for when multi-segment wells have been activated, or zero
otherwise. & NCRDMX & & \\
2-272 to 2-290 & & Undefined. & 0 & & \\
2-291 & & Number of integer IUADs. & NOIUADS & & \\
2-292 & & Number of integer IUAPs. & NOIUAPS & & \\
2-293 to 2-296 & & Undefined. & & & \\
2-297 & Data & & RSEED & & \\
2-298 to 2-410 & & Undefined. & & & \\
2-411 & Data & The third part of the current simulation time in the form
HH:MM:SS, that is SS for seconds. The value should be greater than or
equal to zero and less than or equal to 59. & ISECND & & \\
& LOGIHEAD & Global grid keyword for the RESTART Data file that defines
the logical variables (T for true and F for false) for this time step.
NLHEAD is the number of entries associated with this keyword, and is
dependent on the features being used in the model. Undefined items
should be set to F (False). & Required & & \\
1-1 & Format & LOGIHEAD & NLHEAD & LOGI & \\
2-1 & Data & Dissolved gas is present in the model indicator, set to T
if present in the model else set to F. & DISGAS & & \\
2-2 & Data & Vaporized oil is present in the model indicator, set to T
if present in the model else set to F. & VAPOIL & & \\
2-3 & Data & SATOPTS in the RUNSPEC section - directional relative
permeability curves are active in the model. & DIRECT & & \\
2-4 & Data & SATOPTS keyword in the RUNSPEC section - reversible
relative permeability curves (black-oil) are active in the model. &
IRREVERS & & \\
Radial grid geometry has been used for a compositional model indicator,
set to T or F. Note that currently OPM Flow does not support radial
grids and therefore this item should be set to F. & RADIAL & & & & \\
2-5 & Data & Radial grid geometry has been used for a black-oil model
indicator, set to T or F. Note that currently OPM Flow does not support
radial grids and therefore this item should be set to F. & RADIAL & & \\
Reversible relative permeability (compositional) are active in the
model. & IRREVERS & & & & \\
2-6 & & Undefined. & & & \\
2-7 & Data & Enable hysteresis indicator. & HYSER & & \\
2.8 to 2-14 & & Undefined. & F & & \\
2-15 & Data & The dual porosity option has been activated in the model
indicator, set to T or F. Note that currently OPM Flow does not support
dual porosity grids and therefore this item should be set to F. &
DUALPORO & & \\
2-16 & & Undefined. & F & & \\
2-17 & Data & ENDSCALE keyword in the RUNSPEC section - enable end-point
scaling indicator. & ENDSCALE & & \\
2-18 & Data & ENDSCALE keyword in the RUNSPEC section - directional
end-point scaling indicator. & DIRECT & & \\
2-19 & Data & ENDSCALE keyword in the RUNSPEC section - reversible
end-point scaling indicator. & IRREVERS & & \\
2-20 & Data & SCALECRS keyword in the PROPS section - alternative three
point end-point scaling indicator. & SCALECRS & & \\
2-21 to 2-30 & & Undefined. & F & & \\
2-31 & Data & Coal Bed Methane option in the black-oil model indicator,
set to T or F. Note that currently OPM Flow does not support the Coal
Bed Methane option and therefore this item should be set to F. & F &
& \\
2-32 to 2-38 & & Undefined. & F & & \\
2-39 & Data & Constant oil compressibility (PVCDO) indicator. & ConstCo
& & \\
2-40 to 2-75 & & Undefined. & & & \\
2-76 & Data & Multi-segment wells are present in the model indicator. &
HasMSWells & & \\
2-77 to 2-127 & & Undefined. & & & \\
2-128 & & Coal Bed Methane option in the compositional model indicator,
set to T or F. Note that currently OPM Flow does not support the Coal
Bed Methane option and therefore this item should be set to F. & F &
& \\
& DOUBHEAD & Global grid keyword for the RESTART Data file that defines
the double precision REAL variables for this time step. NDHEAD is the
number of entries associated with this keyword, and is dependent on the
features being used in the model. Undefined values should be set to
zero. & Required & & \\
1-1 & Format & DOUBHEAD & NDHEAD & DOUB & \\
2-1 & Data & The simulation REPORT TIME STEP, expressed in days for all
units, except for the LAB system of units which should be expressed in
hours. & & & \\
2-2 & Data & TUNING keyword in the SCHEDULE section - maximum length of
next time step. & TSINIT & & \\
2-3 & Data & TUNING keyword in the SCHEDULE section - maximum length of
time step after next. & TSMAXZ & & \\
2-4 & Data & TUNING keyword in the SCHEDULE section - minimum length of
all time steps. & TSMINZ & & \\
2-5 & Data & TUNING keyword in the SCHEDULE section - minimum length of
all chopped time steps. & TSMCHP & & \\
2-6 & Data & TUNING keyword in the SCHEDULE section - maximum growth
rate a time step can be increased by, subject to the maximum allowable
time step size set by TSMAXZ. & TSFMAX & & \\
2-7 & Data & TUNING keyword in the SCHEDULE section - minimum decay rate
a time step can be decreased by, subject to the minimum allowable time
step size set by TSMINZ. & TSFMIN & & \\
2-8 & Data & TUNING keyword in the SCHEDULE section - the decay rate a
time step can be decreased by after the number of target iterations has
been exceeded. & TSFCNV & & \\
2-9 & Data & TUNING keyword in the SCHEDULE section - the time
truncation error target. & TRGTTE & & \\
2-10 & Data & TUNING keyword in the SCHEDULE section - thee non-linear
convergence error. & TRGCNV & & \\
2-11 & Data & TUNING keyword in the SCHEDULE section - the target
material balance error. & TRGMBE & & \\
2-12 & Data & TUNING keyword in the SCHEDULE section - the linear
convergence error target. & TRGLCV & & \\
2-13 to 2-16 & & Undefined. & & & \\
2-17 & Data & TUNING keyword in the SCHEDULE section - the maximum time
truncation error. & XXXTTE & & \\
2-18 & Data & TUNING keyword in the SCHEDULE section - the maximum
non-linear convergence error. & XXXCNV & & \\
2-19 & Data & TUNING keyword in the SCHEDULE section - the maximum mass
balance error. & XXXMBE & & \\
2-20 & Data & TUNING keyword in the SCHEDULE section - the maximum
linear convergence error. & XXXLCV & & \\
2-21 & Data & TUNING keyword in the SCHEDULE section - the maximum well
flow convergence error. & XXXWFL & & \\
2-22 to 2-82 & & Undefined. & & & \\
2-83 & Data & TUNING keyword in the SCHEDULE section - target pressure
change within a time step. & TRGDPR & & \\
2-84 & Data & TUNING keyword in the SCHEDULE section - time step growth
factor of the time step after a convergence failure. & TFDIFF & & \\
2-85 & Data & TUNING keyword in the SCHEDULE section - the maximum
pressure change at the last Newton iteration. & DDPLIM & & \\
2.86 & Data & TUNING keyword in the SCHEDULE section - the maximum
saturation change at the last Newton iteration. & DDSLIM & & \\
2.87 & & Undefined. & & & \\
2-88 & Data & GUIDERAT keyword in the SCHEDULE section - guide rate
parameter A. & GRPAR\_A & & \\
2-89 & Data & GUIDERAT keyword in the SCHEDULE section - guide rate
parameter B. & GRPAR\_B & & \\
2-90 & Data & GUIDERAT keyword in the SCHEDULE section - guide rate
parameter C. & GRPAR\_C & & \\
2-91 & Data & GUIDERAT keyword in the SCHEDULE section - guide rate
parameter D. & GRPAR\_D & & \\
2-92 & Data & GUIDERAT keyword in the SCHEDULE section - guide rate
parameter E. & GRPAR\_E & & \\
2-93 & Data & GUIDERAT keyword in the SCHEDULE section - guide rate
parameter F. & GRPAR\_F & & \\
2-94 to 2-97 & & Undefined. & & & \\
2.98 & Data & Guide rate parameter delay interval. & GRPAR\_INT & & \\
2-99 & & Undefined. & & & \\
2-100 & Data & TUNING keyword in the SCHEDULE section - the maximum
throughput ratio over a time step. & THRURPT & & \\
2-101 & Data & TUNING keyword in the SCHEDULE section - maximum
tolerable pressure change within a time step. & XXXDPR & & \\
2-102 & Data & TUNING keyword in the SCHEDULE section - the target fluid
in-place error in Local Grid Refinements. & TRGFIP & & \\
2-103 & Data & TUNING keyword in the SCHEDULE section - target
surfactant change when the Surfactant Model is active in the run. &
TRGSFT & & \\
2-104 to 2-144 & & Undefined. & & & \\
2-145 & Data & Guide rate parameter damping factor. & GRPAR\_DAMP & & \\
2-146 to 2-161 & & Undefined. & & & \\
2-161 & Data & The simulation START time, expressed in days for all
units, except for the LAB system of units which should expressed in
hours. & & & \\
2-162 & Data & Cumulative time from the start of the simulation,
expressed in days for all units, except for the LAB system of units
which should expressed in hours. That is the sum of values (2-1) and
(2-161). & & & \\
2-163 to 2-212 & & Undefined. & 0 & & \\
2-213 & Data & UDQPARAM item number 2 (Permitted range (+/-) of
user-defined quantities). & UDQPAR\_2 & & \\
2-214 & Data & UDQPARAM item number 3 (Value given to undefined elements
when outputting data). & UDQPAR\_3 & & \\
2-215 & Data & UDQPARAM item number 4 (fractional equality tolerance
used in == \textless= etc. functions). & UDQPAR\_4 & & \\
Notes: & & & & & \\
\end{longtable}

\emph{Table A.21: RESTART Data - Header Keywords}

\subsubsection{Example: RESTART Data - Header
Keywords}\label{example-restart-data---header-keywords}

The following example shows a typical formatted RESTART Data - Header
keywords.

\begin{verbatim}
'INTEHEAD'         249 'INTE'
 -955283513      200400           2       -2345       -2345       -2345
      -2345       -2345          10          10           1         100
         10       -2345           7       -2345           0           1
          0           2           2           0           0           0
        110         108         109           3          97          93
      -2345       -2345          19          38          53       -2345
         97          93         146           5           0           1
         15          24           8           5           2           4
          0           0           0           0           0           0
          0           0           0           0           0           0
          0           0           0           0          19          10
       1982           0           0           0           1           0
          0           0           0           0           1          10
          0           0          12           1          25           1
      -2345       -2345           8           8           5           1
          2           1           2           0         100           2
          7           0         -11           1           0           1
          0           0           0          10          10           1
        100           1           1           1           1          10
         10           1          19           1           0           0
          0           0           0           1           1           1
          0           0           0           0           0           0
          0          14          10          10          16           1
          1           1           1           1           2           1
          1           1           1           1           1          31
        108           0           0           0           0           0
          0          50          10           4           5           9
          0           2           8           8          12           1
         25           1      -32767      -32767      -32767      -32767
          0           1           1           1          18          86
          5           1           1           1           1          18
         86      -32767      -32767           0           0         109
         53         146           8           0          19         110
          0           0           1           1          86           0
          0           0           0           0           0           1
          0           0           0           0           1           0
          0           1           0          -1          12           0
          0          10          13           1           0           0
          0           0           2           0           0        3600
          1           6           1          10           1          10
          1           1           1           0          30           3
         18          10           9
'LOGIHEAD'          79 'LOGI'
 T  F  F  T  F  F  F  F  T  F  F  F  F  F  F  F  T  F  T  F  F  F  F  F  F
 F  F  F  F  F  F  F  F  F  F  F  F  F  F  F  F  F  F  F  F  F  F  F  F  F
 F  F  F  F  F  F  F  F  F  F  F  F  F  F  F  F  F  F  F  F  F  F  F  F  F
 F  F  F  F
'DOUBHEAD'         185 'DOUB'
  0.00000000000000D+00   0.10000000000000D+01   0.36500000000000D+03
  0.10000000149012D+00   0.15000000596046D+00   0.30000000000000D+01
  0.30000001192093D+00   0.10000000149012D+00   0.10000000149012D+00
 ………………………….……………..
 -0.10000000200409D+21   0.10000000000000D+01   0.10000000000000D+01
  0.10000000000000D+01   0.00000000000000D+00   0.00000000000000D+00
  0.10000000000000D+01   0.10000000000000D-03
\end{verbatim}

End of Example

The Restart Data -- Well and Group Completion Keywords {[}The RESTART
Data -- Group, Well and Connection keywords are used both for global and
LGR grids with groups, well, and connection data for the global grid,
and wells and connection data for wells contained in the individual LGRs
for the LGR grids. The LGR data is preceded by a series of LGR header
keywords and terminated by an LGR termination keyword.{]} define various
parameters associated with groups, wells, and connection keywords in the
RESTART file for both global and LGR grids.

The structure for this type of keyword is based on the data type as
summarized below:

\begin{verbatim}
- The group data set keywords are specified in Table A.22 and consists of the IGRP, SGRP, XGRP,  and ZGRP keywords.
- Table A.23 contains the multi-segment well keywords that consist of the ISEG, RSEG, ILBS, ILBR, and ICRD keywords.
- For wells the keywords are IWEL, SWEL, XWEL, and ZWEL and Table A.24 describes these keywords.
- Well connection keywords are specified in Table A.25 and consist of the ICON, SCON, and XCON keywords.
- Tracer keywords are listed in Table A.26 and consists of the ICOT, XCOT, IWET, XWET, IGRT, and XGRT keywords.
- Finally, the network keywords INODE, IBRAN, INOBR, RNODE, RBRAN, and ZNODE are specified in Table A.27.
\end{verbatim}

Note that if multi-segment wells, tracers, and networks are absent from
the simulation deck then this data will not be written to the RESTART
Data file.

The following table outlines the keywords for groups.

\begin{longtable}[]{@{}
  >{\raggedright\arraybackslash}p{(\linewidth - 10\tabcolsep) * \real{0.1667}}
  >{\raggedright\arraybackslash}p{(\linewidth - 10\tabcolsep) * \real{0.1667}}
  >{\raggedright\arraybackslash}p{(\linewidth - 10\tabcolsep) * \real{0.1667}}
  >{\raggedright\arraybackslash}p{(\linewidth - 10\tabcolsep) * \real{0.1667}}
  >{\raggedright\arraybackslash}p{(\linewidth - 10\tabcolsep) * \real{0.1667}}
  >{\raggedright\arraybackslash}p{(\linewidth - 10\tabcolsep) * \real{0.1667}}@{}}
\toprule\noalign{}
\begin{minipage}[b]{\linewidth}\raggedright
No.
\end{minipage} & \begin{minipage}[b]{\linewidth}\raggedright
keyword Name
\end{minipage} & \begin{minipage}[b]{\linewidth}\raggedright
RESTART Data -- Group, Well, and Connection keywords Table A.22 - Group
Keywords
\end{minipage} & \begin{minipage}[b]{\linewidth}\raggedright
Status or Value
\end{minipage} & \begin{minipage}[b]{\linewidth}\raggedright
\end{minipage} & \begin{minipage}[b]{\linewidth}\raggedright
\end{minipage} \\
\midrule\noalign{}
\endhead
\bottomrule\noalign{}
\endlastfoot
Keyword & No.~of Entries & Data Type & & & \\
& IGRP & Keyword defines the integer Group, well, and connection data
for this time step. Undefined values should be set to zero. & Required &
& \\
1-1 & Format & IGRP & NSGRPZ x NGMAXZ & INTE & \\
2-1 to 2-NWGMAX & Data & Integer group data array IGRP(NIGRPZ, NGMAXZ)
where NIGRPZ and NGMAXZ are defined on the INTEHEAD keyword. Entries 1
to NWGMAX define the index of each well in the group, if this is a well
group, or the index of each child group in this group if it is a node
group. Note that undefined items in this array may be set to zero. & &
& \\
2 - (NWGMAX + 1) & & & & & \\
2 - (NWGMAX + 2) & & & & & \\
2 - (NWGMAX + 6) & & & & & \\
2 - (NWGMAX + 7) & & & & & \\
2 - (NWGMAX + 8) & & & & & \\
2 - (NWGMAX + 11) & & & & & \\
2 - (NWGMAX + 17) & & & & & \\
2 - (NWGMAX + 18) & & & & & \\
2 - (NWGMAX + 22) & & & & & \\
2 - (NWGMAX + 23) & & & & & \\
2 - (NWGMAX + 27) & & & & & \\
2 - (NWGMAX + 28) & & & & & \\
2 - (NWGMAX + 29) & & & & & \\
2 - (NWGMAX + 34) & & & & & \\
2 - (NWGMAX + 89) & & & & & \\
& SGRP & Keyword defines the real group data array SGRP(NIGRPZ, NGMAXZ)
where NIGRPZ and NGMAXZ are defined on the INTEHEAD keyword. & Required
& & \\
1-1 & Format & SGRP & NSGRPZ x NGMAXZ & REAL & \\
2-1 to 2-6 & & Undefined. & & & \\
2-7 & Data & Group's oil production target/limit. & & & \\
2-8 & Data & Group's water production target/limit. & & & \\
2-9 & Data & Group's gas production target/limit. & & & \\
2-10 & Data & Group's liquid production target/limit. & & & \\
2-11 & Data & Group's oil surface volume injection rate target/limit. &
& & \\
2-12 & Data & Group's oil reservoir volume injection rate target/limit.
& & & \\
2-13 & Data & Group's oil re-injection fraction target/limit. & & & \\
2-14 & Data & Group's oil voidage injection fraction target/limit. & &
& \\
2-15 & & Undefined. & & & \\
2-16 & Data & Group's water surface volume injection rate target/limit.
& & & \\
2-17 & Data & Group's water reservoir volume injection rate
target/limit. & & & \\
2-18 & Data & Group's water re-injection fraction target/limit. & & & \\
2-19 & Data & Group's water voidage injection fraction target/limit. & &
& \\
2-20 & & Undefined. & & & \\
2-21 & Data & Group's gas surface volume injection rate target/limit. &
& & \\
2-22 & Data & Group's gas reservoir volume injection rate target/limit.
& & & \\
2-23 & Data & Group's gas re-injection fraction target/limit. & & & \\
2-24 & Data & Group's gas voidage injection fraction target/limit. & &
& \\
& XGRP & Keyword defines the double precision group data array
XGRP(NIGRPZ, NGMAXZ) where NIGRPZ and NGMAXZ are defined on the INTEHEAD
keyword. & Required & & \\
1-1 & Format & XGRP & NSGRPZ x NGMAXZ & DOUB & \\
2-1 & Data & Group's oil production rate. & & & \\
2-2 & Data & Group's water production rate. & & & \\
2-3 & Data & Group's gas production rate. & & & \\
2-4 & Data & Group's liquid production rate. & & & \\
2-5 & & Undefined. & & & \\
2-6 & Data & Group's water injection rate. & & & \\
2-7 & Data & Group's gas injection rate. & & & \\
2-8 & & Undefined. & & & \\
2-9 & Data & Group's producing water cut. & & & \\
2-10 & Data & Group's producing gas/oil ratio. & & & \\
2-11 & Data & Group's total cumulative oil production. & & & \\
2-12 & Data & Group's total cumulative water production. & & & \\
2-13 & Data & Group's total cumulative gas production. & & & \\
2-14 & Data & Group's total cumulative reservoir voidage production. & &
& \\
2-15 & Data & Undefined. & & & \\
2-16 & Data & Group's total cumulative water injection. & & & \\
2-17 & Data & Group's total cumulative gas injection. & & & \\
2-18 & Data & Group's total cumulative reservoir volume injection. & &
& \\
2-19 & & Undefined. & & & \\
2-20 & & Undefined. & & & \\
2-21 & & Undefined. & & & \\
2-22 & & Undefined. & & & \\
2-23 & Data & Group's oil production potential. & & & \\
2-24 & Data & Group's water production potential. & & & \\
2-25 to 2-86 & & Undefined. & & & \\
2-86 & Data & Group's producer guide rate for oil. & & & \\
2-87 & Data & Group's producer guide rate for water. & & & \\
2-88 & Data & Group's producer guide rate for gas. & & & \\
2-89 & Data & Group's producer guide rate for reservoir voidage volume.
& & & \\
2-90 & Data & Group's injection guide rate for oil. & & & \\
2-91 & & Undefined. & & & \\
2-92 & Data & Group's injection guide rate for water. & & & \\
2-93 & Data & Second copy of group's injection guide rate for water. Not
fully characterized. & & & \\
2-94 & Data & Groups injection guide rate for gas. & & & \\
2-95 to 2-127 & & Undefined. & & & \\
2-128 & Data & Second copy of group's producer guide rate for oil. Not
fully characterized. & & & \\
2-129 & Data & Second copy of group's producer guide rate for water. Not
fully characterized. & & & \\
2-130 & Data & Second copy of group's producer guide rate for gas. Not
fully characterized. & & & \\
2-131 & Data & Second copy of group's producer guide rate for reservoir
voidage volume. Not fully characterized. & & & \\
2-132 to 2-135 & & Undefined. & & & \\
2-136 & Data & Group's total cumulative oil production
(observed/historical rates). & & & \\
2-137 & & Undefined. & & & \\
2-138 & & Undefined. & & & \\
2-139 & & Undefined. & & & \\
2-140 & Data & Group's total cumulative water production
(observed/historical rates). & & & \\
2-141 & Data & Group's total cumulative water injection
(observed/historical rates). & & & \\
2-142 & & Undefined. & & & \\
2-143 & & Undefined. & & & \\
2-144 & Data & Group's total cumulative gas production
(observed/historical rates). & & & \\
2-145 & Data & Group's total cumulative gas injection
(observed/historical rates). & & & \\
& ZGRP & Keyword defines the character GROUP data for this keyword. &
Required & & \\
1-1 & Format & ZGRP & NSGRPZ x NGMAXZ & CHAR & \\
2-1 & Data & Character group data array ZGRP(NIGRPZ, NGMAXZ) where
NIGRPZ and NGMAXZ are defined on the INTEHEAD keyword. & & & \\
\end{longtable}

\emph{Table A.22: RESTART Data -- Group, Well, and Connection Keywords
(Groups)}

The next table (Table A.23) contains the multi-segment well keywords
that consist of the ISEG, RSEG, ILBS, ILBR, and ICRD keywords.

\begin{longtable}[]{@{}
  >{\raggedright\arraybackslash}p{(\linewidth - 10\tabcolsep) * \real{0.1667}}
  >{\raggedright\arraybackslash}p{(\linewidth - 10\tabcolsep) * \real{0.1667}}
  >{\raggedright\arraybackslash}p{(\linewidth - 10\tabcolsep) * \real{0.1667}}
  >{\raggedright\arraybackslash}p{(\linewidth - 10\tabcolsep) * \real{0.1667}}
  >{\raggedright\arraybackslash}p{(\linewidth - 10\tabcolsep) * \real{0.1667}}
  >{\raggedright\arraybackslash}p{(\linewidth - 10\tabcolsep) * \real{0.1667}}@{}}
\toprule\noalign{}
\begin{minipage}[b]{\linewidth}\raggedright
No.
\end{minipage} & \begin{minipage}[b]{\linewidth}\raggedright
keyword Name
\end{minipage} & \begin{minipage}[b]{\linewidth}\raggedright
RESTART Data -- Group, Well, and Connection keywords Table A.23 -
Multi-Segment Well Keywords
\end{minipage} & \begin{minipage}[b]{\linewidth}\raggedright
Status or Value
\end{minipage} & \begin{minipage}[b]{\linewidth}\raggedright
\end{minipage} & \begin{minipage}[b]{\linewidth}\raggedright
\end{minipage} \\
\midrule\noalign{}
\endhead
\bottomrule\noalign{}
\endlastfoot
Keyword & No.~of Entries & Data Type & & & \\
& ISEG & Integer multi-segment well data array ISEG(NISEGZ, NSEGMX,
NSWLMX) where NISEGZ, NSEGMX, and NSWLMX are defined on the INTEHEAD
keyword. & Optional & & \\
1-1 & Format & ISEG & NISEGZ x NSEGMX x NSWLMX & INTE & \\
2-1 & Data & Segment number (one-based). Special ordering -- see code
for details. & & & \\
2-2 & Data & Outlet segment (one-based). Defines the multi-segment
outlet segment number starting with 0 for the segment nearest wellhead
(NISEGZ = 2). & & & \\
2-3 & Data & Inflow segment current branch (one-based). & & & \\
2-4 & Data & Branch number (one-based). Defines the multi-segment branch
for this segment number 1 for main stem and 0 if not active (NISEGZ =
4). & & & \\
2-5 to 2-11 & & Undefined. & & & \\
2-12 & Data & Segment type. & & & \\
2-123 to 2-18 & & Undefined. & & & \\
2-19 & Data & ICD scaling mode. & & & \\
2-20 & Data & ICD open/shut flag. & & & \\
& RSEG & Double precision multi-segment well data array RSEG(NISEGZ,
NSEGMX, NSWLMX) where NISEGZ, NSEGMX, and NSWLMX are defined on the
INTEHEAD keyword. & Optional & & \\
1-1 & Format & RSEG & NISEGZ x NSEGMX x NSWLMX & DOUB & \\
2-1 & Data & Segment's distance to outlet. & & & \\
2-2 & Data & Segment's depth differential to outlet. & & & \\
2-3 & Data & Internal diameter of segment. & & & \\
2-4 & Data & Roughness parameter of segment. & & & \\
2-5 & Data & Cross-sectional area of segment. & & & \\
2-6 & Data & Physical volume of segment. & & & \\
2-7 & Data & Segment's distance to BHP reference node. & & & \\
2-8 & Data & Segment's depth differential to BHP reference node. & &
& \\
2-9 & Data & Normalized total segment flow rate. & & & \\
2-10 & Data & Normalized water flow rate fraction. & & & \\
2-11 & Data & Normalized gas flow rate fraction. & & & \\
2-12 & Data & Segment pressure. & & & \\
2-13 to 2-30 & & Undefined. & & & \\
2-31 & Data & Very close to normalized water flow rate fraction. & &
& \\
2-32 to 2-39 & & Undefined. & & & \\
2-40 & Data & Undefined. & & & \\
2-41 & Data & Length of valve. & & & \\
2-42 & Data & Cross-sectional area of valve. & & & \\
2-43 & Data & Valve's dimensionless flow coefficient. & & & \\
2-44 & Data & Maximal cross-sectional valve area. & & & \\
2-45 to 2-86 & & Undefined. & & & \\
2-87 & Data & Device base strength. & & & \\
2-88 & & Undefined. & & & \\
2-89 & Data & Calibrated fluid density. & & & \\
2-90 & Data & Calibrated fluid viscosity. & & & \\
2-91 & Data & Critical water fraction. & & & \\
2-92 & Data & Transition region width. & & & \\
2-93 & Data & Maximum emulsion ratio. & & & \\
2-94 to 2-97 & & Undefined. & & & \\
2-98 & Data & Maximum valid flow rate. & & & \\
2-99 to 2-102 & & Undefined. & & & \\
2-103 & Data & ICD length. & & & \\
2-104 & Data & Valve area fraction. & & & \\
2-105 to 2-111 & & Undefined. & & & \\
& ILBS & Integer multi-segment well data array for lateral branches
ILBR(NLBRMX, NSWLMX) where NLBRMX and NSWLMX are defined on the INTEHEAD
keyword. & Optional & & \\
1-1 & Format & ILBS & NLBRMX x NSWLMX & INTE & \\
2-1 to no\_branches & Data & The segment number of the first segment in
the branch for branch number two and upwards. & & & \\
& ILBR & Integer multi-segment well data array for chords ILBR(NILBRZ,
NLBRMX, NSWLMX) where NILBRZ, NLBRMX, and NSWLMX are defined on the
INTEHEAD keyword. & Optional & & \\
1-1 & Format & ILBR & NILBRZ x NLBRMX x NSWLMX & INTE & \\
2-1 & Data & Outlet segment number. & & & \\
2-2 & Data & Number of segments in the branch. & & & \\
2.3 & Data & First segment in branch. & & & \\
2-4 & Data & Last segment in branch. & & & \\
2-5 & Data & Branch number minus one. & & & \\
& ICRD & Integer multi-segment well data array for chords ICRD(2,
NCRDMX, NSWLMX) where NCRDMX, and NSWLMX are defined on the INTEHEAD
keyword. & Optional & & \\
1-1 & Format & ILBR & 2 x NCRDMX x NSWLMX & INTE & \\
2-1 & Data & ICRD (2, NCRDMX, NSWLMX) array. & & & \\
\end{longtable}

\emph{Table A.23: RESTART Data -- Group, Well, and Connection Keywords
(Multi-Segment Wells)}

The wells keywords are IWEL, SWEL, XWEL, and ZWEL and Table A.24 below
describes these keywords in detail.

\begin{longtable}[]{@{}
  >{\raggedright\arraybackslash}p{(\linewidth - 10\tabcolsep) * \real{0.1667}}
  >{\raggedright\arraybackslash}p{(\linewidth - 10\tabcolsep) * \real{0.1667}}
  >{\raggedright\arraybackslash}p{(\linewidth - 10\tabcolsep) * \real{0.1667}}
  >{\raggedright\arraybackslash}p{(\linewidth - 10\tabcolsep) * \real{0.1667}}
  >{\raggedright\arraybackslash}p{(\linewidth - 10\tabcolsep) * \real{0.1667}}
  >{\raggedright\arraybackslash}p{(\linewidth - 10\tabcolsep) * \real{0.1667}}@{}}
\toprule\noalign{}
\begin{minipage}[b]{\linewidth}\raggedright
No.
\end{minipage} & \begin{minipage}[b]{\linewidth}\raggedright
keyword Name
\end{minipage} & \begin{minipage}[b]{\linewidth}\raggedright
RESTART Data -- Group, Well, and Connection keywords Table A.24 - Well
Keywords
\end{minipage} & \begin{minipage}[b]{\linewidth}\raggedright
Status or Value
\end{minipage} & \begin{minipage}[b]{\linewidth}\raggedright
\end{minipage} & \begin{minipage}[b]{\linewidth}\raggedright
\end{minipage} \\
\midrule\noalign{}
\endhead
\bottomrule\noalign{}
\endlastfoot
Keyword & No.~of Entries & Data Type & & & \\
& IWEL & Well integer data array IWEL(NIWELZ, NWELLS) where NIWELX and
NWELS are defined on the INTEHEAD keyword. Undefined values should be
set to zero. & & & \\
1-1 & Format & IWEL & NIWELZ x NWELLS & INTE & Required \\
2-1 & Data & Defines the location of the wellhead in the grid for the I
dimension (NIWELZ = 1). & & & \\
2-2 & Data & Defines the location of the wellhead in the grid for the J
dimension (NIWELZ = 2). & & & \\
2-3 & Data & Layer ID (one-based) of top/first connection. & & & \\
2-4 & Data & Layer ID (one-based) of bottom/last connection. & & & \\
2-5 & Data & Number of active cells connected to a well. & & & \\
2-6 & Data & Index (one-based) of well's current group. & & & \\
2-7 & Data & Defines the well type, set 1 for a production well, 2 for
an oil injection well, 3 for a water injection well, or 4 for a gas
injection well (NIWELZ = 7). & & & \\
2-8 & Data & Well's active target control mode (constraint). & & & \\
2-9 & & Undefined. & & & \\
2-10 & & Undefined. & & & \\
2-11 & Data & Defines the well status, set to greater than zero for
producing and less than or equal to zero for being shut (NIWELZ =11). &
& & \\
2-12 & Data & ID (one-based) of well's current VFP table. & & & \\
2-13 & & Undefined. & & & \\
2-14 & & Undefined. & & & \\
2-15 & & Undefined. & & & \\
2-16 & Data & Well's requested control mode from simulation deck
(WCONINJE, WCONPROD). & & & \\
2-17 to 2-22 & & Undefined. & & & \\
2-23 & Data & Cross flow indicator. & & & \\
2-24 to 2-42 & & Undefined. & & & \\
2-43 & Data & Defines the LGR index for a well with local completions in
a LGR (NIWELZ = 43). & & & \\
2-44 to 2-48 & & Undefined. & & & \\
2-49 & Data & Defines the well friction indicator, set to non-zero for
horizontal wells (NIWELZ = 49). & & & \\
2-50 & Data & Well's requested control mode from simulation deck
(WCONHIST, WCONINJH). & & & \\
2-51 to 2-70 & & Undefined. & & & \\
2-71 & Data & Multi-segment well ID; value of one for regular wells, or
zero for multi-segment wells. & & & \\
2-72 & Data & Number of well segments: value of one for regular wells,
or the number of segments for multi-segment wells. & & & \\
2-73 to 2-98 & & Undefined. & & & \\
2-99 & Data & Well's completion ordering scheme, set to: & & & \\
& SWEL & Well real data array SWEL(NSWELZ, NWELLS) where NSWELX and
NWELS are defined on the INTEHEAD record. & & & \\
1-1 & Format & SWEL & NSWELZ x NWELLS & REAL & Required \\
2-1 & Data & Well's current oil rate production target. & & & \\
2-2 & Data & Well's current water rate production target. & & & \\
2-3 & Data & Well's current gas rate production target. & & & \\
2-4 & Data & Well's current liquid rate production target. & & & \\
2-5 & Data & Well's current reservoir voidage rate production target. &
& & \\
2-6 & Data & Well's tubing head pressure target. & & & \\
2-7 & Data & Well's bottom hole pressure target. & & & \\
2-8 & & Undefined. & & & \\
2-9 & & Undefined. & & & \\
2-10 & & Well's reference depth for BHP. & & & \\
2-11 & & Well's artificial lift quantity. & & & \\
2-12 to 2-17 & & Undefined. & & & \\
2-18 & & Well's drainage radius - item seven from the WELSPECS keyword.
& & & \\
2-19 to 2-20 & & Undefined. & & & \\
2-25 & & Item two from WEFAC; this value is repeated at two locations. &
& & \\
2-26 to 2-31 & & Undefined. & & & \\
2-32 & & Item two from WEFAC keyword. & & & \\
2-33 & & Undefined. & & & \\
2-34 & & Well's historical/observed liquid rate target/limit. & & & \\
2-35 to 2-54 & & Undefined. & & & \\
2-55 & & Well's historical/observed gas rate target/limit. & & & \\
2-56 & & Well's historical/observed bottom hole pressure target/limit. &
& & \\
& XWEL & Well's injection and production data array XWEL(NXWELZ,
NWELLS). & & & \\
1-1 & Format & XWEL & NSWELZ x NWELLS & DOUB & Required \\
2-1 & Data & Well's oil production rate. & & & \\
2-2 & Data & Well's water production rate. & & & \\
2-3 & Data & Well's gas production rate. & & & \\
2-4 & Data & Well's liquid production rate. & & & \\
2-5 & Data & Well's reservoir voidage production rate. & & & \\
2-6 & Data & Well's tubing head pressure. & & & \\
2-7 & Data & Well's flowing/producing bottom hole pressure. & & & \\
2-8 & Data & Well's producing water cut. & & & \\
2-9 & Data & Well's producing gas/oil ratio. & & & \\
2-10 to 2-18 & & Undefined. & & & \\
2-19 & Data & Well's total cumulative oil production. & & & \\
2-20 & Data & Well's total cumulative water production. & & & \\
2-21 & Data & Well's total cumulative gas production. & & & \\
2-22 & Data & Well's total cumulative reservoir voidage production. & &
& \\
2-23 & & Undefined. & & & \\
2-24 & Data & Well's total cumulative water injection. & & & \\
2-25 & Data & Well's total cumulative gas injection. & & & \\
2-26 & Data & Well's total cumulative reservoir volume injection. & &
& \\
2-27 to 2-34 & & Undefined. & & & \\
2-35 & Data & Well's producing gas formation volume factor. & & & \\
2-36 & & Undefined. & & & \\
2-37 & & Undefined. & & & \\
2-38 & & Undefined. & & & \\
2-39 to 2-41 & & Undefined. & & & \\
2-42 & Data & Well's current BHP Target/Limit. & & & \\
2-43 to 2-48 & & Undefined. & & & \\
2-49 & Data & Well's ``primary'' guide rate (oil for producers,
preferred phase for injectors). & & & \\
2-50 & Data & Well's producer guide rate for water. & & & \\
2-51 & Data & Well's producer guide rate for gas. & & & \\
2-52 to 2-68 & Data & Undefined. & & & \\
2-69 & Data & Well's producer guide rate for reservoir voidage volume. &
& & \\
2-70 to 2-75 & & Undefined. & & & \\
2-76 & Data & Well's total cumulative oil production
(observed/historical rates). & & & \\
2-77 & Data & Well's total cumulative water production
observed/historical rates). & & & \\
2-78 & Data & Well's total cumulative gas production
(observed(historical rates). & & & \\
2-79 to 2-81 & & Undefined. & & & \\
2-82 & Data & Well's total cumulative water injection
(observed/historical rates). & & & \\
2-83 & Data & Well's total cumulative gas injection (observed/historical
rates). & & & \\
2-84 to 2-91 & & Undefined. & & & \\
2-92 & & Second copy of well's primary guide rate. Not fully
characterized. & & & \\
2-93 & Data & Second copy of well's producer guide rate for water. Not
fully characterized. & & & \\
2-94 & Data & Second copy of well's producer guide rate for gas. Not
fully characterized. & & & \\
2-95 & Data & Second copy of well's producer guide rate for reservoir
voidage. Not fully characterized. & & & \\
2-96 to 2-122 & & Undefined. & & & \\
2-123 & Data & Well's voidage production rate. & & & \\
2-124 & Data & Well's voidage production rate. & & & \\
& ZWEL & Well character data array ZWEL(NXWELZ, NWELLS) where NXWELX and
NWELS are defined on the INTEHEAD keyword. All strings are eight
characters in length. & & & \\
1-1 & Format & ZWEL & NSWELZ x NWELLS & CHAR & Required \\
2-1 & Data & Defines the well name consisting of eight characters
(NXWELZ = 1). & & & \\
2-2 & Data & Defines the name of the Well List the well belongs to, set
to ``blank'' for the well does not belong to Well List (NXWELZ = 2),
that is the ACTIONX name. & & & \\
2-3 & Data & Defines the end of time step action for the well (NXWELZ =
3). & & & \\
\end{longtable}

\emph{Table A.24: RESTART Data -- Group, Well, and Connection Keywords
(Wells)}

The well connection keywords are specified in Table A.25 and consist of
the ICON, SCON, and XCON keywords.

\begin{longtable}[]{@{}
  >{\raggedright\arraybackslash}p{(\linewidth - 10\tabcolsep) * \real{0.1667}}
  >{\raggedright\arraybackslash}p{(\linewidth - 10\tabcolsep) * \real{0.1667}}
  >{\raggedright\arraybackslash}p{(\linewidth - 10\tabcolsep) * \real{0.1667}}
  >{\raggedright\arraybackslash}p{(\linewidth - 10\tabcolsep) * \real{0.1667}}
  >{\raggedright\arraybackslash}p{(\linewidth - 10\tabcolsep) * \real{0.1667}}
  >{\raggedright\arraybackslash}p{(\linewidth - 10\tabcolsep) * \real{0.1667}}@{}}
\toprule\noalign{}
\begin{minipage}[b]{\linewidth}\raggedright
No.
\end{minipage} & \begin{minipage}[b]{\linewidth}\raggedright
keyword Name
\end{minipage} & \begin{minipage}[b]{\linewidth}\raggedright
RESTART Data -- Group, Well, and Connection keywords Table A.25 -
Connection Keywords
\end{minipage} & \begin{minipage}[b]{\linewidth}\raggedright
Status or Value
\end{minipage} & \begin{minipage}[b]{\linewidth}\raggedright
\end{minipage} & \begin{minipage}[b]{\linewidth}\raggedright
\end{minipage} \\
\midrule\noalign{}
\endhead
\bottomrule\noalign{}
\endlastfoot
Keyword & No.~of Entries & Data Type & & & \\
& ICON & Connection integer data array ICON(NICONZ, NCWMAX, NWELLS)
where NICONZ, NCWMAX, and NWELLS are defined on the INTEHEAD keyword.
Data keywords are required for each completion in the well and the
number of data items, N, should be NICONZ x NCWMAX x NWELLS. Undefined
values should be set to zero. & & & \\
1-1 & Format & ICON & N & INTE & \\
2-1 & Data & Defines the well connect index (NICONZ = 1). & 0 & & \\
2-2 & Data & Defines the location of connection in the grid for the I
dimension (NICONZ = 2). & 0 & & \\
2-3 & Data & Defines the location of connection in the grid for the J
dimension (NICONZ = 3). & 0 & & \\
2-4 & Data & Defines the location of connection in the grid for the K
dimension (NICONZ = 4). & 0 & & \\
2-5 & & Undefined. & 0 & & \\
2-6 & Data & Defines the status of the connection, set to less than or
equal to zero for being shut or greater than zero for being open (NICONZ
= 5). & 0 & & \\
2-7 & & Drainage saturation function table number. & 0 & & \\
2-8 & & Undefined. & 0 & & \\
2-9 & & Undefined. & 0 & & \\
2-10 & & Imbibition saturation function table number. & 0 & & \\
2-11 & & Undefined. & 0 & & \\
2-12 & & Undefined. & 0 & & \\
2-13 & & Completion ID (1-based). & 0 & & \\
2-14 & Data & Defines the direction of connection, set to one for the
x-direction, two for the y-direction, or three for the z-direction. The
default value of zero indicates the z-direction (NICONZ = 14). & 0 &
& \\
2-15 & Data & Defines, for multi-segment wells, the segment number the
connect belongs to, for conventional wells the default value of zero
should be used (NICONZ = 15). & & & \\
& SCON & Connection real data array SCON(NSCONZ, NCWMAX, NWELLS) where
NSCONZ, NCWMAX, and NWELLS are defined on the INTEHEAD keyword. Data are
required for each completion in the well and the number of data items,
N, should be NSCONZ x NCWMAX x NWELLS. & Required & & \\
1-1 & Format & ICON & N & REAL & \\
2-1 & Data & Connection transmissibility factor NSCONZ = 1 in the SCON
array. & & & \\
2-2 & Data & Connection center depth. & & & \\
2-3 & Data & Connection diameter. & & & \\
2-4 & Data & Connection KH, that is the net thickness times the
permeability, NSCONZ = 4 in the SCON array. & & & \\
2-5 & Data & Skin factor - item `SKIN' from COMPDAT keyword in the
SCHEDULE section. & & & \\
2-6 to 2-11 & & Undefined. & & & \\
2-12 & Data & Connection transmissibility factor. & & & \\
2-13 to 2-20 & & Undefined. & & & \\
2-21 & Data & Distance to end of connection in segment. & & & \\
2-22 & Data & Distance to start of connection in segment. & & & \\
2-23 to 2-29 & & Undefined. & & & \\
2-30 & Data & Undefined. & & & \\
2-31 & Data & Undefined. & & & \\
2-32 to 2-40 & & Undefined. & & & \\
2-41 & Data & Connection factor defined in input deck (set to one), or
zero for undefined in input deck. & & & \\
& XCON & Connection double precision data array XCON(NXCONZ, NCWMAX,
NWELLS) where NXCONZ, NCWMAX, and NWELLS are defined on the INTEHEAD
keyword. The number of data items, N, should be NXCONZ x NCWMAX x NWELLS
& Required & & \\
1-1 & Format & ICON & N & DOUB & \\
2-1 & Data & Oil surface rate. & & & \\
2-2 & Data & Water surface rate. & & & \\
2-3 & Data & Gas surface rate. & & & \\
2-4 to 2-34 & & Undefined. & & & \\
2-35 & & Connection pressure value. & & & \\
2-36 to 2-49 & & Undefined. & & & \\
2-50 & & Reservoir voidage rate. & & & \\
\end{longtable}

\emph{Table A.25: RESTART Data -- Group, Well, and Connection Keywords
(Connections)}

The tracer keywords are listed in Table A.26 and consists of the ICOT,
XCOT, IWET, XWET, IGRT, and XGRT keywords. Note that this series of
keywords are only available if tracers are present in the run.

\begin{longtable}[]{@{}
  >{\raggedright\arraybackslash}p{(\linewidth - 10\tabcolsep) * \real{0.1667}}
  >{\raggedright\arraybackslash}p{(\linewidth - 10\tabcolsep) * \real{0.1667}}
  >{\raggedright\arraybackslash}p{(\linewidth - 10\tabcolsep) * \real{0.1667}}
  >{\raggedright\arraybackslash}p{(\linewidth - 10\tabcolsep) * \real{0.1667}}
  >{\raggedright\arraybackslash}p{(\linewidth - 10\tabcolsep) * \real{0.1667}}
  >{\raggedright\arraybackslash}p{(\linewidth - 10\tabcolsep) * \real{0.1667}}@{}}
\toprule\noalign{}
\begin{minipage}[b]{\linewidth}\raggedright
No.
\end{minipage} & \begin{minipage}[b]{\linewidth}\raggedright
keyword Name
\end{minipage} & \begin{minipage}[b]{\linewidth}\raggedright
RESTART Data -- Group, Well, and Connection keywords Table A.26 - Tracer
Keywords
\end{minipage} & \begin{minipage}[b]{\linewidth}\raggedright
Status or Value
\end{minipage} & \begin{minipage}[b]{\linewidth}\raggedright
\end{minipage} & \begin{minipage}[b]{\linewidth}\raggedright
\end{minipage} \\
\midrule\noalign{}
\endhead
\bottomrule\noalign{}
\endlastfoot
Keyword & No.~of Entries & Data Type & & & \\
& ICOT & Tracer connection integer data array ICOT(NICOTZ, NSTRA2,
NCWMAX, NWELLS) where NICOTZ, NSTRA2, NCWMAX, and NWELLS are defined on
the INTEHEAD keyword. The number of data items, N, should be NICONZ x
NSTRA2 x NCWMAX x NWELLS. & Optional & & \\
1-1 & Format & ICOT & N & INTE & \\
2-1 & Data & ICOT(NICOTZ, NSTRA2, NCWMAX, NWELLS) array. & & & \\
& XCOT & Tracer connection double precision data array ICOT(NICOTZ,
NSTRA2, NCWMAX, NWELLS) where NICOTZ, NSTRA2, NCWMAX, and NWELLS are
defined on the INTEHEAD keyword. The number of data items, N, should be
NICONZ x NSTRA2 x NCWMAX x NWELLS. & Optional & & \\
1-1 & Format & XCOT & N & DOUB & \\
2-1 & Data & XCOT(NICOTZ, NSTRA2, NCWMAX, NWELLS) array. & & & \\
& IWET & Tracer well integer data array IWET(NIWETZ, NSTRA2, NWMAXZ,
NWELLS) where NIWETZ, NSTRA2, and NWMAXZ are defined on the INTEHEAD
keyword. The number of items, N, should be NIWETNZ x NSTRA2 x NWMAXZ. &
Optional & & \\
1-1 & Format & IWET & N & INTE & \\
2-1 & Data & IWET(NIWETZ, NSTRA2, NWMAXZ, NWELLS) array. & & & \\
& XWET & Tracer connection double precision data array XWET(NXWETZ,
NSTRA2, NWMAXZ) where NXWETZ, NSTRA2, and NWMAXZ are defined on the
INTEHEAD keyword. The number of items, N, should be NXWETNZ x NSTRA2 x
NWMAXZ. & Optional & & \\
1-1 & Format & XWET & N & DOUB & \\
2-1 & Data & XWET(NXWETZ, NSTRA2, NWMAXZ) array. & & & \\
& IGRT & Tracer well integer data array IGRT(NIGRTZ, NSTRA2, NWMAXZ,
NWELLS) where NIGRTZ, NSTRA2, and NWMAXZ are defined on the INTEHEAD
keyword. The number of elements, N, should be NIGRTNZ x NSTRA2 x NWMAXZ.
& Optional & & \\
1-1 & Format & IGRT & N & INTE & \\
2-1 & Data & IGRT(NIGRTZ, NSTRA2, NWMAXZ, NWELLS) array. & & & \\
& XGRT & Tracer connection double precision data array XGRT(NXGRTZ,
NSTRA2, NWMAXZ) where NXGRTZ, NSTRA2, and NWMAXZ are defined in the
INTEHEAD keyword. The number of data items, N, should be NXGRTNZ x
NSTRA2 x NWMAXZ. & Optional & & \\
1-1 & Format & XGRT & N & DOUB & \\
2-1 & Data & XWET(NXWETZ, NSTRA2, NWMAXZ) array. & & & \\
\end{longtable}

\emph{Table A.26: RESTART Data -- Group, Well, and Connection Keywords
(Tracers)}

Finally, the network keywords INODE, IBRAN, INOBR, RNODE, RBRAN, and
ZNODE are specified in Table A.27. Again, this series of keywords are
only available if networks are present in the run.

\begin{longtable}[]{@{}
  >{\raggedright\arraybackslash}p{(\linewidth - 10\tabcolsep) * \real{0.1667}}
  >{\raggedright\arraybackslash}p{(\linewidth - 10\tabcolsep) * \real{0.1667}}
  >{\raggedright\arraybackslash}p{(\linewidth - 10\tabcolsep) * \real{0.1667}}
  >{\raggedright\arraybackslash}p{(\linewidth - 10\tabcolsep) * \real{0.1667}}
  >{\raggedright\arraybackslash}p{(\linewidth - 10\tabcolsep) * \real{0.1667}}
  >{\raggedright\arraybackslash}p{(\linewidth - 10\tabcolsep) * \real{0.1667}}@{}}
\toprule\noalign{}
\begin{minipage}[b]{\linewidth}\raggedright
No.
\end{minipage} & \begin{minipage}[b]{\linewidth}\raggedright
keyword Name
\end{minipage} & \begin{minipage}[b]{\linewidth}\raggedright
RESTART Data -- Group, Well, and Connection keywords Table A.27 -
Network Keywords
\end{minipage} & \begin{minipage}[b]{\linewidth}\raggedright
Status or Value
\end{minipage} & \begin{minipage}[b]{\linewidth}\raggedright
\end{minipage} & \begin{minipage}[b]{\linewidth}\raggedright
\end{minipage} \\
\midrule\noalign{}
\endhead
\bottomrule\noalign{}
\endlastfoot
Keyword & No.~of Entries & Data Type & & & \\
& INODE & Network node integer data array INODE(NINODE, NODMAX) where
NINODE and NODMAX are defined in the INTEHEAD keyword. & Optional & & \\
1-1 & Format & INODE & NINODE x NODMAX & INTE & \\
2-1 & Data & INODE(NINODE, NODMAX) array & & & \\
& IBRAN & Network node branch integer data array IBRAN(NIBRAN, NBRMAX)
where NIBRAN and NBRMAX are defined in the INTEHEAD keyword. & Optional
& & \\
1-1 & Format & IBRAN & NIBRAN x NBRMAX & INTE & \\
2-1 & Data & IBRAN(NIBRAN, NBRMAX) & & & \\
& INOBR & Network node-branch integer data array INOBR(NIOBR) where
NIOBR are defined in the INTEHEAD keyword. & Optional & & \\
1-1 & & INOBR & NIOBR & INTE & \\
2-1 & & INOBR(NIOBR) array. & & & \\
& RNODE & Network node double precision data array RNODE(NRNODE, NODMAX)
where NRNODE and NODMAX are defined in the INTEHEAD keyword. & Optional
& & \\
1-1 & Format & RNODE & NRNODE x NODMAX & DOUB & \\
2-1 & Data & RNODE(NRNODE, NODMAX) array. & & & \\
& RBRAN & Network node branch integer data array RBRAN(NRBRAN, NBRMAX)
where NRBRAN and NBRMAX are defined in the INTEHEAD keyword. & Optional
& & \\
1-1 & Format & RBRAN & NRBRAN x NBRMAX & DOUB & \\
2-1 & Data & RBRAN(NRBRAN, NBRMAX) array. & & & \\
& ZNODE & Network node character data array ZNODE(NZNODE, NODMAX) where
NZNODE and NODMAX are defined in the INTEHEAD keyword. & Optional & & \\
1-1 & Format & ZNODE & NZNODE x NODMAX & CHAR & \\
2-1 & Data & ZNODE(NZNODE, NODMAX) array. & & & \\
Notes: & & & & & \\
\end{longtable}

\emph{Table A.27: RESTART Data -- Group, Well, and Connection Keywords
(Network)}

\subsubsection{Example: RESTART Data -- Group, Well, and Connection
Keywords}\label{example-restart-data-group-well-and-connection-keywords}

The following example shows a typical formatted RESTART Data -- Group,
Well, and Connection keyword data set, showing the IGRP, SGRP, XGRP, and
ZGRP keywords.

\begin{verbatim}
 'IGRP    '        2864 'INTE'
          1           0           0           0           0           0
          0           0           0           0           0           0
          0           0           0           0           0           0
          0           0           0           0           0           0
 ………………………….……………..
          0           0           0           0           0           0
          0           0           0           0           0           0
          0           0           0           0           0           0
          0           0           0           0           0          16
         16           0           0           0           0           0
'SGRP    '        1488 'REAL'
  0.10000000E+21   0.10000000E+21  -0.10000000E+21   0.10000000E+21
  0.00000000E+00   0.00000000E+00   0.10000000E+21   0.10000000E+21
  0.10000000E+21   0.10000000E+21   0.10000000E+21   0.10000000E+21
 ………………………….……………..
  0.00000000E+00   0.00000000E+00   0.10000000E+21   0.10000000E+21
  0.00000000E+00   0.10000000E+21   0.00000000E+00   0.00000000E+00
  0.00000000E+00   0.00000000E+00   0.00000000E+00   0.00000000E+00
  0.00000000E+00   0.00000000E+00   0.00000000E+00   0.10000000E+01
'XGRP    '        2784 'DOUB'
  0.00000000000000D+00   0.00000000000000D+00   0.00000000000000D+00
  0.00000000000000D+00   0.00000000000000D+00   0.00000000000000D+00
  0.00000000000000D+00   0.00000000000000D+00   0.00000000000000D+00
 ………………………….……………..
  0.00000000000000D+00   0.00000000000000D+00   0.00000000000000D+00
  0.00000000000000D+00   0.00000000000000D+00   0.00000000000000D+00
  0.00000000000000D+00   0.00000000000000D+00   0.00000000000000D+00
  0.00000000000000D+00
'ZGRP    '          80 'CHAR'
'MANI-C  ' '        ' '        ' '        ' '        ' 'B1-DUMMY' '        '
'        ' '        ' '        ' 'MANI-D1 ' '        ' '        ' '        '
'        ' 'INJE    ' '        ' '        ' '        ' '        ' 'PROD    '
'        ' '        ' '        ' '        ' 'MANI-B2 ' '        ' '        '
'        ' '        ' 'MANI-B1 ' '        ' '        ' '        ' '        '
'MANI-D2 ' '        ' '        ' '        ' '        ' 'MANI-E1 ' '        '
'        ' '        ' '        ' 'MANI-E2 ' '        ' '        ' '        '
'        ' 'MANI-K1 ' '        ' '        ' '        ' '        ' 'MANI-K2 '
'        ' '        ' '        ' '        ' 'MANI-F  ' '        ' '        '
'        ' '        ' 'WI-GSEG ' '        ' '        ' '        ' '        '
'D2-DUMMY' '        ' '        ' '        ' '        ' 'FIELD   ' '        '
'        ' '        ' '        '
\end{verbatim}

End of Example

The restart data for the UDQ and ACTIONX keywords are written out in a
set of UDQ and ACTIONX arrays. If there are no UDQ or ACTIONX data, the
corresponding arrays are not written to the restart file and only the
relevant arrays are written out.

The structure for this type of keyword is based on the data type as
summarized below:

\begin{verbatim}
- The UDQ data set keywords are specified in Table A.28 and consists of the IUDQ, IUAD, ZUDN, ZUDL, IGPH, IUAP, DUDW, DUDG, and DUDF keywords.
- The ACTIONX data set keywords are specified in Table A.29 and consists of the keywords IACT, SACT, ZACT, ZLACT, ZACN, IACN, and SACN.
\end{verbatim}

\begin{longtable}[]{@{}
  >{\raggedright\arraybackslash}p{(\linewidth - 10\tabcolsep) * \real{0.1667}}
  >{\raggedright\arraybackslash}p{(\linewidth - 10\tabcolsep) * \real{0.1667}}
  >{\raggedright\arraybackslash}p{(\linewidth - 10\tabcolsep) * \real{0.1667}}
  >{\raggedright\arraybackslash}p{(\linewidth - 10\tabcolsep) * \real{0.1667}}
  >{\raggedright\arraybackslash}p{(\linewidth - 10\tabcolsep) * \real{0.1667}}
  >{\raggedright\arraybackslash}p{(\linewidth - 10\tabcolsep) * \real{0.1667}}@{}}
\toprule\noalign{}
\begin{minipage}[b]{\linewidth}\raggedright
No.
\end{minipage} & \begin{minipage}[b]{\linewidth}\raggedright
keyword Name
\end{minipage} & \begin{minipage}[b]{\linewidth}\raggedright
RESTART Data -- UDQ and ACTIONX keywords Table A.28 - UDQ Keywords
\end{minipage} & \begin{minipage}[b]{\linewidth}\raggedright
Status or Value
\end{minipage} & \begin{minipage}[b]{\linewidth}\raggedright
\end{minipage} & \begin{minipage}[b]{\linewidth}\raggedright
\end{minipage} \\
\midrule\noalign{}
\endhead
\bottomrule\noalign{}
\endlastfoot
Keyword & No.~of Entries & Data Type & & & \\
& IUDQ & Keyword defines the integer properties as per the UDQ variable.
& Required & & \\
1-1 & Format & IUDQ & No.~UDQ's x 3 & INTE & \\
2-1 & Data & & & & \\
2-2 & Data & -4 : ASSIGN UDQ (-1) times the number of operators in
expression: DEFINE UDQ (see code for details) & & & \\
2-3 & Data & & & & \\
& IUAD & This keyword defines the various controls for wells, groups,
etc. & Required & & \\
1-1 & Format & IUAD & No.~IUAD's x 5 & INTE & \\
2-1 & Data & Set to: & & & \\
2-2 & Data & Sequence number of UDQ used (from input sequence) for the
actual constraint/target. & & & \\
2-3 & Data & Set to one (1). & & & \\
2-4 & Data & Number of times the UDQ variable is used. & & & \\
2-5 & Data & The sequence number for the first use of the actual UDQ,
equal to one (1) plus the (sum over the earlier UDQ's of the number of
uses of the respective UDQ's). & & & \\
& IUAP & The keyword defines the sequence number of the actual well,
group, etc. for each UDQ use in WCONPROD, GCONPROD, GCONINJE, etc. &
Required & & \\
1-1 & Format & IUAP & No.~IUAP's & INTE & \\
2-1 & Data & Equals the well, group, etc. sequence number for the
UDQ-Variable used in actual target, limit. & & & \\
& IGPH & This keyword defines an integer number per group, and for the
FIELD which is the last group in the list. & Required & & \\
1-1 & Format & IGPH & Max Number Groups in the Field plus one. & INTE
& \\
2-1 & Data & The value of the integer denotes the phase (or component)
for injection group control (GCONINJE): & & & \\
& ZUDN & Keyword pairs of data for each UDQ. & Required & & \\
1-1 & Format & ZUDN & No.~UDQ's x 2 & CHAR & \\
2-1 & Data & UDQ name (maximum of eight characters per UDQ). & & & \\
2-2 & Data & UDQ's units (maximum of eight characters per UDQ). & & & \\
& ZUDL & This keyword contains the EXPRESSION variable (item three) on
the UDQ keyword in the input deck ``Data for operation''. & Required &
& \\
1-1 & Format & ZUDL & No.~UDQ's x 16 & CHAR & \\
2-1 to 2-16 & Data & A character string that defines the EXPRESSION
variable, ``Data for operation'', split into strings of eight characters
each. & & & \\
& DUDW & Keyword containing the values of all the well UDQ's. & Required
& & \\
1-1 & Format & DUDW & No.~Well UDQs x Max No.~of Wells. & DOUBLE & \\
2-1 to no of well UDQ's & Data & The value of the UDQ for the actual
well (repeated for all wells in the well input sequence). & & & \\
& DUDG & Keyword containing the values of all the group UDQ's. &
Required & & \\
1-1 & Format & DUDG & No.~Group\_UDQs x Max No.~of Groups & DOUBLE & \\
2-1 to no of group UDQ's & Data & The value of UDQ for the actual group
(repeated for all groups in the group input sequence). & & & \\
& DUDF & Keyword containing the values of all the field UDQ's. &
Required & & \\
1-1 & Format & DUDF & No.~Field UDQs & DOUBLE & \\
2-1 to no of field UDQ's & Data & Value of UDQ for field UDQs. & & & \\
\end{longtable}

\emph{Table A.28: RESTART Data -- UDQ Keywords}

\begin{longtable}[]{@{}
  >{\raggedright\arraybackslash}p{(\linewidth - 10\tabcolsep) * \real{0.1667}}
  >{\raggedright\arraybackslash}p{(\linewidth - 10\tabcolsep) * \real{0.1667}}
  >{\raggedright\arraybackslash}p{(\linewidth - 10\tabcolsep) * \real{0.1667}}
  >{\raggedright\arraybackslash}p{(\linewidth - 10\tabcolsep) * \real{0.1667}}
  >{\raggedright\arraybackslash}p{(\linewidth - 10\tabcolsep) * \real{0.1667}}
  >{\raggedright\arraybackslash}p{(\linewidth - 10\tabcolsep) * \real{0.1667}}@{}}
\toprule\noalign{}
\begin{minipage}[b]{\linewidth}\raggedright
No.
\end{minipage} & \begin{minipage}[b]{\linewidth}\raggedright
keyword Name
\end{minipage} & \begin{minipage}[b]{\linewidth}\raggedright
RESTART Data -- UDQ and ACTIONX keywords Table A.29 - ACTIONX Keywords
\end{minipage} & \begin{minipage}[b]{\linewidth}\raggedright
Status or Value
\end{minipage} & \begin{minipage}[b]{\linewidth}\raggedright
\end{minipage} & \begin{minipage}[b]{\linewidth}\raggedright
\end{minipage} \\
\midrule\noalign{}
\endhead
\bottomrule\noalign{}
\endlastfoot
Keyword & No.~of Entries & Data Type & & & \\
& IACT & The keyword defines the integer properties for the ACTIONX
variable. & Required & & \\
1-1 & Format & IACT & No.~of Actions x 9 & INTE & \\
2-1 & Data & & & & \\
2-2 & Data & Number of lines of SCHEDULE data. & & & \\
2-3 & Data & & & & \\
2-4 & Data & & & & \\
2-5 & Data & & & & \\
2-6 & Data & & & & \\
2-7 & Data & & & & \\
2-8 & Data & & & & \\
2-9 & Data & & & & \\
& IACN & Keyword defines various parameters for & Required & & \\
1-1 & Format & IACN & No.~Actions x 26 x Max No.~Conditions per Action.
& INTE & \\
2-1 to 2-10 & Data & & & & \\
2-11 & Data & The type of quantity for the condition: & & & \\
2-12 & Data & Quantity type for the Right Hand Side (``RHS'') quantity:
& & & \\
2-13 & Data & Index for relational operator (\textless, =,
\textgreater{} ): & & & \\
2-14 & Data & Relates to operator: & & & \\
2-17 & Data & Defines the operator used in ACTIONX for defined
quantities: & & & \\
2-18 & Data & The value relates to the operator for the triggering
condtion and if the right hand quantity is a constant or not. There is
one value for each condition and the value is determined as follows: If
the previous condition has a constant rhs =\textgreater{} iacn{[}17{]} =
0. If previous rhs is of type \{W,G, F\} and If previous boolean
operator is AND =\textgreater{} iacn{[}17{]} = 1 If previous boolean
operator is OR =\textgreater{} iacn{[}17{]} = 0 Note as mentioned
previously, the No.~column in this table and all tables in APPENDIX A:,
the count is based on one (offset one); however, in C++ the base is zero
(offset zero) and therefore C++ programmers must subtract one from this
column to obtain the correct reference. Hence, the iacn{[}17{]}
reference is the correct reference for C++. & & & \\
& SACT & Keyword defines various parameters for an ACTIONX keyword &
Required & & \\
1-1 & Format & SACT & No.~Actions x 5 & REAL & \\
2-1 to 2-5 & Data & & & & \\
& ZACT & Keyword defines the name of an action & Required & & \\
1-1 & Format & ZACT & No.~Actions x 4 & CHAR & \\
2-1 & Data & The name of an ACTION (maximum eight characters). & & & \\
& ZLACT & Keyword containing the actions defined for an ACTIONX keyword.
& Required & & \\
1-1 & Format & ZLACT & No\_Actions x Max No.~of Input Lines x Max
No.~Characters per Input Line. & CHAR & \\
2-1 to 2-? & Data & Contains the actions defined for an ACTIONX keyword.
& & & \\
& ZACN & Keyword containing the EXPRESSION parameter (item 3) on the UDQ
keyword, the ``Data for operation''. & Required & & \\
1-1 & Format & ZACN & No\_Actions x Max No.~of Conditions x 13 & CHAR
& \\
2-1 & Data & Left hand side quantity if (DAY, MONTH or YEAR). & & & \\
2-2 & Data & Right hand side quantity if (Well or Group). & & & \\
2-3 & Data & Operator. & & & \\
2-4 & Data & Well name -- left hand quantity, if it is a well quantity.
& & & \\
2-5 & Data & Well name -- right hand quantity, if it is a well quantity.
& & & \\
2-6 & Data & Group name -- left hand quantity, if it is a group
quantity. & & & \\
2-7 & Data & Group name -- right hand quantity, if it is a group
quantity. & & & \\
& SACN & & Required & & \\
1-1 & Format & SACN & No.~of Actions x Max No.~Conditions x 16 & DOUBLE
& \\
2-1 & Data & & & & \\
2-2 & Data & & & & \\
2-3 & Data & Value of condition: & & & \\
2-4 & Data & & & & \\
2-5 & Data & & & & \\
2-6 & Data & Value of right hand side. & & & \\
2-7 & Data & Value of left hand side. & & & \\
2-8 & Data & Value of right hand side. & & & \\
2-9 & Data & Value of left hand slide. & & & \\
2-10 & Data & Value of right hand side. & & & \\
\end{longtable}

\emph{Table A.29: RESTART Data -- ACTIONX Keywords}

\subsubsection{Example: RESTART Data -- Group, Well, and Connection
Keywords}\label{example-restart-data-group-well-and-connection-keywords-1}

The following example shows a typical formatted RESTART Data -- UDQ --
data.

\begin{verbatim}
'ZUDN    '          22 'CHAR'
'WUWI3   ' '        ' 'WUPR1   ' '        ' 'WUPR3   ' '        ' 'WUPR4   '
'        ' 'WUPR5   ' '        ' 'GUOPR1  ' '        ' 'GULPR1  ' '        '
'GUOPR2  ' '        ' 'GULPR2  ' '        ' 'GUOPR3  ' '        ' 'GULPR3  '
'        '
'ZUDL    '         176 'CHAR'
'(GOPR 'L' 'OWER' *0' '.931 *0.' '50 *1)  ' '        ' '        ' '        '
'        ' '        ' '        ' '        ' '        ' '        ' '        '
'        ' '        ' '1/(WWCT ' ''OP*' ) ' '        ' '        ' '        '
'        ' '        ' '        ' '        ' '        ' '        ' '        '
'        ' '        ' '        ' '        ' 'SORTA(WU' 'PR1)    ' '        '
'        ' '        ' '        ' '        ' '        ' '        ' '        '
'        ' '        ' '        ' '        ' '        ' '        ' '1/(WWCT '
''OPU*' )' '        ' '        ' '        ' '        ' '        ' '        '
'        ' '        ' '        ' '        ' '        ' '        ' '        '
'        ' 'SORTA(WU' 'PR4)    ' '        ' '        ' '        ' '        '
'        ' '        ' '        ' '        ' '        ' '        ' '        '
'        ' '        ' '        ' '(2500 - ' 'GOPR LOW' 'ER *0.03' ' - GOPR '
'UPPER *0' '.15) * 0' '.899    ' '        ' '        ' '        ' '        '
'        ' '        ' '        ' '        ' '        ' '(4500 - ' 'GOPR LOW'
'ER *0.03' ' - GOPR ' 'UPPER *0' '.15) * 0' '.899    ' '        ' '        '
'        ' '        ' '        ' '        ' '        ' '        ' '        '
'(2500 - ' 'GOPR LOW' 'ER *0.03' ' - GOPR ' 'UPPER *0' '.15) * 0' '.920    '
'        ' '        ' '        ' '        ' '        ' '        ' '        '
'        ' '        ' '(4500 - ' 'GOPR LOW' 'ER *0.03' ' - GOPR ' 'UPPER *0'
'.15) * 0' '.920    ' '        ' '        ' '        ' '        ' '        '
'        ' '        ' '        ' '        ' '(2500 - ' 'GOPR LOW' 'ER *0.03'
' - GOPR ' 'UPPER *0' '.15) * 0' '.880    ' '        ' '        ' '        '
'        ' '        ' '        ' '        ' '        ' '        ' '(4500 - '
'GOPR LOW' 'ER *0.03' ' - GOPR ' 'UPPER *0' '.15) * 0' '.880    ' '        '
'        ' '        ' '        ' '        ' '        ' '        ' '        '
'        '
'IUDQ    '          33 'INTE'
          2          -4           1           2          -1           2
          2          -1           3           2          -1           4
          2          -1           5           2          -4           1
          2          -4           2           2          -4           3
          2          -4           4           2          -4           5
          2          -4           6
\end{verbatim}

End of Example

This set of keywords are used to define the various aquifer arrays
written out by OPM Flow. If there are no aquifers in the model then this
set of keywords are skipped. If an analytical or numerical model aquifer
is present in the model, a complete set of keywords, for a given aquifer
type, should be written out. For example, if a numerical model is
defined in the model then the IAQN and RAQN keywords should be written
to the file.

The structure for this set of keywords is dependent on the aquifer type,
for analytical aquifers the format is outlined in Table A.30 and
numerical aquifers are characterized in Table A.31.

\begin{longtable}[]{@{}
  >{\raggedright\arraybackslash}p{(\linewidth - 10\tabcolsep) * \real{0.1667}}
  >{\raggedright\arraybackslash}p{(\linewidth - 10\tabcolsep) * \real{0.1667}}
  >{\raggedright\arraybackslash}p{(\linewidth - 10\tabcolsep) * \real{0.1667}}
  >{\raggedright\arraybackslash}p{(\linewidth - 10\tabcolsep) * \real{0.1667}}
  >{\raggedright\arraybackslash}p{(\linewidth - 10\tabcolsep) * \real{0.1667}}
  >{\raggedright\arraybackslash}p{(\linewidth - 10\tabcolsep) * \real{0.1667}}@{}}
\toprule\noalign{}
\begin{minipage}[b]{\linewidth}\raggedright
No.
\end{minipage} & \begin{minipage}[b]{\linewidth}\raggedright
keyword Name
\end{minipage} & \begin{minipage}[b]{\linewidth}\raggedright
RESTART Data -- Aquifer Data keywords Table A.30 - Analytical Aquifer
Keywords
\end{minipage} & \begin{minipage}[b]{\linewidth}\raggedright
Status or Value
\end{minipage} & \begin{minipage}[b]{\linewidth}\raggedright
\end{minipage} & \begin{minipage}[b]{\linewidth}\raggedright
\end{minipage} \\
\midrule\noalign{}
\endhead
\bottomrule\noalign{}
\endlastfoot
Keyword & No.~of Entries & Data Type & & & \\
& IAAQ & Analytical aquifers integer data array IAAQ(NIAAQZ, MAAQID)
where NIAAQZ and MAAQIDR are defined in the INTEHEAD keyword. & Optional
& & \\
1-1 & Format & IAAQ & NIAAQZ x MAAQID & INTE & \\
2-1 & Data & Aquifer connection number. & & & \\
2-2 & Data & Water property table. & PVTNUM & & \\
2-3 to 2-9 & & Undefined. & & & \\
2-10 & Data & Analytical aquifer type, set to zero for Fetkovich
aquifers and one for CarterTracy aquifers. & & & \\
2-11 & Data & Type Related 2. & & & \\
& SAAQ & Analytical aquifers real data array SAAQ(NSAAQZ, MAAQID) where
NSAAQZ and MAAQIDR are defined in the INTEHEAD keyword. & Optional &
& \\
1-1 & Format & SAAQ & NSAAQZ x MAAQID & REAL & \\
2-1 & Data & Total Compressibility. & COMP & & \\
2-2 & Data & Fetkovich analytical aquifer pore volume. & PORV & & \\
Carter-Tracy aquifer external radius. & RE & & & & \\
2-3 & Data & Fetkovich analytical aquifer productivity index. & PI &
& \\
Carter-Tracy aquifer permeability. & PERM & & & & \\
2-4 & Data & Fetkovich analytical aquifer time constant. & & & \\
Carter-Tracy aquifer porosity. & PORO & & & & \\
2-5 & Data & Analytical aquifer pressure. & PRESS & & \\
2-6 & Data & Analytical aquifer datum depth. & DATUM & & \\
2-7 & Data & Carter-Tracy aquifer thickness. & DZ & & \\
2-8 & Data & Carter-Tracy aquifer influence angle. & ANGLE & & \\
2-9 & Data & Carter-Tracy aquifer mass density of water. & & & \\
2-10 & Data & Carter-Tracy aquifer water viscosity. & & & \\
& XAAQ & Analytical aquifers double precision data array XAAQ(NXAAQZ,
MAAQID) where NXAAQZ and MAAQIDR are defined in the INTEHEAD keyword. &
Optional & & \\
1-1 & Format & XAAQ & NXAAQZ x MAAQID & DOUB & \\
2-1 & Data & Aquifer flow rate (AAQR: N). & & & \\
2-2 & Data & Dynamic aquifer pressure (AAQP: N). & & & \\
2-3 & Data & Liquid volume produced from the aquifer into the reservoir
(AAQT: N). & & & \\
2-4 & Data & Total aquifer influx coefficient across all aquifer
connections. & & & \\
2-5 & Data & Reciprocal time constant for Carter-Tracy aquifer. & & & \\
2-6 & Data & Influx constant ``beta'' for Carter-Tracy aquifer. & & & \\
2.7 & Data & Unknown. & & & \\
2-8 & Data & Unknown. & & & \\
2-9 & Data & Dimensionless time for Carter-Tracy aquifer (AAQTD: N) & &
& \\
2-10 & Data & Dimensionless pressure for Carter-Tracy aquifer (AAQPD: N)
& & & \\
& IAQL & Analytical aquifers integer aquifer list data array
IAQL(NIAQLX, MXNALI, MXAAQL) where NIAQLX, MXNALI, and MXAAQL are
defined in the INTEHEAD keyword. & Optional & & \\
1-1 & Format & IAQL & NIAQLX x MXNALI x MXAAQL & INTE & \\
2-1 & Data & IAQL(NIAQLX, MXNALI, MXAAQL) array. & & & \\
& ZAQL & Analytical aquifers character aquifer list data array
ZAQL(NZAQLX, MXNALI) where NZAQLX, and MXNALI are defined in the
INTEHEAD keyword. & Optional & & \\
1-1 & Format & ZAQL & NZAQLX x MXNALI & CHAR & \\
2-1 & Data & ZAQL(NZAQLX, MXNALI) array. & & & \\
& ICAQ & Analytical aquifers integer aquifer connection data array
MAAQID arrays, each of size ICAQ(NICAQZ, NGCAUS) where NICAQZ, NGCAUS,
and MAAQID are defined in the INTEHEAD keyword. & Optional & & \\
1-1 & Format & ICAQ & NICAQZ x NGCAUS x MAAQID & INTE & \\
2-1 & Data & Reservoir connecting cell in the I-direction. & I & & \\
2-2 & Data & Reservoir connecting cell in the J-direction. & J & & \\
2-3 & Data & Reservoir connecting cell in the K-direction. & K & & \\
2-4 & Data & Connecting cell status, with active equal to one otherwise
zero. & & & \\
2-5 & Data & Connecting ``face'' of the reservoir cell, where I- = 1, I+
= 2, J- = 3, J+ = 4, K- = 5, and K+ = 6. & AQUFACE & & \\
& SCAQ & Analytical aquifers real aquifer connection data array MAAQID
arrays, each of size SCAQ(NSCAQZ, NGCAUS) where NSCAQZ, NGCAUS, and
MAAQID are defined in the INTEHEAD keyword. & Optional & & \\
1-1 & Format & SCAQ & NICAQZ x NGCAUS x MAAQID & REAL & \\
2-1 & Data & Connection's fraction of total aquifer influx coefficient.
& & & \\
2-2 & Data & Connection's effective face area divided by aquifer's total
influx coefficient. & & & \\
& ACAQ & Analytical aquifers double precision aquifer connection data
array MAAQID arrays, each of size ACAQ(NACAQZ, NGCAUS) where NACAQZ,
NGCAUS, and MAAQID are defined in the INTEHEAD keyword. & Optional &
& \\
1-1 & Format & ACAQ & NACAQZ x NGCAUS x MAAQID & DOUB & \\
2-1 & Data & ACAQ( NSCAQZ, NGCAUS, MAAQID) array. & & & \\
Notes: & & & & & \\
\end{longtable}

\emph{Table A.30: RESTART Data -- Analytical Aquifer Data Keywords}

And for numerical aquifers the format is outlined in Table A.31.

\begin{longtable}[]{@{}
  >{\raggedright\arraybackslash}p{(\linewidth - 10\tabcolsep) * \real{0.1667}}
  >{\raggedright\arraybackslash}p{(\linewidth - 10\tabcolsep) * \real{0.1667}}
  >{\raggedright\arraybackslash}p{(\linewidth - 10\tabcolsep) * \real{0.1667}}
  >{\raggedright\arraybackslash}p{(\linewidth - 10\tabcolsep) * \real{0.1667}}
  >{\raggedright\arraybackslash}p{(\linewidth - 10\tabcolsep) * \real{0.1667}}
  >{\raggedright\arraybackslash}p{(\linewidth - 10\tabcolsep) * \real{0.1667}}@{}}
\toprule\noalign{}
\begin{minipage}[b]{\linewidth}\raggedright
No.
\end{minipage} & \begin{minipage}[b]{\linewidth}\raggedright
keyword Name
\end{minipage} & \begin{minipage}[b]{\linewidth}\raggedright
RESTART Data -- Aquifer Data keywords Table A.31 - Numerical Aquifer
keywords
\end{minipage} & \begin{minipage}[b]{\linewidth}\raggedright
Status or Value
\end{minipage} & \begin{minipage}[b]{\linewidth}\raggedright
\end{minipage} & \begin{minipage}[b]{\linewidth}\raggedright
\end{minipage} \\
\midrule\noalign{}
\endhead
\bottomrule\noalign{}
\endlastfoot
Keyword & No.~of Entries & Data Type & & & \\
& IAQN & Numerical aquifers integer data array IAQN(NIIAQN, NUMAQN)
where NIIAQN and NUMAQN are defined in the INTEHEAD keyword. & Optional
& & \\
1-1 & Format & IAAQ & NIIAQN x NUMAQN & INTE & \\
2-1 & Data & Aquifer identification number. & AQUID & & \\
2-2 & Data & Aquifer coordinate in the I-direction & I & & \\
2-3 & Data & Aquifer coordinate in the J-direction & J & & \\
2-4 & Data & Aquifer coordinate in the K-direction & K & & \\
2-5 & Data & Aquifer Water Property Table. & PVTNUM & & \\
2-6 & Data & Aquifer SATNUM table. & SATNUM & & \\
& RAQN & Numerical aquifers double precision data array RAQN(NIRAQN,
NUMAQN) where NIIAQN and NUMAQN are defined in the INTEHEAD keyword. &
Optional & & \\
1-1 & Format & RAAQ & NIRAQN x NUMAQN & REAL & \\
2-1 & Data & Aquifer inflow area. & AREA & & \\
2-2 & Data & Aquifer length. & LENGTH & & \\
2-3 & Data & Aquifer porosity & PORO & & \\
2-4 & Data & Aquifer permeability. & PERM & & \\
2-5 & Data & Aquifer datum depth. & DATUM & & \\
2-6 & Data & Aquifer pressure. & PRESS & & \\
2-7 & Data & Unknown & & & \\
2-8 & Data & Unknown & & & \\
2-9 & Data & Unknown & & & \\
2-10 & Data & Total aquifer pore-volume. & & & \\
2-11 & Data & Aquifer inflow rate (ANQR: N). & & & \\
2-12 & Data & Total liquid volume produced from aquifer (AQNT: N). & &
& \\
2-13 & Data & Dynamic aquifer pressure (ANQP: N). & & & \\
Notes: & & & & & \\
\end{longtable}

\emph{Table A.31: RESTART Data -- Numerical Aquifer Data Keywords}

\subsubsection{Example: RESTART Data -- Aquifer Data
Keywords}\label{example-restart-data-aquifer-data-keywords}

The following example shows a typical formatted RESTART Data -- Aquifer
Data keyword data set for an analytical aquifer.

\begin{verbatim}
 'IAAQ    '          15 'INTE'
          3           1           0           0           0           0
          0           0           0           1           1           1
          0           0           0
'SAAQ    '          24 'REAL'
  0.29999999E-04   0.24000000E+04   0.20000000E+02   0.10000000E+00
  0.49619409E+04   0.54000000E+04   0.14000000E+03   0.13333334E-01
  0.63151001E+02   0.31000000E+00   0.00000000E+00   0.43846342E+00
  0.00000000E+00   0.00000000E+00   0.00000000E+00   0.00000000E+00
  0.00000000E+00   0.00000000E+00   0.00000000E+00   0.00000000E+00
  0.00000000E+00   0.00000000E+00   0.00000000E+00   0.00000000E+00
'XAAQ    '           8 'DOUB'
  0.63857101611063D+03   0.49619408179004D+04   0.26831697328212D+05
  0.15000000000000D+06   0.23624684055200D-01   0.36097145428688D+02
  0.00000000000000D+00   0.00000000000000D+00
'ICAQNUM '           1 'INTE'
          1
'ICAQ    '          15 'INTE'
          5           1           1          13           2           5
          1           2          14           2           5           1
          3          15           2
'SCAQNUM '           1 'INTE'
          1
'SCAQ    '           6 'REAL'
  0.33333334E+00   0.33333334E+00   0.33333334E+00   0.33333334E+00
  0.33333334E+00   0.33333334E+00
\end{verbatim}

End of Example

This keyword {[}The RESTART Data -- Hidden keyword format is used both
for global and LGR grids.{]} defines solution arrays that have been
exported by OPM Flow that should not be processed by post-processing
software.

The structure for this type of keyword is defined in Table A.32.

\begin{longtable}[]{@{}
  >{\raggedright\arraybackslash}p{(\linewidth - 10\tabcolsep) * \real{0.1667}}
  >{\raggedright\arraybackslash}p{(\linewidth - 10\tabcolsep) * \real{0.1667}}
  >{\raggedright\arraybackslash}p{(\linewidth - 10\tabcolsep) * \real{0.1667}}
  >{\raggedright\arraybackslash}p{(\linewidth - 10\tabcolsep) * \real{0.1667}}
  >{\raggedright\arraybackslash}p{(\linewidth - 10\tabcolsep) * \real{0.1667}}
  >{\raggedright\arraybackslash}p{(\linewidth - 10\tabcolsep) * \real{0.1667}}@{}}
\toprule\noalign{}
\begin{minipage}[b]{\linewidth}\raggedright
No.
\end{minipage} & \begin{minipage}[b]{\linewidth}\raggedright
keyword Name
\end{minipage} & \begin{minipage}[b]{\linewidth}\raggedright
RESTART Data -- Hidden keyword Table A.32
\end{minipage} & \begin{minipage}[b]{\linewidth}\raggedright
Status or Value
\end{minipage} & \begin{minipage}[b]{\linewidth}\raggedright
\end{minipage} & \begin{minipage}[b]{\linewidth}\raggedright
\end{minipage} \\
\midrule\noalign{}
\endhead
\bottomrule\noalign{}
\endlastfoot
Keyword & No.~of Entries & Data Type & & & \\
& HIDDEN & The keyword defines the solution arrays that have not been
exported by the simulator, where N is the number of solution variable
names that have been ``hidden'', each enclosed in single quotes and of
eight characters in length. & Required & & \\
1-1 & Format & HIDDEN & N & CHAR & \\
2-1 & Data & Character list of solution arrays to be hidden from post
processing software. & & & \\
Notes: & & & & & \\
\end{longtable}

\emph{Table A.32: RESTART Data -- Hidden Keyword}

\subsubsection{Example: RESTART Data -- Hidden
Keyword}\label{example-restart-data-hidden-keyword}

The following example shows a typical formatted hidden data set, showing
51 Solution arrays that should be hidden.

\begin{verbatim}
'HIDDEN  '          51 'CHAR'
'ISTHGGI ' 'SWMAX   ' 'SWHY1   ' 'SWHY2   ' 'SWHY3   ' 'ISTHW   ' 'SOMAX   '
'SGHY1   ' 'SGHY2   ' 'SGHY3   ' 'ISTHG   ' 'SGMAX   ' 'SHMAX   ' 'WATQUIES'
'GASQUIES' 'PRESROCC' 'SPDOIL  ' 'SPDWATER' 'SPDGAS  ' 'SWINIT  ' 'SGINIT  '
'POLYMAX ' 'EIP0    ' 'EIP     ' 'TINIT   ' 'IOCN    ' 'SGMIN   ' 'SGSTRD  '
'SWSTRD  ' 'RISTRD  ' 'SGSTRI  ' 'SWSTRI  ' 'RISTRI  ' 'SGTRPD  ' 'SGSTRWD '
'SGSTRWI ' 'SWMIN   ' 'SGSTRWS ' 'RPTHMW  ' 'PPCW    ' 'GRADWELL' 'GRADCONN'
'SWHY4   ' 'SGHY4   ' 'GRADGRUP' 'GRADRESV' 'GRADFIP ' 'GRADSOL ' 'GRADEXP '
'IAQN    ' 'RAQN    '
\end{verbatim}

End of Example

The solution data keywords {[}The RESTART Data - Solution Data keywords
are used both for global and LGR grids SOLUTION arrays. The LGR data is
preceded by a series of LGR head keywords and terminated by an LGR
termination keyword.{]} define the solution for each active cell for
both global and LGR grids. This is then followed by a series of global
data for a global entry or LGR data for a LGR entry.

The structure for this type of keyword is defined in Table A.33 for
standard solution data and Table A.34 for tracers. A list of all
solution names is tabulated in Table A.35.

\begin{longtable}[]{@{}
  >{\raggedright\arraybackslash}p{(\linewidth - 10\tabcolsep) * \real{0.1667}}
  >{\raggedright\arraybackslash}p{(\linewidth - 10\tabcolsep) * \real{0.1667}}
  >{\raggedright\arraybackslash}p{(\linewidth - 10\tabcolsep) * \real{0.1667}}
  >{\raggedright\arraybackslash}p{(\linewidth - 10\tabcolsep) * \real{0.1667}}
  >{\raggedright\arraybackslash}p{(\linewidth - 10\tabcolsep) * \real{0.1667}}
  >{\raggedright\arraybackslash}p{(\linewidth - 10\tabcolsep) * \real{0.1667}}@{}}
\toprule\noalign{}
\begin{minipage}[b]{\linewidth}\raggedright
No.
\end{minipage} & \begin{minipage}[b]{\linewidth}\raggedright
keyword Name
\end{minipage} & \begin{minipage}[b]{\linewidth}\raggedright
RESTART Data - Solution Data keyword Table A.33
\end{minipage} & \begin{minipage}[b]{\linewidth}\raggedright
Status or Value
\end{minipage} & \begin{minipage}[b]{\linewidth}\raggedright
\end{minipage} & \begin{minipage}[b]{\linewidth}\raggedright
\end{minipage} \\
\midrule\noalign{}
\endhead
\bottomrule\noalign{}
\endlastfoot
Keyword & No.~of Entries & Data Type & & & \\
& STARTSOL & This keyword marks the start of the solution variable
section for both global and local grids. Note that there is no data
associated with this keyword. & Required & & \\
1-1 & Format & STARTSOL & 0 & MESG & \\
& LGRNAMES & This keyword defines a list of LGR names for the reporting
time step, where N is the number of LGR names. & & & \\
1-1 & Format & LGRNAMES & N & CHAR & \\
2-1 & Data & A character list of LGR names with each name enclosed in
single quotes and each having a length of eight characters. & & & \\
& SOLUTION & The SOLUTION keyword defines the start of a SOLUTION array
data set, where SOLUTION is the name of the array being written out, as
outlined in Table A.35. For example, PRESSURE, SWAT, SGAS, RS, RV etc.
The SOLUTION keyword (items 1-1 and 2-1) is repeated for each SOLUTION
array being written out. SOLUTION arrays can be in any order and only
the active cells are written out (NACTIV). & Required & & \\
1-1 & Format & SOLUTION & NACTIV & REAL & \\
2-1 & Data & SOLUTION data set. & & & \\
& ENDSOL & This keyword marks the end of the solution variable section
for both global and local grids. Note that there is no data associated
with this keyword. & Required & & \\
1-1 & Format & ENDSOL & 0 & MESG & \\
Notes: & & & & & \\
\end{longtable}

\emph{Table A.33: RESTART Data - Solution Data keyword}

For SOLUTION tracer concentration arrays the ZTRACER keyword must
precede the SOLUTION tracer data array -- see Table A.34 for a
description of ZTRACER keyword.

\begin{longtable}[]{@{}
  >{\raggedright\arraybackslash}p{(\linewidth - 10\tabcolsep) * \real{0.1667}}
  >{\raggedright\arraybackslash}p{(\linewidth - 10\tabcolsep) * \real{0.1667}}
  >{\raggedright\arraybackslash}p{(\linewidth - 10\tabcolsep) * \real{0.1667}}
  >{\raggedright\arraybackslash}p{(\linewidth - 10\tabcolsep) * \real{0.1667}}
  >{\raggedright\arraybackslash}p{(\linewidth - 10\tabcolsep) * \real{0.1667}}
  >{\raggedright\arraybackslash}p{(\linewidth - 10\tabcolsep) * \real{0.1667}}@{}}
\toprule\noalign{}
\begin{minipage}[b]{\linewidth}\raggedright
No.
\end{minipage} & \begin{minipage}[b]{\linewidth}\raggedright
keyword Name
\end{minipage} & \begin{minipage}[b]{\linewidth}\raggedright
RESTART Data -- Solution Data keyword for Tracer Concentration Name
Table A.34
\end{minipage} & \begin{minipage}[b]{\linewidth}\raggedright
Status or Value
\end{minipage} & \begin{minipage}[b]{\linewidth}\raggedright
\end{minipage} & \begin{minipage}[b]{\linewidth}\raggedright
\end{minipage} \\
\midrule\noalign{}
\endhead
\bottomrule\noalign{}
\endlastfoot
Keyword & No.~of Entries & Data Type & & & \\
& ZTRACER & This keyword marks the start of the solution variable
section for both global and local grids. The TRACER in the Data field is
the name of the tracer concentration being processed. & Optional & & \\
1-1 & Format & ZTRACER & 2 & CHAR & \\
2-1 & Data & TRACER & ` ` & & \\
& TRACER & TRACER is the name of the tracer concentration that is
contained on the next keyword. Tracer arrays can be in any order and
only the active cells are written out (NACTIV). & & & \\
1-1 & Format & TRACER & NACTIV & REAL & \\
2-1 & Data & Tracer solution data array. & & & \\
Notes: & & & & & \\
\end{longtable}

\emph{Table A.34: RESTART Data -- Solution Data Keyword for Tracer
Concentration Name}

The SOLUTION keywords are listed in Table A.35.

\begin{longtable}[]{@{}
  >{\raggedright\arraybackslash}p{(\linewidth - 10\tabcolsep) * \real{0.1667}}
  >{\raggedright\arraybackslash}p{(\linewidth - 10\tabcolsep) * \real{0.1667}}
  >{\raggedright\arraybackslash}p{(\linewidth - 10\tabcolsep) * \real{0.1667}}
  >{\raggedright\arraybackslash}p{(\linewidth - 10\tabcolsep) * \real{0.1667}}
  >{\raggedright\arraybackslash}p{(\linewidth - 10\tabcolsep) * \real{0.1667}}
  >{\raggedright\arraybackslash}p{(\linewidth - 10\tabcolsep) * \real{0.1667}}@{}}
\toprule\noalign{}
\begin{minipage}[b]{\linewidth}\raggedright
Model Type
\end{minipage} & \begin{minipage}[b]{\linewidth}\raggedright
Model Option
\end{minipage} & \begin{minipage}[b]{\linewidth}\raggedright
Solution Data Keywords
\end{minipage} & \begin{minipage}[b]{\linewidth}\raggedright
Solution Variable Description
\end{minipage} & \begin{minipage}[b]{\linewidth}\raggedright
OPM Output Status
\end{minipage} & \begin{minipage}[b]{\linewidth}\raggedright
\end{minipage} \\
\midrule\noalign{}
\endhead
\bottomrule\noalign{}
\endlastfoot
Black-Oil & Compositional & & & & \\
& & STANDARD & 1OVERBO & Reciprocal of oil formation volume factor. & \\
& & STANDARD & 1OVERBW & Reciprocal of water formation volume factor.
& \\
& & STANDARD & CNV\_DGAS & Worst cells depending on the gas saturation
increment / Rv increment / Rs increment. & \\
& & STANDARD & CNV\_DPRE & Worst cells depending on the pressure
increment. & \\
& & STANDARD & CNV\_DWAT & Worst cells depending on the water saturation
increment. & \\
& & STANDARD & CNV\_GAS & Worst cells depending on the residual of gas
equation. & \\
& & STANDARD & CNV\_OIL & Worst cells depending on the residual of the
oil equation. & \\
& & STANDARD & CNV\_WAT & Worst cells depending on the residual of the
water equation. & \\
& & STANDARD & CONV\_NEW & Number of Newtons required by each cell in
order to satisfy the solution change convergence criteria at the last
time step. & \\
& & STANDARD & CONV\_PRU & Worst cells depending on the pressure update.
& \\
& & STANDARD & CONV\_VBR & Worst cells depending on the volume balance
residual. & \\
& & STANDARD & DRAINAGE & Drainage region numbers. & \\
& & STANDARD & DRAINMIN & Drainage sink indicator. & \\
& & STANDARD & FIPGAS & Gas fluid-in-place. & \\
& & STANDARD & FIPOIL & Oil fluid-in-place. & \\
& & STANDARD & FIPWAT & Water fluid-in-place. & \\
& & STANDARD & GAS\_DEN & Gas density. & \\
& & STANDARD & GAS\_POTN & Gas potential. & \\
& & STANDARD & GAS\_PRES & Gas phase pressure. & \\
& & STANDARD & GAS\_VISC & Gas viscosity. & \\
& & STANDARD & IOVERBG & Reciprocal of gas formation volume factor. & \\
& & STANDARD & ISTHG & Gas capillary pressure state. & \\
& & STANDARD & ISTHW & Water capillary pressure state. & \\
& & STANDARD & OIL\_DEN & Oil density. & \\
& & STANDARD & OIL\_POTN & Oil potential. & \\
& & STANDARD & OIL\_VISC & Oil viscosity. & \\
& & STANDARD & PBUB & Bubble point pressure. & \\
& & STANDARD & PCOG & Oil-gas capillary pressure. & \\
& & STANDARD & PCOW & Oil-water capillary pressure. & \\
& & STANDARD & PDEW & Dew point pressure. & \\
& & STANDARD & PORV & Pore volume at surface conditions. & \\
& & STANDARD & PRESSURE & Pressure. & \\
& & STANDARD & RFIPGAS & Gas fluid-in-place at reservoir conditions.
& \\
& & STANDARD & RFIPOIL & Oil fluid-in-place at reservoir conditions.
& \\
& & STANDARD & RFIPWAT & Water fluid-in-place at reservoir conditions.
& \\
& & STANDARD & RPORV & Pore volume at reservoir conditions. & \\
& & STANDARD & RS & Gas-oil ratio. & \\
& & STANDARD & RSSAT & Saturated gas-oil ratio. & \\
& & STANDARD & RV & Oil-gas ratio. & \\
& & STANDARD & RVSAT & Saturated oil-gas ratio. & \\
& & STANDARD & SFIPGAS & Gas fluid-in-place at surface/separator
conditions. & \\
& & STANDARD & SFIPOIL & Oil fluid-in-place at surface/separator
conditions. & \\
& & STANDARD & SFIPWAT & Water fluid-in-place at surface/separator
conditions. & \\
& & STANDARD & SGAS & Gas saturation. & \\
& & STANDARD & SGMAX & Maximum gas saturation. & \\
& & STANDARD & SGMIN & Minimum gas saturation. & \\
& & STANDARD & SOIL & Oil saturation. & \\
& & STANDARD & SOMAX & Maximum oil saturation. & \\
& & STANDARD & SSOL & Solvent saturation. & \\
& & STANDARD & STATES & Gas-oil state indicator. & \\
& & STANDARD & SWAT & Water saturation. & \\
& & STANDARD & SWMAX & Maximum water saturation. & \\
& & STANDARD & SWMIN & Minimum water saturation. & \\
& & STANDARD & WAT\_DEN & Water density. & \\
& & STANDARD & WAT\_POTN & Water potential. & \\
& & STANDARD & WAT\_PRES & Water phase pressure. & \\
& & STANDARD & WAT\_VISC & Water viscosity. & \\
& & STANDARD & XMF & Liquid mole fractions. & \\
& & STANDARD & YMF & Vapor mole fractions. & \\
& & STANDARD & ZMF & Total mole fractions. & \\
& & ALKALINE & ALKADS & Alkaline adsorption. & \\
& & ALKALINE & ALKALINE & Alkaline concentration. & \\
& & ALKALINE & ALKMAX & Alkaline maximum historic concentration. & \\
& & ALKALINE & PADMAX & Alkaline polymer maximum historic adsorption.
& \\
& & ALKALINE & PLADALK & Alkaline polymer adsorption multipliers. & \\
& & ALKALINE & SFADALK & Alkaline surfactant adsorption multipliers.
& \\
& & ALKALINE & STMALK & Alkaline water/oil surface tension multipliers.
& \\
& & API & SDENO & Oil surface densities. & \\
& & API & OILAPI & Oil API values. & \\
& & AQUIFER & HYDH & Hydraulic head aquifer. & \\
& & AQUIFER & HYDHFW & Fresh water hydraulic head aquifer. & \\
& & BRINE & CNV\_DSAL & Worst cells depending on the brine concentration
increment. & \\
& & BRINE & CNV\_SAL & Worst cells depending on the residual of the
brine equation (BRINEoption when used with polymer). & \\
& & BRINE & ESALSUR & Effective salinity for surfactant. & \\
& & BRINE & SALT & Brine concentration. & \\
& & COAL & COALGAS & Coal gas concentration for coal bed methane option.
& \\
& & COAL & COALSOLV & Solvent concentration for coal bed methane option.
& \\
& & COAL & GASSATC & Initial coal gas saturated content for coal bed
methane option. & \\
& & COAL & MLANG & Langmuir scaling factors for coal bed methane option.
& \\
& & COAL & MLANGSLV & Langmuir scaling factors for solvent in coal bed
methane option. & \\
& & ENDSCALE & GASKR & Gas relative permeability. & \\
& & ENDSCALE & OILKR & Oil relative permeability. & \\
& & ENDSCALE & WATKR & Water relative permeability. & \\
& & EXCAVATE & EXCAVNUM & Excavation status identifier. & \\
& & EXCAVATE & TRANEXX/Y/Z & Transmissibilities in excavation runs. & \\
& & FOAM & FOAM & Foam concentration. & \\
& & FOAM & FOAM\_HL & Foam half-life. & \\
& & FOAM & FOAMADS & Foam adsorption. & \\
& & FOAM & FOAMCNM & Foam capillary numbers. & \\
& & FOAM & FOAMDCY & Foam decay. & \\
& & FOAM & FOAMMAX & Foam maximum historic concentration. & \\
& & FOAM & FOAMMOB & Foam mobility multiplier. & \\
& & GIMODEL & CNV\_DGGI & Worst cells depending on the dry injection gas
parameter (Gi) increment. & \\
& & GIMODEL & CNV\_GGI & Worst cells depending on the residual of the Gi
Pseudo-Compositional modeL. & \\
& & GIMODEL & GGI & GI injected gas ratio. & \\
& & GIMODEL & RFIPGGI & Dry injection gas in place at reservoir
conditions. & \\
& & GIMODEL & SFIPGGI & Dry injection gas in place at surface
conditions. & \\
& & POLYMER & CABINnnn & This is for use with the PLYTRRFA keyword only.
& \\
& & POLYMER & CNV\_DPLY & Worst cells depending on the polymer
concentration increment. & \\
& & POLYMER & CNV\_PLY & Worst cells depending on the residual of the
polymer equation. & \\
& & POLYMER & PADS & Adsorbed polymer concentrations. & \\
& & POLYMER & POLYMAX & Maximum historic polymer concentration. & \\
& & POLYMER & POLYMER & Polymer concentrations. & \\
& & POLYMER & SFIPPLY,RFIPPLY & Polymer-in-place. & \\
& & POLYMER & SFIPSAL,RFIPSAL & Salt-in-place (BRINE option when used
with polymer). & \\
& & ROCKCOMP & PRESROCC & Rock pressure values used for rock compaction
model. & \\
& & SOLVENT & CNV\_DSOL & Worst cells depending on the solvent
concentration increment. & \\
& & SOLVENT & CNV\_SOL & Worst cells depending on the residual of the
solvent equation. & \\
& & SOLVENT & RFIPSOL & Solvent-in-place at reservoir conditions. & \\
& & SOLVENT & SFIPSOL & Solvent-in-place at surface conditions. & \\
& & SURFACT & CATROCK & Divalent cation concentration associated with
rock. & \\
& & SURFACT & CATSURF & Divalent cation concentration associated with
surfactant. & \\
& & SURFACT & SURFACT & Surface interactions. & \\
& & SURFACT & SURFADS & Adsorbed surfactant concentrations. & \\
& & SURFACT & SURFCNM & Surfactant capillary numbers. & \\
& & SURFACT & SURFMAX & Maximum surfactant concentrations. & \\
& & SURFACT & SURFST & Surface tension in surfactant runs. & \\
& & TEMP & TEMP & Temperature. & \\
OPM & & THERMAL & TEMP & Temperature. & \\
& & VE & EOGC & Maximum oil-gas contact. & \\
& & VE & EOWC & Minimum oil-water contact. & \\
& & VE & GWC & Gas-water contact. & \\
& & VE & OGC & Oil-gas contact. & \\
& & VE & OWC & Oil-water contact. & \\
& & VE & POT\_CORR & Initial contact corrected potential. & \\
Notes: & & & & & \\
\end{longtable}

\emph{Table A.35: RESTART Data - Solution Data Keywords}

\subsubsection{Example: RESTART Data -- Solution Data
Keyword}\label{example-restart-data-solution-data-keyword}

The following example shows a typical formatted global grid data set for
this keyword set.

\begin{verbatim}
 'STARTSOL'           0 'MESS'
 'PRESSURE'       44431 'REAL'
   0.26889725E+03   0.26838983E+03   0.26826810E+03   0.26820352E+03
   0.26814493E+03   0.26809378E+03   0.26807767E+03   0.26806363E+03
   0.26802625E+03   0.26798474E+03   0.26795001E+03   0.26791434E+03
   0.26787915E+03   0.26783920E+03   0.26777118E+03   0.26769208E+03         ………………………….……………..
 'SWAT    '       44431 'REAL'
   0.10500000E+00   0.10500000E+00   0.10500000E+00   0.10500000E+00
   0.10500000E+00   0.10500000E+00   0.10500000E+00   0.14000000E+00
   0.14000000E+00   0.14500000E+00   0.14500000E+00   0.14500000E+00
   0.14500000E+00   0.14500000E+00   0.16000000E+00   0.16000000E+00
………………………….……………..
 'SGAS    '       44431 'REAL'
   0.00000000E+00   0.89499998E+00   0.89499998E+00   0.89499998E+00
   0.89499998E+00   0.89499998E+00   0.89499998E+00   0.86000001E+00
   0.86000001E+00   0.85500002E+00   0.85500002E+00   0.85500002E+00
………………………….……………..
 'RS      '       44431 'REAL'
   0.11697504E+03   0.12030303E+03   0.12023456E+03   0.12019825E+03
   0.12016528E+03   0.12013654E+03   0.12012746E+03   0.12011956E+03
   0.12009855E+03   0.12007520E+03   0.12005566E+03   0.12003561E+03
   0.12001582E+03   0.11999335E+03   0.11995509E+03   0.11991060E+03
   0.11987761E+03   0.11985194E+03   0.11983020E+03   0.11981972E+03
   0.11980798E+03   0.11977385E+03   0.11970672E+03   0.11963727E+03
  ………………………….……………..
 'RV      '       44431 'REAL'
   0.57879315E-04   0.57552861E-04   0.57483696E-04   0.57447018E-04
   0.57413723E-04   0.57384681E-04   0.57375513E-04   0.57367535E-04
   0.57346311E-04   0.57322733E-04   0.57302990E-04   0.57282738E-04
   0.57262747E-04   0.57240050E-04   0.57201407E-04   0.57156467E-04
   0.57123143E-04   0.57097215E-04   0.57075260E-04   0.57064670E-04
   0.57052810E-04   0.57018340E-04   0.56950528E-04   0.56880377E-04
………………………….……………..
 'TRFIELD '          28 'DOUB'
   0.00000000000000D+00   0.00000000000000D+00   0.00000000000000D+00
   0.00000000000000D+00   0.00000000000000D+00   0.00000000000000D+00
   0.00000000000000D+00   0.00000000000000D+00   0.00000000000000D+00
   0.00000000000000D+00   0.00000000000000D+00   0.00000000000000D+00
   0.00000000000000D+00   0.00000000000000D+00   0.00000000000000D+00
   0.00000000000000D+00   0.00000000000000D+00   0.00000000000000D+00
   0.00000000000000D+00   0.00000000000000D+00   0.00000000000000D+00
   0.00000000000000D+00   0.00000000000000D+00   0.00000000000000D+00
   0.00000000000000D+00   0.00000000000000D+00   0.00000000000000D+00
   0.00000000000000D+00
 'ZTRACER '           2 'CHAR'
 'SEAF    ' '        '
 'SEAF    '       44431 'REAL'
   0.00000000E+00   0.00000000E+00   0.00000000E+00   0.00000000E+00
   0.00000000E+00   0.00000000E+00   0.00000000E+00   0.00000000E+00
   0.00000000E+00   0.00000000E+00   0.00000000E+00   0.00000000E+00
   0.00000000E+00   0.00000000E+00   0.00000000E+00   0.00000000E+00
   0.00000000E+00   0.00000000E+00   0.00000000E+00   0.00000000E+00
   0.00000000E+00   0.00000000E+00   0.00000000E+00   0.00000000E+00
   0.00000000E+00   0.00000000E+00   0.00000000E+00   0.00000000E+00
………………………….……………..
\end{verbatim}

End of Example

The file structure for this series of keywords is similar to the global
grid RESTART Data keywords described previously. However, additional
data that the defines the LGR properties (LGR Name for example) are
included in this keyword definition and the LGR keyword is repeated for
each LGR in the model. The keyword description is outlined in Table
A.36.

Note that currently OPM Flow does not support LGR grids and therefore
this keyword format cannot be used by OPM Flow.

\begin{longtable}[]{@{}
  >{\raggedright\arraybackslash}p{(\linewidth - 6\tabcolsep) * \real{0.2500}}
  >{\raggedright\arraybackslash}p{(\linewidth - 6\tabcolsep) * \real{0.2500}}
  >{\raggedright\arraybackslash}p{(\linewidth - 6\tabcolsep) * \real{0.2500}}
  >{\raggedright\arraybackslash}p{(\linewidth - 6\tabcolsep) * \real{0.2500}}@{}}
\toprule\noalign{}
\begin{minipage}[b]{\linewidth}\raggedright
No.
\end{minipage} & \begin{minipage}[b]{\linewidth}\raggedright
keyword Name
\end{minipage} & \begin{minipage}[b]{\linewidth}\raggedright
RESTART Data - LGR Grid Property Data keyword Table A.36 (Repeated for
Each LGR in the Model)
\end{minipage} & \begin{minipage}[b]{\linewidth}\raggedright
Status or Value
\end{minipage} \\
\midrule\noalign{}
\endhead
\bottomrule\noalign{}
\endlastfoot
& LGR & LGR Grid Header keyword. See section A.5.1 RESTART Data - LGR
Grid Header Keywords. & Required \\
LGRHEADI & & & \\
LGRHEADQ & & & \\
LGRHEADD & & & \\
INTEHEAD & Header keyword. See section A.8.3 RESTART Data - Header
Keywords. & & \\
LOGIHEAD & & & \\
DOUBHEAD & & & \\
& IGRP SGRP etc. & Group, well, and connection data status keywords for
this reporting time step. See section A.8.4 RESTART Data -- Group, Well,
and Connection Data Keywords. & Required \\
& IAAQ SAAQ ect. & Aquifer definition arrays, only applicable if an
aquifer is connected to an LGR. See section A.8.6 RESTART Data --
Aquifer Data Keywords. & Optional \\
& HIDDEN & HIDDEN solution data, that is a list of solution arrays NOT
to be read by post processing software. See section A.8.7 RESTART Data
-- Hidden Keyword. & Required \\
& SOLUTION & LGR SOLUTION array section (PRESSURE, SGAS, SOILS, WAT,
etc.). See section A.8.8 RESTART Data - Solution Data keyword. &
Required \\
& ENDSOL & LGR SOLUTION section termination keyword. See section A.8.8
RESTART Data - Solution Data keyword. & Required \\
& ENDLGR & LGR grid header section termination keywords. See section
A.5.1 RESTART Data - LGR Grid Header Keywords. & Required \\
Notes: & & & \\
\end{longtable}

\emph{Table A.36: RESTART Data -- LGR Solution Data Keywords}

This set of keywords are repeated for each LGR in the model with the
last LGR data set terminated by the ENDSOL and ENDLGR keywords to mark
the end of the LGR input - see the example.

\subsubsection{Example: RESTART Data - LGR Solution Data
Keywords}\label{example-restart-data---lgr-solution-data-keywords}

The following example shows a typical formatted data set with two LGR
grids named LGR-1 and LGR-2.

\begin{verbatim}
 'LGR     '           1 'CHAR'
 'LGR-1   '
 'LGRHEADI'          45 'INTE'
           1         100       -2345       -2345       -2345       -2345
………………………….……………..
       -2345       -2345       -2345
 'LGRHEADQ'           5 'LOGI'
  F  F  F  F  F
 'LGRHEADD'           5 'DOUB'
   0.00000000000000D+00  -0.10000000200409D+21  -0.10000000200409D+21
  -0.10000000200409D+21  -0.10000000200409D+21
 'INTEHEAD'         249 'INTE'
………………………….……………..
 'LOGIHEAD'          79 'LOGI'
………………………….……………..
 'DOUBHEAD'         185 'DOUB'
………………………….……………..
 'IGRP    '         358 'INTE'
           0           0           0           0           0           0
………………………….……………..
           0           0           0           0
 'SGRP    '         186 'REAL'
   0.00000000E+00   0.00000000E+00   0.00000000E+00   0.00000000E+00
………………………….……………..
   0.00000000E+00   0.10000000E+01
 'XGRP    '         348 'DOUB'
   0.00000000000000D+00   0.00000000000000D+00   0.00000000000000D+00
………………………….……………..
   0.00000000000000D+00   0.00000000000000D+00   0.00000000000000D+00
 'ZGRP    '          10 'CHAR'
 '        ' '        ' '        ' '        ' '        ' 'FIELD   ' '        '
 '        ' '        ' '        '
 'ICOT    '           0 'INTE'
 'XCOT    '           0 'DOUB'
 'IWET    '           0 'INTE'
 'XWET    '           0 'DOUB'
 'IGRT    '          18 'INTE'
………………………….……………..
 'XGRT    '         180 'DOUB'
   0.00000000000000D+00   0.00000000000000D+00   0.00000000000000D+00
………………………….……………..
   0.10000000200409D+21   0.10000000200409D+21   0.10000000200409D+21
 'DLYTIM  '          30 'DOUB'
   0.00000000000000D+00   0.00000000000000D+00   0.00000000000000D+00
………………………….……………..
   0.00000000000000D+00   0.00000000000000D+00   0.00000000000000D+00
 'HIDDEN  '          51 'CHAR'
 'ISTHGGI ' 'SWMAX   ' 'SWHY1   ' 'SWHY2   ' 'SWHY3   ' 'ISTHW   ' 'SOMAX   '………………………….……………..
 'IAQN    ' 'RAQN    '
 'STARTSOL'           0 'MESS'
 'PRESSURE'           9 'REAL'
   0.26823886E+03   0.26824997E+03   0.26826108E+03   0.26822879E+03
………………………….……………..
 'SWAT    '           9 'REAL'
   0.10500000E+00   0.10500000E+00   0.10500000E+00   0.10500000E+00
………………………….……………..
 'SGAS    '           9 'REAL'
   0.89499998E+00   0.89499998E+00   0.89499998E+00   0.89499998E+00
………………………….……………..
 'RS      '           9 'REAL'
   0.12021812E+03   0.12022437E+03   0.12023062E+03   0.12021246E+03
………………………….……………..
 'RV      '           9 'REAL'
   0.57467099E-04   0.57473411E-04   0.57479720E-04   0.57461380E-04
………………………….……………..
 'ENDSOL  '           0 'MESS'
 'ENDLGR  '           1 'INTE'
           1
 'LGR     '           1 'CHAR'
 'LGR-2   '
 'LGRHEADI'          45 'INTE'
………………………….……………..
 'LGRHEADQ'           5 'LOGI'
………………………….……………..
 'LGRHEADD'           5 'DOUB'
………………………….……………..
 'INTEHEAD'         249 'INTE'
………………………….……………..
 'LOGIHEAD'          79 'LOGI'
………………………….……………..
 'DOUBHEAD'         185 'DOUB'
………………………….……………..
 'IGRP    '         358 'INTE'
………………………….……………..
 'SGRP    '         186 'REAL'
………………………….……………..
 'XGRP    '         348 'DOUB'
………………………….……………..
 'ZGRP    '          10 'CHAR'
………………………….……………..
 'ICOT    '           0 'INTE'
 'XCOT    '           0 'DOUB'
 'IWET    '           0 'INTE'
 'XWET    '           0 'DOUB'
 'IGRT    '          18 'INTE'
………………………….……………..
 'XGRT    '         180 'DOUB'
………………………….……………..
 'DLYTIM  '          30 'DOUB'
………………………….……………..
 'HIDDEN  '          51 'CHAR'
………………………….……………..
 'STARTSOL'           0 'MESS'
 'PRESSURE'           9 'REAL'
………………………….……………..
 'SWAT    '           9 'REAL'
………………………….……………..
 'SGAS    '           9 'REAL'
………………………….……………..
 'RS      '           9 'REAL'
………………………….……………..
 'RV      '           9 'REAL'
………………………….……………..
 'ENDSOL  '           0 'MESS'
 'ENDLGR  '           1 'INTE'
           2
 'SEQNUM  '           1 'INTE'
           4
\end{verbatim}

End of Example

The Repeat Formation Tester (``RFT'') file contains a well's depth
pressure profile at the time the data is requested using the WRFT or
WRFTPLT keywords in the SCHEDULE section. The data is meant to be used
to compare the simulator's pressure profile with the field measured data
collected from one of the open hole wire line logging tools: Formation
Interval Tester (``FIT''), Repeat Formation Tester (``RFT'') tool,
Modular Dynamics Tester (``MDT''), or similar tool. The keywords also
export the fluid saturations for the well connections at the same time.

In addition to the pressure and saturation data, the WRFTPLT keyword can
also active the writing out of each well connection's fluid rates,
connection factors and KH data, etc., collectively known as the
Production Logging Tool (``PLT'') data. The PLT data is used to compare
with measured data from wire line production logging tools. Currently,
output of the PLT data is not supported by OPM Flow.

The file structure for the RFT Data File consists of various keywords: a
Time and Date keyword, a Well and Connection keyword followed by the
requested optional keywords of RFT, PLT, Multi-Segment Well, and River
keywords, that are written out depending on the data requested. A brief
description of the keywords is given in Table A.37. A complete RFT data
set is written to file for each well and time step as requested on the
WRFT and WRFTPLT keywords.

\begin{longtable}[]{@{}
  >{\raggedright\arraybackslash}p{(\linewidth - 6\tabcolsep) * \real{0.2500}}
  >{\raggedright\arraybackslash}p{(\linewidth - 6\tabcolsep) * \real{0.2500}}
  >{\raggedright\arraybackslash}p{(\linewidth - 6\tabcolsep) * \real{0.2500}}
  >{\raggedright\arraybackslash}p{(\linewidth - 6\tabcolsep) * \real{0.2500}}@{}}
\toprule\noalign{}
\begin{minipage}[b]{\linewidth}\raggedright
Reference Section
\end{minipage} & \begin{minipage}[b]{\linewidth}\raggedright
keyword Name
\end{minipage} & \begin{minipage}[b]{\linewidth}\raggedright
keyword Type
\end{minipage} & \begin{minipage}[b]{\linewidth}\raggedright
Description
\end{minipage} \\
\midrule\noalign{}
\endhead
\bottomrule\noalign{}
\endlastfoot
A.9.2 & TIME & Time and Date Date & This data set contains the simulated
time and date information that the date was written out. \\
DATE & & & \\
A.9.3 & WELLETC & Well and Connection Data & The Well and Connection
keywords outlines the type and units being written out, the type of well
and the well connections at the simulated time the data is written out
to the RFT file. \\
CONIPOS & & & \\
CONJPOS & & & \\
CONKPOS & & & \\
HOSTGRID & & & \\
A.9.4 & DEPTH & RFT Solution Data & RFT keyword data consists of the
pressure and saturations for each declared connection in the well at the
simulated time the data is written out to file. \\
PRESSURE & & & \\
SWAT & & & \\
SGAS & & & \\
CPLY & & & \\
CPLAD & & & \\
CBRI & & & \\
A.9.5 & & PLT Solution Data & PLT series of keywords contain the
production logging data which consists of the mid-point depth of each
connection, the tubing length for the connection from the bottom-hole
reference depth, the pressures and fluid rates for the connections. In
addition, this data also includes the well pressures, fluid rates,
connection transmissibility and KH vales, for the given simulated time.
Currently this data set is not supported. \\
A.9.6 & & Multi-Segment Well Data & The Multi-Segment Well data set
consists of pressure and flow rate data for segments and branches,
currently this data set is not supported. \\
A.9.7 & & River Data & Data pertaining to the Rivers option, this option
is not supported by OPM Flow. \\
Notes: & & & \\
\end{longtable}

\emph{Table A.37: RFT Data File Format}

The time and date keywords define the start of an RFT data set and
declares the time and date for the subsequent property data written to
file. The structure for these two keywords is summarized in Table A.38.

\begin{longtable}[]{@{}
  >{\raggedright\arraybackslash}p{(\linewidth - 10\tabcolsep) * \real{0.1667}}
  >{\raggedright\arraybackslash}p{(\linewidth - 10\tabcolsep) * \real{0.1667}}
  >{\raggedright\arraybackslash}p{(\linewidth - 10\tabcolsep) * \real{0.1667}}
  >{\raggedright\arraybackslash}p{(\linewidth - 10\tabcolsep) * \real{0.1667}}
  >{\raggedright\arraybackslash}p{(\linewidth - 10\tabcolsep) * \real{0.1667}}
  >{\raggedright\arraybackslash}p{(\linewidth - 10\tabcolsep) * \real{0.1667}}@{}}
\toprule\noalign{}
\begin{minipage}[b]{\linewidth}\raggedright
No.
\end{minipage} & \begin{minipage}[b]{\linewidth}\raggedright
keyword Name
\end{minipage} & \begin{minipage}[b]{\linewidth}\raggedright
RFT Data File -- Time and Date keywords Table A.38
\end{minipage} & \begin{minipage}[b]{\linewidth}\raggedright
Status or Value
\end{minipage} & \begin{minipage}[b]{\linewidth}\raggedright
\end{minipage} & \begin{minipage}[b]{\linewidth}\raggedright
\end{minipage} \\
\midrule\noalign{}
\endhead
\bottomrule\noalign{}
\endlastfoot
Keyword & No.~of Entries & Data Type & & & \\
& TIME & This keyword marks the start of of an RFT file data set. &
Required & & \\
1-1 & Format & TIME & 1 & REAL & \\
2-1 & Data & Simulation time the data is being written to file. & & & \\
& DATE & Three integer values representing the date of the data being
written out. & Required & & \\
1-1 & Format & DATE & 3 & INTE & \\
2-1 & Data & Calendar day and should range between one and 31. & DAY &
& \\
2-2 & Data & Calendar month and should range between one and 12. & MONTH
& & \\
2-3 & Data & Calendar year in four digits, that is 2020. & YEAR & & \\
Notes: & & & & & \\
\end{longtable}

\emph{Table A.38: RFT Data File -- Time and Date Keywords}

The DATE keyword is immediate followed by the well and connection
keywords.

The well and connection data set of keywords specify the well and the
well connection data as well as the system of units for the pressure and
flow data being written to the RFT file. The structure for this set of
keywords is outlined in Table A.39 and consists of WELLETC, CONIPOS,
CONJPOS, CONKPOS and HOSTGRID keywords.

\begin{longtable}[]{@{}
  >{\raggedright\arraybackslash}p{(\linewidth - 10\tabcolsep) * \real{0.1667}}
  >{\raggedright\arraybackslash}p{(\linewidth - 10\tabcolsep) * \real{0.1667}}
  >{\raggedright\arraybackslash}p{(\linewidth - 10\tabcolsep) * \real{0.1667}}
  >{\raggedright\arraybackslash}p{(\linewidth - 10\tabcolsep) * \real{0.1667}}
  >{\raggedright\arraybackslash}p{(\linewidth - 10\tabcolsep) * \real{0.1667}}
  >{\raggedright\arraybackslash}p{(\linewidth - 10\tabcolsep) * \real{0.1667}}@{}}
\toprule\noalign{}
\begin{minipage}[b]{\linewidth}\raggedright
No.
\end{minipage} & \begin{minipage}[b]{\linewidth}\raggedright
keyword Name
\end{minipage} & \begin{minipage}[b]{\linewidth}\raggedright
RFT Data File -- Well and Connection Data keyword Table A.39
\end{minipage} & \begin{minipage}[b]{\linewidth}\raggedright
Status or Value
\end{minipage} & \begin{minipage}[b]{\linewidth}\raggedright
\end{minipage} & \begin{minipage}[b]{\linewidth}\raggedright
\end{minipage} \\
\midrule\noalign{}
\endhead
\bottomrule\noalign{}
\endlastfoot
Keyword & No.~of Entries & Data Type & & & \\
& WELLETC & This keyword marks the start of the well and connection data
set and comprises of 16 character variables all enclose in single
quotes. & Required & & \\
1-1 & Format & WELLETC & 16 & CHAR & \\
2-1 & Data & Time units. & & & \\
2-2 & Data & Well name. & & & \\
2-3 & Data & Local Grid Refinement (``LGR'') name or left blank if the
well is not in an LGR. LGRs are not currently supported by OPM Flow. &
Blank & & \\
2-4 & Data & Depth units. & & & \\
2-5 & Data & Pressure units. & & & \\
2-6 & Data & Type of data being written out in subsequent keywords and
should be set to R for RFT data, P for PLT data, or S for Segment data.
The Segment data option is not currently supported in OPM Flow. & & & \\
2.7 & Data & Well type data set to STANDARD for a standard well,
FRICTION for a wellbore friction well, and MULTSEG for a multi-segment
well. Only the STANDARD well option is currently supported by OPM Flow.
& & & \\
2-8 & Data & Liquid flow rate units. & & & \\
2-9 & Data & Gas flow rate units. & & & \\
2-10 & Data & Local volumetric flow rate units. & & & \\
2-11 & Data & Flow velocity units. & & & \\
2-12 & Data & Not used and should be a blank character string. & Blank &
& \\
2-13 & Data & Liquid and gas viscosity units. & & & \\
2-14 & Data & Polymer and Brine concentration units. & & & \\
2-15 & Data & Polymer and Brine flow rate units. & & & \\
2-16 & Data & Adsorbed Polymer concentration units. & & & \\
& CONIPOS & This keyword defines the grid index in the x direction (I
direction) for all NCON connections in the well associated with the RFT
data set. & Required & & \\
1-1 & Format & CONIPOS & NCON & INTE & \\
2-1 & Data & Grid index in the I direction for each connection. & & & \\
& CONJPOS & This keyword defines the grid index in the y direction (J
direction) for all NCON connections in the well associated with the RFT
data set. & Required & & \\
1-1 & Format & CONJPOS & NCON & INTE & \\
2-1 & Data & Grid index in the J direction for each connection. & & & \\
& CONKPOS & This keyword defines the grid index in the z direction (K
direction) for all NCON connections in the well associated with the RFT
data set. & Required & & \\
1-1 & Format & CONKPOS & NCON & INTE & \\
2-1 & Data & Grid index in the K direction for each connection. & & & \\
& HOSTGRID & Defines the LGR name for each of the NCON connections.
LGR's have not been implemented in OPM Flow and therefore this keyword
should contain NCON blank character strings. & Blank & & \\
1-1 & Format & HOSTGRID & NCON & CHAR & Blank \\
2-1 & Data & Connections host LGR grid name. & & & \\
Notes: & & & & & \\
\end{longtable}

\emph{Table A.39: RFT Data File -- Well and Connection Data Keyword}

This series of keywords are immediate followed by the RFT solution
keywords: DEPTH, PRESSURE, SWAT, and SGAS.

The RFT solution keywords specify the actual solution data (pressure,
saturation etc.) for each connection in the well, as described in Table
A.40.

\begin{longtable}[]{@{}
  >{\raggedright\arraybackslash}p{(\linewidth - 10\tabcolsep) * \real{0.1667}}
  >{\raggedright\arraybackslash}p{(\linewidth - 10\tabcolsep) * \real{0.1667}}
  >{\raggedright\arraybackslash}p{(\linewidth - 10\tabcolsep) * \real{0.1667}}
  >{\raggedright\arraybackslash}p{(\linewidth - 10\tabcolsep) * \real{0.1667}}
  >{\raggedright\arraybackslash}p{(\linewidth - 10\tabcolsep) * \real{0.1667}}
  >{\raggedright\arraybackslash}p{(\linewidth - 10\tabcolsep) * \real{0.1667}}@{}}
\toprule\noalign{}
\begin{minipage}[b]{\linewidth}\raggedright
No.
\end{minipage} & \begin{minipage}[b]{\linewidth}\raggedright
keyword Name
\end{minipage} & \begin{minipage}[b]{\linewidth}\raggedright
RESTART Data - Solution Data keyword Table A.40
\end{minipage} & \begin{minipage}[b]{\linewidth}\raggedright
Status or Value
\end{minipage} & \begin{minipage}[b]{\linewidth}\raggedright
\end{minipage} & \begin{minipage}[b]{\linewidth}\raggedright
\end{minipage} \\
\midrule\noalign{}
\endhead
\bottomrule\noalign{}
\endlastfoot
Keyword & No.~of Entries & Data Type & & & \\
1-1 & DEPTH & Connection grid block depth. & Required & & \\
1-1 & Format & DEPTH & NCON & REAL & \\
2-1 & Data & Depth data for each connection. & & & \\
& PRESSURE & Connection grid block pressure at the requested time. &
Required & & \\
1-1 & Format & PRESSURE & NCON & REAL & \\
2-1 & Data & Pressure data for each connection. & & & \\
& SWAT & Connection grid block water saturation at the requested time. &
Required & & \\
1-1 & Format & SWAT & NCON & REAL & \\
2-1 & Data & Water saturation data for each connection. & & & \\
2-4 & SGAS & Connection grid block gas saturation at the requested time.
& Required & & \\
1-1 & Format & SGAS & NCON & REAL & \\
2-1 & Data & Gas saturation data for each connection. & & & \\
& CPLY & Connection grid block polymer concentration at the requested
time. & Optional & & \\
1-1 & Format & CPLY & NCON & REAL & \\
2-1 & Data & Polymer concentration data for each connection. & & & \\
& CPLAD & Connection grid block adsorbed polymer concentration at the
requested time. & Optional & & \\
1-1 & Format & CPLAD & NCON & REAL & \\
2-1 & Data & Adsorbed polymer concentration data for each connection. &
& & \\
& CBRI & Connection grid block brine concentration at the requested
time. & Optional & & \\
1-1 & Format & CBRI & NCON & REAL & \\
2-1 & Data & Brine concentration data for each connection. & & & \\
Notes: & & & & & \\
\end{longtable}

\emph{Table A.40: RFT Data File -- RFT Solution Data Keywords}

A complete standard RFT example showing the TIME, DATE, WELLETC,
CONIPOS, CONJPOS, CONKPOS, DEPTH, PRESSURE, SWAT, and SGAS keywords, for
one well at one time step is shown on the following page.

\subsubsection{Example}\label{example-613}

The example shows all the keywords required to define a standard RFT
data set for well B-4H which has 20 connection as of March 29, 1998. The
units declared on the WELLETC keyword indicate the standard metric unit
system.

\begin{verbatim}
'TIME    '           1 'REAL'
  0.14300000E+03
'DATE    '           3 'INTE'
         29           3        1998
'WELLETC '          16 'CHAR'
'  DAYS  ' 'B-4H    ' '        ' ' METRES ' '  BARSA ' 'R       ' 'STANDARD'
' SM3/DAY' ' SM3/DAY' ' RM3/DAY' ' M/SEC  ' '        ' '   CP   ' ' KG/SM3 '
' KG/DAY ' '  KG/KG '
'CONIPOS '          20 'INTE'
         10          10          10          10          10          10
         10          10          10           9           9           9
          9           9           9           9           9           9
          9           9
'CONJPOS '          20 'INTE'
         32          32          32          32          32          32
         32          32          32          32          32          32
         32          32          32          32          32          31
         31          31
'CONKPOS '          20 'INTE'
          1           2           3           5           6           7
          8           9          10          13          14          15
         16          17          18          19          20          20
         21          22
'HOSTGRID'          20 'CHAR'
'        ' '        ' '        ' '        ' '        ' '        ' '        '
'        ' '        ' '        ' '        ' '        ' '        ' '        '
'        ' '        ' '        ' '        ' '        ' '        '
'DEPTH   '          20 'REAL'
  0.25333979E+04   0.25452786E+04   0.25577705E+04   0.25742898E+04
  0.25797493E+04   0.25855930E+04   0.25918203E+04   0.25966951E+04
  0.26071150E+04   0.26206902E+04   0.26333235E+04   0.26461284E+04
  0.26576768E+04   0.26637878E+04   0.26665945E+04   0.26786790E+04
  0.27012888E+04   0.27033108E+04   0.27381892E+04   0.27746746E+04
'PRESSURE'          20 'REAL'
  0.27384171E+03   0.27406891E+03   0.27430649E+03   0.25063243E+03
  0.25072800E+03   0.25085713E+03   0.25126656E+03   0.25138820E+03
  0.25212048E+03   0.25290869E+03   0.25380568E+03   0.25471658E+03
  0.25669318E+03   0.25713290E+03   0.25733768E+03   0.26151834E+03
  0.26378452E+03   0.26398804E+03   0.28049472E+03   0.28405634E+03
'SWAT    '          20 'REAL'
  0.79952940E-01   0.10493817E+00   0.16989987E+00   0.10517924E+00
  0.10518920E+00   0.14012721E+00   0.14357868E+00   0.80211885E-01
  0.80138676E-01   0.80140039E-01   0.90355128E-01   0.19185907E+00
  0.12343995E+00   0.29019111E+00   0.14448708E+00   0.32566309E+00
  0.10000000E+01   0.10000000E+01   0.10000000E+01   0.10000000E+01
'SGAS    '          20 'REAL'
  0.92004704E+00   0.89506185E+00   0.83010012E+00   0.89231288E+00
  0.89236671E+00   0.14750090E+00   0.17341191E-01   0.64456672E-02
  0.00000000E+00   0.00000000E+00   0.00000000E+00   0.00000000E+00
  0.00000000E+00   0.00000000E+00   0.00000000E+00   0.00000000E+00
  0.00000000E+00   0.00000000E+00   0.00000000E+00   0.00000000E+00
\end{verbatim}

End of Example

This data set is currently not supported.

This data set is currently not supported.

This data set is currently not supported.

The SUMMARY files contain the variables requested to be written to the
summary files via the keywords declared in the SUMMARY section. The data
are used to generate line graphs of properties such as oil flow rate
versus time, grid block pressure versus time, etc. See sections Error:
Reference source not found Error: Reference source not found and Error:
Reference source not found Error: Reference source not found that
declare the variables to be written out to the SUMMARY files. The
default behavior is write out the requested variables at each time step.
As this can lead to large files, especially for full field simulation
models, the Error: Reference source not found keyword allows one to
write out the data only at a report time step instead.

There are three types of SUMMARY files, the first SUMMARY file is the
OPM Flow specific Enhanced SUMMARY (*.ESMRY) output file that is
optimized for fast loading of selected vectors by post-processing
applications, including OPM ResInsight. This file type is standalone,
that is there is no separate index file.

The other two files types associated with the time based vector data, is
an index file (SUMMARY Index) that contains a list of variables and
objects (wells, groups, connections etc.) that are to be written out to
summary data file, and the summary data file (SUMMARY Data) that
contains the actual data written out at the required time steps. The
commercial simulator uses the term ``MINISTEP'' to describe the time the
data is written, and as mentioned above, ``MINISTEP'' may be either a
time step or a reporting time step.

Again, this is an OPM FLOW specific Enhanced SUMMARY (\emph{.ESMRY)
output file format that is optimized for fast loading of selected
vectors by post-processing applications. The }.ESMRY file is re-written
for every time step and load times are in the order of 30 times faster
compared with fetching selected vectors using the *.UNSMRY file format
in post processing software. The file type is compatible with OPM
ResInsight and the option is activated via the enable-esmry=true command
line option.

The structure for this series of keywords is outlined in Table A.41 and
depends on the data requested to be written out to the SUMMARY files.

\begin{longtable}[]{@{}
  >{\raggedright\arraybackslash}p{(\linewidth - 10\tabcolsep) * \real{0.1667}}
  >{\raggedright\arraybackslash}p{(\linewidth - 10\tabcolsep) * \real{0.1667}}
  >{\raggedright\arraybackslash}p{(\linewidth - 10\tabcolsep) * \real{0.1667}}
  >{\raggedright\arraybackslash}p{(\linewidth - 10\tabcolsep) * \real{0.1667}}
  >{\raggedright\arraybackslash}p{(\linewidth - 10\tabcolsep) * \real{0.1667}}
  >{\raggedright\arraybackslash}p{(\linewidth - 10\tabcolsep) * \real{0.1667}}@{}}
\toprule\noalign{}
\begin{minipage}[b]{\linewidth}\raggedright
No.
\end{minipage} & \begin{minipage}[b]{\linewidth}\raggedright
keyword Name
\end{minipage} & \begin{minipage}[b]{\linewidth}\raggedright
Enhanced SUMMARY File keywords Table A.41
\end{minipage} & \begin{minipage}[b]{\linewidth}\raggedright
Status or Value
\end{minipage} & \begin{minipage}[b]{\linewidth}\raggedright
\end{minipage} & \begin{minipage}[b]{\linewidth}\raggedright
\end{minipage} \\
\midrule\noalign{}
\endhead
\bottomrule\noalign{}
\endlastfoot
Keyword & No.~of Entries & Data Type & & & \\
& START & Three integer values representing the start date of the
simulation., and four integers representing the hour, minutes, seconds
and micro seconds the start date. & Required & & \\
1-1 & Format & START & 7 & INTE & \\
2-1 & Data & First part of start date of the run, DAY, the value should
be greater than or equal to one and less than or equal to 31 for the
calendar days. & DAY & & \\
2-2 & Data & Second part of start date of the run, MONTH, the value
should be greater than or equal to one and less than or equal to 12 for
the calendar month. & MONTH & & \\
2-3 & Data & Third part of start date of run, YEAR, a positive four
digit integer value of the start year, which must be specified fully by
four digits, for example 1986. & YEAR & & \\
2-4 & Data & The first part of the start time in the form HH:MM:SS:MS,
that is HH for hours. The value should be greater than or equal to zero
and less than or equal to 23. & IHOURZ & & \\
2-5 & Data & The second part of the start time in the form HH:MM:SS:MS,
that is MM for minutes. The value should be greater than or equal to
zero and less than or equal to 59. & IMINTS & & \\
2-6 & Data & The third part of the start time in the form HH:MM:SS:MS,
that is SS for seconds. The value should be greater than or equal to
zero and less than or equal to 59. & ISECND & & \\
2-7 & Data & The fourth part of the start time in the form HH:MM:SS:MS,
that is MS for milliseconds. The value should be greater than or equal
to zero and less than or equal to 999. & & & \\
& RESTART & RESTART is only present for restarted runs and contains the
root file name (absolute or relative path) to the case which this run
was restarted from. & Optional & & \\
1-1 & Format & RESTART & 1 & C0nn & \\
2-1 & Data & The name of file, enclosed in quotes, which contains the
name of the case the current run was re-started from. The length of the
character string is set by the Data Type, C0nn, for example, C125 would
mean that the RESTART filename character length is 125 characters. & &
& \\
& RSTNUM & The RSTNUM keyword is only present for restart runs and
defines the report time step this case was restarted from. & Optional &
& \\
1-1 & Format & RSTNUM & 1 & INTE & \\
2-1 & Data & Restart report time step. number. & & & \\
& KEYCHECK & A list of all the summary variable mnemonics written to the
file at each time step., where NLIST is the number of vectors and C0nn
is the constant fixed length of each variable character string. For
example: `KEYCHECK' 6 `C010' `TIME' `YEARS' `FOPR' `GGOR:PROD' `TCPU'
`WOPR:PROD1' Where each vector name consists of 10 characters. For
vectors associated with objects, groups, wells etc., the vector name is
follows by a ``:'' and then the object name as shown in the above
example. & Required & & \\
1-1 & Format & KEYCHECK & NLIST & C0nn & \\
2-1 & Data & Each mnemonic string should be enclosed in quotes and have
a character length of C0nn, for example `C010' in the above example. & &
& \\
& UNITS & Units & Required & & \\
1-1 & Format & UNITS & NLIST & CHAR & \\
2-1 & Data & Units for the associated summary variable mnemonics. & &
& \\
& RSTEP & An array that holds one integer value for each time step. An
integer value of one indicates time steps that are also report time
steps and a zero indicates time steps only. & Required & & \\
1-1 & Format & RSTEP & NSTEP & INTE & \\
2-1 & Data & One value for each time step and the values should be
either zero or one. & & & \\
& TSTEP & An array holding one integer value for each time step with the
values representing the time step indices. Note that the values starts
from zero, such that the Data will have values of the form 0, 1, 2,
\ldots, NSTEP - 1, that is unless the RPTONLY keyword has been used in
the SUMMARY section to write only the report time steps to the file.
Values should range between 0 to NSTEP - 1 & Required & & \\
1-1 & Format & TSTEP & NSTEP & INTE & \\
2-1 & Data & One value for each time step. with the value ranging from
zero to NSTEP - 1. & & & \\
& VX & One data array (length = NSTEP) for each summary vector specified
in array KEYCHECK. Naming of these arrays should be V0, V1, V2, \ldots{}
etc. & Required & & \\
1-1 & Format & VX & NSTEP & REAL & \\
2-1 & Data & One real value for each time step for one specific summary
vector VX vector, where X ranges from zero to NLIST -- 1, for the NLIST
variables declared via the KEYCHECK keyword. & & & \\
Notes: & & & & & \\
\end{longtable}

\emph{Table A.41: Enhanced SUMMARY File Keywords}

\subsubsection{Example: Enhanced SUMMARY File
Keywords}\label{example-enhanced-summary-file-keywords}

The following example shows a typical formatted Enhanced SUMMARY file
set of keywords.

\begin{verbatim}
'START   '           7 'INTE'
           1          11        2018           0           0           0
           0
'KEYCHECK'           6 'C010'
'TIME      ' 'YEARS     ' 'FOPR      ' 'GGOR:PROD ' 'TCPU      ' 'WOPR:PROD1'
'UNITS   '           6 'CHAR'
'DAYS    ' '        ' 'SM3/DAY ' 'SM3/SM3 ' 'SECONDS ' 'SM3/DAY '
'RSTEP   '           4 'INTE'
           1           1           0           1
'TSTEP   '           4 'INTE'
           0           1           2           3
'V0      '           4 'REAL'
    1.0000000E+00    4.0000000E+00    9.0000000E+00    1.4000000E+01
'V1      '           4 'REAL'
    2.7378509E-03    1.0951404E-02    2.4640657E-02    3.8329910E-02
'V2      '           4 'REAL'
    0.0000000E+00    0.0000000E+00    0.0000000E+00    0.0000000E+00
'V3      '           4 'REAL'
    0.0000000E+00    0.0000000E+00    0.0000000E+00    0.0000000E+00
'V4      '           4 'REAL'
    6.0681427E-01    6.5710169E-01    7.0169014E-01    7.6383120E-01
'V5      '           4 'REAL'
    0.0000000E+00    0.0000000E+00    0.0000000E+00    0.0000000E+00
\end{verbatim}

The example shows the file holding six summary vector and four time
steps, with the first, second and fourth time steps being report time
steps.

The structure for this series of keywords is outlined in Table A.42,
Table A.43, and Table A.44 and depends on the data requested to be
written out to the SUMMARY Data File.

\begin{longtable}[]{@{}
  >{\raggedright\arraybackslash}p{(\linewidth - 10\tabcolsep) * \real{0.1667}}
  >{\raggedright\arraybackslash}p{(\linewidth - 10\tabcolsep) * \real{0.1667}}
  >{\raggedright\arraybackslash}p{(\linewidth - 10\tabcolsep) * \real{0.1667}}
  >{\raggedright\arraybackslash}p{(\linewidth - 10\tabcolsep) * \real{0.1667}}
  >{\raggedright\arraybackslash}p{(\linewidth - 10\tabcolsep) * \real{0.1667}}
  >{\raggedright\arraybackslash}p{(\linewidth - 10\tabcolsep) * \real{0.1667}}@{}}
\toprule\noalign{}
\begin{minipage}[b]{\linewidth}\raggedright
No.
\end{minipage} & \begin{minipage}[b]{\linewidth}\raggedright
keyword Name
\end{minipage} & \begin{minipage}[b]{\linewidth}\raggedright
SUMMARY Index File keywords (Global) Table A.42
\end{minipage} & \begin{minipage}[b]{\linewidth}\raggedright
Status or Value
\end{minipage} & \begin{minipage}[b]{\linewidth}\raggedright
\end{minipage} & \begin{minipage}[b]{\linewidth}\raggedright
\end{minipage} \\
\midrule\noalign{}
\endhead
\bottomrule\noalign{}
\endlastfoot
Keyword & No.~of Entries & Data Type & & & \\
& INTEHEAD & The integer header keyword contains two values indicating
the units system and the simulator used to generate the file. & Required
& & \\
1-1 & Format & INTEHEAD & 2 & INTE & \\
2-1 & Data & Unit system used in the simulation, set to 1 for metric, 2
for field, and 3 for laboratory & & & \\
2-2 & Data & Simulator code used to generate the file, set to: 100 for
Schlumberger's ECLIPSE 100, 300 for Schlumberger's ECLIPSE 300, 500 for
Schlumberger's ECLIPSE 300 Thermal, 700 for Schlumberger's INTERSECT
simulator, 800 for Schlumberger's FrontSim simulator, or a negative
value for other simulators. OPM Flow uses a value of 100. & 100 & & \\
& RESTART & Contains the root file name of the case which this run was
restarted from, if this was a restart run, other wise blank. & Required
& & \\
1-1 & Format & RESTART & 9 & CHAR & \\
2-1 & Data & Maximum of 72 characters divided into eight character
words. For example, if the data file is restarted from
NOR01-PRED01-OPM2004-R00, then this data will be of the form: `NOR01-PR'
`ED01-OPM' `2010-R00' ' ' ' ' ' ' ' ' ' ' ' ' & & & \\
& DIMENS & This keyword contains the number of vectors and objects
written to the SUMMARY Data file at each time step. & Required & & \\
1-1 & Format & DIMENS & 6 & REAL & \\
2-1 & Data & NLIST is the number of values stored at each time step on
the SUMMARY Data File. & & & \\
2-2 & Data & The number of grid blocks in the x-direction (NX). & & & \\
2-3 & Data & The number of grid blocks in the y-direction (NY). & & & \\
2-4 & Data & The number of grid blocks in the z-direction (NZ). & & & \\
2-5 & Data & Not Used. & 0 & & \\
2-6 & Data & Report time step of the case which this run was restated
from, if this was a restart run. & 0 & & \\
& KEYWORDS & A list of all the summary variable mnemonics written to the
SUMMARY Data file at each time step. & Required & & \\
1-1 & Format & KEYWORDS & NLIST & CHAR & \\
2-1 & Data & Each mnemonic string should be enclosed in quotes. & & & \\
& WGNAMES & A list of all well and group objects associated with summary
variable mnemonics & Required & & \\
1-1 & Format & WGNAMES & NLIST & CHAR & \\
2-1 & Data & Each well and group name should be enclosed in quotes & &
& \\
& NUMS & Integer grid cell or region number associated with the summary
variable mnemonics. & Required & & \\
1-1 & Format & NUMS & NLIST & INTE & \\
2-1 & Data & Integer values. & & & \\
& LENGTHS & Horizontal well length. & Optional & & \\
1-1 & Format & LENGTHS & NLIST & REAL & \\
2-1 & Data & The distance from the bottom-hole reference depth to the
completion depth for the associated summary variable mnemonics for
horizontal wells & & & \\
& LENUNITS & Horizontal well length units. & Optional & & \\
1-1 & Format & LENUNITS & 1 & CHAR & \\
2-1 & Data & Horizontal well length units. & & & \\
& MEASRMNT & Measurements associated with each summary variable
mnemonic. & Optional & & \\
1-1 & Format & MEASRMNT & NBLOCK x NLIST & CHAR & \\
2-1 & Data & NBLOCK is equal to the number of elements provided by this
keyword divided by NLIST. & & & \\
& UNITS & Units & Required & & \\
1-1 & Format & UNITS & NLIST & CHAR & \\
2-1 & Data & Units for the associated summary variable mnemonics. & &
& \\
& STARTDAT & Three integer values representing the date and four
integers representing the time of the data being written out. & Required
& & \\
1-1 & Format & STARTDAT & 3 & INTE & \\
2-1 & Data & Calendar day and should range between one and 31. & DAY &
& \\
2-2 & Data & Calendar month should range between one and 12. & MONTH &
& \\
2-3 & Data & Calendar year in four digits, that is 2020. & YEAR & & \\
2-4 & Data & Hour. & HOUR & & \\
2-5 & Data & Minutes. & MINUTE & & \\
2-6 & Data & Seconds in micro seconds, 0-59:999:999 & SECOND & & \\
Notes: & & & & & \\
\end{longtable}

\emph{Table A.42: SUMMARY Index File Keywords (Global)}

In addition, Table A.43 summarizes the SUMMARY Index File keywords
associated with Local Grid Refinements.

\begin{longtable}[]{@{}
  >{\raggedright\arraybackslash}p{(\linewidth - 10\tabcolsep) * \real{0.1667}}
  >{\raggedright\arraybackslash}p{(\linewidth - 10\tabcolsep) * \real{0.1667}}
  >{\raggedright\arraybackslash}p{(\linewidth - 10\tabcolsep) * \real{0.1667}}
  >{\raggedright\arraybackslash}p{(\linewidth - 10\tabcolsep) * \real{0.1667}}
  >{\raggedright\arraybackslash}p{(\linewidth - 10\tabcolsep) * \real{0.1667}}
  >{\raggedright\arraybackslash}p{(\linewidth - 10\tabcolsep) * \real{0.1667}}@{}}
\toprule\noalign{}
\begin{minipage}[b]{\linewidth}\raggedright
No.
\end{minipage} & \begin{minipage}[b]{\linewidth}\raggedright
keyword Name
\end{minipage} & \begin{minipage}[b]{\linewidth}\raggedright
SUMMARY Index File keywords (LGR) Table A.43
\end{minipage} & \begin{minipage}[b]{\linewidth}\raggedright
Status or Value
\end{minipage} & \begin{minipage}[b]{\linewidth}\raggedright
\end{minipage} & \begin{minipage}[b]{\linewidth}\raggedright
\end{minipage} \\
\midrule\noalign{}
\endhead
\bottomrule\noalign{}
\endlastfoot
Keyword & No.~of Entries & Data Type & & & \\
& LGRS & Local Grid Refinement names associated summary variable
mnemonic for LGR grids. & Optional & & \\
1-1 & Format & LGRS & NLIST & CHAR & \\
2-1 & Data & Name of the LGRs. & & & \\
& NUMLX & This keyword defines the LGR grid index in the x direction (I
direction) for the associated summary variable mnemonic for LGR grids. &
Optional & & \\
1-1 & Format & NUMLX & NLIST & INTE & \\
2-1 & Data & LGR I location of local grid block data or the connection
data. & & & \\
& NUMLY & This keyword defines the LGR grid index in the y direction (J
direction) for the associated summary variable mnemonic for LGR grids. &
Optional & & \\
1-1 & Format & NUMLY & NLIST & INTE & \\
2-1 & Data & LGR J location of local grid block data or the connection
data. & & & \\
& NUMLZ & This keyword defines the LGR grid index in the z direction (K
direction) for the associated summary variable mnemonic for LGR grids. &
Optional & & \\
1-1 & Format & NUMLX & NLIST & INTE & \\
2-1 & Data & LGR K location of local grid block data or the connection
data. & & & \\
& LGRNAMES & The names of all the LGRs in this case with the number of
of LGRs in the case given by NLGR. & Optional & & \\
1-1 & Format & LGRNAMES & NLGR & CHAR & \\
2-1 & Data & List of LGR names. & & & \\
& LGRVEC & Integer list of the number of summary variable mnemonics for
each LGR. & Optional & & \\
1-1 & Format & LGRVEC & NLGR & INTE & \\
2-1 & Data & Number of summary variable mnemonics for each LGR. & & & \\
& LGRTIMES & A list of the total number of local time steps for each
LGR. & & & \\
1-1 & Format & LGRTIMES & NLGR & INTE & \\
2-1 & Data & Number of local time steps for each LGR. & & & \\
Notes: & & & & & \\
\end{longtable}

\emph{Table A.43: SUMMARY Index File Keywords (LGR)}

Finally, Table A.44 summarizes the SUMMARY Index File keywords specific
to the commercial simulator and are therefore not described in any
detail.

\begin{longtable}[]{@{}
  >{\raggedright\arraybackslash}p{(\linewidth - 6\tabcolsep) * \real{0.2500}}
  >{\raggedright\arraybackslash}p{(\linewidth - 6\tabcolsep) * \real{0.2500}}
  >{\raggedright\arraybackslash}p{(\linewidth - 6\tabcolsep) * \real{0.2500}}
  >{\raggedright\arraybackslash}p{(\linewidth - 6\tabcolsep) * \real{0.2500}}@{}}
\toprule\noalign{}
\begin{minipage}[b]{\linewidth}\raggedright
No.
\end{minipage} & \begin{minipage}[b]{\linewidth}\raggedright
keyword Name
\end{minipage} & \begin{minipage}[b]{\linewidth}\raggedright
SUMMARY Index File keywords (Commercial Simulator) Table A.44
\end{minipage} & \begin{minipage}[b]{\linewidth}\raggedright
Status or Value
\end{minipage} \\
\midrule\noalign{}
\endhead
\bottomrule\noalign{}
\endlastfoot
& RUNTIMEI & Run time monitoring -- not supported. & \\
& RUNTIMED & Run time monitoring -- not supported. & \\
& STEPRESN & Run time monitoring -- not supported. & \\
& XCOORD & OFM Data -- not supported. & \\
& YCOORD & OFM -- not supported. & \\
& TIMESTMP & Not supported. & \\
Notes: & & & \\
\end{longtable}

\emph{Table A.44: SUMMARY Index File Keywords (Commercial Simulator)}

\subsubsection{Example: SUMMARY Index File
Keywords}\label{example-summary-index-file-keywords}

The following example shows a typical formatted SUMMARY Index file set
of keywords.

\begin{verbatim}
 'INTEHEAD'           2 'INTE'
           1         100
 'RESTART '           9 'CHAR'
 '        ' '        ' '        ' '        ' '        ' '        ' '        '
 '        ' '        '
 'DIMENS  '           6 'INTE'
         164          10          10          10           0           0
 'KEYWORDS'         164 'CHAR'
 'TIME    ' 'YEARS   ' 'FGIP    ' 'FGIPG   ' 'FGIPL   ' 'FGIR    ' 'FGIT    '
 'FGOR    ' 'FGPR    ' 'FGPT    ' 'FOIP    ' 'FOIPG   ' 'FOIPL   ' 'FOIR    '
 'FOIT    ' 'FOPR    ' 'FOPT    ' 'FPR     ' 'FVIR    ' 'FVIT    ' 'FVPR    '
 'FVPT    ' 'FWCT    ' 'FWIP    ' 'FWIR    ' 'FWIT    ' 'FWPR    ' 'FWPT    '
 'GGIR    ' 'GGIT    ' 'GGOR    ' 'GGPR    ' 'GGPT    ' 'GOIR    ' 'GOIT    '
 'GOPR    ' 'GOPT    ' 'GVIR    ' 'GVIT    ' 'GVPR    ' 'GVPT    ' 'GWCT    '
 'GWIR    ' 'GWIT    ' 'GWPR    ' 'GWPT    ' 'SPR     ' 'SPR     ' 'SPR     '
 'SPR     ' 'SPR     ' 'SPR     ' 'SPR     ' 'SPR     ' 'SPR     ' 'SPR     '
  ………………………….……………..
 'WWCT    ' 'WWCT    ' 'WWIR    ' 'WWIR    ' 'WWIT    ' 'WWIT    ' 'WWPR    '

 'WWPR    ' 'WWPT    ' 'WWPT    '
 'WGNAMES '         164 'CHAR'
 ':+:+:+:+' ':+:+:+:+' ':+:+:+:+' ':+:+:+:+' ':+:+:+:+' ':+:+:+:+' ':+:+:+:+'
 ':+:+:+:+' ':+:+:+:+' ':+:+:+:+' ':+:+:+:+' ':+:+:+:+' ':+:+:+:+' ':+:+:+:+'
 ':+:+:+:+' ':+:+:+:+' ':+:+:+:+' ':+:+:+:+' ':+:+:+:+' ':+:+:+:+' ':+:+:+:+'
 ':+:+:+:+' ':+:+:+:+' ':+:+:+:+' ':+:+:+:+' ':+:+:+:+' ':+:+:+:+' ':+:+:+:+'
 'PLATFORM' 'PLATFORM' 'PLATFORM' 'PLATFORM' 'PLATFORM' 'PLATFORM' 'PLATFORM'
 'PLATFORM' 'PLATFORM' 'PLATFORM' 'PLATFORM' 'PLATFORM' 'PLATFORM' 'PLATFORM'
 'PLATFORM' 'PLATFORM' 'PLATFORM' 'PLATFORM' 'OP01    ' 'OP01    ' 'OP01    '
 'OP01    ' 'OP01    ' 'OP01    ' 'OP01    ' 'OP01    ' 'OP01    ' 'OP01    '
 'OP01    ' 'OP01    ' 'OP01    ' 'OP01    ' 'OP01    ' 'OP01    ' 'OP01    '
 'OP01    ' 'OP01    ' 'OP01    ' 'OP01    ' 'OP01    ' 'OP01    ' 'OP01    '
 'OP01    ' 'OP01    ' 'OP01    ' 'OP01    ' 'OP01    ' 'OP01    ' 'OP01    '
 'OP01    ' 'OP01    ' 'OP01    ' 'OP01    ' 'OP01    ' 'OP01    ' 'OP01    '
 'OP01    ' 'OP01    ' 'OP01    ' 'OP01    ' 'OP01    ' 'OP01    ' 'OP01    '
      ………………………….……………..

 'WI01    ' 'OP01    ' 'WI01    '
 'NUMS    '         164 'INTE'
           0           0           0           0           0           0
           0           0           0           0           0           0
           0           0           0           0           0           0
           0           0           0           0           0           0
           0           0           0           0           0           0
           0           0           0           0           0           0
           0           0           0           0           0           0
           0           0           0           0           1           2
           3           4           5           6           7           8
           9          10          11          12          13          14
      ………………………….……………..
         0           0
 'UNITS   '         164 'CHAR'
 'DAYS    ' 'YEARS   ' 'SM3     ' 'SM3     ' 'SM3     ' 'SM3/DAY ' 'SM3     '
 'SM3/SM3 ' 'SM3/DAY ' 'SM3     ' 'SM3     ' 'SM3     ' 'SM3     ' 'SM3/DAY '
 'SM3     ' 'SM3/DAY ' 'SM3     ' 'BARSA   ' 'RM3/DAY ' 'RM3     ' 'RM3/DAY '
 'RM3     ' 'SM3/SM3 ' 'SM3     ' 'SM3/DAY ' 'SM3     ' 'SM3/DAY ' 'SM3     '
      ………………………….……………..
 'SM3/DAY ' 'SM3     ' 'SM3     '
 'STARTDAT'           6 'INTE'
           1           1        2020           0           0           0
\end{verbatim}

End of Example

The structure for this series of keywords is relatively simple compared
to the SUMMARY Index file, in that it only consists of four keyword
types that are written out at the requested time steps, as outlined in
Table A.45.

\begin{longtable}[]{@{}
  >{\raggedright\arraybackslash}p{(\linewidth - 10\tabcolsep) * \real{0.1667}}
  >{\raggedright\arraybackslash}p{(\linewidth - 10\tabcolsep) * \real{0.1667}}
  >{\raggedright\arraybackslash}p{(\linewidth - 10\tabcolsep) * \real{0.1667}}
  >{\raggedright\arraybackslash}p{(\linewidth - 10\tabcolsep) * \real{0.1667}}
  >{\raggedright\arraybackslash}p{(\linewidth - 10\tabcolsep) * \real{0.1667}}
  >{\raggedright\arraybackslash}p{(\linewidth - 10\tabcolsep) * \real{0.1667}}@{}}
\toprule\noalign{}
\begin{minipage}[b]{\linewidth}\raggedright
No.
\end{minipage} & \begin{minipage}[b]{\linewidth}\raggedright
keyword Name
\end{minipage} & \begin{minipage}[b]{\linewidth}\raggedright
SUMMARY Data File keywords Table A.45
\end{minipage} & \begin{minipage}[b]{\linewidth}\raggedright
Status or Value
\end{minipage} & \begin{minipage}[b]{\linewidth}\raggedright
\end{minipage} & \begin{minipage}[b]{\linewidth}\raggedright
\end{minipage} \\
\midrule\noalign{}
\endhead
\bottomrule\noalign{}
\endlastfoot
Keyword & No.~of Entries & Data Type & & & \\
& SEQNUM & This keyword contains a sequence data value corresponding to
an encoded integer representing the time the file was created. Although
required the data is not used. Secondly, this keyword is written at each
report step only and not for each inter leaving mini step, that is the
file format is of the following format: SEQNUM Report step n MINISTEP
PARAMS MINISTEP PARAMS MINISTEP PARAMS INISTEP PARAMS SEQNUM Report step
n + 1 MINISTEP PARAMS MINISTEP PARAMS & Required & & \\
1-1 & Format & SEQNUM & 1 & INTE & \\
2-1 & Data & The value entered is not used by either OPM Flow or OPM
ResInsight & 0 & & \\
& MINSTEP & The time step number starting from zero and incremented by
one for each subsequent time step. & Required & & \\
1-1 & Format & MINISTEP & 1 & INTE & \\
2-1 & Data & Time step. & & & \\
& PARAMS & This keyword contains the actual data for all the vectors and
objects specified in the index file. & Required & & \\
1-1 & Format & PARAMS & NLIST & REAL & \\
2-1 & Data & There should be NLIST values corresponding to the variable
list declared on the corresponding SUMMARY Index file. & & & \\
Notes: & & & & & \\
\end{longtable}

\emph{Table A.45: SUMMARY Data File Keywords}

For LGRs, the global cell data is written to the standard SUMMARY Data
file containing an additional keyword that states the LGR name, as
depicted in Table A.46.

\begin{longtable}[]{@{}
  >{\raggedright\arraybackslash}p{(\linewidth - 10\tabcolsep) * \real{0.1667}}
  >{\raggedright\arraybackslash}p{(\linewidth - 10\tabcolsep) * \real{0.1667}}
  >{\raggedright\arraybackslash}p{(\linewidth - 10\tabcolsep) * \real{0.1667}}
  >{\raggedright\arraybackslash}p{(\linewidth - 10\tabcolsep) * \real{0.1667}}
  >{\raggedright\arraybackslash}p{(\linewidth - 10\tabcolsep) * \real{0.1667}}
  >{\raggedright\arraybackslash}p{(\linewidth - 10\tabcolsep) * \real{0.1667}}@{}}
\toprule\noalign{}
\begin{minipage}[b]{\linewidth}\raggedright
No.
\end{minipage} & \begin{minipage}[b]{\linewidth}\raggedright
keyword Name
\end{minipage} & \begin{minipage}[b]{\linewidth}\raggedright
SUMMARY Data File keywords (LGR) Table A.46
\end{minipage} & \begin{minipage}[b]{\linewidth}\raggedright
Status or Value
\end{minipage} & \begin{minipage}[b]{\linewidth}\raggedright
\end{minipage} & \begin{minipage}[b]{\linewidth}\raggedright
\end{minipage} \\
\midrule\noalign{}
\endhead
\bottomrule\noalign{}
\endlastfoot
Keyword & No.~of Entries & Data Type & & & \\
& SEQHDR & As per Table A.45 & Required & & \\
& LGRNAME & Name of Local Grid Refinement. & Required & & \\
& Format & LGRNAME & 1 & CHAR & \\
& Data & LGR name enclosed in single quotes. & & & \\
& MINSTEP & As per Table A.45. & Required & & \\
& PARAMS & As per Table A.45. & Required & & \\
Notes: & & & & & \\
\end{longtable}

\emph{Table A.46: SUMMARY Data File Keywords (LGR)}

The data for the LGR time steps, as opposed to the global time step, are
written out to a separate SUMMARY Data File, with the suffix of LGR or
FLG for fixed format output. In this case, the MINSTEP keyword states
the ``local'' time step and the keyword is repeated for each local grid
in the model.

\subsubsection{Example: SUMMARY Data Data File
keyword}\label{example-summary-data-data-file-keyword}

The following example shows a typical formatted SUMMARY Data File
keyword data set for one complete report time step and part of a second
report time step.

\begin{verbatim}
'SEQHDR  '           1 'INTE'
          1
'MINISTEP'           1 'INTE'
          0
'PARAMS  '         164 'REAL'
  0.50000000E+00   0.13689254E-02   0.18381236E+11   0.51459348E+10
  0.13235301E+11   0.00000000E+00   0.00000000E+00   0.37649313E+03
  0.95391960E+07   0.47695980E+07   0.10947282E+09   0.30000488E+06
  0.10917282E+09   0.00000000E+00   0.00000000E+00   0.25336973E+05
  0.12668486E+05   0.27468176E+03   0.15571577E+05   0.77857886E+04
  0.63867320E+05   0.31933660E+05   0.12643164E-02   0.60474832E+08
  0.15000000E+05   0.75000000E+04   0.32074501E+02   0.16037251E+02
  0.00000000E+00   0.00000000E+00   0.37649313E+03   0.95391960E+07
  0.47695980E+07   0.00000000E+00   0.00000000E+00   0.25336973E+05
  0.12668486E+05   0.15571577E+05   0.77857886E+04   0.63867320E+05
  0.31933660E+05   0.12643164E-02   0.15000000E+05   0.75000000E+04
  0.32074501E+02   0.16037251E+02   0.26000000E+03   0.26172531E+03
  0.26362421E+03   0.26596732E+03   0.26747589E+03   0.26896045E+03
  ……………………….……………..
'SEQHDR  '           1 'INTE'
          2
'MINISTEP'           1 'INTE'
          1
'PARAMS  '         164 'REAL'
  0.10000000E+01   0.27378509E-02   0.18377710E+11   0.51453399E+10
  0.13232370E+11   0.00000000E+00   0.00000000E+00   0.32637939E+03
 ………………………….……………..
\end{verbatim}

End of Example

This file format is currently not supported by OPM Flow.




\end{document}
